\documentclass[11pt,letterpaper,reqno]{amsart}
\makeatletter
\NeedsTeXFormat{LaTeX2e}[2020/10/01]

\RequirePackage[T1]{fontenc}
\RequirePackage{lmodern}
\RequirePackage{microtype}
\RequirePackage[letterpaper,margin=1in,headheight=15pt]{geometry}
\RequirePackage{amsmath,amsthm,amssymb,mathtools}
\RequirePackage{enumitem}
\RequirePackage{perpage}
\RequirePackage{xcolor}
\definecolor{manuscriptred}{rgb}{0.6,0.1,0.1}
\definecolor{manuscriptblue}{rgb}{0.1,0.3,0.5}
\definecolor{manuscriptmagenta}{HTML}{60445C}
\RequirePackage[breaklinks=true,colorlinks=true,
  linkcolor=manuscriptred,citecolor=manuscriptblue,urlcolor=manuscriptmagenta,
  filecolor=manuscriptmagenta,runcolor=manuscriptmagenta]{hyperref}
\RequirePackage{cleveref}
\hypersetup{pdfauthor={}}

\setcounter{tocdepth}{1}
\renewcommand{\thefigure}{\Roman{figure}}
\renewcommand{\thetable}{\Roman{table}}
\MakePerPage{footnote}
\renewcommand{\thefootnote}{%
  \ifnum\value{footnote}>9\arabic{footnote}\else\fnsymbol{footnote}\fi}

\theoremstyle{definition}
\newtheorem{ms@theorem}{Theorem}
\renewcommand{\thems@theorem}{\Alph{ms@theorem}}
\newtheorem*{ms@theoremstar}{Theorem}
\newtheorem{ms@proposition}{Proposition}[section]
\newtheorem{ms@lemma}[ms@proposition]{Lemma}
\newtheorem{ms@definition}[ms@proposition]{Definition}
\newtheorem{ms@corollary}[ms@proposition]{Corollary}
\newtheorem*{ms@corollarystar}{Corollary}

\newcommand{\ms@majorstart}{%
  \pushQED{\qed}\renewcommand{\qedsymbol}{$\blacksquare\square$}}
\newcommand{\ms@minorstart}{%
  \pushQED{\qed}\renewcommand{\qedsymbol}{$\blacksquare$}}
\newenvironment{theorem}[1][]{\ms@majorstart\ms@theorem[#1]}
  {\popQED\endms@theorem}
\newenvironment{theorem*}[1][]{\ms@majorstart\ms@theoremstar[#1]}
  {\popQED\endms@theoremstar}
\newenvironment{proposition}[1][]{\ms@minorstart\ms@proposition[#1]}
  {\popQED\endms@proposition}
\newenvironment{lemma}[1][]{\ms@minorstart\ms@lemma[#1]}
  {\popQED\endms@lemma}
\newenvironment{definition}[1][]{\ms@minorstart\ms@definition[#1]}
  {\popQED\endms@definition}
\newenvironment{corollary}[1][]{\ms@majorstart\ms@corollary[#1]}
  {\popQED\endms@corollary}
\newenvironment{corollary*}[1][]{\ms@majorstart\ms@corollarystar[#1]}
  {\popQED\endms@corollarystar}

\theoremstyle{remark}
\newtheorem{remark}[ms@proposition]{Remark}
\newtheorem*{remark*}{Remark}

\crefname{ms@theorem}{theorem}{theorems}
\Crefname{ms@theorem}{Theorem}{Theorems}
\crefname{ms@proposition}{proposition}{propositions}
\Crefname{ms@proposition}{Proposition}{Propositions}
\crefname{ms@lemma}{lemma}{lemmas}
\Crefname{ms@lemma}{Lemma}{Lemmas}
\crefname{ms@definition}{definition}{definitions}
\Crefname{ms@definition}{Definition}{Definitions}
\crefname{ms@corollary}{corollary}{corollaries}
\Crefname{ms@corollary}{Corollary}{Corollaries}
\crefname{remark}{remark}{remarks}
\Crefname{remark}{Remark}{Remarks}

\RequirePackage{mathrsfs,alphalph}
\renewcommand{\thems@theorem}{\AlphAlph{\value{ms@theorem}}}
\allowdisplaybreaks[2]
\emergencystretch=1em
\newcommand{\YMMref}[2]{\href{run:YM3_I.pdf\#page=#2}{\textup{YMM1:#1}}}
\pdfstringdefDisableCommands{\def\YMMref#1#2{YMM1:#1}}
\AtBeginDocument{%
 \@ifundefined{r@IV-54-29}{}{%
  \expandafter\let\csname r@IV-30-47\expandafter\endcsname\csname r@IV-54-29\endcsname}}
\makeatother

\title[Curvature observables and continuum limits]{Curvature observables, dynamics, and conditional confinement in three-dimensional lattice Yang--Mills theory}
\date{}
\hypersetup{pdftitle={Curvature observables, dynamics, and conditional confinement in three-dimensional lattice Yang--Mills theory},pdfauthor={}}
\begin{document}
\begin{abstract}
Reflection-positive symmetric-tensor correlations with bounded slab insertions, a power bound in the integrable norm, a scale-independent off-diagonal bound and conservation are Euclidean covariant. We prove this implication and apply it to the curvature bilinears of three-dimensional $SU(2)$ lattice Yang--Mills theory. We also give quantitative construction estimates and the two-term short-distance calculation, and study local nets, dynamics and Wilson loops. Confinement assumes a box inequality uniform along a fixed physical scale and label. The quantitative short-distance law at arbitrary order is an additional hypothesis; its continuum and state-transfer consequences retain that hypothesis. Weighted critical backgrounds are constructed near real regular references, with uniform local response and Hessian estimates. The moved-background and centering results retain a complex scalar gauge-transfer hypothesis. Overlapping-source rates, some centre-flux limits, flowed observables and complete gauge-group transfers retain the additional hypotheses stated for them.
\end{abstract}
\maketitle
\tableofcontents
\section{Introduction}\label{sec:sketches-introduction}
Let $V=\operatorname{Sym}_3(\mathbb R)$ and $V_{\mathbb C}=\operatorname{Sym}_3(\mathbb C)$. An isometry $r:x\mapsto Rx+v$ acts on tensor tests by
\[
 (rH)(x)=R H(r^{-1}x)R^{\rm T}.
\]
For an open set $O$, let $\mathfrak S_V(O)$ consist of finite tuples of continuous compactly supported $V$-valued functions in $O$ with pairwise disjoint supports. Tuples are concatenated by juxtaposition and the empty tuple is denoted by $\varnothing$. Complex combinations and complex tensor values are interpreted by multilinearity. A tuple has type $(d;R)$ if its supports are at mutual distance at least $d$ and each lies in a ball of radius $R$. Write $|H|_\infty=\sup_x\max_{i,j}|H_{ij}(x)|$. For a compactly supported vector field $\eta$ put
\[
 H_\eta=(\nabla\!\cdot\eta)\mathbf1-\nabla\eta-(\nabla\eta)^{\rm T}.
\]
The vanishing of its insertions expresses conservation of $S_{ij}=\frac12\delta_{ij}\operatorname{tr}T-T_{ij}$. Here $T$ denotes the gauge-invariant curvature bilinear, defined below by lattice observables; separate gauge-dependent fields $B_i$ are not assumed to have been constructed. The scalar observable $\frac12\operatorname{tr}T$ is the action density considered in the companion manuscript \cite{YMM1}.


Write $\theta(x_1,x_2,x_3)=(x_1,x_2,-x_3)$ and $\langle[F],[G]\rangle=\sigma(\theta\bar F\cdot G)$ for the reflection-positive form on half-space tuples. We use the tensor action $\rho(R)E=RER^{\rm T}$. Reflection of a test function includes this action on its values. For $o_{ij}=e_ie_j^{\rm T}+e_je_i^{\rm T}$, $i\ne j$, the following abstract theorem gives the principal covariance argument.

\begin{theorem}\label{sketches-main}
Let $\sigma$ be a real, symmetric multilinear functional on the tuples just defined, with $\sigma(\varnothing)=1$ and $\sigma((H))=0$. Suppose that:
\begin{enumerate}[label=\textup{(\roman*)}]
\item For every $n,d,R$ there is $C_n(d;R)<\infty$ such that $|\sigma(H_1,\ldots,H_n)|\le C_n(d;R)\prod_a|H_a|_\infty$ on tuples of type $(d;R)$.
\item The functional is invariant under translations and signed permutations, and is reflection positive in every plane $x_i=c$ and $x_i\pm x_j=c$.
\item In the Hilbert space of reflection positivity at $x_3=0$, with vacuum $\Omega=[\varnothing]$, time translations are $e^{-tH}$, spatial translations are $e^{iv\cdot P}$, and $H\ge |P_1|,|P_2|$.
\item If $h$ is supported in $0<x_3<l$, the insertion
\[
 \mathsf B_l(h)[F]=[(h)\cdot\tau_{le_3}F]
\]
extends to a bounded operator, with norm at most $\mathsf C(\mathcal K,l)|h|_\infty$ for $\operatorname{supp}h\subset\mathcal K\Subset\{0<x_3<l\}$. For some $p<\infty$ and $C<\infty$ it also satisfies
\[
 \|\mathsf B_{2d}(\tau_{de_3}h)\|
 \le C\min\{1,d-d'\}^{-p}\|h\|_{L^1}
\]
when $0\le d'<d$ and $\operatorname{supp}h\subset\{|x_3|\le d'\}$.
\item There is $C_{\rm od}<\infty$ such that
\[
 \sigma(\theta(h o_{ij}),h o_{ij})\le C_{\rm od}|h|_\infty^2
\]
for real $h$ supported in $B(x_0,d)$, $0<d<1$, with $3d/2\le x_{0,3}\le2d$.
\item For every $\eta\in C_c^2(\{x_3>0\};\mathbb R^3)$, $[(H_\eta)]=0$ in that Hilbert space.
\end{enumerate}
Then $\sigma(rX)=\sigma(X)$ for every Euclidean isometry $r$. Moreover,
\[
 \sigma((H_\eta)\cdot X)=0
\]
for every $\eta\in C_c^2(\mathbb R^3;\mathbb R^3)$ and every tuple $X$ whose supports are disjoint from $\operatorname{supp}\eta$.
\end{theorem}

We now specify the lattice observable. Fix the admitted construction constants of \cite{YMM1}, with odd $L\ge13$, running-coupling increment bound $B_\sharp$ and the short-distance ceiling $\gamma_W^{\rm sd}$. Put
\[
 \bar B=\frac{LB_\sharp}{L-1},\qquad
 I(g)=(g^{-2}+\bar B,\,L(g^{-2}+B_\sharp)-\bar B],
 \qquad a_K=L^{-K}.
\]
The admitted torus sides form $\mathfrak L(g)=\{2L^m:m\ge m_{\mathfrak L}(g)\}$, with the construction threshold $m_{\mathfrak L}(g)$. On the lattice of spacing $a_K$, every bond carries an $SU(2)$ matrix and the Wilson weight is $\exp(-\beta_K\sum_pa_p)$, with Haar probability measure on each bond and
$a_p=1-\frac12\operatorname{tr}\partial U(p)$.
For $b=\langle x,x+a_Ke_\mu\rangle$, let $\lambda_b<\lambda'_b$ be the other two directions. Write $W_x(\mu,\pm\alpha)$ for the corresponding based plaquette holonomies, and set
\[
 w(G)=\frac{G-G^{-1}}{2i},\qquad
 u_{b,\alpha}=-\tfrac12\bigl[w(W_x(\mu,\alpha))-w(W_x(\mu,-\alpha))\bigr],
 \qquad q_b=\tfrac12\operatorname{tr}(u_{b,\lambda_b}u_{b,\lambda'_b}).
\]
The gauge-invariant form is
\begin{equation}\label{sf-lat}
 O_K^T(H)=a_K^{-1}\left[
 \sum_pH_{\lambda(p)\lambda(p)}(x(p))a_p
 -\sum_bH_{\lambda_b\lambda'_b}(x_b)q_b\right],
\end{equation}
where $x(p)$ and $x_b$ are the midpoints and $\lambda(p)$ is the plaquette normal. This represents $\frac12\int\operatorname{tr}(HT)$; $O_K^T(f\mathbf1)$ is the action-density observable.

\begin{corollary}[Curvature bilinears in the $SU(2)$ model]\label{sketches-model}
Let $0<g\le\gamma_W^{\rm sd}$, $b\in I(g)$ and $\ell\in\mathfrak L(g)$, and put $\beta_K=L^Kb$.
\begin{enumerate}[label=\textup{(\roman*)}]
\item For $n\ge2$, on $T_\ell$ the joint cumulants of $O_K^T(H_1),\ldots,O_K^T(H_n)$ converge as $K\to\infty$ for every tuple of Lipschitz tensor weights with pairwise disjoint supports.
\item For any sequence $\nu_K$ of Gibbs states at $(a_K,\beta_K)$ invariant under a subgroup of finite index of the lattice translations, there is a subsequence along which
\[
 \nu_K\!\left(\prod_{a=1}^n
 [O_K^T(H_a)-\nu_K(O_K^T(H_a))]\right)
\]
converges for every tuple of continuous compactly supported tensor weights on $\mathbb R^3$ with pairwise disjoint supports. Every such limit $\sigma$ is invariant under all Euclidean isometries, is reflection positive in every plane, and satisfies
\[
 \sigma((H_\eta)\cdot X)=0
\]
for $\eta\in C_c^2(\mathbb R^3;\mathbb R^3)$ and every tuple $X$ with supports disjoint from $\operatorname{supp}\eta$.
\item The functional $\sigma$ is a probability mixture of its canonical components. After recentering, each component reconstructs tempered local Poincar\'e-covariant tensor fields in $1+2$ dimensions, with spectrum in the closed forward cone and a cyclic vacuum unique in that component.
\end{enumerate}
\end{corollary}
\begin{proof}
Theorem~\ref{C-355-15}(E3)(ii) gives (i). Theorem~\ref{C-355-9}(b) gives the subsequence in (ii), and Theorem~\ref{cov-thm}(b),(c) gives conservation and covariance. The positivity and reconstruction statements follow from Theorems~\ref{C-355-10} and~\ref{C-355-13}, with the component decomposition of the latter.
\end{proof}

Theorem~\ref{cov-thm-data} gives the analytic form of this result and its variants. Its proof has two steps. Conservation and the scale-independent off-diagonal estimate eliminate the extreme two-point angular harmonics. Positivity, applied to polygon configurations, and halving then eliminate the higher-point harmonics. The argument requires the analytic continuation of the coefficient kernels; we give that continuation as part of the proof.

For curvature bilinears the hypotheses come from estimates for bond sources and from the lattice stress identity. The latter has a remainder, and its conditional mean must be estimated through the renormalisation steps before conservation passes to the limit. We keep these estimates separate from the abstract covariance theorem.

The later sections concern short-distance coefficients, compact gauge groups, local nets, lattice dynamics, Wilson loops and infrared estimates. Their limiting objects and hypotheses are stated separately. In particular, confinement assumes a box inequality at a fixed physical side and every sufficiently large cutoff along one label. No verification of these additional hypotheses is inferred from the short-distance construction. The weighted critical-background assertion in Theorem~\ref{C-360-14}(a) is proved in the regular-reference setting of Definition~\ref{hyp:weighted-critical-response}; its moved-background and centering clauses retain the scalar gauge-transfer hypothesis of Definition~\ref{hyp:complex-scalar-transfer}. The arbitrary-order quantitative torus law is the additional hypothesis in Definition~\ref{hyp:quantitative-torus}; Theorem~\ref{C-360-2}(f), the free-law conclusion in (g)(i), and (g)(ii) retain it. The comparison estimate in (g)(i) and the separately stated narrower results in (h) do not use that additional input.

The general-group construction assumes the estimates in Definition~\ref{C-368-31}. Proposition~\ref{C-368-27} proves the quotient extensions with their stated separation conditions and local bounds.

We use the scalar construction in the companion manuscript \cite{YMM1}, citing its results and pages at each application. The supplementary arguments are given here. Additional estimates required by the later extensions are stated as hypotheses where they enter.


\section{Curvature bilinears: estimates and continuum limits}\label{C-355-1}

We assemble the six curvature bilinears into a symmetric-matrix-valued field. The principal estimate bounds one centred insertion in a chessboard cell by every power greater than $11/2$ of the inverse cell size. It holds uniformly in the cutoff, the admissible datum, the torus side and the periodic Gibbs state.

This estimate gives subsequential limits and bounded insertion operators. The sharper comparison between consecutive cutoffs gives full convergence on a fixed torus. Conservation and the abstract covariance theorem then yield the tensor fields in real time. We state the quantitative results first; the source estimates, comparison arguments, component decomposition and contact extensions follow.

\subsection{Notation and hypotheses}

Theorem~\YMMref{D}{8} gives the inductive description of the effective densities; Theorem~\YMMref{AJ}{673} supplies its sourced form, and Theorem~\YMMref{AD}{607} gives the pressure bounds. The covariance of the tensor limits under $\mathrm O(3)$ is Proposition~\ref{C-355-28}; we state it with its conditions and give the outline of its proof, from the exact stress identity of the lattice, the smallness of its remainder after the renormalisation steps, and a theorem for data. The bound of the conditional mean of the cubic remainder of the stress identity is Proposition~\ref{C-355-30}~(b); we state it and give the five inputs on which it rests, with proofs and precise references for the construction estimates used.

\emph{Lattices.} We work in dimension $3$. $L\ge13$ is an odd integer, $a_i=L^{-i}$, $\mathbb L_K=a_K\mathbb Z^3$, $T_\ell$ is the torus of side $\ell$, and $\Lambda_{K,\ell}$ is the lattice of spacing $a_K$ on $T_\ell$; the gauge field assigns $U(\mathfrak b)\in\mathrm{SU}(2)$ to the bonds $\mathfrak b$, and $a_p=1-\frac12\operatorname{tr}\partial U(p)$ for a plaquette $p$ with boundary loop $\partial U(p)$. $\langle\cdot\rangle_{K,\ell,\beta}$ is the Wilson mean on $\Lambda_{K,\ell}$ at the bare inverse coupling $\beta$, $\langle\cdot;\dots;\cdot\rangle_{K,\ell,\beta}$ the joint cumulant, and $\mathcal G^{\rm per}_{K,\beta}$ the set of the periodic Gibbs states on $\mathbb L_K$ at $\beta$. An observable is centred by subtracting its mean.

\emph{Constants.} Fix one choice of the real and source construction constants and auxiliary data as in Section~\YMMref{82}{665}, at the source radius required by Theorems~\YMMref{AJ}{673} and~\YMMref{AD}{607}, with the numerical source reserve \eqref{eq:canonical-source-quarter-room}. Write $\varsigma^{\rm S}$ for that choice, $\varsigma$ for its unchanged restriction to the real construction, and $\pi^+$ for this restriction. The real construction keeps the same numerical constants; the additional source floors are imposed in their stated order before the coupling ceiling is chosen. The geometric source choice also retains \eqref{eq:array-borel-radius}. The hypotheses are uniform for $x\ge\gamma^{-2}$. In particular, we retain $\gamma^{-2}\ge2B_\sharp$ and $\log\gamma^{-2}>2\mathfrak p_0$.

The constants that occur below are: $L$; $\mathfrak p_0\in\mathbb N$; $A_0$, $M$, $B_\sharp$, $\bar B=B_\sharp L/(L-1)$ and $\gamma\in(0,\frac12]$; $c'$, $C_{W,0}$ and $c_W=C_{W,0}A_0^2$; the ceilings $\gamma^{\rm sd}_W\le\gamma^{\rm ref}_W\le\gamma$, where $(\gamma^{\rm ref}_W)^{-2}=\mathsf x_W$ is the number of equation~\YMMref{(2326)}{845}; the function $p_0(g')=A_0(\log g'^{-2})^{\mathfrak p_0}$; the function $\mathsf B(\cdot)$ of the definition of a small side, with $\mathsf B(x)\ge\frac92\log x+2$; the function $R(\cdot)$ of the condition (Sd.1), which grows like the power $(\log x)^r$ of the logarithm, $r$ a constant; the constant $C_{\rm s}\ge2$ of the definition of the admissible sides; the function $\theta(\cdot)$ of the definitions preceding \eqref{eq:weighted-history-threshold}, a constant times $\mathsf t^{-1/2}(\log\mathsf t)^{\mathfrak p_0}$, and the function $\mathsf T$ of the definitions preceding \eqref{eq:weighted-history-threshold}; the enlargement fraction $\beta_{\rm s}\in(0,\frac15]$ and the decay rate $\kappa>0$ and the fixed radius $\alpha_0\in(0,1]$ of Definition~\ref{def:weighted-progression-data}, where $M=L^{m'}$ is the cube size, and the number $\mathsf c^\flat$ of \eqref{eq:weighted-progression-numbers}; the decay rate enters only through the numbers $T_r(u)$ and the sizes $\|\cdot\|_{j,\kappa}$ of the families of the progression, and elsewhere the letter $\kappa$ is free.

The numbers $T_r(u)$ used in condition (N.$\theta$) are defined below. A \emph{number} is a quantity formed from the constants and from what is displayed as its argument; it does not depend on the cutoff, the bare coupling, the side, the position or the state, unless this is said.

\emph{Couplings and data.} A reference coupling is a real number $g>0$. We use the first-passage definitions and Proposition~\YMMref{74.1}{615}, its room estimate~\YMMref{(1533)}{615}, and Lemma~\ref{lem:setup-stopped-arithmetic}: $X_g=L(g^{-2}+B_\sharp)$; $x_j(\beta)$ is the sequence of running couplings, written as inverse squares, with $x_0(\beta)=\beta$; $g_j=x_j^{-1/2}$ and $y_j=\log x_j$. A pair $(K,\beta)$ is a \emph{first-passage datum at $g$} if $\beta>X_g$ and $K=K(\beta,g)$; then
\[g^{-2}<x_K(\beta)\le X_g<x_{K-1}(\beta)<\dots<x_0(\beta)=\beta .\]
Its \emph{label} is $b=a_K\beta$; by \YMMref{(1531)}{615}, $g^{-2}<b\le L(g^{-2}+\bar B)$. By \YMMref{(1532)}{615} and Lemma~\ref{lem:setup-stopped-arithmetic}, $\log g^{-2}<y_K\le\log g^{-2}+\log\frac{3L}2$ and $K-k\le y_k/\log L$. $\mathfrak L(g)$ is the set of the admissible sides $\ell=2L^m$, $m\ge m_{\mathfrak L}(g)$, of Proposition~\YMMref{74.1}{615}; the least of them is the \emph{least side} at $g$. For $1\le i\le K-1$ the \emph{stopped coupling} $g_{(i)}$ is given by $g_{(i)}^{-2}=L^{-1}x_{K-i}(\beta)-B_\sharp$; by Corollary~\YMMref{74.3}{616} and Lemma~\ref{lem:setup-stopped-arithmetic}, $g_{(i)}<g$, the pair $(K-i,\beta)$ is a first-passage datum at $g_{(i)}$ with label $L^ib$, $L^i\ell\in\mathfrak L(g_{(i)})$ for $\ell\in\mathfrak L(g)$, and
\begin{equation}\label{C-355-2}
C_{g_{(i)}}\le C_g+\tfrac92\,i\log L+\tfrac18,\qquad C_g=C_{\rm st}+\tfrac92\log X_g ,
\end{equation}
Here $C_g=C_{\rm st}+\frac92\log X_g$. With the vacuum constants
$C_\flat,C_{\mathfrak E}$ of equation~\YMMref{(1408)}{575}, put
$E_+=C_{\rm maj}+2$ and
\[
 C_{\rm st}=\max\{C_\flat,C_{\mathfrak E},E_+,
 C_{\mathfrak E}+E_+,c_{\rm low},c_{\rm low}+C_{\mathfrak E}\}.
\]
The constants $c_{\rm low}$ and $C_{\rm maj}$ are those of
Proposition~\YMMref{70.2}{606}. In particular
$C_{\rm st}\ge c_{\rm low}$, so the lower bound of
Proposition~\YMMref{74.4}{617} applies with $C_g$ at first passage.
The \emph{pressure} of the torus $T_\ell$ at $(K,\beta)$ is
$p_\ell=\ell^{-3}\log Z_{K,\ell,\beta}$. We put
\[G_g=C_{\rm S}+C_g+2C_{\rm tt}C'_{\rm tt}E_0,\]
where, by Theorems~\YMMref{AJ}{673} and~\YMMref{AD}{607} at the source $0$, $-C_g\le p_\ell-E_{K,\ell,\beta}\,\ell^{-3}\le C_{\rm S}$ for every first-passage datum $(K,\beta)$ at $g$ and every side $\ell\in\mathfrak L(g)$, $E_{K,\ell,\beta}$ the normalising energy of Theorems~\YMMref{AJ}{673} and~\YMMref{AD}{607} at the source $0$, and $2C_{\rm tt}C'_{\rm tt}E_0$ bounds the difference of the numbers $E_{K,\ell,\beta}\,\ell^{-3}$ at two such sides; $C_{\rm S}$, $C_{\rm tt}$, $C'_{\rm tt}$, $E_0$ depend on no coupling, cutoff, side or source.

\emph{Boards and the chessboard seminorm.} A board $(i,o)$ of $\mathbb L_K$ is given by an integer $i\le K-1$ and an offset $o\in a_K\mathbb Z^3$; it is the partition of $\mathbb R^3$ into the cubes $\square_{\mathbf k}=o+a_i\mathbf k+[-\frac{a_i}2,\frac{a_i}2]^3$, $\mathbf k\in\mathbb Z^3$, its \emph{cells}, whose faces lie on the \emph{wall planes} $\{x_\mu=o_\mu+a_i(t+\frac12)\}$, $\mu\in\{1,2,3\}$, $t\in\mathbb Z$. As $a_i/a_K=L^{K-i}$ is odd, $o_\mu+a_i(t+\frac12)\in a_K(\mathbb Z+\frac12)$: the wall planes are mid-planes of the lattice, no site lies on them, and the reflections in them are isometries of $\mathbb L_K$. The mirror map from $\square_{\mathbf k}$ to $\square_{\mathbf k'}$, $x\mapsto o+a_i\mathbf k'+\operatorname{diag}(\sigma_1,\sigma_2,\sigma_3)(x-o-a_i\mathbf k)$ with $\sigma_\mu=(-1)^{\mathbf k'_\mu-\mathbf k_\mu}$, is the composition of the $|\mathbf k'-\mathbf k|_1$ reflections in the wall planes between the two cells, and the pair of cells is \emph{even} or \emph{odd} with this number.

The mirror copy in $\square_{\mathbf k'}$ of a bounded complex function $F$ of the bonds inside $\square_{\mathbf k}$ is the image of $F$ under the mirror map if the pair is even, and the complex conjugate of this image if the pair is odd; for a real $F$ it is the image in every cell. For a side $\ell$ with $\ell/a_i$ an even integer, the mean on $T_\ell$ of the product of the mirror copies of $F$ over the $(\ell/a_i)^3$ cells of a period box is $\ge0$, and the \emph{chessboard seminorm} $\|F\|_\ell=\|F\|^{(i,o)}_{{\rm ch},\ell}$ is its $(\ell/a_i)^3$-th root (Proposition~\YMMref{86.3}{681}); the seminorm of a constant is its modulus. We use the chessboard estimate, seminorm properties and volume comparison of Propositions~\YMMref{86.3}{681} and~\YMMref{86.4}{682}, and Corollary~\YMMref{86.6}{683}: $\|\cdot\|_\ell$ is a seminorm, $|\langle F\rangle_{K,\ell,\beta}|\le\|F\|_\ell$, the mean on $T_\ell$ of a product of functions of different cells modulo the side is at most the product of their seminorms, and for two sides $\ell\ge\ell'\ge2a_i$
\[\|F\|_\ell\le e^{a_i^3(p_{\ell'}-p_\ell)}\,\|F\|_{\ell'} ;\]
and the domination of the periodic Gibbs states, Proposition~\YMMref{91.9}{705}~(b): for $\nu\in\mathcal G^{\rm per}_{K,\beta}$, a board, finitely many different cells and functions $\Phi_k$ of the bonds inside them,
\[\Bigl|\nu\Bigl(\prod_k\Phi_k\Bigr)\Bigr|\le\liminf_{\ell\to\infty}\prod_k\|\Phi_k\|_\ell ,\]
the lower limit along the sides of $\mathfrak L(g)$.

\emph{The space of labels.} $V=\mathrm{Sym}_3(\mathbb R)$, $V_{\mathbb C}=\mathrm{Sym}_3(\mathbb C)$, $V_0=\mathbb R\mathbf 1$; $\mathrm O(3)$ acts by $\rho(\mathsf R)A=\mathsf RA\mathsf R^{\rm T}$. We use the two norms $|A|_\infty=\max_{\lambda\lambda'}|A_{\lambda\lambda'}|$ and $|A|=(\operatorname{tr}\bar AA)^{1/2}$, which is invariant; $|A|_\infty\le|A|$. In sup norms $|H|_\infty=\sup_x|H(x)|_\infty$ the norm of $V$ is $|\cdot|_\infty$; in $\|h\|_{L^1}=\int|h(x)|\,dx$ it is the invariant norm unless the contrary is said. $W_3$ is the group of the $48$ signed permutation matrices, and $\theta=\operatorname{diag}(1,1,-1)$. For $1\le i,j\le3$ the \emph{label matrices of the stress tensor} are
\[
\begin{multlined}
\mathsf m_{ij}=\delta_{ij}\mathbf 1-e_ie_j^{\rm T}-e_je_i^{\rm T};\qquad\rho(\mathsf R)\mathsf m_{ij}=\sum_{i',j'}\mathsf R_{i'i}\mathsf R_{j'j}\,\mathsf m_{i'j'}\ \relax \ \relax (\mathsf R\in\mathrm O(3)),\qquad\theta\mathsf m_{ij}\theta=\vartheta_i\vartheta_j\mathsf m_{ij}, \\
|\mathsf m_{ij}|\le\sqrt3,
\end{multlined}
\]
with $\vartheta_1=\vartheta_2=1$, $\vartheta_3=-1$; for a symmetric matrix $T$, $\frac12\operatorname{tr}(\mathsf m_{ij}T)=\frac12\delta_{ij}\operatorname{tr}T-T_{ij}$.

\begin{definition}[Lattice forms]\label{C-355-3}
We use the tensor form \eqref{sf-lat} and write $\mathcal T_K(H)=O_K^T(H)$. For arbitrary plaquette weights $w$ and bond weights $\mathsf w$, put
\[
 \mathcal O(w)=a_K^{-1}\sum_pw(p)a_p,\qquad
 Y[\mathsf w]=-a_K^{-1}\sum_{\mathfrak b}\mathsf w(\mathfrak b)q_{\mathfrak b}.
\]
Thus $\mathcal T_K(H)=\mathcal O(w_H)+Y[\mathsf w_H]$, with $w_H(p)=H_{\lambda(p)\lambda(p)}(x(p))$ and $\mathsf w_H(\mathfrak b)=H_{\lambda_{\mathfrak b}\lambda'_{\mathfrak b}}(x_{\mathfrak b})$. The midpoint and plaquette-normal conventions, $q_{\mathfrak b}$ and $u_\alpha=u_{\mathfrak b,\alpha}$ are those of the introduction. The scalar form $\mathcal T_K(f\mathbf1)=\mathcal O(f)$ is the action-density observable.
\end{definition}

\emph{Sources.} Theorem~\YMMref{D}{8} describes on a small torus the effective densities that the renormalisation steps produce from the lattice measure. With sources on the plaquettes in the exponent, $\sum_a\lambda_a\mathcal O(w_a)$ with complex parameters $\lambda_a$, the steps are carried by the weighted progression of Definition~\ref{def:weighted-progression-data}, in which the source $w_a$ enters at the height $j$ through a family of the size $\mathsf m_{j,a}=12M^3\alpha_0^2L^{K-j}|w_a|_\infty$ of \eqref{eq:weighted-progression-numbers}. A source under a majorant is what this progression carries in the place of a plaquette weight. A \emph{majorant} is a pair $(\mathsf s,S)$ of a number $\mathsf s\ge0$ and a finite set $S$ of plaquettes; its weight is the plaquette weight $\tilde w=\mathsf s\mathbf 1_S$, and its numbers are $\mathsf m_j=12M^3\alpha_0^2L^{K-j}\mathsf s$, $0\le j\le k$, the sizes \eqref{eq:weighted-progression-numbers} at $\tilde w$, $k$ being the last height of the progression.

A \emph{source under the majorant} $(\mathsf s,S)$ is a triple $(F,\mathbb V,(H^{(j)})_{1\le j\le k})$ such that: (So0) $F$ is a continuous gauge-invariant function of the bare configurations with $|F(U)|\le L^K\sum_p\tilde w(p)\,a_p(U)$; (So1) $\mathbb V$ is a family in the fourth place at the height $0$ (Section~\ref{sec:step-input-conventions}) of size at most $\mathsf c^\flat\mathsf m_0$, whose members vanish on the domains whose plaquettes miss $S$, and which sums to $F$ at the configurations of the height $0$, as the family of the action density of a plaquette weight $w$ sums to $\mathcal O(w)$ under the height-zero global sum condition \eqref{eq:fourth-family-global-sum}; (So2) for $1\le j\le k$, $H^{(j)}$ is a family without base point of level $j$ at the fixed radii (Definition~\YMMref{49.5}{375}) with $\|H^{(j)}\|_{j,\kappa}\le\frac12\mathsf m_j$, whose members $H^{(j)}(X,\cdot)$ vanish unless the plaquettes of $X$ meet $S$, and $\sum_{X\in\mathcal D_j}[H^{(j)}(X,U)-H^{(j)}(X,\mathbf 1)]=F(U)$ at every bare configuration $U$ at which the action of the bare arrays of the height $j$ is defined.

The function $F$ is the \emph{first member} of the source, and, the triple being understood, $F$ itself is also called a source under $(\mathsf s,S)$; the source is \emph{real} if $F$ is real. A \emph{source system} is a finite tuple of sources $(F_a,\mathbb V_a,(H^{(j)}_a))$ under majorants $(\mathsf s_a,S_a)$, entered with independent complex parameters $\lambda_a$ in the places of the plaquette sources of Definition~\ref{def:weighted-progression-data}, the exponent being $\sum_a\lambda_aF_a$ and all the numbers of the progression being formed at the weights $\mathsf s_a\mathbf 1_{S_a}$; its first members are the $F_a$. The bounds of Propositions~\ref{C-355-19} and~\ref{C-355-20} read of a source the properties (So0)--(So2) alone.

The \emph{working level} of a datum, the \emph{history} without large-field regions and the weights $\mathsf w_t$ of the first creation of a large field at the step $t$ are those among the construction results; so are the \emph{small sides}, the sides that are small over the running couplings of a datum, among whose conditions is (Sd.4), $\mathsf B(x_{K})\ell^3\le\frac12c'p_0(g_{K})^2$; the least side at $g$ is small for $g\le\gamma^{\rm sd}_W$. The \emph{depth} $s(d)$ of a distance $0<d<1$ is the least $s$ with $L^sd\ge C_W(s)=c_WY_s^{2\mathfrak p_0}$, $Y_s=(s+1)\log L+\log(g^{-2}+\bar B)$, formed at the reference coupling of the datum, with the room constant $C_{W,0}=\max\{1,4L^2M/A_0^2\}$ in $c_W=C_{W,0}A_0^2$, so that $C_W(s)\ge4L^2MY_s^{2\mathfrak p_0}$; we write $\bar x_s=e^{Y_s}$.

For $u>0$ and an integer $r\ge1$ we put
\begin{equation}\label{eq:source-mixed-series}
T_r(u)=\sum_{n'\ge1}(n'+1)L^{-(n'-1)}\Bigl[\bigl(\tfrac43\bigr)^{1/2}L^{n'/2}(1+n')^{\mathfrak p_0}\Bigr]^re^{-4\beta_{\rm s}\kappa[L^{n'-1}u-2]_+},\qquad[v]_+=\max\{v,0\};\end{equation}
the series converges, the quotient of two consecutive terms tending to $0$; $T_r(u)$ does not increase in $u$ and does not decrease in $r$, the bracket being $\ge1$; and $T_2=\frac43L\,\mathsf T$ with the function $\mathsf T$ of the definitions preceding \eqref{eq:weighted-history-threshold}. The condition (N.$\theta$) at $(n,g,s)$ is $2T_n(3)\theta(L^{s-2}g^{-2})\le1$; as $T_r(3)$ does not decrease in $r$, (N.$\theta$) at $(n,g,s)$ gives (N.$\theta$) at $(m,g,s)$ for $2\le m\le n$.

\begin{definition}[Tuples, tensor limits, the reconstructed space]\label{C-355-4}
\begin{enumerate}
\item[(1)] For an open $O\subset\mathbb R^3$, $\mathfrak S_V(O)$ is the set of the tuples $X=(H_1,\dots,H_n)$, $n\ge0$, of continuous $V$-valued functions with compact, pairwise disjoint supports in $O$; $\mathfrak S^1_V(O)$ is the same with $C^1_c$ entries; $\varnothing$ is the empty tuple. $\mathbb C\mathfrak S_V(O)$ is the space of the formal complex combinations of tuples, $F\cdot F'$ the concatenation (formed at disjoint supports), $\bar F$ the conjugate combination, and $(h)$ the tuple with the one entry $h$; an entry $h=h'+ih''$ with values in $V_{\mathbb C}$ stands for $(h')+i(h'')$. An isometry $\mathsf rx=\mathsf Rx+v$ acts by $(\mathsf rH)(x)=\mathsf RH(\mathsf r^{-1}x)\mathsf R^{\rm T}$, entry by entry; $\tau_v$ is the translation and $\theta$ acts as the reflection in $\{x_3=0\}$, on points and on labels. $\mathbb E_+=\{x_3>0\}$.
\item[(2)] A \emph{tensor limit at $g$} is a functional on $\mathfrak S_V(\mathbb R^3)$ of the form
\[\sigma(X)=\lim_{j\to\infty}\rho_j\Bigl(\prod_a\bigl[\mathcal T_{K_j}(H_a)-\rho_j(\mathcal T_{K_j}(H_a))\bigr]\Bigr),\]
the limit existing at every tuple, where $(K_j,\beta_j)$ are first-passage data at $g$ with $K_j\to\infty$, $\rho_j\in\mathcal G^{\rm per}_{K_j,\beta_j}$, and the labels converge, $a_{K_j}\beta_j\to b_*$; $b_*$ is the label of $\sigma$. Its restriction to $\mathfrak S^1_V(\mathbb R^3)$ is written $\sigma^1$.
\item[(3)] Given that $\sigma(\theta\bar F\cdot F)\ge0$ on $\mathbb C\mathfrak S_V(\mathbb E_+)$, $\mathcal H(\sigma)$ is the completion of $\mathbb C\mathfrak S_V(\mathbb E_+)$ modulo the null vectors of $\langle[F'],[F]\rangle=\sigma(\theta\bar F'\cdot F)$; $\Omega=[\varnothing]$; $e^{-tH}[F]=[\tau_{te_3}F]$ and $\mathcal U(v)[F]=[\tau_{(v,0)}F]$ define $H\ge0$ and a unitary representation $\mathcal U$ of $\mathbb R^2$, and $\mathcal P$ is the joint spectral measure on $[0,\infty)\times\mathbb R^2$ with $e^{-tH}\mathcal U(v)=\int e^{-tE+iv\cdot q}d\mathcal P(E,q)$; the momenta are $P_k=\int q_k\,d\mathcal P$, $k=1,2$, so that $\mathcal U(v)=e^{iv\cdot P}$. The \emph{square cone} is $\{E\ge\max(|q_1|,|q_2|)\}$, the \emph{circular cone} $\{E\ge|q|\}$. $\mathcal H_0\subset\mathcal H(\sigma)$ is the closed span of the classes of the tuples with entries in $V_0$. $\mathcal H^{\rm ord}(\sigma)$ is the closed span of $\Omega$ and of the classes of the $X\in\mathfrak S_V(\mathbb E_+)$ whose entries lie in pairwise disjoint slabs $\{t<x_3<t'\}$.
\item[(4)] For $l>0$ and $h$ carried by $\{0<x_3<l\}$, the \emph{slab operator} $\mathsf B_l(h)$ is the bounded operator, if there is one, with $\mathsf B_l(h)[F]=[(h)\cdot\tau_{le_3}F]$ for all $F\in\mathbb C\mathfrak S_V(\mathbb E_+)$.
\item[(5)] Let $\sigma$ be a functional on $\mathfrak S_V(\mathbb R^3)$ that is invariant under the permutations of the entries, $q\ge1$ and $X^1,\dots,X^q\in\mathfrak S_V(\mathbb R^3)$. For $1\le p\le q$ let $\mathsf A^{(p)}$ be the set of the $\underline v=(v_{p+1},\dots,v_q)\in(\mathbb R^3)^{q-p}$ at which the unions of the supports of the entries of $X^1,\tau_{v_{p+1}}X^{p+1},\dots,\tau_{v_q}X^q$ are pairwise disjoint, so that the concatenation $X'_p(\underline v)=X^1\cdot\tau_{v_{p+1}}X^{p+1}\cdots\tau_{v_q}X^q$ is formed, and put $G^{(1)}(\underline v)=\sigma(X'_1(\underline v))$ on $\mathsf A^{(1)}$. Then $\sigma$ \emph{has iterated plateaux} at $(X^1,\dots,X^q)$ if, successively for $p=2,\dots,q$, at every $\underline v\in\mathsf A^{(p)}$ the limit
\[G^{(p)}(\underline v)=\lim_{|v_p|\to\infty}G^{(p-1)}(v_p,\underline v),\qquad v_p\in\mathbb R^3,\]
exists; $\mathsf A^{(q)}$ is a point, and the number $G^{(q)}$ is the \emph{iterated plateau} of $\sigma$ at $(X^1,\dots,X^q)$. At $q=1$ the iterated plateau is $\sigma(X^1)$; at $q=2$ it is the two-fold plateau $\lim_{|v|\to\infty}\sigma(X^1\cdot\tau_vX^2)$.
\end{enumerate}
\end{definition}

\begin{definition}[Data with labels; real time; the frame]\label{C-355-5}
\begin{enumerate}
\item[(1)] A \emph{symmetric functional on labelled tuples} is a map $\sigma$ on $\mathfrak S^1_V(\mathbb R^3)$ with $\sigma(\varnothing)=1$, invariant under the permutations of the entries and $n$-linear on $\prod_aC^1_c(O_a;V)$ for open sets $O_a$ with pairwise disjoint closures; it is continued linearly to $\mathbb C\mathfrak S^1_V$. It is \emph{reconstruction data with labels} with the numbers $C_V,p_V\ge0$, $c_0\in(0,\frac12]$, if: $(\mathrm{Rec1})_V$ $\sigma(\mathsf rX)=\sigma(X)$ for every isometry $\mathsf r$; $(\mathrm{Rec2})_V$ $\sigma(\theta\bar F\cdot F)\ge0$ on $\mathbb C\mathfrak S^1_V(\mathbb E_+)$, the space $\mathcal H(\sigma)$ is separable, and $H$, $\mathcal U$, $\mathcal P$ are formed as in Definition~\ref{C-355-4}~(3); $(\mathrm{Rec3})_V$ $\mathcal P(\{E<|q|\})=0$; $(\mathrm{Rec4})_V$ for $\eta\in(0,1]$, $F,F'\in\mathbb C\mathfrak S^1_V(\{x_3>\eta\})$ and $h\in C^1_c(\{|x_3|<c_0\eta\};V)$,
\[|\sigma(\theta\bar F\cdot(h)\cdot F')|\le C_V\eta^{-p_V}\|h\|_{L^1}\,\|[\tau_{-\eta e_3}F]\|\,\|[\tau_{-\eta e_3}F']\| .\]
(Rec5) is the condition $\ker H=\mathbb C\Omega$; it is not part of the data.
\item[(2)] A point of $\mathbb R^{1+2}$ is $x=(t,\vec x)$; $U(x)=e^{itH}\mathcal U(\vec x)$; $\mathcal D_\infty=\bigcap_j\operatorname{dom}H^j$; $\mathcal S_V=\mathcal S(\mathbb R^{1+2};V_{\mathbb C})$ and $f^*=\rho(\theta)\bar f$. $\iota(t,\vec x)=(\vec x,-it)$; $\mathcal L$ is the restricted Lorentz group of $\mathbb R^{1+2}$, and $\hat\Lambda=\iota\Lambda\iota^{-1}\in\mathrm{SO}(3,\mathbb C)$. For $\mathfrak z\in\mathbb C^3$, $d_\square(\mathfrak z)=\operatorname{Re}\mathfrak z_3-|\operatorname{Im}\mathfrak z_1|-|\operatorname{Im}\mathfrak z_2|$ and $d_\circ(\mathfrak z)=\operatorname{Re}\mathfrak z_3-|(\operatorname{Im}\mathfrak z_1,\operatorname{Im}\mathfrak z_2)|$.
\item[(3)] The \emph{frame} (M1)--(M5) concerns a tensor carrying one label on a space with a scalar field $\phi(f)$, $f\in\mathcal S(\mathbb R^{1+2})$, a unique vacuum and a representation $U$ of $\mathbb R^{1+2}\rtimes\mathcal L$; with $\vartheta_1=\vartheta_2=1$, $\vartheta_3=-1$: (M1) sesquilinear forms $\mathfrak q_{ij}=\mathfrak q_{ji}$ on $\mathcal D_\infty\times\mathcal D_\infty$ with $|\mathfrak q_{ij}(\psi',\psi)|\le c_T\|(1+H)^l\psi'\|\,\|(1+H)^l\psi\|$, $\overline{\mathfrak q_{ij}(\psi',\psi)}=\vartheta_i\vartheta_j\mathfrak q_{ij}(\psi,\psi')$, $\mathfrak q_{ij}(\Omega,\Omega)=0$, and the operators $T^E_{ij}(F)$ they define, with the real-time components $T_{00}=-T^E_{33}$, $T_{0j}=-iT^E_{3j}$, $T_{jk}=T^E_{jk}$; (M2) continuous kernels $\mathsf s_{ij}(x;y_1,\dots,y_m)$ of one tensor entry among $m$ scalar entries, covariant under $\mathrm O(3)$ and the translations, $\mathsf s_{ij}(\mathsf Rx+b;\mathsf Ry+b)=\mathsf R_{ii'}\mathsf R_{jj'}\mathsf s_{i'j'}(x;y)$, related to the forms by the product formula of two point operators at different heights, $\mathsf s^{F',G}_{ij}(x,y)=\mathfrak q_{ij}(\psi'_x,\psi_x)$ as written out in the proof of Lemma~\ref{C-355-31}~(c); (M3) $T^E_{ij}(F_1)\phi(f)\psi=\phi(f)T^E_{ij}(F_1)\psi$ on the span $\mathcal D_0$ of $\Omega$ and the vectors $\phi(f_1)\cdots\phi(f_n)\Omega$, for spacelike separated supports; (M4)~(a) $T_{0\nu}(\partial_0F)\psi=\sum_{k=1,2}T_{k\nu}(\partial_kF)\psi$ on $\mathcal D_0$; (M5) $U(\Lambda)T[a,b](F)\Omega=T[\Lambda a,\Lambda b](F\circ\Lambda^{-1})\Omega$, $T[a,b]=a^\mu b^\nu T_{\mu\nu}$.
\end{enumerate}
\end{definition}

\begin{definition}[The setting of the cumulants on a torus]\label{C-355-6}
Let $n\ge2$, $0<g\le\gamma^{\rm ref}_W$, $(K,\beta)$ a first-passage datum at $g$ with label $b$, and $\ell$ a side. For bounded $H_a:T_\ell\to V_{\mathbb C}$, $\Sigma_a$ is the set of the points $x(p)$, $x_{\mathfrak b}$ at which the component of $H_a$ that enters $\mathcal T_K(H_a)$ is not $0$. The tuple has the property \emph{(Geo) at $d$} if each $\Sigma_a$ lies in a ball of radius $0.87\,d$ and any two of the $\Sigma_a$ are at distance $\ge1.26\,d$. We put $\bar x_0=L(g^{-2}+\bar B)$ and $\mathsf c(\ell)=2\ell^3e^{\mathsf B(\bar x_0)\ell^3}/(c'A_0^2)$.

For the extensions over the diagonals, on $T_\ell$ with the sup distance: for $\lambda>0$ and an integer $m\ge0$, $\|F\|_{(\lambda),m}=\max_{|\gamma|\le m}\lambda^{|\gamma|}\sup|\partial^\gamma F|$; a function $F$ on $T_\ell^n$ is \emph{$\lambda$-supported} if its support lies in a product $\prod_a\bar B_\infty(z_a,\lambda)$ of closed sup balls, the points $z_a\in T_\ell$ coinciding or not. For a real number $\kappa$, the parameter of the power counting, $\mathfrak E_n(\kappa)$ is the class of the distributions $T$ on $T_\ell^n$, acting on the $V_{\mathbb C}^{\otimes n}$-valued test functions, for which there are $C,q,m$ and $h_1>0$ with
\[|T(F)|\le C\lambda^\kappa\bigl(1+\log\tfrac{h_1}\lambda\bigr)^q\|F\|_{(\lambda),m}\qquad\text{for all }0<\lambda\le h_1\text{ and all smooth $\lambda$-supported }F:\]
the power counting $\kappa$, up to logarithms, also where points meet. For a partition $\pi$ of $\{1,\dots,n\}$ into $k$ blocks, $B(a)$ the block of $a$ and $a_B=\min B$, the \emph{coordinates at $\pi$} are $y_B=x_{a_B}\in T_\ell$ and $\xi_a=x_a-x_{a_B}$ for $a\ne a_{B(a)}$, taken where these differences are below $\ell/4$; $\Delta_\pi=\{\xi=0\}$, and $\iota_\pi(y)_a=y_{B(a)}$. For a distribution $T'$ on $T_\ell^k$ and a multi-index $\alpha$ of $\xi$, $(\partial_\xi^\alpha\delta_\pi\otimes T')(F)=(-1)^{|\alpha|}\,T'\bigl((\partial_\xi^\alpha F)\circ\iota_\pi\bigr)$.
\end{definition}

\subsection{Principal results}

\begin{theorem}[The estimate at a cutoff]\label{C-355-7}
Let $0<\delta\le\frac14$.
\begin{enumerate}
\item[(a)] Let $0<g\le\gamma^{\rm sd}_W$. Put
\[
\begin{multlined}
c_E=\log(L^2c_W/\delta)+2\mathfrak p_0\log(4\mathfrak p_0)+\log(X_g+\bar B),\qquad\mathsf Y_i=2[c_E+\log(1/a_i)],\qquad q_E \\
=(4+\mathfrak p_0)(2\mathfrak p_0-1)+10\mathfrak p_0,
\end{multlined}
\]
\[C_E=510\cdot2^{2\mathfrak p_0+1}\,\delta^{-4}e^{G_g+1/8}\,(1+c'A_0^2/9)^{4+\mathfrak p_0}\max\{C^I,C_2\}\,L^5c_W^4\,(X_g+\bar B),\]
where $C^I$, $C_2$ are numbers formed from the constants alone, which enter through Proposition~\ref{C-355-19}. There are numbers $i_E=i_E(g,\delta)\ge1$ and $s_E(i)=s_E(i;g,\delta)$ such that the following holds for every $i\ge i_E$, every first-passage datum $(K,\beta)$ at $g$ with $K\ge K_E(i):=i+s_E(i)$, every board $(i,o)$ of $\mathbb L_K$, every cell $\square$ of it, every $\ell\in\mathfrak L(g)$ and every bounded $H:\mathbb R^3\to V_{\mathbb C}$ that vanishes at the points whose sup distance from the complement of $\square$ is $<\delta a_i$: if $c$ is the mean of $\mathcal T_K(H)$ on a torus $T_{\ell_1}$ with $\ell_1\in\mathfrak L(g)$, or in a state $\nu\in\mathcal G^{\rm per}_{K,\beta}$, then
\begin{equation}\label{C-355-8}
\|\mathcal T_K(H)-c\|^{(i,o)}_{{\rm ch},\ell}\le2C_E\,\mathsf Y_i^{\,q_E}\,a_i^{-11/2}\,b^{-1}|H|_\infty .
\end{equation}
The dependence on $\delta$ is by powers: $i_E$ may be taken with $L^{i_E}\le c_i\delta^{-1}$, and $C_E\mathsf Y_i^{q_E}a_i^{-11/2}\le c_3\delta^{-5}a_i^{-13/2}$ for all $i\ge1$, where $c_i$, $c_3$ are numbers formed from the constants and $g$.
\item[(b)] Let $0<g\le\gamma^{\rm sd}_W$. Put $c'_E=\log(2L(g^{-2}+\bar B))$, $\mathsf Y'_i=2[c'_E+\log(1/a_i)]$, $q'_E=15r+4$,
\[
\begin{multlined}
c_N=\max\bigl\{1,\,8L^6M^3C_{\rm s}^3\,2^{-3r}\bigr\},\qquad c_u=4+\tfrac2{A_\flat}\bigl(2c_N[c_2+\log(8c_N)+3r]+3\log L \\
+\log3\bigr),
\end{multlined}
\]
with the numbers $A_\flat$, $c_2$, $C_1$ and the threshold function $u'_*$ of Proposition~\ref{C-355-19}~(b); let $s_i$ be the least integer $s\ge1$ with $\delta L^s\ge M\,u'_*(\lceil c_N\mathsf Y_i'^{\,3r}\rceil)$, and 
\[
C'_E=480\,e^{G_g+1/8}c_NC_1(LMc_u/\delta)^4L(g^{-2}+\bar B).
\]
Then the bound
\[\|\mathcal T_K(H)-c\|^{(i,o)}_{{\rm ch},\ell}\le2C'_E\,\mathsf Y_i'^{\,q'_E}\,a_i^{-11/2}\,b^{-1}|H|_\infty\]
holds for every integer $i\ge1$ and every first-passage datum with $K\ge i+s_i+2$, with the other quantifiers of (a).
\item[(c)] Let $0<g\le\gamma$. There are numbers $i_A(g,\delta)$ and, for $i\ge0$, $K_0(i)$ and $\mathsf N_i$, formed from the constants, $g$, $\delta$ and $i$, such that: (i) for $i\ge i_A(g,\delta)$, every first-passage datum at $g$ with $K\ge K_E(i)$ (formed at $g$ by the formulas of (a)), every board, cell and side $\ell\in\mathfrak L(g)$, every bounded Borel $H:\mathbb R^3\to V_{\mathbb C}$ vanishing at sup distance $<\delta a_i$ from the complement of the cell, and both centrings of (a),
\[\|\mathcal T_K(H)-c\|^{(i,o)}_{{\rm ch},\ell}\le2e^{1/8}C_E\,\mathsf Y_i^{\,q_E}L^{11i/2}\,b^{-1}|H|_\infty,\qquad b^{-1}<g^2,\]
with $C_E$, $\mathsf Y_i$ formed at $g$ by the formulas of (a); (ii) for every integer $i\ge0$, every first-passage datum at $g$ with $K\ge K_0(i)$, every board $(i,o)$, cell $\square$ and side $\ell\in\mathfrak L(g)$, every real bounded $h$ that vanishes off $\square$ and at the points within sup distance $a_i/32$ of its boundary, and both centrings,
\[\|\mathcal T_K(h)-c\|^{(i,o)}_{{\rm ch},\ell}\le\mathsf N_i\,b^{-1}|h|_\infty\le C_TL^{ip_T}|h|_\infty\qquad\text{for every }p_T>\tfrac{11}2,\quad C_T=g^2\sup_{i\ge0}\mathsf N_iL^{-ip_T}<\infty ;\]
(iii) \emph{(uniform bounds)} for $n\ge1$, $d>0$, $R>0$ let $\delta_n=1/(8n)$, $i^A$ the least $i\ge i_A(g,\delta_n)$ with $2a_i\le d$, $K^A$ the least $K\ge K_E(i^A)$ (formed at $g$ and $\delta_n$ by the formulas of (a)) with $L^{K-i^A}\ge8n$, and
\[C^A_n(d;R)=\Bigl[\bigl(16nR/a_{i^A}+3\bigr)^3\,2e^{1/8}C_E(\delta_n)\,\mathsf Y_{i^A}^{\,q_E}\,a_{i^A}^{-11/2}\Bigr]^n ,\]
$C_E$ and $\mathsf Y_i$ formed at $g$; then for every first-passage datum $(K,\beta)$ at $g$ with $K\ge K^A$, any two states $\nu,\nu'\in\mathcal G^{\rm per}_{K,\beta}$ and all bounded $H_1,\dots,H_n:\mathbb R^3\to V_{\mathbb C}$ whose supports lie in balls of radius $R$ and are pairwise at distance $\ge d$,
\[\Bigl|\nu\Bigl(\prod_a\bigl[\mathcal T_K(H_a)-\nu'(\mathcal T_K(H_a))\bigr]\Bigr)\Bigr|\le C^A_n(d;R)\,b^{-n}\prod_a|H_a|_\infty .\]
\end{enumerate}
All bounds are uniform in the cutoff, the first-passage datum, the board, the cell, the side (hence in the number of the cells) and the state.
\end{theorem}

\begin{theorem}[The tensor limits without the rotations]\label{C-355-9}
Let $0<g\le\gamma^{\rm sd}_W$, and let the numbers of Theorem~\ref{C-355-7}~(a) be taken at the value of $\delta$ that is displayed.
\begin{enumerate}
\item[(a)] \emph{(Uniform bounds.)} For $n\ge1$, $d>0$, $R>0$ put $\delta_n=1/(8n)$; let $i^G$ be the least $i\ge i_E(g,\delta_n)$ with $2a_i\le d$, $K^G$ the least $K\ge K_E(i^G)$ (at $\delta_n$) with $L^{K-i^G}\ge8n$, and
\[C^G_n(d;R)=\Bigl[\bigl(16nR/a_{i^G}+3\bigr)^3\,2C_E(\delta_n)\,\mathsf Y_{i^G}^{\,q_E}\,a_{i^G}^{-11/2}\Bigr]^n .\]
For every first-passage datum $(K,\beta)$ at $g$ with $K\ge K^G$, every $\nu\in\mathcal G^{\rm per}_{K,\beta}$ and all bounded $H_1,\dots,H_n:\mathbb R^3\to V_{\mathbb C}$ whose supports lie in balls of radius $R$ and are pairwise at distance $\ge d$,
\[\Bigl|\nu\Bigl(\prod_a\bigl[\mathcal T_K(H_a)-\nu(\mathcal T_K(H_a))\bigr]\Bigr)\Bigr|\le C^G_n(d;R)\,b^{-n}\prod_a|H_a|_\infty ,\]
and the same holds with the centring by the mean on any torus $T_\ell$, $\ell\in\mathfrak L(g)$. If $d\le2a_{i_E(g,\delta_n)}$, then $C^G_n(d;R)\le\bigl[(32nLR/d+3)^3\,2C_E(\delta_n)\,(2[c_E(\delta_n)+\log(2L/d)])^{q_E}(2L/d)^{11/2}\bigr]^n$.
\item[(b)] \emph{(Limits; the reconstructed space.)} Every sequence of first-passage data $(K_j,\beta_j)$ at $g$ with $K_j\to\infty$, with states $\rho_j\in\mathcal G^{\rm per}_{K_j,\beta_j}$, has a subsequence along which the tensor limit exists at every tuple. Every tensor limit $\sigma$ at $g$ has a label $b_*\in[g^{-2},L(g^{-2}+\bar B)]$; it is real, symmetric and multilinear, $\sigma(\varnothing)=1$, $\sigma((H))=0$, $|\sigma(X)|\le C^G_n(d;R)\,g^{2n}\prod_a|H_a|_\infty$ for the tuples of (a), and $\sigma$ is continuous under the motion of its entries; it is invariant under $\mathbb R^3\rtimes W_3$ and $\sigma(\theta\bar F\cdot F)\ge0$ on $\mathbb C\mathfrak S_V(\mathbb E_+)$. Hence $\mathcal H(\sigma)$, $\Omega$, $H$, $\mathcal U$, $\mathcal P$ are formed; $\mathcal P$ is carried by the square cone; $\mathcal H_0$ reduces $H$, $\mathcal U$, $\mathcal P$.

With $i_1=i_E(g,\frac1{16})$ and $\mathsf c_E(i)=2C_E\mathsf Y_i^{q_E}a_i^{-11/2}$ at $\delta=\frac1{16}$: for $i\ge i_1$, $z\in\mathbb R^3$ with $z_3=a_i/2$ and $h\in C_c(\mathbb R^3;V_{\mathbb C})$ carried by $\{|x-z|_\infty\le\frac38a_i\}$, the slab operator $\mathsf B_{a_i}(h)$ exists and $\|\mathsf B_{a_i}(h)\|\le\mathsf c_E(i)\,b_*^{-1}|h|_\infty$; and wherever slab operators exist, $e^{-tH}\mathsf B_l(h)e^{-t'H}=\mathsf B_{l+t+t'}(\tau_{te_3}h)$ and $\mathcal U(v)\mathsf B_l(h)\mathcal U(v)^{-1}=\mathsf B_l(\tau_{(v,0)}h)$.
\item[(c)] \emph{(The point bound.)} For every $p>\frac{17}2$ there is a number $C_p$ formed from the constants, $g$ and $p$ such that for $0\le d'<d$, $\vartheta=d-d'$ and $h\in C_c(\mathbb R^3;V_{\mathbb C})$ carried by $\{|x_3|\le d'\}$,
\[\|\mathsf B_{2d}(\tau_{de_3}h)\|\le C_p\min\{1,\vartheta\}^{-p}\,b_*^{-1}\|h\|_{L^1},\]
either norm of $V$ being used in $\|h\|_{L^1}$. Consequently $\sigma^1$ satisfies $(\mathrm{Rec4})_V$ with $c_0=\frac12$, every $p_V>\frac{17}2$ and $C_V=2^{p_V}C_{p_V}b_*^{-1}\le2^{p_V}C_{p_V}g^2$: one triple $(2^{p_V}C_{p_V}g^2,p_V,\frac12)$ serves all the tensor limits at $g$. In the sup norm of a cell every power $>\frac{11}2$ is admissible.
\item[(d)] \emph{(Real time.)} Let $p=p_V$, $C_V$ be as in (c) and $l$ the least integer $>p$. (i) $\mathcal T(\mathfrak z)=\int e^{-\mathfrak z_3E+i(\mathfrak z_1q_1+\mathfrak z_2q_2)}d\mathcal P$ are contractions on $\{d_\square\ge0\}$, form a semigroup, and are holomorphic in norm on $\{d_\square>0\}$. (ii) There are bounded point operators $\mathsf b(w|e|w')$, $d_\square(w),d_\square(w')>0$, $e\in V_{\mathbb C}$, holomorphic in norm and linear in $e$, with $\|\mathsf b(w|e|w')\|\le\sqrt2C_Vs^{-p}|e|$ for $0<s\le\min\{1,d_\square(w),d_\square(w')\}$, $\mathcal T(a)\mathsf b(w|e|w')\mathcal T(a')=\mathsf b(w+a|e|w'+a')$, and $\mathsf B_{2\eta}(\tau_{\eta e_3}h)=\int\mathsf b(y+\eta e_3|h(y)|-y+\eta e_3)\,dy$. (iii) The functions
\[\mathbf W_n(\mathfrak z;e)=\bigl\langle\Omega,\mathsf b(w|e_1|\tfrac12\mathfrak z_1)\,\mathsf b(\tfrac12\mathfrak z_1|e_2|\tfrac12\mathfrak z_2)\cdots\mathsf b(\tfrac12\mathfrak z_{n-1}|e_n|w)\Omega\bigr\rangle,\]
which do not depend on $w$, are holomorphic on $\{d_\square>0\}^{n-1}$, $n$-linear in $e$, and $|\mathbf W_n(\mathfrak z;e)|\le(C_V\eta^{-p})^n\prod_k|e_k|$ for $e\in V^n$, $d_\square(\mathfrak z_k)\ge2\eta$, $0<\eta\le1$; $\sigma(X)=\int\mathbf W_n(x_2-x_1,\dots,x_n-x_{n-1};h_1(x_1),\dots,h_n(x_n))\,dx$ for tuples ordered in the third coordinate, and correspondingly for tuples ordered in another coordinate direction; so $|\sigma(X)|\le\prod_kC_V\eta^{-p}\|h_k\|_{L^1}$ for gaps $\ge2\eta$, with no factorial. (iv) For $f\in\mathcal S_V$ there is exactly one operator $T(f):\mathcal D_\infty\to\mathcal D_\infty$ with $\langle\psi,T(f)\psi'\rangle=\int q_0(U(x)^{-1}\psi,U(x)^{-1}\psi';f(x))\,dx$, where $q_0(\psi,\psi';e)=\langle(1+H)^l\psi,\mathsf A_l(e)(1+H)^l\psi'\rangle$ with $\mathsf A_l(e)$ bounded; $\langle\psi,T(f)\psi'\rangle=\langle T(f^*)\psi,\psi'\rangle$, $U(a)T(f)U(a)^{-1}=T(f(\cdot-a))$, $[H,T(f)]=iT(\partial_tf)$, $\|(1+H)^jT(f)\psi\|\le c_j(f)\|(1+H)^{2l+j}\psi\|$, and $T$ is covariant under the eight signed permutations that fix $e_3$; $T(f)$ is determined by $\sigma$ alone. (v) $\langle\Omega,T(f_1)\cdots T(f_n)\Omega\rangle$ are tempered and are the boundary values of the $\mathbf W_n$. (vi) For every $E\in V$ the form $q_0(\cdot,\cdot;E)$ on $\mathcal D_\infty\cap\mathcal H_0$ is bounded by $\|\mathsf A_l\|\,|E|\,\|(1+H)^l\psi'\|\,\|(1+H)^l\psi\|$, vanishes at $(\Omega,\Omega)$ and has $\overline{q_0(\psi',\psi;E)}=q_0(\psi,\psi';\rho(\theta)E)$; so the forms $\mathfrak q_{ij}=q_0(\cdot,\cdot;\mathsf m_{ij})$ at the label matrices of the stress tensor satisfy (M1) with $c_T=\sqrt3\,\|\mathsf A_l\|$, and the operator formula of (M2) holds at different heights for the kernels $\mathsf s_{ij}(x;y)=\mathsf s_{1+m}(x,y;\mathsf m_{ij},\mathbf 1,\dots,\mathbf 1)$, $\mathsf s_{1+m}$ the kernels of (g).
\item[(e)] \emph{(The circular cone.)} $\mathcal P(\{E<|q|\})=0$. The operators $\mathcal T(\mathfrak z)$ are contractions on the closure of $\{d_\circ>0\}$ and holomorphic inside, and $\mathbf W_n$ continues holomorphically to $\{d_\circ>0\}^{n-1}$, $n$-linear in $e$; for $\varepsilon>0$ and the $\mathfrak z$ with $d_\circ(\mathfrak z_k-2\varepsilon e_3)>0$ for all $k$ it is bounded by $\prod_kC_{\rm pt}(\varepsilon)|e_k|_\infty$, where
\[C_{\rm pt}(\varepsilon)=3\sqrt2\,C_V\min\{1,\varepsilon\}^{-p},\]
$C_V$ and $p$ as in (d). Indeed, at $w=\varepsilon e_3$ the relation $\mathcal T(a)\mathsf b(w|e|w')\mathcal T(a')=\mathsf b(w+a|e|w'+a')$ of (d)~(ii) and the semigroup law give $\mathbf W_n(\mathfrak z;e)=\langle\Omega,\mathsf b_1\mathcal T(\mathfrak z_1-2\varepsilon e_3)\mathsf b_2\cdots\mathcal T(\mathfrak z_{n-1}-2\varepsilon e_3)\mathsf b_n\Omega\rangle$, $\mathsf b_k=\mathsf b(\varepsilon e_3|e_k|\varepsilon e_3)$; the right member continues holomorphically to the $\mathfrak z$ with $d_\circ(\mathfrak z_k-2\varepsilon e_3)>0$, where the factors $\mathcal T(\cdot)$ are contractions, and $\|\mathsf b_k\|\le\sqrt2C_V\min\{1,\varepsilon\}^{-p}|e_k|\le3\sqrt2C_V\min\{1,\varepsilon\}^{-p}|e_k|_\infty$ by (d)~(ii) at $s=\min\{1,\varepsilon\}$ and $|e|\le3|e|_\infty$.

In each of the nine pictures $u$ of the coordinate and diagonal mirrors, with the joint spectral measure $\mathcal P_u$ formed as in Definition~\ref{C-355-4}~(3) with $u$ in the place of $e_3$, $\mathcal P_u(\{E<|q|\})=0$, and the operators $\mathcal T_u(\mathfrak z)$ formed with $\mathcal P_u$ as $\mathcal T(\mathfrak z)$ with $\mathcal P$ are contractions on the closure of $\{d_\circ>0\}$ and holomorphic inside; the continuation of $\mathbf W_n$, with its bound, holds in the pictures of $e_1$, $e_2$, $e_3$, the invariance under $W_3$ carrying the picture of $e_3$ to the two others. There are two strongly continuous unitary groups $V_1$, $V_2$ on $\mathcal H(\sigma)$ with $V_k(s)\Omega=\Omega$ and $V_k(s)\mathcal P(B)V_k(s)^{-1}=\mathcal P(\Lambda^{(k)}_{-\pi s/2}B)$ for every Borel set $B$, $\Lambda^{(k)}_\varsigma$ the boost of rapidity $\varsigma$ in the plane of $(E,q_k)$, which multiplies $E+q_k$ by $e^{\varsigma}$ and $E-q_k$ by $e^{-\varsigma}$ and fixes the other momentum; the two boosts generate the restricted Lorentz group $\mathrm{SO}_0(1,2)$ of the form $E^2-|q|^2$, so the null sets and the support of $\mathcal P$ are invariant under it.

This clause holds also for the tensor limits of Theorem~\ref{C-355-10}~(e) at every $0<g\le\gamma^{\rm ref}_W$.
\item[(f)] \emph{(Ordered classes.)} $\mathcal H^{\rm ord}(\sigma)=\mathcal H(\sigma)$; $\Omega$ and the vectors $\mathsf B_{l_1}(h_1)\cdots\mathsf B_{l_n}(h_n)\Omega$ are total, and the span $\mathcal D_0$ of $\Omega$ and of the $T(f_1)\cdots T(f_n)\Omega$ is dense in $\mathcal H(\sigma)$. Two tensor limits that agree on ordered tuples are equal; so $\sigma$ is determined by its functions $\mathbf W_n$.
\item[(g)] \emph{(Kernels.)} Let $\mathsf K_n=(10n)^{\lceil pn\rceil n(n-1)/2}$. For every admissible tuple $M=(h_1,\dots,h_n)$ of continuous compactly supported $V_{\mathbb C}$-valued entries with pairwise disjoint supports in a slab $\{0<x_3<l\}$, entries side by side allowed, there is a bounded operator $\mathsf B_l(M)$ with $[M\cdot\tau_{le_3}F]=\mathsf B_l(M)[F]$ and
\[\|\mathsf B_l(M)\|\le\mathsf K_n\bigl(\sqrt2C_V(3/d(M))^p\bigr)^n\prod_a\|h_a\|_{L^1},\]
$d(M)$ the least of $1$, of the distance of the supports from the two faces, and of the sup distances between two supports. For every $X\in\mathfrak S_V(\mathbb R^3)$, $\sigma(X)=\int\mathsf s_n(x;h_1(x_1),\dots,h_n(x_n))\,dx$ with a kernel $\mathsf s_n$ continuous at all non-coinciding points, $n$-linear in $e$, real-analytic in the coordinates of one direction at fixed values of the others, and
\[|\mathsf s_n(x;e)|\le\mathsf K_n\Bigl(\sqrt2C_V\bigl(3/\min\{1,\min_{a\ne b}|x_a-x_b|_\infty\}\bigr)^p\Bigr)^n\prod_a|e_a| ;\]
hence $|\sigma(X)|\le\mathsf K_n\bigl(\sqrt2C_V(3/\min\{1,d/\sqrt3\})^p\tfrac{4\pi}3R^3\bigr)^n\prod_a\sup|h_a|$ for tuples in balls of radius $R$ at distances $\ge d$, whatever their position. $\sigma$ does not charge the configurations with two equal heights. The scalar part $\sigma_0(f_1,\dots,f_n)=\sigma(f_1\mathbf 1,\dots,f_n\mathbf 1)$, at $C^1_c$ entries, is the limit of the centred products of the action densities along the data $(K_j,\beta_j,\rho_j)$ of $\sigma$, since $\mathcal T_K(f\mathbf 1)=\mathcal O(f)$; these data form a convergent Gibbs sequence at $g$ (Definition~\ref{C-365-5}(2)), so $\sigma_0$ is scalar data (Rec) with the numbers $16\,C^g_{\rm pt},4,\frac12$ (Theorem~\textup{\YMMref{AP}{780}}, Section~\YMMref{107}{765}). With it, the kernel of one tensor entry among $m$ scalar entries is real-analytic at all non-coinciding points, for every $m$ (Proposition~\ref{C-355-27}~(c)).
\item[(h)] \emph{(Growth.)} There is a number $c_4$ formed from the constants and $g$ such that for $n\ge1$, $0<d\le1$, $i=\max\{i_E(g,\frac1{8n}),\ \text{the least $i$ with }a_i\le d/4\}$, every first-passage datum with $K\ge K_E(i)$ (at $\delta=\frac1{8n}$), $K\ge i+1$, $8na_K\le a_i$, any two states $\mathfrak s,\mathfrak s'\in\mathcal G^{\rm per}_{K,\beta}$, and real $H_a$ vanishing off closed sets in balls of radius $1$ at Euclidean distances $\ge d$,
\[\Bigl|\mathfrak s\Bigl(\prod_a\bigl(\mathcal T_K(H_a)-\mathfrak s'(\mathcal T_K(H_a))\bigr)\Bigr)\Bigr|\le\bigl(c_4n^{35/2}d^{-19/2}b^{-1}\bigr)^n\prod_a|H_a|_\infty ;\]
so $|\sigma(X)|\le(c_4g^2e^{35/2})^n(n!)^{35/2}d^{-(19/2)n}\prod_a|H_a|_\infty$ for every tensor limit and every component of Theorem~\ref{C-355-13}. With the upper Riemann sums $\|H\|_\natural=\sum_\square|\square|\sup_\square|H|_\infty$ of $|H|_\infty$ over the boxes $\square=\frac{a_i}{4n}(\mathbf k+[0,1)^3)$, $\mathbf k\in\mathbb Z^3$, of the mesh $a_i/(4n)\le d/(16n)$, $i$ the index of the first display, in the place of the sup norms, $|\sigma(X)|\le(Cn^{16}d^{-9})^n\prod_a\|H_a\|_\natural$ for entries of any size at pairwise sup distances $\ge2d$, ordered or not, $C$ formed from the constants and $g$. For every $n$ there is exactly one linear functional $\mathsf S_n$ on the $V_{\mathbb C}^{\otimes n}$-valued Schwartz functions on $\mathbb R^{3n}$ that vanish with all their derivatives on the diagonals, bounded for one Schwartz norm, with $\mathsf S_n(H_1\otimes\dots\otimes H_n)=\sigma(H_1,\dots,H_n)$ at smooth tuples; it satisfies
\[|\mathsf S_n(\varphi)|\le(C^V_E)^n(n!)^{51/2}|\varphi|_{18n},\]
the growth hypothesis in Proposition~\YMMref{93.5}{719} with a label, with $C^V_E$ formed from the constants and $g$, the same for all tensor limits at $g$ and all their components; the pairs $(17n,\frac{49}2)$ and $(15n,\frac{43}2)$ of the index of the norm and the power of the factorial serve as well, with other numbers; on compact subsets of the non-coinciding configurations it is a distribution of order at most $4n+1$.
\end{enumerate}
The proofs of (a)--(h) use no covariance of $\sigma$ under rotations outside $W_3$ at tuples with nonscalar entries; neither Theorem~\ref{C-355-10}~(a) nor Proposition~\ref{C-355-28} is used. The last two sentences of (g) use only the previously established Euclidean invariance of the scalar limit, through its scalar reconstruction data (Rec); see Theorem~\textup{\YMMref{AP}{780}}.
\end{theorem}

\begin{theorem}[Covariance, reconstruction, locality]\label{C-355-10}
Let $0<g\le\gamma^{\rm sd}_W$ and let $\sigma$ be a tensor limit at $g$ with label $b_*$.
\begin{enumerate}
\item[(a)] $\sigma(\mathsf RX)=\sigma(X)$ for every $\mathsf R\in\mathrm O(3)$ and every $X\in\mathfrak S_V(\mathbb R^3)$.
\item[(b)] $\sigma(\theta_{n,c}\bar F\cdot F)\ge0$ for every unit vector $n$, every $c\in\mathbb R$ and every $F\in\mathbb C\mathfrak S_V(\{n\cdot x>c\})$, where $(\theta_{n,c}H)(x)=\Theta_nH(x-2(n\cdot x-c)n)\Theta_n$ and $\Theta_n=1-2nn^{\rm T}$.
\item[(c)] $\sigma^1$ is reconstruction data with labels at $V=\mathrm{Sym}_3(\mathbb R)$ with the invariant norm, with $c_0=\frac12$, every $p_V>\frac{17}2$ and $C_V=2^{p_V}C_{p_V}b_*^{-1}\le2^{p_V}C_{p_V}g^2$; its invariance holds at all tuples of continuous entries. The map $[F]_{\sigma^1}\mapsto[F]_\sigma$ extends to a unitary map of $\mathcal H(\sigma^1)$ onto $\mathcal H(\sigma)$ that carries $\Omega$, $H$, $\mathcal U$, $\mathcal P$ to $\Omega$, $H$, $\mathcal U$, $\mathcal P$.
\item[(d)] (i) For every $n$ there is exactly one continuous function $\mathsf s_n(x;e)$ on the non-coinciding configurations, $n$-linear in $e\in V^n$, with $\sigma(H_1,\dots,H_n)=\int\mathsf s_n(x;H_1(x_1),\dots,H_n(x_n))\,dx$ for all $X\in\mathfrak S_V(\mathbb R^3)$; it is real-analytic, symmetric when points and labels are permuted together, $\mathsf s_n(\mathsf rx_1,\dots,\mathsf rx_n;\rho(\mathsf R)e_1,\dots,\rho(\mathsf R)e_n)=\mathsf s_n(x;e)$ for every isometry $\mathsf rx=\mathsf Rx+v$, and $|\mathsf s_n(x;e)|\le\bigl(C_V(2n^2/\min\{1,\mathsf d(x)\})^{p_V}\bigr)^n\prod_a|e_a|$ at every point, $\mathsf d(x)=\min_{a\ne b}|x_a-x_b|$; so $|\sigma(X)|\le(C_V(2n^2/d)^{p_V})^n\prod_a\|H_a\|_{L^1}$ at pairwise distances $\ge d$, $0<d\le1$, and $|\mathsf S_n(\varphi)|\le(8\pi^23^{\lceil p_V\rceil+2}(2e^2)^{p_V}C_V)^n(n!)^{2p_V}|\varphi|_{(\lceil p_V\rceil+4)n}$. (ii) There is a strongly continuous unitary representation $U$ of $\mathbb R^{1+2}\rtimes\mathcal L$ on $\mathcal H(\sigma)$ with $U(x,1)=e^{itH}\mathcal U(\vec x)$, $U(a,\Lambda)\Omega=\Omega$, $U(a,\Lambda)\mathcal D_0=\mathcal D_0$, the joint spectrum in the closed forward cone, and
\[U(a,\Lambda)T(f)U(a,\Lambda)^{-1}=T\bigl(\hat\Lambda f(\Lambda^{-1}(\cdot-a))\hat\Lambda^{\rm T}\bigr)\quad\text{on }\mathcal D_0 ,\]
with the transpose, not the adjoint. (iii) $T(f)T(f')\psi=T(f')T(f)\psi$ for $\psi\in\mathcal D_0$ whenever $f,f'\in\mathcal S_V$ have spacelike separated supports: all six components are local with respect to all six.
\item[(e)] Let $0<g\le\gamma$. Every sequence of first-passage data at $g$ with $K_j\to\infty$ and periodic Gibbs states has a subsequence along which the tensor limit exists. Every such limit $\sigma$ satisfies (a), (b) and the conclusion of Theorem~\ref{C-355-11}~(c), and $\sigma^1$ is reconstruction data with labels with $c_0=\frac12$, every $p_V>\frac{17}2$ and
\[C_V=3^{p_V}L^{4z}\max_{k=1,2}2^{p_V}C_{p_V}[g'_k]\,g'^2_k ,\]
where $z$ is the least integer $\ge1$ with $L^zg^{-2}\ge2(\gamma^{\rm sd}_W)^{-2}$, $g'^{-2}_1=L^zg^{-2}/2$, $g'^{-2}_2=L^{z+1}g^{-2}/4$, and $C_p[g']$ is the number $C_p$ of Theorem~\ref{C-355-9}~(c) at the coupling $g'$: a number formed from the constants, $g$ and $p_V$, one triple for all the tensor limits at $g$; so (c) and (d) hold for it.
\end{enumerate}
\end{theorem}

\begin{theorem}[Conservation]\label{C-355-11}
For $\eta\in C^2_c(\mathbb R^3;\mathbb R^3)$ put $H^{\rm cons}[\eta]=(\nabla\cdot\eta)\mathbf 1-\nabla\eta-(\nabla\eta)^{\rm T}$ and $\|\eta\|_{C^2}=\sup|\eta|+\sup|\nabla\eta|+\sup|\nabla^2\eta|$; with the label matrices $\mathsf m_{ij}=\delta_{ij}\mathbf 1-e_ie_j^{\rm T}-e_je_i^{\rm T}$ of the stress tensor one has $\sum_i(\partial_i\zeta)\mathsf m_{ij}=H^{\rm cons}[\zeta e_j]$, and $\frac12\operatorname{tr}(\mathsf m_{ij}T)=\frac12\delta_{ij}\operatorname{tr}T-T_{ij}$ is the stress tensor. Let $0<g\le\gamma$.
\begin{enumerate}
\item[(a)] \emph{(At a cutoff, in every periodic Gibbs state.)} Let $n\ge2$, $1\le n_1\le n$, $R>0$, $0<c<\frac1{28}$, $\phi=(\frac1{28}-c)/(4n-1)$ and $u=\lceil\phi K\rceil$. There are numbers $C$, $p$ (from the constants, $g$, $n$, $R$, $c$) and $K_0$ (from the constants, $g$, $n$, $c$) such that for all first-passage data $(K,\beta)$ at $g$ with $K\ge K_0$, all $\nu\in\mathcal G^{\rm per}_{K,\beta}$, every $d>0$ with $dL^u\ge K^2$, all $\eta_1,\dots,\eta_{n_1}\in C^2_c(\mathbb R^3;\mathbb R^3)$ and all $H_{n_1+1},\dots,H_n\in C_c(\mathbb R^3;V)$ whose $n$ supports lie in a ball of radius $R$ and are pairwise at sup distance $\ge d$, with $X_a=\mathcal T_K(H^{\rm cons}[\eta_a])$ for $a\le n_1$ and $X_a=\mathcal T_K(H_a)$ for $a>n_1$,
\[\Bigl|\nu\Bigl(\prod_{a=1}^n\bigl[X_a-\nu(X_a)\bigr]\Bigr)\Bigr|\le C\,K^pL^{-cn_1K}\prod_{a\le n_1}\|\eta_a\|_{C^2}\prod_{a>n_1}|H_a|_\infty ,\]
and the same bound holds for the joint cumulant of $X_1,\dots,X_n$ under $\nu$. At $c=\frac1{56}$, $n_1=1$ the rate is $K^pL^{-K/56}$.
\item[(b)] \emph{(A second form; tori.)} Write $A^\nu[F]=F-\nu(F)$. For $q\ge2$, $d>0$, $R>0$ there are $C$, $p$, $K_0$, formed from the constants, $g$, $q$, $d$, $R$, such that for all first-passage data $(K,\beta)$ at $g$ with $K\ge K_0$, all $\nu\in\mathcal G^{\rm per}_{K,\beta}$, every $\eta\in C^2_c(\mathbb R^3;\mathbb R^3)$ and all bounded $H_2,\dots,H_q:\mathbb R^3\to V$ whose supports, together with $\operatorname{supp}\eta$, lie in a cube of side $2R$ at pairwise sup distances $\ge d$,
\[\Bigl|\nu\Bigl(A^\nu\bigl[\mathcal T_K(H^{\rm cons}[\eta])\bigr]\prod_{a=2}^qA^\nu\bigl[\mathcal T_K(H_a)\bigr]\Bigr)\Bigr|\le C\,K^pL^{-K/36}\,\|\eta\|_{C^2}\prod_{a=2}^q|H_a|_\infty ,\]
also for the joint cumulant. For the same first-passage data, on every torus $T_\ell$, $\ell\in\mathfrak L(g)$, the corresponding bound holds for $C^2$ vector fields $\eta$ on $T_\ell$ and bounded $H_2,\dots,H_q$ whose supports, together with $\operatorname{supp}\eta$, are pairwise at sup distance $\ge d$ on $T_\ell$:
\[
|\langle\mathcal T_K(H^{\rm cons}[\eta]);\mathcal T_K(H_2);\dots;\mathcal T_K(H_q)\rangle_{K,\ell,\beta}|\le CK^pL^{-K/36}\|\eta\|_{C^2}\prod_{a=2}^q|H_a|_\infty,
\]
with $C$, $p$, $K_0$ depending besides on $\ell$.
\item[(c)] \emph{(In the limit.)} For every tensor limit $\sigma$ at $g$ (in the sense of Theorem~\ref{C-355-10}~(e) if $g>\gamma^{\rm sd}_W$), every $n\ge2$, every $\eta\in C^2_c(\mathbb R^3;\mathbb R^3)$ and all continuous $H_2,\dots,H_n$ with compact, pairwise disjoint supports off $\operatorname{supp}\eta$,
\begin{equation}\label{C-355-12}
\begin{multlined}
\sigma(H^{\rm cons}[\eta],H_2,\dots,H_n)=0,\qquad\text{equivalently}\qquad\sum_i\sigma\bigl((\partial_i\zeta)\mathsf m_{ij},H_2,\dots,H_n\bigr)=0 \\
(\zeta\in C^2_c).
\end{multlined}
\end{equation}
On every torus $T_\ell$, $\ell\in\mathfrak L(g)$, the limits of the cumulants of Theorem~\ref{C-355-15}~(E3) vanish at $(H^{\rm cons}[\eta],H_2,\dots,H_q)$ for Lipschitz $H_a$.
\item[(d)] \emph{(Kernels and operators.)} Let $\sigma'$ be scalar reconstruction data with the frame (M1), (M2) and a unique vacuum. Then the identity \eqref{C-355-12} for the kernels $\mathsf s_{ij}$ of (M2) among scalar entries, $\sum_i\int\partial_i\zeta(x)\mathsf s_{ij}(x;y)\,dx=0$, and (M4)~(a) are equivalent.
\end{enumerate}
\end{theorem}

\begin{theorem}[Components and Wightman tensor fields]\label{C-355-13}
Let $\sigma$ be the tensor limit of a sequence $(K_j,\beta_j,\nu_j)$ of first-passage data and periodic Gibbs states at $g$, and $A^\nu_K(X)=\prod_a[\mathcal T_K(H_a)-\nu(\mathcal T_K(H_a))]$.
\begin{enumerate}
\item[(a)] \emph{(Decomposition; $0<g\le\gamma^{\rm ref}_W$.)} There are a compact metrisable space $\mathsf X$, a Borel probability measure $\lambda$ on $\mathsf X$ with support $\mathsf X$, and real symmetric functionals $\sigma_\xi$, $\xi\in\mathsf X$, on labelled tuples, with $\xi\mapsto\sigma_\xi(X)$ continuous and $\xi\mapsto\sigma_\xi$ injective, such that $\sigma$ has iterated plateaux in the sense of Definition~\ref{C-355-4}~(5) at every $(X^1,\dots,X^q)$, $q\ge1$, $X^i\in\mathfrak S_V(\mathbb R^3)$, and its iterated plateau there is $\int\prod_i\sigma_\xi(X^i)\,d\lambda(\xi)$.
\item[(b)] \emph{(Each component; $0<g\le\gamma^{\rm ref}_W$.)} For every $\xi\in\mathsf X$: (i) there are translation-invariant events $E_j$ with $\nu_j(E_j)>0$ such that $\rho_j=\nu_j(\cdot\,|\,E_j)\in\mathcal G^{\rm per}_{K_j,\beta_j}$ and $\sigma_\xi(X)=\lim_j\rho_j(A^{\nu_j}_{K_j}(X))$ along the full sequence, indeed $\|E_{\rho_j}[A^{\nu_j}_{K_j}(X)\,|\,\mathcal I]-\sigma_\xi(X)\|_{L^\infty(\rho_j)}\to0$, $\mathcal I$ the invariant $\sigma$-field; (ii) $\sigma_\xi((H))=\frac{c_{\pi(\xi)}}3\int\operatorname{tr}H$ with a real number $c_{\pi(\xi)}$, and the recentred functional $\sigma^c_\xi(X)=\sum_{J\subset\{1..n\}}\prod_{a\notin J}(-\sigma_\xi((H_a)))\,\sigma_\xi(X_J)$ is the tensor limit of $(K_j,\beta_j,\rho_j)$; so it has all the properties of a tensor limit, and $[F]\mapsto[\mathsf S_\xi F]$, $\mathsf S_\xi X=\sum_J\prod_{a\notin J}\sigma_\xi((H_a))X_J$, is a unitary map of $\mathcal H(\sigma_\xi)$ onto $\mathcal H(\sigma^c_\xi)$ that carries $\Omega_\xi$ to $\Omega$ and intertwines $H$ and $\mathcal P$; (iii) $\ker H_\xi=\mathbb C\Omega_\xi$; (iv) there is exactly one map $\pi$ of $\mathsf X$ onto the space $\mathsf X_0$ of the canonical decomposition $(\mathsf X_0,\lambda_0,(\sigma_{0,\zeta}))$ of the scalar part $\sigma|_{V_0}$ with $\sigma_\xi(f_1\mathbf 1,\dots,f_n\mathbf 1)=\sigma_{0,\pi(\xi)}(f_1,\dots,f_n)$; it is continuous and onto, and $\pi_*\lambda=\lambda_0$.
\item[(c)] \emph{(Wightman fields; $0<g\le\gamma^{\rm sd}_W$.)} For every $\xi$, the restriction of $\sigma^c_\xi$ to $C^1_c$ entries is reconstruction data with labels with $c_0=\frac12$, every $p_V>\frac{17}2$ and the numbers of Theorem~\ref{C-355-10}~(c), and it has (Rec5). Its reconstruction is a Wightman tensor field of six components in $1+2$ dimensions: the fields $T(f)$, $f\in\mathcal S_V$, on $\mathcal D_0$ are tempered; $\langle\psi,T(f)\psi'\rangle=\langle T(f^*)\psi,\psi'\rangle$, so that $T_{33}$ and $T_{jk}$, $j,k\le2$, are Hermitian and $T_{3j}$ is Hermitian after the factor $i$; $U(a,\Lambda)T(f)U(a,\Lambda)^{-1}=T(\hat\Lambda f(\Lambda^{-1}(\cdot-a))\hat\Lambda^{\rm T})$; the joint spectrum lies in the closed forward cone; $\Omega_\xi$ is the only invariant vector and is cyclic; all components are local with respect to all. The growth condition of Theorem~\ref{C-355-9}~(h) holds with the same numbers for all components. The given limit $\sigma$ is the integral over $\lambda$ of its components in the sense of (a).
\item[(d)] \emph{(The frame over a component; $0<g\le\gamma^{\rm sd}_W$.)} On the scalar subspace of $\mathcal H(\sigma^c_\xi)$ the frame (M1)--(M5) holds for the tensor formed at the label matrices $\mathsf m_{ij}$ of the stress tensor, with the forms $q_0(\cdot,\cdot;\mathsf m_{ij})$, the kernels $\mathsf s_{ij}(x;y)=\mathsf s_{1+m}(x,y;\mathsf m_{ij},\mathbf 1,\dots,\mathbf 1)$, the scalar field $\phi(f)=T(f\mathbf 1)$ and the restriction to the scalar subspace of the representation $U$ of Theorem~\ref{C-355-10}~(d)~(ii) for $\sigma^c_\xi$. The scalar subspace reduces $U(a,\Lambda)$: $\phi(f)$ keeps it; $U(a,\Lambda)\phi(f)U(a,\Lambda)^{-1}=\phi(f(\Lambda^{-1}(\cdot-a)))$ on $\mathcal D_0$, since $\hat\Lambda\mathbf 1\hat\Lambda^{\rm T}=\mathbf 1$; $U(a,\Lambda)\Omega=\Omega$; and $\Omega$ with the vectors $\phi(f_1)\cdots\phi(f_n)\Omega$ spans a dense subspace of it, by Proposition~\ref{C-355-25}~(a)~(v) for the restriction of $\sigma^c_\xi$ to the tuples with entries in $V_0$, whose space is the scalar subspace and whose fields are the restrictions of the $T(f\mathbf 1)$. In particular the real-time stress tensor is conserved, (M4)~(a).
\item[(e)] \emph{(The frame on $\mathcal H_0$; $0<g\le\gamma^{\rm sd}_W$.)} For $\sigma$ itself, (M1) holds, the kernels $\mathsf s_{ij}(x;y)$ of one tensor entry among scalar entries are covariant as in (M2) and continuous at all non-coinciding points, $[T^E_{ij}(F),\phi(f)]=0$ between scalar vectors at spacelike separation as in (M3), and $\langle\Omega,\phi(f_1)\cdots\phi(f_n)T[t](F)\Omega\rangle$ is invariant under $(f,t,F)\mapsto(f\circ\Lambda^{-1},\Lambda t,F\circ\Lambda^{-1})$ as in (M5).
\end{enumerate}
\end{theorem}

\begin{proposition}[Covariance at one label suffices]\label{C-355-14}
Let $0<g\le\gamma^{\rm sd}_W$ and let $\sigma$ be a tensor limit at $g$. If $\sigma(\mathsf RX)=\sigma(X)$ for all $\mathsf R\in\mathrm O(3)$ at the tuples with exactly one entry outside $V_0$, then the kernels of (M2) are covariant at all non-coinciding points and (M3), (M5) hold on $\mathcal H_0$.
\end{proposition}

By Theorem~\ref{C-355-10}~(a) every tensor limit at $g$ satisfies the condition of Proposition~\ref{C-355-14}, and its conclusion is contained in Theorem~\ref{C-355-13}~(e). Its proof uses rotational covariance only for tuples with exactly one entry outside $V_0$. It follows the proofs of Propositions~\YMMref{108.3}{771}, \YMMref{110.2}{777} and~\YMMref{110.1}{776}, with one labelled insertion.

\begin{theorem}[Cumulants on a torus; laws at short distance]\label{C-355-15}
Assume the setting of Definition~\ref{C-355-6}; let $k'_n(g,s)$ and $C_n(g)$, $C(g)$ be the numbers of Proposition~\ref{C-355-20}.
\begin{enumerate}
\item[(E1)] Let $\ell$ be small over the running couplings of $(K,\beta)$, $0<d<1$ with $s(d)\le K$, and (N.$\theta$) at $(n,g,s(d))$. For all $H_a:T_\ell\to V_{\mathbb C}$ with (Geo) at $d$,
\[
\begin{multlined}
|\langle\mathcal T_K(H_1);\dots;\mathcal T_K(H_n)\rangle_{K,\ell,\beta}|\le60^nk'_n(g,s(d))\,b^{-n}\prod_a|H_a|_\infty\le60^nC_n(g)\bigl[2\log\tfrac1d \\
+C(g)\bigr]^{8n\mathfrak p_0}b^{-n}\prod_a|H_a|_\infty .
\end{multlined}
\]
No upper bound on the distances and no smoothness is required. On sides with (Sd.1), in particular on every $\ell\in\mathfrak L(g)$, the bound holds with the factor $\max\{1,\mathsf c(\ell)\}^n$.
\item[(E2)] Let $s_{\rm th}(n)$ be the least integer $\ge3$ from which on (N.$\theta$) holds at $(n,g,\cdot)$, $d''=\min\{\frac34d,L^{-s_{\rm th}(n)}\}$, $s''=s(d'')$, cut $T_\ell$ into boxes $\square$ of sides $\le0.99\,d''$, and let the sets $\Sigma_a$ be pairwise at distance $\ge d$. (i) If $\ell$ is small and $K\ge s''$: $|\langle\mathcal T_K(H_1);\dots;\mathcal T_K(H_n)\rangle|\le60^nk'_n(g,s'')\,b^{-n}\prod_a\sum_\square\sup_\square|H_a|\le60^nk'_n(g,s'')(2\ell/d'')^{3n}b^{-n}\prod_a|H_a|_\infty$. (ii) If $K<s''$, on any side and with or without the separation: the same cumulant is at most $f_n[18\bar x_0(C_W(s'')+1)^4/d'']^nb^{-n}\prod_a\sum_\square\sup_\square|H_a|$, $f_n=\sum_\pi(|\pi|-1)!$ over the partitions $\pi$ of $\{1..n\}$. (iii) On small sides, at every cutoff, it is at most $k_nb^{-n}\prod_a|H_a|_\infty$ with $k_n\le C(n,g,\ell)\,d^{-3n}[2\log\frac1d+C']^{8n\mathfrak p_0}$.
\item[(E3)] (i) \emph{(Scaling.)} Let $g_\star\le\gamma^{\rm ref}_W$, $\beta>X_{g_\star}$, $K_\star=K(\beta,g_\star)$, $z=K-K_\star$, and let $L^z\ell$ be small over the running couplings of $(K_\star,\beta)$. Then $\delta_\star b\le L^{-z}<g_\star^2b$ with $\delta_\star=[L(g_\star^{-2}+\bar B)]^{-1}$, and for $0<d<L^{-z}$ with $s(L^zd)\le K_\star$, (N.$\theta$) and (Geo) at $d$: $|\langle\mathcal T_K(H_1);\dots;\mathcal T_K(H_n)\rangle_{K,\ell,\beta}|\le60^nC_n(g_\star)[2\log\frac bd+C(g_\star)]^{8n\mathfrak p_0}b^{-n}\prod_a|H_a|_\infty$; at $g_\star=\gamma^{\rm sd}_W$ the numbers depend on $n$ alone. (ii) \emph{(The limit of the cutoff on one torus.)} Let $0<g\le\gamma^{\rm ref}_W$, $b\in(g^{-2}+\bar B,X_g-\bar B]$, $\ell\in\mathfrak L(g)$, and let $H_a$ be bounded Borel functions vanishing off closed sets at positive distances, entering through their means over the lattice cubes.

The cumulants $\kappa_K$ at $(K,\ell,L^Kb)$ converge along the full sequence $K\to\infty$; the limit is $n$-linear, symmetric, and translation invariant at uniformly continuous entries; at Lipschitz entries it equals the limit with point values. The rates are those of Proposition~\ref{C-355-36}: with the least sup distance $\mathsf d$ of the sets there are numbers $C$, $p$, formed from the constants, $n$ and $b$, and a cutoff $K_1$, formed besides from $g$, $\ell$ and $\mathsf d$, such that for $K\ge K_1$
\[|\kappa_{K+1}-\kappa_K|\le C\ell^{3(n+1)}K^pL^{-\theta_nK}\prod_a|H_a|_\infty,\qquad\theta_n=\frac1{1260\,n(8n+3)},\]
for the entries through their means over the lattice cubes, and the same with $|H_a|_\infty+{\rm Lip}_\infty H_a$ in the place of $|H_a|_\infty$ for Lipschitz entries read at the points; for one bond form among action densities, with real Lipschitz weights, none of them $0$, read at the points and supports with no common point, whose spread (the definition preceding Lemma~\YMMref{114.6}{805}) is at least $\mathsf d_0>0$, the rate is $K^pL^{-K/120}$ from a cutoff on, with $C$, $p$ and the cutoff formed from the constants, $n$, $\ell$, $b$ and $\mathsf d_0$. If the sets lie in balls of radius $0.86\,d$ at distances $\ge1.27\,d$ and (N.$\theta$) holds at $(n,g,s(d))$, the limit is at most $60^n\max\{1,\mathsf c(\ell)\}^nk'_n(g,s(d))b^{-n}\prod_a|H_a|_\infty$.
\item[(E4)] \emph{(Finer forms.)} (i) Let $g\le\gamma^{\rm sd}_W$, $0<d<1$, $K\ge s(d)$, and let $2\mathsf l$ be the least side at $g$. For two insertions of any components with real weights whose supports lie in balls of radius $d-3a_K$ whose centres are at sup distance $\ge3d+a_K$ (on a torus, modulo its side), in every periodic Gibbs state and on every torus of $\mathfrak L(g)$:
\[|\langle\mathcal T_K(H_1);\mathcal T_K(H_2)\rangle|\le180^2C^\diamond(g)\,b^{-2}|H_1|_\infty|H_2|_\infty,\qquad C^\diamond(g)=\max\bigl\{2\,\mathsf c_\diamond(g)^2,\ C^T_2(\tfrac12;1)\bigr\},\]
with no power of $d$ and no logarithm, and with the factor $4$ for complex weights. Here $C^T_2$ is the number of (E5)~(i), and $\mathsf c_\diamond(g)=e^{G_g\mathsf l^3}\mathsf c'_\diamond(g)$, where $\mathsf c'_\diamond(g)$, formed from the constants and $g$, is the larger of a bound of the eighth root of the mean of the eight mirror copies of a centred form of one direction near a point on the least torus $T_{2\mathsf l}$ and of the crude bound of such a form at the cutoffs below a threshold formed from $g$; the factor $e^{G_g\mathsf l^3}$ is that of the transfer from the least torus to every side and every state (Lemma~\ref{C-355-17}). So $C^\diamond(g)$ is a number of the constants and $g$.

On one torus, from a first depth $s_1(g,\ell)$ on, the same bound holds for two cubes of side $\frac32h$ at sup distance $\ge\frac52h$, with a number depending on $\ell$ in the place of $180^2C^\diamond(g)$. (ii) At $n=3$, on one torus, for supports in three cubes of a common side $\delta\le\delta_\star(g,\ell)$ at pairwise sup distances $\ge2\delta$ and $s(\delta/4)\le K$, the cumulant is at most $64^3\,\Theta(\delta/4)\prod_a\sup|H_a|$ with $\Theta(\delta/4)\le C(g,\ell)[2\log\frac4\delta+C(g)]^{16\mathfrak p_0}b^{-3}$: the power $16\mathfrak p_0$ of the logarithm in the place of $24\mathfrak p_0$. (iii) Let $p\ge2$, $\ell=2L^m\in\mathfrak L(g)$, $0<d<1$, and let $w_1,\dots,w_p$ be real weights of the plaquettes whose supports have their barycentres in sup balls of radius $d$, the balls being pairwise at sup distance $\ge d$ (complex weights cost the factor $2^p$).

Let $y_K\ge\mathsf Y^\natural$, where $\mathsf Y^\natural$ is the threshold of Definition~\ref{def:rarity-choices}, under which Proposition~\ref{prop:local-large-field-rarity} bounds the large-field regions uniformly in the volume with its constant $c^\natural_2$: a ceiling on the running coupling at the unit scale, formed from the constants. Let $T$ be the largest level with $\hat d:=dL^{K-T}\ge M(2+\frac p2y_T)+1$, so that two supports are at least $M(2+\frac p2y_T)$ apart in the units of the level $T$; let $T\ge2$, which holds from a cutoff on that is formed from $d$, $b$ and $p$; and let $d\le b\,e^{-y^\sharp}$, where $y^\sharp\ge1$, formed from the constants, $p$ and $\ell/b$, is the least number such that for all $y\ge y^\sharp$: $e^y\ge x_{\rm J}$, the threshold of the definitions preceding \eqref{eq:weighted-history-threshold}; $\theta(e^y)$ is at most a number formed from the constants and $p$; $2\mathfrak p_0c^\natural_2A_0^2y^{2\mathfrak p_0-1}\ge4+(3+p)p$; $4\cdot24^3(\frac54)^3(\ell/b)^3\exp\bigl((3+p^2)y-c^\natural_2A_0^2y^{2\mathfrak p_0}\bigr)\le1$; and $c_{\rm A}\alpha_0e^{y/2}\ge2MB_3A_0y^{\mathfrak p_0}(\frac y2+1)$, with the number $c_{\rm A}$ of equation~\YMMref{(2262)}{828} and the constant $B_3$.

These conditions give $d_*(T)\ge\frac12y_T$ for the distance 
\[
d_*(T)=\min\{c_{\rm A}\alpha_0/(2MB_3\varepsilon(x_T)),(L^{m+K-T}-4)/(2M)\}-1
\]
 of equation~\YMMref{(2262)}{828} at the level $T$, $\varepsilon(t)=A_0t^{-1/2}(\log t)^{\mathfrak p_0}$. Then
\[
\begin{multlined}
|\langle\mathcal O(w_1);\dots;\mathcal O(w_p)\rangle_{K,\ell,\beta}|\le\mathsf C^{(p)}\prod_{a=1}^p\frac{y_T^{2\mathfrak p_0}\hat d^{\,3}|w_a|_\infty}b\le\mathsf C'^{(p)}\prod_{a=1}^p\frac{y_T^{2\mathfrak p_0+3}|w_a|_\infty}b,\qquad y_T \\
\le\log\tfrac bd+\log\bigl(\tfrac43L(M(2+\tfrac p2y_T)+1)\bigr),
\end{multlined}
\]
with numbers $\mathsf C^{(p)}$, $\mathsf C'^{(p)}$ formed from the constants and $p$: no power of $d$, and for each entry the power $2\mathfrak p_0+3$ of the logarithm of $b/d$ in the place of the power $8\mathfrak p_0$ of (E1).
\item[(E5)] \emph{(Bounds independent of the side and of the state; $0<g\le\gamma^{\rm sd}_W$.)} (i) For $n\ge2$, $d,R>0$ there is a number $C^T_n(d;R)$, formed from the constants, $g$, $n$, $d$, $R$, such that in every periodic Gibbs state and on every torus of $\mathfrak L(g)$, for supports in balls of radius $R$ pairwise at distance $\ge d$, at every cutoff: $|\langle\mathcal T_K(H_1);\dots;\mathcal T_K(H_n)\rangle|\le60^nC^T_n(d;R)b^{-n}\prod_a|H_a|_\infty$. With the numbers of Theorem~\ref{C-355-9}~(a), each formed at its own order: for $K\ge\max_{2\le m\le n}K^G(m)$, $K^G(m)$ the cutoff of that clause at the order $m$, the cumulant in a periodic Gibbs state is at most $\sum_\pi(|\pi|-1)!\prod_{B\in\pi}C^G_{\#B}(d;R)\,b^{-n}\prod_a|H_a|_\infty$, the sum over the partitions $\pi$ of $\{1,\dots,n\}$ without singletons; on a torus of $\mathfrak L(g)$ the same holds, by the chessboard estimate in the place of the domination, with the distances taken on the torus. (ii) Limits through arbitrary periodic Gibbs states, and along growing tori, exist along subsequences and obey these bounds; they are tensors over the scalar limit, invariant under $\mathbb R^3\rtimes W_3$ and reflection positive; on every torus of $\mathfrak L(g)$ the full sequence of cutoffs converges, with bounds independent of the side. (iii) \emph{(Folding.)} Let $n\ge2$, $\ell'=2\mathsf l$ the least side at $g$, $0<d<1$ with $1.74\,d\le\mathsf l/(4n)$ and (N.$\theta$) at $(8n,g,s(d))$.

For $r\ge2$ let $\mathfrak m_r(g,s)=\sum_{\pi'}\prod_{I\in\pi'}k'_{\#I}(g,s)$, the sum over the partitions $\pi'$ of a set of $r$ members into blocks of at least two members, and
\[k^{\rm fold}_n(g,s)=\max_{\mathcal C}\ \sum_\pi(|\pi|-1)!\prod_{B\in\pi}\ \prod_{C\in\mathcal C,\ B\cap C\ne\varnothing}e^{\mathsf l^3G_g}\,\mathfrak m_{8\#(B\cap C)}(g,s)^{1/8},\]
the sum over the partitions $\pi$ of $\{1,\dots,n\}$, the maximum over the partitions $\mathcal C$ of $\{1,\dots,n\}$. Then for every first-passage datum $(K,\beta)$ at $g$ with $s(d)\le K$, in every periodic Gibbs state and on every torus of $\mathfrak L(g)$, and for all $H_a$ with (Geo) at $d$,
\[
\begin{multlined}
|\langle\mathcal T_K(H_1);\dots;\mathcal T_K(H_n)\rangle|\le60^nk^{\rm fold}_n(g,s(d))\,b^{-n}\prod_a|H_a|_\infty,\qquad k^{\rm fold}_n(g,s(d)) \\
\le C'_n(g)\bigl[2\log\tfrac1d+C(g)\bigr]^{8n\mathfrak p_0},
\end{multlined}
\]
$C'_n(g)$ a number formed from the constants, $n$ and $g$: a power of the logarithm alone, free of the side and of the state. The cumulants of every tensor limit at $g$, formed from $\sigma$ by the recursion over the partitions, are at most $60^nk^{\rm fold}_n(g,s(d))\,b_*^{-n}\prod_a|H_a|_\infty$ at the tuples with (Geo) at $d$.
\item[(E6)] \emph{(The free two-point law; $0<g\le\gamma^{\rm sd}_W$.)} Put $K_3(u)_{\lambda\lambda'}=(3\hat u_\lambda\hat u_{\lambda'}-\delta_{\lambda\lambda'})/(4\pi|u|^3)$ and
\[\mathsf W(H_1,H_2)=\tfrac32\int\!\!\int\operatorname{tr}\bigl[H_1(x)K_3(x-y)H_2(y)K_3(x-y)\bigr]\,dx\,dy .\]
There are $d_0(g)>0$ and a function $e_g(d)$, formed from the constants and $g$, with $e_g(d)\to0$ as $d\to0$ and $\sup_d(\log\frac1d)^{\mathfrak p_0/4}e_g(d)<\infty$, such that for every tensor limit $\sigma$ at $g$, all $0<d_-<d<d_0(g)$, all $x_0,y_0$ with $3d\le|x_0-y_0|\le4d$ and all $H_1,H_2\in C^1_c$ with supports in the balls of radius $d_-$ about $x_0$, $y_0$, for all $21$ pairs of components,
\[|\sigma(H_1,H_2)-b_*^{-2}\mathsf W(H_1,H_2)|\le e_g(d)\,b_*^{-2}|H_1|_\infty|H_2|_\infty ;\]
the label $b_*$ is determined by $\sigma$. The same holds on every torus of $\mathfrak L(g)$ at every label of the interval $(g^{-2}+\bar B,X_g-\bar B]$, for Lipschitz weights, and at a finite cutoff in every periodic Gibbs state with an explicit remainder. Consequently, for $0<d<d_0(g)$, $x_0=\frac32d\,e_3$, every $\varphi\in C^1_c(\mathbb R^3)$ with $0\le\varphi\le1$, $\varphi=1$ on $B(x_0,d/128)$ and $\varphi=0$ off $B(x_0,d/64)$, and every $E\in V$,
\[\sigma(\theta(\varphi E),\varphi E)\ge(w_0-e_g(d))\,b_*^{-2}\operatorname{tr}E^2,\qquad w_0=\tfrac1{12}388^{-6}:\]
the two supports lie in the balls of radius $d/64$ about the points $\pm x_0$, whose distance is $3d$; $\mathsf W(\theta(\varphi E),\varphi E)\ge w_0\operatorname{tr}E^2$, which reads $\varphi\ge0$ and $\varphi=1$ on $B(x_0,d/128)$; and the remainder is at most $e_g(d)b_*^{-2}|E|_\infty^2\le e_g(d)b_*^{-2}\operatorname{tr}E^2$, which reads $\varphi\le1$. At the $d$ with $e_g(d)<w_0$, which exist since $e_g(d)\to0$, $\|[(\varphi E)]\|>0$ for $E\ne0$, and no real combination of the six fields vanishes.
\item[(E7)] \emph{(Extension over the diagonals; one torus.)} Let $0<g\le\gamma^{\rm ref}_W$, $b\in(g^{-2}+\bar B,X_g-\bar B]$ and $\ell\in\mathfrak L(g)$, as in (E3)~(ii); let $S_n$ be the cumulant functional of $n$ entries of the limit (E3)~(ii), continued complex linearly to $V_{\mathbb C}$-valued entries; let $h_0\in\ell2^{-\mathbb N}$ with $h_0\le\min\{\ell/(64n),\frac12L^{-s_{\rm th}(n)}\}$, $s_{\rm th}(n)$ as in (E2), and $l(\lambda)=1+\log(h_0/\lambda)$; the norms, the $\lambda$-supported functions, the classes $\mathfrak E_n(\kappa)$ and the coordinates at a partition are those of Definition~\ref{C-355-6}. Then $S_n$ extends to a distribution $\tilde S_n$ on the whole configuration space $T_\ell^n$, of order at most $3n-2$, with
\[|\tilde S_n(F)|\le Cb^{-n}l(\lambda)^{8n\mathfrak p_0+1}\|F\|_{(\lambda),3n-2}\qquad\text{for all }0<\lambda\le h_0\text{ and all $\lambda$-supported }F,\]
$C$ formed from the constants, $n$, $g$ and $\ell$; so $\tilde S_n\in\mathfrak E_n(0)$, the class of the power counting of $S_n$. The subtracted lattice cumulants, the sums of Lemma~\ref{C-355-34} over the Whitney cubes $Q$ with $s(1.52\,h(Q))\le K$, $h(Q)$ the side of $Q$, obey this bound with one number for all the cutoffs $K$, at all $\lambda$ with $8h_K\le\lambda\le h_0$, $h_K$ the least side of these cubes, and converge to $\tilde S_n$. A distribution on $T_\ell^n$ lies in $\mathfrak E_n(0)$ and is carried by the diagonals if and only if it is a sum
\[\sum_{|\pi|\le n-1}\ \sum_{|\alpha|\le3|\pi|}\partial_\xi^\alpha\delta_\pi\otimes T_{\pi,\alpha},\qquad T_{\pi,\alpha}\in\mathfrak E_{|\pi|}(|\alpha|)\text{ on }T_\ell^{|\pi|},\]
over the partitions $\pi$ of the points into at most $n-1$ clusters and the multi-indices $\alpha$ of the coordinates $\xi$ at $\pi$. So two extensions of $S_n$ in $\mathfrak E_n(0)$ differ by such a sum, and an extension plus any such sum is again an extension of $S_n$ in $\mathfrak E_n(0)$; a term with $|\alpha|>3|\pi|$ and a coefficient in $\mathfrak E_{|\pi|}(|\alpha|)$ vanishes. For translation-invariant, real, symmetric two-point contact terms in $\mathfrak E_2(0)$, the invariant space has dimension six under $\mathrm O(3)$ on $\mathbb R^3$, and twelve under the cubic group on $\mathbb R^3$ or $T_\ell$.
\end{enumerate}
\end{theorem}

\subsection{Construction estimates}

\begin{remark}[The input]\label{C-355-16}
\begin{enumerate}
\item[(1)] The real and sourced finite constructions, Theorems~\YMMref{D}{8} and~\YMMref{AJ}{673}, with the small-field and first-large-field estimates recalled below: on a torus that is small over the running couplings of a first-passage datum, the inductive description of the effective densities with sources under majorants, the step map and its domain, the term of the history without large-field regions, the weights $\mathsf w_t$ of the first creation of a large field at the step $t$, which on a small side obey $\mathsf w_t\le2M^{-3}\ell^3L^{3(K-t)}e^{-\frac12c'p_0(g_t)^2}$, the working level of a distance $d$, at which $a_K\le d/C_W(s(d))$ for $s(d)\le K$, the depth being formed with the room constant $C_{W,0}=\max\{1,4L^2M/A_0^2\}$, so that $C_W(s)\ge4L^2MY_s^{2\mathfrak p_0}$, and the floor of $\mathsf B$; on every side with (Sd.1) the same with $\mathsf w_t\le2M^{-3}\ell^3e^{\mathsf B(x_K)\ell^3}L^{3(K-t)}e^{-c'p_0(g_t)^2}$.
\item[(2)] Theorems~\YMMref{AJ}{673} and~\YMMref{AD}{607}, at the source $0$ alone: the bounds $-C_g\le p_\ell-E_{K,\ell,\beta}\,\ell^{-3}\le C_{\rm S}$ with the normalising energy $E_{K,\ell,\beta}$, and the comparison of the normalising energies of two sides.
\item[(3)] The flow and the label, Proposition~\YMMref{74.1}{615}, its room estimate~\YMMref{(1533)}{615}, and Lemma~\ref{lem:setup-stopped-arithmetic} with the flow bound \YMMref{(1530)}{615}, $|x_j-L^{-(j-i)}x_i|\le\bar B\,(1-L^{-(j-i)})$ for $0\le i\le j\le K$, and with \YMMref{(1531)}{615}, \YMMref{(1532)}{615} and Lemma~\ref{lem:setup-stopped-arithmetic}, \ref{lem:setup-pressure-stability}; the stopped sequence, Corollary~\YMMref{74.3}{616} and Lemma~\ref{lem:setup-stopped-arithmetic}; the first-passage windows $\mathcal W_K(g)$ and the interval $I(g)=(g^{-2}+\bar B,X_g-\bar B]$ of Corollary~\YMMref{74.2}{616}, with $L^K\cdot I(g)\subset\mathcal W_K(g)$; and the estimate \eqref{eq:setup-threshold-convexity}: for $0\le t\le k\le K$,
\[p_0(g_t)^2-p_0(g_k)^2\ge2\mathfrak p_0A_0^2\,y_k^{2\mathfrak p_0-1}(k-t)\log\tfrac L2 .\]
\item[(4)] The exact dilation of the lattice model (Propositions~\YMMref{119.1}{815}, \ref{lem:lattice-law-dilation}); the chessboard estimate and the seminorm as the mean of the mirror copies over a period box (Corollary~\ref{IV-18-18}, equation~\eqref{eq:period-box-seminorm}); the domination of the periodic Gibbs states (Proposition~\YMMref{91.9}{705}); their invariance under the lattice isometries, their reflection positivity at the nine families of lattice mirrors, the five-point inequality and the clustering at the rate $ra/|x|_\infty$ on one lattice (Propositions~\YMMref{103.4}{758} and~\YMMref{99.7}{742}, and Lemma~\ref{lem:spectral-steps}, Proposition~\textup{\YMMref{99.7}{742}}); the multiple reflections (Proposition~\YMMref{104.2}{759}); for the action density, the comparison of two consecutive cutoffs on one torus, the convergence of the laws of the block fields and the convergence along the full sequence of cutoffs (Theorem~\YMMref{AR}{810}, Corollaries~\textup{\ref{IV-62-66}}, \YMMref{117.1}{812}, Corollary~\YMMref{118.4}{815}), with the comparison of the background actions at two consecutive cutoffs and the loss of memory of the flows (Proposition~\YMMref{113.9}{795}, Propositions~\YMMref{114.2}{801}, \YMMref{114.3}{803}); the conditional expectations given the invariant $\sigma$-field (Proposition~\YMMref{100.1}{744}) and the clustering uniformly in the cutoff (Proposition~\YMMref{100.1}{744}); for the scalar limits, the canonical decomposition and its uniqueness (Theorem~\textup{\YMMref{AL}{746}}, Proposition~\textup{\YMMref{101.1}{747}}); the limit of a convergent Gibbs sequence at $g$ (Definition~\ref{C-365-5}(2)), which is invariant under all the isometries of $\mathbb R^3$ and is scalar data (Rec) with the numbers $16\,C^g_{\rm pt},4,\frac12$ (Theorem~\textup{\YMMref{AP}{780}}, Section~\YMMref{107}{765}), the scalar part of a tensor limit on $C^1_c$ entries being such a limit; for scalar data (Rec), the kernels, the one holomorphic function on the union of the tubes, the Lorentz covariance, the locality and the reconstruction of the Wightman field (Propositions~\YMMref{108.3}{771}, \YMMref{110.1}{776}, \YMMref{110.2}{777}, Theorem~\textup{\YMMref{AO}{778}}); two mirrors at the angle $\frac\pi4$ (Proposition~\YMMref{94.4}{722}); the kernel theorem on the non-coinciding configurations (Proposition~\YMMref{93.1}{715}); the extension of multilinear forms on separated test functions (Proposition~\YMMref{93.3}{717}) and the growth hypothesis in Proposition~\YMMref{93.5}{719}; the threshold $x_{\rm J}$ of the couplings (the definitions preceding \eqref{eq:weighted-history-threshold}).
\item[(5)] For the stress identity, in the setting (Th) of Section~\ref{sec:cutoff-calibration} with the top level $t=K-\lceil\phi K\rceil$ (Lemma~\YMMref{118.1}{812}) and the small-field chain of Proposition~\textup{\YMMref{115.2}{806}}: the chain with one insertion and with several insertions of disjoint carriers; the conditional means at the level $t$, and the cut of the configurations at a level by the field cuts (the definitions preceding Lemma~\YMMref{113.10}{796}, Lemma~\YMMref{113.10}{796}); the bound of the cumulants of insertions whose carriers are at sup distance $d_\times=\min\{d,2\delta\}-8a_K$, $d$ the separation and $\delta$ the margin of Proposition~\ref{C-355-30}~(c), so that $d_\times L^{K-i}\ge\frac12K^2$ at the levels $i\le t$ once $dL^{K-t}\ge K^2$ and $2\delta L^{\phi K}\ge K^2$, through the polymers that hold boxes of two insertions (\eqref{eq:analytic-domain-tail} and \eqref{eq:analytic-domain-diameter}), with a second member for the configurations with large fields; and the bound of the conditional mean of the cubic remainder, with its five inputs: the linearisation of a step (Proposition~\ref{prop:source-entry}), the size of an insertion at the background of a level, the parity of the births in the fluctuation field, the H\"older continuity of the backgrounds in local gauges (Theorem~\textup{\ref{I-56-65}}), and the comparison of a regular configuration with its mirror image under the inversion in a site (with Proposition~\ref{prop:ward-orbital-cubic}).

These are stated in this section, each with the idea of its proof: the bound of the cubic remainder and its five inputs in Proposition~\ref{C-355-30}~(b) and its sketch, the chains and the bound of the cumulants of insertions with distant carriers in the sketch of Proposition~\ref{C-355-30}~(b), (c).
\end{enumerate}
Of the volume, only the pressure gap of one unit cell at the source $0$ and the lower limit over the sides in the domination enter: no limit in the side is taken, and no clustering of a limit, no gap and no uniqueness of a state is used.
\end{remark}

\subsection{Auxiliary estimates}

\begin{lemma}[The pressure gap; transfer between sides]\label{C-355-17}
Let $0<g'\le\gamma$, let $(K,\beta)$ be a first-passage datum at $g'$, and let $\ell_1\le\ell_2$ lie in $\mathfrak L(g')$. Then $0\le p_{\ell_1}-p_{\ell_2}\le G_{g'}$. For a board of $\mathbb L_K$ of side $\mathsf l$ with $\ell_2\ge\ell_1\ge2\mathsf l$ and every bounded function $F$ of the bonds inside one cell, $\|F\|_{{\rm ch},\ell_2}\le e^{\mathsf l^3G_{g'}}\|F\|_{{\rm ch},\ell_1}$; for every $\nu\in\mathcal G^{\rm per}_{K,\beta}$, finitely many different cells and functions $\Phi_k$ of their bonds, $|\nu(\prod_k\Phi_k)|\le\prod_ke^{\mathsf l^3G_{g'}}\|\Phi_k\|_{{\rm ch},\ell_1}$. Finally, at the stopped coupling, $G_{g_{(i)}}-G_g=C_{g_{(i)}}-C_g\le\frac92i\log L+\frac18$ for $1\le i\le K-1$.
\end{lemma}

\begin{proof}[Proof sketch]
Theorems~\YMMref{AJ}{673} and~\YMMref{AD}{607}, at the source $0$, give $-C_{g'}\le p_\ell-E_{K,\ell,\beta}\,\ell^{-3}\le C_{\rm S}$ at both sides $\ell=\ell_1,\ell_2$, uniformly in the cutoff, $E_{K,\ell,\beta}$ its normalising energy, and $|E_{K,\ell_1,\beta}\,\ell_1^{-3}-E_{K,\ell_2,\beta}\,\ell_2^{-3}|\le2C_{\rm tt}C'_{\rm tt}E_0$; together these give $p_{\ell_1}-p_{\ell_2}\le C_{\rm S}+C_{g'}+2C_{\rm tt}C'_{\rm tt}E_0=G_{g'}$, and the chessboard estimate on the larger torus, applied to the constant function (Corollary~\ref{IV-18-18}~(e)), gives $p_{\ell_1}\ge p_{\ell_2}$. The chessboard estimate converts the gap into the transfer of seminorms, and the domination of the periodic Gibbs states, with the lower limit taken along the sides of $\mathfrak L(g')$, into the bound in the states. Since $C_{\rm S}$, $C_{\rm tt}$, $C'_{\rm tt}$, $E_0$ depend on no coupling, $G_{g_{(i)}}-G_g=C_{g_{(i)}}-C_g$, which \eqref{C-355-2} bounds.
\end{proof}

\begin{lemma}[Change of the length scale]\label{C-355-18}
Let $0<g\le\gamma^{\rm sd}_W$, $(K,\beta)$ a first-passage datum at $g$, $1\le i\le K-1$, $k=K-i$, $\ell\in\mathfrak L(g)$, $0<\delta\le\frac14$.
\begin{enumerate}
\item[(a)] The map $x\mapsto L^ix$ carries $\mathbb L_K$ onto $\mathbb L_k$ and $\Lambda_{K,\ell}$ onto $\Lambda_{k,L^i\ell}$, bonds onto bonds and plaquettes onto plaquettes, the measure at $(K,\beta)$ onto the measure at $(k,\beta)$, and the board $(i,o)$ onto the board $(0,L^io)$ of side $1$, with its cells and mirror maps; with $H'=L^iH(L^{-i}\cdot)$, $\mathcal T_K(H)$ goes over into $\mathcal T_k(H')$, so that $\|\mathcal T_K(H)-c\|^{(i,o)}_{{\rm ch},\ell}$ at $(K,\beta)$ equals $\|\mathcal T_k(H')-c\|^{(0,L^io)}_{{\rm ch},L^i\ell}$ at $(k,\beta)$, and $|H'|_\infty/(L^ib)=|H|_\infty/b$.
\item[(b)] $(k,\beta)$ is a first-passage datum at $g_{(i)}<g$ with label $L^ib$; $X_{g_{(i)}}=x_k$; $L^i\ell\in\mathfrak L(g_{(i)})$; and $L^{i-1}(g^{-2}-\bar B)-B_\sharp\le g_{(i)}^{-2}\le L^{i-1}(X_g+\bar B)-B_\sharp$.
\item[(c)] Let $\ell''$ be the least side at $g_{(i)}$, $N=\ell''^3$, and $s=s(\delta)$ formed at $g_{(i)}$. Then $8\le N\le(c'A_0^2/9)\,y_k^{2\mathfrak p_0-1}\le(c'A_0^2/9)\,\mathsf Y_i^{2\mathfrak p_0-1}$, $\log g_{(i)}^{-2}\le Y_s\le\mathsf Y_i$, $L^s<Lc_W\mathsf Y_i^{2\mathfrak p_0}/\delta$, $\bar x_s\le L^2c_W\delta^{-1}\mathsf Y_i^{2\mathfrak p_0}(X_g+\bar B)L^i$, and $s\le s_E(i)$, the depth $s(\delta)$ formed at the coupling whose inverse square is $L^{i-1}(X_g+\bar B)-B_\sharp$.
\item[(d)] There is a number $\mathsf x_E(\delta)$, formed from the constants and $\delta$ alone, which may be taken equal to $\mathsf x^*/\delta$ with $\mathsf x^*$ formed from the constants alone, such that the thresholds (T1) at $n=N$ and (T2) of Proposition~\ref{C-355-19} hold at $g'=g_{(i)}$, $K'=k$, $\ell''$, $d=\delta$ as soon as $g_{(i)}^{-2}\ge\mathsf x_E(\delta)$; so for $i\ge i_E(g,\delta)$, the least $i\ge1$ with $L^{i-1}(g^{-2}-\bar B)-B_\sharp\ge\mathsf x_E(\delta)$; and $L^{i_E}\le c_i\delta^{-1}$ with $c_i=\max\{L,\,L^2(\mathsf x^*+B_\sharp)/(g^{-2}-\bar B)\}$, a number formed from the constants and $g$.
\end{enumerate}
\end{lemma}

\begin{proof}[Proof sketch]
(a) is the exact dilation of the lattice model, Lemma~\ref{lem:lattice-law-dilation}, a relabelling; since $a_K^{-1}=L^ia_k^{-1}$, the factor $a_K^{-1}$ in Definition~\ref{C-355-3} turns $H$ into $L^iH(L^{-i}\cdot)$. (b) is Corollary~\YMMref{74.3}{616} and Lemma~\ref{lem:setup-stopped-arithmetic}, with the flow bound $|x_K-L^{-i}x_k|\le\bar B(1-L^{-i})\le\bar B$ of Proposition~\YMMref{74.1}{615}, \YMMref{(1530)}{615}. (c) Since $g_{(i)}\le\gamma^{\rm sd}_W$, the least side $\ell''$ is small over the running couplings of $(k,\beta)$; (Sd.4) with the floor of $\mathsf B$ gives $\frac92y_kN\le\mathsf B(x_k)\ell''^3\le\frac12c'A_0^2y_k^{2\mathfrak p_0}$: the number of cells of the least torus is a power of a logarithm, not of the order $a_i^{-3}$. The depth satisfies $L^s\delta<LC_W(s)=Lc_WY_s^{2\mathfrak p_0}$, so $\bar x_s=L^{s+1}(g_{(i)}^{-2}+\bar B)<L^2c_W\delta^{-1}Y_s^{2\mathfrak p_0}(g_{(i)}^{-2}+\bar B)$, and $2\mathfrak p_0\log Y\le\frac12Y+2\mathfrak p_0\log(4\mathfrak p_0)$ gives $Y_s\le\mathsf Y_i$ with (b). (d) By (c), $N$ grows like the power $2\mathfrak p_0-1$ of $Y_s$ while (T1) asks $2Y_s^{2\mathfrak p_0}\ge u_*(N)$, a quantity of the order $N\log N$; the quotient of the two sides tends to infinity with $Y_s\ge\log g_{(i)}^{-2}$. (T2) holds because $L^{3s}$ is a power of a logarithm of $g_{(i)}^{-2}$, hence $\le g_{(i)}^{-40}<x_k^{20}$ from a number on.

The threshold can be taken proportional to $\delta^{-1}$, and then $L^{i_E}\le c_i\delta^{-1}$ by (b). In numbers: put $x=g_{(i)}^{-2}$ and $Y=\log x$; then $x_k=L(x+B_\sharp)>x$, and $y_k\le Y+\log2L$ as soon as $x\ge B_\sharp$. Let $\mathsf Y^*\ge1$ be the least number such that $2Y'^{2\mathfrak p_0}\ge u_*\bigl((c'A_0^2/9)(Y'+\log2L)^{2\mathfrak p_0-1}\bigr)$ for all $Y'\ge\mathsf Y^*$, $u_*$ being read at real arguments by its formula (the right member is of the order $Y'^{2\mathfrak p_0-1}\log Y'$), and put
\[\mathsf c_5=\tfrac{74}{11}\,L^7c_W^4\,(12\mathfrak p_0)^{6\mathfrak p_0}(4\mathfrak p_0)^{2\mathfrak p_0},\qquad\mathsf x^*=\max\bigl\{1,\,B_\sharp,\,e^{\mathsf Y^*},\,L^23^{\mathfrak p_0},\,\mathsf c_5\bigr\},\]
numbers formed from the constants alone. Let $x\ge\mathsf x^*/\delta$. (T1) at $n=N$: by (c), $N\le(c'A_0^2/9)(Y+\log2L)^{2\mathfrak p_0-1}$ and $Y\le Y_s$; as $t\mapsto t[c_1+\log(c^*t)]$ increases on $[1,\infty)$, $u_*$ does not decrease there, and $Y\ge\mathsf Y^*$ gives $u_*(N)\le2Y^{2\mathfrak p_0}\le2Y_s^{2\mathfrak p_0}$. (T2): $y_k>\log x\ge2\log L+\mathfrak p_0\log3$. By (c), $L^{3s}<(Lc_W/\delta)^3\mathsf Y_i^{6\mathfrak p_0}$; the inequality $t^q\le(2q)^qe^{t/2}$ ($t,q>0$) gives $\mathsf Y_i^{6\mathfrak p_0}\le(12\mathfrak p_0)^{6\mathfrak p_0}e^{\mathsf Y_i/2}=(12\mathfrak p_0)^{6\mathfrak p_0}c_0\delta^{-1}L^i$ with $c_0=L^2c_W(4\mathfrak p_0)^{2\mathfrak p_0}(X_g+\bar B)$; by (b), $L^i\le L(x+B_\sharp)/(g^{-2}-\bar B)\le2Lx/(g^{-2}-\bar B)$; and $(X_g+\bar B)/(g^{-2}-\bar B)=L(g^{-2}+\bar B)/(g^{-2}-\bar B)\le\frac{37}{11}L$, because $\bar B\le\frac{13}{24}g^{-2}$ for $g\le\gamma$ (Lemma~\ref{lem:setup-stopped-arithmetic}).

Together $L^{3s}<\mathsf c_5\delta^{-4}x\le x^{20}<x_k^{20}$, since $x^{19}\ge\max\{1,\mathsf c_5\}^{19}\delta^{-19}\ge\mathsf c_5\delta^{-4}$. So $\mathsf x_E(\delta)=\mathsf x^*/\delta$: the reference coupling enters the threshold through the quotient $(X_g+\bar B)/(g^{-2}-\bar B)$ alone, and its bound $\frac{37}{11}L$ is independent of the coupling. By (b), $g_{(i)}^{-2}\ge L^{i-1}(g^{-2}-\bar B)-B_\sharp$, which increases with $i$; if $i_E\ge2$, then $L^{i_E-2}(g^{-2}-\bar B)-B_\sharp<\mathsf x^*/\delta$, whence $L^{i_E}<L^2(\mathsf x^*+B_\sharp)\delta^{-1}/(g^{-2}-\bar B)\le c_i\delta^{-1}$; if $i_E=1$, then $L\le c_i\le c_i\delta^{-1}$.
\end{proof}

\begin{proposition}[As many sources as a small side has unit cubes]\label{C-355-19}
There are numbers $A_\kappa>0$, $A_\flat\in(0,1]$, $c^*\ge4$, $c_2$, $C^I$, $C_2$, $C_1$, formed from the constants alone, with $c_2\ge c_1:=\log[(4/3)^{1/2}L^{1/2}2^{\mathfrak p_0}]$; put
\[
\begin{multlined}
u_*(n)=\max\bigl\{4,\ \relax (2/A_\kappa)(n[c_1+\log(c^*n)]+\log3)\bigr\},\qquad u'_*(n)=\max\bigl\{4,\ \relax (2/A_\flat)(n[c_2 \\
+\log(4n)]+3\log L+\log3)\bigr\},
\end{multlined}
\]
\[\mathsf A_n(s)=2^{2\mathfrak p_0+1}n\max\{C^I,C_2n^{3+\mathfrak p_0}\}(LC_W(s))^3\,\bar x_s\,Y_s^{2\mathfrak p_0}.\]
Let $0<g'\le\gamma^{\rm ref}_W$, $(K',\beta)$ a first-passage datum at $g'$ and $\ell$ a side small over its running couplings. A \emph{source under the majorant} $(\mathsf s,S)$ is carried by the description of Theorem~\YMMref{D}{8} as a triple: a continuous gauge-invariant function $F$ of the bonds with $|F|\le a_{K'}^{-1}\sum_{p\in S}\mathsf s\,a_p$; a family of effective terms of the level $0$; and families of the levels $1,\dots,k$ up to the level $k$ of the chain; each family sums to $F$ up to a constant, vanishes on the polymers that miss $S$, and is bounded by the size that the weight $\mathsf s\mathbf 1_S$ has at its level. $F$ is the \emph{first member} of the source; it is the function that is itself called the source when the triple is understood. A \emph{source system} is a finite family of sources under majorants, entered with independent complex parameters. In (a) let $0<d<1$ with $s=s(d)\le K'$ (formed at $g'$), $k=\max\{K'-s,2\}$, $\upsilon=1/(a_{K'-k}x_k)$, so that $\upsilon^2\le\frac98(a_{K'}\beta)^{-2}$, and $\hat d=d/a_{K'-k}$, the distance $d$ in the units of the level $k$; since $K'-k\in\{s-2,s-1,s\}$ and $C_W(s)\ge4L^2MY_s^{2\mathfrak p_0}$, the definition of the depth gives $4MY_s^{2\mathfrak p_0}\le L^{-2}C_W(s)\le\hat d<LC_W(s)$.
\begin{enumerate}
\item[(a)] Let $2\le n\le\ell^3$ and let $F_1,\dots,F_n$ be the first members of a source system under majorants $(\mathsf s_a,S_a)$ such that every plaquette of $S_a$ has its barycentre within $\frac78d$ of a point $x_a$ and barycentres of plaquettes of different $S_a$ are at distance $\ge\frac54d$. Suppose
\[\text{(T1)}\quad2Y_s^{2\mathfrak p_0}\ge u_*(n);\qquad\text{(T2)}\quad y_{K'}\ge2\log L+\mathfrak p_0\log3\ \text{ and }\ L^{3s}\le x_{K'}^{20}.\]
Then there are complex numbers $c_a$, the first-order terms of the sources at the trivial background summed up to the level $k$, such that for every non-empty $B\subset\{1,\dots,n\}$
\[\Bigl|\Bigl\langle\prod_{a\in B}(F_a-c_a)\Bigr\rangle_{K',\ell,\beta}\Bigr|\le\prod_{a\in B}\mathsf A_n(s)\,\upsilon\,\mathsf s_a .\]
\item[(b)] Let instead $s$ be any integer with $2\le k:=K'-s\le K'-1$, no distance being given, and $1\le n\le\ell^3$; let the plaquettes of each $S_a$ lie in the interior of a cell of its own of a board of side $1$ of the torus, and, if $n\ge2$, let the least sup distance between a corner of a plaquette of one $S_a$ and a corner of a plaquette of another, in the units of the level $k$, be at least $Mu'_*(n)$. Then, with the numbers $c_a$ of (a) formed at this level $k$, for every non-empty $B\subset\{1,\dots,n\}$ with $r=\#B$ members
\[\Bigl|\Bigl\langle\prod_{a\in B}(F_a-c_a)\Bigr\rangle_{K',\ell,\beta}\Bigr|\le\prod_{a\in B}2r\,C_1L^{4s}\,\mathsf s_a\le\prod_{a\in B}2n\,C_1L^{4s}\,\mathsf s_a .\]
\end{enumerate}
\end{proposition}

\begin{proof}[Proof sketch]
Put $\sum_a\lambda_aF_a$ into the exponent. The normalized source sizes and mixed-derivative estimates used below are detailed in the proof of Lemma~\ref{cov-mom}; the moment argument for (a) is given in Lemma~\ref{bil:prod}. By Theorem~\YMMref{D}{8} the effective actions $E^{(j)}_\lambda$ of the levels $j\le k$ are holomorphic in $\lambda$ on a polydisc whose radii are governed by sizes $\hat g_{j,a}$, which grow by at most $L^{1/2}$ per step on the small fields. For $I\subset\{1,\dots,n\}$ with $r=\#I\ge2$, the mixed derivative $\partial_IE^{(j)}_\lambda|_0$ vanishes on every polymer that misses some $S_a$, $a\in I$, and so carries the factor $e^{-A_\kappa[\mathsf D_j/M-2]_+}$, $\mathsf D_j$ the least sup distance between a corner of a plaquette of one support and a corner of a plaquette of another, in the units of the level $j$; by the geometry of the majorants and the working level, $\mathsf D_k\ge(5/(4\sqrt3))\hat d-1$ with $\hat d\ge4MY_s^{2\mathfrak p_0}$, so that $u:=\mathsf D_k/M\ge(5/\sqrt3)Y_s^{2\mathfrak p_0}-1\ge2Y_s^{2\mathfrak p_0}$ ($Y_s\ge7$), and (T1) gives $\mathsf D_j/M=L^{k-j}u\ge u_*(n)$ at all levels $j\le k$. An induction on $r$, and for each $r$ on $j$, with the chain rule of all orders for the step map and Cauchy's estimate at the radius $1/(4r\hat g_{j,a})$, gives
\[\bigl\|\partial_IE^{(j)}_\lambda\big|_0\bigr\|\le c_\theta\,(c^*r)^r\prod_{a\in I}\hat g_{j-1,a},\]
with a number $c_\theta$ formed from the constants alone, the same at every level $j$, for every $n$ and every source system; the induction closes because $u\ge u_*(n)$ makes the series of the polymer sums times $(c^*m)^m$ at most $1$ for $m\le n$, and because on the polydisc the perturbation of the arguments is at most $0.27$ times the radius of the domain of the step map, so that they stay in it. The centred moment is a sum over the partitions of $B$ of products of these derivatives on the history without large-field regions, plus the sum over the step $t$ of the first creation of a large field with the weight $\mathsf w_t$, in which every insertion loses at most $L^4$ per scale. On a small side this loss is paid: by \eqref{eq:setup-threshold-convexity}, (Sd.4) and the floor of $\mathsf B$,
\[\tfrac12c'p_0(g_t)^2\ge\ell^3\Bigl[\tfrac92y_{K'}+2+9\mathfrak p_0(K'-t)\log\tfrac L2\Bigr],\]
and since $r\le n\le\ell^3$ and $9\log\frac L2>4\log L$, the sum over $t$ of $L^{4r(K'-t)}\mathsf w_t$ is at most $L^{3s}x_{K'}^{-20}\le1$ by (T2). In (a) the number $n$ enters at five places: through $u_*(n)$; through the number of partitions, ${\rm Bl}(r)\le r^r\le n^r$; through $r\le\ell^3$; through the domain in $\mathbb C^n$, of which only the derivatives at $0$ are used; and through the sizes of the mixed derivatives of order $r\le n$, summed over the polymers and the levels, whose $r$-th root is at most $C_2r^{3+\mathfrak p_0}$: the factor $(c^*r)^r$ of the Cauchy estimate gives one power of $r$, and the sum over the lag $q$ between the level of a polymer and the level of the moment, $\sum_{q\ge0}L^{(3+r/2)q}(q+1)^{r\mathfrak p_0}e^{-3(L^q-1)}\le5\,(r(2+\mathfrak p_0))^{r(2+\mathfrak p_0)}$, gives $r^{2+\mathfrak p_0}$.

With ${\rm Bl}(r)\le n^r$ these give the factor $n\max\{C^I,C_2n^{3+\mathfrak p_0}\}$ of $\mathsf A_n(s)$, $C^I$ bounding the sizes of the first order. In (b) the level $k=K'-s$ is given, and the threshold $u'_*(n)$ of the distance in its units takes the place of (T1): the polymer sums of the levels below $k$ now carry the factor $L^{3n'}$ of the number of the cubes of the lag $n'$, whence the term $3\log L$ and the exponent $A_\flat$. The sizes are taken crudely, $C_1L^{4(K'-t)}\mathsf s_a$ for a source at the level $t$, and the exponential smallness of a polymer that joins two cells absorbs the factor $(4m)^m$ of the mixed derivatives of order $m$, so that in (b) the number $n$ enters at the first four places alone.

The same sum over $t$ then gives $|\langle\prod_B(F_a-c_a)\rangle|\le{\rm Bl}(r)\,C_1^r\prod_B\mathsf s_a\,[L^{4sr}+\sum_tL^{4(K'-t+1)r}\mathsf w_t]$, the bracket is at most $2L^{4sr}$, and ${\rm Bl}(r)\le r^r$.
\end{proof}

\begin{proposition}[$n$ sources with the sharper number]\label{C-355-20}
Let $n\ge2$ and $0<g\le\gamma^{\rm ref}_W$. For $u>0$ and $r\ge1$ the numbers $T_r(u)$ are those of the paragraph \emph{Sources} of the setting,
\[T_r(u)=\sum_{n'\ge1}(n'+1)L^{-(n'-1)}\bigl[(4/3)^{1/2}L^{n'/2}(1+n')^{\mathfrak p_0}\bigr]^r\,e^{-A_\kappa[L^{n'-1}u-2]_+},\]
with the number $A_\kappa$ of Proposition~\ref{C-355-19}, which is the number $4\beta_{\rm s}\kappa$ in the exponent of the function $\mathsf T$ of the definitions preceding \eqref{eq:weighted-history-threshold}. The series converges, the quotient of two consecutive terms tending to $0$; $T_r(u)$ does not increase in $u$ and does not decrease in $r$, the bracket being $\ge1$; and $T_2=\frac43L\,\mathsf T$. The condition (N.$\theta$) at $(n,g,s)$ is $2T_n(3)\,\theta(L^{s-2}g^{-2})\le1$, $\theta(\cdot)$ the function of the definitions preceding \eqref{eq:weighted-history-threshold}; if it holds at $(n,g,s)$, it holds at $(n',g,s)$ for $2\le n'\le n$. There is a function $s\mapsto k'_n(g,s)\in(0,\infty)$, formed from the constants, $n$, $g$ and $s$ alone, such that:
\begin{enumerate}
\item[(a)] for every first-passage datum $(K,\beta)$ at $g$, $b=a_K\beta$, every side $\ell$ small over its running couplings, every $0<d<1$ with $s=s(d)\le K$ and (N.$\theta$) at $(n,g,s)$, all majorants $(\mathsf s_a,S_a)$ with the geometry of Proposition~\ref{C-355-19}~(a), and every source system under them, real or not, with the first members $F_1,\dots,F_n$,
\[|\langle F_1;\dots;F_n\rangle_{K,\ell,\beta}|\le k'_n(g,s(d))\,b^{-n}\prod_a\mathsf s_a ;\]
\item[(a$'$)] the same holds on every side $\ell=2L^m$ with (Sd.1), $MR(x_K)\le L^m$, in particular on every $\ell\in\mathfrak L(g)$, with a number $k'^{,\ell}_n\le\max\{1,\mathsf c(\ell)\}^nk'_n$ in the place of $k'_n$;
\item[(b)] as $s\to\infty$, at fixed $n$ and $g$, the leading size per entry is $162B_3^2A_0^2Y_s^{2\mathfrak p_0}(LC_W(s))^3$, $B_3$ a constant; there are numbers $C_n(g)$, $C(g)$ with $k'_n(g,s)\le C_n(g)Y_s^{8n\mathfrak p_0}$ for $s\ge3$ and $Y_{s(d)}\le2\log\frac1d+C(g)$; and uniformly in the coupling, at every depth $s=s(d)$ of a distance $0<d<1$: $k'_n(g,s(d))\le C_nY_{s(d)}^{8n\mathfrak p_0}$ with a number $C_n$ formed from the constants and $n$ alone, and $Y_{s(d)}\le2\log(g^{-2}/d)+c$, $c$ a constant;
\item[(c)] (N.$\theta$) holds at $(n,g,s)$ from a depth on, and at $n=2$ for every $s\ge3$.
\end{enumerate}
\end{proposition}

\begin{proof}[Proof sketch]
Let $k=\max\{K-s,2\}$, $y=y_k$, $\upsilon=1/(a_{K-k}x_k)$ and $\hat d=d/a_{K-k}$ be formed as in Proposition~\ref{C-355-19}, and $u=\mathsf D_k/M$, $\mathsf D_k$ the least sup distance of the supports in the units of the level $k$ as there. The proof has two layers. In the first, as in Proposition~\ref{C-355-19}, the centred moments satisfy $|\langle\prod_B(F_a-c_a)\rangle|\le{\rm Bl}(r)[\mathsf A^r+\sum_{t\le k}\mathsf{bb}_{k-t+1}^r\mathsf w_t]\prod_{a\in B}\hat m_a$, with $\hat m_a=\upsilon\mathsf s_a$, ${\rm Bl}(r)$ the number of partitions, $\mathsf A$ the size on the history without large-field regions and $\mathsf{bb}_{n'}=\mathsf c_{\mathsf B}(n)\,\hat d^3x_ky^{2\mathfrak p_0}L^{4n'}(n'+2)^{2\mathfrak p_0}$, $\mathsf c_{\mathsf B}(n)$ a number formed from the constants and $n$; here (N.$\theta$) replaces (T1) in the induction for the mixed derivatives: at a level $j\le k$ the step carries the number $\theta(x_j)$ (the definitions preceding \eqref{eq:weighted-history-threshold}), and $\theta(x_j)\le\theta(L^{s-2}g^{-2})$, since $x_j>L^{s-2}g^{-2}\ge g^{-2}\ge\mathsf x_W\ge x_{\rm J}\ge e^{2\mathfrak p_0}$ and $\theta$ does not increase on $[e^{2\mathfrak p_0},\infty)$, while $\mathsf D_j/M=L^{k-j}u\ge3$; so (N.$\theta$) gives $\theta(x_j)\,T_m(\mathsf D_j/M)\le\frac12$ for $2\le m\le n$, and the series $T_m$ of the polymer sums enter the induction exactly through this product.

At the level $k$ the first-order size of one entry is $162B_3^2A_0^2y^{2\mathfrak p_0}\hat d^3$: of the exponent $8\mathfrak p_0$ per entry, $2\mathfrak p_0$ is the supremum of $x_ka_p$ on the small fields and $6\mathfrak p_0$ comes from the number $\hat d^3<(LC_W(s))^3$ of unit cubes of the working level in a ball of radius $d$; the higher orders are smaller by $e^{-(u-2)}$, $u-2\ge(5/\sqrt3)Y_s^{2\mathfrak p_0}-3$. In the second layer the weights are bounded, $\mathsf w_t\le\mathsf L_0L^{3(K-t)}e^{-c''A_0^2y_t^{2\mathfrak p_0}}$, with $\mathsf L_0=M^{-3}c'A_0^2Y_s^{2\mathfrak p_0}$, $c''=\frac12c'$ on a small side and $\mathsf L_0=2M^{-3}\ell^3e^{\mathsf B(x_K)\ell^3}$, $c''=c'$ on a side with (Sd.1): this is the only place where the side enters.

Since $y_{k-a}\ge y+a\log\frac{3L}4$, and $\hat d<LC_W(s)$, $x_k\le\bar x_s$, $y\le Y_s$, the sum over the step of the first creation is at most
\[\mathsf{BB}^r\,\mathsf L_0\,L^{3s}\,S''_r\,e^{-c''A_0^2(Y_s-\log2L^3)^{2\mathfrak p_0}},\qquad\mathsf{BB}=\mathsf c_{\mathsf B}(n)\,(LC_W(s))^3\,\bar x_s\,Y_s^{2\mathfrak p_0},\]
with $S''_r=\sum_{a\ge0}L^{3a+4r(a+1)}(a+3)^{2r\mathfrak p_0}e^{-\frac12c'A_0^2(a\log\frac{3L}4)^2}$, a number formed from the constants and $r$. As $\bar x_s=e^{Y_s}$ and $L^{3s}\le e^{3Y_s}$, this is $\mathsf L_0$ times a polynomial in $Y_s$ times $e^{(r+3)Y_s-c''A_0^2(Y_s-\log2L^3)^{2\mathfrak p_0}}$, which tends to $0$ as $s\to\infty$ because $2\mathfrak p_0\ge2$; the factor $\mathsf L_0$ is a power of $Y_s$ on a small side, and on a side with (Sd.1) it reads the side and not $s$. Cumulants are obtained from the moments by the recursion over partitions, $\langle F_1;\dots;F_n\rangle=\langle\prod_a(F_a-c_a)\rangle-\sum_{\#\pi\ge2}\prod_{\gamma\in\pi}\langle\cdot\rangle^{\rm cum}_\gamma$, and $\prod_a\hat m_a\le(9/8)^{n/2}b^{-n}\prod_a\mathsf s_a$.

For (b), $Y_s<\log\frac1d+2\log L+\log c_W+\log(g^{-2}+\bar B)+2\mathfrak p_0\log Y_s$ at $s=s(d)$ by the definition of the depth, and $2\mathfrak p_0\log Y\le\frac12Y+2\mathfrak p_0\log(4\mathfrak p_0)$ gives $Y_{s(d)}\le2\log\frac1d+C(g)$; with $\log(g^{-2}+\bar B)\le\log g^{-2}+\log\frac{37}{24}$ (Lemma~\ref{lem:setup-stopped-arithmetic}) it gives $Y_{s(d)}\le2\log(g^{-2}/d)+c$. The bound uniform in the coupling holds at the depths of distances: at $s=s(d)$, $0<d<1$, one has $L^s\ge C_W(s)/d>C_W(s)\ge4L^2MY_s^{2\mathfrak p_0}$, so that $L^{s-1}/2\ge2LMY_s^{2\mathfrak p_0}$; besides, $gL^{-(s-2)/2}\le c\,e^{-Y_s/2}$, $L^{3s}\le e^{3Y_s}$ and $\bar x_s=e^{Y_s}$, $c$ a constant. Besides through $Y_s$, the number $k'_n(g,s)$ reads the coupling only through the correction to the first-order size of one entry, which carries $gL^{-(s-2)/2}$ and $e^{-\mathsf{ds}(g,s)}$ with the number $\mathsf{ds}(g,s)=\min\{c_{\rm d}L^{(s-2)/2}/(A_0\,g\,Y_s^{\mathfrak p_0}),\ L^{s-1}/2-2\}-1$, $c_{\rm d}$ a constant.

This number bounds from below the quantity $\min\{c_{\rm d}\,\varepsilon_k^{-1},\ (\ell/(2a_{K-k})-4)/(2M)\}-1$ of the datum and the side that stands in its place in the size at the level $k$, $\varepsilon_k=A_0y^{\mathfrak p_0}x_k^{-1/2}$ being the size of the field there and $\ell/a_{K-k}$ the side of the torus in the units of that level: $\varepsilon_k\le A_0Y_s^{\mathfrak p_0}gL^{-(s-2)/2}$, as $y\le Y_s$ and $x_k>L^{s-2}g^{-2}$, and $\ell/(2a_{K-k})\ge LM\,L^{s-2}$, as $K-k\ge s-2$ and $\ell=2L^m\ge2LM$. So the size at the level $k$ is at most the expression formed with $\mathsf{ds}(g,s)$ in the place of that quantity. At $s=s(d)$, by $L^{s-1}/2\ge2LMY_s^{2\mathfrak p_0}$ and $gL^{-(s-2)/2}\le c\,e^{-Y_s/2}$, one has $\mathsf{ds}(g,s)\ge\min\{(c_{\rm d}/(cA_0))\,e^{Y_s/2}Y_s^{-\mathfrak p_0},\,2LMY_s^{2\mathfrak p_0}-2\}-1$.

So at $s=s(d)$ every quantity of which $k'_n(g,s)$ is formed is bounded by a function of $Y_s$ alone, which is of the order $Y_s^{8n\mathfrak p_0}$ for $Y_s\ge7$, and $C_n$ is independent of the coupling. For (c), $\theta(t)$ is a constant times $t^{-1/2}(\log t)^{\mathfrak p_0}$, which decreases to $0$; at $n=2$, $T_2(3)=\frac43L\,\mathsf T(3)$ with the function $\mathsf T$ of the definitions preceding \eqref{eq:weighted-history-threshold}, and since $L^{s-2}g^{-2}\ge g^{-2}\ge\mathsf x_W\ge x_{\rm J}$ for $s\ge2$ and $\theta$ does not increase, the condition $\frac43C^\flat_1L\,\mathsf T(3)\,\theta(x)\le1$ of the threshold $x_{\rm J}$ in \eqref{eq:weighted-history-threshold}, with $C^\flat_1\ge2$, gives $2T_2(3)\,\theta(L^{s-2}g^{-2})\le1$.
\end{proof}

\begin{lemma}[The forms as sources]\label{C-355-21}
Let $0<g'\le\gamma^{\rm ref}_W$, $(K',\beta)$ a first-passage datum at $g'$ and $\ell$ a side with (Sd.1), small over the running couplings or not; sources are those of the description of Theorem~\YMMref{D}{8} at these data. Let $w$ be a complex function on the plaquettes and $\mathsf w$ a complex function on the bonds, of finite supports. The plaquette form $\mathcal O(w)$ is the first member of a source under the majorant $(|w|_\infty,\operatorname{supp}w)$, and the bond form $Y[\mathsf w]$ of a source under $(59|\mathsf w|_\infty,S[\mathsf w])$, $S[\mathsf w]$ the plaquettes at the bonds where $\mathsf w\ne0$; both sources are complex linear in the weight; and if $|w|_\infty\le h$ and $|\mathsf w|_\infty\le h$, then $\mathcal O(w)+Y[\mathsf w]$ is the first member of a source under $(60h,\operatorname{supp}w\cup S[\mathsf w])$.

A lattice isometry $\mathsf r$ with linear part $\operatorname{diag}(\sigma_1,\sigma_2,\sigma_3)$, $\sigma_\mu=\pm1$, carries $\mathcal O(w)$ to $\mathcal O(w\circ\mathsf r^{-1})$ and $Y[\mathsf w]$ to $Y[\mathsf r^\flat\mathsf w]$, $(\mathsf r^\flat\mathsf w)(\mathfrak b')=\sigma_{\lambda}\sigma_{\lambda'}\mathsf w(\mathsf r^{-1}\mathfrak b')$, and leaves the torus means invariant. Hence, for every bounded $H:\mathbb R^3\to V_{\mathbb C}$ of bounded support, $\mathcal T_{K'}(H)$ is the first member of a source under $(60|H|_\infty,S_H)$, $S_H=\operatorname{supp}w_H\cup S[\mathsf w_H]$ the plaquettes on which it depends, and so are its mirror copies.
\end{lemma}

\begin{proof}
The plaquette-source assertion is the weighted progression of the construction. Lemma~\ref{C-371-51} supplies the bond-source families, including the action majorant, the height-zero and later-level norms, locality and the exact sum rule. Thus the bond conclusion uses the actual analytic source families, rather than boundedness of the bare polynomial alone. Source linearity follows from the linear dependence of these families on the weights. Increasing a majorant preserves all its defining inequalities, and adding sources adds their sizes and takes the union of their supports; $1+59=60$ gives the mixed-source bound.

Finally, $a_p\circ\mathsf r^{-1}=a_{\mathsf rp}$ and $q_{\mathfrak b}\circ\mathsf r^{-1}=\sigma_\lambda\sigma_{\lambda'}q_{\mathsf r\mathfrak b}$, with the reflected bond taken in its positive orientation. These identities give the stated reflection law and the assertion for mirror copies.
\end{proof}

\begin{proposition}[Two consecutive cutoffs on one torus for the plaquette forms of one plane and the bond forms]\label{C-355-36}
Let $0<g\le\gamma^{\rm ref}_W$, $b\in I(g)$, $\ell\in\mathfrak L(g)$, $n\ge2$ and $\mathsf d>0$; by Corollary~\YMMref{118.4}{815} the pairs $(K,L^Kb)$ are first-passage data at $g$ with label $b$. For a bounded Borel function $\psi$ on $T_\ell$ and $x\in\Lambda_{K,\ell}$ put $\psi^{(K)}(x)=a_K^{-3}\int_{x+a_K[-\frac12,\frac12)^3}\psi$. A \emph{cell insertion} at the cutoff $K$ is one of the following: the plaquette form $\mathcal O(w)$ of one plane $\{\kappa,\kappa'\}$, with $w(p)=\psi^{(K)}(x)$ at the plaquette $p$ spanned at $x$ by $a_Ke_\kappa$, $a_Ke_{\kappa'}$ and $w=0$ at the plaquettes of the two other planes; the sum of the three, the action density with the weight $\psi^{(K)}$; and the bond form $Y[\mathsf w]$ of one direction $\mu$, with $\mathsf w(\langle x,x+a_Ke_\mu\rangle)=\psi^{(K)}(x)$ and $\mathsf w=0$ at the bonds of the two other directions. Let $\mathcal E_1,\dots,\mathcal E_n$ be cell insertions, of any kinds, whose functions $\psi_a$ are real and vanish off closed sets at pairwise sup distances $\ge\mathsf d$, and $\kappa_K=\langle\mathcal E_1;\dots;\mathcal E_n\rangle_{K,\ell,L^Kb}$.
\begin{enumerate}
\item[(a)] There are numbers $C$, $p$, formed from the constants, $n$ and $b$, and a cutoff $K_1$, formed besides from $g$, $\ell$ and $\mathsf d$, such that for $K\ge K_1$
\[|\kappa_{K+1}-\kappa_K|\le C\,\ell^{3(n+1)}K^pL^{-\theta_nK}\prod_a\sup|\psi_a|,\qquad\theta_n=\frac1{1260\,n(8n+3)} .\]
So $\kappa_K$ converges along the full sequence $K\to\infty$, and the limit is $n$-linear and symmetric. The same holds, with $\sup|\psi_a|+{\rm Lip}_\infty\psi_a$ in the place of $\sup|\psi_a|$, for the forms of the three kinds whose weights are real functions $\psi_a$ with ${\rm Lip}_\infty\psi_a<\infty$ read at the barycentres of the plaquettes and at the midpoints of the bonds, $w(p)=\psi_a(x(p))$, $\mathsf w(\mathfrak b)=\psi_a(x_{\mathfrak b})$, under the same separation of the supports; and for such weights the two limits agree.
\item[(b)] In this clause the separation of the head is replaced by a weaker one. Let $\psi_1,\dots,\psi_n$ be real Lipschitz functions on $T_\ell$, none of them $0$, let $\mathcal E_1=Y[\mathsf w]$ be the bond form of one direction $\mu$ with $\mathsf w(\mathfrak b)=\psi_1(x_{\mathfrak b})$ at the bonds $\mathfrak b$ of the direction $\mu$ and $\mathsf w=0$ at the others, let $\mathcal E_a=\mathcal O(w_a)$, $w_a(p)=\psi_a(x(p))$, $2\le a\le n$, be action densities, and let $\kappa_K$ be the joint cumulant of these $n$ forms. If the $n$ supports have no common point and their spread is at least $\mathsf d_0>0$ (the definition preceding Lemma~\YMMref{114.6}{805}), then from a cutoff on $|\kappa_{K+1}-\kappa_K|\le C\,K^pL^{-K/120}\prod_a\|\psi_a\|_\sharp$, $\|\psi\|_\sharp=\max\{\sup|\psi|,40M{\rm Lip}\,\psi\}$, with $C$, $p$ and the cutoff formed from the constants, $n$, $\ell$, $b$ and $\mathsf d_0$.
\item[(c)] If all the $\mathcal E_a$ are action densities with real Lipschitz weights read at the barycentres, $\mathcal E_a=\mathcal O(w_a)$, $w_a(p)=\psi_a(x(p))$, and $\kappa_K$ is their joint cumulant, the convergence is given by Theorem~\YMMref{AR}{810}, equation~\YMMref{(2203)}{810}, Corollary~\YMMref{117.1}{812} and Corollary~\YMMref{118.4}{815}: from a cutoff on, $|\kappa_{K+1}-\kappa_K|\le C\,\ell^{3(n+1)}K^{p}L^{-K/10}\prod_a\|\psi_a\|_\sharp$.
\end{enumerate}
Clauses (a) and (b) carry the convergence of Theorem~\YMMref{AR}{810} over to the plaquette form of one plane, to the bond form and to weights that are means over the lattice cubes; clause (c) is that theorem for the action density.
\end{proposition}

\begin{proof}[Proof sketch]
(a) At the two cutoffs $K$ and $K+1$, at the one label $b$ on the one torus, put
\[\sigma=\lfloor K/(1260\,n)\rfloor,\qquad s=\lfloor\sigma/(8n+3)\rfloor;\]
then $\sigma\le\lceil K/1260\rceil$, $\sigma-s\ge(8n+2)s$, $s\ge\theta_nK-2$ and $4n\sigma\le K/315$. Every insertion is the first member of a source under a majorant at both cutoffs (Lemma~\ref{C-355-21}, which holds on every side with (Sd.1), so on every $\ell\in\mathfrak L(g)$ by Proposition~\YMMref{74.1}{615}): under $(\mathsf s_a,S_a)$ with $\mathsf s_a=\sup|\psi_a|$ for a plaquette form or an action density and $\mathsf s_a=59\sup|\psi_a|$ for a bond form, $S_a$ the plaquettes whose corners lie within $\frac32a_K$ of the closed set off which $\psi_a$ vanishes; a source under a majorant is one under every larger majorant, so that these majorants serve, and corners of plaquettes of different $S_a$ are at sup distance $\ge\mathsf d-3a_K\ge\mathsf d/2$ once $6a_K\le\mathsf d$. The two cumulants are compared through three errors, each bounded by a number formed from the constants and $n$ times the factors written.

(i) \emph{One cutoff.} By Theorem~\YMMref{D}{8}, on the small fields, the chain with the sources is run up to the level $K-s$, at which the $n$ insertions have sizes at most $4\ell^3L^{4s}\mathsf s_a$; each insertion is replaced by its \emph{read-out}, its first member read at the background of the block field of the depth $\sigma$, a function of that block field alone. The cumulant of the $n$ insertions and the cumulant of their $n$ read-outs differ by at most
\[\Bigl[2L^{4n}V_\ell\,L^{-(4n+2)s}+n\,\epsilon(\sigma-s)+2e^{3-L^s\mathsf d/(2M)}\Bigr]\prod_a4\ell^3L^{4s}\mathsf s_a .\]
The first term is the large-field term (Proposition~\textup{\YMMref{115.2}{806}}), $V_\ell=4M^{-3}\ell^3e^{\mathsf B(X_g)\ell^3}$ times a number formed from the constants and $n$. The second is the loss of memory of the progression with sources between the depths $s$ and $\sigma$: the sources switched on at the level $K-\sigma$ in the place of the level $0$ change the effective terms of the level $K-s$ by at most $\epsilon(\sigma-s)$ times their sizes, $\epsilon(\nu)=C_{\rm m}(\nu+1)^{2\mathfrak p_0+2}L^{-\nu/2}$ with a number $C_{\rm m}$ of the constants, because the linearised steps contract on the small fields, as the flows of Propositions~\YMMref{114.2}{801}, \YMMref{114.3}{803} forget their start. The third bounds the polymers that join two supports, whose linear size is at least $L^s\mathsf d/(2M)-2$ in the units of the level $K-s$.

(ii) \emph{The tie at the depth $\sigma$.} The read-outs of the two cutoffs are functions of one block field $V$ of the depth $\sigma$, vanishing where $V$ has large fields, of sizes at most $\frac32B_3^2\ell^3L^{4\sigma}\mathsf s_a$; over one such $V$ they differ by at most
\[\tau_a=C_{\rm t}\,\ell^3L^{4\sigma}K^{1/2}L^{-(K-\sigma)/2}\sup|\psi_a|\]
for weights that are means over the lattice cubes, and by the same with $\sup|\psi_a|+L^{-\sigma}{\rm Lip}_\infty\psi_a$ in the place of $\sup|\psi_a|$ for a weight read at the points of the two lattices, $C_{\rm t}$ a number of the constants; so, under one law of $V$, the cumulants of the read-outs of the two cutoffs differ by at most $\sum_a\tau_a\prod_{b\ne a}\frac32B_3^2\ell^3L^{4\sigma}\mathsf s_b$. The tie is a statement of its own, the counterpart of equation~\YMMref{(2156)}{795} for the products of two plaquette loops at a site. Let $(t,\varepsilon_{\rm r},V)$ be data of the comparison (the setting preceding Lemma~\YMMref{113.8}{794}), $t=K-\sigma$, $U^{\rm A}$, $U^{\rm B}$ the backgrounds over $V$ of the bare lattices $\Lambda_{K,\ell}$, $\Lambda_{K+1,\ell}$, and $N_t=\ell^3L^{3\sigma}$ the number of the sites of the level $t$; for two loops $l$, $l'$ at a site $x$ let $Q_x(l,l')=\langle w(W_{x,l}),w(W_{x,l'})\rangle$, $\langle X,Y\rangle=\frac12\operatorname{tr}XY$, $W_{x,l}$ the plaquette loop of $l$ at $x$. Then for $\xi\in[-1,1]^3$ and a complex function $\phi$ on $T_\ell$ with ${\rm Lip}_\infty\phi<\infty$,
\[\Bigl|L^{t+1}\!\!\sum_{x\in\Lambda_{K+1,\ell}}\!\!\phi(x+a_{K+1}\xi)\,Q_x(l,l')(U^{\rm B})-L^t\!\!\sum_{y\in\Lambda_{K,\ell}}\!\!\phi(y+a_K\xi)\,Q_y(l,l')(U^{\rm A})\Bigr|\le C_{\rm t}\,\|\phi\|_{(t)}\,\varepsilon_{\rm r}^2N_t\,t^{1/2}L^{-t/2},\]
$\|\phi\|_{(t)}$ the norm in equation~\YMMref{(2154)}{795}; and for a weight $\mathsf v_{\rm f}$ on the sites of $\Lambda_{K+1,\ell}$ and its block mean $\mathsf v_{\rm c}(y)=L^{-3}\sum_{x\in B(y)}\mathsf v_{\rm f}(x)$ on $\Lambda_{K,\ell}$, $B(y)$ the $L^3$ sites of the block of $y$, the same holds with $\mathsf v_{\rm f}$, $\mathsf v_{\rm c}$ in the place of the two readings of $\phi$ and $\sup|\mathsf v_{\rm f}|$ in the place of $\|\phi\|_{(t)}$. The proof follows that of Proposition~\YMMref{113.9}{795}: in a gauge on a box about a site of the coarse lattice the background is $e^{ia}$ with a small 1-form $a$ with values in the Lie algebra, the plaquette variables $w(W_{x,l})$ of the fine lattice are the curls of $a$ up to the second order, and those of the coarse lattice their averages (one averaging step with values in the Lie algebra); the variance of the curls about their averages, plaquette by plaquette, is bounded by the defect of Jensen's inequality, whose sum over the plaquettes is $2L^{-2}$ times the difference $L\sum_pa_p-\sum_{p'}a_{p'}$ of the actions of the fine background and of its block average, a difference bounded through the comparison of the two background actions, equation~\YMMref{(2155)}{795}, and the minimality of the background of the coarse lattice; and the background of the coarse lattice is compared with the block average of the fine background, a nearly minimal competitor, through its excess of action.

The bond density is a sum of four such products, $q_{\mathfrak b}=\frac14\sum_{\varepsilon,\varepsilon'=\pm}\varepsilon\varepsilon'\,Q_x((\mu,\varepsilon\lambda),(\mu,\varepsilon'\lambda'))$ at $\xi=\frac12e_\mu$, and the plaquette density of one plane is one such product up to $a_p^2$, by $a_p=\frac12Q_x(l,l)+\frac12a_p^2$ at the loop $l$ of the plane, $a_p^2$ being of the fourth order in the field at the backgrounds; the action density is equation~\YMMref{(2156)}{795} itself. With $a_K^{-1}=L^\sigma L^t$ and $\varepsilon_{\rm r}\le1$, $t\le K$, this gives $\tau_a$; for means over the lattice cubes the weights of the two cutoffs are tied by the block mean, $\psi^{(K)}(y)=L^{-3}\sum_{x\in B(y)}\psi^{(K+1)}(x)$, since the cell of $y$ is the union of the cells of the $x\in B(y)$, and the tie asks no regularity.

(iii) \emph{The laws.} The laws of the block field of the depth $\sigma\le\lceil K/1260\rceil$ at the two cutoffs differ, on the bounded gauge-invariant functions, by at most $\mathsf C_{\rm bf}\ell^3K^{\mathsf p_0}L^{-K/10}$ times their sup norm for $K\ge K_{\rm bf}$ (Corollary~\textup{\ref{IV-62-66}}), with the numbers $\mathsf C_{\rm bf}$, $K_{\rm bf}$ of Definition~\textup{\ref{IV-62-62}}~(5) at the label $b$ and $\mathsf p_0=2\mathfrak p_0+\frac12$; for the cumulant of the read-outs this costs $\mathsf C_{\rm bf}\ell^3K^{\mathsf p_0}L^{-K/10}\prod_a\frac32B_3^2\ell^3L^{4\sigma}\mathsf s_a$.

\emph{The balance.} In (i), $\sigma-s\ge(8n+2)s$ gives $L^{4ns}\epsilon(\sigma-s)\le C_{\rm m}(\sigma+1)^{2\mathfrak p_0+2}L^{-s}$; and from a cutoff on $L^{4n}V_\ell L^{-s}\le1$ and $2e^{3-L^s\mathsf d/(2M)}L^{(4n+1)s}\le1$; so (i) is at most $(4\ell^3)^n[3+nC_{\rm m}(\sigma+1)^{2\mathfrak p_0+2}]L^{-s}\prod_a\mathsf s_a$. In (ii) and (iii), $4n\sigma\le K/315$ gives $L^{4n\sigma}L^{-(K-\sigma)/2}\le L^{-K/3}$ and $L^{4n\sigma}L^{-K/10}\le L^{-K/11}$. As $L^{-s}\le L^2L^{-\theta_nK}$, $\theta_n<\frac1{11}$, $\sigma+1\le K$ and $\mathsf s_a\le59\sup|\psi_a|$, the three errors give (a) with $p=2\mathfrak p_0+2$; the factor $\ell^{3(n+1)}$ comes from the $n$ sizes and the law, and $C$ is formed from the constants, $n$ and $\mathsf C_{\rm bf}$, so from the constants, $n$ and $b$.

The cutoff $K_1$ is the least from which $K\ge K_{\rm bf}$, $s\ge1$, $L^s\mathsf d\ge6M$, $6a_K\le\mathsf d$, $L^\sigma b\ge\Theta_\star$ (the number of equation~\YMMref{(2148)}{794}), the two inequalities above and the thresholds of the chain of $n$ sources at the depth $s$ hold; through $V_\ell$, the distance and $K_{\rm bf}$ it reads besides $g$, $\ell$ and $\mathsf d$, while $C$ and $p$ read neither. The Lipschitz numbers enter at one place, the tie, with the factor $L^{-\sigma}$. The two readings of a Lipschitz weight, at the points and through the means over the lattice cubes, differ by at most $a_K{\rm Lip}_\infty\psi_a$, the size of the weight of a source; the bound of the moments of first members of sources at one fixed depth then makes the two cumulants at the cutoff $K$ differ by at most a number formed from the constants, $n$, $g$, $\ell$ and $\mathsf d$ times $a_K\sum_a{\rm Lip}_\infty\psi_a\prod_{b\ne a}\sup|\psi_b|$, which tends to $0$; so the two limits agree.

Since $\sum_KK^pL^{-\theta_nK}<\infty$, $\kappa_K$ converges along the full sequence; multilinearity and symmetry pass to the limit.

(b) Only the bond form needs a tie: the cumulant is the mixed derivative at $0$, in the parameters of the action densities, of the mean of the bond form under the measure with these densities as sources; along the chain the derivative of the bond form obeys a linear recursion, whose sizes are those of a member of the level $j$ and of what a step creates, so that no disc in a parameter of the bond form is needed; the polymers that meet all the $n$ supports, which make the contact terms, have at a depth $u$ a linear size of at least $\mathsf d_0L^u/M-2$ in the units of their level, as in the polymer estimate in Theorem~\YMMref{AR}{810}, so that the spread, and no pairwise distance, enters; and the comparison of the bond forms of the two backgrounds over one block field, the tie of (ii), gives the rate $L^{-K/120}$. (c) is Theorem~\YMMref{AR}{810}, equation~\YMMref{(2203)}{810} with Corollary~\YMMref{117.1}{812} and Corollary~\YMMref{118.4}{815}.
\end{proof}

\begin{lemma}[A board for $n$ small sets]\label{C-355-22}
Let $n\ge1$, $\mathsf l>0$, $0<h<\mathsf l/(2n)$ with $\mathsf l\in h\mathbb Z$, and let $S_1,\dots,S_n\subset\mathbb R^3$ be non-empty sets of sup diameter $\le\mathsf l/(4n)$. There is $o\in h\mathbb Z^3$ such that every point of every $S_a$ is at distance $>\mathsf l/(8n)$ from every plane $\{x_\mu=o_\mu+\mathsf l(t'+\frac12)\}$, $\mu\in\{1,2,3\}$, $t'\in\mathbb Z$; so each $S_a$ lies in the interior of one cube $o+\mathsf l\mathbf k_a+[-\frac{\mathsf l}2,\frac{\mathsf l}2]^3$.
\end{lemma}

\begin{proof}
Fix $\mu$. The closed $\frac{\mathsf l}{8n}$-neighbourhoods of the projections of the $S_a$ on the $x_\mu$-axis have images on the circle $\mathbb R/\mathsf l\mathbb Z$ that are $n$ closed arcs of total length $\le\frac{\mathsf l}2$. Their complement has measure $\ge\frac{\mathsf l}2$ and at most $n$ components, so one of them is an open arc of length $\ge\frac{\mathsf l}{2n}>h$, and it contains a point of $\frac{\mathsf l}2+h\mathbb Z$; let $o_\mu+\frac{\mathsf l}2$ be such a point. The box of the neighbourhoods is connected and meets no wall plane.
\end{proof}

\begin{lemma}[Dilation of the coupling; refinement]\label{C-355-23}
\begin{enumerate}
\item[(a)] For $0<g\le\gamma$ and $0<\delta\le\frac14$ let $z(g)$ be the least integer $z\ge1$ with $\mathsf G^-_z:=L^{z-1}(g^{-2}-\bar B)-B_\sharp\ge(\gamma^{\rm sd}_W)^{-2}$, and $i_A(g,\delta)=z(g)+\tilde\imath$, $\tilde\imath$ the least integer $\ge1$ with $L^{\tilde\imath-1}(\mathsf G^-_{z(g)}-\bar B)-B_\sharp\ge\mathsf x_E(\delta)$, $\mathsf x_E(\delta)$ the number of Lemma~\ref{C-355-18}~(d), formed from the constants and $\delta$ alone. For a first-passage datum $(K,\beta)$ at $g$ with $K\ge z+1$, $z=z(g)$, the stopped coupling $\tilde g=g_{(z)}$ satisfies $\tilde g\le\gamma^{\rm sd}_W$; $(K-z,\beta)$ is a first-passage datum at $\tilde g$ with label $L^zb$; $L^z\ell\in\mathfrak L(\tilde g)$ for $\ell\in\mathfrak L(g)$; $X_{\tilde g}+\bar B\le L^z(X_g+\bar B)$, $G_{\tilde g}\le G_g+\frac92z\log L+\frac18$; and for $i\ge z+1$, $\tilde\imath=i-z$: $\mathsf Y_{\tilde\imath}[\tilde g]\le\mathsf Y_i[g]$, $C_E[\tilde g]\le e^{1/8}L^{11z/2}C_E[g]$, $K_E(\tilde\imath)[\tilde g]\le K_E(i)[g]-z$ and $i_E(\tilde g,\delta)\le i_A(g,\delta)-z$, the bracket naming the coupling at which a number is formed.
\item[(b)] Let $(K_j,\beta_j,\nu_j)$ be first-passage data at $g$ with periodic Gibbs states and labels converging to $b_*$, $z\ge1$, $\Lambda=L^z$ and $g'=(2L^{z-1}b_*)^{-1/2}$. Then $g'<g$, and along a subsequence $(K_j-z,\beta_j)$ are first-passage data at $g'$, with the dilated states, whose labels converge to $\Lambda b_*$. For a weight $H$ on $\mathbb R^3$ put $H^{(z)}(x)=H(\Lambda^{-1}x)$, for a tuple $X=(H_1,\dots,H_n)$ put $X^{(z)}=(H_1^{(z)},\dots,H_n^{(z)})$, continued linearly to the combinations of tuples, and for a vector field $\eta$ put $\eta^{[z]}(x)=\Lambda\,\eta(\Lambda^{-1}x)$. If they define the tensor limit $\sigma'$, then $\sigma$ exists and $\sigma(X)=\Lambda^n\sigma'(X^{(z)})$ for $X$ with $n$ entries.

The dilation $X\mapsto X^{(z)}$ of the tuples is a bijection, multiplicative for the concatenation, commutes with $\mathrm O(3)$, carries $\tau_v$ to $\tau_{\Lambda v}$ and $\theta_{n,c}$ to $\theta_{n,\Lambda c}$, and carries $H^{\rm cons}[\eta]$ to $H^{\rm cons}[\eta^{[z]}]$, that is $(H^{\rm cons}[\eta])^{(z)}=H^{\rm cons}[\eta^{[z]}]$, $\eta\mapsto\eta^{[z]}$ a bijection of $C^2_c$. If $\sigma'$ satisfies $(\mathrm{Rec4})_V$ with $(C_V,p_V,\frac12)$, then $\sigma$ satisfies it with $(3^{p_V}\Lambda^4C_V,p_V,\frac12)$; if $\sigma'$ satisfies it with $(C_V,p_V,c_0)$, then $\sigma$ satisfies it with $(\Lambda^4C_V,p_V,\Lambda^{-1}c_0)$. All this holds as well with $g'$ replaced by a reference coupling $g''$ such that $\Lambda b_*$ lies in the interior of $I(g'')$; with $g_1'^{-2}=\Lambda g^{-2}/2$ and $g_2'^{-2}=L\Lambda g^{-2}/4$, every label $\Lambda b$, $b\in[g^{-2},L(g^{-2}+\bar B)]$, lies in the interior of $I(g'_1)$ or of $I(g'_2)$.
\item[(c)] Let $(i,o)$ be a board of $\mathbb L_K$, $i+2\le i'\le K-1$, $\mathfrak s$ a law on $\mathbb L_K$ invariant under the lattice isometries, $E\ge0$, and let the eight boards $(i',o+wm)$, $m\in\{0,1\}^3$, $w=(a_{i'}-a_K)/2$, have the property: for every cell of them and every real $h'$ vanishing at sup distance $<a_{i'}/8$ from its complement, $\|\mathcal T_K(h')-\mathfrak s(\mathcal T_K(h'))\|_{{\rm ch},\ell}\le E|h'|_\infty$. Then for every cell of $(i,o)$ and every real $h$ vanishing off it and within sup distance $a_i/32$ of its boundary, $\|\mathcal T_K(h)-\mathfrak s(\mathcal T_K(h))\|^{(i,o)}_{{\rm ch},\ell}\le8(L^{i'-i}+1)^3E|h|_\infty$.
\end{enumerate}
\end{lemma}

\begin{proof}[Proof sketch]
(a) Corollary~\YMMref{74.3}{616} and Lemma~\ref{lem:setup-stopped-arithmetic} at the scale $z$, and the monotony of the explicit formulas of Theorem~\ref{C-355-7}~(a) in the coupling: $\mathsf Y$ and $C_E$ grow with $X_g$, and $G_{\tilde g}-G_g$ is bounded by \eqref{C-355-2}. In detail: the flow bound \YMMref{(1530)}{615} of Proposition~\YMMref{74.1}{615} at the indices $K-z\le K$ reads $|x_K-L^{-z}x_{K-z}|\le\bar B(1-L^{-z})$. With $g^{-2}<x_K\le X_g$ it gives $x_{K-z}>L^z(g^{-2}-\bar B)$, so that $\tilde g^{-2}=L^{-1}x_{K-z}-B_\sharp>\mathsf G^-_z\ge(\gamma^{\rm sd}_W)^{-2}$, and $x_{K-z}+\bar B\le L^z(x_K+\bar B)\le L^z(X_g+\bar B)$; and $X_{\tilde g}=x_{K-z}$ by Corollary~\YMMref{74.3}{616} and Lemma~\ref{lem:setup-stopped-arithmetic}, which give the datum at $\tilde g$, its label and its sides.

So $X_{\tilde g}+\bar B\le L^z(X_g+\bar B)$; hence $\mathsf Y_{\tilde\imath}[\tilde g]\le\mathsf Y_i[g]$, and $C_E[\tilde g]/C_E[g]=e^{G_{\tilde g}-G_g}(X_{\tilde g}+\bar B)/(X_g+\bar B)\le e^{1/8}L^{9z/2}L^z$. The depth $s(\delta)$ does not decrease with the inverse square of the coupling at which it is formed, because $C_W(s)$ grows with it; as $L^{\tilde\imath-1}(X_{\tilde g}+\bar B)-B_\sharp\le L^{i-1}(X_g+\bar B)-B_\sharp$, this gives $s_E(\tilde\imath)[\tilde g]\le s_E(i)[g]$ and the bound of $K_E$. As $\mathsf x_E(\delta)$ is independent of the coupling, $\tilde g^{-2}>\mathsf G^-_z$ gives $L^{j-1}(\tilde g^{-2}-\bar B)-B_\sharp>L^{j-1}(\mathsf G^-_z-\bar B)-B_\sharp$ for every $j\ge1$, whence $i_E(\tilde g,\delta)\le i_A(g,\delta)-z$. (b) The exact dilation of the lattice model (Lemma~\ref{lem:lattice-law-dilation}) gives $\mathcal T_K(H)$ as $\Lambda\,\mathcal T_{K-z}(H^{(z)})$ after transport, for the dilated weight $H^{(z)}$, as in Lemma~\ref{C-355-18}~(a), where $L^iH(L^{-i}\cdot)=\Lambda H^{(z)}$ at $i=z$, and carries periodic Gibbs states to periodic Gibbs states; the labels are placed by the first-passage windows of Corollary~\YMMref{74.2}{616}, at $g'$ and at every $g''$ of the kind named, for the labels $\Lambda b_j$ lie in the interior of $I(g'')$ from an index on; the formula of $q_{\mathfrak b}$ shows that the dilation commutes with the action of $\mathrm O(3)$ on the labels, and the identity for $H^{\rm cons}$ is the chain rule: $\nabla(\eta^{[z]})(x)=\Lambda\,\Lambda^{-1}(\nabla\eta)(\Lambda^{-1}x)=(\nabla\eta)^{(z)}(x)$, so $\nabla\cdot\eta^{[z]}=(\nabla\cdot\eta)^{(z)}$ and $H^{\rm cons}[\eta^{[z]}]=(H^{\rm cons}[\eta])^{(z)}$.

For the two couplings: with $y=\Lambda g^{-2}$ the interiors of $I(g'_1)$, $I(g'_2)$ are $(y/2+\bar B,Ly/2+(L-2)\bar B)$ and $(Ly/4+\bar B,L^2y/4+(L-2)\bar B)$, which overlap and cover $[y,Ly+L\Lambda\bar B]$, as $\bar B\le\frac{13}{24}g^{-2}$. For $(\mathrm{Rec4})_V$: the map $F\mapsto\Lambda^{|F|}F^{(z)}$, $|F|$ the number of entries, is multiplicative, commutes with $\theta$ and the conjugation, carries $\tau_{-\eta e_3}$ to $\tau_{-\Lambda\eta e_3}$ and the classes of $\sigma$ isometrically onto those of $\sigma'$, and $\|\Lambda h^{(z)}\|_{L^1}=\Lambda^4\|h\|_{L^1}$. So for $\eta\le\Lambda^{-1}$ the bound of $\sigma'$ at the distance $\Lambda\eta\le1$ gives that of $\sigma$ with $\Lambda^{4-p_V}C_V$, in either form; for $\Lambda^{-1}<\eta\le1$ and $h$ carried by $\{|x_3|<\Lambda^{-1}c_0\eta\}$ the bound of $\sigma'$ at the distance $1$ applies, and the semigroup $e^{-tH}$ is a contraction: the second triple.

For the first, let $0<\eta\le1$, $h$ be carried by $\{|x_3|<\eta/2\}$ and $\eta_1=\min\{\Lambda^{-1},\eta/3\}$; cut $h$ in $x_3$ by a partition of unity into pieces $h_k$ carried by the slabs of width $\eta_1$ about the heights $y_k=k\eta_1/2$, and translate each piece, with the two tuples beside it, to the height $0$, $\sigma$ being invariant under the translations. The translated tuples lie at heights $>\eta-|y_k|>\eta_1$; the bound at the distance $\eta_1$, the contraction of the semigroup, $\sum_k\|h_k\|_{L^1}=\|h\|_{L^1}$ and $\eta_1^{-1}\le3\Lambda/\eta$ give the number $3^{p_V}\Lambda^4C_V$. (c) Let $\xi_\mu(x)$ be the coordinate of $x$ in the direction $\mu$ measured from a wall plane of the coarse board $(i,o)$; put $m_\mu(x)=0$ if ${\rm dist}(\xi_\mu(x),a_{i'}\mathbb Z)\ge a_{i'}/8$ and $m_\mu(x)=1$ otherwise, and $E_m=\{x:m(x)=m\}$, $m\in\{0,1\}^3$: eight Borel sets that part $\mathbb R^3$.

As $a_i$ is an odd multiple of $a_{i'}$, a mirror map of the coarse board changes each $\xi_\mu$ modulo $a_{i'}$ at most by its sign, so it keeps every $E_m$. In the coordinate $\xi_\mu$ the walls of the fine board $(i',o+wm)$ are the points of $wm_\mu+a_{i'}\mathbb Z$; so the points of $E_m$, and with them their coarse mirror images, lie at sup distance $\ge a_{i'}/8$ from the walls of this board, and the coarse mirror image of a piece $h\mathbf 1_{E_m\cap\square'}$, $\square'$ a cell of the board $(i',o+wm)$, lies in one cell of the same board. Copies of a piece in different coarse cells are at sup distance $\ge a_i/16>a_{i'}$, hence in fine cells that are different modulo the side. The pieces $h\mathbf 1_{E_m\cap\square'}$ that meet the coarse cell, at most $8(L^{i'-i}+1)^3$ of them, add up to $h$; the chessboard estimate on the fine board, for a subset of a period box (Corollary~\ref{IV-18-18}~(b)), bounds the coarse seminorm of each piece by the fine seminorms of its copies, each at most $E|h|_\infty$, and the seminorm is subadditive.
\end{proof}

\begin{lemma}[From the sup norm on a cube to $L^1$; for data]\label{C-355-24}
Let $\mathcal H$ be a Hilbert space, $\mathcal P$ a projection-valued measure on $[0,\infty)\times\mathbb R^2$ carried by the square cone, $e^{-tH}\mathcal U(v)=\int e^{-tE+iv\cdot q}d\mathcal P$. Let $W$ be a complex normed space of finite dimension, $Q=\{|x-z|_\infty\le\rho\}$, $\mathsf b$ a linear map from the continuous $W$-valued functions carried by $Q$ to the bounded operators with $\|\mathsf b(g)\|\le\mathsf c\sup|g|$, $\varepsilon\ge12\rho$, and $\mathsf B^\varepsilon(g)=e^{-\varepsilon H}\mathsf b(g)e^{-\varepsilon H}$. Suppose that $\mathsf B^\varepsilon(\tau_ug)=e^{-(\varepsilon+u_3)H}\mathcal U(\vec u)\mathsf b(g)\mathcal U(-\vec u)e^{-(\varepsilon-u_3)H}$ whenever $|u|_\infty\le2\rho$ and $g$, $\tau_ug$ are carried by $Q$. Then $\|\mathsf B^\varepsilon(h)\|\le4\mathsf c\rho^{-3}\|h\|_{L^1}$ for every continuous $h$ carried by $Q$.
\end{lemma}

\begin{proof}
Let $\chi$ be the product of the tents $\rho^{-1}(1-|t|/\rho)_+$ about the coordinates of $z$; then $\int\chi=1$, $\sup\chi=\rho^{-3}$, and with $g_u=\chi\,h(\cdot+u)$,
\[
\begin{aligned}
h&=\int\tau_ug_u\,du,\\
\langle\psi',\mathsf B^\varepsilon(h)\psi\rangle
&=\int_{|u|_\infty\le2\rho}\langle\psi'_u,\mathsf b(g_u)\psi_u\rangle\,du,
\end{aligned}
\]
$\psi_u=\mathcal U(-\vec u)e^{-(\varepsilon-u_3)H}\psi$, $\psi'_u=\mathcal U(-\vec u)e^{-(\varepsilon+u_3)H}\psi'$. On the square cone $|q_k|\le E$, so 
\[
\|P_1^{\alpha_1}P_2^{\alpha_2}H^{\alpha_3}e^{-\varepsilon H}\|\le|\alpha|!\,\varepsilon^{-|\alpha|};
\]
 this is the only use of the cone. Expanding $\psi_u$, $\psi'_u$ in power series in $u$ (convergent since $|u|_1\le6\rho\le\varepsilon/2$), exchanging sum and integral, and using 
\[
\sup|\int u^\gamma g_u\,du|\le\rho^{-3}(2\rho)^{|\gamma|}\|h\|_{L^1},
\]
 one obtains the bound 
\[
\mathsf c\rho^{-3}\|h\|_{L^1}(\sum_m(6\rho/\varepsilon)^m)^2\le4\mathsf c\rho^{-3}\|h\|_{L^1}.
\]

\end{proof}

\begin{proposition}[Reconstruction for data with labels]\label{C-355-25}
\begin{enumerate}
\item[(a)] Let $\sigma$ be reconstruction data with labels with the numbers $C_V$, $p=p_V$, $c_0$; (Rec5) is not supposed. Then $\mathcal H(\sigma)$ is separable, and there are a strongly continuous unitary representation $U$ of $\mathbb R^{1+2}\rtimes\mathcal L$ with $U(x,1)=e^{itH}\mathcal U(\vec x)$, operators $\phi(f)$, $f\in\mathcal S_V$, on the span $\mathcal D_0$ of $\Omega$ and of the vectors $\phi(f_1)\cdots\phi(f_n)\Omega$, functions $\mathbf W_n(\mathfrak z;e)$ holomorphic on $\{d_\circ>0\}^{n-1}$ and $n$-linear in $e$, and vectors $\Psi_n(z;e)$ with $[X]=\int\Psi_n(x;h_1(x_1),\dots,h_n(x_n))\,dx$ for the tuples in $\mathbb E_+$ ordered in the third coordinate, such that: (i) the spectral measure of $U(\cdot,1)$ is carried by the closed forward cone and $U(a,\Lambda)\Omega=\Omega$; with (Rec5), the invariant vectors are the multiples of $\Omega$; (ii) $\phi(f)\mathcal D_0\subset\mathcal D_0$, $U(a,\Lambda)\mathcal D_0=\mathcal D_0$, $f\mapsto\langle\psi,\phi(f)\psi'\rangle$ is tempered, and $\langle\psi,\phi(f)\psi'\rangle=\langle\phi(f^*)\psi,\psi'\rangle$; (iii) $U(a,\Lambda)\phi(f)U(a,\Lambda)^{-1}=\phi(D(\Lambda)f(\Lambda^{-1}(\cdot-a)))$ on $\mathcal D_0$, where $D(\Lambda)=\rho_{\mathbb C}(\hat\Lambda)$ and $\rho_{\mathbb C}$ is the holomorphic continuation of $\rho|_{\mathrm{SO}(3)}$ to $\mathrm{SO}(3,\mathbb C)$; at $V=\mathrm{Sym}_3(\mathbb R)$, $D(\Lambda)A=\hat\Lambda A\hat\Lambda^{\rm T}$; (iv) $\phi(f)\phi(f')\psi=\phi(f')\phi(f)\psi$ on $\mathcal D_0$ at spacelike separated supports, for all components; (v) $\mathcal D_0$ is dense; (vi) $|\mathbf W_n(\mathfrak z;e)|\le(C_V\eta^{-p})^n\prod_k|e_k|$ for $e\in V^n$, $d_\circ(\mathfrak z_k)\ge2\eta$, $0<\eta\le1$; $\langle\Omega,\phi(f_1)\cdots\phi(f_n)\Omega\rangle=\lim_{\epsilon\downarrow0}\int\mathbf W_n(\iota(x_2-x_1+i\epsilon\alpha_1),\dots;f_1(x_1),\dots,f_n(x_n))\,dx$ for all $\alpha_k$ in the open forward cone; $\sigma(X)=\int\mathbf W_n(x_2-x_1,\dots;h_1(x_1),\dots)\,dx$ at the tuples ordered in the third coordinate; (vii) there is exactly one function $\mathbf S_n$, holomorphic on the union of the tubes $\{d_\circ(\mathsf R^{-1}(z_{\varsigma(k+1)}-z_{\varsigma(k)}))>0,\ 1\le k<n\}$ over $\mathsf R\in\mathrm{SO}(3)$ and the permutations $\varsigma$, which on each of them is $\mathbf W_n$ written in the frame $\mathsf R$ at the labels $\rho(\mathsf R)^{-1}e_{\varsigma(k)}$; its restriction $\mathsf s_n$ to the real non-coinciding configurations is real-analytic, is the only continuous kernel of $\sigma$, is symmetric in points and labels together, and is invariant under every isometry; (viii) $\Omega$ and the vectors $\Psi_n(x;e)$ at real $x$ with $0<x_{13}<\dots<x_{n3}$ are total, and so are $\Omega$ and the classes of the tuples whose entries lie in disjoint slabs.

Clauses (i)--(vi) hold as well for complex symmetric functionals invariant under the orientation-preserving isometries, with $\overline{\sigma(W)}=\sigma(\theta\bar W)$ and with $(\mathrm{Rec2})_V$--$(\mathrm{Rec4})_V$.
\item[(b)] For data as in (a) the two-fold plateau $m_2(X,Y)=\lim_{|v|\to\infty}\sigma(X\cdot\tau_vY)$ exists, and $\sigma$ has a Hermitian canonical decomposition: a compact metrisable $\mathsf X$, a Borel probability measure $\lambda$ with support $\mathsf X$, and complex symmetric functionals $\sigma_\xi$ with $\overline{\sigma_\xi(X)}=\sigma_\xi(\theta X)$, continuous and injective in $\xi$, such that the iterated plateaux of $\sigma$ (Definition~\ref{C-355-4}~(5)) are $\int\prod_i\sigma_\xi(X^i)\,d\lambda$; it is unique up to exactly one homeomorphism, and its components are real, so that it is a canonical decomposition, if and only if $m_2(F,F)\ge0$ for the real combinations $F$ of tuples.

Every component has $(\mathrm{Rec2})_V$--$(\mathrm{Rec4})_V$ with the numbers of $\sigma$ and (Rec5), is reflection positive at every mirror, and $|\sigma_\xi(X\cdot\tau_vY)-\sigma_\xi(X)\sigma_\xi(Y)|\le K(X,Y)|v|^{-1}$ for $|v|\ge r(X,Y)$, with $K$, $r$ formed from $C_V$, $p$, $c_0$ and the two tuples, the same for all $\xi$. Let $V_0$ be a subrepresentation such that the classes of the tuples with entries in $V_0$ are total, and let $(\mathsf X_0,\lambda_0,(\sigma_{0,\zeta}))$ be the Hermitian canonical decomposition of $\sigma|_{V_0}$, which is reconstruction data with labels in $V_0$ with the numbers of $\sigma$. There is exactly one map $\pi:\mathsf X\to\mathsf X_0$ with $\sigma_\xi|_{V_0}=\sigma_{0,\pi(\xi)}$ for every $\xi$; it is continuous and onto, $\pi_*\lambda=\lambda_0$, and $g\mapsto g\circ\pi$ is a unitary map of $L^2(\mathsf X_0,\lambda_0)$ onto $L^2(\mathsf X,\lambda)$.

So for every tuple $X$ there is a bounded Borel function $g_X$ on $\mathsf X_0$ with $\sigma_\xi(X)=g_X(\pi(\xi))$ for $\lambda$-almost every $\xi$, and the iterated plateaux of $\sigma$ are $\int\prod_ig_{X^i}\,d\lambda_0$: the decomposition of $\sigma|_{V_0}$ bears that of $\sigma$ as measure spaces. The map $\pi$ is a homeomorphism if and only if every $g_X$ agrees $\lambda_0$-almost everywhere with a continuous function, in particular if $\dim\ker H<\infty$.
\end{enumerate}
\end{proposition}

\begin{proof}[Proof sketch]
(a) On the circular cone $|e^{-\mathfrak z_3E+i(\mathfrak z_1q_1+\mathfrak z_2q_2)}|\le e^{-d_\circ(\mathfrak z)E}$; this gives the contractions $\mathcal T(\mathfrak z)$, holomorphic inside the tube. By $(\mathrm{Rec4})_V$ the forms $(F,F')\mapsto\sigma(\theta\bar F\cdot(h)\cdot F')$ are bounded, whence insertion operators $B_\eta(h)$ with $\|B_\eta(h)\|\le C_V\eta^{-p}\|h\|_{L^1}$, covariant under the translations, $B_\eta(\tau_yh')=\mathcal T(y+(\eta-\eta')e_3)B_{\eta'}(h')\mathcal T(-y+(\eta-\eta')e_3)$. Since translations act on them Lipschitz-continuously (on the cone each derivative of $e^{-(\eta/2+y_3)E+i\vec y\cdot q}$ is at most $4/(e\eta)$), the limits $\mathsf B(\eta,\eta;e)=\lim_{\epsilon\downarrow0}B_\eta(\delta_\epsilon e)$, $\delta_\epsilon=\epsilon^{-3}\chi(\cdot/\epsilon)$ with $\chi\ge0$ radial of integral $1$, exist in norm, and $\mathsf b(w|e|w')=\mathcal T(w-se_3)\mathsf B(s,s;e)\mathcal T(w'-se_3)$ are holomorphic point operators with $B_\eta(h)=\int\mathsf b(y+\eta e_3|h(y)|-y+\eta e_3)\,dy$.

Ordered tuples are products of insertion operators and semigroup factors; this gives $\mathbf W_n$, $\Psi_n$ and the integral formulas. Any $n$ distinct points are ordered in some frame $\mathsf R\in\mathrm{SO}(3)$ whose cap $\{\mathsf Re_3\cdot\hat v>\frac12\}$ contains a given direction (the zones $\{|u\cdot(x_a-x_b)|<\gamma_N|x_a-x_b|\}$, $\gamma_N=1/(2N^2)$ with $N\ge n$, cover less than a quarter of the sphere); the pieces in two frames agree at the real points of the intersection of two convex tubes, hence on it, and they form one function $\mathbf S_n$. Totality: the map $y\mapsto\Phi(y)$, the limit of the classes of tuples of bumps at $y$, has the real-analytic Gram kernel $\mathsf s_{2n}(\theta y,y')$, so its matrix elements are real-analytic; for a vector orthogonal to the $\Psi_n$ they vanish on the ordered configurations, hence on the connected set of the non-coinciding ones.

The fields are defined by the forms $q_0$ with $\mathsf A_l(e)=\Gamma(l)^{-2}\int\!\!\int(ss')^{l-1}e^{-s-s'}\mathsf B(s,s';e)\,ds\,ds'$, $\mathsf B(s,s';e)=\mathsf b(se_3|e|s'e_3)$, convergent for $l>p$; then $\mathcal T(w)\phi(f)\mathcal T(w')=\int\mathsf b(w+\iota x|f(x)|w'-\iota x)\,dx$ for $w,w'$ in the open tube, and $\|(1+H)^j\phi(f)\psi\|\le c_j(f)\|(1+H)^{2l+j}\psi\|$. \emph{Boundary values} (vi): with $c_k=\iota(\frac i2\epsilon\alpha_k)$, $\alpha_k$ in the open forward cone, the contractions $\mathcal T(2c_k)$ commute with $H$ and tend strongly to $1$ as $\epsilon\downarrow0$; inserted between the fields, they turn $\langle\Omega,\phi(f_1)\cdots\phi(f_n)\Omega\rangle$ into $\int\mathbf W_n(\iota(x_2-x_1+i\epsilon\alpha_1),\dots;f_1(x_1),\dots,f_n(x_n))\,dx$, up to an error that the energy bounds, applied from the right to the left, send to $0$.

\emph{The density of $\mathcal D_0$} (v): let $\psi\perp\mathcal D_0$. For a word $\mathsf w$ in the operators $\phi(f)$ and $\mathcal T(w)$, $\psi\perp\mathsf w\Omega$, by induction on the number of the letters $\mathcal T$: the matrix element is continuous on the closed tube and holomorphic inside in the argument $w$ of the leftmost such letter; it vanishes at $w=\iota a$, $a\in\mathbb R^{1+2}$, because $\mathcal T(\iota a)=U(a)$ moves the rest of the word into a word with translated test functions and one such letter fewer; so it vanishes on the half planes in time by the reflection principle, and on the tube by the identity theorem in the space variables. As $\mathcal T(w)\phi(g_ke)\mathcal T(w')\to\mathsf b(w|e|w')$ in norm for $g_k\ge0$ of integral $1$ whose supports shrink to a point, $\psi\perp\Psi_n(z;e)$, and $\psi=0$ by (viii).

\emph{The representation $D$}: with the covering $\mathsf p:\mathrm{SL}(2,\mathbb C)\to\mathrm{SO}(3,\mathbb C)$, the differential of $\rho\circ\mathsf p$ on $\mathrm{SU}(2)$ continues complex-linearly to $\mathfrak{sl}(2,\mathbb C)$ and integrates to a holomorphic representation of the simply connected $\mathrm{SL}(2,\mathbb C)$, trivial on $\pm1$, hence to a holomorphic representation $\rho_{\mathbb C}$ of $\mathrm{SO}(3,\mathbb C)$ that continues $\rho$; $\Lambda\mapsto\hat\Lambda$ is a homomorphism of $\mathcal L$ into $\mathrm{SO}(3,\mathbb C)$, so $D(\Lambda)=\rho_{\mathbb C}(\hat\Lambda)$ is a representation, and $D(\Lambda)e^*=(D(\Lambda)e)^*$. Every $\Lambda\in\mathcal L$ is $(1\oplus\mathsf R_1)B(s)(1\oplus\mathsf R_2)$ with rotations $\mathsf R_1,\mathsf R_2$ of the space plane and the boost $B(s)$ : a spatial rotation and a boost send the future unit vector $e_0$ to $\Lambda e_0$, and their inverse followed by $\Lambda$ fixes $e_0$ and acts as a rotation of its orthogonal spatial plane; $\widehat{1\oplus\mathsf R}=\mathsf R\oplus1$, and $\widehat{B(s)}$ is the rotation in the plane $(e_1,e_3)$ by the imaginary angle $is$.

\emph{Lorentz covariance}: $\mathbf W_n$ is invariant under the real rotations about $e_3$ and under those in the plane $(e_1,e_3)$, and this last invariance continues holomorphically in the angle to imaginary angles, the boosts, the labels being moved by $e^{is\,d\rho(\omega_1)}$, $\omega_1=E_{13}-E_{31}$, $d\rho$ the differential of $\rho$; no bound on this factor is needed. So $\mathbf W_n(\hat\Lambda\mathfrak z;D(\Lambda)e)=\mathbf W_n(\mathfrak z;e)$, and the boundary values at the cone vectors $\Lambda^{-1}\alpha_k$ give $\langle\Omega,\phi(\Lambda f_1)\cdots\phi(\Lambda f_n)\Omega\rangle=\langle\Omega,\phi(f_1)\cdots\phi(f_n)\Omega\rangle$ with $(\Lambda f)(x)=D(\Lambda)f(\Lambda^{-1}x)$.

As $(\Lambda f)^*=\Lambda(f^*)$, the map $\phi(f_1)\cdots\phi(f_n)\Omega\mapsto\phi(\Lambda f_1)\cdots\phi(\Lambda f_n)\Omega$ is isometric on the dense $\mathcal D_0$ and extends to a unitary operator $U(\Lambda)$; these form a representation with $U(\Lambda)U(a)U(\Lambda)^{-1}=U(\Lambda a)$, whose strong continuity reads of $D$ only that it is locally bounded. With the spectrum of $U(\cdot,1)$, the image of $\mathcal P$, carried by the closed forward cone by $(\mathrm{Rec3})_V$, this gives (i)--(iii). \emph{Locality}: the componentwise argument is that of Proposition~\YMMref{110.2}{777}. Fix the labels and compactly supported scalar tests with two adjacent supports spacelike separated. The height deformation used there stays in tilted tubes with every consecutive margin at least $3\eta$. The Cauchy--Riemann equations and integration by parts move derivatives onto the tests, and Lemma~\YMMref{109.4}{774} gives a finite smeared bound $Q_e(g)$ uniform along the deformation. For $g$ the product test, the integrals for the two adjacent orders differ by at most $2\eta Q_e(D'g)$, where $D'=\frac12(\partial_{t_{k+1}}-\partial_{t_k})$ and $k,k+1$ are the exchanged entries. Their endpoint boundary values therefore agree as $\eta\to0$. Graph-norm continuity extends the equality to Schwartz tests, as in the cited proof.

These are the proofs for scalar data of Propositions~\YMMref{108.3}{771}, \YMMref{109.1}{773}, \YMMref{110.1}{776}, \YMMref{110.2}{777} and Theorem~\textup{\YMMref{AO}{778}} among the construction results, with a label at each entry and with $D$ in the place of the trivial representation; besides them they read the spectral theorem, the identity theorem, the reflection principle in one variable and the covering of $\mathrm{SO}(3,\mathbb C)$ by $\mathrm{SL}(2,\mathbb C)$. They read the reality of $\sigma$ and its invariance under an orientation-reversing isometry at one place only, the adjoint $B_\eta(h)^*=B_\eta(\theta\bar h)$, where $\overline{\sigma(W)}=\sigma(\theta\bar W)$ is all that is used; this gives (i)--(vi) for the complex functionals of the last sentence of (a). (b) For a combination $Y$ of tuples put $\mathsf e(Y)=P_0[\tau_{se_3}Y]$, $P_0$ the projection onto $\ker H$ and $s$ so large that $\tau_{se_3}Y$ lies over $\mathbb E_+$: it does not depend on $s$, it is invariant under the translations, and these vectors are dense in $\ker H$.

In a frame $\mathsf R$, by the cone, $\sigma(A\cdot\tau_wY)\to\langle\mathsf e(\theta\bar A'),\mathsf e(Y')\rangle$, $A'=\mathsf R^{-1}A$, $Y'=\mathsf R^{-1}Y$, uniformly as $|w|\to\infty$ with $\mathsf Re_3\cdot w\ge\frac12|w|$; two frames whose caps meet give the same limit, and the caps cover the sphere; so the two-fold plateau exists, $m_2(A,Y)=\langle\mathsf e(\theta\bar A),\mathsf e(Y)\rangle$. \emph{The rotations are trivial at infinity}: for a unit vector $u$, $\beta_u(F,G)=m_2(\theta_u\bar F,G)$ is a positive Hermitian form; for $u'\perp u$ the involution $\theta_u\theta_{u'}$ acts on its space as a self-adjoint unitary operator $\ge0$, which is $1$; hence $m_2(\theta_uA,Y)$ does not depend on $u$, $m_2(\mathsf rA,Y)=m_2(A,Y)$ and $\mathsf e(\mathsf rY)=\mathsf e(Y)$ for the orientation-preserving isometries $\mathsf r$, and $J\mathsf e(Y)=\mathsf e(\theta Y)$ defines a unitary involution $J$ of $\ker H$ with $m_2(\bar A,Y)=\langle J\mathsf e(A),\mathsf e(Y)\rangle$.

\emph{The operators at infinity}: $M(X)$ is the bounded operator of $\ker H$ with $\langle\mathsf e(Z),M(X)\mathsf e(Y)\rangle=\lim_{|v|\to\infty}m_2(\theta\bar Z,X\cdot\tau_vY)$, the limit of $\sigma(\tau_w\theta\bar Z\cdot X\cdot\tau_vY)$ as $w$ and $v$ go to infinity on opposite sides of $X$; for $X$ ordered in a frame, $M(X)=P_0\mathsf b_X|_{\ker H}$, $\mathsf b_X$ the product of the insertion operators and semigroup factors of the ordered tuple, and a partition of unity reduces every tuple to such ones. Then $M(F)^*=M(\theta\bar F)$, $M(\mathsf rX)=M(X)$ for the orientation-preserving and $M(\mathsf rX)=JM(X)J$ for the other isometries, $M(X\cdot\tau_vY)\to M(X)M(Y)$ strongly as $|v|\to\infty$, and, $\sigma$ being symmetric, the $M(X)$ commute.

They generate a separable commutative $C^*$-algebra; its space of characters is $\mathsf X$, the vacuum gives $\lambda$, whose support is $\mathsf X$ because $\Omega$ is cyclic, hence separating, for the algebra; $\ker H$ is carried unitarily onto $L^2(\mathsf X,\lambda)$, $\sigma_\xi(X)$ is the Gelfand transform of $M(X)$, so that $\overline{\sigma_\xi(X)}=\sigma_\xi(\theta X)$, and the iterated plateaux of $\sigma$ are $\langle\Omega,\prod_iM(X^i)\Omega\rangle$. Uniqueness follows from the density of the polynomials in the $\sigma_\xi(X)$ and their conjugates. \emph{Reality}: $m_2(\bar F,F)=\langle J\mathsf e(F),\mathsf e(F)\rangle$, so $m_2(F,F)\ge0$ for the real combinations $F$ if and only if $J\ge0$, that is $J=1$, that is $M(X)^*=M(X)$ for every tuple, that is every $\sigma_\xi$ is real.

\emph{The components}: an inequality of $\sigma$ at the tuples $F\cdot\tau_{tu}Y$, the spectator $Y$ sent to infinity along $u$, becomes an inequality between elements of the algebra on a dense set of $\ker H$, and, the Gelfand transforms being continuous and $\operatorname{supp}\lambda=\mathsf X$, an inequality at every $\xi$ (test it with $\lambda(O)^{-1/2}\mathbf 1_O$ for small neighbourhoods $O$ of $\xi$); so every $\sigma_\xi$ is reflection positive at every mirror, has $(\mathrm{Rec4})_V$ with the numbers of $\sigma$, and the translations act on its space by contractions and unitary operators, which gives $(\mathrm{Rec2})_V$, one countable family of tuples making all these spaces separable.

The spectral measures disintegrate, $\mu_F=\int\mu^\xi_F\,d\lambda(\xi)$ for $F$ over $\mathbb E_+$; so $\mu^\xi_F(\{E<|q|\})=0$ for $\lambda$-almost every $\xi$, and, this number being lower semicontinuous in $\xi$, for every $\xi$: $(\mathrm{Rec3})_V$. At $\{(0,0)\}$ the disintegration gives $\int\|P^\xi_0[F]_\xi\|^2d\lambda=\|P_0[F]\|^2=\int|\sigma_\xi(F)|^2d\lambda$, while $\|P^\xi_0[F]_\xi\|\ge|\sigma_\xi(F)|$, $P^\xi_0$ the projection onto $\ker H_\xi$: so $\ker H_\xi=\mathbb C\Omega_\xi$ for almost every $\xi$. \emph{The rate}: for data, and for Hermitian data, $|\sigma(X\cdot\tau_vY)-m_2(X,Y)|\le K(X,Y)|v|^{-1}$ for $|v|\ge r(X,Y)$, with $K$, $r$ formed from $C_V$, $p$, $c_0$ and the two tuples alone.

In each frame of a finite family whose caps cover the sphere and in which $X$ and $Y$ are ordered, the difference $f(v)$ is the integral of $e^{-(v'_3-c)E+i\vec v'\cdot q}$, $v'=\mathsf R^{-1}v$, against a complex measure off $(0,0)$ whose variation is at most the product of the numbers $C_V\eta^{-p}\|h\|_{L^1}$ of the entries of $X$ and $Y$. The image $\kappa_v$ of this measure, with the exponential, under $(E,q)\mapsto s=E^2-|q|^2$ has the moments $\Delta^kf(v)$, which are independent of the frame. Its damping gives $\int e^{a\sqrt s}\,d|\kappa_v|(s)<\infty$ for every $a>0$ below the positive time margin, since $0\le\sqrt s\le E$. This square-root exponential condition determines complex measures within this class by their moments. Indeed, the real and imaginary parts of the difference of two such measures have zero moments, including order zero. Their positive and negative parts have equal masses and moments and inherit the same integrability condition; after normalization when the common mass is nonzero, Hardy's criterion applies to them \cite[Theorem~1]{StoyanovLin2013}. Thus $\kappa_v$ is independent of the frame, and its masses $f_{\mathsf B}(v)=\kappa_v(\mathsf B)$ on windows $\mathsf B$ are functions outside a ball $\{|v|<r_2\}$, with $\Delta f_{\mathsf B}(v)=\int_{\mathsf B}s\,\kappa_v(ds)$.

For $\mathsf B=[m_1^2,m_2^2)$, $|(\Delta-m_1^2)\operatorname{Re}(e^{i\varphi}f_{\mathsf B})|\le(m_2^2-m_1^2)\mathsf V_{\mathsf B}$, $\mathsf V_{\mathsf B}$ the variation of the part, and the comparison principle for $\Delta-m_1^2$ outside the ball, with the radial solution $\frac{r_2}re^{-m_1(r-r_2)}$ of the dimension three, gives $|f_{\mathsf B}(v)|\le\mathsf V_{\mathsf B}(\frac{r_2}{|v|}+\frac{m_2^2-m_1^2}{m_1^2})$; the harmonic part at the mass $0$, which tends to $0$ at infinity, is at most $\mathsf V_{\{0\}}\frac{r_2}{|v|}$; the sum over the windows $[(1+\epsilon)^j,(1+\epsilon)^{j+1})$, $j\in\mathbb Z$, with $\epsilon\downarrow0$, gives the rate. One rate serves all the components, which are Hermitian data with the numbers of $\sigma$; with the continuity in $\xi$ and $\operatorname{supp}\lambda=\mathsf X$ the bound by $\sigma_\xi(X)\sigma_\xi(Y)$ holds at every $\xi$, so that $\langle[F]_\xi,e^{-tH_\xi}[F]_\xi\rangle\to|\sigma_\xi(F)|^2$ and $\ker H_\xi=\mathbb C\Omega_\xi$ for every $\xi$.

\emph{The subrepresentation}: the defining limits of the operators at infinity of $\sigma|_{V_0}$ and of $\sigma$ agree at the tuples with entries in $V_0$; when their classes are total, $\sigma|_{V_0}$ and $\sigma$ have one space and one $\ker H$, the operators of these tuples generate a $C^*$-subalgebra with the same unit, for which $\Omega$ is cyclic too, $\pi$ is the restriction of the characters, the only map with $\sigma_\xi|_{V_0}=\sigma_{0,\pi(\xi)}$ because $\zeta\mapsto\sigma_{0,\zeta}$ is injective, $\pi_*\lambda=\lambda_0$ by the vacuum state, $\pi$ is onto because its image is compact of full measure, and $g\mapsto g\circ\pi$ is unitary because both decompositions carry the one space $\ker H$ unitarily onto their $L^2$; $\pi$ is injective if and only if the subalgebra is the whole algebra, by the theorem of Stone and Weierstrass, that is if and only if every $g_X$ agrees almost everywhere with a continuous function; and if $\dim\ker H=N<\infty$, then $\mathsf X$ and $\mathsf X_0$ have $N$ points each.
\end{proof}

\begin{lemma}[The circular cone from five mirrors; for data]\label{C-355-26}
Let $\sigma$ be a real symmetric multilinear functional on labelled tuples, with $\sigma(\varnothing)=1$ and bounds independent of the position, invariant under the translations. For a unit vector $u$ at whose mirror $\{u\cdot x=0\}$ $\sigma$ is reflection positive, let $\mathcal H_u$, $H_u$, $\mathcal P_u$ be formed as in Definition~\ref{C-355-4}~(3) with $u$ in the place of $e_3$.
\begin{enumerate}
\item[(a)] Let $u'$ be a unit vector orthogonal to $u$. If $\sigma$ is reflection positive at the mirrors with the normals $u$ and $(u+u')/\sqrt2$, then $H_u\ge|u'\cdot P_u|$.
\item[(b)] Let $(u,u',u'')$ be an orthonormal basis and $n_\beta=\cos\beta\,u+\sin\beta\,u'$. If $\sigma$ is reflection positive at the four planes through the axis $\mathbb Ru''$ with the normals $n_0,n_{\pi/4},n_{\pi/2},n_{3\pi/4}$ and at one plane with the normal $(n_{\beta_1}+u'')/\sqrt2$, $\beta_1\in\frac\pi4\mathbb Z$, then $\mathcal P_{n_\beta}(\{E<|q|\})=0$ for every $\beta\in\frac\pi4\mathbb Z$.
\item[(c)] \emph{(Eight measures in one plane; analysis alone.)} Let $R>0$, $\kappa>0$, $C<\infty$, $\mathsf V=\{(E,q)\in\mathbb R^2:E\ge|q|\}$, and for each of the eight angles $\beta\in\frac\pi4\{0,1,\dots,7\}$ let $m_\beta$ be a complex Borel measure on $\mathsf V$ with $|m_\beta|(\mathsf V)\le C$; for $\rho>0$ and complex $c$ put
\[f_\beta(c;\rho)=\int_{\mathsf V}\exp\bigl(2RE-\rho[E\cos(c-\beta)-iq\sin(c-\beta)]\bigr)\,dm_\beta(E,q)\]
where the integrand is bounded. Suppose that $f_\beta(t;\rho)=f_{\beta'}(t;\rho)$ for any two of the eight angles $\beta,\beta'$, all $\rho>0$ and all real $t$ with $\rho\cos(t-\beta)>2R$ and $\rho\cos(t-\beta')>2R$; that $\mu=e^{2RE}m_0$ is a finite positive measure; and that for one of the eight angles, $\beta_1$, the measure $m_{\beta_1}$ is carried by $\{E\ge\kappa\}$. Then $\mu(\{(E,q)\in\mathsf V:E^2-q^2<\kappa^2\})=0$.
\end{enumerate}
No invariance under $W_3$, under a reflection or under a rotation is supposed.
\end{lemma}

\begin{proof}[Proof sketch]
(a) is the argument of Proposition~\YMMref{94.4}{722}, applied at the mirrors with the normals $u$ and $(u+u')/\sqrt2$; the labels ride along, the one number $\sigma(\theta_u\bar F\cdot\tau_wF)$ being read in the two pictures. (b) By the symmetry of $\sigma$ (the entries commute) one and the same number $\sigma(A\cdot\tau_xB)$ is a matrix element in every picture whose mirror separates $A$ from $\tau_xB$, also when $x$ has gone around $A$. After a cut $|q\cdot u''|\ge\kappa$ in the momentum normal to the plane of the eight normals $n_\beta$, which is common to the eight pictures, this gives eight complex measures $m_\beta$, $\beta\in\frac\pi4\mathbb Z$, carried by $\mathsf V$ by (a) in each picture, whose Laplace--Fourier transforms $\int\exp(2RE-\rho[E\cos(c-\beta)-iq\sin(c-\beta)])\,dm_\beta$ agree wherever two of them are defined, by the uniqueness of the Fourier transform in the common momentum; the one of the picture $u$ is positive after the weight $e^{2RE}$, and by (a) at the fifth mirror the one of the picture $n_{\beta_1}$ is carried by $\{E\ge|q\cdot u''|\}\subset\{E\ge\kappa\}$.

Clause (c) gives $\mu(\{E^2-q^2<\kappa^2\})=0$ for every $\kappa>0$, and so the circular cone in each picture. (c) Write $c=t+i\gamma$. On $\mathsf V$, $E\cosh\gamma+q\sinh\gamma\ge Ee^{-|\gamma|}$; so on $S_\beta(\rho)=\{\rho\cos(t-\beta)e^{-|\gamma|}>2R\}$ the integrand of $f_\beta$ is bounded by $1$, and $f_\beta(\cdot;\rho)$ is holomorphic there (Morera), bounded by $C$ and of the period $2\pi$. The sets $S_\beta(\rho)$ cover the band $\{|\gamma|<\Gamma(\rho)\}$, $\Gamma(\rho)=\log(\rho\cos\frac\pi8/2R)$, and two of the functions agree on real intervals of the overlaps, hence on the overlaps (identity theorem): they glue to one holomorphic function $f(\cdot;\rho)$ on the band, $|f|\le C$.

Near $\beta_1$, $|f(c;\rho)|\le Ce^{-\kappa\rho_\delta}$ for $|c-\beta_1|\le\delta$, $\rho_\delta=\rho e^{-\delta}\cos\delta-2R$, by the carrier of $m_{\beta_1}$. With $N(c)=\log|2\sin((c-\beta_1)/2)|$, harmonic off $\beta_1+2\pi\mathbb Z$, the function $\log|f|-\log C+\kappa\rho_\delta h$, $h=(\mathsf M-N)/(\mathsf M+\log(2/\delta))$, $\mathsf M=\log(2\sinh(\Gamma'/2))$, $\Gamma'=\Gamma(\rho)-1$, is subharmonic on the band $\{|\gamma|<\Gamma'\}$ without the discs $|c-\beta_1-2\pi k|\le\delta$ and $\le0$ on its boundary; by the maximum principle $|f|\le Ce^{-\kappa\rho_\delta h}$ there. At the points $i\gamma_0$, $f(i\gamma_0;\rho)=\int\exp(-\rho(E\cosh\gamma_0+q\sinh\gamma_0))\,d\mu$ is the Laplace transform of the positive measure $\mu$, and $h(i\gamma_0)\to1$ as $\rho\to\infty$. For fixed $\gamma_0\ne0$, choose $\delta>0$ small enough that $i\gamma_0$ lies outside the excluded discs. First letting $\rho\to\infty$ and then $\delta\downarrow0$ gives $E\cosh\gamma_0+q\sinh\gamma_0\ge\kappa$ almost everywhere. Apply this first on a countable dense subset of $\mathbb R\setminus\{0\}$ and use continuity in $\gamma_0$; the inequality then holds for every real $\gamma_0$ on one set of full measure, and the infimum over $\gamma_0$ of the left member is $0$ if $E=|q|$ and $(E^2-q^2)^{1/2}$ if $E>|q|$.
\end{proof}

\begin{proposition}[Ordered classes and kernels; for data with the square cone]\label{C-355-27}
Let $\sigma$ be a real symmetric multilinear functional on labelled tuples with bounds $C_n(d;R)\prod_a|H_a|_\infty$ independent of the position, invariant under $\mathbb R^3\rtimes W_3$, with $\sigma(\theta\bar F\cdot F)\ge0$ on $\mathbb C\mathfrak S_V(\mathbb E_+)$ and with $\mathcal P$ carried by the square cone.
\begin{enumerate}
\item[(a)] $\mathcal H^{\rm ord}(\sigma)=\mathcal H(\sigma)$; two such functionals that agree on the tuples ordered in the third coordinate are equal.
\item[(b)] If $\sigma^1$ satisfies besides $(\mathrm{Rec4})_V$ with $C_V$, $p$ and $c_0=\frac12$, then the slab operators of tuples with entries side by side, the kernels and the bounds of Theorem~\ref{C-355-9}~(g) hold for $\sigma$, with $\mathsf K_n=(10n)^{\lceil pn\rceil n(n-1)/2}$; $\sigma$ does not charge the configurations with two equal heights.
\item[(c)] If moreover the scalar part of $\sigma$ is scalar reconstruction data, the kernel of one tensor entry among $m$ scalar entries is real-analytic at all non-coinciding points, for every $m$.
\end{enumerate}
\end{proposition}

\begin{proof}[Proof sketch]
(a) By the bounds the class of a tuple depends continuously, in the sup norm, on its entries; so it suffices to show $[X]\in\mathcal H^{\rm ord}(\sigma)$ for $X\in\mathfrak S_V(\mathbb E_+)$ with smooth entries. A partition of unity of a small mesh cuts $X$ into tuples $Y=(g_1,\dots,g_n)$ on \emph{frames}: $g_a$ is carried by a box $Q_a$ with sides at most $2r$, the boxes lie in $\{x_3\ge2t_1\}$ and are pairwise at sup distance $\ge4nr$, for some $r,t_1>0$. A pair $\{a,b\}$ is \emph{bad} if the intervals of the heights of $Q_a$ and $Q_b$ meet; without bad pairs $Y$ is ordered. Induction on the number of the bad pairs. A bad pair $\{a,b\}$ is parted by a vertical cut, a plane $\{x_\mu=c\}$, $\mu\in\{1,2\}$, at a distance $\ge r$ from every box, with $Q_a$ and $Q_b$ on different sides: in the direction $\mu$ in which $Q_a$ and $Q_b$ are $4nr$ apart, the gap between them leaves an interval of length $\ge(4n-2)r$ for $c$, of which the other $n-2$ boxes forbid at most $(4n-8)r$.

Let $Y_{\mathsf L}$, $Y_{\mathsf R}$ be the groups of the entries on the two sides, $Y_{\mathsf R}$ holding $g_b$. The lifted vector $u(\mathsf h)=[Y_{\mathsf L}\cdot\tau_{\mathsf he_3}Y_{\mathsf R}]$, $\mathsf h>-t_1$, lies in $\mathcal H^{\rm ord}(\sigma)$ for large $\mathsf h$ by the induction, as the lifted frame has a bad pair fewer. Its Gram kernel across the diagonal is $\langle u(\mathsf h_0+x-y),u(\mathsf h_0+x+y)\rangle=\sigma(A\cdot\tau_{ye_3}B_{\mathsf h_0+x})$ with $A=\theta Y_{\mathsf L}\cdot Y_{\mathsf L}$ and $B_s=\tau_{-se_3}\theta Y_{\mathsf R}\cdot\tau_{se_3}Y_{\mathsf R}$: as $\theta$ reverses the heights, $y$ moves each group rigidly together with its mirror copy. The plane $\{x_\mu=c\}$ parts $A$ from $B_s$; in the picture of its normal $e_\mu$, which an element of $W_3$ turns into $e_3$, this number is $\langle\Phi_A,e^{-rH}\mathcal U(ye_\mu)\Phi_B\rangle$, the norms of $\Phi_A$, $\Phi_B$ being bounded by the position-free bounds at $4\#\mathsf L$ and $4\#\mathsf R$ entries, and by the square cone $|\partial^j_y\sigma(A\cdot\tau_{ye_3}B_s)|\le(j/er)^j\|\Phi_A\|\|\Phi_B\|$ for every $j$: this is the only use of the cone.

The norm of the Fourier transform of a windowed curve reads the Gram kernel at the frequencies $(p,-p)$ alone, that is across the diagonal; with windows that are $(2q+1)$-fold convolutions of indicator functions, whose derivatives of the orders $j\le2q$ are at most $(24q)^j$ on the unit scale, Lemma~\ref{cov-gram}~(b) gives real-analyticity of $u$ in norm on $(-\frac58t_1,\infty)$. So for $\psi\perp\mathcal H^{\rm ord}(\sigma)$ the real-analytic function $\mathsf h\mapsto\langle\psi,u(\mathsf h)\rangle$ vanishes for large $\mathsf h$, hence at $\mathsf h=0$: $\psi\perp[Y]$. For the second clause let $\sigma'$ be the other functional. Then $\frac12(\sigma+\sigma')$ satisfies what the Proposition supposes: its classes are carried isometrically into $\mathcal H(\sigma)\oplus\mathcal H(\sigma')$ by $[F]\mapsto2^{-1/2}([F]_\sigma,[F]_{\sigma'})$, which intertwines the translations, so that its spectral measure is carried by the square cone.

A combination $G$ of ordered tuples whose class is close to that of $X$ for $\frac12(\sigma+\sigma')$ is close to it in $\mathcal H(\sigma)$ and in $\mathcal H(\sigma')$ at once; $|\sigma(F)|=|\langle\Omega,[F]\rangle|\le\|[F]\|$ for $F$ over $\mathbb E_+$, and $\sigma(G)=\sigma'(G)$; so $\sigma(X)=\sigma'(X)$ for $X$ over $\mathbb E_+$, and for every $X$ by a translation. (b) For points $z_a=(\vec x_a,t_a)$, $t_a$ complex, whose real heights lie in $(0,l)$ and are pairwise different, $\mathsf G_l(z;e)$ is the product, in the order of increasing height, of the point operators $\mathsf b(\cdot|e_a|\cdot)$ with the semigroup factors between them, the first factor taking the lowest point and the last factor $le_3$ minus the highest point; the point operators are those that $(\mathrm{Rec4})_V$ gives on the square tube $\{d_\square>0\}$, as in the sketch of Proposition~\ref{C-355-25}~(a) with the square cone in the place of the circular one, with the bound of Theorem~\ref{C-355-9}~(d)~(ii); so $\mathsf G_l$ is holomorphic on the union of the tubes of the orders in height and bounded by a power of the least gap.

Two fields side by side, at the horizontal sup distance $\gamma=|\vec x_a-\vec x_b|_\infty>0$, are one holomorphic function of their two heights across the wall $\operatorname{Re}t_a=\operatorname{Re}t_b$, for $|\operatorname{Im}(t_a-t_b)|<\gamma$. At a point of the wall, $\operatorname{Re}t_a=\operatorname{Re}t_b=T_0$, put $\zeta=t_a-t_b$, $\tau_0=\operatorname{Im}\zeta\in(-\gamma,\gamma)$, $\kappa=\frac12\min\{1,T_0,l-T_0,\gamma-|\tau_0|\}$, $m_1$ the least integer $\ge2p$, and $\mu(\zeta)=(\zeta-i\tau_0)^2+\kappa^2$. On the circle $|\zeta-i\tau_0|=\kappa$ one has $|\mu(\zeta)|=2\kappa|\operatorname{Re}\zeta|$; so $\mu(\zeta)^{-m_1}$ times the Cauchy integral of $\mu^{m_1}\mathsf G_l$ over this circle is formed, the bound $C|\operatorname{Re}\zeta|^{-2p}$ of $\mathsf G_l$ off the wall being absorbed, and it is holomorphic inside.

It is the continuation: its matrix elements between $\mathcal U(v)[F']$ and $\mathcal U(w)[F]$ are holomorphic in $v,w$ for $|\operatorname{Im}v_1|+|\operatorname{Im}v_2|<\frac\kappa2$, and so are those of $\mathsf G_l$, by the square cone; for real $v,w$ far out along $e_k$ in opposite senses, $k\in\{1,2\}$ the direction with $|x_{a,k}-x_{b,k}|=\gamma$, the spectators and the two fields are ordered along $e_k$, and in the frame of $W_3$ that turns $e_k$ into $e_3$ the kernel is holomorphic across $\operatorname{Re}\zeta=0$ for $|\operatorname{Im}\zeta|<\gamma$, so that the two agree there; the identity theorem in $(v,w)$ carries the agreement to $v=w=0$. \emph{The bound by the least distance.} Let $0<d\le1$ with $d\le\operatorname{Re}t_a\le l-d$ for all $a$ and, for every pair $\{a,b\}$, $|\operatorname{Re}(t_a-t_b)|\ge d$ or $|\vec x_a-\vec x_b|_\infty-|\operatorname{Im}(t_a-t_b)|\ge d$.

On the disc $t_a(\lambda)=t_a+a\lambda$, $|\lambda|<d/(3n)$, on which the $a$-th point moves with the velocity $a$, the walls of the pairs $\{a,b\}$ with $|\operatorname{Re}(t_a-t_b)|<d$ are the lines $\operatorname{Re}\lambda=-r_{ab}$, $r_{ab}=\operatorname{Re}(t_a-t_b)/(a-b)$, which may be supposed pairwise different. With $d_n=d/(4n)$ and $m_2=\lceil pn\rceil$ the multiplier $Q(\lambda)=\prod_{\{a,b\}}(\lambda^2+2r_{ab}\lambda+d_n^2)^{m_2}$, the product over these pairs, has the modulus $\prod(2d_n|\operatorname{Re}\lambda+r_{ab}|)^{m_2}$ on the circle $|\lambda|=d_n$ and $d_n^{2m_2P}$ at $\lambda=0$, $P$ the number of these pairs; $Q\mathsf G_l$ continues across the walls by the clause on two fields, and the maximum principle at $\lambda=0$ gives the factor $\mathsf K_n$.

\emph{The slab operators.} Every family of small sets at mutual distances larger than $n$ times their diameters can be cut by a coordinate plane into two non-empty groups; with this, an induction on $n$ shows $[M\cdot\tau_{le_3}F]=\mathsf I_l(M)[F]$, $\mathsf I_l(M)$ the integral of $\mathsf G_l$ over the configurations with different heights, whose complement is a null set. For a cut by a plane $\{x_3=c_1\}$ the two groups give a product of two operators of the induction. For a cut by $\{x_k=c_1\}$, $k\in\{1,2\}$, let $\eta>0$ with $\operatorname{supp}M\subset\{\eta<x_3<l-\eta\}$: the difference of $\sigma(\theta\overline{\tau_{(v,0)}F'}\cdot M\cdot\tau_{le_3}\tau_{(w,0)}F)$ and $\langle\mathcal U(v)[F'],\mathsf I_l(M)\mathcal U(w)[F]\rangle$ is holomorphic in $v$ for $|\operatorname{Im}v_1|+|\operatorname{Im}v_2|<\eta$ at real $w$, and in $w$ at real $v$, by the square cone; it vanishes when $v$ and $w$ are far out along $e_k$ in opposite senses, where the groups and the spectators are ordered along $e_k$ and the induction applies, in the frame that turns $e_k$ into $e_3$, to the turned groups; the identity theorem, first in $v$ and then in $w$, gives it at $v=w=0$. (c) Fix non-coinciding points $x^0,y^0_1,\dots,y^0_m$.

For each $k$ choose a coordinate plane $\{x_{\mu_k}=c_k\}$ strictly between $x^0$ and $y^0_k$, in the direction $\mu_k$ in which they are farthest apart, at a positive distance from all the points; an element of $\mathbb R^3\rtimes W_3$ turns it into $\{x_3=0\}$ with $y^0_k$ above, and on that side stand scalar points alone. For entries $h,f_1,\dots,f_m$ carried by small cubes about the points, $W(v)=\sigma(h,\tau_{v_1}f_1\mathbf 1,\dots,\tau_{v_m}f_m\mathbf 1)$ is, in the picture of that plane, the scalar product of the class of the entries below it, the tensor entry among them, with the class $[F_v]$ of the scalar entries above it. The latter is an integral of scalar point vectors, which are holomorphic in their points, in norm and uniformly on compact sets, because their Gram kernel $\mathsf s_{2n}(\theta\eta,\eta')$ is real-analytic (Proposition~\YMMref{108.3}{771}, the scalar part being scalar data) and Cauchy's estimates of a Gram kernel bound the derivatives of the vectors; the former is bounded by the position-free sup bounds at two tensor entries.

So $W$ is holomorphic in each $v_k$ apart, the others being real, with one bound $C\sup|h|\prod_k\sup|f_k|$; by the bound of the Chebyshev coefficients (Bernstein), a continuous function on a cube that is holomorphic in each variable apart, on ellipses about the real intervals, with one bound, is holomorphic in all the variables together on a smaller polyellipse. The translations move the tensor point as well. Cauchy's estimate for the derivatives in a shift, with $\int\omega(u)\tau_uH\,du=\omega*H$, replaces the sup norm of one entry by its $L^1$ norm, entry after entry; so the values of $\sigma$ at bumps of shrinking width about the points converge uniformly on a polydisc about $(x^0,y^0)$, and the limit is a real-analytic kernel there.

These neighbourhoods cover the non-coinciding configurations; a partition of unity gives the formula of the kernel for all tuples, and the bumps its uniqueness.
\end{proof}

\begin{proposition}[Covariance of the tensor limits under $\mathrm O(3)$]\label{C-355-28}
\begin{enumerate}
\item[(a)] Let $0<g\le\gamma^{\rm sd}_W$. Every tensor limit $\sigma$ at $g$ satisfies $\sigma(\mathsf RX)=\sigma(X)$ for all $\mathsf R\in\mathrm O(3)$ and all $X\in\mathfrak S_V(\mathbb R^3)$. By Lemma~\ref{C-355-23}~(b), since the dilations commute with $\mathrm O(3)$, the same holds for the tensor limits at every $0<g\le\gamma$.
\item[(b)] \emph{(For data.)} Let $\sigma$ be a real symmetric multilinear functional on labelled tuples with $\sigma(\varnothing)=1$, $\sigma((H))=0$, bounds independent of the position, invariant under $\mathbb R^3\rtimes W_3$, reflection positive at the coordinate and diagonal mirrors, with $H\ge|P_1|,|P_2|$. Suppose: (i) bounded slab operators $\mathsf B_l(h)$ exist; (ii) for every off-diagonal label $o=e_\mu e_\nu^{\rm T}+e_\nu e_\mu^{\rm T}$, $\mu\ne\nu$, every $0<d<1$ and every real continuous $h$ carried by a ball $B(x_0,d)$ with $x_{0,3}\in[\frac32d,2d]$, $\sigma(\theta(ho),(ho))\le C|h|_\infty^2$; (iii) for the scalar part: there are a real $a<4$ and, for every real $\varphi\in C^1_c(\mathbb R)$, a number $C_\varphi$ with $\sigma(\theta(f\mathbf 1),(f\mathbf 1))\le C_\varphi t^{-a}$ for all $t\in(0,1]$ and all $f(x)=\alpha(x_3,x_1)\varphi(x_2)$, $\alpha$ a real $C^1_c$ function carried by $[t,2t]\times[-t,t]$ with $|\alpha|\le3$ and ${\rm Lip}\,\alpha\le12/t$; (iv) the class of the tuple $(H^{\rm cons}[\eta])$ vanishes in $\mathcal H(\sigma)$ for $\eta$ carried by $\mathbb E_+$; (v) for one $p<\infty$, $\|\mathsf B_{2d}(\tau_{de_3}h)\|\le C\min\{1,d-d'\}^{-p}\|h\|_{L^1}$ for $h$ carried by $\{|x_3|\le d'\}$.

Then $\sigma(\mathsf RX)=\sigma(X)$ for all $\mathsf R\in\mathrm O(3)$ and all $X$, and $\sigma((H^{\rm cons}[\eta])\cdot X)=0$ for every $X$ off $\operatorname{supp}\eta$. In (v) a weaker growth suffices. Write $y=(y^0,y^1)=(x_3,x_1)$ for the points of the plane of $e_3$, $e_1$, $\xi^\pm=\xi^0\pm i\xi^1$ for $\xi\in\mathbb C^2$, $d(\xi)=\min\operatorname{Re}\xi^\pm$, ${\rm Tb}=\{\xi:d(\xi)>0\}$, $\mathsf Z(\kappa)=\int e^{-\kappa^0E+i\kappa^1q_1}\,d\mathcal P$ for $\kappa$ in the closure of ${\rm Tb}$, and $\mathsf R_\zeta$ for the rotation about $e_2$ by the angle $\zeta$, continued to complex $\zeta$ by $(\mathsf R_\zeta\xi)^\pm=e^{\pm i\zeta}\xi^\pm$. Under (i) there is, for every label $\lambda=(\varphi,a)$, $\varphi\in C^1_c(\mathbb R)$ real and $a\in V_{\mathbb C}$, exactly one family of bounded operators $\mathsf B_\lambda(\xi,\xi')$, $(\xi,\xi')\in{\rm Tb}^2$, holomorphic in norm, with $\mathsf Z(\kappa)\mathsf B_\lambda(\xi,\xi')\mathsf Z(\kappa')=\mathsf B_\lambda(\kappa+\xi,\xi'+\kappa')$ and $\mathsf B_l(h)=\int\alpha(y)\,\mathsf B_\lambda(y,le_3-y)\,dy$ for the plates $h(x)=\alpha(x_3,x_1)\varphi(x_2)a$ with $\alpha$ carried by $\{0<y^0<l\}$: the \emph{point plates}.

The \emph{plate vectors} are $\Psi(z;\lambda_1,\dots,\lambda_j)=\mathsf B_{\lambda_1}(z_1,\frac12(z_2-z_1))\cdots\mathsf B_{\lambda_j}(\frac12(z_j-z_{j-1}),\kappa)\Omega$, which do not depend on $\kappa\in{\rm Tb}$. It suffices that for every $j$ there is $A_j\ge0$ with $\sup_{s\in\mathbb R}e^{-A_j|s|}\|\Psi(\mathsf R_{is}z;\lambda_1,\dots,\lambda_j)\|<\infty$ uniformly on the compact sets of real $z$ whose $z_1$ and differences $z_{i+1}-z_i$ lie at angles $<\frac{5\pi}{16}$ from $e_3$, the labels running over a basis of $V_{\mathbb C}$ of eigenvectors of the rotations about $e_2$, with $\liminf_j(A_j-j^2)/j=-\infty$; with the diagonal mirrors, $\liminf_jA_j/j^2<2$ suffices. Of the scalar part only the bound (iii) is used, at one step of the proof; its covariance is a case of the conclusion.

The free model formed by the tensor of free photon fields with $\epsilon\,\delta_{\lambda\lambda'}\partial^2_\lambda\chi$ added, $\chi$ a free massive scalar field and $\epsilon\ne0$, has the other properties and fails exactly (iv).
\end{enumerate}
\end{proposition}

\begin{proof}
In (b), the normalization, multilinearity, bounds, cubic invariance, reflection positivity and square spectral cone are precisely (D1)--(D4) of Definition~\ref{cov-def}. Conditions (i), (ii), (iv) and (v) give (S), (O), (N) and (P), respectively. Theorem~\ref{cov-thm-data} proves the two conclusions. For the alternative growth hypotheses, Proposition~\ref{cov-two}(a) first gives covariance at pairs from (O) and (N), and Proposition~\ref{cov-rot} gives covariance at every order. The final conservation assertion follows from the null-vector identity by Lemma~\ref{cov-slide}(c).

For (a) at $g\le\gamma_W^{\rm sd}$, Theorem~\ref{C-355-9} supplies the slab and integrable insertion bounds. Theorem~\ref{C-355-15}(E4)(i) supplies the scale-independent off-diagonal bound. Theorem~\ref{C-355-11}(c), applied to $H^{\rm cons}[\eta]$ and its reflected member, gives (N). Theorem~\ref{C-355-9} and its proof supply (D1)--(D4), so Theorem~\ref{cov-thm-data} applies. Lemma~\ref{C-355-23}(b) gives the assertion at every $g\le\gamma$, since dilations commute with rotations. These inputs use the cutoff estimates and conservation argument, before the Wightman reconstruction of Theorem~\ref{C-355-10}.
\end{proof}

\begin{lemma}[Mirrors]\label{C-355-29}
\begin{enumerate}
\item[(a)] For unit vectors $m,n$, $\mathsf R\in\mathrm O(3)$ with $\mathsf Rm=n$, $c\in\mathbb R$ and $\mathsf gx=\mathsf Rx+cn$: $\mathsf g\{m\cdot x>0\}=\{n\cdot x>c\}$ and, with $\theta_m=\theta_{m,0}$, $\theta_{n,c}=\mathsf g\theta_m\mathsf g^{-1}$ on points, on labelled tuples and on their combinations. $\rho(\Theta_n)$ is $+1$ on $\mathbb Rnn^{\rm T}\oplus\{E:En=0\}$ and $-1$ on $\{nm^{\rm T}+mn^{\rm T}:m\perp n\}$; a mirror outside the nine lattice families mixes the diagonal and the off-diagonal labels, and $\Theta_n\in W_3$ if and only if $n$ is one of the eighteen normals of these families.
\item[(b)] A functional invariant under the translations and under $\mathrm O(3)$ and reflection positive at one plane is reflection positive at every plane.
\item[(c)] Let $\sigma$ be a real symmetric multilinear functional with bounds, invariant under the translations and reflection positive at the nine lattice families of mirrors. A functional reflection positive at a mirror is invariant under its reflection, labels turned; a closed subgroup of $\mathrm O(3)$ that contains $W_3$ and one element more is $\mathrm O(3)$. Hence reflection positivity at a single plane whose normal is not one of the eighteen is equivalent to covariance under $\mathrm O(3)$; the set of the mirrors of $\sigma$ is the nine families or all planes, and there are functionals with exactly the nine.
\end{enumerate}
\end{lemma}

\begin{proof}[Proof sketch]
(a) is linear algebra: $\Theta_nE\Theta_n=E-2n(En)^{\rm T}-2(En)n^{\rm T}+4(n\cdot En)nn^{\rm T}$. (b) follows from (a): $\sigma(\theta_{n,c}\bar F\cdot F)=\sigma(\mathsf g(\theta_m\overline{\mathsf g^{-1}F}\cdot\mathsf g^{-1}F))=\sigma(\theta_m\overline{\mathsf g^{-1}F}\cdot\mathsf g^{-1}F)\ge0$. (c) Positivity at the mirror $\{n\cdot x=c\}$ makes the form $(F,F')\mapsto\sigma(\theta_{n,c}\bar F\cdot F')$ Hermitian; at $F'=\varnothing$, $\sigma$ being real, this gives $\sigma(\theta_{n,c}X)=\sigma(X)$ for the tuples $X$ carried by $\{n\cdot x>c\}$, and for every tuple $X$, with $s$ so large that $\tau_{sn}X$ is carried there, $\sigma(X)=\sigma(\tau_{sn}X)=\sigma(\theta_{n,c}\tau_{sn}X)=\sigma(\tau_{-sn}\theta_{n,c}X)=\sigma(\theta_{n,c}X)$ by the invariance under the translations.

So $\sigma$ is invariant under the reflection at every mirror at which it is reflection positive, labels turned, and with the translations under the reflection $\Theta_n$ in the parallel plane through $0$. The nine reflections of the lattice families generate $W_3$; so the group of the elements of $\mathrm O(3)$ under which $\sigma$ is invariant, which is closed by the bounds, contains $W_3$; the octahedral group $W_3$ is a maximal finite subgroup of $\mathrm O(3)$, so a closed subgroup containing one element more is infinite and contains a one-parameter group of rotations; since $W_3$ acts irreducibly on $\mathbb R^3$, the conjugates of its axis generate $\mathrm{SO}(3)$, and with $-1\in W_3$ the group is $\mathrm O(3)$.

For the last clause let $\sigma$ be a tensor limit at $g\le\gamma^{\rm sd}_W$, covariant by Theorem~\ref{C-355-10}~(a), let $z\ne\pm1$ be real and $Z_z$ the linear map of $V$ that is the identity on the diagonal and the factor $z$ on the off-diagonal matrices. It commutes with $\rho(W_3)$, so $\sigma\circ Z_z$, $Z_z$ applied to every entry, is real, symmetric and multilinear, keeps the bounds and the translations, and is reflection positive at the nine families. Were it covariant, the rotation by $\frac\pi4$ about $e_3$, which carries $o=e_1e_2^{\rm T}+e_2e_1^{\rm T}$ to the diagonal label $e_2e_2^{\rm T}-e_1e_1^{\rm T}$, would give $(z^2-1)\sigma(h_1o,h_2o)=0$ for all real $h_1,h_2$ with disjoint supports, against the free law (E6) of Theorem~\ref{C-355-15}, by which $\sigma(\theta(\varphi o),\varphi o)>0$ for a small bump $\varphi$.

So the mirrors of $\sigma\circ Z_z$ are exactly the nine families.
\end{proof}

\begin{proposition}[The lattice stress identity; the remainder; one cell]\label{C-355-30}
Let $0<g\le\gamma$, let $\ell'=2\mathsf l$ be one side of $\mathfrak L(g)$, fixed once, and let $(K,\beta)$ be first-passage data at $g$.
\begin{enumerate}
\item[(a)] \emph{(Identity.)} For $\eta\in C^2_c(\mathbb R^3;\mathbb R^3)$ there is an exact identity of functions on the configurations of the torus
\[b\,\mathcal T_K(H^{\rm cons}[\eta])=\mathcal N_\eta-\tfrac14\mathfrak N[\mathsf k_\eta]-c_\eta-\sum_{e=1}^9\mathfrak J_e[\eta],\]
where $\mathcal N_\eta$ and $\mathfrak N[\mathsf k_\eta]$ are null terms: every joint cumulant of order $\ge2$ on $T_{\ell'}$ with a null term in one slot and, in the others, bounded functions of the bonds at sup distance $>a_K$ from $\operatorname{supp}\eta$, vanishes; $c_\eta$ is a constant; and the $\mathfrak J_e[\eta]$ are nine local lattice remainders, entries of order $\ge2$ in the field strength, with real weights at most $(3+60L(g^{-2}+\bar B))L^{2K}\|\eta\|_{C^2}$, carried within sup distance $4a_K$ of $\operatorname{supp}\eta$. Among them is the renormalised remainder of the stress identity, with a size $\mathsf H(\eta)$ at most a multiple of $\|\eta\|_{C^2}$.
\item[(b)] \emph{(Remainder.)} Let $0<\phi\le\frac1{84}$, $u=\lceil\phi K\rceil$, $t=K-u$, and let $\omega_t(\mathfrak J)$ be the distance from the constants, in the sup norm on the configurations without large fields, of the conditional mean of $\mathfrak J$ given the block fields at the level $t$. There are $C_\omega$, $p_\omega$ (formed from the constants, $g$, $\ell'$, $\phi$) and, for a margin $0<\delta\le\frac18\mathsf l$, a cutoff $K_{\rm B}$ (formed from these and $\delta$) such that for $K\ge K_{\rm B}$, every board of side $\mathsf l$ whose offset lies in $a_{K-t}\mathbb Z^3$, every cell, every $\eta$ with $\|\eta\|_{C^2}\le1$ carried in it with margin $\delta$, and every mirror copy: $\sum_e\omega_t(\mathfrak J_e[\eta])\le\varpi_{\rm s}=C_\omega K^{p_\omega}L^{3u-K/28}$, while for $\mathcal T_K(H)$ with $|H|_\infty\le1$, $\omega_t\le\varpi_{\rm T}=C_\omega K^{p_\omega}L^{4u}$; and $L^{-K}\le\varpi_{\rm s}\le\varpi_{\rm T}\le L^K$. The renormalised remainder $\mathcal S[b\eta]$ of the stress identity, a density of the third order, has a conditional mean that is a number up to a small error: for $0<\tau\le\frac14$ there are $C_{\rm HR}$, $p_{\rm HR}$ (formed from the constants, $g$, $\ell'$, $\phi$, $\tau$) and a number $\mathsf a[\eta]$, linear in $\eta$, with
\[\bigl|(\mathsf Q_{t\leftarrow0}\mathcal S[b\eta])(V)-\mathsf a[\eta]\bigr|\le C_{\rm HR}\,\mathsf H(\eta)\,K^{p_{\rm HR}}L^{-c_0'K},\qquad c_0'=\min\bigl\{\tfrac1{14}-\tfrac\tau7,\ \tfrac1{14}-3\phi\bigr\},\]
for almost every block field $V$ of the level $t$ without large fields, $\mathsf Q_{t\leftarrow0}$ the conditional mean given the block field of the level $t$ and $\mathsf H(\eta)=\max_\rho\sup|\eta^\rho|+\max_\rho{\rm Lip}\,\eta^\rho$. This holds also when sources enter the chain; at the source $0$, the case here, $\min\{\frac1{10}-\frac\tau5,\frac1{10}-3\phi\}$ may be taken in the place of $c_0'$. At $\tau=\frac14$, $c_0'\ge\frac1{28}$, so that $\omega_t\le C_{\rm HR}\mathsf H(\eta)K^{p_{\rm HR}}L^{-K/28}$.
\item[(c)] \emph{(One cell.)} Let $m\ge1$, $0<\delta\le\frac18\mathsf l$. There are pure numbers $\mathsf C_m$ and a cutoff $K_{\rm P}$ (constants, $g$, $\ell'$, $\phi$, $m$, $\delta$) such that for $K\ge K_{\rm P}$, every $d>0$ with $dL^u\ge K^2$, every board and cell as in (b), and slots $\mathsf X_1,\dots,\mathsf X_m$ in the cell with margin $\delta$, each of them $\mathcal T_K(H_a)$ or $\mathcal T_K(H^{\rm cons}[\eta_a])$ or one of the terms of (a), with supports pairwise at sup distance $\ge d$,
\[\Bigl\|\prod_{a=1}^m\bigl(\mathsf X_a-\langle\mathsf X_a\rangle_{K,\ell',\beta}\bigr)\Bigr\|_{{\rm ch},\ell'}\le\mathsf C_m\prod_{a=1}^m\varpi(\mathsf X_a)\,\mathsf n(\mathsf X_a),\]
$\varpi=\varpi_{\rm T}$ for the tensor slots and $\varpi_{\rm s}$ for the others, $\mathsf n$ the norm $|H_a|_\infty$ or $\|\eta_a\|_{C^2}$.
\end{enumerate}
\end{proposition}

\begin{proof}[Proof sketch]
(a) is the Ward identity of the lattice stress tensor: the Haar measure is invariant under the local variation $U(\mathfrak b)\mapsto e^{tX_{\mathfrak b}}U(\mathfrak b)$ along the field $\mathsf k_\eta$ obtained from $\eta$, and the derivative at $t=0$ of the mean of a product with this variation in one place, integrated by parts on $\mathrm{SU}(2)$ against the Wilson weight, is the null term; the variation of the action is $b\,\mathcal T_K(H^{\rm cons}[\eta])$ up to the remainders, which are of order $\ge2$ in the field strength or carry a factor $a_K$. (b) In the setting (Th) of Section~\ref{sec:cutoff-calibration} with the top level $t$ (Lemma~\YMMref{118.1}{812}), the conditional means $\mathsf Q_{t\leftarrow0}$ given the block fields form a tower of averages along the small-field chain of Proposition~\textup{\YMMref{115.2}{806}}, with the flows of Proposition~\YMMref{114.2}{801}; they contract in the sup norm over the field cuts (the definitions preceding Lemma~\YMMref{113.10}{796}, Lemma~\YMMref{113.10}{796}).

The renormalised remainder $\mathcal S[b\eta]$ is an insertion of the third order; carried through the tower to first order in its parameter, it is the sum of its value at the background of a level $n_0=K-\lceil K/7\rceil$ and of the births of the steps, carried on by chains of averaging operators that are stopped at levels fixed by the count of their members. Five inputs enter. (1) The linearisation of a step: the derivative of the output of one step (Proposition~\ref{prop:source-entry}) with respect to the parameter of the family of the cubic potential, added to the fourth argument of the step, is a pointed family of the next level, linear in the weight; these derivatives are the births. (2) The size of the insertion at the background of the level $n_0$: at most a multiple of $\ell'^3L^{3K-6n_0}$ times the cube of the radius of the field cuts there and the weight. (3) Parity: the part of a birth that is even in the fluctuation field vanishes to the second order, and the odd part enters evenly; this bounds the norm of the birth of the step $k+1$ by $C(\log x_k)^{\mathfrak p_0}g_k^2L^{-3k}$ times the weight. (4) The H\"older continuity of the backgrounds: in a local gauge on a cube about a site (Theorem~\textup{\ref{I-56-65}}) a regular configuration is $e^{ihA}$ with $A$, $\nabla A$ small and $\nabla A$ H\"older continuous. (5) The comparison of a regular configuration with its mirror image under the inversion in a site: the cubic density is odd under this inversion (Proposition~\ref{prop:ward-orbital-cubic}(c)), and in the gauge of (4) the part of the potential that is even under the inversion is of the second order in the ratio of the scales, the odd part smaller by a further power of it, so that the local terms at a site vanish to the second order.

The least of the rates of the parts so bounded is $c_0'$. The nine entries of (a) are bounded one by one, and so are their mirror copies, the mirror maps of a board whose offset lies in $a_{K-t}\mathbb Z^3$ preserving the lattices of the levels up to $t$. The entry $\mathcal S[b\eta]$, of the third order with weights of the order $L^{2K}$, is bounded as just described, at $\tau=\frac14$. Five entries of the second order with weights bounded independently of $K$, a multiple of $a_K$ times a plaquette form among them, and two entries of the third and the fourth order with weights of the order $L^K$, are bounded by power counting along the chain with one insertion: an entry of the second order has $\omega_t\le CK^{2\mathfrak p_0}L^{4u-K}$ times its weight, an entry of a higher order $\omega_t\le CK^{3\mathfrak p_0}L^{4u-3K/2}$ times its weight.

For the last entry, of the second order with a weight of the order $L^K$, power counting gives only the size $L^{4u}$ of a tensor entry; its conditional mean is first cut at the level $n_\Delta=K-y$, $y=\lceil K/14\rceil$, where $\omega_t\le\epsilon_X+CK^{3\mathfrak p_0}L^{3u-K-y/2}$ times the weight, $\epsilon_X$ being its supremum on the regular backgrounds of the level $n_\Delta$. By the H\"older continuity of the backgrounds with the exponent $\frac12$ (input (4)), $\epsilon_X\le CK^{2\mathfrak p_0}L^{7y/2-K/2}b\sup|\nabla\eta|$, and the second member is at most $CK^{3\mathfrak p_0}L^{3u-y/2}b$. As $u\le\frac{13}{28}K$, all these are at most $CK^{p}L^{3u-K/28}\|\eta\|_{C^2}$: the rate $\frac1{28}$ is fixed by this last entry, through $L^{3u-y/2}$, and by $\tau=\frac14$ in the bound of $\mathcal S[b\eta]$.

A tensor entry only grows, by at most $L^4$ per scale over the last $u$ steps: power counting gives $\omega_t\le CK^{2\mathfrak p_0}L^{4u}|H|_\infty$. (c) The eighth power of the seminorm is the mean on $T_{\ell'}$, whose board has eight cells, of the $8m$ mirror copies $Y_{(a,\kappa)}$, all centred. The mean is the sum over the partitions into blocks of at least two members of products of cumulants; in each cumulant every slot may be replaced by its list of entries of (a), since the null terms and constants drop out of cumulants of order $\ge2$ whose other slots depend only on bonds away from their support. The carriers of the entries of different copies are at sup distance at least $d_\times=\min\{d,2\delta\}-8a_K\ge\frac12\min\{d,2\delta\}$, and $d_\times L^{K-i}\ge d_\times L^u\ge\frac12K^2$ at the levels $i\le t$, because $dL^u\ge K^2$ and $K_{\rm P}$ is so large that $2\delta L^{\phi K}\ge K^2$.

A cumulant of $r+1\le N=8m$ such entries is at most $\mathsf{op}_{r+1}\prod_i\omega_t(\cdot)$, with $\mathsf{op}_N=\sum_\pi(|\pi|-1)!$ over the partitions $\pi$ of $N$ members, plus two members: $C^N(K^{3\mathfrak p_0}L^{2K})^N\mathsf t_\times$, with $\mathsf t_\times=\sum_{i\le t}\ell'^3L^{3(K-i)}e^{-d_\times L^{K-i}/M}\le CK\ell'^3L^{3u}e^{-K^2/2M}$, under the condition $32(N+1)C_8\mathsf t_\times\le1$, $C_8$ a number formed from the constants; and $4\,\mathsf{op}_{r+1}\mathsf R_t\prod\sup|\mathfrak J|$, with $\mathsf R_t\le CL^{3u}\exp[-c'A_0^2(\phi K\log L)^{2\mathfrak p_0}]$ the weight of the creation of a large field up to the level $t$, $\mathfrak p_0\ge1$, and $\sup|\mathfrak J|\le CL^{5K}$: the configurations with large fields, which $\omega_t$ does not see, enter here.

This bound comes from the chain with several insertions of disjoint carriers, whose exponents differ from the sum of the exponents with one insertion only by polymers that hold boxes of two insertions; such a polymer has a linear size at least the sup distance of the carriers in the units of its level divided by $M$, less $2$ (\eqref{eq:analytic-domain-diameter}), and the sum over these polymers is bounded by the entropy bound of \eqref{eq:analytic-domain-tail}. The logarithms of both members are $-cK^2+O(K\log K)$ with $c>0$; so for $K\ge K_{\rm P}$ they are at most $\mathsf c_NL^{-2NK}$, $\mathsf c_N$ a pure number. Every copy lies in exactly one block, so the sum is $\mathsf{Op}_{8m}\prod_a\varpi(\mathsf X_a)^8$, $\mathsf{Op}_N$ being the sum over the partitions of $N$ members into blocks of at least two members of the products of the numbers $\mathsf{op}$ of the blocks, up to a term $\le L^{-8mK}\le\prod_a\varpi(\mathsf X_a)^8$; take the eighth root, which gives $\mathsf C_m=(\mathsf{Op}_{8m}+1)^{1/8}$.
\end{proof}

\begin{lemma}[Conservation of kernels and of operators; for data]\label{C-355-31}
Let $\sigma'$ be scalar reconstruction data with a unique vacuum on $\mathcal H$, with a tensor of one label satisfying (M1), (M2), forms $\mathfrak q_{ij}$ and kernels $\mathsf s_{ij}$.
\begin{enumerate}
\item[(a)] A subspace $\mathcal V\subset\mathcal D_\infty$, dense in $\mathcal H$ and invariant under $e^{-\tau H}$, $\tau>0$, is a core of $(1+H)^N$ for every $N$.
\item[(b)] $\mathcal V_1$, the span of the $e^{-\tau H}[F']$, and, if $\mathcal H\ne\mathbb C\Omega$, $\mathcal V_2$, the span of the $\mathcal U(v)\mathsf B_0(s,s')[F]$, $\mathsf B_0$ the scalar point operators, are such subspaces.
\item[(c)] If $\sum_i\int\partial_i\zeta(x)\,\mathsf s_{ij}(x;y)\,dx=0$ for all $m\ge1$, $y\in(\mathbb R^3)^m_{\ne}$ and $\zeta\in C^\infty_c$ off the $y_k$, then the form $CL_j(\psi',\psi)=\mathfrak q_{3j}(H\psi',\psi)-\mathfrak q_{3j}(\psi',H\psi)-i\sum_k[\mathfrak q_{kj}(P_k\psi',\psi)-\mathfrak q_{kj}(\psi',P_k\psi)]$, $P_k=\int q_k\,d\mathcal P$, vanishes on $\mathcal D_\infty\times\mathcal D_\infty$, and $T_{0\nu}(\partial_0F)\psi=\sum_{k=1,2}T_{k\nu}(\partial_kF)\psi$ on $\mathcal D_\infty$. Conversely (M4)~(a) gives the identity of the kernels.
\end{enumerate}
\end{lemma}

\begin{proof}[Proof sketch]
(a) A vector orthogonal to $\mathcal V$ in the graph norm of $(1+H)^N$ gives a finite measure $(1+E)^N\mu$ that annihilates all the functions $e^{-\tau E}$, hence vanishes by the Stone--Weierstrass theorem. (b) The closure of $\mathcal V_2$ reduces $H$ and $\mathcal U$ and contains $\phi(f)\mathcal D_\infty$, hence all the $\phi(f_1)\cdots\phi(f_n)\Omega$; if it were proper, its orthogonal complement would hold an invariant vector, hence $\Omega$ by the unique vacuum, and $\phi(f)\Omega=0$ would follow. (c) Let $a>0$, let $F',G$ be scalar tuples over $\mathbb E_+$ with smooth entries, $F'(z')$, $G(z)$ the products of their entries at the points $z'$, $z$, and $P_k=\int q_k\,d\mathcal P$, so that $\mathcal U(v)=e^{iv\cdot P}$. For $x\ne y$ in $\{|x_3|<a\}$ with $x_3<y_3$ put
\[\mathsf s^{F',G}_{ij}(x,y)=\int\overline{F'(z')}\,G(z)\,\mathsf s_{ij}\bigl(x;y,\theta(z'+ae_3),z+ae_3\bigr)\,dz'\,dz ;\]
the operator formula of (M2), the product of the point operator of $\mathfrak q_{ij}$ at $x$ and of the scalar point operator at $y$ between $[F']$ and $[G]$, reads $\mathsf s^{F',G}_{ij}(x,y)=\mathfrak q_{ij}(\psi'_x,\psi_x)$ with $\psi'_x=\mathcal U(-\vec x)e^{-(a+x_3)H}[F']$ and $\psi_x=\mathcal U(-\vec x)e^{-(h_*-x_3)H}\mathcal U(\vec y)\mathsf B_0(y_3-h_*,a-y_3)\mathcal U(-\vec y)[G]$, for any $h_*\in(x_3,y_3)$. As $\partial_{x_3}\psi'_x=-H\psi'_x$, $\partial_{x_3}\psi_x=H\psi_x$, $\partial_{x_k}\psi'_x=-iP_k\psi'_x$ and $\partial_{x_k}\psi_x=-iP_k\psi_x$ for $k=1,2$, one has $\sum_i\partial_{x_i}\mathsf s^{F',G}_{ij}(x,y)=-CL_j(\psi'_x,\psi_x)$, and $|CL_j(\psi',\psi)|\le6c_T\|(1+H)^{l+1}\psi'\|\,\|(1+H)^{l+1}\psi\|$.

Given $\tau,s,s'>0$, $v\in\mathbb R^2$ and $F$, the choice $a=\frac12(\tau+s+s')$, $x=(0,\tau-a)$, $y=(v,a-s')$, $G=\tau_{(v,0)}F$ gives $\psi'_x=e^{-\tau H}[F']$ and $\psi_x=\mathcal U(v)\mathsf B_0(s,s')[F]$. The identity of the kernels supposed in (c), integrated against $\overline{F'}G$, kills the divergence near $x$; so $CL_j=0$ on $\mathcal V_1\times\mathcal V_2$, and on $\mathcal D_\infty\times\mathcal D_\infty$ by (a), (b) and the bound. Read at $U(b)^{-1}\psi'$, $U(b)^{-1}\psi$, $b\in\mathbb R^{1+2}$, multiplied by the test function of the operator identity and integrated by parts in $b$, this is that identity. Conversely, $\mathcal D_0$ is invariant under the operators $U(b)$, so it is a core of every $(1+H)^N$, as in (a) with the group in the place of the semigroup, the finite measure $(1+E)^N\mu$ having a vanishing Fourier transform; the bounds $\|T_{\mu\nu}(F_1)\psi\|\le C\|(1+H)^{2l}\psi\|$ of the operators of (M1) carry (M4)~(a) from $\mathcal D_0$ to $\mathcal D_\infty$, and, the test function being arbitrary, $CL_j=0$ on $\mathcal D_\infty\times\mathcal D_\infty$.

At a configuration $(x,y_1,\dots,y_m)$ with pairwise different heights take a slab that holds $x$ and the point next to it in height and no other point; the other points are smeared by bumps into $G$ above and, through their mirror images, into $F'$ below; the divergence of $\mathsf s^{F',G}_{ij}$ is $-CL_j(\psi'_x,\psi_x)=0$, and approximate units for the bumps give $\sum_i\int\partial_i\zeta(x)\,\mathsf s_{ij}(x;y)\,dx=0$ for $\zeta$ carried near $x$. A rotation that makes the heights pairwise different reaches every configuration, by the covariance of the kernels in (M2), and a partition of unity every $\zeta$.
\end{proof}

\begin{lemma}[Plateaux on the lattice]\label{C-355-32}
Let $0<g\le\gamma^{\rm ref}_W$, let $(K_j,\beta_j,\nu_j)$ be first-passage data at $g$ with periodic Gibbs states and $K_j\to\infty$, $A_j(X)=A^{\nu_j}_{K_j}(X)$, and $\mathcal I$ the invariant $\sigma$-field. In (c) let moreover $\sigma(X)=\lim_j\nu_j(A_j(X))$ exist for every $X\in\mathfrak S_V(\mathbb R^3)$: $\sigma$ is the tensor limit of the sequence, as in the head of Theorem~\ref{C-355-13}.
\begin{enumerate}
\item[(a)] For $n$, $\delta$, $\varrho$ there are $c_n(\delta;\varrho)$ and $j_n(\delta;\varrho)$ such that for tuples $X,X'$ of $n$ entries carried by one family of cubes of radius $\varrho$ at distances $\ge\delta$, $j\ge j_n$ and $\rho\in\mathcal G^{\rm per}_{K_j,\beta_j}$: $|\rho(A_j(X))-\rho(A_j(X'))|\le c_n\mathsf D(X,X')$, $|E_\rho[A_j(X)|\mathcal I]|\le c_n|X|_\infty$ and $|E_\rho[A_j(X)-A_j(X')|\mathcal I]|\le c_n\mathsf D(X,X')$ almost surely, with $|X|_\infty=\prod_a|H_a|_\infty$ and $\mathsf D(X,X')=\sum_a|H_a-H'_a|_\infty\prod_{b<a}|H'_b|_\infty\prod_{b>a}|H_b|_\infty$.
\item[(b)] For tuples $X,Y$ in a cube $(-\bar R,\bar R)^3$, $j$ large and lattice vectors $v$ with $|v|\ge4\sqrt3\bar R$: $|\rho(A_j(X\cdot\tau_vY))-\rho(E_\rho[A_j(X)|\mathcal I]\,E_\rho[A_j(Y)|\mathcal I])|\le\mathfrak k_0\mathsf N(X)^{1/2}\mathsf N(Y)^{1/2}\bar R/|v|$, with numbers $\mathsf N$ formed from the tuples.
\item[(c)] For $X^1,\dots,X^q$ the means $\nu_j(\prod_i\mathsf e_j(X^i))$, $\mathsf e_j(X)$ a real-valued choice of $E_{\nu_j}[A_j(X)|\mathcal I]$, converge, and the limit is the iterated plateau of $\sigma$ at $(X^1,\dots,X^q)$ (Definition~\ref{C-355-4}~(5)): with $G^{(1)}(v_2,\dots,v_q)=\sigma(X^1\cdot\tau_{v_2}X^2\cdots\tau_{v_q}X^q)$ where this tuple is formed, the limits $G^{(p)}(v_{p+1},\dots,v_q)=\lim_{|v_p|\to\infty}G^{(p-1)}(v_p,v_{p+1},\dots,v_q)$, $p=2,\dots,q$, exist at every $(v_{p+1},\dots,v_q)$ for which $X^1\cdot\tau_{v_{p+1}}X^{p+1}\cdots\tau_{v_q}X^q$ is formed, and the number $G^{(q)}$ is the iterated plateau.
\item[(d)] Let $\rho_j\in\mathcal G^{\rm per}_{K_j,\beta_j}$, and let $\sigma'$ be a functional with $\|E_{\rho_j}[A_j(X)|\mathcal I]-\sigma'(X)\|_{L^\infty(\rho_j)}\to0$ for every $X$, so that $\sigma'(X)=\lim_j\rho_j(A_j(X))$, the products $A_j(X)$ being centred by $\nu_j$ and integrated by $\rho_j$. Then $\lim_{|v|\to\infty}\sigma'(X\cdot\tau_vY)=\sigma'(X)\sigma'(Y)$.
\end{enumerate}
\end{lemma}

\begin{proof}[Proof sketch]
(a) Multilinearity and the uniform bound of Theorem~\ref{C-355-7}~(c)~(iii), which holds at every $0<g\le\gamma$ for any two periodic Gibbs states, the one integrating and the other centring, give the first inequality, $j_n(\delta;\varrho)$ being an index from which on $K_j$ is at least the cutoff $K^A$ of that clause at the order $n$ and the distance $\delta$, which exists because $K_j\to\infty$; the invariant-conditioning argument in the proof of Proposition~\YMMref{100.1}{744}, for the class of the periodic Gibbs states at one cutoff, carries it to the conditional expectations given $\mathcal I$. (b) is the clustering on one lattice at the rate $ra/|x|_\infty$, Proposition~\textup{\YMMref{99.7}{742}}, from reflection positivity at the nine families of mirrors of one lattice, applied to the centred products as in Proposition~\YMMref{100.1}{744}; no mass enters, only a power. (c) follows from (a), (b) by induction on $q$, the plateaux being taken one translation after another. (d) is (b) with the conditional means replaced by constants.
\end{proof}

\begin{lemma}[Identification of the limits; recentring]\label{C-355-33}
\begin{enumerate}
\item[(a)] A sequence of first-passage data $(K_j,\beta_j)$ at $g$, with $K_j\to\infty$ and periodic Gibbs states, whose labels and centred moments on the tuples with $C^1_c$ entries converge has a tensor limit that agrees with its limit on these tuples, with the same label; the restriction of a tensor limit to $C^1_c$ entries is such a limit; a limit along data whose labels do not converge has a subsequence whose labels converge, since they lie in a compact interval. The lattice observables of the plaquette and bond forms are the same in all these limits.
\item[(b)] Let $\sigma^c$ be a limit of periodic Gibbs states and $m(H)=\frac c3\int\operatorname{tr}H$. Then $\sigma^\#(X)=\sum_{S}\prod_{a\notin S}m(H_a)\,\sigma^c(X_S)$ is a symmetric functional; since $m(\mathsf rH)=m(H)$ for every isometry and $m(H^{\rm cons}[\eta])=\frac c3\int\nabla\cdot\eta=0$, covariance and conservation pass from $\sigma^c$ to $\sigma^\#$.
\item[(c)] For a functional $\sigma_\xi$ with $s_\xi(H)=\sigma_\xi((H))=\frac{c_\xi}3\int\operatorname{tr}H$, the map $\mathsf S_\xi X=\sum_J\prod_{a\notin J}s_\xi(H_a)X_J$ is a bijection of the combinations of tuples, multiplicative for the concatenation, commuting with the conjugation and with every isometry, and $\sigma_\xi=\sigma^c_\xi\circ\mathsf S_\xi$ for the recentred $\sigma^c_\xi$; $[F]\mapsto[\mathsf S_\xi F]$ is unitary and intertwines $e^{-tH}$ and $\mathcal U$, and the point operators correspond as $\mathsf B(s,s';e)\cong\mathsf B^c(s,s';e)+\frac{c_\xi}3(\operatorname{tr}e)e^{-(s+s')H}$.
\end{enumerate}
\end{lemma}

\begin{proof}[Proof sketch]
(a) Extraction along a countable family of tuples dense for the uniform norm on each product of compact sets, with the continuity of Theorem~\ref{C-355-9}~(b); the labels lie in $[g^{-2},L(g^{-2}+\bar B)]$. (b), (c) are algebra: $\operatorname{tr}H^{\rm cons}[\eta]=3\nabla\cdot\eta-2\nabla\cdot\eta=\nabla\cdot\eta$, and the inverse of $\mathsf S_\xi$ is the same sum with $-s_\xi$.
\end{proof}

\begin{lemma}[Extension over the diagonals; contact terms]\label{C-355-34}
Let $\mathbb M$ be $\mathbb R^3$ or a torus $T_\ell$, $n\ge2$, and $V$ a real vector space of finite dimension with a norm. A \emph{functional of $n$ entries} is a map $u$ on the tuples $(f_1,\dots,f_n)$ of functions $f_a\in C^\infty_c(\mathbb M;V_{\mathbb C})$ with pairwise disjoint supports that is $n$-linear on $\prod_aC^\infty_c(O_a;V_{\mathbb C})$ for all open sets $O_a$ with pairwise disjoint closures; the entries may as well take their values in $n$ different spaces $V_1,\dots,V_n$, with $\bigotimes_aV_{a,\mathbb C}$ in the place of $V_{\mathbb C}^{\otimes n}$ below, and nothing else changes (this is used in (e)). We write $\|G\|_{(h),r}=\max_{|\gamma|\le r}h^{|\gamma|}\sup|\partial^\gamma G|$, on $\mathbb M$ and on $\mathbb M^n$; $\mathsf d_\infty(x)=\min_{a\ne b}|x_a-x_b|_\infty$ for $x\in\mathbb M^n$; and $\Delta=\{\mathsf d_\infty=0\}$, the union of the diagonals.
\begin{enumerate}
\item[(1)] \emph{Strata.} For a partition $\pi$ of $\{1,\dots,n\}$ with $k=|\pi|$ blocks, $B(a)$ the block of $a$ and $a_B=\min B$: $\Delta_\pi=\{x:x_a=x_b\text{ whenever }B(a)=B(b)\}$, and $\Sigma_k=\bigcup_{|\pi|\le k}\Delta_\pi$, so that $\Sigma_{n-1}=\Delta$. The coordinates at $\pi$ are $y_B=x_{a_B}\in\mathbb M$ and $\xi_a=x_a-x_{a_{B(a)}}$ for $a\ne a_{B(a)}$; $\iota_\pi(y)_a=y_{B(a)}$, so that $\Delta_\pi=\{\xi=0\}=\iota_\pi(\mathbb M^k)$. For a distribution $T'$ on $\mathbb M^k$ and a multi-index $\alpha$ of $\xi$, $(\partial_\xi^\alpha\delta_\pi\otimes T')(F)=(-1)^{|\alpha|}\,T'\bigl((\partial_\xi^\alpha F)\circ\iota_\pi\bigr)$.
\item[(2)] \emph{The bound $(\mathrm W;\kappa,q,r)$}, $\kappa\in\mathbb R$, $q\ge0$, $r\in\mathbb N_0$, with numbers $A,h_0>0$ (on $T_\ell$: $h_0\in\ell2^{-\mathbb N}$, $h_0\le\ell/(64n)$): for every $0<h\le h_0$, all closed cubes $Q^*_1,\dots,Q^*_n\subset\mathbb M$ of the side $\frac32h$, wherever they lie, at pairwise sup distances $\ge\frac52h$, this lower bound being the only condition on their distances, and all $f_a\in C^\infty_c(\operatorname{int}Q^*_a;V_{\mathbb C})$,
\[|u(f_1,\dots,f_n)|\le A\,h^\kappa\,l(h)^q\prod_a\|f_a\|_{(h),r},\qquad l(h)=1+\log(h_0/h).\]
\item[(3)] \emph{The class.} A function $F$ on $\mathbb M^n$ is \emph{$\lambda$-supported} if $\operatorname{supp}F\subset\prod_a\{|x_a-z_a|_\infty\le\lambda\}$ for some $z\in\mathbb M^n$, the $z_a$ coinciding or not. $\mathfrak E_n(\kappa)$ is the set of the $(V_{\mathbb C}^{\otimes n})^*$-valued distributions $T$ on $\mathbb M^n$ for which there are $C$, $q$, $m$, $h_1$ with $|T(F)|\le C\lambda^\kappa(1+\log\frac{h_1}\lambda)^q\|F\|_{(\lambda),m}$ for all $0<\lambda\le h_1$ and all $\lambda$-supported $F$: the \emph{power counting $\kappa$}, up to logarithms, also where the points meet. On $\mathbb M=T_\ell$ the norms $\|\cdot\|_{(\lambda),m}$, the $\lambda$-supported functions, the classes $\mathfrak E_n(\kappa)$, the coordinates at $\pi$ and the distributions $\partial_\xi^\alpha\delta_\pi\otimes T'$ of (1) are those of Definition~\ref{C-355-6}. We put $\omega_k=\lfloor3k-\kappa\rfloor$.
\item[(4)] \emph{Whitney products.} A product $Q=\prod_aQ_a$ of half-open cubes $Q_a\subset\mathbb M$ of one side $h(Q)\le h_0$ is \emph{admissible} if ${\rm dist}_\infty(Q_a,Q_b)\ge3h(Q)$ for $a\ne b$; $c(Q)$ is its centre, and $Q^+$, $Q^*$ are the closed concentric products of the sides $\frac54h(Q)$ and $\frac32h(Q)$, so that the cubes $Q^*_a$ are at pairwise sup distances $\ge\frac52h(Q)$. Let $h_j=2^{-j}h_0$, $j\ge0$, and let $\mathcal Q$ be the set of the admissible products of dyadic cubes of the sides $h_j$ that lie in no larger admissible product of dyadic cubes. For $Q\in\mathcal Q$, $\pi(Q)$ is the partition of $\{1,\dots,n\}$ into the connected components of the graph $a\sim b\iff|c(Q)_a-c(Q)_b|_\infty\le12h(Q)$, and $k(Q)=|\pi(Q)|$.

Then $\mathcal Q$ is a partition of $\mathbb M^n\setminus\Delta$; on $Q$, $\mathsf d_\infty\ge3h(Q)$, and if $h(Q)<h_0$, $\mathsf d_\infty<10h(Q)$ and $k(Q)\le n-1$; there are $\varphi_Q\in C^\infty_c(\operatorname{int}Q^+)$, $\varphi_Q\ge0$, with $\sum_{Q\in\mathcal Q}\varphi_Q=1$ on $\mathbb M^n\setminus\Delta$ and $|\partial^\gamma\varphi_Q|\le c_\gamma h(Q)^{-|\gamma|}$. Fix a smooth $\mu$ on $\mathbb R^3$ carried by $\{|y|_\infty<1\}$ with $\int\mu=1$. For $Q\in\mathcal Q$ with $h=h(Q)$, $\pi=\pi(Q)$, $k=k(Q)\le n-1$ and $\omega_k\ge0$ put $\mu_Q(y')=\prod_{B\in\pi}h^{-3}\mu\bigl((y'_B-c(Q)_{a_B})/h\bigr)$ for $y'\in\mathbb M^k$ and
\[P_QF(x)=\int\mu_Q(y')\sum_{|\gamma|\le\omega_k}\frac{(x-\iota_\pi(y'))^\gamma}{\gamma!}\,\partial^\gamma F(\iota_\pi(y'))\,dy' ,\]
the Taylor polynomial of $F$ of total order $\omega_k$ averaged over the piece of the size $h$ of the stratum $\Delta_\pi$ next to $Q$; $P_Q=0$ for the other $Q$.
\end{enumerate}
Let $u$ be a functional of $n$ entries with $(\mathrm W;\kappa,q,r)$, and put $N_0=\lfloor\frac{3n}2\rfloor+1+nr$ and $m=\max\{\omega_{n-1}+1,N_0\}$.
\begin{enumerate}
\item[(a)] For every admissible product $Q$, dyadic or not, in particular for every $Q\in\mathcal Q$, there is exactly one linear functional $u_Q$ on $C^\infty_c(\operatorname{int}Q^+;V_{\mathbb C}^{\otimes n})$ with $u_Q(g_1\otimes\dots\otimes g_n)=u(g_1,\dots,g_n)$ and
\[|u_Q(G)|\le C_{\rm F}\,A\,h(Q)^\kappa\,l(h(Q))^q\,\|G\|_{(h(Q)),N_0},\]
$C_{\rm F}$ formed from $n$, $r$ and the norm of $V$; it reads $(\mathrm W)$ at the side $h(Q)$ alone, and two of these functionals agree where both are defined. There is exactly one distribution $u^\circ$ on $\mathbb M^n\setminus\Delta$ with $u^\circ(f_1\otimes\dots\otimes f_n)=u(f_1,\dots,f_n)$ for the tuples of the definition, and $u^\circ=u_Q$ on $C^\infty_c(\operatorname{int}Q^+)$: off the diagonals $u$ is given by exactly one distribution.
\item[(b)] The sum
\[\tilde u(F)=\sum_{Q\in\mathcal Q}u_Q\bigl(\varphi_Q(F-P_QF)\bigr)\]
converges absolutely for $F\in C^m_c(\mathbb M^n;V_{\mathbb C}^{\otimes n})$, and $\tilde u$ is a distribution of order at most $m$ that extends $u^\circ$. At $\kappa=0$, $r=0$: $\omega_k=3k$ and $m=3n-2$.
\item[(c)] For $0<\lambda\le h_0$ and every $\lambda$-supported $F$,
\[|\tilde u(F)|\le C_{\rm ext}\,A\,\lambda^\kappa\,l(\lambda)^{q+\epsilon}\,\|F\|_{(\lambda),m},\]
where $\epsilon=1$ if $3k-\kappa$ is an integer $\ge0$ for some $k\in\{1,\dots,n-1\}$, $\epsilon=0$ otherwise, and $C_{\rm ext}$ is formed from $n,\kappa,q,r$, the norm of $V$ and $\mu$, and does not depend on $u$; so $\tilde u\in\mathfrak E_n(\kappa)$. \emph{(A range of sides.)} If $u$ has the inequality of $(\mathrm W;\kappa,q,r)$ at the sides $h_\flat\le h\le h_0$ alone, for some $h_\flat\in(0,h_0]$, then the sum $\tilde u^\flat(F)$ over the $Q\in\mathcal Q$ with $h(Q)\ge h_\flat$ has the same bound, with the same $C_{\rm ext}$, for $8h_\flat\le\lambda\le h_0$ and every $\lambda$-supported $F$.
\item[(d)] If $u$ is invariant under a group generated by permutations of the entries, translations of $\mathbb M$ acting on all the entries at once, and a closed subgroup of $\mathrm O(3)$ that maps $\mathbb M$ onto itself, acting by $f\mapsto\rho(\mathsf R)f(\mathsf R^{-1}\cdot)$, then $u^\circ$ has an extension in $\mathfrak E_n(\kappa)$ of order at most $m$ that is invariant under this group.
\item[(e)] \emph{(Contact terms.)} A distribution $T$ on $\mathbb M^n$ with support in $\Delta$ lies in $\mathfrak E_n(\kappa)$ if and only if
\[T=\sum_{|\pi|\le n-1}\ \sum_{|\alpha|\le\omega_{|\pi|}}\partial_\xi^\alpha\delta_\pi\otimes T_{\pi,\alpha}\qquad\text{with}\quad T_{\pi,\alpha}\in\mathfrak E_{|\pi|}(\kappa+|\alpha|),\]
the sum over the partitions $\pi$ with at most $n-1$ blocks and the multi-indices $\alpha$ of $\xi$ at $\pi$. So two extensions of $u^\circ$ in $\mathfrak E_n(\kappa)$ differ by such a sum, and every such sum may be added to an extension. A term with $|\alpha|>\omega_{|\pi|}$ and a coefficient in $\mathfrak E_{|\pi|}(\kappa+|\alpha|)$ vanishes, and for every $\pi$ and every $\alpha$ with $|\alpha|\le\omega_{|\pi|}$ there are nonzero examples when $V\ne\{0\}$. The coefficients in this representation need not be unique. If $\kappa>3(n-1)$, the extension in $\mathfrak E_n(\kappa)$ is unique.
\item[(f)] \emph{(Two points.)} Let $n=2$, $\kappa=0$, $V=\mathrm{Sym}_3(\mathbb R)$ with $\rho(\mathsf R)H=\mathsf RH\mathsf R^{\rm T}$, and let $T\in\mathfrak E_2(0)$ have its support in $\{x_1=x_2\}$, be invariant under the translations, real, and symmetric under the exchange of the two entries. Then $T(f_1\otimes f_2)=\sum_{|\alpha|\le3}\int\mathsf c_\alpha(\partial^\alpha f_1(x),f_2(x))\,dx$ with bilinear forms $\mathsf c_\alpha$ on $V_{\mathbb C}$. If $\mathbb M=\mathbb R^3$ and $T$ is covariant under $\mathrm O(3)$, the symbol $\hat T(p)(H,H')=\sum_\alpha\mathsf c_\alpha(H,H')(ip)^\alpha$ is a real combination of the six forms
\[\operatorname{tr}H\operatorname{tr}H',\quad\operatorname{tr}HH',\quad p^2\operatorname{tr}H\operatorname{tr}H',\quad p^2\operatorname{tr}HH',\quad(p\cdot Hp)\operatorname{tr}H'+(p\cdot H'p)\operatorname{tr}H,\quad(Hp)\cdot(H'p);\]
if $T$ is covariant under the cubic group $W_3$, on $\mathbb R^3$ or on $T_\ell$, then $\hat T$ is a real combination of twelve forms.
\end{enumerate}
\end{lemma}

\begin{proof}[Proof sketch]
(4) The facts on $\mathcal Q$: if $Q$ is admissible, so is the product of the next generation that holds a given point; $|x_a-x_b|_\infty\ge\mathsf d_\infty(x)-2h_j\ge3h_j$ once $5h_j\le\mathsf d_\infty(x)$, so every $x\notin\Delta$ lies in exactly one member of $\mathcal Q$; if $h(Q)<h_0$ the parent of $Q$ is not admissible, whence $\mathsf d_\infty<10h(Q)$ on $Q$ and two points in one block of $\pi(Q)$ are linked by a chain of centres at most $12h(Q)$ apart. Members of $\mathcal Q$ whose sets $Q^+$ meet have sides in the ratio $\frac12$, $1$ or $2$, and a point lies in at most $3\cdot2^{3n}$ sets $Q^+$; $\varphi_Q$ is a normalised product of cutoffs of the scale $h(Q)$. A count used below: at most ${\rm Bl}(n)(24n)^{3n}(3R/h_j)^{3k}$ members of $\mathcal Q$ have the side $h_j$, $k$ clusters and all the coordinates of their centres within $R\ge h_j$ of those of a given point of $\mathbb M^n$, a partition, one cube for each block and the cubes of the other indices within $12(n-1)h_j$ of it being chosen.

(a) Let $h=h(Q)$, $\mathsf s=\frac32h$, $c=c(Q)$. A $G\in C^\infty_c(\operatorname{int}Q^+;V_{\mathbb C}^{\otimes n})$, continued periodically with the cell $Q^*$, is the sum of its Fourier series in the modes $\mathsf E_{\mathbf k}(x)=\prod_ae^{2\pi i\mathbf k_a\cdot(x_a-c_a)/\mathsf s}$, $\mathbf k\in\mathbb Z^{3n}$, with all derivatives, uniformly; with cutoffs $\chi_a\in C^\infty_c(\operatorname{int}Q^*_a)$, equal to $1$ on $Q^+_a$, fixed once at the scale $h$, $\|\chi_ae^{2\pi i\mathbf k_a\cdot(\cdot-c_a)/\mathsf s}\|_{(h),r}\le C_r(1+|\mathbf k_a|)^r$, and every mode times $\prod_a\chi_a$ is a product of entries carried by the cubes $Q^*_a$, of the side $\frac32h$ and at pairwise sup distances $\ge\frac52h$.

Put $u_Q(G)=\sum_{\mathbf k}\hat G_{\mathbf k}\,u(\chi_1e^{2\pi i\mathbf k_1\cdot(\cdot-c_1)/\mathsf s},\dots,\chi_ne^{2\pi i\mathbf k_n\cdot(\cdot-c_n)/\mathsf s})$, componentwise in a basis of $V$. By $(\mathrm W)$ the terms are at most $Ah^\kappa l(h)^qC_r^n(1+|\mathbf k|)^{nr}|\hat G_{\mathbf k}|$; by the Cauchy--Schwarz inequality, with $\sum_{\mathbf k}(1+|\mathbf k|)^{-2s}<\infty$ for $2s>3n$, and by Parseval, $\sum_{\mathbf k}(1+|\mathbf k|)^{nr}|\hat G_{\mathbf k}|\le C\|G\|_{(h),N_0}$: this is the origin of $N_0=\lfloor\frac{3n}2\rfloor+1+nr$. At a product $G=g_1\otimes\dots\otimes g_n$ the coefficients are products, the series of the entries converge in $\|\cdot\|_{(h),r}$, and $u_Q(G)=u(g_1,\dots,g_n)$.

A function carried by a product of open sets $O_a$ is a limit in every $C^N$, within one compact subset of $\prod_aO_a$, of finite sums of products of functions carried by the $O_a$ (truncated Fourier series times cutoffs); so two functionals that are bounded by a $C^N$ norm and equal to $u$ on products agree, which gives the uniqueness of $u_Q$ and the agreement of $u_Q$, $u_{Q'}$ on $\operatorname{int}Q^+\cap\operatorname{int}Q'^+$. Then $u^\circ(G)=\sum_Qu_Q(\varphi_QG)$ is a finite sum on every compact subset of $\mathbb M^n\setminus\Delta$, a distribution; on a tuple with disjoint supports, a partition of unity of $\mathbb M$ by dyadic cubes of a side $h_j$ with $6h_j$ below the least distance of the supports and the multilinearity of $u$ give $u^\circ=u$, the products of the cubes that meet the supports being admissible; and $u^\circ$ is unique by the same density.

(b) Four facts on $P_Q$, at $Q\in\mathcal Q$ with $h=h(Q)$, $\pi=\pi(Q)$, $k=k(Q)\le n-1$, $\omega=\omega_k\ge0$. (P1) For $x\in Q^+$ and $y'$ in the support of $\mu_Q$, $|x_a-y'_{B(a)}|_\infty\le\frac58h+12(n-1)h+h\le12nh$. (P2) $P_QF=0$ if $F$ vanishes near $\Delta_\pi$: it reads the derivatives of $F$ on $\Delta_\pi$ alone. (P3) For $h\le\lambda$ and $|\gamma|\le N_0$, on $Q^+$, $h^{|\gamma|}|\partial^\gamma(F-P_QF)|\le C(h/\lambda)^{\omega+1}\|F\|_{(\lambda),m}$: since $\int\mu_Q=1$ and $\partial^\gamma$ of the Taylor polynomial of order $\omega$ is the Taylor polynomial of order $\omega-|\gamma|$ of $\partial^\gamma F$, Taylor's formula with (P1) gives $(36n^2h)^{\omega+1-|\gamma|}\lambda^{-\omega-1}\|F\|_{(\lambda),\omega+1}$ for $|\gamma|\le\omega$, and for $|\gamma|>\omega$, $\partial^\gamma P_QF=0$. (P4) If $F$ is $\lambda$-supported at $z$ and $\lambda\le h$, then, in the coordinates at $\pi$ and after an integration by parts in $y'$, $h^{|\gamma|}|\partial^\gamma P_QF|\le C\sum_{j\le\omega}(h/\lambda)^{j-3k}\|F\|_{(\lambda),\omega}$ on $Q^+$, because $\mu_Q$ and its derivatives are of the size $h^{-3k}$ and $(\partial_\xi^\alpha F)\circ\iota_\pi$, of the size $\lambda^{-|\alpha|}$, is carried by a set of volume $\le(2\lambda)^{3k}$; and $P_QF=0$ unless $|c(Q)_a-z_a|_\infty\le12nh$ for all $a$.

Now let $F\in C^m_c$ be carried by a ball of radius $R_F\ge h_0$. By (P3) at $\lambda=h_0$, (a), the count of (4) and $\|\varphi_QG\|_{(h),N_0}\le C\|G\|_{(h),N_0}$ on $Q^+$, the terms of the members of the side $h_j$ with $k$ clusters add up to at most $C\,A\,\|F\|_{(h_0),m}((R_F+12nh_0)/h_j)^{3k}h_j^\kappa l(h_j)^q(h_j/h_0)^{(\omega_k+1)_+}$; the exponent of $h_j$ is $\varepsilon_k=\kappa-3k+(\omega_k+1)_+>0$, equal to $\lfloor3k-\kappa\rfloor+1-(3k-\kappa)\in(0,1]$ if $\omega_k\ge0$ and to $\kappa-3k$ if $\omega_k<0$; for $k=n$ only $h_j=h_0$ occurs. So the series converges absolutely, $\tilde u$ is bounded by a $C^m$ norm on every compact set, and by (P2) $\tilde u=u^\circ$ off $\Delta$. At $\kappa=0$, $r=0$: $\omega_{n-1}+1=3n-2\ge\lfloor\frac{3n}2\rfloor+1$.

(c) Let $F$ be $\lambda$-supported at $z$. Only finitely many members with $h(Q)>\lambda$ give a term, and
\[\tilde u(F)=\underbrace{\sum_{h(Q)\le\lambda}u_Q\bigl(\varphi_Q(F-P_QF)\bigr)}_{\rm I}+\underbrace{u^\circ\Bigl(\sum_{h(Q)>\lambda}\varphi_QF\Bigr)}_{\rm II}-\underbrace{\sum_{h(Q)>\lambda}u_Q(\varphi_QP_QF)}_{\rm III}.\]
In I the members that give a term have their centres within $13n\lambda$ of $z$, and (P3) at the scale $\lambda$ with the count gives $C\,A\,\lambda^\kappa\|F\|_{(\lambda),m}\sum_k\sum_{h_j\le\lambda}(h_j/\lambda)^{\varepsilon_k}l(h_j)^q$, which is at most $C\,A\,\lambda^\kappa l(\lambda)^q\|F\|_{(\lambda),m}$, as $l(2^{-i}\lambda)\le2(1+i)l(\lambda)$. In II the function is smooth, carried where $\mathsf d_\infty\ge\frac{11}4\lambda$, with $\|\cdot\|_{(\lambda),N_0}\le C\|F\|_{(\lambda),N_0}$; a partition of unity by all the products of dyadic cubes of one side $h'\in(\frac\lambda8,\frac\lambda4]$ meets its support only in admissible products, at most $19^{3n}$ of them, and (a) bounds each term by $C\,A\,\lambda^\kappa l(\lambda)^q\|F\|_{(\lambda),N_0}$.

In III, by (P4), at most $(24n+1)^{3n}$ members of each side $h_j>\lambda$ give a term, each at most $C\,A\,\lambda^\kappa l(\lambda)^q(h_j/\lambda)^{\kappa+\omega_k-3k}\|F\|_{(\lambda),\omega_k}$; the exponent $\kappa+\omega_k-3k$ is $\le0$, and it is $0$ exactly when $3k-\kappa$ is an integer; so the sum over $\lambda<h_j\le h_0$ is a convergent geometric series or has at most $1+\log_2(h_0/\lambda)\le2\,l(\lambda)$ equal terms: this is the one logarithm more. \emph{A range of sides:} if $(\mathrm W)$ holds only at the sides $\ge h_\flat$ and $8h_\flat\le\lambda$, I loses its members below $h_\flat$ and was bounded term by term; III is unchanged; and II, now $\sum_{h(Q)>\lambda}u_Q(\varphi_QF)$, is the sum over the products of the side $h'>\frac\lambda8\ge h_\flat$ of the partition of II, by the agreement in (a) at sides $\ge h_\flat$, and is bounded in the same way.

(d) The grids, so $\mathcal Q$ and $\pi(Q)$, and with covariant choices $\varphi_Q$, $\mu_Q$, are invariant under the translations by $h_0\mathbb Z^3$, so $\tilde u$ is; the mean of $\tilde u(\tau_t\,\cdot)$ over $t\in[0,h_0)^3$ is invariant under all the translations, and the mean over the permutations and over the Haar measure of the closed subgroup of $\mathrm O(3)$ keeps this invariance. Each of these maps carries $\lambda$-supported functions to $\sqrt3\lambda$-supported ones and changes $\|\cdot\|_{(\lambda),m}$ by a factor formed from $n$ and $m$; so it maps $\mathfrak E_n(\kappa)$ to itself with uniform numbers and extensions of the invariant $u^\circ$ to extensions, and means of extensions with one bound are such extensions.

(e) \emph{If:} for $\lambda$-supported $F$ the function $g=(\partial_\xi^\alpha F)\circ\iota_\pi$ on $\mathbb M^k$ is $\lambda$-supported, and $\|g\|_{(\lambda),m'}\le n^{m'}\lambda^{-|\alpha|}\|F\|_{(\lambda),m'+|\alpha|}$, as $\partial_{y_B}=\sum_{a\in B}\partial_{x_a}$; so $|T_{\pi,\alpha}(g)|\le C\lambda^{\kappa+|\alpha|}(1+\log\frac{h_1}\lambda)^q\,n^{m'}\lambda^{-|\alpha|}\|F\|_{(\lambda),m'+|\alpha|}$, the power counting $\kappa$. A distribution $T'\in\mathfrak E_k(\kappa')$ with $\kappa'>3k$ is $0$: a test function on $\mathbb M^k$ is the sum of at most $C\lambda^{-3k}$ pieces carried by balls of radius $\lambda$, each of value $\le C\lambda^{\kappa'}$ up to logarithms, and $\lambda\to0$; this is the case $|\alpha|>\omega_k$, where $\kappa+|\alpha|>3k$.

The Lebesgue measure of $\mathbb M^k$ times a fixed nonzero element of the dual lies in $\mathfrak E_k(\kappa')$ exactly when $\kappa'\le3k$, that is, at $\kappa'=\kappa+|\alpha|$, exactly when $|\alpha|\le\omega_k$; this supplies a nonzero example for every allowed pair $(\pi,\alpha)$ when $V\ne\{0\}$. \emph{Only if:} a descent over $k=n-1,\dots,1$ for $T\in\mathfrak E_n(\kappa)$, with its numbers $C,q,m,h_1$, carried by $\Sigma_k$. Off $\Sigma_{k-1}$ the strata $\Delta_\pi$ with $|\pi|=k$ are disjoint; in the tube $U_\pi=\{|\xi|_\infty<\frac14\min\{\mathsf d_\infty(y),h_1\}\}$ about the open stratum, which meets no other of them, and with a cutoff $\chi$ carried by $\{|\xi|_\infty<1\}$ and equal to $1$ on $\{|\xi|_\infty\le\frac12\}$, the coefficients
\[t_{\pi,\alpha}(g)=(-1)^{|\alpha|}\,T\Bigl(\chi(\xi/\varepsilon)\,\frac{\xi^\alpha}{\alpha!}\,g(y)\Bigr),\qquad|\alpha|\le\omega_k,\ \ 0<\varepsilon\le\tfrac14\min\Bigl\{\min_{\operatorname{supp}g}\mathsf d_\infty,h_1\Bigr\},\]
do not depend on $\varepsilon$ and are distributions on $\mathbb M^k\setminus\Delta^{(k)}$ alone, $\Delta^{(k)}$ the diagonals of $\mathbb M^k$. \emph{Claim A:} if $G$ is carried by $U_\pi$ and $\partial_\xi^\alpha G(y,0)=0$ for $|\alpha|\le\omega_k$, then $T(G)=0$; for $(1-\chi(\xi/\varepsilon))G$ vanishes near $\Delta_\pi$ and is carried by $U_\pi$, which meets the support of $T$ in the open stratum alone, so that $T(G)=T(\chi(\xi/\varepsilon)G)$ for small $\varepsilon$, and a partition of unity at the scale $\varepsilon$ in $y$ cuts $\chi(\xi/\varepsilon)G$ into at most $C\varepsilon^{-3k}$ pieces, each $2\varepsilon$-supported with $\|\cdot\|_{(2\varepsilon),m}\le C\varepsilon^{(\omega_k+1)_+}$, so that $|T(G)|\le C\varepsilon^{\kappa-3k+(\omega_k+1)_+}$ times logarithms, which tends to $0$.

\emph{Claim B:} on $C^\infty_c(\mathbb M^n\setminus\Sigma_{k-1})$, by Taylor's formula in $\xi$ and Claim A, $T=\sum_{|\pi|=k}\sum_{|\alpha|\le\omega_k}\partial_\xi^\alpha\delta_\pi\otimes t_{\pi,\alpha}$. The coefficients have a power counting of their own: for cubes $R^*_B$ of the side $\frac32h$ at pairwise sup distances $\ge\frac52h$, $h\le\frac14h_1$, and $g$ carried by $\prod_BR^*_B$, the function $\chi(2\xi/h)\frac{\xi^\alpha}{\alpha!}g(y)$ is $2h$-supported with $\|\cdot\|_{(2h),m}\le Ch^{|\alpha|}\|g\|_{(h),m}$; so $t_{\pi,\alpha}$ is a functional of $k$ entries with $(\mathrm W;\kappa+|\alpha|,q,m)$, with derivatives of the order $m$ in the norms, which is why (a)--(c) are stated for every $r$.

For $k\ge2$, (a)--(c) at $r=m$ extend $t_{\pi,\alpha}$ to $T_{\pi,\alpha}\in\mathfrak E_k(\kappa+|\alpha|)$; for $k=1$ there are no diagonals in $\mathbb M^1$ and $T_{\pi,\alpha}=t_{\pi,\alpha}$ lies in $\mathfrak E_1(\kappa+|\alpha|)$ by the same inequality. By Claim B, $T-\sum_{|\pi|=k}\sum_{|\alpha|\le\omega_k}\partial_\xi^\alpha\delta_\pi\otimes T_{\pi,\alpha}$ vanishes off $\Sigma_{k-1}$, and it lies in $\mathfrak E_n(\kappa)$ by the first part; the descent ends at $\Sigma_0=\varnothing$. If $\kappa>3(n-1)$, all $\omega_k$ with $k\le n-1$ are negative and the sum is empty.

(f) At $n=2$ the one stratum has $k=1$, $y=x_1$, $\xi=x_2-x_1$ and $\omega_1=3$; the coefficients $T_\alpha$ are translation-invariant distributions on $\mathbb M$, hence constants $\mathsf c_\alpha\,dy$, and an integration by parts gives the form; conversely a constant lies in $\mathfrak E_1(\kappa')$ for $\kappa'\le3$. Covariance reads $\hat T(\mathsf Rp)(\rho(\mathsf R)H,\rho(\mathsf R)H')=\hat T(p)(H,H')$, and the exchange $\hat T(-p)(H',H)=\hat T(p)(H,H')$. Since $-1$ lies in both groups and $\rho(-1)=1$, $\hat T$ is even in $p$; so only the terms of the orders $0$ and $2$ in $p$ occur, with real coefficients, $\hat T(p)$ is symmetric in $(H,H')$, and the number of the independent terms is that of the invariants of $\mathrm{Sym}^2V^*\otimes\mathrm{Sym}^j(\mathbb R^3)^*$, $j=0,2$. Under $\mathrm{SO}(3)$, $V=[0]\oplus[2]$, $\mathrm{Sym}^2V=2[0]\oplus2[2]\oplus[4]$ and $\mathrm{Sym}^2\mathbb R^3=[0]\oplus[2]$: $2+4=6$; the six forms are invariant. Separating the homogeneous terms of degrees $0$ and $2$, their independence follows by evaluating at $H=H'\in\{\operatorname{diag}(1,-1,0),\operatorname{diag}(1,0,0)\}$ and $p\in\{e_1,e_2,e_3\}$ within each degree. Under the cubic group $V=A_1\oplus E\oplus T_2$, $\mathrm{Sym}^2V=3A_1\oplus3E\oplus3T_2\oplus T_1$ and $\mathrm{Sym}^2\mathbb R^3\cong V$: $3+9=12$; the characters give the same numbers, and $0$ for the odd orders.
\end{proof}

\begin{remark}[The point bound and the bounds of kernels, for data]\label{C-355-35}
For reconstruction data with labels the point bound gives pointwise bounds of the kernels by powers of the least distance of the points (Proposition~\ref{C-355-25} and Theorem~\ref{C-355-10}~(d)~(i)). The converse fails: at $V=\mathbb R$, $\sigma(f_1,\dots,f_n)=n!\prod_k\int f_k$ is symmetric and has constant kernels, hence every bound of the kind of (E1), but no bound of one insertion between two vectors. The reconstruction uses nothing at coinciding points.
\end{remark}

\subsection{Proofs of the theorems}

\begin{proof}[Proof of Theorem~\ref{C-355-7}]
\emph{Step 1: reduction.} It suffices to treat real $H$; for complex $H$ the real and the imaginary part are treated separately, at the price of a factor $2$. By Lemma~\ref{C-355-18}~(a),
\[\|\mathcal T_K(H)-c\|^{(i,o)}_{{\rm ch},\ell}\ \text{at }(K,\beta)\ =\ \|\mathcal T_k(H')-c\|^{(0,L^io)}_{{\rm ch},L^i\ell}\ \text{at }(k,\beta),\qquad|H'|_\infty/b'=|H|_\infty/b,\]
so the seminorm is taken at the stopped datum $(k,\beta)$, $k=K-i$, at the coupling $g'=g_{(i)}\le\gamma^{\rm sd}_W$, on a board of side $1$, for $H'=L^iH(L^{-i}\cdot)$, which vanishes within sup distance $\delta$ of the walls of its cell; the label is $b'=L^ib$, the least side $\ell''$ at $g'$ is small over the running couplings of $(k,\beta)$, and $s=s(\delta)\le s_E(i)\le k$, so that at the working level $a_k\le\delta/C_W(s)\le\delta/2197$.

\emph{Step 2: pieces.} The points of the cell at sup distance $\ge\delta$ from its walls form a cube of side $1-2\delta$, a union of at most $2\delta^{-3}$ boxes of sides $\le\frac{99}{100}\delta$, each in the ball of radius $0.8574\,\delta$ about its centre. Restricting $w_{H'}$ and $\mathsf w_{H'}$ to the plaquettes and bonds with barycentres and midpoints in a box $\mathfrak b$,
\[\mathcal T_k(H')=\sum_{\mathfrak b}\bigl[\mathcal O(w_{\mathfrak b})+Y[\mathsf w_{\mathfrak b}]\bigr],\]
a sum of at most $4\delta^{-3}$ functions of the bonds of the cell.

\emph{Step 3: one piece and its copies.} Let $F$ be one of these pieces, taken on $T_{\ell''}$, whose board has $N=\ell''^3$ cells, and $F_\kappa$ its mirror copies. By Lemma~\ref{C-355-21} each $F_\kappa$ is the first member of a source under $(\mathsf s,S_\kappa)$ with $\mathsf s\le59|H'|_\infty$. The barycentres of the plaquettes of $S_\kappa$ lie within $0.8574\,\delta+0.71\,\delta/2197\le\frac78\delta$ of the centre of the copy of the box; two points in different cells, each at sup distance $\ge\delta$ from the walls, are parted on the circle by a wall on either side, so their distance is $\ge2\delta$, and the barycentres of different copies are at distance $\ge2\delta-\sqrt2a_k\ge\frac54\delta$.

\emph{Step 4: the moments of the copies.} For $i\ge i_E$ the thresholds (T1) at $n=N$ and (T2) hold (Lemma~\ref{C-355-18}~(d)), and $N=\ell''^3$ is admissible in Proposition~\ref{C-355-19}~(a), whose condition $n\le\ell^3$ is exactly the number of cells. So there are $c_\kappa\in\mathbb C$ with
\[\Bigl|\Bigl\langle\prod_{\kappa\in B}(F_\kappa-c_\kappa)\Bigr\rangle_{k,\ell'',\beta}\Bigr|\le(\mathsf A_N(s)\,\upsilon\,\mathsf s)^{\#B}\qquad\text{for every non-empty set }B\text{ of cells.}\]

\emph{Step 5: one constant.} The torus mean is invariant under the diagonal isometries, so all the copies have one mean $m_F$, and $|m_F-c_\kappa|\le\mathsf A_N(s)\upsilon\mathsf s$ by Step~4 at $B=\{\kappa\}$. Expanding about the $c_\kappa$,
\[\|F-m_F\|^N_{\ell''}=\Bigl\langle\prod_\kappa(F_\kappa-m_F)\Bigr\rangle=\sum_B\Bigl\langle\prod_{\kappa\in B}(F_\kappa-c_\kappa)\Bigr\rangle\prod_{\kappa\notin B}(c_\kappa-m_F)\le(2\mathsf A_N(s)\upsilon\mathsf s)^N .\]

\emph{Step 6: the sum, and every side.} By subadditivity and $\upsilon\le(9/8)^{1/2}/b'$, $\|\mathcal T_k(H')-m\|_{\ell''}\le2\delta^{-3}\cdot2\mathsf A_N(s)\upsilon(1+59)|H'|_\infty\le255\,\delta^{-3}\mathsf A_N(s)|H|_\infty/b$, $m=\sum m_F$. For $\ell\in\mathfrak L(g)$, $L^i\ell\in\mathfrak L(g')$ and $L^i\ell\ge\ell''$; Lemma~\ref{C-355-17} at $g'$, with cells of side $1$, gives
\[\|\cdot\|_{L^i\ell}\le e^{G_{g'}}\|\cdot\|_{\ell''},\qquad G_{g'}\le G_g+\tfrac92i\log L+\tfrac18 ,\]
so $\|\mathcal T_K(H)-m\|_\ell\le\mathsf E_ib^{-1}|H|_\infty$ with $\mathsf E_i=510\,\delta^{-3}e^{G_g+1/8}L^{9i/2}\mathsf A_N(s)$ for complex $H$.

\emph{Step 7: the centrings.} The mean on a torus $T_{\ell_1}$ differs from $m$ by at most the seminorm on $T_{\ell_1}$; the mean in $\nu\in\mathcal G^{\rm per}_{K,\beta}$ differs from $m$ by at most the lower limit of the seminorms, by the domination along the sides $\ell_k$ of $\mathfrak L(g)$:
\[|\nu(\mathcal T_K(H))-m|\le\liminf_k\|\mathcal T_K(H)-m\|_{\ell_k}\le\mathsf E_ib^{-1}|H|_\infty ;\]
with the triangle inequality, $\|\mathcal T_K(H)-c\|_\ell\le2\mathsf E_ib^{-1}|H|_\infty$.

\emph{Step 8: the numbers.} Lemma~\ref{C-355-18}~(c) in the formula of $\mathsf A_N(s)$ gives
\[
\begin{multlined}
N\max\{C^I,C_2N^{3+\mathfrak p_0}\}\le\max\{C^I,C_2\}(1+c'A_0^2/9)^{4+\mathfrak p_0}\mathsf Y_i^{(4+\mathfrak p_0)(2\mathfrak p_0-1)},\quad(LC_W(s))^3 \\
\le L^3c_W^3\mathsf Y_i^{6\mathfrak p_0},\quad\bar x_s\le L^2c_W\delta^{-1}\mathsf Y_i^{2\mathfrak p_0}(X_g+\bar B)L^i,
\end{multlined}
\]
and $Y_s^{2\mathfrak p_0}\le\mathsf Y_i^{2\mathfrak p_0}$; hence $\mathsf E_i\le C_E\mathsf Y_i^{q_E}L^{11i/2}$, which is \eqref{C-355-8}. Of the power $\frac{11}2$, $\frac92$ is the increment of the pressure gap of one unit cell at the stopped coupling, and $1$ is the crude size $\bar x_s\sim L^i$ of one insertion at the working level. With $\mathsf x_E(\delta)=\mathsf x^*/\delta$, $\mathsf x^*$ formed from the constants alone (Lemma~\ref{C-355-18}~(d)), $L^{i_E}\le c_i\delta^{-1}$, and $\mathsf Y_i^{q_E}\le(2q_E)^{q_E}c_0\delta^{-1}L^i$ with $c_0=L^2c_W(4\mathfrak p_0)^{2\mathfrak p_0}(X_g+\bar B)$; so $C_E\mathsf Y_i^{q_E}L^{11i/2}\le c_3\delta^{-5}L^{13i/2}$, $C_E$ being $\delta^{-4}$ times a number free of $\delta$.

\emph{Step 9: the second form (b).} The same steps with Proposition~\ref{C-355-19}~(b) in the place of (a), at the stopped datum $(K-i,\beta)$ and at $s=s_i$, so that the level $k=K-i-s_i\in[2,K-i-1]$ is chosen by the distance alone and the depth $s(\delta)$ plays no part. The whole cell is one source, and by (Sd.1) and the definition of the admissible sides the least torus at $g_{(i)}$, of side $\ell''\le2L^2MC_{\rm s}(\frac12\mathsf Y'_i)^r$, has few cells:
\[N=\ell''^3\le c_N\mathsf Y_i'^{\,3r},\qquad\delta L^{s_i}\ge Mu'_*\bigl(\lceil c_N\mathsf Y_i'^{\,3r}\rceil\bigr)\ge Mu'_*(N);\]
two corners of plaquettes of different copies are at sup distance $\ge\delta$, that is $\ge\delta L^{s_i}$ in the units of the level $k$. No threshold on $i$ is then needed, and the pressure gap enters as in Step~6. As $u'_*(\lceil c_N\mathsf Y_i'^{\,3r}\rceil)\le c_u\mathsf Y_i'^{\,3r+1}$ and, by the minimality of $s_i$, $L^{s_i}\le(LMc_u/\delta)\mathsf Y_i'^{\,3r+1}$, this gives $C'_E$ and $q'_E=15r+4$.

\emph{Step 10: all couplings $\le\gamma$ and all scales (c).} (i) By Lemma~\ref{C-355-23}~(a) the datum stopped at $z=z(g)$ is a first-passage datum at $\tilde g\le\gamma^{\rm sd}_W$, the board $(i,o)$ becomes a board of side $a_{i-z}$, and (a) at $\tilde g$, $\tilde\imath=i-z\ge i_E(\tilde g,\delta)$, with the monotony of the numbers gives
\[\|\mathcal T_K(H)-c\|_\ell\le2C_E[\tilde g]\mathsf Y_{\tilde\imath}[\tilde g]^{q_E}L^{11\tilde\imath/2}(L^zb)^{-1}L^z|H|_\infty\le2e^{1/8}C_E[g]\mathsf Y_i[g]^{q_E}L^{11i/2}b^{-1}|H|_\infty .\]
(ii) For $i\ge i_A(g,\frac1{32})$ this is (i) at $\delta=\frac1{32}$; for the finitely many smaller $i$, Lemma~\ref{C-355-23}~(c) with $i'=\max\{i+2,i_A(g,\frac18)\}$ and (i) at $\delta=\frac18$ on the eight shifted fine boards give $\mathsf N_i=8(L^{i'-i}+1)^3$ times the number of (i) at $i'$. Since $\mathsf N_i\le C\mathsf Y_i^{q_E}L^{11i/2}$ for large $i$ and $b^{-1}<g^2$, $C_T<\infty$ for every $p_T>\frac{11}2$. (iii) Let $i=i^A$. Cut each $H_a$ by a smooth partition of unity of mesh $a_i/(8n)$ into at most $(16nR/a_i+3)^3$ pieces of sup diameter $\le a_i/(4n)$; by multilinearity it suffices to bound one choice of a piece for each entry. Lemma~\ref{C-355-22} at the mesh $a_K$ (possible since $L^{K-i}\ge8n$) gives one board of side $a_i$ in which each of the $n$ pieces lies in a cell at the margin $a_i/(8n)=\delta_na_i$, pieces of different entries in different cells, since two points of one cell are at distance $\le\sqrt3a_i<d$. The domination of the state $\nu$, with (i) at $\delta_n$ for each piece centred by its mean in $\nu'$, bounds the product of the pieces by $\prod_a2e^{1/8}C_E(\delta_n)\mathsf Y_i^{q_E}a_i^{-11/2}b^{-1}|h_a|_\infty$; the sum over the pieces gives $C^A_n(d;R)$.
\end{proof}

\begin{proof}[Proof of Theorem~\ref{C-355-9}]
\emph{Step 1: the uniform bound (a).} Let $i=i^G$, so that $2a_i\le d$. Cut each $H_a$ by a smooth partition of unity of mesh $a_i/(8n)$ into pieces of sup diameter $\le a_i/(4n)$, at most $(16nR/a_i+3)^3$ of them per entry; by multilinearity it suffices to bound one choice of a piece for each entry. Lemma~\ref{C-355-22}, at a mesh of the lattice (possible since $L^{K-i}\ge8n$), gives one board of side $a_i$ in which each of the $n$ pieces lies in a cell at margin $a_i/(8n)=\delta_na_i$; since two points of one cell are at distance $\le\sqrt3a_i<d$, the pieces of different entries lie in different cells. The domination and \eqref{C-355-8} at $\delta=\delta_n$, centred in the state $\nu$, give
\[\Bigl|\nu\Bigl(\prod_a\bigl[\mathcal T_K(h_a)-\nu(\mathcal T_K(h_a))\bigr]\Bigr)\Bigr|\le\prod_a2C_E(\delta_n)\mathsf Y_i^{q_E}a_i^{-11/2}b^{-1}|h_a|_\infty ;\]
no order of the cells is used. On a torus the chessboard estimate takes the place of the domination. For $d\le2a_{i_E}$, $a_{i^G}\ge d/(2L)$.

\emph{Step 2: extraction and the properties (b).} Along a countable family of tuples, dense for the uniform norm on every product of compact sets, the bound (a), uniform in the datum and the state, gives bounded sequences, and a diagonal subsequence converges; multilinearity and (a) give the continuity under the motion of the entries, which carries the convergence to all tuples; the labels lie in the compact interval $[g^{-2},L(g^{-2}+\bar B)]$ by \YMMref{(1531)}{615}. Reality, symmetry and multilinearity pass to the limit, and $\sigma((H))=0$ by the centring. The periodic Gibbs states are invariant under the isometries of the lattice and reflection positive at its mirrors (Proposition~\YMMref{103.4}{758}), the plaquette and bond forms being carried to their images with the signs of Lemma~\ref{C-355-21}; as $a_K\to0$ these isometries approach all of $\mathbb R^3\rtimes W_3$, and the entries are continuous. The five-point inequality, Lemma~\ref{lem:spectral-steps}, from the positivity at the diagonal mirrors,
\[\mathsf S(0)-4\mathsf S(te_3)+\mathsf S(2te_3)+\mathsf S(te_3+te_k)+\mathsf S(te_3-te_k)\ge0,\]
gives $\mathcal P(\{E<\max_k|q_k|\})=0$. $\mathcal H_0$ is invariant under $e^{-tH}$ and $\mathcal U$. The basic inequality comes from the multiple reflections in the periodic Gibbs states (Proposition~\YMMref{104.2}{759}, using the domination inequality~\YMMref{(1832)}{705}),
\[|\rho_j(\overline{\theta_cA'}\,D\,\tau_{a_ie_3}A)|\le\liminf_\ell\|D\|_\ell\,\|A'\|\,\|A\|,\]
for $D$ the centred $\mathcal T_{K_j}(h)$ in a cell over the mirror, whose seminorms \eqref{C-355-8} at $\delta=\frac1{16}$ bounds uniformly in the side (the support is at sup distance $\ge a_i/8-a_{K_j}/2\ge a_i/16$ from the walls); letting $j\to\infty$ gives $|\sigma(\theta\bar F'\cdot(h)\cdot\tau_{a_ie_3}F)|\le\mathsf c_E(i)b_*^{-1}|h|_\infty\|[F']\|\,\|[F]\|$.

\emph{Step 3: the point bound (c).} Let $i=i(\vartheta)$ be the least $i\ge i_1$ with $a_i\le\vartheta$, $\rho=a_i/32$ and $\varepsilon=12\rho$. Cut $\tau_{de_3}h$ by the tent partition of unity on $\rho\mathbb Z^3$; for a piece carried by $Q_z$ put $\mathsf b(g)=\mathsf B_{a_i}(\tau_{-ce_3}g)$, $c=z_3-a_i/2$ and $c'=2d-z_3-a_i/2$. Since $Q_z$ meets the support of $\tau_{de_3}h$, both gaps are at least $\vartheta-a_i/2-\rho\ge15a_i/32>\varepsilon$. Thus $\|\mathsf b(g)\|\le\mathsf c_E(i)b_*^{-1}|g|_\infty$ by Step~2; the relations of the slab operators give the covariance asked in Lemma~\ref{C-355-24}, and $\mathsf B_{2d}(g)=e^{-(c-\varepsilon)H}\mathsf B^\varepsilon(g)e^{-(c'-\varepsilon)H}$ with two contractions. Summing over the pieces,
\[\|\mathsf B_{2d}(\tau_{de_3}h)\|\le4\cdot32^3a_i^{-3}\mathsf c_E(i)\,b_*^{-1}\|h\|_{L^1},\qquad a_i^{-3}\mathsf c_E(i)=2C_E\mathsf Y_i^{q_E}a_i^{-17/2}\le C_p\min\{1,\vartheta\}^{-p}\ (p>\tfrac{17}2).\]
For $(\mathrm{Rec4})_V$: $\sigma(\theta\bar F\cdot(h)\cdot F')=\langle[\tau_{-\eta e_3}F],\mathsf B_{2\eta}(\tau_{\eta e_3}h)[\tau_{-\eta e_3}F']\rangle$ with $\vartheta>\eta/2$, and $|A|_\infty\le|A|$; $b_*\ge g^{-2}$ gives the one triple. The price of the passage to $L^1$ is $a_i^{-3}$: $\frac{17}2=\frac{11}2+3$.

\emph{Step 4: real time (d).} The proofs of Proposition~\ref{C-355-25}~(a) that give the semigroup on the tube, the insertion and point operators, the ordered tuples, $\mathbf W_n$, $\Psi_n$, the fields and their boundary values use, of the isometries, the translations and $\theta$ alone, besides $(\mathrm{Rec2})_V$ and $(\mathrm{Rec4})_V$; they use the cone at five places, and each remains valid with $d_\square$ and the square cone, because on it
\[|\operatorname{Im}\mathfrak z_1\,q_1+\operatorname{Im}\mathfrak z_2\,q_2|\le(|\operatorname{Im}\mathfrak z_1|+|\operatorname{Im}\mathfrak z_2|)\,E :\]
the bound $|e^{-\mathfrak z_3E+i(\mathfrak z_1q_1+\mathfrak z_2q_2)}|\le e^{-d_\square(\mathfrak z)E}$, the estimate $|-\zeta_3E+i(\zeta_1q_1+\zeta_2q_2)|\le2|\zeta|E$ for the holomorphy, $\mathcal P(\{0\}\times\mathbb R^2)=\mathcal P(\{(0,0)\})$, the bound $Ee^{-\eta E/4}$ of the derivatives of the translations, and the convexity of the square tube. The invariance under $W_3$ gives the covariance under the signed permutations fixing $e_3$. For (vi), $\mathcal H_0$ reduces $H$, $\|\mathsf A_l(E)\|\le\|\mathsf A_l\|\,|E|$, $\overline{q_0(\psi,\psi';e)}=q_0(\psi',\psi;e^*)$ with $E^*=\rho(\theta)E$ for a constant real label, so that $\mathsf m_{ij}^*=\vartheta_i\vartheta_j\mathsf m_{ij}$ and $|\mathsf m_{ij}|\le\sqrt3$, and $q_0(\Omega,\Omega;E)$ is an integral of $\mathbf W_1=0$, since $\sigma((h))=0$; the operator formula of (M2) is the product $\mathsf b(x+ae_3|\mathsf m_{ij}|\frac12(y-x))\,\mathsf b(\frac12(y-x)|\mathbf 1|-y+ae_3)$ at $x_3<y_3$, between the vectors $[F']$, $[F]$.

\emph{Step 5: the circular cone (e).} A tensor limit is reflection positive at the coordinate mirrors and at the diagonal mirrors $\{x_\mu\pm x_\nu=c\}$ of the lattice. Lemma~\ref{C-355-26}~(b) at $u=e_3$, $u'=e_1$, $u''=e_2$, with the four mirrors of normals $e_3$, $(e_3\pm e_1)/\sqrt2$, $e_1$ through the axis $\mathbb Re_2$ and the fifth with normal $(e_3+e_2)/\sqrt2$, all of them lattice mirrors, gives the circular cone in the picture $e_3$ (the case $\beta=0$), and likewise in each of the nine pictures, the axis of the four mirrors being the coordinate axis normal to the plane that holds the normal of the picture, and the fifth mirror a diagonal one; in each picture $u$ the contractions $\mathcal T_u(\mathfrak z)$ on the circular tube follow.

In the picture $e_3$ the product formula for $\mathbf W_n$ written in (e), with these contractions, gives the continuation and its bound with $C_{\rm pt}(\varepsilon)=3\sqrt2\,C_V\min\{1,\varepsilon\}^{-p}$; the invariance under $W_3$ carries them to the pictures $e_1$, $e_2$. \emph{The boosts.} $\sigma$ is invariant under the quarter turn $\mathsf qx=(x_3,x_2,-x_1)$ in the plane of $e_1$, $e_3$, an element of $W_3$; it is reflection positive at the mirror $\{x_3=0\}$ and at the diagonal mirror $\{x_1=x_3\}$, of the normal $n_-=(e_3-e_1)/\sqrt2$, whose reflection is $\mathsf q^{-1}\theta$; and $H\ge|P_1|$, so that $H_\pm=H\pm P_1\ge0$. Let $\mathcal C_-$, $\mathcal C_+$ be the sets of the classes of the tuples carried by the quadrants $\{x_1<0,\ x_3>0\}$ and $\{x_1>0,\ x_3>0\}$.

Both are total: if $\psi\perp\mathcal C_-$, then for $F\in\mathbb C\mathfrak S_V(\mathbb E_+)$ and $\varepsilon>0$ the function 
\[
v\mapsto\langle\psi,\mathcal U(ve_1)e^{-\varepsilon H}[F]\rangle=\langle\psi,[\tau_{(v,0,\varepsilon)}F]\rangle
\]
 continues holomorphically to $|\operatorname{Im}v|<\varepsilon$, because $|P_1|\le H$, and vanishes for $v$ far to the left, where $\tau_{(v,0,\varepsilon)}F$ is carried by the quadrant; so it vanishes, and $\varepsilon\downarrow0$ gives $\psi\perp[F]$; likewise for $\mathcal C_+$. For $F,F'$ carried by $\{x_1<0,\ x_3>0\}$, $\mathsf qF$ is carried by $\{x_1>0,\ x_3>0\}$, and, $[\cdot]_{n_-}$ denoting the classes of the picture $n_-$, the invariance of $\sigma$ under $\mathsf q$ gives
\[\langle[F']_{n_-},[F]_{n_-}\rangle=\sigma(\mathsf q^{-1}\theta\bar F'\cdot F)=\sigma(\theta\bar F'\cdot\mathsf qF)=\langle[F'],[\mathsf qF]\rangle .\]
So the map $I:[F]\mapsto[F]_{n_-}$ on the span of $\mathcal C_-$ is well defined, its form $\|I\psi\|^2$ is the form of the positive symmetric operator $[F]\mapsto[\mathsf qF]$, and the form is closable: $I$ is closable, and with its closure $\bar I$, $D=\bar I^*\bar I$ is the Friedrichs extension of that operator, a positive self-adjoint operator with $D[F]=[\mathsf qF]$ on $\mathcal C_-$ and $D\Omega=\Omega$; $\ker D=(\operatorname{ran}D)^\perp=0$, since $\operatorname{ran}D\supset\mathcal C_+$. For $x=(-a,0,t)$, $a,t\ge0$, the translation by $x$ maps $\{x_1<0,\ x_3>0\}$ into itself; $\mathcal T(x)=e^{-tH}\mathcal U(-ae_1)$ is the operator of this translation, and $\bar I\,\mathcal T(x)\supset\mathcal T^{n_-}(x)\,\bar I$ with the translation $\mathcal T^{n_-}(x)$ of the picture $n_-$, a contraction since $n_-\cdot x\ge0$; taking adjoints in this relation at $x'=-\theta_{n_-}x=(-t,0,a)$, which is of the same kind, and using $\mathcal T^{n_-}(x')^*=\mathcal T^{n_-}(x)$ and $\mathcal T(x')^*=\mathcal T(\mathsf qx)$, one gets $D\,\mathcal T(x)\supset\mathcal T(\mathsf qx)\,D$; in the same way $D$ commutes with $\mathcal U(ve_2)$.

With $\Xi(\alpha,\beta)=\exp(-\frac\alpha2H_+-\frac\beta2H_-)$, $\operatorname{Re}\alpha,\operatorname{Re}\beta\ge0$, and $z=t+ia$, this reads $D\,\Xi(z,\bar z)\supset\Xi(-iz,i\bar z)\,D$ for $z$ in the closed first quadrant. For $\phi,\psi\in\operatorname{dom}D$ both members of $\langle D\phi,\Xi(\alpha,\beta)\psi\rangle=\langle\phi,\Xi(-i\alpha,i\beta)D\psi\rangle$ are holomorphic for $\alpha$ in the open first and $\beta$ in the open fourth quadrant, continuous up to the boundary, and they agree on the totally real set $\{\beta=\bar\alpha\}$; so they agree everywhere, and letting $\beta\to0$, $D\,e^{-\alpha H_+/2}\supset e^{i\alpha H_+/2}D$ for $0\le\arg\alpha\le\frac\pi2$. For $\phi,\psi$ in the spectral subspace of $\log D$ of a bounded set and $r>0$, the function
\[h(z)=\bigl\langle D^{-\bar z}\phi,\exp\bigl(-\tfrac r2e^{-i\pi z/2}H_+\bigr)D^z\psi\bigr\rangle\]
is holomorphic and bounded on $|\operatorname{Re}z|\le1$, as $\operatorname{Re}e^{-i\pi z/2}\ge0$ there, and the last relation at $\alpha=re^{-i\pi z/2}$ gives $h(z+1)=h(z)$ for $-1\le\operatorname{Re}z\le0$; so $h$ continues to a bounded entire function and is constant. At $z=is$ this says $D^{is}e^{-rH_+/2}D^{-is}=e^{-re^{\pi s/2}H_+/2}$ for all $r>0$, that is $D^{is}H_+D^{-is}=e^{\pi s/2}H_+$; in the same way, from the relation for $H_-$ with $\beta$ in the closed fourth quadrant, $D^{is}H_-D^{-is}=e^{-\pi s/2}H_-$, and $D^{is}$ commutes with $P_2$. So $V_1(s)=D^{is}$ is a strongly continuous unitary group with $V_1(s)\Omega=\Omega$ and $V_1(s)\mathcal P(B)V_1(s)^{-1}=\mathcal P(\Lambda^{(1)}_{-\pi s/2}B)$, where the boost $\Lambda^{(1)}_\varsigma$ multiplies $E\pm q_1$ by $e^{\pm\varsigma}$ and keeps $q_2$; $V_2$ is its conjugate by the exchange of $e_1$ and $e_2$, an element of $W_3$ that fixes $e_3$, and $\Lambda^{(2)}_\varsigma$ multiplies $E\pm q_2$ by $e^{\pm\varsigma}$. Two boosts in independent directions generate $\mathrm{SO}_0(1,2)$, so the null sets and the support of $\mathcal P$ are invariant under it.

\emph{Step 6: ordered classes and kernels (f), (g).} The suppositions of Proposition~\ref{C-355-27} hold by (a)--(c); its clause (a) gives $\mathcal H^{\rm ord}(\sigma)=\mathcal H(\sigma)$, and the density of $\mathcal D_0$ in $\mathcal H^{\rm ord}(\sigma)$ follows as in Proposition~\ref{C-355-25}: a vector orthogonal to $\mathcal D_0$ is orthogonal to the holomorphic vectors $\Psi_n(z;e)$, which span a dense subspace of $\mathcal H^{\rm ord}$. Its clause (b) gives (g) but for its last sentence; the bound for tuples in balls follows from the bound of the kernels with
\[\|h_a\|_{L^1}\le\tfrac{4\pi}3R^3\sup|h_a|,\qquad|x|_\infty\ge|x|/\sqrt3 .\]
The last sentence of (g) is clause (c), whose supposition holds: the restriction of $\sigma$ to the tuples with entries in $V_0$ is the limit, along the sequence $(K_j,\beta_j,\rho_j)$ of first-passage data at $g$ and periodic Gibbs states that defines $\sigma$, of the centred products of the action densities $\mathcal T_{K_j}(f\mathbf 1)=\mathcal O(f)$ (Definition~\ref{C-355-3}). This sequence is a Gibbs sequence at $g$ in the sense of Definition~\ref{C-365-5}(2), and it converges at the scalar tuples with $C^1_c$ entries, which are the tuples among the construction results; by Theorem~\textup{\YMMref{AP}{780}} its limit, the scalar part of $\sigma$ on $C^1_c$ entries, is invariant under all the isometries of $\mathbb R^3$ and is scalar reconstruction data, (Rec) with the numbers $16\,C^g_{\rm pt}$, $4$, $\frac12$ of that theorem. Of the rotations outside $W_3$, this reads the invariance of the scalar limit alone, on no tensor entry.

\emph{Step 7: growth (h).} Let $i$ be as in (h), so that $a_i\le d/4$; as $L^{i_E(g,\delta_n)}\le c_i\delta_n^{-1}=8nc_i$ (Theorem~\ref{C-355-7}~(a)) and $L^{i(d)}\le4L/d$ for the least $i(d)$ with $a_{i(d)}\le d/4$, and $c_i$ may be taken $\ge1$,
\[L^i\le\max\{8nc_i,\,4L/d\}\le8c_iLn/d,\qquad\mathsf E_i(\delta_n)\le c_3(8n)^5L^{13i/2}.\]
Cut each $H_a$ by a smooth partition of unity of mesh $a_i/(8n)$ into pieces of sup diameter $\le a_i/(4n)$, at most a number times $(n/a_i)^3$ of them, the supports lying in balls of radius $1$. Lemma~\ref{C-355-22} at $h=a_K$, possible as $8na_K\le a_i$, puts the $n$ pieces of one choice, one of each entry, on one board of side $a_i$ in cells at the margin $a_i/(8n)=\delta_na_i$, and in $n$ different cells, since two points of one cell are at distance $\le\sqrt3a_i<d$; the domination with \eqref{C-355-8} at $\delta_n$ bounds the centred product of the pieces in any two states by $(2\mathsf E_i(\delta_n)b^{-1})^n\prod|h_a|_\infty$, without a factorial and without an order of the cells. One entry thus contributes at most a number times $n^3L^{3i}\cdot n^5L^{13i/2}=n^8L^{19i/2}\le c\,n^{35/2}d^{-19/2}$; this is the first bound of (h), and with $n^{35n/2}\le e^{35n/2}(n!)^{35/2}$ and $b_*^{-1}\le g^2$ it gives the bound of $\sigma$, for every tensor limit and, by Theorem~\ref{C-355-13}~(b)~(i), for every component.

\emph{The finer count.} By the formulas of Theorem~\ref{C-355-7}~(a), $C_E(\delta_n)=(8n)^4C_E(1)$, and $\mathsf Y_i=2[c_E(\delta_n)+i\log L]\le c\,(1+\log(n/d))$, as $c_E(\delta_n)=\log(8n)+c_E(1)$ and $L^i\le8c_iLn/d$; so $\mathsf Y_i^{q_E}\le c\,(n/d)^{1/2}$, in the place of the crude bound by a number times $\delta_n^{-1}L^i$ of Step~8 of the proof of Theorem~\ref{C-355-7}, and $\mathsf E_i(\delta_n)\le C_E(\delta_n)\mathsf Y_i^{q_E}L^{11i/2}\le c\,n^4(n/d)^{1/2}L^{11i/2}$. For entries of any size at pairwise sup distances $\ge2d$, cut each $H_a$ into its restrictions to the boxes $\mathfrak b=\frac{a_i}{4n}(\mathbf n+[0,1)^3)$, $\mathbf n\in\mathbb Z^3$, of the mesh $a_i/(4n)\le d/(16n)$; the pieces of two different entries are at sup distance $\ge2d>a_i$, hence in different cells of the board of Lemma~\ref{C-355-22}; and $\sum_{\mathfrak b}\sup_{\mathfrak b}|H_a|_\infty=(4n/a_i)^3\|H_a\|_\natural$ with the upper Riemann sum $\|H\|_\natural=\sum_{\mathfrak b}|\mathfrak b|\sup_{\mathfrak b}|H|_\infty$. One entry then contributes
\[2\mathsf E_i(\delta_n)b^{-1}(4n/a_i)^3\|H_a\|_\natural\le c\,n^7(n/d)^{1/2}L^{17i/2}\,b^{-1}\|H_a\|_\natural\le c'\,n^{16}d^{-9}\,b^{-1}\|H_a\|_\natural ,\]
which, in the limit along the data of $\sigma$ and with $b_*^{-1}\le g^2$, is the bound of (h) with $\|\cdot\|_\natural$, ordered or not. Entries in balls of radius $1$ at Euclidean distances $\ge d$ are at sup distances $\ge d/\sqrt3=2d'$, and the bound at $d'$ with $\|H_a\|_\natural\le27|H_a|_\infty$ gives $|\sigma(X)|\le(c\,g^2n^{16}d^{-9})^n\prod_a|H_a|_\infty$ for $0<d\le1$.

\emph{The extension.} Proposition~\YMMref{93.3}{717}, applied to each of the $6^n$ components of the labels with a bound $\mathsf c_{\mathsf T}\max\{1,d^{-\mathsf a}\}\prod_a\|f_a\|_\sharp$ at radius $\le1$ and distances $\ge d$, gives exactly one functional, bounded for $|\cdot|_{8n+\mathsf a-1}$ with the factor $C^n(n!)^8\mathsf c_{\mathsf T}$. With the first bound, $\mathsf a=\lceil19n/2\rceil$ and $\mathsf c_{\mathsf T}=(c_4g^2e^{35/2})^n(n!)^{35/2}$: the norm $|\cdot|_{18n}$ and the factorial power $\frac{51}2$. With the finer bound, $\mathsf a=9n$ and $\mathsf c_{\mathsf T}=(c\,g^2n^{16})^n\le(c\,g^2e^{16})^n(n!)^{16}$: the norm $|\cdot|_{17n}$ and the factorial power $24\le\frac{49}2$, so that the pair $(17n,\frac{49}2)$ serves.

The pair $(15n,\frac{43}2)$ comes from an extension that reads the bound on small cubes itself. Let $i_n=i_E(g,\delta_n)$, and for an integer $i\ge i_n$ let $P_1,\dots,P_n$ be closed cubes of the side $a_i/(4n)$, with edges along the axes, at pairwise sup distances $\ge a_i$, and $h_a$ continuous, with values in $V_{\mathbb C}$ and carried by $P_a$. At the cutoffs with $K\ge K_E(i)$ (at $\delta_n$) and $8na_K\le a_i$, Lemma~\ref{C-355-22} at $h=a_K$ puts the $P_a$ on one board of side $a_i$, in cells at the margin $\delta_na_i$ and in $n$ different cells, since two points of one cell at this margin are at sup distance $<a_i$; the domination with \eqref{C-355-8} at $\delta_n$, at complex entries as in (a), and $b_*^{-1}\le g^2$ give, for $\sigma$ and, by Theorem~\ref{C-355-13}~(b)~(i), for each of its components,
\[|\sigma(h_1,\dots,h_n)|\le\bigl(\mathsf u\,L^{\mathsf vi}\bigr)^n\prod_a|h_a|_\infty,\qquad\mathsf u=2c_3g^2(8n)^5,\quad\mathsf v=\tfrac{13}2 .\]
With an orthonormal basis $(\varepsilon_\alpha)_{1\le\alpha\le6}$ of $V$ for the trace form, $|\varepsilon_\alpha|_\infty\le|\varepsilon_\alpha|=1$; so each of the $6^n$ components of the labels, $\Lambda(f_1,\dots,f_n)=\sigma(f_1\varepsilon_{\alpha_1},\dots,f_n\varepsilon_{\alpha_n})$, continued linearly to complex $f_k$, has the same bound for every $i\ge i_n$ and all $f_k$ carried by such cubes. A multilinear $\Lambda$ with this bound has exactly one extension $\mathsf S$, bounded for one Schwartz norm, to the complex Schwartz functions on $\mathbb R^{3n}$ vanishing with all their derivatives on the diagonals, and
\[|\mathsf S(\varphi)|\le2\bigl(C_*\,\mathsf u\,n^7L^{(\mathsf v+3)i_n}\bigr)^n\,|\varphi|_{(\lceil\mathsf v\rceil+8)n},\]
$C_*$ a number formed from $\mathsf v$ and two cutoff functions of one variable fixed once. Differences of products over the pairs $a<b$ of cutoffs in $|x_a-x_b|$ at the scales $L^{-s}$, $s\ge0$, cut the non-coinciding configurations into shells on which the least distance of the points lies between $L^{-s}$ and $\sqrt2L^{1-s}$, and is $\ge1$ at $s=0$; a partition of unity of the mesh $w_s=L^{-i_s}/(16n)$, $i_s=\max\{i_n,s+1\}$, cuts each shell into pieces carried by products of boxes of the side $2w_s$ whose concentric cubes of the side $L^{-i_s}/(4n)$ lie at pairwise sup distances $\ge L^{-i_s}$, as $L\ge13$. A piece is the sum of its Fourier series of the period $4w_s$ in each coordinate, and each mode times fixed cutoffs is a product of entries carried by these cubes, of sup norm $\le1$; the operator $\prod_k(1-(2w_s/\pi)^2\Delta_k)^2$, four derivatives in each entry, bounds the sum of the moduli of the coefficients of a piece by $62^n$ times the supremum of its image under the operator, in which Leibniz's formula has at most $(n+1)^{4n}$ terms and the powers of $2w_s/\pi$ pay for the derivatives of the cutoffs.

The sum over the boxes of a shell, with the weight $(1+|x|^2)^{-m/2}$ of the norm $|\cdot|_m$, $m=(\lceil\mathsf v\rceil+8)n$, costs $(16n)^{3n}L^{3i_sn}$, and the bound on the cubes gives $(\mathsf uL^{\mathsf vi_s})^n$: in all $L^{(\mathsf v+3)i_sn}\le L^{(\mathsf v+3)(i_n+s+1)n}$. On the shell $s\ge1$, Taylor's formula of the order $(\lceil\mathsf v\rceil+4)n$ at a point of the diagonals at distance $\le L^{1-s}$, where $\varphi$ vanishes with all its derivatives, gives, up to a number to the power $n$, the factor $L^{-(s-1)(\mathsf v+4)n}$; so the shell $s$ carries $L^{(\mathsf v+3)i_nn}L^{(2\mathsf v+7)n}L^{-sn}$, and the series over $s$ converges. The index is $4n+(\lceil\mathsf v\rceil+4)n=m$, and $(n+1)^{4n}(16n)^{3n}\le16^{4n}n^{7n}$.

At $\mathsf v=\frac{13}2$ the index is $15n$, and with $L^{i_n}\le8nc_i$ one entry carries $n^5\cdot n^7\cdot n^{19/2}=n^{43/2}$; as $n^{43n/2}\le(e^nn!)^{43/2}$, the sum over the $6^n$ components of the labels gives
\[|\mathsf S_n(\varphi)|\le2\bigl(6\,C_*\cdot2c_3g^2\,8^5(8c_i)^{19/2}e^{43/2}\bigr)^n(n!)^{43/2}\,|\varphi|_{15n}\]
for the one functional $\mathsf S_n$ of (h), this sum being bounded for a Schwartz norm and equal to $\sigma$ at the smooth tuples: the pair $(15n,\frac{43}2)$ serves. The kernel theorem, Proposition~\YMMref{93.1}{715}, gives the order $\le4n+1$ on compact sets.
\end{proof}

\begin{proof}[Proof of Theorem~\ref{C-355-10}]
\emph{Step 1: covariance (a)} is Proposition~\ref{C-355-28}~(a): with $(\mathsf RH)(x)=\mathsf RH(\mathsf R^{-1}x)\mathsf R^{\rm T}$,
\[\sigma(\mathsf RH_1,\dots,\mathsf RH_n)=\sigma(H_1,\dots,H_n)\qquad(\mathsf R\in\mathrm O(3)).\]

\emph{Step 2: every mirror (b).} By Theorem~\ref{C-355-9}~(b), $\sigma$ is invariant under the translations and reflection positive at $\{x_3=0\}$; with (a), Lemma~\ref{C-355-29}~(b) gives, for $\mathsf gx=\mathsf Rx+cn$, $\mathsf Re_3=n$,
\[\sigma(\theta_{n,c}\bar F\cdot F)=\sigma\bigl(\theta\,\overline{\mathsf g^{-1}F}\cdot\mathsf g^{-1}F\bigr)\ge0 .\]

\emph{Step 3: data (c).} $(\mathrm{Rec1})_V$ is the composition of the translations (Theorem~\ref{C-355-9}~(b)) with (a); $(\mathrm{Rec2})_V$ is Theorem~\ref{C-355-9}~(b), the space being separable by the bounds and a countable dense family of tuples; $(\mathrm{Rec3})_V$ is Theorem~\ref{C-355-9}~(e), or follows from (a) and $H\ge|P_1|$ by the unitary rotations about $e_3$, which turn $\mathcal P$; $(\mathrm{Rec4})_V$ is Theorem~\ref{C-355-9}~(c). The classes of the $C^1_c$ tuples are dense, because a continuous entry is the uniform limit of its mollifications and $\sigma$ is continuous under such limits; so the map of the two spaces is unitary. Since $g^{-2}\le b_*$, the numbers do not depend on $\sigma$:
\[C_V=2^{p_V}C_{p_V}b_*^{-1}\le2^{p_V}C_{p_V}g^2 .\]

\emph{Step 4: kernels, representation, locality (d).} Proposition~\ref{C-355-25}~(a) at $V=\mathrm{Sym}_3(\mathbb R)$ applies to $\sigma^1$. Its clause (vii) gives the kernels, extended to continuous entries by the continuity of Step~3; for the bound at every point, a direction $u$ with $|u\cdot(x_a-x_b)|\ge|x_a-x_b|/n^2$ for all pairs exists, and the bound of $\mathbf S_n$ in the tube of the frame of $u$ at $\eta=\frac12\min\{1,\mathsf d(x)\}/n^2$ gives
\[|\mathsf s_n(x;e)|\le\bigl(C_V(2n^2/\min\{1,\mathsf d(x)\})^{p_V}\bigr)^n\prod_a|e_a| ;\]
integrating gives the bound of $\sigma$, and $\mathsf S_n(\varphi)=\int\mathsf s_n\varphi$ with the Taylor bound of $\varphi$ near the diagonals gives the growth with the norm $|\cdot|_{(\lceil p_V\rceil+4)n}$. Its clauses (i)--(iii) and (v) give the representation, with $D(\Lambda)A=\hat\Lambda A\hat\Lambda^{\rm T}$, and (iv) the locality. The fields of Proposition~\ref{C-355-25} are the $T(f)$ of Theorem~\ref{C-355-9}~(d): both are determined on $\mathcal D_\infty$ by $\langle\psi,T(f)\psi'\rangle=\int\lim_{s\downarrow0}\langle U(x)^{-1}\psi,\mathsf B(s,s;f(x))U(x)^{-1}\psi'\rangle\,dx$, which involves $\sigma$ alone, and the totality used in Proposition~\ref{C-355-25}~(a)~(viii) is $\mathcal H^{\rm ord}=\mathcal H(\sigma)$, Theorem~\ref{C-355-9}~(f).

\emph{Step 5: all couplings (e).} Let $0<g\le\gamma$. Extract a subsequence with converging labels, $b_j\to b_*$, and let $z$, $g'_1$, $g'_2$ be as in (e); $g'_2<g'_1\le\gamma^{\rm sd}_W$, and by Lemma~\ref{C-355-23}~(b) the label $L^zb_*$ lies in the interior of $I(g')$ for one $g'\in\{g'_1,g'_2\}$. By the same Lemma the dilated data are, from an index on, first-passage data at $g'$, Theorem~\ref{C-355-9}~(b) at $g'$ gives a further subsequence with a tensor limit $\sigma'$, and
\[\sigma(X)=L^{zn}\sigma'(X^{(z)})\]
exists. The dilation commutes with $\mathrm O(3)$, carries mirrors to mirrors and $H^{\rm cons}[\eta]$ to $H^{\rm cons}[\eta^{[z]}]$; so (a), (b) and the conservation pass from $\sigma'$ to $\sigma$, and by Lemma~\ref{C-355-23}~(b) $(\mathrm{Rec4})_V$ holds with $(3^{p_V}L^{4z}C_V[\sigma'],p_V,\frac12)$, where $C_V[\sigma']\le2^{p_V}C_{p_V}[g']\,g'^2$ by Theorem~\ref{C-355-9}~(c) at $g'$; this is at most the number $C_V$ of (e), which does not depend on the sequence. Then (c), (d) follow as in Steps~3, 4.
\end{proof}

\begin{proof}[Proof of Theorem~\ref{C-355-11}]
\emph{Step 1: the identity.} By Proposition~\ref{C-355-30}~(a), for $a\le n_1$,
\[X_a=b^{-1}\Bigl[\mathcal N_{\eta_a}-\tfrac14\mathfrak N[\mathsf k_{\eta_a}]-c_{\eta_a}\Bigr]-b^{-1}\sum_{e=1}^9\mathfrak J_e[\eta_a]:\]
null terms and a constant, and nine small entries. The same holds for each of the terms of the identity in the place of $X_a$.

\emph{Step 2: pieces.} Fix the side $\ell'=2\mathsf l$, put $s=\mathsf l/(8n)$, and cut every datum by a smooth partition of unity $(\chi_i)$ of mesh $s$, into pieces of sup diameter $\le\mathsf l/(4n)$; a ball of radius $R$ meets at most $N_\chi=(2R/s+3)^3$ of the supports, and
\[|\chi_iH|_\infty\le|H|_\infty,\qquad\|\chi_i\eta\|_{C^2}\le C_\chi\|\eta\|_{C^2},\qquad C_\chi=1+2\sup|\nabla\chi_i|+\sup|\nabla^2\chi_i| .\]
By multilinearity it suffices to bound one choice of a piece for each slot.

\emph{Step 3: the board.} Lemma~\ref{C-355-22} with $h=a_{K-t}=L^{-u}<\mathsf l/(2n)$ gives a board of side $\mathsf l$ with offset in $a_{K-t}\mathbb Z^3$ in which each piece lies in one cell at the margin $\delta$,
\[{\rm dist}_\infty\bigl(\operatorname{supp}(\text{piece}),\text{wall planes}\bigr)>\delta=\frac{\mathsf l}{8n};\]
several pieces may share a cell.

\emph{Step 4: one cell, then the state.} For $B\subset\{1,\dots,n\}$ group the factors $X_a-c_a$, $c_a=\langle X_a\rangle_{K,\ell',\beta}$, by cells. Lemma~\ref{C-355-17} with the pressure gap, $\mathsf Q=e^{\mathsf l^3G_g}$, and Proposition~\ref{C-355-30}~(c) in each cell give, for every $\nu\in\mathcal G^{\rm per}_{K,\beta}$,
\[\Bigl|\nu\Bigl(\prod_{a\in B}(X_a-c_a)\Bigr)\Bigr|\le\prod_{a\in B}\mathsf e_a,\qquad\mathsf e_a=\mathsf Q\,\mathsf C_n\,\varpi(X_a)\,\mathsf n(X_a);\]
at $B=\{a\}$, $|\nu(X_a)-c_a|\le\mathsf e_a$.

\emph{Step 5: the centring in $n$ factors.} With $\delta_a=\nu(X_a)-c_a$,
\[\prod_aA^\nu[X_a]=\sum_{B}\prod_{a\in B}(X_a-c_a)\prod_{a\notin B}(-\delta_a),\]
and each of the $2^n$ terms has a $\nu$-mean of modulus $\le\prod_a\mathsf e_a$: every term carries the factor of every slot, the small ones among them. A joint cumulant of order $n\ge2$ does not change when constants are subtracted from its entries; as a sum over the partitions of products of the means of Step~4, it is at most $\mathsf{op}_n\prod_a\mathsf e_a$, $\mathsf{op}_n=\sum_\pi(|\pi|-1)!$.

\emph{Step 6: the exponent.} By Proposition~\ref{C-355-30}~(b), $\prod_a\varpi(X_a)=C_\omega^nK^{np_\omega}L^{E}$ with $u\le\phi K+1$ and
\[E=n_1\bigl(3u-\tfrac K{28}\bigr)+4u(n-n_1)\le-K\bigl[n_1\bigl(\tfrac1{28}-(4n-1)\phi\bigr)\bigr]+4n=-c\,n_1K+4n,\]
since $n-n_1\le n_1(n-1)$. This gives (a), with $p=np_\omega$ and $C=N_\chi^n\max\{2^n,\mathsf{op}_n\}(\mathsf Q\mathsf C_nC_\omega C_\chi)^nL^{4n}$.

\emph{Step 7: the second form (b).} Change the scale by $z=z_C(g,d)$ steps, the least integer $z\ge1$ with $L^zg^{-2}/2\ge\mathsf x_W=(\gamma^{\rm ref}_W)^{-2}$, the number of equation~\YMMref{(2326)}{845}, and $\mathsf l_2\le L^zd/2$, $2\mathsf l_\iota$ the least side at $g'_\iota$, to the couplings
\[g'^{-2}_1=L^zg^{-2}/2,\qquad g'^{-2}_2=L^{z+1}g^{-2}/4,\]
at which the first-passage windows hold the dilated data. In the units of $g$ the eight-cell board of the least torus at $g'_\iota$ then has a side smaller than $d$, and every piece gets a cell of its own; the conservation slot is bounded by Proposition~\ref{C-355-30} at the dilated datum, and the spectators by Proposition~\ref{C-355-20} at the orders $2,\dots,8$, uniformly in the cutoff. The rate becomes $L^{-K/36}$ in the original cutoff index. On a torus $T_\ell$ the same argument is made with the torus mean.

\emph{Step 8: the limit (c).} At $n_1=1$ the bound of (a) tends to $0$ along the data of $\sigma$, while the left side tends to $\sigma(H^{\rm cons}[\eta],H_2,\dots,H_n)$; continuous entries with compact disjoint supports off $\operatorname{supp}\eta$ are at a positive distance from each other, and $dL^u\ge K^2$ holds for large $K$. The second form of \eqref{C-355-12} is the identity of $V$
\[\sum_i(\partial_i\zeta)\mathsf m_{ij}=H^{\rm cons}[\zeta e_j].\]
On a torus, (b) and the convergence of Theorem~\ref{C-355-15}~(E3)~(ii) give the vanishing of the limit.

\emph{Step 9: kernels and operators (d)} is Lemma~\ref{C-355-31}: the identity of the kernels gives, and is given by,
\[T_{0\nu}(\partial_0F)\psi=\sum_{k=1,2}T_{k\nu}(\partial_kF)\psi\qquad(\psi\in\mathcal D_\infty).\]
\end{proof}

\begin{proof}[Proof of Theorem~\ref{C-355-13}]
\emph{Step 1: the space (a).} By Lemma~\ref{C-355-32}~(a) the conditional means $\mathsf e_j(X)$ take values in compact intervals and depend continuously on $X$, uniformly in $j$; for a countable dense family $\mathfrak F$ of tuples, the images of $\nu_j$ under $\omega\mapsto(\mathsf e_j(X)(\omega))_{X\in\mathfrak F}$ are probability measures on a countable product of compact intervals, compact and metrisable, and along a subsequence they converge to a measure; by Lemma~\ref{C-355-32}~(c) its moments are the iterated plateaux of $\sigma$, which do not depend on the subsequence. $\mathsf X$ is its support, $\sigma_\xi(X)$ the coordinate of $\xi$ at $X$, extended by continuity, and
\[m(X^1,\dots,X^q)=\int_{\mathsf X}\prod_i\sigma_\xi(X^i)\,d\lambda(\xi).\]
The continuity of $\xi\mapsto\sigma_\xi(X)$ and the injectivity hold by construction, the coordinates of $\mathfrak F$ separating the points of $\mathsf X$.

\emph{Step 2: the approximants (b)~(i).} For $\xi\in\mathsf X$ and finitely many $X\in\mathfrak F$, the set of the $\omega$ at which $\mathsf e_j(X)(\omega)$ is within $\varepsilon$ of $\sigma_\xi(X)$ is an invariant event of positive $\nu_j$-measure for large $j$, since $\xi$ lies in the support of $\lambda$; a periodic Gibbs state conditioned on an invariant event is again one. Letting the finite set grow and $\varepsilon\downarrow0$ slowly along the full sequence gives $\rho_j=\nu_j(\cdot\,|\,E_j)$ with
\[\bigl\|E_{\rho_j}[A_j(X)\,|\,\mathcal I]-\sigma_\xi(X)\bigr\|_{L^\infty(\rho_j)}\to0 .\]

\emph{Step 3: one entry; recentring (b)~(ii).} $H\mapsto\sigma_\xi((H))$ is linear, continuous for the sup norm on compact sets, and invariant under $\mathbb R^3\rtimes W_3$; so it is $\int\operatorname{tr}(AH)$ with a matrix $A$ invariant under $W_3$, that is $A=\frac{c}3\mathbf 1$; and $c$ is determined by the scalar entries, so it depends on $\pi(\xi)$ alone. Replacing in the lattice entries the centring by $\nu_j$ with the centring by $\rho_j$ changes each factor by a quantity that tends to $\sigma_\xi((H_a))$ uniformly; so $\sigma^c_\xi$ is the tensor limit of $(K_j,\beta_j,\rho_j)$, and Lemma~\ref{C-355-33}~(c) gives the unitary map.

\emph{Step 4: the unique vacuum (b)~(iii).} By Lemma~\ref{C-355-32}~(d), $\lim_{|v|\to\infty}\sigma_\xi(X\cdot\tau_vY)=\sigma_\xi(X)\sigma_\xi(Y)$, so for $F\in\mathbb C\mathfrak S_V(\mathbb E_+)$
\[\|P^\xi_0[F]_\xi\|^2=\lim_{t\to\infty}\langle[F]_\xi,e^{-tH_\xi}[F]_\xi\rangle=|\sigma_\xi(F)|^2=|\langle\Omega_\xi,[F]_\xi\rangle|^2,\]
on a dense set, hence everywhere: $\ker H_\xi=\mathbb C\Omega_\xi$.

\emph{Step 5: the map to the scalar decomposition (b)~(iv).} The restrictions $\sigma_\xi|_{V_0}$ form, with $(\mathsf X,\lambda)$, a decomposition of the scalar part with the plateau formula; by the uniqueness of the canonical decomposition of the scalar limits, Proposition~\textup{\YMMref{101.1}{747}}, $\xi\mapsto\sigma_\xi|_{V_0}$ factors through exactly one map $\pi$ to $\mathsf X_0$, continuous and onto,
\[\sigma_\xi(f_1\mathbf 1,\dots,f_n\mathbf 1)=\sigma_{0,\pi(\xi)}(f_1,\dots,f_n),\]
and the plateau formulas give $\pi_*\lambda=\lambda_0$.

\emph{Step 6: Wightman fields (c).} By (b)~(ii), $\sigma^c_\xi$ is a tensor limit at $g$, so Theorem~\ref{C-355-10}~(c), (d) apply to it with the same numbers; it has (Rec5) by (b)~(iii) and the unitary map of (b)~(ii). Proposition~\ref{C-355-25}~(a)~(i) with (Rec5) gives the uniqueness of the invariant vector; the Hermiticity is
\[\langle\psi,T(f)\psi'\rangle=\langle T(f^*)\psi,\psi'\rangle,\qquad f^*=\rho(\theta)\bar f,\]
under which a component with $r$ indices equal to $3$ is Hermitian after the factor $i^r$; the representation, the spectrum, the locality and the cyclicity are Theorem~\ref{C-355-10}~(d) and Theorem~\ref{C-355-9}~(f) for $\sigma^c_\xi$. The growth is Theorem~\ref{C-355-9}~(h), whose numbers do not depend on the limit.

\emph{Step 7: the frame over a component (d).} The frame is formed at the label matrices $\mathsf m_{ij}$ of the stress tensor, the labels of $\frac12\delta_{ij}\operatorname{tr}T-T_{ij}$; (M1), (M2), (M3), (M5) are linear in the label. (M1) and the operator formula of (M2) are Theorem~\ref{C-355-9}~(d)~(vi) for $\sigma^c_\xi$, and the covariance of the kernels is Theorem~\ref{C-355-10}~(d)~(i) with $\rho(\mathsf R)\mathsf m_{ij}=\sum_{i'j'}\mathsf R_{i'i}\mathsf R_{j'j}\mathsf m_{i'j'}$, at $\mathsf s_{ij}(x;y)=\mathsf s_{1+m}(x,y;\mathsf m_{ij},\mathbf 1,\dots,\mathbf 1)$. The scalar subspace reduces $H$, $\mathcal U$, $\mathcal P$, and it is invariant under the families $\mathcal T(\mathfrak z)$ and $\mathsf b(w|\mathbf 1|w')$, which are closed under the adjoint; so its projection $\Pi_0$ commutes with $\phi(f)=T(f\mathbf 1)$ on $\mathcal D_\infty$, and with $T^E_{ij}(F_1)=\Pi_0T(F_1\mathsf m_{ij})\Pi_0$,
\[T^E_{ij}(F_1)\phi(f)\psi=\Pi_0T(F_1\mathsf m_{ij})\phi(f)\psi=\Pi_0\phi(f)T(F_1\mathsf m_{ij})\psi=\phi(f)T^E_{ij}(F_1)\psi\]
at spacelike separation by the locality: (M3). The scalar subspace reduces $U(a,\Lambda)$ as well: the restriction of $(\sigma^c_\xi)^1$ to the tuples with entries in $V_0$ is reconstruction data with the label space $V_0$, whose space is the scalar subspace, and by Proposition~\ref{C-355-25}~(a)~(v) the span of $\Omega$ and of the vectors $\phi(f_1)\cdots\phi(f_n)\Omega$ is dense in it; as $\hat\Lambda\mathbf 1\hat\Lambda^{\rm T}=\hat\Lambda\hat\Lambda^{\rm T}=\mathbf 1$, Theorem~\ref{C-355-10}~(d)~(ii) gives $U(a,\Lambda)\phi(f)U(a,\Lambda)^{-1}=\phi(f(\Lambda^{-1}(\cdot-a)))$ on $\mathcal D_0$, and with $U(a,\Lambda)\Omega=\Omega$ every $U(a,\Lambda)$ maps this span onto itself; $U$ being a unitary representation, the scalar subspace and its orthogonal complement are invariant, and $\Pi_0$ commutes with every $U(a,\Lambda)$.

So $U$ restricts to a representation on the scalar subspace, and (M5) is $\Pi_0$ applied to the Lorentz covariance on $\Omega$ of Theorem~\ref{C-355-10}~(d)~(ii), with the passage from the Euclidean labels to $T_{\mu\nu}$ of Definition~\ref{C-355-5}~(3). (M4)~(a): conservation passes to $\sigma^c_\xi$ by Lemma~\ref{C-355-33}~(b), so Theorem~\ref{C-355-11}~(c), in the second form of \eqref{C-355-12} with $\zeta$ vanishing near the scalar points, gives $\sum_i\int\partial_i\zeta(x)\,\mathsf s_{ij}(x;y)\,dx=\sigma^c_\xi(H^{\rm cons}[\zeta e_j],\text{the scalar entries})=0$; the scalar subspace has a unique vacuum, by (b)~(iii); and Lemma~\ref{C-355-31} gives (M4)~(a). At the labels $e_ie_j^{\rm T}+e_je_i^{\rm T}=\delta_{ij}\mathbf 1-\mathsf m_{ij}$ the identity of the kernels fails: among one scalar entry their divergence is $\partial_{x_j}\mathsf s_2(x,y;\mathbf 1,\mathbf 1)$, which does not vanish identically by the free law (E6) of Theorem~\ref{C-355-15}; the components $T_{ij}$ themselves are not conserved.

\emph{Step 8: the frame on $\mathcal H_0$ (e).} For $\sigma$ itself the same compression applies without the unique vacuum: (M1) is Theorem~\ref{C-355-9}~(d)~(vi), the covariant kernels come from Theorem~\ref{C-355-10}~(d)~(i), and (M3), (M5) are the compressions of Theorem~\ref{C-355-10}~(d)~(iii), (ii), as identities between scalar vectors, for instance
\[\bigl\langle\Omega,\phi(f_1)\cdots\phi(f_n)T[\Lambda t](F\circ\Lambda^{-1})\Omega\bigr\rangle=\bigl\langle\Omega,\phi(f_1\circ\Lambda)\cdots\phi(f_n\circ\Lambda)T[t](F)\Omega\bigr\rangle .\]
\end{proof}

\begin{proof}[Proof of Proposition~\ref{C-355-14}]
The argument follows the proofs of Propositions~\YMMref{108.3}{771}, \YMMref{110.2}{777} and \YMMref{110.1}{776} for scalar data, with one entry carrying a label; its steps are these.

\emph{Step 1: the scalar part and one label.} Let $e'\in V\setminus V_0$. For a tuple whose entries all lie in $V_0$, $f_1\mathbf 1=f_1(\mathbf 1-e')+f_1e'$ writes it as the sum of two tuples with exactly one entry outside $V_0$, at which $\sigma(\mathsf RX)=\sigma(X)$ is supposed; so the scalar part of $\sigma$ is invariant under $\mathrm O(3)$, and with the translations it is invariant under all the isometries. With Theorem~\ref{C-355-9}~(b), (c), (e), which give the positivity, the point bound and the circular cone, the space being separable by the bounds and a countable dense family of tuples, its restriction to $C^1_c$ entries is scalar reconstruction data (Rec); Theorem~\textup{\YMMref{AP}{780}} gives this as well, as in Step~6 of the proof of Theorem~\ref{C-355-9}. The insertion and point operators with one tensor label and their energy bounds are those of Theorem~\ref{C-355-9}~(d), and Proposition~\ref{C-355-27}~(c) gives the kernel $\mathsf s_{1+m}(x,y;e,\mathbf 1,\dots,\mathbf 1)$ of one tensor entry among $m$ scalar entries, continuous and real-analytic at all non-coinciding points; by continuity, the supposed covariance at one label is $\mathsf s_{1+m}(\mathsf Rx+v,\mathsf Ry+v;\rho(\mathsf R)e,\mathbf 1,\dots,\mathbf 1)=\mathsf s_{1+m}(x,y;e,\mathbf 1,\dots,\mathbf 1)$ for $\mathsf R\in\mathrm O(3)$, $v\in\mathbb R^3$.

\emph{Step 2: one holomorphic function at one label; (M2).} For a frame $\mathsf R_u\in\mathrm{SO}(3)$ with $\mathsf R_ue_3=u$ and an order $\varsigma$ of the $1+m$ points along $u$, the piece on the tube of the configurations ordered along $u$ in the order $\varsigma$ is the function $\mathbf W_{1+m}$ of Theorem~\ref{C-355-9}~(d)~(iii) written in the frame $\mathsf R_u$ at the label $\rho(\mathsf R_u)^{-1}e$; in the components of the label it is
\[\sum_{i',j'}(\mathsf R_u)_{ii'}(\mathsf R_u)_{jj'}\,\mathbf W_{1+m}\bigl(\mathsf R_u^{-1}\,\cdot\,;E_{i'j'},\mathbf 1,\dots,\mathbf 1\bigr),\]
$E_{i'j'}$ the symmetrised matrix units, that is $\rho(\mathsf R_u)\mathbf W_{1+m}(\mathsf R_u^{-1}\,\cdot\,)$ at the one label. At the real points of the intersection of two tubes, the configurations ordered in both frames, both pieces are equal to the kernel of Step~1, by the covariance at one label; the intersection is convex and these real points form an open subset of its real section, so the two pieces agree on it, and they form one holomorphic function on the union of the tubes. As $|(\mathsf R_u)_{ii'}|\le1$, its bound is that of $\mathbf W_{1+m}$ times at most $9$, the number of the terms $(i',j')$. Its restriction to the real non-coinciding configurations is the kernel; so the kernels $\mathsf s_{ij}(x;y)=\mathsf s_{1+m}(x,y;\mathsf m_{ij},\mathbf 1,\dots,\mathbf 1)$ are covariant at all non-coinciding points, and with the product formula of Theorem~\ref{C-355-9}~(d)~(vi) this is (M2).

\emph{Step 3: locality; (M3).} The matrix elements of products of the fields between vectors of $\mathcal D_\infty$ are the boundary values of the one function of Step~2 from the tubes of the orders of the real times. For a pair of spacelike separated points, one of them carrying the label, the tubes of the two orders are joined by the tubes of the frames tilted by an angle of the order $\eta$ towards the direction of the separation, on which the function is one and the same; apply the uniform smeared-bound argument in the proof of Proposition~\ref{C-355-25}~(a) to the holomorphic kernel and polynomial bound of Step~2, with its single tensor label fixed and the remaining labels scalar. The Cauchy--Riemann identities and integration by parts give the same bound $2\eta Q_e(D'g)$ for the difference of the two orders, with $Q_e(D'g)$ independent of $\eta$; hence the boundary values agree as $\eta\to0$. So $\langle\psi',[T(F_1\mathsf m_{ij}),\phi(f)]\psi\rangle=0$ at spacelike separation for $\psi',\psi$ in the span of the vectors $e^{-\tau H}[F]$, $\tau>0$, $F$ scalar, which lie in $\mathcal D_\infty$, are dense in $\mathcal H_0$ and form a core of every power of $1+H$ there (Lemma~\ref{C-355-31}~(a), (b), at the span $\mathcal V_1$); the bound of (M1) and the commutation of $\Pi_0$ with $\phi(f)$ on $\mathcal D_\infty$, as in Step~7 of the proof of Theorem~\ref{C-355-13}, carry this to $T^E_{ij}(F_1)\phi(f)\psi=\phi(f)T^E_{ij}(F_1)\psi$ on $\mathcal D_0$: (M3).

\emph{Step 4: Lorentz covariance; (M5).} The invariance of $\mathbf W_{1+m}$ at one label under the real rotations $\mathsf R_\zeta$ in the plane of $e_1$, $e_3$, the label being turned by $\rho(\mathsf R_\zeta)$, continues holomorphically in the angle $\zeta$ to the imaginary angles $\zeta=is$, which act on the arguments, through $\iota$, as the boosts of rapidity $s$ in the plane of $(t,x_1)$; the label is then moved by $\rho_{\mathbb C}(\mathsf R_{is})$, $\rho_{\mathbb C}$ the holomorphic continuation of $\rho|_{\mathrm{SO}(3)}$, that is by $e\mapsto\hat\Lambda e\hat\Lambda^{\rm T}$, and no bound on this factor is needed, the identity being one of holomorphic functions at fixed arguments. With the real rotations in the plane of $e_1$, $e_2$, which fix $e_3$, this gives the invariance of $\langle\Omega,\phi(f_1)\cdots\phi(f_m)T(F)\Omega\rangle$, the tensor entry at the last place, under $(f_k,F)\mapsto(f_k\circ\Lambda^{-1},\hat\Lambda F(\Lambda^{-1}\cdot)\hat\Lambda^{\rm T})$ for all $\Lambda\in\mathcal L$.

The scalar data of Step~1 carry the representation $U$ of $\mathbb R^{1+2}\rtimes\mathcal L$ on $\mathcal H_0$ with 
\[
U(\Lambda)\phi(f_1)\cdots\phi(f_m)\Omega=\phi(f_1\circ\Lambda^{-1})\cdots\phi(f_m\circ\Lambda^{-1})\Omega
\]
 (Proposition~\ref{C-355-25}~(a) at $V_0$); these vectors are dense in $\mathcal H_0$, and the invariance reads $U(\Lambda)\Pi_0T(F)\Omega=\Pi_0T(\hat\Lambda F(\Lambda^{-1}\cdot)\hat\Lambda^{\rm T})\Omega$. At the label matrices $\mathsf m_{ij}$, with the passage from the Euclidean labels to $T_{\mu\nu}$ of Definition~\ref{C-355-5}~(3), this is (M5).
\end{proof}

\begin{proof}[Proof of Theorem~\ref{C-355-15}]
\emph{Step 1: (E1).} By Lemma~\ref{C-355-21}, which holds for complex weights, $\mathcal T_K(H_a)$ is the first member of a source under $(60|H_a|_\infty,S_{H_a})$ for $H_a$ with values in $V_{\mathbb C}$, and Proposition~\ref{C-355-20}~(a) holds for source systems real or not; so the number $60^n$ holds for complex $H_a$. Since $s(d)\le K$, $a_K\le d/C_W(s(d))\le d/2197$, and a plaquette of $S_{H_a}$ lies within $a_K/\sqrt2$ of the points of $\Sigma_a$; so (Geo) at $d$ gives the geometry of Proposition~\ref{C-355-20}:
\[0.87\,d+0.71\,d/2197\le\tfrac78d,\qquad1.26\,d-1.42\,d/2197\ge\tfrac54d .\]
Proposition~\ref{C-355-20}~(a), (b) give the bound and its logarithmic form, and (a$'$) the factor off the small sides.

\emph{Step 2: (E2).} A box of side $\le0.99\,d''$ lies in a ball of radius $0.8574\,d''$, and boxes that meet different $\Sigma_a$ are at distance $\ge d\ge\frac43d''\ge1.26\,d''$; by multilinearity the cumulant is a sum over the choices of boxes, each bounded by (E1) at $d''$, which gives (i), the number of boxes being at most $(2\ell/d'')^3$. For $K<s''$ a box holds at most $3(d''/a_K+1)^3$ points of plaquettes and bonds, $|a_p|\le2$, $|q_{\mathfrak b}|\le4$, $L^Kd''<C_W(s'')$ and $b\le\bar x_0$:
\[|\mathcal T_K(h)|\le18\,\bar x_0\,(C_W(s'')+1)^4/d''\cdot b^{-1}\sup|h|\quad\text{for $h$ carried by one box,}\]
and the cumulant is a sum over the partitions of products of means, which gives (ii) with $f_n$. For (iii), at low cutoffs an exponential moment $\mu(e^{\beta a_p})\le C$ of one plaquette gives a bound $C^{\rm cr}_n(6\ell^3C_W(s'')^3/d''^3)^n$, and at high cutoffs (i) applies.

\emph{Step 3: (E3).} (i) The dilation of the lattice by $L^z$ (Proposition~\YMMref{119.1}{815}) gives the identity of the cumulants
\[\langle\mathcal T_K(H_1);\dots;\mathcal T_K(H_n)\rangle_{K,\ell,\beta}=L^{zn}\langle\mathcal T_{K_\star}(H_1(L^{-z}\cdot));\dots\rangle_{K_\star,L^z\ell,\beta},\]
at the label $a_{K_\star}\beta=L^zb$, and (E1) at $g_\star$ and the distance $L^zd$; the bounds of the label at $g_\star$ give $\delta_\star b\le L^{-z}<g_\star^2b$, and $\log(1/(L^zd))\le\log(b/d)+C(g_\star)$. The statement is invariant under $(K,\ell,b,d)\mapsto(K-1,L\ell,Lb,Ld)$. (ii) By Corollary~\YMMref{74.2}{616} the pairs $(K,L^Kb)$, $b\in I(g)$, are first-passage data at $g$ with label $b$, and every $\ell\in\mathfrak L(g)$ has (Sd.1) by Proposition~\YMMref{74.1}{615}; the convergence along the full sequence, $n$-linear and symmetric, with the rates of its clauses (a) and (b), is Proposition~\ref{C-355-36}, applied to the six cell insertions of each entry, real and imaginary parts apart, the cumulant being a sum of at most $12^n$ cumulants of cell insertions ($6^n$ at real entries), and the bound is (E1) with (a$'$), the radii $0.86\,d$ and $1.27\,d$ leaving room for the lattice corrections.

\emph{Step 4: (E4).} (i) At two points the leading term of the cumulant at the working level is the Gaussian value of the effective density, which is free of $d$, and the other terms are smaller; it is computed on the least torus with the moments of eight and sixteen mirror copies of the pair, and carried to every side and every state by Lemma~\ref{C-355-17}, which gives, for real weights,
\[|\langle\mathcal T_K(H_1);\mathcal T_K(H_2)\rangle|\le180^2C^\diamond(g)\,b^{-2}|H_1|_\infty|H_2|_\infty ,\]
with the number $C^\diamond(g)=\max\{2\,\mathsf c_\diamond(g)^2,C^T_2(\frac12;1)\}$ of (E4)~(i), and with the factor $4$ for complex weights, their real and imaginary parts being taken apart; on one torus the same holds from the first depth $s_1(g,\ell)$ on. (ii) Cut each of the three cubes into $4^3$ cubes of side $\delta/4$; the cumulant is a sum of $64^3$ cumulants of pieces, and three pieces, one from each cube, are at pairwise sup distances $\ge2\delta=8\cdot\frac\delta4$, so that a coordinate mirror parts one of them from the other two. The third cumulant is the third centred moment, and the Schwarz inequality at this mirror gives $|\kappa|^2\le\langle\text{pair}\cdot\text{its mirror image}\rangle\,\langle\text{single}\cdot\text{its mirror image}\rangle$.

The pair with its mirror image is a centred moment of four entries, the sum over the partitions into blocks of at least two members of products of cumulants, at most $60^4(k'_4+3k_2'^2)b^{-4}$ times the sup norms by Proposition~\ref{C-355-20}, with the power $32\mathfrak p_0$ of the logarithm; the single entry with its mirror image is a two-point function on the one torus, bounded by the second statement of (i), with a number free of the distance and of the logarithm. After the square root the power of the logarithm is $16\mathfrak p_0$, and the bound carries $b^{-3}$. (iii) The history without large-field regions is entered at the level $T$ of (E4)~(iii), in the place of the working level of (E1), and the smallness of the side is replaced by the ceiling $y_K\ge\mathsf Y^\natural$ of Definition~\ref{def:rarity-choices}, under which Proposition~\ref{prop:local-large-field-rarity} bounds the large-field regions uniformly in the volume, with its constant $c^\natural_2$.

By the flow bound \YMMref{(1530)}{615} at the indices $0\le T$, $|x_T-L^{K-T}b|\le\bar B$, and $\bar B\le\frac{13}{24}g^{-2}<\frac{13}{24}b$ (Lemma~\ref{lem:setup-stopped-arithmetic}); with $L^{K-T}=\hat d/d$, $\hat d\ge3$ and $d<1$ this gives $x_T\ge3b/d-\bar B>2b/d$ and $\bar B\le\frac14x_T$, so that $d\le b\,e^{-y^\sharp}$ gives $y_T\ge y^\sharp$, and the conditions that define $y^\sharp$ hold at $y=y_T$. At the levels $j\le T$, $x_j\ge x_T\ge x_{\rm J}$ and $\theta(x_j)\le\theta(e^{y_T})$, a number formed from the constants and $p$; this closes the induction of the mixed derivatives, as (N.$\theta$) does in the proof of Proposition~\ref{C-355-20}. By the choice of $T$ the least sup distance $\mathsf D_T$ of two supports in the units of the level $T$ is at least $M(2+\frac p2y_T)$, so that a polymer that joins two supports has there a linear size of at least $\mathsf D_T/M-2\ge\frac p2y_T$.

The correction to the first-order size of one entry at the level $T$ carries $e^{-d_*(T)}$, $d_*(T)$ being the counterpart at the level $T$ of the quantity $\min\{c_{\rm d}\,\varepsilon_k^{-1},(\ell/(2a_{K-k})-4)/(2M)\}-1$ of the proof of Proposition~\ref{C-355-20}, and $d_*(T)\ge\frac12y_T$ makes it at most $x_T^{-1/2}$. The configurations with a large field enter through the level $t\le T$ of its first creation, with the weight of Proposition~\ref{prop:local-large-field-rarity} per unit cube of that level; with the count of the unit cubes of the torus at the level $T$, $(\ell L^{K-T})^3\le(\frac54)^3(\ell/b)^3e^{3y_T}$, as $L^{K-T}b\le x_T+\bar B\le\frac54x_T$, and with the sizes of the $p$ insertions, they add at most a number formed from the constants and $p$ times
\[e_T=4\cdot24^3\bigl(\tfrac54\bigr)^3(\ell/b)^3\exp\bigl((3+p^2)y_T-c^\natural_2A_0^2y_T^{2\mathfrak p_0}\bigr)\le1\]
times the term of the history without large-field regions, the terms of the levels $t<T$, at which $y_t>y_T$, decreasing geometrically by $2\mathfrak p_0c^\natural_2A_0^2y^{2\mathfrak p_0-1}\ge4+(3+p)p$. At the level $T$ the size of a plaquette source carries $y_T^{2\mathfrak p_0}$, the supremum of $x_Ta_p$ on the small fields, and the count of the unit cubes of that level in a ball of radius $d$ the factor $\hat d^{\,3}$: this is the first form of the bound. By the maximality of $T$, $\hat d<L(M(2+\frac p2y_{T+1})+1)\le L(M(2+\frac p2y_T)+1)$, which gives the second form; and $x_T\le L^{K-T}b+\bar B<\frac43L(M(2+\frac p2y_T)+1)\,b/d$ gives the bound of $y_T$.

\emph{Step 5: (E5).} (i) The cumulant of order $n\ge2$ is the sum over the partitions $\pi$ of $\{1,\dots,n\}$ without singletons of $(-1)^{|\pi|-1}(|\pi|-1)!$ times the products over the blocks $B$ of the moments of the entries of $B$ centred in the state. For $K\ge\max_mK^G(m)$ each of these moments is bounded by Theorem~\ref{C-355-9}~(a) at its own order $\#B$, which gives the second bound; on a torus the chessboard estimate takes the place of the domination, the pieces of different entries lying in cells different modulo the side; for the finitely many smaller cutoffs the crude bound of Step~2 holds. Together these give a number $C^T_n(d;R)$. (ii) Extraction as in Theorem~\ref{C-355-9}~(b); on the scalar entries the limit is the scalar limit; on one torus the full sequence converges by (E3)~(ii), and the bounds of (i) do not depend on the side. (iii) The cumulant does not change when each entry is centred by its mean $m_a$ on the least torus $T_{\ell'}$; by the moment--cumulant formula it is then the sum over the partitions $\pi$ of $\{1,\dots,n\}$ of $(-1)^{|\pi|-1}(|\pi|-1)!$ times the products over the blocks $B$ of the moments of the entries of $B$, centred by these numbers.

Since $s(d)\le K$, $a_K\le d/2197<\mathsf l/(2n)$, the distance $d$ being $\le\mathsf l/(6.96\,n)$, and $\mathsf l=\ell'/2\in a_K\mathbb Z$; so Lemma~\ref{C-355-22} at $h=a_K$ gives a board of side $\mathsf l$ in which the $n$ balls of radius $0.87\,d$ that hold the sets $\Sigma_a$, of sup diameter $\le1.74\,d\le\mathsf l/(4n)$, lie in cells at the margin $\mathsf l/(8n)$; on a torus $T_\ell$, $\ell\in\mathfrak L(g)$, the lemma is applied to lifts, $\ell$ being a multiple of $2\mathsf l$, and the cells are taken modulo the side. Let $\mathcal C$ be the partition of the entries whose sets share a cell, modulo the side on a torus. The moment of a block $B$ is the mean of the product of the functions $\Phi_C=\prod_{a\in B\cap C}(\mathcal T_K(H_a)-m_a)$ of different cells, one for each class $C\in\mathcal C$ that meets $B$, each a function of the bonds inside its cell, as the margin $\mathsf l/(8n)$ exceeds $a_K$; and the domination with the pressure gap, Lemma~\ref{C-355-17} at $\ell_1=\ell'=2\mathsf l$,
\[\Bigl|\nu\Bigl(\prod_C\Phi_C\Bigr)\Bigr|\le\prod_Ce^{\mathsf l^3G_g}\|\Phi_C\|_{{\rm ch},\ell'},\]
bounds it in every periodic Gibbs state, and on every torus of $\mathfrak L(g)$ with the chessboard estimate, the cells being different modulo the side. The board of side $\mathsf l$ has eight cells on $T_{\ell'}$; so the eighth power of $\|\Phi_C\|_{{\rm ch},\ell'}$ is the mean on $T_{\ell'}$ of the product of the $8m$ mirror copies of the $m=\#(B\cap C)$ centred entries of the cell, at complex weights with the conjugations of the reflections, $\overline{\mathcal T_K(H)}=\mathcal T_K(\bar H)$ being a source under the same majorant. The copies are centred, the mean on $T_{\ell'}$ being invariant under the lattice isometries (Lemma~\ref{C-355-21}); so this mean is the sum over the partitions of the copies into blocks of at least two members of products of cumulants.

The copies in one cell are the images of the entries under one isometry, so their sets lie in balls of radius $0.87\,d$ and are $\ge1.26\,d$ apart; the copies in different cells lie at sup distance $\ge\mathsf l/(8n)$ from the walls, hence at distance $\ge\mathsf l/(4n)\ge1.74\,d$ from each other: the $8m$ copies have (Geo) at $d$. The least side at $g\le\gamma^{\rm sd}_W$ is small over the running couplings of every first-passage datum, so Proposition~\ref{C-355-20}~(a) applies on $T_{\ell'}$ at every order $r=\#I\le8n$: (N.$\theta$) at $(8n,g,s(d))$ gives (N.$\theta$) at $(r,g,s(d))$ for $2\le r\le8n$, the numbers $T_r(3)$ not decreasing in $r$, as each term of the series that forms $T_r(3)$ carries the $r$-th power of a number $\ge1$.

With the majorants $60|H_a|_\infty$ of Step~1, at complex weights too, it needs no upper bound on the distances and bounds each cumulant by $60^{\#I}k'_{\#I}(g,s(d))\,b^{-\#I}$ times the sup norms. So
\[\|\Phi_C\|_{{\rm ch},\ell'}\le60^m\,\mathfrak m_{8m}(g,s(d))^{1/8}\,b^{-m}\prod_{a\in B\cap C}|H_a|_\infty ,\]
and the sum over the partitions $\pi$, with the maximum over $\mathcal C$, gives the number $k^{\rm fold}_n$. By Proposition~\ref{C-355-20}~(b), $\mathfrak m_{8m}(g,s)$ is at most a number times $Y_s^{64m\mathfrak p_0}$; so each eighth root carries the power $8\mathfrak p_0$ of $Y_{s(d)}\le2\log\frac1d+C(g)$ per entry of the cell, and the powers add up to $8n\mathfrak p_0$, as the classes $B\cap C$ part $\{1,\dots,n\}$ for each $\pi$. The cumulants of a tensor limit are the limits of the lattice cumulants, which gives the last bound.

\emph{Step 6: (E6).} At the working level of the distance $d$ the effective density is, up to the non-Gaussian part, that of three free photon fields $B_\lambda$; the Wick value of the cumulant of $\langle B_\lambda,B_{\lambda'}\rangle(x)$ and $\langle B_\mu,B_{\mu'}\rangle(y)$ is
\[b^{-2}\cdot\tfrac32\operatorname{tr}\bigl[E_1K_3(x-y)E_2K_3(x-y)\bigr]\quad\text{for the labels }E_1,E_2,\]
since the propagator of the field strength of a free photon in three dimensions is $K_3$ up to a contact term. The non-Gaussian part is smaller by a negative power of $\log\frac1d$, which gives $e_g(d)$. The transfer from the least torus to every side and every state uses the chessboard estimate with eight and sixteen mirror copies. Since $\mathsf W\ne0$, the limit determines $b_*$. For the bump of (E6), $0\le\varphi\le1$, $\varphi=1$ on $B(x_0,d/128)$ and $\varphi=0$ off $B(x_0,d/64)$, $x_0=\frac32d\,e_3$, the explicit evaluation of $\mathsf W(\theta(\varphi E),\varphi E)$ gives at least $w_0\operatorname{tr}E^2$; with the remainder, at most $e_g(d)b_*^{-2}\operatorname{tr}E^2$, this gives the positivity at the $d$ with $e_g(d)<w_0$.

\emph{Step 7: (E7).} Let $\ell\in\mathfrak L(g)$, $b\in I(g)$, and let $u$ be the limit cumulant of (E3)~(ii), continued complex linearly and restricted to smooth entries: a functional of $n$ entries on $\mathbb M=T_\ell$ with $V=\mathrm{Sym}_3(\mathbb R)$, in the sense of Lemma~\ref{C-355-34}. \emph{The bound $(\mathrm W)$.} Closed cubes of the side $\frac32h$ lie in the balls of radius $\frac{3\sqrt3}4h<1.3\,h\le0.86\,d$ about their centres, $d=1.52\,h$, and cubes at pairwise sup distances $\ge\frac52h$ are $\ge1.27\,d$ apart; if $1.52\,h\le L^{-s_{\rm th}(n)}$, then $d<1$ and $L^{s(d)}\ge C_W(s(d))/d\ge L^{s_{\rm th}(n)}$, so that (N.$\theta$) holds at $(n,g,s(d))$, $s_{\rm th}(n)$ being the number of (E2). So the bound of the limit in (E3)~(ii), with Proposition~\ref{C-355-20}~(b) and
\[2\log\tfrac1{1.52\,h}+C(g)\le\bigl[2+2\log\tfrac1{h_0}+C(g)\bigr]\,l(h)\qquad(0<h\le h_0),\]
gives $(\mathrm W;0,8n\mathfrak p_0,0)$ for $u$, with $h_0\in\ell2^{-\mathbb N}$, $h_0\le\min\{\ell/(64n),\frac12L^{-s_{\rm th}(n)}\}$, and
\[A=60^n\max\{1,\mathsf c(\ell)\}^nC_n(g)\bigl[2+2\log\tfrac1{h_0}+C(g)\bigr]^{8n\mathfrak p_0}b^{-n},\]
at complex entries as well, (E1) holding at complex weights (Step~1) and the limit being continued complex linearly. \emph{The extension.} Lemma~\ref{C-355-34}~(a)--(c) at $\kappa=0$, $q=8n\mathfrak p_0$, $r=0$ give
\[\tilde S_n(F)=\sum_{Q\in\mathcal Q}u_Q\bigl(\varphi_Q(F-P_QF)\bigr),\qquad P_QF\ \text{of total order}\le3k\ \text{at a stratum of $k$ clusters},\]
a distribution of order at most $3n-2$, as $\lfloor\frac{3n}2\rfloor+1\le3n-2$; and $\epsilon=1$, every $3k$ being an integer, so that $|\tilde S_n(F)|\le C\,b^{-n}l(\lambda)^{8n\mathfrak p_0+1}\|F\|_{(\lambda),3n-2}$ for $0<\lambda\le h_0$ and every $\lambda$-supported $F$, with $C=C(n,g,\ell)$, which depends on the side through $\mathsf c(\ell)$ and $h_0$. \emph{The lattice cumulants.} Let $u_K(F)$ be the lattice cumulant at $(K,\ell,L^Kb)$ with the weights read at the barycentres of the plaquettes and the midpoints of the bonds, continued to functions $F$ on $T_\ell^n$ as a finite sum of values of $F$. By (E1) with the factor $\max\{1,\mathsf c(\ell)\}^n$ of the sides with (Sd.1), $u_K$ has the inequality of $(\mathrm W;0,8n\mathfrak p_0,0)$, with the same $A$ and $h_0$, at the sides $h\le h_0$ with $s(1.52\,h)\le K$, an interval $[h'_K,h_0]$ since $s(\cdot)$ does not increase; and $(u_K)_Q(G)=u_K(G)$, the Fourier series converging at every point. With $\mathcal Q_K=\{Q\in\mathcal Q:s(1.52\,h(Q))\le K\}$ and $h_K$ the least side in $\mathcal Q_K$, the subtracted lattice cumulant is
\[\tilde u_K(F)=\sum_{Q\in\mathcal Q_K}(u_K)_Q\bigl(\varphi_Q(F-P_QF)\bigr);\]
it differs from $u_K(F)$ by $\sum_{Q\in\mathcal Q_K}u_K(\varphi_QP_QF)$, which reads of $F$ the derivatives of the orders $\le3k$ on the strata of $k$ clusters alone, and by a term that reads $F$ where $\mathsf d_\infty<6h_K$ alone. The clause of Lemma~\ref{C-355-34}~(c) on a range of sides, at $h_\flat=h_K$, gives $|\tilde u_K(F)|\le C\,b^{-n}l(\lambda)^{8n\mathfrak p_0+1}\|F\|_{(\lambda),3n-2}$ for all $K$, all $8h_K\le\lambda\le h_0$ and all $\lambda$-supported $F$: one bound for all the cutoffs. For one $Q$, $(u_K)_Q(G)$ is the series over the Fourier modes of $G$ of the values of $u_K$ at products of smooth entries carried by the cubes $Q^*_a$; each of these values converges to that of $u$, by (E3)~(ii) at Lipschitz entries read at the points, and the series is dominated, uniformly in $K$, by $C\,A\,l(h(Q))^{8n\mathfrak p_0}\sum_{\mathbf k}|\hat G_{\mathbf k}|<\infty$, $\hat G_{\mathbf k}$ the Fourier coefficients of the proof of Lemma~\ref{C-355-34}~(a); so $(u_K)_Q(G)\to u_Q(G)$.

The sum over $Q$ is dominated, uniformly in $K$, by the series that gives the absolute convergence in Lemma~\ref{C-355-34}~(b); hence $\tilde u_K(F)\to\tilde S_n(F)$ for every smooth $F$. \emph{The contact terms.} The list is Lemma~\ref{C-355-34}~(e) at $\kappa=0$: $\omega_k=3k$, the coefficients $T_{\pi,\alpha}\in\mathfrak E_{|\pi|}(|\alpha|)$, $|\alpha|\le3|\pi|$, and two extensions of the cumulant functional in the class $\mathfrak E_n(0)$ differ exactly by such a sum. The translation-invariant, real, symmetric two-point contact terms are classified in Lemma~\ref{C-355-34}~(f): the invariant space has dimension six under $\mathrm O(3)$ on $\mathbb R^3$, and twelve under the cubic group on $\mathbb R^3$ or $T_\ell$.
\end{proof}

\subsection{Scope of the conclusions}

The theorems concern the components $\operatorname{tr}F_{ij}F_{kl}$ at separated points; at coinciding points they give the extensions of Theorem~\ref{C-355-15}~(E7), on one torus. The Wightman property and the unique vacuum are asserted for each component of the canonical decomposition. Free of the side of the torus and of the state are the bounds of Theorem~\ref{C-355-7}, of Theorem~\ref{C-355-9}~(a) and (h), of Theorem~\ref{C-355-11}~(a) and of the first display of its clause (b), and, in Theorem~\ref{C-355-15}, those of the first statement of (E4)~(i), of (E5), and of (E6) with its numbers $d_0(g)$ and $e_g(d)$. The numbers of (E1), of the first member of (E2)~(i) and of (E3)~(i) contain no side, the side being one that is small over the running couplings, and (E2)~(ii) holds on any side with a number free of it. The side enters the numbers of (E1) off the small sides, through the factor $\max\{1,\mathsf c(\ell)\}^n$; of the second member of (E2)~(i), through $(2\ell/d'')^{3n}$; of (E2)~(iii), through $C(n,g,\ell)$; of (E3)~(ii), through the numbers of Proposition~\ref{C-355-36} and the factor of the bound of the limit; of the second statement of (E4)~(i), through the first depth $s_1(g,\ell)$; of (E4)~(ii), through $C(g,\ell)$ and $\delta_\star(g,\ell)$; of (E7), whose number is $C=C(n,g,\ell)$; and of the torus form of Theorem~\ref{C-355-11}~(b).

\section{Conservation of the stress tensor}
\label{stc:sec}

The curvature bilinears form a symmetric tensor $T$. The associated stress tensor is
\[
 S_{ij}=\tfrac12\delta_{ij}\operatorname{tr}T-T_{ij}.
\]
We prove its conservation inside correlations by estimating the remainder in an exact lattice identity. The argument separates two tasks: making the orbital remainder small under the renormalisation steps, and identifying the remaining insertion with the divergence of $S$.

The main estimates hold at finite cutoff in every periodic Gibbs state and on a fixed torus. Theorems~\ref{stc:cutoff} and~\ref{stc:scale} give the rates and their dependence on the separation of the supports. Theorem~\ref{stc:limit} passes the identity to every convergent tensor sequence. We retain the alternative null-form argument because it isolates the collapse estimate needed for the reflection-positive null vector.

\subsection{Notation and hypotheses}

\emph{Constants, couplings, sides.} Take the admitted constant choice of Proposition~\YMMref{69.1}{595} and Section~\YMMref{82}{665}. Its constants include the blocking factor $L$ (odd, $L\ge13$), the cube size $M$, an exponent $\mathfrak p_0\ge1$, the bound $\gamma\in(0,\frac12]$ of the reference couplings, $B_\sharp\ge0$, $\bar B=B_\sharp L/(L-1)$, and the further letters fixed there. A number is a \emph{function} of some data if it depends on these data and on the constants alone; $C$ and $p$ denote such numbers, which may change from line to line. The cutoff is $K\in\mathbb N$, $a_K=L^{-K}$, $\mathbb L_K=a_K\mathbb Z^3$; $\beta>0$ is the inverse coupling of the Wilson weight $\rho_0=e^{-\beta\sum_pa_p}$, where $a_p=1-\frac12{\rm tr}\,\partial U(p)$ is the action density of the plaquette $p$ (equations~\YMMref{(6)--(7)}{13}); $b=a_K\beta$ is the \emph{label}.

For $0<g\le\gamma$, $\mathfrak d(g)$ is the set of first-passage data $(K,\beta)$ at $g$ (\YMMref{74.2}{616}, with the canonical data of \YMMref{74}{614}), $\mathfrak L(g)$ the set of sides at $g$ (\YMMref{(1528)}{615}), all of the form $2L^m$, $X_g=L(g^{-2}+B_\sharp)$, $\mathsf b_+=X_g+\bar B=L(g^{-2}+\bar B)$ (\YMMref{74}{614}, with $\bar B=B_\sharp L/(L-1)$, \YMMref{(1526)}{615}), and $I(g)=(g^{-2}+\bar B,X_g-\bar B]$. Lemmas~\ref{stc:cmean}\,(c), (d) and \ref{stc:reg} and Propositions~\ref{stc:rm10} and \ref{stc:cubic} take a number $\mathsf b_+\ge\gamma^{-2}$, and some of them a number $\gamma^{-2}\le\mathsf b_-\le\mathsf b_+$, as free data; there $\mathsf b_+$ is any such number, and first-passage data at $g$ are the case $\mathsf b_-=\gamma^{-2}$, $\mathsf b_+=X_g+\bar B$ (Lemma~\ref{stc:frame}\,(a)).

A side $\ell'=2\mathsf l\in\mathfrak L(g)$, $\mathsf l=L^{m_1}$, is fixed once and for all; the least side at $g$ will do. $\gamma^{\rm ref}_W\le\gamma$ is the ceiling of the reference couplings of equation~\YMMref{(2326)}{845}, and $\gamma^{\rm sd}_W\le\gamma^{\rm ref}_W$ is the smaller ceiling of equation~\YMMref{(2327)}{846}; both are fixed in the choices in equations~\YMMref{(2326)--(2327)}{845}.

\emph{States.} For a side $\ell$, the torus $T_\ell=\mathbb R^3/\ell\mathbb Z^3$ carries the lattice of spacing $a_K$, the torus measure $\langle\cdot\rangle_{K,\ell,\beta}$ with the Wilson weight, its partition function $Z_{K,\ell,\beta}$ and the pressure $p_\ell=\ell^{-3}\log Z_{K,\ell,\beta}$; a function on the configurations of $\mathbb L_K$ that depends on the bonds of a cube of side $<\ell$ is regarded as a function on the configurations of $T_\ell$ through its periodisation. $\mathcal G^{\rm per}_{K,\beta}$ is the set of periodic Gibbs states on $\mathbb L_K$; they are invariant under the reflections in the lattice mirrors of the nine families $\{\mathsf n\cdot x=c\}$, $\mathsf n\in\{e_\mu,e_\mu\pm e_\lambda\}$, and reflection positive at these mirrors (\textup{\YMMref{103.4}{758}}); products of two such reflections give translations of the lattice. Torus measures are invariant under all isometries of the lattice of the torus. For a state $\nu$, $A^\nu[X]=X-\nu(X)$, $\langle X;Y\rangle_\nu=\nu(A^\nu[X]A^\nu[Y])$, and $\langle X_1;\dots;X_n\rangle_\nu$ is the joint cumulant. Over the partitions $\pi$ of $\{1,\dots,n\}$ we put $\mathsf{op}_n=\sum_\pi(|\pi|-1)!$ and $\mathsf{op}'_n=\sum_\pi|\pi|!$; $\mathsf{Bl}_n$ is the number of these partitions.

\emph{Test objects.} $\eta\in C^2_c(\mathbb R^3;\mathbb R^3)$ (or a $C^2$ vector field on a torus), with $\|\eta\|=\sup|\eta|+\sup|\nabla\eta|+\sup|\nabla^2\eta|$ (largest entries); $H:\mathbb R^3\to V={\rm Sym}_3(\mathbb R)$ bounded Borel (or $H:T_\ell\to V$), $|H|_\infty=\sup_x\max_{\lambda,\lambda'}|H_{\lambda\lambda'}(x)|$. Distances and diameters are sup distances. With $\Theta={\rm diag}(1,1,-1)$: $(\theta\eta)(x)=\Theta\eta(\Theta x)$, $(\theta H)(x)=\Theta H(\Theta x)\Theta$.

\emph{The Lie algebra.} $\mathfrak g$ is the space of Hermitian $2\times2$ matrices of trace $0$, with $\langle X,Y\rangle=\frac12{\rm tr}\,XY$ and $|X|^2=\langle X,X\rangle$; $\mathbf w(V)=\frac1{2i}(V-V^{-1})$ is the vector part of $V\in SL(2,\mathbb C)$. On matrices $|\cdot|$ is the operator norm; $|V^{-1}-\mathbf 1|=|V-\mathbf 1|$ for $V\in SL(2,\mathbb C)$.

\emph{Petals, bonds, differences.} For a site $x$ and signed directions $\alpha,\alpha'\in\{\pm1,\pm2,\pm3\}$ with $|\alpha|\ne|\alpha'|$, $W_x(\alpha,\alpha')$ is the holonomy of the plaquette loop from $x$, first along $e_\alpha$; the pair $\pi=(\alpha,\alpha')$ is a \emph{petal} at $x$, $m(\pi)=\{|\alpha|,|\alpha'|\}$ is its plane and $\mathsf s(\pi)=\pm1$ its orientation, and $q_x(\pi,\pi')=\langle\mathbf w(W_x(\pi)),\mathbf w(W_x(\pi'))\rangle$. For a positively oriented bond $\mathfrak b=\langle x,x+a_Ke_\mu\rangle$, $\mu=\mu_{\mathfrak b}$, with midpoint $x_{\mathfrak b}$, and a direction $\rho\ne\mu$, $p(\mathfrak b,\pm\rho)$ are the two plaquettes that contain $\mathfrak b$ in the plane of $e_\mu,e_\rho$; $W_{\mathfrak b,p}$ is the holonomy of the plaquette $p\ni\mathfrak b$ from $x$, starting with $\mathfrak b$, so that $W_{\mathfrak b,p(\mathfrak b,\pm\rho)}=W_x(\mu,\pm\rho)$; $\mathbf u_{\mathfrak b,\rho}=-\frac12[\mathbf w(W_x(\mu,\rho))-\mathbf w(W_x(\mu,-\rho))]$ and $\mathbf y_{\mathfrak b,\rho}=\mathbf w(W_x(\mu,\rho))+\mathbf w(W_x(\mu,-\rho))$; $\lambda_{\mathfrak b}<\lambda'_{\mathfrak b}$ are the two directions normal to $\mathfrak b$, and $q_{\mathfrak b}=\langle\mathbf u_{\mathfrak b,\lambda_{\mathfrak b}},\mathbf u_{\mathfrak b,\lambda'_{\mathfrak b}}\rangle$.

For a plaquette $p$, $x(p)$ is its centre, $\rho(p)$ its normal direction, $\mathsf P(p)$ the two directions of its plane; for a direction $\rho$, $p_\pm=p\pm a_Ke_\rho$ wherever $p_\pm$ stands beside the index $\rho$, as in $\mathfrak r^\sharp_\rho(p)$ and $\mathfrak e_\rho(p)$, so that at $\rho=\rho(p)$ it is the shift along the normal; and $p_m(x)$ is the plaquette of the plane $m$ with the lowest corner $x$. Differences: $(\nabla^\pm_\rho f)(y)=\pm a_K^{-1}[f(y\pm a_Ke_\rho)-f(y)]$, $(\nabla^{\rm c}_\rho f)(y)=a_K^{-1}[f(y+\frac12a_Ke_\rho)-f(y-\frac12a_Ke_\rho)]$, $(\nabla^{\rm s}_\rho f)(y)=(2a_K)^{-1}[f(y+a_Ke_\rho)-f(y-a_Ke_\rho)]$, $\nabla^{\rm s}\!\cdot\eta=\sum_\rho\nabla^{\rm s}_\rho\eta^\rho$, and $(\Delta_\rho f)(y)=4a_K^{-2}[f(y+\frac12a_Ke_\rho)-2f(y)+f(y-\frac12a_Ke_\rho)]$.

\begin{definition}[Lattice fields]\label{stc:fields}
The derivations $\mathfrak v_{\mathbf X}$ along fields of directions $\mathbf X$ on the bonds and their divergence ${\rm div}\,\mathfrak v_{\mathbf X}$, the orbital potential $\Psi_\rho(p)$, the orbital remainder $\mathfrak r_\rho(p)$ and its part of third order $\mathfrak r^\sharp_\rho(p)=\mathfrak r_\rho(p)-\Psi_\rho(p)+\Psi_\rho(p_-)$ are those of equation~\YMMref{(2103)}{781}, \ref{sec:ward-general-calculus}, \ref{sec:ward-general-calculus} and equations~\YMMref{(2104)--(2105)}{781}, Section~\ref{sec:ward-orbital-complement}; $\mathfrak v_\eta$ is the displacement along $\eta$.
\begin{enumerate}
\item[(1)] For a weight $w$ of the plaquettes $\mathcal O(w)=a_K^{-1}\sum_pw(p)a_p$; for a weight $\mathsf w$ of the positively oriented bonds $Y[\mathsf w]=-a_K^{-1}\sum_{\mathfrak b}\mathsf w(\mathfrak b)q_{\mathfrak b}$; and for a bounded Borel $H$
\[
\mathcal T_K(H)=\mathcal O(w_H)+Y[\mathsf w_H],\qquad w_H(p)=H_{\rho(p)\rho(p)}(x(p)),\quad\mathsf w_H(\mathfrak b)=H_{\lambda_{\mathfrak b}\lambda'_{\mathfrak b}}(x_{\mathfrak b}).
\]
Formally $\mathcal T_K(H)=\frac12\int{\rm Tr}(HT)$ with $T_{\lambda\lambda'}=\langle B_\lambda,B_{\lambda'}\rangle$, $B_\lambda=\frac12\varepsilon_{\lambda\mu\nu}F_{\mu\nu}$; and $\mathcal T_K(f\mathbf 1)=\mathcal O(w_{f\mathbf 1})$ is the action density field of the main theorems.
\item[(2)] $H^{\rm cons}[\eta]=(\nabla\!\cdot\!\eta)\mathbf 1-\nabla\eta-(\nabla\eta)^{\rm T}$, $(\nabla\eta)_{\nu\rho}=\partial_\nu\eta^\rho$, $H^{\rm c}_\eta=-H^{\rm cons}[\eta]$, and $\mathcal T^{\rm c}[\eta]=\mathcal T_K(H^{\rm cons}[\eta])$; formally $\mathcal T^{\rm c}[\eta]=-\int\eta_j\partial_iS_{ij}$.
\item[(3)] The insertion field is $\mathsf\Theta_\beta[\eta]=\beta a_K^{-1}\sum_p\sum_{\mathfrak b\subset\partial p}\sum_\rho\bigl(\eta^\rho(x_{\mathfrak b})-\eta^\rho(x(p))\bigr)\langle\mathbf u_{\mathfrak b,\rho},\mathbf w(W_{\mathfrak b,p})\rangle$, with $\mathbf u_{\mathfrak b,\mu_{\mathfrak b}}=0$. The general-field formula \eqref{eq:ward-general-displacement} at the constant weight $\beta$ splits $\mathfrak v_\eta(\beta\sum_pa_p)=\mathfrak K[\eta]+\mathsf{St}[\eta]$ into the \emph{orbital remainder} $\mathfrak K[\eta]$, the coupling of $\beta\eta$ to the orbital remainders $\mathfrak r_\rho(p)$, and the \emph{stress insertion} $\mathsf{St}[\eta]=\mathsf\Theta_\beta[\eta]-b\,\mathcal O(\nabla^{\rm s}\!\cdot\eta)$. The \emph{defect} is $\Delta[\eta]=\mathsf{St}[\eta]-b\,\mathcal T^{\rm c}[\eta]$.
\item[(4)] $\mathsf k^\rho_\eta=\nabla^+_\rho\eta^\rho$ and $w_\eta(p)=\sum_{\rho\in\mathsf P(p)}(\nabla^{\rm c}_\rho\eta^\rho)(x(p))$. For a complex weight $\varkappa$ of the plaquettes $\mathcal S_\varkappa=\sum_p\varkappa(p)\,\mathfrak r^\sharp_{\rho(p)}(p)$; a \emph{field} $\Xi$ is a triple of real functions on the centres of the plaquettes, $\mathsf H(\Xi)=\max_\rho\sup|\Xi_\rho|+\max_\rho{\rm Lip}\,\Xi_\rho$, and $\mathcal S[\Xi]=\mathcal S_\varkappa$ at $\varkappa(p)=L^{2K}\Xi_{\rho(p)}(x(p))$. The \emph{cubic half} is $\mathcal S[b\eta]=\sum_pL^{2K}b\,\eta^{\rho(p)}(x(p))\mathfrak r^\sharp_{\rho(p)}(p)$, the \emph{bilinear half} $\mathfrak J^{\rm II}[\mathsf k]=\beta\sum_{p,\rho}\mathsf k^\rho(x(p))\Psi_\rho(p)$.
\item[(5)] For a field of directions $\mathbf X$ that vanishes off a finite set $\Lambda$ of bonds, $\mathsf N_{\mathbf X}=\beta\,\mathfrak v_{\mathbf X}(\sum_pa_p)-{\rm div}\,\mathfrak v_{\mathbf X}$ is a \emph{null term}; $\mathcal N_\eta$ is the null term of $\mathfrak v_\eta$, and $\mathfrak N[\mathsf k]$ that of the field of directions $\sum_\rho\mathsf k^\rho(x_{\mathfrak b})\mathbf y_{\mathfrak b,\rho}$.
\item[(6)] The \emph{remainder} is $\mathsf{Rm}[\eta]=\mathcal S[b\eta]-\sum_{v=1}^5\mathfrak J_{0,v}[\mathsf k_\eta]+\frac32a_K\mathcal O(w_\eta)$, with the insertions $\mathfrak J_{0,v}$ of Lemma~\ref{stc:null}.
\end{enumerate}
\end{definition}

\begin{definition}[Insertions, rough insertions, oscillations]\label{stc:ins}
\begin{enumerate}
\item[(1)] A \emph{box} is $Q=\{x+a_K\sum_\mu n_\mu e_\mu:0\le n_\mu\le r_\mu\}$, $r\in\{0,1,2\}^3$ with at least two $r_\mu\ge1$; $x$ is its lowest site. A \emph{density} on $Q$ is a polynomial $\mathsf j$ in the entries of complex configurations $\mathbf U(\mathfrak b)\in SL(2,\mathbb C)$, $\mathbf U(\bar{\mathfrak b})=\mathbf U(\mathfrak b)^{-1}$, on the bonds of $Q$, invariant under the gauge transformations of $Q$; it has \emph{height} $\hat c$ if $|\mathsf j|\le\hat c$ whenever all $|\mathbf U(\mathfrak b)|\le2$, and \emph{order} $m$ with \emph{constant} $c$ if moreover $|\mathsf j|\le c\,r^m$ whenever all plaquette variables of $Q$, taken from one corner and in one sense, satisfy $|\partial\mathbf U(q)-\mathbf 1|\le r\le1$. An \emph{insertion} of type $(m,c,\hat c,N)$ is $\mathfrak J=\sum_\iota\varkappa_\iota\mathsf j_\iota(U|_{Q_\iota})$ with such densities on boxes $Q_\iota$, complex weights $\varkappa_\iota$, $|\varkappa|_\infty=\max_\iota|\varkappa_\iota|$, and at most $N$ indices with one lowest site ($N$ is the \emph{counting number}); indices of weight $0$ are struck before the type is formed, the \emph{carrier} is the set of the sites of the remaining boxes, and $\mathfrak J$ \emph{meets} a set $X$ if one of these boxes lies in $X$.
\item[(2)] A \emph{rough insertion} is $F(k,v)=a_K^{-1}\sum_x\sum_{\pi,\pi'}k_{\pi\pi'}(x)q_x(\pi,\pi')+\mathcal O(v)$, the sums over the sites and the petals at them, with $|k|_1=\sup_x\sum_{\pi,\pi'}|k_{\pi\pi'}(x)|$. Its \emph{symbol} at $x$ is the quadratic form in $F_{12},F_{13},F_{23}\in\mathfrak g$ obtained by putting $\mathsf s(\pi)F_{m(\pi)}$ for $\mathbf w(W_x(\pi))$ and $\frac12|F_m|^2$ for $a_{p_m(x)}$; $F(k,v)$ is \emph{null} if its symbol vanishes at every site.
\item[(3)] In the small-field setting (Th) of \ref{sec:cutoff-calibration} on a torus $T_\ell$, $\ell=2L^m$, with the data $(K,\mathsf b',m,t_*)$, a level $t$ with $1\le t\le t_*$ (in most uses the top level $t=t_*$) and the line $x_k=L^{K-k}\mathsf b'$: $\mathcal V_k$ is the space of the configurations of the block lattice of level $k$, $\mathcal V_0$ that of the configurations of the bare lattice, $\mathcal M^k:\mathcal V_0\to\mathcal V_k$ is the $k$-fold block averaging, $\mathsf{Cut}_k\supset\mathsf{In}_k$ are the field cuts (\YMMref{113.10}{796}), with the radii $\mathsf e^{\rm c}_k=c_\rho A_0x_k^{-1/2}(\log x_k)^{\mathfrak p_0}$ of the definitions preceding Lemma~\YMMref{113.10}{796}, $c_\rho$ the constant defined in Section~\ref{sec:cutoff-calibration} of the field cuts. For $0\le j<t_*$, $\mathfrak G_j[x_j]$ is the gauge-fixing factor, the \emph{gauge factor}, and $\chi^{\rm sf}_{x_j}$ the fluctuation test, the \emph{small-field indicator}, of the step of height $j$ at the inverse coupling $x_j$ (Section~\ref{sec:sf-chart-interface}), written $\mathfrak G_j$ and $\chi^{\rm sf}$ for short: functions on $\mathcal V_j$, the first bounded Borel with values in $[0,\infty)$, the second Borel with the values $0$ and $1$. The \emph{small-field weights} of the chain up to the level $k$, $0\le k\le t_*$, are
\[
\mathsf p_k(U)=\mathbf 1_{\mathsf{Cut}_k}(V_k)\prod_{j=0}^{k-1}\bigl[\mathbf 1_{\mathsf{Cut}_j}\,\mathfrak G_j[x_j]\,\chi^{\rm sf}_{x_j}\bigr](V_j),\qquad V_j=\mathcal M^jU,
\]
the product over $j<k$ of the indicator of $\mathsf{Cut}_j$, of the gauge factor and of the small-field indicator at the level $j$, times the indicator of $\mathsf{Cut}_k$: bounded Borel functions $\ge0$ on $\mathcal V_0$, with $\mathsf p_0=\mathbf 1_{\mathsf{Cut}_0}$ and $\mathsf p_{k+1}=(\mathbf 1_{\mathsf{Cut}_{k+1}}\circ\mathcal M^{k+1})\,((\mathfrak G_k[x_k]\chi^{\rm sf}_{x_k})\circ\mathcal M^k)\,\mathsf p_k$. The measures $\nu_k=(\mathcal M^k)_*(\mathsf p_k\rho_0dU)$ are those of the chain at zero source, given by the measure recursion in Proposition~\YMMref{114.2}{801} at $t_0=0$: $\nu_0=\mathbf 1_{\mathsf{Cut}_0}\rho_0dU$ and $\nu_{k+1}=\mathbf 1_{\mathsf{Cut}_{k+1}}\mathbf T(\mathfrak G_k[x_k]\chi^{\rm sf}_{x_k}\nu_k)$, $\mathbf T$ the image of complex measures under one block averaging $\mathcal M:\mathcal V_k\to\mathcal V_{k+1}$ (Proposition~\YMMref{114.2}{801}); this follows from the recursion of the $\mathsf p_k$, by induction on $k$ with the rules of images of measures, as in the proof of Proposition~\YMMref{115.2}{806}, its image-identity proof.

For $0\le k\le t_*$ let $\mathbf B_k$ be the determining set of the constant sequence $T_\ell=\Omega_0=\dots=\Omega_k$, whose averaging map is $\mathcal M^k$, and for $V\in\mathcal V_k$ let $\mathcal C(\mathbf B_k,V)$ be the set of the configurations $U$ with $\mathcal M^kU=V$ that are critical for the action under this constraint and regular (Definitions~\YMMref{7.1}{29}, \YMMref{9.3}{40} and~\YMMref{9.5}{41}). For $V\in\mathsf{Cut}_k$ this set is one orbit of the residual gauge group, formed by the gauge transformations of the bare lattice that equal $\mathbf 1$ at the sites of the block lattice of level $k$ (\YMMref{113.10}{796}; for $k=0$ it is $\{V\}$, by the definition preceding Proposition~\YMMref{114.2}{801}): $\mathcal C(\mathbf B_k,V)$ is the minimal orbit over $V\in\mathsf{Cut}_k$ (Section~\ref{sec:analytic-core-complements}), its elements are the \emph{backgrounds over} $V$, and $U_k(V)$, the background of $V$ (Section~\YMMref{9}{37}, immediately before equation~\YMMref{(77)}{42}), is any of them.

A configuration $U$ is \emph{over} $V$, and \emph{over} $\mathsf{Cut}_k$, if $V\in\mathsf{Cut}_k$ and $U\in\mathcal C(\mathbf B_k,V)$; so a configuration over $\mathsf{Cut}_k$ is a background over some $V\in\mathsf{Cut}_k$, and a \emph{background of level $k$} is a configuration over $\mathsf{Cut}_k$. A function of the bare configuration that is invariant under its gauge transformations takes one value at the configurations over $V$, that is, on the orbit, and so is a function of $V$ there. $\mathsf Q_{k\leftarrow0}Y$ is the conditional mean of a bounded Borel $Y$ given the block field of the level $k$ (Lemma~\ref{stc:cmean}). The \emph{small-field functional} is $\Lambda^{\rm sf}(F)=\int F\,(\mathbf 1_{\mathsf{In}_t}\circ\mathcal M^t)\,\mathsf p_t\,\rho_0\,dU$ for bounded Borel $F$; as $\mathsf{In}_t\subset\mathsf{Cut}_t$ \YMMref{(2159)}{796}, its weight $(\mathbf 1_{\mathsf{In}_t}\circ\mathcal M^t)\,\mathsf p_t$ is the weight $\mathfrak w^{\rm sf}_t$ of level $t$ of the small-field chain defined before Proposition~\YMMref{115.2}{806}; the image of the measure $F\,(\mathbf 1_{\mathsf{In}_t}\circ\mathcal M^t)\,\mathsf p_t\,\rho_0\,dU$ under $\mathcal M^t$ is $\mathbf 1_{\mathsf{In}_t}(\mathsf Q_{t\leftarrow0}F)\,\nu_t$ (Lemma~\ref{stc:cmean}\,(a)), and cumulants under the small-field functional are those of the normalised functional $\Lambda^{\rm sf}(\cdot)/\Lambda^{\rm sf}(1)$.

For a pointed family $\mathsf f$ among the construction results of level $j\le k$ (with the norm $\|\cdot\|$ of rate $\kappa+1$ in which Proposition~\YMMref{45.1}{353} and Lemma~\ref{lem:analytic-domain-geometry} control the sums over polymers: on the domain of the local terms its member at a polymer $X$ of level $j$ is at most $\|\mathsf f\|e^{-(\kappa+1)d_j(X)}$, $d_j(X)$ the linear size of $X$), $\mathsf{Sum}[\mathsf f](U)=\sum_X\sum_z\mathsf f(X,U,z)$ is the sum of its members at $U$, and $\Sigma_k[\mathsf f](V)=\mathsf{Sum}[\mathsf f](U_k(V))$ on $\mathsf{Cut}_k$, $0$ off it. The \emph{oscillation} of an insertion is $\omega_t(\mathfrak J)=\inf_{c\in\mathbb C}\|\mathsf Q_{t\leftarrow0}\mathfrak J-c\|_{L^\infty(\mathbf 1_{\mathsf{In}_t}\nu_t)}$.
\end{enumerate}
\end{definition}

\begin{definition}[Tensor limits]\label{stc:tlim}
Let $(K_j,\beta_j)\in\mathfrak d(g)$, $K_j\to\infty$, and $\rho_j\in\mathcal G^{\rm per}_{K_j,\beta_j}$. A functional $\sigma$ is a \emph{tensor limit at $g$ along $(K_j,\beta_j,\rho_j)$} if for all $n\ge1$ and all $H_1,\dots,H_n\in C_c(\mathbb R^3;V)$ with pairwise disjoint supports
\[
\sigma(H_1,\dots,H_n)=\lim_j\rho_j\Bigl(\prod_{a=1}^nA^{\rho_j}[\mathcal T_{K_j}(H_a)]\Bigr),
\]
the limits being assumed to exist. For the conservation statements in this section, tensor limit refers only to this moment convergence; convergence of the labels $a_{K_j}\beta_j$ is not required.
\end{definition}

\subsection{Results}

\begin{theorem}[Conservation among $n$ entries at finite cutoff]\label{stc:cutoff}
Let $0<g\le\gamma$, $n\ge2$, $1\le n_1\le n$, $R>0$. For $a\le n_1$ let $\eta_a\in C^2_c(\mathbb R^3;\mathbb R^3)$ and let $\mathsf X_a$ be one of $\mathcal T^{\rm c}[\eta_a]$, $\mathsf{St}[\eta_a]$, $\mathfrak K[\eta_a]$, $\Delta[\eta_a]$, $\mathsf{Rm}[\eta_a]$; for $n_1<a\le n$ let $H_a:\mathbb R^3\to V$ be bounded Borel and $\mathsf X_a=\mathcal T_K(H_a)$. Let the $n$ supports lie in a ball of radius $R$, pairwise at sup distance $\ge\mathsf d>0$, and put $\Pi=\prod_{a\le n_1}\|\eta_a\|\prod_{a>n_1}|H_a|_\infty$.
\begin{enumerate}
\item[(i)] For every $0<c<\frac1{14}$ there are $C,p,K_0$, functions of $g,\ell',n,\mathsf d,R,c$, such that for all $(K,\beta)\in\mathfrak d(g)$ with $K\ge K_0$ and all $\nu\in\mathcal G^{\rm per}_{K,\beta}$, if each $\mathsf X_a$, $a\le n_1$, is one of the first three,
\[
\Bigl|\nu\Bigl(\prod_{a=1}^nA^\nu[\mathsf X_a]\Bigr)\Bigr|\le CK^pL^{-cn_1K}\,\Pi .
\]
\item[(ii)] The bound of (i) holds also when $\Delta[\eta_a]$ or $\mathsf{Rm}[\eta_a]$ are among the $\mathsf X_a$, $a\le n_1$, and for the joint cumulant $\langle\mathsf X_1;\dots;\mathsf X_n\rangle_\nu$.
\item[(iii)] Let $0<c<\frac1{28}$ and $\phi=(\frac1{28}-c)/(4n-1)$. There are $C,p$, functions of $g,\ell',n,R,c$, and $K_0$, a function of $g,\ell',n,c$, such that the bounds of (i), (ii) hold for all $(K,\beta)\in\mathfrak d(g)$ with $K\ge K_0$, all $\nu\in\mathcal G^{\rm per}_{K,\beta}$ and all data as above whose supports are pairwise at sup distance $\ge\mathsf d$ for some $\mathsf d>0$ with $\mathsf d\,L^{\lceil\phi K\rceil}\ge K^2$.
\item[(iv)] Under the hypotheses of (i), also $|\nu(\mathfrak K[\eta_1])|\le CK^pL^{-cK}\|\eta_1\|$ and, if $n_1\ge2$, also $|\nu(\mathfrak K[\eta_1]\mathfrak K[\eta_2])|\le CK^pL^{-2cK}\|\eta_1\|\|\eta_2\|$.
\end{enumerate}
\end{theorem}

\begin{corollary}[Two points]\label{stc:two}
Let $0<g\le\gamma$, $0<c<\frac1{14}$, $R,\mathsf d>0$; let $\eta,\eta'\in C^2_c(\mathbb R^3;\mathbb R^3)$ and a bounded Borel $H':\mathbb R^3\to V$ have their supports in a ball of radius $R$, with ${\rm supp}\,\eta$ at sup distance $\ge\mathsf d$ from ${\rm supp}\,\eta'$ and from ${\rm supp}\,H'$. There are $C,p,K_0$, functions of $g,\ell',\mathsf d,R,c$, such that for all $(K,\beta)\in\mathfrak d(g)$ with $K\ge K_0$ and all $\nu\in\mathcal G^{\rm per}_{K,\beta}$:
\begin{enumerate}
\item[(a)] $|\langle\mathsf{St}[\eta'];\mathsf{St}[\eta]\rangle_\nu|=|\langle\mathsf{Rm}[\eta'];\mathsf{Rm}[\eta]\rangle_\nu|\le CK^pL^{-2cK}\|\eta'\|\|\eta\|$; $|\nu(\mathfrak K[\eta'])|\le CK^pL^{-cK}\|\eta'\|$; and $|\nu(\mathfrak K[\eta']\mathfrak K[\eta])|+|\nu(\mathfrak K[\eta'])\nu(\mathfrak K[\eta])|\le CK^pL^{-2cK}\|\eta'\|\|\eta\|$;
\item[(b)] $|\langle\mathcal T^{\rm c}[\eta'];\mathcal T^{\rm c}[\eta]\rangle_\nu|\le CK^pL^{-2cK}\|\eta'\|\|\eta\|$;
\item[(c)] (the collapse) 
\[
|\langle\Delta[\eta];\mathcal T_K(H')\rangle_\nu|\le CK^pL^{-cK}\|\eta\||H'|_\infty,
\]
 and $|\langle\Delta[\eta];Z\rangle_\nu|\le CK^pL^{-2cK}\|\eta\|\|\eta'\|$ for $Z=\mathsf{St}[\eta']$ and for $Z=\Delta[\eta']$;
\item[(d)] $|\langle\mathcal T^{\rm c}[\eta];\mathcal T_K(H')\rangle_\nu|\le CK^pL^{-cK}\|\eta\||H'|_\infty$.
\end{enumerate}
At $\eta'=\theta\eta$, $\eta\in C^2_c(\{x_3>0\};\mathbb R^3)$, the last bound of (a) bounds the mirrored two-point function $\langle\mathfrak K[\theta\eta];\mathfrak K[\eta]\rangle_\nu$ of the orbital remainder, which gives (N) by Theorem~\ref{stc:limit}\,(iii); and (b) is the lattice form of (N) below.
\end{corollary}

\begin{theorem}[Conservation among bounded tensor fields, by a change of scale]\label{stc:scale}
Let $0<g\le\gamma$, $q\ge2$, $\mathsf d>0$.
\begin{enumerate}
\item[(S)] Let $R>0$. There are $C,p,K_0$, functions of $g,q,\mathsf d,R$, such that for all $(K,\beta)\in\mathfrak d(g)$ with $K\ge K_0$, all $\nu\in\mathcal G^{\rm per}_{K,\beta}$, every $\eta\in C^2_c(\mathbb R^3;\mathbb R^3)$ and all bounded Borel $H_2,\dots,H_q:\mathbb R^3\to V$ whose $q$ supports lie in a ball of radius $R$, pairwise at sup distance $\ge\mathsf d$,
\[
\Bigl|\nu\Bigl(A^\nu[\mathcal T^{\rm c}[\eta]]\prod_{a=2}^qA^\nu[\mathcal T_K(H_a)]\Bigr)\Bigr|\le CK^pL^{-K/36}\|\eta\|\prod_{a=2}^q|H_a|_\infty ,
\]
and the same bound holds for the joint cumulant of the $q$ functions under $\nu$.
\item[(T)] Let $\ell\in\mathfrak L(g)$. There are $C,p,K_0$, functions of $g,q,\mathsf d,\ell$, such that for all $(K,\beta)\in\mathfrak d(g)$ with $K\ge K_0$, every $C^2$ vector field $\eta$ on $T_\ell$ and all bounded Borel $H_2,\dots,H_q:T_\ell\to V$ whose $q$ supports are pairwise at sup distance $\ge\mathsf d$ on $T_\ell$,
\[
|\langle\mathcal T_K(H^{\rm cons}[\eta]);\mathcal T_K(H_2);\dots;\mathcal T_K(H_q)\rangle_{K,\ell,\beta}|\le CK^pL^{-K/36}\|\eta\|\prod_{a=2}^q|H_a|_\infty .
\]
\end{enumerate}
\end{theorem}

\begin{proposition}[The remainder at every rate below $\frac1{10}$]\label{stc:rm10}
Let $m_1\in\mathbb N$, $\mathsf l=L^{m_1}$, $\gamma^{-2}\le\mathsf b_-\le\mathsf b_+$ with $MR(\mathsf b_+)\le L^{m_1}$ ($R$ the cube-size function of \YMMref{74}{614}), $r>0$, and let parameters $0<\psi<\frac29$, $0\le\kappa_\star<1$, $0<\tau<\frac12$, $0<\phi<\frac16\psi$ be given with
\[
c''=\min\bigl\{1-\tfrac92\psi,\ \kappa_\star-(3+\kappa_\star)\psi,\ (\tfrac12-\tau)\psi,\ \tfrac12\psi-3\phi\bigr\}>0 .
\]
\begin{enumerate}
\item[(a)] There are $C,p,K_\star$, functions of these data, such that for all $K\ge K_\star$, all $\beta=L^K\mathsf b'$ with $\mathsf b'\in[\mathsf b_-,\mathsf b_+]$, every board of eight cells of side $\mathsf l$ on $T_{2\mathsf l}$ with an offset in $a_K\mathbb Z^3$, every cell of it and every $\eta\in C^2(\mathbb R^3;\mathbb R^3)$ carried by the points of that cell at sup distance $\ge r$ from its walls, with $c_0=\langle\mathsf{Rm}[\eta]\rangle_{K,2\mathsf l,\beta}$,
\[
0\le\Bigl\langle\prod_\kappa\bigl(\mathsf{Rm}[\eta]-c_0\bigr)\circ\mathsf r_\kappa^{-1}\Bigr\rangle_{K,2\mathsf l,\beta}=\|\mathsf{Rm}[\eta]-c_0\|_{{\rm ch},2\mathsf l}^8\le\bigl(CK^pL^{-c''K}\|\eta\|\bigr)^8 ,
\]
the product over the eight mirror maps $\mathsf r_\kappa$ of the board. Moreover, let $u=\lceil\phi K\rceil$ and $t=K-u$, so that, $K_\star$ being taken $\ge K_\phi$, (Th) holds on $T_{2\mathsf l}$ with $(K,\mathsf b',m_1,t)$ (Lemma~\ref{stc:frame}\,(b), as $MR(\mathsf b')\le MR(\mathsf b_+)\le L^{m_1}$, $R$ being non-decreasing), and attach to the seven insertions of $\mathsf{Rm}[\eta]$ (Definition~\ref{stc:fields}\,(6)) the numbers
\[
\epsilon_1=C\mathsf b_+\mathsf H(\eta)K^pL^{-c''K},\qquad\epsilon_2=CK^{2\mathfrak p_0}L^{4u-K}\|\eta\|,\qquad\epsilon_3=CK^{3\mathfrak p_0}L^{4u-K/2}\|\eta\| ,
\]
$\epsilon_1$ to $\mathcal S[b\eta]$, $\epsilon_2$ to $\mathfrak J_{0,1}[\mathsf k_\eta]$, $\mathfrak J_{0,2}[\mathsf k_\eta]$, $\mathfrak J_{0,3}[\mathsf k_\eta]$ and $\frac32a_K\mathcal O(w_\eta)$, $\epsilon_3$ to $\mathfrak J_{0,4}[\mathsf k_\eta]$ and $\mathfrak J_{0,5}[\mathsf k_\eta]$, with $C\ge1$ and $p$ those of the display and $\mathsf H(\eta)=\max_\rho\sup|\eta^\rho|+\max_\rho{\rm Lip}\,\eta^\rho$ (Definition~\ref{stc:fields}\,(4)); these numbers depend neither on the board, its offset and its cell nor on the mirror map. For every mirror copy $\mathfrak J\circ\mathsf r_\kappa^{-1}$ of one of the seven insertions $\mathfrak J$, with the number $\epsilon$ attached to $\mathfrak J$, there is a number $\mathsf a$ with $|\mathsf Q_{t\leftarrow0}(\mathfrak J\circ\mathsf r_\kappa^{-1})-\mathsf a|\le\epsilon$ almost everywhere for $\mathbf 1_{\mathsf{In}_t}\nu_t$ (Definition~\ref{stc:ins}\,(3)); and every joint cumulant under $\langle\cdot\rangle_{K,2\mathsf l,\beta}$ of $2$ to $8$ such mirror copies, in pairwise different cells, is at most $1+\mathsf{op}'_8\mathsf{Bl}_8$ times the product of the numbers $\epsilon$ attached to the insertions of which they are copies.
\item[(b)] Without the cubic half: for $0<\phi<\frac18$ there are $C,K_\star$ with $\|\mathsf{Rm}[\eta]-\mathcal S[b\eta]-c_0\|_{{\rm ch},2\mathsf l}\le CK^{3\mathfrak p_0}L^{-(\frac12-4\phi)K}\|\eta\|$ in the situation of (a), the torus mean of $\mathcal S[b\eta]$ being $0$; and $\|\mathcal S[\Xi]\|_{{\rm ch},2\mathsf l}\le C\mathsf H(\Xi)K^pL^{-c''K}$ for every field $\Xi$ carried by a cell at sup distance $\ge r$ from its walls.
\item[(c)] Let $0<g\le\gamma$, $2\mathsf l\in\mathfrak L(g)$, $\mathsf b_-=\gamma^{-2}$, $\mathsf b_+=X_g+\bar B$. Then (a), (b) hold for all $(K,\beta)\in\mathfrak d(g)$ with $K\ge K_\star$; at $\psi=\frac15$, $\kappa_\star=\frac9{10}$, $c''=\frac1{10}-\max\{\frac15\tau,3\phi\}$, so every $c''<\frac1{10}$ is reached. For every $c<\frac1{10}$ and $R,\mathsf d>0$ there are $C,p,K_0$, functions of $g,c,R,\mathsf d$, such that for all $(K,\beta)\in\mathfrak d(g)$ with $K\ge K_0$, all $\nu\in\mathcal G^{\rm per}_{K,\beta}$ and all $\eta\in C^2_c(\{x_3>0\};\mathbb R^3)$ with support in the ball of radius $R$ at sup distance $\ge\frac12\mathsf d$ from $\{x_3=0\}$:
\[
|\langle\mathsf{St}[\theta\eta];\mathsf{St}[\eta]\rangle_\nu|=|\langle\mathsf{Rm}[\theta\eta];\mathsf{Rm}[\eta]\rangle_\nu|\le CK^pL^{-2cK}\|\eta\|^2,\qquad|\langle\mathfrak K[\theta\eta];\mathfrak K[\eta]\rangle_\nu|\le CK^pL^{-2cK}\|\eta\|^2 .
\]
\end{enumerate}
\end{proposition}

\begin{theorem}[The limit]\label{stc:limit}
Let $0<g\le\gamma$, let $\sigma$ be a tensor limit at $g$ along $(K_j,\beta_j,\rho_j)$, and $b_j=a_{K_j}\beta_j$.
\begin{enumerate}
\item[(i)] $\sigma(H^{\rm cons}[\eta],H_2,\dots,H_n)=0$ for every $n\ge2$, every $\eta\in C^2_c(\mathbb R^3;\mathbb R^3)$ and all $H_2,\dots,H_n\in C_c(\mathbb R^3;V)$ whose supports are pairwise disjoint and disjoint from ${\rm supp}\,\eta$.
\item[(ii)] In particular: (C) $\sigma(H^{\rm cons}[\eta],H')=0$ whenever ${\rm supp}\,\eta\cap{\rm supp}\,H'=\emptyset$; (N) 
\[
\sigma(\theta H^{\rm cons}[\eta],H^{\rm cons}[\eta])=0
\]
 for $\eta\in C^2_c(\{x_3>0\};\mathbb R^3)$.
\item[(iii)] For $\eta\in C^2_c(\{x_3>0\};\mathbb R^3)$, with $\gamma^{-2}\le b_j\le\mathsf b_+$,
\[
b_j^2\,\langle\mathcal T^{\rm c}[\theta\eta];\mathcal T^{\rm c}[\eta]\rangle_{\rho_j}-\langle\mathfrak K[\theta\eta];\mathfrak K[\eta]\rangle_{\rho_j}\to0 .
\]
So $\langle\mathfrak K[\theta\eta];\mathfrak K[\eta]\rangle_{\rho_j}\to0$, which holds by Corollary~\ref{stc:two}\,(a), gives (N) by itself; and (N) for $\sigma$ gives $\langle\mathfrak K[\theta\eta];\mathfrak K[\eta]\rangle_{\rho_j}\to0$ and $\langle\mathsf{St}[\theta\eta];\mathsf{St}[\eta]\rangle_{\rho_j}\to0$ along $(K_j,\beta_j,\rho_j)$. The comparison does not use Proposition~\ref{stc:cubic}.
\end{enumerate}
No positivity, covariance or growth of $\sigma$ is used.
\end{theorem}

\begin{theorem}[One torus; the cutoff removed first]\label{stc:torus}
\begin{enumerate}
\item[(i)] Let $0<g\le\gamma$, $\ell\in\mathfrak L(g)$, $q\ge2$, $\mathsf d>0$, $0<c<\frac1{28}$. There are $C,p,K_0$, functions of $g,\ell,q,\mathsf d,c$, such that for all $(K,\beta)\in\mathfrak d(g)$ with $K\ge K_0$, every $C^2$ vector field $\eta$ on $T_\ell$ and all bounded Borel $H_2,\dots,H_q:T_\ell\to V$ such that ${\rm supp}\,\eta,{\rm supp}\,H_2,\dots,{\rm supp}\,H_q$ are pairwise at sup distance $\ge\mathsf d$ on $T_\ell$,
\[
|\langle\mathcal T_K(H^{\rm cons}[\eta]);\mathcal T_K(H_2);\dots;\mathcal T_K(H_q)\rangle_{K,\ell,\beta}|\le CK^pL^{-cK}\|\eta\|\prod_{a=2}^q|H_a|_\infty .
\]
\item[(ii)] Let $0<g\le\gamma^{\rm ref}_W$, $\ell\in\mathfrak L(g)$ and $b\in I(g)$. Then $(K,L^Kb)\in\mathfrak d(g)$ for every $K$, and for $\eta,H_2,\dots,H_q$ as in (i) the cumulants of (i) at $\beta=L^Kb$ tend to $0$ as $K\to\infty$. For Lipschitz $H_2,\dots,H_q$ whose supports are pairwise disjoint and disjoint from ${\rm supp}\,\eta$ this says that the torus limit functional $\mathcal S^{\ell,b}_q$, the limit of the joint cumulants $\langle\mathcal T_K(H_1);\dots;\mathcal T_K(H_q)\rangle_{K,\ell,L^Kb}$ at Lipschitz weights with pairwise disjoint supports (statement (iii) of the subsection on the input), vanishes:
\[
\mathcal S^{\ell,b}_q(H^{\rm c}_\eta,H_2,\dots,H_q)=0,\qquad H^{\rm c}_\eta=\nabla\eta+(\nabla\eta)^{\rm T}-(\nabla\!\cdot\!\eta)\mathbf 1 .
\]
In particular, with scalar spectators, $\mathcal S^{\ell,b}_{m+1}(H^{\rm c}_\eta,h_1\mathbf 1,\dots,h_m\mathbf 1)=0$ for every $m\ge1$ and all real Lipschitz $h_1,\dots,h_m$ on $T_\ell$ whose supports are pairwise disjoint and disjoint from ${\rm supp}\,\eta$, at every pair $(\ell,b)$ with $\ell\in\mathfrak L(g)$, $b\in I(g)$.
\end{enumerate}
\end{theorem}

\begin{corollary}[Tensor limits at small coupling]\label{stc:small}
Let $0<g\le\gamma^{\rm ref}_W$.
\begin{enumerate}
\item[(i)] Every sequence $(K_j,\beta_j,\nu_j)$ with $(K_j,\beta_j)\in\mathfrak d(g)$, $K_j\to\infty$, $\nu_j\in\mathcal G^{\rm per}_{K_j,\beta_j}$, has a subsequence along which $a_{K_j}\beta_j$ converges to some $b_*\in[g^{-2},L(g^{-2}+\bar B)]$ and which defines a tensor limit. A subsequence of a defining sequence defines the same tensor limit; so every tensor limit at $g$ is a limit along a sequence with convergent labels. Every tensor limit $\sigma$ at $g$ is symmetric in its entries and vanishes at a single entry; it is covariant, $\sigma(\mathsf rH_1,\dots,\mathsf rH_n)=\sigma(H_1,\dots,H_n)$ for every $\mathsf rx=\mathsf Rx+v$, $\mathsf R$ in the cubic group, $v\in\mathbb R^3$, $(\mathsf rH)(x)=\mathsf RH(\mathsf r^{-1}x)\mathsf R^{\rm T}$; for weights that vanish off closed subsets of a ball of radius $R_0$ at pairwise sup distances $\ge2\delta_0$, $|\sigma(H_1,\dots,H_n)|\le C_n\prod_a|H_a|_\infty$ with $C_n$ a function of $g,n,\delta_0,R_0$, and the same holds with a number that depends on $g,n,\delta_0$ and the sup diameters of the single supports, not on $R_0$; and $\sigma(f_1\mathbf 1,\dots,f_n\mathbf 1)$ is the limit, along the same sequence, of the centred products of the action density fields $\mathcal O(w_{f_a\mathbf 1})$ of the main theorems.

The limits at $g$ of statement (iv) of the subsection on the input are tensor limits in the sense of Definition~\ref{stc:tlim}.
\item[(ii)] Every tensor limit $\sigma$ at $g$ is reflection positive at the nine families of mirrors: $\sigma(\theta_{\mathsf n,c}\bar F\cdot F)\ge0$ for $\mathsf n\in\{e_\mu,e_\mu\pm e_\lambda\}$, $c\in\mathbb R$, $\theta_{\mathsf n,c}$ the reflection in $\{\mathsf n\cdot x=c\}$ acting on arguments and indices, and $F$ a finite complex combination of tuples of $C^1_c$ tensor fields with pairwise disjoint supports in $\{\mathsf n\cdot x>c\}$. The form $\langle[F'],[F]\rangle=\sigma(\theta_{e_3,0}\bar F'\cdot F)$ gives a separable Hilbert space $\mathcal H(\sigma)$ with a unit vector $\Omega=[\emptyset]$, a strongly continuous semigroup $e^{-tH}[F]=[\tau_{te_3}F]$, $t\ge0$, of self-adjoint contractions, $H\ge0$, $H\Omega=0$, and a strongly continuous unitary representation $e^{i(v_1P_1+v_2P_2)}[F]=[\tau_{(v,0)}F]$ of the translations parallel to the mirror, which commutes with $H$ and fixes $\Omega$; and $H\ge|P_1|$, $H\ge|P_2|$. Here $\tau_v$ translates every entry of a tuple by $v$.
\item[(iii)] For $\eta\in C^2_c(\{x_3>0\};\mathbb R^3)$ and every such $F$ in $\{x_3>0\}$ whose supports do not meet ${\rm supp}\,\eta$, the vector $[(H^{\rm cons}[\eta])\cdot F]\in\mathcal H(\sigma)$ is zero; at $F=\emptyset$ this is (N). In particular $\sigma(\theta_{e_3,0}\bar F'\cdot(H^{\rm cons}[\eta])\cdot F)=0$ for every such $F'$; for tuples separated by the plane $\{x_3=0\}$ this is contained in Theorem~\ref{stc:limit}\,(i), which holds in general position and without the Hilbert space.
\end{enumerate}
\end{corollary}

\begin{proposition}[A second road: null forms]\label{stc:second}
For weights $v=(v_\rho(\mathfrak b))_{\rho\ne\mu_{\mathfrak b}}$ and $z=(z_{\rho\tau}(\mathfrak b))_{\rho,\tau\ne\mu_{\mathfrak b}}$ of the positively oriented bonds and $w$ of the plaquettes, with $a_\pm=a_{p(\mathfrak b,\pm\rho)}$ and $\rho'$ the direction normal to $e_{\mu_{\mathfrak b}}$ and $e_\rho$, put
\[
\begin{multlined}
\mathcal Y[v]=a_K^{-1}\sum_{\mathfrak b}\sum_{\rho\ne\mu_{\mathfrak b}}v_\rho(\mathfrak b)|\mathbf y_{\mathfrak b,\rho}|^2,\quad\mathcal A[v]=a_K^{-1}\sum_{\mathfrak b,\rho}v_\rho(\mathfrak b)(a_+^2+a_-^2),\quad\mathcal A^2[w]=a_K^{-1}\sum_pw(p)a_p^2, \\
\mathcal R[z]=a_K^{-1}\sum_{\mathfrak b}\sum_{\rho,\tau}z_{\rho\tau}(\mathfrak b)\langle\mathbf u_{\mathfrak b,\rho},\mathbf y_{\mathfrak b,\tau}\rangle ;
\end{multlined}
\]
the \emph{planar null form} is $\mathcal P_K[v,w]=\mathcal Y[v]+\mathcal A^2[w]$, with $|(v,w)|_\infty=\max\{\sup|v|,\sup|w|\}$; for six bounded real functions $\kappa=(\kappa_{\mu\nu})_{\mu\ne\nu}$, $\mathfrak Y[\kappa]=\frac\beta4\sum_{\mathfrak b}\sum_{\nu\ne\mu_{\mathfrak b}}\kappa_{\mu_{\mathfrak b}\nu}(x_{\mathfrak b})|\mathbf y_{\mathfrak b,\nu}|^2=\mathcal P_K[\frac b4v^\kappa,0]$, $v^\kappa_\nu(\mathfrak b)=\kappa_{\mu_{\mathfrak b}\nu}(x_{\mathfrak b})$. For a vector field $\eta$: $v^\eta_\rho(\mathfrak b)=(\nabla^{\rm c}_\rho\eta^\rho)(x_{\mathfrak b})$, $z^\eta_{\rho\tau}(\mathfrak b)=\frac{a_K}8(\Delta_\tau\eta^\rho)(x_{\mathfrak b})$, $\mathfrak Y[\eta]=\mathfrak Y[\kappa]$ at $\kappa_{\mu\nu}=\nabla^{\rm c}_\nu\eta^\nu$; and $H^{\rm cons}_K[\eta]$ is the lattice symbol whose diagonal entry $\lambda\lambda$ at $x(p)$, $\lambda=\rho(p)$, is $(\nabla^{\rm s}\!\cdot\eta-2\nabla^{\rm s}_\lambda\eta^\lambda)(x(p))$ and whose entry $\lambda_{\mathfrak b}\lambda'_{\mathfrak b}$ at $x_{\mathfrak b}$ is $-(\nabla^{\rm c}_{\lambda'_{\mathfrak b}}\eta^{\lambda_{\mathfrak b}}+\nabla^{\rm c}_{\lambda_{\mathfrak b}}\eta^{\lambda'_{\mathfrak b}})(x_{\mathfrak b})$, inserted in $\mathcal T_K$ as in Definition~\ref{stc:fields}\,(1); $H''_K[\eta]=H^{\rm cons}_K[\eta]-H^{\rm cons}[\eta]$ has $|H''_K[\eta]|_\infty\le\frac32a_K\sup|\nabla^2\eta|$.
\begin{enumerate}
\item[(a)] (An exact identity.) For every $K$, $\beta>0$, every vector field $\eta$ of bounded support (or on a torus) and every configuration,
\[
\mathsf{St}[\eta]=b\,\mathcal T_K(H^{\rm cons}_K[\eta])-b\bigl\{\tfrac14\mathcal Y[v^\eta]+\tfrac12\mathcal A[v^\eta]+\mathcal R[z^\eta]\bigr\}.
\]
\item[(b)] (One bond.) Let $\nu$ be a Gibbs state at $(a_K,\beta)$, $z$ a weight of bounded support, $Y$ a bounded Borel local function that does not depend on the bonds $\mathfrak b$ with $z(\mathfrak b)\ne0$; let $\zeta_\rho(\mathfrak b)=z_{\rho\rho}(\mathfrak b)-z_{\rho\rho'}(\mathfrak b)$, $\zeta'_\rho(\mathfrak b)=z_{\rho\rho'}(\mathfrak b)$, and $(\delta\xi)(p)=\sum_{(\mathfrak b,\rho):\,p(\mathfrak b,\rho)=p}\xi_\rho(\mathfrak b)-\sum_{(\mathfrak b,\rho):\,p(\mathfrak b,-\rho)=p}\xi_\rho(\mathfrak b)$ for a weight $\xi$ of the pairs. Then
\[
\langle\mathcal R[z];Y\rangle_\nu=\bigl\langle\mathcal O\bigl(\delta[-\zeta+\tfrac3{2\beta}\zeta']\bigr);Y\bigr\rangle_\nu+\tfrac12\langle\mathcal A^2[\delta\zeta];Y\rangle_\nu .
\]
\item[(c)] (Eight copies of a null form.) Let $0<g\le\gamma$, $\ell\in\mathfrak L(g)$, $r>0$, and let parameters $0<\kappa_\star<1$, $0<\phi<\psi<1$, $\phi\le\frac1{84}$, have $c_{\rm Y}=\min\{2\kappa_\star-(3+2\kappa_\star)\psi,\ \frac12\psi-4\phi\}>0$. There are $C,K_0$, functions of $g,\ell,r$ and these parameters, such that for all $(K,\beta)\in\mathfrak d(g)$ with $K\ge K_0$, every board of eight cells of side $\frac12\ell$ on $T_\ell$ with an offset in $a_K\mathbb Z^3$ and all real weights $(v,w)$ carried by the points of one cell at sup distance $\ge r$ from its walls, with $F=\mathcal P_K[v,w]$ and $c=\langle F\rangle_{K,\ell,\beta}$: every joint cumulant of $q\ge2$ mirror copies of $F$ in different cells is at most $(\frac32+\mathsf{op}'_q\mathsf{Bl}_q)(CK^{2\mathfrak p_0}L^{-c_{\rm Y}K}|(v,w)|_\infty)^q$, and
\[
0\le\Bigl\langle\prod_\kappa(F-c)\circ\mathsf r_\kappa^{-1}\Bigr\rangle_{K,\ell,\beta}\le\bigl(CK^{2\mathfrak p_0}L^{-c_{\rm Y}K}|(v,w)|_\infty\bigr)^8 .
\]
For every $c'<\frac2{11}$ there are parameters with $c_{\rm Y}>c'$; at $(\kappa_\star,\psi,\phi)=(\frac12,\frac18,\frac1{256})$, $c_{\rm Y}=\frac3{64}$. For $\mathfrak Y[\kappa]$ the same holds with $\frac14\mathsf b_+\sup|\kappa|$ for $|(v,w)|_\infty$.
\item[(d)] (The collapse in the limit.) Let $0<g\le\gamma^{\rm sd}_W$, and take for the fixed side $\ell'=2\mathsf l$ the least side at $g$. For $0<r\le\frac1{16}\mathsf l$ there are $C,K_0$, functions of $g,r$, such that for all $(K,\beta)\in\mathfrak d(g)$ with $K\ge K_0$, all $\nu\in\mathcal G^{\rm per}_{K,\beta}$, every direction $\mu_0$, every $c\in a_K(\mathbb Z+\frac12)$, $\theta_c$ the reflection in $\{x_{\mu_0}=c\}$, and all real weights $(v,w)$ carried in a cube of side $r$ in $\{c+2r-2a_K\le x_{\mu_0}<c+3r\}$ or in its image under $\theta_c$,
\[
\|\mathcal P_K[v,w]\|_{\nu,c}:=\langle\theta_c\mathcal P_K[v,w];\mathcal P_K[v,w]\rangle_\nu^{1/2}\le CK^{2\mathfrak p_0}L^{-3K/64}|(v,w)|_\infty ;
\]
and for every tensor limit $\sigma$ at $g$ along $(K_j,\beta_j,\rho_j)$ with $b_j\to b_*$, every $\eta\in C^2_c(\mathbb R^3;\mathbb R^3)$ and every $H'\in C_c(\mathbb R^3;V)$ with disjoint supports,
\[
\lim_j\langle\mathsf{St}[\eta];\mathcal T_{K_j}(H')\rangle_{\rho_j}=b_*\,\sigma(H^{\rm cons}[\eta],H').
\]
Proposition~\ref{stc:cubic} is not used for (d).
\item[(e)] ((N) at a lattice mirror.) Let $0<g\le\gamma^{\rm sd}_W$, $h_0,R_0>0$ and $c<\frac1{10}$. There are $C,p,K_0$, functions of $g,h_0,R_0,c$, such that for all $(K,\beta)\in\mathfrak d(g)$ with $K\ge K_0$, all $\rho\in\mathcal G^{\rm per}_{K,\beta}$ and all $\eta\in C^2(\mathbb R^3;\mathbb R^3)$ with support in $\{x_3\ge h_0\}\cap\{|x|_\infty\le R_0\}$, $\theta_K$ being the reflection in the bond mirror $\{x_3=\frac12a_K\}$,
\[
0\le\bigl\langle\theta_K\mathcal T_K(H^{\rm cons}_K[\eta]);\mathcal T_K(H^{\rm cons}_K[\eta])\bigr\rangle_\rho\le b^{-2}\bigl(CK^pL^{-cK}\|\eta\|\bigr)^2 .
\]
This gives (N) again for every tensor limit at $g$.
\end{enumerate}
\end{proposition}

\begin{proposition}[Two instances]\label{stc:free}
\begin{enumerate}
\item[(a)] (The free value.) Let $K_3(u)=\nabla\nabla\frac1{4\pi|u|}$ be the matrix of second derivatives, $y_1,y_2\in\mathbb R^3$, $B$ a real $3\times3$ matrix and $M(x)=K_3(x-y_1)\,B\,K_3(x-y_2)$. On $\mathbb R^3\setminus\{y_1,y_2\}$
\[
\sum_\nu\partial_\nu(M_{\nu\rho}+M_{\rho\nu})=\partial_\rho\,{\rm tr}\,M .
\]
Hence $\int{\rm tr}[H^{\rm cons}[\eta](x)M(x)]\,dx=0$ for $\eta\in C^2_c(\mathbb R^3\setminus\{y_1,y_2\};\mathbb R^3)$, and for three free photon fields, $\langle B^{\mathsf a}_\alpha(x)B^{\mathsf b}_\gamma(y)\rangle=\delta^{\mathsf a\mathsf b}K_3(x-y)_{\alpha\gamma}$ for $x\ne y$, and for bounded $H_1,\dots,H_q$ with compact supports that are pairwise disjoint, the value of the joint cumulant of $\frac12\int{\rm Tr}(H_aT)$, $a=1,\dots,q$,
\[
\tfrac32\sum_{\pi\in S_{q-1}}\int\!\!\cdots\!\!\int{\rm tr}\bigl[H_1(x_1)K_3(x_1-x_{\pi2})H_{\pi2}(x_{\pi2})K_3(x_{\pi2}-x_{\pi3})\cdots H_{\pi q}(x_{\pi q})K_3(x_{\pi q}-x_1)\bigr]\,dx_1\cdots dx_q ,
\]
in which every integrand is bounded, vanishes at $H_1=H^{\rm cons}[\eta]$ for every $q\ge2$, every $\eta\in C^2_c(\mathbb R^3;\mathbb R^3)$ and all bounded $H_2,\dots,H_q$ with compact supports that are pairwise disjoint and disjoint from ${\rm supp}\,\eta$.
\item[(b)] (Stress bounds without (N).) On $\mathbb L_K$ let $(U,\chi_K)$ have the law $\nu\otimes\gamma_K$, $\nu\in\mathcal G^{\rm per}_{K,\beta}$, $\gamma_K$ the free lattice scalar field of mass $1$, and $\mathcal T^{\rm X}_K(H)=\mathcal T_K(H)+\epsilon\,\mathsf X_K[H]$, $\mathsf X_K[H]=a_K^3\sum_x\sum_\lambda H_{\lambda\lambda}(x)(\nabla^-_\lambda\nabla^+_\lambda\chi_K)(x)$, $\epsilon\ne0$. The functions $\mathsf{St}[\eta]$, $\mathsf{Rm}[\eta]$, $\mathfrak K[\eta]$ depend on $U$ alone, so Corollary~\ref{stc:two}\,(a) holds word for word under $\nu\otimes\gamma_K$; but $b\,\mathcal T^{\rm X}_K(H^{\rm cons}[\eta])-\mathsf{St}[\eta]=-\Delta[\eta]+b\epsilon\,\mathsf X_K[H^{\rm cons}[\eta]]$, and the last member is no function of the gauge field. At $\eta=he_3$, $0\le h\in C^\infty_c(\{x_3>0\})$, $h\ne0$, the mirrored two-point function $\gamma_K(\mathsf X_K[H^{\rm cons}[\theta\eta]]\,\mathsf X_K[H^{\rm cons}[\eta]])$ tends, as $K\to\infty$, to
\[
\begin{multlined}
\int|\hat h(p)|^2E(p)^2\bigl(2E(p)^2-1\bigr)^2\,d\rho(p)>0,\qquad E(p)=(1+|p|^2)^{1/2},\quad d\rho(p) \\
=\frac{d^2p}{(2\pi)^2\,2E(p)},\quad\hat h(p)=\int e^{-E(p)x_3+ip\cdot(x_1,x_2)}h(x)\,dx .
\end{multlined}
\]
So along every sequence of first-passage data and states as in Corollary~\ref{stc:two} the mirrored two-point function of $\mathcal T^{\rm X}_K(H^{\rm cons}[\eta])$ tends to $\epsilon^2$ times this positive number: the dressed field violates (N), while the bounds for $\mathsf{St}$, $\mathsf{Rm}$, $\mathfrak K$ hold.
\end{enumerate}
\end{proposition}

\subsection{Input}

Lemma~\YMMref{118.1}{812} supplies the constants satisfying (Th1), (Th2) of the calibrated small-field setting (Th) of \ref{sec:cutoff-calibration}, whose admissible top levels are those of \ref{sec:cutoff-calibration}, with the number $c_{\log}=\log L/\log\gamma^{-2}$ of Section~\ref{sec:cutoff-calibration}. In this setting we use one renormalisation step with its locality (Proposition~\ref{prop:source-entry}, with the admitted local data and the radius $\varepsilon_{\rm src}$ of the sources \eqref{eq:source-entry-domain}) and with its parity: if the four families of a step depend holomorphically on a parameter $\mu$, the other data of the step not, and the reflection $B\mapsto-B$ of the fluctuation field carries the four families at $\mu$ into those at $-\mu$, the output of the step is holomorphic, bounded and even in $\mu$ (Proposition~\ref{prop:source-entry-parity}), and the three interaction families of a step are families in the fourth place of size at most $K_P\delta(x_k)$ (\eqref{eq:source-internal-size}), which bounds their even and odd parts; the potentials of local functions (equation~\YMMref{(870)}{410}), with their sum rule and decay (Proposition~\YMMref{51.15}{410}, equations~\YMMref{(872)}{411} and~\YMMref{(871)}{410}); the chain on one torus with its runs and its numbers (\YMMref{114}{799}, \YMMref{(2163)}{798}, \YMMref{(2176)}{803}), its exponent and measure recursion in Proposition~\YMMref{114.2}{801}, the bound in Proposition~\YMMref{114.2}{801} on its families, with the step identity and the actions on a field cut there, the field cuts (\YMMref{113.10}{796}, \YMMref{113.10}{796}, the definitions preceding Lemma~\YMMref{113.10}{796}), the gauge-fixing factor and the fluctuation test of a step (Section~\ref{sec:sf-chart-interface}), the small-field weights, the weight defined before Proposition~\YMMref{115.2}{806} of the small-field chain and the chart integrals, and the discarding of large fields (\textup{\YMMref{115.2}{806}}) with the number $\mathsf R_t$ of \YMMref{115.2}{806}; the sums over polymers of Proposition~\YMMref{45.1}{353} and Lemma~\ref{lem:analytic-domain-geometry}; minimal orbits and backgrounds (Definitions~\YMMref{9.3}{40} and~\YMMref{9.5}{41}; Section~\ref{sec:analytic-core-complements}); local terms and their evaluation in a local gauge (Section~\ref{sec:regular-background-complements}), with the usable radius \YMMref{(933)}{427}; the number $\varepsilon_{\mathcal W}$ (Proposition~\YMMref{52.3}{433}), regular configurations and the current of a level (Definition~\ref{def:level-regularity}); local regular gauges (Theorem~\YMMref{F}{199}, with its constants $c_{*,\mathrm{rg}}=c_*$ and $B_*$) and the gauge of a regular configuration on a cube about a polymer (Lemma~\ref{lem:regular-cube-gauge}); and local gauges with a H\"older bound on cubes (Theorem~\textup{\ref{I-56-65}}).

Theorems~\YMMref{AJ}{673} and~\YMMref{AD}{607} enter only through the transfer statement \ref{lem:canonical-source-system} with \eqref{eq:source-system-pressure}, taken at zero source: it bounds the difference of the pressures of two sides at $g$ uniformly in the cutoff and in the larger side. We use besides the labels and the sides at a reference coupling (Corollary~\YMMref{74.2}{616}, \YMMref{(1535)}{616}, the side definition~\YMMref{(1528)}{615}, at the common choice of Section~\YMMref{82}{665}); the lattice fields of equation~\YMMref{(2103)}{781}, \ref{sec:ward-general-calculus}, \ref{sec:ward-general-calculus} and equations~\YMMref{(2104)--(2105)}{781}, Section~\ref{sec:ward-orbital-complement}, with the densities \eqref{eq:ward-orbital-normal-potential}, \eqref{eq:ward-orbital-inplane}; the general-field formula \eqref{eq:ward-general-displacement}, the calculus \ref{lem:ward-general-calculus} and the divergence \eqref{eq:ward-general-divergence} of the derivations, the summation by parts in the orbital remainder and its covariance (Proposition~\ref{prop:ward-orbital-cubic}), and the vanishing of its torus means (\ref{C-360-8}(b)); the equations of Dobrushin, Lanford and Ruelle for Gibbs states (\YMMref{(1819)}{698}); the chessboard seminorm on tori (Section~\ref{sec:periodic-chessboard}), the chessboard estimate and the comparison of two sides (\ref{IV-18-18}), and the domination of periodic Gibbs states (Proposition~\YMMref{91.9}{705}); the reflection positivity and invariance of periodic Gibbs states (\textup{\YMMref{103.4}{758}}); and the invariance of torus measures under lattice isometries.

For the chain with insertions, its sources and its covariance under lattice isometries we use further: the bound of the action of a bare array at a background field (Proposition~\YMMref{113.13}{799}), with the number $C_{\rm old}$ of equation~\YMMref{(2163)}{798}; the loss of memory of a run, that is, its Lipschitz bound in its early members (the proof of Proposition~\YMMref{114.3}{803}); the bound $\frac{15}{16}E_0$ on the family produced by every run whose early members have norms $\le E_0$ (the proof of Proposition~\YMMref{114.2}{801}), with the constant \YMMref{(2173)}{801} of a step and the step identity \YMMref{(2174)}{801}; the classes of covariance, with their heredity, their window and their invariance (Lemma~\ref{lem:sf-class-window}); the covariance of the chart, of its Gaussian normaliser and of the minimal orbits under the isometries of the coarser lattices (the chart and push-forward calculation in the proof of Proposition~\YMMref{52.4}{435}, the minimal-orbit argument in Proposition~\YMMref{52.3}{433}, and Lemma~\ref{lem:setup-lattice-covariance}), with the chart integral \eqref{eq:sf-chart-integral} and the Gaussian normaliser of Section~\ref{sec:sf-chart-interface}; and the radius and the sets of the sources (\eqref{eq:source-entry-domain}, \eqref{eq:source-entry-domain}, the half-radius and rate conditions following \eqref{eq:source-entry-domain}), with the size \eqref{eq:source-entry-norm} of the families in the fourth place.

The square cone of Corollary~\ref{stc:small}\,(ii) uses the five-point inequality of periodic Gibbs states (Lemma~\ref{lem:spectral-steps}), which rests on the lattice energy--momentum inequality (Proposition~\textup{\YMMref{99.3}{739}} and Lemma~\ref{lem:spectral-steps}), and the lemma of slope one at small steps (Lemma~\ref{lem:spectral-steps}).

From Section~\ref{sf-sec} on the six fields ${\rm tr}\,F_{ij}F_{kl}$, their continuum limits and a Wightman tensor field, in which the six fields form one tensor field and whose statements rest on Theorems~\YMMref{D}{8}, \YMMref{AJ}{673} and~\YMMref{AD}{607}, we use the following three statements, restated in the notation of this section; the lattice field of that section, $O^T_K(H)=a_K^{-1}\bigl[\sum_pH_{\rho(p)\rho(p)}(x(p))\,a_p-\sum_{\mathfrak b}H_{\lambda_{\mathfrak b}\lambda'_{\mathfrak b}}(x_{\mathfrak b})\,q_{\mathfrak b}\bigr]$, is $\mathcal T_K(H)$ of Definition~\ref{stc:fields}\,(1), its bond form being $q_{\mathfrak b}=\langle\mathbf u_{\mathfrak b,\lambda_{\mathfrak b}},\mathbf u_{\mathfrak b,\lambda'_{\mathfrak b}}\rangle=\frac14\sum_{\epsilon,\epsilon'=\pm}\epsilon\epsilon'\langle\mathbf w(W_x(\mu,\epsilon\lambda_{\mathfrak b})),\mathbf w(W_x(\mu,\epsilon'\lambda'_{\mathfrak b}))\rangle$ at the bond $\mathfrak b=\langle x,x+a_Ke_\mu\rangle$.
\begin{itemize}
\item[(i)] (The lemma on exact symmetries and separated insertions of that section, part~(b).) Let $0<g\le\gamma^{\rm ref}_W$, $\ell\in\mathfrak L(g)$, $q\ge2$ and $\delta>0$. There are a depth $s_q(\delta)$, a function of $g$, $q$ and $\delta$, and a number $\mathsf K^{\rm ub}_q(g,\ell,\delta)$, neither of which depends on the cutoff or on $\beta$, such that for all $(K,\beta)\in\mathfrak d(g)$ with $K\ge s_q(\delta)+2$ and $3a_K\le\delta$ and all bounded Borel $H_1,\dots,H_q:T_\ell\to V$ vanishing off closed sets that are pairwise at sup distance $\ge2\delta$ on $T_\ell$,
\[
|\langle\mathcal T_K(H_1);\dots;\mathcal T_K(H_q)\rangle_{K,\ell,\beta}|\le\mathsf K^{\rm ub}_q(g,\ell,\delta)\prod_{a=1}^q60|H_a|_\infty .
\]
\item[(ii)] (The theorem on periodic Gibbs states and limits on $\mathbb R^3$ of that section, part~(a).) Let $0<g\le\gamma^{\rm ref}_W$, $n\ge2$ and $\delta_0,\varrho_0>0$. There are a number $C^{\rm lc}_n(g,\delta_0,\varrho_0)$ and a threshold, functions of $g$, $n$, $\delta_0$ and $\varrho_0$, such that for all $(K,\beta)\in\mathfrak d(g)$ beyond the threshold, all $\nu,\nu'\in\mathcal G^{\rm per}_{K,\beta}$ and all real bounded Borel $H_1,\dots,H_n:\mathbb R^3\to V$ vanishing off closed sets of sup diameter $\le\varrho_0$ that are pairwise at sup distance $\ge2\delta_0$,
\[
\Bigl|\nu\Bigl(\prod_{a=1}^n\bigl(\mathcal T_K(H_a)-\nu'(\mathcal T_K(H_a))\bigr)\Bigr)\Bigr|\le C^{\rm lc}_n(g,\delta_0,\varrho_0)\prod_{a=1}^n|H_a|_\infty ;
\]
the number depends neither on the cutoff, the states and the positions of the supports nor on their distances beyond $2\delta_0$.
\item[(iii)] (The theorem on the tori of that section, part~(a).) Let $0<g\le\gamma^{\rm ref}_W$, $\ell\in\mathfrak L(g)$, $b\in I(g)$ and $q\ge2$. For Lipschitz $H_1,\dots,H_q:T_\ell\to V$ with pairwise disjoint supports the joint cumulants $\langle\mathcal T_K(H_1);\dots;\mathcal T_K(H_q)\rangle_{K,\ell,L^Kb}$ converge as $K\to\infty$; the limit is the torus limit functional $\mathcal S^{\ell,b}_q(H_1,\dots,H_q)$.
\end{itemize}
Besides, Corollary~\ref{stc:small} uses two further statements of that section, for $0<g\le\gamma^{\rm ref}_W$: the limits at $g$ with their reflection positivity at the nine families of mirrors, (iv), and the Hilbert space of such a limit with its forward cone, (v).
\begin{itemize}
\item[(iv)] (The theorem on periodic Gibbs states and limits on $\mathbb R^3$ of that section, part~(c).) Let $0<g\le\gamma^{\rm ref}_W$. A \emph{Gibbs sequence} at $g$ is $(K_j,\beta_j,\nu_j)$ with $(K_j,\beta_j)\in\mathfrak d(g)$, $K_j\to\infty$ and $\nu_j\in\mathcal G^{\rm per}_{K_j,\beta_j}$; a \emph{limit} at $g$ is the functional
\[
\sigma(H_1,\dots,H_n)=\lim_j\nu_j\Bigl(\prod_{a=1}^nA^{\nu_j}[\mathcal T_{K_j}(H_a)]\Bigr),
\]
for all $n\ge0$ and all continuous $H_1,\dots,H_n:\mathbb R^3\to V$ of compact support with pairwise disjoint supports, along a Gibbs sequence whose labels $a_{K_j}\beta_j$ converge, these limits existing. Every Gibbs sequence at $g$ has a subsequence along which $a_{K_j}\beta_j$ converges to a number in $[g^{-2},L(g^{-2}+\bar B)]$ and all these limits exist. A limit at $g$ is real, symmetric and multilinear, $\sigma(\emptyset)=1$, and $\sigma$ vanishes at a single entry; it obeys the bounds of (ii); it is invariant under every $\mathsf rx=\mathsf Rx+v$, $\mathsf R$ in the cubic group, $v\in\mathbb R^3$, acting on points and indices; and $\sigma(\theta_{\mathsf n,c}\bar F\cdot F)\ge0$ for the reflection $\theta_{\mathsf n,c}$ in every plane $\{\mathsf n\cdot x=c\}$, $\mathsf n\in\{e_\mu,e_\mu\pm e_\lambda\}$, $c\in\mathbb R$, acting on points and indices, and every finite complex combination $F$ of such tuples with supports in $\{\mathsf n\cdot x>c\}$, where $\sigma$ is extended linearly, $\bar F$ has the conjugate coefficients and $F'\cdot F$ is formed by concatenating the tuples. Its trace part $\sigma(f_1\mathbf 1,\dots,f_n\mathbf 1)$ is the scalar limit of the action density of the main theorems along the same sequence.
\item[(v)] (The theorem on invariance, the Wightman field and the decomposition of that section, part~(a).) Let $\sigma$ be a limit at $0<g\le\gamma^{\rm ref}_W$ as in (iv), $\theta$ the reflection in $\{x_3=0\}$ and $\tau_v$ the translation of every entry of a tuple by $v$. The form $\langle[F'],[F]\rangle=\sigma(\theta\bar F'\cdot F)$ on the finite complex combinations of tuples with supports in $\{x_3>0\}$ gives, by quotient and completion, a separable Hilbert space $\mathcal H(\sigma)$ with the unit vector $\Omega=[\emptyset]$; $e^{-tH}[F]=[\tau_{te_3}F]$, $t\ge0$, and $\mathcal U(v)[F]=[\tau_{(v,0)}F]$, $v\in\mathbb R^2$, define a self-adjoint $H\ge0$ with $H\Omega=0$ and a unitary group $\mathcal U(v)=e^{i(v_1P_1+v_2P_2)}$ that commutes with it; and the joint spectral measure $\mathcal P$ of $(H,P_1,P_2)$ is carried by the forward cone, $\mathcal P(\{(E,q):E<|q|\})=0$, which contains $H\ge|P_1|$ and $H\ge|P_2|$. The same holds with every mirror of the nine families in the place of $\{x_3=0\}$.
\end{itemize}
The covariance of tensor limits under $O(3)$ in the Remarks uses two further inputs, which rest, like the scalar theory, on Theorems~\YMMref{D}{8}, \YMMref{AJ}{673} and~\YMMref{AD}{607}. A scalar limit is invariant under all isometries of $\mathbb R^3$ (Theorem~\textup{\YMMref{AP}{780}}). And for $0<g\le\gamma^{\rm sd}_W$ Section~\ref{C-371-1} on reflection positivity at every plane, the forward light cone and Wightman fields for the tensor ${\rm tr}\,F_{ij}F_{kl}$ gives a bound on off-diagonal entries near a mirror that depends neither on the scale nor on the cutoff (Proposition~\ref{C-371-24}\,(e)): there is $C_\infty(g)$, a function of $g$, such that every tensor limit $\sigma$ at $g$ has $\sigma(\theta h,h)\le C_\infty(g)\,b_*^{-2}|h|_\infty^2$, $\theta$ the reflection in $\{x_3=0\}$ and $b_*$ the limit of the labels along a defining sequence of $\sigma$ with convergent labels, for $0<d<1$ and every real $h=h_0(E_{\mu\nu}+E_{\nu\mu})$, $E_{\mu\nu}$ the matrix units, $\mu\ne\nu$, with $h_0$ supported in the ball of radius $d$ about a point $x_0$ with $x_{0,3}\in[\frac32d,2d]$.

\subsection{Estimates and proofs}

\begin{lemma}[Null terms and the remainder]\label{stc:null}
Let $K\in\mathbb N$, $\beta>0$, and let $\eta$ be a vector field of bounded support on $\mathbb R^3$, evaluated at the points of $\mathbb L_K$, or a vector field on a torus.
\begin{enumerate}
\item[(a)] Put $\mathsf e^{\mathsf k}_{\mathsf P}(y)=\sum_{\rho\in\mathsf P}[\mathsf k^\rho(y+\frac12a_Ke_\rho)+\mathsf k^\rho(y-\frac12a_Ke_\rho)]$ and, with $\rho=\rho(p)$, $\mathfrak e_\rho(p)=\sum_{\mathfrak b\subset\partial p}\langle\mathbf y_{\mathfrak b,\rho},\mathbf w(W_{\mathfrak b,p})\rangle+2[\Psi_\rho(p)+\Psi_\rho(p_-)]+[a_{p_+}+a_{p_-}-2a_p]$, and let the insertions $\mathfrak J_{0,v}[\mathsf k]$, linear in a vector field $\mathsf k$, be
\[
\begin{aligned}
\mathfrak J_{0,1}[\mathsf k]
&=\sum_p\Bigl[\tfrac34\mathsf e^{\mathsf k}_{\mathsf P(p)}(x(p))
-\tfrac14b\,a_K\sum_\rho(\nabla^-_\rho\nabla^+_\rho\mathsf k^\rho)(x(p))\Bigr]a_p,\\
\mathfrak J_{0,2}[\mathsf k]
&=\tfrac b4\sum_{p,\rho}\sum_{\mathfrak b\subset\partial p}
\frac{\mathsf k^\rho(x_{\mathfrak b})-\mathsf k^\rho(x(p))}{a_K}\\
&\qquad\langle\mathbf y_{\mathfrak b,\rho},\mathbf w(W_{\mathfrak b,p})\rangle,\\
\mathfrak J_{0,3}[\mathsf k]
&=-\tfrac b2\sum_{p,\rho}(\nabla^+_\rho\mathsf k^\rho)(x(p))\,\Psi_\rho(p),\\
\mathfrak J_{0,4}[\mathsf k]
&=\tfrac\beta4\sum_p\mathsf k^{\rho(p)}(x(p))\,\mathfrak e_{\rho(p)}(p),\\
\mathfrak J_{0,5}[\mathsf k]
&=-\tfrac\beta2\sum_p\Bigl(\sum_{\rho\in\mathsf P(p)}\mathsf k^\rho(x(p))\Bigr)a_p^2 .
\end{aligned}
\]
There is a number $c_{\mathsf k}$ such that, identically in the configuration,
\[
\mathfrak J^{\rm II}[\mathsf k]=-\tfrac14\mathfrak N[\mathsf k]+\sum_{v=1}^5\mathfrak J_{0,v}[\mathsf k]-\tfrac14c_{\mathsf k},\qquad
\mathfrak K[\eta]=\mathcal S[b\eta]-\mathfrak J^{\rm II}[\mathsf k_\eta]=\mathcal N_\eta-\mathsf{St}[\eta]-\tfrac32a_K\mathcal O(w_\eta);
\]
so
\[
\mathsf{St}[\eta]=\mathcal N_\eta-\tfrac14\mathfrak N[\mathsf k_\eta]-\mathsf{Rm}[\eta]-\tfrac14c_{\mathsf k_\eta},\qquad\mathfrak K[\eta]=\mathsf{Rm}[\eta]-\tfrac32a_K\mathcal O(w_\eta)+\tfrac14\mathfrak N[\mathsf k_\eta]+\tfrac14c_{\mathsf k_\eta}.
\]
If $\|\eta\|\le1$, the weights of $\mathfrak J_{0,v}[\mathsf k_\eta]$ are at most $3+5b$ for $v\le3$ and at most $L^Kb$ for $v=4,5$, and those of $\frac32a_K\mathcal O(w_\eta)$ are at most $3$. The field of directions of $\mathcal N_\eta$ is $a_K^{-1}\sum_\rho\eta^\rho(x_{\mathfrak b})\mathbf u_{\mathfrak b,\rho}$ and that of $\mathfrak N[\mathsf k_\eta]$ is $\sum_\rho\mathsf k^\rho_\eta(x_{\mathfrak b})\mathbf y_{\mathfrak b,\rho}$; both vanish off the bonds with ${\rm dist}(x_{\mathfrak b},{\rm supp}\,\eta)\le a_K$. All these functions depend only on the bond variables within $4a_K$ of ${\rm supp}\,\eta$.
\item[(b)] Let $\mathsf N_{\mathbf X}$ be a null term, $\nu$ a Gibbs state at $(a_K,\beta)$ or a torus measure, and $Y$ a bounded Borel function that does not depend on the bonds of $\Lambda$. Then $\nu(\mathsf N_{\mathbf X}Y)=0$. Hence every joint cumulant of $\mathsf N_{\mathbf X}$ with functions that do not depend on these bonds vanishes; so does every joint cumulant of order $\ge2$ with a constant entry.
\item[(c)] If $Y_1,\dots,Y_r$ do not depend on the bonds within $2a_K$ of ${\rm supp}\,\eta$, the joint cumulants under $\nu$ of $\mathsf{St}[\eta]$ and of $-\mathsf{Rm}[\eta]$ with $Y_1,\dots,Y_r$ coincide, and $\nu\bigl(A^\nu[\mathsf{St}[\eta]]\prod_iA^\nu[Y_i]\bigr)=-\nu\bigl(A^\nu[\mathsf{Rm}[\eta]]\prod_iA^\nu[Y_i]\bigr)$. The same holds, without the sign, for $\mathfrak K[\eta]$ and $\mathsf{Rm}[\eta]-\frac32a_K\mathcal O(w_\eta)$.
\end{enumerate}
\end{lemma}

\begin{proof}[Proof sketch]
(a) The general-field formula \eqref{eq:ward-general-displacement} at the constant weight $\beta$ gives $\mathsf{St}[\eta]=\mathfrak v_\eta(\beta\sum_pa_p)-\mathfrak K[\eta]$; the divergence formula \eqref{eq:ward-general-divergence} gives ${\rm div}\,\mathfrak v_\eta=-\frac32\sum_pa_pw_\eta(p)=-\frac32a_K\mathcal O(w_\eta)$, so $\mathfrak v_\eta(\beta\sum_pa_p)=\mathcal N_\eta-\frac32a_K\mathcal O(w_\eta)$; and the summation by parts in the orbital remainder Proposition~\ref{prop:ward-orbital-cubic}(c) at the vector weight $\beta\eta$ gives $\mathfrak K[\eta]=\mathcal S[b\eta]-\mathfrak J^{\rm II}[\mathsf k_\eta]$, as $a_K^{-1}\beta=L^{2K}b$. The first identity is obtained by computing $\mathfrak N[\mathsf k]$ from its definition with the calculus of the derivations \ref{lem:ward-general-calculus} and ordering the result by densities; it is an identity between polynomials in finitely many bond variables, and inserting it gives the last two.

On $\mathbb L_K$ the finitely many bonds on which these functions depend may be realised on a torus so large that nothing wraps. The weights: $\mathsf k^\rho_\eta(y)=\int_0^1\partial_\rho\eta^\rho(y+sa_Ke_\rho)\,ds$, so $|\mathsf k_\eta|\le\sup|\nabla\eta|$, ${\rm Lip}\,\mathsf k_\eta\le3\sup|\nabla^2\eta|$, $|w_\eta|\le2\sup|\nabla\eta|$; $|\mathsf e^{\mathsf k}_{\mathsf P}|\le4\sup|\mathsf k|$ and $a_K|\nabla^-_\rho\nabla^+_\rho\mathsf k^\rho|\le2\,{\rm Lip}\,\mathsf k$ give $3+\frac92b$ for $v=1$, the half lattice spacing between $x_{\mathfrak b}$ and $x(p)$ gives $\frac b8{\rm Lip}\,\mathsf k$ for $v=2$, and $\frac b2{\rm Lip}\,\mathsf k$, $\frac\beta4\sup|\mathsf k|$, $\beta\sup|\mathsf k|$ for $v=3,4,5$, with $\beta=L^Kb$.

The field of directions of $\mathfrak v_\eta$ is given in \YMMref{(2103)}{781}; $\mathsf k^\rho_\eta(x_{\mathfrak b})\ne0$ only if $x_{\mathfrak b}$ or $x_{\mathfrak b}+a_Ke_\rho$ lies in ${\rm supp}\,\eta$; a weight at an index that belongs to $p$ depends on the values of $\mathsf k$ within $a_K$ of $x(p)$, and the sites of its box lie within $\frac32a_K$ of $x(p)$ (Lemma~\ref{stc:rsharp}).
(b) Haar measure is invariant under left translations, so $\int\partial^{\mathsf a}_{\mathfrak b}\Phi\,dU(\mathfrak b)=0$ for every $\Phi$ smooth in $U(\mathfrak b)$. Take $\Phi=\mathbf X^{\mathsf a}_{\mathfrak b}e^{H_\Lambda}$, $H_\Lambda=-\beta\sum a_p$ over the plaquettes that contain a bond of $\Lambda$, and sum over $\mathfrak b\in\Lambda$ and $\mathsf a$:
\[
\int\Bigl({\rm div}\,\mathfrak v_{\mathbf X}-\beta\,\mathfrak v_{\mathbf X}\sum_pa_p\Bigr)e^{H_\Lambda}\prod_{\mathfrak b\in\Lambda}dU(\mathfrak b)=0
\]
at every value of the other bonds, and $Y$ is constant under this integral. So the Gibbs kernel of $\Lambda$ annihilates $\mathsf N_{\mathbf X}Y$, and the equations of Dobrushin, Lanford and Ruelle (\YMMref{(1819)}{698}) give $\nu(\mathsf N_{\mathbf X}Y)=0$; on a torus the same holds with the Wilson weight. A joint cumulant is a sum over partitions of products of means, and the block that contains $\mathsf N_{\mathbf X}$ contributes $0$; of the other entries nothing is used but that they do not depend on the bonds of $\Lambda$, so tensor entries are admitted like scalar ones. (c) follows from (a) and (b), constants dropping out of $A^\nu$ and of cumulants of order $\ge2$.
\end{proof}

\begin{lemma}[The densities; the cubic density]\label{stc:rsharp}
\begin{enumerate}
\item[(a)] A polynomial in the entries of finitely many $SU(2)$ matrices and their inverses has exactly one polynomial continuation to $SL(2,\mathbb C)$, and gauge invariance passes to it. The functions $a_p$, $\langle\mathbf y_{\mathfrak b,\rho},\mathbf w(W_{\mathfrak b,p})\rangle$ ($\mathfrak b\subset\partial p$, $\rho\ne\mu_{\mathfrak b}$) and $\Psi_\rho(p)$ are densities of order $2$; $\mathfrak e_{\rho(p)}(p)$ is one of order $3$; $a_p^2$ is one of order $4$; each on the smallest box that holds the bonds on which it depends, all of whose sites lie within $\frac32a_K$ of $x(p)$, with constants and heights that are pure numbers. All are real on $SU(2)$ configurations, and every insertion built from them has a counting number that is a pure number.
\item[(b)] $\mathfrak r^\sharp_\rho(p)=0$ for $\rho$ in the plane of $p$. For $\rho=\rho(p)$ it is a gauge invariant polynomial in the twenty bond variables of the two cubes with common face $p$, on the box $Q(p)$ of their twelve sites. Proposition~\ref{prop:ward-orbital-cubic} proves the real bound $|\mathfrak r^\sharp_\rho(p)|\le c_{\rm orb}\mu^3$ when their eleven plaquette variables are within $\mu$ of $\mathbf1$, with a fixed $c_{\rm orb}$, and the complex bound of order $3$ with a pure constant $c_\sharp$ and height $\hat c_\sharp$. A site is the lowest site of three boxes $Q(p)$, so $\mathcal S_\varkappa$ is an insertion of type $(3,c_\sharp,\hat c_\sharp,3)$ with $|\varkappa|_\infty=\max|\varkappa|$.
\item[(c)] For a diagonal isometry $\mathsf r:x\mapsto{\rm diag}(s_1,s_2,s_3)x+v$ of the lattice, $\mathfrak r^\sharp_\rho(\mathsf rp)(U)=s_\rho\,\mathfrak r^\sharp_\rho(p)(\mathsf r^{-1}U)$; so $\mathcal S_\varkappa\circ\mathsf r^{-1}=\mathcal S_{\varkappa'}$, $\varkappa'(q)=s_{\rho(q)}\varkappa(\mathsf r^{-1}q)$, and $\mathcal S[\Xi]\circ\mathsf r^{-1}=\mathcal S[\Xi^{\mathsf r}]$ with $\Xi^{\mathsf r}_\rho(y)=s_\rho\Xi_\rho(\mathsf r^{-1}y)$, $\mathsf H(\Xi^{\mathsf r})=\mathsf H(\Xi)$. In particular $\mathfrak r^\sharp$ is odd under inversions, and the torus mean of $\mathcal S[\Xi]$ vanishes.
\end{enumerate}
\end{lemma}

\begin{proof}[Proof sketch]
(a) A polynomial that vanishes on $SU(2)^n$ vanishes on $SL(2,\mathbb C)^n$ ($SU(2)$ is a real form). The densities \eqref{eq:ward-orbital-normal-potential}, \eqref{eq:ward-orbital-inplane} are polynomials in the vector parts and the traces of plaquette holonomies at a corner ($\langle\mathbf w(V),\mathbf w(V)\rangle=a(2-a)$, $a=1-\frac12{\rm tr}\,V$, and $|a_p|\le\frac12r^2$ as in Lemma~\ref{stc:defect}); they vanish to the second order at the trivial configuration for $a_p$, $\langle\mathbf y,\mathbf w(W)\rangle$, $\Psi_\rho(p)$, to the third order for $\mathfrak e_\rho(p)$, whose terms of second order cancel, and to the fourth for $a_p^2$. A petal holonomy is conjugate, by a product of at most three bond matrices of norm $\le2$, to a plaquette variable or its inverse, so it is within $2^6r$ of $\mathbf 1$ when the plaquette variables are within $r$; the orders follow, and the heights at $|U(\mathfrak b)|\le2$ in the same way. (b) The local polynomial, the real cubic estimate and its complex version follow from Proposition~\ref{prop:ward-orbital-cubic}. Its proof uses a tree on the two cubes and cancellation of the quadratic Taylor polynomial. (c) The covariance in that proposition gives the stated transformation of the weights. The isometry preserves plaquette centres and sup distances, hence also $\mathsf H$. Reflection through the plane of $p$ reverses its normal component and fixes $p$, so invariance of the torus measure gives the zero mean.
\end{proof}

\begin{lemma}[The defect is a null petal form]\label{stc:defect}
Let $K\in\mathbb N$, $\beta>0$, $\eta\in C^2_c(\mathbb R^3;\mathbb R^3)$ and $H:\mathbb R^3\to V$ bounded.
\begin{enumerate}
\item[(a)] $\Delta[\eta]=F(k^0,v^0)+F(k^1,v^1)$, where $(k^0,v^0)$ is null, $|k^0|_1+|v^0|_\infty\le22\,b\sup|\nabla\eta|$, $|k^1|_1+|v^1|_\infty\le60\,b\,a_K\sup|\nabla^2\eta|$, and the weights vanish at the sites farther than $2a_K$ from ${\rm supp}\,\eta$. Moreover $\mathcal T_K(H)=F(k,v)$ with $|k|_1\le3|H|_\infty$, $|v|_\infty\le|H|_\infty$.
\item[(b)] A rough insertion $F(k,v)$ is an insertion of order $2$ on the boxes $x+a_K\{-1,0,1\}^3$, with pure numbers as constant and height, counting number $\le600$, and $|\varkappa|_\infty\le L^K\max\{|k|_1,|v|_\infty\}$.
\end{enumerate}
\end{lemma}

\begin{proof}[Proof sketch]
(a) Each bond $\mathfrak b=\langle x,x+a_Ke_\mu\rangle$ lies in four plaquettes whose loops are petals at $x$, $W_{\mathfrak b,p(\mathfrak b,\alpha')}=W_x(\mu,\alpha')$, and $\mathbf u_{\mathfrak b,\rho}=-\frac12\sum_{\epsilon=\pm}\epsilon\,\mathbf w(W_x(\mu,\epsilon\rho))$; so $\mathsf\Theta_\beta[\eta]$, $Y$ and $\mathcal O$ are petal forms: $\mathsf\Theta_\beta[\eta]=a_K^{-1}\sum_x\sum_\mu\sum_{\alpha'}\sum_{\rho\ne\mu}\sum_{\epsilon=\pm}k^\Theta q_x((\mu,\epsilon\rho),(\mu,\alpha'))$ with $k^\Theta=-\frac\epsilon2ba_K^{-1}[\eta^\rho(x_{\mathfrak b})-\eta^\rho(x_{\mathfrak b}+\frac12a_Ke_{\alpha'})]$, $48$ terms at a site. Expanding the values of $\eta$ about the site $x$ splits the weights into a part linear in $\nabla\eta(x)$, which is $(k^0,v^0)$, and a remainder of size $a_K\sup|\nabla^2\eta|$; counting the terms gives $|k^0|_1\le18b\sup|\nabla\eta|$, $|v^0|_\infty\le4b\sup|\nabla\eta|$ and the bound for $(k^1,v^1)$. With $G=\nabla\eta(x)$, $T_{\lambda\lambda'}=\langle B_\lambda,B_{\lambda'}\rangle$, $H=({\rm tr}\,G)\mathbf 1-G-G^{\rm T}$ and $\sum_\mu\langle F_{\mu\rho},F_{\mu\nu}\rangle=\delta_{\rho\nu}{\rm tr}\,T-T_{\rho\nu}$, which follows from $\sum_\mu\varepsilon_{\mu\rho\lambda}\varepsilon_{\mu\nu\lambda'}=\delta_{\rho\nu}\delta_{\lambda\lambda'}-\delta_{\rho\lambda'}\delta_{\lambda\nu}$, the symbol of $(k^0,v^0)$ at $x$ is
\[
b\,[({\rm tr}\,G)\,{\rm tr}\,T-{\rm Tr}(GT)]-\tfrac b2({\rm tr}\,G)\,{\rm tr}\,T-\tfrac b2{\rm Tr}(HT)=0,
\]
the three members coming from $\mathsf\Theta_\beta$, from $-b\,\mathcal O(\nabla\!\cdot\!\eta)$ and from $-b\,\mathcal T_K(H^{\rm cons}[\eta])$, because $\frac12{\rm Tr}(HT)=\frac12({\rm tr}\,G)\,{\rm tr}\,T-{\rm Tr}(GT)$ for symmetric $T$. This is the one place where it is used that the six fields are petal forms of the gauge field itself. (b) All the densities at the site $x$ are densities on the one box $x+a_K\{-1,0,1\}^3$: $576$ forms $q_x(\pi,\pi')$ and three $a_{p_m(x)}$; they are gauge invariant since both loops are based at $x$. A petal holonomy is conjugate, by a product of at most three bond matrices of norm $\le2$, to a plaquette variable or its inverse; so $|\mathbf w(W_x(\pi))|\le2^6r$, $|q_x(\pi,\pi')|\le2^{12}r^2$ and $|a_p|\le\frac12r^2$ (as ${\rm tr}(V-\mathbf 1)=-\det(V-\mathbf 1)$ on $SL(2,\mathbb C)$) when $|U(\mathfrak b)|\le2$ and $|\partial U(q)-\mathbf 1|\le r\le1$; the heights follow in the same way.
\end{proof}

\begin{lemma}[The frame]\label{stc:frame}
Let $0<g\le\gamma$.
\begin{enumerate}
\item[(a)] Every $(K,\beta)\in\mathfrak d(g)$ has $g^{-2}+B_\sharp-\bar B/L<b\le\mathsf b_+$; so $b\in[\gamma^{-2},\mathsf b_+]$ and $b\ge4$. For every side $2L^m\in\mathfrak L(g)$ and every $1<\mathsf b'\le\mathsf b_+$, $MR(\mathsf b')\le L^m$, with the non-decreasing cube-size function $R$ of \YMMref{74}{614}.
\item[(b)] Let $0<\phi<1$, and let $K_\phi$ be the least integer $K\ge2$ with $(1-\phi)K\ge2$, $L^{\phi K}\gamma^{-2}\ge\mathsf t_\natural$ and $\phi K>16\mathfrak p_0^2c_{\log}/(\log L)^2$, where $\mathsf t_\natural$, $c_{\log}$ are numbers of the setting (\ref{sec:cutoff-calibration}, Section~\ref{sec:cutoff-calibration}). For $K\ge K_\phi$, every $\mathsf b'\ge\gamma^{-2}$ and every $m$ with $MR(\mathsf b')\le L^m$, the level $t=K-\lceil\phi K\rceil$ is an admissible top level: (Th) holds on $T_{2L^m}$ with $(K,\mathsf b',m,t)$, and the numbers $\theta_k$ of the chain satisfy
\[
L\theta_k\le\Theta_KL^{-(K-k)/2}\quad(0\le k\le t),\qquad\Theta_K=\max\{1,\tfrac{13}{72}(L-1)(1+Kc_{\log})^{\mathfrak p_0}\}\le CK^{\mathfrak p_0},
\]
so that $\sum_{k<n}L\theta_k\le\Theta_KL^{-(K-n)/2}$ for $1\le n\le t$.
\item[(c)] If $0<\phi\le\frac1{84}$ and $K\ge\max\{K_\phi,5\}$, then $u=\lceil\phi K\rceil\le\lceil\frac K{84}\rceil\le\lceil\frac K{14}\rceil$, $u\le\frac K{84}+1$, $4u\le\frac K{21}+4\le K$ and $t\ge\frac34K$.
\end{enumerate}
\end{lemma}

\begin{proof}[Proof sketch]
(a) By \YMMref{(1535)}{616} and Corollary~\YMMref{74.2}{616}, applied at the common constants of Section~\YMMref{82}{665} and at $g\le\gamma$, the label lies in $(g^{-2}+B_\sharp-L^{-1}\bar B,L(g^{-2}+\bar B)]$; by \YMMref{74}{614}, with $\bar B=B_\sharp L/(L-1)$, $B_\sharp-L^{-1}\bar B=B_\sharp\frac{L-2}{L-1}\ge0$, and $\gamma\le\frac12$. The bound $MR(\mathsf b')\le L^m$ is \YMMref{(1528)}{615} with $R$ non-decreasing. (b) Of a top level (Th) asks only $1\le t\le K-1$ and $L^{K-t}\mathsf b'\ge\mathsf t_\natural$ (\ref{sec:cutoff-calibration}), besides $K\ge2$, $\mathsf b'\ge\gamma^{-2}$ and $MR(\mathsf b')\le L^m$; the first two conditions on $K_\phi$ give these. By \YMMref{(2176)}{803}, $\mathsf s^{\rm abs}_{K-y}\le\mathsf s_\gamma L^{-y/2}(1+yc_{\log})^{\mathfrak p_0}$; with $\log(1+a)\le2\sqrt a$ and the third condition the logarithm of the factor is negative for $y\ge\phi K$, so the runs of the chain are admissible at all levels $k\le t$, and $L\theta_k\le\frac{13}{72}(L-1)L^{-(K-k)/2}(1+(K-k)c_{\log})^{\mathfrak p_0}$; summing the geometric series ($L\ge13$) gives the last bound. (c) $\phi K\le\frac K{84}$, and the ceiling is non-decreasing: this gives the first two bounds. $\lceil x\rceil<x+1$ gives $u\le\frac K{84}+1$, so $4u\le\frac K{21}+4$ and $t=K-u\ge\frac{83}{84}K-1$; for $K\ge5$, $\frac K{21}+4\le K$ and $\frac{83}{84}K-1\ge\frac34K$, both being $K\ge\frac{21}5$.
\end{proof}

\begin{proposition}[The chain with one insertion, on the whole cut set]\label{stc:chain}
Let the constants satisfy (Th1), (Th2), and let (Th) hold on $T_\ell$, $\ell=2L^m$, with $(K,\mathsf b',m,t_*)$; the line is $x_j=L^{K-j}\mathsf b'$, $x_0=\beta$. Let $\mathfrak J$ be an insertion of type $(m_{\mathfrak J},c,\hat c,N)$, $m_{\mathfrak J}\ge2$, $\varsigma=2m_{\mathfrak J}-3$, and put
\[
\begin{multlined}
\mathsf m_j=8NM^3c\,\alpha_0^{m_{\mathfrak J}}|\varkappa|_\infty L^{-\varsigma j},\quad\mathsf n^{\rm f}=16NM^3\hat c\,\tfrac{\mathsf C_{\mathcal A}}{r^\circ_P}\,\delta(x_0)|\varkappa|_\infty,\quad\mathsf d_1=K^{\rm d}\mathsf n^{\rm f},\quad\mathsf n_j=\sum_{i\le j}L^{i-j}\mathsf d_i, \\
\mathsf d_{j+1}=\theta_j(\mathsf n_j+\mathsf m_j),
\end{multlined}
\]
\[
\mathsf o_k=N_k\bigl[Nc\,(\mathsf e^{\rm c}_k)^{m_{\mathfrak J}}|\varkappa|_\infty L^{-\varsigma k}+C_{\rm old}\mathsf n_k\bigr],\qquad r_{\mathfrak J}=\min\Bigl\{\frac{w_\star}{32\,\mathsf n^{\rm f}},\frac{E_0}{16(\mathsf d_1+\mathsf m_1)}\Bigr\},
\]
with the numbers $\alpha_0$, $\mathsf C_{\mathcal A}$, $r^\circ_P$, $\delta(\cdot)$, $K^{\rm d}$, $\theta_j$, $\mathsf e^{\rm c}_k$, $N_k=(\ell L^{K-k})^3$, $C_{\rm old}$, $w_\star$, $E_0$ of the chain (\YMMref{114}{799}, \YMMref{(2163)}{798}). For complex $|s|,|s'|<r_{\mathfrak J}$ there are pointed families $E^{(i)}[s]$ of level $i$, $1\le i\le t_*$, with $E^{(i)}[0]=E^{(i)}$ the families of the chain, and numbers $\mathsf c_i(s)$, such that:
\begin{enumerate}
\item[(i)] $\|E^{(i)}[s]\|<\frac{15}{16}E_0+\mathbf 1_{\{i=1\}}|s|\mathsf d_1$ and $\|E^{(i)}[s]-E^{(i)}[s']\|\le|s-s'|\mathsf d_i$; the members are holomorphic in $s$;
\item[(ii)] (locality) if $\mathfrak J_0,\dots,\mathfrak J_r$ have pairwise disjoint carriers and $\mathfrak J_\sigma=\sum_a\sigma_a\mathfrak J_a$, $\sigma\in\{0,1\}^{r+1}$, then for $|s|$ below a common radius the member at a polymer $X$ of the families of $\mathfrak J_\sigma$ depends only on the $\sigma_a$ for which $\mathfrak J_a$ meets $X$;
\item[(iii)] for $1\le k\le t_*$,
\[
(\mathcal M^k)_*(\mathsf p_k\,e^{s\mathfrak J}\rho_0\,dU)=\exp\Bigl[\sum_{i\le k}\mathsf c_i(s)\Bigr]\mathbf 1_{\mathsf{Cut}_k}e^{\Phi^s_k}\,dV,
\]
a bounded complex Borel density, where $\Phi^s_k(V)=\Phi_k(V)+\sum_{i\le k}\sum_{X,z}[\delta E^{(i)}_s(X,U,z)-\delta E^{(i)}_s(X,\mathbf 1,z)]+s\mathfrak J(U)$, $\delta E^{(i)}_s=E^{(i)}[s]-E^{(i)}$, $\Phi_k$ is the exponent of the chain (the definition of $\Phi_k^w$ in Proposition~\YMMref{114.2}{801} at zero source) and $U$ is any element of the minimal orbit over $V$, that is, any configuration over $V$ (Definition~\ref{stc:ins}\,(3)), on which $\Phi^s_k$ does not depend: $\Phi_k$ is a function of $V$ (Proposition~\YMMref{114.2}{801}); the double sum is the difference of the actions at $U$ of the two bare arrays formed by the families $E^{(i)}[s]$, resp.\ $E^{(i)}$, $i\le k$, and the action of a bare array is defined at every element of the minimal orbit over $V$ and takes one value there (Propositions~\YMMref{114.2}{801} and~\YMMref{113.13}{799}); and $\mathfrak J$ is gauge invariant, while the minimal orbit is one orbit of gauge transformations of the bare lattice. The numbers are $\mathsf c_i(s)=\mathsf c_i(0)+\sum_{X,z}\delta E^{(i)}_s(X,\mathbf 1,z)$, with $\mathsf c_i(0)$ the constants of the chain (from the constant \YMMref{(2173)}{801} of a step). At the level $k=t$ the display gives $\Lambda^{\rm sf}(e^{s\mathfrak J})=\exp[\sum_{i\le t}\mathsf c_i(s)]\int\mathbf 1_{\mathsf{In}_t}e^{\Phi^s_t}\,dV$;
\item[(iv)] $|\Phi^s_k-\Phi^{s'}_k|\le|s-s'|\,\mathsf o_k$ on $\mathsf{Cut}_k$;
\item[(v)] (several insertions; a source) for finitely many insertions $\mathfrak J_a$ of orders $\ge2$, with their numbers $\mathsf d_{i,a}$, $\mathsf m_{j,a}$, $\mathsf o_{k,a}$, $r_a=r_{\mathfrak J_a}$, and strengths $\mathbf s=(s_a)$ in the parameter domain $\{\sum_a|s_a|/r_a<1\}$, the same construction with $\sum_as_a\mathfrak J_a$ gives families $E^{(i)}[\mathbf s]$ with $\|E^{(i)}[\mathbf s]\|<\frac{15}{16}E_0+\mathbf 1_{\{i=1\}}\sum_a|s_a|\mathsf d_{1,a}$, $\|E^{(i)}[\mathbf s]-E^{(i)}[\mathbf s']\|\le\sum_a|s_a-s'_a|\mathsf d_{i,a}$, $|\Phi^{\mathbf s}_k-\Phi^{\mathbf s'}_k|\le\sum_a|s_a-s'_a|\mathsf o_{k,a}$ on $\mathsf{Cut}_k$, members at $X$ that depend only on the $s_a$ with $\mathfrak J_a$ meeting $X$ and are jointly holomorphic, and (iii) for $\sum_as_a\mathfrak J_a$; the same holds with a complex source $w$ of the action density in the class, $\rho_0e^{\mathcal O(w)}$ taking the place of $\rho_0$.
\end{enumerate}
No largeness of $K$ enters: (i)--(v) hold at every admissible top level.
\end{proposition}

\begin{proof}
The one-step estimates used here are Lemma~\YMMref{77.1}{631} and
Proposition~\YMMref{77.2}{632}, together with the density identity in Theorem~\YMMref{W}{423}.
They are applied on the full analytic background cores, with reference
coupling $\mathsf t=x_j$ and $\tau=1$; the source and insertion parameters
do not change this Gaussian reference. In their fourth-family norm the
fluctuation variable is $B$, the physical field is $g_jB$, and the
quadratic weight is $1+M^{-3}\sum_b|B(b)|^2$. The sum runs over the
fluctuation bonds in the polymer. The source class and the array ceiling
are those of Proposition~\YMMref{114.2}{801}, with $E_0^{\rm s}=E_0$ and
$K^{\rm d}=(351/2)L^3K_\Pi$. Its hypotheses include the fixed smallness
budgets and the whole-chart conditions needed for the pushed-forward
density. Its proof gives the output bound $15E_0/16$ from every
admissible array; Proposition~\YMMref{114.3}{803} gives the weighted difference
estimate used below.

The families are defined by recursion on the height, as for the chain, with two additions. At height $0$ the first step receives, besides its families, the family $s\mathbb W_{\mathfrak J}$ in the fourth place, $\mathbb W_{\mathfrak J}(Y;\cdot)=\sum_{\iota:X_0(Q_\iota)=Y}\varkappa_\iota[\mathsf j_\iota(e^{iB'}U)-\mathsf j_\iota(U)]$ ($X_0(Q)$ the union of the cubes of level $0$ that hold a site of $Q$), of norm $\le|s|\mathsf n^{\rm f}<w_\star/32$; the source-inclusive unperturbed fourth family has norm at most $7w_\star/32$, so the one-step hypotheses of Lemma~\YMMref{77.1}{631} hold, and its output $E^{(1)}[s]$ is a pointed family. At height $j\ge1$ the early member of level $j$ is replaced by $E^{(j)}[s]+sH^{(j)}_{\mathfrak J}$, where $H^{(j)}_{\mathfrak J}(X,\cdot,z)$ distributes over the sites of $X$ the boxes $Q_\iota$ with $X_j(Q_\iota)=X$; it has norm $\le\mathsf m_j$, it vanishes at the trivial configuration, and it sums to $\mathfrak J$ at every background over a field cut, that is, at every element $U$ of the minimal orbit $\mathcal C(\mathbf B_j,V)$ over a block field $V\in\mathsf{Cut}_j$. (Here and below the minimal orbit $\mathcal C(\mathbf B_k,V)$ over a block field $V$ of level $k$ and its elements, the \emph{backgrounds} over $V$, are those of Section~\ref{sec:analytic-core-complements}; a gauge invariant function, such as $\mathfrak J$ or the sum $\mathsf{Sum}[\mathsf f]$ of the members of a pointed family, takes one value on the orbit; a configuration \emph{over} $\mathsf{Cut}_k$ is a background over some $V\in\mathsf{Cut}_k$.) For every run whose early members have norms $\le E_0$, Proposition~\YMMref{114.2}{801} gives an output family of norm $<15E_0/16$. The three contributions to this bound are at most $E_0/2$, $E_0/4$ and $3E_0/16$; the next run admits every early member of norm at most $E_0$.

With $|s|(\mathsf d_1+\mathsf m_1)<\frac1{16}E_0$ the perturbed early members, of norms $<\frac{15}{16}E_0+|s|(\mathsf d_1+\mathsf m_1)$ (as $\mathsf m_j\le\mathsf m_1$), are admissible, and the recursion of Proposition~\YMMref{114.2}{801} gives $E^{(j+1)}[s]$, of norm $<\frac{15}{16}E_0$: this is the first half of (i) at $i\ge2$. At $i=1$ the Lipschitz bound of one step in its family in the fourth place (Proposition~\YMMref{77.2}{632}, whose coefficient is at most $54(3L^3-3)K_\Pi<K^{\rm d}$ in these size budgets) gives $\|E^{(1)}[s]-E^{(1)}[s']\|\le K^{\rm d}|s-s'|\mathsf n^{\rm f}=|s-s'|\mathsf d_1$, and with $s'=0$ the first half of (i) at $i=1$.

For $j\ge1$ the second half of (i) at $i=j+1$ follows from the Lipschitz bound of a run in its early members, the loss of memory of Proposition~\YMMref{114.3}{803}: for two runs from the level $j$ whose early members of the levels $i\le j$ differ in norm by at most $\Delta_i$, the families they produce differ by $\Delta_{j+1}\le\theta_j\sum_{i\le j}L^{-(j-i)}\Delta_i$, $\theta_j=K^{\rm d}K_\Lambda\mathsf s^{\rm abs}_j$. The runs at $s$ and at $s'$ have $\Delta_i\le|s-s'|\mathsf d_i$ for $i<j$ and $\Delta_j\le|s-s'|(\mathsf d_j+\mathsf m_j)$, so $\Delta_{j+1}\le\theta_j|s-s'|(\mathsf n_j+\mathsf m_j)=|s-s'|\mathsf d_{j+1}$. Holomorphy is that of one step in its data. (ii) follows from the locality clauses of Lemma~\YMMref{77.1}{631} and Proposition~\YMMref{114.2}{801} (two families in the fourth place, or two bare arrays, that agree near $X$ give the same output at $X$), by induction on the level.

For (iii), the chart identity is Theorem~\YMMref{W}{423} on the same whole-chart domain; $\mu_k=(\mathcal M^k)_*(\mathsf p_ke^{s\mathfrak J}\rho_0dU)$ satisfies $\mu_{k+1}=\mathbf 1_{\mathsf{Cut}_{k+1}}\mathbf T(\mathfrak G_k\chi^{\rm sf}\mu_k)$ by the rules of images of complex measures ($\mathbf T$ the image under one block averaging, $\mathfrak G_k$ the gauge factor, $\chi^{\rm sf}$ the small-field indicator); the chart integral of the first step is $\exp[\mathsf c_1(s)+\Phi^s_1]$, and by induction $e^{\Phi^s_k}$ is on $\mathsf{Cut}_k$ the integrand of the run started at level $k$. The numbers of (iii) are the constants of these steps (Proposition~\YMMref{114.2}{801}): in the normalizing constant of a step only the values at the trivial configuration of the family that the step produces depend on its data, so that
\[
\mathsf c_i(s)=\mathsf c_i(0)+\sum_{X,z}\delta E^{(i)}_s(X,\mathbf 1,z),\qquad\delta E^{(i)}_s=E^{(i)}[s]-E^{(i)},
\]
with $\mathsf c_i(0)$ the constants of the chain. The sums are continuous on the open set $\mathsf{Cut}_k$, hence Borel, and gauge invariant. (iv): at a background $U$ over $V\in\mathsf{Cut}_k$,
\[
\Phi^s_k(V)-\Phi^{s'}_k(V)=(s-s')\,\mathfrak J(U)+S[\mathfrak g](U),\qquad S[\mathfrak g](U)=\sum_{i\le k}\sum_{X,z}\bigl[\mathsf g_i(X,U,z)-\mathsf g_i(X,\mathbf 1,z)\bigr],
\]
the action of the bare array $\mathfrak g$ with the families $\mathsf g_i=E^{(i)}[s]-E^{(i)}[s']$, of norms $\le|s-s'|\mathsf d_i$ by (i), with vanishing counterterms and empty slot. The bound of the action of a bare array at a background field, Proposition~\YMMref{113.13}{799}, states: for $1\le k\le t_*$, $V\in\mathsf{Cut}_k$, $U\in\mathcal C(\mathbf B_k,V)$ and pointed families $\mathsf g_i$ of the levels $i\le k$,
\[
|S[\mathfrak g](U)|\le C_{\rm old}\,N_k\sum_{i\le k}L^{-(k-i)}\|\mathsf g_i\|,
\]
with the number $C_{\rm old}$ of equation~\YMMref{(2163)}{798}, a function of the constants, at every admissible top level; no largeness of $K$ enters. Here it gives $|S[\mathfrak g](U)|\le C_{\rm old}N_k\mathsf n_k|s-s'|$. The idea of that bound: $U$ is regular of level $k$ at the radius $\alpha_0$ (\YMMref{113.10}{796}), hence regular of level $i$ at the radius $\alpha_0L^{-2(k-i)}$, and a member at a polymer $X$ of level $i$ minus its value at $\mathbf 1$ is of second order in this radius as long as $X$ is not too large (the local estimate in the proof of Lemma~\YMMref{113.12}{798}); so the polymers of moderate size through the $N_i=L^{3(k-i)}N_k$ sites of level $i$ contribute $N_iL^{-4(k-i)}=N_kL^{-(k-i)}$ times the norm, and the larger polymers, and all polymers at the few levels $i$ next to $k$, are paid for by the decay $e^{-(\kappa+1)d_i(X)}$ of the members and the sums of \eqref{eq:analytic-domain-tail}, the numbers $L^{4l}e^{-c_{\rm A}L^l}$, $l\ge0$, being bounded ($c_{\rm A}$ the number of equation~\YMMref{(2163)}{798}).

Besides, $|\mathfrak J(U)|\le N_kNc(\mathsf e^{\rm c}_k)^{m_{\mathfrak J}}|\varkappa|_\infty L^{-\varsigma k}$: the plaquette variables of $U$ are within $\mathsf m_k(U)L^{-2k}<\mathsf e^{\rm c}_kL^{-2k}\le1$ of $\mathbf 1$ ($\mathsf{Cut}_k=\mathcal C^\flat_k(\mathsf e^{\rm c}_k)$, \YMMref{113.10}{796}, with the regularity number $\mathsf m_k(U)=L^{2k}\max_p|\partial U(p)-\mathbf 1|$ of Definition~\ref{def:sf-chart-data}), so that a density of order $m_{\mathfrak J}$ is at most $c(\mathsf e^{\rm c}_kL^{-2k})^{m_{\mathfrak J}}$ there, and $\mathfrak J$ has at most $N(\ell L^K)^3=NN_kL^{3k}$ indices. The two bounds add up to $|s-s'|\,\mathsf o_k$. (v): the statements of one step are worded for parameters in a point or an open subset of some $\mathbb C^p$; every estimate above is linear in the insertion; so the proof goes through with $\mathbf s\in\mathbb C^{\#a}$, and with the source among the parameters.
\end{proof}

\begin{lemma}[Conditional means; the numbers of one insertion]\label{stc:cmean}
In the setting of Proposition~\ref{stc:chain}:
\begin{enumerate}
\item[(a)] $(\mathcal M^k)_*(\mathsf p_kY\rho_0dU)=(\mathsf Q_{k\leftarrow0}Y)\,\nu_k$ for bounded Borel $Y$, with $\mathsf Q_{k\leftarrow0}=\mathsf Q_{k-1}\cdots\mathsf Q_0$ and $\mathsf Q_kh=\mathsf I[hF_k]/\mathsf I[F_k]$ on $\mathsf{Cut}_{k+1}$, $0$ off it ($\mathsf I$ the chart integral, $F_k$ the integrand of the step). $\mathsf Q_kh$ is a bounded Borel function, linear in $h$, that depends on $h$ on $\mathsf{Cut}_k$ alone; $\mathsf Q_k1=\mathbf 1_{\mathsf{Cut}_{k+1}}$; $\nu_k\ge0$ and $\|\mathsf Q_kh-a_0\|_{L^\infty(\nu_{k+1})}\le\|h-a_0\|_{L^\infty(\nu_k)}$ for $a_0\in\mathbb C$. For $k\ge1$ and a pointed family $\mathsf g$ of level $i'\le k$ (or a formal sum of such), $\mathsf Q_k\Sigma_k[\mathsf g]=\Sigma_{k+1}[\mathsf g+\mathsf S_k\mathsf g]$ at every point, with a linear map $\mathsf S_k$, $\|\mathsf S_k\mathsf g\|\le\theta_kL^{-(k-i')}\|\mathsf g\|$, whose member at $X$ depends on the members of $\mathsf g$ at the polymers inside $X$ alone.
\item[(b)] There is a number $\dot c_k$ with $\mathsf Q_{k\leftarrow0}\mathfrak J=\dot c_k+\partial_sA_k(0,\cdot)$ $\nu_k$-almost everywhere on $\mathsf{Cut}_k$, $A_k(s,\cdot)=\Phi^s_k-\Phi_k$; $|\partial_sA_k(0,\cdot)|\le\mathsf o_k$, and $|A_k(s,V)/s-\partial_sA_k(0,V)|\le2|s|\mathsf o_k/r_{\mathfrak J}$ for $0<|s|\le\frac12r_{\mathfrak J}$. Hence $\omega_t(\mathfrak J)\le\mathsf o_t$; $\omega_t$ is a seminorm.
\item[(c)] Let $m\in\mathbb N$, $\ell=2L^m$, $\gamma^{-2}\le\mathsf b_-\le\mathsf b_+$ with $MR(\mathsf b_+)\le L^m$, $0<\phi\le\frac14$, and fix a type. Let $K\ge K_\phi$, $u=\lceil\phi K\rceil$, $t=K-u$, $\mathsf b'\in[\mathsf b_-,\mathsf b_+]$ and $\beta=L^K\mathsf b'$, so that (Th) holds on $T_\ell$ with $(K,\mathsf b',m,t)$ (Lemma~\ref{stc:frame}\,(b)). With $C$ a function of $m$, $\mathsf b_+$, $\phi$ and the type: $\mathsf o_t\le CK^{2\mathfrak p_0}L^{4u-K}|\varkappa|_\infty$ at order $2$, $\mathsf o_t\le CK^{3\mathfrak p_0}L^{4u-\frac32K}|\varkappa|_\infty$ at orders $\ge3$, and $1/r_{\mathfrak J}\le C|\varkappa|_\infty$. If $0<\phi\le\frac1{84}$ and $K\ge\max\{K_\phi,5\}$, these bounds hold with $C$ a function of $m$, $\mathsf b_+$ and the type alone, not of $\phi$, and with such a $C$ also $\max_i\mathsf d_i\le C|\varkappa|_\infty$ and, at order $2$, $\sum_{i\le n_0}L^i\mathsf d_i\le CK^{\mathfrak p_0}L^{-(K-n_0)/2}|\varkappa|_\infty$ for $1\le n_0\le t$. First-passage data $(K,\beta)\in\mathfrak d(g)$ on a side $\ell\in\mathfrak L(g)$ are the case $\mathsf b_-=\gamma^{-2}$, $\mathsf b_+=X_g+\bar B$, $\mathsf b'=b$ (Lemma~\ref{stc:frame}\,(a)).
\item[(d)] For the data of (c), now with any $0<\phi<1$, $\sup|\mathfrak J|\le N\hat c\,\ell^3L^{3K}|\varkappa|_\infty$ at real configurations, and, with the numbers $c'$, $A_0$, $\mathsf C^\Sigma_{\rm maj}$, $C_{\rm low}(x)=c_{\rm low}+\frac92\log x$ of the definition \YMMref{115.2}{806} of $\mathsf R_t$,
\[
\mathsf R_t\le C_{\rm R}L^{3u}\exp\bigl[-c'A_0^2(\phi K\log L)^{2\mathfrak p_0}\bigr],\qquad C_{\rm R}=2M^{-3}\ell^3\exp\bigl[(\mathsf C^\Sigma_{\rm maj}+2+C_{\rm low}(\mathsf b_+))\ell^3\bigr],
\]
uniformly in $\mathsf b'\in[\gamma^{-2},\mathsf b_+]$; for $\phi K\log L\ge1$ the exponent is $\le-c'A_0^2(\phi\log L)^2K^2$, so $\mathsf R_t$ lies below every power of $L^{-K}$ for large $K$.
\end{enumerate}
\end{lemma}

\begin{proof}[Proof sketch]
(a) If $\mu_k=h\nu_k$, then $\mu_{k+1}=(\mathsf Q_kh)\nu_{k+1}$, by the expression of the step through the chart integral (which depends on its function at the points of $\mathsf{Cut}_k$ alone, and is a positive kernel at zero source); induction on $k$. At zero source $\mathsf Q_k$ averages against a finite positive measure, whence the contraction ($|\int_E(\mathsf Q_kh-a_0)\,d\nu_{k+1}|\le\|h-a_0\|_\infty\nu_{k+1}(E)$ for Borel $E$). For the last identity, perturb the early member of level $i'$ by $\mu\mathsf g$, admissible for $|\mu|\|\mathsf g\|<\frac1{16}E_0$ by the room between $\frac{15}{16}E_0$ and $E_0$, and take the derivative in $\mu$ at $0$ of the identity of the step for the run with this early member: $\mathsf S_k\mathsf g$ is the derivative of the next family, bounded by the Lipschitz bound of a run (the one-step estimate in the proof of Proposition~\YMMref{114.3}{803}, with $\Delta_{i'}=|\mu|\|\mathsf g\|$ and $\Delta_i=0$ at the other levels $i\le k$), and local by the locality of one step. (b) Integrate the display of Proposition~\ref{stc:chain}\,(iii) against a bounded Borel $h$: both sides are holomorphic in $s$; by (iv), $|A_k|\le|s|\mathsf o_k$, so Cauchy's estimates give $|\partial_sA_k(0)|\le\mathsf o_k$ and $|A_k/s-\partial_sA_k(0)|\le2|s|\mathsf o_k/r_{\mathfrak J}$; differentiate at $s=0$ under the integrals (dominated convergence) and compare with (a). (c) Insert $x_t=L^u\mathsf b'$, $(\mathsf e^{\rm c}_t)^2=c_\rho^2A_0^2(\log x_t)^{2\mathfrak p_0}/x_t\le CK^{2\mathfrak p_0}L^{-u}$ (as $\log x_t\le K\log L+\log\mathsf b_+$ and $x_t\ge\gamma^{-2}L^u$), $N_t=\ell^3L^{3u}$, $\delta(x_0)\le CK^{\mathfrak p_0}L^{-K/2}$ into the definitions; the recursion gives $L^t\mathsf n_t\le C(\mathsf d_1+\sum_j\theta_jL^j\mathsf m_j)$ and $L^{j+1}\mathsf d_{j+1}=L\theta_j(L^j\mathsf n_j+L^j\mathsf m_j)$, and Lemma~\ref{stc:frame}\,(b) sums the $L\theta_j$.

At orders $\ge3$, $\mathsf e^{\rm c}_t\le1$ and $L^{-3t}\le L^{-t-K/2}$ because $t\ge(1-\phi)K-1\ge\frac K4$, as $\phi\le\frac14$ and $K\ge2$. (d) The height at $|U(\mathfrak b)|=1$ and at most $N(\ell L^K)^3$ indices; for $\mathsf R_t$ insert $N_t=\ell^3L^{3u}$, $\log x_t\ge u\log L\ge\phi K\log L$ ($\mathsf b'>1$), $C_{\rm low}(\mathsf b')\le C_{\rm low}(\mathsf b_+)$ and $\mathfrak p_0\ge1$ into \YMMref{115.2}{806}.
\end{proof}

\begin{proposition}[Cumulants of distant insertions on one torus]\label{stc:cum}
\begin{enumerate}
\item[(a)] Let (Th) hold on $T_\ell$ with top level $t$, let $\mathsf R_t\le\frac12$, and let $\mathfrak J_0,\dots,\mathfrak J_r$, $r\ge1$, be insertions of orders $\ge2$ whose carriers are pairwise at sup distance $\ge\mathsf d_\times>0$ on $T_\ell$. Let $r_*$ and $\bar{\mathsf d}$ be the numbers $r_{\mathfrak J}$ and $\max_i\mathsf d_i$ of Proposition~\ref{stc:chain} formed at order $2$ with the largest constants and the largest weight among the $\mathfrak J_a$, $r_1=\min\{\frac12r_*,(16\sum_a\mathsf o_t(\mathfrak J_a))^{-1}\}$ and $\mathsf t_\times=e^2C_{{\rm E},0}\sum_{i\le t}N_i\exp[-\mathsf d_\times L^{K-i}/M]$, $C_{{\rm E},0}$ the number of \eqref{eq:analytic-domain-tail}. If $32(r+2)r_1\bar{\mathsf d}\,\mathsf t_\times\le1$, then
\[
|\langle\mathfrak J_0;\dots;\mathfrak J_r\rangle_{K,\ell,\beta}|\le\mathsf{op}_{r+1}\prod_a\omega_t(\mathfrak J_a)+9\cdot2^{r+1}(r+2)\,\bar{\mathsf d}\,\mathsf t_\times r_1^{-r}+4\,\mathsf{op}'_{r+1}\mathsf R_t\prod_a\sup|\mathfrak J_a| .
\]
\item[(b)] (A form with any bounds of the conditional means.) In (Th), $1\le t\le t_*$, $\mathsf R_t\le\frac12$, let $\mathfrak J_a$, $a\in\mathsf E$, $\mathsf r=\#\mathsf E\ge2$, be insertions of orders $\ge2$ with carriers pairwise at sup distance $\ge\mathsf d_\times$ on $T_\ell$, $\mathsf y=(\mathsf d_\times-a_K)L^{K-t}/M\ge4$, $\mathsf t'_\times=4e^2C_{{\rm E},0}N_te^{-\mathsf y}$, and let $\epsilon_a>0$, $\mathsf a_a\in\mathbb C$ satisfy $|\mathsf Q_{t\leftarrow0}\mathfrak J_a-\mathsf a_a|\le\epsilon_a$ almost everywhere for $\mathbf 1_{\mathsf{In}_t}\nu_t$, and $2E_0\mathsf t'_\times\prod_{a\in I}(\mathsf r/r_{\mathfrak J_a})\le\prod_{a\in I}\epsilon_a$ for all $I\subset\mathsf E$ with $\#I\ge2$. Then the joint cumulant of the $\mathfrak J_a$ under the small-field functional (Definition~\ref{stc:ins}\,(3)) is at most $(\frac12+\mathsf{op}'_{\mathsf r}\mathsf{Bl}_{\mathsf r})\prod_a\epsilon_a$ in modulus, and the joint cumulant under $\langle\cdot\rangle_{K,\ell,\beta}$ differs from it by at most $4\,\mathsf{op}'_{\mathsf r}\mathsf R_t\prod_a\sup|\mathfrak J_a|$.
\end{enumerate}
\end{proposition}

\begin{proof}[Proof sketch]
(a) Let $Z_\sigma(s)$ be the small-field functional at the top level $t$ applied to $e^{s\mathfrak J_\sigma}$. Since $\sum_\sigma(-1)^{r+1-|\sigma|}\prod_{a\in B}\sigma_a$ is $1$ for $B=\{0,\dots,r\}$ and $0$ otherwise,
\[
(r+1)!\,{\rm cum}^{\rm sf}(\mathfrak J_0,\dots,\mathfrak J_r)=\sum_{\sigma\in\{0,1\}^{r+1}}(-1)^{r+1-|\sigma|}\,\partial_s^{r+1}\log Z_\sigma(0).
\]
By Proposition~\ref{stc:chain}\,(ii), $\Phi_\sigma-\Phi_0-\sum_a\sigma_aA_a$, where $A_a$ is the change of the exponent for $\mathfrak J_a$ alone, is a sum over polymers that contain boxes of two insertions, and so is the same combination of the constants $\sum_{i\le t}(\mathsf c_i(s)-\mathsf c_i(0))=\sum_{i\le t}\sum_{X,z}\delta E^{(i)}_s(X,\mathbf 1,z)$ (proof of Proposition~\ref{stc:chain}\,(iii)); by \eqref{eq:analytic-domain-diameter} such a polymer of level $i$ has linear size $\ge\mathsf d_\times L^{K-i}/M-2$, and by (i) and \eqref{eq:analytic-domain-tail} the sum is at most $2(r+2)|s|\bar{\mathsf d}\,\mathsf t_\times$. With $\lambda\propto\mathbf 1_{\mathsf{In}_t}e^{\Phi_0}dV$ a probability measure, Proposition~\ref{stc:chain}\,(iii) gives, on $|s|\le r_1$, $\log(Z_\sigma/Z_0)=(\text{linear in }\sigma)+\log\lambda(\exp\sum_a\sigma_aA_a)+O(\bar{\mathsf d}\mathsf t_\times)$.

The alternating sum annihilates the linear part and leaves $s^{r+1}{\rm cum}_\lambda(\dot A_0,\dots,\dot A_r)+O(s^{r+2})$; Cauchy's estimate on $|s|=r_1$ bounds the rest by the second member. By Lemma~\ref{stc:cmean}\,(b), $\dot A_a=\mathsf Q_{t\leftarrow0}\mathfrak J_a-{\rm const}$ $\lambda$-almost everywhere, and a cumulant of $r+1$ bounded functions, from which constants drop out, is at most $\mathsf{op}_{r+1}$ times the product of their sup norms. The discarding statement among the construction results (\textup{\YMMref{115.2}{806}}: the small-field functional lies between the term without large fields and the bare density, for all bounded Borel gauge invariant functions) gives $|Z\langle F\rangle-\Lambda^{\rm sf}(F)|\le\sup|F|\,\mathsf R_tZ$ for real gauge invariant $F$, and the expression of a cumulant through means gives the last member. (b) The generating function $T(\mathbf s)=\Lambda^{\rm sf}(e^{\sum_as_a\mathfrak J_a})$ is entire, and the cumulant is $\partial^{\mathsf E}_{\mathbf s}\log T(0)$, $\partial^{\mathsf E}_{\mathbf s}=\prod_{a\in\mathsf E}\partial_{s_a}$.

Let $E^{(i)}[\mathbf s]$ be the families of Proposition~\ref{stc:chain}\,(v) for the $\mathfrak J_a$, $\delta E^{(i)}_{\mathbf s}=E^{(i)}[\mathbf s]-E^{(i)}$, and put
\[
\begin{multlined}
\mathsf C(\mathbf s)=\sum_{i\le t}\bigl(\mathsf c_i(\mathbf s)-\mathsf c_i(0)\bigr)=\sum_{i\le t}\sum_{X,z}\delta E^{(i)}_{\mathbf s}(X,\mathbf 1,z),\qquad\mathsf D_{\mathbf s}(V)=\Phi^{\mathbf s}_t(V)-\Phi_t(V) \\
=\sum_as_a\mathfrak J_a(U)+\sum_{i\le t}\sum_{X,z}\bigl[\delta E^{(i)}_{\mathbf s}(X,U,z)-\delta E^{(i)}_{\mathbf s}(X,\mathbf 1,z)\bigr]
\end{multlined}
\]
for $V\in\mathsf{In}_t$ and a background $U$ over $V$, the first identity by the formula for the numbers $\mathsf c_i(\mathbf s)$ in the proof of Proposition~\ref{stc:chain}\,(iii). Multiplied by $\mathbf 1_{\mathsf{In}_t}$, the display of Proposition~\ref{stc:chain}\,(iii), which holds for $\sum_as_a\mathfrak J_a$ by (v), gives, on the parameter domain $\{\sum_a|s_a|/r_{\mathfrak J_a}<1\}$,
\[
T(\mathbf s)=T(0)\,e^{\mathsf C(\mathbf s)}\,\langle e^{\mathsf D_{\mathbf s}}\rangle ,
\]
where $\langle\cdot\rangle$ is the mean under the probability measure $\mathbf 1_{\mathsf{In}_t}\nu_t/\nu_t(\mathsf{In}_t)$, proportional to $\mathbf 1_{\mathsf{In}_t}e^{\Phi_t}dV$, and $T(0)=\nu_t(\mathsf{In}_t)\ge(1-\mathsf R_t)Z_{K,\ell,\beta}\ge\frac12Z_{K,\ell,\beta}>0$ by the discarding statement (\textup{\YMMref{115.2}{806}}) and $\mathsf R_t\le\frac12$. Hence the cumulant is $\partial^{\mathsf E}_{\mathbf s}\mathsf C(0)+\partial^{\mathsf E}_{\mathbf s}{\rm Log}\langle e^{\mathsf D_{\mathbf s}}\rangle|_{\mathbf s=0}$. The term of $(i,X,z)$ in $\mathsf C$ and in $\mathsf D_{\mathbf s}$ depends on the $s_a$ with $\mathfrak J_a$ meeting $X$ alone (locality, Proposition~\ref{stc:chain}\,(v)), and $\sum_as_a\mathfrak J_a$ is linear in $\mathbf s$; so a jet $\partial^I_{\mathbf s}$, $\#I\ge2$, receives contributions only from polymers $X$ that meet the carriers of two insertions. Such a polymer of level $i$ holds two points at sup distance $\ge\mathsf d_\times$, so $d_i(X)\ge\mathsf d_\times L^{K-i}/M-2\ge\mathsf yL^{t-i}-2$ (\eqref{eq:analytic-domain-diameter}), and by \eqref{eq:analytic-domain-tail}, over the $N_i=L^{3(t-i)}N_t$ sites of level $i$,
\[
\sum_{i\le t}\ \sum_{z}\ \sum_{X\ni z\ \text{meeting two carriers}}e^{-(\kappa+1)d_i(X)}\le e^2C_{{\rm E},0}N_t\sum_{l\ge0}L^{3l}e^{-\mathsf yL^l}\le2e^2C_{{\rm E},0}N_te^{-\mathsf y}=\tfrac12\mathsf t'_\times ,
\]
the quotient of two successive terms being at most $L^3e^{-4(L-1)}\le\frac12$ as $\mathsf y\ge4$. On this parameter domain $\|E^{(i)}[\mathbf s]\|<E_0$ by (v) and the choice of the $r_{\mathfrak J_a}$, so $|\delta E^{(i)}_{\mathbf s}(X,\cdot,z)|\le2E_0e^{-(\kappa+1)d_i(X)}$ at the background and at $\mathbf 1$. $\mathsf D_{\mathbf s}$ holds two values of each such term (at $U$ and at $\mathbf 1$), $\mathsf C$ one; so the parts of $\mathsf D_{\mathbf s}(V)$ and of $\mathsf C$ that $\partial^I_{\mathbf s}$ sees are holomorphic and of moduli at most $2E_0\mathsf t'_\times$ and $E_0\mathsf t'_\times$. This domain contains the polydisc of radii $r_{\mathfrak J_a}/\mathsf r$, and Cauchy's estimate for one derivative in each variable of $I$ gives
\[
\bigl|\partial^I_{\mathbf s}\mathsf D_{\mathbf s}(V)|_{\mathbf s=0}\bigr|\le2E_0\mathsf t'_\times\prod_{a\in I}\frac{\mathsf r}{r_{\mathfrak J_a}}\le\prod_{a\in I}\epsilon_a,\qquad\bigl|\partial^{\mathsf E}_{\mathbf s}\mathsf C(0)\bigr|\le E_0\mathsf t'_\times\prod_{a\in\mathsf E}\frac{\mathsf r}{r_{\mathfrak J_a}}\le\tfrac12\prod_{a\in\mathsf E}\epsilon_a ;
\]
the second is the term $\frac12$ of the bound. The first jets are $\partial_{s_a}\mathsf D_{\mathbf s}|_{\mathbf s=0}=\mathsf Q_{t\leftarrow0}\mathfrak J_a-\partial_{s_a}\mathsf C(0)$ almost everywhere for $\mathbf 1_{\mathsf{In}_t}\nu_t$ (Lemma~\ref{stc:cmean}\,(b) for $\mathfrak J_a$, with $\dot c_t=\partial_{s_a}\mathsf C(0)$), so within $\epsilon_a$ of the constant $\mathsf a'_a=\mathsf a_a-\partial_{s_a}\mathsf C(0)$. With $\tilde{\mathsf D}_{\mathbf s}=\mathsf D_{\mathbf s}-\sum_as_a\mathsf a'_a$, ${\rm Log}\langle e^{\mathsf D_{\mathbf s}}\rangle={\rm Log}\langle e^{\tilde{\mathsf D}_{\mathbf s}}\rangle+\sum_as_a\mathsf a'_a$ near $0$, and the linear function is killed by $\partial^{\mathsf E}_{\mathbf s}$ as $\mathsf r\ge2$. Moments by partitions, $\partial^I_{\mathbf s}e^{\tilde{\mathsf D}_{\mathbf s}}|_{\mathbf s=0}=\sum_{\pi}\prod_{I'\in\pi}\partial^{I'}_{\mathbf s}\tilde{\mathsf D}_{\mathbf s}|_{\mathbf s=0}$ over the partitions $\pi$ of $I$, bounded by $\mathsf{Bl}_{\#I}\prod_{a\in I}\epsilon_a$ almost everywhere, and logarithms by partitions give $|\partial^{\mathsf E}_{\mathbf s}{\rm Log}\langle e^{\tilde{\mathsf D}_{\mathbf s}}\rangle|_{\mathbf s=0}|\le\mathsf{op}_{\mathsf r}\mathsf{Bl}_{\mathsf r}\prod_a\epsilon_a\le\mathsf{op}'_{\mathsf r}\mathsf{Bl}_{\mathsf r}\prod_a\epsilon_a$. The discarding is as in (a).
\end{proof}

\begin{lemma}[Regular configurations; a family at a background]\label{stc:reg}
Let (Th) hold on $T_\ell$, $\ell=2L^m$, with $(K,\mathsf b',m,t_*)$ and $\gamma^{-2}\le\mathsf b'\le\mathsf b_+$.
\begin{enumerate}
\item[(a)] For $1\le j\le n\le t_*$, $V\in\mathsf{Cut}_n$ and $U\in\mathcal C(\mathbf B_n,V)$: $\varepsilon_n:=2\mathsf e^{\rm c}_n/\varepsilon_{\mathcal W}\le\alpha_0$ ($\varepsilon_{\mathcal W}$ the number of Proposition~\YMMref{52.3}{433}, which enters the field cuts \YMMref{113.10}{796}, \YMMref{113.10}{796}), and $U$ is regular of level $n$ at the radius $\varepsilon_n$ on $T_\ell$ in the sense of Definition~\ref{def:level-regularity}: $|\partial U(p)-\mathbf 1|<\varepsilon_nL^{-2n}$ at every bare plaquette and $|J^{(n)}(U)(\mathfrak b)|<\varepsilon_n$ at every positively oriented bare bond, $J^{(n)}$ the current of level $n$ (Definition~\ref{def:level-regularity}). Moreover $U\in\mathcal C(\mathbf B_j,\mathcal M^jU)$ and $\mathcal M^jU\in\mathsf{Cut}_j$; the same holds for $\mathsf rU$, $\mathsf r$ any isometry of the block lattice of level $j$, which is again regular of level $n$ at the radius $\varepsilon_n$; every member of every pointed family of level $j$ is formed at $U$, at $\mathsf rU$ and at $\mathbf 1$. Finally $\varepsilon_n^2\le\bar{\mathsf c}_\varepsilon(K\log L+\log\mathsf b_+)^{2\mathfrak p_0}L^{-(K-n)}$, $\bar{\mathsf c}_\varepsilon=4c_\rho^2A_0^2\gamma^2/\varepsilon_{\mathcal W}^2$.
\item[(b)] Let $\mathsf f$ be a pointed family of level $j$, $1\le j\le n\le t_*$, $V\in\mathsf{Cut}_n$, $U\in\mathcal C(\mathbf B_n,V)$. Then
\[
\bigl|\mathsf{Sum}[\mathsf f](U)-\mathsf{Sum}[\mathsf f](\mathbf 1)\bigr|\le\|\mathsf f\|\Bigl\{32B_*^2M^2C_{{\rm E},2}r_\alpha^{-2}\,\varepsilon_n^2\,N_nL^{-(n-j)}+2C_{{\rm E},0}N_je^{-[D_*-1]_+}\Bigr\},
\]
where $B_*\ge1$ and $c_{*,{\rm rg}}\in(0,1]$ are the constants of Theorem~\YMMref{F}{199}, with $c_{*,\mathrm{rg}}=c_*$ as in Section~\ref{sec:regular-background-complements}, $r_\alpha=a^{\rm us}(\alpha_0,\alpha_1)$ is the usable radius \YMMref{(933)}{427}, $C_{{\rm E},p}$ are the numbers of Proposition~\YMMref{45.1}{353} and Lemma~\ref{lem:analytic-domain-geometry}, and, with $\alpha=\varepsilon_nL^{-2(n-j)}$,
\[
D_*=D_*(j,n)=\min\Bigl\{\frac{c_{*,{\rm rg}}}{4M\alpha},\ \frac{r_\alpha}{8B_*M\alpha},\ \frac{L^{m+K-j}-5}{2M}\Bigr\}:
\]
for a polymer $X$ of level $j$ with $d_j(X)+1\le D_*$, Lemma~\ref{lem:regular-cube-gauge}, on a cube of side $2(d_j(X)+1)M$ of level $j$ that holds $X$, gives a gauge in which $U=e^{iL^{-j}A}$ on $X$ with $|A|,|\nabla A|<4B_*(d_j(X)+1)M\alpha\le\frac12r_\alpha$. If $t_*\le K-\phi K$, then $K-j\ge K-n\ge\phi K$ and $\alpha\le\varepsilon_n$, so by (a) the first two members of $D_*$ are at least $cL^{\phi K/2}(K\log L+\log\mathsf b_+)^{-\mathfrak p_0}$, $c$ a function of the constants, and the third is at least $(L^{\phi K}-5)/(2M)$; hence $D_*>7K\log L+1$ for $K$ large in terms of $\mathsf b_+$ and $\phi$. So for $0<\phi<1$, $t_*\le K-\phi K$ and $K\ge K_2$, a function of $\mathsf b_+$ and $\phi$,
\[
\begin{multlined}
\bigl|\mathsf{Sum}[\mathsf f](U)-\mathsf{Sum}[\mathsf f](\mathbf 1)\bigr|\le\mathsf B_nL^j\|\mathsf f\|,\qquad\mathsf B_n=\mathsf B'_0L^{3(K-n)-K},\quad\mathsf B'_0 \\
=\ell^3\bigl[32B_*^2M^2C_{{\rm E},2}r_\alpha^{-2}\bar{\mathsf c}_\varepsilon(K\log L+\log\mathsf b_+)^{2\mathfrak p_0}+2C_{{\rm E},0}\bigr]\le CK^{2\mathfrak p_0}.
\end{multlined}
\]
\item[(c)] (Petals in a H\"older gauge.) Let $0<\kappa_{\rm H}<1$, $h=L^{-n}$, and let $U$ be regular of level $n$ at a radius $\varepsilon\le1$ with $6\rho_{\rm H}<L^{m+K-n}$, $\rho_{\rm H}\varepsilon\le a_{\rm g}$, $h\,\mathsf c_{\rm g}\rho_{\rm H}\varepsilon\le\frac14$, where $\rho_{\rm H}$, $a_{\rm g}$, $\mathsf c_{\rm g}$, $\mathsf c_{\rm g}(\kappa_{\rm H})$ are the numbers of Theorem~\textup{\ref{I-56-65}}; put $e_1=\mathsf c_{\rm g}\rho_{\rm H}\varepsilon$, $e_2=\mathsf c_{\rm g}(\kappa_{\rm H})\rho_{\rm H}\varepsilon$. At every site $x$ there are $F_{12}(x),F_{13}(x),F_{23}(x)\in\mathfrak g$ with, for all petals $\pi,\pi'$ at $x$,
\[
\begin{multlined}
|q_x(\pi,\pi')-\mathsf s(\pi)\mathsf s(\pi')h^4\langle F_{m(\pi)},F_{m(\pi')}\rangle|\le D(2\varepsilon h^2+D),\qquad|a_{p_m(x)}-\tfrac12h^4|F_m|^2| \\
\le\tfrac12D(2\varepsilon h^2+D)+\tfrac18\varepsilon^4h^8,
\end{multlined}
\]
$D=4e_2h^{2+\kappa_{\rm H}}+(21+29e_1)e_1^2h^3$; and if moreover $\mathsf c_{\rm g}\rho_{\rm H}\varepsilon\le1$, then 
\[
|\mathbf y_{\mathfrak b,\rho}(U)|\le48\,\mathsf c_{\rm g}(\kappa_{\rm H})\rho_{\rm H}\,\varepsilon\,h^{2+\kappa_{\rm H}}
\]
 at every pair $(\mathfrak b,\rho)$.
\end{enumerate}
\end{lemma}

\begin{proof}[Proof sketch]
(a) For a block field $V$ of level $n$ the largest plaquette defect is $\max_p|\partial V(p)-\mathbf 1|$, the maximum over the plaquettes $p$ of the block lattice of level $n$ (the plaquette size $\varpi(V)$). On $\mathsf{Cut}_n$ it is $<2\mathsf e^{\rm c}_n$ (\YMMref{113.10}{796}), and $2\mathsf e^{\rm c}_n\le\varepsilon_{\mathcal W}\alpha_0$, $\mathsf e^{\rm c}_n\le\frac12c^{\rm SO}_2$ (\YMMref{(2134)}{790}); so $\varepsilon_n\le\alpha_0\le1$, and $2\mathsf e^{\rm c}_n=\varepsilon_{\mathcal W}\varepsilon_n$. The field cut is $\mathsf{Cut}_n=\mathcal C^\flat_n(\mathsf e^{\rm c}_n)$ (\YMMref{113.10}{796}), and the window of the classes of covariance, Lemma~\ref{lem:sf-class-window}(1) at the radii $(b_0,b_1)=(\varepsilon_n,1)$ and $\mathsf e=\mathsf e^{\rm c}_n$ (as $2\mathsf e=\varepsilon_{\mathcal W}b_0$), gives $\mathsf{Cut}_n=\mathcal C_n(\varepsilon_n,1;\mathsf e^{\rm c}_n)$: the pairs of every $U\in\mathcal C(\mathbf B_n,V)$ lie in the cores at the radii $(\varepsilon_n,1)$, so that $\mathsf m_n(U)<\varepsilon_n$ and $|J^{(n)}(U)(\mathfrak b)|<\varepsilon_n$ at every positively oriented bare bond (Lemma~\ref{lem:sf-class-window}(1)): $U$ is regular of level $n$ at the radius $\varepsilon_n$, current included.

Heredity, Lemma~\ref{lem:sf-class-window}(2): as $x_n=L^{-(n-j)}x_j$ gives $\mathsf e^{\rm c}_n\le L^{(n-j)/2}\mathsf e^{\rm c}_j$, hence $L^{-2(n-j)}\mathsf e^{\rm c}_n\le\mathsf e^{\rm c}_j$, the background over $\mathcal M^jU$ exists, $U\in\mathcal C(\mathbf B_j,\mathcal M^jU)$, and $\mathcal M^jU\in\mathcal C^\flat_j(\mathsf e^{\rm c}_j)=\mathsf{Cut}_j$. An isometry $\mathsf r$ of the block lattice of level $j$ commutes with block averaging up to the level $j$, $\mathcal M^j(\mathsf rU)=\mathsf r\mathcal M^jU$, and carries the minimal orbits of the level $j$ and the sets $\mathcal C^\flat_j(\mathsf e)$ into themselves (Lemma~\ref{lem:setup-lattice-covariance}, Proposition~\YMMref{52.3}{433}, and Lemma~\ref{lem:sf-class-window}(3)); so $\mathcal M^j(\mathsf rU)\in\mathsf{Cut}_j$ and $\mathsf rU\in\mathcal C(\mathbf B_j,\mathcal M^j(\mathsf rU))$; it keeps plaquette defects and currents bond by bond (Lemma~\ref{lem:setup-lattice-covariance}), so $\mathsf rU$ is regular of level $n$ at the radius $\varepsilon_n$.

As $2\mathsf e^{\rm c}_j\le\varepsilon_{\mathcal W}\alpha_0$, the window at the radii $(\alpha_0,\alpha_1)$ gives $\mathsf{Cut}_j=\mathcal C_j(\alpha_0,\alpha_1;\mathsf e^{\rm c}_j)$: the pairs of $U$ and of $\mathsf rU$ at every polymer $X$ of level $j$ lie in the cores $\mathcal U^\circ_j(X;\alpha_0,\alpha_1)$ on which the members of the pointed families of level $j$ are local terms, and so does the trivial pair (the usable-ball inclusion in Section~\ref{sec:analytic-core-complements}); so every such member is formed at $U$, at $\mathsf rU$ and at $\mathbf 1$. Finally $x_n=L^{K-n}\mathsf b'\ge\gamma^{-2}L^{K-n}$, $\log x_n\le K\log L+\log\mathsf b_+$, and $\varepsilon_n^2=4c_\rho^2A_0^2\varepsilon_{\mathcal W}^{-2}(\log x_n)^{2\mathfrak p_0}/x_n$ (the definitions preceding Lemma~\YMMref{113.10}{796}). (b) Fix a polymer $X$ of level $j$, $d=d_j(X)$, and a member $\mathsf E$, $\|\mathsf E\|_X\le\|\mathsf f\|e^{-(\kappa+1)d}$.

If $d+1\le D_*$, the gauge of Lemma~\ref{lem:regular-cube-gauge} on $X$ gives $U^u=e^{iL^{-j}A}$ and $v=(A,\operatorname{Ad}(u)J^{(j)}(U))$ of size $\le4B_*(d+1)M\varepsilon_n\vartheta^2\le\frac12r_\alpha$, $\vartheta=L^{-(n-j)}$, and the local term is formed at $U$ with $\mathsf E(U)=\mathsf E(v)$, $\mathsf E(\mathbf 1)=\mathsf E(0)$, by pair dependence and gauge invariance (Definition~\YMMref{49.4}{375}; Section~\ref{sec:regular-background-complements}). $v\mapsto\mathsf E(v)$ is holomorphic on the ball of radius $r_\alpha$ and its differential at $0$ vanishes: it is invariant under the constant gauge transformations $g\in\{\mathbf 1,i\sigma_1,i\sigma_2,i\sigma_3\}$, and $\sum_ggYg^{-1}=0$ for $Y$ of trace $0$; by the maximum principle $|\mathsf E(v)-\mathsf E(0)|\le2\|\mathsf E\|_X(\|v\|/r_\alpha)^2$. If $d+1>D_*$ the difference is $\le2\|\mathsf E\|_X$.

Sum with \eqref{eq:analytic-domain-tail} over the $N_j$ sites; $N_j\vartheta^4=N_nL^{-(n-j)}$. The second form follows from (a), $N_n=\ell^3L^{3(K-n)}$ and $N_je^{-[D_*-1]_+}\le\ell^3L^{-4K}$. (c) Theorem~\textup{\ref{I-56-65}} gives, on a cube of side $2\rho_{\rm H}$ about $x$ in the units of level $n$, a gauge in which $U=e^{ihA}$, $A$ a $1$-form with values in $\mathfrak g$, with $|A|,|\nabla A|<e_1$ and $\nabla A$ H\"older continuous of exponent $\kappa_{\rm H}$ with constant $e_2$ (forward differences of spacing $h$). In this gauge put, for a plane $m=\{\mu<\nu\}$ and a site $x''$ of the cube,
\[
F_m(x''):=(\nabla_\mu A_\nu-\nabla_\nu A_\mu)(x'')+i[A_\mu(x''),A_\nu(x'')],
\]
the curvature in the gauge: it is Hermitian and of trace $0$, so $F_m(x'')\in\mathfrak g$, and the $F_m(x)$ of the statement are these at $x''=x$, in the gauge about $x$. Expanding the four exponentials of the plaquette variable $P_m(x'')=P_{\mu\nu}(x'')=W_{x''}(\mu,\nu)$ of $p_m(x'')$, with $\mathbf a=ihA_\mu(x'')$ and $\mathbf b=ihA_\nu(x''+he_\mu)$,
\[
P_{\mu\nu}(x'')-\mathbf 1=ih^2(\nabla_\mu A_\nu-\nabla_\nu A_\mu)(x'')+[\mathbf a,\mathbf b]+R,\qquad|R|\le(5+29e_1)e_1^2h^3,
\]
and $[\mathbf a,\mathbf b]=-h^2[A_\mu(x''),A_\nu(x'')]+R'$, $|R'|\le2e_1^2h^3$; the commutator, of size up to $2e_1^2h^2$, is the part $ih^2\cdot i[A_\mu,A_\nu]$ of $ih^2F_m$. So $|P_{\mu\nu}(x'')-\mathbf 1-ih^2F_m(x'')|\le(7+29e_1)e_1^2h^3$, and, taking the Hermitian part of $(P_{\mu\nu}(x'')-\mathbf 1)/i$, $|\mathbf w(P_{\mu\nu}(x''))-h^2F_m(x'')|\le(7+29e_1)e_1^2h^3$. The holonomy of a petal $\pi$ at $x$ is $W_x(\pi)=T\,P_{m(\pi)}(x'')^{\mathsf s(\pi)}\,T^{-1}$ for a plaquette $p_{m(\pi)}(x'')$ with $|x''-x|_1\le2h$, $T$ a product of at most three bond variables; so $|\mathbf w(W_x(\pi))-\mathsf s(\pi)\mathbf w(P_{m(\pi)}(x''))|\le2|T-\mathbf 1|\,\varepsilon h^2\le6e_1^2h^3$, as $\varepsilon\le e_1$. The H\"older bound of $\nabla A$ and $|A(x'')-A(x)|\le2he_1$ give $|F_m(x'')-F_m(x)|\le2e_2(2h)^{\kappa_{\rm H}}+8e_1^2h$. Together
\[
|\mathbf w(W_x(\pi))-\mathsf s(\pi)h^2F_{m(\pi)}(x)|\le4e_2h^{2+\kappa_{\rm H}}+(21+29e_1)e_1^2h^3=D,
\]
and $|\mathbf w(W_x(\pi))|<\varepsilon h^2$, a petal holonomy being, at a real configuration, a unitary conjugate of a plaquette variable or of its inverse. Both bounds follow by $|\langle\mathbf w,\mathbf w'\rangle-\langle\mathbf v,\mathbf v'\rangle|\le|\mathbf w-\mathbf v||\mathbf w'|+|\mathbf v||\mathbf w'-\mathbf v'|$, and by $a_p=\frac12|\mathbf w(\partial U(p))|^2+\frac12a_p^2$ with $0\le a_p<\frac12\varepsilon^2h^4$. $\mathbf y_{\mathfrak b,\rho}$ is the covariant difference of the curvatures of the two neighbouring plaquettes of one plane at $\mathfrak b$; in the gauge it is at most $2e_2h^{2+\kappa_{\rm H}}+46e_1^2h^3$.
\end{proof}

\begin{lemma}[A regular configuration against its mirror image]\label{stc:mirror}
Let $0\le\kappa_\star<1$, $r_1=(6e^3)^{-1}$, and in the notation of Lemma~\ref{stc:reg}\,(c), with $\kappa_\star$ for $\kappa_{\rm H}$, let $e_3=12e_2+2304e_1^2r_1^{-3}$.
\begin{enumerate}
\item[(a)] Let $1\le j\le n\le K$, $\vartheta=L^{-(n-j)}$, $6\rho_{\rm H}<L^{m+K-n}$, $U$ regular of level $n$ at the radius $\varepsilon$ on $T_\ell$ with $\rho_{\rm H}\varepsilon\le a_{\rm g}$, $e_1\le r_1^2$, $z$ a site, $\mathsf r_z$ the inversion through $z$, and $1\le R_\circ\le\vartheta^{-1}$. In the units of level $j$, on the cube $Q$ of radius $R_\circ$ about $z$, there are a gauge transformation $u$ and a real $1$-form $a'$ with $U^u=e^{iL^{-j}a'}$ on the bonds of $Q$ and $(\mathsf r_zU)^{\mathsf r_zu}=e^{iL^{-j}\mathsf r_z^*a'}$, $(\mathsf r_z^*a')_\mu(x)=-a'_\mu(2z-x-L^{-j}e_\mu)$; and with $\bar a=\frac12(a'+\mathsf r_z^*a')$, $w=\frac12(a'-\mathsf r_z^*a')$,
\[
|\bar a|,\ |\nabla\bar a|\le45R_\circ\vartheta^2e_1,\qquad|w|,\ |\nabla w|\le R_\circ^2\vartheta^{2+\kappa_\star}e_3,\qquad|J^{(j)}(U)|,\ |J^{(j)}(\mathsf r_zU)|<\vartheta^3\varepsilon .
\]
\item[(b)] Let (Th) hold, $1\le j<n_0<t_*$, $K-j+1\ge y_{\rm par}$ (the number of Lemma~\ref{stc:births}\,(e)), $\vartheta=L^{-(n_0-j)}$, $\varepsilon=\varepsilon_{n_0}$, with $6\rho_{\rm H}<L^{m+K-n_0}$, $\rho_{\rm H}\varepsilon\le a_{\rm g}$, $e_1\le r_1^2$, $180e_1\le r_\alpha$, $4e_3\le r_\alpha$. For a site $z$ of level $j$ let $\theta_z(p)=\prod_\lambda\vartheta_0((x(p)-z)_\lambda/(L^ja_K))$, $\vartheta_0=1$ on $(-\frac12,\frac12)$, $\frac12$ at $\pm\frac12$, $0$ elsewhere, so that $\theta_z\ge0$, $\sum_z\theta_z=1$ and $\theta_z\circ\mathsf r_z=\theta_z$; for a weight $\varkappa$ of the plaquettes put $\Sigma_z[\varkappa]=\mathsf{Sum}[\mathsf D^{(j)}[\theta_z\varkappa]]$ with the births of Lemma~\ref{stc:births}. Then for $V\in\mathsf{Cut}_{n_0}$ and $U\in\mathcal C(\mathbf B_{n_0},V)$
\[
\sum_z\bigl|\Sigma_z[\varkappa](U)-\Sigma_z[\varkappa](\mathsf r_zU)\bigr|\le C\,N_j\bar{\mathsf p}_j\max|\varkappa|\Bigl\{\frac{e_3(720e_1+11e_3)}{r_\alpha^2}\vartheta^{4+\kappa_\star}+\frac{\varepsilon^2}{r_\alpha^2}\vartheta^5+e^{-[D_*-1]_+}\Bigr\},
\]
$C$ a function of the constants, $D_*=D_*(j,n_0)$ the number of Lemma~\ref{stc:reg}\,(b), $\bar{\mathsf p}_j$ that of Lemma~\ref{stc:births}\,(e).
\end{enumerate}
\end{lemma}

\begin{proof}[Proof sketch]
(a) Theorem~\textup{\ref{I-56-65}} gives the H\"older gauge on the cube of side $2\rho_{\rm H}$ about $z$ in the units of level $n$; in the units of level $j$, $a=\vartheta A$ has $|a|<\vartheta e_1$, $|\nabla a|<\vartheta^2e_1$, and $\nabla a$ has the H\"older constant $\vartheta^{2+\kappa_\star}e_2$. A second gauge $e^{i\lambda}$, $\lambda(x)=\sum_\nu(x-z)_\nu c_\nu$ with $c_\mu$ the mean of $a_\mu$ at $z$ and $z-L^{-j}e_\mu$, is odd under $\mathsf r_z$; the new $1$-form is $a'_\mu(x)=L^j\Psi(\lambda(x);L^{-j}a_\mu(x),L^{-j}c_\mu)$, where $e^{i\Psi(X;Y,Z)}=e^{iX}e^{iY}e^{-i(X+Z)}$ defines a holomorphic $\Psi$ with $|\Psi-(Y-Z)|$ of second order and with the parity $G(-X,Y_o,-Y_e,-Z)=-G(X,Y_o,Y_e,Z)$ for $G(X,Y_o,Y_e,Z)=\Psi(X;Y_e+Y_o,Z)-\Psi(-X;-Y_e+Y_o,-Z)$.

The bound for $\bar a$ follows from $|\nabla a|$ and the Taylor estimate for $\Psi$. To see the extra gain for $w$, put $h=L^{-j}$, $x^*=2z-x-he_\mu$, $X=\lambda(x)$, $Y=ha_\mu(x)$, $Y^*=ha_\mu(x^*)$, $Z=hc_\mu$, $Y_o=(Y-Y^*)/2$ and $Y_e=(Y+Y^*)/2$. Since $\lambda(x^*)=-X-Z$, the inverse identity gives $2hw_\mu(x)=G(X,Y_o,Y_e,Z)$. The Taylor expansion of this function is $G=2(Y_e-Z)+i[2X+Z,Y_o]+R_{\ge3}$, where $R_{\ge3}$ starts at degree three and each of its monomials contains $Y$, $Y^*$ or $Z$, because $\Psi(X;0,0)=0$. The paired paths from $z$ to $x$ and from $z-he_\mu$ to $x^*$ cancel the constant gradient. The H\"older estimate therefore bounds $h^{-1}(Y_e-Z)$ and its forward gradient by $12R_\circ^2\vartheta^{2+\kappa_\star}e_2$. Moreover $|X|\le5R_\circ\vartheta e_1$, $|Y_o|\le4hR_\circ\vartheta^2e_1$ and $|Y|,|Y^*|,|Z|\le h\vartheta e_1$. Thus the quadratic term contains the smaller variable $Y_o$, and every higher term contains a bond variable. Cauchy's estimate in scalar directions at the radius $r_1$, with split radii for the first derivatives, using $e_1\le r_1^2$ and $R_\circ\vartheta\le1$, bounds their contributions to $w$ and $\nabla w$ by $800r_1^{-3}R_\circ^2\vartheta^3e_1^2$. Finally $\vartheta^3\le\vartheta^{2+\kappa_\star}$, which gives the claimed bound with $e_3$. The currents scale as $J^{(j)}=\vartheta^3J^{(n)}$. (b) Fix $z$, a polymer $X$ of level $j$ whose plaquettes meet the support of $\theta_z$, and a member $\mathsf E$ of $\mathsf D^{(j)}[\theta_z\varkappa]$ at $X$, $\|\mathsf E\|_X\le\bar{\mathsf p}_j\max|\varkappa|e^{-(\kappa+1)d}$. With $R_d=M(d+2)+3$: if $R_d\le\vartheta^{-1}$, by (a) the configurations $U$ and $\mathsf r_zU$, written in the gauge of (a), are $v=\bar v+w$ and $v^*=\bar v-w$ with $\|\bar v\|\le45R_d\vartheta^2e_1$, $\|w\|\le R_d^2\vartheta^{2+\kappa_\star}e_3$, and a local term with vanishing differential at $0$ (as in Lemma~\ref{stc:reg}\,(b)) changes by at most $\|\mathsf E\|_X\|w\|(16\|\bar v\|+11\|w\|)/r_\alpha^2$ between them; if $R_d>\vartheta^{-1}$ but $d+1\le D_*$, both values are within $2\|\mathsf E\|_X(\cdot/r_\alpha)^2$ of $\mathsf E(0)$ as in Lemma~\ref{stc:reg}\,(b), and $e^{-d}\le C\vartheta$ there; if $d+1>D_*$, the difference is $\le2\|\mathsf E\|_X$.

At most $R_d^3$ sites $z$ meet a polymer; sum with \eqref{eq:analytic-domain-tail}.
\end{proof}

\begin{lemma}[A cut at a level; the null part of the defect]\label{stc:cut}
Let $\ell=2L^m\in\mathfrak L(g)$, $(K,\beta)\in\mathfrak d(g)$, $0<\phi\le\frac1{84}$, $u=\lceil\phi K\rceil$, $t=K-u$, and let $K\ge\max\{K_\phi,K_2,K_\Theta\}$, $K_2$ the number of Lemma~\ref{stc:reg}\,(b) and $K_\Theta$ the least integer from which on $\Theta_KL^{-\phi K/2}\le1$.
\begin{enumerate}
\item[(a)] Let $1\le n_0\le t$, $\mathfrak J$ an insertion of order $2$, and $\epsilon_X=\sup|\mathfrak J(U)|$ over the $U\in\mathcal C(\mathbf B_{n_0},V)$, $V\in\mathsf{Cut}_{n_0}$. Then
\[
\omega_t(\mathfrak J)\le\epsilon_X+e\,\mathsf B'_0L^{3u-K}\sum_{i\le n_0}L^i\mathsf d_i\le\epsilon_X+CK^{3\mathfrak p_0}L^{3u-K-(K-n_0)/2}|\varkappa|_\infty .
\]
\item[(b)] Let $\eta$ be a $C^2$ vector field on $T_\ell$ (or $\eta\in C^2_c(\mathbb R^3;\mathbb R^3)$ carried by the interior of a cell of side $\frac12\ell$ and regarded on $T_\ell$), $0<\kappa_{\rm H}<1$, $1\le n\le t$, $y=K-n$, $\mathsf r$ an isometry of the block lattice of level $n$, and let $K$ be so large that the data of Lemma~\ref{stc:reg}\,(c) hold at $\varepsilon=\varepsilon_n$. Then, the supremum being over the $U\in\mathcal C(\mathbf B_n,V)$, $V\in\mathsf{Cut}_n$,
\[
\sup|F(k^0,v^0)(\mathsf r^{-1}U)|\le CK^{2\mathfrak p_0}L^{(3+\kappa_{\rm H})y-\kappa_{\rm H}K}\,b\sup|\nabla\eta| .
\]
\end{enumerate}
\end{lemma}

\begin{proof}[Proof sketch]
(a) Let $0<s\le\frac12r_{\mathfrak J}$ be real. On $\mathsf{Cut}_{n_0}$, by Proposition~\ref{stc:chain}\,(iii), $s^{-1}(\Phi^s_{n_0}-\Phi_{n_0})=\mathfrak J\circ U_{n_0}+\sum_{i\le n_0}(\Sigma_{n_0}[\mathsf f^{(i)}_s]-a[\mathsf f^{(i)}_s])$, with $\mathsf f^{(i)}_s=s^{-1}(E^{(i)}[s]-E^{(i)})$ of norm $\le\mathsf d_i$ and $a[\cdot]$ the value at the trivial configuration. Apply $\mathsf Q_{t\leftarrow n_0}$. A conditional mean is not larger than a supremum, so the first term stays $\le\epsilon_X$. By Lemma~\ref{stc:cmean}\,(a) a family of level $i$ is carried to level $t$ as $\Sigma_t[\mathsf T_{n_0}\mathsf f]$, $\mathsf T_{n_0}\mathsf f=\sum S_{k_m}\cdots S_{k_1}\mathsf f$ over the chains $n_0\le k_1<\dots<k_m<t$, at the price $\prod_k(1+L\theta_k)\le\exp[\Theta_KL^{-u/2}]\le e$; and at level $t$ a pointed family of level $\le t$ differs at a background from its value at the trivial configuration by at most $\mathsf B_tL^i\|\mathsf f\|$, $\mathsf B_t=\mathsf B'_0L^{3u-K}$ (Lemma~\ref{stc:reg}\,(b)).

Constants are carried to constants. Let $s\to0$ with Lemma~\ref{stc:cmean}\,(b) at $k=n_0$, and use (c) of that Lemma. What is born below $n_0$ is born at a small coupling; above $n_0$ nothing is born. (b) By Lemma~\ref{stc:reg}\,(a) the configuration $\mathsf r^{-1}U$ is regular of level $n$ at $\varepsilon_n$, with $\varepsilon_n^2\le CK^{2\mathfrak p_0}L^{-y}$. The main terms of Lemma~\ref{stc:reg}\,(c) add up at each site to $h^4$ times the symbol of $(k^0,v^0)$ at $(F_m(x))_m$, which vanishes: the symbol of the null pair $(k^0,v^0)$ is zero at every triple $(F_{12},F_{13},F_{23})\in\mathfrak g^3$ (Lemma~\ref{stc:defect}\,(a), Definition~\ref{stc:ins}\,(2)), whatever the curvatures $F_m(x)$ of Lemma~\ref{stc:reg}\,(c) are; what is left is at most $L^K(\ell L^K)^3(|k^0|_1+3|v^0|_\infty)\,C\varepsilon_n^2h^{4+\kappa_{\rm H}}$, $h=L^{-n}$, and $4K-y-(4+\kappa_{\rm H})(K-y)=(3+\kappa_{\rm H})y-\kappa_{\rm H}K$.
\end{proof}

\begin{lemma}[Births of the cubic half]\label{stc:births}
Let (Th) hold on $T_\ell$, at zero source, let $0\le k<t_*$, and let $\varkappa$ be a complex weight of the plaquettes.
\begin{enumerate}
\item[(o)] At the height $k$ there is a family $\mathbb V^\sharp_\varkappa$ in the fourth place, the \emph{cubic potential}: for $k=0$ the family $\mathbb W$ of the insertion $\mathcal S_\varkappa$ of the proof of Proposition~\ref{stc:chain}, for $k\ge1$ the sum over the polymers $X_1$ of level $k$ of the potentials $\mathbb V[f_{X_1}]$ (equation~\YMMref{(870)}{410}) of the local functions $f_{X_1}=\sum_{p:X_k(Q(p))=X_1}\varkappa(p)\mathfrak r^\sharp_{\rho(p)}(p)$. It is linear in $\varkappa$, and its member at $Y$ depends on $\varkappa$ on the plaquettes of $Y$ alone. It is formed and holomorphic on the sets $\mathfrak D^{\rm src}(Y)$ of \eqref{eq:source-entry-domain}, on which the fluctuation field $\mathcal A=g_kB$ of the step is bounded, bond by bond, by the radius $\varepsilon_{\rm src}$ of \eqref{eq:source-entry-domain}, a number that does not depend on $x_k$ ($\varepsilon_{\rm src}=r^\circ_P$ at $k=0$); it vanishes at $\mathcal A=0$; and there
\[
|\mathbb V^\sharp_\varkappa(Y;P,\mathcal A)|\le CL^{-3k}\max|\varkappa|\,e^{-\lambda_Vd_k(Y)},
\]
$\lambda_V$ the rate of the size $\mathsf n^{\rm f}$ of the families in the fourth place \eqref{eq:source-entry-norm}. Its size is at most $C\delta(x_k)L^{-3k}\max|\varkappa|$, and its sum is the variation of $\mathcal S_\varkappa$ from the zero-fluctuation background to the configuration of the chart integral. For $k\ge1$ this is $\mathcal S_\varkappa(U_k(V_W[\mathcal A]))-\mathcal S_\varkappa(U_{k+1}(W))$ in the chart over $W$; for $k=0$ it is the shifted value minus the background value in the definition of $\mathbb W$.
\end{enumerate}
Let $E_\mu$ be the output of the step of height $k$ with the family $\mathbb V'_k+\mu\mathbb V^\sharp_\varkappa$ in the fourth place, $\mathbb V'_k$ that of the chain and $|\mu|\,\mathsf n^{\rm f}(\mathbb V^\sharp_\varkappa)\le\frac1{32}w_\star$, and let $\mathsf D^{(k+1)}[\varkappa]=\partial_\mu E_\mu|_{\mu=0}$, value by value (the \emph{births}). Put $\mathsf s_0=\mathcal S_\varkappa$, and $\mathsf s_j=\mathcal S_\varkappa\circ U_j$ on $\mathsf{Cut}_j$, $0$ off it.
\begin{enumerate}
\item[(a)] $\mathsf D^{(k+1)}[\varkappa]$ is a pointed family of level $k+1$, complex linear in $\varkappa$, with $\|\mathsf D^{(k+1)}[\varkappa]\|\le K^{\rm d}\mathsf n^{\rm f}(\mathbb V^\sharp_\varkappa)$.
\item[(b)] Its member at $(X,z)$ depends on $\varkappa$ on the plaquettes of $X$ alone, and vanishes if $\varkappa=0$ there.
\item[(c)] For $k+1<t_*$: $\mathsf Q_k\mathsf s_k=\mathsf s_{k+1}+\Sigma_{k+1}[\mathsf D^{(k+1)}[\varkappa]]$ at every point of $\mathcal V_{k+1}$.
\item[(d)] $\sup_{\mathsf{Cut}_n}|\mathsf s_n|\le3c_{\rm orb}\,\ell^3L^{3K-6n}(\mathsf e^{\rm c}_n)^3\max|\varkappa|$ for $1\le n\le t_*$.
\item[(e)] (Parity.) There are $y_{\rm par}$ and $\mathsf C'_{\rm par}$, functions of the constants, such that for $K-k\ge y_{\rm par}$
\[
\|\mathsf D^{(k+1)}[\varkappa]\|\le\bar{\mathsf p}_{k+1}\max|\varkappa|,\qquad\bar{\mathsf p}_{k+1}=\mathsf C'_{\rm par}(\log x_k)^{\mathfrak p_0}\,g_k^2\,L^{-3k},\quad g_k^2=x_k^{-1}.
\]
\item[(f)] (Covariance.) Let $k+1<t_*$, $\mathsf r:x\mapsto{\rm diag}(s_1,s_2,s_3)x+v$ an isometry of the block lattice of level $k+1$, and $(\mathsf r^\sharp\varkappa)(p)=s_{\rho(p)}\varkappa(\mathsf rp)$ (for an inversion, $-\varkappa(\mathsf rp)$). For $W\in\mathsf{Cut}_{k+1}$ and $U\in\mathcal C(\mathbf B_{k+1},W)$: $\mathsf rW\in\mathsf{Cut}_{k+1}$, $\mathsf rU\in\mathcal C(\mathbf B_{k+1},\mathsf rW)$, and $\mathsf{Sum}[\mathsf D^{(k+1)}[\varkappa]](\mathsf rU)=\mathsf{Sum}[\mathsf D^{(k+1)}[\mathsf r^\sharp\varkappa]](U)$.
\end{enumerate}
\end{lemma}

\begin{proof}[Proof sketch]
(o) For $k\ge1$ each $f_{X_1}$ is a local term of level $k$: by Lemma~\ref{stc:rsharp}\,(b) a summand is $\le C\alpha_0^3L^{-6k}|\varkappa(p)|$ on the domain of the local terms, and at most $24M^3L^{3k}$ plaquettes have one region $X_k(Q(p))$. The potential of a term on its region is formed and holomorphic on the sets $\mathfrak D(Y)$ of Section~\YMMref{51.8}{410}, which at the radius $\varepsilon=\varepsilon_{\rm src}$ are the sets $\mathfrak D^{\rm src}(Y)$ (\eqref{eq:source-entry-domain}), and it vanishes at $\mathcal A=0$, by Proposition~\YMMref{51.15}{410}; it sums to the variation of the term by equation~\YMMref{(872)}{411}, at the admitted chart and background points of that proposition, and it decays in the size of the polymer by equation~\YMMref{(871)}{410}, as $e^{-(\bar\kappa-1)\mathfrak m_k(X_1,Y)}$ in the number $\mathfrak m_k(X_1,Y)$ of the cubes of $Y$ outside $X_1$.

A region $X_1=X_k(Q(p))$ has linear size $0$, so $d_k(Y)\le\mathfrak m_k(X_1,Y)$ (\eqref{eq:analytic-domain-extension}); with $0<\lambda_V\le\bar\kappa-2$ (the rate condition following \eqref{eq:source-entry-domain}, at the places of the chain) the remaining decay $e^{-\mathfrak m_k(X_1,Y)}$ pays for the at most $8$ regions per cube of $Y$. Summing over the regions gives the family, linear in $\varkappa$, local, formed and holomorphic on $\mathfrak D^{\rm src}(Y)$, zero at $\mathcal A=0$, with $|\mathbb V^\sharp_\varkappa(Y;P,\mathcal A)|\le\mathsf h_\varkappa e^{-\lambda_Vd_k(Y)}$ there, $\mathsf h_\varkappa=CL^{-3k}\max|\varkappa|$. For $k=0$ the same holds for $\mathbb W$, at the radius $r^\circ_P$ and with $\mathsf h_\varkappa=C\max|\varkappa|$: by its definition in the proof of Proposition~\ref{stc:chain} it is a sum, over the at most $8NM^3$ boxes of one region, of differences $\mathsf j_\iota(e^{iB'}U)-\mathsf j_\iota(U)$, which vanish at $\mathcal A=0$, where $B'=0$; at fluctuation fields below the radius $r^\circ_P$ the bond variables of $e^{iB'}U$ and of $U$ have norms $<2$, so each value is at most the height of the density; and the member at $Y$ vanishes unless $d_0(Y)=0$.

The sizes follow by Schwarz's lemma. On the field polydisc of the step $|\mathcal A(b)|=g_k|B(b)|<\mathsf C_{\mathcal A}\delta(x_k)\le\frac12\varepsilon_{\rm src}$ (condition (Sr), the half-radius condition following \eqref{eq:source-entry-domain}). Fix $Y$, $P$ and $B\ne0$, and put $\beta=\sup_b|B(b)|$: the function $\zeta\mapsto\mathbb V^\sharp_\varkappa(Y;P,\zeta g_kB)$ is holomorphic on the disc $|\zeta|<\varepsilon_{\rm src}/(g_k\beta)$, of radius $>2$, bounded there by $\mathsf h_\varkappa e^{-\lambda_Vd_k(Y)}$ and zero at $0$, and its even part is even. Schwarz's lemma at $\zeta=1$ bounds the family and its part odd under $B\mapsto-B$ by $\mathsf h_\varkappa e^{-\lambda_Vd_k(Y)}g_k\beta/\varepsilon_{\rm src}$, and its even part by $\mathsf h_\varkappa e^{-\lambda_Vd_k(Y)}(g_k\beta/\varepsilon_{\rm src})^2$.

With $g_k\beta<\mathsf C_{\mathcal A}\delta(x_k)$, $\beta\le\frac12M^{3/2}(1+\mathsf q)$ and $\beta^2\le M^3(1+\mathsf q)$, $\mathsf q=M^{-3}\sum_b|B(b)|^2$ the weight in the size \eqref{eq:source-entry-norm}, the sizes are at most $\mathsf C_{\mathcal A}\varepsilon_{\rm src}^{-1}\delta(x_k)\mathsf h_\varkappa$ for the family, $\frac12M^{3/2}\varepsilon_{\rm src}^{-1}g_k\mathsf h_\varkappa$ for its odd part and $M^3\varepsilon_{\rm src}^{-2}g_k^2\mathsf h_\varkappa$ for its even part. (a)--(c) The step of height $k$ with the further family $\mu\mathbb V^\sharp_\varkappa$, which sums to $\mu$ times the variation of $\mathcal S_\varkappa$ from its background value, gives, at every $W\in\mathsf{Cut}_{k+1}$ and $U\in\mathcal C(\mathbf B_{k+1},W)$, the chart integral of $F_ke^{\mu\mathsf s_k}$ as $\mathsf I[F_k](W)\exp[\sum_{X,z}(E_\mu-E_0)(X,U,z)+\mu\mathcal S_\varkappa(U)]$ (Proposition~\ref{prop:source-entry}); differentiating at $\mu=0$ gives $\mathsf I[\mathsf s_kF_k](W)=\mathsf I[F_k](W)\{\mathsf{Sum}[\mathsf D^{(k+1)}[\varkappa]](U)+\mathcal S_\varkappa(U)\}$, which is (c).

The derivative of a family that is holomorphic in $\mu$ is a pointed family, with the Lipschitz bound of the step in its family in the fourth place; linearity follows over a bidisc and by the chain rule; locality from the locality clauses of the two Theorems. (d) On $\mathsf{Cut}_n$ the plaquette variables are within $\mathsf e^{\rm c}_nL^{-2n}\le1$ of $\mathbf 1$; Lemma~\ref{stc:rsharp}\,(b) at $\mu=\mathsf e^{\rm c}_nL^{-2n}$, $3\ell^3L^{3K}$ plaquettes, the fixed factor $3c_{\rm orb}$. (e) Split $\mathbb V^\sharp_\varkappa=\mathbb V^{\rm ev}+\mathbb V^{\rm od}$ into its parts even and odd under the reflection $B\mapsto-B$ of the fluctuation field; both are families in the fourth place, and $\mathsf D^{(k+1)}=\mathsf D^{\rm ev}+\mathsf D^{\rm od}$, the derivatives of the output in the two directions.

The even part vanishes to the second order at $B=0$: by Schwarz's lemma on the disc of radius $\varepsilon_{\rm src}/(g_k\sup|B|)>2$, as in (o), $\varepsilon_{\rm src}$ the radius of the sources \eqref{eq:source-entry-domain}, its size is at most $CM^3\varepsilon_{\rm src}^{-2}g_k^2L^{-3k}\max|\varkappa|$, and by (a) the norm of $\mathsf D^{\rm ev}$ is at most $K^{\rm d}$ times this. The odd part has, by (o), the size $\mathsf n^{\rm f}(\mathbb V^{\rm od})\le CM^{3/2}\varepsilon_{\rm src}^{-1}g_kL^{-3k}\max|\varkappa|$ only, and its contribution is bounded by a parity argument in two parameters, in which the odd parts of the step's own families enter; these are small: the base family $\mathbb V'_k$ in the fourth place has even and odd parts of size at most a number $\nu_k$ of the chain, and the three interaction families have even and odd parts of size at most $K_P\delta(x_k)$ (\eqref{eq:source-internal-size}); both numbers are proportional to $\delta(x_k)$ and at most $\frac18w_\star$.

Put
\[
\rho=\min\Bigl\{\frac{w_\star}{8K_P\delta(x_k)}-1,\ \frac{w_\star}{64\nu_k}\Bigr\},\qquad\varpi_0=\frac{w_\star}{64\rho\,\mathsf n^{\rm f}(\mathbb V^{\rm od})},
\]
and for $|\mu|<\rho$, $|\varpi|<\varpi_0$ give the step, in the fourth place, the family $(\mathbb V'_k)^{\rm ev}+\mu(\mathbb V'_k)^{\rm od}+\varpi\mu\mathbb V^{\rm od}$, and, for each of the three interaction families, its even part plus $\mu$ times its odd part. These data are admissible: the interaction families have sizes at most $(1+\rho)K_P\delta(x_k)\le\frac18w_\star$, and the family in the fourth place has size at most $(1+\rho)\nu_k+\varpi_0\rho\,\mathsf n^{\rm f}(\mathbb V^{\rm od})\le\frac14w_\star$. Keep the Gaussian structure, cutoff and quadratic data fixed, and set
\[
 \Phi(\mu,\varpi)=\sum_b h_z(b)
       E^\bullet(X,b;\mathfrak E_{\mu,\varpi}),
\]
the interaction-dependent pointed term at $(X,z)$. The normalizer and
the fixed cutoff term are independent of $\mu,\varpi$, so their
$\varpi$ derivatives vanish. The change of variables $B\mapsto-B$
preserves the Gaussian, the cutoff, the field polydiscs and the quadratic
pairing, and carries all four families at $\mu$ to those at $-\mu$.
The orthogonal covariance in Proposition~\YMMref{48.4}{372} and Lemma~\YMMref{77.1}{631} therefore
gives $\Phi(-\mu,\varpi)=\Phi(\mu,\varpi)$. Those same localized-term
estimates give holomorphy on the parameter bidisc and
\[
 |\Phi(\mu,\varpi)|\le
 3L^3K_\Pi(32\nu_{\rm int}+9\nu')
       e^{-(\kappa+1)d_{k+1}(X)}
 \le Cw_\star e^{-(\kappa+1)d_{k+1}(X)},
\]
where $\nu_{\rm int}\le w_\star/8$ bounds each internal family and
$\nu'\le w_\star/4$ bounds the fourth family. The tent weights are
fixed throughout.

Hence $\psi(\mu)=\partial_\varpi\Phi(\mu,0)$ is even, and Cauchy's estimate gives $|\psi|\le Cw_\star e^{-(\kappa+1)d_{k+1}(X)}/\varpi_0=C\rho\,\mathsf n^{\rm f}(\mathbb V^{\rm od})e^{-(\kappa+1)d_{k+1}(X)}$ on $|\mu|<\rho$. At $\mu=0$ none of the four families reads $\varpi$, so $\psi(0)=0$, and $\psi$, being even, vanishes to the second order at $0$; the maximum principle for $\psi(\mu)/\mu^2$ on $|\mu|<\rho$ gives $|\psi(1)|\le C\rho^{-1}\mathsf n^{\rm f}(\mathbb V^{\rm od})e^{-(\kappa+1)d_{k+1}(X)}$. At $\mu=1$ the data are those of the step with $\mathbb V'_k+\varpi\mathbb V^{\rm od}$ in the fourth place, so the member of $\mathsf D^{\rm od}$ at $X$ is $\psi(1)$. For $K-k\ge y_{\rm par}$, $\rho>2$ and $\rho^{-1}\le\mathsf c'_\rho\delta(x_k)=\mathsf c'_\rho A_1g_k(\log x_k)^{\mathfrak p_0}$, $\mathsf c'_\rho$ a function of the constants; so the norm of $\mathsf D^{\rm od}$ is at most $CM^{3/2}\varepsilon_{\rm src}^{-1}(\log x_k)^{\mathfrak p_0}g_k^2L^{-3k}\max|\varkappa|$, the factor $\rho^{-1}$ supplying the second power of $g_k$.

Together: $\bar{\mathsf p}_{k+1}$; the factor $g_k^2$, not $g_k$, is what makes the label disappear in $g_k^2\mathsf b'=L^{-(K-k)}$. (f) Let $0\le i\le k$. An isometry $\mathsf r$ of the block lattice of level $k+1$ is one of every block lattice of a level $i+1\le k+1$ (Lemma~\ref{lem:setup-lattice-covariance}). The chart of the step of height $i$ is covariant under $\mathsf r$ (the chart calculation in the proof of Proposition~\YMMref{52.4}{435}): the push-forward operator $\mathsf O=\mathsf O_{\mathsf r,W}$ of the isometry (defined on the positive nonpivot bonds at the beginning of that proof) is a linear bijection of the fluctuation fields that keeps $dB$, the quadratic form $(\cdot,\cdot)_1$ and the range $\mathsf{Rg}_{x_i}$; $V^{\rm bf}_i(\mathsf rW)=\mathsf rV^{\rm bf}_i(W)$ and $V_{\mathsf rW}(\mathsf OB)=\mathsf rV_W(B)$; the exponent and the Jacobian are unchanged, $G_{\mathsf rW}(\mathsf OB)=G_W(B)$, $\mathsf J_{\mathsf rW}(\mathsf OB)=\mathsf J_W(B)$; and the elimination factor $z^{(i)}_{V^{\rm bf}_i(W)}(\Gamma)$ of the chart integral \eqref{eq:sf-chart-integral} is unchanged, as the Gaussian normaliser $Z^{(i)}(W;\mathsf A^\circ)$ of Section~\ref{sec:sf-chart-interface}, which is this factor times a Gaussian integral that $B\mapsto\mathsf OB$ carries into itself, is unchanged.

The substitution $B\mapsto\mathsf OB$ in \eqref{eq:sf-chart-integral} gives
\[
\mathsf I[F](\mathsf rW)=\mathsf I[F\circ\mathsf r](W)
\]
for bounded Borel $F$ on $\mathcal V_i$ and $W\in\mathsf{Cut}_{i+1}$; and $\mathsf rW\in\mathsf{Cut}_{i+1}$, the minimal orbits being carried along (Lemma~\ref{lem:sf-class-window}(3); Lemma~\ref{stc:reg}\,(a)). The integrands of the chain are invariant, by induction on $i$: $F_0=\rho_0$ on $\mathsf{Cut}_0$ (the definition preceding Proposition~\YMMref{114.2}{801}), and the lattice isometries keep $\rho_0$ and $\mathsf{Cut}_0$; if $F_i\circ\mathsf r=F_i$ on $\mathsf{Cut}_i$, then for $W\in\mathsf{Cut}_{i+1}$ the points $V_W(B)$ and $\mathsf rV_W(B)=V_{\mathsf rW}(\mathsf OB)$ lie in $\mathsf{Cut}_i$ (the downward stability \YMMref{113.10}{796}), and the step identity $\mathsf I[F_i]=e^{\mathsf c_{i+1}}F_{i+1}$ on $\mathsf{Cut}_{i+1}$, with a number $\mathsf c_{i+1}$ (Proposition~\YMMref{114.2}{801}, equation~\YMMref{(2174)}{801}), gives $e^{\mathsf c_{i+1}}F_{i+1}(\mathsf rW)=\mathsf I[F_i\circ\mathsf r](W)=\mathsf I[F_i](W)=e^{\mathsf c_{i+1}}F_{i+1}(W)$. So, for $W\in\mathsf{Cut}_{k+1}$,
\[
(\mathsf Q_kh)(\mathsf rW)=\frac{\mathsf I[(h\circ\mathsf r)(F_k\circ\mathsf r)](W)}{\mathsf I[F_k\circ\mathsf r](W)}=\frac{\mathsf I[(h\circ\mathsf r)F_k](W)}{\mathsf I[F_k](W)}=(\mathsf Q_k(h\circ\mathsf r))(W).
\]
By Lemma~\ref{stc:rsharp}\,(c), $\mathcal S_\varkappa(\mathsf rU)=\mathcal S_{\mathsf r^\sharp\varkappa}(U)$, and the minimal orbits are carried along, so $\mathsf s^\varkappa_n\circ\mathsf r=\mathsf s^{\mathsf r^\sharp\varkappa}_n$ on $\mathsf{Cut}_n$ ($n=k+1$, and $n=k$ at the points $V_W(B)$); inserted in (c) at $\mathsf rW$ and at $W$ this gives $\mathsf{Sum}[\mathsf D^{(k+1)}[\varkappa]](\mathsf rU)=(\mathsf Q_k\mathsf s^\varkappa_k)(\mathsf rW)-\mathsf s^\varkappa_{k+1}(\mathsf rW)=(\mathsf Q_k\mathsf s^{\mathsf r^\sharp\varkappa}_k)(W)-\mathsf s^{\mathsf r^\sharp\varkappa}_{k+1}(W)=\mathsf{Sum}[\mathsf D^{(k+1)}[\mathsf r^\sharp\varkappa]](U)$.
\end{proof}

\begin{lemma}[Carrying families up; stopping]\label{stc:stop}
Let $t\ge1$ and $L\ge2$. For $1\le j\le t$ let $\mathfrak F_j$ be complex vector spaces with seminorms $\|\cdot\|_j$, $\mathsf c(\mathsf f)=L^j\|\mathsf f\|_j$ for $\mathsf f\in\mathfrak F_j$, $\mathfrak F_{\le n}=\bigoplus_{j\le n}\mathfrak F_j$; for $1\le n\le t$ let $\mathscr B_n$ be complex vector spaces with seminorms $|\cdot|_n$ and elements $\mathbf 1_n$; let $\Sigma_n:\mathfrak F_{\le n}\to\mathscr B_n$ be linear, $a$ a linear form on $\mathfrak F_{\le t}$, $S_i:\mathfrak F_{\le i}\to\mathfrak F_{i+1}$ and $\mathsf Q_i:\mathscr B_i\to\mathscr B_{i+1}$ linear ($1\le i<t$), and $\theta_i,\mathsf B_n\ge0$, with
\begin{enumerate}
\item[(N)] $\|S_i\mathsf f\|_{i+1}\le\theta_iL^{-(i-i')}\|\mathsf f\|_{i'}$ for $\mathsf f\in\mathfrak F_{i'}$, $i'\le i$ (so $\mathsf c(S_i\mathsf f)\le L\theta_i\mathsf c(\mathsf f)$);
\item[(A)] $\mathsf Q_i\mathbf 1_i=\mathbf 1_{i+1}$, $|\mathsf Q_ih|_{i+1}\le|h|_i$, $\mathsf Q_i\Sigma_i[\mathsf g]=\Sigma_{i+1}[\mathsf g+S_i\mathsf g]$;
\item[(B)] $|\Sigma_n[\mathsf f]-a[\mathsf f]\mathbf 1_n|_n\le\mathsf B_n\mathsf c(\mathsf f)$ for $\mathsf f\in\mathfrak F_j$, $j\le n$.
\end{enumerate}
Put $\mathsf T_t={\rm id}$ and $\mathsf T_n\mathsf g=\mathsf T_{n+1}(\mathsf g+S_n\mathsf g)$.
\begin{enumerate}
\item[(a)] $\Sigma_t[\mathsf T_n\mathsf g]=\mathsf Q_{t-1}\cdots\mathsf Q_n\Sigma_n[\mathsf g]$; $|\Sigma_t[\mathsf T_n\mathsf g]-a_0\mathbf 1_t|_t\le|\Sigma_n[\mathsf g]-a_0\mathbf 1_n|_n$ for every $a_0\in\mathbb C$; and $\mathsf T_n\mathsf g=\mathsf T_{n'}\mathsf g+\sum_{k=n}^{n'-1}\mathsf T_{k+1}S_k\mathsf g$ for $n\le n'\le t$.
\item[(b)] Let $1\le j\le n_0\le t$ and $\mathsf D\in\mathfrak F_j$. A \emph{chain} is $\underline k=(k_1<\dots<k_m)$, $m\ge1$, $j\le k_1$, $k_m<t$, and $\mathsf f_{\underline k}=S_{k_m}\cdots S_{k_1}\mathsf D$. Let $\mathsf n$ map the chains to integers with $k_m<\mathsf n(\underline k)\le t$; a chain is \emph{live} if $k_1<n_0$ and $k_{l+1}<\mathsf n(k_1,\dots,k_l)$ for $1\le l<m$. Then $\mathsf T_j\mathsf D=\mathsf T_{n_0}\mathsf D+\sum_{\rm live}\mathsf T_{\mathsf n(\underline k)}\mathsf f_{\underline k}$, and for every $a_0\in\mathbb C$
\[
\Bigl|\Sigma_t[\mathsf T_j\mathsf D]-\Bigl(a_0+\sum_{\rm live}a[\mathsf f_{\underline k}]\Bigr)\mathbf 1_t\Bigr|_t\le|\Sigma_{n_0}[\mathsf D]-a_0\mathbf 1_{n_0}|_{n_0}+\mathsf c(\mathsf D)\sum_{\underline k\ {\rm live}}\mathsf B_{\mathsf n(\underline k)}\prod_lL\theta_{k_l}.
\]
\item[(c)] Let $K'\ge t$, $q\ge1$, $0<\tau<\frac12$, $\Theta_0\ge1$, $\mathsf B_0>0$ with $L\theta_k\le\Theta_0L^{-(K'-k)/2}$ ($1\le k<t$) and $\mathsf B_n\le\mathsf B_0L^{q(K'-n)-K'}$ ($1\le n\le t$); let $u=K'-t$, $y_*=K'-n_0$, $m_{q,\tau}=2+\lfloor\log(q/\tau)/\log(1+\tau/q)\rfloor$, $Y(\underline k)=\sum_l(K'-k_l)$ and $\mathsf n(\underline k)=\max\{k_m+1,\min\{t,\lceil K'-\tau Y(\underline k)/q\rceil\}\}$. Then every live chain has at most $m_{q,\tau}$ members, and
\[
\sum_{\underline k\ {\rm live}}\mathsf B_{\mathsf n(\underline k)}\prod_lL\theta_{k_l}\le\mathsf B_0L^{-K'}m_{q,\tau}(t\Theta_0)^{m_{q,\tau}}\max\bigl\{L^{-(\frac12-\tau)y_*},L^{qu-y_*/2}\bigr\}.
\]
\end{enumerate}
\end{lemma}

\begin{proof}[Proof sketch]
(a) Downward induction on $n$ with the third member of (A); $a_0\mathbf 1_t=\mathsf Q_{t-1}\cdots\mathsf Q_n(a_0\mathbf 1_n)$ and the second member of (A); $\mathsf T_n\mathsf g=\mathsf T_{n+1}\mathsf g+\mathsf T_{n+1}S_n\mathsf g$ and induction on $n'-n$. (b) (a) at $(j,n_0)$ gives $\mathsf T_{n_0}\mathsf D$ and the live chains of length $1$; for a live chain of length $m$, (a) at $(k_m+1,\mathsf n(\underline k))$ splits $\mathsf T_{k_m+1}\mathsf f_{\underline k}$ into $\mathsf T_{\mathsf n(\underline k)}\mathsf f_{\underline k}$ and the live chains of length $m+1$ that prolong it; a chain has at most $t-1$ members. The bound: (a) for each summand, (B) at the level $\mathsf n(\underline k)$, and (N). (c) The price of a chain is $\le\mathsf B_0L^{-K'}\Theta_0^mL^{q(K'-\mathsf n)-Y/2}$, $\mathsf n=\mathsf n(\underline k)$, a bound that does not increase with $\mathsf n$.

Put $\mathsf n^\circ=\min\{t,\lceil K'-\tau Y/q\rceil\}$, so that $\mathsf n=\max\{k_m+1,\mathsf n^\circ\}\ge\mathsf n^\circ$ and $\mathsf n\le t$. If $\mathsf n^\circ$ is the ceiling, then $K'-\mathsf n\le K'-\mathsf n^\circ\le\tau Y/q$ and the exponent is $\le-(\frac12-\tau)Y$; otherwise $\mathsf n^\circ=t$, so $\mathsf n=t$, and the exponent is $qu-Y/2$. (When $\mathsf n=k_m+1>\mathsf n^\circ$, then $\mathsf n^\circ<t$ is the ceiling, and the first case applies.) For a live chain $Y>y_*$. If $(k_1,\dots,k_{m+1})$ is live, then $k_{m+1}<K'-\tau Y_m/q$, so $y_{m+1}>\tau Y_m/q$ with $y_l=K'-k_l$; hence $Y_{l+1}>(1+\tau/q)Y_l$, and $y_{m+1}<y_1$ gives $(1+\tau/q)^{m-1}y_1\le Y_m<(q/\tau)y_1$. There are at most $t^m$ chains of length $m$.
\end{proof}

\begin{proposition}[The conditional mean of the cubic half]\label{stc:cubic}
Let $\ell=2L^m$, $\gamma^{-2}\le\mathsf b_-\le\mathsf b_+$ with $MR(\mathsf b_+)\le L^m$, and let parameters $0<\psi<\frac29$, $0\le\kappa_\star<1$, $0<\tau<\frac12$, $0<\phi<\frac16\psi$ be given; put
\[
c''=\min\bigl\{1-\tfrac92\psi,\ \kappa_\star-(3+\kappa_\star)\psi,\ (\tfrac12-\tau)\psi,\ \tfrac12\psi-3\phi\bigr\},\qquad\bar m=2+\lfloor\log(3/\tau)/\log(1+\tau/3)\rfloor .
\]
There are $K_3$, a function of $\mathsf b_+$ and the parameters, and $C_3,p_3$, functions of $m$, $\mathsf b_+$ and the parameters, such that for all $K\ge K_3$, $\mathsf b'\in[\mathsf b_-,\mathsf b_+]$, $\beta=L^K\mathsf b'$ and $t=K-\lceil\phi K\rceil$:
\begin{enumerate}
\item[(a)] (Th) holds on $T_\ell$ with $(K,\mathsf b',m,t)$, so that Proposition~\ref{stc:chain} and Lemma~\ref{stc:cmean} hold with $t_*=t$;
\item[(b)] for every real field $\eta=(\eta^\rho)_\rho$ on $T_\ell$, with $\mathsf H(\eta)=\max_\rho\sup|\eta^\rho|+\max_\rho{\rm Lip}\,\eta^\rho$, there is a number $\mathsf a[\eta]$, real linear in $\eta$ and a sum of values at the trivial configuration, with
\[
|(\mathsf Q_{t\leftarrow0}\mathcal S[\mathsf b'\eta])(V)-\mathsf a[\eta]|\le C_3\,\mathsf H(\eta)\,K^{p_3}L^{-c''K}\qquad\text{for $\nu_t$-almost all }V\in\mathsf{Cut}_t .
\]
\end{enumerate}
At $\psi=\frac17$, $\kappa_\star=\frac34$, $0<\tau\le\frac14$, $0<\phi<\frac1{42}$:
\[
c''=c_0':=\min\bigl\{\tfrac1{14}-\tfrac\tau7,\ \tfrac1{14}-3\phi\bigr\},\qquad p_3=3\mathfrak p_0+1+\bar m(1+\mathfrak p_0),
\]
$K_3$ a function of $\mathsf b_+$ and $\phi$, $C_3$ a function of $m$, $\mathsf b_+$ and $\tau$; $c_0'\ge\frac1{28}$ for $\tau=\frac14$, $\phi\le\frac1{84}$, and $c_0'$ does not fall as $\phi\to0$. At $\psi=\frac15$, $\kappa_\star=\frac9{10}$: $c''=\frac1{10}-\max\{\frac15\tau,3\phi\}$.
\end{proposition}

\begin{proof}[Proof sketch]
Put $\varkappa(p)=L^{2K}\mathsf b'\eta^{\rho(p)}(x(p))$, so that $\mathcal S[\mathsf b'\eta]=\mathcal S_\varkappa$ is an insertion of order $3$ (Lemma~\ref{stc:rsharp}\,(b)); power counting alone (Lemma~\ref{stc:cmean}\,(c)) gives it only $L^{4u+K/2}$, and the gain comes from parity and covariance. Let $u=\lceil\phi K\rceil$, $t=K-u$, $y_*=\lceil\psi K\rceil$, $n_0=K-y_*$. $K_3$ is so large that: Lemma~\ref{stc:frame}\,(b) applies; $1\le n_0<t$; $\psi K\ge y_{\rm par}$; $K\ge K_2$ (Lemma~\ref{stc:reg}\,(b)); and the data of Lemmas~\ref{stc:reg}\,(c) and \ref{stc:mirror} hold at $\varepsilon_{n_0}$, with $D_*\ge7K\log L+1$. Each of these holds from a $K$ on, a power of $L^K$ against a power of $K$; this is (a).

\emph{1. The births.} By Lemma~\ref{stc:births}\,(c), $\mathsf Q_{j-1}\mathsf s_{j-1}=\mathsf s_j+\Sigma_j[\mathsf D^{(j)}]$, $\mathsf D^{(j)}=\mathsf D^{(j)}[\varkappa]$, for $1\le j\le n_0$; by induction and linearity
\[
\mathsf Q_{t\leftarrow0}\mathcal S_\varkappa=\mathsf Q_{t\leftarrow n_0}\mathsf s_{n_0}+\sum_{j=1}^{n_0}\mathsf Q_{t\leftarrow j}\Sigma_j[\mathsf D^{(j)}].
\]
\emph{2. The scheme.} The pointed families of level $j$ with their norms, the bounded functions on $\mathcal V_n$ vanishing off $\mathsf{Cut}_n$ with the essential supremum for $\nu_n$, $\mathbf 1_n=\mathbf 1_{\mathsf{Cut}_n}$, the maps $\Sigma_n$, the value $a[\mathsf f]=\mathsf{Sum}[\mathsf f](\mathbf 1)$ at the trivial configuration, the maps $\mathsf S_i$ and $\mathsf Q_i$ of Lemma~\ref{stc:cmean}\,(a), the numbers $\theta_i$ and $\mathsf B_n$ of Lemma~\ref{stc:reg}\,(b) satisfy (N), (A), (B) of Lemma~\ref{stc:stop}, and its hypotheses (c) hold with $K'=K$, $q=3$, $\Theta_0=\Theta_K$ (Lemma~\ref{stc:frame}\,(b)) and $\mathsf B_0=\mathsf B'_0$; $m_{3,\tau}=\bar m$. So $\mathsf Q_{t\leftarrow j}\Sigma_j[\mathsf D^{(j)}]=\Sigma_t[\mathsf T_j\mathsf D^{(j)}]$.

\emph{3. The entry.} A conditional mean does not exceed a supremum: 
\[
|\mathsf Q_{t\leftarrow n_0}\mathsf s_{n_0}|\le3c_{\rm orb}\ell^3L^{3K-6n_0}(\mathsf e^{\rm c}_{n_0})^3\max|\varkappa|
\]
 (Lemma~\ref{stc:births}\,(d)), with $(\mathsf e^{\rm c}_{n_0})^3\mathsf b'\le CK^{3\mathfrak p_0}L^{-\frac32y_*}$ and $\max|\varkappa|\le L^{2K}\mathsf b'\mathsf H(\eta)$: the exponent of $L$ is $\frac92y_*-K$.

\emph{4. The norms.} By Lemma~\ref{stc:births}\,(e), with $g_k^2\mathsf b'=L^{-(K-k)}$: $\mathsf c_j:=L^j\|\mathsf D^{(j)}\|\le CK^{\mathfrak p_0}L^{K-j}\mathsf H(\eta)$; the label has disappeared.

\emph{5. Covariance at the levels $j<n_0$.} With the partition of unity $\theta_z$ of Lemma~\ref{stc:mirror}\,(b) and $\Sigma_z[\varkappa]=\mathsf{Sum}[\mathsf D^{(j)}[\theta_z\varkappa]]$, linearity gives $\mathsf{Sum}[\mathsf D^{(j)}[\varkappa]]=\sum_z\Sigma_z[\varkappa]$, and Lemma~\ref{stc:births}\,(f) for the inversion $\mathsf r_z$, with $\theta_z\circ\mathsf r_z=\theta_z$, gives $\Sigma_z[\varkappa](\mathsf r_zU)=-\Sigma_z[\varkappa\circ\mathsf r_z](U)$. So
\[
\Sigma_z[\varkappa](U)=\tfrac12\bigl(\Sigma_z[\varkappa](U)-\Sigma_z[\varkappa](\mathsf r_zU)\bigr)+\tfrac12\Sigma_z[\varkappa-\varkappa\circ\mathsf r_z](U).
\]
The first parts are bounded by Lemma~\ref{stc:mirror}\,(b): with $N_j\vartheta^4=N_{n_0}\vartheta$, $\varepsilon_{n_0}^2\le CK^{2\mathfrak p_0}L^{-y_*}$ and $\bar{\mathsf p}_jL^{2K}\mathsf b'\le CK^{\mathfrak p_0}L^{K-2j+2}$, their sum over $z$ is at most $CK^{3\mathfrak p_0}\mathsf H(\eta)L^{3y_*}L^{-j}\vartheta^{\kappa_\star}$, a gain $\vartheta^{\kappa_\star}$ over the size of one family at a background: the even part of the gauge potential of a regular configuration costs $\vartheta^2$, its odd part $\vartheta^{2+\kappa_\star}$, and a local term vanishes quadratically, so the difference is of order $\varepsilon_{n_0}^2\vartheta^{4+\kappa_\star}$. The second parts are $\frac12\mathsf{Sum}[\mathsf D^{(j)}[\tilde\varkappa_j]]$, $\tilde\varkappa_j=\sum_z\theta_z(\varkappa-\varkappa\circ\mathsf r_z)$, a weight with $\max|\tilde\varkappa_j|\le L^{2K}\mathsf b'\mathsf H(\eta)L^{j-K}$ (an inversion keeps $\rho(p)$ and moves $x(p)$ by at most $L^{j-K}$ where $\theta_z(p)\ne0$); by Lemma~\ref{stc:reg}\,(b) at $n=n_0$ they differ from their values at $\mathbf 1$ by $CK^{3\mathfrak p_0}\mathsf H(\eta)L^{3y_*}L^{-j}L^{-(K-j)}$. Together, for $1\le j<n_0$,
\[
\varsigma'_j:=\sup_{V\in\mathsf{Cut}_{n_0}}\bigl|\mathsf{Sum}[\mathsf D^{(j)}](U_{n_0}(V))-\tfrac12\mathsf{Sum}[\mathsf D^{(j)}[\tilde\varkappa_j]](\mathbf 1)\bigr|\le C\mathsf H(\eta)K^{3\mathfrak p_0}L^{3y_*}L^{-j-\kappa_\star(n_0-j)} .
\]

\emph{6. Stopping.} Lemma~\ref{stc:stop}\,(b) with $a_0=\frac12\mathsf{Sum}[\mathsf D^{(j)}[\tilde\varkappa_j]](\mathbf 1)$ for $j<n_0$ (where $|\Sigma_{n_0}[\mathsf D^{(j)}]-a_0\mathbf 1_{n_0}|\le\varsigma'_j$), and (B) at $n=j=n_0$, with the stopping rule of Lemma~\ref{stc:stop}\,(c); a chain of the maps $\mathsf S_k$ costs a factor $L\theta_k$ at each step, a chain is stopped at the level $\mathsf n(\underline k)$, where its value at the background is compared with its value at the trivial configuration at the price $\mathsf B_{\mathsf n}=\mathsf B'_0L^{3(K-\mathsf n)-K}$, and the live chains have at most $\bar m$ members. With
\[
\mathsf a[\eta]=\sum_{j<n_0}\Bigl(\tfrac12\mathsf{Sum}[\mathsf D^{(j)}[\tilde\varkappa_j]](\mathbf 1)+\sum_{\underline k\ {\rm live}}\mathsf{Sum}[\mathsf S_{k_m}\cdots\mathsf S_{k_1}\mathsf D^{(j)}](\mathbf 1)\Bigr)+\mathsf{Sum}[\mathsf D^{(n_0)}](\mathbf 1),
\]
real linear in $\eta$ (the live chains depend on $K,t,n_0,j,\tau$ alone), for $\nu_t$-almost all $V\in\mathsf{Cut}_t$
\[
\begin{multlined}
|(\mathsf Q_{t\leftarrow0}\mathcal S_\varkappa)(V)-\mathsf a[\eta]|\le\sup_{\mathsf{Cut}_{n_0}}|\mathsf s_{n_0}|+\sum_{j<n_0}\varsigma'_j+\mathsf B_{n_0}\mathsf c_{n_0} \\
+\sum_{j<n_0}\mathsf B'_0\mathsf c_jL^{-K}\bar m(t\Theta_K)^{\bar m}\max\bigl[L^{-(\frac12-\tau)y_*},L^{3u-y_*/2}\bigr].
\end{multlined}
\]

\emph{7. The exponents.} With $\psi K\le y_*\le\psi K+1$ and $u\le\phi K+1$, the four groups give the exponents $1-\frac92\psi$ (the entry), $1-4\psi$ (the births at the level $n_0$, $\mathsf B_{n_0}\mathsf c_{n_0}\le CK^pL^{4y_*-K}\mathsf H(\eta)$), $\kappa_\star-(3+\kappa_\star)\psi$ (covariance, as $j+\kappa_\star(n_0-j)\ge\kappa_\star n_0$) together with $1-3\psi$ (the part in ${\rm Lip}\,\eta$ and the tails), and $(\frac12-\tau)\psi$, $\frac12\psi-3\phi$ (the families carried up from the levels $j<n_0$ and stopped, as $\sum_jL^{-j}\le1$). As $1-4\psi$ and $1-3\psi$ exceed $1-\frac92\psi$, the least is $c''$; all powers of $K$ are at most $K^{p_3}$. At $\psi=\frac17$, $\kappa_\star=\frac34$ the six numbers are $\frac5{14}$, $\frac37$, $\frac3{14}$, $\frac47$, $\frac1{14}-\frac17\tau$, $\frac1{14}-3\phi$. The label $\mathsf b'$ has entered only through $x_i\in[\gamma^{-2}L^{K-i},\mathsf b_+L^{K-i}]$, $g_k^2\mathsf b'=L^{-(K-k)}$ and $\mathsf b'^{-1/2}\le\gamma$, so the bound is uniform in $\mathsf b'\in[\mathsf b_-,\mathsf b_+]$.
\end{proof}

\emph{Boards.} A board of $\mathbb L_K$ with cells of side $\mathsf l$ and offset $o$ has the wall planes $x_\mu=o_\mu+\mathsf l(t'+\frac12)$, $t'\in\mathbb Z$; on $T_{\ell'}$ it has eight cells, exchanged by the mirror maps $\mathsf r_\kappa$, $\kappa\in(\mathbb Z/2)^3$. A function of a cell is a bounded Borel function of the bonds with both end points in it, and its chessboard seminorm is given by $\|F\|^8_{{\rm ch},\ell'}=\langle\prod_\kappa F\circ\mathsf r_\kappa^{-1}\rangle_{K,\ell',\beta}$, after periodisation (Section~\ref{sec:periodic-chessboard}); it is subadditive and homogeneous (\ref{IV-18-18}).

\begin{lemma}[Boards for small sets]\label{stc:board}
\begin{enumerate}
\item[(a)] (Two sets.) Let $0<\delta_0\le\mathsf l$, $0<h\le\frac1{64}\delta_0$ with $\mathsf l\in h\mathbb Z$, and let $S_1,S_2\subset\mathbb R^3$ have sup diameters $\le\frac18\delta_0$, with $\sup_{S_1}x_{\mu_0}+\delta_0\le\inf_{S_2}x_{\mu_0}$ for one direction $\mu_0$. There is $o\in h\mathbb Z^3$ such that $S_1$, $S_2$ lie in two different cells of the board with offset $o$, each at sup distance $\ge\frac1{16}\delta_0$ from the walls of its cell.
\item[(b)] ($n$ sets.) Let $n\ge1$, $0<h<\mathsf l/(2n)$ with $\mathsf l\in h\mathbb Z$, and let $S_1,\dots,S_n\subset\mathbb R^3$ be non-empty sets of sup diameter $\le\mathsf l/(4n)$. There is $o\in h\mathbb Z^3$ such that every point of every $S_a$ is farther than $\mathsf l/(8n)$ from every wall plane of the board with offset $o$; so each $S_a$ lies in one cell. Different $S_a$ may lie in the same cell.
\end{enumerate}
\end{lemma}

\begin{proof}[Proof sketch]
(a) For a direction $\mu$ let $I_1$, $I_2$ be the closed $\frac1{16}\delta_0$-neighbourhoods of the convex hulls of the projections of $S_1$, $S_2$ (the projections themselves need not be intervals): intervals of length $\le\frac14\delta_0$. Wanted is $c_\mu\in\frac{\mathsf l}2+h\mathbb Z$ with $(c_\mu+\mathsf l\mathbb Z)\cap(I_1\cup I_2)=\emptyset$, and for $\mu=\mu_0$ besides $\sup I_1<c_{\mu_0}<\inf I_2$. The bad set $(I_1\cup I_2)+\mathsf l\mathbb Z$ has measure $\le\frac12\delta_0$ in a period, in at most two intervals. For $\mu_0$ the gap $(\sup I_1,\sup I_1+\frac78\delta_0)$ lies between $I_1$ and $I_2$; minus the bad set it has measure $\ge\frac38\delta_0$ in at most four intervals, one of length $\ge\frac3{32}\delta_0>h$, which holds a point of $\frac{\mathsf l}2+h\mathbb Z$.

For $\mu\ne\mu_0$ a period minus the bad set has measure $\ge\mathsf l-\frac12\delta_0\ge\frac12\delta_0$ in at most three intervals. A connected set that meets no wall plane lies in one cell, and a wall plane of the direction $\mu_0$ parts $S_1$ from $S_2$. (b) Fix a direction. The closed $\frac{\mathsf l}{8n}$-neighbourhoods of the convex hulls of the projections of the $S_a$ are intervals of length $\le\frac{\mathsf l}{2n}$; on the circle $\mathbb R/\mathsf l\mathbb Z$ they form $n$ arcs of total length $\le\frac{\mathsf l}2$. The complement has measure $\ge\frac{\mathsf l}2$ and at most $n$ components, so one of them is an open arc of length at least $\frac{\mathsf l}{2n}>h$, which contains a point of $\frac{\mathsf l}2+h\mathbb Z$; take it as a wall position.

A box that meets no wall plane lies in one cell.
\end{proof}

\begin{lemma}[Families, mirror copies, oscillations]\label{stc:fam}
Fix $c_*\in\{\frac1{28},\frac1{14}\}$ and $0<\phi\le\frac1{84}$, put $u=\lceil\phi K\rceil$, $t=K-u$, and let $(\tau,y_{\rm c},\kappa_{\rm H})=(\frac14,\lceil\frac K{14}\rceil,\frac12)$ if $c_*=\frac1{28}$, $(\tau,y_{\rm c},\kappa_{\rm H})=(21\phi,\lceil\frac K6\rceil,\frac34)$ if $c_*=\frac1{14}$. To each of the functions $\mathsf X=\mathcal T_K(H)$, $a_K\mathcal O(w_\eta)$, $\mathsf{Rm}[\eta]$, $\Delta[\eta]$, $\mathsf{St}[\eta]$, $\mathfrak K[\eta]$, $\mathcal T^{\rm c}[\eta]$ we attach the following family of at most nine insertions, whose sum differs from $\mathsf X$ by null terms and a constant (Lemmas~\ref{stc:null}, \ref{stc:defect}):
\begin{itemize}
\item $\mathcal T_K(H)$ itself (a rough insertion) for $\mathcal T_K(H)$, and $a_K\mathcal O(w_\eta)$ itself for $a_K\mathcal O(w_\eta)$; difference $0$;
\item $\mathcal S[b\eta]$, $-\mathfrak J_{0,v}[\mathsf k_\eta]$ ($v\le5$), $\frac32a_K\mathcal O(w_\eta)$ for $\mathsf{Rm}[\eta]$; difference $0$;
\item $F(k^0,v^0)$, $F(k^1,v^1)$ for $\Delta[\eta]$; difference $0$;
\item the negatives of the seven insertions of $\mathsf{Rm}[\eta]$ for $\mathsf{St}[\eta]$; difference $\mathcal N_\eta-\frac14\mathfrak N[\mathsf k_\eta]-\frac14c_{\mathsf k_\eta}$;
\item $\mathcal S[b\eta]$, $-\mathfrak J_{0,v}[\mathsf k_\eta]$ for $\mathfrak K[\eta]$; difference $\frac14\mathfrak N[\mathsf k_\eta]+\frac14c_{\mathsf k_\eta}$;
\item $-b^{-1}$ times the seven insertions of $\mathsf{Rm}[\eta]$ and the two of $\Delta[\eta]$ for $\mathcal T^{\rm c}[\eta]$; difference $b^{-1}(\mathcal N_\eta-\frac14\mathfrak N[\mathsf k_\eta]-\frac14c_{\mathsf k_\eta})$.
\end{itemize}
The support of $\mathsf X$ is that of $H$, resp.\ $\eta$; $\mathsf n(\mathsf X)=|H|_\infty$, resp.\ $\|\eta\|$; $\mathsf X$ is \emph{small} unless $\mathsf X=\mathcal T_K(H)$. Let $0<\delta\le\frac18\mathsf l$. There are $C_\omega,p_\omega$, functions of $g,\ell',c_*,\phi$, and $K_{\rm B}$, a function of these and $\delta$, such that for all $(K,\beta)\in\mathfrak d(g)$ with $K\ge K_{\rm B}$, every board with offset in $a_{K-t}\mathbb Z^3$, every such $\mathsf X$ with support in a cell at distance $\ge\delta$ from its walls, and every $\kappa$:
\begin{enumerate}
\item[(a)] $\mathsf r_\kappa$ maps the block lattices of all levels $\le t$ onto themselves. The mirror copy $\mathfrak J\circ\mathsf r_\kappa^{-1}$ of an insertion of type $(m,c,\hat c,N)$ is an insertion of order $m$ with constant $2^{6m}\max\{c,\hat c\}$, height $\hat c$, the same weights, carrier the image of the carrier, and counting number $\le20N$; for rough insertions the constants are unchanged; and $\mathcal S[b\eta]\circ\mathsf r_\kappa^{-1}=\mathcal S[b\eta^{\mathsf r}]$ with $(\eta^{\mathsf r})^{\rho'}(\mathsf r_\kappa y)=\pm\eta^\rho(y)$, $\mathsf H(\eta^{\mathsf r})=\mathsf H(\eta)$.
\item[(b)] $\mathsf X$ and the members of its family are real functions of the cell that depend only on the bond variables within $4a_K$ of the support, with weights $\le(3+60\mathsf b_+)L^{2K}\mathsf n(\mathsf X)$; the members are insertions whose orders, constants, heights and counting numbers are pure numbers.
\item[(c)] The oscillations of the copies of the family add up to $\sum_e\omega_t(\mathfrak J_e\circ\mathsf r_\kappa^{-1})\le\varpi(\mathsf X)\,\mathsf n(\mathsf X)$, with $\varpi=\varpi_{\rm T}=C_\omega K^{p_\omega}L^{4u}$ if $\mathsf X=\mathcal T_K(H)$ and $\varpi=\varpi_{\rm s}=C_\omega K^{p_\omega}L^{3u-c_*K}$ if $\mathsf X$ is small; and $L^{-K}\le\varpi_{\rm s}\le\varpi_{\rm T}\le L^K$.
\end{enumerate}
\end{lemma}

\begin{proof}[Proof sketch]
(a) Twice a wall coordinate lies in $a_{K-t}\mathbb Z$, as $o\in a_{K-t}\mathbb Z^3$ and $\mathsf l\in a_{K-t}\mathbb Z$, so the reflections preserve $a_{K-n}\mathbb Z^3$ for $n\le t$. The plaquette variables of the reflected configuration, from the standard corner and sense, are $g\,\partial U(\mathsf rq)^{\pm1}g^{-1}$ with $|g^{\pm1}|\le2^3$, hence within $2^6r$ of $\mathbf 1$ if those of $U$ are within $r$; the density is then $\le c(2^6r)^m$ if $2^6r\le1$, and $\le\hat c\le\hat c(2^6r)^m$ otherwise. A reflected box has the same shape; for a given lowest site there is one possible original lowest site for each of the $20$ admissible shapes, whence $20N$. For petal forms the bound $|q_x|\le2^{12}r^2$ holds for all petals at once.

The identity for the cubic half is Lemma~\ref{stc:rsharp}\,(c). (b) Lemmas~\ref{stc:null} and \ref{stc:defect}; the weights of the members are those listed there, and for $\mathcal T^{\rm c}$ they carry the factor $b^{-1}\le\gamma^2$; the types are those of Lemmas~\ref{stc:rsharp}\,(a), (b) and \ref{stc:defect}\,(b). (c) Let $\mathsf n(\mathsf X)=1$, so $\mathsf H(\eta)\le3$. For every copy: $\mathcal S[b\eta]$, by (a) and Proposition~\ref{stc:cubic} applied to $\eta^{\mathsf r}$ at $(\tau,\phi)$ and $\psi=\frac17$, $\kappa_\star=\frac34$: $\omega_t\le3C_3K^{p_3}L^{-c_0'K}$, and $c_0'\ge c_*-3\phi$; the four insertions of order $2$ of $\mathsf{Rm}$, of weight $\le3+5\mathsf b_+$, $F(k^1,v^1)$, of weight $\le60\mathsf b_+$, and $a_K\mathcal O(w_\eta)$: by Lemmas~\ref{stc:defect}\,(b) and \ref{stc:cmean}\,(b), (c), $\le CK^{2\mathfrak p_0}L^{4u-K}$; $\mathfrak J_{0,4},\mathfrak J_{0,5}$, of orders $3,4$ and weight $\le L^K\mathsf b_+$: $\le CK^{3\mathfrak p_0}L^{4u-K/2}$; $F(k^0,v^0)\circ\mathsf r_\kappa^{-1}$, of weight $\le22\mathsf b_+L^K$: Lemma~\ref{stc:cut}\,(a) at $n_0=K-y_{\rm c}\le t$ with (b) at $\kappa_{\rm H}$, $\le CK^{2\mathfrak p_0}L^{(3+\kappa_{\rm H})y_{\rm c}-\kappa_{\rm H}K}+CK^{3\mathfrak p_0}L^{3u-y_{\rm c}/2}$, the two exponents being $\le-\frac K4+\frac72$ and $3u-\frac K{28}$ for $c_*=\frac1{28}$, and $\le-\frac K8+\frac{15}4$ and $3u-\frac K{12}$ for $c_*=\frac1{14}$; $\mathcal T_K(H)=F(k,v)$, of weight $\le3L^K|H|_\infty$: $\le CK^{2\mathfrak p_0}L^{4u}$.

These bounds depend on the type and the weight alone, so they hold for every copy. Since $u\le\frac K{84}+1$, every exponent of a small kind is $\le3u-c_*K+4$ for $K\ge5$, and a family has at most nine members. Finally $3u-c_*K\ge-K$ and $4u\le K/21+4$.
\end{proof}

\begin{proposition}[One torus, one cell: products]\label{stc:prod}
In the situation of Lemma~\ref{stc:fam}, let $m\ge1$. There is $K_{\rm P}$, a function of $g,\ell',c_*,\phi,m,\delta$, such that for all $(K,\beta)\in\mathfrak d(g)$ with $K\ge K_{\rm P}$, every $\mathsf d>0$ with $\mathsf dL^u\ge K^2$, every board with offset in $a_{K-t}\mathbb Z^3$, and all $\mathsf X_1,\dots,\mathsf X_m$ with supports in one cell, at distance $\ge\delta$ from its walls and pairwise at sup distance $\ge\mathsf d$,
\[
\Bigl\|\prod_{a=1}^m\bigl(\mathsf X_a-\langle\mathsf X_a\rangle_{K,\ell',\beta}\bigr)\Bigr\|_{{\rm ch},\ell'}\le\mathsf C_m\prod_{a=1}^m\varpi(\mathsf X_a)\,\mathsf n(\mathsf X_a),\qquad\mathsf C_m=\max_{m'\le m}(\mathsf{Op}_{8m'}+1)^{1/8},
\]
where $\mathsf{Op}_N=\sum_\pi\prod_{I\in\pi}\mathsf{op}_{\#I}$ over the partitions $\pi$ of $\{1,\dots,N\}$ without blocks of one element.
\end{proposition}

\begin{proof}[Proof sketch]
Both sides are homogeneous in each datum; let $\mathsf n(\mathsf X_a)=1$. The eighth power of the left side is the mean of the $8m$ functions $Y_{(a,\kappa)}=\mathsf X_a\circ\mathsf r_\kappa^{-1}-\langle\mathsf X_a\rangle$, which are centred by the invariance of the torus measure under lattice isometries; so it equals $\sum_\pi\prod_{I\in\pi}\langle Y_i;i\in I\rangle$ over the partitions without singletons. Two copies of one datum are $\ge2\delta$ apart, two data of one cell are $\ge\mathsf d$ apart in every copy, and by Lemma~\ref{stc:fam}\,(b) the carriers of different entries are $\ge\mathsf d_\times=\min\{\mathsf d,2\delta\}-8a_K\ge\frac12\min\{\mathsf d,2\delta\}$ apart. In a cumulant each $\mathsf X_a$ may be replaced by its family: its difference with the family consists of null terms, on whose bonds the other entries do not depend, so that their joint cumulants with the other entries vanish, and of a constant, which drops out of joint cumulants of order $\ge2$ (Lemma~\ref{stc:null}\,(b), after moving the copy back by an isometry).

So $\langle Y_i;i\in I\rangle$ is a sum of at most $9^{\#I}$ cumulants of insertions, each at most $\mathsf{op}_{\#I}\prod\omega_t+\epsilon_\circ$ by Proposition~\ref{stc:cum}\,(a), $\epsilon_\circ$ a bound of its last two members. By Lemma~\ref{stc:cmean}\,(c), (d), the weights $\le(3+60\mathsf b_+)L^{2K}$, $\mathsf dL^u\ge K^2$ and $\mathfrak p_0\ge1$,
\[
\mathsf t_\times\le e^2C_{{\rm E},0}K\ell'^3L^{3u}e^{-K^2/2M},\qquad\mathsf R_t\le C_{\rm R}L^{3u}\exp[-c'A_0^2(\phi\log L)^2K^2]
\]
for $\phi K\log L\ge1$; both members of $\epsilon_\circ$ are $\exp[-{\rm const}\,K^2+O(K\log K)]$, so $9^{8m}\epsilon_\circ\le L^{-16mK}$ times a pure number for $K\ge K_{\rm P}$, and the condition of Proposition~\ref{stc:cum}\,(a) holds. Each $i$ lies in exactly one block; summing, $\|\cdot\|^8\le\mathsf{Op}_{8m}\prod_a\varpi_a^8+L^{-8mK}\le(\mathsf{Op}_{8m}+1)\prod_a\varpi_a^8$, because $\varpi_a\ge L^{-K}$.
\end{proof}

\begin{corollary}[Seminorms of single functions]\label{stc:one}
Let $0<\delta\le\frac18\mathsf l$, and for a cutoff $K$ put $u=\lceil K/400\rceil$, $t=K-u$. There are $C,p,K_0$, functions of $g,\ell',\delta$, such that for all $(K,\beta)\in\mathfrak d(g)$ with $K\ge K_0$, every board of eight cells of side $\mathsf l$ on $T_{\ell'}$ with an offset in $a_{K-t}\mathbb Z^3$, and all data with support in one cell at distance $\ge\delta$ from its walls, with $\langle\cdot\rangle=\langle\cdot\rangle_{K,\ell',\beta}$:
\begin{enumerate}
\item[(a)] let $m^\circ,c^\circ,N^\circ$ be numbers; there is $K^\circ_0$, a function of $g,\ell',\delta,m^\circ,c^\circ,N^\circ$, such that for $K\ge K^\circ_0$: if $F$ is a real function of the cell that is a sum of at most seven insertions of orders $\ge2$, of types $(m,c,\hat c,N)$ with $m\le m^\circ$, $c\le c^\circ$, $\hat c\le c^\circ$ and $N\le N^\circ$, with weights $\le L^{2K}\mathsf b_+$ and carriers within $4a_K$ of the support, and if for every mirror copy the oscillations of the copies of these insertions add up to at most $\omega\le1$, then $\|F-\langle F\rangle\|_{{\rm ch},\ell'}\le C_8(\omega+L^{-K})$, with a pure number $C_8$. In (b)--(d), (a) is applied to insertions whose orders, constants, heights and counting numbers are pure numbers (Lemma~\ref{stc:fam}\,(b)), and $K_0$ is taken not smaller than $K^\circ_0$ at these numbers;
\item[(b)] $\|\mathsf{St}[\eta]-\langle\mathsf{St}[\eta]\rangle\|_{{\rm ch},\ell'}=\|\mathsf{Rm}[\eta]-\langle\mathsf{Rm}[\eta]\rangle\|_{{\rm ch},\ell'}\le CK^pL^{-K/28}\|\eta\|$, and the same bound holds for $\mathfrak K[\eta]-\frac14\mathfrak N[\mathsf k_\eta]$, whose torus mean is $0$;
\item[(c)] $\|\Delta[\eta]-\langle\Delta[\eta]\rangle\|_{{\rm ch},\ell'}\le CK^pL^{-K/36}\|\eta\|$, and the same holds for $\mathcal T^{\rm c}[\eta]$;
\item[(d)] $\|\mathcal T_K(H')-\langle\mathcal T_K(H')\rangle\|_{{\rm ch},\ell'}\le CK^{2\mathfrak p_0}L^{K/100}|H'|_\infty$.
\end{enumerate}
\end{corollary}

\begin{proof}[Proof sketch]
(a) The eighth power of the left side is the mean of the eight centred copies, which is the sum, over the partitions of the eight copies without blocks of one element, of products of their joint cumulants. Copies in different cells have carriers at sup distance $\ge2\delta-8a_K\ge\delta$. The copies of the insertions are insertions with the same weights, of orders $\le m^\circ$, constants $\le2^{6m^\circ}c^\circ$, heights $\le c^\circ$ and counting numbers $\le20N^\circ$ (Lemma~\ref{stc:fam}\,(a)). A joint cumulant of $r+1\ge2$ copies is a sum of at most $7^{r+1}$ joint cumulants of insertions, constants dropping out, each bounded by Proposition~\ref{stc:cum}\,(a); its condition holds for large $K$, as $r_1\bar{\mathsf d}\le\frac12r_*\bar{\mathsf d}$ is at most a number of the constants and $\mathsf t_\times\to0$.

In its last two members $\bar{\mathsf d}\le CL^{2K}$, $1/r_1\le CK^{3\mathfrak p_0}L^{3K}$, $\mathsf t_\times\le e^2C_{{\rm E},0}K\ell'^3L^{3u}e^{-\delta L^u/M}$, $\sup|\mathfrak J|\le CL^{5K}$, with $C$ a function of $\ell'$, $\mathsf b_+$, $m^\circ$, $c^\circ$, $N^\circ$ (Lemma~\ref{stc:cmean}\,(c), (d), at these types: the numbers $\mathsf d_i$, $\mathsf o_k$ and $1/r_{\mathfrak J}$ of Proposition~\ref{stc:chain} increase with the constant, the height and the counting number, and the orders are finitely many), and $\mathsf R_t$ lies below every power of $L^{-K}$ (Lemma~\ref{stc:cmean}\,(d)); so these members are $\le7^{-8}L^{-8K}$ for $K\ge K^\circ_0$, a function of $g$, $\ell'$, $\delta$, $m^\circ$, $c^\circ$, $N^\circ$.

Summing over the choices of the insertions, the products of their oscillations add up to at most $\mathsf{op}_{r+1}\omega^{r+1}$; so a block $I$ of a partition contributes at most $\mathsf{op}_{\#I}\omega^{\#I}+L^{-8K}\le(\mathsf{op}_{\#I}+1)(\omega+L^{-K})^{\#I}$, and summing over the partitions gives the bound with $C_8^8=\sum_\pi\prod_{I\in\pi}(\mathsf{op}_{\#I}+1)$, the sum over the partitions $\pi$ of eight elements without blocks of one element. (b) The mirror copy of a null term is a null term, and the torus measure is invariant under lattice isometries; so Lemma~\ref{stc:null}\,(a), (b), applied to one copy at a time, show that a joint cumulant of $r+1$ copies of $\mathsf{St}[\eta]$ is $(-1)^{r+1}$ times the joint cumulant of the corresponding copies of $\mathsf{Rm}[\eta]$; over a partition of the eight copies the signs multiply to $+1$, and the two eighth moments agree.

The family of $\mathsf{Rm}[\eta]$ has seven members (Lemma~\ref{stc:fam}), whose copies have oscillations $\le CK^pL^{-K/28}\|\eta\|$ by Proposition~\ref{stc:cubic} at $\tau=\frac14$, $\phi=\frac1{400}$, where $c_0'=\frac1{28}$, and by Lemma~\ref{stc:cmean}\,(c); (a) applies after division by a multiple of $\|\eta\|$. By Lemma~\ref{stc:null}\,(a), $\mathfrak K[\eta]-\frac14\mathfrak N[\mathsf k_\eta]=\mathsf{Rm}[\eta]-\frac32a_K\mathcal O(w_\eta)+\frac14c_{\mathsf k_\eta}$ is a constant plus six of these members; its torus mean vanishes by \ref{C-360-8}(b) and by Lemma~\ref{stc:null}\,(b) at $Y=1$. (c) The two members of the family of $\Delta[\eta]$ have oscillations $\le CK^pL^{3u-K/28}\|\eta\|\le CK^pL^{-K/36}\|\eta\|$ (Lemma~\ref{stc:fam}\,(c) at $c_*=\frac1{28}$, as $\frac3{400}+\frac1{36}<\frac1{28}$); and $b\,\mathcal T^{\rm c}[\eta]=\mathsf{St}[\eta]-\Delta[\eta]$ with $b\ge\gamma^{-2}\ge4$, the seminorm being subadditive and homogeneous (\ref{IV-18-18}). (d) $\mathcal T_K(H')$ is one rough insertion, with oscillation $\le CK^{2\mathfrak p_0}L^{4u}|H'|_\infty$ (Lemma~\ref{stc:fam}\,(c)), and $4u\le\frac K{100}+4$.
\end{proof}

\begin{lemma}[From one torus to every periodic Gibbs state]\label{stc:boards}
There is $\mathsf Q=\mathsf Q(g,\ell')\ge1$ such that for all $(K,\beta)\in\mathfrak d(g)$:
\begin{enumerate}
\item[(a)] for every $\nu\in\mathcal G^{\rm per}_{K,\beta}$, every board with cells of side $\mathsf l$ (any offset), every finite set $\mathcal K_0$ of cells and all real functions $F_{\mathbf k}$ of the cells $\mathbf k\in\mathcal K_0$,
\[
\Bigl|\nu\Bigl(\prod_{\mathbf k\in\mathcal K_0}F_{\mathbf k}\Bigr)\Bigr|\le\prod_{\mathbf k\in\mathcal K_0}\mathsf Q\,\|F_{\mathbf k}\|_{{\rm ch},\ell'} ;
\]
in particular $|\nu(F)-c|\le\mathsf Q\|F-c\|_{{\rm ch},\ell'}$, and $|\langle F_1;F_2\rangle_\nu|\le2\mathsf Q^2\|F_1-c_1\|_{{\rm ch},\ell'}\|F_2-c_2\|_{{\rm ch},\ell'}$ for functions of two different cells and real $c,c_1,c_2$;
\item[(b)] the same bound holds for the torus measure $\langle\cdot\rangle_{K,\ell,\beta}$ of every side $\ell\in\mathfrak L(g)$, $\ell\ge\ell'$, in the place of $\nu$, if the cells of $\mathcal K_0$ lie in one period box of $T_\ell$.
\end{enumerate}
\end{lemma}

\begin{proof}[Proof sketch]
(a) Let $\ell_k=2L^{m_1+k}\in\mathfrak L(g)$. The domination of periodic Gibbs states (Proposition~\YMMref{91.9}{705}) gives $|\nu(\prod F_{\mathbf k})|\le\liminf_k\prod\|F_{\mathbf k}\|_{{\rm ch},\ell_k}$; the comparison of two sides (\ref{IV-18-18}\,(e)) gives $\|F\|_{{\rm ch},\ell_k}\le\|F\|_{{\rm ch},\ell'}e^{\mathsf l^3(p_{\ell'}-p_{\ell_k})}$, $p_\ell=\ell^{-3}\log Z_{K,\ell,\beta}$; and the transfer statement \ref{lem:canonical-source-system}, with \eqref{eq:source-system-pressure}, taken at zero source, bounds the factor $e^{\mathsf l^3(p_{\ell'}-p_{\ell_k})}=Z_{K,\ell',\beta}^{1/8}Z_{K,\ell_k,\beta}^{-(\mathsf l/\ell_k)^3}$ by a number $\mathsf Q$ that depends on none of $K,\beta,\ell_k$. Here Theorems~\YMMref{AJ}{673} and~\YMMref{AD}{607} enter, as a bound of stability; no reflection positivity of $\nu$, and nothing of $\nu$ beyond periodicity, is used. The two-cell forms follow from $\langle F_1;F_2\rangle_\nu=\nu((F_1-c_1)(F_2-c_2))-\nu(F_1-c_1)\nu(F_2-c_2)$. (b) The chessboard estimate on $T_\ell$ for cells of one period box (\ref{IV-18-18}\,(b)), the comparison of the two sides $\ell$, $\ell'$, and the same bound of the pressures.
\end{proof}

\begin{lemma}[A change of scale]\label{stc:dil}
Let $z\ge1$ be an integer, $\Lambda=L^z$, $D^z:\mathbb L_K\to\mathbb L_{K-z}$, $x\mapsto\Lambda x$, $(D^zU)(D^z\mathfrak b)=U(\mathfrak b)$, $H^{(z)}(x)=H(\Lambda^{-1}x)$ and $\eta^{[z]}(x)=\Lambda\eta(\Lambda^{-1}x)$.
\begin{enumerate}
\item[(a)] $\mathcal T_K(H)=\Lambda\,\mathcal T_{K-z}(H^{(z)})\circ D^z$; $(H^{\rm cons}[\eta])^{(z)}=H^{\rm cons}[\eta^{[z]}]$; $\|\eta^{[z]}\|\le\Lambda\|\eta\|$ for $\Lambda\ge1$, $|H^{(z)}|_\infty=|H|_\infty$, and supports are multiplied by $\Lambda$.
\item[(b)] $D^z_*$ maps $\mathcal G^{\rm per}_{K,\beta}$ onto $\mathcal G^{\rm per}_{K-z,\beta}$, and the torus measure of $T_\ell$ at $(K,\beta)$ to that of $T_{\Lambda\ell}$ at $(K-z,\beta)$.
\item[(c)] Let $0<g\le\gamma$, and let $g'_1,g'_2$ be the reference couplings with $g'^{-2}_1=\Lambda g^{-2}/2$, $g'^{-2}_2=L\Lambda g^{-2}/4$. For every $(K,\beta)\in\mathfrak d(g)$ with $K>z$ there is $\iota\in\{1,2\}$ with $(K-z,\beta)\in\mathfrak d(g'_\iota)$.
\item[(d)] If $2\mathsf l_\iota$ is the least side at $g'_\iota$, then $\mathsf l_\iota/\Lambda\to0$ and $g'_\iota\to0$ as $z\to\infty$; so $g'_\iota\le\gamma^{\rm sd}_W\le\gamma^{\rm ref}_W$ for large $z$.
\end{enumerate}
\end{lemma}

\begin{proof}[Proof sketch]
(a) $x(D^zp)=\Lambda x(p)$, $x_{D^z\mathfrak b}=\Lambda x_{\mathfrak b}$, plaquette and bond variables are carried, and $a_{K-z}^{-1}=\Lambda^{-1}a_K^{-1}$; insert in Definition~\ref{stc:fields}\,(1). The derivatives of orders $0,1,2$ of $\eta^{[z]}$ get the factors $\Lambda,1,\Lambda^{-1}$, and $\partial\eta^{[z]}=(\partial\eta)^{(z)}$. (b) $D^z$ is a bijection of the bonds and plaquettes that keeps incidence; the Gibbs kernel of a finite set of bonds depends on $\beta$ and the plaquettes alone, not on the spacing; $D^z\tau_v=\tau_{\Lambda v}D^z$ carries periodicity to periodicity. (c) By Lemma~\ref{stc:frame}\,(a), $g^{-2}<b\le L(g^{-2}+\bar B)$. As $g'^{-2}_2>g'^{-2}_1\ge\frac{13}2g^{-2}$, $g'_2<g'_1<g\le\gamma$ are reference couplings.

With $LB_\sharp=(L-1)\bar B$: $I(g'_1)=(\Lambda g^{-2}/2+\bar B,L\Lambda g^{-2}/2+(L-2)\bar B]$ and $I(g'_2)=(L\Lambda g^{-2}/4+\bar B,L^2\Lambda g^{-2}/4+(L-2)\bar B]$; since $\bar B<\gamma^{-2}$ (equation~\YMMref{(1523)}{614}), $\gamma^{-2}\le g^{-2}$ and $L\ge13$, these two intervals overlap and cover $(\Lambda g^{-2},L\Lambda(g^{-2}+\bar B)]\ni\Lambda b$. For $b'\in I(g')$ and every $K'$ the pair $(K',L^{K'}b')$ is a first-passage datum at $g'$ with label $b'$ (Corollary~\YMMref{74.2}{616}, at the reference coupling $g'\le\gamma$); take $K'=K-z$, $b'=\Lambda b$: $L^{K-z}\Lambda b=\beta$. (d) By the definition of the sides (\YMMref{(1528)}{615}) through the cube-size function $R$ of \YMMref{74}{614}, which grows like a power of $\log x$, the least side at $g'_\iota$ grows like a power of $z\log L+\log(g^{-2}+\bar B)$, while $\Lambda=L^z$; and $g'^{-2}_\iota\ge\Lambda g^{-2}/2\to\infty$.
\end{proof}

\begin{lemma}[Schwarz at a lattice mirror; uniform boundedness]\label{stc:schwarz}
\begin{enumerate}
\item[(a)] Let $\nu\in\mathcal G^{\rm per}_{K,\beta}$, $\mu_0$ a direction, $c\in a_K(\mathbb Z+\frac12)$, $\theta_c$ the reflection in $\{x_{\mu_0}=c\}$, and $F$, $G$ real bounded Borel local functions of the bonds in $\{x_{\mu_0}>c\}$, resp.\ in $\{x_{\mu_0}<c\}$. Then $\langle\theta_cF;F\rangle_\nu\ge0$, $\langle\theta_cG;G\rangle_\nu\ge0$ and $|\langle F;G\rangle_\nu|\le\|F\|_{\nu,c}\|G\|_{\nu,c}$, where $\|F\|_{\nu,c}=\langle\theta_cF;F\rangle_\nu^{1/2}$.
\item[(b)] Let $\rho_j$ be probability measures on the configurations of $\mathbb L_{K_j}$ such that 
\[
\lim_j\langle\mathcal T_{K_j}(H);\mathcal T_{K_j}(H')\rangle_{\rho_j}
\]
 exists for all $H,H'\in C_c(\mathbb R^3;V)$ with disjoint supports. For disjoint compact sets $A,A'$ there is $C_2(A,A')<\infty$ with 
\[
|\langle\mathcal T_{K_j}(H);\mathcal T_{K_j}(H')\rangle_{\rho_j}|\le C_2(A,A')|H|_\infty|H'|_\infty
\]
 for all $j$ and all continuous $H,H'$ that vanish off $A$, $A'$.
\end{enumerate}
\end{lemma}

\begin{proof}[Proof sketch]
(a) $\nu$ is invariant under $\theta_c$ and reflection positive at the bond mirror $\{x_{\mu_0}=c\}$ (\textup{\YMMref{103.4}{758}}); so $(A,B)=\nu(\overline{\theta_cA}\,B)$ is a positive semidefinite Hermitian form on the bounded local functions of the closed upper side. With $F_0=F-\nu(F)$ and $G_0=G-\nu(G)$: $\langle F;G\rangle_\nu=(\theta_cG_0,F_0)$, $(F_0,F_0)=\langle\theta_cF;F\rangle_\nu$, $(\theta_cG_0,\theta_cG_0)=\langle\theta_cG;G\rangle_\nu$; Schwarz's inequality. (b) On the Banach spaces $E_A$, $E_{A'}$ of continuous $V$-valued functions vanishing off $A$, $A'$ with the sup norm, the bilinear forms $B_j(H,H')=\langle\mathcal T_{K_j}(H);\mathcal T_{K_j}(H')\rangle_{\rho_j}$ are bounded ($|\mathcal T_K(H)|\le a_K^{-1}(2\#\{p:x(p)\in A\}+\#\{\mathfrak b:x_{\mathfrak b}\in A\})|H|_\infty$) and converge at every point, so they are bounded at every point; the principle of uniform boundedness, applied twice, gives $\sup_j\|B_j\|<\infty$.
\end{proof}

\subsection{Proofs of the results}

\begin{proof}[Proof of Theorem~\ref{stc:cutoff}]
\emph{Step 1: parameters.} For (i), (ii) take $c_*=\frac1{14}$, for (iii) $c_*=\frac1{28}$, and $\phi=(c_*-c)/(4n-1)\le\frac1{84}$, with $u$, $t$ and the other parameters of Lemma~\ref{stc:fam}.

\emph{Step 2: pieces.} Let $\chi_i(x)=\prod_\mu\psi(x_\mu/r_0-i_\mu)$, $i\in\mathbb Z^3$, $r_0=\mathsf l/(8n)$, with $0\le\psi\in C^\infty(\mathbb R)$, ${\rm supp}\,\psi\subset(-1,1)$, $\sum_{j\in\mathbb Z}\psi(\cdot-j)=1$: a smooth partition of unity with ${\rm supp}\,\chi_i$ of sup diameter $\le2r_0$. All the functions $\mathsf X_a$ are linear in their data, $A^\nu$ is linear and cumulants are multilinear; so each left side is a sum of at most $(2R/r_0+3)^{3n}$ terms in which each datum is replaced by a piece $\chi_i\eta_a$ or $\chi_iH_a$, with $\|\chi_i\eta\|\le C_\chi\|\eta\|$ and $|\chi_iH|_\infty\le|H|_\infty$. The pieces are pieces of the data: $H^{\rm cons}[\chi_i\eta]$, not $\chi_iH^{\rm cons}[\eta]$. In one term the supports have sup diameter $\le\mathsf l/(4n)$ and are pairwise $\ge\mathsf d$ apart; no relation between the size of the pieces and $\mathsf d$ is asked.

\emph{Step 3: board.} Lemma~\ref{stc:board}\,(b) with $h=a_{K-t}=L^{-u}<\mathsf l/(2n)$ gives an offset in $a_{K-t}\mathbb Z^3$ for which every $\mathsf X_a$ has its support in one cell, at distance $\ge\delta=\mathsf l/(8n)$ from its walls; no wall needs to separate two supports.

\emph{Step 4: subsets.} Let $c_a=\langle\mathsf X_a\rangle_{K,\ell',\beta}$ and $\mathsf e_a=\mathsf Q\,\mathsf C_n\varpi(\mathsf X_a)\mathsf n(\mathsf X_a)$. Grouping the factors by cells, Lemma~\ref{stc:boards}\,(a) and Proposition~\ref{stc:prod} in each cell give, for every $B\subset\{1,\dots,n\}$,
\[
\Bigl|\nu\Bigl(\prod_{a\in B}(\mathsf X_a-c_a)\Bigr)\Bigr|\le\prod_{a\in B}\mathsf e_a .
\]

\emph{Step 5: centring.} $A^\nu[\mathsf X_a]=(\mathsf X_a-c_a)-\delta_a$, $\delta_a=\nu(\mathsf X_a)-c_a$, and $|\delta_a|\le\mathsf e_a$ by Step 4 at $B=\{a\}$. Expanding the product gives $2^n$ terms, each of modulus $\le\prod_{a=1}^n\mathsf e_a$: every term carries the factor of every entry, the small ones among them.

\emph{Step 6: cumulants.} A joint cumulant of order $n\ge2$ is unchanged when constants are subtracted from its entries, and it is a sum over partitions of products of moments; so $|\langle\mathsf X_1;\dots;\mathsf X_n\rangle_\nu|\le\mathsf{op}_n\prod_a\mathsf e_a$.

\emph{Step 7: exponent.} By Lemma~\ref{stc:fam}\,(c), $\prod_a\varpi(\mathsf X_a)=C_\omega^nK^{np_\omega}L^E$, and since $u\le\phi K+1$ and $n-n_1\le n_1(n-1)$,
\[
E=n_1(3u-c_*K)+4u(n-n_1)\le-n_1K\bigl(c_*-(4n-1)\phi\bigr)+4n=-cn_1K+4n .
\]

\emph{Step 8: constants.} Both bounds hold with $p=np_\omega$ and 
\[
C=(2R/r_0+3)^{3n}\max\{2^n,\mathsf{op}_n\}(\mathsf Q\mathsf C_nC_\omega C_\chi)^nL^{4n}
\]
, for $K\ge K_0=\max_{m\le n}K_{\rm P}(m,\mathsf l/(8n))$, enlarged so that $L^{-u}<\mathsf l/(2n)$, whenever $\mathsf dL^u\ge K^2$. The distance has entered only through this condition. For (iii) this is the assertion; for (i), (ii), $K_0$ is enlarged so that $\mathsf dL^{\lceil\phi K\rceil}\ge K^2$ for all $K\ge K_0$. The cases of (ii) follow also from (i): $\Delta=\mathsf{St}-b\,\mathcal T^{\rm c}$ with $b\le\mathsf b_+$ in a multilinear product; $\mathsf{St}[\eta]$ and $-\mathsf{Rm}[\eta]$ have equal centred products and cumulants with distant entries (Lemma~\ref{stc:null}\,(c)); and a joint cumulant is a sum over partitions of products of the moments of Step 4. (iv) is Step 4 for $\mathsf X_a=\mathfrak K[\eta_a]$, $B=\{1\},\{1,2\}$, $c_a=\langle\mathfrak K[\eta_a]\rangle_{K,\ell',\beta}=0$ (\ref{C-360-8}(b)).
\end{proof}

\begin{proof}[Proof of Corollary~\ref{stc:two}]
The bounds in (b), (c), (d) and the first bound in (a) are Theorem~\ref{stc:cutoff}\,(i), (ii) at $n=2$: with $n_1=2$ for (a), (b) and for (c) with $Z=\mathsf{St}[\eta'],\Delta[\eta']$; with $n_1=1$ for (c) with $Z=\mathcal T_K(H')$ and for (d). The equality in (a) is Lemma~\ref{stc:null}\,(c), applied twice in $\nu$, first with $Y=A^\nu[\mathsf{St}[\eta']]$ and then with $Y=A^\nu[\mathsf{Rm}[\eta]]$ ($\mathsf d>8a_K$ for large $K$). For the mean and the products in (a), Step 4 of the proof above is applied to $\mathsf X_a=\mathfrak K[\cdot]$ with $c_a=\langle\mathfrak K[\cdot]\rangle_{K,\ell',\beta}=0$, the orbital remainder having mean zero under every torus measure (\ref{C-360-8}(b)): at $B=\{1\}$ it bounds $|\nu(\mathfrak K[\eta'])|$, at $B=\{1,2\}$ it bounds $|\nu(\mathfrak K[\eta']\mathfrak K[\eta])|$, and the product of the means is bounded by the product of the two bounds at $B=\{1\}$.
\end{proof}

\begin{proof}[Proof of Theorem~\ref{stc:scale}]
\emph{Step 0: the number $z$.} With the notation of Lemma~\ref{stc:dil}, let $z\ge1$ be the least integer such that $g'_\iota\le\gamma^{\rm ref}_W$ for $\iota=1,2$ and $\mathsf l_2\le\Lambda\mathsf d/2$, $2\mathsf l_\iota$ the least side at $g'_\iota$; it exists by Lemma~\ref{stc:dil}\,(d), and $\mathsf l_1\le\mathsf l_2$. It is a function of $g$ and $\mathsf d$.

\emph{Step 1: the labels.} Let $(K,\beta)\in\mathfrak d(g)$, $K>z$. By Lemma~\ref{stc:dil}\,(c) there is $\iota$ with $(K',\beta)\in\mathfrak d(g')$, $K'=K-z$, $g'=g'_\iota\le\gamma^{\rm ref}_W$. Put $\mathsf l=\mathsf l_\iota$, $\ell_0=2\mathsf l$, $\delta'=\mathsf l/(8q)$.

\emph{Step 2: the change of scale.} By Lemma~\ref{stc:dil}\,(a), (b), the left members of (S), (T) are $\Lambda^q$ times the same expressions at $(K',\beta)$, for the state $\mathfrak s=D^z_*\nu\in\mathcal G^{\rm per}_{K',\beta}$, resp.\ the torus measure of $T_{\Lambda\ell}$, the field $\eta'=\eta^{[z]}$ and the weights $H'_a=H_a^{(z)}$; $\|\eta'\|\le\Lambda\|\eta\|$, $|H'_a|_\infty=|H_a|_\infty$, and the supports are pairwise at sup distance $\ge\Lambda\mathsf d\ge2\mathsf l$. In (T), two points of $T_\ell$ have sup distance $\le\ell/2$, so $\mathsf d\le\ell/2$, $\mathsf l\le\Lambda\ell/4$, and $\Lambda\ell=2L^{m+z}$ is a side at $g'$.

\emph{Step 3: pieces.} A partition of unity of mesh $\mathsf l/(8q)$ for $\eta'$ and a partition into boxes of the same size for the $H'_a$ cut each datum into at most $P=(2\Lambda R/\delta'+3)^3$, resp.\ $(\Lambda\ell/\delta')^3$, pieces; by multilinearity it is enough to bound one choice of pieces, $\eta'_1$ and $H''_a$, of sup diameters $\le\mathsf l/(4q)$ and pairwise at sup distance $\ge2\mathsf l$, and to multiply by $P^q$.

\emph{Step 4: a board.} Lemma~\ref{stc:board}\,(b) with $h=a_{K'-t'}$, $t'=K'-\lceil K'/400\rceil$, gives a board of side $\mathsf l$ with an offset in $a_{K'-t'}\mathbb Z^3$ in which every piece lies in one cell at distance $\ge\delta'$ from its walls; a cell has diameter $\mathsf l<2\mathsf l$, so the cells of different pieces are different, in (T) also modulo the period of $T_{\Lambda\ell}$.

\emph{Step 5: the seminorms on the least torus at $g'$.} Let $\langle\cdot\rangle_0=\langle\cdot\rangle_{K',\ell_0,\beta}$, $\Phi_1=\mathcal T_{K'}(H^{\rm cons}[\eta'_1])-\langle\cdot\rangle_0$ and $\Phi_a=\mathcal T_{K'}(H''_a)-\langle\cdot\rangle_0$, $a\ge2$. At the coupling $g'$, with the fixed side $\ell_0$, Corollary~\ref{stc:one}\,(c) (that is, Lemma~\ref{stc:fam}\,(c) at $c_*=\frac1{28}$, $\phi=\frac1{400}$, $u'=\lceil K'/400\rceil$, and Proposition~\ref{stc:prod} at $m=1$) gives $\|\Phi_1\|_{{\rm ch},\ell_0}\le CK'^pL^{3u'-K'/28}\|\eta'_1\|\le CK'^pL^{-K'/36}\|\eta'_1\|$, since $\frac3{400}+\frac1{36}<\frac1{28}$. For a spectator, $\|\Phi_a\|^8_{{\rm ch},\ell_0}$ is the mean of the eight centred copies of $\mathcal T_{K'}(H''_a)$ in the eight cells, carried by closed sets pairwise at sup distance $\ge2\delta'$ on $T_{\ell_0}$; the mean of a product of centred functions is the sum, over the partitions without singletons, of products of joint cumulants, and the uniform bound of joint cumulants on a torus (item (i) of the subsection on the input, Lemma~\ref{sf-sym}\,(b)), at the coupling $g'\le\gamma^{\rm ref}_W$, the side $\ell_0\in\mathfrak L(g')$ and $\delta'$, bounds every joint cumulant of $2\le r\le8$ of these copies by $\mathsf K^{\rm ub}_r(g',\ell_0,\delta')\prod60|H''_a|_\infty$ for $K'\ge s_r(\delta')+2$ and $3a_{K'}\le\delta'$, a number that depends neither on $K'$ nor on $\beta$.

So $\|\Phi_a\|_{{\rm ch},\ell_0}\le60\mathsf M_8^{1/8}|H_a|_\infty$ with a number $\mathsf M_8$ that does not depend on $K'$.

\emph{Step 6: from $T_{\ell_0}$ to $\mathfrak s$.} Lemma~\ref{stc:boards}, at the coupling $g'$, (a) for $\mathfrak s=D^z_*\nu$ and (b) for the torus measure of $T_{\Lambda\ell}$, gives for every non-empty $B\subset\{1,\dots,q\}$: $|\mathfrak s(\prod_{a\in B}\Phi_a)|\le\prod_{a\in B}\mathsf Q'\|\Phi_a\|_{{\rm ch},\ell_0}$.

\emph{Step 7: the two polynomials.} The joint cumulant of the $q$ functions is a sum over the partitions of $\{1,\dots,q\}$ of products of the moments of Step 6, and the centred product is a sum of $2^q$ products of such moments; in each term the index $1$ stands in exactly one factor. So both are at most $\max\{2^q,\mathsf{op}_q\}\,\epsilon\prod_{a\ge2}\mathsf m_a$, with $\epsilon=\mathsf Q'CK'^pL^{-K'/36}\|\eta'_1\|$ and $\mathsf m_a=60\mathsf Q'\mathsf M_8^{1/8}|H_a|_\infty$.

\emph{Step 8: back to $(K,\beta)$.} The left members of (S), (T) are at most 
\[
\Lambda^qP^q\max\{2^q,\mathsf{op}_q\}\mathsf Q'C\Lambda C_\chi K^pL^{z/36}L^{-K/36}\|\eta\|\prod_{a\ge2}60\mathsf Q'\mathsf M_8^{1/8}|H_a|_\infty
\]
; the numbers depend on $g'_\iota$, $\mathsf l_\iota$, $q$, $z$ and $R$, resp.\ $\ell$, hence on $g,q,\mathsf d$ and $R$, resp.\ $\ell$; take the larger over $\iota=1,2$. In (T), $\mathcal T_K$ of a weight on $T_\ell$ carried by one cell is $\mathcal T_K$ of its lift, taken through the periodisation. A spectator costs here a number independent of the cutoff, and every piece has a cell of its own: these are the two points of the change of scale.
\end{proof}

\begin{proof}[Proof of Proposition~\ref{stc:rm10}]
(a) By Lemmas~\ref{stc:null}\,(a) and \ref{stc:rsharp}, $\mathsf{Rm}[\eta]$ is a sum of seven insertions: $\mathcal S[b\eta]$, of order $3$ and weight $\le L^{2K}b\sup|\eta|$; four of order $2$ with weights $\le(3+5b)\|\eta\|$; two of orders $3$ and $4$ with weights $\le L^Kb\|\eta\|$; their sites lie within $\frac72a_K$ of ${\rm supp}\,\eta$. The mirror copies are insertions of the same orders and weights (Lemma~\ref{stc:fam}\,(a), whose part on insertions holds for every diagonal isometry of the lattice), and $\mathcal S[b\eta]\circ\mathsf r_\kappa^{-1}=\mathcal S[b\eta^{\mathsf r}]$. In the frame of Proposition~\ref{stc:cubic}\,(a) at the given parameters, the oscillation of every copy of $\mathcal S[b\eta]$ is $\le C\mathsf b_+\mathsf H(\eta)K^pL^{-c''K}$ (Proposition~\ref{stc:cubic}\,(b) for the field $\eta^{\mathsf r}$), and those of the other six are $\le CK^{2\mathfrak p_0}L^{4u-K}$ and $\le CK^{3\mathfrak p_0}L^{4u-K/2}$ times $\|\eta\|$ (Lemma~\ref{stc:cmean}\,(b), (c), at a depth $\phi\le\frac14$); as $\phi<\frac16\psi$ and $\psi<\frac29$, $\frac12-4\phi>c''$.

These bounds depend on the type, the weight and $\mathsf H$ alone, not on the cell or the offset. The carriers of copies in different cells are $\ge2r-8a_K$ apart; Proposition~\ref{stc:cum}\,(b), whose conditions hold for large $K$, since $\mathsf t'_\times\le CL^{3u}e^{-rL^u/M}$ lies below every power of $L^{-K}$ and $\mathsf R_t\le\frac12$ from a cutoff on that depends on $m_1$, $\mathsf b_+$ and $\phi$ (Lemma~\ref{stc:cmean}\,(d)), bounds every joint cumulant of $q$ such copies, $2\le q\le8$, by $(\frac12+\mathsf{op}'_q\mathsf{Bl}_q)\prod_a\epsilon_a$ plus the discarding term, which is smaller for large $K$ (Lemma~\ref{stc:cmean}\,(d)). Indeed, every number $\epsilon$ of the statement is at least $L^{-K}$ times $\mathsf b_+\mathsf H(\eta)$, resp.\ $\|\eta\|$, as $C\ge1$, the powers of $K$ there are $\ge1$ and every exponent of $L$ there is $\ge-K$, while every copy $\mathfrak J$ has $1/r_{\mathfrak J}\le CL^{2K}$ and $\sup|\mathfrak J|\le CL^{5K}$ times the same norm (Lemma~\ref{stc:cmean}\,(c), (d), with the weights above), and $\mathsf t'_\times$ and $\mathsf R_t$ lie below every power of $L^{-K}$; so for large $K$, $2E_0\mathsf t'_\times\prod_{a\in I}(q/r_{\mathfrak J_a})\le\prod_{a\in I}\epsilon_a$ for every set $I$ of at least two of the $q$ copies, which is the last condition of Proposition~\ref{stc:cum}\,(b), and $4\,\mathsf{op}'_q\mathsf R_t\prod_a\sup|\mathfrak J_a|\le\frac12\prod_a\epsilon_a$, which absorbs the discarding term; as $\mathsf{op}'_q$ and $\mathsf{Bl}_q$ increase with $q$, the bound $(\frac12+\mathsf{op}'_q\mathsf{Bl}_q)\prod_a\epsilon_a+\frac12\prod_a\epsilon_a$ gives the factor $1+\mathsf{op}'_8\mathsf{Bl}_8$ of the statement.

The eighth moment of the centred copies is the sum, over the $715$ partitions of eight elements without singletons (each with at most four blocks), of products of joint cumulants of copies, each a sum of at most $7^{\#I}$ cumulants of insertions; this gives the bound. It is $\ge0$ by reflection positivity of the torus measure at a wall (Section~\ref{sec:periodic-chessboard}).
(b) The same without the cubic member, at any $0<\phi<\frac18$: the six exponents are $4\phi-1$ and $4\phi-\frac12$. For $\mathcal S[\Xi]$: the copies $\mathcal S[\Xi^{\mathsf r}]$, Proposition~\ref{stc:cubic}\,(b), Proposition~\ref{stc:cum}\,(b) and the moments as in (a); its mean is $0$ by Lemma~\ref{stc:rsharp}\,(c).
(c) For first-passage data the labels lie in $[\gamma^{-2},X_g+\bar B]$ and $MR(\mathsf b_+)\le L^{m_1}$ for every side (Lemma~\ref{stc:frame}\,(a)). At $\psi=\frac15$, $\kappa_\star=\frac9{10}$ the four numbers are $\frac1{10}$, $\frac3{25}$, $\frac1{10}-\frac15\tau$, $\frac1{10}-3\phi$. For the mirrored pair let $c<c''<\frac1{10}$, $\delta_0=\min\{\frac12\mathsf d,\mathsf l\}$, and cut $\eta$ into at most $C(R/\delta_0)^3$ pieces of sup diameter $\le\frac18\delta_0$, and $\theta\eta$ into their mirror images. A piece of $\theta\eta$ and a piece of $\eta$ are $\ge\mathsf d\ge2\delta_0$ apart, so in one direction their projections are $\ge\delta_0$ apart, and Lemma~\ref{stc:board}\,(a) with $h=a_K$ puts them in two different cells of a board with an offset in $a_K\mathbb Z^3$, at distance $\ge\frac1{16}\delta_0$ from the walls.

By (a), the chessboard seminorm of the centred remainder of each piece is $\le CK^pL^{-c''K}\|\eta\|$, and Lemma~\ref{stc:boards}\,(a) for two cells bounds the covariance of the two pieces by $2\mathsf Q^2$ times the product. The equality with the stress insertions is Lemma~\ref{stc:null}\,(c), used twice. For $\mathfrak K$: $\mathfrak K[\eta]=\mathsf{Rm}[\eta]-\frac32a_K\mathcal O(w_\eta)+\frac14\mathfrak N[\mathsf k_\eta]+\frac14c_{\mathsf k_\eta}$; the null terms and the constant drop out of the covariance (Lemma~\ref{stc:null}\,(b)), and the seminorm of the centred $a_K\mathcal O(w_\eta)$ of a piece is $\le CK^{2\mathfrak p_0}L^{4u-K}\|\eta\|$ as in (b); the seminorm is subadditive.
\end{proof}

\begin{proof}[Proof of Theorem~\ref{stc:limit}]
(i) The supports are compact and pairwise disjoint, hence pairwise at a positive sup distance $\mathsf d$, and they lie in a ball of radius $R$. By Definition~\ref{stc:tlim} and Theorem~\ref{stc:cutoff}\,(i) with $n_1=1$, $\mathsf X_1=\mathcal T^{\rm c}[\eta]$,
\[
\begin{multlined}
|\sigma(H^{\rm cons}[\eta],H_2,\dots,H_n)|=\lim_j\Bigl|\rho_j\Bigl(A^{\rho_j}[\mathcal T^{\rm c}[\eta]]\prod_{a\ge2}A^{\rho_j}[\mathcal T_{K_j}(H_a)]\Bigr)\Bigr| \\
\le\lim_jCK_j^pL^{-cK_j}\|\eta\|\prod_{a\ge2}|H_a|_\infty=0 .
\end{multlined}
\]
(ii) Take $n=2$; and $\partial_i(\theta\eta)^j(x)=\Theta_{ii}\Theta_{jj}(\partial_i\eta^j)(\Theta x)$ gives $H^{\rm cons}[\theta\eta]=\theta H^{\rm cons}[\eta]$, with support in $\{x_3<0\}$.
(iii) By Lemma~\ref{stc:null}\,(a) and $\mathsf{St}[\eta]=b\,\mathcal T^{\rm c}[\eta]+\Delta[\eta]$, identically in the configuration, $\mathfrak K[\eta]=\mathcal N_\eta-b\,\mathcal T^{\rm c}[\eta]-\Delta[\eta]-\frac32a_K\mathcal O(w_\eta)$, and likewise for $\theta\eta$; the only null term is $\mathcal N_\eta$, there is no constant, and the null terms drop out of the covariance by Lemma~\ref{stc:null}\,(b). Hence $\langle\mathfrak K[\theta\eta];\mathfrak K[\eta]\rangle_{\rho_j}-b_j^2\langle\mathcal T^{\rm c}[\theta\eta];\mathcal T^{\rm c}[\eta]\rangle_{\rho_j}$ is a sum of covariances each of which has a factor $\Delta$ or $a_K\mathcal O(w)$. The family of such a factor has oscillations $\le CK^pL^{3u-c_*K}$ by Lemmas~\ref{stc:defect}, \ref{stc:cmean} and \ref{stc:cut} alone; the other factor is either of the same kind or $b\,\mathcal T^{\rm c}[\cdot]=\mathcal T_K(bH^{\rm cons}[\cdot])$, treated as a tensor field at the price $CK^pL^{4u}$.

Steps 2--7 of the proof of Theorem~\ref{stc:cutoff} at $n=2$, $n_1=1$, with $\phi<c_*/7$, bound each covariance by $CK_j^pL^{(7\phi-c_*)K_j}\to0$. With $\gamma^{-2}\le b_j\le\mathsf b_+$ (Lemma~\ref{stc:frame}\,(a)): if $\langle\mathfrak K[\theta\eta];\mathfrak K[\eta]\rangle_{\rho_j}\to0$, then $\langle\mathcal T^{\rm c}[\theta\eta];\mathcal T^{\rm c}[\eta]\rangle_{\rho_j}\to0$, which by Definition~\ref{stc:tlim} is (N); conversely, (N) gives $\langle\mathfrak K[\theta\eta];\mathfrak K[\eta]\rangle_{\rho_j}\to0$, and $\langle\mathsf{St}[\theta\eta];\mathsf{St}[\eta]\rangle_{\rho_j}\to0$ because $\mathsf{St}[\eta]=\mathcal N_\eta-\mathfrak K[\eta]-\frac32a_K\mathcal O(w_\eta)$ identically (Lemma~\ref{stc:null}\,(a)), the null term dropping out and the cross terms with $a_K\mathcal O(w)$ tending to $0$ by the same count.
\end{proof}

\begin{proof}[Proof of Theorem~\ref{stc:torus}]
(i) Any side of $\mathfrak L(g)$ may serve as the fixed side: take $\ell'=\ell$, $c_*=\frac1{28}$, $\phi=(\frac1{28}-c)/(4q-1)\le\frac1{196}$; Lemmas~\ref{stc:frame}, \ref{stc:cmean}, \ref{stc:cut} and Proposition~\ref{stc:cubic} hold on every side at $g$. A partition of unity of $T_\ell$ into pieces of sup diameter $\le\ell/8$ and multilinearity leave a bounded number of cumulants of pieces of $\eta$; each piece lifts to a datum on $\mathbb R^3$ in a cell of side $\ell/2$, at distance $\ge\ell/16$ from its walls, and, the formulas being local, its lattice functions on $T_\ell$ are the periodised ones, so that Lemmas~\ref{stc:null}, \ref{stc:defect} and \ref{stc:fam}\,(c), without mirror copies, hold on $T_\ell$. By Lemma~\ref{stc:null}\,(b) under $\langle\cdot\rangle_{K,\ell,\beta}$, the entry $\mathcal T_K(H^{\rm cons}[\eta])$ of the cumulant may be replaced by its family of nine insertions, the difference being null terms whose bonds the other entries do not depend on and a constant; a spectator $\mathcal T_K(H_a)$ is one rough insertion, of weight $\le3L^K|H_a|_\infty$, whose oscillation bound $CK^{2\mathfrak p_0}L^{4u}|H_a|_\infty$ depends on the weight and the type alone. The cumulant becomes a sum of at most nine cumulants of $q$ insertions with carriers pairwise $\ge\mathsf d-8a_K$ apart, and Proposition~\ref{stc:cum}\,(a) on $T_\ell$ bounds each by
\[
\bigl(\mathsf{op}_qC_\omega^qK^{qp_\omega}L^{3u-K/28+4u(q-1)}+L^{-qK}\bigr)\|\eta\|\prod_{a\ge2}|H_a|_\infty,\qquad3u-\tfrac K{28}+4u(q-1)\le-cK+4q .
\]
No board is needed.
(ii) For $b\in I(g)$, Corollary~\YMMref{74.2}{616} gives $(K,L^Kb)\in\mathfrak d(g)$, with the label $b$, for every $K$. So (i) holds along $\beta=L^Kb$, and the cumulants tend to $0$; for Lipschitz weights with pairwise disjoint supports these are the joint cumulants $\langle\mathcal T_K(H_1);\dots;\mathcal T_K(H_q)\rangle_{K,\ell,L^Kb}$ at $H_1=H^{\rm cons}[\eta]$, whose limit as $K\to\infty$ is the torus limit functional $\mathcal S^{\ell,b}_q$ (item (iii) of the subsection on the input, Theorem~\ref{sf-tori}\,(a), for $0<g\le\gamma^{\rm ref}_W$), and $H^{\rm c}_\eta=-H^{\rm cons}[\eta]$. At $H_a=h_a\mathbf 1$, $\mathcal T_K(h_a\mathbf 1)=\mathcal O(w_{h_a\mathbf 1})$.
\end{proof}

\begin{proof}[Proof of Corollary~\ref{stc:small}]
(i) Let $O_1,\dots,O_n$ be bounded open sets with pairwise disjoint closures; the closures are compact, of bounded sup diameters and pairwise at a positive sup distance. So the uniform bound of centred products in periodic Gibbs states (item (ii) of the subsection on the input, Theorem~\ref{sf-gibbs}\,(a), at $\nu'=\nu$: beyond a threshold in the cutoff, for real bounded weights vanishing off closed sets of sup diameter $\le\varrho_0$ that are pairwise at sup distance $\ge2\delta_0$, the bound $C^{\rm lc}_n(g,\delta_0,\varrho_0)\prod_a|H_a|_\infty$, uniform in the states and the positions) gives $|\sigma_j(H_1,\dots,H_n)|\le C_n\prod_a|H_a|_\infty$ uniformly in $j$, for $K_j$ beyond the threshold of that statement, for all real bounded weights $H_a$ vanishing off $O_a$, where $\sigma_j(H_1,\dots,H_n)=\nu_j(\prod_aA^{\nu_j}[\mathcal T_{K_j}(H_a)])$.

A diagonal argument over a countable family of tuples of $C^1_c$ fields, dense in the sup norm, and the compactness of the range of the labels (Lemma~\ref{stc:frame}\,(a)) give a subsequence along which $a_{K_j}\beta_j\to b_*$ and $\sigma_j$ converges at all tuples of $C^1_c$ fields with disjoint supports; $\sigma_j$ is $n$-linear, and $C^1_c$ is sup-norm dense in the continuous functions vanishing off such a set, so $(\sigma_j)$ is a Cauchy sequence at every $C_c$ tuple. A subsequence of a defining sequence defines the same $\sigma$, and its labels converge along a further subsequence. The bound passes to $\sigma$; as the number of item (ii) depends on the sup diameters of the single supports and on a lower bound of their distances, not on their positions, the same statement gives the second form, for tuples of any extent, and this form passes to $\sigma$ in the same way.

A lattice isometry $\mathsf r$ carries $\mathcal T_K(H)$ to $\mathcal T_K(\mathsf rH)$: it permutes the plaquettes and the bonds, and $q_{\mathfrak b}$ changes by the signs that $\mathsf R$ gives to the two directions normal to $\mathfrak b$. The $\nu_j$ are invariant under the reflections in the lattice mirrors of the nine families; products of two of them give translations of the lattice, and those through a site generate the cubic group. For $v\in\mathbb R^3$ take lattice vectors $v_j\to v$ and use the uniform bound and the uniform continuity of the weights; this gives the covariance. The single entry gives $\nu_j(A^{\nu_j}[\,\cdot\,])=0$, and the trace part is the definition, $\mathcal T_K(f\mathbf 1)=\mathcal O(w_{f\mathbf 1})$.

The limits at $g$ of item (iv) of the subsection on the input are tensor limits in the sense of Definition~\ref{stc:tlim}: the lattice field $O^T_K(H)$ displayed in the subsection on the input, which weights the action density $a_p$ of each plaquette with $H_{\rho(p)\rho(p)}$ at its centre and the bond form $q_{\mathfrak b}$ with $H_{\lambda_{\mathfrak b}\lambda'_{\mathfrak b}}$ at the midpoint of the bond, is $\mathcal O(w_H)+Y[\mathsf w_H]=\mathcal T_K(H)$, with the weights read at the same points and the normalisations of Definition~\ref{stc:fields}\,(1), those limits are limits of products of centred entries along first-passage data at $g$ with $K_j\to\infty$, and the states are periodic Gibbs states.
(ii) Reflection positivity at the nine families of mirrors is in item (iv) of the subsection on the input (Theorem~\ref{sf-gibbs}\,(c)), and the Hilbert space with $H\ge0$ and the cone $\mathcal P(\{E<|q|\})=0$ of the joint spectral measure $\mathcal P$ of $(H,P_1,P_2)$, which contains $H\ge|P_1|$ and $H\ge|P_2|$, in item (v) there (Theorem~\ref{sf-wight}\,(a)); we recall the argument. Reflection positivity of the $\nu_j$ at the nine families of lattice mirrors and their invariance under the reflections in these mirrors (\textup{\YMMref{103.4}{758}}) give positivity at lattice mirrors $\{\mathsf n\cdot x=c_j\}$, $c_j\to c$; a lattice reflection acts on $\mathcal T_K(H)$ as on its weight; the uniform bound of (i) and the uniform continuity of $C^1_c$ weights pass this to the limit; the covariance of $\sigma$ under all translations is that of (i).

The Hilbert space is obtained by quotient and completion. For a lattice translation by $k$ steps normal to the mirror, iterated Schwarz inequalities in the positive form, $|u(k)|\le u(0)^{1-2^{-n}}|u(2^nk)|^{2^{-n}}$ for $u(k)=\langle A,\tau_{ka_Ke_3}A\rangle$, give $|\sigma(\theta\bar F'\cdot\tau_{te_3}F)|\le\|[F']\|\,\|[F]\|$; so $e^{-tH}$ is well defined, a semigroup of self-adjoint contractions, weakly continuous by the uniform bound, hence $e^{-tH}$ with $H\ge0$, $H\Omega=0$; the translations parallel to the mirror act unitarily and commute with it, and a countable dense family gives separability. The square cone, $k\in\{1,2\}$, comes from the five-point inequality of periodic Gibbs states (Lemma~\ref{lem:spectral-steps}, at the directions $3$, $k$ and the lattice mirror $\{x_3=0\}$): for $\nu\in\mathcal G^{\rm per}_{K,\beta}$, every bounded Borel local function $A$ of the bonds in $\{x_3\ge0\}$, $\theta_0$ the reflection in $\{x_3=0\}$, $u_\nu(w)=\nu(\overline{\theta_0A}\,\tau_wA)$ for $w\in a_K\mathbb Z^3$, and every $\mathsf t\in a_K\mathbb N$,
\[
u_\nu(0)-4u_\nu(\mathsf te_3)+u_\nu(2\mathsf te_3)+u_\nu(\mathsf te_3+\mathsf te_k)+u_\nu(\mathsf te_3-\mathsf te_k)\ge0 ;
\]
it rests on the reflection positivity of $\nu$ at $\{x_3=0\}$, at $\{x_3=\frac12a_K\}$ and at the diagonal mirror $\{x_3+x_k=0\}$ and on its invariance under $a_K\mathbb Z^3$ (the lattice energy--momentum inequality, Proposition~\textup{\YMMref{99.3}{739}} and Lemma~\ref{lem:spectral-steps}). Let $F=\sum_ic_iX_i$ be a finite complex combination of tuples $X_i=(H_{i,a})_a$ of $C^1_c$ fields with pairwise disjoint supports in $\{x_3>0\}$, and take $\nu=\nu_j$, $A=A_j=\sum_ic_i\prod_aA^{\nu_j}[\mathcal T_{K_j}(H_{i,a})]$, for large $j$ a local function of the bonds in $\{x_3>0\}$, and $\mathsf t=\mathsf t_j\in a_{K_j}\mathbb N$ with $\mathsf t_j\to t>0$. As $\nu_j$ is invariant under the lattice reflections and translations, which act on $\mathcal T_K(H)$ as on its weight, $u_{\nu_j}(w)=\sum_{i,l}\bar c_ic_l\,\sigma_j(\theta X_i\cdot\tau_wX_l)$, with $\sigma_j$ as in (i); by the uniform bound of (i) and the uniform continuity of the weights, $u_{\nu_j}(w_j)\to u(w)=\sigma(\theta\bar F\cdot\tau_wF)$ for $w_j\to w$. So $u(0)-4u(te_3)+u(2te_3)+u(te_3+te_k)+u(te_3-te_k)\ge0$ for every $t>0$; with $\psi=[F]$, $e^{-tH}[F]=[\tau_{te_3}F]$ and $e^{itP_k}[F]=[\tau_{te_k}F]$ this reads
\[
\bigl\langle\psi,\bigl\{1-4e^{-tH}+e^{-2tH}+e^{-tH}\bigl(e^{itP_k}+e^{-itP_k}\bigr)\bigr\}\psi\bigr\rangle=\int4e^{-tE}\Bigl(\sinh^2\frac{tE}2-\sin^2\frac{tq_k}2\Bigr)\,d\langle\psi,\mathcal P\psi\rangle\ge0 ,
\]
$\mathcal P$ the joint spectral measure of $(H,P_1,P_2)$. The classes $[F]$ are dense in $\mathcal H(\sigma)$; so the lemma of slope one at small steps (Lemma~\ref{lem:spectral-steps}), applied to the image of $\mathcal P$ under $(E,q_1,q_2)\mapsto(E,q_k)$, gives $\mathcal P(\{E<|q_k|\})=0$: $E\ge|q_k|$ on the support of $\mathcal P$, that is $H\ge|P_k|$.
(iii) $\|[(H^{\rm cons}[\eta])\cdot F]\|^2$ is a sum of values $\bar c_kc_l\,\sigma(\theta\bar X_k\cdot\theta H^{\rm cons}[\eta]\cdot H^{\rm cons}[\eta]\cdot X_l)$ at tuples with pairwise disjoint supports in which $H^{\rm cons}[\eta]$ stands beside the other entries ($\theta H^{\rm cons}[\eta]=H^{\rm cons}[\theta\eta]$ has its support in $\{x_3<0\}$); each vanishes by Theorem~\ref{stc:limit}\,(i). The last assertion is $\langle[F'],[(H^{\rm cons}[\eta])\cdot F]\rangle=0$.
\end{proof}

\begin{proof}[Proof of Proposition~\ref{stc:second}]
(a) Order the sum that defines $\mathsf\Theta_\beta[\eta]$ by the bonds. The four plaquettes that contain $\mathfrak b$ are $p(\mathfrak b,\pm\tau)$, $x(p(\mathfrak b,\pm\tau))=x_{\mathfrak b}\pm\frac12a_Ke_\tau$, and $\mathbf w(W_{\mathfrak b,p(\mathfrak b,\pm\tau)})=\frac12\mathbf y_{\mathfrak b,\tau}\mp\mathbf u_{\mathfrak b,\tau}$. For $f=\eta^\rho$ and $D_\pm=f(x_{\mathfrak b})-f(x_{\mathfrak b}\pm\frac12a_Ke_\tau)$: $D_++D_-=-\frac{a_K^2}4(\Delta_\tau f)(x_{\mathfrak b})$ and $D_--D_+=a_K(\nabla^{\rm c}_\tau f)(x_{\mathfrak b})$, so that, with $\beta=ba_K^{-1}$,
\[
\mathsf\Theta_\beta[\eta]=b\,a_K^{-1}\sum_{\mathfrak b}\sum_{\rho,\tau}(\nabla^{\rm c}_\tau\eta^\rho)(x_{\mathfrak b})\langle\mathbf u_{\mathfrak b,\rho},\mathbf u_{\mathfrak b,\tau}\rangle-b\,\mathcal R[z^\eta].
\]
For $\rho\ne\tau$ the two terms give $-H^{\rm cons}_K[\eta]_{\lambda_{\mathfrak b}\lambda'_{\mathfrak b}}(x_{\mathfrak b})q_{\mathfrak b}$, which sum to $b\,Y[\mathsf w_{H^{\rm cons}_K[\eta]}]$ exactly. For $\rho=\tau$, $|\mathbf w(W_x(\mu,\pm\rho))|^2=a_\pm(2-a_\pm)$ gives $|\mathbf u_{\mathfrak b,\rho}|^2=(a_++a_-)-\frac12(a_+^2+a_-^2)-\frac14|\mathbf y_{\mathfrak b,\rho}|^2$; the parts with $a_+^2+a_-^2$ and $|\mathbf y|^2$ are $-b\{\frac12\mathcal A[v^\eta]+\frac14\mathcal Y[v^\eta]\}$, and the part with $a_++a_-$, ordered by the plaquettes with $(\nabla^{\rm c}_\tau f)(y+\frac12a_Ke_\tau)+(\nabla^{\rm c}_\tau f)(y-\frac12a_Ke_\tau)=2(\nabla^{\rm s}_\tau f)(y)$, has at $a_K^{-1}a_p$ the weight $2(\nabla^{\rm s}\!\cdot\eta-\nabla^{\rm s}_\lambda\eta^\lambda)(x(p))$, $\lambda=\rho(p)$; taking off $b\,\mathcal O(\nabla^{\rm s}\!\cdot\eta)$ leaves $b\,\mathcal O(w_{H^{\rm cons}_K[\eta]})$.
(b) With $\partial^{\mathsf a}_{\mathfrak b}$ the derivative along $e^{is\sigma_{\mathsf a}}U(\mathfrak b)$ and $J_{\mathfrak b}=\sum_{p\ni\mathfrak b}\mathbf w(W_{\mathfrak b,p})=\mathbf y_{\mathfrak b,\rho}+\mathbf y_{\mathfrak b,\rho'}$, the calculus of the derivations (\ref{lem:ward-general-calculus}) gives $\partial^{\mathsf a}_{\mathfrak b}\sum_{p\ni\mathfrak b}a_p=\langle\sigma_{\mathsf a},J_{\mathfrak b}\rangle$ and $\sum_{\mathsf a}\partial^{\mathsf a}_{\mathfrak b}\langle\sigma_{\mathsf a},\mathbf u_{\mathfrak b,\rho}\rangle=\frac32(a_+-a_-)$. Invariance of Haar measure and the equation of the Gibbs state for the single bond $\mathfrak b$ (\YMMref{(1819)}{698}) give $\beta\langle\langle\mathbf u_{\mathfrak b,\rho},J_{\mathfrak b}\rangle;Y\rangle_\nu=\frac32\langle a_+-a_-;Y\rangle_\nu$. Write $\sum_\tau z_{\rho\tau}\langle\mathbf u_{\mathfrak b,\rho},\mathbf y_{\mathfrak b,\tau}\rangle=(z_{\rho\rho}-z_{\rho\rho'})\langle\mathbf u_{\mathfrak b,\rho},\mathbf y_{\mathfrak b,\rho}\rangle+z_{\rho\rho'}\langle\mathbf u_{\mathfrak b,\rho},J_{\mathfrak b}\rangle$, use $\langle\mathbf u_{\mathfrak b,\rho},\mathbf y_{\mathfrak b,\rho}\rangle=-(a_+-a_-)+\frac12(a_+^2-a_-^2)$, and sum by parts over the plaquettes, which defines $\delta$.
(c) $|\mathbf y_{\mathfrak b,\rho}|^2=4(a_++a_-)-(a_++a_-)^2-2a_{R(\mathfrak b,\rho)}$, $a_{R}$ the action of the holonomy around the $1\times2$ rectangle $R(\mathfrak b,\rho)=p(\mathfrak b,\rho)\cup p(\mathfrak b,-\rho)$, is a density of order $2$ on that rectangle, and $a_p^2$ one of order $4$; so $\mathcal P_K[v,w]$ is an insertion with weights $L^K|(v,w)|_\infty$ whose carrier lies within $a_K$ of the carrier of the weights; a mirror copy of a planar null form is a planar null form with the same sup norm of the weights (the trace of a holonomy around a box depends neither on the corner nor on the sense). Let $u=\lceil\phi K\rceil$, $t=K-u$, $n_0=K-\lceil\psi K\rceil$. At a background $U$ over $V\in\mathsf{Cut}_{n_0}$, $U\in\mathcal C(\mathbf B_{n_0},V)$, regular of level $n_0$ at $\varepsilon_{n_0}$ (Lemma~\ref{stc:reg}\,(a)), Lemma~\ref{stc:reg}\,(c) gives $|\mathbf y_{\mathfrak b,\rho}|\le C\varepsilon_{n_0}L^{-(2+\kappa_\star)n_0}$ and $a_p\le\frac12\varepsilon_{n_0}^2L^{-4n_0}$; with at most $9\ell^3L^{3K}$ boxes, $\epsilon_X\le C\ell^3L^{4K}\varepsilon_{n_0}^2L^{-(4+2\kappa_\star)n_0}|(v,w)|_\infty\le CK^{2\mathfrak p_0}L^{-[2\kappa_\star-(3+2\kappa_\star)\psi]K}|(v,w)|_\infty$.

Lemma~\ref{stc:cut}\,(a) at $n_0$ adds $CK^{3\mathfrak p_0}L^{3u-\psi K/2}|(v,w)|_\infty$, which is $\le CK^{2\mathfrak p_0}L^{4u-\psi K/2}|(v,w)|_\infty$ for large $K$, as $K^{\mathfrak p_0}\le L^{\phi K}$; so every copy has oscillation $\le CK^{2\mathfrak p_0}L^{-c_{\rm Y}K}|(v,w)|_\infty$. Proposition~\ref{stc:cum}\,(b) bounds the joint cumulants of $2\le q\le8$ copies in different cells, whose carriers are $\ge r$ apart, and the eighth moment of the centred copies is the sum over the $715$ partitions without singletons; it is $\ge0$ by Section~\ref{sec:periodic-chessboard}. The supremum of $c_{\rm Y}$ is approached at $\kappa_\star\uparrow1$, $\phi\downarrow0$, $\psi=\frac4{11}$.
(d) Take the board of the least torus whose walls normal to $e_{\mu_0}$ stand at $c+\mathsf l\mathbb Z$, the other walls chosen so that the cube lies in one cell at distance $\ge r$ from its walls ($\mathsf l\ge16r$); $F=\mathcal P_K[v,w]$ and $\theta_cF$ are functions of two neighbouring cells, and the eight copies of $\theta_c(F-c_0)$ are those of $F-c_0$, $c_0$ the torus mean. With $\nu(\theta_cF)=\nu(F)$ and Lemma~\ref{stc:boards}\,(a) for two cells, $\|F\|^2_{\nu,c}=\nu(\theta_c(F-c_0)(F-c_0))-(\nu(F)-c_0)^2\le\mathsf Q^2\|F-c_0\|^2_{{\rm ch},\ell'}$, and (c) at the parameters $(\kappa_\star,\psi,\phi)=(\frac12,\frac18,\frac1{256})$ gives the first assertion. For the limit let $d_0$ be the distance of the supports, $r\le\frac1{24}d_0$, and cut $\eta$ and $H'$ into pieces of sup diameter $<\frac12r$.

By (a), $\mathsf{St}[\eta]=b\,\mathcal T_K(H^{\rm cons}[\eta])+b\,\mathcal T_K(H''_K[\eta])-b\,\mathcal N_K$, $\mathcal N_K=\frac14\mathcal Y[v^\eta]+\frac12\mathcal A[v^\eta]+\mathcal R[z^\eta]$. By Definition~\ref{stc:tlim} and $b_j\to b_*$, $\langle b_j\mathcal T_{K_j}(H^{\rm cons}[\eta]);\mathcal T_{K_j}(H')\rangle_{\rho_j}\to b_*\sigma(H^{\rm cons}[\eta],H')$; the term with $H''_K$, of size $\frac32a_K\sup|\nabla^2\eta|$ and carried by a fixed compact set disjoint from ${\rm supp}\,H'$, tends to $0$ by Lemma~\ref{stc:schwarz}\,(b). By (b), $\langle\mathcal N_K;\mathcal T_K(H')\rangle_{\rho_j}=\langle F;\mathcal T_K(H')\rangle_{\rho_j}+\langle\mathcal T_K(H^\natural_K);\mathcal T_K(H')\rangle_{\rho_j}$ with the planar null form $F=\mathcal P_K[\frac14v^\eta,\frac12\sigma v^\eta+\frac12\delta\zeta^\eta]$ ($\sigma v$ the sum of the weights of a plaquette over its four bonds) and a diagonal weight $H^\natural_K$ of size $\le Ca_K(1+\beta^{-1})\sup|\nabla^2\eta|$, whose term tends to $0$ by Lemma~\ref{stc:schwarz}\,(b).

A lattice mirror $c_j\in a_{K_j}(\mathbb Z+\frac12)$ at distance about $2r$ from ${\rm supp}\,\eta$ parts the carrier of $F$, which lies in a cube of side $r$ as in the first assertion, from the bonds on which $\mathcal T_K(H')$ depends; Lemma~\ref{stc:schwarz}\,(a) gives $|\langle F;\mathcal T_K(H')\rangle_{\rho_j}|\le\|F\|_{\rho_j,c_j}\|\mathcal T_K(H')\|_{\rho_j,c_j}$; the second factor is $\langle\mathcal T_K(\theta_{c_j}H');\mathcal T_K(H')\rangle_{\rho_j}^{1/2}$, bounded uniformly in $j$ by Lemma~\ref{stc:schwarz}\,(b), and the first tends to $0$ by the first assertion.
(e) By reflection positivity at the bond mirror, $\langle F,G\rangle_\rho=\langle\theta_KF;G\rangle_\rho$ is a positive semidefinite form on the bounded functions of the bonds in $\{x_3>\frac12a_K\}$, with seminorm $\|F\|_\rho$, and $\langle\theta_K\mathcal T_K(H^{\rm cons}_K[\eta]);\mathcal T_K(H^{\rm cons}_K[\eta])\rangle_\rho=\|\mathcal T_K(H^{\rm cons}_K[\eta])\|_\rho^2\ge0$. By (a) and Lemma~\ref{stc:null}\,(a), $b\,\mathcal T_K(H^{\rm cons}_K[\eta])-\mathfrak Y[\eta]=\mathcal Z-\frac14c_{\mathsf k_\eta}+X''[\eta]$, $\mathcal Z=\mathcal N_\eta-\frac14\mathfrak N[\mathsf k_\eta]$, $X''[\eta]=-\mathsf{Rm}[\eta]+b(\frac12\mathcal A[v^\eta]+\mathcal R[z^\eta])$, since $\mathfrak Y[\eta]=\frac b4\mathcal Y[v^\eta]$. $\rho(\mathcal ZY)=0$ for every bounded function $Y$ of the lower side (Lemma~\ref{stc:null}\,(b)), and $\theta_K\mathcal Z$ is a null term of the lower side; so the null terms drop out of the form, and $\|b\,\mathcal T_K(H^{\rm cons}_K[\eta])-\mathfrak Y[\eta]\|_\rho=\|X''[\eta]\|_\rho$.

Cut $\eta$ by a partition of unity into pieces of sup diameter $\le\frac18\mathsf l$, $\ell'=2\mathsf l$ the least side at $g$. For a function $F_1$ of the upper side carried by one piece, a board $\mathsf B$ of the least torus with an offset in $a_K\mathbb Z^3$ puts the piece and its mirror image in two cells, both at sup distance $\ge r=\min\{\frac14h_0,\frac1{32}\mathsf l\}$ from the walls (either the bond mirror is a wall of the board, or a wall stands between the piece and its image). Let $F_2=\theta_KF_1$ and $c_1$ the torus mean of $F_1$. As $\rho$ is invariant under $\theta_K$ (\textup{\YMMref{103.4}{758}}), $\rho(F_2)=\rho(F_1)$, so $\|F_1\|_\rho^2=\rho((F_2-c_1)(F_1-c_1))-(\rho(F_1)-c_1)^2$, and Lemma~\ref{stc:boards}\,(a) for two cells gives $\|F_1\|_\rho^2\le\mathsf Q^2\|F_2-c_1\|_{{\rm ch},\ell'}\|F_1-c_1\|_{{\rm ch},\ell'}$, both seminorms formed on $\mathsf B$.

The reflection $\theta_K$ induces an isometry of the lattice of the least torus; it carries $\mathsf B$ to a board $\theta_K\mathsf B$ with an offset in $a_K\mathbb Z^3$ (if the walls of $\mathsf B$ normal to $e_3$ stand at $W+\mathsf l\mathbb Z$, those of $\theta_K\mathsf B$ stand at $a_K-W+\mathsf l\mathbb Z$), conjugates the mirror maps of $\mathsf B$ to those of $\theta_K\mathsf B$ and keeps the torus measure. So $F_2$ has the torus mean $c_1$, and the seminorm of $F_2-c_1$ formed on $\mathsf B$ is the seminorm of $F_1-c_1$ formed on $\theta_K\mathsf B$, in which the piece lies in one cell at sup distance $\ge r$ from the walls. The bounds used next hold for every board with an offset in $a_K\mathbb Z^3$ and every cell, the function being carried at sup distance $\ge r$ from its walls; so both factors fall under one bound.

The seminorm $\|\cdot\|_\rho$ is subadditive, and the pieces add up. The chessboard seminorms are bounded by Proposition~\ref{stc:rm10}\,(a) for $\mathsf{Rm}$ with a rate $c''>c$, by power counting for $b(\frac12\mathcal A[v^\eta]+\mathcal R[z^\eta])$ (orders $4$ and $2$, weights $\le CL^K\|\eta\|$ and $\le C\|\eta\|$, rates $\frac12-4\phi$ and $1-4\phi$), and by (c) for $\mathfrak Y[\eta]$ with parameters such that $c_{\rm Y}>c$ (possible as $c<\frac1{10}<\frac2{11}$). So $\|b\,\mathcal T_K(H^{\rm cons}_K[\eta])\|_\rho\le\|X''[\eta]\|_\rho+\|\mathfrak Y[\eta]\|_\rho\le CK^pL^{-cK}\|\eta\|$. For (N): $\mathcal T_K(H^{\rm cons}_K[\eta])-\mathcal T_K(H^{\rm cons}[\eta])=\mathcal T_K(H''_K[\eta])$ with $|H''_K[\eta]|_\infty\le\frac32a_K\|\eta\|$, and the bound of item (ii) of the subsection on the input (Theorem~\ref{sf-gibbs}\,(a)) at $n=2$ on the covariance of two tensor entries with disjoint supports, by a number times the product of the sup norms of the weights, uniformly in the cutoff and in the periodic Gibbs state, gives $\|\mathcal T_K(H^{\rm cons}_K[\eta])\|^2_\rho-\|\mathcal T_K(H^{\rm cons}[\eta])\|^2_\rho\to0$; $\theta_K\mathcal T_K(H)=\mathcal T_K(\theta_KH)$ with $\theta_KH$ the translate by $a_Ke_3$ of $\theta H$, which tends to $\theta H$ uniformly; so along the defining sequence of a tensor limit $\sigma(\theta H^{\rm cons}[\eta],H^{\rm cons}[\eta])=\lim_j\|\mathcal T_{K_j}(H^{\rm cons}[\eta])\|^2_{\rho_j}=0$ (with $b_j^{-1}\le\gamma^2$).
\end{proof}

\begin{proof}[Proof of Proposition~\ref{stc:free}]
(a) $K_3=\nabla\nabla\varphi$, $\varphi(u)=\frac1{4\pi|u|}$ harmonic off $0$, is symmetric, $\sum_\nu\partial_\nu(K_3)_{\nu\rho}=\partial_\rho\Delta\varphi=0$, and $\partial_\nu(K_3)_{\beta\rho}=\partial_\rho(K_3)_{\beta\nu}$. With $K^i=K_3(\cdot-y_i)$: $\sum_\nu\partial_\nu M_{\nu\rho}=K^1_{\nu\alpha}B_{\alpha\beta}\partial_\rho K^2_{\beta\nu}$ and $\sum_\nu\partial_\nu M_{\rho\nu}=(\partial_\rho K^1_{\nu\alpha})B_{\alpha\beta}K^2_{\beta\nu}$, summed over the repeated indices; the sum is $\partial_\rho(K^1_{\nu\alpha}B_{\alpha\beta}K^2_{\beta\nu})=\partial_\rho{\rm tr}\,M$. Integrating by parts, $\int{\rm tr}[H^{\rm cons}[\eta]M]=\sum_\rho\int\eta^\rho[-\partial_\rho{\rm tr}\,M+\sum_\nu\partial_\nu(M_{\nu\rho}+M_{\rho\nu})]=0$.

The joint cumulant of the quadratic forms $\frac12\int{\rm Tr}(H_aT)$ of Gaussian fields is $\frac12$ times the sum over the cyclic orders of the traces of the products of weights and covariances, times $3$ for the three fields; as the supports are compact and pairwise disjoint, consecutive points of a term lie in different supports, at sup distances bounded below, so every integrand is bounded and compactly supported. In each term the integral over $x_1$ is $\int{\rm tr}[H_1(x_1)M(x_1)]\,dx_1$ with $y_1=x_{\pi2}$, $y_2=x_{\pi q}$ and $B$ the product of the other factors ($K_3$ being even); for $x_2,\dots,x_q$ in the supports of $H_2,\dots,H_q$ the points $y_1,y_2$ lie off ${\rm supp}\,\eta$, so $\eta\in C^2_c(\mathbb R^3\setminus\{y_1,y_2\};\mathbb R^3)$, and at $H_1=H^{\rm cons}[\eta]$ the integral over $x_1$ vanishes term by term.
(b) Under $\nu\otimes\gamma_K$ the functions of $U$ alone have the laws they have under $\nu$, and $\mathsf X_K$ is linear in $\chi_K$ with $\gamma_K(\mathsf X_K)=0$; so the mirrored covariance of $\mathcal T^{\rm X}_K(H^{\rm cons}[\cdot])$ at $\theta\eta$ and $\eta$ is that of $\mathcal T^{\rm c}$ under $\nu$, which tends to $0$ by Corollary~\ref{stc:two}\,(b), plus $\epsilon^2\gamma_K(\mathsf X_K[H^{\rm cons}[\theta\eta]]\mathsf X_K[H^{\rm cons}[\eta]])$. At $\eta=he_3$ the diagonal entries of $H^{\rm cons}[\eta]$ are $\partial_3h\,(1,1,-1)$; its off-diagonal entries, $-\partial_1h$ at $(1,3)$, $(3,1)$ and $-\partial_2h$ at $(2,3)$, $(3,2)$, do not enter $\mathsf X_K$, which reads the entries $H_{\lambda\lambda}$ alone.

So in the limit $K\to\infty$, in which the free lattice field of mass $1$ gives the continuum one, $\mathsf X_K[H^{\rm cons}[\eta]]\to\int\partial_3h\,(\partial_1^2+\partial_2^2-\partial_3^2)\chi=\int h\,(2\partial_3^3-\partial_3\Delta)\chi$. For $f$ supported in $\{x_3>0\}$ the free covariance gives $\int\!\!\int f(\Theta x)G(x-y)f(y)\,dx\,dy=\int|\hat f(p)|^2d\rho(p)$, $\hat f(p)=\int e^{-E(p)x_3+ip\cdot(x_1,x_2)}f(x)\,dx$; with $\partial_3\mapsto-E$ and $\Delta\mapsto E^2-|p|^2=1$ on the mass shell this is the integral of $|\hat h|^2E^2(2E^2-1)^2$. For $0\le h\ne0$, $\hat h(0)=\int he^{-x_3}dx>0$ and $\hat h$ is continuous, so the integral is positive. The added member is no function of the gauge field, and no null term takes its divergence away: the collapse of Lemma~\ref{stc:defect} fails for the dressed field, and with it (N).
\end{proof}

\subsection{Scope of the conclusions}

The bounds are bounds in the cutoff at a fixed distance of the supports, or, in Theorem~\ref{stc:cutoff}\,(iii), at distances that may shrink like $K^2L^{-\lceil\phi K\rceil}$; no mass gap, uniqueness of the vacuum or clustering is used. In every statement the supports of the entries are pairwise disjoint, and a tensor limit is defined only on such tuples. Theorem~\ref{stc:limit}\,(iii) concerns covariances: (N) for one tensor limit gives the vanishing of the mirrored two-point functions of $\mathfrak K$ and $\mathsf{St}$ along its own defining sequence, and nothing on the means $\nu(\mathfrak K)$ apart. The bounds for $\mathsf{St}$, $\mathsf{Rm}$ and $\mathfrak K$ contain no tensor field and give (N) only through the collapse (Lemma~\ref{stc:defect}), as the instance of Proposition~\ref{stc:free}\,(b) shows.

For $\gamma^{\rm ref}_W<g\le\gamma$ tensor limits at $g$ exist along subsequences as well, by the following route. Let $z\ge1$ be the least integer for which the reference couplings $g'_1,g'_2$ of Lemma~\ref{stc:dil}\,(c), $g'^{-2}_1=\Lambda g^{-2}/2$, $g'^{-2}_2=L\Lambda g^{-2}/4$, $\Lambda=L^z$, satisfy $g'_1,g'_2\le\gamma^{\rm ref}_W$; it exists by Lemma~\ref{stc:dil}\,(d) and is a function of $g$, and $D^z$ is the dilation $x\mapsto\Lambda x$ of $\mathbb L_K$ onto $\mathbb L_{K-z}$, $(D^zU)(D^z\mathfrak b)=U(\mathfrak b)$. By Lemma~\ref{stc:dil}\,(a)--(c), $\mathcal T_K(H)=L^z\mathcal T_{K-z}(H^{(z)})\circ D^z$, $D^z_*$ carries periodic Gibbs states to periodic Gibbs states, and every $(K,\beta)\in\mathfrak d(g)$ with $K>z$ has $(K-z,\beta)\in\mathfrak d(g'_1)\cup\mathfrak d(g'_2)$.

So item (ii) of the subsection on the input (Theorem~\ref{sf-gibbs}\,(a)), read at $g'_1$ and at $g'_2$ for the dilated weights, bounds $\nu(\prod_aA^\nu[\mathcal T_K(H_a)])$ by $C\prod_a|H_a|_\infty$ for real bounded weights $H_a$ vanishing off fixed bounded open sets with pairwise disjoint closures, uniformly in $(K,\beta)\in\mathfrak d(g)$ with large $K$ and in $\nu\in\mathcal G^{\rm per}_{K,\beta}$, the number $C$ containing the factor $L^{nz}$. The labels lie in the compact interval $[\gamma^{-2},\mathsf b_+]$ (Lemma~\ref{stc:frame}\,(a)), and the diagonal extraction of the proof of Corollary~\ref{stc:small}\,(i) gives a subsequence along which the labels converge and which defines a tensor limit at $g$.

Conservation is an input of the covariance of tensor limits under all of $O(3)$. For $0<g\le\gamma^{\rm sd}_W$, and by the change of scale of Lemma~\ref{stc:dil} for $0<g\le\gamma^{\rm ref}_W$, every tensor limit $\sigma$ at $g$ has $\sigma(\mathsf RH_1,\mathsf RH_2)=\sigma(H_1,H_2)$ for every $\mathsf R\in O(3)$, with $(\mathsf RH)(x)=\mathsf RH(\mathsf R^{-1}x)\mathsf R^{\rm T}$, and all $H_1,H_2\in C^1_c(\mathbb R^3;V)$ with disjoint supports (the theorem on the tensor limits of Section~\ref{C-371-1} on reflection positivity at every plane, the forward light cone and Wightman fields for the tensor ${\rm tr}\,F_{ij}F_{kl}$, part~(3), at two entries; the seven conditions (h1)--(h7) below are derived for tensor limits in Step~3 of its proof).

The argument is that of Proposition~\ref{C-371-27} of that section, a statement on a real symmetric bilinear functional $\sigma_2$ on the pairs of continuous weights with disjoint supports under seven conditions (h1)--(h7). Let first $0<g\le\gamma^{\rm sd}_W$, and $\sigma_2=\sigma$. (h1) $\sigma$ is invariant under the translations and the cubic group, and (h2) $|\sigma(\theta\bar h,h')|\le c(\mathcal K)|h|_\infty|h'|_\infty$ for $h,h'$ supported in a compact set $\mathcal K\subset\{x_3>0\}$ (Corollary~\ref{stc:small}\,(i)). (h3) In the Hilbert space of Corollary~\ref{stc:small}\,(ii), with $v(h)=[(h)]$, the joint spectral measure $\mathcal P$ of $(H,P_1,P_2)$ is carried by the square cone $\{E\ge\max(|q_1|,|q_2|)\}$, and $e^{-a_3H}e^{i(a_1P_1+a_2P_2)}v(h)=v(\tau_ah)$ for $a_3\ge0$; so the mirrored two-point function is the Laplace transform of a matrix of measures carried by this cone. (h4) $\sigma$ is positive at the diagonal mirror $\{x_1+x_3=0\}$ (Corollary~\ref{stc:small}\,(ii)). (h5) is (C) of Theorem~\ref{stc:limit}\,(ii); in the helicity basis of the rotations about one axis it gives three linear relations between these measures.

The last two conditions are inputs that rest, like the scalar theory, on Theorems~\YMMref{D}{8}, \YMMref{AJ}{673} and~\YMMref{AD}{607}. (h6) The scalar two-point function is invariant under the rotations: by Corollary~\ref{stc:small}\,(i) the trace part $\sigma(f_1\mathbf 1,f_2\mathbf 1)$ is the limit, along the defining sequence of $\sigma$, of the centred products of the action density fields, that is the scalar limit of this sequence, and a scalar limit is invariant under all isometries of $\mathbb R^3$ (Theorem~\textup{\YMMref{AP}{780}}). (h7) A bound $\sigma(\theta h,h)\le Cd^{-a}|h|_\infty^2$ with $a<4$ at the off-diagonal bumps $h$ of size $d$ near the mirror; here it holds with $a=0$: there is $C_\infty(g)$, a function of $g$, such that $\sigma(\theta h,h)\le C_\infty(g)\,b_*^{-2}|h|_\infty^2$, $b_*$ the limit of the labels along a defining sequence of $\sigma$ with convergent labels (Corollary~\ref{stc:small}\,(i)), for $0<d<1$ and every real $h=h_0(E_{\mu\nu}+E_{\nu\mu})$, $E_{\mu\nu}$ the matrix units, $\mu\ne\nu$, with $h_0$ supported in the ball of radius $d$ about a point $x_0$ with $x_{0,3}\in[\frac32d,2d]$ (Proposition~\ref{C-371-24}\,(e), stated for $0<g\le\gamma^{\rm sd}_W$ alone, and recalled in the subsection on the input).

It is the bound at short distance of the covariance of two off-diagonal bond fields, whose means vanish, in the periodic Gibbs states at the first-passage data at $g$, with a number that depends neither on the scale $d$ nor on the cutoff, carried to the limit: on the least torus at $g$ the renormalisation group of Theorem~\YMMref{D}{8} gives it at the scale $d$, and the chessboard estimate with the bound of Theorems~\YMMref{AJ}{673} and~\YMMref{AD}{607} on differences of pressures carries it to every periodic Gibbs state. It is had for $0<g\le\gamma^{\rm sd}_W$, and this sets the first range. With (h1)--(h7) only the harmonics of a covariant kernel remain: the imaginary rotations about $e_2$ act on the measures of (h3) as the boosts of $(E,q_1)$; by (h7) and conservation the boosted measures have finitely many harmonics, of orders at most $6+a<10$; (h6) makes the measure of the trace invariant under the boosts; a wrong harmonic of odd order is positive and negative by (h4), so it vanishes, and one of even order is emptied by (h5) and by its scaling under the boosts; so the kernel is covariant under the rotations about $e_2$, which generate $O(3)$ with the cubic group.

For $\gamma^{\rm sd}_W<g\le\gamma^{\rm ref}_W$ take $z$ with $g'_1,g'_2\le\gamma^{\rm sd}_W$ (Lemma~\ref{stc:dil}\,(d)), and pass to a subsequence of the defining sequence $(K_j,\beta_j,\rho_j)$ of $\sigma$ along which the index $\iota$ of Lemma~\ref{stc:dil}\,(c) is one and the same; a subsequence defines the same $\sigma$. By Lemma~\ref{stc:dil}\,(a)--(c), $(K_j-z,\beta_j,D^z_*\rho_j)$ defines the tensor limit $\sigma_z(H'_1,\dots,H'_n)=\Lambda^{-n}\sigma(H'_1(\Lambda\,\cdot),\dots,H'_n(\Lambda\,\cdot))$ at $g'_\iota\le\gamma^{\rm sd}_W$, so $\sigma(H_1,H_2)=\Lambda^2\sigma_z(H_1^{(z)},H_2^{(z)})$; as $(\mathsf RH)^{(z)}=\mathsf RH^{(z)}$, the covariance of $\sigma_z$ gives that of $\sigma$.

\section{Point plates and holomorphic kernels}\label{C-353-1}

We construct bounded point-plate operators from reflection-positive functionals with a finite slab-insertion bound. The construction first permits an arbitrary bound as a function of the slab width. A power bound gives quantitative growth of the holomorphic kernels and point fields.

The tensor indices are part of the reflection and covariance conventions. We give the operator construction, the gluing of quarter-turned tubes and the converses on the closed spans of the resulting vectors. Gaussian examples show why a finite insertion bound alone need not be a power bound. The final application identifies the precise lattice estimates required for curvature bilinears.

\subsection{Notation and hypotheses}

We work in $\mathbb R^3\ni x=(x_1,x_2,x_3)$ with the standard basis $e_1,e_2,e_3$; the coordinate $x_3$ is Euclidean time, and all distances are Euclidean. Scalar products are antilinear in the left argument, and products of operators are taken in the order of their indices. $M\ge1$ is the constant in the norm $\|\cdot\|_\sharp$ of Definition~\ref{C-353-2}(2); only $M\ge1$ is used here. The reference coupling $g$, its two ceilings $\gamma^{\rm sd}_W\le\gamma^{\rm ref}_W$, first-passage data, periodic Gibbs states, the spacing $a_K$ of the lattice at the cutoff $K$ and the constants $L$ and $\bar B$ are the construction data fixed in Section~\ref{C-355-1}; they enter Theorem~\ref{C-353-14} alone.

\begin{definition}[Values, functions, isometries]\label{C-353-2}
\begin{enumerate}
\item[(1)] $V={\rm Sym}_3(\mathbb R)$ and $V_{\mathbb C}={\rm Sym}_3(\mathbb C)$ carry the norm $|E|_2=(\operatorname{tr}E^*E)^{1/2}$; besides, $|E|_\infty:=\max_{\lambda,\lambda'}|E_{\lambda\lambda'}|\le|E|_2$. Every $E\in V_{\mathbb C}$ is $E=E^0+iE^1$ with $E^0=\frac12(E+\bar E)$ and $E^1=\frac1{2i}(E-\bar E)$ in $V$, and $|E|_2^2=|E^0|_2^2+|E^1|_2^2$. For $r\in{\rm O}(3)$ we put $\varrho(r)E:=rEr^{\mathsf T}$; then $|\varrho(r)E|_2=|E|_2$ and $\overline{\varrho(r)E}=\varrho(r)\bar E$.
\item[(2)] For $O\subset\mathbb R^N$ open and $\mathbb K\in\{V,V_{\mathbb C},\mathbb C\}$, $C^1_c(O;\mathbb K)$ is the space of the $C^1$ maps $\mathbb R^N\to\mathbb K$ whose support is compact and contained in $O$, and $C_c(O;\mathbb K)$ the same with continuous maps. With $|\cdot|_2$ on the values, $\|f\|_\infty=\sup|f(x)|_2$, ${\rm Lip}\,f=\sup_{x\ne x'}|f(x)-f(x')|_2/|x-x'|$, $|\nabla f|^2=\sum_\mu|\partial_\mu f|_2^2$, $\|f\|_{C^1}=\max\{\|f\|_\infty,\sup|\nabla f|\}$ and $\|f\|_{L^1}=\int|f(x)|_2\,dx$; at $N=3$, $\|f\|_\sharp=\max\{\|f\|_\infty,40M\,{\rm Lip}\,f\}$. The functions $f^0$, $f^1$, $\bar f$ are taken pointwise, and $f$ is real if its values lie in $V$ (or in $\mathbb R$).
\item[(3)] An isometry $r\colon x\mapsto r^\circ x+v$ of $\mathbb R^3$ acts on functions by the tensor action $(rf)(x)=\varrho(r^\circ)f(r^{-1}x)$. We denote by $\tau_v$ the translation by $v$, by $\theta_c$ the reflection $x\mapsto x-2(x_3-c)e_3$, by $\theta=\theta_0$ both this reflection and the matrix ${\rm diag}(1,1,-1)$, so that $(\theta_cf)(x)=\varrho(\theta)f(\theta_cx)$; $E_+=\{x_3>0\}$; $W_3$ is the group of the $48$ signed permutation matrices, $\mathbb R^3\rtimes W_3$ the group of the isometries $x\mapsto wx+v$, $w\in W_3$, and $B(y,r)$ the open ball of radius $r$ about $y$.
\end{enumerate}
\end{definition}

\begin{definition}[Tuples, complex members, tensor functionals]\label{C-353-3}
\begin{enumerate}
\item[(1)] For $O\subset\mathbb R^3$ open, $\mathfrak S^V(O)$ consists of $\emptyset$ and of the tuples $X=(H_1,\dots,H_n)$ of functions $H_a\in C^1_c(O;V)$ with pairwise disjoint supports; the function $0$ is admitted as a member. $X$ is of type $(\mathsf d;R)$ if each support lies in a closed ball of radius $R$ and any two supports are at a distance $\ge\mathsf d$. Isometries act member by member; $X\cdot Y$ is the concatenation, formed when no support of a member of $X$ meets one of $Y$. $\mathbb C\mathfrak S^V(O)$ is the free complex vector space over $\mathfrak S^V(O)$; for $F=\sum\zeta_iX_i$ we put $\bar F=\sum\bar\zeta_iX_i$, and isometries and concatenation are extended linearly and bilinearly.
\item[(2)] For $H_1,\dots,H_n\in C^1_c(\mathbb R^3;V_{\mathbb C})$ with pairwise disjoint supports,
\[(H_1,\dots,H_n)_{\mathbb C}:=\sum_\varepsilon i^{|\varepsilon|}(H_1^{\varepsilon_1},\dots,H_n^{\varepsilon_n})\in\mathbb C\mathfrak S^V(\mathbb R^3),\qquad|\varepsilon|=\varepsilon_1+\dots+\varepsilon_n,\]
the sum running over the $\varepsilon\in\{0,1\}^n$ with $H_a^1\ne0$ wherever $\varepsilon_a=1$; thus $(H)_{\mathbb C}=(H)$ for real $H$.
\item[(3)] A tensor functional is a map $\sigma\colon\mathfrak S^V(\mathbb R^3)\to\mathbb R$ with $\sigma(\emptyset)=1$ which is $n$-linear on $\prod_aC^1_c(O_a;V)$ whenever $O_1,\dots,O_n\subset\mathbb R^3$ are open sets with pairwise disjoint closures; it is extended linearly to $\mathbb C\mathfrak S^V(\mathbb R^3)$. It is centred if $\sigma((H))=0$ for every $H\in C^1_c(\mathbb R^3;V)$.
\end{enumerate}
\end{definition}

\begin{definition}[The class]\label{C-353-4}
Let $\sigma$ be a tensor functional.
\begin{itemize}
\item[(T1$^\circ$)] $\sigma$ is symmetric and $\sigma(\tau_vX)=\sigma(X)$ for all $v\in\mathbb R^3$ and $X$.
\item[(T1)] $\sigma$ is symmetric and $\sigma(rX)=\sigma(X)$ for all $r\in\mathbb R^3\rtimes W_3$ and $X$.
\item[(T2)] There are numbers $C^\sigma_n(\mathsf d;R)\in[0,\infty)$, one for each $n\ge1$, $\mathsf d>0$, $R>0$, with $|\sigma(H_1,\dots,H_n)|\le C^\sigma_n(\mathsf d;R)\prod_a\|H_a\|_\sharp$ at all tuples of type $(\mathsf d;R)$.
\item[(T3)] $\sigma(\theta\bar F\cdot F)\ge0$ for all $F\in\mathbb C\mathfrak S^V(E_+)$.
\end{itemize}
For $\sigma$ with (T3), $\langle F,G\rangle_\sigma:=\sigma(\theta\bar F\cdot G)$ is a positive semi-definite Hermitian form on $\mathbb C\mathfrak S^V(E_+)$ (Lemma~\ref{C-353-19}); $\mathcal H(\sigma)$ is the completion of the quotient by its null space, $[F]$ is the class of $F$, and $\Omega:=[\emptyset]$. For $\sigma$ with (T1$^\circ$), (T2) and (T3), Lemma~\ref{C-353-19} provides the operators $\mathcal T_t[F]=[\tau_{te_3}F]$, $t\ge0$, and $\mathcal U(v)[F]=[\tau_vF]$, $v\in e_3^\perp$, and their joint spectral measure $\mathcal P$ on $[0,\infty)\times e_3^\perp\ni(E,q)$. In these terms:
\begin{itemize}
\item[(T4)] $\mathcal P(\{(E,q):E<|q_1|\})=0$.
\item[(T5$_N$)] For a number $c_I\in(0,\frac1{16}]$ and a function $N\colon(0,1]\to(0,\infty)$, which may be any function: for all $\eta\in(0,1]$, all real $H\in C^1_c(\mathbb R^3;V)$ with ${\rm supp}\,H\subset B(y,c_I\eta)$ for some $y$ with $|y_3|\le\eta/8$, and all $F,G\in\mathbb C\mathfrak S^V(\{x_3>\eta/4\})$,
\[|\sigma(\theta\bar F\cdot(H)\cdot G)|\le N(\eta)\,\|H\|_\sharp\,\|[\tau_{-(\eta/4)e_3}F]\|\,\|[\tau_{-(\eta/4)e_3}G]\|.\]
\item[(T5)] The insertion bound: (T5$_N$) holds with $N(\eta)=C_I\eta^{-\mathsf p_I}$ for numbers $C_I>0$ and $\mathsf p_I\ge0$.
\end{itemize}
The class is the set of the tensor functionals with (T1), (T2), (T3), (T4) and (T5$_N$), for some $c_I\in(0,\frac1{16}]$ and some function $N$ as in (T5$_N$); a member of the class has the insertion bound with a power if it has (T5) besides. (T1$^\circ$), which follows from (T1), takes the place of (T1) in (A)--(C) of Theorem~\ref{C-353-7} and in Theorem~\ref{C-353-8}. We say that $\sigma$ has (T5$_N$), or (T5), in the sup norm if the bound holds with $\|H\|_\infty$ in place of $\|H\|_\sharp$; as $\|H\|_\infty\le\|H\|_\sharp$, it then has (T5$_N$), or (T5), with the same numbers. In (T2) the numbers $C^*_n(\mathsf d;R):=\sup|\sigma(H_1,\dots,H_n)|/\prod_a\|H_a\|_\sharp$, the supremum over the tuples of type $(\mathsf d;R)$ without a member $0$, serve as well, and they do not increase with $\mathsf d$; so the class does not change if one asks in (T2) for numbers that do not increase with $\mathsf d$.
\end{definition}

\begin{definition}[The plane, the tubes, labels, profiles]\label{C-353-5}
\begin{enumerate}
\item[(1)] We put $x=(y,x_2)$ with $y=(y^0,y^1)=(x_3,x_1)$, $n_0=(1,0)$, $\tau_y:=\tau_{y^0e_3+y^1e_1}$, $\theta y=(-y^0,y^1)$ and $\mathbb R^2_+=\{y^0>0\}$. A point $\xi=(\xi^0,\xi^1)\in\mathbb C^2$ is identified with $(\xi^1,0,\xi^0)\in\mathbb C^3$.
\item[(2)] $d_{\rm Tb}(\xi)={\rm Re}\,\xi^0-|{\rm Im}\,\xi^1|$, ${\rm Tb}=\{d_{\rm Tb}>0\}$, $\xi^*=(\overline{\xi^0},-\overline{\xi^1})$, $\xi^\pm=\xi^0\pm i\xi^1$, and ${\rm Tb}^{(n)}=\{x\in(\mathbb C^2)^n:x_1\in{\rm Tb},\ x_{k+1}-x_k\in{\rm Tb}\}$.
\item[(3)] $d_3(\mathfrak z)={\rm Re}\,\mathfrak z_3-|{\rm Im}\,\mathfrak z_1|-|{\rm Im}\,\mathfrak z_2|$, ${\rm Tb}_3=\{d_3>0\}$, $\mathfrak z^*=(-\bar{\mathfrak z}_1,-\bar{\mathfrak z}_2,\bar{\mathfrak z}_3)$, and ${\rm Tb}_3^{(n)}$ is formed as in (2).
\item[(4)] $R_t$, $t\in\mathbb R$, is the rotation about $e_2$ with $R_te_3=\cos t\,e_3+\sin t\,e_1$ and $R_te_1=-\sin t\,e_3+\cos t\,e_1$. On $\mathbb C^2$ and for $\zeta\in\mathbb C$, $R_\zeta$ is the linear map with $(R_\zeta\xi)^\pm=e^{\pm i\zeta}\xi^\pm$; for real $\zeta$ it is the rotation $R_\zeta$ acting on $y$. We put $n_\upsilon:=R_\upsilon n_0=(\cos\upsilon,\sin\upsilon)$.
\item[(5)] A label is an $\mathfrak a\in\mathcal L_{\mathbb C}:=C^1_c(\mathbb R;V_{\mathbb C})$; it is real if it is $V$-valued. We put $\mathfrak a^\theta:=\varrho(\theta)\bar{\mathfrak a}$; $\lambda_{\mathfrak a}$ is the length of the smallest closed interval that contains ${\rm supp}\,\mathfrak a$ ($\lambda_0=0$); $\nu(\mathfrak a):=(\lambda_{\mathfrak a}+1)\|\mathfrak a\|_{C^1}$ and $\nu_\infty(\mathfrak a):=(\lambda_{\mathfrak a}+1)\|\mathfrak a\|_\infty$. For $r\in{\rm O}(3)$, $\varrho(r)\mathfrak a$ is taken pointwise; for a tuple $\mathbf a=(\mathfrak a_1,\dots,\mathfrak a_n)$ of labels, $\bar{\mathbf a}$ and $\varrho(r)\mathbf a$ are taken label by label.
\item[(6)] A profile is an $\alpha\in C^1_c(\mathbb R^2;\mathbb C)$, and $(\alpha\otimes\mathfrak a)(x):=\alpha(y)\mathfrak a(x_2)$. Profiles $\alpha_1,\dots,\alpha_n$ have time-ordered supports if ${\rm supp}\,\alpha_k\subset\{g_{k-1}<y^0<g_k\}$ for real numbers $g_0<\dots<g_n$, and are time-ordered in $\mathbb R^2_+$ if moreover $g_0\ge0$; then $\sigma(\alpha_1\otimes\mathfrak a_1,\dots,\alpha_n\otimes\mathfrak a_n):=\sigma((\alpha_1\otimes\mathfrak a_1,\dots,\alpha_n\otimes\mathfrak a_n)_{\mathbb C})$. We put $Y=(y_1,\dots,y_n)$ and $dY=dy_1\cdots dy_n$.
\item[(7)] Helicity. $\epsilon_\pm:=e_3\pm ie_1$; $E_{\pm2}:=\epsilon_\pm\epsilon_\pm^{\mathsf T}$, $E_{\pm1}:=\frac12(e_2\epsilon_\pm^{\mathsf T}+\epsilon_\pm e_2^{\mathsf T})$, $E_0:=e_2e_2^{\mathsf T}$ and $E_{0'}:=\frac12(\epsilon_+\epsilon_-^{\mathsf T}+\epsilon_-\epsilon_+^{\mathsf T})=e_3e_3^{\mathsf T}+e_1e_1^{\mathsf T}$ form a basis of $V_{\mathbb C}$; ${\rm hel}(\pm2)=\pm2$, ${\rm hel}(\pm1)=\pm1$, ${\rm hel}(0)={\rm hel}(0')=0$. A helicity label is a label $\varphi E_s$ with a real $\varphi\in C^1_c(\mathbb R)$; for a tuple $\mathbf a$ of helicity labels $\varphi_kE_{s_k}$ we put $S(\mathbf a):=\sum_k{\rm hel}(s_k)$.
\end{enumerate}
\end{definition}

\begin{definition}[Operators and numbers]\label{C-353-6}
Let $\sigma$ have (T1$^\circ$), (T2) and (T3).
\begin{enumerate}
\item[(1)] Under (T4): $T(\xi):=\int_{\{E\ge|q_1|\}}e^{-\xi^0E+i\xi^1q_1}\,d\mathcal P(E,q)$ for $\xi\in\overline{\rm Tb}$.
\item[(2)] Under (T1) and (T4): $T_3(\mathfrak z):=\int_{\{E\ge|q_1|,\ E\ge|q_2|\}}e^{-\mathfrak z_3E+i\mathfrak z_1q_1+i\mathfrak z_2q_2}\,d\mathcal P(E,q)$ for $\mathfrak z\in\overline{{\rm Tb}_3}$.
\item[(3)] Under (T5$_N$): for $d>0$ and $f\in C^1_c(\{|x_3|<d\};V_{\mathbb C})$, the centred slab operator $B^c_d(f)$ is the bounded operator with $B^c_d(f)[F]=[(\tau_{de_3}f)_{\mathbb C}\cdot\tau_{2de_3}F]$ for $F\in\mathbb C\mathfrak S^V(E_+)$ (Lemmas~\ref{C-353-21} and~\ref{C-353-22}); then $\langle[\tau_{-de_3}F],B^c_d(f)[\tau_{-de_3}F']\rangle=\sigma(\theta\bar F\cdot(f)_{\mathbb C}\cdot F')$ for $F,F'\in\mathbb C\mathfrak S^V(\{x_3>d\})$.
\item[(4)] ${\rm mv}_0=9\pi/(5\pi^2-8)$, ${\rm mv}_1=\frac{3\pi}4{\rm mv}_0$, ${\rm mv}_2=\sqrt2\,{\rm mv}_0{\rm mv}_1+2{\rm mv}_0^2$ and ${\rm mv}_3={\rm mv}_0^3+\sqrt3\,{\rm mv}_0^2{\rm mv}_1$ are the construction constants; $c_{\rm plt}:=480\cdot8^3{\rm mv}_2c_I^{-4}M$ and $c_{\rm pt}:=320\cdot8^4{\rm mv}_3c_I^{-4}M$; under (T5), $\wp:=\mathsf p_I+4$, $\wp':=\mathsf p_I+3$, $C_{\rm plt}:=c_{\rm plt}C_I$, $C_{\rm pt}:=c_{\rm pt}C_I$, $C_\flat:=2^{\wp+1}C_{\rm plt}$ and $C'_\flat:=2^{\wp'+1}\cdot12\cdot8^2{\rm mv}_0^2c_I^{-3}C_I$.
\end{enumerate}
\end{definition}

\subsection{Principal results}

\begin{theorem}[Point plates, kernels and vectors with matrix labels; point fields with a matrix index]\label{C-353-7}
Let $\sigma$ be a tensor functional. The numbers of Definition~\ref{C-353-6} are formed with the constants $c_I$ and $N$, or $c_I$, $C_I$ and $\mathsf p_I$, with which $\sigma$ has (T5$_N$), or (T5); none of them depends on $n$, on the labels or on the numbers of (T2).
\begin{enumerate}
\item[(A)] \emph{Point plates.} Suppose that $\sigma$ has (T1$^\circ$), (T2), (T3), (T4) and (T5$_N$), with $N$ arbitrary. Then there are bounded operators $B_{\mathfrak a}(\xi,\xi')$ on $\mathcal H(\sigma)$, for $\mathfrak a\in\mathcal L_{\mathbb C}$ and $\xi,\xi'\in{\rm Tb}$, holomorphic in norm in $(\xi,\xi')$ and complex linear in $\mathfrak a$, such that:
\begin{enumerate}
\item[(i)] $\|B_{\mathfrak a}(\xi,\xi')\|\le2c_{\rm plt}\,\nu(\mathfrak a)\,N(\eta)\eta^{-4}$ for every $\eta\in(0,\min\{1,d_{\rm Tb}(\xi),d_{\rm Tb}(\xi')\}]$, and the factor $2$ may be omitted for real $\mathfrak a$;
\item[(ii)] $T(\kappa)B_{\mathfrak a}(\xi,\xi')T(\kappa')=B_{\mathfrak a}(\kappa+\xi,\xi'+\kappa')$ for $\kappa,\kappa'\in\overline{\rm Tb}$, and $B_{\mathfrak a}(\xi,\xi')^*=B_{\mathfrak a^\theta}(\xi'^*,\xi^*)$;
\item[(iii)] $B^c_d(\alpha\otimes\mathfrak a)=\int\alpha(y)\,B_{\mathfrak a}(dn_0+y,dn_0-y)\,dy$ for $d>0$ and profiles $\alpha$ with support in $\{|y^0|<d\}$; if ${\rm supp}\,\alpha\subset\{|y^0|\le d'\}$ with $0\le d'<d$, then $\|B^c_d(\alpha\otimes\mathfrak a)\|\le2c_{\rm plt}\,\nu(\mathfrak a)\,N(\eta)\eta^{-4}\|\alpha\|_{L^1}$ for every $\eta\in(0,\min\{1,d-d'\}]$;
\item[(iv)] a weakly continuous family of bounded operators on ${\rm Tb}\times{\rm Tb}$ that has the shift rule of (ii) and satisfies the identity of (iii) in matrix elements, for all $d>0$ and all real profiles with support in $\{|y^0|<d\}$, is the family $B_{\mathfrak a}$; in particular $B_{\mathfrak a}$ depends on $\sigma$ and $\mathfrak a$ alone, not on $c_I$ or $N$;
\item[(v)] for $\xi,\xi',\kappa_1,\kappa_2\in{\rm Tb}$, $\kappa_0\in\overline{\rm Tb}$ and $\mathfrak a'\in\mathcal L_{\mathbb C}$, the operator $B_{\mathfrak a}(\xi,\kappa_1)T(\kappa_0)B_{\mathfrak a'}(\kappa_2,\xi')$ depends on $(\kappa_1,\kappa_0,\kappa_2)$ only through $\varkappa=\kappa_1+\kappa_0+\kappa_2$ and equals $B_{\mathfrak a}(\xi,\frac\varkappa2)B_{\mathfrak a'}(\frac\varkappa2,\xi')$; the vector $B_{\mathfrak a}(\xi,\kappa)\Omega$ and the numbers $\langle\Omega,B_{\mathfrak a}(\kappa,\xi)\psi\rangle$, $\psi\in\mathcal H(\sigma)$, do not depend on $\kappa\in{\rm Tb}$.
\end{enumerate}
\item[(B)] \emph{Kernels and vectors.} Let $\sigma$ be as in (A), $n\ge1$, $\mathbf a=(\mathfrak a_1,\dots,\mathfrak a_n)\in\mathcal L_{\mathbb C}^n$ and $\kappa\in{\rm Tb}$. Put
\[K_n(\xi;\mathbf a):=\langle\Omega,B_{\mathfrak a_1}(\kappa,\tfrac{\xi_1}2)B_{\mathfrak a_2}(\tfrac{\xi_1}2,\tfrac{\xi_2}2)\cdots B_{\mathfrak a_n}(\tfrac{\xi_{n-1}}2,\kappa)\Omega\rangle,\qquad\xi\in{\rm Tb}^{n-1},\]
\[\Psi(x;\mathbf a):=B_{\mathfrak a_1}(x_1,\tfrac{x_2-x_1}2)B_{\mathfrak a_2}(\tfrac{x_2-x_1}2,\tfrac{x_3-x_2}2)\cdots B_{\mathfrak a_n}(\tfrac{x_n-x_{n-1}}2,\kappa)\Omega,\qquad x\in{\rm Tb}^{(n)},\]
with $K_1(\mathfrak a_1)=\langle\Omega,B_{\mathfrak a_1}(\kappa,\kappa)\Omega\rangle$ and $\Psi(x_1;\mathfrak a_1)=B_{\mathfrak a_1}(x_1,\kappa)\Omega$; we call the $\Psi(x;\mathbf a)$ plate vectors. They do not depend on $\kappa$, they are complex linear in each label, and:
\begin{enumerate}
\item[(KV1$_N$)] $K_n(\cdot;\mathbf a)$ is holomorphic on ${\rm Tb}^{n-1}$, and $|K_n(\xi;\mathbf a)|\le\prod_k\bigl(2c_{\rm plt}\,\nu(\mathfrak a_k)\bigr)\cdot(N(\eta)\eta^{-4})^n$ for every $\eta\in(0,\min\{1,\frac12\min_kd_{\rm Tb}(\xi_k)\}]$; for $n=1$, for every $\eta\in(0,1]$;
\item[(KV2)] if the profiles $\alpha_1,\dots,\alpha_n$ have time-ordered supports, then
\[\sigma(\alpha_1\otimes\mathfrak a_1,\dots,\alpha_n\otimes\mathfrak a_n)=\int\prod_k\alpha_k(y_k)\,K_n(y_2-y_1,\dots,y_n-y_{n-1};\mathbf a)\,dY,\]
and if they are time-ordered in $\mathbb R^2_+$, then 
\[
[(\alpha_1\otimes\mathfrak a_1,\dots,\alpha_n\otimes\mathfrak a_n)_{\mathbb C}]=\int\prod_k\alpha_k(y_k)\,\Psi(Y;\mathbf a)\,dY;
\]

\item[(KV3)] $\Psi(\cdot;\mathbf a)$ is holomorphic in norm on ${\rm Tb}^{(n)}$; $T(\kappa')\Psi(x;\mathbf a)=\Psi(x+(\kappa',\dots,\kappa');\mathbf a)$ for $\kappa'\in\overline{\rm Tb}$; 
\[
\langle\Omega,\Psi(z;\mathbf a')\rangle=K_{n'}(z_2-z_1,\dots,z_{n'}-z_{n'-1};\mathbf a')
\]
 and 
\[
\langle\Psi(x;\mathbf a),\Omega\rangle=K_n((x_n-x_{n-1})^*,\dots,(x_2-x_1)^*;\mathfrak a_n^\theta,\dots,\mathfrak a_1^\theta);
\]
and, for $\mathbf a'\in\mathcal L_{\mathbb C}^{n'}$ and $z\in{\rm Tb}^{(n')}$,
\[
\begin{multlined}
\langle\Psi(x;\mathbf a),\Psi(z;\mathbf a')\rangle=K_{n+n'}\bigl((x_n-x_{n-1})^*,\dots,(x_2-x_1)^*,x_1^*+z_1,z_2-z_1,\dots,z_{n'} \\
-z_{n'-1};\mathfrak a_n^\theta,\dots,\mathfrak a_1^\theta,\mathfrak a'_1,\dots,\mathfrak a'_{n'}\bigr);
\end{multlined}
\]
\item[(KV4)] $\overline{K_n(\xi;\mathbf a)}=K_n(\bar\xi;\bar{\mathbf a})$, the bars taken coordinate by coordinate and label by label; if $\sigma$ is centred, then $K_1=0$;
\item[(KV5)] a continuous function on $(\mathbb R^2_+)^{n-1}$ that satisfies the first formula of (KV2) for all real profiles with time-ordered supports is $K_n(\cdot;\mathbf a)$ there; hence $K_n(\cdot;\mathbf a)$ is the only holomorphic function on ${\rm Tb}^{n-1}$ with that formula, $\Psi(\cdot;\mathbf a)$ is the only map on ${\rm Tb}^{(n)}$, holomorphic in norm, with the second formula, and both depend on $\sigma$ and $\mathbf a$ alone.
\end{enumerate}
\item[(C)] \emph{The power.} Suppose that $\sigma$ has (T1$^\circ$), (T2), (T3), (T4) and (T5). Then, besides (B), for all $n\ge1$, $\mathbf a\in\mathcal L_{\mathbb C}^n$ and $\xi\in{\rm Tb}^{n-1}$,
\[\text{(KV1)}\qquad|K_n(\xi;\mathbf a)|\le\prod_k\bigl(C_\flat\,\nu(\mathfrak a_k)\bigr)\cdot\min\{1,\min_kd_{\rm Tb}(\xi_k)\}^{-\wp n},\]
the right side being $C_\flat\nu(\mathfrak a_1)$ for $n=1$, and with $2^\wp C_{\rm plt}$ in place of $C_\flat$ if all labels are real. More precisely, $|K_n(\xi;\mathbf a)|\le\prod_{k=1}^n\bigl(2C_{\rm plt}\,\nu(\mathfrak a_k)\min\{1,\frac12d_{\rm Tb}(\xi_{k-1}),\frac12d_{\rm Tb}(\xi_k)\}^{-\wp}\bigr)$, the terms with $\xi_0$ and $\xi_n$ left out, and $\|B_{\mathfrak a}(\xi,\xi')\|\le2C_{\rm plt}\,\nu(\mathfrak a)\min\{1,d_{\rm Tb}(\xi),d_{\rm Tb}(\xi')\}^{-\wp}$. If $\sigma$ has (T5) in the sup norm, then (KV1) holds also with $\wp'$, $\nu_\infty(\mathfrak a_k)$ and $C'_\flat$ in place of $\wp$, $\nu(\mathfrak a_k)$ and $C_\flat$.
\item[(D)] \emph{Point fields.} Suppose that $\sigma$ has (T1), (T2), (T3), (T4) and (T5$_N$). Then $\sigma$ is reflection positive at each of the mirrors $\{x_\mu=0\}$, $\mu=1,2,3$; $\mathcal P$ is carried by $\{E\ge|q_1|,\ E\ge|q_2|\}$; and $T_3$ is a strongly continuous family of contractions on $\overline{{\rm Tb}_3}$, holomorphic in norm on ${\rm Tb}_3$, with $T_3(\mathfrak z)T_3(\mathfrak z')=T_3(\mathfrak z+\mathfrak z')$, $T_3(\mathfrak z)^*=T_3(\mathfrak z^*)$, $T_3(\mathfrak z)\Omega=\Omega$, $T_3(x)[F]=[\tau_xF]$ for real $x$ with $x_3\ge0$, and $T(\xi)=T_3(\xi^1,0,\xi^0)$. There are bounded operators $B_E(\mathfrak z,\mathfrak z')$ on $\mathcal H(\sigma)$, for $E\in V_{\mathbb C}$ and $\mathfrak z,\mathfrak z'\in{\rm Tb}_3$, holomorphic in norm in $(\mathfrak z,\mathfrak z')$ and complex linear in $E$, such that:
\begin{enumerate}
\item[(i)] $\|B_E(\mathfrak z,\mathfrak z')\|\le2c_{\rm pt}|E|_2N(\eta)\eta^{-4}$ for every $\eta\in(0,\min\{1,d_3(\mathfrak z),d_3(\mathfrak z')\}]$, and the factor $2$ may be omitted for $E\in V$;
\item[(ii)] $T_3(\kappa)B_E(\mathfrak z,\mathfrak z')T_3(\kappa')=B_E(\kappa+\mathfrak z,\mathfrak z'+\kappa')$ for $\kappa,\kappa'\in\overline{{\rm Tb}_3}$, and $B_E(\mathfrak z,\mathfrak z')^*=B_{\varrho(\theta)\bar E}(\mathfrak z'^*,\mathfrak z^*)$;
\item[(iii)] $B^c_d(H)=\int B_{H(x)}(de_3+x,de_3-x)\,dx$ for $d>0$ and $H\in C^1_c(\{|x_3|<d\};V_{\mathbb C})$;
\item[(iv)] for each $E$, a weakly continuous family of bounded operators on ${\rm Tb}_3\times{\rm Tb}_3$ that has the shift rule of (ii) and satisfies (iii) in matrix elements at $H=hE$, for all $d>0$ and all real $h\in C^1_c(\{|x_3|<d\};\mathbb R)$, is the family $B_E$; for $\kappa_1,\kappa_2\in{\rm Tb}_3$ and $\kappa_0\in\overline{{\rm Tb}_3}$ the operator $B_E(\mathfrak z,\kappa_1)T_3(\kappa_0)B_{E'}(\kappa_2,\mathfrak z')$ depends on $(\kappa_1,\kappa_0,\kappa_2)$ only through $\varkappa=\kappa_1+\kappa_0+\kappa_2$ and equals $B_E(\mathfrak z,\frac\varkappa2)B_{E'}(\frac\varkappa2,\mathfrak z')$; and $B_E(\mathfrak z,\kappa)\Omega$ and $\langle\Omega,B_E(\kappa,\mathfrak z)\psi\rangle$ do not depend on $\kappa\in{\rm Tb}_3$.
\end{enumerate}
\item[(E)] \emph{Plates from point fields; point bound and insertion bound.} Let $\sigma$ be as in (D).
\begin{enumerate}
\item[(i)] $B_{\mathfrak a}(\xi,\xi')=\int B_{\mathfrak a(s)}(\xi+se_2,\xi'-se_2)\,ds$ for $\mathfrak a\in\mathcal L_{\mathbb C}$ and $\xi,\xi'\in{\rm Tb}$; hence $\|B_{\mathfrak a}(\xi,\xi')\|\le2c_{\rm pt}\|\mathfrak a\|_{L^1}N(\eta)\eta^{-4}$ for $\eta$ as in (A)(i). For $H\in C^1_c(\{|x_3|<d\};V_{\mathbb C})$ with support in $\{|x_3|\le d'\}$, $0\le d'<d$, and for $\eta\in(0,\min\{1,d-d'\}]$, $\|B^c_d(H)\|\le2c_{\rm pt}N(\eta)\eta^{-4}\|H\|_{L^1}$ and $\|B^c_d(H)\|\le c_{\rm pt}N(\eta)\eta^{-4}\sum_c\|\operatorname{tr}(b_cH)\|_{L^1}$ for every basis $(b_c)$ of $V$ orthonormal for $(E,E')\mapsto\operatorname{tr}EE'$. If $\sigma$ has (T5), then (KV1) holds also with $2^{\wp+1}C_{\rm pt}\|\mathfrak a_k\|_{L^1}$ in place of $C_\flat\nu(\mathfrak a_k)$.
\item[(ii)] Let $\mathsf C_{\rm pt}(\eta):=\sup\{\|B_E(\mathfrak z,\mathfrak z')\|:E\in V,\ |E|_2=1,\ d_3(\mathfrak z),d_3(\mathfrak z')\ge\eta\}$ for $\eta\in(0,1]$. Then $\mathsf C_{\rm pt}(\eta)\le c_{\rm pt}N(\eta)\eta^{-4}$, and $\mathsf C_{\rm pt}(\eta)\le8^4{\rm mv}_0^3c_I^{-3}N(\eta)\eta^{-3}$ if $\sigma$ has (T5$_N$) in the sup norm. For $\eta$, $H$, $F$, $G$ as in (T5$_N$),
\[|\sigma(\theta\bar F\cdot(H)\cdot G)|\le\mathsf C_{\rm pt}(\tfrac\eta{16})\,\|H\|_{L^1}\,\|[\tau_{-(\eta/4)e_3}F]\|\,\|[\tau_{-(\eta/4)e_3}G]\|,\qquad\|H\|_{L^1}\le\tfrac{4\pi}3(c_I\eta)^3\|H\|_\infty.\]
Consequently, if $\sigma$ has (T5) in the sup norm with $C_I$ and $\mathsf p_I$, then $\mathsf C_{\rm pt}(\eta)\le8^4{\rm mv}_0^3c_I^{-3}C_I\eta^{-\mathsf p_I-3}$; if $\mathsf C_{\rm pt}(\eta)\le C\eta^{-\mathsf p}$ for $\eta\in(0,1]$, then $\sigma$ has (T5) in the sup norm with the same $c_I$, with $\max\{\mathsf p-3,0\}$ for $\mathsf p_I$ and $\frac{4\pi}3\,16^{\mathsf p}c_I^3C$ for $C_I$; in particular (T5) with $\mathsf p_I$ gives (T5) in the sup norm with $\mathsf p_I+1$.
\end{enumerate}
\item[(F)] \emph{Reflections; other spaces of values.} (a) If $\sigma$ has (T1$^\circ$) and (T3), then $\sigma(\theta_cX)=\sigma(X)$ for every $X\in\mathfrak S^V(\mathbb R^3)$ none of whose supports meets $\{x_3=c\}$. (b) Let $V'$ be a real Euclidean space of finite dimension and $\vartheta$ an orthogonal involution of $V'$; define tuples, tensor functionals and the conditions (T1$^\circ$), (T2), (T3), (T4), (T5$_N$), (T5) as above, with $V'$ in place of $V$, its Euclidean norm in place of $|\cdot|_2$, translations acting on the argument alone, $(\theta_cf)(x):=\vartheta f(\theta_cx)$, and $\mathfrak a^\theta:=\vartheta\bar{\mathfrak a}$. Then (A), (B) and, under (T5), (KV1) hold for $V'$ and $\vartheta$, with the same constants.
\end{enumerate}
\end{theorem}

\begin{theorem}[A converse on the closed span of the plate vectors]\label{C-353-8}
Let $\sigma$ have (T1$^\circ$), (T2), (T3), (T4) and (T5$_N$), with $N$ arbitrary, and let $\mathcal H_{\rm pl}\subset\mathcal H(\sigma)$ be the closed span of $\Omega$ and of all plate vectors $\Psi(x;\mathbf a)$, $n\ge1$, $\mathbf a\in\mathcal L_{\mathbb C}^n$, $x\in{\rm Tb}^{(n)}$. Suppose that there are numbers $C'>0$ and $\varpi\ge0$ with $|K_n(\xi;\mathbf a)|\le\prod_k(C'\nu(\mathfrak a_k))\cdot\min\{1,\min_kd_{\rm Tb}(\xi_k)\}^{-\varpi n}$ for all $n\ge1$, $\xi\in{\rm Tb}^{n-1}$ and $\mathbf a\in\mathcal L_{\mathbb C}^n$. Then $B_{\mathfrak b}(\xi,\xi')\mathcal H_{\rm pl}\subset\mathcal H_{\rm pl}$ and
\[\|B_{\mathfrak b}(\xi,\xi')\Phi\|\le C'\nu(\mathfrak b)\min\{1,2d_{\rm Tb}(\xi),d_{\rm Tb}(\xi')\}^{-\varpi}\|\Phi\|\qquad(\Phi\in\mathcal H_{\rm pl},\ \mathfrak b\in\mathcal L_{\mathbb C},\ \xi,\xi'\in{\rm Tb}).\]
Conversely, if $\|B_{\mathfrak b}(\xi,\xi')\Phi\|\le C''\nu(\mathfrak b)\min\{1,d_{\rm Tb}(\xi),d_{\rm Tb}(\xi')\}^{-\varpi}\|\Phi\|$ for all such $\Phi$, $\mathfrak b$, $\xi$, $\xi'$, then the bound on the $K_n$ holds with $C'=2^\varpi C''$.
\end{theorem}

\begin{corollary}[The density formula with a matrix]\label{C-353-9}
Let $\sigma$ have (T1), (T2), (T3), (T4) and (T5$_N$), let $d>0$ and $F,F'\in\mathbb C\mathfrak S^V(\{x_3>d\})$, and put, for $|x_3|<d$ and $E\in V_{\mathbb C}$,
\[\mathsf v_d(x;E):=\langle[\tau_{-de_3}F],B_E(de_3+x,de_3-x)[\tau_{-de_3}F']\rangle.\]
Then $\mathsf v_d(x;E)$ is complex linear in $E$; $\mathsf v_d(\cdot;E)$ is the restriction to the real points of a holomorphic function on $\{w\in\mathbb C^3:|{\rm Re}\,w_3|+|{\rm Im}\,w_1|+|{\rm Im}\,w_2|<d\}$; $\sigma(\theta\bar F\cdot(H)_{\mathbb C}\cdot F')=\int\mathsf v_d(x;H(x))\,dx$ for every $H\in C^1_c(\{|x_3|<d\};V_{\mathbb C})$; and $|\mathsf v_d(x;E)|\le2c_{\rm pt}|E|_2N(\eta)\eta^{-4}\|[\tau_{-de_3}F]\|\,\|[\tau_{-de_3}F']\|$ for every $\eta\in(0,\min\{1,d-|x_3|\}]$, the factor $2$ omitted for $E\in V$.
\end{corollary}

\begin{theorem}[One holomorphic function on all quarter-turned tubes]\label{C-353-10}
Let $n\ge2$, and let $\sigma$ have (T1), (T2), (T3), (T4) and (T5$_N$). For a permutation $\mathfrak s$ of $\{1,\dots,n\}$, $Z=(z_1,\dots,z_n)\in(\mathbb C^2)^n$ and $\mathbf a\in\mathcal L_{\mathbb C}^n$ put $\mathfrak sZ:=(z_{\mathfrak s(1)},\dots,z_{\mathfrak s(n)})$, $\mathfrak s\mathbf a:=(\mathfrak a_{\mathfrak s(1)},\dots,\mathfrak a_{\mathfrak s(n)})$ and $\partial Z:=(z_2-z_1,\dots,z_n-z_{n-1})$; for $\upsilon\in\frac\pi2\mathbb Z$ let $\mathcal Z(\mathfrak s,\upsilon):=\{Z:\partial(\mathfrak sZ)\in(R_\upsilon{\rm Tb})^{n-1}\}$, $\mathcal Z_n:=\bigcup_{\mathfrak s,\upsilon}\mathcal Z(\mathfrak s,\upsilon)$, and on $\mathcal Z(\mathfrak s,\upsilon)$
\[K^{\mathfrak s}_\upsilon(Z;\mathbf a):=K_n\bigl(R_{-\upsilon}\partial(\mathfrak sZ);\mathfrak s\,\varrho(R_{-\upsilon})\mathbf a\bigr),\]
$R_{-\upsilon}$ acting on each difference. Then:
\begin{enumerate}
\item[(a)] $K^{\mathfrak s}_\upsilon$ depends on $\upsilon$ modulo $2\pi$ only and is holomorphic on $\mathcal Z(\mathfrak s,\upsilon)$;
\item[(b)] for non-zero profiles $\alpha_1,\dots,\alpha_n$ with ${\rm supp}\,\alpha_1\times\dots\times{\rm supp}\,\alpha_n\subset\mathcal Z(\mathfrak s,\upsilon)$, $\sigma(\alpha_1\otimes\mathfrak a_1,\dots,\alpha_n\otimes\mathfrak a_n)=\int\prod_k\alpha_k(y_k)\,K^{\mathfrak s}_\upsilon(Y;\mathbf a)\,dY$;
\item[(c)] $K^{\mathfrak s}_\upsilon=K^{\mathfrak s'}_{\upsilon'}$ on $\mathcal Z(\mathfrak s,\upsilon)\cap\mathcal Z(\mathfrak s',\upsilon')$; so there is exactly one holomorphic function $\mathcal K_n(\cdot;\mathbf a)$ on $\mathcal Z_n$ that equals $K^{\mathfrak s}_\upsilon$ on every $\mathcal Z(\mathfrak s,\upsilon)$, and it is complex linear in each label;
\item[(d)] $\mathcal K_n(\mathfrak s'Z;\mathfrak s'\mathbf a)=\mathcal K_n(Z;\mathbf a)=\mathcal K_n(Z+(v,\dots,v);\mathbf a)$ for every permutation $\mathfrak s'$ and $v\in\mathbb C^2$; $\mathcal K_n(R_{\pi/2}Z;\mathbf a)=\mathcal K_n(Z;\varrho(R_{-\pi/2})\mathbf a)$; and $\overline{\mathcal K_n(Z;\mathbf a)}=\mathcal K_n(\bar Z;\bar{\mathbf a})$;
\item[(e)] if $\sigma$ has (T5), then $|\mathcal K_n(Z;\mathbf a)|\le\prod_k(C_\flat\nu(\mathfrak a_k))\cdot\min\{1,\min_kd_{\rm Tb}((R_{-\upsilon}\partial(\mathfrak sZ))_k)\}^{-\wp n}$ on $\mathcal Z(\mathfrak s,\upsilon)$;
\item[(f)] for real $x\in{\rm Tb}^{(n_1)}$ and $z\in{\rm Tb}^{(n_2)}$, real $s$, labels $\mathbf a\in\mathcal L_{\mathbb C}^{n_1}$, $\mathbf a'\in\mathcal L_{\mathbb C}^{n_2}$ and $X:=(\theta x_{n_1},\dots,\theta x_1,z_1,\dots,z_{n_2})$,
\[\langle\Psi(R_{is}x;\mathbf a),\Psi(R_{is}z;\mathbf a')\rangle=\mathcal K_{n_1+n_2}\bigl(R_{is}X;\mathfrak a_{n_1}^\theta,\dots,\mathfrak a_1^\theta,\mathfrak a'_1,\dots,\mathfrak a'_{n_2}\bigr);\]
\item[(g)] for helicity labels, $\varrho(R_t)E_s=e^{-i\,{\rm hel}(s)t}E_s$ and $(\varphi E_s)^\theta=(-1)^{{\rm hel}(s)}\varphi E_s$; hence, for a tuple $\mathbf a$ of helicity labels, $K^{\mathfrak s}_\upsilon(Z;\mathbf a)=e^{iS(\mathbf a)\upsilon}K_n(R_{-\upsilon}\partial(\mathfrak sZ);\mathfrak s\mathbf a)$, $\mathcal K_n(R_{\pi/2}Z;\mathbf a)=e^{iS(\mathbf a)\pi/2}\mathcal K_n(Z;\mathbf a)$, and the Gram formula of Theorem~\ref{C-353-7}\,(KV3) takes the form 
\[
\langle\Psi(x;\mathbf a),\Psi(z;\mathbf a')\rangle=(-1)^{S(\mathbf a)}K_{n+n'}(\dots;\mathfrak a_n,\dots,\mathfrak a_1,\mathfrak a'_1,\dots,\mathfrak a'_{n'}),
\]
with the same arguments;
\item[(h)] let $\mathbf a$ be a tuple of helicity labels, $S=S(\mathbf a)$, and $X$ a configuration of $n$ pairwise different points of $\mathbb R^2$ on one line. Then $R_\zeta X\in\mathcal Z_n$ for every $\zeta\in\mathbb C$, and $\zeta\mapsto e^{-iS\zeta}\mathcal K_n(R_\zeta X;\mathbf a)$ is entire and of period $\frac\pi2$; so
\[\mathcal K_n(R_\zeta X;\mathbf a)=\sum_{m\in\mathbb Z}c_m(X;\mathbf a)\,e^{i(S+4m)\zeta},\qquad c_m(R_tX;\mathbf a)=e^{i(S+4m)t}c_m(X;\mathbf a)\ (t\in\mathbb R),\]
the series converging locally uniformly. With the weight $k(E)$ of an $E\in V_{\mathbb C}$ defined by $\varrho(R_t)E=e^{ik(E)t}E$ one has $k(E_s)=-{\rm hel}(s)$, so the frequencies are the multiples of $4$ diminished by $\kappa=\sum_kk(E_{s_k})=-S$. If $\sigma$ has (T5), then $c_m(X;\mathbf a)=0$ unless $4|m|\le\wp n+|S|$, in particular unless $4|m|\le(\wp+2)n$.
\end{enumerate}
\end{theorem}

\begin{theorem}[A converse on the closed span of the point vectors]\label{C-353-11}
Let $\sigma$ have (T1), (T2), (T3), (T4) and (T5$_N$). For $n\ge1$, $\mathbf E=(E_1,\dots,E_n)\in V_{\mathbb C}^n$, $\kappa\in{\rm Tb}_3$, $\mathfrak z\in{\rm Tb}_3^{n-1}$ and $\mathfrak x\in{\rm Tb}_3^{(n)}$ put
\[
\begin{multlined}
W_n(\mathfrak z;\mathbf E):=\langle\Omega,B_{E_1}(\kappa,\tfrac{\mathfrak z_1}2)B_{E_2}(\tfrac{\mathfrak z_1}2,\tfrac{\mathfrak z_2}2)\cdots B_{E_n}(\tfrac{\mathfrak z_{n-1}}2,\kappa)\Omega\rangle,\qquad\Psi_3(\mathfrak x;\mathbf E):=B_{E_1}(\mathfrak x_1, \\
\tfrac{\mathfrak x_2-\mathfrak x_1}2)\cdots B_{E_n}(\tfrac{\mathfrak x_n-\mathfrak x_{n-1}}2,\kappa)\Omega,
\end{multlined}
\]
which do not depend on $\kappa$ by Theorem~\ref{C-353-7}\,(D)(iv), and let $\mathcal H_{\rm pt}$ be the closed span of $\Omega$ and of all point vectors $\Psi_3(\mathfrak x;\mathbf E)$. Suppose that there are $C'>0$ and $\varpi\ge0$ with $|W_n(\mathfrak z;\mathbf E)|\le\prod_k(C'|E_k|_2)\cdot\min\{1,\min_kd_3(\mathfrak z_k)\}^{-\varpi n}$ for all $n$, $\mathfrak z$, $\mathbf E$. Then $B_E(\mathfrak z,\mathfrak z')\mathcal H_{\rm pt}\subset\mathcal H_{\rm pt}$,
\[\|B_E(\mathfrak z,\mathfrak z')\Phi\|\le C'|E|_2\min\{1,2d_3(\mathfrak z),d_3(\mathfrak z')\}^{-\varpi}\|\Phi\|\qquad(\Phi\in\mathcal H_{\rm pt}),\]
and for $\eta\in(0,1]$, real $H\in C^1_c(\mathbb R^3;V)$ with support in $B(y,c_I\eta)$, $|y_3|\le\eta/8$, $F\in\mathbb C\mathfrak S^V(\{x_3>\eta/4\})$ and $G\in\mathbb C\mathfrak S^V(\{x_3>\eta/4\})$ with $[\tau_{-(\eta/4)e_3}G]\in\mathcal H_{\rm pt}$,
\[|\sigma(\theta\bar F\cdot(H)\cdot G)|\le\tfrac{4\pi}3\,c_I^3\,16^\varpi C'\,\eta^{3-\varpi}\,\|H\|_\infty\,\|[\tau_{-(\eta/4)e_3}F]\|\,\|[\tau_{-(\eta/4)e_3}G]\|:\]
on $\mathcal H_{\rm pt}$, the insertion bound in the sup norm holds with the power $\max\{\varpi-3,0\}$.
\end{theorem}

\begin{theorem}[Three Gaussian examples]\label{C-353-12}
\begin{enumerate}
\item[(a)] Let $\sigma_0$ be a scalar functional, that is, a functional on the tuples of real scalar $C^1_c$ functions on $\mathbb R^3$ with pairwise disjoint supports, with $\sigma_0(\emptyset)=1$ and multilinear in the same sense, and suppose that $\sigma_0$ has conditions (R1)--(R5) of Definition~\YMMref{95.1}{723} (symmetry and invariance under $\mathbb R^3\rtimes W_3$, regularity, reflection positivity, the condition $E\ge|q_1|$, and the insertion bound with numbers $c_I$, $C_I$, $\mathsf p_I$). Then $\sigma(H_1,\dots,H_n):=\sigma_0(\operatorname{tr}H_1,\dots,\operatorname{tr}H_n)$ is a tensor functional with (T1)--(T5), with the numbers $c_I$, $\sqrt3\,C_I$ and $\mathsf p_I$. If $\sigma_0$ has (R1)--(R4) and the insertion bound with a function $N_0$ in place of $C_I\eta^{-\mathsf p_I}$, then $\sigma$ has (T1)--(T4) and (T5$_N$) with $N=\sqrt3\,N_0$.
\item[(b)] Let $\phi^1,\phi^2,\phi^3$ be independent Gaussian free fields of mass zero on $\mathbb R^3$, with covariance $G(x)=(4\pi|x|)^{-1}$, let $\phi^c_\varepsilon:=j^{(3)}_\varepsilon*\phi^c$, and for $H\in C^1_c(\mathbb R^3;V)$ put $T_\varepsilon(H):=\sum_{c=1}^3\sum_{a,b=1}^3\int H_{ab}(x){:}\partial_a\phi^c_\varepsilon(x)\,\partial_b\phi^c_\varepsilon(x){:}\,dx$, the colons denoting the Wick order. For $(H_1,\dots,H_n)\in\mathfrak S^V(\mathbb R^3)$ the limit $\sigma(H_1,\dots,H_n):=\lim_{\varepsilon\to0}\mathbb E[T_\varepsilon(H_1)\cdots T_\varepsilon(H_n)]$ exists, and with $\sigma(\emptyset)=1$ it is a centred tensor functional, invariant under all isometries of $\mathbb R^3$ in the tensor action, with (T1)--(T4) and with (T5) in the sup norm with $c_I=\frac1{16}$, $\mathsf p_I=0$ and some $C_I$. The same holds when $\nabla\phi^c$ is replaced by the Hodge duals of the field strengths of three independent free photon fields with the action $\frac14\int F_{ij}F_{ij}$: at disjoint supports this multiplies $\sigma(H_1,\dots,H_n)$ by $(-1)^n$.
\item[(c)] Let $S(x):=\int_0^\infty e^{\sqrt m}e^{-m|x|}(4\pi|x|)^{-1}\,dm$ for $x\ne0$, and
\[\sigma(H_1,\dots,H_n):=\sum_\pi\prod_{\{k,l\}\in\pi}\iint\operatorname{tr}H_k(x)\,S(x-y)\,\operatorname{tr}H_l(y)\,dx\,dy,\]
the sum over the partitions $\pi$ of $\{1,\dots,n\}$ into pairs (it is empty for odd $n$), and $\sigma(\emptyset)=1$. Then $\sigma$ is a centred tensor functional with (T1)--(T4) and with (T5$_N$) for $c_I=\frac1{16}$ and a finite function $N$, so that (A), (B) and (D) of Theorem~\ref{C-353-7}, and (E)(i) but for its last sentence, hold for it. But for every real $\varphi\in C^1_c(\mathbb R)$ with $\varphi\ge0$, $\varphi\ne0$, and the unit matrix $\mathbf 1$, $t^qK_2((t,0);\varphi\mathbf 1,\varphi\mathbf 1)\to\infty$ as $t\downarrow0$ for every $q\ge0$; hence (KV1) holds for no $C_\flat$ and no $\wp$, and $\sigma$ has (T5) for no $C_I$ and no $\mathsf p_I$.
\end{enumerate}
\end{theorem}

\begin{remark}[A fixed matrix does not reduce the theorem to the scalar case]\label{C-353-13}
For $E\in V$ with $\varrho(\theta)E=E$ and $\sigma$ as in Theorem~\ref{C-353-7}\,(C), the scalar functional $\sigma_E(h_1,\dots,h_n):=\sigma(h_1E,\dots,h_nE)$ is reflection positive, since $\theta(hE)=(h\circ\theta)E$; the map $[F]_{\sigma_E}\mapsto[FE]_\sigma$, every member $h$ replaced by $hE$, is then an isometry that intertwines the translations, so that $\sigma_E$ has the regularity bound with $|E|_2^nC^\sigma_n$, the cone condition and the insertion bound with $|E|_2C_I$, because $\|hE\|_\sharp=|E|_2\|h\|_\sharp$. The scalar reconstruction hypotheses in
Definition~\YMMref{95.1}{723} include invariance under
$\mathbb R^3\rtimes W_3$, whereas (T1$^\circ$) gives $\sigma_E$
invariance under translations. If $\sigma$ has (T1), then $\sigma_E(wX)=\sigma_{\varrho(w)^{-1}E}(X)$ for $w\in W_3$, $w$ acting on the scalar tuple $X$ by $h\mapsto h\circ w^{-1}$, since $(h\circ w^{-1})E=w\bigl(h\,\varrho(w)^{-1}E\bigr)$ in the tensor action and (T1) holds at $w$; and $\varrho(w)E=E$ for all $w\in W_3$ only for $E\in\mathbb R\mathbf 1$.

Moreover, polarisation in $E$ of the scalar kernels of the $\sigma_E$ yields only the sums 
\[
\sum_\pi K_n(\xi;\varphi_1E_{\pi(1)},\dots,\varphi_nE_{\pi(n)})
\]
 over the permutations $\pi$, not the single kernels. Thus applying scalar reconstruction separately to the
$\sigma_E$ does not prove Theorem~\ref{C-353-7}. Its proof below
retains a matrix in every tuple member and every label.
\end{remark}

\begin{theorem}[The limits of the fields $\operatorname{tr}F_{ij}F_{kl}$]\label{C-353-14}
Let $g$ be a reference coupling with $0<g\le\gamma^{\rm ref}_W$. A Gibbs sequence at $g$ is a sequence $(K_j,\beta_j,\nu_j)_{j\ge1}$ in which $(K_j,\beta_j)$ are first-passage data at $g$, $K_j\to\infty$, and $\nu_j$ is a periodic Gibbs state at $(K_j,\beta_j)$. For a bounded $V$-valued function $H$ on $\mathbb R^3$ with bounded support put
\[\mathcal O^{\mathsf T}_K(H):=a_K^{-1}\sum_pH_{\lambda(p)\lambda(p)}(x(p))\,a_p-a_K^{-1}\sum_bH_{\lambda_b\lambda_b'}(x_b)\,q_b,\]
the sums running over the plaquettes $p$ and the positively oriented bonds $b$ of $a_K\mathbb Z^3$. Here $a_p:=1-\frac12\operatorname{tr}U(\partial p)$, with $U(\partial p)$ the product of the values of the configuration $U$ around $p$, $x(p)$ is the centre of $p$ and $\lambda(p)$ the direction normal to $p$; for $b=\langle x,x+a_Ke_\mu\rangle$, $x_b$ is the midpoint of $b$, $\lambda_b<\lambda_b'$ are the two directions other than $\mu$, and
\[q_b:=\tfrac14\sum_{\varepsilon,\varepsilon'=\pm1}\varepsilon\varepsilon'\,\bigl\langle w(W_x(\mu,\varepsilon\lambda_b)),w(W_x(\mu,\varepsilon'\lambda_b'))\bigr\rangle,\]
where, for $\alpha,\alpha'\in\{\pm1,\pm2,\pm3\}$ with $|\alpha|\ne|\alpha'|$ and $e_{-\mu}:=-e_\mu$, $W_x(\alpha,\alpha')$ is the product, in the order of the path $x\to x+a_Ke_\alpha\to x+a_Ke_\alpha+a_Ke_{\alpha'}\to x+a_Ke_{\alpha'}\to x$, of the values of $U$ on its four bonds, a bond run against its orientation carrying the inverse value; $w(V):=(V-V^{-1})/2i$ and $\langle X,Y\rangle:=\frac12\operatorname{tr}XY$. This is the lattice field $\operatorname{tr}F_{ij}F_{kl}$ at the cutoff $K$ smeared with $H$, as in Section~\ref{sf-sec} on the six fields $\operatorname{tr}F_{ij}F_{kl}$. It is real and linear in $H$, and $\mathcal O^{\mathsf T}_K(f\mathbf 1)$ is the action-density field smeared with $f$. The sequence converges if
\[\sigma(H_1,\dots,H_n):=\lim_{j\to\infty}\nu_j\Bigl(\prod_{a=1}^n\bigl(\mathcal O^{\mathsf T}_{K_j}(H_a)-\nu_j(\mathcal O^{\mathsf T}_{K_j}(H_a))\bigr)\Bigr)\]
exists for every $(H_1,\dots,H_n)\in\mathfrak S^V(\mathbb R^3)$; with $\sigma(\emptyset):=1$, $\sigma$ is its limit. A tensor limit at $g$ is the limit, in the same sense but at every tuple of continuous $V$-valued functions with compact, pairwise disjoint supports, of a Gibbs sequence at $g$ along which $a_{K_j}\beta_j$ converges to a number $b_*$.
\begin{enumerate}
\item[(a)] Every Gibbs sequence at $g$ has a subsequence that converges at every tuple of continuous $V$-valued functions with compact, pairwise disjoint supports and along which $a_{K_j}\beta_j$ converges to a number $b_*\in[g^{-2},L(g^{-2}+\bar B)]$; its limit is a tensor limit at $g$ with the number $b_*$. In particular every Gibbs sequence at $g$ has a convergent subsequence, and tensor limits at $g$ exist.
\item[(b)] Let $\sigma$ be the limit of a convergent Gibbs sequence at $g$. Then:
\begin{enumerate}
\item[(i)] $\sigma$ is a centred tensor functional with (T1), (T2), (T3) and (T4), and $\mathcal P$ is carried by $\{E\ge\max(|q_1|,|q_2|)\}$;
\item[(ii)] for every $\eta\in(0,1]$ the least number $\mathsf c^\circ_\sigma(\eta)$ with 
\[
|\sigma(\theta\bar F\cdot(h)\cdot F')|\le\mathsf c^\circ_\sigma(\eta)\|h\|_{L^1}\|[\tau_{-\eta e_3}F]\|\,\|[\tau_{-\eta e_3}F']\|
\]
 for all $F,F'\in\mathbb C\mathfrak S^V(\{x_3>\eta\})$ and all real $h\in C^1_c(\{|x_3|<\eta/2\};V)$ is finite, and $\sigma$ has (T5$_N$) with $c_I=\frac1{16}$ and $N(\eta)=\frac{4\pi}3\,16^{-3}\eta^3\,\mathsf c^\circ_\sigma(\eta/8)$;
\item[(iii)] the trace part $\sigma_0(f_1,\dots,f_n):=\sigma(f_1\mathbf 1,\dots,f_n\mathbf 1)$ is the centred-moment limit of the same Gibbs sequence for the action-density field, and $[F]_{\sigma_0}\mapsto[F\mathbf 1]_\sigma$, every member $f$ replaced by $f\mathbf 1$, is an isometry of $\mathcal H(\sigma_0)$ onto a closed subspace of $\mathcal H(\sigma)$ that reduces $\mathcal T_t$, $\mathcal U(v)$ and $\mathcal P$;
\item[(iv)] consequently (A), (B) and (D) of Theorem~\ref{C-353-7}, (E)(i) but for its last sentence, Corollary~\ref{C-353-9}, and Theorem~\ref{C-353-10} but for (e) and the last assertion of (h), hold for $\sigma$.
\end{enumerate}
\item[(c)] Let $0<g\le\gamma^{\rm sd}_W$, let $\sigma$ be a tensor limit at $g$ with the number $b_*$, and let $\sigma^1$ be its restriction to $\mathfrak S^V(\mathbb R^3)$. Then $\sigma^1$ is the limit of a convergent Gibbs sequence at $g$, so that (b) applies to it, and for every $\mathsf p>\frac{17}2$ there is a number $C_{\mathsf p}$, depending only on $g$, $\mathsf p$ and the fixed construction constants, such that:
\begin{enumerate}
\item[(i)] $|\sigma^1(\theta\bar F\cdot(h)\cdot G)|\le C_{\mathsf p}\min\{1,d-d'\}^{-\mathsf p}b_*^{-1}\|h\|_{L^1}\|[\tau_{-de_3}F]\|\,\|[\tau_{-de_3}G]\|$ for $0\le d'<d$, real $h\in C^1_c(\mathbb R^3;V)$ with support in $\{|x_3|\le d'\}$, and $F,G\in\mathbb C\mathfrak S^V(\{x_3>d\})$;
\item[(ii)] $\sigma^1$ has (T5) in the sup norm with $c_I=\frac1{16}$, $\mathsf p_I=\mathsf p-3$ and $C_I=\frac{4\pi}3\,16^{\mathsf p-3}C_{\mathsf p}/b_*$; so all of Theorem~\ref{C-353-7}, Corollary~\ref{C-353-9} and Theorem~\ref{C-353-10} apply to $\sigma^1$, and (KV1) holds in the form for the sup norm with $\wp'=\mathsf p$;
\item[(iii)] $\|B_{\mathfrak a}(\xi,\xi')\|\le\sqrt2\,C_{\mathsf p}b_*^{-1}\|\mathfrak a\|_{L^1}\min\{1,d_{\rm Tb}(\xi),d_{\rm Tb}(\xi')\}^{-\mathsf p}$ for $\mathfrak a\in\mathcal L_{\mathbb C}$, without the factor $\sqrt2$ for real $\mathfrak a$, and $|K_n(\xi;\mathbf a)|\le\prod_k\bigl(\sqrt2\cdot2^{\mathsf p}C_{\mathsf p}b_*^{-1}\|\mathfrak a_k\|_{L^1}\bigr)\cdot\min\{1,\min_kd_{\rm Tb}(\xi_k)\}^{-\mathsf pn}$;
\item[(iv)] $\|B_E(\mathfrak z,\mathfrak z')\|\le\sqrt2\,C_{\mathsf p}b_*^{-1}|E|_2\min\{1,d_3(\mathfrak z),d_3(\mathfrak z')\}^{-\mathsf p}$ for $E\in V_{\mathbb C}$, without the factor $\sqrt2$ for $E\in V$; $\mathcal U_wB_E(\mathfrak z,\mathfrak z')\mathcal U_w^{-1}=B_{\varrho(w)E}(w\mathfrak z,w\mathfrak z')$ for the eight $w\in W_3$ with $we_3=e_3$, where $\mathcal U_w[F]:=[wF]$; the functions $W_n$ of Theorem~\ref{C-353-11} are holomorphic on ${\rm Tb}_3^{n-1}$, $n$-linear in $\mathbf E$, and $|W_n(\mathfrak z;\mathbf E)|\le\prod_k\bigl(\sqrt2\cdot2^{\mathsf p}C_{\mathsf p}b_*^{-1}|E_k|_2\bigr)\cdot\min\{1,\min_kd_3(\mathfrak z_k)\}^{-\mathsf pn}$; and
\[\sigma^1(h_1,\dots,h_n)=\int W_n\bigl(x_2-x_1,\dots,x_n-x_{n-1};h_1(x_1),\dots,h_n(x_n)\bigr)\,dx_1\cdots dx_n\]
whenever $\sup_{{\rm supp}\,h_k}x_3<\inf_{{\rm supp}\,h_{k+1}}x_3$ for all $k$; for $w\in W_3$, the same holds with $W_n(w^{-1}(x_2-x_1),\dots;\varrho(w^{-1})h_1(x_1),\dots)$ whenever $w^{-1}(h_1,\dots,h_n)$ is so ordered, that is, whenever the supports are ordered in the coordinate direction $we_3$;
\item[(v)] for a compact set $\mathcal C$ of real points of ${\rm Tb}^{(j)}$ and $c_{\mathcal C}:=\min\{1,\min_{z\in\mathcal C}\min_i\min\{z_1^0,\frac12(z_{i+1}-z_i)^0\}\}$: $\|\Psi(R_{is}z;\mathbf a)\|\le\prod_i\bigl(\sqrt2\,C_{\mathsf p}b_*^{-1}\|\mathfrak a_i\|_{L^1}\bigr)\cdot c_{\mathcal C}^{-\mathsf pj}e^{\mathsf pj|s|}$ for $z\in\mathcal C$, $s\in\mathbb R$ and $\mathbf a\in\mathcal L_{\mathbb C}^j$.
\end{enumerate}
\end{enumerate}
\end{theorem}

\subsection{Construction estimates}

Theorems \ref{C-353-7}--\ref{C-353-12} use neither Theorem \YMMref{D}{8} nor Theorems \YMMref{AJ}{673} and \YMMref{AD}{607} among the construction results, nor the continuum-limit and reconstruction conclusions of Theorems \YMMref{A}{5} and \YMMref{AP}{780}; no lattice, coupling, cutoff or volume enters them. They rest on the following statements among the construction results.
\begin{itemize}
\item The abstract reconstruction theorem: a vector space with a positive semi-definite Hermitian form and a distinguished vector, a semigroup of maps symmetric for the form with a bound of $\langle F,\mathcal T^\circ_tF\rangle$ uniform in $t$, and commuting maps that preserve the form, all weakly continuous, give a strongly continuous semigroup of self-adjoint contractions and a strongly continuous unitary group, commuting and fixing the vector.
\item The statement on the joint spectral measure of such a pair, on $[0,\infty)\times\mathbb R^2$, with $\mathcal P(\{0\})$ fixing an invariant vector; the rule $\int\mathbf 1_B\,d\mathcal P=\mathcal P(B)$ and the spectral integrals $g(\mathcal P)$ with the complex measures $\langle\psi,\mathcal P(\cdot)\psi'\rangle$.
\item The facts of analysis (the identity theorem from real points; the fundamental lemma; for scalar $C^1_c$ functions ${\rm Lip}\le\sup|\nabla\cdot|$, the continuity of translations in $\|\cdot\|_{C^1}$, the convergence in $\|\cdot\|_{C^1}$ of the Riemann sums $\sum_k\varepsilon^N\alpha(\varepsilon k)\tau_{\varepsilon k}j$ to $\alpha*j$ with supports in ${\rm supp}\,\alpha+{\rm supp}\,j$, and $\alpha*j_\delta\to\alpha$ in $\|\cdot\|_{C^1}$; a weakly holomorphic, locally bounded operator function is holomorphic in norm; a $\frac\pi2$-periodic entire function $f$ with $|f(t+i\tau)|\le Ce^{A|\tau|}$ equals $\sum_{4|m|\le A}c_me^{4im\zeta}$).
\item The scalar partitions of unity (a function $\chi\in C^1_c(\mathbb R)$ with $0\le\chi\le1$, ${\rm supp}\,\chi=[-1,1]$, ${\rm Lip}\,\chi=\frac\pi2$, $\sum_k\chi(s-k)=1$, and the products $\chi_{\mathsf h,k}(x)=\prod_i\chi(x_i/\mathsf h-k_i)$, $k\in\mathbb Z^3$, with $\sum_k\chi_{\mathsf h,k}=1$, ${\rm supp}\,\chi_{\mathsf h,k}\subset\{|x-\mathsf hk|_\infty\le\mathsf h\}$ and ${\rm Lip}\,\chi_{\mathsf h,k}\le\sqrt3\pi/2\mathsf h$).
\item The calculus of operator functions (products and compositions with affine maps of functions holomorphic in norm are holomorphic in norm; the Cauchy estimate, by which a function holomorphic in norm and bounded by $b$ on a polydisc of radius $\rho$ in $\mathbb C^m$ has a Lipschitz constant at most $c_mb/\rho$ on the concentric polydisc of radius $\rho/2$; integrals of norm-continuous operator functions of compact support, with matrix elements, the norm bound by the integral of the norm, linearity, and passage to uniform limits).
\item The mean-value formula with its numbers and mollifiers (for $N\in\{2,3\}$ and $r>0$ there are real functions $\alpha^{(N)}_{r,t}\in C^1_c(\mathbb R^N)$, $t\in(-r,r)^N$, supported in $[-r,r]^N$, with $\sup|\alpha^{(N)}_{r,t}|\le{\rm mv}_0^Nr^{-2N}$ and $\sup|\partial_j\alpha^{(N)}_{r,t}|\le{\rm mv}_0^{N-1}{\rm mv}_1r^{-2N-1}$, such that $u(0)=\int_{(-r,r)^N}dt\int u(y+it)\alpha^{(N)}_{r,t}(y)\,dy$ for every $u$ holomorphic near the closed polydisc of radius $r$ about $0$; and mollifiers $j^{(N)}_\delta\ge0$, real, with integral $1$, supported in $\{|y|\le\delta/2\}$ and even under the reflection of each coordinate).
\end{itemize}
For Theorem~\ref{C-353-12} we use, in addition, Wick's theorem, the reflection positivity of Gaussian free fields and of the Gaussian fields with the covariances $S_\Lambda(x)=\int_0^\Lambda e^{\sqrt m}e^{-m|x|}(4\pi|x|)^{-1}dm$, $\Lambda<\infty$, and the description of the reconstructed spaces of these fields as Fock spaces on which the time and space translations act by $e^{-t\mathbf H_0}$ and $e^{iv\cdot\mathbf P_0}$, with the fields in Fock-space form. Each $S_\Lambda$ is a superposition with positive weights of the covariances of free fields of masses $m<\Lambda$; the proof of Theorem~\ref{C-353-12}\,(c) uses these fields at $\Lambda<\infty$, the functional of (c) being defined by the pairing formula of $S$.

For the free field $\phi$ of mass $m\ge0$ on $\mathbb R^3$ (and in the same way for superpositions of such fields over $m$ with a positive weight, such as the fields with the covariances $S_\Lambda$, the one-particle space carrying the weight in $m$): the covariance $\mathbb E[\phi(x)\phi(y)]=\int e^{ip\cdot(\vec x-\vec y)-|x_3-y_3|\omega_m(p)}\,d\mu_m(p)$, with $\vec x=(x_1,x_2)$, $\omega_m(p)=(|p|^2+m^2)^{1/2}$ and $d\mu_m(p)=(2\pi)^{-2}(2\omega_m(p))^{-1}d^2p$ on $\mathbb R^2$; the identity $\mathbb E\bigl[\prod_i{:}e^{\phi(f_i)}{:}\bigr]=\exp\sum_{i<k}\mathbb E[\phi(f_i)\phi(f_k)]$ for the Wick exponentials of finitely many smeared fields; the class of ${:}e^{\phi(g)}{:}$, for real $g\in C^1_c(E_+;\mathbb C)$, which is the coherent vector $e^{a^*(u_g)}\Omega$ of the Fock space over $L^2(\mathbb R^2,\mu_m)$, $u_g(p)=\int g(x)e^{-x_3\omega_m(p)+ip\cdot\vec x}\,dx$; and the matrix elements $\langle e^{a^*(u)}\Omega,{:}O(a^*,a){:}\,e^{a^*(v)}\Omega\rangle=O(\bar u,v)\,e^{\langle u,v\rangle}$ of a normal-ordered polynomial in the creation and annihilation operators between coherent vectors.

Theorem~\ref{C-353-14} rests in addition on three statements among the construction results, which give Gibbs sequences at $g$: the choice preceding equation~\YMMref{(2326)}{845}, by which $\gamma^{\rm ref}_W\le\gamma$ for the constant $\gamma$ among the construction results; Corollary~\YMMref{74.2}{616}, by which, for $0<g\le\gamma$, the interval $I(g)$ among the construction results, defined in~(\YMMref{(1527)}{615}), is not empty and the pairs $(K,L^K\mathsf b)$, $\mathsf b\in I(g)$, $K\in\mathbb N$, are first-passage data at $g$; and Proposition~\YMMref{91.2}{699}, by which at each first-passage datum at $g$ the states of the tori form a non-empty set of periodic Gibbs states. It rests further on three results of two other sections, which in turn rest on Theorems~\YMMref{D}{8}, \YMMref{AJ}{673} and~\YMMref{AD}{607} among the construction results.
\begin{itemize}
\item[(I)] (Section~\ref{sf-sec} on the six fields $\operatorname{tr}F_{ij}F_{kl}$ on the tori and on $\mathbb R^3$.) Let $0<g\le\gamma^{\rm ref}_W$. Every Gibbs sequence $(K_j,\beta_j,\nu_j)$ at $g$ has a subsequence that converges at every tuple of $\mathfrak S^V(\mathbb R^3)$, and along a further subsequence $a_{K_j}\beta_j$ converges to a number in $[g^{-2},L(g^{-2}+\bar B)]$. Moreover, for open sets $O_1,\dots,O_n$ that lie in a bounded set and are at pairwise positive distances there are a number $C$, depending on $g$, $n$ and the sets, and an index $j_0$ such that for $j\ge j_0$ the centred moments $\nu_j\bigl(\prod_a(\mathcal O^{\mathsf T}_{K_j}(H_a)-\nu_j(\mathcal O^{\mathsf T}_{K_j}(H_a)))\bigr)$ are bounded by $C\prod_a\sup_x|H_a(x)|_\infty$ for all bounded Borel functions $H_a\colon\mathbb R^3\to V$ that vanish off $O_a$, $a=1,\dots,n$; they are $n$-linear in such functions.
\item[(II)] (The same section.) Let $\sigma$ be the limit of a convergent Gibbs sequence at $0<g\le\gamma^{\rm ref}_W$. Then $\sigma$ is symmetric and centred; for every $n\ge1$, $\delta>0$ and $R>0$ there is a number $C_n(\delta,R)$ with $|\sigma(H_1,\dots,H_n)|\le C_n(\delta,R)\prod_a\sup_x|H_a(x)|_\infty$ at all tuples of type $(\delta;R)$; $\sigma(rX)=\sigma(X)$ for $r\in\mathbb R^3\rtimes W_3$, acting on arguments and indices; $\sigma(\theta_{n,c}\bar F\cdot F)\ge0$ for every $n\in\{e_\mu,\ e_\mu\pm e_\lambda\}$, every real $c$, the reflection $\theta_{n,c}$ in $\{n\cdot x=c\}$ acting on indices too, and every $F\in\mathbb C\mathfrak S^V(\{n\cdot x>c\})$; the joint spectral measure of the reconstructed translations is carried by $\{E\ge\max(|q_1|,|q_2|)\}$; for every $\eta\in(0,1]$ there is a finite number bounding $|\sigma(\theta\bar F\cdot(h)\cdot F')|$ by itself times $\int|h(x)|_\infty dx\,\|[\tau_{-\eta e_3}F]\|\,\|[\tau_{-\eta e_3}F']\|$, for $F,F'\in\mathbb C\mathfrak S^V(\{x_3>\eta\})$ and real $h\in C^1_c(\{|x_3|<\eta/2\};V)$; and the statement on the trace part in Theorem~\ref{C-353-14}\,(b)(iii) holds.
\item[(III)] (Section~\ref{C-355-1} on the point bound for the fields $\operatorname{tr}F_{ij}F_{kl}$.) Let $0<g\le\gamma^{\rm sd}_W$ and let $\sigma$ be a tensor limit at $g$ with the number $b_*$. For every $\mathsf p>\frac{17}2$ there is a number $C_{\mathsf p}$, depending only on $g$, $\mathsf p$ and the fixed construction constants, such that the centred slab operators of $\sigma$, formed on the space reconstructed from tuples of continuous members, satisfy $\|B^c_d(h)\|\le C_{\mathsf p}\min\{1,d-d'\}^{-\mathsf p}b_*^{-1}\int|h(x)|_\infty\,dx$ for $0\le d'<d$ and every real $h\in C_c(\mathbb R^3;V)$ with support in $\{|x_3|\le d'\}$.
\end{itemize}
The passage from (I) to the tuples of continuous members, on the whole range $0<g\le\gamma^{\rm ref}_W$, is a step of this section, made in item (d) of the remark below from the convergence and the bound of (I); Theorem~\ref{C-353-14}\,(a) rests on (I) and on this step. Accordingly the tensor limits at $g$ of Theorem~\ref{C-353-14} are defined in this section for every $0<g\le\gamma^{\rm ref}_W$; for $0<g\le\gamma^{\rm sd}_W$ they are the tensor limits that Section~\ref{C-355-1} on the point bound treats and that (III) reads.

\begin{remark}[How (I) and (II) are obtained]
In Section~\ref{sf-sec} on the six fields, (I) and (II) come from bounds for products of centred fields that read only the number of the fields, the sup norms of their weights, the least distance of the supports and the diameters of the supports, together with the passage of the lattice symmetries and of reflection positivity to the limit; item (d) ends with the step of this section that passes the convergence of (I) to the tuples of continuous members. For a Gibbs sequence $(K_j,\beta_j,\nu_j)$ at $g$ we write $\sigma_j(H_1,\dots,H_n):=\nu_j\bigl(\prod_a(\mathcal O^{\mathsf T}_{K_j}(H_a)-\nu_j(\mathcal O^{\mathsf T}_{K_j}(H_a)))\bigr)$, so that the limit of a convergent Gibbs sequence is $\sigma=\lim_j\sigma_j$.
\begin{itemize}
\item[(a)] \emph{One torus.} At a first-passage datum $(K,\beta)$ at $g$ and a side $\ell$ among the construction results at $g$, the joint cumulants of $q$ fields $\mathcal O^{\mathsf T}_K(H_a)$ whose weights vanish off closed sets at pairwise sup distances $\ge2\delta$ are bounded by a number of $g$, $q$, $\ell$ and $\delta$ times $\prod_a\sup_x|H_a(x)|_\infty$, as soon as $3a_K\le\delta$ and $K$ exceeds a threshold that reads $g$, $q$ and $\delta$. Here Theorem~\YMMref{D}{8} enters.
\item[(b)] \emph{One cell.} Let $\mathsf l=L^{-i}$, $0<\delta\le\mathsf l/4$, and let $\ell'$ be a side at $g$ with $\ell'\ge2\mathsf l$. Let $H_1,\dots,H_n$ vanish off closed subsets of one cell, of side $\mathsf l$, of a board of the lattice $a_K\mathbb Z^3$, at sup distances $\ge\delta$ from the walls of the cell and $\ge2\delta$ from one another, and let $\mathsf D:=\prod_a(\mathcal O^{\mathsf T}_K(H_a)-c_a)$, with $c_a$ the mean of $\mathcal O^{\mathsf T}_K(H_a)$ in the torus measure of side $\ell'$. The $N$-th power of the chessboard seminorm of $\mathsf D$ on the torus of side $\ell'$, $N=(\ell'/\mathsf l)^3$, is the mean of the product of the $N$ mirror copies of $\mathsf D$. The mirror maps have diagonal linear parts, and a lattice symmetry $r$ carries $\mathcal O^{\mathsf T}_K(H)$ to $\mathcal O^{\mathsf T}_K(rH)$; so the copies are products of $nN$ centred fields of the same kind, with weights of the same sup norms at pairwise sup distances $\ge2\delta$. The mean of a product of centred fields is the sum, over the partitions without blocks of one element, of the products of the joint cumulants of the blocks, and (a) bounds each of them. So the chessboard seminorm of $\mathsf D$ is at most a number of $g$, $n$, $\mathsf l$, $\ell'$ and $\delta$ times $\prod_a\sup_x|H_a(x)|_\infty$.
\item[(c)] \emph{All volumes and all periodic Gibbs states.} By the chessboard estimate among the construction results, the modulus of the mean of $\mathsf D$ in the torus measure of every side $\ell\ge\ell'$ at $g$ is at most $e^{\mathsf l^3(p_{\ell'}-p_\ell)}$ times this seminorm, $p_\ell$ being the pressure on the torus of side $\ell$; by the domination among the construction results along a sequence of sides tending to infinity, the same holds in every periodic Gibbs state. The pressure gap $0\le p_{\ell'}-p_\ell\le G_g$, with a number $G_g$ of $g$ and the constants among the construction results, rests on Theorems~\YMMref{AJ}{673} and~\YMMref{AD}{607} and turns the factor into $e^{G_g\mathsf l^3}$.

For weights in several cells of one board, the mean of the product of the cell products is at most the product of the cell bounds, and no number reads the positions of the cells. Cutting the weights into cubes of side $\mathsf l/4n$ of one grid and choosing, for each product of pieces, a board whose walls keep a distance from all the pieces, one obtains a bound by a number times $\prod_a\sup_x|H_a(x)|_\infty$ for weights that vanish off closed sets of given diameters at a given least distance, the number reading $g$, $n$, the least distance and the diameters alone; the expansion of $\prod_a(\mathcal O^{\mathsf T}_K(H_a)-\nu(\mathcal O^{\mathsf T}_K(H_a)))$ into the products $\prod_{a\in J}(\mathcal O^{\mathsf T}_K(H_a)-c_a)$, $J\subset\{1,\dots,n\}$, whose coefficients $\prod_{a\notin J}(c_a-\nu(\mathcal O^{\mathsf T}_K(H_a)))$ are bounded by the case of one weight, passes it to the $\sigma_j$.

The bound holds uniformly in the cutoff from a threshold on, in the first-passage datum and in the periodic Gibbs state. These are the bounds of (II) and, at the $\sigma_j$, the bound of (I).
\item[(d)] \emph{Extraction.} The bounds of (c) hold for bounded weights, and the $\sigma_j$ are multilinear. For each finite union $O$ of open cubes with rational corners choose a countable subset of $C^1_c(O;V)$, dense in $C^1_c(O;V)$ for the sup norm; the tuples of such members, for sets $O$ with pairwise disjoint closures, form a countable family, and a diagonal subsequence converges at each of its tuples. For a tuple $X\in\mathfrak S^V(\mathbb R^3)$, and a tuple $X'$ of the family whose members are $\varepsilon$-close to those of $X$ in the sup norm and lie in the same sets $O$, the bound (c) and multilinearity give $|\sigma_j(X)-\sigma_j(X')|\le C\varepsilon$ for all large $j$, with $C$ independent of $j$ and $\varepsilon\le1$; so the subsequence converges at $X$. The numbers $a_{K_j}\beta_j$ lie in $(g^{-2},L(g^{-2}+\bar B)]$, and a further subsequence makes them converge. This is (I).

\emph{The step of this section: continuous members.} Let $0<g\le\gamma^{\rm ref}_W$, take the subsequence of (I) along which $a_{K_j}\beta_j$ converges to a number $b_*$, and let $X=(H_1,\dots,H_n)$ be a tuple of continuous $V$-valued functions with compact, pairwise disjoint supports. Choose open sets $O_a\supset{\rm supp}\,H_a$ in a bounded set with pairwise disjoint closures. As $C^1_c(O_a;V)$ is dense in $C_c(O_a;V)$ for the sup norm (by mollification), there are $H'_a\in C^1_c(O_a;V)$ with $\sup_x|H_a(x)-H'_a(x)|_\infty\le\varepsilon\le1$, and $X'=(H'_1,\dots,H'_n)\in\mathfrak S^V(\mathbb R^3)$. By multilinearity,
\[\sigma_j(X)-\sigma_j(X')=\sum_{a=1}^n\sigma_j(H'_1,\dots,H'_{a-1},H_a-H'_a,H_{a+1},\dots,H_n),\]
each term having weights that vanish off $O_1,\dots,O_n$; so the bound of (I) gives $|\sigma_j(X)-\sigma_j(X')|\le nC\varepsilon\prod_a(\sup_x|H_a(x)|_\infty+1)$ for $j\ge j_0$, a bound free of $j$. As $(\sigma_j(X'))_j$ converges along the subsequence, $(\sigma_j(X))_j$ is a Cauchy sequence along it. So the subsequence converges at every tuple of continuous members, its limit is a tensor limit at $g$ with the number $b_*$, and together with (I) this is Theorem~\ref{C-353-14}\,(a).
\item[(e)] \emph{Symmetry, covariance, positivity.} Symmetry and centring pass to the limit. The periodic Gibbs states are invariant under the symmetries $a_K\mathbb Z^3\rtimes W_3$ of the lattice, and such a symmetry $r$ carries $\mathcal O^{\mathsf T}_K(H)$ to $\mathcal O^{\mathsf T}_K(rH)$, $r$ acting on the argument and on the indices of $H$; so $\sigma_j(rX)=\sigma_j(X)$. A real translation is the limit of lattice translations, with an error that the bounds of (c) and the Lipschitz constants of the members control. The periodic Gibbs states are reflection positive at the lattice mirrors $\{n\cdot x=c\}$, $c\in a_K\mathbb Z$, of the nine families $n\in\{e_\mu,\ e_\mu\pm e_\lambda\}$; the reflection acts on the field as on the weight, on the argument and on the indices, and a real $c$ is the limit of such lattice values.
\item[(f)] \emph{The cone.} Reflection positivity at the mirrors of the families $e_3$ and $e_3\pm e_k$, $k=1,2$, gives, by Lemma~\ref{lem:spectral-steps}, for $u(w):=\sigma_j(\theta\bar F\cdot\tau_wF)$, $F\in\mathbb C\mathfrak S^V(E_+)$, and every $t\in a_{K_j}\mathbb N$,
\[u(0)-4u(te_3)+u(2te_3)+u(te_3+te_k)+u(te_3-te_k)\ge0.\]
In the limit, with lattice points $t_j\to t$, and by density, for every $\psi\in\mathcal H(\sigma)$ and every $t>0$, $\int f_t\,d\langle\psi,\mathcal P\psi\rangle\ge0$ with
\[f_t(E,q)=(1-e^{-tE})^2-e^{-tE}|1-e^{itq_k}|^2=4e^{-tE}\bigl(\sinh^2\tfrac{tE}2-\sin^2\tfrac{tq_k}2\bigr).\]
At $\psi=\mathcal P(\Delta)\psi'$, for all Borel sets $\Delta$, this gives $f_t\ge0$ $\mathcal P$-almost everywhere, for all rational $t>0$ at once; there $\sinh\frac{tE}2\ge|\sin\frac{tq_k}2|$, and $t\downarrow0$ gives $E\ge|q_k|$. So $\mathcal P(\{E<\max(|q_1|,|q_2|)\})=0$.
\item[(g)] \emph{Slab and point operators.} Let $\mathsf l=L^{-i}$, $Q:=(v,0)+(-\frac{\mathsf l}2,\frac{\mathsf l}2)^2\times(0,\mathsf l)$, and let $X=(h_1,\dots,h_n)$ be a tuple of members with supports in $Q$ at sup distances $\ge\delta$ from $\partial Q$ and $\ge2\delta$ from one another. The cell-insertion bound of Proposition~\YMMref{104.2}{759},
applied to a board one of whose cells is $Q$ up to a shift of
the order of $a_{K_j}$, bounds $|\sigma_j(\theta\bar F'\cdot X\cdot\tau_{\mathsf le_3}F)|$, for $F,F'\in\mathbb C\mathfrak S^V(E_+)$, by the chessboard seminorm of the centred product of the fields of $X$ times $\sigma_j(\theta\bar F'\cdot F')^{1/2}\sigma_j(\theta\bar F\cdot F)^{1/2}$.

By (b) and (c), in the limit there is exactly one bounded operator $\mathsf B_{\mathsf l}(X)$ on $\mathcal H(\sigma)$ with $\mathsf B_{\mathsf l}(X)[F]=[X\cdot\tau_{\mathsf le_3}F]$, and $\|\mathsf B_{\mathsf l}(X)\|\le\mathsf C^{\rm sl}_n(g;\mathsf l,\delta)\prod_a\sup_x|h_a(x)|_\infty$, with a number $\mathsf C^{\rm sl}_n(g;\mathsf l,\delta)$ formed as in (c). Let $n=1$ and $\delta=\mathsf l/16$. The translations act covariantly on these operators, and by the square cone of (f) the matrix elements of $\mathsf B_{\mathsf l}((\tau_yh))$ between vectors in the range of $\mathcal T_\delta$ are the restrictions to real $y$ of functions holomorphic on $\{|{\rm Re}\,y_3|+|{\rm Im}\,y_1|+|{\rm Im}\,y_2|<\delta\}$.

The Cauchy estimate for the derivatives $\partial_S$, $S\subset\{1,2,3\}$, at $y=0$ and an integration by parts in $\partial_1\partial_2\partial_3$ show that the form $h\mapsto\langle\mathcal T_\delta\psi',\mathsf B_{\mathsf l}((h))\mathcal T_\delta\psi\rangle$, on members with support in the points of $Q$ at sup distance $>3\delta$ from $\partial Q$, has a continuous density and is bounded by $6(80/\mathsf l)^3\mathsf C^{\rm sl}_1(g;\mathsf l,\mathsf l/16)\|\psi'\|\,\|\psi\|\int|h(x)|_\infty\,dx$. These densities are the matrix elements of point operators $\mathsf B(s,s';e)$, $s,s'>0$, $e\in V$, with $\|\mathsf B(s,s';e)\|\le\mathsf C^{\rm fin}(\min\{s,s'\})|e|_\infty$ for the finite function
\[\mathsf C^{\rm fin}(\eta):=6(80/\mathsf l_\eta)^3\,\mathsf C^{\rm sl}_1(g;\mathsf l_\eta,\mathsf l_\eta/16),\qquad\mathsf l_\eta\ \text{the largest power of $L$ with}\ \tfrac9{16}\mathsf l_\eta\le\eta;\]
and they give $|\sigma(\theta\bar F\cdot(h)\cdot F')|\le\mathsf C^{\rm fin}(\eta/2)\int|h(x)|_\infty\,dx\,\|[\tau_{-\eta e_3}F]\|\,\|[\tau_{-\eta e_3}F']\|$ for $\eta$, $F$, $F'$ and $h$ as in (II), the last bound of (II).
\item[(h)] \emph{The trace part.} $\mathcal O^{\mathsf T}_K(f\mathbf 1)$ is the action-density field smeared with $f$, so $\sigma(f_1\mathbf 1,\dots,f_n\mathbf 1)$ is the limit of the centred moments of the scalar field. The map $F\mapsto F\mathbf 1$ preserves the forms, as $\varrho(\theta)\mathbf 1=\mathbf 1$, and the closed span of the classes $[F\mathbf 1]$ is invariant under the self-adjoint operators $\mathcal T_t$ and the unitary group $\mathcal U(v)$; so it reduces them and $\mathcal P$.
\end{itemize}
\end{remark}

\begin{remark}[Where the exponent $\frac{17}2$ and the factor $b_*^{-1}$ come from]\label{C-353-15}
In Section~\ref{C-355-1} on the point bound, (III) is obtained from an estimate at a cutoff of one centred smeared field in one cell of side $a_i=L^{-i}$: for $i$ past a threshold its chessboard norm is bounded by the power $a_i^{-11/2}$ times a power of $\log(1/a_i)$ times $(a_K\beta)^{-1}$ times the sup norm of the weight, at every cutoff past a threshold that reads $i$, uniformly in the cutoff, the first-passage datum, the cell, the side of the torus and the periodic Gibbs state; it uses a dilation to unit scale, a moment of mirror copies on the least torus, and the pressure gap. The factor $(a_K\beta)^{-1}$ comes from the bound on the moments of the sources at the working level (equation~\YMMref{(2322)}{844}), which carries a factor at most $(9/8)^{1/2}(a'\beta)^{-1}$ per field, $a'$ the spacing of the dilated lattice; as the dilation multiplies the weight and the spacing by the same power of $L$, this is the factor $(a_K\beta)^{-1}$ per unit of the sup norm of the original weight.

As $a_{K_j}\beta_j\to b_*$ along a tensor limit, this factor is the origin of the factor $b_*^{-1}$ in (III) and in Theorem~\ref{C-353-14}\,(c). Multiple reflections in periodic Gibbs states turn the estimate into a bound, by the sup norm and with the factor $b_*^{-1}$, of the slab operator of height $a_i$ for weights carried by a cube of side comparable to $a_i$; and an abstract lemma, which uses only the covariance under small translations and the spectral condition $E\ge\max(|q_1|,|q_2|)$, turns a sup-norm bound on cubes of half-side $\rho$ into a bound by the $L^1$ norm at the price $\rho^{-3}$. So $\frac{17}2=\frac{11}2+3$, and the volume $\frac{4\pi}3(\eta/16)^3$ of the balls of (T5) gives back three powers: $\mathsf p_I=\mathsf p-3>\frac{11}2$.
\end{remark}

\subsection{Estimates and proofs}

\begin{lemma}[Functions with matrix values]\label{C-353-16}
Let $f$ take values in $V_{\mathbb C}$; the norms are those of Definition~\ref{C-353-2}.
\begin{enumerate}
\item[(a)] ${\rm supp}\,f^\varepsilon\subset{\rm supp}\,f$, and $\|f^\varepsilon\|\le\|f\|$ for $\varepsilon=0,1$ and each of $\|\cdot\|_\infty$, ${\rm Lip}$, $\|\cdot\|_{C^1}$, $\|\cdot\|_\sharp$, $\|\cdot\|_{L^1}$.
\item[(b)] ${\rm Lip}\,f\le\sup|\nabla f|$; so $\|f\|_\sharp\le40M\|f\|_{C^1}$ on $\mathbb R^3$, and $\|\mathfrak a\|_{C^1}=\max\{\|\mathfrak a\|_\infty,{\rm Lip}\,\mathfrak a\}$ for a label $\mathfrak a$.
\item[(c)] For $f\in C^1_c(\mathbb R^N;V_{\mathbb C})$, $u\mapsto\tau_uf$ is continuous for $\|\cdot\|_{C^1}$.
\item[(d)] $\|\chi f\|_\sharp\le(\|\chi\|_\infty+40M\,{\rm Lip}\,\chi)\|f\|_\sharp$ for every bounded Lipschitz $\chi\colon\mathbb R^3\to\mathbb C$.
\item[(e)] $f\mapsto rf$ is complex linear; $(rr')f=r(r'f)$, $\overline{rf}=r\bar f$, $(rf)^\varepsilon=rf^\varepsilon$, ${\rm supp}\,rf=r\,{\rm supp}\,f$ and $\|rf\|_\sharp=\|f\|_\sharp$. Hence every identity between products of isometries, as maps of points, holds for their actions on functions, on tuples and on $\mathbb C\mathfrak S^V(\mathbb R^3)$.
\item[(f)] For a profile $\alpha$ and a label $\mathfrak a$, $\alpha\otimes\mathfrak a\in C^1_c(\mathbb R^3;V_{\mathbb C})$ is complex bilinear in $(\alpha,\mathfrak a)$ and real if both are; $\tau_y(\alpha\otimes\mathfrak a)=(\tau_y\alpha)\otimes\mathfrak a$; $\theta\overline{(\alpha\otimes\mathfrak a)}=(\bar\alpha\circ\theta)\otimes\mathfrak a^\theta$; $\|\alpha\otimes\mathfrak a\|_\infty\le\|\alpha\|_\infty\|\mathfrak a\|_\infty$, ${\rm Lip}(\alpha\otimes\mathfrak a)\le{\rm Lip}\,\alpha\,\|\mathfrak a\|_\infty+\|\alpha\|_\infty{\rm Lip}\,\mathfrak a$ and $\|\alpha\otimes\mathfrak a\|_{C^1}\le\sqrt2\|\alpha\|_{C^1}\|\mathfrak a\|_{C^1}$. The map $\mathfrak a\mapsto\mathfrak a^\theta$ is an antilinear involution with $\lambda_{\mathfrak a^\theta}=\lambda_{\mathfrak a}$, $\|\mathfrak a^\theta\|_{C^1}=\|\mathfrak a\|_{C^1}$ and $\|\mathfrak a^\theta\|_{L^1}=\|\mathfrak a\|_{L^1}$.
\item[(g)] Let $\chi$ be the function of the scalar partition of unity, $r>0$, and $\mathfrak a_k(s):=\chi(s/4r-k)\mathfrak a(s)$, $k\in\mathbb Z$. Then $\mathfrak a=\sum_k\mathfrak a_k$, ${\rm supp}\,\mathfrak a_k\subset[4r(k-1),4r(k+1)]$, $\|\mathfrak a_k\|_\infty\le\|\mathfrak a\|_\infty$, ${\rm Lip}\,\mathfrak a_k\le(1+\frac\pi{8r})\|\mathfrak a\|_{C^1}$, each $\mathfrak a_k$ is real if $\mathfrak a$ is, and fewer than $\lambda_{\mathfrak a}/4r+3$ of the $\mathfrak a_k$ are non-zero.
\end{enumerate}
\end{lemma}

\begin{proof}
Items (a), (b), (d) and (e) follow from $|E^\varepsilon|_2\le|E|_2$, $|\varrho(r)E|_2=|E|_2$, $\overline{\varrho(r)E}=\varrho(r)\bar E$ and from $f(x)-f(x')=\int_0^1\sum_\mu(x-x')_\mu\,\partial_\mu f(x'+s(x-x'))\,ds$; in (e), $((rr')f)(x)=\varrho(r^\circ)\varrho(r'^\circ)f(r'^{-1}r^{-1}x)=(r(r'f))(x)$, so the tensor action is a left action. Item (c) is the continuity of translations of scalar $C^1_c$ functions applied to the nine entries, together with $\|g\|_{C^1}\le\sum_{\lambda\lambda'}\|g_{\lambda\lambda'}\|_{C^1}$. In (f), $\theta$ keeps $x_2$, so
\[\bigl(\theta\overline{\alpha\otimes\mathfrak a}\bigr)(x)=\varrho(\theta)\overline{\alpha(\theta y)\mathfrak a(x_2)}=\bar\alpha(\theta y)\,\mathfrak a^\theta(x_2);\]
this is the one rule that is new against the scalar case. Further $|\nabla(\alpha\otimes\mathfrak a)|^2=|\nabla\alpha|^2|\mathfrak a|_2^2+|\alpha|^2|\mathfrak a'|_2^2\le2\|\alpha\|_{C^1}^2\|\mathfrak a\|_{C^1}^2$, and $\varrho(\theta)^2=1$ while $\varrho(\theta)$ and the bar keep $|\cdot|_2$. In (g), $0\le\chi\le1$, ${\rm supp}\,\chi=[-1,1]$, ${\rm Lip}\,\chi=\pi/2$ and $\sum_k\chi(s-k)=1$; the Lipschitz bound is that of (d) in one variable; and if $\mathfrak a_k\ne0$, the open interval $(4r(k-1),4r(k+1))$ meets an interval of length $\lambda_{\mathfrak a}$ containing ${\rm supp}\,\mathfrak a$, which confines $k$ to an open interval of length $\lambda_{\mathfrak a}/4r+2$.
\end{proof}

\begin{lemma}[Tubes]\label{C-353-17}
\begin{enumerate}
\item[(a)] $d_{\rm Tb}(\xi+\xi')\ge d_{\rm Tb}(\xi)+d_{\rm Tb}(\xi')$, $d_{\rm Tb}(c\xi)=c\,d_{\rm Tb}(\xi)$ for $c>0$, $d_{\rm Tb}(\xi+cn_0)=d_{\rm Tb}(\xi)+c$ for real $c$, and $d_{\rm Tb}(\xi+w)\ge d_{\rm Tb}(\xi)-|{\rm Re}\,w^0|-|{\rm Im}\,w^1|$. So ${\rm Tb}$ is an open convex cone, ${\rm Tb}+\overline{\rm Tb}\subset{\rm Tb}$ and $\overline{\rm Tb}=\{d_{\rm Tb}\ge0\}$. Moreover $d_{\rm Tb}(\xi^*)=d_{\rm Tb}(\bar\xi)=d_{\rm Tb}(\xi)\le d_{\rm Tb}({\rm Re}\,\xi)$; the star is additive and commutes with real factors, $n_0^*=n_0$ and $\xi^{**}=\xi$; $d_{\rm Tb}(y)=y^0$ at real $y$, so ${\rm Tb}\cap\mathbb R^2=\mathbb R^2_+$ and $d_{\rm Tb}(dn_0\pm y)=d\pm y^0$.
\item[(b)] The same holds for $d_3$, ${\rm Tb}_3$ and $\mathfrak z^*$, with 
\[
d_3(\mathfrak z+w)\ge d_3(\mathfrak z)-|{\rm Re}\,w_3|-|{\rm Im}\,w_1|-|{\rm Im}\,w_2|
\]
 and $d_3(de_3\pm x)=d\pm x_3$ at real $x$; besides $d_3(\xi\pm se_2)=d_{\rm Tb}(\xi)$ for $\xi\in\mathbb C^2$ and real $s$.
\item[(c)] $d_{\rm Tb}(\xi)=\min\{{\rm Re}\,\xi^+,{\rm Re}\,\xi^-\}$ and $(\xi^*)^\pm=\overline{\xi^\pm}$. For real $\tau$, 
\[
d_{\rm Tb}(R_{i\tau}\xi)=\min\{e^{-\tau}{\rm Re}\,\xi^+,e^\tau{\rm Re}\,\xi^-\}\ge e^{-|\tau|}d_{\rm Tb}(\xi)
\]
 on $\overline{\rm Tb}$; hence $R_{i\tau}{\rm Tb}={\rm Tb}$, and $d_{\rm Tb}(R_{i\tau}\xi)=e^{-|\tau|}\xi^0$ for real $\xi$ with $\xi^0>0$. Further $(R_{i\tau}\xi)^*=R_{i\tau}\xi^*$, $\xi^*=-\theta\xi$ for real $\xi$, $(R_{-\upsilon}w)^0=n_\upsilon\cdot w$ for real $w$, and $R_tR_\upsilon{\rm Tb}=R_{t+\upsilon}{\rm Tb}$ for real $t,\upsilon$.
\end{enumerate}
\end{lemma}

\begin{proof}
The real part is additive, $|{\rm Im}(u+v)|\le|{\rm Im}\,u|+|{\rm Im}\,v|$ and ${\rm Im}(-\bar u)={\rm Im}\,u$; this gives (a) and (b). For (c), ${\rm Re}\,\xi^\pm={\rm Re}\,\xi^0\mp{\rm Im}\,\xi^1$, and $R_{i\tau}$ multiplies $\xi^\pm$ by $e^{\mp\tau}$; on $\overline{\rm Tb}$ both real parts are $\ge0$, which gives the inequality. The remaining identities are direct computations with $(R_\zeta\xi)^\pm=e^{\pm i\zeta}\xi^\pm$.
\end{proof}

\begin{lemma}[Complex members]\label{C-353-18}
Let $\sigma$ be a tensor functional.
\begin{enumerate}
\item[(a)] If $X\cdot(H)\cdot Y$ and $X\cdot(H')\cdot Y$ are formed and $c\in\mathbb R$, then $\sigma(X\cdot(H+cH')\cdot Y)=\sigma(X\cdot(H)\cdot Y)+c\,\sigma(X\cdot(H')\cdot Y)$; $\sigma$ vanishes at every tuple with a member $0$; and $\overline{\sigma(F)}=\sigma(\bar F)$ for $F\in\mathbb C\mathfrak S^V(\mathbb R^3)$.
\item[(b)] $(H_1,\dots,H_n)_{\mathbb C}=(H_1)_{\mathbb C}\cdots(H_n)_{\mathbb C}$ and $r(H_1,\dots,H_n)_{\mathbb C}=(rH_1,\dots,rH_n)_{\mathbb C}$; under $\sigma$ the sum that defines $(H_1,\dots,H_n)_{\mathbb C}$ may be taken over all $\varepsilon\in\{0,1\}^n$.
\item[(c)] $(H_a)\mapsto\sigma((H_1,\dots,H_n)_{\mathbb C})$ is $n$-linear over $\mathbb C$ on $\prod_aC^1_c(O_a;V_{\mathbb C})$ for open $O_a$ with pairwise disjoint closures.
\item[(d)] $\overline{\sigma(A\cdot(H)_{\mathbb C}\cdot A')}=\sigma(\bar A\cdot(\bar H)_{\mathbb C}\cdot\bar A')$.
\end{enumerate}
\end{lemma}

\begin{proof}
(a) The set ${\rm supp}\,H\cup{\rm supp}\,H'$ and the supports of the members of $X$ and $Y$ are finitely many compact, pairwise disjoint sets, at mutual distances $\ge\mathsf d$ for some $\mathsf d>0$; their open $\mathsf d/3$-neighbourhoods have pairwise disjoint closures, and the multilinearity of $\sigma$ applies; $\sigma$ is real on tuples. (b) follows from Lemma~\ref{C-353-16}\,(e) and from (a). (c) The map is additive and real homogeneous by (a), and $(iH)^0=-H^1$, $(iH)^1=H^0$ give $u((iH)^0)+i\,u((iH)^1)=i\,(u(H^0)+i\,u(H^1))$ for the value $u$ of $\sigma$ as a function of the $a$-th member. (d) If $H^1=0$, both sides equal $\sigma(\bar A\cdot(H)\cdot\bar A')$; if $H^1\ne0$, the left side is $\sigma(\bar A\cdot((H^0)-i(H^1))\cdot\bar A')$ and $(\bar H)_{\mathbb C}=(H^0)+i(-H^1)$, and (a) applies. The restriction of the sum in Definition~\ref{C-353-3}\,(2) matters in $\mathbb C\mathfrak S^V$ itself only, where a tuple with a member $0$ is a basis vector and not the zero vector.
\end{proof}

\begin{lemma}[The form, the space, the reconstruction]\label{C-353-19}
\begin{enumerate}
\item[(a)] Let $\sigma$ have (T3). The form $\langle F,G\rangle_\sigma=\sigma(\theta\bar F\cdot G)$ on $\mathbb C\mathfrak S^V(E_+)$ is positive semi-definite and Hermitian, and $\|\Omega\|=1$. For open sets $O_1,\dots,O_n$ whose closures are compact, pairwise disjoint and contained in $E_+$, $(H_a)\mapsto[(H_1,\dots,H_n)_{\mathbb C}]$ is $n$-linear over $\mathbb C$ on $\prod_aC^1_c(O_a;V_{\mathbb C})$; the class of a tuple with a member $0$ is $0$.
\item[(b)] Let $\sigma$ have (T1$^\circ$), (T2) and (T3). There are exactly one family of bounded operators $\mathcal T_t$, $t\ge0$, and exactly one family $\mathcal U(v)$, $v\in e_3^\perp$, with $\mathcal T_t[F]=[\tau_{te_3}F]$ and $\mathcal U(v)[F]=[\tau_vF]$; $(\mathcal T_t)$ is a strongly continuous semigroup of self-adjoint operators with $0\le\mathcal T_t\le1$, $(\mathcal U(v))$ is a strongly continuous unitary group, all of them commute and fix $\Omega$, and there is exactly one projection-valued measure $\mathcal P$ on $[0,\infty)\times e_3^\perp$ with $\mathcal T_t\,\mathcal U(v)=\int e^{-tE+iv\cdot q}\,d\mathcal P(E,q)$; moreover $\mathcal P(\{0\})\Omega=\Omega$.
\item[(c)] Let $\sigma$ have (T1$^\circ$) and (T3). Then $\sigma(\theta_cX)=\sigma(X)$ for every $X\in\mathfrak S^V(\mathbb R^3)$ none of whose supports meets $\{x_3=c\}$.
\end{enumerate}
\end{lemma}

\begin{proof}
(a) For $X,Y\in\mathfrak S^V(E_+)$ the tuple $\theta X\cdot Y$ is formed, as $\theta$ maps $E_+$ onto $\{x_3<0\}$. The form is sesquilinear, and its quadratic form $\mathsf q(F)=\langle F,F\rangle_\sigma$ is real and $\ge0$ by (T3); by polarisation, $\langle F,G\rangle_\sigma=\frac14\sum_{k=0}^3i^{-k}\mathsf q(F+i^kG)$, so it is Hermitian. $\|\Omega\|^2=\sigma(\emptyset)=1$. For multilinearity let $X[u]$ be the tuple whose $a$-th member is $u$, $u_1=H_a+cH'_a$, $u_2=H_a$, $u_3=H'_a$, $(\gamma_i)=(1,-1,-c)$ and $Z=\sum_i\gamma_iX[u_i]$; then $\|[Z]\|^2=\sum_{i,k}\gamma_i\gamma_k\,\sigma(\theta X[u_i]\cdot X[u_k])$, and $(u,u')\mapsto\sigma(\theta X[u]\cdot X[u'])$ is bilinear because the $2n$ open sets $\theta O_b$, $O_b$ have pairwise disjoint closures and $u\mapsto\theta u$ is linear; so $\|[Z]\|^2=0$. Complex members are treated as in Lemma~\ref{C-353-18}\,(c).

(b) We apply the abstract reconstruction theorem to $\mathbb C\mathfrak S^V(E_+)$ with the form of (a), the vector $\emptyset$, the maps $\tau_{te_3}$, $t\ge0$, and $\tau_v$, $v\in e_3^\perp$. As maps of points $\tau_{te_3}\theta\tau_{te_3}=\theta$ and $\theta\tau_v=\tau_v\theta$; with Lemma~\ref{C-353-16}\,(e) and (T1$^\circ$) the maps are symmetric, respectively isometric, for the form. For $F=\sum\zeta_iX_i$ there are $R<\infty$ and $\mathsf d>0$ such that every tuple $\theta X_i\cdot\tau_{te_3}X_k$, $t\ge0$, is of type $(\mathsf d;R)$, so that
\[|\langle F,\tau_{te_3}F\rangle_\sigma|\le\sum_{i,k}|\zeta_i||\zeta_k|\max\{1,C^\sigma_{n_i+n_k}(\mathsf d;R)\}\prod_{H\in X_i}\|H\|_\sharp\prod_{H\in X_k}\|H\|_\sharp,\]
a bound that does not see $t$. For weak continuity, with $Y=(f_1,\dots,f_n)$,
\[\sigma(\theta X\cdot\tau_uY)-\sigma(\theta X\cdot Y)=\sum_a\sigma\bigl(\theta X\cdot(\tau_uf_1,\dots,\tau_uf_{a-1},\tau_uf_a-f_a,f_{a+1},\dots,f_n)\bigr)\]
by Lemma~\ref{C-353-18}\,(a); for $|u|<\mathsf d/3$ each tuple on the right is of type $(\mathsf d/3;R+\mathsf d/3)$, and (T2) with $\|\tau_uf_a-f_a\|_\sharp\le40M\|\tau_uf_a-f_a\|_{C^1}\to0$ (Lemma~\ref{C-353-16}\,(b), (c)) gives the continuity. The theorem gives $\mathcal T_t$ and $\mathcal U(v)$; the statement on the joint spectral measure gives $\mathcal P$ and $\mathcal P(\{0\})\Omega=\Omega$.

(c) By translation invariance, with $\theta_c=\tau_{ce_3}\theta\tau_{-ce_3}$, we may take $c=0$. Let $X=(H_1,\dots,H_n)$. The support of each $H_a$ is compact and does not meet the closed plane $\{x_3=0\}$, so it lies at a positive distance from it, and $H_a$ vanishes near the plane; hence $H_a^\pm:=H_a\mathbf 1_{\{\pm x_3>0\}}$ lie in $C^1_c(\mathbb R^3;V)$, ${\rm supp}\,H_a^\pm\subset{\rm supp}\,H_a$, $H_a=H_a^++H_a^-$ and $\theta H_a=\theta H_a^++\theta H_a^-$. So all the tuples $(H_1^{s_1},\dots,H_n^{s_n})$, $s\in\{+,-\}^n$, and their images under $\theta$ are formed, and Lemma~\ref{C-353-18}\,(a), applied member after member, gives
\[\sigma(X)=\sum_{s\in\{+,-\}^n}\sigma(H_1^{s_1},\dots,H_n^{s_n}),\qquad\sigma(\theta X)=\sum_{s\in\{+,-\}^n}\sigma(\theta H_1^{s_1},\dots,\theta H_n^{s_n}).\]
In each term every member has its support in one of the open half-spaces $\{x_3>0\}$ and $\{x_3<0\}$ (the member $0$ is admitted); so, by symmetry, the term of $X$ is $\sigma(\theta A\cdot B)$ with real $A,B\in\mathfrak S^V(E_+)$, and the term of $\theta X$ with the same $s$ is $\sigma(\theta(\theta A\cdot B))$. As the form is Hermitian and $\sigma$ is real on tuples,
\[\sigma(\theta A\cdot B)=\langle A,B\rangle_\sigma=\langle B,A\rangle_\sigma=\sigma(\theta B\cdot A)=\sigma(A\cdot\theta B)=\sigma(\theta(\theta A\cdot B)),\]
the last two equalities by symmetry and $\theta\theta=1$; so the two sums agree term by term, and $\sigma(\theta X)=\sigma(X)$.
\end{proof}

\begin{lemma}[The semigroups of the cone]\label{C-353-20}
\begin{enumerate}
\item[(a)] Let $\sigma$ have (T1$^\circ$), (T2), (T3) and (T4). Then $\mathcal P$ is carried by $\mathsf S_1:=\{E\ge|q_1|\}$, on which $|e^{-\xi^0E+i\xi^1q_1}|\le e^{-d_{\rm Tb}(\xi)E}$. The operators $T(\xi)$, $\xi\in\overline{\rm Tb}$, satisfy $\|T(\xi)\|\le1$, $T(\xi)T(\xi')=T(\xi+\xi')$, $T(0)=1$, $T(\xi)^*=T(\xi^*)$ and $T(\xi)\Omega=\Omega$; $\xi\mapsto T(\xi)$ is strongly continuous on $\overline{\rm Tb}$ and holomorphic in norm on ${\rm Tb}$; and $T(y)=\mathcal T_{y^0}\mathcal U(y^1e_1)$, $T(y)[F]=[\tau_yF]$ for real $y$ with $y^0\ge0$.
\item[(b)] Let $\sigma$ have (T1), (T2), (T3) and (T4). For $w\in W_3$ with $we_3=e_3$, $\mathcal U_w[F]:=[wF]$ is a unitary operator that fixes $\Omega$, commutes with $\mathcal T_t$ and satisfies $\mathcal U_w\mathcal U(v)\mathcal U_w^{-1}=\mathcal U(wv)$; hence $\mathcal U_w\mathcal P(B)\mathcal U_w^{-1}=\mathcal P(\{(E,wq):(E,q)\in B\})$, $\mathcal P(\{E<|q_2|\})=0$, and $\mathcal P$ is carried by $\mathsf S:=\{E\ge|q_1|,\ E\ge|q_2|\}$, on which $|e^{-\mathfrak z_3E+i\mathfrak z_1q_1+i\mathfrak z_2q_2}|\le e^{-d_3(\mathfrak z)E}$. The operators $T_3(\mathfrak z)$, $\mathfrak z\in\overline{{\rm Tb}_3}$, have the properties of (a), with $\mathfrak z^*$, $d_3$, and $T_3(x)[F]=[\tau_xF]$ for real $x$ with $x_3\ge0$; $T(\xi)=T_3(\xi^1,0,\xi^0)$; and $\mathcal U_wT_3(\mathfrak z)\mathcal U_w^{-1}=T_3(w\mathfrak z)$. For $w\in W_3$ with $we_3=e_\mu$, $\sigma$ is reflection positive at the mirror $\{x_\mu=0\}$, and $J_w[F]:=[wF]^{(\mu)}$ is a unitary map of $\mathcal H(\sigma)$ onto the space $\mathcal H^{(\mu)}(\sigma)$ built with that mirror, which carries $\Omega$, $\mathcal T_t$, $\mathcal U(v)$ and $\mathcal P$ to the vacuum, the time and space translations and the joint spectral measure of $\mathcal H^{(\mu)}(\sigma)$, the last through $(E,q)\mapsto(E,wq)$.
\end{enumerate}
\end{lemma}

\begin{proof}
(a) By (T4), $\mathsf S_1$ carries $\mathcal P$, and on $\mathsf S_1$, $|e^{-\xi^0E+i\xi^1q_1}|=e^{-{\rm Re}\,\xi^0E-{\rm Im}\,\xi^1q_1}\le e^{-d_{\rm Tb}(\xi)E}\le1$ for $\xi\in\overline{\rm Tb}$. The algebraic rules follow from the spectral calculus and the rule $\int\mathbf 1_B\,d\mathcal P=\mathcal P(B)$, the adjoint from $\overline{e^{-\xi^0E+i\xi^1q_1}}=e^{-\overline{\xi^0}E-i\overline{\xi^1}q_1}$, and $T(\xi)\Omega=\Omega$ from $\Omega=\mathcal P(\{0\})\Omega$. Strong continuity follows by dominated convergence from $\|(T(\xi_j)-T(\xi))\psi\|^2=\int_{\mathsf S_1}|e_{\xi_j}-e_\xi|^2\,d\langle\psi,\mathcal P\psi\rangle$, $e_\xi(E,q)=e^{-\xi^0E+i\xi^1q_1}$, the integrand being $\le4$; holomorphy in norm, because the matrix elements are holomorphic and the family is bounded; the real points by Lemma~\ref{C-353-19}\,(b).

(b) As $w$ commutes with $\theta$ and maps $E_+$ onto itself, (T1) gives $\sigma(\theta\overline{wF}\cdot wG)=\sigma(w(\theta\bar F\cdot G))=\sigma(\theta\bar F\cdot G)$, so $\mathcal U_w$ is well defined and isometric, and onto because $w^{-1}$ is of the same kind; the relations with $\mathcal T_t$ and $\mathcal U(v)$ hold on classes, and the uniqueness of the joint spectral measure gives the transformation of $\mathcal P$. Thus $B\mapsto\mathcal U_w\mathcal P(B)\mathcal U_w^{-1}=\mathcal P(\{(E,wq):(E,q)\in B\})$ is the image of $\mathcal P$ under $(E,q)\mapsto(E,w^{-1}q)$, and for every bounded Borel function $g$
\[\mathcal U_w\,g(\mathcal P)\,\mathcal U_w^{-1}=\int g(E,w^{-1}q)\,d\mathcal P(E,q).\]
For the $w$ that exchanges $e_1$ and $e_2$, $\{E<|q_2|\}$ is the image of $\{E<|q_1|\}$. The assertions on $T_3$ that repeat those of (a) are (a) in three variables. For the rule $\mathcal U_wT_3(\mathfrak z)\mathcal U_w^{-1}=T_3(w\mathfrak z)$ take $g=\mathbf 1_{\mathsf S}e_{\mathfrak z}$, $e_{\mathfrak z}(E,q):=e^{-\mathfrak z_3E+i\mathfrak z_1q_1+i\mathfrak z_2q_2}$, $\mathfrak z\in\overline{{\rm Tb}_3}$, $q=q_1e_1+q_2e_2$. As $w$ is a real orthogonal matrix with $we_3=e_3$, it maps $e_3^\perp$ onto itself, $(w\mathfrak z)_3=\mathfrak z_3$ and $\mathfrak z\cdot(w^{-1}q)=(w\mathfrak z)\cdot q$, the dot being the bilinear product; so $e_{\mathfrak z}(E,w^{-1}q)=e_{w\mathfrak z}(E,q)$. The set $\mathsf S$ is stable under $(E,q)\mapsto(E,wq)$, and $d_3(w\mathfrak z)=d_3(\mathfrak z)$, because $(w\mathfrak z)_1,(w\mathfrak z)_2$ are $\mathfrak z_1,\mathfrak z_2$ up to order and signs; so $w\mathfrak z\in\overline{{\rm Tb}_3}$ and $\mathcal U_wT_3(\mathfrak z)\mathcal U_w^{-1}=\int\mathbf 1_{\mathsf S}e_{w\mathfrak z}\,d\mathcal P=T_3(w\mathfrak z)$ on all of $\overline{{\rm Tb}_3}$. For $we_3=e_\mu$, $\theta^{(\mu)}:=w\theta w^{-1}$ is the reflection in $\{x_\mu=0\}$, and by (T1) $\sigma(\theta^{(\mu)}\bar F\cdot F)=\sigma(\theta\overline{w^{-1}F}\cdot w^{-1}F)\ge0$ for $F$ over $\{x_\mu>0\}$; the same computation shows that $J_w$ is isometric and onto.
\end{proof}

\begin{lemma}[Slab operators]\label{C-353-21}
Let $\sigma$ have (T1$^\circ$), (T2), (T3) and (T5$_N$), and let $l>0$.
\begin{enumerate}
\item[(a)] For every $H\in C^1_c(\{0<x_3<l\};V_{\mathbb C})$ there is exactly one bounded operator $B_l(H)$ on $\mathcal H(\sigma)$ with $B_l(H)[F]=[(H)_{\mathbb C}\cdot\tau_{le_3}F]$ for $F\in\mathbb C\mathfrak S^V(E_+)$. $H\mapsto B_l(H)$ is complex linear, and $\|B_l(H)\|\le C(\mathcal K,l)\|H\|_\sharp$ if ${\rm supp}\,H\subset\mathcal K$, $\mathcal K\subset\{0<x_3<l\}$ compact, with a finite number $C(\mathcal K,l)$ that depends only on $\mathcal K$, $l$, $M$, $c_I$ and $N$.
\item[(b)] $B_l(H)=\mathcal T_cB_{c'-c}(\tau_{-ce_3}H)\mathcal T_{l-c'}$ for $0\le c<c'\le l$ and ${\rm supp}\,H\subset\{c<x_3<c'\}$; $\mathcal U(v)B_l(H)\mathcal U(-v)=B_l(\tau_vH)$ for $v\in e_3^\perp$.
\item[(c)] $B_l(H)^*=B_l(\overline{\theta_{l/2}H})$, with the tensor action of $\theta_{l/2}$.
\item[(d)] $\|B_{\eta/2}(H)\|\le N(\eta)\|H\|_\sharp$ for $\eta\in(0,1]$ and real $H$ with support in a ball $B(y,c_I\eta)$, $\eta/8\le y_3\le3\eta/8$.
\item[(e)] $[(H_1,\dots,H_n)_{\mathbb C}]=B_{c_1}(H_1)B_{c_2-c_1}(\tau_{-c_1e_3}H_2)\cdots B_{c_n-c_{n-1}}(\tau_{-c_{n-1}e_3}H_n)\Omega$ for $0=c_0<c_1<\dots<c_n$ and ${\rm supp}\,H_a\subset\{c_{a-1}<x_3<c_a\}$.
\end{enumerate}
\end{lemma}

\begin{proof}
(d) Put $H_\circ:=\tau_{-(\eta/4)e_3}H$, whose support lies in a ball $B(y',c_I\eta)$ with $|y'_3|\le\eta/8$, and $b(F,G):=\sigma(\theta\overline{\tau_{(\eta/4)e_3}F}\cdot(H_\circ)\cdot\tau_{(\eta/4)e_3}G)$ for $F,G\in\mathbb C\mathfrak S^V(E_+)$. By (T5$_N$), $|b(F,G)|\le N(\eta)\|H\|_\sharp\|[F]\|\,\|[G]\|$; this is the one place of the whole argument at which the insertion bound is used. So $b$ is a bounded form on the classes and is given by a bounded operator $A$ with $\|A\|\le N(\eta)\|H\|_\sharp$. By translation invariance and $\tau_{(\eta/4)e_3}\theta\tau_{(\eta/4)e_3}=\theta$, $b(F,G)=\langle[F],[(H)\cdot\tau_{(\eta/2)e_3}G]\rangle$; so $A=B_{\eta/2}(H)$.

(a) Uniqueness is clear, the classes being dense. For real $H$ with support in $\mathcal K$ let $\delta:=\min\{\min_{\mathcal K}x_3,l-\max_{\mathcal K}x_3\}$, $\eta:=\min\{1,4\delta\}$, $\mathsf h:=c_I\eta/2$, and let $\chi_{\mathsf h,k}$ be the scalar partition of unity at scale $\mathsf h$. For each of the finitely many $k$ with ${\rm supp}\,\chi_{\mathsf h,k}\cap\mathcal K\ne\emptyset$ choose $c_k\in[0,l-\eta/2]$ with $|\mathsf hk_3-\eta/4-c_k|\le\mathsf h$; then $\tau_{-c_ke_3}(\chi_{\mathsf h,k}H)$ is real, with support in a ball of radius $\sqrt3\mathsf h<c_I\eta$ whose centre has its third coordinate in $[\eta/8,3\eta/8]$, and the identity $(H_k)\cdot\tau_{le_3}G=\tau_{c_ke_3}[(\tau_{-c_ke_3}H_k)\cdot\tau_{(\eta/2)e_3}(\tau_{(l-c_k-\eta/2)e_3}G)]$ in $\mathbb C\mathfrak S^V(E_+)$, together with (a) of Lemma~\ref{C-353-19}, gives
\[B_l(H)=\sum_k\mathcal T_{c_k}B_{\eta/2}(\tau_{-c_ke_3}\chi_{\mathsf h,k}H)\,\mathcal T_{l-c_k-\eta/2},\qquad\|B_l(H)\|\le\sum_kN(\eta)\|\chi_{\mathsf h,k}H\|_\sharp,\]
which is bounded by a multiple of $\|H\|_\sharp$ by (d), Lemma~\ref{C-353-16}\,(d) and ${\rm Lip}\,\chi_{\mathsf h,k}\le\sqrt3\pi/2\mathsf h$. For complex $H$ put $B_l(H):=B_l(H^0)+iB_l(H^1)$ and use Lemma~\ref{C-353-16}\,(a).

(b) and (e) follow from identities in $\mathbb C\mathfrak S^V(E_+)$: $(H)_{\mathbb C}\cdot\tau_{le_3}G=\tau_{ce_3}[(\tau_{-ce_3}H)_{\mathbb C}\cdot\tau_{(c'-c)e_3}(\tau_{(l-c')e_3}G)]$, and, for (e), induction on $n$ with $[(H_1,\dots,H_n)_{\mathbb C}]=B_{c_1}(H_1)[\tau_{-c_1e_3}(H_2,\dots,H_n)_{\mathbb C}]$.

(c) Let $B=(H)_{\mathbb C}\cdot\tau_{le_3}G$. The form is Hermitian, so $\langle F,B\rangle_\sigma=\overline{\langle B,F\rangle_\sigma}$, and by one translation, $\tau_{le_3}\theta=\theta_{l/2}$, $\tau_{le_3}\theta\tau_{le_3}=\theta$, symmetry and Lemma~\ref{C-353-18}\,(d),
\[\langle B,F\rangle_\sigma=\sigma\bigl(\theta\overline{(H)_{\mathbb C}}\cdot\theta\tau_{le_3}\bar G\cdot F\bigr)=\sigma\bigl(\theta_{l/2}\overline{(H)_{\mathbb C}}\cdot\theta\bar G\cdot\tau_{le_3}F\bigr)=\sigma\bigl(\theta\bar G\cdot(\overline{\theta_{l/2}H})_{\mathbb C}\cdot\tau_{le_3}F\bigr),\]
which is $\langle[G],B_l(\overline{\theta_{l/2}H})[F]\rangle$. No invariance of $\sigma$ under a reflection is used: the matrix index is moved only by the $\theta$ of the form.
\end{proof}

\begin{lemma}[Centred slab operators]\label{C-353-22}
Let $\sigma$ be as in Lemma~\ref{C-353-21}, $d>0$ and $f\in C^1_c(\{|x_3|<d\};V_{\mathbb C})$; then $B^c_d(f)=B_{2d}(\tau_{de_3}f)$.
\begin{enumerate}
\item[(a)] $f\mapsto B^c_d(f)$ is complex linear, and $B^c_d(f)^*=B^c_d(\overline{\theta f})$.
\item[(b)] $\mathcal T_{c+x_3}\mathcal U(\vec x)B^c_d(f)\,\mathcal U(-\vec x)\mathcal T_{c-x_3}=B^c_{d+c}(\tau_xf)$ for $x\in\mathbb R^3$, $\vec x=x_1e_1+x_2e_2$, $c\ge|x_3|$; so $T(cn_0+y)B^c_d(f)T(cn_0-y)=B^c_{d+c}(\tau_yf)$ for $|y^0|\le c$ under (T4), and $T_3(ce_3+x)B^c_d(f)T_3(ce_3-x)=B^c_{d+c}(\tau_xf)$ for $|x_3|\le c$ under (T1) and (T4).
\item[(c)] $\|B^c_d(h)\|\le40M\,C(de_3+\mathcal K,2d)\|h\|_{C^1}$ if ${\rm supp}\,h\subset\mathcal K$, $\mathcal K\subset\{|x_3|<d\}$ compact.
\item[(d)] $\|B^c_{\eta/4}(h)\|\le N(\eta)\|h\|_\sharp$ for $\eta\in(0,1]$ and real $h$ with support in a ball $B(v,c_I\eta)$, $v_3=0$.
\item[(e)] $\langle[\tau_{-de_3}F],B^c_d(f)[\tau_{-de_3}F']\rangle=\sigma(\theta\bar F\cdot(f)_{\mathbb C}\cdot F')$ for $F,F'\in\mathbb C\mathfrak S^V(\{x_3>d\})$.
\end{enumerate}
\end{lemma}

\begin{proof}
(a) is Lemma~\ref{C-353-21}\,(c) at $l=2d$, with $\theta_d\tau_{de_3}=\tau_{de_3}\theta$. (b) By Lemma~\ref{C-353-21}\,(b) the left side is $\mathcal T_{c+x_3}B_{2d}(\tau_{\vec x}\tau_{de_3}f)\mathcal T_{c-x_3}$, which the same clause, with $l=2(d+c)$ and the slab $\{c+x_3<x_3'<c+x_3+2d\}$, identifies with $B_{2(d+c)}(\tau_{(d+c)e_3}\tau_xf)=B^c_{d+c}(\tau_xf)$; the forms with $T$ and $T_3$ follow from the values at real points in Lemma~\ref{C-353-20}. (c) is Lemma~\ref{C-353-21}\,(a) with Lemma~\ref{C-353-16}\,(b). (d) is Lemma~\ref{C-353-21}\,(d) after the translation by $(\eta/4)e_3$, as $\eta/4\in[\eta/8,3\eta/8]$. (e) follows from one translation by $-de_3$ and $\tau_{-de_3}\theta\tau_{-de_3}=\theta$.
\end{proof}

\begin{proposition}[Regularised point plates; a bound that does not depend on the regularisation]\label{C-353-23}
Let $\sigma$ have (T1$^\circ$), (T2), (T3), (T4) and (T5$_N$). Let $j_\delta:=j^{(2)}_\delta$, and for $\mathfrak a\in\mathcal L_{\mathbb C}$, $\delta>0$ and $d_{\rm Tb}(\xi),d_{\rm Tb}(\xi')\ge\delta$ put
\[A_\delta(\mathfrak a):=B^c_\delta(j_\delta\otimes\mathfrak a),\qquad B^\delta_{\mathfrak a}(\xi,\xi'):=T(\xi-\delta n_0)\,A_\delta(\mathfrak a)\,T(\xi'-\delta n_0).\]
\begin{enumerate}
\item[(a)] $B^\delta_{\mathfrak a}$ is complex linear in $\mathfrak a$ and holomorphic in norm on $\{d_{\rm Tb}>\delta\}^2$; $T(\kappa)B^\delta_{\mathfrak a}(\xi,\xi')T(\kappa')=B^\delta_{\mathfrak a}(\kappa+\xi,\xi'+\kappa')$ for $\kappa,\kappa'\in\overline{\rm Tb}$; $A_\delta(\mathfrak a)^*=A_\delta(\mathfrak a^\theta)$ and $B^\delta_{\mathfrak a}(\xi,\xi')^*=B^\delta_{\mathfrak a^\theta}(\xi'^*,\xi^*)$; $B^\delta_{\mathfrak a}(dn_0+y,dn_0-y)=B^c_d((\tau_yj_\delta)\otimes\mathfrak a)$ for $d\ge\delta+|y^0|$; and $\int\alpha(y)B^\delta_{\mathfrak a}(dn_0+y,dn_0-y)\,dy=B^c_d((\alpha*j_\delta)\otimes\mathfrak a)$ for continuous $\alpha$ with compact support in $\{|y^0|<d-\delta\}$.
\item[(b)] For $0<\eta\le\min\{1,d_{\rm Tb}(\xi),d_{\rm Tb}(\xi')\}$ and $0<\delta\le c_I\eta/8$, $\|B^\delta_{\mathfrak a}(\xi,\xi')\|\le2c_{\rm plt}\,\nu(\mathfrak a)N(\eta)\eta^{-4}$, without the factor $2$ for real $\mathfrak a$; if $\sigma$ has (T5$_N$) in the sup norm, also 
\[
\|B^\delta_{\mathfrak a}(\xi,\xi')\|\le2\cdot12\cdot8^2{\rm mv}_0^2c_I^{-3}\nu_\infty(\mathfrak a)N(\eta)\eta^{-3}.
\]

\item[(c)] For $\varepsilon>0$, $\eta_0:=\frac12\min\{1,\varepsilon\}$, $b_\varepsilon:=2c_{\rm plt}\nu(\mathfrak a)N(\eta_0)\eta_0^{-4}$ and $0<\delta,\delta'\le c_I\eta_0/8$: 
\[
\|B^\delta_{\mathfrak a}(\varepsilon n_0,\varepsilon n_0)-B^{\delta'}_{\mathfrak a}(\varepsilon n_0,\varepsilon n_0)\|\le c\,b_\varepsilon\varepsilon^{-1}(\delta+\delta')
\]
 with an absolute number $c$. Hence $B_{\mathfrak a}(\xi,\xi'):=\lim_{\delta\to0}B^\delta_{\mathfrak a}(\xi,\xi')$ exists in norm on ${\rm Tb}\times{\rm Tb}$, uniformly on $\{d_{\rm Tb}\ge\varepsilon\}^2$, and has the properties (i)--(iv) of Theorem~\ref{C-353-7}\,(A); the bounds of (b) pass to it.
\end{enumerate}
\end{proposition}

\begin{proof}
(a) By Lemma~\ref{C-353-16}\,(f), $\theta\overline{(j_\delta\otimes\mathfrak a)}=(j_\delta\circ\theta)\otimes\mathfrak a^\theta=j_\delta\otimes\mathfrak a^\theta$, so Lemma~\ref{C-353-22}\,(a) gives $A_\delta(\mathfrak a)^*=A_\delta(\mathfrak a^\theta)$ in place of the scalar rule $A_\delta^*=A_\delta$; with $T(\xi)^*=T(\xi^*)$ and $(\xi-\delta n_0)^*=\xi^*-\delta n_0$ this is the adjoint rule. The shift rule and holomorphy come from Lemma~\ref{C-353-20}\,(a) and the calculus of operator functions. The value at $(dn_0+y,dn_0-y)$ is Lemma~\ref{C-353-22}\,(b) with $c=d-\delta$; the integral identity follows from it, the convergence of Riemann sums of $\alpha*j_\delta$ in $\|\cdot\|_{C^1}$, and Lemma~\ref{C-353-22}\,(c).

(b) Let $\mathfrak a$ be real, $r:=c_I\eta/8$, $\psi,\psi'$ unit vectors, and $u(z):=\langle\psi,B^\delta_{\mathfrak a}(\xi+z,\xi'-z)\psi'\rangle$ for $|z^0|,|z^1|<2r$; there $d_{\rm Tb}(\xi\pm z)>\eta-4r>\delta$, so $u$ is holomorphic near the closed polydisc of radius $r$. By the mean-value formula, $u(0)=\int_{(-r,r)^2}dt\int u(y+it)\alpha^{(2)}_{r,t}(y)\,dy$. For fixed $t$ let $\kappa_t:=\xi-\frac\eta4n_0+it$ and $\kappa'_t:=\xi'-\frac\eta4n_0-it$, which lie in ${\rm Tb}$; by the shift rule $u(y+it)=\langle\psi_t,B^\delta_{\mathfrak a}(\frac\eta4n_0+y,\frac\eta4n_0-y)\psi'_t\rangle$ with $\psi_t:=T(\kappa_t)^*\psi$, $\psi'_t:=T(\kappa'_t)\psi'$ of norms $\le1$, and by (a) at $d=\eta/4$
\[\int u(y+it)\,\alpha^{(2)}_{r,t}(y)\,dy=\langle\psi_t,B^c_{\eta/4}(\gamma_t\otimes\mathfrak a)\psi'_t\rangle,\qquad\gamma_t:=\alpha^{(2)}_{r,t}*j_\delta.\]
The profile $\gamma_t$ is real, supported in $[-2r,2r]^2$, with $\|\gamma_t\|_\infty\le{\rm mv}_0^2r^{-4}$ and ${\rm Lip}\,\gamma_t\le\sqrt2\,{\rm mv}_0{\rm mv}_1r^{-5}$. We cut the label as in Lemma~\ref{C-353-16}\,(g), at the scale $r$; at most $3(\lambda_{\mathfrak a}+1)/(c_I\eta)$ pieces $\mathfrak a_k$ are non-zero, and each $\gamma_t\otimes\mathfrak a_k$ is real, supported in the closed ball of radius $\sqrt{24}\,r<8r=c_I\eta$ about $4rke_2$, with $\|\gamma_t\otimes\mathfrak a_k\|_\sharp\le40M\,{\rm mv}_2r^{-5}\|\mathfrak a\|_{C^1}$, since ${\rm Lip}(\gamma_t\otimes\mathfrak a_k)\le\sqrt2\,{\rm mv}_0{\rm mv}_1r^{-5}\|\mathfrak a\|_{C^1}+{\rm mv}_0^2r^{-4}(1+\frac\pi{8r})\|\mathfrak a\|_{C^1}\le{\rm mv}_2r^{-5}\|\mathfrak a\|_{C^1}$. Lemma~\ref{C-353-22}\,(d) bounds each piece, and
\[
\begin{multlined}
|u(0)|\le(2r)^2\cdot\frac{3(\lambda_{\mathfrak a}+1)}{c_I\eta}\cdot N(\eta)\cdot40M\,{\rm mv}_2r^{-5}\|\mathfrak a\|_{C^1}=480\cdot8^3{\rm mv}_2M(c_I\eta)^{-4}N(\eta)\,\nu(\mathfrak a) \\
=c_{\rm plt}\,\nu(\mathfrak a)N(\eta)\eta^{-4}.
\end{multlined}
\]
The supremum over $\psi,\psi'$ is the norm. A complex label is split as $\mathfrak a=\mathfrak a^0+i\mathfrak a^1$, whence the factor $2$. In the sup norm the slope of $\gamma_t$ is not paid: $\|B^c_{\eta/4}(\gamma_t\otimes\mathfrak a_k)\|\le N(\eta){\rm mv}_0^2r^{-4}\|\mathfrak a\|_\infty$, and the same sum gives $12\cdot8^2{\rm mv}_0^2(c_I\eta)^{-3}N(\eta)\nu_\infty(\mathfrak a)$.

(c) For $|z^0|,|z^1|<\varepsilon/4$, $d_{\rm Tb}(\varepsilon n_0\pm z)>\varepsilon/2\ge\eta_0$, so by (b) at the one number $\eta_0$ the holomorphic function $z\mapsto B^\delta_{\mathfrak a}(\varepsilon n_0+z,\varepsilon n_0-z)$ is bounded by $b_\varepsilon$ there, and the Cauchy estimate gives $\|B^\delta_{\mathfrak a}(\varepsilon n_0+y,\varepsilon n_0-y)-B^\delta_{\mathfrak a}(\varepsilon n_0,\varepsilon n_0)\|\le c\,b_\varepsilon\varepsilon^{-1}|y|$ for real $|y|\le\varepsilon/8$. Integrating against $j_{\delta'}$, which has mass $1$ and support in $\{|y|\le\delta'/2\}$, and using (a), $\|B^c_\varepsilon((j_{\delta'}*j_\delta)\otimes\mathfrak a)-B^\delta_{\mathfrak a}(\varepsilon n_0,\varepsilon n_0)\|\le c\,b_\varepsilon\varepsilon^{-1}\delta'$; the same with $\delta$ and $\delta'$ exchanged, and $j_{\delta'}*j_\delta=j_\delta*j_{\delta'}$, give the estimate.

As the space of bounded operators is complete, $B_{[\varepsilon]}:=\lim_\delta B^\delta_{\mathfrak a}(\varepsilon n_0,\varepsilon n_0)$ exists, and for $d_{\rm Tb}(\xi),d_{\rm Tb}(\xi')\ge\varepsilon$ the shift rule gives $B^\delta_{\mathfrak a}(\xi,\xi')=T(\xi-\varepsilon n_0)B^\delta_{\mathfrak a}(\varepsilon n_0,\varepsilon n_0)T(\xi'-\varepsilon n_0)$, so $B_{\mathfrak a}(\xi,\xi')=T(\xi-\varepsilon n_0)B_{[\varepsilon]}T(\xi'-\varepsilon n_0)$, with $\|B^\delta_{\mathfrak a}(\xi,\xi')-B_{\mathfrak a}(\xi,\xi')\|\le\|B^\delta_{\mathfrak a}(\varepsilon n_0,\varepsilon n_0)-B_{[\varepsilon]}\|$. Holomorphy, linearity, the bound (i), the shift rule and the adjoint (ii) pass to the limit. For (iii), let ${\rm supp}\,\alpha\subset\{|y^0|\le d'\}$, $d'<d$: on ${\rm supp}\,\alpha$, $d_{\rm Tb}(dn_0\pm y)\ge d-d'$, so the left side of the integral identity of (a) converges in norm to $\int\alpha(y)B_{\mathfrak a}(dn_0+y,dn_0-y)dy$, while $\alpha*j_\delta\to\alpha$ in $\|\cdot\|_{C^1}$ with supports in a fixed compact subset of $\{|y^0|<d\}$, so that the right side converges to $B^c_d(\alpha\otimes\mathfrak a)$ by Lemma~\ref{C-353-22}\,(c); the bound in (iii) follows from the norm bound of the integral.

For (iv), if $B'$ is such a family, $g(y):=\langle\psi,(B'-B_{\mathfrak a})(dn_0+y,dn_0-y)\psi'\rangle$ is continuous on $\{|y^0|<d\}$ and $\int\alpha g=0$ for all real profiles there; by the fundamental lemma $g=0$, in particular $B'(dn_0,dn_0)=B_{\mathfrak a}(dn_0,dn_0)$ for all $d>0$, and the shift rule gives
\[B'(\xi,\xi')=T(\xi-\varepsilon n_0)B'(\varepsilon n_0,\varepsilon n_0)T(\xi'-\varepsilon n_0)=B_{\mathfrak a}(\xi,\xi'),\qquad\varepsilon=\min\{d_{\rm Tb}(\xi),d_{\rm Tb}(\xi')\}.\]
\end{proof}

\begin{lemma}[Rules for the point plates]\label{C-353-24}
Under the hypotheses of Theorem~\ref{C-353-7}\,(A), the family $B_{\mathfrak a}$ of Proposition~\ref{C-353-23} has property (v) of Theorem~\ref{C-353-7}\,(A).
\end{lemma}

\begin{proof}
For $0<\varepsilon\le\min\{d_{\rm Tb}(\kappa_1),d_{\rm Tb}(\kappa_2)\}$ the shift rule gives $B_{\mathfrak a}(\xi,\kappa_1)=B_{\mathfrak a}(\xi,\varepsilon n_0)T(\kappa_1-\varepsilon n_0)$ and $B_{\mathfrak a'}(\kappa_2,\xi')=T(\kappa_2-\varepsilon n_0)B_{\mathfrak a'}(\varepsilon n_0,\xi')$, so that
\[B_{\mathfrak a}(\xi,\kappa_1)T(\kappa_0)B_{\mathfrak a'}(\kappa_2,\xi')=B_{\mathfrak a}(\xi,\varepsilon n_0)\,T(\varkappa-2\varepsilon n_0)\,B_{\mathfrak a'}(\varepsilon n_0,\xi');\]
for two triples with the same sum we take $\varepsilon$ below the four distances, and $(\frac\varkappa2,0,\frac\varkappa2)$ is such a triple by Lemma~\ref{C-353-17}\,(a). Further $B_{\mathfrak a}(\xi,\kappa)\Omega=B_{\mathfrak a}(\xi,\varepsilon n_0)T(\kappa-\varepsilon n_0)\Omega=B_{\mathfrak a}(\xi,\varepsilon n_0)\Omega$ and $\langle\Omega,B_{\mathfrak a}(\kappa,\xi)\psi\rangle=\langle T((\kappa-\varepsilon n_0)^*)\Omega,B_{\mathfrak a}(\varepsilon n_0,\xi)\psi\rangle=\langle\Omega,B_{\mathfrak a}(\varepsilon n_0,\xi)\psi\rangle$ for $\varepsilon\le\min\{d_{\rm Tb}(\kappa),d_{\rm Tb}(\kappa')\}$ and the comparison with $\kappa'$.
\end{proof}

\begin{proposition}[Point fields]\label{C-353-25}
Let $\sigma$ have (T1), (T2), (T3), (T4) and (T5$_N$). With $j^{(3)}_\delta$ put $A_\delta(E):=B^c_\delta(j^{(3)}_\delta E)$ and $B^\delta_E(\mathfrak z,\mathfrak z'):=T_3(\mathfrak z-\delta e_3)A_\delta(E)T_3(\mathfrak z'-\delta e_3)$ for $E\in V_{\mathbb C}$ and $d_3(\mathfrak z),d_3(\mathfrak z')\ge\delta$. Then $A_\delta(E)^*=A_\delta(\varrho(\theta)\bar E)$, $B^\delta_E$ has the shift rule with $T_3$ and the adjoint $B^\delta_E(\mathfrak z,\mathfrak z')^*=B^\delta_{\varrho(\theta)\bar E}(\mathfrak z'^*,\mathfrak z^*)$, $B^\delta_E(de_3+x,de_3-x)=B^c_d((\tau_xj^{(3)}_\delta)E)$ for $d\ge\delta+|x_3|$, and for $0<\eta\le\min\{1,d_3(\mathfrak z),d_3(\mathfrak z')\}$ and $0<\delta\le c_I\eta/8$,
\[\|B^\delta_E(\mathfrak z,\mathfrak z')\|\le2c_{\rm pt}|E|_2N(\eta)\eta^{-4},\]
without the factor $2$ for $E\in V$, and $\|B^\delta_E(\mathfrak z,\mathfrak z')\|\le2\cdot8^4{\rm mv}_0^3c_I^{-3}|E|_2N(\eta)\eta^{-3}$ if $\sigma$ has (T5$_N$) in the sup norm. The limits $B_E(\mathfrak z,\mathfrak z'):=\lim_{\delta\to0}B^\delta_E(\mathfrak z,\mathfrak z')$ exist in norm and have the properties (i)--(iv) of Theorem~\ref{C-353-7}\,(D).
\end{proposition}

\begin{proof}
The proof of Proposition~\ref{C-353-23} is repeated with $\mathbb R^3$, $x$, $\mathfrak z$, $T_3$, $e_3$, $d_3$ and $\beta E$ in place of $\mathbb R^2$, $y$, $\xi$, $T$, $n_0$, $d_{\rm Tb}$ and $\beta\otimes\mathfrak a$, with Lemmas~\ref{C-353-17}\,(b), \ref{C-353-20}\,(b) and~\ref{C-353-22}\,(b), and without the cut of a label. For real $E$, $r=c_I\eta/8$ and $u(w):=\langle\psi,B^\delta_E(\mathfrak z+w,\mathfrak z'-w)\psi'\rangle$, $|w_i|<2r$: there $d_3(\mathfrak z\pm w)>\eta-6r>\delta$; $\kappa_t:=\mathfrak z-\frac\eta4e_3+it$ has $d_3(\kappa_t)>\frac34\eta-2r>0$; the mean-value formula in three variables gives the profiles $\gamma_t=\alpha^{(3)}_{r,t}*j^{(3)}_\delta$, real, supported in $[-2r,2r]^3$, hence in the ball of radius $2\sqrt3\,r<c_I\eta$ about $0$, with $\|\gamma_t\|_\infty\le{\rm mv}_0^3r^{-6}$, ${\rm Lip}\,\gamma_t\le\sqrt3\,{\rm mv}_0^2{\rm mv}_1r^{-7}$, so $\|\gamma_tE\|_\sharp\le40M\,{\rm mv}_3r^{-7}|E|_2$, and
\[|u(0)|\le(2r)^3\cdot N(\eta)\cdot40M\,{\rm mv}_3r^{-7}|E|_2=320\cdot8^4{\rm mv}_3c_I^{-4}M\,\eta^{-4}N(\eta)|E|_2;\]
in the sup norm $(2r)^3N(\eta){\rm mv}_0^3r^{-6}|E|_2=8^4{\rm mv}_0^3c_I^{-3}\eta^{-3}N(\eta)|E|_2$. The limit is taken at $(\varepsilon e_3,\varepsilon e_3)$ on the polydisc $|w_i|<\varepsilon/6$, where $d_3(\varepsilon e_3\pm w)>\varepsilon/2$, by the Cauchy estimate and the comparison through $B^c_\varepsilon((j^{(3)}_{\delta'}*j^{(3)}_\delta)E)$. The slab formula (iii) is first obtained for $H=\beta E$ with scalar $\beta\in C^1_c(\{|x_3|\le d'\};\mathbb C)$, $d'<d$, as the limit $\delta\to0$ of $\int\beta(x)B^\delta_E(de_3+x,de_3-x)\,dx=B^c_d((\beta*j^{(3)}_\delta)E)$, and then for every $H=\sum_ch_cb_c$ by linearity in $E$. Uniqueness and the rules are proved as in Propositions~\ref{C-353-23}\,(c) and Lemma~\ref{C-353-24}, with the fundamental lemma in three variables.
\end{proof}

\begin{proof}[Proof of Theorem~\ref{C-353-7}]
\emph{Step 1 ((A)).} Lemmas~\ref{C-353-19}, \ref{C-353-20}\,(a), \ref{C-353-21} and~\ref{C-353-22} provide $\mathcal H(\sigma)$, $T(\xi)$ and the slab operators; (A)(i)--(iv) is Proposition~\ref{C-353-23}\,(c), and (A)(v) is Lemma~\ref{C-353-24}.

\emph{Step 2 (kernels and vectors; (KV1$_N$), (KV3)).} By (A)(v), $K_n$ and $\Psi$ do not depend on $\kappa$; holomorphy comes from the calculus of operator functions, since the arguments are affine functions of $\xi$ or $x$ with values in ${\rm Tb}$, and ${\rm Tb}^{(n)}$ is the preimage of ${\rm Tb}^n$ under a linear bijection. At $\kappa=n_0$ every argument has $d_{\rm Tb}\ge\min\{1,\frac12\min_kd_{\rm Tb}(\xi_k)\}$, and (A)(i), applied factor by factor, gives (KV1$_N$). The shift in (KV3) is the shift rule at the first factor. By (A)(ii) adjoints reverse the order, star the arguments and replace $\mathfrak a_k$ by $\mathfrak a_k^\theta$:
\[\langle\Psi(x;\mathbf a),\Psi(z;\mathbf a')\rangle=\langle\Omega,B_{\mathfrak a_n^\theta}(\kappa^*,p_n^*)\cdots B_{\mathfrak a_1^\theta}(q_1^*,x_1^*)\,B_{\mathfrak a'_1}(z_1,q'_1)\cdots B_{\mathfrak a'_{n'}}(p'_{n'},\kappa)\Omega\rangle,\]
with $q_k^*=p_{k+1}^*=\frac12(x_{k+1}-x_k)^*$ and $q'_k=p'_{k+1}=\frac12(z_{k+1}-z_k)$; by (A)(v) the arguments $x_1^*$ and $z_1$ may be replaced by $\frac12(x_1^*+z_1)$ and $\kappa^*$ by $\kappa$, which is the product that defines $K_{n+n'}$ at the stated arguments and labels. The two edge formulas are the cases $n=0$ and $n'=0$.

\emph{Step 3 (the formulas (KV2)).} For supports time-ordered in $\mathbb R^2_+$ put $c_a:=g_a$ ($c_0=0$), $H_a:=\alpha_a\otimes\mathfrak a_a$, $d_a:=\frac12(c_a-c_{a-1})$, $m_a:=\frac12(c_a+c_{a-1})$. The product formula of Lemma~\ref{C-353-21}\,(e) writes $[(H_1,\dots,H_n)_{\mathbb C}]$ as a product of the factors $B^c_{d_a}((\tau_{-m_an_0}\alpha_a)\otimes\mathfrak a_a)$ applied to $\Omega$; by (A)(iii) each factor is $\int\alpha_a(y)B_{\mathfrak a_a}(y-c_{a-1}n_0,c_an_0-y)\,dy$, so
\[
\begin{multlined}
[(H_1,\dots,H_n)_{\mathbb C}]=\int\prod_a\alpha_a(y_a)\,B_{\mathfrak a_1}(y_1,c_1n_0-y_1)B_{\mathfrak a_2}(y_2-c_1n_0,c_2n_0-y_2)\cdots B_{\mathfrak a_n}(y_n \\
-c_{n-1}n_0,c_nn_0-y_n)\Omega\,dY.
\end{multlined}
\]
At each junction the two neighbouring arguments add up to $y_{a+1}-y_a$, and by (A)(v) both may be replaced by $\frac12(y_{a+1}-y_a)$; the last argument may be replaced by $\kappa$; this gives the vector formula. The first formula follows from $\langle\Omega,[F]\rangle=\sigma(F)$ and the edge formula of Step 2; the general time-ordered case follows by one translation in time and the substitution $y_a\mapsto y_a+cn_0$, which keeps the differences.

\emph{Step 4 (conjugation, centring, uniqueness: (KV4), (KV5)).} For real labels and $n\ge2$, $\gamma(Y):={\rm Im}\,K_n(y_2-y_1,\dots;\mathbf a)$ is continuous on $O=\{y_1^0<\dots<y_n^0\}$, and for real profiles with $\prod_k{\rm supp}\,\alpha_k\subset O$ the supports are time-ordered and the left side of (KV2) is real; by the fundamental lemma $\gamma=0$. The function $\xi\mapsto\overline{K_n(\bar\xi;\mathbf a)}$ is holomorphic on ${\rm Tb}^{n-1}$ and equals $K_n$ on the real points, a non-empty open subset of $\mathbb R^{2(n-1)}$ of the connected set ${\rm Tb}^{n-1}$; by the identity theorem it is $K_n$. For complex labels,
\[K_n(\xi;\mathbf a)=\sum_\varepsilon i^{|\varepsilon|}K_n(\xi;\mathfrak a_1^{\varepsilon_1},\dots,\mathfrak a_n^{\varepsilon_n}),\qquad\overline{K_n(\xi;\mathbf a)}=\sum_\varepsilon(-i)^{|\varepsilon|}K_n(\bar\xi;\mathfrak a_1^{\varepsilon_1},\dots)=K_n(\bar\xi;\bar{\mathbf a}).\]
If $\sigma$ is centred, (KV2) at $n=1$ with a real profile of integral $1$ gives $K_1(\mathfrak a)=\sigma((\alpha\otimes\mathfrak a^0))+i\sigma((\alpha\otimes\mathfrak a^1))=0$. Uniqueness: the difference of two continuous densities is annihilated by all products of real profiles with $\prod{\rm supp}\,\alpha_k\subset O$, hence vanishes by the fundamental lemma, and every point of $(\mathbb R^2_+)^{n-1}$ is a tuple of differences of a point of $O$; then the identity theorem on ${\rm Tb}^{n-1}$. For $\Psi$ the same argument applies to $x\mapsto\langle\chi,\Psi(x;\mathbf a)-\Psi'(x)\rangle$ on the real points of ${\rm Tb}^{(n)}$, for every vector $\chi$.

\emph{Step 5 (the power, (C)).} (T5) is (T5$_N$) at $N(\eta)=C_I\eta^{-\mathsf p_I}$; at $\eta=\frac12\min\{1,\rho\}$, $\rho=\min_kd_{\rm Tb}(\xi_k)$,
\[c_{\rm plt}N(\eta)\eta^{-4}=C_{\rm plt}\eta^{-\wp}=C_{\rm plt}2^\wp\min\{1,\rho\}^{-\wp},\]
and (KV1$_N$) gives (KV1) with $C_\flat=2^{\wp+1}C_{\rm plt}$, also at $n=1$ ($\eta=\frac12$). The finer bound and the plate bound are (A)(i) under (T5) at $\kappa=n_0$, factor by factor. The exponent $\wp=\mathsf p_I+4$ arises in Proposition~\ref{C-353-23}\,(b) as $-5$ from the norm $\|\cdot\|_\sharp$ of $\gamma_t\otimes\mathfrak a_k$ (height $r^{-4}$, slope $r^{-1}$), $+2$ from the square of the imaginary parts $t$, and $-1$ from the number of pieces of the label. In the sup norm the slope is not paid, and the bound of Proposition~\ref{C-353-23}\,(b) in the sup norm gives (KV1) with $\wp'$, $\nu_\infty$ and $C'_\flat$.

\emph{Step 6 ((D)).} (D) is Lemma~\ref{C-353-20}\,(b) and Proposition~\ref{C-353-25}.

\emph{Step 7 (plates from point fields; the point bound, (E)).} (i) The right side of $B_{\mathfrak a}(\xi,\xi')=\int B_{\mathfrak a(s)}(\xi+se_2,\xi'-se_2)ds$ is weakly continuous, has the shift rule because $T(\kappa)=T_3(\kappa^1,0,\kappa^0)$ and $d_3(\xi\pm se_2)=d_{\rm Tb}(\xi)$, and satisfies the identity of (A)(iii), by (D)(iii) at $H=\alpha\otimes\mathfrak a$ and Fubini, since $de_3\pm x=(dn_0\pm y)\pm x_2e_2$; by (A)(iv) it is $B_{\mathfrak a}$. The bounds follow from (D)(i), $|\mathfrak a(s)|_2$ and $|H(x)|_2$ under the integral, or $H=\sum_c\operatorname{tr}(b_cH)\,b_c$ with real $b_c$. (ii) The first two bounds on $\mathsf C_{\rm pt}$ are those of Proposition~\ref{C-353-25} at $\eta\le1$. For $\eta$, $H$, $F$, $G$ as in (T5$_N$), ${\rm supp}\,H\subset\{|x_3|<3\eta/16\}$, and by Lemma~\ref{C-353-22}\,(e) and (D)(iii)
\[|\sigma(\theta\bar F\cdot(H)\cdot G)|\le\int\bigl\|B_{H(x)}(\tfrac\eta4e_3+x,\tfrac\eta4e_3-x)\bigr\|dx\,\|[\tau_{-(\eta/4)e_3}F]\|\,\|[\tau_{-(\eta/4)e_3}G]\|\le\mathsf C_{\rm pt}(\tfrac\eta{16})\|H\|_{L^1}\|\cdot\|\,\|\cdot\|,\]
as $d_3(\frac\eta4e_3\pm x)\ge\eta/16$ on ${\rm supp}\,H$; and $\|H\|_{L^1}\le\frac{4\pi}3(c_I\eta)^3\|H\|_\infty$. The consequences are arithmetic: $\mathsf C_{\rm pt}(\eta)\le C\eta^{-\mathsf p}$ gives $|\sigma(\theta\bar F\cdot(H)\cdot G)|\le C16^{\mathsf p}\frac{4\pi}3c_I^3\eta^{3-\mathsf p}\|H\|_\infty\|\cdot\|\,\|\cdot\|$, and (T5) gives $\mathsf C_{\rm pt}(\eta)\le C_{\rm pt}\eta^{-\mathsf p_I-4}$.

\emph{Step 8 (reflections; other spaces of values, (F)).} (a) is Lemma~\ref{C-353-19}\,(c). (b) In the parts of Definitions~\ref{C-353-2}--\ref{C-353-6} that concern the plane, in Lemmas~\ref{C-353-16}--\ref{C-353-19}, \ref{C-353-20}\,(a), \ref{C-353-21}, \ref{C-353-22} (but for the form of (b) with $T_3$) and~\ref{C-353-24}, in Proposition~\ref{C-353-23} and in Steps~1--5 above, of $V$, $|\cdot|_2$ and $\varrho$ only the following enters: $V$ is a real Euclidean space of finite dimension, $V_{\mathbb C}$ its complexification with $|E|^2=|E^0|^2+|E^1|^2$ and the conjugation, translations act on the argument alone, and $\varrho(\theta)$ is a real linear isometry $\vartheta$ with $\vartheta^2=1$; the isometries that act form the group generated by the translations and $\theta$, and
\[\tau_v\mapsto1,\qquad\tau_v\theta\mapsto\vartheta\]
is a homomorphism, which is all that Lemma~\ref{C-353-16}\,(e) asks (Lemma~\ref{C-353-16}\,(c) uses a basis of $V'$). So the proofs of (A), (B) and, under (T5), of (KV1) hold for $V'$ and $\vartheta$ as they stand.
\end{proof}

\begin{proof}[Proof of Theorem~\ref{C-353-8}]
\emph{Step 1 (plate vectors to plate vectors).} Let $A:=B_{\mathfrak b}(\xi,\xi')$; then $A^*=B_{\mathfrak b^\theta}(\xi'^*,\xi^*)$, and by Theorem~\ref{C-353-7}\,(A)(v)
\[A\,\Psi(z;\mathbf a')=\Psi\bigl((\xi,\xi+\xi'+z_1,\dots,\xi+\xi'+z_{n'});\mathfrak b,\mathbf a'\bigr),\qquad A\Omega=\Psi(\xi;\mathfrak b),\]
the point lying in ${\rm Tb}^{(n'+1)}$; so $A\mathcal H_{\rm pl}\subset\mathcal H_{\rm pl}$, and likewise for $A^*$.

\emph{Step 2 (one matrix element).} For plate vectors $\Psi_i=\Psi(x;\mathbf a)$, $\Psi_j=\Psi(z;\mathbf a')$ and $m\ge1$, the Gram formula and the rules show that $\langle\Psi_i,(A^*A)^{2m}\Psi_j\rangle$ is $K_{n+4m+n'}$ at the sums of neighbouring arguments: $(x_{k+1}-x_k)^*$; $x_1^*+\xi'^*$; in turn $\xi^*+\xi$ and $\xi'+\xi'^*$; $\xi'+z_1$; $z_{k+1}-z_k$. By Lemma~\ref{C-353-17}\,(a) these have $d_{\rm Tb}\ge d_{\rm Tb}(\xi')$, $2d_{\rm Tb}(\xi)$, $2d_{\rm Tb}(\xi')$ and $d_{\rm Tb}(\xi')$, so if all inner differences of $\Psi_i$, $\Psi_j$ have $d_{\rm Tb}\ge\rho_0:=\min\{1,2d_{\rm Tb}(\xi),d_{\rm Tb}(\xi')\}$, the hypothesis and Lemma~\ref{C-353-16}\,(f) give
\[|\langle\Psi_i,(A^*A)^{2m}\Psi_j\rangle|\le c_ic_j\beta^{4m},\qquad\beta:=C'\nu(\mathfrak b)\rho_0^{-\varpi},\quad c_i:=\prod_k\bigl(C'\nu(\mathfrak a_k)\rho_0^{-\varpi}\bigr).\]

\emph{Step 3 (Schwarz's inequality, repeated).} For a finite combination $\Phi=\sum\zeta_i\Psi_i$ of such vectors and of $\Omega$, and $B:=A^*A\ge0$, $\|B^{2^k}\Phi\|^2=\langle\Phi,B^{2^{k+1}}\Phi\rangle\le\|\Phi\|\,\|B^{2^{k+1}}\Phi\|$ gives by induction $\|B\Phi\|\le\|\Phi\|^{1-2^{-k}}\|B^{2^k}\Phi\|^{2^{-k}}$; with $m=2^k$ and Step 2, $\|B^m\Phi\|^{1/m}\le(\sum|\zeta_i|c_i)^{1/m}\beta^2\to\beta^2$. Hence $\|B\Phi\|\le\beta^2\|\Phi\|$ and $\|A\Phi\|^2=\langle\Phi,B\Phi\rangle\le\beta^2\|\Phi\|^2$.

\emph{Step 4 (density).} If $\chi\in\mathcal H_{\rm pl}$ is orthogonal to $\Omega$ and to all plate vectors whose inner differences have $d_{\rm Tb}>\rho_0$, then $x\mapsto\langle\chi,\Psi(x;\mathbf a)\rangle$ is holomorphic on the convex open set ${\rm Tb}^{(n)}$ and vanishes on the real points with $y_1^0>0$, $y_{k+1}^0-y_k^0>\rho_0$, a non-empty open subset of $(\mathbb R^2)^n$; by the identity theorem it vanishes, so $\chi=0$. As $A$ is bounded, $\|A\Phi\|\le\beta\|\Phi\|$ extends to $\mathcal H_{\rm pl}$.

\emph{Step 5 (the converse direction).} At $\kappa=n_0$ the kernel is $\langle\Omega,B_{\mathfrak a_1}(n_0,\frac{\xi_1}2)\cdots B_{\mathfrak a_n}(\frac{\xi_{n-1}}2,n_0)\Omega\rangle$, every partial product applied to $\Omega$ lies in $\mathcal H_{\rm pl}$ by Step 1, and every factor is bounded on $\mathcal H_{\rm pl}$ by $C''\nu(\mathfrak a_k)\min\{1,\frac12\min_kd_{\rm Tb}(\xi_k)\}^{-\varpi}\le2^\varpi C''\nu(\mathfrak a_k)\min\{1,\min_kd_{\rm Tb}(\xi_k)\}^{-\varpi}$.
\end{proof}

\begin{proof}[Proof of Corollary~\ref{C-353-9}]
By Lemma~\ref{C-353-22}\,(e) and Theorem~\ref{C-353-7}\,(D)(iii), 
\[
\sigma(\theta\bar F\cdot(H)_{\mathbb C}\cdot F')=\langle[\tau_{-de_3}F],\int B_{H(x)}(de_3+x,de_3-x)\,dx\,[\tau_{-de_3}F']\rangle=\int\mathsf v_d(x;H(x))\,dx,
\]
matrix elements passing under the integral. The function 
\[
w\mapsto\langle[\tau_{-de_3}F],B_E(de_3+w,de_3-w)[\tau_{-de_3}F']\rangle
\]
 is holomorphic where $d_3(de_3\pm w)>0$, that is on $\{|{\rm Re}\,w_3|+|{\rm Im}\,w_1|+|{\rm Im}\,w_2|<d\}$, by composition with an affine map (Lemma~\ref{C-353-17}\,(b)). Linearity and the bound are those of (D)(i), as $d_3(de_3\pm x)\ge d-|x_3|$ at real $x$.
\end{proof}

\begin{proof}[Proof of Theorem~\ref{C-353-10}]
\emph{Step 1 (the tubes).} $R_\upsilon{\rm Tb}$ is an open convex cone that is stable under the conjugation and under ${\rm Re}$ (Lemma~\ref{C-353-17}); so each $\mathcal Z(\mathfrak s,\upsilon)$, the preimage of $(R_\upsilon{\rm Tb})^{n-1}$ under a linear map, is open and convex and contains ${\rm Re}\,Z$, $\bar Z$, $R_{is}Z$ and $Z+(v,\dots,v)$ with $Z$; its real points are the configurations with $n_\upsilon\cdot y_{\mathfrak s(1)}<\dots<n_\upsilon\cdot y_{\mathfrak s(n)}$, since $R_\upsilon{\rm Tb}\cap\mathbb R^2=\{n_\upsilon\cdot w>0\}$. As $R_{2\pi}=1$ and the labels are turned by $\varrho(R_{-\upsilon})$, (a) holds, holomorphy coming from composition with linear maps.

\emph{Step 2 (the $\sigma$-formula on one tube, (b)).} Since $\upsilon\in\frac\pi2\mathbb Z$, $R_{-\upsilon}\in W_3$, and $R_{-\upsilon}(\alpha\otimes\mathfrak a)=(\alpha\circ R_\upsilon)\otimes\varrho(R_{-\upsilon})\mathfrak a$, the rotation keeping $x_2$; this is the only point at which the tensor action meets a member. For profiles as in (b) the profiles $\alpha_{\mathfrak s(k)}\circ R_\upsilon$ have time-ordered supports, the largest time on one support being smaller than the least time on the next; so symmetry, (T1) at $R_{-\upsilon}$ and (KV2) give
\[
\begin{multlined}
\sigma(\alpha_1\otimes\mathfrak a_1,\dots)=\sigma\bigl(R_{-\upsilon}(\alpha_{\mathfrak s(1)}\otimes\mathfrak a_{\mathfrak s(1)},\dots)\bigr)=\int\prod_k(\alpha_{\mathfrak s(k)}\circ R_\upsilon)(y_k)\,K_n(\partial Y;\mathfrak s\varrho(R_{-\upsilon})\mathbf a)\,dY \\
=\int\prod_k\alpha_k(y_k)K^{\mathfrak s}_\upsilon(Y;\mathbf a)\,dY.
\end{multlined}
\]

\emph{Step 3 (one function, (c)).} If $\mathcal Z(\mathfrak s,\upsilon)\cap\mathcal Z(\mathfrak s',\upsilon')$ is not empty, it is convex, hence connected, and it contains real points, which form a non-empty open subset of $(\mathbb R^2)^n$. There both functions are continuous densities of $\sigma$ by Step 2, so they agree by the fundamental lemma, and on the whole intersection by the identity theorem.

\emph{Step 4 (symmetries and bounds, (d), (e)).} Invariance under the permutations of points and labels follows in the same way from the symmetry of $\sigma$, and translation invariance from the dependence on $\partial Z$. For $Z\in\mathcal Z(\mathfrak s,\upsilon)$, $R_{\pi/2}Z\in\mathcal Z(\mathfrak s,\upsilon+\frac\pi2)$ and
\[K^{\mathfrak s}_{\upsilon+\pi/2}(R_{\pi/2}Z;\mathbf a)=K_n\bigl(R_{-\upsilon}\partial(\mathfrak sZ);\mathfrak s\varrho(R_{-\upsilon})\varrho(R_{-\pi/2})\mathbf a\bigr)=K^{\mathfrak s}_\upsilon(Z;\varrho(R_{-\pi/2})\mathbf a).\]
The conjugation rule follows from (KV4), as $R_{-\upsilon}$ and $\varrho(R_{-\upsilon})$ are real. (e) is (KV1) for the labels $\varrho(R_{-\upsilon})\mathfrak a_{\mathfrak s(k)}$, with $\nu(\varrho(R)\mathfrak a)=\nu(\mathfrak a)$.

\emph{Step 5 (the Gram formula along boosts, (f)).} By (KV3) at $R_{is}x$, $R_{is}z$, with $(R_{is}w)^*=R_{is}w^*$ and $w^*=-\theta w$ for real $w$ (Lemma~\ref{C-353-17}\,(c)), the arguments of $K_{n_1+n_2}$ are $R_{is}$ applied to the consecutive differences of $X$: $(R_{is}(x_{k+1}-x_k))^*=R_{is}(\theta x_k-\theta x_{k+1})$ and $(R_{is}x_1)^*+R_{is}z_1=R_{is}(z_1-\theta x_1)$. These lie in ${\rm Tb}$, so $R_{is}X\in\mathcal Z({\rm id},0)$ and the value is $\mathcal K_{n_1+n_2}(R_{is}X;\dots)$.

\emph{Step 6 (helicity labels, (g)).} $R_t\epsilon_\pm=e^{\mp it}\epsilon_\pm$ and $R_te_2=e_2$, and each $E_s$ is bilinear in two of $\epsilon_+,\epsilon_-,e_2$, so $\varrho(R_t)E_s=e^{-i\,{\rm hel}(s)t}E_s$; $\bar\epsilon_\pm=\epsilon_\mp$ and $\theta\epsilon_\pm=-\epsilon_\mp$, so $\varrho(\theta)\bar E_s=(-1)^{{\rm hel}(s)}E_s$, and for real $\varphi$, $(\varphi E_s)^\theta=(-1)^{{\rm hel}(s)}\varphi E_s$. The three formulas follow by linearity in each label from $\varrho(R_{-\upsilon})E_s=e^{i\,{\rm hel}(s)\upsilon}E_s$, from (d) and from (KV3).

\emph{Step 7 (harmonics, (h)).} Let the points of $X$ lie on a line with direction $n_\beta$, with coordinates $r_1,\dots,r_n$ along it. For real $t$ choose $\upsilon\in\frac\pi2\mathbb Z$ with $|\beta+t-\upsilon|\le\frac\pi4$ modulo $2\pi$, so that $\cos(\beta+t-\upsilon)\ge2^{-1/2}$, and $\mathfrak s$ ordering the points by increasing $r_a$; then $n_\upsilon\cdot(R_ty_{\mathfrak s(k+1)}-R_ty_{\mathfrak s(k)})\ge c_X:=\min\{1,2^{-1/2}\min_{a\ne b}|r_a-r_b|\}$, so $R_tX\in\mathcal Z(\mathfrak s,\upsilon)$, and $R_{t+i\tau}X=R_{i\tau}R_tX$ lies there too. So $\zeta\mapsto\mathcal K_n(R_\zeta X;\mathbf a)$ is entire, and by (g), $e^{-iS(\zeta+\pi/2)}\mathcal K_n(R_{\pi/2}R_\zeta X;\mathbf a)=e^{-iS\zeta}\mathcal K_n(R_\zeta X;\mathbf a)$: the function $f_X(\zeta):=e^{-iS\zeta}\mathcal K_n(R_\zeta X;\mathbf a)$ is entire and $\frac\pi2$-periodic, hence a locally uniformly convergent series $\sum_mc_m(X)e^{4im\zeta}$, with $f_{R_tX}(\zeta)=e^{iSt}f_X(\zeta+t)$, which gives the rule for $c_m(R_tX)$. $k(E_s)=-{\rm hel}(s)$ is the identity $\varrho(R_t)E_s=e^{-i\,{\rm hel}(s)t}E_s$, put in the form $e^{ik(E_s)t}E_s$. Under (T5), by (e) and Lemma~\ref{C-353-17}\,(c), with $\zeta=t+i\tau$,
\[|f_X(t+i\tau)|\le e^{|S||\tau|}\prod_kC_\flat\nu(\mathfrak a_k)\cdot(e^{-|\tau|}c_X)^{-\wp n}=Ce^{(\wp n+|S|)|\tau|},\]
Put $A=\wp n+|S|$. The contour shift used in the proof of
Proposition~\YMMref{96.3}{728} gives, for every real $s$,
\[
 c_m=\frac2\pi\int_0^{\pi/2}
       f_X(t+is)e^{-4im(t+is)}\,dt,
 \qquad |c_m|\le C e^{A|s|+4ms}.
\]
The vertical sides cancel by periodicity. If $4|m|>A$, take
$s\to-\infty$ for $m>0$ and $s\to+\infty$ for $m<0$;
then $c_m=0$. Finally $|{\rm hel}|\le2$ gives $|S|\le2n$.
\end{proof}

\begin{proof}[Proof of Theorem~\ref{C-353-11}]
\emph{Step 1 (point vectors and point kernels).} By (D)(iv) of Theorem~\ref{C-353-7}, $W_n$ and $\Psi_3$ do not depend on $\kappa$, and by (D)(ii) and (D)(iv), exactly as in Step 2 of the proof of Theorem~\ref{C-353-7},
\[\langle\Psi_3(\mathfrak x;\mathbf E),\Psi_3(\mathfrak x';\mathbf E')\rangle=W_{n+n'}\bigl((\mathfrak x_n-\mathfrak x_{n-1})^*,\dots,\mathfrak x_1^*+\mathfrak x'_1,\mathfrak x'_2-\mathfrak x'_1,\dots;\varrho(\theta)\bar E_n,\dots,\varrho(\theta)\bar E_1,E'_1,\dots\bigr),\]
and $A:=B_E(\mathfrak z,\mathfrak z')$ maps $\Psi_3(\mathfrak x';\mathbf E')$ to $\Psi_3((\mathfrak z,\mathfrak z+\mathfrak z'+\mathfrak x'_1,\dots);E,\mathbf E')$ and $\Omega$ to $\Psi_3(\mathfrak z;E)$; so $A\mathcal H_{\rm pt}\subset\mathcal H_{\rm pt}$.

\emph{Step 2 (one matrix element).} As in Step 2 of the proof of Theorem~\ref{C-353-8}, $\langle\Psi_i,(A^*A)^{2m}\Psi_j\rangle$ is $W$ at arguments with $d_3\ge\rho_0:=\min\{1,2d_3(\mathfrak z),d_3(\mathfrak z')\}$, where $d_3(\mathfrak z+\mathfrak z^*)=2d_3(\mathfrak z)$, and $|\varrho(\theta)\bar E|_2=|E|_2$; so it is bounded by $c_ic_j(C'|E|_2\rho_0^{-\varpi})^{4m}$.

\emph{Step 3 (Schwarz's inequality, repeated, and density).} Steps 3 and 4 of the proof of Theorem~\ref{C-353-8} carry over, with the identity theorem on ${\rm Tb}_3^{(n)}$, whose real points with positive time differences form a non-empty open set; this gives the bound of $B_E(\mathfrak z,\mathfrak z')$ on $\mathcal H_{\rm pt}$.

\emph{Step 4 (the arguments of the insertion).} For $H$ as stated, ${\rm supp}\,H\subset\{|x_3|\le3\eta/16\}$; for $x$ there, $\mathfrak z=\frac\eta4e_3+x$ and $\mathfrak z'=\frac\eta4e_3-x$ have $\min\{1,2d_3(\mathfrak z),d_3(\mathfrak z')\}\ge\eta/16$.

\emph{Step 5 (the insertion bound on $\mathcal H_{\rm pt}$).} By Lemma~\ref{C-353-22}\,(e), (D)(iii) and Step 3, with $\psi:=[\tau_{-(\eta/4)e_3}G]\in\mathcal H_{\rm pt}$,
\[
\begin{multlined}
|\sigma(\theta\bar F\cdot(H)\cdot G)|\le\|[\tau_{-(\eta/4)e_3}F]\|\int\|B_{H(x)}(\tfrac\eta4e_3+x,\tfrac\eta4e_3-x)\psi\|\,dx \\
\le16^\varpi C'\eta^{-\varpi}\|H\|_{L^1}\|[\tau_{-(\eta/4)e_3}F]\|\,\|\psi\|,
\end{multlined}
\]
and $\|H\|_{L^1}\le\frac{4\pi}3(c_I\eta)^3\|H\|_\infty$.
\end{proof}

\begin{lemma}[Bounds relative to the energy in a Fock space]\label{C-353-26}
Let $(\mathsf Y,\mu)$ be a $\sigma$-finite measure space, $\omega\colon\mathsf Y\to(0,\infty)$ measurable, $\mathscr F$ the symmetric Fock space over $L^2(\mathsf Y,\mu)$ with the annihilation and creation forms $a(p)$, $a^*(p)$, and $\mathbf H_0:=d\Gamma(\omega)$. Put $\|g\|_\omega:=(\int|g|^2\omega^{-1}d\mu)^{1/2}$ and $\|K\|_\omega:=(\iint|K(p,p')|^2\omega(p)^{-1}\omega(p')^{-1}d\mu\,d\mu)^{1/2}$. Then, for vectors $\chi,\psi$ in the relevant domains,
\[\|a(g)\psi\|\le\|g\|_\omega\|\mathbf H_0^{1/2}\psi\|,\qquad\Bigl\|\iint K(p,p')a(p)a(p')\,d\mu\,d\mu\;\psi\Bigr\|\le\|K\|_\omega\|\mathbf H_0\psi\|,\]
\[\Bigl|\iint K(p,p')\,\bigl\langle a(p)\chi,A\,a(p')\psi\bigr\rangle\,d\mu\,d\mu\Bigr|\le\|K\|_\omega\|A\|\,\|\mathbf H_0^{1/2}\chi\|\,\|\mathbf H_0^{1/2}\psi\|\]
for every bounded operator $A$; at $A=1$ the left side is $|\langle\chi,\iint K(p,p')a^*(p)a(p')\,d\mu\,d\mu\;\psi\rangle|$. Further, $a(p)f(\mathbf H_0)=f(\mathbf H_0+\omega(p))a(p)$.
\end{lemma}

\begin{proof}
In the $m$-particle space, $(a(p)\psi)(k)=\sqrt m\,\psi(p,k)$; by Schwarz's inequality 
\[
|\int\bar g(p)\psi(p,k)d\mu(p)|^2\le\|g\|_\omega^2\int\omega(p)|\psi(p,k)|^2d\mu(p),
\]
and $m\int\omega(p)|\psi(p,k)|^2d\mu(p)\,dk=\langle\psi,\mathbf H_0\psi\rangle$ by symmetry. In the same way, with two variables, 
\[
m(m-1)\int\omega(p)\omega(p')|\psi(p,p',k)|^2\le\langle\psi,\mathbf H_0^2\psi\rangle,
\]
because $\sum_{i\ne j}\omega_i\omega_j\le(\sum_i\omega_i)^2$; the mixed term is bounded by $\|A\|\iint|K|\,\|a(p)\chi\|\,\|a(p')\psi\|$ and Schwarz's inequality, with $\int\omega(p)\|a(p)\chi\|^2d\mu(p)=\langle\chi,\mathbf H_0\chi\rangle$. The commutation rule holds sector by sector.
\end{proof}

\begin{proof}[Proof of Theorem~\ref{C-353-12}]
\emph{Step 1 (the trace construction, (a)).} $\operatorname{tr}\varrho(r)E=\operatorname{tr}E$ for $r\in{\rm O}(3)$, so $\operatorname{tr}(rH)$ is the scalar action of $r$ on $\operatorname{tr}H$, and (T1) follows from (R1); $|\operatorname{tr}E|\le\sqrt3|E|_2$ gives $\|\operatorname{tr}H\|_\sharp\le\sqrt3\|H\|_\sharp$ and (T2) with $3^{n/2}C^{\sigma_0}_n$; $\operatorname{tr}(\theta H)=(\operatorname{tr}H)\circ\theta$ gives (T3) from (R3). Moreover
\[\langle F,G\rangle_\sigma=\langle\operatorname{tr}F,\operatorname{tr}G\rangle_{\sigma_0},\]
$\operatorname{tr}$ applied member by member, so $[F]_\sigma\mapsto[\operatorname{tr}F]_{\sigma_0}$ is an isometry that intertwines the translations; hence the spectral measure of each $[F]_\sigma$ equals that of $[\operatorname{tr}F]_{\sigma_0}$, which is carried by $\{E\ge|q_1|\}$, and (T4) follows, the classes being dense. The insertion bound of $\sigma_0$ at $\operatorname{tr}H$ gives (T5), or (T5$_N$), with the factor $\sqrt3$.

\emph{Step 2 (the free fields, (b): definition).} By Wick's theorem, $\mathbb E[T_\varepsilon(H_1)\cdots T_\varepsilon(H_n)]$ is a finite sum over the pairings of the $2n$ factors $\partial_a\phi^c_\varepsilon$ in which no pair lies at one vertex; each pairing contributes an integral of a product of $n$ covariances $\delta_{cc'}G^\varepsilon_{ab}(x_k-x_l)$, $G^\varepsilon_{ab}=-\partial_a\partial_b(j^{(3)}_\varepsilon*j^{(3)}_\varepsilon*G)$, with $k\ne l$. On the supports of a tuple of $\mathfrak S^V(\mathbb R^3)$, $|x_k-x_l|\ge\mathsf d>0$, and $G^\varepsilon_{ab}$ converges uniformly on $\{|x|\ge\mathsf d\}$ to
\[G_{ab}(x)=-\partial_a\partial_bG(x)=\frac{\delta_{ab}|x|^2-3x_ax_b}{4\pi|x|^5}.\]
So the limit exists, is multilinear, vanishes for $n=1$, and is bounded by $c_n(c\mathsf d^{-3})^n\prod_k\|H_k\|_{L^1}\le c_n(c\mathsf d^{-3})^n(\frac{4\pi}3R^3)^n\prod_k\|H_k\|_\infty$ at tuples of type $(\mathsf d;R)$: (T2).

\emph{Step 3 ((T1) and (T3) for (b)).} $G_{ab}(r^\circ x)=(r^\circ G(x)r^{\circ\mathsf T})_{ab}$ for $r^\circ\in{\rm O}(3)$, so the limit is invariant under all isometries in the tensor action. For reflection positivity let $F=\sum\zeta_iX_i$ with $X_i\in\mathfrak S^V(E_+)$ and $F_\varepsilon:=\sum_i\zeta_i\prod_{H\in X_i}T_\varepsilon(H)$; for small $\varepsilon$, $F_\varepsilon$ is a square-integrable function of the fields in $E_+$, and the reflection of the fields, $\phi\mapsto\phi\circ\theta$, carries $T_\varepsilon(H)$ to $T_\varepsilon(\theta H)$, since $j^{(3)}_\varepsilon$ is even and $\partial_a(\phi\circ\theta)=\theta_{aa}(\partial_a\phi)\circ\theta$. Reflection positivity of the free fields gives $\mathbb E[\overline{\Theta F_\varepsilon}\,F_\varepsilon]\ge0$, and as $\varepsilon\to0$ this tends to $\sigma(\theta\bar F\cdot F)$, because only the moments at the disjoint tuples $\theta X_i\cdot X_k$ enter.

\emph{Step 4 ((T4) for (b)).} The reconstructed space of the free fields is the Fock space over $L^2(\mathbb R^2,\mu)\otimes\mathbb C^3$, $d\mu(p)=(2\pi)^{-2}(2|p|)^{-1}d^2p$, with $\omega(p):=|p|$, $\mathbf H_0=d\Gamma(\omega)$ and $\mathbf P_0=d\Gamma(p)$, in the following form. For real $f\in C_c(E_+)$ put $jf(p):=\int f(y)\,e^{ip\cdot y_\perp-\omega(p)y_3}\,dy$, $y_\perp=(y_1,y_2)$. As $G(z)=\int e^{ip\cdot z_\perp-\omega(p)|z_3|}\,d\mu(p)$ for $z_3\ne0$,
\[\langle jf,jg\rangle=\iint f(y)\,G(\theta y-y')\,g(y')\,dy\,dy'=\mathbb E[\phi^c(\theta f)\,\phi^c(g)],\]
so, by Wick's theorem, the map $\mathsf W$ defined on the Wick monomials of fields over $E_+$ by
\[\mathsf W\bigl({:}\phi^{c_1}(f_1)\cdots\phi^{c_m}(f_m){:}\bigr):=a^{c_1*}(jf_1)\cdots a^{c_m*}(jf_m)\Omega\]
satisfies $\langle\mathsf WA,\mathsf WB\rangle=\mathbb E[\overline{\Theta A}\,B]$ and $\|\mathsf WA\|^2\le\mathbb E[|A|^2]$. It therefore extends to all square-integrable functions of the fields in $E_+$, whatever the order in time of the arguments of the fields, and $\mathsf W(\tau_{te_3}A)=e^{-t\mathbf H_0}\mathsf WA$, $\mathsf W(\tau_vA)=e^{iv\cdot\mathbf P_0}\mathsf WA$ for $t\ge0$ and $v\in e_3^\perp$, the translations acting on the arguments of the fields, because $j(\tau_{te_3}f)=e^{-t\omega}jf$ and $j(\tau_vf)=e^{ip\cdot v}jf$. The vectors $\mathsf V_\varepsilon(F):=\mathsf W(F_\varepsilon)$ of the $F_\varepsilon$ of Step 3 converge as $\varepsilon\to0$, since $\|\mathsf V_\varepsilon(F)-\mathsf V_{\varepsilon'}(F)\|^2$ is a combination of moments at the disjoint tuples $\theta X_i\cdot X_k$ whose limits do not depend on the regularisation; so $\mathsf V[F]:=\lim\mathsf V_\varepsilon(F)$ is an isometry of $\mathcal H(\sigma)$ into the Fock space with $\mathsf V\mathcal T_t=e^{-t\mathbf H_0}\mathsf V$ and $\mathsf V\mathcal U(v)=e^{iv\cdot\mathbf P_0}\mathsf V$. The image is closed and reduces $\mathbf H_0$ and $\mathbf P_0$, so $\mathcal P$ is the restriction of their joint spectral measure, which is carried by $\{E\ge|q|\}$, because on $m$ particles $E=\sum_j|p_j|\ge|\sum_jp_j|=|q|$.

\emph{Step 5 ((T5) with $\mathsf p_I=0$ for (b)).} Put $u(p):=(-ip_1,-ip_2,\omega(p))$, so that $|u(p)|^2=2\omega(p)^2$, and, on the Fock space of Step~4 and for $E\in V_{\mathbb C}$,
\[B^c_a:=\int u_a(p)\bigl(a^c(p)-a^{c*}(p)\bigr)\,d\mu(p),\qquad{:}Q_E{:}:=\sum_c\sum_{a,b}E_{ab}\,{:}B^c_aB^c_b{:},\]
the colons denoting the normal order. With the time-zero fields $\phi^c_0(0):=\int\bigl(a^c(p)+a^{c*}(p)\bigr)\,d\mu(p)$ and $\phi^c_0(\vec x):=e^{i\vec x\cdot\mathbf P_0}\phi^c_0(0)e^{-i\vec x\cdot\mathbf P_0}$, the forms $B^c_1$, $B^c_2$ are the spatial derivatives $\partial_1\phi^c_0$, $\partial_2\phi^c_0$ at $\vec x=0$, and
\[B^c_3=-[\mathbf H_0,\phi^c_0(0)]=[\phi^c_0(0),\mathbf H_0];\]
the sign is that of $\partial_{x_3}\bigl[e^{-(d+x_3)\mathbf H_0}\phi^c_0(0)e^{-(d-x_3)\mathbf H_0}\bigr]=e^{-(d+x_3)\mathbf H_0}\bigl(-[\mathbf H_0,\phi^c_0(0)]\bigr)e^{-(d-x_3)\mathbf H_0}$, $|x_3|<d$: in this sandwich, the form in which the fields enter below, the lower heights stand on the left. In normal order, ${:}Q_E{:}$ is the sum over $c$ of the forms $\iint K\,a^ca^c$, $\iint K\,a^{c*}a^{c*}$ and $-2\iint K\,a^{c*}a^c$ with the kernel $K(p,p'):=u(p)^{\mathsf T}Eu(p')=K(p',p)$, and $|K(p,p')|\le|E|_2\,|u(p)|\,|u(p')|=2|E|_2\,\omega(p)\omega(p')$. Moving half of $e^{-s\mathbf H_0}$ and of $e^{-s'\mathbf H_0}$ through the annihilation and creation operators by $a(p)f(\mathbf H_0)=f(\mathbf H_0+\omega(p))a(p)$, and using Lemma~\ref{C-353-26}, $\int\omega e^{-s\omega}d\mu=(4\pi s^2)^{-1}$ and $\|\mathbf H_0^\alpha e^{-s\mathbf H_0}\|\le(\alpha/es)^\alpha$, one bounds the three terms of $e^{-s\mathbf H_0}{:}Q_E{:}e^{-s'\mathbf H_0}$ by absolute multiples of $|E|_2s'^{-3}$, $|E|_2s^{-3}$ and $|E|_2(ss')^{-3/2}$; in the last term the factor $e^{-(s+s')\mathbf H_0/2}$ stays between $a^*$ and $a$ and is the operator $A$ of Lemma~\ref{C-353-26}. So
\[\|e^{-s\mathbf H_0}\,{:}Q_E{:}\,e^{-s'\mathbf H_0}\|\le c_0|E|_2\min\{s,s'\}^{-3}\]
with an absolute number $c_0$, the power $3$ being the scaling of a field of dimension three. The same bound holds, with $s-\varepsilon$ and $s'-\varepsilon$ in place of $s$ and $s'$, if the kernel $K$ is replaced by $K(p,p')m(p)m(p')$ with $|m(p)|\le e^{\varepsilon\omega(p)/2}$ and $0<\varepsilon<\min\{s,s'\}$.

Let $0\le d'<d$, $F,G\in\mathbb C\mathfrak S^V(\{x_3>d\})$ and $H\in C^1_c(\{|x_3|<d\};V)$ with ${\rm supp}\,H\subset\{|x_3|\le d'\}$. We show
\[\sigma(\theta\bar F\cdot(H)\cdot G)=\int\bigl\langle\mathsf V[\tau_{-de_3}F],e^{-(d+x_3)\mathbf H_0}e^{i\vec x\cdot\mathbf P_0}\,{:}Q_{H(x)}{:}\,e^{-i\vec x\cdot\mathbf P_0}e^{-(d-x_3)\mathbf H_0}\mathsf V[\tau_{-de_3}G]\bigr\rangle\,dx,\]
in which only the bounded operators $e^{-s\mathbf H_0}{:}Q_E{:}e^{-s'\mathbf H_0}$ with $s,s'\ge d-d'$ enter. Let $\varepsilon>0$ be smaller than $\frac12(d-d')$ and than the distance of the supports of the members of $F$ and $G$ from $\{x_3=d\}$, and put $m_\varepsilon(p):=\int j^{(3)}_\varepsilon(z)e^{ip\cdot z_\perp-\omega(p)z_3}dz$. As $j^{(3)}_\varepsilon$ is even in each coordinate and supported in $\{|z|\le\varepsilon/2\}$, $m_\varepsilon$ is real, $|m_\varepsilon(p)|\le e^{\varepsilon\omega(p)/2}$, and $m_\varepsilon(p)\to1$ as $\varepsilon\to0$. For real $f,g\in C_c(\{x_3>d\})$ and $|x_3|\le d'$, the formula for $G$ in Step~4 gives the two contractions of a leg of the insertion:
\[
\begin{multlined}
\mathbb E[\partial_a\phi^c_\varepsilon(x)\,\phi^{c'}(g)]=\delta_{cc'}\int u_am_\varepsilon\,e^{-ip\cdot\vec x+\omega x_3}\,jg\,d\mu,\qquad\mathbb E[\partial_a\phi^c_\varepsilon(x)\,\phi^{c'}(\theta f)]= \\
-\delta_{cc'}\int u_am_\varepsilon\,e^{ip\cdot\vec x-\omega x_3}\,\overline{jf}\,d\mu.
\end{multlined}
\]
With $B^{c,\varepsilon}_a:=\int u_am_\varepsilon(a^c-a^{c*})\,d\mu$, the first is the matrix element of the annihilation part of $e^{i\vec x\cdot\mathbf P_0}B^{c,\varepsilon}_ae^{-i\vec x\cdot\mathbf P_0}$ between $\Omega$ and $e^{-(d-x_3)\mathbf H_0}\mathsf W(\phi^{c'}(\tau_{-de_3}g))$, and the second is that of its creation part between $e^{-(d+x_3)\mathbf H_0}\mathsf W(\phi^{c'}(\tau_{-de_3}f))$ and $\Omega$: a leg contracted upwards acts as an annihilation operator, a leg contracted downwards as a creation operator. By Wick's product formula, $F_\varepsilon$ and $G_\varepsilon$ are finite sums of Wick monomials in fields smeared, with smooth kernels, over $\{x_3>d\}$. For two such monomials $A$ and $B$, Wick's theorem writes $\mathbb E[\overline{\Theta A}\,{:}\partial_a\phi^c_\varepsilon\partial_b\phi^c_\varepsilon{:}(x)\,B]$ as the sum over the pairings in which each of the two legs is paired with a factor of $\overline{\Theta A}$ or of $B$ and the remaining factors of $\overline{\Theta A}$ are paired with those of $B$, the colons excluding the pair of the two legs.

The normal order gives the same sum for the matrix element of 
\[
e^{-(d+x_3)\mathbf H_0}e^{i\vec x\cdot\mathbf P_0}{:}B^{c,\varepsilon}_aB^{c,\varepsilon}_b{:}e^{-i\vec x\cdot\mathbf P_0}e^{-(d-x_3)\mathbf H_0}
\]
 between $\mathsf W(\tau_{-de_3}A)$ and $\mathsf W(\tau_{-de_3}B)$, as the matrix elements of normal-ordered polynomials between coherent vectors, expanded in their parameters, show: once the operators on either side of the normal product are moved onto the two vectors, a pair of a factor $\phi^c(\theta f)$ of $\overline{\Theta A}$ and a factor $\phi^c(g)$ of $B$ gives $\langle jf,jg\rangle=\mathbb E[\phi^c(\theta f)\phi^c(g)]$, whatever $x$. Summing and integrating against $H(x)$, with $(\tau_{-de_3}F)_\varepsilon=\tau_{-de_3}F_\varepsilon$,
\[\mathbb E[\overline{\Theta F_\varepsilon}\,T_\varepsilon(H)\,G_\varepsilon]=\int\bigl\langle\mathsf V_\varepsilon(\tau_{-de_3}F),e^{-(d+x_3)\mathbf H_0}e^{i\vec x\cdot\mathbf P_0}\,{:}Q^\varepsilon_{H(x)}{:}\,e^{-i\vec x\cdot\mathbf P_0}e^{-(d-x_3)\mathbf H_0}\mathsf V_\varepsilon(\tau_{-de_3}G)\bigr\rangle\,dx,\]
where ${:}Q^\varepsilon_E{:}$ is formed with the $B^{c,\varepsilon}_a$, that is, with the kernel $K(p,p')m_\varepsilon(p)m_\varepsilon(p')$; by the bound above at $s-\varepsilon$, $s'-\varepsilon$, only bounded operators enter. Now let $\varepsilon\to0$. The left side tends to $\sigma(\theta\bar F\cdot(H)\cdot G)$, by Step~2 at the tuples $\theta X_i\cdot(H)\cdot Y_k$, where $F=\sum\zeta_iX_i$ and $G=\sum\zeta'_kY_k$, and by $\Theta T_\varepsilon(H')=T_\varepsilon(\theta H')$ (Step~3). On the right, $\mathsf V_\varepsilon(\tau_{-de_3}F)\to\mathsf V[\tau_{-de_3}F]$ and $\mathsf V_\varepsilon(\tau_{-de_3}G)\to\mathsf V[\tau_{-de_3}G]$ in norm by Step~4, and $e^{-s\mathbf H_0}({:}Q^\varepsilon_E{:}-{:}Q_E{:})e^{-s'\mathbf H_0}\to0$ in norm, uniformly for $s,s'\ge d-d'$ and $|E|_2\le1$, because the kernels $K(p,p')(m_\varepsilon(p)m_\varepsilon(p')-1)$, weighted as in the bound above, tend to $0$ in $\|\cdot\|_\omega$ by dominated convergence. So each factor of the integrand converges in norm, uniformly for $x\in{\rm supp}\,H$, and the identity follows.

At $d=\eta/4$ and for $H$ as in (T5$_N$) with $c_I=\frac1{16}$, $\min\{d+x_3,d-x_3\}\ge\eta/16$, so $|\sigma(\theta\bar F\cdot(H)\cdot G)|\le c_016^3\eta^{-3}\frac{4\pi}3(\eta/16)^3\|H\|_\infty\|\cdot\|\,\|\cdot\|=\frac{4\pi}3c_0\|H\|_\infty\|\cdot\|\,\|\cdot\|$. For the duals of the field strengths of free photons the covariance off the diagonal is $+\partial_a\partial_bG$, the negative of $G_{ab}$; every pairing has $n$ covariances, which gives the factor $(-1)^n$, and $[X]\mapsto(-1)^{n(X)}[X]$ carries all the assertions over.

\emph{Step 6 ((c): the class).} $\sigma=\sigma_0\circ\operatorname{tr}$, where $\sigma_0(f_1,\dots,f_n):=\sum_\pi\prod_{\{k,l\}\in\pi}\iint f_k(x)S(x-y)f_l(y)\,dx\,dy$, the sum over the partitions $\pi$ of $\{1,\dots,n\}$ into pairs, for real scalar $C^1_c$ functions $f_1,\dots,f_n$ with pairwise disjoint supports, and $\sigma_0(\emptyset)=1$; the integrals are finite because $S$ is smooth on $\mathbb R^3\setminus\{0\}$. As $S$ is radial and bounded on $\{|x|\ge\mathsf d\}$ for every $\mathsf d>0$, the properties (R1) and (R2) of Definition~\YMMref{95.1}{723} can be read off this formula. For $\Lambda<\infty$ let $S_\Lambda(x):=\int_0^\Lambda e^{\sqrt m}e^{-m|x|}(4\pi|x|)^{-1}dm\le\Lambda e^{\sqrt\Lambda}(4\pi|x|)^{-1}$, and let $\sigma_{0,\Lambda}$ be formed as $\sigma_0$ with $S_\Lambda$ in place of $S$. $S_\Lambda$ is locally integrable and a superposition with positive weights of the covariances of free fields of mass $m<\Lambda$; so there is a Gaussian field $\chi_\Lambda$ with covariance $S_\Lambda$, and $\sigma_{0,\Lambda}$ is the functional of its moments.

Steps~3 and~4, with one scalar field in place of the derivatives and without regularisation, apply to it, with the one-particle space $L^2((0,\Lambda)\times\mathbb R^2,\mu_\Lambda)$, $d\mu_\Lambda(m,p)=e^{\sqrt m}dm\,(2\pi)^{-2}d^2p/2\omega_m(p)$, $\omega_m(p)=(|p|^2+m^2)^{1/2}\ge|p|$, $jf(m,p):=\int f(y)e^{ip\cdot y_\perp-\omega_m(p)y_3}dy$, $\mathbf H_0=d\Gamma(\omega_m(p))$ and $\mathbf P_0=d\Gamma(p)$, as $S_\Lambda(z)=\int e^{ip\cdot z_\perp-\omega_m(p)|z_3|}d\mu_\Lambda(m,p)$ for $z_3\ne0$: $\sigma_{0,\Lambda}$ is reflection positive, its reconstructed space is mapped isometrically into the Fock space by a map $\mathsf V_\Lambda$ that intertwines the translations, and its joint spectral measure $\mathcal P_\Lambda$ is carried by $\{E\ge|q|\}$, since on $k$ particles $E=\sum_j\omega_{m_j}(p_j)\ge|\sum_jp_j|$.

On $\{|x|\ge\mathsf d\}$, $0\le S(x)-S_\Lambda(x)\le\int_\Lambda^\infty e^{\sqrt m-m\mathsf d}dm\,(4\pi\mathsf d)^{-1}$, as $e^{-m|x|}/|x|$ decreases in $|x|$; so $S_\Lambda\to S$ uniformly there, and the moments of $\sigma_{0,\Lambda}$ at tuples with pairwise disjoint supports converge to those of $\sigma_0$. Hence (R3), an inequality between such moments, passes to $\sigma_0$.

(R4) passes as follows. Let $F$ be a complex combination of scalar tuples over $E_+$, $t>0$, and $u_\Lambda(w):=\sigma_{0,\Lambda}(\theta\bar F\cdot\tau_wF)=\int e^{-w_3E+iw_\perp\cdot q}\,d\langle[F],\mathcal P_\Lambda[F]\rangle$ for $w_3\ge0$. The function
\[f_t(E,q):=(1-e^{-tE})^2-e^{-tE}|1-e^{itq_1}|^2=4e^{-tE}\bigl(\sinh^2\tfrac{tE}2-\sin^2\tfrac{tq_1}2\bigr)\]
is $\ge0$ on $\{E\ge|q_1|\}$, which carries $\mathcal P_\Lambda$, as $\sinh\frac{tE}2\ge\frac{tE}2\ge\frac{t|q_1|}2\ge|\sin\frac{tq_1}2|$ there; so
\[u_\Lambda(0)-4u_\Lambda(te_3)+u_\Lambda(2te_3)+u_\Lambda(te_3+te_1)+u_\Lambda(te_3-te_1)=\int f_t\,d\langle[F],\mathcal P_\Lambda[F]\rangle\ge0.\]
The left side is a combination of moments at disjoint supports; so the inequality passes to $\Lambda=\infty$, that is, $\langle\psi,f_t(\mathcal P)\psi\rangle\ge0$ for the classes $\psi=[F]$ of $\sigma_0$, and, by density, for all $\psi\in\mathcal H(\sigma_0)$, $\mathcal P$ being the joint spectral measure of $\sigma_0$. At $\psi=\mathcal P(\Delta)\psi'$, for all Borel sets $\Delta$, this gives $f_t\ge0$ almost everywhere for $\langle\psi',\mathcal P\psi'\rangle$, for all rational $t>0$ at once; there $\sinh\frac{tE}2\ge|\sin\frac{tq_1}2|$, and $t\downarrow0$ gives $E\ge|q_1|$. So $\mathcal P(\{E<|q_1|\})=0$, which is (R4) of Definition~\YMMref{95.1}{723}.

The insertion bound is obtained at $\Lambda<\infty$ and passed to the limit. For the time-zero field $\chi_0(0)=a(1)+a^*(1)$ on the Fock space over $L^2((0,\Lambda)\times\mathbb R^2,\mu_\Lambda)$, Lemma~\ref{C-353-26} gives $\|e^{-s\mathbf H_0}\chi_0(0)e^{-s'\mathbf H_0}\|\le\mathsf n_\Lambda(\min\{s,s'\})\le\mathsf n(\min\{s,s'\})$, where $\mathsf n_\Lambda$ is given by the integral below over $m<\Lambda$ and
\[\mathsf n(s)^2:=c\int_0^\infty e^{\sqrt m}dm\int d^2p\,\frac{e^{-s\omega_m(p)}}{\omega_m(p)}\Bigl(1+\frac1{s\,\omega_m(p)}\Bigr)<\infty\qquad(s>0),\]
with an absolute number $c$; $\mathsf n$ decreases, and it is finite because $e^{-s\omega_m(p)}\le e^{-sm/2}e^{-s|p|/2}$, $e^{\sqrt m-sm/2}$ is integrable and $\int_0^1dm\int_{|p|\le1}d^2p\,\omega_m(p)^{-2}<\infty$. As in Step~5, with one leg in place of two and without regularisation (the leg contracted upwards an annihilation, downwards a creation), for $0\le d'<d$, scalar tuples $F,G$ over $\{x_3>d\}$ and real scalar $h\in C^1_c(\mathbb R^3)$ with support in $\{|x_3|\le d'\}$,
\[\sigma_{0,\Lambda}(\theta\bar F\cdot(h)\cdot G)=\int h(x)\bigl\langle\mathsf V_\Lambda[\tau_{-de_3}F],e^{-(d+x_3)\mathbf H_0}e^{i\vec x\cdot\mathbf P_0}\chi_0(0)e^{-i\vec x\cdot\mathbf P_0}e^{-(d-x_3)\mathbf H_0}\mathsf V_\Lambda[\tau_{-de_3}G]\bigr\rangle\,dx,\]
so that
\[|\sigma_{0,\Lambda}(\theta\bar F\cdot(h)\cdot G)|\le\mathsf n(d-d')\,\|h\|_{L^1}\,\|[\tau_{-de_3}F]\|_{\sigma_{0,\Lambda}}\|[\tau_{-de_3}G]\|_{\sigma_{0,\Lambda}}.\]
The left side and the squares of the two norms are combinations of moments at disjoint supports, so the bound passes to $\sigma_0$. At $d=\eta/4$ and $d'=3\eta/16$ it gives the scalar insertion bound, (R5) of Definition~\YMMref{95.1}{723} with $c_I=\frac1{16}$ and with the finite function $N_0(\eta)=\frac{4\pi}3(\eta/16)^3\mathsf n(\eta/16)$ in place of $C_I\eta^{-\mathsf p_I}$. Then (a) gives (T1)--(T4) and (T5$_N$); centring holds because odd moments vanish.

\emph{Step 7 ((c): the growth of $K_2$).} At real $\xi\in\mathbb R^2_+$, by (KV5) the kernel is the continuous density of the two-point $\sigma$-formula, $K_2(\xi;\mathfrak a,\mathfrak a')=\iint S(\xi,s'-s)\operatorname{tr}\mathfrak a(s)\operatorname{tr}\mathfrak a'(s')\,ds\,ds'$, with $(\xi,u)$ the point of $\mathbb R^3$ with $(x_3,x_1)=\xi$ and $x_2=u$. For $\mathfrak a=\mathfrak a'=\varphi\mathbf 1$, $\varphi\ge c_\varphi>0$ on an interval, and small $t>0$,
\[K_2((t,0);\varphi\mathbf 1,\varphi\mathbf 1)\ge9\min_{|u|\le t}S(t,0,u)\iint_{|s-s'|\le t}\varphi(s)\varphi(s')\,ds\,ds'\ge c\,t\cdot t^{-1}e^{1/(4\sqrt2\,t)-\sqrt2\,t},\]
since $S$ is radial and decreasing, $|(t,0,u)|\le\sqrt2\,t$ for $|u|\le t$, and $\sqrt m-\sqrt2\,tm$ has the maximum $1/(4\sqrt2\,t)$ at $m=1/(8t^2)$ and stays above it minus $\sqrt2\,t$ on the next interval of length $1$. (KV1) would give $|K_2((t,0);\cdot)|\le C\min\{1,t\}^{-2\wp}$; so it holds for no $\wp$, and by Theorem~\ref{C-353-7}\,(C) $\sigma$ has (T5) for no $\mathsf p_I$.
\end{proof}

\begin{lemma}[An $L^1$ slab bound gives the insertion bound in the sup norm]\label{C-353-27}
Let $\sigma$ have (T1$^\circ$), (T2) and (T3), and let $C>0$, $\mathsf p\ge3$ be such that $|\sigma(\theta\bar F\cdot(h)\cdot G)|\le C\min\{1,d-d'\}^{-\mathsf p}\|h\|_{L^1}\|[\tau_{-de_3}F]\|\,\|[\tau_{-de_3}G]\|$ for all $0\le d'<d$, all real $h\in C^1_c(\mathbb R^3;V)$ with support in $\{|x_3|\le d'\}$ and all $F,G\in\mathbb C\mathfrak S^V(\{x_3>d\})$. Then $\sigma$ has (T5) in the sup norm with $c_I=\frac1{16}$, $\mathsf p_I=\mathsf p-3$ and $C_I=\frac{4\pi}3\,16^{\mathsf p-3}C$.
\end{lemma}

\begin{proof}
For $\eta\in(0,1]$ and $H$ with support in $B(y,\eta/16)$, $|y_3|\le\eta/8$, take $d=\eta/4$ and $d'=3\eta/16$, so that $d-d'=\eta/16\le1$, and use $\|H\|_{L^1}\le\frac{4\pi}3(\eta/16)^3\|H\|_\infty$:
\[|\sigma(\theta\bar F\cdot(H)\cdot G)|\le C(\tfrac\eta{16})^{-\mathsf p}\tfrac{4\pi}3(\tfrac\eta{16})^3\|H\|_\infty\|\cdot\|\,\|\cdot\|=\tfrac{4\pi}3\,16^{\mathsf p-3}C\,\eta^{-(\mathsf p-3)}\|H\|_\infty\|\cdot\|\,\|\cdot\|.\qedhere\]
\end{proof}

\begin{lemma}[A finite slab bound gives (T5$_N$)]\label{C-353-28}
Let $\sigma$ have (T1$^\circ$), (T2) and (T3), and suppose that for every $\eta\in(0,1]$ there is a finite number $\mathsf c^\circ(\eta)$ with $|\sigma(\theta\bar F\cdot(h)\cdot F')|\le\mathsf c^\circ(\eta)\|h\|_{L^1}\|[\tau_{-\eta e_3}F]\|\,\|[\tau_{-\eta e_3}F']\|$ for all $F,F'\in\mathbb C\mathfrak S^V(\{x_3>\eta\})$ and all real $h\in C^1_c(\{|x_3|<\eta/2\};V)$. Then $\sigma$ has (T5$_N$) with $c_I=\frac1{16}$ and $N(\eta)=\frac{4\pi}3\,16^{-3}\eta^3\,\mathsf c^\circ(\eta/8)$.
\end{lemma}

\begin{proof}
Let $H$ have support in $B(y,\eta/16)$, $|y_3|\le\eta/8$, and $F,G\in\mathbb C\mathfrak S^V(\{x_3>\eta/4\})$. Put $h':=\tau_{-y_3e_3}H$, supported in $\{|x_3|<\eta/16\}$, $F':=\tau_{y_3e_3}F$ and $F'':=\tau_{-y_3e_3}G$, both over $\{x_3>\eta/8\}$. By translation invariance and $\tau_{-y_3e_3}\theta=\theta\tau_{y_3e_3}$, $\sigma(\theta\bar F\cdot(H)\cdot G)=\sigma(\theta\bar F'\cdot(h')\cdot F'')$, which the hypothesis at $\eta/8$ bounds by $\mathsf c^\circ(\eta/8)\|h'\|_{L^1}\|[\tau_{-(\eta/8)e_3}F']\|\,\|[\tau_{-(\eta/8)e_3}F'']\|$. Now
\[[\tau_{-(\eta/8)e_3}F']=\mathcal T_{\eta/8+y_3}[\tau_{-(\eta/4)e_3}F],\qquad[\tau_{-(\eta/8)e_3}F'']=\mathcal T_{\eta/8-y_3}[\tau_{-(\eta/4)e_3}G],\]
with $\eta/8\pm y_3\ge0$ and $\|\mathcal T_t\|\le1$; and $\|h'\|_{L^1}\le\frac{4\pi}3(\eta/16)^3\|H\|_\infty\le\frac{4\pi}3(\eta/16)^3\|H\|_\sharp$.
\end{proof}

\begin{proof}[Proof of Theorem~\ref{C-353-14}]
\emph{Step 1 ((a)).} By (I), every Gibbs sequence at $g$ has a subsequence that converges at every tuple of $\mathfrak S^V(\mathbb R^3)$, with a further subsequence along which $a_{K_j}\beta_j$ converges to a number $b_*\in[g^{-2},L(g^{-2}+\bar B)]$. By the step of this section that ends item (d) of the remark on how (I) and (II) are obtained (the density of $C^1_c$ in $C_c$ for the sup norm, the bound of (I), multilinearity and a telescoping sum), this further subsequence converges at every tuple of continuous $V$-valued functions with compact, pairwise disjoint supports, and its limit is a tensor limit at $g$ with the number $b_*$; as a subsequence of the first, it converges in particular at every tuple of $\mathfrak S^V(\mathbb R^3)$.

Gibbs sequences at $g$ exist. Indeed $g\le\gamma^{\rm ref}_W\le\gamma$, with the constant $\gamma$ (the choice preceding equation~\YMMref{(2326)}{845}); so the interval $I(g)$, defined in~(\YMMref{(1527)}{615}), is not empty, and for every $\mathsf b\in I(g)$ the pairs $(K,L^K\mathsf b)$, $K\in\mathbb N$, are first-passage data at $g$, with $a_KL^K\mathsf b=\mathsf b$ (Corollary~\YMMref{74.2}{616}). At each of them the set of the states of the tori is not empty and consists of periodic Gibbs states (Proposition~\YMMref{91.2}{699}); choosing one such state $\nu_K$ at each $K$ gives the Gibbs sequence $(K,L^K\mathsf b,\nu_K)_{K\ge1}$ at $g$, and what precedes, applied to it, gives a tensor limit at $g$.

\emph{Step 2 ((b)(i), (iii)).} Each $\mathcal O^{\mathsf T}_K(H)$ is a real function of the configuration and linear in $H$. So, for open sets $O_1,\dots,O_n$ with pairwise disjoint closures, the numbers $\nu_j\bigl(\prod_a(\mathcal O^{\mathsf T}_{K_j}(H_a)-\nu_j(\mathcal O^{\mathsf T}_{K_j}(H_a)))\bigr)$ are real and $n$-linear in $(H_a)\in\prod_aC^1_c(O_a;V)$, and so is their limit $\sigma$; with $\sigma(\emptyset)=1$, $\sigma$ is a tensor functional. By (II), $\sigma$ is symmetric and centred; the bounds of (II) with $\sup|H_a|_\infty\le\|H_a\|_\infty\le\|H_a\|_\sharp$ give (T2); the covariance of (II) is (T1); positivity at $\{x_3=0\}$ is (T3). The joint spectral measure of (II) is the measure $\mathcal P$ of Lemma~\ref{C-353-19}\,(b), by its uniqueness, so (T4) holds and $\mathcal P$ is carried by $\{E\ge\max(|q_1|,|q_2|)\}$. (iii) is part of (II).

\emph{Step 3 ((b)(ii)).} As $\int|h(x)|_\infty dx\le\|h\|_{L^1}$, (II) makes $\mathsf c^\circ_\sigma(\eta)$ finite, and Lemma~\ref{C-353-28} gives (T5$_N$) with
\[c_I=\tfrac1{16},\qquad N(\eta)=\tfrac{4\pi}3\,16^{-3}\eta^3\,\mathsf c^\circ_\sigma(\eta/8).\]

\emph{Step 4 ((b)(iv)).} By Steps 2 and 3, $\sigma$ has (T1)--(T4) and (T5$_N$), which are the hypotheses of the clauses named.

\emph{Step 5 ((c)(i), (ii)).} A tensor limit converges at every tuple of continuous members, in particular at every tuple of $\mathfrak S^V(\mathbb R^3)$, so $\sigma^1$ is the limit of a convergent Gibbs sequence. The classes of tuples of $C^1$ members span a closed subspace of the space reconstructed from tuples of continuous members, with the same scalar products; so the bound (III) on the centred slab operators gives, in matrix elements, $|\sigma^1(\theta\bar F\cdot(h)\cdot G)|\le\|B^c_d(h)\|\,\|[\tau_{-de_3}F]\|\,\|[\tau_{-de_3}G]\|$, and with $\int|h|_\infty\le\|h\|_{L^1}$ this is (i). Lemma~\ref{C-353-27} with $C=C_{\mathsf p}/b_*$ gives (ii); then (T1)--(T5) hold, and (C) of Theorem~\ref{C-353-7} in the sup norm gives (KV1) with $\wp'=\mathsf p_I+3=\mathsf p$.

\emph{Step 6 ((c)(iii)).} By Proposition~\ref{C-353-23} and Lemma~\ref{C-353-22}\,(b) with $c=\varepsilon-\delta$, $B^\delta_{\mathfrak a}(\varepsilon n_0,\varepsilon n_0)=B^c_\varepsilon(j_\delta\otimes\mathfrak a)$, and for real $\mathfrak a$, by (i) with $d'=\delta/2$ and $\|j_\delta\otimes\mathfrak a\|_{L^1}=\|\mathfrak a\|_{L^1}$,
\[\|B_{\mathfrak a}(\varepsilon n_0,\varepsilon n_0)\|=\lim_{\delta\to0}\|B^c_\varepsilon(j_\delta\otimes\mathfrak a)\|\le C_{\mathsf p}b_*^{-1}\min\{1,\varepsilon\}^{-\mathsf p}\|\mathfrak a\|_{L^1}.\]
The shift rule with contractions extends this to $d_{\rm Tb}(\xi),d_{\rm Tb}(\xi')\ge\varepsilon$; a complex label is split, and $\|\mathfrak a^0\|_{L^1}+\|\mathfrak a^1\|_{L^1}\le\sqrt2\|\mathfrak a\|_{L^1}$. The kernel bound follows at $\kappa=n_0$, factor by factor, with $\min\{1,d_{\rm Tb}(p_k),d_{\rm Tb}(q_k)\}\ge\frac12\min\{1,\min_kd_{\rm Tb}(\xi_k)\}$.

\emph{Step 7 ((c)(iv)).} The point bound is obtained as in Step 6 from Proposition~\ref{C-353-25} and $B^\delta_E(\varepsilon e_3,\varepsilon e_3)=B^c_\varepsilon(j^{(3)}_\delta E)$. For $w\in W_3$ with $we_3=e_3$, $\mathcal U_wB^c_d(H)\mathcal U_w^{-1}=B^c_d(wH)$, and the family $\mathcal U_w^{-1}B_{\varrho(w)E}(w\mathfrak z,w\mathfrak z')\mathcal U_w$ is weakly continuous, has the shift rule by Lemma~\ref{C-353-20}\,(b), and satisfies the slab formula at $hE$, because $w(hE)=(h\circ w^{-1})\varrho(w)E$, $|\det w|=1$ and $de_3+wx=w(de_3+x)$; by (D)(iv) it is $B_E$. The functions $W_n$ do not depend on $\kappa$ and are holomorphic and $n$-linear by (D) and the calculus; the bound follows at $\kappa=e_3$, factor by factor. The $\sigma$-formula is proved as (KV2) in Step 3 of the proof of Theorem~\ref{C-353-7}, with the product formula, the slab formula (D)(iii) in each factor and the rules (D)(iv):
\[[(h_1,\dots,h_n)]=\int B_{h_1(x_1)}(x_1,c_1e_3-x_1)B_{h_2(x_2)}(x_2-c_1e_3,c_2e_3-x_2)\cdots\Omega\,dx_1\cdots dx_n\]
for heights $0=c_0<c_1<\dots$ separating the supports, after a translation in time. For $w\in W_3$, $\sigma^1(X)=\sigma^1(w^{-1}X)$ by (T1), $(w^{-1}h)(x)=\varrho(w^{-1})h(wx)$, and the substitution $x_k\mapsto w^{-1}x_k$ gives the second form, as $(w^{-1}x)_3=x\cdot we_3$.

\emph{Step 8 ((c)(v)).} At $\kappa=n_0$ the factors of $\Psi(R_{is}z;\mathbf a)$ have the arguments $R_{is}z_1$, $\frac12R_{is}(z_{i+1}-z_i)$ and $n_0$, whose $d_{\rm Tb}$ are $e^{-|s|}z_1^0$, $\frac12e^{-|s|}(z_{i+1}-z_i)^0$ and $1$ by Lemma~\ref{C-353-17}\,(c), all $\ge e^{-|s|}c_{\mathcal C}$; the plate bound of (iii) gives
\[\|\Psi(R_{is}z;\mathbf a)\|\le\prod_i\bigl(\sqrt2\,C_{\mathsf p}b_*^{-1}\|\mathfrak a_i\|_{L^1}\bigr)\,(e^{-|s|}c_{\mathcal C})^{-\mathsf pj}.\qedhere\]
\end{proof}

\subsection{Scope}

Theorems~\ref{C-353-7}--\ref{C-353-12} are statements on functionals, and $\Omega$ enters them only through $\|\Omega\|=1$ and $T(\xi)\Omega=\Omega$. The trivial functional, with $\sigma(X)=0$ for $X\ne\emptyset$, $\mathcal H(\sigma)=\mathbb C\Omega$ and $K_n=0$, belongs to the class. All the bounds are short-distance bounds; at distances $\ge1$ they are constants. The formulas (KV2), and the formula for $W_n$ in Theorem~\ref{C-353-14}, hold at tuples whose members are ordered in time (for $W_n$, in a coordinate direction), with the members of (KV2) of the form $\alpha_k\otimes\mathfrak a_k$. Of the rotations, Theorem~\ref{C-353-10} uses the quarter turns about $e_2$.

\section{Conservation and rotation covariance}\label{cov-sec}

The lattice supplies cubic symmetry, reflection positivity and a square spectral cone. Full rotation covariance requires more: conservation removes the two-point angular growth, and positivity propagates this fact to correlations with any number of insertions.

We first state the model and abstract results. The lattice estimates verify the abstract hypotheses; the analytic part uses point plates and their angular Fourier coefficients. The polygon argument controls the extreme coefficient, while the halving argument relates its degree to the two-point degree. The latter is zero by conservation and the off-diagonal bound.

\subsection{Notation and hypotheses}\label{cov-set}

\emph{Constants and data.} The radius, the constants ($M,\alpha_0,\mathfrak p_0,\dots$) and the fixed data are those fixed in the opening notation and hypotheses, with the common construction choice of Section~\YMMref{82}{665}; $L\ge13$ is the odd blocking factor; $\gamma^{\rm sd}_W\le\gamma^{\rm ref}_W\le\gamma$ are the two ceilings of the choices in equations~\YMMref{(2326)--(2327)}{845}, functions of the constants alone. A \emph{reference coupling} is a number $0<g\le\gamma$ (Section~\YMMref{74}{614}). We write $a_K=L^{-K}$, $\mathbb L_K=a_K\mathbb Z^3$, and $\mathcal G^{\rm per}_{K,\beta}$ for the set of the Gibbs states of the Wilson weight $e^{-\beta\sum_pa_p}$ on $\mathbb L_K$ that are invariant under a subgroup of finite index of $a_K\mathbb Z^3$, the \emph{periodic Gibbs states}.

For a reference coupling $g$, $\mathfrak d(g)$ is the set of the first-passage data $(K,\beta)$ at $g$ (Corollary~\YMMref{74.2}{616}), and $b=a_K\beta=L^{-K}\beta$ is the \emph{label} of a datum. By Proposition~\YMMref{74.1}{615} and Corollary~\YMMref{74.2}{616} (the flow and the windows of first passage), $g^{-2}<b\le L(g^{-2}+\bar B)$ for $(K,\beta)\in\mathfrak d(g)$, and $(K,\beta)\in\mathfrak d(g)$ whenever $b\in(g^{-2}+\bar B,X_g-\bar B]$, where $X_g=L(g^{-2}+B_\sharp)$ and $\bar B=LB_\sharp/(L-1)$; so there are data at every $K$. The \emph{run} of a datum $(K,\beta)\in\mathfrak d(g)$ is the flow $(x_0(\beta),\dots,x_K(\beta))$ of Proposition~\YMMref{74.1}{615}: $x_0(\beta)=\beta$ and $x_j(\beta)=L^{-1}x_{j-1}(\beta)-\mathsf e_j(\beta)$ with increments $|\mathsf e_j(\beta)|\le B_\sharp$, so that, by \YMMref{(1530)}{615} and \YMMref{(1529)}{615},
\[
\begin{multlined}
\bigl|x_j(\beta)-L^{-(j-i)}x_i(\beta)\bigr|\le\bar B\bigl(1-L^{-(j-i)}\bigr)\quad(0\le i\le j\le K),\qquad g^{-2}<x_K(\beta) \\
\le X_g<x_{K-1}(\beta)<\dots<x_0(\beta);
\end{multlined}
\]
the label $b=L^{-K}x_0(\beta)$ lies within $\bar B$ of $x_K(\beta)$. The run is not the line $x_k=L^{K-k}\mathsf b'$ of a label $\mathsf b'$, which enters the setting (Th) of Subsection~\ref{cov-cut}. A side $\ell=2L^m$ is \emph{small over the run} if it satisfies the conditions (Sd.1)--(Sd.4) of equation~\YMMref{(2312)}{841}: the first three make it admissible at every level, and the fourth makes it so small that the suppression of large fields at the last level beats the factor $e^{\mathsf B(x_K)\ell^3}$ of that definition. For $1\le i\le K-1$ the \emph{stopped coupling} $g_{(i)}$ of \YMMref{(1536)}{616} is given by
\[g_{(i)}^{-2}=L^{-1}x_{K-i}(\beta)-B_\sharp,\qquad\text{so that}\qquad X_{g_{(i)}}=x_{K-i}(\beta);\]
$g_{(i)}^{-2}$ is the least value of $g'^{-2}$ at which $(K-i,\beta)$ is a first-passage datum at $g'$, and $g_{(i)}<g$; the run of $(K-i,\beta)$ is the cut $(x_0(\beta),\dots,x_{K-i}(\beta))$ of the run of $(K,\beta)$, and its label is $L^ib$ (Corollary~\YMMref{74.3}{616} and Lemma~\ref{lem:setup-stopped-arithmetic}). The sides at $g$ form the unbounded set $\mathfrak L(g)=\{2L^m:m\ge m_{\mathfrak L}(g)\}$ of equation~\YMMref{(1528)}{615}, with the least element $\ell'$. For $g\le\gamma^{\rm sd}_W$ the least side $\ell'$ is small over the run of every first-passage datum at $g$ (Lemma~\YMMref{126.3}{846}). $\langle\cdot\rangle_{K,\ell,\beta}$ is the Wilson measure on the torus $T_\ell$ of side $\ell$ at the spacing $a_K$, and $p_\ell=\ell^{-3}\log Z_{K,\ell,\beta}$ its pressure. Boards of $\mathbb L_K$, their cells and walls and the mirror copies $F^{[\kappa]}$ of a function $F$ of one cell are those of Section~\ref{sec:periodic-chessboard}, and the chessboard seminorms $\|F\|_{{\rm ch},\ell}$ those of equation~\YMMref{(1756)}{681} and Corollary~\ref{IV-18-18}; on the least torus with cells of side $\ell'/2$ the board has eight cells, and $\|F\|^8_{{\rm ch},\ell'}=\langle\prod_\kappa F^{[\kappa]}\rangle_{K,\ell',\beta}$. We use the following statement, which collects Propositions~\YMMref{103.4}{758} and~\YMMref{91.9}{705}:

\smallskip
\noindent\textup{(Per)} \emph{The periodic Gibbs states form a non-empty set, are the Gibbs states invariant under $a_K\mathbb Z^3$, and every $\nu\in\mathcal G^{\rm per}_{K,\beta}$ is invariant under $a_K\mathbb Z^3\rtimes W_3$, reflection positive at the three families of coordinate lattice mirrors and the six families of diagonal lattice mirrors, and dominated at finitely many cells by chessboard seminorms along every sequence of sides $\ell_k\to\infty$ in $\mathfrak L(g)$: $|\nu(\prod_{k\in\mathcal K_0}F_k)|\le\liminf_k\prod_{k\in\mathcal K_0}\|F_k\|_{{\rm ch},\ell_k}$ for functions $F_k$ of finitely many distinct cells of one board.}

\smallskip
\emph{The lattice tensor.} We use the forms of Definition~\ref{C-355-3}, with the same midpoint convention:
\begin{equation}\label{cov-eq-tensor}
 \mathcal T_K(H)=O_K^T(H)=\mathcal O(w_H)+Y[\mathsf w_H].
\end{equation}
Besides the differences $u_{b,\alpha}$ in the introduction, we use the sums
\[
 \mathbf y_{b,\alpha}=w(W_x(\mu,\alpha))+w(W_x(\mu,-\alpha)),
 \qquad \mathbf u_{b,\alpha}=u_{b,\alpha}.
\]
The same definitions apply on a torus.

\emph{Tuples and limits.} An isometry $\mathsf r:x\mapsto\mathsf Rx+v$ acts by $(\mathsf rH)(x)=\rho(\mathsf R)H(\mathsf r^{-1}x)$. For an open set $O$, $\mathfrak S_V(O)$ is the set of the tuples $X=(H_1,\dots,H_n)$, $n\ge0$, of members $H_a\in C_c(O;V)$ with pairwise disjoint supports, and $\mathbb C\mathfrak S_V(O)$ the set of their complex combinations; a tuple with members in $C_c(O;V_{\mathbb C})$ stands for its expansion into real and imaginary parts. Isometries act member by member, $X\cdot Y$ is the concatenation, $(h)$ the tuple with the one member $h$, and $X$ is of \emph{type $(\mathsf d;R)$} if each support lies in a closed ball of radius $R$ and any two supports are at distance $\ge\mathsf d$. For a state $\nu$ put $A^\nu_K(X)=\prod_a[\mathcal T_K(H_a)-\nu(\mathcal T_K(H_a))]$. Let $(K_j,\beta_j)\in\mathfrak d(g)$ with $K_j\to\infty$, $\rho_j\in\mathcal G^{\rm per}_{K_j,\beta_j}$ and $a_{K_j}\beta_j\to b_*$. If
\begin{equation}\label{cov-eq-limit}
\sigma(X)=\lim_{j\to\infty}\rho_j\bigl(A^{\rho_j}_{K_j}(X)\bigr)
\end{equation}
exists for every $X\in\mathfrak S_V(\mathbb R^3)$, we call $\sigma$ a \emph{tensor limit at $g$} with the \emph{label} $b_*\in[g^{-2},L(g^{-2}+\bar B)]$, and $\sigma_0$, its restriction to the tuples with members in $C_c(\mathbb R^3;V_0)$, its \emph{scalar part}. For $\eta\in C^2_c(\mathbb R^3;\mathbb R^3)$ let
\[H_\eta=(\nabla\!\cdot\!\eta)\mathbf 1-\nabla\eta-(\nabla\eta)^{\rm T},\qquad(\nabla\eta)_{ij}=\partial_i\eta_j .\]
Then $\frac12\operatorname{Tr}(H_\eta T)=\sum_{i,j}\partial_i\eta_j(\frac12\delta_{ij}\operatorname{tr}T-T_{ij})$, so that $\int$ of it vanishes for all $\eta$ exactly when $\frac12\delta_{ij}\operatorname{tr}T-T_{ij}$ is conserved, and $\mathsf rH_\eta=H_{\mathsf r\eta}$ with $(\mathsf r\eta)(x)=\mathsf R\eta(\mathsf r^{-1}x)$; $\|\eta\|_{C^2}$ is the sum of the suprema of the entries of $\eta$, $\nabla\eta$ and $\nabla^2\eta$.

\emph{Euclidean data.} $\theta$ is the reflection $x_3\mapsto-x_3$, acting on points and, through $\rho$, on values; $\theta_c$ is the reflection in the plane $x_3=c$; $\mathbb E_+=\{x_3>0\}$; $W_3$ is the group of the $48$ signed permutation matrices; $\mathsf R_t$ is the rotation about $e_2$ with $\mathsf R_te_3=\cos t\,e_3+\sin t\,e_1$; $\tau_v$ is the translation by $v$; $\Sigma$ is the diagonal plane $x_1+x_3=0$ and $\theta_\Sigma\in W_3$ the reflection in it.

\begin{definition}\label{cov-def}
A real functional $\sigma$ on $\mathfrak S_V(\mathbb R^3)$, continued linearly to complex combinations and to members with values in $V_{\mathbb C}$, is \emph{labelled data} if
\begin{itemize}
\item[(D1)] $\sigma(\varnothing)=1$ and $\sigma((H))=0$; $\sigma$ is symmetric and $n$-linear on $\prod_aC_c(O_a;V)$ for open sets $O_a$ with pairwise disjoint closures; and $|\sigma(X)|\le C_n(\mathsf d;R)\prod_a|H_a|_\infty$ for $X$ of type $(\mathsf d;R)$ with $n$ members;
\item[(D2)] $\sigma(\mathsf rX)=\sigma(X)$ for $\mathsf r\in\mathbb R^3\rtimes W_3$;
\item[(D3)] $\sigma(\theta_{\Sigma'}\bar F\cdot F)\ge0$ for every plane $\Sigma'$ of the form $x_\mu=c$ or $x_\mu\pm x_\nu=c$, with the reflection $\theta_{\Sigma'}$, and every $F\in\mathbb C\mathfrak S_V$ carried by one open side of $\Sigma'$;
\item[(D4)] in the Hilbert space $\mathcal H(\sigma)$ of the form $\langle[F'],[F]\rangle=\sigma(\theta\bar F'\cdot F)$ on $\mathbb C\mathfrak S_V(\mathbb E_+)$, with $\Omega=[\varnothing]$, $e^{-tH}[F]=[\tau_{te_3}F]$, $e^{iv\cdot P}[F]=[\tau_{(v,0)}F]$ and the joint spectral measure $\mathcal P$ of $H$ and $P$: $H\ge|P_1|$ and $H\ge|P_2|$.
\end{itemize}
Labelled data may have the following properties.
\begin{itemize}
\item[(S)] For $l>0$ and $h\in C_c(\{0<x_3<l\};V_{\mathbb C})$ there is a bounded operator $\mathsf B_l(h)$ on $\mathcal H(\sigma)$ with $\mathsf B_l(h)[F]=[(h)\cdot\tau_{le_3}F]$ and $\|\mathsf B_l(h)\|\le\mathsf C(\mathcal K,l)|h|_\infty$ when $\operatorname{supp}h$ lies in a compact set $\mathcal K$; nothing is supposed of the numbers $\mathsf C(\mathcal K,l)$.
\item[(O)] There is $C_{\rm od}$ with $\sigma(\theta(h\,o),h\,o)\le C_{\rm od}|h|_\infty^2$ for every $o=o_{\mu\nu}=e_\mu e_\nu^{\rm T}+e_\nu e_\mu^{\rm T}$, $\mu\ne\nu$, every $0<\mathsf d<1$, and every real $h\in C_c(B(x_0,\mathsf d))$ with $x_{0,3}\in[\frac32\mathsf d,2\mathsf d]$.
\item[(O$_a$)] For one $a<4$ there are $C_a$ and a sequence $1>\mathsf d_n\downarrow0$ with $\sigma(\theta(h\,o),h\,o)\le C_a\mathsf d_n^{-a}|h|_\infty^2$ for all $n$, all $o$ as in (O) and all real $h\in C_c(B(2\mathsf d_ne_3,\mathsf d_n))$.
\item[(N)] $\sigma(\theta H_\eta,H_\eta)=0$ for every $\eta\in C^2_c(\mathbb E_+;\mathbb R^3)$; that is, $[(H_\eta)]=0$ in $\mathcal H(\sigma)$.
\item[(P)] For one $p<\infty$ there is $C$ with $\|\mathsf B_{2d}(\tau_{de_3}h)\|\le C\min\{1,d-d'\}^{-p}\|h\|_{L^1}$ whenever $0\le d'<d$ and $h\in C_c(\{|x_3|\le d'\};V_{\mathbb C})$, where $\|h\|_{L^1}=\int|h(x)|_\infty\,dx$.
\item[(I$_0$)] $\sigma(\mathsf R_tX)=\sigma(X)$ for all real $t$ and all tuples $X$ with members in $C_c(\mathbb R^3;V_0)$.
\end{itemize}
\end{definition}

\subsection{Principal results}\label{cov-res}

The numbers called constants below are functions of the fixed construction data specified in the opening notation and hypotheses.

\begin{theorem}\label{cov-thm}
Let $0<g\le\gamma$.
\begin{enumerate}
\item[(a)] \emph{(Conservation at a cutoff.)} Let $q\ge2$ and $\mathsf d,R>0$. There are $C$, $p$ and $K_0$, functions of the constants, $g$, $q$, $\mathsf d$ and $R$ alone, such that for all $(K,\beta)\in\mathfrak d(g)$ with $K\ge K_0$, all $\nu\in\mathcal G^{\rm per}_{K,\beta}$, all $\eta\in C^2_c(\mathbb R^3;\mathbb R^3)$ and all bounded $H_2,\dots,H_q:\mathbb R^3\to V$ whose supports, with $\operatorname{supp}\eta$, lie in a cube of side $2R$ at pairwise sup distances $\ge\mathsf d$,
\begin{equation}\label{cov-eq-cons}
\bigl|\nu\bigl(A^\nu_K(H_\eta,H_2,\dots,H_q)\bigr)\bigr|\le CK^pL^{-K/36}\,\|\eta\|_{C^2}\prod_{a\ge2}|H_a|_\infty ;
\end{equation}
the same bound holds for the joint cumulant of $\mathcal T_K(H_\eta),\mathcal T_K(H_2),\dots,\mathcal T_K(H_q)$ in $\nu$, and, with numbers that depend also on the side, for the joint cumulant in $\langle\cdot\rangle_{K,\ell,\beta}$, $\ell\in\mathfrak L(g)$, of members on $T_\ell$ at pairwise sup distances $\ge\mathsf d$. For $q=2$ and $H_2=H_{\eta'}$ the right side may be replaced by $CK^pL^{-K/20}\|\eta\|_{C^2}\|\eta'\|_{C^2}$.
\item[(a$'$)] \emph{(The mirror form.)} Let $h_0,R_0>0$ and $c<\frac1{10}$. There are $C$, $p$, $K_0$, functions of the constants, $g$, $h_0$, $R_0$ and $c$, such that for $(K,\beta)\in\mathfrak d(g)$, $K\ge K_0$, $\nu\in\mathcal G^{\rm per}_{K,\beta}$, the reflection $\theta_K$ in $x_3=\frac12a_K$ and $\eta\in C^2_c(\mathbb R^3;\mathbb R^3)$ with support in $\{x_3\ge h_0,\ |x|_\infty\le R_0\}$,
\[0\le\nu\bigl(A^\nu_K(\theta_KH^K_\eta,H^K_\eta)\bigr)\le b^{-2}\bigl(CK^pL^{-cK}\|\eta\|_{C^2}\bigr)^2,\]
$H^K_\eta$ being $H_\eta$ with the lattice differences of Lemma~\ref{cov-var}.
\item[(b)] \emph{(Conservation in the limit.)} Every tensor limit $\sigma$ at $g$ satisfies $\sigma((H_\eta)\cdot X)=0$ for every $\eta\in C^2_c(\mathbb R^3;\mathbb R^3)$ and every $X\in\mathfrak S_V(\mathbb R^3)$ whose members have supports disjoint from $\operatorname{supp}\eta$.
\item[(c)] \emph{(Covariance.)} Every tensor limit $\sigma$ at $g$ satisfies $\sigma(\mathsf RX)=\sigma(X)$ for all $\mathsf R\in\mathrm O(3)$ and all $X\in\mathfrak S_V(\mathbb R^3)$.
\end{enumerate}
\end{theorem}

Let $K_3(u)=(3\hat u\hat u^{\rm T}-\mathbf 1)/(4\pi|u|^3)$, $\hat u=u/|u|$, and
\[
\begin{multlined}
\mathcal W(H_1,H_2)=\\
\tfrac32\iint\operatorname{Tr}\bigl[H_1(x)\,K_3(x-y)\,H_2(y)\,K_3(x-y)\bigr]\,dx\,dy .
\end{multlined}
\]
For three free photon fields $B^{\mathsf a}$ with 
\[
\langle B^{\mathsf a}_\lambda(x)B^{\mathsf b}_{\lambda'}(y)\rangle=\mathsf b^{-1}\delta^{\mathsf a\mathsf b}K_3(x-y)_{\lambda\lambda'}
\]
 and $T_{\lambda\lambda'}=\sum_{\mathsf a}{:}B^{\mathsf a}_\lambda B^{\mathsf a}_{\lambda'}{:}$, the covariance of $\frac12\int\operatorname{Tr}(H_1T)$ and $\frac12\int\operatorname{Tr}(H_2T)$ is $\mathsf b^{-2}\mathcal W(H_1,H_2)$; at $H_a=\phi_a\mathbf 1$ it is 
\[
\frac9{16\pi^2}\mathsf b^{-2}\iint\phi_1(x)\phi_2(y)|x-y|^{-6},
\]
 and $|\mathcal W(H_1,H_2)|\le24|H_1|_\infty|H_2|_\infty$ at the positions of Theorem~\ref{cov-thm-nz}\,(i).

\begin{theorem}\label{cov-thm-nz}
Let $0<g\le\gamma^{\rm sd}_W$. There are $\mathsf d_0(g)>0$ and a function $e_g$ on $(0,\mathsf d_0(g))$, formed with the constants and $g$ alone, with $e_g(\mathsf d)\to0$ as $\mathsf d\to0$ and $\sup_{\mathsf d}(\log\frac1{\mathsf d})^{\mathfrak p_0/4}e_g(\mathsf d)<\infty$, such that the following holds for every tensor limit $\sigma$ at $g$ with label $b_*$.
\begin{enumerate}
\item[(i)] For $0<\mathsf d_-<\mathsf d<\mathsf d_0(g)$, $x_0,y_0$ with $3\mathsf d\le|x_0-y_0|\le4\mathsf d$, and $H_1,H_2\in C^1_c(\mathbb R^3;V)$ with supports in the closed balls of radius $\mathsf d_-$ about $x_0$ and $y_0$,
\[\bigl|\sigma(H_1,H_2)-b_*^{-2}\mathcal W(H_1,H_2)\bigr|\le e_g(\mathsf d)\,b_*^{-2}|H_1|_\infty|H_2|_\infty .\]
\item[(ii)] The label $b_*$ is determined by $\sigma$.
\item[(iii)] Let $w_0=\frac1{12}\,388^{-6}$, $e_g(\mathsf d)<w_0$, $x_0=\frac32\mathsf de_3$, and $\varphi\in C^1_c(\mathbb R^3)$ with $\mathbf 1_{B(x_0,\mathsf d/128)}\le\varphi\le\mathbf 1_{B(x_0,\mathsf d/64)}$. Then $\sigma(\theta(\varphi E),\varphi E)\ge(w_0-e_g(\mathsf d))\,b_*^{-2}\operatorname{tr}E^2$ for every $E\in V$: the vectors $[(\varphi E)]$ are nonzero for $E\ne0$, and no nonzero real combination of the six fields vanishes in any tensor limit.
\item[(iv)] With the same $e_g$, the bound (i) holds for the limit functionals of Proposition~\ref{cov-torus} on every torus $T_\ell$, $\ell\in\mathfrak L(g)$, at every label $\mathsf b\in(g^{-2}+\bar B,X_g-\bar B]$, for Lipschitz members, with $\mathsf b$ in the place of $b_*$ and $\mathcal W$ formed with lifts of the members to balls about lifts of $x_0,y_0$ at the distance $|x_0-y_0|$.
\end{enumerate}
\end{theorem}

The next two theorems concern abstract data; no lattice occurs in them, and the kernels $\mathcal K_n$, the plate vectors $\Psi$, the tubes $\mathcal Z_n$, the weight labels and the imaginary rotations $\mathsf R_{is}$ in them are those of Subsection~\ref{cov-abs} and Lemma~\ref{cov-plate}.

\begin{theorem}\label{cov-thm-data}
Let $\sigma$ be labelled data with the properties (S), (O), (N) and (P). Then:
\begin{enumerate}
\item[(i)] $\sigma(\mathsf RX)=\sigma(X)$ for all $\mathsf R\in\mathrm O(3)$ and all $X\in\mathfrak S_V(\mathbb R^3)$;
\item[(ii)] $\sigma((H_\eta)\cdot X)=0$ for every $\eta\in C^2_c(\mathbb R^3;\mathbb R^3)$ and every $X\in\mathfrak S_V(\mathbb R^3)$ with members off $\operatorname{supp}\eta$;
\item[(iii)] $\mathcal K_n(\mathsf R_{is}Z;\underline\lambda)=e^{\kappa(\underline\lambda)s}\mathcal K_n(Z;\underline\lambda)$ for all $n\ge2$, all tuples $\underline\lambda$ of weight labels, all $Z\in\mathcal Z_n$ and all real $s$;
\item[(iv)] $e^{-\kappa(\underline\lambda)s}\|\Psi(\mathsf R_{is}z;\underline\lambda)\|=\|\Psi(z;\underline\lambda)\|$ for all real configurations $z$ in the domain of $\Psi$, all tuples of weight labels and all real $s$; and there is a unitary group $U(s)$ on the closed span of $\Omega$ and the plate vectors with $U(s)\Omega=\Omega$ and $U(s)\Psi(w;\underline\lambda)=e^{-\kappa(\underline\lambda)s}\Psi(\mathsf R_{is}w;\underline\lambda)$.
\end{enumerate}
The same holds if (P) is replaced by any growth of the plate vectors along the imaginary rotations with exponents $A_j$ such that $\liminf_jA_j/j^2<2$ (Proposition~\ref{cov-rot}), and (O) may be replaced by (O$_a$) when (I$_0$) holds at pairs. Under (P), or under a growth with $\liminf_j(A_j-j^2)/j=-\infty$, it suffices that (D3) hold at the plane $x_3=0$ and, at tuples of one member, at the plane $\Sigma$.
\end{theorem}

\begin{theorem}\label{cov-thm-rider}
Let $\sigma$ be labelled data with the properties (S), (O), (N) and (I$_0$). Then $\sigma(\mathsf RX)=\sigma(X)$ for all $\mathsf R\in\mathrm O(3)$ and all tuples $X=(H_1,H_2,f_1\mathbf 1,\dots,f_n\mathbf 1)\in\mathfrak S_V(\mathbb R^3)$, $n\ge0$. The same holds for labelled data with (S) and (I$_0$) that are covariant under the $\mathsf R_t$ at all pairs.
\end{theorem}

Tensor limits at $0<g\le\gamma^{\rm sd}_W$ have (S), (O) and (N) by Propositions~\ref{cov-lim}, \ref{cov-od} and Theorem~\ref{cov-thm}\,(a), and their scalar part is invariant under $\mathbb R^3\rtimes\mathrm O(3)$ by Theorem~\YMMref{AN}{763} (limit functionals along periodic Gibbs states), whose proof rests on Theorem~\YMMref{AK}{737} (the rotation theorem for the action density); so Theorem~\ref{cov-thm-rider} gives Theorem~\ref{cov-thm}\,(c) at these tuples without the power of Proposition~\ref{cov-P}.

\begin{corollary}\label{cov-cor-real}
Let $0<g\le\gamma^{\rm sd}_W$, $\sigma$ a tensor limit at $g$ with label $b_*$, $\sigma^1$ its restriction to the tuples with members in $C^1_c(\mathbb R^3;V)$, and $p>\frac{17}2$.
\begin{enumerate}
\item[(A)] $\sigma^1$ is invariant under all isometries of $\mathbb R^3$; it is reflection positive at $x_3=0$, with $\mathcal H(\sigma^1)=\mathcal H(\sigma)$ and the same $\Omega$, $H$, $P$, $\mathcal P$; the spectral measure is carried by the cone $\{E\ge|q|\}$, that is $\mathcal P(\{E<|q|\})=0$; and $|\sigma(\theta\bar F\cdot(h)\cdot F')|\le C_V\eta^{-p}\|h\|_{L^1}\|[\tau_{-\eta e_3}F]\|\,\|[\tau_{-\eta e_3}F']\|$ for $\eta\in(0,1]$, $F,F'$ over $\{x_3>\eta\}$, $h\in C^1_c(\{|x_3|<\frac12\eta\};V)$, with $C_V=2^pC_pb_*^{-1}\le2^pC_pg^2$, $C_p$ the number of Proposition~\ref{cov-P}; so the number $2^pC_pg^2$ serves for all tensor limits at $g$.
\item[(B)] For every $n\ge2$ there is exactly one continuous function $\mathsf s_n(x;e)$ on the configurations of $n$ pairwise different points of $\mathbb R^3$ and $e\in V^n$, $n$-linear in $e$, with $\sigma(H_1,\dots,H_n)=\int\mathsf s_n(x;H_1(x_1),\dots,H_n(x_n))\,dx$ for all $X\in\mathfrak S_V(\mathbb R^3)$; it is real-analytic in $x$, symmetric under simultaneous permutations of points and values, invariant under $\mathsf s_n(\mathsf rx;\rho(\mathsf R)e)=\mathsf s_n(x;e)$ for every isometry $\mathsf r$, and $|\mathsf s_n(x;e)|\le(\mathsf c\,C_V)^n\max\{1,2n^2/\delta(x)\}^{pn}\prod_k|e_k|$, $\delta(x)$ the least distance of two points, $\mathsf c$ a pure number. The Schwinger functions satisfy the growth hypothesis in Proposition~\YMMref{93.5}{719} with the constants $(\lceil p\rceil,1,\mathsf C_E,2p)$, so at $p=9$: $|\mathsf S_n(\varphi)|\le\mathsf C_E^n(n!)^{18}|\varphi|_{9n}$.
\item[(C)] By the reconstruction theorem among the construction results in its labelled form (Theorem~\textup{\YMMref{AO}{778}} for members with values in $V$, with the changes listed at the end of Subsection~\ref{cov-start}), applied to $\sigma^1$, which by (A) has the data (Rec1)--(Rec4) so read: the span $\mathcal D_0$ of $\Omega$ and of the vectors $T(f_1)\cdots T(f_n)\Omega$ of the reconstructed fields is dense in $\mathcal H(\sigma)$; there is a strongly continuous unitary representation $U$ of the proper orthochronous Poincar\'e group of $\mathbb R^{1+2}$ on $\mathcal H(\sigma)$ with $U(x,1)=e^{itH}e^{i\vec x\cdot P}$, $U\Omega=\Omega$, spectrum in the closed forward cone, and $U(a,\Lambda)T(f)U(a,\Lambda)^{-1}=T(\hat\Lambda f(\Lambda^{-1}(\cdot-a))\hat\Lambda^{\rm T})$ on $\mathcal D_0$, $\hat\Lambda=\iota\Lambda\iota^{-1}\in\mathrm{SO}(3,\mathbb C)$, $\iota(t,\vec x)=(\vec x,-it)$; and $T(f)T(f')\psi=T(f')T(f)\psi$ for $\psi\in\mathcal D_0$ and test tensors $f,f'$ with spacelike separated supports, all six components with all six.
\end{enumerate}
\end{corollary}

\subsection{Construction estimates}\label{cov-start}

Besides Theorems~\YMMref{D}{8}, \YMMref{AJ}{673} and~\YMMref{AD}{607}, the argument uses the statements among the construction results listed below, and, for Corollary~\ref{cov-cor-real}\,(C), the reconstruction theorem among the construction results in the labelled form stated at the end of this subsection.

\emph{The chain of renormalisation steps on one torus:} the setting (Th) (Section~\ref{sec:cutoff-calibration}, Lemma~\YMMref{118.1}{812}); the cuts $\mathsf{Cut}_k$ (the definitions preceding Lemma~\YMMref{113.10}{796}, Lemma~\YMMref{113.10}{796}) and the classes of covariance, which give the radius of regularity of the backgrounds over the cuts (Section~\ref{sec:sf-class-window} and Lemma~\ref{lem:sf-class-window}); one step as an identity at every point of the next cut (Lemma~\ref{lem:sf-chart-integral}); the one-step theorem with a complex plaquette weight of small size and small oscillation (Proposition~\ref{prop:source-entry}); the flows from a level for every system of early members (Proposition~\YMMref{114.2}{801}, Proposition~\YMMref{114.2}{801}), with the running numbers of equations~\YMMref{(2176)--(2177)}{803}; the Lipschitz form of forgetting (Proposition~\YMMref{114.3}{803}); parity in the fluctuation field (Proposition~\ref{prop:source-entry-parity}); the orbital potential (Section~\ref{sec:ward-orbital-complement}, Proposition~\ref{prop:ward-orbital-cubic}); the laws of the block fields at two consecutive cutoffs (Corollary~\textup{\ref{IV-62-66}}).

\emph{Two consecutive cutoffs on one torus:} the setting of Section~\YMMref{113}{786} with its threshold $\Theta_\star$, the data of the comparison (the setting preceding Lemma~\YMMref{113.8}{794}), the lower and the upper halves (Lemma~\YMMref{113.8}{794}), the comparison of the background actions at the cutoffs $K$ and $K+1$ (Proposition~\YMMref{113.9}{795}), and the competitors and their excess (Proposition~\YMMref{113.5}{792}, Lemma~\YMMref{113.6}{792}, Proposition~\YMMref{113.7}{793}); the setting (E), the two depths with their conditions (k0)--(k6), and the block fields with their laws (Definitions~\ref{sec:cutoff-calibration}, \ref{sec:cutoff-calibration}, \textup{\ref{IV-62-62}}); the rarity of large-field regions on one torus (Proposition~\ref{prop:ward-fixed-volume-rarity}).

\emph{One small torus:} the data, the numbers, the parameter sets and the step map of the weighted progression (Section~\ref{sec:weighted-progression-interface} and equations~\YMMref{(2251)--(2252)}{825}), its height zero (the height-zero part of Proposition~\ref{prop:weighted-finite-system}), the small-field progression with sources, the weighted progression of Proposition~\textup{\YMMref{122.1}{826}}, and its first derivatives at regular configurations (equation~\YMMref{(2264)}{828}); the setting of one lattice and the objects of the construction in it (Definition~\ref{def:weighted-history-setting}, Proposition~\ref{prop:weighted-canonical-history}), one exact step and the first step (Lemma~\ref{lem:weighted-array-step}), the term of the history without large-field regions (Proposition~\YMMref{123.1}{829}), the quotients of the sources along it (equation~\YMMref{(2267)}{829} and Proposition~\YMMref{123.1}{829}), the quotients of a step and their tower (the quotient definitions in Section~\ref{sec:weighted-covariance-interface} and equation~\YMMref{(2273)}{830}), and the one-step covariance of two sources (Lemma~\ref{lem:weighted-step-covariance}); the sum over the level of the first large-field region (equation~\eqref{eq:ward-first-creation}) and the rarity of large fields on a small side (Lemma~\YMMref{125.3}{841}); the stopped datum (Corollary~\YMMref{74.3}{616} and Lemma~\ref{lem:setup-stopped-arithmetic}); the small sides (equation~\YMMref{(2312)}{841}, Lemma~\YMMref{126.3}{846}); the depth fitted to a scale (equation~\YMMref{(2321)}{843} and the argument immediately following it), with the value of the room constant (equation~\YMMref{(2320)}{843}); the working level and its arithmetic (equation~\YMMref{(2322)}{844}, Lemma~\YMMref{126.1}{844}).

\emph{The first order and the two-point law on a small torus:} a pointed family inserted to first order in one step and the families carried from it to the higher levels (equations~\YMMref{(2276)--(2277)}{831}, \YMMref{(2280)}{831} and Proposition~\YMMref{123.3}{832}); one level at a regular configuration (Lemma~\ref{lem:weighted-regular-family}); the flat top (Proposition~\YMMref{123.3}{832}, equation~\YMMref{(2283)}{832}); the plaquette angles at a small point of the chart and the plaquette terms to the second order (Section~\YMMref{124}{832}, Lemma~\YMMref{124.4}{835}); the step couplings (equation~\YMMref{(2299)}{837}); the Gaussian numbers of the slices of the covariance and the errors of the flat top (equations~\YMMref{(2303)--(2305)}{838}); the carried family at the flat top (Proposition~\YMMref{125.1}{838}); the value of the truncated two-point function, with its errors, its sizes and its tails (Theorem~\YMMref{AT}{839}, Lemmas~\YMMref{125.2}{840}, \YMMref{125.4}{841}, \YMMref{125.5}{842}), at the levels attached to the working level (equation~\YMMref{(2316)}{842}), in the settings that the construction satisfies (Lemma~\YMMref{126.2}{845}); the envelopes and the error functions of the depth (equations~\YMMref{(2328)--(2336)}{846}, with the notation in \eqref{eq:short-distance-envelope-names}, the depth estimates in equations~\YMMref{(2332)}{846} and~\YMMref{(2336)}{847} and their proofs); the two-point law on a torus of small side (Theorem~\YMMref{AU}{847}); the background kernel of a depth and the slices of the free curvature kernel, with their telescope and their decay (Section~\YMMref{120}{817}, \YMMref{120.5}{821} and~\YMMref{120.6}{821}, Lemma~\YMMref{120.1}{818}, Propositions~\YMMref{120.5}{821} and~\YMMref{120.6}{821}).

\emph{The source zero:} the canonical system of localised source energies (Lemma~\ref{lem:canonical-source-system}), which gives the gap of the pressures between sides.

\emph{The lattice:} the lattice calculus of fields of directions (Lemma~\ref{lem:ward-general-calculus}, the general field \eqref{eq:ward-general-displacement}), the lattice Ward identity with its cubic density (Proposition~\YMMref{112.1}{782}) and the summation by parts of the orbital remainder (Proposition~\ref{prop:ward-orbital-cubic}); local gauges on cubes, with the H\"older regularity of backgrounds (Theorem~\textup{\ref{I-56-65}}); the covariance of lattice objects under lattice isometries (Lemma~\ref{lem:setup-lattice-covariance}).

\emph{Periodic Gibbs states and tori:} (Per) of Subsection~\ref{cov-set}; the chessboard estimate (Proposition~\YMMref{86.3}{681}, Corollary~\ref{IV-18-18}); multiple reflections (Proposition~\YMMref{104.2}{759}); the lattice energy-momentum inequality (Proposition~\YMMref{99.3}{739}) and its five-point form (Lemma~\ref{lem:spectral-steps}); the admissible labels on the tori of $\mathfrak L(g)$ (Lemma~\ref{lem:cutoff-admissible-labels}); the limit functionals along periodic Gibbs states (Theorem~\YMMref{AN}{763}) and the rotation theorem for the action density (Theorem~\YMMref{AK}{737}); the statement on two points for scalar plates (the proof of Theorem~\YMMref{AK}{737}).

\emph{Functionals on separated tuples:} the test spaces $\mathcal S_{\ne}$ and their norms (Section~\YMMref{93}{713}), the extension proposition (Proposition~\YMMref{93.3}{717}, with Lemma~\YMMref{93.2}{716} for uniqueness), the growth hypothesis in Proposition~\YMMref{93.5}{719}, and the reconstruction theorem for Euclidean data (Theorem~\textup{\YMMref{AO}{778}}), with the data (Rec) of Section~\YMMref{107}{765} and the Wightman properties stated in Theorem~\YMMref{AO}{778}.

\emph{The reconstruction theorem in its labelled form.} This is Theorem~\textup{\YMMref{AO}{778}} for members with values in $V$, with the following changes. In the data (Rec): the isometries of $\mathbb R^3$, the reflection $\theta$ among them, act on a member also through $\rho$ on its values, $(\mathsf rH)(x)=\rho(\mathsf R)H(\mathsf r^{-1}x)$, in (Rec1) and in the positivity of (Rec2); the norm $\|h\|_{L^1}$ of (Rec4) is formed with a norm of $V$ that $\rho$ keeps, for instance $(\operatorname{tr}A^2)^{1/2}\ge|A|_\infty$; and the data are (Rec1)--(Rec4), so read. In the Wightman clauses: the test functions $f$ take their values in $V_{\mathbb C}$, and $f^*=\rho(\theta)\bar f$ takes the place of $\bar f$ in the Hermiticity identity; a Lorentz transformation $\Lambda$ acts on the values through $D(\Lambda)A=\hat\Lambda A\hat\Lambda^{\rm T}$, $\hat\Lambda=\iota\Lambda\iota^{-1}\in\mathrm{SO}(3,\mathbb C)$, $\iota(t,\vec x)=(\vec x,-it)$, a homomorphism of the proper orthochronous Lorentz group into $\mathrm{GL}(V_{\mathbb C})$ with $D(1\oplus\mathsf R)=\rho(\mathsf R\oplus1)$ at a rotation $\mathsf R$ of the space coordinates, $D(B(s))=\exp(is\,d\rho(\omega_1))$ at the boost $B(s)$ in the plane of $x^0,x^1$, $\omega_1=e_1e_3^{\rm T}-e_3e_1^{\rm T}$, and $D(\Lambda)e^*=(D(\Lambda)e)^*$ for $e^*=\rho(\theta)\bar e$; Lorentz covariance reads $U(a,\Lambda)T(f)U(a,\Lambda)^{-1}=T(D(\Lambda)f(\Lambda^{-1}(\cdot-a)))$; locality and cyclicity are required for all the components; the unitary representation and spectrum conditions are retained, and uniqueness of the vacuum is imposed only when stated separately.

The statements among the construction results on the point operators, the functions $\mathbf W_n$, the kernels, the one holomorphic function on the union of the tubes, the field operators, their Lorentz covariance and their locality (Propositions~\YMMref{107.2}{767}, \YMMref{108.2}{770}, \YMMref{108.3}{771}, \YMMref{109.1}{773}, \YMMref{110.1}{776}, \YMMref{110.2}{777}) hold in the labelled form with the same proofs: these use of a member only its linearity and the action of the isometries on it, and the labels enter through $\rho$ and $D$ alone.

\subsection{Bounds on the lattice}\label{cov-lat}

\begin{lemma}\label{cov-cell}
Let $0<g\le\gamma$, $(K,\beta)\in\mathfrak d(g)$, and $G_g=C_{\rm S}+C_g+2C_{\rm tt}C'_{\rm tt}E_0$, formed with the numbers of the canonical system of Lemma~\ref{lem:canonical-source-system} at the source zero.
\begin{enumerate}
\item[(a)] $|p_{\ell_1}-p_{\ell_2}|\le G_g$ for all $\ell_1,\ell_2\in\mathfrak L(g)$.
\item[(b)] $\|F\|_{{\rm ch},\ell}\le e^{G_g\mathsf l^3}\|F\|_{{\rm ch},\ell'}$ for every function $F$ of a cell of side $\mathsf l$ with $2\mathsf l\le\ell'$ and every $\ell\in\mathfrak L(g)$ with $\ell\ge\ell'$.
\item[(c)] For $\nu\in\mathcal G^{\rm per}_{K,\beta}$, a board with cells of side $\ell'/2$ and any offset in $a_K\mathbb Z^3$, distinct cells $\square_1,\dots,\square_q$, real functions $F_k$ of $\square_k$ and real numbers $c_k$: $|\nu(\prod_k(F_k-c_k))|\le\prod_ke^{G_g(\ell'/2)^3}\|F_k-c_k\|_{{\rm ch},\ell'}$; in particular $|\nu(F_1)-c_1|\le e^{G_g(\ell'/2)^3}\|F_1-c_1\|_{{\rm ch},\ell'}$.
\item[(d)] \emph{(Multiple reflections at several cells of one layer.)} Let a board of side $\mathsf l$ have the wall plane $\{x_3=\gamma'\}$, let $\square_k$, $k\in\mathcal K_0$, be finitely many cells in the layer $\{\gamma'\le x_3\le\gamma'+\mathsf l\}$, $D_k$ local functions of $\square_k$, $D_k^\sharp$ the conjugate of the reflection of $D_k$ in the middle plane of the layer, with $\|D_k^\sharp\|_{{\rm ch},\ell_k}\le d_k$ along a sequence $\ell_k\to\infty$ in $\mathfrak L(g)$. Then for $\nu\in\mathcal G^{\rm per}_{K,\beta}$, local functions $X,X'$ of $\{x_3>\gamma'\}$ and $t,t'\in a_K\{0,1,2,\dots\}$,
\[\Bigl|\nu\Bigl(\overline{\theta_{\gamma'}\tau_{te_3}X'}\cdot\prod_{k\in\mathcal K_0}D_k\cdot\tau_{(\mathsf l+t')e_3}X\Bigr)\Bigr|\le\prod_kd_k\;\|X'\|_{\nu,\gamma'}\,\|X\|_{\nu,\gamma'},\qquad\|X\|^2_{\nu,\gamma'}=\nu\bigl(\overline{\theta_{\gamma'}X}\,X\bigr).\]
\end{enumerate}
\end{lemma}
\begin{proof}
Lemma~\ref{lem:canonical-source-system} at the source zero gives $p_\ell=\ell^{-3}E_{K,\ell,\beta}+\vartheta_\ell$ with $-C_g\le\vartheta_\ell\le C_{\rm S}$, and $|\ell_1^{-3}E_{K,\ell_1,\beta}-\ell_2^{-3}E_{K,\ell_2,\beta}|\le2C_{\rm tt}C'_{\rm tt}E_0$: this is (a). The comparison of two sides in Corollary~\ref{IV-18-18}\,(e), $\|F\|_{{\rm ch},\ell}\le\|F\|_{{\rm ch},\ell'}e^{\mathsf l^3(p_{\ell'}-p_\ell)}$ for $\ell\ge\ell'\ge2\mathsf l$, gives (b). The domination in (Per), along $\ell_k\to\infty$, and (b) give (c); a constant has the seminorm equal to its modulus. For (d) one follows the proof of Proposition~\YMMref{104.2}{759} (multiple reflections): Schwarz' inequality at the wall planes $\theta_{\gamma'}$ and $\theta_{\gamma'+\mathsf l}$ separates $X'$ and $X$ from the layer, the convexity step is taken before the quotient so that the product over all the cells of the layer is reflected at once, and the chessboard estimate in the layer bounds the reflected product by $\prod_k\|D_k^\sharp\|_{{\rm ch},\ell_k}$.
\end{proof}

Every estimate below is of short distance; Lemma~\ref{cov-cell} makes it uniform in the side and in the state.

\emph{Sources under majorants.} The small-field progression with sources carries weights of the plaquettes (\eqref{eq:weighted-progression-recursion}); its proofs read a source only through the properties that the next definition names. Let data (Wp1)--(Wp5) of Definition~\ref{def:weighted-progression-data} be given on a torus $T_\ell$ at the spacing $a_K$, with a last height $k\ge1$. $\mathcal P(X)$ is the set of the plaquettes with their four corners in a region $X$; the regions, the triples of the cores, the pointed families, the families without base point, the families in the fourth place, the bare arrays and their sizes are those that Section~\ref{sec:weighted-progression-interface} and Section~\YMMref{122}{825} name.

\begin{definition}\label{cov-maj}
\begin{enumerate}
\item[(1)] A \emph{majorant} is a pair $(\mathsf s,S)$ of a number $\mathsf s>0$ and a non-empty set $S$ of plaquettes. Its \emph{weight} is the real source $\mathsf s\mathbf 1_S$ on the plaquettes, and the numbers $\mathsf m_j=12M^3\alpha_0^2L^{K-j}\mathsf s$ of \eqref{eq:weighted-progression-numbers} are formed at it.
\item[(2)] A \emph{source under the majorant} $(\mathsf s,S)$ is a triple $\mathfrak s=(F,\mathbb V^\natural,(H^{(j)})_{1\le j\le k})$ with the following properties.
\begin{itemize}
\item[(So0)] The \emph{first member} $F$ is a continuous complex function of the configurations, invariant under the gauge transformations of the bare lattice, with $|F(U)|\le a_K^{-1}\mathsf s\sum_{p\in S}a_p(U)$ at every $\mathrm{SU}(2)$-valued configuration $U$.
\item[(So1)] (\emph{The height zero.}) $\mathbb V^\natural$ is a family in the fourth place at the height $0$ (Section~\ref{sec:step-input-conventions}), of size at most $\mathsf c^\flat\mathsf m_0$, $\mathsf c^\flat=81C_P^2\mathsf C_{\mathcal A}^2/K_\Lambda$ the number of \eqref{eq:weighted-progression-numbers}; its member at a region $Y$ vanishes unless $\mathcal P(Y)$ meets $S$; and it sums to $F$, in the sense of \eqref{eq:fourth-family-global-sum}, at every configuration $W$ of the first level satisfying the background and chart conditions of Section~\ref{sec:step-input-conventions} at the height $0$, as the bracketed action of a weight $w$ of the plaquettes sums to $\mathcal O(w)$ (the height-zero sum rule in Proposition~\ref{prop:weighted-finite-system}).
\item[(So2)] (\emph{The heights $j\ge1$.}) $H^{(j)}$ is a family without base point of level $j$ at the radii $(\alpha_0,\alpha_1)$, of size at most $\frac12\mathsf m_j$; its member at a region $X$ vanishes unless $\mathcal P(X)$ meets $S$; and $\sum_X\bigl(H^{(j)}(X,U)-H^{(j)}(X,\mathbf 1)\bigr)=F(U)$ wherever the action of the bare arrays of the height $j$ is defined. It takes the place of the weight member in the slot of the bare array of the height $j$ (\eqref{eq:weighted-progression-recursion}).
\end{itemize}
The source is \emph{real} if its first member is real.
\item[(3)] A \emph{system of sources} is a tuple of sources $\mathfrak s_a=(F_a,\mathbb V^\natural_a,(H^{(j)}_a)_j)$ under majorants $(\mathsf s_a,S_a)$, $1\le a\le p$, $p\ge1$. All the numbers and sets of equations~\eqref{eq:weighted-progression-numbers}--\eqref{eq:weighted-progression-domain}, among them $\mathsf m_{j,a}$, $\mathsf d_{j,a}$, $\mathsf n_{j,a}$, $N^{[j]}_\star$ and the parameter sets $\Lambda^{\rm w}_j$, and the numbers of Sections~\ref{sec:weighted-covariance-interface} and~\YMMref{123--124}{828} and~\YMMref{125--126}{837} that are formed with the sources, are formed at the weights $\mathsf s_a\mathbf 1_{S_a}$. For $\lambda\in\mathbb C^p$ put $F_\lambda=\sum_a\lambda_aF_a$. The \emph{progression} $(E'^{(j)}_\lambda)_j$ of the system is the weighted progression of \eqref{eq:weighted-progression-recursion}, built with the same provisions, numbers of a step and counterterms, with $\sum_a\lambda_a\mathbb V^\natural_a$ in the place of the bracketed action at the height $0$ and with $\sum_a\lambda_aH^{(j)}_a$ in the place of the weight member in the slot of the height $j$; $E^{(j)}=E'^{(j)}_0$ is the progression without sources.
\item[(4)] The system is \emph{admissible} if the following holds.
\begin{itemize}
\item[(Ad1)] For every $j$ and every $\lambda\in\Lambda^{\rm w}_{j-1}$ the family $E'^{(j)}_\lambda$ is a pointed family of level $j$ at $(\alpha_0,\alpha_1)$; each of its values is holomorphic in $\lambda$ and in the fields; and $\|E'^{(j)}_\lambda-E'^{(j)}_{\lambda'}\|_{j,(1+4\beta_{\rm s})\kappa}\le\sum_a|\lambda_a-\lambda'_a|\,\mathsf d_{j,a}$.
\item[(Ad2)] The member of $E'^{(j)}_\lambda$ at a region $X$ does not depend on $\lambda_a$ unless $\mathcal P(X)$ meets $S_a$.
\item[(Ad3)] The progression is given by the step map: $E'^{(j+1)}_\lambda=E^{(j+1)}+\mathsf{St}^{[j]}(x^{[j]},\mathfrak f'_{\lambda,j})-\mathsf{St}^{[j]}(x^{[j]},\mathfrak f'_{0,j})$, where $\mathfrak f'_{\lambda,j}$ is the bare array of the families $E'^{(i)}_\lambda$, $i\le j$, with the counterterms of the regular progression and with the member $\sum_a\lambda_aH^{(j)}_a$ in its slot.
\item[(Ad4)] Along the history without large-field regions up to the level $t$, the term of $\rho_0e^{F_\lambda}$ is $\mathsf r'^\lambda_t\cdot\mathsf m_t$, $\mathsf r'^\lambda_t$ the exponential of the relative exponent in equation~\YMMref{(2265)}{829} formed with $F_\lambda$ and the $E'^{(i)}_\lambda$ in the places of $\mathcal O(w_\lambda)$ and the $E^{(i)}_\lambda$, as in Proposition~\YMMref{123.1}{829}.
\item[(Ad5)] The irrelevant part $\mathrm{Irr}^{(t)}_a$ of the first derivative (the definitions following equation~\YMMref{(2267)}{829}, with $\dot E'^{(i)}_a=\partial_{\lambda_a}E'^{(i)}_\lambda|_{\lambda=0}$) obeys the bound of equation~\YMMref{(2264)}{828} formed at the weight $\mathsf s_a\mathbf 1_{S_a}$.
\end{itemize}
The \emph{constants of the progression at the level $t$} are $\mathsf c^{(t)}_a=\sum_{i\le t}\sum_{(X,z)\in\mathcal Q_i}\dot E'^{(i)}_a(X,\mathbf 1,z)$, the sums over the regions and sites $\mathcal Q_i$ of the height $i$ (the finite index sets defined in Section~\ref{sec:weighted-covariance-interface}).
\item[(5)] A \emph{kernel under the majorant} $(\mathsf s,S)$ is a real symmetric function $\mathsf t$ on the pairs of plaquettes that vanishes unless both plaquettes belong to $S$, with $2\sum_{p'}|\mathsf t(p,p')|\le\mathsf s$ for every $p$. A real source $\mathfrak s$ under $(\mathsf s,S)$ together with a kernel $\mathsf t$ under $(\mathsf s,S)$ is a \emph{quadratic source} if
\begin{itemize}
\item[(So3)] at every height $i$, every real point $(v,\mathcal A)$ of the chart of equations~\YMMref{(2285)--(2286)}{833} and all real numbers $r\ge\|v\|_{T_\ell}$, $a\ge\sup|\mathcal A|$ with $\mathsf{rt}=\min\{\varrho_{\rm ch}/r,\varrho_{\rm fl}/a\}\ge2$, as in Lemma~\YMMref{124.4}{835},
\[\Bigl|a_KF\bigl(e^{i\eta_iA^{\rm lw}(v,\mathcal A)}\bigr)-\sum_{p,p'}\mathsf t(p,p')\bigl\langle\mathsf{ang}_p(v,\mathcal A),\mathsf{ang}_{p'}(v,\mathcal A)\bigr\rangle\Bigr|\le\mathsf M_{\rm pl}^2\,\eta_i^4\,\mathsf{rt}^{-3}\,(\#S)\,\mathsf s\,:\]
the second order of the first member at a small point of the chart is the quadratic form of the kernel in the plaquette angles of Section~\YMMref{124}{832}.
\end{itemize}
\end{enumerate}
\end{definition}

\begin{lemma}\label{cov-src}
Let $g'\le\gamma^{\rm ref}_W$ be a reference coupling, $(K,\beta)\in\mathfrak d(g')$, and $T_\ell$ a torus of a side $\ell\in\mathfrak L(g')$. For a complex function $\mathsf w$ of the bonds, not identically zero, let $S[\mathsf w]$ be the set of the plaquettes that contain a bond $b$ with $\mathsf w(b)\ne0$.
\begin{enumerate}
\item[(a)] For all data (Wp1)--(Wp5) on $T_\ell$ with a last height $k\ge1$, the data of the weighted progression with at least one level (Proposition~\textup{\YMMref{122.1}{826}}, Proposition~\YMMref{123.1}{829}), the bond form $Y[\mathsf w]$ is the first member of a source $\mathfrak s[\mathsf w]$ under the majorant $(59|\mathsf w|_\infty,S[\mathsf w])$, complex linear in $\mathsf w$ and real if $\mathsf w$ is real; in particular $|Y[\mathsf w](U)|\le2a_K^{-1}|\mathsf w|_\infty\sum_{p\in S[\mathsf w]}a_p(U)$ at every $\mathrm{SU}(2)$-valued configuration $U$. For a complex weight $w$ of the plaquettes, not identically zero, $\mathcal O(w)$ is the first member of a source under $(|w|_\infty,\operatorname{supp}w)$, with the bracketed action at the height $0$ and the weight members in the slots. A source under a majorant $(\mathsf s,S)$ is one under every $(\mathsf s',S')$ with $\mathsf s'\ge\mathsf s$, $S'\supset S$; sums of sources, member by member, are sources under $(\mathsf s+\mathsf s',S\cup S')$, and, for $c\ne0$, $c\,\mathfrak s$ is one under $(|c|\mathsf s,S)$. So $\mathcal T_K(H)=\mathcal O(w_H)+Y[\mathsf w_H]$ is the first member of a source under the majorant $(60|H|_\infty,S_H)$, $S_H=\operatorname{supp}w_H\cup S[\mathsf w_H]$, for every bounded $H:\mathbb R^3\to V_{\mathbb C}$ with $S_H\ne\varnothing$.
\item[(b)] In the setting of one lattice of Definition~\ref{def:weighted-history-setting}, which the construction on $T_\ell$ along the run of $(K,\beta)$ satisfies (Proposition~\ref{prop:weighted-canonical-history}, Lemma~\YMMref{126.2}{845}), every system of sources under majorants (Definition~\ref{cov-maj}\,(3)), of any number and with complex first members, is admissible (Definition~\ref{cov-maj}\,(4)). So every finite tuple of forms $Y[\mathsf w_a]$, $\mathcal O(w_a)$, $\mathcal T_K(H_a)$, and of sums and multiples of them, is the tuple of the first members of an admissible system of sources under the majorants of (a).
\item[(c)] Let $\mathsf w$ be real. For a bond $b$ of the direction $\mu$ and a direction $\lambda\ne\mu$, the linear part of $\mathbf u_{b,\lambda}$ in the angles $\Phi_p$ of the plaquettes, $\Phi$ a function of the plaquettes with values in the hermitian traceless $2\times2$ matrices, is $\mathsf u_{b,\lambda}(\Phi)=\sum_pc_{b,\lambda}(p)\Phi_p$, with $c_{b,\lambda}(p)=\pm\frac12$ at the two plaquettes of the plane of $e_\mu,e_\lambda$ that contain $b$, the signs being those of their orientations, and $c_{b,\lambda}(p)=0$ at the other plaquettes. Put
\[\mathsf t[\mathsf w](p,p')=-\tfrac12\sum_b\mathsf w(b)\bigl[c_{b,\lambda_b}(p)\,c_{b,\lambda'_b}(p')+c_{b,\lambda_b}(p')\,c_{b,\lambda'_b}(p)\bigr],\]
so that $\sum_{p,p'}\mathsf t[\mathsf w](p,p')\langle\Phi_p,\Phi_{p'}\rangle=-\sum_b\mathsf w(b)\langle\mathsf u_{b,\lambda_b}(\Phi),\mathsf u_{b,\lambda'_b}(\Phi)\rangle$. Then $\mathsf t[\mathsf w]$ is real, symmetric and carried by $S[\mathsf w]$, $|\mathsf t[\mathsf w](p,p')|\le\frac18|\mathsf w|_\infty$, $2\sum_{p'}|\mathsf t[\mathsf w](p,p')|\le2|\mathsf w|_\infty$ and $\sum_{p,p'}|\mathsf t[\mathsf w](p,p')|\le\sum_b|\mathsf w(b)|$; and, for the standing data of Definition~\ref{def:weighted-history-setting} and the actual step setting of Section~\YMMref{124}{832}, $\mathfrak s[\mathsf w]$ with the kernel $\mathsf t[\mathsf w]$ is a quadratic source under $(59|\mathsf w|_\infty,S[\mathsf w])$. For a real weight $w$ of the plaquettes the source of $\mathcal O(w)$ is quadratic under $(|w|_\infty,\operatorname{supp}w)$ with the kernel $\frac12w(p)\delta_{pp'}$. Sums and real multiples of quadratic sources are quadratic sources; so for a real $H$ the source of $\mathcal T_K(H)$ is quadratic under $(60|H|_\infty,S_H)$ with the kernel $\mathsf t_H(p,p')=\mathsf t[\mathsf w_H](p,p')+\frac12w_H(p)\delta_{pp'}$, and $2\sum_{p'}|\mathsf t_H(p,p')|\le3|H|_\infty$.
\end{enumerate}
\end{lemma}
\begin{proof}
(a) \emph{(So0).} At $\mathrm{SU}(2)$-valued fields $|w(W)|^2=2a_p-a_p^2\le2a_p$ for the loop $W$ of a plaquette $p$, so $|\mathbf u_{b,\lambda}|^2\le a_p+a_{p'}$ for the two plaquettes $p,p'$ of the plane of $e_{\mu_b},e_\lambda$ at $b$, and $q_b$ is real with $|q_b|\le\frac12(|\mathbf u_{b,\lambda_b}|^2+|\mathbf u_{b,\lambda'_b}|^2)\le\frac12\sum_{p\ni b}a_p$. Summing over the four bonds of each plaquette, $|Y[\mathsf w](U)|\le a_K^{-1}\sum_pa_p(U)\,\frac12\sum_{b\subset\partial p}|\mathsf w(b)|\le2a_K^{-1}|\mathsf w|_\infty\sum_{p\in S[\mathsf w]}a_p(U)$, and $2\le59$; $Y[\mathsf w]$ is a gauge invariant polynomial in the entries of the bond variables. \emph{(So1).} For a bond $b$ let $X_0(b)$ be the union of the cubes of the partition of the height $0$ that hold a corner of one of the four plaquettes at $b$: at most six cubes with a common corner, a region of the height $0$ whose set $\mathcal P(X_0(b))$ holds these four plaquettes; one region is $X_0(b)$ for at most $18M^3$ bonds.

The family at the height $0$ is the variation of the bond form along the fluctuation of the first step, $\mathbb V^\natural(Y;(P,B))=-a_K^{-1}\sum_{b:X_0(b)=Y}\mathsf w(b)\,[q_b(e^{iB'}\mathbf U)-q_b(\mathbf U)]$, formed at a triple $P$ of the core and a scaled fluctuation $B$ with the configuration $\mathbf U$ and the $1$-form $B'$ with which the bracketed action of a weight of the plaquettes is formed in the height-zero proof of Proposition~\YMMref{122.1}{826}; it is holomorphic, gauge covariant and vanishes at $B=0$ as that action does, $q_b$ being a gauge invariant polynomial in the variables of the four plaquettes at $b$. At the triples of the core, with $\rho\le2C_P\mathsf C_{\mathcal A}\mathsf s^{[0]}\alpha_0\le\frac1{128}$ the size of $B'$ and $\zeta_0=\alpha_0L^{-2}$ that of the plaquette variables of the triple, the variation of $q_b$ is at most $784\rho^2+132\rho\zeta_0\le3138\,C_P^2\mathsf C_{\mathcal A}^2\mathsf s^{[0]}\alpha_0^2$, with $\mathsf s^{[0]}\le1$ of \eqref{eq:weighted-progression-numbers} and the numbers $C_P$, $\mathsf C_{\mathcal A}$ of $\mathsf c^\flat$; so the size of $\mathbb V^\natural$ is at most $18\cdot3138\,M^3C_P^2\mathsf C_{\mathcal A}^2\alpha_0^2a_K^{-1}|\mathsf w|_\infty/K_\Lambda$, while (So1) allows $\mathsf c^\flat\mathsf m_0=81\cdot12\,M^3C_P^2\mathsf C_{\mathcal A}^2\alpha_0^2a_K^{-1}\mathsf s/K_\Lambda$ at a majorant $(\mathsf s,S)$.

As $18\cdot3138=56484<57348=81\cdot12\cdot59$, the majorant $(59|\mathsf w|_\infty,S[\mathsf w])$ serves: whence the number $59$. Each bond occurs at the one region $X_0(b)$, so the members of $\mathbb V^\natural$ at the restrictions of $(P_W,gB)$ add up to the variation $Y[\mathsf w](U_0(V_W(B)))-Y[\mathsf w](U_1(W))$, $q_b$ being constant on the gauge orbits: this is the sum rule \eqref{eq:fourth-family-global-sum}. \emph{(So2).} With $X_j(b)$ formed alike from the cubes of the partition of the height $j$, the slot member is $H^{(j)}(X,P)=-a_K^{-1}\sum_{b:X_j(b)=X}\mathsf w(b)\,q_b(\mathbf U)$. At the triples of the core the bond variables have norm at most $e^{\eta_j\alpha_1}$ and the plaquette variables lie within $\alpha_0\eta_j^2$ of $\mathbf 1$, so that $|q_b|\le4e^{4\eta_j\alpha_1}\alpha_0^2\eta_j^4$ there; with at most $18M^3\eta_j^{-3}$ bonds at one region, the size of $H^{(j)}$ at the radii $(\alpha_0,\alpha_1)$ is at most $72e^{4/13}M^3\alpha_0^2L^{K-j}|\mathsf w|_\infty<98M^3\alpha_0^2L^{K-j}|\mathsf w|_\infty$, while $\frac12\mathsf m_j=354M^3\alpha_0^2L^{K-j}|\mathsf w|_\infty$ at $\mathsf s=59|\mathsf w|_\infty$; its sum is $Y[\mathsf w]$, as $q_b(\mathbf 1)=0$.

For $\mathcal O(w)$ the three members are those of \eqref{eq:weighted-progression-recursion}, with the properties of the height-zero part of Proposition~\ref{prop:weighted-finite-system}. A larger majorant loosens each of (So0)--(So2), the numbers $\mathsf m_j$ growing with $\mathsf s$. Sums and multiples: (So0)--(So2) are linear conditions with subadditive sizes, $\mathsf s\mathbf 1_S+\mathsf s'\mathbf 1_{S'}\le(\mathsf s+\mathsf s')\mathbf 1_{S\cup S'}$, and $|w_H|_\infty+59|\mathsf w_H|_\infty\le60|H|_\infty$.

(b) The proofs of Proposition~\textup{\YMMref{122.1}{826}} (with the height-zero part of Proposition~\ref{prop:weighted-finite-system}), of Proposition~\YMMref{123.1}{829} and of the first-derivative bounds in equations~\YMMref{(2264)}{828} and~\YMMref{(2268)}{829} read a source only through (So0)--(So2) and through numbers formed at its weight: at the height $0$ through the family in the fourth place, its size, its locality and its sum rule; at the heights $j\ge1$ through the member in the slot of the bare array, its size, its locality and its action; and through the domination (So0), which bounds the source at the backgrounds of small fields as the action density of its weight is bounded there. The statements that these proofs cite besides are worded for every family in the fourth place, every bare array of the domain of the step map (equations~\YMMref{(2251)--(2252)}{825}), every bare array of $\mathsf{Arr}^{[j]}$ (Lemma~\ref{lem:weighted-array-step}) and every bounded complex Borel function. Read with a system of sources in the place of the weights, these proofs give the progression of the system on the parameter sets $\Lambda^{\rm w}_j$ with (Ad1) and (Ad3), the locality (Ad2) from that of the members, (Ad4) from the first step and from one exact step at each height (Lemma~\ref{lem:weighted-array-step}), and (Ad5) as in equation~\YMMref{(2264)}{828}. None of these readings depends on the number of the sources or on their reality.

(c) \emph{The coefficients.} For $b=\langle x,x+a_Ke_\mu\rangle$, $\mathbf u_{b,\lambda}=-\frac12[w(W_x(\mu,\lambda))-w(W_x(\mu,-\lambda))]$, and to the first order $w(W_x(\mu,\pm\lambda))$ is the angle of the plaquette of the loop with the sign of its orientation; this gives $c_{b,\lambda}$. \emph{The kernel.} Two different plaquettes have at most one bond in common; so $\mathsf t[\mathsf w](p,p')\ne0$ only if $p$ and $p'$ contain one bond $b$ with $\mathsf w(b)\ne0$ and lie in the two different planes through $b$, and then just one of the two products in the bracket is not zero: $|\mathsf t[\mathsf w](p,p')|=\frac18|\mathsf w(b)|$. A plaquette has four bonds, at each of them two partners: $2\sum_{p'}|\mathsf t[\mathsf w](p,p')|\le2\cdot8\cdot\frac18|\mathsf w|_\infty=2|\mathsf w|_\infty\le59|\mathsf w|_\infty$; a bond gives eight ordered pairs: $\sum_{p,p'}|\mathsf t[\mathsf w](p,p')|\le\sum_b|\mathsf w(b)|$.

\emph{(So3).} $a_KY[\mathsf w]=-\sum_b\mathsf w(b)q_b$. At a small real point of the chart the vector parts of the plaquette loops at $b$ are, up to terms of the third order, the plaquette angles with the signs of the orientations, as for $a_p$ in Lemma~\YMMref{124.4}{835}; this gives $|q_b(e^{i\eta_iA^{\rm lw}(v,\mathcal A)})-\langle\mathsf u_{b,\lambda_b}(\mathsf{ang}),\mathsf u_{b,\lambda'_b}(\mathsf{ang})\rangle|\le12\,\mathsf M_{\rm pl}^2\eta_i^4\mathsf{rt}^{-3}$, $\mathsf{ang}=\mathsf{ang}(v,\mathcal A)$. Times $|\mathsf w(b)|$ and summed: as every bond lies in four plaquettes and every plaquette has four bonds, at most $\#S[\mathsf w]$ bonds carry $\mathsf w$, and the remainder is at most $\mathsf M_{\rm pl}^2\eta_i^4\mathsf{rt}^{-3}(\#S[\mathsf w])\,12|\mathsf w|_\infty$, with $12\le59$.

For $\mathcal O(w)$, (So3) is the first inequality of Lemma~\YMMref{124.4}{835}, times $|w(p)|$, summed over the support. Sums and real multiples: the conditions on kernels are linear with subadditive bounds, and $(\#S)\mathsf s+(\#S')\mathsf s'\le\#(S\cup S')(\mathsf s+\mathsf s')$; finally $2\sum_{p'}|\mathsf t_H(p,p')|\le|H|_\infty+2|H|_\infty$.
\end{proof}

\begin{lemma}\label{cov-mom}
Let $g'$ and $(K,\beta)$ be as in Lemma~\ref{cov-src}, with $K\ge3$, let $T_\ell$ be a torus whose side is small over the run of $(K,\beta)$, and let $s$ be a depth with $1\le s\le K-2$; $\mathsf k=K-s\in[2,K-1]$ is the working level. Let $F_1,\dots,F_n$, $1\le n\le\ell^3$, be the first members of an admissible system of sources under majorants $(\mathsf s_a,S_a)$ (Definition~\ref{cov-maj}) such that the corners of the plaquettes of each $S_a$ lie in the interior of one unit cell of a board of $T_\ell$ of side $1$, and, if $n\ge2$, any two of the sets $S_a$ are at least $M\,u_*(n)$ units of the level $\mathsf k$ apart, where
\[u_*(m)=\max\Bigl\{4,\ \frac2{A_\flat}\bigl(m\,[c_2+\log(4m)]+3\log L+\log3\bigr)\Bigr\}\qquad(m\ge1),\]
with numbers $A_\flat\in(0,1]$ and $c_2$ of the constants, an increasing function of the order $m\log m$. Then, with the constants $c_a=\mathsf c^{(\mathsf k)}_a$ of the progression at the level $\mathsf k$ (Definition~\ref{cov-maj}\,(4)), for every non-empty $B\subset\{1,\dots,n\}$
\[\Bigl|\Bigl\langle\prod_{a\in B}(F_a-c_a)\Bigr\rangle_{K,\ell,\beta}\Bigr|\le\prod_{a\in B}2n\,\mathsf C_1L^{4s}\,\mathsf s_a ,\]
$\mathsf C_1$ a number of the constants that depends on none of $g'$, $K$, $\beta$, $\ell$, $n$, $s$ and the sources.
\end{lemma}
\begin{proof}
Let $r=\#B$. The mean of the product $G_B=\prod_{a\in B}(F_a-c_a)$ is a sum over the level $t$ of the first creation of a large field (equation~\eqref{eq:ward-first-creation}): $|\langle G_B\rangle|\le\mathsf q_{\mathsf k}+\sum_{1\le t\le\mathsf k}\mathsf q_{t-1}\mathsf w_t$, where $\mathsf q_0\ge\sup|G_B|$, $\mathsf q_t$ bounds the density of $G_B$ at the level $t$ along the histories without large fields, and $\mathsf w_t$ is the weight of a first creation at $t$. Along the histories without large fields the density is $\sum_{\pi}\prod_{\beta\in\pi}\mathsf D^\beta_t$ over the partitions $\pi$ of $B$, by (Ad4) and the formula $\partial_Be^\Psi=e^\Psi\sum_\pi\prod_{\beta\in\pi}\partial_\beta\Psi$: the blocks of one element are the conditional means $\mathsf c^{(t)}_a-c_a+F_a(U_t)+\mathrm{Irr}^{(t)}_a$ of the progression, and those of several elements are the sums over the regions and sites of the heights $i\le t$ of the mixed derivative families $\partial_\beta E'^{(i)}_\lambda|_{\lambda=0}$, which live on polymers joining the supports.

For these, put $c_\theta=\theta^{[j]}N^{[j]}_\star$, a number free of $j$, and $\hat g_{0,a}=\mathsf c^\flat\mathsf m_{0,a}/N^{[0]}_\star$, $\hat g_{j,a}=\mathsf n_{j,a}/N^{[j]}_\star$ for $j\ge1$: the size of the source $a$ in the array of the height $j$ over the room $N^{[j]}_\star$ of the domain of the step map (\eqref{eq:weighted-progression-numbers}). The chain rule of all orders for the step map (equations~\YMMref{(2251)--(2252)}{825} give it for two variables; its proof, which reads the holomorphy and the Lipschitz bound of the step map on its domain, gives it for any number of variables), with Cauchy's estimate on the polydiscs of the radii $1/(4r\hat g_{j,a})$ in the variables $\lambda_a$, bounds the size of the mixed derivative family of a set $\beta$ of $r\ge2$ sources at the height $j$ by $c_\theta(4r)^r\prod_{a\in\beta}\hat g_{j-1,a}$, times $e^{-4\beta_{\rm s}\kappa[\mathsf D_j/M-2]_+}$ in the distance $\mathsf D_j$ of the supports in units of the level $j$; these are exponentially small in the distance, and the separation $M\,u_*(n)$ makes them pay the combinatorial factors.

The sizes are $\sup|G_B|\le\prod\mathsf C_1L^{4K}\mathsf s_a$ and, at the level $t$, $\operatorname{Bl}(r)\prod\mathsf C_1L^{4(K-t)}\mathsf s_a$, $\operatorname{Bl}(r)\le r^r$ the number of partitions; for the blocks of several elements this uses $L^{\frac52(K-t)+1}\le L^{4(K-t)}$, that is $K-t\ge K-\mathsf k=s\ge1$. On a small side the rarity of large fields (Lemma~\YMMref{125.3}{841}) gives $\sum_tL^{4(K-t+1)r}\mathsf w_t\le1$ for every $r\le\ell^3$: the large fields pay for as many sources as the side has unit cells.
\end{proof}

For $0<\mathsf d<1$ and a reference coupling $g'$ the \emph{depth fitted to $\mathsf d$} is the witness depth $s(\mathsf d)$ of equation~\YMMref{(2321)}{843}, the least integer $s\ge0$ with $L^s\mathsf d\ge C_W(s)$, where $C_W(s)=\mathsf c_WY_s^{2\mathfrak p_0}$, $Y_s=\log\bar x_s$, $\bar x_s=L^{s+1}(g'^{-2}+\bar B)$ and $\mathsf c_W=C_{W,0}A_0^2$, formed with the value $C_{W,0}=\max\{1,4L^2M/A_0^2\}$ of the room constant of equation~\YMMref{(2320)}{843}, which serves throughout the section, so that $\mathsf c_W\ge4L^2M$; then $L^{s(\mathsf d)}\mathsf d<LC_W(s(\mathsf d))$ by the minimality argument following equation~\YMMref{(2321)}{843}. The next lemma bounds moments of eight sources at their natural size; unlike Lemma~\ref{cov-mom}, it measures the sizes at the working level $\mathsf k=\max\{K-s,2\}$ of equation~\YMMref{(2322)}{844}, and not by a factor $L^{4(K-t)}$ per source and level.

\begin{lemma}\label{cov-mom8}
Let $0<g'\le\gamma^{\rm sd}_W$, $(K,\beta)\in\mathfrak d(g')$, $b=a_K\beta$, $T_\ell$ the least torus at $g'$, and $0<\mathsf d'<1$ with $s=s(\mathsf d')\le K$, the depth fitted to $\mathsf d'$ at $g'$, and so small that $2Y_s^{2\mathfrak p_0}\ge u_*(8)$. Let $\mathsf k=\max\{K-s,2\}$ be the working level and $s_{\mathsf k}=K-\mathsf k=\min\{s,K-2\}$ its depth (equation~\YMMref{(2322)}{844}), so that $s\ge3$, $2\le\mathsf k\le K-1$ and $s-2\le s_{\mathsf k}\le s$ (Lemma~\YMMref{126.1}{844}), and put
\[\upsilon=\frac{L^{s_{\mathsf k}}}{x_{\mathsf k}(\beta)}=\bigl(a_{s_{\mathsf k}}x_{\mathsf k}(\beta)\bigr)^{-1},\]
the step coupling $\upsilon_{\mathsf k+1}$ of equation~\YMMref{(2299)}{837} (Lemma~\YMMref{126.1}{844}, equation~\YMMref{(2324)}{844}): the inverse of the coupling of the run at the working level, read in the units of that level. Let $F_1,\dots,F_8$ be the first members of an admissible system of sources under majorants $(\mathsf s_a,S_a)$ such that the barycentres of the plaquettes of each $S_a$ lie in a ball of radius $\frac78\mathsf d'$ and barycentres of plaquettes of two different sets $S_a$ are at least $\frac54\mathsf d'$ apart. Then, with the constants $c_a=\mathsf c^{(\mathsf k)}_a$ of the progression at the level $\mathsf k$, for every non-empty $B\subset\{1,\dots,8\}$
\[\Bigl|\Bigl\langle\prod_{a\in B}(F_a-c_a)\Bigr\rangle_{K,\ell,\beta}\Bigr|\le\prod_{a\in B}\mathsf A_8(s)\,\upsilon\,\mathsf s_a,\qquad\mathsf A_n(s)=2^{2\mathfrak p_0+1}n\max\{\mathsf C^{\rm I},C_2n^{3+\mathfrak p_0}\}\bigl(LC_W(s)\bigr)^3\bar x_sY_s^{2\mathfrak p_0},\]
where $\mathsf C^{\rm I}$, $C_2$ are numbers of the constants; $\upsilon$ is the natural size of a source of unit majorant at the working level, and
\[|b\,\upsilon-1|\le\bar Bg'^2L^{2-s}\le\tfrac1{24},\qquad\upsilon^2\le\tfrac98b^{-2},\qquad|b^2\upsilon^2-1|\le3\bar Bg'^2L^{2-s}.\]
In particular $\mathsf A_8(s(\mathsf d'))\le C\mathsf d'^{-1}\mathsf Y(\mathsf d')^{10\mathfrak p_0}$, with a number $C$ of the constants and $g'$ and a function $\mathsf Y$ that is affine in $\log\frac1{\mathsf d'}$.
\end{lemma}
\begin{proof}
As in the proof of Lemma~\ref{cov-mom}, the mean of $G_B=\prod_{a\in B}(F_a-c_a)$ is a sum over the level $t$ of the first creation of a large field (equation~\eqref{eq:ward-first-creation}), and along the histories without large fields the density at the level $t$ is a sum over the partitions of $B$, the blocks of one element being the conditional means of the progression and those of several elements the terms of its exponent on polymers joining the supports. \emph{The distance.} Two corners of plaquettes of different sets $S_a$ are at a sup distance $\ge\frac5{4\sqrt3}\mathsf d'-a_K$, so at least $\mathsf D_{\mathsf k}\ge\frac5{4\sqrt3}\hat{\mathsf d}'_{\mathsf k}-1$ units of the level $\mathsf k$, $\hat{\mathsf d}'_{\mathsf k}=L^{s_{\mathsf k}}\mathsf d'$.

By Lemma~\YMMref{126.1}{844}, equation~\YMMref{(2323)}{844}, which reads the property $\mathsf c_W\ge4L^2M$ of the room constant, $4MY_s^{2\mathfrak p_0}\le L^{-2}C_W(s)\le\hat{\mathsf d}'_{\mathsf k}<LC_W(s)$; so $\mathsf D_{\mathsf k}\ge2.88\,MY_s^{2\mathfrak p_0}-1\ge2MY_s^{2\mathfrak p_0}\ge M\,u_*(8)$, and at a level $j<\mathsf k$ the distance is $L^{\mathsf k-j}$ times larger: the terms on polymers joining the supports pay the combinatorial factors. \emph{The sizes.} A source is now bounded by its natural size: at the working level it is a form in the field strength over a ball of the radius $\frac78\hat{\mathsf d}'_{\mathsf k}<LC_W(s)$ units of the level, at the coupling $x_{\mathsf k}(\beta)\le\bar x_s$ of the run (Lemma~\YMMref{126.1}{844}, equation~\YMMref{(2324)}{844}), so that each block is bounded at that level by a power of $\max\{\mathsf C^{\rm I},C_2n^{3+\mathfrak p_0}\}(LC_W(s))^3\bar x_sY_s^{2\mathfrak p_0}$ times $\prod\upsilon\mathsf s_a$, and at a level $t<\mathsf k$ by $L^{4(\mathsf k-t)}(\mathsf k-t+2)^{2\mathfrak p_0}$ times more.

On the least side, which is small over the run (Lemma~\YMMref{126.3}{846}), the rarity of large fields (Lemma~\YMMref{125.3}{841}) pays these factors, and the sum over $t$ gives $\mathsf A_8(s)$. \emph{The number $\upsilon$.} By the flow bound of Subsection~\ref{cov-set}, $|x_{\mathsf k}(\beta)-L^{-\mathsf k}\beta|\le\bar B$; multiplied by $a_{s_{\mathsf k}}$, as $a_{s_{\mathsf k}}L^{-\mathsf k}\beta=a_K\beta=b$, this is $|\upsilon^{-1}-b|\le\bar Ba_{s_{\mathsf k}}\le\bar BL^{2-s}$; and $x_{\mathsf k}(\beta)>L^{s_{\mathsf k}}g'^{-2}$ gives $\upsilon<g'^2$ (Lemma~\YMMref{126.1}{844}, equation~\YMMref{(2324)}{844}). So $|b\upsilon-1|\le\bar Bg'^2L^{2-s}\le\frac{13}{24L}\le\frac1{24}$, as $\bar Bg'^2\le\frac{13}{24}$ and $s\ge3$; hence $(b\upsilon)^2\le(\frac{25}{24})^2\le\frac98$ and $|(b\upsilon)^2-1|\le3\bar Bg'^2L^{2-s}$.

For the last assertion: $L^s\mathsf d'<LC_W(s)$ gives $\bar x_s<L^2\mathsf c_WY_s^{2\mathfrak p_0}(g'^{-2}+\bar B)/\mathsf d'$, and so $Y_s<\log(L^2\mathsf c_W/\mathsf d')+2\mathfrak p_0\log Y_s+\log(g'^{-2}+\bar B)$; with $2\mathfrak p_0\log Y_s\le\frac12Y_s+2\mathfrak p_0\log(4\mathfrak p_0)$ this bounds $Y_s$ by a function $\mathsf Y(\mathsf d')$ affine in $\log\frac1{\mathsf d'}$, and $\mathsf A_8(s)$ is at most a number times $Y_s^{6\mathfrak p_0}\bar x_sY_s^{2\mathfrak p_0}\le C\mathsf d'^{-1}\mathsf Y(\mathsf d')^{10\mathfrak p_0}$.
\end{proof}

\begin{proposition}\label{cov-sd}
Let $0<g\le\gamma^{\rm sd}_W$ and $0<\delta\le\frac14$. Put $\bar x_0=L(g^{-2}+\bar B)$, $c_E=\log(2\bar x_0)$, $\mathsf Y_i=2[c_E+\log\frac1{a_i}]$, $q_E=15r+4$, $r$ the exponent of the cube-size function $R$ of Section~\YMMref{74}{614}, $R(x)$ being the least power $L^j$, $j\ge0$, with $L^j\ge(\log x)^r$, $c_N=\max\{1,8L^6M^3C_{\rm s}^32^{-3r}\}$, $\bar N_i=c_N\mathsf Y_i^{3r}$, $s_i$ the least $s$ with $\delta L^s\ge M\,u_*(\lceil\bar N_i\rceil)$, $K_E(i)=i+s_i+2$, and $C_E=480\,e^{G_g+1/8}c_N\mathsf C_1(LMc_u/\delta)^4\bar x_0$ with $C_{\rm s}$ the number of Section~\YMMref{74}{614}, $\mathsf C_1$ that of Lemma~\ref{cov-mom}, and a number $c_u$ of the constants with $u_*(\lceil\bar N_i\rceil)\le c_u\mathsf Y_i^{3r+1}$. For every integer $i\ge1$, every $(K,\beta)\in\mathfrak d(g)$ with $K\ge K_E(i)$, every board of $\mathbb L_K$ with cells of side $a_i$, every cell $\square$, every side $\ell\in\mathfrak L(g)$, every bounded $H:\mathbb R^3\to V_{\mathbb C}$ that vanishes at the points at sup distance $<\delta a_i$ from the complement of $\square$, and for $c$ the mean of $\mathcal T_K(H)$ in a $\nu\in\mathcal G^{\rm per}_{K,\beta}$ or on a torus $T_{\ell_1}$, $\ell_1\in\mathfrak L(g)$,
\begin{equation}\label{cov-eq-sd}
\|\mathcal T_K(H)-c\|_{{\rm ch},\ell}\le\mathsf c_E(i)\,b^{-1}|H|_\infty,\qquad\mathsf c_E(i)=2C_E\,\mathsf Y_i^{\,q_E}a_i^{-11/2}.
\end{equation}
In particular $\mathsf c_E(i)\le C_\flat a_i^{-6}$ with $C_\flat=2C_E\sup_i\mathsf Y_i^{q_E}L^{-i/2}<\infty$.
\end{proposition}
\begin{proof}
\emph{Pieces.} It suffices to treat four pieces of $H$: the real and the imaginary parts of the diagonal entries, which enter through $\mathcal O$, and of the off-diagonal ones, which enter through $Y$.
\emph{Dilation.} The dilation $x\mapsto L^ix$ carries $\mathbb L_K$ onto $\mathbb L_{K'}$, $K'=K-i$, the board onto a board with cells of side $1$, the Wilson measures at $(K,\beta)$ onto those at $(K',\beta)$, the seminorms at the side $\ell$ onto those at the side $L^i\ell$, and $\mathcal T_K(H)$ onto $L^i\mathcal T_{K'}(H(L^{-i}\cdot))$.
\emph{The stopped datum.} $(K',\beta)$ is a first-passage datum at the stopped coupling $g_{(i)}<g$, with label $L^ib$; its least side $\ell''$ is small over the run, $L^i\ell\in\mathfrak L(g_{(i)})$, and $N=\ell''^3\le c_N\mathsf Y_i^{3r}$.
\emph{The copies.} On $T_{\ell''}$ the $N$ mirror copies of a piece are the first members of an admissible system of $N$ sources (Lemma~\ref{cov-src}\,(a), (b)), each carried by a whole unit cell with margin $\delta$; two points in different cells, each at distance $\ge\delta$ from the walls, are $\ge2\delta$ apart, so the corners of the plaquettes of different copies are $\ge2(\delta-2a_{K'})\ge\delta$ apart, that is $\ge\delta L^s\ge M\,u_*(N)$ in units of the level $K'-s$, $s=s_i$; here $1\le s\le K'-2$, as $\delta<4\le Mu_*(\lceil\bar N_i\rceil)$ and $K\ge K_E(i)$. By Lemma~\ref{cov-mom}, $|\langle\prod_{a\in B}(F_a-c_a)\rangle_{K',\ell'',\beta}|\le(2N\mathsf C_1L^{4s}\mathsf s)^{\#B}$ for every set $B$ of copies, $\mathsf s$ the weight of the piece. With $m$ the common mean of the copies, $\prod_\kappa(F_\kappa-m)=\sum_B\prod_{\kappa\in B}(F_\kappa-c_\kappa)\prod_{\kappa\notin B}(c_\kappa-m)$ and $|c_\kappa-m|\le2N\mathsf C_1L^{4s}\mathsf s$; so the seminorm on $T_{\ell''}$ of a piece less its mean is $\le4N\mathsf C_1L^{4s}\mathsf s$.
\emph{The pressure gap.} Lemma~\ref{cov-cell}\,(b) at $g_{(i)}$, with the bound \eqref{eq:setup-stopped-pressure}, $C_{g_{(i)}}\le C_g+\frac92i\log L+\frac18$, carries the seminorm of a function of one unit cell from $T_{\ell''}$ to any side $L^i\ell$ at the price $e^{G_g+1/8}L^{9i/2}$.
\emph{Collecting.} The dilation contributes $L^i$ and the pressure gap $L^{9i/2}$, together $a_i^{-11/2}$; $N$ and $L^{4s_i}$ are powers of $\mathsf Y_i$. So, with $m'=L^im$, $m$ the common mean of the copies of the dilated piece on $T_{\ell''}$, $\|F-m'\|_{{\rm ch},\ell}$ is bounded, for a piece $F$ and every side $\ell\in\mathfrak L(g)$, by the product of these factors, a bound uniform in the side. The change of the centring to the mean $c$ of the piece in $\nu$ or on $T_{\ell_1}$ costs a factor $2$, the means of the four pieces adding up to that of $\mathcal T_K(H)$: by the domination in (Per) at the one cell, along the sides $\ell_j\to\infty$ of $\mathfrak L(g)$, $|\nu(F)-m'|\le\liminf_j\|F-m'\|_{{\rm ch},\ell_j}$, and by the chessboard estimate on $T_{\ell_1}$, $|\langle F\rangle_{K,\ell_1,\beta}-m'|\le\|F-m'\|_{{\rm ch},\ell_1}$ (Corollary~\ref{IV-18-18}\,(c)); both are at most the bound just obtained, and, the seminorm of a constant being its modulus, $\|F-c\|_{{\rm ch},\ell}\le\|F-m'\|_{{\rm ch},\ell}+|m'-c|$ is at most twice that bound. Each member stands once, in a cell of its own: no tensor source is put into an exponent. The power $\frac{11}2=1+\frac92$ comes from the dilation and from the gap of the pressures; the natural power is $0$.
\end{proof}

\begin{remark}\label{cov-rem-stab}
The stability clause of Theorems~\YMMref{AJ}{673} and~\YMMref{AD}{607} bounds the sourced expectation after removing a holomorphic normalising factor, the sourced vacuum constant. It does not provide a cutoff-uniform bound for the raw centred exponential. The free-photon analogue illustrates the distinction: on the orthogonal complement of the kernel of $d$, a constant diagonal source $H=c\mathbf 1$, $0<c<b$, has centred logarithmic generating function $3(N^3-1)[-\log(1-c/b)-c/b]$, $N=\ell/a_K$. This grows like $N^3$ for fixed $c/b$.
\end{remark}

\begin{proposition}\label{cov-lim}
Let $0<g\le\gamma^{\rm sd}_W$.
\begin{enumerate}
\item[(a)] \emph{(The uniform bound.)} For $n\ge1$ and $\mathsf d,R>0$ let $i$ be the least integer $\ge1$ with $2a_i\le\mathsf d$, $K^G$ the least integer $\ge K_E(i)$ with $L^{K^G-i}\ge8n$, and $C^G_n(\mathsf d;R)=[(16nR/a_i+3)^3\mathsf c_E(i)]^n$, both at $\delta=\frac1{8n}$. For all $(K,\beta)\in\mathfrak d(g)$ with $K\ge K^G$, all $\nu\in\mathcal G^{\rm per}_{K,\beta}$ and all $X$ of type $(\mathsf d;R)$ with $n$ members, $|\nu(A^\nu_K(X))|\le C^G_n(\mathsf d;R)\,b^{-n}\prod_a|H_a|_\infty$; the same holds with every member centred at its mean on a torus of a side in $\mathfrak L(g)$. The number $C^G_n(\mathsf d;R)$ depends on none of the cutoff, the datum, the side, the state and the positions; for small $\mathsf d$ it is a power of $1/\mathsf d$ up to logarithms.
\item[(b)] Every sequence $(K_j,\beta_j)\in\mathfrak d(g)$, $K_j\to\infty$, $\rho_j\in\mathcal G^{\rm per}_{K_j,\beta_j}$, has a subsequence along which $a_{K_j}\beta_j$ converges and a tensor limit exists.
\item[(c)] Every tensor limit $\sigma$ at $g$ is labelled data, with $C_n(\mathsf d;R)=C^G_n(\mathsf d;R)b_*^{-n}$. \emph{(Continuity in the members.)} If $X_j\to X$ member by member in the sup norm, with supports in one compact set and at mutual distances $\ge\mathsf d>0$, then $\rho_j(A^{\rho_j}_{K_j}(X_j))\to\sigma(X)$. The scalar part $\sigma_0$, taken at members $f\mathbf 1$, $f\in C^1_c$, is the centred limit functional of the action density along $(K_j,\beta_j,\rho_j)$ of Theorem~\YMMref{AN}{763}, invariant under $\mathbb R^3\rtimes\mathrm O(3)$.
\item[(d)] \emph{(Slab operators.)} $\sigma$ has the property (S). For $i\ge1$, $z_3=\frac12a_i$, $h\in C_c(\{|x-z|_\infty\le\frac38a_i\};V_{\mathbb C})$ and $F,F'\in\mathbb C\mathfrak S_V(\mathbb E_+)$,
\begin{equation}\label{cov-eq-basic}
|\sigma(\theta\bar F'\cdot(h)\cdot\tau_{a_ie_3}F)|\le\mathsf c_E(i)\,b_*^{-1}|h|_\infty\|[F']\|\,\|[F]\|,
\end{equation}
with $\mathsf c_E(i)$ at $\delta=\frac1{16}$; with members $h_1,\dots,h_m$ in the place of $h$, in the cubes $o+a_i\mathbf k_a+[-\frac38a_i,\frac38a_i]^3$, $o=(t_1,t_2,\frac12a_i)$, $\mathbf k_a\in\mathbb Z^3$ pairwise different with third entries $0$, the right side becomes $\prod_a\mathsf c_E(i)b_*^{-1}|h_a|_\infty\|[F']\|\|[F]\|$. If $\operatorname{supp}h$ lies in a ball of radius $R$ and $a_i$, $i\ge1$ least, is at most the distance of $\operatorname{supp}h$ from the planes $x_3=0$ and $x_3=l$, then $\|\mathsf B_l(h)\|\le((16/3)R/a_i+3)^3\mathsf c_E(i)b_*^{-1}|h|_\infty$. Further $e^{-tH}\mathsf B_l(h)e^{-t'H}=\mathsf B_{l+t+t'}(\tau_{te_3}h)$, $e^{iv\cdot P}\mathsf B_l(h)e^{-iv\cdot P}=\mathsf B_l(\tau_{(v,0)}h)$, $\mathsf B_l(h)^*=\mathsf B_l(h^\star)$ with $h^\star(x)=\rho(\theta)\overline{h(x_1,x_2,l-x_3)}$, the products $\mathsf B_{l_1}(h_1)\cdots\mathsf B_{l_n}(h_n)\Omega$ are the classes of the tuples whose members lie in pairwise disjoint slabs, the eight elements of $W_3$ that fix $e_3$ act by unitary operators $[F]\mapsto[\mathsf RF]$ commuting with $H$, and the closed span $\mathcal H_0$ of the classes of the scalar tuples is $\mathcal H(\sigma_0)$ and reduces $H$, $P$ and $\mathcal P$.
\item[(e)] \emph{(The insertion bound with a power, in all six directions.)} Put $\mathsf a=\frac1{4L}$, $c_{\rm I}=\frac38\mathsf a$. For $\eta\in(0,1]$, $h\in C_c(\mathbb R^3;V_{\mathbb C})$ in a closed ball $\bar B(y,c_{\rm I}\eta)$ with $|y_3|\le\frac\eta8$, and $F,G\in\mathbb C\mathfrak S_V(\{x_3>\frac\eta4\})$,
\[|\sigma(\theta\bar F\cdot(h)\cdot G)|\le2C_E\mathsf a^{-11/2}\Bigl(2c_E+2\log\frac1{\mathsf a\eta}\Bigr)^{q_E}g^2\,\eta^{-11/2}|h|_\infty\,\|[\tau_{-(\eta/4)e_3}F]\|\,\|[\tau_{-(\eta/4)e_3}G]\| ;\]
here $C_E$ is taken at $\delta=\frac1{16}$; so the left side is $\le C_{\rm I}\eta^{-\mathsf p_{\rm I}}|h|_\infty\|[\tau_{-(\eta/4)e_3}F]\|\|[\tau_{-(\eta/4)e_3}G]\|$ for every $\mathsf p_{\rm I}>\frac{11}2$, with $C_{\rm I}$ a function of the constants, $g$ and $\mathsf p_{\rm I}$ that depends neither on $\sigma$ nor on its label.
\item[(f)] \emph{(Towers of one member.)} Put $\mathsf a'=\frac1{16L^2}$ and $\mathsf n_*=(32L^2+3)^3$. For $\eta'\in(0,\frac12]$, $h$ with support in the closed ball of radius $\frac{\eta'}{16}$ about a point $z$ with $z_3\in[\frac{\eta'}4,\frac{3\eta'}4]$, with $h^\star$ as in (d) at $l=\eta'$, the tower $\mathbb T_k=(h^\star)\cdot\tau_{\eta'e_3}(h)\cdot\tau_{2\eta'e_3}(h^\star)\cdots\tau_{(2k-1)\eta'e_3}(h)$ of $2k$ floors, and the reflection $\vartheta_c$ in $x_1=c=z_1+\frac{\eta'}8$,
\[0\le\sigma\bigl(\vartheta_c\overline{\mathbb T_k}\cdot\mathbb T_k\bigr)\le\Bigl(\mathsf n_*\,2C_E\Bigl(2c_E+2\log\frac1{\mathsf a'\eta'}\Bigr)^{q_E}(\mathsf a'\eta')^{-11/2}\,g^2\,|h|_\infty\Bigr)^{4k},\]
with $C_E$ at $\delta=\frac1{16}$; so the left side is at most $(C_{\mathsf p}\eta'^{-\mathsf p}|h|_\infty)^{4k}$ for every $\mathsf p>\frac{11}2$.
\end{enumerate}
\end{proposition}
\begin{proof}
(a) Put $r=a_i/(8n)$ and $\lambda=L^{K-i}=a_i/a_K$, an odd integer $\ge8n$ as $K\ge K^G$. Cut each member $H_a$ by the tent partition of unity of the scale $r$, $H_a=\sum_{z\in r\mathbb Z^3}\chi^r_zH_a$, the tent $\chi^r_z$ being carried by the cube $\{|x-z|_\infty<r\}$ of half-side $r$; as the support of $H_a$ lies in a closed ball of radius $R$, at most $(2R/r+3)^3=(16nR/a_i+3)^3$ of the pieces are not zero. Fix one piece of each member, carried by the cube about $z_a$. In each direction $\mu$ the wall planes of the boards with cells of side $a_i$ lie at one of the $\lambda$ positions $a_K(m+\frac12)+a_i\mathbb Z$, $0\le m<\lambda$; a piece excludes the positions that have a wall plane in the closed interval $[z_{a,\mu}-2r,z_{a,\mu}+2r]$ of length $4r=a_i/(2n)$, at most $\lambda/(2n)+1$ of them, and the $n$ pieces exclude at most $\lambda/2+n<\lambda$ positions.

A free position in each direction gives one board on which the cube of every piece lies in the interior of one cell, at a sup distance $>r$ from its walls, so that the piece vanishes at the points at sup distance $<a_i/(8n)$ from the complement of its cell; different members lie in different cells because $\mathsf d\ge2a_i>\sqrt3a_i$. Domination of $\nu$ at several cells of one board (Per) and Proposition~\ref{cov-sd} at $\delta=\frac1{8n}$, uniform in the side, bound each term by $\prod_a\mathsf c_E(i)b^{-1}|H_a|_\infty$, and there are at most $(16nR/a_i+3)^{3n}$ terms; the centring on a torus is allowed by Proposition~\ref{cov-sd}. (b) The bound (a) gives equicontinuity on each type; extract diagonally along a countable dense set of tuples and along the labels, which lie in a compact interval.
(c) (D1) and the continuity in the members follow from (a) and multilinearity; $\sigma$ is real as a limit of real numbers. (D2): $\mathsf r\mathcal T_K(H)=\mathcal T_K(\mathsf rH)$ for lattice isometries, by the covariance of lattice objects (Lemma~\ref{lem:setup-lattice-covariance}), $\nu$ is invariant under $a_K\mathbb Z^3\rtimes W_3$ (Per), and the lattices are nested since $L$ is odd, so that a translation or a mirror of one cutoff is one of every finer cutoff; with the continuity in the members, all translations follow. (D3): reflection positivity of periodic Gibbs states at the coordinate and diagonal lattice mirrors (Per), which keep $\mathcal T_K$, and the continuity in the members. (D4): $e^{-tH}$ and $e^{iv\cdot P}$ exist by the iterated Schwarz inequality and (D1); the inequality
\[S_F(0)-4S_F(te_3)+S_F(2te_3)+S_F(te_3+te_k)+S_F(te_3-te_k)\ge0,\qquad S_F(w)=\sigma(\theta\bar F\cdot\tau_wF),\ k=1,2,\]
passes from the five-point inequality \textup{(\ref{IV-54-29})} of Lemma~\ref{lem:spectral-steps}, a form of the lattice energy-momentum inequality, to the limit and gives $\mathcal P(\{E<\max(|q_1|,|q_2|)\})=0$. The scalar part: $\mathcal T_K(f\mathbf 1)=\mathcal O(f)$, and Theorem~\YMMref{AN}{763} applies to the sequence $(K_j,\beta_j,\rho_j,\rho_j)$.
(d) Multiple reflections (Lemma~\ref{cov-cell}\,(d)) in the layer of cells of height $a_i$ above the bond mirror $x_3=\frac12a_{K_j}$ bound $\rho_j$ of the reflected $F'$, the centred cell observable and the translated $F$ by the chessboard seminorm of the cell observable, uniformly in the side, times the two mirror norms; Proposition~\ref{cov-sd} at $\delta=\frac1{16}$ (the support is $\ge\frac18a_i-\frac12a_{K_j}\ge\frac1{16}a_i$ from the walls once $a_{K_j}\le\frac18a_i$) and the continuity in the members give \eqref{cov-eq-basic}, and the same with several cells of one layer. A partition of unity at the one scale $a_i$, at most $((16/3)R/a_i+3)^3$ pieces, gives the bound of $\mathsf B_l(h)$, and the operator itself exists by F.~Riesz's theorem. The relations are identities between classes, which follow from (D2); the dihedral operators are unitary by (D2), and $\mathcal H_0$ is invariant under the semigroup, the translations and the dihedral group, which are generated on it by the scalar tuples.
(e) Let $i$ be the least integer $\ge1$ with $a_i\le\frac\eta4$; then $\mathsf a\eta\le a_i\le\frac\eta4$. Inside the slab of width $\frac\eta2$ about the ball there is a cell of the scale $a_i$ about the ball as in (d); two contractions $e^{-cH}$, $c\ge0$, bridge the remaining heights, and (d) gives the bound, with $b_*^{-1}<g^2$.
(f) Let $i_2$ be the least integer $\ge1$ with $a_{i_2}\le\frac{\eta'}{16}$, $i=i_2+1$ and $\rho=\frac1{16}a_i$. Then $a_{i_2-1}>\frac{\eta'}{16}$: for $i_2\ge2$ by the choice of $i_2$, and for $i_2=1$ because $a_0=1>\frac1{32}\ge\frac{\eta'}{16}$. So $a_i=L^{-2}a_{i_2-1}>\frac{\eta'}{16L^2}=\mathsf a'\eta'$ and $\frac{2\eta'}{a_i}<32L^2$; with $\mathsf Y_i\le2c_E+2\log\frac1{\mathsf a'\eta'}$ this bounds $\mathsf c_E(i)$, at $\delta=\frac1{16}$, by $2C_E(2c_E+2\log\frac1{\mathsf a'\eta'})^{q_E}(\mathsf a'\eta')^{-11/2}$, and $\mathsf c_E(i_2)\le\mathsf c_E(i)$. \emph{Claim:} for $u,v\in C_c(\mathbb R^3;V_{\mathbb C})$ with supports in $\bar B(\zeta,\frac{\eta'}{16})$ and $\bar B(\zeta+\frac{\eta'}4e_1,\frac{\eta'}{16})$, $\zeta_3\in[\frac{\eta'}4,\frac{3\eta'}4]$, and $F,F'\in\mathbb C\mathfrak S_V(\mathbb E_+)$,
\[|\sigma(\theta\bar F'\cdot(u)\cdot(v)\cdot\tau_{\eta'e_3}F)|\le\mathsf n_*^2\,\mathsf c_E(i)^2\,b_*^{-2}|u|_\infty|v|_\infty\,\|[F']\|\,\|[F]\| .\]
Cut $u$ and $v$ by tents of the scale $\rho$; each has at most $(\frac{2\eta'}{a_i}+3)^3\le\mathsf n_*$ pieces, centred at heights in $[\frac{3\eta'}{16}-\rho,\frac{13\eta'}{16}+\rho]$, and the pieces of $v$ lie at least $\frac{\eta'}8-2\rho$ further along $e_1$ than those of $u$. For a pair of pieces whose heights differ by at least $a_i$, \eqref{cov-eq-basic} at the scale $i$, applied twice, each time to one piece moved by a translation along $e_3$ into the layer of (d), with the heights left free above and below it at least $\frac12a_i$, bounds the pair by $\mathsf c_E(i)^2b_*^{-2}|u|_\infty|v|_\infty\|[F']\|\|[F]\|$. For a pair whose heights differ by less than $a_i$, the two pieces lie, after a translation along $e_3$, in a slab of half-width $\frac9{16}a_i\le\frac38a_{i_2}$; an offset in the first two coordinates of its own for this pair puts them in the inner cubes of two cells of one layer of the board of side $a_{i_2}$, distinct because $\frac{\eta'}8-2\rho>\frac34a_{i_2}$, with the heights left free above and below the layer at least $\frac12a_{i_2}$; the bound of (d) for several cells of one layer at the scale $i_2$ bounds the pair alike.

Summing over the at most $\mathsf n_*^2$ pairs gives the claim. \emph{The column.} Let $u_m=h^\star$ for even $m$ and $u_m=h$ for odd $m$, $0\le m<2k$, be the untranslated floor profiles, and let $v_m=\vartheta_c\bar u_m$; $\vartheta_c$ moves $z_1$ to $z_1+\frac{\eta'}4$ and commutes with the translations along $e_3$, and $|v_m|_\infty=|h|_\infty$. Put $Y_{2k}=\varnothing$ and $Y_m=(u_m)\cdot(v_m)\cdot\tau_{\eta'e_3}Y_{m+1}$, a tuple in $\mathbb E_+$. The claim at $F'=Y_m$, $F=Y_{m+1}$ gives $\|[Y_m]\|^2\le\mathsf M\|[Y_m]\|\,\|[Y_{m+1}]\|$ with $\mathsf M=\mathsf n_*^2\mathsf c_E(i)^2b_*^{-2}|h|_\infty^2$, so $\|[Y_0]\|\le\mathsf M^{2k}$; $Y_0$ is $\vartheta_c\overline{\mathbb T_k}\cdot\mathbb T_k$ in another order, and $|\sigma(Y_0)|=|\langle\Omega,[Y_0]\rangle|\le\mathsf M^{2k}$, with $b_*^{-1}\le g^2$.

As $\mathbb T_k$ lies in $\{x_1<c\}$, (D3) at the plane $x_1=c$ gives the sign.
\end{proof}

The bound on off-diagonal members rests on two lemmas: a decay of the background kernel, which fixes the number $C_{\mathsf k}$, and the factorisation of the moments of close pairs of bond forms that are far from one another.

\begin{lemma}\label{cov-bg}
Let $(\ell,K)$ be a lattice, $0\le s\le K$, and let $p,q$ be plaquettes with $\mathsf r=|x(p)-x(q)|_\infty>0$. The background kernel $C^{\rm bg}_{K,s}$ of the depth $s$ (Section~\YMMref{120}{817}) satisfies
\[\bigl|C^{\rm bg}_{K,s}(p,q)\bigr|\le\bigl[C_{\rm tp}\mathsf S_3\mathsf r^{-3}+C^{\rm bg}_{\rm tp}\bigr]a_K^3,\qquad\mathsf S_3=\sup_{t>0}\sum_{n\in\mathbb Z}(tL^n)^3e^{-\delta_{\rm tp}tL^n}<\infty,\]
with the numbers $C_{\rm tp}$, $C^{\rm bg}_{\rm tp}$, $\delta_{\rm tp}$ of Propositions~\YMMref{120.5}{821} and~\YMMref{120.6}{821}. So, for $0<\mathsf d<1$ and $|x(p)-x(q)|\ge\mathsf d$,
\[\bigl|C^{\rm bg}_{K,s}(p,q)\bigr|\le C_{\mathsf k}\Bigl(\frac{a_K}{\mathsf d}\Bigr)^3,\qquad C_{\mathsf k}=3^{3/2}C_{\rm tp}\mathsf S_3+C^{\rm bg}_{\rm tp},\]
a number of the constants.
\end{lemma}
\begin{proof}
By the telescope \YMMref{(2227)}{818} at the depths $s$ and $0$, $C^{\rm bg}_{K,s}=C^{\rm bg}_{K,0}+\sum_{K-s<j\le K}\mathcal C_{(j)}$. By Proposition~\YMMref{120.6}{821} at $n=K-j+1\in\{1,\dots,s\}$, $|\mathcal C_{(j)}(p,q)|\le C_{\rm tp}a_K^3L^{3n}e^{-\delta_{\rm tp}\mathsf rL^n}=C_{\rm tp}a_K^3\mathsf r^{-3}(\mathsf rL^n)^3e^{-\delta_{\rm tp}\mathsf rL^n}$, and by Proposition~\YMMref{120.5}{821} at the depth $0$, $|C^{\rm bg}_{K,0}(p,q)|\le C^{\rm bg}_{\rm tp}a_K^3$. The series of $\mathsf S_3$ does not change when $t$ is multiplied by $L$, and for $t\in[1,L]$ it converges uniformly, its terms being at most $(tL^n)^3$ for $n<0$ and decaying like $e^{-\delta_{\rm tp}L^n}$ for $n>0$. For $|x(p)-x(q)|\ge\mathsf d$, $\mathsf r\ge\mathsf d/\sqrt3$, and $a_K^3\le(a_K/\mathsf d)^3$ as $\mathsf d<1$.
\end{proof}

\begin{lemma}\label{cov-pairs}
Let $0<g\le\gamma^{\rm ref}_W$, $(K,\beta)\in\mathfrak d(g)$, $b=a_K\beta$, $T_\ell$ a torus whose side is small over the run of $(K,\beta)$, and $0<\mathsf d<1$ with $s=s(\mathsf d)\le K$; let $\mathsf k$, $s_{\mathsf k}$ and $\upsilon$ be as in Lemma~\ref{cov-mom8}, at $g$ and $\mathsf d$, and $C^{\rm bg}=C^{\rm bg}_{K,s_{\mathsf k}}$ the background kernel of the depth $s_{\mathsf k}$. Let $n\ge1$, and number the members of $n$ pairs $a=(\tau,\iota)$, $\tau=1,\dots,n$, $\iota=1,2$. For each member let $z_a\in T_\ell$ be a point and $\mathsf w_a$ a real weight of the bonds of one direction, not identically zero, with $|x_b-z_a|\le\mathsf d-2a_K$ where $\mathsf w_a(b)\ne0$, such that $3\mathsf d\le|z_{(\tau,1)}-z_{(\tau,2)}|\le4\mathsf d$ for each pair, and that corners of plaquettes of $S[\mathsf w_a]$ and of $S[\mathsf w_{a'}]$ are at a sup distance $\ge11M$ whenever $a$ and $a'$ belong to different pairs.

Put $F_a=Y[\mathsf w_a]$, the first member of the quadratic source under $(\mathsf s_a,S_a)=(59|\mathsf w_a|_\infty,S[\mathsf w_a])$ with the kernel $\mathsf t_a=\mathsf t[\mathsf w_a]$ (Lemma~\ref{cov-src}), $\hat m_a=\upsilon\mathsf s_a$, $P_\tau=F_{(\tau,1)}F_{(\tau,2)}$, $\hat m_\tau=\hat m_{(\tau,1)}\hat m_{(\tau,2)}$, and, for a non-empty set $J$ of pairs, $P_J=\prod_{\tau\in J}P_\tau$ and $\hat m_J=\prod_{\tau\in J}\hat m_\tau=59^{2\#J}\upsilon^{2\#J}\prod_{\tau\in J}|\mathsf w_{(\tau,1)}|_\infty|\mathsf w_{(\tau,2)}|_\infty$; let
\[G_\tau=6\,\upsilon^2\sum_{p,p',q,q'}\mathsf t_{(\tau,1)}(p,p')\,\mathsf t_{(\tau,2)}(q,q')\,C^{\rm bg}(p,q)\,C^{\rm bg}(p',q')\]
be the Gaussian value of the pair $\tau$. Then $|G_\tau|\le106\,C_{\mathsf k}^2\upsilon^2|\mathsf w_{(\tau,1)}|_\infty|\mathsf w_{(\tau,2)}|_\infty$, with the number $C_{\mathsf k}$ of Lemma~\ref{cov-bg}. And there are an integer $s_{2n}(g)$ and a function $\omega_{g,2n}$ of the depth, formed with the constants, $g$ and $n$, with $\omega_{g,2n}(s)\to0$ as $s\to\infty$, such that for $s_{2n}(g)\le s(\mathsf d)\le K$ and every non-empty set $J$ of pairs
\[\Bigl|\langle P_J\rangle_{K,\ell,\beta}-\prod_{\tau\in J}G_\tau\Bigr|\le\omega_{g,2n}(s(\mathsf d))\,\hat m_J .\]
\end{lemma}
\begin{proof}
\emph{The data.} At the $2n$ majorants the working level, the step couplings and the upper level $T$ attached to the working level (equations~\YMMref{(2299)}{837}, \YMMref{(2316)}{842} and~\YMMref{(2322)}{844}) are formed, and the settings of Sections~\YMMref{125--126}{837} hold for the construction on $T_\ell$ along the run of $(K,\beta)$ (Lemmas~\YMMref{126.1}{844}, \YMMref{126.2}{845}); the forms $F_a$ are the first members of an admissible system of quadratic sources (Lemma~\ref{cov-src}). \emph{The bound of $G_\tau$.} The barycentres of plaquettes of the two members of a pair are at least $3\mathsf d-2(\mathsf d-\frac32a_K)>\mathsf d$ apart, so $|C^{\rm bg}(p,q)|\le C_{\mathsf k}(a_K/\mathsf d)^3$ at the plaquettes that count (Lemma~\ref{cov-bg}); $\sum_{p,p'}|\mathsf t_a(p,p')|\le\sum_b|\mathsf w_a(b)|$ (Lemma~\ref{cov-src}\,(c)), and at most $\frac{4\pi}3(\mathsf d/a_K)^3$ bonds of one direction have their midpoints in a ball of radius $\mathsf d-2a_K$; so $|G_\tau|\le6(\frac{4\pi}3)^2C_{\mathsf k}^2\upsilon^2|\mathsf w_{(\tau,1)}|_\infty|\mathsf w_{(\tau,2)}|_\infty$, and $6(\frac{4\pi}3)^2=105.3<106$.

So $|G_\tau|\le\mathsf g_0\hat m_\tau$, $\mathsf g_0=106C_{\mathsf k}^2/59^2$. \emph{The moments.} The mean of $P_J$ is a sum over the level of the first creation of a large field (equation~\eqref{eq:ward-first-creation}), and along the histories without large fields the quotient of a product of the forms at the level $t$ is the sum over the partitions of the products of its blocks, as in the proof of Lemma~\ref{cov-mom}, now with $c_a=0$. Four kinds of terms occur.
(i) \emph{Blocks that join two pairs.} Their supports are at least $11ML^{K-i}$ units of the level $i$ apart, and the mixed derivative families decay exponentially in this distance (proof of Lemma~\ref{cov-mom}); such a block is at most $\mathsf{far}_{g,2n}(s)\prod_a\hat m_a$, with a function $\mathsf{far}_{g,2n}$ of the depth, formed with the constants, $g$ and $n$, that is a power of $L^s$ times $e^{-11L^{s-2}/\mathsf{Lm}_s}$, $\mathsf{Lm}_s$ the number in \eqref{eq:short-distance-envelope-names}, and so tends to $0$.
(ii) \emph{A member alone.} The mean of $q_b$ on the torus vanishes, by the reflection through $b$ in the direction $\lambda_b$; so the mean of $F_a$ vanishes, and by the identity of the first creation at $F_a$ (equation~\eqref{eq:ward-first-creation}), with the sizes and the tails of Lemmas~\YMMref{125.2}{840} and~\YMMref{125.4}{841} at the weight of the majorant, the quotient of $F_a$ at the upper level $T$ is at most $\mathsf e_g(s)\hat m_a$, with a function $\mathsf e_g$ of the depth that tends to $0$.
(iii) \emph{The block of one pair.} At the upper level it is $G_\tau$ up to $\epsilon_g(s)\hat m_\tau$, $\epsilon_g$ the error function in equation~\YMMref{(2330)}{846}, with the coefficient fixed in Lemma~\ref{lem:short-distance-collected-constant}. Step by step: the conditional covariance of the two members in one step is the total of the pointed family $\mathsf K^{(j)}$ of the level $j$ that the mixed derivative of the step map in the two slot members gives (Lemma~\ref{lem:weighted-step-covariance}, read for the slot members of the two sources, which enter the array of the step as the weight members of plaquette sources do); it vanishes off the regions that meet both supports. It is carried to the upper level by the first-order statement (equations~\YMMref{(2276)--(2277)}{831}, \YMMref{(2280)}{831} and Proposition~\YMMref{123.3}{832}) and by the tower of the quotients (equation~\YMMref{(2273)}{830}); at the upper level the carried family is flat on the small fields (Lemma~\ref{lem:weighted-regular-family} and the carried-family proof of Proposition~\YMMref{123.3}{832}); and its total at the trivial configuration is the Gaussian number of the slices of the covariance, formed with the kernels $\mathsf t_{(\tau,1)}$, $\mathsf t_{(\tau,2)}$ and the step couplings $\upsilon_j$, up to the error of the flat top of Proposition~\YMMref{125.1}{838}, with the same constants.

Indeed the proof of that Proposition reads a source only through its second order at a small point of the chart, which is (So3) with a kernel of two indices in the place of the diagonal kernel $\frac12w(p)\delta_{pp'}$ of an action density, and through the sizes formed at the weights of the majorants. The Gaussian numbers of the slices $\mathcal C_{(j)}$ of the steps above the working level, with the couplings $\upsilon_j$, add up, by the telescope \YMMref{(2227)}{818} and up to the errors of Theorem~\YMMref{AT}{839} at the attached levels (Lemma~\YMMref{125.5}{842}), to $G_\tau$, the Gaussian value at the background kernel of the working level.
(iv) \emph{The large fields.} On the small side the rarity of large fields (Lemma~\YMMref{125.3}{841}) bounds the contributions of the levels below the upper level by $\mathsf{tl}_{g,2n}(s)\prod_a\hat m_a$, with a function $\mathsf{tl}_{g,2n}$ of the depth that tends to $0$.
\emph{The sum over the partitions}, with $\delta_g=\epsilon_g+\mathsf e_g^2$ the error of a pair block together with its two singletons, gives the bound with
\[\omega_{g,2n}(s)=n\,\delta_g(s)\bigl(1+\mathsf g_0+\delta_g(s)\bigr)^{n-1}+\operatorname{Bl}(2n)\,\mathsf{far}_{g,2n}(s)\bigl(1+\mathsf g_0+\epsilon_g(s)\bigr)^n+\mathsf{sm}_g(s)(1+\mathsf g_0)^n+\mathsf{tl}_{g,2n}(s),\]
$\mathsf{sm}_g$ the function in \eqref{eq:short-distance-envelope-names} and $\operatorname{Bl}(2n)$ the number of the partitions of $2n$ elements. Each of these functions tends to $0$ as $s\to\infty$, by the asymptotics of the envelopes in the depth (the depth estimates in equations~\YMMref{(2332)}{846} and~\YMMref{(2336)}{847} and their proofs); $s_{2n}(g)$ is the least depth, not below the least depth of the two-point law of Theorem~\YMMref{AU}{847}, from which the thresholds of the proof, $\mathsf{far}_{g,2n}(s)\le1$, $\mathsf e_g(s)\le1$, the smallness of $s\,\mathsf{sm}_g(s)$ and the condition on the mixed derivative families of $2n$ sources, hold at every larger depth.
\end{proof}

\begin{proposition}\label{cov-od}
Let $0<g\le\gamma^{\rm sd}_W$ and $\mathsf l=\ell'/2$. Let $C_{\mathsf k}$ be the number of Lemma~\ref{cov-bg}, $s_8(g)$ the integer and $\omega_{g,8}$ the function of Lemma~\ref{cov-pairs} at four pairs, and $\bar\omega_{g,8}=\sup_{s\ge s_8(g)}\omega_{g,8}(s)<\infty$; these are formed with the constants and $g$. With the depth $s(\mathsf d)$ fitted to $\mathsf d$ at the coupling $g$, put $C_8=(\frac98)^4[(106C_{\mathsf k}^2)^4+59^8\bar\omega_{g,8}]$, $\mathsf d_\flat=L^{-s_8(g)}$ and
\[C_\infty(g)=\max\bigl\{e^{2G_g\mathsf l^3}C_8^{1/4},\ \bar C(g,\mathsf d_\flat)\bigr\},\]
where $\bar C(g,\mathsf d_\flat)$ is the number $C^G_2(\mathsf d_\flat;1)$ of Proposition~\ref{cov-lim}\,(a), raised, with the bound $|Y[\mathsf w]|\le2a_K^{-1}|\mathsf w|_\infty\sum_{p\in S[\mathsf w]}a_p$ of Lemma~\ref{cov-src}, so as to cover also the finitely many cutoffs below the threshold of that proposition. Then the following holds for all $(K,\beta)\in\mathfrak d(g)$, all $0<\mathsf d<1$ with $s(\mathsf d)\le K$, and all real weights $\mathsf w$ carried by bonds of one direction with $|x_b-x_0|\le\mathsf d-2a_K$ on the support.
\begin{enumerate}
\item[(a)] For every lattice mirror $x_3=c$, $c\in a_K(\mathbb Z+\frac12)$, with $x_{0,3}-c\in[\frac32\mathsf d,2\mathsf d]$, and every $\nu\in\mathcal G^{\rm per}_{K,\beta}$:
\[0\le\nu\bigl(\theta_cY[\mathsf w]\cdot Y[\mathsf w]\bigr)\le C_\infty(g)\,b^{-2}|\mathsf w|_\infty^2 .\]
On every torus $T_\ell$ of a side $\ell\in\mathfrak L(g)$, $0\le\langle\theta_cY[\mathsf w];Y[\mathsf w]\rangle_{K,\ell,\beta}\le C_\infty(g)b^{-2}|\mathsf w|_\infty^2$ for the reflection $\theta_c$ of $T_\ell$ in the lattice planes $x_\mu=c$ and $x_\mu=c+\frac\ell2$, $c\in a_K(\mathbb Z+\frac12)$, in each of the three directions $\mu$, whenever $x_{0,\mu}-c$ lies in $[\frac32\mathsf d,\frac\ell2-\frac32\mathsf d]$ modulo $\ell$.
\item[(b)] On every torus $T_\ell$, $\ell\in\mathfrak L(g)$: for weights $\mathsf w_1,\mathsf w_2$ as above about two centres $x_1,x_2$ at a sup distance $|x_1-x_2|_\infty\ge3\mathsf d+a_K$ in $T_\ell$, $|\langle Y[\mathsf w_1];Y[\mathsf w_2]\rangle_{K,\ell,\beta}|\le C_\infty(g)b^{-2}|\mathsf w_1|_\infty|\mathsf w_2|_\infty$. In periodic Gibbs states the statement is the mirror form of (a), at the heights $x_{0,3}-c\in[\frac32\mathsf d,2\mathsf d]$.
\item[(c)] Every tensor limit at $g$ has (O) with $C_{\rm od}=C_\infty(g)b_*^{-2}$.
\item[(d)] \emph{(All six components, with a power.)} There are $\mathsf d'_\flat>0$, $C_\flat(g)$ and an affine function $\mathsf Y_g$ of $\log\frac1{\mathsf d}$ such that for $0<\mathsf d\le\mathsf d'_\flat$, $\frac32\mathsf d\le x_{0,3}-c\le\frac12\mathsf l$, every $\nu\in\mathcal G^{\rm per}_{K,\beta}$ at large $K$, and every bounded $H:\mathbb R^3\to V_{\mathbb C}$ vanishing off $\bar B(x_0,\mathsf d)$,
\[0\le\nu\bigl(A^\nu_K(\theta_c\bar H,H)\bigr)\le E_g(\mathsf d)^2b^{-2}|H|_\infty^2,\qquad E_g(\mathsf d)\le C_\flat(g)\,\mathsf d^{-1}\mathsf Y_g(\mathsf d)^{10\mathfrak p_0} .\]
So every tensor limit at $g$ has (O$_a$) for every $a\in(2,4)$.
\end{enumerate}
\end{proposition}
\begin{proof}
(a) The lower bound is reflection positivity at the bond mirror (Per), and on $T_\ell$ that of the torus measure (Corollary~\ref{IV-18-18}\,(a)). Take the board with cells of side $\mathsf l$ that has $x_3=c$ as a wall plane and $x_0$ in the middle of a cell in the two other directions; $Y[\mathsf w]$ and $\theta_cY[\mathsf w]$ are functions of the two cells that $\theta_c$ exchanges, and by Lemma~\ref{cov-cell}\,(c) $\nu(\theta_cY[\mathsf w]\cdot Y[\mathsf w])\le e^{2G_g\mathsf l^3}\|Y[\mathsf w]\|^2_{{\rm ch},\ell'}$. The eighth power $M_8$ of the seminorm is the mean on $T_{\ell'}$ of the eight mirror copies of $Y[\mathsf w]$; as the period of the torus is $2\mathsf l$, they form four close pairs, each a copy and its image across a wall, at distances in $[3\mathsf d,4\mathsf d]$.

The copies are $Y[\mathsf w_a]$, $\mathsf w_a$ the mirror images of $\mathsf w$, real weights of the bonds of one direction with $|\mathsf w_a|_\infty=|\mathsf w|_\infty$; corners of plaquettes of copies of different pairs are at a sup distance $\ge\mathsf l-2\mathsf d\ge LM-2\ge11M$, as $\mathsf l\ge LM$ for a small side (the second inequality of equation~\YMMref{(2312)}{841}). So Lemma~\ref{cov-pairs} applies at four pairs, on the least torus, which is small over the run (Lemma~\YMMref{126.3}{846}). With $J$ the set of the four pairs, $M_8=\langle P_J\rangle_{K,\ell',\beta}$; $|G_\tau|\le106C_{\mathsf k}^2\upsilon^2|\mathsf w|_\infty^2=(106C_{\mathsf k}^2/59^2)\hat m_\tau$, uniformly in $\mathsf d$, as natural scaling demands: the covariance of two quadratic forms in the field strength at a distance of the order $\mathsf d$, integrated over volumes of the order $\mathsf d^3$, is dimensionless.

So at $s(\mathsf d)\ge s_8(g)$, $|\langle P_J\rangle|\le((106C_{\mathsf k}^2/59^2)^4+\omega_{g,8}(s(\mathsf d)))\hat m_J$, and $\hat m_J=59^8\upsilon^8|\mathsf w|_\infty^8\le59^8(\frac98)^4b^{-8}|\mathsf w|^8_\infty$ (Lemma~\ref{cov-mom8}); hence $0\le M_8\le C_8b^{-8}|\mathsf w|_\infty^8$, which gives the constant $e^{2G_g\mathsf l^3}C_8^{1/4}$. For $\mathsf d_\flat\le\mathsf d<1$ the two members $\theta_cY[\mathsf w]$ and $Y[\mathsf w]$, which are $\mathcal T_K$ of bounded off-diagonal members, are of the type $(\mathsf d_\flat;1)$, and the proof of Proposition~\ref{cov-lim}\,(a), which applies to bounded members, gives $\bar C(g,\mathsf d_\flat)$. On a torus $T_\ell$, $\ell\in\mathfrak L(g)$, the same two arguments, with the chessboard estimate on $T_\ell$ (Corollary~\ref{IV-18-18}\,(b)) and Lemma~\ref{cov-cell}\,(b) in the place of Lemma~\ref{cov-cell}\,(c) and of the domination in (Per), give the bound when $x_{0,\mu}-c\in[\frac32\mathsf d,2\mathsf d]$, in each direction by the invariance under $W_3$; the means of the $q_b$ vanish, so these forms are the truncated ones.

For the other heights let $Y_h$ be the translate of $Y[\mathsf w]$ along $e_\mu$ whose centre lies at the height $h$ over the plane $x_\mu=c$, $h\in x_{0,\mu}-c+a_K\mathbb Z$. By the invariance of the torus measure under the translations along $e_\mu$, $\langle\theta_cY_u\cdot Y_v\rangle_{K,\ell,\beta}$ depends on $u+v$ alone; Schwarz' inequality for the positive form $(F,G)\mapsto\langle\overline{\theta_cF}\,G\rangle_{K,\ell,\beta}$ on the functions of $\{c<x_\mu<c+\frac\ell2\}$ then gives, for $f(h)=\langle\theta_cY_h\cdot Y_h\rangle_{K,\ell,\beta}\ge0$,
\[f(h)^2=\langle\theta_cY_{h-a_K}\cdot Y_{h+a_K}\rangle_{K,\ell,\beta}^2\le f(h-a_K)\,f(h+a_K)\qquad\bigl(\tfrac32\mathsf d<h<\tfrac\ell2-\tfrac32\mathsf d\bigr).\]
So the maximum of $f$ over the heights in $[\frac32\mathsf d,\frac\ell2-\frac32\mathsf d]$ is attained at the lowest or at the highest of them: at the leftmost interior point $h$ of a maximum $M>0$ one would have $M^2\le f(h-a_K)f(h+a_K)<M^2$. These two heights lie in $[\frac32\mathsf d,2\mathsf d]$ over the plane $x_\mu=c$, resp. in $[\frac32\mathsf d,2\mathsf d]$ under the plane $x_\mu=c+\frac\ell2$, which $\theta_c$ fixes as well, and there the bound holds. (b) Let $\mu$ be a direction with $|x_{1,\mu}-x_{2,\mu}|\ge3\mathsf d+a_K$ on the circle $\mathbb R/\ell\mathbb Z$. On the shorter arc between $x_{1,\mu}$ and $x_{2,\mu}$ there is $c\in a_K(\mathbb Z+\frac12)$ at distances $\ge\frac32\mathsf d$ from both, and these distances are at most $\frac\ell2-\frac32\mathsf d$; the two weights lie on the two sides of the reflection $\theta_c$ of $T_\ell$. As the means of the $q_b$ vanish, Schwarz' inequality for $\theta_c$ gives
\[\bigl|\langle Y[\mathsf w_1]\,Y[\mathsf w_2]\rangle_{K,\ell,\beta}\bigr|\le\langle Y[\mathsf w_1]\cdot\theta_cY[\mathsf w_1]\rangle_{K,\ell,\beta}^{1/2}\,\langle\theta_cY[\mathsf w_2]\cdot Y[\mathsf w_2]\rangle_{K,\ell,\beta}^{1/2},\]
and (a) on $T_\ell$ bounds both factors, the form being the same for a centre on either side of the plane by the invariance under $W_3$. (c) $\mathcal T_K(h\,o)=Y[\mathsf w]$ with $\mathsf w(b)=h(x_b)$ on the bonds normal to the plane of $e_\mu,e_\nu$; $\nu(q_b)=0$ by the reflection through $b$ in the direction $\lambda_b$, an element of $a_K\mathbb Z^3\rtimes W_3$ (Per); and the continuity of Proposition~\ref{cov-lim}\,(c).
(d) The cell of side $\mathsf l$ of the least torus with the wall plane $x_3=c$ and the centre above $x_0$ holds $\bar B(x_0,\mathsf d)$ at distance $\ge\frac12\mathsf d$ from its walls. Cut $\mathcal T_K(H)$ into the at most $54$ diagonal and off-diagonal forms of $27$ boxes of side $\frac23\mathsf d$; put $\mathsf d'=\frac34\mathsf d$. The eight mirror copies of each form are the first members of an admissible system of real sources (Lemma~\ref{cov-src}\,(a), (b)) under the majorants $(\mathsf s,S_\kappa)$ with $\mathsf s\le59|H|_\infty$; the barycentres of the plaquettes of each $S_\kappa$ lie within $\frac78\mathsf d'$ of the image of the centre of its box, and those of two different copies are at least $\frac54\mathsf d'$ apart, since two points of different cells, each at a distance $\ge\frac23\mathsf d'$ from the walls, are $\ge\frac43\mathsf d'$ apart.

Lemma~\ref{cov-mom8} at the depth $s=s(\mathsf d')$ gives numbers $c_\kappa$ with $|\langle\prod_{\kappa\in B}(F_\kappa-c_\kappa)\rangle_{K,\ell',\beta}|\le(\mathsf A_8(s)\upsilon\mathsf s)^{\#B}$ for every set $B$ of copies; with $m$ the common mean of the copies, $|c_\kappa-m|\le\mathsf A_8(s)\upsilon\mathsf s$, and expanding $\prod_\kappa(F_\kappa-m)$ as in the proof of Proposition~\ref{cov-sd} bounds the seminorm on $T_{\ell'}$ of each form less its mean by $2\mathsf A_8(s)\upsilon\mathsf s$, with $\mathsf A_8(s)\le C\mathsf d^{-1}\mathsf Y_g(\mathsf d)^{10\mathfrak p_0}$ by Lemma~\ref{cov-mom8} and $\upsilon^2\le\frac98b^{-2}$. Lemma~\ref{cov-cell} carries the bound to all sides and to every periodic Gibbs state; the reflection is an element of $a_K\mathbb Z^3\rtimes W_3$, and reflection positivity gives the sign.

The power $\mathsf d^{-2}$ up to logarithms is $\le C_a\mathsf d^{-a}$ for every $a>2$.
\end{proof}

\begin{lemma}\label{cov-smooth}
Let $\mathcal P$ be a projection-valued measure on $[0,\infty)\times\mathbb R^2$ carried by $\{E\ge\max(|q_1|,|q_2|)\}$, $e^{-tH}\mathcal U(v)=\int e^{-tE+iv\cdot q}\,d\mathcal P$, $Q=\{|x-z|_\infty\le\varrho\}$, $\varepsilon\ge12\varrho$, and $\mathsf b$ a linear map from the continuous functions with support in $Q$ to the bounded operators, with $\|\mathsf b(\mathsf f)\|\le\mathsf c\sup|\mathsf f|$ and
\[e^{-\varepsilon H}\mathsf b(\tau_u\mathsf f)e^{-\varepsilon H}=e^{-(\varepsilon+u_3)H}\mathcal U(\vec u)\mathsf b(\mathsf f)\mathcal U(-\vec u)e^{-(\varepsilon-u_3)H}\]
whenever $|u|_\infty\le2\varrho$ and $\mathsf f$, $\tau_u\mathsf f$ have support in $Q$; here $\vec u=(u_1,u_2)$. Then $\|e^{-\varepsilon H}\mathsf b(h)e^{-\varepsilon H}\|\le4\mathsf c\,\varrho^{-3}\|h\|_{L^1}$.
\end{lemma}
\begin{proof}
Let $\chi$ be the tent of mass one on $Q$ and $\mathsf f_u(y)=\chi(y)h(y+u)$; then $h=\int\tau_u\mathsf f_u\,du$, and the covariance gives $\langle\psi',e^{-\varepsilon H}\mathsf b(h)e^{-\varepsilon H}\psi\rangle=\int\langle\psi'_u,\mathsf b(\mathsf f_u)\psi_u\rangle du$ with $\psi_u=\mathcal U(-\vec u)e^{-(\varepsilon-u_3)H}\psi$, and likewise $\psi'_u$. On the carrier of $\mathcal P$ one has $|q_k|\le E$, so $\psi_u=\sum_\alpha\frac{u^\alpha}{\alpha!}A^\alpha e^{-\varepsilon H}\psi$ with $\|A^\alpha e^{-\varepsilon H}\|\le\sup_EE^{|\alpha|}e^{-\varepsilon E}\le|\alpha|!\,\varepsilon^{-|\alpha|}$, a series that converges for $|u|_1<\varepsilon$. Exchanging sums and integral, $\int u^\gamma\mathsf b(\mathsf f_u)du=\mathsf b(k_\gamma)$ with $\sup|k_\gamma|\le\varrho^{-3}(2\varrho)^{|\gamma|}\|h\|_{L^1}$, and with $\sum_{|\alpha|=n}n!/\alpha!=3^n$ the double series is $(\sum_n(6\varrho/\varepsilon)^n)^2\le4$. So a point mass is as good as a bump once $e^{-\varepsilon H}$ makes the translates analytic, at the price of three powers of the scale.
\end{proof}

\begin{proposition}\label{cov-P}
Let $0<g\le\gamma^{\rm sd}_W$. For every $p>\frac{17}2$ there is $C_p$, a function of the constants, $g$ and $p$, such that every tensor limit at $g$ has (P) with $C=C_pb_*^{-1}$. Equivalently: for $\eta\in(0,1]$, $F,F'\in\mathbb C\mathfrak S_V(\{x_3>\eta\})$ and $h\in C_c(\{|x_3|<\frac12\eta\};V_{\mathbb C})$, $|\sigma(\theta\bar F\cdot(h)\cdot F')|\le2^pC_pb_*^{-1}\eta^{-p}\|h\|_{L^1}\|[\tau_{-\eta e_3}F]\|\,\|[\tau_{-\eta e_3}F']\|$.
\end{proposition}
\begin{proof}
Let $\vartheta=d-d'$, $i$ the least integer $\ge1$ with $a_i\le\vartheta$, $\varrho=a_i/32$, $\varepsilon=12\varrho=\frac38a_i$. A tent partition of unity at the scale $\varrho$ cuts $\tau_{de_3}h$ into pieces on cubes $Q$ of side $2\varrho$ whose $L^1$ norms add up to $\|h\|_{L^1}$. For a cube centred at the height $z_3$, Lemma~\ref{cov-smooth} applies to $\mathsf b(\mathsf f)=\mathsf B_{a_i}(\tau_{(\frac12a_i-z_3)e_3}\mathsf f)$, by Proposition~\ref{cov-lim}\,(d) and the relations $e^{-tH}\mathsf B_l(h)e^{-t'H}=\mathsf B_{l+t+t'}(\tau_{te_3}h)$, $e^{iv\cdot P}\mathsf B_l(h)e^{-iv\cdot P}=\mathsf B_l(\tau_{(v,0)}h)$; the remaining distances to the planes $x_3=0$ and $x_3=2d$ are $\ge\frac{15}{32}a_i>\varepsilon$ and are bridged by contractions $e^{-tH}$. So
\[\|\mathsf B_{2d}(\tau_{de_3}h)\|\le4\cdot32^3a_i^{-3}\mathsf c_E(i)b_*^{-1}\|h\|_{L^1},\qquad a_i^{-3}\mathsf c_E(i)=2C_E\mathsf Y_i^{q_E}a_i^{-17/2}\le C_p\min\{1,\vartheta\}^{-p}\]
for $p>\frac{17}2$, as $a_i\ge\min\{1,\vartheta\}/L$. The equivalent form follows with $d=\eta$, $d'=\frac12\eta$ and two contractions.
\end{proof}

\begin{proposition}\label{cov-growth}
Let $0<g\le\gamma^{\rm sd}_W$ and $\sigma$ a tensor limit at $g$.
\begin{enumerate}
\item[(a)] There is $\mathsf c_4$, a function of the constants and $g$, such that for $n\ge1$, $0<\mathsf d\le1$ and $X\in\mathfrak S_V(\mathbb R^3)$ with $n$ members whose supports lie in balls of radius $1$ at pairwise distances $\ge\mathsf d$,
\[|\sigma(X)|\le(\mathsf c_4g^2e^{35/2})^n(n!)^{35/2}\mathsf d^{-19n/2}\prod_a|H_a|_\infty .\]
\item[(b)] For every $n\ge1$ there is exactly one complex linear functional $\mathsf S_n$ on the functions $\varphi:\mathbb R^{3n}\to V_{\mathbb C}^{\otimes n}$ whose components lie in the space $\mathcal S_{\ne}(\mathbb R^{3n})$ of Section~\YMMref{93}{713}, of Schwartz functions flat off the configurations of pairwise different points, that is bounded for one of the norms $|\varphi|_k$ (the largest of the norms $|\cdot|_k$ of Section~\YMMref{93}{713} of the $6^n$ components) and has $\mathsf S_n(H_1\otimes\dots\otimes H_n)=\sigma(H_1,\dots,H_n)$ at smooth tuples with compact, pairwise disjoint supports; and $|\mathsf S_n(\varphi)|\le\mathsf C^n(n!)^{51/2}|\varphi|_{18n}$, $\mathsf C$ a function of the constants and $g$: $\sigma$ satisfies the growth hypothesis in Proposition~\YMMref{93.5}{719} with the constants $(18,1,\mathsf C,\frac{51}2)$, in its form for members with values in $V$.
\end{enumerate}
\end{proposition}
\begin{proof}
(a) At a cutoff, Proposition~\ref{cov-lim}\,(a) at the scale $i$ at which both the bound of Proposition~\ref{cov-sd} holds at the margin $\delta_n=\frac1{8n}$ and $a_i\le\frac14\mathsf d$, applied to the mean in one periodic Gibbs state of a product centred in another, gives $(\mathsf c_4n^{35/2}\mathsf d^{-19/2}b^{-1})^n\prod|H_a|_\infty$: the number of pieces, the factor $\delta_n^{-4}$ of $C_E$ and $a_i^{-11/2}\mathsf Y_i^{q_E}$ are powers of $n$ and of $1/\mathsf d$. With $b^{-1}<g^2$, $n^n\le e^n n!$ and the continuity of Proposition~\ref{cov-lim}\,(c) this passes to $\sigma$. (b) The extension proposition, Proposition~\YMMref{93.3}{717}, gives, for a multilinear map with $|\mathsf T(X)|\le\mathsf c_{\mathsf T}\max\{1,\mathsf d^{-\mathsf a}\}\prod\|f_i\|_\sharp$ at tuples with $R(X)\le1$ and separation $\mathsf d(X)\ge\mathsf d$, exactly one extension $\tilde{\mathsf T}$ to $\mathcal S_{\ne}(\mathbb R^{3n})$ bounded for one of the norms, with $|\tilde{\mathsf T}(\varphi)|\le\mathsf C_{\ne}^{n+\mathsf a}M^n(n!)^8\mathsf c_{\mathsf T}|\varphi|_{8n+\mathsf a-1}$.

Applied component by component with $\mathsf a=10n$, as $\min\{1,\mathsf d\}^{-19n/2}\le\max\{1,\mathsf d^{-10n}\}$, with $\mathsf c_{\mathsf T}=(\mathsf c_4g^2e^{35/2})^n(n!)^{35/2}$, and with $|f|_\infty\le\|f\|_\sharp$, this gives the factorial power $8+\frac{35}2=\frac{51}2$ and the index $8n+10n-1\le18n$. Uniqueness is that of Lemma~\YMMref{93.2}{716}, component by component.
\end{proof}

The convergence on a torus along the full sequence of cutoffs rests on a comparison of two consecutive cutoffs at one label. For a side $\ell=2L^m$ and a cutoff $K$, $\Lambda_{K,\ell}=a_K\mathbb Z^3/\ell\mathbb Z^3$ is the lattice of $T_\ell$ at the spacing $a_K$, and $\mathcal M$ is the block averaging map of one renormalisation step, from the configurations on $\Lambda_{k+1,\ell}$ to those on $\Lambda_{k,\ell}$ (Definition~\textup{\ref{IV-62-62}}\,(3)); for a configuration $U$ on $\Lambda_{K,\ell}$ and a depth $0\le d\le K$, $\mathcal M^{K-d}U$ is its \emph{block field of the depth $d$}, a configuration on $\Lambda_{d,\ell}$. A \emph{weight on the sites} is a complex function $\mathsf v$ on the sites of $\Lambda_{K,\ell}$.

For $\{\lambda,\lambda',\mu\}=\{1,2,3\}$ put $\mathcal T_{\lambda\lambda'}[\mathsf v]:=Y[\mathsf w]$ with $\mathsf w(\langle x,x+a_Ke_\mu\rangle):=\mathsf v(x)$ and $\mathsf w:=0$ at the bonds of the two other directions, and $\mathcal T_{\lambda\lambda}[\mathsf v]:=2\mathcal O(w)$ with $w(p_{\kappa\kappa'}(x)):=\mathsf v(x)$ at the plaquettes of the plane $\{\kappa,\kappa'\}$ normal to $e_\lambda$, $p_{\kappa\kappa'}(x)$ the plaquette of that plane with the lowest corner $x$, and $w:=0$ at the other plaquettes. Then $\mathcal T_K(H)=\frac12\sum_{\lambda,\lambda'}\mathcal T_{\lambda\lambda'}[\mathsf v^{\rm pt}_{\lambda\lambda'}]$ with the \emph{point readings} $\mathsf v^{\rm pt}_{\lambda\lambda'}(x)=H_{\lambda\lambda'}(x+\frac12a_Ke_\mu)$ for $\lambda\ne\lambda'$ and $\mathsf v^{\rm pt}_{\lambda\lambda}(x)=H_{\lambda\lambda}(x+\frac12a_K(e_\kappa+e_{\kappa'}))$.

The \emph{cell reading} of a bounded Borel function $\psi$ on $T_\ell$ is the weight $\psi^{(K)}(x):=a_K^{-3}\int_{x+a_K[-\frac12,\frac12)^3}\psi$, and the \emph{cell average} of a bounded Borel $H$ is $\mathcal T^{\rm av}_K(H):=\frac12\sum_{\lambda,\lambda'}\mathcal T_{\lambda\lambda'}[H^{(K)}_{\lambda\lambda'}]$. The sites of $\Lambda_{K+1,\ell}$ in the cell $y+a_K[-\frac12,\frac12)^3$ of a site $y$ of $\Lambda_{K,\ell}$ form its \emph{block} $B(y)$ of $L^3$ sites, and the \emph{block mean} of a weight $\mathsf v_{\rm f}$ on the sites of $\Lambda_{K+1,\ell}$ is the weight $\mathsf v_{\rm A}(y):=L^{-3}\sum_{x\in B(y)}\mathsf v_{\rm f}(x)$ on the sites of $\Lambda_{K,\ell}$. As $L$ is odd, the cell of $y$ is the disjoint union of the cells of the $x\in B(y)$; so $\psi^{(K)}$ is the block mean of $\psi^{(K+1)}$.

\begin{lemma}\label{cov-twocut}
Let $0<g\le\gamma^{\rm ref}_W$, let $\mathsf b$ and $\ell=2L^m$ be as in Proposition~\ref{cov-torus} below, so that the setting (E) of Section~\ref{sec:cutoff-calibration} holds at the label $\mathsf b$ on the side $\ell$, and let $q\ge2$ and $\mathsf d>0$. We compare the lattice $\mathrm A=\Lambda_{K,\ell}$ at $\beta=L^K\mathsf b$ with the lattice $\mathrm B=\Lambda_{K+1,\ell}$ at $\beta=L^{K+1}\mathsf b$, $\langle\cdot\rangle_K=\langle\cdot\rangle_{K,\ell,L^K\mathsf b}$; their block fields of a depth $d$, the images $\mathcal M^{K-d}U$ and $\mathcal M^{K+1-d}U$ of the configurations, are configurations on one lattice $\Lambda_{d,\ell}$.
\begin{enumerate}
\item[(a)] \emph{(The laws of the block fields.)} Let $K_{\rm bf}$ and $\mathsf C_{\rm bf}$ be the numbers of Definition~\textup{\ref{IV-62-62}}\,(5) at the interval of labels $\{\mathsf b\}$, and $\mathsf p_0=2\mathfrak p_0+\frac12$. For every cutoff $K\ge K_{\rm bf}$, every depth $0\le d\le\lceil K/1260\rceil$ and every bounded gauge invariant Borel function $F$ of the configurations on $\Lambda_{d,\ell}$, $|\langle F(\mathcal M^{K+1-d}U)\rangle_{K+1}-\langle F(\mathcal M^{K-d}U)\rangle_K|\le\mathfrak o_K\sup|F|$ with $\mathfrak o_K\le\mathsf C_{\rm bf}\ell^3K^{\mathsf p_0}L^{-K/10}$.
\item[(b)] \emph{(Tied sources.)} Let $K\ge K_{\rm bf}$, and let $1\le s<\varpi$ be depths with $\varpi\le\min\{K-1,\lceil K/1260\rceil\}$ and $L^\varpi\mathsf b\ge\Theta_\star$, the threshold of equation~\YMMref{(2148)}{794}, and with $s$ so large that $L^s\mathsf d\ge3M$ and that the data of the weighted progression of the depth $s$ are given on both lattices (Proposition~\textup{\YMMref{122.1}{826}}, read along the run up to the level $K-s$, resp.\ $K+1-s$; the two smallness conditions on the numbers of the run that this asks hold for all large $s$). A \emph{small block field of the depth $\varpi$} is a configuration $V$ on $\Lambda_{\varpi,\ell}$ with $|\partial V(p)-\mathbf 1|<\varepsilon_{\rm r}:=2c_\rho\,\varepsilon(L^\varpi\mathsf b)$ at every plaquette, where $\varepsilon(x)=A_0x^{-1/2}(\log x)^{\mathfrak p_0}$ is the threshold of \YMMref{(634)}{339} and $c_\rho$, set in Section~\ref{sec:cutoff-calibration}, is the number of the constants in the radii of the cuts (the definitions preceding Lemma~\YMMref{113.10}{796}); over it the background $U^{\rm A}(V)$ of $t:=K-\varpi$ renormalisation steps in $\mathrm A$ and the background $U^{\rm B}(V)$ of $t+1$ steps in $\mathrm B$ exist (the setting preceding Lemma~\YMMref{113.8}{794}, equation~\YMMref{(2149)}{794}), and a gauge invariant function takes one value on each of their minimal orbits.

Let $F^{\rm A}_a$ on $\mathrm A$ and $F^{\rm B}_a$ on $\mathrm B$, $1\le a\le q$, be the first members of sources under majorants $(\mathsf s_a,S^{\rm A}_a)$, resp.\ $(\mathsf s_a,S^{\rm B}_a)$ (Definition~\ref{cov-maj}), with one number $\mathsf s_a>0$ for the two lattices, for these data (in the sense of Lemma~\ref{cov-src}, which holds at them on the side $\ell$: see the proof), and so, on each lattice, the first members of an admissible system of sources (Lemma~\ref{cov-src}\,(a), (b)), the corners of the plaquettes of $S^\bullet_a$ and of $S^\bullet_{a'}$, $a\ne a'$, being at sup distances $\ge\mathsf d$; and let them be \emph{tied at the depth $\varpi$ with the defects $\tau_a$}: $|F^{\rm B}_a(U^{\rm B}(V))-F^{\rm A}_a(U^{\rm A}(V))|\le\tau_a$ at every small block field $V$ of the depth $\varpi$.

At these backgrounds $|F^\bullet_a|\le\mathsf S_a:=\frac32B_3^2\ell^3L^{4\varpi}\mathsf s_a$, where $B_3\ge1$ is the constant of Section~\YMMref{52.7}{432} through which equation~\YMMref{(2149)}{794} bounds the backgrounds. Then, for every real $\mathfrak b\ge4q$,
\[\bigl|\langle F^{\rm B}_1;\dots;F^{\rm B}_q\rangle_{K+1}-\langle F^{\rm A}_1;\dots;F^{\rm A}_q\rangle_K\bigr|\le\mathsf E_q\prod_a\mathsf s_a+\mathfrak c_q\sum_a\tau_a\prod_{b\ne a}\mathsf S_b+q\,\mathfrak c_q\,\mathfrak o_K\prod_a\mathsf S_a ,\]
where $\mathfrak c_q=\sum_\pi(|\pi|-1)!$, the sum over the partitions $\pi$ of a set of $q$ elements, and
\[\mathsf E_q:=2\,\mathfrak K'_q\Bigl[2L^{4q}\,\mathsf V_\ell(\mathfrak b)\,L^{-(\mathfrak b+1)s}+q\,\epsilon_{\rm fg}(\varpi-s)+2e^{3-L^s\mathsf d/M}\Bigr]\bigl(4\ell^3L^{4s}\bigr)^q\]
is the sum of the weight of the large fields, of the error of forgetting from the depth $\varpi$ to the depth $s$, $\epsilon_{\rm fg}(\nu)=C_{\rm fg}(\nu+1)^{2\mathfrak p_0+2}L^{-\nu/2}$, and of the terms on polymers joining two supports, which are exponentially small in $L^s\mathsf d$, each times the size $(4\ell^3L^{4s})^q\prod_a\mathsf s_a$ of the sources at the depth $s$, the factor $2$ counting the two lattices. Here $\mathfrak K'_q$ is a number of the constants, $q$ and $L^{4q}\mathsf V_\ell(\mathfrak b)L^{-(\mathfrak b+1)s}$, bounded by a number of the constants and $q$ when this last is at most $1$, and $\mathsf V_\ell(\mathfrak b)$, a number of the constants, $g$, $\ell$ and $\mathfrak b$ that carries the side through a factor $\ell^3e^{\mathsf B(X_g)\ell^3}$, bounds the weights $\mathsf w_t$ of a first creation of a large field at the level $t$: $\sum_{t\le K-s}L^{4q(K-s-t)}\mathsf w_t\le L^{-(\mathfrak b+1)s}\mathsf V_\ell(\mathfrak b)$, by the rarity of large-field regions on one torus (Proposition~\ref{prop:ward-fixed-volume-rarity}), read along the run on the side $\ell$.
\item[(c)] \emph{(Cell averages.)} For a weight $\mathsf v_{\rm f}$ on the sites of $\mathrm B$, its block mean $\mathsf v_{\rm A}$ and every pair $(\lambda,\lambda')$, the pair formed by $\mathcal T_{\lambda\lambda'}[\mathsf v_{\rm A}]$ on $\mathrm A$ and $\mathcal T_{\lambda\lambda'}[\mathsf v_{\rm f}]$ on $\mathrm B$ is tied at every depth $\varpi$ as in (b), with the defect $\tau\le\mathsf K_{\rm av}\ell^3L^{4\varpi}K^{1/2}L^{-(K-\varpi)/2}\sup|\mathsf v_{\rm f}|$, $\mathsf K_{\rm av}$ a number of the constants. Let $\varpi:=\lfloor K/(1260q)\rfloor$, $s:=\lfloor\varpi/(8q+3)\rfloor$, $\mathfrak b:=4q+1$ and $\vartheta_q:=1/(1260q(8q+3))$. There are numbers $\mathsf C^{\rm av}$, $\mathsf p$ of the constants, the label $\mathsf b$ and $q$, and an integer $K_1$ of the constants, $g$, $\mathsf b$, $\ell$, $q$ and $\mathsf d$, such that for $K\ge K_1$ and all real bounded Borel functions $\psi_1,\dots,\psi_q$ on $T_\ell$ that vanish off closed sets at pairwise sup distances $\ge\mathsf d$, the right side of (b) for the cell entries $\mathcal T_{\lambda_a\lambda_a'}[\psi_a^{(K)}]$ on $\mathrm A$ and $\mathcal T_{\lambda_a\lambda_a'}[\psi_a^{(K+1)}]$ on $\mathrm B$ is at most $2^q\mathsf C^{\rm av}\ell^{3(q+1)}K^{\mathsf p}L^{-\vartheta_qK}\prod_a\sup|\psi_a|$; so, by multilinearity, the joint cumulants of the cell averages of $\mathcal T_K(H_a)$ and of $\mathcal T_{K+1}(H_a)$, for bounded Borel $H_a$ with values in $V_{\mathbb C}$ that vanish off such sets, differ by at most $18^q\mathsf C^{\rm av}\ell^{3(q+1)}K^{\mathsf p}L^{-\vartheta_qK}\prod_a|H_a|_\infty$.
\end{enumerate}
\end{lemma}
\begin{proof}
(a) is Corollary~\textup{\ref{IV-62-66}}\,(a): at $n=1$ the depth $u$ of Section~\ref{sec:cutoff-calibration} is $\lceil K/1260\rceil$, and at $K\ge K_{\rm bf}$ the conditions under which that Corollary is stated hold.

(b) \emph{Sources on every side.} The admission of Lemma~\ref{cov-src} holds for all data of the weighted progression with at least one level, on every torus: a source enters the progression only through its domination by the plaquette action with its weight, its slot members, its family in the fourth place at the height zero and its sum rules, and none of these reads the side, so that the admission of the bond forms holds for all such data. So, at the data of the depth $s$ on every side $\ell\in\mathfrak L(g)$, $g\le\gamma^{\rm ref}_W$, $\mathcal O(w)$ is the first member of a source under $(|w|_\infty,\operatorname{supp}w)$ and $Y[\mathsf w]$ of one under $(59|\mathsf w|_\infty,S[\mathsf w])$; sums and multiples of sources are sources; and a source under $(\mathsf s,S)$ is one under every $(\mathsf s',S')$ with $\mathsf s'\ge\mathsf s$, $S'\supset S$. A side that is small over the run is read only by the rarity of large fields, in Lemmas~\ref{cov-mom} and~\ref{cov-mom8}; here the rarity on the side $\ell$ enters through the number $\mathsf V_\ell(\mathfrak b)$.
\emph{One lattice.} In each of the two lattices the truncated function of the sources is, up to $\frac12\mathsf E_q\prod_a\mathsf s_a$, that of the functions $\mathsf X^\bullet_a$ of the block field of the depth $\varpi$ that take the values $F^\bullet_a(U^\bullet(V))$ at the small block fields $V$ and vanish at the others. The moments are sums over the level of the first creation of a large field (equation~\eqref{eq:ward-first-creation}); the histories in which one is created cost $2L^{4q}\mathsf V_\ell(\mathfrak b)L^{-(\mathfrak b+1)s}$; along the histories without large fields, the progression of the depth $s$ and the one in which the sources enter at the depth $\varpi$ give conditional means that differ by $q\,\epsilon_{\rm fg}(\varpi-s)$, the error of the Lipschitz form of forgetting (Proposition~\YMMref{114.3}{803}); and the terms of the exponent on polymers joining two supports are at most $2e^{3-L^s\mathsf d/M}$, exponentially small in the distance in units of the level $K-s$; each of the three times the size of the sources at the depth $s$, with the number $\mathfrak K'_q$ for the passage from moments to cumulants.

At a background of a small block field $a_p<\frac12B_3^2\varepsilon_{\rm r}^2\eta_t^4$, $\eta_t=L^{-t}$, at every bare plaquette (equation~\YMMref{(2149)}{794}); with $\varepsilon_{\rm r}\le1$ the domination by the plaquette action gives $|\mathsf X^\bullet_a|\le\mathsf S_a$.
\emph{Two backgrounds, two laws.} A joint cumulant of $q$ bounded functions is multilinear, a polynomial in their mixed moments with coefficients of $q$ alone, and at most $\mathfrak c_q\prod\sup|\cdot|$. Under one law, replacing $\mathsf X^{\rm A}_a$ by $\mathsf X^{\rm B}_a$, one index $a$ at a time, changes the cumulant by at most $\mathfrak c_q\tau_a\prod_{b\ne a}\mathsf S_b$, by the tie. The products $\prod_{a\in B}\mathsf X^{\rm A}_a$ are bounded gauge invariant Borel functions of the block field of the depth $\varpi\le\lceil K/1260\rceil$, and by (a) their means under the laws of the block fields of $\mathrm B$ and of $\mathrm A$ differ by at most $\mathfrak o_K\prod_{a\in B}\mathsf S_a$; through the formula of the cumulant by moments this costs $q\,\mathfrak c_q\,\mathfrak o_K\prod_a\mathsf S_a$.

(c) \emph{The tie.} Let $V$ be a small block field of the depth $\varpi$, $t=K-\varpi$, $\eta_t=L^{-t}$, $N_t=\ell^3L^{3\varpi}$ the number of its sites, $U^{\rm A}=U^{\rm A}(V)$, $U^{\rm B}=U^{\rm B}(V)$, and $U':=\mathcal M(U^{\rm B})$, a configuration on $\Lambda_{K,\ell}$. It is a competitor at the level $t$ of $\mathrm A$ over the same block field, $\mathcal M^tU'=V$, and $\mathsf{ab}^{\rm A}\le A^{\eta_t}(U')\le(1+\mathsf j_t)^2\mathsf{ab}^{\rm B}$, $\mathsf{ab}^\bullet$ the actions of the two backgrounds and $\mathsf j_t\le\frac25$ (Lemma~\YMMref{113.8}{794}); with $|\mathsf{ab}^{\rm B}-\mathsf{ab}^{\rm A}|\le\mathsf K_u\varepsilon_{\rm r}^2\eta_tN_t$ (equation~\YMMref{(2155)}{795}) the energy that the one averaging step from $\mathrm B$ to $\mathrm A$ loses at the background of $\mathrm B$ is small: the sum $\mathsf{Jd}$ over the plaquettes of $\mathrm A$ of the defects of Jensen's inequality of the averaging, with errors of the third order, is at most $c\,\varepsilon_{\rm r}^2\eta_t^2N_t$, one power of $\eta_t$ below the energy $\sum_pa_p(U^{\rm B})\le\frac32B_3^2\varepsilon_{\rm r}^2\eta_{t+1}N_t$.
\emph{(i) The fine form against the coarse one, over one block field.} About a site $y$ of $\Lambda_{K,\ell}$, on the box formed by the $27$ blocks $B(y')$, $y'-y\in a_K\{-1,0,1\}^3$, a gauge brings $U^{\rm B}$ to $e^{ia}$ with a real $1$-form $a$, $\sup|a|\le3\pi L\varpi_{\rm r}$, $\varpi_{\rm r}:=B_3\varepsilon_{\rm r}\eta_{t+1}^2$; let $g_p:=a(\partial p)$ be the fine curvatures and $\bar G_l(y)$ the fine curvature at the coarse site $y$ in the plane of a loop $l$. The vector part of a loop of a plaquette is $\pm$ its curvature up to $c_{\rm f}\varpi_{\rm r}^2$; the coarse curvatures are the sums $G'_{p'}=\sum_p\mathsf W(p',p)g_p$ of the fine ones under the kernel $\mathsf W$ of the block averaging, $\sum_p\mathsf W(p',p)=L^2$, up to $c_{\rm c}\varpi_{\rm r}^2$; and the variance $\sum_p\mathsf W(p',p)|g_p-L^{-2}G'_{p'}|^2$ is at most the defect of Jensen's inequality at $p'$ plus an error of the third order. So $D^l_x:=g_{p(x,l)}-\bar G_l(y)$ has $\sum_{x\in B(y)}|D^l_x|^2\le12L^4V_y$, $V_y$ the variances of the four coarse plaquettes at $y$ in the plane of $l$, with $\sum_yV_y\le4\,\mathsf{Jd}$: the variation of the curvature inside a block is bounded by the lost energy, and not by the energy. With $Q_x(l,l')=\langle\mathbf w(W_{x,l}),\mathbf w(W_{x,l'})\rangle$ the product of the vector parts of the loops $l$, $l'$ at $x$ and $\sigma(l)=\pm1$ their orientations,
\[\bigl|Q_x(l,l')(U^{\rm B})-\sigma(l)\sigma(l')\langle\bar G_l,\bar G_{l'}\rangle(y)\bigr|\le2\varpi_{\rm r}\bigl(|D^l_x|+|D^{l'}_x|+2c_{\rm f}\varpi_{\rm r}^2\bigr)\qquad(x\in B(y)),\]
and likewise for $L^{-4}Q_y(l,l')(U')$ at the coarse site, through the same $\mathsf{Jd}$. Cauchy's and Schwarz' inequality over the $\eta_t^{-3}L^3N_t$ fine sites gives $\sum_x|D^l_x|\le(\eta_t^{-3}L^3N_t)^{1/2}(48L^4\mathsf{Jd})^{1/2}$; and when $\mathsf v_{\rm A}$ is the block mean of $\mathsf v_{\rm f}$, $L\sum_{x\in B(y)}\mathsf v_{\rm f}(x)=L^4\mathsf v_{\rm A}(y)$, so that the weights cancel exactly at the common term $\langle\bar G_l,\bar G_{l'}\rangle(y)$ ($a_{K+1}^{-1}L^3=a_K^{-1}L^4$). Hence
\[\Bigl|\eta_{t+1}^{-1}\sum_x\mathsf v_{\rm f}(x)Q_x(l,l')(U^{\rm B})-\eta_t^{-1}\sum_y\mathsf v_{\rm A}(y)Q_y(l,l')(U')\Bigr|\le\mathsf K^{\rm cf}_0\sup|\mathsf v_{\rm f}|\,\varepsilon_{\rm r}^2N_t\,\eta_t^{1/2},\]
$\mathsf K^{\rm cf}_0$ a number of the constants.
\emph{(ii) The competitor against the background of $\mathrm A$, over one block field.} A nearly minimal competitor has nearly the loops of the background. By the level-by-level comparison in the proof of Lemma~\YMMref{113.6}{792}, the configurations $\hat U$, $\check U$ of two consecutive levels of the competitor are related by $\hat U=(e^{iA'}\check U)^u$ with a real $1$-form $A'$ and one gauge transformation $u$, so that the loops at every corner differ by at most $\rho(p)$, $\eta_k^{-1}\sum_p\rho(p)^2\le2C_{\rm nrg}\|\mathcal A_k\|^2$; Cauchy's and Schwarz' inequality over the plaquettes and over the $t$ levels, $L^tA(\cdot)\le A^{\eta_t}(U')$ at the backgrounds of the levels, and the bound \YMMref{(2144)}{792} of the fields $\mathcal A_k$ by the excess give, as in the proof of Proposition~\YMMref{113.7}{793},
\[\Bigl|\eta_t^{-1}\sum_y\mathsf v(y)\bigl[Q_y(l,l')(U')-Q_y(l,l')(U^{\rm A})\bigr]\Bigr|\le12\sup|\mathsf v|\bigl(C_{\rm nrg}\gamma_{\rm r}^{-1}\,t\,\mathsf{exc}_t(U')\,A^{\eta_t}(U')\bigr)^{1/2},\]
with $\mathsf{exc}_t(U')=A^{\eta_t}(U')-\mathsf{ab}^{\rm A}\le2\mathsf K_u\varepsilon_{\rm r}^2\eta_tN_t$ and $A^{\eta_t}(U')\le3B_3^2\varepsilon_{\rm r}^2N_t$ (Lemma~\YMMref{113.8}{794}): this is $\varepsilon_{\rm r}^2N_t(t\eta_t)^{1/2}$ up to a number, and here the factor $t^{1/2}\le K^{1/2}$ enters.
\emph{(iii) The components.} $q_b=\frac14\sum_{\varepsilon,\varepsilon'=\pm}\varepsilon\varepsilon'Q_x((\mu,\varepsilon\lambda),(\mu,\varepsilon'\lambda'))$ at $b=\langle x,x+a_Ke_\mu\rangle$, the loops being named by their two directions; so $\mathcal T_{\lambda\lambda'}[\mathsf v]$, $\lambda\ne\lambda'$, is a mean of four forms $a_K^{-1}\sum_x\mathsf v(x)Q_x(l,l')$; and $2a_p=Q_x(l_m,l_m)+a_p^2$ for the loop $l_m$ of the plaquette $p=p_{\kappa\kappa'}(x)$, with $a_p^2<\frac14B_3^4\varepsilon_{\rm r}^4\eta_t^8$ at the backgrounds, does the same for $\mathcal T_{\lambda\lambda}[\mathsf v]$. In the units of the depth $\varpi$, $a_K^{-1}=L^\varpi\eta_t^{-1}$ and $a_{K+1}^{-1}=L^\varpi\eta_{t+1}^{-1}$; with $\varepsilon_{\rm r}\le1$, $t\le K$ and $\eta_t\le(t\eta_t)^{1/2}$, (i) and (ii) give the defect $\mathsf K_{\rm av}\ell^3L^{4\varpi}K^{1/2}L^{-(K-\varpi)/2}\sup|\mathsf v_{\rm f}|$.
\emph{The rate.} At $\varpi=\lfloor K/(1260q)\rfloor$, $s=\lfloor\varpi/(8q+3)\rfloor$, $\mathfrak b=4q+1$: $\varpi\le\lceil K/1260\rceil$, $\varpi-s\ge(8q+2)s$, $s\ge\vartheta_qK-2$ and $4q\varpi\le K/315$. Let $K_1\ge K_{\rm bf}$ be the least cutoff from which on $s\ge1$, $L^s\mathsf d\ge6M$, the data of the depth $s$ are given, $L^{4q}\mathsf V_\ell(\mathfrak b)L^{-s}\le1$, $2e^{3-L^s\mathsf d/(2M)}L^{(4q+1)s}\le1$, $L^\varpi\mathsf b\ge\Theta_\star$ and $6a_K\le\mathsf d$. The cell entries are first members of sources under the majorants $(59\sup|\psi_a|,S^\bullet_a)$, $S^\bullet_a$ the plaquettes all of whose corners are within $\frac32a_{K^\bullet}$ of the closed set off which $\psi_a$ vanishes, at corner distances $\ge\mathsf d-3a_K\ge\frac12\mathsf d$; so (b) applies at $\frac12\mathsf d$, with the ties of the first part of (c), $\psi^{(K)}$ being the block mean of $\psi^{(K+1)}$. Then $q\,\epsilon_{\rm fg}(\varpi-s)L^{4qs}\le qC_{\rm fg}K^{2\mathfrak p_0+2}L^{-s}$, and the three parts of $\mathsf E_q$ are at most a number times $K^{2\mathfrak p_0+2}\ell^{3q}L^{-s}$, $L^{-s}\le L^2L^{-\vartheta_qK}$; the ties contribute $(\ell^3L^{4\varpi})^qL^{-(K-\varpi)/2}\le\ell^{3q}L^{-K/3}$ times $K^{1/2}$, and the laws $(\ell^3L^{4\varpi})^qL^{-K/10}\le\ell^{3q}L^{-K/11}$ times $K^{\mathsf p_0}$; and $\vartheta_q<\frac1{11}$.
\end{proof}

\begin{proposition}\label{cov-torus}
Let $0<g\le\gamma^{\rm ref}_W$, $\mathsf b\in(g^{-2}+\bar B,X_g-\bar B]$, $\ell\in\mathfrak L(g)$ and $\langle\cdot\rangle_K=\langle\cdot\rangle_{K,\ell,L^K\mathsf b}$, so that $(K,L^K\mathsf b)\in\mathfrak d(g)$ for every $K$. Let $n\ge2$ and $\mathsf d>0$.
\begin{enumerate}
\item[(a)] There are $K_1$ and $\mathsf C^\flat_n(\mathsf d)$, functions of the constants, $g$, $\mathsf b$, $\ell$, $n$ and $\mathsf d$, such that for $K\ge K_1$ the joint cumulants of $\mathcal T_K(H_1),\dots,\mathcal T_K(H_n)$, for $H_a\in C(T_\ell;V_{\mathbb C})$ with supports at pairwise sup distances $\ge\mathsf d$, are at most $18^n\mathsf C^\flat_n(\mathsf d)\prod_a|H_a|_\infty$.
\item[(b)] The limit $\mathfrak S_n(H_1,\dots,H_n)=\lim_{K\to\infty}\langle\mathcal T_K(H_1);\dots;\mathcal T_K(H_n)\rangle_K$ exists along the full sequence of cutoffs, with a rate: the difference at $K$ is at most \[\begin{multlined}18^n\mathsf C^{\rm av}\ell^{3(n+1)}\sum_{K'\ge K}K'^{\mathsf p}L^{-\vartheta_nK'}\prod_a|H_a|_\infty\\{}+18^n\mathsf C^\flat_n(\mathsf d)\sum_a\omega_a(a_K)\prod_{b\ne a}|H_b|_\infty,\end{multlined}\] with the numbers $\mathsf C^{\rm av},\mathsf p>0$ and $\vartheta_n=1/(1260n(8n+3))$ of Lemma~\ref{cov-twocut} and $\omega_a$ the modulus of continuity of $H_a$. It is symmetric, $n$-linear, real at real members, bounded by $18^n\mathsf C^\flat_n(\mathsf d)\prod|H_a|_\infty$, and it extends to bounded Borel members vanishing off closed sets at pairwise distances $\ge\mathsf d$, continuously under bounded pointwise convergence. For the normalized plaquette and corner forms of Theorem~\ref{sf-tori}\,(c), the limit depends only on their common quadratic symbol, with the planes and orientation signs fixed there. In indices, the corresponding curvature bilinears tend to $\sum_{\lambda\lambda'}\varepsilon_{ij\lambda}\varepsilon_{kl\lambda'}T_{\lambda\lambda'}$, antisymmetric in $(ij)$ and in $(kl)$ and symmetric under their exchange.
\item[(c)] The centred functional $\sigma_\ell(X)=\sum_\pi\prod_{J\in\pi}\mathfrak S_{|J|}(X_J)$, the sum over the partitions into blocks of at least two members, is $\lim_K\langle A_K(X)\rangle_K$, $A_K$ centred in $\langle\cdot\rangle_K$; it is invariant under the translations of $T_\ell$ and under $W_3$, acting on values through $\rho$, and reflection positive at every coordinate plane of $T_\ell$.
\item[(d)] If $z\ge0$ is an integer such that $(L^z\mathsf b,L^z\ell)$ is admitted at some reference coupling, then $\mathfrak S^{\mathsf b,\ell}_n(H_\bullet)=L^{zn}\mathfrak S^{L^z\mathsf b,L^z\ell}_n(H_\bullet(L^{-z}\cdot))$.
\item[(e)] \emph{(Conservation.)} For every $C^2$ vector field $\eta$ on $T_\ell$ and all Lipschitz $H_2,\dots,H_n:T_\ell\to V$ whose supports are pairwise disjoint and disjoint from $\operatorname{supp}\eta$,
\[\mathfrak S_n(H_\eta,H_2,\dots,H_n)=0 ;\]
by (d), the same holds at every label and side reached by a dilation.
\end{enumerate}
\end{proposition}
\begin{proof}
The cell entries of bounded Borel functions are first members of sources under majorants (Definition~\ref{cov-maj}) at the data of the depth $s$ on $T_\ell$, and so first members of an admissible system of sources (Lemma~\ref{cov-src}\,(a), (b), which hold at these data on every side $\ell\in\mathfrak L(g)$, as in the proof of Lemma~\ref{cov-twocut}\,(b)), and the cell readings of a function at the cutoffs $K'$ and $K'+1$ are tied, the one at $K'$ being the block mean of the one at $K'+1$. \emph{A uniform bound for weights of no regularity.} Let $K_1$ be the integer of Lemma~\ref{cov-twocut}\,(c) at $(q,\mathsf d)=(n,\frac12\mathsf d)$, and
\[\mathsf C^\flat_n(\mathsf d):=\mathfrak c_n\bigl(2\ell^3L^{4K_1}\bigr)^n+\mathsf C^{\rm av}\ell^{3(n+1)}\sum_{K'\ge K_1}K'^{\mathsf p}L^{-\vartheta_nK'} .\]
For $K\ge K_1$, all pairs $(\lambda_a,\lambda'_a)$ and all weights $\mathsf v_a$ on the sites of $\Lambda_{K,\ell}$ whose sets $\{\mathsf v_a\ne0\}$ are at pairwise sup distances $\ge\frac12\mathsf d+a_K$,
\[\bigl|\langle\mathcal T_{\lambda_1\lambda'_1}[\mathsf v_1];\dots;\mathcal T_{\lambda_n\lambda'_n}[\mathsf v_n]\rangle_K\bigr|\le4^n\,\mathsf C^\flat_n(\mathsf d)\prod_a\sup|\mathsf v_a|.\qquad\text{(UB)}\]
Indeed, for real $\mathsf v_a$, the weight $\mathsf v_a$ at the cutoff $K$ is the cell reading at $K$ of the step function $\psi_a:=\sum_x\mathsf v_a(x)\mathbf 1_{x+a_K[-\frac12,\frac12)^3}$, with $\sup|\psi_a|=\sup|\mathsf v_a|$, and $\psi_a$ vanishes off the union of the closed cells of the sites of $\{\mathsf v_a\ne0\}$, closed sets at pairwise sup distances $\ge\frac12\mathsf d$. From $K$ down to $K_1$ the joint cumulants of the cell entries of the $\psi_a$ change from $K'$ to $K'+1$ by at most $\mathsf C^{\rm av}\ell^{3(n+1)}K'^{\mathsf p}L^{-\vartheta_nK'}\prod_a\sup|\psi_a|$ (for the cell entries $\mathcal O(w)$ of one plane and $Y[\mathsf w]$ of one direction), by Lemma~\ref{cov-twocut}\,(c), which reads no regularity of the functions.

At $K_1$ the crude bounds $|\mathcal O(w)|\le2\ell^3L^{4K_1}|w|_\infty$ for a weight carried by the plaquettes of one plane ($a_p\le2$, $\ell^3L^{3K_1}$ such plaquettes) and $|Y[\mathsf w]|\le\ell^3L^{4K_1}|\mathsf w|_\infty$ for a weight carried by the bonds of one direction ($|q_b|\le1$) hold, a cumulant of $n$ bounded functions being at most $\mathfrak c_n\prod\sup$. The factor $2$ of a diagonal component and the splitting of a complex weight into its real and imaginary parts give $4^n$. \emph{(a), (b).} $\mathcal T_K(H)=\frac12\sum_{\lambda,\lambda'}\mathcal T_{\lambda\lambda'}[\mathsf v^{\rm pt}_{\lambda\lambda'}]$, nine terms with the factor $\frac12$, and $\mathcal T^{\rm av}_K(H)$ is the same with the weights $H^{(K)}_{\lambda\lambda'}$; both weights vanish at the sites at sup distance $>\frac12a_K$ from $\operatorname{supp}H$, so that at $12a_K\le\mathsf d$, which holds for $K\ge K_1$, the sites of two members are $\ge\mathsf d-a_K\ge\frac12\mathsf d+a_K$ apart; and $|\mathsf v^{\rm pt}_{\lambda\lambda'}-H^{(K)}_{\lambda\lambda'}|\le\omega(a_K)$, $\omega$ the modulus of continuity of $H$ for the sup distance, as a point of the cell of $x$ and the point of the point reading are within $a_K$.

By multilinearity and (UB) the joint cumulants of the $\mathcal T_K(H_a)$ and of the $\mathcal T^{\rm av}_K(H_a)$ are at most $18^n\mathsf C^\flat_n(\mathsf d)\prod_a|H_a|_\infty$, which is (a), and they differ by at most $18^n\mathsf C^\flat_n(\mathsf d)\sum_a\omega_a(a_K)\prod_{b\ne a}|H_b|_\infty$; the cumulants of the cell averages converge, by Lemma~\ref{cov-twocut}\,(c) applied to the real and imaginary parts, with the rate $18^n\mathsf C^{\rm av}\ell^{3(n+1)}\sum_{K'\ge K}K'^{\mathsf p}L^{-\vartheta_nK'}\prod_a|H_a|_\infty$. This gives the existence and the rate in (b). For bounded Borel members the cumulants of the cell averages converge in the same way, with the bound of (a); at a fixed cutoff they are polynomials in the finitely many cell readings, which converge under bounded pointwise convergence by dominated convergence, and with the rate, uniform in such sequences, the limit is continuous under bounded pointwise convergence.

For the last assertion of (b), use Theorem~\ref{sf-tori}\,(c) and its proof. Each class insertion is reduced to the canonical corner loops: a shift costs $a_K$ times the Lipschitz constant, loops in the same planes are exchanged with their orientation signs by Lemma~\ref{sf-cut}\,(c), and $a_p=(Q_x(l_\lambda,l_\lambda)+a_p^2)/2$, with $a_p^2$ of fourth order at the backgrounds. At the same site weights there is no shift cost. In particular, the difference of two such lattice forms, with a complex weight $\mathsf v$ on the sites, has at one cutoff joint cumulants with $n-1$ sources under majorants $(\mathsf s_a,S_a)$ far from it of at most $C\ell^{3n}K^{2\mathfrak p_0+2}L^{-\vartheta_nK}\sup|\mathsf v|\prod_{a\ge2}\mathsf s_a$, for weights of any regularity. (c) The formula for centred moments by cumulants; lattice isometries of $T_\ell$ keep the Wilson measure and map $\mathcal T_K$ covariantly, with translations of $T_\ell$ reached by limits; reflection positivity of the Wilson measure at the lattice planes of the torus passes to the limit, the planes of one cutoff being planes of every finer cutoff ($L$ odd). (d) The dilation of the lattice is exact: $\mathcal T_K(H)$ on the torus of side $\ell$ is $L^z\mathcal T_{K-z}(H(L^{-z}\cdot))$ on that of side $L^z\ell$, the two being one abstract lattice. (e) The supports are compact and pairwise disjoint, hence at a positive sup distance $\mathsf d$; $H_\eta$ is Lipschitz.

By Theorem~\ref{cov-thm}\,(a) on $T_\ell$ the joint cumulants of $\mathcal T_K(H_\eta),\mathcal T_K(H_2),\dots,\mathcal T_K(H_n)$ in $\langle\cdot\rangle_K$ are at most $CK^pL^{-K/36}\|\eta\|_{C^2}\prod_{a\ge2}|H_a|_\infty$, and by (b) they converge to $\mathfrak S_n(H_\eta,H_2,\dots,H_n)$. At other labels and sides, $H_\eta(L^{-z}\cdot)=H_{\eta^{[z]}}$ with $\eta^{[z]}(x)=L^z\eta(L^{-z}x)$, and (d) applies.
\end{proof}

\subsection{Conservation at a cutoff}\label{cov-cut}

Here $0<g\le\gamma$. The setting (Th) of Section~\ref{sec:cutoff-calibration} is used on a torus $T_\ell$, $\ell=2L^m$, with data $(K,\mathsf b',m,t)$, $\gamma^{-2}\le\mathsf b_-\le\mathsf b'\le\mathsf b_+$, $MR(\mathsf b_+)\le L^m$; $x_k=L^{K-k}\mathsf b'$ is the line of the label $\mathsf b'$ of (Th3), which is not the run of a datum of Subsection~\ref{cov-set}, and $\rho_0\,dU$ is the bare measure. $\mathcal V_k$ is the set of the configurations of the level $k$, on the lattice of $T_\ell$ of the spacing $a_{K-k}$, and $\mathcal M^k:\mathcal V_0\to\mathcal V_k$ is the block map of $k$ renormalisation steps, $\mathcal M$ that of one step. \emph{The cuts} (the definitions preceding Lemma~\YMMref{113.10}{796}, Lemma~\YMMref{113.10}{796}): with the radius $\mathsf e^{\rm c}_k:=c_\rho\,\varepsilon(x_k)$, $\varepsilon(x)=A_0x^{-1/2}(\log x)^{\mathfrak p_0}$ the threshold of \YMMref{(634)}{339} and $c_\rho$ a number of the constants, $\mathsf{Cut}_0$ is the set of the bare configurations $V$ with $\max_p|\partial V(p)-\mathbf 1|<\mathsf e^{\rm c}_0$, and for $1\le k\le t$, $\mathsf{Cut}_k\subset\mathcal V_k$ is the set of the $V$ over which the background exists and every $U$ of its minimal orbit has $L^{2k}\max_p|\partial U(p)-\mathbf 1|<\mathsf e^{\rm c}_k$, the maximum over the bare plaquettes.

These sets are open and gauge invariant, and they are stable downward: for $W\in\mathsf{Cut}_{k+1}$ and every fluctuation field $B$ in its range, the configuration $V_W(B)$ of the level $k$ lies in $\mathsf{Cut}_k$. Every $U$ of the minimal orbit over $V\in\mathsf{Cut}_k$, $k\ge1$, is \emph{regular of level $k$ at the radius}
\[\varepsilon_k:=2\,\mathsf e^{\rm c}_k/\varepsilon_{\mathcal W}\le\alpha_0 ,\]
that is $|\partial U(p)-\mathbf 1|<\varepsilon_kL^{-2k}$ and $|J^{(k)}(U)|<\varepsilon_k$ at all bare plaquettes and bonds, $J^{(k)}$ the current at the spacing of the level $k$ (Definition~\ref{def:level-regularity}; Lemma~\ref{lem:sf-class-window}(1), $\varepsilon_{\mathcal W}$ a number of the constants), and $U$ lies on the minimal orbit over $\mathcal M^jU\in\mathsf{Cut}_j$ for $1\le j\le k$ (Lemma~\ref{lem:sf-class-window}(2)); besides, $\varepsilon_k^2\le CK^{2\mathfrak p_0}L^{-(K-k)}$, $C$ a number of the constants and $\mathsf b_+$. The weight of the chain is
\[p_k(U):=\mathbf 1_{\mathsf{Cut}_k}(V_k)\prod_{j<k}\bigl[\mathbf 1_{\mathsf{Cut}_j}\,\mathfrak G_j[x_j]\,\chi^{\rm sf}_{x_j}\bigr](V_j),\qquad V_j=\mathcal M^jU,\]
with the bounded Borel functions $\mathfrak G_j[x_j]\ge0$ and the small-field indicators $\chi^{\rm sf}_{x_j}$ of the steps (the definitions preceding Proposition~\YMMref{115.2}{806}, where the small-field weight of the level $t$ is $(\mathbf 1_{\mathsf{In}_t}\circ\mathcal M^t)\,p_t$). For the least torus at $g$ these data are admissible for all labels of $\mathfrak d(g)$, by Lemma~\ref{lem:cutoff-admissible-labels} on the tori of $\mathfrak L(g)$. A \emph{source} $w$ is a complex weight of the plaquettes in the class of sources $\mathsf{Src}$ of (Th1); $\rho^w_0=\rho_0e^{\mathcal O(w)}$.

\begin{definition}\label{cov-loc}
A \emph{local observable} of order $m\ge2$ is a finite sum $\mathfrak J(U)=\sum_\iota\varkappa_\iota\mathsf j_\iota(U|_{Q_\iota})$ over boxes $Q_\iota$ of at most $3\times3\times3$ sites that extend in at least two directions, where $\varkappa_\iota\in\mathbb C$, $\mathsf j_\iota$ is a gauge invariant polynomial in the entries of the bond matrices of $Q_\iota$, taken in $\mathrm{SL}(2,\mathbb C)$, with $|\mathsf j_\iota|\le\hat c$ where these matrices have norm $\le2$ and $|\mathsf j_\iota|\le c\,\delta^m$ where besides all plaquette variables of $Q_\iota$ lie within $\delta\le1$ of $\mathbf 1$, and every site is the lowest site of at most $N$ boxes; $(m,c,\hat c,N)$ is its \emph{type} and $|\varkappa|_\infty=\max_\iota|\varkappa_\iota|$. The \emph{conditional mean after $t$ steps} is $\mathsf Q_{t\leftarrow0}=\mathsf Q_{t-1}\circ\dots\circ\mathsf Q_0$, where $\mathsf Q_k$ sends a bounded complex Borel function $h$ of the field of the level $k$ to the function that vanishes off $\mathsf{Cut}_{k+1}$ and is, at every point of $\mathsf{Cut}_{k+1}$, the integral over the fluctuation field of the step $k$ of $h$ times the integrand of that step at the source $w=0$, divided by the integral of that integrand alone; the denominators do not vanish on the cut, and Borel functions go to Borel functions.
\end{definition}

The cubic density with the orbital potential taken off, $\mathfrak r^\sharp_\lambda(p)$ of Lemma~\ref{cov-cubic} below, is a local observable of order $3$ on a box of $2\times2\times3$ sites, of the type $(3,2^{36}\hat c_\sharp,\hat c_\sharp,N)$ with $\hat c_\sharp=23185$ and a pure number $N$: the bound $\hat c_\sharp$ at bond matrices of norm at most $2$ is the sum of the bounds of the terms of its closed form there. The families in the fourth place that represent it at the levels of the chain have sizes bounded through the pure number $48M^3\hat c_\sharp$.

\begin{lemma}\label{cov-chain}
Let $0<\phi<\frac1{42}$, $K_\phi$ the least integer $K\ge2$ with $(1-\phi)K\ge2$ and $L^{\phi K}\gamma^{-2}\ge\mathsf t_\natural$, $\mathsf t_\natural$ the calibration of Section~\ref{sec:cutoff-calibration}, formed from the numbers of (Th1). For $K\ge K_\phi$, $\mathsf b'\in[\mathsf b_-,\mathsf b_+]$, $u=\lceil\phi K\rceil$ and $t=K-u$ the setting (Th) holds with $(K,\mathsf b',m,t)$, and the running numbers $\theta_k=K^{\rm d}K_\Lambda\mathsf s^{\rm abs}_k$ of the setting (equations~\YMMref{(2176)--(2177)}{803}) satisfy $L\theta_k\le\Theta_KL^{-(K-k)/2}$ and $\prod_k(1+L\theta_k)\le C_{\rm fg}$, $\Theta_K=\max\{1,\frac{13}{72}(L-1)(1+Kc_{\log})^{\mathfrak p_0}\}$. Let $\mathfrak J_1,\dots,\mathfrak J_q$ be local observables of orders $\ge2$.
\begin{enumerate}
\item[(a)] \emph{(The chain with observables, at every source.)} For one observable $\mathfrak J$ put 
\[
\mathsf m_j=8NM^3c\,\alpha_0^{m}|\varkappa|_\infty L^{-(2m-3)j},
\]
 the size of $\mathfrak J$ as a pointed family at the level $j\ge1$, $\mathsf n^{\rm f}_{\mathfrak J}$ its size as a family in the fourth place at the height zero, 
\[
\mathsf d_1=K^{\rm d}\mathsf n^{\rm f}_{\mathfrak J},
\]
 $\mathsf n_j=\sum_{i\le j}L^{-(j-i)}\mathsf d_i$, $\mathsf d_{j+1}=\theta_j(\mathsf n_j+\mathsf m_j)$, and 
\[
r_{\mathfrak J}=\min\{w_\star/(32\mathsf n^{\rm f}_{\mathfrak J}),E_0/(16(\mathsf d_1+\mathsf m_1))\},
\]
 with the numbers $K^{\rm d}$ and $w_\star$ of \eqref{eq:source-weight-room} and $E_0$ of (Th1). There are numbers $\mathsf c_i(w,s)$ and functions $\Phi^{w,s}_k$ on $\mathsf{Cut}_k$, defined by a recursion that carries $s\mathfrak J$ as one more early member, such that for every source $w$, complex sources included, every $|s|<r_{\mathfrak J}$ and $1\le k\le t$,
\[(\mathcal M^k)_*\bigl(p_k\,e^{s\mathfrak J}\rho^w_0\,dU\bigr)=\exp\Bigl[\sum_{i\le k}\mathsf c_i(w,s)\Bigr]\mathbf 1_{\mathsf{Cut}_k}e^{\Phi^{w,s}_k}dV,\]
a bounded complex Borel density; the families $E^{(i)}[w,s]$ of the recursion are systems of early members with $\|E^{(i)}[w,s]\|<\frac{15}{16}E_0+\mathbf 1_{\{i=1\}}|s|\mathsf d_1$, Lipschitz in $s$ with the numbers $\mathsf d_i$, local in $w$ and in $\mathfrak J$, and holomorphic in $(w,s)$ and in the fields. The same holds for several observables, $s\mathfrak J$ replaced by $\sum_as_a\mathfrak J_a$ with $\sum_a|s_a|/r_a<1$, holomorphically in all variables. At the source $w=0$, $|\Phi^s_t-\Phi^{s'}_t|\le\sum_a|s_a-s'_a|\,\mathsf o_t(\mathfrak J_a)$ on $\mathsf{Cut}_t$, with $\mathsf o_t(\mathfrak J)\le CK^{2\mathfrak p_0}L^{4u-K}|\varkappa|_\infty$ at order $2$ and $\mathsf o_t(\mathfrak J)\le CK^{3\mathfrak p_0}L^{4u-3K/2}|\varkappa|_\infty$ at higher orders, $C$ a function of the constants, $m$, $\mathsf b_+$ and the type.
\item[(b)] $\mathsf Q_{t\leftarrow0}\mathfrak J$ differs from a constant by at most $\mathsf o_t(\mathfrak J)$ on $\mathsf{Cut}_t$.
\item[(c)] \emph{(Cumulants from oscillations.)} If $2\le q\le8$ and $|\mathsf Q_{t\leftarrow0}\mathfrak J_a-\mathsf a_a|\le\epsilon_a$ on $\mathsf{Cut}_t$, the joint cumulant of $\mathfrak J_1,\dots,\mathfrak J_q$ under $\langle\cdot\rangle_{K,\ell,L^K\mathsf b'}$ is at most $c_q\prod_a\epsilon_a+4\,\mathsf{op}_q\mathsf R_t\prod_a\sup|\mathfrak J_a|$ plus a term that is exponentially small in the distance of the carriers in units of the level $t$; $c_q=1+3.1^q\mathsf{op}_q\operatorname{Bl}_q$ and $\mathsf{op}_q$ are pure numbers, $\operatorname{Bl}_q$ the number of partitions of $q$ elements, and $\mathsf R_t\le C_{\rm R}L^{3u}e^{-c'A_0^2(\phi K\log L)^{2\mathfrak p_0}}$ is the large-field weight of the frame, smaller than every power of $L^{-K}$.
\item[(d)] A lattice isometry with diagonal linear part carries a local observable of type $(m,c,\hat c,N)$ to one of type $(m,64^m\max\{c,\hat c\},\hat c,27N)$.
\end{enumerate}
\end{lemma}
\begin{proof}
The frame: (Th3) of Section~\ref{sec:cutoff-calibration} asks an integer $t$ with $1\le t\le K-1$ and $x_t\ge\mathsf t_\natural$, and $x_t=L^u\mathsf b'\ge L^{\phi K}\gamma^{-2}$; the bounds on $\theta_k$ are those of equations~\YMMref{(2176)--(2177)}{803}, with $1+x\le e^x$. (a) An observable of order $m$ is a pointed family at the heights $\ge1$, of the size $\mathsf m_j$ at the level $j$, and a family in the fourth place at the height zero. The term $s\mathfrak J$ is carried as one more early member: the flow from a level exists for every system of early members (Proposition~\YMMref{114.2}{801}), the one-step theorem holds for every family in the fourth place of size $\le w_\star/4$ (Proposition~\ref{prop:source-entry}), one step is an identity at every point of the next cut (Lemma~\ref{lem:sf-chart-integral}), and $|s|<r_{\mathfrak J}$ keeps all norms in the domains.

At the first step the family in the fourth place is the one of the source, less its weight, plus $s$ times that of $\mathfrak J$. The Lipschitz numbers in $s$ propagate by $\mathsf d_{j+1}=\theta_j(\mathsf n_j+\mathsf m_j)$, by the Lipschitz form of forgetting (the proof of Proposition~\YMMref{114.3}{803}), and their contribution at $t$ is $\mathsf o_t$; locality in the observable gives several observables at once. (b) is the derivative of (a) at $s=0$, with Cauchy's estimate. (c) By polarisation the cumulant under the small-field functional is an alternating sum of derivatives of the logarithms of total masses; the exponents split into parts of single observables and terms that live on polymers joining two carriers, exponentially small in their distance, and the main term is a cumulant under a probability measure of the functions $\mathsf Q_{t\leftarrow0}\mathfrak J_a-\mathsf a_a$; the Wilson measure differs from the small-field functional by the weight $\mathsf R_t$. (d) The image of a box of at most $3\times3\times3$ sites is such a box with another lowest site, so that a site is the lowest site of at most $27N$ image boxes; the plaquette variables at the new base points are conjugate to the old ones by products of at most three bond matrices of norm $\le2$, so they lie within $64\delta$ of $\mathbf 1$ when the old ones lie within $\delta$, and where $64\delta>1$ the bound $\hat c$ serves.
\end{proof}

\begin{lemma}\label{cov-cubic}
For a field $\Xi$ on $T_\ell$, three real functions on the barycentres of the plaquettes, let $\mathsf H(\Xi)$ be the largest of their suprema plus the largest of their Lipschitz constants in the sup distance of the torus, and
\[\mathcal S[\Xi]=\sum_pL^{2K}\,\Xi_{\lambda(p)}(x(p))\,\mathfrak r^\sharp_{\lambda(p)}(p),\qquad\mathfrak r^\sharp_\lambda(p)=\mathfrak r_\lambda(p)-\Psi_\lambda(p)+\Psi_\lambda(p-a_Ke_\lambda),\]
with the cubic density $\mathfrak r_\lambda$ of the lattice Ward identity (Proposition~\YMMref{112.1}{782}) and the orbital potential $\Psi_\lambda$ of Section~\ref{sec:ward-orbital-complement}. Let $0<\phi<\frac1{42}$. There are $C_{\rm HR}$, $p_{\rm HR}$ and $K_{\rm HR}$, functions of the constants, $m$, $\mathsf b_+$ and $\phi$, such that for all integers $K\ge K_{\rm HR}$, all $\mathsf b'\in[\gamma^{-2},\mathsf b_+]$, $t=K-\lceil\phi K\rceil$, and every field $\Xi$ there is a number $\mathsf a$ with
\[\bigl|(\mathsf Q_{t\leftarrow0}\mathcal S[\Xi])(V)-\mathsf a\bigr|\le C_{\rm HR}\mathsf H(\Xi)K^{p_{\rm HR}}L^{-(1/14-3\phi)K}\qquad\text{for $\nu_t$-almost every }V\in\mathsf{Cut}_t ;\]
for $0<\phi<\frac1{30}$ the same holds with $\frac1{10}$ in the place of $\frac1{14}$. Both rates concern the conditional mean $\mathsf Q_{t\leftarrow0}$ of Definition~\ref{cov-loc}, which is formed at the source zero.
\end{lemma}
\begin{proof}
Apply Proposition~\ref{stc:cubic} to the same zero-source chain with $\eta=\Xi/\mathsf b'$, so that $\mathcal S[\mathsf b'\eta]=\mathcal S[\Xi]$ and $\mathsf H(\eta)\le\gamma^2\mathsf H(\Xi)$. The hypotheses on the torus and the label interval are the standing hypotheses here. With $(\psi,\kappa_\star)=(1/7,3/4)$ choose $\tau=\min\{1/4,21\phi\}$. The rate in that proposition is then $\min\{1/14-\tau/7,1/14-3\phi\}=1/14-3\phi$.

For the second rate take $(\psi,\kappa_\star)=(1/5,9/10)$ and $\tau=\min\{1/4,15\phi\}$, giving $1/10-\max\{\tau/5,3\phi\}=1/10-3\phi$. The choices of $\tau$ explain the permitted $\phi$ dependence of the constants and of the power of $K$. The localized inversion split, the H\"older gain and the stopping of the carried families are given in the proof of Proposition~\ref{stc:cubic}, using Lemmas~\ref{stc:mirror} and~\ref{stc:births}.


\end{proof}

\begin{remark}\label{cov-rem-oneloop}
The cubic density has a Gaussian illustration at fixed parameters. Fix a finite torus, a regulator $\mathsf m>0$, and a gauge fixing of the following class: $\mathsf D_{\bar U}$ and $\mathsf G_{\bar U}\ge0$ are real analytic in the background, $\mathsf D_{\mathbf1}=d$, and $\mathsf G_{\mathbf1}$ is a convolution with a real even function, the same for the three colours. For a real background direction $\bar A$ put $\bar U_\lambda=e^{i\lambda\bar A}$ and
\[\langle A,\mathsf H_{\bar U}A\rangle=\mathrm{Hs}_{\bar U}(A)+\langle\mathsf D_{\bar U}^*A,\mathsf G_{\bar U}\mathsf D_{\bar U}^*A\rangle+\mathsf m^2\|A\|^2,\qquad \mathrm{Hs}_{\bar U}(A)=\left.\frac{d^2}{d\mathsf g^2}\right|_{\mathsf g=0}\sum_qa_q(e^{i\mathsf gA}\bar U).\]
At $\lambda=0$ this form is positive definite, and for sufficiently small real $\lambda$ it defines a centred Gaussian with covariance $C_\lambda=\mathsf H_{\bar U_\lambda}^{-1}$. The size of this neighbourhood may depend on all the fixed parameters. Here $\mathsf g=\beta^{-1/2}$ serves only as an expansion parameter.

Let $T_p(\lambda)$ be the coefficient of $\mathsf g^2$ in the Gaussian mean of $\mathfrak r^\sharp_{\lambda(p)}(p)(e^{i\mathsf gA}\bar U_\lambda)$, and let $\Phi_p^{(r,s)}$ be its density term of degree $r$ in $A$ and degree $s$ in $\bar A$. The cancellation proved in Proposition~\ref{prop:ward-orbital-cubic} shows that the density has no term of total degree at most two. Its cubic term is $\frac14\sum_{\sigma=\pm}\sigma\sum_{j\in\mathbb Z_4}(\mathsf z_j^\sigma\times\mathsf z_{j+1}^\sigma)\cdot\mathsf x$, where $\mathsf x$ is the field strength at $p$ and $\mathsf z_j^\sigma$ are the legs at the four side faces above and below $p$.

Write $C_\lambda=C_0+\lambda C_1+O(\lambda^2)$, with $C_1=-C_0\mathsf H_1C_0$ and $\mathsf H_1=\left.\partial_\lambda\mathsf H_{\bar U_\lambda}\right|_0$. The matrix $C_0$ is diagonal in the colours. Thus $\Phi_p^{(2,0)}=0$, and every term of $\langle\Phi_p^{(2,1)}\rangle_{C_0}$ contracts an antisymmetric colour tensor with that diagonal matrix, so it is zero. Consequently

\[
T_p(\lambda)=\lambda^2\mathsf q_p[\bar A]+O(\lambda^3),\qquad \mathsf q_p=\langle\Phi_p^{(2,1)}\rangle_{C_1}+\langle\Phi_p^{(2,2)}\rangle_{C_0}.
\]
Brackets with $C_1$ denote linear Gaussian contraction with that symmetric matrix; $C_1$ need not be positive. All parameters except $\lambda$ are fixed in this expansion.



\end{remark}

\begin{lemma}\label{cov-var}
For all $K$, all $\beta>0$ and every $\eta\in C^2_c(\mathbb R^3;\mathbb R^3)$, on $\mathbb L_K$ (and, with any $C^2$ vector field, on a torus), identically in $U$,
\begin{equation}\label{cov-eq-var}
b\,\mathcal T_K(H^K_\eta)=\mathcal Z_\eta+X''[\eta]+\mathfrak Y[\eta],
\end{equation}
where the signs are the following. \emph{Lattice differences} of a function $\varphi$, with $a=a_K$: the centred one $(\nabla^{\rm c}_\rho\varphi)(y)=a^{-1}[\varphi(y+\frac a2e_\rho)-\varphi(y-\frac a2e_\rho)]$, the symmetric one $(\nabla^{\rm s}_\rho\varphi)(y)=(2a)^{-1}[\varphi(y+ae_\rho)-\varphi(y-ae_\rho)]$, the forward and the backward ones $(\nabla^\pm_\rho\varphi)(y)=\pm a^{-1}[\varphi(y\pm ae_\rho)-\varphi(y)]$, and $(\delta^2_\nu\varphi)(y)=\varphi(y+\frac a2e_\nu)+\varphi(y-\frac a2e_\nu)-2\varphi(y)$; $\nabla^{\rm s}\!\cdot\eta=\sum_\kappa\nabla^{\rm s}_\kappa\eta^\kappa$. \emph{The lattice tensor.} $H^K_\eta$ is $H_\eta$ with the entries that $\mathcal T_K$ reads replaced by lattice ones: at the barycentre of a plaquette $p$ with the normal $\lambda=\lambda(p)$ the diagonal entry $(\nabla^{\rm s}\!\cdot\eta-2\nabla^{\rm s}_\lambda\eta^\lambda)(x(p))$, and at the midpoint of a bond $b$ the off-diagonal entry $-(\nabla^{\rm c}_{\lambda'}\eta^\lambda+\nabla^{\rm c}_\lambda\eta^{\lambda'})(x_b)$, $\{\lambda,\lambda'\}=\{\lambda_b,\lambda'_b\}$; as $|\nabla^{\rm s}_\rho\varphi-\partial_\rho\varphi|\le\frac a2\sup|\partial^2_\rho\varphi|$ and $|\nabla^{\rm c}_\rho\varphi-\partial_\rho\varphi|\le\frac a4\sup|\partial^2_\rho\varphi|$, $|H^K_\eta-H_\eta|_\infty\le\frac32a_K\|\nabla^2\eta\|_\infty$.

\emph{The null terms.} For a field of directions $\mathbf X$, $\mathfrak v_{\mathbf X}$ is the derivation along it and its null term is $\beta\mathfrak v_{\mathbf X}\sum_pa_p-\operatorname{div}\mathfrak v_{\mathbf X}$; $\mathcal N_\eta$ is the null term of $\mathbf X_\eta(b)=a^{-1}\sum_\rho\eta^\rho(x_b)\mathbf u_{b,\rho}$, and $\mathcal N_{\mathsf k}$, for a vector field $\mathsf k$, that of $\mathbf X^{\mathsf k}(b)=\sum_\rho\mathsf k^\rho(x_b)\mathbf y_{b,\rho}$, with $\mathbf u_{b,\mu_b}=\mathbf y_{b,\mu_b}:=0$. With $\mathsf k^\rho_\eta:=\nabla^+_\rho\eta^\rho$ and the number $c_{\mathsf k}:=6\sum_b\sum_{\rho\ne\mu_b}\mathsf k^\rho(x_b)$, free of $U$, $\mathcal Z_\eta=\mathcal N_\eta-\frac14\mathcal N_{\mathsf k_\eta}-\frac14c_{\mathsf k_\eta}$ is the number $-\frac14c_{\mathsf k_\eta}$ plus null terms.

\emph{The nine observables.} 
\[
X''[\eta]=\sum_{v=1}^9\mathfrak J_v[\eta],
\]
 a sum of local observables of the orders $3,2,2,2,2,2,3,4,4$:
\[
\begin{multlined}
\mathfrak J_1=-\mathcal S[b\eta],\quad(\mathfrak J_2,\mathfrak J_3,\mathfrak J_4,\mathfrak J_7,\mathfrak J_8)=(\mathfrak J_{0,1},\mathfrak J_{0,2},\mathfrak J_{0,3},\mathfrak J_{0,4},\mathfrak J_{0,5})[\mathsf k_\eta],\quad\mathfrak J_5=-\tfrac32a\,\mathcal O(w_\eta),\quad\mathfrak J_6= \\
-E_2[\eta],\quad\mathfrak J_9=-E_4[\eta],
\end{multlined}
\]
where $w_\eta(p)=\sum_{\rho\in\mathsf P(p)}(\nabla^{\rm c}_\rho\eta^\rho)(x(p))$, $\mathsf P(p)$ the two directions of the plane of $p$;
\[
\begin{multlined}
E_4[\eta]=-\frac\beta2\sum_b\sum_{\nu\ne\mu_b}(\nabla^{\rm c}_\nu\eta^\nu)(x_b)\bigl(a_{p(b,\nu)}^2+a_{p(b,-\nu)}^2\bigr),\qquad E_2[\eta]= \\
-\frac\beta{2a}\sum_b\sum_{\rho,\nu\ne\mu_b}(\delta^2_\nu\eta^\rho)(x_b)\,\langle\mathbf u_{b,\rho},\mathbf y_{b,\nu}\rangle ,
\end{multlined}
\]
$p(b,\pm\nu)$ the plaquette of the loop $W_x(\mu_b,\pm\nu)$, $b=\langle x,x+ae_{\mu_b}\rangle$; and, for a vector field $\mathsf k$, the five pieces of the orbital bilinear
\[
\begin{aligned}
\mathfrak J_{0,1}[\mathsf k]
&=\sum_p\Bigl[\tfrac34\mathsf e^{\mathsf k}_{\mathsf P(p)}(x(p))
-\tfrac14\beta a^2\sum_\rho(\nabla^-_\rho\nabla^+_\rho\mathsf k^\rho)(x(p))\Bigr]a_p,\\
\mathfrak J_{0,2}[\mathsf k]
&=\tfrac\beta4\sum_{p,\rho}\sum_{b\subset\partial p}
\bigl(\mathsf k^\rho(x_b)-\mathsf k^\rho(x(p))\bigr)\\
&\qquad\langle\mathbf y_{b,\rho},\mathbf w(W_{b,p})\rangle,\\
\mathfrak J_{0,3}[\mathsf k]
&=-\tfrac{\beta a}2\sum_{p,\rho}(\nabla^+_\rho\mathsf k^\rho)(x(p))\,\Psi_\rho(p),\\
\mathfrak J_{0,4}[\mathsf k]
&=\tfrac\beta4\sum_p\mathsf k^{\lambda(p)}(x(p))\,\mathfrak e_{\lambda(p)}(p),\\
\mathfrak J_{0,5}[\mathsf k]
&=-\tfrac\beta2\sum_p\Bigl(\sum_{\rho\in\mathsf P(p)}\mathsf k^\rho(x(p))\Bigr)a_p^2,
\end{aligned}
\]
with $\mathsf e^{\mathsf k}_{\mathsf P}(y)=\sum_{\rho\in\mathsf P}[\mathsf k^\rho(y+\frac a2e_\rho)+\mathsf k^\rho(y-\frac a2e_\rho)]$, $W_{b,p}$ the loop of $p$ that starts along $b$ ($W_x(\mu_b,\pm\nu)$ at $p=p(b,\pm\nu)$), the orbital potential $\Psi_\rho(p)$ of Section~\ref{sec:ward-orbital-complement}, $p_\pm=p\pm ae_\rho$, and the density 
\[
\mathfrak e_\rho(p):=\sum_{b\subset\partial p}\langle\mathbf y_{b,\rho},\mathbf w(W_{b,p})\rangle+2[\Psi_\rho(p)+\Psi_\rho(p_-)]+[a_{p_+}+a_{p_-}-2a_p],
\]
which is of the order $3$ at the normal direction. When $b=a_K\beta\le\mathsf b_+$ the weights of the nine observables are at most $(3+5\mathsf b_+)L^{\omega K}\|\eta\|_{C^2}$, $\omega=2,0,0,0,0,0,1,1,1$, and every site of a box with a weight $\ne0$ lies within $4a_K$ of $\operatorname{supp}\eta$. The \emph{orbital remainder} of the lattice Ward identity is $X'[\eta]:=X''[\eta]+E_4[\eta]+E_2[\eta]$, the sum of the seven observables other than $\mathfrak J_6$ and $\mathfrak J_9$. And
\[\mathfrak Y[\kappa]=\frac\beta4\sum_b\sum_{\nu\ne\mu_b}\kappa_{\mu_b\nu}(x_b)\,|\mathbf y_{b,\nu}|^2,\qquad\mathfrak Y[\eta]=\mathfrak Y[\kappa]\text{ at }\kappa_{\mu\nu}=\nabla^{\rm c}_\nu\eta^\nu .\]
For every Gibbs state $\nu$ at $(a_K,\beta)$ and every torus measure, $\nu(\mathcal NY)=0$ for every bounded Borel $Y$ that does not depend on the bonds at which the field of directions of $\mathcal N$ is not zero; so every joint cumulant with the member $\mathcal N$ and further members that do not depend on those bonds vanishes. The coefficient of $\mathcal T_K$ in \eqref{cov-eq-var} is exactly $b$.
\end{lemma}
\begin{proof}
All the identities are identities between polynomials in the matrices at finitely many bonds; they are proved on a torus, and they hold on $\mathbb L_K$ for a field of compact support, whose bonds may be placed on a torus without wrapping. Write $A=\sum_pa_p$, $\mathfrak d_{b,p,\rho}=\langle\mathbf u_{b,\rho},\mathbf w(W_{b,p})\rangle$, and recall $|\mathbf w(W)|^2=2a_q-a_q^2$ for a loop $W$ of a plaquette $q$: for $W=c\mathbf 1+is\mathsf n$ with $c^2+s^2=1$ and a unit vector $\mathsf n$ of the Lie algebra, $\mathbf w(W)=s\mathsf n$ and $a_q=1-c$.

\emph{1. The variation along $\eta$.} The lattice calculus of Lemma~\ref{lem:ward-general-calculus}, with the general field \eqref{eq:ward-general-displacement} at the constant weight $\beta$ and the divergence \eqref{eq:ward-general-divergence}, gives $\beta\mathfrak v_{\mathbf X_\eta}A=\mathfrak K^{\rm v}[\beta\eta]-b\,\mathcal O(\nabla^{\rm s}\!\cdot\eta)+\mathsf\Theta_\beta[\eta]$, with $\mathsf\Theta_\beta[\eta]=\beta\sum_pa^{-1}\sum_{b\subset\partial p}\sum_\rho(\eta^\rho(x_b)-\eta^\rho(x(p)))\,\mathfrak d_{b,p,\rho}$, and $\operatorname{div}\mathfrak v_{\mathbf X_\eta}=-\frac32a\,\mathcal O(w_\eta)$. So the \emph{stress insertion} $\mathsf{St}_K[\eta]:=\mathsf\Theta_\beta[\eta]-b\,\mathcal O(\nabla^{\rm s}\!\cdot\eta)$ is $\mathcal N_\eta-\mathfrak K^{\rm v}[\beta\eta]-\frac32a\,\mathcal O(w_\eta)$. The summation by parts \eqref{eq:ward-orbital-summation} at the vector weight $\beta\eta$, with $\beta a^{-1}=bL^{2K}$, gives $\mathfrak K^{\rm v}[\beta\eta]=\mathcal S[b\eta]-\mathfrak J^{\rm II}[\mathsf k_\eta]$, where
\[\mathfrak J^{\rm II}[\mathsf k]:=\beta\sum_p\sum_\rho\mathsf k^\rho(x(p))\,\Psi_\rho(p)\]
is the \emph{orbital bilinear}: an observable of the order $2$ with the weight $\beta=L^Kb$, which the power counting of Lemma~\ref{cov-chain} does not make small.

\emph{2. The orbital bilinear is a null term up to observables that power counting bears.} For every vector field $\mathsf k$, identically in $U$,
\[\mathfrak J^{\rm II}[\mathsf k]=-\tfrac14\mathcal N_{\mathsf k}+\sum_{v=1}^5\mathfrak J_{0,v}[\mathsf k]-\tfrac14c_{\mathsf k}.\]
By Lemma~\ref{lem:ward-general-calculus}, $\partial^{\mathsf a}_ba_p=\langle\sigma_{\mathsf a},\mathbf w(W_{b,p})\rangle$ for $p\ni b$ and $\partial^{\mathsf a}_b\langle\sigma_{\mathsf a},\mathbf w(W_{b,p})\rangle=1-a_p$; so $\mathfrak v_{\mathbf X^{\mathsf k}}a_p=\sum_{b\subset\partial p}\langle\mathbf X^{\mathsf k}(b),\mathbf w(W_{b,p})\rangle$ and
\[\operatorname{div}\mathfrak v_{\mathbf X^{\mathsf k}}=3\sum_b\sum_{\rho\ne\mu_b}\mathsf k^\rho(x_b)\bigl[2-a_{p(b,\rho)}-a_{p(b,-\rho)}\bigr]=c_{\mathsf k}-3\sum_pa_p\,\mathsf e^{\mathsf k}_{\mathsf P(p)}(x(p)),\]
the last sum ordered by the plaquettes: $p=p(b,\mp\rho)$ exactly for $\rho\in\mathsf P(p)$ and the bond $b$ of $p$ with $x_b=x(p)\pm\frac a2e_\rho$. In $\mathfrak v_{\mathbf X^{\mathsf k}}A$ write $\mathsf k^\rho(x_b)=\mathsf k^\rho(x(p))+(\mathsf k^\rho(x_b)-\mathsf k^\rho(x(p)))$; the first part is $\sum_{p,\rho}\mathsf k^\rho(x(p))\{-2[\Psi_\rho(p)+\Psi_\rho(p_-)]-[a_{p_+}+a_{p_-}-2a_p]+\mathfrak e_\rho(p)\}$, by the definition of $\mathfrak e_\rho$. If $\rho$ lies in the plane of $p$, then $\mathfrak e_\rho(p)=-2a_p^2$: at the two bonds $b_1$, $b_2$ of $p$ of the other direction, $\langle\mathbf w+\mathbf w',\mathbf w\rangle=\frac12|\mathbf w+\mathbf w'|^2+\frac12|\mathbf w|^2-\frac12|\mathbf w'|^2$ with $\mathbf w=\mathbf w(W_{b,p})$, $|\mathbf w(W)|^2=2a_q-a_q^2$, and the first form of \eqref{eq:ward-orbital-inplane}, $|\mathbf y_{b_2,\rho}|^2=-4\Psi_\rho(p)-a_p^2-a_{p_+}^2$ and $|\mathbf y_{b_1,\rho}|^2=-4\Psi_\rho(p_-)-a_{p_-}^2-a_p^2$, give $\sum_{b\subset\partial p}\langle\mathbf y_{b,\rho},\mathbf w(W_{b,p})\rangle=-2[\Psi_\rho(p)+\Psi_\rho(p_-)]-[a_{p_+}+a_{p_-}-2a_p]-2a_p^2$.

Shifting $p\mapsto p\pm ae_\rho$, $\sum_p\mathsf k^\rho(x(p))\Psi_\rho(p_-)=\sum_p\mathsf k^\rho(x(p)+ae_\rho)\Psi_\rho(p)$ and $\sum_p\mathsf k^\rho(x(p))[a_{p_+}+a_{p_-}-2a_p]=a^2\sum_p(\nabla^-_\rho\nabla^+_\rho\mathsf k^\rho)(x(p))\,a_p$. Solving for $\sum_{p,\rho}\mathsf k^\rho(x(p))\Psi_\rho(p)$, multiplying by $\beta=ba^{-1}$, writing $\beta\mathfrak v_{\mathbf X^{\mathsf k}}A=\mathcal N_{\mathsf k}+\operatorname{div}\mathfrak v_{\mathbf X^{\mathsf k}}$ and using the value of $\mathfrak e_\rho$ in the plane gives the identity with the five pieces of the statement. Their densities $a_p$, $\langle\mathbf y_{b,\rho},\mathbf w(W_{b,p})\rangle$, $\Psi_\rho(p)$ are of the order $2$, $\mathfrak e_{\lambda(p)}(p)$ is of the order $3$ (its part of the total order at most two in the logarithms of the bond variables cancels) and $a_p^2$ is of the order $4$, each on a box within $\frac32a$ of $x(p)$; their weights are at most $3\|\mathsf k\|_\infty+\frac32b\operatorname{Lip}_\infty\mathsf k$, $\frac18b\operatorname{Lip}_\infty\mathsf k$, $\frac12b\operatorname{Lip}_\infty\mathsf k$, $\frac14\beta\|\mathsf k\|_\infty$, $\beta\|\mathsf k\|_\infty$, and a weight at an index of $p$ vanishes unless $x(p)$ is within $a$ of $\operatorname{supp}\mathsf k$.

\emph{3. Two forms of the stress insertion.} Steps 1 and 2 give $\mathsf{St}_K[\eta]=\mathcal Z_\eta+X'[\eta]$. On the other hand, order $\mathsf\Theta_\beta$ by the bonds: $b$ lies in the four plaquettes $p(b,\pm\nu)$, $\nu\ne\mu_b$, with $x(p(b,\pm\nu))=x_b\pm\frac a2e_\nu$, and the vector parts $\mathbf w_\pm=\mathbf w(W_x(\mu_b,\pm\nu))$ are $\frac12\mathbf y_{b,\nu}\mp\mathbf u_{b,\nu}$; so
\[\mathsf\Theta_\beta[\eta]=\beta\sum_b\sum_{\rho,\nu\ne\mu_b}(\nabla^{\rm c}_\nu\eta^\rho)(x_b)\langle\mathbf u_{b,\rho},\mathbf u_{b,\nu}\rangle+E_2[\eta].\]
The terms with $\rho\ne\nu$ are $b\,Y[\mathsf w]$ at the off-diagonal entries of $H^K_\eta$, because $\langle\mathbf u_{b,\lambda_b},\mathbf u_{b,\lambda'_b}\rangle=q_b$. In the terms with $\rho=\nu$ the parallelogram law and $|\mathbf w(W)|^2=2a_q-a_q^2$ give $|\mathbf u_{b,\nu}|^2=(a_{p(b,\nu)}+a_{p(b,-\nu)})-\frac12(a_{p(b,\nu)}^2+a_{p(b,-\nu)}^2)-\frac14|\mathbf y_{b,\nu}|^2$; the second and the third part give $E_4[\eta]$ and $-\mathfrak Y[\eta]$, and the first, ordered by the plaquettes with $(\nabla^{\rm c}_\nu\eta^\nu)(x(p)\pm\frac a2e_\nu)=(\nabla^\pm_\nu\eta^\nu)(x(p))$ and $\nabla^++\nabla^-=2\nabla^{\rm s}$, is $b\,\mathcal O(w)$ with $w(p)=2\sum_{\nu\in\mathsf P(p)}(\nabla^{\rm s}_\nu\eta^\nu)(x(p))$; less $b\,\mathcal O(\nabla^{\rm s}\!\cdot\eta)$ this is $b\,\mathcal O$ at the diagonal entries of $H^K_\eta$. So $\mathsf{St}_K[\eta]=b\,\mathcal T_K(H^K_\eta)-\mathfrak Y[\eta]+E_4[\eta]+E_2[\eta]$.

\emph{4.} The two forms together give \eqref{cov-eq-var}, with $X''=X'-E_4-E_2$; the coefficient $b$ of $\mathcal T_K$ is that of the variation of the action itself. The orders and the weights of $\mathfrak J_1$, $\mathfrak J_5$, $E_2$, $E_4$ are read off their formulas ($|\delta^2_\nu\varphi|\le\frac{a^2}4\sup|\partial^2_\nu\varphi|$; $\langle\mathbf u_{b,\rho},\mathbf y_{b,\nu}\rangle$ is a density of the order $2$), and those of the five pieces at $\mathsf k_\eta$ follow from $\|\mathsf k_\eta\|_\infty\le\|\nabla\eta\|_\infty$, $\operatorname{Lip}_\infty\mathsf k_\eta\le3\|\nabla^2\eta\|_\infty$ and $b\le\mathsf b_+$. \emph{The null property.} Let the field of directions $\mathbf X$ of $\mathcal N$ vanish off a finite set $\Lambda$ of bonds, and let $Y$ not read the bonds of $\Lambda$: invariance of Haar measure (Lemma~\ref{lem:ward-general-calculus}) gives $\int(\operatorname{div}\mathfrak v_{\mathbf X}-\beta\mathfrak v_{\mathbf X}A)\,Y\,e^{-\beta\sum_{p\cap\Lambda\ne\emptyset}a_p}\prod_{b\in\Lambda}dU(b)=0$, so the Gibbs kernel of $\Lambda$ annihilates $\mathcal NY$, and the DLR equations give $\nu(\mathcal NY)=0$.
\end{proof}

\begin{proposition}\label{cov-eight}
Let $r>0$, and consider on the least torus $T_{\ell'}$ a board with cells of side $\ell'/2$; let $t_0=K-\lceil K/400\rceil$. There are $C$, $p$, $K_8$, functions of the constants, $g$, $r$ and the rate $c$ below, such that for all $(K,\beta)\in\mathfrak d(g)$ with $K\ge K_8$, with $\langle\cdot\rangle=\langle\cdot\rangle_{K,\ell',\beta}$, and for objects carried by one cell at sup distance $\ge r$ from its walls:
\begin{enumerate}
\item[(a)] $\|X''[\eta]-\langle X''[\eta]\rangle\|_{{\rm ch},\ell'}\le CK^pL^{-cK}\|\eta\|_{C^2}$ for every $c<\frac1{10}$, at every offset of the board in $a_K\mathbb Z^3$; the same for the orbital remainder $X'[\eta]=X''[\eta]+E_4[\eta]+E_2[\eta]$ of the lattice Ward identity, the sum of the seven pieces of $X''[\eta]$ other than $\mathfrak J_6=-E_2[\eta]$ and $\mathfrak J_9=-E_4[\eta]$ (Lemma~\ref{cov-var}); and the joint cumulants of $2\le q'\le8$ of the nine pieces $\mathfrak J_v[\eta]$ in different cells are at most $(CK^pL^{-cK}\|\eta\|_{C^2})^{q'}$;
\item[(b)] $\|\mathfrak Y[\kappa]-\langle\mathfrak Y[\kappa]\rangle\|_{{\rm ch},\ell'}\le CK^pL^{-cK}\sup|\kappa|$ for every $c<\frac2{11}$ and six bounded real functions $\kappa_{\mu\nu}$;
\item[(c)] $\|\mathcal T_K(H')-\langle\mathcal T_K(H')\rangle\|_{{\rm ch},\ell'}\le CK^{2\mathfrak p_0}L^{K/100}|H'|_\infty$ for every bounded $H'$, at offsets in $a_{K-t_0}\mathbb Z^3$;
\item[(d)] $\|\mathcal T_K(H_\eta)-\langle\mathcal T_K(H_\eta)\rangle\|_{{\rm ch},\ell'}\le CK^pL^{-cK}\|\eta\|_{C^2}$ for every $c<\frac1{10}$, at offsets in $a_{K-t_0}\mathbb Z^3$.
\end{enumerate}
\end{proposition}
\begin{proof}
The eight copies of an observable carried by one cell are local observables: the mirror maps are diagonal lattice isometries, and Lemma~\ref{cov-chain}\,(d) gives the copies of a local observable of the type $(m,c,\hat c,N)$ the type $(m,64^m\max\{c,\hat c\},\hat c,27N)$. Their carriers are $\ge2(r-4a_K)\ge r$ apart once $8a_K\le r$, as the boxes of a piece reach at most $4a_K$ beyond the support of $\eta$ (Lemma~\ref{cov-var}), that is $\ge rL^{K-t}$ in units of the level $t$. The mean of the product of eight centred copies is a sum over the $715$ partitions of the eight copies into blocks of at least two elements of products of joint cumulants, which Lemma~\ref{cov-chain}\,(c) bounds by oscillations of conditional means, the large-field term and the term on polymers joining two carriers being negligible. (a) The cubic term is bounded by Lemma~\ref{cov-cubic} with $\Xi=b\eta$, $\mathsf H(b\eta)\le4\mathsf b_+\|\eta\|_{C^2}$, at a $\phi$ so small that $\frac1{10}-3\phi>c$, in each of the eight copies: for a diagonal lattice isometry $\mathsf r:x\mapsto\operatorname{diag}(s_1,s_2,s_3)x+v$, the covariance Proposition~\ref{prop:ward-orbital-cubic}(c) of $\mathfrak r^\sharp$ under the lattice isometries gives $\mathsf r\,\mathcal S[\Xi]=\mathcal S[\Xi^{\mathsf r}]$ with $\Xi^{\mathsf r}_\rho=s_\rho\,\Xi_\rho\circ\mathsf r^{-1}$, and $\mathsf H(\Xi^{\mathsf r})=\mathsf H(\Xi)$; so Lemma~\ref{cov-cubic}, which holds for every field, applies to each copy.

The other eight terms have orders $m\ge2$ and weights $\le L^{\omega K}$ with $\omega=1$ only for $m\ge3$, and Lemma~\ref{cov-chain}\,(a), (b), at the types of the copies, give oscillations $\le CK^{3\mathfrak p_0}L^{4u-K/2}$. (b) The mirror copies of $\mathfrak Y[\kappa]$ are forms of the same kind: for a diagonal lattice isometry $\mathsf r$ and $b^{\mathsf r}$ the positively oriented one of the bonds $\mathsf rb$, $\overline{\mathsf rb}$, $\mathsf r|\mathbf y_{b,\nu}|^2=|\mathbf y_{b^{\mathsf r},\nu}|^2$, by the covariance of the holonomies (Lemma~\ref{lem:setup-lattice-covariance}), so that $\mathsf r\,\mathfrak Y[\kappa]=\mathfrak Y[\kappa^{\mathsf r}]$ with a weight $\kappa^{\mathsf r}$ of the same sup norm, and the bound below holds for every bounded weight. $\mathfrak Y[\kappa]$ is an observable of order $2$ with weights of the size $\beta=L^Kb$; but at every $U$ of the minimal orbit over $V\in\mathsf{Cut}_n$, which is regular of level $n$ at the radius $\varepsilon_n$ of Subsection~\ref{cov-cut}, each $\mathbf y_{b,\nu}$, the sum of the vector parts of the two plaquettes opposite at $b$, is a lattice derivative of the field strength and is one power of the scale smaller than each part:
\[|\mathbf y_{b,\nu}(U)|\le48\,\mathsf c_{\rm g}(\beta_{\rm H})\rho_{\rm H}\,\varepsilon_nL^{-(2+\beta_{\rm H})n},\]
with the H\"older exponent $0\le\beta_{\rm H}<1$ and the numbers $\mathsf c_{\rm g}(\beta_{\rm H})$, $\rho_{\rm H}$, $a_{\rm g}$ of the local gauges on cubes (Theorem~\textup{\ref{I-56-65}}), once $\rho_{\rm H}\varepsilon_n\le a_{\rm g}$, $\mathsf c_{\rm g}(\beta_{\rm H})\rho_{\rm H}\varepsilon_n\le1$ and $6\rho_{\rm H}<L^{m+K-n}$, which hold for large $K$. Indeed, on the cube of the side $2\rho_{\rm H}$, in units of the level $n$, about a site there is a gauge in which $U=e^{iL^{-n}A}$ with $|A|,|\nabla A|<\mathsf c_{\rm g}\rho_{\rm H}\varepsilon_n$ and a H\"older bound of the exponent $\beta_{\rm H}$ and the number $\mathsf c_{\rm g}(\beta_{\rm H})\rho_{\rm H}\varepsilon_n$ on $\nabla A$; there the sum of the vector parts of the two loops opposite at $b$ is $L^{-2n}$ times the difference of the field strength at two neighbouring plaquettes, up to $O(\varepsilon_n^2L^{-3n})$.

Cutting the chain at the level $n_0=K-\lceil\psi K\rceil$, where $\varepsilon_{n_0}^2\le CK^{2\mathfrak p_0}L^{-\psi K}$, what is born below $n_0$ is born at a small coupling, and above $n_0$ the conditional mean is a mean of the value at $n_0$; the oscillation decays at the rate $\min\{\frac\psi2-4\phi,\frac12-4\phi,2\beta_{\rm H}-(3+2\beta_{\rm H})\psi\}$, which is $\frac1{14}$ at $(\beta_{\rm H},\psi,\phi)=(\frac12,\frac16,\frac1{336})$ and approaches $\frac2{11}$ as $\phi\to0$, $\beta_{\rm H}\to1$, $\psi=\frac4{11}$. (c) $\mathcal T_K(H')$ is an observable of order $2$ with weights $\le L^K|H'|_\infty$; Lemma~\ref{cov-chain} at $\phi=\frac1{400}$ gives the oscillation $CK^{2\mathfrak p_0}L^{4u}|H'|_\infty$, $4u\le\frac K{100}+4$. (d) Insert \eqref{cov-eq-var}.

In the mean of the eight centred copies every term that contains a null term of one copy vanishes, since the other copies do not depend on its bonds; so (a), (b) and (c) for $\mathcal T_K(H_\eta-H^K_\eta)$, whose sup norm is $\le\frac32a_K\|\nabla^2\eta\|_\infty$, give (d), as $b>g^{-2}$.
\end{proof}

The \emph{stress insertion} $\mathsf{St}_K[\eta]=\mathsf\Theta_\beta[\eta]-b\,\mathcal O(\nabla^{\rm s}\!\cdot\eta)$ of the lattice Ward identity (the proof of Lemma~\ref{cov-var}) is $\mathcal Z_\eta+X'[\eta]=b\,\mathcal T_K(H^K_\eta)-\mathfrak Y[\eta]+(X'[\eta]-X''[\eta])$; it differs from $b\,\mathcal T_K(H^K_\eta)$ by $X'[\eta]-X''[\eta]-\mathfrak Y[\eta]=E_4[\eta]+E_2[\eta]-\mathfrak Y[\eta]$. By Proposition~\ref{cov-eight}\,(a), (b) this difference has the eight-copy seminorm $\le CK^pL^{-cK}\|\eta\|_{C^2}$. With Proposition~\ref{cov-eight}\,(c) and the two-cell argument of the proof of Theorem~\ref{cov-thm}, and with $A^\nu[F]:=F-\nu(F)$: for $\mathsf d,R>0$ there are $C$, $p$, $K_0$, functions of the constants, $g$, $\mathsf d$ and $R$, such that for $K\ge K_0$ and in every periodic Gibbs state $|\nu(A^\nu[\mathsf{St}_K[\eta]-b\mathcal T_K(H_\eta)]\,A^\nu[Z])|\le CK^pL^{-K/60}\|\eta\|_{C^2}\mathsf z$, for $Z=\mathcal T_K(H')$, $\mathsf z=|H'|_\infty$, and for $Z$ a further stress insertion, $\mathsf z=\|\eta'\|_{C^2}$, with supports in a cube of the side $2R$ at sup distance $\ge\mathsf d$. So the limits of the correlations of the stress insertion are those of $b_*$ times the tensor at $H_\eta$: the coefficient is exactly one, and no other combination of the components enters.

\subsection{Abstract data}\label{cov-abs}

In this subsection no lattice occurs. We write $y=(y^0,y^1)=(x_3,x_1)$ for the plane of the rotations $\mathsf R_t$ and $x=(y,x_2)$; for $\xi\in\mathbb C^2$, $\xi^\pm=\xi^0\pm i\xi^1$, $(\mathsf R_\zeta\xi)^\pm=e^{\pm i\zeta}\xi^\pm$ for $\zeta\in\mathbb C$, $d(\xi)=\min\operatorname{Re}\xi^\pm$, $\mathsf{Tb}=\{d>0\}$, whose real points form the half plane $\mathbb R^2_+=\{y\in\mathbb R^2:y^0>0\}$, $\xi^*=(\overline{\xi^0},-\overline{\xi^1})$, $n_\alpha=(\cos\alpha,\sin\alpha)$, and $\mathsf T(\xi)=\int e^{-\xi^0E+i\xi^1q_1}d\mathcal P$, a contraction for $\xi\in\overline{\mathsf{Tb}}$, holomorphic in norm on $\mathsf{Tb}$, with $\mathsf T(\xi)\mathsf T(\xi')=\mathsf T(\xi+\xi')$, $\mathsf T(\xi)^*=\mathsf T(\xi^*)$, $\mathsf T(\xi)\Omega=\Omega$ and $\mathsf T(y)[F]=[\tau_yF]$ at real $y$ with $y^0\ge0$; for real $s$, $d(\mathsf R_{is}\xi)\ge e^{-|s|}d(\xi)$, with equality $d(\mathsf R_{is}y)=e^{-|s|}y^0$ at real $y$ with $y^0>0$. The \emph{weight basis}
\[\varepsilon_{\mp2}=u_\pm u_\pm^{\rm T},\quad\varepsilon_{\mp1}=u_\pm e_2^{\rm T}+e_2u_\pm^{\rm T},\quad\varepsilon_0=e_3e_3^{\rm T}+e_1e_1^{\rm T},\quad\varepsilon'_0=e_2e_2^{\rm T},\qquad u_\pm=e_3\pm ie_1,\]
of $V_{\mathbb C}$ has $\rho(\mathsf R_t)a=e^{ik(a)t}a$ with $k(\varepsilon_{\mp2})=\mp2$, $k(\varepsilon_{\mp1})=\mp1$, $k(\varepsilon_0)=k(\varepsilon'_0)=0$, and $a^\dagger:=\rho(\theta)\bar a=(-1)^{k(a)}a$ on it; $\mathbf1=\varepsilon_0+\varepsilon'_0$. A \emph{label} is $\lambda=(\varphi,a)$, $\varphi\in C^1_c(\mathbb R)$ real, $a\in V_{\mathbb C}$; it is a \emph{weight label} if $\pm a$ belongs to the weight basis, and then $k(\lambda)=k(a)$; a \emph{scalar label} is a label $(\varphi,\mathbf1)$, written $\varphi$, of weight $0$, with the point operators $\mathsf B_\varphi=\mathsf B_{(\varphi,\mathbf1)}$ of Lemma~\ref{cov-plate}; $(\alpha\otimes\lambda)(x)=\alpha(y)\varphi(x_2)a$ is a \emph{plate} with the profile $\alpha$; $\lambda^\dagger=(\varphi,\rho(\theta)\bar a)$, $\bar\lambda=(\varphi,\bar a)$, and $\kappa(\underline\lambda)=\sum_kk(\lambda_k)$.

At a weight label $\mathsf R_t(\alpha\otimes\lambda)=e^{ik(\lambda)t}(\alpha\circ\mathsf R_{-t})\otimes\lambda$. For a permutation $\mathfrak s$ and a \emph{frame} $\upsilon\in\frac\pi2\mathbb Z$ the \emph{tube} $\mathcal Z(\mathfrak s,\upsilon)$ is the set of $Z\in(\mathbb C^2)^n$ whose consecutive differences after the permutation, $\partial(\mathfrak sZ)$, lie in $(\mathsf R_\upsilon\mathsf{Tb})^{n-1}$, and $\mathcal Z_n$ is the union of the tubes. The tubes are open convex cones, invariant under common complex translations, under $\mathsf R_{is}$ for real $s$ and under conjugation; a real configuration lies in $\mathcal Z_n$ if and only if its times, or its abscissae, are pairwise different.

A \emph{fan} of $j$ points is a real $z$ with $z_1=r_0n_{\alpha_0}$, $z_{i+1}-z_i=r_in_{\alpha_i}$, $r_i>0$, $|\alpha_i|<\frac{5\pi}{16}$; a \emph{pencil} is a real configuration whose consecutive differences make angles $<\frac\pi{16}$ with $n_0$.

\begin{lemma}\label{cov-plate}
Let $\sigma$ be labelled data with (S).
\begin{enumerate}
\item[(a)] For every label $\lambda$ there is exactly one family of bounded operators $\mathsf B_\lambda(\xi,\xi')$, $(\xi,\xi')\in\mathsf{Tb}^2$, continuous in norm, with $\mathsf T(\kappa)\mathsf B_\lambda(\xi,\xi')\mathsf T(\kappa')=\mathsf B_\lambda(\kappa+\xi,\xi'+\kappa')$ for $\kappa,\kappa'\in\overline{\mathsf{Tb}}$ and $\mathsf B_l(\alpha\otimes\lambda)=\int\alpha(y)\mathsf B_\lambda(y,ln_0-y)\,dy$ for $\alpha\in C_c(\{0<y^0<l\})$. It is holomorphic in norm, complex linear in $a$, real linear in $\varphi$; $\mathsf B_\lambda(\xi,\xi')^*=\mathsf B_{\lambda^\dagger}(\xi'^*,\xi^*)$; $\|\mathsf B_\lambda(\xi,\xi')\|$ depends on $\operatorname{Re}\xi^\pm$, $\operatorname{Re}\xi'^\pm$ alone and does not increase in them; and for $\epsilon>0$, $d(\xi),d(\xi')>4\epsilon$,
\[\mathsf B_\lambda(\xi,\xi')=\int_{|b|_\infty\le\epsilon}\mathsf T(\xi+ib-3\epsilon n_0)\,\mathsf B_{6\epsilon}\bigl((\tau_{3\epsilon n_0}\chi_{\epsilon,b})\otimes\lambda\bigr)\,\mathsf T(\xi'-ib-3\epsilon n_0)\,db,\]
so that $\|\mathsf B_\lambda(\xi,\xi')\|\le4c_0\epsilon^{-2}\mathsf C(\mathcal K,6\epsilon)|\varphi|_\infty|a|_\infty$, $\mathcal K=\{2\epsilon\le x_3\le4\epsilon,\ |x_1|\le\epsilon,\ x_2\in I\}$, $\operatorname{supp}\varphi\subset I$; here $\chi_{\epsilon,b}(a')=\chi_\epsilon(a'+ib)$, $\chi_\epsilon(w)=c_\chi\epsilon^{-4}\mathsf X(|w^+|^2/\epsilon^2)\mathsf X(|w^-|^2/\epsilon^2)$, with a fixed smooth $\mathsf X\ge0$ supported in $[0,1]$ and $c_\chi$ fixed, is a bump of total mass one on the polydisc $\{|w^\pm|\le\epsilon\}$ of $\mathbb C^2$, and $|\chi_{\epsilon,b}|\le c_0\epsilon^{-4}$.
\item[(b)] The \emph{plate vectors} $\Psi(x;\underline\lambda)=\mathsf B_{\lambda_1}(x_1,\frac12(x_2-x_1))\cdots\mathsf B_{\lambda_n}(\frac12(x_n-x_{n-1}),\kappa)\Omega$, for $x_1,x_{k+1}-x_k\in\mathsf{Tb}$, and the \emph{kernels} $\mathsf K_n(\xi;\underline\lambda)=\langle\Omega,\mathsf B_{\lambda_1}(\kappa,\frac12\xi_1)\cdots\mathsf B_{\lambda_n}(\frac12\xi_{n-1},\kappa)\Omega\rangle$, $\xi\in\mathsf{Tb}^{n-1}$, do not depend on $\kappa\in\mathsf{Tb}$ and are holomorphic and multilinear; $\mathsf K_1=0$; $\mathsf K_n$ is the density of $\sigma$ at plates with time-ordered supports, $\sigma(\alpha_1\otimes\lambda_1,\dots,\alpha_n\otimes\lambda_n)=\int\prod\alpha_a(y_a)\mathsf K_n(y_2-y_1,\dots;\underline\lambda)\,dY$; and
\[\langle\Psi(x;\underline\lambda),\Psi(z;\underline\lambda')\rangle=\mathsf K_{n+n'}\bigl((x_n-x_{n-1})^*,\dots,(x_2-x_1)^*,x_1^*+z_1,z_2-z_1,\dots;\lambda_n^\dagger,\dots,\lambda_1^\dagger,\underline\lambda'\bigr).\]
\item[(c)] For weight labels there is exactly one holomorphic function $\mathcal K_n(\cdot;\underline\lambda)$ on $\mathcal Z_n$ that equals $e^{-i\kappa\upsilon}\mathsf K_n(\mathsf R_{-\upsilon}\partial(\mathfrak sZ);\mathfrak s\underline\lambda)$ on every tube $\mathcal Z(\mathfrak s,\upsilon)$; it is invariant under permutations of points with their labels and under common translations, $\mathcal K_n(\mathsf R_{\pi/2}Z;\underline\lambda)=e^{-i\kappa\pi/2}\mathcal K_n(Z;\underline\lambda)$, $\overline{\mathcal K_n(Z;\underline\lambda)}=\mathcal K_n(\bar Z;\underline{\bar\lambda})$, and for real $s$ and real $x$, $z$ in the domains of the plate vectors
\[\langle\Psi(\mathsf R_{is}x;\underline\mu),\Psi(\mathsf R_{is}z;\underline\mu')\rangle=\mathcal K_{j+j'}\bigl(\mathsf R_{is}(\theta x_j,\dots,\theta x_1,z_1,\dots,z_{j'});\mu_j^\dagger,\dots,\mu_1^\dagger,\underline\mu'\bigr),\]
$\underline\mu$ and $\underline\mu'$ being tuples of $j$ and $j'$ weight labels and $\theta$ acting on points of the plane by $y^0\mapsto-y^0$.
\item[(d)] If $\sigma$ has (P), then 
\[
\|\mathsf B_\lambda(\xi,\xi')\|\le C|a|_\infty\|\varphi\|_{L^1}\min\{1,d(\xi),d(\xi')\}^{-p};
\]
for every compact set $\mathcal C$ of fans of $j$ points and every tuple of $j$ weight labels, \[\sup_{s\in\mathbb R,\,z\in\mathcal C}e^{-pj|s|}\|\Psi(\mathsf R_{is}z;\underline\lambda)\|<\infty.\] The kernels satisfy 
\[
|\mathsf K_n(\xi;\underline\lambda)|\le\prod_k\bigl(2^pC|a_k|_\infty\|\varphi_k\|_{L^1}\bigr)\min\{1,\min_kd(\xi_k)\}^{-pn}.
\]

\item[(e)] If instead $\sigma$ has only an insertion bound in the sup norm, 
\[
|\sigma(\theta\bar F\cdot(h)\cdot G)|\le N(\eta)|h|_\infty\|[\tau_{-(\eta/4)e_3}F]\|\,\|[\tau_{-(\eta/4)e_3}G]\|
\]
 at the positions of Proposition~\ref{cov-lim}\,(e), then 
\[
|\mathsf K_n(\xi;\underline\lambda)|\le\prod_k\bigl(2c_{\rm plt}(\lambda_{\varphi_k}+1)\|\varphi_k\|_{C^1}|a_k|_\infty\bigr)(N(\eta)\eta^{-4})^n
\]
 for every $\eta\in(0,\frac12\min\{1,\min_kd(\xi_k)\}]$, $\lambda_\varphi$ the length of the least interval holding $\operatorname{supp}\varphi$ and $c_{\rm plt}$ a number of the constants and $c_{\rm I}$; with $N(\eta)=C_{\rm I}\eta^{-\mathsf p_{\rm I}}$ this is the power $\mathsf p_{\rm I}+4$.
\end{enumerate}
\end{lemma}
\begin{proof}
(a) The obstacle is that (S) bounds a plate by the sup norm of its profile, which blows up as the profile shrinks to a point; but the slab operator of a translate of a plate is that of the plate multiplied by $\mathsf T(\kappa)$ on the left and $\mathsf T(\kappa')$ on the right, also at complex $\kappa$, $\kappa'$, and a holomorphic function is its own mean over a polydisc of $\mathbb C^2$, a set of four real dimensions of which two are taken by the profile and two by the complex translations. So put $\mathsf M_\epsilon(\beta;\zeta,\zeta')=\mathsf T(\zeta-3\epsilon n_0)\mathsf B_{6\epsilon}((\tau_{3\epsilon n_0}\beta)\otimes\lambda)\mathsf T(\zeta'-3\epsilon n_0)$; it is linear and bounded in $\beta$, holomorphic in $(\zeta,\zeta')$, and $\mathsf M_\epsilon(\tau_x\beta;\zeta,\zeta')=\mathsf M_\epsilon(\beta;\zeta+x,\zeta'-x)$ for small real $x$, by the relation $\mathsf T(\kappa)\mathsf B_l(h)\mathsf T(\kappa')=\mathsf B_{l+\kappa^0+\kappa'^0}(\tau_\kappa h)\mathsf T((0,\kappa^1+\kappa'^1))$.

The displayed integral defines $\mathsf B^\epsilon(\xi,\xi')$; the mean value property on the polydisc shows that $\int\beta(x)\mathsf B^\epsilon(\xi+x,\xi'-x)dx=\mathsf M_\epsilon(\beta;\xi,\xi')$ for small $\beta$, so the families for different $\epsilon$ agree on an open set and, by the identity theorem, everywhere; a partition of unity gives the integral formula. The adjoint comes from (D2) at the reflection $x_3\mapsto l-x_3$, and the monotony from the semigroup law, the imaginary parts acting by unitaries. Uniqueness: the integral formula and continuity fix $\mathsf B_\lambda(dn_0,dn_0)$, and $\mathsf B_\lambda(\xi,\xi')=\mathsf T(\xi-dn_0)\mathsf B_\lambda(dn_0,dn_0)\mathsf T(\xi'-dn_0)$ at $d=\min\{d(\xi),d(\xi')\}$.
(b) Products depend on their inner arguments through their sums alone; the density and the Gram formula follow from the integral formula, (S) applied from the right, and (D2); $\mathsf K_1=0$ is $\sigma((H))=0$.
(c) By (D2) at the quarter turn $\mathsf R_{\pi/2}\in W_3$ and the symmetry of $\sigma$, each of the functions displayed is the density of $\sigma$ at the plates whose supports lie in its tube; two of them agree at the real points of the convex intersection of their tubes, hence on it, and $\mathcal Z_n$ is connected. Conjugation: $\sigma$ is real at real tuples, so $\overline{\mathsf K_n(\xi;\underline\lambda)}=\mathsf K_n(\xi;\underline{\bar\lambda})$ at real points, and $\xi\mapsto\overline{\mathsf K_n(\bar\xi;\underline\lambda)}$ is holomorphic. The scalar products are the Gram formula of (b) at boosted points, taken in the tube $\mathcal Z(\mathrm{id},0)$.
(d) Approximate $\mathsf B_\lambda(dn_0,dn_0)$ by $\mathsf B_{2d}(\alpha_\epsilon\otimes\lambda)$, with bumps $\alpha_\epsilon\ge0$ of mass one shrinking to $dn_0$, and apply (P) with $d'=\epsilon$; other points by (a) with two contractions. As $d(\mathsf R_{is}\xi)\ge e^{-|s|}d(\xi)$ and $\Psi(\mathsf R_{is}z;\underline\lambda)$ is a product of $j$ point operators at rotated arguments from a compact set of pairs of points of $\mathsf{Tb}$, applied to $\Omega$, the growth follows; in $\mathsf K_n$ take $\kappa=n_0$: every argument has $d\ge\frac12\min\{1,\min_kd(\xi_k)\}$.
(e) The same construction, with the insertion bound in the place of (P): a point plate is the mean of slab operators of bumps against the semigroup, a bump of size $\eta$ costs $\eta^{-4}$ in the sup norm of its profile on the polydisc, and the label, of class $C^1$, is cut into at most $\lambda_\varphi/\eta+3$ pieces.
\end{proof}

\begin{lemma}\label{cov-C1}
Let $\sigma$ be labelled data and $\sigma^1$ its restriction to the tuples with members in $C^1_c(\mathbb R^3;V)$. (a) The classes of the tuples of $C^1$ members are dense in $\mathcal H(\sigma)$; the identification of $\mathcal H(\sigma^1)$ with $\mathcal H(\sigma)$ carries $\Omega$, the semigroup, the translations and the spectral measure into one another, and an insertion bound of $\sigma$ holds for $\sigma^1$ with the same numbers. (b) If $\sigma^1(\mathsf rX)=\sigma^1(X)$ for one $\mathsf r\in\mathrm O(3)$ and all tuples of $C^1$ members, then $\sigma(\mathsf rX)=\sigma(X)$ for all $X\in\mathfrak S_V(\mathbb R^3)$.
\end{lemma}
\begin{proof}
Mollify the members with a radial mollifier $j_\epsilon$: $(\mathsf rH)*j_\epsilon=\mathsf r(H*j_\epsilon)$ for $\mathsf r\in\mathrm O(3)$, the supports stay in fixed neighbourhoods, and the bound (D1) at one type gives the convergence of both sides of each identity, and of the classes in $\mathcal H(\sigma)$.
\end{proof}

\begin{lemma}\label{cov-tower}
Let $\sigma$ be labelled data; (S) is not supposed. Let $l>0$, $c'\in\mathbb R$, $G=(g_1,\dots,g_m)$ a tuple with complex members in the slab $\{0<x_3<l\}$, at any heights, with supports in $\{x_1<c'\}$; $G^\star=\theta\overline{\tau_{-le_3}G}$ its mirror image in the middle plane of the slab, $g^\star(x)=\rho(\theta)\overline{g(x_1,x_2,l-x_3)}$; $\mathbb T_n(G)=G^\star\cdot\tau_{le_3}G\cdot\tau_{2le_3}G^\star\cdots\tau_{(2n-1)le_3}G$ the tower of $2n$ floors; $\vartheta_{c'}$ the reflection in $x_1=c'$, and
\[\beta_G=\limsup_{n\to\infty}\sigma\bigl(\vartheta_{c'}\overline{\mathbb T_n(G)}\cdot\mathbb T_n(G)\bigr)^{1/4n}\in[0,\infty] .\]
If $\beta_G<\infty$, there is exactly one bounded operator $\mathbb B(G;l)$ on $\mathcal H(\sigma)$ with $\mathbb B(G;l)[F]=[G\cdot\tau_{le_3}F]$ for all $F\in\mathbb C\mathfrak S_V(\mathbb E_+)$, $\|\mathbb B(G;l)\|\le\beta_G$, and $\mathbb B(G;l)^*[F]=[G^\star\cdot\tau_{le_3}F]$.
\end{lemma}
\begin{proof}
For $F$ carried by the quarter space $\{x_3>0,\ x_1>c'\}$ put $\mathsf s_n=\sigma(\theta\bar F\cdot\mathbb T_n(G)\cdot\tau_{2nle_3}F)$, so that $\mathsf s_0=\|[F]\|^2$ and $\mathsf s_1=\|[G\cdot\tau_{le_3}F]\|^2$. Positivity at $x_3=0$ and Schwarz' inequality, iterated as in the proof of the chessboard estimate with the convexity step taken before any quotient, give $\mathsf s_1^n\le\mathsf s_0^{n-1}\mathsf s_n$, although no operator is known to exist yet. Schwarz' inequality at the second mirror $x_1=c'$, positive by (D3), gives $\mathsf s_n\le\sigma(\vartheta_{c'}\overline{\mathbb T_n(G)}\cdot\mathbb T_n(G))^{1/2}\sigma(\vartheta_{c'}\overline{Y_n}\cdot Y_n)^{1/2}$ with $Y_n=\theta\bar F\cdot\tau_{2nle_3}F$; the last factor is a value at four groups of copies of $F$ that stay at mutual distances bounded below independently of $n$, so it is at most a number $C_F$ free of $n$, by (D1). Hence $\mathsf s_1\le\mathsf s_0\beta'^2(C_F^{1/2}/\mathsf s_0)^{1/n}$ for every $\beta'>\beta_G$ and large $n$, and $\|[G\cdot\tau_{le_3}F]\|\le\beta_G\|[F]\|$. The classes over the quarter space are dense by the cone (D4), and the adjoint identity, applied to all pairs, carries the operator to all classes.
\end{proof}

For a tensor limit and one member $h$ as in Proposition~\ref{cov-lim}\,(f), with $l=\eta'$ and $c'=z_1+\frac{\eta'}8$, the lemma gives $\|\mathbb B((h);\eta')\|\le C_{\mathsf p}\eta'^{-\mathsf p}|h|_\infty$ for every $\mathsf p>\frac{11}2$, at every $\eta'\in(0,\frac12]$.

A curve in a Hilbert space is smooth, or real-analytic, in norm as soon as its Gram kernel is so across the diagonal alone.

\begin{lemma}\label{cov-gram}
Let $\mathcal H$ be a Hilbert space, $N\ge1$ an integer, $\varsigma,\delta,\mathsf M>0$, $I=(-\varsigma,\varsigma)^N$, $v:I\to\mathcal H$ continuous in norm, and $f(x,y):=\langle v(x-y),v(x+y)\rangle$ at the points $(x,y)\in\mathbb R^N\times\mathbb R^N$ with $x\pm y\in I$. No regularity of $f$ in $x$ is supposed.
\begin{enumerate}
\item[(a)] Let $m\ge1$, $k=2(N+m+1)$, and let the derivatives $\partial^c_{y_j}f$, $1\le j\le N$, $0\le c\le k$, exist, be continuous and satisfy $|\partial^c_{y_j}f|\le\mathsf M^2(c/e\delta)^c$. Then $v$ is $m$ times continuously differentiable in norm on $(-\frac58\varsigma,\frac58\varsigma)^N$, and there $\|\partial^\gamma v\|\le c_{N,m}(\varsigma,\delta)\,\mathsf M$ for $|\gamma|\le m$, $c_{N,m}(\varsigma,\delta)$ a number of $N$, $m$, $\varsigma$ and $\delta$.
\item[(b)] If $N=1$ and the bounds of (a) on the derivatives in $y$ hold for all $c\ge0$, then on $(-\frac58\varsigma,\frac58\varsigma)$, with $\mathsf y=\frac{48}\varsigma+\frac2{e\delta}$, $\|v^{(m)}\|\le\frac{4\varsigma}\pi\mathsf M\bigl(\frac12(m+2)\mathsf y\bigr)^{m+1}$ for every $m\ge0$, and $v$ is real-analytic in norm: its Taylor series at $s_0$ converges in norm to $v(s)$ for $|s-s_0|<2/(e\mathsf y)$.
\end{enumerate}
\end{lemma}
\begin{proof}
Let $w\in C^k_c(I)$ be a window with values in $[0,1]$, equal to $1$ on $(-\frac58\varsigma,\frac58\varsigma)^N$. The norm square of the Fourier transform $\hat V(p)=\int w(s)e^{-ip\cdot s}v(s)\,ds$ reads the Gram kernel at the frequency $(p,-p)$ alone, that is across the diagonal: with $s=x-y$, $s'=x+y$,
\[\|\hat V(p)\|^2=\iint w(s)w(s')e^{ip\cdot(s-s')}\langle v(s),v(s')\rangle\,ds\,ds'=2^N\iint w(x-y)\,w(x+y)\,f(x,y)\,e^{-2ip\cdot y}\,dx\,dy .\]
At fixed $x$, $k$ integrations by parts in the coordinate $y_j$ with $|p_j|=|p|_\infty$ and Leibniz' rule give $\|\hat V(p)\|^2\le(2\varsigma)^{2N}\mathsf M^2\,c_w\,(2|p|_\infty)^{-k}$, $c_w$ a number of the bounds of the derivatives of $w$ up to the order $k$ and of the numbers $(c/e\delta)^c$, $c\le k$. As $k-2m=2N+2$, the functions $p^\gamma\hat V(p)$, $|\gamma|\le m$, are integrable in norm, and Fourier inversion, applied to the scalar functions $\langle\psi,wv\rangle$, $\psi\in\mathcal H$, shows that $wv$ is $m$ times continuously differentiable in norm, with $\partial^\gamma(wv)(s)=(2\pi)^{-N}\int(ip)^\gamma e^{ip\cdot s}\hat V(p)\,dp$; this is (a). For (b) take, at the order $m$, $q=m+2$ and the window $w_q=\mathbf 1_{[-\frac34\varsigma,\frac34\varsigma]}*\beta_q^{*(2q+1)}$, $\beta_q$ the normalised indicator of an interval of the length $\varsigma/(4(2q+1))$ about $0$: it lies in $C^{2q}$, has values in $[0,1]$, equals $1$ on $[-\frac58\varsigma,\frac58\varsigma]$, vanishes off $[-\frac78\varsigma,\frac78\varsigma]$, and $|w_q^{(c)}|\le(24q/\varsigma)^c$ for $c\le2q$, one factor $\beta_q$ being spent for each derivative.

The multinomial theorem then gives $\|\hat V_q(p)\|\le2\varsigma\mathsf M\min\{1,(q\mathsf y/2|p|)^q\}$, so that 
\[
(2\pi)^{-1}\int|p|^m\|\hat V_q(p)\|\,dp\le\frac{4\varsigma}\pi\mathsf M(\frac12q\mathsf y)^{m+1},
\]
 which bounds $v^{(m)}$ on the inner interval; Taylor's formula with these bounds and $\mathsf m!\ge(\mathsf m/e)^{\mathsf m}$ gives the convergence of the Taylor series.
\end{proof}

\begin{lemma}\label{cov-ordered}
Let $\sigma$ be labelled data. For every $X\in\mathfrak S_V(\mathbb E_+)$ and $\epsilon>0$ there are finitely many tuples $X_1,\dots,X_q\in\mathfrak S_V(\mathbb E_+)$, each with as many members as $X$, whose members lie in pairwise disjoint slabs $\{c<x_3<c'\}$, the $a$-th member of each being $\psi\cdot(H_a*\phi')$ with a mollifier $\phi'$ and a function $0\le\psi\le1$ in $C^\infty_c(\mathbb R^3)$, such that $\|[X]-\sum_i[X_i]\|<\epsilon$. So $\Omega$ and the classes of the tuples in pairwise disjoint slabs are total in $\mathcal H(\sigma)$, and two labelled data that agree at all such tuples over $\mathbb E_+$ are equal. If $\sigma$ has (S), $\Omega$ and the vectors $\mathsf B_{l_1}(h_1)\cdots\mathsf B_{l_n}(h_n)\Omega$ are total.
\end{lemma}
\begin{proof}
\emph{Smooth members, small pieces.} By (D1) the classes are multilinear in the members and bounded by their sup norms, with numbers that read no position; so, mollifying the members to $H_a*\phi_\varepsilon$ and telescoping, it suffices to treat smooth members. Let $\mathsf d_X$ be the least sup distance of two supports and $t_X$ the least height of a point of a support, choose $r>0$ with $4(n+1)r\le\mathsf d_X$ and $4r\le t_X$, put $t_1=\frac14t_X$, and cut every member by the smooth partition of unity $\chi_z(x)=\prod_\mu\chi((x_\mu-z_\mu)/r)$, $z\in r\mathbb Z^3$, $0\le\chi\le1$, $\operatorname{supp}\chi_z\subset z+[-\frac58r,\frac58r]^3$. Then $[X]$ is a finite sum of classes of tuples $Y=(g_1,\dots,g_n)$ on \emph{frames}: $g_a$ is carried by a box $Q_a$ with sides of lengths $\le2r$ in $\{x_3\ge2t_1\}$, and $\operatorname{dist}_\infty(Q_a,Q_b)\ge4nr$ for $a\ne b$. A pair $\{a,b\}$ is \emph{bad} if the projections of $Q_a$ and $Q_b$ to the $x_3$-axis meet; a tuple on a frame without bad pairs is ordered. We show by induction on the number of bad pairs, at fixed $n$, $r$, $t_1$, that $[Y]$ is a limit of sums of ordered classes.
\emph{A cut.} Let $\{a,b\}$ be bad. As $Q_a$, $Q_b$ are $4nr$ apart in the sup distance and their third sides meet, they are $4nr$ apart in a direction $\mu\in\{1,2\}$; among the planes $x_\mu=c$ between them, an interval of the length $\ge(4n-2)r$, the $n-2$ other boxes forbid intervals of the total length $\le(4n-8)r$. So there is $c$ such that every box lies in $\{x_\mu\le c-r\}$ or in $\{x_\mu\ge c+r\}$, $Q_a$ in the first half space and $Q_b$ in the second; let $Y_{\mathsf L}$, $Y_{\mathsf R}$ be the two groups of members.
\emph{The two groups move smoothly against one another in the height.} For $s\in(-t_1,t_1)^2$ put $v(s)=[\tau_{s_1e_3}Y_{\mathsf L}\cdot\tau_{s_2e_3}Y_{\mathsf R}]$. As $\theta\tau_{te_3}=\tau_{-te_3}\theta$ and $\sigma$ is symmetric, the Gram kernel at $s=x-y$, $s'=x+y$ is $\langle v(s),v(s')\rangle=\sigma(A_{x_1}\cdot\tau_{(y_2-y_1)e_3}B_{x_2})$, with $A_t=\tau_{-te_3}\theta Y_{\mathsf L}\cdot\tau_{te_3}Y_{\mathsf L}$ and $B_t=\tau_{-te_3}\theta Y_{\mathsf R}\cdot\tau_{te_3}Y_{\mathsf R}$: the variable $y$ moves each group together with its mirror copy, rigidly, and no regularity in $x$ is needed. The plane $x_\mu=c$ parts $A_t$ from $B_t$ with the margin $r$. By (D2), with the element of $W_3$ that exchanges $e_\mu$ and $e_3$ and a translation, and by (D4), $\sigma(A\cdot\tau_{\eta e_3}B)=\langle\Phi_A,e^{-rH}e^{i\eta P_\mu}\Phi_B\rangle$ for two vectors with $\|\Phi_A\|^2=\sigma(\vartheta A\cdot A)$, $\vartheta$ the reflection in the plane $x_\mu=c-\frac r2$, and likewise for $B$ with the plane $x_\mu=c+\frac r2$; as $|P_\mu|\le H$ on the carrier of the spectral measure, the derivatives of the order $c$ in $\eta$ are at most $(c/er)^c\|\Phi_A\|\|\Phi_B\|$. $\|\Phi_{A_t}\|^2$ is a value of $\sigma$ at $4\#\mathsf L$ members with the sup norms $|g_e|_\infty$, of the type $(\min\{r,2t_1\};\sqrt3r)$, so by (D1) it is bounded by a number free of $t$ and of the positions; and the same for $B_t$.

So Lemma~\ref{cov-gram}\,(a) at $N=2$ makes $v$ twice continuously differentiable in norm, with bounds by the sup norms of the members, and
\[
\begin{multlined}
\partial_{s_1}v(0)=-\sum_{e\in\mathsf L}[Y[e\!:\!\partial_3g_e]],\qquad\partial_{s_2}v(0)=-\sum_{e'\in\mathsf R}[Y[e'\!:\!\partial_3g_{e'}]],\qquad\partial_{s_1}\partial_{s_2}v(0) \\
=\sum_{e\in\mathsf L,\,e'\in\mathsf R}[Y[e\!:\!\partial_3g_e,e'\!:\!\partial_3g_{e'}]],
\end{multlined}
\]
$Y[e\!:\!g]$ being $Y$ with $g$ in the place of $g_e$. The same argument about every $\varrho_0\ge0$, with Lemma~\ref{cov-gram}\,(b), makes the lift $u(\varrho)=[Y_{\mathsf L}\cdot\tau_{\varrho e_3}Y_{\mathsf R}]$ real-analytic in norm on $(-\frac58t_1,\infty)$.
\emph{$L^1$ in the height, in two members at once.} The derivative formulas bound $[Y]$, with $g_a$ replaced by $\partial_3G$, or $g_b$ by $\partial_3G'$, or both, by the sup norms of $G$, $G'$ and of the other members and of their derivatives $\partial_3g_e$. Writing $g_a=\partial_3G+\kappa_a\int g_a\,dx_3$ with a fixed bump $\kappa_a\ge0$ of the integral $1$ in the third side of $Q_a$, so that $|G|_\infty\le2\|g_a\|_{\infty,1}$, $\|g\|_{\infty,1}:=\sup_{x_1,x_2}\int|g(x_1,x_2,t)|\,dt$, and the same for $g_b$, one gets
\[\|[Y]\|\le\mathsf K\,\|g_a\|_{\infty,1}\|g_b\|_{\infty,1}\prod_{e\ne a,b}\max\{|g_e|_\infty,|\partial_3g_e|_\infty\},\]
$\mathsf K$ a number of $n$, $r$, $t_1$, the two bumps and the numbers of (D1), free of the members.
\emph{Layers.} Cut $g_a$ and $g_b$ by the partition of unity $\chi(Nx_3-k)$, $k\in\mathbb Z$, of the height: $[Y]=\sum_{k,k'}[Y_{kk'}]$, with at most $2rN+3$ values of $k$ and of $k'$. For $|k-k'|\ge2$ the pair $\{a,b\}$ of $Y_{kk'}$, a tuple on a frame with smaller boxes, is good, and no new bad pair arises; by the induction $[Y_{kk'}]$ is a limit of sums of ordered classes. For $|k-k'|\le1$, at most $3(2rN+3)$ terms, each layer has $\|\cdot\|_{\infty,1}\le|g|_\infty/N$, so that their sum has the norm $\le3(2rN+3)\,\mathsf K\,N^{-2}|g_a|_\infty|g_b|_\infty\prod_{e\ne a,b}\max\{|g_e|_\infty,|\partial_3g_e|_\infty\}$, which tends to $0$ as $N\to\infty$. So $[Y]$ is a limit of sums of ordered classes, and the approximants of $[X]$ are sums, with the coefficients $1$, of ordered tuples whose $a$-th member is $H_a*\phi_\varepsilon$ times a finite product of the functions $\chi_z$ and $\chi(N\cdot-k)$, each with values in $[0,1]$; they are formed with $X$, $\varepsilon$, $r$ and the numbers $N$ of the stages alone.

This is the first assertion, and $\Omega$ and the ordered classes are total. The lift gives the totality too: for $\varrho$ large the frame of $Y_{\mathsf L}\cdot\tau_{\varrho e_3}Y_{\mathsf R}$ has fewer bad pairs, so by the induction $u(\varrho)$ lies in the closed span of the ordered classes, and the scalar product of $u(\varrho)$ with a vector orthogonal to that span is a real-analytic function of $\varrho$ on $(-\frac58t_1,\infty)$ that vanishes for large $\varrho$, hence at $\varrho=0$.
\emph{Equality of data.} $|\sigma(F)|=|\langle\Omega,[F]\rangle|\le\|[F]\|$. Taking at every stage $\varepsilon$ so small and $N$ so large that the errors are small in the norms of both data, which the approximants allow as they are formed with $X$, $\varepsilon$, $r$ and $N$ alone, $|\sigma(X)-\sigma'(X)|$ is at most twice the error when $\sigma$ and $\sigma'$ agree at the ordered tuples over $\mathbb E_+$; by (D2) at a translation along $e_3$ this reaches every $X\in\mathfrak S_V(\mathbb R^3)$. The last assertion follows from (S) applied from the right.
\end{proof}

\begin{lemma}\label{cov-slide}
Let $\sigma$ be labelled data.
\begin{enumerate}
\item[(a)] If a slab $\{|x_\mu-c|<\delta\}$ parts the supports of $Y_I$ from those of $Y_J$ and $u\perp e_\mu$ is a unit vector, then $s\mapsto\sigma(Y_I\cdot\tau_{su}Y_J)$ is the restriction to $\mathbb R$ of a function holomorphic on $\{|\operatorname{Im}s|<\sqrt2\,\delta\}$, and on $\{|\operatorname{Im}s|<2\delta\}$ if $u$ is a coordinate vector.
\item[(b)] Let $\delta H=\frac d{dt}\mathsf R_tH|_{t=0}$ and $\mathsf D(X)=\sum_a\sigma(\dots,\delta H_a,\dots)$ for tuples of $C^1$ members; then $\frac d{dt}\sigma(\mathsf R_tX)=\mathsf D(\mathsf R_tX)$. If $\mathsf D(X)=0$ at every tuple $X$ of $n$ members of class $C^1$ whose projections to the $x_3$-axis are pairwise disjoint, then $\mathsf D=0$ at all tuples of $n$ members of class $C^1$, and $\sigma(\mathsf R_tX)=\sigma(X)$ for all real $t$ and all $X\in\mathfrak S_V(\mathbb R^3)$ of $n$ members.
\item[(c)] If $\sigma$ has (N), then $\sigma((H_\eta)\cdot X)=0$ for every $\eta\in C^2_c(\mathbb R^3;\mathbb R^3)$ and every $X$ with members off $\operatorname{supp}\eta$, among any number of members; in particular $\sigma(H_\eta,H')=0$ at disjoint supports.
\end{enumerate}
\end{lemma}
\begin{proof}
(a) By (D2), with an element of $W_3$ that turns $e_\mu$ into $e_3$ and $u$ into a unit vector $u'=(u'_1,u'_2,0)$, and a translation, the function is $\langle[\tau_{-\delta e_3}\theta Y_I],e^{-2\delta H}e^{is(u'_1P_1+u'_2P_2)}[\tau_{-\delta e_3}Y_J]\rangle=\int e^{-2\delta E+is(u'_1q_1+u'_2q_2)}dm(E,q)$ for a finite complex measure $m$ carried by $\{E\ge\max(|q_1|,|q_2|)\}$ (D4). There $|u'_1q_1+u'_2q_2|\le(|u'_1|+|u'_2|)E\le\sqrt2E$, and $\le E$ if $u'$ is a coordinate vector. (b) The derivative formula is the product rule under (D1). Cut the members into small pieces and argue by induction on the number of pairs of pieces with overlapping $x_3$-projections. Such a pair is parted by a slab normal to $e_1$ or $e_2$; sliding the group on one side along $e_3$ by $s$ lowers the number for large $s$, and since $\delta(\tau_{se_3}H)=\tau_{se_3}[\delta H-s\,\partial_1H]$ the derivative at the slid tuple is $f_0(s)+sf_1(s)$ with $f_0,f_1$ real-analytic by (a): it vanishes for large $s$, hence at $s=0$.

Continuous members follow by mollification and (D1). (c) Put $\Lambda(\eta;X)=\sigma((H_\eta)\cdot X)$; it is linear in $\eta$ and in each member, by (D1). \emph{Step 0.} If $\operatorname{supp}\eta\subset\{x_3>c\}$ and the members of $X$ lie in $\{x_3<c\}$, then $\Lambda(\eta;X)=0$: by (D2) let $c=0$; $F=\theta X$ is a real tuple over $\mathbb E_+$ with $\theta\bar F=X$, and $\Lambda(\eta;X)=\langle[F],[(H_\eta)]\rangle=0$ by (N). \emph{Step 1.} If the $x_3$-intervals of $\operatorname{supp}\eta$ and of the supports of the members are pairwise disjoint, then $\Lambda(\eta;X)=0$. Let $J$ be the set of the members whose $x_3$-interval lies above that of $\operatorname{supp}\eta$; if $J$ is empty, this is Step 0.

Else a slab $\{|x_3-c|<\delta\}$ parts the members of $J$, above it, from $\operatorname{supp}\eta$ and the other members, below it; put $Y_I=(H_\eta)\cdot(H_a)_{a\notin J}$ and $Y_J=(H_a)_{a\in J}$. By (a) at $\mu=3$, $u=e_1$, the function $f(s)=\sigma(Y_I\cdot\tau_{se_1}Y_J)$ is real-analytic on $\mathbb R$. For $s>s_0$ the $x_1$-intervals of the members of $\tau_{se_1}Y_J$ lie to the right of all those of $Y_I$, with a gap; fix such an $s$. Then a slab normal to $e_1$ parts $\tau_{se_1}Y_J$ from $Y_I$, and by (a) at $\mu=1$, $u=e_3$, the function $g(s')=\sigma(Y_I\cdot\tau_{-s'e_3}\tau_{se_1}Y_J)$ is real-analytic on $\mathbb R$; for $s'>s'_0$ all the members of $\tau_{-s'e_3}\tau_{se_1}Y_J$ lie below $\operatorname{supp}\eta$, as the members outside $J$ do, and $g(s')=0$ by Step 0.

So $g=0$, and $f(s)=g(0)=0$ for every $s>s_0$; so $f=0$, and $\Lambda(\eta;X)=f(0)=0$. \emph{Step 2.} Let the supports of $\eta$ and of the $n$ members be at mutual distances $\ge\mathsf d>0$, let $(n+3)\epsilon<\mathsf d/\sqrt3$, and cut $\eta$ and the members by a smooth partition of unity $(\chi_k)$ with supports in cubes of side $\epsilon$: $H_\eta=\sum_kH_{\chi_k\eta}$ with $\chi_k\eta\in C^2_c$, and $\Lambda(\eta;X)$ is the sum of the values at the tuples of pieces, one piece of $\eta$ and one of each member; a tuple with a piece $0$ gives $0$. In such a tuple two pieces whose $x_3$-intervals meet have, in the direction $e_1$ or $e_2$, intervals at a distance $>(n+1)\epsilon$; the other $n-1$ intervals in that direction, of lengths $\le\epsilon$, leave in this gap a level $x_\mu=c$ at a positive distance from all of them, which splits the pieces into two groups.

Argue by induction on the number of the pairs of pieces with meeting $x_3$-intervals: slide the group on one side of such a level along $e_3$ by $s$; the value is real-analytic in $s$ by (a); the slid tuple is again a tuple of pieces of this kind, as $\tau_{se_3}H_{\chi\eta}=H_{\tau_{se_3}(\chi\eta)}$ and the intervals in the directions $e_1$, $e_2$ do not move; and for large $s$ it has fewer such pairs, so that the value vanishes there, hence at $s=0$. At the end of the induction the $x_3$-intervals are pairwise disjoint: Step 1.
\end{proof}

\begin{proposition}\label{cov-two}
\begin{enumerate}
\item[(a)] Labelled data with (O) and (N) satisfy $\sigma(\mathsf RH,\mathsf RH')=\sigma(H,H')$ for all $\mathsf R\in\mathrm O(3)$ and all pairs with disjoint supports. If besides (S) holds, then $\mathcal K_2(\mathsf R_{is}Z;\lambda,\lambda')=e^{(k(\lambda)+k(\lambda'))s}\mathcal K_2(Z;\lambda,\lambda')$ on $\mathcal Z_2$ for all weight labels and real $s$, and at weight labels $\mathsf K_2(rn_\vartheta;(\varphi,a),(\psi,b))=\ell_0(r)e^{-i(k(a)+k(b))\vartheta}$, $|\vartheta|<\frac\pi2$: every harmonic $m\ne0$ vanishes at every pair of labels.
\item[(b)] The same conclusions hold for labelled data with (N), (O$_a$) and (I$_0$) at pairs, and it suffices that (D3) hold at the plane $x_3=0$ and, for tuples of one member, at the plane $\Sigma$.
\item[(c)] If labelled data are covariant at pairs under the rotations about $e_3$, the spectral measure of every vector $[(h)]$, $h\in C_c(\mathbb E_+;V_{\mathbb C})$, is carried by the cone $\{E\ge|q|\}$.
\end{enumerate}
\end{proposition}
\begin{proof}
(a) \emph{Kernel.} By (D3) at $x_3=0$ and (D4) there is a matrix $\mathsf M_{ab}$ of complex measures on the cone $\{E\ge\max(|q_1|,|q_2|)\}$, antilinear in the first label, with $(\mathsf M_{a_ia_j})_{ij}$ a positive matrix of measures, and $\sigma(H,H')=\iint\mathcal G(x'-x)(H(x),H'(x'))\,dx\,dx'$ at disjoint supports, where $\mathcal G(u)(a,b)=\int e^{-u_3E+i(u_1q_1+u_2q_2)}d\mathsf M_{a^\dagger b}$ at $u_3>0$; by (D2) the kernel continues to a real-analytic function on $\mathbb R^3\setminus\{0\}$ with $\mathcal G(\mathsf wu)(\rho(\mathsf w)a,\rho(\mathsf w)b)=\mathcal G(u)(a,b)$, $\mathsf w\in W_3$. With $\Gamma_a(t)=\int e^{-tE}d\mathsf M_{aa}$, non-increasing, and $\Gamma_*(t)=\max\{\Gamma_{\rho(\mathsf w)a}(t),\Gamma_{\rho(\mathsf w)b}(t):\mathsf w\in W_3\}$: on the polydisc $|v_k-u_k|\le|u|/6\sqrt3$ the function $\mathcal G(\cdot)(a,b)$ is holomorphic and bounded by $\Gamma_*(|u|/2\sqrt3)$, after a turn by an element of $W_3$ that makes the largest coordinate of $u$ the third; so $\mathcal G$ is bounded on $|u|\ge1$.
\emph{Conservation.} By (N) and Lemma~\ref{cov-slide}\,(c), $\sigma(H'',H_\eta)=0$ at disjoint supports; integrating by parts, $2\partial_i\mathcal G_{ij}(u;b)=\partial_j\operatorname{Tr}\mathcal G(u;b)$, $\mathcal G_{ij}(u;b)$ the kernel at the first label $e_ie_j^{\rm T}$, symmetrised.
\emph{Growth.} (O) and a test function whose transform has no zero on the cone give $\Gamma_o(t)\le c\,C_{\rm od}t^{-6}$, $t\le1$, at off-diagonal $o$, and $W_3$ maps off-diagonal labels to off-diagonal ones up to signs; Cauchy's estimate on the polydisc gives $|\partial_\mu\partial_\nu\mathcal G(u)(o,o')|\le C|u|^{-8}$ for $|u|\le2$. With $\bar a_\mu=2e_\mu e_\mu^{\rm T}-\mathbf1$ and $o_{\mu\nu}=e_\mu e_\nu^{\rm T}+e_\nu e_\mu^{\rm T}$, the $o$ of (O), conservation gives $\partial_\mu\mathcal G(\bar a_\mu,b)=-\sum_{\nu\ne\mu}\partial_\nu\mathcal G(o_{\mu\nu},b)$, hence $\partial_\mu^2\mathcal G(\bar a_\mu,\bar a_\mu)=\sum_{\nu,\nu'\ne\mu}\partial_\nu\partial_{\nu'}\mathcal G(o_{\mu\nu},o_{\mu\nu'})$, and two integrations in $u_\mu$ give $|\mathcal G(u)(\bar a_\mu,\bar a_\mu)|\le C(1+|u|^{-6})$.

Put $\Gamma_{\varphi a}(t)=\int e^{-tE}|\hat\varphi(q_2)|^2d\mathsf M_{aa}=\iint\varphi(s)\varphi(s')\mathcal G((0,s'-s,t))(a^\dagger,a)\,ds\,ds'$; with $\int(t^2+\tau^2)^{-3}d\tau=\frac{3\pi}8t^{-5}$ this gives $\Gamma_{\varphi a}(t)\le C_\varphi\max\{1,t^{-1}\}^5$ at $a=\bar a_\mu$ and at off-diagonal $a$; these span $V_{\mathbb C}$ ($\mathbf1=-\sum_\mu\bar a_\mu$), and $\mathsf M_{a+a',a+a'}\le2\mathsf M_{aa}+2\mathsf M_{a'a'}$, so the bound holds at every label. The scalar label $\mathbf1=\varepsilon_0+\varepsilon'_0$, of weight $0$, is reached by conservation and the off-diagonal bound alone.
\emph{Harmonics.} For weight labels $a,b$, $\kappa=k(a)+k(b)$, and real $\varphi,\psi\in C^1_c(\mathbb R)$, the function $\Upsilon(\zeta)=\iint\varphi(s)\psi(s')\mathcal G(\mathsf R_\zeta(rn_0),s'-s)(a,b)\,ds\,ds'$, the first argument of $\mathcal G$ being the point of the plane $y$ and the second the coordinate $x_2$, continues to an entire function with $\Upsilon(\zeta+\frac\pi2)=e^{-i\kappa\pi/2}\Upsilon(\zeta)$ and $|\Upsilon(\vartheta+is)|\le Ce^{5|s|}$; so
\[\Upsilon(\zeta)=\sum_{|4m-\kappa|\le5}\ell_m(r)\,e^{i(4m-\kappa)\zeta} .\]
\emph{Two mirrors.} For labels $(\varphi_i,a_i)$ put $\ell^{ij}_m=\ell_m(r;\varphi_ia_i^\dagger,\varphi_ja_j)$. The matrix $(\sum_m\ell^{ij}_mx^m)_{ij}$ is positive semi-definite for $x>0$, by (D3) at $x_3=0$, and for $x<0$, by (D3) at $\Sigma$ for tuples of one member, a positive definite function that continues to a strip being the Laplace transform of a positive measure there. A Laurent polynomial that is $\ge0$ on $\mathbb R\setminus\{0\}$ has even extreme powers; so at the diagonal entries only $m\in\{0,\mp1,\mp2\}$ survive at the labels of weight $\mp2$ and only $m=0$ at the others. Off the diagonal $|\sum_m\ell^{ij}_mx^m|^2\le(\sum_m\ell^{ii}_mx^m)(\sum_m\ell^{jj}_mx^m)$ on $\mathbb R\setminus\{0\}$; comparing the orders as $x\to0$ and as $x\to\infty$, the powers of an off-diagonal entry lie between the half sum of the least powers and the half sum of the greatest powers of the two diagonal entries. Hence at two labels of weights of modulus $\le1$ only $m=0$ survives, at a label of weight $\mp2$ and one of weight of modulus $\le1$ only $m\in\{0,\mp1\}$, and at two labels of weight $\mp2$ only $m\in\{0,\mp1,\mp2\}$; and by positivity all the $\ell_m$ with $m\ne0$ vanish once the extreme ones do.
\emph{The remaining harmonics vanish.} With $z=u_3+iu_1$, conservation contracted with $u_+$ reads $\partial_z\mathcal G(\varepsilon_{-2},b)=\partial_{\bar z}\mathcal G(\varepsilon'_0,b)-\frac12\partial_2\mathcal G(\varepsilon_{-1},b)$. Here $z=re^{i\vartheta}$ at the point $rn_\vartheta$ of the plane $y$, and $\partial_z=\frac12e^{-i\vartheta}(\partial_r-\frac ir\partial_\vartheta)$ lowers the frequency $4m-\kappa$ of a harmonic by one, $\partial_{\bar z}$ raises it by one, and $\partial_2$ keeps it. For $|k(b)|\le1$ only $m=0$ survives on the right side, at $(\varepsilon'_0,b)$ and at $(\varepsilon_{-1},b)$, with the frequency $1-k(b)$, while $m\in\{0,-1\}$ on the left side; so the frequency $-3-k(b)$, of $m=-1$, occurs on the left side alone, and gives
\[\ell_{-1}'-\frac{2+k(b)}r\,\ell_{-1}=0,\qquad\ell_{-1}=c\,r^{2+k(b)},\]
bounded as $r\to\infty$ because $\mathcal G$ is bounded on $|u|\ge1$; so $\ell_{-1}=0$. Then at $b=\varepsilon_{-2}$ only $m\in\{0,-1\}$ survive on the right side, at $(\varepsilon'_0,\varepsilon_{-2})$ and at $(\varepsilon_{-1},\varepsilon_{-2})$, with the frequencies $3$ and $-1$, while the left side has the frequencies $3$, $-1$, $-5$; so the frequency $-5$, of $m=-2$, occurs on the left side alone and gives $\ell_{-2}=c\,r^4=0$; the least power of the diagonal entry at $\varepsilon_{-2}$ being even, its harmonic $m=-1$ vanishes too; and likewise for $\varepsilon_{+2}$. So all $\ell_m$ with $m\ne0$ vanish, $\mathcal G(\mathsf R_tu)(\rho(\mathsf R_t)a,\rho(\mathsf R_t)b)=\mathcal G(u)(a,b)$, and $\sigma$ is covariant at pairs under the $\mathsf R_t$; with $W_3$ they generate $\mathrm O(3)$. Under (S) the function $e^{i\kappa\zeta}\mathsf K_2(\mathsf R_\zeta w)$, $w=rn_{\vartheta_0}$, is holomorphic on the strip $|\operatorname{Re}\zeta+\vartheta_0|<\frac\pi2$ and constant on its real interval, hence on the strip, which gives the law of $\mathcal K_2$ along $\mathsf R_{is}$.
(b) Here the momentum picture is used. Let $\mathsf{Tb}^\square=\{u\in\mathbb C^3:\operatorname{Re}u_3>|\operatorname{Im}u_1|+|\operatorname{Im}u_2|\}$, the tube of the square cone, and let $\beta_\tau(E,q_1,q_2)=(E\cosh\tau+q_1\sinh\tau,\,q_1\cosh\tau+E\sinh\tau,\,q_2)$ be the boost of the momenta; under it $p_\pm=E\pm q_1$ are multiplied by $e^{\pm\tau}$, and $E\circ\beta_\tau\ge e^{-|\tau|}E$ on $\{E\ge|q_1|\}$. The matrix $\mathsf M(e,e')$ of measures on the cone, antilinear in $e$, with $\langle[(ke)],\mathcal P(B)[(k'e')]\rangle=\int_B\overline{F_k}F_{k'}\,d\mathsf M(e,e')$ for $k,k'\in C_c(\mathbb E_+;\mathbb C)$, $F_k(E,q)=\int k(x)e^{-x_3E+i(x_1q_1+x_2q_2)}dx$, exists without a point bound: the sup norms of (D1) are unchanged by a modulation $h\mapsto he^{-i\kappa\cdot x}$, and, averaged over $\kappa$ against $e^{-\mathsf d(|\kappa_1|+|\kappa_2|)}$ at a bump $h$ of the size $\mathsf d$, they become bounds of $\int e^{-7\mathsf dE}d\mathsf M(e,e)$ at every $\mathsf d>0$ (the modulation bound).

The kernel $\mathcal G(u)(a,b)=\int e^{-u_3E+i(u_1q_1+u_2q_2)}d\mathsf M(a^\dagger,b)$ is holomorphic on $\mathsf{Tb}^\square$; its transforms by the elements of $W_3$ agree where their domains meet, so that it is one holomorphic function on the union of the images of $\mathsf{Tb}^\square$ under $W_3$, whose real points are $\mathbb R^3\setminus\{0\}$; and an imaginary rotation acts on $\mathsf M$ as the boost: $\mathcal G(\mathsf R_{i\tau}u)(a^\dagger,b)=\int e^{-u_3E+i(u_1q_1+u_2q_2)}d(\beta_\tau)_*\mathsf M(a,b)$ at real $u$ with $u_3>0$. Conservation, contracted with $u_+$, $u_-$ and $e_2$, becomes three linear relations between the entries of $\mathsf M$ with the coefficients $p_\pm$ and $q_2$:
\[-p_-\mathsf M(e,\varepsilon_{-2})+iq_2\mathsf M(e,\varepsilon_{-1})+p_+\mathsf M(e,\varepsilon'_0)=0,\qquad-p_+\mathsf M(e,\varepsilon_2)+iq_2\mathsf M(e,\varepsilon_1)+p_-\mathsf M(e,\varepsilon'_0)=0,\]
\[p_-\mathsf M(e,\varepsilon_{-1})+p_+\mathsf M(e,\varepsilon_1)=2iq_2\bigl[\mathsf M(e,\varepsilon'_0)-\mathsf M(e,\varepsilon_0)\bigr];\]
and (I$_0$) at pairs makes the scalar entry $\mathsf M(\mathbf1,\mathbf1)$ invariant under the boosts. \emph{The growth.} For a height $t>0$ let $\tau_n=\log(t/7\mathsf d_n)$, so that $te^{-\tau_n}=7\mathsf d_n$ and $\tau_n\to\infty$. (O$_a$) at $\mathsf d_n$ and the modulation bound give $\int e^{-7\mathsf d_nE}d\mathsf M(o,o)\le C\mathsf d_n^{-6-a}$ at every off-diagonal $o$, so $\int e^{-tE\circ\beta_{\pm\tau_n}}d\mathsf M(o,o)\le\int e^{-7\mathsf d_nE}d\mathsf M(o,o)\le C(t)e^{(6+a)\tau_n}$; the scalar entry does not move under the boosts; conservation passes the bound to $e_3e_3^{\rm T}$, and to $e_1e_1^{\rm T}$ off $\{|q_1|<\frac12E\}$, while on $\{|q_1|\le\frac12E\}$ a boost does not lower the energy below its half, $E\circ\beta_\tau\ge\frac12E$; and $e_2e_2^{\rm T}=\mathbf1-e_1e_1^{\rm T}-e_3e_3^{\rm T}$.

So $\int e^{-tE\circ\beta_{\pm\tau_n}}d\mathsf M(e,e)\le C(t)|e|^2e^{(6+a)\tau_n}$ at every $e$. This bound along the sequence $\pm\tau_n$ alone suffices: for real $u$ with $(u_3,u_1)\ne0$ the function $f(\zeta)=\mathcal G(\mathsf R_\zeta u)(e,e')$ is entire and of period $2\pi$, its Fourier coefficients $f_\nu=\frac1{2\pi}\int_0^{2\pi}f(t'+is)e^{-i\nu(t'+is)}dt'$ do not depend on $s$, and $|f_\nu|\le Ce^{(6+a)\tau_n\mp\nu\tau_n}\to0$ for $\pm\nu>6+a$. So $f$ is a trigonometric polynomial with $|\nu|\le6+a$, and by Vandermonde's determinant the boosted matrix is a finite sum $(\beta_\tau)_*\mathsf M=\sum_{|\nu|\le6+a}e^{-\nu\tau}\mathsf M_\nu$, with $(\beta_\tau)_*\mathsf M_\nu=e^{-\nu\tau}\mathsf M_\nu$.

\emph{The harmonics.} By the quarter turn $\mathsf R_{\pi/2}\in W_3$, at weight labels $\mathsf M_\nu(e,e')=0$ unless $\nu=4m-k(e)-k(e')$ with an integer $m$. The \emph{harmonic of number $m$} is the matrix $\mathsf M^{(m)}=(\mathsf M_{4m-k(e)-k(e')}(e,e'))_{e,e'}$ over the weight basis, with $\mathbf1$ in the place of $\varepsilon_0$; its number is the $m$ of (a), of the frequency $4m-\kappa$; it is \emph{wrong} if $m\ne0$; and $\nu=4m-k(e)-k(e')$ is the index of its entry at $(e,e')$. Each $\mathsf M^{(m)}$ satisfies the three relations: $p_\pm$ carry the boost weights $\pm1$, so that a relation taken at one power $e^{-\lambda\tau}$ links the indices $\lambda+1$, $\lambda$, $\lambda-1$, as in $-p_-\mathsf M_{\lambda+1}(e,\varepsilon_{-2})+iq_2\mathsf M_\lambda(e,\varepsilon_{-1})+p_+\mathsf M_{\lambda-1}(e,\varepsilon'_0)=0$, and at $\lambda=4m-k(e)+1$ these are entries of one $\mathsf M^{(m)}$.

The lowest and the highest $\mathsf M^{(m)}$ that are not zero are positive matrices (multiply the positive matrix $(e^{-(k(e)+k(e'))\tau}(\beta_\tau)_*\mathsf M(e,e'))_{e,e'}=(\sum_me^{-4m\tau}\mathsf M^{(m)}(e,e'))_{e,e'}$ by $e^{4m\tau}$ and let $\tau\to\pm\infty$), and the diagonal mirror makes $\sum_m(-1)^m\mathsf M^{(m)}(e,e)$ positive at each weight label $e$. Let the lowest wrong harmonic have the number $-k$, $k\ge1$ (the highest is treated alike). If $k$ is odd, it is both positive and negative at each label, hence zero. If $k$ is even, $k\ge2$, its diagonal entries at the labels $\varepsilon_1$, $\varepsilon_2$ carry the indices $-4k-2$, $-4k-4$, of modulus $\ge10$, outside the window because $6+a<10$, and vanish, and so does its diagonal entry at the scalar label $\mathbf1$, the scalar entry having no wrong harmonic; by positivity the rows of these labels vanish; the three relations then push the remaining entries onto $\{p_-=0\}$ and the entry at $\varepsilon'_0$ onto $\{0\}$, where the boost scales it by $e^{4k\tau}$; and a positive measure on $\{p_-=0\}$ whose mass on a smaller set is $e^{|\nu|\tau}$ times its mass on a larger one, $\nu\ne0$, is zero.

So all wrong harmonics vanish.
(c) A rotation $\mathsf S$ about $e_3$ keeps $\mathbb E_+$ and commutes with $\theta$; covariance at pairs makes $[(h)]\mapsto[(\mathsf Sh)]$ unitary on the closed span of the vectors of one member, with $e^{-tH}e^{iv\cdot P}$ turned into $e^{-tH}e^{i(\mathsf Sv)\cdot P}$; so the spectral measure of $[(\mathsf Sh)]$ is the image of that of $[(h)]$ under $(E,q)\mapsto(E,\mathsf Sq)$. As $[(\mathsf Sh)]$ is again a vector of one member, its spectral measure is carried by the square cone (D4); so that of $[(h)]$ is carried by the preimage of the square cone under $(E,q)\mapsto(E,\mathsf Sq)$, for every $\mathsf S$, and hence by the intersection of these preimages, the cone $\{E\ge|q|\}$: a point with $E<|q|$ lies outside the preimage for an open set of angles of $\mathsf S$, so that countably many $\mathsf S$ suffice.
\end{proof}

\begin{lemma}\label{cov-boost}
Let $\sigma$ be labelled data with (S).
\begin{enumerate}
\item[(a)] If $\sigma(\mathsf R_tX)=\sigma(X)$ for all real $t$ at the pairs of plates, then for weight labels $\lambda,\lambda'$, $\xi,\xi'\in\mathsf{Tb}$ and real $s$: $\langle\Psi(\mathsf R_{is}\xi;\lambda),\Psi(\mathsf R_{is}\xi';\lambda')\rangle=e^{(k(\lambda)+k(\lambda'))s}\langle\Psi(\xi;\lambda),\Psi(\xi';\lambda')\rangle$.
\item[(b)] If $\sigma(\mathsf R_tX)=\sigma(X)$ for all real $t$ at all tuples of plates, there is a unitary group $U(s)$ on the closed span $\mathcal H_{\rm pl}$ of $\Omega$ and the plate vectors with $U(s)\Omega=\Omega$, $U(s)\Psi(w;\underline\lambda)=e^{-\kappa(\underline\lambda)s}\Psi(\mathsf R_{is}w;\underline\lambda)$, and $\mathsf B_\lambda(\mathsf R_{is}\xi,\mathsf R_{is}\xi')=e^{k(\lambda)s}U(s)\mathsf B_\lambda(\xi,\xi')U(-s)$ on $\mathcal H_{\rm pl}$; so $e^{-\kappa s}\|\Psi(\mathsf R_{is}z;\underline\lambda)\|=\|\Psi(z;\underline\lambda)\|$.
\end{enumerate}
\end{lemma}
\begin{proof}
At a strictly time-ordered $X$ let $I_X\ni0$ be the open interval of the real $t$ for which $\mathsf R_tX$ is time-ordered. For $\zeta=t+is$, $d(\mathsf R_\zeta\xi)=e^{-|s|}(\mathsf R_t\xi)^0$ at real $\xi$ with $(\mathsf R_t\xi)^0>0$, so $\mathsf R_\zeta\partial X\in\mathsf{Tb}^{n-1}$ for $\zeta\in I_X+i\mathbb R$, and $g(\zeta)=e^{i\kappa\zeta}\mathsf K_n(\mathsf R_\zeta\partial X;\underline\lambda)=e^{i\kappa\zeta}\mathcal K_n(\mathsf R_\zeta X;\underline\lambda)$ is holomorphic on $I_X+i\mathbb R$. By the hypothesis at plates carried by small discs about the points of $X$, $g$ is constant on $I_X$, hence on $I_X+i\mathbb R$: $\mathcal K_n(\mathsf R_{is}X)=e^{\kappa s}\mathcal K_n(X)$ for all real $s$. By Lemma~\ref{cov-plate}\,(c) the Gram matrix of $\Omega$ and the vectors $e^{-\kappa s}\Psi(\mathsf R_{is}x;\underline\mu)$ at real points does not depend on $s$ ($\langle\Omega,\Psi(z;\lambda)\rangle=\mathsf K_1=0$); at pairs this is (a), after holomorphic continuation in $\xi,\xi'$. The vectors at real points span a dense subspace of $\mathcal H_{\rm pl}$ (a vector orthogonal to them gives a holomorphic function vanishing at real points), so the isometry they define extends to a unitary group, which acts as stated at complex points by holomorphy; the formula for $\mathsf B_\lambda$ follows on the total set of plate vectors.
\end{proof}

\begin{lemma}\label{cov-scalar}
Let $\sigma$ be labelled data with (S) and (I$_0$). A \emph{scalar label} is a label $(\varphi,\mathbf1)$, written $\varphi$, of weight $0$, with the point operators $\mathsf B_\varphi=\mathsf B_{(\varphi,\mathbf1)}$; $\mathbb R^2_+=\{y\in\mathbb R^2:y^0>0\}$ is the set of the real points of $\mathsf{Tb}$. For $n\ge1$, real scalar labels $\varphi_1,\dots,\varphi_n$ and a compact set $\mathcal V\subset(\mathbb R^2_+)^{n+1}$ there is $c_{\mathcal V}<\infty$ with
\[\bigl\|\mathsf A_{\underline\varphi}(\mathsf R_{is}v_0,\dots,\mathsf R_{is}v_n)\,\Psi(\xi;\lambda)\bigr\|\le c_{\mathcal V}\|\Psi(\xi;\lambda)\|\]
for all $\underline v\in\mathcal V$, real $s$, labels $\lambda$ and $\xi\in\mathsf{Tb}$, where $\mathsf A_{\underline\varphi}(\xi_0,\dots,\xi_n)=\mathsf B_{\varphi_1}(\xi_0,\frac12\xi_1)\cdots\mathsf B_{\varphi_n}(\frac12\xi_{n-1},\xi_n)$ is the product of the scalar point operators.
\end{lemma}
\begin{proof}
Put $\mathsf A_s=\mathsf A_{\underline\varphi}(\mathsf R_{is}\underline v)$ and $P_s=\mathsf A_s^*\mathsf A_s$, a bounded positive operator. Let $R=1+\sup_{\underline v\in\mathcal V}\sum_j|v_j^1|$, let $\lambda,\lambda'$ be weight labels, and let $\eta,\eta'$ be real points of $\mathbb R^2_+$ with $R<\eta^1<\eta'^1$. \emph{The stack.} $P_s^k$ is a product of $2kn$ scalar point operators along a vertical stack $Z_k=(q_1,\dots,q_{2kn})$ of points with the period $2l\,n_0$, $l=\sum_jv_j^0$, and with abscissae $|q_i^1|<R$, so that, by Lemma~\ref{cov-plate}\,(b), (c), $m_k=\langle\Psi(\mathsf R_{is}\eta;\lambda),P_s^k\Psi(\mathsf R_{is}\eta';\lambda')\rangle$ is the kernel of the labels $\lambda^\dagger$, the $2kn$ scalar labels of the stack and $\lambda'$ at the time-ordered configuration $(\theta\eta,q_1,\dots,q_{2kn},2kl\,n_0+\eta')$, boosted by $\mathsf R_{is}$; its two tensor points lie beyond the vertical line $y^1=R$, the stack before it.

\emph{The doubled stack costs $\beta$ per point.} Reflect the periodic stack in the line $y^1=R$. For all but finitely many small angles $\epsilon$ the doubled stack, turned by $\mathsf R_\epsilon$, has its points at time gaps $\ge2\tau>0$; by continuity and the compactness of $\mathcal V$ one $\tau$ and finitely many angles serve all of $\mathcal V$, and then by (S) the scalar kernel of $M'$ points of the turned doubled stack is at most $\beta^{M'}$, $\beta=\max\{1,\max_j\|\mathsf B_{\varphi_j}(\tau n_0,\tau n_0)\|\}$. \emph{The frame $e_1$.} Take the abscissae as times in $m_k$, after small shifts of the points of the stack that make their abscissae pairwise different and $<R$: in the order of the abscissae come the points of the stack, then $\theta\eta$, then $2kl\,n_0+\eta'$, and by the Gram formula in this frame and Schwarz' inequality $|m_k|$ is at most the product of the norms of two plate vectors.

The first is made of the $2kn$ scalar plates of the stack; its squared norm is the scalar kernel of the stack together with its mirror image in $y^1=R$, which does not move with $s$, since by (I$_0$) the scalar kernels are invariant under the $\mathsf R_t$ and so, by holomorphy, under the $\mathsf R_{is}$; turned by one of the angles $\epsilon$ it is at most $\beta^{4kn}$. The second is $\mathsf B_{\lambda^\dagger}(\mathsf R_{is}u',\mathsf R_{is}\delta)\,\mathsf T(\mathsf R_{is}(0,-\Lambda_k))\,\Psi(\mathsf R_{is}\delta;\lambda')$, with $u'=(\eta^1-R,\eta^0)$ and $\delta=(\frac12(\eta'^1-\eta^1),0)$, both in $\mathbb R^2_+$ because $\eta^1>R$ and $\eta'^1>\eta^1$, and $\Lambda_k=2kl+\eta^0+\eta'^0$: a point operator of $\lambda^\dagger$, a unitary, and a one-member vector of $\lambda'$; its norm is a number $C_1$ that may depend on $s$, $\eta$, $\eta'$ in any way, but not on $k$.

So $|m_k|\le C_1\beta^{2nk}$. \emph{The root.} Let $E_s$ be the spectral measure of $P_s$; 
\[
\mu=\langle\Psi(\mathsf R_{is}\eta;\lambda),E_s(\cdot)\Psi(\mathsf R_{is}\eta';\lambda')\rangle
\]
 is a finite complex measure on $[0,\|P_s\|]$ with the moments $m_k$, and a finite complex measure on $[0,M]$ whose moments are at most $Cc^k$, $k\ge0$, has no mass on $(c,M]$; the $k$-th root removes the number $C_1$, which may depend on $s$. So 
\[
\langle\Psi(\mathsf R_{is}\eta;\lambda),\mathsf E\,\Psi(\mathsf R_{is}\eta';\lambda')\rangle=0
\]
 for $\mathsf E=E_s((\beta^{2n},\infty))$ and all real $\eta,\eta'$ with $R<\eta^1<\eta'^1$. The function 
\[
\begin{multlined}(\eta,\eta')\mapsto\\ \langle\Psi(\mathsf R_{is}\bar\eta;\lambda),\mathsf E\,\Psi(\mathsf R_{is}\eta';\lambda')\rangle\end{multlined},
\]
 $\bar\eta$ the conjugate point, is holomorphic on $\mathsf{Tb}\times\mathsf{Tb}$ and vanishes on a non-empty open set of real points; so it vanishes, and at $\lambda'=\lambda$, $\eta'=\xi$, $\eta=\bar\xi$ it gives $\mathsf E\,\Psi(\mathsf R_{is}\xi;\lambda)=0$.

As $\mathsf R_{is}$ maps $\mathsf{Tb}$ onto itself and every label is a combination of weight labels, $\mathsf E\,\Psi(\xi;\lambda)=0$ for all $\xi\in\mathsf{Tb}$ and all labels $\lambda$, and 
\[
\|\mathsf A_s\Psi(\xi;\lambda)\|^2=\langle\Psi(\xi;\lambda),P_s\Psi(\xi;\lambda)\rangle\le\beta^{2n}\|\Psi(\xi;\lambda)\|^2
\]
Thus $c_{\mathcal V}=\beta^n$.
\end{proof}

\begin{proposition}\label{cov-rot}
Let $\sigma$ be labelled data with (S), and suppose that for every $j\ge1$, every tuple $\underline\lambda$ of $j$ weight labels and every compact set $\mathcal C$ of fans of $j$ points, $\sup_{s\in\mathbb R,\,z\in\mathcal C}e^{-A_j|s|}\|\Psi(\mathsf R_{is}z;\underline\lambda)\|<\infty$, nothing being asked of the suprema, where $\liminf_jA_j/j^2<2$. If $\liminf_j(A_j-j^2)/j=-\infty$ (for instance $\liminf_jA_j/j^2<1$, or $A_j=pj$), it suffices below that (D3) hold at the diagonal plane $\Sigma$ for tuples of one member. If $\sigma(\mathsf R_tH,\mathsf R_tH')=\sigma(H,H')$ for all real $t$ and all pairs, then $\mathcal K_n(\mathsf R_{is}Z;\underline\lambda)=e^{\kappa(\underline\lambda)s}\mathcal K_n(Z;\underline\lambda)$ for all $n\ge2$, weight labels, $Z\in\mathcal Z_n$ and real $s$, and $\sigma(\mathsf RX)=\sigma(X)$ for all $\mathsf R\in\mathrm O(3)$ and $X\in\mathfrak S_V(\mathbb R^3)$.
\end{proposition}
\begin{proof}[Proof of Proposition~\ref{cov-rot}]
We first obtain a finite angular expansion. Its extreme coefficient is a positive kernel at a coordinate line and satisfies a modulus Schwarz inequality at every line. A regular polygon forces that coefficient to vanish unless its frequency satisfies a divisibility condition. The Schwarz inequality propagates this vanishing from the polygon to an arc with half as many vertices. Real analyticity then removes the other points. Iterating this argument gives the divisibility condition at every order. The resulting convex integer sequence is linear under either growth hypothesis in the statement.

\emph{Finite angular expansions.}
We may replace $A_j$ by $\max\{A_j,0\}$; either growth hypothesis in the statement is preserved. Put
\[
 \mathsf A_n=A_{\lfloor n/2\rfloor}+A_{\lceil n/2\rceil}+2n,
 \qquad A_0=0.
\]
Call $X=(y_1,\ldots,y_n)$ a pencil if
$y_{j+1}-y_j=r_jn_{\vartheta_j}$ with $r_j>0$ and
$|\vartheta_j|<\pi/16$. These configurations form a nonempty open set. If
$\zeta=t+is$, choose $\upsilon\in(\pi/2)\mathbb Z$ with
$|t-\upsilon|\le\pi/4$. In the light-cone coordinates of the tube,
\[
 \mathrm{Re}\bigl(\mathsf R_{\zeta-\upsilon}(y_{j+1}-y_j)\bigr)^\pm
 =r_je^{\mp s}\cos(\vartheta_j+t-\upsilon)>0.
\]
Thus $\mathsf R_\zeta X\in\mathcal Z_n$ for every complex $\zeta$.
For a fixed tuple of weight labels, write $\kappa=\sum_jk(\lambda_j)$. The function
\[
 h_X(\zeta)=e^{i\kappa\zeta}
       \mathcal K_n(\mathsf R_\zeta X;\underline\lambda)
\]
is entire and, by quarter-turn covariance, has period $\pi/2$.

For $|t|\le\pi/4$, cut $\mathsf R_tX$ between its $\lfloor n/2\rfloor$-th and
next points, at their midpoint. The points after the cut and the reflected points
before it are fans. Their angles have modulus less than $5\pi/16$, and as $t$
varies they range over compact sets of fans. The Gram formula in
Lemma~\ref{cov-plate}(c), the growth assumption and Schwarz's inequality give
\[
 |h_X(t+is)|\le C_Xe^{\mathsf A_n|s|}.
\]
Here $|\kappa|\le2n$ accounts for the last term in $\mathsf A_n$.
The Fourier coefficients may be computed on any horizontal line. Shifting that
line to imaginary height tending to either infinity shows that the coefficients
with $4|m|>\mathsf A_n$ vanish. Hence
$h_X(\zeta)=\sum_{4|m|\le\mathsf A_n}\mathsf c_m(X)e^{4im\zeta}$.

This expansion extends from pencils to every tube by a finite interpolation.
Let $\mathcal M_n=\{m\in\mathbb Z:4|m|\le\mathsf A_n\}$ and
$\mu=|\mathcal M_n|$. The matrix
$(e^{-4m\ell})_{0\le\ell<\mu,\,m\in\mathcal M_n}$ is a Vandermonde matrix.
Applying its inverse to
$e^{-\kappa\ell}\mathcal K_n(\mathsf R_{i\ell}Z;\underline\lambda)$ defines
holomorphic functions $\mathsf c_m$ on $\mathcal Z_n$, since each imaginary
rotation preserves each tube. On pencils they are the preceding coefficients.
Every tube contains an open set of real configurations that a permutation and
a quarter turn carry to pencils. The identity theorem on that tube therefore gives
\begin{equation}\label{cov-rot:expansion}
 e^{-\kappa s}\mathcal K_n(\mathsf R_{is}Z;\underline\lambda)
 =\sum_{m\in\mathcal M_n}\mathsf c_m(Z;\underline\lambda)e^{-4ms}
 \qquad (Z\in\mathcal Z_n,\ s\in\mathbb R).
\end{equation}
The coefficients are invariant under joint permutations of points and labels
and under translations.

\emph{Exact phases on the real tube union.}
Let $\mathcal O_n=\mathcal Z_n\cap(\mathbb R^2)^n$. Suppose that $X$ belongs
to the real part of the tube indexed by $(\mathfrak s,\upsilon)$. Write
\[
 \mathsf R_{-\upsilon}(y_{\mathfrak s(j+1)}-y_{\mathfrak s(j)})
 =r_jn_{\vartheta_j},\qquad r_j>0,\quad |\vartheta_j|<\pi/2,
\]
and put
\[
 I_X=\bigcap_j(-\pi/2-\vartheta_j,\pi/2-\vartheta_j).
\]
This is an open interval containing zero. The preceding light-cone computation
shows that $\mathsf R_{t+is}X$ stays in this same tube for
$t\in I_X$ and every real $s$. The two sides of
\eqref{cov-rot:expansion}, with $s$ replaced by $w$, are holomorphic on the
connected strip $-\mathrm{Im}\,w\in I_X$ and agree for real $w$. At $w=s-it$,
coefficient comparison in their finite exponential sums gives
\begin{equation}\label{cov-rot:phase}
 \mathsf c_m(\mathsf R_tX;\underline\lambda)
 =e^{i(4m-\kappa)t}\mathsf c_m(X;\underline\lambda)
 \qquad(t\in I_X).
\end{equation}

We need this identity when both endpoints belong to $\mathcal O_n$, even if
the intervening real orbit leaves a particular tube. Consider the configurations
with distinct points and with no two distinct difference vectors perpendicular.
They are rotation invariant and dense. Indeed their complement is the zero set
of the polynomial
\[
 Q(X)=\prod_{a<b}|y_a-y_b|^2
       \prod_{\{a,b\}\ne\{c,d\}}(y_a-y_b)\cdot(y_c-y_d),
\]
where each unordered pair occurs once in the second product. This polynomial is
not zero: at $y_a=(a,a^2)$, every scalar product in the second product is
$(a-b)(c-d)[1+(a+b)(c+d)]\ne0$.
For such a configuration, every real rotation belongs to $\mathcal O_n$.
Otherwise one pair would have equal times and another equal abscissae; those
difference vectors would be perpendicular, or the coincident pair would consist
of equal points. By \eqref{cov-rot:phase},
$t\mapsto e^{-i(4m-\kappa)t}\mathsf c_m(\mathsf R_tX)$ is locally constant
on the whole real line, and is therefore constant. Approximating any $X$ with
$X,\mathsf R_tX\in\mathcal O_n$ by this dense set proves
\eqref{cov-rot:phase} for those two endpoints.

\emph{Extension to all distinct real configurations.}
For a configuration $X$ of pairwise different points, choose $t$ so that
$\mathsf R_tX$ has pairwise different time coordinates. Such an angle exists:
each difference forbids only one angle modulo $\pi$. Define
\begin{equation}\label{cov-rot:extension}
 \widetilde{\mathsf c}_m(X;\underline\lambda)
 =e^{-i(4m-\kappa)t}\mathsf c_m(\mathsf R_tX;\underline\lambda).
\end{equation}
The endpoint phase identity makes this independent of the choice of $t$.
A single choice works in a neighborhood of any given $X$, where the right-hand
side is the restriction of a holomorphic function. Thus the extension is real
analytic on the entire collision-free real configuration space. It agrees with
$\mathsf c_m$ on $\mathcal O_n$, is invariant under translations and joint
permutations, and has phase $e^{i(4m-\kappa)t}$ under rotation by $t$ about
any centre. In particular this extension includes polygons, configurations with
ties in both coordinates, and configurations with some points fixed. We now
omit the tilde.

\emph{The extreme coefficient and Schwarz's inequality.}
Let $\bar m(n)$ be the largest $|m|$ for which a coefficient at some tuple of
$n$ labels is nonzero, taking $\bar m(n)=0$ when all such coefficients vanish, and put $\bar m(0)=\bar m(1)=0$. Reality and conjugation
of the labels exchange the extreme positive and negative frequencies. For
$a\ge1$, put $\nu(a)=\bar m(2a)$ and let
$\mathsf t_a=\mathsf c_{-\nu(a)}$ at $2a$ points.
The Gram formula, after multiplying each vector by its factor
$e^{-\kappa(S)s}$, reads
\[
 \bigl\langle e^{-\kappa(S)s}\Psi(\mathsf R_{is}S),
                e^{-\kappa(T)s}\Psi(\mathsf R_{is}T)\bigr\rangle
 =\sum_m\mathsf c_m(\theta S\cup T)e^{-4ms}.
\]
Here reflection carries the reflected conjugate labels. Multiplication by
$e^{-4\nu(a)s}$ and passage to $s\to+\infty$ make
$\mathsf G_a(S,T)=\mathsf t_a(\theta S\cup T)$ a positive Hermitian kernel
on the configurations of $a$ points in the positive half plane. Continuity
allows ties in time. At any other line $\mathsf L$, the exact phase formula
identifies the corresponding values with those at the coordinate line, up to
factors of modulus one. Consequently
\begin{equation}\label{cov-rot:schwarz}
 |\mathsf t_a(\theta_{\mathsf L}S\cup T)|^2
 \le |\mathsf t_a(\theta_{\mathsf L}S\cup S)|\,
       |\mathsf t_a(\theta_{\mathsf L}T\cup T)|.
\end{equation}
We use this modulus inequality, without asserting positivity after discarding
the phases at a tilted line. Each reflected conjugate weight label is a
nonzero scalar multiple of the original weight label; those scalar factors
can be removed in every vanishing assertion below by multilinearity.

Schwarz's inequality before taking the extreme coefficient also gives
\begin{equation}\label{cov-rot:convexity}
 2\bar m(a+b)\le\bar m(2a)+\bar m(2b).
\end{equation}
Indeed a nonzero coefficient of frequency $-\bar m(a+b)$ grows as
$e^{4\bar m(a+b)s}$, whereas the two squared norms on the right of Schwarz
have growth at most $e^{4\bar m(2a)s}$ and $e^{4\bar m(2b)s}$.
In particular $\nu(0)=0$ and
$2\nu(a)\le\nu(a-1)+\nu(a+1)$.

\emph{Halving a prescribed subconfiguration.}
For a labelled configuration $S$ with at most $2a$ points, let
$\mathcal N_a(S)$ mean that $\mathsf t_a$ vanishes on every collision-free
configuration of $2a$ labelled points containing $S$. Let $S'$ consist of
$j'\le a$ points on one open side of a line $\mathsf L$. If
$S\subset\theta_{\mathsf L}S'\cup S'$ and $\mathcal N_a(S)$ holds, we claim
that both $\mathcal N_a(S')$ and
$\mathcal N_a(\theta_{\mathsf L}S')$ hold.

Fix the reflected points $P=\theta_{\mathsf L}S'$. Add $a-j'$ variable points
$y$ and $a$ variable points $z$, with arbitrary fixed weight labels. Their
collision-free parameter space $D$ is open and connected: its complement is a
finite union of affine subspaces of real codimension two. On $D$ the function
\[
 (y,z)\longmapsto\mathsf t_a(P\cup y\cup z)
\]
is real analytic by \eqref{cov-rot:extension}. On the nonempty open subset
where $y$ is on the side of $P$ and $z$ is on the opposite side, set
$T=S'\cup\theta_{\mathsf L}y$. The configuration
$\theta_{\mathsf L}T\cup T$ contains $S$, so its top coefficient is zero.
Applying \eqref{cov-rot:schwarz} to $T$ and $z$ gives
$\mathsf t_a(P\cup y\cup z)=0$ on this open subset. The real-analytic
identity theorem gives zero on all of $D$. Since the added labels were arbitrary,
this proves $\mathcal N_a(P)$. Exchanging the two sides proves
$\mathcal N_a(S')$. This argument uses the extension on the full collision-free
domain; analyticity only on the original tube union would not suffice.

\emph{Polygons and divisibility.}
Suppose, towards a contradiction, that $a$ does not divide $2\nu(a)$, and
fix any weight label $\lambda$ of weight $k$. On a regular polygon of $2a$
vertices carrying this label, rotation by $\pi/a$ fixes the labelled
configuration. Since its total weight is $2ak$, the phase of the extreme
coefficient under this rotation is
\[
 \exp\bigl(i(-4\nu(a)-2ak)\pi/a\bigr)
 =\exp\bigl(-4\pi i\nu(a)/a\bigr)\ne1.
\]
The coefficient therefore vanishes on every such polygon.

Write its vertices as $P_j=p+r n_{\omega+j\pi/a}$. Suppose that
$\mathcal N_a$ holds for every arc of $j$ consecutive vertices and put
$j'=\lceil j/2\rceil$. An arbitrary arc
$S'=\{P_{k_0+1},\ldots,P_{k_0+j'}\}$ lies on one side of the line through
$p$ with direction $n_{\omega+(k_0+1/2)\pi/a}$. Reflection in that line sends
$P_{k_0+q}$ to $P_{k_0+1-q}$. Thus
$\theta_{\mathsf L}S'\cup S'$ is an arc of $2j'\ge j$ vertices containing
an arc for which $\mathcal N_a$ holds. The halving argument proves
$\mathcal N_a(S')$. Starting at $j=2a$ and iterating proves
$\mathcal N_a$ for every one-vertex arc. Every point is such an arc for a
suitable polygon. Hence $\mathsf t_a$ vanishes on every configuration containing
the label $\lambda$. Since $\lambda$ was arbitrary, it vanishes identically,
contrary to the choice of the extreme frequency. We have proved
\begin{equation}\label{cov-rot:divisibility}
 a\mid2\nu(a)\qquad(a\ge1).
\end{equation}

\emph{The integer sequence and the growth threshold.}
Set $\mu(a)=2\nu(a)/a$ for $a\ge1$. These are nonnegative integers by
\eqref{cov-rot:divisibility}. Convexity of $\nu$, with $\nu(0)=0$, makes
$\mu$ nondecreasing. If $\Delta_a=\mu(a+1)-\mu(a)$, the same convexity is
\[
 (a+1)\Delta_a\ge(a-1)\Delta_{a-1}\qquad(a\ge2).
\]
If $\mu$ is not constant, let $a'$ be its first index of increase. From that
index on, every $\Delta_a$ is a positive integer, and therefore at least one.
It follows that
\[
 \nu(a)\ge a\nu(1)+\tfrac12a(a-a')\qquad(a\ge a').
\]
But \eqref{cov-rot:expansion} gives $\nu(a)\le A_a/2+a$. The nonlinear
alternative would imply
\[
 \frac{A_a-a^2}{a}\ge2\nu(1)-a'-2\qquad(a\ge a'),
\]
contradicting $\liminf_a(A_a-a^2)/a=-\infty$. Under that growth hypothesis,
therefore, $\nu(a)=a\nu(1)$.

For the weaker hypothesis $\liminf_aA_a/a^2<2$, use (D3) at the diagonal
plane for tuples of every length. At the diagonal mirror, the angular expansion
has the additional factor $(-1)^m$, apart from label factors of modulus one
attached to each configuration. Positivity there therefore makes the kernel
\[
 \sum_m(-1)^m\mathsf c_m(\theta S\cup T)
\]
positive. This positivity continues to the complex points used in the boosts.
To see the continuation, use light-cone coordinates, in which the domain is a
product of right half planes. Every finite set of points in a right half plane
lies in a disc centred at a positive real point and contained in that half
plane. The Taylor expansion about such a real point expresses the continued
Gram matrices as limits of quadratic forms in derivative Gram matrices.
The latter are limits of finite-difference Gram matrices at real points and
are positive. Applying this in all coordinates gives positivity of
\[
 \left(\sum_m(-1)^me^{-4ms}
       \mathsf c_m(\theta S_i\cup S_{i'})\right)_{i,i'}
 \qquad(s\in\mathbb R).
\]
The same positive normalisations of the vectors as above do not affect the
conclusion. Its extreme limit is $(-1)^{\nu(a)}\mathsf G_a$. If $\nu(a)>0$, the kernel
$\mathsf G_a$ is nonzero: otherwise the extreme coefficient would vanish on
the open set of configurations split into $a$ points on each side, hence
identically by real analyticity. Since both $\mathsf G_a$ and
$(-1)^{\nu(a)}\mathsf G_a$ are positive, $\nu(a)$ is even. This is also true
when $\nu(a)=0$.

In the nonlinear alternative, $\mu(a)$ is consequently a multiple of four at
odd $a$: indeed $a\mu(a)=2\nu(a)$ is then a multiple of four and $a$ is odd.
Its increments are positive integers from $a'$ on. Between consecutive odd
indices it therefore increases by at least four. The possible first and last
single steps cost at most three in this lower bound, giving
\[
 \nu(a)\ge a\nu(1)+a(a-a'-\tfrac32)\qquad(a>a').
\]
Combining this with $\nu(a)\le A_a/2+a$ gives
$A_a\ge2a^2-(2a'+5)a$ and hence $\liminf_a A_a/a^2\ge2$, a contradiction.
Thus in both cases
\[
 \bar m(2a)=a\bar m(2),\qquad
 \bar m(n)\le\tfrac n2\bar m(2),
\]
the second inequality following from \eqref{cov-rot:convexity} with
$a=\lfloor n/2\rfloor$ and $b=\lceil n/2\rceil$.

\emph{Two points and arbitrary tuples.}
By the assumed covariance at pairs, the twisted two-point kernel is constant
on a real interval of rotation angles. Holomorphy continues this identity to
the corresponding angular strip. Coefficient comparison in
\eqref{cov-rot:expansion} gives $\bar m(2)=0$. Therefore every $\bar m(n)$
is zero. The coefficients with $m\ne0$ vanish on the real tube union, hence
on each complex tube, and \eqref{cov-rot:expansion} gives
\[
 \mathcal K_n(\mathsf R_{is}Z;\underline\lambda)
 =e^{\kappa s}\mathcal K_n(Z;\underline\lambda).
\]
For real time-ordered configurations, the exact phase formula with $m=0$
gives covariance under small real rotations. Integrating against plate profiles
proves covariance at time-ordered plates. Finite sums of plates approximate
weights in pairwise disjoint slabs by Weierstrass's theorem; (D1) passes the
identity to their limits. Lemma~\ref{cov-slide}(b) then gives covariance for
all tuples. The invariance group contains $W_3$ and the rotations
$\mathsf R_t$. Conjugation by $W_3$ supplies rotations about all three axes,
which generate $\mathrm{SO}(3)$, and $-\mathbf1\in W_3$ supplies the remaining
component of $\mathrm O(3)$.

The argument may also be restricted to a set $\Lambda$ of weight labels closed
under $\lambda\mapsto-\lambda$, $\lambda\mapsto\lambda^\dagger$ and
$\lambda\mapsto\bar\lambda$. Define the largest frequencies using labels
in $\Lambda$ only. The Gram argument uses the reflected conjugate labels,
the polygon argument uses one label at a time, and all vanishing assertions
are unchanged by the reflection phases. Thus the same conclusion holds with
the exponents $A_j^\Lambda$. For a finite set $\Phi$ of real transverse
profiles, this applies to $\Lambda_\Phi=\{(\varphi,\pm b):\varphi\in\Phi,
\ b\text{ in the weight basis}\}$. Each finite plate approximation uses only
one such set, so these restricted conclusions suffice for the assertion about
all tuples.
\end{proof}

\begin{corollary}\label{cov-equiv}
Let $\sigma$ be labelled data with (S) that is covariant at all pairs under the $\mathsf R_t$.
\begin{enumerate}
\item[(a)] The following are equivalent: $\sigma$ is $\mathrm O(3)$-covariant; the twisted plate vectors are bounded along the boosts, $\sup_{s,\,z\in\mathcal C}e^{-\kappa(\underline\lambda)s}\|\Psi(\mathsf R_{is}z;\underline\lambda)\|<\infty$ on compact sets $\mathcal C$ of fans; the growth hypothesis of Proposition~\ref{cov-rot} holds for some exponents with $\liminf A_j/j^2<1$; it holds with $A_j=2j$; the point operators grow at most exponentially along the boosts, $\sup_s\sup_{\mathcal K}e^{-\gamma|s|}\|\mathsf B_\lambda(\mathsf R_{is}\xi,\mathsf R_{is}\xi')|_{\mathcal H_{\rm pl}}\|<\infty$ on compact sets $\mathcal K$ of pairs of real points $rn_\alpha$, $|\alpha|<\frac{5\pi}{16}$, for some $\gamma<\infty$; it holds with $\gamma=2$. The same follows from (P) with any $p<\infty$, and from a growth $\gamma_\lambda<\infty$ of its own at every label.
\item[(b)] If $\sigma$ is not covariant and the growth hypothesis holds with some exponents, there is $a'\ge1$ with $\bar m(2a)\ge\frac12a(a-a')$, so $A_a\ge a(a-a'-2)$, for all $a\ge a'$; if $\bar m(4)\ge1$, then $a'=1$ and $\bar m(2a)\ge\binom a2$: the defect grows at least quadratically.
\end{enumerate}
\end{corollary}
\begin{proof}
By Lemma~\ref{cov-boost}\,(b), covariance gives the twisted bound with equality and the operator bound with $\gamma=2$, as $|k|\le2$; the twisted bound gives the growth with $A_j=2j$, as $|\kappa|\le2j$; such a growth gives covariance by Proposition~\ref{cov-rot}; an exponential growth $\gamma$ of the point operators gives $A_j=\gamma j$, a product of $j$ point operators making up a plate vector; and (P) gives $\gamma=p$ by Lemma~\ref{cov-plate}\,(d). A growth $\gamma_\lambda<\infty$ of its own at every label gives covariance through Proposition~\ref{cov-rot} relative to a set of labels: for a finite set $\Phi$ the set $\Lambda_\Phi$ is finite, a plate vector with $j$ labels in $\Lambda_\Phi$ is a product of $j$ point operators applied to $\Omega$, at pairs of arguments in a compact set of pairs of real points $rn_\alpha$, $|\alpha|<\frac{5\pi}{16}$, when its fan ranges over a compact set; so the growth hypothesis holds at the tuples of labels in $\Lambda_\Phi$ with $A^\Phi_j=j\gamma_\Phi$, $\gamma_\Phi=\max_{\lambda\in\Lambda_\Phi}\gamma_\lambda$, and $(A^\Phi_j-j^2)/j\to-\infty$; covariance at pairs gives $\bar m_{\Lambda_\Phi}(2)=0$; so the identity for $\mathcal K_n$ holds at all tuples of labels in $\Lambda_\Phi$, for every finite $\Phi$, and $\sigma$ is $\mathrm O(3)$-covariant. (b) is the second case of the integer argument in the proof of Proposition~\ref{cov-rot}, with $\nu(1)=\bar m(2)=0$; if $\nu(2)\ge1$, then $\mu(2)=\nu(2)>\mu(1)=0$, the differences $\Delta_a$ are $\ge1$ from $a=1$ on, and $\nu(a)\ge\binom a2$.
\end{proof}

\begin{remark}\label{cov-ex}
Of the hypotheses of Theorem~\ref{cov-thm-data}, neither (N), nor (O), nor positivity at the diagonal planes can be dropped, and covariance under $W_3$, reflection positivity and bounds do not suffice; the instances below are data, built from free fields with the standard properties of Wick polynomials (reflection positivity at every plane, Euclidean covariance, bounds at non-coinciding points), and they say nothing about the theory. (1) Let $T^{\rm W}$ be the tensor of three free photon fields and $\chi$ an independent free scalar field of mass $1$; $T^{\rm W}+2\epsilon\,\delta_{\lambda\lambda'}\partial_\lambda^2\chi$ gives labelled data with (S), (O), the properties of the scalar part, growth and an insertion bound with a power, and even the short-distance value of Theorem~\ref{cov-thm-nz}\,(i) with the error $c\epsilon^2b^2\mathsf d$; but not (N), since $\|[(H_{he_3})]\|^2=\epsilon^2\int|\tilde h|^2E^2(2E^2-1)^2d\rho_\chi>0$, and it is not covariant: the top harmonic $2$ sits at the labels of weight $-2$.

So the statement on two points for scalar plates, the two-point conclusion in the proof of Theorem~\YMMref{AK}{737}, does not carry over to labelled data: covariance under $W_3$, reflection positivity, the bounds and even the short-distance value do not give covariance at pairs; in Proposition~\ref{cov-two} conservation takes its place. With a massless $\chi$ conservation among scalar members and covariance with one tensor member among scalars hold, while (N) and covariance at pairs fail. (2) A covariant base with decay plus a constant random tensor that is diagonal, traceless and invariant in law under $W_3$ has all the hypotheses but positivity at the diagonal mirror, and is not covariant (the top harmonic is $1$): the diagonal mirror is really needed. (3) Let $\Theta(p)=\sum_\beta\mu_\beta(p)(p\times e_\beta)(p\times e_\beta)^{\rm T}$, $\mu_\beta(p)=2p^2p_\beta^2-(p^4-\sum_\lambda p_\lambda^4)$: symmetric, homogeneous of the order $6$, with $\sum_\nu p_\nu\Theta_{\nu\rho}=0$, $\operatorname{tr}\Theta=0$, covariant under $W_3$ and not under $\mathrm{SO}(3)$.

Then $T^{\rm W}+\epsilon\,\Theta(-i\partial)\chi$ gives labelled data with (D1)--(D4), (N), a covariant scalar part and a point bound with a power, but not (O): $\sigma(\theta h,h)$ grows like $\mathsf d^{-7}|h|_\infty^2$ at off-diagonal $h$ of the scale $\mathsf d$, as the leading part at short distance of the kernel at the off-diagonal label of the plane of $e_1,e_2$ has the transform $\Theta_{12}(p)^2/p^2$, which is no polynomial. Nor is it covariant at pairs: covariance under $\mathsf R$ would make the difference of two kernels a distribution carried by the origin, so that $N(p)/(p^2+1)$ would be a polynomial for $N=\Theta\otimes\Theta-\Theta^{\mathsf R}\otimes\Theta^{\mathsf R}$, homogeneous of the order $12$, $\Theta^{\mathsf R}=\mathsf R^{\rm T}\Theta(\mathsf R\,\cdot)\mathsf R$; comparing the lowest and the highest homogeneous parts gives $N=0$, so $\Theta^{\mathsf R}=\pm\Theta$, and as $\mathrm{SO}(3)$ is connected $\Theta$ would be covariant under it.

So the bound on the off-diagonal members cannot be dropped. The threshold $a<4$ of Proposition~\ref{cov-two}\,(b) is sharp up to the interval $4\le a<5$. Let $S(p)$ be the symmetric matrix with $S_{\beta\beta}=0$, $S_{12}=p_3(p_1^2-p_2^2)$, $S_{23}=p_1(p_2^2-p_3^2)$, $S_{31}=p_2(p_3^2-p_1^2)$, let $[p]_\times$ be the matrix of $v\mapsto p\times v$, and $\Theta_5=[p]_\times S[p]_\times^{\rm T}$. It is symmetric and homogeneous of the order $5$; $\sum_\nu p_\nu(\Theta_5)_{\nu\rho}=0$, as $p^{\rm T}[p]_\times=0$; $\operatorname{tr}\Theta_5=|p|^2\operatorname{tr}S-p^{\rm T}Sp=0$, as $\operatorname{tr}S=0$ and $p^{\rm T}Sp=2p_1p_2p_3\sum_{\rm cyc}(p_1^2-p_2^2)=0$; on the generators of $W_3$ one checks $S(\mathsf wp)=\det\mathsf w\,\rho(\mathsf w)S(p)$, and with $[\mathsf wp]_\times=\det\mathsf w\,\mathsf w[p]_\times\mathsf w^{\rm T}$ this gives $\Theta_5(\mathsf wp)=\det\mathsf w\,\rho(\mathsf w)\Theta_5(p)$ for $\mathsf w\in W_3$; and $\Theta_5$ is not covariant under the $\mathsf R_t$.

Put $c_e(p)=-\frac12\operatorname{Tr}(e\,\Theta_5(p))$. The two-point kernel $c_e(\partial)c_{e'}(\partial)(4\pi|u|)^{-1}$ is, up to the sign $(-1)^5$ of the odd order, the covariance of the derivatives $c_e(\partial)\chi$ of a massless free field $\chi$; under an element $\mathsf w$ of $W_3$ the two members bring the factor $(\det\mathsf w)^2=1$, and at a mirror the factor $\det\theta=-1$ that the reflection puts on the reflected member compensates that sign, so that the kernel gives pair data with covariance under $\mathbb R^3\rtimes W_3$ and positivity at every coordinate and diagonal mirror; it has (N) at pairs, as $\Theta_5$ is conserved and traceless; its scalar entries vanish, so (I$_0$) holds at pairs; and the kernel is homogeneous of the order $-11$, so that $\sigma(\theta(h\,o),h\,o)\le C\mathsf d^{-5}|h|_\infty^2$ at the scale $\mathsf d$, the bound (O$_a$) with $a=5$.

Its harmonics $m=\pm1,\pm2$ do not vanish. So the conclusion of Proposition~\ref{cov-two}\,(b) holds for every $a<4$ and fails at $a=5$. (4) $\partial_i\partial_j\phi$, $\phi$ a free field of mass $m>0$, is covariant and not conserved: (N) is not necessary.
\end{remark}

\subsection{Proofs of the theorems}\label{cov-pf}

\begin{proof}[Proof of Theorem~\ref{cov-thm-data}]
\emph{Step 1 (conservation).} Lemma~\ref{cov-slide}\,(c) gives (ii), and in particular $\sigma(H_\eta,H')=0$ at disjoint supports.
\emph{Step 2 (two points).} Proposition~\ref{cov-two}\,(a) gives $\sigma(\mathsf RH,\mathsf RH')=\sigma(H,H')$ for all $\mathsf R\in\mathrm O(3)$; under (O$_a$) and (I$_0$) at pairs, Proposition~\ref{cov-two}\,(b).
\emph{Step 3 (growth).} Lemma~\ref{cov-plate}\,(d) gives the growth of the plate vectors with $A_j=pj$, and $(A_j-j^2)/j=p-j\to-\infty$.
\emph{Step 4 (all points).} Proposition~\ref{cov-rot} gives (i) and (iii), under (P) or under any of the growths named after the theorem.
\emph{Step 5 (unitary boosts).} By (i), $\sigma$ is invariant under the $\mathsf R_t$ at all tuples of plates, and Lemma~\ref{cov-boost}\,(b) gives the unitary group $U(s)$ and the norm identity of (iv).
\end{proof}

\begin{proof}[Proof of Theorem~\ref{cov-thm-rider}]
\emph{Step 1 (one member).} By Proposition~\ref{cov-two}\,(a), $\sigma$ is covariant at pairs; by Lemma~\ref{cov-boost}\,(a), $\langle\Psi(\mathsf R_{is}\xi;\lambda),\Psi(\mathsf R_{is}\xi';\lambda')\rangle=e^{(k(\lambda)+k(\lambda'))s}\langle\Psi(\xi;\lambda),\Psi(\xi';\lambda')\rangle$ at weight labels, and at scalar labels by linearity, as $\mathbf1=\varepsilon_0+\varepsilon'_0$. For the second assertion of the theorem only this step changes: covariance at pairs is assumed.
\emph{Step 2 (a bound free of the boost).} Let $\underline\lambda=(\lambda,\varphi_1,\dots,\varphi_n,\lambda')$, two weight labels and $n$ scalar ones, $w=k(\lambda)+k(\lambda')$, and let $\mathcal K$ be the kernel of these labels, continuous on the real points of $\mathcal Z_{n+2}$. Near a fixed configuration in which the two tensor points lie far apart on the time axis and the scalar points on a short segment between them at the angle $\frac\pi4$, every rotated configuration $\mathsf R_tY$ is time-ordered after a quarter turn, with consecutive differences in a compact set $\mathcal V_0$ of $\mathbb R^2_+$; so $g_Y(\zeta)=e^{iw\zeta}\mathcal K(\mathsf R_\zeta Y)$ is entire of period $\frac\pi2$, and
\[|g_Y(t+is)|=e^{-ws}\bigl|\bigl\langle\Psi(\tfrac12\xi_0^*;\lambda^\dagger),\ \mathsf A_{\underline\varphi}(\tfrac12\xi_0,\xi_1,\dots,\xi_{n-1},\tfrac12\xi_n)\,\Psi(\tfrac12\xi_n;\lambda')\bigr\rangle\bigr|,\qquad\xi_j=\mathsf R_{is}v_j,\ v_j\in\mathcal V_0 .\]
By Lemma~\ref{cov-scalar} and Step~1 this is bounded uniformly in $t$ and $s$.
\emph{Step 3 (the kernel is covariant).} By Liouville's theorem $g_Y$ is constant; so the real-analytic function $\frac d{dt}|_{t=0}e^{iwt}\mathcal K(\mathsf R_tY)$ vanishes on an open set of real configurations, hence on the connected set of the real points of $\mathcal Z_{n+2}$, the complement of finitely many subspaces of codimension $2$; across the gaps of the set of $t$ with $\mathsf R_tY$ in it one passes by the density of the configurations at which no two differences are orthogonal. So $\mathcal K(\mathsf R_tY)=e^{-iwt}\mathcal K(Y)$ whenever $Y$ and $\mathsf R_tY$ are real points of the tubes.
\emph{Step 4 (all configurations).} $\mathcal K$ extends to exactly one real-analytic function $\tilde{\mathcal K}$ on all configurations of pairwise different points with the same law under rotations, and $\sigma$ at plates with pairwise disjoint supports is the integral of $\tilde{\mathcal K}$: the difference of the two is real-analytic in a slide along $e_1$ of the plates above a level in $y^0$ (Lemma~\ref{cov-slide}\,(a)), and vanishes once the abscissae are parted.
\emph{Step 5 (tuples).} Hence $\sigma(\mathsf R_tX)=\sigma(X)$ at tuples of plates of the class; tent partitions of unity make the members limits of sums of plates in the sup norm, with (D1), which gives the tuples whose planar projections are pairwise disjoint. For the other tuples of the class, cut the members into small pieces, so that two pieces whose planar projections meet have disjoint $x_2$-ranges, and argue by induction on the number of such pairs. For one such pair take a level $x_2=c$ between their $x_2$-ranges and outside those of all pieces, write $X=X_I\cdot X_J$ with the pieces of $X_I$ below and those of $X_J$ above the level, and put
\[F(s)=\sigma\bigl(\mathsf R_t(X_I\cdot\tau_{se_1}X_J)\bigr)-\sigma(X_I\cdot\tau_{se_1}X_J).\]
As $\mathsf R_t$ keeps $x_2$, $\mathsf R_t(X_I\cdot\tau_{se_1}X_J)=\mathsf R_tX_I\cdot\tau_{s\mathsf R_te_1}\mathsf R_tX_J$, with the unit vector $\mathsf R_te_1\perp e_2$; so Lemma~\ref{cov-slide}\,(a) at $\mu=2$, at the unit vectors $\mathsf R_te_1$ and $e_1$, makes both terms real-analytic in $s$. For large $s$ fewer pairs have meeting planar projections, so $F(s)=0$ by the induction; hence $F(0)=0$, which gives all tuples of the class.
\emph{Step 6 (the group).} The tuples $(H_1,H_2,f_1\mathbf1,\dots,f_n\mathbf1)$ are carried into such tuples by every $\mathsf R\in\mathrm O(3)$, as $\rho(\mathsf R)\mathbf1=\mathbf1$; the group of the $\mathsf R$ that keep $\sigma$ at them holds $W_3$ and the $\mathsf R_t$, hence $\mathrm O(3)$.
\end{proof}

\begin{lemma}\label{cov-dil}
Let $0<g\le\gamma$, let $\sigma$ be a tensor limit at $g$ along $(K_j,\beta_j,\rho_j)$ with label $b_*$, let $z\ge1$ be an integer with $L^zg^{-2}>2\max\{(\gamma^{\rm sd}_W)^{-2},|\bar B|\}$, and $g'=(L^zb_*/2)^{-1/2}$. Then $g'<g$, $g'<\gamma^{\rm sd}_W$, $(K_j-z,\beta_j)\in\mathfrak d(g')$ for large $j$, the images $\rho'_j$ of the $\rho_j$ under $x\mapsto L^zx$ are periodic Gibbs states, and $\sigma'(X)=\lim_j\rho'_j(A^{\rho'_j}_{K_j-z}(X))$ is a tensor limit at $g'$ with label $L^zb_*$ and
\[\sigma(X)=L^{nz}\sigma'(X^{(z)}),\qquad X^{(z)}=\bigl(H_1(L^{-z}\cdot),\dots,H_n(L^{-z}\cdot)\bigr).\]
\end{lemma}
\begin{proof}
As $b_*\ge g^{-2}$, $g'^{-2}=L^zb_*/2\ge L^zg^{-2}/2$ exceeds $g^{-2}$, $(\gamma^{\rm sd}_W)^{-2}$ and $|\bar B|$. The labels $L^za_{K_j}\beta_j$ tend to $2g'^{-2}$, which lies inside $(g'^{-2}+\bar B,X_{g'}-\bar B]$, because $X_{g'}-\bar B-2g'^{-2}=(L-2)(g'^{-2}+\bar B)>0$; so $(K_j-z,\beta_j)\in\mathfrak d(g')$ for large $j$. The dilation carries $\mathbb L_{K_j}$ onto $\mathbb L_{K_j-z}$, bonds to bonds and plaquettes to plaquettes; the Wilson weight does not depend on the spacing, so Gibbs states and invariance under lattice translations go over; and as $a_K^{-1}=L^za_{K-z}^{-1}$, $\mathcal T_K(H)$ goes over into $L^z\mathcal T_{K-z}(H(L^{-z}\cdot))$. Hence $\rho_j(A^{\rho_j}_{K_j}(X))=L^{nz}\rho'_j(A^{\rho'_j}_{K_j-z}(X^{(z)}))$, and $X\mapsto X^{(z)}$ is a bijection of $\mathfrak S_V(\mathbb R^3)$.
\end{proof}

\begin{proof}[Proof of Theorem~\ref{cov-thm}]
\emph{Step 1 (two members, at $g$).} Let $0<g\le\gamma$ and $\delta_0=\min\{\mathsf d,\frac12\ell'\}$. Cut $\eta$ and $H_2$ by partitions of unity into pieces of sup diameter $\le\frac18\delta_0$; $H_\eta=\sum_kH_{\chi_k\eta}$, as $H_\eta$ is linear in $\eta$, $\|\chi_k\eta\|_{C^2}\le C\|\eta\|_{C^2}$, and the number of pairs of pieces depends on $\mathsf d$, $R$ and $\ell'$ alone. Two pieces at sup distance $\ge\mathsf d$ are $\ge\delta_0$ apart in one coordinate, and counting the offsets in $a_{K-t_0}\mathbb Z^3$, $t_0$ as in Proposition~\ref{cov-eight}, gives a board of $\mathbb L_K$ with cells of side $\frac12\ell'$ on which they lie in two different cells, each at distance $\ge\frac1{16}\delta_0$ from the walls.

By Lemma~\ref{cov-cell}\,(c), with $Q=e^{G_g(\ell'/2)^3}$, $|\nu(A^\nu[F_1]A^\nu[F_2])|\le2Q^2\|F_1-c_1\|_{{\rm ch},\ell'}\|F_2-c_2\|_{{\rm ch},\ell'}$ for functions $F_i$ of the two cells and real $c_i$. For $F_1=\mathcal T_K(H_{\chi\eta})$ Proposition~\ref{cov-eight}\,(d) at $c=\frac1{15}$ gives the seminorm $CK^pL^{-K/15}\|\eta\|_{C^2}$, and for $F_2=\mathcal T_K(\chi'H_2)$ Proposition~\ref{cov-eight}\,(c) gives $CK^{2\mathfrak p_0}L^{K/100}|H_2|_\infty$. As $\frac1{15}-\frac1{100}>\frac1{60}$, $|\nu(A^\nu_K(H_\eta,H_2))|\le CK^pL^{-K/60}\|\eta\|_{C^2}|H_2|_\infty$, the case $q=2$ of (a) at a weaker rate; and for $H_2=H_{\eta'}$ both factors are of the first kind, and $\frac2{15}>\frac1{20}$.
\emph{Step 2 (any number of members).} Let $z$ be the least integer $\ge1$ such that the couplings $g'_1$, $g'_2$ with $g_1'^{-2}=L^zg^{-2}/2$ and $g_2'^{-2}=L^{z+1}g^{-2}/4$ are $\le\gamma^{\rm sd}_W$ and their least sides $\ell'(g'_\iota)$ are $\le\frac12L^z\mathsf d$; the least sides grow like a power of $z$, so $z$ exists. By the windows of first passage (Corollary~\YMMref{74.2}{616}) the dilation $x\mapsto L^zx$ carries $\mathfrak d(g)$ into $\mathfrak d(g'_1)\cup\mathfrak d(g'_2)$, and it carries members, states, tori, boards and seminorms exactly: the lattice, the Gibbs states and the members by Lemma~\ref{cov-dil}, the tori, the boards and the chessboard seminorms as in the step \emph{Dilation} of the proof of Proposition~\ref{cov-sd} and in Proposition~\ref{cov-torus}\,(d).

Let $g'$ be the one of $g'_1,g'_2$ at which the dilated datum $(K',\beta)$, $K'=K-z$, lies. Pieces of the dilated members of sup diameter $\le\ell'(g')/(8q)$ get, by an offset in $a_{K'-t'_0}\mathbb Z^3$, $t'_0=K'-\lceil K'/400\rceil$, distinct cells of side $\frac12\ell'(g')$, at distance $\ge\ell'(g')/(16q)$ from the walls; Lemma~\ref{cov-cell}\,(c) at $g'$ bounds a mean of centred pieces by $e^{qG_{g'}(\ell'(g')/2)^3}$ times their seminorms: $CK^pL^{-K/15}\|\eta\|_{C^2}$ for the conserved piece (Proposition~\ref{cov-eight}\,(d) at $g'$), and a number free of $K$ for each other piece. For the latter, cut the piece further, by a partition of unity, into small pieces of sup diameter $\le\frac14$, whose number depends on $\ell'(g')$, $q$, $\mathsf d$ and $R$ alone; the chessboard seminorm is subadditive, so the seminorm of the piece less its mean is at most the sum of the seminorms of the small pieces less their means.

Fix a small piece, a member $H_a$ of which it is a piece, and a direction $\mu$, and let $c+\frac12\ell'(g')\mathbb Z$, $c\in a_{K'}(\mathbb Z+\frac12)$, be the walls of the board of side $\frac12\ell'(g')$ normal to $e_\mu$. Of the two families of planes $c+\mathbb Z$ and $c+\frac12+\mathbb Z$ one keeps a sup distance $\ge\frac18$ from the projection of the small piece to the axis of $e_\mu$, an interval of length $\le\frac14$. If it is $c+\mathbb Z$, take $c+\mathbb Z$ as the walls normal to $e_\mu$: the reflections in the walls of the coarse board keep this family, so that all the eight mirror copies of the small piece keep the distance $\ge\frac18$ from it. Otherwise the projection lies within $\frac38$ of a point of $c+\mathbb Z$, and so do those of all the copies; then take $c+\frac12(1-a_{K'})+\mathbb Z$ as the walls normal to $e_\mu$, which lie in $a_{K'}(\mathbb Z+\frac12)$, as walls must, while the planes $c+\frac12+\mathbb Z$ lie in $a_{K'}\mathbb Z$, $L^{K'}$ being odd: all the copies keep a distance $\ge\frac18-\frac12a_{K'}$ from them.

This gives one board of $T_{\ell'(g')}$ with cells of side $1$ in which the corners of the plaquettes of each copy lie in the interior of one cell, at a distance $\ge\frac18-\frac32a_{K'}$ from its walls; two copies may lie in one cell. Two different copies differ in a direction in which the projection of the small piece lies at a distance $\ge\ell'(g')/(16q)$ from both walls of its cell of side $\frac12\ell'(g')$; so on the circle of length $\ell'(g')$ the reflected copy is at a distance $\ge\ell'(g')/(8q)$, and the corners of the plaquettes of different copies are at least $\ell'(g')/(8q)-2a_{K'}\ge\ell'(g')/(16q)$ apart once $a_{K'}\le\ell'(g')/(32q)$. Let $s$ be the least integer $s\ge1$ with $(\ell'(g')/(16q))L^s\ge Mu_*(8)$, a number of $g$, $\mathsf d$ and $q$, and let $K$ be so large that $K'\ge s+2$ and $a_{K'}\le\min\{\frac1{16},\ell'(g')/(32q)\}$: then $K'\ge3$ and $1\le s\le K'-2$, as Lemma~\ref{cov-mom} asks at $(K',\beta)$, and the corners of the plaquettes of different copies, at least $\ell'(g')/(16q)$ apart, are at least $(\ell'(g')/(16q))L^s\ge Mu_*(8)$ units of the level $K'-s$ apart.

The least torus at $g'\le\gamma^{\rm sd}_W\le\gamma^{\rm ref}_W$ is small over the run of $(K',\beta)$ (Lemma~\YMMref{126.3}{846}), and on it the eight copies are the first members of an admissible system of real sources under the majorants $(60|H_a|_\infty,S_\kappa)$, $S_\kappa$ the set of their plaquettes (Definition~\ref{cov-maj}, Lemma~\ref{cov-src}\,(a), (b)); as $8\le\ell'(g')^3$, Lemma~\ref{cov-mom} at $(K',\beta)$, $n=8$ and the depth $s$ gives numbers $c_\kappa$, the constants of the progression at the level $K'-s$, with $|\langle\prod_{\kappa\in B}(F_\kappa-c_\kappa)\rangle_{K',\ell'(g'),\beta}|\le(16\mathsf C_1L^{4s}\cdot60|H_a|_\infty)^{\#B}$ for every set $B$ of copies; with $m$ the common mean of the copies, $|c_\kappa-m|\le16\mathsf C_1L^{4s}\cdot60|H_a|_\infty$, and expanding $\prod_\kappa(F_\kappa-m)$ as in the step \emph{The copies} of the proof of Proposition~\ref{cov-sd} bounds the seminorm on $T_{\ell'(g')}$ of the small piece less its mean by $32\mathsf C_1L^{4s}\cdot60|H_a|_\infty$, a number free of $K$.

Centred products and joint cumulants are sums of products of such means with the conserved piece in one factor, and $\frac1{15}>\frac1{36}$, the shift of the cutoff by $z$ costing a constant. On $T_\ell$: $L^z\ell\in\mathfrak L(g')$, Lemma~\ref{cov-cell}\,(a), and the chessboard estimate on the torus itself, with numbers that depend on $\ell$. This is (a).
\emph{Step 3 (the mirror form).} For $\eta$ as in (a$'$), cut $\eta$ by a partition of unity into pieces $\chi_k\eta$ of sup diameter $\le\frac18\min\{h_0,\ell'\}$; the covariance in (a$'$) is the sum over the pairs $(k,k')$ of the covariances of $\theta_K\mathcal T_K(H^K_{\chi_k\eta})$ and $\mathcal T_K(H^K_{\chi_{k'}\eta})$, and the number of pairs depends on $h_0$, $R_0$ and $\ell'$ alone. The two members of a pair lie on the two sides of the plane $x_3=\frac12a_K$, at distance $\ge2h_0-a_K$ from each other, and counting the offsets gives a board of $\mathbb L_K$ with cells of side $\frac12\ell'$ on which they lie in distinct cells, at a distance from the walls fixed by $h_0$ and $\ell'$. By \eqref{cov-eq-var}, which is linear in $\eta$, the null property and Proposition~\ref{cov-eight}\,(a), (b), the eight-copy seminorm of $b\,\mathcal T_K(H^K_{\chi_k\eta})$ less its mean is $\le CK^pL^{-cK}\|\eta\|_{C^2}$ for every $c<\frac1{10}$, and so is that of its image under the lattice isometry $\theta_K$; Lemma~\ref{cov-cell}\,(c) bounds each covariance by $2e^{2G_g(\ell'/2)^3}$ times the product of the two seminorms, which gives the upper bound. The covariance is $\nu(\theta_KF\cdot F)$ with $F=\mathcal T_K(H^K_\eta)-\nu(\mathcal T_K(H^K_\eta))$, as $\nu$ is invariant under $\theta_K$, and reflection positivity at the bond mirror $x_3=\frac12a_K$ (Per) gives the lower bound.
\emph{Step 4 (the four properties).} Let $g\le\gamma^{\rm sd}_W$. By Propositions~\ref{cov-lim}, \ref{cov-od}, \ref{cov-P} a tensor limit $\sigma$ at $g$ is labelled data with (S), (O) and (P) for every $p>\frac{17}2$. For $\eta\in C^2_c(\mathbb E_+;\mathbb R^3)$ the supports of $\theta H_\eta=H_{\theta\eta}$ and $H_\eta$ are at positive distance, and (a) at $q=2$ gives (N): $\sigma(\theta H_\eta,H_\eta)=\lim_j\rho_j(A^{\rho_j}_{K_j}(\theta H_\eta,H_\eta))=0$.
\emph{Step 5 ((b), (c) at $g\le\gamma^{\rm sd}_W$).} Theorem~\ref{cov-thm-data}\,(i), (ii).
\emph{Step 6 (all reference couplings).} Let $g\le\gamma$ and $\sigma'$, $g'$, $z$ as in Lemma~\ref{cov-dil}. Rotations commute with the dilation, $(\mathsf RX)^{(z)}=\mathsf RX^{(z)}$, and $H_\eta(L^{-z}\cdot)=H_{\eta^{[z]}}$ with $\eta^{[z]}(x)=L^z\eta(L^{-z}x)$; by Step~5 at $g'$,
\[\sigma(\mathsf RX)=L^{nz}\sigma'(\mathsf RX^{(z)})=L^{nz}\sigma'(X^{(z)})=\sigma(X),\qquad\sigma((H_\eta)\cdot X)=L^{(n+1)z}\sigma'\bigl((H_{\eta^{[z]}})\cdot X^{(z)}\bigr)=0 .\]
\end{proof}

\subsection{The limits are not zero}\label{cov-scope}

Theorem~\ref{cov-thm-nz} shows that the tensor limits at $0<g\le\gamma^{\rm sd}_W$, to which Theorems~\ref{cov-thm}, \ref{cov-thm-data} and~\ref{cov-thm-rider} and Corollary~\ref{cov-cor-real} apply, are not zero, and that each of the six fields, and every real combination of them, survives in every such limit. By Lemma~\ref{cov-dil} the same holds for every tensor limit at every reference coupling $0<g\le\gamma$, where Theorem~\ref{cov-thm} applies. Let $\sigma$ be one at $g$, with the label $b_*$; let $z$, $g'=(L^zb_*/2)^{-1/2}<\gamma^{\rm sd}_W$ and $\sigma'$, a tensor limit at $g'$ with the label $L^zb_*$, be as in Lemma~\ref{cov-dil}; choose $0<\mathsf d'<\mathsf d_0(g')$ with $e_{g'}(\mathsf d')<w_0$, the bump $\varphi'$ of Theorem~\ref{cov-thm-nz}\,(iii) at $\mathsf d'$, and $E\in V$, $E\ne0$; and put $H=(\varphi'E)(L^z\,\cdot)$, a member carried by $\mathbb E_+$ with $H^{(z)}=\varphi'E$. As $\theta$ commutes with the dilation, $(\theta H)^{(z)}=\theta(\varphi'E)$, and Theorem~\ref{cov-thm-nz}\,(iii) at $g'$ gives
\[
\begin{multlined}
\|[(H)]\|^2=\sigma(\theta H,H)=L^{2z}\sigma'(\theta(\varphi'E),\varphi'E)\ge L^{2z}\bigl(w_0-e_{g'}(\mathsf d')\bigr)(L^zb_*)^{-2}\operatorname{tr}E^2=\bigl(w_0 \\
-e_{g'}(\mathsf d')\bigr)b_*^{-2}\operatorname{tr}E^2>0 ;
\end{multlined}
\]
here $g'$ and $\mathsf d'$ depend on $b_*$.

\begin{proof}[Proof of Theorem~\ref{cov-thm-nz}]
(i) Let $s_{16}(g)$, an integer formed with the constants and $g$, be the least depth from which on the thresholds of the factorisation at sixteen copies (below), that of Lemma~\ref{cov-pairs} at eight pairs with the copies centred at their means, hold, and let $\mathsf d_0(g)=L^{-s_{16}(g)}$; for $0<\mathsf d<\mathsf d_0(g)$ the depth $s=s(\mathsf d)$ fitted to $\mathsf d$ exceeds $s_{16}(g)$. Let $\ell'$ be the least side at $g$ and $\mathsf l=\ell'/2$. We work first at a cutoff: let $(K,\beta)\in\mathfrak d(g)$ with $K\ge s+2$ and $4a_K\le\mathsf d-\mathsf d_-$, and $b=a_K\beta$. The working level is $\mathsf k=K-s$ (equation~\YMMref{(2322)}{844}), and $\upsilon=(a_sx_{K-s}(\beta))^{-1}$, the number $\upsilon$ of Lemma~\ref{cov-mom8} at $g$ and $\mathsf d$, with $s_{\mathsf k}=s$ as $K\ge s+2$, is the coupling of the working level in the physical scale, the natural size there of a source of unit majorant; by Lemmas~\YMMref{126.1}{844} and~\ref{cov-mom8}, $|\upsilon^{-1}-b|\le\bar Ba_s$, $\upsilon^2\le\frac98b^{-2}$, $|\upsilon^2-b^{-2}|\le3\bar Bg^2L^{2-s}b^{-2}$, $8a_K\le\mathsf d$, $a_s\le\mathsf d$, $20\mathsf d\le\ell'$ and $Y_s\ge s\log L\ge\log\frac1{\mathsf d}$.

Put $F_\iota=\mathcal T_K(H_\iota)$, and on the least torus $c_\iota=\langle F_\iota\rangle_{K,\ell',\beta}$, $\tilde F_\iota=F_\iota-c_\iota$, and the \emph{pair value} $m=\langle\tilde F_1\tilde F_2\rangle_{K,\ell',\beta}=\langle F_1;F_2\rangle_{K,\ell',\beta}$.
\emph{The sixteen copies.} Consider the board of the least torus with eight cells of side $\mathsf l$ and an offset within $\frac12a_K$ of the midpoint of $x_0$ and $y_0$: the two members lie in the middle of one cell, every corner of a plaquette on which $F_\iota$ depends lying within $3\mathsf d<\frac12\mathsf l$ of its centre, and their sixteen mirror copies $F_{(\kappa,\iota)}$, $\kappa\in\{0,1\}^3$, with the majorants $(60|H_\iota|_\infty,S_{H_\iota})$ of Lemma~\ref{cov-src}, are eight close pairs $\tau=\kappa$: the two copies of one pair are $3\mathsf d$ to $4\mathsf d$ apart, and the corners of the plaquettes of different pairs are at sup distances $\ge\mathsf l-6\mathsf d\ge11M$. By the invariance of the torus measure under the reflections every copy of $F_\iota$ has the mean $c_\iota$, and every pair the value $m$.
\emph{Quadratic sources; the Gaussian value.} For real $H$ let $\mathsf t_H=\mathsf t[w_H]+\mathsf t[\mathsf w_H]$, the kernel of Lemma~\ref{cov-src}\,(c), a real symmetric kernel on the pairs of bare plaquettes, where
\[\mathsf t[w](p,p')=\tfrac12w(p)\delta_{pp'},\qquad\mathsf t[\mathsf w](p,p')=-\tfrac12\sum_b\mathsf w(b)\bigl[c_{b,\lambda_b}(p)c_{b,\lambda'_b}(p')+c_{b,\lambda_b}(p')c_{b,\lambda'_b}(p)\bigr],\]
$c_{b,\lambda}(p)\in\{0,\pm\frac12\}$ being the coefficients of $\mathbf u_{b,\lambda}$, to the first order, as a linear form in the vector parts of the plaquettes; they are carried by the two plaquettes of the plane of $\lambda$ and $\mu_b$ that contain $b$. The kernel vanishes off $S_H\times S_H$, and $2\sum_{p'}|\mathsf t_H(p,p')|\le3|H|_\infty$. With this kernel $\mathcal T_K(H)$ is a \emph{quadratic source} under $(60|H|_\infty,S_H)$ (Definition~\ref{cov-maj}\,(5), Lemma~\ref{cov-src}\,(c)): besides the properties of a source under a majorant that Lemma~\ref{cov-src} names, at a small point of the chart of every step its second order is the quadratic form $\sum_{p,p'}\mathsf t_H(p,p')\langle\mathsf{ang}_p,\mathsf{ang}_{p'}\rangle$ in the angles of the plaquettes, up to a remainder of a higher order in the radius times $(\#S_H)\,60|H|_\infty$.

For $\mathcal O(w)$ this is the expansion of $a_p$ to the second order; for $Y[\mathsf w]$, $q_b=\langle\mathbf u_{b,\lambda_b},\mathbf u_{b,\lambda'_b}\rangle$ is, to the second order, the product of the two linear forms. The carried family at the flat top (Proposition~\YMMref{125.1}{838}) is worded for the plaquette sources, with the kernels $\frac12w(p)\delta_{pp'}$, and its proof reads a source only through this quadratic form and the row sums of its kernel; so it holds for kernels with two indices, with the same constants and the same errors, read at the majorants. With it the law of two points (Proposition~\YMMref{126.4}{846}) holds for two real quadratic sources under majorants $(\mathsf s_\iota,S_\iota)$ whose plaquettes lie within $\mathsf d-a_K$ of $x_0$ and of $y_0$: $|\langle F_1;F_2\rangle_{K,\ell,\beta}-G|\le\epsilon^\natural_g(s)\,\upsilon^2\mathsf s_1\mathsf s_2$ on every side $\ell$ that is small over the run, with a function $\epsilon^\natural_g$ of the constants, $g$ and $s$ that tends to $0$, where
\[G=6\upsilon^2\sum_{p,p',q,q'}\mathsf t_1(p,p')\,\mathsf t_2(q,q')\,C(p,q)\,C(p',q')\]
is the \emph{Gaussian value} of the two sources with the kernels $\mathsf t_1$, $\mathsf t_2$. Here $6=2(N^2-1)$ at $N=2$, and $C=C^{\rm free}_K-C^{\rm fl}_{K,s}$, the background kernel $C^{\rm bg}_{K,s}$ of the depth $s$ of the working level (Section~\YMMref{120}{817}, Lemma~\ref{cov-bg}), is the covariance of the background of the chain at the depth $s$, on pairs of bare plaquettes and in the bare units: with the orientation signs $\epsilon_p=\pm1$ of the plaquettes and $u=x(p)-x(q)$, the free part differs from $a_K^3K_3(u)_{\lambda(p)\lambda(q)}$, after the signs $\epsilon_p\epsilon_q$, by at most $C_{\rm lat}a_K^3[a_K^2|u|^{-5}+\ell'^{-3}]$, and the part of the fluctuations is at most $C^{\rm fl}_*a_K^3a_s^{-3}e^{-\delta_{\rm tp}L|u|/(\sqrt3a_s)}$, with numbers $C_{\rm lat}$, $C^{\rm fl}_*$, $\delta_{\rm tp}$ of the constants. At the plaquettes of our two members $|C(p,q)|\le C_{\mathsf k}(a_K/\mathsf d)^3$, $C_{\mathsf k}$ a number of the constants, that of Lemma~\ref{cov-bg}, and $\#S_{H_\iota}\le81(\mathsf d/a_K)^3$; so the Gaussian value $G_\tau$ of every pair $\tau$ has $|G_\tau|\le\mathsf g'_0\hat m$, with $\mathsf g'_0=\frac32\,81^2C_{\mathsf k}^2$ and $\hat m=3600\upsilon^2|H_1|_\infty|H_2|_\infty$; the pair of the members themselves has $G_0=G$ at $\mathsf t_\iota=\mathsf t_{H_\iota}$.
\emph{Factorisation, centred at the means.} For a pair $\tau$ let $\tilde P_\tau=\tilde F_{(\tau,1)}\tilde F_{(\tau,2)}$, and for a set $J$ of pairs $\tilde P_J=\prod_{\tau\in J}\tilde P_\tau$. The mean on the least torus of a product of centred copies is, as in the proof of Lemma~\ref{cov-pairs}, a sum over the level of the first creation of a large field (equation~\eqref{eq:ward-first-creation}); along the histories without large fields the conditional density at the upper level is a sum over the partitions of the copies, the blocks of one element being the conditional means less the true means, those of the two copies of one pair their conditional covariance, and those that link two pairs terms of the exponent on polymers joining them, exponentially small in their distance in units of the working level; a block of two or more elements reads no centring constant.

In the place of a vanishing mean one reads: the conditional mean of $F_\iota$ at the upper level is the constant $c'_\iota$ of the progression up to a small number, and the sum over the level of the first large field, read for $F_\iota-c'_\iota$, puts the true mean $c_\iota$ within the same small number of $c'_\iota$; so a block of one element is at most $\mathsf e^{\rm c}_g(s)\,60\upsilon|H_\iota|_\infty$, with $\mathsf e^{\rm c}_g(s)\to0$. A block of a close pair is its Gaussian value up to a small multiple of $\hat m$, by the flat top for kernels with two indices. On the least side, which is small over the run (Lemma~\YMMref{126.3}{846}), the large fields are rare (Lemma~\YMMref{125.3}{841}) and pay for the levels below the upper one.

So, for every set $J$ of pairs,
\[\Bigl|\langle\tilde P_J\rangle_{K,\ell',\beta}-\prod_{\tau\in J}G_\tau\Bigr|\le\omega'_{g,16}(s)\,\hat m^{\#J},\qquad\bigl|\langle\tilde P_J\rangle_{K,\ell',\beta}-m^{\#J}\bigr|\le\omega'^\sharp_{g,16}(s)\,\hat m^{\#J},\]
the second because every pair has the value $m$; and for the eight copies of one member, one of each pair, $|\langle\prod_\kappa\tilde F_{(\kappa,\iota)}\rangle_{K,\ell',\beta}|\le(\mathsf{mn}_g(s)\,60\upsilon|H_\iota|_\infty)^8$. The numbers $\omega'_{g,16}(s)\le\omega'^\sharp_{g,16}(s)$ and $\mathsf{mn}_g(s)$ are formed with the constants, $g$ and $s\ge s_{16}(g)$; they tend to $0$ as $s\to\infty$, and $Y_s^{2\mathfrak p_0}\omega'^\sharp_{g,16}(s)$ is bounded.
\emph{The lattice sum.} As $C$ is $K_3$ up to the lattice and volume corrections, and $a_K^3\sum_{p,p'}|\mathsf t_H(p,p')|\le45\mathsf d^3|H|_\infty$,
\[
\begin{multlined}
\Bigl|6\sum_{p,p',q,q'}\mathsf t_{H_1}(p,p')\mathsf t_{H_2}(q,q')C(p,q)C(p',q')-\mathcal W(H_1,H_2)\Bigr|\le\Bigl[12150E_1\bigl(\tfrac1\pi+E_1\bigr) \\
+7400\frac{a_K}{\mathsf d}\Bigr]|H_1|_\infty|H_2|_\infty+\mathsf{Ri}_K,
\end{multlined}
\]
with $E_1=C_{\rm lat}[(a_K/\mathsf d)^2+(\mathsf d/\ell')^3]+C^{\rm fl}_*\check{\mathsf d}^3e^{-\delta_{\rm tp}L\check{\mathsf d}/\sqrt3}$, $\check{\mathsf d}=L^s\mathsf d$, and the Riemann remainder $\mathsf{Ri}_K\le310[\omega_1(a_K)|H_2|_\infty+|H_1|_\infty\omega_2(a_K)]+12800\frac{a_K}{\mathsf d}|H_1|_\infty|H_2|_\infty$, $\omega_\iota$ the modulus of continuity of $H_\iota$; and $|\mathcal W(H_1,H_2)|\le24|H_1|_\infty|H_2|_\infty$.
\emph{All states and sides.} Let $\mathfrak s$ be any $\nu\in\mathcal G^{\rm per}_{K,\beta}$, or the torus measure of any side $\ell\in\mathfrak L(g)$. For every real function $D$ of the cell, $|\mathfrak s(D)|\le e^{G_g\mathsf l^3}\|D\|_{{\rm ch},\ell'}$ (Lemma~\ref{cov-cell}\,(b), (c)), and $\|D\|^8_{{\rm ch},\ell'}$ is the mean on the least torus of the eight mirror copies of $D$. At $D=F_\iota-c_\iota$ these are the eight centred copies of one member, one of each pair: $|\mathfrak s(F_\iota)-c_\iota|\le e^{G_g\mathsf l^3}\mathsf{mn}_g(s)\,60\upsilon|H_\iota|_\infty$. At $D=(F_1-c_1)(F_2-c_2)-m$: $\|D\|^8_{{\rm ch},\ell'}=\sum_J(-m)^{8-\#J}\langle\tilde P_J\rangle_{K,\ell',\beta}$, the sum over the sets $J$ of pairs; with $\langle\tilde P_J\rangle_{K,\ell',\beta}=m^{\#J}+r_J$, $|r_J|\le\omega'^\sharp_{g,16}(s)\hat m^{\#J}$, the terms $m^8\sum_J(-1)^{8-\#J}=m^8(1-1)^8$ cancel, and $|m|\le(\mathsf g'_0+\omega'_{g,16}(s))\hat m$ bounds the rest by $\omega'^\sharp_{g,16}(s)(1+\mathsf g'_0+\omega'_{g,16}(s))^8\hat m^8$; the eighth root gives $|\mathfrak s((F_1-c_1)(F_2-c_2))-m|\le\mathsf{Df}'_g(s)\,\hat m$ with $\mathsf{Df}'_g(s)=e^{G_g\mathsf l^3}\omega'^\sharp_{g,16}(s)^{1/8}(1+\mathsf g'_0+\omega'_{g,16}(s))$. As $\mathfrak s(F_1;F_2)=\mathfrak s((F_1-c_1)(F_2-c_2))-\mathfrak s(F_1-c_1)\,\mathfrak s(F_2-c_2)$,
\[\bigl|\mathfrak s(F_1;F_2)-m\bigr|\le\bigl[\mathsf{Df}'_g(s)+e^{2G_g\mathsf l^3}\mathsf{mn}_g(s)^2\bigr]\hat m .\]
\emph{The law at a cutoff.} With $|m-\upsilon^2X|\le\omega'_{g,16}(s)\hat m$ for the sum $X=6\sum\mathsf t_{H_1}\mathsf t_{H_2}CC$ (the case $J=\{0\}$), the lattice sum, $\upsilon^2\le\frac98b^{-2}$, $|\upsilon^2-b^{-2}|\le3\bar Bg^2L^{2-s}b^{-2}$ and $|\mathcal W|\le24|H_1|_\infty|H_2|_\infty$:
\[
\begin{multlined}
\bigl|\mathfrak s(F_1;F_2)-b^{-2}\mathcal W(H_1,H_2)\bigr|\le b^{-2}\Bigl\{e_g(\mathsf d)+\tfrac98\Bigl[12150\Bigl(E_1\bigl(\tfrac1\pi+E_1\bigr)-E^\infty_1\bigl(\tfrac1\pi+E^\infty_1\bigr)\Bigr) \\
+7400\frac{a_K}{\mathsf d}\Bigr]\Bigr\}|H_1|_\infty|H_2|_\infty+\tfrac98b^{-2}\mathsf{Ri}_K,
\end{multlined}
\]
where $E^\infty_1=E_1-C_{\rm lat}(a_K/\mathsf d)^2$ and
\[e_g(\mathsf d)=\tfrac98\Bigl\{3600\bigl[\omega'_{g,16}(s)+\mathsf{Df}'_g(s)+e^{2G_g\mathsf l^3}\mathsf{mn}_g(s)^2\bigr]+12150E^\infty_1\bigl(\tfrac1\pi+E^\infty_1\bigr)\Bigr\}+72\,\bar Bg^2L^{2-s},\qquad s=s(\mathsf d).\]
The function $e_g$ is formed with the constants and $g$ alone: it reads none of $K$, $\beta$, the state, the side and the members. The other terms read $K$ and the members, and they tend to $0$ as $K\to\infty$ at continuous members.
\emph{The rate.} As $Y_s^{2\mathfrak p_0}\omega'^\sharp_{g,16}(s)$ is bounded, $\omega'_{g,16}(s)\le\omega'^\sharp_{g,16}(s)\le cY_s^{-2\mathfrak p_0}$, and the eighth root gives $\mathsf{Df}'_g(s)\le c'Y_s^{-\mathfrak p_0/4}$; $\mathsf{mn}_g(s)^2$ decays faster in $Y_s$, and so do $L^{2-s}$ and $E^\infty_1$, as $\mathsf d<LC_W(s)L^{-s}$ and $\check{\mathsf d}\ge C_W(s)\to\infty$. With $Y_s\ge\log\frac1{\mathsf d}$ this gives $\sup_{\mathsf d}(\log\frac1{\mathsf d})^{\mathfrak p_0/4}e_g(\mathsf d)<\infty$, and $e_g(\mathsf d)\to0$ as $\mathsf d\to0$, the depth $s(\mathsf d)$ tending to infinity.
\emph{The limit.} Let $\sigma$ be a tensor limit along $(K_j,\beta_j,\rho_j)$, with the label $b_*$. For large $j$, $K_j\ge s+2$ and $4a_{K_j}\le\mathsf d-\mathsf d_-$, and the bound at $\mathfrak s=\rho_j$ bounds $\rho_j(A^{\rho_j}_{K_j}(H_1,H_2))=\rho_j(F_1;F_2)$; as $j\to\infty$, $a_{K_j}\beta_j\to b_*$, $E_1\to E^\infty_1$, $a_{K_j}/\mathsf d\to0$ and $\mathsf{Ri}_{K_j}\to0$, and the definition of a tensor limit gives (i).
(ii) $\mathcal W$ is invariant under the dilations of the members; so for fixed $H_1,H_2$ with $\mathcal W(H_1,H_2)\ne0$, (i) at the dilated members $H_\iota(\cdot/\epsilon)$ gives $b_*^{-2}=\lim_{\epsilon\to0}\sigma(H_1(\cdot/\epsilon),H_2(\cdot/\epsilon))/\mathcal W(H_1,H_2)$.
(iii) A direct computation gives $\mathcal W(\theta(\varphi E),\varphi E)\ge w_0\operatorname{tr}E^2$ for the bump $\varphi$ and the mirror image at the distance $3\mathsf d$; (i) at $H_1=\theta(\varphi E)$, $H_2=\varphi E$, with $|E|_\infty^2\le\operatorname{tr}E^2$, gives the bound, and the left side is $\|[(\varphi E)]\|^2$.
(iv) The same argument on $T_\ell$, with the torus functionals of Proposition~\ref{cov-torus}, at every label $\mathsf b\in(g^{-2}+\bar B,X_g-\bar B]$, for which $(K,L^K\mathsf b)\in\mathfrak d(g)$ at every $K$.
\end{proof}

\subsection{The fields in real time}\label{cov-rt}

\begin{proof}[Proof of Corollary~\ref{cov-cor-real}]
(A) Invariance under all isometries: Theorem~\ref{cov-thm}\,(c), (D2), and Lemma~\ref{cov-C1}, which also identifies $\mathcal H(\sigma^1)$ with $\mathcal H(\sigma)$. A rotation $\mathsf S$ about $e_3$ keeps $\mathbb E_+$ and commutes with $\theta$, so $[F]\mapsto[\mathsf SF]$ is unitary on $\mathcal H(\sigma)$ and turns $e^{-tH}e^{iv\cdot P}$ into $e^{-tH}e^{i(\mathsf Sv)\cdot P}$; so, with $U_{\mathsf S}[F]=[\mathsf SF]$ and $\hat{\mathsf S}(E,q)=(E,\mathsf Sq)$, $U_{\mathsf S}\mathcal P(B)U_{\mathsf S}^{-1}=\mathcal P(\hat{\mathsf S}B)$ for every Borel set $B$. $\mathcal P$ is carried by the square cone (D4), hence by its images under all the $\hat{\mathsf S}$, and so by their intersection, the cone $\{E\ge|q|\}$, as in Proposition~\ref{cov-two}\,(c). The insertion bound is the equivalent form of Proposition~\ref{cov-P}, with $b_*\ge g^{-2}$.
(B) The construction of Lemma~\ref{cov-plate}\,(a) in three variables, with the semigroup $\mathsf T_3(\mathfrak z)=\int e^{-\mathfrak z_3E+i(\mathfrak z_1q_1+\mathfrak z_2q_2)}d\mathcal P$ on the tube $\mathsf{Tb}^\square=\{\operatorname{Re}\mathfrak z_3>|\operatorname{Im}\mathfrak z_1|+|\operatorname{Im}\mathfrak z_2|\}$ of the proof of Proposition~\ref{cov-two}\,(b) and the mean value property on a polydisc of $\mathbb C^3$, gives point operators of the members $e\,\delta_x$, which (P) bounds by $C|e|\min\{1,\cdot\}^{-p}$ as in Lemma~\ref{cov-plate}\,(d); so at time-ordered configurations with consecutive time differences $\ge2\eta$ the kernel $\mathbf W_n$ of $\sigma$ is continuous and $|\mathbf W_n|\le(C_V\eta^{-p})^n\prod|e_k|$.

For $n$ pairwise different points the directions $u\in S^2$ with $|u\cdot(x_i-x_j)|\ge n^{-2}|x_i-x_j|$ for all $i\ne j$ fill more than half of the sphere; a rotation that turns such a $u$ into $e_3$ makes the configuration time-ordered with gaps $\ge n^{-2}\delta(x)$, and by (A) the value $\mathsf s_n(x;e)=\mathbf W_n$ at the rotated configuration does not depend on the rotation chosen. This gives the bound; real-analyticity comes from the holomorphy of the kernels in the tubes, and the representation of $\sigma$ at all tuples from the slides alone: Lemma~\ref{cov-slide}\,(a), with the induction on the pairs of pieces with overlapping heights of the proof of Lemma~\ref{cov-slide}\,(b), as in Step~4 of the proof of Theorem~\ref{cov-thm-rider}.

For $\varphi$ with components in $\mathcal S_{\ne}(\mathbb R^{3n})$, flat on the coincidence set, Taylor's formula at that set pays for the factor $\max\{1,2n^2/\delta(x)\}^{pn}$, at the cost of $(n!)^{2p}$ and of $\lceil p\rceil n$ derivatives; this is the growth condition with the constants $(\lceil p\rceil,1,\mathsf C_E,2p)$, $\mathsf C_E=6\pi^22^pe^{2p}3^{\lceil p\rceil}\mathsf c\,C_V$.
(C) The reconstruction theorem among the construction results, Theorem~\textup{\YMMref{AO}{778}}, is used in its form for members with values in $V$: the labelled form stated, with the full list of its changes, at the end of Subsection~\ref{cov-start}. It is applied to $\sigma^1$, which has the data (Rec1)--(Rec4) by (A). Its data are a real symmetric functional $\sigma^1$ on the tuples of members in $C^1_c(\mathbb R^3;V)$ with pairwise disjoint supports, $\sigma^1(\varnothing)=1$, $n$-linear on the products of the spaces $C^1_c(O_a;V)$ for open sets with pairwise disjoint closures, with (Rec1)$_V$ $\sigma^1(\mathsf rX)=\sigma^1(X)$ for every isometry $\mathsf r$ of $\mathbb R^3$, acting on points and values; (Rec2)$_V$ positivity at $x_3=0$, a separable space $\mathcal H(\sigma^1)$, and $H\ge0$, the translations $e^{iv\cdot P}$ and the joint spectral measure $\mathcal P$ as in (D4); (Rec3)$_V$ $\mathcal P(\{E<|q|\})=0$; (Rec4)$_V$ the insertion bound $|\sigma^1(\theta\bar F\cdot(h)\cdot F')|\le C_V\eta^{-p_V}\|h\|_{L^1}\|[\tau_{-\eta e_3}F]\|\,\|[\tau_{-\eta e_3}F']\|$ for $\eta\in(0,1]$, $F,F'$ over $\{x_3>\eta\}$ and $h\in C^1_c(\{|x_3|<c_0\eta\};V)$, the norm $\|h\|_{L^1}$ being formed with a norm of $V$ that $\rho$ keeps.

For such data the theorem gives: a strongly continuous unitary representation $U$ of the proper orthochronous Poincar\'e group of $\mathbb R^{1+2}$ with $U(x,1)=e^{itH}e^{i\vec x\cdot P}$, $U\Omega=\Omega$ and spectrum in the closed forward cone; fields $T(f)$, $f\in\mathcal S(\mathbb R^{1+2};V_{\mathbb C})$, linear in $f$, on the span $\mathcal D_0$ of $\Omega$ and of the vectors $T(f_1)\cdots T(f_n)\Omega$, tempered, with $\langle\psi,T(f)\psi'\rangle=\langle T(f^*)\psi,\psi'\rangle$, $f^*=\rho(\theta)\bar f$; $U(a,\Lambda)T(f)U(a,\Lambda)^{-1}=T(D(\Lambda)f(\Lambda^{-1}(\cdot-a)))$ on $\mathcal D_0$; locality of all components with all components; $\mathcal D_0$ dense; and functions $\mathbf W_n(\mathfrak z;e)$, holomorphic on the tubes and $n$-linear in $e\in V_{\mathbb C}^n$, bounded by $(C_V\eta^{-p_V})^n\prod|e_k|$ where the consecutive differences have the margin $2\eta$, $0<\eta\le1$, whose boundary values are the Wightman functions and which at real configurations ordered in $x_3$ are the density of $\sigma^1$.

Here $D(\Lambda)A=\hat\Lambda A\hat\Lambda^{\rm T}$, $A\in V_{\mathbb C}$, $\hat\Lambda=\iota\Lambda\iota^{-1}$: $\hat\Lambda^{\rm T}\hat\Lambda=\mathbf1$, $\det\hat\Lambda=1$ and $\widehat{\Lambda\Lambda'}=\hat\Lambda\hat\Lambda'$, so that $D$ is a continuous homomorphism of the Lorentz group into $\mathrm{GL}(V_{\mathbb C})$; $D(1\oplus\mathsf R)=\rho(\mathsf R\oplus1)$ for $\mathsf R\in\mathrm{SO}(2)$; $D(B(s))=\exp(is\,d\rho(\omega_1))$ for the boost $B(s)$ in the plane of $x^0$ and $x^1$, $\omega_1=e_1e_3^{\rm T}-e_3e_1^{\rm T}$, $d\rho(\omega)A=\omega A+A\omega^{\rm T}$; and $D(\Lambda)e^*=(D(\Lambda)e)^*$ with $e^*=\rho(\theta)\bar e$, as $\theta\omega_1\theta^{-1}=-\omega_1$. The proof is that among the construction results, index by index: the point operators become linear in a value $e$, with $\mathsf B(s,s';e)^*=\mathsf B(s',s;e^*)$, and the rotations $\mathsf R$ about $e_3$ act on them by $e\mapsto\rho(\mathsf R)e$; the kernels are permuted together with their values, and $\mathsf s_n(\mathsf rx;\rho(\mathsf R)e)=\mathsf s_n(x;e)$ for $\mathsf rx=\mathsf Rx+v$; the bounds keep their constants because $\rho(\mathsf R)$ is orthogonal for the norm of $V$ that $\rho$ keeps; the kernels satisfy $\mathbf W_n(\hat\Lambda\mathfrak z;D(\Lambda)e)=\mathbf W_n(\mathfrak z;e)$, at real rotations by the covariance of $\mathsf s_n$ and at the boosts by holomorphic continuation, no bound of $D(\Lambda)$ being read; in the Lorentz covariance of the fields $f\circ\Lambda^{-1}$ becomes $D(\Lambda)f\circ\Lambda^{-1}$ and $\bar f$ becomes $f^*$, and the one new fact is $(\Lambda f)^*=\Lambda(f^*)$ for $(\Lambda f)(x)=D(\Lambda)f(\Lambda^{-1}x)$; locality is proved index by index.

The theorem is used in its form without (Rec5), that is, except the uniqueness of the vacuum. Its hypotheses are (A), with $p_V=p$, $c_0=\frac12$ and the number $C_V$ of (A); the norm $|h(x)|_\infty$ in the bound of (A) is at most the norm $(\operatorname{tr}h(x)^2)^{1/2}$, which $\rho$ keeps. It gives the fields $T(f)$ on the smooth vectors, with the kernels of (B) as Euclidean Green's functions; the boosts are the analytic continuation of the rotations $\mathsf R_t$, which on the plate vectors is the unitary group of Lemma~\ref{cov-boost}\,(b); the spectral condition comes from the cone; locality from the symmetry of the kernels of (B) at all configurations of pairwise different points; and the density of $\mathcal D_0$ from Lemma~\ref{cov-ordered}.

Before invoking rotation covariance, the construction already gives the semigroup on the tube, the point operators, the functions $\mathbf W_n$ and the vectors, the forms and the operators $T(f)$, and the boundary values. These steps use no rotations and apply with the square cone: with $d_\square(\mathfrak z)=\operatorname{Re}\mathfrak z_3-|\operatorname{Im}\mathfrak z_1|-|\operatorname{Im}\mathfrak z_2|$ in the place of $\operatorname{Re}\mathfrak z_3-|\operatorname{Im}\vec{\mathfrak z}|$, the tube $\mathsf{Tb}^\square=\{d_\square>0\}$, and $V^\square_+=\{\alpha^0>|\alpha^1|+|\alpha^2|\}$ in the place of the forward cone; the required spectral estimate is (on $\{E\ge\max(|q_1|,|q_2|)\}$, $|e^{-\mathfrak z_3E+i(\mathfrak z_1q_1+\mathfrak z_2q_2)}|\le e^{-d_\square(\mathfrak z)E}$ and $|q_k|\le E$), and the symmetries used are translations, $\theta$ and the eight elements of $W_3$ that fix $e_3$.

With the inputs (D1)--(D4), (S) and (P) of Propositions~\ref{cov-lim} and~\ref{cov-P} they give fields $T(f)$ on the smooth vectors, tempered, with $f^*=\rho(\theta)\bar f$, covariant under the translations and under the eight elements of $W_3$ that fix $e_3$, with Wightman functions on the tubes of the square cone bounded by $(C\eta^{-p})^n$, $0<\eta\le1$, and with spectrum in $\{E\ge\max(|q_1|,|q_2|)\}$.
\end{proof}

\subsection{Scope of the conclusions}\label{cov-rem}

The theorems are statements on tuples with continuous members and pairwise disjoint supports. Each tensor limit is treated by itself, with its own label. The rates of Theorem~\ref{cov-thm}\,(a) are rates in the cutoff at fixed $\mathsf d$ and $R$, not in the distance of the supports; neither covariance nor conservation is a statement on the decay of correlations, and at large distances the argument needs only the boundedness of the two-point kernels. Theorem~\ref{cov-thm-nz} is not used in the proofs of Theorems~\ref{cov-thm}, \ref{cov-thm-data}, \ref{cov-thm-rider} and of Corollary~\ref{cov-cor-real}. The rotation theorem for the action density, Theorem~\YMMref{AK}{737}, is not used in the proof of Theorem~\ref{cov-thm}\,(c). Of the positivity (D3), Theorem~\ref{cov-thm-data} under (P) and Proposition~\ref{cov-two}\,(b) need only the plane $x_3=0$ and, at tuples of one member, the plane $\Sigma$.

\section{Reflection positivity, spectral cones and reconstruction}\label{C-371-1}

Reflections act both on the arguments and on the tensor labels. We compare the Hilbert spaces reconstructed in the different lattice pictures and obtain the forward spectral cone before using full rotation covariance. This also makes precise which positivity statements survive for growing tori.

The lattice-source admission argument below treats bond bilinears directly. It supplies the tensor input that cannot be obtained merely by applying a scalar-source theorem by name. The remaining results concern the passage to components, their tensor lifts and the labelled reconstruction in real time.

\subsection{Notation and hypotheses}

\emph{Constants and lattice notions.} The constants and the data are those of the main theorems; a number depends on the constants if it is formed from them alone. $L\ge13$ is the odd block factor, $a_K=L^{-K}$, and $a_i=L^{-i}$ for every integer $i$. Recall the reference coupling $g$ and its ceilings $\gamma^{\rm sd}_W\le\gamma^{\rm ref}_W\le\gamma$; the first-passage data $(K,\beta)$ at $g$, whose labels $a_K\beta$ lie in the window $(g^{-2},L(g^{-2}+\bar B)]$, $\bar B$ a constant, so that limits of labels lie in the closed interval $[g^{-2},L(g^{-2}+\bar B)]$, and the interval $I(g)=(g^{-2}+\bar B,L(g^{-2}+\bar B)-2\bar B]$, which is sufficient for first passage: for $\mathsf b\in I(g)$ and every $K$ the pair $(K,L^K\mathsf b)$ is a first-passage datum at $g$ (Proposition~\YMMref{74.1}{615}, Corollary~\YMMref{74.2}{616}); the sets $\mathcal G^{\rm per}_{K,\beta}$ of periodic Gibbs states on $a_K\mathbb Z^3$ (Section~\YMMref{91}{697}); the sides $\ell$ admissible at $g$, the least of them $\ell'$, the tori $T_\ell$ and their Wilson measures $\tilde\mu_{K,\ell,\beta}$; boards of cells of side $a_i$ with offsets in $a_K\mathbb Z^3$, the mirror copies $F^{[\mathbf k']}$ of a function $F$ of the bonds of one cell, and the chessboard seminorms: $\|F\|_{{\rm ch},\ell}^N$ is the mean under $\tilde\mu_{K,\ell,\beta}$ of the product of the $N=(\ell/a_i)^3$ mirror copies of $F$; the pressures $p_\ell=\ell^{-3}\log Z_{K,\ell,\beta}$ (Section~\YMMref{86}{680}) and the pressure gap $|p_\ell-p_{\ell_1}|\le G_g$ for admissible sides, $G_g$ a number depending on the constants and $g$ (Lemma~\ref{lem:canonical-source-system}); and the exponent $\mathfrak p_0\ge1$ of the logarithms in the small-field bounds. $W_3$ is the group of the $48$ signed permutation matrices, $e_1,e_2,e_3$ the standard basis, $|x|$ the Euclidean norm and $|x|_\infty=\max_\mu|x_\mu|$; sup diameters and sup distances are formed with $|\cdot|_\infty$. $\mathsf R_t$ is the rotation about $e_2$ with $\mathsf R_te_3=\cos t\,e_3+\sin t\,e_1$.

\emph{Labels, weights and tuples.} $V=\mathrm{Sym}_3(\mathbb R)$, $V_{\mathbb C}=\mathrm{Sym}_3(\mathbb C)$, with $\rho(R)E=RER^{\rm t}$ for $R\in\mathrm O(3)$ and $|E|_\infty=\max_{\lambda\lambda'}|E_{\lambda\lambda'}|$; $E_{\lambda\lambda'}$ denotes also a matrix unit, and $o_{\mu\nu}=E_{\mu\nu}+E_{\nu\mu}$ for $\mu\ne\nu$. As $F_{\mu\nu}=\varepsilon_{\mu\nu\lambda}B_\lambda$, the array $\operatorname{tr}F_{ij}F_{kl}$ is the tensor $T_{\lambda\lambda'}=\langle B_\lambda,B_{\lambda'}\rangle$, $\langle X,Y\rangle=\frac12\operatorname{tr}XY$, whose components are indexed by $V$. The decomposition $V=V_A\oplus V_E\oplus V_T$ into the multiples of $\mathbf 1$, the traceless diagonal matrices and the matrices with zero diagonal is invariant under $\rho(W_3)$, with projections $P_A,P_E,P_T$, and $Z_z=P_A+P_E+zP_T$, $z$ real, multiplies the off-diagonal labels by $z$.

A weight is an $H\in C^1_c(\mathbb R^3;V)$; where continuous weights $H\in C_c(\mathbb R^3;V)$ are admitted we say so. $|H|_\infty=\sup_x|H(x)|_\infty$, $\|H\|_{L^1}=\int|H(x)|_\infty\,dx$, $H^{(z)}(x)=H(L^{-z}x)$ for $z\in\mathbb N$. An isometry $\iota x=Rx+v$ acts by $(\iota H)(x)=\rho(R)H(\iota^{-1}x)$, and on vector fields by $(\iota\eta)(x)=R\eta(\iota^{-1}x)$; $\tau_v$ is the translation by $v$. For open $O\subset\mathbb R^3$, $\mathfrak S_V(O)$ is the set of the unordered tuples $X=(H_1,\dots,H_n)$, $n\ge0$, of weights with pairwise disjoint supports in $O$; $X\cdot Y$ is the concatenation of two tuples with disjoint supports, $\mathbb C\mathfrak S_V(O)$ the space of the complex combinations $F=\sum c_iX_i$, and $\bar F=\sum\bar c_iX_i$.

Isometries and linear maps of $V$, applied pointwise, act on tuples entry by entry and on combinations linearly. An entry with values in $V_{\mathbb C}$ stands for the expansion of $H=\operatorname{Re}H+i\operatorname{Im}H$ in that entry. For a vector field $\eta\in C^1(\mathbb R^3;\mathbb R^3)$, with $(\nabla\eta)_{ij}=\partial_i\eta_j$,
\[H^{\rm cons}[\eta]=(\nabla\cdot\eta)\mathbf 1-\nabla\eta-(\nabla\eta)^{\rm t},\qquad\iota H^{\rm cons}[\eta]=H^{\rm cons}[\iota\eta];\]
formally $\frac12\int\operatorname{tr}(H^{\rm cons}[\eta]T)=-\int\eta_\nu\partial_\mu\Theta_{\mu\nu}$ with the stress tensor $\Theta=\frac12(\operatorname{tr}T)\mathbf 1-T$, $\Theta_{\mu\nu}=F_{\mu\alpha}F_{\nu\alpha}-\frac14\delta_{\mu\nu}F_{\alpha\beta}F_{\alpha\beta}$.

\emph{Lattice observables and limits.} A configuration on $a_K\mathbb Z^3$, or on a torus $T_\ell$, assigns to every oriented bond $b$ a matrix $U(b)\in\mathrm{SU}(2)$, with $U(\bar b)=U(b)^{-1}$ at the reversed bond. For a plaquette $p$, its plaquette variable $\partial U(p)$ is the product of $U$ around $p$, defined up to a conjugation that the functions below do not see; $x(p)$ is its centre, $\lambda(p)$ its normal direction, and $a_p=1-\frac12\operatorname{tr}\partial U(p)\in[0,2]$ its plaquette square. With $e_{-\mu}=-e_\mu$, for $\alpha,\alpha'\in\{\pm1,\pm2,\pm3\}$ with $|\alpha|\ne|\alpha'|$ let $W_x(\alpha,\alpha')$ be the product of $U$ along the loop $x\to x+a_Ke_\alpha\to x+a_Ke_\alpha+a_Ke_{\alpha'}\to x+a_Ke_{\alpha'}\to x$, and $\mathbf w(V)=(V-V^{-1})/(2i)$. For a positively oriented bond $b=\langle x,x+a_Ke_\mu\rangle$ with midpoint $x_b$, and $\lambda_b<\lambda'_b$ the two directions other than $\mu$, put
\[u_{b,\alpha}=-\tfrac12\bigl[\mathbf w(W_x(\mu,\alpha))-\mathbf w(W_x(\mu,-\alpha))\bigr]\quad(\alpha\in\{\lambda_b,\lambda'_b\}),\qquad q_b=\langle u_{b,\lambda_b},u_{b,\lambda'_b}\rangle:\]
$u_{b,\alpha}$ is formed from the two plaquettes of the plane of $e_\mu,e_\alpha$ that contain $b$, read from $x$, and $q_b$ is the bond form of $b$. A weight $w$ of the plaquettes, or $\mathsf w$ of the bonds, is a bounded complex function on the plaquettes, or on the positively oriented bonds, and
\[O(w)=a_K^{-1}\sum_pw(p)a_p,\qquad Y[\mathsf w]=-a_K^{-1}\sum_b\mathsf w(b)q_b .\]
A bounded weight is a bounded Borel function $H:\mathbb R^3\to V$, or $V_{\mathbb C}$, of bounded support; weights are bounded weights. For a bounded weight $H$
\[O^T_K(H)=O(w_H)+Y[\mathsf w_H],\qquad w_H(p)=H_{\lambda(p)\lambda(p)}(x(p)),\quad\mathsf w_H(b)=H_{\lambda_b\lambda'_b}(x_b),\]
and the same formulas define $O^T_K(H)$ on a torus, for $H$ on $T_\ell$ or for $H$ on $\mathbb R^3$ with support in a period box. So $O^T_K(H)$ weights the square of each plaquette with $H_{\lambda\lambda}$ at its centre, $\lambda$ normal to the plaquette, and the bond form of the directions $\lambda\ne\lambda'$ with $H_{\lambda\lambda'}$ at the midpoint of the bond, and $O^T_K(f\mathbf 1)=O(f)$ is the action density weighted with $f$. The plaquette squares and the bond forms enter with one normalisation: for the configuration $U(b)=e^{ia_KA_\mu(x_b)}$ of a smooth field $A$ with curvature $F$, $F_{\mu\nu}=\varepsilon_{\mu\nu\lambda}B_\lambda$, one has $a_p=\frac12a_K^4\langle F_{\mu\nu},F_{\mu\nu}\rangle+o(a_K^4)$ at a plaquette of the plane of $e_\mu,e_\nu$, and $q_b=a_K^4\langle F_{\mu\lambda_b},F_{\mu\lambda'_b}\rangle+o(a_K^4)=-a_K^4\langle B_{\lambda_b},B_{\lambda'_b}\rangle+o(a_K^4)$; so formally $O^T_K(H)=\frac12\int\operatorname{tr}(H(x)T(x))\,dx$, the minus sign in $Y$ compensating that of $\langle F_{\mu\lambda},F_{\mu\lambda'}\rangle=-\langle B_\lambda,B_{\lambda'}\rangle$. A Gibbs sequence at $g$ is a sequence $(K_j,\beta_j,\nu_j)$ of first-passage data at $g$ with $K_j\to\infty$ and of states $\nu_j\in\mathcal G^{\rm per}_{K_j,\beta_j}$. A tensor limit at $g$ is a functional on $\mathfrak S_V(\mathbb R^3)$ given, for some Gibbs sequence at $g$, by
\begin{equation}\label{C-371-2}
\sigma(X)=\lim_j\nu_j\Bigl(\prod_{a=1}^n\bigl(O^T_{K_j}(H_a)-\nu_j(O^T_{K_j}(H_a))\bigr)\Bigr),\qquad X=(H_1,\dots,H_n),
\end{equation}
the limit existing for every $X$; a limit at continuous weights is the same at all tuples of continuous weights. The label of the sequence is $b=\lim_ja_{K_j}\beta_j$ when this limit exists; the trace part of $\sigma$ is $\sigma_0(f_1,\dots,f_n)=\sigma(f_1\mathbf 1,\dots,f_n\mathbf 1)$. Let $\xi$ be a component of the canonical decomposition of the scalar limit of the sequence and $\rho_j\in\mathcal G^{\rm per}_{K_j,\beta_j}$ its approximants, the states $\nu_j$ conditioned on translation-invariant events, as in the theorem on the scalar limit among the construction results (Theorems~\YMMref{AL}{746} and~\YMMref{AP}{780}). A lift $\sigma^\sharp_\xi$ is a limit \eqref{C-371-2} with the outer state $\nu_j$ replaced by $\rho_j$, along a subsequence; its re-centred functional $\sigma^c_\xi$ is the limit with $\rho_j$ in both places.

\emph{Mirrors and pictures.} A mirror is a pair $(n,c)\in S^2\times\mathbb R$, with the reflection $\theta_{n,c}x=x-2(n\cdot x-c)n$, $\Theta_n=1-2nn^{\rm t}$, $\theta_n=\theta_{n,0}$; it acts on weights by $(\theta_{n,c}H)(x)=\Theta_nH(\theta_{n,c}x)\Theta_n$. We put $\mathcal D_{n,c}=\mathbb C\mathfrak S_V(\{n\cdot x>c\})$ and $\mathcal D_n=\mathcal D_{n,0}$. A functional $\sigma$ is positive at $(n,c)$ if $\sigma(\theta_{n,c}\bar F\cdot F)\ge0$ for all $F\in\mathcal D_{n,c}$. $\mathfrak N_{18}$ is the set of the vectors $\pm e_\mu$, $(\pm e_\mu\pm e_\lambda)/\sqrt2$, $\mu\ne\lambda$, the normals of the nine families of lattice planes. For an orthonormal frame $(u,u',u'')$ the cubic group of the frame consists of the $48$ orthogonal maps that permute $u,u',u''$ up to signs; its $18$ normals are $\pm u,\pm u',\pm u''$ and $(\pm v\pm v')/\sqrt2$, $v\ne v'$ in the frame.

If $\sigma$ satisfies (D0)--(D2) of Definition~\ref{C-371-3} and is positive at $(u,0)$, the picture $u$ (Lemma~\ref{C-371-7}) consists of the Hilbert space $\mathcal H_u$ of the form $\sigma(\theta_u\bar F'\cdot F)$ on $\mathcal D_u$ with the classes $[F]_u$, the vacuum $\Omega_u=[\emptyset]_u$, the energy $H_u\ge0$ with $e^{-tH_u}[F]_u=[\tau_{tu}F]_u$, the momenta $P_u$ generating $[F]_u\mapsto[\tau_vF]_u$, $v\perp u$, and the joint spectral measure $\mathcal P_u$ of $(H_u,P_u)$ on $[0,\infty)\times u^\perp$; for $G,G'\in\mathcal D_u$, $m^u_{G',G}=\langle[G']_u,\mathcal P_u(\cdot)[G]_u\rangle$. At $u=e_3$ we write $\mathcal H(\sigma)$, $[F]$, $\Omega$, $H$, $P=(P_1,P_2)$, $\mathcal P$, $\theta=\theta_{e_3}$, and $U(t,\vec x)=e^{itH}e^{i(x_1P_1+x_2P_2)}$, $\vec x=(x_1,x_2)$.

The forward cone is $\{(E,q):E\ge|q|\}$, $|q|$ Euclidean, and the square cone $\{E\ge\max(|q_1|,|q_2|)\}$; the circular tube is $\mathsf{Tb}=\{\zeta\in\mathbb C^3:\operatorname{Re}\zeta_3>|(\operatorname{Im}\zeta_1,\operatorname{Im}\zeta_2)|\}$, the square tube $\mathsf{Tb}^\square=\{\zeta:\operatorname{Re}\zeta_3>|\operatorname{Im}\zeta_1|+|\operatorname{Im}\zeta_2|\}$, and $\mathsf E(\zeta)=\int e^{-\zeta_3E+i(\zeta_1q_1+\zeta_2q_2)}\,d\mathcal P(E,q)$. $\Lambda^{(1)}_\varsigma(E,q_1,q_2)=(E\cosh\varsigma+q_1\sinh\varsigma,\,q_1\cosh\varsigma+E\sinh\varsigma,\,q_2)$ is the boost in the plane $(E,q_1)$, and $\Lambda^{(2)}_\varsigma$ the same in $(E,q_2)$.

\begin{definition}\label{C-371-3}
For $\sigma:\mathfrak S_V(\mathbb R^3)\to\mathbb R$ we consider the conditions:
\begin{itemize}
\item[(D0)] $\sigma(\emptyset)=1$; $\sigma$ is symmetric and $n$-linear on $\prod_aC^1_c(O_a;V)$ for open sets $O_a$ with pairwise disjoint closures, and it is extended linearly to combinations;
\item[(D1)] for $n\ge1$ and $\delta,\varrho_0>0$ there is $C_n(\delta,\varrho_0)$ with $|\sigma(H_1,\dots,H_n)|\le C_n(\delta,\varrho_0)\prod_a|H_a|_\infty$ whenever the supports have sup diameters at most $\varrho_0$ and pairwise sup distances at least $2\delta$; the bound does not depend on the position of the supports, and for complex entries it holds with the factor $2^n$;
\item[(D2)] $\sigma\circ\tau_v=\sigma$ for $v\in\mathbb R^3$; \quad(D3) $\sigma\circ w=\sigma$ for $w\in W_3$;
\item[(D4)] $\sigma$ is positive at $(n,c)$ for all $n\in\mathfrak N_{18}$ and $c\in\mathbb R$.
\end{itemize}
We put $G_\sigma=\{R\in\mathrm O(3):\sigma(RX)=\sigma(X)\text{ for all }X\}$; $\sigma$ is covariant under $\mathrm O(3)$ if $G_\sigma=\mathrm O(3)$. A set $\mathfrak K$ of tuples is full if $\iota X\in\mathfrak K$ for every $X\in\mathfrak K$ and every isometry $\iota$, and then $G_\sigma[\mathfrak K]=\{R\in\mathrm O(3):\sigma(RX)=\sigma(X)\text{ for all }X\in\mathfrak K\}$. $\mathfrak K_N$ is the set of the tuples with at most $N$ entries and $\mathfrak T_k$ that of the tuples with at most $k$ entries not of the form $f\mathbf 1$; both are full, as $\rho(R)\mathbf 1=\mathbf 1$. $\sigma$ is covariant at order $N$ if $G_\sigma[\mathfrak K_N]=\mathrm O(3)$. For $m\in\mathbb N\cup\{\infty\}$, $\mathfrak M^{(m)}(\sigma)$ is the set of the $u\in S^2$ with $\sigma(\theta_u\bar F\cdot F)\ge0$ for every $F\in\mathcal D_u$ that combines tuples with at most $m$ entries, and $\mathfrak M(\sigma)=\mathfrak M^{(\infty)}(\sigma)$. Positivity is meant with complex coefficients throughout.
\end{definition}

\subsection{Principal results}

\begin{theorem}[Mirrors and the forward cone for labelled functionals]\label{C-371-4}
Let $\sigma:\mathfrak S_V(\mathbb R^3)\to\mathbb R$.
\begin{enumerate}
\item[(a)] \emph{The mirror on the labels.} For every mirror $(n,c)$ and every $E\in V$,
\[(\Theta_nE\Theta_n)_{\lambda\lambda'}=E_{\lambda\lambda'}-2n_\lambda(En)_{\lambda'}-2(En)_\lambda n_{\lambda'}+4(n\cdot En)\,n_\lambda n_{\lambda'};\]
$\rho(\Theta_n)$ is $+1$ on $\mathbb Rnn^{\rm t}\oplus\{E:En=0\}$, of dimension $1+3$, and $-1$ on $V_n^-=\{nm^{\rm t}+mn^{\rm t}:m\perp n\}$, of dimension $2$; at $n=e_\mu$ a component with $k$ indices equal to $\mu$ acquires the factor $(-1)^k$; on the four indices of $\operatorname{tr}F_{ij}F_{kl}$ the mirror acts as $\Theta_n^{\otimes4}$. $\Theta_n\in W_3$ if and only if $n\in\mathfrak N_{18}$, and for $n\notin\mathfrak N_{18}$ the map $\rho(\Theta_n)$ mixes $V_A\oplus V_E$ with $V_T$; at $n=(1,1,1)/\sqrt3$, for instance, $(\Theta_nE\Theta_n)_{\lambda\lambda'}=E_{\lambda\lambda'}-\frac23(r_\lambda+r_{\lambda'})+\frac49s$, with the row sums $r_\lambda$ and the sum $s$ of the entries of $E$. Under (D0), positivity at $(n,c)$ and at $(-n,-c)$ are equivalent; under (D0), (D2) both are equivalent to positivity at $(n,0)$. The action on the labels is forced: if $\sigma$ satisfies (D0)--(D2) and is positive at $(u,0)$, and if $A:V\to V$ commutes with $\rho(\Theta_u)$ and $\sigma((A\theta_u\bar F)\cdot F)\ge0$ for all $F\in\mathcal D_u$, then $[AF]_u=[F]_u$ for all $F$; in particular, for the reflection applied to the arguments alone, $A=\rho(\Theta_u)$, one gets $[(h)]_u=0$ for every weight $h$ with values in $V_u^-$.
\item[(b)] \emph{A positive mirror is a symmetry; nine families or all planes.} If $\sigma$ satisfies (D0), (D2) and $u\in\mathfrak M(\sigma)$, then $\Theta_u\in G_\sigma$; it suffices that the numbers $\sigma(\theta_u\bar F\cdot F)$, $F\in\mathcal D_u$, be real. If $\sigma$ satisfies (D0), (D2), $m\in\mathfrak M(\sigma)$, $R\in G_\sigma$ and $Rm=n$, then $\sigma$ is positive at $(n,c)$ for every $c$. If $\sigma$ satisfies (D0)--(D4), then $G_\sigma$ is a closed subgroup of $\mathrm O(3)$ containing $W_3$, and the following are equivalent: (i) $\sigma$ is covariant under $\mathrm O(3)$; (ii) $G_\sigma$ contains an element outside $W_3$; (iii) $\sigma$ is positive at one plane whose normal is not in $\mathfrak N_{18}$; (iv) $\sigma$ is positive at every plane. In (ii) one reflection $\Theta_n$, $n\notin\mathfrak N_{18}$, may be taken, and (iii) may be weakened to: the numbers $\sigma(\theta_n\bar F\cdot F)$, $F\in\mathcal D_n$, are real for one $n\notin\mathfrak N_{18}$. Hence $\mathfrak M(\sigma)=\mathfrak N_{18}$ or $\mathfrak M(\sigma)=S^2$.
\item[(c)] \emph{How much covariance one plane uses.} Let $\sigma$ satisfy (D0)--(D4). For every full set $\mathfrak K$, $G_\sigma[\mathfrak K]$ is $W_3$ or $\mathrm O(3)$. If $\mathfrak F$ and $\mathfrak K$ are full, $X\cdot Y\in\mathfrak K$ for all $X,Y\in\mathfrak F$ with disjoint supports, and $G_\sigma[\mathfrak K]=\mathrm O(3)$, then $\sigma(\theta_{n,c}\bar F\cdot F)\ge0$ at every mirror for every combination $F$ of tuples of $\mathfrak F$ on the side of the mirror; this applies to $(\mathfrak F,\mathfrak K)=(\mathfrak K_m,\mathfrak K_{2m})$, $(\mathfrak T_k,\mathfrak T_{2k})$ and $(\mathfrak T_0\cup\mathfrak K_1,\mathfrak T_1\cup\mathfrak K_2)$. For $m\in\mathbb N\cup\{\infty\}$: covariance at order $2m$ implies $\mathfrak M^{(m)}(\sigma)=S^2$, which implies $\mathfrak M^{(m)}(\sigma)\ne\mathfrak N_{18}$, which implies covariance at order $m$; for $m=\infty$ the four statements are equivalent.

Covariance at order $1$ holds always: $\sigma((H))=c_\sigma\int\operatorname{tr}H$ with a number $c_\sigma$. If $\sigma(H,H')=\iint\langle H(x)\otimes H'(y),s(y-x)\rangle\,dx\,dy$ with a continuous $s:\mathbb R^3\setminus\{0\}\to V\otimes V$, then covariance at order $2$, $\mathfrak M^{(1)}(\sigma)=S^2$ and $\mathfrak M^{(1)}(\sigma)\ne\mathfrak N_{18}$ are equivalent; if $s$ is real analytic, covariance at order $2$ follows from $\sigma(RH,RH')=\sigma(H,H')$ for one $R\notin W_3$ and all $H,H'$ supported in two fixed balls with disjoint closures.
\item[(d)] \emph{The forward cone from three mirrors.} Let $\sigma$ satisfy (D0)--(D2), let $(u,u',u'')$ be an orthonormal frame, and let $\sigma$ be positive at $(u,0)$, $((u+u')/\sqrt2,0)$ and $((u+u'')/\sqrt2,0)$. Then $\sigma$ is invariant under the cubic group of the frame and positive at $(n,c)$ for its $18$ normals $n$ and all $c$, and
\[\mathcal P_n(\{(E,q):E<|q|\})=0\qquad\text{for each of the $18$ normals $n$.}\]
In the picture $u$, $\mathsf E(\zeta)$ is a contraction for $\zeta\in\overline{\mathsf{Tb}}$, holomorphic in norm on $\mathsf{Tb}$, with $\mathsf E(\zeta)\mathsf E(\zeta')=\mathsf E(\zeta+\zeta')$ and $\mathsf E((v,t))=e^{-tH_u}e^{iv\cdot P_u}$ at real points. In the coordinates $e_3=u$, $e_1=-u'$, $e_2=u''$ there is a positive self-adjoint injective operator $D$ on $\mathcal H_u$ with $D\Omega_u=\Omega_u$, $D^{is}(H_u\pm P_1)D^{-is}=e^{\pm\pi s/2}(H_u\pm P_1)$ and $D^{is}\mathcal P_u(B)D^{-is}=\mathcal P_u(\Lambda^{(1)}_{-\pi s/2}B)$; for every $\Lambda\in\mathrm{SO}_0(1,2)$ there is a unitary $V_\Lambda$ with $V_\Lambda\Omega_u=\Omega_u$ and $V_\Lambda\mathcal P_u(B)V_\Lambda^{-1}=\mathcal P_u(\Lambda B)$.

So the null sets and the support of $\mathcal P_u$ are invariant under $\mathrm{SO}_0(1,2)$, and the support is the union of $\{0\}$ with some of the orbits $\{E=|q|>0\}$ and $\{E^2-|q|^2=M^2,\ E>0\}$, $M>0$. Between the pictures of any two of the $18$ normals there is a unitary map that carries vacuum, energy, momenta and spectral measure into one another; so the set of the masses $M$ that occur is the same in the $18$ pictures. For the vectors of the wedge $\{x\cdot u>|x\cdot u'|\}$ the pictures $u$ and $(u+u')/\sqrt2$ give one law to the pair (mass squared in the plane of the two mirrors, momentum along $u''$), in the sense of Lemma~\ref{C-371-19}. All this holds with $V$ replaced by any finite-dimensional real representation of $\mathrm O(3)$.

The two diagonal mirrors may be replaced by square cones: if $\sigma$ satisfies (D0)--(D2), $u,u'\in\mathfrak M(\sigma)$ and $\mathcal P_u(\{E<|q\cdot u'|\})=\mathcal P_{u'}(\{E<|q\cdot u|\})=\mathcal P_u(\{E<|q\cdot u''|\})=0$, then the forward cone holds in the pictures $\pm u$, $\pm u'$.
\item[(e)] \emph{From rotations.} If $\sigma$ satisfies (D0)--(D2), $e_3\in\mathfrak M(\sigma)$, $H\ge|P_1|$, and $G_\sigma$ contains the rotations about $e_3$ by a dense set of angles, then $\mathcal P(\{E<|q|\})=0$. If $\sigma$ satisfies (D0)--(D2), $e_3\in\mathfrak M(\sigma)$ and $G_\sigma=\mathrm O(3)$, then $\mathfrak M(\sigma)=S^2$; for $Re_3=u$ the map $[F]\mapsto[RF]_u$ is unitary from $\mathcal H(\sigma)$ onto $\mathcal H_u$ and carries $(H,P,\mathcal P)$ to $(H_u,P_u,\mathcal P_u)$, the momenta turned by $R$; $\mathcal P_u(\{E<|q|\})=0$ for every $u$; and the maps $[F]\mapsto[RF]$, $Re_3=e_3$, form a strongly continuous unitary representation of $\mathrm O(2)$ on $\mathcal H(\sigma)$.
\item[(f)] \emph{Examples.} (i) Let $\sigma^{\rm fr}(H_1,\dots,H_n)$ be the correlation of the Wick products 
\[
\frac12\int\sum_{\lambda\lambda'}H_a(x)_{\lambda\lambda'}{:}\langle B_\lambda,B_{\lambda'}\rangle(x){:}\,dx
\]
 for three independent Gaussian fields $B^{\mathsf a}=(B^{\mathsf a}_\alpha)_{\alpha\le3}$, $\mathsf a\in\{1,2,3\}$, with covariance $\delta^{\mathsf a\mathsf b}K_3(x-y)_{\alpha\gamma}$ at $x\ne y$, 
\[
K_3(u)=\partial\partial(4\pi|u|)^{-1}=(4\pi|u|^3)^{-1}(3\hat u\hat u^{\rm t}-\mathbf 1)
\]
 for $u\ne0$, $\hat u=u/|u|$, taken at tuples with disjoint supports and without self-contraction, so that $K_3$ is read off $0$ alone. Then $\sigma^{\rm fr}$ satisfies (D0)--(D4) and is covariant under $\mathrm O(3)$, while $\sigma^{\rm fr}\circ Z_2$ satisfies (D0)--(D4), has the forward cone in the nine pictures, and $\mathfrak M(\sigma^{\rm fr}\circ Z_2)=\mathfrak N_{18}$. (ii) The diagonal mirrors cannot be omitted in (d): for $m,\lambda>0$ the Gaussian functional $\sigma_D$ of the traces $\operatorname{tr}H_a$ with the pair covariance 
\[
\hat S_D(p)=[m^2+|p|^2+\lambda(p_1^2p_2^2+p_1^2p_3^2+p_2^2p_3^2)]^{-1}
\]
 satisfies (D0)--(D3), is positive at every coordinate mirror, and $\mathcal P(\{E<c|q_1|\})\ne0$ for every $c>0$; it is positive at no diagonal mirror. (iii) The third mirror cannot be omitted in (d): for $0<c<1$ and $m>0$ the Gaussian functional $\sigma_c$ of the traces with $\hat S_c(p)=[p_1^2+c^2p_2^2+p_3^2+m^2]^{-1}$ satisfies (D0)--(D2) and is positive at $(u,0)$ for every unit vector $u\perp e_2$, so at $(e_3,0)$ and $((e_3+e_1)/\sqrt2,0)$; on the vectors of one entry its spectral measure is carried by $\{E^2=q_1^2+c^2q_2^2+m^2\}$, and $\mathcal P(\{E<|q|\})\ne0$.
\end{enumerate}
\end{theorem}

\begin{theorem}[The tensor limits]\label{C-371-5}
Let $0<g\le\gamma^{\rm ref}_W$.
\begin{enumerate}
\item[(1)] \emph{The limits are data.} Every Gibbs sequence at $g$ has a subsequence with a tensor limit. If a Gibbs sequence has a tensor limit $\sigma$, its label $b$ exists, $g^{-2}\le b\le L(g^{-2}+\bar B)$, and $\sigma$ determines it: $b^{-2}=\lim_{d\downarrow0}\sigma(\varphi^d_1E,\varphi^d_2E)/\mathcal W(\varphi^d_1E,\varphi^d_2E)$ for $E=o_{12}$ and the bumps of Proposition~\ref{C-371-24}. The trace part of $\sigma$ is the scalar limit of the sequence. Every tensor limit at $g$ and every lift satisfies (D0)--(D4), the number $C_n(\delta,\varrho_0)$ of (D1) being the number $C^{\rm lc}_n$ of Proposition~\ref{C-371-22}, which depends on the constants, $g$, $n$, $\delta$ and $\varrho_0$ alone; tensor limits have $\sigma((H))=0$, lifts $\sigma^\sharp_\xi((H))=\frac{c_\xi}3\int\operatorname{tr}H$ with a real number $c_\xi$. The re-centred functional $\sigma^c_\xi$ is a tensor limit at $g$, that of $(K_j,\beta_j,\rho_j)$, and
\[\sigma^\sharp_\xi(X)=\sum_{S\subset\{1,\dots,n\}}\ \prod_{a\notin S}\Bigl(\frac{c_\xi}3\int\operatorname{tr}H_a\Bigr)\,\sigma^c_\xi(X_S),\qquad\sigma^c_\xi(X)=\sum_{S\subset\{1,\dots,n\}}\ \prod_{a\notin S}\Bigl(-\frac{c_\xi}3\int\operatorname{tr}H_a\Bigr)\,\sigma^\sharp_\xi(X_S),\]
$X_S=(H_a)_{a\in S}$.
\item[(2)] \emph{The forward cone.} For every tensor limit and every lift, $\mathcal P_n(\{E<|q|\})=0$ in the picture of each $n\in\mathfrak N_{18}$, together with the contractions $\mathsf E(\zeta)$ on $\overline{\mathsf{Tb}}$, the boosts $D^{is}$, the unitaries $V_\Lambda$, the invariance of the support under $\mathrm{SO}_0(1,2)$ and the unitary maps between the pictures of Theorem~\ref{C-371-4}(d).
\item[(3)] \emph{Covariance.} For every tensor limit at every $0<g\le\gamma$, and for every lift, $\sigma(RX)=\sigma(X)$ for all $R\in\mathrm O(3)$ and all $X\in\mathfrak S_V(\mathbb R^3)$. At two entries the covariance follows from conservation, the diagonal mirror, the invariance of the scalar pair and the scale-free bound on the off-diagonal fields alone (Proposition~\ref{C-371-27}), with no bound by a power of the distance.
\item[(4)] \emph{Positivity at every plane.} Every tensor limit at every $0<g\le\gamma$, and every lift, is positive at every mirror $(n,c)\in S^2\times\mathbb R$:
\[\sigma(\theta_{n,c}\bar F\cdot F)\ge0\quad(F\in\mathcal D_{n,c}),\qquad(\theta_{n,c}H)(x)=\Theta_nH(x-2(n\cdot x-c)n)\Theta_n.\]
So $\mathfrak M(\sigma)=S^2$ at every $0<g\le\gamma$. For $0<g\le\gamma^{\rm ref}_W$, for every tensor limit and every lift, every $u\in S^2$ has its picture, unitarily equivalent to the picture $e_3$ and with the forward cone, and the rotations about $e_3$ act by a strongly continuous unitary representation of $\mathrm O(2)$ on $\mathcal H(\sigma)$. Turning the labels is necessary: for a tensor limit at $0<g\le\gamma^{\rm ref}_W$, every $E\in V$ and the bumps $\varphi\ge0$ of Proposition~\ref{C-371-24}(b), $\|[(\varphi E)]\|^2=\sigma(\theta(\varphi E),\varphi E)\ge(w_0-\epsilon_g(d))\,b^{-2}\operatorname{tr}E^2$, by the two-point law of Proposition~\ref{C-371-24}, which is carried from $g\le\gamma^{\rm sd}_W$ to $\gamma^{\rm sd}_W<g\le\gamma^{\rm ref}_W$ by the dilation, with $d_0(g)=L^{-z}\min_sd_0(g'_s)$; so no nonzero real combination of the six fields vanishes in $\mathcal H(\sigma)$; at $0\ne E\in V^-_{e_3}$ the form $\sigma((\varphi\circ\theta)E,\varphi E)$ of the reflection acting on the points alone is negative, and by the covariance the same holds at every plane.
\item[(5)] \emph{The hypotheses of the reconstruction.} Every tensor limit $\sigma$ at $g$ satisfies (R1)--(R4) of Proposition~\ref{C-371-16} for every $p>17/2$, with $C_V$ depending on the constants, $g$ and $p$; for $g\le\gamma^{\rm sd}_W$ one may take $C_V=2^pC_pb^{-1}\le2^pC_pg^2$, $C_p$ the number of Proposition~\ref{C-371-15}, and for $\gamma^{\rm sd}_W<g\le\gamma^{\rm ref}_W$ the same with one further factor $L^{4z}$, $z$ the dilation of Lemma~\ref{C-371-14}. The kernels $\mathbf W_n$ of the product formula are holomorphic on $\mathsf{Tb}^{n-1}$ (for every lift too) and
\[|\mathbf W_n(\zeta;e)|\le(C_V\eta^{-p})^n\prod_k|e_k|_\infty\quad\text{if }\operatorname{Re}\zeta^k_3-|(\operatorname{Im}\zeta^k_1,\operatorname{Im}\zeta^k_2)|\ge2\eta\text{ for all }k,\ 0<\eta\le1.\]
\item[(6)] \emph{Wightman fields.} Every tensor limit $\sigma$ at $g$ has the conclusions (i)--(iv) of Proposition~\ref{C-371-16} with $p=9$: for each $n\ge2$ one continuous kernel $s_n$ on the configurations of $n$ different points, real analytic, symmetric and covariant under the isometries acting on points and labels, which integrates to $\sigma$ at every tuple of $C^1$ or of continuous weights, with $|s_n(x;e)|\le(cC_V)^n\max\{1,2n^2/\delta(x)\}^{9n}\prod_k|e_k|_\infty$, $|\sigma(X)|\le(cC_V)^n(2n^2/d)^{9n}\prod_a\|H_a\|_{L^1}$ for tuples at mutual distances $\ge d$, $0<d\le1$, and $|S_n(\varphi)|\le C^n(n!)^{18}|\varphi|_{9n}$ on test functions flat at the coincidences; the truncated functions have the kernels that the formula of the cumulants writes in the $s_m$; six tempered fields $T(f)$ with $T(f)^*\supset T(\rho(\Theta_{e_3})\bar f)$, so that $T_{33}$ and $T_{jk}$, $j,k\le2$, are Hermitian and $T_{3j}$ is Hermitian after the factor $i$; a strongly continuous unitary representation of the Poincar\'e group of $\mathbb R^{1+2}$ under which the fields are covariant, the vacuum is invariant and the spectrum lies in the closed forward cone; locality; cyclicity of the vacuum; the Wightman functions are tempered boundary values of the $\mathbf W_n$ from every direction of the forward cone.

Every lift $\sigma^\sharp_\xi$ has such kernels, with $2\max\{1,|c_\xi|\}\max\{1,cC_V\}$ in the place of $cC_V$ and $S_1(\varphi)=\frac{c_\xi}3\int\operatorname{tr}\varphi$.
\item[(7)] \emph{Components.} Every tensor limit $\sigma$ at $g$ has a decomposition $(\mathsf X,\lambda,(\sigma_\chi)_{\chi\in\mathsf X})$ with labels, the canonical decomposition among the construction results (Theorem~\YMMref{AL}{746}) applied to labelled tuples: $\mathsf X$ compact metrisable, $\lambda$ a Borel probability measure of support $\mathsf X$, $\chi\mapsto\sigma_\chi$ injective and continuous at each tuple, and $\int\prod_i\sigma_\chi(X^i)\,\lambda(d\chi)$ the iterated plateau of $\sigma$ at $(X^1,\dots,X^q)$ for every $q\ge1$ and all tuples $X^1,\dots,X^q$. The iterated plateau is that of Definition~\ref{C-355-4}(5), read at labelled tuples: $G^{(1)}(v_2,\dots,v_q)=\sigma(X^1\cdot\tau_{v_2}X^2\cdots\tau_{v_q}X^q)$ where the supports of $X^1,\tau_{v_2}X^2,\dots,\tau_{v_q}X^q$ are pairwise disjoint, $G^{(p)}(v_{p+1},\dots,v_q)=\lim_{|v_p|\to\infty}G^{(p-1)}(v_p,v_{p+1},\dots,v_q)$ for $p=2,\dots,q$, the limits taken one after the other, and the iterated plateau is the number $G^{(q)}$; for $q=2$ it is the plateau $\lim_{|v|\to\infty}\sigma(X^1\cdot\tau_vX^2)$ of $\sigma$ at $(X^1,X^2)$.

That all these limits exist is part of the assertion, and along the sequence of $\sigma$ one has $G^{(q)}=\lim_j\nu_j\bigl(\prod_iE_{\nu_j}[A_j(X^i)\mid\mathcal I]\bigr)$, $A_j(X)$ the product of the entries of $X$ centred in $\nu_j$ and $\mathcal I$ the translation-invariant events. Let $(\mathsf X_0,\lambda_0,(\sigma_{0,\zeta}))$ be the canonical decomposition of the trace part. Exactly one map $\pi:\mathsf X\to\mathsf X_0$ has $\pi(\chi)$ the trace part of $\sigma_\chi$; it is continuous and onto, and $\pi_*\lambda=\lambda_0$. Each $\sigma_\chi$ is the lift of $\pi(\chi)$ formed with its own approximants along the full sequence, with $\sigma_\chi((H))=\frac{c_{\pi(\chi)}}3\int\operatorname{tr}H$, and it clusters: $\sigma_\chi(X\cdot\tau_vY)\to\sigma_\chi(X)\sigma_\chi(Y)$ as $|v|\to\infty$. $\ker H=\mathbb C\Omega$ on $\mathcal H(\sigma_\chi)$ and on $\mathcal H(\sigma^c_\chi)$; so the six fields of $\sigma^c_\chi$ satisfy all the Wightman axioms, uniqueness of the vacuum included.
\item[(8)] \emph{A functional with exactly nine families.} Let $0<g\le\gamma^{\rm sd}_W$ and $\sigma$ be a tensor limit at $g$. Then $\sigma\circ Z_2$ satisfies (D0)--(D4), with $2^nC_n(\delta,\varrho_0)$ in (D1), has the trace part of $\sigma$, the operators, the slab and point operators, the product formula and the fields of $\sigma$ up to a unitary map, with $Z_2$ applied to the labels, and the forward cone; and $\mathfrak M(\sigma\circ Z_2)=\mathfrak N_{18}$. Among the functionals $\sigma\circ Z_z$, $z\ne0$, the two-point law of Proposition~\ref{C-371-24}, at all pairs of weights and with some label $b'>0$ in the place of $b$, holds for $z=1$ alone; at the pairs of off-diagonal weights alone, with the label $b$ of $\sigma$, it holds if and only if $z^2=1$.

At the pairs $(\varphi^d_1\mathbf 1,\varphi^d_2\mathbf 1)$ of diagonal weights, with the bumps of Proposition~\ref{C-371-24}(a), where $\sigma\circ Z_z=\sigma$ and the kernel of $\mathcal W$ at $u=x-y$ is $\frac32\operatorname{tr}K_3(u)^2=9(4\pi|u|^3)^{-2}>0$, the law forces $b'=b$; at the pairs $(\varphi^d_1o_{12},\varphi^d_2o_{12})$ of Proposition~\ref{C-371-24}(a), placed along $e_3$, where $\sigma\circ Z_z=z^2\sigma$, it then forces $z^2=1$; at the mixed pairs $(\varphi^d_1\mathbf 1,\varphi^d_2o_{12})$ with the same bumps placed along $(e_1+e_2)/\sqrt2$, where $\sigma\circ Z_z=z\sigma$ and the kernel of $\mathcal W$ is $\frac32\operatorname{tr}[K_3(u)o_{12}K_3(u)]=9\hat u_1\hat u_2(4\pi|u|^3)^{-2}>0$, so that $\mathcal W$ is a positive number independent of $d$, it forces $z=1$.

And $\sigma'\circ Z_2$, $\sigma'$ the limit at continuous weights of Lemma~\ref{C-371-14}(d) that agrees with $\sigma$ at the tuples of weights, has every hypothesis of Proposition~\ref{C-371-28} except conservation at two points.
\end{enumerate}
\end{theorem}

\subsection{Construction estimates}

\begin{itemize}
\item[(S1)] Theorems~\YMMref{D}{8}, \YMMref{AJ}{673} and~\YMMref{AD}{607} among the construction results, and the following statements among the construction results built on them: the small-field renormalisation progression with plaquette weights as sources (Proposition~\ref{prop:running-source-chain}, including its zero-source specialization), with its block maps and its small-field sets of each level (Theorem~\YMMref{D}{8}); the never-large histories, that is the histories without large-field regions, and the decomposition of a mean by the scale at which a large field is first created, an exact sum rule (Section~\YMMref{123}{828}, equation~\eqref{eq:ward-first-creation}); the rarity of large fields (Proposition~\ref{prop:ward-fixed-volume-rarity}, Proposition~\ref{prop:local-large-field-rarity}); the sides of the torus that are small over the run of a first-passage datum, its canonical sequence of running couplings, the least admissible side among them when $g\le\gamma^{\rm sd}_W$, and the weights of the first creation of a large field on such a side (equation~\YMMref{(2312)}{841}, Lemmas~\YMMref{126.3}{846} and~\YMMref{125.3}{841}); the working level of a depth (equation~\YMMref{(2322)}{844}) and the carried family at the flat top of the second order, with its Gaussian value (Proposition~\YMMref{125.1}{838}); the chessboard estimates (Proposition~\YMMref{86.3}{681}), the domination of every periodic Gibbs state by the tori (Proposition~\YMMref{91.9}{705}) and that of every larger torus by the least admissible torus (Corollary~\YMMref{86.6}{683}); the pressure gap (Lemma~\ref{lem:canonical-source-system}); the invariance of the periodic Gibbs states under the lattice isometries and their reflection positivity at the nine families of lattice planes (Proposition~\textup{\YMMref{103.4}{758}}); the lattice steps to the cone, that is the energy-momentum inequality and the five-point inequality of the periodic Gibbs states (Proposition~\YMMref{99.3}{739} and Lemma~\ref{lem:spectral-steps}); the dilation of lattices, Gibbs states, tori, boards and chessboard seminorms (Proposition~\YMMref{119.1}{815} and Lemma~\ref{lem:lattice-law-dilation}), with the windows of labels of first-passage data (Proposition~\YMMref{74.1}{615}, Corollary~\YMMref{74.2}{616}); the existence of Gibbs sequences at every reference coupling, obtained from a label in $I(g)$, its first-passage line and the periodic torus states, as in the proof of Theorem~\ref{C-353-14}.
\item[(S2)] The theorem on the scalar limit (Theorem~\textup{\YMMref{AP}{780}}): the canonical decomposition of a scalar limit (the specification preceding Theorem~\YMMref{AL}{746}), with approximants obtained by conditioning on translation-invariant events, and the scalar lifts (Theorems~\YMMref{AL}{746} and~\YMMref{AP}{780}); the invariance of a scalar limit under all isometries, its reflection positivity, its forward cone, its Wightman theory (Theorem~\YMMref{AP}{780}), and its bound $\sigma_0(f\circ\theta,f)\le C_\varphi t^{-a_\Gamma}$, the plate growth~\YMMref{(1896)}{723} of property (R7) of Definition~\YMMref{95.1}{723}, which the scalar limit has by Proposition~\YMMref{105.2}{762}, with $a_\Gamma=3<4$, for $t\in(0,1]$ and $f(x)=\alpha(x_3,x_1)\varphi(x_2)$, $\varphi\in C^1_c(\mathbb R)$ real, $\alpha$ real with support in $[t,2t]\times[-t,t]$, $|\alpha|\le3$, $\operatorname{Lip}\alpha\le12/t$.

Three statements among the construction results beneath it are used through their proofs, which run unchanged at labelled tuples: the existence of a canonical decomposition of the limit of a frame (Theorem~\textup{\YMMref{AL}{746}}), with the iterated plateaux of Definition~\ref{C-355-4}(5), in Proposition~\ref{C-371-30}; the reconstruction of a Wightman theory from data (Theorem~\textup{\YMMref{AO}{778}}) in Proposition~\ref{C-371-16}; and the extension of multilinear forms bounded at separated tuples (Proposition~\YMMref{93.3}{717}), with the growth hypothesis in Proposition~\YMMref{93.5}{719}, in Proposition~\ref{C-371-29}(c).
\item[(S3)] The following abstract facts will also be used. \emph{Reconstruction:} the operator argument in the proof of Proposition~\YMMref{94.1}{720} gives the following statement: let the complex vector space $\mathcal D$ carry a positive semidefinite Hermitian form and a unit vector $\mathbf 1$, and let $(T_t)_{t\ge0}$, $(U(v))_{v\in\mathbb R^N}$ be a semigroup and a group of commuting linear maps fixing $\mathbf 1$, $T_t$ symmetric, $U(v)$ isometric, weakly continuous at $0$, with $c_F=\sup_t|\langle F,T_tF\rangle|<\infty$; then $\|T_tF\|\le\|F\|^{1-2^{-k}}c_F^{2^{-k-1}}$ for all $k$, the $T_t$ give a strongly continuous semigroup of self-adjoint contractions with zero kernel, the $U(v)$ a unitary group, and one projection-valued measure $\mathcal P$ on $[0,\infty)\times\mathbb R^N$ has $T_tU(v)=\int e^{-tE+iv\cdot q}\,d\mathcal P$. \emph{One measure in two pictures} (Lemma~\YMMref{94.3}{722}): if $\mu\ge0$ is finite on $[0,\infty)\times\mathbb R$, $w=(E-iq)/\sqrt2$, and $\int e^{-sw-s'\bar w}d\mu=\langle a,e^{-sH_+}e^{-is'Q}b\rangle$ for $s,s'\ge0$, with $H_+\ge0$ and $Q$ self-adjoint, not necessarily commuting, then $\int_{w\ne0}e^{2\kappa\arg w}d\mu\le\|a\|\|b\|e^{|\kappa|\pi/2}$ for real $\kappa$, so $\mu(\{|q|>E\})=0$. \emph{Uniqueness:} finite complex Borel measures on $[0,\infty)\times\mathbb R^N$ with the same integrals of all $e^{-tE+iv\cdot q}$ are equal. Indeed, for their difference $\nu$, polynomial
approximation in $e^{-E}$ shows that
$\int\varphi(E)e^{iv\cdot q}\,d\nu=0$ for every continuous
$\varphi$ on the compactification $[0,\infty]$.
Fourier uniqueness in $\mathbb R^N$ then gives
$\int\varphi(E)\psi(q)\,d\nu=0$ for every bounded continuous
$\psi$, hence $\nu=0$. This is the measure-uniqueness argument
in the proof of Proposition~\YMMref{96.1}{726}, with
$q\in\mathbb R^N$. \emph{Subgroups:} a subgroup of $\mathrm O(3)$ containing $W_3$ and a nontrivial one-parameter group of rotations is $\mathrm O(3)$ (Lemma~\YMMref{112.4}{785}).
\end{itemize}
From the books we use the identity theorem, the theorems of Morera, Weierstrass, Stone--Weierstrass and Liouville, the maximum principle for subharmonic functions, the spectral theorem, von Neumann's theorem on $A^*A$ and the polar decomposition, the generation of $\mathrm{SO}_0(1,2)$ by two boosts, the reflection positivity of the free massless field at every plane and its passage to Wick polynomials. The statements on the six lattice fields that the argument needs beyond (S1)--(S3) (the admission of the bond forms, in any number, as sources of the progression of (S1), Lemma~\ref{C-371-51}; the passage to $\mathbb R^3$, the reconstructed operators, the two-point law, the power bound, conservation and covariance) are proved below.

\subsection{Mirrors, pictures and the forward cone}

\begin{lemma}[Mirrors and labels]\label{C-371-6}
\begin{enumerate}
\item[(a)] Isometries act as a group on weights, tuples and combinations. If $R\in\mathrm O(3)$, $Rm=n$, $c\in\mathbb R$ and $\iota x=Rx+cn$, then $\iota\{m\cdot x>0\}=\{n\cdot x>c\}$ and $\theta_{n,c}=\iota\theta_m\iota^{-1}$ on points, weights and combinations. The formula for $\Theta_nE\Theta_n$, its eigenspaces, the action $\Theta_n^{\otimes4}$, the characterisation of $\Theta_n\in W_3$, the mixing off $\mathfrak N_{18}$ and the equivalences of positivity in Theorem~\ref{C-371-4}(a) hold.
\item[(b)] If $\sigma$ satisfies (D0), (D2) and the numbers $\sigma(\theta_u\bar F\cdot F)$, $F\in\mathcal D_u$, are real, then $\Theta_u\in G_\sigma$.
\item[(c)] If $\sigma$ satisfies (D0), (D2), $m\in\mathfrak M(\sigma)$, $R\in G_\sigma$ and $Rm=n$, then $\sigma$ is positive at $(n,c)$ for every $c$.
\item[(d)] Let $\sigma$ satisfy (D0)--(D2) and be positive at $(u,0)$, and let $A:V\to V$ commute with $\rho(\Theta_u)$ and satisfy $\sigma((A\theta_u\bar F)\cdot F)\ge0$ for $F\in\mathcal D_u$. Then $[AF]_u=[F]_u$.
\item[(e)] If $\sigma$ satisfies (D0), (D1), then $G_\sigma$ is closed, and so is $G_\sigma[\mathfrak K]$ for every full $\mathfrak K$. A closed subgroup of $\mathrm O(3)$ that contains $W_3$ and one element more is $\mathrm O(3)$.
\end{enumerate}
\end{lemma}
\begin{proof}
(a) $\rho$ is a homomorphism; $n\cdot\iota x-c=m\cdot x$ and $\theta_{n,c}(\iota x)=R(x-2(m\cdot x)m)+cn=\iota(\theta_mx)$. Expanding $(1-2nn^{\rm t})E(1-2nn^{\rm t})$ with $(nn^{\rm t}E)_{\lambda\lambda'}=n_\lambda(En)_{\lambda'}$ gives the formula, and $\Theta_nn=-n$, $\Theta_nm=m$ for $m\perp n$ give the eigenspaces; at $n=(1,1,1)/\sqrt3$ one has $(En)_\lambda=r_\lambda/\sqrt3$ and $n\cdot En=s/3$. From $\varepsilon_{abc}R_{ia}R_{jb}R_{\lambda c}=\det R\,\varepsilon_{ij\lambda}$ the array $\varepsilon_{ij\lambda}\varepsilon_{kl\lambda'}E_{\lambda\lambda'}$ turns with $R^{\otimes4}$, the two determinants cancelling. The entries of $\Theta_n$ are $\delta_{\lambda\lambda'}-2n_\lambda n_{\lambda'}$, while a signed permutation matrix has one nonzero entry in each row: this excludes three nonzero coordinates of $n$ and forces $n_\mu^2=n_\lambda^2=\frac12$ if two are nonzero. Off $\mathfrak N_{18}$, $\Theta_nE_{\mu\mu}\Theta_n$ has the entry $-2(1-2n_\mu^2)n_\mu n_\lambda$ at $(\mu,\lambda)$, or $4n_\mu^2n_\lambda n_\nu$ at $(\lambda,\nu)$, which is not zero. For the equivalences put $G=\theta\bar F$, so that $\theta\bar G\cdot G=\theta\bar F\cdot F$, and use $\theta_{n,c}=\tau_{cn}\theta_n\tau_{-cn}$ with (D2).

(b) A sesquilinear form on a complex space with real diagonal is Hermitian; for real tuples $X',X$ on the side of $u$ this gives $\sigma(\theta_uX'\cdot X)=\sigma(\theta_uX\cdot X')=\sigma(\theta_u(\theta_uX'\cdot X))$. At $X'=\emptyset$ we get $\sigma(\theta_uX)=\sigma(X)$; a translation along $u$ carries every tuple to the side of $u$, and $\theta_u\tau_{su}=\tau_{-su}\theta_u$ with (D2) gives the identity for all $X$; and $\theta_u$ is $\Theta_u$ as a linear isometry.

(c) With $\iota$ as in (a), $\sigma\circ\iota=\sigma$ by (D2) and $R\in G_\sigma$, and for $F\in\mathcal D_{n,c}$ and $G=\iota^{-1}F\in\mathcal D_m$,
\[\sigma(\theta_{n,c}\bar F\cdot F)=\sigma(\iota(\theta_m\bar G\cdot G))=\sigma(\theta_m\bar G\cdot G)\ge0.\]

(d) As in (b), $\sigma(AY)=\sigma(Y)$ for all $Y$; so $[F]_u\mapsto[AF]_u$ extends to an isometry $\mathsf A$ with $\langle\mathsf A\psi,\psi\rangle\ge0$, which is self-adjoint with $\mathsf A^2=1$ and spectrum in $\{1\}$: $\mathsf A=1$. For $A=\rho(\Theta_u)$ and $h$ with values in $V_u^-$ it gives $[(h)]_u=[(-h)]_u=-[(h)]_u$.

(e) If $R_j\to R$, the supports of the $R_jH_a$ stay in fixed disjoint neighbourhoods and $|R_jH_a-RH_a|_\infty\to0$; (D0), (D1) give $\sigma(R_jX)\to\sigma(RX)$, and for full $\mathfrak K$ the set $G_\sigma[\mathfrak K]$ is a group because $\mathfrak K$ is stable under $\mathrm O(3)$. Let $G$ be the subgroup; $G_+=G\cap\mathrm{SO}(3)$ properly contains the $24$ cubic rotations. If $G_+$ is infinite, it has elements $g_j\ne1$, $g_j\to1$, rotations by angles $t_j\to0$ about axes $n_j\to n$, and $g_j^{\lfloor t/t_j\rfloor}$ tends to the rotation by $t$ about $n$; so $G$ contains a one-parameter group, and the lemma on subgroups (S3) applies. If $G_+$ is finite of order $N>24$, the count of the poles of its rotations, $2-2/N=\sum_i(1-1/r_i)$ over three orbits with $1/r_1+1/r_2+1/r_3>1$, the stabilisers at $\pm e_\mu$ being of order divisible by $4$, has the only solution $(4,2,3)$, with $N=24$.
\end{proof}

\begin{lemma}[The picture of a positive normal]\label{C-371-7}
Let $\sigma$ satisfy (D0)--(D2) and $u\in\mathfrak M(\sigma)$. Then the picture $u$ exists: the classes $[F]_u$ are dense, $e^{-tH_u}$ is a strongly continuous semigroup of self-adjoint contractions with zero kernel, $[F]_u\mapsto[\tau_vF]_u$, $v\perp u$, a strongly continuous unitary group, and $\mathcal P_u$ is the unique projection-valued measure on $[0,\infty)\times u^\perp$ with
\begin{equation}\label{C-371-8}
\sigma(\theta_u\bar G'\cdot\tau_wG)=\int e^{-(w\cdot u)E+iw\cdot q}\,dm^u_{G',G}(E,q)\qquad(w\cdot u\ge0);
\end{equation}
$m^u_{G,G}\ge0$ and $|m^u_{G',G}|(B)\le\|\mathcal P_u(B)[G']_u\|\,\|\mathcal P_u(B)[G]_u\|$. For $R\in G_\sigma$, $Ru\in\mathfrak M(\sigma)$ and $[F]_u\mapsto[RF]_{Ru}$ is unitary; it carries $\Omega_u$ and $H_u$ to $\Omega_{Ru}$ and $H_{Ru}$, the translation by $v$ to that by $Rv$, and $\mathcal P_u$ to $\mathcal P_{Ru}$ with $q$ turned by $R$.
\end{lemma}
\begin{proof}
Apply the reconstruction lemma of (S3) to the form on $\mathcal D_u$ with $T_t=\tau_{tu}$ and $U(v)=\tau_v$. $T_t$ is symmetric because $\theta_u\tau_{tu}=\tau_{-tu}\theta_u$ and (D2). In $\theta_uX_i\cdot\tau_{tu}X_k$ the supports keep their diameters and stay at distance at least $\min\{d_1,2d_0\}$, $d_0$ the distance of the supports from the plane and $d_1$ their mutual distance, and $|\Theta_uE\Theta_u|_\infty\le3|E|_\infty$; as (D1) does not depend on the position of the supports, $c_F<\infty$ at every mirror, and the iterated inequality of Schwarz, $\|T_tF\|\le\|F\|^{1-2^{-k}}c_F^{2^{-k-1}}$, gives the contractions as $k\to\infty$. Continuity follows from $|\tau_wH-H|_\infty\le|w|\operatorname{Lip}H$ with (D0), (D1); uniqueness of $\mathcal P_u$ is the uniqueness statement of (S3); the bound on $|m^u_{G',G}|$ is the inequality of Schwarz for the positive form $\langle\cdot,\mathcal P_u(B)\cdot\rangle$. For the transport, $\theta_{Ru}R=R\theta_u$.
\end{proof}

\begin{lemma}[Two mirrors at the angle $\pi/4$]\label{C-371-9}
Let $\sigma$ satisfy (D0)--(D2), $u\perp u'$, and $u,(u+u')/\sqrt2\in\mathfrak M(\sigma)$. Then $\mathcal P_u(\{(E,q):E<|q\cdot u'|\})=0$.
\end{lemma}
\begin{proof}
Let $e_\pm=(u\pm u')/\sqrt2$, $\theta_+=\theta_{e_+}$, and let $F$ be supported in the wedge $\{x\cdot u>|x\cdot u'|\}$; then $F\in\mathcal D_u\cap\mathcal D_{e_+}$, $\theta_+\theta_uF\in\mathcal D_{e_+}$ and $\theta_+(\theta_+\theta_u\bar F)=\theta_u\bar F$. As $se_++s'e_-=\frac{s+s'}{\sqrt2}u+\frac{s-s'}{\sqrt2}u'$, \eqref{C-371-8} in the pictures $u$ and $e_+$ gives for $s,s'\ge0$
\[\int e^{-sw-s'\bar w}\,d\mu=\sigma(\theta_u\bar F\cdot\tau_{se_++s'e_-}F)=\langle[\theta_+\theta_uF]_{e_+},e^{-sH_{e_+}}e^{-is'Q}[F]_{e_+}\rangle,\]
where $w=(E-iq)/\sqrt2$, $\mu\ge0$ is the image of $m^u_{F,F}$ under $(E,q)\mapsto(E,q\cdot u')$, and $Q=-e_-\cdot P_{e_+}$. The lemma on one measure in two pictures (S3) gives $\int_{w\ne0}e^{2\kappa\arg w}d\mu\le C e^{|\kappa|\pi/2}$ for all real $\kappa$; letting $\kappa\to\pm\infty$ gives $|\arg w|\le\pi/4$ $\mu$-almost everywhere, that is $\mu(\{|q|>E\})=0$. Every $F'\in\mathcal D_u$ is carried into the wedge by some $\tau_{ru}$, and $e^{-rH_u}$ has zero kernel.
\end{proof}

\begin{lemma}[A lemma on a cylinder]\label{C-371-10}
Let $r,\kappa>0$, $C<\infty$, and let $B$ be a finite set of angles containing $0$ and $\beta_1$ such that every real $t$ lies at a distance less than $\pi/2$ from $B+2\pi\mathbb Z$; for instance $B_8=\frac\pi4\{0,\dots,7\}$ or $B_4=\frac\pi2\{0,1,2,3\}$. For $\beta\in B$ let $m_\beta$ be complex measures on $\{(E,q)\in\mathbb R^2:E\ge|q|\}$ of total variation at most $C$, and
\[f_\beta(c;\rho)=\int\exp\bigl(2rE-\rho[E\cos(c-\beta)-iq\sin(c-\beta)]\bigr)\,dm_\beta(E,q)\]
where the integrand is bounded. Suppose (i) $f_\beta(t;\rho)=f_{\beta'}(t;\rho)$ for $\beta,\beta'\in B$, $\rho>0$ and real $t$ with $\rho\cos(t-\beta)>2r$ and $\rho\cos(t-\beta')>2r$; (ii) $\mu=e^{2rE}m_0$ is a finite positive measure; (iii) $m_{\beta_1}$ is carried by $\{E\ge\kappa\}$. Then $\mu(\{E^2-q^2<\kappa^2\})=0$.
\end{lemma}
\begin{proof}
Fix $\gamma_0\ne0$ and $0<\delta<\min\{1/4,|\gamma_0|/2\}$. For $c=t+i\gamma$ the real part of $E\cos(c-\beta)-iq\sin(c-\beta)$ is $\cos(t-\beta)(E\cosh\gamma+q\sinh\gamma)\ge\cos(t-\beta)Ee^{-|\gamma|}$ on the cone; so $f_\beta$ is holomorphic and bounded by $C$ where $\rho\cos(t-\beta)e^{-|\gamma|}>2r$. By (i) and the identity theorem the $f_\beta$ are restrictions of one bounded $2\pi$-periodic holomorphic function $f$ on a band $|\operatorname{Im}c|<\Gamma(\rho)$, $\Gamma(\rho)=\log(\rho\cos\vartheta_B/2r)$, $\vartheta_B<\pi/2$ the largest distance of a real $t$ from $B$. By (iii), $|f|\le Ce^{-\kappa\rho(e^{-\delta}\cos\delta-2r/\rho)}$ on the disc of radius $\delta$ about $\beta_1$. With $N(c)=\log|2\sin\frac{c-\beta_1}2|$, harmonic off $\beta_1+2\pi\mathbb Z$ with $e^{2N(c)}=2\cosh\gamma-2\cos(t-\beta_1)$, $\Gamma'=\Gamma(\rho)-1$, $M=\log(2\sinh(\Gamma'/2))$ and $h=(M-N)/(M+\log(2/\delta))$, one has $h\le0$ on $|\gamma|=\Gamma'$ and $h\le1$ on the boundary of the disc, and $h(i\gamma_0)\to1$ as $\rho\to\infty$; the maximum principle for the subharmonic function $\log|f|-\log C+\kappa\rho(e^{-\delta}\cos\delta-2r/\rho)h$ on the band less the discs about $\beta_1+2\pi\mathbb Z$ gives
\[\int e^{-\rho(E\cosh\gamma_0+q\sinh\gamma_0)}\,d\mu=f(i\gamma_0;\rho)\le Ce^{-\kappa\rho(e^{-\delta}\cos\delta-o(1))}\qquad(\rho\to\infty).\]
For fixed $\gamma_0\ne0$, first let $\rho\to\infty$ and then $\delta\downarrow0$. Positivity of $\mu$ gives $E\cosh\gamma_0+q\sinh\gamma_0\ge\kappa$ almost everywhere. Apply this on a countable dense subset of $\mathbb R\setminus\{0\}$; continuity in $\gamma_0$ then gives the inequality for every real $\gamma_0$ on one set of full measure. Its infimum over $\gamma_0$ is $(E^2-q^2)^{1/2}$ for $E>|q|$ and $0$ for $E=|q|$, proving the assertion.
\end{proof}

\begin{proposition}[The forward cone from three mirrors]\label{C-371-11}
\begin{enumerate}
\item[(a)] Let $\sigma$ satisfy (D0)--(D2), let $(u,u',u'')$ be an orthonormal frame, and let $\sigma$ be positive at $(u,0)$, $((u+u')/\sqrt2,0)$ and $((u+u'')/\sqrt2,0)$. Then $\sigma$ is invariant under the cubic group of the frame, positive at $(n,c)$ for its $18$ normals and all $c$, and $\mathcal P_n(\{E<|q|\})=0$ at each of them.
\item[(b)] Let $\sigma$ satisfy (D0)--(D2), let $(u,u',u'')$ be an orthonormal frame, let $u,u'\in\mathfrak M(\sigma)$, and $\mathcal P_u(\{E<|q\cdot u'|\})=\mathcal P_{u'}(\{E<|q\cdot u|\})=\mathcal P_u(\{E<|q\cdot u''|\})=0$. Then $\mathcal P_{\pm u}(\{E<|q|\})=\mathcal P_{\pm u'}(\{E<|q|\})=0$; no diagonal mirror is used.
\end{enumerate}
\end{proposition}
\begin{proof}
\emph{Symmetries.} By Lemma~\ref{C-371-6}(b) the three reflections lie in $G_\sigma$. With $\mathsf p=\Theta_{(u+u')/\sqrt2}$, $\mathsf p'=\Theta_{(u+u'')/\sqrt2}$ one has $\mathsf p\Theta_u\mathsf p=\Theta_{u'}$, $\mathsf p'\Theta_u\mathsf p'=\Theta_{u''}$, and $\Theta_u\Theta_{u'}\mathsf p$, $\Theta_u\Theta_{u''}\mathsf p'$ exchange $u$ with $u'$, resp. with $u''$. These elements generate the cubic group of the frame, which carries the three normals to all $18$; Lemma~\ref{C-371-6}(c) gives positivity there, Lemma~\ref{C-371-6}(a) at every $c$.

\emph{Eight pictures.} Let $n_\beta=\cos\beta\,u+\sin\beta\,u'$, $n'_\beta=n_{\beta+\pi/2}$, $\beta\in B_8$. Since $(n_\beta+n'_\beta)/\sqrt2=n_{\beta+\pi/4}$ is a positive normal, Lemma~\ref{C-371-9} gives $E\ge|q\cdot n'_\beta|$ on the carrier of $\mathcal P_{n_\beta}$; at $(u,u'')$ it gives $E\ge|q\cdot u''|$ in the picture $u$.

\emph{One number in eight pictures.} Let $F\in\mathcal D_u$ have its supports in $\{|x|<r\}$, and put $\mathsf U(x)=\sigma(\theta_u\bar F\cdot\tau_xF)$ for $|x|\ge2r$. With $G'_\beta=\theta_{n_\beta}\tau_{-rn_\beta}\theta_uF$ and $G_\beta=\tau_{rn_\beta}F$, both in $\mathcal D_{n_\beta}$, (D2) and \eqref{C-371-8} give, whenever $x\cdot n_\beta\ge2r$,
\begin{equation}\label{C-371-12}
\mathsf U(x)=\int e^{-(x\cdot n_\beta-2r)E+ix\cdot q}\,dm^\beta(E,q),\qquad m^\beta=m^{n_\beta}_{G'_\beta,G_\beta},\quad m^0=e^{-2rE}m^u_{F,F}.
\end{equation}
The left side does not depend on $\beta$: $\sigma$ is symmetric and the tuple $\theta_u\bar F\cdot\tau_xF$ is unordered, also when $x$ has gone around the support of $F$. This symmetry takes the place of the rotations.

\emph{The cut.} The momentum $q\cdot u''$ is common to the eight pictures. Let $\kappa>0$; restrict $m^\beta$ to $\{|q\cdot u''|\ge\kappa\}$ and carry it to the variables $(E,q\cdot n'_\beta)$. At $x=\rho n_t+su''$ the uniqueness of the Fourier transform in $s$ turns \eqref{C-371-12} into (i) of Lemma~\ref{C-371-10}; (ii) comes from $m^0$, and (iii), with $\beta_1=0$, from $E\ge|q\cdot u''|\ge\kappa$. The lemma gives $m^u_{F,F}(\{|q\cdot u''|\ge\kappa,\ E^2-(q\cdot u')^2<\kappa^2\})=0$ for every $\kappa>0$; if $E\ge|q\cdot u'|$ and $E^2<(q\cdot u')^2+(q\cdot u'')^2$ there is a rational $\kappa$ with $(E^2-(q\cdot u')^2)^{1/2}<\kappa\le|q\cdot u''|$; so $m^u_{F,F}(\{E<|q|\})=0$. The classes are dense.

\emph{All $18$ normals.} The hypotheses (eight positive normals $n_\beta$, the cone $E\ge|q\cdot n'_\beta|$ in each picture $n_\beta$, the cone $E\ge|q\cdot u''|$ in one of them) do not change if $(u,u')$ is replaced by $(n_{\beta_0},n_{\beta_0+\pi/2})$, $\beta_0\in B_8$; then the picture $u$, which carries $E\ge|q\cdot u''|$, sits at the angle $\beta_1=-\beta_0$, which Lemma~\ref{C-371-10} admits. So the cone holds in the eight pictures $n_\beta$; the frames $(u,u'',u')$ and $(u',u'',u)$ give the remaining normals.

(b) By Lemma~\ref{C-371-6}(a), $-u,-u'\in\mathfrak M(\sigma)$. $\Theta_u\in G_\sigma$ (Lemma~\ref{C-371-6}(b)), and Lemma~\ref{C-371-7} at $R=\Theta_u$ gives a unitary map $\mathcal H_u\to\mathcal H_{-u}$ that carries $\mathcal P_u(B)$ to $\mathcal P_{-u}(B)$, as $\Theta_uq=q$ on $u^\perp$; so the square cones hold in the pictures $-u$, $-u'$ too, for the four angles of $B_4$. The parts \emph{one number}, \emph{the cut} of the proof of (a) use nothing else of the pictures; Lemma~\ref{C-371-10} holds with $B_4$; for $\pm u'$ replace the frame as there, with $\beta_0\in B_4$.
\end{proof}

\begin{remark}\label{C-371-31}
The proof uses all real translations, through \eqref{C-371-12} at every $x$ with $x\cdot n_\beta\ge2r$. Discrete translations would not do: on a lattice of unit spacing the free field has the joint spectrum $\cosh E=1+\sum_k(1-\cos q_k)$, which meets $\{E<|q|\}$, e.g.\ at $q=(\pi,0)$, where $E=\operatorname{arccosh}3<\pi$.
\end{remark}

\begin{proposition}[Boosts and the identification of pictures]\label{C-371-13}
Let $\sigma$ satisfy (D0)--(D2), be invariant under the quarter turn $\mathsf q$ about $e_2$ with $\mathsf qe_3=e_1$, and be positive at $(e_3,0)$ and at $(n_-,0)$, $n_\pm=(e_3\pm e_1)/\sqrt2$, with $\mathcal P(\{E<|q_1|\})=0$. These hypotheses hold if $\sigma$ is positive at the three mirrors $(e_3,0)$, $(n_-,0)$, $((e_3+e_2)/\sqrt2,0)$; then also $\mathcal P(\{E<|q_2|\})=0$. Let $T(x)=e^{-x_3H+ix_1P_1}$ for $x=(x_1,0,x_3)$, $x_3\ge0$.
\begin{enumerate}
\item[(a)] The classes of the tuples in the quadrant $\{x_1<0<x_3\}$ are dense, and so are those of $\{x_1>0,x_3>0\}$. $I_0:[F]\mapsto[F]_{n_-}$, $F$ supported in the quadrant $\{x_1<0<x_3\}$, is a well defined closable operator from $\mathcal H(\sigma)$ into $\mathcal H_{n_-}$, with closure $\bar I$; $D=\bar I^*\bar I$ is positive, self-adjoint and injective, $D[F]=[\mathsf qF]$ and $\langle[F],D[F]\rangle=\sigma(\theta_{n_-}\bar F\cdot F)$ for these $F$, and $D\Omega=\Omega$.
\item[(b)] $DT(x)=T(\mathsf qx)D$ on $\operatorname{dom}D$ for $x_1\le0\le x_3$; $D$ commutes with the translations along $e_2$; and for real $s$
\[D^{is}(H\pm P_1)D^{-is}=e^{\pm\pi s/2}(H\pm P_1),\qquad D^{is}\mathcal P(B)D^{-is}=\mathcal P(\Lambda^{(1)}_{-\pi s/2}B).\]
\item[(c)] If besides $\mathcal P(\{E<|q_2|\})=0$, then $\mathcal P(\{E<|q|\})=0$. If $\sigma$ is invariant under the exchange of $e_1$ and $e_2$, then for every $\Lambda\in\mathrm{SO}_0(1,2)$ there is a unitary $V_\Lambda$ with $V_\Lambda\Omega=\Omega$ and $V_\Lambda\mathcal P(B)V_\Lambda^{-1}=\mathcal P(\Lambda B)$; the null sets and the support of $\mathcal P$ are invariant under $\mathrm{SO}_0(1,2)$.
\item[(d)] $\bar I=WD^{1/2}$ with a unitary $W:\mathcal H(\sigma)\to\mathcal H_{n_-}$ that carries $\Omega$, $e^{-tH}$, the translation by $v_1e_1+v_2e_2$ and $\mathcal P$ to $\Omega_{n_-}$, $e^{-tH_{n_-}}$, the translation by $v_1n_++v_2e_2$ and $\mathcal P_{n_-}$, the momenta of the picture $n_-$ being taken in the basis $(n_+,e_2)$.
\end{enumerate}
\end{proposition}
\begin{proof}
Under the three mirrors the reflections lie in $G_\sigma$, $\mathsf q^{-1}=\Theta_{n_-}\Theta_{e_3}$, and Lemma~\ref{C-371-9} gives $E\ge|q_1|$ and $E\ge|q_2|$ on the carrier of $\mathcal P$. As $\theta_{n_-}=\theta_{e_3}\mathsf q=\mathsf q^{-1}\theta_{e_3}$, for $F,G$ in the quadrant
\[\langle[G]_{n_-},[F]_{n_-}\rangle=\sigma(\theta_{n_-}\bar G\cdot F)=\langle[\mathsf qG],[F]\rangle=\langle[G],[\mathsf qF]\rangle,\]
which gives (a) by von Neumann's theorem on $\bar I^*\bar I$: if $[F]=0$, then $\|[F]_{n_-}\|^2=\langle[\mathsf qF],[F]\rangle=0$; a limit point of $I_0$ at $0$ is orthogonal to the range, in whose closure it lies; for $\eta\in\operatorname{dom}\bar I$ the identity, extended by closure, reads $\langle\bar I\eta,I_0[F]\rangle=\langle\eta,[\mathsf qF]\rangle$, so $I_0[F]\in\operatorname{dom}\bar I^*$ and $D[F]=[\mathsf qF]$; and injectivity follows from the density of the classes of $\{x_1>0,x_3>0\}=\mathsf q\{x_1<0<x_3\}$. The quadrant classes are dense because $v\mapsto\langle\psi,e^{ivP_1}e^{-\eta H}[F]\rangle$ is holomorphic in the strip $|\operatorname{Im}v|<\eta$, by $E\ge|q_1|$, and vanishes for large $v$ of one sign if $\psi$ is orthogonal to them.

(b) Write $x=(x_1,0,x_3)=an_-+bn_+$. The closed quadrant $\{x_1\le0\le x_3\}$ is $\{|b|\le a\}$, so it holds $x^\dagger=\Theta_{n_+}x=an_--bn_+$ with $x$; and $\Theta_{e_1}x^\dagger=\mathsf qx$. In the picture $n_-$, with its momenta in the basis $(n_+,e_2)$, $T^-(x)=e^{-aH_{n_-}}e^{ib\,n_+\cdot P_{n_-}}$, $a\ge0$, is a contraction with $T^-(x)[F]_{n_-}=[\tau_xF]_{n_-}$ (Lemma~\ref{C-371-7} at $n_-$) and $T^-(x)^*=T^-(x^\dagger)$; likewise $T(y)^*=T(\Theta_{e_1}y)$. Let $x$ lie in the closed quadrant. ($\alpha$) The quadrant is stable under $\tau_x$, so $I_0T(x)=T^-(x)I_0$ on the quadrant classes; both operators being bounded and $\bar I$ closed, $T(x)\operatorname{dom}\bar I\subset\operatorname{dom}\bar I$ and $\bar IT(x)=T^-(x)\bar I$ there. ($\beta$) For $\xi\in\operatorname{dom}D$ and $\eta\in\operatorname{dom}\bar I$, by ($\alpha$) at $x$ and at $x^\dagger$,
\[\langle\bar I\eta,\bar IT(x)\xi\rangle=\langle T^-(x^\dagger)\bar I\eta,\bar I\xi\rangle=\langle\bar IT(x^\dagger)\eta,\bar I\xi\rangle=\langle T(x^\dagger)\eta,D\xi\rangle=\langle\eta,T(\mathsf qx)D\xi\rangle;\]
so $\bar IT(x)\xi\in\operatorname{dom}\bar I^*$, that is $T(x)\xi\in\operatorname{dom}D$, and $DT(x)\xi=T(\mathsf qx)D\xi$. ($\gamma$) The translations $\mathsf U$, $\mathsf U^-$ by $ve_2$ in the two pictures are unitary, keep the quadrant and satisfy $I_0\mathsf U=\mathsf U^-I_0$; so $\bar I\mathsf U=\mathsf U^-\bar I$, and $\mathsf U^{-1}D\mathsf U=\bar I^*(\mathsf U^-)^{-1}\mathsf U^-\bar I=D$: $D$ commutes with the translations along $e_2$. Put $T(w,w')=\int\exp(-[w(E-q_1)+w'(E+q_1)]/2)\,d\mathcal P$ for $\operatorname{Re}w,\operatorname{Re}w'\ge0$, a contraction by $E\ge|q_1|$; then $T(x)=T(w_x,\bar w_x)$ with $w_x=x_3+ix_1$, and $T(\mathsf qx)=T(iw_x,-i\bar w_x)$. For $\xi,\eta\in\operatorname{dom}D$ the functions $\langle D\eta,T(w,w')\xi\rangle$ and $\langle\eta,T(iw,-iw')D\xi\rangle$ are holomorphic on the sector $-\pi/2<\arg w<0<\arg w'<\pi/2$; by ($\beta$) they agree at its points $(w_x,\bar w_x)$, $x_1<0<x_3$, which form an open set of the real plane of the coordinates $(\frac12(w+w'),\frac1{2i}(w-w'))$; so they agree on the sector, and, $D$ being self-adjoint, $T(w,w')\operatorname{dom}D\subset\operatorname{dom}D$ and $DT(w,w')=T(iw,-iw')D$ on $\operatorname{dom}D$ for $(w,w')$ in the sector.

For such $(w,w')$ and for $\varphi,\psi$ in spectral subspaces of $\log D$ for bounded intervals, $G(z)=\langle D^{i\bar z}\varphi,T(e^{-\pi z/2}w,e^{\pi z/2}w')D^{iz}\psi\rangle$ is bounded and holomorphic on a neighbourhood of $-1\le\operatorname{Im}z\le0$; at $z=r-i$, $r$ real, one has $D^{iz}\psi=D(D^{ir}\psi)$ with $D^{ir}\psi\in\operatorname{dom}D$, the arguments of $T$ are $(iw(r),-iw'(r))$ with $(w(r),w'(r))=(e^{-\pi r/2}w,e^{\pi r/2}w')$ in the sector, and the last identity gives $G(r-i)=G(r)$. So $G$ extends to a bounded entire function, constant by Liouville's theorem:
\[D^{is}T(w,w')D^{-is}=T(e^{-\pi s/2}w,e^{\pi s/2}w'),\]
on the sector and then, by continuity, for $\operatorname{Re}w,\operatorname{Re}w'\ge0$. The uniqueness statement (S3) turns this, with ($\gamma$), into the formulas of (b), $E\mp q_1$ being multiplied by $e^{\mp\pi s/2}$. (c) $\{E<|q|\}\subset\{E<|q_1|\}\cup\bigcup_{\varsigma\in\mathbb Q}\Lambda^{(1)}_\varsigma\{E<|q_2|\}$: if $E\ge|q_1|$ and $(E-q_1)(E+q_1)<q_2^2$, the energy $\frac12((E+q_1)e^{-\varsigma}+(E-q_1)e^{\varsigma})$ of $\Lambda^{(1)}_{-\varsigma}(E,q)$ falls below $|q_2|$ for some rational $\varsigma$. The exchange of $e_1$ and $e_2$ yields, through Lemma~\ref{C-371-7}, a second group that boosts in $(E,q_2)$; the two generate $\mathrm{SO}_0(1,2)$. (d) is the polar decomposition $\bar I=WD^{1/2}$: $W$ is unitary as $\ker D=0$ and the range of $I_0$ is dense in $\mathcal H_{n_-}$ (translate along $n_-$ into the quadrant).

For $x_1<0<x_3$ the identity of (b), continued to $s=-i/2$ on the spectral subspaces of $\log D$ for bounded intervals, which form a core of $D^{1/2}$, gives $D^{1/2}T(x)=T(x^\flat)D^{1/2}$ on $\operatorname{dom}D^{1/2}=\operatorname{dom}\bar I$, $w_{x^\flat}=e^{i\pi/4}w_x$. With ($\alpha$), $T^-(x)WD^{1/2}=\bar IT(x)=WD^{1/2}T(x)=WT(x^\flat)D^{1/2}$, and the range of $D^{1/2}$ is dense; so $T^-(x)W=WT(x^\flat)$. For $x=an_-+bn_+$, $|b|<a$, one has $x^\flat=ae_3+be_1$, hence $W\,e^{-aH}e^{ibP_1}W^{-1}=e^{-aH_{n_-}}e^{ib\,n_+\cdot P_{n_-}}$; at $b=0$ this is $We^{-aH}W^{-1}=e^{-aH_{n_-}}$, and as $e^{-aH_{n_-}}$ has zero kernel, $We^{ibP_1}W^{-1}=e^{ib\,n_+\cdot P_{n_-}}$ for $|b|<a$, so for all $b$.

By ($\gamma$) and the functional calculus $D^{1/2}$ commutes with the translations along $e_2$, and $\bar I\mathsf U=\mathsf U^-\bar I$ gives $W\mathsf U=\mathsf U^-W$; the spectral measures correspond by uniqueness.
\end{proof}

\begin{lemma}[One correlation in two pictures]\label{C-371-19}
Let $\sigma$ satisfy (D0)--(D2), let $u,n\in\mathfrak M(\sigma)$ with $u\cdot n=\cos\gamma$, $0<\gamma<\pi/2$, $u'=(n-\cos\gamma\,u)/\sin\gamma$, $m=\sin\gamma\,u-\cos\gamma\,u'$, $u''\perp u,u'$; take the momenta of the picture $u$ in the basis $(u',u'')$ and those of the picture $n$ in $(m,u'')$, $q=(q_1,q_2)$, and suppose $\mathcal P_u(\{E<|q_1|\})=\mathcal P_n(\{E<|q_1|\})=0$. Let $F$ be supported in $W_\gamma=\{x:u\cdot x>\tan\gamma\,|u'\cdot x|\}$, $\mathsf m=m^u_{F,F}\ge0$ and $\mathsf n=m^n_{\theta_n\theta_uF,F}$.
\begin{enumerate}
\item[(a)] For real $w=w_0u+w_1u'+w_2u''$ with $w_0\ge0$ and $w\cdot n\ge0$:
$\int e^{-w_0E+i(w_1q_1+w_2q_2)}d\mathsf m=\sigma(\theta_u\bar F\cdot\tau_wF)=\int e^{-(w\cdot n)E+i((w\cdot m)q_1+w_2q_2)}d\mathsf n$.
\item[(b)] For a bounded Borel $g\ge0$ let $\mu_g$ be the image of $g(q_2)\mathsf m$ under $(E,q)\mapsto(E-q_1,E+q_1)$ and $\nu_g$ that of $g(q_2)\mathsf n$ under $(E,q)\mapsto(E+q_1,E-q_1)$, measures on $[0,\infty)^2$. Then for $\gamma-\pi/2<\arg z<\pi/2$ and $-\pi/2<\arg z'<\pi/2-\gamma$
\[\int e^{-(za+z'b)/2}\,d\mu_g(a,b)=\int e^{-(e^{-i\gamma}za+e^{i\gamma}z'b)/2}\,d\nu_g(a,b).\]
\item[(c)] For real $\kappa_1,\kappa_2$: $\int_{(0,\infty)^2}a^{-i\kappa_1}b^{-i\kappa_2}\,d\mu_g=e^{-(\kappa_1-\kappa_2)\gamma}\int_{(0,\infty)^2}a^{-i\kappa_1}b^{-i\kappa_2}\,d\nu_g$, with the like identities on the two half axes and $\mu_g(\{(0,0)\})=\nu_g(\{(0,0)\})$: at a fixed value of $ab=E^2-q_1^2$ and of $q_2$, $\nu_g$ is the continuation of $\mu_g$ in the rapidity $\frac12\log(b/a)$ by $i\gamma$.
\item[(d)] $\int\Phi(E^2-q_1^2,q_2)\,d\mathsf m=\int\Phi(E^2-q_1^2,q_2)\,d\mathsf n$ for every bounded Borel $\Phi$. So, with $\mathsf T(E,q)=(E^2-q_1^2,q_2)$, $\mathcal P_u(\mathsf T^{-1}B)=0$ if and only if $\mathcal P_n(\mathsf T^{-1}B)=0$, for every Borel $B\subset[0,\infty)\times\mathbb R$; in particular the forward cone holds in the picture $u$ if and only if it holds in the picture $n$.
\end{enumerate}
\end{lemma}
\begin{proof}
(a) On $W_\gamma$, $n\cdot x\ge\cos\gamma(u\cdot x-\tan\gamma|u'\cdot x|)>0$ and $n\cdot\theta_ux<0$; so $F$ and $\theta_n\theta_uF$ lie in $\mathcal D_n$, and $\theta_n(\theta_n\theta_u\bar F)=\theta_u\bar F$; \eqref{C-371-8} in the two pictures. (b) Both outer members of (a) are, as functions of $w_2$, Fourier transforms of measures in $q_2$; their equality lets one insert $g(q_2)$. With $z=w_0+iw_1$, $z'=w_0-iw_1$ one has $-w_0E+iw_1q_1=-[z(E-q_1)+z'(E+q_1)]/2$ and $-(w\cdot n)E+i(w\cdot m)q_1=-[e^{-i\gamma}z(E+q_1)+e^{i\gamma}z'(E-q_1)]/2$; so the identity holds at $z'=\bar z$, $\gamma-\pi/2\le\arg z\le\pi/2$. Both members are holomorphic on the product of the two sectors (the cones in $q_1$ make the integrands bounded), and agree on a non-empty open set of its real points: they agree. (c) Integrate (b) at $(z,z')=(2\varrho,2\varrho')$ against $d_{\eta,\mathsf z}(\varrho)d_{\eta,\mathsf z'}(\varrho')$, where $d_{\eta,\mathsf z}=c_{\eta,\mathsf z}-\eta c_{\eta,\mathsf z+1}$, $c_{\eta,\mathsf z}(\varrho)=\Gamma(\mathsf z)^{-1}\varrho^{\mathsf z-1}e^{-\eta\varrho}$, whose Laplace transform $\zeta(\zeta+\eta)^{-\mathsf z-1}$ vanishes at $\zeta=0$ and is bounded by $\max\{1,\eta^{-1}\}e^{|\operatorname{Im}\mathsf z|\gamma}$ on $|\arg\zeta|\le\gamma$; let $\operatorname{Re}\mathsf z\downarrow0$, then $\eta\downarrow0$, by dominated convergence, and use $(e^{\mp i\gamma}x)^{-i\kappa}=x^{-i\kappa}e^{\mp\kappa\gamma}$; the half axes come from $\varrho\to\infty$ in one variable. (d) At $\kappa_1=\kappa_2$ the images of the two measures under $\log ab$ have one Fourier transform, and the masses on $\{ab=0\}$ agree; rectangles generate the Borel sets.

If $\mathcal P_n(\mathsf T^{-1}B)=0$, then $\mathsf m(\mathsf T^{-1}B)=\|\mathcal P_u(\mathsf T^{-1}B)[F]_u\|^2=0$ for all $F$ in $W_\gamma$, whose classes are dense in $\mathcal H_u$. Conversely, (a) holds by polarisation for pairs $F',F$ in $W_\gamma$, and the classes $[\theta_n\theta_uF']_n$ and $[F]_n$, $F,F'$ in $W_\gamma$, are total in $\mathcal H_n$: the translations along $u=\cos\gamma\,n+\sin\gamma\,m$ and along $\theta_n\theta_uu=\cos\gamma\,n-\sin\gamma\,m$, which in the picture $n$ are $e^{-s\cos\gamma H_n}$ times unitaries, carry every compact set into $W_\gamma$, resp. into $\theta_n\theta_uW_\gamma$, for large $s$, and for $\psi\in\mathcal H_n$, $G\in\mathcal D_n$ the function $s\mapsto\langle\psi,e^{-s\cos\gamma H_n}e^{\pm is\sin\gamma\,m\cdot P_n}[G]_n\rangle$ is holomorphic on the sector $|\arg s|<\pi/2-\gamma$ by the cone in $q_1$; if it vanishes for large real $s$, it vanishes at $s\downarrow0$.
\end{proof}

\begin{lemma}[How much covariance one plane uses]\label{C-371-20}
Let $\sigma$ satisfy (D0)--(D4).
\begin{enumerate}
\item[(a)] For a full $\mathfrak K$, $G_\sigma[\mathfrak K]$ is a closed subgroup of $\mathrm O(3)$ containing $W_3$, so it is $W_3$ or $\mathrm O(3)$: covariance on $\mathfrak K$ under one $R\notin W_3$ is covariance on $\mathfrak K$ under $\mathrm O(3)$.
\item[(b)] Let $\mathfrak F$, $\mathfrak K$ be full, with $X\cdot Y\in\mathfrak K$ for all $X,Y\in\mathfrak F$ with disjoint supports, and $G_\sigma[\mathfrak K]=\mathrm O(3)$. Then $\sigma(\theta_{n,c}\bar F\cdot F)\ge0$ for every mirror and every combination $F$ of tuples of $\mathfrak F$ on its side. For one mirror $(u,c)$ and one $F=\sum\zeta_iX_i$ the proof uses $\sigma(gY)=\sigma(Y)$ at one isometry $gx=Rx+cu$, $Rn=u$, $n\in\mathfrak N_{18}$, and at the tuples $Y=g^{-1}(\theta_{u,c}X_i\cdot X_k)$ alone.
\item[(c)] Instances of (b) are $(\mathfrak F,\mathfrak K)=(\mathfrak K_m,\mathfrak K_{2m})$, $(\mathfrak T_k,\mathfrak T_{2k})$, $(\mathfrak T_0\cup\mathfrak K_1,\mathfrak T_1\cup\mathfrak K_2)$. In particular covariance at order $2$ gives positivity at every mirror on the vectors of at most one entry, and covariance at order $2$ together with $G_\sigma[\mathfrak T_1]=\mathrm O(3)$ gives positivity at every mirror on the combinations of scalar tuples and of single tensor entries; $G_\sigma[\mathfrak T_0]=\mathrm O(3)$ is the invariance of the trace part.
\item[(d)] Let $\mathfrak F$ be full with $\emptyset\in\mathfrak F$, $u\in S^2$, and let $\sigma(\theta_{u,c}\bar F\cdot F)\ge0$ for all real $c$ and all combinations $F$ of tuples of $\mathfrak F$ on the side of the plane. Then $\sigma(\Theta_uX)=\sigma(X)$ for $X\in\mathfrak F$, and for $X=X'\cdot X''$ with $X',X''\in\mathfrak F$ on the two sides of a plane normal to $u$.
\item[(e)] $\sigma((H))=c_\sigma\int\operatorname{tr}H$ with a number $c_\sigma$. If $\sigma(H,H')=\iint s(y-x;H(x),H'(y))\,dx\,dy$ with $s$ continuous on $(\mathbb R^3\setminus\{0\})\times V\times V$, then positivity of the two-point form at the planes normal to one $u\notin\mathfrak N_{18}$ gives covariance at order $2$; if $s$ is real analytic and $\sigma(RH,RH')=\sigma(H,H')$ for one $R\notin W_3$ and all $H\in C^1_c(B;V)$, $H'\in C^1_c(B';V)$, $B,B'$ open balls with disjoint closures, then $\sigma$ is covariant at order $2$.
\end{enumerate}
\end{lemma}
\begin{proof}
(a) Lemma~\ref{C-371-6}(e) and (D3). (b) As in Lemma~\ref{C-371-6}(c): $Y_i=g^{-1}X_i\in\mathfrak F$ lies on the side of $n$, $Z_{ik}=\theta_nY_i\cdot Y_k\in\mathfrak K$, $\theta_{u,c}X_i\cdot X_k=gZ_{ik}$, $\sigma(gZ_{ik})=\sigma(Z_{ik})$, and $\sum\bar\zeta_i\zeta_k\sigma(Z_{ik})\ge0$ by (D4). (c) A concatenation of two tuples with at most $m$ entries has at most $2m$, and the same for the entries outside $\mathbb R\mathbf 1$; at most one non-scalar entry in each of two tuples gives at most two. (d) As in Lemma~\ref{C-371-6}(b), with $F'=\emptyset$, and with $F'=\theta X'$, $F=X''$. (e) $\ell(H)=\sigma((H))$ is linear, invariant under the translations and bounded in sup norm for supports in a fixed ball; so $f\mapsto\ell(fE)$ is a translation-invariant Radon measure, a multiple of Lebesgue measure, and $\ell(H)=\int\operatorname{tr}(\mathsf cH)$ with $\rho(w)\mathsf c=\mathsf c$ for $w\in W_3$: the sign changes kill the entries of $\mathsf c$ off the diagonal and the permutations make the diagonal constant.

For the second assertion, (d) at $\mathfrak F=\mathfrak K_1$ says $s(\Theta_uv;\rho(\Theta_u)e,\rho(\Theta_u)e')=s(v;e,e')$ for $u\cdot v\ne0$, by continuity for all $v\ne0$; so $\Theta_u\in G_\sigma[\mathfrak K_2]$, as two continuous kernels of one bilinear form agree, and $\Theta_u\notin W_3$; (a). If $s$ is real analytic, $s(Rv;\rho(R)e,\rho(R)e')-s(v;e,e')$ is real analytic on the connected set $\mathbb R^3\setminus\{0\}$ and vanishes on the open set $B'-B$; for a general pair cut both entries by partitions of unity into pieces so small that a coordinate plane parts any two of them.
\end{proof}

\subsection{Examples}

\begin{proposition}[Functionals with exactly nine families, and the inputs of the cone]\label{C-371-21}
\begin{enumerate}
\item[(a)] The functional $\sigma^{\rm fr}$ of Theorem~\ref{C-371-4}(f)(i) satisfies (D0)--(D4) and is covariant under $\mathrm O(3)$; at two entries it is 
\[
\sigma^{\rm fr}(H_1,H_2)=\mathcal W(H_1,H_2)=\frac32\iint\operatorname{tr}[H_1(x)K_3(x-y)H_2(y)K_3(x-y)]\,dx\,dy
\]
. The functional $\sigma^{\rm fr}\circ Z_2$ satisfies (D0)--(D4), with $2^n$ times the numbers of (D1), is not covariant at order $2$, has the forward cone in the nine pictures, and $\mathfrak M(\sigma^{\rm fr}\circ Z_2)=\mathfrak N_{18}$.
\item[(b)] The functional 
\[
\sigma_D(H_1,\dots,H_n)=\sum_{\rm pairings}\prod_{\{a,b\}}S_D(\operatorname{tr}H_a,\operatorname{tr}H_b),
\]
$S_D(f,f')=(2\pi)^{-3}\int\overline{\hat f(p)}\hat f'(p)\hat S_D(p)\,dp$, with $\hat S_D$ of Theorem~\ref{C-371-4}(f)(ii), has (D0)--(D3) and is positive at every coordinate mirror; on the vectors $[(h)]$ the joint spectrum of $(H,P_1,P_2)$ is the surface $E=\omega(q)$, $\omega(q)^2=(m^2+|q|^2+\lambda q_1^2q_2^2)/(1+\lambda|q|^2)$, with $\omega(q_1,0)^2\le\max\{m^2,1/\lambda\}$; so $\mathcal P(\{E<c|q_1|\})\ne0$ for every $c>0$, and $\sigma_D$ is positive at no diagonal mirror.
\item[(c)] The functional $\sigma_c$ formed in the same way with $\hat S_c(p)=[p_1^2+c^2p_2^2+p_3^2+m^2]^{-1}$, $0<c<1$, $m>0$, has (D0)--(D2), is positive at $(u,0)$ for every unit $u\perp e_2$, and on the vectors $[(h)]$ its spectral measure is carried by $\{E^2=q_1^2+c^2q_2^2+m^2\}$ and charges $\{E<|q|\}$.
\end{enumerate}
\end{proposition}
\begin{proof}
(a) (D0) is Wick's rule: at disjoint supports $\sigma^{\rm fr}$ is a finite sum over the pairings without self-contraction of integrals of products of $n$ kernels $K_3$; $|K_3(u)|_\infty\le(2\pi|u|^3)^{-1}$ gives (D1), and $K_3(Ru)=RK_3(u)R^{\rm t}$ gives (D2), (D3) and the covariance. Off $0$, $-K_3$ is the covariance of the gradient of the free massless scalar field $\phi$: $\langle\partial_\alpha\phi(x)\,\partial_\gamma\phi(y)\rangle=\partial^x_\alpha\partial^y_\gamma(4\pi|x-y|)^{-1}=-K_3(x-y)_{\alpha\gamma}$ for $x\ne y$. The free field is reflection positive at every plane; so for every $n\in S^2$ the reflected covariance of the gradient, a vector field that turns with $\Theta_n$, is the positive kernel $-\Theta_nK_3(\theta_nx-y)$ on $\{n\cdot x>0\}$.

The field $B$, of covariance $+K_3$ off $0$, is a pseudovector: it turns with $-\Theta_n$, so its reflected covariance is the same positive kernel $-\Theta_nK_3(\theta_nx-y)$, and in $\langle B_\lambda,B_{\lambda'}\rangle$ the two signs cancel, so that the labels turn by $\rho(\Theta_n)$. Equivalently, the Wick squares of the gradient give $(-1)^n\sigma^{\rm fr}$ at $n$ entries, each term having $n$ kernels, and $F=\sum c_iX_i\mapsto\sum c_i(-1)^{|X_i|}X_i$ keeps positivity. Positivity passes from a Gaussian field to its Wick polynomials, after a mollification that is removed at separated supports. For $E=o_{12}$ and $u$ parallel to $e_3$, $\frac32\operatorname{tr}[EK_3(u)EK_3(u)]=3K_{11}K_{22}=3(4\pi|u|^3)^{-2}>0$, so $\sigma^{\rm fr}(\varphi_1E,\varphi_2E)\ne0$ for nonnegative bumps $\varphi_1,\varphi_2$ about two points on a vertical line.

Proposition~\ref{C-371-17} makes $\sigma^{\rm fr}\circ Z_2$ a functional with (D0)--(D4) that is not covariant at order $2$; Theorem~\ref{C-371-4}(b) gives $\mathfrak M=\mathfrak N_{18}$, and Proposition~\ref{C-371-11} the cone.

(b) $\hat S_D$ is invariant under the signed permutations of $p$, and $\operatorname{tr}(\rho(R)E)=\operatorname{tr}E$: (D0), (D2), (D3). With $q=(p_1,p_2)$, $A=m^2+|q|^2+\lambda q_1^2q_2^2$, $B'=1+\lambda|q|^2$ and $\omega=(A/B')^{1/2}$, the integral over $p_3$ gives
\[S_D(z)=(2\pi)^{-2}\int e^{iq\cdot\vec z-\omega(q)|z_3|}\bigl(2(AB')^{1/2}\bigr)^{-1}d^2q\qquad(z_3\ne0),\]
which converges uniformly for $|z_3|\ge2\delta$. Indeed, on $|q_1|\ge\max\{1,|q_2|\}$ one has $(AB')^{-1/2}\le[\sqrt\lambda\,q_1^2\sqrt{1+\lambda q_2^2}]^{-1}$; integration in $q_1$ leaves an integrable multiple of $[\max\{1,|q_2|\}\sqrt{1+\lambda q_2^2}]^{-1}$. The other unbounded wedge follows by symmetry, and the compact part is bounded since $m>0$. Thus the kernel is bounded for $|z|_\infty\ge2\delta$ by (D3): (D1). For $h,h'$ in $\{x_3>0\}$, $f=\operatorname{tr}h$, $\tilde f(q)=\int f(x)e^{-iq\cdot\vec x-\omega(q)x_3}dx$: $\sigma_D(\theta h',\tau_{(v,t)}h)=(2\pi)^{-2}\int\overline{\tilde f'(q)}\tilde f(q)e^{-iq\cdot v-\omega(q)t}(2(AB')^{1/2})^{-1}d^2q$; at $h'=h$, $v=0$, $t=0$ it is $\ge0$. A Gaussian functional is positive at a mirror if its two-point form is: sorting the pairings of $\theta X'\cdot X$ by the pairs that cross the mirror, $\sigma_D(\theta X'\cdot X)=\langle\Phi(X'),\Phi(X)\rangle$ with $\Phi(X)=\sum_{A\subset X}c(X\setminus A)(|A|!)^{1/2}\mathrm{Sym}\bigotimes_{a\in A}[f_a]$ in the symmetric Fock space over the one-entry space, $c(Y)$ the sum over the pairings of $Y$; so $\sigma_D(\theta\bar F\cdot F)=\|\sum\zeta_i\Phi(X_i)\|^2$; (D3) gives the other coordinate mirrors. By uniqueness of the transform the spectral measure of $[(h)]$ is the image of $|\tilde f|^2(8\pi^2(AB')^{1/2})^{-1}d^2q$ under $q\mapsto(\omega(q),-q)$; $\tilde f$ is real analytic and not $0$ for $f\ge0$, $f\ne0$, and the open set $\{\omega(q)<c|q_1|\}$ is not empty. By Lemma~\ref{C-371-9} with (D3), positivity at one diagonal mirror would give $E\ge|q_1|$.

(c) With $\mathsf s_c(x)=(x_1,x_2/c,x_3)$, $\sigma_c$ is, up to the factor $c^n$, the Gaussian functional of the free field of mass $m$ evaluated at $f_a\circ\mathsf s_c^{-1}$; $\mathsf s_c$ commutes with $\theta_u$ for $u\perp e_2$, so positivity passes from the free field; the transform in $p_3$ is $(2E)^{-1}e^{-|x_3|E}$, $E=(q_1^2+c^2q_2^2+m^2)^{1/2}$, which is $<|q|$ where $(1-c^2)q_2^2>m^2$.
\end{proof}

\subsection{The lattice input: bond forms as sources, passage, operators, the two-point law, dilation and a power bound}

\begin{lemma}[The bond forms as sources]\label{C-371-51}
Let $0<g\le\gamma^{\rm ref}_W$, let $(K,\beta)$ be a first-passage datum at $g$ and $\ell$ an admissible side; plaquettes and bonds are those of $T_\ell$ at the spacing $a_K$. For a weight $\mathsf w$ of the bonds let $S[\mathsf w]$ be the set of the plaquettes that contain a bond $b$ with $\mathsf w(b)\ne0$. A majorant is a pair $(s,S)$ of a number $s\ge0$ and a set $S$ of plaquettes; with the constants $\mathsf c^\flat$, $M=L^{m'}$ and $\alpha_0$ of the progression of (S1) put $\mathsf m_j=12M^3\alpha_0^2L^{K-j}s$ for $j\ge0$. A continuous gauge-invariant function $F$ of the configuration is a source under $(s,S)$ if: (i) $|F|\le a_K^{-1}s\sum_{p\in S}a_p$; (ii) at height $0$ a family represents the variation of $F$ from its zero-fluctuation background value in the first chart integral; at each level $j\ge1$, $F$ is the total of a family of local terms without base point, with the same sum rule as for a plaquette weight; these families vanish at the regions that hold no plaquette of $S$, and their sizes are at most those that the progression grants to the plaquette weight $s\mathbf 1_S$: $\mathsf c^\flat\mathsf m_0$ at the height $0$ and $\frac12\mathsf m_j$ at the level $j$; (iii) if $F$ is real: at the small points of the chart of each level, $a_KF$ differs from a quadratic form $\sum_{p,p'}t(p,p')\langle\Phi_p,\Phi_{p'}\rangle$ in the plaquette angles, $\partial U(p)=e^{i\Phi_p}$, by at most the error that the progression admits there for $s\mathbf 1_S$, with a real symmetric kernel $t$ that vanishes off $S\times S$ and has $2\sum_{p'}|t(p,p')|\le s$ for every $p$.
\begin{enumerate}
\item[(a)] $q_b$ is a polynomial in the entries of the bond variables of the four plaquettes that contain $b$ and of their inverses; it is gauge invariant, real on $\mathrm{SU}(2)$, and $|q_b|\le\frac12\sum_{p\ni b}a_p$. So $Y[\mathsf w]$ is gauge invariant and $|Y[\mathsf w]|\le2a_K^{-1}\sup|\mathsf w|\sum_{p\in S[\mathsf w]}a_p$. For a weight $w$ of the plaquettes, $O(w)$ is a source under $(|w|_\infty,\operatorname{supp}w)$, with $t(p,p')=\frac12w(p)\delta_{pp'}$ if $w$ is real; for a weight $\mathsf w$ of the bonds, $Y[\mathsf w]$ is a source under $(59\sup|\mathsf w|,S[\mathsf w])$, with
\[t(p,p')=-\tfrac12\sum_b\mathsf w(b)\bigl[c_{b,\lambda_b}(p)c_{b,\lambda'_b}(p')+c_{b,\lambda_b}(p')c_{b,\lambda'_b}(p)\bigr]\]
if $\mathsf w$ is real, $c_{b,\alpha}(p)\in\{0,\pm\frac12\}$ the coefficients of $u_{b,\alpha}=\sum_pc_{b,\alpha}(p)\Phi_p+\dots$ to the first order in the angles; if $F$, $F'$ are sources under $(s,S)$, $(s',S')$ and $c\in\mathbb C$ ($c$ real for (iii)), then $F+F'$ is one under $(s+s',S\cup S')$ and $cF$ one under $(|c|s,S)$. So $O^T_K(H)$ is a source under $(60|H|_\infty,S_H)$, $S_H$ the set of the plaquettes that carry a weight of $H$ or contain a bond that carries one; for real $H$ its kernel $t_H$, the sum of the kernels of $O(w_H)$ and of $Y[\mathsf w_H]$, has $2\sum_{p'}|t_H(p,p')|\le3|H|_\infty$.
\item[(b)] Let $F_1,\dots,F_{n'}$ be sources under majorants $(s_a,S_a)$, in any finite number $n'$ and mixed with plaquette weights. The progression of Theorem~\YMMref{D}{8} with the source $\sum_a\lambda_aF_a$, $\lambda$ in a polydisc about $0$ in $\mathbb C^{n'}$, is formed, with all its numbers taken at the plaquette weights $s_a\mathbf 1_{S_a}$, and it has the finite-source properties of Proposition~\ref{prop:weighted-finite-system}: its families are holomorphic and Lipschitz in $\lambda$; they are local, a family on a region reading only the sources whose sets meet the region; the density of a never-large history up to the level $t$ is $e^{D^\lambda_t}$ times that at $\lambda=0$, with $D^\lambda_t$ holomorphic in $\lambda$ and $D^0_t=0$; the first derivatives in the $\lambda_a$ at $\lambda=0$, read at a small-field configuration of the level $t$ and less their values at the trivial one, are bounded, summed over the polymers of the levels up to $t$, by a number of the constants times $L^{K-t}s_a$, times the number of the cubes of the level $t$ that meet $S_a$, times the radius of the small fields of the level $t$; and the mixed derivatives at $\lambda=0$ in $r\ge2$ different parameters vanish on the polymers that do not meet the $r$ sets, and are bounded by a number of $r$ and the constants times the product of the $r$ sizes, the number being a constant times $(4r)^r$ when the sets lie far apart in the units of the level.

The decomposition of a mean by the level at which a large field is first created and the quotients of products at the never-large histories hold as in (S1), with all their numbers formed at the plaquette weights $s_a\mathbf 1_{S_a}$: the quotient of $e^{\sum_a\lambda_aF_a}$ at the never-large history up to the level $t$ is $e^{D^\lambda_t}$, and that of a product $\prod_a(F_a-c_a)$ of centred sources is the sum, over the partitions of the sources, of the products over the blocks of the derivatives of $D^\lambda_t-\sum_ac_a\lambda_a$ at $\lambda=0$ in the parameters of the block.
\item[(c)] Consequently, for $n\ge2$ and $\delta>0$ there are $K_n(g,\delta)$ and $k_n(g,\ell,\delta)$, depending on the constants, $g$, $n$, $\ell$ and $\delta$ alone, such that for $K\ge K_n(g,\delta)$ with $3a_K\le\delta$ and for closed sets $\Sigma_1,\dots,\Sigma_n$ at pairwise sup distances $\ge2\delta$ the following holds. Let each $F_a$ be one of $O^T_K(H_a)$, $O(w_a)$, $Y[\mathsf w_a]$, or a sum of two of these, with a bounded weight $H_a$ that vanishes off $\Sigma_a$ and weights $w_a$, $\mathsf w_a$ that vanish at the plaquettes and bonds whose centres and midpoints lie off $\Sigma_a$, and let $s_a$ be $60|H_a|_\infty$, $|w_a|_\infty$, $59\sup|\mathsf w_a|$, or the sum of the two numbers: so $F_a$ is a source under $(s_a,S_a)$ with the corners of the plaquettes of $S_a$ within the sup distance $a_K$ of $\Sigma_a$. Then the joint cumulant under the Wilson measure of $T_\ell$ satisfies
\[|\langle F_1;\dots;F_n\rangle_{K,\ell,\beta}|\le k_n(g,\ell,\delta)\prod_as_a .\]
If $\ell$ is small over the run of $(K,\beta)$, its canonical sequence of running couplings (equation~\YMMref{(2312)}{841}), as the least admissible side is when $g\le\gamma^{\rm sd}_W$ (Lemma~\YMMref{126.3}{846}), then for $K\ge K^\circ_n(g,\delta)$ the bound holds in the form $k^\circ_n(g,\delta)(a_K\beta)^{-n}\prod_as_a$, with numbers that depend on the constants, $g$, $n$ and $\delta$ alone and read neither the side nor the datum. On such a side, for $0<d<d^\circ_2(g)$ and $K\ge K'_2(g,d)$, numbers formed with the constants, $g$ and $d$, two real sources $F_1$, $F_2$ under majorants $(s_1,S_1)$, $(s_2,S_2)$, with the kernels $t_1$, $t_2$ of (iii) and with the centres of the plaquettes of $S_1$ and of $S_2$ within the distance $d-a_K$ of two points at a distance between $3d$ and $4d$, satisfy
\[\Bigl|\langle F_1;F_2\rangle_{K,\ell,\beta}-6\upsilon^2\sum_{p_1,p'_1,p_2,p'_2}t_1(p_1,p'_1)\,t_2(p_2,p'_2)\,C^{\rm bg}(p_1,p_2)\,C^{\rm bg}(p'_1,p'_2)\Bigr|\le\epsilon'_g(d)(a_K\beta)^{-2}s_1s_2,\]
with $\epsilon'_g(d)$ formed with the constants, $g$ and $d$ and tending to $0$ as $d\to0$; here $C^{\rm bg}$ is the covariance of the background field at the working level of the depth of $d$ (equation~\YMMref{(2322)}{844}) and $\upsilon\le(\frac98)^{1/2}(a_K\beta)^{-1}$ the running coupling of that level. This gives the two-point law of Proposition~\ref{C-371-24} and the bound of its clause (e).
\end{enumerate}
\end{lemma}
\begin{proof}
(a) Under a gauge transformation the loops $W_x(\mu,\pm\alpha)$, read from $x$, are conjugated by its value at $x$, and $\langle\cdot,\cdot\rangle$ is invariant under conjugation. For $V\in\mathrm{SU}(2)$ with $\operatorname{tr}V=2\cos\vartheta$, $\mathbf w(V)$ is Hermitian and $\langle\mathbf w(V),\mathbf w(V)\rangle=\sin^2\vartheta\le2(1-\cos\vartheta)$; so $\langle u_{b,\alpha},u_{b,\alpha}\rangle\le a_p+a_{p'}$ for the two plaquettes $p,p'$ of the plane of $e_\mu,e_\alpha$ at $b$, and $|q_b|\le\frac12(\langle u_{b,\lambda_b},u_{b,\lambda_b}\rangle+\langle u_{b,\lambda'_b},u_{b,\lambda'_b}\rangle)$. A plaquette has four bonds.

For $O(w)$ the three properties are those of the plaquette weights in the progression of (S1), and (iii) is the expansion $a_p=\frac12\langle\Phi_p,\Phi_p\rangle+\dots$. For $Y[\mathsf w]$, (i) is the bound just shown. At the height $0$ the member of the family at a region is the change of $-a_K^{-1}\sum_b\mathsf w(b)q_b$, over the bonds $b$ whose four plaquettes have their corners in the region, under the fluctuation field of the first block step; the region of a bond consists of at most six blocks of the level $0$ with a common corner and holds its four plaquettes, and at most $18M^3$ bonds have one and the same region. At configurations with $|U(b')^{\pm1}|\le e^{1/13}$ whose plaquette variables lie within $\zeta$ of $\mathbf 1$, $e^{2/13}\zeta\le1$, a fluctuation $e^{iB'}$ with $\max|B'|=\rho\le\frac1{128}$ changes $q_b$ by at most $784\rho^2+132\rho\zeta$, by the expansion of $\mathbf w$ to the second order.

At the radii of the first block step this, with the count of the bonds that have a given region, bounds the size of the family of the height $0$ by $18\cdot3138=56484$ units of a quantity proportional to $\sup|\mathsf w|$, of which $\mathsf c^\flat\mathsf m_0$ at $s=59\sup|\mathsf w|$ is $81\cdot12\cdot59=57348$ units; so the family of the height $0$ stays below the bound for the plaquette weight $59\sup|\mathsf w|\,\mathbf 1_S$, $S=S[\mathsf w]$, and this is the origin of the number $59$. At the levels $j\ge1$ the bond forms of the bonds whose blocks of the level $j$ lie in a region, read on the analyticity cores of that level, form a family of local terms without base point whose sizes stay below the bounds for $59\sup|\mathsf w|\,\mathbf 1_S$: at most $98M^3\alpha_0^2L^{K-j}\sup|\mathsf w|$, against $\frac12\mathsf m_j=354M^3\alpha_0^2L^{K-j}\sup|\mathsf w|$ at $s=59\sup|\mathsf w|$.

To the first order $u_{b,\alpha}$ is $\frac12$ times the sum, with signs, of the angles of the two plaquettes of the plane of $e_\mu,e_\alpha$ at $b$, so the second order of $q_b$ is the form of the kernel $t$; two different plaquettes share at most one bond, so $|t(p,p')|\le\frac18\sup|\mathsf w|$, and a plaquette has four bonds with two partners at each, so $2\sum_{p'}|t(p,p')|\le2\sup|\mathsf w|$. The three properties are additive or subadditive in the source and monotone in $(s,S)$, and $|w_H|_\infty+59\sup|\mathsf w_H|\le60|H|_\infty$; for real $H$, $2\sum_{p'}|t_H(p,p')|\le|w_H(p)|+2\sup|\mathsf w_H|\le3|H|_\infty$.

(b) The proofs of the progression with sources, of the decomposition of a mean by the level of the first large field and of the quotients read a source through (i)--(iii) alone; every other step of them is worded for every family of local terms or every bounded Borel function of the configuration, or concerns numbers formed from $s$ and $S$. So they hold for sources under majorants, in any number and mixed with plaquette weights, with the numbers formed at the majorants. The derivatives of all orders in the parameters come from the chain rule for the step of the progression in any number of variables. A family on a region reads the parameters of the sources whose sets meet the region alone, so a mixed derivative vanishes on the polymers that do not meet all its sets; on the polydisc of the radii $1/(4r)$, in the units in which the sizes are measured against the room of the progression, the step stays within that room, and Cauchy's estimate gives the factor $(4r)^r$, the decay of the families in the size of the polymers that join sets far apart paying the combinatorial factor. $D^\lambda_t$ is the difference of the actions of the never-large history at $\lambda$ and at $0$, the source and the families of the progression up to the level $t$ being read at its configuration; so the quotient of $e^{\sum_a\lambda_aF_a}$ is $e^{D^\lambda_t}$, that of $\prod_a(F_a-c_a)$ is the derivative $\partial_{\lambda_1}\cdots\partial_{\lambda_n}$ at $\lambda=0$ of $e^{D^\lambda_t-\sum_ac_a\lambda_a}$, and this derivative of the exponential of a function that vanishes at $0$ is the sum over the partitions of the products of its derivatives in the blocks.

(c) The mean of a product of centred sources is the sum over the level of the first creation of a large field of integrals, along the never-large history, of the quotients of the product; by (b) these are sums over the partitions of the sources of products of mixed derivative families of the progression, and a family of a block of at least two sources lives on regions that meet every $S_a$ of the block, so it is small in the separation; the large fields are paid by their rarity. A joint cumulant is a polynomial in such means; this gives the first bound. On a small side, the least side at $g\le\gamma^{\rm sd}_W$ among them, the progression is run to the working level of the depth fixed by $\delta$, the bound comes in the units $\upsilon s_a$, $\upsilon$ the running coupling of that level, with $\upsilon^2\le\frac98(a_K\beta)^{-2}$, and the weights of the first creation of a large field (Lemma~\YMMref{125.3}{841}) are summable against the sizes of the deeper levels; a small side is at most a number formed with the constants and $g$, by the fourth inequality of equation~\YMMref{(2312)}{841} and the bound of the last coupling of the run, so these numbers read no side.

For two real sources the second cumulant at the never-large history is the second order of the progression: there the two sources are the quadratic forms of their kernels (iii) in the plaquette angles, whose covariance at the working level is the Gaussian term $6\upsilon^2\sum t_1t_2C^{\rm bg}C^{\rm bg}$ displayed in (c), compared with the value of the free field as for plaquette weights; the bound of the carried family at the flat top (Proposition~\YMMref{125.1}{838}), which reads a source only through its kernel and its majorant, carries over with its constants and errors, and the rest is small with the depth of $d$. At the kernels of the bond forms this is the law of Proposition~\ref{C-371-24} at the off-diagonal labels, and the bound of its clause (e) with a number that does not depend on the distance.
\end{proof}

\begin{proposition}[The passage to $\mathbb R^3$]\label{C-371-22}
Let $0<g\le\gamma^{\rm ref}_W$.
\begin{enumerate}
\item[(a)] Let $n\ge1$, $\delta,\varrho_0>0$ and $\delta'=\min\{\delta,1/(8n)\}$. For $r\le n$ let $C^\flat_r$ be a number with $\|D\|_{{\rm ch},\ell'}\le C^\flat_r\prod_a|H_a|_\infty$ for every product $D=\prod_{a\le r}(O^T_K(H_a)-c_a)$, $c_a=c(H_a)$, of centred entries carried by one cell of side $1$, at sup distances $\ge\delta'$ from its walls and $\ge2\delta'$ from one another, on the least admissible torus, uniformly in the first-passage data at $g$ with $K\ge K^G$. Here $c(H)=e_{\ell'}a_K^{-1}\sum_pH_{\lambda(p)\lambda(p)}(x(p))$ for a bounded weight $H$, $e_{\ell'}$ the mean of a plaquette square $a_p$ under the Wilson measure of the least admissible torus, which is the same at every plaquette; $c$ is linear in $H$, and for an entry carried by one cell $c(H)$ is the mean of $O^T_K(H)$ under that measure, the bond forms having mean $0$. $K^G=K^G(n,\delta')$, depending on the constants, $g$, $n$ and $\delta'$ alone, is a cutoff with $3a_K\le\delta'$ for $K\ge K^G$ beyond which the joint cumulants of at most $n\ell'^3$ centred entries of this kind on the least admissible torus are bounded uniformly in the cutoff (Lemma~\ref{C-371-51}(c) on the least side, with $s_a=60|H_a|_\infty$); and put
\[C^{\rm loc}_m=(4n\varrho_0+2)^{3m}\Bigl(\max_{1\le r\le n}\max\{1,e^{G_g}C^\flat_r\}\Bigr)^m,\qquad C^{\rm lc}_n=\sum_{k=0}^n\tbinom nk(C^{\rm loc}_1)^{n-k}C^{\rm loc}_k .\]
Such numbers exist and depend on the constants, $g$, $n$, $\delta$, $\varrho_0$ alone, and for every first-passage datum $(K,\beta)$ at $g$ with $K\ge K^G$, all $\nu,\nu'\in\mathcal G^{\rm per}_{K,\beta}$ and all bounded weights $H_1,\dots,H_n$ that vanish off closed sets of sup diameters $\le\varrho_0$ at pairwise sup distances $\ge2\delta$:
\[\Bigl|\nu\Bigl(\prod_a\bigl(O^T_K(H_a)-\nu'(O^T_K(H_a))\bigr)\Bigr)\Bigr|\le C^{\rm lc}_n\prod_a|H_a|_\infty;\]
the same holds under the Wilson measures of the admissible tori whose period boxes hold the supports with room. No position, cutoff, volume or state enters the numbers.
\item[(b)] Every Gibbs sequence at $g$ has a subsequence along which \eqref{C-371-2} exists at every tuple of $C^1$ weights and at every tuple of continuous weights, and along which the labels converge, to $b\in[g^{-2},L(g^{-2}+\bar B)]$.
\item[(c)] Such a limit $\sigma$ is real, symmetric, $\sigma((H))=0$, and has (D0)--(D4) with $C_n=C^{\rm lc}_n$; for every $d>0$, $|\sigma(X\cdot\tau_vY)|$ is bounded uniformly over translations $v$ for which all supports in $X\cdot\tau_vY$ remain at pairwise sup distance at least $d$; its truncated functions are the limits of the lattice ones; its trace part is the scalar limit of the subsequence; and if $X_j\to X$ entry by entry in the sup norm, with supports in a fixed compact set at pairwise distances $\ge d>0$, the lattice correlations at $X_j$ tend to $\sigma(X)$.
\item[(d)] For every component $\xi$ of the decomposition of the scalar limit, with approximants $\rho_j$, a further subsequence has a lift $\sigma^\sharp_\xi$; it has (D0)--(D4) with the numbers of (a), its trace part is the component $\sigma_{0,\xi}$, and $\sigma^\sharp_\xi((H))=\frac{c_\xi}3\int\operatorname{tr}H$, $c_\xi$ the number with $\rho_j(O^T_{K_j}(f\mathbf 1))-\nu_j(O^T_{K_j}(f\mathbf 1))\to c_\xi\int f$ of the theorem on the scalar limit (Theorem~\textup{\YMMref{AP}{780}}).
\item[(e)] Along growing tori, with the Wilson measures of admissible sides $\ell_j\to\infty$ in the place of the $\nu_j$, (b) and (c) hold, except that of (D4) only the three coordinate families of mirrors are asserted.
\end{enumerate}
\end{proposition}
\begin{proof}
(a) Put $E_a=O^T_K(H_a)$. The Wilson measure of the least torus is invariant under the lattice isometries $\hat r$, which act transitively on the plaquettes and satisfy $\hat rO^T_K(H)=O^T_K(\hat rH)$, $(\hat rH)(x)=RH(\hat r^{-1}x)R^{\rm t}$ with $R$ the linear part of $\hat r$; so $a_p$ has one mean $e_{\ell'}$, and the reflection in the plane through $b$ normal to $e_{\lambda_b}$ keeps $b$ and changes the sign of $q_b$, so $q_b$ has mean $0$: $c(H)$ is the mean of $O^T_K(H)$ for an entry carried by one cell, and $c(\hat rH)=c(H)$. For a product $D$ of $r$ entries $E_a-c(H_a)$ in one cell of a board of side $l$, $\|D\|_{{\rm ch},\ell'}^N$, $N=(\ell'/l)^3$, is the mean on the least torus of the $N$ mirror copies of $D$; a copy is $\prod_a(O^T_K(\hat rH_a)-c(\hat rH_a))$ with a lattice isometry $\hat r$ and $|\hat rH_a|_\infty=|H_a|_\infty$, so the copies are centred entries of the same kind at pairwise sup distances $\ge2\delta'$.

The mean of a product of centred variables is the sum over the partitions without singletons of the products of their joint cumulants, each bounded uniformly in the cutoff on the least torus by Lemma~\ref{C-371-51}(c), with $s_a=60|H_a|_\infty$; the $N$-th root gives $C^\flat_r$. In a torus of side $\ell\ge\ell'$ the chessboard estimates give $|\tilde\mu_{K,\ell,\beta}(D)|\le\|D\|_{{\rm ch},\ell}\le e^{l^3(p_{\ell'}-p_\ell)}\|D\|_{{\rm ch},\ell'}\le e^{G_gl^3}\|D\|_{{\rm ch},\ell'}$, and the domination of every periodic Gibbs state by the tori gives the same bound in $\nu$; products over several cells are bounded by the products of the bounds, whatever the positions of the cells. For entries of any extent cut $\mathbb R^3$ into the half-open cubes $P$ of side $l/(4n)$ of one grid.

As $H\mapsto O^T_K(H)-c(H)$ is linear, $\prod_a(E_a-c(H_a))$ is the sum, over the choices of one cube $P_a$ for each entry, of the products $\prod_a(O^T_K(H_a\mathbf 1_{P_a})-c(H_a\mathbf 1_{P_a}))$, whose pieces $H_a\mathbf 1_{P_a}$ are bounded weights; at most $(4n\varrho_0/l+2)^{3n}$ of these products are not $0$. For each, in every coordinate the offsets of the board that put a wall within $\delta'$ of one of the $n$ pieces have a measure at most $l/2$ in a period, so an offset in $a_K\mathbb Z^3$ puts every piece into one cell at distance $\ge\delta'$ from the walls, and two pieces of different entries lie $\ge2\delta'$ apart; the bounds above, in the occupied cells, give $C^{\rm loc}_n\prod_a|H_a|_\infty$ for the mean of $\prod_a(E_a-c(H_a))$ under $\nu$.

At $l=1$, with $c_a=c(H_a)$, $\prod_a(E_a-\nu'(E_a))=\sum_J\prod_{a\notin J}(c_a-\nu'(E_a))\prod_{a\in J}(E_a-c_a)$ and $|c_a-\nu'(E_a)|=|\nu'(E_a-c_a)|\le C^{\rm loc}_1|H_a|_\infty$, this gives $C^{\rm lc}_n$.
(b) Let $\mathfrak O$ be the finite unions of open cubes with rational corners, and for $O\in\mathfrak O$ let $\mathfrak D_O$ be a countable sup-dense subset of $C^1_c(O;V)$. A diagonal subsequence converges at every tuple with entries in the $\mathfrak D_{O_a}$, $O_a\in\mathfrak O$ with disjoint closures, and has convergent labels, which lie in the bounded window and converge in its closure; by (a) and multilinearity the lattice correlations at any tuple are Cauchy. The same for continuous weights.
(c) Symmetry, multilinearity, $\sigma(\emptyset)=1$, $\sigma((H))=0$ and the bound pass to the limit. $\nu_j$ is invariant under $a_{K_j}\mathbb Z^3\rtimes W_3$, and real translations are reached by moving the entries, $|\tau_vH-\tau_{v_j}H|_\infty\le\operatorname{Lip}H|v-v_j|$. $\nu_j$ is reflection positive at the lattice planes $\{n\cdot x=c'\}$ of the nine families with $c'\in n\cdot(a_{K_j}\mathbb Z^3)$, that is $c'\in a_{K_j}\mathbb Z$ at the coordinate normals and $c'\in a_{K_j}\mathbb Z/\sqrt2$ at the diagonal ones, and at the bond mirrors $\{x_\mu=c'\}$, $c'\in a_{K_j}(\mathbb Z+\frac12)$, among them the plane $x_3=\frac12a_{K_j}$ (Proposition~\textup{\YMMref{103.4}{758}}); the lattice observables turn with the reflection as the weights do; taking $c_j\to c$ in $n\cdot(a_{K_j}\mathbb Z^3)$ gives positivity at every $c$. $O^T_K(f\mathbf 1)$ is the action density, so the trace part is the scalar limit.
(d) The proofs of (a)--(c) use of $\nu_j$ only that it is a periodic Gibbs state, and of the centring constants only that they are invariant under the lattice isometries and bounded as in (a); so they hold for $\rho_j$. The one-point function is $\rho_j(O^T_K(H))-\nu_j(O^T_K(H))=(e_{\rho_j}-e_{\nu_j})a_K^{-1}\sum_pH_{\lambda(p)\lambda(p)}(x(p))$, $e$ the means of the plaquette squares $a_p$, the bond forms having mean $0$; at $H=f\mathbf 1$ it tends to $c_\xi\int f$, and $a_K^3\sum_{p:\lambda(p)=\lambda}H_{\lambda\lambda}(x(p))\to\int H_{\lambda\lambda}$.
(e) The bounds hold on the tori with $\ell\ge\ell'$, whose Wilson measures are invariant under the lattice isometries and reflection positive at the coordinate planes of the torus.
\end{proof}

\begin{proposition}[The reconstructed operators]\label{C-371-23}
Let $\sigma$ be a tensor limit at $0<g\le\gamma^{\rm ref}_W$.
\begin{enumerate}
\item[(a)] $\mathcal H(\sigma)$ is separable; $e^{-tH}$, $e^{iv\cdot P}$ and $\mathcal P$ are as in Lemma~\ref{C-371-7}, $H\Omega=0$; and $\mathcal P(\{E<\max(|q_1|,|q_2|)\})=0$.
\item[(b)] The eight $w\in W_3$ with $we_3=e_3$ act by unitaries $[F]\mapsto[wF]$, a representation of the dihedral group of order eight that commutes with $H$, fixes $\Omega$ and turns the translations. The closed span $\mathcal H_0$ of the classes of scalar tuples is $\mathcal H(\sigma_0)$ through $[F]_0\mapsto[F\mathbf 1]$; it reduces $H$, $P$, $\mathcal P$ and the dihedral group, and on it the forward cone and the Poincar\'e representation of the scalar theory hold.
\item[(c)] \emph{Slab operators of tuples.} Let $i\in\mathbb Z$, $l=a_i$, $v\in\mathbb R^2$, $Q=(v,0)+(-\frac l2,\frac l2)^2\times(0,l)$, and $X=(h_1,\dots,h_n)\in\mathfrak S_V(Q)$ with supports at sup distances $\ge\delta$ from $\partial Q$ and $\ge2\delta$ from one another. Then $B_l(X)[F]=[X\cdot\tau_{le_3}F]$ is bounded, $\|B_l(X)\|\le C^{\rm sl}_n(g;l,\delta)\prod_a|h_a|_\infty$, where
\[C^{\rm sl}_n(g;l,\delta)=\sum_{k=0}^n\tbinom nk\bigl(e^{G_gl^3}C^\flat_1\bigr)^{n-k}e^{G_gl^3}C^\flat_k,\qquad C^\flat_0=1,\]
and, for $1\le r\le n$, $C^\flat_r=C^\flat_r(g;l,\delta)$ is a number with $\|D\|_{{\rm ch},\ell_l}\le C^\flat_r\prod_a|H_a|_\infty$ for every product $D$ of $r$ entries $O^T_K(H_a)$ centred at their means on $T_{\ell_l}$ and carried by one cell of side $l$, at sup distances $\ge\frac12\delta$ from its walls and $\ge\delta$ from one another, uniformly in the first-passage data at $g$ beyond a cutoff; here $\ell_l$ is the least admissible side with $\ell_l\ge2l$ (the least admissible side $\ell'$ if $l\le\frac12\ell'$), and $C^\flat_r$ is formed as the number $C^\flat_r$ of Proposition~\ref{C-371-22}(a), with $l$, $\frac12\delta$, $\ell_l$ in the places of $1$, $\delta'$, $\ell'$. In particular $C^{\rm sl}_1(g;l,\delta)=(e^{2G_gl^3}+e^{G_gl^3})C^\flat_1\le2e^{2G_gl^3}C^\flat_1$. Moreover $B_l(X)\Omega=[X]$, $B_l(X)^*=B_l(\theta_{e_3,l/2}X)$, and $e^{iv'\cdot P}B_l(X)e^{-iv'\cdot P}=B_l(\tau_{(v',0)}X)$.
\item[(d)] \emph{Slab operators of one entry.} For every $l>0$ and $h\in C_c(\{0<x_3<l\};V_{\mathbb C})$ (a complex entry taken by its expansion; a continuous one through the limit at continuous weights of Lemma~\ref{C-371-14}(d) and the unitary map there, or, what is the same, as the norm limit of the operators of $C^1$ entries that tend to it in sup norm with supports in a fixed compact set), $B_l(h)[F]=[(h)\cdot\tau_{le_3}F]$ is bounded, $\|B_l(h)\|\le C^{\rm sl}(\mathcal K,l)\,b^{-1}|h|_\infty$ for $\operatorname{supp}h\subset\mathcal K$ compact, with a number $C^{\rm sl}(\mathcal K,l)$ depending on the constants, $g$, $\mathcal K$ and $l$; $\langle[F'],B_{l_1}(h_1)\cdots B_{l_n}(h_n)[F]\rangle=\sigma(\theta\bar F'\cdot(h_1)\cdot(\tau_{l_1e_3}h_2)\cdots(\tau_{(l_1+\dots+l_{n-1})e_3}h_n)\cdot\tau_{(l_1+\dots+l_n)e_3}F)$; $e^{-tH}B_l(h)e^{-t'H}=B_{l+t+t'}(\tau_{te_3}h)$; $B_l(h)^*=B_l(h^\star)$, $h^\star(x)=\rho(\Theta_{e_3})\overline{h(x_1,x_2,l-x_3)}$, so that a component with $r$ indices equal to $3$ acquires the sign $(-1)^r$; $B_l(h)$ turns with the translations along $e_1,e_2$ and with the dihedral group.
\item[(e)] \emph{Point operators.} For $s,s'>0$ and $e\in V$ there is exactly one bounded operator $B(s,s';e)$, linear in $e$, with
\[\sigma(\theta(\tau_{se_3}\bar G')\cdot(h)\cdot\tau_{s'e_3}G)=\int\langle[G'],e^{i\vec x\cdot P}B(s+x_3,s'-x_3;h(x))e^{-i\vec x\cdot P}[G]\rangle\,dx\]
for $G,G'\in\mathcal D_{e_3}$ and $h\in C^1_c(\{-s<x_3<s'\};V)$; $\|B(s,s';e)\|\le C_{\rm pt}(\min\{s,s'\})|e|_\infty$ with $C_{\rm pt}(\eta)=6(80/l_\eta)^3C^{\rm sl}_1(g;l_\eta,l_\eta/16)$, $l_\eta$ the largest power $L^{-i}$ with $\frac9{16}L^{-i}\le\eta$; $e^{-tH}B(s,s';e)e^{-t'H}=B(s+t,s'+t';e)$, $B(s,s';e)^*=B(s',s;\rho(\Theta_{e_3})e)$, $(s,s')\mapsto B(s,s';e)$ is norm continuous, and the dihedral group turns $e$ by $\rho(w)$. So $|\sigma(\theta\bar F\cdot(h)\cdot F')|\le C_{\rm pt}(\eta/2)\|h\|_{L^1}\|[\tau_{-\eta e_3}F]\|\|[\tau_{-\eta e_3}F']\|$ for $\eta\in(0,1]$, $F,F'\in\mathcal D_{e_3,\eta}$, $h\in C^1_c(\{|x_3|<\eta/2\};V)$.
\item[(f)] \emph{The product formula.} For $n\ge1$, $e\in V^n$, $\zeta\in(\mathsf{Tb}^\square)^{n-1}$ and $\varepsilon>0$ with $\zeta^k-2\varepsilon e_3\in\mathsf{Tb}^\square$ for all $k$,
\[\mathbf W_n(\zeta;e)=\langle\Omega,B(\varepsilon,\varepsilon;e_1)\mathsf E(\zeta^1-2\varepsilon e_3)B(\varepsilon,\varepsilon;e_2)\cdots\mathsf E(\zeta^{n-1}-2\varepsilon e_3)B(\varepsilon,\varepsilon;e_n)\Omega\rangle\]
does not depend on $\varepsilon$; it is holomorphic in $\zeta$, $n$-linear in $e$, $|\mathbf W_n(\zeta;e)|\le\prod_kC_{\rm pt}(\varepsilon)|e_k|_\infty$, and $\sigma(h_1,\dots,h_n)=\int\mathbf W_n(x^2-x^1,\dots,x^n-x^{n-1};h_1(x^1),\dots,h_n(x^n))\,dx$ for every tuple whose supports are ordered in the third coordinate.
\end{enumerate}
The lifts have (a)--(c), (e), (f) with the same numbers, and (d) with the bound $\|B_l(h)\|\le C^{\rm sl}(\mathcal K,l)\,b^{-1}|h|_\infty+\frac{|c_\xi|}3\bigl|\int\operatorname{tr}h\bigr|$, $c_\xi$ the number of Theorem~\ref{C-371-5}(1).
\end{proposition}
\begin{proof}
(a) The form is positive by (D4) and Hermitian; separability comes from the countable dense family of the proof of Proposition~\ref{C-371-22}(b). Lemma~\ref{C-371-7} gives the operators. For the square cone, the lattice steps to the cone, here the five-point inequality of Lemma~\ref{lem:spectral-steps}, give, for $\nu_j$, $t_j\in a_{K_j}\mathbb N$ and $k\in\{1,2\}$, $u(0)-4u(t_je_3)+u(2t_je_3)+u(t_je_3+t_je_k)+u(t_je_3-t_je_k)\ge0$ with $u(w)=\sigma_j(\theta\bar F\cdot\tau_wF)$; in the limit $\int f_t\,d\langle\psi,\mathcal P\psi\rangle\ge0$ with $f_t(E,q)=4e^{-tE}(\sinh^2(tE/2)-\sin^2(tq_k/2))$, first for $\psi=[F]$ and by density for all $\psi$; at $\psi=\mathcal P(\Delta)\psi'$ this gives $f_t\ge0$ $\mathcal P$-almost everywhere for all rational $t>0$ at once, and $t\downarrow0$ gives $E\ge|q_k|$. (b) (D3) and $w\theta=\theta w$; $\rho(\Theta_{e_3})\mathbf 1=\mathbf 1$ makes the map isometric, and a closed subspace invariant under a self-adjoint semigroup and a unitary group reduces them. (c) The number $C^\flat_r$ exists: the mean on $T_{\ell_l}$ of the $N=(\ell_l/l)^3$ mirror copies of $D$ is the sum over the partitions without singletons of the products of joint cumulants of centred entries of the same kind, at sup distances $\ge\delta$ from one another, and these are bounded uniformly in the cutoff on every admissible torus by Lemma~\ref{C-371-51}(c), with $\frac12\delta$ in the place of its $\delta$, $s_a=60|H_a|_\infty$ and a number that reads the side; on the least torus this is the bound of the proof of Proposition~\ref{C-371-22}(a).

Let $\sigma_j(X)=\nu_j\bigl(\prod_a(O^T_{K_j}(H_a)-\nu_j(O^T_{K_j}(H_a)))\bigr)$ be the functional under the limit \eqref{C-371-2}, extended linearly to combinations. In $\nu_j$ the multiple reflections in the walls of a board of cells of side $l$ with a wall in the plane $x_3=\frac12a_{K_j}$, one cell of which is $(v_j,0)+(-\frac l2,\frac l2)^2\times(\frac12a_{K_j},\frac12a_{K_j}+l)$, $v_j\in a_{K_j}\mathbb Z^2$, $v_j\to v$, and holds the entries of $X$ at sup distances $\ge\frac12\delta$ from its walls for large $j$, give
\[|\sigma_j(\theta_{e_3,a_{K_j}/2}\bar F'\cdot X\cdot\tau_{le_3}F)|\le\|D_j\|_{\rm ch}\,\sigma_j(\theta_{e_3,a_{K_j}/2}\bar F'\cdot F')^{1/2}\sigma_j(\theta_{e_3,a_{K_j}/2}\bar F\cdot F)^{1/2},\]
$D_j$ the product of the entries $E_a=O^T_{K_j}(h_a)$ of $X$ centred in $\nu_j$, and $\|D_j\|_{\rm ch}$ the lower limit of its chessboard seminorms along admissible sides that tend to infinity (the domination of the periodic Gibbs states by the tori, Proposition~\YMMref{91.9}{705}). With $c_a$ the mean of $E_a$ on $T_{\ell_l}$,
\[D_j=\sum_{J\subset\{1,\dots,n\}}\ \prod_{a\notin J}\bigl(c_a-\nu_j(E_a)\bigr)\prod_{a\in J}(E_a-c_a);\]
the chessboard seminorm is subadditive, homogeneous, and equal to $|c|$ at a constant $c$; on every admissible side $\ell\ge\ell_l$ the chessboard estimates with the pressure gap give $\|\prod_{a\in J}(E_a-c_a)\|_{{\rm ch},\ell}\le e^{G_gl^3}C^\flat_{\#J}\prod_{a\in J}|h_a|_\infty$, and the domination gives $|c_a-\nu_j(E_a)|\le e^{G_gl^3}C^\flat_1|h_a|_\infty$. So $\|D_j\|_{\rm ch}\le C^{\rm sl}_n\prod_a|h_a|_\infty$. In the limit (Proposition~\ref{C-371-22}(c)) the moving mirror $\theta_{e_3,a_{K_j}/2}$ becomes $\theta$, and the lattice correlations become the values of $\sigma$; the adjoint follows from $\theta_{e_3,l/2}=\tau_{le_3}\theta$. (d) A compact $\mathcal K$ in the slab is covered by at most $27(\operatorname{diam}_\infty\mathcal K/a_i+2)^3$ cubes $\{|x-z|_\infty\le\frac12a_i-\delta\}$, $\delta=\frac18a_i$, at the least $i$ with $a_i\le\min\{l,\ell'/2,8\operatorname{dist}(\mathcal K,\{x_3\in\{0,l\}\})\}$; a partition of unity cuts $h$ into pieces carried by these cubes, of sup norms at most $|h|_\infty$, and the translations and the semigroup move each cube to the height $z_3=\frac12a_i$.

For one piece: if $g\le\gamma^{\rm sd}_W$, the reflections of (c) with one cell give the basic inequality $|\sigma(\theta\bar F'\cdot(h)\cdot\tau_{a_ie_3}F)|\le c(i,\delta)b^{-1}|h|_\infty\|[F']\|\|[F]\|$, where $c(i,\delta)$ is $236e^{G_ga_i^3}$ times the $N$-th root, $N=(\ell'/a_i)^3$, of the sum over the partitions of the $N$ mirror copies of the products of the bounds $2^r60^rk^\circ_r(g,\delta)$, $k^\circ_r(g,\delta)(a_K\beta)^{-r}\prod_as_a$ with $s_a=60|h|_\infty$ being the bound of the joint cumulants of $r$ such copies on the least torus, in the form of Lemma~\ref{C-371-51}(c) on a small side, which the least side is for $g\le\gamma^{\rm sd}_W$; so $C^{\rm sl}(\mathcal K,l)=27(\operatorname{diam}_\infty\mathcal K/a_i+2)^3c(i,\frac18a_i)$.

If $\gamma^{\rm sd}_W<g\le\gamma^{\rm ref}_W$, (c) at $n=1$ on the box of side $a_i$ about the cube, with the factor $2$ of the real and the imaginary part and $1\le L(g^{-2}+\bar B)b^{-1}$, gives the same with $C^{\rm sl}(\mathcal K,l)=54(\operatorname{diam}_\infty\mathcal K/a_i+2)^3C^{\rm sl}_1(g;a_i,\frac18a_i)L(g^{-2}+\bar B)$. The bound is in the sup norm, so it passes to continuous entries; the rules are identities between values of $\sigma$, at continuous entries values of the limit of Lemma~\ref{C-371-14}(d), which agrees with $\sigma$ at the tuples of weights. (e) Fix $l=a_i$, $\delta=l/16$, $Q=(-\frac l2,\frac l2)^2\times(0,l)$ and $\omega,\omega'\in\mathcal H(\sigma)$; $\mu(h)=\langle e^{-\delta H}\omega',B_l((h))e^{-\delta H}\omega\rangle$ is a bounded linear form of $h$ in the sup norm, and $\mu(\tau_yh)=\langle e^{-(\delta+y_3)H}e^{-i\vec y\cdot P}\omega',B_l((h))e^{-(\delta-y_3)H}e^{-i\vec y\cdot P}\omega\rangle$, which by the square cone continues holomorphically to $y\in\mathbb C^3$ with $|\operatorname{Re}y_3|+|\operatorname{Im}y_1|+|\operatorname{Im}y_2|<\delta$, bounded by $C=C^{\rm sl}_1\|\omega'\|\|\omega\|$.

Cauchy's estimates bound $\mu(\partial_Sh)$ by $(3/\delta)^{|S|}C|h|_\infty$, and the identity 
\[
\chi\partial_1\partial_2\partial_3f=\sum_S(-1)^{3-|S|}\partial_S[(\partial_{S^c}\chi)f]
\]
 for a cut-off $\chi$ shows that $\partial_1\partial_2\partial_3(\chi\mu)$ is a measure: $\mu$ has a bounded Lipschitz density near the centre of $Q$, and $|\mu(h)|\le6(5/\delta)^3C\int|h|_\infty$. Its value at the centre, transported by the translations and the semigroup, defines $B(s,s';e)$; the formula, the bound and the relations follow. The $L^1$ bound is the formula at $s=s'=\eta$ with 
\[
\|B(\eta\pm x_3,\eta\mp x_3;e)\|\le\|B(\eta/2,\eta/2;e)\|.
\]
 (f) For $h\in C^1_c(\{0<x_3<c\};V)$ the formula gives $[(h)\cdot\tau_{ce_3}G]=\mathsf A_c(h)[G]$ with 
\[
\mathsf A_c(h)=\int e^{i\vec x\cdot P}B(x_3,c-x_3;h(x))e^{-i\vec x\cdot P}dx;
\]
by induction $[X]$ is a product of such operators applied to $\Omega$, and $\sigma(X)=\langle\Omega,[X]\rangle$; $B(s,s';e)=e^{-(s-\varepsilon)H}B(\varepsilon,\varepsilon;e)e^{-(s'-\varepsilon)H}$ collects the factors into $\mathsf E(x^{k+1}-x^k-2\varepsilon e_3)$, and $\|\mathsf E\|\le1$ on the square tube by the square cone.
\emph{Lifts.} For a lift the proofs of (a)--(c), (e), (f) run in its approximants $\rho_j$, which are periodic Gibbs states, with the entries centred in $\nu_j$, the bounds used being those of every periodic Gibbs state. For (d): by the identity of Theorem~\ref{C-371-5}(1), $\sigma^\sharp_\xi=\sigma^c_\xi\circ\mathsf S_\xi$ with $\mathsf S_\xi(X)=\sum_J\prod_{a\notin J}s_\xi(H_a)X_J$, $s_\xi(H)=\frac{c_\xi}3\int\operatorname{tr}H$, and $\sigma^c_\xi$ is a tensor limit at $g$; $\mathsf S_\xi$ is multiplicative and commutes with the bar and with the isometries, so $[F]\mapsto[\mathsf S_\xi F]$ is a unitary map of $\mathcal H(\sigma^\sharp_\xi)$ onto $\mathcal H(\sigma^c_\xi)$ that intertwines $e^{-tH}$, the translations and the dihedral group, and it carries $B_l(h)$ to $B^c_l(h)+s_\xi(h)e^{-lH}$, as 
\[
\mathsf S_\xi((h)\cdot\tau_{le_3}F)=(h)\cdot\tau_{le_3}\mathsf S_\xi F+s_\xi(h)\,\tau_{le_3}\mathsf S_\xi F;
\]
(d) for $\sigma^c_\xi$ gives the bound. At the components of the decomposition this algebra is Proposition~\ref{C-371-30}(d).
\end{proof}

\begin{proposition}[The two-point law at short distance]\label{C-371-24}
Let $0<g\le\gamma^{\rm ref}_W$. There are $d_0(g)>0$ and a function $\epsilon_g(d)\to0$ as $d\to0$, with $\sup_{0<d<d_0(g)}(\log\frac1d)^{\mathfrak p_0/4}\epsilon_g(d)<\infty$, depending on the constants and $g$ alone, such that for every tensor limit $\sigma$ at $g$ with label $b$, every $0<d<d_0(g)$, all $x_0,y_0$ with $3d\le|x_0-y_0|\le4d$ and all weights $H_1,H_2$ supported in closed balls of radius less than $d$ about $x_0$, $y_0$:
\[
\begin{multlined}
|\sigma(H_1,H_2)-b^{-2}\mathcal W(H_1,H_2)|\le\epsilon_g(d)\,b^{-2}|H_1|_\infty|H_2|_\infty,\qquad\mathcal W(H_1,H_2)=\tfrac32\iint\operatorname{tr}[H_1(x)K_3(x \\
-y)H_2(y)K_3(x-y)]\,dx\,dy,
\end{multlined}
\]
$K_3(u)=\partial\partial(4\pi|u|)^{-1}=(4\pi|u|^3)^{-1}(3\hat u\hat u^{\rm t}-\mathbf 1)$ for $u\ne0$, $\hat u=u/|u|$, read off $0$ alone as in Theorem~\ref{C-371-4}(f)(i). The same holds at a finite cutoff with two further errors, one for all the volumes: for every first-passage datum $(K,\beta)$ at $g$ with $K\ge K_0(g,d)$, $b=a_K\beta$, in every periodic Gibbs state and under the Wilson measure of every admissible torus $T_\ell$ (the weights read on $T_\ell$, and $\mathcal W$ formed at their lifts to balls about lifts of $x_0$, $y_0$ at the distance $|x_0-y_0|$), and for continuous $H_1,H_2$ carried by closed balls of a radius $d_-<d$ with $4a_K\le d-d_-$,
\[\bigl|\langle O^T_K(H_1);O^T_K(H_2)\rangle-b^{-2}\mathcal W(H_1,H_2)\bigr|\le b^{-2}\bigl[(\epsilon_g(d)+C_g(a_K/d)^2)|H_1|_\infty|H_2|_\infty+r_K(H_1,H_2)\bigr],\]
with a number $C_g$ depending on the constants and $g$, and a number $r_K(H_1,H_2)\ge0$ that reads $K$, $d$ and the moduli of continuity of $H_1$, $H_2$ alone and tends to $0$ as $K\to\infty$; none of the errors reads $\beta$, the side or the state. Consequently:
\begin{enumerate}
\item[(a)] The label exists along every Gibbs sequence with a tensor limit, and $\sigma$ determines it: with a fixed bump $\varphi\ge0$ supported in the ball of radius $\frac12$ about $0$, $\varphi^d_i(x)=\varphi((x-x_i)/d)$ and $x_2-x_1=\frac72de_3$, $\mathcal W(\varphi^d_1o_{12},\varphi^d_2o_{12})$ is a positive number independent of $d$ ($\mathcal W$ is invariant under dilations), and $b^{-2}=\lim_{d\downarrow0}\sigma(\varphi^d_1o_{12},\varphi^d_2o_{12})/\mathcal W(\varphi^d_1o_{12},\varphi^d_2o_{12})$.
\item[(b)] With $w_0=\frac1{12}388^{-6}$, for $d$ with $\epsilon_g(d)<w_0$, the reflection $\theta'$ in a plane $\{x_\mu=c\}$ with $|c-x_{0,\mu}|=\frac32d$, and $\varphi\in C^1_c$, $0\le\varphi\le1$, equal to $1$ on $B(x_0,d/128)$ and to $0$ off $B(x_0,d/64)$: $\sigma(\theta'(\varphi E),\varphi E)\ge(w_0-\epsilon_g(d))b^{-2}\operatorname{tr}E^2$ for every $E\in V$.
\item[(c)] $|\sigma(rH_1,rH_2)-\sigma(H_1,H_2)|\le10\,\epsilon_g(d)b^{-2}|H_1|_\infty|H_2|_\infty$ for every rotation $r$ about $x_0$.
\item[(d)] If $Z=aP_A+eP_E+zP_T$ and $\sigma(ZH_1,ZH_2)$ obeys the same law with the same $b$ and an error $\epsilon'(d)\to0$, then $a=e=z\in\{1,-1\}$.
\end{enumerate}
Moreover, for $0<g\le\gamma^{\rm sd}_W$:
\begin{enumerate}
\item[(e)] for $0<d<1$ and real off-diagonal $h=h_0o_{\mu\nu}$ with $h_0$ supported in $B(x_0,d)$, $x_{0,3}\in[\frac32d,2d]$: $\|[(h)]\|^2=\sigma(\theta h,h)\le C_\infty(g)b^{-2}|h|_\infty^2$, $C_\infty(g)$ depending on the constants and $g$ alone;
\item[(f)] for every weight $h$ supported in a ball of radius $d$ at height at least $\frac32d$: $\|[(h)]\|^2\le C(g)d^{-2}(\log\frac1d)^{20\mathfrak p_0}b^{-2}|h|_\infty^2$ for small $d$.
\end{enumerate}
\end{proposition}
\begin{proof}
Let first $g\le\gamma^{\rm sd}_W$, so that the least admissible side is small. On the least admissible torus, Theorem~\YMMref{D}{8} gives the moments of several close pairs of quadratic observables (plaquette squares and bond forms with bounded coefficients, the bond forms entering the progression as sources, Lemma~\ref{C-371-51}(a), (b)), centred at their means, that are far from one another: the renormalisation progression shows that the conditional mean given the block field is a constant plus something small, and the sum rule among the construction results, the exact decomposition of a mean by the level at which a large-field region is first created (equation~\eqref{eq:ward-first-creation}), shows that the true mean is as close to that constant; the value on the least torus is the Gaussian term of Lemma~\ref{C-371-51}(c) at the kernels $t_{H_1}$, $t_{H_2}$ of Lemma~\ref{C-371-51}(a), the plaquettes of $S_{H_1}$ and of $S_{H_2}$ lying within the distance $d-2a_K$ of $x_0$ and of $y_0$ once $4a_K\le d-d_-$, $d_-<d$ the radius of the balls that carry the weights, and comparing the lattice Gaussian term with $\mathcal W$ gives $b^{-2}\mathcal W$ up to errors that tend to $0$ with $d$, besides two errors of a finite cutoff, $C_g(a_K/d)^2$ from the lattice kernel and $r_K$ from the moduli of continuity of the weights at the scale $a_K$, which tend to $0$ as $K\to\infty$ and read neither $\beta$ nor the side.

The chessboard estimate then bounds, in any periodic Gibbs state or larger torus, the mean of a function of one cell of the board of eight cells of side $l=\frac12\ell'$ by $e^{G_gl^3}$ times the eighth root of the mean, on the least torus, of the product of its eight mirror images; applied to the centred product minus its mean, the products of means cancel and what is left tends to $0$ with $d$, at the rate of the eighth root, $(\log\frac1d)^{-\mathfrak p_0/4}$; the difference of the means between volumes is small by the same estimate applied to one centred field (Theorems~\YMMref{AJ}{673} and~\YMMref{AD}{607} enter through $G_g$). This gives the law for $g\le\gamma^{\rm sd}_W$, at a finite cutoff in every volume and, as $K\to\infty$, for $\sigma$, with $d_0(g)$ a negative power of $L$ fixed by the depth at which the progression is stopped; above, the dilation of Lemma~\ref{C-371-14}(c) carries it over with $d_0(g)=L^{-z}\min_sd_0(g'_s)$, as $\mathcal W(H_1^{(z)},H_2^{(z)})=\mathcal W(H_1,H_2)$ and $\sigma(X)=L^{2z}\sigma_z(X^{(z)})$ at two entries; at a finite cutoff it carries $(K,\beta)$ to $(K-z,\beta)$, $T_\ell$ to the admissible torus $T_{L^z\ell}$ and the periodic Gibbs states to periodic Gibbs states, with $a_{K-z}/(L^zd)=a_K/d$. (a) Two limit points of the bounded labels give two laws for one $\sigma$; at $E=o_{12}$ and $x_2-x_1\parallel e_3$ the kernel of $\mathcal W$ is $3(4\pi|x-y|^3)^{-2}>0$. (b) is the law with the computation $\mathcal W(\theta'(\varphi E),\varphi E)\ge w_0\operatorname{tr}E^2$. (c) $\mathcal W$ is invariant under $\mathrm O(3)$, and two applications of the law. (d) Compare the two laws on pairs within each of the three blocks and across them. (e) is the two-point bound at short distance in every periodic Gibbs state at off-diagonal labels, whose means vanish, carried to the limit: at the off-diagonal weights the bond forms enter the small-field progression and the term of the never-large histories as sources under the majorant $(59\sup|\mathsf w|,S[\mathsf w])$, $S[\mathsf w]$ the plaquettes that contain a bond where $\mathsf w\ne0$, and their second order at small fields as a quadratic form in the plaquette angles of $S[\mathsf w]$ with a kernel whose row sums lie under the same majorant (Lemma~\ref{C-371-51}(a), (b)); the second cumulant of Lemma~\ref{C-371-51}(c) gives the bound with a number that does not depend on $d$. (f) follows from the law.

Let $d<d_0(g)$, and cut $h$ by a partition of unity into a fixed number $m_0$ of pieces carried by balls of radius $\frac14d$ with centres in the ball of radius $d$ that carries $h$; the centres are at heights $y_3\ge\frac12d$. As $e^{-tH}$ is a contraction, the norm of the class of a piece is at most that of the class of the same piece moved down to the height $y'_3=\min\{y_3,d\}\in[\frac12d,d]$. There the law applies to the piece $k$ and its mirror image $\theta k$, at the parameter $d''=\frac47y'_3<d_0(g)$: the radius $\frac14d$ is less than $d''$, and the centres are $2y'_3=\frac72d''$ apart. For weights carried by balls of radius less than $d''$ whose centres are at least $3d''$ apart, $|K_3(u)_{\lambda\lambda'}|\le(2\pi|u|^3)^{-1}$ gives $|\mathcal W(H_1,H_2)|\le\frac32\cdot81(2\pi d''^3)^{-2}(\frac43\pi d''^3)^2|H_1|_\infty|H_2|_\infty=54|H_1|_\infty|H_2|_\infty$.

So $\|[(k)]\|^2=\sigma(\theta k,k)\le(54+\bar\epsilon_g)b^{-2}|h|_\infty^2$, $\bar\epsilon_g=\sup_{0<d<d_0(g)}\epsilon_g(d)$, and $\|[(h)]\|^2\le m_0^2(54+\bar\epsilon_g)b^{-2}|h|_\infty^2$, at all six labels, which is below the bound of (f) for small $d$.
\end{proof}

\begin{lemma}[Exact dilation]\label{C-371-14}
Let $z\in\mathbb N$.
\begin{enumerate}
\item[(a)] For $K\ge z+1$ and bounded $H$ of bounded support, $O^T_K(H)(U)=L^zO^T_{K-z}(H^{(z)})(U^{(z)})$, $U^{(z)}$ the configuration carried to $a_{K-z}\mathbb Z^3$ by $x\mapsto L^zx$. The map $X\mapsto X^{(z)}$ is a bijection of $\mathfrak S_V(O)$ onto $\mathfrak S_V(L^zO)$, keeps sup norms, commutes with $\mathrm O(3)$, and $(\tau_vH)^{(z)}=\tau_{L^zv}H^{(z)}$, $(\theta_{n,c}H)^{(z)}=\theta_{n,L^zc}H^{(z)}$, $\|H^{(z)}\|_{L^1}=L^{3z}\|H\|_{L^1}$, $(H^{\rm cons}[\eta])^{(z)}=H^{\rm cons}[\eta^{[z]}]$ with $\eta^{[z]}(x)=L^z\eta(L^{-z}x)$; on tori the cumulants of $n$ entries scale by $L^{nz}$, the side by $L^z$.
\item[(b)] Let $0<g\le\gamma$ and $\sigma$ be a tensor limit at $g$ with label $b$. Then $g_z=(L^zb/2)^{-1/2}<g$; along a subsequence the dilated sequence $(K_j-z,\beta_j,\nu_j^{(z)})$ is a Gibbs sequence at $g_z$ with a tensor limit $\sigma_z$ of label $L^zb$, and
\[\sigma(X)=L^{nz}\sigma_z(X^{(z)})\qquad(X\text{ with $n$ entries});\]
$[F]_\sigma\mapsto[S_zF]_{\sigma_z}$, with $S_z$ multiplying every entry by $L^z$ after dilating it, is unitary and carries $e^{-tH}$ to $e^{-L^ztH_z}$. Moreover $g_z\le\gamma^{\rm sd}_W$ once $z$ exceeds a number depending on the constants and $g$. The same holds at continuous weights.
\item[(c)] Let $0<g\le\gamma^{\rm ref}_W$, $t_1=L^zg^{-2}-\bar B$, $t_2=L^{z+1}g^{-2}-3\bar B$ and $g'_s=t_s^{-1/2}$. For every first-passage datum $(K,\beta)$ at $g$ with $K\ge z+1$, $(K-z,\beta)$ is a first-passage datum at $g'_1$ or at $g'_2$; $L^z\ell$ is admissible at $g'_s$ for every side $\ell$ admissible at $g$; and $g'_s\le\gamma^{\rm sd}_W$ once $z$ exceeds a number depending on the constants and $g$.
\item[(d)] Let $0<g\le\gamma^{\rm sd}_W$. If the limit \eqref{C-371-2} of a Gibbs sequence at $g$ exists at all tuples of weights, a subsequence has a label and a limit $\sigma'$ at all tuples of continuous weights, which agrees with the first limit at the tuples of weights, with the same label; the restriction of a limit at continuous weights to the tuples of weights is a tensor limit; and $[F]_{\sigma}\mapsto[F]_{\sigma'}$ extends to a unitary map of the two Hilbert spaces that carries $\Omega$, $H$, the translations and $\mathcal P$ into one another. Above $\gamma^{\rm sd}_W$ the same holds through (c).
\end{enumerate}
\end{lemma}
\begin{proof}
(a) The dilation carries $a_K\mathbb Z^3$ onto $a_{K-z}\mathbb Z^3$, plaquettes, bonds, Haar measure and action onto themselves, and the normalisation $a_K^{-1}$ to $L^za_{K-z}^{-1}$; the rest is the chain rule: $\nabla\eta^{[z]}(x)=(\nabla\eta)(L^{-z}x)$. (b) By the dilation statements (Propositions~\YMMref{119.1}{815} and~\ref{lem:lattice-law-dilation}), periodic Gibbs states go to periodic Gibbs states; the datum $(K_j-z,\beta_j)$ passes first at $g_z$ for large $j$, because its label $L^za_{K_j}\beta_j$ tends to $L^zb$, which is an inner point of the interval $I(g_z)=(g_z^{-2}+\bar B,L(g_z^{-2}+\bar B)-2\bar B]$, $g_z^{-2}=L^zb/2$, sufficient for first passage at $g_z$ (Corollary~\YMMref{74.2}{616} and Lemma~\ref{lem:setup-label-coverage}); and $b\ge g^{-2}$ gives $g_z^{-2}\ge L^zg^{-2}/2>g^{-2}$, that is $g_z<g$.

The formula follows from (a), every entry carrying the factor $L^z$, and the unitarity from $(\theta H)^{(z)}=\theta H^{(z)}$. (c) A pair whose label lies in $I(g')$ is a first-passage datum at $g'$ (Corollary~\YMMref{74.2}{616}). Here $I(g'_1)=(L^zg^{-2},L^{z+1}g^{-2}-2\bar B]$ and $I(g'_2)=(L^{z+1}g^{-2}-2\bar B,L^{z+2}g^{-2}-(2L+2)\bar B]$, and their union holds $L^z\mathsf b$ for every $\mathsf b$ in the window $(g^{-2},L(g^{-2}+\bar B)]$, as $(L^{z+1}+2L+2)\bar B\le L^{z+1}(L-1)g^{-2}$ by $g^{-2}\ge4\bar B$, which holds for $g\le\gamma^{\rm ref}_W$; the label of $(K-z,\beta)$ is $L^za_K\beta$. For the sides: $\ell=2L^m$ is admissible at $g$ if and only if $m\ge m_{\mathfrak L}(g)=1+m'+\max\{\log_LR(L(g^{-2}+\bar B)),\lceil\log_LC_{\rm s}\rceil\}$, where $R(x)$ is the least power of $L$ that is $\ge(\log x)^r$ and $m'$, $C_{\rm s}$, $r\ge1$ are constants (equation~\YMMref{(1528)}{615}); and $L^z\ell$ is admissible at $g'$ if and only if $m+z\ge m_{\mathfrak L}(g')$ (equation~\YMMref{(1528)}{615}).

So it suffices that $R(L(t_s+\bar B))\le L^zR(L(g^{-2}+\bar B))$. Now $L(t_1+\bar B)=L^{z+1}g^{-2}$ and $L(t_2+\bar B)=L^{z+2}g^{-2}-2L\bar B$ are at most $L^{z+1}\cdot L(g^{-2}+\bar B)$; with $y_g=\log(L(g^{-2}+\bar B))$,
\[\bigl((z+1)\log L+y_g\bigr)^r\le y_g^r\,L^{r(z+1)/y_g}\le y_g^r\,L^z,\]
because $r(z+1)\le2rz\le y_gz$: at $g\le\gamma^{\rm ref}_W$ one has $y_g\ge\log g^{-2}\ge2\mathfrak p_0>3r+2$, by the logarithmic lower bounds in the reference choice preceding equation~\YMMref{(2326)}{845} and the exponent choice $r\ge2$, $\mathfrak p_0\ge6r+1$ in Theorem~\YMMref{D}{8}. So $L^zR(L(g^{-2}+\bar B))$ is a power of $L$ that is $\ge(\log L(t_s+\bar B))^r$, and it is $\ge R(L(t_s+\bar B))$. The least side grows like a power of $\log g^{-2}$, the dilation like $L^z$. Finally $t_2>t_1\ge(L^z-1)g^{-2}$ gives $g'_2<g'_1<g$ and the last assertion. (d) The bound of Proposition~\ref{C-371-22}(a) holds for bounded weights uniformly in the cutoff; a diagonal extraction over a countable family of continuous tuples, dense in sup norm with controlled supports, gives the subsequence, and at the tuples of weights the full sequence converges already; identities and inequalities between values pass, and so do the Hilbert spaces.
\end{proof}

\begin{proposition}[A power bound for one tensor entry]\label{C-371-15}
Let $0<g\le\gamma^{\rm ref}_W$, $p_T>11/2$ and $0<\delta\le\frac14$. There are $i_0$, $K_0(i)$ and $C_T$, depending on the constants, $g$, $p_T$ and $\delta$, such that for every integer $i\ge i_0$, every first-passage datum $(K,\beta)$ at $g$ with $K\ge K_0(i)$, every board of cells of side $a_i$, every cell $\square$, every side $\ell$ admissible at $g$, and every bounded $H:\mathbb R^3\to V_{\mathbb C}$ vanishing at the points within sup distance $\delta a_i$ of the complement of $\square$,
\[\|O^T_K(H)-c\|_{{\rm ch},\ell}\le C_TL^{p_Ti}|H|_\infty,\]
$c$ the mean of $O^T_K(H)$ on an admissible torus or in a periodic Gibbs state. For $g\le\gamma^{\rm sd}_W$ this holds for every $i\ge1$ and $K\ge K_E(i)$ in the form $\|O^T_K(H)-c\|_{{\rm ch},\ell}\le2C_E\mathsf Y_i^{q_E}L^{(11/2)i}b^{-1}|H|_\infty$, $b=a_K\beta$, $\mathsf Y_i=2(c_E+i\log L)$, where $c_E=\log(2L(g^{-2}+\bar B))$ and $q_E$, the latter formed from the constants alone, do not depend on $\delta$, $C_E$ is $\delta^{-4}$ times a number depending on the constants and $g$, of the order $\delta^{-4}e^{G_g}$, and besides that $\delta$ enters only $K_E(i)$, which depends on the constants, $g$ and $\delta$. Consequently, for $g\le\gamma^{\rm sd}_W$:
\begin{enumerate}
\item[(i)] for entries $H_a$ carried by different cells of one board of side $a_i$, at the margin $\delta a_i$, and $\nu,\nu'\in\mathcal G^{\rm per}_{K,\beta}$: $|\nu(\prod_a(O^T_K(H_a)-\nu'(O^T_K(H_a))))|\le\prod_a2e^{G_g}C_E\mathsf Y_i^{q_E}L^{(11/2)i}b^{-1}|H_a|_\infty$, with no factorial and no order of the cells;
\item[(ii)] every tensor limit at $g$ has $\|B_{a_i}(h)\|\le c_E(i)b^{-1}|h|_\infty$ for $h$ supported in $\{|x-z|_\infty\le\frac38a_i\}$, $z_3=\frac12a_i$, with $c_E(i)=2C_E\mathsf Y_i^{q_E}a_i^{-11/2}$ at $\delta=\frac1{16}$;
\item[(iii)] every tensor limit at $g$ has, for every $p>17/2$, $0\le d'<d$ and $h\in C_c(\mathbb R^3;V_{\mathbb C})$ supported in $\{|x_3|\le d'\}$ (a complex entry taken by its expansion, a continuous one through Lemma~\ref{C-371-14}(d)),
\[\|B_{2d}(\tau_{de_3}h)\|\le C_p\min\{1,d-d'\}^{-p}\,b^{-1}\|h\|_{L^1},\]
$C_p$ depending on the constants, $g$ and $p$ alone; so (R4) of Proposition~\ref{C-371-16} holds with $C_V=2^pC_pb^{-1}\le2^pC_pg^2$; and for $0<\eta\le1$, $C_{\rm pt}(\eta)\le12\cdot80^3e^{2G_g}C'_T(9L/16)^{3+p_T}\eta^{-3-p_T}$, where $C'_T$ is the larger of the number $C_T$ of the first display at $\delta=\frac1{32}$ and of the finitely many numbers $L^{-ip_T}C^\flat_1(g;a_i,\frac1{16}a_i)$, $0\le i<i_0$, of Proposition~\ref{C-371-23}(c).
\end{enumerate}
For $\gamma^{\rm sd}_W<g\le\gamma^{\rm ref}_W$, (iii) holds with $C_p$ depending on the constants, $g$ and $p$, which includes the factor $L^{4z}$ of the dilation.
\end{proposition}
\begin{proof}
Let first $g\le\gamma^{\rm sd}_W$ and $i\ge1$. Let $x_\bullet$ be the run of $(K,\beta)$, $K'=K-i$, and let $g_{(i)}$ be the stopped coupling, $g_{(i)}^{-2}=L^{-1}x_{K'}-B_\sharp$ (equation~\YMMref{(1536)}{616}): then $g_{(i)}<g$, $(K',\beta)$ is a first-passage datum at $g_{(i)}$, $L^i\ell$ is admissible at $g_{(i)}$ for every side $\ell$ admissible at $g$, and $G_{g_{(i)}}\le G_g+\frac92i\log L+\frac18$, one constant alone of the pressure gap reading the coupling (Corollary~\YMMref{74.3}{616} and Lemma~\ref{lem:setup-stopped-arithmetic}). The exact dilation by $L^i$ (Lemma~\ref{C-371-14}(a)) carries the board of side $a_i$ at $(K,\beta)$ to a board of side $1$ at $(K',\beta)$, the entry $O^T_K(H)$ to $L^iO^T_{K'}(H')$ with $H'$ vanishing within $\delta$ of the complement of its unit cell, and the chessboard seminorm on $T_\ell$ at $(K,\beta)$ to that on $T_{L^i\ell}$ at $(K',\beta)$.

Let $F'$ be one of the four real pieces of $O^T_{K'}(H')$ (the plaquette and the bond part of the real and of the imaginary part of $H'$), a source under a majorant $(\mathsf s,S)$ with $\mathsf s\le59|H|_\infty$ (Lemma~\ref{C-371-51}(a)). Let $\ell''$ be the least admissible side at $g_{(i)}$; it is small over the run of $(K',\beta)$ (Lemma~\YMMref{126.3}{846}, as $g_{(i)}<g\le\gamma^{\rm sd}_W$), and it has $N=\ell''^3$ cells of side $1$, with $N\le C(\log\frac1{a_i})^{2\mathfrak p_0-1}$ by the condition (Sd.4) of a small side. The chessboard seminorm of $F'-m_0$ on $T_{\ell''}$, $m_0$ the mean of $F'$ there, is the $N$-th root of the mean of the product of the $N$ mirror copies $F_\kappa-m_0$, and the copies $F_\kappa$ are sources under majorants $(\mathsf s,S_\kappa)$ whose sets lie in $N$ different cells, at sup distances $\ge\delta$ from one another.
\emph{The moments.} For every set $B$ of copies,
\[\Bigl|\Bigl\langle\prod_{\kappa\in B}(F_\kappa-c_\kappa)\Bigr\rangle_{K',\ell'',\beta}\Bigr|\le(\mathsf A\mathsf s)^{\#B},\qquad\mathsf A=2N\mathsf C_1L^{4s},\]
with numbers $c_\kappa$, a number $\mathsf C_1$ depending on the constants alone, and $s$ the depth below $K'$ of the level $\mathsf k=K'-s$ at which two different sets are $M\mathsf u(N)$ units of the level apart, $\mathsf u(N)$ a number of the constants and $N$ that grows like $N\log N$; so $L^s$ is at most a power of $\log\frac1{a_i}$ divided by $\delta$. Indeed, by the sum rule of the first creation (equation~\eqref{eq:ward-first-creation}) the mean is the term of the never-large history up to the level $\mathsf k$ plus the sum, over the level $t\le\mathsf k$ of the first creation of a large field, of its weight $\mathsf w_t$ times an integral along the never-large history of the quotient of the product, the conditional mean given the block field; by Lemma~\ref{C-371-51}(b) the quotient of $e^{\sum\lambda_\kappa F_\kappa}$ is $e^{D^\lambda_t}$, so that of the centred product is a sum over the partitions of $B$ of products of derivatives of $D^\lambda_t$ at $\lambda=0$; a block of one element is at most the size of one source at the level $t$, $\mathsf C_1L^{4(K'-t)}\mathsf s$, a block of $m\ge2$ elements is a sum of mixed derivative families on polymers that join $m$ different cells, which the separation of the cells makes small enough to pay the factor $(4m)^m$, and a set of $r$ elements has at most $r^r$ partitions.

The large fields are paid by their rarity on the small side: $\mathsf w_t\le2M^{-3}\ell''^3L^{3(K'-t)}e^{-\frac12c'p_0(g_t)^2}$ (Lemma~\YMMref{125.3}{841}, which rests on the rarity of large fields, Proposition~\ref{prop:ward-fixed-volume-rarity} and Proposition~\ref{prop:local-large-field-rarity}), $p_0(g_t)$ the large-field threshold at the running coupling $g_t$ of the level $t$, and the condition $(\frac92\log x_{K'}+2)\ell''^3\le\frac12c'p_0(g_{K'})^2$ of a small side gives $\sum_{t\le K'}L^{4(K'-t+1)r}\mathsf w_t\le1$ for every $r\le\ell''^3$: the rarity pays the sizes of as many sources as the side has unit cells. This is Theorem~\YMMref{D}{8}; Theorems~\YMMref{AJ}{673} and~\YMMref{AD}{607} do not enter the moments.
\emph{The seminorm.} At $B$ of one element $|m_0-c_\kappa|\le\mathsf A\mathsf s$, and the expansion of $\prod_\kappa(F_\kappa-m_0)$ about the $c_\kappa$ gives $\|F'-m_0\|_{{\rm ch},\ell''}\le2\mathsf A\mathsf s$: the moment contributes a power of $\log\frac1{a_i}$ alone. Back to every admissible side $L^i\ell$ at $g_{(i)}$ the chessboard estimate pays the pressure gap of one unit cell at the stopped coupling, $e^{G_{g_{(i)}}}\le e^{G_g+\frac18}L^{\frac92i}$ (Theorems~\YMMref{AJ}{673} and~\YMMref{AD}{607}, at zero source); the dilation pays $L^i$; the re-centring at the mean of the entry on an admissible torus or in a periodic Gibbs state pays $2$ (the domination of the periodic Gibbs states by the tori); and $1\le L(g^{-2}+\bar B)b^{-1}$.

So the exponent is $\frac{11}2=1+\frac92$, $1$ from the dilation and $\frac92$ from the pressure gap, and the four pieces, of sizes at most $1$, $59$, $1$, $59$ times $|H|_\infty$, give the factor $480=4(1+59+1+59)$ in $C_E$, which with $\mathsf A\le c\,\mathsf Y_i^{q_E}\delta^{-4}$ is of the order $\delta^{-4}e^{G_g}$. (i) The domination of the periodic Gibbs states by the tori bounds a product over several cells by the product of the seminorms. (ii) The multiple reflections of Proposition~\ref{C-371-23}(c) with one cell. (iii) Let $\vartheta=d-d'$ and $i(\vartheta)$ the least $i\ge i_0$ with $a_i\le\vartheta$. A sup bound becomes an $L^1$ bound on a cube $Q=\{|x-z|_\infty\le\varrho\}$: if $\mathsf b$ maps $C_c(Q)$ linearly into bounded operators with $\|\mathsf b(g)\|\le c\sup|g|$, $\varepsilon\ge12\varrho$, and $\mathsf B^\varepsilon(g)=e^{-\varepsilon H}\mathsf b(g)e^{-\varepsilon H}$ turns with the translations by $|u|_\infty\le2\varrho$ as the point operators do, then writing $h$ as an integral of translates of pieces $\chi(\cdot)h(\cdot+u)$ and using the square cone to control the factors $e^{\pm u_3H}$ against the margin $e^{-\varepsilon H}$ gives $\|\mathsf B^\varepsilon(h)\|\le4c\varrho^{-3}\|h\|_{L^1}$.

For $\vartheta<a_{i_0}$ the minimality gives $a_{i(\vartheta)}>\vartheta/L$; otherwise $i(\vartheta)=i_0$. Thus the preceding estimate gives $\|B_{2d}(\tau_{de_3}h)\|\le4\cdot32^3a_{i(\vartheta)}^{-3}c_E(i(\vartheta))b^{-1}\|h\|_{L^1}$, and $a_{i(\vartheta)}^{-3-11/2}\mathsf Y^{q_E}\le C_p\min\{1,\vartheta\}^{-p}$ for every $p>17/2=\frac{11}2+3$; with $d=\eta$, $d'=\eta/2$ this is (R4). The bound of $C_{\rm pt}$: for $0<\eta\le1$, $l_\eta=a_i$ with $i\ge0$, so $e^{G_gl_\eta^3}\le e^{G_g}$ and $C^{\rm sl}_1(g;l_\eta,\frac1{16}l_\eta)\le2e^{2G_g}C^\flat_1(g;l_\eta,\frac1{16}l_\eta)$, the number $C^\flat_1$ being formed at the margin $\frac1{32}l_\eta$; the first display at $\delta=\frac1{32}$ for $i\ge i_0$, and the finitely many numbers of Proposition~\ref{C-371-23}(c) for $0\le i<i_0$, give $C^\flat_1(g;a_i,\frac1{16}a_i)\le C'_TL^{p_Ti}$; and $l_\eta^{-1}<9L/(16\eta)$, as $\frac9{16}Ll_\eta>\eta$.

For $\gamma^{\rm sd}_W<g\le\gamma^{\rm ref}_W$, Lemma~\ref{C-371-14}(c) carries data, boards and seminorms at $g$ to the couplings $g'_1,g'_2$, and the dilation of an entry costs $L^{z}$ and of the $L^1$ norm $L^{3z}$.
\end{proof}

\subsection{Conservation and covariance}

\begin{lemma}[The cubic remainder after the small scales]\label{C-371-26}
Let $0<g\le\gamma^{\rm sd}_W$, $0<\tau\le\frac14$, $0<\phi<\frac1{42}$ and $c'=\min\{\frac1{14}-\frac\tau7,\frac1{14}-3\phi\}$; for $\phi<\frac1{30}$ one may take $c'=\min\{\frac1{10}-\frac\tau5,\frac1{10}-3\phi\}$. Fix an admissible torus, let $K\ge K_0$, let $b'$ lie in the window of labels at $g$, $\beta=L^Kb'$, and $t=K-\lceil\phi K\rceil$. Let $S[\Xi]=\sum_pL^{2K}\Xi^{\rho(p)}(x(p))\,\mathfrak r^\sharp_{\rho(p)}(p)$ be the cubic lattice entry of a field $\Xi=(\Xi^1,\Xi^2,\Xi^3)$ of directions on the torus, three real functions of the centres of the plaquettes, and $\mathsf H(\Xi)=\max\sup|\Xi|+\max\operatorname{Lip}\Xi$. Here $\mathfrak r^\sharp$ is defined by \eqref{eq:ward-orbital-remainder}. By Proposition~\ref{prop:ward-orbital-cubic}, it is a gauge invariant polynomial on the two adjacent cubes, is cubic in the plaquette defects on the complex bond domain, and satisfies $|\mathfrak r^\sharp_\rho(p)|\le c_{\rm orb}\mu^3$ on real configurations whose eleven plaquette matrices are within $\mu$ of the identity.

A \emph{box} is a set $Q=\{x+a_K\sum_\mu n_\mu e_\mu:0\le n_\mu\le r_\mu\}$ of bare sites with $r\in\{0,1,2\}^3$ and two $r_\mu\ge1$; a \emph{density} on $Q$ is a polynomial $\mathsf j$ in the entries of the bond variables of $Q$, read also at $SL(2,\mathbb C)$-valued configurations and invariant under the gauge transformations on $Q$; it has the \emph{height} $\hat c$ if $|\mathsf j|\le\hat c$ where the bond variables of $Q$ have operator norms $\le2$, and the \emph{order} $m$ with the constant $c$ if besides $|\mathsf j|\le c\,\varrho^m$ where $|\partial U(q)-\mathbf 1|\le\varrho\le1$ at the plaquettes $q$ of $Q$. A \emph{local entry} of the type $(m,c,\hat c,N)$, of the order $m$, is a finite family $\mathfrak J=(Q_\iota,\mathsf j_\iota,\varkappa_\iota)_\iota$ of boxes, densities on them of the height $\hat c$ and the order $m$ with the constant $c$, and complex coefficients, at most $N$ boxes having one lowest site; it is read as $\mathfrak J(U)=\sum_\iota\varkappa_\iota\mathsf j_\iota(U|_{Q_\iota})$, and $|\varkappa|_\infty=\max_\iota|\varkappa_\iota|$.

Small $|s|$ means $|s|<r_{\mathfrak J}$, a number formed with the type, $|\varkappa|_\infty$ and the first coupling $x_0=\beta$ of the run, of the order $1/|\varkappa|_\infty$, for which $s\mathfrak J$ stays within the room of the progression. So $S[\Xi]$ is a local entry of the order three, the boxes being the pairs of cubes with a common face. There are $K_0$, $C$ and $p$, depending on the constants, $g$, $\tau$, $\phi$ and the side, such that:
\begin{enumerate}
\item[(a)] for every local entry $\mathfrak J$ of order at least two and small $|s|$, the image under the block maps of level $k\le t$ of the Wilson measure, restricted to the small-field histories and weighted with $e^{s\mathfrak J}$, has on the small-field set of level $k$ the density $\exp[\sum_{i\le k}c_i(s)]\,e^{\Phi^s_k}$, with numbers $c_i(s)$ and a bounded function $\Phi^s_k$: the entry enters the renormalisation progression of Theorem~\YMMref{D}{8} as a source;
\item[(b)] for every $\Xi$ there is a number $\mathsf a(\Xi)$ with $|(\mathsf Q_{t\leftarrow0}S[\Xi])(V)-\mathsf a(\Xi)|\le C\,\mathsf H(\Xi)K^pL^{-c'K}$ for $\nu_t$-almost every small-field configuration $V$ of level $t$, $\mathsf Q_{t\leftarrow0}$ the conditional expectation onto the block field of level $t$ along the small-field progression.
\end{enumerate}
\end{lemma}
\begin{proof}
(a) The progression of Theorem~\YMMref{D}{8} is set up with sources. A local entry of order at least two is a sum of local terms on the regions of level zero that hold its boxes, and it enters the first step as a further family of local terms at the height $0$, as the bond forms do in Lemma~\ref{C-371-51}(a): the first block step makes it into its value at the background and one pointed family of level $1$, that is, a family of local terms $E(X,\cdot,z)$ on the polymers $X$ of the level, each with a marked site $z\in X$, measured in a norm with the weight $e^{\kappa d_1(X)}$, $d_1(X)$ the size of $X$ in the units of the level; every further step carries the pointed families of the chain, the families $E^{(i)}$ of the small-field progression, with the families of the entry among them; its order makes these families small at small fields, and for $|s|<r_{\mathfrak J}$ they stay within the room of the progression.

The normalisations $c_i(s)$ are the constants that the steps produce, and the image of the measure is the displayed density. (b) Let $\mathsf Q_kh=\mathsf I[hF_k]/\mathsf I[F_k]$ on the small-field set of level $k+1$, $F_k$ the density of the progression at the level $k$ and $\mathsf I$ the integral over the fluctuation field of the step: the conditional expectation of one step; the tower is $\mathsf Q_{n\leftarrow k}=\mathsf Q_{n-1}\cdots\mathsf Q_k$, and each $\mathsf Q_k$ is a contraction in the sup norm on the small-field set that carries constants to themselves. Put $\varkappa(p)=L^{2K}\Xi^{\rho(p)}(x(p))$, so that $S[\Xi]=\sum_p\varkappa(p)\mathfrak r^\sharp_{\rho(p)}(p)$, and let $\mathsf s_n$ be $S[\Xi]$ read at the background of level $n$ on the small-field set.

To the first order in the source the tower is linear in the weight: $\mathsf Q_k\mathsf s_k=\mathsf s_{k+1}+\Sigma_{k+1}[\mathsf D^{(k+1)}[\varkappa]]$, where $\mathsf D^{(k+1)}[\varkappa]$ is the pointed family born at the step, linear in $\varkappa$, its member on $X$ reading $\varkappa$ on the plaquettes of $X$ alone, and $\Sigma_{k+1}$ is the sum of a pointed family over the polymers and their sites, read at the background. Fix a middle level $n_0=K-\lceil\psi K\rceil$, $\phi<\psi<1$, and a H\"older exponent $0\le\beta_{\rm H}<1$. Then
\[\mathsf Q_{t\leftarrow0}S[\Xi]=\mathsf Q_{t\leftarrow n_0}\mathsf s_{n_0}+\sum_{j\le n_0}\mathsf Q_{t\leftarrow j}\Sigma_j\bigl[\mathsf D^{(j)}[\varkappa]\bigr].\]
\emph{The entry at $n_0$.} The density is of the order three, and at a background of level $n_0$ the plaquette variables are within $\varepsilon_{n_0}L^{-2n_0}$ of $\mathbf 1$, $\varepsilon_n$ the small-field radius of the level $n$; so $\sup|\mathsf s_{n_0}|$ is at most a power of $K$ times $L^{-(1-\frac92\psi)K}\mathsf H(\Xi)$.
\emph{The norm of a birth.} The Gaussian integral over the fluctuation field of a step is invariant under $B\mapsto-B$; split into its even and its odd part in $B$, the birth has an even part of the second order in the fluctuation, of the size $g_k^2$, and an odd part that meets only the odd part of the base family, of the size $g_k(\log x_k)^{\mathfrak p_0}$; so $\|\mathsf D^{(k+1)}[\varkappa]\|\le C(\log x_k)^{\mathfrak p_0}g_k^2L^{-3k}\max|\varkappa|$, $g_k$ the running coupling and $x_k=g_k^{-2}$.
\emph{A family at a background.} The differential of a local term at the trivial configuration vanishes, by its invariance under the constant gauge transformations; so a pointed family of level $j$, read at a background of level $n\ge j$, differs from its value at the trivial configuration by at most $\varepsilon_n^2L^{-4(n-j)}$ times its norm at each of its sites. With the norm of a birth this leaves, for the level $j$, $L^{3(K-n)-j}$ up to logarithms, which is not small at $n=t$; the smallness comes from the inversions.
\emph{The inversions.} The progression is covariant under the lattice isometries of the levels, and the inversion $\mathsf r_z$ through a site $z$ of level $j$ changes the sign of the cubic density, so that the birth of the weight $\varkappa$, read at $\mathsf r_zU$, is the birth of $-\varkappa\circ\mathsf r_z$ read at $U$. Cut $\varkappa$ into the pieces $\theta_z\varkappa$ about the sites $z$ of level $j$, with $\sum_z\theta_z=1$ and $\theta_z\circ\mathsf r_z=\theta_z$, and let $\Sigma_z[\varkappa]$ be the sum of the birth of $\theta_z\varkappa$; this gives the exact split
\[\Sigma_z[\varkappa](U)=\tfrac12\bigl(\Sigma_z[\varkappa](U)-\Sigma_z[\varkappa](\mathsf r_zU)\bigr)+\tfrac12\Sigma_z[\varkappa-\varkappa\circ\mathsf r_z](U).\]
In the first part a background of level $n_0$ is compared with its mirror image: in a H\"older gauge of exponent $\beta_{\rm H}$ about $z$ their difference is of the relative order $\vartheta^{\beta_{\rm H}}$, $\vartheta=L^{-(n_0-j)}$, and the family moves by $\varepsilon_{n_0}^2\vartheta^{4+\beta_{\rm H}}$ times its norm. In the second part the weight $\varkappa-\varkappa\circ\mathsf r_z$ is at most $L^{2K}$ times the Lipschitz constant of $\Xi$ times $L^{j-K}$, and its family at the background is bounded as above, its value at the trivial configuration being kept in the constant. Together, the births of the levels $j<n_0$, read at $n_0$ and less the values of their second parts at the trivial configuration, are at most a power of $K$ times $\mathsf H(\Xi)L^{3(K-n_0)-j-\beta_{\rm H}(n_0-j)}$, of the rate $\beta_{\rm H}-(3+\beta_{\rm H})\psi$, the second parts alone of the rate $1-3\psi$; and the birth of the level $n_0$ has the rate $1-4\psi$.
\emph{From $n_0$ to $t$.} The tower is a contraction in the sup norm, but every step feeds parts of the families that it carries into the families of the chain, each feeding at a level $k$ smaller by $L^{-(K-k)/2}$ up to logarithms. The fed parts are followed along chains of levels $k_1<\dots<k_m$, and each chain is stopped at the level $\lceil K-\frac13\tau Y\rceil$, $Y=\sum_l(K-k_l)$, where its family is replaced by its value at the trivial configuration and the error of a family at a background; the chains that live have at most a number of $\tau$ members, and their errors have the rates $(\frac12-\tau)\psi$ and $\frac12\psi-3\phi$.
\emph{The constant and the rate.} $\mathsf a(\Xi)$ is the sum of the values at the trivial configuration of the second parts of the births of the levels $j<n_0$, of the families of the stopped chains and of the birth of the level $n_0$; it is real linear in $\Xi$ and reads no configuration. The error is a power of $K$ times $\mathsf H(\Xi)L^{-c'K}$, $c'$ the least of the six rates $1-\frac92\psi$, $1-4\psi$, $\beta_{\rm H}-(3+\beta_{\rm H})\psi$, $1-3\psi$, $(\frac12-\tau)\psi$, $\frac12\psi-3\phi$. At $(\psi,\beta_{\rm H})=(\frac17,\frac34)$ they are $\frac5{14}$, $\frac37$, $\frac3{14}$, $\frac47$, $\frac1{14}-\frac\tau7$, $\frac1{14}-3\phi$, and the least is $\min\{\frac1{14}-\frac\tau7,\frac1{14}-3\phi\}$; at $(\frac15,\frac9{10})$ and $\phi<\frac1{30}$ they are $\frac1{10}$, $\frac15$, $\frac3{25}$, $\frac25$, $\frac1{10}-\frac\tau5$, $\frac1{10}-3\phi$, and the least is $\min\{\frac1{10}-\frac\tau5,\frac1{10}-3\phi\}$: the two numbers $c'$ of the statement. Theorems~\YMMref{AJ}{673} and~\YMMref{AD}{607} are not used.
\end{proof}

\begin{proposition}[Conservation of the tensor]\label{C-371-25}
\begin{enumerate}
\item[(a)] Let $0<g\le\gamma^{\rm sd}_W$, $h_0,R_0>0$ and $c<\frac1{10}$. There are $C$, $p$, $K_0$, depending on the constants, $g$, $h_0$, $R_0$ and $c$, such that for every first-passage datum $(K,\beta)$ at $g$ with $K\ge K_0$, $b=a_K\beta$, every $\rho\in\mathcal G^{\rm per}_{K,\beta}$ and every $\eta\in C^2(\mathbb R^3;\mathbb R^3)$ with support in $\{x_3\ge h_0,|x|_\infty\le R_0\}$:
\[0\le\rho\bigl(\theta_K\mathcal T^{\rm lat}_K(\eta);\mathcal T^{\rm lat}_K(\eta)\bigr)\le b^{-2}\bigl(CK^pL^{-cK}\|\eta\|_{C^2}\bigr)^2,\]
and $|\rho(\mathcal T^{\rm lat}_K(\eta);Y)|\le b^{-1}CK^pL^{-cK}\|\eta\|_{C^2}\,\rho(\theta_KY;Y)^{1/2}$ for every real bounded local function $Y$ of the bonds in $\{x_3<\frac12a_K\}$. Moreover, for every $r>0$ there are $C$, $p$, $K_0$, depending on the constants, $g$, $c$ and $r$, such that for every first-passage datum $(K,\beta)$ at $g$ with $K\ge K_0$, on the least admissible torus with the board of eight cells of side $\frac12\ell'$ and any offset in $a_K\mathbb Z^3$, and for every $C^2$ vector field $\eta$ on $T_{\ell'}$ carried by the points of one cell at sup distances $\ge r$ from its walls,
\[\bigl\|O^T_K(H^{\rm cons}[\eta])-m_\eta\bigr\|_{{\rm ch},\ell'}\le CK^pL^{-cK}\|\eta\|_{C^2},\qquad m_\eta\text{ the mean on }T_{\ell'} .\]
Here $\theta_K$ is the reflection in the plane $x_3=\frac12a_K$, $\rho(F;G)=\rho(FG)-\rho(F)\rho(G)$, and $\mathcal T^{\rm lat}_K(\eta)=O^T_K(H^{\rm cons}[\eta]+H''_K)$ is the conservation observable formed with the symmetric lattice differences $(\nabla^{\rm s}_\rho\varphi)(y)=(2a_K)^{-1}[\varphi(y+a_Ke_\rho)-\varphi(y-a_Ke_\rho)]$ and $(\nabla^{\rm c}_\rho\varphi)(y)=a_K^{-1}[\varphi(y+\frac12a_Ke_\rho)-\varphi(y-\frac12a_Ke_\rho)]$, $\nabla^{\rm s}\cdot\eta=\sum_\rho\nabla^{\rm s}_\rho\eta^\rho$: its plaquette weights are $(\nabla^{\rm s}\cdot\eta)(x(p))-2(\nabla^{\rm s}_{\lambda(p)}\eta^{\lambda(p)})(x(p))$ at the centres and its bond weights $-(\nabla^{\rm c}_{\lambda'}\eta^\lambda+\nabla^{\rm c}_\lambda\eta^{\lambda'})(x_b)$ at the midpoints, $\lambda<\lambda'$ the two directions other than that of the bond $b$, the lattice forms of the entries $\nabla\cdot\eta-2\partial_\lambda\eta^\lambda$ and $-(\partial_\lambda\eta^{\lambda'}+\partial_{\lambda'}\eta^\lambda)$ of $H^{\rm cons}[\eta]$; they differ from these entries by at most $\frac32a_K\|\nabla^2\eta\|_\infty$ at the plaquettes and $\frac12a_K\|\nabla^2\eta\|_\infty$ at the bonds, so that $|H''_K|_\infty\le\frac32a_K\|\nabla^2\eta\|_\infty$.
\item[(b)] Let $0<g\le\gamma$, $q\ge2$, $d>0$ and $R>0$. There are $C$, $p$, $K_0$, depending on the constants, $g$, $q$, $d$ and $R$, such that for every first-passage datum $(K,\beta)$ at $g$ with $K\ge K_0$, every $\nu\in\mathcal G^{\rm per}_{K,\beta}$, every $\eta\in C^2_c(\mathbb R^3;\mathbb R^3)$ and all bounded weights $H_2,\dots,H_q$ such that the $q$ supports of $\eta,H_2,\dots,H_q$ lie in a cube of side $2R$ at pairwise sup distances $\ge d$,
\[\Bigl|\nu\Bigl(A^\nu[O^T_K(H^{\rm cons}[\eta])]\prod_{a\ge2}A^\nu[O^T_K(H_a)]\Bigr)\Bigr|\le CK^pL^{-K/36}\|\eta\|_{C^2}\prod_{a\ge2}|H_a|_\infty,\qquad A^\nu[F]=F-\nu(F),\]
and the joint cumulant $\bigl\langle O^T_K(H^{\rm cons}[\eta]);O^T_K(H_2);\dots;O^T_K(H_q)\bigr\rangle$ under $\nu$ obeys the same bound. On a torus $T_\ell$ of an admissible side $\ell$ at $g$, the joint cumulant under its Wilson measure obeys the same bound for $K\ge K_0$, for every $C^2$ vector field $\eta$ and all bounded weights on $T_\ell$ whose supports have pairwise sup distances $\ge d$ on $T_\ell$, with $C$, $p$, $K_0$ depending on the constants, $g$, $q$, $d$ and $\ell$.
\item[(c)] Every tensor limit $\sigma$ at every $0<g\le\gamma$ has $\sigma(\theta H^{\rm cons}[\eta],H^{\rm cons}[\eta])=0$ for $\eta\in C^2_c(\{x_3>0\};\mathbb R^3)$, that is $[(H^{\rm cons}[\eta])]=0$ in $\mathcal H(\sigma)$, and $\sigma((H^{\rm cons}[\eta])\cdot X)=0$ for every $\eta\in C^2_c(\mathbb R^3;\mathbb R^3)$ and every tuple $X$ whose supports do not meet $\operatorname{supp}\eta$: the stress tensor $\Theta$ is conserved, $\partial_\mu\Theta_{\mu\nu}=0$, against every tuple.
\end{enumerate}
\end{proposition}
\begin{proof}
(a) Let $\mathfrak v_\eta$ be the derivation of the functions of the configuration along the left translations of the bond variables by the field of directions that $\eta$ induces on the bonds, and $A=\sum_pa_p$ the Wilson action. Computing the variation $\beta\mathfrak v_\eta A$ of the action and the divergence $\operatorname{div}\mathfrak v_\eta$, ordering the sums by the bonds and by the plaquettes, and using the parallelogram law $\frac14|\mathbf w_+-\mathbf w_-|^2=\frac12(|\mathbf w_+|^2+|\mathbf w_-|^2)-\frac14|\mathbf w_++\mathbf w_-|^2$ for the $\mathbf w$ of the plaquette variables of the two plaquettes of a plane at a bond, and $|\mathbf w(W)|^2=2a_q-a_q^2$ for the plaquette variable $W$ of a plaquette $q$, one obtains, identically in the configuration, as an identity between polynomials in the bond variables at finitely many bonds, the lattice Ward identity
\begin{equation}\label{C-371-34}
b\,\mathcal T^{\rm lat}_K(\eta)=\mathcal N_\eta-\tfrac14\mathfrak N[\mathsf k_\eta]+X''[\eta]+\mathfrak Y[\eta]-\tfrac14c_\eta .
\end{equation}
Here $\mathcal N_\eta=\beta\mathfrak v_\eta A-\operatorname{div}\mathfrak v_\eta$, and $\mathfrak N[\mathsf k_\eta]$ is formed in the same way with the field of the lattice derivative $\mathsf k_\eta$ of $\eta$; they are carried by the bonds at which $\eta$, resp. $\mathsf k_\eta$, does not vanish. They are null terms: the Haar measure of one bond is invariant under the left translations of its variable, so an integration by parts over the finitely many bonds at which the field does not vanish, the other bonds being fixed, shows that the conditional mean of $\mathcal N_\eta Y$, resp. $\mathfrak N[\mathsf k_\eta]Y$, given the other bonds, vanishes for every polynomial $Y$ in the bonds at which the field vanishes; with the equations of a Gibbs state, $\nu(\mathcal N_\eta Y)=\nu(\mathfrak N[\mathsf k_\eta]Y)=0$ for every Gibbs state $\nu$, and the same holds under the Wilson measure of every torus. $X''[\eta]$ is a sum of nine local entries, linear in $\eta$: the cubic entry $-S[b\eta]$ of Lemma~\ref{C-371-26}, of order three, with weights at most $L^{2K}b\|\eta\|_\infty$; five entries formed with the lattice derivative $\mathsf k_\eta$, $\mathsf k^\rho_\eta(y)=a_K^{-1}[\eta^\rho(y+a_Ke_\rho)-\eta^\rho(y)]$, of the orders two, two, two, three and four, the one of order three with weights at most $\frac14L^Kb\|\nabla\eta\|_\infty$; the entry $-\frac32\sum_pw_\eta(p)a_p$ of order two, $w_\eta(p)$ the sum of $(\nabla^{\rm c}_\rho\eta^\rho)(x(p))$ over the two directions $\rho$ of the plane of $p$; an entry of order two formed with the second lattice differences of $\eta$; and an entry of order four formed with the $a_p^2$. $\mathfrak Y[\eta]$ is a form in the squares of the in-plane differences of plaquette variables at the bonds, with the coefficients $\nabla^{\rm c}_\nu\eta^\nu$, and $c_\eta$ is a number.

The plaquette squares and the bond forms enter \eqref{C-371-34} with one normalisation, that of $O^T_K$: this is what makes the left member the lattice form of $H^{\rm cons}[\eta]$. With $\theta_K$, $\langle F,G\rangle_\rho=\rho(\theta_KF;G)$ is a positive semidefinite form on the real local functions of the bonds above the mirror, by the reflection positivity of $\rho$; the null terms are null in it, because $\theta_K$ of a function above the mirror is a polynomial in bonds below it, where $\eta$ vanishes. So $\|b\,\mathcal T^{\rm lat}_K(\eta)\|_\rho\le\|X''[\eta]\|_\rho+\|\mathfrak Y[\eta]\|_\rho$. Cut $\eta$ by a partition of unity into at most $(2R_0/\delta+3)^3$ pieces $\eta_i$ in cubes of side $\frac18\mathsf l'$, $\mathsf l'=\frac12\ell'$, $\delta=\frac1{16}\mathsf l'$; for each piece there is a board of the least torus with an offset in $a_K\mathbb Z^3$ in which the piece and its mirror image lie in two cells, at distance $\ge r=\min\{\frac14h_0,\frac1{32}\mathsf l'\}$ from the walls.

By the domination of $\rho$ (the chessboard estimate with the pressure gap) $\|F_1\|^2_\rho\le e^{2G_g\mathsf l'^3}\bar{\mathsf m}^2$, where $\bar{\mathsf m}^8$ bounds the mean on the least torus of the eight mirror copies of the centred piece $F_1$. For $F_1=X''[\eta_i]$ the copies are entries carried by different cells at distance $\ge r$; by Lemma~\ref{C-371-26}(a) they enter the progression as sources, a joint cumulant of entries that, integrated over the scales up to $t$, are constant up to $\epsilon_a$ on the small-field set is at most a number times $\prod\epsilon_a$, and the large fields are negligible at the depth $t$; the eight entries other than $-S[b\eta]$ are small by power counting, with the exponent $\frac12-4\phi$, which exceeds $\frac{11}{30}>\frac1{10}>c'$ as $\phi<\frac1{30}$, and $-S[b\eta]$ is small by the estimate on the cubic remainder, Lemma~\ref{C-371-26}(b), with the rate $c'=\min\{\frac1{10}-\frac\tau5,\frac1{10}-3\phi\}$ that the Lemma admits for $\phi<\frac1{30}$; so $\bar{\mathsf m}\le CK^pL^{-c'K}\|\eta_i\|_{C^2}$.

For $F_1=\mathfrak Y[\eta_i]$ a pointwise bound at the background, where the in-plane differences of the plaquette variables are H\"older small, gives the mean of the eight copies $\le(CK^{2\mathfrak p_0}L^{-c_{\rm Y}K}\|\nabla\eta_i\|_\infty)^8$, with
\[c_{\rm Y}=\min\bigl\{\tfrac12\psi-4\phi',\ \tfrac12-4\phi',\ 2\beta_{\rm H}-(3+2\beta_{\rm H})\psi\bigr\}\]
for any $0<\beta_{\rm H}<1$ and $0<\phi'<\psi<1$ with $\phi'\le\frac14$ and $c_{\rm Y}>0$, the number $C$ depending on them besides: $\beta_{\rm H}$ is the H\"older exponent of the backgrounds from a plaquette to its neighbour, the pointwise bound is taken at the level $K-\lceil\psi K\rceil$, and the progression is stopped at the level $K-\lceil\phi'K\rceil$. At $(\beta_{\rm H},\psi,\phi')=(\frac12,\frac16,\frac1{336})$ one has $c_{\rm Y}=\frac1{14}$, at $(\frac34,\frac14,\frac1{160})$ one has $c_{\rm Y}=\frac1{10}$, and $\sup c_{\rm Y}=\frac2{11}$, not reached. The seminorm is subadditive; summing the pieces gives the first display with the rate $\min\{c',c_{\rm Y}\}$. At the choice $c_{\rm Y}=\frac1{10}>c'$ this rate is $c'=\min\{\frac1{10}-\frac\tau5,\frac1{10}-3\phi\}$, and every $c<\frac1{10}$ is reached by the choice of $\tau$ and $\phi$.

The second display is the inequality of Schwarz for the form. For the chessboard bound: by \eqref{C-371-34}, $b\,\mathcal T^{\rm lat}_K(\eta)-\bar c=Z+R$, $\bar c$ its mean on $T_{\ell'}$, with $Z=\mathcal N_\eta-\frac14\mathfrak N[\mathsf k_\eta]$, of mean $0$, and $R=X''[\eta]+\mathfrak Y[\eta]-\bar c'$, $\bar c'$ the mean of $X''[\eta]+\mathfrak Y[\eta]$. In the mean on $T_{\ell'}$ of the product of the eight mirror copies of $Z+R$, every term with a factor $Z^{[\kappa]}$ vanishes: by the invariance of the Wilson measure under the reflections in the walls it is the mean of $Z$ times a polynomial in the bonds of the other cells, where $\eta$ and $\mathsf k_\eta$ vanish, and the null terms are null under the Wilson measure of the torus.

So $\|b\,\mathcal T^{\rm lat}_K(\eta)-\bar c\|_{{\rm ch},\ell'}=\|R\|_{{\rm ch},\ell'}\le\|X''[\eta]-c_1\|_{{\rm ch},\ell'}+\|\mathfrak Y[\eta]-c_2\|_{{\rm ch},\ell'}$, $c_1$, $c_2$ the means, the seminorm being subadditive; these are the two numbers $\bar{\mathsf m}$ bounded above, at the rates $c'$ and $c_{\rm Y}$, the copies being entries in different cells, at sup distances $\ge r$ from their walls; $b^{-1}\le g^2$; and $O^T_K(H^{\rm cons}[\eta])=\mathcal T^{\rm lat}_K(\eta)-O^T_K(H''_K)$ with $\|O^T_K(H''_K)-c_3\|_{{\rm ch},\ell'}\le\frac32C^\flat_1(g;\frac12\ell',r)a_K\|\nabla^2\eta\|_\infty$, the bound of one centred entry of Proposition~\ref{C-371-23}(c). The rate is again every $c<\frac1{10}$.
(b) Let $\Lambda=L^z$, $z\ge1$ an integer, and let $g'_1,g'_2$ be the couplings with $g'^{-2}_1=\Lambda g^{-2}/2$ and $g'^{-2}_2=L\Lambda g^{-2}/4$. As $g'^{-2}_2>g'^{-2}_1>g^{-2}$, they are reference couplings for every $0<g\le\gamma$, and they read $g$ alone, not the label. Put $x=g^{-2}$. The labels of the first-passage data at $g$ lie in the window $(x,L(x+\bar B)]$, which the dilation carries to $(\Lambda x,L\Lambda(x+\bar B)]$, and $I(g'_1)=(\frac12\Lambda x+\bar B,\frac12L\Lambda x+(L-2)\bar B]$, $I(g'_2)=(\frac14L\Lambda x+\bar B,\frac14L^2\Lambda x+(L-2)\bar B]$. Here $\bar B<\gamma^{-2}\le x$: the constants in Section~\YMMref{74}{614} satisfy $B_\sharp\le\gamma^{-2}/2$ and $\bar B=LB_\sharp/(L-1)<\gamma^{-2}$. So $\bar B\le\frac{L-4}4x$, as $L\ge13$, and the three comparisons of ends
\[\tfrac12\Lambda x+\bar B\le\Lambda x,\qquad\tfrac14L\Lambda x+\bar B\le\tfrac12L\Lambda x+(L-2)\bar B,\qquad L\Lambda(x+\bar B)\le\tfrac14L^2\Lambda x+(L-2)\bar B\]
show that $I(g'_1)\cup I(g'_2)$ holds the dilated window, at every $z\ge1$. By the windows of labels (Corollary~\YMMref{74.2}{616}) every first-passage datum $(K,\beta)$ at $g$ with $K\ge z+1$ thus dilates to a first-passage datum $(K-z,\beta)$ at $g'_1$ or at $g'_2$, of label $\Lambda a_K\beta$. Let $2\mathsf l_\iota$ be the least admissible side at $g'_\iota$, and let $z$ be the least integer $z\ge1$ with $\Lambda g^{-2}/2\ge(\gamma^{\rm sd}_W)^{-2}$, so that $g'_2<g'_1\le\gamma^{\rm sd}_W$, and with $\mathsf l_2\le\Lambda d/2$; such integers exist, the least side at $g'_\iota$ growing like a power of $\log g'^{-2}_\iota$, that is of $z$, and $\Lambda d$ like $L^z$. The dilation of Lemma~\ref{C-371-14}(a) carries entries, states and tori at $g$ to those at $g'_\iota$, and the supports to sets at pairwise sup distances $\ge\Lambda d\ge2\mathsf l_2$.

Use the board of cells of side $\mathsf l_\iota$ on the dilated lattice. Cutting each entry into small pieces and choosing the offset as in the proof of Theorem~\ref{cov-thm}, the selected pieces occupy distinct cells, each with margin $\mathsf l_\iota/(8q)$. Evaluate each cell's eight-copy seminorm separately on the least torus at $g'_\iota$. The cell of $\eta$ carries the chessboard bound of (a) on the least torus at $g'_\iota\le\gamma^{\rm sd}_W$, at the margin $\mathsf l_\iota/(8q)$: $\|O^T(H^{\rm cons}[\eta'])-\langle\cdot\rangle\|_{\rm ch}\le CK^pL^{-cK}\|\eta'\|_{C^2}$ for every $c<\frac1{10}$, in particular at $c=\frac1{36}$, $\eta'$ the dilated piece of $\eta$; every other cell carries the eighth root of the mean of the eight mirror copies of its centred entry, which the cumulant bounds of the orders $2,\dots,8$ on the least torus bound, as in the proof of Proposition~\ref{C-371-22}(a): the least side at $g'_\iota\le\gamma^{\rm sd}_W$ is small, and the entries, bond forms included, are sources (Lemma~\ref{C-371-51}(a), (c)).

The chessboard estimates, the pressure gap and the domination of the periodic Gibbs states and of the larger tori by the least torus bound the mean of the product of the centred entries by the product of these numbers; the joint cumulant is a polynomial in such means, each monomial with one factor that holds a piece of $\eta$. The powers of $\Lambda$ that the dilation brings to the entries and to $\|\eta\|_{C^2}$, $(H^{\rm cons}[\eta])^{(z)}=H^{\rm cons}[\eta^{[z]}]$ with $\|\eta^{[z]}\|_{C^2}\le\Lambda\|\eta\|_{C^2}$, and the factor $L^{z/36}$ of the shift of the cutoff enter $C$; on a torus $T_\ell$ the same argument runs on $T_{\Lambda\ell}$, with numbers that depend on $\ell$.
(c) Let first $g\le\gamma^{\rm sd}_W$ and $\eta\in C^2_c(\{x_3>0\};\mathbb R^3)$. By (a) in $\rho=\nu_j$, with $b_j^{-1}\le g^2$, $\|\mathcal T^{\rm lat}_{K_j}(\eta)\|_{\nu_j}\to0$; by Proposition~\ref{C-371-22}(a) the part $O^T_K(H''_K)$ contributes $O(a_K)$, and by Proposition~\ref{C-371-22}(c) with the moving mirror $\theta_{K_j}H\to\theta H$, $\nu_j(\theta_{K_j}O^T(H);O^T(H))\to\sigma(\theta H,H)$ for $H=H^{\rm cons}[\eta]$. So $\sigma(\theta H^{\rm cons}[\eta],H^{\rm cons}[\eta])=0$. The second assertion follows from (b) in the limit, or from the first by Proposition~\ref{C-371-28}(c). Above $\gamma^{\rm sd}_W$ use Lemma~\ref{C-371-14}(b) and $(H^{\rm cons}[\eta])^{(z)}=H^{\rm cons}[\eta^{[z]}]$.
\end{proof}

\begin{proposition}[Covariance at two points, for labelled functionals]\label{C-371-27}
Let $\sigma_2$ be real, symmetric and bilinear on the pairs of continuous weights with disjoint supports, bilinear on $C_c(O;V)\times C_c(O';V)$ for open $O,O'$ with disjoint closures and extended bilinearly to $V_{\mathbb C}$, with:
\begin{itemize}
\item[(h1)] $\sigma_2(\iota H,\iota H')=\sigma_2(H,H')$ for $\iota\in\mathbb R^3\rtimes W_3$;
\item[(h2)] for every compact $\mathcal K\subset\{x_3>0\}$ a number $c(\mathcal K)$ with $|\sigma_2(\theta\bar h,h')|\le c(\mathcal K)|h|_\infty|h'|_\infty$ for $h,h'\in C_c(\{x_3>0\};V_{\mathbb C})$ supported in $\mathcal K$;
\item[(h3)] a Hilbert space with a linear map $v$ from $C_c(\{x_3>0\};V_{\mathbb C})$, $\langle v(h),v(h')\rangle=\sigma_2(\theta\bar h,h')$, and a projection-valued measure $\mathcal P$ on $[0,\infty)\times\mathbb R^2$ carried by the square cone, with $\int e^{-a_3E+i(a_1q_1+a_2q_2)}d\mathcal P\,v(h)=v(\tau_ah)$ for $a_3\ge0$;
\item[(h4)] $\sigma_2(\theta_\Sigma\bar H,H)\ge0$ for $H\in C_c(\{x_1+x_3>0\};V_{\mathbb C})$, $\theta_\Sigma$ the reflection in $\{x_1+x_3=0\}$;
\item[(h5)] $\sigma_2(H^{\rm cons}[\eta],H')=0$ for $\eta\in C^2_c(\mathbb R^3;\mathbb R^3)$ and $H'$ with support disjoint from $\operatorname{supp}\eta$;
\item[(h6)] $\sigma_2((f\circ\mathsf R_{-t})\mathbf 1,(f'\circ\mathsf R_{-t})\mathbf 1)=\sigma_2(f\mathbf 1,f'\mathbf 1)$ for real $t$ and $f,f'$ with disjoint supports;
\item[(h7)] there are $a<4$, $C_{\rm od}$ and $1>d_n\downarrow0$ with $\sigma_2(\theta h,h)\le C_{\rm od}d_n^{-a}|h|_\infty^2$ for all $n$, $\mu\ne\nu$ and $h=h_0o_{\mu\nu}$, $h_0\in C_c(B(2d_ne_3,d_n))$ real.
\end{itemize}
Then the kernel $G(u)(e,e')$ of $\sigma_2$ is real analytic on $\mathbb R^3\setminus\{0\}$ with $G(Ru)(\rho(R)e,\rho(R)e')=G(u)(e,e')$, and $\sigma_2(RH,RH')=\sigma_2(H,H')$ for all $R\in\mathrm O(3)$ and all pairs with disjoint supports.
\end{proposition}
\begin{proof}
\emph{The measure of one entry.} Write $p=(E,q_1,q_2)$, $p_\pm=E\pm q_1$, $\pi(p)=(iq_1,iq_2,-E)$ and $\langle u,p\rangle=u_3E-i(u_1q_1+u_2q_2)$ for $u\in\mathbb C^3$, so that $\partial_{u_k}e^{-\langle u,p\rangle}=\pi_ke^{-\langle u,p\rangle}$. Modulating an entry, $h\mapsto he^{-i\kappa\cdot x}$, does not change its sup norm; from (h2), (h3) one obtains one family of complex Radon measures $M(e,e')$ on the square cone, antilinear in $e$, linear in $e'$, positive at $e=e'$, with $\langle v(ke),\mathcal P(B)v(k'e')\rangle=\int_B\bar F_kF_{k'}dM(e,e')$, $F_k(p)=\int k(x)e^{-\langle x,p\rangle}dx=\int k(x)e^{-x_3E+i(x_1q_1+x_2q_2)}dx$, and $\int e^{-tE}dM(e,e)<\infty$ for $t>0$. For the last, fix a continuous $g_1:\mathbb R\to[0,1]$ with support in $(-\frac12,\frac12)$ and $\max g_1=1$, put $c_1=\int g_1$, $\hat g_1(\kappa)=\int g_1(x)e^{ix\kappa}dx$ and $c_g=\int_{\mathbb R}|\hat g_1(\kappa)|^2e^{-|\kappa|}d\kappa$, and for $d>0$ and $\kappa\in\mathbb R^2$ let $h^d_\kappa(x)=g_1(x_1/d)g_1(x_2/d)g_1(x_3/d-2)e^{-i(\kappa_1x_1+\kappa_2x_2)}$ be the modulated bumps of size $d$: $|h^d_\kappa|\le1$, their supports lie in $\mathcal K_d=B(2de_3,d)$, and $F_{h^d_\kappa}(p)=\mathsf L(E)\,d\hat g_1(d(q_1-\kappa_1))\,d\hat g_1(d(q_2-\kappa_2))$ with $\mathsf L(E)=\int g_1(x_3/d-2)e^{-x_3E}dx_3\ge c_1d\,e^{-5dE/2}$.

If $\mu\ge0$ is carried by the square cone and $\int|F_{h^d_\kappa}|^2d\mu\le A$ for all $\kappa\in\mathbb R^2$, then integrating over the modulation against $e^{-d(|\kappa_1|+|\kappa_2|)}d\kappa$, with $|q_1|+|q_2|\le2E$ on the cone, gives $\int e^{-7dE}d\mu\le4A(c_1c_g)^{-2}d^{-6}$; at $\mu=M(e,e)$, (h2) gives this with $A=c(\mathcal K_d)|e|_\infty^2$. \emph{The kernel.} $G(u)(e,e')=\int e^{-\langle u,p\rangle}dM(e^\theta,e')$, $e^\theta=\rho(\Theta_{e_3})\bar e$, is holomorphic on the square tube and represents $\sigma_2$ on pairs parted by a plane $x_3=c$; by (h1) the kernels of the frames of $W_3$ agree on overlaps and give one function on the union of the rotated square tubes, whose real points are $\mathbb R^3\setminus\{0\}$; for a real $u$ with $(u_3,u_1)\ne0$, $\zeta\mapsto G(\mathsf R_\zeta u)$ is entire and $2\pi$-periodic, and an imaginary rotation $\mathsf R_{i\tau}$ is the boost $\beta_\tau(E,q_1,q_2)=(E\cosh\tau+q_1\sinh\tau,q_1\cosh\tau+E\sinh\tau,q_2)$ of $(E,q_1)$ acting on $M$: $\langle\mathsf R_{i\tau}u,p\rangle=\langle u,\beta_\tau p\rangle$ and $p_\pm\circ\beta_\tau=e^{\pm\tau}p_\pm$.

\emph{Conservation is algebra.} With $v\vee w=\frac12(vw^{\rm t}+wv^{\rm t})$ for $v,w\in\mathbb C^3$ one has $H^{\rm cons}[\eta]=(\nabla\cdot\eta)\mathbf 1-2\sum_{k,l}\partial_k\eta_l\,e_k\vee e_l$, and (h5) gives, as measures, $2\sum_k\pi_kM(e,e_k\vee e_l)=\pi_lM(e,\mathbf 1)$; in the helicity basis $E_{\pm2}=\varepsilon_\pm\vee\varepsilon_\pm$, $E_{\pm1}=e_2\vee\varepsilon_\pm$, $E_0=e_2e_2^{\rm t}$, $\mathbf 1$, $\varepsilon_\pm=e_3\pm ie_1$, whose helicities $h_s=2,1,0,0,-1,-2$ for $E_2,E_1,E_0,\mathbf 1,E_{-1},E_{-2}$ are fixed by $\rho(\mathsf R_t)E_s=e^{-ih_st}E_s$ (so $h_s$ is minus the helicity $k$ of a label of the plates below, which is fixed by $\rho(\mathsf R_t)\varepsilon_k=e^{ikt}\varepsilon_k$): $-p_-M(e,E_2)+2iq_2M(e,E_1)+p_+M(e,E_0)=0$, $-p_+M(e,E_{-2})+2iq_2M(e,E_{-1})+p_-M(e,E_0)=0$, $2iq_2M(e,E_0)-p_-M(e,E_1)-p_+M(e,E_{-1})=iq_2M(e,\mathbf 1)$.

\emph{The trace.} (h6) makes $M(\mathbf 1,\mathbf 1)$ invariant under the boosts. \emph{Growth.} By (h7), through the modulation bound with $A=4C_{\rm od}d_n^{-a}$ (the factor $4$ for the real and imaginary parts of $h^{d_n}_\kappa$), $\int e^{-7d_nE}dM(o,o)\le16C_{\rm od}(c_1c_g)^{-2}d_n^{-6-a}$ at the off-diagonal labels $o=o_{\mu\nu}$, and the boosted off-diagonal measures grow at most like $e^{(6+a)|\tau|}$ along $\tau_n=\log(t/7d_n)$; conservation passes this growth to the diagonal labels, except near $q_1\approx0$, where a boost does not lower the energy below its half. So only the harmonics $|\nu|\le6+a$ occur: $(\beta_\tau)_*M(e,e')=\sum_\nu e^{-\nu\tau}M_\nu(e,e')$, and by the quarter turn $M_\nu(E_s,E_{s'})=0$ unless $\nu\equiv h_s+h_{s'}$ modulo $4$, $h_s$ the helicities.

The extreme wrong harmonic $M^{(m)}$, $m\ne0$, is itself a positive matrix of measures, conserved, and it scales under the boosts. \emph{Odd and even.} If its index $m$ is odd, positivity at the diagonal mirror (h4) says that it is also negative: it vanishes. If $m$ is even, $|m|\ge2$, the entries at the helicities $\mp1$, $\mp2$ have the frequencies $\mp(2+4|m|)$, $\mp(4+4|m|)$, beyond $6+a<10$, so vanish, and with the trace three of the six rows are empty; conservation then empties the rest outside $\{p_\mp=0\}$, puts the helicity-$0$ entry at $p=0$, where the scaling under boosts kills it, and on $\{p_\mp=0\}$ a measure whose mass on a smaller set is $e^{|\nu|\tau}$ times its mass on the larger one is not finite on bounded sets unless it vanishes.

So all harmonics but $\nu=h_s+h_{s'}$ vanish, $G(\mathsf R_tu)(\rho(\mathsf R_t)e,\rho(\mathsf R_t)e')=G(u)(e,e')$, and the group generated by $W_3$ and the $\mathsf R_t$ is $\mathrm O(3)$.
\end{proof}

\begin{remark}\label{C-371-32}
The exponent enters through $10>6+a$; $a<4$ cannot be replaced by $a=5$. Let $S(p)$ be the symmetric matrix with zero diagonal and $S_{12}=p_3(p_1^2-p_2^2)$, $S_{23}=p_1(p_2^2-p_3^2)$, $S_{31}=p_2(p_3^2-p_1^2)$, $[p]_\times$ the matrix of $q\mapsto p\times q$, and $\Upsilon(p)=[p]_\times S(p)[p]_\times^{\rm t}$: it is symmetric, traceless, conserved ($p^{\rm t}\Upsilon(p)=0$), homogeneous of order five, $\Upsilon(wp)=\det w\,\rho(w)\Upsilon(p)$ for $w\in W_3$, and not covariant under the $\mathsf R_t$. The pair functional with the kernel $G(u)(e,e')=c_e(\partial)c_{e'}(\partial)(4\pi|u|)^{-1}$, $c_e(p)=-\frac12\operatorname{tr}(e\Upsilon(p))$, built on a free massless field and positive at the planes of the nine families by the positivity of that field, has (h1)--(h6) and (h7) with $a=5$, and it is not covariant.
\end{remark}

\emph{Plates and boosts.} For Proposition~\ref{C-371-28} write $y=(x_3,x_1)$ for the coordinates of the plane of the rotations $\mathsf R_t$ and $x=(y,x_2)$. With $u_\pm=e_3\pm ie_1$, the matrices $\varepsilon_{\mp2}=u_\pm u_\pm^{\rm t}$, $\varepsilon_{\mp1}=u_\pm e_2^{\rm t}+e_2u_\pm^{\rm t}$, $\varepsilon_0=e_3e_3^{\rm t}+e_1e_1^{\rm t}$, $\varepsilon_0'=e_2e_2^{\rm t}$ form a basis of $V_{\mathbb C}$ with $\rho(\mathsf R_t)\varepsilon_k=e^{ikt}\varepsilon_k$, $\rho(\mathsf R_t)\varepsilon'_0=\varepsilon'_0$; $k$ is the helicity of the label, the opposite of the helicity $h$ of the proof of Proposition~\ref{C-371-27}: the matrix $u_+u_+^{\rm t}=\varepsilon_{-2}$ is the $E_2$ of that proof, of helicity $h=2$.

A plate is a weight $\alpha(y)\varphi(x_2)\epsilon$ with $\alpha\in C_c(\mathbb R^2)$, $\varphi\in C^1_c(\mathbb R)$ real and $\epsilon$ in this basis; $\lambda=(\varphi,\epsilon)$ is its label, $k(\lambda)$ its helicity, $\lambda^\dagger=(\varphi,\rho(\Theta_{e_3})\bar\epsilon)$ its reflected label, and $\kappa(\lambda_1,\dots,\lambda_n)=\sum k(\lambda_a)$; for $\xi\in\mathbb C^2$, $\xi^*=(\bar\xi^0,-\bar\xi^1)$. In the coordinates $\xi^\pm=\xi^0\pm i\xi^1$ of $\mathbb C^2$ an imaginary rotation acts by $(\mathsf R_{is}\xi)^\pm=e^{\mp s}\xi^\pm$: it is a boost of the plane. For data with bounded slab operators the slab operator of a plate has a density in the point of the plane, the point plate $B_\lambda(\xi,\xi')$, holomorphic in norm for $\xi,\xi'$ in $\{\operatorname{Re}\xi^\pm>0\}$; the plate vectors $\Psi(x;\lambda_1,\dots,\lambda_n)$ are the products of point plates applied to $\Omega$, and the kernels $\mathcal K_n(Z;\lambda_1,\dots,\lambda_n)$ of the plates are their scalar products, continued to one holomorphic function on the union $\mathcal Z_n$ of the tubes $\{Z=(z_1,\dots,z_n):\operatorname{Re}(\mathsf R_{-\upsilon}(z_{\mathfrak s(c+1)}-z_{\mathfrak s(c)}))^\pm>0,\ 1\le c<n\}$ of the frames $\upsilon\in\frac\pi2\mathbb Z$ and the permutations $\mathfrak s$ of the points.

With $n_\alpha=(\cos\alpha,\sin\alpha)$, a real configuration $X=(y_1,\dots,y_n)$ lies in the tube of $(\mathfrak s,\upsilon)$ if and only if $n_\upsilon\cdot y_{\mathfrak s(1)}<\dots<n_\upsilon\cdot y_{\mathfrak s(n)}$; so the real points of $\mathcal Z_n$ are the configurations whose times, or whose abscissae, are pairwise different. A fan of $j$ points is a real $z$ with $z_1=r_0n_{\alpha_0}$, $z_{i+1}-z_i=r_in_{\alpha_i}$, $r_i>0$, $|\alpha_i|<\frac{5\pi}{16}$. $\mathcal H_{\rm pl}$ is the closed span of $\Omega$ and of the plate vectors.

\begin{proposition}[Conservation and a power bound give covariance]\label{C-371-28}
Let $\sigma$ be a real functional on the tuples of continuous weights with:
\begin{itemize}
\item[(C1)] $\sigma(\emptyset)=1$, $\sigma((H))=0$, $\sigma$ symmetric and $n$-linear on $\prod C_c(O_a;V)$ for open $O_a$ with disjoint closures, with sup-norm bounds as in (D1) that do not depend on the position; invariance under $\mathbb R^3\rtimes W_3$; positivity at the nine families; the square cone $\mathcal P(\{E<\max(|q_1|,|q_2|)\})=0$;
\item[(C2)] bounded slab operators of one entry: for $l>0$ and $h\in C_c(\{0<x_3<l\};V_{\mathbb C})$, $B_l(h)[F]=[(h)\cdot\tau_{le_3}F]$ is bounded by $C(\mathcal K,l)|h|_\infty$ for $\operatorname{supp}h\subset\mathcal K$ compact, with no condition on the growth of $C(\mathcal K,l)$;
\item[(C3)] the scale-free bound $\sigma(\theta h,h)\le C_{\rm od}|h|_\infty^2$ for $0<d<1$ and real $h=h_0o_{\mu\nu}$, $h_0\in C_c(B(x_0,d))$, $x_{0,3}\in[\frac32d,2d]$;
\item[(C4)] the bound $\sigma(\theta(f\mathbf 1),f\mathbf 1)\le C_\varphi t^{-a_\Gamma}$, $a_\Gamma<4$, of (S2) at the plates $f$ of scale $t\in(0,1]$;
\item[(C5)] conservation at two points: $\sigma(\theta H^{\rm cons}[\eta],H^{\rm cons}[\eta])=0$ for $\eta\in C^2_c(\{x_3>0\};\mathbb R^3)$;
\item[(C6)] for one $p<\infty$, $\|B_{2d}(\tau_{de_3}h)\|\le C\min\{1,d-d'\}^{-p}\|h\|_{L^1}$ for $0\le d'<d$ and $h\in C_c(\mathbb R^3;V_{\mathbb C})$ supported in $\{|x_3|\le d'\}$.
\end{itemize}
Then:
\begin{enumerate}
\item[(a)] $\sigma(RX)=\sigma(X)$ for all $R\in\mathrm O(3)$ and all tuples;
\item[(b)] $\mathcal K_n(\mathsf R_{is}Z;\underline\lambda)=e^{\kappa(\underline\lambda)s}\mathcal K_n(Z;\underline\lambda)$ for all $n\ge2$, all weight labels $\underline\lambda$ of plates and all real $s$, and $e^{-\kappa(\underline\lambda)s}\|\Psi(\mathsf R_{is}z;\underline\lambda)\|=\|\Psi(z;\underline\lambda)\|$ at real $z$;
\item[(c)] $\sigma((H^{\rm cons}[\eta])\cdot X)=0$ for every $\eta\in C^2_c(\mathbb R^3;\mathbb R^3)$ and every tuple $X$ whose supports do not meet $\operatorname{supp}\eta$.
\end{enumerate}
A growth with the sequence $A_\bullet=(A_j)_{j\ge1}$, $A_j\ge0$, is the bound $\|\Psi(\mathsf R_{is}z;\underline\lambda)\|\le Ce^{A_j|s|}$ for all real $s$, all labels $\underline\lambda$ of $j$ plates and all $z$ in a compact set of fans of $j$ points, $C$ depending on the labels and the compact set, for every $j\ge1$. The same holds if (C6) is replaced by a growth with $\liminf_j(A_j-j^2)/j=-\infty$, or, the diagonal mirror being used at all tuples, with $\liminf_jA_j/j^2<2$. If $\sigma$ has (C1), (C2) and is covariant at two entries under the rotations $\mathsf R_t$, the following are equivalent: covariance under $\mathrm O(3)$; the growth with $A_j=2j$; a growth with $\liminf_jA_j/j^2<1$; a bound $\sup_{s\in\mathbb R}\sup_{(\xi,\xi')\in\mathcal C}e^{-\gamma|s|}\|B_\lambda(\mathsf R_{is}\xi,\mathsf R_{is}\xi')|_{\mathcal H_{\rm pl}}\|<\infty$ of the boosted point plates on $\mathcal H_{\rm pl}$, with one $\gamma<\infty$, for every label $\lambda$ and every compact set $\mathcal C$ of pairs of real $\xi=rn_\alpha$, $\xi'=r'n_{\alpha'}$, $r,r'>0$, $|\alpha|,|\alpha'|<\frac{5\pi}{16}$; and (C6) with any $p<\infty$ implies them.

If $\sigma$ has (C1), (C2), the covariance at two entries under the $\mathsf R_t$ and a growth with some sequence $A_\bullet$, of any size, then the top frequency $\bar m(n)$ of the kernels of $n$ plates is finite, $\bar m(n)\le\frac14(A_{\lfloor n/2\rfloor}+A_{\lceil n/2\rceil}+2n)$, and if covariance fails, there is an $a'\ge1$ with $\bar m(2a)\ge\frac12a(a-a')$, so $A_a\ge a(a-a'-2)$, for all $a\ge a'$.
\end{proposition}
\begin{proof}
Conditions (C1)--(C3), (C5) and (C6) are (D1)--(D4), (S), (O), (N) and (P) of Definition~\ref{cov-def}. Theorem~\ref{cov-thm-data} gives (a)--(c). Under the alternative growth hypotheses, Proposition~\ref{cov-two}(a) gives covariance at two points and Proposition~\ref{cov-rot} gives it at all orders. Its proof includes the coefficient continuation, the modulus Schwarz inequality at tilted lines and the polygon-halving argument. Lemma~\ref{cov-boost} gives the exact norm formula for the plate vectors.

The equivalences: under covariance at two points about one axis $\bar m(2)=0$ as soon as a growth is had, the kernels of two points being then constant along the boosts; a growth with $\liminf_jA_j/j^2<1$ gives covariance by the convexity; covariance gives unitary boosts $U(s)$ on $\mathcal H_{\rm pl}$, $U(s)\Psi(w;\underline\lambda)=e^{-\kappa(\underline\lambda)s}\Psi(\mathsf R_{is}w;\underline\lambda)$, $U(s)B_\lambda(\xi,\xi')U(-s)=e^{-k(\lambda)s}B_\lambda(\mathsf R_{is}\xi,\mathsf R_{is}\xi')$ on $\mathcal H_{\rm pl}$, so the growth $A_j=2j$ and the bound of the boosted point plates with $\gamma=2$; that bound with any $\gamma$ gives, by the product of the point plates, the growth $A_j=\gamma j$, so covariance; and under a growth with any sequence, if covariance fails, the second case of the integer sequences holds with $\nu(1)=0$ and $\nu(a)\le\frac12A_a+a$.

\end{proof}

\begin{remark}\label{C-371-33}
Conservation cannot be omitted from Proposition~\ref{C-371-28}. For a tensor limit $\sigma$ at $g\le\gamma^{\rm sd}_W$ let $\sigma'$ be the limit at continuous weights of Lemma~\ref{C-371-14}(d). Then $\sigma'\circ Z_2$ has (C1)--(C4) and (C6): $Z_2$ commutes with $\rho(W_3)$, is the identity on $V_A\oplus V_E$, keeps the operators up to a unitary map, multiplies the bound of (C3) by $4$ and the $L^1$ norm in (C6) by at most $2$. It is not covariant (Proposition~\ref{C-371-17} with Proposition~\ref{C-371-24}(a) and Theorem~\ref{C-371-5}(3)). So it violates (C5): conservation, in which the plaquette squares and the bond forms enter with one normalisation, is the property of the theory that its rescaled forms $\sigma\circ Z_z$, $z\ne1$, do not have: for $z^2\ne1$ by the argument above, and at $z=-1$ because, with $R$ the rotation by $\pi/4$ about $e_3$, $\rho(R)o_{12}$ lies in $V_E$, where $Z_{-1}$ is the identity, so that covariance of $\sigma'\circ Z_{-1}$ would give $-\sigma=\sigma$ at the mixed pairs $(\varphi^d_1\mathbf 1,\varphi^d_2o_{12})$ of Proposition~\ref{C-371-17}, against the law of Proposition~\ref{C-371-24} with $\mathcal W\ne0$ there; and $\sigma'\circ Z_{-1}$ has (C1)--(C4) and (C6) as $\sigma'\circ Z_2$ does.
\end{remark}

\subsection{Reconstruction, rescaling and components}

\begin{proposition}[Wightman fields from labelled functionals]\label{C-371-16}
Let $\sigma$ be a real functional on $\mathfrak S_V(\mathbb R^3)$ with (D0), (D1) and $\sigma((H))=0$, and suppose:
\begin{itemize}
\item[(R1)] $\sigma\circ\iota=\sigma$ for every isometry $\iota$ of $\mathbb R^3$;
\item[(R2)] $\sigma$ is positive at $(e_3,0)$, and $\mathcal H(\sigma)$ is separable;
\item[(R3)] $\mathcal P(\{E<|q|\})=0$;
\item[(R4)] there are $C_V$ and $p$ with $|\sigma(\theta\bar F\cdot(h)\cdot F')|\le C_V\eta^{-p}\|h\|_{L^1}\|[\tau_{-\eta e_3}F]\|\|[\tau_{-\eta e_3}F']\|$ for $\eta\in(0,1]$, $F,F'\in\mathcal D_{e_3,\eta}$ and $h\in C^1_c(\{|x_3|<\eta/2\};V)$.
\end{itemize}
Then:
\begin{enumerate}
\item[(i)] for $n\ge2$ there is exactly one continuous function $s_n$ on $\{x\in(\mathbb R^3)^n:x_i\ne x_j\ (i\ne j)\}\times V^n$, $n$-linear in $e$, with $\sigma(H_1,\dots,H_n)=\int s_n(x;H_1(x_1),\dots,H_n(x_n))\,dx$ for all tuples, and for the tuples of continuous weights if $\sigma$ is given there; it is real analytic, symmetric, invariant under the isometries acting on points and labels, and
\[|s_n(x;e)|\le(cC_V)^n\max\{1,2n^2/\delta(x)\}^{pn}\prod_k|e_k|_\infty,\qquad\delta(x)=\min_{i\ne j}|x_i-x_j|,\]
$c$ a number depending on $(V,\rho)$ and its norm; so $|\sigma(X)|\le(cC_V)^n(2n^2/d)^{pn}\prod_a\|H_a\|_{L^1}$ if the supports are at mutual distances $\ge d$, $0<d\le1$. For $\varphi:\mathbb R^{3n}\to V_{\mathbb C}^{\otimes n}$ whose $6^n$ components are Schwartz functions vanishing with all their derivatives where two of the $n$ points coincide, $S_n(\varphi)=\int s_n\varphi$ converges absolutely and
\[|S_n(\varphi)|\le(c'C_V)^n(n!)^{2p}|\varphi|_{k_En},\qquad k_E=\max\{\lceil p\rceil,4\},\quad c'=6\pi^22^pe^{2p}3^{k_E}c,\]
$|\varphi|_k$ the greatest of the norms $\sup(1+|x|^2)^{k/2}|\partial^\alpha\varphi_\iota|$, $|\alpha|\le k$, over the components; $S_n(H_1\otimes\dots\otimes H_n)=\sigma(H_1,\dots,H_n)$ at smooth entries with compact disjoint supports, and $S_n$ is the only linear functional bounded for one of the norms with these values. The truncated functions have the kernels that the formula of the cumulants writes in the $s_m$, $m\le n$.
\item[(ii)] $\Omega$ and the classes of the tuples whose entries lie in pairwise disjoint slabs $\{t<x_3<t'\}$ are total in $\mathcal H(\sigma)$.
\item[(iii)] On $\mathcal D_\infty=\bigcap_k\operatorname{dom}H^k$ there are operators $T(f)$, $f\in\mathcal S(\mathbb R^{1+2};V_{\mathbb C})$, linear in $f$, with $\langle\psi,T(f)\psi'\rangle=\int\lim_{s\downarrow0}\langle U(x)^{-1}\psi,B(s,s;f(x))U(x)^{-1}\psi'\rangle\,dx$, $B$ the point operators that (R4) provides, and $\langle\psi,T(f)\psi'\rangle=\langle T(\rho(\Theta_{e_3})\bar f)\psi,\psi'\rangle$, so that a component with $r$ indices equal to $3$ is Hermitian after the factor $i^r$; and there is a strongly continuous unitary representation $U$ of the Poincar\'e group of $\mathbb R^{1+2}$ extending $x\mapsto U(x)$, with $U\Omega=\Omega$, the spectrum of the translations in the closed forward cone, and $U(a,\Lambda)T(f)U(a,\Lambda)^{-1}=T(\rho_{\mathbb C}(\hat\Lambda)f(\Lambda^{-1}(\cdot-a)))$ on the span $\mathcal D_0$ of $\Omega$ and of the vectors $T(f_1)\cdots T(f_n)\Omega$.

Here $\hat\Lambda=\iota\Lambda\iota^{-1}\in\mathrm{SO}(3,\mathbb C)$ for $\Lambda\in\mathrm{SO}_0(1,2)$, with $\iota(t,\vec x)=(\vec x,-it)$: a time $t$ sits at the Euclidean height $x_3=-it$, as $U(t,\vec x)=\mathsf E((\vec x,-it))$; the boost of $(t,x_1)$ by $s$ goes to $\mathsf R_{is}$, and a rotation $1\oplus R$, $R\in\mathrm{SO}(2)$ acting on $\vec x$, to $R\oplus1$; and $\rho_{\mathbb C}(\hat\Lambda)E=\hat\Lambda E\hat\Lambda^{\rm t}$, so that $\rho_{\mathbb C}(\mathsf R_{is})=\exp(is\,d\rho(E_{13}-E_{31}))$, $d\rho(\omega)E=\omega E+E\omega^{\rm t}$.
\item[(iv)] $T(f)T(f')=T(f')T(f)$ on $\mathcal D_0$ when the supports of $f$ and $f'$ are spacelike separated, all six components with all six; $\mathcal D_0$ is dense; the Wightman functions $\langle\Omega,T(f_1)\cdots T(f_n)\Omega\rangle$ are tempered and equal $\lim_{\epsilon\downarrow0}\int\mathbf W_n(\iota(x_2-x_1+i\epsilon\alpha_1),\dots;f_1(x_1),\dots)\,dx$ for all $\alpha_k$ in the forward cone, with the $\iota$ of (iii).
\item[(v)] The kernels $\mathbf W_n$ of the product formula are holomorphic on $\mathsf{Tb}^{n-1}$, with $|\mathbf W_n(\zeta;e)|\le(C_V\eta^{-p})^n\prod_k|e_k|_\infty$ for $e\in V^n$ where $\operatorname{Re}\zeta^k_3-|(\operatorname{Im}\zeta^k_1,\operatorname{Im}\zeta^k_2)|\ge2\eta$ for all $k$, $0<\eta\le1$; $\hat\Lambda$ of (iii) maps $\mathsf{Tb}=\iota(\mathbb R^{1+2}+iV_+)$, $V_+=\{\alpha:\alpha^0>|\vec\alpha|\}$, onto itself, and $\mathbf W_n(\hat\Lambda\zeta^1,\dots,\hat\Lambda\zeta^{n-1};\rho_{\mathbb C}(\hat\Lambda)e_1,\dots,\rho_{\mathbb C}(\hat\Lambda)e_n)=\mathbf W_n(\zeta;e)$ for every $\Lambda\in\mathrm{SO}_0(1,2)$. The $s_n$ continue to one holomorphic function $\mathbf S_n$ on the union of the tubes $\{z\in\mathbb C^{3n}:R^{-1}(z_{\varsigma(k+1)}-z_{\varsigma(k)})\in\mathsf{Tb},\ 1\le k<n\}$ of the real frames $R\in\mathrm{SO}(3)$ and the permutations $\varsigma$, on the tube of $(R,\varsigma)$ equal to $\mathbf W_n(R^{-1}(z_{\varsigma(2)}-z_{\varsigma(1)}),\dots,R^{-1}(z_{\varsigma(n)}-z_{\varsigma(n-1)});\rho(R)^{-1}e_{\varsigma(1)},\dots,\rho(R)^{-1}e_{\varsigma(n)})$; $\mathbf S_n$ is symmetric when points and labels are permuted together, $\mathbf S_n(Rz_1+a,\dots,Rz_n+a;\rho(R)e)=\mathbf S_n(z;e)$ for the real isometries $x\mapsto Rx+a$, $R\in\mathrm O(3)$, and $|\mathbf S_n(z;e)|\le(cC_V\eta^{-p})^n\prod_k|e_k|_\infty$ for $e\in V^n$, $c$ the number of (i), where for some $(R,\varsigma)$ the differences $R^{-1}(z_{\varsigma(k+1)}-z_{\varsigma(k)})$ have the margin $2\eta$ of the first bound.
\end{enumerate}
If besides $\ker H=\mathbb C\Omega$, every vector invariant under the translations is a multiple of $\Omega$, and all the Wightman axioms hold.
\end{proposition}
\begin{proof}
\emph{Point operators and kernels.} By (R4) the form $h\mapsto\sigma(\theta\bar F\cdot(h)\cdot F')$ is bounded in $\|h\|_{L^1}$; with the covariance under translations this gives point operators with $\|B(s,s;e)\|\le C_Vs^{-p}|e|_\infty$ and the product formula of Proposition~\ref{C-371-23}(f), whose kernels are holomorphic on $\mathsf{Tb}^{n-1}$ by (R3), $\mathsf E(\zeta)$ being a contraction on $\overline{\mathsf{Tb}}$; the first bound of (v) is $\prod\|B(\eta,\eta;e_k)\|$. By (R1) the same holds in every frame $R\in\mathrm{SO}(3)$: there the heights are taken along $Re_3$, the kernel is $\mathbf W_n(R^{-1}\cdot\,;\rho(R)^{-1}e)$, the labels being turned back by $\rho(R)^{-1}$, and the bound holds with $\prod_k|\rho(R)^{-1}e_k|_\infty\le c^n\prod_k|e_k|_\infty$; $c=3$ serves, as each row of $R$ has $\sum_l|R_{kl}|\le\sqrt3$. The kernels of two frames, and of two orders of the points, agree at the common real points of their tubes, as the proof of (i) shows; the intersection of two tubes is convex and holds the real parts of its points, so they agree on it by the identity theorem, and they form the one function $\mathbf S_n$ of (v), symmetric as $\sigma$ is. Its covariance under the real isometries holds at the real points by (i), and on each tube, which an isometry carries onto a tube, by the identity theorem.
(i) For a unit vector $v$ the map $u\mapsto u\cdot v$ carries the surface measure of $S^2$ to the uniform law on $[-1,1]$, so the band $\{|u\cdot v|<n^{-2}\}$ has the fraction $n^{-2}$; for pairwise different $x_1,\dots,x_n$ the $u$ with $|u\cdot(x_i-x_j)|\ge n^{-2}|x_i-x_j|$ for all $i\ne j$ fill at least the fraction $1-(n-1)/(2n)>\frac12$ of the sphere. In a frame whose third axis is such a $u$ the points are ordered in height with gaps at least $n^{-2}\delta(x)$, and the kernel of that frame is real analytic near $x$ with the bound, at $\eta=\min\{1,\delta(x)/(2n^2)\}$. For two frames that both order $x$, both kernels are continuous on a product of small balls about the $x_i$ and have the same integrals against all products $f_1(x_1)\cdots f_n(x_n)$; so they are equal.

The same reason gives uniqueness, symmetry and covariance. For a general tuple, cover the product of the supports by finitely many such products of balls, cut each entry by a partition of unity, and sum by the $n$-linearity on neighbourhoods with disjoint closures. For $S_n$: at a point $x$ with $\delta(x)<1$, $\varphi$ is flat at the point where the two nearest points coincide, and Taylor's formula of order $m=\lceil pn\rceil$ along the segment gives $\max\{1,2n^2/\delta(x)\}^{pn}|\varphi(x)|\le(2n^2)^{pn}3^{k}|\varphi|_k(1+|x|^2)^{-k/2}$, $k=k_En\ge3n+1$; $\int(1+|x|^2)^{-k/2}dx\le\pi^{2n}$ and $n^n\le e^nn!$.
(ii) The vectors $\Phi(y)$ of $n$ points of the half space have the Gram kernel $s_{2n}$ at $(\theta y,y')$, real analytic by (i); so for $\psi$ orthogonal to the classes of (ii), $y\mapsto\langle\psi,\Phi(y)\rangle$ is real analytic on the connected set of the non-coinciding $n$-tuples of the half space, and vanishes where the heights are ordered.
(iii), (iv), (v) The fields are built from the point operators and $U(x)$, and the Wightman functions are the boundary values of the $\mathbf W_n$ from every direction $\alpha$ of the open cone $V_+$, tempered by the power $p$; this follows the reconstruction of Theorem~\textup{\YMMref{AO}{778}}, with labels; (ii) gives cyclicity, and $B(s,s';e)^*=B(s',s;\rho(\Theta_{e_3})e)$ gives the adjoints. \emph{The invariance of the $\mathbf W_n$ in (v).} $\hat\Lambda\iota(x+i\alpha)=\iota(\Lambda x+i\Lambda\alpha)$, and $\Lambda V_+=V_+$: so $\hat\Lambda\mathsf{Tb}=\mathsf{Tb}$. For a rotation $\Lambda=1\oplus R$ of $\vec x$, $\hat\Lambda=R\oplus1$ is a real rotation that keeps the heights; at real $\zeta\in\mathsf{Tb}^{n-1}$, that is at configurations ordered in height, $\mathbf W_n=s_n$, and the covariance of $s_n$ gives the identity, which the identity theorem carries to $\mathsf{Tb}^{n-1}$.

For the boost $B(s)$ of $(t,x_1)$, $\hat\Lambda=\mathsf R_{is}$: fix a real $\zeta\in\mathsf{Tb}^{n-1}$ and $e$; there is $\delta>0$ such that $\mathsf R_\tau\zeta^k\in\mathsf{Tb}$ for all $k$ and all complex $\tau$ at a distance $<\delta$ from the segment $\{is':|s'|\le|s|\}$, a convex open set on which $\gamma(\tau)=\mathbf W_n(\mathsf R_\tau\zeta^1,\dots,\mathsf R_\tau\zeta^{n-1};\rho_{\mathbb C}(\mathsf R_\tau)e_1,\dots,\rho_{\mathbb C}(\mathsf R_\tau)e_n)$ is holomorphic; at real $\tau$, $|\tau|<\delta$, the turned points are still ordered in height, and the covariance of $s_n$ under the real rotation $\mathsf R_\tau$ gives $\gamma(\tau)=\gamma(0)$; so $\gamma$ is constant, and at $\tau=is$ this is the identity at real $\zeta$, which the identity theorem carries to $\mathsf{Tb}^{n-1}$.

The boosts $B(s)$ and the rotations $1\oplus R$ generate $\mathrm{SO}_0(1,2)$, and $\Lambda\mapsto\hat\Lambda$ is a homomorphism: the identity holds for every $\Lambda$. \emph{(iii).} By the substitution $x_k=\Lambda y_k$ and $\iota(\Lambda y+i\epsilon\alpha)=\hat\Lambda\iota(y+i\epsilon\Lambda^{-1}\alpha)$, with $\Lambda^{-1}\alpha\in V_+$, the invariance of the $\mathbf W_n$ makes the Wightman functions invariant under $f_k\mapsto\Lambda_af_k=\rho_{\mathbb C}(\hat\Lambda)f_k(\Lambda^{-1}(\cdot-a))$, the translations being had already; as $\bar\iota=\Theta_{e_3}\iota$ and $\Lambda$ is real, $\rho_{\mathbb C}(\hat\Lambda)$ commutes with $e\mapsto\rho(\Theta_{e_3})\bar e$, so the adjoints of (iii) go over too, and $U(a,\Lambda)T(f_1)\cdots T(f_n)\Omega:=T(\Lambda_af_1)\cdots T(\Lambda_af_n)\Omega$ is well defined and isometric on $\mathcal D_0$, which is dense; it extends to a unitary representation, strongly continuous as the fields are continuous in $f$, which fixes $\Omega$ and has the law of (iii) on $\mathcal D_0$.

\emph{(iv).} As $\mathcal D_0$ is dense it suffices that a Wightman function does not change when two neighbours $f_k$, $f_{k+1}$ with spacelike separated supports are exchanged, first for $f_j$ of compact support, to which cut-offs that do not enlarge the supports reduce the general case. Let $x$ run through the product of the supports. At each $x$ let $\hat\xi$ be the unit vector of $\vec x_{k+1}-\vec x_k$, and tilt $e_3$ towards it: $u_\vartheta=\cos\vartheta\,e_3+\sin\vartheta(\hat\xi,0)$, $\sin\vartheta=A\eta$, a real unit vector. The points $z_j=\iota(x_j)+\eta h^\lambda_je_3$, $|\lambda|\le1$, $0<\eta\le\eta_0$, with heights $h^\lambda_j$ formed from the supports alone, increasing with $j$ at $\lambda=1$ and in the order with $k,k+1$ exchanged at $\lambda=-1$, lie in the tube of the real frame with third axis $u_\vartheta$, its differences having the margin $3\eta$: across the pair $k,k+1$ the spacelike separation, $\hat\xi\cdot(\vec x_{k+1}-\vec x_k)>|t_{k+1}-t_k|$, pays for the exchanged heights, and $|\operatorname{Im}\mathfrak z-(u_\vartheta\cdot\operatorname{Im}\mathfrak z)u_\vartheta|=|t_{j+1}-t_j|\sin\vartheta$ for $\mathfrak z=z_{j+1}-z_j$.

There $\mathbf S_n$, symmetric in the points with their labels, is bounded by $(cC_V\eta^{-p})^n\prod_k|e_k|_\infty$ by (v). The integral $I(\lambda,\eta)$ of $\mathbf S_n(z;f_1(x_1),\dots,f_n(x_n))$ over $x$ has $\partial_\eta I$ and $\partial_\lambda I/\eta$ equal to integrals of the same kind with derivatives of the $f_j$ in the times, after an integration by parts. The one-label smoothing and integration-by-parts argument in the proof of Proposition~\ref{C-355-25}, applied to each tensor component, gives Taylor's formula with the integral remainder and bounds all of them uniformly in $\eta\in(0,\eta_0]$. So $|I(1,\eta)-I(-1,\eta)|\le C\eta\to0$; and $I(1,\eta)$, $I(-1,\eta)$ tend to the Wightman functions of the two orders, the labels being permuted with the points. \emph{The last sentence.} A vector $\psi$ invariant under the translations has $e^{itH}\psi=U(t,\vec 0)\psi=\psi$ for every real $t$, by (iii), so $H\psi=0$ by Stone's theorem; here $\ker H=\mathbb C\Omega$ enters and gives $\psi\in\mathbb C\Omega$.

With this uniqueness of the vacuum, (iii), (iv) and the separability of $\mathcal H(\sigma)$ in (R2) give all the Wightman axioms.
\end{proof}

\begin{proposition}[What is had without rotations]\label{C-371-29}
\begin{enumerate}
\item[(a)] Let $\sigma$ be a real symmetric functional on tuples of continuous, or of $C^1$, weights, multilinear as in (D0), with the sup-norm bounds of (D1), invariant under $\mathbb R^3\rtimes W_3$, positive at $(e_3,0)$, with the square cone. Then $\Omega$ and the classes of the tuples whose entries lie in pairwise disjoint slabs $\{t<x_3<t'\}$ are total in $\mathcal H(\sigma)$; more exactly, every $[X]$ is a limit of sums of classes of ordered tuples with as many entries, formed from mollified and cut entries of $X$. Two such functionals that agree at all the tuples ordered in the third coordinate are equal.
\item[(b)] Let $\sigma$ be a tensor limit at $0<g\le\gamma^{\rm sd}_W$ and $C_V$, $p$ as in Proposition~\ref{C-371-15}. Without covariance under rotations: there are operators $T(f)$, $f\in\mathcal S(\mathbb R^{1+2};V_{\mathbb C})$, on $\mathcal D_\infty$, tempered, with $\langle\psi,T(f)\psi'\rangle=\langle T(\rho(\Theta_{e_3})\bar f)\psi,\psi'\rangle$, $U(a)T(f)U(a)^{-1}=T(f(\cdot-a))$, $[H,T(f)]=iT(\partial_tf)$, and covariant under the dihedral group of the eight $w\in W_3$ with $we_3=e_3$; the $\mathbf W_n$ are holomorphic on $(\mathsf{Tb}^\square)^{n-1}$ with $|\mathbf W_n(\zeta;e)|\le(C_V\eta^{-p})^n\prod_k|e_k|_\infty$ for $e\in V^n$ where $\operatorname{Re}\zeta^k_3-|\operatorname{Im}\zeta^k_1|-|\operatorname{Im}\zeta^k_2|\ge2\eta$, $0<\eta\le1$; the Wightman functions are their tempered boundary values; $\sigma$ has one kernel at the configurations whose heights differ in one coordinate direction; and $\Omega$ is cyclic for the $T(f)$.
\item[(c)] Let $\sigma$ be a tensor limit at $0<g\le\gamma^{\rm ref}_W$, let $n\ge2$, let $i\ge i_n$ be an integer, where $i_n=1$ if $g\le\gamma^{\rm sd}_W$ and $i_n=z+1$ if $\gamma^{\rm sd}_W<g\le\gamma^{\rm ref}_W$, $z$ the dilation of Lemma~\ref{C-371-14}(c), chosen with $g'_1,g'_2\le\gamma^{\rm sd}_W$, which depends on the constants and $g$ and not on $n$, and let $\mathsf l=L^{-i}$. Then for every tuple $X=(H_1,\dots,H_n)$ of weights at pairwise sup distances $\ge2\mathsf l$, ordered in some direction or not,
\[|\sigma(X)|\le\prod_{a=1}^n2\,\mathsf g_{n,i}\Bigl(\frac{4n}{\mathsf l}\Bigr)^3\|H_a\|_\natural,\qquad\mathsf g_{n,i}=\mathsf c_g\,n^4\,\mathsf Y_i^{q_E}L^{(11/2)i},\]
where $\|H\|_\natural=\sum_{\mathfrak b}|\mathfrak b|\sup_{\mathfrak b}|H|_\infty$, the sum over the boxes $\mathfrak b=\frac{\mathsf l}{4n}(m+[0,1)^3)$, $m\in\mathbb Z^3$, is the upper Riemann sum of $\int|H|_\infty$ at the mesh $\mathsf l/(4n)$, $\mathsf Y_i=2(c_E+i\log L)$ and $q_E$ are the numbers of Proposition~\ref{C-371-15} at the margin $\delta=1/(8n)$, of which $c_E=\log(2L(g^{-2}+\bar B))$ and $q_E$, the latter formed from the constants alone, do not depend on the margin (above $\gamma^{\rm sd}_W$ they are taken at the couplings $g'_1,g'_2$ of Lemma~\ref{C-371-14}(c) and at the index $i-z$, the larger of the two), and $\mathsf c_g$ depends on the constants and $g$. No rotation, no ordering of the supports and no analyticity enters. There is a number $\mathsf C_G$, depending on the constants and $g$, such that for all $n\ge2$, all $0<d\le1$ and all tuples of $n$ weights at pairwise sup distances $\ge2d$, ordered or not,
\[|\sigma(H_1,\dots,H_n)|\le(\mathsf C_Gn^{16}d^{-9})^n\prod_a\|H_a\|_\natural\le(e^{16}\mathsf C_G)^n(n!)^{16}d^{-9n}\prod_a\|H_a\|_\natural,\]
the sums $\|\cdot\|_\natural$ taken at the mesh $L^{-i}/(4n)$, $i$ the least integer with $L^{-i}\le d$ and $i\ge i_n$; and $\sigma$ satisfies the growth hypothesis in Proposition~\YMMref{93.5}{719}, with labels, with the constants $(17,1,\mathsf C_E,24)$, $\mathsf C_E$ depending on the constants and $g$: for every $n$ there is exactly one complex linear functional $\mathsf S^\sigma_n$ on the test functions $\varphi:\mathbb R^{3n}\to V_{\mathbb C}^{\otimes n}$ of Proposition~\ref{C-371-16}(i), whose $6^n$ components are Schwartz functions vanishing with all their derivatives where two of the points coincide, that is bounded for one of the norms $|\cdot|_k$ and has the value $\sigma(H_1,\dots,H_n)$ at $H_1\otimes\dots\otimes H_n$ for smooth entries with compact, pairwise disjoint supports; and $|\mathsf S^\sigma_n(\varphi)|\le\mathsf C_E^n(n!)^{24}|\varphi|_{17n}$, so that each component is a distribution of order at most $17n$ on the configurations of $n$ different points. The $L^1$ form of these bounds in several entries, $|\sigma(X)|\le C\prod_a\|H_a\|_{L^1}$ near the configurations whose heights no coordinate direction orders, is the bound of Theorem~\ref{C-371-5}(6), which uses the rotations.
\end{enumerate}
\end{proposition}
\begin{proof}
(a) Use the ordered-tuple argument of Lemma~\ref{cov-ordered}. Its proof uses only the sup-norm bounds, translation and cubic symmetry, coordinate reflection positivity and the square cone, which are the hypotheses here. It gives the stated approximations by mollified and cut members, using the two-member height-$L^1$ bound to remove thin layers and analytic continuation to lift one group past another. The same approximants work for two functionals, giving the last assertion; the classes determine $\sigma(X)=\langle\Omega,[X]\rangle$.
(b) The proof of Proposition~\ref{C-371-16} runs with the square tube $\mathsf{Tb}^\square$ in the place of $\mathsf{Tb}$, the square cone of Proposition~\ref{C-371-23}(a) in the place of (R3), the dihedral group in the place of $\mathrm O(3)$ and Proposition~\ref{C-371-15}(iii) for (R4); cyclicity comes from (a). With the forward cone of Theorem~\ref{C-371-5}(2) the boundary values are taken from every direction of the forward cone.
(c) Let $\delta_n=1/(8n)$, and let first $g\le\gamma^{\rm sd}_W$, so that $i_n=1$. In a periodic Gibbs state, and on an admissible torus, cut every entry into its pieces on the boxes of side $\mathsf l/(4n)$. For one choice of a piece of each entry the $n$ pieces lie at pairwise sup distances $>\mathsf l$, and an offset in $a_K\mathbb Z^3$ of a board of cells of side $\mathsf l$ puts them into $n$ different cells, each at the margin $\delta_n\mathsf l$ from the walls (in each coordinate the offsets that bring a wall near one of the $n$ pieces fill less than a period). Proposition~\ref{C-371-15} at the margin $\delta_n$ holds for every $i\ge1$ and at all weights. Of the margin $\delta$ its proof reads only the depth $s$ at which the renormalisation progression is stopped on the least torus of the scale $L^{-i}$, chosen so that the mirror copies of a piece, about $\delta$ apart in units of the cell, are a number of blocks of the stopped level apart that is formed from the constants, $g$ and $i$ alone; so $L^s$ is $\delta^{-1}$ times such a number, and the moment bound pays $L^{4s}$: this gives the factor $\delta^{-4}$ in $C_E$, which is $\delta^{-4}$ times a number free of $\delta$, and the cutoffs $K_E(i)$, while $c_E$ and $q_E$ are formed without $\delta$.

At the off-diagonal labels it reads the bond form of a weight $\mathsf w$ of the bonds as a source of the small-field progression with sources (Propositions~\ref{prop:source-entry} and~\ref{prop:weighted-finite-system}, with the local representation specified here), and of the term of the never-large histories in the decomposition of a mean by the level at which a large field is first created (Section~\YMMref{123}{828}, equation~\eqref{eq:ward-first-creation}), under the majorant $(59\sup|\mathsf w|,S[\mathsf w])$, $S[\mathsf w]$ the plaquettes that contain a bond where $\mathsf w\ne0$ (Lemma~\ref{C-371-51}(a), (b)): the bond form has the properties of a source that these statements read, so that they hold for it as for a weight of the plaquettes of the sup norm $59\sup|\mathsf w|$ carried by $S[\mathsf w]$.

Proposition~\ref{C-371-15} at the margin $\delta_n$ and its consequence (i), with $b^{-1}\le g^2$, bound the mean of the product of the pieces, centred in the state, by $\prod_a2\mathsf g_{n,i}\sup|H_a|_\infty$, the suprema over their boxes, with no factorial. By multilinearity the mean of the product of the centred entries is the sum of these terms over the choices of one box for each entry, and $(4n/\mathsf l)^3\|H\|_\natural=\sum_{\mathfrak b}\sup_{\mathfrak b}|H|_\infty$; the bound passes to the limit. Above $\gamma^{\rm sd}_W$ the dilation of Lemma~\ref{C-371-14} carries it over, a board of the index $i\ge z+1$ going to one of the index $i-z\ge1$ at $g'_1$ or $g'_2$, its factors entering $\mathsf c_g$.

Nothing in this uses a rotation, an ordering of the supports or analyticity. For the count, take $i$ the least integer with $L^{-i}\le d$ and $i\ge i_n$; then $L^i\le L^{i_n}/d$, and $L^{i_n}$ depends on the constants and $g$ alone; $2\mathsf g_{n,i}(4n/\mathsf l)^3=128\,\mathsf c_gn^7\mathsf Y_i^{q_E}L^{(17/2)i}$, and, as $t^q\le(4q)^qe^{t/4}$ for $t,q>0$ and $c_E$, $q_E$ do not read $n$, $\mathsf Y_i^{q_E}\le(4q_E)^{q_E}e^{c_E/2}L^{i/2}$ (above $\gamma^{\rm sd}_W$ the same with the numbers at $g'_s$ and the index $i-z$): one power $d^{-1/2}$ absorbs it, and each factor is at most $Cn^7d^{-9}\le\mathsf C_Gn^{16}d^{-9}$; and $n^n\le e^nn!$. The growth condition follows from this bound by the extension of multilinear forms bounded at separated tuples (Proposition~\YMMref{93.3}{717}) with the exponent $9n$ of $d^{-1}$, applied to each of the $6^n$ forms $(f_1,\dots,f_n)\mapsto\sigma(f_1E_{\iota_1},\dots,f_nE_{\iota_n})$, $E_\iota$ running through the basis $E_{\lambda\lambda}$, $o_{\lambda\lambda'}$ $(\lambda<\lambda')$ of $V$: each has exactly one extension to the scalar test functions flat at the coinciding points that is bounded for one of the norms; $\mathsf S^\sigma_n(\varphi)$ is the sum of these extensions at the components $\varphi_\iota$ of $\varphi$ in that basis; it has the value $\sigma(H_1,\dots,H_n)$ at $H_1\otimes\dots\otimes H_n$ by multilinearity; and for a second functional $\mathsf S'$ with these properties, $\psi\mapsto\mathsf S'(\psi E_{\iota_1}\otimes\dots\otimes E_{\iota_n})$ is an extension of the same form, bounded for one of the norms, so it is that extension, and $\mathsf S'=\mathsf S^\sigma_n$.

For weights in balls of radius $1$, $\|H\|_\natural\le27|H|_\infty$, and $(n!)^8(n!)^{16}=(n!)^{24}$, $|\varphi|_{17n-1}\le|\varphi|_{17n}$.
\end{proof}

\begin{proposition}[Rescaling the off-diagonal labels]\label{C-371-17}
Let $\sigma$ satisfy (D0)--(D4), $z\ne0$ real, and $\sigma^z=\sigma\circ Z_z$. Then $\sigma^z$ satisfies (D0)--(D4), with $\max\{1,|z|\}^nC_n(\delta,\varrho_0)$ in (D1), and has the trace part of $\sigma$; $[F]_{\sigma^z}\mapsto[Z_zF]_\sigma$ is a unitary map of $\mathcal H(\sigma^z)$ onto $\mathcal H(\sigma)$ that commutes with the translations and the semigroup, under which the slab and point operators and the kernels of $\sigma^z$ are those of $\sigma$ with $Z_z$ applied to the labels. If $\sigma(\varphi_1E,\varphi_2E)\ne0$ for $E=o_{12}$ and some $\varphi_1,\varphi_2$ with disjoint supports, and $z^2\ne1$, then $\sigma$ and $\sigma^z$ are not both covariant at order $2$. For a tensor limit at $g\le\gamma^{\rm sd}_W$ the two-point law of Proposition~\ref{C-371-24}, at all pairs of weights, holds for $\sigma^z$ with some label $b'>0$ in the place of $b$ if and only if $z=1$; at the pairs of off-diagonal weights alone it holds with the label $b$ of $\sigma$ if and only if $z^2=1$.
\end{proposition}
\begin{proof}
$Z_z$ is the identity on $V_A$ and commutes with $\rho(W_3)$, in particular with $\rho(\Theta_n)$ for $n\in\mathfrak N_{18}$; this gives the first part. Let $R$ be the rotation by $\pi/4$ about $e_3$, so that $\rho(R)E=E_{22}-E_{11}\in V_E$. If both were covariant at order $2$, then
\[z^2\sigma(\varphi_1E,\varphi_2E)=\sigma^z(\varphi_1E,\varphi_2E)=\sigma^z(R(\varphi_1E),R(\varphi_2E))=\sigma(R(\varphi_1E),R(\varphi_2E))=\sigma(\varphi_1E,\varphi_2E).\]
For the law: at $z=1$, $\sigma^z=\sigma$. Let $\sigma^z$ obey the law with an error $\epsilon'(d)\to0$, at all pairs with a label $b'$, or at the pairs of off-diagonal weights alone with the label $b$. In the first case, at the pairs of diagonal weights, which $Z_z$ leaves alone, its label is that of $\sigma$: there $\sigma^z=\sigma$, and at $(\varphi^d_1\mathbf 1,\varphi^d_2\mathbf 1)$, with the bumps of Proposition~\ref{C-371-24}(a), the kernel of $\mathcal W$ is $\frac32\operatorname{tr}K_3(u)^2=9(4\pi|u|^3)^{-2}>0$, $u=x-y$, so that $\mathcal W(\varphi^d_1\mathbf 1,\varphi^d_2\mathbf 1)$ is a positive number independent of $d$, and $b'^{-2}$ and $b^{-2}$ are both the limit of $\sigma(\varphi^d_1\mathbf 1,\varphi^d_2\mathbf 1)/\mathcal W(\varphi^d_1\mathbf 1,\varphi^d_2\mathbf 1)$ as $d\downarrow0$: $b'=b$.

With the label $b$, in both cases, at off-diagonal entries $\sigma^z(H_1,H_2)=z^2\sigma(H_1,H_2)$, and $|z^2-1|\,b^{-2}|\mathcal W(H_1,H_2)|\le(z^2\epsilon_g(d)+\epsilon'(d))b^{-2}|H_1|_\infty|H_2|_\infty$ for all small $d$ forces $z^2=1$ by the lower bound of $\mathcal W$ (Proposition~\ref{C-371-24}(a)); conversely at $z=\pm1$ these pairs are unchanged. At the mixed pair $H_1=\varphi^d_1\mathbf 1$, $H_2=\varphi^d_2o_{12}$, with the bumps of Proposition~\ref{C-371-24}(a) placed with $x_2-x_1=\frac72d(e_1+e_2)/\sqrt2$, one has $\sigma^z(H_1,H_2)=z\,\sigma(H_1,H_2)$, and the kernel of $\mathcal W$ is $\frac32\operatorname{tr}[K_3(u)o_{12}K_3(u)]=3(K_3(u)^2)_{12}=9\hat u_1\hat u_2(4\pi|u|^3)^{-2}$, with $K_3(u)=(3\hat u\hat u^{\rm t}-\mathbf 1)(4\pi|u|^3)^{-1}$, $\hat u=u/|u|$, $u$ the difference of the two points; on the supports $\hat u$ lies within the angle $\arcsin\frac27<\frac\pi4$ of $\pm(e_1+e_2)/\sqrt2$, so $\hat u_1\hat u_2>0$, and $\mathcal W(H_1,H_2)=w_\times>0$ does not depend on $d$, $\mathcal W$ being invariant under dilations.

So $|z-1|\,b^{-2}w_\times\le(|z|\epsilon_g(d)+\epsilon'(d))b^{-2}|\varphi|_\infty^2$ for all small $d$, and $z=1$. This is Proposition~\ref{C-371-24}(d) at $a=e=1$.
\end{proof}

\begin{proposition}[The decomposition with labels]\label{C-371-30}
Let $\sigma$ be a tensor limit at $0<g\le\gamma^{\rm ref}_W$ of the Gibbs sequence $(K_j,\beta_j,\nu_j)$, $\sigma_0$ its trace part and $(\mathsf X_0,\lambda_0,(\sigma_{0,\zeta}))$ the canonical decomposition of $\sigma_0$.
\begin{enumerate}
\item[(a)] $\sigma$ has a canonical decomposition $(\mathsf X,\lambda,(\sigma_\chi)_{\chi\in\mathsf X})$: $\mathsf X$ compact metrisable, $\lambda$ a Borel probability measure with support $\mathsf X$, each $\sigma_\chi$ a real symmetric functional on labelled tuples, $\chi\mapsto\sigma_\chi(X)$ continuous and $\chi\mapsto\sigma_\chi$ injective, and the iterated plateaux of $\sigma$ at $(X^1,\dots,X^q)$ equal to $\int\prod_i\sigma_\chi(X^i)\,\lambda(d\chi)$. For every $\chi$ there are translation-invariant events $E_j$ with $\nu_j(E_j)>0$ such that $\rho_j=(\nu_j)_{1_{E_j}}\in\mathcal G^{\rm per}_{K_j,\beta_j}$ have $\sigma_\chi(X)=\lim_j\rho_j(A_j(X))$ along the full sequence, $A_j(X)$ the product of the entries centred in $\nu_j$; with $\mathcal I$ the invariant sigma-algebra, $\|\mathbb E_{\rho_j}[A_j(X)\mid\mathcal I]-\sigma_\chi(X)\|_{L^\infty(\rho_j)}\to0$. The bounds of Proposition~\ref{C-371-22}(a) hold for every $\sigma_\chi$ with numbers free of $\chi$.
\item[(b)] Every $\sigma_\chi$ clusters, $\sigma_\chi(X\cdot\tau_vY)\to\sigma_\chi(X)\sigma_\chi(Y)$ as $|v|\to\infty$, and $\ker H=\mathbb C\Omega$ on $\mathcal H(\sigma_\chi)$.
\item[(c)] Exactly one map $\pi:\mathsf X\to\mathsf X_0$ has $\sigma_\chi(f_1\mathbf 1,\dots,f_n\mathbf 1)=\sigma_{0,\pi(\chi)}(f_1,\dots,f_n)$; it is continuous and onto, and $\pi_*\lambda=\lambda_0$. The $\rho_j$ of (a) are approximants of $\pi(\chi)$, $\sigma_\chi$ is the lift $\sigma^\sharp_{\pi(\chi)}$ formed with them, along the full sequence, and $\sigma_\chi((H))=\frac{c_{\pi(\chi)}}3\int\operatorname{tr}H$.
\item[(d)] With $s_\chi(H)=\sigma_\chi((H))$, $(K_j,\beta_j,\rho_j)$ is a Gibbs sequence at $g$ with the tensor limit $\sigma^c_\chi(X)=\sum_J\prod_{a\notin J}(-s_\chi(H_a))\sigma_\chi(X_J)$, along the full sequence. The linear map $\mathsf S_\chi(X)=\sum_J\prod_{a\notin J}s_\chi(H_a)X_J$ of combinations is a bijection, multiplicative, commutes with the bar and with every isometry, and $\sigma_\chi=\sigma^c_\chi\circ\mathsf S_\chi$; $[F]\mapsto[\mathsf S_\chi F]$ is a unitary map of $\mathcal H(\sigma_\chi)$ onto $\mathcal H(\sigma^c_\chi)$ that carries $\Omega$ to $\Omega$ and intertwines $e^{-tH}$, the translations and the dihedral group; under it $B_l(X)$ corresponds to $\sum_J\prod_{a\notin J}s_\chi(h_a)B^c_l(X_J)$ and $B(s,s';e)$ to $B^c(s,s';e)+\frac{c_{\pi(\chi)}}3(\operatorname{tr}e)e^{-(s+s')H}$. In particular $\ker H=\mathbb C\Omega$ on $\mathcal H(\sigma^c_\chi)$.
\end{enumerate}
\end{proposition}
\begin{proof}
(a) The construction of the canonical decomposition in the proof of Theorem~\textup{\YMMref{AL}{746}} uses of the scalar functional only bounds of the kind of Proposition~\ref{C-371-22}(a), uniform along the sequence and in every periodic Gibbs state; with them it applies to labelled tuples: at a fixed physical size the conditional expectations of the centred products onto the translation-invariant events have plateaux that do not depend on the cutoff in the limit, iterated plateaux commute with $j\to\infty$, and the plateau functions take their values in a countable product of compact intervals; Stone--Weierstrass and Riesz give $\mathsf X$ and $\lambda$, and conditioning on invariant events of small probability near a point $\chi$ gives the $\rho_j$. (b) The plateau of a component is the product: $\sigma_\chi(X\cdot\tau_vY)\to\sigma_\chi(X)\sigma_\chi(Y)$.

So $\langle[F'],e^{-tH}[F]\rangle=\sigma_\chi(\theta\bar F'\cdot\tau_{te_3}F)\to\langle[F'],\Omega\rangle\langle\Omega,[F]\rangle$ as $t\to\infty$, and the projection onto $\ker H$ is $|\Omega\rangle\langle\Omega|$. (c) The trace parts of the components decompose the trace part, and the canonical decomposition of $\sigma_0$ is unique; the conditioned states are approximants of the scalar component $\pi(\chi)$, and Proposition~\ref{C-371-22}(d) applies to them, along the full sequence because the limit $\sigma_\chi$ exists. (d) The expansion of Theorem~\ref{C-371-5}(1), inverted; $\mathsf S_\chi$ has the inverse with $-s_\chi$; the operators follow from the identity of values.
\end{proof}

\subsection{Proof of Theorem~\ref{C-371-4}}

\emph{Step 1: clause (a).} Lemma~\ref{C-371-6}(a) gives all but the assertion that the action is forced, which is Lemma~\ref{C-371-6}(d); for $A=\rho(\Theta_u)$ and $h$ with values in $V_u^-$ it gives $[(h)]_u=[(-h)]_u=-[(h)]_u$.

\emph{Step 2: clause (b).} That a positive mirror, or one with real numbers $\sigma(\theta_u\bar F\cdot F)$, is a symmetry is Lemma~\ref{C-371-6}(b); positivity at $(n,c)$ from $m$ and $R$ is Lemma~\ref{C-371-6}(c); closedness is Lemma~\ref{C-371-6}(e), and $W_3\subset G_\sigma$ is (D3). For the equivalences: (iii)$\Rightarrow$(ii) because the plane may be moved through $0$ (Lemma~\ref{C-371-6}(a)), $\Theta_n\in G_\sigma$ by Lemma~\ref{C-371-6}(b), which uses only the reality of the numbers, and $\Theta_n\notin W_3$; (ii)$\Rightarrow$(i) by Lemma~\ref{C-371-6}(e); (i)$\Rightarrow$(iv) by Lemma~\ref{C-371-6}(c) with $m=e_3$, positive by (D4), and $Re_3=n$:
\[\sigma(\theta_{n,c}\bar F\cdot F)=\sigma(\theta_{e_3}\bar G\cdot G)\ge0,\qquad G=\iota^{-1}F,\ \iota x=Rx+cn.\]
Since $\mathfrak N_{18}\subset\mathfrak M(\sigma)$ by (D4), $\mathfrak M(\sigma)\ne\mathfrak N_{18}$ gives (iii), hence (iv).

\emph{Step 3: clause (c).} The first assertions are Lemma~\ref{C-371-20}(a)--(c). The chain: covariance at order $2m$ gives $\mathfrak M^{(m)}(\sigma)=S^2$ by Lemma~\ref{C-371-20}(b), (c) at $(\mathfrak K_m,\mathfrak K_{2m})$; the second arrow is $\mathfrak N_{18}\ne S^2$; for the third, a $u\in\mathfrak M^{(m)}(\sigma)\setminus\mathfrak N_{18}$ gives, by Lemma~\ref{C-371-20}(d) at $\mathfrak F=\mathfrak K_m$, with all real $c$ by Lemma~\ref{C-371-6}(a), whose proof keeps the number of entries, $\Theta_u\in G_\sigma[\mathfrak K_m]$, $\Theta_u\notin W_3$, and Lemma~\ref{C-371-20}(a). At $m=\infty$ the orders $m$ and $2m$ agree. The assertions on one and two entries are Lemma~\ref{C-371-20}(e).

\emph{Step 4: clause (d).} Proposition~\ref{C-371-11}(a) gives the symmetries, the positivity and the cone. On the forward cone
\[\bigl|e^{-\zeta_3E+i(\zeta_1q_1+\zeta_2q_2)}\bigr|\le e^{-E(\operatorname{Re}\zeta_3-|(\operatorname{Im}\zeta_1,\operatorname{Im}\zeta_2)|)},\]
so $\|\mathsf E(\zeta)\|\le1$ on $\overline{\mathsf{Tb}}$, and the bounds $E^ke^{-Ed}\le(k/ed)^k$ on the carrier give holomorphy in norm; the product law and the real points are the calculus of $\mathcal P_u$. With $e_3=u$, $e_1=-u'$, $e_2=u''$ one has $n_-=(u+u')/\sqrt2$, and the three mirrors are those of Proposition~\ref{C-371-13}; the exchange of $e_1$ and $e_2$ lies in the cubic group of the frame. Proposition~\ref{C-371-13}(b), (c), (d) give $D$, the $V_\Lambda$, the invariance of the null sets and of the support, and the unitary map to the picture of $(u+u')/\sqrt2$; a closed subset of the forward cone invariant under $\mathrm{SO}_0(1,2)$ is a union of orbits, and the orbits are $\{0\}$, the light cone without its tip and the mass shells. The cubic group of the frame, through Lemma~\ref{C-371-7}, connects the pictures of the six normals $\pm u,\pm u',\pm u''$ and of the twelve others among themselves. Lemma~\ref{C-371-19} at $\gamma=\pi/4$, $n=(u+u')/\sqrt2$, gives the law; its hypothesis is Lemma~\ref{C-371-9} at $(u,u')$ and at $(n,m)$, as $(n+m)/\sqrt2=u$. The proofs use of $\rho$ only that it is a homomorphism and, in Lemma~\ref{C-371-7}, a bound $|\rho(\Theta)E|\le c_\rho|E|$. The last assertion is Proposition~\ref{C-371-11}(b).

\emph{Step 5: clause (e).} By Lemma~\ref{C-371-7} a rotation $D_t$ about $e_3$ in $G_\sigma$ acts unitarily on $\mathcal H(\sigma)$ and turns $q$ in $\mathcal P$, so $\{E<|q\cdot D_te_1|\}$ is a null set together with $\{E<|q_1|\}$; for a countable dense set $A_0$ of angles
\[\{E<|q|\}=\textstyle\bigcup_{t\in A_0}\{E<|q\cdot D_te_1|\},\]
since $E<|q|$ gives $q\cdot m>E$ for all unit $m$ near $q/|q|$. Under covariance, Lemma~\ref{C-371-6}(c) gives $\mathfrak M(\sigma)=S^2$, Lemma~\ref{C-371-9} at $(e_3,e_1)$ gives $H\ge|P_1|$, and Lemma~\ref{C-371-7} transports the picture; strong continuity of the representation: $\|[RF]-[F]\|^2=2\|[F]\|^2-2\operatorname{Re}\sigma(\theta\bar F\cdot RF)\to0$ as $R\to1$, by (D0), (D1) as in the proof of Lemma~\ref{C-371-6}(e), and the classes are dense.

\emph{Step 6: clause (f).} Proposition~\ref{C-371-21}.

\subsection{Proof of Theorem~\ref{C-371-5}}

\emph{Step 1: clause (1).} Existence, (D0)--(D4) with $C_n=C^{\rm lc}_n$, the trace part and the lifts are Proposition~\ref{C-371-22}(b)--(d); the label exists and is determined by $\sigma$ by Proposition~\ref{C-371-24}(a). For a lift, $\rho_j(O_a)-\nu_j(O_a)\to\frac{c_\xi}3\int\operatorname{tr}H_a$, $O_a=O^T_{K_j}(H_a)$; expanding
\[\prod_a\bigl(O_a-\nu_j(O_a)\bigr)=\sum_S\prod_{a\notin S}\bigl(\rho_j(O_a)-\nu_j(O_a)\bigr)\prod_{a\in S}\bigl(O_a-\rho_j(O_a)\bigr)\]
and inverting the sum shows that $\sigma^c_\xi$ exists along the same subsequence and gives both identities; $(K_j,\beta_j,\rho_j)$ is a Gibbs sequence, $\rho_j$ being a periodic Gibbs state.

\emph{Step 2: clause (2).} By (1), (D4) gives positivity at the three mirrors of each frame $(e_\mu,e_\lambda,e_\nu)$, and Theorem~\ref{C-371-4}(d) applies; the $18$ normals of the cubic group of the standard frame are $\mathfrak N_{18}$.

\emph{Step 3: clause (3).} Let first $g\le\gamma^{\rm sd}_W$, $\sigma$ a tensor limit at $g$, and $\sigma'$ the limit at continuous weights of Lemma~\ref{C-371-14}(d), which agrees with $\sigma$ at the tuples of weights. $\sigma'$ has the hypotheses of Proposition~\ref{C-371-28}: (C1) by Propositions~\ref{C-371-22}(c) and~\ref{C-371-23}(a); (C2) by Proposition~\ref{C-371-23}(d); (C3) by Proposition~\ref{C-371-24}(e); (C4) because its trace part is the scalar limit, which has the bound of (S2); (C5) by Proposition~\ref{C-371-25}(c); (C6) by Proposition~\ref{C-371-15}(iii). At the off-diagonal labels (C3), (C6) and (h7) below rest on the bond forms entering the renormalisation progression as sources under the majorant $(59\sup|\mathsf w|,S[\mathsf w])$ (Lemma~\ref{C-371-51}(a), (b)), as in the proof of Proposition~\ref{C-371-29}(c); for (C3) and (h7) their second order at small fields enters besides, a quadratic form in the plaquette angles of $S[\mathsf w]$ with a kernel whose row sums lie under the same majorant, and this gives the two-point bound of Proposition~\ref{C-371-24}(e) at these labels with a number free of $d$ (Lemma~\ref{C-371-51}(c)).

So $\sigma'$ is covariant, and so is $\sigma$. At two entries Proposition~\ref{C-371-27} applies instead: (h1), (h2) by Proposition~\ref{C-371-22}(c), (h3) by Proposition~\ref{C-371-23}(a) with $v(h)=[(h)]$, (h4) by (D4), (h5) by Proposition~\ref{C-371-25}(c), (h6) by the invariance of the scalar limit (S2), (h7) with $a=0$ by Proposition~\ref{C-371-24}(e); no bound with a power enters. For $\gamma^{\rm sd}_W<g\le\gamma$, pass to a subsequence along which the bounded labels converge, choose $z$ with $g_z\le\gamma^{\rm sd}_W$ and use Lemma~\ref{C-371-14}(b):
\[\sigma(RX)=L^{nz}\sigma_z((RX)^{(z)})=L^{nz}\sigma_z(RX^{(z)})=L^{nz}\sigma_z(X^{(z)})=\sigma(X).\]
For a lift, the identity of (1), the covariance of the tensor limit $\sigma^c_\xi$ and $\int\operatorname{tr}(RH)=\int\operatorname{tr}H$ give the claim.

\emph{Step 4: clause (4).} For $g\le\gamma^{\rm ref}_W$, clause (3) is (i) of Theorem~\ref{C-371-4}(b) for tensor limits and lifts, which with (D4) gives (iv); Theorem~\ref{C-371-4}(e) gives the pictures and the representation of $\mathrm O(2)$. For a tensor limit at $\gamma^{\rm ref}_W<g\le\gamma$, with $z$ as in Step 3, $(\theta_{n,c}H)^{(z)}=\theta_{n,L^zc}H^{(z)}$ and Lemma~\ref{C-371-14}(b) make $\sigma(\theta_{n,c}\bar F\cdot F)$ the form of $\sigma_z$ at $(n,L^zc)$ at the combination $S_zF$, and $\sigma_z$ is positive there; this gives $\mathfrak M(\sigma)=S^2$ at every $0<g\le\gamma$. For the labels, at $0<g\le\gamma^{\rm ref}_W$, where the two-point law of Proposition~\ref{C-371-24} holds (above $\gamma^{\rm sd}_W$ through the dilation in its proof), Proposition~\ref{C-371-24}(b) at $\theta'=\theta$ and $x_0=\frac32de_3$ gives $\|[(\varphi E)]\|^2=\sigma(\theta(\varphi E),\varphi E)\ge(w_0-\epsilon_g(d))b^{-2}\operatorname{tr}E^2$; for $E\in V^-_{e_3}$, $\theta(\varphi E)=-(\varphi\circ\theta)E$, so the form of the reflection acting on the points alone takes the negative value $-\|[(\varphi E)]\|^2$; for $u=Re_3$ the covariance carries $\varphi E$ to $R(\varphi E)$, with values in $V^-_u$, and the value is the same; the translations give every $c$.

\emph{Step 5: clause (5).} (R1) is clause (3) with (D2); (R2) is Proposition~\ref{C-371-23}(a); (R3) is clause (2); (R4) is Proposition~\ref{C-371-15}(iii). For the kernels, take $\varepsilon=\eta$ in the product formula of Proposition~\ref{C-371-23}(f): then $\zeta^k-2\eta e_3\in\overline{\mathsf{Tb}}$, the $\mathsf E(\zeta^k-2\eta e_3)$ are contractions, holomorphic in norm by clause (2), and (R4), with $h$ concentrating at a point of $\{x_3=0\}$, gives $\|B(\eta,\eta;e)\|\le C_V\eta^{-p}|e|_\infty$; hence
\[|\mathbf W_n(\zeta;e)|\le\prod_k\|B(\eta,\eta;e_k)\|\le(C_V\eta^{-p})^n\prod_k|e_k|_\infty.\]
For a lift the holomorphy is the same argument with its point operators.

\emph{Step 6: clause (6).} By clause (1) and clause (5), Proposition~\ref{C-371-16} applies; at $p=9$ one has $k_E=9$ and $(n!)^{2p}=(n!)^{18}$. At continuous weights the kernel of the limit $\sigma'$ of Lemma~\ref{C-371-14}(d) serves. Above $\gamma^{\rm sd}_W$, $s_n(x;e)=L^{4nz}s'_n(L^zx;e)$ with the kernel $s'_n$ at the smaller coupling and $\delta(L^zx)\ge\delta(x)$. For a lift the kernel is $\sum_S\prod_{a\notin S}\frac{c_\xi}3\operatorname{tr}e_a\,s^c_{|S|}(x_S;e_S)$, $s^c_0=1$, $s^c_1=0$, by the identity of clause (1); $|\frac{c_\xi}3\operatorname{tr}e|\le|c_\xi||e|_\infty$, $\delta(x_S)\ge\delta(x)$, and there are $2^n$ subsets.

\emph{Step 7: clause (7).} Proposition~\ref{C-371-30}(a)--(c) give the decomposition, the map $\pi$, the lifts and the clustering, and $\ker H=\mathbb C\Omega$ on $\mathcal H(\sigma_\chi)$; Proposition~\ref{C-371-30}(d) gives the same on $\mathcal H(\sigma^c_\chi)$; $\sigma^c_\chi$ is a tensor limit at $g$, so clauses (5), (6) apply to it, and with $\ker H=\mathbb C\Omega$ the last sentence of Proposition~\ref{C-371-16} gives all the Wightman axioms.

\emph{Step 8: clause (8).} By clause (3) $\sigma$ is covariant. By Proposition~\ref{C-371-24}(a), with the bumps $\varphi^d_i$ there, $b^2|\sigma(\varphi^d_1o_{12},\varphi^d_2o_{12})|\ge\mathcal W(\varphi^d_1o_{12},\varphi^d_2o_{12})-\epsilon_g(d)|\varphi|_\infty^2>0$ for small $d$. Proposition~\ref{C-371-17} at $z=2$ gives the properties of $\sigma\circ Z_2$, its operators and fields through the unitary map, and shows that it is not covariant at order $2$; Theorem~\ref{C-371-4}(b) gives $\mathfrak M(\sigma\circ Z_2)=\mathfrak N_{18}$, and Theorem~\ref{C-371-4}(d) the cone. The law holds for $\sigma\circ Z_z$, with some label in the place of $b$, at $z=1$ alone, and at the pairs of off-diagonal weights alone, with the label $b$, exactly at $z^2=1$, by Proposition~\ref{C-371-17}: the pairs of diagonal weights fix the label, and the mixed pair of a diagonal and an off-diagonal weight excludes $z=-1$ (Proposition~\ref{C-371-24}(d) at $a=e=1$); the last assertion is Remark~\ref{C-371-33}.

\begin{remark}\label{C-371-18}
Uniqueness of the vacuum holds in the components of clause (7) of Theorem~\ref{C-371-5}. Near coinciding points the bounds of clause (6) give a growth of the kernels $s_n$ at most like $\delta(x)^{-9n}$ and the extension $S_n$ of $\sigma$ to the test functions flat at the coincidences, with numbers $C_V$, $C$ that depend on the constants and $g$; the limits are taken through periodic Gibbs states, and along growing tori (D0)--(D3) and positivity at the three coordinate families of mirrors hold (Proposition~\ref{C-371-22}(e)). Clause (8) and Theorem~\ref{C-371-4}(f) concern functionals with (D0)--(D4) only: $\sigma\circ Z_2$ is $\sigma$ written in the basis $Z_2$ of $V$, covariant and positive at every plane for the representation $R\mapsto Z_2^{-1}\rho(R)Z_2$, and the tensor limits are covariant by clause (3). For the tensor limits clause (6) gives a representation of the Poincar\'e group that acts covariantly on the fields.
\end{remark}

\section{Two short scales and point splitting}\label{C-352-1}

A pair of nearby action-density insertions has a tensor-valued leading term. We compare the close pair with this tensor at a second, larger distance, retaining the dependence on both scales. The three-point comparison, its extension to plane fields, the fixed-distance point-splitting limit and the four-point law have different hypotheses and are stated separately.

The proof begins with the local kernels and their Gaussian contractions, then controls their remainders through the renormalisation steps. The treatment of narrow test functions is needed to distinguish continuous continuum kernels from the almost-everywhere classes obtained on a fixed torus.

\subsection{Notation and hypotheses}

\emph{Constants and conventions.} Fix one common choice of the real and sourced construction constants,
as in Proposition~\YMMref{69.1}{595} and Section~\YMMref{82}{665},
at a fixed source radius $2\varrho>0$. The real construction uses the
restriction of this choice, with no change in the value of any retained
constant. In particular, the hypotheses of Theorems~\YMMref{AJ}{673}
and~\YMMref{AD}{607} are imposed at this same choice: the blocking factor $L\ge13$ is odd, the cube size is $M=L^{m'}$, $a_j=\eta_j=L^{-j}$, and $A_0,A_1$, $\mathfrak p_0\ge1$, $\bar B$, $B_\sharp$ (with $\bar B=B_\sharp L/(L-1)$), $\alpha_0$, $\beta_{\mathrm s}$, $\kappa$, together with the remaining constants of the real and sourced steps, have these fixed values. By $C,c,c'$, with or without an index, we denote positive numbers made from these constants only, and by $C(g)$, $C(g,\mathsf b)$, $C(g,\ell)$ numbers that may also depend on the displayed arguments; such a number depends on nothing that is not displayed, in particular not on $K$, $\beta$, $\ell$, a level, a length, a position, a direction, a sign or a weight. We put $\delta(t)=A_1(\log t)^{\mathfrak p_0}t^{-1/2}$, $\mathfrak b(t)=A_1(\log t)^{\mathfrak p_0}$, $\varepsilon(t)=A_0(\log t)^{\mathfrak p_0}t^{-1/2}$; $\delta$ and $\varepsilon$ decrease on $[e^{2\mathfrak p_0},\infty)$.

\emph{Couplings, data, sides.} The reference coupling satisfies $0<g\le\gamma$, with the ceiling fixed by Proposition~\YMMref{82.2}{668}. A first-passage datum $(K,\beta)$ at $g$ consists of a bare inverse coupling $\beta$ and the number $K$ of renormalisation steps after which, for the first time, the running inverse coupling that starts at $\beta$ is at most $L(g^{-2}+B_\sharp)$, as in Proposition~\YMMref{74.1}{615}. Its run is $x_0=\beta>x_1>\dots>x_K>g^{-2}$; we write $y_j=\log x_j$, $g_j=x_j^{-1/2}$, $\delta_j=\delta(x_j)$, $\mathfrak b_j=\mathfrak b(x_j)$, $\varepsilon_j=\varepsilon(x_j)$. The label $\mathsf b=a_K\beta$ is a length, $g^{-2}<\mathsf b\le\bar{\mathsf b}:=L(g^{-2}+\bar B)$. We put $I(g)=(g^{-2}+\bar B,\,L(g^{-2}+B_\sharp)-\bar B]$; for $\mathsf b\in I(g)$ the pair $(K,L^K\mathsf b)$ is a first-passage datum at $g$ with the label $\mathsf b$ for every $K\ge1$. $\mathfrak L(g)$ is the set of the sides at $g$, each of the form $\ell=2L^m$, and $\ell_g=\min\mathfrak L(g)\ge26$ is the least side; $T_\ell$ is the torus of side $\ell$ with the lattice $a_K\mathbb Z^3$, and the lattice $(\ell,K)$ is this torus with the coarser lattices of its levels.

\emph{States and limits.} $\mu_{K,\ell,\beta}$ is the measure of the torus $T_\ell$ at the cutoff $K$ and the bare coupling $\beta$, $\langle\,\cdot\,\rangle_{K,\ell,\beta}$ its expectation, $\tilde\mu_{K,\ell,\beta}$ its periodic extension to $\mathbb L_K=a_K\mathbb Z^3$, and $\mathcal G^{\mathrm{per}}_{K,\beta}$ is the set of periodic Gibbs states (equation~\YMMref{(1819)}{698} and the definition of periodic states following it). By a state at $(K,\beta)$ we mean either a periodic extension $\tilde\mu_{K,\ell,\beta}$ with $\ell\in\mathfrak L(g)$ or a periodic Gibbs state in $\mathcal G^{\mathrm{per}}_{K,\beta}$. The test pairs $\mathfrak T_{T_\ell}(r;\vartheta)$ of scale $r$ and shape $\vartheta$ are those of Lemma~\ref{lem:ward-test-pairs}.

Along first-passage data $(K_j,\beta_j)$ at $g$ with
$K_j\to\infty$ and $a_{K_j}\beta_j\to\mathsf b>0$, the torus
connected functions converge to $S_n^{(\ell,\mathsf b)}$ on real
Lipschitz tuples with pairwise disjoint supports, independently of
the defining sequence, by Theorem~\YMMref{A}{5}. The first-passage
labels lie in $(g^{-2},\bar{\mathsf b}]$ by
Proposition~\YMMref{74.1}{615}; in particular their limiting labels
are positive. A law on $\mathbb R^3$ is a functional on separated
tuples of real $C_c^1$ functions with the properties of the
growing-torus or volume-first limits in
Propositions~\YMMref{92.1}{710}, \YMMref{92.3}{712} and
\YMMref{92.4}{713}. Limits of periodic-Gibbs sequences and the
components of their canonical decompositions are examples, by
Theorems~\YMMref{AP}{780} and~\YMMref{AL}{746}.

For limits of periodic-Gibbs sequences and their canonical
components, Theorem~\YMMref{AP}{780} and
Proposition~\YMMref{108.3}{771} give unique continuous moment
kernels, real analytic and invariant under Euclidean isometries.
The connected kernels have the same properties by the finite
moment--cumulant formula. For an arbitrary growing-torus law,
the inherited bounds initially give distributions on separated
configurations. Full rotation invariance has the additional Ward
criterion of Theorem~\YMMref{AQ}{785}; the analytic-kernel conclusion
is available when the reconstruction hypotheses \textup{(Rec)} of
Proposition~\YMMref{108.3}{771} hold. The pointwise assertions
below for such a law explicitly assume continuity of the kernels
they evaluate. The finite-cutoff, integrated and almost-everywhere
assertions retain their stated hypotheses.

\emph{Fields.} For a plaquette $p$, $x(p)$ is its barycentre, $\lambda(p)\in\{1,2,3\}$ its dual direction (the direction orthogonal to its plane), and $a_p=1-\frac12\operatorname{tr}\partial U(p)$. For a real weight $w$ on the plaquettes with finite support, $\mathcal O(w)=a_K^{-1}\sum_pw(p)a_p$, $|w|_\infty=\max_p|w(p)|$ and $\Sigma(w)=\{x(p):w(p)\ne0\}$; for a bounded function $f$, $\mathcal O(f)=\mathcal O(f\circ x)$. The plane fields are $\mathcal O_\lambda(\varphi)=\mathcal O(\mathbf 1_{\{\lambda(p)=\lambda\}}\,\varphi\circ x)$, with plane weights, so that $\mathcal O=\sum_\lambda\mathcal O_\lambda$; $T_\lambda=\mathcal O_\lambda-\frac13\mathcal O$ are the diagonal elements of the traceless tensor.

Off the diagonal stand the bond forms. For $V\in SU(2)$ put $\mathrm w(V)=(V-V^{-1})/2i$, and $\langle X,Y\rangle=\frac12\operatorname{tr}XY$. A bond $b=\langle x,x+a_Ke_\mu\rangle$ has the base point $x$, the direction $\mu_b=\mu$ and the midpoint $x_b=x+\frac12a_Ke_\mu$. For signed directions $\alpha,\alpha'$ with $\alpha\ne\pm\alpha'$ let $W_x(\alpha,\alpha')$ be the product of $U$ along the closed path $x\to x+a_Ke_\alpha\to x+a_Ke_\alpha+a_Ke_{\alpha'}\to x+a_Ke_{\alpha'}\to x$, where $e_{-\alpha}=-e_\alpha$. For a direction $\alpha\ne\mu_b$ the loops $W_x(\mu_b,\alpha)$ and $W_x(\mu_b,-\alpha)$ run around the two plaquettes of the plane spanned by $\mu_b$ and $\alpha$ that contain $b$, both from the base point $x$ along $b$ first, and we put $u_{b,\alpha}=-\frac12[\mathrm w(W_x(\mu_b,\alpha))-\mathrm w(W_x(\mu_b,-\alpha))]$; with $\lambda_b<\lambda'_b$ the two directions other than $\mu_b$, $q_b=\langle u_{b,\lambda_b},u_{b,\lambda'_b}\rangle$.

For a real weight $\mathsf w$ on the bonds, $Y[\mathsf w]=-a_K^{-1}\sum_b\mathsf w(b)q_b$, and $\mathcal O^{\mathrm{od}}_{\lambda\lambda'}(\varphi)=Y[\mathsf w]$ at $\mathsf w(b)=\varphi(x_b)\mathbf 1_{\{\mu_b=\mu\}}$, where $\{\lambda,\lambda',\mu\}=\{1,2,3\}$: the bond weight is taken at the midpoint. Written in the angles $F_p$ of the plaquette variables, the linear part of $u_{b,\lambda}$ is $\sum_pu_{b,\lambda}(p)F_p$, with $u_{b,\lambda}(p)=\frac12\epsilon_{\lambda\mu_b\lambda'}\epsilon_p$ at the two plaquettes with the barycentres $x_b\pm\frac12a_Ke_\lambda$ and $u_{b,\lambda}(p)=0$ at all other plaquettes; here $\{\lambda,\mu_b,\lambda'\}=\{1,2,3\}$, $\epsilon_{\lambda\mu_b\lambda'}$ is the sign of this permutation and $\epsilon_p=\pm1$ is the orientation sign of $p$.

For a unit vector $n$,
\[
T_n(\varphi)=\sum_\lambda n_\lambda^2\,T_\lambda(\varphi)+\sum_{\lambda<\lambda'}n_\lambda n_{\lambda'}\,\mathcal O^{\mathrm{od}}_{\lambda\lambda'}(\varphi).
\]
A mixed source is a pair $\mathfrak s=(w,\mathsf w)$ of real weights on the plaquettes and on the bonds, $\mathcal O^T(\mathfrak s)=\mathcal O(w)+Y[\mathsf w]$, $|\mathfrak s|_\infty=|w|_\infty+8|\mathsf w|_\infty$.

\emph{Connected functions and natural sizes.} For weights $w_1,\dots,w_n$ and $w_\lambda=\sum_a\lambda_aw_a$, $S^T_n=S^T_{n;K,\ell,\beta}(w_1,\dots,w_n)=\partial_{\lambda_1}\cdots\partial_{\lambda_n}\log\langle e^{\mathcal O(w_\lambda)}\rangle_{K,\ell,\beta}|_{\lambda=0}$; for functions $f_a$ it is taken at $w_a=f_a\circ x$; $\langle X;Y\rangle$ denotes $S^T_2$, and $S^T_n[\mathfrak s]$ the same in a state $\mathfrak s$. For a function $f$ on $\mathbb R^3$, $f^{\langle\ell\rangle}$ is its periodisation. We write $\mathsf u_a=\mathsf b^{-1}|w_a|_\infty$. Three bumps of width $\rho$, two of them at the near distance $r$ and the third at the far distance $\mathsf d$, are measured against $\mathrm{nat}_3=\rho^9r^{-3}\mathsf d^{-6}\mathsf b^{-3}\prod_{a\le3}|w_a|_\infty$, and two such pairs against $\mathrm{nat}_4=\rho^{12}r^{-6}\mathsf d^{-6}\mathsf b^{-4}\prod_{a\le4}|w_a|_\infty$ (at functions $\|f_a\|_\infty$ replaces $|w_a|_\infty$); a kernel of three points is measured against $\mathsf b^{-3}r_*^{-3}\mathsf d_*^{-6}$ and a kernel of two points against $\mathsf b^{-2}|x-y|^{-6}$. These are the natural sizes.

\emph{Placed weights.} Real weights $w_1,w_2$ are placed at $(\rho,r)$ if there are centres $x_a$ with $\Sigma(w_a)\subset\mathbb B(x_a,\rho)$ (closed balls), $\operatorname{dist}(\Sigma(w_1),\Sigma(w_2))\ge r$ and $|x_1-x_2|\le8r$. Weights $w_1,w_2,w_3$ are placed at $(\rho,r,\mathsf d)$ if besides $\Sigma(w_3)\subset\mathbb B(x_3,\rho)$, $\operatorname{dist}(\Sigma(w_a),\Sigma(w_3))\ge\mathsf d$ and $|x_a-x_3|\le18\mathsf d$ for $a=1,2$. Weights $w_1,\dots,w_4$ form two pairs at $(\rho,r,\mathsf d)$ if $(w_1,w_2)$ and $(w_3,w_4)$ are placed at $(\rho,r)$ and $\operatorname{dist}(\Sigma(w_a),\Sigma(w_b))\ge\mathsf d$, $|x_a-x_b|\le18\mathsf d$ for $a\le2<b$. A set $\beta$ of sources crosses if it meets both clusters, $\{1,2\}$ and $\{3\}$ (or $\{3,4\}$), and $c(\beta)$ is the number of whole pairs in $\beta$.

\emph{Free objects.} $\mathbb K_3(u)=(4\pi|u|^3)^{-1}(3\hat u\hat u^{\mathsf T}-1)=(4\pi)^{-1}\nabla\nabla|u|^{-1}$, $\hat u=u/|u|$; $k_3(x,y,z)=\operatorname{tr}\mathbb K_3(x-y)\mathbb K_3(y-z)\mathbb K_3(z-x)$; $k_4(x_1,\dots,x_4)=\frac12\sum_\varsigma\operatorname{tr}\prod_l\mathbb K_3(x_{\varsigma(l)}-x_{\varsigma(l+1)})$, the sum over the six cyclic orders of $\{1,2,3,4\}$ (three rings, each once); $\mathcal K_{\lambda\mu}(\varphi_1,\varphi_2)=\iint\varphi_1(x)\varphi_2(y)[\mathbb K_3(x-y)_{\lambda\mu}]^2\,dx\,dy$; and for a unit vector $n$ and $X\ne0$
\[
\mathsf t^\circ(n,X)=\tfrac32\mathsf b^{-2}\bigl[\langle n,\mathbb K_3(X)^2n\rangle-\tfrac13\operatorname{tr}\mathbb K_3(X)^2\bigr]=\frac3{32\pi^2}\,\mathsf b^{-2}|X|^{-6}\bigl(3(n\cdot\hat X)^2-1\bigr).
\]
$C^{\mathrm{free}}_K$ is the free curvature kernel on the bare plaquettes of the lattice $(\ell,K)$, $\mathcal C_{(j)}$, $1\le j\le K$, are its slices by level, and $\upsilon_j=(a_{K-j+1}x_{j-1})^{-1}$. Over a base level $k_\flat$ we put $\mathcal S_i=\sum_{k_\flat<j\le i}\upsilon_j\mathcal C_{(j)}$; for a nonempty set $\beta$ of sources with $m=|\beta|$, $R_{i;\beta}=\sum w_{a_1}\mathcal S_iw_{a_2}\cdots\mathcal S_iw_{a_m}$, the sum over the orderings $(a_1,\dots,a_m)$ of $\beta$, the weights acting as diagonal matrices; $q_{i;\beta}(F)=\frac1{2a_K}\sum_{p,q}R_{i;\beta}(p,q)\langle F_p,F_q\rangle$ for $\mathfrak{su}(2)$-valued functions $F$ of the plaquettes; and $g_{i;\beta}=\frac32\sum_\varsigma\operatorname{Tr}\prod_l(w_{\varsigma(l)}\mathcal S_i)$, the sum over the cyclic orders $\varsigma$ of $\beta$. At $\beta=\{1,2\}$ we write $q_n=q_{n;12}$, so $q_n(F)=a_K^{-1}\sum_{p,q}w_1(p)\mathcal S_n(p,q)w_2(q)\langle F_p,F_q\rangle$, and $\mathcal Q_i=q_{i+1}-q_i$. The lattice ring of three weights is $3\operatorname{Tr}\prod_a(w_a\mathsf b^{-1}C^{\mathrm{free}}_K)$. Finally $\chi^\rho_x=\frac3{\pi\rho^3}(1-|\cdot-x|/\rho)_+$ is the radial tent of radius $\rho$ and mass $1$.

\emph{The renormalisation group.} We use the following signs. The level $j$ is the lattice of spacing $a_{K-j}$, level $0$ being the bare lattice; for a length $\varrho$, $\hat\varrho_j=\varrho/a_{K-j}$. The chain classes $\mathsf{Cc}_n(r_n)$ have the radii $r_K=0$, $r_i=L^{-1}r_{i+1}+\frac{11}{10}\mathsf c_{\mathfrak h}\delta_i$ ($\mathsf c_{\mathfrak h}\ge1$); for $A\in\mathsf{Cc}_n(r_n)$ we put $\tilde U=e^{i\eta_nA}$ and $\mathsf{ang}^{[n]}_p(A)=\eta_n(\mathrm dA)(p)$, and $\mathbf 1$ is the trivial configuration. The witness depth is $s(\varrho)=\min\{s\ge0:L^s\varrho\ge C_W(s)\}$, where $C_W(s)=C_{W,0}A_0^2(\log[L^{s+1}(g^{-2}+\bar B)])^{2\mathfrak p_0}$, with the room constant $C_{W,0}=\max\{1,4L^2M/A_0^2\}$, so that $C_{W,0}A_0^2\ge4L^2M$ (equations~\YMMref{(2320)}{843} and~\YMMref{(2322)}{844}); the working level is $\mathsf k(\varrho)=\max\{K-s(\varrho),2\}$; $Y_s=(s+1)\log L+\log(g^{-2}+\bar B)$, so that $C_W(s)=C_{W,0}A_0^2Y_s^{2\mathfrak p_0}$, $Y_\varrho=Y_{s(\varrho)}$, and $\mathsf{Lm}_s$ is a constant times $Y_s^{5\mathfrak p_0}$.

For sources $w_1,\dots,w_p$ the weighted progression $E^{(j)}_\lambda$, $1\le j\le K$, is the sequence of families of local terms of level $j$ that the renormalisation steps make with the source $w_\lambda$, for $\lambda$ in a neighbourhood of $0$; $E^{(j)}=E^{(j)}_0$ and $E^{(j)}_\beta=\partial_\beta E^{(j)}_\lambda|_0$. The relative exponent is
\[
D^\lambda_i(U)=\mathcal O(w_\lambda)(U)+\sum_{i'\le i}\mathsf{Sum}[E^{(i')}_\lambda-E^{(i')}](U),\qquad\mathsf{Sum}[\mathsf f](U)=\sum_{(X,z)}\mathsf f(X,U,z),
\]
the sum over the domains $X$ of the level and their base points $z$; $G_{i;\beta}(A)=\partial_\beta D^\lambda_i(e^{i\eta_iA})|_0$, the pair functional on configurations is defined by $\mathsf F_n(e^{i\eta_nA})=G_{n;12}(A)$, the jets are $\mathsf J_n=G_{K;[n]}(0)$, and $\Delta[\mathsf f]=\mathsf{Sum}[\mathsf f](\tilde U)-\mathsf{Sum}[\mathsf f](\mathbf 1)$. A family $\mathsf f$ of level $i$ has the norms $\|\mathsf f\|_{i,\kappa'}=\sup|\mathsf f(X,P,z)|e^{\kappa'd_i(X)}$, $d_i$ the linear size of a domain, and $\lambda_{\mathrm{out}}=(1+4\beta_{\mathrm s})\kappa$; $\pi_i$ is the partition into cubes of side $Ma_{K-i}$; for a weight $w_a$, $\mathsf{Cr}_a$ is the set of the corners of the plaquettes of $\operatorname{supp}w_a$, $\mathsf N_i(w_a)$ the number of the cubes of $\pi_i$ that hold a site of $\mathsf{Cr}_a$, and $\mathsf D^{ab}_i$ the least sup distance of $\mathsf{Cr}_a$ from $\mathsf{Cr}_b$ in level-$i$ units.

The step map $\mathsf{St}^{[j]}$ is the renormalisation step of the height $j$, viewed as a function from bare arrays $\mathfrak f=((G^{(i)},e_i)_{i\le j},(H))$ of level $j$ to families of level $j+1$; $N(\mathfrak f)$ is the size of an array and $\mathfrak n[\mathfrak f]\le N(\mathfrak f)$ the size of the family that it puts in the fourth place of the one-step theorem; $N^{[j]}_\star$ and $\theta^{[j]}=K_{\mathbf M}(1+\beta_{\mathrm s})\delta_j/\alpha_0$ are numbers of the step $j$, and $\mathsf W_w$ is the weight member of a weight $w$. Further construction objects are defined where they enter.

\begin{definition}[Ceilings, thresholds, numbers]\label{C-352-2}
\begin{enumerate}
\item The ceiling $(\Gamma1)$: $0<g\le\gamma_{\mathrm{jt}}$; $g^{-2}\ge\max\{x_{\mathrm J},e^{2\mathfrak p_0},e^{y_{\mathsf p}}\}$; $\delta(g^{-2})\le\min\{1/2c_{\mathrm{rt}},1/c_\epsilon,1/2C_{\mathrm{cu}}\}$; $\theta(g^{-2}):=K_{\mathbf M}(1+\beta_{\mathrm s})\delta(g^{-2})/\alpha_0\le\tau^\natural$. The ceiling $(\Gamma2)$: $0<g\le\gamma^{\mathrm{sd}}_W$, $g^{-2}\ge x_{\mathrm J}$, $\log g^{-2}\ge\mathsf Y_{\mathrm{lt}}(\frac12)$; data under $(\Gamma2)$ have $K\ge2$. $(\Gamma)$ is $(\Gamma1)$ together with $(\Gamma2)$. Here $\gamma_{\mathrm{jt}}$ is defined in \eqref{eq:pair-threshold}, $\gamma^{\mathrm{sd}}_W$ in equation~\YMMref{(2327)}{846}, and $x_{\mathrm J}$ in \eqref{eq:weighted-history-threshold}; Proposition~\ref{prop:weighted-canonical-history} verifies the required actual-step hypotheses for every finite number of sources. The least-torus threshold $\mathsf Y_{\mathrm{lt}}$ is specified below; $y_{\mathsf p},c_{\mathrm{rt}},c_\epsilon,C_{\mathrm{cu}},\tau^\natural$ are numbers formed from the constants in Lemmas~\ref{C-352-10}, \ref{C-352-12} and~\ref{C-352-20}.

For $0<\eth<1$, put $c'=(1-\eth)c_\star$ and
$C_{\rm sd}=8L^6(\mathsf C_{\rm maj}^{\Sigma}+2+c_{\rm low}+9/2)$.
Choose $\mathsf Y_{\rm lt}(\eth)\ge1$ so that all the history thresholds
in Proposition~\ref{prop:ward-fixed-volume-rarity} hold above it and
\begin{equation}\label{eq:pair-least-torus-threshold}
 2C_{\rm sd}y[MR(e^y)]^3\le\tfrac12c'A_0^2y^{2\mathfrak p_0}
 \qquad(y\ge\mathsf Y_{\rm lt}(\eth)).
\end{equation}
Such a finite threshold exists: $R(e^y)<Ly^r$ and
$2\mathfrak p_0>3r+2$, as in the proof of
Lemma~\YMMref{126.3}{846}. On the least admissible side that proof gives
\[
 \mathsf B(x_K)\ell_g^3\le C_{\rm sd}y_K[MR(e^{y_K})]^3.
\]
Thus $\log g^{-2}\ge\mathsf Y_{\rm lt}(\eth)$ implies
$2\mathsf B(x_K)\ell_g^3\le c'p_K^2/2$, the stronger small-side
condition in Proposition~\ref{prop:ward-fixed-volume-rarity}.

Since $\delta$ decreases on $[e^{2\mathfrak p_0},\infty)$, each member is a condition ``$g\le$ a number formed from the constants''. $(\Gamma^{(4)})$ is $(\Gamma)$ with $c^{(4)}_\epsilon$ of Lemma~\ref{C-352-20} in the place of $c_\epsilon$. As $c_\epsilon\le c^{(4)}_\epsilon$ (Lemma~\ref{C-352-20}(a), where $C_{\mathrm{cu}}$, $c_{\mathrm{cu}}$ and the numbers $c^{(n)}_\epsilon$ are formed), $(\Gamma1)$ with $c^{(4)}_\epsilon$ in the place of $c_\epsilon$ contains $(\Gamma1)$, and $(\Gamma^{(4)})$ contains $(\Gamma)$: every $g$ under $(\Gamma^{(4)})$ satisfies $(\Gamma)$. For the bond forms we use the ceiling $\gamma_{\mathrm N}=x_{\mathrm N}^{-1/2}$, where $x_{\mathrm N}$ is the least real number $\ge x_{\mathrm{jt}}=\gamma_{\mathrm{jt}}^{-2}$ such that every $t\ge x_{\mathrm N}$ satisfies $t\ge\max\{x_{\mathrm J},e^{2\mathfrak p_0},e^{y_{\mathsf p}}\}$, $\delta(t)\le\min\{1/2c_{\mathrm{rt}},1/c_\epsilon\}$, $\theta(t)=K_{\mathbf M}(1+\beta_{\mathrm s})\delta(t)/\alpha_0\le\tau^\natural$, $\mathfrak b(t)\ge1$ and $c_{\mathrm A}\alpha_0\ge5M\mathsf c_{\mathfrak h}\delta(t)(1+\log t)$, with $c_{\mathrm A}$ of Lemma~\ref{C-352-14}.

So $\gamma_{\mathrm N}\le\gamma_{\mathrm{jt}}$, and every $0<g\le\gamma_{\mathrm N}$ satisfies $(\Gamma1)$ but for its member $\delta(g^{-2})\le1/2C_{\mathrm{cu}}$.
\item $r_1(\mathsf b)>0$ is the number of Lemma~\ref{C-352-12}, formed from the constants and $\mathsf b$; it does not increase in $\mathsf b$. $\mathsf d_4(g)>0$ is the number fixed in Proposition~\ref{C-352-24}, formed from the constants and $g$, with $\mathsf d_4(g)\le\min\{e^{-2(15\mathfrak p_0+2)},r_1(\bar{\mathsf b})\}$; $\mathsf d_5(g)=\min\{2\mathsf d_4(g),\,\min\{1,16r_1(\bar{\mathsf b})\},\,\frac1{32}\ell_g\}$. For two pairs, $\mathsf d^{(4)}_4(g)$ is fixed in the same Proposition, and $\mathsf d^{(4)}_5(g)=\min\{2\mathsf d^{(4)}_4(g),\frac1{48}\ell_g\}$.
\item $\bar{\mathsf M}^{(24)}=1+2L^{24}\mathsf C_{\mathrm{ag}}(25,\frac12)$, where $\mathsf C_{\mathrm{ag}}$ is defined by \eqref{eq:ward-volume-constants}; $\Lambda_3\ge1$ is the number of Lemma~\ref{C-352-26}; $C_{\mathrm{lat}}$ is the absolute constant in equations~\YMMref{(2245)--(2246)}{824}, comparing $C^{\mathrm{free}}_K$ with $\mathbb K_3$ on every side; $C^{(3)}=3\cdot81^3(1+3C_{\mathrm{lat}})^2C_{\mathrm{lat}}+2.1\cdot10^7$, $C^{(4)}=3[4\cdot81^4C_{\mathrm{lat}}(1+3C_{\mathrm{lat}})^3+1.49\cdot10^8]$, $C_{\mathrm{Rs}}=7400+729C_{\mathrm{lat}}(2+3C_{\mathrm{lat}})$, the constants of the lattice sums of Lemma~\ref{C-352-27}; $C_{\mathrm{PS}}=62+1089C_{\mathrm{lat}}(2+3C_{\mathrm{lat}})$ and $C^{\mathrm{od}}_{\mathrm{PS}}=62+729(2+5C_{\mathrm{lat}})(1+C_{\mathrm{lat}})$, the constants of the lattice sums at two balls on every torus (Lemma~\ref{C-352-36a}); $C^{\mathrm{tr}3}_3(g)$, $C^{\mathrm{tr}3}_4(g)$, $C^{\mathrm{tr}}_2(g)$ are the numbers of Proposition~\ref{C-352-33}, formed from the constants and $g$.
\item With the number $C(g)$ of Proposition~\ref{C-352-24}: $\epsilon^{\mathrm{nar}}(g,t)=C(g)(\log\frac1t)^{39\mathfrak p_0+2}t^{1/2}$, $\epsilon^{\mathrm{nar}}_2(g,t)=C(g)(\log\frac1t)^{19\mathfrak p_0}t^{1/2}$, $\epsilon^{\mathrm{nar}}_4(g,t)=C(g)(\log\frac1t)^{47\mathfrak p_0+2}t^{1/2}$; and
\begin{align*}
\epsilon_{\mathrm A}(g,t)&=512\bigl[\epsilon^{\mathrm{nar}}(g,\tfrac t2)+3C^{(3)}(5t/\ell_g)^3+13\Lambda_3^3(\bar{\mathsf M}^{(24)})^3\bar{\mathsf b}^3t^{3/2}\bigr],\\
\mathsf E_3(g,t)&=512\bigl[\epsilon^{\mathrm{nar}}(g,\tfrac t2)+3C^{(3)}(5t/\ell_g)^3+\bigl(13\Lambda_3^3(\bar{\mathsf M}^{(24)})^3+C^{\mathrm{tr}3}_3(g)\bigr)\bar{\mathsf b}^3t^{3/2}\bigr],\\
\epsilon_{\mathrm P}(g,t)&=64\bigl[\epsilon^{\mathrm{nar}}_2(g,\tfrac t2)+3\Lambda_3^2(\bar{\mathsf M}^{(24)})^2\bar{\mathsf b}^2(\tfrac t2)^{3/2}+\tfrac32C_{\mathrm{Rs}}(t/2\ell_g)^3\bigr],\\
\mathsf E_2(g,t)&=\epsilon^{\mathrm{nar}}_2(g,t)+\bigl[3\Lambda_3^2(\bar{\mathsf M}^{(24)})^2+C^{\mathrm{tr}}_2(g)\bigr]\bar{\mathsf b}^2t^{3/2}+\tfrac32C_{\mathrm{Rs}}(t/\ell_g)^3,\qquad\epsilon_\sigma(g,t)=64\sup_{0<t'\le t}\mathsf E_2(g,\tfrac{t'}2).
\end{align*}
So $\epsilon_{\mathrm A},\mathsf E_3\le C(g)[(\log\frac1t)^{39\mathfrak p_0+2}t^{1/2}+(t/\ell_g)^3]$ and $\epsilon_{\mathrm P},\epsilon_\sigma\le C(g)[(\log\frac1t)^{19\mathfrak p_0}t^{1/2}+(t/\ell_g)^3]$ for small $t$; all of them are functions of $(g,t)$ alone.
\item \emph{The threshold $r_\alpha(g,\mathsf b)$ of Theorem~\ref{C-352-7}.} For a label $\mathsf b$ and $0<r\le L^{-4}$ consider every first-passage datum $(K,\beta)$ at $g$ with $a_K\beta=\mathsf b$, $128a_K\le r$ and $s(\frac r2)\le K$, and the near pair $w_{1,2}=\chi^{r/8}_{c\pm\frac r2e_\lambda}\circ x$ of that Theorem, a test pair of the scale $\frac r2$; $n_1$ and $n_0$ are the least levels with $a_{K-n_1}\ge(\frac r2)^{3/4}$ and $a_{K-n_0}\ge r^{1/2}$. The conditions are:
\begin{enumerate}
\item[\textup{($\alpha$1)}] $\mathsf k(\frac r2)<n_1<n_0<K$, $2(r+a_K)\le a_{K-n_0}$ and $\frac{27}{8M^3}(r/2a_{K-n_1})^3\le1$;
\item[\textup{($\alpha$2)}] for $\mathsf k(\frac r2)\le i<n_1$, $n_0\le n<K$ and $m=n-i-1$: $(4\mathsf h_i+4)a_{K-i-1}\le\ell_g$ and $7\mathsf h_iL^{-2m}r_n\le\delta_i/2\Lambda_i$, where $\mathsf h_i$ and $\Lambda_i$ are the numbers of Proposition~\ref{C-352-19} for the near pair at $(i,n)$, formed with $R_w=\frac54\hat r_{i+1}+1$ and $\bar{\mathsf e}=2\bar{\mathsf w}_1\bar{\mathsf w}_2\mathsf E_i$, $\bar{\mathsf w}_a=81\hat r_i^3|w_a|_\infty/a_{K-i}$;
\item[\textup{($\alpha$3)}] $\frac r2\le\min\{r_0(\mathsf b),r_1(\mathsf b)\}$;
\item[\textup{($\alpha$4)}] $s(\frac r2)\ge3$ and $\mathsf{Lm}_{s(r/2)}\le L^{s(r/2)-3}$;
\item[\textup{($\alpha$5)}] $2\sqrt3\,\delta_{\mathrm{tp}}L(L-1)M(\log\frac2r)^{2\mathfrak p_0}\ge3\log L+\log2$, so that Lemma~\ref{C-352-22}(b) applies at the level $\mathsf k(\frac r2)$ to two plaquettes of the two supports of the near pair.
\end{enumerate}
$r_\alpha(g,\mathsf b)$ is the largest number $r_\star\in(0,L^{-4}]$ such that \textup{($\alpha$1)}--\textup{($\alpha$5)} hold at every $0<r\le r_\star$ and every such datum. It is positive. \textup{($\alpha$1)} holds for $r\le L^{-4}$ at every datum. \textup{($\alpha$2)}, \textup{($\alpha$4)} and \textup{($\alpha$5)} hold for $r$ below a number formed from the constants and $g$: by the working level (Lemma~\ref{C-352-10}(c)) and step 6 of the proof of Proposition~\ref{C-352-19}, $\mathsf h_i\le C(1+Y)^{2\mathfrak p_0}$ and $\Lambda_i\le C(1+Y)$ with $Y=Y_{r/2}\le2\log\frac2r+C(g)$; $r_n\le\frac85\mathsf c_{\mathfrak h}(\frac54)^{1/2}L^{(n-i)/2}\delta_i$; and $L^{-(n-i)}=a_{K-i}/a_{K-n}<(\frac r2)^{3/4}r^{-1/2}\le r^{1/4}$; none of these bounds depends on $\mathsf b$ or on $K$. \textup{($\alpha$3)} holds for $\frac r2\le\min\{\inf_{\mathsf b'}r_0(\mathsf b'),r_1(\bar{\mathsf b})\}$, the infimum over $g^{-2}<\mathsf b'\le\bar{\mathsf b}$ being positive by (B.1) and $r_1$ not increasing. So $r_\alpha(g,\mathsf b)\ge r^-_\alpha(g)$ for all $\mathsf b\in(g^{-2},\bar{\mathsf b}]$, with a positive number $r^-_\alpha(g)$ formed from the constants and $g$.
\item \emph{The ceiling $(\Gamma3)$ and the threshold $r^{\mathcal W}_\alpha(g,\mathsf b)$.} $(\Gamma3)$ asks: $c_{\mathrm A}\alpha_0\ge8MB_3\varepsilon(t)(1+\log t)$ for all $t\ge g^{-2}$; $g\le2C_{\mathrm A}C_{\mathrm E,3}M^2B_3^2A_0^2/C_{\mathrm E,1}\alpha_0^2$; and $7488B_3\varepsilon(g^{-2})\le\epsilon_{\mathrm G}$, $C_{\mathcal W}\varepsilon(g^{-2})\le\frac1{120}$ with $C_{\mathcal W}=\max\{720C^\sharp_{\mathrm G},2\}B_3$. Here $c_{\mathrm A}$, $C_{\mathrm A}$ and $C_{\mathrm E,1}$ are the numbers of Lemma~\ref{C-352-14} and of its proof, $C_{\mathrm E,3}$ is the number in Lemma~\ref{lem:source-support-sum} at $q=2$, $2B_3\varepsilon_n$ is the radius at which the small-field backgrounds are regular (Lemma~\ref{C-352-25}(d)), and $\epsilon_{\mathrm G}$, $C^\sharp_{\mathrm G}$ are the numbers of the lemma on a regular gauge on a cube, which precedes Lemma~\ref{C-352-38}.

As $\varepsilon(t)(1+\log t)\to0$ and $\varepsilon$ decreases on $[e^{2\mathfrak p_0},\infty)$, each member is a condition ``$g\le$ a number formed from the constants and $\beta_{\mathrm H}$''. $r^{\mathcal W}_\alpha(g,\mathsf b)$ is the largest number $r_\star\in(0,r_\alpha(g,\mathsf b)]$ such that, at every $0<r\le r_\star$ and every datum as in (5), the numbers $\mathsf h_i$ and $\Lambda_i$ of \textup{($\alpha$2)}, formed with $C_{\mathcal W}\varepsilon_n$ in the place of $r_n$, satisfy for $\mathsf k(\frac r2)\le i<n_1$, $n_0\le n\le K$ and $m=n-i-1$: $(4\mathsf h_i+4)a_{K-i-1}\le\ell_g$, $7\mathsf h_iL^{-2m}C_{\mathcal W}\varepsilon_n\le\delta_i/2\Lambda_i$ and $(\mathsf h_i+1)L^{-m}\le3$, and besides $Lr^{1/2}\le(2\bar{\mathsf M}^{(24)})^{-1/24}$.

The bounds of (5), with $\varepsilon_n=(A_0/A_1)\delta_n$, give $r^{\mathcal W}_\alpha(g,\mathsf b)\ge r^{\mathcal W-}_\alpha(g)$, a positive number formed from the constants, $\beta_{\mathrm H}$ and $g$.
\end{enumerate}
\end{definition}

\subsection{Principal results}

\begin{theorem}[Three points at two scales]\label{C-352-4}
Let $g$ satisfy $(\Gamma)$.
\begin{enumerate}
\item[(A.1)] \emph{At a cutoff, on the least torus.} Let $(K,\beta)$ be a first-passage datum at $g$, $\mathsf b=a_K\beta$, $\ell=\ell_g$; let $0<\rho\le r\le\mathsf d\le\mathsf d_4(g)$ and $s(\rho)\le K$, and let $w_1,w_2,w_3$ be real weights placed at $(\rho,r,\mathsf d)$, for instance $w_a=f_a\circ x$ with bounded real Borel functions $f_a$. Then
\[
\Bigl|S^T_3(w_1,w_2,w_3)-3\operatorname{Tr}\prod_a(w_a\mathsf b^{-1}C^{\mathrm{free}}_K)\Bigr|\le\bigl[\epsilon^{\mathrm{nar}}(g,\mathsf d)+13\Lambda_3^3\mathsf M'^3\mathsf b^3\mathsf d^{3/2}\bigr]\mathrm{nat}_3,
\]
where $\mathsf M'=\int\alpha^{-24}d\mathbb P\le\bar{\mathsf M}^{(24)}$ (Lemma~\ref{C-352-25}). The bound holds at every ratio $r/\mathsf d\in(0,1]$.
\item[(A.2)] \emph{The kernel on the least torus.} Let $\mathsf b$ be a limit of labels of first-passage data at $g$ (every $\mathsf b\in I(g)$ is one), and $\mathcal D=\{(x,y,z)\in T_{\ell_g}^3:0<|x-y|=r_*\le\mathsf d_*=\operatorname{dist}(z,\{x,y\})\le\mathsf d_5(g)\}$. At every triple of $\mathcal D$, for every $0<\rho<\min\{\frac14r_*,r_1(\bar{\mathsf b})\}$ and all real Lipschitz $f_1,f_2,f_3$ with supports in balls of radius $\rho$ inside the open balls of radius $\frac14r_*$ about $x,y,z$,
\[
\Bigl|S^{(\ell_g,\mathsf b)}_3(f_1,f_2,f_3)-3\mathsf b^{-3}\!\int k_3\prod_af_a\Bigr|\le\epsilon_{\mathrm A}(g,\mathsf d_*)\,\mathsf b^{-3}r_*^{-3}\mathsf d_*^{-6}\,\rho^9\prod_a\|f_a\|_\infty .
\]
Consequently $S^{(\ell_g,\mathsf b)}_3$ has on $\mathcal D$ a Borel kernel $\mathsf s_3$, and at every Lebesgue point of $\mathsf s_3$ in $\mathcal D$, so almost everywhere on $\mathcal D$,
\[
|\mathsf s_3-3\mathsf b^{-3}k_3|\le\epsilon_{\mathrm A}(g,\mathsf d_*)\,\mathsf b^{-3}r_*^{-3}\mathsf d_*^{-6},
\]
uniformly in $r_*/\mathsf d_*\in(0,1]$. Choose the closest pair first and call it $(x,y)$. The triple lies in $\mathcal D$ whenever the resulting far distance $\mathsf d_*=\operatorname{dist}(z,\{x,y\})$ is at most $\mathsf d_5(g)$.
\item[(A.3)] \emph{Means at every triple.} At every triple of $\mathcal D$ and for every $0<\rho<\min\{\frac14r_*,r_1(\bar{\mathsf b})\}$ the mean of $\mathsf s_3-3\mathsf b^{-3}k_3$ over the product of the three balls of radius $\rho$ about its points is at most $(\frac3{4\pi})^3\epsilon_{\mathrm A}(g,\mathsf d_*)\mathsf b^{-3}r_*^{-3}\mathsf d_*^{-6}$ in modulus; with the radial tents $\chi^\rho$ about the three points in place of the normalised indicators of the balls the same holds with $(\frac3\pi)^3$ in place of $(\frac3{4\pi})^3$.
\item[(A.4)] \emph{Every state; limits and components.} (i) Let $(K,\beta)$ be a first-passage datum at $g$, $\mathfrak s$ a state at $(K,\beta)$, $0<\rho\le r\le\mathsf d\le1$, $\rho\le\frac12$, $s(\rho)\le K$, $\rho\le r_1(\mathsf b)$, and let $f_1,f_2,f_3$ be bounded real Borel functions on $\mathbb R^3$ with $\operatorname{supp}f_a\subset\mathbb B(x_a,\rho)$, $\operatorname{dist}(\operatorname{supp}f_1,\operatorname{supp}f_2)\ge r$, $\operatorname{dist}(\operatorname{supp}f_a,\operatorname{supp}f_3)\ge\mathsf d$ for $a=1,2$, all supports lying in one ball of radius at most $\frac18\ell_g-1$. Then
\[
\bigl|S^T_3[\mathfrak s](f_1,f_2,f_3)-S^T_{3;K,\ell_g,\beta}\bigl(f^{\langle\ell_g\rangle}_1,f^{\langle\ell_g\rangle}_2,f^{\langle\ell_g\rangle}_3\bigr)\bigr|\le C^{\mathrm{tr}3}_3(g)\,\mathsf b^3\mathsf d^{3/2}\,\mathrm{nat}_3;
\]
so (A.1) holds for $S^T_3[\mathfrak s]$ with $C^{\mathrm{tr}3}_3(g)\mathsf b^3\mathsf d^{3/2}$ more in the bracket. (ii) Let $(K_j,\beta_j)$ be first-passage data at $g$ with $K_j\to\infty$ and $a_{K_j}\beta_j\to\mathsf b$, $\nu_j$ a state at $(K_j,\beta_j)$, and $\sigma'$ a law on $\mathbb R^3$ with $S^T_3[\nu_j]\to S^{\sigma'}_3$ on separated triples. Assume that $S^{\sigma'}_3$ has a continuous kernel on distinct
triples, as it does for a periodic-Gibbs limit or one of its canonical
components. Then that kernel $\mathsf s_3$ satisfies, at every $(x,y,z)\in(\mathbb R^3)^3$ with $0<|x-y|=r_*\le\mathsf d_*=\operatorname{dist}(z,\{x,y\})\le\mathsf d_5(g)$,
\[
|\mathsf s_3(x,y,z)-3\mathsf b^{-3}k_3(x,y,z)|\le\mathsf E_3(g,\mathsf d_*)\,\mathsf b^{-3}r_*^{-3}\mathsf d_*^{-6},
\]
uniformly in $r_*/\mathsf d_*$. (iii) Every Gibbs sequence at $g$ that converges, with its limit, and the approximants of every component of a canonical decomposition of that limit, with the component, are such data along every subsequence on which the labels converge.
\item[(A.5)] \emph{Gibbs sequences at couplings above the ceilings.} Let $0<g_0\le\gamma$ be a reference coupling with $g\le g_0$. There is an integer $z_+\ge0$ formed from the constants, $g_0$ and $g$ such that for every Gibbs sequence at $g_0$ that converges, and every subsequence along which its labels converge, to $\mathsf b_0$ say, the continuous kernel $\mathsf s_3$ of the limit, and that of every component of a canonical decomposition of it, satisfy at every triple with $0<r_*\le\mathsf d_*\le L^{-z_+}\mathsf d_5(g)$
\[
|\mathsf s_3-3\mathsf b_0^{-3}k_3|\le\max_{0\le z'\le z_+}\mathsf E_3(g,L^{z'}\mathsf d_*)\;\mathsf b_0^{-3}r_*^{-3}\mathsf d_*^{-6}.
\]
\end{enumerate}
\end{theorem}

\begin{theorem}[A close pair seen from the coarser levels]\label{C-352-3}
Let $0<g\le\gamma$ be a reference coupling, $(K,\beta)$ a first-passage datum at $g$ with its run, $\mathsf b=a_K\beta$, and $\ell\in\mathfrak L(g)$. For a level $n$ and $A\in\mathsf{Cc}_n(r_n)$ put $\tilde U=e^{i\eta_nA}$.
\begin{enumerate}
\item[(B.1)] \emph{A pair as wide as its distance.} There are numbers $\gamma_{\mathrm O},C,c,c'>0$ formed from the constants, and a number $r_0(\mathsf b)>0$ formed from these and $\mathsf b$, with $\inf\{r_0(\mathsf b'):g^{-2}<\mathsf b'\le\bar{\mathsf b}\}>0$, such that the following holds. Let $g\le\min\{\gamma_{\mathrm{jt}},\gamma_{\mathrm O}\}$; $0<r\le r_0(\mathsf b)$; $s=s(r)$ with $3\le s\le K$ and $\mathsf{Lm}_s\le L^{s-3}$; $\vartheta\in(0,1)$; $(f_1,f_2)\in\mathfrak T_{T_\ell}(r;\vartheta)$ with $w_a=f_a\circ x\not\equiv0$; $k_0=\max\{K-s,2\}$; and let $q_n$ be formed over the base level $k_0$. Then for every $k_0<n\le K$ and every $A\in\mathsf{Cc}_n(r_n)$
\[
\bigl|\mathsf F_n(\tilde U)-\mathsf F_n(\mathbf 1)-q_n(\mathsf{ang}^{[n]}(A))\bigr|\le\Bigl\{C\mathfrak b_n^2\delta_n+C(\log x_{k_0})^c\bigl[(r/\mathsf b)^{1/2}+e^{-c'Y_s^{2\mathfrak p_0}}\bigr]\Bigr\}\,\hat r_n^3\,\mathsf u_1\mathsf u_2 .
\]
None of $\gamma_{\mathrm O},C,c,c'$ depends on $\vartheta$, $K$, $\ell$ or the level.
\item[(B.2)] \emph{A narrow pair.} Let $g$ satisfy $(\Gamma1)$; let $0<\rho\le r\le e^{-2\sqrt3}$, $s(\rho)\le K$, $\rho\le r_1(\mathsf b)$; let $w_1,w_2\not\equiv0$ be real weights placed at $(\rho,r)$, and let $q_n$ be formed over the base level $\mathsf k(r)$. Then for every $\mathsf k(r)\le n\le K$ and every $A\in\mathsf{Cc}_n(r_n)$
\[
\begin{multlined}
\bigl|\mathsf F_n(\tilde U)-\mathsf F_n(\mathbf 1)-q_n(\mathsf{ang}^{[n]}(A))\bigr|\le\bigl[C_{\mathrm N}\mathfrak b_n^2\delta_n+\epsilon'_{\mathrm N}(r)\bigr]\Bigl(\frac\rho r\Bigr)^6\hat r_n^3\,\mathsf u_1\mathsf u_2,\qquad\epsilon'_{\mathrm N}(r) \\
=C'_{\mathrm N}Y_r^{15\mathfrak p_0+2}\bigl[g\,r^{1/2}+e^{-Y_r^{2\mathfrak p_0}/\sqrt3}\bigr].
\end{multlined}
\]
$C_{\mathrm N},C'_{\mathrm N}$ are formed from the constants, $\epsilon'_{\mathrm N}$ depends on $g$ and $r$ alone; nothing in the bracket depends on $\rho$, $K$, $\ell$, the position, the direction, the sign or the shape of the weights, and $\epsilon'_{\mathrm N}(r)(\log\frac1r)^N\to0$ as $r\to0$ for every $N$.
\item[(B.3)] \emph{The explicit form.} In both cases $\sup_{A\in\mathsf{Cc}_n(r_n)}|q_n(\mathsf{ang}^{[n]}(A))|\le C\,\mathfrak b_n^2(\rho/r)^6\hat r_n^3\mathsf u_1\mathsf u_2$, with $\rho=r$ in (B.1).
\end{enumerate}
\end{theorem}

\begin{theorem}[The pair against the plane fields]\label{C-352-5}
Let $g$ satisfy $(\Gamma)$, let $\mathsf b$ be as in (A.2), and let $\mathsf s_3$ be the kernel of (A.2). For points $c,z$ put $\mathsf d=|c-z|$, let $0<r\le\frac12\mathsf d$, $\mathsf d\le\frac{16}{17}\mathsf d_5(g)$, $\bar\epsilon_{\mathrm A}=\sup\{\epsilon_{\mathrm A}(g,t):\frac34\mathsf d\le t\le\frac{17}{16}\mathsf d\}$, and write $(c\pm\frac r2n,z)$ for the triple $(c+\frac r2n,c-\frac r2n,z)$; $e_\lambda$ is the unit vector of the direction $\lambda$, and $\mathsf M_\rho$ denotes the mean over the balls of radius $\rho$ about the arguments.
\begin{enumerate}
\item[(C.1)] \emph{Against the free harmonic.} (i) For almost every $(c,h,z)\in(T_{\ell_g})^3$, with $h=rn$, $r=|h|$ in the range above,
\[
\bigl|\tfrac{2\pi}3\mathsf br^3\mathsf s_3(c\pm\tfrac r2n,z)-\mathsf t^\circ(n,c-z)\bigr|\le\bigl[12\bar\epsilon_{\mathrm A}+2.7(r/\mathsf d)^2\bigr]\mathsf b^{-2}\mathsf d^{-6}.
\]
(ii) For every $(c,n,r,z)$ in the range above the same holds with $\mathsf s_3(c\pm\frac r2n,z)$ replaced by $\limsup_{\rho\to0}$ of its mean $\mathsf M_\rho\mathsf s_3(c\pm\frac r2n,z)$; and, for every $0<\rho<\min\{\frac14r,r_1(\bar{\mathsf b})\}$, with $\mathsf s_3(c\pm\frac r2n,z)$ replaced by $\int\mathsf s_3\,\chi^\rho_{c+\frac r2n}\otimes\chi^\rho_{c-\frac r2n}\otimes\chi^\rho_z$ and $7900\rho/r$ added in the bracket. (iii) For a law $\sigma'$ as in (A.4)(ii), in particular for every periodic-Gibbs limit at $g$ and every
component of its canonical decomposition, the bound of (i) holds for its continuous kernel at every $(c,r,z)$ and every unit vector $n$, with $12\sup\{\mathsf E_3(g,t):\frac34\mathsf d\le t\le\frac{17}{16}\mathsf d\}$ in place of $12\bar\epsilon_{\mathrm A}$.
\item[(C.2)] \emph{The plane laws.} (i) At a cutoff: let $(K,\beta)$ be a first-passage datum at $g$, $\mathsf b=a_K\beta$, $0<\rho\le\mathsf d\le\mathsf d_4(g)$, $s(\rho)\le K$, $\vartheta\in(0,1]$, and let $\varphi_1,\varphi_2$ be real Lipschitz functions on $T_{\ell_g}$ with $\operatorname{supp}\varphi_a\subset\mathbb B(x_a,\rho)$, $\operatorname{Lip}\varphi_a\le\|\varphi_a\|_\infty/\vartheta\rho$, the balls $\mathbb B(x_a,\rho)$ at a distance $\ge\mathsf d$ and $|x_1-x_2|\le18\mathsf d$; no sign condition. Then for all $\lambda,\mu$
\[
\begin{multlined}
\bigl|\langle\mathcal O_\lambda(\varphi_1);\mathcal O_\mu(\varphi_2)\rangle_{K,\ell_g,\beta}-\tfrac32\mathsf b^{-2}\mathcal K_{\lambda\mu}(\varphi_1,\varphi_2)\bigr|\le\Bigl[\epsilon^{\mathrm{nar}}_2(g,\mathsf d)+3\Lambda_3^2\mathsf M'^2\mathsf b^2\mathsf d^{3/2}+\tfrac32C_{\mathrm{Rs}}\Bigl(\frac{a_K}{\vartheta\rho} \\
+\Bigl(\frac{\mathsf d}{\ell_g}\Bigr)^3\Bigr)\Bigr]\Bigl(\frac\rho{\mathsf d}\Bigr)^6\mathsf b^{-2}\|\varphi_1\|_\infty\|\varphi_2\|_\infty .
\end{multlined}
\]
(ii) Let $(K_j,\beta_j)$ be first-passage data at $g$ with $K_j\to\infty$, $a_{K_j}\beta_j\to\mathsf b$, and $\Omega_{\mathrm P}=\{(x,y)\in T_{\ell_g}^2:0<|x-y|<2\mathsf d_4(g)\}$. There is a subsequence along which the nine forms $\langle\mathcal O_\lambda(\varphi_1);\mathcal O_\mu(\varphi_2)\rangle_{K_j,\ell_g,\beta_j}$ converge for all real Lipschitz $\varphi_1,\varphi_2$ with $\operatorname{supp}\varphi_1\times\operatorname{supp}\varphi_2\subset\Omega_{\mathrm P}$; the nine limits $T_{\lambda\mu}$ along one such subsequence are a limit point of the plane laws. Every limit point has Borel kernels $\tau_{\lambda\mu}$ on $\Omega_{\mathrm P}$ with
\[
\bigl|\tau_{\lambda\mu}(x,y)-\tfrac32\mathsf b^{-2}[\mathbb K_3(x-y)_{\lambda\mu}]^2\bigr|\le\epsilon_{\mathrm P}(g,|x-y|)\,\mathsf b^{-2}|x-y|^{-6}\quad\text{almost everywhere on }\Omega_{\mathrm P};
\]
before a kernel is named, at every $(x_0,y_0)\in\Omega_{\mathrm P}$ with $\mathsf d_0=|x_0-y_0|$ and for Lipschitz $\varphi_1,\varphi_2$ supported in balls of radius $\rho$ inside the balls of radius $\frac18\mathsf d_0$ about $x_0,y_0$: $|T_{\lambda\mu}(\varphi_1,\varphi_2)-\frac32\mathsf b^{-2}\mathcal K_{\lambda\mu}(\varphi_1,\varphi_2)|\le\epsilon_{\mathrm P}(g,\mathsf d_0)\mathsf b^{-2}\mathsf d_0^{-6}\rho^6\|\varphi_1\|_\infty\|\varphi_2\|_\infty$. The kernel $t_\lambda=\sum_\mu(\tau_{\lambda\mu}-\frac13\sum_{\lambda'}\tau_{\lambda'\mu})$ of the limit of $\langle T_\lambda;\mathcal O\rangle$ satisfies $|t_\lambda(x,y)-\mathsf t^\circ(e_\lambda,x-y)|\le6\epsilon_{\mathrm P}(g,|x-y|)\mathsf b^{-2}|x-y|^{-6}$ almost everywhere.
\item[(C.3)] \emph{The pair against the kernel of $\langle T_\lambda;\mathcal O\rangle$.} Take one limit point as in (C.2)(ii), along a subsequence of $(K_j,\beta_j)$; along it $S^T_{3;K_j,\ell_g,\beta_j}\to S^{(\ell_g,\mathsf b)}_3$. For every $c,z,r,\lambda$ in the range above
\[
\limsup_{\rho\to0}\bigl|\tfrac{2\pi}3\mathsf br^3\mathsf M_\rho\mathsf s_3(c\pm\tfrac r2e_\lambda,z)-\mathsf M_\rho t_\lambda(c,z)\bigr|\le\bigl[12\bar\epsilon_{\mathrm A}+6\epsilon_{\mathrm P}(g,\mathsf d)+2.7(r/\mathsf d)^2\bigr]\mathsf b^{-2}\mathsf d^{-6},
\]
and the same with the radial tents of mass $1$ in place of the means. For almost every $(c,h,z)$, $h=rn$, in the range above,
\[
\bigl|\tfrac{2\pi}3\mathsf br^3\mathsf s_3(c\pm\tfrac r2n,z)-t_\lambda(c,z)\bigr|\le\bigl[12\bar\epsilon_{\mathrm A}+6\epsilon_{\mathrm P}(g,\mathsf d)+2.7(r/\mathsf d)^2+0.057|n-e_\lambda|\bigr]\mathsf b^{-2}\mathsf d^{-6}.
\]
\item[(C.4)] \emph{Limits and components on $\mathbb R^3$.} Let $(K_j,\beta_j)$, $\nu_j$, $\sigma'$ be as in (A.4)(ii), now with $S^T_n[\nu_j]\to S^{\sigma'}_n$ for $n=2,3$, assume that $S^{\sigma'}_3$ and $S^{\sigma'}_2$ have continuous
kernels $\mathsf s_3$ and $\mathsf s^T_2$ on distinct configurations, and $\Omega_\sigma=\{(x,y)\in(\mathbb R^3)^2:0<|x-y|<\mathsf d_5(g)\}$. (a) There is a subsequence along which $\langle\mathcal O_\lambda(\varphi_1);\mathcal O_\mu(\varphi_2)\rangle_{\nu_j}$ converges for all $\lambda,\mu$ and all real Lipschitz $\varphi_1,\varphi_2$ with compact supports, $\operatorname{supp}\varphi_1\times\operatorname{supp}\varphi_2\subset\Omega_\sigma$. (b) For every such subsequence the limits have Borel kernels $\tau^\sigma_{\lambda\mu}$ on $\Omega_\sigma$ with $|\tau^\sigma_{\lambda\mu}(x,y)-\frac32\mathsf b^{-2}[\mathbb K_3(x-y)_{\lambda\mu}]^2|\le(\frac3\pi)^2\epsilon_\sigma(g,|x-y|)\mathsf b^{-2}|x-y|^{-6}$ almost everywhere; $\tau^\sigma_{\lambda\mu}(x,y)$ depends on $x-y$ alone (almost everywhere); and $\sum_{\lambda,\mu}\tau^\sigma_{\lambda\mu}=\mathsf s^T_2$ almost everywhere, so that $|\mathsf s^T_2(x,y)-\frac9{16\pi^2}\mathsf b^{-2}|x-y|^{-6}|\le9\epsilon_\sigma(g,|x-y|)\mathsf b^{-2}|x-y|^{-6}$ at every $(x,y)\in\Omega_\sigma$. (c) With $t^\sigma_\lambda=\sum_\mu(\tau^\sigma_{\lambda\mu}-\frac13\sum_{\lambda'}\tau^\sigma_{\lambda'\mu})$ there is one null set $\mathcal N\subset\Omega_\sigma$, fixed before $r$ is named, such that for every $(c,z)\notin\mathcal N$ with $0<\mathsf d=|c-z|\le\frac12\mathsf d_5(g)$, every $0<r\le\frac12\mathsf d$ and every $\lambda$
\[
\begin{multlined}
\bigl|\tfrac{2\pi}3\mathsf br^3\mathsf s_3(c\pm\tfrac r2e_\lambda,z)-t^\sigma_\lambda(c,z)\bigr|\le\bigl[\epsilon'_\sigma(g,\mathsf d)+2.7(r/\mathsf d)^2\bigr]\mathsf b^{-2}\mathsf d^{-6},\qquad\epsilon'_\sigma(g,\mathsf d) \\
=12\!\sup_{\frac34\mathsf d\le t\le\frac{17}{16}\mathsf d}\!\mathsf E_3(g,t)+6\epsilon_\sigma(g,\mathsf d),
\end{multlined}
\]
and $\epsilon'_\sigma(g,\mathsf d)\le C(g)(\log\frac1{\mathsf d})^{39\mathfrak p_0+2}\mathsf d^{1/2}$; at every $(c,z)$ in this range the same holds with $\sum_\mu(\mathsf t_{\lambda\mu}-\frac13\sum_{\lambda'}\mathsf t_{\lambda'\mu})(\chi^\rho_c,\chi^\rho_z)$ in place of $t^\sigma_\lambda(c,z)$, $\mathsf t_{\lambda\mu}$ the limit forms, for every $0<\rho<\frac18\mathsf d$, with $21\rho/\mathsf d$ added in the bracket. (d) At a cutoff, for every first-passage datum $(K,\beta)$ at $g$, every state $\mathfrak s$ at $(K,\beta)$, points $x_0,y_0$ with $0<\mathsf d_0=|x_0-y_0|\le2\mathsf d_4(g)$, $\mathsf d=\frac12\mathsf d_0$, $s(\rho)\le K$, and bounded real Borel $f_1,f_2$ on $\mathbb R^3$ supported in balls of radius $\rho$ inside the closed balls of radius $\frac18\mathsf d_0$ about $x_0,y_0$:
\[
\begin{multlined}
\Bigl|\langle\mathcal O_\lambda(f_1);\mathcal O_\mu(f_2)\rangle_{\mathfrak s}-\tfrac32\mathsf b^{-2}\!\!\sum_{\lambda(p)=\lambda,\,\lambda(q)=\mu}\!\!f_1(x(p))f_2(x(q))\,C^{\mathrm{free}}_K(\mathrm{pr}\,p,\mathrm{pr}\,q)^2\Bigr|\le\Bigl[\mathsf E_2(g,\mathsf d) \\
-\tfrac32C_{\mathrm{Rs}}\Bigl(\frac{\mathsf d}{\ell_g}\Bigr)^3\Bigr]\Bigl(\frac\rho{\mathsf d}\Bigr)^6\mathsf b^{-2}\|f_1\|_\infty\|f_2\|_\infty,
\end{multlined}
\]
with $C^{\mathrm{free}}_K$ of the lattice $(\ell_g,K)$ and $\mathrm{pr}$ the projection onto $T_{\ell_g}$.
\end{enumerate}
\end{theorem}

\begin{theorem}[The traceless tensor in every direction]\label{C-352-6}
Let $0<g\le\gamma$ be a reference coupling and $\ell\in\mathfrak L(g)$. There are an integer $K_\circ(g)$ and a number $\mathsf d_{\mathrm{cl}}(g)>0$ formed from the constants and $g$, and a finite function $\epsilon^{\mathrm{cl}}(g,\ell,\cdot)$ on $(0,\frac18\ell]$ formed from the constants, $g$ and $\ell$, with $\epsilon^{\mathrm{cl}}(g,\ell,\mathsf d')\le C(g,\ell)(\log\frac1{\mathsf d'})^{19\mathfrak p_0}\mathsf d'^{1/2}$ for $\mathsf d'\le\mathsf d_{\mathrm{cl}}(g)$, such that the following hold.
\begin{enumerate}
\item[(D.1)] \emph{The bond forms against the plane fields, on every torus.} Let $(K,\beta)$ be a first-passage datum at $g$ with $K\ge K_\circ(g)$, $\mathsf b=a_K\beta$; $0<16r\le\mathsf d'\le\frac18\ell$ with $s(r)\le K$; $x_1,x_2\in T_\ell$ with $|x_1-x_2|=\mathsf d'$; $\vartheta\in(0,1]$; $\varphi_1,\varphi_2$ real Lipschitz functions on $T_\ell$ with $\operatorname{supp}\varphi_a\subset\mathbb B(x_a,r)$ and $\operatorname{Lip}\varphi_a\le\|\varphi_a\|_\infty/\vartheta r$. Then for all $\lambda\ne\lambda'$ and all $\gamma,\mu$
\begin{align*}
&\Bigl|\langle\mathcal O^{\mathrm{od}}_{\lambda\lambda'}(\varphi_1);\mathcal O_\gamma(\varphi_2)\rangle_{K,\ell,\beta}\\
&\qquad-3\mathsf b^{-2}\!\iint\!\varphi_1(x)\varphi_2(y)\mathbb K_3(x-y)_{\lambda\gamma}\mathbb K_3(x-y)_{\lambda'\gamma}\,dx\,dy\Bigr|\\
&\quad\le\Bigl[8\epsilon^{\mathrm{cl}}(g,\ell,\mathsf d')+3C^{\mathrm{od}}_{\mathrm{PS}}\frac{a_K}{\vartheta r}
+3C^{\mathrm{od}}_{\mathrm{PS}}\Bigl(\frac{\mathsf d'}\ell\Bigr)^3\Bigr]\\
&\qquad\Bigl(\frac r{\mathsf d'}\Bigr)^6\mathsf b^{-2}\|\varphi_1\|_\infty\|\varphi_2\|_\infty,\\
&\bigl|\langle\mathcal O_\lambda(\varphi_1);\mathcal O_\mu(\varphi_2)\rangle_{K,\ell,\beta}
-\tfrac32\mathsf b^{-2}\mathcal K_{\lambda\mu}(\varphi_1,\varphi_2)\bigr|\\
&\quad\le\Bigl[\epsilon^{\mathrm{cl}}(g,\ell,\mathsf d')+\tfrac32C_{\mathrm{PS}}\frac{a_K}{\vartheta r}
+\tfrac32C_{\mathrm{PS}}\Bigl(\frac{\mathsf d'}\ell\Bigr)^3\Bigr]\\
&\qquad\Bigl(\frac r{\mathsf d'}\Bigr)^6\mathsf b^{-2}\|\varphi_1\|_\infty\|\varphi_2\|_\infty .
\end{align*}
\item[(D.2)] \emph{The kernels of $\langle T_n;\mathcal O\rangle$ for all directions at once.} Let $(K_j,\beta_j)$ be first-passage data at $g$ with $K_j\to\infty$ and $a_{K_j}\beta_j\to\mathsf b$; $\Omega^\circ=\{(x,y)\in T_\ell^2:0<|x-y|<\frac79\mathsf d_{\mathrm{cl}}(g)\}$; 
\[
\bar\epsilon^{\mathrm{cl}}(t)=\sup\{8\epsilon^{\mathrm{cl}}(g,\ell,\mathsf d'')+3\max\{C_{\mathrm{PS}},C^{\mathrm{od}}_{\mathrm{PS}}\}(\mathsf d''/\ell)^3:\frac79t\le\mathsf d''\le\frac97t\},
\]
so that $\bar\epsilon^{\mathrm{cl}}(t)\le C(g,\ell)[(\log\frac1t)^{19\mathfrak p_0}t^{1/2}+(t/\ell)^3]$. Along a subsequence all the forms $\langle\mathcal O_\lambda(f);\mathcal O_\mu(h)\rangle_{K_j,\ell,\beta_j}$ and $\langle\mathcal O^{\mathrm{od}}_{\lambda\lambda'}(f);\mathcal O_\gamma(h)\rangle_{K_j,\ell,\beta_j}$ converge, at all real Lipschitz $f,h$ with $\operatorname{supp}f\times\operatorname{supp}h\subset\Omega^\circ$.

For every such subsequence the limits have Borel kernels on $\Omega^\circ$ that lie, almost everywhere, within $4.2\,\bar\epsilon^{\mathrm{cl}}(|x-y|)\mathsf b^{-2}|x-y|^{-6}$ of the free kernels $\frac32\mathsf b^{-2}[\mathbb K_3(x-y)_{\lambda\mu}]^2$ and $3\mathsf b^{-2}\mathbb K_3(x-y)_{\lambda\gamma}\mathbb K_3(x-y)_{\lambda'\gamma}$; and outside one null set of $\Omega^\circ$, for all unit vectors $n$ at once, the kernel $t_n$ of the limit of $\langle T_n(f);\mathcal O(h)\rangle$ satisfies
\[
|t_n(x,y)-\mathsf t^\circ(n,x-y)|\le38\,\bar\epsilon^{\mathrm{cl}}(|x-y|)\,\mathsf b^{-2}|x-y|^{-6}.
\]
\item[(D.3)] \emph{The pair against $t_n$.} Let moreover $g$ satisfy $(\Gamma)$ and $\ell=\ell_g$, and take $\mathsf s_3$ of (A.2) and a limit point of (D.2). For every $c,z,r,n$ with $0<r\le\frac12\mathsf d$, $\mathsf d=|c-z|\le\frac{16}{17}\mathsf d_5(g)$ and $\mathsf d<\frac79\mathsf d_{\mathrm{cl}}(g)$,
\[
\limsup_{\rho\to0}\bigl|\tfrac{2\pi}3\mathsf br^3\mathsf M_\rho\mathsf s_3(c\pm\tfrac r2n,z)-\mathsf M_\rho t_n(c,z)\bigr|\le\bigl[12\bar\epsilon_{\mathrm A}+38\bar\epsilon^{\mathrm{cl}}(\mathsf d)+2.7(r/\mathsf d)^2\bigr]\mathsf b^{-2}\mathsf d^{-6},
\]
the same with the radial tents of mass $1$ in place of the means, and the same bound holds for $|\frac{2\pi}3\mathsf br^3\mathsf s_3(c\pm\frac r2n,z)-t_n(c,z)|$ at almost every $(c,rn,z)$ in this range.
\end{enumerate}
\end{theorem}

\begin{theorem}[Pair minus tensor field; the limit $r\to0$]\label{C-352-7}
Fix $\beta_{\mathrm H}\in(0,1)$. There is a number $\gamma_{\mathrm R}\in(0,\gamma_{\mathrm{jt}}]$ formed from the constants and $\beta_{\mathrm H}$ such that every $0<g\le\gamma_{\mathrm R}$ satisfies $(\Gamma)$ and $g\le\gamma_{\mathrm O}$ of (B.1), and besides the two conditions $\frac32C_\natural c_\gamma g\le\log2$, where $C_\natural=(L^8+8c_3)L^4$ with $c_3$ of Lemma~\ref{C-352-12}(c) and $c_\gamma$ of Lemma~\ref{C-352-12}(f), and $c_{\mathrm A}\alpha_0\ge32M\mathsf c_{\mathfrak h}\delta(t)(1+\log t)$ for all $t\ge g^{-2}$, which contains the condition $c_{\mathrm A}\alpha_0\ge5M\mathsf c_{\mathfrak h}\delta(t)(1+\log t)$ for all $t\ge g^{-2}$, the last member of the list of $\gamma_{\mathrm N}$ in Definition~\ref{C-352-2}(1), together with $g\le19C_{\mathrm A}C_{\mathrm E,3}M^2\mathsf c_{\mathfrak h}^2A_1^2/2C_{\mathrm E,1}\alpha_0^2$, with the numbers of Definition~\ref{C-352-2}(6); each of them is a condition ``$g\le$ a number formed from the constants'', and together they contain the ceilings of all the lemmas used in the proof.

The following holds for $0<g\le\gamma_{\mathrm R}$ on the least torus $\ell=\ell_g$. Let $(K,\beta)$ be a first-passage datum at $g$, $\mathsf b=a_K\beta$, $\lambda$ a direction, $\mathsf d>0$ with $6\mathsf d\le\min\{1,\mathsf b\}$, and $\varphi_1,\varphi_2$ bounded real functions on $T_{\ell_g}$ with $\operatorname{supp}\varphi_a\subset\mathbb B(y_a,\mathsf d)$ and $\operatorname{dist}(\operatorname{supp}\varphi_1,\operatorname{supp}\varphi_2)\ge\mathsf d$. Let $r_\alpha(g,\mathsf b)>0$ be the number of Definition~\ref{C-352-2}(5), formed from the constants, $g$ and $\mathsf b$. For every $r$ with
\[
\begin{multlined}
128a_K\le r\le\min\{r_\alpha(g,\mathsf b),(\mathsf d/5LM)^2\},\qquad3\le s(\tfrac r2)\le K,\qquad\tfrac r2\le\min\{r_0(\mathsf b),r_1(\mathsf b)\}, \\
\mathsf{Lm}_{s(r/2)}\le L^{s(r/2)-3},
\end{multlined}
\]
the following holds. For $c\in\operatorname{supp}\varphi_1$ let $w_{1,2}=\chi^{r/8}_{c\pm\frac r2e_\lambda}\circ x$ and $w_T(p)=(\delta_{\lambda(p)\lambda}-\frac13)\chi^{r/8}_c(x(p))$, so that $\mathcal O(w_T)=T_\lambda(\chi^{r/8}_c)$; let $\mathsf c=\frac{2\pi}3\mathsf br^3$; let $w_3$ be $\varphi_2\circ x$ restricted to the plaquettes with barycentres in a cube of side $\frac r4$ whose centre $z$ has $|c-z|\le6\mathsf d$, and $\mu_3=r^3\|\varphi_2\|_\infty$. Let $\mathsf Y_\beta$ be the remainders of Lemma~\ref{C-352-25} for the lists $(w_1,w_2,w_3)$ and $(w_T,w_3)$, $\mathsf Y_\Delta=\mathsf c\mathsf Y_{12}-\mathsf Y_T$ and $\mathsf Y_{\Delta3}=\mathsf c\mathsf Y_{123}-\mathsf Y_{T3}$.
\begin{enumerate}
\item[(E.1)] \emph{The part that is not explicit.} There is a finite function $\mathsf C_{\mathrm R}(g,\mathsf b;\cdot)$ on $(0,1]$, formed from the constants, $\beta_{\mathrm H}$, $g$ and $\mathsf b$, such that
\[
\Bigl|\int\mathsf Y_{\Delta3}\,d\mathbb P+\kappa_{\mathbb P}(\mathsf Y_\Delta,\mathsf Y_3)\Bigr|\le\mathsf C_{\mathrm R}(g,\mathsf b;\mathsf d)\bigl[r^{\beta_{\mathrm H}/72}+(a_K/r)^{1/18}\bigr]\mathsf b^{-2}\mathsf d^{-6}\mu_3 .
\]
If $g$ satisfies besides $(\Gamma3)$ of Definition~\ref{C-352-2}(6), one may take $\mathsf C_{\mathrm R}(g,\mathsf b;\mathsf d)=C(g,\mathsf b)\,\mathsf d^{-5/2}$, and the pair minus the tensor field is nearly constant up to the last level: with $\Delta_n=\mathsf cG_{n;12}-G_{n;T}$ and $\Gamma_n=\mathsf cG_{n;123}-G_{n;T3}$, formed with one progression of the four sources, with $\varepsilon^{\mathcal W}=C(g,\mathsf b)[r^{\beta_{\mathrm H}/72}+(a_K/r)^{1/18}]$, and for $r\le r^{\mathcal W}_\alpha(g,\mathsf b)$, there are numbers $\mathsf c_n$ such that at every level $n\le K$ with $a_{K-n}\ge r^{1/2}$, for $\mathsf m_n$-almost every small-field background $W$ of the level $n$ (Lemma~\ref{C-352-25}),
\[
\begin{multlined}
|\Delta_n(U_n(W))-\mathsf c_n|\le\varepsilon^{\mathcal W}\varepsilon_n^2a_{K-n}^{-4},\qquad|\mathsf c_n-\mathsf c_K|\le\varepsilon^{\mathcal W}\varepsilon_n^2a_{K-n}^{-4},\qquad|\Gamma_n(U_n(W)) \\
-\Gamma_K(\mathbf 1)|\le C(\mathsf b)\,\mathsf d^{-5/2}\varepsilon^{\mathcal W}\mathsf b^{-2}\mathsf d^{-6}\mu_3 .
\end{multlined}
\]
\item[(E.2)] \emph{The true functions.} There are numbers $C_{\mathrm A}(g,\mathsf b)$, formed from the constants, $\beta_{\mathrm H}$, $g$ and $\mathsf b$, and $r_{\mathcal W}(g,\mathsf b,\mathsf d)>0$, formed besides with $\mathsf d$, such that for $r\le r_{\mathcal W}(g,\mathsf b,\mathsf d)$
\[
\bigl|\mathsf c\,S^T_3(w_1,w_2,w_3)-S^T_2(w_T,w_3)\bigr|\le C_{\mathrm A}(g,\mathsf b)\,\mathsf d^{-18}\bigl[r^{\beta_{\mathrm H}/72}+(a_K/r)^{1/18}\bigr]\mu_3 .
\]
\item[(E.3)] \emph{At a fixed label: the limit in the cutoff, and $r\to0$.} Let $\varphi_1,\varphi_2$ be as above, Lipschitz, with $|y_1-y_2|\le3\mathsf d$, let $f^c_\pm=\chi^{r/8}_{c\pm\frac r2e_\lambda}$, and for a first-passage datum $(K,\beta)$ at $g$ with the label $\mathsf b'=a_K\beta$ put
\[
\mathsf A_{r,K}=\tfrac{2\pi}3\mathsf b'r^3\int dc\,\varphi_1(c)\,S^T_{3;K,\ell_g,\beta}(f^c_+,f^c_-,\varphi_2).
\]
If either $\varphi_1$ or $\varphi_2$ is zero, the assertions are immediate by multilinearity; assume both are nonzero in the normalized estimates that follow. At every label $\mathsf b'\in(g^{-2},\bar{\mathsf b}]$ there are numbers $r_\sharp>0$ and $\mathsf C(\mathsf d)$, formed from the constants, $\beta_{\mathrm H}$, $g$, $\mathsf b'$, $\mathsf d$ and $\mathsf d\operatorname{Lip}\varphi_1/\|\varphi_1\|_\infty$, such that: (i) for $0<r\le r_\sharp$ and every first-passage datum $(K,\beta)$ at $g$ with the label $\mathsf b'$, $128a_K\le r$ and $s(\frac r2)\le K$, $|\mathsf A_{r,K}-\langle T_\lambda(\varphi_1);\mathcal O(\varphi_2)\rangle_{K,\ell_g,\beta}|\le\mathsf C(\mathsf d)[r^{\beta_{\mathrm H}/72}+(a_K/r)^{1/18}]\mathsf b'^{-2}\|\varphi_1\|_\infty\|\varphi_2\|_\infty$. Let now $\mathsf b\in I(g)$ and $\beta_K=L^K\mathsf b$, so that $(K,\beta_K)$ is a first-passage datum at $g$ with the label $\mathsf b$ for every $K\ge1$; let $\mathsf A_{r,K}$ be formed at $(K,\beta_K)$, and $r_\sharp$, $\mathsf C(\mathsf d)$ at the label $\mathsf b$.

Then (ii) the limit $t_\lambda(\varphi_1,\varphi_2)=\lim_{K\to\infty}\langle T_\lambda(\varphi_1);\mathcal O(\varphi_2)\rangle_{K,\ell_g,\beta_K}$ exists; (iii) 
\[
\mathsf A_r=\lim_K\mathsf A_{r,K}=\frac{2\pi}3\mathsf br^3\int dc\,\varphi_1(c)S^{(\ell_g,\mathsf b)}_3(f^c_+,f^c_-,\varphi_2)
\]
 exists and 
\[
|\mathsf A_r-t_\lambda(\varphi_1,\varphi_2)|\le\mathsf C(\mathsf d)r^{\beta_{\mathrm H}/72}\mathsf b^{-2}\|\varphi_1\|_\infty\|\varphi_2\|_\infty
\]
 for $0<r\le r_\sharp$; so $\mathsf A_r\to t_\lambda(\varphi_1,\varphi_2)$ as $r\to0$ at a fixed $\mathsf d$.
\item[(E.4)] \emph{Labels that tend to $\mathsf b$.} For $\mathsf b\in I(g)$, the $\varphi_a$ of (E.3) and every sequence of first-passage data $(K_j,\beta_j)$ at $g$ with $K_j\to\infty$ and $a_{K_j}\beta_j\to\mathsf b$: $\langle T_\lambda(\varphi_1);\mathcal O(\varphi_2)\rangle_{K_j,\ell_g,\beta_j}\to t_\lambda(\varphi_1,\varphi_2)$.
\end{enumerate}
\end{theorem}

\begin{theorem}[Four points at two scales]\label{C-352-8}
Let $g$ satisfy $(\Gamma^{(4)})$, let $(K,\beta)$ be a first-passage datum at $g$, $\mathsf b=a_K\beta$, $0<\rho\le r\le\mathsf d\le\mathsf d^{(4)}_4(g)$, $s(\rho)\le K$, and let $w_1,\dots,w_4$ form two pairs at $(\rho,r,\mathsf d)$.
\begin{enumerate}
\item[(a)] On every side $\ell\in\mathfrak L(g)$ with $64\mathsf d\le\ell$,
\[
\Bigl|\mathsf J_4-\tfrac32\sum_\varsigma\operatorname{Tr}\prod_l\bigl(w_{\varsigma(l)}\mathsf b^{-1}C^{\mathrm{free}}_K\bigr)\Bigr|\le\epsilon^{\mathrm{nar}}_4(g,\mathsf d)\,\mathrm{nat}_4,
\]
the sum over the six cyclic orders of $\{1,2,3,4\}$; on the least torus 
\[
|S^T_4-\mathsf J_4|\le77\Lambda_3^4\mathsf M'^4\mathsf b^4\mathsf d^{3/2}\mathrm{nat}_4;
\]
 and in every state at $(K,\beta)$, for bounded real Borel functions on $\mathbb R^3$ in the geometry of (A.4)(i) for two pairs, the connected function differs from that of the least torus at the periodised functions by at most 
\[
C^{\mathrm{tr}3}_4(g)\mathsf b^4\mathsf d^{3/2}\mathrm{nat}_4.
\]

\item[(b)] If $w_a=f_a\circ x$ with real Lipschitz $f_a$ on $T_{\ell_g}$ such that $\operatorname{supp}f_a\subset\mathbb B(x_a,\rho)$, $x_a$ the centres with which the $w_a$ form the two pairs (so that $|x_1-x_2|,|x_3-x_4|\le8r$ and $|x_a-x_b|\le18\mathsf d$ for $a\le2<b$), the balls $\mathbb B(x_a,\rho)$ are at least $r$ apart in $\{1,2\}$ and in $\{3,4\}$ and at least $\mathsf d$ apart across, $\operatorname{Lip}f_a\le\|f_a\|_\infty/\vartheta\rho$, $\vartheta\in(0,1]$, and all supports lie in one ball of radius $\frac12\mathsf d^+$ with $4\mathsf d^+<\ell_g$, then
\[
\Bigl|S^T_4(f)-3\mathsf b^{-4}\!\int k_4\prod_af_a\Bigr|\le\Bigl[\epsilon^{\mathrm{nar}}_4(g,\mathsf d)+77\Lambda_3^4\mathsf M'^4\mathsf b^4\mathsf d^{3/2}+3C^{(4)}\Bigl(\frac{a_K}{\vartheta\rho}+\Bigl(\frac{\mathsf d^+}{\ell_g}\Bigr)^3\Bigr)\Bigr]\mathrm{nat}_4 .
\]
Here $8a_K\le\rho$, by the working level at $\varrho=\rho$, and $\mathsf d\le\mathsf d^+$, the supports of $f_1$ and $f_3$ being at least $\mathsf d$ apart in a ball of diameter $\mathsf d^+$; if an $f_a$ vanishes, both members are $0$.
\item[(c)] Let $\mathsf b$ be a limit of labels of first-passage data at $g$ and $\mathcal D_4=\{(x_1,\dots,x_4):r_*\le|x_1-x_2|,|x_3-x_4|\le2r_*,\ r_*\le\mathsf d_*=\operatorname{dist}(\{x_1,x_2\},\{x_3,x_4\})\le\mathsf d^{(4)}_5(g)\}$. Then $S^{(\ell_g,\mathsf b)}_4$ has on $\mathcal D_4$ a Borel kernel $\mathsf s_4$ with, almost everywhere on $\mathcal D_4$ and uniformly in $r_*/\mathsf d_*\in(0,1]$,
\[
\begin{multlined}
|\mathsf s_4-3\mathsf b^{-4}k_4|\le\epsilon_4(g,\mathsf d_*)\,\mathsf b^{-4}r_*^{-6}\mathsf d_*^{-6},\qquad\epsilon_4(g,t)=4096\bigl[\epsilon^{\mathrm{nar}}_4(g,\tfrac t2)+77\Lambda_3^4(\bar{\mathsf M}^{(24)})^4\bar{\mathsf b}^4t^{3/2} \\
+3C^{(4)}(12t/\ell_g)^3\bigr].
\end{multlined}
\]
\end{enumerate}
\end{theorem}

\begin{remark}[The ceilings]\label{C-352-9}
Let $x_W$ be the threshold of Section~\YMMref{126}{843},
equation~\YMMref{(2326)}{845}, and put
\begin{equation}\label{eq:pair-threshold}
 x_{\rm jt}=\max\left\{x_W,
 \left(\frac{20M\mathsf c_{\mathfrak h}A_1
                   (4\mathfrak p_0)^{\mathfrak p_0}}
                  {c_A\alpha_0}\right)^4\right\},
 \qquad\gamma_{\rm jt}=x_{\rm jt}^{-1/2}.
\end{equation}
The notation $x_{\mathrm{jt}}$, $\gamma_{\mathrm{jt}}$ denotes these
same numbers. The defining inequalities for $x_W$ include those for
$x_{\rm J}$, so
$x_{\rm jt}\ge x_W\ge\max\{\gamma^{-2},x_{\rm J}\}$.
For $t>1$, the inequality $\log u\le u$, applied to
$u=t^{1/(4\mathfrak p_0)}$, gives
\[
 \delta(t)=A_1t^{-1/2}(\log t)^{\mathfrak p_0}
       \le A_1(4\mathfrak p_0)^{\mathfrak p_0}t^{-1/4}.
\]
Consequently $t\ge x_{\rm jt}$ implies
$5\mathsf c_{\mathfrak h}\delta(t)\le c_A\alpha_0/(4M)$.
Also $\gamma_W^{\rm ref}=x_W^{-1/2}$ and
$\gamma_W^{\rm sd}\le\gamma_W^{\rm ref}$ by
Section~\YMMref{126}{843}, equation~\YMMref{(2327)}{846}.
Hence $\gamma_{\rm jt}\le\gamma_W^{\rm ref}\le\gamma$ and
$\gamma_W^{\rm sd}\le\gamma$.
Every coupling under $(\Gamma1)$, $(\Gamma2)$, $(\Gamma)$ or
$(\Gamma^{(4)})$, or below $\gamma_{\mathrm N}$ or
$\gamma_{\mathrm R}$, therefore satisfies $0<g\le\gamma$.
For $g\le\gamma_{\rm jt}$ one also has
$g^{-2}\ge e^{2\mathfrak p_0}$. All these ceilings are positive
numbers depending only on the fixed constants.
\end{remark}

\subsection{Construction estimates}

We use the construction results at their stated hypotheses, including Theorems~\YMMref{D}{8}, \YMMref{AJ}{673} and~\YMMref{AD}{607}. The finite construction of Theorem~\YMMref{D}{8} and its sourced form, Theorem~\YMMref{AJ}{673}, enters wherever the objects of the theory (the run of a first-passage datum, the effective actions of all levels with their bounds, and the weighted progression with any number of sources) satisfy the actual-step hypotheses of Proposition~\ref{prop:weighted-canonical-history}; this holds for $0<g\le\gamma$ on every side of $\mathfrak L(g)$, with $g^{-2}\ge x_{\mathrm J}$ where sources are present. Proposition~\ref{prop:weighted-finite-system} constructs the weighted progression for sources on the bare plaquettes; mixed sources enter through clauses (5) and (6) of the lemma on bond sources, which follows Lemma~\ref{C-352-36}: at $0<g\le\gamma_{\mathrm N}$, on every side of $\mathfrak L(g)$, the progression and the representation of the torus expectations are formed by the same construction, with the weight member of the mixed source in the source slot of the bare arrays; the frame of the pair functional is read with the hulls of the sources in the place of their supports; and the tower identity holds for a pair of mixed sources in that frame.

Theorems~\YMMref{AJ}{673} and~\YMMref{AD}{607} enter through the two-sided stability bound of the torus measures, which together with their reflection positivity gives the chessboard bound of Lemma~\ref{C-352-30}. Through this bound they are used in (A.4), (A.5), (C.1)(iii), (C.4) and for the states in Theorem~\ref{C-352-8}(a); they are not used in Theorem~\ref{C-352-3}, in (A.1)--(A.3), (C.1)(i), (ii), (C.2), (C.3), in Theorem~\ref{C-352-6}, or in Theorem~\ref{C-352-8} on the tori. Theorem~\ref{C-352-7}, like every statement of this section, assumes the construction inputs just listed.

We use the fixed construction setting just specified, the torus limits of Theorem~\YMMref{A}{5}, and the periodic-Gibbs limits and decompositions of Theorems~\YMMref{AP}{780} and~\YMMref{AL}{746}. On one renormalisation step (whose provision is that of Section~\YMMref{52.4}{422}) and on the progression: Proposition~\ref{prop:source-entry}, with the fourth-family norm \eqref{eq:source-entry-norm}; the bracket localisation of the families with its count; the locality and linear-weight dependence of Proposition~\ref{prop:weighted-finite-system}; the chart and normalized real measures of Lemmas~\ref{lem:sf-chart-integral} and~\ref{lem:weighted-step-integrals}, with the cumulant bounds of Lemma~\ref{C-352-20}; the expansion $|a_p(e^{i\eta_iA})-\frac12|\mathsf{ang}^{[i]}_p(A)|^2|\le\mathsf M_{\mathrm{pl}}^2\eta_{i-1}^4(r_i/\varrho_{\mathrm{ch}})^3$ on the chain classes, and $|\mathsf{ang}^{[i]}_p(A)|\le2\eta_i^2r_i$; the regular-field evaluation and domain sums of Lemmas~\ref{lem:weighted-regular-family} and~\ref{lem:source-support-sum}; the lowering maps of the chain classes, by which $A'\in\mathsf{Cc}_{i+1}(r_{i+1})$ and a fluctuation field $B$ give $A^{\mathrm{lw}}(B)\in\mathsf{Cc}_i(r_i)$, and the pair tower identity of Proposition~\ref{prop:weighted-pair-tower}, extended to several sources in the lemma preceding Proposition~\ref{C-352-21}; the history decomposition \eqref{eq:ward-first-creation}, its sourced quotients in Proposition~\ref{prop:running-source-density}, and the tail bounds of Proposition~\ref{prop:ward-fixed-volume-rarity}.

On the free kernel, Proposition~\YMMref{120.6}{821} and equations~\YMMref{(2245)--(2246)}{824} give the decay of the slices $|\mathcal C_{(j)}(p,q)|\le C_{\mathrm{tp}}a_K^3\mathsf r^{-3}t_j^3e^{-\delta_{\mathrm{tp}}t_j}$, $\mathsf r=|x(p)-x(q)|_\infty$, $t_j=\mathsf r/a_{K-j+1}$; and the comparison $|\epsilon_p\epsilon_qC^{\mathrm{free}}_K(p,q)-a_K^3\mathbb K_3(u)_{\lambda(p)\lambda(q)}|\le C_{\mathrm{lat}}a_K^3[a_K^2|u|^{-5}+\ell^{-3}]$, $|C^{\mathrm{free}}_K(p,q)|\le(1+3C_{\mathrm{lat}})a_K^3|u|^{-3}$, for $p\ne q$, $u$ the least lift of $x(p)-x(q)$, $|u|\ge2^{-1/2}a_K$, $\epsilon_p=\pm1$ the orientation signs. On the working level (Lemma~\YMMref{126.1}{844}): for $0<\varrho<1$ with $s=s(\varrho)\le K$ and $k=\mathsf k(\varrho)$, $Y=Y_s$: $s\ge3$, $2\le k\le K-1$, $4MY^{2\mathfrak p_0}\le\hat\varrho_k<LC_{W,0}A_0^2Y^{2\mathfrak p_0}$, $Y\ge\log\frac1\varrho$, $y_k\le Y$, $K-k\le s\le Y/\log L$; in particular $a_K\le a_{K-k}=\varrho/\hat\varrho_k\le\varrho/4M$, and $Y\ge2\sqrt3$ when $\varrho\le e^{-2\sqrt3}$.

The lower bounds come from the definition of $s(\varrho)$ and from $C_{W,0}A_0^2\ge4L^2M$: $L^s\varrho\ge C_W(s)\ge4L^2MY^{2\mathfrak p_0}>L^2$ and $\varrho<1$ give $s\ge3$; $K-k=\min\{s,K-2\}\ge s-2$ gives $\hat\varrho_k=L^{K-k}\varrho\ge L^{-2}C_W(s)\ge4MY^{2\mathfrak p_0}$; and $L^s\varrho\ge C_W(s)\ge1$ gives $Y\ge s\log L\ge\log\frac1\varrho$. Besides, $L^m\ge LMy_K$ for the sides $\ell=2L^m\in\mathfrak L(g)$. Finally the dilation of lattices: $x\mapsto L^zx$ carries the lattice $(\ell,K)$ to the lattice $(L^z\ell,K-z)$ and $\mu_{K,\ell,\beta}$ to $\mu_{K-z,L^z\ell,\beta}$.

\subsection{Estimates and proofs}

The argument has eight parts: the calculus of the derivative families (Lemmas~\ref{C-352-10}--\ref{C-352-16}); the kernel families (Lemma~\ref{C-352-17} to Proposition~\ref{C-352-19}); the induction over the levels and the jets (Lemma~\ref{C-352-20} to Proposition~\ref{C-352-24}); the connected functions on the least torus (Lemmas~\ref{C-352-25}--\ref{C-352-29}); every state (Lemma~\ref{C-352-30} to Lemma~\ref{C-352-35}); the bond forms (Lemma~\ref{C-352-36} to Proposition~\ref{C-352-37}); pair minus tensor field (the lemma on a regular gauge on a cube, and Lemmas~\ref{C-352-38} to~\ref{C-352-41}); and the proofs of the theorems.

\paragraph{The calculus of the derivative families.}

\begin{lemma}[The run]\label{C-352-10}
Let $0<g\le\gamma$ with $g^{-2}\ge x_{\mathrm J}$, which holds for every $0<g\le\gamma_{\mathrm{jt}}$ (Remark~\ref{C-352-9}), and $(K,\beta)$ a first-passage datum at $g$ with its run and label $\mathsf b$.
\begin{enumerate}
\item[(a)] $x_K>g^{-2}\ge x_{\mathrm J}\ge\max\{e^{2\mathfrak p_0},eL,4\bar B\}$, so $y_K\ge\max\{2\mathfrak p_0,1+\log L\}$. For $0\le i\le j\le K$: $\frac45\mathsf b\le a_{K-j}x_j\le\frac43\mathsf b$; $\frac34L^{j-i}x_j\le x_i\le\frac54L^{j-i}x_j$, in particular $x_j\ge\frac45L^{-(j-i)}x_i$; $y_j\le y_i$; $\delta_j\le(\frac54)^{1/2}L^{(j-i)/2}\delta_i$ and $\delta_i\le\delta_j$; $g^{-2}<\mathsf b\le\bar{\mathsf b}$; and $\frac{11}{10}\mathsf c_{\mathfrak h}\delta_i\le r_i\le\frac85\mathsf c_{\mathfrak h}\delta_i$ for $1\le i<K$, so $x_nr_n^2\le\frac{64}{25}\mathsf c_{\mathfrak h}^2\mathfrak b_n^2$. Of these, $g^{-2}<\mathsf b\le\bar{\mathsf b}$ holds, together with the flow bound $|x_j-L^{-(j-i)}x_i|\le\bar B$ $(0\le i\le j\le K)$, for the label and the run of every first-passage datum at every $0<g\le\gamma$.
\item[(b)] Let $g^{-2}\ge\max\{e^{2\mathfrak p_0},e^{y_{\mathsf p}}\}$ and $\delta(g^{-2})\le\min\{1/2c_{\mathrm{rt}},1/c_\epsilon\}$, where $c_\delta=(\frac54L)^{1/2}$, $c_{\mathrm{rt}}=\max\{8\mathsf c_{\mathfrak h}c_\delta/5\varrho_{\mathrm{ch}},1/\varrho_{\mathrm{fl}}\}$, $c_\epsilon=c_{\mathrm{cu}}+16\mathsf c_{\mathbf C}\mathsf c_{\mathfrak h}c_\delta/5r_\alpha$, $y_{\mathsf p}=\max\{1,(LM)^3A_1^{-2}(1+2\log^+(c_{\mathsf p}/A_1))\}$. Here $\varrho_{\mathrm{ch}}$ and $\varrho_{\mathrm{fl}}$ are the radii
denoted by $\rho_{\rm ch}$ and $\rho_{\rm fl}$ in
equation~\YMMref{(2285)}{833}, and
$r_\alpha=a^{\rm us}(\alpha_0,\alpha_1)$ is the usable chart radius
of equation~\YMMref{(933)}{427}. The covariance and fluctuation-tail
constants are
\[
 \mathsf c_{\mathbf C}=2/\gamma_r,\qquad
 c_{\mathsf p}=C_{\rm ex}((L+2)M)^3e^{4\lambda_1},
\]
as in equations~\YMMref{(2294)}{836} and~\YMMref{(2296)}{836}.
For this regular one-step expansion,
$\gamma_2=2(LM)^{-3}$ by equation~\YMMref{(889)}{415}; hence its
tail factor $e^{-\gamma_2\mathfrak b_i^2/2}$ is
$e^{-(LM)^{-3}\mathfrak b_i^2}$.
Finally, $c_{\mathrm{cu}}$ is the constant of Lemma~\ref{C-352-20}.

Then at all heights $1\le i<K$: $2c_{\mathrm{rt}}\delta_i\le1$, $c_\epsilon\delta_i\le1$, $y_i\ge y_{\mathsf p}$; hence $2r_i\le\varrho_{\mathrm{ch}}$, $r_{i+1}\le\frac85\mathsf c_{\mathfrak h}c_\delta\delta_i$, $\mathsf p_{\star,i}\le\delta_i$ and $\epsilon'_i\le c_\epsilon\delta_i\le1$, where $\mathsf p_{\star,i}=c_{\mathsf p}\exp[-(LM)^{-3}\mathfrak b_i^2]$ and $\epsilon'_i$ is the number of Lemma~\ref{C-352-20}; if besides $\delta(g^{-2})\le1/2C_{\mathrm{cu}}$, then $\epsilon^{\mathrm{cu}}_i\le2C_{\mathrm{cu}}\delta_i\le1$. By the conditions on the constants $\mathfrak b_j\ge(LM)^{3/2}\ge1$, so $g_j\le\delta_j$ and $x_j^{-1}\le\delta_j^2\le r_j^2$ $(j<K)$; and $\theta^{[j]}\le\theta^{[K-1]}\le\theta(g^{-2})$.
\item[(c)] For $0<\varrho<1$ with $s(\varrho)\le K$: $\log\frac1\varrho\le Y_\varrho\le2\log\frac1\varrho+C(g)$, $a_{K-\mathsf k(\varrho)}\le\varrho/4M$, and $\delta_{\mathsf k(\varrho)}\le A_1Y_\varrho^{\mathfrak p_0}(\varrho/\mathsf b)^{1/2}$.
\end{enumerate}
\end{lemma}

\begin{proof}
(a) At every $0<g\le\gamma$ the run of a first-passage datum has $x_0>x_1>\dots>x_K>g^{-2}$ and $|x_j-L^{-(j-i)}x_i|\le\bar B$ for $i\le j$, and its label has $g^{-2}<\mathsf b\le\bar{\mathsf b}$ (Proposition~\YMMref{74.1}{615}). The threshold $x_{\mathrm J}$ is at least $\max\{e^{2\mathfrak p_0},eL,4\bar B\}$ by its definition (the definitions preceding \eqref{eq:weighted-history-threshold}); so $g^{-2}\ge x_{\mathrm J}$ gives $x_j\ge x_K>g^{-2}\ge4\bar B$, hence $\bar B\le\frac14x_j$ and $\frac34x_j\le L^{-(j-i)}x_i\le\frac54x_j$, which at $i=0$ gives the bounds of $a_{K-j}x_j$ against $\mathsf b=a_Kx_0$. As $y$ decreases along the run and $x_i\le\frac54L^{j-i}x_j$, $\delta_j/\delta_i=(y_j/y_i)^{\mathfrak p_0}(x_i/x_j)^{1/2}\le(\frac54)^{1/2}L^{(j-i)/2}$.

As $x_i\ge x_j>g^{-2}\ge e^{2\mathfrak p_0}$ lie in $[e^{2\mathfrak p_0},\infty)$, where $\delta$ decreases, $\delta_i\le\delta_j$. The radii follow from their recursion and these bounds. (b) $x_i>g^{-2}\ge e^{2\mathfrak p_0}$ and $\delta$ decreases there, so every condition stated at $g^{-2}$ holds at every height; $2r_i\le\varrho_{\mathrm{ch}}$ and $r_{i+1}\le\frac85\mathsf c_{\mathfrak h}c_\delta\delta_i$ by (a). For $y\ge y_{\mathsf p}$: $(LM)^{-3}A_1^2y^{2\mathfrak p_0}\ge\frac y2+\log^+(c_{\mathsf p}/A_1)$, so $\mathsf p_{\star,i}\le A_1e^{-y_i/2}\le A_1y_i^{\mathfrak p_0}x_i^{-1/2}=\delta_i$; then $\epsilon'_i\le c_{\mathrm{cu}}\delta_i+2\mathsf c_{\mathbf C}r_{i+1}/r_\alpha\le c_\epsilon\delta_i$ and $\epsilon^{\mathrm{cu}}_i\le C_{\mathrm{cu}}(\delta_i+\mathsf p_{\star,i})\le2C_{\mathrm{cu}}\delta_i$. (c) As $K-s\le\mathsf k(\varrho)\le K-s+2$, the working level gives $\varrho L^s\le L^2\hat\varrho_k<L^3C_{W,0}A_0^2Y^{2\mathfrak p_0}$, whence
\[
Y=s\log L+\log L(g^{-2}+\bar B)<\log\tfrac1\varrho+2\mathfrak p_0\log Y+\log[L^4C_{W,0}A_0^2(g^{-2}+\bar B)],
\]
and $2\mathfrak p_0\log Y\le\frac12Y+2\mathfrak p_0\log4\mathfrak p_0$; further $a_{K-k}=\varrho/\hat\varrho_k\le\varrho/4M$ and $\delta_k=A_1y_k^{\mathfrak p_0}x_k^{-1/2}\le A_1Y^{\mathfrak p_0}(5a_{K-k}/4\mathsf b)^{1/2}$.
\end{proof}

The frame used below includes the identity of the step data with the
canonical construction. It also records the numerical room needed to
lower a field through the chain. Keeping both requirements explicit
allows us to apply the estimates at an arbitrary lower level.

\begin{definition}\label{def:pair-functional-frame}
Let $\mathfrak w=(\ell,K)$, $\ell=2L^m$, and let
$x_0,\ldots,x_K$ be a finite running sequence. The pair setting consists
of the following conditions.
\begin{itemize}
\item[(Jt.1)] Fix $2\le k_0<K$, $d\ge a_K$ and real Lipschitz
functions $f_1,f_2$, each supported in a closed ball of radius $d$.
Put $w_a(p)=f_a(x(p))$ and $s_0=K-k_0$.
\item[(Jt.2)] The setting \textup{(J)} of
Definition~\ref{def:weighted-history-setting} holds through the upper
level $K$, with $p=2$ and the calibration equal to this running sequence.
Thus its index $k$ is $K-1$, and its actual step data, including
\textup{(Pv)}, are required. Both sampled sources are nonzero, and
\[
 \widehat d_{k_0}:=d/a_{K-k_0}\ge M,
 \qquad D_{k_0}\ge3M.
\]
The running and weighted bounds are
\begin{equation}\label{eq:pair-frame-running}
\begin{gathered}
 |x_j-L^{-(j-i)}x_i|\le\bar B\quad(0\le i\le j\le K),
 \qquad x_K\ge4\bar B,
 \qquad y_K\ge\max\{2\mathfrak p_0,1+\log L\},\\
 C_0^\flat\mathsf S_A\theta^{[K-1]}\le1,
 \qquad\varepsilon_j\le c_A\alpha_0/(8MB_3)\quad(1\le j\le K),
 \qquad\varepsilon_j\le\varepsilon_\flat\quad(0\le j<K).
\end{gathered}
\end{equation}
The weighted data restrict to every earlier height. At heights
$1\le j<k_0$ they satisfy
\[
 s^{[j]}\le s^\star/2,
 \qquad\tfrac43 C_1^\flat L\theta^{[j]}\mathsf T(D_j/M)\le1.
\]
These are the hypotheses for the mixed derivative bound at the lower
cut, in addition to the running bounds just stated.

The flat-top radii are
$r_K=0$ and
$r_i=L^{-1}r_{i+1}+\frac{11}{10}\mathsf c_{\mathfrak h}\delta_i$.
Set
\[
 c_R=\max\left\{\tfrac85\sqrt{5L/4}\,
       \mathsf c_{\mathfrak h}/\varrho_{\rm ch},
       \varrho_{\rm fl}^{-1}\right\},\qquad
 e_i^{\rm KF}=\epsilon_i^{\rm cu}
       +\tfrac{16}5\sqrt{5L/4}\,
       \mathsf c_{\mathbf C}\mathsf c_{\mathfrak h}\delta_i/r_\alpha.
\]
Require $\ell\ge N_{\mathsf H}$ and, for $k_0\le i<K$,
\begin{equation}\label{eq:pair-frame-lowering}
\begin{gathered}
 c_R\delta_i\le\tfrac12,\qquad e_i^{\rm KF}\le1,\\
 [\tfrac{48}5\sqrt{5/4}L^{-3/2}\mathsf c_{\mathfrak h}
          +8C_{\rm db}]\delta_i<\varepsilon_{\mathcal W}\alpha_0.
\end{gathered}
\end{equation}
\item[(Jt.3)] The two-source weighted progression is the one of
Definition~\ref{def:weighted-progression-data}, with $k=K-1$ and
$t_j=x_j$. It uses the same canonical data as \textup{(Jt.2)}, when
that condition is imposed. Its source is
$w_\lambda=\lambda_1w_1+\lambda_2w_2$,
$w^{[j]}_\lambda=L^{K-j}w_\lambda$, and
$C_0^\flat\mathsf S_A\theta^{[K-1]}\le1$.
\item[(Jt.4)] Both sources are nonzero,
$s^{[j]}\le s^\star/2$ for $0\le j<k_0$, and
\[
 \tfrac43 C_1^\flat L\theta^{[j]}\mathsf T(D_j/M)\le1
 \quad(1\le j<k_0).
\]
\item[(Jt.5)] An intermediate integer $T$ is fixed with $k_0<T\le K$.
\item[(Jt.6)] For $k_0\le i<K$,
\begin{equation}\label{eq:pair-frame-curvature}
 5\mathsf c_{\mathfrak h}\delta_i\le c_A\alpha_0/(4M).
\end{equation}
\end{itemize}
The \emph{frame at $k_0$} means \textup{(Jt.1)--(Jt.3)} and
\textup{(Jt.6)}. The full pair setting includes \textup{(Jt.4)} and
\textup{(Jt.5)} as well.
\end{definition}

\begin{lemma}[The frame of the pair functional; tents]\label{C-352-11}
Let $0<g\le\gamma_{\mathrm{jt}}$, $g^{-2}\ge x_{\mathrm J}$, $(K,\beta)$ a first-passage datum at $g$ and $\ell\in\mathfrak L(g)$.
\begin{enumerate}
\item[(a)] Let $2\le k<K$ and let $f,h$ be real Lipschitz functions on $T_\ell$, neither $f\circ x$ nor $h\circ x$ identically $0$, each with its support in a closed ball of a radius $\bar{\mathsf d}\ge a_K$ with $\bar{\mathsf d}/a_{K-k}\ge M$, such that the corners of the plaquettes of the two supports are at least $3M$ level-$k$ spacings apart in the sup distance. Then $(f\circ x,h\circ x)$ satisfies, with the lower level $k$, the four conditions \textup{(Jt.1)}--\textup{(Jt.3)} and \textup{(Jt.6)} of Definition~\ref{def:pair-functional-frame}: the data, the flat top at $K$, the weighted progression of the two sources, and the smallness of $\delta_i$ at the heights $k\le i<K$; we call these four conditions the frame at $k$.

The two further conditions of that setting, \textup{(Jt.4)} at the levels below the lower level and \textup{(Jt.5)} on an intermediate level, hold, together with the other four, for a test pair $(f_1,f_2)\in\mathfrak T_{T_\ell}(r;\vartheta)$, $0<r<1$, $\vartheta\in(0,1)$, with $f_a\circ x\not\equiv0$, at the lower level $\max\{K-s,2\}$ when $3\le s=s(r)\le K$ and $\mathsf{Lm}_s\le L^{s-3}$ with the intermediate level specified in the proof. The frame at $k$ holds in the same way for the two sampled linear combinations of $n$ functions with pairwise disjoint supports falling into two clusters, provided both combinations are nonzero, the two unions of supports lie in such balls, and the corners of the two clusters are $3M$ level-$k$ spacings apart.
\item[(b)] Let real weights $w_1,w_2$ be placed at $(\rho,r)$, $0<\rho\le r\le e^{-2\sqrt3}$, $s(\rho)\le K$. Then $a_K\le\rho/4M$; with the tent $\tau$ of height $1$ and radius $\frac13a_K$, the functions $f_a=\sum_pw_a(p)\tau(\cdot-x(p))$ are Lipschitz, $f_a\circ x=w_a$, $\operatorname{supp}f_a\subset\mathbb B(x_a,2r)$, and (a) applies at $k=\mathsf k(r)$ with $\bar{\mathsf d}=2r$. For $0\le i\le K$: (P1) $\sum_p|w_a(p)|\le81(\rho/a_K)^3|w_a|_\infty$; (P2) $\mathsf N_i(w_a)\le(2\hat\rho_i/M+3)^3\le108(1+t_i)$, $t_i=\frac{27}{8M^3}\hat\rho_i^3$; (P3) $\mathsf D^{12}_i\ge\hat r_i/\sqrt3-1$; and every corner of a plaquette of $\operatorname{supp}w_1\cup\operatorname{supp}w_2$ lies within the sup distance $R_w\le11\hat r_{i+1}$, in level-$(i+1)$ units, of a bare site next to $x_1$. For weights placed at $(\rho,r,\mathsf d)$, or forming two pairs at $(\rho,r,\mathsf d)$, with $\mathsf d\le e^{-2\sqrt3}$, the same holds, (P3) holds across the clusters with $\hat{\mathsf d}_i$ for $\hat r_i$, and (a) applies to the two clusters at $k=\mathsf k(\mathsf d)$, each lying in a ball of radius $10\mathsf d$; for two weights with $\Sigma(w_a)\subset\mathbb B(x_a,\rho)$ at a distance $\ge\mathsf d\ge\rho$, (a) applies at $k=\mathsf k(\mathsf d)$.
\end{enumerate}
\end{lemma}

\begin{proof}
(a) First consider a pair satisfying the four geometric hypotheses
in the statement. The first-passage construction gives
$x_0>\cdots>x_K>g^{-2}$ and the running comparison in
\eqref{eq:pair-frame-running}; see Proposition~\YMMref{74.1}{615}.
Since $\ell\in\mathfrak L(g)$, it also gives
$MR(x_K)\le L^m$. The lattice is therefore admissible over every
initial segment of the run. In particular $L^m\ge LM$ and
$\ell\ge N_{\mathsf H}$, by the proof of
Lemma~\YMMref{126.2}{845}.

Proposition~\ref{prop:weighted-canonical-history}, applied with
$k=K-1$, supplies the actual setting \textup{(J)}, its weighted
progression and the identification \textup{(Pv)}. This is the same
canonical datum at every height, with the counterterms of this run.
The inequalities defining $x_{\rm J}$ apply at every $x_j$:
they give the regular-background radius, the half-ceiling and
$C_0^\flat\mathsf S_A\theta^{[K-1]}\le1$.
The cutoff bound $\varepsilon_j\le\varepsilon_\flat$ is part of the
actual history verification in that proposition.

The supports satisfy $\widehat d_k\ge M$ and $D_k\ge3M$ by
hypothesis. Their separation at earlier levels is
$D_j=L^{k-j}D_k$, so the decreasing function $\mathsf T$ obeys
$\mathsf T(D_j/M)\le\mathsf T(3)$ for $j<k$.
Equation~\eqref{eq:weighted-history-threshold} therefore gives
all the mixed derivative hypotheses at the lower cut. The threshold
inequalities in equation~\YMMref{(2326)}{845} give
\eqref{eq:pair-frame-lowering}. The factor $2$ in the corresponding
lowering-domain inequality makes its conclusion strict. They also
give $y_K\ge\max\{2\mathfrak p_0,1+\log L\}$ and
$x_K\ge4\bar B$.
Thus \textup{(Jt.1)--(Jt.3)} hold. Finally
$x_i>g^{-2}\ge x_{\rm jt}$ and the threshold computation in
Remark~\ref{C-352-9} give \textup{(Jt.6)}.

This verification uses only the four geometric properties in the
statement. It uses neither the filling nor the sign condition of a
test pair, nor the formula for the working level. For $n$ functions
in two clusters, apply it to the two sums over the clusters, with
real coefficients for which both sampled sums are nonzero.
Shrinking a support preserves the radius and separation conditions.

For a test pair at $k_0=\max\{K-s,2\}$, Lemma~\YMMref{126.1}{844}
gives $d\ge a_K$, $\widehat d_{k_0}\ge M$ and $D_{k_0}\ge3M$.
The preceding argument gives its frame. At all heights below $k_0$,
\eqref{eq:weighted-history-threshold} gives
$s^{[j]}\le s^\star/2$, including $j=0$, and the same separation
argument proves \textup{(Jt.4)}. Choose
\[
 n_T=\left\lceil\log_L
      (\widehat d_{k_0}y_{k_0}^{3\mathfrak p_0})\right\rceil,
 \qquad T=k_0+n_T.
\]
The room bound $\mathsf{Lm}_s\le L^{s-3}$ and the final assertion of
Lemma~\YMMref{126.1}{844} give $k_0<T\le K-2$.
This proves \textup{(Jt.5)} and the claimed full pair setting.

(b) The working level at $\varrho=\rho$ gives $a_K\le a_{K-\mathsf k(\rho)}\le\rho/4M$. Two distinct barycentres are at least $2^{-1/2}a_K$ apart, so the tents are disjoint and $f_a(x(p))=w_a(p)$.

A corner is at the distance $2^{-1/2}a_K$ from its barycentre: corners of the two supports are $\ge r-\sqrt2a_K$ apart, $\ge(r-\sqrt2a_K)/\sqrt3$ in the sup distance, that is $\ge\hat r_i/\sqrt3-1$ level-$i$ units. At $\varrho=r$ the working level gives $2\le\mathsf k(r)\le K-1$, $2\hat r_{\mathsf k(r)}\ge M$ and $\mathsf D^{12}_{\mathsf k(r)}\ge\frac4{\sqrt3}MY_r^{2\mathfrak p_0}-1\ge3M$, as $Y_r\ge2\sqrt3$. (P1): the barycentres of one orientation form a translate of $a_K\mathbb Z^3$, with at most $(2\rho/a_K+1)^3\le27(\rho/a_K)^3$ points in a ball of radius $\rho\ge a_K$, and there are three orientations. (P2): an interval of length $2\hat\rho_i+\eta_i$ meets at most $2\hat\rho_i/M+3$ intervals of length $M$, and $(2y+3)^3\le4(8y^3+27)$. $R_w$: all points of $\Sigma(w_1)\cup\Sigma(w_2)$ are within $9r$ of $x_1$.

The three-point and four-point cases are the same with $\mathsf d$ for $r$ across the clusters, and $9r+a_K\le10\mathsf d$.
\end{proof}

\begin{lemma}[The calculus of the derivative families]\label{C-352-12}
Let $0<g\le\gamma$ with $g^{-2}\ge x_{\mathrm J}$, $(K,\beta)$ a first-passage datum at $g$ with $K\ge2$, $\ell\in\mathfrak L(g)$, and $w_1,\dots,w_p$ real weights, none identically $0$, $p\ge1$ arbitrary. For $\emptyset\ne\beta\subset[p]$, $m=|\beta|$, $0\le j\le K-1$, put
\[
\begin{multlined}
\mathsf m_{j,a}=12M^3\alpha_0^2L^{K-j}|w_a|_\infty,\quad\varkappa_j=\frac{M^{3/2}}{2\mathsf C_{\mathcal A}\mathfrak b_j},\quad\gamma_j=\theta^{[j]}\varkappa_j,\quad\mathsf U_{j;\beta} \\
=\varkappa_j^{m-1}(N^{[j]}_\star)^{1-m}\prod_{a\in\beta}2\mathsf m_{j,a},\quad\bar\nu_{j;\beta}=\gamma_j\mathsf U_{j;\beta},
\end{multlined}
\]
with the constant $\mathsf C_{\mathcal A}$ of the one-step theorem; the bare arrays of the weighted progression, $\mathfrak f_{\lambda,j}=((E^{(i)}_\lambda,e_i)_{i\le j},(\mathsf W_{L^{K-j}w_\lambda}))$ for $\lambda$ near $0$, where the counterterms $e_i$ of the regular progression of the run are real numbers that depend neither on $\lambda$ nor on the sources, and the weight member $\mathsf W_w$ is linear in $w$, so that $E^{(j+1)}_\lambda-E^{(j+1)}=\mathsf{St}^{[j]}(\mathfrak f_{\lambda,j})-\mathsf{St}^{[j]}(\mathfrak f_{0,j})$ (equation~\eqref{eq:weighted-progression-recursion}); their value at $\lambda=0$, $\mathfrak f_{0,j}=((E^{(i)},e_i)_{i\le j},(0))$; and their derivatives $\mathfrak f_{\beta,j}=\partial_\beta\mathfrak f_{\lambda,j}|_{\lambda=0}=((E^{(i)}_\beta,0)_{i\le j},(H_\beta))$ with $H_{\{a\}}=\mathsf W_{L^{K-j}w_a}$ and $H_\beta=0$ for $m\ge2$; and $\bar{\mathsf h}_{j;\beta}=\mathfrak n[\mathfrak f_{\beta,j}]/\mathsf U_{j;\beta}$, $\bar\Psi_{j;\beta}=\|E^{(j+1)}_\beta\|_{j+1,\lambda_{\mathrm{out}}}/\bar\nu_{j;\beta}$. Then for $1\le j\le K-1$:
\begin{enumerate}
\item[(a)] \emph{Chain rule.} $E^{(j+1)}_\beta=\sum_{\pi\in\mathrm{Part}(\beta)}\bigl(\prod_{\beta'\in\pi}\partial_{s_{\beta'}}\bigr)\mathsf{St}^{[j]}\bigl(\mathfrak f_{0,j}+\sum_{\beta'\in\pi}s_{\beta'}\mathfrak f_{\beta',j}\bigr)\big|_{s=0}$; $E^{(j)}_\beta$ is $m$-linear in $(w_a)_{a\in\beta}$, depends on no other source, and is the same in every progression that holds the sources of $\beta$.
\item[(b)] \emph{Derivatives of the step.} $N(\mathfrak f_{0,j})\le\frac12N^{[j]}_\star$, and for bare arrays $\mathfrak f,\mathfrak g_1,\dots,\mathfrak g_q$ with $N(\mathfrak f)\le\frac12N^{[j]}_\star$
\[
\bigl\|\partial_{s_1}\cdots\partial_{s_q}\mathsf{St}^{[j]}\bigl(\mathfrak f+\textstyle\sum_ls_l\mathfrak g_l\bigr)\big|_0\bigr\|_{j+1,\lambda_{\mathrm{out}}}\le(2(q-1))^{q-1}\,\theta^{[j]}\varkappa_j^q\,(N^{[j]}_\star)^{1-q}\prod_l\mathfrak n[\mathfrak g_l].
\]
\item[(c)] \emph{Lags.} If $\mathfrak g=((G^{(i)},0)_{i\le j},(H))$ and each $G^{(i)}(X,\cdot,z)$ vanishes unless $X$ meets a set $\mathsf{Cr}$ of bare sites, then $\mathfrak n[\mathfrak g]\le\|H\|_{j,\kappa}+\sum_{i\le j}\omega^{\mathrm{sp}}_{j-i}(i)\|G^{(i)}\|_{i,\kappa}$, where $\omega^{\mathrm{sp}}_0=1$, $\omega^{\mathrm{sp}}_l(i)=\varpi_l+L^{-4l}\min\{L^{3l},c_3\mathsf n^{\mathrm{cb}}_i(\mathsf{Cr})\}$, $0\le\varpi_l\le L^{8-5l}$, $\mathsf n^{\mathrm{cb}}_i(\mathsf{Cr})$ is the number of the cubes of $\pi_i$ that hold a site of $\mathsf{Cr}$, and $c_3$ is a constant.
\item[(d)] \emph{Support.} $E^{(i)}_\beta(X,\cdot,z)=0$ unless every $\operatorname{supp}w_a$, $a\in\beta$, has a plaquette with its corners in $X$; and then $d_i(X)\ge[\max_{a,b\in\beta}\mathsf D^{ab}_i/M-2]_+$.
\item[(e)] \emph{Start.} $\bar\Psi_{0;\beta}\le m^m(\mathsf c^\flat/2)^m$ with a constant $\mathsf c^\flat$; and if $\sum_p|w_a(p)|\le81(\rho/a_K)^3|w_a|_\infty$ for a width $\rho\ge a_K$, then $\|\mathsf W_{L^{K-j}w_a}\|_{j,\kappa}\le\varsigma_j\mathsf m_{j,a}$, $\varsigma_j=\min\{1,t_j\}$.
\item[(f)] \emph{Units.} $\gamma_j=c_\gamma x_j^{-1/2}$, $\bar\nu_{j;\beta}=\bar c_mx_j^{-m/2}\prod_{a\in\beta}(|w_a|_\infty/a_{K-j})$; for $i\le j$: $(\frac45)^{m/2}L^{(j-i)m/2}\bar\nu_{j;\beta}\le\bar\nu_{i;\beta}\le(\frac43)^{m/2}L^{(j-i)m/2}\bar\nu_{j;\beta}$, $\gamma_i\le(\frac43)^{1/2}L^{-(j-i)/2}\gamma_j$, $\sum_{i<j}\gamma_i\le0.45\gamma_j$; and $\varkappa_j\le\frac16L^{-3/2}$.
\item[(g)] \emph{Recursions.} $\bar\Psi_{j;\beta}\le\sum_{\pi\in\mathrm{Part}(\beta)}(2(|\pi|-1))^{|\pi|-1}\prod_{\beta'\in\pi}\bar{\mathsf h}_{j;\beta'}$; and for $m\ge2$ and every $a_*\in\beta$, with $\omega^{\mathrm{sp}}$ formed with $\mathsf n^{\mathrm{cb}}_i=\mathsf N_i(w_{a_*})$,
\[
\begin{multlined}
\bar{\mathsf h}_{j;\beta}\le\gamma_j(\tfrac43)^{m/2}\sum_{i\le j}\omega^{\mathrm{sp}}_{j-i}(i)\,\sigma_{i;\beta}\,L^{(j-i+1)m/2}\,\bar\Psi_{i-1;\beta},\qquad\sigma_{i;\beta}=e^{-c_\sigma[\max_{a,b\in\beta}\mathsf D^{ab}_i/M-2]_+}, \\
c_\sigma=\min\{\tfrac12,4\beta_{\mathrm s}\kappa\}.
\end{multlined}
\]
\end{enumerate}
No assumption of (a)--(g) depends on $p$; the order $m$ enters through the displayed numbers alone. There are numbers $\tau^\natural,\gamma^\natural>0$ formed from the constants such that $\theta^{[K-1]}\le\tau^\natural$ gives $\gamma_j\le\gamma^\natural$ for all $j<K$ and $2C_\omega\mathsf c^\flat\theta^{[0]}\le1$, $C_\omega$ the constant of Lemma~\ref{C-352-15}; and a number $r_1(\mathsf b)>0$, formed from the constants and $\mathsf b$ and not increasing in $\mathsf b$, such that every width $0<\rho\le r_1(\mathsf b)$ with $s(\rho)\le K$ has $\rho\ge4Ma_K$ and $E_\sigma(\hat\rho_1/2)\le\varkappa_0^4$, where $E_\sigma(u)=\exp(-c_\sigma[u/\sqrt3M-3]_+)$.
\end{lemma}

\begin{proof}
(a) The progression is $E^{(j+1)}_\lambda-E^{(j+1)}=\mathsf{St}^{[j]}(\mathfrak f_{\lambda,j})-\mathsf{St}^{[j]}(\mathfrak f_{0,j})$ (equation~\eqref{eq:weighted-progression-recursion}), and the multilinear derivatives at $\lambda=0$ of a composition are sums over partitions; the dependence through $w_\lambda$ alone gives the multilinearity. (b) Within the holomorphic, gauge-invariant input class of Proposition~\ref{prop:source-entry}, the fourth family must vanish at the block field $0$ and have finite norm \eqref{eq:source-entry-norm} on the polydisc $|B(b)|<\mathsf C_{\mathcal A}\mathfrak b_j$. A family that vanishes at $B=0$ and is bounded on that polydisc is at most its supremum times $|B|/\mathsf C_{\mathcal A}\mathfrak b_j$ (Schwarz's lemma); so $\mathsf{St}^{[j]}(\mathfrak f)$ is formed for $\varkappa_j\mathfrak n[\mathfrak f]\le N^{[j]}_\star$, with
\[
|\mathsf{St}^{[j]}(\mathfrak f)-\mathsf{St}^{[j]}(\mathfrak f')|(X,P,z)\le\theta^{[j]}\varkappa_j\,\mathfrak n[\mathfrak f-\mathfrak f']\,e^{-\lambda_{\mathrm{out}}d_{j+1}(X)},
\]
and it is holomorphic in the parameters of $\mathfrak f+\sum s_l\mathfrak g_l$ on $\{\varkappa_j\mathfrak n[\cdot]<N^{[j]}_\star\}$. Dividing the Lipschitz bound by $|s|$ and letting $s\to0$ gives $\|\mathsf D\mathsf{St}^{[j]}(\mathfrak h)[\mathfrak g_1]\|_{j+1,\lambda_{\mathrm{out}}}\le\theta^{[j]}\varkappa_j\mathfrak n[\mathfrak g_1]$ wherever $\varkappa_j\mathfrak n[\mathfrak h]<N^{[j]}_\star$; this is (b) at $q=1$. For $q\ge2$, as $\varkappa_j\mathfrak n[\mathfrak f]\le N(\mathfrak f)\le\frac12N^{[j]}_\star$, the point $\mathfrak f+\sum_{l\ge2}s_l\mathfrak g_l$ stays in that set on the polydisc of the radii $N^{[j]}_\star/2(q-1)\varkappa_j\mathfrak n[\mathfrak g_l]$ in $s_2,\dots,s_q$; Cauchy's inequality in these $q-1$ variables, applied to the first derivative in the direction $\mathfrak g_1$, gives the factor $\prod_{l\ge2}2(q-1)\varkappa_j\mathfrak n[\mathfrak g_l]/N^{[j]}_\star$, that is (b), with $(2(q-1))^{q-1}=1,2,16,216$ for $q=1,\dots,4$.

This factor $1/\mathfrak b_j$ for every order in the sources is the whole of the loss $\mathfrak b_j^2$ that a bound through the size $N$ would cost. $N(\mathfrak f_{0,j})\le\frac12N^{[j]}_\star$ is the half-ball estimate in the proof of Proposition~\YMMref{122.2}{827}: the admitted bound $s_j=\theta^{[j]}/K_{\mathbf M}\le s_\omega^\star/2$ gives $N_j(f_{0,j})\le N_\star\le w_\star/(4K_\Lambda s_\omega^\star)\le w_\star/(8K_\Lambda s_j)=N_{\star,j}/2$. Proposition~\ref{prop:weighted-canonical-history} verifies this ceiling at the actual run. (c) Use the wide-domain and octet decomposition in the proof of
Proposition~\YMMref{51.16}{412}, with zero counterterms and all relative
localization radii equal to $1$. Put $l=j-i\ge1$ and
$D_l=\max\{2L^{l-1}-2,0\}$. The wide domains have size at least $D_l$.
Their rooted sum gives the factor
$\varpi_l=L^{3l}e^{-D_l}$ of equation~\YMMref{(884)}{414}.
For $l=1$ this is $L^3$; for $l\ge2$,
\[
 D_l\ge2(l-1)(L-1)\ge8(l-1)\log L,
 \qquad\varpi_l\le L^{8-5l},
\]
where the second inequality uses $L\ge13$.
For a narrow domain, equation~\YMMref{(864)}{408} bounds the octet
term minus its constant by
\[
 C_{\rm oc}\|G^{(i)}\|_{i,\kappa}
      (d_i(X)+2)^2L^{4-4l}e^{-\kappa d_i(X)}.
\]
Sum first over domains meeting $\mathsf{Cr}$ and then over their
base points. The rooted count in Lemma~\ref{lem:source-support-sum}
applies with the marked sites of $\mathsf{Cr}$ in place of plaquette
corners: each relevant domain contains one of the
$\mathsf n_i^{\rm cb}(\mathsf{Cr})$ marked cubes, and
$\#(X\cap\mathbb T^{(i)})\le8M^3(d_i(X)+1)$.
The remaining sum is bounded by the same convergent series of a
polynomial in $d_i(X)$ times its exponential decay.
This replaces the $L^{3l}$ possible positions in the unrestricted
count by $c_3\mathsf n_i^{\rm cb}(\mathsf{Cr})$.
Taking the better of the two counts, with the normalization of
equation~\YMMref{(885)}{414}, gives the narrow contribution
$L^{-4l}\min\{L^{3l},c_3\mathsf n_i^{\rm cb}(\mathsf{Cr})\}$.
The fixed constants, including $L^4$, are those in the bracket
normalization and in $c_3$. There is no growing-radius factor here:
the relative radii are $1$. The level-$j$ family and $H$ have no lag
and retain their ordinary norm bounds. This proves (c). (d) is the locality of the progression with the linear size of a domain. (e) The one-step theorem applied at the height $0$, where the source stands in the fourth place itself, with Schwarz's lemma as in (b); the weight member: $|a_p|<\frac12\alpha_0^2\eta_j^4$ on the cores, and the count of plaquettes. (f) is arithmetic along the run (Lemma~\ref{C-352-10}). (g) follows from (a), (b), (e) in the units, and from (c) with $\mathsf{Cr}=\mathsf{Cr}_{a_*}$, (d) and (f).

The number $\gamma^\natural$, formed from the constants, is the smallness that Lemma~\ref{C-352-15} reads: $\gamma^\natural\le(3.6\,C_\omega)^{-1}$, so that $\epsilon^{(1)}_j=C_\omega(0.9\gamma_j+\mathsf c^\flat\theta^{[0]})\le\frac34$ in Lemma~\ref{C-352-15}(a) once $2C_\omega\mathsf c^\flat\theta^{[0]}\le1$; $\gamma^\natural$ is at most the number under which the feedback $\sum_{i'<j}\gamma_{i'}\le0.45\gamma_j$ is absorbed in the recursion of the proof of Lemma~\ref{C-352-15}(b); and $\gamma^\natural\le(5C_1L^{7/2}\hat C)^{-1}$, which the proof of Lemma~\ref{C-352-15}(e) reads, where $\hat C=\sup_{v>0}(1+v^3)\max\{1,v^{12}\}E_\sigma(v/2)$ is finite, $E_\sigma$ decaying exponentially, and $C_1$ is the constant of the bound $\omega^{\mathrm{sp}}_l(i)\le C_1L^{-4l}(1+t_i)$ in that proof.

Then $\tau^\natural=\min\{6L^{3/2}\gamma^\natural,(2C_\omega\mathsf c^\flat)^{-1}\}$: for $j<K$, $\delta_j\le\delta_{K-1}$ (Lemma~\ref{C-352-10}(a)) gives $\theta^{[j]}\le\theta^{[K-1]}$, in particular $\theta^{[0]}\le\theta^{[K-1]}$, and with $\varkappa_j\le\frac16L^{-3/2}$ of (f), $\theta^{[K-1]}\le\tau^\natural$ gives $\gamma_j=\theta^{[j]}\varkappa_j\le\gamma^\natural$ and $2C_\omega\mathsf c^\flat\theta^{[0]}\le1$. For $r_1$, let $C_7=\sup_{u\ge1}[4\mathfrak p_0\log(1+\log u)-c_\sigma u/4\sqrt3LM]$, finite, let $\Lambda_1(\mathsf b)\ge1$ be the least number such that for all $\Lambda\ge\Lambda_1(\mathsf b)$
\[
\frac{c_\sigma}{\sqrt3L}\Lambda^{2\mathfrak p_0}\ge3c_\sigma+C_7+4\log^+\bigl(2\mathsf C_{\mathcal A}A_1M^{-3/2}\bigr)+4\mathfrak p_0\bigl[\log(1+\log^+\mathsf b)+\log(1+\Lambda)\bigr],
\]
and $r_1(\mathsf b)=e^{-\Lambda_1(\mathsf b)}$. The left member grows like $\Lambda^{2\mathfrak p_0}$, $\mathfrak p_0\ge1$, the right one like $\log\Lambda$, so $\Lambda_1(\mathsf b)$ is finite; it is formed from the constants and $\mathsf b$, it holds no $K$, and since the right member increases with $\mathsf b$, $r_1$ does not increase in $\mathsf b$. Let $0<\rho\le r_1(\mathsf b)$ with $s(\rho)\le K$, $\Lambda=\log\frac1\rho\ge\Lambda_1(\mathsf b)\ge1$ and $u=\rho/a_K=\hat\rho_0$. The working level at $\varrho=\rho<1$ gives $u\ge\hat\rho_{\mathsf k(\rho)}\ge4MY_\rho^{2\mathfrak p_0}\ge4M\Lambda^{2\mathfrak p_0}\ge4M$, so $\rho\ge4Ma_K$. Next, $1/\varkappa_0=2\mathsf C_{\mathcal A}A_1M^{-3/2}y_0^{\mathfrak p_0}$, and $x_0=\beta=\mathsf b/a_K=\mathsf bu/\rho$ gives $y_0=\log\mathsf b+\log u+\Lambda$, whence $\log y_0\le\log(1+\log^+\mathsf b)+\log(1+\log u)+\log(1+\Lambda)$, as $x+y+z\le(1+x)(1+y)(1+z)$ for $x,y,z\ge0$. As $\hat\rho_1=u/L$, the bound $E_\sigma(\hat\rho_1/2)\le\varkappa_0^4$ follows from
\[
c_\sigma\Bigl(\frac u{2\sqrt3LM}-3\Bigr)\ge4\log^+\bigl(2\mathsf C_{\mathcal A}A_1M^{-3/2}\bigr)+4\mathfrak p_0\log y_0 .
\]
Halve the first term: by $u\ge4M\Lambda^{2\mathfrak p_0}$ one half is at least $\frac{c_\sigma}{\sqrt3L}\Lambda^{2\mathfrak p_0}$, and the other half is at least $4\mathfrak p_0\log(1+\log u)-C_7$; with the bound of $\log y_0$, what is left is the display that defines $\Lambda_1(\mathsf b)$, at $\Lambda=\log\frac1\rho$. The term $\log(1+\log u)$, through which $y_0$ grows with $K$, is beaten by the half of the exponent, so no number in the argument depends on $K$.
\end{proof}

\begin{lemma}[The kernel family and the rest family]\label{C-352-13}
In the data of Lemma~\ref{C-352-12} with $p\ge2$, let $\dot{\mathfrak e}_{a,j}=((E^{(i)}_a,0)_{i\le j},(0))$ and $\mathfrak u_{a,j}=((0,0)_{i\le j},(\mathsf W_{L^{K-j}w_a}))$, so that $\mathfrak f_{\{a\},j}=\dot{\mathfrak e}_{a,j}+\mathfrak u_{a,j}$, and $\mathsf B_j(\mathfrak g_1,\mathfrak g_2)=\partial_{s_1}\partial_{s_2}\mathsf{St}^{[j]}(\mathfrak f_{0,j}+s_1\mathfrak g_1+s_2\mathfrak g_2)|_0$. With the kernel family $\mathsf K^{(j+1)}=\mathsf B_j(\mathfrak u_{1,j},\mathfrak u_{2,j})$, for $1\le j\le K-1$
\[
E^{(j+1)}_{12}=\mathsf K^{(j+1)}+\mathsf R^{(j+1)},\qquad\mathsf R^{(j+1)}=\mathsf B_j(\dot{\mathfrak e}_{1,j},\mathfrak u_{2,j})+\mathsf B_j(\mathfrak u_{1,j},\dot{\mathfrak e}_{2,j})+\mathsf B_j(\dot{\mathfrak e}_{1,j},\dot{\mathfrak e}_{2,j})+\mathsf D\mathsf{St}^{[j]}(\mathfrak f_{0,j})[\mathfrak f_{12,j}],
\]
and $\mathsf R^{(j+1)}$ is a pointed family of level $j+1$ that vanishes off $\mathsf{Cr}_1$ and off $\mathsf{Cr}_2$; $\mathsf K^{(j+1)}(X,\cdot,z)=0$ unless a plaquette of $\operatorname{supp}w_1$ and one of $\operatorname{supp}w_2$ have their corners in $X$.
\end{lemma}

\begin{proof}
Lemma~\ref{C-352-12}(a) at $\beta=\{1,2\}$ has two partitions. $\mathsf B_j$ is additive in each entry: for bare arrays $\mathfrak a,\mathfrak b,\mathfrak g$ the value $h(t)=\mathsf{St}^{[j]}(\mathfrak f_{0,j}+t_1\mathfrak a+t_2\mathfrak b+t_3\mathfrak g)(X,P,z)$ is holomorphic on a polydisc about $0$, the size being subadditive, and $\partial_{s_1}\partial_{s_2}h(s_1,s_1,s_2)|_0=(\partial_{t_1}+\partial_{t_2})\partial_{t_3}h(0)$. The first and second derivatives of the step are pointed families with the locality of the step, and Lemma~\ref{C-352-12}(d) gives where they vanish.
\end{proof}

\begin{lemma}[One level at a regular configuration]\label{C-352-14}
Let $g$ satisfy $(\Gamma1)$, $(K,\beta)$ be a first-passage datum at $g$ and $\ell\in\mathfrak L(g)$.
\begin{enumerate}
\item[(a)] Let the frame at a level $k_0$ hold for a pair of real Lipschitz functions (Lemma~\ref{C-352-11}(a)), and let $k_0\le n<K$, $1\le i'\le n$, $A\in\mathsf{Cc}_n(r_n)$, $\mathsf{Cr}$ a set of bare sites, $D'\ge0$, and $\mathsf f$ a pointed family of level $i'$ that vanishes off $\mathsf{Cr}$ and on the domains $X$ with $d_{i'}(X)<D'$. Then
\[
|\Delta[\mathsf f]|\le C_\Delta\,\mathsf n^{\mathrm{cb}}_{i'}(\mathsf{Cr})\,\|\mathsf f\|_{i',\lambda_{\mathrm{out}}}\,r_n^2\,L^{-4(n-i')}e^{-D'/2}.
\]
\item[(b)] Let $G$ be a family of local terms of level $i$ that vanishes unless $\mathcal P(X)$, the set of the plaquettes with their corners in $X$, meets a set $\mathcal S$ of plaquettes and $d_i(X)\ge D_i$, and let $\mathsf N_i$ be the number of the cubes of $\pi_i$ that hold a corner of a plaquette of $\mathcal S$. At every configuration $U$ whose pairs of level $i$ lie in the cores, $\sum_{(X,z)}|G(X,U,z)|\le8M^3C_{\mathrm E,1}\mathsf N_ie^{-D_i}\|G\|_{i,\lambda_{\mathrm{out}}}$; if $U$ is besides regular of a level $t'\ge i$ at a radius $\varepsilon$ with $4M\varepsilon\le c_{\mathrm A}\alpha_0$, then $\sum_{(X,z)}|G(X,U,z)-G(X,\mathbf 1,z)|\le2C_{\mathrm f}L^{-4(t'-i)}\mathsf N_ie^{-D_i/2}\|G\|_{i,\lambda_{\mathrm{out}}}$.
\end{enumerate}
The numbers $C_\Delta,C_{\mathrm f},C_{\mathrm E,1},c_{\mathrm A}$ are formed from the constants. Of $(\Gamma1)$ the proof of (a) uses only, besides the frame at $k_0$, the two bounds $2r_n\le\varrho_{\mathrm{ch}}$ and $x_n^{-1}\le r_n^2$ at the heights concerned, and the proof of (b) only $g\le\gamma$ and $g^{-2}\ge x_{\mathrm J}$; both use besides the bound $L^m\ge LMy_K$ of the sides $\ell=2L^m\in\mathfrak L(g)$. So (a) holds whenever the frame at $k_0$ and these two bounds hold, and (b) holds whenever $g\le\gamma$ and $g^{-2}\ge x_{\mathrm J}$, with no frame.
\end{lemma}

\begin{proof}
(a) Let $k_0\le n<K$. Then $r_n>0$ and $2r_n\le\varrho_{\mathrm{ch}}$ (Lemma~\ref{C-352-10}(a), (b)); in the frame at $k_0$, Lemma~\ref{lem:weighted-chain-background}, at the radii $r_n<r'=2r_n\le\varrho_{\mathrm{ch}}$, gives that $\tilde U$ is regular of level $n$ at the radius $\varepsilon:=6r_n$ on $T_\ell$ and that the actions of the bare arrays over the absolute system of height $n$ are formed at it; this is where the frame is read. For one domain $X$ of size $d$ and $\mathsf E=\mathsf f(X,\cdot,z)$, the evaluation of a local term at a regular configuration gives either $|\mathsf E(\tilde U)-\mathsf E(\mathbf 1)|\le C_{\mathrm A}\|\mathsf f\|e^{-\lambda_{\mathrm{out}}d}(\varepsilon/\alpha_0)^2L^{-4(n-i')}M^2(1+d)^2$, or $d\ge[d_*]_+$ and $|\mathsf E(\tilde U)-\mathsf E(\mathbf 1)|\le2\|\mathsf f\|e^{-\lambda_{\mathrm{out}}d}$, with the depth $d_*=\min\{c_{\mathrm A}\alpha_0\tau/6M,(L^{m+K-i'}-4)/2M\}-1$, $\tau=L^{2(n-i')}/r_n$.

A contributing $X$ holds one of the $\mathsf n^{\mathrm{cb}}_{i'}$ cubes, has $d\ge D'$ and at most $8M^3(d+1)$ base points; the sum over domains through a cube is bounded by \eqref{eq:analytic-domain-tail}. In the tail, $e^{-[d_*]_+/2}\le e^{1/2}[e^{-c_{\mathrm A}\alpha_0\tau/12M}+e^{-(L^{m+K-i'}-4)/4M}]$; the first part is $\le C\tau^{-2}$; for the second, $L^m\ge LMy_K$ on the sides of $\mathfrak L(g)$, $x_n\le\frac53L^{K-n}x_K$ and $x_n^{-1}\le r_n^2$ give $e^{-(L^{m+K-i'}-4)/4M}\le\frac53e^{1/M}r_n^2L^{-4(n-i')}$. So the bound holds on every side $\ell\in\mathfrak L(g)$. (b) is the same count at a configuration whose pairs lie in the cores, and, for the difference, at a configuration regular of level $t'$, with $\vartheta=L^{-(t'-i)}$.
\end{proof}

\begin{lemma}[Sizes at three lengths]\label{C-352-15}
Let $g$ satisfy $(\Gamma1)$, $(K,\beta)$ be a first-passage datum at $g$ and $\ell\in\mathfrak L(g)$, and let real weights be placed at $(\rho,r)$, or at $(\rho,r,\mathsf d)$, or form two pairs at $(\rho,r,\mathsf d)$, with $0<\rho\le r\le\mathsf d\le e^{-2\sqrt3}$, $s(\rho)\le K$, $\rho\le r_1(\mathsf b)$. Put $q=\rho/r$, $u_i=\hat r_i$, $v_i=\hat{\mathsf d}_i$, $\Xi^r_j=q^6\min\{1,u_j\}^3$, $\mathsf z(\beta)=q^{3c(\beta)}(\rho/\mathsf d)^6$ for crossing $\beta$, $\mathsf D(u)=[u/\sqrt3M-3]_+$, $C_\sigma=\sup_u\max\{1,u^6\}E_\sigma(u)$ and $\bar\nu_j=\bar\nu_{j;12}$. There are numbers $C_\omega,\mathsf A_{\mathrm N},C_{\mathrm B},\mathsf A^\times_m,C^\Psi_m,C_\times$ formed from the constants such that for $1\le j\le K-1$:
\begin{enumerate}
\item[(a)] $\mathfrak n[\mathfrak u_{a,j}]\le\varsigma_j\mathsf m_{j,a}$, $\mathfrak n[\dot{\mathfrak e}_{a,j}]\le\epsilon^{(1)}_j\varsigma_j\mathsf m_{j,a}$ with $\epsilon^{(1)}_j=C_\omega(0.9\gamma_j+\mathsf c^\flat\theta^{[0]})\le\frac34$, and $\bar{\mathsf h}_{j;a},\bar\Psi_{j;a}\le\varsigma_j$;
\item[(b)] for a pair, $\bar{\mathsf h}_{j;12}\le\mathsf A_{\mathrm N}\gamma_j\Xi^r_j$ and $\bar\Psi_{j;12}\le2\varsigma_j^2+\mathsf A_{\mathrm N}\gamma_j\Xi^r_j$;
\item[(c)] $\|\mathsf K^{(j+1)}\|_{j+1,\lambda_{\mathrm{out}}}\le\frac12\varsigma_j^2\bar\nu_j$ and $\|\mathsf R^{(j+1)}\|_{j+1,\lambda_{\mathrm{out}}}\le[\frac32\epsilon^{(1)}_j\varsigma_j^2+\mathsf A_{\mathrm N}\gamma_j\Xi^r_j]\bar\nu_j$;
\item[(d)] below the level of the distance: if $1\le n\le K-1$, $u_n\ge\sqrt3M$, and the pairs of all levels $\le n$ of $U$ lie in the cores, then $|\sum_{i\le n}\mathsf{Sum}[E^{(i)}_{12}](U)|\le C_{\mathrm B}q^6e^{-u_n/2\sqrt3M}\bar\nu_n$;
\item[(e)] for crossing $\beta$, $m=|\beta|\le4$, with $\Xi^3_j(\beta)=\mathsf z(\beta)\min\{1,v_j\}^{3/2}$: $\bar{\mathsf h}_{j;\beta}\le\mathsf A^\times_m\gamma_j\Xi^3_j(\beta)$, and $\mathsf N_iE_\sigma(v_i)\bar\Psi_{i-1;\beta}\le108\,C^\Psi_m\mathsf z(\beta)\min\{1,v_i\}^{3/2}E_\sigma(v_i/6)$ for $1\le i\le K$; and at $n=\mathsf k(\mathsf d)$, with $P=v_n$ and $A\in\mathsf{Cc}_n(r_n)$, $|G_{n;\beta}(A)|\le C_\times e^{-P/2\sqrt3M}\mathsf z(\beta)\bar\nu_{n;\beta}$;
\item[(f)] \emph{without the placing:} the bounds (a)--(d), and those of (e) but the last one, at the level $\mathsf k(\mathsf d)$, hold with the same numbers for real weights, none identically $0$, that form one pair, or a pair and a third source, or two pairs, with numbers $0<\rho\le r\le\mathsf d\le1$ such that for $0\le i\le K$: $a_K\le\rho$; (P1) and (P2); the corners of the plaquettes of the supports are at least $u_i/\sqrt3-1$ level-$i$ spacings apart in a cluster and at least $v_i/\sqrt3-1$ across; $\rho\ge\frac23Ma_K$; $E_\sigma(\hat\rho_1/2)\le\varkappa_0^4$; and $\gamma_{K-1}\le\gamma^\natural$. No bound on the distances of the centres and no bound $\mathsf d\le e^{-2\sqrt3}$ is asked in (f); these serve only the frame, the working level, the number $R_w$ and the rings.
\end{enumerate}
\end{lemma}

\begin{proof}
All of it is Lemma~\ref{C-352-12} with (P1)--(P3) of Lemma~\ref{C-352-11}, $\rho\ge\frac23Ma_K$ (so $\varsigma_0=1$), $E_\sigma(\hat\rho_1/2)\le\varkappa_0^4$, $\gamma_j\le\gamma^\natural$, and three inequalities of arithmetic, $\omega^{\mathrm{sp}}_l(i)L^{l+1}\varsigma_{i-1}\le C_\omega\varsigma_{i+l}$, $\omega^{\mathrm{sp}}_l(i)\le C_1L^{-4l}(1+t_i)$ and $(1+t_i)\varsigma_{i-1}^2\le2L^6t_i^2$, which use only $a_K\le\rho$ and (P2). (a) By induction on $j$, with (c) of Lemma~\ref{C-352-12} at $\mathsf{Cr}=\mathsf{Cr}_a$, (e) and the first of these inequalities; then (b) of Lemma~\ref{C-352-12} at $q=1$. (b) In the recursion (g) of Lemma~\ref{C-352-12} the support rule and (P3) give $\sigma_{i;12}\le E_\sigma(u_i)$; the kernel parts are $(1+t_i)2\varsigma_{i-1}^2\le4L^6q^6u_i^6$, and
\[
\sum_{i\le j}L^{-3(j-i)}\min\{u_i^6,C_\sigma\}\le2.002\,C_\sigma\min\{1,u_j\}^3
\]
($u_i=L^{j-i}u_j$: two geometric sums); the term $i=1$ uses $E_\sigma(u_1)\le E_\sigma(u_1/2)E_\sigma(\hat\rho_1/2)\le E_\sigma(u_1/2)\varkappa_0^4$; the feedback $\sum_{i'<j}\gamma_{i'}\le0.45\gamma_j$ is absorbed under $\gamma_j\le\gamma^\natural$. (c) $\mathsf K^{(j+1)}=\mathsf B_j(\mathfrak u_1,\mathfrak u_2)$ and (b) of Lemma~\ref{C-352-12} at $q=2$; for $\mathsf R^{(j+1)}$, Lemma~\ref{C-352-13} with (a), (b) and $\frac12(2+\epsilon)\epsilon\le\frac32\epsilon$. (d) By the support rule $E^{(i)}_{12}$ vanishes on the domains with $d_i(X)<\mathsf D(u_i)$; Lemma~\ref{C-352-14}(b) with (b), $\bar\nu_{i-1}\le\frac43L^{n-i+1}\bar\nu_n$, and $\sum_lL^{l+1}e^{-\frac58\mathsf D(L^lu_n)}\le Ce^{-u_n/2\sqrt3M}$. (e) Induction on $m$ and, for one $\beta$, on the level, through one inequality on the natural sizes $\mu_i(\beta')=\hat\rho_i^{3|\beta'|}u_i^{-3c(\beta')}v_i^{-6[\beta'\text{ crosses}]}$, which is $\hat\rho_i^3$ at a point, $q^6u_i^3$ at a whole pair and $\mathsf z(\beta')$ at a crossing set, at every level (there $3|\beta'|=3c(\beta')+6$): for crossing $\beta$ and a partition $\pi$ of it with $|\pi|\ge2$,
\[
\prod_{\beta'\in\pi}\mu_i(\beta')\le\max\{1,v_i^{12}\}\,\mathsf z(\beta);
\]
the quotient is $u_i^{3s}v_i^{6-6k_\times}$, $s$ the number of whole pairs of $\beta$ that $\pi$ cuts and $k_\times$ the number of its crossing blocks (with $k$ blocks that do not cross, $3k+6k_\times=6+3s$), and $u_i\le v_i$, $k_\times\le2$, and $k_\times=2$ forces $s=2$: at $k_\times=0$ the quotient is $\le v_i^{3s+6}$, at $k_\times=1$ (then $s=k\ge1$) it is $u_i^{3s}\le\max\{1,v_i^6\}$, at $k_\times=2$ it is $(u_i/v_i)^6\le1$. At $v_i<1$ the display gives the bound $\mathsf z(\beta)$ alone (at $k_\times=2$ the quotient is $1$ when $r=\mathsf d$); the factor $\min\{1,v_i\}^{3/2}$ comes, at $k_\times\ge1$, from the crossing blocks, whose factor of the induction is kept, and at $k_\times=0$ from the quotient.

\emph{The blocks.} At a level $i-1\ge1$: a point by (a), $\bar{\mathsf h}_{i-1;a}\le\varsigma_{i-1}\le t_{i-1}=L^3t_i\le L^3\mu_i(\{a\})$; a whole pair $c$ by (b), $\bar{\mathsf h}_{i-1;c}\le\mathsf A_{\mathrm N}\gamma^\natural q^6\min\{1,Lu_i\}^3\le\mathsf A_{\mathrm N}\gamma^\natural L^3\mu_i(c)$; a crossing proper subset $\beta'$ by the first bound of (e) at $|\beta'|<m$, $\bar{\mathsf h}_{i-1;\beta'}\le\mathsf A^\times_{|\beta'|}\gamma^\natural\mathsf z(\beta')\min\{1,Lv_i\}^{3/2}\le\mathsf A^\times_{|\beta'|}\gamma^\natural L^{3/2}\mu_i(\beta')\min\{1,v_i\}^{3/2}$. With $\mathsf a_{(2)}=\max\{1,\mathsf A_{\mathrm N}\gamma^\natural\}$ and $\mathsf a_{(m+1)}=\max\{\mathsf a_{(m)},\mathsf A^\times_m\gamma^\natural\}$, every partition $\pi$ of $\beta$ with $|\pi|\ge2$ has
\[
\prod_{\beta'\in\pi}\bar{\mathsf h}_{i-1;\beta'}\le(L^3\mathsf a_{(m)})^m\min\{1,v_i\}^{3/2}\max\{1,v_i^{12}\}\,\mathsf z(\beta):
\]
if $\pi$ has a crossing block, the factor $\min\{1,v_i\}^{3/2}$ of one of them is kept and the display bounds the rest; if it has none, $k_\times=0$ and the quotient is $\le v_i^{3s+6}\le\min\{1,v_i\}^{3/2}\max\{1,v_i^{12}\}$, as $s\le2$. The block $\pi=\{\beta\}$ is the first bound of (e) for $\beta$ itself at the level $i-1$: $\bar{\mathsf h}_{i-1;\beta}\le\mathsf A^\times_m\gamma^\natural L^{3/2}\mathsf z(\beta)\min\{1,v_i\}^{3/2}$. \emph{The second bound.} Put $T_i=(1+t_i)E_\sigma(v_i)\bar\Psi_{i-1;\beta}$ and $\hat C=\sup_v(1+v^3)\max\{1,v^{12}\}E_\sigma(v/2)$. As $t_i\le\hat\rho_i^3\le v_i^3$, $E_\sigma(v)\le E_\sigma(v/2)^2$ and $E_\sigma(v/2)\le E_\sigma(v/6)^3\le1$, we have $(1+t_i)\max\{1,v_i^{12}\}E_\sigma(v_i)\le\hat C\,E_\sigma(v_i/6)$; so the first recursion of (g) of Lemma~\ref{C-352-12} gives for $2\le i\le K$
\[
\begin{multlined}
T_i\le\hat C\bigl[L^{3/2}\mathsf A^\times_m\gamma^\natural+\mathsf b_m(L^3\mathsf a_{(m)})^m\bigr]\mathsf z(\beta)\min\{1,v_i\}^{3/2}E_\sigma(v_i/6),\qquad\mathsf b_m=\sum_{\pi\in\mathrm{Part}(\beta),\,|\pi|\ge2}(2(|\pi| \\
-1))^{|\pi|-1}=2,22,326 .
\end{multlined}
\]
At $i=1$, (e) of Lemma~\ref{C-352-12} gives $\bar\Psi_{0;\beta}\le\mathsf a^\flat_m:=m^m(\mathsf c^\flat/2)^m$; as $\rho\ge\frac23Ma_K$ and $M\ge L$, $v_1\ge u_1\ge\hat\rho_1\ge\frac23$, so that $\mathsf z(\beta)=\hat\rho_1^{3m}u_1^{-3c(\beta)}v_1^{-6}\ge(\hat\rho_1/v_1)^{3m}\ge(\frac23)^{12}\max\{1,v_1^{12}\}^{-1}$ and $\min\{1,v_1\}^{3/2}\ge(\frac23)^{3/2}$; hence $T_1\le\mathsf a^\flat_m(\frac32)^{27/2}\hat C\,\mathsf z(\beta)\min\{1,v_1\}^{3/2}E_\sigma(v_1/6)$. With $C^\Psi_m=2\hat C[\mathsf b_m(L^3\mathsf a_{(m)})^m+\mathsf a^\flat_m(\frac32)^{27/2}]$, $\mathsf A^\times_m=2.5\,C_1L^2C^\Psi_m$ and $\gamma^\natural\le(5C_1L^{7/2}\hat C)^{-1}$, a smallness of the same kind as the one that absorbs the feedback in (b) and part of the choice of $\gamma^\natural$, the feedback $\hat CL^{3/2}\mathsf A^\times_m\gamma^\natural$ is $\le\frac12C^\Psi_m$, and $T_i\le C^\Psi_m\mathsf z(\beta)\min\{1,v_i\}^{3/2}E_\sigma(v_i/6)$ for $1\le i\le K$; with $\mathsf N_i\le108(1+t_i)$ (P2) this is the second bound of (e). \emph{The first bound.} By the second recursion of (g), with $\sigma_{i;\beta}\le E_\sigma(v_i)$ (the corners across are $v_i/\sqrt3-1$ apart, and $[(v_i/\sqrt3-1)/M-2]_+\ge\mathsf D(v_i)$), $\omega^{\mathrm{sp}}_{j-i}(i)\le C_1L^{-4(j-i)}(1+t_i)$, $v_{j-l}=L^lv_j$ and $\min\{1,L^lv_j\}^{3/2}\le L^{3l/2}\min\{1,v_j\}^{3/2}$:
\[
\bar{\mathsf h}_{j;\beta}\le(\tfrac43)^{\frac m2}C_1L^{\frac m2}\gamma_j\sum_{l=0}^{j-1}L^{(-4+\frac m2)l}T_{j-l}\le(\tfrac43)^{\frac m2}C_1L^{\frac m2}C^\Psi_m\,\gamma_j\,\mathsf z(\beta)\min\{1,v_j\}^{3/2}\sum_{l\ge0}L^{(-4+\frac m2+\frac32)l};
\]
the exponent $-4+\frac m2+\frac32\le-\frac12$ ($m\le4$) gives $\sum_l\le1.39$, and $\bar{\mathsf h}_{j;\beta}\le\mathsf A^\times_m\gamma_j\Xi^3_j(\beta)$. The induction runs over $m=2,3,4$ and, for one $\beta$, over the level: $T_i$ reads the first bounds at the levels $\le i-1$ only. The bound of $G_{n;\beta}$ at the level of $\mathsf d$ is Lemma~\ref{C-352-14}(b) with the factors $e^{-\frac12\mathsf D(v_i)}$, summed over $i\le n$. (f): the proofs of (a)--(d) and of the first two bounds of (e) use of the weights only the properties listed in (f), the separation of the corners standing for (P3); the last bound of (e) uses besides the working level $\mathsf k(\mathsf d)$, which asks $\mathsf d\le e^{-2\sqrt3}$.
\end{proof}

\begin{lemma}[One source of any width]\label{C-352-16}
Let $g$ satisfy $(\Gamma1)$, let the frame at $k_0$ hold for a pair that holds the source $w_a$, and let $w_a$ satisfy (P1), (P2) at a width $\rho$ with $\frac23Ma_K\le\rho$. Then for $k_0\le n\le K$ and $A\in\mathsf{Cc}_n(r_n)$, with $q_{n;a}(F)=\frac1{2a_K}\sum_pw_a(p)|F_p|^2$,
\[
\bigl|G_{n;a}(A)-G_{n;a}(0)-q_{n;a}(\mathsf{ang}^{[n]}(A))\bigr|\le C_{\mathrm S}\,\mathfrak b_n^2\delta_n\,\hat\rho_n^3\,\mathsf u_a .
\]
\end{lemma}

\begin{proof}
At $n=K$, $A=0$ and both members vanish; let $k_0\le n<K$. $G_{n;a}(A)=\mathcal O(w_a)(\tilde U)+\sum_{i\le n}\mathsf{Sum}[E^{(i)}_a](\tilde U)$ and $a_p(\mathbf 1)=0$. The weight member: the expansion of $a_p$ on the chain class and (P1) give $|\mathcal O(w_a)(\tilde U)-q_{n;a}(\mathsf{ang}^{[n]}(A))|\le C r_n^3\hat\rho_n^3x_n\mathsf u_a\le C\mathfrak b_n^2\delta_n\hat\rho_n^3\mathsf u_a$ (Lemma~\ref{C-352-10}). The families: Lemma~\ref{C-352-14}(a), in the frame at $k_0$, at $\mathsf{Cr}_a$, $D'=0$, with $\|E^{(i)}_a\|\le\varsigma_{i-1}\mathsf m_{i-1,a}\mathsf g_i$, $\mathsf g_i=2\gamma_{i-1}$ ($i\ge2$), $\mathsf g_1=\mathsf c^\flat\theta^{[0]}$ (Lemma~\ref{C-352-15}(a)), and $\mathsf N_i\varsigma_{i-1}\le216L^3t_i$; then every level gives the same natural size
\[
t_i\,\mathsf m_{i-1,a}L^{-4(n-i)}=\tfrac{81}2\alpha_0^2L\,\hat\rho_n^3|w_a|_\infty/a_{K-n}
\]
times $\mathsf g_i$, and $\sum_{i\le n}\mathsf g_i\le C\delta_n$, $|w_a|_\infty r_n^2/a_{K-n}\le C\mathfrak b_n^2\mathsf u_a$.
\end{proof}

\paragraph{The kernel families.}

\begin{lemma}[From a real segment to the second order]\label{C-352-17}
Let $R>0$, $0<\mathsf e\le\mathsf M$, and let $\psi$ be holomorphic on $\{|\zeta|<R\}$ with $|\psi|\le\mathsf M$ there, $|\psi(t)|\le\mathsf e$ for real $t\in(-R,R)$, and $\psi(0)=\psi'(0)=0$. With $\Lambda=\max\{1,\log(\mathsf M/\mathsf e)\}$, $|\psi(\zeta)|\le4e\,\mathsf e\,\Lambda^2|\zeta|^2/R^2$ for $|\zeta|\le R/2\Lambda$.
\end{lemma}

\begin{proof}
On the upper half disc $h(\zeta)=\frac2\pi\arg\frac{R+\zeta}{R-\zeta}$ is harmonic, $0$ on the diameter and $1$ on the arc; $\log|\psi|-\log\mathsf e-h\log(\mathsf M/\mathsf e)$ is subharmonic and bounded above, with nonpositive boundary values but at two points, so it is $\le0$ (the two-constants theorem); the same on the lower half disc. As $h(\zeta)\le2|\zeta|/R$, $|\psi|\le e\,\mathsf e$ on $|\zeta|\le R/2\Lambda$; $\psi(\zeta)/\zeta^2$ is holomorphic, and the maximum principle on that disc gives the claim.
\end{proof}

\begin{lemma}[A local gauge]\label{C-352-18}
In units of a level, let the bare lattice have the spacing $\xi\le\frac1{13}$, let $x_0$ be a bare site, $N_s=\{x:|x-x_0|_\infty\le s\}$, $\rho\ge2$ with $4\rho+4$ at most the side of the torus, and let $v=(A,\mathbf J)$ be a real pair of $\mathfrak{su}(2)$-valued $1$-forms on the bare bonds with $|A|\le\alpha\le\frac1{60}$, $|\nabla^\xi A|\le\gamma_\circ$, $|\mathbf J|\le\iota$. With $c_\mu=A_\mu(x_0)$, $X(x)=\sum_\mu(x-x_0)_\mu c_\mu$ and $u=e^{iX}$ on $N_{2\rho+1}$, there is a real pair $v'=(A',\mathbf J')$ which at the bonds with base point in $N_\rho$ is the gauge transform of $v$ under $u$, with $\|v'\|\le\max\{\rho(6\gamma_\circ+19\alpha^2),\iota\}$ in the norm of the chart, $|A'(b)|\le3s\gamma_\circ+(9s+8\xi)\alpha^2$ at the bonds with base point in $N_s$, $s\le\rho$, and $|\nabla^\xi A'|\le\frac{100}{99}(\gamma_\circ+3\alpha^2)$ on $N_\rho$. In the sharp first-difference bound, the whole difference stencil must lie in $N_\rho$; the chart-norm bound is global.

The same holds in a local form, with data on $N_{\rho+1}$ alone: let $U$ be a configuration, $\mathsf g\colon N_{\rho+1}\to SU(2)$ and $A$ a real $\mathfrak{su}(2)$-valued $1$-form on the bonds $b=\langle b_-,b_+\rangle$ with both end points in $N_{\rho+1}$ such that $\mathsf g(b_-)U(b)\mathsf g(b_+)^{-1}=e^{i\xi A(b)}$, $|A|\le\alpha\le\frac1{60}$ and $|\nabla^\xi A|\le\gamma_\circ$ there, the current $\mathbf J(U)$ of $U$ being at most $\iota$ there; with $c_\mu$ and $X$ as above and $h=e^{iX}\mathsf g$ on $N_{\rho+1}$, there is a real pair $v'$ on all the bare bonds which, at the bonds with base point in $N_\rho$, is the image of $U$ and of $\mathbf J(U)$ under the gauge transformation $h$, and which obeys the three bounds above; and if $\|v'\|$ is less than the radius $r_\alpha$ of the ball of the chart, every local term of the level whose domain lies in $N_{\rho-1}$, formed at $U$, has the same value at $P[v']$ restricted to that domain.
\end{lemma}

\begin{proof}
For hermitian $Y,a,c$ with $|a|+|c|\le\frac14$ let $\Psi(Y,a,c)=-i\log(e^{iY}e^{ia}e^{-i(Y+c)})$; matrix calculus gives $|\Psi|\le|a-c|+3|Y||c|+2(|a|+|c|)^2$ and a Lipschitz bound in $(Y,a)$. Put $A^u_\mu(x)=\xi^{-1}\Psi(X(x),\xi A_\mu(x),\xi c_\mu)$, so that $e^{i\xi A^u_\mu(x)}=u(x)e^{i\xi A_\mu(x)}u(x+\xi e_\mu)^{-1}$, a real form; $|X|\le3s\alpha$ and $|A_\mu-c_\mu|\le3s\gamma_\circ$ on $N_s$ along three segments. Cut off with $\chi(x)=\min\{1,\max\{0,2-|x-x_0|_\infty/\rho\}\}$: $A'=\chi A^u$, $\mathbf J'=\chi u\mathbf Ju^{-1}$. The point is that all first difference quotients of $A$ are bounded; lowered by $m$ levels, $|A|$ falls like $L^{-m}$ and $|\nabla A|$ like $L^{-2m}$, and only the constant part is too large, which a gauge linear in $x$ removes at the price $|c|^2$.

In the local form $A^u$ is formed by the same formula from the potential $A$ of the transform of $U$ under $\mathsf g$, so that $e^{i\xi A^u(b)}=h(b_-)U(b)h(b_+)^{-1}$ at the bonds with both end points in $N_{\rho+1}$, and the cut-off $\chi(x)=\min\{1,\max\{0,(\rho+1-2\xi-|x-x_0|_\infty)/(1-2\xi)\}\}$, of the width $1-2\xi$ and equal to $1$ on $N_\rho$, takes the place of the cut-off of the width $\rho$: $A'=\chi A^u$, $\mathbf J'=\chi h\mathbf J(U)h^{-1}$. Both vanish at the bonds whose base point lies at a distance $\ge\rho+1-2\xi$ from $x_0$, and the other bonds have both end points in $N_{\rho+1}$, where $A$, $h$ and the bound $\iota$ are given. As $\nabla^\xi_\nu(\chi A^u)(x)=\chi(x+\xi e_\nu)\nabla^\xi_\nu A^u(x)+A^u(x)(\chi(x+\xi e_\nu)-\chi(x))/\xi$, $|\chi(x+\xi e_\nu)-\chi(x)|\le\frac{13}{11}\xi$, and, on $N_{\rho+1}$, $|\nabla^\xi A^u|\le\frac{100}{99}(\gamma_\circ+3\alpha^2)$ and $|A^u|\le3(\rho+1)\gamma_\circ+(9\rho+9.62)\alpha^2$ (the third and the second bound, at $s=\rho+1$, by the Lipschitz bound of $\Psi$ and the three segments in $N_{\rho+1}$, with $8\xi\le0.62$), for $\rho\ge2$
\[
|\nabla^\xi A'|\le\tfrac{100}{99}(\gamma_\circ+3\alpha^2)+\tfrac{13}{11}\bigl[3(\rho+1)\gamma_\circ+(9\rho+9.62)\alpha^2\bigr]\le\rho(6\gamma_\circ+19\alpha^2),\qquad|A'|\le\rho(6\gamma_\circ+19\alpha^2),
\]
which with $|\mathbf J'|\le\iota$ gives the global chart bound; on stencils contained in $N_\rho$, the cutoff equals $1$ and the sharper local bounds are unchanged. The bonds that a local term with its domain in $N_{\rho-1}$ reads have their end points in $N_\rho$, where $P[v']$ is the image of $U$ under $h$, and a local term is a gauge invariant function of the pair on its domain.
\end{proof}

\begin{proposition}[The kernel families]\label{C-352-19}
Let $g$ satisfy $(\Gamma1)$, let the frame at $k_0$ hold for $(w_1,w_2)$, let the forms $q_n$ be formed over a base level $k_\flat\le k_0$, and let $k_0\le i<n\le K$, $m=n-i-1$, $A\in\mathsf{Cc}_n(r_n)$. As $k_\flat\le i$, $\mathcal S_{i+1}-\mathcal S_i=\upsilon_{i+1}\mathcal C_{(i+1)}$, so that $\mathcal Q_i=q_{i+1}-q_i$ is the slice $\mathcal Q_i(F)=a_K^{-1}\upsilon_{i+1}\sum_{p,q}w_1(p)\mathcal C_{(i+1)}(p,q)w_2(q)\langle F_p,F_q\rangle$. At $n=K$, $A=0$ and the member $\mathsf{LM}$ below is $0$; the numbers below are formed for $n<K$. Put $\mathsf w_a=a_K^{-1}\eta_i^4\sum_p|w_a(p)|$; let $R_w$ be such that every corner of a plaquette of $\operatorname{supp}w_1\cup\operatorname{supp}w_2$ lies within the sup distance $R_w$, in level-$(i+1)$ units, of one bare site; $\alpha=L^{-m}r_n$, $\gamma_\circ=L^{-2m}r_n$, $\mathsf c_i=C_{\mathrm{tp}}e^{-\delta_{\mathrm{tp}}[\mathsf D^{12}_i-1]_+}$, $\mathsf M^{\mathsf K}_i=8M^3C_{\mathrm E,1}\mathsf N_{i+1}(w_1)\|\mathsf K^{(i+1)}\|_{i+1,\lambda_{\mathrm{out}}}$, $\mathsf M_i=2\mathsf M^{\mathsf K}_i+4L^{-4}\mathsf c_ig_i^2\delta_i^2\mathsf w_1\mathsf w_2$, and, with $\varrho_\wedge=\min\{\varrho_{\mathrm{ch}},\varrho_{\mathrm{fl}}\}$, $\epsilon=\epsilon^{\mathrm{cu}}_i$ and $\mathsf c_{\mathbf C}$ of Lemma~\ref{C-352-20}(a), $\epsilon'=\epsilon+2\mathsf c_{\mathbf C}\delta_i/r_\alpha$, the number $C^{\mathrm{row}}_{\mathrm{tp}}$ of Proposition~\YMMref{120.6}{821} that bounds the row sums $\sum_b|\mathsf T^{\mathrm{cv}}_i(p,b)|\le C^{\mathrm{row}}_{\mathrm{tp}}\eta_i^2$ of the matrix $\mathsf T^{\mathrm{cv}}_i$ of step 3 (Proposition~\YMMref{120.6}{821}, equation~\YMMref{(2239)}{821}), $\bar X=(4L^{-2}\mathsf c_{\mathfrak h}+2\mathsf c_{\mathfrak h}^2)\delta_i^2$ and $\bar\rho=\mathsf M_{\mathrm{pl}}^2\delta_i^3/\varrho_\wedge^3$,
\[
\begin{multlined}
\mathsf E_i=12L^{-4}(C^{\mathrm{row}}_{\mathrm{tp}})^2g_i^2\delta_i^2\epsilon'+6\sqrt3L^{-2}(C^{\mathrm{row}}_{\mathrm{tp}})^3g_i^3\delta_i\epsilon(1+2\mathsf c_{\mathbf C}+2\epsilon)+\tfrac94(C^{\mathrm{row}}_{\mathrm{tp}})^4g_i^4\bigl[\epsilon+4\epsilon^2+4(\mathsf c_{\mathbf C} \\
+\epsilon)\epsilon^2+2\epsilon'(2\mathsf c_{\mathbf C}+\epsilon)\bigr]+4\bar X\bar\rho+2\bar\rho^2,
\end{multlined}
\]
and for $\delta_i\le\mathsf r\le\frac12\varrho_{\mathrm{ch}}$ let $\mathsf E_i(\mathsf r)$ be the same number with $g_i^2\mathsf r^2$, $g_i^3\mathsf r$ in the places of $g_i^2\delta_i^2$, $g_i^3\delta_i$, with $\epsilon+2\mathsf c_{\mathbf C}\mathsf r/r_\alpha$ in the place of $\epsilon'$, $4L^{-2}\mathsf c_{\mathfrak h}\mathsf r\delta_i+2\mathsf c_{\mathfrak h}^2\delta_i^2$ in that of $\bar X$ and $\mathsf M_{\mathrm{pl}}^2\max\{\mathsf r/\varrho_{\mathrm{ch}},\delta_i/\varrho_{\mathrm{fl}}\}^3$ in that of $\bar\rho$. All summands are $\ge0$ and the last one is $2\bar\rho^2$, so $2\mathsf M_{\mathrm{pl}}^4\varrho_\wedge^{-6}\delta_i^6\le\mathsf E_i$. As $g_i\le\delta_i$ (Lemma~\ref{C-352-10}(b)) and $\epsilon^{\mathrm{cu}}_i\le1$, $\mathsf E_i$ and, for $\mathsf r\le\delta_i+r_{i+1}$, $\mathsf E_i(\mathsf r)$ are at most $C_{\mathsf E}\delta_i^4(\delta_i+\epsilon^{\mathrm{cu}}_i)$, with a number $C_{\mathsf E}$ formed from the constants; and as $\epsilon^{\mathrm{cu}}_i\le C_{\mathrm{cu}}(\delta_i+\mathsf p_{\star,i})$ (Lemma~\ref{C-352-20}(a)) and $\mathsf p_{\star,i}\le\max\{1,(c_{\mathsf p}/A_1)e^{y_{\mathsf p}/2}\}\delta_i$ at every height (the quotient $\mathsf p_{\star,i}/\delta_i=(c_{\mathsf p}/A_1)y_i^{-\mathfrak p_0}\exp[\frac{y_i}2-(LM)^{-3}A_1^2y_i^{2\mathfrak p_0}]$ is $\le1$ for $y_i\ge y_{\mathsf p}$ by the proof of Lemma~\ref{C-352-10}(b), and $\le(c_{\mathsf p}/A_1)e^{y_{\mathsf p}/2}$ for $1\le y_i<y_{\mathsf p}$; $y_i>\log g^{-2}\ge\log\gamma^{-2}>1$ since $\gamma\le\frac12$), they are at most $C\delta_i^5$.

Let $\bar{\mathsf e}$ be any number $\ge2\mathsf w_1\mathsf w_2\mathsf E_i$; $\Lambda_i=\max\{1,\log(\mathsf M_i/\bar{\mathsf e})\}$, $\mathsf h_i=\max\{2,R_w+1+M(2+\Lambda_i+2\log^+(\delta_i/\gamma_\circ))\}$, $\mathsf s=\mathsf h_i\gamma_\circ(6+19r_n)$. Then $\mathsf{LM}=|\Delta[\mathsf K^{(i+1)}]-\mathcal Q_i(\mathsf{ang}^{[n]}(A))|$ satisfies:
\begin{enumerate}
\item[(i)] always $\mathsf{LM}\le2\mathsf w_1\mathsf w_2\mathsf E_i(\max\{\delta_i,\alpha\})$, where $\max\{\delta_i,\alpha\}\le\min\{\frac12\varrho_{\mathrm{ch}},\delta_i+r_{i+1}\}$ ($\alpha\le r_{i+1}$, as $r_j\ge L^{-1}r_{j+1}$; $2\delta_i\le\varrho_\wedge$ and $2r_{i+1}\le\varrho_{\mathrm{ch}}$ in the frame, and under $(\Gamma1)$ by Lemma~\ref{C-352-10}(b)), so that $\mathsf E_i(\max\{\delta_i,\alpha\})\le C\delta_i^5$;
\item[(ii)] if $\mathsf s\le\delta_i/2\Lambda_i$: $\mathsf{LM}\le8e\,\bar{\mathsf e}\Lambda_i^2\mathsf s^2/\delta_i^2+L^{-4}\mathsf c_ig_i^2\mathsf w_1\mathsf w_2[18(\alpha^2+\alpha'^2)(\gamma_\circ+\gamma')+48R_w\alpha\gamma_\circ^2]$, with $\alpha'=(R_w+1)(3\gamma_\circ+10\alpha^2)$, $\gamma'=\frac{100}{99}(\gamma_\circ+3\alpha^2)$.
\end{enumerate}
The weights enter only through $\mathsf w_a$, $R_w$, $\mathsf D^{12}_i$, $\mathsf N_{i+1}(w_1)$ and $\|\mathsf K^{(i+1)}\|$; no test pair, filling or sign is asked. For weights placed at $(\rho,r)$ as in (B.2), with $k_0=\mathsf k(r)$, this gives for $\mathsf k(r)\le i<n\le K$
\[
\mathsf{LM}\le\bigl\{C_{\mathrm K}\mathfrak b_i^4\delta_i\Lambda_\sharp^2\mathsf h_\sharp^2u_i^3+C'_{\mathrm K}\mathsf c_i(1+u_i)^2u_i^3\,\mathfrak b_n^2r_n\bigr\}(\rho/r)^6\hat r_n^3\mathsf u_1\mathsf u_2,\qquad\Lambda_\sharp\le CY_r,\quad\mathsf h_\sharp\le CY_r^{2\mathfrak p_0},
\]
and neither $\Lambda_\sharp$ nor $\mathsf h_\sharp$ depends on $\rho$. Of $(\Gamma1)$ the proofs of (i) and (ii) use only the frame at $k_0$ and, at the heights concerned, $2r_n\le\varrho_{\mathrm{ch}}$ (so $r_n\le\frac1{120}$, as $\varrho_{\mathrm{ch}}\le\frac1{60}$ by equation~\YMMref{(2285)}{833}), $\epsilon^{\mathrm{cu}}_i\le1$ and $x_n^{-1}\le r_n^2$, besides the bound $\ell\ge2LMy_K$ of the sides $\ell\in\mathfrak L(g)$; (i) and (ii), with the bounds of $\mathsf E_i$ and $\mathsf E_i(\cdot)$ stated above, hold whenever these hold. The narrow display holds under the data of (B.2), where step 6 reads besides Lemma~\ref{C-352-15}(c) and $\delta_i\le\delta_n$.
\end{proposition}

\begin{proof}
1. \emph{Lowering.} $L^{-m}A$, viewed at the level $i+1$, is a real point $v$ of the chart with $|A_m|\le\alpha$, $|\nabla A_m|\le\gamma_\circ$, current $\le\gamma_\circ$, and $\mathsf{ang}^{[i+1]}(A_m)=\mathsf{ang}^{[n]}(A)$; $\mathsf{Sum}[\mathsf K^{(i+1)}](\tilde U)=\mathcal K_i(v):=\mathsf{Sum}[\mathsf K^{(i+1)}](P[v])$, and $\mathcal Q_i(v)$ is the third term.

2. \emph{Holomorphy and no first order.} $\mathcal K_i$ is holomorphic on the ball of radius $r_\alpha$ of the chart, $|\mathcal K_i|\le\mathsf M^{\mathsf K}_i$, and it is invariant under constant gauge rotations; as $\sum_{k=1}^3\operatorname{Ad}(i\sigma_k)=-\mathrm{id}$, its first derivative at $0$ vanishes.

3. \emph{The real ball.} At the real points $\|v\|\le\delta_i$, $\mathcal K_i(v)$ is a covariance under the probability measure $\nu_{i,v}$ of one step of the plaquette terms at the lowered configuration; this rests on a sum rule at every point of the chart, which (a) and (b) give.

(a) \emph{The sum rule at every point.} Proposition~\YMMref{51.15}{410} proves the sum rule of one local term at the two admitted background points, in equation~\YMMref{(872)}{411}, and from it the sum rule of the totals at the chart points $v_A$ of the chain classes (Proposition~\YMMref{124.2}{834}, equation~\YMMref{(2289)}{834}). The finite-vertex argument in that proof uses only membership in the localization domain $\mathfrak D(T_\ell)$ until it identifies the two endpoint values with background values; so for a term $f$ on a domain $X$ whose radius exceeds $1$, as do those of the localisation of bare arrays, and every point $\mathfrak u=(P,\mathcal A)\in\mathfrak D(T_\ell)$, critical or not,
\[
\sum_{Y'\in\mathcal D_i:\,Y'\supset X^\uparrow}\mathbb V[f]\bigl(Y';\mathfrak u|_{Y'}\bigr)=f\bigl(\rho_XP\dotplus\mathfrak h_{T_\ell}(\mathbf 1;\mathfrak u)\bigr)-f(\rho_XP).
\]
Indeed, with the interpolating function $F(s)=\Phi_{T_\ell}[f](s,1;\mathfrak u)$ of equation~\YMMref{(869)}{410}, including its distinct rescaling factors for the potential and the current, the resummation of brackets (\eqref{eq:analytic-bracket-resummation}) writes $F(\mathbf 1)-F(0)$ as the sum of the brackets $F^*[\mathcal K|\mathsf N_0]$ over the sets $\mathcal K$ of fine cubes off the cubes $\mathsf N_0$ of $X^\uparrow$; a bracket for which $\mathcal K\cup\mathsf N_0$ has more than one wall-component vanishes, $F$ being local for the roots $\mathsf N_0$ (Proposition~\YMMref{51.15}{410} and Lemma~\ref{lem:analytic-rooted-bracket}); if it has one, its union is a region $Y'\supset X^\uparrow$ and the bracket is $\mathbb V[f](Y';\mathfrak u|_{Y'})$ by the zero-extension consistency in Proposition~\YMMref{51.15}{410}, and $\mathcal K\mapsto Y'$ is a bijection onto these regions; finally $F(\mathbf 1)=f(\rho_XP\dotplus\mathfrak h_{T_\ell}(\mathbf 1;\mathfrak u))$ by the definition of $F$, and $F(0)=f(\rho_XP)$ by the zero-value statement in Proposition~\YMMref{51.15}{410}. Identifying the endpoint values with the two background values gives equation~\YMMref{(872)}{411}; the finite-bracket identity itself requires only the stated domain membership.

(b) \emph{The totals at every point of the chart.} The weight member of a weight $w$ is, on a domain $X$ of level $i$, $\eta_i^{-1}\sum_{p:X(p)=X}w(p)a_p$, every bare plaquette $p$ having one domain $X(p)$; let $w=w^{[i]}_a=L^{K-i}w_a$, so that $\eta_i^{-1}L^{K-i}=L^K$, and let $\mathfrak f_i[\mathsf W_{w^{[i]}_a}]$ be the bare array whose only member is this weight member. Applied to its members at the point $\mathfrak u=(P[v],g_iB)$, with the pair identity in equation~\YMMref{(2286)}{833} for the configuration of the substituted pair, (a) gives at every point $v=(A,\mathbf J)$ of the chart of level $i+1$ with $\|v\|<\varrho_{\mathrm{ch}}$, real or not, in a chain class or not, and every $B\in\mathsf{Rg}_{x_i}$ (as $\delta_i<\varrho_{\mathrm{fl}}$)
\[
\mathbb V'\bigl[\mathfrak f_i[\mathsf W_{w^{[i]}_a}]\bigr]^{\mathrm{tot}}(P[v];B)=L^K\sum_pw_a(p)\Bigl[a_p\bigl(e^{i\eta_iA^{\mathrm{lw}}(v,g_iB)}\bigr)-a_p\bigl(e^{i\eta_{i+1}A}\bigr)\Bigr],
\]
This proves at every such chart point the algebraic identity whose real-background form is equation~\YMMref{(2290)}{834}.

(c) \emph{The covariance.} By Lemma~\ref{lem:weighted-step-integrals} and equation~\YMMref{(2298)}{837}, $\mathsf{Sum}[\mathsf K^{(i+1)}](P)$ is, at every triple $P$ of the core, the covariance under $\nu_i[P]$ of the two totals of (b); the terms $a_p(e^{i\eta_{i+1}A})$ are constants. So at every real $v$ with $\|v\|\le\delta_i$
\[
\mathcal K_i(v)=L^{2K}\sum_{p,q}w_1(p)w_2(q)\bigl\langle a_p(\tilde U^{\mathrm{lw}});a_q(\tilde U^{\mathrm{lw}})\bigr\rangle_{\nu_{i,v}},\qquad\tilde U^{\mathrm{lw}}(B)=e^{i\eta_iA^{\mathrm{lw}}(v,g_iB)},
\]
where $\nu_{i,v}$ is a probability measure on $\mathsf{Rg}_{x_i}$ and $|\langle F;F'\rangle_{\nu_{i,v}}|\le2\sup|F|\sup|F'|$ (Proposition~\YMMref{124.5}{836}).

(d) \emph{The plaquette terms.} $(v,g_iB)$ is a real point with $\|v\|\le\delta_i$ and $\sup|g_iB|<g_i\mathfrak b_i=\delta_i$. The linear part of the lowered angle in the fluctuation is $\mathsf{sg}_i\sum_b\mathsf T^{\mathrm{cv}}_i(p,b)\mathcal A(b)$, where $\mathsf T^{\mathrm{cv}}_i=(\mathsf T^{\mathrm{cv}}_i(p,b))$ is the real matrix, indexed by the bare plaquettes and the positively oriented fluctuation bonds, of the curvature response of one step composed with its elimination map (the definitions preceding Lemma~\YMMref{124.3}{835}), $\mathsf{sg}_i\in\{1,-1\}$ is the sign of Lemma~\YMMref{124.3}{835}, and $\sum_b|\mathsf T^{\mathrm{cv}}_i(p,b)|\le C^{\mathrm{row}}_{\mathrm{tp}}\eta_i^2$ (Proposition~\YMMref{120.6}{821}, equation~\YMMref{(2239)}{821}). Lemma~\YMMref{124.4}{835} at $r=a=\delta_i$, $\mathsf{rt}=\varrho_\wedge/\delta_i\ge2$, gives
\[
a_p(\tilde U^{\mathrm{lw}}(B))=\tfrac12|F_p|^2+\langle F_p,\xi_p\rangle+\tfrac12\langle\xi_p,\xi_p\rangle+\mathsf o_p(B),\qquad\xi_p=\mathsf{sg}_ig_i\sum_b\mathsf T^{\mathrm{cv}}_i(p,b)B(b),
\]
with $F_p=\mathsf{ang}^{[i+1]}_p(A)$, $|\mathsf o_p|\le\mathsf M_{\mathrm{pl}}^2\eta_i^4\mathsf{rt}^{-3}=\eta_i^4\bar\rho$, $|F_p|\le2L^{-2}\eta_i^2\delta_i$ and $|\xi_p|\le2\mathsf c_{\mathfrak h}\eta_i^2\delta_i$. So $\mathsf X_p=\langle F_p,\xi_p\rangle+\frac12\langle\xi_p,\xi_p\rangle$ has $\sup|\mathsf X_p|\le\eta_i^4\bar X$, and $\langle a_p;a_q\rangle_{\nu_{i,v}}$ differs from $\langle\mathsf X_p;\mathsf X_q\rangle_{\nu_{i,v}}$ by at most $\eta_i^8(4\bar X\bar\rho+2\bar\rho^2)$.

(e) \emph{The Gaussian part.} Under $\nu_{i,v}$ the joint cumulants of the orders $1$, $3$, $4$ of the components of $B$ are at most $\epsilon=\epsilon^{\mathrm{cu}}_i$, and those of the order $2$ lie within $\epsilon'=\epsilon+2\mathsf c_{\mathbf C}\delta_i/r_\alpha$ of $\mathbf C_0$, the covariance of the step at $v=0$ (Proposition~\YMMref{124.5}{836}, equations~\YMMref{(2296)}{836}--\YMMref{(2297)}{836}, at $\|v\|\le\delta_i$); and $\mathbf C_0=(\mathsf A^{\mathrm{cv}}_i)^{-1}\otimes1$ with $\mathsf T^{\mathrm{cv}}_i(\mathsf A^{\mathrm{cv}}_i)^{-1}(\mathsf T^{\mathrm{cv}}_i)^{\mathsf T}=\eta_i\mathcal C_{(i+1)}$, $\mathsf A^{\mathrm{cv}}_i$ the real symmetric matrix of one step at the trivial configuration on the fluctuation bonds (equation~\YMMref{(2297)}{836}).

A covariance of products is the sum, over the partitions that are no refinements of the two factors, of the products of the cumulants of the blocks. With $\mathbf C_0$ in the place of the second cumulants, the part of $\langle\langle F_p,\xi_p\rangle;\langle F_q,\xi_q\rangle\rangle$ is $g_i^2\eta_i\mathcal C_{(i+1)}(p,q)\langle F_p,F_q\rangle$, that of $\frac14\langle\langle\xi_p,\xi_p\rangle;\langle\xi_q,\xi_q\rangle\rangle$ is $\frac32g_i^4\eta_i^2\mathcal C_{(i+1)}(p,q)^2$ (three colours, two pairings), the mixed covariances $\langle\langle F_p,\xi_p\rangle;\frac12\langle\xi_q,\xi_q\rangle\rangle$ have none, and the rest is at most $\eta_i^8$ times the first three summands of $\mathsf E_i$, by the row sums.

Multiplied by $L^{2K}w_1(p)w_2(q)$ and summed ($L^{2K}\eta_i^8\sum_p|w_1(p)|\sum_q|w_2(q)|=\mathsf w_1\mathsf w_2$, $L^{2K}g_i^2\eta_i=a_K^{-1}\upsilon_{i+1}$), this gives, with $\mathsf k_i=\frac32g_i^4\eta_i^2L^{2K}\sum_{p,q}w_1(p)w_2(q)\mathcal C_{(i+1)}(p,q)^2$ and $\mathcal Q_i(v)=\mathcal Q_i(\mathsf{ang}^{[i+1]}(A))$,
\[
|\mathcal K_i(v)-\mathsf k_i-\mathcal Q_i(v)|\le\mathsf w_1\mathsf w_2\mathsf E_i .
\]

(f) \emph{Why a ball.} A ball of the chart is closed under $v\mapsto\zeta v$, under cut-offs and under small gauge transformations, which steps 4 and 5 use; a chain class is closed under none of the three, and this is why the totals are taken at every point of the chart. The radius $\delta_i$ enters only the expansion, the bound of $|F_p|$ and the ratio $\|v\|/r_\alpha$; read at a radius $\delta_i\le\mathsf r\le\frac12\varrho_{\mathrm{ch}}$ ($a=\delta_i$, $\mathsf{rt}=\min\{\varrho_{\mathrm{ch}}/\mathsf r,\varrho_{\mathrm{fl}}/\delta_i\}\ge2$), the same argument gives $|\mathcal K_i(v)-\mathsf k_i-\mathcal Q_i(v)|\le\mathsf w_1\mathsf w_2\mathsf E_i(\mathsf r)$ at the real points with $\|v\|\le\mathsf r$; at $\mathsf r=\max\{\delta_i,\alpha\}$, at the point $v$ of step 1 and at $0$, this is (i).

4. \emph{Two constants.} For real $v'$ with $\|v'\|\le\delta_i/2\Lambda_i$ put $\psi(\zeta)=\mathcal K_i(\zeta v')-\mathcal K_i(0)-\zeta^2\mathcal Q_i(v')$: holomorphic for $|\zeta|<\delta_i/\|v'\|$, $\psi(0)=\psi'(0)=0$, $|\psi|\le\mathsf M_i$, $|\psi|\le\bar{\mathsf e}$ on the real segment (step 3 at $tv'$ and at $0$); Lemma~\ref{C-352-17} gives $|\psi(1)|\le4e\bar{\mathsf e}\Lambda_i^2\|v'\|^2/\delta_i^2$ (if $\bar{\mathsf e}>\mathsf M_i$, then $|\psi|\le\mathsf M_i<\bar{\mathsf e}$ on the disc and Schwarz's lemma of the order two gives more). Of the complex ball only the bound $\mathsf M_i$ is used, through $\log\mathsf M_i$; the smallness comes from the real ball, where the step measure is a probability measure, and any $\bar{\mathsf e}\ge2\mathsf w_1\mathsf w_2\mathsf E_i$ serves.

5. \emph{The gauge.} Let first $4\mathsf h_i+4$ be at most the side $\ell L^{K-i-1}$ of the torus in level-$(i+1)$ units. Lemma~\ref{C-352-18} on the box $N_{\mathsf h_i}$ about the pair gives $v'$ with $\|v'\|\le\mathsf s$; the domains in the box take the same values at $v$ and at $v'$, those outside have $d_{i+1}(X)\ge(\mathsf h_i-1-R_w)/M-2$ and cost $2\mathsf M^{\mathsf K}_ie^{-[\cdot]_+}\le\bar{\mathsf e}\gamma_\circ^2/\delta_i^2$ by the choice of $\mathsf h_i$; replacing $\mathcal Q_i(v')$ by $\mathcal Q_i(v)$ costs the last term of (ii), since $|\mathcal C_{(i+1)}|\le\mathsf c_i\eta_i^3$ on the supports. If $4\mathsf h_i+4>\ell L^{K-i-1}$, then $L^m\le L^{K-i-1}<(4\mathsf h_i+4)/\ell\le\mathsf h_i$, as $\ell\ge8$ and $\mathsf h_i\ge2$; so $\|v\|\le\max\{\alpha,\gamma_\circ\}=\alpha=L^m\gamma_\circ\le\mathsf h_i\gamma_\circ\le\mathsf s$, and step 4 applies at $v'=v$ itself, with no gauge and no domains outside a box: (ii) holds without its last term.

6. \emph{The narrow pair.} Take $\bar{\mathsf e}=2\bar{\mathsf w}_1\bar{\mathsf w}_2\mathsf E_i$ with $\bar{\mathsf w}_a=81\hat\rho_i^3|w_a|_\infty/a_{K-i}\ge\mathsf w_a$. By Lemma~\ref{C-352-15}(c), $\|\mathsf K^{(i+1)}\|\le\frac12t_i^2\bar\nu_i=C x_i^{-1}\bar{\mathsf w}_1\bar{\mathsf w}_2$: the factor $\hat\rho_i^6$ cancels in $\mathsf M_i/\bar{\mathsf e}\le C(1+\mathsf N_{i+1})x_i^{-1}\delta_i^{-6}\le C(1+u_{\mathsf k(r)}^3)x_{\mathsf k(r)}^2$ (by the lower bound $2\mathsf M_{\mathrm{pl}}^4\varrho_\wedge^{-6}\delta_i^6\le\mathsf E_i$), so $\Lambda_i\le CY_r$ by the working level; $R_w\le11u_{i+1}$ and $\log^+(\delta_i/\gamma_\circ)\le2m\log L\le2Y_r$ give $\mathsf h_i\le CY_r^{2\mathfrak p_0}$. If $\mathsf s\le\delta_i/2\Lambda_i$, (ii) with $\bar{\mathsf w}_1\bar{\mathsf w}_2x_i^{-1}L^{-4m}r_n^2\le C\mathfrak b_n^2u_i^3(\rho/r)^6\hat r_n^3\mathsf u_1\mathsf u_2$; else $L^{3(n-i)}\le C\mathsf h_i^2\Lambda_i^2$, and (i) gives the first term.
\end{proof}

\paragraph{The induction over the levels; the jets.}

\begin{lemma}[One step: cumulants, quadratic polynomials, the Gaussian step]\label{C-352-20}
\begin{enumerate}
\item[(a)] In the data of one renormalisation step at a height $i$, for a real point $v$ of the chart of the level $i+1$ with $\|v\|<r_\alpha$, every $n\ge1$ and components $B^{\alpha_1}(b_1),\dots,B^{\alpha_n}(b_n)$ of the fluctuation field, repetitions allowed, the joint cumulants under the measure $\nu_{i,v}$ of the step are formed and satisfy $|\kappa(B^{\alpha_1}(b_1),\dots)|\le\epsilon^{\mathrm{cu},n}_i$ for $n\ne2$ and $|\kappa(B^{\alpha_1}(b_1),B^{\alpha_2}(b_2))-\mathbf C_v((b_1,\alpha_1),(b_2,\alpha_2))|\le\epsilon^{\mathrm{cu},2}_i$, where $\epsilon^{\mathrm{cu},n}_i=C_{\mathrm E,0}(K_\Pi\mathsf n^{\mathbb V}_i+K^\Pi_P\mathsf p_{\star,i})(n\mathsf n_{\mathrm f}/w_\star)^n$ with construction constants, $\mathsf n^{\mathbb V}_i\le c_{\mathbb V}\delta_i$; $\mathbf C_v$ is the covariance of the step, $|\mathbf C_v|\le\mathsf c_{\mathbf C}$ and $|\mathbf C_v-\mathbf C_0|\le2\mathsf c_{\mathbf C}\|v\|/r_\alpha$ entry by entry, $\mathbf C_0=(\mathsf A^{\mathrm{cv}}_i)^{-1}\otimes1$, and $\mathsf T^{\mathrm{cv}}_i(\mathsf A^{\mathrm{cv}}_i)^{-1}(\mathsf T^{\mathrm{cv}}_i)^{\mathsf T}=\eta_i\mathcal C_{(i+1)}$ (at $n\le4$ the bounds on the cumulants are those of the proof of Proposition~\YMMref{124.5}{836}, and the statements on $\mathbf C_v$ and $\mathbf C_0$ are in Proposition~\YMMref{124.5}{836}).

Here $\mathsf A^{\mathrm{cv}}_i$ is the real symmetric matrix of one step at the trivial configuration on the positively oriented fluctuation bonds $b\in\mathfrak X^+$, positive definite as $\mathbf C_0$ is (Proposition~\YMMref{124.5}{836} at $v=0$); $\mathsf T^{\mathrm{cv}}_i=(\mathsf T^{\mathrm{cv}}_i(p,b))$ is the real matrix, indexed by the bare plaquettes and the bonds of $\mathfrak X^+$, of the curvature response of one step composed with its elimination map (the definitions preceding Lemma~\YMMref{124.3}{835}); and there are a sign $\mathsf{sg}_i\in\{1,-1\}$ and the number $C^{\mathrm{row}}_{\mathrm{tp}}$ of Proposition~\YMMref{120.6}{821} such that the linear part of the lowered angle in a fluctuation $\mathcal A$ is $\mathsf{sg}_i\sum_b\mathsf T^{\mathrm{cv}}_i(p,b)\mathcal A(b)$ and $\sum_b|\mathsf T^{\mathrm{cv}}_i(p,b)|\le C^{\mathrm{row}}_{\mathrm{tp}}\eta_i^2$ (Lemma~\YMMref{124.3}{835} and Proposition~\YMMref{120.6}{821}, equation~\YMMref{(2239)}{821}, on a side $\ell$ at least the constant $N_{\mathsf H}$ fixed in Section~\YMMref{120.3}{820}, as in the frame).

At a real point $(v,g_iB)$, $v=(A,\mathbf J)$, $B\in\mathsf{Rg}_{x_i}$, with $r\ge\|v\|$ and $\mathsf{rt}=\min\{\varrho_{\mathrm{ch}}/r,\varrho_{\mathrm{fl}}/\delta_i\}\ge2$, Lemma~\YMMref{124.4}{835} (at $a=\delta_i$) splits the angle of the lowered configuration:
\[
\mathsf{ang}^{[i]}_p\bigl(A^{\mathrm{lw}}(v,g_iB)\bigr)=\mathsf{ang}^{[i+1]}_p(A)+\psi_p(B)+\zeta_p(B),\qquad\psi_p(B)=\mathsf{sg}_ig_i\sum_b\mathsf T^{\mathrm{cv}}_i(p,b)B(b),
\]
$|\psi_p|\le2\mathsf c_{\mathfrak h}\eta_i^2\delta_i$, $|\zeta_p|\le4\mathsf M_{\mathrm{gr}}\eta_i^2\mathsf{rt}^{-2}$, with the number $\mathsf M_{\mathrm{gr}}$ of equation~\YMMref{(2285)}{833}; under the centred Gaussian with covariance $\mathbf C_0$, $\psi$ has the covariance $g_i^2\eta_i\mathcal C_{(i+1)}=a_K\upsilon_{i+1}\mathcal C_{(i+1)}$, as $\mathsf{sg}_i^2=1$ and $g_i^2\eta_i=a_K\upsilon_{i+1}$. At $v=v_{A'}$, $A'\in\mathsf{Cc}_{i+1}(r_{i+1})$, one takes $r=r_{i+1}$; under $(\Gamma1)$ Lemma~\ref{C-352-10}(b) gives $\mathsf{rt}\ge2$ and $|\zeta_p|\le C\eta_i^2\delta_i^2$. We put $\epsilon^{\mathrm{cu}}_i=\epsilon^{\mathrm{cu},4}_i$; with $C_\star=C_{\mathrm E,0}\max\{K_\Pi c_{\mathbb V},K^\Pi_P\}$, so that $\epsilon^{\mathrm{cu},n}_i\le C_\star(\delta_i+\mathsf p_{\star,i})(n\mathsf n_{\mathrm f}/w_\star)^n$, we take $C_{\mathrm{cu}}=C_\star(4\mathsf n_{\mathrm f}/w_\star)^4$, $c^{(n)}_{\mathrm{cu}}=2C_\star(2n\mathsf n_{\mathrm f}/w_\star)^{2n}$ and $c_{\mathrm{cu}}=c^{(3)}_{\mathrm{cu}}$, so that $\epsilon^{\mathrm{cu}}_i\le C_{\mathrm{cu}}(\delta_i+\mathsf p_{\star,i})$ and $\epsilon^{\mathrm{cu},2n}_i\le c^{(n)}_{\mathrm{cu}}(\delta_i+\mathsf p_{\star,i})/2$; we put $\epsilon'^{(n)}_i=\epsilon^{\mathrm{cu},2n}_i+2\mathsf c_{\mathbf C}r_{i+1}/r_\alpha$ and $\epsilon'_i=\epsilon'^{(3)}_i$, and $c^{(n)}_\epsilon$ is the number $c_\epsilon$ of Lemma~\ref{C-352-10}(b) with $c^{(n)}_{\mathrm{cu}}$ in place of $c_{\mathrm{cu}}$.

We have $4\mathsf n_{\mathrm f}/w_\star>1$ ($\mathsf n_{\mathrm f}=\frac12M^{3/2}$ and $w_\star<1$: equation~\YMMref{(2294)}{836} and the proof of Proposition~\YMMref{124.5}{836}); so with $x=\mathsf n_{\mathrm f}/w_\star$, $(nx)^n\le(n'x)^n\le(n'x)^{n'}$ for $n\le n'$, $n'\ge4$, and $\epsilon^{\mathrm{cu},n}_i\le\epsilon^{\mathrm{cu},2m}_i$ for $n\le2m$, $m=2,3,4$: at $\|v\|\le r_{i+1}$ the joint cumulants of all the orders $\le2m$ lie within $\epsilon'^{(m)}_i$ of those of the centred Gaussian with covariance $\mathbf C_0$ (at the order two through $|\mathbf C_v-\mathbf C_0|\le2\mathsf c_{\mathbf C}r_{i+1}/r_\alpha$). Besides $\epsilon'^{(2)}_i\le\epsilon'_i\le\epsilon'^{(4)}_i$ and $c^{(2)}_\epsilon\le c_\epsilon\le c^{(4)}_\epsilon$; so $(\Gamma1)$ with $c^{(4)}_\epsilon$ in place of $c_\epsilon$ contains $(\Gamma1)$, and $(\Gamma^{(4)})$ contains $(\Gamma)$.
\item[(b)] Let $I$ be finite, $\nu$ a probability measure with bounded support on $\mathbb R^I$, $\mathbf C$ a real symmetric positive semidefinite matrix on $I$ with $|\mathbf C(\mathsf i,\mathsf j)|\le\mathsf c$, $\mathsf c\ge1$, $0\le\epsilon\le1$, and let the cumulants of $\nu$ of orders $n\le6$ be within $\epsilon$ of those of the centred Gaussian $\mathsf G$ with covariance $\mathbf C$. For quadratic polynomials $P_1,P_2,P_3$ in the coordinates, with $\|P\|_{\mathrm{pol}}$ the sum of the moduli of their coefficients of orders one and two, and $1\le m\le3$: $|\kappa_\nu(P_1,\dots,P_m)-\kappa_{\mathsf G}(P_1,\dots,P_m)|\le c_6\epsilon\prod_l\|P_l\|_{\mathrm{pol}}$, $c_6=1218(\mathsf c+1)^5$. If only the cumulants of the orders $n\le4$ are within $\epsilon$ of those of $\mathsf G$, the same holds for $m\le2$ with $c_4=60(\mathsf c+1)^3$; if those of the orders $n\le8$ are, it holds for $m\le4$ (four sources) with $c_8=33120(\mathsf c+1)^7$; and $c_4\le c_6\le c_8$.
\item[(c)] Let $k_\flat\le i<K$, $\mathsf s=\upsilon_{i+1}\mathcal C_{(i+1)}$, and let $\psi$ be a centred Gaussian vector of $\mathfrak{su}(2)$-valued functions of the plaquettes with $\mathbb E\psi^\alpha_p\psi^{\alpha'}_q=a_K\delta_{\alpha\alpha'}\mathsf s(p,q)$. For every $F'$ and every nonempty set $\beta$ of sources,
\[
\sum_{\pi\in\mathrm{Part}(\beta)}\kappa\bigl(q_{i;\beta'}(F'+\psi):\beta'\in\pi\bigr)=(g_{i+1;\beta}-g_{i;\beta})+q_{i+1;\beta}(F').
\]
\end{enumerate}
\end{lemma}

\begin{proof}
(a) The bound on the cumulants of one step (the proof of Proposition~\YMMref{124.5}{836}) is proved with $n$ insertions of one sort, and depends on the number $n$ only through the factor $(n\mathsf n_{\mathrm f}/w_\star)^n$; so it holds at every order. The remaining statements in (a) follow from the results cited there, and the comparison of the orders is the arithmetic stated in (a). (b) The cumulant of products $X_{S_1},\dots,X_{S_m}$ is the sum, over the partitions $\sigma$ of $S=\bigsqcup S_l$ with $\sigma\vee\{S_1,..,S_m\}=\{S\}$, of the products of the cumulants of the blocks (from the expansions of $\exp$ and $\log$ in $\mathbb C[z]/(z_l^2)$). Replacing the factors one by one costs $|\sigma|\epsilon(\mathsf c+1)^{|\sigma|-1}\le2m\epsilon(\mathsf c+1)^{2m-1}$, as $|S|\le2m$, and a set of $2m$ elements has $B_{2m}$ partitions, $B_4=15$, $B_6=203$, $B_8=4140$: so $c_{2m}=2m\,B_{2m}(\mathsf c+1)^{2m-1}$ serves, which is $c_6=6\cdot203(\mathsf c+1)^5$ for $m\le3$, $c_4$ for $m\le2$ and $c_8$ for $m\le4$. (c) Put $w_\lambda=\sum_a\lambda_aw_a$, $R^\lambda_i=w_\lambda(1-\mathcal S_iw_\lambda)^{-1}$, $g^\lambda_i=-\frac32\operatorname{Tr}\log(1-\mathcal S_iw_\lambda)$ and $q^\lambda_i(F)=\frac1{2a_K}\langle F,R^\lambda_iF\rangle$; their coefficients of $\lambda^\beta=\prod_{a\in\beta}\lambda_a$ are $R_{i;\beta}$, $g_{i;\beta}$ and $q_{i;\beta}$.

The Gaussian integral, with the factor $3$ from the dimension of $\mathfrak{su}(2)$, is
\[
\log\mathbb E\,e^{q^\lambda_i(F'+\psi)}=-\tfrac32\operatorname{Tr}\log(1-\mathsf sR^\lambda_i)+\tfrac1{2a_K}\langle F',R^\lambda_i(1-\mathsf sR^\lambda_i)^{-1}F'\rangle,
\]
and $1-\mathsf sR^\lambda_i=(1-\mathcal S_{i+1}w_\lambda)(1-\mathcal S_iw_\lambda)^{-1}$, $R^\lambda_i(1-\mathsf sR^\lambda_i)^{-1}=R^\lambda_{i+1}$, as $\mathcal S_{i+1}=\mathcal S_i+\mathsf s$. The coefficient of $\lambda^\beta$ of the left member is the sum over the partitions of the joint cumulants; that of the right member is $(g_{i+1;\beta}-g_{i;\beta})+q_{i+1;\beta}(F')$. This is Wick's rule for the photon fields of one step.
\end{proof}

\begin{lemma}[The tower identity at several sources]\label{C-352-20a}
Let $0<g\le\gamma_{\mathrm{jt}}$, $g^{-2}\ge x_{\mathrm J}$, $(K,\beta)$ a first-passage datum at $g$, $\ell\in\mathfrak L(g)$, $n\ge2$, and let $w_a=f_a\circ x\not\equiv0$, $a\in[n]$, with real Lipschitz functions $f_a$ on $T_\ell$ with pairwise disjoint supports (for instance the tents of Lemma~\ref{C-352-11}(b)) that fall into two nonempty clusters $[n]=I_1\sqcup I_2$ lying in the frame at a level $k$. Let $k\le i<K$, $A'\in\mathsf{Cc}_{i+1}(r_{i+1})$, $\tilde U=e^{i\eta_{i+1}A'}$, $v_{A'}=(A',J^{(i+1)}(\tilde U))$ the real point of the chart of level $i+1$ that $A'$ gives, with $\|v_{A'}\|\le r_{i+1}$, $\nu=\nu_{i,A'}:=\nu_{i,v_{A'}}$ the measure of the step, and $B\mapsto A^{\mathrm{lw}}(B)=A^{\mathrm{lw}}(v_{A'},g_iB)$ the lowering map (Proposition~\YMMref{124.2}{834}). Then for every nonempty $\beta\subset[n]$
\[
G_{i+1;\beta}(A')=\sum_{\pi\in\mathrm{Part}(\beta)}\kappa_\nu\bigl(G_{i;\beta'}\circ A^{\mathrm{lw}}:\beta'\in\pi\bigr).
\]
The same holds for real mixed sources $\mathfrak s_a\ne0$, $a\in[n]$, at $0<g\le\gamma_{\mathrm N}$ with $K\ge2$, the progression being formed with their weight members in the source slot (the lemma on bond sources, which follows Lemma~\ref{C-352-36}), when the two unions of the hulls of the clusters $I_1$, $I_2$ satisfy the conditions of the frame at $k$ of clause~(6) of that lemma. For mixed sources the frame is that of clause~(6), read with the hulls in the place of the supports, and the tower identity for a pair is the one stated in that clause.
\end{lemma}

\begin{proof}
For $\lambda$ near $0$ in $\mathbb C^n$ put $\Phi(\lambda)=\exp D^\lambda_{i+1}(\tilde U)-\int\exp D^\lambda_i(e^{i\eta_iA^{\mathrm{lw}}(B)})\,d\nu(B)$, formed with the progression of $(w_1,\dots,w_n)$. Each relative exponent is a finite sum of functions that are holomorphic in $\lambda$, and the integrand is bounded uniformly in $B$ for $\lambda$ near $0$, as the progression with any number of sources lies in the setting of one step (Theorem~\YMMref{D}{8}); so $\Phi$ is holomorphic near $0$. Let $t\in\mathbb R^n$ be nonzero on $I_1$ and on $I_2$, and put $f^t=\sum_{a\in I_1}t_af_a$, $h^t=\sum_{a\in I_2}t_af_a$. These are real Lipschitz functions; $f^t\circ x$ and $h^t\circ x$ do not vanish identically, the supports of the $w_a$ being disjoint; and their supports lie in the two unions of supports.

So $(f^t,h^t)$ lies in the frame at $k$, and the tower identity among the construction results for a pair (equation~\eqref{eq:weighted-pair-tower}) gives $\exp D^\mu_{i+1}[f^t,h^t](\tilde U)=\int\exp D^\mu_i[f^t,h^t](e^{i\eta_iA^{\mathrm{lw}}})\,d\nu$ for $\mu$ near $0$ in $\mathbb C^2$, with the same measure $\nu$ and the same lowering map, neither of which depends on the sources. The progression depends on the parameters and the sources only through $w_\lambda$; at $\lambda(\mu)=(\mu_1t|_{I_1},\mu_2t|_{I_2})$ one has $w_{\lambda(\mu)}=\mu_1f^t\circ x+\mu_2h^t\circ x$, so $D^{\lambda(\mu)}_\cdot=D^\mu_\cdot[f^t,h^t]$, and at $\mu=(s,s)$: $\Phi(st)=0$ for all small real $s$. With the expansion $\Phi=\sum_N\Phi_N$ in homogeneous polynomials, $\sum_Ns^N\Phi_N(t)=0$ for small $s$; so $\Phi_N(t)=0$ for every $N$ and every $t$ of a dense subset of $\mathbb R^n$, hence $\Phi_N=0$ and $\Phi=0$ near $0$.

As $D^0_{i+1}=0$ and $D^0_i=0$, the multilinear jets of the two members of $\exp D^\lambda_{i+1}(\tilde U)=\int\exp D^\lambda_i\,d\nu$ agree in $\mathbb C[\lambda_1,\dots,\lambda_n]/(\lambda_a^2)$, and the logarithm of a moment series is the cumulant series; differentiation under the integral is allowed by Cauchy's formula in $\lambda$. The coefficient of $\lambda^\beta$ is the claim. At $n=2$ the lemma is the identity among the construction results. For mixed sources the argument is the same, with the weight member of $\mathfrak s_\lambda=\sum_a\lambda_a\mathfrak s_a$, which is linear in $\mathfrak s_\lambda$ (clause~(1) of the lemma on bond sources), in the source slot. The $t\in\mathbb R^n$ at which $\sum_{a\in I_1}t_a\mathfrak s_a$ or $\sum_{a\in I_2}t_a\mathfrak s_a$ is $0$ form two proper linear subspaces, every $\mathfrak s_a$ being nonzero; off them, a dense set, the two sums are nonzero mixed sources whose hulls lie in the two unions of the hulls of the clusters, so they are in the frame of clause~(6), and the identity for a pair stated there, with the same measure $\nu$ and the same lowering map, stands in the place of equation~\eqref{eq:weighted-pair-tower}.
\end{proof}

\begin{proposition}[The induction over the levels]\label{C-352-21}
Let $n\in\{2,3,4\}$ and let $g$ satisfy $(\Gamma1)$, with $c^{(4)}_\epsilon$ in place of $c_\epsilon$ if $n=4$; let $(K,\beta)$ be a first-passage datum at $g$, $\ell\in\mathfrak L(g)$, and let real weights $w_1,\dots,w_n$ with the two clusters $\{1\},\{2\}$, or $\{1,2\},\{3\}$, or $\{1,2\},\{3,4\}$ lie in the frame at a level $k$; let $k\le k_1<K$, $0\le k_\flat\le k_1$, and form the free objects over the base level $k_\flat$. Let numbers $\mathsf v_\beta>0$, $c^\circ_\beta$ and $\epsilon_\circ\ge0$, $\emptyset\ne\beta\subset[n]$, satisfy
\begin{enumerate}
\item[(I1)] $|G_{k_1;\beta}(A)-c^\circ_\beta-q_{k_1;\beta}(\mathsf{ang}^{[k_1]}(A))|\le\epsilon_\circ\mathsf v_\beta$ for $A\in\mathsf{Cc}_{k_1}(r_{k_1})$;
\item[(I2)] $\varrho_{i;\beta}:=\eta_i^4\sum_{p,q}|R_{i;\beta}(p,q)|/(a_Kx_i)\le L^{-3(i-k_1)}\mathsf v_\beta$ for $k_1\le i<K$;
\item[(I3)] $\prod_{\beta'\in\pi}\mathsf v_{\beta'}\le\mathsf v_\beta$ for every partition $\pi$ of $\beta$.
\end{enumerate}
Put $Z=Z_0(1+\mathfrak b_{k_1})^2$ and $\varpi=\epsilon_\circ+\sum_{i=k_1}^{K-1}L^{-3(i-k_1)}(\delta_i+c_{2n}\epsilon'^{(n)}_i)$, with a number $Z_0$ formed from the constants and the numbers $c_{2n}$ of Lemma~\ref{C-352-20}(b), and let $(K-k_1)\varpi\le1$. Then $\varpi\le\epsilon_\circ+\frac65(1+c_{2n}c^{(n)}_\epsilon)\delta_{k_1}$, and for $k_1\le i\le K$, $A\in\mathsf{Cc}_i(r_i)$ and every $\beta$, $m=|\beta|$,
\[
\bigl|G_{i;\beta}(A)-c^\circ_\beta-(g_{i;\beta}-g_{k_1;\beta})-q_{i;\beta}(\mathsf{ang}^{[i]}(A))\bigr|\le\mathsf a_mZ^m\varpi\,\mathsf v_\beta,\qquad(\mathsf a_1,\mathsf a_2,\mathsf a_3)=(1,14,285),
\]
$\mathsf a_4$ a number; in particular $|\mathsf J_n-c^\circ_{[n]}-g_{K;[n]}+g_{k_1;[n]}|\le\mathsf a_nZ^n\varpi\mathsf v_{[n]}$. Of the base level only the relation $\mathcal S_{i+1}-\mathcal S_i=\upsilon_{i+1}\mathcal C_{(i+1)}$, $i\ge k_1$, is used.
\end{proposition}

\begin{proof}
Let $\mathsf e_{i;\beta}$ be the supremum over $\mathsf{Cc}_i(r_i)$ of the left member. 1. \emph{The tower.} For $A'\in\mathsf{Cc}_{i+1}(r_{i+1})$ and the measure $\nu=\nu_{i,A'}$ of the step, $G_{i+1;\beta}(A')=\sum_{\pi}\kappa_\nu(G_{i;\beta'}\circ A^{\mathrm{lw}}:\beta'\in\pi)$ (the lemma on the tower identity at several sources, which precedes this proposition, for the two clusters of the frame). 2. \emph{The lowered angle.} The point $v_{A'}$ of the chart is real, with $\|v_{A'}\|\le r_{i+1}<\varrho_{\mathrm{ch}}\le r_\alpha$ by the definition of the chain classes (equation~\eqref{eq:weighted-chain-class}); $\nu=\nu_{i,v_{A'}}$ is a probability measure carried by the real $B$ with $|B(b)|<\mathfrak b_i$ (Proposition~\YMMref{124.5}{836}), and $A^{\mathrm{lw}}(B)\in\mathsf{Cc}_i(r_i)$ (Proposition~\YMMref{124.2}{834} with the radius recurrence~\YMMref{(2291)}{834}).

By the splitting of Lemma~\ref{C-352-20}(a) at $r=r_{i+1}$, $\mathsf{ang}^{[i]}_p(A^{\mathrm{lw}}(B))=F'_p+\psi_p(B)+\zeta_p(B)$, with $F'=\mathsf{ang}^{[i+1]}(A')$, $\psi_p(B)=\mathsf{sg}_ig_i\sum_b\mathsf T^{\mathrm{cv}}_i(p,b)B(b)$ linear in the fluctuation field $B$, $|\psi_p|\le2\mathsf c_{\mathfrak h}\eta_i^2\delta_i$ and $|\zeta_p|\le C\eta_i^2\delta_i^2$ by the smallness of Lemma~\ref{C-352-10}(b); besides $|F'_p|\le2\eta_{i+1}^2r_{i+1}\le\mathsf c_{\mathfrak h}\eta_i^2\delta_i$. So $Y_\beta=G_{i;\beta}\circ A^{\mathrm{lw}}$ is within $\mathsf e_{i;\beta}+Z\delta_iL^{-3(i-k_1)}\mathsf v_\beta$ of the quadratic polynomial $P_\beta=c^\circ_\beta+g_{i;\beta}-g_{k_1;\beta}+q_{i;\beta}(F'+\psi)$ in the coordinates $B^\alpha(b)$: the term with $\zeta$ is at most $Ca_K^{-1}\eta_i^4\delta_i^3\sum_{p,q}|R_{i;\beta}(p,q)|=C\mathfrak b_i^2\delta_i\varrho_{i;\beta}$.

The linear part of $P_\beta$ has the coefficients $\mathsf{sg}_ig_ia_K^{-1}\sum_{p,q}R_{i;\beta}(p,q)\mathsf T^{\mathrm{cv}}_i(q,b)F'^\alpha_p$, its quadratic part the coefficients $\frac12g_i^2a_K^{-1}\sum_{p,q}R_{i;\beta}(p,q)\mathsf T^{\mathrm{cv}}_i(p,b)\mathsf T^{\mathrm{cv}}_i(q,b')$ of $B^\alpha(b)B^\alpha(b')$; with the row sums $\sum_b|\mathsf T^{\mathrm{cv}}_i(p,b)|\le C^{\mathrm{row}}_{\mathrm{tp}}\eta_i^2$, $x_ig_i\delta_i=\mathfrak b_i$ and $x_ig_i^2=1$,
\[
\|P_\beta\|_{\mathrm{pol}}\le\bigl[2\mathsf c_{\mathfrak h}C^{\mathrm{row}}_{\mathrm{tp}}\mathfrak b_i+\tfrac32(C^{\mathrm{row}}_{\mathrm{tp}})^2\bigr]\varrho_{i;\beta};
\]
so, by (I2), $\mathfrak b_i\le\mathfrak b_{k_1}$ and $|B(b)|<\mathfrak b_i$, $\|P_\beta\|_{\mathrm{pol}}$ and $\sup|P_\beta-\text{const}|$ are $\le ZL^{-3(i-k_1)}\mathsf v_\beta$, $Z_0$ being at least $2\mathsf c_{\mathfrak h}C^{\mathrm{row}}_{\mathrm{tp}}+\frac32(C^{\mathrm{row}}_{\mathrm{tp}})^2$ and the constant of the term with $\zeta$. 3. \emph{The Gaussian value.} Under the Gaussian with covariance $\mathbf C_0$, $\psi=\mathsf{sg}_ig_i\mathsf T^{\mathrm{cv}}_iB$ has the covariance $g_i^2\eta_i\mathcal C_{(i+1)}=a_K\upsilon_{i+1}\mathcal C_{(i+1)}$ (Lemma~\ref{C-352-20}(a): $\mathsf{sg}_i^2=1$, $\mathsf T^{\mathrm{cv}}_i(\mathsf A^{\mathrm{cv}}_i)^{-1}(\mathsf T^{\mathrm{cv}}_i)^{\mathsf T}=\eta_i\mathcal C_{(i+1)}$, $g_i^2\eta_i=a_K\upsilon_{i+1}$), and Lemma~\ref{C-352-20}(c) gives $\sum_\pi\kappa_{\mathsf G}(P_{\beta'})=c^\circ_\beta+g_{i+1;\beta}-g_{k_1;\beta}+q_{i+1;\beta}(F')$. 4.

\emph{From $\mathsf G$ to $\nu$} by Lemma~\ref{C-352-20}(a), (b), with $\mathsf c=\max\{1,\mathsf c_{\mathbf C}\}$: as $\|v_{A'}\|\le r_{i+1}$, all the cumulants of the orders $\le2n$ lie within $\epsilon=\epsilon'^{(n)}_i$ of those of $\mathsf G$, and $\epsilon'^{(n)}_i\le c^{(n)}_\epsilon\delta_i\le1$ (under $(\Gamma1)$, as $c^{(2)}_\epsilon\le c_\epsilon$, and with $c^{(4)}_\epsilon$ at $n=4$); the number is $c_{2n}$; \emph{from $P$ to $Y$} by $|\kappa(X,Y)|\le\|X\|^\sim\|Y\|^\sim$, $|\kappa(X,Y,W)|\le2\|X\|^\sim\|Y\|^\sim\|W\|^\sim$, $\|X\|^\sim=\inf_c\sup|X-c|$. 5. \emph{The recursion.} With (I3) every product of sizes over a partition is at most $Z^m\mathsf v_\beta$ times pure numbers, and $\varepsilon_{i;\beta}=\mathsf e_{i;\beta}/Z^m\mathsf v_\beta$ obeys $\varepsilon_{i+1;a}\le\varepsilon_{i;a}+\tau_i$, $\varepsilon_{i+1;ab}\le\varepsilon_{i;ab}+2\tau_i+4\varpi^2+4\theta_i\varpi$ and the like at three sources, $\tau_i=L^{-3(i-k_1)}(\delta_i+c_{2n}\epsilon'^{(n)}_i)$, $\theta_i=L^{-3(i-k_1)}$; summed over at most $K-k_1$ heights with $(K-k_1)\varpi\le1$ this gives $1$, $14$ and $285$.

At $i=K$, $\mathsf{Cc}_K=\{0\}$ and $q_{K;\beta}(0)=0$. The bound of $\varpi$: $\epsilon'^{(n)}_i\le c^{(n)}_\epsilon\delta_i$ (as $\epsilon'_i\le c_\epsilon\delta_i$ in Lemma~\ref{C-352-10}(b)), $\delta_i\le(\frac54)^{1/2}L^{(i-k_1)/2}\delta_{k_1}$, $(\frac54)^{1/2}\sum_kL^{-5k/2}\le\frac65$. For two sources Steps 1--5 use only \eqref{eq:weighted-pair-tower-derivatives} and the cumulants up to order four.
\end{proof}

\begin{lemma}[The slices]\label{C-352-22}
Let $(K,\beta)$ be a first-passage datum at $g\le\gamma$, $p\ne q$ plaquettes, $\mathsf r=|x(p)-x(q)|_\infty$, $c^{(N)}_L=\sup_\tau\sum_{j\in\mathbb Z}(L^j\tau)^Ne^{-\delta_{\mathrm{tp}}L^j\tau}$ and $\mathsf m_K=\frac54C_{\mathrm{tp}}c^{(3)}_L$. (a) For every set $\mathcal J\subset\{1,\dots,K\}$: $|\sum_{j\in\mathcal J}\upsilon_j\mathcal C_{(j)}(p,q)|\le\mathsf m_K\mathsf b^{-1}a_K^3\mathsf r^{-3}$. (b) If $t=\mathsf r/a_{K-k}$ satisfies $\delta_{\mathrm{tp}}L(L-1)t\ge3\log L+\log2$, then
\[
\Bigl|\sum_{k<j\le K}\upsilon_j\mathcal C_{(j)}(p,q)-\mathsf b^{-1}C^{\mathrm{free}}_K(p,q)\Bigr|\le\mathsf b^{-1}a_K^3\bigl[\tfrac54C_{\mathrm{tp}}c^{(2)}_L\bar B\,\mathsf r^{-2}\mathsf b^{-1}+C^{\mathrm{bg}}_{\mathrm{tp}}+2L^3C_{\mathrm{tp}}\mathsf r^{-3}t^3e^{-\delta_{\mathrm{tp}}Lt}\bigr],
\]
with $C^{\mathrm{bg}}_{\mathrm{tp}}$ from Proposition~\YMMref{120.5}{821}, evaluated at $s=0$ in equation~\YMMref{(2236)}{821}.
\end{lemma}

\begin{proof}
(a) $\upsilon_j=(a_{K-j+1}x_{j-1})^{-1}\le\frac{24}{23}\mathsf b^{-1}\le\frac54\mathsf b^{-1}$ at every $0<g\le\gamma$: the flow bound and $g^{-2}<\mathsf b$, which Lemma~\ref{C-352-10}(a) states for every first-passage datum at every $0<g\le\gamma$, give $|a_{K-j+1}x_{j-1}-\mathsf b|\le a_{K-j+1}\bar B\le\bar B/L$, and $\bar B\le\frac{13}{24}g^{-2}<\frac{13}{24}\mathsf b$ at $g\le\gamma$ (Lemma~\ref{lem:setup-stopped-arithmetic} gives this bound), so $a_{K-j+1}x_{j-1}\ge\frac{23}{24}\mathsf b$ as $L\ge13$; and the decay of the slices: $\sum_j(\mathsf r/a_{K-j+1})^3e^{-\delta_{\mathrm{tp}}\mathsf r/a_{K-j+1}}\le c^{(3)}_L$, a lacunary sum. (b) Equation~\YMMref{(2227)}{818} gives $\sum_{k<j\le K}\mathcal C_{(j)}=C^{\mathrm{free}}_K-C^{\mathrm{bg}}_{K,0}-\sum_{j\le k}\mathcal C_{(j)}$, and Proposition~\YMMref{120.5}{821} bounds $|C^{\mathrm{bg}}_{K,0}(p,q)|$ by $C^{\mathrm{bg}}_{\mathrm{tp}}a_K^3$; $|\upsilon_j-\mathsf b^{-1}|=|\mathsf b-a_{K-j+1}x_{j-1}|\upsilon_j\mathsf b^{-1}\le\frac54\bar Ba_{K-j+1}\mathsf b^{-2}$ by the same bounds of the run, which with the decay gives the first term; the slices of the levels $j\le k$ are at $t_j\ge Lt$ and give the last term.
\end{proof}

\begin{lemma}[The explicit sizes; the weights of the induction]\label{C-352-23}
Let weights be placed at $(\rho,r,\mathsf d)$, or form two pairs at $(\rho,r,\mathsf d)$, with $\mathsf d\le e^{-2\sqrt3}$, $s(\rho)\le K$; let $k_\flat=\mathsf k(r)$, $k_1=\mathsf k(\mathsf d)$, $u=\hat r_{k_1}$, $\hat\rho=\hat\rho_{k_1}=qu$, $P=\hat{\mathsf d}_{k_1}$, and $I_1,I_2$ the clusters. Put $\mathsf w(\emptyset)=1$, $\mathsf w(\{a\})=\hat\rho^3$, $\mathsf w(\text{a whole pair})=q^6u^3\max\{1,u^3\}$, $c_*=\max\{102,4\cdot81\cdot3^{3/2}(1+\mathsf m_K)\}$, and
\[
\mathsf v_\beta=c_*^{|\beta|}\,\mathsf w(\beta\cap I_1)\,\mathsf w(\beta\cap I_2)\,\mathsf u_\beta,\qquad\mathsf u_\beta=\prod_{a\in\beta}\mathsf u_a .
\]
Then (I2) and (I3) of Proposition~\ref{C-352-21} hold; $\mathsf v_{[3]}\le c_*^3P^9\mathrm{nat}_3$ and $\mathsf v_{[4]}\le c_*^4P^{12}\mathrm{nat}_4$; for crossing $\beta$, $\mathsf z(\beta)\mathsf u_\beta\le P^{-6}\mathsf v_\beta$, and if $\mathsf e_{\mathrm{far}}=8e^{-\delta_{\mathrm{tp}}LP/2\sqrt3}\le1$, then $\varrho_{k_1;\beta}\le\mathsf e_{\mathrm{far}}\mathsf v_\beta$ and $|g_{k_1;[n]}|\le\frac32(n-1)!(81\cdot3^{3/2}\mathsf m_K)^n\mathsf e_{\mathrm{far}}^2\mathrm{nat}_n$. The factor $\max\{1,u^3\}$ cannot be dropped: $u$ may be a power of $\log\frac1{\mathsf d}$.
\end{lemma}

\begin{proof}
By Lemma~\ref{C-352-22}(a), $|\mathcal S_i(p,q)|\le\mathsf m_K\mathsf b^{-1}a_K^3\mathsf r^{-3}$ with $\mathsf r\ge r/\sqrt3$ in a cluster and $\ge\mathsf d/\sqrt3$ across. An ordering of $\beta$ is a path with $m-1$ edges, at most $c(\beta)$ of them inside a whole pair; with (P1),
\[
\varrho_{i;\beta}\le\tfrac54m!\,81^m(3^{3/2}\mathsf m_K)^{m-1}\rho^{3m}r^{-3c}\mathsf d^{-3(m-1-c)}a_{K-i}^{-3}\mathsf u_\beta,
\]
and with $a=a_{K-k_1}$ the powers of $a$ cancel: $\rho^{3m}r^{-3c}\mathsf d^{-3(m-1-c)}a^{-3}=\hat\rho^{3m}u^{-3c}P^{-3(m-1-c)}\le\mathsf w(\beta\cap I_1)\mathsf w(\beta\cap I_2)$, $a_{K-i}^{-3}=L^{-3(i-k_1)}a^{-3}$: (I2). (I3): either $\beta\cap I_\iota$ lies in one block, or it is a whole pair cut into two points, and $\hat\rho^6=q^6u^6\le q^6u^3\max\{1,u^3\}$. The top sets: $q^6u^3\hat\rho^3=P^6\rho^9r^{-3}\mathsf d^{-6}$. At the level $k_1$ the slices of the levels $\le k_1$ at distances $\ge\mathsf d/\sqrt3$ are exponentially small, $|\mathcal S_{k_1}(p,q)|\le\mathsf e_{\mathrm{far}}\mathsf m_K\mathsf b^{-1}a_K^3\mathsf r^{-3}$; a path through a crossing set has a far edge, a ring through all sources two.
\end{proof}

\begin{proposition}[The jets of narrow bumps]\label{C-352-24}
There are numbers $C(g)$ and $\mathsf d_4(g),\mathsf d^{(4)}_4(g)\in(0,\min\{e^{-2(15\mathfrak p_0+2)},r_1(\bar{\mathsf b})\}]$, formed from the constants and $g$, such that the following holds for $g$ under $(\Gamma1)$, every first-passage datum $(K,\beta)$ at $g$, every side $\ell\in\mathfrak L(g)$ with $64\mathsf d\le\ell$, and $s(\rho)\le K$.
\begin{enumerate}
\item[(a)] For real weights placed at $(\rho,r,\mathsf d)$ with $0<\rho\le r\le\mathsf d\le\mathsf d_4(g)$: $|\mathsf J_3-3\operatorname{Tr}\prod_a(w_a\mathsf b^{-1}C^{\mathrm{free}}_K)|\le\epsilon^{\mathrm{nar}}(g,\mathsf d)\,\mathrm{nat}_3$.
\item[(b)] For real weights $w_1,w_2$ with $\Sigma(w_a)\subset\mathbb B(x_a,\rho)$, $\operatorname{dist}(\Sigma(w_1),\Sigma(w_2))\ge\mathsf d$, $|x_1-x_2|\le18\mathsf d$, $0<\rho\le\mathsf d\le\mathsf d_4(g)$: $|\mathsf J_2-\frac32\mathsf b^{-2}\sum_{p,q}w_1(p)w_2(q)C^{\mathrm{free}}_K(p,q)^2|\le\epsilon^{\mathrm{nar}}_2(g,\mathsf d)(\rho/\mathsf d)^6\mathsf u_1\mathsf u_2$.
\item[(c)] Under $(\Gamma1)$ with $c^{(4)}_\epsilon$ in place of $c_\epsilon$, for two pairs at $(\rho,r,\mathsf d)$ with $\mathsf d\le\mathsf d^{(4)}_4(g)$: $|\mathsf J_4-\frac32\sum_\varsigma\operatorname{Tr}\prod_l(w_{\varsigma(l)}\mathsf b^{-1}C^{\mathrm{free}}_K)|\le\epsilon^{\mathrm{nar}}_4(g,\mathsf d)\,\mathrm{nat}_4$.
\end{enumerate}
The normalized error factors and constants are uniform in $\rho/r$, $r/\mathsf d$, $K$, $\ell$, position and direction; in particular they hold for $w_a=f_a\circ x$ with bounded Borel $f_a$, and for plane weights.
\end{proposition}

\begin{proof}
1. \emph{Levels and data.} $k_1=\mathsf k(\mathsf d)$, $k_\flat=\mathsf k(r)\le k_1$, $P=\hat{\mathsf d}_{k_1}$; by the working level $4MY_{\mathsf d}^{2\mathfrak p_0}\le P<CY_{\mathsf d}^{2\mathfrak p_0}$, $K-k_1\le Y_{\mathsf d}/\log L$, $\delta_{k_1}\le A_1Y_{\mathsf d}^{\mathfrak p_0}(\mathsf d/\mathsf b)^{1/2}$. Lemma~\ref{C-352-11}(b) puts the two clusters in the frame at $k_1$; the smallness holds at all heights (Lemma~\ref{C-352-10}); (I2), (I3) hold with the weights of Lemma~\ref{C-352-23}.

2. \emph{(I1).} At a single source, $c^\circ_a=G_{k_1;a}(0)$ and Lemma~\ref{C-352-16} give the tolerance $C\mathfrak b_{k_1}^2\delta_{k_1}\mathsf v_a$. At the whole pair, $c^\circ=G_{k_1;12}(0)$: by the multilinearity $G_{k_1;12}$ is the pair functional of the pair alone, and (B.2) at $n=k_1\ge\mathsf k(r)$ (its proof, given below, rests on Lemmas~\ref{C-352-10}--\ref{C-352-19} alone) gives $[C_{\mathrm N}\mathfrak b_{k_1}^2\delta_{k_1}+\epsilon'_{\mathrm N}(r)]q^6u^3\mathsf u_{12}\le[\dots]\mathsf v_{12}$. At a crossing set, $c^\circ_\beta=0$: Lemma~\ref{C-352-15}(e) and Lemma~\ref{C-352-23} give $|G_{k_1;\beta}|\le Ce^{4Y_{\mathsf d}-\frac2{\sqrt3}Y_{\mathsf d}^{2\mathfrak p_0}}\mathsf v_\beta$, and $|q_{k_1;\beta}|\le C\mathfrak b_{k_1}^2\mathsf e_{\mathrm{far}}\mathsf v_\beta$. So
\[
\epsilon_\circ=C_\circ\bigl\{\mathfrak b_{k_1}^2\delta_{k_1}+\epsilon'_{\mathrm N}(r)+e^{4Y_{\mathsf d}-\frac2{\sqrt3}Y_{\mathsf d}^{2\mathfrak p_0}}+\mathfrak b_{k_1}^2\mathsf e_{\mathrm{far}}\bigr\}.
\]
Here $C_\circ$ is a number formed from the constants that dominates the constants of these four bounds (of Lemma~\ref{C-352-16}, of (B.2), of Lemma~\ref{C-352-15}(e) with Lemma~\ref{C-352-23}, and of the bound of $q_{k_1;\beta}$), and those of the two bounds of (I1) in step 6. The tolerance of (B.2) falls together with its yardstick; for this reason no relation between $r$ and $\mathsf d$ is needed, and the bounds hold at every ratio $r/\mathsf d\in(0,1]$.

3. \emph{The induction} (Proposition~\ref{C-352-21}) at $i=K$: $|\mathsf J_n-g_{K;[n]}+g_{k_1;[n]}|\le\mathsf a_nZ^n\varpi c_*^nP^{3n}\mathrm{nat}_n$.

4. \emph{The rings.} Lemma~\ref{C-352-22}(b) at $k=\mathsf k(r)$ for $p,q$ in different supports ($t\ge\hat r_{\mathsf k(r)}/\sqrt3$) gives $|\mathcal S_K-\mathsf b^{-1}C^{\mathrm{free}}_K|\le\mathrm{br}\,\mathsf b^{-1}a_K^3\mathsf r^{-3}$ with
\[
\mathrm{br}=25C_{\mathrm{tp}}c^{(2)}_L\bar B\,\mathsf d/\mathsf b+8000C^{\mathrm{bg}}_{\mathrm{tp}}\mathsf d^3+2L^3C_{\mathrm{tp}}\sup_{t\ge\frac4{\sqrt3}MY_r^{2\mathfrak p_0}}t^3e^{-\delta_{\mathrm{tp}}Lt};
\]
telescoping a product of $n$ kernels, $|g_{K;[n]}-\frac32\sum_\varsigma\operatorname{Tr}\prod_l(w_{\varsigma(l)}\mathsf b^{-1}C^{\mathrm{free}}_K)|\le C\,\mathrm{br}\,\mathrm{nat}_n$; at $n=3$ the subtrahend is $3\operatorname{Tr}\prod_a(w_a\mathsf b^{-1}C^{\mathrm{free}}_K)$, two cyclic orders with one trace. With Lemma~\ref{C-352-23} for $g_{k_1;[n]}$:
\[
\epsilon^{\mathrm{nar}}_{(n)}=\mathsf a_nZ^nc_*^nP^{3n}\varpi+C\,[\mathsf e_{\mathrm{far}}^2+\mathrm{br}].
\]

5. \emph{The threshold $\mathsf d_4(g)$.} It is fixed by six conditions: (i) $\mathsf d\le e^{-2(15\mathfrak p_0+2)}$, so that by Lemma~\ref{C-352-10}(c) $\epsilon'_{\mathrm N}(r)\le C(g)(\log\frac1{\mathsf d})^{15\mathfrak p_0+2}\mathsf d^{1/2}$ for all $r\le\mathsf d$; (ii) $\frac4{\sqrt3}M\delta_{\mathrm{tp}}L(L-1)(\log\frac1{\mathsf d})^2\ge3\log L+\log2$; (iii) $\log\frac1{\mathsf d}\ge\max\{5,\sqrt3/\delta_{\mathrm{tp}}LM\}$, so that $e^{4Y_{\mathsf d}-\frac2{\sqrt3}Y_{\mathsf d}^2}\le\mathsf d$, $\mathsf e_{\mathrm{far}}\le8\mathsf d^2$ and the last member of $\mathrm{br}$ is $\le C\mathsf d$; (iv) $8\mathsf d^2\le1$, $\overline{\mathrm{br}}\le1$; (v) $(Y_{\mathsf d}/\log L)[\bar\epsilon_\circ+\frac65(1+c_6c_\epsilon)A_1Y_{\mathsf d}^{\mathfrak p_0}g\mathsf d^{1/2}]\le1$; (vi) $\mathsf d\le r_1(\bar{\mathsf b})$, so that $\rho\le r_1(\mathsf b)$ holds of itself. In (iv) and (v) the numbers $\mathrm{br}$ and $\epsilon_\circ$, which read $r$ and the datum, stand replaced by majorants that read $g$ and $\mathsf d$ alone. As $r\le\mathsf d$, $s(r)\ge s(\mathsf d)$ and $Y_r\ge Y_{\mathsf d}$; with $\mathsf d/\mathsf b<g^2\mathsf d$,
\[
\mathrm{br}\le\overline{\mathrm{br}}=25C_{\mathrm{tp}}c^{(2)}_L\bar B\,g^2\mathsf d+8000C^{\mathrm{bg}}_{\mathrm{tp}}\mathsf d^3+2L^3C_{\mathrm{tp}}\sup_{t\ge\frac4{\sqrt3}MY_{\mathsf d}^{2\mathfrak p_0}}t^3e^{-\delta_{\mathrm{tp}}Lt}.
\]
In $\epsilon_\circ$, $\epsilon'_{\mathrm N}(r)$ is at most the bound of (i); $\mathfrak b_{k_1}\le A_1Y_{\mathsf d}^{\mathfrak p_0}$ and $\delta_{k_1}\le A_1Y_{\mathsf d}^{\mathfrak p_0}g\mathsf d^{1/2}$ (step 1, $\mathsf b>g^{-2}$); and $\mathsf e_{\mathrm{far}}\le8\mathsf d^2$ by (iii), as $P\ge4MY_{\mathsf d}^{2\mathfrak p_0}$. So
\[
\epsilon_\circ\le\bar\epsilon_\circ=C_\circ\bigl\{A_1^3Y_{\mathsf d}^{3\mathfrak p_0}g\mathsf d^{1/2}+C(g)(\log\tfrac1{\mathsf d})^{15\mathfrak p_0+2}\mathsf d^{1/2}+e^{4Y_{\mathsf d}-\frac2{\sqrt3}Y_{\mathsf d}^{2\mathfrak p_0}}+8A_1^2Y_{\mathsf d}^{2\mathfrak p_0}\mathsf d^2\bigr\}.
\]
$Y_{\mathsf d}$ is a function of $g$ and $\mathsf d$, with $\log\frac1{\mathsf d}\le Y_{\mathsf d}\le2\log\frac1{\mathsf d}+C(g)$ (Lemma~\ref{C-352-10}(c)). So (i)--(vi) are conditions on $(g,\mathsf d)$ alone, and they hold for all small $\mathsf d$, every majorised member tending to $0$ with $\mathsf d$. $\mathsf d_4(g)$ is half the supremum of the numbers $\mathsf d_*$ such that (i)--(vi) hold at every $0<\mathsf d\le\mathsf d_*$: a positive number formed from the constants and $g$, which reads no $K$, $\beta$, $\ell$, $\rho$, $r$; at every $\mathsf d\le\mathsf d_4(g)$, every $r\le\mathsf d$, every datum at $g$ and every label $\mathsf b\in(g^{-2},\bar{\mathsf b}]$ the conditions hold with $\mathrm{br}$ and $\epsilon_\circ$ themselves.

As $\varpi\le\epsilon_\circ+\frac65(1+c_6c_\epsilon)\delta_{k_1}$ at $n=3$ (Proposition~\ref{C-352-21}) and $K-k_1\le Y_{\mathsf d}/\log L$, (v) secures $(K-k_1)\varpi\le1$; at $n=2$ (step 6) the same (v) secures it, as $c_4c^{(2)}_\epsilon\le c_6c_\epsilon$ (Lemma~\ref{C-352-20}) and $C_\circ$ dominates the constants of step 6. $\mathsf d^{(4)}_4(g)$ is formed in the same way with $n=4$, that is with $c_8c^{(4)}_\epsilon$ in (v) and the tolerance $\epsilon_\circ$ of two pairs, under $(\Gamma1)$ with $c^{(4)}_\epsilon$. Then $Z^n\le CY_{\mathsf d}^{2n\mathfrak p_0}$, $P^{3n}\le CY_{\mathsf d}^{6n\mathfrak p_0}$, $\mathfrak b^2_{k_1}\delta_{k_1}\le CY_{\mathsf d}^{3\mathfrak p_0}g\mathsf d^{1/2}$, and the exponents are $8n\mathfrak p_0+15\mathfrak p_0+2$: $39\mathfrak p_0+2$ at $n=3$, $47\mathfrak p_0+2$ at $n=4$.

6. \emph{Two sources (b).} The same with the base level equal to the start level $k_1=\mathsf k(\mathsf d)$, $\mathsf v_a=c_*\hat\rho^3\mathsf u_a$, $\mathsf v_{12}=c_*^2\hat\rho^6\mathsf u_{12}$; at the pair $\mathcal S_{k_1}=0$, so $q_{k_1;12}=g_{k_1;12}=0$, and Lemma~\ref{C-352-15}(d) at $n=k_1$ gives $|G_{k_1;12}|\le Ce^{2Y_{\mathsf d}-\frac2{\sqrt3}Y_{\mathsf d}^{2\mathfrak p_0}}\mathsf v_{12}$; then $g_{K;12}=\frac32\sum w_1(p)w_2(q)\mathcal S_K(p,q)^2$ and Lemma~\ref{C-352-22}(b) at $k=k_1$; the logarithms count $4\mathfrak p_0+12\mathfrak p_0+3\mathfrak p_0=19\mathfrak p_0$.
\end{proof}

\paragraph{The connected functions on the least torus.}

\begin{lemma}[The mixture; the conditioning stopped at any level]\label{C-352-25}
Let $0<g\le\gamma$ with $g^{-2}\ge x_{\mathrm J}$, $(K,\beta)$ a first-passage datum at $g$ with $K\ge2$, $\ell\in\mathfrak L(g)$, and $w_1,\dots,w_p$ real weights, $p\ge1$ arbitrary. For $0\le k\le K$ let $\Omega_k=\mathcal V_k\sqcup\bigsqcup_{t=1}^k\mathcal V_{t-1}$ and $\mathbb P_k=Z_k^{-1}[\mathsf m_k\oplus\bigoplus_{t\le k}\mathsf g^{\mathrm w}_t\mathsf m_{t-1}]$, where $\mathcal V_t$ is the space of the block fields of the level $t$, $\mathsf m_t$ the measure of the never-large history up to the level $t$, and $\mathsf g^{\mathrm w}_t\ge0$ the weight of a first large field at the step $t$ in \eqref{eq:ward-first-creation}; $t'=k$ on the first part and $t'=t-1$ on the part $t$; $\alpha=a_{K-t'}$; $U_0(W)=W$, and for $t'\ge1$ let $U_{t'}(W)$ be the background configuration of $W$ in the small-field region $\mathcal V^\chi_{t'}$. Write $\Omega=\Omega_K$, $\mathbb P=\mathbb P_K$, and $\mathsf Y_\beta(\omega)=\partial_\beta[D^\lambda_{t'}(U_{t'}(W))-D^\lambda_K(\mathbf 1)]|_{\lambda=0}$ on $\Omega$. Then:
\begin{enumerate}
\item[(a)] $\mathbb P_k$ is a Borel probability measure formed with no source; its parts with $t\le k$ are parts of $\mathbb P$; the part $t$ of $\mathbb P$ has the mass $\mathsf w_t$, and $\int\alpha^{-b}d\mathbb P=\mathbb P(\mathcal V_K)+\sum_tL^{b(K-t+1)}\mathsf w_t$ for every real $b$.
\item[(b)] $\langle e^{\mathcal O(w_\lambda)}\rangle_{K,\ell,\beta}=\int\exp[D^\lambda_{t'}(U_{t'}(W))]\,d\mathbb P_k$ for $\lambda$ near $0$; coefficientwise, $\langle\prod_{\iota\in\gamma}\mathcal O(w_\iota)\rangle_{K,\ell,\beta}=\int[\exp\mathbb W]_\gamma\,d\mathbb P$ in $\mathbb C[\lambda_\iota]/(\lambda_\iota^2)$, with $\mathbb W$ the multilinear jet of $\lambda\mapsto D^\lambda_{t'}(U_{t'}(W))$; each $[\mathbb W]_\gamma$ is a bounded Borel function that depends on the sources of $\gamma$ alone.
\item[(c)] $S^T_n-\mathsf J_n=\sum_{\pi\in\mathrm{Part}([n])}\kappa_{\mathbb P}(\mathsf Y_{\beta'}:\beta'\in\pi)$; and with $G_\beta(\omega)=G_{t';\beta}(U_{t'}(W))$ on $\Omega_k$, $S^T_n=\sum_\pi\kappa_{\mathbb P_k}(G_{\beta'}:\beta'\in\pi)$.
\item[(d)] Where $t'\ge1$, $U_{t'}(W)$ has its pairs of all levels $\le t'$ in the cores and is regular of the level $t'$ at the radius $2B_3\varepsilon_{t'}$, with $8MB_3\varepsilon_{t'}\le c_{\mathrm A}\alpha_0$.
\item[(e)] On the least torus, under $(\Gamma2)$: $\int\alpha^{-b}d\mathbb P\le1+2L^b\mathsf C_{\mathrm{ag}}(b+1,\frac12)$ for every $b\ge0$, so $\mathsf M'=\int\alpha^{-24}d\mathbb P\le\bar{\mathsf M}^{(24)}$; and $\int\alpha^{-b}d\mathbb P_k\le a_{K-k}^{-b}+\mathsf M'$, $\int_{\{\alpha<\varrho\}}\alpha^{-b}d\mathbb P_k\le\varrho^{24-b}\mathsf M'$ for $0\le b\le24$, $\varrho\le a_{K-k}$. On a side $\ell\in\mathfrak L(g)$: $\int\alpha^{-b}d\mathbb P\le1+4L^bM^{-3}\ell^3e^{\mathsf B(X_g)\ell^3}\mathsf C_{\mathrm{ag}}(b+1,\eth)$, finite and growing with $\ell$, with the function $\mathsf B$ and the letter $\eth$ and $X_g=L(g^{-2}+B_\sharp)$.
\end{enumerate}
\end{lemma}

\begin{proof}
Equation~\eqref{eq:ward-first-creation} writes the torus expectation of a bounded gauge invariant function $F$ as $Z^{-1}$ times the integral against $\mathsf m_k$ of its quotient at the level $k$ plus the integrals against $\mathsf g^{\mathrm w}_t\mathsf m_{t-1}$ of its quotients at the levels $t-1$, $t\le k$, at every level $k\le K$. At $F=1$ this makes $\mathbb P_k$ a probability measure; at $F=e^{\mathcal O(w_\lambda)}$ the quotient at a level $t$ is, on $\mathcal V^\chi_t$, $\exp D^\lambda_t(U_t(W))$, by Proposition~\ref{prop:running-source-density}, applied to the finite progression of Proposition~\ref{prop:weighted-finite-system}. Coefficients of the multilinear jets give (b); the logarithm of the moment series is the cumulant series, and constants drop out of the cumulants of orders $\ge2$, which gives (c), at every level $k$. (d) is Lemma~\ref{lem:weighted-passing-fields}, in the setting verified by Proposition~\ref{prop:weighted-canonical-history} at $x_K\ge x_{\mathrm J}$. (e) The mass $\mathsf w_t$ is at most the probability that the first large field arises at the step $t$; on the least torus Proposition~\ref{prop:ward-fixed-volume-rarity}, under $\log g^{-2}\ge\mathsf Y_{\mathrm{lt}}(\frac12)$, gives $\sum_tL^{b(K-t+1)}\mathsf w_t\le2L^b\mathsf C_{\mathrm{ag}}(b+1,\frac12)$ for every real $b\ge0$; on a larger torus the number of places costs $\ell^3e^{\mathsf B(X_g)\ell^3}$. The parts $t\le k$ of $\mathbb P_k$ are those of $\mathbb P$, and on the first part $\alpha=a_{K-k}$.
\end{proof}

\begin{lemma}[The remainder]\label{C-352-26}
Let $g$ satisfy $(\Gamma)$, $(K,\beta)$ be a first-passage datum at $g$, $\ell=\ell_g$, $n\in\{3,4\}$, and let real weights satisfy the conditions of weights placed at $(\rho,r,\mathsf d)$ (if $n=3$), or of two pairs at $(\rho,r,\mathsf d)$ (if $n=4$), but for the upper bounds on the distances of the centres, with $0<\rho\le r\le\mathsf d\le1$, $\rho\le\frac12$, $s(\rho)\le K$, $\rho\le r_1(\mathsf b)$, all supports in one ball of radius $\le\frac18\ell_g-1$. There are numbers $\Lambda_3\ge1$ and $C_{\mathrm J}$ formed from the constants such that at $\mathbb P$-almost every $\omega$, for every nonempty $\beta$,
\[
\begin{multlined}
|\mathsf Y_\beta(\omega)|\le\Lambda_3\,\alpha^{-6}\,\mathsf s(\beta)\prod_{a\in\beta}|w_a|_\infty,\qquad\mathsf s(\beta)=\rho^3\ \relax (\text{a point}),\quad\rho^6r^{-3}\ \relax (\text{a whole pair}), \\
(\rho/r)^{3c(\beta)}\rho^6\mathsf d^{-9/2}\ \relax (\beta\text{ crossing});
\end{multlined}
\]
hence $|S^T_n-\mathsf J_n|\le\mathsf k_n\Lambda_3^n\mathsf M'^n\mathsf b^n\mathsf d^{3/2}\mathrm{nat}_n$ with $\mathsf k_3=13$, $\mathsf k_4=77$, and $|\mathsf J_n|\le C_{\mathrm J}(\mathsf b/\mathsf d)^{n/2}\mathrm{nat}_n$, uniformly in $\rho\le r\le\mathsf d$ and in $K$. For two weights $w_1,w_2$ with $\Sigma(w_a)\subset\mathbb B(x_a,\rho)$ whose supports are at a distance $\ge\mathsf d$, with $0<\rho\le\mathsf d\le1$, $\rho\le\frac12$, $s(\rho)\le K$, $\rho\le r_1(\mathsf b)$ and both supports in one ball of radius $\le\frac18\ell_g-1$, at $\mathbb P$-almost every $\omega$
\[
\begin{multlined}
|\mathsf Y_a(\omega)|\le\Lambda_3\,\alpha^{-6}\rho^3|w_a|_\infty,\qquad|\mathsf Y_{12}(\omega)|\le\Lambda_3\,\alpha^{-6}\rho^6\mathsf d^{-3}|w_1|_\infty|w_2|_\infty \\
\le\Lambda_3\,\alpha^{-6}\rho^6\mathsf d^{-9/2}|w_1|_\infty|w_2|_\infty,
\end{multlined}
\]
and 
\[
|S^T_2-\mathsf J_2|\le3\Lambda_3^2\mathsf M'^2\rho^6\mathsf d^{-9/2}|w_1|_\infty|w_2|_\infty;
\]
 if besides the balls are $\ge\mathsf d$ apart, 
\[
|\mathsf J_2|\le C_{\mathrm J}(\mathsf b/\mathsf d)\rho^6\mathsf d^{-6}\mathsf b^{-2}|w_1|_\infty|w_2|_\infty.
\]

\end{lemma}

\begin{proof}
For $|\beta|\ge2$, $\mathsf Y_\beta=\sum_{i\le t'}(\mathsf{Sum}[E^{(i)}_\beta](U)-\mathsf{Sum}[E^{(i)}_\beta](\mathbf 1))-\sum_{t'<i\le K}\mathsf{Sum}[E^{(i)}_\beta](\mathbf 1)$, $U=U_{t'}(W)$; at a point $\mathcal O(w_a)(U)-\mathcal O(w_a)(\mathbf 1)$ comes besides, and for $t'\ge1$ it is $\le162B_3^2\varepsilon_{t'}^2\rho^3\alpha^{-4}|w_a|_\infty$ on the regular configuration of Lemma~\ref{C-352-25}(d). The first sum is bounded by Lemma~\ref{C-352-14}(b) at the regular background of Lemma~\ref{C-352-25}(d), the second at $\mathbf 1$, both with the sizes of Lemma~\ref{C-352-15}(f) at the width $\rho$, whose conditions hold: (P1), (P2) and the separation of the corners follow, as in the proof of Lemma~\ref{C-352-11}(b), from the balls that hold the supports and from the distances of the supports alone, and $\rho\ge\frac23Ma_K$, $E_\sigma(\hat\rho_1/2)\le\varkappa_0^4$ from $\rho\le r_1(\mathsf b)$, $s(\rho)\le K$ (Lemma~\ref{C-352-12}); the sums over the levels on both sides of $t'$ are geometric.

For $t'\ge1$, with $j=t'-1$, $a_{K-j}=\alpha/L$ and $x_j>1$: $\bar\nu_{j;\beta}\le\bar c_m(L/\alpha)^m\prod|w_a|_\infty$, $\Xi^r_j\le\rho^6r^{-3}L^3\alpha^{-3}$, and $\mathsf z(\beta)\min\{1,v_j\}^{3/2}\le\mathsf s(\beta)L^{3/2}\alpha^{-3/2}$; so $|\mathsf Y_\beta|\le C\alpha^{-6}\mathsf s(\beta)\prod|w_a|_\infty$ ($\alpha\le1$, $m\le4$). The part $t=1$ of $\Omega$ has $t'=0$, $\alpha=a_K$, and there $W\in\mathcal V_0$ is any configuration and $U_0(W)=W$: $D^\lambda_0=\mathcal O(w_\lambda)$, so the first sum is empty; at a point $0\le a_p\le2$ and (P1) give $|\mathcal O(w_a)(W)|\le2a_K^{-1}\sum_p|w_a(p)|\le162\rho^3a_K^{-4}|w_a|_\infty=162\rho^3\alpha^{-4}|w_a|_\infty$, and $\mathcal O(w_a)(\mathbf 1)=0$; the second sum runs over $1\le i\le K$ at $\mathbf 1$ and is led by $i=1$, with the sizes above at $j=0$: $\bar\nu_{0;\beta}\le\bar c_m\alpha^{-m}\prod|w_a|_\infty$ ($x_0>1$), $\Xi^r_0\le\rho^6r^{-3}\alpha^{-3}$, $\mathsf z(\beta)\min\{1,v_0\}^{3/2}\le\mathsf s(\beta)\alpha^{-3/2}$.

So the bound holds at every $\omega$ of this part, with $\Lambda_3\ge162$. The weights $\mathsf v_1=\mathsf v_2=\rho^3r^{-3/2}$, $\mathsf v_3=\rho^3\mathsf d^{-9/2}$ (and $\mathsf v_a=\rho^3r^{-3/2}\mathsf d^{-9/8}$ at $n=4$) have $\mathsf s(\beta)\le\mathsf v^\beta$ and $\mathsf v^{[n]}=\mathsf s([n])$, so every product over a partition is $\le\mathsf s([n])$; Lemma~\ref{C-352-25}(c), the bounds $|\kappa_2|\le2$, $|\kappa_3|\le6$ times the mean of the product of the bounds, and the moments of $\alpha^{-1}$ give $13=1+3\cdot2+6$ at $n=3$; finally $\mathsf s([3])\prod|w_a|_\infty=\mathsf b^3\mathsf d^{3/2}\mathrm{nat}_3$. The bound of $\mathsf J_n$ is Lemma~\ref{C-352-14}(b) at $\mathbf 1$ with the bounds of Lemma~\ref{C-352-15}(e) at all the levels, in the form (f).

At two sources whose supports are at a distance $\ge\mathsf d$ the two weights form one pair at the widths $(\rho,\mathsf d)$ in the sense of Lemma~\ref{C-352-15}(f): (P1), (P2) and the separation of the corners follow from the balls $\mathbb B(x_a,\rho)$ and from the distance $\ge\mathsf d$ of the supports, and $\rho\ge\frac23Ma_K$, $E_\sigma(\hat\rho_1/2)\le\varkappa_0^4$ from $\rho\le r_1(\mathsf b)$, $s(\rho)\le K$, as above. The proof of the bound above, run for the list of the two weights, gives at a point and at the whole pair, read at $r=\mathsf d$, the display at two sources, as $\rho^6\mathsf d^{-3}\le\rho^6\mathsf d^{-9/2}$ for $\mathsf d\le1$; with $S^T_2-\mathsf J_2=\int\mathsf Y_{12}d\mathbb P+\kappa_{\mathbb P}(\mathsf Y_1,\mathsf Y_2)$ (Lemma~\ref{C-352-25}(c)), $|\kappa_2|\le2$ and the moments of $\alpha^{-1}$ it gives $3=1+2$.

If besides the balls are $\ge\mathsf d$ apart, put a dummy bump $w_2'$, the indicator of the plaquettes with barycentres in $\mathbb B(x_1+3\rho e,\rho)$, $e$ the unit vector from $x_2$ to $x_1$; the three are at the lengths $(\rho,\rho,\mathsf d)$, and by multilinearity the objects of the pair are those of the list with the dummy: $S^T_2-\mathsf J_2=\int\mathsf Y_{12}d\mathbb P+\kappa_{\mathbb P}(\mathsf Y_1,\mathsf Y_2)$, with $|\mathsf Y_a|\le\Lambda_3\alpha^{-6}\rho^3|w_a|_\infty$ and $|\mathsf Y_{12}|\le\Lambda_3\alpha^{-6}\rho^6\mathsf d^{-9/2}|w_1|_\infty|w_2|_\infty$; and the bound of $\mathsf J_n$ above, at the crossing set that the two weights form in this list, is the bound of $\mathsf J_2$.
\end{proof}

\begin{lemma}[Lattice sums and continuum integrals]\label{C-352-27}
For every side $\ell$ and every cutoff $K$, that is on the torus $T_\ell$ with the lattice $a_K\mathbb Z^3$, and with the least lifts of the differences, the following hold.
\begin{enumerate}
\item[(a)] For $4a_K\le\rho\le r\le\mathsf d\le\mathsf d^+$, $4\mathsf d^+<\ell$, $\vartheta\in(0,1]$ and Lipschitz $f_1,f_2,f_3$ with $\operatorname{supp}f_a\subset\mathbb B(x_a,\rho)$, $\operatorname{Lip}f_a\le\|f_a\|_\infty/\vartheta\rho$, $\operatorname{dist}(\operatorname{supp}f_1,\operatorname{supp}f_2)\ge r$, $\operatorname{dist}(\operatorname{supp}f_a,\operatorname{supp}f_3)\ge\mathsf d$, all supports in one ball of radius $\frac12\mathsf d^+$:
\[
\Bigl|\operatorname{Tr}\prod_a\bigl((f_a\circ x)C^{\mathrm{free}}_K\bigr)-\int k_3\prod_af_a\Bigr|\le C^{(3)}\Bigl[\frac{a_K}{\vartheta\rho}+\Bigl(\frac{\mathsf d^+}\ell\Bigr)^3\Bigr]\rho^9r^{-3}\mathsf d^{-6}\prod_a\|f_a\|_\infty .
\]
\item[(b)] For two pairs with the same data ($8a_K\le\rho$, the balls $\ge r$ apart in a pair and $\ge\mathsf d$ apart across): $|\frac12\sum_\varsigma\operatorname{Tr}\prod_l((f_{\varsigma(l)}\circ x)C^{\mathrm{free}}_K)-\int k_4\prod f_a|\le C^{(4)}[\frac{a_K}{\vartheta\rho}+(\frac{\mathsf d^+}\ell)^3]\rho^{12}r^{-6}\mathsf d^{-6}\prod\|f_a\|_\infty$.
\item[(c)] For $8a_K\le\rho\le\mathsf d$, $64\mathsf d\le\ell$, Lipschitz $\varphi_a$ in $\mathbb B(x_a,\rho)$ with $\operatorname{Lip}\varphi_a\le\|\varphi_a\|_\infty/\vartheta\rho$, the balls $\ge\mathsf d$ apart and $|x_1-x_2|\le18\mathsf d$:
\[
\Bigl|\sum_{\lambda(p)=\lambda,\,\lambda(q)=\mu}\varphi_1(x(p))\varphi_2(x(q))C^{\mathrm{free}}_K(p,q)^2-\mathcal K_{\lambda\mu}(\varphi_1,\varphi_2)\Bigr|\le C_{\mathrm{Rs}}\Bigl[\frac{a_K}{\vartheta\rho}+\Bigl(\frac{\mathsf d}\ell\Bigr)^3\Bigr]\Bigl(\frac\rho{\mathsf d}\Bigr)^6\|\varphi_1\|_\infty\|\varphi_2\|_\infty .
\]
\end{enumerate}
\end{lemma}

\begin{proof}
\emph{Step A.} In a ring, and in a square, every sign $\epsilon_p$ stands twice, so $C^{\mathrm{free}}_K(p_l,p_{l+1})$ may be replaced by $\epsilon\epsilon C^{\mathrm{free}}_K$ and then, by equations~\YMMref{(2245)--(2246)}{824}, by $a_K^3\mathbb K_3(u_l)_{\lambda(p_l)\lambda(p_{l+1})}$; both are $\le(1+3C_{\mathrm{lat}})a_K^3|u|^{-3}$, and telescoping costs $C_{\mathrm{lat}}a_K^3[a_K^2|u_l|^{-5}+\ell^{-3}]$ times the other factors, with at most $81(\rho/a_K)^3$ plaquettes for each source. \emph{Step B.} For each choice of the directions the remaining sum is a Riemann sum over the cubes of side $a_K$ about the barycentres, which for one orientation form a translate of $a_K\mathbb Z^3$; summed over the directions, $\sum\prod_l\mathbb K_3(\cdot)_{\lambda_l\lambda_{l+1}}=\operatorname{tr}\prod_l\mathbb K_3$. On the cubes that count the differences are at least $0.56$ times their least values, $\|\mathbb K_3(u)\|\le\frac1{2\pi}|u|^{-3}$, $\|\partial\mathbb K_3(u)\|\le\frac6{4\pi}|u|^{-4}$ (Lemma~\ref{C-352-29}(b)), and the Lipschitz bounds of the $f_a$ give the error $a_K/\vartheta\rho$ times the natural size; the numbers $C^{(3)},C^{(4)},C_{\mathrm{Rs}}$ are these counts.
\end{proof}

\begin{lemma}[Kernels of forms bounded on small bumps]\label{C-352-28}
Let $U_1,\dots,U_p$ be open subsets of $\mathbb R^3$ or of a torus, $\rho_0>0$, $\mathsf A\ge0$, and $T$ a $p$-linear form on $\prod_a\mathrm{Lip}_c(U_a)$ with $|T(f_1,\dots,f_p)|\le\mathsf A\rho^{3p}\prod_a\|f_a\|_\infty$ whenever $\rho\le\rho_0$ and each $f_a$ has its support in a ball of radius $\rho$ in $U_a$. Then: (a) $T(f)=\int k\prod_af_a$ with a Borel function $k$, $|k|\le(\frac3\pi)^p\mathsf A$ almost everywhere, and at every point where $T$ has a continuous kernel; (b) $T$ extends to bounded Borel functions with compact supports in the $U_a$, and for balls $\mathbb B(z_a,\rho)\subset U_a$, $\rho\le\rho_0$, the mean of $k$ over $\prod_a\mathbb B(z_a,\rho)$ is at most $(\frac3{4\pi})^p\mathsf A$ in modulus, at every choice of the centres; with the radial tents, $|\int k\prod_a\chi^\rho_{z_a}|\le(\frac3\pi)^p\mathsf A$.
\end{lemma}

\begin{proof}
Let $k_\rho(x_1,\dots,x_p)=T(\chi^\rho_{x_1},\dots,\chi^\rho_{x_p})$; $\|\chi^\rho\|_\infty=3/\pi\rho^3$ gives $|k_\rho|\le(\frac3\pi)^p\mathsf A$. On functions with supports in fixed compact sets $T$ is bounded in the sup norms (a Lipschitz partition of unity of width $\rho_0$), and the Riemann sums of $f_a*\chi^\rho$ converge uniformly, so $\int k_\rho\prod f_a=T(f_1*\chi^\rho,\dots)\to T(f)$; a weak-$*$ limit point of $k_\rho$ in $L^\infty$ is a kernel. (b) By dominated convergence $T$ extends with the same bound; at $f_a=\mathbf 1_{\mathbb B(z_a,\rho)}$, $|\int_{\prod\mathbb B(z_a,\rho)}k|\le\mathsf A\rho^{3p}$, which is the bound of the mean.
\end{proof}

\begin{lemma}[The free template]\label{C-352-29}
Let $X\ne0$, $\mathsf d=|X|$, $\mathsf b>0$, $n,n'$ unit vectors.
\begin{enumerate}
\item[(a)] $\mathbb K_3(X)^2=(16\pi^2\mathsf d^6)^{-1}(3\hat X\hat X^{\mathsf T}+1)$; $\frac32\mathsf b^{-2}\sum_\mu\{[\mathbb K_3(X)_{\lambda\mu}]^2-\frac13\sum_{\lambda'}[\mathbb K_3(X)_{\lambda'\mu}]^2\}=\mathsf t^\circ(e_\lambda,X)$; $\frac32\mathsf b^{-2}\sum_{\lambda,\mu}[\mathbb K_3(X)_{\lambda\mu}]^2=\frac9{16\pi^2}\mathsf b^{-2}\mathsf d^{-6}$; $\sum_\lambda n_\lambda^2\mathsf t^\circ(e_\lambda,X)+\sum_{\lambda<\lambda'}n_\lambda n_{\lambda'}3\mathsf b^{-2}[\mathbb K_3(X)^2]_{\lambda\lambda'}=\mathsf t^\circ(n,X)$; and $|\mathsf t^\circ(n,X)-\mathsf t^\circ(n',X)|\le0.057|n-n'|\mathsf b^{-2}\mathsf d^{-6}$.
\item[(b)] In the operator norm, $\|\mathbb K_3(y)\|=\frac2{4\pi}|y|^{-3}$, $\|(h\cdot\nabla)\mathbb K_3(y)\|\le\frac6{4\pi}|h||y|^{-4}$, $\|(h\cdot\nabla)^2\mathbb K_3(y)\|\le\frac{24}{4\pi}|h|^2|y|^{-5}$.
\item[(c)] For $0<r\le\frac12\mathsf d$ and $c-z=X$: $|\frac{2\pi}3\mathsf br^3\cdot3\mathsf b^{-3}k_3(c+\frac r2n,c-\frac r2n,z)-\mathsf t^\circ(n,X)|\le2.66\,(r/\mathsf d)^2\mathsf b^{-2}\mathsf d^{-6}$; and for $0<\rho<\frac14r$ the same function averaged against the radial tents $\chi^\rho$ about the three points differs from its value by at most $7900\,(\rho/r)\mathsf b^{-2}\mathsf d^{-6}$.
\item[(d)] For $0<\rho\le\frac18\mathsf d$: $|\mathcal K_{\lambda\mu}(\chi^\rho_c,\chi^\rho_z)-[\mathbb K_3(X)_{\lambda\mu}]^2|\le2.28\,(\rho/\mathsf d)\mathsf d^{-6}$.
\end{enumerate}
\end{lemma}

\begin{proof}
(a) With $E=\hat X\hat X^{\mathsf T}$, $E^2=E$: $(3E-1)^2=3E+1$; $\sum_\mu[\mathbb K_3(X)_{\lambda\mu}]^2=[\mathbb K_3(X)^2]_{\lambda\lambda}$ and $\operatorname{tr}\mathbb K_3(X)^2=6(16\pi^2\mathsf d^6)^{-1}$; the identity in $n$ is the expansion of $\langle n,\mathbb K_3(X)^2n\rangle$; the last bound is $\frac{18}{32\pi^2}\le0.057$ times $|n-n'|$. (b) $\mathbb K_3=(4\pi)^{-1}\nabla\nabla|y|^{-1}$; the third and fourth derivatives of $|y|^{-1}$ contracted with $h$, in a frame with $\hat y=e_1$, have eigenvalues bounded by $6$ and $24$ times $|h|^k|y|^{-3-k}$. (c) With $D(n)=3nn^{\mathsf T}-1$, $h=\frac r2n$ and $\phi(t)=\operatorname{tr}[D(n)\mathbb K_3(X-th)\mathbb K_3(X+th)]$: $\mathbb K_3(rn)=(4\pi r^3)^{-1}D(n)$, so the first member is $\frac12\mathsf b^{-2}\phi(1)$, while $\frac12\mathsf b^{-2}\phi(0)=\mathsf t^\circ(n,X)$ by (a). $\phi$ is even, since $\operatorname{tr}[DAB]=\operatorname{tr}[DBA]$ for symmetric matrices, so $|\phi(1)-\phi(0)|\le\frac12\sup|\phi''|$; with $|X\pm th|\ge\frac34\mathsf d$, the trace norm $4$ of $D(n)$ and (b),
\[
|\phi''|\le\frac4{(4\pi)^2}(24\cdot2+2\cdot6\cdot6+2\cdot24)\frac{r^2}4\Bigl(\frac43\Bigr)^8\mathsf d^{-8}=10.63\,(r/\mathsf d)^2\mathsf d^{-6},
\]
and $\frac14\cdot10.63\le2.66$. The tents: the mean value theorem with (b) for the three factors of $k_3$. (d) $|\nabla(\mathbb K_{3,\lambda\mu}^2)|\le\frac3{2\pi^2}|u|^{-7}$ on a segment in $\{|u|\ge\frac34\mathsf d\}$, and $\frac3{2\pi^2}(\frac43)^7\cdot2\le2.28$.
\end{proof}

\paragraph{Every state.}

\begin{lemma}[The chessboard bound]\label{C-352-30}
Let $0<g\le\gamma$, $(K,\beta)$ a first-passage datum at $g$, $\ell'\in\mathfrak L(g)$, and consider the board of $\mathbb L_K$ into cubes of side $\frac12\ell'$, the cells, with any offset in $a_K\mathbb Z^3$; let $\mathsf Q_0=\exp[\frac18\ell'^3(2C_{\mathrm S}+C_g)]$ with the constants $C_{\mathrm S},C_g$ of the two-sided stability bound. For every $\ell\in\mathfrak L(g)$ with $\ell\ge\ell'$, every $\nu\in\mathcal G^{\mathrm{per}}_{K,\beta}$, every cell and every bounded complex Borel function $F$ of the variables of that cell,
\[
|\tilde\mu_{K,\ell,\beta}(F)|,\ |\nu(F)|\le\mathsf Q_0\|F\|_{\mathrm{ch}},\qquad\|F\|_{\mathrm{ch}}^8=\tilde\mu_{K,\ell',\beta}\Bigl(\prod_{\text{the eight cells of }T_{\ell'}}F^{[\cdot]}\Bigr)\ge0,
\]
$F^{[\cdot]}$ are the reflected copies of $F$, conjugated when an odd number of coordinate reflections is used.
\end{lemma}

\begin{proof}
The bound is a chain of four links. The cells are cubes of side $\mathsf l=\frac12\ell'$; the torus $T_\ell$ holds $\mathsf N_\ell=(2\ell/\ell')^3$ of them and $T_{\ell'}$ eight. For $F$ as in the statement let $F^{\mathrm{all}}_\ell$ be the product of the reflected copies of $F$, conjugated with the parity convention in the statement, over all the cells of $T_\ell$, and $Z_{K,\ell,\beta}[F^{\mathrm{all}}_\ell]$ the integral of the density of $T_\ell$ with this insertion. (1) \emph{Reflection positivity on $T_\ell$.} The torus measures are reflection positive in the planes of the board, and the chessboard estimate (Corollary~\ref{IV-18-18}(b), (c)) gives $|\tilde\mu_{K,\ell,\beta}(F)|\le\|F\|_{\mathrm{ch},\ell}$, where $\|F\|_{\mathrm{ch},\ell}^{\mathsf N_\ell}=Z_{K,\ell,\beta}[F^{\mathrm{all}}_\ell]/Z_{K,\ell,\beta}\ge0$; on $T_{\ell'}$ this is the number of the statement, $\|F\|_{\mathrm{ch},\ell'}=\|F\|_{\mathrm{ch}}$. (2) \emph{Two sides, with the insertions.} As $\ell/\ell'$ is a power of $L$, the subadditivity in the volume of the integrals with the reflected copies inserted (Proposition~\YMMref{86.5}{682}, Corollary~\YMMref{86.6}{683}) gives $Z_{K,\ell,\beta}[F^{\mathrm{all}}_\ell]^{1/\mathsf N_\ell}\le Z_{K,\ell',\beta}[F^{\mathrm{all}}_{\ell'}]^{1/8}$; with $p_\ell=\ell^{-3}\log Z_{K,\ell,\beta}$, $Z_{K,\ell,\beta}^{1/\mathsf N_\ell}=e^{\mathsf l^3p_\ell}$ and $Z_{K,\ell',\beta}^{1/8}=e^{\mathsf l^3p_{\ell'}}$ this is clause~(e) of the same corollary:
\[
\|F\|_{\mathrm{ch},\ell}\le\|F\|_{\mathrm{ch},\ell'}\,e^{\mathsf l^3(p_{\ell'}-p_\ell)} .
\]
(3) \emph{Stability, with no insertion.} The two-sided stability bound among the construction results, which rests on Theorems~\YMMref{AJ}{673} and~\YMMref{AD}{607} (Lemma~\ref{lem:canonical-source-system}), concerns the partition functions without insertions: $Z_{K,\ell,\beta}=e^{E_{K,\ell,\beta}}\widetilde Z_\ell$ with the vacuum constant $E_{K,\ell,\beta}$ and a number $e^{-C_g\ell^3}\le\widetilde Z_\ell\le e^{C_{\mathrm S}\ell^3}$, and $|E_{K,\ell,\beta}\ell^{-3}-E_{K,\ell',\beta}\ell'^{-3}|\le C_{\mathrm S}$ for two sides $\ell,\ell'\in\mathfrak L(g)$. So
\[
p_{\ell'}-p_\ell\le\bigl(E_{K,\ell',\beta}\ell'^{-3}+C_{\mathrm S}\bigr)-\bigl(E_{K,\ell,\beta}\ell^{-3}-C_g\bigr)\le2C_{\mathrm S}+C_g ,
\]
and (1)--(3), with $\mathsf l^3=\frac18\ell'^3$, give $|\tilde\mu_{K,\ell,\beta}(F)|\le\mathsf Q_0\|F\|_{\mathrm{ch}}$. (4) \emph{Periodic Gibbs states.} For $\nu\in\mathcal G^{\mathrm{per}}_{K,\beta}$ the domination theorem (Proposition~\YMMref{91.9}{705}, for every board of $\mathbb L_K$ and every offset) gives $|\nu(F)|\le\liminf_k\|F\|_{\mathrm{ch},\ell_k}$ along every sequence $\ell_k\to\infty$ in $\mathfrak L(g)$, and (2), (3) at $\ell=\ell_k$ bound the right member by $\mathsf Q_0\|F\|_{\mathrm{ch}}$. No link reads what $F$ is made of, and $\mathsf Q_0$ depends neither on $\ell\ge\ell'$ nor on $K$.
\end{proof}

\begin{lemma}[The fold]\label{C-352-31}
Let $\mathfrak N_I=\mathbb C[\lambda_\iota:\iota\in I]/(\lambda_\iota^2)$, $[x]_\gamma$ the coefficient of $\lambda^\gamma$, $\mathfrak m=\{[x]_\emptyset=0\}$, and $\|x\|_u=\sum_\gamma|[x]_\gamma|/u^\gamma$ for numbers $u_\iota>0$.
\begin{enumerate}
\item[(a)] $\|\cdot\|_u$ is submultiplicative; for $y\in\mathfrak m$, $\|e^y-1\|_u\le(1+\|y\|_u)^{|I|}-1$ and $\|\log(1+y)\|_u\le\|y\|_u(1+\|y\|_u)^{|I|-1}$; the logarithm of the moment series is the cumulant series.
\item[(b)] In the data of Lemma~\ref{C-352-30} let $f_1,\dots,f_n$ be bounded real functions that vanish off the interior of one cell, $w_\iota$, $\iota\in I=[n]\times\{\text{the eight cells}\}$, the eight mirror images of the $f_a\circ x$ on $T_{\ell'}$, and $(\Omega,\mathbb P)$ a probability space with a measurable $\mathbb W\colon\Omega\to\mathfrak m$ such that $\langle\prod_{\iota\in\gamma}\mathcal O(w_\iota)\rangle_{K,\ell',\beta}=\int[\exp\mathbb W]_\gamma d\mathbb P$ for all $\gamma\subset I$. Let $u_{(a,\cdot)}=u_a>0$, $\mathsf c=\sum\mathsf c^{(\cdot)}$ a sum of constants of $\mathfrak m$ each in the variables of one cell, of any size, and $\mathsf Z\ge0$ measurable with $\|\mathbb W-\mathsf c\|_u\le\mathsf Z$ almost everywhere. Put $\mathsf y=\int[(1+\mathsf Z)^n-1]d\mathbb P$, $\mathsf h=\mathsf y(1+\mathsf y)^{n-1}$, $\eta_0=[\int(1+\mathsf Z+8\mathsf h)^{8n}d\mathbb P]^{1/8}$, $\eta=(2^n-1)\mathsf Q_0\eta_0$. Then for every nonempty $\beta\subset[n]$, every $\ell\in\mathfrak L(g)$ with $\ell\ge\ell'$, and $\mathfrak s=\tilde\mu_{K,\ell,\beta}$ or $\mathfrak s\in\mathcal G^{\mathrm{per}}_{K,\beta}$:
\[
\bigl|S^T_{|\beta|}[\mathfrak s]((f_a)_{a\in\beta})-S^T_{|\beta|}[\tilde\mu_{K,\ell',\beta}]((f_a)_{a\in\beta})\bigr|\le\eta(1+\eta)^{n-1}u^\beta .
\]
\end{enumerate}
The bound holds at every size of $\mathsf Z$, of the constants $\mathsf c$ and of $\eta$.
\end{lemma}

\begin{proof}
(a) is the algebra of $\exp$ and $\log$ in an algebra with $|I|$ nilpotent generators. (b) Let $\mathsf k=\log\sum_\beta\tilde\mu_{K,\ell',\beta}(\prod_{a\in\beta}\mathcal O(f_a))\lambda^\beta$ and $\Xi_\beta=[\exp(\sum_a\lambda_a\mathcal O(f_a)-\mathsf k)]_\beta$, a function of the variables of the cell. The mirror images of $\Xi_\beta$ are the same expressions in the mirrored sources, so by the representation $\|\Xi_\beta\|^8_{\mathrm{ch}}=\int[\exp(\mathbb W-\mathsf k^{\mathrm{all}})]_{\beta\times\{\text{cells}\}}d\mathbb P$, $\mathsf k^{\mathrm{all}}$ the sum of the eight copies of $\mathsf k$. Each copy has $e^{\mathsf k^{(\cdot)}}=\int\exp(\pi_{\text{cell}}\mathbb W)d\mathbb P$, so $\mathsf k^{(\cdot)}-\mathsf c^{(\cdot)}=\log(1+\mathsf y^{(\cdot)})$ with $\|\mathsf y^{(\cdot)}\|_u\le\mathsf y$ and $\|\mathsf k^{(\cdot)}-\mathsf c^{(\cdot)}\|_u\le\mathsf h$; the constants of any size cancel, $\|\mathbb W-\mathsf k^{\mathrm{all}}\|_u\le\mathsf Z+8\mathsf h$, and $|[\exp(\cdot)]_{\beta\times\{\text{cells}\}}|\le(1+\mathsf Z+8\mathsf h)^{8n}(u^\beta)^8$.

So $\|\Xi_\beta\|_{\mathrm{ch}}\le\eta_0u^\beta$, and Lemma~\ref{C-352-30} gives $|\mathfrak s(\Xi_\beta)|\le\mathsf Q_0\eta_0u^\beta$. With $\mathsf b'=\sum_\beta\mathfrak s(\Xi_\beta)\lambda^\beta=\mathsf M_{\mathfrak s}e^{-\mathsf k}$, $\mathsf M_{\mathfrak s}$ the moment series of $\mathfrak s$: $\|\mathsf b'-1\|_u\le\eta$ and $\|\log\mathsf M_{\mathfrak s}-\mathsf k\|_u=\|\log\mathsf b'\|_u\le\eta(1+\eta)^{n-1}$.
\end{proof}

\begin{lemma}[Mirror copies]\label{C-352-32}
Zero sources are omitted before forming the weighted jet norm; coefficients involving their directions vanish by multilinearity. Assume the remaining source weights are nonzero. Let $g$ satisfy $(\Gamma)$, $\ell'=\ell_g$, and let $f_1,\dots,f_n$, $n\in\{2,3,4\}$, be as in (A.4)(i) (for $n=4$, two pairs; for $n=2$, two bounded real Borel functions with $\operatorname{supp}f_a\subset\mathbb B(x_a,\rho)$, the two supports at a distance $\ge\mathsf d$, $0<\rho\le\mathsf d\le1$, $\rho\le\frac12$, $s(\rho)\le K$, $\rho\le r_1(\mathsf b)$, both supports in one ball of radius $\le\frac18\ell_g-1$), placed by the offset in the interior of one cell, so that two mirror copies are at least $\frac14\ell_g\ge6.5$ apart. Let $\mathbb W$ be the multilinear jet of $\lambda\mapsto D^\lambda_{t'}(U_{t'}(W))$ for the $8n$ sources $w_\iota$ on the space of Lemma~\ref{C-352-25}, $[\mathsf c]_\gamma=\partial^\gamma D^\lambda_K(\mathbf 1)|_0$ for $\gamma$ in one cell and $[\mathsf c]_\gamma=0$ for $\gamma$ that link two cells, and $u_a=\mathsf v_a|w_a|_\infty$ with $\mathsf v_1=\mathsf v_2=\rho^3r^{-3/2}$, $\mathsf v_3=\rho^3\mathsf d^{-9/2}$ ($n=3$), $\mathsf v_a=\rho^3r^{-3/2}\mathsf d^{-9/8}$ ($n=4$), $\mathsf v_a=\rho^3\mathsf d^{-9/4}$ ($n=2$).

Then $\|\mathbb W(\omega)-\mathsf c\|_u\le\mathsf Z_n\alpha(\omega)^{-6}$ almost everywhere, $\mathsf Z_n=8(2^n-1)\Lambda_3+2^{8n}\max_m\Lambda^{\mathrm L}_m$, where the numbers $\Lambda^{\mathrm L}_m$, formed from the constants, bound the linking jets: for $\gamma$ linking two copies, $|\sum_{i\le t}\mathsf{Sum}[E^{(i)}_\gamma](U)|\le\Lambda^{\mathrm L}_m\mathsf b^{-m/2}\prod_{\text{copies}}\mathsf s(\gamma\cap\text{copy})\prod_{\iota\in\gamma}|w_\iota|_\infty$ at every $U$ whose pairs of the levels $\le t$ lie in the cores.
\end{lemma}

\begin{proof}
Within one cell the coefficients of $\mathbb W-\mathsf c$ are the $\mathsf Y_\beta$ of the copy (by multilinearity the mixture names no source), bounded by Lemma~\ref{C-352-26} and $\mathsf s(\beta)\le\mathsf v^\beta$; there are $8(2^n-1)$ such sets. At $n=2$ the two bumps are taken as one pair at the widths $(\rho,\mathsf d)$ in the sense of Lemma~\ref{C-352-15}(f), since (P1), (P2) and the separation of the corners follow from the balls $\mathbb B(x_a,\rho)$ and from the distance $\ge\mathsf d$ of the supports, as in the proof of Lemma~\ref{C-352-26}; that proof, with the size $\rho^6\mathsf d^{-3}$ of a whole pair at the distance $\mathsf d$ in the place of $\rho^6r^{-3}$, gives at $\mathbb P$-almost every $\omega$
\[
|\mathsf Y_a|\le\Lambda_3\alpha^{-6}\rho^3|w_a|_\infty,\qquad|\mathsf Y_{12}|\le\Lambda_3\alpha^{-6}\rho^6\mathsf d^{-3}|w_1|_\infty|w_2|_\infty\le\Lambda_3\alpha^{-6}\rho^6\mathsf d^{-9/2}|w_1|_\infty|w_2|_\infty\quad(\mathsf d\le1),
\]
and $\rho^3\le\mathsf v_a$, $\rho^6\mathsf d^{-9/2}=\mathsf v_1\mathsf v_2$. For a set $\gamma$ that links two copies, the calculus of Lemma~\ref{C-352-12} at $8n$ sources, with the copies as cells and $\mathsf D^{\iota\iota'}_i\ge\hat{\mathsf D}^\times_i/\sqrt3-1$ across, where $\mathsf D^\times\ge\frac14\ell_g$ is the least distance on $T_{\ell_g}$ of the supports of two different mirror copies and $\hat{\mathsf D}^\times_i=\mathsf D^\times/a_{K-i}$, gives by the recursion (g), with $\sigma_{i;\gamma}\le E_\sigma(\hat{\mathsf D}^\times_i)$, bounds on $\bar{\mathsf h}_{j;\gamma}$ and $\bar\Psi_{j;\gamma}$ by products of the one-cell sizes; with Lemma~\ref{C-352-14}(b) and the geometric sums this is the bound of the linking jets. The products of natural sizes over a partition do not exceed that of the whole set. The numbers $\Lambda^{\mathrm L}_m$ are finite and formed from the constants, two copies being at least $\frac14\ell_g\ge6.5$ apart; no condition on $g$ beyond $(\Gamma)$ is needed.
\end{proof}

\begin{proposition}[Transfer to every state]\label{C-352-33}
Let $g$ satisfy $(\Gamma)$ and $(K,\beta)$ be a first-passage datum at $g$.
\begin{enumerate}
\item[(a)] For $n\in\{3,4\}$ and the data of (A.4)(i) (for $n=4$ two pairs, all supports in one ball of radius $\le\frac18\ell_g-1$), every $\ell\in\mathfrak L(g)$ and every state $\mathfrak s$ at $(K,\beta)$: $|S^T_n[\mathfrak s](f)-S^T_{n;K,\ell_g,\beta}(f^{\langle\ell_g\rangle})|\le C^{\mathrm{tr}3}_n(g)\mathsf b^n\mathsf d^{3/2}\mathrm{nat}_n$, where $C^{\mathrm{tr}3}_n(g)=\eta(1+\eta)^{n-1}$, $\eta=(2^n-1)\mathsf Q_0(g)\eta_0$, $\eta_0=(1+\mathsf Z_n+8\mathsf h)^n(\bar{\mathsf M}^{(48n)})^{1/8}$, $\mathsf h=\mathsf y(1+\mathsf y)^{n-1}$, $\mathsf y=(1+\mathsf Z_n)^n\bar{\mathsf M}^{(6n)}$, $\bar{\mathsf M}^{(b)}=1+2L^b\mathsf C_{\mathrm{ag}}(b+1,\frac12)$ and $\mathsf Q_0(g)=\exp[\frac18\ell_g^3(2C_{\mathrm S}+C_g)]$.
\item[(b)] For two bumps: let $0<\rho\le\mathsf d\le1$, $\rho\le\frac12$, $s(\rho)\le K$, $\rho\le r_1(\mathsf b)$, $f_1,f_2$ bounded real Borel functions on $\mathbb R^3$ with $\operatorname{supp}f_a\subset\mathbb B(x_a,\rho)$, the two supports at a distance $\ge\mathsf d$, both in one ball of radius $\le\frac18\ell_g-1$, $t_a\colon\{1,2,3\}\to[-1,1]$, $w_a(p)=t_a(\lambda(p))f_a(x(p))$, and $\bar w_a$ the same with $f^{\langle\ell_g\rangle}_a$ on $T_{\ell_g}$. Then for every state $\mathfrak s$ at $(K,\beta)$, $|\langle\mathcal O(w_1);\mathcal O(w_2)\rangle_{\mathfrak s}-\langle\mathcal O(\bar w_1);\mathcal O(\bar w_2)\rangle_{K,\ell_g,\beta}|\le C^{\mathrm{tr}}_2(g)(\rho/\mathsf d)^6\mathsf d^{3/2}\|f_1\|_\infty\|f_2\|_\infty$, where $C^{\mathrm{tr}}_2(g)=\eta(1+\eta)$ is the number of (a) formed at $n=2$, with $\mathsf Z_2$ of Lemma~\ref{C-352-32}.
\end{enumerate}
\end{proposition}

\begin{proof}
Every periodic Gibbs state is invariant under lattice translations by Proposition~\YMMref{103.4}{758}; choose the offset of the board so that the tuple lies in the interior of one cell. Lemma~\ref{C-352-31} with $\ell'=\ell_g$, the representation of Lemma~\ref{C-352-25}(b) at the $8n$ mirrored sources, $\mathsf Z=\mathsf Z_n\alpha^{-6}$ by Lemma~\ref{C-352-32}, and the moments of Lemma~\ref{C-352-25}(e): $\mathsf y\le(1+\mathsf Z_n)^n\bar{\mathsf M}^{(6n)}$ and $\eta_0\le(1+\mathsf Z_n+8\mathsf h)^n(\bar{\mathsf M}^{(48n)})^{1/8}$, so $\eta$ is formed from the constants and $g$. With the weights of Lemma~\ref{C-352-32}, $u^{[n]}=\mathsf s([n])\prod|w_a|_\infty=\mathsf b^n\mathsf d^{3/2}\mathrm{nat}_n$. For (b) the same argument runs at $n=2$, with Lemma~\ref{C-352-32} at $n=2$, whose bounds read only the balls of radius $\rho$ and the distance $\ge\mathsf d$ of the supports; a mirror map carries $a_p$ to $a_{\mathsf rp}$ with $x(\mathsf rp)=\mathsf rx(p)$ and $\lambda(\mathsf rp)=\lambda(p)$, so plane weights go to plane weights, and $u_1u_2=(\rho/\mathsf d)^6\mathsf d^{3/2}|w_1|_\infty|w_2|_\infty$.
\end{proof}

\begin{lemma}[Limits and components]\label{C-352-34}
\begin{enumerate}
\item[(a)] Every Gibbs sequence at $g$ that converges has subsequences on which the labels converge (they lie in $(g^{-2},\bar{\mathsf b}]$); along every such subsequence the members are states at their data, the limit $\sigma$ is a law on $\mathbb R^3$ with $S^T_n[\nu_j]\to S^\sigma_n$ on separated tuples, and for every component $\sigma_\xi$ of a canonical decomposition of $\sigma$ its approximants $\rho_j$ are periodic Gibbs states at the same data with $S^T_n[\rho_j]\to S^{\sigma_\xi}_n$.
\item[(b)] For a law with continuous moment kernels through order
three, the connected kernel is
$\mathsf s_3=\mathsf s^{\mathrm m}_3-\mathsf s_2\otimes\mathsf s_1
-\dots+2\mathsf s_1\otimes\mathsf s_1\otimes\mathsf s_1$,
and $\mathsf s^T_2=\mathsf s_2-\mathsf s_1\otimes\mathsf s_1$.
These kernels are continuous. For the periodic-Gibbs limits and
components in (a), they are also real analytic and Euclidean
invariant, by Proposition~\YMMref{108.3}{771}.
\item[(c)] Let $\eta(t)=1$ $(t\le0)$, $1-3t^2+2t^3$ $(0\le t\le1)$, $0$ $(t\ge1)$, $\psi_\rho(v)=\eta(5|v|/\rho-4)$ and $f^\rho_a=\psi_\rho(\cdot-z_a)/\int\psi_\rho$. These are $C^1$ probability densities supported in $\mathbb B(z_a,\rho)$ with $\rho^9\prod_{a\le3}\|f^\rho_a\|_\infty\le0.102$; if $\mathsf s_3$ is continuous at $(z_1,z_2,z_3)$ and
$|\int(\mathsf s_3-3\mathsf b^{-3}k_3)\prod f^\rho_a|\le\mathsf A$
for all small $\rho$, then $|\mathsf s_3-3\mathsf b^{-3}k_3|(z_1,z_2,z_3)\le\mathsf A$.
\end{enumerate}
\end{lemma}

\begin{proof}
(a) follows from the component approximation in Theorem~\YMMref{AP}{780} and the canonical decomposition of Theorem~\YMMref{AL}{746}; a joint cumulant of order $\ge2$ does not change if constants are added to its entries, so the centring of the approximants does not enter. (b) Apply the finite moment--cumulant formula to the assumed
continuous moment kernels. For the laws in (a), their analytic
and invariance properties follow from Theorem~\YMMref{AP}{780}
and Proposition~\YMMref{108.3}{771}. For a self-centred limit,
$S^\sigma_1=0$. (c) $\psi_\rho$ is constant on $|v|\le\frac45\rho$, and $\mathsf s_3-3\mathsf b^{-3}k_3$ is continuous at the point.
\end{proof}

\begin{lemma}[One subsequence for all pairs of bumps]\label{C-352-35}
Let $\Omega$ be an open subset of $\{(x,y):0<|x-y|<\mathsf d_\Omega\}$ in $T_\ell^2$ or in $(\mathbb R^3)^2$, and let bilinear forms $\mathsf t^{(j)}_{\lambda\mu}$, $\lambda,\mu\in\{1,2,3\}$, $j\in\mathbb N$, on bounded real Borel functions be such that for every $(x_0,y_0)\in\Omega$, with closed balls $\bar U,\bar V$ about $x_0,y_0$ of a radius $\varrho(x_0,y_0)>0$ for which $\bar U\times\bar V\subset\Omega$, there is $\mathsf B<\infty$ with $|\mathsf t^{(j)}_{\lambda\mu}(f,h)|\le\mathsf B\|f\|_\infty\|h\|_\infty$ for all large $j$ and all $f,h$ supported in $\bar U,\bar V$. Then there is a subsequence along which $\mathsf t^{(j)}_{\lambda\mu}(\varphi_1,\varphi_2)$ converges for all $\lambda,\mu$ and all real Lipschitz $\varphi_1,\varphi_2$ with compact supports, $\operatorname{supp}\varphi_1\times\operatorname{supp}\varphi_2\subset\Omega$; the limits are bilinear. If on every such product of balls the limit forms satisfy the assumption of Lemma~\ref{C-352-28}, they have Borel kernels on $\Omega$, unique almost everywhere.
\end{lemma}

\begin{proof}
$\mathrm{Lip}_c(U)$ is a subset of the separable space $C_0(U)$, so it has a countable dense subset in the sup norm. Take a countable dense set of centres in $\Omega$ and a diagonal subsequence along which the countably many bounded sequences $\mathsf t^{(j)}_{\lambda\mu}(f,h)$, $f,h$ in these dense sets, converge; by the uniform bound and bilinearity the sequences converge at all Lipschitz $f,h$ in the balls. A pair $\varphi_1,\varphi_2$ as in the claim is decomposed by a Lipschitz partition of unity into cubes of a side small against the distances of $\operatorname{supp}\varphi_1\times\operatorname{supp}\varphi_2$ to the diagonal and to the boundary of $\Omega$: finitely many pairs, each in one product of balls. Two kernels that represent one form on an open product agree almost everywhere (mollification), and a numbering of the countable cover glues the local kernels into one Borel function.
\end{proof}

\paragraph{The bond forms.}

\begin{lemma}[Dilation to two couplings below the ceiling]\label{C-352-36}
Let $0<g\le\gamma$, let $z=z(g)\ge1$ be the least integer with $L^{z-1}(g^{-2}+\bar B)-\bar B\ge x_{\mathrm N}=\gamma_{\mathrm N}^{-2}$ (Definition~\ref{C-352-2}(1)), and define $g'_1,g'_2$ by $g_1'^{-2}=L^z(g^{-2}+\bar B)-\bar B$, $g_2'^{-2}=L^{-1}(g_1'^{-2}+2B_\sharp)-B_\sharp$.
\begin{enumerate}
\item[(a)] $0<g'_1<g'_2\le\gamma_{\mathrm N}$, and both depend on $g$ and the constants alone.
\item[(b)] For a first-passage datum $(K,\beta)$ at $g$ with $K\ge z+1$ there is $\iota\in\{1,2\}$ such that $(K-z,\beta)$ is a first-passage datum at $g'_\iota$; its run is $(x_0,\dots,x_{K-z})$, and its label is $L^z\mathsf b$.
\item[(c)] $L^z\ell\in\mathfrak L(g'_\iota)$ for $\ell\in\mathfrak L(g)$; if $s(t)\ge z$, then $s'_\iota(L^zt)\le s(t)-z$ for the witness depth $s'_\iota$ at $g'_\iota$, and if $s(t)<z$, then $L^zt>1$.
\item[(d)] $U\mapsto U^{(z)}$, $U^{(z)}(L^zb)=U(b)$, carries $\mu_{K,\ell,\beta}$ to $\mu_{K-z,L^z\ell,\beta}$, with $a_{L^zp}(U^{(z)})=a_p(U)$ and $q_{L^zb}(U^{(z)})=q_b(U)$. So $\mathcal O^T(\mathfrak s)(U)=L^z\mathcal O'^T(\mathfrak s^{(z)})(U^{(z)})$ and $\langle\mathcal O^T(\mathfrak s_1);\mathcal O^T(\mathfrak s_2)\rangle_{K,\ell,\beta}=L^{2z}\langle\mathcal O'^T(\mathfrak s^{(z)}_1);\mathcal O'^T(\mathfrak s^{(z)}_2)\rangle_{K-z,L^z\ell,\beta}$, primes marking the objects of the torus $T_{L^z\ell}$ at the cutoff $K-z$ and $\mathfrak s^{(z)}(L^z\cdot)=\mathfrak s(\cdot)$; and the numbers $\|\varphi\|_\infty$, $r\operatorname{Lip}\varphi$, $a_K/r$, $\mathsf d'/\ell$, $r/\mathsf d'$ and the integrals $\iint\varphi_1\varphi_2\mathbb K_3(x-y)_{\lambda\gamma}\mathbb K_3(x-y)_{\lambda'\gamma'}$ are the same for $(\varphi_a;r,\mathsf d',\ell,K)$ and for $(\varphi^{(z)}_a;L^zr,L^z\mathsf d',L^z\ell,K-z)$.

Hence each display of (D.1) at $(K,\ell,\beta)$ is the same display at $(K-z,L^z\ell,\beta)$, with the same tolerance, both members being multiplied by $L^{-2z}$. The identity $\mathcal O^T(\mathfrak s)(U)=L^z\mathcal O'^T(\mathfrak s^{(z)})(U^{(z)})$ holds at every configuration $U$ of $\mathbb L_K$ as well, the primes now marking the fields of $\mathbb L_{K-z}$; in particular $\mathcal O(f)(U)=L^z\mathcal O'(f^{(z)})(U^{(z)})$ with $f^{(z)}=f(L^{-z}\cdot)$. So for every law $\nu$ on the configurations of $\mathbb L_K$ and its image $\nu^{(z)}$ under the dilation, $S^T_n[\nu](f_1,\dots,f_n)=L^{nz}S^T_n[\nu^{(z)}](f^{(z)}_1,\dots,f^{(z)}_n)$ for bounded real functions $f_a$ with compact supports, and the joint cumulants of the fields $\mathcal O^T(\mathfrak s_a)$ of mixed sources scale in the same way; and $\nu\mapsto\nu^{(z)}$ is a bijection of $\mathcal G^{\mathrm{per}}_{K,\beta}$ onto $\mathcal G^{\mathrm{per}}_{K-z,\beta}$.

These statements of (d) hold with every integer $0\le z'<K$ in the place of $z$.
\end{enumerate}
\end{lemma}

\begin{proof}
(a), (b) are arithmetic of the inverse couplings. Put $G=g^{-2}+\bar B$ and $X_{g'}=L(g'^{-2}+B_\sharp)$ for a coupling $g'$. One step of the run is $x_{j-1}=L(x_j+b_j)$ with $|b_j|\le B_\sharp$ (the flow bound in Section~\ref{sec:setup-first-passage-arithmetic}); as $\bar B=B_\sharp L/(L-1)$ sums the geometric series, $z$ steps give $L^z(x_K-\bar B)+\bar B\le x_{K-z}\le L^z(x_K+\bar B)-\bar B$, and with $g^{-2}<x_K\le X_g$ and $X_g+\bar B=LG$:
\[
L^z(g^{-2}-\bar B)+\bar B<x_{K-z}\le L^{z+1}G-\bar B .
\]
(a) $g_2'^{-2}=L^{z-1}G-\bar B+2B_\sharp/L\ge L^{z-1}G-\bar B\ge x_{\mathrm N}$ by the choice of $z$, and $g_1'^{-2}-g_2'^{-2}=(L-1)L^{z-1}G-2B_\sharp/L>0$; both read $g$ and the constants alone. (b) $X_{g'_1}=L^{z+1}G-\bar B$ is the upper end above, and $X_{g'_2}=g_1'^{-2}+2B_\sharp$; so the half-open intervals $I_\iota=(g_\iota'^{-2}+2B_\sharp,X_{g'_\iota}]$, $\iota=1,2$, abut, are disjoint, and $I_2\cup I_1=(g_2'^{-2}+2B_\sharp,X_{g'_1}]$. For the lower end, $g_2'^{-2}+2B_\sharp\le L^{z-1}G+2\bar B$ and
\[
\bigl[L^z(g^{-2}-\bar B)+\bar B\bigr]-\bigl[L^{z-1}G+2\bar B\bigr]=L^{z-1}\bigl[(L-1)g^{-2}-(L+1)\bar B\bigr]-\bar B\ge(L^{z-1}-1)\bar B\ge0,
\]
since $\bar B\le\frac{13}{24}g^{-2}$ at $g\le\gamma$ (Lemma~\ref{lem:setup-stopped-arithmetic} gives this bound) gives $(L-1)g^{-2}\ge\frac{24}{13}(L-1)\bar B\ge(L+2)\bar B$, as $11L\ge50$. So $x_{K-z}$ lies in exactly one $I_\iota$. There $x_{K-z}\le X_{g'_\iota}$ and $x_{K-z-1}=L(x_{K-z}+b_{K-z})\ge L(x_{K-z}-B_\sharp)>L(g_\iota'^{-2}+B_\sharp)=X_{g'_\iota}$, and the run decreases: $K-z$ is the first index at which the run that starts at $\beta$ is at most $X_{g'_\iota}$, which is the first passage at the coupling $g'_\iota\le\gamma_{\mathrm N}\le\gamma$. The run of $(K-z,\beta)$ is the same sequence cut at $K-z$, and its label is $a_{K-z}\beta=L^z\mathsf b$.

The margin $2B_\sharp$ in the lower ends of the $I_\iota$ is what gives $x_{K-z-1}>X_{g'_\iota}$, the second half of the first passage. (c) Put $G=g^{-2}+\bar B$. Then $g_1'^{-2}+\bar B=L^zG$ and $g_2'^{-2}+\bar B\le L^zG$. The radius-scaling inequality~\YMMref{(1534)}{616}, applied to $X=LG$ and $X'=L(g_\iota'^{-2}+\bar B)$, gives $R(X')\le L^zR(X)$. The admissibility formula~\YMMref{(1528)}{615} therefore holds with $m+z$ at $g'_\iota$ whenever it holds with $m$ at $g$; its fixed $C_{\mathrm s}$ threshold is also preserved. Moreover $C_W^{g'_\iota}(s-z)\le C_W^g(s)$ for $s\ge z$, directly from the displayed definition of $C_W$. Thus a witness at depth $s(t)$ transfers to depth $s(t)-z$. If $s(t)<z$, then $L^{s(t)}t\ge C_W^g(s(t))\ge1$, whence $L^zt>1$. (d) is the dilation; $\mathsf b'=L^z\mathsf b$, so $\mathsf b^{-2}=L^{2z}\mathsf b'^{-2}$, and $\mathbb K_3(L^zu)=L^{-3z}\mathbb K_3(u)$ makes the integrals invariant. The identity of the fields reads of the configuration only $a_{L^zp}(U^{(z)})=a_p(U)$ and $q_{L^zb}(U^{(z)})=q_b(U)$, and of the lattice only $x(L^zp)=L^zx(p)$, $x_{L^zb}=L^zx_b$ and $a_{K-z}=L^za_K$; so it holds at every configuration of $\mathbb L_K$, and the joint cumulants under a law are multilinear in the fields.

The dilation maps the bonds, the plaquettes and the finite sets of bonds of $\mathbb L_K$ onto those of $\mathbb L_{K-z}$. The Gibbs kernel of a finite set of bonds has the exponent $-\beta\sum_pa_p$ over the plaquettes that contain a bond of the set, which reads $\beta$ and the incidences of the bonds and the plaquettes alone; so the image of a Gibbs state at $(a_K,\beta)$ satisfies the equations of the kernels at $(a_{K-z},\beta)$, and invariance under a subgroup of finite index of $a_K\mathbb Z^3$ becomes invariance under its image under $x\mapsto L^zx$, a subgroup of finite index of $a_{K-z}\mathbb Z^3$. The inverse dilation carries the periodic Gibbs states at $(K-z,\beta)$ back in the same way, so $\nu\mapsto\nu^{(z)}$ is a bijection of $\mathcal G^{\mathrm{per}}_{K,\beta}$ onto $\mathcal G^{\mathrm{per}}_{K-z,\beta}$ (equation~\YMMref{(1819)}{698} and the definition of periodic states following it).

The same arguments hold at every integer $0\le z'<K$.
\end{proof}

\begin{lemma}[Bond sources]\label{C-352-36b}
Let $0<g\le\gamma_{\mathrm N}$, $(K,\beta)$ a first-passage datum at $g$ with $K\ge2$, $\ell\in\mathfrak L(g)$, and let $\mathfrak s=(w,\mathsf w)$ be a real mixed source whose barycentres $x(p)$, $w(p)\ne0$, and midpoints $x_b$, $\mathsf w(b)\ne0$, lie in a ball of a radius $\rho\ge a_K$. Put $\mathsf m_j=12M^3\alpha_0^2L^{K-j}|\mathfrak s|_\infty$ and $\mathsf u=\mathsf b^{-1}|\mathfrak s|_\infty$, let $\varsigma_j$ be as in Lemma~\ref{C-352-12}(e) at the width $\rho$, and let
\[
\mathsf W_{\mathfrak s}(p,q)=w(p)\delta_{pq}-\sum_b\mathsf w(b)\bigl[u_{b,\lambda_b}(p)u_{b,\lambda'_b}(q)+u_{b,\lambda'_b}(p)u_{b,\lambda_b}(q)\bigr],
\]
a symmetric matrix that joins two plaquettes only through a bond of the support, with the coefficients $u_{b,\lambda}(p)$ of the paragraph on the fields. There are numbers $C$ formed from the constants such that:
\begin{enumerate}
\item[(1)] the weight member of the source $L^{K-j}\mathfrak s$ at the level $j$ is linear in $\mathfrak s$, vanishes on the domains that hold no corner of a plaquette $p$ with $w(p)\ne0$ and no site of a bond $b$ with $\mathsf w(b)\ne0$, and its norm in $\|\cdot\|_{j,\kappa}$ is at most $C\varsigma_j\mathsf m_j$;
\item[(2)] at the first level, $\|E^{(1)}_{\mathfrak s}\|_{1,\lambda_{\mathrm{out}}}\le C\theta^{[0]}\mathsf m_0$;
\item[(3)] for $1\le n\le K$, $A\in\mathsf{Cc}_n(r_n)$, $\tilde U=e^{i\eta_nA}$ and $F=\mathsf{ang}^{[n]}(A)$: 
\[
\bigl|\mathcal O^T(\mathfrak s)(\tilde U)-\frac1{2a_K}\sum_{p,q}\mathsf W_{\mathfrak s}(p,q)\langle F_p,F_q\rangle\bigr|\le C\mathfrak b_n^2\delta_n\hat\rho_n^3\mathsf u,
\]
and $\eta_n^4\sum_{p,q}|\mathsf W_{\mathfrak s}(p,q)|/(a_Kx_n)\le C\hat\rho_n^3\mathsf u$;
\item[(4)] Lemma~\ref{C-352-20}(c) and Proposition~\ref{C-352-21} hold with the matrices $\mathsf W_{\mathfrak s_a}$ in the place of the diagonal matrices $w_a$ in $R_{i;\beta}$, $q_{i;\beta}$ and $g_{i;\beta}$, for two mixed sources in the frame of (6), the tower identity of (6) standing in Step~1 of the proof of Proposition~\ref{C-352-21};
\item[(5)] the weighted progression and the representation of the torus expectations through the never-large histories, which name no source, are formed for mixed sources at $g\le\gamma_{\mathrm N}$ on every side $\ell\in\mathfrak L(g)$, with the weight member of (1) in the source slot of the bare arrays; the lemma on the tower identity at several sources (it precedes Proposition~\ref{C-352-21}), for clusters in the frame of (6), the representation of Lemma~\ref{C-352-25} and the almost everywhere bound of Lemma~\ref{C-352-26} on the remainders hold for mixed sources, the moments of $\alpha^{-1}$ being those of Lemma~\ref{C-352-25}(e) on that side;
\item[(6)] \emph{the frame and the tower identity.} The hull of a mixed source $\mathfrak s=(w,\mathsf w)$ is the set of the plaquettes $p$ with $w(p)\ne0$ and of the four plaquettes through each bond $b$ with $\mathsf w(b)\ne0$. For $2\le k<K$, two mixed sources, neither $0$, are in the frame at $k$, that is they satisfy the four conditions \textup{(Jt.1)}--\textup{(Jt.3)} and \textup{(Jt.6)} of the frame at $k$ of Lemma~\ref{C-352-11}(a), read with the hulls in the place of the supports and with the members of (1) in the slot, if each hull has its barycentres in a closed ball of a radius $\bar{\mathsf d}\ge a_K$ with $\bar{\mathsf d}/a_{K-k}\ge M$ and the corners of the plaquettes of the two hulls are at least $3M$ level-$k$ spacings apart in the sup distance; this holds at $k=\mathsf k(\mathsf d)$ for hulls in balls of radius $\le\mathsf d$ at a distance $\ge\mathsf d$, $0<\mathsf d\le e^{-2\sqrt3}$, $s(\mathsf d)\le K$, as in Lemma~\ref{C-352-11}(b). In the frame, for $k\le i<K$, $A'\in\mathsf{Cc}_{i+1}(r_{i+1})$, $\tilde U=e^{i\eta_{i+1}A'}$ and $\lambda$ near $0$,
\[
\exp D^\lambda_{i+1}(\tilde U)=\int\exp D^\lambda_i\bigl(e^{i\eta_iA^{\mathrm{lw}}(B)}\bigr)\,d\nu_{i,A'}(B),
\]
with the measure $\nu_{i,A'}$ of the step and the lowering map; and the one-source Lemma~\ref{C-352-16} holds for each source $\mathfrak s$ of such a pair, with $\frac1{2a_K}\sum_{p,q}\mathsf W_{\mathfrak s}(p,q)\langle F_p,F_q\rangle$ in the place of $q_{n;a}(F)$, its frame read with the hulls and its count (P1) replaced by the second bound of (3).
\end{enumerate}
\end{lemma}

\begin{proof}
(1) $\mathcal O^T(\mathfrak s)=a_K^{-1}\sum_pw(p)a_p-a_K^{-1}\sum_b\mathsf w(b)q_b$, and $q_b$ is a gauge invariant function of the variables of the four plaquettes through $b$; so the source is a local term of the bare level, linear in $\mathfrak s$, which lives at the sites of its support. On the cores the plaquette variables are within $C\alpha_0\eta_j^2$ of $\mathbf 1$, so that $|q_b|\le|u_{b,\lambda_b}||u_{b,\lambda'_b}|\le C\alpha_0^2\eta_j^4$ there, as $|a_p|<\frac12\alpha_0^2\eta_j^4$; at most $81(\rho/a_K)^3$ bonds have their midpoints in a ball of the radius $\rho\ge a_K$; with $|\mathfrak s|_\infty=|w|_\infty+8|\mathsf w|_\infty$ the count of Lemma~\ref{C-352-12}(e) gives the bound. (2) The one-step theorem at the height $0$, with this member in the fourth place, and Schwarz's lemma, as in Lemma~\ref{C-352-12}(e). (3) On the chain class the angle of a plaquette variable of $\tilde U$ is $\pm F_p$ up to $C\eta_n^2r_n^2$, and $\mathrm w(e^{i\theta})=\theta+O(|\theta|^3)$; so $u_{b,\lambda}=\sum_pu_{b,\lambda}(p)F_p$ up to $C\eta_n^2r_n^2$, and, as $|F_p|\le2\eta_n^2r_n$, $|q_b-\langle\sum_pu_{b,\lambda_b}(p)F_p,\sum_qu_{b,\lambda'_b}(q)F_q\rangle|\le C\eta_n^4r_n^3$, the bound that the expansion of $a_p$ on the chain class has.

Summed against $\mathsf w$, this gives the bond part of $\frac1{2a_K}\langle F,\mathsf W_{\mathfrak s}F\rangle$ with the error $Ca_K^{-1}(\rho/a_K)^3|\mathsf w|_\infty\eta_n^4r_n^3$; with the plaquette part as in the proof of Lemma~\ref{C-352-16}, $a_{K-n}x_n\ge\frac45\mathsf b$ and $x_nr_n^2\le C\mathfrak b_n^2$, this is the first bound. Each $u_{b,\lambda}$ has two coefficients of modulus $\frac12$, so $\sum_{p,q}|\mathsf W_{\mathfrak s}(p,q)|\le\sum_p|w(p)|+2\sum_b|\mathsf w(b)|\le81(\rho/a_K)^3|\mathfrak s|_\infty$, and $\eta_n^4a_K^{-4}=a_{K-n}^{-4}$ gives the second. (4) In the proof of Lemma~\ref{C-352-20}(c) the sources enter only through $R^\lambda_i=\mathsf W_\lambda(1-\mathcal S_i\mathsf W_\lambda)^{-1}$, $\mathsf W_\lambda=\sum_a\lambda_a\mathsf W_{\mathfrak s_a}$, which is symmetric as $\mathsf W_\lambda$ and $\mathcal S_i$ are, and through the Gaussian integral of a quadratic form with a symmetric matrix; the identity $1-\mathsf sR^\lambda_i=(1-\mathcal S_{i+1}\mathsf W_\lambda)(1-\mathcal S_i\mathsf W_\lambda)^{-1}$ does not use that $\mathsf W_\lambda$ is diagonal.

The proof of Proposition~\ref{C-352-21} uses the sources only through the tower identity of its Step~1, which (6) gives for two mixed sources in the frame, through (I1)--(I3), through the sums $\sum_{p,q}|R_{i;\beta}(p,q)|$, which the second bound of (3) controls as (P1) does for plaquette weights, and through the lowered angles at the plaquettes of the hulls. (5) Propositions~\ref{prop:weighted-finite-system} and~\ref{prop:running-source-density}, together with \eqref{eq:ward-first-creation}, give the weighted progression and the history representation for sources on the bare plaquettes, on every side of $\mathfrak L(g)$ with $g^{-2}\ge x_{\mathrm J}$ (Theorem~\YMMref{D}{8}). For mixed sources the same construction is made with the weight member of (1) in the source slot of the bare arrays; the other data of the progression name no source.

The proofs among the construction results read a source only through that member, and of the member they read these properties: the entry of the objects into the settings of one step, for any number of sources, $K\ge2$ and every side, reads that it is a local term of the bare level with the bound of (1); the never-large histories and the measures of the mixture are formed with no source; the representation reads that $\mathcal O^T(\mathfrak s)$ is a bounded gauge invariant function, linear in $\mathfrak s$ and $0$ at $\mathbf 1$, whose member is that of (1), so that the quotient at a level $t$ of $e^{\mathcal O^T(\mathfrak s_\lambda)}$ is $\exp D^\lambda_t$; and the tower identity reads the member in the slot of the step map, which takes any member there with the bound of (1).

By (1) a mixed source has these properties. The proofs of the lemma on the tower identity at several sources, which precedes Proposition~\ref{C-352-21}, and of Lemma~\ref{C-352-25} use the sources only through the progression and the tower identity for a pair, which (6) gives for mixed sources; the almost everywhere bound of Lemma~\ref{C-352-26} uses them only through (1), (2), the support rule of Lemma~\ref{C-352-12}(d) and the sizes of Lemma~\ref{C-352-15}(f), which hold for mixed sources by (1)--(3) with the hulls in the place of the supports; on a side $\ell$ the moments of $\alpha^{-1}$ are bounded by Lemma~\ref{C-352-25}(e). (6) The frame at $k$ is formed by the four conditions \textup{(Jt.1)}--\textup{(Jt.3)} and \textup{(Jt.6)} of the setting of the pair functional among the construction results, read with the hulls in the place of the supports: the proof of Lemma~\ref{C-352-11}(a), using Definition~\ref{def:pair-functional-frame}, reads of the two sources only the four geometric properties listed there, now of the hulls, which hold all the plaquettes that the members of the sources read, and the members in the slot.

The tower identity for a pair is that among the construction results, read with the member of (1) in the slot: its proof forms the measure $\nu_{i,A'}$ and the lowering map with no source and reads the sources only through the relative exponents. Hulls at a distance $\ge\mathsf d$ have corners at least $(\mathsf d-\sqrt2a_K)/\sqrt3$ apart in the sup distance, which is at least $3M$ level-$\mathsf k(\mathsf d)$ spacings by the working level, as in the proof of Lemma~\ref{C-352-11}(b); a ball of radius $\mathsf d$ has $\mathsf d/a_{K-\mathsf k(\mathsf d)}\ge4M$. In Lemma~\ref{C-352-16} the member of $\mathfrak s$ is expanded on the chain class by the first bound of (3), and the families by the count of its proof with the second bound of (3).
\end{proof}

\begin{lemma}[Lattice sums at two balls on every torus]\label{C-352-36a}
On every torus $T_\ell$ with the lattice $a_K\mathbb Z^3$, with the least lifts of the differences, let $4a_K\le r$, $16r\le\mathsf d'$, $8\mathsf d'\le\ell$, $x_1,x_2\in T_\ell$ with $|x_1-x_2|=\mathsf d'$, $\vartheta\in(0,1]$, and let $\varphi_1,\varphi_2$ be real Lipschitz functions with $\operatorname{supp}\varphi_a\subset\mathbb B(x_a,r)$ and $\operatorname{Lip}\varphi_a\le\|\varphi_a\|_\infty/\vartheta r$. Put $\mathsf e_{\mathrm{PS}}=[a_K/\vartheta r+(\mathsf d'/\ell)^3](r/\mathsf d')^6\|\varphi_1\|_\infty\|\varphi_2\|_\infty$.
\begin{enumerate}
\item[(a)] For all $\lambda,\mu$: $\bigl|\sum_{\lambda(p)=\lambda,\,\lambda(q)=\mu}\varphi_1(x(p))\varphi_2(x(q))\,C^{\mathrm{free}}_K(p,q)^2-\mathcal K_{\lambda\mu}(\varphi_1,\varphi_2)\bigr|\le C_{\mathrm{PS}}\,\mathsf e_{\mathrm{PS}}$.
\item[(b)] Let $\{\lambda,\lambda',\mu\}=\{1,2,3\}$, $\gamma$ a direction, $\mathsf w(b)=\varphi_1(x_b)\mathbf 1_{\{\mu_b=\mu\}}$, $w_2(q)=\varphi_2(x(q))\mathbf 1_{\{\lambda(q)=\gamma\}}$, and $C_\alpha(b,q)=\sum_pu_{b,\alpha}(p)C^{\mathrm{free}}_K(p,q)$ for a direction $\alpha\ne\mu$. Then
\[
\Bigl|-\sum_{b,q}\mathsf w(b)w_2(q)\,C_\lambda(b,q)C_{\lambda'}(b,q)-\iint\varphi_1(x)\varphi_2(y)\,\mathbb K_3(x-y)_{\lambda\gamma}\mathbb K_3(x-y)_{\lambda'\gamma}\,dx\,dy\Bigr|\le C^{\mathrm{od}}_{\mathrm{PS}}\,\mathsf e_{\mathrm{PS}} .
\]
\end{enumerate}
\end{lemma}

\begin{proof}
Points of the two balls are at least $\mathsf d'-2r\ge\frac78\mathsf d'$ and at most $\mathsf d'+2r\le\frac98\mathsf d'<\frac14\ell$ apart, so all the least lifts are taken in one ball of radius $\frac14\ell$. (a) Steps A and B of the proof of Lemma~\ref{C-352-27} at these distances. In a square every sign $\epsilon_p$ stands twice, so equations~\YMMref{(2245)--(2246)}{824} replace $C^{\mathrm{free}}_K(p,q)^2$ by $a_K^6[\mathbb K_3(u)_{\lambda\mu}]^2$ at the cost $C_{\mathrm{lat}}(2+3C_{\mathrm{lat}})a_K^6|u|^{-3}[a_K^2|u|^{-5}+\ell^{-3}]$ for each pair, and a ball of radius $r\ge a_K$ holds at most $27(r/a_K)^3$ barycentres of one orientation; with $|u|\ge\frac78\mathsf d'$ and $a_Kr\le\frac14r^2\le\frac1{1024}\mathsf d'^2$ this costs at most $729(\frac87)^3C_{\mathrm{lat}}(2+3C_{\mathrm{lat}})\,\mathsf e_{\mathrm{PS}}\le1089\,C_{\mathrm{lat}}(2+3C_{\mathrm{lat}})\,\mathsf e_{\mathrm{PS}}$.

The Riemann sums over the cubes of side $a_K$ about the barycentres, which for one orientation form a translate of $a_K\mathbb Z^3$, with $\|\mathbb K_3(u)\|\le\frac1{2\pi}|u|^{-3}$, $\|\partial\mathbb K_3(u)\|\le\frac6{4\pi}|u|^{-4}$ (Lemma~\ref{C-352-29}(b)) and the Lipschitz bounds of the $\varphi_a$, give the rest of $C_{\mathrm{PS}}\mathsf e_{\mathrm{PS}}$. (b) $u_{b,\alpha}(p)\ne0$ only at the two plaquettes $p_\pm$ with the barycentres $x_b\pm\frac12a_Ke_\alpha$, which lie in the plane of $\mu$ and $\alpha$, and there $u_{b,\alpha}(p_\pm)=\frac12\epsilon_{\alpha\mu\alpha'}\epsilon_{p_\pm}$, $\{\alpha,\mu,\alpha'\}=\{1,2,3\}$. By equation~\YMMref{(2245)}{824}, $C_\lambda(b,q)$ is $\epsilon_{\lambda\mu\lambda'}\epsilon_qa_K^3$ times the mean of $\mathbb K_3(\cdot)_{\lambda'\gamma}$ at the two differences $x(p_\pm)-x(q)$, and $C_{\lambda'}(b,q)$ is $\epsilon_{\lambda'\mu\lambda}\epsilon_qa_K^3$ times the like mean of $\mathbb K_3(\cdot)_{\lambda\gamma}$, up to the terms of the comparison.

As $\epsilon_{\lambda\mu\lambda'}\epsilon_{\lambda'\mu\lambda}=-1$ and $\epsilon_q^2=1$, $-C_\lambda(b,q)C_{\lambda'}(b,q)$ is $a_K^6\,\mathbb K_3(x_b-x(q))_{\lambda\gamma}\mathbb K_3(x_b-x(q))_{\lambda'\gamma}$ up to these terms and to the shifts by $\pm\frac12a_Ke_\lambda$, $\pm\frac12a_Ke_{\lambda'}$, which cost $a_K\|\partial\mathbb K_3\|$. The midpoints of the bonds of one direction form a translate of $a_K\mathbb Z^3$, as the barycentres of the plaquettes of one orientation do, and steps A and B as in (a) give $C^{\mathrm{od}}_{\mathrm{PS}}\mathsf e_{\mathrm{PS}}$. The numbers $C_{\mathrm{PS}}$ and $C^{\mathrm{od}}_{\mathrm{PS}}$ of Definition~\ref{C-352-2}(3) are these counts.
\end{proof}

\begin{proposition}[Two mixed sources on every side]\label{C-352-37}
For $0<g\le\gamma_{\mathrm N}$ let $t_{\mathrm N}(g)\in(0,\frac12e^{-2\sqrt3}]$ be the threshold, formed from the constants and $g$, below which the conditions of the induction at two sources with the base level equal to the start level hold at $g$ and at every label in $(g^{-2},\bar{\mathsf b}]$ (the conditions of step~5 of the proof of Proposition~\ref{C-352-24} that its step~6 reads, with the members that $g\le\gamma_{\mathrm N}$ gives), and put $r_g=\frac12\min\{r_1(\bar{\mathsf b}),t_{\mathrm N}(g)\}$. For $\ell\in\mathfrak L(g)$ there is a finite function $\epsilon^{\mathrm{ev}}_\sharp(g,\ell,\cdot)$ on $(0,\frac18\ell]$, formed from the constants, $g$ and $\ell$ (it is written out at the end of the proof), with $\epsilon^{\mathrm{ev}}_\sharp(g,\ell,\mathsf d')\le C(g,\ell)(\log\frac1{\mathsf d'})^{19\mathfrak p_0}\mathsf d'^{1/2}$ for $\mathsf d'\le\frac76t_{\mathrm N}(g)$, such that the following holds. Let $(K,\beta)$ be a first-passage datum at $g$, $\mathsf b=a_K\beta$, $x_1,x_2\in T_\ell$ with $|x_1-x_2|=\mathsf d'$, $0<16r\le\mathsf d'\le\frac18\ell$, $s(\min\{r,r_g\})\le K$. For real mixed sources $\mathfrak s_a=(w_a,\mathsf w_a)$ with $\{x(p):w_a(p)\ne0\}\cup\{x_b:\mathsf w_a(b)\ne0\}\subset\mathbb B(x_a,r)$,
\[
\Bigl|\langle\mathcal O^T(\mathfrak s_1);\mathcal O^T(\mathfrak s_2)\rangle_{K,\ell,\beta}-\tfrac32\mathsf b^{-2}\operatorname{Tr}\bigl(\mathsf W_{\mathfrak s_1}C^{\mathrm{free}}_K\mathsf W_{\mathfrak s_2}C^{\mathrm{free}}_K\bigr)\Bigr|\le\epsilon^{\mathrm{ev}}_\sharp(g,\ell,\mathsf d')\Bigl(\frac r{\mathsf d'}\Bigr)^6\mathsf b^{-2}|\mathfrak s_1|_\infty|\mathfrak s_2|_\infty,
\]
where $\mathsf W_{\mathfrak s}(p,q)=w(p)\delta_{pq}-\sum_b\mathsf w(b)[u_{b,\lambda_b}(p)u_{b,\lambda'_b}(q)+u_{b,\lambda'_b}(p)u_{b,\lambda_b}(q)]$ and $u_{b,\alpha}(p)$ are the coefficients of the plaquette angles in the linear part of $u_{b,\alpha}$, given in the paragraph on the fields (it is the matrix of the lemma on bond sources, which follows Lemma~\ref{C-352-36}). Consequently both displays of (D.1) hold at $g\le\gamma_{\mathrm N}$, with $\epsilon^{\mathrm{ev}}_\sharp$ in place of $\epsilon^{\mathrm{cl}}$, for every $K$ with $s(\min\{r,r_g\})\le K$.
\end{proposition}

\begin{proof}
A bond weight enters the progression through its weight member, a local term on the plaquettes of the two planes that contain the bond. Put $\rho=\min\{r,r_g\}$; the working level at $\rho$ ($\rho\le\frac14e^{-2\sqrt3}$, $s(\rho)\le K$) gives $a_K\le\rho/4M\le\frac1{52}\rho$. The hull of $\mathfrak s_a$ (clause (6) of the lemma on bond sources, which follows Lemma~\ref{C-352-36}) has its barycentres in $\mathbb B(x_a,r+\frac12a_K)$, so that for $p,q$ in the two hulls
\[
\tfrac67\mathsf d'\le|x(p)-x(q)|\le\tfrac87\mathsf d'=:\mathsf d^+,
\]
all these points lying in a ball of radius $<\frac14\ell$, where the least lifts are taken. Put $\mathsf d=\min\{\frac67\mathsf d',t_{\mathrm N}(g)\}$, the floor of the distances that is read.

\emph{The case $r<r_g$.} The width $r_T=r+\frac12a_K\le\frac{105}{104}r$ has $r_T\le\mathsf d$, $s(r_T)\le K$, and $r_T\le\frac12\cdot1.01\,r_1(\bar{\mathsf b})\le r_1(\mathsf b)$, as $r_1$ does not increase and $\mathsf b\le\bar{\mathsf b}$; so the two sources are in the frame at $\mathsf k(\mathsf d)$ (clause (6) of the lemma on bond sources). By the lemma on bond sources the calculus of Lemma~\ref{C-352-12}, the sizes of Lemma~\ref{C-352-15}(f), the one-source Lemma~\ref{C-352-16} and the induction at two sources (Proposition~\ref{C-352-21}, base level equal to start level $\mathsf k(\mathsf d)$, as in step~6 of the proof of Proposition~\ref{C-352-24}) apply to the mixed sources, with the explicit form $\frac1{2a_K}\langle F,\mathsf W_{\mathfrak s}F\rangle$ of a source, the orientation signs entering through $u_{b,\lambda}(p)$. Of $(\Gamma1)$ these use only the members that $g\le\gamma_{\mathrm N}$ gives: the member $\delta(g^{-2})\le1/2C_{\mathrm{cu}}$ is used only for the kernel families (Proposition~\ref{C-352-19}), which do not occur at two sources whose distance fixes the base level. The difference of the two members of the display is at most the sum of three members.

(i) \emph{The remainder.} $S^T_2-\mathsf J_2=\int\mathsf Y_{12}\,d\mathbb P+\kappa_{\mathbb P}(\mathsf Y_1,\mathsf Y_2)$ (Lemma~\ref{C-352-25}(c) for mixed sources, by clause (5) of the lemma on bond sources). The two sources form one pair at the widths $(r_T,\mathsf d)$ in the sense of Lemma~\ref{C-352-15}(f), so the proof of Lemma~\ref{C-352-26}, with the sizes of clauses (1)--(3) of the lemma on bond sources, gives at $\mathbb P$-almost every $\omega$ the bounds $|\mathsf Y_a|\le\Lambda_T\alpha^{-6}r_T^3|\mathfrak s_a|_\infty$ and $|\mathsf Y_{12}|\le\Lambda_T\alpha^{-6}r_T^6\mathsf d^{-9/2}|\mathfrak s_1|_\infty|\mathfrak s_2|_\infty$, with a number $\Lambda_T\ge1$ formed from the constants. With the moments of $\alpha^{-1}$ on the side $\ell$, $\int\alpha^{-24}d\mathbb P\le\bar{\mathsf M}_\ell:=1+4L^{24}M^{-3}\ell^3e^{\mathsf B(X_g)\ell^3}\mathsf C_{\mathrm{ag}}(25,\eth)$ (Lemma~\ref{C-352-25}(e)), which grows with $\ell$, this member is at most $3\Lambda_T^2\bar{\mathsf M}_\ell^2\bar{\mathsf b}^2\mathsf d^{3/2}\,(r_T/\mathsf d)^6\mathsf b^{-2}|\mathfrak s_1|_\infty|\mathfrak s_2|_\infty$.

(ii) \emph{The induction.} The Gaussian step of Lemma~\ref{C-352-20}(c) with $\mathsf W_{\mathfrak s}$ in place of $w$ gives at the top the value $g_{K;12}=\frac32\operatorname{Tr}(\mathsf W_{\mathfrak s_1}\mathcal S_K\mathsf W_{\mathfrak s_2}\mathcal S_K)$, $\mathcal S_K=\sum_{\mathsf k(\mathsf d)<j\le K}\upsilon_j\mathcal C_{(j)}$, Wick's rule for two quadratic forms of the photon field. Proposition~\ref{C-352-21} and Lemma~\ref{C-352-15}(d), as in step~6 of the proof of Proposition~\ref{C-352-24}, bound $|\mathsf J_2-g_{K;12}|$ by $\mathsf E_{\mathrm J}(g,\mathsf d)(r_T/\mathsf d)^6\mathsf b^{-2}|\mathfrak s_1|_\infty|\mathfrak s_2|_\infty$, with a function $\mathsf E_{\mathrm J}$ of $(g,\mathsf d)$ alone that is at most $C(g)(\log\frac1{\mathsf d})^{19\mathfrak p_0}\mathsf d^{1/2}$ for $\mathsf d\le t_{\mathrm N}(g)$.

(iii) \emph{The free kernel.} For $p,q$ in the two hulls, $\mathsf r=|x(p)-x(q)|_\infty\in[\mathsf d/\sqrt3,\mathsf d^+]$ and $t=\mathsf r/a_{K-\mathsf k(\mathsf d)}\ge\hat{\mathsf d}_{\mathsf k(\mathsf d)}/\sqrt3$, the floor of the working level, at which the condition of Lemma~\ref{C-352-22}(b) holds by the conditions that fix $t_{\mathrm N}(g)$ (as condition (ii) of step~5 of the proof of Proposition~\ref{C-352-24}). So Lemma~\ref{C-352-22}(b) at $k=\mathsf k(\mathsf d)$, multiplied by $\mathsf r^3$, gives $|\mathcal S_K-\mathsf b^{-1}C^{\mathrm{free}}_K|(p,q)\le\mathsf b^{-1}a_K^3\mathsf r^{-3}\,\mathsf e_{\mathrm{fr}}(g;\mathsf d,\mathsf d^+)$, with $\mathsf b^{-1}<g^2$ and
\[
\mathsf e_{\mathrm{fr}}(g;\mathsf d,\mathsf d^+)=\tfrac54C_{\mathrm{tp}}c^{(2)}_L\bar Bg^2\,\mathsf d^++C^{\mathrm{bg}}_{\mathrm{tp}}(\mathsf d^+)^3+2L^3C_{\mathrm{tp}}\sup_{t\ge\hat{\mathsf d}_{\mathsf k(\mathsf d)}/\sqrt3}t^3e^{-\delta_{\mathrm{tp}}Lt}.
\]
Telescoping the two kernels in $\operatorname{Tr}(\mathsf W_{\mathfrak s_1}\cdot\,\mathsf W_{\mathfrak s_2}\cdot\,)$, with Lemma~\ref{C-352-22}(a), $|C^{\mathrm{free}}_K(p,q)|\le(1+3C_{\mathrm{lat}})a_K^3|u|^{-3}$, $\mathsf r^{-3}\le3^{3/2}\mathsf d^{-3}$ and the second bound of clause (3) of the lemma on bond sources for $\sum_{p,q}|\mathsf W_{\mathfrak s_a}(p,q)|$, this member is at most $C\,\mathsf e_{\mathrm{fr}}(g;\mathsf d,\mathsf d^+)(r_T/\mathsf d)^6\mathsf b^{-2}|\mathfrak s_1|_\infty|\mathfrak s_2|_\infty$: a member polynomial in $\mathsf d^+$, besides the far term at the floor.

The members (i) and (ii), which carry the renormalisation group, read of the distances the floor $\mathsf d$ alone; the largest distance $\mathsf d^+$ enters the one explicit member (iii), the comparison of the sum of the slices with the free kernel, through a bound polynomial in $\mathsf d^+$.

\emph{The case $r\ge r_g$.} Then $\rho=r_g$ and $s(r_g)\le K$. Tile the lift by the half-open cubes $Q$ of side $r_g$ and cut $\mathfrak s_a=\sum_Q\mathfrak s^Q_a$, $\mathfrak s^Q_a$ the restriction of $\mathfrak s_a$ to the plaquettes and the bonds whose barycentres or midpoints lie in $Q$; $|\mathfrak s^Q_a|_\infty\le|\mathfrak s_a|_\infty$, and at most $(2r/r_g+2)^3\le64(r/r_g)^3$ pieces of each source are not $0$. The hull of a piece is a part of the hull of $\mathfrak s_a$ and has its barycentres within $\frac{\sqrt3}2r_g+\frac12a_K\le r_g$ of the centre of its cube; so every pair of pieces, neither $0$, is in the data of the first case at the widths $(r_g,\mathsf d)$, with $r_g\le\mathsf d$, $s(r_g)\le K$, $r_g\le r_1(\mathsf b)$, the centres of the cubes in the place of $x_1,x_2$, and the same bounds $\frac67\mathsf d'\le|x(p)-x(q)|\le\mathsf d^+$ on the hulls. Both members of the display are bilinear in $(\mathfrak s_1,\mathfrak s_2)$, as $\mathcal O^T$ and $\mathsf W_{\mathfrak s}$ are linear in $\mathfrak s$, and the sum over the pairs of pieces costs $64^2(r/r_g)^6(r_g/\mathsf d)^6=4096(r/\mathsf d)^6$.

\emph{Together.} $(\mathsf d'/\mathsf d)^6=(\frac76)^6\max\{1,6\mathsf d'/7t_{\mathrm N}(g)\}^6$ and $4096(\frac76)^6(\frac{105}{104})^6\le11060$ give the display with
\[
\begin{multlined}
\epsilon^{\mathrm{ev}}_\sharp(g,\ell,\mathsf d')=11060\max\Bigl\{1,\frac{6\mathsf d'}{7t_{\mathrm N}(g)}\Bigr\}^6\Bigl[\mathsf E_{\mathrm J}(g,\mathsf d)+C\,\mathsf e_{\mathrm{fr}}\bigl(g;\mathsf d,\tfrac87\mathsf d'\bigr)+3\Lambda_T^2\bar{\mathsf M}_\ell^2\bar{\mathsf b}^2\mathsf d^{3/2}\Bigr],\qquad\mathsf d \\
=\min\{\tfrac67\mathsf d',t_{\mathrm N}(g)\},
\end{multlined}
\]
a finite function on $(0,\frac18\ell]$ formed from the constants, $g$ and $\ell$. For $\mathsf d'\le\frac76t_{\mathrm N}(g)$ one has $\mathsf d=\frac67\mathsf d'$ and $\mathsf d^+=\frac43\mathsf d$, and each of the three members is at most $C(g,\ell)(\log\frac1{\mathsf d'})^{19\mathfrak p_0}\mathsf d'^{1/2}$: the first by step~6 of the proof of Proposition~\ref{C-352-24}, the second as the member $\mathrm{br}$ of its step~4 under the conditions of its step~5, the third as $\mathsf d^{3/2}\le\mathsf d^{1/2}$.

\emph{The plane and the bond forms.} For Lipschitz bumps the lattice form is turned into the integrals of (D.1) by Lemma~\ref{C-352-36a}, whose data hold: $4a_K\le r$, $16r\le\mathsf d'$, $8\mathsf d'\le\ell$. For the plane fields take $\mathfrak s_a=(\mathbf 1_{\{\lambda(\cdot)=\lambda_a\}}\varphi_a\circ x,0)$, $\lambda_1=\lambda$, $\lambda_2=\mu$, so that $|\mathfrak s_a|_\infty\le\|\varphi_a\|_\infty$ and the main term is $\frac32\mathsf b^{-2}\sum_{p,q}w_1(p)w_2(q)C^{\mathrm{free}}_K(p,q)^2$; Lemma~\ref{C-352-36a}(a) times $\frac32\mathsf b^{-2}$ gives the second display of (D.1). For the bond form take $\mathfrak s_1=(0,\mathsf w)$, $\mathsf w(b)=\varphi_1(x_b)\mathbf 1_{\{\mu_b=\mu\}}$, so that $Y[\mathsf w]=\mathcal O^{\mathrm{od}}_{\lambda\lambda'}(\varphi_1)$ and $|\mathfrak s_1|_\infty\le8\|\varphi_1\|_\infty$, and $\mathfrak s_2=(\mathbf 1_{\{\lambda(\cdot)=\gamma\}}\varphi_2\circ x,0)$.

For the symmetric matrix $C=C^{\mathrm{free}}_K$, $\operatorname{Tr}(\mathsf W_{\mathfrak s_1}C\mathsf W_{\mathfrak s_2}C)=-2\sum_{b,q}\mathsf w(b)w_2(q)C_\lambda(b,q)C_{\lambda'}(b,q)$ with $C_\alpha(b,q)=\sum_pu_{b,\alpha}(p)C(p,q)$, so the main term is $-3\mathsf b^{-2}\sum_{b,q}\mathsf w(b)w_2(q)C_\lambda(b,q)C_{\lambda'}(b,q)$, and Lemma~\ref{C-352-36a}(b) times $3\mathsf b^{-2}$ gives the first display, with $8\epsilon^{\mathrm{ev}}_\sharp$ in the place of $8\epsilon^{\mathrm{cl}}$.
\end{proof}

\paragraph{Pair minus tensor field.}

\begin{lemma}[A regular gauge on a cube; the curvature is H\"older continuous]\label{C-352-37a}
Fix $\beta_{\mathrm H}\in(0,1)$. There are numbers $\epsilon_{\mathrm G}\in(0,1]$ and $C^\sharp_{\mathrm G},C_{\mathrm H}\ge1$, formed from the constants and $\beta_{\mathrm H}$, such that the following holds for $1\le n\le K$; lengths are in level-$n$ units, $\xi=\eta_n$, and $Q_s(x_0)$ is the set of the bare sites at the sup distance at most $s$ from a bare site $x_0$.
\begin{enumerate}
\item[(a)] Let $12a_{K-n}\le\ell$ and let $A$ be a real $\mathfrak{su}(2)$-valued $1$-form on the bonds of $Q_2(x_0)$ such that $|A|$, $|\nabla^\xi A|$ and the current of $e^{i\xi A}$ are at most $\mathsf r_\circ\le\frac1{120}$ there. Then the linear angles $F_p=\xi(\mathrm dA)(p)$ satisfy $|F_p|\le2\xi^2\mathsf r_\circ$ and $|F_{p_{\mu\nu}(x)}-F_{p_{\mu\nu}(x')}|\le C_{\mathrm H}\,\xi^2\mathsf r_\circ|x-x'|^{\beta_{\mathrm H}}$ for the plaquettes $p_{\mu\nu}(x)$, $p_{\mu\nu}(x')$ of one orientation with $x,x'\in Q_{1/2}(x_0)$.
\item[(b)] Let $W\in\mathcal V^\chi_n$, $U=U_n(W)$, and let $\varrho\ge1$ with $(8\varrho+2)a_{K-n}<\ell$ and $(4\varrho)^2(1+4\varrho)\,4B_3\varepsilon_n\le\epsilon_{\mathrm G}$. Then there are $\mathsf g\colon Q_{4\varrho}(x_0)\to SU(2)$ and a real $1$-form $A$ on the bonds of $Q_{2\varrho}(x_0)$ with $\mathsf g(b_-)U(b)\mathsf g(b_+)^{-1}=e^{i\xi A(b)}$ there and, on $Q_\varrho(x_0)$, $|A|\le80C^\sharp_{\mathrm G}B_3\varrho^2\varepsilon_n$ and $|\nabla^\xi A|\le80C^\sharp_{\mathrm G}B_3\varrho\,\varepsilon_n$; the plaquette variables and the current of $e^{i\xi A}$ are those of $U$ conjugated by $\mathsf g$, and the current has a modulus $<2B_3\varepsilon_n$.
\end{enumerate}
None of these numbers depends on $n$, $K$, $\ell$, the spacing, $x_0$ or the configuration.
\end{lemma}

\begin{proof}
(a) The $2$-form $f_{\mu\nu}=\xi^{-2}F_{p_{\mu\nu}}=\nabla^\xi_\mu A_\nu-\nabla^\xi_\nu A_\mu$ is antisymmetric and closed, $\mathrm df=0$, and its lattice codifferential differs from the current of $e^{i\xi A}$ by at most $40|A||\nabla^\xi A|+32|A|^3$ in the sup bounds, by the expansion of the lattice current; so $|f|\le2\mathsf r_\circ$ and $|\mathrm d^*f|\le2\mathsf r_\circ$. For an antisymmetric lattice $2$-form $f$ on a cube $Q_R(x_0)$, $R\ge8\xi$, with $|f|\le m_0$ and $|\mathrm df|,|\mathrm d^*f|\le m_1$, the lattice Hodge identity $\Delta(\chi f)=\mathrm d^*\mathrm d(\chi f)+\mathrm d\mathrm d^*(\chi f)$, for a cut-off $\chi$ equal to $1$ on $Q_{R/2}(x_0)$ and written as $\chi f=G_\xi*\Delta(\chi f)$ with the lattice Green function $G_\xi$ of $\xi\mathbb Z^3$, and the bounds $|\nabla G_\xi(z)|\le C(|z|+\xi)^{-2}$, $|\nabla^2G_\xi(z)|\le C(|z|+\xi)^{-3}$, which are uniform in $\xi$, give $|f(x)-f(x')|\le C(\beta_{\mathrm H})(m_0+Rm_1)(|x-x'|/R)^{\beta_{\mathrm H}}$ for $x,x'\in Q_{R/2}(x_0)$. At $R=1$, $m_0=m_1=2\mathsf r_\circ$ this is (a).
(b) Put $\alpha=4B_3\varepsilon_n$ and lift the sites of $Q_{4\varrho}(x_0)$ to the box $\Omega=x_0+\xi\{-M,\dots,M\}^3$, $M=\lfloor4\varrho/\xi\rfloor$, of side $s=(2M+1)\xi\le9\varrho$. The room hypothesis makes the lift injective. By Lemma~\ref{C-352-25}(d), $U$ is regular of level $n$ at radius $\alpha_0=2B_3\varepsilon_n$. The current-to-matrix-divergence calculation in Lemma~\YMMref{9.4}{40} gives $|D_U^{\mathsf T}\partial U|\le\xi^{-1}(\alpha_0\xi^2)^2+\xi^2\alpha_0<\alpha\xi^2$ at every full box bond; its plaquette bound is also $<\alpha\xi^2$. Choose $\epsilon_{\mathrm G}\le\min\{1,c_\square\}$ and $C^\sharp_{\mathrm G}\ge\max\{1,B_\square\}$, where $c_\square,B_\square$ are the box-gauge constants in the proof of Theorem~\YMMref{F}{199}. The stated smallness implies $s\alpha\le c_\square$.

That box estimate gives a real gauge and potential on $\Omega$ with $|A|<B_\square s\alpha\le36B_\square B_3\varrho\varepsilon_n$, and the same bound for first differences at base sites of depth at least $3/2$. Every positive bond with an endpoint in $Q_\varrho$ has base depth at least $3\varrho-2\xi>3/2$. Restricting the potential to $Q_{2\varrho}$ gives both displayed bounds, since $\varrho\ge1$. Plaquette and current moduli are preserved by the real gauge.
\end{proof}

In Lemmas~\ref{C-352-38}--\ref{C-352-41} the data are those of Theorem~\ref{C-352-7}; one progression holds the sources $w_1,w_2,w_T,w_3$, so that by multilinearity the objects of the two lists $(w_1,w_2,w_3)$ and $(w_T,w_3)$ are obtained from it. We write $\Delta_n=\mathsf cG_{n;12}-G_{n;T}$ and $\Gamma_n=\mathsf cG_{n;123}-G_{n;T3}$, functions of the configurations of the level $n$ whose pairs lie in the cores, and $\mathrm{nat}^{\mathcal W}_n=\varepsilon_n^2a_{K-n}^{-4}$.

\begin{lemma}[Pair minus tensor field on all small-field backgrounds]\label{C-352-38}
Let $r\le r_\alpha(g,\mathsf b)$; let $n_1$ and $n_0$ be the least levels with $a_{K-n_1}\ge(\frac r2)^{3/4}$ and $a_{K-n_0}\ge r^{1/2}$, $N_0=n_0+m_1$, where $m_1\ge1$ is the least integer with $L^{m_1/2}\ge C_7A_0/(0.85\,\mathsf c_{\mathfrak h}A_1)$, $C_7=2B_3+69\cdot1280\,C^\sharp_{\mathrm G}B_3$, an integer formed from the constants and $\beta_{\mathrm H}$ (through $C^\sharp_{\mathrm G}$ of the lemma on a regular gauge on a cube, which precedes this lemma), $\epsilon=r^{\beta_{\mathrm H}/36}+(a_K/r)^{1/9}$, and let $x_\sharp\ge e^{2\mathfrak p_0}$ be the least number such that every $x\ge x_\sharp$ has $17408B_3\varepsilon(x)\le\epsilon_{\mathrm G}$ and $1280C^\sharp_{\mathrm G}B_3\varepsilon(x)\le\frac1{100}$, with the numbers $\epsilon_{\mathrm G}$, $C^\sharp_{\mathrm G}$ of the lemma on a regular gauge on a cube.

Let $L^{N_0-n_1}\ge2M$, and let $\epsilon^{\mathrm e}_i=C_\omega[1.45\gamma_{i-1}+\frac12\mathsf c^\flat\theta^{[0]}]\le1$ for $\mathsf k(\frac r2)\le i<n_1$; this number bounds the single-source arrays of the near pair at the level $i$, $\mathfrak n[\dot{\mathfrak e}_{a,i}]\le2\epsilon^{\mathrm e}_i\varsigma_i\mathsf m_{i,a}$ ($a=1,2$, $\varsigma_i$ formed at the width $\frac r8$), as in Lemma~\ref{C-352-15}(a). For $N_0\le n^\sharp<K$ with $x_{n^\sharp}\ge x_\sharp$ there are real numbers $\mathsf c_n$, $N_0\le n\le n^\sharp$, such that for $\mathsf m_n$-almost every $W\in\mathcal V^\chi_n$
\[
|\Delta_n(U_n(W))-\mathsf c_n|\le\varepsilon^{\mathcal W}\,\mathrm{nat}^{\mathcal W}_n,\qquad|\mathsf c_n-\mathsf c_{n'}|\le2.01\,\tilde\varepsilon\,\mathrm{nat}^{\mathcal W}_n\quad(n\le n'\le n^\sharp),
\]
with $\varepsilon^{\mathcal W}\le\tilde\varepsilon\le C(g,\mathsf b)\,\epsilon$.
\end{lemma}

\begin{proof}
(i) We use the regular-field estimate of Lemma~\ref{lem:weighted-regular-family} and the insertion identity of Lemma~\ref{lem:weighted-step-integrals} on the small-field backgrounds. The first is the evaluation of a local term at a configuration that is regular of a level, with the factor $\varepsilon_n^2$: for $1\le i\le n\le K$, $W\in\mathcal V^\chi_n$ and a pointed family $\mathsf f$ of level $i$ that vanishes unless one of $\mathsf N$ given cubes of $\pi_i$ is a cube of the domain, and on the domains with $d_i(X)<D$,
\[
\sum_{(X,z)}\bigl|\mathsf f(X,U_n(W),z)-\mathsf f(X,\mathbf 1,z)\bigr|\le C\,\mathsf N\,\|\mathsf f\|_{i,\lambda_{\mathrm{out}}}\,\varepsilon_n^2\,L^{-4(n-i)}e^{-D/2} .
\]
It is \eqref{eq:weighted-regular-family} at the radius $\varepsilon=2B_3\varepsilon_n$ at which $U_n(W)$ is regular of the level $n$ (Lemma~\ref{C-352-25}(d), $4M\varepsilon\le c_{\mathrm A}\alpha_0$), which gives for one domain of size $d$ the factor $(\varepsilon/\alpha_0)^2L^{-4(n-i)}M^2(1+d)^2$, together with the tails of the proof of Lemma~\ref{C-352-14}(a): the domains beyond the depth $d_*=\min\{c_{\mathrm A}\alpha_0L^{2(n-i)}/M\varepsilon,(L^{m+K-i}-4)/2M\}-1$ have $e^{-[d_*]_+/2}\le C\varepsilon_n^2L^{-4(n-i)}$, by $e^{-s/2}\le8s^{-2}$, by $L^{m+K-n}\ge13My_n$ on the sides $\ell=2L^m$ of $\mathfrak L(g)$ (from $L^m\ge LMy_K$ and Lemma~\ref{C-352-10}(a)) and by $x_n^{-1}\le\varepsilon_n^2/A_0^2$; no ceiling beyond those of the data is read. (Lemma~\ref{C-352-14}(b) is the same count without the factor $\varepsilon_n^2$.) The second is the averaging identity of one step: the quotient $\mathsf{Qt}_{j+1}$ of one step, a positive normalised mean that takes bounded Borel functions of the block fields of the level $j$ to functions of those of the level $j+1$, takes the action $S[\mathfrak g](U_j(\cdot))$ of a bare array $\mathfrak g$ of the level $j$ to $(S[\mathfrak g]+\mathsf{Sum}[\mathsf D\mathsf{St}^{[j]}(\mathfrak f_{0,j})[\mathfrak g]])(U_{j+1}(\cdot))$, $\mathsf m_{j+1}$-almost everywhere; so the seminorm $\inf_c\operatorname{ess\,sup}_{\mathsf m_n}|\Phi(U_n(W))-c|$ of such actions does not grow from one level to the next. (ii) \emph{The split.} For $n_0\le j<K$ let $\mathfrak x_j$ be the bare array whose families are $\mathsf c$ times the kernel families $\mathsf K^{(i)}$ of the near pair at the levels $\mathsf k(\frac r2)<i\le n_1$ and whose slot holds minus the weight member of $w_T$; its action $X^{\mathrm{fz}}=S[\mathfrak x_j]$, the frozen part, does not depend on $j$.

Let $\mathfrak r_j$ hold all the other families of $\Delta$ of the levels $\le n_0$. Then $\Delta_n=\Delta_n(\mathbf 1)+X_n+S_n$, where $X_n$ adds to $X^{\mathrm{fz}}$ the families that the first derivatives of the steps above $n_0$ make from $\mathfrak x$ alone, and $S_n$ is the action of $\mathfrak r_n$ together with the families $\mathsf s^{(i)}$, $n_0<i\le n$, that the first derivatives of the steps make from $\mathfrak r$ and from the $\mathsf s^{(i')}$ of the lower levels, and with the second derivatives $\mathsf c\,\mathsf B_{i-1}(\mathfrak f_{\{1\},i-1},\mathfrak f_{\{2\},i-1})$ of the steps (Lemma~\ref{C-352-13}). The evaluation of (i), the lags of Lemma~\ref{C-352-12}(c) and the sizes of Lemma~\ref{C-352-15} give $|S_n|\le C\varepsilon_S\,\mathrm{nat}^{\mathcal W}_n$ at every small-field background of the level $n$, with $\varepsilon_S=2\varepsilon_{\mathfrak r}+C(r/a_{K-n_0})^3$ and the number $\varepsilon_{\mathfrak r}$ of step (i) of the proof of the lemma on the chain classes, which follows this lemma; that step uses the same split on the chain classes. (iii) The explicit parts cancel.

With $D(n)=3n\otimes n-1$, the pair, multiplied by $\mathsf c$, carries the explicit form with the kernel $\mathsf b^{-1}C^{\mathrm{free}}_K\approx\mathsf b^{-1}a_K^3\mathbb K_3(re_\lambda)=\mathsf b^{-1}a_K^3(4\pi r^3)^{-1}D(e_\lambda)$, while the tensor source carries $\frac1{2a_K}\sum_pw_T(p)|F_p|^2$ with $\delta_{\lambda(p)\lambda}-\frac13=\frac13D(e_\lambda)_{\lambda(p)\lambda(p)}$; the moment condition $\sum_k\mathsf a_kD(n_k)=\mathsf a_TD(e_\lambda)$, here with one pair of direction $e_\lambda$ and $\mathsf a_1=\mathsf a_T=1$, makes them equal. (iv) \emph{The frozen part on the backgrounds.} Let $N_0\le n\le n^\sharp$, $W\in\mathcal V^\chi_n$, $U=U_n(W)$, $x_0$ a bare site with $|x_0-c|_\infty\le a_K$, and let lengths be in level-$n$ units, $\xi=\eta_n$.

\emph{The gauge.} As $x_n\ge x_\sharp$ and $34a_{K-n}<\ell$, clause (b) of the lemma on a regular gauge on a cube, which precedes this lemma, applies at $\varrho=4$ and gives $\mathsf g$ on $Q_{16}(x_0)$ and a potential $A$ on the bonds of $Q_8(x_0)$ with $\mathsf g(b_-)U(b)\mathsf g(b_+)^{-1}=e^{i\xi A(b)}$, and on $Q_4(x_0)$ the bounds $|A|,|\nabla^\xi A|\le C_6\varepsilon_n\le\frac1{100}$, $C_6=1280C^\sharp_{\mathrm G}B_3$, and a current of modulus $<2B_3\varepsilon_n$. \emph{The cut-off.} Let $h=\prod_\mu\chi(|x_\mu-x_{0,\mu}|)$ with $\chi=1$ on $[0,3]$, $\chi=0$ on $[\frac{15}4,\infty)$, $|\chi'|\le2$ and $\operatorname{Lip}\chi'\le11$; put $\hat A=hA$ on the bonds of $Q_4(x_0)$, $\hat A=0$ on all other bonds of $T_\ell$, and $\hat U=e^{i\xi\hat A}$.

Then $\hat U$ is the transform of $U$ by $\mathsf g$ on every bond with an end in $Q_3(x_0)$, and the product rules for $hA$ with the expansion of the lattice current give
\[
\max\bigl\{\sup|\hat A|,\ \sup|\nabla^\xi\hat A|,\ \sup|J(\hat U)|\bigr\}\le(2B_3+69\,C_6)\,\varepsilon_n=C_7\varepsilon_n\qquad\text{on all of }T_\ell .
\]
\emph{The shift of the level.} Seen from the level $n^\flat=n-m_1\ge n_0$, the same configuration has the potential $\hat A^\flat=L^{-m_1}\hat A$, with $|\hat A^\flat|\le L^{-m_1}C_7\varepsilon_n$, differences $\le L^{-2m_1}C_7\varepsilon_n$ and current $\le L^{-3m_1}C_7\varepsilon_n$ in level-$n^\flat$ units; with $\varepsilon_n=(A_0/A_1)\delta_n$, $\delta_n\le(\frac54)^{1/2}L^{m_1/2}\delta_{n^\flat}$ (Lemma~\ref{C-352-10}(a)) and the choice of $m_1$,
\[
L^{-m_1}C_7\varepsilon_n\le(\tfrac54)^{1/2}\cdot0.85\,\mathsf c_{\mathfrak h}\delta_{n^\flat}<\tfrac{11}{10}\mathsf c_{\mathfrak h}\delta_{n^\flat}\le r_{n^\flat} .
\]
This is the norm clause of the chain class of radius $r_{n^\flat}$, and it is all that the steps below read of the field. \emph{The pair.} Step (ii) of the proof of the lemma on the chain classes is read at the level $n^\flat$ with $U$ in the place of $\tilde U$, with two changes. (1) The gauge $u=e^{iX}$ of Lemma~\ref{C-352-18} at the pair is formed for $\hat A^\flat$ on the bonds within two level-$n^\flat$ units of $x_0$, where $\hat U$ is the transform of $U$; it gives the angles $F'$, H\"older continuous as in that step, through the estimate of the proof of clause (a) of the lemma on a regular gauge on a cube on $Q_1(x_0)$ at $m_0=2.08\,r_{n^\flat}$ and $m_1=3.7\,r_{n^\flat}$, and the plaquette term of the weight member of $w_T$ in this gauge, as $a_p(U)=a_p(\hat U)$ at the plaquettes of $\operatorname{supp}w_T$, which lie in $Q_{1/2}(x_0)$. (2) For the kernel families of the levels $\mathsf k(\frac r2)\le i<n_1$, $m=n^\flat-i-1$: Lemma~\ref{C-352-18} for the real pair of the level $i+1$ that $\hat A^\flat$ and its current give on all of $T_\ell$ (the potential multiplied by $L^{-m}$, the current by $L^{-3m}$), on the box $N_{\rho'}$ about $x_0$, $\rho'=\min\{\mathsf h_i,3L^{n-i-1}-1\}$, with the same $u$, gives a real point $v'$ of the chart with $\|v'\|\le\mathsf h_iL^{-2m}r_{n^\flat}(6+19r_{n^\flat})\le\delta_i/2\Lambda_i$, by \textup{($\alpha$2)} at $n^\flat\ge n_0$, the pair having $|L^{-m}\hat A^\flat|\le L^{-m}r_{n^\flat}$ and differences and current $\le L^{-2m}r_{n^\flat}$.

A domain $X$ of the level $i+1$ with all its sites within $\rho'-1$ of $x_0$ lies in $Q_3(x_0)$, where $\hat U$ is the transform of $U$, so that $\mathsf K^{(i+1)}(X,U,z)=\mathsf K^{(i+1)}(X,P[v']|_X,z)$ by the gauge invariance and the locality of a local term; the other domains have $d_{i+1}(X)\ge2L^{n-i-1}/M-2$, or the depth of step~5 of the proof of Proposition~\ref{C-352-19} if $\rho'=\mathsf h_i$, and are paid by the norm (Lemma~\ref{C-352-14}(b) at $U$ and at $P[v']$). Steps 2--4 of the proof of Proposition~\ref{C-352-19} (holomorphy, the real ball, two constants), which read of the field only the real point $v'$ and its norm, give at $v'$ the kernel step with the angles $F'$. As $L^{n-i-1}\ge L^{N_0-n_1}\ge2M$ and $\epsilon^{\mathrm e}_i\le1$, the domains that leave the box cost at most $C\mathsf b\,r\,\mathrm{nat}^{\mathcal W}_n$ together.

With $r_{n^\flat}^2a_{K-n^\flat}^{-4}\le C\,\mathrm{nat}^{\mathcal W}_n$ (Lemma~\ref{C-352-10}(a), $a_{K-n^\flat}=L^{-m_1}a_{K-n}$) this bounds the frozen part: $|X^{\mathrm{fz}}(U)|\le\varepsilon^{\mathcal W}_X\mathrm{nat}^{\mathcal W}_n$, with $\varepsilon^{\mathcal W}_X\le C[\varepsilon_X+\mathsf b\,r]$ and the number $\varepsilon_X$ of step (ii) of the proof of the lemma on the chain classes. (v) Carried upward with (i) and the lags of Lemma~\ref{C-352-12}(c), $\|X_n\|\le[\varepsilon^{\mathcal W}_X+C\gamma_{N_0-1}+C\gamma_n^{8/9}(\varepsilon^{\mathcal W}_X)^{1/9}]\mathrm{nat}^{\mathcal W}_n$; the exponents $\frac89,\frac19$, and $\frac12$ below, come from interpolations between the oscillation bound and the bound of the sizes, and the choice of the levels, $\mathsf k(\frac r2)<n_1<n_0$ with $a_{K-n_1}\ge(\frac r2)^{3/4}$, costs the exponent $\frac14$.

Together $\varepsilon^{\mathcal W}\le\tilde\varepsilon\le C(g,\mathsf b)[r^{\beta_{\mathrm H}/36}+(a_K/r)^{1/9}]$ for $r\le r_\alpha(g,\mathsf b)$, the levels being ordered by \textup{($\alpha$1)}--\textup{($\alpha$5)} of Definition~\ref{C-352-2}(5); the constants $\mathsf c_n$ of two levels differ by at most these oscillations plus the part of the steps between that the sizes of the sources at those levels bound, which is of the same order.
\end{proof}

\begin{lemma}[Pair minus tensor field on the chain classes]\label{C-352-38a}
Let $r\le r_\alpha(g,\mathsf b)$, let $n_1$ and $n_0$ be the levels of Lemma~\ref{C-352-38}, and $\epsilon=r^{\beta_{\mathrm H}/36}+(a_K/r)^{1/9}$. There are numbers $\mathsf c'_n$, $n_0\le n\le K$, and a number $\varepsilon_\alpha\le C(g,\mathsf b)\,\epsilon$ such that
\[
\sup_{A\in\mathsf{Cc}_n(r_n)}\bigl|\Delta_n(e^{i\eta_nA})-\mathsf c'_n\bigr|\le\varepsilon_\alpha\,r_n^2a_{K-n}^{-4}\qquad(n_0\le n\le K).
\]
\end{lemma}

\begin{proof}
(i) \emph{The split} is that of step (ii) of the proof of Lemma~\ref{C-352-38}. For $n_0\le j<K$, $\mathfrak x_j$ is the array whose families are $\mathsf c$ times the kernel families $\mathsf K^{(i)}$ of the near pair at the levels $\mathsf k(\frac r2)<i\le n_1$ and whose slot holds minus the weight member of $w_T$; its action $X^{\mathrm{fz}}$, the frozen part, does not depend on $j$. $\mathfrak r_j$ holds all the other families of $\Delta$ of the levels $\le n_0$. Lemma~\ref{C-352-12} and the sizes of Lemma~\ref{C-352-15} give $\mathfrak n[\mathfrak x_j]\le Ca_{K-j}^{-4}$ and $\mathfrak n[\mathfrak r_j]\le\varepsilon_{\mathfrak r}a_{K-j}^{-4}$, where $\varepsilon_{\mathfrak r}$ is a constant times $\gamma_{n_0}+s(\frac r2)\theta^{[0]}+(r/a_{K-n_1})^3$ and terms exponentially small in $\hat r_{\mathsf k(r/2)}$ and in $r/a_{K-1}$, so that $\varepsilon_{\mathfrak r}\le C(g,\mathsf b)[r^{1/4}+(a_K/r)^{1/4}]$.

On the chain classes $\Delta_n=\Delta_n(\mathbf 1)+X_n+S_n$, where $X_n$ adds to $X^{\mathrm{fz}}$ the families that the first derivatives of the steps above $n_0$ make from $\mathfrak x$ alone, and $S_n$ collects the rest: the action of $\mathfrak r_n$, the families that the first derivatives make from $\mathfrak r$ and from them, and the second derivatives of the steps; Lemma~\ref{C-352-14}(a), in the frame at $\mathsf k(\frac r2)$ in which Lemma~\ref{C-352-11}(a) puts the near pair, a test pair of the scale $\frac r2$, under the conditions $3\le s(\frac r2)\le K$ and $\mathsf{Lm}_{s(r/2)}\le L^{s(r/2)-3}$ of the data (and $\mathsf k(\frac r2)<n_1<n_0\le n<K$ by \textup{($\alpha$1)}), and the lags of Lemma~\ref{C-352-12}(c) give $|S_n|\le C\varepsilon_Sr_n^2a_{K-n}^{-4}$ with $\varepsilon_S=2\varepsilon_{\mathfrak r}+C(r/a_{K-n_0})^3$.
(ii) \emph{The frozen part.} Let $n_0\le n\le K$, $A\in\mathsf{Cc}_n(r_n)$, $x_0$ a bare site with $|x_0-c|_\infty\le a_K$, let $u=e^{iX}$, $X$ linear, be the gauge of Lemma~\ref{C-352-18} at the pair, about $x_0$, $A'$ the potential of $\tilde U=e^{i\eta_nA}$ in this gauge and $F'=\eta_n\mathrm dA'$ its angles. Near the pair $|A'|$, $|\nabla A'|$ and the current are at most a constant times $r_n$ (Lemma~\ref{C-352-18}, at $\alpha=\gamma_\circ=r_n$, in the units of the level $n$): on $N_2$, $|A'|\le6r_n+18.7r_n^2\le6.2\,r_n$ and $|\nabla^{\eta_n}A'|\le\frac{100}{99}(r_n+3r_n^2)\le1.04\,r_n$, and the current, conjugated by the gauge, is at most $r_n$. So the $2$-form $f=\eta_n^{-2}F'$ is closed, $|f|\le2.08\,r_n$, and, as its codifferential differs from the current by at most $40|A'||\nabla^{\eta_n}A'|+32|A'|^3$, $|\mathrm d^*f|\le3.7\,r_n$ on $Q_1(x_0)$, since $2r_n\le\varrho_{\mathrm{ch}}$ (Lemma~\ref{C-352-10}(b)) and $\varrho_{\mathrm{ch}}\le\frac1{60}$ by equation~\YMMref{(2285)}{833}.

The lemma on a regular gauge on a cube, which precedes Lemma~\ref{C-352-38}, is read through the estimate of the proof of its clause (a) on $Q_1(x_0)$, at $m_0=2.08\,r_n$ and $m_1=3.7\,r_n$: it makes $F'$ H\"older continuous with the exponent $\beta_{\mathrm H}$ on $Q_{1/2}(x_0)$, on the scale of the pair. Proposition~\ref{C-352-19}, in the frame of step (i), at the levels $\mathsf k(\frac r2)\le i<n_1$, where \textup{($\alpha$2)} of Definition~\ref{C-352-2}(5) gives its condition in (ii), taken before the last replacement of step 5 of its proof, so that $\mathcal Q_i$ stands at $F'$ and the last term of (ii) does not occur, replaces $\mathsf c\sum_i(\mathsf{Sum}[\mathsf K^{(i+1)}](\tilde U)-\mathsf{Sum}[\mathsf K^{(i+1)}](\mathbf 1))$ by $\mathsf c\,q_{n_1}(F')$ with an error at most $C\sum_i\hat r_i^3\mathfrak b_i^2(\delta_i+\epsilon^{\mathrm{cu}}_i)\Lambda_i^2\mathsf h_i^2\,r_n^2a_{K-n}^{-4}\le C(g,\mathsf b)\,r^{1/4}r_n^2a_{K-n}^{-4}$; the weight member of $w_T$ is expanded in the gauge $u$ at the pair.

With $\xi=\eta_n$, the plaquette angles of $\tilde U$ in this gauge are $F'_p+\Psi_p$ with the commutator part $|\Psi_p|\le24\xi^2|A'|^2$; Lemma~\ref{C-352-18} at $\alpha=\gamma_\circ=r_n$ gives $|A'(b)|\le3sr_n+(9s+8\xi)r_n^2\le4sr_n+8\xi r_n^2$ at the bonds with base point within $s$ of $x_0$, and $|F'_p|\le2.08\,\xi^2r_n$; as $a_p=\frac12|\theta|^2+O(|\theta|^4)$ in the plaquette angle $\theta$,
\[
\bigl|a_p(\tilde U)-\tfrac12|F'_p|^2\bigr|\le|F'_p||\Psi_p|+\tfrac12|\Psi_p|^2+C|F'_p+\Psi_p|^4\le C\xi^4r_n^3(s+\xi)^2
\]
at the plaquettes within $s\le\frac12$ of $x_0$. The plaquettes of $\operatorname{supp}w_T$ lie within $s\le r/a_{K-n}$ of $x_0$, so $s+\xi\le2r/a_{K-n_0}$; summed against $w_T$, with $\sum_p|w_T(p)|\le Ca_K^{-3}$, this replaces the tensor source by $\frac1{2a_K}\sum_pw_T(p)|F'_p|^2$ up to $C(r/a_{K-n_0})^2r_n^2a_{K-n}^{-4}\le C\,r\,r_n^2a_{K-n}^{-4}$, as $a_{K-n_0}\ge r^{1/2}$: in the gauge at the pair the angles are small at the scale $s+\xi$ of the tensor source. By the H\"older continuity and the comparison of $\mathcal S_{n_1}$ with $\mathsf b^{-1}C^{\mathrm{free}}_K$ and of $C^{\mathrm{free}}_K$ with $a_K^3\mathbb K_3$ (Lemma~\ref{C-352-22}(b), with \textup{($\alpha$5)}), both main terms equal $\frac1{6a_K^4}\sum_{\mu,\nu}D(e_\lambda)_{\mu\nu}\langle F^\circ_\mu,F^\circ_\nu\rangle$, $F^\circ_\mu$ the angle of the orientation $\mu$ at the centre, up to $C[(r/a_{K-n_1})^{\beta_{\mathrm H}}+a_K/r]\,r_n^2a_{K-n}^{-4}$ and terms of the same order; the moment condition of step (iii) of the proof of Lemma~\ref{C-352-38} cancels the main terms. So $\sup_{\mathsf{Cc}_n(r_n)}|X^{\mathrm{fz}}|\le\varepsilon_Xr_n^2a_{K-n}^{-4}$ with $\varepsilon_X\le C(g,\mathsf b)[r^{\beta_{\mathrm H}/4}+a_K/r]$.
(iii) \emph{Carrying upward.} On the chain classes the averaging identity of one step holds at every point: for $n_0\le m<K$, $A'\in\mathsf{Cc}_{m+1}(r_{m+1})$, the measure $\nu=\nu_{m,A'}$ of the step and every bare array $\mathfrak g$ of the level $m$ of size below $N^{[m]}_\star$,
\[
\int S[\mathfrak g]\bigl(e^{i\eta_mA^{\mathrm{lw}}(B)}\bigr)\,d\nu(B)=\bigl(S[\mathfrak g]+\mathsf{Sum}[\mathsf D\mathsf{St}^{[m]}(\mathfrak f_{0,m})[\mathfrak g]]\bigr)\bigl(e^{i\eta_{m+1}A'}\bigr)+c_{\mathfrak g},
\]
with a number $c_{\mathfrak g}$ that does not depend on $A'$. It is the first derivative, in the direction of $\mathfrak g$, of the exact step integral in Lemma~\ref{lem:weighted-step-integrals} at the real chart point $A'$, which writes the action of the level $m+1$ through the mean under $\nu_{m,A'}$ of the exponential of the action of the level $m$ at the lowered configurations; the measure $\nu_{m,A'}$ and the lowering map name no array, and the step takes any bare array of that size in its fourth place (Lemma~\ref{C-352-12}(b)). As $A^{\mathrm{lw}}(B)\in\mathsf{Cc}_m(r_m)$ and $\nu$ is a probability measure, the oscillation of the right member over $\mathsf{Cc}_{m+1}(r_{m+1})$ is at most the oscillation of $S[\mathfrak g]$ over $\mathsf{Cc}_m(r_m)$.

A family that the steps above $n_0$ make from $\mathfrak x$, first created at the level $m+1$, is bounded in two ways: by the norm of the calculus, with the factor $\gamma_m$ (Lemma~\ref{C-352-12}(b), (c)), and, through this identity at $\mathfrak g=\mathfrak x_m$, by the oscillation of $X^{\mathrm{fz}}$ at the levels $m$ and $m+1$, at most $2\varepsilon_XL^{4(n-m)}r_n^2a_{K-n}^{-4}$. The smaller of the two, summed over $m$ with $\sum_{l\ge0}\min\{aL^{-l/2},bL^{4l}\}\le2.4a^{8/9}b^{1/9}$ ($L\ge13$), bounds the oscillation of $X_n$ by $[\varepsilon_X+C\gamma_{K-1}^{8/9}\varepsilon_X^{1/9}]r_n^2a_{K-n}^{-4}$. Together the oscillation of $\Delta_n$ on $\mathsf{Cc}_n(r_n)$ is at most $\varepsilon_\alpha r_n^2a_{K-n}^{-4}$ with $\varepsilon_\alpha=\varepsilon_X+C\gamma_{K-1}^{8/9}\varepsilon_X^{1/9}+C\varepsilon_S\le C(g,\mathsf b)\,\epsilon$, as $\gamma_{K-1}\le1$; $\mathsf c'_n$ is the centre of this oscillation.

At $n=K$ the chain class is $\{0\}$ and $r_K=0$.
\end{proof}

\begin{lemma}[Pair minus tensor field up to the last level]\label{C-352-38b}
Let $g$ satisfy besides $(\Gamma3)$ of Definition~\ref{C-352-2}(6), let $r\le r^{\mathcal W}_\alpha(g,\mathsf b)$, let $n_0$ be the level of Lemma~\ref{C-352-38}, and $\varepsilon^{\mathcal W}=C(g,\mathsf b)[r^{\beta_{\mathrm H}/72}+(a_K/r)^{1/18}]$ with a suitable number $C(g,\mathsf b)\ge1$. There are numbers $\mathsf c_n$, $n_0\le n\le K$, such that for $\mathsf m_n$-almost every $W\in\mathcal V^\chi_n$
\[
|\Delta_n(U_n(W))-\mathsf c_n|\le\varepsilon^{\mathcal W}\,\mathrm{nat}^{\mathcal W}_n,\qquad|\mathsf c_n-\mathsf c_K|\le\varepsilon^{\mathcal W}\,\mathrm{nat}^{\mathcal W}_n,
\]
and, for $w_3\not\equiv0$ and $\mathsf d'=|c-z|$, $|\Gamma_n(U_n(W))|+|\Gamma_K(\mathbf 1)|\le C(\mathsf b)\,\mathsf d'^{-5/2}\varepsilon^{\mathcal W}\mathsf b^{-2}\mathsf d'^{-6}\mu_3$.
\end{lemma}

\begin{proof}
The proof of Lemma~\ref{C-352-38} is made at every level $n_0\le n\le K$, with the radius $C_{\mathcal W}\varepsilon_n$ of the backgrounds in the place of the radius of a chain class and without the shift of the level. (i) By the first two members of $(\Gamma3)$ the domains that leave the cores at the radius $2B_3\varepsilon_n$ have a depth that makes the tails of Lemma~\ref{C-352-14}(b) at most $C\varepsilon_n^2L^{-4(n-i)}$; so a pointed family of level $i$ that vanishes off a set of sites changes on the backgrounds of the level $n$ by at most a constant times its norm, the number of the cubes that hold a site of the set, $L^{-4(n-i)}$ and $\mathrm{nat}^{\mathcal W}_n$. (ii) Lemma~\ref{C-352-37a}(b) at $\varrho=3$, which applies as $7488B_3\varepsilon_n\le\epsilon_{\mathrm G}$ by $(\Gamma3)$ and $26a_{K-n}<\ell_g$ (as $\ell_g\ge2LMy_K\ge338$), gives about a bare site $x_0$ next to $c$ a gauge $\mathsf g$ on $Q_{12}(x_0)$ and a potential $A$ on the bonds of $Q_6(x_0)$ with $\mathsf g(b_-)U(b)\mathsf g(b_+)^{-1}=e^{i\eta_nA(b)}$ there, $U=U_n(W)$; on $Q_3(x_0)$ the potential, its differences and the current are at most $\mathsf r_\circ:=C_{\mathcal W}\varepsilon_n\le\frac1{120}$.

Lemma~\ref{C-352-18} is used in a local form, with data on $N_{\rho+1}$: \emph{in the units of a level, with the bare spacing $\xi\le\frac1{13}$, let $\rho\ge2$ with $4\rho+4$ at most the side of the torus, and let $\mathsf g\colon N_{\rho+1}\to SU(2)$ and a real $1$-form $A$ on the bonds with both end points in $N_{\rho+1}$ satisfy $\mathsf g(b_-)U(b)\mathsf g(b_+)^{-1}=e^{i\xi A(b)}$, $|A|\le\alpha\le\frac1{60}$ and $|\nabla^\xi A|\le\gamma_\circ$ there, the current of $U$ being at most $\iota$ there; with $c_\mu=A_\mu(x_0)$, $X(x)=\sum_\mu(x-x_0)_\mu c_\mu$ and $h=e^{iX}\mathsf g$ on $N_{\rho+1}$ there is a real pair $v'$ on all the bare bonds which, at the bonds with base point in $N_\rho$, is the image of $U$ and of its current under the gauge transformation $h$, and which obeys the three bounds of Lemma~\ref{C-352-18}; if $\|v'\|$ is less than the radius of the ball of the chart, every local term of the level on a domain $X\subset N_{\rho-1}$ that is formed at $U$ has the same value at $P[v']|_X$.} Its proof is that of Lemma~\ref{C-352-18}, with the transformed potential formed on $N_{\rho+1}$ and the cut-off $\chi(x)=\min\{1,\max\{0,(\rho+1-2\xi-|x-x_0|_\infty)/(1-2\xi)\}\}$, of the width $1-2\xi$ and equal to $1$ on $N_\rho$, in the place of the cut-off of the width $\rho$: as $|\chi(x+\xi e_\nu)-\chi(x)|\le\frac{13}{11}\xi$, $|\nabla^\xi A'|\le\frac{100}{99}(\gamma_\circ+3\alpha^2)+\frac{13}{11}[3(\rho+1)\gamma_\circ+(9\rho+9.62)\alpha^2]\le\rho(6\gamma_\circ+19\alpha^2)$ for $\rho\ge2$; and the bonds that a local term on $X\subset N_{\rho-1}$ reads have their end points in $N_\rho$, where $P[v']$ is the image of $U$ under $h$, while a local term is a gauge invariant function of the pair on its domain.

Viewed from the level $i+1$, $\mathsf k(\frac r2)\le i<n_1$, $m=n-i-1$, the data on $Q_3(x_0)$ have $|A|\le\alpha=L^{-m}\mathsf r_\circ$, $|\nabla A|\le\gamma_\circ=L^{-2m}\mathsf r_\circ$ and a current at most $L^{-3m}\mathsf r_\circ\le\gamma_\circ$, and the cube $N_{\mathsf h_i+1}$ of the level $i+1$ lies in $Q_3(x_0)$, as $(\mathsf h_i+1)L^{-m}\le3$ (Definition~\ref{C-352-2}(6)). The local form at $\rho=\mathsf h_i$, where $(4\mathsf h_i+4)a_{K-i-1}\le\ell_g$, gives a real point $v'$ of the chart with $\|v'\|\le\mathsf h_i\gamma_\circ(6+19\mathsf r_\circ)\le7\mathsf h_iL^{-2m}C_{\mathcal W}\varepsilon_n\le\delta_i/2\Lambda_i$ (Definition~\ref{C-352-2}(6)), inside the ball of the chart; the linear gauge $X$ is the same function at all the levels $i$, so that one $h$ serves them all.

Steps 2--4 of the proof of Proposition~\ref{C-352-19} at the real point $v'$, the domains in $N_{\mathsf h_i-1}$ taking the same values at $U$ and at $P[v']$, and the others, of a linear size at least $D_*=(\mathsf h_i-1-R_w)/M-2$, costing $2\mathsf M^{\mathsf K}_ie^{-[D_*]_+}\le\bar{\mathsf e}\gamma_\circ^2/\delta_i^2$ by the choice of $\mathsf h_i$, replace $\mathsf c\sum_i(\mathsf{Sum}[\mathsf K^{(i+1)}](U)-\mathsf{Sum}[\mathsf K^{(i+1)}](\mathbf 1))$ by $\mathsf c\,q_{n_1}(F')$, $F'$ the linear angles in the gauge $h$, up to $C\sum_i\hat r_i^3\mathfrak b_i^2(\delta_i+\epsilon^{\mathrm{cu}}_i)\Lambda_i^2\mathsf h_i^2\,C_{\mathcal W}^2\mathrm{nat}^{\mathcal W}_n$. In the gauge $h$, in the units of the level $n$, $|A'|\le4s\mathsf r_\circ+8\eta_n\mathsf r_\circ^2$ at the bonds with base point within $s\le2$ of $x_0$, $|\nabla A'|\le1.04\,\mathsf r_\circ$ and $|F'_p|\le2.08\,\eta_n^2\mathsf r_\circ$ there, and $|a_p(U)-\frac12|F'_p|^2|\le C\eta_n^4\mathsf r_\circ^3(s+\eta_n)^2$ at the plaquettes within $s\le\frac12$ of $x_0$; this replaces the tensor source by $\frac1{2a_K}\sum_pw_T(p)|F'_p|^2$ up to $C(r/a_{K-n_0})^2C_{\mathcal W}^2\mathrm{nat}^{\mathcal W}_n$.

Lemma~\ref{C-352-37a}(a) is read through the estimate of its proof on $Q_1(x_0)$, at $m_0=2.08\,\mathsf r_\circ$ and $m_1=3.7\,\mathsf r_\circ$ (on $N_2$, $|A'|\le6.2\,\mathsf r_\circ$, and the codifferential of $\eta_n^{-2}F'$ differs from the current by at most $40|A'||\nabla A'|+32|A'|^3$): it makes $F'$ H\"older continuous with the exponent $\beta_{\mathrm H}$ on the scale of the pair. The comparison of $\mathcal S_{n_1}$ with $\mathsf b^{-1}C^{\mathrm{free}}_K$ and of $C^{\mathrm{free}}_K$ with $a_K^3\mathbb K_3$ (Lemma~\ref{C-352-22}(b), with \textup{($\alpha$5)}) brings both main terms to $\frac1{6a_K^4}\sum_{\mu\nu}D(e_\lambda)_{\mu\nu}\langle F^\circ_\mu,F^\circ_\nu\rangle$, $F^\circ$ the angles at the centre, up to $C[(r/a_{K-n_1})^{\beta_{\mathrm H}}+a_K/r]C_{\mathcal W}^2\mathrm{nat}^{\mathcal W}_n$; the moment condition cancels them, and the frozen part on the backgrounds is at most $\varepsilon^{\mathcal W}_X\mathrm{nat}^{\mathcal W}_n$. (iii) From one level to the next the averaging identity of one step, written for arbitrary bare arrays, gives $\mathsf m_{n+1}$-almost everywhere $\Delta_{n+1}=\mathsf{Qt}_{n+1}[\Delta_n]+\mathsf b_n$, where $\mathsf{Qt}_{n+1}$ is the quotient of the step, a positive normalised mean, and $\mathsf b_n$ is $\mathsf c$ times the covariance, under the step, of the quotients of the two single sources of the near pair.

The one-source bound $|G_{n;a}-G_{n;a}(\mathbf 1)|\le Cr^3|w_a|_\infty\mathrm{nat}^{\mathcal W}_n$ on the backgrounds gives $|\mathsf b_n|\le CA_0^2y_{n_0}^{2\mathfrak p_0}(r/a_{K-n_0})^3\mathrm{nat}^{\mathcal W}_n\le C(\mathsf b)\,r\,\mathrm{nat}^{\mathcal W}_n$. A quotient does not increase the oscillation, so the centres $\mathsf c_n$ of the oscillations of two consecutive levels differ by at most twice the oscillation bound plus $C(\mathsf b)\,r\,\mathrm{nat}^{\mathcal W}_n$; as $\mathrm{nat}^{\mathcal W}$ falls by the factor $\frac54L^{-3}$ from one level to the next, the sum over the levels gives the second bound. (iv) For $\Gamma_n$ the averaging identity of one step, in its form for two and for three arrays, gives $\mathsf m_{i+1}$-almost everywhere, for $n_0\le i<K$, $\Gamma_{i+1}=\mathsf{Qt}_{i+1}[\Gamma_i]+\mathsf{Qt}_{i+1}[\Delta_i;G_{i;3}]+\mathsf c\,\mathsf{Sum}[\mathsf C^{(i+1)}]$; here $\mathsf{Qt}_{i+1}[\,\cdot\,;\cdot\,]$ is the covariance of the quotients of the step, and the family $\mathsf C^{(i+1)}$ is made of the mixed second derivatives of the step in the arrays of the sets $\{1,3\}$ and $\{2\}$ and in those of $\{2,3\}$ and $\{1\}$, and of its third derivative in the arrays of the three single sources.

A quotient does not raise the essential supremum, so the bound of $\Gamma_n$ is that of $\Gamma_{n_0}$ plus the sum over $n_0\le i<n$ of the bounds of the two other terms. The covariance is at most the smaller of the bound by the oscillations of $\Delta_i$ and of the far source on the backgrounds and the bound by the norms, $\mathsf C^{(i+1)}$ is bounded by the norms, and $\Gamma_{n_0}$ by the first display of Lemma~\ref{C-352-39} at $U_{n_0}(W)$, whose pairs of all levels $\le n_0$ lie in the cores (Lemma~\ref{C-352-25}(d)); the sum over the levels is that of the proof of Lemma~\ref{C-352-40}, with the bound of the oscillation of $\Delta$ on the backgrounds in the place of $\varepsilon_\alpha$.

All the numbers so formed are at most $C(g,\mathsf b)[r^{\beta_{\mathrm H}/36}+(a_K/r)^{1/9}]$ for $r\le r^{\mathcal W}_\alpha$; with $(\frac{\mathsf b}{\mathsf d'})^{3/2}(\frac r{\mathsf d'})^3\le\mathsf b^{3/2}\mathsf d'^{-2}r^{1/2}$ and $(\frac{\mathsf b}{\rho_0})^{3/2}e^{-\mathsf d'/8M\rho_0}\le C\mathsf b^{3/2}\mathsf d'^{-5/2}r^{1/2}$, $\rho_0=a_{K-n_0}$, the claims follow.
\end{proof}

\begin{lemma}[Crossing functionals at levels below the distance]\label{C-352-39}
For $\beta\in\{13,23,123\}$ of the list $(w_1,w_2,w_3)$ and $\beta=T3$ of the list $(w_T,w_3)$, $m=|\beta|$, $1\le i\le K$, and every configuration $U$ whose pairs of the level $i$ lie in the cores,
\[
\sum_{(X,z)}|E^{(i)}_\beta(X,U,z)|\le C_\times\bigl(\tfrac{8r}{3\mathsf d}\bigr)^6\bar\nu_{i-1;\beta}\,e^{-D^\sharp_i/2},\qquad D^\sharp_i=[\mathsf d/2Ma_{K-i}-2]_+ .
\]
So if $8Ma_{K-n}\le\mathsf d$, then for every such $U$ of the level $n$: 
\[
|G_{n;\beta}(U)|\le C'_\times\mathsf b^{-m/2}a_{K-n}^{-m/2}e^{-\mathsf d/4Ma_{K-n}}(\frac{8r}{3\mathsf d})^6\prod_{a\in\beta}|w_a|_\infty.
\]

\end{lemma}

\begin{proof}
By the support rule (Lemma~\ref{C-352-12}(d)) a crossing family lives on domains that reach from the pair to the piece, of linear size at least $D^\sharp_i$; the near supports lie in $\mathbb B(c,\frac58r)$, the piece in $\mathbb B(z,\frac14r)$, so the corners of the two kinds are at least $\frac12\mathsf d$ apart in the sup distance; its norm is bounded by the sizes of Lemma~\ref{C-352-15}(e), in the form (f), at the width $r$ and a crossing distance $\mathsf d_\times\ge\frac38\mathsf d$, with $(r/\mathsf d_\times)^6\le(\frac{8r}{3\mathsf d})^6$; Lemma~\ref{C-352-14}(b) at $U$ gives the sum with the factor $e^{-D^\sharp_i/2}$, and $\bar\nu_{i-1;\beta}\le\bar c_m(5L/4\mathsf b)^{m/2}a_{K-i}^{-m/2}\prod|w_a|_\infty$; the sum over $i\le n$ is led by $i=n$.
\end{proof}

\begin{lemma}[The difference of the jets]\label{C-352-40}
Let $w_3\not\equiv0$. Then the centre $z$ of its cube has $(1-\frac1{64})\mathsf d\le|c-z|\le6\mathsf d$, and
\[
\bigl|\mathsf c\,\mathsf J_3(w_1,w_2,w_3)-\mathsf J_2(w_T,w_3)\bigr|\le C(g,\mathsf b)\,\mathsf d^{-12}\bigl[r^{\beta_{\mathrm H}/72}+(a_K/r)^{1/18}\bigr]\mathsf b^{-2}\mathsf d^{-6}\mu_3 .
\]
\end{lemma}

\begin{proof}
A barycentre in the cube lies in $\operatorname{supp}\varphi_2$, at a distance $\ge\mathsf d$ from $c$ and $\le\frac{\sqrt3}8r\le\frac1{64}\mathsf d$ from $z$, as $r\le(\mathsf d/5LM)^2$. Put $\mathsf d'=|c-z|$ and $\rho_0=a_{K-n_0}$; the data give $16r\le\mathsf d'\le\min\{1,\mathsf b\}$, $4M\rho_0\le\mathsf d'$ and $64\mathsf d'\le\ell_g$ ($\mathsf d'\le6\mathsf d$, $r\le(\mathsf d/5LM)^2$, $\rho_0<Lr^{1/2}$, $\ell_g\ge2LMy_K\ge338$). The corners of the plaquettes of the near supports and of the piece are at least $\frac12\mathsf d'$ apart in the sup distance, so Lemma~\ref{C-352-39} holds with $\mathsf d'$ in the place of $\mathsf d$. Lemma~\ref{C-352-15} is used in its form (f), which bounds the crossing distance by $1$ in the place of the bound $e^{-2\sqrt3}$ of the placing: at the width $r$ and the crossing distance $\mathsf d_\times=\frac38\mathsf d'\in[6r,1]$, with the corners of the two kinds at least $\frac12\mathsf d'\ge\mathsf d_\times/\sqrt3$ apart, its conditions hold as in the proof of Lemma~\ref{C-352-39}. By the support rule (Lemma~\ref{C-352-12}(d)) a domain of the level $i+1$ that holds corners of both kinds has a linear size at least $D_{i+1}=[\mathsf d'/2Ma_{K-i-1}-2]_+$.

\emph{The tower.} Let $n_0\le i<K$, $A'\in\mathsf{Cc}_{i+1}(r_{i+1})$ and $\nu=\nu_{i,A'}$. The tower identity at several sources (step 1 of the proof of Proposition~\ref{C-352-21}; the two clusters, the near sources and the piece, lie in the frame at the level $n_0-1$, as $4M\rho_0\le\mathsf d'$), at the five partitions of $\{1,2,3\}$ and the two of $\{T,3\}$, with every $G_{i;\beta}$ taken at $A^{\mathrm{lw}}(B)$, gives
\[
\Gamma_{i+1}(A')=\int\Gamma_i\,d\nu+\kappa_\nu(\Delta_i,G_{i;3})+\mathsf c\bigl[\kappa_\nu(G_{i;13},G_{i;2})+\kappa_\nu(G_{i;23},G_{i;1})+\kappa_\nu(G_{i;1},G_{i;2},G_{i;3})\bigr].
\]
As $A^{\mathrm{lw}}(B)\in\mathsf{Cc}_i(r_i)$ and $\nu$ is a probability measure, the mean is at most $\sup_{\mathsf{Cc}_i(r_i)}|\Gamma_i|$; as $\mathsf{Cc}_K(r_K)=\{0\}$, $|\Gamma_K(\mathbf 1)|=|\mathsf c\mathsf J_3-\mathsf J_2|$ is at most $\sup_{\mathsf{Cc}_{n_0}(r_{n_0})}|\Gamma_{n_0}|$ plus the sum over $n_0\le i<K$ of the suprema of the four cumulants.

\emph{The covariance of $\Delta$ and of the far source.} By the oscillations: the explicit parts of $\Delta_i$ cancel by the moment condition of step (iii) of the proof of Lemma~\ref{C-352-38}, and Lemma~\ref{C-352-38a} bounds the oscillation of $\Delta_i$ on $\mathsf{Cc}_i(r_i)$ by $2\varepsilon_\alpha r_i^2a_{K-i}^{-4}$; Lemma~\ref{C-352-16} and $|\mathsf{ang}^{[i]}_p|\le2\eta_i^2r_i$ bound that of $G_{i;3}$ by $C\mu_3r_i^2a_{K-i}^{-4}$; with $x_ir_i^2\le C\mathfrak b_i^2$ and $a_{K-i}x_i\ge\frac45\mathsf b$, $|\kappa_\nu(\Delta_i,G_{i;3})|\le C\varepsilon_\alpha\mathfrak b_i^4\mathsf b^{-2}a_{K-i}^{-6}\mu_3$. By the norms: $\kappa_\nu(\Delta_i,G_{i;3})$ is the value of $\mathsf{Sum}$ of the mixed second derivative of the step in the array of $\Delta_i$, of the size at most $Ca_{K-i}^{-4}$ (Lemma~\ref{C-352-12}, and Lemma~\ref{C-352-15} in the form (f)), and in the array of the far source; this family vanishes on the domains with $d_{i+1}(X)<D_{i+1}$, and Lemma~\ref{C-352-12}(b), the sizes of Lemma~\ref{C-352-15}(a), in the form (f), and Lemma~\ref{C-352-14}(b) give $|\kappa_\nu(\Delta_i,G_{i;3})|\le C\mathsf b^{-1}a_{K-i}^{-7}e^{-D_{i+1}}\mu_3$.

The smaller of the two bounds is at most their geometric mean, $C\varepsilon_\alpha^{1/2}\mathfrak b_i^2\mathsf b^{-3/2}a_{K-i}^{-13/2}e^{-D_{i+1}/2}\mu_3$; with $\mathfrak b_i=A_1y_i^{\mathfrak p_0}$ and $y_i\le\log(4\mathsf b/3a_{K-i})$ the sum over $i$ is led by the levels with $a_{K-i}$ about $\mathsf d'$ and gives the first term below.

\emph{The crossing terms.} $\kappa_\nu(G_{i;1},G_{i;2},G_{i;3})$ is the value of $\mathsf{Sum}$ of the third derivative of the step in the arrays of the three single sources, and $\kappa_\nu(G_{i;13},G_{i;2})$, $\kappa_\nu(G_{i;23},G_{i;1})$ are the values of $\mathsf{Sum}$ of its mixed second derivatives in the arrays of a crossing set and of a single source; these families vanish on the domains with $d_{i+1}(X)<D_{i+1}$. Lemma~\ref{C-352-12}(b) at $q=3$, the sizes of Lemma~\ref{C-352-15}(a), in the form (f), and Lemma~\ref{C-352-14}(b) give $\mathsf c|\kappa_\nu(G_{i;1},G_{i;2},G_{i;3})|\le C\mathsf c\,\bar\nu_{i;123}\varsigma_i^3e^{-D_{i+1}}\le C\mathsf b^{-1/2}r^3a_{K-i}^{-21/2}e^{-D_{i+1}}\mu_3$, the widths of the three sources being of the order of $r$; the two others are smaller, by a factor $\gamma_i$ (Lemma~\ref{C-352-15}(e), in the form (f)). The sum over $i$, led by the levels with $a_{K-i}$ about $\mathsf d'$, gives the second term below.

\emph{The start.} At the configurations of $\mathsf{Cc}_{n_0}(r_{n_0})$, whose pairs lie in the cores, 
\[
\Gamma_{n_0}=\sum_{i\le n_0}\mathsf{Sum}[\mathsf cE^{(i)}_{123}-E^{(i)}_{T3}];
\]
 the first display of Lemma~\ref{C-352-39}, with 
\[
\bar\nu_{i-1;\beta}\le\bar c_m(5L/4\mathsf b)^{m/2}a_{K-i}^{-m/2}\prod_{a\in\beta}|w_a|_\infty
\]
 and 
\[
e^{-D^\sharp_i/2}\le e\cdot e^{-\mathsf d'/4Ma_{K-i}},
\]
 summed over $i\le n_0$, where the quotient of two consecutive terms is at most $L^{3/2}e^{-(L-1)}<1$ as $4M\rho_0\le\mathsf d'$, gives 
\[
|\Gamma_{n_0}|\le C(\mathsf b/\rho_0)^{3/2}e^{-\mathsf d'/4M\rho_0}\mathsf b^{-2}\mathsf d'^{-6}\mu_3,
\]
 at most the third term below.

The result is
\[
\begin{multlined}
|\Gamma_K(\mathbf 1)|\le C\Bigl[\Bigl(\frac{\mathsf b}{\mathsf d'}\Bigr)^{1/2}\mathsf L_{\mathsf d'}\varepsilon_\alpha^{1/2}+\Bigl(\frac{\mathsf b}{\mathsf d'}\Bigr)^{3/2}\Bigl(\frac r{\mathsf d'}\Bigr)^3+\Bigl(\frac{\mathsf b}{\rho_0}\Bigr)^{3/2}e^{-\mathsf d'/8M\rho_0}\Bigr]\mathsf b^{-2}\mathsf d'^{-6}\mu_3,\qquad\mathsf L_{\mathsf d'}=\Bigl(1 \\
+\log^+\frac{\mathsf b}{\mathsf d'}\Bigr)^{2\mathfrak p_0},
\end{multlined}
\]
where $\varepsilon_\alpha\le C(g,\mathsf b)\epsilon$ is the number of Lemma~\ref{C-352-38a}. Then $\varepsilon_\alpha^{1/2}\le C(g,\mathsf b)[r^{\beta_{\mathrm H}/72}+(a_K/r)^{1/18}]$, $(r/\mathsf d')^3\le\mathsf d'^{-3}r^{\beta_{\mathrm H}/72}$, $e^{-s}\le2s^{-2}$ at $s=\mathsf d'/8M\rho_0$ with $\rho_0<Lr^{1/2}$, and $(1-\frac1{64})\mathsf d\le\mathsf d'\le6\mathsf d\le1$ give the claim.
\end{proof}

\begin{lemma}[Uniformity in the label]\label{C-352-41}
Let $0<g\le\gamma_{\mathrm R}$ and let $\mathcal K\subset(g^{-2},\bar{\mathsf b}]$ be compact. Then $C_{\mathrm A}(g,\mathsf b)$, $\mathsf C_{\mathrm R}(g,\mathsf b;\mathsf d)$ and $\mathsf C(\mathsf d)$ of Theorem~\ref{C-352-7} are bounded, and $r_\alpha(g,\mathsf b)$, $r_{\mathcal W}(g,\mathsf b,\mathsf d)$, $r_\sharp$, $r_0(\mathsf b)$, $r_1(\mathsf b)$ are bounded away from $0$, uniformly in $\mathsf b\in\mathcal K$; here $r_\sharp$ and $\mathsf C(\mathsf d)$ are the numbers that steps 2 and 3 of the proof of (E.3) form at a first-passage datum at $g$ with the label $\mathsf b$, for every $\mathsf b\in(g^{-2},\bar{\mathsf b}]$.
\end{lemma}

\begin{proof}
\emph{The thresholds.} $r_1$ does not increase in $\mathsf b$, so $r_1(\mathsf b)\ge r_1(\bar{\mathsf b})>0$, and the infimum of $r_0$ over the labels at $g$ is positive by (B.1). The threshold $r_\alpha(g,\mathsf b)$ is the largest number below which the inequalities \textup{($\alpha$1)}--\textup{($\alpha$5)} of Definition~\ref{C-352-2}(5) hold, and the bound $r\le r_\alpha(g,\mathsf b)$ enters Lemma~\ref{C-352-38} and the lemma on the chain classes through these inequalities; by the bounds given there each of them holds as soon as $r$ is below a number formed from the constants and $g$, or, for \textup{($\alpha$3)}, below $2\min\{\inf_{\mathsf b'}r_0(\mathsf b'),r_1(\bar{\mathsf b})\}$; so $r_\alpha(g,\mathsf b)\ge r^-_\alpha(g)>0$ for every label at $g$. $r_{\mathcal W}(g,\mathsf b,\mathsf d)$ adds to $r_\alpha$ the inequalities of step 0 of the proof of (E.2), which depend on $\mathsf b$ only through $a_\sharp=4\mathsf b/5x_\sharp\ge4g^{-2}/5x_\sharp$; and $r_\sharp$ is the least of these thresholds and of those of the data. So all the thresholds are bounded below by positive numbers formed from the constants, $\beta_{\mathrm H}$, $g$ and $\mathsf d$, for all labels in $(g^{-2},\bar{\mathsf b}]$.

\emph{The constants.} The label enters the proofs of Lemmas~\ref{C-352-38}--\ref{C-352-40} and of Theorem~\ref{C-352-7} through five kinds of quantities. (1) The run: by Lemma~\ref{C-352-10}(a), $\frac45\mathsf b\le a_{K-j}x_j\le\frac43\mathsf b$, so that $x_j$, $y_j$, $\delta_j$, $\varepsilon_j$, $\mathfrak b_j$, $\gamma_j=c_\gamma x_j^{-1/2}$, $\theta^{[j]}$ and $r_j$ lie between two functions of $\mathsf b/a_{K-j}$, each continuous in $\mathsf b$, and $\frac3{4\mathsf b}\le\upsilon_j\le\frac5{4\mathsf b}$. (2) The prefactor $\mathsf c=\frac{2\pi}3\mathsf br^3$ and the units $\mathsf u_a=\mathsf b^{-1}|w_a|_\infty$. (3) The slices, through $\mathsf b^{-1}$ and $\bar B/\mathsf b$ (Lemma~\ref{C-352-22}). (4) The number $a_\sharp=4\mathsf b/5x_\sharp$, through $a_*^{-1}\le L\max\{8M,5x_\sharp/4\mathsf b\}\ell_\epsilon/\mathsf d$. (5) The thresholds, bounded below as above.

Member by member: $\gamma_{n_0}\le c_\gamma(5Lr^{1/2}/4\mathsf b)^{1/2}$; $s(\frac r2)\theta^{[0]}$, with $x_0=\mathsf b/a_K$; $\sum_i\hat r_i^3\mathfrak b_i^2(\delta_i+\epsilon^{\mathrm{cu}}_i)\Lambda_i^2\mathsf h_i^2$, with $x_i\ge\frac45\mathsf b/a_{K-i}$, while $\Lambda_i\le C(1+Y_{r/2})$ and $\mathsf h_i\le C(1+Y_{r/2})^{2\mathfrak p_0}$ are formed with $g$ and $r$; the cost, at most $C\mathsf b\,r\,\mathrm{nat}^{\mathcal W}_n$, of the domains that leave the cube of the gauge, and the part of the steps between two levels, of the same order (proof of Lemma~\ref{C-352-38}); $y_{n_0}\le\log(4\mathsf b/3\rho_0)$; the bracket of the proof of Lemma~\ref{C-352-40}, through $(\mathsf b/\mathsf d')^{1/2}$, $(\mathsf b/\mathsf d')^{3/2}$, $(\mathsf b/\rho_0)^{3/2}$ and $\mathsf L_{\mathsf d'}$; the size $b_\Delta=(C\mathsf b+C)\Lambda_3$ of $\mathsf Y_\Delta$, the terms $T_1$--$T_4$ and the moments $\mathsf M_b$ of the proof of (E.2); the number $\epsilon_0(g,\mathsf b)$, the reciprocal of the constant in $\sum_i\epsilon^{\mathrm e}_i\le C(g,\mathsf b)\,\epsilon$; the quotients that form $\mathsf C_{\mathrm R}$; and the constant $C(\mathsf b,\mathsf d)$ of step 3 of the proof of (E.3).

Each majorant so formed is a sum of products of a factor that does not contain $\mathsf b$ and of a function of $\mathsf b$ made by finitely many products, powers, logarithms and maxima from $\mathsf b$, $\mathsf b^{-1}$, the constants, $\beta_{\mathrm H}$, $g$, $\mathsf d$ and the lower bounds of the thresholds; such a function is continuous on the compact interval $[g^{-2},\bar{\mathsf b}]$, where $g^{-2}\ge e^{2\mathfrak p_0}>1$, and so bounded there. Hence $C_{\mathrm A}(g,\mathsf b)$, $\mathsf C_{\mathrm R}(g,\mathsf b;\mathsf d)$ and $\mathsf C(\mathsf d)$ are bounded on all of $(g^{-2},\bar{\mathsf b}]$ by numbers formed from the constants, $\beta_{\mathrm H}$, $g$, $\mathsf d$ and $\mathsf d\operatorname{Lip}\varphi_1/\|\varphi_1\|_\infty$, in particular on $\mathcal K$.
\end{proof}

\paragraph{Proofs of the theorems.}

\begin{proof}[Proof of Theorem~\ref{C-352-3}]
\emph{(B.2).} 1. \emph{The data.} By Lemma~\ref{C-352-11}(b) the tents put the pair in the frame at $k_0=\mathsf k(r)$; $(\Gamma1)$ gives $g\le\gamma_{\mathrm{jt}}$ and $g^{-2}\ge x_{\mathrm J}$, so Lemma~\ref{C-352-12} applies, with (P1)--(P3), and $\rho\le r_1(\mathsf b)$ gives $\rho\ge4Ma_K$ and $E_\sigma(\hat\rho_1/2)\le\varkappa_0^4$; the smallness holds at all heights (Lemma~\ref{C-352-10}(b)). At $n=K$, $A=0$ and both members vanish; let $\mathsf k(r)\le n<K$.

2. \emph{The decomposition.} Differentiate the defining formula for $D_n^\lambda$ twice: the source term is linear in $\lambda$, so $\mathsf F_n=\sum_{i\le n}\mathsf{Sum}[E^{(i)}_{12}]$; with Lemma~\ref{C-352-13} and $q_n=\sum_{i=\mathsf k(r)}^{n-1}\mathcal Q_i$, $q_{\mathsf k(r)}=0$,
\[
\mathsf F_n(\tilde U)-\mathsf F_n(\mathbf 1)-q_n(\mathsf{ang}^{[n]}(A))=\underbrace{\sum_{i=1}^{\mathsf k(r)}\Delta[E^{(i)}_{12}]}_{\mathrm I}+\underbrace{\sum_{j=\mathsf k(r)}^{n-1}\Delta[\mathsf R^{(j+1)}]}_{\mathrm{II}}+\underbrace{\sum_{i=\mathsf k(r)}^{n-1}\bigl(\Delta[\mathsf K^{(i+1)}]-\mathcal Q_i(\mathsf{ang}^{[n]}(A))\bigr)}_{\mathrm{III}} .
\]

3. \emph{The passage to the yardstick} $\mathsf Y_n=(\rho/r)^6\hat r_n^3\mathsf u_1\mathsf u_2$. For a pointed family $\mathsf f$ of level $j+1$ that vanishes off $\mathsf{Cr}_1$ and on the domains with $d_{j+1}(X)<D'$, with $\|\mathsf f\|\le\Phi\bar\nu_j$, Lemma~\ref{C-352-14}(a) in the frame at $k_0=\mathsf k(r)$ of step 1, $\bar\nu_j\le\frac{25}{16}\bar c_2x_j\mathsf u_1\mathsf u_2$, $x_j\le\frac53L^{n-j}x_n$ and $x_nr_n^2\le C\mathfrak b_n^2$ give
\[
|\Delta[\mathsf f]|\le C_{\mathrm U}\,\mathfrak b_n^2\,\mathsf N_{j+1}(w_1)\,\Phi\,u_j^{-3}e^{-D'/2}\,u_n^3\mathsf u_1\mathsf u_2 .
\]

4. \emph{I, the levels at and below the working level, are cut by the distance.} By the support rule $D'=\mathsf D(u_i)$; Lemma~\ref{C-352-15}(b) and (P2) give $\mathsf N_i\bar\Psi_{i-1;12}\le Cq^6(1+u_i^6)$; with $u_i=L^lu_{\mathsf k(r)}$ and $u_{\mathsf k(r)}/\sqrt3M\ge\frac4{\sqrt3}Y_r^{2\mathfrak p_0}$ the sum is led by $i=\mathsf k(r)$: $|\mathrm I|\le C\mathfrak b_n^2e^{-Y_r^{2\mathfrak p_0}/\sqrt3}\mathsf Y_n$.

5. \emph{II, the rest families.} Step 3 with $D'=0$ and Lemma~\ref{C-352-15}(c); $\mathsf N_{j+1}\varsigma_j^2u_j^{-3}\le216L^3q^6u_{j+1}^3$, $\mathsf N_{j+1}\Xi^r_ju_j^{-3}\le216q^6$, $\sum_{j<n}\gamma_j\le0.45\gamma_n\le C\delta_n/\mathfrak b_n$, and $u_{\mathsf k(r)}^3[\gamma_{\mathsf k(r)}+\theta^{[0]}]\le CY_r^{6\mathfrak p_0}\delta_{\mathsf k(r)}$, $\delta_{\mathsf k(r)}\le A_1Y_r^{\mathfrak p_0}(r/\mathsf b)^{1/2}$ (Lemma~\ref{C-352-10}(c)):
\[
|\mathrm{II}|\le\bigl\{C\mathfrak b_n\delta_n+C\mathfrak b_n^2Y_r^{7\mathfrak p_0}(r/\mathsf b)^{1/2}\bigr\}\mathsf Y_n .
\]

6. \emph{III, the kernel families.} Proposition~\ref{C-352-19}, narrow form; $\mathsf c_i\le Ce^{-\delta_{\mathrm{tp}}u_i/\sqrt3}$, so $\sum_i\mathsf c_i(1+u_i)^2u_i^3\le C$, and $\sum_{i\ge\mathsf k(r)}\mathfrak b_i^4\delta_iu_i^3\le1.12\,\mathfrak b_{\mathsf k(r)}^4u_{\mathsf k(r)}^3\delta_{\mathsf k(r)}$:
\[
|\mathrm{III}|\le\bigl\{CY_r^{15\mathfrak p_0+2}(r/\mathsf b)^{1/2}+C\mathfrak b_n^2r_n\bigr\}\mathsf Y_n .
\]

7. \emph{The sum.} $\mathfrak b_n\le\mathfrak b_{\mathsf k(r)}\le A_1Y_r^{\mathfrak p_0}$, $r_n\le\frac85\mathsf c_{\mathfrak h}\delta_n$ and $(r/\mathsf b)^{1/2}\le g\,r^{1/2}$ give the bound. The factor $(\rho/r)^6$ costs nothing: $\mathsf F_n$ is bilinear in $(w_1,w_2)$, every estimate depends on the weights only through $\sum_p|w_a(p)|$ or through a norm that (P1) bounds, and the one place where a sup norm could enter, $\Lambda_i$, sees a quotient (Proposition~\ref{C-352-19}, step 6). By Lemma~\ref{C-352-10}(c), $\log\frac1r\le Y_r\le2\log\frac1r+C(g)$, whence $\epsilon'_{\mathrm N}(r)(\log\frac1r)^N\to0$.

\emph{(B.1).} 1. Lemma~\ref{C-352-11}(a) and its proof put the test pair in all six conditions of Definition~\ref{def:pair-functional-frame} at $k_0=\max\{K-s,2\}$, under $3\le s\le K$ and $\mathsf{Lm}_s\le L^{s-3}$, as required in (B.1); there the number $\theta^{[j]}/K_{\mathbf M}$ of every step is at most half of the largest value that the one-step theorem admits, $\hat r_{k_0}\ge M$, so $\varsigma_i=1$ for $i\le k_0$, and $\mathsf D^{12}_i\ge\hat r_i/\sqrt3-1$. 2. Below $k_0$, Proposition~\YMMref{122.2}{827} gives $\|E^{(j)}_{12}\|\le C\nu_{j-1}$, with $\nu_j$ defined in \eqref{eq:weighted-mixed-numbers}; its running, half-ceiling and separation hypotheses are those verified in the pair frame of Lemma~\ref{C-352-11}(a); from $k_0$ on, Lemma~\ref{C-352-12}(b), (c) and the start at $k_0$ give the rest families with $\|\mathsf R^{(j+1)}\|_{j+1,\lambda_{\mathrm{out}}}\le\epsilon^{(2)}_j\varsigma_j\bar\nu_j$,
\[
\epsilon^{(2)}_j\le C\bigl[x_j^{-1/2}(1+\mathsf e_\downarrow)+\theta^{[k_0]}\bigr],\qquad\mathsf e_\downarrow=\sum_{i\le k_0}e^{-4\beta_{\mathrm s}\kappa[\mathsf D^{12}_i/M-2]_+}\varkappa_{i-1}^{-2}:
\]
no number in it grows with $\log\frac1r$ or with $n-k_0$, as the derivatives of the step gain the factor $1/\mathfrak b_j$ per order in the sources (Lemma~\ref{C-352-12}(b)) and the lags of Lemma~\ref{C-352-12}(c) fall like $L^{-4l}$ and sum geometrically over the levels. 3. The decomposition of step 2 above, over the base level $k_0$, with Lemma~\ref{C-352-14}(a), in the frame at $k_0$, for I and II and Proposition~\ref{C-352-19} at the test pair for III, both in the form stated after them, whose conditions are the frame, the smallness at the heights concerned ($2r_n\le\varrho_{\mathrm{ch}}$, so that $r_n\le\frac1{120}$ as $\varrho_{\mathrm{ch}}\le\frac1{60}$; $\epsilon^{\mathrm{cu}}_i\le1$; $x_n^{-1}\le r_n^2$) and the side bound $\ell\ge2LMy_K$: for a test pair at $k_0=\max\{K-s,2\}$ these hold for $g\le\min\{\gamma_{\mathrm{jt}},\gamma_{\mathrm O}\}$ and $r\le r_0(\mathsf b)$, the frame by Lemma~\ref{C-352-11}(a), the side bound as $\ell\in\mathfrak L(g)$, and the smallness at the heights $\ge k_0$ being among the conditions that fix $\gamma_{\mathrm O}$ and $r_0(\mathsf b)$.

The shape $\vartheta$ does not enter, since any $\bar{\mathsf e}\ge2\mathsf w_1\mathsf w_2\mathsf E_i$ is admissible. 4. In the yardstick $\hat r_n^3\mathsf u_1\mathsf u_2$ the main part of II is $\le C\mathfrak b_n^2[\gamma_n(1+\mathsf e_\downarrow)+(n-k_0)\theta^{[k_0]}]$, III is $\le C\hat r_{k_0}^3\mathfrak b_{k_0}^4(\delta_{k_0}+\epsilon^{\mathrm{cu}}_{k_0})\Lambda^2\mathsf h^2+C\mathfrak b_n^2r_n$, by the bound $\mathsf E_i\le C_{\mathsf E}\delta_i^4(\delta_i+\epsilon^{\mathrm{cu}}_i)$ of Proposition~\ref{C-352-19}, and I is $\le C\mathfrak b_n^2e^{-c'Y_s^{2\mathfrak p_0}}$. With $\theta^{[k_0]}\le C(r/\mathsf b)^{1/2}y_{k_0}^{\mathfrak p_0}$, $K-k_0\le s\le Cy_{k_0}$ and $\mathsf e_\downarrow\le Cy_{k_0}^{2\mathfrak p_0}e^{-c'Y_s^{2\mathfrak p_0}}$, and the smallness that the recursions need holding for $g\le\gamma_{\mathrm O}$ and $r\le r_0(\mathsf b)$, this is the bound of (B.1); its last term $C\mathfrak b_n^2r_n\le C\mathfrak b_n^2\delta_n$ is the price of writing the form with the linear angles of one chart. 5.

The conditions that fix $r_0(\mathsf b)$ depend on $\mathsf b$ through $(r/\mathsf b)^{1/2}$, through $y_{k_0}\le\log\frac{\mathsf b}r+C\log\log\frac{\mathsf b}r$ and through $x_{k_0}\le\frac{24}{23}\mathsf b\hat r_{k_0}/r$; on $(g^{-2},\bar{\mathsf b}]$ these are bounded by their values at the two ends, so $r_0$ has a positive infimum over the labels at $g$.

\emph{(B.3).} $|\mathsf{ang}^{[n]}_p(A)|\le2\eta_n^2r_n$, Lemma~\ref{C-352-22}(a) with $\mathsf r\ge r/\sqrt3$, and (P1):
\[
|q_n(\mathsf{ang}^{[n]}(A))|\le C\,r_n^2\,\mathsf b\,a_{K-n}^{-4}\rho^6r^{-3}\mathsf u_1\mathsf u_2=C\,r_n^2\,\mathsf b\,a_{K-n}^{-1}(\rho/r)^6\hat r_n^3\mathsf u_1\mathsf u_2\le C\mathfrak b_n^2(\rho/r)^6\hat r_n^3\mathsf u_1\mathsf u_2,
\]
as $\mathsf b\,a_{K-n}^{-1}\le\frac54x_n$ and $x_nr_n^2\le C\mathfrak b_n^2$.
\end{proof}

\begin{proof}[Proof of Theorem~\ref{C-352-4}]
1. \emph{(A.1).} Proposition~\ref{C-352-24}(a) bounds $\mathsf J_3$ against the lattice ring ($64\mathsf d\le\ell_g$ as $\ell_g\ge26$ and $\mathsf d\le e^{-34}$; $\rho\le r_1(\mathsf b)$ by $\mathsf d\le\mathsf d_4(g)\le r_1(\bar{\mathsf b})\le r_1(\mathsf b)$), and Lemma~\ref{C-352-26} bounds $S^T_3-\mathsf J_3$ by $13\Lambda_3^3\mathsf M'^3\mathsf b^3\mathsf d^{3/2}\mathrm{nat}_3$ (all supports lie within $20\mathsf d\le\frac18\ell_g-1$ of $x_1$).

2. \emph{The geometry of (A.2).} For a triple $(z_1,z_2,z_3)=(x,y,z)$ of $\mathcal D$ put $r=\frac12r_*$, $\mathsf d=\frac12\mathsf d_*$; bumps $f_a$ with supports in balls of radius $\rho<\min\{\frac14r_*,r_1(\bar{\mathsf b})\}$ inside the open balls of radius $\frac14r_*$ about the $z_a$ are placed at $(\rho,r,\mathsf d)$: their supports are $\ge r_*-\frac12r_*=r$ apart in the pair and $\ge\mathsf d_*-\frac12r_*\ge\mathsf d$ across, with $|x_1-x_2|\le3r$, $|x_a-x_3|\le5\mathsf d$, all in $\mathbb B(z_1,\frac94\mathsf d_*)$; this holds for every ratio $r_*/\mathsf d_*\in(0,1]$, and $\mathsf d_*\le\mathsf d_5(g)$ gives $\mathsf d\le\mathsf d_4(g)$.

3. \emph{Lattice to continuum.} Lemma~\ref{C-352-27}(a) with $\mathsf d^+=5\mathsf d_*$ and a Lipschitz number $\vartheta$ of the $f_a$:
\[
|3\operatorname{Tr}\prod_a((f_a\circ x)\mathsf b^{-1}C^{\mathrm{free}}_K)-3\mathsf b^{-3}\!\int k_3\prod f_a|\le3C^{(3)}\bigl[a_K/\vartheta\rho+(5\mathsf d_*/\ell_g)^3\bigr]\mathrm{nat}_3 .
\]

4. \emph{The limit.} Along first-passage data with $K_j\to\infty$ and labels $\to\mathsf b$, $S^T_3(f)\to S^{(\ell_g,\mathsf b)}_3(f)$ (Theorem~\YMMref{A}{5}), $a_{K_j}/\vartheta\rho\to0$, and the brackets are bounded with $\bar{\mathsf b}$ and $\bar{\mathsf M}^{(24)}$ (Lemma~\ref{C-352-25}(e)); with $r^{-3}\mathsf d^{-6}=512\,r_*^{-3}\mathsf d_*^{-6}$ this is the bound of the form in (A.2).

5. \emph{Kernels.} Lemma~\ref{C-352-28}(a) on $\prod_aU_a$, with $(\frac3\pi)^3\le1$: a Borel kernel, bounded at its Lebesgue points by the bound of the form at the centre; the cover of $\mathcal D$ by such products and the glueing of Lemma~\ref{C-352-35} give one kernel on $\mathcal D$; the closest pair of a triple is $(x,y)$ after a relabelling. \emph{(A.3)} is Lemma~\ref{C-352-28}(b).

6. \emph{(A.4).} (i) is Proposition~\ref{C-352-33}(a) at $n=3$. (ii) For a triple with $\mathsf d_*\le\mathsf d_5(g)$, the bumps of Lemma~\ref{C-352-34}(c) about it are in the data of (i) at every large $j$ (the cube of side $\frac92\mathsf d_*\le\frac12\ell_g$ about $x$ holds them, so the projection keeps all distances and the integrals of $k_3$); (i), step 4 at the label $\mathsf b$, and $S^T_3[\nu_j]\to S^{\sigma'}_3$ give $|\int(\mathsf s_3-3\mathsf b^{-3}k_3)\prod f^\rho_a|\le[\epsilon_{\mathrm A}+512C^{\mathrm{tr}3}_3\mathsf b^3\mathsf d_*^{3/2}]\mathsf b^{-3}r_*^{-3}\mathsf d_*^{-6}\cdot0.102\le\mathsf E_3(g,\mathsf d_*)\mathsf b^{-3}r_*^{-3}\mathsf d_*^{-6}$, and Lemma~\ref{C-352-34}(c) gives the bound at the point. (iii) is Lemma~\ref{C-352-34}(a).

7. \emph{(A.5).} Let $(K_j,\beta_j,\nu_j)$ be a Gibbs sequence at $g_0$; its members are periodic Gibbs states, $\nu_j\in\mathcal G^{\mathrm{per}}_{K_j,\beta_j}$. For large $j$, $\beta_j$ exceeds $L(g^{-2}+B_\sharp)$; let $K'_j$ be the first-passage index of $\beta_j$ at $g$ and $z_j=K_j-K'_j$; as $g\le g_0$, $0\le z_j\le z_+$ with $z_+$ formed from $g_0$, $g$ and the constants. (1) \emph{The states.} For $0\le z<K$ the dilation $U\mapsto U^{(z)}$, $U^{(z)}(L^zb)=U(b)$, maps the bonds, the plaquettes and the finite sets of bonds of $\mathbb L_K$ onto those of $\mathbb L_{K-z}$, with $a_{L^zp}(U^{(z)})=a_p(U)$. The Gibbs kernel of a finite set of bonds has the exponent $-\beta\sum_pa_p$ over the plaquettes that contain a bond of the set, which reads $\beta$ and the lattice and not its spacing; so the image $\nu^{(z)}$ of a Gibbs state at $(a_K,\beta)$ satisfies the equations of the kernels at $(a_{K-z},\beta)$, and invariance under a subgroup $\Gamma'$ of finite index of $a_K\mathbb Z^3$ becomes invariance under $L^z\Gamma'$, of finite index in $a_{K-z}\mathbb Z^3$.

So the dilation carries $\mathcal G^{\mathrm{per}}_{K,\beta}$ onto $\mathcal G^{\mathrm{per}}_{K-z,\beta}$ (equation~\YMMref{(1819)}{698} and the definition of periodic states following it), and $\nu^{(z_j)}_j$ is a periodic Gibbs state at the first-passage datum $(K'_j,\beta_j)$ at $g$, with the label $L^{z_j}a_{K_j}\beta_j$. (2) \emph{The fields.} At every configuration $U$ of $\mathbb L_K$, $x(L^zp)=L^zx(p)$ and $a_{K-z}=L^za_K$ give $\mathcal O(f)(U)=L^z\mathcal O'(f^{(z)})(U^{(z)})$, $f^{(z)}=f(L^{-z}\cdot)$, the prime marking the field at the cutoff $K-z$, as in Lemma~\ref{C-352-36}(d), which is the dilation; so in every law $\nu$ on $\mathbb L_K$, $S^T_3[\nu](f)=L^{3z}S^T_3[\nu^{(z)}](f^{(z)})$. (3) \emph{The approximants} of a component of a canonical decomposition of the limit are periodic Gibbs states at the same data, by the definition of the components and their approximants in Lemma~\ref{C-352-34}(a); (1) and (2) apply to them. (4) \emph{The limit.} On a subsequence $z_j=z$ is constant and the labels converge, to $\mathsf b_0$ say; along it $S^T_3[\nu^{(z)}_j](f)=L^{-3z}S^T_3[\nu_j](f(L^z\cdot))\to L^{-3z}S^\sigma_3(f(L^z\cdot))$ on separated triples, $\sigma$ the limit, and this functional has the continuous kernel $\mathsf s^{(z)}_3=L^{-12z}\mathsf s_3(L^{-z}\cdot)$; likewise for a component and its approximants.

The proof of (A.4)(ii) uses of the law $\sigma'$ the convergence $S^T_3[\nu_j]\to S^{\sigma'}_3$ on separated triples and the continuity of the kernel of $S^{\sigma'}_3$ at the triple (Lemma~\ref{C-352-34}(c)); so (A.4)(ii) at $g$ applies to the dilated sequence, whose labels tend to $L^z\mathsf b_0$, at the triples with $L^z\mathsf d_*\le\mathsf d_5(g)$. In kernels $\mathsf s_3(x)=L^{12z}\mathsf s^{(z)}_3(L^zx)$, and the natural size is invariant: $L^{12z}(L^z\mathsf b_0)^{-3}(L^zr_*)^{-3}(L^z\mathsf d_*)^{-6}=\mathsf b_0^{-3}r_*^{-3}\mathsf d_*^{-6}$, $k_3(L^z\cdot)=L^{-9z}k_3$. The approximants of a component are periodic Gibbs states at the same couplings and are dilated alike.
\end{proof}

\begin{proof}[Proof of Theorem~\ref{C-352-5}]
1. \emph{(C.1).} At $x,y=c\pm\frac r2n$: $r_*=r$; the two distances $|c-z\pm\frac r2n|$ lie in $[\mathsf d-\frac r2,\mathsf d+\frac r2]$ and the sum of their squares is $2\mathsf d^2+\frac12r^2$, so $\mathsf d_*\in[\frac34\mathsf d,\frac{17}{16}\mathsf d]$, $r\le\mathsf d_*\le\mathsf d_5(g)$. (A.2) multiplied by $\frac{2\pi}3\mathsf br^3$ gives $\frac{2\pi}3(\frac43)^6\le12$ times $\bar\epsilon_{\mathrm A}\mathsf b^{-2}\mathsf d^{-6}$, and Lemma~\ref{C-352-29}(c) gives $2.66\le2.7$; ``almost everywhere'' is in $(c,h,z)$, a linear change of the variables $(x,y,z)$. (ii) (A.3) at every triple; $k_3$ is continuous off the diagonals, so its means tend to its values; with tents, (A.3) in the tent form and the second part of Lemma~\ref{C-352-29}(c). (iii) (A.4)(ii) at every triple and Lemma~\ref{C-352-29}(c).

2. \emph{(C.2)(i).} Proposition~\ref{C-352-24}(b) bounds $\mathsf J_2$; Lemma~\ref{C-352-26} at two sources gives $|S^T_2-\mathsf J_2|\le3\Lambda_3^2\mathsf M'^2\mathsf b^2\mathsf d^{3/2}(\rho/\mathsf d)^6\mathsf b^{-2}|w_1|_\infty|w_2|_\infty$; at the plane weights $w_1=\mathbf 1_{\{\lambda(p)=\lambda\}}\varphi_1\circ x$, $w_2=\mathbf 1_{\{\lambda(p)=\mu\}}\varphi_2\circ x$, Lemma~\ref{C-352-27}(c) turns the lattice square into $\mathcal K_{\lambda\mu}$ ($8a_K\le\rho$ by $s(\rho)\le K$; the assumption $|x_1-x_2|\le18\mathsf d$ is the one that this step and the rings need).

3. \emph{(C.2)(ii).} For $(x_0,y_0)\in\Omega_{\mathrm P}$ and bounded Borel $f,h$ supported in the closed balls of radius $\frac18\mathsf d_0$ about $x_0,y_0$, the jets, the remainder, $|C^{\mathrm{free}}_K(p,q)|\le(1+3C_{\mathrm{lat}})a_K^3|u|^{-3}$ and the count of plaquettes give
\[
|\langle\mathcal O_\lambda(f);\mathcal O_\mu(h)\rangle_{K_j,\ell_g,\beta_j}|\le\mathsf B(g,\mathsf d_0)\|f\|_\infty\|h\|_\infty,\qquad\mathsf B=4^{-6}g^4\bigl[\tfrac32\cdot729(1+3C_{\mathrm{lat}})^2+\mathsf E_2(g,\tfrac12\mathsf d_0)\bigr],
\]
with no Lipschitz number (the Riemann sums of step 2 involve $a_K/\vartheta\rho$ and give no bound in the sup norms). Lemma~\ref{C-352-35} gives one subsequence for all pairs of balls; along it step 2, with $a_{K_j}/\vartheta\rho\to0$, gives the bump form at every centre: the balls of radius $\frac18\mathsf d_0$ are $\ge\frac12\mathsf d_0=:\mathsf d$ apart, $|x_1-x_2|\le\frac52\mathsf d$, and $(\rho/\mathsf d)^6=64(\rho/\mathsf d_0)^6$. Lemma~\ref{C-352-28} at $p=2$ gives the kernels and the almost everywhere bound, and Lemma~\ref{C-352-29}(a) gives $t_\lambda-\mathsf t^\circ(e_\lambda,\cdot)=\sum_\mu(\delta_{\lambda\mu}-\frac13\sum_{\lambda'}\delta_{\lambda'\mu})$ with $\delta_{\lambda\mu}=\tau_{\lambda\mu}-\frac32\mathsf b^{-2}[\mathbb K_{3,\lambda\mu}]^2$: three terms and a third of nine, so the factor $6$.

4. \emph{(C.3).} Along the subsequence of the limit point $S^T_3\to S^{(\ell_g,\mathsf b)}_3$, since that limit is taken along all data with labels $\to\mathsf b$. By (A.3) and by Lemma~\ref{C-352-28}(b) applied to the bump form of (C.2)(ii) at the centre $(c,z)$, the means over balls of radius $\rho$ of $\mathsf s_3-3\mathsf b^{-3}k_3$ and of $\tau_{\lambda\mu}-\frac32\mathsf b^{-2}[\mathbb K_{3,\lambda\mu}]^2$ are bounded at every centre, with the factors $(\frac3{4\pi})^3,(\frac3{4\pi})^2\le1$; $k_3$ and $[\mathbb K_3]^2$ are continuous off the diagonals; adding with Lemma~\ref{C-352-29}(a), (c) gives the display, and the tents likewise. The cone form: (C.1)(i) at $n$, (C.2)(ii) at $(c,z)$ almost everywhere (a null set of $(c,z)$ is a null set of $(c,h,z)$), and $|\mathsf t^\circ(n,X)-\mathsf t^\circ(e_\lambda,X)|\le0.057|n-e_\lambda|\mathsf b^{-2}|X|^{-6}$.

5. \emph{(C.4).} (d) is Proposition~\ref{C-352-33}(b) with step 2 on the least torus (the projection is injective and keeps distances on $\mathbb B(x_0,\frac32\mathsf d_0)$, so the integrals and the norms are those of the periodised functions); it gives the bound $\mathsf B(g,\mathsf d_0)$ of step 3 along $\nu_j$ for $j$ large. (a) Lemma~\ref{C-352-35} on $\Omega_\sigma$. (b) The bump form at every centre of $\Omega_\sigma$ with $64\mathsf E_2(g,\frac12\mathsf d_0)\mathsf b^{-2}\mathsf d_0^{-6}\rho^6$; Lemma~\ref{C-352-28}; the bound at a point by centres $\zeta_m\to(x,y)$ of distances $\mathsf d_{\zeta_m}<|x-y|$, whence the supremum in $\epsilon_\sigma$. Every $\nu_j$ is invariant under the translations of $\mathbb L_{K_j}$ by Proposition~\YMMref{103.4}{758}, so the limit forms are invariant under all translations, and by Fubini's theorem $\tau^\sigma_{\lambda\mu}(x,y)$ is a function of $x-y$ almost everywhere. $\sum_{\lambda,\mu}\langle\mathcal O_\lambda;\mathcal O_\mu\rangle=S^T_2$, which converges along the whole sequence to $S^{\sigma'}_2$ with the continuous kernel $\mathsf s^T_2$; with Lemma~\ref{C-352-29}(a) this gives the bound of $\mathsf s^T_2$ almost everywhere, and at every point by continuity, $\epsilon_\sigma(g,\cdot)$ being nondecreasing. (c) Let $\mathcal N$ be the union of the null sets of (b); it is fixed before $r$ is named.

Off $\mathcal N$, (b) and Lemma~\ref{C-352-29}(a) give $|t^\sigma_\lambda-\mathsf t^\circ(e_\lambda,\cdot)|\le6(\frac3\pi)^2\epsilon_\sigma\mathsf b^{-2}\mathsf d^{-6}\le6\epsilon_\sigma\mathsf b^{-2}\mathsf d^{-6}$, and (C.1)(iii) at $n=e_\lambda$ holds at every point, since $\mathsf s_3$ is continuous; $12\sup\mathsf E_3+6\epsilon_\sigma=\epsilon'_\sigma$. With tents of radius $\rho<\frac18\mathsf d$ at $(c,z)$: the bump form of (b) at the centre $(c,z)$, and Lemma~\ref{C-352-29}(d) times $\frac32$, six times, give $21\rho/\mathsf d$.
\end{proof}

\begin{proof}[Proof of Theorem~\ref{C-352-6}]
1. \emph{(D.1) below $\gamma_{\mathrm N}$} is Proposition~\ref{C-352-37}. 2. \emph{(D.1) at every $g\le\gamma$.} Let $z,g'_1,g'_2$ be as in Lemma~\ref{C-352-36} and $K_\circ(g)=z+\max_\iota s'_\iota(r_{g'_\iota})$. For $K\ge K_\circ(g)$, $(K-z,\beta)$ is a first-passage datum at some $g'_\iota\le\gamma_{\mathrm N}$, on the side $L^z\ell\in\mathfrak L(g'_\iota)$, and the conditions of Proposition~\ref{C-352-37} hold at $(L^zr,L^z\mathsf d')$; by Lemma~\ref{C-352-36}(d) each display at $(K,\ell,\beta)$ is that display at $(K-z,L^z\ell,\beta)$. So (D.1) holds with
\[
\epsilon^{\mathrm{cl}}(g,\ell,\mathsf d')=\max_{\iota=1,2}\epsilon^{\mathrm{ev}}_\sharp(g'_\iota,L^z\ell,L^z\mathsf d'),\qquad\mathsf d_{\mathrm{cl}}(g)=\tfrac76L^{-z}\min_\iota t_{\mathrm N}(g'_\iota),
\]
a finite function of $(g,\ell,\mathsf d')$, with the rate for $\mathsf d'\le\mathsf d_{\mathrm{cl}}(g)$. On the dilated torus, at the distance $L^z\mathsf d'$, the members of $\epsilon^{\mathrm{ev}}_\sharp$ that carry the renormalisation group, the remainder and the induction, read the distances through the floor $\min\{\frac67L^z\mathsf d',t_{\mathrm N}(g'_\iota)\}$ alone; the largest distance $\frac87L^z\mathsf d'$ enters the one explicit member, the comparison of the sum of the slices with the free kernel (Lemma~\ref{C-352-22}(b)), which is polynomial in that distance; and the passage from the floor to $L^z\mathsf d'$ costs the factor $\max\{1,6L^z\mathsf d'/7t_{\mathrm N}(g'_\iota)\}^6$ of $\epsilon^{\mathrm{ev}}_\sharp$ (proof of Proposition~\ref{C-352-37}). 3.

\emph{(D.2).} The forms are bounded in the sup norms at every pair of small balls in $\Omega^\circ$, as in step 3 of the proof of Theorem~\ref{C-352-5}; Lemma~\ref{C-352-35} gives one subsequence for the plane and the bond forms; (D.1) with $a_{K_j}/\vartheta r\to0$ gives the bump forms at every centre; Lemma~\ref{C-352-28}, with centres in the window $[\frac79t,\frac97t]$ and $(\frac3\pi)^2(\frac97)^6\le4.2$, gives the kernels. For $T_n$: the kernel of $\langle T_\lambda;\mathcal O\rangle$ is within $3(4.2+4.2)\bar\epsilon^{\mathrm{cl}}=25.2\,\bar\epsilon^{\mathrm{cl}}$ of $\mathsf t^\circ(e_\lambda,\cdot)$, that of $\langle\mathcal O^{\mathrm{od}}_{\lambda\lambda'};\mathcal O\rangle=\sum_\gamma\langle\mathcal O^{\mathrm{od}}_{\lambda\lambda'};\mathcal O_\gamma\rangle$ within $12.6\,\bar\epsilon^{\mathrm{cl}}$ of $3\mathsf b^{-2}[\mathbb K_3^2]_{\lambda\lambda'}$; with $\sum_\lambda n_\lambda^2=1$, $\sum_{\lambda<\lambda'}|n_\lambda n_{\lambda'}|\le1$ and Lemma~\ref{C-352-29}(a), $25.2+12.6\le38$.

The kernel $t_n$ is a quadratic polynomial in $n$ whose coefficients are the nine kernels, so one null set serves all $n$. 4. \emph{(D.3)} is (C.1) and (D.2) by the triangle inequality: almost everywhere in $(c,h,z)$ by Fubini's theorem; for the means and the tents, the bump forms of (D.1) at the centre $(c,z)$ itself, extended to indicators of balls by Lemma~\ref{C-352-28}(b), bound the means of the kernels with factors below those of the almost everywhere bound.
\end{proof}

\begin{proof}[Proof of Theorem~\ref{C-352-7}]
\emph{(E.2).} 0. \emph{The stopping level.} Let $\epsilon=r^{\beta_{\mathrm H}/36}+(a_K/r)^{1/9}$, $\ell_\epsilon=\max\{1,\log\frac1\epsilon\}$, $a^\flat=\min\{\mathsf d/8M\ell_\epsilon,a_\sharp\}$ with $a_\sharp=4\mathsf b/5x_\sharp$, so that $a_{K-n}\le a_\sharp$ gives $x_n\ge x_\sharp$ (Lemma~\ref{C-352-10}(a)), $n_*$ the largest level with $a_{K-n_*}\le a^\flat$, and $a_*=a_{K-n_*}\in(a^\flat/L,a^\flat]$. The conditions that define $r_{\mathcal W}(g,\mathsf b,\mathsf d)$, namely $r\le\min\{r_\alpha,(\mathsf d/5LM)^2,e^{-1}\}$, $8ML^{m_1+2}r^{1/2}(1+\log\frac1r)\le\mathsf d$, $L^{m_1+2}r^{1/2}\le a_\sharp$ and $r^{-1/4}\ge2M$, secure $N_0<n_*$, $x_{n_*}\ge x_\sharp$ and $L^{N_0-n_1}\ge2M$, as $L^{N_0-n_1}=L^{m_1}a_{K-n_0}/a_{K-n_1}\ge r^{-1/4}$ (so Lemma~\ref{C-352-38} holds with $n^\sharp=n_*$ when $\epsilon^{\mathrm e}_i\le1$), $8Ma_*\le\mathsf d$, $e^{-\mathsf d/4Ma_*}\le\epsilon^2$ and $a_*^{-1}\le C\ell_\epsilon/\mathsf d$.

1. \emph{The identity.} Lemma~\ref{C-352-25}(c) at $k=n_*$, $\mathbb P_*=\mathbb P_{n_*}$, for the two lists, $G_3$ being the same function in both:
\[
\begin{multlined}
\mathsf cS^T_3(w_1,w_2,w_3)-S^T_2(w_T,w_3)=\int\Gamma\,d\mathbb P_*+\kappa_{\mathbb P_*}(\Delta,G_3)+\mathsf c[\kappa_{\mathbb P_*}(G_{13},G_2)+\kappa_{\mathbb P_*}(G_{23},G_1)] \\
+\mathsf c\,\kappa_{\mathbb P_*}(G_1,G_2,G_3)=:T_1+T_2+T_3+T_4,
\end{multlined}
\]
with $\Delta$, $\Gamma$, $G_\beta$ taken at $U_{t'}(W)$.

2. \emph{Sizes.} $r^3|w_{1,2}|_\infty\le489$, $|w_T|_\infty\le\frac23|w_1|_\infty$, $\mathsf c|w_1|_\infty|w_2|_\infty r^3\le C\mathsf b$, $|w_3|_\infty\le\mu_3r^{-3}$; the bounds of Lemma~\ref{C-352-26} hold at every point of the small-field region, as their proof uses of $W$ only that the background is regular.

3. \emph{$T_1$, $T_3$: the crossings.} At $t'\le n_*$, $8Ma_{K-t'}\le\mathsf d$, and Lemma~\ref{C-352-39} with $e^{-\mathsf d/4Ma_*}\le\epsilon^2$ gives $|T_1|\le C(\mathsf b)a_*^{-3/2}\mathsf d^{-6}\epsilon^2\mu_3$ and $|T_3|\le Ca_*^{-1}r^3\mathsf d^{-6}\epsilon^2\mu_3\mathsf M_6$, $\mathsf M_b=\int\alpha^{-b}d\mathbb P_*$: no negative power of $r$ is left.

4. \emph{$T_4$.} $|T_4|\le C\mathsf br^3\Lambda_3^3\mathsf M_{18}\mu_3$.

5. \emph{$T_2$: pair minus tensor field.} $T_2=\int(\Delta-\mathsf c_{n_*})(G_3-\int G_3)d\mathbb P_*$, and $|G_3-\int G_3|\le\Lambda_3\mu_3(\alpha^{-6}+\mathsf M_6)$ by step 2. On the parts with $t'\ge N_0$ each measure is absolutely continuous with respect to $\mathsf m_{t'}$, and Lemma~\ref{C-352-38} gives $|\Delta-\mathsf c_{n_*}|\le3.01\,\tilde\varepsilon\,\alpha^{-4}$ almost everywhere. The number $\mathsf c_{n_*}$: the first part of $\mathbb P_*$, on which $t'=n_*$ and $\alpha=a_*$, has positive mass, as the density of $\mathsf m_{n_*}$ on $\mathcal V_{n_*}^\chi$
is a positive exponential by \eqref{eq:weighted-history-chart} and
\eqref{eq:weighted-array-step}. This region is open, its defining
tests are strict, and they hold at the trivial field, where all
curvature and fluctuation defects vanish. It therefore has positive
Haar measure; at almost every point of it $|\Delta-\mathsf c_{n_*}|\le\varepsilon^{\mathcal W}a_*^{-4}$ (Lemma~\ref{C-352-38}) and $|\mathsf Y_\Delta|=|\Delta-\Delta_K(\mathbf 1)|\le b_\Delta a_*^{-6}$, $b_\Delta=(C\mathsf b+C)\Lambda_3$ (Lemma~\ref{C-352-26} at every point, step 2), so that $|\mathsf c_{n_*}-\Delta_K(\mathbf 1)|\le(\varepsilon^{\mathcal W}+b_\Delta)a_*^{-6}$.

On the parts with $t'<N_0$, $\alpha<a_{K-N_0}=L^{m_1}a_{K-n_0}<L^{m_1+1}r^{1/2}<a_*$, so that there $|\Delta-\mathsf c_{n_*}|\le|\mathsf Y_\Delta|+|\Delta_K(\mathbf 1)-\mathsf c_{n_*}|\le(2b_\Delta+\varepsilon^{\mathcal W})\alpha^{-6}$ and $\mathsf M_6\le(1+\mathsf M')\alpha^{-6}$; and $\int_{\{\alpha<\varrho\}}\alpha^{-12}d\mathbb P_*\le\varrho^{12}\mathsf M'$ (Lemma~\ref{C-352-25}(e)) at $\varrho=a_{K-N_0}$ gives a factor $Cr^6$.

6. \emph{Together.} $\mathsf M_b\le a_*^{-b}+\mathsf M'\le C(g,\mathsf b)\ell_\epsilon^b\mathsf d^{-b}$, $\epsilon^{1/2}\ell_\epsilon^{10}\le C$ and $r^{3-\beta_{\mathrm H}/72}(1+\log\frac1r)^{18}\le C$ give $|\mathsf cS^T_3-S^T_2|\le C_{\mathrm A}(g,\mathsf b)\mathsf d^{-18}[r^{\beta_{\mathrm H}/72}+(a_K/r)^{1/18}]\mu_3$. If the smallness $\epsilon^{\mathrm e}_i\le1$ of the single-source arrays of the near pair, which Lemma~\ref{C-352-38} asks at the levels $\mathsf k(\frac r2)\le i<n_1$, fails, then, as $\sum_{\mathsf k(r/2)\le i<n_0}\epsilon^{\mathrm e}_i\le C_\omega[2.2\gamma_{n_0}+\frac12s(\frac r2)\mathsf c^\flat\theta^{[0]}]\le C(g,\mathsf b)\,\epsilon$ (with $\gamma_{n_0}\le C(\mathsf b)r^{1/4}$ and $s(\frac r2)\theta^{[0]}\le C(\mathsf b)(a_K/r)^{1/4}$), the number $\epsilon$ exceeds a number $\epsilon_0(g,\mathsf b)>0$, the bracket is $\ge(\frac12\epsilon_0)^{1/2}$, and the bounds of Lemma~\ref{C-352-26} give the claim with a larger $C_{\mathrm A}$.

\emph{(E.1).} By Lemma~\ref{C-352-25}(c) for the two lists,
\[
\begin{multlined}
[\mathsf cS^T_3-S^T_2]-[\mathsf c\mathsf J_3-\mathsf J_2]=\int\mathsf Y_{\Delta3}d\mathbb P+\kappa_{\mathbb P}(\mathsf Y_\Delta,\mathsf Y_3)+\mathsf E,\qquad\mathsf E=\mathsf c[\kappa_{\mathbb P}(\mathsf Y_{13},\mathsf Y_2)+\kappa_{\mathbb P}(\mathsf Y_{23},\mathsf Y_1) \\
+\kappa_{\mathbb P}(\mathsf Y_1,\mathsf Y_2,\mathsf Y_3)],
\end{multlined}
\]
and Lemma~\ref{C-352-26} gives $|\mathsf E|\le C\Lambda_3^3\mathsf M'\mathsf b^3\mathsf d^{3/2}r^3\cdot\mathsf b^{-2}\mathsf d^{-6}\mu_3$. For $r\le r_{\mathcal W}$ the first bracket is bounded by (E.2) and the second by Lemma~\ref{C-352-40} (if $w_3\equiv0$ all members vanish): together $C(g,\mathsf b)\mathsf d^{-12}[r^{\beta_{\mathrm H}/72}+(a_K/r)^{1/18}]\mathsf b^{-2}\mathsf d^{-6}\mu_3$, and $|\mathsf E|$ is of this form as $r^3\le r^{\beta_{\mathrm H}/72}$. For $r>r_{\mathcal W}$, Lemma~\ref{C-352-26} bounds the left member of (E.1) by $C(\mathsf b)\bar{\mathsf M}^{(24)}[\mathsf d^{-9/2}+1]\mu_3$, while the bracket is $\ge r_{\mathcal W}^{\beta_{\mathrm H}/72}$; the larger of the two quotients so formed is $\mathsf C_{\mathrm R}(g,\mathsf b;\mathsf d)$.

\emph{(E.1) under $(\Gamma3)$.} The stronger bounds stated in (E.1) are Lemma~\ref{C-352-38b}, at the levels $n\ge n_0$, that is at the levels with $a_{K-n}\ge r^{1/2}$, with $\mathsf d'=|c-z|\in[(1-\frac1{64})\mathsf d,6\mathsf d]$. Let $r\le r^{\mathcal W}_\alpha(g,\mathsf b)$. On the part of $\Omega$ with $t'=n\ge n_0$ the measure $\mathbb P$ is absolutely continuous with respect to $\mathsf m_n$, $\mathsf Y_{\Delta3}=\Gamma_n(U_n(W))-\Gamma_K(\mathbf 1)$ and, with $\mathsf c_*=\mathsf c_K-\Delta_K(\mathbf 1)$, $\mathsf Y_\Delta-\mathsf c_*=\Delta_n(U_n(W))-\mathsf c_K$, so that Lemma~\ref{C-352-38b} bounds both by $\varepsilon^{\mathcal W}$ times their natural sizes, $\mathrm{nat}^{\mathcal W}_n\le\alpha^{-4}$; $\kappa_{\mathbb P}(\mathsf Y_\Delta,\mathsf Y_3)=\int(\mathsf Y_\Delta-\mathsf c_*)(\mathsf Y_3-\int\mathsf Y_3d\mathbb P)\,d\mathbb P$, and $|\mathsf c_*|\le C(g,\mathsf b)$, since $\mathbb P(\alpha<\alpha_1)\le\frac12$ for $\alpha_1=(2\bar{\mathsf M}^{(24)})^{-1/24}$ (Lemma~\ref{C-352-25}(e)), while $a_{K-n_0}<Lr^{1/2}\le\alpha_1$, so that on the set $\{\alpha\ge\alpha_1\}$, of positive measure, both $|\mathsf Y_\Delta-\mathsf c_*|\le2\varepsilon^{\mathcal W}\alpha_1^{-4}$ and the bound of $\mathsf Y_\Delta$ by Lemma~\ref{C-352-26} hold.

The parts with $t'<n_0$ have $\alpha<a_{K-n_0}<Lr^{1/2}$, where nothing cancels and the moments of Lemma~\ref{C-352-25}(e) give the factors $r^6$ and $r^9$; for $r>r^{\mathcal W}_\alpha$ Lemma~\ref{C-352-26} bounds the left member while the bracket is $\ge(r^{\mathcal W}_\alpha)^{\beta_{\mathrm H}/72}$. Together the left member of (E.1) is at most $C(g,\mathsf b)\mathsf d^{-5/2}[r^{\beta_{\mathrm H}/72}+(a_K/r)^{1/18}]\mathsf b^{-2}\mathsf d^{-6}\mu_3$.

\emph{(E.3).} 1. For $\mathsf b\in I(g)$ every $(K,L^K\mathsf b)$ is a first-passage datum at $g$ with the label $\mathsf b$. 2. \emph{The pieces.} Cut $\varphi_2$ by the indicators of the cubes of side $\frac r4$; at most $9^3(\mathsf d/r)^3$ of them meet its support, and as $|y_1-y_2|\le3\mathsf d$ the centre $z$ of each has $|c-z|\le5\mathsf d+\frac{\sqrt3}8r\le6\mathsf d$ for $c\in\operatorname{supp}\varphi_1$, so every piece is a $w_3$ of the data. By (E.2), summed over the pieces and integrated against $\varphi_1(c)\,dc$, $\int|\varphi_1|\le\frac{4\pi}3\mathsf d^3\|\varphi_1\|_\infty$:
\[
\Bigl|\mathsf A_{r,K}-\int dc\,\varphi_1(c)\,S^T_2(w^c_T,\varphi_2)\Bigr|\le C(g,\mathsf b)\mathsf d^{-12}\bigl[r^{\beta_{\mathrm H}/72}+(a_K/r)^{1/18}\bigr]\|\varphi_1\|_\infty\|\varphi_2\|_\infty .
\]
3. \emph{The smeared tensor.} $\int dc\,\varphi_1(c)\mathcal O(w^c_T)=T_\lambda(\varphi_1*\chi^{r/8})$ exactly, $\|\varphi_1*\chi^{r/8}-\varphi_1\|_\infty\le\frac r8\operatorname{Lip}\varphi_1$, and the two-point bound of Lemma~\ref{C-352-26} in the sup norms (supports at a distance $\ge\frac12\mathsf d$, after cutting $\varphi_1*\chi^{r/8}-\varphi_1$ into pieces) gives $|\langle T_\lambda(\varphi_1*\chi^{r/8}-\varphi_1);\mathcal O(\varphi_2)\rangle|\le C(\mathsf b,\mathsf d)\,r\operatorname{Lip}\varphi_1\|\varphi_2\|_\infty$; this is (i), with $r_\sharp$ the least of the thresholds of (E.2) and of the data. Steps 2 and 3 read the datum through (E.2) and Lemma~\ref{C-352-26} at it; so (i) holds, with the same proof, at every first-passage datum $(K,\beta)$ at $g$ whose label $\mathsf b'=a_K\beta$ lies in $(g^{-2},\bar{\mathsf b}]$, with $\mathsf A_{r,K}$ formed at $(K,\beta)$ with the prefactor $\frac{2\pi}3\mathsf b'r^3$ and with the numbers $r_\sharp$ and $\mathsf C(\mathsf d)$ formed at $\mathsf b'$. 4.

\emph{(ii).} At a fixed $r$, $\mathsf A_{r,K}\to\mathsf A_r$ as $K\to\infty$: for each $c$, $S^T_{3;K,\ell_g,\beta_K}(f^c_+,f^c_-,\varphi_2)\to S^{(\ell_g,\mathsf b)}_3(f^c_+,f^c_-,\varphi_2)$ by Theorem~\YMMref{A}{5} (Lipschitz functions with disjoint supports), with a bound uniform in $c$ and $K$ (Lemma~\ref{C-352-26} and Proposition~\ref{C-352-24}), and dominated convergence in $c$. Hence, by (i),
\[
\limsup_{K,K'\to\infty}\bigl|\langle T_\lambda(\varphi_1);\mathcal O(\varphi_2)\rangle_{K,\ell_g,\beta_K}-\langle T_\lambda(\varphi_1);\mathcal O(\varphi_2)\rangle_{K',\ell_g,\beta_{K'}}\bigr|\le2\mathsf C(\mathsf d)r^{\beta_{\mathrm H}/72}\mathsf b^{-2}\|\varphi_1\|_\infty\|\varphi_2\|_\infty,
\]
and $r\to0$ shows that the sequence is a Cauchy sequence; $s(\frac r2)\le K$ and $128a_K\le r$ hold for all large $K$, so the conditions of (i) on the cutoff hold along the sequence. 5. \emph{(iii)}: $K\to\infty$ in (i). As $g$ satisfies $(\Gamma)$, the limit points of (C.2) along data with $\beta_j=L^{K_j}\mathsf b$ all give the form $t_\lambda$ on these pairs.

\emph{(E.4).} Let $\mathsf b_j=a_{K_j}\beta_j\to\mathsf b\in I(g)$, and let $\mathsf A^{(j)}_r$ be $\mathsf A_{r,K}$ formed at $(K_j,\beta_j)$ with $\frac{2\pi}3\mathsf b_jr^3$. As $g^{-2}<\mathsf b<\bar{\mathsf b}$, for $j$ large the labels lie in a compact subset of $(g^{-2},\bar{\mathsf b}]$. The bound (i), in its form at every label of $(g^{-2},\bar{\mathsf b}]$ (step 3 above; a label $\mathsf b_j$ may lie outside $I(g)$), holds at $(K_j,\beta_j)$ with the prefactor $\frac{2\pi}3\mathsf b_jr^3$, and by Lemma~\ref{C-352-41} its constants and thresholds are uniform in $j$. At a fixed $r$, $\mathsf A^{(j)}_r\to\mathsf A_r$: the prefactor $\frac{2\pi}3\mathsf b_jr^3$ tends to $\frac{2\pi}3\mathsf br^3$; for each $c$, $S^T_{3;K_j,\ell_g,\beta_j}(f^c_+,f^c_-,\varphi_2)\to S^{(\ell_g,\mathsf b)}_3(f^c_+,f^c_-,\varphi_2)$ by Theorem~\YMMref{A}{5}, one limit along all data whose labels tend to $\mathsf b$; and the bounds of Lemma~\ref{C-352-26} on $|S^T_3-\mathsf J_3|$ and on $|\mathsf J_3|$, summed over the pieces of $\varphi_2$, are uniform in $c$ and in $j$, as $\mathsf b_j\le\bar{\mathsf b}$, so that dominated convergence applies in $c$. So
\[
\begin{multlined}
\limsup_j\bigl|\langle T_\lambda(\varphi_1);\mathcal O(\varphi_2)\rangle_{K_j,\ell_g,\beta_j}-t_\lambda(\varphi_1,\varphi_2)\bigr|\le\limsup_j\bigl(|\langle T_\lambda;\mathcal O\rangle_j-\mathsf A^{(j)}_r|+|\mathsf A^{(j)}_r-\mathsf A_r|\bigr)+|\mathsf A_r \\
-t_\lambda|\le C\,r^{\beta_{\mathrm H}/72},
\end{multlined}
\]
and $r\to0$.
\end{proof}

\begin{proof}[Proof of Theorem~\ref{C-352-8}]
(a) Proposition~\ref{C-352-24}(c), Lemma~\ref{C-352-26} at $n=4$, and Proposition~\ref{C-352-33}(a) at $n=4$. (b) The assumptions of (b) give $\operatorname{supp}f_a\subset\mathbb B(x_a,\rho)$, with the balls $\mathbb B(x_a,\rho)$ at least $r$ apart in each pair and at least $\mathsf d$ apart across, as Lemma~\ref{C-352-27}(b) asks; the $w_a=f_a\circ x$ form two pairs at $(\rho,r,\mathsf d)$ with the centres $x_a$, with the bounds $|x_1-x_2|,|x_3-x_4|\le8r$ and $|x_a-x_b|\le18\mathsf d$, $a\le2<b$, of the data; and $8a_K\le\rho$ by $s(\rho)\le K$ (Lemma~\ref{C-352-11}(b)) and $\mathsf d\le\mathsf d^+$, as (b) notes. The bounds of (a) on $|S^T_4-\mathsf J_4|$ and on $|\mathsf J_4-\frac32\sum_\varsigma\operatorname{Tr}\prod_l(w_{\varsigma(l)}\mathsf b^{-1}C^{\mathrm{free}}_K)|$, the identity $\frac32\sum_\varsigma\operatorname{Tr}\prod_l(w_{\varsigma(l)}\mathsf b^{-1}C^{\mathrm{free}}_K)=3\mathsf b^{-4}\cdot\frac12\sum_\varsigma\operatorname{Tr}\prod_l((f_{\varsigma(l)}\circ x)C^{\mathrm{free}}_K)$ and Lemma~\ref{C-352-27}(b) give (b). (c) For a point of $\mathcal D_4$ take the open balls of radius $\frac14r_*$ about the $x_a$, $r=\frac12r_*$, $\mathsf d=\frac12\mathsf d_*$, and bumps $f_a$ supported in balls $\mathbb B(x'_a,\rho)$ of radius $\rho$ that lie inside them: these balls are $\ge r$ apart in a pair and, since $r_*\le\mathsf d_*$, $\ge\mathsf d$ apart across, as (b) asks, and the $f_a\circ x$ form two pairs at $(\rho,r,\mathsf d)$ with the centres $x'_a$ (in a pair the centres are $\le\frac52r_*=5r$ apart, across $\le11\mathsf d$), all inside a ball of radius $6\mathsf d_*=\frac12\mathsf d^+$; then (b), the limit $K\to\infty$ at the label $\mathsf b$ (Theorem~\YMMref{A}{5}), Lemma~\ref{C-352-28} at $p=4$ with $(\frac3\pi)^4<1$, and $r^{-6}\mathsf d^{-6}=4096\,r_*^{-6}\mathsf d_*^{-6}$.

The rings $(1234)$, $(1243)$ of $k_4$ have the size $r_*^{-6}\mathsf d_*^{-6}$, the ring $(1324)$ the size $\mathsf d_*^{-12}$.
\end{proof}

\subsection{Scope}

Every statement of this section is made at short distances, on one torus or in states at a cutoff and in their limits, and the constants that depend on the side grow with it. The side of the torus enters through conditions that ask it to be large against the lengths in play, and through constants that grow with it, among them: the conditions that the supports lie in a ball small against it; the conditions that the boxes and the cubes of the local gauges do not wrap around the torus, namely $4\rho+4$ at most the side in Lemma~\ref{C-352-18}, in the units of the level, the case distinction of step 5 of the proof of Proposition~\ref{C-352-19} on whether its box fits in the torus, which reads $\ell\ge8$, $(4\mathsf h_i+4)a_{K-i-1}\le\ell_g$ in \textup{($\alpha$2)} and in Definition~\ref{C-352-2}(6), and $12a_{K-n}\le\ell$, $(8\varrho+2)a_{K-n}<\ell$ for the regular gauge on a cube, read at $\varrho=4$ in Lemma~\ref{C-352-38} and at $\varrho=3$ up to $n=K$ in the proof of the bounds of (E.1) under $(\Gamma3)$, where $26<\ell_g$; the bound $L^m\ge LMy_K$ of the sides of $\mathfrak L(g)$, which gives $\ell_g\ge338$, in the tails of Lemma~\ref{C-352-14}, in Proposition~\ref{C-352-19}, in Lemma~\ref{C-352-40} and for $26<\ell_g$; the distance $\frac14\ell_g\ge6.5$ of two mirror copies in Lemma~\ref{C-352-32}; the membership $\ell\in\mathfrak L(g)$ wherever the objects of the theory are entered, as in the frame of Lemma~\ref{C-352-11}; the terms in $\mathsf d/\ell$ of the comparison of $C^{\mathrm{free}}_K$ with $\mathbb K_3$; the moments of Lemma~\ref{C-352-25}(e); and the constant $\mathsf Q_0$ of Lemma~\ref{C-352-30}.

In the bond forms the members that carry the renormalisation group, the remainder and the induction of the proof of Proposition~\ref{C-352-37}, read the distances through the floor $\min\{\frac67\mathsf d',t_{\mathrm N}(g)\}$; the largest distance $\frac87\mathsf d'$ enters the one explicit member, the comparison of the sum of the slices with $C^{\mathrm{free}}_K$ (Lemma~\ref{C-352-22}(b)), which is polynomial in that distance and a part of $\epsilon^{\mathrm{ev}}_\sharp$. Theorems~\YMMref{AJ}{673} and~\YMMref{AD}{607} enter through the two-sided stability bound of the torus measures and the chessboard bound of Lemma~\ref{C-352-30}, in (A.4), (A.5), (C.1)(iii), (C.4) and for the states of Theorem~\ref{C-352-8}(a); Theorem~\ref{C-352-7}, like every statement of this section, uses Theorems~\YMMref{D}{8}, \YMMref{AJ}{673} and~\YMMref{AD}{607}.

The crossing functionals of Lemma~\ref{C-352-39} are joint cumulants of the steps given the block fields of a scale below $\mathsf d$. The comparisons with $\mathsf t^\circ$ and with the kernels $t_\lambda$ and $t_n$ are additive, in the unit $\mathsf b^{-2}\mathsf d^{-6}$; $\mathsf t^\circ(n,X)$ vanishes on the cone $3(n\cdot\hat X)^2=1$. On a torus a kernel is a class modulo null sets, while at fixed $(c,z)$ the triples $(c\pm\frac r2n,z)$ form a null set; there (C.1)--(C.3) and (D.3) are statements almost everywhere, or through means, tents and cones of directions. For the laws on $\mathbb R^3$ satisfying the continuity
hypotheses in (A.4)(ii) and (C.4), including the periodic-Gibbs
limits and their canonical components, $\mathsf s_3$ and
$\mathsf s^T_2$ are continuous, and the kernels $\tau^\sigma_{\lambda\mu}$ and $t^\sigma_\lambda$ are classes modulo null sets: (C.1)(iii) holds at every point and for every unit vector $n$; (C.4)(c) holds for $\mathsf s_3$ itself, without means, for every $0<r\le\frac12\mathsf d$ and every $\lambda$, at every $(c,z)$ off one null set fixed before $r$ is named, and at every $(c,z)$ with the limit forms at the tents $\chi^\rho_c,\chi^\rho_z$, $0<\rho<\frac18\mathsf d$, in the place of $t^\sigma_\lambda(c,z)$ and $21\rho/\mathsf d$ added in the bracket.

The kernels $\tau_{\lambda\mu}$, $t_\lambda$, $t_n$ of (C.2)--(C.4) and (D.2) belong to limit points, each chosen once for all pairs of bumps; at a label of $I(g)$, for $g\le\gamma_{\mathrm R}$ and on the least torus, the forms $\langle T_\lambda(\varphi_1);\mathcal O(\varphi_2)\rangle$ of (E.3) have one limit along all cutoffs.

\section{Plane fields, stress and anisotropic variations}\label{C-360-1}

We compare the curvature bilinears with plane insertions and the stress tensor obtained from anisotropic changes of the lattice action. Uniform estimates permit the removal of the cutoff; the collapse and remainder estimates identify the conserved stress insertion.

The Ward identities retain their boundary terms and quantitative hypotheses. The two-scale comparisons use the results of Section~\ref{C-352-1}; they do not identify arbitrary coincident-point products without a specified extension.

\subsection{Notation and hypotheses}

\emph{Model, constants and data.} We work in the common real and sourced construction setting of Proposition~\YMMref{69.1}{595} and Section~\YMMref{82}{665}: the Wilson action of ${\rm SU}(2)$ on $T_\ell=\mathbb R^3/\ell\mathbb Z^3$, $\ell=2L^m$, at the spacing $a=a_K=L^{-K}$, $L\ge13$ odd; $a_p=1-\frac12{\rm tr}\,\partial U(p)$ and $A=\sum_pa_p$; $\langle\cdot\rangle_{K,\ell,\beta}$ and $\langle\cdot\,;\dots;\cdot\rangle_{K,\ell,\beta}$ are the means and the joint cumulants of $Z_{K,\ell,\beta}^{-1}e^{-\beta A}dU$, and $S^T_{n;K,\ell,\beta}$ are the connected functions of $\mathcal O$. For a plaquette $p$, $x(p)$ is its barycentre, $\lambda(p)$ its normal direction and $P(p)$ the set of its two directions; for a bond $b$, $x_b$ is its midpoint and $\mu_b$ its direction.

The constants are those of the list of constants and conditions among the construction results, fixed in its order, $\gamma$ last. For a reference coupling $0<g\le\gamma$, $\mathfrak L(g)$ is the set of the sides at $g$ (equation~\YMMref{(1528)}{615}), $\ell_g$ the least of them, and $\mathfrak D(g)$ the set of the first-passage data $(K,\beta)$ at $g$ (Corollary~\YMMref{74.2}{616}), with the label $\mathsf b=a_K\beta$; $X_g=L(g^{-2}+B_\sharp)$, $\bar B=B_\sharp L/(L-1)$, $I(g)=(g^{-2}+\bar B,X_g-\bar B]$. Every label of $\mathfrak D(g)$ lies in $[\gamma^{-2},X_g+\bar B]$; $MR(X_g+\bar B)\le L^m$ on every side $2L^m$ in $\mathfrak L(g)$, $R(x)$ being the least power of $L$ above $(\log x)^r$ with the exponent $r$ among the construction results; and for $\mathsf b\in I(g)$ every pair $(K,L^K\mathsf b)$, $K\in\mathbb N$, is a first-passage datum at $g$.

The bounds $\gamma^{\rm sd}_W\le\gamma^{\rm ref}_W=x_W^{-1/2}\le\gamma$ on the reference coupling (the choices in equations~\YMMref{(2326)--(2327)}{845}), the threshold $x_{\rm J}$, the radii $\alpha_0,\alpha_1$ of the cores, the numbers $\mathfrak p_0$, $A_0$, $A_1$, $c'$, $\kappa$, $M$, $C_{\rm lat}$ and the function $B(\cdot)$ are those among the construction results, and so is the depth function: $s(\rho)$ is the least integer $s\ge0$ with $L^s\rho\ge C_W(s)$, where $C_W(s)=C_{W,0}A_0^2Y_s^{2\mathfrak p_0}$ and $Y_s=(s+1)\log L+\log(g^{-2}+\bar B)$. The choice preceding equation~\YMMref{(2326)}{845} satisfies every inequality defining $x_{\rm J}$ in \eqref{eq:weighted-history-threshold}. Since $x_{\rm J}$ is the least such threshold, $x_W\ge\max\{\gamma^{-2},x_{\rm J}\}$, and hence $g^{-2}\ge x_{\rm J}$ for $0<g\le\gamma^{\rm ref}_W$. $\mathcal G^{\rm per}_{K,\beta}$ is the set of the periodic Gibbs states at $(a_K,\beta)$; a \emph{state} at $(K,\beta)$ is one of them or the measure of a torus with a side in $\mathfrak L(g)$; for a state $\nu$ we write $A^\nu[F]=F-\nu(F)$.

We write $p_\ell=\ell^{-3}\log Z_{K,\ell,\beta}$ for the pressure of $T_\ell$ per unit of physical volume, formed with the plain partition function $Z_{K,\ell,\beta}$ of the measure above (Section~\YMMref{86}{680}); at zero source Theorems~\YMMref{AJ}{673} and~\YMMref{AD}{607} write $Z_{K,\ell,\beta}=e^{E_{(K)}[\ell]}\widetilde Z_\ell(0)$, with the vacuum constant $E_{(K)}[\ell]$ of the side $\ell$ and the normalised partition function $\widetilde Z_\ell(0)$ among the construction results. For $(K,\beta)\in\mathfrak D(g)$ and $\ell\ge\ell'$ in $\mathfrak L(g)$, $0\le p_{\ell'}-p_\ell\le G_g$, with the number $G_g=C_{\rm S}+C_g+C_{\rm gap}$ among the construction results.

The lower bound is reflection positivity: $Z_{K,\ell,\beta}\le Z_{K,\ell',\beta}^{(\ell/\ell')^3}$, Corollary~\ref{IV-18-18}(e). The upper bound follows from $|E_{(K)}[\ell]\ell^{-3}-E_{(K)}[\ell']\ell'^{-3}|\le C_{\rm gap}$, $\widetilde Z_{\ell'}(0)\le e^{C_{\rm S}\ell'^3}$ and $\widetilde Z_\ell(0)\ge e^{-C_g\ell^3}$ (Propositions~\YMMref{74.4}{617} and~\YMMref{74.5}{617}, through Corollary~\ref{IV-18-18}, Proposition~\YMMref{89.2}{692} and Lemma~\ref{lem:canonical-source-system}, which compare two sides; $C_{\rm S}$ and $C_{\rm gap}=2C_{\rm tt}C'_{\rm tt}E_0$ are the numbers of Section~\ref{sec:source-system-complements}, and $C_g$ that of \eqref{eq:setup-pressure-constants}). Distances are Euclidean unless they are called sup distances; ${\rm Lip}_\infty$ is the Lipschitz constant for the sup distance, $\|f\|_\sharp=\max\{\sup|f|,40M\,{\rm Lip}f\}$, $\|\eta\|_{C^2}$ is the sum of the suprema of $|\eta|$, $|\nabla\eta|$, $|\nabla^2\eta|$, and $\|\varphi\|_{C^{1,1}}=\max\{\sup|\varphi|,\sup|\nabla\varphi|,{\rm Lip}\,\nabla\varphi\}$.

Constants in the statements depend only on the constants among the construction results and on the arguments shown, never on $K$, $\beta$, the state or the test functions.

\begin{lemma}[A bound on separated torus cumulants]
\label{lem:torus-cumulant-bound}
Fix the common construction choice and $0<g\le\gamma$. Put
$\delta_e=\kappa-\log9216>0$ and
\[
 S(\kappa)=\frac{2^{12}}{1-e^{-\delta_e}},\qquad
 C'_\Sigma=\frac{27}{26}\,5^3+
       \left(\frac6{e\delta_e}\right)^3
                       \frac{e^{2\delta_e}}{1-e^{-\delta_e}},
 \qquad c_V=C_{\rm S}(2+\log g^{-2})+1.
\]
For $n\ge2$ and $d>0$, define
\[
 C_n(d)=\left(\frac n{2\varrho}\right)^n
    \left[4^n\log2+81\sqrt3\,C_{\rm S}S(\kappa)C'_\Sigma d^{-3}\right].
\]
At every first-passage datum $(K,\beta)$ at $g$, on every
$\ell\in\mathfrak L(g)$, real Lipschitz functions $f_1,\ldots,f_n$
with two supports at Euclidean distance at least $d$ satisfy
\begin{equation}\label{eq:torus-cumulant-bound}
 |S^T_{n;K,\ell,\beta}(f_1,\ldots,f_n)|
 \le C_n(d)e^{nc_V\ell^3}\prod_i\|f_i\|_\sharp.
\end{equation}
The constants are independent of $K$, $\beta$, $\ell$ and the
functions; $C_n(d)$ depends only on $n,d,\varrho$ and the fixed
construction constants, while $c_V$ may also depend on $g$.
If each support is in a sup-ball of radius $R>0$ and
$\ell\ge16n(R+1)$, Proposition~\YMMref{89.4}{693} instead
gives a bound independent of $\ell$.
\end{lemma}
\begin{proof}
Proposition~\YMMref{88.3}{688}, with
$Q_\ell=\exp[C_{\rm S}(2+\log g^{-2})\ell^3]$, gives
\[
 |S^T_n(f)|\le\left(\frac n{2\varrho}\right)^n
 \left[\log2\,(2Q_\ell+2)^n+
  81\sqrt3 C_{\rm S}S(\kappa)C'_\Sigma\ell^3d^{-3}\right]
                       \prod_i\|f_i\|_\sharp.
\]
Here $C'_\Sigma$ is the scale-sum constant in equation
\YMMref{(1776)}{686}. Since $Q_\ell\ge1$,
$(2Q_\ell+2)^n\le4^nQ_\ell^n$; and
$\ell^3\le e^{n\ell^3}$. Absorbing these two factors proves
\eqref{eq:torus-cumulant-bound}. Symmetry of a cumulant allows
the separated pair to occupy any two positions. The per-ball
assertion is exactly Proposition~\YMMref{89.4}{693}, under the
additional support and side assumptions stated above.
\end{proof}

In this section $c_V$ denotes the displayed enlarged constant.
Every use below is through \eqref{eq:torus-cumulant-bound}; no
identity with a pressure or a vacuum-energy coefficient is used.


\emph{Plane fields and rings.} A \emph{weight} is a real or complex function $w$ of the plaquettes; $\mathcal O(w)=a^{-1}\sum_pw(p)a_p$, $\mathcal O(f)=\mathcal O(f\circ x)$, $\mathcal O_\lambda(f)=a^{-1}\sum_{\lambda(p)=\lambda}f(x(p))a_p$, so that $\mathcal O=\mathcal O_1+\mathcal O_2+\mathcal O_3$, and $\check D_m=\mathcal O-2\mathcal O_m$. $C^{\rm free}_K(p,q)$ is the matrix of the orthogonal projection onto the exact lattice curvatures, and
\[
\mathcal K^{(n)}_K(w_1,\dots,w_n)=\tfrac12\sum_{\varsigma}\sum_{p_1,\dots,p_n}\prod_aw_a(p_a)\prod_{i\in\mathbb Z/n}C^{\rm free}_K(p_{\varsigma(i)},p_{\varsigma(i+1)}),
\]
the sum over the $(n-1)!$ cyclic orders $\varsigma$; $\mathcal K^{(n)}_{K,\ell'}$ is the same sum on $T_{\ell'}$. With $\mathbb K_3(u)=(3\hat u\hat u^{\rm T}-\mathbf 1)/(4\pi|u|^3)$, the continuum ring $\mathcal K^{(n)}_\lambda(\phi)$ is the same expression with $\int\prod_a\phi_a(y_a)\,dy_a$ for the sum over the plaquettes and $\mathbb K_3(y_{\varsigma(i)}-y_{\varsigma(i+1)})_{\lambda_{\varsigma(i)}\lambda_{\varsigma(i+1)}}$ for $C^{\rm free}_K$. We put $\mathcal K_{\lambda\mu}(\phi_1,\phi_2)=\iint\phi_1(x)\phi_2(y)[\mathbb K_3(x-y)_{\lambda\mu}]^2\,dx\,dy$, $k_3(x,y,z)={\rm tr}[\mathbb K_3(x-y)\mathbb K_3(y-z)\mathbb K_3(z-x)]$ and, for a unit vector $n$, $\mathsf t^\circ(n,X)=\frac3{32\pi^2}\mathsf b^{-2}|X|^{-6}(3(n\cdot\hat X)^2-1)$.

Since $\mathbb K_3(X)^2=(16\pi^2|X|^6)^{-1}(3\hat X\hat X^{\rm T}+\mathbf 1)$, we have $\frac32\mathsf b^{-2}\sum_{\lambda,\mu}[\mathbb K_3(X)_{\lambda\mu}]^2=\frac9{16\pi^2}\mathsf b^{-2}|X|^{-6}$, the free two-point kernel of $\mathcal O$, and $\frac32\mathsf b^{-2}\sum_\mu\{[\mathbb K_3(X)_{\lambda\mu}]^2-\frac13\sum_{\lambda'}[\mathbb K_3(X)_{\lambda'\mu}]^2\}=\mathsf t^\circ(e_\lambda,X)$. For plaquettes $p\ne q$ and $u$ a least lift of $x(p)-x(q)$, equations~\YMMref{(2245)--(2246)}{824} bound the lattice kernel by $|C^{\rm free}_K(p,q)|\le(1+3C_{\rm lat})a^3|u|^{-3}$ and compare it with $\mathbb K_3$, $|\epsilon_p\epsilon_qC^{\rm free}_K(p,q)-a^3\mathbb K_3(u)_{\lambda(p)\lambda(q)}|\le C_{\rm lat}a^3(a^2|u|^{-5}+\ell^{-3})$, $\epsilon_p=\pm1$ the orientations (\YMMref{(2246)}{824} and \YMMref{(2245)}{824} of equations~\YMMref{(2245)--(2246)}{824}).

Its differences along one plane obey
\[
\begin{multlined}
|C^{\rm free}_K(p,q)-C^{\rm free}_K(p',q)|\le C_\nabla\,a^3|x(p)-x(p')|\,|u|^{-4}\qquad\text{for parallel, equally oriented }p, \\
p'\text{ with }|x(p)-x(p')|\le\tfrac14|u|\text{ and }|u|\ge8a,
\end{multlined}
\]
$C_\nabla=256\sqrt3\,C_{\rm h}(1+3C_{\rm lat})$, $C_{\rm h}$ the constant of the interior gradient estimate for discrete harmonic functions on $\mathbb Z^3$. Indeed the Laplacian $dd^*+d^*d$ of the cubic lattice acts on $2$-forms plane by plane as the scalar lattice Laplacian, and it commutes with the projection; so it carries the column $C^{\rm free}_K(\cdot,q)$, the projection of the unit form $\delta_q$, to $dd^*\delta_q$, which vanishes at every plaquette that shares no bond with $q$, in particular at every plaquette whose barycentre is more than $a$ away from $x(q)$. On each plane the column is therefore discrete harmonic at the distance $\ge2a$ from $x(q)$. A lattice path from $x(p)$ to $x(p')$ of length $\le\sqrt3|x(p)-x(p')|$ keeps the distance $\frac12|u|$ from $x(q)$; on the balls of radius $\frac14|u|$ about its points the first bound bounds the column by $4^3(1+3C_{\rm lat})a^3|u|^{-3}$, and the gradient estimate on these balls bounds each step by $4^4C_{\rm h}(1+3C_{\rm lat})a^4|u|^{-4}$.

\emph{Lattice calculus.} $\mathbf w(V)=(V-V^{-1})/2i$ and $\langle X,Y\rangle=\frac12{\rm tr}XY$. For a positively oriented bond $b$ and a plaquette $p\ni b$, $W_{b,p}$ is the holonomy of $p$ from the initial point of $b$, along $b$; $p(b,\pm\rho)$, $\rho\ne\mu_b$, are the plaquettes through $b$ in the directions $\pm\rho$; $\mathbf u_{b,\rho}=-\frac12[\mathbf w(W_{b,p(b,\rho)})-\mathbf w(W_{b,p(b,-\rho)})]$, $\mathbf u_{b,\mu_b}=0$, the displacement of $b$, and $\mathfrak d_{b,p,\rho}=\langle\mathbf u_{b,\rho},\mathbf w(W_{b,p})\rangle$, the elementary densities; $\mathbf u_{b,\rho}$ is the half sum of the field strengths of the two plaquettes at $b$ on the sides $\pm e_\rho$, and $G_\rho(b)=\mathbf w(W_{b,p(b,\rho)})+\mathbf w(W_{b,p(b,-\rho)})$ is the transported difference of these field strengths. With $p_\pm=p\pm ae_\rho$,
\[
\begin{multlined}
\mathfrak r_\rho(p)=\sum_{b\subset\partial p}\mathfrak d_{b,p,\rho}-\tfrac12(a_{p_+}-a_{p_-}),\qquad\mathfrak s_{l\rho}(p)=\sum_{b\subset\partial p}\frac{(x_b-x(p))_l}a\,\mathfrak d_{b,p,\rho},\qquad\mathfrak r^\sharp_\rho(p)=\mathfrak r_\rho(p) \\
-\Psi_\rho(p)+\Psi_\rho(p_-),
\end{multlined}
\]
the orbital remainder, the spin densities and the cubic remainder, where the orbital potential is $\Psi_\rho(p)=\frac12a_{p\cup p_+}-a_p-a_{p_+}+\frac12a_pa_{p_+}$ for $\rho\in P(p)$, $a_{p\cup p_+}$ being the plaquette term of the $1\times2$ rectangle, and $\Psi_\rho(p)=-\frac1{16}\sum_y\langle Y_y,Y_y\rangle$, $Y_y=\mathbf w(W^y_p)-{\rm Ad}(U(y,y+ae_\rho))\mathbf w(W^{y+ae_\rho}_{p_+})$, for $\rho=\lambda(p)$, the sum running over the corners $y$ of $p$ and $W^y_p$ being the holonomy of $p$ from $y$. For a weight $v$ and a vector field $\eta$ the stress insertion is
\[
\Theta_v[\eta]=\sum_pv(p)\,a^{-1}\sum_{b\subset\partial p}\sum_\rho\bigl(\eta^\rho(x_b)-\eta^\rho(x(p))\bigr)\mathfrak d_{b,p,\rho},
\]
and $\Theta_\beta[\eta]$ is its value at $v\equiv\beta$. A \emph{field} is a triple $\eta$ of real functions on the barycentres; $\mathsf H(\eta)=\max_\rho\sup|\eta^\rho|+\max_\rho{\rm Lip}_\infty\eta^\rho$, and
\[
\mathfrak K^{\rm v}[\eta]=a^{-1}\sum_p\sum_\rho\eta^\rho(x(p))\,\mathfrak r_\rho(p),\qquad\mathfrak K^\sharp[\eta]=a^{-1}\sum_p\eta^{\lambda(p)}(x(p))\,\mathfrak r^\sharp_{\lambda(p)}(p)
\]
are the orbital insertion and the cubic insertion. We write ${\rm div}^s\eta(y)=\sum_\rho(2a)^{-1}[\eta^\rho(y+ae_\rho)-\eta^\rho(y-ae_\rho)]$, $\nabla^+_\rho\varphi(y)=a^{-1}[\varphi(y+ae_\rho)-\varphi(y)]$, $\nabla^c_\rho\varphi(y)=a^{-1}[\varphi(y+\frac a2e_\rho)-\varphi(y-\frac a2e_\rho)]$, $e^\eta(p)=[\nabla^c_\mu\eta^\mu+\nabla^c_\nu\eta^\nu](x(p))$ with $\{\mu,\nu\}=P(p)$, and $\mathrm{St}_K[\eta]=\Theta_\beta[\eta]-\mathsf b\,\mathcal O({\rm div}^s\eta)$. With $v_\eta$ the derivation that displaces $U$ along $\eta$ and ${\rm div}\,v_\eta$ its divergence, the identities (Lemma~\ref{lem:ward-general-calculus}, with the integration by parts \ref{lem:ward-general-calculus}) are $v_\eta A(\beta,\cdot)=\mathfrak K^{\rm v}[\beta\eta]-\mathsf b\,\mathcal O({\rm div}^s\eta)+\Theta_\beta[\eta]$ and ${\rm div}\,v_\eta=-\frac32a\,\mathcal O(e^\eta)$. So the null term $N_\beta[\eta]=v_\eta A(\beta,\cdot)-{\rm div}\,v_\eta$ of Lemma~\ref{C-360-8}(a) is $\mathrm{St}_K[\eta]+\mathfrak K^{\rm v}[\beta\eta]+\frac32a\,\mathcal O(e^\eta)$; and since, identically in $U$, $\mathfrak K^{\rm v}[\beta\eta]=-\beta\sum_{p,\rho}(\nabla^+_\rho\eta^\rho)(x(p))\Psi_\rho(p)+\beta\mathfrak K^\sharp[\eta]$ (Lemma~\ref{C-360-8}(b)), $N_\beta[\eta]=\mathbb T_\beta[\eta]+\beta\mathfrak K^\sharp[\eta]$ with the \emph{stress functional}
\[
\mathbb T_\beta[\eta]=\mathrm{St}_K[\eta]-\beta\sum_{p,\rho}(\nabla^+_\rho\eta^\rho)(x(p))\,\Psi_\rho(p)+\tfrac32a\,\mathcal O(e^\eta),
\]
which reads $\eta$ through its differences alone and vanishes at a constant field; the cubic insertion $\beta\mathfrak K^\sharp[\eta]$, which completes it to the null term, is its \emph{cubic part}.

\emph{Tensor entries.} $V={\rm Sym}_3(\mathbb R)$ and $|H|_\infty=\sup_x\max_{l,r}|H_{lr}(x)|$. For $H\in C_c(\mathbb R^3;V)$ we put $\mathcal T_K(H)=\mathcal O(w_H)-a^{-1}\sum_bH_{\lambda_b\lambda'_b}(x_b)\langle\mathbf u_{b,\lambda_b},\mathbf u_{b,\lambda'_b}\rangle$, $w_H(p)=H_{\lambda(p)\lambda(p)}(x(p))$, $\lambda_b<\lambda'_b$ being the directions normal to $b$; formally $\mathcal T_K(H)=\frac12\int{\rm Tr}(HT)$ with $T_{lr}={\rm tr}B_lB_r$, and $\mathcal T_K(f\mathbf 1)=\mathcal O(f)$, $\mathcal T_K(f(\mathbf 1-2E_{mm}))=\check D_m(f)$. We put $H^{\rm cons}[\eta]=({\rm div}\,\eta)\mathbf 1-\nabla\eta-(\nabla\eta)^{\rm T}$ and $\Delta_K[\eta]=\mathrm{St}_K[\eta]-\mathsf b\,\mathcal T_K(H^{\rm cons}[\eta])$.

At smooth fields the density of $\mathrm{St}_K[\eta]$ against $\partial_l\eta^r$ is $\mathsf b[\frac12\delta_{lr}|B|^2-\langle B_l,B_r\rangle]$, so that of $\mathsf b\,\mathcal T_K(H^{\rm cons}[\eta])$ is the same and the classical density of $\Delta_K[\eta]$ vanishes. A \emph{tensor limit} at $g$ along $(K_j,\beta_j,\rho_j)$, with $(K_j,\beta_j)\in\mathfrak D(g)$, $K_j\to\infty$ and $\rho_j\in\mathcal G^{\rm per}_{K_j,\beta_j}$, is a functional with $\sigma(H_1,\dots,H_n)=\lim_j\rho_j(\prod_aA^{\rho_j}[\mathcal T_{K_j}(H_a)])$ for all $H_a\in C_c(\mathbb R^3;V)$ with pairwise disjoint supports; the labels need not converge. The stress entries are $T_{lr}(\varphi)=\beta\sum_p\varphi(x(p))[\mathfrak s_{lr}(p)-\delta_{lr}a_p]$.

\emph{Rotations.} For $\omega\in\mathfrak{so}(3)$, $|\omega|=\max|\omega_{\rho\lambda}|$, and $\omega y$ is taken in the chart $(-\frac\ell2,\frac\ell2)^3$; $\nabla^s_\omega\varphi(y)=\sum_\rho(\omega y)_\rho(2a)^{-1}[\varphi(y+ae_\rho)-\varphi(y-ae_\rho)]$, and $\mathfrak S_\omega[v]=\sum_pv(p)\sum_{\rho,\lambda}\omega_{\rho\lambda}\mathfrak s_{\lambda\rho}(p)$ is the spin insertion. A cut-off function at $K$ is a real function $\chi$ on $\mathbb R^3$ that vanishes outside $(-\frac\ell2+2a,\frac\ell2-2a)^3$, regarded as a function on $T_\ell$ (Section~\YMMref{112}{781}); it has the margin $\delta_0\ge2a$ if it vanishes outside $(-\frac\ell2+\delta_0,\frac\ell2-\delta_0)^3$; it covers $f_1,\dots,f_n$ if $\chi=1$ at the barycentre and at the bond midpoints of every plaquette at which some $f_i\ne0$; ${\rm Var}(\chi)$ is the closed set of the points near which $\chi$ is not constant. The cut-off variation density and the flux insertion of Section~\YMMref{112}{781} are
\[
\begin{multlined}
h_{\omega,\chi}(p)=a^{-1}\sum_{b\subset\partial p}\sum_\rho\bigl[\chi(x_b)-\chi(x(p))\bigr](\omega x_b)_\rho\,\mathfrak d_{b,p,\rho},\qquad\mathfrak B_{\omega,\chi}=\mathsf b\,\mathcal O(\nabla^s_\omega\chi)-\beta\sum_ph_{\omega,\chi}(p) \\
-\tfrac32a\,\mathcal O(v_\chi),
\end{multlined}
\]
with $v_\chi(p)=\sum_{\rho\in P(p)}(\omega x(p))_\rho\nabla^c_\rho\chi(x(p))$; $\mathfrak B_{\omega,\chi}$ depends on $U$ near ${\rm Var}(\chi)$ only, and $\beta\sum_ph_{\omega,\chi}(p)=\Theta_\beta[\chi\,\omega y]-\beta\,\mathfrak S_\omega[\chi]$. With $\mathcal O(f_\cdot)=\mathcal O(f_1);\dots;\mathcal O(f_n)$, the bulk term and the cubic remainder are
\[
\mathfrak U_{\omega,\chi}(f)=-\beta\langle\mathfrak S_\omega[\chi]+\mathfrak K^{\rm v}[\chi\,\omega y];\mathcal O(f_\cdot)\rangle+a^{-1}\sum_i\langle\mathfrak S_\omega[f_i]+\mathfrak K^{\rm v}[f_i\,\omega y];(\mathcal O(f_l))_{l\ne i}\rangle,
\]
\[
\mathfrak U^\sharp[\eta](f)=-\beta\langle\mathfrak K^\sharp[\eta];\mathcal O(f_\cdot)\rangle+a^{-1}\sum_i\langle\mathfrak K^\sharp[f_i\eta];(\mathcal O(f_l))_{l\ne i}\rangle,
\]
and, for $\chi$ covering the $f_i$, the rotational Ward identity at finite cutoff is
\[
\sum_i\langle\mathcal O(f_1);\dots;\mathcal O(\nabla^s_\omega f_i);\dots;\mathcal O(f_n)\rangle=\mathfrak U_{\omega,\chi}(f)+\langle\mathfrak B_{\omega,\chi};\mathcal O(f_\cdot)\rangle .
\]
The spread of sets $S_0,\dots,S_n$ is $\inf\max_{i<i'}|y_i-y_{i'}|_\infty$ over $y_i\in S_i$; it is at least the sup distance of any two of them.

\subsection{Principal results}

\begin{definition}[Additional quantitative torus-law hypothesis]\label{hyp:quantitative-torus}
For a fixed integer $n\ge2$, hypothesis $\mathrm{(Q_n)}$ is the
following assertion about the interacting Wilson measure. There
exist a coupling threshold $\gamma_n$, depth thresholds $s_n(g)$
and $\mathsf d_n(g,\ell)$, and errors $\varepsilon_n(g,s)$ and
$\omega_n(g,s;c)$, uniform in the cutoff, label, position and
weights, with the properties below. We require
$\varepsilon_n(g,s)\to0$ as $s\to\infty$ and use the explicit
large-field majorant
\[\omega_n(g,s;c)=L^{3s}(\mathsf Z^\vee_s)^n\big[L^{4n}c_\Sigma+s\big]\exp\big[-cA_0^2(Y^*_s)^{2\mathfrak p_0}\big],\qquad c_\Sigma=\sum_{l\ge0}L^{-2l}(l+2)^{2n\mathfrak p_0},\]
$Y^*_s=\log(\frac45L^{-t_1(s)}x^-_s)$, $x^-_s=\frac45L^{s-2}g^{-2}(L^2C_W(s)Y_s^{3\mathfrak p_0})^{-1}$, $Y_s=\log\bar x_s$, and $\mathsf Z^\vee_s\ge1$, $t_1(s)\ge1$ are numbers formed with the basic constants, $n$, $g$ and $s$: $\mathsf Z^\vee_s$ bounds, in the units of the levels above the working level of depth $s$, the sizes of the $n$ sources carried there by Theorem~\YMMref{D}{8}, and is at most a power of $\bar x_s$ times a polynomial in $Y_s$, and $L^{t_1(s)}$ is at most a polynomial in $Y_s$; so $\mathsf L_w\omega_n(g,s;c)$ bounds the large-field histories summed against their weights $\mathsf w_t$, each with the $n$-th power of the sizes of the sources at its level.

Let $n\ge2$, $0<g\le\gamma_n$, $(K,\beta)\in\mathfrak D(g)$, $\ell\in\mathfrak L(g)$, $0<\mathsf d<1$, $s_n(g)\le s=s(\mathsf d)\le K$, and let $c,\mathsf L_w>0$ satisfy $\mathsf w_t\le\mathsf L_wL^{3(K-t)}e^{-cp_0(g_t)^2}$ $(1\le t\le K)$, $\varsigma_K(c)\ge4n+5$ and $\mathsf L_w\omega_n(g,s;c)\le\frac18$; here $\mathsf w_t$ is the probability of a first creation of a large field at level $t$, $p_0(g_t)=A_0(\log g_t^{-2})^{\mathfrak p_0}$, and $\varsigma_K(c)=2c\mathfrak p_0A_0^2(\log g_K^{-2})^{2\mathfrak p_0-1}$. For real weights $w_a$ and points $x_a$ with $|x(p)-x_a|<\mathsf d$ where $w_a(p)\ne0$ and $|x(p)-x(q)|\ge\mathsf d$ where $w_a(p)\ne0\ne w_b(q)$, $a\ne b$,
\[
\begin{multlined}
\bigl|\langle\mathcal O(w_1);\dots;\mathcal O(w_n)\rangle_{K,\ell,\beta}-3\mathsf b^{-n}\mathcal K^{(n)}_K(w)\bigr|\le\mathsf e_n\,\mathsf b^{-n}\prod_a\sup|w_a|,\qquad\mathsf e_n=\varepsilon_n(g,s)+2(1 \\
+\bar Bg^2L^{2-s})^n\mathsf L_w\omega_n(g,s;c).
\end{multlined}
\]
The hypotheses hold with $c=c'$ and $\mathsf L_w=2M^{-3}\ell^3e^{B(X_g)\ell^3}$ for $\mathsf d\le\mathsf d_n(g,\ell)$ and every $K\ge s(\mathsf d)$, and then $\mathsf e_n\to0$ as $\mathsf d\to0$ (for $n\ge6$ provided $\log g^{-2}\ge(\frac{4n+5}{2c'\mathfrak p_0A_0^2})^{1/(2\mathfrak p_0-1)}$).
\end{definition}

The arbitrary-order small-field remainder in
Definition~\ref{hyp:quantitative-torus} is an additional hypothesis.
The Gaussian cyclic calculation and the large-field majorant alone
do not establish it. The clauses below that use $\mathrm{(Q_n)}$
are stated conditionally; the other clauses retain their own hypotheses.

\begin{theorem}[Plane fields]\label{C-360-2}
Let $0<g\le\gamma^{\rm ref}_W$ and $\ell'\in\mathfrak L(g)$.
\begin{enumerate}
\item[(a)] \emph{One torus.} For $r\ge2$ and $0<\mathsf d<1$ there are $K_0(r,\mathsf d)$, depending on $g,r,\mathsf d$, and $k(r,\mathsf d)$, depending besides on $\ell'$, such that
\[
|\langle\mathcal O(w_1);\dots;\mathcal O(w_r)\rangle_{K,\ell',\beta}|\le k(r,\mathsf d)\prod_a\sup|w_a|
\]
for all $(K,\beta)\in\mathfrak D(g)$ with $K\ge K_0(r,\mathsf d)$ and all real weights whose sets of barycentres are mutually at distance $\ge\mathsf d$; the weights need no regularity and no sign, and they may depend on the plane; with $2^rk(r,\mathsf d)$ the same holds for complex weights. For every $p>7$, $k(r,\mathsf d)\le C_{r,p}\mathsf d^{-pr}$. Besides, the growth in the order is controlled: for every $0<g\le\gamma$ with $g^{-2}\ge x_{\rm J}$, every side $\ell\in\mathfrak L(g)$, every $q\ge2$, every $\mathsf d\ge a_K$, all planes and all bounded complex functions $u_1,\dots,u_q$ on $T_\ell$ whose supports are at pairwise sup distances $\ge2\mathsf d$, and for $K$ above a threshold depending on $g,q,\mathsf d$,
\[
\begin{multlined}
|\langle\mathcal O_{\lambda_1}(u_1);\dots;\mathcal O_{\lambda_q}(u_q)\rangle_{K,\ell,\beta}|\le A_\ell^{\,q}\,(q!)^{\frac{21}2+\mathfrak p_0}\,e^{36c_Lq^2}\max\{1,\mathsf d^{-4q}\}\prod_a\sup|u_a|,\qquad c_L \\
=\frac{(\log L)^2}{2c'A_0^2(\log\frac L2)^2},
\end{multlined}
\]
with $A_\ell$ depending on $g$ and $\ell$.
\item[(b)] \emph{Every state.} For $n\ge2$ and $\mathsf d_0,R>0$ there is $C_n(\mathsf d_0;R)$, depending on $g,\ell',n,\mathsf d_0,R$, such that for every $(K,\beta)\in\mathfrak D(g)$, all planes and all bounded real $f_i$ supported in closed balls of radius $R$ at mutual distances $\ge\mathsf d_0$,
\[
|\langle\mathcal O_{\lambda_1}(f_1);\dots;\mathcal O_{\lambda_n}(f_n)\rangle_s|\le C_n(\mathsf d_0;R)\prod_i\sup|f_i|,
\]
where $s$ is any member of $\mathcal G^{\rm per}_{K,\beta}$, or the measure of $T_\ell$ with $\ell=2L^m\ge L\ell'$ in $\mathfrak L(g)$ and the supports in $(-\frac\ell4,\frac\ell4)^3$. The same holds for bounded weights, and $|\nu'(\mathcal O_\lambda(f))-\nu(\mathcal O_\lambda(f))|\le C_1(R)\sup|f|$ for any two such states $\nu,\nu'$.
\item[(c)] \emph{Limit points.} Let $(K_j,\beta_j)\in\mathfrak D(g)$, $K_j\to\infty$, and let $s_j$ be either any periodic Gibbs state at $(K_j,\beta_j)$, or the measure of $T_{\ell_j}$ with $\ell_j\in\mathfrak L(g)$, $\ell_j\ge L\ell'$, $\ell_j\to\infty$, the cumulants being formed once the supports lie in $(-\frac{\ell_j}4,\frac{\ell_j}4)^3$. Every subsequence has a further subsequence along which the cumulants of (b) converge for all $n\ge2$, all planes and all real $f_i\in C^1_c(\mathbb R^3)$ with disjoint supports. The limits $S_n$ are real, multilinear and symmetric; they obey the bound of (b) and $|S_n|\le B^n(n!)^{4+p}\prod_i\sup|f_i|$ for every $p>7$, with $B$ depending on $g,\ell',\mathsf d_0,R,p$ and not on $n$; they are invariant under $x\mapsto rx+v$, $r\in W_3$, $v\in\mathbb R^3$, the planes being permuted as the axes, and their sum over the planes is the limit of $S^T_n$.

With $\Pi(J)=\sum_\pi\prod_{I\in\pi}S_{\#I}(I)$, $\pi$ running over the partitions of a finite family $J$ of pairs (plane, function) into blocks of at least two members, $\sum_{\alpha,\alpha'}\bar z_\alpha z_{\alpha'}\Pi(\theta J_\alpha\cup J_{\alpha'})\ge0$ for families $J_\alpha$ on one side of a mirror $\theta$: at every coordinate mirror, the planes being fixed, in both cases; and, for limits of periodic Gibbs states, also at the diagonal mirrors with normal $e_\mu\pm e_\nu$, which exchange the planes $\mu,\nu$ in the reflected entries.
\item[(d)] \emph{The full sequence on a torus.} Let $0<g\le\gamma$ with $g^{-2}\ge x_{\rm J}$, $\mathsf b\in I(g)$, $\ell\in\mathfrak L(g)$, $n\ge2$, $\mathsf d>0$ and $\theta_n=1/(1260n(8n+3))$. There are $C$, depending on $n$ and $\mathsf b$, and $K_1$, depending on $g,\ell,\mathsf d,\mathsf b,n$, such that for $K\ge K_1$, all planes and all real $\phi_a$ with ${\rm Lip}_\infty\phi_a<\infty$ and supports at pairwise sup distances $\ge\mathsf d$, $G_K=\langle\mathcal O_{\lambda_1}(\phi_1);\dots;\mathcal O_{\lambda_n}(\phi_n)\rangle_{K,\ell,L^K\mathsf b}$ obeys
\[
|G_{K+1}-G_K|\le C\,\ell^{3(n+1)}K^{2\mathfrak p_0+2}L^{-\theta_nK}\prod_a\bigl(\sup|\phi_a|+{\rm Lip}_\infty\phi_a\bigr).
\]
So $\lim_KG_K$ exists with the summed rate; for continuous $\phi_a$ with supports at pairwise sup distances $\ge2\mathsf d$ it exists without a rate; and the same holds with $\mathsf b\check D_m$ for an entry. With one plane field among action densities the rate is better: for real Lipschitz $k,f_1,\dots,f_n$ with supports pairwise $\ge\mathsf d$ apart, $K\ge K_1$ and a $C$ depending on $g,\ell,n,\mathsf d,\mathsf b$,
\[
\bigl|\langle\mathcal O_\lambda(k);\mathcal O(f_\cdot)\rangle_{K+1,\ell,L^{K+1}\mathsf b}-\langle\mathcal O_\lambda(k);\mathcal O(f_\cdot)\rangle_{K,\ell,L^K\mathsf b}\bigr|\le CK^{2\mathfrak p_0(n+2)+2}L^{-K/80}\|k\|_\sharp\prod_i\|f_i\|_\sharp .
\]
The Lipschitz constants cannot be dropped from the bound on $|G_{K+1}-G_K|$: it fails with $\prod_a\sup|\phi_a|$ alone for bounded functions sampled at the barycentres, for instance for the indicator of the points of a ball that are barycentres at an even cutoff and at no smaller cutoff, whereas the bound on $|G_K|$ itself needs sup norms only, by (a). One pair $(C,K_1)$ serves a compact interval of labels in $I(g)$; every $\mathsf b\ge2\max\{\gamma^{-2},x_{\rm J}\}$ is admitted at $g=(\mathsf b/2)^{-1/2}$, and by the dilation $G_K(\phi;\mathsf b,\ell)=L^{zn}G_{K-z}(\phi(L^{-z}\cdot);L^z\mathsf b,L^z\ell)$ the same holds at every label $\mathsf b>0$ on every side $2L^m$.
\item[(e)] \emph{Two points at one scale.} Let $0<g\le\gamma^{\rm sd}_W$, $\ell=\ell_g$, $\vartheta\in(0,1)$, $(K,\beta)\in\mathfrak D(g)$, $0<\mathsf d\le\mathsf d^{\rm ref}_0(g,\vartheta)$ with $s(\mathsf d)\le K$, $\mathsf d^{\rm ref}_0(g,\vartheta)$ being the threshold in equation~\YMMref{(2335)}{847}, and let $(f,h)$ be a test pair of $\mathfrak T_{T_\ell}(\mathsf d;\vartheta)$ (Lemma~\ref{lem:ward-test-pairs}). Then, uniformly in $K$,
\[
\Bigl|\langle\mathcal O_\lambda(f);\mathcal O_{\lambda'}(h)\rangle_{K,\ell,\beta}-\tfrac32\mathsf b^{-2}\mathcal K_{\lambda\lambda'}(f,h)\Bigr|\le C(g,\vartheta)\bigl(\log\tfrac1{\mathsf d}\bigr)^{-2\mathfrak p_0}\mathsf b^{-2}\|f\|_\infty\|h\|_\infty .
\]
Since $\mathcal K_{\lambda\lambda'}(f,h)\ge\frac6{175}(\frac\vartheta{20})^{10}\|f\|_\infty\|h\|_\infty$, the nine two-point functions are positive for small $\mathsf d$. So is the two-point function of $\mathcal O_1-\mathcal O_2$ at the test pairs with $\hat u_1^2+\hat u_2^2\le\frac14$ for all $u\in{\rm supp}\,f-{\rm supp}\,h$, for instance at those carried by balls of radius $\mathsf d$ whose centres are $4\mathsf de_3$ apart: there
\[
\begin{multlined}
[\mathbb K_3(u)]^2_{11}-2[\mathbb K_3(u)]^2_{12}+[\mathbb K_3(u)]^2_{22}\\
=\frac{9(\hat u_1^2-\hat u_2^2)^2-6(\hat u_1^2+\hat u_2^2)+2}{16\pi^2|u|^6}
\ge\frac1{32\pi^2|u|^6},
\end{multlined}
\]
and, since $|u|\le10\mathsf d$ on ${\rm supp}\,f-{\rm supp}\,h$ (Lemma~\ref{lem:ward-test-pairs}(a)) and the masses of $f,h$ are at least $\vartheta\mathsf d^3\|f\|_\infty$, $\vartheta\mathsf d^3\|h\|_\infty$ (Lemma~\ref{lem:ward-test-pairs}), its main term is at least 
\[
\frac32\mathsf b^{-2}(32\pi^2)^{-1}10^{-6}\vartheta^2\|f\|_\infty\|h\|_\infty,
\]
 against $4$ times the error of the display; so no plane field is a multiple of $\mathcal O$. The law holds with $9$ times the error for $\mathcal O-2\mathcal O_m$. For two entries $\mathsf b[\mathcal O-2\mathcal O_m]$, both the main term and the error are multiplied by $\mathsf b^2$. The two-point law on a torus of small side, Theorem~\YMMref{AU}{847} in the form of Proposition~\YMMref{126.4}{846} against the lattice kernel, holds for all complex weights $w_1,w_2$ carried by balls of radius $\mathsf d-a_K$ about points $x_0,y_0$ with $3\mathsf d\le|x_0-y_0|\le4\mathsf d$, with $16$ times its tolerance ($4$ times for real weights), whatever the dependence of the weights on the plane.
\item[(f)] \emph{Free photons and the conditional continuum comparison.}
For three free lattice photon fields the cyclic expression
$3\mathsf b^{-n}\mathcal K_K^{(n)}(w)$ is the exact cumulant.
For the interacting plane fields, assume $\mathrm{(Q_n)}$
of Definition~\ref{hyp:quantitative-torus}, with its stated data
and error $\mathsf e_n$. For real Lipschitz $\phi_a$ with ${\rm supp}\,\phi_a\subset B(x_a,\mathsf d)\subset B(c,\frac\ell4)$, ${\rm Lip}\,\phi_a\le\|\phi_a\|_\infty/(\vartheta_0\mathsf d)$, $0<\vartheta_0\le1$, and supports $\ge\mathsf d$ apart, the plane fields $\mathcal O_{\lambda_a}(\phi_a)$ obey the law with $\mathcal K^{(n)}_\lambda(\phi)$ for $\mathcal K^{(n)}_K(w)$ and with the bracket $\mathsf e_n+C'_n(\frac{a_K}{\vartheta_0\mathsf d}+\frac{\mathsf d^3}{\ell^3})$, where $C'_n=\frac32n![6^nC_{\rm lat}(1+3C_{\rm lat})^{n-1}+43\cdot11^n]$; with $3^nC'_n$ the same holds for the combinations $\sum_\lambda k_{a\lambda}\mathcal O_\lambda$, $k_a\in[-1,1]^3$, such as $\mathcal O-2\mathcal O_m$. Multiplying any of these entries by $\mathsf b$ multiplies both the main term and its error by $\mathsf b$ for that entry.
\item[(g)] \emph{The law in every state.} (i) Let $0<g\le\gamma^{\rm sd}_W$, $n\ge1$ and $\sigma\in(0,1)$. There are $\mathsf d_2>0$ and $C,c$, depending on $g,n,\sigma$, such that for $(K,\beta)\in\mathfrak D(g)$, $0<\mathsf d\le\mathsf d_2$ with $s(\mathsf d)\le K$, all bounded real $f_a$ vanishing outside closed balls of radius $\mathsf d$ inside one ball of radius $9\mathsf d$, at mutual distances $\ge\mathsf d$, and every state $s$ at $(K,\beta)$,
\[
|S^T_n[s](f)-S^T_{n;K,\ell_g,\beta}(f)|\le\mathsf e'_n(g,\mathsf d)\,\mathsf b^{-n}\prod_a\|f_a\|_\infty,\qquad\mathsf e'_n(g,\mathsf d)\le C\bigl(\log\tfrac1{\mathsf d}\bigr)^{cn}\mathsf d^{\,3(1-\sigma)-[\frac n2-1]_+};
\]
If $\mathrm{(Q_n)}$ holds on $T_{\ell_g}$ and $g\le\min\{\gamma_n,\gamma_W^{\rm sd}\}$, the preceding comparison and (f) give the free law in every state for $2\le n\le7$, under the depth hypotheses of $\mathrm{(Q_n)}$. (ii) Assume $\mathrm{(Q_m)}$ on $T_{\ell'}$ for every $2\le m\le8n$. Let $n\ge2$, $0<g\le\min_{2\le m\le8n}\gamma_m$ with $\log g^{-2}\ge(\frac{32n+5}{2c'\mathfrak p_0A_0^2})^{1/(2\mathfrak p_0-1)}$, and $\ell'\in\mathfrak L(g)$. For $0<\mathsf d\le\min_{2\le m\le8n}\mathsf d_m(g,\ell')$, every $(K,\beta)\in\mathfrak D(g)$ with $K\ge s(\mathsf d)$, every $s\in\mathcal G^{\rm per}_{K,\beta}$ and the measure $s$ of every torus $T_\ell$, $\ell\ge\ell'$ in $\mathfrak L(g)$, and real plaquette weights with the support and separation conditions of Definition~\ref{hyp:quantitative-torus}, inside a cube of side $\frac14\ell'$,
\[
\bigl|\kappa_s(\mathcal O(w_1),\dots,\mathcal O(w_n))-3\mathsf b^{-n}\mathcal K^{(n)}_{K,\ell'}(w)\bigr|\le\mathsf e^\sharp_n(g,\ell',\mathsf d)\,\mathsf b^{-n}\prod_a\sup|w_a|,
\]
where $\mathsf e^\sharp_n$ depends on no state, side or cutoff and tends to $0$ with $\mathsf d$. Under the same hypotheses, the following comparison holds for every non-empty $B\subset\{1,\dots,n\}$. It also holds for $n=1$ assuming $\mathrm{(Q_m)}$ for $2\le m\le8$ and their coupling and depth conditions:
\[
\bigl|\kappa_s\bigl((\mathcal O(w_a))_{a\in B}\bigr)-\kappa_{T_{\ell'}}\bigl((\mathcal O(w_a))_{a\in B}\bigr)\bigr|\le\eta^\sharp(1+\eta^\sharp)^{n-1}\mathsf b^{-\#B}\prod_{a\in B}\sup|w_a|,
\]
$\kappa_{T_{\ell'}}$ being the joint cumulant under the measure of $T_{\ell'}$, with $\eta^\sharp\to0$ as $\mathsf d\to0$. The same holds for the plane fields and their combinations $\mathcal O-2\mathcal O_m$, with the corresponding three-plane coefficient factors. For $\mathsf b[\mathcal O-2\mathcal O_m]$, multiply the comparison bound by $\mathsf b$ for every such entry.
\item[(h)] \emph{Short distance in every state and in every limit.} Let $0<g\le\gamma_\sigma$, where $\gamma_\sigma\le\gamma^{\rm sd}_W$ is a constant formed with the constants, $\mathsf d_\sigma(g)>0$ a threshold and $\Omega_\sigma=\{(x,y):0<|x-y|<\mathsf d_\sigma(g)\}$; put $E_2(g,t)$, $E_3(g,t)$ for the tolerances at the distance $t$ of the laws of two and of three points at narrow bumps, which Step~4 proves on the least torus and (i) and (iii) below state, functions of $g$ and $t$ formed with the constants, $E_2(g,t)\le C(g)(\log\frac1t)^{19\mathfrak p_0}t^{1/2}$, $E_3(g,t)\le C(g)[(\log\frac1t)^{39\mathfrak p_0+2}t^{1/2}+(t/\ell_g)^3]$.
\begin{enumerate}
\item[(i)] At a cutoff, in every state: for $(K,\beta)\in\mathfrak D(g)$, every state $s$ at $(K,\beta)$, points $x_0,y_0$ with $\mathsf d_0=|x_0-y_0|\le\mathsf d_\sigma(g)$, $\rho>0$ with $s(\rho)\le K$, and bounded real $f_1,f_2$ supported in balls of radius $\rho$ inside $B(x_0,\frac18\mathsf d_0)$, $B(y_0,\frac18\mathsf d_0)$,
\[
\begin{multlined}
\Bigl|\langle\mathcal O_\lambda(f_1);\mathcal O_\mu(f_2)\rangle_s-\tfrac32\mathsf b^{-2}\!\!\sum_{\lambda(p)=\lambda}\sum_{\lambda(q)=\mu}\!f_1(x(p))f_2(x(q))C^{\rm free}_K(p,q)^2\Bigr|\le E_2(g, \\
\tfrac12\mathsf d_0)\Bigl(\frac{2\rho}{\mathsf d_0}\Bigr)^6\mathsf b^{-2}\|f_1\|_\infty\|f_2\|_\infty,
\end{multlined}
\]
$C^{\rm free}_K$ being that of $T_{\ell_g}$; for real Lipschitz $\phi_a$ with ${\rm Lip}\,\phi_a\le\|\phi_a\|_\infty/(\vartheta\rho)$ the same holds with $\frac32\mathsf b^{-2}\mathcal K_{\lambda\mu}(\phi_1,\phi_2)$ and with the bracket $E_2+C\frac{a_K}{\vartheta\rho}$. By the change of units $x\mapsto L^zx$ this narrow form holds at every $0<g\le\gamma$, on every side of $\mathfrak L(g)$, for distances $\mathsf d_0$ below a threshold depending on $g$.
\item[(ii)] Let $(K_j,\beta_j)\in\mathfrak D(g)$, $K_j\to\infty$, with labels $\mathsf b_j\to\mathsf b$, and let $s_j$ be states at $(K_j,\beta_j)$. Along a subsequence $\langle\mathcal O_\lambda(\phi_1);\mathcal O_\mu(\phi_2)\rangle_{s_j}$ converges for all real Lipschitz $\phi_a$ with ${\rm supp}\,\phi_1\times{\rm supp}\,\phi_2\subset\Omega_\sigma$, and every such limit has Borel kernels $\tau_{\lambda\mu}(x,y)=\hat\tau_{\lambda\mu}(x-y)$ with
\[
\begin{multlined}
\bigl|\tau_{\lambda\mu}(x,y)-\tfrac32\mathsf b^{-2}[\mathbb K_3(x-y)_{\lambda\mu}]^2\bigr|\\ \le(\tfrac3\pi)^2\epsilon_\sigma(g,|x-y|)\,\mathsf b^{-2}|x-y|^{-6} \\
\text{almost everywhere},\\ \epsilon_\sigma(g,t)=64\sup_{0<t'\le t}E_2(g,\tfrac12t')\le C(g)\bigl(\log\tfrac1t\bigr)^{19\mathfrak p_0}t^{1/2};
\end{multlined}
\]
before a kernel is named, at every centre and for bumps of radius $\rho$ in $B(x_0,\frac18\mathsf d_0)$, $B(y_0,\frac18\mathsf d_0)$, the limit differs from $\frac32\mathsf b^{-2}\mathcal K_{\lambda\mu}(\phi_1,\phi_2)$ by at most 
\[
64E_2(g,\frac12\mathsf d_0)\mathsf b^{-2}\mathsf d_0^{-6}\rho^6\|\phi_1\|_\infty\|\phi_2\|_\infty.
\]
 If besides, for all $n\ge2$, the $S^T_n[s_j]$ converge on separated tuples to the connected functions $S^{\sigma'}_n$ of reconstruction data $\sigma'$ on $\mathbb R^3$, let $\mathsf s^T_n$ be the continuous kernel of $S^{\sigma'}_n$ $(n=2,3)$: with the continuous kernels $\mathsf s_1,\mathsf s_2,\mathsf s^{\rm m}_3$ of the moments of $\sigma'$ (Proposition~\YMMref{108.3}{771}), $\mathsf s^T_2=\mathsf s_2-\mathsf s_1\otimes\mathsf s_1$ and 
\[
\begin{multlined}
\mathsf s^T_3(x,y,z)=\mathsf s^{\rm m}_3(x,y,z)-\mathsf s_2(x,y)\mathsf s_1(z)-\mathsf s_2(x,z)\mathsf s_1(y)-\mathsf s_2(y,z)\mathsf s_1(x) \\
+2\mathsf s_1(x)\mathsf s_1(y)\mathsf s_1(z)
\end{multlined}
.\]
Then $\sum_{\lambda,\mu}\tau_{\lambda\mu}=\mathsf s^T_2$ almost everywhere on $\Omega_\sigma$, so that $|\mathsf s^T_2(x,y)-\frac9{16\pi^2}\mathsf b^{-2}|x-y|^{-6}|\le9\epsilon_\sigma(g,|x-y|)\mathsf b^{-2}|x-y|^{-6}$ on $\Omega_\sigma$.
\item[(iii)] Three points at two scales: in the latter case, at every $(x,y,z)$ with $0<|x-y|=r_*\le\mathsf d_*={\rm dist}(z,\{x,y\})\le\mathsf d_\sigma(g)$,
\[
|\mathsf s^T_3(x,y,z)-3\mathsf b^{-3}k_3(x,y,z)|\le E_3(g,\mathsf d_*)\,\mathsf b^{-3}r_*^{-3}\mathsf d_*^{-6},
\]
uniformly in $r_*/\mathsf d_*\in(0,1]$.
\item[(iv)] Point splitting: off one null set of pairs $(c,z)$ with $0<\mathsf d=|c-z|\le\frac12\mathsf d_\sigma(g)$, for all $0<r\le\frac12\mathsf d$ and all $\lambda$,
\[
\begin{multlined}
\bigl|\tfrac{2\pi}3\mathsf br^3\mathsf s^T_3(c+\tfrac r2e_\lambda,c-\tfrac r2e_\lambda,z)-t_\lambda(c,z)\bigr|\le\bigl[\epsilon'_\sigma(g,\mathsf d)+2.7(r/\mathsf d)^2\bigr]\mathsf b^{-2}\mathsf d^{-6},\qquad\epsilon'_\sigma(g, \\
\mathsf d)\le C(g)\bigl(\log\tfrac1{\mathsf d}\bigr)^{39\mathfrak p_0+2}\mathsf d^{1/2},
\end{multlined}
\]
$t_\lambda=\sum_\mu(\tau_{\lambda\mu}-\frac13\sum_{\lambda'}\tau_{\lambda'\mu})$ being the kernel of the limit of $\langle\mathcal O_\lambda-\frac13\mathcal O;\mathcal O\rangle$; and at every such $(c,z)$, for every unit vector $n$, 
\[
\begin{multlined}
|\frac{2\pi}3\mathsf br^3\mathsf s^T_3(c+\frac r2n,c-\frac r2n,z)-\mathsf t^\circ(n,c-z)|\le[12\sup_{\frac34\mathsf d\le t\le\frac{17}{16}\mathsf d}E_3(g,t) \\
+2.7(r/\mathsf d)^2]\mathsf b^{-2}\mathsf d^{-6}
\end{multlined}
.\]

\item[(v)] On the least torus at a fixed label $\mathsf b\in I(g)$, for $g$ small enough in terms of the constants. For real Lipschitz $\phi_1,\phi_2$ with disjoint supports, $\lim_K\langle\mathcal O_\lambda(\phi_1)-\frac13\mathcal O(\phi_1);\mathcal O(\phi_2)\rangle_{K,\ell_g,L^K\mathsf b}$ exists. Let moreover $\mathsf d>0$ with $6\mathsf d\le\min\{1,\mathsf b\}$ and $0<\vartheta_0\le1$, let $\phi_1,\phi_2$ be real Lipschitz functions each carried by a ball of radius $\mathsf d$, with supports $\ge\mathsf d$ apart, $|c-z|\le6\mathsf d$ for $c\in{\rm supp}\,\phi_1$ and $z\in{\rm supp}\,\phi_2$, and ${\rm Lip}\,\phi_1\le\|\phi_1\|_\infty/(\vartheta_0\mathsf d)$; for $r>0$ put $\rho=\frac r8$ and let $f^c_\pm(y)=\frac3{\pi\rho^3}\bigl(1-|y-c\mp\frac r2e_\lambda|/\rho\bigr)_+$ be the radial tents of radius $\rho$ and mass $1$ about $c\pm\frac r2e_\lambda$, of sup norm $\frac3{\pi\rho^3}=\frac{1536}{\pi r^3}$. There are $r_*>0$, depending on $g,\mathsf b,\mathsf d$, and $C(g,\mathsf b,\mathsf d,\vartheta_0)$ such that for $0<r\le r_*$ and every $K$ with $128a_K\le r$, that is $16a_K\le\rho$, and $s(\rho)\le K$
\[
\begin{multlined}
\Bigl|\tfrac{2\pi}3\mathsf br^3\!\int\!\phi_1(c)\,S^T_{3;K,\ell_g,L^K\mathsf b}(f^c_+,f^c_-,\phi_2)\,dc-\bigl\langle\mathcal O_\lambda(\phi_1)-\tfrac13\mathcal O(\phi_1);\mathcal O(\phi_2)\bigr\rangle_{K,\ell_g,L^K\mathsf b}\Bigr| \\
\le C(g,\mathsf b,\mathsf d,\vartheta_0)\bigl[r^{\beta_H/72}+(a_K/r)^{1/18}\bigr]\mathsf b^{-2}\|\phi_1\|_\infty\|\phi_2\|_\infty ,
\end{multlined}
\]
$\beta_H$ being the H\"older exponent. As the depth function is non-increasing, $s(\rho)\le K$ gives $s(\frac r2)\le K$; and $r_*$ is taken so that for $0<r\le r_*$ the depth $s=s(\frac r2)$ obeys $s\ge3$ and $L^2C_W(s)Y_s^{3\mathfrak p_0}\le L^{s-3}$ (both hold for small $r$, since $s(\frac r2)\to\infty$ as $r\to0$ and the second holds at all large $s$) and $\frac r2\le\min\{r_0(\mathsf b),r_1(\mathsf b)\}$, the radii, depending on $g$ and $\mathsf b$, up to which Step~4 forms the induction over the levels and the frame of the pair. Here $C(g,\mathsf b,\mathsf d,\vartheta_0)=C(g,\mathsf b,\vartheta_0)\mathsf d^{-9/2}$, of which the part of the three-point function beyond its explicit Gaussian part takes $C(g,\mathsf b)\mathsf d^{-5/2}$.

The explicit Gaussian parts of the two members, which take the rest, are compared on the lattice of $T_{\ell_g}$: the near leg, between the two tents, against $a_K^3\mathbb K_3(re_\lambda)$, by the comparison of the lattice kernel with $\mathbb K_3$ in the setting, at the cost $(a_K/r)^2+(r/\ell_g)^3$, the mean of $\mathbb K_3$, whose entries are harmonic off $0$, over the two radial tents being its value at $re_\lambda$, and $\frac{2\pi}3r^3\cdot3\mathbb K_3(re_\lambda)=\frac32(E_{\lambda\lambda}-\frac13\mathbf 1)$; the far legs by the bound on the differences of $C^{\rm free}_K$ along one plane of the setting; and the shift of the argument of $\phi_1$ by at most $r$ costs $r\,{\rm Lip}\,\phi_1\le r\|\phi_1\|_\infty/(\vartheta_0\mathsf d)$.

Hence, at this fixed $\mathsf d$, the limit in $K$ is the limit as $r\to0$ of 
\[
\frac{2\pi}3\mathsf br^3\int\phi_1(c)\,S^{(\ell_g,\mathsf b)}_3(f^c_+,f^c_-,\phi_2)\,dc.
\]
 The limit in $K$ exists for all real Lipschitz $\phi_1,\phi_2$ with disjoint supports, and it is identified in this way for all $\phi_1,\phi_2$ as above.
\end{enumerate}
\end{enumerate}
\end{theorem}

\begin{theorem}[One stress insertion; stress entries; conservation]\label{C-360-3}
\begin{enumerate}
\item[(a)] \emph{Bounds in every geometry.} Let $0<g\le\gamma^{\rm ref}_W$, $\ell'\in\mathfrak L(g)$, $n\ge1$, $0<\rho<1$, $\omega\in\mathfrak{so}(3)$, and let $\chi$ be a cut-off function at $K$, Lipschitz on $\mathbb R^3$. For every $(K,\beta)\in\mathfrak D(g)$,
\[
\Bigl|\Bigl\langle\beta\sum_ph_{\omega,\chi}(p);\mathcal O(f_\cdot)\Bigr\rangle_{K,\ell',\beta}\Bigr|\le\mathsf b\,C^\flat_n(\rho)\,|\omega|\,{\rm Lip}_\infty(\chi)\prod_i\|f_i\|,
\]
(i) with $\|f_i\|=\sup|f_i|$ for bounded real $f_i$ (or real weights), ${\rm Var}(\chi)$ and the ${\rm supp}\,f_i$ being pairwise $\ge\rho$ apart; (ii) with $\|f_i\|=\|f_i\|_\sharp$ for real Lipschitz $f_i$ in any mutual position, ${\rm Var}(\chi)$ being at distance $\ge\rho$ from each ${\rm supp}\,f_i$. $C^\flat_n(\rho)=468\ell'k^\flat_n(\rho)$, where $k^\flat_n(\rho)$ depends on $g,\ell',n,\rho$, not on $\chi$, on the $f_i$ or on the place of the shell. In (ii) the number may be taken free of the side and of the state: on every torus $T_\ell$, $\ell\ge\ell'$ in $\mathfrak L(g)$, and in every $\nu\in\mathcal G^{\rm per}_{K,\beta}$, for $K$ above a threshold depending on $g,\ell',n,\rho$, the left member is at most $156\,\mathsf b\,C^e_n(\rho)|\omega|\,{\rm Lip}(\chi)R_1(\chi)N_\chi\prod_iN_i\|f_i\|_\sharp$, where $R_1(\chi)\le\frac12+\frac32\ell$ is the largest $|x_b|_1$ at the bonds where $\chi$ varies, $N_\chi$ and $N_i$ are the numbers of the cubes of a mesh $h(n,\rho)$ that meet the shell and the supports, and $C^e_n(\rho)$ depends on $g,\ell',n,\rho$; at every label $\mathsf b>0$ on every side $2L^m$ the same holds for all $K$ with a number depending on $\mathsf b,\ell,n,\rho$.

Both bounds are cases of one rule for a rough entry $F$ of Lemma~\ref{C-360-6} whose carrier lies, with respect to the $f_i$, as ${\rm Var}(\chi)$ in (i), resp.\ (ii): the left member with $F$ in the place of $\beta\sum_ph_{\omega,\chi}(p)$ is at most $k^\flat_n(\rho)\|F\|\prod_i\|f_i\|$ on $T_{\ell'}$, and, in (ii), at most $C^e_n(\rho)N_F\|F\|\prod_iN_i\|f_i\|_\sharp$ free of the side and of the state, $N_F$ being the number of the cubes of the mesh that meet the carrier of $F$; here $-\beta\sum_ph_{\omega,\chi}(p)=\mathsf b\,{\rm Pf}(k^h)$ with $\|{\rm Pf}(k^h)\|\le156|\omega|\,{\rm Lip}(\chi)R_1(\chi)\le468\ell'|\omega|\,{\rm Lip}_\infty(\chi)$. The same rule gives the bounds for the whole flux insertion $\mathfrak B_{\omega,\chi}=\mathsf b\,{\rm Pf}(k^h)+\mathcal O(v^{\mathfrak B})$, of norm $\le(158\mathsf b+3a)|\omega|\,{\rm Lip}(\chi)R'(\chi)$, $R'(\chi)=\max\{R_1(\chi),R'_1(\chi)\}$, where $R'_1(\chi)\le\frac32\ell$ is the number of Lemma~\ref{C-360-6}, the largest $|x(p)|_1$ at the plaquettes at which $\chi$ is not constant on the points $x(p)\pm ae_\mu$, $x(p)\pm\frac a2e_\mu$, $\mu=1,2,3$, among which are those where $\nabla^s_\omega\chi$ or $v_\chi$ does not vanish, which hold every plaquette at which $v^{\mathfrak B}\ne0$, so that the largest $|x(p)|_1$ at these may stand for $R'_1(\chi)$ throughout: on $T_{\ell'}$ with $\frac{158\mathsf b+3a}{156}C^\flat_n(\rho)$ in the place of $\mathsf b\,C^\flat_n(\rho)$, and free of the side and of the state with $(158\mathsf b+3a)C^e_n(\rho)|\omega|\,{\rm Lip}(\chi)R'(\chi)N_\chi$ in the place of $156\,\mathsf b\,C^e_n(\rho)|\omega|\,{\rm Lip}(\chi)R_1(\chi)N_\chi$.

For ${\rm supp}\,\eta$ at distance $\ge\rho$ from each ${\rm supp}\,f_i$ (and, in (i), the ${\rm supp}\,f_i$ pairwise $\ge\rho$ apart) it gives them for $\Theta_\beta[\eta]=\mathsf b\,{\rm Pf}(k^\eta)$, of norm $\le156\,\mathsf b\max_\rho{\rm Lip}\,\eta^\rho$, and for $\mathfrak K^{\rm v}[\beta\eta]$, which beside fields that do not read the bonds within $2a$ of ${\rm supp}\,\eta$ is the rough entry $-\Theta_\beta[\eta]+\mathsf b\,\mathcal O({\rm div}^s\eta)-\frac32a\,\mathcal O(e^\eta)$ of norm $\le(159\mathsf b+3a)\max_\rho{\rm Lip}\,\eta^\rho$ (the null term of Lemma~\ref{C-360-8}(a) drops): on $T_{\ell'}$ with $156\,\mathsf b\,k^\flat_n(\rho)\max_\rho{\rm Lip}_\infty\eta^\rho$, resp.\ $(159\mathsf b+3a)k^\flat_n(\rho)\max_\rho{\rm Lip}_\infty\eta^\rho$, in the place of $\mathsf b\,C^\flat_n(\rho)|\omega|\,{\rm Lip}_\infty(\chi)$, and free of the side and of the state with $156\,\mathsf b\,C^e_n(\rho)\max_\rho{\rm Lip}\,\eta^\rho N_\eta$, resp.\ $(159\mathsf b+3a)C^e_n(\rho)\max_\rho{\rm Lip}\,\eta^\rho N_\eta$, in the place of $156\,\mathsf b\,C^e_n(\rho)|\omega|\,{\rm Lip}(\chi)R_1(\chi)N_\chi$, $N_\eta$ being the number of the cubes of the mesh that meet the $\frac\rho2$-neighbourhood of ${\rm supp}\,\eta$.

Bounded fields, or real weights that depend on the plane, may be added in sup norms if they are $\ge\rho$ apart from the shell, from one another and from all the other fields.
\item[(b)] \emph{Stress entries.} $T_{lr}(\varphi)$ is $\mathsf b$ times a rough entry (Lemma~\ref{C-360-6}) of norm at most an absolute constant times $\sup|\varphi|$. So for $0<g\le\gamma^{\rm ref}_W$ the joint cumulants on one torus of $r$ entries $T_{l_ar_a}(\varphi_a)$ and $\mathcal O(w_a)$ with supports pairwise $\ge\mathsf d$ apart are bounded uniformly in $K\ge K'_0(r,\mathsf d)$, for merely bounded $\varphi_a,w_a$; for two entries this holds at every $0<g\le\gamma$: $|\langle T_{lr}(\varphi);\mathcal O(f)\rangle|\le C\mathsf b\sup|\varphi|\sup|f|$ and $|\langle T_{lr}(\varphi);T_{l'r'}(\varphi')\rangle|\le C\mathsf b^2\sup|\varphi|\sup|\varphi'|$, $C$ depending on $g$, the side and the distance. Along subsequences $\langle T;\mathcal O\rangle$ and $\langle T;T\rangle$ converge for continuous entries with disjoint supports, and the limits are symmetric in $(l,r)$, translation invariant, covariant under $W_3$ with the signs of a tensor, and reflection positive.
\item[(c)] \emph{The value, Lipschitz norms.} Let $\mathsf b\ge2x_W$, $g=(\mathsf b/2)^{-1/2}$, $\ell'\in\mathfrak L(g)$, $\beta=L^K\mathsf b$, $n\ge1$, $0<\mathsf d<1$ and $c<\frac1{14}$. There are $C,p,K_*$, depending on $\mathsf b,\ell',n,\mathsf d,c$, such that for $K\ge K_*$, every real $C^2$ vector field $\eta$ and all real Lipschitz $f_i$, the $n+1$ supports being pairwise at sup distance $\ge\mathsf d$,
\[
|\langle\mathfrak K^{\rm v}[\beta\eta];\mathcal O(f_\cdot)\rangle|+|\langle\Theta_\beta[\eta]-\mathsf b\,\mathcal O({\rm div}^s\eta);\mathcal O(f_\cdot)\rangle|\le CK^pL^{-cK}\|\eta\|_{C^2}\prod_i\|f_i\|_\sharp .
\]
By $x\mapsto L^zx$ the same holds at every label $\mathsf b>0$ on every side $2L^m$.
\item[(d)] \emph{The value, sup norms.} In the setting of (c): for every field $\eta$ and all real weights $w_a$, the $n+1$ carriers being pairwise at sup distance $\ge\mathsf d$, $|\beta\langle\mathfrak K^\sharp[\eta];\mathcal O(w_1);\dots;\mathcal O(w_n)\rangle|\le CK^pL^{-cK}\mathsf H(\eta)\prod_a\sup|w_a|$, and the same with rough entries or block forms (Lemmas~\ref{C-360-6} and~\ref{C-360-8}) in the places of the $\mathcal O(w_a)$; for $C^2$ fields $\eta$ vanishing within sup distance $\mathsf d$ of the carriers, the two members of (c) with $\mathcal O(w_a)$ for $\mathcal O(f_a)$ are $\le CK^pL^{-cK}\|\eta\|_{C^2}\prod_a\sup|w_a|$. At $n=1$, for every real vector field $\eta$, Lipschitz for the sup distance, and every bounded real $h$ (a function or a weight) with $\eta=0$ within sup distance $\mathsf d$ of $\{x(p):h(p)\ne0\}$,
\[
\begin{multlined}
|\langle\Theta_\beta[\eta];\mathcal O(h)\rangle-\mathsf b\langle\mathcal O({\rm div}^s\eta);\mathcal O(h)\rangle|+|\langle\mathfrak K^{\rm v}[\beta\eta];\mathcal O(h)\rangle|\le CK^pL^{-K/200}\bigl(\mathsf H(\eta) \\
+\max_\rho{\rm Lip}_\infty\eta^\rho\bigr)\sup|h| .
\end{multlined}
\]
For two fields with carriers $\ge\mathsf d$ apart, $|\beta^2\langle\mathfrak K^\sharp[\eta];\mathfrak K^\sharp[\eta']\rangle|\le CK^{2p}L^{-2cK}\mathsf H(\eta)\mathsf H(\eta')$, and, with $\mathsf H^+=\mathsf H+\max_\rho{\rm Lip}_\infty$, 
\[
|\langle\mathfrak K^{\rm v}[\beta\eta];\mathfrak K^{\rm v}[\beta\eta']\rangle|\le CK^pL^{-K/200}\mathsf H^+(\eta)\mathsf H^+(\eta')
\]
 and 
\[
|\langle\mathrm{St}_K[\eta];\mathrm{St}_K[\eta']\rangle|\le CK^pL^{-K/200}\mathsf b\,\mathsf H^+(\eta)\max_\rho{\rm Lip}_\infty\eta'^\rho
.\]
All this holds at every fixed label $\mathsf b>0$ on every side $2L^m$. At every first-passage datum $(K,\beta)\in\mathfrak D(g)$ on every side of $\mathfrak L(g)$, with numbers that depend on $g$ and on no label, the clauses on the cubic insertion alone (one $\beta\mathfrak K^\sharp[\eta]$ beside $n$ weights, rough entries or block forms, and two of them) hold for $0<g\le\gamma$ with $g^{-2}\ge x_{\rm J}$, and the other clauses, which also read the bilinear half of the orbital insertion and the term $a\,\mathcal O(e^\eta)$, for $0<g\le\gamma^{\rm ref}_W$, where $g^{-2}\ge x_{\rm J}$ holds of itself.
\item[(e)] \emph{Conservation and collapse in every periodic Gibbs state.} Let $0<g\le\gamma$ and $R,\mathsf d>0$. There are $C,p,K_0$, depending on $g,\mathsf d,R$, such that for all $(K,\beta)\in\mathfrak D(g)$ with $K\ge K_0$, all $\nu\in\mathcal G^{\rm per}_{K,\beta}$, all $\eta,\eta'\in C^2_c(\mathbb R^3;\mathbb R^3)$ and $H'\in C_c(\mathbb R^3;V)$ supported in a ball of radius $R$, with ${\rm supp}\,\eta$ at sup distance $\ge\mathsf d$ from ${\rm supp}\,\eta'$ and from ${\rm supp}\,H'$, and with $A=A^\nu$:
\[
\begin{multlined}
|\nu(A[\mathrm{St}_K[\eta']]A[\mathrm{St}_K[\eta]])|\le CK^pL^{-K/14}\|\eta'\|_{C^2}\|\eta\|_{C^2},\qquad|\nu(\mathfrak K^{\rm v}[\beta\eta])| \\
\le CK^pL^{-K/28}\|\eta\|_{C^2},
\end{multlined}
\]
\[
\begin{multlined}
|\nu(A[\mathcal T_K(H^{\rm cons}[\eta'])]A[\mathcal T_K(H^{\rm cons}[\eta])])|\le CK^pL^{-K/20}\|\eta'\|_{C^2}\|\eta\|_{C^2}, \\
|\nu(A[\mathcal T_K(H^{\rm cons}[\eta])]A[\mathcal T_K(H')])|\le CK^pL^{-K/60}\|\eta\|_{C^2}|H'|_\infty,
\end{multlined}
\]
and $|\nu(A[\Delta_K[\eta]]A[Z])|\le CK^pL^{-K/60}\|\eta\|_{C^2}\mathsf z$ for $Z=\mathcal T_K(H')$, $\mathsf z=|H'|_\infty$, and for $Z=\mathrm{St}_K[\eta']$ or $\Delta_K[\eta']$, $\mathsf z=\|\eta'\|_{C^2}$. The last bound says that the stress insertion collapses onto $\mathsf b\,\mathcal T_K(H^{\rm cons}[\eta])$ in every periodic Gibbs state.
\item[(f)] \emph{Any number of entries.} Let $0<g\le\gamma$, $n\ge2$, $1\le n_1\le n$, $R>0$, $0<c<\frac1{28}$ and $u=\lceil\frac{1/28-c}{4n-1}K\rceil$. There are $C,p$, depending on $g,n,R,c$, and $K_0$, depending on $g,n,c$, such that for all $(K,\beta)\in\mathfrak D(g)$ with $K\ge K_0$, all $\nu\in\mathcal G^{\rm per}_{K,\beta}$, every $\mathsf d$ with $\mathsf dL^u\ge K^2$, all $\eta_a\in C^2_c(\mathbb R^3;\mathbb R^3)$ $(a\le n_1)$ and $H_a\in C_c(\mathbb R^3;V)$ $(a>n_1)$ with supports in a ball of radius $R$ at pairwise sup distances $\ge\mathsf d$, and for $X_a$ any of $\mathcal T_K(H^{\rm cons}[\eta_a])$, $\mathrm{St}_K[\eta_a]$, $\Delta_K[\eta_a]$, $\mathrm{Rm}_K[\eta_a]$ (Lemma~\ref{C-360-8}(c)), $\mathfrak K^{\rm v}[\beta\eta_a]$ $(a\le n_1)$ and $X_a=\mathcal T_K(H_a)$ $(a>n_1)$,
\[
\Bigl|\nu\Bigl(\prod_aA^\nu[X_a]\Bigr)\Bigr|\le CK^pL^{-cn_1K}\prod_{a\le n_1}\|\eta_a\|_{C^2}\prod_{a>n_1}|H_a|_\infty ,
\]
and the same holds for the joint cumulant under $\nu$; with $H_a=h_a\mathbf 1$ the spectators are action densities. At $n_1=1$, $c=\frac1{56}$ and a fixed $\mathsf d>0$ this holds for $K\ge K_0(\mathsf d)$.
\item[(g)] \emph{Every tensor limit is conserved.} Let $0<g\le\gamma$ and let $\sigma$ be a tensor limit at $g$. Then $\sigma(H^{\rm cons}[\eta],H_2,\dots,H_n)=0$ for every $n\ge2$, $\eta\in C^2_c(\mathbb R^3;\mathbb R^3)$ and $H_a\in C_c(\mathbb R^3;V)$ whose supports are pairwise disjoint and disjoint from ${\rm supp}\,\eta$. In particular $\sigma(H^{\rm cons}[\theta\eta],H^{\rm cons}[\eta])=0$ for ${\rm supp}\,\eta\subset\{x_3>0\}$, $\theta\eta=\theta\circ\eta\circ\theta$, $\theta$ the mirror $x_3\mapsto-x_3$, so that in the reconstruction of $\sigma$ the vector of $\mathcal T(H^{\rm cons}[\eta])$ vanishes; and along the sequence $\mathsf b_j^2\rho_j(A[\mathcal T(H^{\rm cons}[\theta\eta])]A[\mathcal T(H^{\rm cons}[\eta])])-\rho_j(A[\mathfrak K^{\rm v}[\beta_j\theta\eta]]A[\mathfrak K^{\rm v}[\beta_j\eta]])\to0$.

For $0<g\le\gamma^{\rm sd}_W$ tensor limits exist: by Theorem~\ref{C-355-9}(a) of Section~\ref{C-355-1}, the centred products of any number of tensor entries at separated supports are bounded in every periodic Gibbs state by numbers that depend on no cutoff, side, state or position, so every sequence of first-passage data at $g$ with $K_j\to\infty$ and of periodic Gibbs states has a subsequence along which all of them converge. Tensor limits are covariant as tensors under the translations and $W_3$, reflection positive at the nine families of mirrors, and their trace part is the scalar limit, $\sigma(f_1\mathbf 1,\dots,f_n\mathbf 1)$ being the limit of the centred products of the $\mathcal O(f_a)$.
\end{enumerate}
\end{theorem}

\begin{remark}[Normalisation; the role of the collapse]\label{C-360-15}
Conservation at disjoint supports does not fix the normalisation of a tensor: $X_{lr}=|x|^{-2}[F_ln_r+F_rn_l-\delta_{lr}F\cdot n+n_ln_rF\cdot n]$, $n=x/|x|$, is symmetric, traceless and conserved off $0$, with the flux $\frac{16\pi}3F$. And Theorem~\ref{C-360-3}(g) rests on the collapse in (e): the decorrelation of the orbital insertion alone (the second bound of (e)) does not make the six tensor fields conserved; it is the identification of the stress insertion with $\mathsf b\,\mathcal T_K(H^{\rm cons}[\eta])$, up to $\Delta_K[\eta]$, that does.
\end{remark}

\begin{theorem}[Ward identities on tori]\label{C-360-4}
Let $0<g\le\gamma$, $\ell\in\mathfrak L(g)$, $n\ge1$, $\mathsf d>0$, $0<\delta_0<\frac\ell2$ and $c<\frac1{14}$.
\begin{enumerate}
\item[(a)] \emph{Rotations on a fixed torus.} There are $C,p$, depending on $g,\ell,n,\mathsf d,\delta_0,c$, such that for all $K$, $\mathsf b'\in I(g)$, $\beta=L^K\mathsf b'$, $\omega\in\mathfrak{so}(3)$, every cut-off function $\chi$ at $K$ of class $C^{1,1}$ with the margin $\delta_0$, and all real $f_i$ vanishing outside $(-\frac\ell2+\delta_0,\frac\ell2-\delta_0)^3$ and covered by $\chi$:
\begin{enumerate}
\item[(i)] if $n\ge2$, $f_i\in C^{1,1}$ and the spread of the supports is $\ge\mathsf d$ (the supports may touch if no point is common to all of them), $|\mathfrak U_{\omega,\chi}(f)|\le C|\omega|K^pL^{-cK}\|\chi\|_{C^{1,1}}\prod_i\|f_i\|_{C^{1,1}}$;
\item[(ii)] if $n=1$, $f_1$ is Lipschitz and $\chi=1$ on the sup $\mathsf d$-neighbourhood of ${\rm supp}\,f_1$, $|\mathfrak U_{\omega,\chi}(f_1)|\le C|\omega|K^pL^{-cK}\|\chi\|_{C^{1,1}}\|f_1\|_\sharp$;
\item[(iii)] whatever the supports, if $\chi$ and $\chi'$ both equal $1$ on the $\mathsf d$-neighbourhood of $\bigcup_i{\rm supp}\,f_i$, $|\langle\mathfrak B_{\omega,\chi}-\mathfrak B_{\omega,\chi'};\mathcal O(f_\cdot)\rangle|\le C|\omega|K^pL^{-cK}(\|\chi\|_{C^{1,1}}+\|\chi'\|_{C^{1,1}})\prod_i\|f_i\|_\sharp$.
\end{enumerate}
If $n\ge2$, the $f_i$ are of class $C^3$, ${\rm dist}_\infty({\rm supp}\,f_1,{\rm supp}\,f_2)\ge\mathsf d$ and a fixed $\chi$ covers the $f_i$ at all large $K$, then $\mathfrak U_{\omega,\chi}(f)+\langle\mathfrak B_{\omega,\chi};\mathcal O(f_\cdot)\rangle\to\sum_iS^{(\ell,\mathsf b')}_n(f_1,\dots,(\omega y)\cdot\nabla f_i,\dots,f_n)$ at the rate $K^pL^{-K/10}+L^{-K}$, and the flux alone converges to the same limit at the rate $K^pL^{-cK}$; the limit does not depend on $\chi$. With nothing asked of the bulk term, for $\omega=\omega_{13}$, the generator of the turns in the plane of $e_1,e_3$, $f^t=f\circ e^{-t\omega}$, $f_i\in C^3$ supported in a ball inside the margin with ${\rm dist}({\rm supp}\,f_1,{\rm supp}\,f_2)\ge\sqrt3\mathsf d$, and $\chi$ covering all the $f^t$, $0\le t\le\frac\pi2$:
\[
\begin{multlined}
\Bigl|\int_0^{\pi/2}\bigl[\mathfrak U_{\omega,\chi}(f^t)+\langle\mathfrak B_{\omega,\chi};\mathcal O(f^t_\cdot)\rangle\bigr]dt\Bigr|\le Ca_K\prod_i\|f_i\|_{C^3},\qquad\int_0^{\pi/2}\sum_iS^{(\ell,\mathsf b')}_n(f^t_1,\dots,(\omega y) \\
\cdot\nabla f^t_i,\dots,f^t_n)\,dt=0 .
\end{multlined}
\]
\item[(b)] \emph{Rotations along growing tori.} Let $n\ge2$, $R,\mathsf d_0>0$, $\chi_R(y)=\chi_1(|y|/R)$ with $\chi_1\in C^\infty$, $\chi_1=1$ on $[0,\frac43]$ and $\chi_1=0$ on $[\frac53,\infty)$. There are $C,p,K_0$, depending on $g,n,R,\mathsf d_0,c$ and on no side, such that
\[
|\mathfrak U_{\omega,\chi_R}(f)_{K,\ell,\beta}|\le C|\omega|K^pL^{-cK}\prod_i\|f_i\|_{C^{1,1}}
\]
for all $(K,\beta)\in\mathfrak D(g)$ with $K\ge K_0$, every $\ell\in\mathfrak L(g)$ with $2R+2a<\frac\ell2$ and all smooth $f_i$ supported in $\{|y|<\frac R2\}$ at mutual distances $\ge\mathsf d_0$; without its cubic member $\mathfrak U^\sharp[\chi_R\,\omega y](f)$ the bulk term is $\le C|\omega|K^{1+\mathfrak p_0}L^{-c_*K}\prod_i\|f_i\|_{C^{1,1}}$ for every $c_*<\frac12$. The same bounds hold under the weak limit points of the torus laws. Along every sequence $(K_j,\ell_j,\beta_j)$ with $K_j,\ell_j\to\infty$ along which the connected functions converge to $S_n$, $\langle\mathfrak B_{\omega,\chi_R};\mathcal O(f_\cdot)\rangle\to-\frac d{dt}S_n(f(e^{-t\omega}\cdot))|_{t=0}$ for every $R$; so the limit is invariant under rotations if and only if this flux tends to $0$.

For the volume-first scalar limits, which are invariant under ${\rm O}(3)$, the flux through every shell tends to $0$. With the limits taken one after the other: at a label $\mathsf b'\in I(g)$, on each torus $T_\ell$, $\ell\in\mathfrak L(g)$, the limit in $K$ of the flux is $\sum_iS^{(\ell,\mathsf b')}_n(\dots,(\omega y)\cdot\nabla f_i,\dots)$ through every enclosing shell, by (a). If along sides $\ell_k\in\mathfrak L(g)$, $\ell_k\to\infty$, the limits $S^\circ_n(X)=\lim_kS^{(\ell_k,\mathsf b')}_n(X)$ exist for all $n\ge2$ and all tuples $X$ of real Lipschitz functions with compact supports at positive distances, then for smooth $f_i$ with compact supports at positive mutual distances and all cut-off functions $\chi_k$ on $T_{\ell_k}$ of class $C^{1,1}$ that vanish near the seam and cover the $f_i$ (for instance $\chi_R$ at a fixed $R$, or $\chi_{R_k}$ with $R_k\to\infty$ and $2R_k<\frac12\ell_k-1$, a far shell)
\[
\lim_k\lim_K\langle\mathfrak B_{\omega,\chi_k};\mathcal O(f_\cdot)\rangle_{K,\ell_k,L^K\mathsf b'}=\sum_iS^\circ_n\bigl(f_1,\dots,(\omega y)\cdot\nabla f_i,\dots,f_n\bigr)=-\tfrac d{dt}S^\circ_n\bigl(f(e^{-t\omega}\cdot)\bigr)\big|_{t=0},
\]
through any shell; and there are integers $K^*_k$ such that $S^\circ_n$ is the limit of the connected functions along $(K_k,\ell_k,L^{K_k}\mathsf b')$ for every sequence $K_k\ge K^*_k$. For each fixed order $n\ge2$, there is $c_\ell(n,g)>0$ such that
the same limit is obtained by taking the spacing limit first on each
torus whenever $\ell_j\to\infty$ and
$\ell_j\le c_\ell(n,g)K_j^{2\mathfrak p_0/3}$.
Thus $S_n=\lim_jS^{(\ell_j,\mathsf b_j)}_n$ at any first-passage
labels, and raising the cutoffs at fixed sides and labels does not
change this limit. The constant is specified in the proof. For every $K\ge1$, the labels of the first-passage data satisfy
\[
 I(g)\subset
 \{L^{-K}\beta:(K,\beta)\in\mathfrak D(g)\}
 \subset
 \bigl(g^{-2}+B_\sharp-L^{-1}\bar B,\,
       L(g^{-2}+\bar B)\bigr].
\]
\item[(c)] \emph{One stress entry: trace and diagonal.} Let $0<\mathsf b_-\le\mathsf b\le\mathsf b_+$, $\beta=L^K\mathsf b$, $\ell=2L^m$ any side, and let two of the $n+1$ supports be $\ge\mathsf d$ apart, the others may touch; $\lambda_L=\log(L/2)/\log L$ and $c_n=\lambda_L/(4(6+3n))$. For $\gamma$ small enough in terms of the other constants, and with $C,p$ depending on $\ell,\mathsf b_\pm,n,\mathsf d$,
\[
|\langle T_{mm}(\varphi)-\mathsf b\check D_m(\varphi);\mathcal O(f_\cdot)\rangle|+|\langle(T_{11}+T_{22}+T_{33})(\varphi)-\mathsf b\,\mathcal O(\varphi);\mathcal O(f_\cdot)\rangle|\le CK^pL^{-c_nK}\|\varphi\|_\sharp\prod_i\|f_i\|_\sharp,
\]
so that $-T_{33}=\mathsf b[\mathcal O_3-\mathcal O_1-\mathcal O_2]$ inside these functions. For every $\epsilon>0$, $\langle(T_{lr}-T_{rl})(\varphi);\mathcal O(f_\cdot)\rangle=O(a^{1-\epsilon}){\rm Lip}\,\varphi$ and $|\langle T_{lr}(\varphi);\mathcal O(f_\cdot)\rangle|\le C_\epsilon a^{-\epsilon}\sup|\varphi|\prod_i\|f_i\|_\sharp$; and $\langle\mathfrak K^{\rm v}[\beta\eta]-\beta\mathfrak K^\sharp[\eta];\mathcal O(f_\cdot)\rangle=O(K^pL^{-c_nK}){\rm Lip}\,\eta$ for Lipschitz vector fields $\eta$, in contact with the $f_i$ or not. At pairwise separated supports and a fixed label $\lim_K\langle T_{mm}(\varphi);\mathcal O(f_\cdot)\rangle$ exists.
\item[(d)] \emph{All Ward identities with one stress entry.} Let $\gamma$ be as in (c), $\gamma^{-2}\le\mathsf b_-\le\mathsf b\le\mathsf b_+$, $MR(\mathsf b_+)\le L^m$, $\ell=2L^m$; let $\xi=e_r$, $\omega y$ or $y$, where for $\xi=\omega y$ and $\xi=y$, taken in $(-\frac\ell2,\frac\ell2)^3$, the functions $\chi$ and $f_i$ vanish within a fixed distance of $\partial(-\frac\ell2,\frac\ell2)^3$; and either $n\ge2$, $f_i\in C^2$ with pairwise disjoint supports, $I\subset\{1,\dots,n\}$ and $\chi\in C^2$ equal to $1$ near ${\rm supp}\,f_i$ for $i\in I$ and to $0$ near the others, or $n=1$, $I=\emptyset$ and $\chi=0$ near ${\rm supp}\,f_1$. Then
\[
\begin{multlined}
\lim_K\Bigl\langle\sum_{l,r}T_{lr}(\xi^r\partial_l\chi);\mathcal O(f_\cdot)\Bigr\rangle_{K,\ell,L^K\mathsf b}=-\sum_{i\in I}S^{(\ell,\mathsf b)}_n(f_1,\dots,\xi\cdot\nabla f_i,\dots,f_n) \\
+\delta_{\xi,y}\bigl(|I|S^{(\ell,\mathsf b)}_n(f)-\mathsf b\,S^{(\ell,\mathsf b)}_{n+1}(\chi,f)\bigr),
\end{multlined}
\]
and at the cutoff the two members differ by $\mathfrak U^\sharp[\chi\xi](f)$ up to $O(K^pL^{-c_nK})$ and up to the convergence of the right member. Since $|\mathfrak U^\sharp[\chi\xi](f)|\le CK^pL^{-cK}$ for every $c<\frac1{14}$ (Lemma~\ref{C-360-10}, the spread of $\{\chi\xi\ne0\}$ and of the supports being positive), the right member converges at the rate $K^pL^{-K/10}$ and $c_n<\frac1{36}$, the left member converges at the rate $K^pL^{-c_nK}$, the slowest of the three. For $I=\emptyset$ this is the conservation of the stress tensor beside action densities; for $\xi=e_r$ and $\chi$ a periodic function of $y_3$ alone, $\langle T_{3r}(\chi');\mathcal O(f_\cdot)\rangle\to-\sum_{i\in I}S^{(\ell,\mathsf b)}_n(\dots,\partial_rf_i,\dots)$. Thus, inside these cumulants on one torus, $T_{3r}$ integrated over the planes $\{y_3={\rm const}\}$ acts on the enclosed entries as $-\partial_r$, so that $-T_{33}$, the Euclidean energy density, acts as $\partial_3$.
\item[(e)] \emph{Translations.} Let $(K,\beta)\in\mathfrak D(g)$, $n\ge2$, $\rho$ a direction and $\mathsf d'>0$. For real $f_i\in C^2$ with Lipschitz $\partial^2_\rho f_i$ and ${\rm dist}({\rm supp}\,f_1,{\rm supp}\,f_2)\ge\mathsf d'+2a$,
\[
\begin{multlined}
\Bigl|\beta\langle\mathfrak K^\sharp[e_\rho];\mathcal O(f_\cdot)\rangle-a^{-1}\sum_i\langle\mathfrak K^{\rm v}[f_ie_\rho];(\mathcal O(f_l))_{l\ne i}\rangle\Bigr|=\Bigl|\sum_iS^T_{n;K,\ell,\beta}(f_1,\dots,\nabla^s_\rho f_i,\dots,f_n)\Bigr| \\
\le(n+1)2^nC_n(\mathsf d')e^{nc_V\ell^3}a_K^2\prod_iN_i,
\end{multlined}
\]
$N_i=\max\{\|f_i\|_\sharp,\|\partial_\rho f_i\|_\sharp,\frac12\|\partial^2_\rho f_i\|_\sharp\}$, with $C_n(\mathsf d')$ and $c_V$ as in
Lemma~\ref{lem:torus-cumulant-bound}; at $n=2$ the left member vanishes identically. For $I\subset\{1,\dots,n\}$, the $f_i$ with $i\in I$ of class $C^3$ with Lipschitz $\partial^3_\rho f_i$, the others Lipschitz, ${\rm dist}({\rm supp}\,f_1,{\rm supp}\,f_2)\ge\mathsf d'+a$, and $\chi$ a real function on $T_\ell$ with $\chi(x(p))=\chi(x_b)=1$ at every plaquette $p$ at which some $f_i$, $i\in I$, does not vanish and at its four bonds $b$, and $\chi(x(p))=\chi(x_b)=0$ at every plaquette $p$ at which some $f_i$, $i\notin I$, does not vanish and at its four bonds $b$,
\[
\langle N_\beta[\chi e_\rho];\mathcal O(f_\cdot)\rangle-a^{-1}\sum_{i\in I}\langle\mathfrak K^{\rm v}[f_ie_\rho];(\mathcal O(f_l))_{l\ne i}\rangle=-\sum_{i\in I}S^T_{n;K,\ell,\beta}(f_1,\dots,\partial_\rho f_i,\dots,f_n)+\mathsf r,
\]
$|\mathsf r|\le\frac{a_K^2}6C_n(\mathsf d')e^{nc_V\ell^3}\sum_{i\in I}\|\partial^3_\rho f_i\|_\sharp\prod_{l\ne i}\|f_l\|_\sharp$, $N_\beta[\chi e_\rho]=\mathbb T_\beta[\chi e_\rho]+\beta\mathfrak K^\sharp[\chi e_\rho]$ the null term of Lemma~\ref{C-360-8}(a). Thus the stress functional $\mathbb T_\beta[\chi e_\rho]$ of a cut-off translation, taken together with its cubic part $\beta\mathfrak K^\sharp[\chi e_\rho]$ and with the contact terms $-a^{-1}\sum_{i\in I}\langle\mathfrak K^{\rm v}[f_ie_\rho];(\mathcal O(f_l))_{l\ne i}\rangle$, generates the translation of the enclosed entries up to $O(a^2)$, with connected functions bounded uniformly in $K$. If the closed $a$-neighbourhoods of the supports lie in balls of radius $R$ and $\ell\ge16n(R+1)$, the numbers may be taken free of $\ell$.
\end{enumerate}
\end{theorem}

\begin{remark}[The spread is needed]\label{C-360-16}
For $n\ge3$ the bound (a)(i) fails at supports with a common point. The symmetric difference is no derivation: by the invariance of the measure under the lattice translations, $\sum_iS^T_n(\dots,\nabla^s_\rho f_i,\dots)$ is $a^2$ times connected functions of difference quotients of the $f_i$, at $n=3$ exactly
\[
\sum_{i=1}^3S^T_3(\dots,\nabla^s_\rho f_i,\dots)=\tfrac{a^2}2\,S^T_3(\nabla^+_\rho f_1,\nabla^+_\rho f_2,\nabla^+_\rho f_3),
\]
and connected functions of action densities at coinciding points are of the order $a^{-3}$ (Proposition~\ref{C-360-21}(c)). For free photons, let the real $f_i$ have compact supports
inside the margin and bounded derivatives through order five,
and let $\chi=1$ on a neighbourhood of their supports. Then
$a\,\mathfrak U_{\omega,\chi}(f)\to3\mathsf b^{-3}\int\sum_\rho(\omega y)_\rho\,\partial_\rho f_1\,\partial_\rho f_2\,\partial_\rho f_3$ at $n=3$, while the identity with the true generator survives at supports with a common point: $\sum_iS^T_n(\dots,(\omega y)\cdot\nabla f_i,\dots)-\langle\mathfrak B_{\omega,\chi};\mathcal O(f_\cdot)\rangle$ is $O(a)$ at $n=2$ and $O(a\log\frac1a)$ at $n=3$. For ${\rm SU}(2)$, at every $g\le\gamma$, on every side $\ell\in\mathfrak L(g)$ and for every $n\ge3$ there are smooth $f_1,\dots,f_n$ with a common point, a rotation $\omega$ and a cut-off function $\chi$ such that
\[
450\,(n-1)!\,\mathsf b'^{-n}\le a_K|\mathfrak U_{\omega,\chi}(f)|\le2.3\cdot10^5\,(n-1)!\,\mathsf b'^{-n}
\]
for all large $K$, uniformly in the labels $\mathsf b'\in I(g)$. At $n=1$ the neighbourhood in (a)(ii) is needed:
there are a real Lipschitz $f_1$, a rotation $\omega$ and a
cutoff $\chi\in C^{1,1}$ covering $f_1$, equal to one on
its support but on no neighbourhood of it, for which the
free-photon model has
$\liminf_K\mathfrak U_{\omega,\chi}(f_1)>0$.
\end{remark}

\begin{theorem}[One plane field between two times]\label{C-360-5}
Let $0<g\le\gamma^{\rm ref}_W$, $\ell'\in\mathfrak L(g)$, $(K_j,\beta_j)\in\mathfrak D(g)$ with $K_j\to\infty$, $\nu_j\in\mathcal G^{\rm per}_{K_j,\beta_j}$, and $(X)_j=\prod_i[\mathcal O_{K_j}(f_i)-\nu_j(\mathcal O_{K_j}(f_i))]$. Every subsequence has a further subsequence along which the following hold.
\begin{enumerate}
\item[(i)] $\sigma(X)=\lim_j\nu_j((X)_j)$ and $\sigma^{(m)}(u;X)=\lim_j\nu_j((X)_j[\check D_{m,j}(u)-\nu_j(\check D_{m,j}(u))])$ exist for all real $u,f_i\in C^1_c(\mathbb R^3)$ with disjoint supports; they are real and multilinear, $\sum_m\sigma^{(m)}(u;X)=\sigma(u\cdot X)$, the product with one more scalar factor, they are translation invariant and covariant under $W_3$: a quarter turn that carries $e_3$ to $e_i$ carries $\sigma^{(3)}$ to $\sigma^{(i)}$.
\item[(ii)] On the Hilbert space $\mathcal H(\sigma)$ of the reconstruction data of the limit, with the Hamiltonian $H\ge0$, the reflection $\theta$, the translations $U(\vec x)$, whose generators $P_1,P_2$ commute with $H$, and the classes $[G]$, the joint spectrum lies in the cones $H\ge|P_k|$, $k=1,2$: every periodic Gibbs state obeys the five-point inequality \textup{(\ref{IV-54-29})} of Lemma~\ref{lem:spectral-steps} at the mirror normal to $e_3$ and the direction $e_k$, which passes to the limit and puts the spectral measure of every class on $\{\cosh(tE)+\cos(tq_k)\ge2\}$ for a dense set of $t>0$, hence on $\{E\ge|q_k|\}$. So $H\ge0$ and $H\ge|P_k|$ are the hypotheses of Lemma~\ref{C-360-12} in the direction $x_3$ and, in its last sentence, in the directions $x_1,x_2$; and there is a unique family of bounded operators $B^{(m)}(s,s')$, $s,s'>0$, with $B^{(m)}(s,s')^*=B^{(m)}(s',s)$, $e^{-tH}B^{(m)}(s,s')e^{-t'H}=B^{(m)}(s+t,s'+t')$ and
\[
\sigma^{(m)}\bigl(u;\theta\overline{\tau_{se_3}G'}\cdot\tau_{s'e_3}G\bigr)=\int u(x)\,\bigl\langle[G'],U(\vec x)B^{(m)}(s+x_3,s'-x_3)U(-\vec x)[G]\bigr\rangle\,dx
\]
for $u$ supported in $\{-s<x_3<s'\}$ and $G,G'$ supported in $\{x_3>0\}$.
\item[(iii)] $\|B^{(m)}(s,s')\|$ is bounded by a finite function of $\min\{s,s',1\}$, $g$ and $\ell'$, and for $u$ supported in $\{\gamma'-s\le x_3\le s'-\gamma'\}$ the left member of (ii) is at most $c(\gamma')\|u\|_{L^1}\|[G']\|\,\|[G]\|$. For $0<g\le\gamma^{\rm sd}_W$ every such limit with converging labels arises, along a further subsequence, from a tensor limit $\sigma_T$ (Theorem~\ref{C-360-3}(g)): $\sigma(X)=\sigma_T(f_1\mathbf 1,\dots,f_n\mathbf 1)$ and $\sigma^{(m)}(u;X)=\sigma_T(u(\mathbf 1-2E_{mm}),f_1\mathbf 1,\dots,f_n\mathbf 1)$. Let $\mathcal H(\sigma_T)$ be the reconstruction of $\sigma_T$ at $\theta$, spanned by the classes $[F]$ of the tuples $F$ of entries in $C_c(\{x_3>0\};V)$ with $\langle[F'],[F]\rangle=\sigma_T(\theta\bar F'\cdot F)$, with its Hamiltonian and its translations, whose joint spectrum lies in the same cones by the same five-point inequality; its \emph{scalar subspace} $\mathcal H_0$, the closed span of the classes of the tuples of scalar entries $f_a\mathbf 1$, is $\mathcal H(\sigma)$ and reduces the Hamiltonian and the translations, and $\Pi_0$ is the projection onto it.

For $e\in V$ the \emph{point operator} of the tensor entry at $e$ is the family $B^e(s,s')$ of bounded operators on $\mathcal H(\sigma_T)$ that Lemma~\ref{C-360-12} gives for $u\mapsto\sigma_T(ue,\cdot)$, with $\sigma_T(ue,\theta\overline{\tau_{se_3}F'}\cdot\tau_{s'e_3}F)=\int u(x)\langle[F'],U(\vec x)B^e(s+x_3,s'-x_3)U(-\vec x)[F]\rangle\,dx$ as in (ii); it is linear in $e$, and $\|B^e(s,s')\|\le C|e|\min\{s,s',1\}^{-p}$ for every $p>\frac{17}2$, $|e|$ the Euclidean norm of $e$ and $C$ depending on $g$, $p$ and the limit label and not on the side, by the bound on the chessboard seminorm of a centred tensor entry. Then $B^{(m)}(s,s')=\Pi_0B^{\mathbf 1-2E_{mm}}(s,s')\Pi_0$ on $\mathcal H_0$, the compression to the scalar subspace of the point operator at $\mathbf 1-2E_{mm}$, so that $\|B^{(m)}(s,s')\|\le C\min\{s,s',1\}^{-p}$; and for $Z=(u;X)$ and $R_\alpha$ the turn by the angle $\alpha$ in the plane of $e_i,e_3$,
\[
\sigma^{(3)}(Z\circ R_\alpha)+\sigma^{(3)}(Z\circ R_{-\alpha})=2\cos^2\alpha\,\sigma^{(3)}(Z)+2\sin^2\alpha\,\sigma^{(i)}(Z).
\]
Thus the $\check D_m$ behave under turns as the diagonal of a symmetric tensor; moreover $\alpha\mapsto\sigma^{(3)}(Z\circ R_\alpha)$ is of the form $c_0+c_1\cos2\alpha+c_2\sin2\alpha$ with numbers $c_0,c_1,c_2$, and its second derivative at $\alpha=0$ is $2[\sigma^{(i)}(Z)-\sigma^{(3)}(Z)]$.
\end{enumerate}
\end{theorem}

One step of the renormalisation group, at every level and with its never-large term, runs with three complex plane couplings, about its prescribed background with free centres (Theorem~\ref{C-360-13}); on the labels of histories with large-field regions the moved backgrounds, the centres of the steps and the counterterms are those of Theorem~\ref{C-360-14}. For a triple $P=(U,A',J)$ at the level $n$, $\mathbf U_P=e^{i\eta_nA'}U$, $\eta_n=L^{-n}$, is its complex configuration (Definition~\YMMref{49.1}{374}); on a region $X$ let ${\rm hol}_n(P;X)$ be the H\"older-$\frac12$ number, transports included, of $L^{2n}(\partial\mathbf U_P-1)$ at distances $\le1$ in level-$n$ units; the \emph{core with the third radius} $\mathcal U_n(X)$ is the set of triples in the hereditary core of Definition~\YMMref{49.1}{374} at the radii $(\alpha_0,\alpha_1)$, with all the tower conditions retained as in Section~\ref{sec:analytic-core-complements}, with ${\rm hol}_n(P;X)<\alpha_2$; $\rho^{k\to j}_Z(U,A',J)=(U,\lambda A',\lambda^3J)|_Z$, $\lambda=L^{-(k-j)}$; $\mathcal D_n$ and $d_n$ are as in Section~\YMMref{45}{351}.

A region $Z\in\mathcal D_j$ is \emph{short} at $k$ if $d_j(Z)+4\le L^{k-j}$. The \emph{box} $Z^{\rm box}$ of a short $Z$ is the least box of level-$j$ cubes that holds a lift of $Z$ (a wall-connected union of cubes of the covering that the covering map carries one to one, cube by cube, onto $Z$; such unions exist, $Z$ being wall-connected, and any two of them differ by a period), enlarged by one level-$j$ cube on each of its six sides and carried back by the covering map: it does not depend on the lift, its sides are at most $d_j(Z)+4\le L^{k-j}$ level-$j$ cubes, it lies in $\mathcal D_j$, holds $Z$, and holds both end points of every bond with an end point in $Z$. Then $X_Z$ is the union of the level-$k$ cubes that hold a level-$j$ cube of $Z^{\rm box}$: one or two along each axis, so at most eight, with $X_Z\in\mathcal D_k$ and $d_k(X_Z)=0$.

For every $Z\in\mathcal D_j$, $Z^{\rm up}$ is the union of the level-$k$ cubes that hold a level-$j$ cube of $Z$. A \emph{term} on $\mathcal U_j(Z)$ is a bounded function on it that is holomorphic on its slices, a function of the pair $(\mathbf U_P,J)$, and gauge invariant.

\emph{The extended list.} Theorems~\ref{C-360-13} and~\ref{C-360-14} are stated at the choices of the following extension of the list $\mathcal T^{\rm S}[2\varrho]$ of Theorems~\YMMref{AJ}{673} and~\YMMref{AD}{607} (the common real and sourced choice of Section~\YMMref{82}{665}, including \eqref{eq:canonical-source-quarter-room}). The base choice also retains \eqref{eq:array-borel-radius} and the
radius, decay and block-number conditions of
Lemma~\ref{lem:ordinary-holder-kernels}. These are imposed when
choosing the response parameters in the proof of
Proposition~\YMMref{69.1}{595}, before the polymer-window size $M$:
the finitely many local constants here depend only on $L$, so the
response block number $\mathfrak M$ can be raised and its radii
and decay rate reduced at that stage. For Lemma~\ref{C-360-22}
take this same response block number at least $M_{\rm ps}$ and
$M_H$, the floors in Theorem~\YMMref{I}{225} and
Lemma~\ref{lem:ordinary-holder-kernels}. Enlarge $w_{\rm lin}$
to at least $\max\{g_{\rm p},w_S\}$, where $w_S$ is from
Proposition~\YMMref{29.1}{231}, and reduce $\alpha_{\rm lin}$
so that $2\alpha_{\rm lin}\le a_Q$ and the four local-system
radii are at most $a_H$. These choices depend only on $L$.
The separated setting has $R_1\ge L$ and terminal torus side at
least $6R_1\mathfrak M$, hence at least $6LM_{\rm ps}$.
On a real regular label, $F'(0)=J(U)$ at full bonds; the scalar
part of $D_U^{\mathsf T}\partial U$ is bounded by
$\xi^{-1}\max_{p\ni b}|\partial U(p)-I|^2$.
Thus equation~\YMMref{(579)}{288} gives the full matrix-divergence
bound at twice the regularity radius. The ceiling on
$C_m\varepsilon$ in part (c) supplies this radius and the four
local-system radii, while its width ceiling supplies the enlarged
$w_{\rm lin}$. These are conditions on the common response
choice; they are not new analytic assumptions.
It adds $108$ entries, in this order: four lower bounds on chosen constants among the construction results, each by a function of earlier constants; $85$ further determined constants, each a positive explicit formula in earlier constants, placed before $\gamma$ (four abbreviations are used among the formulas); and $15$ further upper bounds on $\gamma$. No condition is laid on $L$, on $E_0$ or on a window constant.

The following constants retain their meanings in the real construction and its sourced extension (the proof of Proposition~\YMMref{69.1}{595} and Section~\YMMref{82}{665}): 
\[
\begin{multlined}
\mathsf C_{\rm ex},\mathsf C_{\rm cs},\mathsf C_{\rm mn},\lambda_{\rm s},\gamma_{\rm r},c_{\rm r},\mathsf u_{\rm G},\lambda_{\rm r,gs},\bar c_{\rm P},\beta_{\rm s},\kappa,\kappa'',\bar r_H,\bar\kappa,\lambda_V,\vartheta_\ast,\mathsf c_{\rm os},\mathsf c_{\mathfrak h},\varepsilon_{\mathfrak h},r^\circ_P,\mathsf C_{\mathcal A},\mathsf s^\star,K_\Lambda,N_\star,K_P, \\
w_\star,\varepsilon_{\mathcal W},\mathsf o,B^{(\mathrm{cg})},A_0,A_1,\mathfrak p_0,\alpha_0,\alpha_1,B_3,E_0,M,L
\end{multlined}
,\]
the plaquette constant $\mathsf c_{\rm pl}=600$ and the number $C_P$ of equation~\YMMref{(799)}{390}, $K_{\rm c}=96(1+2\mathsf c_{\rm pl})^2C_{\mathrm E}$, $C_{\mathrm E}=4096$, the numbers $\mathsf b_{\mathrm D}$, $\mathsf B_\sharp(0)$, $\mathsf B_\sharp(1/2)$ and $\mathsf B_\sharp=\mathsf B_\sharp(0)$ of Proposition~\ref{prop:ordinary-holder-response}, and $a^\sharp$, $K^\sharp$ of the tower hypotheses of Section~\ref{sec:analytic-core-complements}.

Here $\lambda_{\rm r,gs}$ denotes the canonical history Gaussian gap,
written $\lambda_{\rm r}$ in equation~\YMMref{(1137)}{495}
with the specialization of equation~\YMMref{(1161)}{501}:
\[
 \lambda_{\rm r,gs}
 =\min\left\{\frac{\gamma_0}{2},
       \left(1+2(B_{\rm nm}^2/\gamma_0)^2\right)^{-1}\right\}.
\]
The symbols $B_{\rm nm}$ and $\gamma_0$ have the values in that
specialization. This history gap is distinct in notation from
the plain-step gap $\lambda_{\rm r}^{\rm pv}$ defined below.

The entries are:
\begin{enumerate}
\item[\textup{1.}] the lower bound $\mathsf C_{\rm ex}\ge\max\{2\mathsf C_{\rm cs},8\}$;
\item[\textup{2.}] directly behind $\bar c_{\rm P}$: $\lambda^{\rm pv}_{\rm r}:=\min\{\frac12\gamma_{\rm r},(1+2c_{\rm r}^2)^{-1}\}$ and $c_q:=\max\bigl\{\bar c_{\rm P}+\mathsf C_{\rm mn}^3/\lambda_{\rm r,gs}+2,\ \mathsf C_{\rm mn}^3(\frac12\mathsf u_{\rm G}+c_{\rm r}+3)/\lambda^{\rm pv}_{\rm r}+2\bigr\}$;
\item[\textup{3.}] with the abbreviation $E_q:=1+2\lambda_{\rm s}+\mathsf C_{\rm cs}(c_q+2)$, the lower bounds $(1+4\beta_{\rm s})\kappa\ge E_q$ and $\frac12\kappa\ge2\lambda_{\rm s}+8\log2+\frac{\mathsf C_{\rm cs}}LE_q$;
\item[\textup{4.}] directly behind $\bar r_H$: $r^q_H:=\kappa''+\lambda_{\rm s}+\mathsf C_{\rm cs}(c_q+2)$, $\kappa':=\kappa''+1$, $r^q_{\rm w}:=\max\{\frac{\mathsf C_{\rm cs}}Lr^q_H+\lambda_{\rm s},\frac{8\kappa'}L\}$, $r_t:=\lambda_{\rm s}$, $r^q:=r^q_{\rm w}+\lambda_{\rm s}$, $r_1:=\kappa''-8\log2$, $r'_1:=\kappa''-16\log2$;
\item[\textup{5.}] the lower bounds $\bar\kappa\ge1+\mathsf C_{\rm cs}r^q_H+\log\mathsf C_{\rm cs}$ and $\bar\kappa\ge3+16\kappa'+2\lambda_{\rm s}+2\lambda_V$;
\item[\textup{6.}] the lower bound $\mathsf C_{\rm mn}(LM)^3\ge c_q$;
\item[\textup{7.}] directly before $\gamma$, in this order: $c_K:=2/\gamma_{\rm r}+4(c_{\rm r}+\frac12)^2/\gamma_{\rm r}^2$; $B':=6\cdot10^5c_K(LM)^3e^{16\kappa'}$; $S':=4096e^{8\kappa'}$; $b_A:=6.1201^2M^3$; $K^{\rm pl}_T:=\max\{12L^4M^3K_{\rm c},2b_A,112b_A+\frac12M^3\}$; $a_Z:=\mathsf c_{\rm os}\alpha_0$; $r^{\rm pl}:=\min\{r^\circ_P,\varepsilon_{\mathfrak h}/\sqrt2,1/(3\mathsf c_{\mathfrak h}),a_Z/(2\mathsf c_{\mathfrak h})\}$; $\delta':=\min\{1,\varepsilon_{\mathfrak h}/\mathsf C_{\mathcal A},\mathsf s^\star\alpha_0/(1+\beta_{\rm s}),A_1(2\mathfrak p_0/e)^{\mathfrak p_0}\}$; $r^{\rm pl}_2:=\mathsf C_{\mathcal A}\delta'$; $r^{\rm pl}_3:=\min\{r^{\rm pl},r^{\rm pl}_2\}$; $\bar u_f:=K_\Lambda\mathsf s^\star N_\star$; $\varrho^{\rm pl}:=\min\{\frac14r^{\rm pl},(r^{\rm pl})^3/(256B'K^{\rm pl}_TS')\}$; $\sigma_3:=\min\{(r^{\rm pl})^2/(64B'K^{\rm pl}_TS'),\ r^{\rm pl}\varrho^{\rm pl}/(4B'K^{\rm pl}_TS')\}$; $\varepsilon_1:=[8\mathsf C_{\rm cs}^2L^3e^{\mathsf C_{\rm cs}(\kappa''+c_q+r^q_H)}]^{-1}$; $v_0:=(2^{17}L^3)^{-1}$; $v_d:=(2^{17}L^3)^{-1}$; $\hat\rho:=\varepsilon_1/(8e)$; $\eta:=1$; $\gamma^q_2:=(LM)^{-3}$; $\theta_m:=\frac12\vartheta_\ast$; $\varepsilon^{\rm pl}_m:=\min\{\frac12\vartheta_\ast,(LM)^{-3/2}\}$; $\theta'_1:=\frac12\theta_m$; $\theta''_1:=\frac14(1/\varepsilon^{\rm pl}_m+4\varepsilon^{\rm pl}_m/\theta_m^2)^{-1}$; $v_{\rm ce}:=\mathsf C_{\mathcal A}$; $v^+:=\sqrt2\,\mathsf C_{\mathcal A}+1$; $a_{\rm cap}:=\varepsilon_1v_0/(96e)$; $c^+_1:=\min\{\frac14r^{\rm pl}_3,\ a_{\rm cap}(r^{\rm pl})^3/(128K^{\rm pl}_TM^3)\}$; $K_Z:=4B'K^{\rm pl}_TS'/r^{\rm pl}$; $B_0:=6\cdot10^5c_K(LM)^3$; $K_6:=1296M^3\max\{1,2C_P\}B_3L^{-2}$; with the abbreviation $\theta_*:=\min\{\theta'_1,\theta''_1\}$, 
\[
\begin{multlined}
\sigma_e:=\min\{\sigma_3,\ \relax a_{\rm cap}(r^{\rm pl})^2/(16K^{\rm pl}_TM^3),\ \relax c^+_1/(2K_Z),\ \relax a_Z\theta_*/(4\sqrt2K_Z\mathsf c_{\mathfrak h}v^+L), \\
\ \relax \theta_*A_1/(4B_0K_6(2L)^{1/2}A_0)\}
\end{multlined}
;\]
$\bar K_{\rm or}:=\mathsf C_{\rm cs}(LM)^3e^{\mathsf C_{\rm cs}(c_q+r^q_H+\kappa'')}$; $Q^{\rm pl}:=e\bar K_{\rm or}/v_d$; $Q':=e^{\kappa-r^q}$; $\bar K_{\rm w}:=\frac12(6(LM)^3)^{1/2}$; $c^{\rm pl}:=\mathsf C_{\rm mn}^{-1}\gamma^q_2(1-\theta_m)^2A_1^2/4$;
\item[\textup{8.}] then, in this order: $\theta_{\rm w}:=\min\{1,E_0e^{-(\kappa''-\kappa)}\}$; $c^H_{\mathfrak h}:=\max\{54C_P\mathsf b_{\mathrm D}e^{8\bar\kappa}\mathsf B_\sharp(\frac12)/\mathsf B_\sharp,\ 9\mathsf c_{\mathfrak h}\}$; $c''_0:=5200(c^H_{\mathfrak h}+\mathsf c_{\mathfrak h})$; $\alpha^{\rm pl}_2:=2c''_0\varepsilon_{\mathfrak h}$, the third radius $\alpha_2$ of the cores above; $c''_{\mathfrak h}:=c''_0+102\,\mathsf c_{\mathfrak h}L^{-5/2}\alpha^{\rm pl}_2$; $K^0_H:=1.02\cdot746$; $K_H:=\max\{2K^0_HB_3,10\mathsf c_{\mathfrak h}M^3\}$; $C_{\rm gs}:=\max\{\alpha_0/\varepsilon_{\mathfrak h},2\mathsf c_{\mathfrak h}/\mathsf c_{\rm os},1200\mathsf c_{\mathfrak h}\}$; with the abbreviations $m_0:=\min\{\alpha^{\rm pl}_2,\alpha_0\}$ and $m_1:=\max\{\alpha^{\rm pl}_2,\alpha_0\}$: $c'':=\max\{C_{\rm gs},c''_{\mathfrak h}\}$; $c^a_3:=L^{-5/2}+c''r^{\rm pl}/m_0$; $c^a_4:=\max\{c''/m_0,1/r^{\rm pl}\}$; $C^{\rm pl}_V:=2^{10}LC_{\mathrm E}(1-1/L)^{-1}$; 
\[
K^0_J:=\max\{1,\ \relax 10\mathsf c_{\mathfrak h}M^3L^{-5/2}m_1,\ \relax 8(1+84C_P(M/L)^3)L^{-5/2}m_1\}
;\]
$K_J:=2K^0_J$; $K^a_{\rm c}:=2B'S'K_J$; 
\[
K_{\rm ta}:=2^8\cdot24(LM)^3[K_JK^a_{\rm c}+2K^{\rm pl}_T(K^a_{\rm c})^2/(r^{\rm pl})^2]
;\]
$a'_Z:=\min\{a_Z,\mathsf c_{\mathfrak h}\varepsilon_{\mathfrak h}\}$; $\sigma'_5:=a'_Z\theta_*/(4\sqrt2K_Z\mathsf c_{\mathfrak h}v^+L)$; $\sigma^a_*:=\min\{\sigma_3,\sigma_e,\sigma'_5,1\}$; $c^{\rm pl}_g:=2\max\{1/\varepsilon_{\mathcal W},K_H\}$; 
\[
\begin{multlined}
w_5:=\min\bigl\{K^0_J/(1.01\,C^{\rm pl}_Vc^a_4K_{\rm ta}),\ \relax \min\{\varrho^{\rm pl},\frac18r^{\rm pl},K_Z\sigma^a_*\}/K^a_{\rm c},\ \relax \sigma^a_*K^{\rm pl}_T/(6144M^3\alpha_0^2), \\
\ \relax (\sigma^a_*K^{\rm pl}_T/(2C^{\rm pl}_V(c^a_3)^2K_{\rm ta}))^{1/2},\ \relax \sigma^a_*,\ \relax m_0\theta_*A_1/(2K^a_{\rm c}c^{\rm pl}_g(2L)^{1/2}A_0)\bigr\}
\end{multlined}
;\]
$\sigma_5:=\frac12w_5$; $\tau_5:=\frac14w_5$; $b_r:=\mathsf oLMB^{(\mathrm{cg})}\alpha_0$; $a_3:=8.4(\alpha_1(1+2b_r)+b_r)$; $C^Q_2:=29.4a_3^2+14.7a_3$; $\bar C_2:=\max\{C^Q_2,1\}$; $a_4:=7.4a_3+C^Q_2$; $C^Q_H:=18$; $c^Q_s:=\min\{1,\mathsf oL,1/(12Ma_3)\}$, the constant $c_s$ of Theorem~\ref{C-360-13}(a); $c^Q_3:=\min\{\frac14a^\sharp,\alpha_0/(4K^\sharp),\alpha_1,\frac14,\alpha^{\rm pl}_2/18\}$; 
\[
\lambda_Q:=\min\{1/(4a_3),\ \relax a_3/C^Q_2,\ \relax \alpha_1/a_4,\ \relax a^\sharp/(2\alpha_0+a_4),\ \relax [\alpha_0/(K^\sharp(2\alpha_0+a_4^2))]^{1/2}\}
;\]
 
\[
C_Q:=\max\{2\lambda_Q^{-4},\ \relax 32(c^Q_s)^{-4},\ \relax 72M^2\bar C_2^2(c^Q_3)^{-2}\}
;\]
$K_{\rm Lf}:=2C_{\mathrm E}(140C_Q+1)$; $u^c:=2^9\cdot129\,K_{\rm Lf}L/(L-1)$; $\bar u^c_f:=\bar u_f+u^c$; $c^c_J:=2^{21}L^3K^a_{\rm c}(7M^3/r^\circ_P+\bar u^c_f/r^{\rm pl}_3)$; $w_6:=\min\{w_5,1/c^c_J\}$, the constant $c_{\rm w}$ of Theorem~\ref{C-360-13}(c); $\sigma_6:=\frac12w_6$; $\tau_6:=\frac14w_6$;
\item[\textup{9.}] fifteen upper bounds on $\gamma$, each of the form $\sup_{x\ge\gamma^{-2}}F(x)\le\Phi$ with $F(x)\to0$ as $x\to\infty$ and $\Phi>0$ formed with earlier constants; with the thresholds $\delta(x)=A_1x^{-1/2}(\log x)^{\mathfrak p_0}$, $\varepsilon(x)=A_0x^{-1/2}(\log x)^{\mathfrak p_0}$ of \YMMref{(634)}{339} and 
\[
\xi(x):=\exp(-\mathsf C_{\rm mn}^{-1}\gamma^q_2(1-\theta_m)^2A_1^2(\log x)^{2\mathfrak p_0})
\]
 they are: $3K_P\delta(x)\le a_{\rm cap}$; $\bar u^c_f\mathsf C_{\mathcal A}\delta(x)/r^{\rm pl}_3\le a_{\rm cap}$; $\frac{64}3(7M^3+\bar u^c_f)c^+_1\mathsf C_{\mathcal A}\delta(x)/(r^{\rm pl}_3)^2\le\frac13a_{\rm cap}$; $\mathsf C_{\mathcal A}\delta(x)\le\frac12r^{\rm pl}_3$; $(\theta'_1+\theta''_1)\delta(x)\le\frac12c^+_1$; $6(LM)^3\xi(x)\le1$; $2(LM)^3\xi(x)^{1/2}\le\varepsilon_1L^3$; $[(1-\theta_m)A_1(\log x)^{\mathfrak p_0}]^{-1}\le1$; $2Q^{\rm pl}Q'L^{-2c^{\rm pl}\log(x/2)}\le1$; $Q'\bar K_{\rm w}e^{-c^{\rm pl}(\log(x/2))^2}\le\hat\rho v_d$; $\mathsf c_{\mathfrak h}v^+\delta(x)\le\frac12a'_Z$; $2B_3\varepsilon(x)\le1$; $2\varepsilon(x)\le\varepsilon_{\mathcal W}\alpha_0$; $K_H\varepsilon(x)\le\frac12\alpha^{\rm pl}_2$; $8K_\Lambda N_\star(1+\beta_{\rm s})\delta(x)/\alpha_0\le w_\star$.
\end{enumerate}
Each formula names only constants that stand before it, and none reads $K$, $m$, a level, a sequence of couplings, a weight or a field. The list so extended is an order table of admissible shape; choices of it exist; the restriction of a choice to the constants of $\mathcal T^{\rm S}[2\varrho]$ is a choice of that list, and its restriction to the constants of the list $\mathcal T^{\rm F+}$ of Theorem~\YMMref{D}{8} (the real choices in Proposition~\YMMref{69.1}{595}) is a choice of that list. Theorem~\ref{C-360-14} lowers $\gamma$ further where it says so.

\begin{theorem}[Three plane couplings]\label{C-360-13}
Let the constants be a choice of the extended list stated before this theorem, which adds to the list of Theorems~\YMMref{AJ}{673} and~\YMMref{AD}{607} four lower bounds on chosen constants, each by a function of earlier constants, further determined constants placed before $\gamma$ and further upper bounds on $\gamma$, and no new condition on $L$ itself; its choices exist, they restrict to choices of the lists of Theorems~\YMMref{D}{8}, \YMMref{AJ}{673} and~\YMMref{AD}{607}, and at each of them the standing conditions of the frame of Definition~\ref{def:running-source-history} hold, as do the locality for all pairs of regions and the gauge covariance at all decoupling parameters of the Gaussian structure among the construction results, the tower hypotheses of Section~\ref{sec:analytic-core-complements} and the inequality $\alpha_0+\alpha_1\le a^\sharp$ with the constant $a^\sharp$ in the radius conditions \YMMref{(867)}{410}. At every choice $L>4\cdot10^{19}$, as a consequence of the condition $L\ge2\mathsf C_{\rm ex}$ among the construction results and of the lower bound $\mathsf C_{\rm ex}\ge2\mathsf C_{\rm cs}$ of the extended list, with $\mathsf C_{\rm cs}\ge24C_{\mathrm E}^2e^{24}$.
\begin{enumerate}
\item[(a)] \emph{Lifted local terms.} For all $K,m$ and $1\le j<k<K$ with $2L^{m+K-k}\ge C_LM$, that is, for which the lattice $(\ell,K)$, $\ell=2L^m$, is allowed at the height $k$ (Section~\YMMref{50.4}{385}), $C_L\ge8\mathsf n_{\rm e}L$ being the room and $\mathsf n_{\rm e}\ge1$ the halo width of Section~\YMMref{50.4}{385}, so that $m+K-k\ge m'+2$, $\mathcal D_j$ and $\mathcal D_k$ are defined and a box of at most two cubes of the level $k$ along each axis does not wrap around the torus, which has at least $8L$ such cubes along each axis, every $Z\in\mathcal D_j$ short at $k$, every $P'\in\mathcal U_k(X_Z)$ and every term $E$ on $\mathcal U_j(Z)$,
\[
|E(\rho^{k\to j}_ZP')-E(\mathbf 1)|\le C_Q\sup|E|\,(d_j(Z)+2)^4L^{-4(k-j)},
\]
and the same with $(d_j(Z)+2)^2$ in the place of $(d_j(Z)+2)^4$ and with $72M^2\bar C_2^2(c^Q_3)^{-2}\le C_Q$ in the place of $C_Q$ when $L^{-(k-j)}\le\lambda_Q$ and $(d_j(Z)+4)L^{-(k-j)}\le c_s$, with the constants $\lambda_Q$, $c_s=c^Q_s$, $\bar C_2$ and $c^Q_3$ of item~8 of the list; $C_Q$ is a function of constants placed before $\gamma$, it lays no condition on any constant, and it depends on no $j,k,K,m,Z,P',E$. The same holds on every larger base in $\mathcal D_k$ that holds the box of $Z$, for cores at any positive radii, and on the plain cores for its local terms of level $j$. The power $4$ is attained: for every short $Z$ and every plaquette with its bonds in the boundary of $Z$ there are a term $E$ with $\sup|E|\le\frac12$ and a triple $P'$ with $E(\rho^{k\to j}_ZP')-E(\mathbf 1)\ge c_QL^{-4(k-j)}$, $c_Q>0$ a constant. For the abelian group ${\rm U}(1)^3$ the inequality fails.
\item[(b)] \emph{The lift.} Let $K,m$ and $1\le k<K$ be as in (a), $2L^{m+K-k}\ge C_LM$, and let $\kappa_{\rm lf}\ge4\log18432$ be a rate. For terms $g^{(j)}(Z,\cdot)$ on $\mathcal U_j(Z)$, $1\le j\le k$, $Z\in\mathcal D_j$, with $B_j=\sup_Ze^{\kappa_{\rm lf}d_j(Z)}\sup|g^{(j)}(Z,\cdot)|<\infty$, the lift
\[
{\rm Lf}[g](X;P')=\sum_{j\le k}\ \sum_{Z\in\mathcal D_j:\,X_Z=X}\bigl[g^{(j)}(Z;\rho^{k\to j}_ZP')-g^{(j)}(Z;\mathbf 1)\bigr],\qquad X\in\mathcal D_k,\ P'\in\mathcal U_k(X),
\]
($X_Z=Z^{\rm up}$, the union of the cubes of the level $k$ that hold a cube of $Z$, for the $Z$ that are not short, among them all $Z\in\mathcal D_k$, for which $Z^{\rm up}=Z$) is a term on $\mathcal U_k(X)$, linear in $g$, depending on $g$ at the $Z$ in $X$ alone, and
\[
|{\rm Lf}[g](X;P')|\le2C_E(140C_Q+1)\sum_{j\le k}L^{-(k-j)}B_j\,e^{-\kappa_{\rm lf}d_k(X)},\qquad C_E=4096,
\]
a bound that is uniform in the level.
\item[(c)] \emph{One step at all levels.} Let $K,m$ with $L^m\ge MR(x_K)$ and let $x_0>\dots>x_K\ge\gamma^{-2}$ be an admissible sequence of couplings in the frame of Definition~\ref{def:running-source-history}; by \eqref{eq:running-source-history-room} the frame has $C_L\le2LR(\gamma^{-2})\le2LR(x_K)$, so that $2L^{m+K-k}\ge2LMR(x_K)\ge C_LM$ at every $k<K$, and (a) and (b) apply at every level. Let $\sigma,\tau\ge0$ with $\sigma+2\tau\le c_{\rm w}$, $c_{\rm w}=w_6>0$ the constant of item~8, placed before $\gamma$, and let $\omega$ be a complex plaquette function whose plane weights $w_\kappa$, $\kappa\in\{12,13,23\}$, obey $\sup|w_\kappa|\le\sigma(x_K-\bar B)$ and $M{\rm Lip}_\infty w_\kappa\le\tau x_K$, with no tie between the planes.

Put $v_n=-L^{K-n}\omega/x_n$, $g_n=x_n^{-1/2}$, $y_n=\log(x_n/2)$, $\mathsf b_n=A_1(\log x_n)^{\mathfrak p_0}$ the threshold of the level $n$ (equation~\YMMref{(634)}{339}), and, for a weight $w$, $A_n(w;U)=A^{\eta_n}(w,U)=L^n\sum_pw(p)a_p(U)$, the weighted action of \eqref{eq:running-weighted-action}. For $1\le n\le K$, $X\in\mathcal D_n$ and $\mathsf a=(\mathsf a_1,\mathsf a_2)\in(0,\alpha_0]\times(0,\alpha_2]$, $\alpha_2=\alpha^{\rm pl}_2$ the third radius of item~8, the \emph{core of the size} $\mathsf a$, $\mathcal U_n(X;\mathsf a)$, is the set of triples $P$ in the same hereditary core of Definition~\YMMref{49.1}{374} at the radii $(\alpha_0,\alpha_1)$ with $|\partial\mathbf U_P(p)-\mathbf 1|<\mathsf a_1L^{-2n}$ at the plaquettes formed on $X$ and ${\rm hol}_n(P;X)<\mathsf a_2$; $\mathcal U_n(X)=\mathcal U_n(X;(\alpha_0,\alpha_2))$ is the core with the third radius, and $N_n(\mathsf a)=\sigma\mathsf a_2/\alpha_2+2\tau L^{-(K-n)}\mathsf a_1/\alpha_0$.

The following objects are formed by recursion on the step $k=0,1,\dots,K-1$, which integrates the fluctuation field of the level $k$ at the coupling $x_k$ over the fibre of a block field $W$ of the level $k+1$, at the triple $Q=P_W$. \emph{The classical action and the trees:} the classical action $\Psi^{(k)}(Q;\cdot)$ of the step is its Gaussian form, with the covariance $C_*(Q)$ of the step, together with its linear term and its higher potentials, the weight $v_k$ and the earlier trees being entered as potentials, so that the action of the step carries no explicit weight; $c^{(k)}_*$ is its stationary point; the tree ${\rm ta}_k(Z';\omega,\cdot)$, $Z'\in\mathcal D_{k+1}$, is the value $\Psi^{(k)}(Q;c^{(k)}_*)$, written as a sum of local terms and re-sorted into terms on the cores $\mathcal U_{k+1}(Z')$; and the sum of the trees is ${\rm Tr}_n(P)=\sum_{j=1}^n(x_{j-1}/x_n)\sum_{Z\in\mathcal D_j}{\rm ta}_{j-1}(Z;\omega,\rho_ZP)$, ${\rm Tr}_0=0$, so that $x_{k+1}{\rm Tr}_{k+1}(Q)=x_k[{\rm Tr}_k(\rho Q)+\Psi^{(k)}(Q;c^{(k)}_*)]$.

\emph{The non-classical array:} the non-classical term $\delta_{k+1}(Z';\omega,\cdot)$ consists of the potential $V^J$ of the Jacobians of the chart (the localised exponent $\mathsf J_W$ of the Jacobians in the weight ${\rm wt}_W$ of the chart integral below, one of the potentials of the step) and of the potential $V'$ of the arrays of the step, both taken at $c^{(k)}_*$, and of the ideal term $E^{\rm id}_k$ below at the centre $g_k^{-1}c^{(k)}_*$ less the same object at $\omega=0$, all re-sorted in the same way; ${\rm Lf}_k={\rm Lf}[g]$ is the lift of (b) at $g^{(j)}=\delta_j$; $S^c_n[\omega](U)=S[\mathfrak f_n](U)+\sum_{j\le n}\sum_{Z\in\mathcal D_j}[\delta_j(Z;\omega,p^{(j)}_Z(U))-\delta_j(Z;\omega,\mathbf 1)]$, where $S[\mathfrak f_n]$ is the action of the bare arrays of the never-large term at the zero source (Proposition~\ref{prop:running-source-density}, at $w=0$) and $p^{(j)}_Z(U)$ the plain pair of $U$ on $Z$, the configuration and its current of the level $j$ on ${\rm Bd}(Z)$; and $E_n[\omega]=E_{(n)}+\sum_{i\le n}\sum_{X\in\mathcal D_i}\delta_i(X;\omega,\mathbf 1)$, $E_{(n)}$ the number of that term.

\emph{The centres, the ideal term and the tiny terms:} the true centres are $\nu^{(n)}_j(V)=g_j^{-1}c^{(j)}_*$, read at the triple $\rho P_V$ of the level $j+1$, $0\le j<n$, and the centre of the step is $\mathsf m=\nu^{(k+1)}_k(W)$; for a system $V$ of potentials $V^o(Y;\cdot)$, $Y\in\mathcal D_k$, and a centre $\mathsf n$, $\Xi_{\mathsf n}[V](Q)=\int d\mathsf G_{\mathsf n}(B)\,\mathbf 1\{B\in{\rm Rg}_{x_k}\}\exp\sum_{Y,o}V^o(Y;B-\mathsf n)$, $\mathsf G_{\mathsf n}$ being the Gaussian measure of the covariance $C_*(Q)$ and the mean $\mathsf n$; $V_{\mathsf m}$ is the system of the three potentials of the step about the centre $\mathsf m$ (the potentials of the step, the potential of the arrays, and the remainder of the shift of the centre), without the carried tiny terms; $E^{\rm id}_k(X';Q,\mu)$, $X'\in\mathcal D_{k+1}$, is the cluster sum at the centre $0$ of the system $V_\mu$ on $X'$, so that $\Xi_0[V_\mu](Q)=\exp\sum_{X'}E^{\rm id}_k(X';Q|_{X'},\mu|_{X'})$; and ${\rm Rf}^{(k+1)}(X';Q,\nu)$ is the cluster sum at the centre $\nu_k$ of the system $V_{\nu_k}$ enlarged, for $k\ge1$, by the carried term ${\rm Rf}^{(k)}$ as one further potential, less the cluster sum at the centre $0$ of $V_{\nu_k}$ without it.

\emph{The classical factor of the step:} $c_k(W)=n_k(W)\exp\bigl[\mathcal E_k(W)-x_k\Psi^{(k)}(Q;c^{(k)}_*)+\sum_{Y\in\mathcal D_k}(V^J+V')(Y;Q|_Y;c^{(k)}_*|_Y)\bigr]$, where $\mathcal E_k(W)$ is the exponent of $\Phi_k[\omega]$ below at $V_W(0)$ ($\Phi_0[\omega]=\rho_0$), $V^J$ and $V'$ are the potential of the Jacobians of the chart and the potential of the arrays of the step, as in $\delta_{k+1}$, and $n_k(W)=\mathsf{pf}_k(x_k)\,\mathsf Z^{(k)}$ is the normalisation of the step, with the prefactor equation~\YMMref{(907)}{419} and the Gaussian integral in Proposition~\YMMref{50.2}{386} at the triple of $W$; so ${\rm wt}_W(B)\Phi_k[\omega](V_W(B))$ is $c_k(W)$ times the density of $\mathsf G_{\mathsf m}$ at $B$ times $\exp\sum_{Y,o}V^o_{\mathsf m}(Y;B-\mathsf m)$.

Then for $1\le n\le K$ the never-large term of $\rho_0=\exp[-A(x_0-L^K\omega,\cdot)]$ at the level $n$ has the density
\[
f_n(V)=\Phi_n[\omega](V)\exp\sum_{X\in\mathcal D_n}{\rm Rf}^{(n)}\bigl(X;P_V,\nu^{(n)}(V)\bigr)\ \text{ on }\mathcal V^\chi_n,\qquad f_n=0\ \text{ off it},
\]
$\mathcal V^\chi_n$ being the set of the configurations of the level $n$ that pass all the small-field tests (Definition~\ref{def:running-source-history}), and
\[
\Phi_n[\omega]=\exp\bigl[E_n[\omega]-x_nA_n(1+v_n;U_n(V))-x_n{\rm Tr}_n(P_V)+S^c_n[\omega](U_n(V))\bigr],\qquad v_n=-L^{K-n}\omega/x_n,
\]
where ${\rm Tr}_n$, the sum of the trees (the classical remainders) of the earlier steps, the non-classical array $S^c_n[\omega]$ and the number $E_n[\omega]$ are those above, and $P_V$ is the frozen triple of the background $U_n(V)$ of Proposition~\ref{prop:running-source-density} and Theorem~\YMMref{E}{10} in the setting of Section~\YMMref{52.7}{432}. Each step $0\le k<K$, at every $W\in\mathcal V^\chi_{k+1}$ and with $f_0=\rho_0$, is the identity
\[
\begin{multlined}
\int_{{\rm Rg}_{x_k}}{\rm wt}_W(B)f_k(V_W(B))\,dB=c_k(W)\,\Xi_0[V_{\mathsf m}](Q)\exp\sum_{X'\in\mathcal D_{k+1}}{\rm Rf}^{(k+1)}\bigl(X';Q, \\
\nu^{(k+1)}(W)\bigr),\qquad c_k(W)\,\Xi_0[V_{\mathsf m}](Q)=\Phi_{k+1}[\omega](W),
\end{multlined}
\]
whose left member is $f_{k+1}(W)$, the chart integral $\mathsf I[f_k](W)$ at $\mathsf t=x_k$ (Definition~\ref{def:sf-chart-data} and equation~\eqref{eq:sf-chart-integral}, \eqref{eq:sf-chart-integral}): ${\rm Rg}_{x_k}$ is the set of the scaled fields below the threshold $\mathfrak b(x_k)$ (Section~\ref{sec:sf-chart-interface}), $V_W(B)$ the chart of the fibre over $W$ (Lemma~\ref{lem:sf-chart-integral}), and ${\rm wt}_W(B)=\mathsf{pf}_k(x_k)\,z^{(k)}_{V_0}(\Gamma)\exp[-x_kG_W(B)+\mathsf J_W(B)]$, $V_0=V^{\rm bf}_k(W)$, the weight of the chart integral; $Q=P_W$, $\mathsf m=\nu^{(k+1)}_k(W)$ is the centre of the step, and $c_k(W)$, $\Xi_0$ and $V_{\mathsf m}$ are those above; so each step is classical, times ideal, times the exponential of a tiny term.

The ideal term is $\exp\sum_{X'}E^{\rm id}_k(X';\cdot)$ with $|E^{\rm id}_k|\le\frac14e^{-\kappa''d_{k+1}(X')}$, so it does not vanish. ${\rm Rf}^{(n)}(X;P,\nu)$ is holomorphic on $\mathcal U_n(X)$ times the polydisc of the centres, one $1$-form $\nu_j$ for every earlier level $j<n$, with $|{\rm Re}\,\nu_j(b)|<\vartheta_{j,n}\theta'_1\mathsf b_j$ and $|{\rm Im}\,\nu_j(b)|<\vartheta_{j,n}\theta''_1\mathsf b_j$, $\vartheta_{j,n}=\frac12(1+2^{-(n-j-1)})$, gauge invariant when the centres are carried along, and $|{\rm Rf}^{(n)}|\le\bar K_{\rm w}e^{-c^{\rm pl}y_n^2}e^{-\kappa d_n(X)}$, $y_n=\log(x_n/2)$, with the constants $\bar K_{\rm w}$ and $c^{\rm pl}$ of item~7; the true centres obey $|\nu^{(n)}_j(V)(b)|\le\frac{\theta_*}2L^{-2(n-j-1)}\mathsf b_j$, $\mathsf b_j=A_1(\log x_j)^{\mathfrak p_0}$ the threshold of the level $j$ (equation~\YMMref{(634)}{339}); the trees obey $e^{r_1d_n(Z)}|{\rm ta}_{n-1}(Z;\omega,P)|\le K_{\rm ta}N_n(\mathsf a)^2$ on the cores $\mathcal U_n(Z;\mathsf a)$ of the size $\mathsf a$, at every $\mathsf a\in(0,\alpha_0]\times(0,\alpha_2]$, with one $K_{\rm ta}$ and one rate $r_1$ at all levels; the non-classical terms obey $e^{r_1d_n(Z)}|\delta_n(Z;\omega,P)|\le129\min\{1,(\sigma+2\tau)/c_{\rm w}\}$ on $\mathcal U_n(Z)$; and the lift of the non-classical array, which is (b) at $g^{(j)}=\delta_j$, $B_j\le129$ and the rate $\kappa_{\rm lf}=r_1$, obeys $e^{r_1d_k(X)}|{\rm Lf}_k(X;\omega,P')|\le129K_{\rm Lf}L/(L-1)$, $K_{\rm Lf}=2C_E(140C_Q+1)$, as $\sum_{j\le k}L^{-(k-j)}\le L/(L-1)$.

The rate $r_1$ is admissible in (b): $r_1=\kappa''-8\log2$ with $\kappa''=(1+4\beta_{\rm s})\kappa+1+\lambda_{\rm s}$ and $\lambda_{\rm s}\ge1$ (Section~\YMMref{82}{665}, where $\lambda_{\rm s}=\log18432$ and the formula for $\kappa''$ is recalled), so that by item~3, $r_1\ge E_q+2-8\log2>\mathsf C_{\rm cs}>4\log18432$. Every number is a constant placed before $\gamma$ or an absolute constant. Everything is holomorphic in $\omega$ on the open set of the plaquette functions that lie, for some $\sigma,\tau\ge0$ with $\sigma+2\tau\le c_{\rm w}$, in the class with strict inequalities $|v_n|_\infty<\sigma$ and ${\rm sw}_n(v_n)<2\tau L^{-(K-n)}$ at every $0\le n\le K$, ${\rm sw}_n(v)$ being the largest difference $|v(p_{\mu\nu}(x'))-v(p_{\mu\nu}(x))|$ over the parallel plaquettes with $x'-x$ along $e_\mu$ or $e_\nu$ and of length at most one spacing of the level $n$.

At $\omega=0$ the trees, the $\delta_n$ and ${\rm Tr}_n$ vanish, the true centres are $\nu^{(n)}_j(V)=0$, and ${\rm Rf}^{(n)}(X;P,0)=0$ for every $P\in\mathcal U_n(X)$, at the zero centres; hence $f_n=\Phi_n[0]$ is the never-large term. The background is $U_n(V)$ from Proposition~\ref{prop:running-source-density}, and the centres of the Gaussian integrals are free arguments; the cores with the third radius coincide with the plain cores at the levels $n\le(\frac L2+29.39)/\log L$, so at the levels $n\le13$ for every odd $L\ge13$, and are proper subsets of them from some level on; at $k=0$, where the step is written with bond labels, its centre is at most $K_Z\sigma$ on the whole plain core, $K_Z$ a constant placed before $\gamma$.
\end{enumerate}
\end{theorem}

\begin{theorem}[Weighted backgrounds]\label{C-360-14}
In the setting of Theorem~\ref{C-360-13}, part (a) uses the regular-reference and geometric setting of Definition~\ref{hyp:weighted-critical-response}; its analytic response and Hessian estimates are proved in Proposition~\ref{prop:weighted-critical-response}. Parts (c) and (d) additionally assume the complex scalar gauge-transfer estimate of Definition~\ref{hyp:complex-scalar-transfer}. Evaluation of the algebraic identity in (e) on those moved backgrounds has the same assumptions; the identity and its bounds for already given families do not require it.
\begin{enumerate}
\item[(a)] \emph{Critical backgrounds.}
Use the real regular reference, natural-pair and separated-region
setting of Definition~\ref{hyp:weighted-critical-response}.
For the plane and symmetric corner weights $\mathbf h=(h,v)$,
Proposition~\ref{prop:weighted-critical-response} gives a unique
critical background in the slice
$RD^{\mathsf T}A=0$, $\mathcal Q(A)=b$, among fields with
$\sup_y\|A\|^{\beta_H}_y\le r_*$, where this norm includes
the scaled first-gradient H\"older seminorm defined there.
It is holomorphic in the row datum and the weights on
$\sup|b|<\rho_*$, $\|\mathbf h\|<z_*$, real at real data,
and obeys the localized response bounds
\eqref{eq:weighted-critical-row-response}--\eqref{eq:weighted-critical-weight-response}.
Its Hessian in the block datum satisfies
\eqref{eq:weighted-critical-hessian-floor} on $Y_0+iY_0$,
with the longer-sequence hypotheses in that proposition.
The unweighted comparison is at the same reference background.
The H\"older norm of the weights cannot be replaced by their
supremum in a uniform H\"older estimate for this response.
There is a separate finite-lattice assertion at $U_0=1$, when
all regions are the whole torus. For real plane weights
$v_{\mu\nu}=\epsilon k_{\mu\nu}$, $|k|\le1$, or for a single
real mixed tensor term
$\epsilon\varkappa\langle f_{\mu\nu},f_{\mu\lambda}\rangle$,
$|\varkappa|\le1$, with $0\le\epsilon<1$, the critical backgrounds
near the zero block datum form a real-analytic family and their
block Hessian at that datum obeys
$(1-\epsilon)Z_0\le Z_\epsilon\le(1+\epsilon)Z_0$.
This assertion uses each of these two perturbations separately.

A common constant $c\ne-1$ in the plane weights factors out:
the weights $(h,v+c)$ have the same critical fields as
$\mathbf h'=(h,v)/(1+c)$ and their multipliers and Hessians
are multiplied by $1+c$. The smallness conditions are those
on $\mathbf h'$. For real $c>-1$ the positive Hessian floor
is multiplied by $1+c$ as well.

After lowering $\gamma$ by the ceilings specified in the proof,
the following holds on every history. At each step $t\ge2$,
the background and row datum of the run satisfy the hypotheses
above, and for real weights with $\|\mathbf h\|\le z_N^+$ the
block Hessian is at least
$\frac12\gamma^{\rm pc}\|B\|^2$ on real $B\in Y_0$.
At the first step the real part of the second variation of the
action itself is at least that amount for
$B\in Y_0+iY_0$ and complex weights with
$\|h\|_\infty+\sup|v|\le z_N^1$.
The latter assertion uses the two-region natural-pair hypothesis
and the explicit background and block-field ceilings in the proof.

\item[(b)] \emph{Walk terms.} Under the natural-pair and separated-region hypotheses of Proposition~\ref{prop:weighted-system-L2}, there are numbers $r_\tau,r'_\tau>0$ and a power $M^\tau$ of $L$, depending only on $L$, such that all four weighted systems are invertible for $|v|_\infty<r_\tau$, radii at most $r'_\tau$, and block number at least $M^\tau$. Their inverses satisfy the mean-square divergence-source estimate \eqref{eq:weighted-system-L2-decay}, and their local walk terms satisfy \eqref{eq:weighted-L2-walk}, at a common exponential rate. The Hessian changes by the local plaquette form of Lemma~\ref{lem:weighted-local-hessian}. Grouping a walk by its visited pieces gives zero for a disconnected set $S$ and a bound with the factor $\epsilon_*^{\#S-1}$ for a connected set, after a further increase of $M^\tau$ for the prescribed $\epsilon_*$, as proved after Proposition~\ref{prop:weighted-system-L2}. The pieces and their adjacency are those specified there.
\item[(c)] \emph{The moved background on a label.} On the label of a history with large-field regions, in every frame of Theorems~\YMMref{AJ}{673} and~\YMMref{AD}{607} (every such frame satisfies the standing conditions under which these labels are formed), after lowering $\gamma$ by three further upper bounds of the shape of the list, $\sup_{x\ge\gamma^{-2}}\max\{w_{\rm G},w_{\rm lin}(L)\}(\log x)^{-r}\le1$ with $w_{\rm G}=4c_U^{-1}(7\log L+\log2)$, which gives $MR(\gamma^{-2})\ge\max\{w_{\rm G},w_{\rm lin}(L)\}$, $\sup_{x\ge\gamma^{-2}}C_m\varepsilon(x)\le\alpha_\star$ and $\sup_{x\ge\gamma^{-2}}2K'_5C_m\varepsilon(x)\le b_2$, where $C_m$ is fixed before $\gamma$ and depends on $L$ and the previously chosen construction constants, including $B_3$, and is the constant with which a label supplies its background, regular at the radius $C_m\varepsilon_k$ with the envelope $e(x)=C_m\varepsilon_{j(x)}$, $\varepsilon_j$ being the radius of regularity of the layer $j$; $\alpha_\star$ is a positive number of $L$ with $\alpha_\star\le\min\{\alpha_{\rm lin}(L),e_0(L)\}$, the numbers of Lemma~\ref{C-360-22}, so that the second ceiling gives the radius of regularity $\alpha\le\alpha_{\rm lin}(L)$ and the bound $\sup e\le e_0(L)$ on the envelope under which Lemma~\ref{C-360-22} applies; $K'_5$ is a positive number formed with $L$ through the constant $K_{\rm lin}$ of Lemma~\ref{C-360-22} and the constant $K_5$ below; and $b_2>0$ is the radius with which the substitution with three pairs of radii among the construction results (Proposition~\ref{prop:analytic-three-radii}) carries the cores at half their radii into the cores at three quarters of them, and for anisotropies ${\rm an}$ of the plaquettes with $t({\rm an})<t_0$, $t_0>0$ a number of $L$ with $t_0\le t_{\rm lin}(L)\le\frac14$, $t_{\rm lin}(L)$ the number of Lemma~\ref{C-360-22}, $t$ being the supremum plus the scaled H\"older-$\frac12$ quotients: the moved background $W=e^{i\eta_kA}U$, $\eta_k=L^{-k}$ the spacing of the level, exists, is critical for $\sum_p(1+{\rm an}(p))a_p$ on the surface of the prescribed block averages up to a complex gauge, and is holomorphic in ${\rm an}$; if $U=\tilde U$ at the bonds with an end point within the weighted distance $\mathsf R$ of a point $x_1$, then, with $e'(x)=e(x)\exp[-c_U(\mathsf R-d_B(x,x_1))_+]$,
\[
l(b)\,|\eta_kA[U](b)-\eta_kA[\tilde U](b)|\le K_U\,t({\rm an})\,e'(b_-)\qquad\text{at every bond }b,
\]
$b_-$ the initial point of $b$, with no gauge transformation between; the same bound holds, at every plaquette $p$, site $x$ and row $c$ of the block constraints, for the curvatures, $l(p)^2\eta_k^2|D_UA[U](p)-D_{\tilde U}A[\tilde U](p)|\le K_Ut({\rm an})e'(x_p)$, $x_p$ the base point of $p$, for the divergences, $l(x)^2\eta_k^2|D^*_UA[U](x)-D^*_{\tilde U}A[\tilde U](x)|\le K_Ut({\rm an})e'(x)$, and for the multipliers $\omega$ of the pair, $|\omega[U](c)-\omega[\tilde U](c)|\le K_Ut({\rm an})e'(c_-)$, $c_-$ the base point of $c$; if $\mathsf R\ge w_*$, then at every point $x_*$ of a layer $n\ge2$ with $d_B(x_*,x_1)\le\mathsf R-w_*$, in one gauge that is a member of gauge systems of both $U$ and $\tilde U$, the H\"older-$\frac12$ quotients of $\eta_k^2(D_UA[U]-D_{\tilde U}A[\tilde U])$, over the plaquettes of one orientation in the cube of the bare sites within $\frac14l_n$ bare spacings of $x_*$, and of $\eta_k(A[U]-A[\tilde U])$, over the bonds of one direction based within $\frac18l_n-2$ bare spacings of $x_*$, obey $(\frac14l_n)^{1/2}[\cdot]_{1/2}\le K_Ut({\rm an})e'(x_*)l_n^{-2}$, resp.\ $l_n^{1/2}[\cdot]_{1/2}\le K_Ut({\rm an})e'(x_*)l_n^{-1}$; and the multiplier $\omega_0$ of the plain background obeys $|\omega_0[U](c)-\omega_0[\tilde U](c)|\le K_Oe'(c_-)$.

Here $l(b),l(p),l(x)$ are the scales of the layers, $l_n=L^n$, $D_U$ is the covariant lattice exterior derivative and $D^*_U$ its adjoint, and $d_B$ is the weighted distance of Definition~\YMMref{12.3}{50}; $c_U$, $K_U$, $K_O$ depend on $L$ alone, $c_U\le1/(6L)$, $K_U=7K_{\rm lin}^2(1+2L^3)^2e^{c_Uw_*}$ with the constant $K_{\rm lin}$ of Lemma~\ref{C-360-22} and $w_*=w+21L$, $w=w(L)$ being the width of the cut-off used in the proof. The map $({\rm an},U)\mapsto A$ is jointly continuous, and the moved background, its triples and its currents are Borel functions of the configuration; every critical point on the fibre of norm $\le r_2(L)$ is carried to $W$ by a unique small gauge transformation that induces the identity at every top of the sequence; the full covariant gradient obeys $l(x)^2|\eta_k^2(\nabla_\kappa A_\nu)(x)|\le K_5\,t({\rm an})\,e(x)$; $W$ lies in the cores at three quarters of their radii; the construction commutes with the translations by vectors of the coarsest lattice; and $\sum_p|(1+{\rm an}(p))a_p(W)-a_p(U)|\le K|{\rm an}|_\infty\sum_pa_p(U)$.
\item[(d)] \emph{The centres on a label.} Let $0<\theta_3\le1$ be fixed before $\gamma$. In the setting of (c), let $\gamma$ be lowered further by upper bounds of the shape of the list formed with $\theta_3$, $A_0$, $A_1$, $\mathfrak p_0$, $r$, $\varepsilon_{\rm v}$ and numbers of $L$, among them $C_\chi(1+2\theta_3)\delta(x)<\frac18r_c$ for $x\ge\gamma^{-2}$, with two numbers $C_\chi,r_c$ of $L$; this keeps the box $|{\rm Re}\,B_1(b)|+|{\rm Im}\,B_1(b)|<(1+2\theta_3)\mathsf b_k$ of the complex block data of the step, of the unscaled radius $(1+2\theta_3)\delta(x_k)$, and with it the polydisc $|\mathsf n(b)|<\theta_3\mathsf b_k$ below, of the unscaled radius $\theta_3\delta(x_k)$, inside the ball, of a radius of $L$ alone, in which the system of the centre has exactly one solution.

Let the envelope of (c) be taken with $C'/L$ in the place of $C_m$, $C'=\max\{LC_m,2LC_\chi(1+2\theta_3)A_1/(\rho_{\rm u}A_0)\}$, $\rho_{\rm u}=\rho_{\rm u}(L)$ being the radius of that ball; $C'$ is the one number that is enlarged. Among the upper bounds on $\gamma$ above are the ceilings $\sup_{x\ge\gamma^{-2}}C'\varepsilon(x)\le c(L)$, $c(L)>0$ numbers of $L$, and the block-number and small-radius conditions of Proposition~\ref{prop:weighted-system-L2}, under which Corollary~\ref{cor:weighted-energy-L2} gives the global energy bound for the linear system; the nonlinear field-energy bound is supplied by Lemma~\ref{lem:weighted-field-energy}. Let the anisotropies obey moreover $t({\rm an})<t_\theta$, where $t_\theta=\min\{t_0,t_{\rm u}(L),c_{\rm ce}\theta_3A_1/A_0\}$, $t_{\rm u}(L)$ being the number of $L$ below which the moved backgrounds with a complex block datum are formed and $c_{\rm ce}>0$ a number formed with $L$ and $C'$; so $t_\theta\in(0,t_0]$ is formed with $\theta_3$ and with constants placed before $\gamma$.

At every step $0\le k<K$ and every regular base of the step (a regular background of the coarser level together with its block datum), the \emph{centre} of the step is defined by two conditions: it is a complex stationary point $\mathsf n$ of the weighted action on the fibre of the block field, and it is in the block axial gauge. There is exactly one $\mathsf n$ with $|\mathsf n(b)|<\theta_3\mathsf b_k$ at all bonds $b$ of the step that obeys the two conditions; it is $\mathsf n=g_k^{-1}a''$, $a''$ the $1$-form, in the axial gauge, of the finer block average of $W[{\rm an};b_0]$ read about the finer block average of $W[0;b_0]$. Here both backgrounds are formed over the same regular reference background $U_{\rm b}$ on the longer sequence, with the same complex block datum $b_0$ obtained from the base of the step; $|\mathsf n(b)|\le C_\theta t({\rm an})\mathsf b_k\le\frac12\theta_3\mathsf b_k$, $C_\theta$ depending on $L$, $\theta_3$, $A_1/A_0$ and $C_m$. At $b_0=0$, the base being $U_{\rm b}$ with its own block averages, $\sum_b|\mathsf n(b)|^2\le C^2|{\rm an}|_\infty^2A(\beta_t,U_{\rm b})$, $C$ a number of $L$, the sum running over all the bonds of the step for $1\le k<K$ and over its external bonds at $k=0$, $A(\beta_t,U_{\rm b})$ being the action of the reference background with the couplings of the step; and $\mathsf n$ is holomorphic in ${\rm an}$, real at real ${\rm an}$ and real $b_0$, and $0$ when ${\rm an}$ is constant.

The linearised system of the centre (the second variation of the constrained weighted action with its multipliers, the averaging constraint and the block axial gauge) is invertible, in $\ell^2$ and in a weighted maximum norm, with bounds depending on $L$ alone. The centres $\mathsf n^{(i)}$ of the earlier steps $0\le i<k$, formed by the same rule about the moved background of the label, vanish at ${\rm an}=0$, are holomorphic in ${\rm an}$ and obey $|g_i\mathsf n^{(i)}(b)|\le Ct({\rm an})\varepsilon_i$, that is $|\mathsf n^{(i)}(b)|\le Ct({\rm an})(A_0/A_1)\mathsf b_i\le\frac12\theta_3\mathsf b_i$, with one number $C$ of $L$ for all levels, $\varepsilon_i=\varepsilon(x_i)$ the radius of regularity of the level $i$ and $g_i=x_i^{-1/2}$; and their squares, summed over the stored bonds, which are the external bonds of these steps, obey $\sum_{i<k}\sum_b|\mathsf n^{(i)}(b)|^2\le C^2|{\rm an}|_\infty^2A(\beta_k,U)$, $A(\beta_k,U)$ being the action of the background with the couplings of the level $k$.
\item[(e)] \emph{The counterterms on a label.} On the label of a history with large-field regions the part of every coupling counterterm among the construction results that sees the anisotropy moves exactly into the local terms and the boundary terms (Lemma~\ref{C-360-23}): the anisotropy stands in the main weight and the local terms, and in the moved background whenever (c) is applied, and in no counterterm; the counterterms, the values at the trivial triple and the trivial values of the sums are unchanged; at a fixed triple, with the original families and counterterms held fixed, the changes in the re-filed local and boundary entries are linear in ${\rm an}$ with explicit Lipschitz constants. The following size bounds require a margin. Let $0<\theta<1$; let the families $E^{(j)}$ before the re-filing obey $\|E^{(j)}\|_{j,\kappa+1}\le(1-\theta)E_0$ at every level and every source; let $\|{\rm an}\|$ be the norm of the anisotropies in the format of the families, with $|{\rm an}|_\infty\le\sigma_3\|{\rm an}\|$, the anisotropies ranging over $\|{\rm an}\|<1$, and, if $B_\sharp>0$, let $\sigma_3\le\theta E_0e^{-7(\kappa+1)}/(\frac32B_\sharp\alpha_0^2)$.

Then the re-filed families $\hat E^{(j)}$ obey $\|\hat E^{(j)}\|_{j,\kappa+1}\le E_0$ and $\|\hat E^{(j)}-E^{(j)}\|_{j,\kappa+1}\le\theta E_0\|{\rm an}\|$; if $\|E^{(j)}-E^{(j)}_0\|_{j,\kappa+1}\le D$, $E^{(j)}_0$ the family at zero source, then $\|\hat E^{(j)}-E^{(j)}_0\|_{j,\kappa+1}\le D+\theta E_0\|{\rm an}\|$, and for the re-filed families this bound stands in the place of the quadratic bound among the construction results; in a frame of Theorems~\YMMref{AJ}{673} and~\YMMref{AD}{607} the added parts of the boundary entries are at most $\theta M^3E_0e^{-7}\|{\rm an}\|e^{-\kappa d_j(X)}$; the trivial values are unchanged.
\end{enumerate}
\end{theorem}

\begin{definition}[Weighted critical backgrounds]
\label{hyp:weighted-critical-response}
Fix $0<\beta_H<1$. Use a separated sequence through level
$1\le k<K$, with bare spacing $\eta=L^{-k}$, the interior and
outward determining rows, and no inward or edge rows. Retain the
natural-pair, block-number, separation and torus-size hypotheses
of the regular response construction in
Propositions~\YMMref{31.1}{252} and~\YMMref{34.2}{267}.
In particular the terminal torus has at least $6R_1\mathfrak M$
sites per direction. The common block number and width are chosen
large enough for the ordinary H\"older kernels and the positive
fluctuation form used below: $\mathfrak M\ge M^{\rm pc}$,
$R_1\mathfrak M\ge\mathfrak r^{\rm pr}$ and
$\delta_*R_1\mathfrak M\ge512\log L$.
Here $M^{\rm pc}$ is enlarged to make the cutoff defect in the
proof at most $1/8$, and the regularity ceiling $a^{\rm pc}$
is decreased to give the local inverse comparison there.

Let $U_0$ be real, critical for the unweighted action on the fibre
of its block averages. Write $\mathcal Q$ for the straightened
row map and $a_{\mathcal Q}$ for its fixed analytic potential
radius, chosen so that this map and its local Taylor bounds
hold on $\sup_y\|A\|_y<2a_{\mathcal Q}$.
Write $K_A=\mathcal Q'(A)$, $K_0=K_A|_{A=0}$, and
\[
 F(A)=\eta^{-1}\sum_pa_p(e^{i\eta A}U_0),\qquad
 F'(0)=K_0^{\mathsf T}\omega_0.
\]
The fine pairing has weight $\eta^3$, and the dimensionless row
pairing has weight $\xi_c^{-1}$. Suppose $U_0$ is regular at
radius $\alpha$: its plaquette defects are at most
$\alpha\eta^2\xi_y^{-2}$ and $\xi_y^3|F'(0)|\le\alpha$.
Require the pair $(U_0,F'(0))$ to belong to the natural class of
that sequence with radii $(\alpha,\alpha,\alpha)$ and real
factor $U_0$. This retains the local real gauges in that class.

The weights are $\mathbf h=(h,v)$. The plane weight $v$ is a
function of inner plaquettes. At each bare site $x$, the symmetric
matrix $h_{ww'}(x)$ is indexed by the $24$ oriented corners
$w=(\iota,\iota')$, $|\iota|\ne|\iota'|$, with
$\iota,\iota'\in\{\pm1,\pm2,\pm3\}$. If $f_w$ is the field
strength of the corresponding plaquette loop based at $x$, put
\[
 T_{\mathbf h}(A)=T_h(A)+T_v(A),\quad
 T_h(A)=\eta^3\sum_{x,w,w'}h_{ww'}(x)\langle f_w,f_{w'}\rangle,
 \quad T_v(A)=\eta^{-1}\sum_pv(p)a_p.
\]
All loops are contained in the ambient region. Set
$|h(x)|=\max_w\sum_{w'}2|h_{ww'}(x)|$.
The norm $\|\mathbf h\|$ is the supremum over tops $y$ of the
sum of the local suprema of $|h|$ and $|v|$ and their
$C^{\beta_H}$ seminorms multiplied by $\xi_y^{\beta_H}$.
The pairs in these seminorms have one base point in the cell of
$y$ and distance at most $\xi_y/4$; plane weights are compared
in the same plane.

For potentials use
\[
 \|A\|_y=\max\{\xi_y|A|_y,\xi_y^2|\nabla A|_y\},\qquad
 \|A\|^{\beta_H}_y
 =\max\{\|A\|_y,\xi_y^{2+\beta_H}[\nabla A]_{\beta_H,y}\}.
\]
The last seminorm uses the transports of $U_0$ along inner combs,
over the same distances. With $N_y$ denoting the appropriate
potential, row or weight norm, write
\[
 \|u\|_{\theta,y_0}=\sup_y e^{\theta d_{\mathbf B}(y,y_0)}N_y(u).
\]
For the Hessian assertion add the next region $\Omega_{k+1}$ and
require the same natural-pair bounds on this longer sequence.
Let $Y_0$ be the real top-layer fields about $\Omega_{k+1}$ whose
next linearized averages vanish and which vanish on the block
axial trees. Its norm and pairing are the top-row ones.
\end{definition}

\begin{proposition}[Weighted critical response]
\label{prop:weighted-critical-response}
In Definition~\ref{hyp:weighted-critical-response}, there are
positive radii $a_*,r_*,\rho_*,z_*$ and constants
$B_*,\Theta,C_T$, depending only on $L$ and $\beta_H$, with
$a_*\le\min\{a^{\rm pc},a_{\mathcal Q},L^{-2}/4\}$ and the
following property. If $\alpha\le a_*$, $\sup|b|<\rho_*$ and
$\|\mathbf h\|<z_*$, there is exactly one critical background
$e^{i\eta A}U_0$ with
\[
 RD^{\mathsf T}A=0,\qquad \mathcal Q(A)=b,
 \qquad \sup_y\|A\|^{\beta_H}_y\le r_*.
\]
Its pair $(A,\omega)$ is holomorphic in $(b,\mathbf h)$ and real
at real data, and equals $(0,\omega_0)$ at $(0,0)$.
It satisfies $\sup_y\|A\|^{\beta_H}_y<r_*$. For a row $c_0$
and a datum $e_{c_0}$ supported there,
\begin{align}
 \|d_bA\,e_{c_0}\|^{\beta_H}_y
 &\le2B_*e^{-\delta_*d_{\mathbf B}(y,y_{c_0})/4}|e|,
       \label{eq:weighted-critical-row-response}\\
 |d_b\omega\,e_{c_0}|_y
 &\le2\Theta B_*e^{-\delta_*d_{\mathbf B}(y,y_{c_0})/4}|e|.
       \notag
\end{align}
For every weight direction $\mathbf k$ and
$0\le\theta\le\delta_*/4$,
\begin{equation}\label{eq:weighted-critical-weight-response}
 \max\{\|d_{\mathbf h}A\,\mathbf k\|^{\beta_H}_{\theta,y_0},
       \Theta^{-1}\|d_{\mathbf h}\omega\,\mathbf k\|_{\theta,y_0}\}
 \le2C_T(\alpha+\sup_y\|A\|^{\beta_H}_y)
                        \|\mathbf k\|_{\theta,y_0}.
\end{equation}
Here $\delta_*$ is the fixed propagator rate, with the response
kernels taken at rate $3\delta_*/8$ before their weighted summation.
With the longer-sequence hypotheses, let $Z=d_b\omega$ be the
Hessian of the weighted critical value and put
$s=\max\{\sup|b|/\rho_*,\|\mathbf h\|/z_*\}$. After decreasing
the radii, there are $\gamma^{\rm pc},e_0,\kappa_0>0$ such that
\begin{equation}\label{eq:weighted-critical-hessian-floor}
 \operatorname{Re}\langle\bar B,ZB\rangle
 \ge(\gamma^{\rm pc}-e_0\alpha-\kappa_0s)\|B\|^2,
 \qquad B\in Y_0+iY_0.
\end{equation}
The reference for the perturbation in $s$ is the unweighted
Hessian at $b=0$ over this same $U_0$.
\end{proposition}

\begin{proof}
We keep the exact constrained response at the reference. Its
Hessian includes the multiplier correction. Put
$\Delta=F''(0)$, and define $W(A,\omega)$ by
$\langle X,W(A,\omega)Y\rangle
=\langle\omega,\mathcal Q''(A)[X,Y]\rangle$.
Thus $W_0=W(0,\omega_0)$. Set
\[
 G_1=(\Delta-W_0+DRD^{\mathsf T}+K_0^{\mathsf T}aK_0)^{-1},
 \quad H=G_1K_0^{\mathsf T}(K_0G_1K_0^{\mathsf T})^{-1},
 \quad \widetilde G=G_1-HK_0G_1.
\]
The physical and dimensionless row conversions are those in the
proof of Proposition~\YMMref{34.2}{267}; $a$ is the positive row
coefficient of that construction, including its boundary extension.
This extension cancels on $N(K_0)$. Elimination of the auxiliary
variables in Proposition~\YMMref{31.1}{252} gives
\[
 K_0H=1,\quad K_0\widetilde G=0,\quad
 \widetilde G(\Delta-W_0+DRD^{\mathsf T})=1-HK_0,
 \quad RD^{\mathsf T}H=0.
\]
The ordinary response kernels and their first differences decay
exponentially. Their H\"older bounds follow as in
Lemma~\ref{lem:ordinary-holder-kernels}, with the fixed exponents
$\beta_H$ and $(1+\beta_H)/2$ in that proof. The exponent gap in
its ordinary-input argument is positive. For a divergence source
of exponent $\beta_H$, use instead the same-exponent parametrix
in the proof of Lemma~\ref{C-360-22}, with that exponent in the
Green-kernel cancellation estimate of Lemma~\YMMref{19.4}{115}.
The local multiplier term $W_0$ has size $C\alpha$ and is one of
the lower-order terms in that parametrix. Thus all constants are
uniform in the region sequence and its depth.

Write
\[
 C(A)=\mathcal Q(A)-K_0A,\qquad
 F'(A)=F'(0)+\Delta A+\mathcal V(A).
\]
The critical equation with the slice term and its constraint is
equivalent to the following two equations:
\begin{align*}
A={}&H(b-C(A))+\widetilde G\bigl(
 C'(A)^{\mathsf T}\omega-F'(0)-\mathcal V(A)
                  -W_0A-T'_{\mathbf h}(A)\bigr),\\
 \omega={}&H^{\mathsf T}\bigl(
 F'(A)+T'_{\mathbf h}(A)+DRD^{\mathsf T}A
                         -C'(A)^{\mathsf T}\omega\bigr).
\end{align*}
Denote this map by $\Phi^{\mathbf h}$.
The exact identities above recover both equations from its fixed
point. Conversely a solution gives a fixed point by the same
elimination. The slice term vanishes at a real fixed point by
gauge invariance and the scalar projection identity; holomorphy
extends that identity to the complex parameter domain.

Here is the uniform smallness estimate behind the contraction.
The Taylor kernels of the row map are confined to the double
block beneath a row. The local action kernels are confined to a
plaquette star. In the scaled local norms their bounds are
\[
 \|C'(A)\|+\|C'(A)^{\mathsf T}\|\le Cr,
 \quad \|\mathcal V'(A)\|\le Cr,
 \quad \|W(A,\omega)-W_0\|
       \le C(|\omega-\omega_0|+r|\omega_0|)
\]
on $\sup\|A\|_y\le r$, while $|\omega_0|\le C\alpha$.
These are the local estimates used in
Proposition~\YMMref{33.1}{259}; their passage to exponentially
weighted norms costs a fixed summation constant. The first
component of $\partial_X\Phi^0$ is therefore small. In the
second component, the only term without a small factor is
$H^{\mathsf T}(D^{\mathsf T}D+DRD^{\mathsf T})\delta A$,
which has a fixed bound $K_8\|\delta A\|$.
Choose $\Theta\ge4K_8$ and measure pairs by
\[
 \|X\|_{\theta,y_0}
 =\max\{\|A\|^{\beta_H}_{\theta,y_0},
                   \Theta^{-1}\|\omega\|_{\theta,y_0}\}.
\]
First reducing $a_*,r_*$ makes the two remaining derivative
bounds at most $1/8$ each. The product inequality
$[fg]_{\beta_H}\le|f|[g]_{\beta_H}+[f]_{\beta_H}|g|$
gives the same bounds for the H\"older terms. The added weighted
action contributes at most $C\|\mathbf h\|$ to the derivative:
its leading term is the covariant divergence of the weight
times the field strength, and its other terms are the local
curvature commutators. The divergence estimate just described
uses the $C^{\beta_H}$ norm of that product. Reduce $z_*$ to
absorb this contribution. We obtain
\begin{equation}\label{eq:weighted-critical-contraction}
 \|\partial_X\Phi^{\mathbf h}\|_{\theta,y_0}\le\tfrac12,
 \qquad 0\le\theta\le\delta_*/4.
\end{equation}
This contraction proves uniqueness in the stated norm, including
its first-gradient H\"older term.

At $X_0=(0,\omega_0)$ the displacement of the map is bounded
by $C(\sup|b|+\alpha\|\mathbf h\|)$. Reducing $\rho_*,z_*$
puts this displacement within half the chosen ball. Iteration
gives the fixed point and its holomorphy, uniformly on smaller
parameter balls. Conjugation gives reality. If a critical point
in the slice has $\sup\|A\|^{\beta_H}_y\le r_*$, the row equation
$K_AH=1+C'(A)H$ has an inverse bounded by its Neumann series.
It gives the multiplier bound needed to place that point in the
same pair ball. Thus it is the fixed point. The actual block
field is obtained from the reduced datum by the accumulated
averaging gauge, as in Corollary~\YMMref{41.3}{317}.

Differentiate the fixed-point equation in a single row. It gives
\[
 (1-\partial_X\Phi^{\mathbf h})d_bX\,e_{c_0}
                  =(He_{c_0},0).
\]
The right side has norm at most $B_*|e|$ in the exponential norm
centred at $y_{c_0}$. The inverse on the left has norm at most
$2$, proving \eqref{eq:weighted-critical-row-response}.
In a weight direction the right side is formed from
$T'_{\mathbf k}(A)$. This first variation is linear in the field
strength and is bounded by
$C_T(\alpha+\sup\|A\|^{\beta_H}_y)\|\mathbf k\|_{\theta,y_0}$.
The same inverse gives \eqref{eq:weighted-critical-weight-response}.

We next record the unweighted lower bound used for the Hessian.
The longer-sequence hypothesis is needed here. Let $T^{\rm ax}$
be the fluctuation form in nonpivot axial coordinates. In the
local inversion proof of Theorem~\YMMref{L}{298}, the flat axial
form $T^0$ has the gap in equation~\YMMref{(593)}{300}.
The finite support of the pivot-coordinate map and the exponential
bound for the row inverse give an upper bound
$T^0\le a_{\rm pc}I$ by the row and column sum test. Hence
\[
 a_{\rm pc}^{-1}I\le (T^0)^{-1}\le\gamma_{\rm c}^{-1}I.
\]
The natural class supplies a real local gauge on each comparison
cube in which the potential and its first differences are small.
In that gauge, the local system differs from the flat one by the
local defect bounded in the proof of Theorem~\YMMref{L}{298}.
The resolvent expansion bounds the difference of the inverse
blocks by $C\alpha e^{-\delta d(b,b')}$.
Summing its rows and columns, and decreasing the regularity
radius, gives
\[
 \|C_\kappa-(T^0)^{-1}\|\le(2a_{\rm pc})^{-1}.
\]
Here $C_\kappa$ is the block of the local inverse from the
nonpivot-field source to the nonpivot field. Elimination of its
other components identifies it with the inverse of the local
reduced form. At real data it is real and symmetric. The last
bound therefore gives
\[
 (2a_{\rm pc})^{-1}I\le C_\kappa\le p_{\rm pc}I,\qquad
 p_{\rm pc}=\gamma_{\rm c}^{-1}+(2a_{\rm pc})^{-1}.
\]
Returning to the original local gauge is unitary and preserves
these inequalities. With the partition from that same local
inversion proof, put
\[
 P=\sum_\kappa h_\kappa C_\kappa h_\kappa,
 \qquad \sum_\kappa h_\kappa^2=1.
\]
The local inverse bounds give
$(2a_{\rm pc})^{-1}I\le P\le p_{\rm pc}I$.
To compare $PT^{\rm ax}$ with the identity, lift a nonpivot
field $B$ to the auxiliary solution $\Phi$ in the reduced system.
Its equation has only the last source component $T^{\rm ax}B$,
and the ordinary inverse estimate gives
$\|\Phi\|_\infty\le C_{\rm pc}\|B\|_\infty$.
On the enlarged comparison cube the local and global equations
agree. Multiply by a cutoff which equals one on the support of
$h_\kappa$ and apply the local inverse. The main term is $h_\kappa B$.
The error is the inverse applied to the cutoff commutator.
Every difference of the cutoff costs $\mathfrak M^{-1}$; the
enlarged cubes have bounded overlap. Multiplying by $h_\kappa$
and summing consequently gives
\[
 PT^{\rm ax}=1-E,\qquad
 \|E\|_\infty\le
 2N_GB_{\rm pc}C'_{\rm cut}C_{\rm cut}C_{\rm pc}
                         \mathfrak M^{-1}\le\tfrac18.
\]
Here $N_G$ is the overlap bound of the enlarged comparison
cubes, $B_{\rm pc}$ the local inverse bound,
$C_{\rm cut}$ the bound for the scaled cutoff differences,
and $C'_{\rm cut}$ the commutator bound per such difference.
The constant $C_{\rm pc}$ is the lifting bound for $\Phi$
above. All depend only on $L$.
The last inequality is imposed by the common block floor.
Although $PT^{\rm ax}$ need not be symmetric, it is similar to
the real symmetric matrix $P^{1/2}T^{\rm ax}P^{1/2}$.
Its eigenvalues lie in $[7/8,9/8]$. Consequently
$T^{\rm ax}\ge7/(8p_{\rm pc})$.
For $B\in Y_0$, the next-average constraint recovers its pivot
entries from the nonpivot entries: $B=\mathsf C\check B$.
Here $\mathsf C$ solves the linear next-average equations
for the pivot bonds in the axial gauge. Its entries are bounded
by $C_K$, and each bond shares a coordinate block with at most
$n_{\rm pc}=36L^3$ positively oriented bonds.
The finite block-coordinate estimate gives
$\|B\|^2\le2(C_Kn_{\rm pc})^2\|\check B\|^2$.
Hence the unweighted fluctuation form has a positive floor
$\gamma^{\rm pc}$ on $Y_0$. The difference between that
Lagrangian form and the Hessian of the critical value is the
second derivative of the row map paired with the reference
current. Its local kernel bound gives a loss at most
$e_0\alpha\|B\|^2$. Thus at $(b,\mathbf h)=(0,0)$,
\[
 \langle B,Z(0,0;U_0)B\rangle
       \ge(\gamma^{\rm pc}-e_0\alpha)\|B\|^2.
\]

The derivative of the multiplier is the block Hessian, because
the critical equation cancels the field variation in the first
derivative of the critical value. The analytic response and its
localized derivative estimates above bound the rows and columns
of this Hessian by a summable exponential kernel on the top
layer. Cauchy's estimate on the common parameter balls, followed
by that same row and column sum, gives
\[
 \|Z(b,\mathbf h;U_0)-Z(0,0;U_0)\|_{Y_0^{\mathbb C}
                    \to(Y_0^{\mathbb C})^*}\le\kappa_0s.
\]
Reducing the stated radii once leaves the larger analytic balls
needed for Cauchy's estimate. Taking the real part proves
\eqref{eq:weighted-critical-hessian-floor}. The unweighted
comparison throughout is at the same reference $U_0$.
\end{proof}



\subsection{Construction estimates}

We use the following results among the construction results. Theorem~\YMMref{D}{8}, and Theorems~\YMMref{AJ}{673} and~\YMMref{AD}{607}. The weighted progression and the never-large term with a source, with the statements on one term and on families: Section~\YMMref{51.9}{411}, Proposition~\YMMref{51.15}{410} and Theorem~\YMMref{V}{416}; and the sections on sources, equation~\eqref{eq:running-source-one-step} and Proposition~\ref{prop:running-source-density}. The rarity of a first creation of a large field on one torus, Proposition~\ref{prop:ward-fixed-volume-rarity}, and the rarity of the large-field region uniformly in the volume, Proposition~\ref{prop:local-large-field-rarity}. The law of the field and of the history, Proposition~\YMMref{44.1}{342}, with the measures in equation~\YMMref{(647)}{342}.

The count of polymers, Lemma~\ref{lem:analytic-domain-geometry}, whose extension to any number of sources is Lemma~\ref{C-360-24}. The frame, Definition~\ref{def:running-source-history}; its background and the history with no large-field region, Proposition~\ref{prop:running-source-density}; and the local step with a source, Proposition~\ref{prop:source-entry}. The hereditary cores, Definition~\YMMref{49.1}{374} and Section~\ref{sec:analytic-core-complements}; the H\"older bound for real backgrounds, Theorem~\textup{\ref{I-56-65}}; and the ordinary H\"older response, Proposition~\ref{prop:ordinary-holder-response}. The boards of the lattices, Section~\ref{sec:periodic-chessboard}; reflection positivity at the site and bond mirrors, which are the walls of the boards, Proposition~\YMMref{91.4}{700}; the chessboard estimate, Proposition~\YMMref{86.3}{681} and, for the torus measures, Corollary~\ref{IV-18-18}; the domination of periodic Gibbs states by the torus measures through their conditional laws, Proposition~\YMMref{91.9}{705}; and the bound $0\le p_{\ell'}-p_\ell\le G_g$ of the setting on the pressures formed with the plain partition functions of Section~\YMMref{86}{680}, Corollary~\ref{IV-18-18} with Proposition~\YMMref{89.2}{692} and Lemma~\ref{lem:canonical-source-system}.

The bound on connected functions of $\mathcal O$ on a torus at separated supports, $|S^T_{n;K,\ell,\beta}(f)|\le C_n(\mathsf d)e^{nc_V\ell^3}\prod_i\|f_i\|_\sharp$ for real Lipschitz $f_i$ two of whose supports are $\ge\mathsf d$ apart, as proved in Lemma~\ref{lem:torus-cumulant-bound}, and its
per-ball form, Proposition~\YMMref{89.4}{693}, for supports in
sup-balls of radius $R$ and $\ell\ge16n(R+1)$. The two-point law on a torus of small side, Theorem~\YMMref{AU}{847} with Proposition~\YMMref{126.4}{846}, and the kernel bounds $|C^{\rm free}_K(p,q)|\le(1+3C_{\rm lat})a^3|u|^{-3}$ and $|\epsilon_p\epsilon_qC^{\rm free}_K(p,q)-a^3\mathbb K_3(u)_{\lambda(p)\lambda(q)}|\le C_{\rm lat}a^3(a^2|u|^{-5}+\ell^{-3})$, $u$ the least lift of $x(p)-x(q)$ and $\epsilon_p=\pm1$ the orientations, of equations~\YMMref{(2245)--(2246)}{824}.

The comparison of two lattices at equal label, Theorem~\YMMref{AR}{810}; the laws of the block fields, Definition~\textup{\ref{IV-62-62}} and Corollary~\textup{\ref{IV-62-66}}; and the limits $S^{(\ell,\mathsf b)}_n$ of $S^T_{n;K,\ell,L^K\mathsf b}$ at a fixed label, with the rate $K^pL^{-K/10}$, Theorem~\YMMref{AS}{813}, with the all-side extension in Corollary~\YMMref{119.2}{816}. Reflection positivity of the periodic Gibbs states at all the families of mirrors, Proposition~\textup{\YMMref{103.4}{758}}, and the reconstruction data of the limits of Gibbs sequences, Theorem~\YMMref{AP}{780}. The large-field part: the histories and their labels defined in Section~\YMMref{44}{339}, with the finite exploration in equations~\YMMref{(643)--(644)}{341}; the separated layers, physical row spaces and unweighted constrained systems of Sections~\YMMref{23}{176} and~\YMMref{27}{203}; their weighted mean-square extension in Proposition~\ref{prop:weighted-system-L2} and Corollary~\ref{cor:weighted-energy-L2}; and the layer-weighted distance of Lemma~\YMMref{12.3}{50}. The same-exponent H\"older estimate for the moved background is proved in Lemma~\ref{C-360-22}, with its normed spaces and envelope class.

The lattice calculus of Section~\ref{sec:ward-general-calculus}, \ref{sec:ward-general-calculus} and equations~\YMMref{(2104)--(2105)}{781} and~\YMMref{112}{781}, the integration by parts \ref{lem:ward-general-calculus} of Lemma~\ref{lem:ward-general-calculus}, and the rotational Ward identity at finite cutoff, Proposition~\YMMref{112.1}{782}. For the step with plane couplings and the backgrounds: the lattices allowed at a height, with the room and the width of the halos, Section~\YMMref{50.4}{385}; the scaled fields below the threshold, Section~\ref{sec:sf-chart-interface}, the chart of the fibre, Lemma~\ref{lem:sf-chart-integral}, and the chart integral, Definition~\ref{def:sf-chart-data} and equation~\eqref{eq:sf-chart-integral} with \eqref{eq:sf-chart-integral}, its prefactor equation~\YMMref{(907)}{419} and the Gaussian integral in Proposition~\YMMref{50.2}{386}; the weighted action \eqref{eq:running-weighted-action} and the density \eqref{eq:running-source-density} at zero source; the common real and sourced choices in Proposition~\YMMref{69.1}{595} and Section~\YMMref{82}{665}; and the substitution with three pairs of radii, Proposition~\ref{prop:analytic-three-radii}.

For the arbitrary-order short-distance comparison we assume
$\mathrm{(Q_n)}$ of Definition~\ref{hyp:quantitative-torus}.
Its uniform interacting small-field remainder is part of that
additional input. The free Gaussian contribution is computed
by expanding $-\frac32\log\det(1-\mathcal S W_\lambda)$:
the mixed coefficients are the cyclic contractions of the
covariance $\mathcal S$. The explicit first-creation weights
give the large-field term in $\omega_n$. These two calculations
do not supply the remaining arbitrary-order interacting error.
From the other sections we use: for $0<g\le\gamma^{\rm sd}_W$, the bound on centred products of any number of tensor entries at separated supports in every periodic Gibbs state, uniformly in the cutoff, the side, the state and the position, Theorem~\ref{C-355-9}(a), and the bound $\|\mathcal T_K(H)-c\|_{\rm ch}\le2\mathsf C_TY_i^{\mathsf q_T}a_i^{-11/2}\mathsf b^{-1}|H|_\infty$, $Y_i=2[\mathsf c_T+\log\frac1{a_i}]$, on the chessboard seminorm of a centred tensor entry in a cell of side $a_i$, \eqref{C-355-8} of Theorem~\ref{C-355-7}(a), with $\mathsf C_T,\mathsf q_T,\mathsf c_T$ the constants $C_E,q_E,c_E$ of that bound (Section~\ref{C-355-1}); and, for $0<g\le\gamma^{\rm sd}_W$, the invariance $\sigma(RX)=\sigma(X)$, $R\in{\rm O}(3)$, of tensor limits at tuples with one non-scalar entry among scalar ones (Section~\ref{cov-sec}, Theorem~\ref{cov-thm-rider}, which gives Theorem~\ref{cov-thm}(c) at these tuples).

\subsection{Auxiliary estimates}

\begin{lemma}[Several sources]\label{C-360-24}
Let $0<g\le\gamma$ with $g^{-2}\ge x_{\rm J}$, which holds for every $0<g\le\gamma^{\rm ref}_W$; let $(K,\beta)\in\mathfrak D(g)$ with $K\ge3$, $x_0=\beta>x_1>\dots>x_K$ its run, $g^{-2}<x_K\le X_g$, and $\ell=2L^m$ with $MR(x_K)\le L^m$, which holds on every side $\ell\in\mathfrak L(g)$; let $s\ge1$ be a depth and $k=\max\{K-s,2\}$ the working level, so that $2\le k\le K-1$; and let $(F_a,o_a,w_a)$, $a=1,\dots,p$, be dominated local sources with the comparison weights $w_a$. A \emph{dominated local source} $(F,o,w)$ consists of a real function $F$ of the configurations of the bare bonds that is a gauge invariant polynomial in the entries of the bond variables; a local representation $o$ of $F$, that is, families of local terms at the heights $1\le j\le k$ whose sums, less their values at the trivial configuration, are $F$, and a family at the height $0$ that sums to the fluctuation variation $F(e^{iB'}\mathbf U)-F(\mathbf U)$ over the background $\mathbf U$ of the level-$1$ block field; and a real weight $w\ge0$ of the plaquettes with $|F|\le\mathcal O(w)$ on all configurations, with sizes of $o$ at most those of a source of the first size $|w|_\infty$, $m_0(o)\le12\mathsf c^\flat M^3\alpha_0^2L^K|w|_\infty$ and $m_j(o)\le12M^3\alpha_0^2L^{K-j}|w|_\infty$ for $1\le j\le k$, $m_0(o)$ being the size of the family at the height $0$, $m_j(o)$ twice the norm $\|\cdot\|_{j,\kappa}$ of the family at the height $j$ and $\mathsf c^\flat=81C_P^2\mathsf C_{\mathcal A}^2/K_\Lambda$, and such that the corners of the plaquettes of ${\rm supp}\,w$ carry $o$: its members vanish on the regions that hold none of these corners. Zero sources are omitted before the parameter radii are formed.
\begin{enumerate}
\item[(a)] Let $\Lambda_0=\{\lambda\in\mathbb C^p:\sum_a|\lambda_a|m_{0,a}<N^{[0]}_\star\}$ and, for $1\le j\le k$, $\Lambda_j=\{\lambda\in\Lambda_{j-1}:\sum_a|\lambda_a|n_{j,a}<N^{[j]}_\star-(\sum_{n\le j}\omega_n)(E_0+\frac32\alpha_0^2B_\sharp)\}$, where $m_{j,a}=m_j(o_a)$, $n_{j,a}=m_{j,a}+\sum_{i=1}^j\omega_{j+1-i}d_{i,a}$, $d_{1,a}=\theta(x_0)m_{0,a}$, $d_{j+1,a}=\theta(x_j)n_{j,a}$, $\omega_n=(n+1)L^{-(n-1)}$, $\theta$ is the size of a source of (b), and $N^{[j]}_\star=w_\star\alpha_0/(4K_\Lambda(1+\beta_{\rm s})\delta(x_j))$ is the bound on the size of the bare array that the step of the level $j$ carries (equations~\eqref{eq:weighted-progression-numbers} and~\eqref{eq:weighted-progression-domain}). These are open convex neighbourhoods of $0$, decreasing in $j$; their thresholds read no source and not $p$, so the sources share the room of each step, and a polydisc $\{|\lambda_a|<\rho_a\}$ lies in $\Lambda_{t-1}$ when $\sum_a\rho_am_{0,a}\le N^{[0]}_\star$ and $\sum_a\rho_an_{j,a}$ is at most the right member of the condition of $\Lambda_j$ for $1\le j<t$; in particular $\Lambda_k=\bigcap_{0\le j\le k}\Lambda_j$ holds the polydisc of the radii $\rho_a=p^{-1}\min_{0\le j\le k}N_j/n_{j,a}$, $N_j$ being the right member of the condition of $\Lambda_j$ and $n_{0,a}=m_{0,a}$: these radii are $p^{-1}$ times numbers that depend on the sizes of the sources alone.

For $1\le t\le k+1$ and $\lambda\in\Lambda_{t-1}$, $\tau_{h_t}[\rho_0e^{\sum_a\lambda_aF_a}]=e^{D^\lambda_t(U_t(\cdot))}\mathsf m_t$ on the small block fields of the level $t$, with $D^\lambda_t$ holomorphic in $\lambda$ and formed from the members $E^{(i)}_\lambda(X)$ of the weighted progression; $\partial_BE^{(i)}_\lambda(X)=0$ unless $X$ meets the carrier of every $F_a$, $a\in B$, and the mixed derivatives at $\lambda=0$ of the level $i$ carry the factor $\exp(-[L^{t-i}\mathsf D_t/M-2]_+)$, $\mathsf D_t$ the least sup distance of the carriers in level-$t$ units. The centred moments at the never-large history are $\Phi[\prod_{a\in B}(F_a-c_a)]=\sum_\pi\prod_{\beta\in\pi}\kappa_\beta$, $\pi$ over the partitions of $B$ into blocks $\beta$, with $\kappa_\beta=\partial_\beta D^\lambda_t|_{\lambda=0}$ for $\#\beta\ge2$ and $\kappa_{\{a\}}=\partial_aD^\lambda_t|_{\lambda=0}-c_a$ at the one-element blocks.

Here $\rho_0=e^{-\beta A}$ is the bare density, $h_t$ the never-large history of the level $t$ and $\tau_{h_t}[\cdot]$ its term, $\mathsf m_t=\tau_{h_t}[\rho_0]$ a finite positive measure on the block fields of the level $t$, $U_t$ the background, $c_a$ the constant of $F_a$ at the level $t$, and $\Phi[Y]$ the bounded density of $\tau_{h_t}[Y\rho_0]$ with respect to $\mathsf m_t$.
\item[(b)] Let $p\ge2$, $0<g\le\gamma^{\rm ref}_W$, $0<\mathsf d<1$ with $s=s(\mathsf d)\le K$ and $2\mathsf T_p(3)\theta(L^{s-2}g^{-2})\le1$, and let the weights $w_a\ne0$ be carried by balls $B(x_a,\frac78\mathsf d)$ at mutual distances $\ge\frac54\mathsf d$. Then $|\langle F_1;\dots;F_p\rangle_{K,\ell,\beta}|\le k^{\rm loc}_p(g,s;\ell)\,\mathsf b^{-p}\prod_a|w_a|_\infty$, with $k^{\rm loc}_p$ depending on $g,p,\ell,s$ alone. Here $\mathsf T_p(u)=\sum_{n\ge1}(n+1)L^{-(n-1)}\bigl[(\frac43)^{1/2}L^{n/2}(1+n)^{\mathfrak p_0}\bigr]^pe^{-4\beta_s\kappa[L^{n-1}u-2]_+}$ is the series of the mixed derivatives of the order $p$ and $\theta(x)=K_{\mathbf M}(1+\beta_s)A_1(\log x)^{\mathfrak p_0}x^{-1/2}/\alpha_0$ the size of a source at the coupling $x$, with the constants $\beta_s,K_{\mathbf M}$.
\item[(c)] In (b), let moreover $2\max\{\tilde{\mathsf T}_p(3),(\frac43)^{1/2}L\,\mathsf S_{\rm nat}\}\,\theta(L^{s-2}g^{-2})\le1$, where $\tilde{\mathsf T}_p$ is the series $\mathsf T_p$ with the further factor $L^{n/4}$ in its $n$-th term and $\mathsf S_{\rm nat}=\sum_{n\ge1}(n+1)^{1+\mathfrak p_0}L^{-n/4}$; this condition implies that of (b), and $\tilde s(p)$ denotes the least integer $s\ge3$ with it. If the sizes of the local representation of $F_1$ at the heights $1\le j\le k$, and its values on the backgrounds of the levels $1\le t\le k$, are smaller than those of a source of the weight $w_1$ by the factors $L^{-j/4}$, resp.\ $L^{-t/4}$, that is, $m_j(o_1)\le12M^3\alpha_0^2L^{K-j}L^{-j/4}|w_1|_\infty$ and $|F_1(U_t(W))|\le L^{-t/4}\,2B_3^2\varepsilon(x_t)^2L^{K-4t}\sum_pw_1(p)$ at the backgrounds of the configurations $W\in\mathcal V^\chi_t$, then
\[
|\langle F_1;\dots;F_p\rangle_{K,\ell,\beta}|\le L^{-(k-1)/4}\,\tilde k_p(g,s;\ell)\,\mathsf b^{-p}\prod_a|w_a|_\infty,
\]
with a number $\tilde k_p(g,s;\ell)\ge k^{\rm loc}_p(g,s;\ell)$ that depends on $g,p,\ell,s$ alone.
\item[(d)] Let $c=\sum_i\kappa_im_i$ be a block form of Lemma~\ref{C-360-8}(b): each monomial $m$ is a product of $n_2(m)$ factors $\langle\mathbf w,\mathbf w'\rangle$, $n_3(m)$ factors $\det(\mathbf w,\mathbf w',\mathbf w'')$ and $n_a(m)$ factors $1-\frac12{\rm tr}\,V$, where $V$ runs over the holonomies of the lassos of the block of $27$ sites centred at a site $x$, the holonomies of its plaquettes carried to $x$ along paths of at most three steps in the block, and $\mathbf w=\mathbf w(V)$; the order of $m$ is $2n_2(m)+3n_3(m)+2n_a(m)$, and $\|c\|=\sum_i|\kappa_i|2^{n_a(m_i)}$ is the norm of $c$ of that lemma. For a real function $k$ of the sites, the source $\mathrm{Fm}^c(k)=a^{-1}\sum_xk(x)c_x(U)$ of Lemma~\ref{C-360-8}(b) is a dominated local source with the comparison weight $\varpi\mathbf 1_{{\rm Pl}(k)}$ for every $\varpi\ge76\|c\||k|_\infty$, ${\rm Pl}(k)$ being the set of the plaquettes of the blocks of the sites with $k\ne0$, and $|\mathrm{Fm}^c(k)|\le\mathcal O(2\|c\||k|_\infty\mathbf 1_{{\rm Pl}(k)})$ on all configurations; if all its monomials have an order $\ge3$, the hypothesis of (c) holds for it with the comparison weight $\Gamma^\circ_c\|c\||k|_\infty\mathbf 1_{{\rm Pl}(k)}$, $\Gamma^\circ_c=\max\{76,\ 34\max\{1,5.04\alpha_0\}^{{\rm Ord}(c)-2},\ 2\max\{1,2B_3\}^{{\rm Ord}(c)-2}\}$, ${\rm Ord}(c)$ the largest order of its monomials, even with $L^{-2j}$ and $L^{-2t}$ in the places of $L^{-j/4}$ and $L^{-t/4}$.
\item[(e)] Let $\ell\in\mathfrak L(g)$, let $(F_a,o_a,w_a)$, $a=1,2$, be two dominated local sources whose carriers are at the sup distance $\mathsf d'$, and let $2\le t\le k$, $k$ the working level, $s'=K-t$ (so that $1\le s'\le K-2$, as $k\le K-1$), with $L^{s'}\mathsf d'\ge3M$, so that $\mathsf D_t\ge3M$, and $\mathsf h_2(s')\le1$, where $\mathsf h_2(s')=\frac43\mathsf k_2L\,\theta(\frac34L^{s'+1}g^{-2})\,\mathsf T^\circ_2(L^{s'}\mathsf d'/M)$ with the series $\mathsf T^\circ_2(u)=\sum_{n\ge1}(n+1)^{2\mathfrak p_0+1}e^{-4\beta_s\kappa[L^{n-1}u-2]_+}$ and the number $\mathsf k_2=\max\{\frac12(\mathsf c^\flat)^2,5\}$. Then, with $y_a\le4\ell^3L^{4(K-t)}|w_a|_\infty$: the real part of the density $D_{t,2}$ of $F_2$ at the never-large history differs from the real part of its constant $c_2$ by at most $C_1y_2$, and the mixed part $D_{t,12}$ of the density of $F_1F_2$ is at most $C_2e^{3-\mathsf D_t/M}y_1y_2$.
\end{enumerate}
\end{lemma}

\begin{proof}[Proof sketch]
$MR(x_K)\le L^m$ on the sides of $\mathfrak L(g)$: \[
L^{m_{\mathfrak L}(g)}\ge LM\,R(X_g+\bar B)
\]
by equation~\YMMref{(1528)}{615}, and $x_K\le X_g+\bar B$. (a) is the derivative calculus for the never-large term with one source (Proposition~\ref{prop:source-entry} and equation~\eqref{eq:running-source-one-step}), carried out in $p$ parameters. Each source enters each step as one more local term, and the $p$ sources share the room of each step, the bound of the step on the size of what it carries: with $m_{0,a}$ the size of $F_a$ at the height $0$ and $n_{j,a}=m_{j,a}+\sum_{i=1}^j\omega_{j+1-i}d_{i,a}$ the size at the level $j$ of $F_a$ together with all that the steps below the level $j$ have made of it, as in the statement, the parameter $\lambda$ ranges at the level $t$, $1\le t\le k+1$, over the set $\Lambda_{t-1}=\bigcap_{0\le j<t}\Lambda_j$, the sets $\Lambda_j$ being those of the statement,
\[
\Lambda_0=\Bigl\{\lambda:\sum_a|\lambda_a|\,m_{0,a}<N_0\Bigr\},\qquad\Lambda_j=\Bigl\{\lambda\in\Lambda_{j-1}:\sum_a|\lambda_a|\,n_{j,a}<N_j\Bigr\}\quad(1\le j\le k),
\]
with $N_0=N^{[0]}_\star$ and $N_j=N^{[j]}_\star-(\sum_{n\le j}\omega_n)(E_0+\frac32\alpha_0^2B_\sharp)$ for $1\le j\le k$, $\omega_n=(n+1)L^{-(n-1)}$ as in the statement: $N^{[j]}_\star$ is the bound on the size of what the step of the level $j$ carries, and for $1\le j\le k$ the array of the level $j$ without the sources has a size at most $(\sum_{n\le j}\omega_n)(E_0+\frac32\alpha_0^2B_\sharp)\le\frac12N^{[j]}_\star$ (the weighted-lag array bound in the proof of Proposition~\ref{prop:weighted-finite-system}, together with $8K_\Lambda N_\star s_j\le w_\star$ from item~9 before Theorem~\ref{C-360-13}), to which the sources add at most $\sum_a|\lambda_a|\,n_{j,a}$; so $N_j\ge\frac12N^{[j]}_\star$ is the room that the step of the level $j$ leaves to the sources, a number that reads no $p$. The sets $\Lambda_j$ are open and convex and decrease with $j$, and $\Lambda_{t-1}$ holds the polydisc of the radii $\min_{0\le j<t}N_j/(p\,n_{j,a})$, $n_{0,a}=m_{0,a}$, since on it $\sum_a|\lambda_a|\,n_{j,a}<N_j$ for $0\le j<t$; so $\Lambda_k$, which lies in every $\Lambda_{t-1}$, holds the polydisc of the radii $\min_{0\le j\le k}N_j/(p\,n_{j,a})$.

These radii are $p^{-1}$ times numbers that depend on the sizes of the sources alone, while the thresholds $N_j$ read no $p$. On $\Lambda_{t-1}$ the members of the progression up to the level $t$ are holomorphic in $\lambda$; the derivative in $\lambda_a$ of a member vanishes unless its polymer meets the carrier of $F_a$, and the count of polymers in Lemma~\ref{lem:analytic-domain-geometry}, applied at every level $1\le j\le k+1$ to the polymers that meet a given set of sites, gives the decay in the distance of the carriers: a polymer of the level $i$ that meets the carriers of all the $F_a$, $a\in B$, has a linear size at least $L^{t-i}\mathsf D_t/M-2$, whence the factor $\exp(-[L^{t-i}\mathsf D_t/M-2]_+)$ of the statement; the formula for the centred moments is the moment--cumulant formula for the quotient of the never-large term with the sources by the term without them, in which the centring moves the first cumulant alone: $\kappa_{\{a\}}=\partial_aD^\lambda_t|_{\lambda=0}-c_a$, and $\kappa_\beta=\partial_\beta D^\lambda_t|_{\lambda=0}$ for $\#\beta\ge2$. (b) The cumulant is a mixed derivative at $0$ of the logarithm of the normalised partition function with the sources.

The histories with a first creation of a large field at a level $t$ weigh at most $\mathsf w_t$. Along the never-large history the mixed derivative families of the order $p$ are bounded in the norm of the rate $\kappa^+=(1+4\beta_s)\kappa$; read at the rate $\kappa$ they keep the factor $e^{-4\beta_s\kappa[D_j/M-2]_+}$, $D_j$ the least sup distance of the carriers in level-$j$ units. On the way from the finer levels up to the working level $k$, a family of the level $i$ enters the step of the level $j\ge i$ with the weight $(n+1)L^{-(n-1)}$ of the lag $n=j+1-i$, while the size of a source against the bound of the step is at the level $i-1$ at most $(\frac43)^{1/2}L^{n/2}(1+n)^{\mathfrak p_0}$ times that at the level $j$; so the series $\mathsf T_p$ of the statement sums the mixed derivatives of the order $p$ over the levels finer than the working level,
\[
\mathsf T_p(u)=\sum_{n\ge1}(n+1)L^{-(n-1)}\Bigl[\bigl(\tfrac43\bigr)^{1/2}L^{n/2}(1+n)^{\mathfrak p_0}\Bigr]^p\,e^{-4\beta_s\kappa[L^{n-1}u-2]_+},
\]
and the induction over the levels up to $k+1$ closes if $\theta(x_j)\,\mathsf T_p(D_j/M)\le\frac12$ for $1\le j\le k$. Since $D_j/M\ge3$, $\theta(x_j)\le\theta(L^{s-2}g^{-2})$ and $\mathsf T_p$ does not increase, this holds under the condition $2\mathsf T_p(3)\theta(L^{s-2}g^{-2})\le1$ of the statement. Cauchy's estimates in each parameter then give the factor $\mathsf b^{-p}\prod_a|w_a|_\infty$. (c) is a recursion of numbers. If the sizes of $F_1$ fall like $L^{-j/4}$, the recursion of (a) gives $n_{j,1}\le\Gamma L^{-j/4}$ times the size at the level $j$ of a source of the weight $w_1$, $\Gamma$ a number formed with the constants: against the induction hypothesis at the level $j$, the lag $n$ contributes at most $(\frac43)^{1/2}L\,\theta(x_j)(n+1)^{1+\mathfrak p_0}L^{-n/4}$, and the induction closes with the sum $\mathsf S_{\rm nat}=\sum_{n\ge1}(n+1)^{1+\mathfrak p_0}L^{-n/4}$ of the statement.

In the induction of (b) the falling sizes of the first source compare the level $i-1$ with the level $j$ with the further factor $L^{n/4}$, so that the series $\tilde{\mathsf T}_p$ of the statement, which is $\mathsf T_p$ with the further factor $L^{n/4}$ in its $n$-th term, takes the place of $\mathsf T_p$. Both inductions close under the condition $2\max\{\tilde{\mathsf T}_p(3),(\frac43)^{1/2}L\,\mathsf S_{\rm nat}\}\,\theta(L^{s-2}g^{-2})\le1$ of (c), which implies the one of (b) at the same $s$; then the bound of (b) holds with $L^{-(k-1)/4}\tilde k_p(g,s;\ell)$ in the place of $k^{\rm loc}_p(g,s;\ell)$, $\tilde k_p\ge k^{\rm loc}_p$ a number of the same kind formed with these series, for the sizes at the levels $1\le j\le k$ and the values on the backgrounds of the levels $1\le t\le k$.

So the recursion closes under the condition of (c); the least integer $s\ge3$ at which it holds with $p=r$ is the number $\tilde s(r)$ of the statement, and we write $\tilde k_r$ for the number $\tilde k_r(g,s;\ell)$ of the bound of (c). (d) Let $c=\sum_i\kappa_im_i$ be a block form of Lemma~\ref{C-360-8}(b), $\|c\|=\sum_i|\kappa_i|2^{n_a(m_i)}$ its norm as in the statement, with $m_i$ written also for the order of the monomial $m_i$, and $k$ a real function of the sites; the source is $\mathrm{Fm}^c(k)=a^{-1}\sum_xk(x)c_x(U)$. At a height $j\ge1$ it is represented by the family $\eta_j^{-1}\sum_xL^{K-j}k(x)c_x$, $\eta_j=L^{-j}$, so that $\eta_j^{-1}L^{K-j}=a^{-1}$, the sum over the sites $x$ whose block lies in a given region of the level $j$; such a region has at most $8$ cubes and the linear size $0$, and at most $8M^3\eta_j^{-3}$ sites have one region (each of the at most eight cubes contains $(M/\eta_j)^3$ bare sites).

On the cores of radii $(b_0,b_1)$ in Definition~\YMMref{49.1}{374} the vector part of the holonomy of a lasso is at most $5.04\,b_0\eta_j^2$ in level-$j$ units and its plaquette term at most $(5.04\,b_0\eta_j^2)^2$, so a monomial of the order $m$ is at most $(5.04\,b_0)^m\eta_j^{2m}$, and the family has the size $8M^3L^{K-j}|k|_\infty\sum_i|\kappa_i|(5.04\,b_0)^{m_i}\eta_j^{2m_i-4}$. The sizes are taken at $(b_0,b_1)=(\alpha_0,\alpha_1)$: as $5.04\,\alpha_0L^{-2}\le1$, $(5.04\,\alpha_0)^{m_i}\eta_j^{2m_i-4}\le25.41\,\alpha_0^2$, and with the factor $2$ of the size of a representation at a height $j\ge1$ this is at most $34\|c\||k|_\infty$ in the unit $12M^3\alpha_0^2L^{K-j}$ of the size of a source of the weight $1$.

At the height $0$ the representation is the family in the fourth place $a^{-1}\sum_xk(x)[c_x(e^{iB'}U)-c_x(U)]$, $B'$ the fluctuation field with its values at the pivot bonds, which sums to the fluctuation variation of the source; a lasso has at most ten factors, and telescoping in them gives the size $911\,\mathsf c^\flat M^3\alpha_0^2L^K\|c\||k|_\infty$, that is $\frac{911}{12}\|c\||k|_\infty\le76\|c\||k|_\infty$ in the unit $12\mathsf c^\flat M^3\alpha_0^2L^K$ of the height $0$. So $76\|c\||k|_\infty$ bounds the sizes at all the heights, being at least $34\|c\||k|_\infty$ at the heights $j\ge1$ and $\frac{911}{12}\|c\||k|_\infty$ at the height $0$. On the real configurations one factor of a monomial is at most $a_p+a_{p'}$, $p,p'$ plaquettes of its lassos: a factor $\langle w,w'\rangle$ or $\det(w,w',w'')$, since $|w|^2\le2a_p$ and $|w|\le1$, or a plaquette term itself; every other factor is at most $1$, or $2$ if it is a plaquette term; so $|\mathrm{Fm}^c(k)|\le\mathcal O(2\|c\||k|_\infty\mathbf 1_{{\rm Pl}(k)})$, ${\rm Pl}(k)$ the plaquettes of the blocks of the sites with $k\ne0$, and the set $\{k\ne0\}$ carries all the representations: $\mathrm{Fm}^c(k)$ is a dominated local source with the comparison weight $\varpi\mathbf 1_{{\rm Pl}(k)}$ for every $\varpi\ge76\|c\||k|_\infty$.

If all the orders are $\ge3$, the factor $(5.04\,\alpha_0\eta_j^2)^{m_i-2}\le\max\{1,5.04\,\alpha_0\}^{{\rm Ord}(c)-2}L^{-2j}$, ${\rm Ord}(c)$ the largest order, gives the fall $L^{-2j}$ of the sizes; on the backgrounds of the level $t$ the plaquette terms are below $\bar a=2B_3^2\varepsilon_t^2L^{-4t}$, $\varepsilon_t$ the radius of regularity of the level $t$, a factor of the order $q\in\{2,3\}$ is at most $(2\bar a)^{q/2}$, and the values fall like $L^{-2t}$. With the comparison weight $\Gamma^\circ_c\|c\||k|_\infty\mathbf 1_{{\rm Pl}(k)}$ of the statement, $\Gamma^\circ_c=\max\{76,\,34\max\{1,5.04\,\alpha_0\}^{{\rm Ord}(c)-2},\,2\max\{1,2B_3\}^{{\rm Ord}(c)-2}\}$, this is the hypothesis of (c), with $L^{-2j}$ and $L^{-2t}$. (e) The density of a source at the never-large history of the level $t$ is its first derivative in (a); its oscillation is bounded by the size of the source at that level, which is $y_2$ up to a constant, and the mixed part is the second mixed derivative, which carries the factor of the distance from (a), with $\mathsf D_t\ge L^{s'}\mathsf d'\ge3M$.

The second mixed derivatives are summed over the levels finer than $t$ as in (b) at $p=2$, $\frac43L\,\mathsf T^\circ_2$ being the series $\mathsf T_2$, and the condition $\mathsf h_2(s')\le1$ closes that induction up to the level $t$.
\end{proof}

\begin{lemma}[Rough entries]\label{C-360-6}
A petal at a site $x$ is a pair $\pi=(\alpha,\alpha')$ of signed directions with $|\alpha|\ne|\alpha'|$, $W_x(\pi)$ the holonomy from $x$ of the plaquette that it spans, and $q_x(\pi,\pi')=\langle\mathbf w(W_x(\pi)),\mathbf w(W_x(\pi'))\rangle$. For real $k=(k_{\pi\pi'}(x))$ put ${\rm Pf}(k)=a^{-1}\sum_x\sum_{\pi,\pi'}k_{\pi\pi'}(x)q_x(\pi,\pi')$ and $|k|_1=\sup_x\sum_{\pi,\pi'}|k_{\pi\pi'}(x)|$. A \emph{rough entry} is $F={\rm Pf}(k)+\mathcal O(v)$ with a real weight $v$; its norm is $\|F\|=13|k|_1+|v|_\infty$, and it is carried by the barycentres of the plaquettes of ${\rm supp}\,v$ and of those with a corner at which $k\ne0$. Let $0<g\le\gamma^{\rm ref}_W$, $\ell'\in\mathfrak L(g)$, $r\ge2$ and $0<\mathsf d<1$. There are $K'_0(r,\mathsf d)$ and $k'(r,\mathsf d)\le C_{r,p}\mathsf d^{-pr}$ $(p>7)$ such that
\[
|\langle F_1;\dots;F_r\rangle_{K,\ell',\beta}|\le k'(r,\mathsf d)\prod_a\|F_a\|
\]
for all $(K,\beta)\in\mathfrak D(g)$ with $K\ge K'_0(r,\mathsf d)$ and all rough entries whose carriers are pairwise $\ge\mathsf d$ apart. Moreover $\Theta_\beta[\eta]=\mathsf b\,{\rm Pf}(k^\eta)$ with $|k^\eta|_1\le12\max_\rho{\rm Lip}_\infty\eta^\rho$, so $\|{\rm Pf}(k^\eta)\|\le156\max_\rho{\rm Lip}_\infty\eta^\rho$; $-\beta\sum_ph_{\omega,\chi}(p)=\mathsf b\,{\rm Pf}(k^h)$ with $|k^h|_1\le12|\omega|\,{\rm Lip}(\chi)R_1(\chi)\le36\ell|\omega|\,{\rm Lip}_\infty\chi$; $\mathfrak B_{\omega,\chi}=\mathsf b\,{\rm Pf}(k^h)+\mathcal O(v^{\mathfrak B})$ with $|v^{\mathfrak B}|_\infty\le(2\mathsf b+3a)|\omega|\,{\rm Lip}(\chi)R'_1(\chi)$, where $R'_1(\chi)\le\frac32\ell$ is the number of Theorem~\ref{C-360-3}(a), the largest $|x(p)|_1$, taken in the chart $(-\frac\ell2,\frac\ell2)^3$, at the plaquettes where $\nabla^s_\omega\chi$ or $v_\chi$ does not vanish (at their barycentres $x(p)$, $\chi$ is not constant on the points $x(p)\pm ae_\rho$, $x(p)\pm\frac a2e_\rho$, $\rho=1,2,3$); the entries $\mathsf b\,{\rm Pf}(k^\eta)$, $\mathsf b\,{\rm Pf}(k^h)$ and $\mathfrak B_{\omega,\chi}$ are carried within $2a$ of ${\rm supp}\,\eta$, resp.\ ${\rm Var}(\chi)$; and $T_{lr}(\varphi)=\mathsf b\,({\rm Pf}(k^{lr})+\mathcal O(v^{lr}))$ with $13|k^{lr}|_1+|v^{lr}|_\infty\le C_T\sup|\varphi|$, $C_T$ an absolute constant.
\end{lemma}

\begin{proof}[Proof sketch]
A rough entry is a dominated local source in the sense of Lemma~\ref{C-360-24}: a gauge invariant polynomial with a local representation at every height, with $|F|\le\mathcal O(w)$ for a weight $w$ of size $\|F\|$ on its carrier. Put $\sum_a\lambda_aF_a$ into the weighted progression (Lemma~\ref{C-360-24}(a)). Its families are holomorphic in $\lambda$ near $0$ with sizes that depend on the $\sup|w_a|$ alone, and they are local: $\partial_BE^{(i)}_\lambda(X)=0$ unless $X$ meets every carrier in $B$, so that $d_i(X)\ge D_i/M-2$, $D_i$ being the least sup distance of these carriers in level-$i$ units, and a factor $e^{-\kappa^+[D_i/M-2]_+}$ appears. We split the mean over the level $t$ of the first creation of a large field. Up to that level the history is the never-large one, and
\[
\begin{multlined}
\tau_{h_t}\bigl[\rho_0e^{\sum_a\lambda_aF_a}\bigr]=e^{D^\lambda_t(U_t(\cdot))}\,\mathsf m_t,\qquad\Phi\Bigl[\prod_{a\in B}(F_a-c_a)\Bigr]=\sum_\pi\prod_{\beta\in\pi}\kappa_\beta,\quad\kappa_\beta=\partial_\beta D^\lambda_t|_{\lambda=0}\ \relax (\#\beta \\
\ge2),\quad\kappa_{\{a\}}=\partial_aD^\lambda_t|_{\lambda=0}-c_a,
\end{multlined}
\]
with the signs of Lemma~\ref{C-360-24}(a); each $\kappa_\beta$ with $\#\beta\ge2$ is bounded by the mixed derivatives above. Behind the level $t$ only a positive weight enters, of the mass $\mathsf w_t\le2M^{-3}\ell'^3e^{B(X_g)\ell'^3}L^{3(K-t)}e^{-c'p_0(g_t)^2}$, which is smaller than every power of $L^{-(K-t)}$, while the centred insertions carried to the level $t-1$ are at most powers of $L^{K-t}$. The mixed derivatives of the orders $\le r$ converge at the depth $s(\mathsf d'')$, $\mathsf d''=\min\{\frac45\mathsf d,L^{-\bar s(r)}\}$, where $\bar s(r)$ is the least integer $s\ge3$ with $2\mathsf T_r(3)\theta(L^{s-2}g^{-2})\le1$, with the series $\mathsf T_r$ and the size $\theta$ of Lemma~\ref{C-360-24}(b); at every $s\ge\bar s(r)$ the condition holds, since $\theta$ does not increase.

Apply Proposition~\ref{C-355-20}(a$'$), (b) to the dominated local sources just constructed. Its summability condition is the one defining $\bar s(r)$ above, and it permits the fixed side $\ell'$: the local cumulant constant may be chosen with $k_r(g,s;\ell')\le C_r(g,\ell')Y_s^{8r\mathfrak p_0}$. Cutting every weight into at most $(2\ell'/\mathsf d'')^3$ pieces and using $\mathsf b^{-r}\le g^{2r}$ therefore gives $k'(r,\mathsf d)\le C_r(g,\ell')g^{2r}(2\ell'/\mathsf d'')^{3r}Y_{s(\mathsf d'')}^{8r\mathfrak p_0}$. The same proposition bounds $Y_{s(\mathsf d'')}\le2\log(1/\mathsf d'')+C(g)$; since $\mathsf d''=\min\{4\mathsf d/5,L^{-\bar s(r)}\}$, this proves the asserted $C_{r,p}\mathsf d^{-pr}$ bound for every $p>7$. The same decomposition, at the least depth $s$ with $L^s\mathsf d\ge3M$ at which the mixed derivatives of the order $q$ are summable, gives the second bound of Theorem~\ref{C-360-2}(a): $L^s\le\max\{1,LM\mathsf u_q/\mathsf d\}$ with $\mathsf u_q=O(q\log q)$, the sums over partitions are at most $C^q(q!)^{3/2+\mathfrak p_0}$, and the rarity of a first creation, at which all the $q$ insertions are bounded by their suprema, costs $\exp[c_L(4q+4)^2]$. The formulas for $k^\eta$, $k^h$ and $k^{lr}$ are the definitions written in petals: for $b=\langle x,x+ae_\mu\rangle$ and $p=p(b,\alpha'')$, $\mathfrak d_{b,p,\rho}=-\frac12[q_x((\mu,\rho),(\mu,\alpha''))-q_x((\mu,-\rho),(\mu,\alpha''))]$, and the shell lies in $|x|_1\le\frac12+\frac32\ell$.

For the whole flux insertion, $\mathfrak B_{\omega,\chi}-\mathsf b\,{\rm Pf}(k^h)=\mathcal O(v^{\mathfrak B})$ with $v^{\mathfrak B}(p)=\mathsf b\,\nabla^s_\omega\chi(x(p))-\frac32a\,v_\chi(p)$; as $\sum_\rho|(\omega y)_\rho|\le2|\omega||y|_1$ and the difference quotients of $\chi$ in $\nabla^s_\omega\chi$ and in $v_\chi$ are at most ${\rm Lip}(\chi)$, $|v^{\mathfrak B}(p)|\le(2\mathsf b+3a)|\omega|\,{\rm Lip}(\chi)|x(p)|_1$, and $v^{\mathfrak B}(p)\ne0$ only at the plaquettes where $\nabla^s_\omega\chi$ or $v_\chi$ does not vanish, at whose barycentres $\chi$ is not constant on the points $x(p)\pm ae_\rho$, $x(p)\pm\frac a2e_\rho$: this gives the factor $R'_1(\chi)$ of Theorem~\ref{C-360-3}(a).
\end{proof}

\begin{lemma}[Chessboard bound]\label{C-360-7}
Let $0<g\le\gamma$, $\ell'\in\mathfrak L(g)$, $(K,\beta)\in\mathfrak D(g)$, and let a board of cubic cells of side $\mathsf l'$ dividing $\frac{\ell'}2$ be placed on $T_{\ell'}$, on every $T_\ell$ with $\ell\ge\ell'$ in $\mathfrak L(g)$ and on $a_K\mathbb Z^3$; $N'=(\ell'/\mathsf l')^3$. For a bounded Borel function $F$ of the bonds inside one cell put $\|F\|_{\rm ch}=|\langle\prod_\kappa F^{[\kappa]}\rangle_{K,\ell',\beta}|^{1/N'}$, the product running over the mirror copies of $F$, conjugated on odd reflection parity, in the $N'$ cells of $T_{\ell'}$. Then
\[
|s(F_1\cdots F_q)|\le\prod_ie^{G_g\mathsf l'^3}\|F_i\|_{\rm ch}
\]
for functions $F_1,\dots,F_q$ of distinct cells, $s$ being a periodic Gibbs state at $(K,\beta)$ or the measure of $T_\ell$, $\ell\in\mathfrak L(g)$, $\ell\ge\ell'$. In particular, for real $F_1,F_2$ of two distinct cells and real $c_1,c_2$: $|s(F_1)-c_1|\le e^{G_g\mathsf l'^3}\|F_1-c_1\|_{\rm ch}$ and $|s(A^s[F_1]A^s[F_2])|\le2e^{2G_g\mathsf l'^3}\|F_1-c_1\|_{\rm ch}\|F_2-c_2\|_{\rm ch}$.
\end{lemma}

\begin{proof}[Proof sketch]
On $T_\ell$, $\ell\ge\ell'$, reflection positivity at the walls gives the chessboard estimate $|\langle F_1\cdots F_q\rangle_{K,\ell,\beta}|\le\prod_i\langle\prod_\kappa F_i^{[\kappa]}\rangle_{K,\ell,\beta}^{1/N_\ell}$, $N_\ell=(\ell/\mathsf l')^3$. The product of the mirror copies over all the cells of $T_\ell$ is a periodisation of the product over the cells of $T_{\ell'}$, and the chessboard seminorms on the two sides compare by the inequality $\langle\prod_\kappa F^{[\kappa]}\rangle_{K,\ell,\beta}^{1/N_\ell}\le\langle\prod_\kappa F^{[\kappa]}\rangle_{K,\ell',\beta}^{1/N'}\,e^{\mathsf l'^3(p_{\ell'}-p_\ell)}$ (Corollary~\ref{IV-18-18}(e)). It is reflection positivity at the walls applied once more, to the periodised product: the pattern of the copies has the period $2\mathsf l'$, which divides $\ell'$, so that a slab of that length is a half slab and its mirror image, the transfer operator across it is positive, and ${\rm Tr}\,T^n\le({\rm Tr}\,T)^n$ in each of the three directions; this bounds the unnormalised $N_\ell$-th root by the unnormalised $N'$-th root, and the ratio $Z_{K,\ell',\beta}^{1/N'}/Z_{K,\ell,\beta}^{1/N_\ell}$ of the plain partition functions per cell is $e^{\mathsf l'^3(p_{\ell'}-p_\ell)}\le e^{G_g\mathsf l'^3}$, the pressures $p_\ell=\ell^{-3}\log Z_{K,\ell,\beta}$ being formed with the plain partition functions (Section~\YMMref{86}{680}).

A periodic Gibbs state is dominated by the torus measures through its conditional laws (Proposition~\YMMref{91.9}{705}), and this domination carries the bound to every periodic Gibbs state. The last two bounds are the cases $q=1$ and $q=2$, with $s(A^s[F_1]A^s[F_2])=s((F_1-c_1)(F_2-c_2))-(s(F_1)-c_1)(s(F_2)-c_2)$.
\end{proof}

\begin{lemma}[Exact identities]\label{C-360-8}
\begin{enumerate}
\item[(a)] Let $\mathbf X$ be a polynomial field of directions that vanishes off a finite set $\Lambda$ of bonds: at each positively oriented bond $b$ a $\mathfrak g$-valued polynomial $\mathbf X(b)$ in the plaquette variables, $0$ for $b\notin\Lambda$; let $v_{\mathbf X}F=\sum_b\sum_{\mathsf a}\langle\sigma_{\mathsf a},\mathbf X(b)\rangle\partial^{\mathsf a}_bF$ and ${\rm div}\,v_{\mathbf X}=\sum_b\sum_{\mathsf a}\partial^{\mathsf a}_b\langle\sigma_{\mathsf a},\mathbf X(b)\rangle$, $\sigma_1,\sigma_2,\sigma_3$ the Pauli matrices and $\partial^{\mathsf a}_b$ the derivative along $e^{i\epsilon\sigma_{\mathsf a}}U(b)$ at $\epsilon=0$. Then $\nu((\beta v_{\mathbf X}A-{\rm div}\,v_{\mathbf X})Y)=0$ for every Gibbs state or torus measure $\nu$ at $(a_K,\beta)$ and every bounded Borel $Y$ that does not depend on the bonds of $\Lambda$; so every joint cumulant under $\nu$ with the entry $\beta v_{\mathbf X}A-{\rm div}\,v_{\mathbf X}$ and further entries not depending on $\Lambda$ vanishes. $N_\beta[\eta]=\Theta_\beta[\eta]+\mathfrak K^{\rm v}[\beta\eta]-\mathsf b\,\mathcal O({\rm div}^s\eta)+\frac32a\,\mathcal O(e^\eta)$ is such a null term.

With the contact terms, for all weights $f_i$ and every vector field $\eta$, $\langle N_\beta[\eta];\mathcal O(f_\cdot)\rangle+\sum_i\langle{\rm Ct}[f_i;\eta];(\mathcal O(f_l))_{l\ne i}\rangle=0$, where ${\rm Ct}[f;\eta]=\mathcal O({\rm div}^s(f\eta))-a^{-1}\Theta_f[\eta]-a^{-1}\mathfrak K^{\rm v}[f\eta]$.
\item[(b)] Identically in $U$, $\mathfrak K^{\rm v}[\eta]=-\sum_p\sum_\rho(\nabla^+_\rho\eta^\rho)(x(p))\Psi_\rho(p)+\mathfrak K^\sharp[\eta]$; for $\rho\in P(p)$, $\mathfrak r^\sharp_\rho(p)=0$ and $-(a_p+a_{p+ae_\rho})\le\Psi_\rho(p)\le0$; $\mathfrak r^\sharp$ transforms as a vector under the coordinate mirrors of the lattice: $\mathfrak r^\sharp_\rho$ is odd under the mirror normal to $\rho$ and even under the two mirrors that contain $e_\rho$, so that it is odd, for every $\rho$, under the inversion $y\mapsto2z-y$ through a site $z$, the product of the three mirrors through $z$. Hence $\langle\mathfrak K^\sharp[\eta]\rangle_{K,\ell,\beta}=0$, and by translation invariance also $\langle\mathfrak K^{\rm v}[\eta]\rangle_{K,\ell,\beta}=0$; and if a field $\zeta$ has $\zeta^\rho(y+se_\rho)=\zeta^\rho(y)$ for all $y$, $s$ and $\rho$, then $\langle\mathfrak K^\sharp[\zeta];\mathcal O(f)\rangle_{K,\ell,\beta}=0$ for every weight $f$.

A \emph{block form} is a combination $c=\sum_i\kappa_im_i$ of monomials in the lassos of the block ${\rm Bk}(x)=x+a\{-1,0,1\}^3$ of $27$ sites centred at a site $x$: a lasso is a path of at most three steps from $x$ inside the block to a corner $y$ of a plaquette of the block, the holonomy of that plaquette from $y$, and the path back; with $w=\mathbf w(W)$ and $a_{\rm lz}=1-\frac12{\rm tr}\,W$ the vector part and the plaquette term of the holonomy $W$ of a lasso, a monomial is a product of $n_2$ factors $\langle w,w'\rangle$, $n_3$ factors $\det(w,w',w'')$ and $n_a$ factors $a_{\rm lz}$, and $2n_2+3n_3+2n_a$ is its order in the field strength. The same lassos are taken at every site, $c_x$ is the form at the block of $x$, the order of $c$ is the least order of its monomials, and its norm is the number $\|c\|=\sum_i|\kappa_i|2^{n_a(m_i)}$ of Lemma~\ref{C-360-24}(d).

For a real function $k$ of the sites, $\mathrm{Fm}^c(k)=a^{-1}\sum_xk(x)c_x(U)$ is a dominated local source of size $\le76\|c\||k|_\infty$, carried by $\{k\ne0\}$, with $|\mathrm{Fm}^c(k)|\le\mathcal O(2\|c\||k|_\infty\mathbf 1_{{\rm Pl}(k)})$, ${\rm Pl}(k)$ the plaquettes of the blocks of the sites with $k\ne0$; and if all the monomials of $c$ have order $\ge3$, as for the commutator terms, the defects of the linear Bianchi identity, $a_p^2$ and $a_pa_{p_+}$, its sizes fall like $L^{-2j}$ with the level $j$ and its values on the backgrounds like $L^{-2t}$ with the level $t$, with the comparison weight $\Gamma^\circ_c\|c\||k|_\infty\mathbf 1_{{\rm Pl}(k)}$ of Lemma~\ref{C-360-24}(d). $\Psi_{\lambda(p)}(p)$ is a block form of order $2$ and norm $1$, and $\mathfrak r^\sharp_\rho(p)$, $\rho=\lambda(p)$, a block form of order $2$ and norm $\le8$, both at the block of the lowest corner of $p$; so for a real weight $k$ the sums $a^{-1}\sum_pk(p)\Psi_{\lambda(p)}(p)$ and $a^{-1}\sum_pk(p)\mathfrak r^\sharp_{\lambda(p)}(p)$ over the plaquettes of one normal are sources $\mathrm{Fm}^c(\tilde k)$, $\tilde k(x)$ the value of $k$ at the plaquette of that normal with the lowest corner $x$, and they are dominated local sources of the sizes $\le76|k|_\infty$ and $\le608|k|_\infty$; as a function $|\mathfrak r^\sharp_\rho(p)|\le C\bar a^{3/2}$, $\bar a$ the largest $a_q$ near $p$, and $\beta\mathfrak K^\sharp$ is treated by parity (Lemmas~\ref{C-360-9}, \ref{C-360-10}).
\item[(c)] For a vector field $\mathsf k$, $\beta\sum_{p,\rho}\mathsf k^\rho(x(p))\Psi_\rho(p)=-\frac14\mathfrak N[\mathsf k]+\sum_{v=1}^5\mathfrak J_v[\mathsf k]-\frac14c_{\mathsf k}$, where $\mathfrak N[\mathsf k]=\beta v_{\mathbf X^{\mathsf k}}A-{\rm div}\,v_{\mathbf X^{\mathsf k}}$ is the null term of (a) for the field of directions $\mathbf X^{\mathsf k}(b)=\sum_{\rho\ne\mu_b}\mathsf k^\rho(x_b)G_\rho(b)$, $c_{\mathsf k}=6\sum_b\sum_{\rho\ne\mu_b}\mathsf k^\rho(x_b)$ is a number, and $\mathfrak J_1,\dots,\mathfrak J_5$ are sums of local densities of the orders $2,2,2,3,4$ with weights at most $3|\mathsf k|_\infty+\frac32\mathsf b{\rm Lip}_\infty\mathsf k$, $\frac18\mathsf b{\rm Lip}_\infty\mathsf k$, $\frac12\mathsf b{\rm Lip}_\infty\mathsf k$, $\frac14\beta|\mathsf k|_\infty$ and $\beta|\mathsf k|_\infty$, vanishing beyond the sup distance $a$ of ${\rm supp}\,\mathsf k$. Consequently, with $\mathsf k^\rho_\eta=\nabla^+_\rho\eta^\rho$, $\mathrm{St}_K[\eta]=N_\beta[\eta]-\frac14\mathfrak N[\mathsf k_\eta]-\mathrm{Rm}_K[\eta]-\frac14c_{\mathsf k_\eta}$, where $\mathrm{Rm}_K[\eta]=\beta\mathfrak K^\sharp[\eta]-\sum_v\mathfrak J_v[\mathsf k_\eta]+\frac32a\,\mathcal O(e^\eta)$ is the remainder of the stress insertion.
\item[(d)] $\Delta_K[\eta]=F^0+F^1$, rough entries carried within $2a$ of ${\rm supp}\,\eta$, with $\|F^0\|\le286\mathsf b\sup|\nabla\eta|$ and $\|F^1\|\le780\mathsf b\,a\sup|\nabla^2\eta|$; $F^0$ is null: its classical density, a quadratic form in the field strength with coefficients linear in $\nabla\eta$, vanishes identically.
\item[(e)] For a positively oriented bond $b$ and directions $\lambda,\rho$: $\sum_{p\ni b}a^{-1}(x_b-x(p))_\lambda\mathfrak d_{b,p,\rho}=\langle\mathbf u_{b,\lambda},\mathbf u_{b,\rho}\rangle$, so that $\sum_p\mathfrak s_{\lambda\rho}(p)=\sum_b\langle\mathbf u_{b,\lambda},\mathbf u_{b,\rho}\rangle$ is a symmetric positive semidefinite matrix; for $\lambda\ne\mu_b$ and $p_\pm=p(b,\pm\lambda)$, $|\mathbf u_{b,\lambda}|^2=(a_{p_+}+a_{p_-})-\frac14(a_{p_+}^2+a_{p_-}^2)+\Psi_\lambda(p_-)$ and $\Psi_\lambda(p_-)=-\frac14|G_\lambda(b)|^2-\frac14(a_{p_+}^2+a_{p_-}^2)$; for $\nu\in P(p)$, $\mathfrak s_{\nu\nu}(p)=\frac32a_p+\frac14(a_{p+ae_\nu}+a_{p-ae_\nu})-\frac12a_p^2+\frac12(\Psi_\nu(p)+\Psi_\nu(p-ae_\nu))$, and $\mathfrak s_{\lambda\rho}(p)=0$ if $\lambda=\lambda(p)$. Consequently, identically in $U$,
\[
T_{mm}(\varphi)-\mathsf b\check D_m(\varphi)=\beta\sum_{p:\,m\in P(p)}\varphi(x(p))\Bigl[\tfrac12\bigl(\Psi_m(p)+\Psi_m(p-ae_m)\bigr)-\tfrac12a_p^2\Bigr]+\mathsf b\,\mathcal O(\tilde\varphi_m),
\]
$\tilde\varphi_m(p)=\frac14[\varphi(x(p)+ae_m)-2\varphi(x(p))+\varphi(x(p)-ae_m)]\mathbf 1\{m\in P(p)\}$, so $|\tilde\varphi_m|\le\frac12a\,{\rm Lip}\,\varphi$; and $\sum_mT_{mm}(\varphi)-\mathsf b\,\mathcal O(\varphi)=\sum_m(T_{mm}(\varphi)-\mathsf b\check D_m(\varphi))$.
\item[(f)] \emph{The null term against sources that meet its carrier.} For a complex vector field $\mathsf k$ and a complex function $f$ on $T_\ell$ put $\mathsf e^{\mathsf k}(p)=\sum_{\rho\in P(p)}[\mathsf k^\rho(x(p)+\frac a2e_\rho)+\mathsf k^\rho(x(p)-\frac a2e_\rho)]$ and
\[
\begin{multlined}
\mathfrak C'[\mathsf k;f]=\sum_p\sum_{b\subset\partial p}\sum_{\rho\ne\mu_b}a^{-1}\bigl(f(x(p))-f(x_b)\bigr)\mathsf k^\rho(x_b)\,\langle G_\rho(b),\mathbf w(W_{b,p})\rangle,\qquad\mathfrak R[\mathsf k;f] \\
=\mathsf b^{-1}{\rm div}\,v_{\mathbf X^{f\mathsf k}}+\mathfrak C'[\mathsf k;f],
\end{multlined}
\]
with $\mathbf X^{\mathsf k}$ and $\mathfrak N[\mathsf k]$ as in (c). Then: (N1) $v_{\mathbf X^{\mathsf k}}\mathcal O(f)=\mathsf b^{-1}\mathfrak N[f\mathsf k]+\mathfrak R[\mathsf k;f]$, identically in $U$; (N2) $\int\mathfrak N[\mathsf k]\,e^{\mathcal O(w)}e^{-\beta A}dU=\int\bigl(v_{\mathbf X^{\mathsf k}}\mathcal O(w)\bigr)e^{\mathcal O(w)}e^{-\beta A}dU$ for every complex function $w$ on $T_\ell$; (N3) for a finite set $\mathsf A$ and complex functions $f_l$, $l\in\mathsf A$,
\[
\bigl\langle\mathfrak N[\mathsf k];(\mathcal O(f_l))_{l\in\mathsf A}\bigr\rangle_{K,\ell,\beta}=\sum_{j=1}^{\#\mathsf A}\ \sum_{(q_1,\dots,q_j)}\mathsf b^{-(j-1)}\Bigl\langle\mathfrak R\bigl[f_{q_1}\cdots f_{q_{j-1}}\mathsf k;f_{q_j}\bigr];(\mathcal O(f_l))_{l\in\mathsf A\setminus\{q_1,\dots,q_j\}}\Bigr\rangle_{K,\ell,\beta},
\]
the inner sum over the sequences of $j$ distinct elements of $\mathsf A$ (for $\mathsf A=\emptyset$ the left member is the mean, and it is $0$); (N4) $\mathfrak R[\mathsf k;f]=\mathsf b^{-1}c_{f\mathsf k}-3\mathsf b^{-1}\sum_pa_p\,\mathsf e^{f\mathsf k}(p)+\mathfrak C'[\mathsf k;f]$ is a number, the density $a_p$ with weights $\le12\,\mathsf b^{-1}\sup|f|\sup|\mathsf k|$ and the densities $\langle G_\rho(b),\mathbf w(W_{b,p})\rangle$ of the order $2$ with weights $\le\frac12{\rm Lip}_\infty f\,\sup|\mathsf k|$: no factor $\beta$ enters. Its weight at a plaquette $p$ vanishes unless a bond $b$ of $p$ has $\mathsf k(x_b)\ne0$ and $f\ne0$ at $x_b$ or at $x(p)$, so that its densities are carried within the sup distance $2a$ of ${\rm supp}\,\mathsf k$ and of ${\rm supp}\,f$, and the number vanishes unless $\mathsf k(x_b)f(x_b)\ne0$ at some bond.

Consequently, if the spread $D$ of ${\rm supp}\,\mathsf k$ and the ${\rm supp}\,f_l$, $l\in\mathsf A$, is $>4a$, the terms of (N3) with $j=\#\mathsf A$ vanish, the numbers drop from the other terms, and in each of them the entry is carried by a set whose spread with the supports of the remaining $f_l$ is $\ge D-2a$. This rests on the merging of sets: if $S_0,\dots,S_n$ have the spread $D$, $Q\subset\{1,\dots,n\}$, $2r<D$, and $T$ lies within the sup distance $r$ of $S_0$ and of every $S_q$, $q\in Q$, then $T=\emptyset$ if $Q=\{1,\dots,n\}$, and otherwise the spread of $T$ and the $S_l$, $l\notin Q$, is $\ge D-r$.
\end{enumerate}
\end{lemma}

\begin{proof}[Proof sketch]
(a) The Haar measure is invariant; with $\partial^{\mathsf a}_b\langle\sigma_{\mathsf a},\mathbf w(W_{b,p})\rangle=1-a_p$, an integration by parts in each bond variable against $e^{-\beta A}$ gives $\nu((\beta v_{\mathbf X}A-{\rm div}\,v_{\mathbf X})Y)=0$ in finite volume, and in a Gibbs state through its conditional laws. With sources the same integration by parts against $e^{-A(\beta-a^{-1}w_z,\cdot)}$, $w_z=\sum_iz_if_i$, followed by the mixed derivative in $z$ at $0$, gives the contact terms, since $v_\eta A(v,\cdot)$ is linear in the weight $v$. (b) For $q_\pm=p(b,\pm\rho)$ in one plane, $\langle\mathbf w(W_{b,q_-}),\mathbf w(W_{b,q_+})\rangle=-2\Psi_\rho(q_-)-a_{q_-}-a_{q_+}$; so $\mathfrak r_\rho(p)=\Psi_\rho(p)-\Psi_\rho(p_-)$ in the plane, and the split is a summation by parts.

Under the mirror $\mathsf r$ normal to $\rho$, $\mathfrak d_{\mathsf rb,\mathsf rp,\rho}(\mathsf rU)=-\mathfrak d_{b,p,\rho}(U)$ and $\Psi_\rho(\mathsf rp)(\mathsf rU)=\Psi_\rho(p_-)(U)$, and the measure of a torus is invariant under $\mathsf r$. Under a coordinate mirror that contains $e_\rho$ these quantities are carried to the same quantities at the image bond and plaquette: the holonomy $W_{b,p}$ is taken from the initial point of $b$ along $b$, and where the mirror reverses the orientation of $b$ both factors of $\mathfrak d_{b,p,\rho}$ change their sign together; so $\mathfrak r^\sharp_\rho$ is even under it, and the three coordinate mirrors through a site compose to the inversion through it.

For $\zeta$: $\langle\mathfrak K^\sharp[\zeta];\mathcal O(f)\rangle=a^{-1}\sum_{p'}f(p')\langle\mathfrak K^\sharp[\zeta];a_{p'}\rangle$, and for each $p'$ and $\rho$ the mirror normal to $\rho$ through $x(p')$ keeps $a_{p'}$, the measure and $\zeta^\rho$, carries the plaquettes normal to $\rho$ among themselves, and changes the sign of $\mathfrak r^\sharp_\rho$ on them. A block form has local representations at every height of size $76\|c\||k|_\infty$ (the proof of Lemma~\ref{C-360-24}(d), with the lassos from the centre of the block, which reach its $27$ sites in at most three steps); the field strength is $O(L^{-2j})$ in level-$j$ units on the backgrounds, which gives the fall for monomials of order $\ge3$, with the comparison weight $\Gamma^\circ_c\|c\||k|_\infty\mathbf 1_{{\rm Pl}(k)}$ of Lemma~\ref{C-360-24}(d). (c) Take in (a) $\mathbf X(b)=\sum_\rho\mathsf k^\rho(x_b)[\mathbf w(W_{b,p(b,\rho)})+\mathbf w(W_{b,p(b,-\rho)})]$; then $v_{\mathbf X}a_p=\sum_{b\subset\partial p}\langle\mathbf X(b),\mathbf w(W_{b,p})\rangle$, ${\rm div}\,v_{\mathbf X}$ is the number $c_{\mathsf k}$ minus $3\sum_pa_p$ times a weight formed from $\mathsf k$, and the in-plane remainder is $-2a_p^2$, of order $4$. (d) At a constant field strength, with $G=\nabla\eta$ and $T_{lr}={\rm tr}B_lB_r$, the classical density of $\Delta_K[\eta]$ is $\mathsf b[{\rm tr}G\,{\rm tr}T-{\rm Tr}\,GT]-\frac{\mathsf b}2{\rm tr}G\,{\rm tr}T-\frac{\mathsf b}2{\rm Tr}(H^{\rm cons}T)=0$, since ${\rm Tr}(H^{\rm cons}T)={\rm tr}G\,{\rm tr}T-2{\rm Tr}\,GT$; the Taylor remainder of $\eta$ gives $F^1$. (e) $\sum_{p\ni b}a^{-1}(x_b-x(p))_\lambda\mathbf w(W_{b,p})=\mathbf u_{b,\lambda}$ and the form is bilinear; with $w_\pm=\mathbf w(W_{b,p_\pm})$, $|w_\pm|^2=a_{p_\pm}(2-a_{p_\pm})$ and $\Psi_\lambda(p_-)=-\frac14|w_++w_-|^2-\frac14(a_{p_-}^2+a_{p_+}^2)$, and the parallelogram law $4|\mathbf u_{b,\lambda}|^2=|w_+-w_-|^2=2|w_+|^2+2|w_-|^2-|w_++w_-|^2$ gives the first identity; the formula for $\mathfrak s_{\nu\nu}$ is the same computation at the two bonds of $p$ of the direction other than $\nu$, and no vector $x_b-x(p)$ has a component normal to $p$.

Inserting these in $T_{mm}(\varphi)=\beta\sum_p\varphi(x(p))[\mathfrak s_{mm}(p)-a_p]$ and $\mathsf b\check D_m(\varphi)=\beta\sum_p\varphi(x(p))a_p(1-2\cdot\mathbf 1\{\lambda(p)=m\})$, the plaquettes normal to $m$ cancel, and moving the shifts $a_{p\pm ae_m}$ onto $\varphi$ gives $\tilde\varphi_m$. Finally $\sum_m\check D_m=3\mathcal O-2\mathcal O=\mathcal O$. (f) $\mathcal O(f)=a^{-1}\sum_pf(x(p))a_p$ and $v_{\mathbf X}a_p=\sum_{b\subset\partial p}\langle\mathbf X(b),\mathbf w(W_{b,p})\rangle$; with $f(x(p))=f(x_b)+(f(x(p))-f(x_b))$ and $f(x_b)\mathbf X^{\mathsf k}(b)=\mathbf X^{f\mathsf k}(b)$, the first part is $a^{-1}v_{\mathbf X^{f\mathsf k}}A=\mathsf b^{-1}\beta v_{\mathbf X^{f\mathsf k}}A=\mathsf b^{-1}(\mathfrak N[f\mathsf k]+{\rm div}\,v_{\mathbf X^{f\mathsf k}})$ and the second is $\mathfrak C'[\mathsf k;f]$: this is (N1). (N2) is the integration by parts of (a) along $\mathbf X^{\mathsf k}$ against $e^{\mathcal O(w)}e^{-\beta A}$.

At the source $w_z=\sum_{l\in\mathsf A}z_lf_l$, (N2) says that the mean of $\mathfrak N[\mathsf k]-\sum_qz_q\,v_{\mathbf X^{\mathsf k}}\mathcal O(f_q)$ under the measure with the source vanishes, and this function is affine in $z$; the mixed derivative in $z$ at $0$ gives $\langle\mathfrak N[\mathsf k];(\mathcal O(f_l))_{l\in\mathsf A}\rangle=\sum_{q\in\mathsf A}\langle v_{\mathbf X^{\mathsf k}}\mathcal O(f_q);(\mathcal O(f_l))_{l\ne q}\rangle$, and inserting (N1), whose first part is a null term of the same kind with $f_q\mathsf k$ for $\mathsf k$, an induction on $\#\mathsf A$ gives (N3). (N4): ${\rm div}\,v_{\mathbf X^{\mathsf k}}=c_{\mathsf k}-3\sum_pa_p\,\mathsf e^{\mathsf k}(p)$, by $\partial^{\mathsf a}_b\langle\sigma_{\mathsf a},\mathbf w(W_{b,p})\rangle=1-a_p$; $|\mathsf e^{f\mathsf k}|\le4\sup|f\mathsf k|$, $|x(p)-x_b|_\infty=\frac a2$, and the boxes of the densities lie within the sup distance $\frac32a$ of $x(p)$.

The merging: if $Q=\{1,\dots,n\}$, points of $S_0,\dots,S_n$ within $r$ of one point of $T$ are pairwise at most $2r<D$ apart, which the spread forbids; otherwise, for $t\in T$, $y_l\in S_l$ $(l\notin Q)$ and points $y_0\in S_0$, $y_q\in S_q$ within $r$ of $t$, some pair of the family $y_0,(y_q),(y_l)$ is at least $D$ apart, it is no pair among $y_0,(y_q)$, and so two of the $y_l$, or one $y_l$ and $t$, are at least $D-r$ apart.
\end{proof}

\begin{lemma}[The chain with an entry; the cubic insertion]\label{C-360-9}
Let $\gamma^{-2}\le\mathsf b_-\le\mathsf b_+$, $MR(\mathsf b_+)\le L^m$, $\ell=2L^m$, $0<\tau\le\frac14$, $0<\phi<\frac1{42}$, $u=\lceil\phi K\rceil$, $t=K-u$ and $c'_0=\min\{\frac1{14}-\frac\tau7,\frac1{14}-3\phi\}$. Let $K_\phi$ be the least integer $K\ge2$ with $(1-\phi)K\ge2$, $L^{\phi K}\gamma^{-2}\ge\mathsf t_\natural$ and $\phi K>16\mathfrak p_0^2c_{\log}/(\log L)^2$, $\mathsf t_\natural$ being the calibration of Section~\ref{sec:cutoff-calibration} and $c_{\log}=\log L/\log\gamma^{-2}$; it depends on the constants and on $\phi$ alone. For $K\ge K_\phi$ and $\mathsf b'\in[\mathsf b_-,\mathsf b_+]$ let $\rho_0dU$ be the measure of $T_\ell$ at $\beta=L^K\mathsf b'$; the data $(K,\mathsf b',m)$ with the top level $t$ form the setting (Th) (Section~\ref{sec:cutoff-calibration}), and the small-field chain is formed along the couplings $x_j=L^{K-j}\mathsf b'$ up to the level $t$. Let $\mathcal M^k$ be the block map to the level $k$, ${\rm Cut}_k$ the field cut and ${\rm In}_k\subset{\rm Cut}_k$ the inner ball of the level $k$ (the definitions preceding Lemma~\YMMref{113.10}{796}), $\mathfrak G_k[x_k]\ge0$ the gauge-fixing factor and $\chi^{\rm sf}_{x_k}\in\{0,1\}$ the fluctuation test of the step from the level $k$ (Section~\ref{sec:sf-chart-interface}), and
\[
\mathsf p_k(U)=\mathbf 1_{{\rm Cut}_k}(\mathcal M^kU)\prod_{j<k}\bigl[\mathbf 1_{{\rm Cut}_j}\,\mathfrak G_j[x_j]\,\chi^{\rm sf}_{x_j}\bigr](\mathcal M^jU),
\]
the weight of the small-field chain up to the level $k$: a bounded Borel function with values in $[0,\infty)$, with $\mathsf p_0=\mathbf 1_{{\rm Cut}_0}$ and $\mathsf p_{k+1}=(\mathbf 1_{{\rm Cut}_{k+1}}\circ\mathcal M^{k+1})\,\bigl((\mathfrak G_k[x_k]\chi^{\rm sf}_{x_k})\circ\mathcal M^k\bigr)\,\mathsf p_k$; and $(\mathbf 1_{{\rm In}_t}\circ\mathcal M^t)\,\mathsf p_t$ is the weight $\mathfrak w^{\rm sf}_t$ of the small-field chain of the definitions preceding Proposition~\YMMref{115.2}{806}. Let $\nu_k=(\mathcal M^k)_*(\mathsf p_k\rho_0dU)$, the measure of the level $k$ of the chain among the construction results at the source $0$ (Proposition~\YMMref{114.2}{801} and the proof of Proposition~\textup{\YMMref{115.2}{806}}), a finite positive measure with $\nu_{k+1}=\mathbf 1_{{\rm Cut}_{k+1}}\mathbf T(\mathfrak G_k[x_k]\chi^{\rm sf}_{x_k}\nu_k)$, $\mathbf T$ the image under the block averaging; let $Q_kh$ be the density of $\mathbf 1_{{\rm Cut}_{k+1}}\mathbf T(\mathfrak G_k[x_k]\chi^{\rm sf}_{x_k}h\,\nu_k)$ with respect to $\nu_{k+1}$, the mean of $h$ over the range ${\rm Rg}_{x_k}$ of the chart of the fibre of the step (Lemma~\ref{lem:sf-chart-integral}), and $Q_{t\leftarrow0}=Q_{t-1}\cdots Q_0$, so that $(\mathcal M^t)_*(\mathsf p_tY\rho_0dU)=(Q_{t\leftarrow0}Y)\,\nu_t$ for bounded Borel $Y$: $(Q_{t\leftarrow0}Y)\circ\mathcal M^t$ is the conditional mean of $Y$ given $\mathcal M^tU$ under $\mathsf p_t\rho_0dU$.

An \emph{entry of the type} $(m_J,c,\hat c,N)$ is a finite family $J=(\mathsf Q_\iota,j_\iota,\varkappa_\iota)_\iota$ of boxes $\mathsf Q_\iota=\{x+a\sum_\mu n_\mu e_\mu:0\le n_\mu\le r_\mu\}$, $(r_1,r_2,r_3)\in\{0,1,2\}^3$ with at least two $r_\mu\ge1$, $x$ the lowest site of the box, of densities $j_\iota$ on them and of complex weights $\varkappa_\iota$, at most $N$ indices having one lowest site. A density on a box is a polynomial in the entries of the complex configurations $\mathbf U(b)\in{\rm SL}(2,\mathbb C)$ at its bonds, $\mathbf U(\bar b)=\mathbf U(b)^{-1}$, invariant under the gauge transformations of the box; it has the height $\hat c$ if $|j|\le\hat c$ wherever $|\mathbf U(b)|\le2$ at all the bonds, and the order $m_J$ with the constant $c$ if $|j|\le c\,\epsilon^{m_J}$ wherever, besides, $|\partial\mathbf U(q)-\mathbf 1|\le\epsilon$ at all the plaquettes $q$ of the box, $0<\epsilon\le1$.

Zero insertions and zero scalar sources are omitted before inverse norms and source radii are formed; the multilinear estimates are immediate when an insertion vanishes. We write $J(U)=\sum_\iota\varkappa_\iota j_\iota(U|_{\mathsf Q_\iota})$ and $|\varkappa|_\infty=\max_\iota|\varkappa_\iota|$; the carrier of $J$ is the set of the sites of the boxes with $\varkappa_\iota\ne0$.
\begin{enumerate}
\item[(a)] For $K\ge K_\phi$, an entry of the type $(m_J,c,\hat c,N)$ with $m_J\ge2$, $|s|<r_J$ with $r_J^{-1}\le C|\varkappa|_\infty$, and $1\le k\le t$:
\[
(\mathcal M^k)_*\bigl(\mathsf p_ke^{sJ}\rho_0dU\bigr)=\exp\Bigl[\sum_{i\le k}c_i(s)\Bigr]\mathbf 1_{{\rm Cut}_k}e^{\Phi^s_k}\,dV,
\]
the numbers $c_i(s)$ and the action $\Phi^s_k$ being holomorphic in $s$ and sums of local families, each of which depends on $\varkappa$ in its domain only. $Q_{t\leftarrow0}J$ is a number plus $\partial_s\Phi^s_t|_{s=0}$ $\nu_t$-almost everywhere on ${\rm Cut}_t$, and its oscillation is at most $CK^{2\mathfrak p_0}L^{4u-K}|\varkappa|_\infty$ for $m_J=2$ and $CK^{3\mathfrak p_0}L^{4u-\frac32K}|\varkappa|_\infty$ for $m_J\ge3$. Here $C$ depends on the constants, $\ell$, $\mathsf b_+$, $\phi$ and the type of the entry.
\item[(b)] There are $K_2\ge K_\phi$, $C_2$, $p_2$, depending on these data alone, such that for $K\ge K_2$ and every $\mathsf b'\in[\mathsf b_-,\mathsf b_+]$ the following holds. For every field $\Xi$ there is a number $\mathsf a$ with $|Q_{t\leftarrow0}(\beta\mathfrak K^\sharp[\Xi])-\mathsf a|\le C_2\mathsf H(\Xi)K^{p_2}L^{-c'_0K}$ $\nu_t$-almost everywhere on ${\rm Cut}_t$. The same holds at real sources: let $n\ge1$, $f_1,\dots,f_n$ real Lipschitz functions, $\mathsf M_t\le CK^{2\mathfrak p_0}\ell^3L^{3u}$ the inverse radius of the sources of the chain at the level $t$, and $x$ real in the polydisc $\mathbb D_t$ of the radii $r_i=(8n\mathsf M_t\|f_i\|_\sharp)^{-1}$; then $w_x=\sum_ix_if_i$ lies in the class of sources of the setting (Th), and with the measure $\rho_0e^{\mathcal O(w_x)}dU$ in the place of $\rho_0dU$, its chain measure $\nu^{w_x}_t$ and its conditional mean $Q^{w_x}_{t\leftarrow0}$, and the entry with the moved weight $\mathfrak K^\sharp[(\beta-a^{-1}w_x)\Xi]$ in the place of $\beta\mathfrak K^\sharp[\Xi]$, the same bound holds with a number $\mathsf a(x)$, the numbers $C_2,p_2,K_2$ being free of $x$ and of $\mathsf b'$.

The number $\mathsf a(x)$ is the sum, over the levels $j\le t$, of the values at the configuration $\mathbf 1$ of local families of the level $j$ whose member at a polymer $X$ reads $w_x$ and $\Xi$ on the plaquettes of $X$ alone and vanishes if $\Xi$ vanishes at the barycentres of these plaquettes. For free photons the cubic piece vanishes identically.
\end{enumerate}
\end{lemma}

\begin{proof}[Proof sketch]
(a) The entry is typed anew at each level. At a height $j\ge1$ the sum of the $\varkappa_\iota j_\iota$ over the boxes whose cubes of the level $j$ form a given polymer is a pointed family of the level $j$: at most $8N(ML^j)^3$ boxes have one polymer, and on the cores each density is at most $c\,\alpha_0^{m_J}L^{-2m_Jj}$ by its order, so that the family has the size $8NM^3c\,\alpha_0^{m_J}|\varkappa|_\infty L^{-(2m_J-3)j}$; this is the power counting, of order $|\varkappa|_\infty L^{(3-2m_J)j}$ in level-$j$ units. At the height $0$ the family in the fourth place $\sum_\iota\varkappa_\iota[j_\iota(e^{iB'}\mathbf U)-j_\iota(\mathbf U)]$, $B'$ the fluctuation field with its values at the pivot bonds, sums to $J(e^{iB'}\mathbf U)-J(\mathbf U)$ and has the size of a number times $N\hat c\,\delta(x_0)|\varkappa|_\infty$, by the height and Schwarz's lemma in the fluctuation field.

So for $|s|<r_J$, $r_J^{-1}\le C|\varkappa|_\infty$ with $C$ reading $N$, $c$ and $\hat c$, the entry enters each step as one more local term, inside the radii of the one-step theorem with sources (Proposition~\ref{prop:source-entry}); the members of the families are holomorphic in $s$ and read $s$ only on the polymers that the entry meets. The identity of the measures follows by induction on $k$ from the recursion of $\mathsf p_k$: the images $\nu^s_k=(\mathcal M^k)_*(\mathsf p_ke^{sJ}\rho_0dU)$ obey $\nu^s_{k+1}=\mathbf 1_{{\rm Cut}_{k+1}}\mathbf T(\mathfrak G_k[x_k]\chi^{\rm sf}_{x_k}\nu^s_k)$, the small-field pushforward in Lemma~\ref{lem:sf-chart-integral}, a chart integral over ${\rm Rg}_{x_k}$ that the one-step theorem evaluates; one function $e^{\Phi^s_k}$ is both the density reached at the level $k$ and the integrand of the step from the level $k$.

The derivative of the exponent in $s$ at $0$ is the conditional mean up to the number $\dot c_t=\sum_i\partial_sc_i(0)$. On ${\rm Cut}_t$ the plaquette variables of the background are within $\mathsf e^{\rm c}_tL^{-2t}$ of $\mathbf 1$, $\mathsf e^{\rm c}_t$ the radius of the field cut, so that the entry at the background is at most $N(\ell L^K)^3c\,(\mathsf e^{\rm c}_t)^{m_J}L^{-2m_Jt}|\varkappa|_\infty$; with what the families of the earlier levels carry, $|\Phi^s_t-\Phi^{s'}_t|\le|s-s'|\,\mathsf o_t$ with $\mathsf o_t$ of the order of the stated oscillation, which is at most $2\mathsf o_t$. The threshold $K\ge K_\phi$ makes $t$ a top level of the setting (Th): $(1-\phi)K\ge2$ gives $1\le t\le K-1$, $L^{\phi K}\gamma^{-2}\ge\mathsf t_\natural$ gives $x_t=L^u\mathsf b'\ge\mathsf t_\natural$, and the third condition keeps the size of a source at the coupling $x_k$, $k\le t$, below the bound of the step; this size $\theta_k$ obeys $L\theta_k\le\Theta_KL^{-(K-k)/2}$, $\Theta_K$ a power of $K$. (b) Let $S=\beta\mathfrak K^\sharp[\Xi]=\sum_p\varkappa(p)\mathfrak r^\sharp_{\lambda(p)}(p)$, $\varkappa(p)=L^{2K}\mathsf b'\,\Xi^{\lambda(p)}(x(p))$, so that $\max|\varkappa|\le L^{2K}\mathsf b_+\mathsf H(\Xi)$; the density vanishes to the third order at $\mathbf 1$ (Lemma~\ref{C-360-8}(b)).

At a real source the weight is moved, $\varkappa(p)=L^{2K}(\mathsf b'-w_x(x(p)))\Xi^{\lambda(p)}(x(p))$, and the chain carries the source; we write the source $0$, the other case being the same with the index $w_x$. Let $s_0=S$ and $s_j=S\circ U_j$ on ${\rm Cut}_j$, the entry read at the background of the level $j$. For a pointed family $\mathsf f$ of a level $\le n$ let $\Sigma_n[\mathsf f](V)$, $V\in{\rm Cut}_n$, be the sum over the polymers of the members of $\mathsf f$ at the background $U_n(V)$, and ${\rm Sum}[\mathsf f](\mathbf 1)$ the same sum at the configuration $\mathbf 1$, a number. Let $S_n\mathsf f$ be the first-order response of the step from the level $n$ to a family $\mathsf f$ of a level $i\le n$, a family of the level $n+1$ with $\|S_n\mathsf f\|_{n+1}\le\theta_nL^{-(n-i)}\|\mathsf f\|_i$, so that $Q_n\Sigma_n[\mathsf f]=\Sigma_{n+1}[\mathsf f+S_n\mathsf f]$ on ${\rm Cut}_{n+1}$; and let $D^{(j)}=D^{(j)}[w;\varkappa]$ be the family born at the level $j$ from the entry at the source $w$ and the weight $\varkappa$: $Q_{j-1}s_{j-1}-s_j=\Sigma_j[D^{(j)}]$ on ${\rm Cut}_j$.

With $T_t={\rm id}$ and $T_n\mathsf g=T_{n+1}(\mathsf g+S_n\mathsf g)$, the carriage of a family from the level $n$ to the level $t$ with all that it feeds, $Q_{t-1}\cdots Q_n\Sigma_n[\mathsf g]=\Sigma_t[T_n\mathsf g]$; so, with $n_0=K-\lceil\frac K7\rceil$ and $Q_{t\leftarrow n_0}=Q_{t-1}\cdots Q_{n_0}$,
\[
Q_{t\leftarrow0}S=Q_{t\leftarrow n_0}s_{n_0}+\sum_{j\le n_0}\Sigma_t[T_jD^{(j)}]\qquad\text{on }{\rm Cut}_t .
\]
The births are small by the parity in the fluctuation field: the part of the variation of the density that is odd in the fluctuation field has the size $g_{j-1}=x_{j-1}^{-1/2}$ and the Gaussian mean $0$, so that it meets only the odd part of the interaction, of the size $g_{j-1}(\log x_{j-1})^{\mathfrak p_0}$, while the even part has the size $g_{j-1}^2$; so $\|D^{(j)}\|\le C(\log x_{j-1})^{\mathfrak p_0}g_{j-1}^2L^{-3(j-1)}\max|\varkappa|$ once $K-j$ exceeds a fixed construction threshold, and $L^j\|D^{(j)}\|\le CK^{\mathfrak p_0}L^{K-j}\mathsf H(\Xi)$. $T_jD^{(j)}$ is $D^{(j)}$ plus the sum of the chains $S_{k_m}\cdots S_{k_1}D^{(j)}$, $j\le k_1<\dots<k_m<t$, of the families that it feeds; each chain is stopped at a level $n$ of its own, chosen with the free fraction $\tau$, at which the variation ${\rm Sum}[\mathsf f](U_n(V))-{\rm Sum}[\mathsf f](\mathbf 1)$ of a family $\mathsf f$ of the level $i$ is at most $\beta_nL^i\|\mathsf f\|_i$, $\beta_n=CK^pL^{3(K-n)-K}$, the square of the radius of regularity of the level $n$ entering. $Q_{t\leftarrow n}$ does not enlarge an oscillation, and each feeding costs $L\theta_k\le\Theta_KL^{-(K-k)/2}$, so that what is fed is forgotten.

The variation of ${\rm Sum}[D^{(j)}]$ over ${\rm Cut}_{n_0}$, $j<n_0$, is where the parity of the cubic density enters. The step is covariant under the isometries of the lattice of its level, and $\mathfrak r^\sharp$ transforms as a vector (Lemma~\ref{C-360-8}(b)); so under the inversion $r_z:y\mapsto2z-y$ through a site $z$ of the lattice of the level $j$, ${\rm Sum}[D^{(j)}[w;\varkappa]](r_zU)=-{\rm Sum}[D^{(j)}[r_z^*w;r_z^*\varkappa]](U)$, $(r_z^*v)(p)=v(r_zp)$. With a partition of unity $\psi_z\ge0$ over these sites, $\psi_z\circ r_z=\psi_z$, $\psi_z(p)=0$ unless $|x(p)-z|_\infty\le\frac12L^{j-K}$, and $\Sigma_z[w;\varkappa]={\rm Sum}[D^{(j)}[w;\psi_z\varkappa]]$, this gives the exact split
\[
\begin{multlined}
\Sigma_z[w;\varkappa](U)=\tfrac12\bigl(\Sigma_z[w;\varkappa](U)-\Sigma_z[w;\varkappa](r_zU)\bigr)+\tfrac12\Sigma_z[w;\varkappa-r_z^*\varkappa](U)+\tfrac12\bigl(\Sigma_z[w;r_z^*\varkappa](U) \\
-\Sigma_z[r_z^*w;r_z^*\varkappa](U)\bigr).
\end{multlined}
\]
(I) The first part compares the background with its image under $r_z$. In the H\"older gauge (Theorem~\textup{\ref{I-56-65}} for real backgrounds) the curvature, a $2$-form, is even under the inversion, so that the linear part of the potential is kept and the potential and its image differ by the odd part of the H\"older remainder, at most $R^2\vartheta^{2+b}e_3$ on a cube of the radius $R$ about $z$, $\vartheta=L^{-(n_0-j)}$, $b=\frac34$, $e_3$ of the order of the radius of regularity of the level $n_0$: a gain $\vartheta^b$. (II) The second part carries the asymmetry of the weight, $|\varkappa-r_z^*\varkappa|\le L^{2K}{\rm Lip}(\Xi)L^{j-K}$ where $\psi_z\ne0$. (III) The third part is the asymmetry of the source in the step: it vanishes at the source $0$, and at a real source it is of the second order at the trivial triple, of the size ${\rm Lip}(w_x)\,\delta(x_{j-1})$, $\delta(x_{j-1})$ of the order of $L^{-(K-j)/2}$ being the threshold of the level.

Summed over $j<n_0$, the parts (I), (II) and (III) are $O(L^{-3K/14})$, $O(L^{-4K/7})$ and $O(L^{-K/14})$; the stopped entry contributes $O(L^{-5K/14})$ and the family born at $n_0$ $O(L^{-3K/7})$; forgetting leaves $O(L^{-(\frac1{14}-\frac\tau7)K})$ of what is fed, and against the factor $L^{3u}$ of the levels above $t$ it leaves $O(L^{-(\frac1{14}-3\phi)K})$; all this up to powers of $K$. The least of these exponents is $c'_0$; at the source $0$ the last two meet at $\frac1{14}$ as $\tau,\phi\to0$, and at a real source the part (III) meets them there. The number $\mathsf a$, the sum of the values at the configuration $\mathbf 1$ of the families born and fed up to their levels of stopping, depends on no $V$, and each of these families reads the source and the weight on its polymer alone and vanishes where the weight vanishes.

The numbers are uniform on $[\mathsf b_-,\mathsf b_+]$ because the frame of Definition~\ref{def:running-source-history} holds for the chain along $x_j=L^{K-j}\mathsf b'$ at every label of the interval and $\mathsf b'$ enters only through $x_i\in[\gamma^{-2}L^{K-i},\mathsf b_+L^{K-i}]$ and $\max|\varkappa|$, each time in a monotone way; on the polydisc $\mathbb D_t$ the source stays in the class of the setting (Th) and the mass of the small-field chain stays away from $0$, and at real $x$ the measures of the chain are positive.
\end{proof}

\begin{lemma}[The cubic remainder]\label{C-360-10}
In the setting of Lemma~\ref{C-360-9}, let $n\ge1$, $\mathsf d>0$ and $c<\frac1{14}$. There are $C,p$, free of $K$ and of $\mathsf b'$, such that for all $K\ge2$, $\mathsf b'\in[\mathsf b_-,\mathsf b_+]$, $\beta=L^K\mathsf b'$, all real Lipschitz $f_i$ and all fields $\eta$, the spread of $\{\eta\ne0\},{\rm supp}\,f_1,\dots,{\rm supp}\,f_n$ being $\ge\mathsf d$,
\[
|\mathfrak U^\sharp[\eta](f)|\le CK^pL^{-cK}\mathsf H(\eta)\prod_i\|f_i\|_\sharp .
\]
The spread condition holds if two of the sets are at sup distance $\ge\mathsf d$, and if $\eta=0$ within the sup distance $\mathsf d$ of one support. By $x\mapsto L^zx$ the same holds at every label and on every side.
\end{lemma}

\begin{proof}[Proof sketch]
Moving the weight of the entry with the sources produces the contact terms:
\[
\mathfrak U^\sharp[\eta](f)=-\partial_s\partial_{z_1}\cdots\partial_{z_n}\log\Bigl\langle\exp\Bigl(s\,\mathfrak K^\sharp\bigl[(\beta-a^{-1}\textstyle\sum_iz_if_i)\eta\bigr]+\sum_iz_i\mathcal O(f_i)\Bigr)\Bigr\rangle\Big|_{s=z=0}.
\]
After discarding the histories with large fields, Lemma~\ref{C-360-9}(b), applied to the chain with the real sources $x$ in its polydisc $\mathbb D_t$ and to the field $\eta$, makes the conditional mean of the entry with the moved weight equal, up to $\varepsilon=K^pL^{-c'_0K}$, to its number $\mathsf a(x)$; and $\mathsf a(x)$ is a sum $\sum_ic_i(x)$ of functions each of which does not depend on $x_i$, up to a remainder below every power of $L^{-K}$. Indeed, let $c_i(x)$ collect the members of the families of $\mathsf a(x)$ at the polymers on whose plaquettes $f_i$ is the first of the sources that vanishes at all the barycentres; such a member reads the sources on its polymer alone, so it does not read $x_i$; the members at the other polymers, which meet $\{\eta\ne0\}$ and all the ${\rm supp}\,f_i$, are left over; under the spread such a polymer of the level $j$ has the linear size $\ge\mathsf dL^{K-j}/M-3$, so that these members cost $e^{-\mathsf dL^{K-j}/M}$ against the sizes.

Such sums have no mixed derivative. At complex sources the mean is crudely at most $B=C\ell^3L^{3K}\max|\varkappa|$, $\varkappa$ the weight of the entry. Interpolating between the two bounds, the mixed derivative is at most $2^n\varepsilon^{1-\sigma}B^\sigma(\frac{4n}{\pi\sigma})^n\prod_ir_i^{-1}$ with $r_i^{-1}\le CK^{2\mathfrak p_0}\ell^3L^{3u}\|f_i\|_\sharp$, and $\sigma,\phi,\tau$ small give every $c<\frac1{14}$. The finitely many $K<K_2$ are paid by the crude bound of the cumulant by the suprema of its entries, which goes into $C$. For $n\ge2$ the bound holds also where $\{\eta\ne0\}$ covers the supports, as in the uses below, since the spread of the $n+1$ sets is at least that of the supports among themselves.
\end{proof}

\begin{lemma}[Chessboard bound in one cell]\label{C-360-11}
Let $0<g\le\gamma$, $\ell'=2\mathsf l\in\mathfrak L(g)$, $0<\phi\le\frac1{84}$, $u=\lceil\phi K\rceil$, $t=K-u$, $0<\delta\le\frac18\mathsf l$ and $m\ge1$. There are $C,p,K_3$, depending on $g,\ell',\phi,m,\delta$, such that for $(K,\beta)\in\mathfrak D(g)$ with $K\ge K_3$, every board of side $\mathsf l$ with offset in $a_{K-t}\mathbb Z^3=L^{-u}\mathbb Z^3$, every cell, every $\mathsf d$ with $\mathsf dL^u\ge K^2$, and $m$ slots $X_1,\dots,X_m$ supported in the cell at sup distance $\ge\delta$ from its walls and pairwise $\ge\mathsf d$ apart,
\[
\Bigl\|\prod_a\bigl(X_a-\langle X_a\rangle_{K,\ell',\beta}\bigr)\Bigr\|_{\rm ch}\le C\prod_a\varpi(X_a)\,\mathsf n(X_a),
\]
where a slot is $\mathcal T_K(H)$, with $\varpi=K^pL^{4u}$ and $\mathsf n=|H|_\infty$, or one of 
\[
\mathcal T_K(H^{\rm cons}[\eta]),\mathrm{St}_K[\eta],\Delta_K[\eta],\mathrm{Rm}_K[\eta],\mathfrak K^{\rm v}[\beta\eta]
,\]
with $\varpi=K^pL^{3u-K/28}$ and $\mathsf n=\|\eta\|_{C^2}$. For $m=1$ and $\phi=\frac1{400}$ we have $\|\mathrm{St}_K[\eta]-\langle\cdot\rangle\|_{\rm ch}\le CK^pL^{-K/28}\|\eta\|_{C^2}$ and 
\[
\|\Delta_K[\eta]-\langle\cdot\rangle\|_{\rm ch}+\|\mathcal T_K(H^{\rm cons}[\eta])-\langle\cdot\rangle\|_{\rm ch}\le CK^pL^{-K/36}\|\eta\|_{C^2}
;\]
for a plain tensor entry $\|\mathcal T_K(H')-\langle\cdot\rangle\|_{\rm ch}\le CK^{2\mathfrak p_0}L^{K/100}|H'|_\infty$. At $m=1$ and $\phi=\frac1{400}$, moreover, $\|\mathrm{Rm}_K[\eta]-\langle\cdot\rangle\|_{\rm ch}=\|\mathrm{St}_K[\eta]-\langle\cdot\rangle\|_{\rm ch}$, and the bound $CK^pL^{-K/28}\|\eta\|_{C^2}$ holds also for $\|F_\eta\|_{\rm ch}$, where
\[
F_\eta=\mathfrak K^{\rm v}[\beta\eta]-\tfrac14\mathfrak N[\mathsf k_\eta]=\beta\mathfrak K^\sharp[\eta]-\sum_{v=1}^5\mathfrak J_v[\mathsf k_\eta]+\tfrac14c_{\mathsf k_\eta},\qquad\langle F_\eta\rangle_{K,\ell',\beta}=0,
\]
with $\mathsf k^\rho_\eta=\nabla^+_\rho\eta^\rho$ and the objects of Lemma~\ref{C-360-8}(c); for every state $s$ of Lemma~\ref{C-360-7}, $s(\mathfrak K^{\rm v}[\beta\eta])=s(F_\eta)$, so that $|s(\mathfrak K^{\rm v}[\beta\eta])|\le e^{G_g\mathsf l^3}\|F_\eta\|_{\rm ch}$ by Lemma~\ref{C-360-7} at $q=1$, $c_1=0$.
\end{lemma}

\begin{proof}[Proof sketch]
By Lemma~\ref{C-360-8}(c), $\mathrm{St}_K[\eta]=N_\beta[\eta]-\frac14\mathfrak N[\mathsf k_\eta]-\mathrm{Rm}_K[\eta]-\frac14c_{\mathsf k_\eta}$, so that the slots $\mathrm{St}_K[\eta]$ and $-\mathrm{Rm}_K[\eta]$ differ by null terms and a number; inside the mean of the $8m$ mirror copies on $T_{\ell'}$ the null terms drop, since no other copy depends on their bonds. A product of centred slots in one cell expands into joint cumulants over the partitions into blocks of at least two members. Apply Proposition~\ref{stc:cum}(a), with its common insertion radius $r_1$, size $\bar{\mathsf d}$ and polymer sum $\mathsf t_\times$. It bounds the higher mixed jets of the exponent as well as the separate conditional means and gives
\[
|\langle J_0;\dots;J_r\rangle|\le{\rm op}_{r+1}\prod_a\omega_t(J_a)+9\cdot2^{r+1}(r+2)\bar{\mathsf d}\,\mathsf t_\times r_1^{-r}+4\,{\rm op}'_{r+1}\mathsf R_t\prod_a\sup|J_a|,\qquad{\rm op}_n=\sum_{\pi}(|\pi|-1)!,
\]
$\pi$ over the partitions of $n$ points, with ${\rm op}'_n$ as in Proposition~\ref{stc:cum}; $\omega_t$ is the oscillation of the conditional mean (Lemma~\ref{C-360-9}) and $\mathsf R_t$ the weight of the histories with large fields. After allowing the carrier collars, take $\mathsf d_\times\ge c\min\{\mathsf d,\delta\}$. The hypothesis $\mathsf dL^u\ge K^2$ gives $\mathsf d_\times L^u\ge cK^2$ for all sufficiently large $K$, uniformly in $\mathsf d$. The prefactor estimates in the proof of Proposition~\ref{stc:prod} bound the second term by a polynomial in $K$ times $L^{AK}e^{-cK^2}$; thus the condition of Proposition~\ref{stc:cum}(a) holds and this term is below any prescribed $L^{-BK}$. The large-field term is controlled there separately. Six of the seven pieces of $\mathrm{Rm}_K$ are small by power counting, the seventh is the cubic one, of oscillation $K^pL^{-c'_0K}$. The principal part of $\Delta_K$ is null (Lemma~\ref{C-360-8}(d)): on the backgrounds of the level $K-\lceil\frac K{14}\rceil$ its value is at most $CK^{2\mathfrak p_0}L^{\frac72\lceil K/14\rceil-\frac K2}\mathsf b\sup|\nabla\eta|$, what is born below that level is born at a small coupling and is forgotten, and a conditional mean is never larger than a supremum.

Each crude tensor entry costs $L^{4u}$. At $m=1$: by (a)--(c) of Lemma~\ref{C-360-8}, $\mathfrak K^{\rm v}[\beta\eta]=\frac14\mathfrak N[\mathsf k_\eta]+\beta\mathfrak K^\sharp[\eta]-\sum_v\mathfrak J_v[\mathsf k_\eta]+\frac14c_{\mathsf k_\eta}$, so that $F_\eta$ consists of six of the seven pieces of $\mathrm{Rm}_K[\eta]$ and a number. The cubic piece has the oscillation $K^pL^{-K/28}\|\eta\|_{C^2}$ by Lemma~\ref{C-360-9}(b) at $\tau=\frac14$, where $c'_0=\frac1{28}$ for $\phi\le\frac1{84}$; the five pieces $\mathfrak J_v$, of the orders $2,2,2$ with weights of the order of $\|\eta\|_{C^2}$ and of the orders $3,4$ with weights at most $\beta\|\eta\|_{C^2}$, have oscillations at most $K^{3\mathfrak p_0}L^{4u-K/2}\|\eta\|_{C^2}$ by Lemma~\ref{C-360-9}(a); so no factor $L^{3u}$ enters the bound of $F_\eta$; in the bounds of the slots of the kind of $\eta$ it enters through $\Delta_K$ alone.

The means of $\mathfrak K^{\rm v}[\beta\eta]$ and of $\mathfrak N[\mathsf k_\eta]$ on $T_{\ell'}$ are $0$ (Lemma~\ref{C-360-8}(a), (b)), so is that of $F_\eta$, and $s(\mathfrak N[\mathsf k_\eta])=0$ in every state $s$ by Lemma~\ref{C-360-8}(a).
\end{proof}

\begin{lemma}[From sup norms to an integrable density]\label{C-360-12}
Let $H\ge0$ be self-adjoint, $T_t=e^{-tH}$, and for every $\eta>0$ let $\rho\mapsto S_\eta[\rho]$ be linear from $C_c((-\eta,\eta))$ to the bounded operators, with $\|S_\eta[\rho]\|\le A_J\sup|\rho|$ for ${\rm supp}\,\rho$ in a compact $J\subset(-\eta,\eta)$, and $S_\eta[\rho(\cdot-c)]=T_{\delta+c}S_{\eta-\delta}[\rho]T_{\delta-c}$ for $|c|\le\delta<\eta$. Then there is a unique family of bounded operators $B(s,s')$, $s,s'>0$, with $T_tB(s,s')T_{t'}=B(s+t,s'+t')$ and $S_\eta[\rho]=\int\rho(c)B(\eta+c,\eta-c)\,dc$; for $0<\gamma'<\eta$ and ${\rm supp}\,\rho\subset[-(\eta-\gamma'),\eta-\gamma']$, $\|S_\eta[\rho]\|\le\|S_{\gamma'/2}[{\rm tent}_{\gamma'/8}]\|\,\|\rho\|_{L^1}$; and $B(s,s')^*=B(s',s)$ if $S_\eta[\rho]^*=S_\eta[\bar\rho(-\cdot)]$. The same holds in the space directions for commuting momenta $P$ with $H\ge|P_k|$ in the direction $k$.
\end{lemma}

\begin{proof}[Proof sketch]
For $G$ holomorphic and bounded on the strip $\{|{\rm Re}\,z|<h'\}$, $h'>2h$, integrating $\Gamma(z)\sin^{-2}(\pi z/2h)$, $\Gamma'=G$, over the boundary of the strip gives
\[
G(0)=\int\mu_2(ds)\int{\rm tent}_{2h}(c)\,G(c+is)\,dc,\qquad\mu_2=\mu_1*\mu_1,\quad\mu_1(ds)=\tfrac\pi{4h}{\rm sech}^2\bigl(\tfrac{\pi s}{2h}\bigr)ds .
\]
Thus an average in real time inverts a tent smearing in Euclidean time. Apply it to $z\mapsto\langle T_{\delta+\bar z}\psi_2,S[\rho]T_{\delta-z}\psi_1\rangle$ for vectors $\psi_1,\psi_2$, holomorphic by the covariance, and use $\rho_0*\rho_1=\rho_1*\rho_0$: the $\mu_2$-average of $e^{-isH}S_{\gamma'/2}[{\rm tent}]e^{isH}$ is the operator $\mathbb B$ of the density, $S_\eta[\rho]=\int\rho(c)T_{\eta-\gamma'/2+c}\mathbb BT_{\eta-\gamma'/2-c}\,dc$. In space, $G(0)=\int h_{d'}(x)\frac12[G(x+id')+G(x-id')]\,dx$ with $h_{d'}=(2d')^{-1}{\rm sech}(\pi x/2d')$, with the imaginary translations that $H\ge|P|$ allows.
\end{proof}

\begin{lemma}[Functions of the block field]\label{C-360-17}
For a history $h$ let $\mathfrak w_h\ge0$ be its weight in the law of the field and of the history and $\mathsf t(h)=\int\mathfrak w_h\rho_0\,dU$ its mass, so that $\sum_h\mathsf t(h)=Z$ (equations~\YMMref{(646)--(648)}{342} and Proposition~\YMMref{44.1}{342}).
\begin{enumerate}
\item[(i)] In the setting of Lemma~\ref{C-360-9}, let $1\le t\le K-1$ be a level up to which the chain is formed (a top level of the chain at the label $\mathsf b'$), $h_t$ the history without large-field regions of the level $t$, ${\rm In}_t=\{V:\varpi(V)<2\varepsilon(x_t)\}\subset{\rm Cut}_t$ the inner ball of the level $t$ (the definitions preceding Lemma~\YMMref{113.10}{796}; $\varpi(V)$ is the largest $|\partial V(p')-\mathbf 1|$ over the plaquettes $p'$ of the level $t$), and $\mathfrak w^{\rm sf}_t$ the weight of the level $t$ of the small-field chain with its number $\mathsf R_t$ (the definitions preceding Proposition~\YMMref{115.2}{806}, the definition preceding Proposition~\YMMref{115.2}{806}, \YMMref{(2187)}{806}): $\mathfrak w^{\rm sf}_t(U)=\mathbf 1_{{\rm In}_t}(\mathcal M^tU)\prod_{i<t}[\mathbf 1_{{\rm Cut}_i}\mathfrak G_i[x_i]\chi^{\rm sf}_{x_i}](\mathcal M^iU)$, with the gauge-fixing factors $\mathfrak G_i[x_i]$ and the fluctuation tests $\chi^{\rm sf}_{x_i}$ among the construction results, so that $\mathfrak w^{\rm sf}_t=(\mathbf 1_{{\rm In}_t}\circ\mathcal M^t)\,\mathsf p_t$ with $\mathsf p_t(U):=\mathbf 1_{{\rm Cut}_t}(\mathcal M^tU)\prod_{i<t}[\mathbf 1_{{\rm Cut}_i}\mathfrak G_i[x_i]\chi^{\rm sf}_{x_i}](\mathcal M^iU)$.

The measure $\nu_t$ of the level $t$ of the chain is $(\mathcal M^t)_*(\mathsf p_t\rho_0\,dU)$ (Proposition~\YMMref{114.2}{801}, the measure recursion in Proposition~\YMMref{114.2}{801}), and $Q_{t\leftarrow0}$ is the conditional mean given $\mathcal M^tU$ under $\mathsf p_t\rho_0\,dU$; the tower of Lemma~\ref{C-360-9} is read with this $\mathsf p_t$. By Proposition~\textup{\YMMref{115.2}{806}}, $\mathfrak w^{\rm sf}_t$ is a bounded Borel function with values in $[0,\infty)$, $(\mathcal M^t)_*(\mathfrak w^{\rm sf}_t\rho_0\,dU)=\mathbf 1_{{\rm In}_t}\nu_t$ \YMMref{(2188)}{806}, $0\le\mathfrak w_{h_t}\le\mathfrak w^{\rm sf}_t$ at every point \YMMref{(2189)}{806}, and $0\le Z-\int\mathfrak w^{\rm sf}_t\rho_0\,dU\le Z-\mathsf t(h_t)\le\mathsf R_tZ$ \YMMref{(2190)}{806}: $\mathsf R_t$ bounds the weight of the histories with a large field up to the level $t$.

For every bounded real Borel gauge invariant $Y$, real $a_0$ and $\epsilon\ge0$ with $|Q_{t\leftarrow0}Y-a_0|\le\epsilon$ $\nu_t$-almost everywhere on ${\rm In}_t$, and every bounded Borel function $\psi$ of the block fields of the level $t$, $|\langle Y;\psi\circ\mathcal M^t\rangle|\le2[\epsilon+2(\sup|Y|+|a_0|)\mathsf R_t]\sup|\psi|$; if besides $\mathsf R_t\le\frac12$, then $|a_0|\le\sup|Y|+\epsilon$.
\item[(ii)] In the setting of Lemma~\ref{C-360-24} at $0<g\le\gamma$ with $g^{-2}\ge x_{\rm J}$ (a first-passage datum $(K,\beta)\in\mathfrak D(g)$, a side $\ell\in\mathfrak L(g)$, a depth $s$ and the working level $k$), let $2\le T\le k$ be a level, $h_T$ the history without large-field regions of the level $T$, $\mathsf m_T=\tau_{h_T}[\rho_0]$ (Lemma~\ref{C-360-24}(a)), and $W_T<1$ the weight of a first creation of a large field at a level $\le T$, so that $Z-\mathsf t(h_T)=W_TZ$. A dominated local source $F_2$ acts on a dominated local source $F_1$ at the sup distance $\mathsf d'$, with $L^{K-T}\mathsf d'\ge3M$ and $\mathsf h_2(K-T)\le1$ (the conditions of Lemma~\ref{C-360-24}(e) at $t=T$), through a bounded function of the block field of the level $T$: let $r_2$ be the real part of the bounded density $D_{T,2}$ of $\tau_{h_T}[F_2\rho_0]$ with respect to $\mathsf m_T$, $c$ real, $E$ the essential supremum of $|r_2-c|$ for $\mathsf m_T$, and $\psi=(r_2-c)\mathbf 1\{|r_2-c|\le E\}$, a bounded real Borel function with $\sup|\psi|\le E$ and $\psi=r_2-c$ $\mathsf m_T$-almost everywhere; then
\[
|\langle F_1;F_2\rangle-\langle F_1;\psi\circ\mathcal M^T\rangle|\le C_2e^{3-\mathsf D_T/M}y_1y_2+2\sup|F_1|\bigl(\sup|F_2|+\sup|\psi|+|c|\bigr)W_T,
\]
with $y_a\le4\ell^3L^{4(K-T)}$ times the sizes of the sources and $\mathsf D_T\ge L^{K-T}\mathsf d'$ the least sup distance of the carriers in level-$T$ units; at $c={\rm Re}\,c_2$, $c_2$ the constant of $F_2$ at the level $T$, $E\le C_1y_2$ and $|c|\le\sup|F_2|+C_1y_2$. The display is additive in $F_1$.
\end{enumerate}
\end{lemma}

\begin{proof}[Proof sketch]
The law of the field and of the history, Proposition~\YMMref{44.1}{342} with the measures in equation~\YMMref{(647)}{342}, gives \[
\begin{aligned}
 \mathbf T^k(\Psi\rho_0\,dU)&=\sum_h\tau_h[\Psi\rho_0],\\
 \int F\,d\tau_h[\Psi\rho_0]
 &=\int F(\mathcal M^kU)\,\mathfrak w_h(U)\Psi(U)\rho_0(U)\,dU
\end{aligned}
\] for every bounded gauge invariant $\Psi$, when $\varepsilon_j\le\varepsilon_\flat=c_2^{\rm SO}$ for $j<k$; hence $Z\langle\Psi\cdot\psi\circ\mathcal M^k\rangle=\sum_{h}\int\psi\,d\tau_h[\Psi\rho_0]$, and, as $\mathfrak w_h\ge0$ and $\sum_h\mathsf t(h)=Z$, \[
 \left|Z\langle\Psi\cdot\psi\circ\mathcal M^k\rangle-
 \int\psi\,d\tau_{h_k}[\Psi\rho_0]\right|
 \le\sup|\Psi|\sup|\psi|\,(Z-\mathsf t(h_k)).
\] (i) At the level $t$ replace $\mathfrak w_{h_t}$ by $\mathfrak w^{\rm sf}_t$: the difference $\mathfrak w^{\rm sf}_t-\mathfrak w_{h_t}$ is $\ge0$ at every point, of the mass $\int\mathfrak w^{\rm sf}_t\rho_0\,dU-\mathsf t(h_t)\le Z-\mathsf t(h_t)\le\mathsf R_tZ$; with the histories other than $h_t$ this gives $|Z\langle\Psi\cdot\psi\circ\mathcal M^t\rangle-\int\Psi\,\psi(\mathcal M^tU)\,\mathfrak w^{\rm sf}_t\rho_0\,dU|\le2\sup|\Psi|\sup|\psi|\mathsf R_tZ$, the factor $2$ counting the two parts.

By the tower property, $\mathfrak w^{\rm sf}_t=(\mathbf 1_{{\rm In}_t}\circ\mathcal M^t)\mathsf p_t$ and ${\rm In}_t\subset{\rm Cut}_t$, $\int\psi(\mathcal M^tU)\,(Y-a_0)\,\mathfrak w^{\rm sf}_t\rho_0\,dU=\int_{{\rm In}_t}\psi\,(Q_{t\leftarrow0}Y-a_0)\,d\nu_t$, of a modulus at most $\epsilon\int_{{\rm In}_t}|\psi|\,d\nu_t\le\epsilon\sup|\psi|\,Z$, since $\nu_t\ge0$ and $\nu_t({\rm In}_t)=\int\mathfrak w^{\rm sf}_t\rho_0\,dU\le Z$. Applied at $\Psi=Y-a_0$, for $\psi$ and for $1$, with $\langle Y;\psi\circ\mathcal M^t\rangle=\langle(Y-a_0)\,\psi\circ\mathcal M^t\rangle-\langle Y-a_0\rangle\langle\psi\circ\mathcal M^t\rangle$, this is the bound of (i). If $\mathsf R_t\le\frac12$: $|Q_{t\leftarrow0}Y|\le\sup|Y|$ $\nu_t$-almost everywhere and $\nu_t({\rm In}_t)\ge(1-\mathsf R_t)Z>0$, so ${\rm In}_t$ holds points at which $|a_0|\le|Q_{t\leftarrow0}Y|+\epsilon\le\sup|Y|+\epsilon$. (ii) By the first display at the level $T$, where $Z-\mathsf t(h_T)=W_TZ$, the means $\langle F_1F_2\rangle$, $\langle F_1\,\psi\circ\mathcal M^T\rangle$, $\langle F_1\rangle$, $\langle F_2\rangle$ and $\langle\psi\circ\mathcal M^T\rangle$ are the integrals against $Z^{-1}\mathsf m_T$ of the bounded densities of these functions at $h_T$, up to the weight $W_T$ times their suprema, and $Z^{-1}\mathsf m_T(1)=1-W_T$.

The densities of $F_1$, $F_2$ and $F_1F_2$ are $D_{T,1}$, $D_{T,2}$ and $D_{T,1}D_{T,2}+D_{T,12}$; the real part $r_2$ of $D_{T,2}$ is again a density of the real source $F_2$, two bounded densities of one function agree $\mathsf m_T$-almost everywhere, and $r_2=\psi+c$ there; so the density of $F_1F_2$ is $D_{T,12}+D_{T,1}(\psi+c)$ almost everywhere, and subtracting, $\langle F_1;F_2\rangle-\langle F_1;\psi\circ\mathcal M^T\rangle$ is $Z^{-1}\int D_{T,12}\,d\mathsf m_T$, at most $C_2e^{3-\mathsf D_T/M}y_1y_2$ by Lemma~\ref{C-360-24}(e), plus two brackets of errors that together are at most $2\sup|F_1|(\sup|F_2|+\sup|\psi|+|c|)W_T$. At $c={\rm Re}\,c_2$: $|r_2-c|\le|D_{T,2}-c_2|\le C_1y_2$ on the small block fields of the level $T$, which carry $\mathsf m_T$ (Lemma~\ref{C-360-24}(e)), and $|r_2|\le\sup|F_2|$ there.

Every member but $D_{T,12}$ is formed with $F_1$ linearly, and $D_{T,12}$ is additive in $F_1$.
\end{proof}

\begin{lemma}[One entry against separated sources]\label{C-360-18}
Let $0<g\le\gamma$, $\ell\in\mathfrak L(g)$, $n\ge1$, $\mathsf d>0$, $0<c_*<\frac12$, and let $J=\sum_\iota\varkappa_\iota j_\iota$ be an entry of the order $m\ge2$ (Lemma~\ref{C-360-9}) of the type $(m,c,\hat c,N)$: at complex configurations with $|\mathbf U(b)|\le2$ at all the bonds of its box, $|j_\iota|\le c\,r^m$ where the plaquette variables of the box are within $r\le1$ of $\mathbf 1$, and $|j_\iota|\le\hat c$ everywhere; and at most $N$ indices have boxes with one lowest site. There are $C,K_1$, depending on $g$, $\ell$, $n$, $\mathsf d$, $c_*$ and the type, such that for $K\ge K_1$, every label $\mathsf b'\in[\gamma^{-2},X_g+\bar B]$ (so at every first-passage datum $(K,\beta)\in\mathfrak D(g)$, and at $\beta=L^K\mathsf b'$ for every $\mathsf b'\in I(g)$) and all real Lipschitz $f_q$, if the spread of the $n+1$ sets formed by the carrier of $J$ and the ${\rm supp}\,f_q$ is $\ge\mathsf d$, as it is when two of them are at sup distance $\ge\mathsf d$,
\[
\begin{multlined}
|\langle J;\mathcal O(f_1);\dots;\mathcal O(f_n)\rangle_{K,\ell,L^K\mathsf b'}|\le C|\varkappa|_\infty L^{-\mathsf g_mK}K^{(2n+1)\mathfrak p_0}L^{(\frac12-c_*)K}\prod_q\|f_q\|_\sharp,\qquad\mathsf g_2=1,\ \relax \mathsf g_m \\
=\tfrac32\ \relax (m\ge3).
\end{multlined}
\]
On all the sides at once: let besides $R>0$ and $\ell'\in\mathfrak L(g)$ with $\ell'\ge2R+2\mathsf d$, and let the weights $\varkappa_\iota$ be real and the densities real on the $G$-valued configurations. There are $C',K'_1$, depending on $g$, $\ell'$, $n$, $\mathsf d$, $R$, $c_*$ and the type, and on no side, such that for all $(K,\beta)\in\mathfrak D(g)$ with $K\ge K'_1$, every side $\ell\ge\ell'$ in $\mathfrak L(g)$, every entry of the type whose boxes lie in one sup ball of radius $R$, and all real Lipschitz $f_q$ on $\mathbb R^3$ that vanish outside that ball, read on $T_\ell$, with the carrier of the entry at sup distance $\ge\mathsf d$ from every ${\rm supp}\,f_q$ (the sources lie anywhere among themselves),
\[
|\langle J;\mathcal O(f_1);\dots;\mathcal O(f_n)\rangle_{K,\ell,\beta}|\le C'|\varkappa|_\infty L^{-\mathsf g_mK}K^{1+\mathfrak p_0}L^{(\frac12-c_*)K}\prod_q\|f_q\|_\sharp ;
\]
the same holds under every periodic Gibbs state at $(K,\beta)$ and under every weak limit point of the torus laws as $\ell\to\infty$ in $\mathfrak L(g)$, and, with the further factor $2^{n+1}$, for complex weights and complex Lipschitz $f_q$.
\end{lemma}

\begin{proof}[Proof sketch]
Put $\phi=(\frac12-c_*)/(3n+4)$, $u=\lceil\phi K\rceil$, $t=K-u$. Since $M\,R(X_g+\bar B)\le L^m$ and $R$ increases, the frame of Definition~\ref{def:running-source-history} holds for the chain along $x_j=L^{K-j}\mathsf b'$ at every label of $[\gamma^{-2},X_g+\bar B]$, and from a $K_1$ on $t$ is a top level of it with $\mathsf R_t\le\frac12$ and $8a\le\mathsf d$.

\emph{The cut.} Let $D\ge\mathsf d$ be the spread. For every site $y_0$ of the carrier there is a source $q$ with ${\rm dist}_\infty(y_0,{\rm supp}\,f_q)\ge\frac12D$: otherwise there are points $y_q\in{\rm supp}\,f_q$ with $|y_0-y_q|_\infty<\frac12D$, and $y_0,y_1,\dots,y_n$ are pairwise closer than $D$, against the definition of the spread. Give each index $\iota$ of $J$ the least such $q$ for the lowest site of its box; this cuts $J$ into at most $n$ entries $J_q$ of the same type, the carrier of $J_q$ being at sup distance $\ge\frac12D-2a\ge\frac14\mathsf d$ from ${\rm supp}\,f_q$, while the other sources lie anywhere. The cumulant is linear in its first entry and symmetric in the others; in the cumulant of $J_q$ the far source is $f_q$.

\emph{The chain with the entry and a complex source.} The chain of Lemma~\ref{C-360-9}(a) is formed with the entry $sJ_q$ and with the source $w=\sum_q\zeta_qf_q$ at once, for $|s|<r_J$ and \[
\|w\|_\flat=\max\{\sup|w|,\frac{40}{\sqrt3}M{\rm Lip}_\infty w\}<2\varrho
\] (the polydisc of the sources of Theorems~\YMMref{AJ}{673} and~\YMMref{AD}{607}): its families $E^{(i)}[w,s]$ are pointed families of the levels $i$, holomorphic in $(w,s)$ and local, the member at a domain $X$ reading $w$ on the plaquettes of $X$ alone, and $s$ only if a box of the entry lies in $X$; with them $(\mathcal M^t)_*(\mathfrak w^{\rm sf}_te^{sJ_q}\rho^w_0\,dU)=\exp[\sum_{i\le t}c_i(w,s)]\mathbf 1_{{\rm In}_t}e^{\Phi^{w,s}_t}dV$, $\rho^w_0=\rho_0e^{\mathcal O(w)}$, and on ${\rm Cut}_t$
\[
|\Phi^{w,s}_t-\Phi^{w,s'}_t|\le|s-s'|\,\mathsf o_t,\qquad|\Phi^{w,0}_t-\Phi^{0,0}_t|\le N_t\mathsf C_t\|w\|_\flat ,
\]
wherever the sources lie, with $N_t=\ell^3L^{3u}$ the number of the sites of the level $t$ and $\mathsf C_t\le CK^{2\mathfrak p_0}$. The entry enters each step as one more local term, as in Lemma~\ref{C-360-9}(a); bounding the radius $\mathsf e^{\rm c}_t\le1$ of the field cut, raised to the order of the entry, by $\mathsf e^{\rm c}_t$ itself, the size of the entry at the level $t$ and what the steps made of it give $\mathsf o_t\le CK^{\mathfrak p_0}L^{4u-\mathsf g_mK}|\varkappa|_\infty$.

\emph{The bound.} Let $T(s,\zeta)=\int e^{sJ_q+\sum_{q'}\zeta_{q'}\mathcal O(f_{q'})}\,\mathfrak w^{\rm sf}_t\rho_0\,dU$. The cumulant under the positive measure $\mathfrak w^{\rm sf}_t\rho_0\,dU$ is $\partial_s\partial_{\zeta_1}\cdots\partial_{\zeta_n}\log T(0,0)$, and it differs from the cumulant under the measure of the torus by at most $4\,{\rm op}_{n+1}\mathsf R_t\sup|J_q|\prod_{q'}6\ell^3L^{4K}\sup|f_{q'}|$ (Proposition~\YMMref{115.2}{806}, its cumulant bound), which is smaller than every power of $L^{-K}$. By the identity of measures, $\partial_s\log T(0,\zeta)=\partial_s\mathsf C(0,\zeta)+\langle\partial_s\mathsf D\rangle_\zeta$, where $\mathsf C(s,\zeta)=\sum_{i\le t}(c_i(w,s)-c_i(w,0))$ collects the values of the families at the trivial triple, $\mathsf D=\Phi^{w,s}_t-\Phi^{w,0}_t$, and $\langle\cdot\rangle_\zeta$ is the mean under $e^{\Phi^{w,0}_t-\Phi^{0,0}_t}$ times the image of $\mathfrak w^{\rm sf}_t\rho_0\,dU$, normalised.

On the polydisc $|\zeta_{q'}|<(2nN_t\mathsf C_t\|f_{q'}\|_\sharp)^{-1}$ the exponent $\Phi^{w,0}_t-\Phi^{0,0}_t$ is at most $\frac12$, so $|\langle G\rangle_\zeta|\le3.1\sup|G|$, and the second member is at most $3.1\,\mathsf o_t$. In the first, a member at a domain that does not hold both a box of $J_q$ and a plaquette of the far source $f_q$ does not depend on $\zeta_q$ (locality), and has no mixed derivative; the members at the domains $X$ of the levels $i$ that hold both have $d_i(X)\ge(\frac14\mathsf d-a)L^{K-i}/M-2$, and with the count of the domains their sum is at most a power of $L^K$ times $e^{-\mathsf dL^u/8M}$. Cauchy's estimates in the $n$ sources on this polydisc cost $\prod_{q'}(nCK^{2\mathfrak p_0}L^{3u}\|f_{q'}\|_\sharp)$.

Together, $K^{(2n+1)\mathfrak p_0}L^{(3n+4)u-\mathsf g_mK}$, and $(3n+4)u\le(\frac12-c_*)K+3n+4$.

\emph{On all the sides.} Let $\mathsf l=L^{-i}$ be the largest power of $L$ with $\mathsf l\le\frac12\mathsf d$; then $\ell'\ge2R+4\mathsf l$. On $T_{\ell'}$ and on every $T_\ell$, $\ell\ge\ell'$, a board of side $\mathsf l$ puts the boxes of the entry and the supports of the sources, which lie in one sup ball of radius $R$ inside one period box of side $\ell'$, in different cells. The cumulant on $T_\ell$ is a mixed derivative at $0$ of a quotient whose coefficients are means on $T_\ell$ of products of functions of different cells of that period box; the chessboard bound (Lemma~\ref{C-360-7}) bounds each by a product of chessboard seminorms on $T_{\ell'}$, at the price $e^{\mathsf l^3(p_{\ell'}-p_\ell)}\le e^{G_g\mathsf l^3}$ for a cell; this ratio $Z_{K,\ell',\beta}^{1/N'}/Z_{K,\ell,\beta}^{1/N_\ell}$ of the plain partition functions per cell, $N'$ and $N_\ell$ being the numbers of the cells of $T_{\ell'}$ and of $T_\ell$, is the one number that depends on $\ell$, and its bound is free of $\ell$.

After a subtraction of Wick type in the functions of the cells, which removes the cumulants of the families inside one cell, the chessboard seminorm of the function of the cell of the entry, the mean on $T_{\ell'}$ of the product of its mirror copies, is by the moment--cumulant formula a sum over the partitions of the copies all of whose blocks join different cells; each such cumulant is small by the bound on one torus, read at a free depth $u$ with $C_{\rm u}K^{1/(2\mathfrak p_0)}\le u\le\frac12K$, $C_{\rm u}\ge1$ a number formed with the data and the type: a factor $|\varkappa|_\infty K^{1+\mathfrak p_0}L^{4u-\mathsf g_mK}$ for each copy of the entry and a factor of Cauchy's estimate for each copy of a source.

The number of the cells that the ball of radius $R$ meets, at most $(2R/\mathsf l+3)^3$, enters $C'$, and at $u=\lceil\frac14(\frac12-c_*)K\rceil$ the bound follows. Every periodic Gibbs state is dominated by the torus measures through its conditional laws (Proposition~\YMMref{91.9}{705}), and the weak limit points of the torus laws inherit the bound; for complex data the real and imaginary parts are taken apart.
\end{proof}

\begin{lemma}[The cubic remainder uniformly in the side]\label{C-360-19}
Let $0<g\le\gamma$, $\ell'\in\mathfrak L(g)$, $n\ge2$, $\mathsf d>0$, $c<\frac1{14}$, and let $Q$ be a cube of side $\frac14\ell'$. There are $C,p,K_0$, depending on $g,\ell',n,\mathsf d,c$ and independent of the larger side $\ell$, such that for all $\ell\ge\ell'$ in $\mathfrak L(g)$, all $(K,\beta)\in\mathfrak D(g)$ with $K\ge K_0$, all Lipschitz fields $\eta$ and real Lipschitz $f_q$ carried by $Q$ with the supports of the $f_q$ pairwise $\ge\mathsf d$ apart,
\[
|\mathfrak U^\sharp[\eta](f)_{T_\ell}|+|\mathfrak U^\sharp[\eta](f)_{T_\ell}-\mathfrak U^\sharp[\eta](f)_{T_{\ell'}}|\le CK^pL^{-cK}\mathsf H(\eta)\prod_q\|f_q\|_\sharp ;
\]
the finitely many sides of $\mathfrak L(g)$ below $\ell'$ are covered by Lemma~\ref{C-360-10}.
\end{lemma}

\begin{proof}[Proof sketch]
Every label of $\mathfrak D(g)$ lies in $[\gamma^{-2},X_g+\bar B]$, and $MR(X_g+\bar B)\le L^m$ on every side of $\mathfrak L(g)$; so the frame of Definition~\ref{def:running-source-history} holds for the chain along $x_j=L^{K-j}\mathsf b$ at the label of every first-passage datum, and Lemma~\ref{C-360-9}(b) holds with one triple of numbers for all of them, also with real sources $\sum_ix_if_i$ in its polydisc. The chain is formed with several entries and complex sources as well, as in the proof of Lemma~\ref{C-360-18} (Lemma~\ref{C-360-9}(a), polarisation and locality in each entry).

\emph{The reduction.} With $\mathcal S[\Xi]=a^{-2}\sum_p\Xi^{\lambda(p)}(x(p))\,\mathfrak r^\sharp_{\lambda(p)}(p)$ for a field $\Xi$ we have $\beta\mathfrak K^\sharp[\eta]=\mathcal S[\mathsf b\eta]$ and $a^{-1}\mathfrak K^\sharp[f_i\eta]=\mathcal S[f_i\eta]$, so $\mathfrak U^\sharp[\eta](f)=-\langle\mathcal S[\mathsf b\eta];\mathcal O(f_1);\dots;\mathcal O(f_n)\rangle+\sum_i\langle\mathcal S[f_i\eta];(\mathcal O(f_l))_{l\ne i}\rangle$, with $\mathsf H(\mathsf b\eta)\le(X_g+\bar B)\mathsf H(\eta)$ and $\mathsf H(f_i\eta)\le2\|f_i\|_\sharp\mathsf H(\eta)$. It suffices therefore to compare $G_\ell=\langle\mathcal S[\Xi];\mathcal O(h_1);\dots;\mathcal O(h_{n'})\rangle_{K,\ell,\beta}$ with $G_{\ell'}$, $n'\in\{n-1,n\}$, $n'\ge1$, for Lipschitz fields $\Xi$ and real Lipschitz $h_q$ carried by $Q$ whose supports are pairwise $\ge\mathsf d$ apart, nothing being asked of $\Xi$ relative to them.

\emph{The fold with one factor for each copy.} Place $Q$ in the middle of one cell of the board of side $\mathsf l=\frac12\ell'$; it lies at sup distance $\ge\frac18\ell'\ge1$ from the walls. On every $T_\ell$, $\ell\ge\ell'$, $\mathcal S[\Xi]$ and the $\mathcal O(h_q)$ are the realisations of one model function each, and their mirror copies in the eight cells of $T_{\ell'}$ are again a cubic entry $\mathcal S[\Xi^\kappa]$, $\mathsf H(\Xi^\kappa)\le\mathsf H(\Xi)$, and sources $\mathcal O(h_{(q,\kappa)})$ with $\|h_{(q,\kappa)}\|_\sharp=\|h_q\|_\sharp$; copies in different cells are $\ge1$ apart. The chessboard bound (Lemma~\ref{C-360-7}) carries the cumulant on $T_\ell$ to the mean on $T_{\ell'}$ of the product of the eight mirror copies, at the price $e^{G_g\mathsf l^3}$ for each function, and by the moment--cumulant formula of the copies (an identity of formal power series, as in the proof of Theorem~\ref{C-360-2}(g)(ii)) the difference $G_\ell-G_{\ell'}$ reads the joint cumulants on $T_{\ell'}$ of the families $B$ of copies that join different cells alone, the families inside one cell entering nowhere. If each of these is at most $\Gamma\prod_{(i,\kappa)\in B}\omega_i$, where $\Gamma\ge1$, with one factor $\omega_0$ for each copy of the entry and one factor $\omega_q$ for each copy of the source $h_q$, then
\[
|G_\ell-G_{\ell'}|\le\mathsf c_{n'}\,e^{(n'+1)G_g\mathsf l^3}\,\Gamma^{(n'+1)/2}\prod_{i=0}^{n'}\omega_i ,
\] To see the exponent, use the within-cell normalization and logarithm identities of Lemma~\ref{C-352-31}. In a fold of $q$ entries, every remaining block joins cells and has at least two members; a partition of the $8q$ copies has at most $4q$ such blocks. Its bound is at most $\Gamma^{4q}\prod_i\omega_i^8$, and the eighth root gives $\Gamma^{q/2}\prod_i\omega_i$. The subsequent logarithm partitions preserve this exponent, since their subset sizes sum to $q=n'+1$. This proves the display, including the stated pressure factor.
Here $\mathsf c_{n'}$ is a pure number. With $u=\lceil\phi K\rceil$, $t=K-u$, the factors $\omega_0=K^pL^{-c_1K}\mathsf H(\Xi)$ and $\omega_q=K^{2\mathfrak p_0}L^{3u}\|h_q\|_\sharp$ serve, in three cases. (i) No copy of the entry: the family holds two copies of sources in different cells, $\ge1$ apart, and the bound \eqref{eq:torus-cumulant-bound} on $T_{\ell'}$ at separated supports applies. (ii) One copy $\mathcal S[\Xi^\kappa]$: moving the weight of the entry with the sources writes the cumulant as $-\mathfrak U^\sharp[\Xi^\kappa/\mathsf b](h)$ plus, when two or more sources are present, the cubic remainders $-\mathfrak U^\sharp[h_\alpha\Xi^\kappa/\mathsf b^2]((h_{\alpha'})_{\alpha'\ne\alpha})$ (the supports of the $h$'s being disjoint), and with one source the contact member is a mean of $\mathfrak K^\sharp$, which vanishes; a source of another cell is $\ge1$ from the entry, and $h_\alpha\Xi^\kappa$ is carried by ${\rm supp}\,h_\alpha$, at sup distance $\ge\min\{\mathsf d,1\}$ from the other supports: the spread condition holds, and Lemma~\ref{C-360-10} on $T_{\ell'}$ at the rate $c_1$ applies. (iii) Two to eight copies of the entry, in different cells, the sources anywhere: the cumulant is a mixed derivative of the logarithm of the generating function of the chain with these entries and the complex source $w_\zeta=\sum\zeta_\alpha h_\alpha$, at the level $t$.
For case (iii), retain the higher mixed jets as in the proof of Proposition~\ref{stc:cum}(b). Proposition~\ref{stc:chain}(v) permits the fixed complex source $w$ and the several insertion parameters $\mathbf s$ together. Write their small-field generating function as $T_w(\mathbf s)=T_w(0)e^{\mathsf C_w(\mathbf s)}\langle e^{\mathsf D_{w,\mathbf s}}\rangle_w$, with the normalized block-field functional $\langle\cdot\rangle_w$. For $q$ copies of the entry and $I$ with $|I|\ge2$, locality and the Cauchy estimate in that proof give $\max\{|\partial_I\mathsf C_w(0)|,\sup|\partial_I\mathsf D_{w,\mathbf s}|_{\mathbf s=0}\}\le2E_0\mathsf t'_\times\prod_{a\in I}(q/r_a)$, where $r_a$ are the insertion radii and $\mathsf t'_\times$ is the polymer sum of that proposition. Copies of the entry lie in distinct cells with a fixed positive margin. Their inverse radii are at most $CL^{2K}\mathsf H(\Xi^a)$, whereas $\mathsf t'_\times\le CK^pL^{AK}e^{-cL^u}$. Hence the mixed-jet bound is below $\prod_{a\in I}\epsilon_a$ for large $K$, uniformly on a fixed fraction of the displayed, $K$-dependent source polydisc. Zero entries are omitted. At real sources the block-field functional is positive; at complex sources it has norm at most $3.1$ on the smaller polydisc, by the exponent bound in the proof of Lemma~\ref{C-360-18}. The moment and logarithm partition formulas in Proposition~\ref{stc:cum}(b) therefore apply in both cases, with finite constants adjusted in the complex case.
At a real source $x$ in the polydisc of the radii $(CK^{2\mathfrak p_0}L^{3u}\|h_\alpha\|_\sharp)^{-1}$, Lemma~\ref{C-360-9}(b), applied to the field $\Xi^\kappa/(\mathsf b-w_x)$, makes the conditional mean of each copy given the block field of the level $t$ a number up to $\epsilon_a=C\mathsf H(\Xi)K^{p_2}L^{-c'_0K}$. Subtracting these constants from the first jets, and using the higher mixed-jet bound above, the same partition formulas bound the cumulant by a constant times $\prod_a\epsilon_a$. At a complex source, the norm bound for the normalized functional and the mixed-jet estimate instead give a constant times the product of the crude bounds $\hat{\mathsf o}_a=C\mathsf H(\Xi)K^{\mathfrak p_0}L^{4u+\frac12K}$ of the oscillations of the cubic entries, of the order $3$ and of the weight $a^{-2}$. The estimate with two constants for a function holomorphic on a polydisc of the radii $r_\alpha$, bounded by $B$ there and by $\epsilon$ at its real points, $|\partial_{\zeta_1}\cdots\partial_{\zeta_{n''}}(0)|\le\epsilon^{1-\sigma}B^\sigma(\frac{4n''}{\pi\sigma})^{n''}\prod_\alpha r_\alpha^{-1}$, gives one factor $\epsilon_a^{1-\sigma}\hat{\mathsf o}_a^\sigma\le C\mathsf H(\Xi)K^pL^{-c_1K}$ for each copy of the entry, $c_1=(1-\sigma)c'_0-\sigma(\frac12+4\phi)$.

In all three cases the histories with a large field up to the level $t$ cost on $T_{\ell'}$ less than every power of $L^{-K}$. With $\delta=\frac1{14}-c$, $\tau=\delta$, $\phi=\delta/(12(n+1))$ and $\sigma=\frac14\delta$ one has $c_1-3n\phi>c$, and the product $\omega_0\prod_q\omega_q$ gives the bound of the difference with the rate $c$ and with numbers independent of the larger side $\ell$; on $T_{\ell'}$ itself the bound is Lemma~\ref{C-360-10}.
\end{proof}

\begin{lemma}[The flow with an insertion]\label{C-360-20}
Let $\ell=2L^m$ be any side, $0<\mathsf b_-\le\mathsf b\le\mathsf b_+$, $\beta=L^K\mathsf b$, $n\ge1$, $\mathsf d>0$ and $f_i$ real Lipschitz; let $S_\varkappa=\sum_l\varkappa(l)\,{\rm dn}_l$ be an insertion of brick densities ${\rm dn}_l$ (gauge invariant polynomials on boxes of at most five sites along each axis, each holding a plaquette, vanishing at $\mathbf 1$) of the order $r$, and let two of the $n+1$ sets formed by the carrier of $S_\varkappa$ and the ${\rm supp}\,f_i$ be at sup distance $\ge\mathsf d$, the others in any position. For $\gamma$ small enough in terms of the other constants there are $C,p,K_0$, depending on $\ell,\mathsf b_\pm,n,\mathsf d$ and the type, such that for $K\ge K_0$:
\begin{enumerate}
\item[(a)] if $r\ge3$ and $|\varkappa|_\infty\le\beta h$: $|\langle S_\varkappa;\mathcal O(f_\cdot)\rangle|\le CK^pL^{-\lambda_LK/4}h\prod_i\|f_i\|_\sharp$; no norm of $\varkappa$ but its supremum enters;
\item[(b)] if $r=2$, $0\le\sigma\le1$, $|\varkappa|_\infty\le L^{\sigma K}h$ and $0<\epsilon\le1$: $|\langle S_\varkappa;\mathcal O(f_\cdot)\rangle|\le C_\epsilon K^pL^{-(1-\sigma-\epsilon)K}h\prod_i\|f_i\|_\sharp$; in particular $|\langle\mathcal O(w);\mathcal O(f_\cdot)\rangle|\le C_\epsilon a^{-\epsilon}\max|w|\prod_i\|f_i\|_\sharp$ for every real weight $w$, in contact with the $f_i$ or not;
\item[(c)] if $r=2$, the insertion is real with $|\varkappa|_\infty\le\beta h$, and its densities are of derivative kind, that is, $|G_\nu(b)|^2$ or $\Psi_{\lambda(p)}(p)$: $|\langle S_\varkappa;\mathcal O(f_\cdot)\rangle|\le CK^pL^{-c_nK}h\prod_i\|f_i\|_\sharp$, $c_n=\lambda_L/(4(6+3n))$; for every $0<\theta_1<1$ and every real insertion of the order $2$ with $|\varkappa|_\infty\le\beta h$, $\langle S_\varkappa-J_{t_1};\mathcal O(f_\cdot)\rangle$ obeys the same bound with $L^{-\lambda_L\theta_1K/4}$, the numbers $C,p,K_0$ depending besides on $\theta_1$, where $t_1=K-\lceil\theta_1K\rceil$ and $J_{t_1}$ is the classical value of the insertion at the background of the level $t_1$, a function of the block field.
\end{enumerate}
\end{lemma}

\begin{proof}[Proof sketch]
Run the chain with the insertion, as in Lemma~\ref{C-360-9}(a), down to a level $t$ at which the large fields are rare. A transfer inequality bounds the cumulant by
\[
2\delta_t(\varkappa)\prod_i\bigl(8nM_t\|f_i\|_\sharp\bigr)+{\rm Con}_t+4\,{\rm op}_{n+1}\mathsf R_t\sup|S_\varkappa|\prod_i\bigl(6\ell^3L^{4K}\sup|f_i|\bigr),
\]
with $\delta_t$ the oscillation of the conditional mean of the insertion given the block field of the level $t$, $M_t$ the inverse radius of the sources at that level, and ${\rm Con}_t$ the mixed derivative of the constants of the chain, which is small by locality when a pair of sets is separated (when the separated pair is the insertion and one source, the constants of the sources that meet the insertion are handled by a slow cut-off at the scale $\mathsf d$, at the price $(2+480M/\mathsf d)^{n-1}$). By power counting an insertion of the order $r$ has the size $|\varkappa|_\infty L^{(3-2r)i}$ at the level $i$, and what the steps make of it is forgotten at the rate $L^{-\lambda_L/2}$ per level, $(L/2)^{-x}=L^{-\lambda_Lx}$; this gives (a) and (b).

For (c) two levels are used. Let $0<\theta_1<1$, $\tau_1=\lceil\theta_1K\rceil$, $t_1=K-\tau_1$, and $t=K-\tau$, $\tau=\lceil\theta K\rceil$, $\theta=\lambda_L\theta_1/(4(7+3n))<\theta_1$. The insertion is put into the chain only up to the level $t_1$, as early members; there it equals its classical value $J_{t_1}$ on the background of the level $t_1$ plus what the steps up to $t_1$ made of it, which in level-$t_1$ terms has the size $C|\varkappa|_\infty L^{-(t_1-1)}L^{-\lambda_L\tau_1/2}$ (what a step makes of an insertion of the order $2$ is forgotten at the rate $L^{-\lambda_L/2}$ per level). From the level $t_1$ to the level $t$ this part is forgotten by the factor $L^{-(t-t_1)}$ of the run, on the field cut of the level $t$; at the level $t$ the histories with a large field are discarded, at a cost below every power of $L^{-K}$, and the volume $N_t=\ell^3L^{3\tau}$ and Cauchy's estimates in the $n$ sources cost $L^{3(n+1)\tau}$ up to powers of $K$; the mixed derivatives of the constants of the chain in the sources are small by locality, two of the sets being separated.

The derivative in the insertion is taken as a difference quotient at a small step $s$, its second derivative being bounded crudely. Together the exponent is $(7+3n)\tau+1-\frac12\lambda_L\tau_1\le-\frac14\lambda_L\theta_1K+8+3n$: this is the second clause of (c), for every $\theta_1$. For the first clause the classical part is bounded besides: $J_{t_1}$ is a function of the block field of the level $t_1$, its cumulant with the sources is bounded by Cauchy's estimates in the sources at that level, at the price $L^{3n\tau_1}$, times its supremum, and for densities of derivative kind the supremum is at most $\beta\ell^3L^{3K}\,C\varepsilon_{t_1}^2(1+K\log L)^2L^{-6t_1}$ by the estimate below; the exponent is $K+3K-6(K-\tau_1)+3n\tau_1=-2K+(6+3n)\tau_1$, which at $\theta_1=1/(6+3n)$ is at most $-K+6+3n$.

With the second clause at this $\theta_1$ the rate is $\min\{1,\frac14\lambda_L\theta_1\}=c_n$. For a lattice field $B$ on a box of radius $R$ with $|B|\le A_0$ and $|{\rm div}^+B|,|{\rm curl}^-B|\le\epsilon_1$, the div--curl estimate $|\partial^\pm_jB_i|\le C_*[A_0/R+\epsilon_1(1+\log R)]$ on the quarter box, applied with the Bianchi identity on a cube, gives $|G_\nu(b)|^2,|\Psi_{\lambda(p)}(p)|\le C\varepsilon_i^2\eta_i^6(1+i\log L)^2$ at the regular backgrounds of the level $i$, $\eta_i=L^{-i}$, which is the order that (c) needs.
\end{proof}

\begin{proposition}[The mean plaquette]\label{C-360-21}
Let $0<g\le\gamma$, $\ell\in\mathfrak L(g)$, $n\ge1$, and $E(\beta)=\langle a_p\rangle$ on a periodic lattice of $N_s$ sites along each axis (it does not depend on $p$).
\begin{enumerate}
\item[(a)] For every $R\ge3$, $N_s\ge R$ and $\beta$ with $116R^3(\beta/2)^{-1/2}\le\frac15$ and $\bar\eta(R,\beta)\le\frac14$, $|\beta E(\beta)-1|\le2\bar\eta(R,\beta)^{1/2}$, with $\bar\eta(R,\beta)\le\frac{9\log(2\beta)+55}{R^3}+\frac{18\log(1+256R^3)+15}R+1.21\cdot10^9R^9\beta^{-1/2}$: at $R=\lceil\log\beta\rceil$ this tends to $0$ uniformly in all volumes $N_s\ge\log\beta$.
\item[(b)] On the side $\ell$, with $N_s=\ell L^K$: there are $C,K_0$, depending on $g,\ell,n$, such that $\sup_{\mathsf b'\in I(g)}|\beta^nE^{(n-1)}(\beta)-(-1)^{n-1}(n-1)!|\le CK^{2n+3}L^{-K/20}$ for $\beta=L^K\mathsf b'$ and $K\ge K_0$.
\item[(c)] For all real Lipschitz $f_1,\dots,f_n$ on $T_\ell$, with nothing asked of the supports, $|a_K^3S^T_{n;K,\ell,L^K\mathsf b'}(f)-3(n-1)!\,\mathsf b'^{-n}\int_{T_\ell}\prod_if_i|\le CK^{2n+3}L^{-K/20}\prod_i\max\{\sup|f_i|,\frac{40}{\sqrt3}M{\rm Lip}_\infty f_i\}$.
\end{enumerate}
\end{proposition}

\begin{proof}[Proof sketch]
(a) Cut the lattice into blocks of $R$ to $2R$ sites along each axis, and in each block take a comb tree. Bounding the partition function with plaquette weights $\beta_p\in[\beta_-,\beta_+]$ from above (dropping a few plaquettes) and from below (freezing a few bonds), both bounds become the same Gaussian integral, $\sum_B(2|B|-2)\log\mathsf g(\beta_B)-\mathsf D$ with $\mathsf g(\beta)=(2/\pi)^{1/2}\beta^{-3/2}$ and a number $\mathsf D$ of the geometry, up to an error $N\eta$ that is small per site; this uses weak coupling on blocks of $R\le\beta^{1/20}$ spacings alone. The convexity of $\log Z$ in $\beta$ then pins down its first derivative. (b) For the higher derivatives analyticity in $\beta$ is needed.

On a torus of a fixed physical side, Theorems~\YMMref{AJ}{673} and~\YMMref{AD}{607} write $Z[\zeta f]=\exp(E+\sum_{j,X}e_j(X;\zeta f))\widetilde Z(\zeta f)$ for real $f$ with $\max\{\sup|f|,\frac{40}{\sqrt3}M{\rm Lip}_\infty f\}\le1$, with $|e_j(X;w)|\le C_{\rm S}e^{-\kappa d_j(X)}$, $|\widetilde Z|\le e^{C_{\rm S}\ell^3}$ and $\widetilde Z(0)\ge e^{-C_g\ell^3}$; the factor $\widetilde Z$ may vanish. Put $N=(\ell L^K)^3$, $h_f(\zeta)=\widetilde Z(\zeta f)/\widetilde Z(0)$ and $D_f(\zeta)=(3N)^{-1}[\sum_{j,X}e_j(X;\zeta f)+\sum_p\log(1-\zeta f(x(p))/\mathsf b')]$, the per-site exponent; by the count of polymers $|D_f|\le M_0$ on a disc $|\zeta|<\rho_1$, with $M_0$ and $\rho_1$ free of $K$. Then $h_f(0)=1$, $h_f>0$ on the real axis and $|h_f|\le e^\Gamma$ with $\Gamma=(C_{\rm S}+C_g)\ell^3$, a number free of $K$.

First, Cauchy's estimate on circles of radius $\frac12\rho_1$ gives $|h_f-1|\le\frac12$ on $|\zeta|\le\frac14\rho_1e^{-\Gamma}$, so that $\log h_f$ is defined there and $|(\log h_f)^{(n)}(0)|\le n!\,(4e^\Gamma/\rho_1)^n$: this member does not grow with the cutoff, and it is of the order $1/N$ per site. Second, the Poisson--Jensen formula and Harnack's inequality, applied to the non-negative function $\Gamma-\log|h_f|$, give $\int_{-\rho_1/2}^{\rho_1/2}|\log|h_f(x)||\,dx\le9\rho_1\Gamma$; with this, the two-sided real-axis bound of (a), at the plaquette weights $\beta(1-xf(x(p))/\mathsf b')$, controls the per-site exponent in $L^1$ on the segment: $\int_{-\rho_1/2}^{\rho_1/2}|D_f(x)|\,dx\le\rho_1\delta$ with $\delta=\frac13(2\bar\eta(R,\beta)+Ra_K+R^{-3})+3(C_{\rm S}+C_g)L^{-3K}$.

Third, A.~Markov's inequality for the derivatives of the Taylor polynomial of $D_f$ of the order $\mathsf m=\max\{n,\lceil\log_2(M_0/\delta)\rceil\}$, with Cauchy's estimate for the rest of the series, gives $|D_f^{(n)}(0)|\le4\,(2/\rho_1)^n\mathsf m^{2n}(\mathsf m+1)^2\delta$, and at $R=\lfloor\beta^{1/20}\rfloor$ this is the rate. By Jensen's formula $h_f$ has at most $\Gamma/\log2$ zeros in $|\zeta|<\frac12\rho_1$, a number that does not grow with the cutoff; the bound on $(\log h_f)^{(n)}(0)$ above does not need it. At $f\equiv1$, where $S^T_n(1,\dots,1)=a_K^{-n}(-1)^{n-1}3N\,E^{(n-1)}(\beta)$, the identity of (c) below gives (b). (c) is the same argument at a general $f$: $S^T_n(f,\dots,f)=(n-1)!\mathsf b'^{-n}\sum_pf(x(p))^n+3ND^{(n)}_f(0)+(\log h_f)^{(n)}(0)$, with the per-site exponent $D_f$ and the factor $h_f$ of Theorems~\YMMref{AJ}{673} and~\YMMref{AD}{607}; polarisation gives distinct entries.
\end{proof}

Use the separated regions and the actual averaging constraints of
Proposition~\ref{prop:weighted-system-L2}, with its block-number and
small-triple hypotheses in addition to the width condition below.
Require also the geometric and regularity hypotheses of
Theorem~\YMMref{I}{225} and Proposition~\YMMref{29.1}{231},
with their block-number, torus-size and gap floors. The matrix-valued
plaquette-divergence bound in that theorem follows from the regular
curvature and current bounds by equation~\YMMref{(579)}{288},
with a factor $2$ absorbed in the regularity ceiling. These are
numerical floors and small-radius conditions on the data, not
assumptions about the inverse estimate to be proved.
Also retain the four local-system radii at most $a_H$, the block
floor $\mathfrak M\ge M_H$, and the permitted decay rate of
Lemma~\ref{lem:ordinary-holder-kernels}.
Write $\xi=\eta_k$ and $\ell_n=\xi L^n$ for the physical layer scale.
For clarity distinguish $D_0$, the covariant differential from sites
to bonds, from $D_1$, the covariant differential from bonds to
plaquettes. Let $U$ be a real background, critical for the unweighted
action on the kernel of the averaging constraint, regular at radius
$\alpha\le\alpha_{\rm lin}(L)\le(16L^2)^{-1}$, with curvature bounded
by the envelope $e(x)=\varepsilon_{j(x)}$ in the scaled norm below,
and $\sup e\le e_0(L)$. The layer radii satisfy
$0<\varepsilon_j\le\varepsilon_{j+1}\le L\varepsilon_j$.
Require
$R_1\mathfrak M\ge w_{\rm lin}(L)$.

At each bare site $x_*$ of layer $n\ge2$, fix a real adjoint gauge
$g_{x_*}$ on the cube of physical sup radius $\ell_n/4$ such that
\[
 |g_{x_*}(b_-)\operatorname{Ad}U(b)g_{x_*}(b_+)^{-1}-I|
       \le L^2\alpha/L^n.
\]
The gauge hypotheses are part of the setting. All H\"older quotients
below compare components of the same orientation in these gauges.
The cube for a potential is the concentric cube of radius
$\ell_n/8-2\xi$; those for a curvature and a source two-form have radius $\ell_n/4$. Quotients are taken over distinct admissible base points
in the indicated cube, with the usual physical Euclidean distance.

We use the following scaled spaces. For a field of inverse
length degree $r$, put
\[
 |f|_{r,h}=\sup_i\frac{\ell(i)^r|f(i)|}{h(x_i)},\qquad
 [f]_{r,1/2,h}
 =\sup_{n\ge2,x_*}\frac{\ell_n^{r+1/2}}
                              {h(x_*)}[g_{x_*}f]_{1/2;C(x_*)}.
\]
For an index $i$, let $\ell(i)$ be the largest layer scale of the
sites in the stencil defining that field value or first difference.
For a row use its prescribed scale $\xi_c$. The envelope is read at
the indicated base point $x_i$. This fixes the norm convention also
for indices meeting a layer boundary. The adjacent-layer comparison
of Lemma~\YMMref{21.2}{141} makes it equivalent, with constants
depending only on $L$, to the base-point scale convention. Set
$\|f\|_{r,1/2,h}=|f|_{r,h}+[f]_{r,1/2,h}$.
For the first two layers only the supremum term is required; their
ratios $\ell_n/\xi\le L$ are uniformly bounded. The curvature
hypothesis above means $\|F_U\|_{2,1/2,e}\le1$.

The row units are dimensionless. If $Q_{\rm ph}$ is the physical
averaging derivative with row pairing of weight $\xi_c^3$, put
\[
 QA(c)=\xi_c Q_{\rm ph}A(c),\qquad
 \omega(c)=\xi_c^3\omega_{\rm ph}(c).
\]
Use the row pairing of weight $\xi_c^{-1}$ for $Q^{\mathsf T}$.
Then $Q^{\mathsf T}\omega=Q_{\rm ph}^{\mathsf T}\omega_{\rm ph}$,
since both pairings with $A$ are
$\sum_c\frac12\operatorname{tr}(Q_{\rm ph}A(c)\omega(c))$.
This conversion changes units only; it retains the same row set and
averaging germ. Row data $b$ and multipliers have degree zero.

Let $h=e\theta$, where $\theta>0$ and
$|\log\theta(x)-\log\theta(x')|\le c_Ud(x,x')$.
Choose $c_U$ at most one eighth of each decay rate used below and
at most $1/(6L)$. Thus, for fixed constants
$C_h=C_h(L)$ and $0<\nu\le\delta_2/2$,
\begin{equation}\label{eq:linear-input-envelope}
       h(z)/h(y)\le C_h e^{\nu d(y,z)}.
\end{equation}
Take the same condition on the local comparison cubes. Define
\begin{align*}
 \|(A,\omega)\|_{\mathcal X_h}
  &=\|A\|_{1,1/2,h}+\|D_1A\|_{2,1/2,h}
       +|D_0^{\mathsf T}A|_{2,h}+|\omega|_{0,h},\\
 \|(Y,f,b)\|_{\mathcal F_h}
  &=\|Y\|_{2,1/2,h}+|f|_{3,h}+|b|_{0,h}.
\end{align*}
The ordinary and divergence sources are read only through
$f+D_1^{\mathsf T}Y$. For a complex anisotropy ${\rm an}$, its size
$t({\rm an})$ is its supremum plus the scaled H\"older-$1/2$
seminorm on the same plaquette cubes.

\begin{lemma}[Linear response at a variable anisotropy]\label{C-360-22}
There are positive numbers
$w_{\rm lin},\allowbreak\alpha_{\rm lin},\allowbreak e_0,\allowbreak t_{\rm lin},\allowbreak K_{\rm lin}$,
depending only on $L$, such that the following holds in the setting
above. If $t({\rm an})\le t_{\rm lin}\le1/4$, the system
\begin{equation}\label{eq:uniform-linear-response-system}
 QA=b,\qquad
 T_\tau A=f+D_1^{\mathsf T}Y+Q^{\mathsf T}\omega,
 \qquad \tau=1+{\rm an},
\end{equation}
has a unique solution, linear in the data, with
\begin{equation}\label{eq:uniform-linear-response-bound}
 \|(A,\omega)\|_{\mathcal X_h}
      \le K_{\rm lin}\|(Y,f,b)\|_{\mathcal F_h}.
\end{equation}
Here
$T_\tau=D_1^{\mathsf T}\tau cD_1+V_0'+D_0RD_0^{\mathsf T}$
is the second variation of the weighted action at the reference
field, with the slice term. The coefficient $c_p$ is determined by
$\partial U(p)=c_pI+i\xi^2F_U(p)$, and $V_0'$ contains the
curvature commutator terms of that second variation. The nonlinear
block-constraint remainder and its multiplier are kept in the
forcing when the lemma is applied to the nonlinear equation.
The row and slice conventions are those of the setting. All
constants are uniform in the volume, the number of layers and the
permitted envelopes. For the mean-scheme rows of a label, rotated
into the comb frame, the same estimate holds with $2K_{\rm lin}$.
\end{lemma}

\begin{proof}[Proof of the linear estimate]
We first treat the unweighted operator. The block constraint must be
retained in this argument: the inverse of the gauge-fixed operator
without the constraint is not the required Green operator. Put
\[
 T=T_\tau\big|_{{\rm an}=0},\qquad
 G=(T+Q^{\mathsf T}Q)^{-1},\qquad
 M=QGQ^{\mathsf T},
\]
and define
\begin{equation}\label{eq:endpoint-constrained-inverses}
 H=GQ^{\mathsf T}M^{-1},\qquad
 \widetilde G=G-HQG.
\end{equation}
Thus $QH=1$, $Q\widetilde G=0$ and
$\widetilde GT=1-HQ$. The rows are the interior and outward
determining bonds; inward bonds are discarded, so there are no edge
rows. On a row of level $j$, the derivative of the iterated comb
average is
\[
 Q A(c)=\xi\,
   \bigl(\mathbf Q(W_{j-1})\cdots\mathbf Q(W_0)A\bigr)(c).
\]
The factors read only the double block beneath that row. Separation
keeps these blocks inside the domains on which the averages are
defined. The row conversion above therefore identifies this $Q$
with the physical averaging derivative in the regular constrained
system, without changing the set of constraints.

The regular current bound supplies the matrix-divergence hypothesis
used for that system. Indeed, at a full bond the quaternion
decomposition gives
\[
 D_U^{\mathsf T}\partial U=aI+i\xi^2J(U),\qquad
 |a|\le\xi^{-1}\max_{p\ni b}|\partial U(p)-I|^2.
\]
Since the adjacent plaquette scale is at least the bond scale and
the latter is at least $\xi$, the regularity bounds give
$|D_U^{\mathsf T}\partial U|\le2\alpha\xi^2\ell(b)^{-3}$.
This is equation~\YMMref{(579)}{288}; our choice
$2\alpha_{\rm lin}\le a_Q$ permits the application below.

Here are the unweighted bounds that will be used. Write
$|f|_{(r)}=\sup_i\ell(i)^r|f(i)|$ and
$\|A\|_X=\max\{|A|_{(1)},|\nabla A|_{(2)}\}$. For sufficiently
small regularity radius, there are $K_1,K_7\ge1$, depending only
on $L$, such that
\begin{equation}\label{eq:endpoint-six-bounds}
\begin{gathered}
 \|\widetilde Gf\|_X\le K_7|f|_{(3)},\qquad
 \|Hb\|_X\le K_1\sup|b|,\qquad
 \sup|H^{\mathsf T}f|\le K_1|f|_{(3)},\\
 \sup|(D_1H)^{\mathsf T}Y|\le3K_1|Y|_{(2)},\qquad
 \sup|(D_0^{\mathsf T}H)^{\mathsf T}\lambda|
       \le3K_1|\lambda|_{(2)},\qquad
 |R\lambda|_{(2)}\le(1+K_1^3)|\lambda|_{(2)}.
\end{gathered}
\end{equation}
These estimates are unchanged by passing to the units of any layer.
The averaging rows are dimensionless, their pairing has weight
$\ell(c)^{-1}$, and their transposes are taken for that pairing.
They are precisely the bounds of
Proposition~\YMMref{29.1}{231}, applied after the row identification
just made. Its proof also records their exponential versions.
For the transpose estimates, test against a datum supported at one
row $c$. The row pairing contributes a factor $\ell(c)$, while
the number of bonds or plaquettes in a cell of scale $\ell_y$ is
at most $3(\ell_y/\xi)^3$. The resulting kernel sum has factor
$\ell(c)/\ell_y$, which is absorbed by the exponential decay.
The same calculation with the first differences gives the fourth
and fifth estimates, including the zero-extension terms at the
boundary. These are equations~\YMMref{(459)}{235} and the
following divergence estimate. The scalar estimate follows from
$R=1-G'Q'^{\mathsf T}(Q'G'^2Q'^{\mathsf T})^{-1}Q'G'$.
At a bond meeting the complement of the ambient region, the
zero-extension difference has local scale $\xi$ and is bounded by
$|A|_{(1)}$. Thus these estimates also include boundary differences.
For an empty ambient region all spaces are zero. The operators
preserve real fields, so the estimates extend to complex fields by
their real and imaginary parts, with constants enlarged by a fixed
factor.

The same operators have exponential decay in the weighted distance.
Their first-difference H\"older estimates follow from the ordinary
kernel estimates in Lemma~\ref{lem:ordinary-holder-kernels}.
Multiplying the input at $z$ and output at $y$ by the envelope ratio
costs $h(z)/h(y)$. The permitted growth of this ratio consumes only
part of the exponential decay. The remaining kernel is summable by
Lemma~\YMMref{21.2}{141}. This gives the envelope versions of
\eqref{eq:endpoint-six-bounds}, with constants independent of the
torus and the number of layers.

For $q=f+D_1^{\mathsf T}Y$, elimination of the multiplier in
$TA=q+Q^{\mathsf T}\omega$, $QA=b$, gives the exact formulas
\begin{equation}\label{eq:endpoint-elimination}
 A=\widetilde Gq+Hb,\qquad
 \omega=H^{\mathsf T}(THb-q).
\end{equation}
Indeed $(T+Q^{\mathsf T}Q)A=q+Q^{\mathsf T}(\omega+b)$,
and applying $QG$ determines $\omega+b$. The identity
$M^{-1}=H^{\mathsf T}TH+1$ gives the displayed formula for
$\omega$. The six bounds control the ordinary forcing, the lift,
and all multiplier terms. The remaining issue is the first-slot
divergence $D_1^{\mathsf T}Y$.

Cover each layer of index at least two by cubes of size comparable
to its layer scale. In the local gauges, write
$U=e^{i\xi a}$, with $a$ and its first differences bounded in the
units of that layer. A partition of unity reduces the divergence
source to the flat equation on these cubes. Convolution with the
lattice Laplace Green kernel gives a local potential; its first
differences involve second differences of that kernel. These
preserve the H\"older exponent by Lemma~\YMMref{19.4}{115}.
The cancellation in that lemma is essential: a bound on the
absolute kernel alone would give a logarithmic loss.

Cut off these local potentials and carry them back to the original
gauge. Their sum $\Phi$ satisfies
\[
 T\Phi=D_1^{\mathsf T}Y+J,\qquad
 |J|_{3,h}\le C\|Y\|_{2,1/2,h}.
\]
Here $J$ consists of cutoff commutators, the terms containing the
local potential and curvature, the lower-order term $V_0'\Phi$,
the term $-D_0(1-R)D_0^{\mathsf T}\Phi$, and the untreated
sources in layers zero and one. The first differences of the local
gauge potential bound the coefficient terms; the cutoff derivatives
have the inverse powers of the local scale. In the first two layers
the scale-to-spacing ratio is bounded by $L$, so a bare difference
is bounded by the corresponding scaled supremum. These are
ordinary sources of degree three. The parametrix identity is
\begin{equation}\label{eq:endpoint-parametrix}
 \widetilde GD_1^{\mathsf T}Y
       =\Phi-HQ\Phi-\widetilde GJ.
\end{equation}
The row datum $Q\Phi$ is bounded in the dimensionless row norm:
each average reads a double block on which the scale and envelope
are comparable, and the factor $\ell(c)$ in the physical path
average cancels the scale $\ell^{-1}$ of $\Phi$.
The remainder $J$ has the degree-three norm displayed above.
Thus both correction terms lie in the ordinary-input classes of
Lemma~\ref{lem:ordinary-holder-kernels}. Its estimates are read
at $s=1$, with the same rows and their row pairing. This proves
the required potential and curvature bounds. The first-difference
H\"older bounds for the correction terms are the ordinary kernel
bounds; those for $\Phi$ are the same-exponent Green-kernel
bounds just used. Their supremum estimates also bound the
divergence. To pass from the comparison cubes
of radius $\ell/56$ to the cubes in the norm, cover the latter
by comparison cubes and use the local gauge transition bounds.
On an overlap those transitions have differences bounded by
$C(L)\alpha/\ell$; the H\"older product inequality and the
supremum bound therefore control the change of gauge. The
bounded overlap gives a constant depending only on $L$.
No H\"older estimate is
obtained here by applying an estimate with a positive loss of
exponent at zero loss.

For completeness, the potential H\"older norm can also be read
directly from its curvature and divergence. On a cube $C_R$ in
bare units, let the adjoint transports differ from the identity by
at most $v/R$. If $6v\le\sigma_H$, then
\[
 \sup_{C_{R/2-2}}|E|+(R/2-2)^{1/2}[E]_{1/2;C_{R/2-2}}
 \le 2C_H\bigl((R-1)^{-3/2}\|E\|_2
                 +R|DE|+R|D_0^*E|\bigr).
\]
To obtain it, apply Proposition~\YMMref{18.1}{102} to each
component, whose flat Laplacian is the divergence of its flat
curvature and divergence. Replacing the covariant quantities by
the flat ones adds at most $Cv$ times the supremum on the next
larger cube. On radii increasing from $3R/4$ by
$R/8,R/16,\ldots$, choose $v$ so that this factor is at most
$1/12$. Iteration absorbs the supremum; the terminal error is at
most $12^{8-\log_2R}\|E\|_2$. A second application on cubes
of radius $R/4$ gives the H\"older bound. For $R<64$ the
supremum is bounded by the finite-volume $\ell^2$ norm. At
$R=L^n/4$, the condition follows by reducing the regularity
radius so that $\alpha\le\sigma_H/(8L^2)$.

We have proved the estimate at zero anisotropy with a constant
$K^{(0)}$. At a nonzero anisotropy, the same solution solves the
unweighted system with data
\[
 (Y-{\rm an}\,cD_1A,\ f-V_0'[{\rm an}]A,\ b).
\]
The H\"older product inequality bounds the first change by
$Ct({\rm an})\|D_1A\|_{2,1/2,h}$; the second is bounded by
$C|{\rm an}|_\infty\alpha|A|_{1,h}$. Choose
$K^{(0)}Ct_{\rm lin}\le1/2$ and absorb these terms. The bound
holds with $K_{\rm lin}=2K^{(0)}$.

The real part of the form of $T_\tau$ on $N(Q)$ is coercive
at these radii, as proved below. Hence a homogeneous solution has
$A=0$, by pairing its equation with $\bar A$, and then
$Q^{\mathsf T}\omega=0$ gives $\omega=0$. The finite square
system is therefore invertible. This proves existence and
uniqueness as well as the estimate.

Finally, the mean-scheme rows of a label, rotated into the comb
frame, have the form $Q+E$. Their comparison kernel is supported
in the double block of the row and obeys
\[
 |EA|_{0,h}\le c_E\alpha|A|_{1,h},\qquad
 |E^{\mathsf T}\omega|_{3,h}\le c_E\alpha|\omega|_{0,h}.
\]
The system with these rows is the one just solved with data
$(Y,f+E^{\mathsf T}\omega,b-EA)$. Taking
$K_{\rm lin}c_E\alpha\le1/2$ proves its bound, with
$2K_{\rm lin}$, by a contraction. This is the row comparison
needed when the estimate is applied to a label.
\end{proof}

\begin{lemma}[The energy of the displacement]
\label{lem:weighted-field-energy}
In the setting of Lemma~\ref{C-360-22}, decrease the regularity
radius and the envelope ceiling by constants depending only on
$L$. Put
\[
 \mathcal E(A)=\xi^3\sum_p|D_1A(p)|^2
                +\xi^3\sum_b\ell(b)^{-2}|A(b)|^2,
 \qquad |b|_{\rm row}^2=\sum_c\ell(c)^{-1}|b(c)|^2.
\]
For the linear solution with data $(Y,f,b)$,
\begin{equation}\label{eq:weighted-field-energy-linear}
 \mathcal E(A)\le K^{(2)}\left(
       \xi^3\sum_p|Y(p)|^2+
       \xi^3\sum_b\ell(b)^2|f(b)|^2+|b|_{\rm row}^2\right).
\end{equation}
For the small fixed point with zero block datum,
\begin{equation}\label{eq:weighted-field-energy-nonlinear}
 \mathcal E(A[{\rm an};0])
      \le C|{\rm an}|_\infty^2\xi^3\sum_p|F_U(p)|^2.
\end{equation}
At a permitted complex block datum $b_0$, the difference
$A'=A[{\rm an};b_0]-A[0;b_0]$ instead satisfies
\begin{equation}\label{eq:weighted-field-energy-difference}
 \mathcal E(A')\le C|{\rm an}|_\infty^2
       \left(\xi^3\sum_p|F_U(p)|^2+|b_0|_{\rm row}^2\right).
\end{equation}
The constants are independent of the number of layers and the
volume. The last two statements concern fixed points in the small
ball of the linear estimate.
\end{lemma}

\begin{proof}
We need an estimate on the field energy, including the forcing from
the nonlinear block constraint. A contraction in a stronger energy
norm of the field and multiplier is unnecessary.

At zero anisotropy the operator $T+Q^{\mathsf T}Q$ is real,
symmetric and positive by Theorem~\YMMref{I}{225}. Its
restriction to $N(Q)$ is $T$, so the constrained Green operator
is positive on that subspace. Apply the first estimate of
\eqref{eq:endpoint-six-bounds} to an eigenvector of the largest
eigenvalue of $\ell^{-1}\widetilde G\ell^{-1}$. If that
eigenvalue is $\lambda$, it gives
\[
 \lambda\sup_b\ell(b)^2|g(b)|
       \le K_7\sup_b\ell(b)^2|g(b)|.
\]
Thus the quadratic form of $T$ on $N(Q)$ controls the weighted
potential square sum with constant $K_7^{-1}$. Its expression as
a curvature form, a slice term and $V_0'$ gives the second lower
bound. The plaquette coefficient satisfies $c_p\ge1/2$, while
the form of $V_0'$ is bounded by $C_V\alpha$ times the weighted
potential square sum. Since $|{\rm an}|_\infty\le1/4$, reducing
$\alpha$ so that $3C_VK_7\alpha\le1$ gives
\begin{equation}\label{eq:weighted-field-energy-floor}
 \operatorname{Re}(\bar A,T_\tau A)_\xi
    \ge c^{(2)}\mathcal E(A)+\tfrac12\|RD_0^*A\|_\xi^2,
 \qquad A\in N(Q),
\end{equation}
where $c^{(2)}=\min\{1/8,(4K_7)^{-1}\}$. Pairing the equation
with $\bar A$ removes the multiplier. Cauchy--Schwarz bounds the
remaining pairing by $\mathcal E(A)^{1/2}$ times the square root
of the first two data sums in
\eqref{eq:weighted-field-energy-linear}.

For nonzero row data use the lift $A^\circ=Hb$. Then
$A-A^\circ\in N(Q)$ and its equation has the additional forcing
$-T_\tau A^\circ$. The bounds on $H$, its first differences and
their transposes in \eqref{eq:endpoint-six-bounds}, applied by the
weighted Schur test, give
\[
 \mathcal E(A^\circ)+\|D_0^*A^\circ\|_\xi^2
                   \le C_H|b|_{\rm row}^2.
\]
The row weight is $\ell(c)^{-1}$; using an unweighted row sum
would not give a constant uniform in the layers. Apply
\eqref{eq:weighted-field-energy-floor} to $A-A^\circ$ and use
this lift estimate. This proves
\eqref{eq:weighted-field-energy-linear}.

At the nonlinear fixed point the data are
\[
 \bigl(-{\rm an}F_U-Y_3(A),\
       -f_3(A)+(dC_2(U,A))^{\mathsf T}(\omega_0+\omega),\
       -C_2(U,A)\bigr).
\]
Here $Y_3,f_3$ are the terms of order at least two in the
displacement in the first variation, and $C_2$ is the nonlinear
remainder of the block constraint. The last entry is not zero.
The small-ball bound on $A$ and the multiplier implies that the
square-sum norm of the nonlinear data is at most
$c_{\rm ms}\sup e\,\mathcal E(A)^{1/2}$.
To see the row estimate in the same units, write $l(c)=\ell(c)/\xi$.
The Taylor kernel of the block average and Cauchy--Schwarz give
\[
 |C_2(U,A)(c)|\le C l(c)^{-1}
                  \sum_{b\ {\rm beneath}\ c}|\xi A(b)|^2.
\]
There are at most $Cl(c)^3$ bonds beneath the row, their scales
are comparable to $l(c)$, and each bond lies beneath at most a
constant number of rows. Hence the row square sum is bounded by
$C(\sup_b\ell(b)|A(b)|)^2\mathcal E(A)$. The same bilinear
kernel bounds the transposed derivative paired with
$\omega_0+\omega$, whose scaled supremum is at most $C\sup e$.
The plaquette and bond remainders are treated by their local
product bounds, using one small factor and one energy factor.

Choose the envelope ceiling so that
$K^{(2)}c_{\rm ms}^2(\sup e)^2\le1/16$.
The nonlinear data can then be absorbed in
\eqref{eq:weighted-field-energy-linear}. The remaining forcing
is $-{\rm an}F_U$, which proves
\eqref{eq:weighted-field-energy-nonlinear}.

At ${\rm an}=0$ and datum $b_0$ the same argument gives
$\mathcal E(A[0;b_0])\le4K^{(2)}|b_0|_{\rm row}^2$.
Subtract the two fixed-point equations at the same $b_0$.
The terms linear in ${\rm an}$ are bounded by its supremum times
the square root of
$\xi^3\sum|F_U|^2+\mathcal E(A[0;b_0])$.
The remaining differences have the factor $C\sup e$; the
multiplier difference is bounded by testing against the same row
lift. One further reduction of the envelope ceiling absorbs
these differences. This proves
\eqref{eq:weighted-field-energy-difference}.
\end{proof}


\begin{definition}[Complex scalar gauge transfer]\label{hyp:complex-scalar-transfer}
Consider a small complex potential $a$ and the small gauge
$e^{i\lambda}$ relative to its real reference configuration. Define
$N$ and $q$ by the exact identities
\[
 \begin{split}
 D^*a^{e^{i\lambda}}&=D^*a-D^*D\lambda+N(a,\lambda),\\
 T^{-1}i^{-1}\log\mathcal R(a;e^{i\lambda})
        &=Q'\lambda+q(a,\lambda).
 \end{split}
\]
Here $\mathcal R$ is the induced gauge at the retained block tops,
and $T$ is the orthogonal transport in its linearization; these are
the maps in the proof of Lemma~\YMMref{30.4}{242}. Their linear
scalar system is precisely the compression $\mathbb K^{[2]}$ of
Section~\YMMref{23.2}{179}, in the covariant form of
Section~\YMMref{25}{189}:
\[
 \mathbb K^{[2]}(\lambda,w,\sigma)
  =(\Delta'_a\lambda+w,\ \Delta'_aw-Q'^*\sigma,\ Q'\lambda),
 \qquad \Delta'_a=\Delta_D+Q'^*\mathsf a'Q'.
\]
Here $\Delta_D$ is the scalar Dirichlet Laplacian; the boundary
identity is $RD^*D=R\Delta_D$. Let $\mathcal S(f,b)$ be the first
component of the solution with data
$(f+Q'^*\mathsf a'b,0,b)$. Its first and third rows give
$w=f-\Delta_D\lambda$, while its second row says $Rw=0$.
Consequently the two desired gauge conditions are exactly the fixed
point $\lambda=\mathcal S(D^*a+N(a,\lambda),-q(a,\lambda))$.
The mass term $Q'^*\mathsf a'b$ is necessary when $b\ne0$.
In the norm
$\|\lambda\|_{\mathcal Z_h}
 :=|\lambda|_{0,h}+\|(D\lambda,0)\|_{\mathcal X_h}$, assume
\begin{equation}\label{eq:complex-scalar-transfer}
 \begin{split}
 \|\mathcal S(N(a,\lambda)-N(a,\tilde\lambda),
                  -q(a,\lambda)+q(a,\tilde\lambda))\|_{\mathcal Z_h}
       &\le\tfrac12\|\lambda-\tilde\lambda\|_{\mathcal Z_h},\\
 \|(a^{e^{i\lambda}}-a,0)\|_{\mathcal X_h}
       &\le C\|\lambda\|_{\mathcal Z_h}.
 \end{split}
\end{equation}
These bounds are required on one small complex ball, invariant under
the scalar iteration, for each admitted envelope $h$, including the
exponentially decaying envelopes used for a cutoff repair. They
include the same bounds after subtracting a given solution and
localizing its scalar system. Thus the initial forcing in a cutoff
repair is the commutator forcing supported where the cutoff varies.
The scalar inverse itself is supplied by Theorem~\YMMref{G}{204};
\eqref{eq:complex-scalar-transfer} is the additional uniform complex
nonlinear estimate, not a claim about that theorem.

\end{definition}

\begin{lemma}[The re-filing of the counterterms]\label{C-360-23}
Let $1\le j\le K$, $m'<m+K-j$, $b_0,b_1,\kappa'>0$ and let $c'$ be a complex function of the plaquettes. For a site $z$ of the level $j$ let $\square_z$ be the cube of $\pi_j$ (of side $M$ in level-$j$ units) that holds $z$, and $X_z=\square_z^\sim$ the block of the $3\times3\times3$ cubes of $\pi_j$ about it, the box of the site $z$: $X_z\in\mathcal D_j$, $z\in X_z$, $d_j(X_z)\le7$ for the linear size $d_j$ of \YMMref{72}{607} (the definitions of $\mathcal D_j$ and $d_j$ at its beginning) (Section~\YMMref{72}{607} (the definitions of $\mathcal D_j$ and $d_j$ at its beginning)), every plaquette $p$ with $h^{(j)}_z(p)\ne0$ has its four corners in $X_z$ and its four bonds in ${\rm Bd}(X_z)$, and at most $M^3$ sites $z'$ have $X_{z'}=X_z$, namely those of $\square_z$ (the tent domains and support assertions are those preceding equation~\YMMref{(762)}{376} and in Lemma~\YMMref{49.7}{376}).

Here $h^{(j)}_z$ is the tent of the level $j$ at $z$, the product of the hat functions $\max\{1-|u_\mu|,0\}$ of the coordinates in level-$j$ units, read at the barycentres (equation~\YMMref{(760)}{375}): $0\le h^{(j)}_z\le1$, $h^{(j)}_z(y)\ne0$ exactly when $|y-z|_\infty<1$ in level-$j$ units, $\sum_zh^{(j)}_z=1$, and $\sum_ph^{(j)}_z(p)=3\eta_j^{-3}$ (Proposition~\YMMref{49.6}{375} and the plaquette-count calculation in the proof of Lemma~\YMMref{49.7}{376}). For a complex weight $c$ and a complex bare configuration $\mathbf U$ put $A^\eta(c,\mathbf U)=\eta^{-1}\sum_pc(p)a_p(\mathbf U)$ (equations~\YMMref{(6)--(7)}{13}), and for a triple $P=(U,A',J)$ of a level $q$ let $\mathbf U_P=e^{i\eta_qA'}U$ be its configuration. The tent family at the weight $c'$,
\[
\begin{multlined}
A^t[c'](X,P,z)=A^{\eta_j}\bigl(c'h^{(j)}_z,\mathbf U_P\bigr)=\eta_j^{-1}\sum_pc'(p)h^{(j)}_z(p)a_p(\mathbf U_P)\ \relax \ \relax \text{if }X=X_z,\qquad A^t[c'](X,P, \\
z)=0\ \relax \ \relax \text{for the other }X\in\mathcal D_j,
\end{multlined}
\]
$z\in X$ a site of the level $j$, is, on the cores $\mathcal U^\circ_j(X;b_0,b_1)$, a pointed family of the level $j$ at $(b_0,b_1)$ (Definitions~\YMMref{49.4}{375} and~\YMMref{49.5}{375}), complex linear in $c'$ and $0$ at the trivial triple; its size \YMMref{(759)}{375}, the supremum of $e^{\kappa'd_j(X)}|A^t[c'](X,P,z)|$ over the $X\in\mathcal D_j$, the sites $z\in X$ and the triples $P$ of the core, obeys $\|A^t[c']\|_{j,\kappa'}\le\frac32|c'|_\infty b_0^2e^{7\kappa'}$; at $c'=1$ it is the tent family of equation~\YMMref{(762)}{376}; and $A^{\eta_j}(c',\mathbf U)=\sum_zA^{\eta_j}(c'h^{(j)}_z,\mathbf U)$ for every complex $\mathbf U$. Consequently, let $1\le n\le K$ with $m'<m+K-n$, and let $E^{(i)}$, $1\le i\le n$, be pointed families of the levels $i$ at $(\alpha_0,\alpha_1)$ with complex counterterms $e_i$, let ${\rm an}$ be a complex weight with $|{\rm an}|_\infty=s$, and $\hat E^{(i)}=E^{(i)}-e_iA^t_i[{\rm an}]$, $A^t_i$ the tent family of the level $i$; let $P$ be a triple of one level $k\ge n$ on the whole torus whose restrictions $\rho_XP=\rho^{k\to i}_XP$ to the domains $X$ of the levels $i\le n$ lie in the cores; since $\eta_iL^{-(k-i)}=\eta_k$, the configuration of $\rho_XP$ is $\mathbf U_P$ on $X$. Then, identically in $P$,
\[
\begin{multlined}
\sum_i\Bigl[\sum_z\sum_{X\ni z}\bigl(E^{(i)}(X,\rho_XP,z)-E^{(i)}(X,\mathbf 1,z)\bigr)-e_iA^{\eta_i}(1+{\rm an},\mathbf U_P)\Bigr]=\sum_i\Bigl[\sum_z\sum_{X\ni z}\bigl(\hat E^{(i)}(X, \\
\rho_XP,z)-\hat E^{(i)}(X,\mathbf 1,z)\bigr)-e_iA^{\eta_i}(1,\mathbf U_P)\Bigr],
\end{multlined}
\]
$\hat E^{(i)}(X,\mathbf 1,z)=E^{(i)}(X,\mathbf 1,z)$, and $\|\hat E^{(i)}-E^{(i)}\|_{i,\kappa'}\le\frac32s|e_i|\alpha_0^2e^{7\kappa'}$.

On the label $\mathcal S$ of a history $h\in\mathcal H_k$ with regular regions $\Lambda_j$, assume $m'<m+K-j$ at every level used. The same identity then holds in the following form. Let $\varphi_j$ be the profile of $\Lambda_j$ (equation~\YMMref{(666)}{349} and the definition of $\varphi_j$ following it): $0\le\varphi_j\le1$, $\varphi_jh^{(j)}_z=h^{(j)}_z$ at the sites $z$ of $\Lambda_j^{\sim-1}$, which is $\Lambda_j$ less one layer of its cells, and $\varphi_j(x(p))\ne0$ only if all the corners of $p$ lie in $\Lambda_j$ (the tent identity in Section~\YMMref{44.4}{349}; the corner assertion also follows since a nonzero profile is more than one quarter of a cell width from $Z_j$, while every plaquette corner is less than one bare spacing from its center). The exponent on the label (equation~\YMMref{(1476)}{600}) holds, besides its Wilson term and its vacuum constant, the regular sums in equation~\YMMref{(1474)}{600}, over the set $\mathsf{Ix}^{\rm rs}_j(\mathcal S)$ of the places $(z,X)$ with $X\subset\Lambda_j$ and $z\in\Lambda_j^{\sim-1}$, with the counterterms $-e_jA^{\eta_j}(\varphi_j,\cdot)$, and the boundary sums in equation~\YMMref{(1475)}{600}, which hold the made boundary entries $\mathscr B(j,X,{\rm mk})$ at the filing places $X\in\mathsf{Ix}^{\rm mk}_j(\mathcal S)$ and the edge parts of equation~\YMMref{(1388)}{571} of the families, the members with base points outside $\Lambda_j^{\sim-1}$, at the places $X\in\mathsf{Ix}^{\rm ed}_j(\mathcal S)$ (the index sets and domains preceding Proposition~\YMMref{62.2}{571}); each member is filed at exactly one place (by the separate level, domain and origin indexing in Proposition~\YMMref{62.2}{571}); the entries are boundary terms of their filing places, additive in the exponent, of the sizes $\le\mathsf B\,e^{-\kappa d_j(X)}$ of a well filed system (Proposition~\YMMref{62.2}{571}, on its stated layered output spaces); with the anisotropy the counterterms are $-e_jA^{\eta_j}((1+{\rm an})\varphi_j,\mathbf U_P)$.

For a site $y$ of the level $j$ outside $\Lambda_j^{\sim-1}$ put $\epsilon_y=1$ if $X_y\subset\Lambda_j$ and $\epsilon_y=0$ otherwise, and
\[
r_{(j,y)}=-e_j\eta_j^{-1}\sum_p{\rm an}(p)\bigl[\varphi_j(x(p))-\epsilon_y\bigr]h^{(j)}_y(p)\,a_p(\mathbf U_P).
\]
Then the exponent formed with the families $E^{(j)}$, the entries $\mathscr B$ and the counterterms $-e_jA^{\eta_j}((1+{\rm an})\varphi_j,\mathbf U_P)$ equals, identically, the exponent formed with the families $\hat E^{(j)}$ at the regular places and in the edge parts, the entries $\mathscr B(j,X,{\rm mk})+\sum\{r_{(j,y)}:y\notin\Lambda_j^{\sim-1},\ X_y=X\}$ and the counterterms $-e_jA^{\eta_j}(\varphi_j,\mathbf U_P)$. Here $r_{(j,y)}=0$ unless the tent of $y$ meets the set where $\varphi_j\ne\epsilon_y$, a collar of the boundary of $\Lambda_j$, and then $X_y\in\mathsf{Ix}^{\rm mk}_j(\mathcal S)$; $r_{(j,y)}$ is linear in ${\rm an}$, reads $\mathbf U_P$ in $X_y$ alone and is $0$ at the trivial triple; and on the space of the filing place $(j,X_y,{\rm mk})$, namely the layered output space at level $j$ specified before Proposition~\YMMref{62.2}{571}, where the plaquettes concerned have their plaquette variables within $r_0\eta_j^2$ of $\mathbf 1$, $r_0=\mathsf r_0((j,{\rm out});j)$ being the prescribed output-layer curvature radius at level $j$ (the same radius used for the edge-tent estimate in the proof of Proposition~\YMMref{58.22}{541}), $|r_{(j,y)}|\le\frac32s|e_j|r_0^2$; in a frame of Theorems~\YMMref{AJ}{673} and~\YMMref{AD}{607} the radius of the level is at most $\alpha_0/(1+\beta_{\rm s})<\alpha_0$. The values at the trivial triple are unchanged.
\end{lemma}

\begin{proof}
The room condition gives $2L^{m+K-j-m'}\ge2L$ cubes of
$\pi_j$ in each coordinate direction. The three consecutive cubes
in each projection of $X_z$ therefore determine its middle cube.
Thus $X_{z'}=X_z$ exactly for the $M^3$ level-$j$ sites
$z'\in\square_z$.

On the cores $|\partial\mathbf U_P(p)-\mathbf 1|<b_0\eta_j^2$ at the plaquettes with their bonds in ${\rm Bd}(X_z)$, and $a_p=\frac12\det(\partial\mathbf U_P(p)-\mathbf 1)$ for a matrix of determinant $1$, so $|a_p|\le\frac12b_0^2\eta_j^4$; with $h^{(j)}_z\ge0$ and $\sum_ph^{(j)}_z(p)=3\eta_j^{-3}$ this gives $|A^t[c'](X_z,P,z)|\le\eta_j^{-1}|c'|_\infty\cdot3\eta_j^{-3}\cdot\frac12b_0^2\eta_j^4=\frac32|c'|_\infty b_0^2$, and $d_j(X_z)\le7$ gives the factor $e^{7\kappa'}$. Each member is a polynomial in the entries of the $U(b)^{\pm1}$ and $e^{\pm i\eta_jA'(b)}$ at the bonds of ${\rm Bd}(X_z)$, entire in $(A',J)$, a function of the pair, gauge invariant, and $0$ at the trivial triple (Lemma~\YMMref{49.7}{376}).

The tents form a partition of unity, so $A^{\eta_j}(c',\mathbf U)=\sum_zA^{\eta_j}(c'h^{(j)}_z,\mathbf U)$; at $\mathbf U=\mathbf U_P$ the counterterm $e_iA^{\eta_i}({\rm an},\mathbf U_P)$ is the sum over the sites $z$ of the members $e_iA^t_i[{\rm an}](X_z,\rho_{X_z}P,z)$, which vanish at the trivial triple, and moving each from the counterterm into the family of its box is the display, term by term: the two members differ by $\sum_ie_i\eta_i^{-1}\sum_p{\rm an}(p)[a_p(\mathbf U_P)-a_p(\mathbf U_P)]=0$, the counterterm and the families being read at the one configuration $\mathbf U_P$. On a label, by the partition of unity and $\varphi_jh^{(j)}_z=h^{(j)}_z$ on $\Lambda_j^{\sim-1}$, for every rule $\epsilon_y\in\{0,1\}$,
\[
A^{\eta_j}({\rm an}\,\varphi_j,\mathbf U)=\sum_{z\in\Lambda_j^{\sim-1}}A^{\eta_j}({\rm an}\,h^{(j)}_z,\mathbf U)+\sum_{y\notin\Lambda_j^{\sim-1}}\epsilon_yA^{\eta_j}({\rm an}\,h^{(j)}_y,\mathbf U)+\sum_{y\notin\Lambda_j^{\sim-1}}A^{\eta_j}\bigl({\rm an}\,(\varphi_j-\epsilon_y)h^{(j)}_y,\mathbf U\bigr).
\]
For $z\in\Lambda_j^{\sim-1}$ the place $(z,X_z)$ is a regular place, and for $y\notin\Lambda_j^{\sim-1}$ the domain $X_y$ is an edge place exactly when $X_y\subset\Lambda_j$; the exponent with $\hat E^{(j)}$ holds $-e_jA^t_j[{\rm an}]$ at the regular places and in the edge parts, which are the first sum and the members of the second with $\epsilon_y=1$: this is why $\epsilon_y=1$ at the edge places and $\epsilon_y=0$ elsewhere, and $-e_j$ times the third sum is $\sum_yr_{(j,y)}$, the contribution to the made boundary entries. This identity also holds at $e_j=0$. If ${\rm an}(p)[\varphi_j(x(p))-\epsilon_y]h^{(j)}_y(p)\ne0$, the corners of $p$ lie in $X_y$, and in $\Lambda_j$ (by the profile corner property just stated if $\varphi_j(x(p))\ne0$; and if $\varphi_j(x(p))=0$, then $\epsilon_y=1$ and $X_y\subset\Lambda_j$); as $\Lambda_j\subset\Omega_j$, its bonds lie in the bonds of the label in $X_y$, and its base point, a corner, is a site of $X_y$ and of $\Omega_j$, so that $X_y$ and $\Omega_j$ have a bare cube in common; and as $y\in X_y$ lies outside $\Lambda_j^{\sim-1}$, $X_y\in\mathsf{Ix}^{\rm mk}_j(\mathcal S)$: this index set consists precisely of domains sharing a bare cube with $\Omega_j$ and not contained in $\Lambda_j^{\sim-1}$, as defined before Proposition~\YMMref{62.2}{571}.

The size: $|\varphi_j-\epsilon_y|\le1$, the mass $3\eta_j^{-3}$ of the tent and $|a_p|\le\frac12r_0^2\eta_j^4$ on the space of the filing place give $\frac32s|e_j|r_0^2$; in a frame of Theorems~\YMMref{AJ}{673} and~\YMMref{AD}{607}, $r_0\le\alpha_0/(1+\beta_{\rm s})$. The tent families and $\varphi_j$, $\epsilon_y$ read no history beyond the label, and the values at the trivial triple are unchanged, the tent families vanishing there.
\end{proof}

\subsection{Proofs of the theorems}

\begin{proof}[Proof of Theorems~\ref{C-360-2}--\ref{C-360-5}, \ref{C-360-13} and \ref{C-360-14}, and of the remarks]
\emph{Step 1: Theorem~\ref{C-360-2}(a)--(c).} (a) is Lemma~\ref{C-360-6} for the rough entries $F_a=\mathcal O(w_a)$, and its second bound is the second decomposition in the sketch of that lemma, applied to the weights $u_a\mathbf 1\{\lambda(\cdot)=\lambda_a\}$; for a weight that vanishes the cumulant is $0$. (b) Let $\ell'=2\mathsf l$, $\mathsf d=\min\{\mathsf d_0,\frac{\mathsf l}{4n},\frac12\}$, and let $K_1$ be the least $K\ge K_0(r,\mathsf d)$, $2\le r\le8n$, with $a_K\le\frac{\mathsf l}{4n}$. For $K\ge K_1$ cut each $f_i$ into at most $(\frac R\rho+2)^3$ pieces in balls of radius $\rho=\frac{\mathsf l}{8n}$; each choice of pieces lies, for a suitable offset (a pigeonhole on the $n$ centres), in cells of a board of side $\mathsf l$ at sup distance $\ge\frac{\mathsf d}2$ from the walls, so that the mirror copies stay $\ge\mathsf d$ apart, and Lemma~\ref{C-360-7} with (a) gives
\[
\Bigl\|\prod_{a\le r}(\mathcal O(w_a)-c_a)\Bigr\|_{\rm ch}\le e^{G_g\mathsf l^3}\Bigl[\sum_\pi\prod_{I\in\pi}k(\#I,\mathsf d)\Bigr]^{1/8}\prod_a\sup|w_a|,
\]
$\pi$ over the partitions of $8r$ points into blocks of at least two; the moment--cumulant formula, $\langle F_1;\dots;F_n\rangle=\sum_\pi(-1)^{|\pi|-1}(|\pi|-1)!\prod_{B\in\pi}\langle\prod_{a\in B}(F_a-c_a)\rangle$, gives a bound $T_n(\mathsf d)$. For $K<K_1$ we use $|\mathcal O(f)|\le m_K(R)\sup|f|$, $m_K(R)=6a_K^{-1}(\frac{2R}{a_K}+1)^3$, increasing in $K$; so $C_n=\max\{(\frac R\rho+2)^{3n}T_n(\mathsf d),\,m_{K_1}(R)^n\sum_\pi(\#\pi-1)!\}$. On $T_\ell$ the same board is placed with $\ell\ge L\ell'$, and a periodic Gibbs state is dominated through its conditional laws (Lemma~\ref{C-360-7}). The bound on $\nu'-\nu$ is the case $n=1$ with the centring at the mean of the torus. (c) The bounds of (b) give equicontinuity in the sup norm, and a diagonal argument over a countable dense family of tuples gives a convergent subsequence; the lattice symmetries and the reflection positivity of the states pass to the limit, at the three coordinate mirrors in both cases and, for periodic Gibbs states, also at the six diagonal mirrors; the translations by all $v\in\mathbb R^3$ follow through $v_j\in a_{K_j}\mathbb Z^3$, $v_j\to v$, and the Lipschitz continuity of the limits; the growth in $n$ follows by cutting into $\lceil\frac{2\mathsf l}{\mathsf d_0}\rceil^3$ cells and counting, with $B=2e^{3+p}(\frac{8R}{\mathsf l}+2)^3\mathsf d_*^{-p}\max_{r\le P}[e^{G_g\mathsf l^3}{\rm Bl}_{8r}^{1/8}]^{1/r}$, $\mathsf d_*=\min\{\mathsf d_0,\frac{\mathsf l}4,\frac12\}$, $P=\lceil\frac{2\mathsf l}{\mathsf d_0}\rceil^3$; the sum over the planes is the scalar cumulant by linearity.

\emph{Step 2: Theorem~\ref{C-360-2}(d).} Put $\sigma=\lfloor K/(1260n)\rfloor$. (i) Let $s=\lfloor\sigma/(8n+3)\rfloor$ be the coarse depth, $T=K-s$, and $\nu=\sigma-s\ge(8n+2)s$ the number of the levels between the fine depth $\sigma$ and the coarse one. By forgetting, with the never-large term and the sources switched on at the level $K-\sigma$, each $\mathcal O(w_a)$ is exchanged for the weighted energy, plane by plane, of the background over the block field $\mathcal M^{K-\sigma}U$: what a source creates at the levels between $K-\sigma$ and $T$ is forgotten by $L^{-\nu/2}$ over these $\nu$ levels. For the $n$ sources, at mutual corner distances $\ge\mathsf d'=\frac12\mathsf d$, the cumulants of the $\mathcal O(w_a)$ and of the energies differ by at most
\[
\mathsf K_n\bigl[2L^{4n}\mathsf V_\ell(b)L^{-(b+1)s}+n\,C_{\rm fg}(\nu+1)^{2\mathfrak p_0+2}L^{-\nu/2}+2e^{3-L^s\mathsf d'/M}\bigr]\prod_ay_a,\qquad y_a\le4\ell^3L^{4s}\sup|w_a|,
\]
for every real $b\ge4n$, with numbers $\mathsf K_n$, $C_{\rm fg}$ and $\mathsf V_\ell(b)$ that read no $K$: the sup norms of the sources are paid at the coarse depth, which brings the factor $L^{4ns}$, and the histories with a first creation of a large field up to the level $T$ bring $\mathsf V_\ell(b)L^{-(b+1)s}$; the member $2e^{3-L^s\mathsf d'/M}$ is the factor of the distance of the mixed parts at the level $T$ (Lemma~\ref{C-360-24}(a), (e)), read under the conditions $L^s\mathsf d'\ge3M$ and $\mathsf h_2(s)\le1$ of Lemma~\ref{C-360-24}(e) at $t=T$, which hold for $K\ge K_1$, since $s$ grows with $K$ and $\mathsf h_2(s)\to0$ as $s\to\infty$. At $b=4n+1$, $L^{4ns}L^{-(4n+1)s}=L^{-s}$; and $L^{4ns}L^{-\nu/2}\le L^{-s}$ since $\nu\ge(8n+2)s$; so this error is at most $C\ell^{3n}K^{2\mathfrak p_0+2}L^{-s}\prod_a\sup|w_a|$, the factor $\ell^{3n}$ of $\ell^{3(n+1)}$ in the statement coming from $\prod_ay_a$ (the further $\ell^3$ comes from (iii)), and $s\ge\theta_nK-2$. (ii) For one block field at the level $t$ and the cutoffs $K$, $K+1$ at equal label the plane-wise energies of the two backgrounds, formed with the coefficient function $\mathsf c$ of the weights, differ per block site by $Ct^{1/2}L^{-t/2}(\sup|\mathsf c|+L^{t-K}{\rm Lip}_\infty\mathsf c)$.

For the sum of the three planes this is equation~\YMMref{(2156)}{795}. For one plane, retain the parallel-plaquette support in the proof of Proposition~\YMMref{15.1}{79}. More precisely, for a set $S$ of coarse plaquettes and a finer real field $W$ with $\varpi=\varpi_{\mathcal N(S)}(W)\le\varpi'_L$,
\[
\Bigl(\sum_{p'\in S}a_{p'}(\mathcal M(W))\Bigr)^{1/2}\le L^{1/2}\Bigl(\sum_{p\in\mathcal N^{\mathsf W}(S)}a_p(W)\Bigr)^{1/2}+\sqrt{12}\,C^{\rm SO}_0\,\varpi\Bigl(\sum_{p\in\mathcal N(S)}a_p(W)\Bigr)^{1/2},
\]
where $\mathcal N(S)$ consists of the fine plaquettes in the cones of $S$, and $\mathcal N^{\mathsf W}(S)$ consists of those parallel to a corresponding plaquette of $S$. To verify the estimate, put $f_W(p)=|\partial W(p)-1|$. Lemma~\YMMref{14.9}{68} bounds $f_{\mathcal M(W)}(p')$ by $(\mathsf Wf_W)(p')+C^{\rm SO}_0\varpi_{p'}^2$. The kernel in Lemma~\YMMref{14.8}{68} uses only parallel plaquettes, so its squared-norm bound is $L\|f_W\|^2_{\mathcal N^{\mathsf W}(S)}$. Equation~\YMMref{(146)}{80} gives
\[
 \sum_{p'\in S}\varpi_{p'}^4
 \le12\varpi^2\|f_W\|^2_{\mathcal N(S)}.
\]
The triangle inequality and $f_W(p)^2=2a_p(W)$ give the displayed energy bound. Applied to the plaquettes of one plane at which the coefficient exceeds a level $\lambda$, and integrated over $\lambda$, this bounds the plane-wise energy of the block average by that of the finer background up to the factor $1+C_{\rm en}\varpi$ and up to the variation of the coefficient over a block; this is the one place where the Lipschitz constants enter, with the factor $L^{-\sigma}$. (iii) The laws of the block fields of depth at most $\lceil K/1260\rceil$ at the two cutoffs at equal label differ, on every bounded Borel gauge invariant function $F$, by at most $C_{\rm bf}\ell^3K^{p'}L^{-K/10}\sup|F|$: this is Corollary~\textup{\ref{IV-62-66}}(a), \textup{(\ref{IV-62-68})}, with the number $C_{\rm bf}$ of Definition~\textup{\ref{IV-62-62}}(5) and the exponent $p'=2\mathfrak p_0+\frac12$, the exponent defined before Theorem~\YMMref{AR}{810} at $n=0$.

With $\sigma=\lfloor K/(1260n)\rfloor$ and $s=\lfloor\sigma/(8n+3)\rfloor$ the three errors are at most $L^{-s}\le L^2L^{-\theta_nK}$, $L^{4n\sigma}L^{-(K-\sigma)/2}\le L^{-K/3}$ and $L^{4n\sigma}L^{-K/10}\le L^{-K/11}$, which gives $\theta_n$ and the exponent $\max\{2\mathfrak p_0+2,p',1\}$ of $K$; it equals $2\mathfrak p_0+2$, since $p'=2\mathfrak p_0+\frac12$ and $\mathfrak p_0\ge1$. The comparison of two lattices at equal label, Theorem~\YMMref{AR}{810} in the setting (E) of Section~\ref{sec:cutoff-calibration}, applies: it asks $\mathsf b\ge\gamma^{-2}$, which holds because $\mathsf b>g^{-2}+\bar B>\gamma^{-2}$, and $MR(\mathsf b)\le L^{m-1}$, which holds because $L^{m_{\mathfrak L}(g)}\ge LM\,R(L(g^{-2}+\bar B))$ by the definition of the sides (equation~\YMMref{(1528)}{615}), $\mathsf b\le L(g^{-2}+\bar B)$ and $R$ increases.

For $\mathsf b\ge2\max\{\gamma^{-2},x_{\rm J}\}$ and $g=(\mathsf b/2)^{-1/2}$ we have $g^{-2}=\frac{\mathsf b}2\ge x_{\rm J}$ and $\mathsf b\in I(g)$. On a compact interval of labels the numbers are monotone in $\mathsf b$; the dilation is the change of units $x\mapsto L^zx$, which carries $\mu_{K,\ell,\beta}$ to $\mu_{K-z,L^z\ell,\beta}$ and multiplies the label by $L^z$. For continuous $\phi_a$ with supports at sup distances $\ge2\mathsf d$, uniform approximation by Lipschitz functions with supports at distances $\ge\mathsf d$ and the bound of (b), uniform in $K$, make $(G_K)$ a Cauchy sequence. With one plane field among action densities, (ii) is needed for the plane field alone, the action densities being compared at the two cutoffs by the comparison of two lattices at equal label (Theorem~\YMMref{AR}{810}).

The plane field is a connected function at both cutoffs: with the tent $\theta_K(y)=\max\{0,1-(2/a_{K+1}){\rm dist}_\infty(y,{\rm Bc}_\lambda)\}$ about the barycentres ${\rm Bc}_\lambda$ of the plaquettes normal to $\lambda$ at the cutoff $K+1$, $k_0=k\,\theta_K$ is Lipschitz, ${\rm supp}\,k_0\subset{\rm supp}\,k$, $\sup|k_0|\le\sup|k|$, and $k_0(x(p))=k(x(p))\mathbf 1\{\lambda(p)=\lambda\}$ at the plaquettes of both cutoffs, $L$ being odd; so $\langle\mathcal O_\lambda(k);\mathcal O(f_\cdot)\rangle=S^T_{n+1}(f_1,\dots,f_n,k_0)$ at both cutoffs. The comparison statements on two cutoffs at equal label (Theorem~\YMMref{AR}{810} with the lemmas on which it rests, Definition~\textup{\ref{IV-62-62}} and Corollary~\textup{\ref{IV-62-66}}) hold, with the same proofs, at any two depths $u<s$ in the places of the depths there: their proofs read of the depths only conditions on the cutoff, which hold for $K\ge K_0$ when $u$ and $s$ are fixed fractions of $K$.

Take $u=\lceil K/(80(3n+8))\rceil$, at whose level $t_1=K-u$ the sources are carried by the run, and $s=\lceil K/20\rceil$, at whose level $K-s$ the plane field is compared. Three members are left, each with the factor $L^{3nu}$ of Cauchy's estimates in the $n$ sources at the level $t_1$. What the plane field makes below the level $K-s$ is forgotten over the $s-u$ levels up to $t_1$, by $L^{-(s-u)}$, against the volume $L^{3u}$ of the level $t_1$ and the size $L^{s/2}$ of its first variation at the level $K-s$: the exponent $(3n+4)u-\frac12s\le-\frac1{80}K+3n+4$. The forgetting of the run between the two depths, against the size of the first variation at the level $t_1$, gives the exponent $(3n+\frac{15}2)u-s\le-\frac3{80}K+3n+8$.

The classical comparison of the plane-wise energies of the two backgrounds at the level $K-s$, as in (ii), with the factor $K^{1/2}L^{-(K-s)/2}(\sup|k|+L^{-s}{\rm Lip}_\infty k)$ against the volume $L^{3s}$, bounds the difference of the two first variations at the real points of the polydisc of the sources; at its complex points the first variations are bounded crudely, and the estimate with two constants for holomorphic functions, with the exponents $\frac34$ and $\frac14$, gives the exponent $3nu+\frac{23}4s-\frac38K\le-\frac6{80}K+3n+6$. The histories with a first creation of a large field up to the level $t_1$, and the mixed derivatives of the constants of the run in the sources, which are small by locality, cost less than every power of $L^{-K}$.

The largest of the three exponents gives the rate $L^{-K/80}$, and the powers of $K$ of Cauchy's estimates in the $n$ sources give $K^{2\mathfrak p_0(n+2)+2}$. For the sharpness: $L$ being odd, the barycentres of $a_K\mathbb Z^3$ are barycentres of $a_{K+1}\mathbb Z^3$ in the same plane, of density $L^{-3}$ among them; so the indicator of the points of a ball that are barycentres at an even cutoff and at no smaller cutoff samples to a weight of density close to $1$ in the ball at the even cutoffs and close to $L^{-3}$ at the odd ones, and for suitable smooth spectators the cumulants along the even and along the odd cutoffs have different limits.

\emph{Step 3: Theorem~\ref{C-360-2}(e)--(g).} (e) $\mathcal O_\lambda(f)=\mathcal O(f\chi_\lambda)$, $\chi_\lambda$ the indicator of the plaquettes normal to $\lambda$. Proposition~\YMMref{126.4}{846}, the two-point law on a torus of small side against the lattice kernel, is stated for every test pair of $\mathfrak T_{T_\ell}(\mathsf d;\vartheta)$ and every shape parameter $\vartheta\in(0,1)$, with a tolerance $\epsilon_g(s(\mathsf d))$ that depends neither on $\vartheta$ nor on $a_K/\mathsf d$. A pair of non-negative Lipschitz functions, neither identically $0$, carried by balls of radius $\mathsf d$ whose supports are between $\mathsf d$ and $6\mathsf d$ apart, is a test pair for every $\vartheta\in(0,1)$ below four positive numbers formed from their masses, their sup norms and their Lipschitz constants (Lemma~\ref{lem:ward-test-pairs}).

For non-negative weights $w_1,w_2$ carried by balls of radius $\mathsf d-a_K$ about $x_0,y_0$ with $3\mathsf d\le|x_0-y_0|\le4\mathsf d$, the functions $\sum_pw_a(p)\,{\rm pk}_p$, ${\rm pk}_p$ the tent of height $1$ and sup radius $\frac12a_K$ at $x(p)$, are Lipschitz, take the values $w_a$ at the barycentres, so that $\mathcal O(\sum_pw_a(p){\rm pk}_p)=\mathcal O(w_a)$, and form such a pair; so Proposition~\YMMref{126.4}{846} applies to them with its tolerance; real weights are differences of two non-negative ones (four terms), complex weights have real and imaginary parts (sixteen terms). This applies to the weights $f\chi_\lambda$, $h\chi_{\lambda'}$, which depend on the plane, and the kernel bounds of equations~\YMMref{(2245)--(2246)}{824}, with the bounds that enter the template sums of Proposition~\YMMref{121.2}{824} and Lemma~\ref{lem:ward-fluctuation-template}, compare the lattice kernel with $\mathbb K_3$ plane by plane at the Lipschitz pair $(f,h)$.

The lower bound on $\mathcal K_{\lambda\lambda'}$ follows from the masses and the Lipschitz constants of test pairs. For $\mathcal O_1-\mathcal O_2$: with $\hat u=u/|u|$ and $\mathbb K_3(u)_{\lambda\mu}=(3\hat u_\lambda\hat u_\mu-\delta_{\lambda\mu})/(4\pi|u|^3)$,
\[
\begin{multlined}
[\mathbb K_3(u)_{11}]^2-2[\mathbb K_3(u)_{12}]^2+[\mathbb K_3(u)_{22}]^2=\frac{(3\hat u_1^2-1)^2-18\hat u_1^2\hat u_2^2+(3\hat u_2^2-1)^2}{16\pi^2|u|^6} \\
=\frac{9(\hat u_1^2-\hat u_2^2)^2-6(\hat u_1^2+\hat u_2^2)+2}{16\pi^2|u|^6},
\end{multlined}
\]
which is at least $(32\pi^2|u|^6)^{-1}$ where $\hat u_1^2+\hat u_2^2\le\frac14$, the bracket being then at least $2-\frac32$; this holds at all $u\in{\rm supp}\,f-{\rm supp}\,h$ for the test pairs of the statement, for instance for those carried by balls of radius $\mathsf d$ whose centres are $4\mathsf de_3$ apart. With $|u|\le10\mathsf d$ there (Lemma~\ref{lem:ward-test-pairs}(a)) and the masses $\ge\vartheta\mathsf d^3\|f\|_\infty$, $\vartheta\mathsf d^3\|h\|_\infty$ of the non-negative pair (Lemma~\ref{lem:ward-test-pairs}), the main term of the two-point function of $\mathcal O_1-\mathcal O_2$ is at least $\frac32\mathsf b^{-2}(32\pi^2)^{-1}10^{-6}\vartheta^2\|f\|_\infty\|h\|_\infty$, and its error is at most $4$ times the error of the display, one for each of the four two-point functions of the planes $1,2$; so it is positive for small $\mathsf d$, whereas a plane field that is a multiple of $\mathcal O$ would make $\mathcal O_1-\mathcal O_2$ vanish.

The combination $\mathcal O-2\mathcal O_m$ has coefficients of modulus $1$ in three planes, whence the factor $9$ for two entries. Multilinearity gives one further factor $\mathsf b$ in both the main term and the error for each entry multiplied by $\mathsf b$. (f) For free photons, expansion of
$\log\det(1-\mathsf b^{-1}M_\lambda)^{-3/2}$, where
$M_\lambda=\sum_a\lambda_aQ^{\rm T}{\rm diag}(w_a)Q$,
gives $3\mathsf b^{-n}\mathcal K_K^{(n)}$ exactly.
For the interacting weights the lattice comparison is
$\mathrm{(Q_n)}$. The kernel comparison and bound
\YMMref{(2245)--(2246)}{824}, applied to the factors of each
ring and telescoped plane by plane, give the displayed
$C'_n$ error; multilinearity gives the factors for plane
combinations and rescaled entries. (g)(i) Lemma~\ref{C-360-7} on $T_{\ell_g}$, applied to the Wick monomials of the tuple, whose eight mirror copies have as mean a mixture over histories, with jets that are nearly constant at the working level and with large fields that are rare on the small side.

A real multilinear form that is invariant under the lattice translations and bounded by $\mathsf a$ at every choice of signs of a packing of $N_c\sim\mathsf d^{-3(1-\sigma)}$ translates of the tuple in one cube of side $\sim\mathsf d^\sigma$ is bounded by $\mathsf a/N_c$; this gives the factor $\mathsf d^{3(1-\sigma)}$. Each derivative in a source is paid by Cauchy's estimate on the polydisc of the height, which costs $(\frac43\bar{\mathsf b}/\mathsf d)^{[\frac n2-1]_+}$, $\bar{\mathsf b}=L(g^{-2}+\bar B)$; the sizes at the working depth give a factor $\Lambda(\mathsf d)\le C(\log\frac1{\mathsf d})^c$; so $\mathsf e'_n=2^{n-1}(\frac{42L}M)^3(\frac43\bar{\mathsf b})^{[\frac n2-1]_+}[c^\circ_nQ_0\Lambda(\mathsf d)]^n\mathsf d^{3(1-\sigma)-[\frac n2-1]_+}$ with $c^\circ_n=n2^{3n}(1+n4^{n-1})$ and $Q_0$ a constant of the stability bounds on $T_{\ell_g}$.

Every periodic Gibbs state is a state at $g$. (ii) For real weights, the mean of the product of the eight mirror copies of the Wick polynomials $\Xi_B$, centred at the cumulants of $T_{\ell'}$, is exactly the sum over the partitions of the copies all of whose blocks link different copies of the products of the cumulants of the blocks (an identity of formal power series). Hypotheses $\mathrm{(Q_m)}$, $2\le m\le8n$, apply to the real plaquette weights without an upper bound on their mutual distances. Together with the following Gaussian ring bound, they show that a linking cumulant is at most $\mathsf e_m+\mathsf v_m$, where $\mathsf v_m=\frac32(m-1)!(81(1+3C_{\rm lat}))^m(\frac{4\mathsf d}{\ell'})^6$ bounds the linking lattice rings, two copies in different cells being $\ge\frac14\ell'$ apart; with ${\rm Bl}_N(k)=\sum_\pi\prod_{I\in\pi}k_{\#I}$ over the partitions of $N$ points into blocks of at least two, Lemma~\ref{C-360-7} gives the second display with $\eta^\sharp=2^ne^{G_g\mathsf l^3}\max_{m\le n}{\rm Bl}_{8m}(\mathsf e+\mathsf v)^{1/8}$ and $\mathsf e^\sharp_n=\mathsf e_n+\eta^\sharp(1+\eta^\sharp)^{n-1}$.

\emph{Step 4: Theorem~\ref{C-360-2}(h).} (i) On the least torus a pair of narrow bumps, of width $\rho$ at the distance $\mathsf d$, is treated at the level $k_1=\max\{K-s(\mathsf d),2\}$ of the distance $\mathsf d$, the working level of Lemma~\ref{C-360-24} at the depth $s(\mathsf d)$: there the jet of each bump at the background is a constant plus an explicit quadratic form of the curvature, up to errors that the bounds of Theorem~\YMMref{D}{8} make summable over the levels. The induction over the levels with the two sources runs upward from the level $k_1$ to the last level $K$, at which the block field is trivial and the quadratic form vanishes, and carries this to the Gaussian value, the explicit part $\frac32\mathsf b^{-2}\sum f_1f_2(C^{\rm free}_K)^2$, with the error $\epsilon^{\rm nar}_2(g,\mathsf d)(\rho/\mathsf d)^6$ in natural sizes, $\epsilon^{\rm nar}_2\le C(g)(\log\frac1{\mathsf d})^{19\mathfrak p_0}\mathsf d^{1/2}$.

The remainder of the two-point function beyond the Gaussian value is bounded by placing a dummy bump at the distance $\rho$ from the first bump, applying the bound on the remainder of three sources of (iii) at the lengths $(\rho,\rho,\mathsf d)$, and setting the parameter of the dummy bump to $0$ in the mixture over the histories; it is $\le C\mathsf b^2\mathsf d^{3/2}(\rho/\mathsf d)^6$; the chessboard bound, read at two lengths (the fold at two lengths), carries it to every state with a further error $\le C(g)(\rho/\mathsf d)^6\mathsf d^{3/2}$; the comparison of $\sum_{\lambda(p)=\lambda,\lambda(q)=\mu}\varphi_1\varphi_2(C^{\rm free}_K)^2$ with $\mathcal K_{\lambda\mu}$ costs $C[\frac{a_K}{\vartheta\rho}+(\frac{\mathsf d}{\ell_g})^3](\rho/\mathsf d)^6$.

The change of units $x\mapsto L^zx$ multiplies the label, the side and the distances by $L^z$, and for $z$ large it carries a first-passage datum at $g\le\gamma$ to one at a coupling below $\gamma_\sigma$; this gives the narrow form at every $0<g\le\gamma$. (ii) For bumps of radius $\rho$ the bounds of (i) are $A\rho^{6}\prod\|\phi_a\|_\infty$, uniform in $j$; a bilinear form with such bounds at all small radii has a Borel kernel with $|k|\le(\frac3\pi)^2A$; the subsequence is taken at all pairs through partitions of unity; every periodic Gibbs state at $\mathfrak D(g)$ is invariant under the lattice translations, and the tori too, so the kernels depend on $x-y$; $\sum_\lambda\mathcal O_\lambda=\mathcal O$ gives $\sum\tau_{\lambda\mu}=\mathsf s^T_2$. (iii) With three sources at two scales, an induction over the levels shows: let $k_1=\max\{K-s(\mathsf d),2\}$ be the level of the distance $\mathsf d$, $\mathsf b_{k_1}=A_1(\log x_{k_1})^{\mathfrak p_0}$ its threshold and $\delta_{k_1}=\delta(x_{k_1})=A_1x_{k_1}^{-1/2}(\log x_{k_1})^{\mathfrak p_0}\le A_1Y_{s(\mathsf d)}^{\mathfrak p_0}(\mathsf d/\mathsf b)^{1/2}$ (equation~\YMMref{(634)}{339}), and $Z=Z_0(1+\mathsf b_{k_1})^2$ with a number $Z_0$ formed with the constants among the construction results, so that $Z^m\le C\,Y_{s(\mathsf d)}^{2m\mathfrak p_0}$; if at the level $k_1$ each jet of the sources is a constant plus an explicit quadratic form of the curvature, up to $\epsilon_\circ\mathsf v_\beta$ with weights $\mathsf v_\beta$ whose products over partitions do not exceed $\mathsf v_\beta$, then at every later level, up to the last level $K$, it is the Gaussian value, up to $\mathsf a_mZ^m\varpi\mathsf v_\beta$ with $(\mathsf a_m)=(1,14,285)$ and $\varpi\le\epsilon_\circ+C\delta_{k_1}$, provided $(K-k_1)\varpi\le1$, which holds for $\mathsf d$ below a threshold depending on $g$.

As $Y_{s(\mathsf d)}\le2\log\frac1{\mathsf d}+C(g)$, the factors $Z^m$ and $\mathsf b_{k_1}^2\delta_{k_1}$ give the powers of $\log\frac1{\mathsf d}$ below. The explicit Gaussian part is the ring $3\mathsf b^{-3}\mathcal K^{(3)}_K$, with the error $C(g)(\log\frac1{\mathsf d})^{39\mathfrak p_0+2}\mathsf d^{1/2}$ times the natural size $\rho^9r^{-3}\mathsf d^{-6}\mathsf b^{-3}\prod\|f_a\|_\infty$; the remainder $S^T_3-\mathsf J_3$ is at most $C\mathsf b^3\mathsf d^{3/2}$ times the natural size, with no power of $r$; the passage over the least torus costs $(\mathsf d/\ell_g)^3$, and the comparison of the lattice and continuum rings at two scales is explicit. As in (i), the fold at two lengths carries the three-point law at narrow bumps from the least torus to every state, with the further error $C^{\rm tr3}_3(g)\,\mathsf d^{3/2}\rho^9r^{-3}\mathsf d^{-6}\prod_a\|f_a\|_\infty$, $C^{\rm tr3}_3(g)$ a number formed with $g$ and the constants among the construction results; $E_3$ contains this term, the labels at $g$ lying in a compact interval and $\mathsf d^{3/2}\le\mathsf d^{1/2}$, and the law of (h)(iii) along the states $s_j$ rests on it. (iv) follows from (ii), (iii) and the free identity
\[
\bigl|\tfrac{2\pi}3\mathsf br^3\cdot3\mathsf b^{-3}k_3\bigl(c+\tfrac r2n,c-\tfrac r2n,z\bigr)-\mathsf t^\circ(n,c-z)\bigr|\le2.66\,(r/\mathsf d)^2\mathsf b^{-2}\mathsf d^{-6},
\]
with $|t_\lambda(c,z)-\mathsf t^\circ(e_\lambda,c-z)|\le6\epsilon_\sigma(g,|c-z|)\mathsf b^{-2}|c-z|^{-6}$ almost everywhere by (ii) and the identity for $\mathsf t^\circ$ of the setting; the null set lies in $(c,z)$ alone because $\mathsf s^T_3$ is continuous. (v) The limit in $K$ exists at disjoint supports by (d): its clause on one plane field among action densities gives the limit of $\langle\mathcal O_\lambda(\phi_1);\mathcal O(\phi_2)\rangle$ at the fixed label, and $\langle\mathcal O(\phi_1);\mathcal O(\phi_2)\rangle$ converges by the convergence of $S^T_2$ at a fixed label. For the identification, the narrow bumps are the $f^c_\pm$ of the statement, the radial tents of radius $\rho=\frac r8$, of mass $1$ and of height $3/(\pi\rho^3)=1536/(\pi r^3)$ about $c\pm\frac r2e_\lambda$; the threshold on $K$ is $128a_K\le r$, that is $16a_K\le\rho$, and $s(\rho)\le K$, which gives $s(\frac r2)\le K$, the depth function being non-increasing.

The number $r_*$ is taken so that for $0<r\le r_*$ the depth $s=s(\frac r2)$ obeys $s\ge3$ and $L^2C_W(s)Y_s^{3\mathfrak p_0}\le L^{s-3}$, the condition under which the family near the pair is formed, and $\frac r2\le\min\{r_0(\mathsf b),r_1(\mathsf b)\}$, the radii up to which the induction of (iii) and the frame of the pair are formed, which read $g$ and $\mathsf b$; all three hold for small $r$, since $s(\frac r2)\to\infty$ as $r\to0$ and the second holds at all large $s$; and $r_*\le\frac16\mathsf d$. Write $S^T_3=\mathsf J_3+(S^T_3-\mathsf J_3)$, $\mathsf J_3$ the explicit Gaussian part, cut $\phi_2$ into pieces carried by cubes of side $\frac14r$, of the natural size $\mathsf b^{-2}\mathsf d^{-6}r^3\|\phi_2\|_\infty$, and compare, at each $c\in{\rm supp}\,\phi_1$ and for each piece, the pair of narrow bumps at $c\pm\frac r2e_\lambda$, weighted by $\frac{2\pi}3\mathsf br^3$, with the plane field $\mathcal O_\lambda-\frac13\mathcal O$ at $c$.

For the remainder: for almost every block field, at every level $n$ with $a_{K-n}\ge r^{1/2}$, the difference jet at the background, the jet of the weighted pair minus that of the plane field, is a constant up to $C[r^{\beta_H/72}+(a_K/r)^{1/18}]\varepsilon_n^2a_{K-n}^{-4}$, $\varepsilon_n$ the radius of regularity of the level $n$ and $\beta_H$ the H\"older exponent; this comes from the H\"older regularity, with the exponent $\beta_H$, of the backgrounds at the scale of the pair, carried over the levels by the induction of (iii). Integrating over the mixture of the histories on the least torus, whose moments of the scale of regularity are bounded uniformly in $K$, gives for the remainder the difference estimate $C(g,\mathsf b)\mathsf d^{-5/2}[r^{\beta_H/72}+(a_K/r)^{1/18}]$ times the natural size.

The explicit Gaussian parts are compared on the lattice of $T_{\ell_g}$, with $C(g,\mathsf b,\vartheta_0)\mathsf d^{-9/2}$ and the same bracket. In the ring of the left member the near leg, between the two tents, whose plaquettes are between $\frac34r$ and $\frac54r$ apart, is $a_K^3\mathbb K_3(re_\lambda)$ up to the cost $(a_K/r)^2+(r/\ell_g)^3$, by the comparison \YMMref{(2245)}{824} of the lattice kernel with $\mathbb K_3$ in the setting. The two far legs $\mathsf a$, $\mathsf a'$, from the two tents to the support of $\phi_2$, enter through $\mathsf a\mathsf a'=\frac12(\mathsf a^2+\mathsf a'^2)-\frac12(\mathsf a-\mathsf a')^2$: the term $(\mathsf a-\mathsf a')^2$ costs $(r/\mathsf d)^2$ by the bound on the differences of $C^{\rm free}_K$ along one plane of the setting, its two plaquettes being parallel, equally oriented and at most $\frac54r\le\frac14(\mathsf d-\frac58r)$ apart; and in each square the far leg reads one tent alone.

There the lattice sums over the tents are Riemann sums, at the cost $a_K/\rho=8a_K/r$; the entries of $\mathbb K_3$ are harmonic off $0$, so its mean over the two radial tents, over the one through the sum over its plaquettes and over the other through the integral over $c$, is its value at $re_\lambda$; and the integral over $c$ moves the argument of $\phi_1$ by at most $r$, at the cost $r\,{\rm Lip}\,\phi_1\le r\|\phi_1\|_\infty/(\vartheta_0\mathsf d)$. Finally $\frac{2\pi}3r^3\cdot3\mathbb K_3(re_\lambda)=\frac32(E_{\lambda\lambda}-\frac13\mathbf 1)$, so the weighted pair becomes the plane field $\mathcal O_\lambda-\frac13\mathcal O$ at $c$ in the explicit part of the two-point function of the right member, whose legs are lattice kernels as well.

Each of these costs is a number formed with $g,\mathsf b,\vartheta_0,\mathsf d$ times the bracket, as $a_K/r\le\frac1{128}$ and $r\le1$. Summing over the at most $9^3(\mathsf d/r)^3$ pieces and integrating over $c$, against the mass $\le\frac{4\pi}3\mathsf d^3\|\phi_1\|_\infty$, gives the display of (h)(v), uniformly in $K$. Letting $K\to\infty$ at a fixed $r$, which is allowed since $128a_K\le r$ and $s(\rho)\le K$ from some $K$ on, and then $r\to0$ gives the identification at the fixed $\mathsf d$.

\emph{Step 5: Theorem~\ref{C-360-3}(a), (b).} By Lemma~\ref{C-360-6}, $-\beta\sum_ph_{\omega,\chi}=\mathsf b\,{\rm Pf}(k^h)$, with $\|{\rm Pf}(k^h)\|\le156|\omega|\,{\rm Lip}(\chi)R_1(\chi)\le468\ell'|\omega|\,{\rm Lip}_\infty(\chi)$, and (i) is Lemma~\ref{C-360-6} at $r=n+1$, with $C^\flat_n(\rho)=468\ell'k^\flat_n(\rho)$, $k^\flat_n(\rho)=\max\{k'(n+1,\frac\rho2),\text{the crude bound below }K'_0\}$. Both clauses follow the one rule of the statement for a rough entry $F$: its norm $\|F\|$ times a number, and, free of the side and of the state, times the counts $N_F$ and $N_i$ of the cubes of the mesh that meet its carrier and the supports. For (ii) we put the rough entry alone in its cell of a board and the Lipschitz fields in other cells, in any mutual position.

For groups $G_j$ of entries, $R_G=\partial_G[e^{\sum_{e\in G}z_eX_e}/\nu(e^{\sum_{e\in G}z_eX_e})]_{z=0}$ has $\nu(\prod_jR_{G_j})=\sum_\pi\prod_{I\in\pi}\kappa(I)$ over the partitions all of whose blocks meet two groups; and, for integers $n\ge0$, $p\ge1$ and numbers $s_1,\dots,s_n>0$, $\rho_*>0$, $Q\ge1$, $\alpha_1,\dots,\alpha_p\ge0$: if $A_1,\dots,A_p,B$ are holomorphic on the domain $\mathsf D=\{z\in\mathbb C^n:\sum_r|z_r|s_r<\rho_*\}$, with $B(0)=1$, $|A_j(z)|\le\alpha_jQ^{\#z}$ and $|B(z)|\le Q^{\#z}$ on $\mathsf D$, $\#z$ being the number of the non-zero coordinates of $z$, then $A_1\cdots A_pB^{-p}$ is holomorphic near $0$ and $|\partial_{[n]}(A_1\cdots A_pB^{-p})(0)|\le\binom{n+p-1}nn!\prod_j\alpha_j\prod_r(2eQs_r/\rho_*)$.

The chessboard bound of Lemma~\ref{C-360-7} and the bounds of Theorems~\YMMref{AJ}{673} and~\YMMref{AD}{607} on tori for the scalar fields, through their cumulant polynomial, with $e^{nc_V\ell'^3}$, conclude. On the larger sides the chessboard seminorm of a centred rough entry is at most $e^{G_g\mathsf l^3}k_{N_0}(\mathsf d)\|F\|$, $k_N(\mathsf d)=[\sum_\pi\prod_{I\in\pi}k'(\#I,\mathsf d)]^{1/N}$ over the partitions of $N$ points into blocks of at least two, by its $N_0$ mirror copies on the reference side; a partition of unity by tents of the mesh $h(n,\rho)$ cuts the shell and the fields, which gives the counts $N_\chi,N_i$; the domination of the periodic Gibbs states (Lemma~\ref{C-360-7}) gives the states, and the change of units $x\mapsto L^zx$ gives every label and side.

The whole flux insertion is the rough entry $\mathsf b\,{\rm Pf}(k^h)+\mathcal O(v^{\mathfrak B})$, of the norm $13\mathsf b|k^h|_1+|v^{\mathfrak B}|_\infty\le(158\mathsf b+3a)|\omega|\,{\rm Lip}(\chi)R'(\chi)$, $R'(\chi)=\max\{R_1(\chi),R'_1(\chi)\}$ (Lemma~\ref{C-360-6}), and the same rule bounds it. For $\Theta_\beta[\eta]$ we use $\Theta_\beta[\eta]=\mathsf b\,{\rm Pf}(k^\eta)$; for $\mathfrak K^{\rm v}[\beta\eta]$, Lemma~\ref{C-360-8}(a) gives $\mathfrak K^{\rm v}[\beta\eta]=N_\beta[\eta]-\Theta_\beta[\eta]+\mathsf b\,\mathcal O({\rm div}^s\eta)-\frac32a\,\mathcal O(e^\eta)$, and the null term drops beside fields that do not depend on the bonds near ${\rm supp}\,\eta$. (b) The spin densities written in petals give $T_{lr}(\varphi)=\mathsf b({\rm Pf}(k^{lr})+\mathcal O(v^{lr}))$; for two entries the two-source mixed derivatives of Lemma~\ref{C-360-24}(a), (e) give the bound first at $0<g\le\gamma$ with $g^{-2}\ge x_{\rm J}$, the setting of that lemma, with the factor $(2\ell/\mathsf d'')^6k_2(g,s(\mathsf d''))g^4$.

Every $0<g\le\gamma$ is then reached by the change of units $x\mapsto L^zx$, which carries $\mu_{K,\ell,\beta}$ to $\mu_{K-z,L^z\ell,\beta}$ and multiplies the label by $L^z$: take $z$ with $L^z\gamma^{-2}\ge2\max\{\gamma^{-2},x_{\rm J}\}$. The labels of $\mathfrak D(g)$ lie in the compact interval $[\gamma^{-2},X_g+\bar B]$, whose image under the multiplication by $L^z$ lies in $[2\max\{\gamma^{-2},x_{\rm J}\},\infty)$; every $\mathsf b'$ there lies inside $I(g')$ at $g'=(\mathsf b'/2)^{-1/2}\le\gamma$, where $g'^{-2}=\frac12\mathsf b'\ge x_{\rm J}$, so finitely many intervals $I(g'_i)$ cover the image, and at the cutoffs $K>z$ the bound at the couplings $g'_i$ applies; the cutoffs $K\le z$ are bounded crudely, by the suprema of the entries.

The limits are formed as in Step~1. For the symmetry in $(l,r)$ of the limits: for Lipschitz $\varphi$, by Lemma~\ref{C-360-8}(e), $(T_{lr}-T_{rl})(\varphi)=\beta\sum_b\sum_{p\ni b}(\varphi(x(p))-\varphi(x_b))a^{-1}[(x_b-x(p))_l\mathfrak d_{b,p,r}-(x_b-x(p))_r\mathfrak d_{b,p,l}]$, which is $\mathsf b$ times a rough entry of norm $\le15a\,{\rm Lip}\,\varphi$; so, by the two-entry bounds just proved, at every first-passage datum and every $0<g\le\gamma$, both $\langle(T_{lr}-T_{rl})(\varphi);\mathcal O(f)\rangle$ and $\langle(T_{lr}-T_{rl})(\varphi);T_{l'r'}(\varphi')\rangle$ are $O(a)$ at separated supports. The limits are continuous in the sup norms of the entries, by the same bounds, and the Lipschitz functions are dense in the continuous ones in the sup norm; so every limit point is symmetric in $(l,r)$.

\emph{Step 6: Theorem~\ref{C-360-3}(c), (d).} (c) By Lemma~\ref{C-360-8}(a) the null term drops, and $a\,\mathcal O(e^\eta)$ contributes $O(a)$ by Theorem~\ref{C-360-2}(a). By Lemma~\ref{C-360-8}(b), $\mathfrak K^{\rm v}[\beta\eta]=-\mathsf b\sum_\rho P_\rho(\nabla^+_\rho\eta^\rho)+\beta\mathfrak K^\sharp[\eta]$, $P_\rho(\mathsf k)=a^{-1}\sum_p\mathsf k(x(p))\Psi_\rho(p)$: a bilinear half and a cubic half. In the bilinear half an exact identity in $U$, direction by direction and with one integration by parts, is used. For a real weight $W$ of the plaquettes and a direction $\rho$, let $(\mu,\nu,\rho)$ be the cyclic permutation of $(1,2,3)$ that ends with $\rho$; for a bare site $y$ let $p(y;\alpha,\alpha')$ be the plaquette with the corners $y$, $y+ae_\alpha$, $y+ae_\alpha+ae_{\alpha'}$, $y+ae_{\alpha'}$, and $w_0(y)=W(p(y;\mu,\nu))$, the weight of the plaquette normal to $\rho$ at $y$; let $g_\mu(y)=G_\rho(\langle y,y+ae_\mu\rangle)$ and $g_\nu(y)=G_\rho(\langle y,y+ae_\nu\rangle)$, the transported differences of the field strengths of the setting, and $E_\mu(y)=\sum_{p\ni b}\mathbf w(W_{b,p})$, the sum over the four plaquettes that hold the bond $b=\langle y,y+ae_\mu\rangle$, the lattice equation of motion at $b$, and $E_\nu(y)$ in the same way. Then
\[
-4\sum_pW(p)\Psi_\rho(p)=\sum_yw_0(y)\bigl[\langle g_\mu,E_\mu\rangle+\langle g_\nu,E_\nu\rangle\bigr]+{\rm Bil}+{\rm Cub}+{\rm Q4},
\]
where ${\rm Bil}$ comes from the differences of the weight $W$, ${\rm Cub}$ from the defect of the linear Bianchi identity and from commutators (factors ${\rm Ad}(W_{b,p})-\mathbf 1$ of plaquette holonomies), and ${\rm Q4}$ from $a_p^2+a_{p_+}^2$. For each plaquette $p=p(y;\mu,\alpha')$ through the bond $b=\langle y,y+ae_\mu\rangle$, the one-bond term $\beta\langle\mathbf w(W_{b,p}),E_\mu(y)\rangle-3(1-a_p)$ has vanishing cumulants with every entry that does not depend on $U(b)$ (an integration by parts in $U(b)$, as in Lemma~\ref{C-360-8}(a)), and $g_\mu(y)$ is the sum of the $\mathbf w(W_{b,p})$ over the two plaquettes through $b$ in the directions $\pm\rho$; so $\beta$ times the first sum is, up to a number, $-3a\,\mathcal O(v)$ inside cumulants whose other entries avoid the bonds concerned, with the weight $v=w_0(y)+w_0(y+ae_\rho)$ on the plaquettes $p(y;\rho,\mu)$ and $p(y;\rho,\nu)$ and $v=0$ on the plaquettes normal to $\rho$; $|v|\le2\sup|W|$ gives the factor $\frac32$ below; $\beta\,{\rm Bil}$ is $a\mathsf b\,{\rm Lip}_\infty\mathsf k$ times a bounded petal form; ${\rm Cub}$ and ${\rm Q4}$ are block forms all of whose monomials have an order $\ge3$, whose sizes fall like $L^{-2j}$ with the level $j$ and give $L^{-(K-s-1)/4}$ at the working depth $s$ (Lemma~\ref{C-360-24}(c), (d)).

So, for $r\ge1$ spectators $F_q$ of the sizes $\omega_q$ (block forms with weights, or $\mathcal O(v_q)$) whose carriers are at Euclidean distances $\ge\mathsf d$ from one another and from $\{\mathsf k\ne0\}$, at all first-passage data with $36a\le\mathsf d$, $K\ge s(\mathsf d_1'')$ and $K\ge s(\mathsf d_1''')$, $\mathsf d_1=\frac34\mathsf d$,
\[
|\mathsf b\langle P_\rho(\mathsf k\circ x);F_1;\dots;F_r\rangle|\le\Bigl\{a\bigl[\tfrac32\sup|\mathsf k|+190\,\mathsf b\,{\rm Lip}_\infty\mathsf k\bigr]e_2+50\,\mathsf b\,\Gamma L^{-(K-s-1)/4}e_3\sup|\mathsf k|\Bigr\}\prod_q\omega_q,
\]
where $e_2=e_2(r+1,\mathsf d_1)$ and $e_3=e_3(r+1,\mathsf d_1)$, the cumulant having $r+1$ entries; for a number $r'\ge2$ of entries and a distance $\mathsf d_1$, $e_2(r',\mathsf d_1)=(3\ell'/\mathsf d_1'')^{3r'}k^{\rm loc}_{r'}(g,s(\mathsf d_1'');\ell')g^{2r'}$, $\mathsf d_1''=\min\{\frac45\mathsf d_1,L^{-\bar s(r')}\}$, with the depth $\bar s(r')$ of the proof of Lemma~\ref{C-360-6}, the least integer $s\ge3$ with the condition $2\mathsf T_{r'}(3)\theta(L^{s-2}g^{-2})\le1$ of Lemma~\ref{C-360-24}(b), read with the series $\mathsf T_{r'}$ of the mixed derivatives of the order $r'$ and the size $\theta$ stated there, and the number $k^{\rm loc}_{r'}$ of Lemma~\ref{C-360-24}(b); $e_3(r',\mathsf d_1)=(3\ell'/\mathsf d_1''')^{3r'}\tilde k_{r'}(g,s(\mathsf d_1''');\ell')g^{2r'}$ is the like number at the second depth $\mathsf d_1'''=\min\{\frac34\mathsf d_1,L^{-\tilde s(r')}\}$, $\tilde s(r')$ the least integer $s\ge3$ with the condition of Lemma~\ref{C-360-24}(c) at $p=r'$, and $\tilde k_{r'}$ the number of the bound of Lemma~\ref{C-360-24}(c); $\Gamma=\max\{\Gamma^\circ_{\rm Cub},\Gamma^\circ_{\rm Q4}\}$, the larger of the comparison weights $\Gamma^\circ_c$ of Lemma~\ref{C-360-24}(d) of the two block forms ${\rm Cub}$ and ${\rm Q4}$; and $s=s(\mathsf d_1''')$ at $r'=r+1$.

In (c) this is used with $r=n$ spectators $\mathcal O(f_i)$ and $\frac12\mathsf d$ in the place of $\mathsf d$, so that the numbers are formed at $(n+1,\frac38\mathsf d)$. These numbers depend on $g,\ell',r,\mathsf d$ and on no $K$ or $\beta$. The cubic half is Lemma~\ref{C-360-10}, the spread being $\ge\mathsf d$; the norm $\|\cdot\|_\sharp$ and the rate $c<\frac1{14}$ come from it alone. (d) For the cubic half in sup norms we couple on one probability space the chain of the first-passage datum and the chain along $x_j=L^{K-j}\mathsf b$ of Lemma~\ref{C-360-9}, which share only the block map, the Haar measure and the bare density, and use the law of total covariance given the block field of the level $t=K-\lceil\phi K\rceil$:
\[
\langle X;F\rangle=\int(r_{XF}-r_Xr_F)\,d\mathsf m_t+\int(r_X-\mathsf a)(r_F-f_0)\,d\mathsf m_t+(\text{large fields}),
\]
$r_X,r_F,r_{XF}$ the densities at the never-large history. The first term is a mixed derivative of two separated sources, $\le C_2e^{3-L^u\mathsf d'/M}y_Xy_F$, $\mathsf d'=\frac12\mathsf d$, $y_X=4\ell'^3L^{4u}\cdot608\,\mathsf bL^K\sup|\eta|$, by Lemma~\ref{C-360-24}(e) at the level $t$, whose conditions $L^u\mathsf d'\ge3M$ and $\mathsf h_2(u)\le1$ hold for large $K$; the second is the oscillation of Lemma~\ref{C-360-9}(b) times $3C_1y_F$, $y_F=4\ell'^3L^{4u}\sup|w|$, and $4\phi<c'_0$; the rest is at most $\mathsf V_{\ell'}(\mathsf b)L^{-11K}$. In words (Lemma~\ref{C-360-17}): a spectator acts on the cubic insertion only through a bounded function of the block field, and the cubic insertion is uncorrelated with every function of the block field up to its oscillation, at the price $L^{4\phi K}$, which a small $\phi$ pays.

For $n$ spectators the mixed derivatives of every order of the never-large term are bounded by Cauchy's estimates, locality and a count, $e^{-L^u\mathsf d/2M}$ paying every power of $L^K$. In the bilinear half in sup norms, $\mathsf b\,a^{-1}\sum_pw(p)\Psi_{\lambda(p)}(p)$, summed over the plaquettes of one normal, is a dominated local source of the weight $76\mathsf b\sup|w|$ (Lemma~\ref{C-360-8}(b)), and a source cut at a level, with the H\"older bound~\textup{\ref{I-56-65}}, gives $L^{-K/200}$. For two cubic insertions the same argument is used twice, and the factor $L^{4u}$ does not enter; two orbital insertions and two stress insertions follow with the bilinear half and the null identity.

The numbers of Lemma~\ref{C-360-9}(b) are uniform on $[\gamma^{-2},X_g+\bar B]$, which holds every label of $\mathfrak D(g)$. The several-source statements used above hold clause by clause in their own ranges: Lemma~\ref{C-360-24}(a), (e), which the clauses on the cubic insertion alone read, at $0<g\le\gamma$ with $g^{-2}\ge x_{\rm J}$; Lemma~\ref{C-360-24}(b), (c), which the bilinear half reads, at $0<g\le\gamma^{\rm ref}_W$, and so does Theorem~\ref{C-360-2}(a), which bounds the term $a\,\mathcal O(e^\eta)$. So at every first-passage datum at $g$ on every side of $\mathfrak L(g)$, with numbers that depend on $g$ and on no label, the clauses on the cubic insertion alone hold for $0<g\le\gamma$ with $g^{-2}\ge x_{\rm J}$, and the other clauses for $0<g\le\gamma^{\rm ref}_W$, where $g^{-2}\ge x_{\rm J}$ holds of itself; and $x\mapsto L^zx$ gives every fixed label and side.

\emph{Step 7: Theorem~\ref{C-360-3}(e)--(g).} In the mean of the mirror copies $\mathrm{St}_K[\eta]$ may be replaced by $-\mathrm{Rm}_K[\eta]$ (Lemma~\ref{C-360-8}(a) and the proof of Lemma~\ref{C-360-11}). A partition of unity cuts $\eta,\eta',H'$ into at most $C(R/\mathsf d)^3$ pieces of sup diameter $\le\frac18\mathsf d$, at the cost $C\mathsf d^{-2}$ in the $C^2$ norms. For two sets of sup diameter $\le\frac18\delta_0$ that are $\delta_0$ apart along one direction there is an offset in $h\mathbb Z^3$, $h\le\frac1{64}\delta_0$, of a board of side $\frac12\ell_g$ for which they lie in different cells at sup distance $\ge\frac1{16}\delta_0$ from the walls; Lemmas~\ref{C-360-7} and~\ref{C-360-11} give $|\nu(A^\nu[F_1]A^\nu[F_2])|\le2e^{2G_g\mathsf l^3}\|F_1-c_1\|_{\rm ch}\|F_2-c_2\|_{\rm ch}$.

Two stress insertions give $L^{-K/28}$ each, hence $L^{-K/14}$; for two conserved tensors $L^{-K/36}$ each gives $L^{-K/18}\le L^{-K/20}$; against a plain tensor entry, which costs $L^{4u}=L^{K/100}$, the factor $L^{-K/36}$ leaves $L^{-K/60}$; and $\mathsf b\,\mathcal T_K(H^{\rm cons}[\eta])=\mathrm{St}_K[\eta]-\Delta_K[\eta]$. For the mean of the orbital insertion, cut into the pieces of the partition of unity, $\nu(\mathfrak K^{\rm v}[\beta\eta])=\nu(F_\eta)$ with $F_\eta=\mathfrak K^{\rm v}[\beta\eta]-\frac14\mathfrak N[\mathsf k_\eta]$, $\mathsf k^\rho_\eta=\nabla^+_\rho\eta^\rho$, since $\nu(\mathfrak N[\mathsf k_\eta])=0$ by Lemma~\ref{C-360-8}(a); Lemma~\ref{C-360-11} at $m=1$, $\phi=\frac1{400}$ gives $\|F_\eta\|_{\rm ch}\le CK^pL^{-K/28}\|\eta\|_{C^2}$, $F_\eta$ having the mean $0$ on $T_{\ell'}$, and Lemma~\ref{C-360-7} at $q=1$, $c_1=0$ gives $|\nu(\mathfrak K^{\rm v}[\beta\eta])|\le e^{G_g\mathsf l^3}CK^pL^{-K/28}\|\eta\|_{C^2}$. (f) $n$ sets of sup diameter $\le\frac{\mathsf l}{4n}$ lie at sup distance $>\frac{\mathsf l}{8n}$ from all the walls of some board with offset on a grid of mesh $<\frac{\mathsf l}{2n}$ (a pigeonhole); several slots may share a cell (Lemma~\ref{C-360-11}); Lemma~\ref{C-360-7} holds with any number of cells; recentring from the torus mean to $\nu$ costs $2^n$ terms; and $4\phi(n-1)<\min\{c'_0,\frac1{28}-3\phi\}$ pays the factors $L^{4u}$.

At $n_1=1$, $c=\frac1{56}$, a fixed $\mathsf d>0$ satisfies $\mathsf dL^u\ge K^2$ for $K\ge K_0(\mathsf d)$. (g) $\sigma(H^{\rm cons}[\eta],H_2,\dots,H_n)$ is the limit of the left member of (f) at $n_1=1$, which tends to $0$ at every fixed $\mathsf d>0$, and disjoint compact supports have a positive sup distance. By reflection positivity of $\sigma$ at the mirror $\theta$, $\sigma(H^{\rm cons}[\theta\eta],H^{\rm cons}[\eta])$ is the square of the norm of the vector of $\mathcal T(H^{\rm cons}[\eta])$ in the reconstruction, which therefore vanishes. The equivalence along the sequence is Lemma~\ref{C-360-8}(a) with the bound on $\Delta_K$ in (e), the labels lying in $[\gamma^{-2},X_g+\bar B]$.

For $0<g\le\gamma^{\rm sd}_W$ the cited bound on centred products of tensor entries and a diagonal argument give the convergent subsequences; the covariance and the reflection positivity pass to the limits from the periodic Gibbs states, and $\mathcal T_K(f\mathbf 1)=\mathcal O(f)$ gives the trace part.

\emph{Step 8: Theorem~\ref{C-360-4}(a), (b).} (a) We write $\mathfrak U_{\omega,\chi}=(\mathfrak U_{\omega,\chi}-\mathfrak U^\sharp[\chi\,\omega y])+\mathfrak U^\sharp[\chi\,\omega y]$, where $\chi\,\omega y$ is a field with $\mathsf H(\chi\,\omega y)\le(2+3\ell)|\omega|\|\chi\|_{C^{1,1}}$; since $\chi$ covers the $f_i$, $f_i\chi=f_i$, and the cubic member of the bulk term is $\mathfrak U^\sharp[\chi\,\omega y](f)$. The first part has $2(n+1)$ members, a spin member and a bilinear orbital member for each $\varphi\in\{\chi,f_1,\dots,f_n\}$. The spin members are $-\beta\langle\mathfrak S_\omega[\chi];\mathcal O(f_\cdot)\rangle$, with $\beta\,\mathfrak S_\omega[\chi]=\mathsf b\,a^{-1}\mathfrak S_\omega[\chi]$, and $\langle a^{-1}\mathfrak S_\omega[f_i];(\mathcal O(f_l))_{l\ne i}\rangle$. By Lemma~\ref{C-360-8}(e), $\sum_{p\ni b}a^{-1}(x_b-x(p))_\lambda\mathfrak d_{b,p,\rho}$ is symmetric in $(\lambda,\rho)$ and $\omega$ is antisymmetric, so that
\[
a^{-1}\mathfrak S_\omega[\varphi]=\sum_b\sum_{p\ni b}\sum_\rho k^\varphi_{b,p,\rho}\,\mathfrak d_{b,p,\rho},\qquad k^\varphi_{b,p,\rho}=a^{-1}\bigl(\varphi(x(p))-\varphi(x_b)\bigr)\sum_\lambda\omega_{\rho\lambda}\,a^{-1}(x_b-x(p))_\lambda :
\]
an entry of the order $2$ with real weights $|k^\varphi|\le\frac14|\omega|{\rm Lip}\,\varphi$, which vanish unless $\varphi(x(p))\ne\varphi(x_b)$; $\beta\,\mathfrak S_\omega[\chi]$ is $\mathsf b$ times such an entry. By Lemma~\ref{C-360-8}(b), and since $\nabla^+_\rho(\varphi\,(\omega y)_\rho)=(\omega y)_\rho\nabla^+_\rho\varphi$, $(\omega y)_\rho$ being free of $y_\rho$, the bilinear members are $\langle\mathfrak J^{\rm bl}[\mathsf k^\chi];\mathcal O(f_\cdot)\rangle$ and $-\mathsf b^{-1}\langle\mathfrak J^{\rm bl}[\mathsf k^{f_i}];(\mathcal O(f_l))_{l\ne i}\rangle$, where $\mathfrak J^{\rm bl}[\mathsf k]=\beta\sum_{p,\rho}\mathsf k^\rho(x(p))\Psi_\rho(p)$ and $\mathsf k^{\varphi,\rho}(y)=(\omega y)_\rho\nabla^+_\rho\varphi(y)$; these vector fields are Lipschitz, with norm $\le5|\omega|(1+\ell)\|\varphi\|_{C^{1,1}}$, which is where the norms $\|f_i\|_{C^{1,1}}$ enter, and supported in the $a$-neighbourhood of ${\rm supp}\,\varphi$.

By Lemma~\ref{C-360-8}(c), $\mathfrak J^{\rm bl}[\mathsf k]=-\frac14\mathfrak N[\mathsf k]+\sum_{v=1}^5\mathfrak J_v[\mathsf k]-\frac14c_{\mathsf k}$: a null term, entries of the orders $2,2,2,3,4$ with weights of the order of the norm of $\mathsf k$ and, for the orders $3,4$, of $\beta$ times it, and a number, which drops from the cumulants. In (i) let $D\ge\mathsf d$ be the spread of the supports. Two rules on spreads are used. A family of sets has a spread at least that of each of its subfamilies of two or more sets. And if $2r'<D$ and a set $T$ lies within the sup distance $r'$ of one of the sets and of each set of an index set $Q$, then $T$ is empty when these are all the sets, and otherwise $T$ and the remaining sets have a spread $\ge D-r'$: points of $T$ and of the remaining sets that were closer would give points of all the sets that are pairwise closer than $D$.

With $\chi$ in the first entry, the carrier of the entry and the supports have a spread $\ge D$. With $f_i$ in the first entry, the carrier lies within $\frac32a$ of ${\rm supp}\,f_i$ and the cumulant is taken with the $\mathcal O(f_l)$, $l\ne i$: the spread is $\ge D-\frac32a$. The null terms meet the sources: the covering fixes $\chi$ at the barycentres and at the bond midpoints of the plaquettes of the supports, while $\mathsf k^\chi$ reads $\chi$ at the points $y+ae_\rho$ as well, so that ${\rm supp}\,\mathsf k^\chi$ may meet a support; and ${\rm supp}\,\mathsf k^{f_i}$ may meet ${\rm supp}\,f_l$. They are taken apart by the contact formula of Lemma~\ref{C-360-8}(f): by (N3) the cumulant of $\mathfrak N[\mathsf k]$ with the sources of a set $\mathsf A$ is a sum of cumulants of the entries $\mathfrak R[f_{q_1}\cdots f_{q_{j-1}}\mathsf k;f_{q_j}]$ with the remaining sources, and by (N4) each such entry is a number, the density $a_p$ and densities of the order $2$, with no factor $\beta$, carried within the sup distance $2a$ of ${\rm supp}\,\mathsf k$ and of the supports that it absorbs.

The spread of ${\rm supp}\,\mathsf k^\chi$ and the supports is $\ge D$, and that of ${\rm supp}\,\mathsf k^{f_i}$ and the ${\rm supp}\,f_l$, $l\ne i$, is $\ge D-a$; once these exceed $4a$, the terms that absorb all the sources vanish, the numbers drop from the others, and by the merging of sets of Lemma~\ref{C-360-8}(f) every remaining entry has with the remaining supports a spread $\ge D-2a$, resp.\ $\ge D-3a$. So in every term that is left the spread is $\ge\mathsf d-3a$, which is $\ge\frac12\mathsf d$ once $6a\le\mathsf d$, and Lemma~\ref{C-360-18} gives $L^{-K+(\frac12-c_*)K}\le L^{-c_*K}$ for the orders $2$ and $L^K\cdot L^{-\frac32K+(\frac12-c_*)K}=L^{-c_*K}$ for the orders $3,4$. The cubic part is Lemma~\ref{C-360-10}, whose spread condition holds because the support of the field $\chi\,\omega y$ is added to the supports, whose spread is $D$.

In (ii), $n=1$: the members with $f_1$ in the first entry are means, of the spin insertion and of the bilinear density $\mathfrak K^{\rm v}-\mathfrak K^\sharp$, and vanish by the mirror symmetries of the measure (Lemma~\ref{C-360-8}(b)); in the members with $\chi$ the weights vanish unless $\chi$ varies across the plaquette or the bond, so, $\chi$ being $1$ on the $\mathsf d$-neighbourhood of ${\rm supp}\,f_1$, their carriers are at sup distance $\ge\mathsf d-2a$ from ${\rm supp}\,f_1$, the null term meets no source, and Lemma~\ref{C-360-18} applies to the pair. The cubic member is $\mathfrak U^\sharp[\chi\,\omega y](f_1)=-\beta\langle\mathfrak K^\sharp[\chi\,\omega y];\mathcal O(f_1)\rangle$, the contact member $a^{-1}\langle\mathfrak K^\sharp[f_1\chi\,\omega y]\rangle$ being a mean of $\mathfrak K^\sharp$, which vanishes (Lemma~\ref{C-360-8}(b)).

Here ${\rm supp}\,f_1$ lies in the set where the field does not vanish, and a periodic field is subtracted: let $\mathsf s$ be the $\ell$-periodic continuous function with $\mathsf s(\tau)=\tau$ for $|\tau|\le\frac\ell2-\delta_0$ that is affine on $[\frac\ell2-\delta_0,\frac\ell2+\delta_0]$, and $\zeta^\rho(y)=\sum_\lambda\omega_{\rho\lambda}\mathsf s(y_\lambda)$. Then $\zeta=\omega y$ on the closure of the cube of the margin, $\zeta^\rho$ does not depend on $y_\rho$, $\sup|\zeta^\rho|\le|\omega|\ell$ and ${\rm Lip}_\infty\zeta^\rho\le|\omega|\ell/\delta_0$. As $\zeta^\rho(y+se_\rho)=\zeta^\rho(y)$ for all $y$, $s$ and $\rho$, the clause of Lemma~\ref{C-360-8}(b) on such fields gives $\langle\mathfrak K^\sharp[\zeta];\mathcal O(f_1)\rangle=0$: for each direction $\rho$ and each plaquette $q$, the mirror normal to $\rho$ through the barycentre of $q$ keeps the measure, $a_q$ and $\zeta^\rho$, and changes the sign of $\mathfrak K^\sharp[\zeta^\rho e_\rho]$.

The field $\eta'=\chi\,\omega y-\zeta$ vanishes on the $\mathsf d$-neighbourhood of ${\rm supp}\,f_1$, where $\chi=1$ and $\zeta=\omega y$; so $f_1\eta'=0$, $\mathfrak U^\sharp[\chi\,\omega y](f_1)=\mathfrak U^\sharp[\eta'](f_1)$, and Lemma~\ref{C-360-10} applies, $\eta'$ vanishing within $\mathsf d$ of the support, with $\mathsf H(\eta')\le|\omega|[(\ell+2)\sup|\chi|+\ell\,{\rm Lip}_\infty\chi+\ell+\ell/\delta_0]$: here the margin $\delta_0$ enters $C$. In (iii), $\chi-\chi'$ vanishes within $\mathsf d$ of all the supports, so the exact identity gives $\langle\mathfrak B_{\omega,\chi}-\mathfrak B_{\omega,\chi'};\mathcal O(f_\cdot)\rangle=\beta\langle\mathfrak S_\omega[\chi-\chi']+\mathfrak K^{\rm v}[(\chi-\chi')\omega y];\mathcal O(f_\cdot)\rangle$, an insertion that vanishes near the sources, and no condition on the supports is needed.

The limit of the flux follows from the exact identity, $\nabla^s_\omega f=(\omega y)\cdot\nabla f+D_f$ with $\|D_f\|_\sharp\le C_0M|\omega|(1+\ell)a\|f\|_{C^3}$ by Taylor's formula, the bound \eqref{eq:torus-cumulant-bound} on connected functions at separated supports, and the convergence of $S^T_n$ at a fixed label with the rate $K^pL^{-K/10}$ (Corollary~\YMMref{117.1}{812}). For the quarter turn: $e^{\frac\pi2\omega_{13}}$ is a signed permutation of the axes, an exact symmetry of the lattice measure on $T_\ell$, so $S^T_n(f^{\pi/2})=S^T_n(f)$; since $\frac d{dt}f^t=-(\omega y)\cdot\nabla f^t$, $\int_0^{\pi/2}\sum_iS^T_n(\dots,(\omega y)\cdot\nabla f^t_i,\dots)dt=0$ at every $K$; with the exact identity and $\nabla^s_\omega f^t-(\omega y)\cdot\nabla f^t=O(a)$ this is the first display, and in the limit the second. (b) In $\mathfrak U-\mathfrak U^\sharp$ the entries sit where $\chi_R$ varies, in $\{\frac43R\le|y|\le\frac53R\}$ up to $2a$, or near one of the supports, which lie in $\{|y|<\frac R2\}$ at mutual distances $\ge\mathsf d_0$; all the boxes and supports lie in the sup ball of radius $2R$ about $0$, and in each cumulant the carrier of the entry is at a sup distance $\ge\min\{\frac{5R}{6\sqrt3},\frac{\mathsf d_0}{\sqrt3}\}-2a\ge\mathsf d_{\rm b}:=\frac12\min\{\frac{5R}{6\sqrt3},\frac{\mathsf d_0}{\sqrt3}\}$ from every source of that cumulant, for $K$ large, the cumulants with $f_i$ in the first entry being taken with the other sources.

So Lemma~\ref{C-360-18} in its form on all the sides, at the radius $2R$, the distance $\mathsf d_{\rm b}$ and on the sides $\ell\ge\ell'$ with $\ell'\ge4R+2\mathsf d_{\rm b}$, gives the bracket with a number free of the side, the finitely many smaller sides being covered by the form on one torus; the spin part comes with $L^{-(\frac12+c_*)K}$. The cubic member $\mathfrak U^\sharp[\chi_R\,\omega y](f)$ has its entry on top of the sources, so no wall parts them; it is Lemma~\ref{C-360-19}, with Lemma~\ref{C-360-10} for the finitely many smaller sides. Weak limit points of the torus laws inherit the bounds. Along a sequence with converging connected functions, the exact identity gives $\mathfrak U^{(j)}+\langle\mathfrak B;\mathcal O(f_\cdot)\rangle^{(j)}=\sum_qS^T_n(f_1,\dots,\nabla^s_\omega f_q,\dots)\to-\frac d{dt}S_n(f(e^{-t\omega}\cdot))|_{t=0}$, and $\mathfrak U^{(j)}\to0$.

For the volume-first limits Theorem~\YMMref{AP}{780} gives invariance under ${\rm O}(3)$, so the flux tends to $0$. With the limits one after the other: at a label $\mathsf b'\in I(g)$, on each torus $T_\ell$, $\ell\in\mathfrak L(g)$, (a) gives the limit in $K$ of the flux, $\sum_iS^{(\ell,\mathsf b')}_n(f_1,\dots,(\omega y)\cdot\nabla f_i,\dots,f_n)$, through every enclosing shell. Along the sides $\ell_k\to\infty$ the convergence $S^{(\ell_k,\mathsf b')}_n\to S^\circ_n$ is the hypothesis of the clause; the functions $(\omega y)\cdot\nabla f_i$ are Lipschitz, with the supports of the $f_i$, so the limit in $k$ is $\sum_iS^\circ_n(f_1,\dots,(\omega y)\cdot\nabla f_i,\dots,f_n)$; and this is $-\frac d{dt}S^\circ_n(f(e^{-t\omega}\cdot))|_{t=0}$, as $S^\circ_n$ is multilinear and bounded per ball in the norms $\|\cdot\|_\sharp$, the per-ball bound of Proposition~\YMMref{91.13}{709} passing to the limit once $\ell\ge16n(R+1)$ and the fixed supports lie in the central cube, and $t\mapsto f_i(e^{-t\omega}\cdot)$ is differentiable in these norms with the derivative $-(\omega y)\cdot\nabla f_i$ at $t=0$.

The integers $K^*_k\ge k$ are chosen so that for $K\ge K^*_k$ and $n\le k$ the connected functions on $T_{\ell_k}$ at the label $\mathsf b'$ lie within $\frac1k$ of $S^{(\ell_k,\mathsf b')}_n$ at all tuples of norms $\|\cdot\|_\sharp\le1$ with supports in the ball of radius $R_k=\min\{k,\frac18\ell_k\}$ about $0$ at mutual distances $\ge\frac1k$, by the rate $K^pL^{-K/10}$ of Corollary~\YMMref{117.1}{812}; every tuple of the clause lies, up to a factor, in these classes from some $k$ on, so along every sequence $K_k\ge K^*_k$ the connected functions converge to $S^\circ_n$. For given $f_i$ and $\chi_k$, raising $K^*_k$ makes the bulk term on $T_{\ell_k}$ at most $\frac1k$ as well, by (a), and the flux converges along such sequences to the limit of the display.

For slowly growing tori, the threshold in the cutoff comparison must
also be uniform in the side. Fix $n\ge2$ and a positive separation
$\mathsf d$. Take the label interval
\[
 J=[\mathsf b_-,\mathsf b_+],\qquad
 \mathsf b_-=g^{-2}+B_\sharp-L^{-1}\bar B\ge g^{-2},\qquad
 \mathsf b_+=L(g^{-2}+\bar B).
\]
It contains all first-passage labels by Corollary~\YMMref{74.2}{616}.
For large sides $\ell=2L^m$, $MR(\mathsf b_+)\le L^m$, so
the interval is admitted by Theorem~\YMMref{AR}{810}.
Its cutoff comparison uses
\[
 \phi_1=\frac1{180(3n+4)},\qquad
 u=\lceil\phi_1K\rceil,\qquad s=\lceil K/9\rceil.
\]
All its smallness conditions are products of powers of $\ell$ and
quantities decaying exponentially or faster in $K$, except for the
volume factor in the discarded mass. The proof of that theorem
writes the latter majorant as
\[
 \bar{\mathsf R}=2M^{-3}\ell^3L^{3u}
 \exp\!\left[-c'A_0^2(\log(L^u\mathsf b_-))^{2\mathfrak p_0}
                     +C_3\ell^3\right],\qquad
 C_3=\mathsf C^\Sigma_{\rm maj}+2+C_{\rm low}(\mathsf b_+).
\]
Put
\[
 c_\ell(n,g)=
 \left[\frac{c'A_0^2(\phi_1\log L)^{2\mathfrak p_0}}{2C_3}\right]^{1/3}.
\]
If $\ell\le c_\ell(n,g)K^{2\mathfrak p_0/3}$, the positive
volume exponent uses at most half of the negative exponent, since
$\mathsf b_-\ge\gamma^{-2}\ge1$ and
$\log(L^u\mathsf b_-)\ge\phi_1K\log L$.
As $2\mathfrak p_0>1$, the remaining bound decays faster than every
exponential in $K$, uniformly over these sides. The other comparison
conditions then hold beyond a threshold independent of the side;
they continue to hold at every $K'\ge K$ with the side fixed.
Theorem~\YMMref{AR}{810} and Corollary~\YMMref{117.1}{812} give
\[
 \sum_{K'\ge K}
 |S^T_{n;K'+1,\ell}-S^T_{n;K',\ell}|
 \le C\ell^{3(n+1)}K^{\mathsf p_n}L^{-K/10}
                         \prod_i\|f_i\|_\sharp,
 \qquad \mathsf p_n=2\mathfrak p_0(n+1)+\tfrac12.
\]
Under the stated side bound this is at most
$C'K^{2\mathfrak p_0(n+1)+\mathsf p_n}L^{-K/10}
\prod_i\|f_i\|_\sharp$, which tends to zero.
Thus the finite-cutoff connected functions and their spacing limits
on those tori have the same limit, and increasing the cutoffs changes
the values by at most the same tail.

For completeness, the label bounds follow directly from the canonical
flow, as in Proposition~\YMMref{74.1}{615} and
Corollary~\YMMref{74.2}{616}. Put
$\mathcal W_K(g)=\{\beta:(K,\beta)\in\mathfrak D(g)\}$ and
$X_g=L(g^{-2}+B_\sharp)$. The flow has increments of modulus at most
$B_\sharp$ and satisfies
\[
 |x_j(\beta)-L^{-j}\beta|
 \le\bar B(1-L^{-j}).
\]
At its first passage, $x_K(\beta)\le X_g<x_{K-1}(\beta)$.
Thus, for $\mathsf b=L^{-K}\beta$,
\[
 L\mathsf b>X_g-\bar B,
 \qquad
 \mathsf b\le X_g+\bar B-\bar B L^{-K}
                    \le X_g+\bar B.
\]
These imply the outer inclusion because
$L^{-1}(X_g-\bar B)=g^{-2}+B_\sharp-L^{-1}\bar B$ and
$X_g+\bar B=L(g^{-2}+\bar B)$.
To prove the inner inclusion, take
$\mathsf b\in I(g)=(L^{-1}(X_g+\bar B),X_g-\bar B]$
and put $\beta=L^K\mathsf b$. The first passage exists at some
$K'\ge1$, since $\beta>X_g$. Its label
$\mathsf b'=L^{K-K'}\mathsf b$ obeys the preceding two inequalities.
If $K'\ge K+1$, then
$L\mathsf b'\le\mathsf b\le X_g-\bar B$, a contradiction.
If $K'\le K-1$, then
$\mathsf b'\ge L\mathsf b>X_g+\bar B$, also a contradiction.
Hence $K'=K$. Existence and uniqueness of the first passage also
show that the sets $\mathcal W_K(g)$ are disjoint and have union
$(X_g,\infty)$.

\emph{Step 9: Theorem~\ref{C-360-4}(c)--(e).} (c) By Lemma~\ref{C-360-8}(e), $T_{mm}(\varphi)-\mathsf b\check D_m(\varphi)$ is $\beta$ times orbital potentials in the plane, which are $-\frac14|G_m(b)|^2$ up to densities of the order $4$, plus $\beta$ times densities of the order $4$ with weights $\le\sup|\varphi|$, plus $\mathsf b\,\mathcal O(\tilde\varphi_m)$ with $|\tilde\varphi_m|\le\frac12a\,{\rm Lip}\,\varphi$. Lemma~\ref{C-360-20}(c) bounds the first by $CK^pL^{-c_nK}\sup|\varphi|\prod_i\|f_i\|_\sharp$, Lemma~\ref{C-360-20}(a) the second by $CK^pL^{-\lambda_LK/4}$, and Lemma~\ref{C-360-20}(b) the third by $C_\epsilon a^{1-\epsilon}{\rm Lip}\,\varphi$. The sum over $m$ gives the trace. By Lemma~\ref{C-360-8}(e), $\sum_{p\ni b}a^{-1}(x_b-x(p))_\lambda\mathfrak d_{b,p,\rho}$ is symmetric in $(\lambda,\rho)$, so $(T_{lr}-T_{rl})(\varphi)=\beta\sum_b\sum_{p\ni b}(\varphi(x(p))-\varphi(x_b))a^{-1}[(x_b-x(p))_l\mathfrak d_{b,p,r}-(x_b-x(p))_r\mathfrak d_{b,p,l}]$, an insertion of the order $2$ with weights $\le\beta a\,{\rm Lip}\,\varphi$, and Lemma~\ref{C-360-20}(b) at $\sigma=0$ gives $O(a^{1-\epsilon}){\rm Lip}\,\varphi$; at $\sigma=1$ it gives $C_\epsilon a^{-\epsilon}\sup|\varphi|$ for $T_{lr}$ itself.

By Lemma~\ref{C-360-8}(b), $\mathfrak K^{\rm v}[\beta\eta]-\beta\mathfrak K^\sharp[\eta]=-\mathsf b\sum_\rho P_\rho(\nabla^+_\rho\eta^\rho)$, whose in-plane potentials are $-\frac14|G|^2$ up to the order $4$ and whose normal potentials are $\Psi_{\lambda(p)}(p)$: Lemma~\ref{C-360-20}(a), (c). The limit at a fixed label follows from the clause of Theorem~\ref{C-360-2}(d) on one plane field among action densities, with the convergence of the connected functions of $\mathcal O$ at a fixed label, since $\check D_m=\mathcal O-2\mathcal O_m$. (d) One lattice step along $e_r$ is $U(b)\mapsto W(\pm r,\mu_b)U(b)$ up to a gauge, and for every gauge invariant $\Phi$
\[
a\,v_{e_r}\Phi=\tfrac12({\rm Sh}^+_r-{\rm Sh}^-_r)\Phi+\mathcal P_r\Phi-{\rm Sh}^-_r\mathcal P_r\Phi+\mathcal R_r\Phi,\qquad|\mathcal P_r\Phi|\le C\bar a,\quad|\mathcal R_r\Phi|\le C\bar a^{3/2},
\]
${\rm Sh}^\pm_r$ the shifts by $\pm ae_r$ and $\bar a$ the largest plaquette term near the support of $\Phi$. With the null identity of Lemma~\ref{C-360-8}(a), its contact form and (c), every member of the Ward identity for the field $\chi\xi$ but $\mathfrak U^\sharp[\chi\xi](f)$ is its stated limit up to $O(K^pL^{-c_nK})$, and $\mathfrak U^\sharp[\chi\xi](f)$ is bounded by Lemma~\ref{C-360-10}, the spread being positive (the field $\chi\xi$ vanishes near the supports of the $f_i$ with $i\notin I$, and near the boundary of the chart for $\xi=\omega y,y$). For $\xi=y$, $\partial_ly^r=\delta_{lr}$, so the contraction with $T_{lr}$ contains the trace, which by (c) is $\mathsf b\,\mathcal O$; this gives the terms $|I|S_n-\mathsf bS_{n+1}(\chi,f)$. For $\xi=e_r$ and $\chi$ a function of $y_3$ alone, $\sum_lT_{lr}(\xi^r\partial_l\chi)=T_{3r}(\chi')$. (e) The equality is the exact identity for the constant field $e_\rho$ (Lemma~\ref{C-360-8}(a) with the contact terms; the stress functional of a constant field vanishes). Put $(\tau_\pm f)(y)=f(y\pm ae_\rho)$, $g_i=\nabla^s_\rho f_i$, $f^e_i=\frac12(\tau_+f_i+\tau_-f_i)$ and $h_i=(f^e_i-f_i)/a^2$. The measure is invariant under the lattice translations, so $S^T_n(\tau_+f)=S^T_n(\tau_-f)$; with $\tau_\pm f_i=f^e_i\pm ag_i$ this gives $0=2\sum_{|J|\ {\rm odd}}a^{|J|}S^T_n(g_J,f^e_{J^c})$, and with $f^e=f+a^2h$,
\[
\sum_iS^T_n(g_i,(f_l)_{l\ne i})=-\sum_{|J|\ge3\ {\rm odd}}a^{|J|-1}S^T_n(g_J,f^e_{J^c})-\sum_i\sum_{J'\ne\emptyset}a^{2|J'|}S^T_n(g_i,h_{J'},f_{\rm rest}),
\]
at most $2^n+n2^n$ terms, each with a factor $a^k$, $k\ge2$. The functions $g_i$ and $h_i$ are averages of translates of $\partial_\rho f_i$ and of $\frac12\partial^2_\rho f_i$, so their norms $\|\cdot\|_\sharp$ are bounded by those of these derivatives, and their supports lie in the $a$-neighbourhoods of the ${\rm supp}\,f_i$; the bound \eqref{eq:torus-cumulant-bound} on connected functions at separated supports bounds each term. At $n=2$, $\langle\mathcal O(\tau_\pm f_1);\mathcal O(f_2)\rangle=\langle\mathcal O(f_1);\mathcal O(\tau_\mp f_2)\rangle$ by the invariance, and the two members cancel exactly: the statement has content for $n\ge3$ and expresses the invariance under lattice translations and the failure of the Leibniz rule for $\nabla^s_\rho$. For the second display, the contact form of Lemma~\ref{C-360-8}(a) at the field $\chi e_\rho$, where $\Theta_{f_i}[\chi e_\rho]=0$ because $\chi(x_b)=\chi(x(p))$ at every plaquette $p$ at which $f_i$ does not vanish and at its four bonds $b$, and the vector weight $f_i\chi e_\rho$ is $f_ie_\rho$ for $i\in I$ and $0$ otherwise, gives the left member as $-\sum_{i\in I}S^T_n(\dots,\nabla^s_\rho f_i,\dots)$; and $\nabla^s_\rho f-\partial_\rho f$ is an average of translates of $\partial^3_\rho f$ with the total weight $\frac{a^2}6$.

\emph{Step 10: Remark~\ref{C-360-16}.} For a trilinear form that is invariant under the lattice translations along $e_\rho$, write the functions through their Fourier coefficients at $k_1+k_2+k_3=0$ and put $z_i=e^{iak_i}$; then $\nabla^s_\rho$ and $\nabla^+_\rho$ act by $(z-z^{-1})/2a$ and $(z-1)/a$, and since $z_1z_2z_3=1$,
\[
\prod_{i=1}^3(z_i-1)=\sum_iz_i-\sum_{i<j}z_iz_j=\sum_i(z_i-z_i^{-1}),
\]
which is the displayed identity. Proposition~\ref{C-360-21}(c) gives
$a^3S^T_3(g_1,g_2,g_3)\to6\mathsf b'^{-3}\int g_1g_2g_3$
for Lipschitz functions, without a separation assumption. Thus the
two sides can have order $a^{-1}$ when the supports meet. In the
rotational identity the flux through a shell separated from the
supports stays bounded, by Theorem~\ref{C-360-3}(a), whereas
the difference between the symmetric and continuous generators
contains three derivatives.

We first give the free-photon calculation. Its local part requires
a difference estimate for the curvature kernel. For equally
oriented parallel plaquettes $p,p'$, put $u=x(p)-x(q)$, using a
least lift. If $|x(p)-x(p')|\le|u|/4$ and $|u|\ge8a$, then
\begin{equation}\label{eq:free-curvature-difference}
 |C^{\rm free}_K(p,q)-C^{\rm free}_K(p',q)|
 \le C_\nabla a^3|x(p)-x(p')|\,|u|^{-4},
 \qquad C_\nabla=256\sqrt3 C_{\rm h}(1+3C_{\rm lat}).
\end{equation}
Here $C_{\rm h}$ is an absolute constant in the interior
gradient estimate for a scalar discrete harmonic function;
Theorem~6.3.8(a) of \cite{LawlerLimic} gives this estimate,
with the finitely many small radii absorbed into the constant.
Indeed, $dd^*+d^*d$ acts on each plane component as the scalar
lattice Laplacian and commutes with the projection onto exact
$2$-forms. Applied to the column $C^{\rm free}_K(\cdot,q)$,
it gives $dd^*\delta_q$, which vanishes beyond the plaquettes
sharing a bond with $q$. The column is therefore harmonic at
distance at least $2a$ from $x(q)$. Join $x(p)$ to $x(p')$
by a lattice path of length at most $\sqrt3|x(p)-x(p')|$.
This path stays at distance at least $|u|/2$ from $x(q)$.
On balls of radius $|u|/4$ about its points,
equation~\YMMref{(2246)}{824} bounds the column by
$4^3(1+3C_{\rm lat})a^3|u|^{-3}$.
The gradient estimate bounds each lattice step by
$4^4C_{\rm h}(1+3C_{\rm lat})a^4|u|^{-4}$.
Summing along the path proves \eqref{eq:free-curvature-difference}.

Let the real functions $f_i$ have compact supports inside the
margin and bounded derivatives through order five, and let
$\chi=1$ on a neighbourhood of these supports. Write $\theta_p$
for the three-colour Gaussian curvature of covariance
$\beta^{-1}C^{\rm free}_K$. In this model
\[
 a_p=\tfrac12|\theta_p|^2,\qquad
 \mathbf w(W_{b,p})=\pm\theta_p,\qquad
 \Psi_\rho(p)=-\tfrac14|\theta_{p_+}-\theta_p|^2.
\]
The sign records the plaquette orientation. The spin densities
and the other insertions are quadratic. The linear Bianchi
identity gives $\mathfrak r^\sharp=0$, hence
$\mathfrak U^\sharp=0$, and Wick's formula gives the connected
functions as rings. The exact Ward identity remains valid after
removing the terms $\frac32a\,\mathcal O(\cdot)$ coming from
the curvature of $SU(2)$; against separated sources these terms
are $O(a)$. The expression with the continuous generator is
therefore, up to $O(a)$,
\begin{equation}\label{eq:free-continuous-ward-remainder}
 \mathfrak U_{\omega,\chi}(f)
 -\sum_iS^T_n\bigl(\ldots,
       \nabla^s_\omega f_i-(\omega y)\cdot\nabla f_i,\ldots\bigr).
\end{equation}

The terms with $\chi$ in the first slot have their quadratic
insertion separated from the supports. The ring formula makes
them $O(a)$. The terms with $f_i$ in that slot are local.
For the bilinear part their kernel is
\[
 M_\rho(p,q)
   =[C^{\rm free}_K(p_+,q)-C^{\rm free}_K(p,q)]^2.
\]
Equation~\eqref{eq:free-curvature-difference} bounds it by
$Ca^8|u|^{-8}$ for $|u|\ge8a$; at the remaining distances
use the kernel bound and the projection bound on the diagonal.
The spin part has one additional factor $a$ and kernels
$C^{\rm free}_K(q,q')C^{\rm free}_K(q',p)$, where the
plaquettes $q$ meet a bond of $p$.

Expand the test functions about the midpoint of $x(p),x(p_+)$
in the first case and about $x_b$ in the second. In lattice units
these kernels have uniformly bounded moments through orders
four and two, respectively; the third-order tail in the second
case is even and has degree $-6$. Together with
\[
 \nabla^s_\omega f-(\omega y)\cdot\nabla f
   =\frac{a^2}{6}\sum_\rho(\omega y)_\rho\partial_\rho^3f
        +O(a^4),
\]
the expansion writes \eqref{eq:free-continuous-ward-remainder}
as local forms with factors $a^{-3},a^{-2},a^{-1},a^0$, plus
an $O(a)$ remainder for $n=2$ and an
$O(a\log(1/a))$ remainder for $n=3$.
The infinite-lattice coefficients are independent of $a$ by
scaling. The torus correction changes each term by $O(a)$.
At order three, configurations with two nearby entries and a
third at distance $|u|$ use the fourth moment of a kernel of
degree $-7$. Their contribution is bounded by $a$ times
the sum of $a^3|u|^{-3}$ over $a\le|u|\le1$, which gives
the logarithm. Five derivatives suffice for these Taylor
remainders and for the $O(a^4)$ error in replacing the lattice
sum of $(\omega y)\cdot\nabla\prod_i f_i$ by its integral.
The constants use the suprema of derivatives through order five.

All four local terms vanish. Inversion through the centre of
the chart preserves the Gaussian measure and the expression,
and removes the odd orders in this local expansion,
namely the coefficients of $a^{-2}$ and $a^0$.
The coefficient of $a^{-3}$ is a multiple of
\[
 \int(\omega y)\cdot\nabla\prod_i f_i=0,
\]
since $\operatorname{div}(\omega y)=0$.
For the coefficient of $a^{-1}$, use the other expression
for the same quantity: the variation of $S^T_n$ under
continuous rotations minus the bounded shell flux.
The potentially divergent terms are rotation derivatives of
cubic-invariant local forms with at most two derivatives.
A cubic-invariant symmetric tensor of rank two is a scalar
multiple of the identity, so these forms are rotation invariant.
At $n=2$ the terms can also be written as combinations of
\[
 \int(\omega y)_\rho
        \partial_\rho(\partial_\rho f_1\partial_\rho f_2)=0,
 \qquad
 \int\bigl[(\omega y\cdot\nabla f_1)\Delta f_2
            +(\omega y\cdot\nabla f_2)\Delta f_1\bigr]=0.
\]
The first identity uses that $(\omega y)_\rho$ is independent
of $y_\rho$. The first cubic-invariant local form that need
not be rotation invariant has four derivatives. Thus only the
stated remainders remain. Finally,
$a^3S^T_3\to6\mathsf b^{-3}\int$ in the Gaussian model and
\[
 \sum_i\int(\omega y)_\rho\partial_\rho^3f_i
                        \prod_{l\ne i}f_l
   =3\int(\omega y)_\rho\prod_i\partial_\rho f_i.
\]
Multiplying the difference of the generators by $a$ gives
the asserted limit of $a\mathfrak U_{\omega,\chi}$.

At order one the neighbourhood condition has a concrete role.
Take
\[
 S=[-1,1]\times[-2,2]\times[-1,1],\qquad
 \kappa(y)=3-y_1y_2,\qquad
 f_1=\kappa\,\operatorname{dist}(\cdot,\partial S)
           \ \hbox{on }S,
 \qquad f_1=0\ \hbox{off }S.
\]
The faces of $S$ are lattice planes at every cutoff, and
$f_1$ is Lipschitz with a kink at its boundary. Let
$\chi=\varsigma(\operatorname{dist}(\cdot,S))$, where
$\varsigma\in C^{1,1}$ equals $1-\mu s^2/2$ for
$0\le s\le s_0$, with $\mu>0$, and decreases to zero
inside the margin. Then $\chi=1$ on $S$ but on no
neighbourhood of it. It covers $f_1$: the barycentre and
bond midpoints of every contributing plaquette lie in $S$.
Choose $\omega_{21}=-\omega_{12}=1$, with the other entries zero.
The contact terms are means of the spin and vector insertions
and vanish by reflection, as in Lemma~\ref{C-360-8}(b). Hence
\[
 \mathfrak U_{\omega,\chi}(f_1)
 =-\beta\langle\mathfrak S_\omega[\chi]
       +\mathfrak K^{\rm v}[\chi\omega y];\mathcal O(f_1)\rangle.
\]
Its bilinear part is
\[
 -\tfrac34\mathsf b^{-1}
   \sum_{p,q,\rho}(\omega x(p))_\rho
      \nabla^+_\rho\chi(x(p))f_1(x(q))M_\rho(p,q),
 \qquad M_\rho\ge0.
\]
Away from the edges of a face with normal $n_F=\pm e_\nu$,
only $\rho=\nu$ contributes. The factor
$-(n_F)_\nu\nabla^+_\nu\chi$ is $\mu a$ times a nonnegative
function of the exterior lattice distance; $f_1$ is $\kappa a$
times the interior lattice distance. Summing the kernel over
pairs across the plane and then over the $a^{-2}$ plaquettes
per unit area gives
\[
 \tfrac34\mathsf b^{-1}\mu c_{\rm lat}
       \int_{\partial S}(\omega y)\cdot n_F\,\kappa\,dS.
\]
The positive constant $c_{\rm lat}$ is the pair sum per unit
lattice face. Cubic symmetry makes it the same for every face;
its convergence is bounded by
$\int_0^\infty\int_0^\infty st(s+t+1)^{-6}\,ds\,dt$.
Pairs a fixed distance apart contribute $O(a^2)$.
The contribution near the edges vanishes with the width of
the edge neighbourhood; the spin part is $O(a\log(1/a))$
by its $a^6|u|^{-6}$ kernel, and the torus correction is $O(a)$.
The divergence theorem gives
\[
 \int_{\partial S}(\omega y)\cdot n_F\,\kappa\,dS
   =\int_S(\omega y)\cdot\nabla\kappa
   =\int_S(y_2^2-y_1^2)=16.
\]
Consequently
$\mathfrak U_{\omega,\chi}(f_1)\to
12\mathsf b^{-1}\mu c_{\rm lat}>0$.

For the $SU(2)$ example, choose an even smooth function
$\psi:\mathbb R\to[0,1]$, equal to one on $[-1,1]$ and
zero off $[-2,2]$, and put $\mathsf c=100$ and
\[
 \varPi_{(r_1,r_2,r_3)}(y)=
 \psi((y_1-\mathsf c)/r_1)\psi(y_2/r_2)\psi(y_3/r_3).
\]
Set $g_3=\varPi_{(1,1,1)}$, $g_2=\varPi_{(3,3,3)}$,
$g_1=\varPi_{(7,7,7)}$, and
\[
 G_i(y)=\int_{-\infty}^{y_2}g_i(y_1,t,y_3)\,dt,
 \qquad f_i=G_i(\cdot+30e_2)+G_i-2G_i(\cdot-30e_2)
       \quad(i\le3).
\]
For $3<j\le n$ take $f_j=\varPi_{(15,45,15)}$ and take
$\chi=\varPi_{(31,91,31)}$, with the same $\omega$ as above
and $\delta_0=\mathsf d=1$. These fixed smooth data have
compact supports with a common point. The cutoff is one on
their closed unit neighbourhood and vanishes outside
$\{|y|_\infty\le182\}$, inside the margin of every admitted
side, since its exponent $m$ is at least seven.

Reflection in $y_1=\mathsf c$ removes the rotation about
the axis through $\mathsf c e_1$; the remaining vector field
is $\mathsf c e_2$. Translation invariance therefore writes
the left side of the exact identity as
\[
 \frac{\mathsf c}{2a}\sum_{|J|\ge3}(|J|-2)a^{|J|}
 S^T_n\bigl((\nabla^+_2f_i)_{i\in J},(f_j)_{j\notin J}\bigr).
\]
The bulk term is this expression minus the flux. The latter
has a quadratic insertion, weights $O(\beta)$ and distance
at least $1/2$ from the supports. Lemma~\ref{C-360-18}, at
$c_*=1/4$, bounds it by
$CK^{(2n+1)\mathfrak p_0}L^{K/4}$.
If $J$ contains an index above three, its derivative entry
vanishes near the support of $f_3$, so the same bound applies.
For $J=\{1,2,3\}$ the three derivatives are translates of
$g_i^K=a^{-1}\int_0^ag_i(\cdot+te_2)\,dt$ with coefficients
$(1,1,-2)$. Replace the spectators by one, discard triples
with different translates, and replace $g_1^K,g_2^K$ by one.
Each error has an entry separated by at least $1/2$ and is
again covered by that lemma. The equal translates sum to
$1+1-8=-6$. With
$I_K=(a^3/3)\sum_pg_3^K(x(p))$, translation invariance gives
\[
 S^T_n(g_3^K,1,\ldots,1)
    =3a^{-n-3}I_K(-1)^{n-1}E^{(n-1)}(\beta).
\]
Thus, up to $CK^{(2n+1)\mathfrak p_0}L^{-3K/4}$,
\[
 a_K\mathfrak U_{\omega,\chi}(f)
   =900 I_K(-1)^n\mathsf b'^{-n}\beta^nE^{(n-1)}(\beta).
\]
The inner and outer boxes of $g_3^K$ give $1\le I_K\le125$,
and its Lipschitz approximation gives $I_K\to\int g_3$.
Proposition~\ref{C-360-21}(b) gives
$(-1)^n\beta^nE^{(n-1)}(\beta)=-(n-1)!+
O(K^{2n+3}L^{-K/20})$, uniformly in $\mathsf b'\in I(g)$.
Finally $900/2=450$ and $900\cdot125\cdot2\le2.3\cdot10^5$
give the asserted two-sided bounds.

\emph{Step 11: Theorem~\ref{C-360-5}.} Let a board of side $a_i\le\frac{\ell'}2$ be placed, $\{x_3=c\}$ one of its wall planes, $\theta_c$ the mirror at it, $\tau_{a_ie_3}$ the translation by $a_ie_3$, and let $\tilde w$ be a real weight carried by the cell between $\{x_3=c\}$ and its translate by $a_i$ that vanishes within $\frac12\mathsf d$ of the walls. The centred entry $D=\mathcal O_K(\tilde w)-c_0$, $c_0$ its mean on $T_{\ell'}$, has the chessboard seminorm $\|D\|_{\rm ch}\le\mathsf q_i(\mathsf d)\sup|\tilde w|$ for $K\ge K^\#(i,\mathsf d)$, with $\mathsf q_i=e^{G_ga_i^3}{\rm Bl}_{N'_i}(k(\cdot,\mathsf d))^{1/N'_i}$, $N'_i=(\ell'L^i)^3$ and ${\rm Bl}_N(k)=\sum_\pi\prod_{I\in\pi}k(\#I,\mathsf d)$ over the partitions of $N$ points into blocks of at least two: the mean on $T_{\ell'}$ of the product of its $N'_i$ mirror copies is, by the moment--cumulant formula, the sum over these partitions of products of cumulants of copies that are $\ge\mathsf d$ apart, each bounded by Lemma~\ref{C-360-6} at the orders $\le N'_i$ (Theorem~\ref{C-360-2}(a)), and the passage between the sides costs $e^{G_ga_i^3}$ (Lemma~\ref{C-360-7}). A law $\nu$ of the lattice that is invariant under its symmetries, reflection positive at the wall planes of the boards, and dominated by the torus measures through its conditional laws obeys the reflection inequality
\[
|\nu(\overline{\theta_cF'}\,D\,\tau_{a_ie_3}F)|\le\|D\|_{\rm ch}\,\|F'\|\,\|F\|,\qquad\|F\|=\nu(\overline{\theta_cF}F)^{1/2},
\]
for bounded local functions $F,F'$ of the bonds in the half space $\{x_3>c\}$: the form $(F,G)\mapsto\nu(\overline{\theta_cF}G)$ is positive, and the Cauchy--Schwarz inequality for it, repeated at the walls of the board on the dominating tori (the chessboard estimate, Proposition~\YMMref{86.3}{681}, with the Cauchy--Schwarz inequalities of Proposition~\YMMref{91.4}{700}), bounds the left member by the chessboard seminorm of $D$ read along these tori. Every periodic Gibbs state has the three properties: it is invariant under the lattice symmetries, it is reflection positive at all the families of mirrors (Proposition~\textup{\YMMref{103.4}{758}}), and it is dominated by the torus measures (Proposition~\YMMref{91.9}{705}). Recentring from $c_0$ to $\nu(\mathcal O_K(\tilde w))$ costs $\mathsf q_i(\mathsf d)\sup|\tilde w|\,\|F'\|\,\|F\|$ more, by Lemma~\ref{C-360-7} at $q=1$ and $|\nu(\overline{\theta_cF'}\,\tau_{a_ie_3}F)|\le\|F'\|\,\|F\|$; whence the factor $2$:
\[
|\nu(\overline{\theta_cF'}\,[\mathcal O_K(\tilde w)-\nu(\cdot)]\,\tau_{a_ie_3}F)|\le2\mathsf q_i(\mathsf d)\sup|\tilde w|\,\|F'\|\,\|F\|\qquad(K\ge K^\#(i,\mathsf d));
\]
eight shifted boards cover all positions. Along subsequences the limits exist (Theorem~\ref{C-360-2}(c)), and in the reconstruction the bound is what Lemma~\ref{C-360-12} asks, and that lemma gives $B^{(m)}$, the density and the $L^1$ bound. The reconstruction of the limits of Gibbs sequences at $0<g\le\gamma$ has $H\ge0$ (Theorem~\YMMref{AP}{780}), the hypothesis of Lemma~\ref{C-360-12} in the direction $x_3$; and it has $H\ge|P_k|$, $k=1,2$, for the commuting momenta $P_k$ that generate the translations $U(\vec x)$ in the directions orthogonal to $e_3$, the hypothesis of the last sentence of that lemma in the directions $x_1,x_2$: every periodic Gibbs state obeys the five-point inequality \textup{(\ref{IV-54-29})} of Lemma~\ref{lem:spectral-steps} at the mirror normal to $e_3$ and the direction $e_k$, which passes to the limit and puts the spectral measure of every class on $\{\cosh(tE)+\cos(tq_k)\ge2\}$ for a dense set of $t>0$, hence on $\{E\ge|q_k|\}$, as $\cosh(tE)+\cos(tq_k)=2+\frac12t^2(E^2-q_k^2)+O(t^4)$.

The same holds for the reconstruction of a tensor limit. So Lemma~\ref{C-360-12} applies in the three directions, which gives the density in all three coordinates and the $L^1$ bound in them. For the bound on $\|B^{(m)}\|$: for $g\le\gamma^{\rm sd}_W$ the chessboard seminorm of a centred tensor entry in a cell of side $a_i$ is at most $2\mathsf C_TY_i^{\mathsf q_T}a_i^{-11/2}\mathsf b^{-1}|H|_\infty$, uniformly in the cutoff, the datum, the board, the side and the state (Section~\ref{C-355-1}, \eqref{C-355-8} of Theorem~\ref{C-355-7}(a)); with Lemma~\ref{C-360-12} every tensor limit has the insertion bound with every $p>\frac{17}2$, and $B^{(m)}$ is the compression to the scalar subspace of the point operator at $e=\mathbf 1-2E_{mm}$, $|e|=\sqrt3$; that every limit with converging labels arises from a tensor limit along a further subsequence is the extraction of Theorem~\ref{C-360-3}(g).

For the behaviour under turns: by Section~\ref{cov-sec}, every tensor limit at $0<g\le\gamma^{\rm sd}_W$ obeys $\sigma(RX)=\sigma(X)$, $R\in{\rm O}(3)$, at tuples with one non-scalar entry among scalar ones (Theorem~\ref{cov-thm-rider}, which gives Theorem~\ref{cov-thm}(c) at these tuples); $R(\mathbf 1-2E_{33})R^{\rm T}=\mathbf 1-2nn^{\rm T}$ with $n=Re_3$, and the sum over $\pm\alpha$ cancels the off-diagonal part of $nn^{\rm T}$, leaving $2\cos^2\alpha(\mathbf 1-2E_{33})+2\sin^2\alpha(\mathbf 1-2E_{ii})$; the entries of $nn^{\rm T}$ are combinations of $1$, $\cos2\alpha$ and $\sin2\alpha$, which gives the form $c_0+c_1\cos2\alpha+c_2\sin2\alpha$ and the second derivative.

\emph{Step 12: Theorem~\ref{C-360-13}.} (a) Put $\lambda=L^{-(k-j)}$, $d=d_j(Z)$ and $S=\sup|E|$; the left member is at most $2S$, and the numbers $a_3$, $C^Q_2$, $\bar C_2=\max\{C^Q_2,1\}$, $c^Q_3$, $c_s=c^Q_s$ and $\lambda_Q$ are those of the extended list. The trivial cases come first: if $\lambda>\lambda_Q$, then $2S\le2\lambda_Q^{-4}\lambda^4S$; if $(d+4)\lambda>c_s$, then $2S\le2((d+4)\lambda/c_s)^4S\le32c_s^{-4}(d+2)^4\lambda^4S$, as $d+4\le2(d+2)$. Let now $\lambda\le\lambda_Q$ and $(d+4)\lambda\le c_s$. Then $\lambda\le1/(4a_3)$, and on a cube that holds the box of $Z$, whose sites lie within the $\ell^1$-distance $l=3(d+4)M$ (level-$j$ units) of its least site, the definition of the cores gives a gauge in which the real part of the potential obeys $|A|\le\lambda$ and $|\nabla A|\le\lambda^2$ in level-$j$ units; with the clause on the complex part this gives a complex potential $A^0$ of the sizes $\lambda$ and $\lambda^2$, up to the complex gauge transformations $g_t=e^{itm}$, $0\le t\le1$, with $|m|\le l\lambda a_3\le3Ma_3c_s\le\frac14$, as $c_s\le1/(12Ma_3)$, so of the size $e^{1/4}$.

The constant part, of the first order, is removed by $g_1$: a term is invariant under real gauge transformations and holomorphic, so it is invariant under every complex gauge transformation that is reached by a path inside its domain, and the path $t\mapsto g_t$ stays there, its triples lying in the cores and below the third radius. What is left, $\tilde A$, vanishes at the least site of the cube and has the gradient $|\nabla\tilde A|\le C^Q_2\lambda^2\le\bar C_2\lambda^2$, so that $|\tilde A|\le l\bar C_2\lambda^2=3(d+4)M\bar C_2\lambda^2$ on the cube; the disc $\zeta\mapsto(\mathbf 1,\zeta\tilde A,\zeta\tilde J)$ of the radius $R=c^Q_3/(l\bar C_2\lambda^2)=c^Q_3/(3(d+4)M\bar C_2\lambda^2)$, of the order $((d+2)\lambda^2)^{-1}$, stays in the domain, because the H\"older number depends only on the curl of the potential, up to squares, and scales with the power $\frac52>2$ under $\rho^{k\to j}_Z$; at the trivial field the first derivative of a gauge invariant term vanishes, since the adjoint representation of ${\rm SU}(2)$ has no invariant vector; and the Schwarz lemma gives $|E((\mathbf 1,\tilde A,\tilde J))-E(\mathbf 1)|\le2SR^{-2}=2l^2\bar C_2^2(c^Q_3)^{-2}\lambda^4S\le72M^2\bar C_2^2(c^Q_3)^{-2}(d+2)^2\lambda^4S$, since $l\le6(d+2)M$ (if $R\le1$ the bound $2S\le2SR^{-2}$ holds already).

This is the clause with $(d+2)^2$, and with the trivial cases it gives 
\[
C_Q=\max\{2\lambda_Q^{-4},32c_s^{-4},72M^2\bar C_2^2(c^Q_3)^{-2}\}
\]
 and the power $(d+2)^4$. For ${\rm U}(1)^3$ the field strength scales like $\lambda^2$ and the linear term survives. The power $4$ is attained by a term of the second order in the field strength at one plaquette, against a triple with a constant field strength of the largest admitted size. (b) A block of the level $k$ holds $L^{3(k-j)}$ regions of the level $j$, each constant up to $L^{-4(k-j)}$ by (a); the count of regions among the construction results with the rate $\kappa_{\rm lf}\ge4\log18432$ of (b) gives the bound, uniformly in $k$. (c) The complex centre of each Gaussian integral is kept as a free argument, one $1$-form for every earlier level, in a polydisc.

The first order of a plane weight gives $\sup|w|L^{-5/2}{\rm hol}$, hence the third radius, which is small at the true points and stable under the shift system among the construction results; the weight and the trees enter as potentials, so that the action of the step carries no explicit weight; a tree of the level $j$ weighs $L^{k-j+1}$ seen from the level $k$, there are $L^{3(k-j)}$ of them and each is smaller by $L^{-5(k-j)}$; the non-classical terms stay at their levels and are lifted afresh at every step by (b), at the rate $r_1$, which (b) admits since $r_1>\mathsf C_{\rm cs}>4\log18432$. Each step is then the identity of the statement, classical times ideal times the exponential of a tiny term, with $c_k(W)\Xi_0[V_{\mathsf m}](Q)=\Phi_{k+1}[\omega](W)$; at $k=0$ the step is written with bond labels.

On the plain cores the H\"older number is at most $\alpha_0$ times a factor that grows like $L^{n/2}$ with the level, while $\alpha_2/\alpha_0\ge4.99\cdot10^6e^{L/4}$; so the condition of the third radius follows from the plain core at the levels $n\le(\frac L2+29.39)/\log L$. At $k=0$ the weight enters linearly through the bond labels, and the centre, the solution of the linear system of the step with a source of the size $\sigma$, is at most $K_Z\sigma$. The extended list, stated before Theorem~\ref{C-360-13}, places each new constant after the constants on which it depends and before $\gamma$; its four new lower bounds are laid on the chosen constants $\mathsf C_{\rm ex}$, $\kappa$, $\bar\kappa$ and $M$, which are free upward, each by a function of earlier constants, its new conditions on $\gamma$ are ceilings whose left members tend to $0$, and it lays no condition on $L$ itself; so it keeps the admissible shape, its choices exist, and they restrict to choices of the lists of Theorems~\YMMref{D}{8}, \YMMref{AJ}{673} and~\YMMref{AD}{607}.

At every choice $L>4\cdot10^{19}$: the retained floor $L\ge L_\Pi=2\mathsf C_{\rm ex}(1+4\beta_{\rm s})/(1-2\delta_{\rm w})\ge2\mathsf C_{\rm ex}$ (equation~\YMMref{(904)}{418}, with $\beta_{\rm s}>0$ and $0<\delta_{\rm w}<\frac12$), the new lower bound $\mathsf C_{\rm ex}\ge2\mathsf C_{\rm cs}$, and $\mathsf C_{\rm cs}\ge24C_{\mathrm E}^2e^{24}$ by the choice preceding equation~\YMMref{(1332)}{552}, $C_{\mathrm E}=4096$, give $L\ge4\mathsf C_{\rm cs}\ge96C_{\mathrm E}^2e^{24}>4\cdot10^{19}$. The further conditions named in the statement are of three kinds. The tower hypotheses of Section~\ref{sec:analytic-core-complements} and the inequality $\alpha_0+\alpha_1\le a^\sharp$, even $2(\alpha_0+\alpha_1)\le a^\sharp$, are closed inequalities between constants that hold at every choice of the list of Theorem~\YMMref{D}{8}, and so at every choice of the extended list, which restricts to such a choice.

The standing conditions of the frame of Definition~\ref{def:running-source-history} follow, condition by condition, from members of the lists of Theorems~\YMMref{D}{8}, \YMMref{AJ}{673} and~\YMMref{AD}{607}, to which the choices of the extended list restrict. The locality for all pairs of regions and the gauge covariance at all decoupling parameters are properties of the Gaussian structure among the construction results as it is formed, at all admitted constants: its structures form a local system, and they are covariant under real gauge transformations at all decoupling parameters $s$ of the closed polydisc (Theorem~\YMMref{U}{378}(2) and the covariance of the operators in \YMMref{(765)}{378}: restrictions commute with the unitary gauge actions, as do inversion and the functional calculus used to form the covariance and coupling).

\emph{Step 13: Theorem~\ref{C-360-14}.}
(a) Proposition~\ref{prop:weighted-critical-response} proves the
uniform critical-background and Hessian assertions. We give the
remaining comparisons and check the actual step data.

At $U_0=1$ with all regions the whole torus, the quadratic form
for plane weights is
\[
 \langle B,Z_\epsilon B\rangle
 =\min_{Qd=B}
 \{(\partial d,\mathsf a\,\partial d)+\|RD^{\mathsf T}d\|^2\},
 \qquad 1-\epsilon\le\mathsf a\le1+\epsilon.
\]
Each form before minimization lies between $1-\epsilon$ and
$1+\epsilon$ times the unweighted form, including its
nonnegative slice term. Taking the infimum gives the claimed
comparison of block Hessians. For the single mixed tensor term,
$2|\langle f_{\mu\nu},f_{\mu\lambda}\rangle|
\le|f_{\mu\nu}|^2+|f_{\mu\lambda}|^2$ gives the same form
comparison in the stated normalization. Adding the square of the
linearized block averages preserves it. Thus the linearization
of the constrained system is invertible at zero. On each fixed
finite lattice the analytic implicit-function theorem gives
the stated family near zero.

The reason a uniform H\"older estimate needs the H\"older norm
of the weight can be seen in this same quadratic problem.
Fix $B\in Y_0$, $B\ne0$, and write
$h_\epsilon=d_bA\cdot B$ at $b=0$.
Take the torus side $N=2L^{m+K-k}$ in level units, as in the
construction; it has the even number $N/\eta_k=2L^{m+K}$
of bare sites in each direction.
Choose a plane $(\mu\nu)$ in which
$F^0=\partial h_0$ has largest component norm, and put
$v_{\mu\nu}(x)=\epsilon(-1)^{x_\mu/\eta_k}$,
with all other plane weights zero. Then
\[
 Qh_\epsilon=B,\quad RD^{\mathsf T}h_\epsilon=0,\quad
 \partial^{\mathsf T}\mathsf a\,\partial h_\epsilon
                              =Q^{\mathsf T}Z_\epsilon B.
\]
The supremum of the source on the right is at most
$(1+\epsilon)C_{\rm tr}c_Z\|B\|$, with
$c_Z\ge\|Z_0\|$ and the row-to-supremum norm $C_{\rm tr}$ of
$Q^{\mathsf T}$ depending only on $L$.
Since $\delta=h_\epsilon-h_0$ has $Q\delta=0$, subtracting
the two variational equations and testing by $\delta$ gives
\[
 \|\partial\delta\|^2+\|RD^{\mathsf T}\delta\|^2
 =-\epsilon\bigl((-1)^{x_\mu/\eta_k}
       (\partial h_\epsilon)_{\mu\nu},
                         (\partial\delta)_{\mu\nu}\bigr).
\]
It follows that
\[
 \|(\partial h_\epsilon)_{\mu\nu}\|
 \ge(1+\epsilon)^{-1}\|F^0_{\mu\nu}\|
 \ge(1+\epsilon)^{-1}(\gamma^{\rm pc}/3)^{1/2}\|B\|.
\]
On the torus of side $N$ in level units, its supremum
$f_\epsilon$ is therefore at least
$(1+\epsilon)^{-1}(\gamma^{\rm pc}/(3N^3))^{1/2}\|B\|$.
At two neighbouring sites in direction $\mu$, use the bond
equation at the site carrying the factor $1+\epsilon$.
It expresses $2\epsilon(\partial h_\epsilon)_{\mu\nu}$
as $\eta_k$ times the source minus two first differences of
curvature components, with coefficients at most $1+\epsilon$
and $1$. Consequently
\[
 2\epsilon f_\epsilon\le
 (2+\epsilon)\operatorname{osc}(\partial h_\epsilon)
       +\eta_k(1+\epsilon)C_{\rm tr}c_Z\|B\|.
\]
Here $\operatorname{osc}$ is the largest modulus of the change
of one curvature component between neighbouring bare sites.
For $\eta_k\le c_\eta\epsilon N^{-3/2}$ this gives
\[
 [\partial h_\epsilon]_{\beta_H}
 \ge\frac16\Bigl(\frac{\gamma^{\rm pc}}{3N^3}\Bigr)^{1/2}
                       \epsilon\eta_k^{-\beta_H}\|B\|.
\]
Here the seminorm includes neighbouring sites. At fixed $N$
this diverges as $k\to\infty$, although the weight has supremum
$\epsilon$. Since the background vanishes at zero block datum,
the example concerns its derivative there. Cauchy's estimate
would bound that derivative if the backgrounds had a uniform
H\"older bound on a common block-datum ball.

The constant-weight assertion follows from the exact identity
\[
 F+T_{(h,v+c)}=(1+c)(F+T_{(h,v)/(1+c)}).
\]
Multiplying the critical equation multiplies its multiplier by
$1+c$ and leaves its constraint and gauge unchanged.
Uniqueness in Proposition~\ref{prop:weighted-critical-response}
then gives the asserted identity of critical fields.

For the history application use the collared sequence before the
step and its completion by the next large-field region.
The background of the run is its real unweighted critical
configuration at the prescribed averages. The block field $A_k$
is read about its finer average in the axial gauge, and its row
datum is $b=JA_k$, extended by zero to the remaining rows.
The regular-background and chart estimates for these actual
step fields, Lemma~\YMMref{37.5}{287},
Proposition~\YMMref{53.12}{457}, and the fibre-chart bounds in
Section~\ref{sec:sf-chart-interface}, give
\[
 \alpha\le2C_{\mathfrak P}\mathfrak M B_3\varepsilon(x_t),
 \qquad
 \sup|b|\le2C_\chi\mathsf c_{\rm fl}\delta(x_k),
 \quad C_{\mathfrak P}=504L^2B^{\rm p}_1 .
\]
For the natural-class assertion, apply Lemma~\YMMref{37.5}{287}
with the separated witness $m=\mathfrak M$, increasing the
common block floor if necessary. The layerwise regularity
radius of the real background is at most $2B_3\varepsilon(x_t)$,
so that lemma gives the displayed natural-class radius on the
completed sequence. The chart converts the fluctuation field
to the row datum with the displayed factor $2C_\chi$.
Choose
\[
 a_N\le\min\{a_*,a^{\rm pc},\gamma^{\rm pc}/(4e_0)\},
 \quad \rho_N^+\le\frac{\gamma^{\rm pc}\rho_*}{4\kappa_0},
 \quad z_N^+\le\frac{\gamma^{\rm pc}z_*}{4\kappa_0},
\]
also keeping $\rho_N^+<\rho_*$ and $z_N^+<z_*$.
Impose the ceilings, for all $x\ge\gamma^{-2}$,
\[
 2C_{\mathfrak P}\mathfrak M B_3\varepsilon(x)\le a_N,
 \qquad 2C_\chi\mathsf c_{\rm fl}\delta(x)\le\rho_N^+ .
\]
Their left sides tend to zero with $\gamma$.
Thus $e_0\alpha\le\gamma^{\rm pc}/4$ and
$\kappa_0s\le\gamma^{\rm pc}/4$ at every step $t\ge2$,
which gives the asserted half-floor.

At the first step $k=0$, the lattice spacing is one, the block
averages are the identity and the slice projection vanishes.
There is no minimization at this step. At the real regular
reference the fluctuation floor differs from the action Hessian
by the second derivative of the block constraint paired with
the current. Its loss is at most $8c_{\rm pc}\alpha$, where
$c_{\rm pc}=\sqrt2C_Kn_{\rm pc}$ and $n_{\rm pc}=36L^3$.
For a complex block field $A^{(0)}$ with
$r_A=\sup|A^{(0)}|$, differentiating the four link factors in
each plaquette and summing their incident bonds gives the
following three operator-norm bounds: the change in the plain
action Hessian is at most $128r_Ae^{4r_A}$; the tensor term
has norm at most $512e^{8r_A}\|h\|_\infty$; and the plane term
has norm at most $32e^{4r_A}\sup|v|$.
Therefore
\begin{align*}
 \operatorname{Re}(F+T_{\mathbf h})''(JA^{(0)})[J\bar B,JB]
 \ge\bigl(&\gamma^{\rm pc}-8c_{\rm pc}\alpha
       -128r_Ae^{4r_A}\\
       &-512e^{8r_A}\|h\|_\infty
       -32e^{4r_A}\sup|v|\bigr)\|B\|^2 .
\end{align*}
This uses the real natural pair on the two-region sequence
$(\Omega_0,\Omega_1)$ and $\mathfrak M\ge M^{\rm pc}$,
exactly as in the unweighted positivity argument above.
Take
\[
 a_N^1\le\min\{a^{\rm pc},\varpi_\chi,
                       \gamma^{\rm pc}/(64c_{\rm pc})\},
 \quad r_N^1=\min\{1/4,\gamma^{\rm pc}/(1024e)\},
 \quad z_N^1=\gamma^{\rm pc}/(2048e^2).
\]
The ceilings
\[
 2C_{\mathfrak P}\mathfrak M B_3\varepsilon(x)\le a_N^1,
 \qquad
 2C_\chi\mathsf c_{\rm fl}\delta(x)
                         \le\min\{r_N^1,r_\chi\},
 \qquad x\ge\gamma^{-2},
\]
put the first-step data in these ranges.
The losses through $\alpha$, through $r_A$, and through the
weights are at most $\gamma^{\rm pc}/8$,
$\gamma^{\rm pc}/8$, and $\gamma^{\rm pc}/4$ respectively.
This leaves $\gamma^{\rm pc}/2$ for the stated complex data.

(b) Lemma~\ref{lem:weighted-local-hessian} identifies the perturbation as $D^{(1)\mathsf T}vD^{(1)}+Z_{1,v}$. The complete mean-square argument is Proposition~\ref{prop:weighted-system-L2}: its proof first controls the divergence-input system with the actual averaging rows, inserts $vD^{(1)}A$ as a source, and then localizes with a divergence-form cutoff. It proves both inverse existence and the termwise walk bounds without taking a difference of $v$.

The grouping by visited pieces after Proposition~\ref{prop:weighted-system-L2} gives the connected-set factor in (b). (c) The label supplies its background regular at the radius $C_m\varepsilon_k$ with the envelope $e(x)=C_m\varepsilon_{j(x)}$. Two of the three ceilings of the statement are the label's: $\sup_{x\ge\gamma^{-2}}C_m\varepsilon(x)\le\alpha_\star$ bounds the radius of regularity and the envelope at every level by the number $\alpha_\star$ of $L$, and $\sup_{x\ge\gamma^{-2}}\max\{w_{\rm G},w_{\rm lin}\}(\log x)^{-r}\le1$ gives $MR(\gamma^{-2})\ge\max\{w_{\rm G},w_{\rm lin}\}$, $w_{\rm G}=4c_U^{-1}(7\log L+\log2)$; under these, and with $t({\rm an})<t_0\le t_{\rm lin}(L)\le\frac14$, Lemma~\ref{C-360-22} and the locality below are applied; the third, the membership ceiling, is read below.

Use Lemma~\ref{C-360-22} and assume the complex scalar transfer bound \eqref{eq:complex-scalar-transfer}. The pair $(A,\omega)$ is the fixed point of the critical-point contraction built on \eqref{eq:uniform-linear-response-bound} and the local H\"older bounds. To localize it, write $G'=(\Delta_D+Q'^*\mathsf a'Q')^{-1}$ in the scalar notation and retain $Q'Q'^*=1$. Since $\ker R=\operatorname{ran}(G'Q'^*)$, the slice condition is equivalently the local equation $(1-Q'^*Q')\Delta_DD^*A=0$. Cut the far field off with a slow cutoff ${\rm ch}$ of width $w(L)$ and bounded scaled first and second differences. In the augmented scalar system the cutoff error is supported where ${\rm ch}$ varies. Its inverse is the actual compression $\mathbb K^{[2]}$ described in Definition~\ref{hyp:complex-scalar-transfer}; the two nonlinear remainders are precisely $N$ and $q$ there. The additional estimate \eqref{eq:complex-scalar-transfer} therefore makes the scalar correction a contraction in the exponentially weighted norm. It repairs the slice and the induced block gauge and retains exponential decay away from the cutoff. The repaired field nearly solves the critical-point system about $U$. Applying its contraction gives the stated locality estimate with $c_U$ and $K_U=7K_{\rm lin}^2(1+2L^3)^2e^{c_Uw_*}$, after the permitted enlargement of $K_{\rm lin}$, where $w_*=w+21L$.

Continuity: a uniform limit of iterates, with the error $\frac23\,4^{-n}$. Uniqueness is in the normalized class: the gauge functions preserve the block averages and induce the identity at the tops. If $v_j[A']$ denotes the accumulated averaging gauge, this means
\[
 u(y)=v_j[A'^u](y)v_j[A'](y)^{-1}
 \quad\text{at every top }y\text{ of level }j.
\]
The scalar transfer carries a sufficiently small critical point into the slice. Its multiplier is absorbed in the fixed-point equation using the same small-radius condition, so the resulting pair lies in the contraction ball and equals the constructed pair. Two normalized gauges with this property agree at every top. Their quotient preserves $W$, hence is covariantly constant along its bonds. It is therefore the identity along every pillar from a top, and those pillars cover the bare sites. This proves the stated uniqueness. Energy: apply \eqref{eq:weighted-field-energy-nonlinear} to the fixed point, with its actual row forcing $-C_2(U,A)$. Since $|\xi^2F_U(p)|^2\le2a_p(U)$, the local plaquette product estimate gives $\sum_p|\partial U(p)^{-1}\partial W(p)-1|^2\le C(L)|{\rm an}|_\infty^2\sum_pa_p(U)$. Cauchy--Schwarz with the unperturbed plaquette defects gives the stated action comparison. Full covariant gradient: let $x_*$ lie in a layer $n\ge2$, of the scale $l_n$, and pass to the axial gauge in the cube of radius $\frac14l_n$ about $x_*$, built from the transports of $U$ along coordinate paths; the transports differ from $\mathbf 1$ there by $O(e(x_*)l_n^{-1})$ in modulus and in H\"older-$\frac12$ quotient.

In this gauge the potential $A$ solves a flat system for its curl and its divergence, $dA=\Phi$, $d^*A=\varpi$, where $\Phi$ is the curvature and $\varpi$ the divergence of $A$, each up to products of $A$ with the transports minus $\mathbf 1$; by the curvature norm of Lemma~\ref{C-360-22} and the slice estimate for the divergence below, $\Phi$ and $\varpi$ are bounded and H\"older-$\frac12$ in the scale of the layer, with the sizes $Kt({\rm an})e(x_*)l_n^{-2}$ and, for the H\"older quotients, $Kt({\rm an})e(x_*)l_n^{-5/2}$, $K$ a constant of $L$ alone. The divergence is even Lipschitz: on the slice $RD^*A=0$, so $D^*A=G'Q'^*m$ for some $m$, and the range of $G'Q'^*$ has first differences bounded in the scale of the layer, $|\eta_k^3(DD^*A)(b)|\le Kt({\rm an})e(x_b)l(b)^{-3}$.

After a cut-off at the radius $\frac18l_n$ whose first and second differences are $O(l_n^{-1})$ and $O(l_n^{-2})$, the identity $\Delta=d^*d+dd^*$ writes each first difference of $A$ as second differences of the lattice Green function of the Laplacian applied to the cut data; an interior estimate of Schauder type for these kernels against H\"older-$\frac12$ data, on a cube of radius $\sim l_n$, bounds all nine covariant first differences, $l(x)^2|\eta_k^2(\nabla_\kappa A_\nu)(x)|\le K_5\,t({\rm an})\,e(x)$ with $K_5$ a function of $L$ alone; in the layers $0,1$ the bound $|\nabla_\kappa A_\nu(x)|\le|A_\nu(x+e_\kappa)|+|A_\nu(x)|$ suffices. Membership: the plain triple of $U$ lies in the cores at half the radii (Corollary~\YMMref{58.4}{508}, under its support, radius and smallness hypotheses); the moved triple is the plain triple with the displacement substituted, and in the units of each layer $\lambda$ the displacement, its first differences, the change of the plaquette variables and the change of the current are at most $K'_5\,t({\rm an})\,C_m\varepsilon_\lambda$, by Lemma~\ref{C-360-22} and the gradient estimate, where $C_m\varepsilon_{j(x)}$ is the envelope with which the label supplies its background, $K'_5$ is a number formed with $L$ through $K_{\rm lin}$ and $K_5$, and $t({\rm an})<t_0\le t_{\rm lin}(L)\le\frac14$, Lemma~\ref{C-360-22} being applied at the anisotropy ${\rm an}$.

These are small against the radii, as the substitution asks, under the third ceiling of the statement,
\[
\sup_{x\ge\gamma^{-2}}2K'_5C_m\,\varepsilon(x)\le b_2,\qquad\varepsilon(x)=A_0x^{-1/2}(\log x)^{\mathfrak p_0}\ \text{(equation~\YMMref{(634)}{339})},
\]
$b_2$ being the radius of Proposition~\ref{prop:analytic-three-radii} among the construction results, with which the substitution with three pairs of radii carries the cores at half their radii into the cores at three quarters of them: under this ceiling they are at most $\frac18b_2$, and the ceiling is of the admissible shape, its left member tending to $0$ as $\gamma\to0$. Then this substitution puts the moved triple in the cores at three quarters of the radii, all the tower clauses included. Frames retain the following six admissions, supplied by the ordered choice of Proposition~\YMMref{69.1}{595} and its source continuation in Section~\YMMref{82}{665}: (1) the geometric and background hypotheses used in that choice; (2) consistent data and a coupling sequence with flow bound $B_\sharp$; (3) at every level $n$, $x_n>\bar B$, $y_n>2\mathfrak p_0$, $g_n\le\gamma$ and $2\varepsilon_n\le\bar\varepsilon$; (4) $\varepsilon_q\le\varepsilon_\flat$ at every level $q$; (5) at every level $t$,
\[
 C_{\rm reg}\varepsilon_t\le\varepsilon_{\rm v},\qquad
 L^2C_{\rm reg}\varepsilon_t\le c^{\rm SO}_2,\qquad
 2L^2C_{\rm reg}\varepsilon_t\le\bar\varepsilon;
\]
and (6) $x_t\ge2B_\sharp$ at every level. The layer and support conclusions used here are those of Corollary~\YMMref{58.4}{508}, with these actual admissions; the pivot evaluation additionally retains the ceilings preceding Lemma~\YMMref{66.1}{583}.

The real and sourced constructions use the same coupling sequence,
tests and canonical real domains, by Lemmas~\YMMref{82.4}{670}
and~\YMMref{82.5}{670}. Condition (1) is part of the geometric
choice in Proposition~\YMMref{69.1}{595}; condition (2) is the
admissible-sequence and room hypothesis of Theorem~\YMMref{D}{8}.
The proof of Proposition~\YMMref{69.1}{595} imposes
$x\ge2B_\sharp$, $x>\bar B$ and $\log x>2\mathfrak p_0$
uniformly for $x\ge\gamma^{-2}$. The additional ceilings in the
proof of Lemma~\YMMref{82.5}{670} give
$2\varepsilon_t\le\bar\varepsilon$,
$\varepsilon_t\le\varepsilon_\flat$ and
$L^2C_{\rm reg}\varepsilon_t\le c_2^{\rm SO}$ at every level,
as well as the remaining graded-field conditions there.
These are (3)--(6). (d) The Hessian of the Lagrangian at a critical background annihilates, against every field, a gauge direction whose gauge transformation preserves the block averages of every configuration. For a gauge direction whose linearized block average vanishes only at the background, it annihilates the fields in $N(Q)$; against other fields the remaining term lies in the range of $Q^{\mathsf T}$ and is absorbed by changing the multiplier. Thus Lemma~\ref{C-360-22}, applied on the longer sequence in its own gauge, inverts the linearized system in the axial gauge. Its weighted maximum bound is that lemma. Its square-sum bound follows from Lemma~\ref{lem:weighted-field-energy}; the multiplier is bounded by testing against a row lift $H\varsigma$ in
\[
 \langle\omega,\varsigma\rangle
   =\langle T_\tau A-f-D_1^{\mathsf T}Y,H\varsigma\rangle.
\]
All constants depend only on $L$.

With Cauchy's inequalities for the remainder of the system this gives exactly one centre in the polydisc $|\mathsf n(b)|<\theta_3\mathsf b_k$ of the statement, $\mathsf b_k=A_1(\log x_k)^{\mathfrak p_0}$ the threshold of the level $k$, and it lies in the half polydisc, $|\mathsf n(b)|\le\frac12\theta_3\mathsf b_k$, once the following hold. The envelope of the step is that of (c) with $C'/L$ in the place of $C_m$, where $C'=\max\{LC_m,2LC_\chi(1+2\theta_3)A_1/(\rho_{\rm u}A_0)\}$ is the number of the statement and $\rho_{\rm u}=\rho_{\rm u}(L)$ the radius of the ball in which the centre is unique: the first member keeps the envelope at least that of (c), and the second puts every complex block datum of the box of the step, read in the chart of the fibre, inside that ball in the norm of the envelope; this is the one number that is enlarged, and it is read by the two conditions that follow.

The size of the anisotropies is below the number of the statement, $t_\theta=\min\{t_0,t_{\rm u}(L),c_{\rm ce}\theta_3A_1/A_0\}$, where $t_{\rm u}(L)$ is the number below which the moved backgrounds with a complex block datum are formed and $c_{\rm ce}>0$ is formed with $L$ and $C'$; so $t_\theta\in(0,t_0]$ is formed with $\theta_3$ and with constants placed before $\gamma$. And $\gamma$ is below ceilings of the shape $\sup_{x\ge\gamma^{-2}}C'\varepsilon(x)\le c(L)$, of the admissible shape, among them the ceiling $C_\chi(1+2\theta_3)\delta(x)<\frac18r_c$, $x\ge\gamma^{-2}$, of the statement, with the numbers $C_\chi$, $r_c$ of $L$: it keeps the box $|{\rm Re}\,B_1(b)|+|{\rm Im}\,B_1(b)|<(1+2\theta_3)\mathsf b_k$ of the complex block data of the step, of the unscaled radius $(1+2\theta_3)\delta(x_k)$, and with it the polydisc $|\mathsf n(b)|<\theta_3\mathsf b_k$, of the unscaled radius $\theta_3\delta(x_k)$, inside the ball of the radius $r_c$, in the unscaled norm, in which the system of the centre has exactly one solution.

The centre is $\mathsf n=g_k^{-1}a''$, $g_k=x_k^{-1/2}$, where $a'$ is the finer block average of $W[{\rm an};b_0]$ read about that of $W[0;b_0]$, the two backgrounds having the same row datum, and $a''$ is this relative field turned into the block axial gauge; $\mathsf n$ is holomorphic in ${\rm an}$, real at real ${\rm an}$ and real $b_0$, and $0$ if ${\rm an}$ is one number. By the bound of (c) on the displacement in the units of the level, read with the envelope $(C'/L)\varepsilon_j$, and the passage to the block axial gauge, $|a''(b)|\le K''t({\rm an})(C'/L)\varepsilon_k$ with a number $K''$ of $L$; as $x_k^{1/2}\varepsilon_k=(A_0/A_1)\mathsf b_k$, this is $|\mathsf n(b)|\le C_\theta t({\rm an})\mathsf b_k$ with $C_\theta=K''C'A_0/(LA_1)$, a number of $L$, $\theta_3$, $A_1/A_0$ and $C_m$, and $C_\theta t({\rm an})\le\frac12\theta_3$ at $t({\rm an})<t_\theta$ once $c_{\rm ce}\le L/(2K''C')$.

The square sum for $1\le k<K$, over all the bonds of the step, follows in mean square, in three inequalities. The passage to the axial gauge: $|a''(c)|^2$ is at most a number of $L$ times $|a'(c)|^2$ plus the sum of $|a'(c')|^2$ over the bonds $c'$ inside the two blocks that $c$ meets, a bond $c'$ being counted at most $12L^3$ times. The averages in the mean square: the square sum of the block averages $a'$ over a set of bonds is at most a number of $L$ times the square sum of the displacement $A'$ over the bonds beneath them, each in the units of its level. The displacement is the difference $A'=A[{\rm an};b_0]-A[0;b_0]$ at the same complex row datum. Its energy is bounded by \eqref{eq:weighted-field-energy-difference}. At $b_0=0$ this is at most $C|{\rm an}|_\infty^2\xi^3\sum_p|F_U(p)|^2$; in bare units it is the corresponding bound by $C|{\rm an}|_\infty^2\sum_p|\tilde F_U(p)|^2$, where $\tilde F_U=\xi^2F_U$ and $|\tilde F_U(p)|^2\le2a_p(U)$.

At $b_0=0$, since $|\mathsf n|^2=x_k|a''|^2$ and, in a frame, $x_kL^k\le\frac{37}{11}\beta_t(y)$ at every point $y$, this gives $\sum_b|\mathsf n(b)|^2\le C^2|{\rm an}|_\infty^2A(\beta_t,U)$, $C$ a number of $L$; at $k=0$ the same holds at the external bonds of the step. The centres $\mathsf n^{(i)}$ of the earlier steps $0\le i<k$ are formed by the same rule about the moved background of the label, so that $g_i\mathsf n^{(i)}$, $g_i=x_i^{-1/2}$, is the block average at the level $i$ of its displacement, turned into the block axial gauge; by the bound of (c) on the displacement in the units of each layer, each obeys $|g_i\mathsf n^{(i)}(b)|\le Ct({\rm an})\varepsilon_i$, that is $|\mathsf n^{(i)}(b)|\le Ct({\rm an})(A_0/A_1)\mathsf b_i$, as $x_i^{1/2}\varepsilon_i=A_0(\log x_i)^{\mathfrak p_0}=(A_0/A_1)\mathsf b_i$, with one $C$ of $L$ for all the levels; this is at most $\frac12\theta_3\mathsf b_i$ at $t({\rm an})<t_\theta$ once $c_{\rm ce}\le1/(2C)$, and $c_{\rm ce}=\min\{L/(2K''C'),1/(2C)\}$ is a number formed with $L$ and $C'$.

The stored bonds, which are the external bonds of these steps, lie for the level $i$ within $2L^i$ (bare units) of the layer $i$; the square sums of the block averages of all the levels over the bonds near their own layers are bounded together, with one number of $L$, by the square sum of the displacement weighted by the scales of the layers, because the layers are wide and the bonds near the own layers of different levels overlap boundedly; and the $\ell^2$ bound of the linear system, as for the centre of the step, then gives $\sum_{i<k}\sum_b|\mathsf n^{(i)}(b)|^2\le C^2|{\rm an}|_\infty^2A(\beta_k,U)$, $b$ running over the stored bonds, with one $C$ of $L$, $A(\beta_k,U)$ being the action of the background with the couplings of the level $k$. (e) The room hypothesis of Lemma~\ref{C-360-23} holds at
every level: $L^m\ge MR(x_K)\ge LM$ gives $m>m'$, hence
$m'<m+K-j$ for $j\le K$. The identity is that lemma on a label, for any list of families with counterterms, with the explicit boundary terms of its clause on a label; the trivial values do not change because the tent families vanish at the trivial triple, and, at a fixed triple with the original families and counterterms held fixed, the linear changes in the re-filed local and boundary entries are those of the tent families.

The size bounds require the stated margin: by Lemma~\ref{C-360-23} with $\kappa'=\kappa+1$ and $|e_j|\le B_\sharp$, $\|\hat E^{(j)}-E^{(j)}\|_{j,\kappa+1}\le\frac32|{\rm an}|_\infty B_\sharp\alpha_0^2e^{7(\kappa+1)}\le\frac32\sigma_3\|{\rm an}\|B_\sharp\alpha_0^2e^{7(\kappa+1)}\le\theta E_0\|{\rm an}\|$, and the triangle inequality, with the bound $(1-\theta)E_0$ of the families before the re-filing, gives $\|\hat E^{(j)}\|_{j,\kappa+1}\le E_0$ and the distance from the family at zero source; the boundary terms are at most $\frac32M^3|{\rm an}|_\infty B_\sharp\alpha_0^2e^{7\kappa}e^{-\kappa d_j(X)}\le\theta M^3E_0e^{-7}\|{\rm an}\|e^{-\kappa d_j(X)}$ in a frame of Theorems~\YMMref{AJ}{673} and~\YMMref{AD}{607}, where the radius of the level is at most $\alpha_0$.
\end{proof}

\subsection{Scope of the conclusions}

Theorems~\ref{C-360-2}(d), \ref{C-360-3}(c), (d) and~\ref{C-360-4}(a), (c), (d) live on one torus, with constants that depend on its side; so does Theorem~\ref{C-360-4}(e), whose numbers are free of the side when the closed $a$-neighbourhoods of the supports lie in balls of radius $R$ and $\ell\ge16n(R+1)$. On $\mathbb R^3$ the statements concern limit points and tensor limits, at separated or disjoint supports. Theorem~\ref{C-360-4}(d) concerns cumulants with one stress entry among action densities on one torus, at a fixed label. The kernels of Theorem~\ref{C-360-2}(h) and the operators of Theorem~\ref{C-360-5} are those of each limit, taken along the subsequence that gives it. Theorems~\ref{C-360-13} and~\ref{C-360-14} concern one step, the never-large term, and the backgrounds, centres and counterterms of the construction.

Along growing tori and in the periodic Gibbs states the section asserts the following. The cumulants of the plane fields have limit points along every sequence of first-passage data at $g$ with $K_j\to\infty$, in periodic Gibbs states and on tori $T_{\ell_j}$ with $\ell_j\to\infty$, whatever the relation between $K_j$ and $\ell_j$, with the bounds, the invariance and the reflection positivity of Theorem~\ref{C-360-2}(c). The bulk term of the rotational Ward identity vanishes with a rate and with numbers that read no side, so that a limit along tori is invariant under rotations exactly when the flux through a shell about the supports tends to $0$; this flux tends to $0$ for the volume-first scalar limits; and, with the limits taken one after the other, at a label $\mathsf b\in I(g)$ the limit in $K$ of the flux on each torus $T_\ell$ is $\sum_iS^{(\ell,\mathsf b)}_n(f_1,\dots,(\omega y)\cdot\nabla f_i,\dots,f_n)$ through every enclosing shell, and along sides $\ell_k\to\infty$ in $\mathfrak L(g)$ on which the limits $S^\circ_n=\lim_kS^{(\ell_k,\mathsf b)}_n$ exist, the limit in $k$ of these limits, through any shell that covers the supports, a fixed one or a far one, is $-\frac d{dt}S^\circ_n(f(e^{-t\omega}\cdot))|_{t=0}$, and there are integers $K^*_k$ such that $S^\circ_n$ is the limit of the connected functions along $(K_k,\ell_k,L^{K_k}\mathsf b)$ for every sequence $K_k\ge K^*_k$ (Theorem~\ref{C-360-4}(b)).

For sides up to $a^{-\theta}$, $\theta<\frac1{30}$, a limit along tori at first-passage labels is the limit taken spacing first, then volume, and raising the cutoffs at fixed sides and labels does not change it (Theorem~\ref{C-360-4}(b)). With numbers that read no side, in every periodic Gibbs state the stress insertion collapses onto its classical tensor and the conserved tensor decorrelates from any number of tensor entries at separated supports (Theorem~\ref{C-360-3}(e), (f)); every tensor limit is conserved at disjoint supports (Theorem~\ref{C-360-3}(g)); and the translation identity of Theorem~\ref{C-360-4}(e), with its null term and its contact terms, holds per ball with numbers free of the side.

\section{The two-term short-distance law on the torus}\label{C-350-1}

For the action density, the first correction to the free two-point law is
of relative order $d/b$. More precisely, the two-term approximation is
\[
c_2b^{-2}\bigl[\mathcal K_6+c_1b^{-1}\mathcal K_5\bigr](f,h),\qquad c_2=\frac9{16\pi^2},\quad c_1=\frac8{3\pi},
\]
The error is
$[\bar\varepsilon(d)b^{-3}d+C_{\rm lat}b^{-2}a_K/d]\|f\|_\infty\|h\|_\infty$.
Theorem~\ref{C-350-3} states the cutoff and label ranges, the improvement
to a remainder of order $d^{3/2}$, and the averaged logarithmic correction.
Corollary~\ref{C-350-28} passes the law to every periodic state and to the
continuum limits and their components.

The proof first expresses the covariance as an average of source-energy
jets over translated pairs. A complex field scale separates the free
term, of order four, from the correction, of order six. A composed
block-field chart identifies the latter with a constrained lattice
Laplace coefficient. Comparison of two bare tori and removal of the flat
block constraint then connect it to the infinite-lattice coefficient.
The lattice calculation is proved in Section~\ref{IV-73-0}.

That calculation first takes the weak-coupling limit in a fixed box and
then exhausts the lattice. Its box-dependent remainder is not a bound in
the simultaneous continuum and large-volume limit. The renormalisation
and finite-torus comparisons below provide the separate passage needed
here.

\subsection{Parameters and test pairs}

We use the Wilson $G=\mathrm{SU}(2)$ model in three dimensions.
Its Lie algebra consists of the hermitian traceless $2\times2$ matrices
with $\langle X,Y\rangle=\frac12\operatorname{tr}XY$, extended
complex bilinearly. The trace is the ordinary matrix trace, and $|X|$
is the operator norm on $\mathbb C^2$. Near the identity we use
$\log X=-\sum_{n\ge1}(I-X)^n/n$, for $|X-I|<1$.
For $U\in\mathrm{SU}(2)\setminus\{-I\}$ the intrinsic branch is
$\log U=iY$, where $Y$ is Hermitian and traceless and $|Y|<\pi$;
these branches agree on their common domain.
Fix a mollifier, a source radius $2\varrho>0$, and one common choice
of the constants in Proposition~\YMMref{69.1}{595} and
Section~\YMMref{82}{665}. In particular $L\ge13$ is odd and $a_j=L^{-j}$.
We retain all the real-step conditions when imposing the sourced-step
conditions: restricting to the real construction leaves every retained
constant unchanged. The two radius conditions are
\[
 2\varrho\le\sigma^{\rm S}(x-\bar B),\qquad
 2\varrho x^{-1/2}(\log x)^{\mathfrak p_0}
       \le\sigma^{\rm S}_{\rm L}\qquad(x\ge\gamma^{-2}).
\]
They are imposed by Proposition~\YMMref{82.2}{668}, after the
radius-independent choices of Section~\YMMref{82}{665}. Thus all
uses of the real flow, the sourced progression and the sourced chain
below refer to the same numerical assignment and the same canonical
data.

The running couplings and the data at a reference coupling, the sides, the torus measures and the states, the weighted progression, the sourced chain on the least torus and its histories are therefore formed with this one choice, and every statement of this section is made at it; where Theorems~\YMMref{AJ}{673} and~\YMMref{AD}{607} enter, they are read at the radius $\varrho$ and at the choice of the longer list. Among the constants are $B_\sharp>0$, with $\bar B=B_\sharp L/(L-1)$, the powers $M$, $C_L$ of $L$ in the geometric setting of Section~\YMMref{52.7}{432}, $A_0,A_1>0$ and $p_0\in\mathbb N$. The real and source hypotheses are those imposed jointly in Section~\YMMref{82}{665}. A number is \emph{determined by the constants} if it depends on them alone.

For $\ell\in\mathfrak L=\{2L^m:m\in\mathbb N\}$ the torus $T_\ell=\mathbb R^3/\ell\mathbb Z^3$ carries the lattice of spacing $a_K$. The \emph{lattice} $(\ell,K)$, a side with a cutoff, carries the lattices $\mathbb T^{(j)}=a_{K-j}\mathbb Z^3/\ell\mathbb Z^3$ of the levels $0\le j\le K$, of spacing $D_j=a_{K-j}$; the level $0$ is the bare lattice, which in bare units is $\Lambda_N=\mathbb Z^3/N\mathbb Z^3$, $N=\ell L^K$. The \emph{height} $i$ is the averaging step from the level $i$ to the level $i+1$. Blocks, the bonds of the coarser lattices and the non-pivot bonds are those of Section~\YMMref{4}{16} and the pivot definition preceding Lemma~\YMMref{40.1}{297}, contours those of Section~\YMMref{4}{16} (coordinate-order contours), and $\mathcal M$ is the block-averaging map (equations~\YMMref{(37)--(39)}{20}).

For a plaquette $p$, $a_p(U)=1-\frac12\operatorname{tr}\partial U(p)$ and $x(p)$ is its barycentre; $A=\sum_pa_p$, $\mu_{K,\ell,\beta}(dU)=Z^{-1}e^{-\beta A(U)}dU$, $\mathcal O(f)=a_K^{-1}\sum_pf(x(p))a_p$, for a complex function $w$ on the bare plaquettes, a \emph{source}, $\mathcal O(w)=a_K^{-1}\sum_pw(p)a_p$, so that $\mathcal O(f)=\mathcal O(f\circ x)$, and $S^T_{2;K,\ell,\beta}(f,h)$ is the covariance of $\mathcal O(f)$ and $\mathcal O(h)$ under $\mu_{K,\ell,\beta}$. The lattice isometries preserve $\mu_{K,\ell,\beta}$ and carry $\mathcal O(f)$ to $\mathcal O$ of the carried function (Lemma~\ref{lem:setup-lattice-covariance}: plaquette holonomies are conjugated or inverted, and the positive-bond variables are permuted or inverted under product Haar measure); so $S^T_{2;K,\ell,\beta}(\tau_vf,\tau_vh)=S^T_{2;K,\ell,\beta}(f,h)$ for every $v\in a_K\mathbb Z^3$ and all real Lipschitz $f,h$, whatever their supports, $\tau_v$ being the translation by $v$.

For a reference coupling $0<g\le\gamma$ ($\gamma$ the last constant) put $X_g=L(g^{-2}+B_\sharp)$, $b_+=L(g^{-2}+\bar B)$ and $I(g)=(g^{-2}+\bar B,X_g-\bar B]$. The sides at $g$ form $\mathfrak L(g)=\{2L^m:m\ge m_{\mathfrak L}(g)\}$ (equation~\YMMref{(1528)}{615}), with least member $\ell_g$; so $\ell_gL^s\in\mathfrak L(g)$ for every integer $s\ge0$. A \emph{datum} at $g$ is a first-passage datum $(K,\beta)$ at $g$ (Corollary~\YMMref{74.2}{616}): its running couplings $x_0=\beta>x_1>\dots>x_K$ obey $x_{j+1}=L^{-1}x_j-e_{j+1}$ with $|e_{j+1}|\le B_\sharp$, and $K$ is the least index with $x_K\le X_g$; its label is $b=a_K\beta$. The running couplings depend on $(K,\beta)$ alone.

Every label at $g$ lies in $[g^{-2},b_+]$ \YMMref{(1531)}{615}, and for $b\in I(g)$ the pair $(K,L^Kb)$ is a datum at $g$ for every $K\in\mathbb N$ (Corollary~\YMMref{74.2}{616}). The \emph{witness depth} is $s(d)=\min\{s\ge0:L^sd\ge C_W(s)\}$, $C_W(s)=C_{W,0}A_0^2(\log[L^{s+1}(g^{-2}+\bar B)])^{2p_0}$, $C_{W,0}=\max\{1,4L^2M/A_0^2\}$, as in equation~\YMMref{(2321)}{843} and Theorem~\YMMref{AV}{851}; it depends on $g$, $d$ and the constants, and on no cutoff, label or side. As $C_W(s)\ge4L^2M$, for $d\le L^{-1}$ and $K>s(d)$ one has $L^{s(d)}\ge4L^3M$, $s(d)\ge2$, $K\ge3$ and $a_K\le L^{-1-s(d)}<d/13$; moreover $L^{s(d)}\ge1/d$, so $C_W(s(d))\ge(\log\frac1d)^2$ and $d/a_K\ge LC_W(s(d))\ge L(\log\frac1d)^2$.

Finally $\gamma^{\rm sd}_W\le\gamma^{\rm ref}_W\le\gamma$ are the coupling thresholds of the choices in equations~\YMMref{(2326)--(2327)}{845}; $(\gamma^{\rm sd}_W)^{-2}\ge4\bar B$.

For $\vartheta\in(0,1)$, $d>0$, and $\mathbb E=\mathbb R^3$ or $\mathbb E=T_\ell$ with $20d\le\ell$, $\mathfrak T_{\mathbb E}(d;\vartheta)$ is the following class of test pairs: Lipschitz functions $f,h\ge0$ on $\mathbb E$, neither $0$, with supports in closed balls of radius $d$ at a distance in $[d,6d]$, $\int f\ge\vartheta\|f\|_\infty d^3$, $\operatorname{Lip}f\le\|f\|_\infty/(\vartheta d)$, and the same for $h$. Then both supports lie in one ball of radius $5d$, and $d\le|x-y|\le10d$ for $x\in\operatorname{supp}f$, $y\in\operatorname{supp}h$. The conditions involve distances alone. So the class is invariant under translations; and for $36d<\ell$ the periodisation of a pair of $\mathfrak T_{\mathbb R^3}(d;\vartheta)$ is a pair of $\mathfrak T_{T_\ell}(d;\vartheta)$ with the same sup norms and the same kernels $\mathcal K_q$ below, every pair of $\mathfrak T_{T_\ell}(d;\vartheta)$ being the periodisation of one such pair, its \emph{lift}, since all the distances that occur are at most $10d<\ell/2$.

For two sides $\ell\le\breve\ell$, the \emph{pair with the same lift} as a pair $(f,h)$ on $T_\ell$ is the periodisation on $T_{\breve\ell}$ of the lift of $(f,h)$.

We put $\eta=\|f\|_\infty\|h\|_\infty$, $\mathcal K_q(f,h)=\iint f(x)h(y)|x-y|^{-q}\,dx\,dy$, $\varkappa=c_1c_2=\frac3{2\pi^3}$ and $L_b^+=c_2b^{-2}[\mathcal K_6+c_1b^{-1}\mathcal K_5]$. On $\mathbb Z^3$, $\mathsf B_1^{\rm lat}(p,q)$ is the coefficient of $\beta^{-3}$ in the covariance of $a_p$ and $a_q$ (Definition~\textup{\ref{IV-73-22}}), and $\mathsf P_1^{\rm lat}(f,h;a)=a\sum_{p,q}f(ax(p))h(ax(q))\mathsf B^{\rm lat}_1(p,q)$. There are a pure number $C_{\rm T}$ and a constant $\kappa^{\rm lat}$ with
\[
|\mathsf P_1^{\rm lat}(f,h;a)-\kappa^{\rm lat}\mathcal K_5(f,h)|\le C_{\rm T}\vartheta^{-1}a\lg(d/a)\,\eta\qquad\bigl((f,h)\in\mathfrak T_{\mathbb R^3}(d;\vartheta),\ 0<a\le d\bigr),
\]
$\lg t=1+|\log t|$ (Corollary~\textup{\ref{IV-73-73}}), and $\kappa^{\rm lat}=\frac3{2\pi^3}=c_1c_2$ (Theorems~\textup{\ref{IV-73-98}}, \textup{\ref{IV-73-100}}). For real Lipschitz $f,h$ on $T_\ell$, $\ell\in\mathfrak L(g)$, with disjoint supports and $b\in\overline{I(g)}$, $S_2^{(\ell,b)}(f,h)=\lim_{K\to\infty}S^T_{2;K,\ell,L^Kb}(f,h)$ is the limit of Theorem~\YMMref{A}{5}. That theorem gives existence for every $b>0$, continuity in $b$ on $(0,\infty)$, and bilinearity on real Lipschitz functions supported in two fixed disjoint closed sets. A function on $\mathbb R^3$ with support in a ball of radius $<\ell/4$ is read on $T_\ell$ by periodisation, and $f_\rho:=f(\cdot/\rho)$.

\begin{lemma}[Kernel bounds for test pairs]\label{lem:test-pair-kernels}
For $(f,h)\in\mathfrak T_{\mathbb E}(d;\vartheta)$ and $q>0$,
\[
 10^{-q}\vartheta^2d^{6-q}\|f\|_\infty\|h\|_\infty
 \le\mathcal K_q(f,h)
 \le(4\pi/3)^2d^{6-q}\|f\|_\infty\|h\|_\infty.
\]
Moreover $d\le\mathcal K_5(f,h)/\mathcal K_6(f,h)\le10d$.
\end{lemma}
\begin{proof}
All cross-distances of the two supports lie between $d$ and
$10d$. Positivity gives
\[
 (10d)^{-q}\int f\int h
 \le\mathcal K_q(f,h)\le d^{-q}\int f\int h.
\]
The defining mass bounds give the lower estimate, and the
volumes of the two radius-$d$ balls give the upper estimate.
On a torus the balls have the same volumes because $20d\le\ell$.
Finally $|x-y|^{-5}=|x-y|\,|x-y|^{-6}$ on the product of the
supports; integration gives the ratio bound. Its denominator
is positive by the lower estimate.
\end{proof}

\begin{definition}\label{C-350-2}
Let $N_3$ be the threefold convolution of $\mathbf 1_{[-1/2,1/2)}$, $\Phi(x)=\prod_{\mu=1}^3N_3(x_\mu)$, $\zeta=1$ on $[0,7]$, $\zeta(t)=1-3(t-7)^2+2(t-7)^3$ on $[7,8]$, $\zeta=0$ on $[8,\infty)$, $P(y)=|y|^2-\frac34$ and $\Xi(y)=P(y)\zeta(|y|)$. For odd $q\ge1$ let $J_q=\{-\frac{q-1}2,\dots,\frac{q-1}2\}$, $c^q_j=\prod_{\mu}q^{-2}\#\{\iota\in J_q^3:\iota_1+\iota_2+\iota_3=j_\mu\}$ for $j\in\mathbb Z^3$, and
\[
\Upsilon^{(q)}(y)=q^{-1}\sum_{j\in\mathbb Z^3}c^q_j\,\Xi\bigl((y+j)/q\bigr)-\Xi(y).
\]
With $e(u)=\int\Phi(x)\Xi(x+u)\,dx$ let $\mathsf k_6=\int_0^\infty\bigl(4\pi-t^{-2}\int_{S^2}e(t\omega)\,d\omega\bigr)t^{-2}\,dt$.
\end{definition}

\subsection{The quantitative two-term law}

\begin{theorem}\label{C-350-3}
There is $\gamma_{\rm e}\in(0,\frac12]$, determined by the constants, with the following property. Let $0<g\le\gamma_{\rm e}$, $\vartheta\in(0,1)$, $\ell=\ell_g$, $P\ge P_\sharp:=3601$, and put $C_{\rm lat}=\frac32C_{\rm tm}\vartheta^{-1}$ and $C^\flat_{\rm lat}=\frac92C_{\rm tm}\vartheta^{-1}+1$, where $C_{\rm tm}$, determined by the constants, is the constant of the template sums of test pairs (Proposition~\YMMref{121.2}{824} and Lemma~\ref{lem:ward-fluctuation-template}). There are $d_{\rm e}>0$, $c_{1/2}<\infty$, $c_{\rm r}<\infty$ and a non-decreasing $\bar\varepsilon:(0,d_{\rm e}]\to[0,\infty)$ with $\bar\varepsilon(d)\to0$ as $d\to0$, depending on $g$, $\vartheta$, $P$ and the constants, and on no cutoff, label, pair or position, such that for $0<d\le d_{\rm e}$ and $(f,h)\in\mathfrak T_{T_\ell}(d;\vartheta)$ the following hold.
\begin{enumerate}
\item[(a)] At every datum $(K,\beta)$ at $g$ with $s(d)<K$ and $L^K\le d^{-P}$, with $b=a_K\beta$,
\begin{equation}\label{C-350-4}
|S^T_{2;K,\ell,\beta}-L_b^+|(f,h)\le\bigl[\bar\varepsilon(d)\,b^{-3}d+C_{\rm lat}\,b^{-2}a_K/d\bigr]\eta .
\end{equation}
\item[(b)] For every $b\in I(g)$, \eqref{C-350-4} holds at $(K,L^Kb)$ for every $K>s(d)$, and $|S_2^{(\ell,b)}-L_b^+|(f,h)\le\min\{\bar\varepsilon(d),c_{1/2}d^{1/2}\}\,b^{-3}d\,\eta$.
\item[(c)] If moreover $g\le L^{-1/4}\min\{\gamma_{\rm e},\gamma^{\rm sd}_W\}$, there are $d'_{\rm e}>0$ and $\bar\varepsilon'$ with $\bar\varepsilon'(d)\to0$, depending on $g$, $\vartheta$, $P$, the radius $\varrho$ and the choice of the longer list of which the constants are the restriction, and on no cutoff, label, pair or position, such that \eqref{C-350-4}, with $\bar\varepsilon'$ in the place of $\bar\varepsilon$, holds for $0<d\le d'_{\rm e}$, $(f,h)\in\mathfrak T_{T_\ell}(d;\vartheta)$ and every datum $(K,\beta)$ at $g$ with $K>s(d)$.
\item[(d)] Let $k_0\in\mathbb Z$ with $44L^{-k_0}<\ell_g$, $r_0=L^{-k_0}$, and for $n\ge1$ let $r_1=L^{-n}r_0$ and
\[
\Delta_n(b)=r_1^{-1}S_2^{(\ell,b)}\bigl(\Phi_{r_1},\Upsilon^{(L^n)}_{r_1}\bigr)-c_2b^{-2}\mathsf k_6\,(r_1^{-1}-r_0^{-1})-\frac6{\pi^2}\,b^{-3}\log\frac{r_0}{r_1}
\]
(the pair has disjoint supports, so that the limit is defined at it). Then $\Delta_n$ converges uniformly on $I(g)$ to a bounded continuous function $\Delta_\infty$, and there are $C_{\rm co}<\infty$ and $n_{\rm co}\in\mathbb N$, depending on $g$, $k_0$ and the constants, with $\sup_{I(g)}|\Delta_n|\le C_{\rm co}$ for all $n\ge1$ and $\sup_{I(g)}|\Delta_n-\Delta_\infty|\le C_{\rm co}L^{-n/2}$ for $n\ge n_{\rm co}$.
\item[(e)] For every $b\in I(g)$ and every $K>s(d)$,
\begin{equation}\label{C-350-30}
|S^T_{2;K,\ell,L^Kb}-L_b^+|(f,h)\le\bigl[c_{\rm r}\,b^{-3}d^{3/2}+C^\flat_{\rm lat}\,b^{-2}a_K/d\bigr]\eta ;
\end{equation}
and \eqref{C-350-30}, with $\beta$ for $L^Kb$, holds at every datum $(K,\beta)$ at $g$ with $s(d)<K\le K_\sharp(d)$ and $b=a_K\beta$, where $K_\sharp(d)$ is the cutoff of Proposition~\ref{C-350-16} at the side $\ell_g$.
\item[(f)] Under the assumptions of (c) there are $d''_{\rm e}>0$ and $c'_{\rm r}<\infty$, depending on $g$, $\vartheta$, the radius $\varrho$ and the choice of the longer list of which the constants are the restriction, and on no $P$, which enters them as $P_\sharp$, and no cutoff, label, pair or position, such that \eqref{C-350-30}, with $c'_{\rm r}$ for $c_{\rm r}$ and $\beta$ for $L^Kb$, holds for $0<d\le d''_{\rm e}$, $(f,h)\in\mathfrak T_{T_\ell}(d;\vartheta)$ and every datum $(K,\beta)$ at $g$ with $K>s(d)$, $b=a_K\beta$.
\end{enumerate}
\end{theorem}

As $|x-y|\le10d$ on the supports, $\mathcal K_5(f,h)\ge(10d)^{-5}\int f\int h\ge10^{-5}\vartheta^2d\,\eta$ by the nonnegativity and integral floors in Lemma~\ref{lem:ward-test-pairs}. So clause (b) determines the number:
\[
\Bigl|\frac{(S_2^{(\ell,b)}-c_2b^{-2}\mathcal K_6)(f,h)}{c_2b^{-3}\mathcal K_5(f,h)}-\frac8{3\pi}\Bigr|\le10^5c_2^{-1}\vartheta^{-2}\min\{\bar\varepsilon(d),c_{1/2}d^{1/2}\}.
\]
The continuum limit on $T_{\ell_g}$ at a label of $I(g)$ determines first the label, since 
\[
S_2^{(\ell,b)}(f,h)/(c_2\mathcal K_6(f,h))\to b^{-2}
\]
 along pairs of the class with $d\to0$, and then the number $\frac8{3\pi}$. The number enters the proof at one place, the last link of the chain of Step~3 below, through $\kappa^{\rm lat}=c_1c_2$; it depends on no $L$, $g$, side or label.

\begin{corollary}[Every state; every side; the limits on $\mathbb R^3$]\label{C-350-28}
Let $g$, $\vartheta$ and $P$ be as in clause (c) of Theorem~\ref{C-350-3}, with $d'_{\rm e},\bar\varepsilon'$ of (c) and $d''_{\rm e},c'_{\rm r}$ of (f); let $C^{\rm T}_2(g)$ and $\mathsf d_\star(g,\vartheta)$ be the numbers of Proposition~\ref{C-350-21}, and put
\[
d_\infty=\min\bigl\{d'_{\rm e},d''_{\rm e},\tfrac1{180},\mathsf d_\star(g,\vartheta),L^{-1}\bigr\},\qquad\bar\varepsilon_\infty(d)=\min\{\bar\varepsilon'(d),c'_{\rm r}d^{1/2}\}+C^{\rm T}_2(g)\vartheta^{-1}b_+d .
\]
The numbers $d_\infty$ and $\bar\varepsilon_\infty(d)$ depend on $g$, $\vartheta$, $P$, the radius $\varrho$ and the choice of the longer list of which the constants are the restriction, and on no cutoff, label, side, state, sequence, component, pair or position. Let $0<d\le d_\infty$ and $(\tilde f,\tilde h)\in\mathfrak T_{\mathbb R^3}(d;\vartheta)$, with $\eta=\|\tilde f\|_\infty\|\tilde h\|_\infty$.
\begin{enumerate}
\item[(i)] For every state $\mathfrak s$ at $g$ at a datum $(K,\beta)$ at $g$ with $K>s(d)$ (Proposition~\ref{C-350-21}), $b=a_K\beta$:
\[
|S^T_2[\mathfrak s]-L_b^+|(\tilde f,\tilde h)\le\bigl[(\bar\varepsilon'(d)+C^{\rm T}_2(g)\vartheta^{-1}b_+d)\,b^{-3}d+C_{\rm lat}b^{-2}a_K/d\bigr]\eta,
\]
and the same with $c'_{\rm r}d^{1/2}$ in the place of $\bar\varepsilon'(d)$ and $C^\flat_{\rm lat}$ in the place of $C_{\rm lat}$. In particular this holds for the torus measure of every side of $\mathfrak L(g)$ at the periodisations of the pair, and for every periodic Gibbs state at $(a_K,\beta)$.
\item[(ii)] For every side $\ell'\in\mathfrak L(g)$, every $b\in I(g)$ and every $(f,h)\in\mathfrak T_{T_{\ell'}}(d;\vartheta)$: $|S_2^{(\ell',b)}-L_b^+|(f,h)\le\bar\varepsilon_\infty(d)\,b^{-3}d\,\eta$.
\item[(iii)] Let $\mathfrak s_j$ be states at $g$ at data $(K_j,\beta_j)$ at $g$ with $K_j\to\infty$ and labels $b_j=a_{K_j}\beta_j\to b$; the torus measures of sides $\ell_j\in\mathfrak L(g)$ that grow with $j$ are among them. Then $b\in[g^{-2},b_+]$, and every limit point $S$ of the sequence $S^T_2[\mathfrak s_j](\tilde f,\tilde h)$ satisfies $|S-L_b^+(\tilde f,\tilde h)|\le\bar\varepsilon_\infty(d)\,b^{-3}d\,\eta$. The labels are determined by the two-point functions: if $\mathfrak s_j$ are states at $g$ at data $(K_j,\beta_j)$ at $g$ with $K_j\to\infty$, and $S^T_2[\mathfrak s_j](F)$ converges at every pair $F$ of a set that holds, for arbitrarily small $d\in(0,d_\infty]$, a pair of $\mathfrak T_{\mathbb R^3}(d;\vartheta)$, then the labels $a_{K_j}\beta_j$ converge, to a number $b\in[g^{-2},b_+]$.

Consequently the labels of every convergent Gibbs sequence at $g$ (a sequence $(K_j,\beta_j,\nu_j)$ of data at $g$ with $K_j\to\infty$ and periodic Gibbs states $\nu_j$ at $(a_{K_j},\beta_j)$) converge, to a number $b$; and the connected two-point function $S^\sigma_2$ of the action density of its limit $\sigma$, and that of every component of a canonical decomposition of $\sigma$, satisfy $|S_2^\sigma-L_b^+|(\tilde f,\tilde h)\le\bar\varepsilon_\infty(d)\,b^{-3}d\,\eta$ whenever $\tilde f$ and $\tilde h$ are of class $C^1$. Such a pair is a separated tuple, at which these functionals are defined, and $\mathfrak T_{\mathbb R^3}(d;\vartheta)$ holds pairs of class $C^1$ for every $d>0$ and $\vartheta\in(0,1)$.

Every volume-first limit at $g$ has one label $b$, the same along all its defining sequences, and the same inequality holds for it and for its ergodic components (Section~\ref{sec:state-limit-conventions} and Definition~\ref{C-365-5}(2)) at the pairs of $\mathfrak T_{\mathbb R^3}(d;\vartheta)$ of class $C^1$. Every $b\in I(g)$ is the label of such limits.
\item[(iv)] In the cases (ii) and (iii), with $S$ the two-point function in question at a pair at which (ii) or (iii) states the inequality, $|(S-c_2b^{-2}\mathcal K_6)/(c_2b^{-3}\mathcal K_5)-\frac8{3\pi}|\le10^5c_2^{-1}\vartheta^{-2}\bar\varepsilon_\infty(d)$: one number $\frac8{3\pi}$ holds for all the limits on the tori of the sides at $g$, the limits of states at $g$ and their components, and each of these limits determines its label and the number. In the case (i), at a state at a finite cutoff, the same quotient differs from $\frac8{3\pi}$ by at most $10^5c_2^{-1}\vartheta^{-2}\bigl[\bar\varepsilon'(d)+C^{\rm T}_2(g)\vartheta^{-1}b_+d+C_{\rm lat}\,b\,a_K/d^2\bigr]$, and by at most the same with $c'_{\rm r}d^{1/2}$ in the place of $\bar\varepsilon'(d)$ and $C^\flat_{\rm lat}$ in the place of $C_{\rm lat}$.
\end{enumerate}
\end{corollary}

\subsection{Renormalisation and coefficient inputs}

Theorem~\YMMref{D}{8} stands behind all the statements on the renormalisation group that we use, and Theorems~\YMMref{AJ}{673} and~\YMMref{AD}{607} behind the existence of the limit on a torus, the rate of the cutoffs and the agreement of states, into which it enters through the bounds on the partition functions of the tori at the source $0$. We use the following construction results: the flow of the coupling and its data (Section~\YMMref{74}{614}, \YMMref{74}{614}, Section~\ref{sec:complex-field-scale}, Proposition~\ref{prop:source-entry}, equation~\YMMref{(1528)}{615}, \YMMref{74.2}{616}, Section~\YMMref{82}{665}, Proposition~\YMMref{74.1}{615}, Corollary~\YMMref{74.2}{616}, the ordered real choice in Proposition~\YMMref{69.1}{595} and the common source choice of Section~\YMMref{82}{665}, the coupling-sequence quantities in equation~\YMMref{(634)}{339} and the function $R$ of equation~\YMMref{(634)}{339}); the averaging chart and gauge factor of Section~\ref{sec:sf-chart-interface}, with the measure identity of Lemma~\ref{lem:sf-chart-integral} and the cube-form sum in equation~\YMMref{(828)}{395}, and the free response on tori (the construction and periodization in Section~\YMMref{52.6}{426}, and the curvature bound~\YMMref{(2234)}{820}); the localisation domains of a torus and of the infinite lattice, their near lifts and their linear sizes (Section~\YMMref{72}{607} (the definitions of $\mathcal D_j$ and $d_j$ at its beginning), Lemma~\ref{lem:analytic-domain-geometry}, Lemma~\ref{lem:near-domain-geometry}); the sizes \YMMref{(759)}{375} of families (Definition~\YMMref{49.5}{375}) and the two sizes $|v|_\infty$ and $\operatorname{sw}_k(v)$ of a plaquette weight in Section~\ref{sec:source-entry-calculus}; the bracket localisation of a bare array and the octet property (Propositions~\YMMref{51.14}{408} and~\YMMref{51.16}{412}, with equation~\YMMref{(877)}{412}); the one-step theorem (Theorem~\YMMref{V}{416}); the naturality of one step under the charts between two lattices with one cutoff (Section~\YMMref{52.5}{425}, in particular equation~\YMMref{(928)}{426}); the regular progression on all the lattices of a level (Theorem~\YMMref{Y}{438}); the weighted progression, its steps at a complex field scale and its chain classes (Definition~\ref{def:weighted-progression-data} and equation~\eqref{eq:weighted-progression-recursion}, \ref{def:weighted-history-setting}, the definitions preceding \eqref{eq:weighted-history-threshold}, the definitions preceding equation~\YMMref{(2276)}{831}, Section~\YMMref{124}{832}, \eqref{eq:weighted-chain-class}, \ref{IV-71-26}, the bracket-localization constants in (Wp4), \eqref{eq:weighted-progression-numbers}, \eqref{eq:weighted-progression-numbers}, \eqref{eq:weighted-running-numbers}, \eqref{eq:weighted-mixed-numbers}, \eqref{eq:weighted-history-threshold}, equation~\YMMref{(2286)}{833}, Proposition~\textup{\YMMref{122.1}{826}}, equation~\YMMref{(2258)}{827}, Lemma~\ref{lem:source-support-sum}, equations~\YMMref{(2286)--(2287)}{833}, \YMMref{124.4}{835}, \YMMref{124.3}{835}, equation~\YMMref{(2291)}{834}, Propositions~\YMMref{122.2}{827}, \YMMref{123.1}{829}, \ref{prop:weighted-canonical-history}, \YMMref{124.2}{834}, \YMMref{124.5}{836}, \ref{prop:weighted-pair-tower}, \ref{IV-71-37}, Corollary~\ref{cor:weighted-canonical-data}, Lemma~\YMMref{126.2}{845}), with the restricted Gaussian measures of one step (Lemma~\ref{lem:ward-box-cumulants}); the level-wise sourced density and packed-source amplitudes
(Proposition~\ref{prop:levelwise-source-density} and
Lemma~\ref{lem:packed-levelwise-source}), the first-creation
decomposition and its coefficient bounds
(Proposition~\ref{prop:first-creation-envelopes},
Corollary~\ref{cor:packed-source-coefficients} and
Proposition~\ref{prop:ward-fixed-volume-rarity}), and the
least-torus comparison proved in Proposition~\ref{C-350-21}; the free template and the template sums (Section~\YMMref{120}{817}, Section~\YMMref{120}{817}, \YMMref{(2231)}{819}, Lemma~\YMMref{120.1}{818}, the test-pair bounds~\eqref{C-351-2}, Proposition~\YMMref{120.3}{819}, \YMMref{120.5}{821} and~\YMMref{120.6}{821}, \YMMref{121.2}{824} and Lemma~\ref{lem:ward-fluctuation-template}, Proposition~\YMMref{120.3}{819}); Laplace's method and the lattice coefficient (Definitions~\textup{\ref{IV-73-15}}, \textup{\ref{IV-73-22}}, Lemmas~\textup{\ref{IV-73-12}}, \textup{\ref{IV-73-21}}, \textup{\ref{IV-73-26}}, \textup{\ref{IV-73-51}}, \textup{\ref{IV-73-53}}, \textup{\ref{IV-73-55}}, \textup{\ref{IV-73-75}}, \textup{(\ref{IV-73-24})}, \textup{(\ref{IV-73-25})}, \textup{(\ref{IV-73-56})}, Corollary~\textup{\ref{IV-73-73}}, Theorems~\textup{\ref{IV-73-98}}, \textup{\ref{IV-73-100}}); the rate of the cutoffs on a torus (the constants and thresholds in Theorem~\YMMref{AR}{810}, \YMMref{(2206)}{811}, the two threshold inequalities at the end of the proof of Theorem~\YMMref{AR}{810}, Corollary~\YMMref{117.1}{812}); the limit at every label on a torus (Theorem~\textup{\YMMref{AS}{813}}, Proposition~\YMMref{118.3}{814}); the separated tuples and the test pairs of class $C^1$ (Section~\ref{sec:state-limit-conventions}, the convention in equation~\YMMref{(2343)}{849}, Lemma~\ref{lem:ward-test-pairs} and equation~\YMMref{(2247)}{824}); the sequences of pairs of laws, their convergence and the connected functions of their limit functionals (Proposition~\YMMref{92.1}{710} and its moment--cumulant identity~\YMMref{(1856)}{711}); the Gibbs sequences, their limits and the components of these (Definition~\ref{C-365-5}(2), Theorem~\textup{\YMMref{AP}{780}}); and the common construction choice of Section~\YMMref{82}{665}, with the torus limits of Theorem~\YMMref{A}{5} and the growing-torus and volume-first limits of Proposition~\YMMref{92.1}{710}.

\subsection{Source-energy jets and the field-scale expansion}

\emph{Notation.} Throughout, $D=(K,\beta)$ is a datum at $g$ with label $b$, $\ell=2L^m\in\mathfrak L(g)$, and $f,h$ are real bounded functions on $T_\ell$; $w_1=f\circ x$ and $w_2=h\circ x$ are the lattice sources on the bare plaquettes. We write $N=\ell L^K$, $\mathsf N=N+\beta$, $r=d/a_K$, $l=\max\{1,\log(2b/d)\}$, $\mathfrak b(x)=A_1(\log x)^{p_0}$, $\delta(x)=\mathfrak b(x)x^{-1/2}$ and $\langle z\rangle=\max\{1,|z|\}$. With $E^{(j)}_\mu$ the local terms of the level $j$ of the weighted progression (\eqref{eq:weighted-progression-recursion}) at the datum, with the weight $\mu_1w_1+\mu_2w_2$ of the plaquettes, and $\mathcal D_j$ its domains of the level $j$, the second jet of the level $j$ is
\[
\mathsf n_j=\sum_{X\in\mathcal D_j}\ \sum_{z\in X\cap\mathbb T^{(j)}}\partial_{\mu_1}\partial_{\mu_2}E^{(j)}_\mu(X,\mathbf 1,z)\big|_{\mu=0},\qquad\mathsf J_2=\sum_{j=1}^K\mathsf n_j .
\]
The same construction at a complex field scale $\lambda$ (every fluctuation field multiplied by $\lambda$, the coefficient of the action at the height $j$ moved to $\xi_j(\lambda)=L^{-j}x_0+\lambda^2(x_j-L^{-j}x_0)$) gives the local terms $E^{(j),\lambda}_\mu$ of the level $j$ at the field scale $\lambda$, with $E^{(j),1}_\mu=E^{(j)}_\mu$, and their second jets $\mathsf n^\lambda_j=\sum_{X\in\mathcal D_j}\sum_{z\in X\cap\mathbb T^{(j)}}\partial_{\mu_1}\partial_{\mu_2}E^{(j),\lambda}_\mu(X,\mathbf 1,z)|_{\mu=0}=\sum_{m\ge0}\mathsf n_{j,[2m]}\lambda^{2m}$ (Proposition~\ref{C-350-7}); $\mathsf h^{(n)}_{2m}=\sum_{j\le n}\mathsf n_{j,[2m]}$, $\mathsf h_{2m}=\mathsf h^{(K)}_{2m}$, and $\mathsf H^{(n)}[D]=\beta^3a_K^2\mathsf h^{(n)}_6$. As $\beta^{-3}a_K^{-2}=b^{-3}a_K$ and $a_Kr=d$, a size $r\eta$ of $\mathsf H$ is a size $b^{-3}d\eta$ of $\mathsf h_6$. Where the side matters we write $\mathsf J_2[\ell]$, $\mathsf h_6[\ell]$, $\mathsf H^{(n)}[D,\ell]$.

\emph{The free template} (Section~\YMMref{120}{817}) is the free abelian model of the curvature: cochains, plaquette averages and nested projections. At the height $0\le i<K$, $\mathsf A^{\rm cv}_i=\mathsf B^{\rm cv}_i+\hat{\mathsf G}$ is its one-step quadratic form, a sum of two real symmetric matrices $\ge0$ on the non-pivot bonds of $\mathbb T^{(i)}$, and $\mathsf A^{\rm cv}_i\ge\gamma^{\rm tr}$ with a lower bound $\gamma^{\rm tr}>0$ (Proposition~\YMMref{120.3}{819}); $\mathsf T_i$ is its curvature map onto the bare plaquettes (Proposition~\YMMref{120.3}{819}). $\Pi^{\rm tp}_0$ is the orthogonal projection onto the exact bare curvatures and $\Pi^{\rm tp}_n$, $n\ge1$, the projection onto those that the plaquette averages of the level $n$ determine, so that $\Pi^{\rm tp}_0-\Pi^{\rm tp}_n$ is the orthogonal projection onto the curvatures whose plaquette averages of the level $n$ vanish; $\mathcal C_{(j)}=\Pi^{\rm tp}_{j-1}-\Pi^{\rm tp}_j$ are the curvature slices, and $C^{\rm fl}_{K,K-n}=\Pi^{\rm tp}_0-\Pi^{\rm tp}_n=\sum_{1\le j\le n}\mathcal C_{(j)}$ is the fluctuation kernel, the conditional covariance of the free lattice curvature given its plaquette averages of the level $n$; on a lattice with $t$ levels we write $\Pi^{\rm fl}_t=\Pi^{\rm tp}_0-\Pi^{\rm tp}_t$.

The numbers $\delta_{\rm tp}\in(0,\frac1{12}]$ and $C_{\rm tp}\ge1$ of the decay of the slices (Propositions~\YMMref{120.5}{821} and~\YMMref{120.6}{821}), and the number $C^{\rm row}_{\rm tp}$ of the row sums of the curvature maps,
\[
\sum_b|\mathsf T_i(p,b)|\le C^{\rm row}_{\rm tp}L^{-2i}\qquad\text{for every bare plaquette }p,
\]
uniformly in the volume, at the height $i=0$ and at every height at which the side $\ell L^{K-i}$ of the level $i$, in its units, is at least $N_{\mathsf H}$, so at every height if $\ell\ge N_{\mathsf H}$ (Proposition~\YMMref{120.6}{821}), depend on $L$ and on the least admissible side $N_{\mathsf H}$ fixed with $C_{\mathsf H}$ and $\delta_{\mathsf H}$ in Section~\YMMref{52.6}{426} and the beginning of Section~\YMMref{120.3}{820}; they, $\gamma^{\rm tr}$ and $N_{\mathsf H}$ are determined by the constants.

\emph{The Gaussian measures of the heights.} Put $\tau_i=L^{-i}x_0/x_i$ and $\mathsf M_i=\tau_i\mathsf B^{\rm cv}_i+\hat{\mathsf G}$. $\nu^\circ_i$ is the centred Gaussian measure with the precision $\mathsf M_i\otimes1$ on the real $\mathfrak g$-valued $1$-forms $B_i$ on the non-pivot bonds of $\mathbb T^{(i)}$, and $\nu^{\rm r}_i$ is $\nu^\circ_i$ conditioned on the box $\{|B_i(b)|<\mathfrak b(x_i)\ \forall b\}$, $|\cdot|$ being the Euclidean norm of $\mathfrak g\cong\mathbb R^3$: the boxes are products of Euclidean balls. As $\tau_i\in[\frac34,\frac43]$ (Proposition~\ref{C-350-6}) and both matrices are $\ge0$, $\mathsf M_i\ge\min\{\tau_i,1\}\mathsf A^{\rm cv}_i\ge\frac34\gamma^{\rm tr}$, so every variance in $\nu^\circ_i$ is at most $\frac4{3\gamma^{\rm tr}}$.

The measures $\nu^\circ_i$ are tied to the steps by Lemma~\ref{C-350-23}(iii): at $\lambda=0$, at the trivial configuration, the density of the step of the height $i$ is proportional to $\exp[-\frac12(B_i,(\mathsf M_i\otimes1)B_i)]$ on the box of the height, the height $0$ included, where $\tau_0=1$ and $\mathsf M_0=\mathsf A^{\rm cv}_0$. For polynomials $X_{a,m}$ ($a=1,2$, $m=0,1,2$), $W_1$, $W_2$ in these fields and a probability measure $\nu$, with the joint cumulants $\langle\cdot;\dots;\cdot\rangle$ taken in $\nu$,
\begin{align}
\gamma_2[\nu]={}&\langle X_{1,2};X_{2,0}\rangle+\langle X_{1,0};X_{2,2}\rangle+\langle X_{1,1};X_{2,1}\rangle-\langle X_{1,1};X_{2,0};W_1\rangle\nonumber\\
&-\langle X_{1,0};X_{2,1};W_1\rangle-\langle X_{1,0};X_{2,0};W_2\rangle+\tfrac12\langle X_{1,0};X_{2,0};W_1^2\rangle ;\label{C-350-5}
\end{align}
$\gamma^{\rm r}_2=\gamma_2[\bigotimes_{0\le i<n}\nu^{\rm r}_i]$ and $\gamma^\circ_2=\gamma_2[\bigotimes_{0\le i<n}\nu^\circ_i]$, with the polynomials of Proposition~\ref{C-350-11}. If $\nu$ is invariant under $B\mapsto-B$, the $X_{a,0}$ are even and $W_1$ is odd, then $\frac12\langle X_{1,0};X_{2,0};W_1^2\rangle=\frac12\langle X_{1,0};X_{2,0};W_1;W_1\rangle$: a cumulant with a product entry is $\langle X;Y;W^2\rangle=\langle X;Y;W;W\rangle+2\langle X;W\rangle\langle Y;W\rangle+2\langle W\rangle\langle X;Y;W\rangle$, and the last two terms vanish by parity.

\emph{Laplace's method} (Lemma~\textup{\ref{IV-73-26}}). Let $\phi,\psi,X,Y$ be real and smooth near $0\in\mathbb R^m$, $\phi(0)=0$, $\nabla\phi(0)=0$, with a positive definite Hessian $H$ at $0$, and let $X$, $Y$ vanish at $0$ with their gradients. Let $\phi_j,\psi_j,X'_j,Y'_j$ be their homogeneous Taylor polynomials of order $j$ at $0$, and $W_j=\phi_{j+2}+\psi_j$, $X_{1,j}=X'_{j+2}$, $X_{2,j}=Y'_{j+2}$. Then the covariance of $X$ and $Y$ under the density proportional to $e^{-\beta\phi-\psi}$ on a small window about $0$ is
\begin{equation}\label{C-350-31}
\beta^{-2}\gamma^{\rm L}_0+\beta^{-3}\gamma^{\rm L}_2+O(\beta^{-4})\qquad(\beta\to\infty),
\end{equation}
where $\gamma^{\rm L}_0=\langle X_{1,0};X_{2,0}\rangle$ and $\gamma^{\rm L}_2$ is \eqref{C-350-5}, both in the centred Gaussian measure of covariance $H^{-1}$. In the variable $\beta^{1/2}x$ the density is the Gaussian one times $e^{-\beta^{-1/2}W_1-\beta^{-1}W_2-\dots}$ and the observables are $\beta^{-1}(X_{a,0}+\beta^{-1/2}X_{a,1}+\beta^{-1}X_{a,2}+\dots)$; the coefficient of $\beta^{-3}$ of the covariance, expanded in joint cumulants, is \eqref{C-350-5}, written with the parity just noted, and the half-integer orders vanish by parity. We write $\gamma^{\rm L}_{m'}[\phi,\psi,X,Y]$, $m'=0,2$. The numbers $\mathsf g_{m'}[N,t](p,q)$ of a flat top are set in Proposition~\ref{C-350-12}; for functions $w_1,w_2$ on the plaquettes of $\Lambda_N$, $\mathsf T[N,t](w)=\sum_{p,q}w_1(p)w_2(q)\mathsf g_2[N,t](p,q)$, and for functions $\tilde w_1,\tilde w_2$ of finite support on the plaquettes of $\mathbb Z^3$, $\mathsf T^{\rm lat}(\tilde w)=\sum\tilde w_1(\tilde p)\tilde w_2(\tilde q)\mathsf B_1^{\rm lat}(\tilde p,\tilde q)$.

If $\tilde w_1(\tilde p)=\tilde f(a_Kx(\tilde p))$, $\tilde w_2(\tilde p)=\tilde h(a_Kx(\tilde p))$, then $a_K\mathsf T^{\rm lat}(\tilde w)=\mathsf P_1^{\rm lat}(\tilde f,\tilde h;a_K)$.

\begin{proposition}[The setting of the expansion]\label{C-350-6}
There is $\gamma_{\rm s}\in(0,\min\{\gamma^{\rm ref}_W,\frac12\}]$, determined by the constants, such that for $0<g\le\gamma_{\rm s}$, every datum $(K,\beta)$ at $g$ with $K\ge3$, every $\ell=2L^m\in\mathfrak L(g)$, and real bounded $f,h$ with supports in balls of radius $d$ at a distance $\ge d$, $8a_K\le d$, at every position of $(f,h)$:
\begin{enumerate}
\item[(i)] $2L^{m+K-j+1}\ge C_LM$ for $1\le j\le K$, in particular $\ell L\ge C_LM$; and $\ell\ge N_{\mathsf H}$;
\item[(ii)] $x_0>x_1>\dots>x_K>g^{-2}$, $x_{K-1}>Lg^{-2}\ge4L$, and for $0\le j<K$: $\frac34b/D_j\le x_j\le2b/D_j$ and $\frac34\le\tau_j\le\frac43$;
\item[(iii)] at every level $j<K$ the constants, the lattice of the step $j\to j+1$ and $x_j$ are data under (Im) of Definition~\ref{def:sf-chart-data};
\item[(iv)] the data of the weighted progression (Definition~\ref{def:weighted-progression-data}) are given with $k=K-1$, the sources $w_1,w_2$ and the running couplings as thresholds; the setting (J) of Definition~\ref{def:weighted-history-setting} holds at $k=K-1$; $\log x_{K-1}\ge\max\{2p_0,1+\log L\}$; and $\delta(x_{K-1})$ lies below the threshold of the chain classes (equations~\YMMref{(2285)}{833} and~\eqref{eq:weighted-tower-smallness}) and below the smallness that Proposition~\ref{C-350-7} asks of its discs.
\end{enumerate}
No condition is put on the witness depth, on the signs or on the Lipschitz constants of $f$ and $h$.
\end{proposition}

\begin{proof}[Proof sketch]
(i) For a side at $g$ the lattice of a datum is large at every level, by the flow of the coupling (Section~\YMMref{74}{614}, \YMMref{74}{614}) and Lemma~\YMMref{126.2}{845}; at $j=K$ this is $\ell L\ge C_LM$. The condition on $N_{\mathsf H}$ and its two companion numbers is a member of the collective condition, which holds at the fixed constants and under which the geometric hypotheses of Section~\YMMref{52.7}{432} hold; it is a condition on these three numbers, and no side occurs in it. The bound $\ell\ge N_{\mathsf H}$ is a property of the sides at $g$: for $\ell=2L^m\in\mathfrak L(g)$ the flow of the coupling gives $LM\,R(x_K)\le L^m$ (Proposition~\YMMref{74.1}{615}; $R$ is the function of equation~\YMMref{(634)}{339}), so $M\,R(x_K)\le L^m$, and with this Lemma~\YMMref{126.2}{845} gives $\ell\ge N_{\mathsf H}$, together with the frames of the datum on $T_\ell$; the setting (J) at $k=K-1$ follows in (iv).

(ii) The run equation gives $|x_j-L^{-j}x_0|\le B_\sharp\sum_{i\ge0}L^{-i}=\bar B$, and $L^{-j}x_0=b/D_j$. As $K$ is the first passage, $x_j\ge x_{K-1}>X_g\ge Lg^{-2}$ for $j<K$, and $x_K=L^{-1}x_{K-1}-e_K>g^{-2}$; $\gamma_{\rm s}$ is so small that $X_g\ge12\bar B$. Hence $|\tau_j-1|\le\bar B/x_j\le\frac1{12}$, and $\frac{12}{13}b/D_j\le x_j\le b/D_j+\bar B\le2b/D_j$. (iii), (iv) $\gamma_{\rm s}\le\gamma^{\rm ref}_W$ is the least of two thresholds determined by the constants. For $g\le\gamma_{\rm s}$, Corollary~\ref{cor:weighted-canonical-data}, Proposition~\ref{prop:weighted-canonical-history} and Lemma~\YMMref{126.2}{845} give the data of the progression and the setting (J) at $k=K-1$; in that setting, at every level, the constants, the lattice of the step and $x_j$ are data under (Im), so that in particular the ratio $A_0/A_1$ lies above the number that (Im) asks.

As $x_{K-1}>L\gamma_{\rm s}^{-2}$ and $\delta$ decreases for large arguments, $\delta(x_{K-1})\le\delta(L\gamma_{\rm s}^{-2})$, which lies below both thresholds for $\gamma_{\rm s}$ small, and $\log x_{K-1}\ge\log(L\gamma_{\rm s}^{-2})\ge\max\{2p_0,1+\log L\}$. The sources enter through their supports alone.
\end{proof}

\begin{proposition}[One field scale]\label{C-350-7}
Under the assumptions of Proposition~\ref{C-350-6} the progression is formed at one complex field scale for all its steps. There are $R_j>1$ ($0\le j<K$) with $R_j^{-2}\le c_0(\log x_j)^{2p_0}D_j/b$, and $\rho\in(0,\frac12]$ with $(\rho R_j)^2\ge2$, such that, with $c_0,\rho,C_{\mathsf n},C_{\rm E},C_{{\rm E},m},C_{\rm R},C_{\rm J}$ determined by the constants (they depend on no cutoff, side, run, scale, position or pair), the following hold.
\begin{enumerate}
\item[(i)] $\lambda\mapsto\mathsf n^\lambda_{j+1}$ is holomorphic and even on $|\lambda|<R_j$, $\mathsf n^1_{j+1}=\mathsf n_{j+1}$, its Taylor coefficients $\mathsf n_{j+1,[0]}$, $\mathsf n_{j+1,[2]}$ of the orders $0$ and $2$ vanish, and for $|\lambda|\le\rho R_j$
\[
\begin{multlined}
|\mathsf n^\lambda_{j+1}|\le C_{\mathsf n}|\lambda|^2R_j^{-2}D_j^{-2}\min\{1,(d/D_j)^{5/2}\}\,\Pi_{j+1}\,\eta,\qquad\Pi_{j+1}=\max\{1, \\
(d/D_{j+1})^3\}e^{-[d/(4MD_{j+1})-2]_+}.
\end{multlined}
\]
\item[(ii)] $|\mathsf J_2-\mathsf h_4-\mathsf h_6|\le C_{\rm E}b^{-4}d^2l^{8p_0}\eta$ and $|\mathsf h_{2m}|\le C_{{\rm E},m}b^{-m}d^{m-2}l^{2mp_0}\eta$ for $m=2,3$.
\item[(iii)] For $0\le n\le K-1$ with $D_n\ge d$ and $m=2,3$: 
\[
\sum_{j=n}^{K-1}|\mathsf n_{j+1,[2m]}|\le C_{\rm R}b^{-m}l^{2mp_0}d^{m-2}(d/D_n)^{\frac52-(m-2)}\eta
.\]
In particular, for $n\le n'\le K$ with $L^n\ge r$: $|\mathsf H^{(n')}[D]-\mathsf H^{(n)}[D]|\le C_{\rm R}l^{6p_0}(r/L^n)^{3/2}r\,\eta$.
\item[(iv)] Let $\dot E^{(i+1),\lambda}_a=\partial_{\mu_a}E^{(i+1),\lambda}_\mu|_{\mu=0}$ ($a=1,2$) and $\ddot E^{(i+1),\lambda}=\partial_{\mu_1}\partial_{\mu_2}E^{(i+1),\lambda}_\mu|_{\mu=0}$ be the first and the second jets of the level $i+1$ at the field scale $\lambda$, pointed families of the level $i+1$, and let $\psi_i=\min\{1,(d/D_i)^3\}$, $\chi_i=\min\{1,(d/D_i)^{5/2}\}$. At every $\lambda$ at which they are formed, the member of $\dot E^{(i+1),\lambda}_a$ on a domain $X$ vanishes unless $X$ holds a plaquette of the support of $w_a$, and the member of $\ddot E^{(i+1),\lambda}$ on $X$ vanishes unless $X$ holds a plaquette of each of the two supports. For $0\le i\le j<K$ and $|\lambda|\le\rho R_j$, their sizes \YMMref{(759)}{375} at the rate $\kappa^+=(1+4\beta_{\rm s})\kappa$, at which the parts of the local terms that depend on the sources decay in the linear size of the domain (Proposition~\YMMref{122.1}{826}), obey
\[
\|\dot E^{(i+1),\lambda}_a\|_{i+1,\kappa^+}\le C_{\rm J}|\lambda|^2R_i^{-2}D_i^{-1}\psi_i\|w_a\|_\infty,\qquad\|\ddot E^{(i+1),\lambda}\|_{i+1,\kappa^+}\le C_{\rm J}|\lambda|^2R_i^{-2}D_i^{-2}\chi_i\,\eta .
\]
\end{enumerate}
If $w_a$ vanishes identically, all the jets in the direction $w_a$ vanish; so (i)--(iii) hold trivially if $w_1$ or $w_2$ vanishes, and so do the bounds of (iv) that involve a vanishing source.
\end{proposition}

\begin{proof}[Proof sketch]
One step at a complex field scale is Proposition~\ref{IV-71-37}, applied with a parameter set (Definition~\ref{IV-71-26}) that holds the field scales of the earlier steps; the scale is put on the diagonal afterwards. A lower calibration $\mathsf t_\flat$, fixed with the constants, gives room for the size of the deformed bare array, the bound of the step being uniform in the parameter. The counterterms $e_i$ are kept, and the coefficient of the action is moved along $\xi_j(\lambda)$, which stays in the admitted disc about $x_j$ for $|\lambda|<R_j:=\delta(\mathsf t_\flat)/(2\delta(x_j))$, because $|\xi_j(\lambda)-x_j|=|\lambda^2-1|\,|x_j-L^{-j}x_0|<2R_j^2\bar B$. To make the last bound smaller than $\sigma_*x_j$, decrease $\gamma_{\rm s}$, if necessary, so that
\[
 \bar B\delta(\mathsf t_\flat)^2
       <2\sigma_*A_1^2(\log x_j)^{2p_0}
\]
for every coupling of the run. This is possible by the lower bound on $x_j$ in Proposition~\ref{C-350-6}; substitution of the formula for $R_j$ gives $2R_j^2\bar B<\sigma_*x_j$. Thus the simultaneous substitution in the field scale and the quadratic coefficient stays within the product domain of Proposition~\ref{IV-71-37}.

By Proposition~\ref{C-350-6}(ii), $R_j^{-2}=4\delta(x_j)^2/\delta(\mathsf t_\flat)^2\le\frac{16}3A_1^2\delta(\mathsf t_\flat)^{-2}(\log x_j)^{2p_0}D_j/b$. Induction over the heights gives the holomorphy, the evenness, $\mathsf n^1_{j+1}=\mathsf n_{j+1}$ and the locality in the weight of (iv): the member of the level $i+1$ on a domain $X$ is formed from the members of the lower levels on the domains in $X$ and from the weight on the plaquettes of $X$, at every field scale as at the field scale $1$ (Proposition~\textup{\YMMref{122.1}{826}}(c)). The derivatives in the weights obey two linear recursions over the levels, by the chain rule and Cauchy's estimate in the direction of the bare array: one step at the field scale $\lambda$ multiplies the size of the derivative in the direction of a bare array by at most $c\,|\lambda|^2R_i^{-2}$ (clauses (cx3), (cx8) of Proposition~\ref{IV-71-37}, uniformly in the parameter), and the weight member of the level $i$ in the direction $w_a$ has the size at most $36M^3\alpha_0^2D_i^{-1}\psi_i\|w_a\|_\infty$, its source meeting at most $35(d/a_K)^3$ plaquettes; so the first derivatives follow the profile $\psi_i=\min\{1,(d/D_i)^3\}$, the second the profile $\chi_i=\min\{1,(d/D_i)^{5/2}\}$.

They close on these profiles for $|\lambda|\le\rho R_j$ by the locality in the weight, the bare array of a source with support of diameter $\le2d$ meeting at most $77\max\{1,(d/D_i)^3\}$ cubes of the level $i$, the cubes of side $MD_i$ of the partition $\pi_i$ of the torus whose wall-connected unions are the domains of the level $i$ (the level-$i$ localization partition of Section~\YMMref{45}{351}); $\rho$ is fixed by the two inequalities that close the induction, and the closed recursions are the two bounds of (iv). A domain that meets both supports holds two points at a distance at least $d/4$, so that its linear size in the units of its level is at least $[d/(4MD_{j+1})-2]_+$; with the count of such domains this gives the display of (i).

At the order $2$ the earlier members are taken at the field scale $0$, where they do not depend on the sources, and the coefficient is a covariance of two linear forms $\ell_Y$ of the field at the trivial configuration. A constant gauge rotation $u$ fixes the trivial configuration, so $\ell_Y(\operatorname{Ad}(u)B)=\ell_Y(B)$, and $\int_G\operatorname{Ad}(u)Y\,du=0$ for $Y\in\mathfrak g$ gives $\ell_Y=0$.

(ii), (iii): an even holomorphic $H$ with $|H|\le\mathsf B$ on $|\lambda|<R$ has $|H_{[2m]}|\le\mathsf BR^{-2m}$ and $|H(1)-\sum_{m\le p}H_{[2m]}|\le\mathsf BR^{-2p-2}/(1-R^{-2})$. Applied on $|\lambda|<\rho R_j$, where $(\rho R_j)^2\ge2$, with the bound of (i), each level gives at most $c\,b^{-m}(\log x_j)^{2mp_0}D_j^{m-2}\chi_j\Pi_{j+1}\eta$ at the order $2m$, and $m=4$ for the remainder. The levels with $D_j\ge d$ have $\Pi_{j+1}=1$, $\log x_j\le l$ and $D_j^{m-2}\chi_j=d^{m-2}(d/D_j)^{\frac52-(m-2)}$: a geometric series over the levels. The other levels are summed against $\Pi_{j+1}$, which is summable against every power of the logarithm. For (iii) at $m=3$: $\mathsf h^{(n')}_6-\mathsf h^{(n)}_6=\sum_{n\le j<n'}\mathsf n_{j+1,[6]}$, $D_n\ge d$ if and only if $L^n\ge r$, $d/D_n=r/L^n$ and $\beta^3a_K^2b^{-3}d=r$.
\end{proof}

\begin{proposition}[The front on the least torus]\label{C-350-8}
There are $\gamma_{\rm C}>0$, determined by the constants, and for $0<g\le\gamma_{\rm C}$ numbers $d_{\rm A}(g)>0$ and $C_{\rm A}(g)<\infty$, depending on $g$ and the constants and on no $K$, $\beta$, $b$, $\vartheta$, $d$, pair or position, such that for every datum $(K,\beta)$ at $g$, $\ell=\ell_g$, $0<d\le d_{\rm A}(g)$ with $s(d)\le K$, $\vartheta\in(0,1)$ and $(f,h)\in\mathfrak T_{T_\ell}(d;\vartheta)$: $K\ge3$ and $a_K\le d$; the set
\[
\mathsf{Pk}=\bigl\{v\in\delta\mathbb Z^3:\ |v_\mu|\le0.575-9d\ (\mu=1,2,3)\bigr\},\qquad\delta=a_K\lceil18d/a_K\rceil,
\]
which lies in $a_K\mathbb Z^3$ and depends on $a_K$ and $d$ alone, consists of vectors of length $<1$, has $\#\mathsf{Pk}\ge(30d)^{-3}$, and $(\tau_vf,\tau_vh)\in\mathfrak T_{T_\ell}(d;\vartheta)$ for $v\in\mathsf{Pk}$; and
\[
\Bigl|S^T_{2;K,\ell,\beta}(f,h)-\frac1{\#\mathsf{Pk}}\sum_{v\in\mathsf{Pk}}\mathsf J_2(\tau_vf,\tau_vh)\Bigr|\le C_{\rm A}(g)\vartheta^{-1}d^2b^{-2}\eta .
\]
\end{proposition}

\begin{proof}
Decrease $\gamma_{\rm C}\le\gamma_{\rm s}$ to satisfy the
source and radius conditions of Section~\ref{sec:levelwise-first-creation}.
These are conditions on the fixed constants and on the lower
coupling threshold. Choose $d_{\rm A}(g)$ small enough for
the packing and amplitude bounds below.

If $f\circ x=0$, every lattice translate of this bare
weight is zero. Both the covariance and its source-energy
mixed derivative therefore vanish; the assertion follows.
The same applies to $h\circ x$. Henceforth neither is zero.
Choose $x_1$ with $\operatorname{supp}f\subset\mathbb B(x_1,d)$.
Both supports lie in $\mathbb B(x_1,9d)$. The balls
$\mathbb B(x_1+v,9d)$, $v\in\mathsf{Pk}$, have disjoint
interiors, since $\delta\ge18d$, and lie in
$\mathbb B(x_1,1)$. Indeed,
$|v|\le\sqrt3(0.575-9d)$ and
$\sqrt3(0.575-9d)+9d<1$.
The condition $s(d)\le K$ gives
$L^Kd\ge C_W(s(d))\ge4L^2M$, so $a_K\le d$;
for $d$ small it also gives $K\ge3$.
Now $\delta\le18d+a_K\le19d$. Thus, for $d$ small,
\[
 \#\mathsf{Pk}=
  \left(2\left\lfloor\frac{0.575-9d}{\delta}\right\rfloor+1\right)^3
  \ge(30d)^{-3}.
\]
The translates preserve the support separation, amplitude
and Lipschitz conditions defining the test class.

Use the norm \eqref{eq:levelwise-source-norm} and the
holomorphic total source energy $\mathcal E$ of
Proposition~\ref{prop:levelwise-source-density}.
Lemma~\ref{lem:packed-levelwise-source} bounds a disjoint
sum $w=(\sum_\iota\zeta_\iota u_\iota)\circ x$ by
\[
 \|w\|_{\rm lw}\le
   \max_\iota|\zeta_\iota|\mathsf A_\iota/\mathsf A_\star(d),
 \qquad
 \mathsf A_\star(d)=\min\left\{\sigma_{\rm src}b^\vee,
   \frac{\mu_\star}{K_{\rm L}^\vee}
   \frac{b^\vee}{(b^\wedge)^{1/2}}
       \left(\frac{\vartheta d}{4\sqrt3M}\right)^{1/2}\right\},
\]
where $\|u_\iota\|_\infty\le\mathsf A_\iota$,
$\operatorname{Lip}u_\iota\le\mathsf A_\iota/(\vartheta d)$,
$b^\vee=g^{-2}-\bar B$ and $b^\wedge=b_+$.
The total-energy identity in
Proposition~\ref{prop:levelwise-source-density} shows that
$\mathsf J_2$ is the mixed derivative of $\mathcal E$ at
zero. Therefore $S^T_2-\mathsf J_2$ is the mixed derivative
of $\log(Z[\cdot]/Z[0])-\mathcal E$.

For signs $\epsilon\in\{\pm1\}^{\mathsf{Pk}}$ put
\[
 F^\epsilon=\sum_v\epsilon_v\tau_vf,\qquad
 H^\epsilon=\sum_v\epsilon_v\tau_vh,\qquad
 w_z=(z_1F^\epsilon+z_2H^\epsilon)\circ x,
 \qquad
 \Omega_\epsilon(z)=Z[w_z]e^{-\mathcal E(w_z)}/Z[0].
\]
The nonzero sets of the individual bumps are pairwise
disjoint. With $u_1=\|f\|_\infty/\mathsf A_\star(d)$ and
$u_2=\|h\|_\infty/\mathsf A_\star(d)$, the amplitude bound
gives $\|w_z\|_{\rm lw}\le\max_a|z_a|u_a$.
Thus $\Omega_\epsilon$ is holomorphic on
$\mathbb P=\{|z_1|u_1<1,\ |z_2|u_2<1\}$ and
$\Omega_\epsilon(0)=1$.

Split the sourced partition by its first large-field
creation, using \eqref{eq:first-creation-split}:
\[
 \Omega_\epsilon=\Omega_\epsilon^\emptyset
                  +\sum_{t=1}^K\Omega_\epsilon^t,\qquad
 \Omega_\epsilon^\emptyset=
        e^{-\mathcal E(w_z)}m_K^{w_z}(\mathcal V_K)/Z[0],
 \qquad
 \Omega_\epsilon^t=
        e^{-\mathcal E(w_z)}\mathsf{Cr}_t(w_z)/Z[0].
\]
Every source here is supported in the one ball
$\mathbb B(x_1,1)$. Proposition~\ref{prop:first-creation-envelopes}
therefore gives
\[
 |\Omega_\epsilon^\emptyset|\le e^{\mathsf G_K}
       \quad\hbox{on }\mathbb P,\qquad
 |\Omega_\epsilon^t|\le
 e^{\|w_z\|_{\rm lw}(\mathsf G_{t-1}+\mathsf H_{t-1})}\mathsf w_t.
\]
Choose $\mathsf C_{\rm cnt}=C_{\rm cnt}(g)\ge1$ in
\eqref{eq:first-creation-count}, and set
$\mathsf{Cnt}_t=\mathsf C_{\rm cnt}L^{4(K-t+1)}$.
On $\mathbb P_t=\{|z_a|u_a\mathsf{Cnt}_t<1, a=1,2\}$
the second bound is at most $e\mathsf w_t$.
For $\emptyset\ne B\subset\{1,2\}$ write
$[z^B]\Omega=(\prod_{a\in B}\partial_{z_a})\Omega(0)$ and
$u^B=\prod_{a\in B}u_a$. Cauchy's formula on these polydiscs
gives
\[
 |[z^B]\Omega_\epsilon^\emptyset|\le e^{\mathsf G_K}u^B,
 \qquad
 |[z^B]\Omega_\epsilon^t|
    \le e\mathsf w_t\mathsf{Cnt}_t^{|B|}u^B
    \le e(\mathsf C_{\rm cnt}L^4)^2
                         L^{8(K-t)}\mathsf w_tu^B.
\]
By \eqref{eq:first-creation-rarity}, the latter coefficients
sum to at most $e(\mathsf C_{\rm cnt}L^4)^2C_8(g)u^B$.
Choose $\mathsf Q(g)\ge1$ above this coefficient bound plus
$e^{\sup\mathsf G_K}$. It is finite and independent of
the datum, the positions and the signs. Thus
\[
 |[z^B]\Omega_\epsilon|\le\mathsf Q(g)u^B,
 \qquad
 |\partial_{12}\log\Omega_\epsilon(0)|
 \le(\mathsf Q+\mathsf Q^2)u_1u_2
 \le2\mathsf Q^2\mathsf A_\star(d)^{-2}\eta.
\]
The logarithm is needed only near zero, where
$\Omega_\epsilon(0)=1$ ensures its existence.

By bilinearity the last mixed derivative is
\[
 \sum_{v,v'}\epsilon_v\epsilon_{v'}
       (S^T_2-\mathsf J_2)(\tau_vf,\tau_{v'}h).
\]
Average over the independent signs. The terms $v\ne v'$
vanish and the diagonal terms remain. Translation
invariance of $S^T_2$ and the packing count then give
\[
 \left|S^T_2(f,h)-\frac1{\#\mathsf{Pk}}
                   \sum_v\mathsf J_2(\tau_vf,\tau_vh)\right|
 \le(30d)^3\,2\mathsf Q^2\mathsf A_\star(d)^{-2}\eta.
\]
For $d\le4\sqrt3Mg^{-2}(K_{\rm L}^\vee/K_{\rm Q}^\vee)^2$
the second entry in the minimum defining $\mathsf A_\star$
is the smaller one. Substitution bounds the last expression
by $C(g)\vartheta^{-1}d^2\eta$. Since the label satisfies
$b\le b_+$, enlarging this fixed-$g$ constant by $b_+^2$
gives the asserted $C_{\rm A}(g)\vartheta^{-1}d^2b^{-2}\eta$.
The finitely many ceilings on $d$ determine $d_{\rm A}(g)$.
\end{proof}

\begin{proposition}[The order four]\label{C-350-9}
Let $0<g\le\gamma_{\rm s}$, $(K,\beta)$ a datum at $g$, $\ell\in\mathfrak L(g)$, $0<d\le L^{-1}$, $s(d)<K$, $20d\le\ell$, $\vartheta\in(0,1)$ and $(f,h)\in\mathfrak T_{T_\ell}(d;\vartheta)$, at any position. Then $K\ge3$, $8a_K\le d$, $d/a_K\ge L(\log\frac1d)^2$, and
\[
|\mathsf h_4-c_2b^{-2}\mathcal K_6|(f,h)\le\bigl[\bar\varepsilon_4(g,d)\,b^{-3}d+C_{\rm lat}\,b^{-2}a_K/d\bigr]\eta,
\]
where, with $l'=\max\{1,\log\frac{2b'}d\}$,
\[
\bar\varepsilon_4(g,d)=\sup_{g^{-2}<b'\le b_+}\frac{b'}d\Bigl\{3C_{\rm tm}d^3+C_{\rm Bx}\bigl[d^{3/2}+e^{-\mathsf c_\chi(\log\frac{3b'}{4d^{1/2}})^2}\bigr]+C_{\rm Bt}\,l'^{\,4p_0}e^{-\mathsf c_\flat L(\log\frac1d)^2}\Bigr\},
\]
$C_{\rm Bx},C_{\rm Bt},\mathsf c_\chi>0$ being determined by the constants and $\mathsf c_\flat=\min\{\frac1{8ML},\frac{\delta_{\rm tp}}{2\sqrt3}\}$. The function $\bar\varepsilon_4(g,\cdot)$ depends on $g$ and the constants alone and is the same for all sides and data at $g$; and $\bar\varepsilon_4(g,d)\le C_4(g)d^{1/2}$, with $C_4(g)$ depending on $g$ and the constants.
\end{proposition}

\begin{proof}[Proof sketch]
The first three assertions are those of the setting on the witness depth. Let $2\le n\le K$ with $D_n\ge d$.

(1)

\emph{The tower at a real field scale.} Let $0<\lambda\le1$, $g_i=x_i^{-1/2}$, and let $c_{\mathfrak h}$ be the constant of the shift systems of the steps. The chain of the levels at the field scale $\lambda$ is formed downwards from the flat top of the level $n$: $A^\lambda_n:=0$ and, for $i=n-1,\dots,1$,
\[
A^\lambda_i:=A^{\rm lw}\bigl(v^{[i+1]}_{A^\lambda_{i+1}},\lambda g_iB_i\bigr),\qquad\tilde U^\lambda_i:=e^{iL^{-i}A^\lambda_i},
\]
with the lowered $1$-form $A^{\rm lw}$ of equation~\YMMref{(2286)}{833} and the chart point $v^{[i+1]}_A$ of a chain class (equation~\eqref{eq:weighted-chain-class}). A \emph{chart point} of a level $k$ is a pair $v=(A,\mathbf J)$ of $1$-forms on the bare bonds; $P[v]=(\mathbf 1,A,\mathbf J)$ is its triple in the chart at the trivial pair (Section~\YMMref{124}{832}), and $P[0]$ is the trivial triple; its norm $\|v\|$ is the largest of the numbers $|A(b)|$, $|\mathbf J(b)|$ over the bare bonds $b$ and of the moduli of the difference quotients of $A$ at the spacing $L^{-k}$; and $\mathsf{Cc}_k(r)$ is the chain class of the level $k$ and the radius $r$. The $A^\lambda_i$ are real $1$-forms on the bare bonds, and $\tilde U^\lambda_i$, a configuration on the bare bonds, is the \emph{background of the level $i$}, as in the proof of Proposition~\ref{C-350-11}(i): up to a gauge transformation it is a background of the level $i$ in the sense of the actual-background condition in \eqref{eq:weighted-chain-class}; $\mathcal O(w_a)$ and the $a_p$ are taken at it.

By Proposition~\YMMref{124.2}{834}, equation~\YMMref{(2288)}{834}, read at the field $\lambda B_i$, which lies in the box, as the box is a product of Euclidean balls, and by the bounds of equation~\YMMref{(2287)}{834} and its domain verification at $\sup|\lambda g_iB_i|\le\lambda\delta(x_i)$, $A^\lambda_i\in\mathsf{Cc}_i(\lambda r_i)$, with $r_n:=0$ and $r_i:=L^{-1}r_{i+1}+\frac{11}{10}c_{\mathfrak h}\delta(x_i)\le\frac85c_{\mathfrak h}\delta(x_i)$; the conditions of that Proposition hold at every height, and $r_1\le\frac12r_\alpha$, $r_\alpha$ being the usable radius of the fixed radii (equations~\YMMref{(933)}{427} and~\YMMref{(2285)}{833}), as $\delta(x_{K-1})$ lies below the threshold of the chain classes (Proposition~\ref{C-350-6}(iv)).

The same Proposition gives $v^{[i]}_{A^\lambda_i}=v^{\rm lw}(v^{[i+1]}_{A^\lambda_{i+1}},\lambda g_iB_i)$, the lowered point of equation~\YMMref{(2286)}{833}; the lowered points are holomorphic (the holomorphy following equation~\YMMref{(2286)}{833}) and obey the bounds of equation~\YMMref{(2287)}{834} and its domain verification at complex points as well, so the same recursion, with the lowered points, continues the chain to the complex field scales $|\lambda|\le1$, where the $A^\lambda_i$ take complex values and the $\tilde U^\lambda_i$ lie in $\mathrm{SL}(2,\mathbb C)$; one name, $\tilde U^\lambda_i$, serves for both. For $\mu\in\mathbb C^2$ near $0$ put $w_\mu=\mu_1w_1+\mu_2w_2$ and
\[
D^{\lambda,\mu}_i=\mathcal O(w_\mu)+\sum_{i'\le i}\mathsf{Sum}\bigl[E^{(i'),\lambda}_\mu-E^{(i'),\lambda}_0\bigr],
\]
$\mathsf{Sum}[\cdot]$ being the sum of a pointed family at a configuration or at a triple (the definitions preceding equation~\YMMref{(2276)}{831}). Each step at a height $1\le i<n$ obeys an identity of its own at the field scale $\lambda$, whose case $\lambda=1$ uses the quotient argument of Proposition~\ref{prop:weighted-pair-tower}: for $r\le r_{i+1}$, $A\in\mathsf{Cc}_{i+1}(r)$, $v=v^{[i+1]}_A$, $\tilde U=e^{iL^{-i-1}A}$ and $\tilde U^{\rm lw}(B)=e^{iL^{-i}A^{\rm lw}(v,\lambda g_iB)}$,
\[
\exp D^{\lambda,\mu}_{i+1}(\tilde U)=\int\exp D^{\lambda,\mu}_i\bigl(\tilde U^{\rm lw}(B)\bigr)\,d\nu^\lambda_i[v](B).
\]
Here $\nu^\lambda_i[v]$ is the complex measure of mass $1$ on the box of the height $i$ whose density is the integrand of the integral form of the step at the field scale $\lambda$, at the source $0$ and at the triple $P[v]$ (clause (cx10) of Proposition~\ref{IV-71-37}), divided by its integral, which is not $0$: the Gaussian density of the covariance $\tau_i(\lambda)^{-1}\mathbf C_\ast(P[v])$, with $\tau_i(\lambda)=\xi_i(\lambda)/x_i$, positive for real $|\lambda|\le1$, and the whole covariance $\mathbf C_\ast$ of the step, times $\exp[-\frac{1-\tau_i(\lambda)}2\sum_\square(B,\hat{\mathsf G}_\square(P[v]|_\square)B)]$, with the gauge-fixing operators $\hat{\mathsf G}_\square$ of the fine cubes (equation~\YMMref{(823)}{394}), times the exponential of the potentials of the step at the field scale $\lambda$. At complex $\lambda$ this is the holomorphic continuation supplied by the same one-step theorem, with its nonvanishing normalizer; no positivity is asserted there.

The identity is proved as follows. The integral form of the step at the triple $P[v]$, at the source $\mu$ and at the source $0$ (Proposition~\ref{IV-71-37}(cx10), at the data of Proposition~\ref{C-350-7}), and the quotient of the two give
\[
\exp\mathsf{Sum}\bigl[E^{(i+1),\lambda}_\mu-E^{(i+1),\lambda}_0\bigr](P[v])=\int\exp\bigl[\mathbb V'[\mathfrak g]^{\rm tot}(P[v];\lambda B)\bigr]\,d\nu^\lambda_i[v](B),
\]
because $\mu$ stands in the rule of the step alone and the rule is linear in the bare array, $\mathfrak g$ being the difference of the bare arrays of the height $i$ at $\mu$ and at $0$, whose members are the $E^{(i'),\lambda}_\mu-E^{(i'),\lambda}_0$, $i'\le i$, and the weight member of $w_\mu$. Proposition~\YMMref{124.2}{834}, equation~\YMMref{(2289)}{834}, which holds for every bare array and every field of the box, read at $\lambda B$, in the plaquette form of equation~\YMMref{(2290)}{834} and by linearity, gives $\mathbb V'[\mathfrak g]^{\rm tot}(P[v];\lambda B)=D^{\lambda,\mu}_i(\tilde U^{\rm lw}(B))-D^{\lambda,\mu}_i(\tilde U)$; and the proof of Proposition~\YMMref{124.2}{834} gives $\mathsf{Sum}[\cdot](P[v])=\mathsf{Sum}[\cdot](\tilde U)$ for the family of the level $i+1$. Multiplying by $\exp D^{\lambda,\mu}_i(\tilde U)$ gives the identity. Iterating it from the height $n-1$, where $A=0$ and $\tilde U=\mathbf 1$, down to the height $1$ along the chain, and taking the mixed derivative in $\mu$ at $0$ of the logarithms, gives
\[
\sum_{j=2}^n\mathsf n^\lambda_j=\bigl\langle\mathcal F^\lambda_1(\tilde U^\lambda_1);\mathcal F^\lambda_2(\tilde U^\lambda_1)\bigr\rangle_{\mathbb Q^\lambda_n}+\int\bigl[\mathsf F^\lambda_1(\tilde U^\lambda_1)-\mathsf F^\lambda_1(\mathbf 1)\bigr]d\mathbb Q^\lambda_n,
\]
where $\mathbb Q^\lambda_n(dB)=\prod_{1\le i<n}\nu^\lambda_i\bigl[v^{[i+1]}_{A^\lambda_{i+1}}\bigr](dB_i)$ is a complex measure of mass $1$ on the product of the boxes, with a bounded continuous density, each factor being formed at the chart point that the chain gives above its height; $\mathcal F^\lambda_a=\mathcal O(w_a)+\mathsf{Sum}[\dot E^{(1),\lambda}_a]$ is $\mathcal O(w_a)$ plus the first jet of the level $1$ in the direction $w_a$, and $\mathsf F^\lambda_1=\mathsf{Sum}[\ddot E^{(1),\lambda}]$ is the second jet of the level $1$, with $\mathsf F^\lambda_1(\mathbf 1)=\mathsf n^\lambda_1$ (the jets of Proposition~\ref{C-350-7}(iv)).

(2)

\emph{The limit $\lambda\to0$.} The lattice being fixed, let $\lambda\downarrow0$ in the identity of (1) divided by $\lambda^4$. Its left side tends to $\sum_{j=2}^n\mathsf n_{j,[4]}$, by Proposition~\ref{C-350-7}(i).

\emph{The measures.} The chart points of the chain have $\sup_B\|v^{[i+1]}_{A^\lambda_{i+1}}\|\le\lambda r_{i+1}\to0$, and $\tau_i(\lambda)\to\tau_i(0)=\tau_i$. Two facts tie the step at the trivial configuration to the free template: the whole covariance of the step at the trivial triple is $(\mathsf A^{\rm cv}_i)^{-1}\otimes1$ (equation~\YMMref{(2297)}{836}; Proposition~\YMMref{120.3}{819}); and the gauge-fixing operators of the fine cubes at the trivial triple add up to the gauge-fixing operator at the block field $\mathbf 1$, which is $\hat{\mathsf G}\otimes1$, as the fine cubes partition the sites of the level $i+1$ and the operator of a cube is the gauge-fixing form restricted to the blocks of that cube (equation~\YMMref{(823)}{394}; Proposition~\YMMref{120.3}{819}).

The covariance and the gauge-fixing operators at the triple $P[v]$ tend to these, entry by entry, as $v\to0$ (equation~\YMMref{(2296)}{836}; the kernels of the $\hat{\mathsf G}_\square$ are holomorphic in the chart point), and the potentials at the field scale $\lambda$ are at most $c\lambda$ times a function bounded on the box (\eqref{eq:complex-input-bounds} of Definition~\ref{IV-71-26}). So the density of $\mathbb Q^\lambda_n$ tends, uniformly on the product of the boxes, to the product over the heights of the densities proportional to
\[
e^{-\frac12(B_i,(\tau_i\mathsf A^{\rm cv}_i+(1-\tau_i)\hat{\mathsf G})B_i)}=e^{-\frac12(B_i,\mathsf M_iB_i)} ;
\]
as the boxes have finite volume and the limits of the normalisations are positive, $\mathbb Q^\lambda_n\to\mathbb Q^0_n:=\bigotimes_{1\le i<n}\nu^{\rm r}_i$ in total variation.

\emph{The plaquette terms.} At every height the linear part of one step of the chain in the field $B_i$, read in the plaquette angle, is $\mathsf{sg}_i(\mathsf T_i\otimes1)$ with a sign $\mathsf{sg}_i\in\{1,-1\}$ (Lemma~\YMMref{124.3}{835}, which reads $\ell\ge N_{\mathsf H}$ of Proposition~\ref{C-350-6}(i); its map $\mathsf T^{\rm cv}_i$ is $\mathsf T_i$, the definitions preceding Lemma~\YMMref{124.3}{835}). With the expansion of the plaquette terms and of the plaquette angles to the second order at one step (Lemma~\YMMref{124.4}{835}), summed along the chain, this gives, uniformly on the boxes and for every bare plaquette $p$,
\[
\bigl|a_p(\tilde U^\lambda_1)-\tfrac12\lambda^2|\Theta_p(B)|^2\bigr|\le c_\Theta\lambda^3,\qquad\Theta_p(B):=\sum_{1\le i<n}\mathsf{sg}_i\,x_i^{-1/2}(\mathsf T_iB_i)(p),
\]
with $c_\Theta$ determined by the constants. Hence $\lambda^{-2}\mathcal O(w_a)(\tilde U^\lambda_1)\to q_a(B):=\frac12a_K^{-1}\sum_pw_a(p)|\Theta_p(B)|^2$ uniformly on the boxes.

\emph{The level $1$.} The local terms of the level $1$ at $\tilde U^\lambda_1$ take the values of their members at the restrictions of the triple $P[v]$, $v=v^{[1]}_{A^\lambda_1}$, as in the proof of Proposition~\YMMref{124.2}{834}. For a pointed family $\mathsf f$ of the level $1$, a domain $X$ of the level $1$, a site $z$ of $X$ and $\kappa'>0$, $d_1(X)$ being the linear size of $X$ in the units of the level $1$,
\[
\bigl|\mathsf f\bigl(X,P[v]|_X,z\bigr)-\mathsf f(X,\mathbf 1,z)\bigr|\le2\|\mathsf f\|_{1,\kappa'}e^{-\kappa'd_1(X)}\lambda^2(r_1/r_\alpha)^2 :
\]
the function $\zeta\mapsto\mathsf f(X,P[\zeta v]|_X,z)$ is holomorphic on $|\zeta|<r_\alpha/\|v\|$, where $r_\alpha/\|v\|\ge r_\alpha/(\lambda r_1)\ge2$, and bounded there by $\|\mathsf f\|_{1,\kappa'}e^{-\kappa'd_1(X)}$; its derivative at $0$ vanishes, because a constant gauge rotation $u$ fixes the trivial configuration and acts on the chart point by $\operatorname{Ad}(u)$, the local terms are invariant, and $\int_G\operatorname{Ad}(u)Y\,du=0$ for $Y\in\mathfrak g$; Cauchy's estimate bounds the rest of the series. At $\mathsf f=\dot E^{(1),\lambda}_a$, whose size at the rate $\kappa^+$ is at most $C_{\rm J}\lambda^2R_0^{-2}a_K^{-1}\|w_a\|_\infty$ (Proposition~\ref{C-350-7}(iv) at $i=0$, where $\psi_0=1$), this shows that the first jet of the level $1$ changes along the chain by $O(\lambda^4)$, uniformly on the boxes; as a covariance does not change when constants are added, $\lambda^{-4}\langle\mathcal F^\lambda_1(\tilde U^\lambda_1);\mathcal F^\lambda_2(\tilde U^\lambda_1)\rangle_{\mathbb Q^\lambda_n}\to\langle q_1;q_2\rangle_{\mathbb Q^0_n}$.

At $\mathsf f=\ddot E^{(1),\lambda}$, whose size is at most $C_{\rm J}\lambda^2R_0^{-2}a_K^{-2}\eta$ and whose members vanish on the domains that do not meet both supports (Proposition~\ref{C-350-7}(iv) at $i=0$), the count of the domains of the level $1$ that meet both supports, weighted with $e^{-\kappa^+d_1(X)}$ (Lemma~\ref{lem:source-support-sum}), brings the factor $\Pi_1$, as in the proof of Proposition~\ref{C-350-7}(i), and gives
\[
\bigl|\mathsf F^\lambda_1(\tilde U^\lambda_1)-\mathsf F^\lambda_1(\mathbf 1)\bigr|\le\lambda^4\,\mathrm{Er}_\flat,\qquad\mathrm{Er}_\flat=c\,R_0^{-2}a_K^{-2}(r_1/r_\alpha)^2\,\Pi_1\,\eta .
\]
So the limit $\mathsf e_\flat$ of $\lambda^{-4}\int[\mathsf F^\lambda_1(\tilde U^\lambda_1)-\mathsf F^\lambda_1(\mathbf 1)]\,d\mathbb Q^\lambda_n$ exists, $|\mathsf e_\flat|\le\mathrm{Er}_\flat$, the total variation of $\mathbb Q^\lambda_n$ tending to $1$, and
\[
\sum_{j=2}^n\mathsf n_{j,[4]}=\langle q_1;q_2\rangle_{\mathbb Q^0_n}+\mathsf e_\flat .
\]
The covariance $\langle q_1;q_2\rangle_{\mathbb Q^0_n}$ does not depend on the signs $\mathsf{sg}_i$: every $\nu^{\rm r}_i$ is invariant under $B_i\mapsto-B_i$, its box being symmetric, so that the law of $(\mathsf{sg}_iB_i)_i$ under $\mathbb Q^0_n$ is $\mathbb Q^0_n$. By Proposition~\ref{C-350-7}(i) at $j=0$ and Cauchy's estimate on $|\lambda|=\rho R_0$, $|\mathsf n_{1,[4]}|\le C_{\mathsf n}\rho^{-2}R_0^{-4}a_K^{-2}\Pi_1\eta$. As $x_0=\beta=b/a_K$, $R_0^{-2}\le c_0(\log x_0)^{2p_0}a_K/b$, $r_1\le\frac85c_{\mathfrak h}\delta(x_1)$ and $\delta(x_1)^2\le\frac43A_1^2(\log x_0)^{2p_0}La_K/b$ (Proposition~\ref{C-350-6}(ii)), both $|\mathsf e_\flat|$ and $|\mathsf n_{1,[4]}|$ are at most $c\,b^{-2}(\log\frac b{a_K})^{4p_0}\Pi_1\eta$.

With $u_0=d/a_K\ge8$, $\log\frac b{a_K}\le l\,(1+\log u_0)$ and $\Pi_1\le\max\{1,u_0^3\}e^{2-u_0/(4ML)}$; as $(1+\log u)^{4p_0}\max\{1,u^3\}e^{2-u/(8ML)}$ is bounded on $u\ge8$ and $\mathsf c_\flat\le\frac1{8ML}$, this is below $C\,b^{-2}l^{4p_0}e^{-\mathsf c_\flat d/a_K}\eta$.

(3)

\emph{The boxes out at the order four.} The cumulants of the orders $2$ and $4$ of the components $B^\alpha_i(b)$ in $\nu^{\rm r}_i$ and in $\nu^\circ_i$ differ by at most $C_\chi e^{-\mathsf c_\chi(\log x_i)^{2p_0}}$, uniformly in the volume, with $C_\chi$, $\mathsf c_\chi$ determined by the constants (Lemma~\ref{lem:ward-box-cumulants}, read with the one quadratic system of the height and the threshold $\mathfrak b(x_i)$); the odd cumulants vanish in both measures by the symmetry $B\mapsto-B$, and the Gaussian ones are $\delta_{\alpha\beta}\mathsf M_i^{-1}(b,b')$ at the order $2$ and $0$ at the order $4$.

(4)

\emph{Wick's rule and the slices.} The $B_i$ are independent in $\mathbb Q^0_n$ and each $\nu^{\rm r}_i$ is even. With $F^\alpha_i(p)=\mathsf{sg}_i\,x_i^{-1/2}\sum_b\mathsf T_i(p,b)B^\alpha_i(b)$, so that $|\Theta_p|^2=\sum_\alpha\bigl(\sum_iF^\alpha_i(p)\bigr)^2$ in an orthonormal basis $(t_\alpha)$ of $\mathfrak g$, $c^{\alpha\beta}_i(p,q)=\langle F^\alpha_i(p)F^\beta_i(q)\rangle_{\nu^{\rm r}_i}$ and $\kappa^{\alpha\beta}_i(p,q)$ the joint fourth cumulant of $F^\alpha_i(p),F^\alpha_i(p),F^\beta_i(q),F^\beta_i(q)$,
\[
\langle q_1;q_2\rangle_{\mathbb Q^0_n}=\tfrac14a_K^{-2}\sum_{p,q}w_1(p)w_2(q)\sum_{\alpha,\beta}\Bigl[2\Bigl(\sum_ic^{\alpha\beta}_i(p,q)\Bigr)^2+\sum_i\kappa^{\alpha\beta}_i(p,q)\Bigr].
\]
By the slice identity $\mathsf T_i\mathsf M_i^{-1}\mathsf T_i^{\mathsf T}=L^{-i}\tau_i^{-1}\mathcal C_{(i+1)}$ (Proposition~\YMMref{120.3}{819}, valid at every $\tau_i\ne0$) and
\[
x_i^{-1}L^{-i}\tau_i^{-1}=x_0^{-1},
\]
the running couplings drop out: in the Gaussian measures $c^{\alpha\beta}_i=\delta_{\alpha\beta}x_0^{-1}\mathcal C_{(i+1)}$. As $a_Kx_0=b$, the Gaussian value of $\langle q_1;q_2\rangle$ is $\frac32b^{-2}\sum_{p,q}w_1(p)w_2(q)\Sigma(p,q)^2$ with
\[
\Sigma=\sum_{1\le i<n}\mathcal C_{(i+1)}=C^{\rm fl}_{K,K-n}-\mathcal C_{(1)} .
\]
The differences of (3) are carried to the $F^\alpha_i$ through the row sums of $\mathsf T_i$, which are at most $C^{\rm row}_{\rm tp}L^{-2i}$ at every height, uniformly in the volume, as $\ell\ge N_{\mathsf H}$ (Proposition~\YMMref{120.6}{821}; Proposition~\ref{C-350-6}(i)). By multilinearity, with $\epsilon_i=C_\chi e^{-\mathsf c_\chi(\log x_i)^{2p_0}}$, $\mathsf{sg}_i^2=1$ and $x_i^{-1}L^{-4i}=\tau_ix_0^{-1}L^{-3i}\le\frac43x_0^{-1}L^{-3i}$,
\[
\bigl|c^{\alpha\beta}_i(p,q)-\delta_{\alpha\beta}x_0^{-1}\mathcal C_{(i+1)}(p,q)\bigr|\le\bar e_i:=\tfrac43(C^{\rm row}_{\rm tp})^2x_0^{-1}L^{-3i}\epsilon_i,\qquad|\kappa^{\alpha\beta}_i(p,q)|\le\tfrac{16}9(C^{\rm row}_{\rm tp})^4x_0^{-2}L^{-6i}\epsilon_i .
\]
There are at most $35^2(d/a_K)^6$ pairs $(p,q)$ with $w_1(p)w_2(q)\ne0$, and for them $|x(p)-x(q)|_\infty\ge d/\sqrt3$; with $u_i=d/D_i$ one has $L^{-3i}(d/a_K)^3=u_i^3$, and Proposition~\YMMref{120.6}{821} gives $(d/a_K)^3(|\Sigma|+|\mathcal C_{(1)}|)(p,q)\le C_{\rm tp}S_\Sigma$, $S_\Sigma=\sup_{u>0}\sum_{s\in\mathbb Z}(uL^s)^3e^{-\delta_{\rm tp}uL^s/\sqrt3}$. Expanding the square, with $\mathsf E_k=\sum_{1\le i<n}u_i^k\epsilon_i$ for $k=3,6$,
\[
\begin{multlined}
\Bigl|\langle q_1;q_2\rangle_{\mathbb Q^0_n}-\tfrac32b^{-2}\sum_{p,q}w_1(p)w_2(q)\Sigma(p,q)^2\Bigr|\le\tfrac{35^2}4\,b^{-2}\bigl[16C_{\rm tp}S_\Sigma(C^{\rm row}_{\rm tp})^2\mathsf E_3+32(C^{\rm row}_{\rm tp})^4\mathsf E_3^2 \\
+16(C^{\rm row}_{\rm tp})^4\mathsf E_6\bigr]\eta .
\end{multlined}
\]
The numbers $\mathsf E_k$ are bounded in three ranges of levels. As $\log x_i\ge\log\frac{3b}{4D_i}\ge0$ (Proposition~\ref{C-350-6}(ii), and $\frac34b\ge1\ge D_i$), $\log x_i\ge1$ and $2p_0\ge2$, one has $\epsilon_i\le C_\chi e^{-\mathsf c_\chi(\log\frac{3b}{4D_i})^2}$. At the levels with $D_i<d$, $\log\frac{3b}{4D_i}=\Lambda_b+t_i$ with $\Lambda_b=\log\frac{3b}{4d}\ge\log\frac{3b}{4d^{1/2}}$ and $t_i=\log u_i>0$, the $t_i$ being spaced by $\log L$, and $(\Lambda_b+t)^2\ge\Lambda_b^2+t^2$, so that $u_i^k\epsilon_i\le C_\chi e^{-\mathsf c_\chi\Lambda_b^2}e^{6t_i-\mathsf c_\chi t_i^2}$ for $k\le6$; so these levels give at most $C_\chi S_\chi e^{-\mathsf c_\chi\Lambda_b^2}$, with $S_\chi=\sup_{t_0\ge0}\sum_{s\ge0}e^{6(t_0+s\log L)-\mathsf c_\chi(t_0+s\log L)^2}$. At the levels with $d\le D_i\le d^{1/2}$, $\sum u_i^3\le2$ and $\epsilon_i\le C_\chi e^{-\mathsf c_\chi(\log\frac{3b}{4d^{1/2}})^2}$. At the levels with $D_i>d^{1/2}$, $\sum u_i^3\le2d^{3/2}$ and $\epsilon_i\le C_\chi$. So
\[
\mathsf E_3,\ \mathsf E_6\le C_\chi(S_\chi+2)\bigl[d^{3/2}+e^{-\mathsf c_\chi(\log\frac{3b}{4d^{1/2}})^2}\bigr]\le2C_\chi(S_\chi+2),
\]
and the bracket above gives the member $C_{\rm Bx}b^{-2}[d^{3/2}+e^{-\mathsf c_\chi(\log(3b/(4d^{1/2})))^2}]\eta$, with $C_{\rm Bx}$ formed with $C_{\rm tp}$, $S_\Sigma$, $C^{\rm row}_{\rm tp}$, $C_\chi$ and $S_\chi$, and so determined by the constants. The missing first slice is small: by Proposition~\YMMref{120.6}{821}, $(d/a_K)^3|\mathcal C_{(1)}(p,q)|\le C_{\rm tp}u_0^3e^{-\delta_{\rm tp}u_0/\sqrt3}$, $u_0=d/a_K$, at $|x(p)-x(q)|\ge d/\sqrt3$; with the terms of the level $1$ of (2) it gives the member $C_{\rm Bt}l^{4p_0}e^{-\mathsf c_\flat d/a_K}b^{-2}\eta$.

(5)

\emph{The free value.} Lemma~\ref{lem:ward-fluctuation-template} (the template sums) gives $|\sum_{p,q}w_1w_2(C^{\rm fl}_{K,K-n})^2-\frac3{8\pi^2}\mathcal K_6|\le C_{\rm tm}[\frac{a_K}{\vartheta d}+(\frac d\ell)^3+(\frac d{D_n})^3]\eta$, and $\frac32\cdot\frac3{8\pi^2}=c_2$. Take $n=K$, $D_K=1$: $(d/\ell)^3+d^3\le2d^3$, the member $\frac32C_{\rm tm}\vartheta^{-1}b^{-2}a_K/d$ is $C_{\rm lat}b^{-2}a_K/d$, and $b^{-2}=\frac bd\,b^{-3}d$; the supremum over the labels at $g$ and $d/a_K\ge L(\log\frac1d)^2$ give $\bar\varepsilon_4$. Finally $\log\frac{3b'}{4d^{1/2}}\ge\frac12\log\frac1d$, as $\frac34b'\ge1$; so $\frac{b'}de^{-\mathsf c_\chi(\cdot)^2}\le b_+e^{\log\frac1d-\frac14\mathsf c_\chi(\log\frac1d)^2}$, $\frac{b'}dC_{\rm Bx}d^{3/2}\le C_{\rm Bx}b_+d^{1/2}$, and the other members are smaller: $\bar\varepsilon_4(g,d)\le C_4(g)d^{1/2}$.
\end{proof}

\begin{remark}\label{C-350-32}
With Proposition~\ref{C-350-7}(ii) at $m=3$, Proposition~\ref{C-350-9} gives, on every side $\ell\in\mathfrak L(g)$ and at every position,
\[
|\mathsf J_2-c_2b^{-2}\mathcal K_6|(f,h)\le\bigl[(C_{{\rm E},3}l^{6p_0}+C_{\rm E}b^{-1}d\,l^{8p_0}+\bar\varepsilon_4(g,d))\,b^{-3}d+C_{\rm lat}b^{-2}a_K/d\bigr]\eta ,
\]
and on $T_{\ell_g}$, under the assumptions of Proposition~\ref{C-350-8} besides ($g\le\gamma_{\rm C}$, $d\le d_{\rm A}(g)$), by that Proposition, the same for $S^T_{2;K,\ell_g,\beta}(f,h)$ with $C_{\rm A}(g)\vartheta^{-1}bd$ added to the coefficient of $b^{-3}d$. The member $\bar\varepsilon_4(g,d)\,b^{-3}d\,\eta$ holds the two finite-size terms of the free value, $\frac32C_{\rm tm}b^{-2}[(d/\ell)^3+(d/D_K)^3]\eta$, of the torus and of the constraint at the top level $D_K=1$, which the proof of Proposition~\ref{C-350-9} bounds by $3C_{\rm tm}b^{-2}d^3\eta=3C_{\rm tm}bd^2\cdot b^{-3}d\,\eta$: they are at most $3C_{\rm tm}\,b^{-3}d\,\eta$ for $d^2\le1/b$, and of the order $bd^2$ relative to the first correction beyond.

The other members of $\bar\varepsilon_4(g,d)\,b^{-3}d\,\eta$ are at most $C_4(g)\,d^{1/2}\,b^{-3}d\,\eta$ (Proposition~\ref{C-350-9}), and $C_{\rm E}b^{-1}d\,l^{8p_0}\le C_{\rm E}c_{\mathfrak p}\,l^{6p_0}$ with $c_{\mathfrak p}=\sup_{0<t\le1}t\max\{1,\log\frac2t\}^{2p_0}$. So, beside the lattice term $C_{\rm lat}b^{-2}a_K/d$ and the two finite-size terms, the deviation from the free law is at most $C\,l^{6p_0}\,b^{-3}d\,\eta$, with a $C$ formed with $g$ and the constants (and with $\vartheta$ on $T_{\ell_g}$): of the order of the first correction, up to the factor $l^{6p_0}$; for $d^2\le1/b$ the two finite-size terms are of that order too. The coefficient of the first correction is fixed by the order six.
\end{remark}

\begin{lemma}[Gaussian measures conditioned on a convex set]\label{C-350-10}
Let $\nu$ be a centred Gaussian measure on $\mathbb R^I$, $I$ finite, with variances $\le\sigma_i^2$, and $Q\supset\prod_i(-c_i,c_i)$ open, bounded, convex and symmetric; put $t_i=c_i/\sigma_i$ and $\epsilon=\sum_ie^{-t_i^2/2}$. Let $F_1,\dots,F_r$ ($r=2,3$) be complex polynomials of orders $\le k_j$ with $|F_j|\le S_j$ on $Q$, and $k=\sum_jk_j$. If $t_i^2\ge\max\{4,2k\}$ for all $i$ and $2^k\epsilon\le1$, the joint cumulants of $F_1,\dots,F_r$ in $\nu(\cdot\,|\,Q)$ and in $\nu$ differ by at most $q_r(1+2^k)\epsilon\prod_jS_j$, with $q_2=5$, $q_3=37$.
\end{lemma}

\begin{proof}[Proof sketch]
For $z\notin Q$ let $s=\max_i|z_i|/c_i\ge1$. The segment from $-z/s$ to $z/s$ lies in $\bar Q$, so $\tau\mapsto F_j(\tau z/s)$ is a polynomial of order $\le k_j$ bounded by $S_j$ on $[-1,1]$, and the extremal property of Chebyshev's polynomials gives $|F_j(z)|\le S_j(2s)^{k_j}$. On $Q^{\rm c}$, $s^k\le\sum_i(|z_i|/c_i)^k\mathbf 1\{|z_i|\ge c_i\}$, and for a centred Gaussian variable of variance $\le\sigma_i^2$, $\mathbb E[(|z_i|/c_i)^k;|z_i|\ge c_i]\le e^{-t_i^2/2}$ once $t_i^2\ge\max\{4,2k\}$; so
\[
\nu(Q^{\rm c})\le\epsilon,\qquad\int_{Q^{\rm c}}\prod_j|F_j|\,d\nu\le2^k\epsilon\prod_jS_j .
\]
Let $F$ be a product of some of the $F_j$, $S$ the product of their $S_j$, and $m$, $m^{\rm r}$ its moments in $\nu$ and in $\nu(\cdot\,|\,Q)$. Then $m^{\rm r}-m=m^{\rm r}\nu(Q^{\rm c})-\int_{Q^{\rm c}}F\,d\nu$ and $|m^{\rm r}|\le S$, so $|m^{\rm r}-m|\le(1+2^k)\epsilon S$; and $|m|\le S+2^k\epsilon S\le2S$. A joint cumulant of order $r$ is a sum, over the partitions into $p$ blocks, of $(-1)^{p-1}(p-1)!$ times products of $p$ moments; telescoping each product, with the factors bounded by $S$ and by $2S$, bounds its change by $p\,2^{p-1}(1+2^k)\epsilon\prod_jS_j$. So $q_r=\sum_p\mathrm S(r,p)\,p!\,2^{p-1}$, $\mathrm S$ the Stirling numbers of the second kind: $q_2=1+4$, $q_3=1+12+24$.
\end{proof}

\begin{lemma}[The coefficients of a logarithm]\label{C-350-22}
Let $G$ be holomorphic on $|\lambda|<1$ with $G(0)=0$ and $\operatorname{Re}G\le\mathsf E$. Then every Taylor coefficient $G_{[m]}$, $m\ge1$, has $|G_{[m]}|\le2\mathsf E$. Consequently, if $F$ is holomorphic on $|\lambda|<1$, $F(0)=1$, $F\ne0$, $|F|\le e^{\mathsf E}$, and $F=1+\lambda\mathsf R_1+\lambda^2\mathsf R_2+O(\lambda^3)$, then $W_1=-\mathsf R_1$ and $W_2=\frac12\mathsf R_1^2-\mathsf R_2$ have $|W_1|,|W_2|\le2\mathsf E$. The condition $F\ne0$ cannot be dropped.
\end{lemma}

\begin{proof}[Proof sketch]
For $0<\rho<1$ and $m\ge1$, $G_{[m]}\rho^m=-\pi^{-1}\int_0^{2\pi}(\mathsf E-\operatorname{Re}G(\rho e^{i\phi}))e^{-im\phi}\,d\phi$; the factor $\mathsf E-\operatorname{Re}G$ is $\ge0$ and has the mean $\mathsf E$ over the circle, as $\operatorname{Re}G(0)=0$; let $\rho\to1$. For $F$, the function $G=\log F$ is defined on the disc, as $F$ has no zero there; $\operatorname{Re}G=\log|F|\le\mathsf E$ and $G=-\lambda W_1-\lambda^2W_2+O(\lambda^3)$. For $F=1+\mathsf c\lambda$ with $\mathsf c\ge3$, which has a zero in the disc, $|\mathsf R_1|=\mathsf c>2\log(1+\mathsf c)=2\mathsf E$.
\end{proof}

\begin{proposition}[The order six at one level]\label{C-350-11}
Under the assumptions of Proposition~\ref{C-350-6}, let $2\le n\le K$ and $\mathsf{Bx}=\prod_{0\le i<n}\{B_i:|B_i(b)|<\mathfrak b(x_i)\ \forall b\}$, a product of Euclidean balls in $\mathfrak g$, one for each non-pivot bond of the heights $0\le i<n$.
\begin{enumerate}
\item[(i)] There are functions $\mathsf O^\lambda_1$, $\mathsf O^\lambda_2$, $\mathsf D^\lambda$ on $\mathsf{Bx}$, continuous and bounded on $\{|\lambda|\le1\}\times\mathsf{Bx}$ and holomorphic in $|\lambda|<1$, with $\mathsf O^\lambda_a(B)=\mathsf o_a(\lambda B)$ for a function $\mathsf o_a$ holomorphic near $0$, $\mathsf O^{-\lambda}_a(-B)=\mathsf O^\lambda_a(B)$, $\mathsf D^{-\lambda}(-B)=\mathsf D^\lambda(B)$, $\int_{\mathsf{Bx}}\mathsf D^\lambda dB=1$, $\mathsf D^0dB=\bigotimes_{0\le i<n}\nu^{\rm r}_i$, and
\[
\sum_{j\le n}\mathsf n^\lambda_j=\bigl\langle\mathsf O^\lambda_1;\mathsf O^\lambda_2\bigr\rangle_{\mathsf D^\lambda dB}\qquad(|\lambda|<1).
\]
Here $\mathsf D^\lambda$ is the product over the heights $0\le i<n$ of the normalised densities of the steps along the chain of the backgrounds, which reads $(\lambda,B)$ through $\lambda B$; for $0\le i<n$, the height $0$ included, the density of the step of the height $i$ at $\lambda=0$, at the trivial configuration, is proportional to $\exp[-\frac12(B_i,(\mathsf M_i\otimes1)B_i)]$ on the box of the height, with the real symmetric $\mathsf M_i\ge\frac34\gamma^{\rm tr}$ of the measures $\nu^\circ_i$ (Lemma~\ref{C-350-23}(iii)); and for real $0<\lambda\le1$, $\log(\mathsf D^\lambda/\mathsf D^0)$ is the sum of a term that does not depend on $B$, of terms $\lambda^{-2}F(\lambda,\lambda B)$ and of terms $F(\lambda,\lambda B)$, with functions $F(\lambda,\zeta)$ holomorphic near $(0,0)$, those of the first kind with $F=O(|\zeta|^2)$.

Uniformly on $\mathsf{Bx}$, $\mathsf O^\lambda_a=\lambda^2(X_{a,0}+\lambda X_{a,1}+\lambda^2X_{a,2})+O(\lambda^5)$ and $\mathsf D^\lambda=\mathsf D^0\exp[-\lambda W_1-\lambda^2W_2+O(\lambda^3)]$, where $X_{a,m}$ is homogeneous of order $m+2$, $X_{a,0}=q_a+q^{(0)}_a$ with the $q_a$ of the proof of Proposition~\ref{C-350-9}, formed with the signs $\mathsf{sg}_i$ of Lemma~\YMMref{124.3}{835} there, and a quadratic form $q^{(0)}_a$ that vanishes where $B_0=0$, $W_1$ is odd of order $\le3$, and $W_2$ is, up to the constant that the term without $B$ gives, even of order $\le4$. These are polynomials on the whole space, with Taylor coefficients at the trivial configuration formed with $L$ and the running couplings, so that the cumulants of \eqref{C-350-5} are formed in $\bigotimes_{0\le i<n}\nu^\circ_i$ as well.

Hence $\mathsf h^{(n)}_4=\langle X_{1,0};X_{2,0}\rangle$ and $\mathsf h^{(n)}_6=\gamma^{\rm r}_2$, the cumulants being taken in $\bigotimes_{0\le i<n}\nu^{\rm r}_i$.
\item[(ii)] If $\gamma^{\rm tr}\mathfrak b(x_{n-1})^2\ge80$ and $10240N^3e^{-\gamma^{\rm tr}\mathfrak b(x_{n-1})^2/8}\le1$, then
\[
\beta^3a_K^2|\gamma^{\rm r}_2-\gamma^\circ_2|\le\mathcal U(x_{n-1},\mathsf N)\|w_1\|_\infty\|w_2\|_\infty,\qquad\mathcal U(x,\mathsf N)=C_{\rm X}\mathsf N^{36}\exp\bigl[-\tfrac18\gamma^{\rm tr}A_1^2(\log x)^{2p_0}\bigr],
\]
with $C_{\rm X}$ determined by the constants.
\end{enumerate}
\end{proposition}

\begin{proof}[Proof sketch]
(i)

\emph{The chain of the backgrounds.} For $|\lambda|\le1$ and $B\in\mathsf{Bx}$ put $v^\lambda_n:=0$ and, downwards for $i=n-1,\dots,1$, let $v^\lambda_i$ be the lowered point (equation~\YMMref{(2286)}{833}) that the step $i+1$ forms of the point $v^\lambda_{i+1}$ of the level $i+1$ and of the fluctuation $\lambda x_i^{-1/2}B_i$. Its first member $A^\lambda_i$ is a $1$-form on the bare bonds, and $\tilde U^\lambda_i:=\exp(iL^{-i}A^\lambda_i)$, a configuration on the bare bonds, is the \emph{background of the level $i$}. The maps of the step are holomorphic on their domain, at real and at complex points (the holomorphy following equation~\YMMref{(2286)}{833}); by their bounds (equation~\YMMref{(2287)}{834} and its domain verification), as $\sup|\lambda x_i^{-1/2}B_i|\le|\lambda|\delta(x_i)$ on the box, $\|v^\lambda_i\|\le|\lambda|r_i$ at every height, with the radii $r_n=0$, $r_i=L^{-1}r_{i+1}+\frac{11}{10}c_{\mathfrak h}\delta(x_i)\le\frac85c_{\mathfrak h}\delta(x_i)$ of equation~\YMMref{(2291)}{834}, which keep the chain in the domain of the maps, as $\delta(x_{K-1})$ lies below the threshold of the chain classes (Proposition~\ref{C-350-6}(iv)).

So the chain is formed for $|\lambda|\le1$, it reads $(\lambda,B)$ through $\lambda B$, and for real $\lambda\in[0,1]$ it stays in the chain classes (Proposition~\YMMref{124.2}{834}, equation~\YMMref{(2288)}{834}). At complex $\lambda$ the background has values in $\mathrm{SL}(2,\mathbb C)$, and $a_p$ is read there as $1-\frac12\operatorname{tr}$ of the product of the variables around $p$; on $\mathrm{SU}(2)$ this is $a_p$. The observable is $\mathsf O^\lambda_a(B):=a_K^{-1}\sum_pw_a(p)\,a_p(\tilde U^\lambda_1)+\mathsf w^\lambda_a(B_0)$: its first member is $\mathcal O(w_a)$ at the background of the level $1$, which carries the fields of the heights $1\le i<n$, and its second, the source term of the height $0$, is the sum over the domains of the level $0$ of the action term of the source $w_a$ at the point $v^\lambda_1$ and the field $\lambda x_0^{-1/2}B_0$; the field of the height $0$ enters $\mathsf O^\lambda_a$ through this term alone.

The density is $\mathsf D^\lambda(B):=\prod_{0\le i<n}\rho^\lambda_i/Z^\lambda_i$, with $\rho^\lambda_i$ the density of the step of the height $i$ at the field scale $\lambda$ at the point $v^\lambda_{i+1}$, the exponential of the quadratic forms of the step in $B_i$ and of its potentials, and $Z^\lambda_i$ its integral over the box of the height.

\emph{The identity.} At real $0<\lambda\le1$ each step at a height $1\le i<n$ obeys the identity of one step at the field scale $\lambda$ of step (1) of the proof of Proposition~\ref{C-350-9}, whose case $\lambda=1$ uses the quotient argument of Proposition~\ref{prop:weighted-pair-tower}, and at the height $0$ the source enters linearly, through the action term of the sources. Iterating from the height $n-1$ down,
\[
\exp D^{\lambda,\mu}_n=\int_{\mathsf{Bx}}\exp\bigl[\mu_1\mathsf O^\lambda_1+\mu_2\mathsf O^\lambda_2\bigr]\mathsf D^\lambda\,dB,
\]
$D^{\lambda,\mu}_n$ the exponent of the progression at the trivial configuration with the sources $\mu_1w_1+\mu_2w_2$. Both sides are holomorphic in $\mu$ near $0$, and holomorphic in $|\lambda|<1$, as the maps of the chain are holomorphic at complex points and the normalisations $Z^\lambda_i$ have no zero (Lemma~\ref{C-350-23}(ii)); they agree on $(0,1]$. The mixed derivative in $\mu$ at $0$ of the logarithms is the covariance.

\emph{At $\lambda=0$.} The chain is trivial and the potentials vanish. At the heights $1\le i<n$ the density of the step is $\exp[-\frac12(B_i,((\tau_i\mathsf A^{\rm cv}_i+(1-\tau_i)\hat{\mathsf G})\otimes1)B_i)]=\exp[-\frac12(B_i,(\mathsf M_i\otimes1)B_i)]$, and at the height $0$, where $\tau_0=1$, it is $\exp[-\frac12(B_0,(\mathsf A^{\rm cv}_0\otimes1)B_0)]$, $\mathsf M_0=\mathsf A^{\rm cv}_0$; the sources of these two identities are given with Lemma~\ref{C-350-23}(iii). So $\mathsf D^0dB=\bigotimes_{0\le i<n}\nu^{\rm r}_i$.

\emph{The orders.} The orders $0$ and $1$ of $\mathsf O^\lambda_a$ vanish, the order $1$ by the constant gauge rotations, as in Proposition~\ref{C-350-7}. As $\mathsf O^\lambda_a=\mathsf o_a(\lambda B)$ with $\mathsf o_a$ holomorphic near $0$, the coefficient of $\lambda^{m+2}$ is a homogeneous polynomial of order $m+2$. For real $0<\lambda\le1$ the integral identity of the step at a complex field scale (Proposition~\ref{IV-71-37}), read at the points of the chain, writes each $Z^\lambda_i$ as $\det^{1/2}$ of the covariance of the step times the exponential of a sum of local terms of the step; in the product of the quotients the families of the progression drop out, by the sum rule of the chain classes (Proposition~\YMMref{124.2}{834}, equation~\YMMref{(2289)}{834}), which telescopes over the heights, and $\mathsf D^\lambda(B)$ is a number $c_n(\lambda)$, which does not depend on $B$, times the exponential of a sum of terms $\lambda^{-2}F(\lambda,\lambda B)$ with $F=O(|\zeta|^2)$, namely the quadratic forms of the steps, $\lambda^{-2}$ times forms in $\lambda B_i$ whose coefficients read $\lambda B$ through the chain, and the potentials of the marks of the gauge fixing and of the action, $x_i\lambda^{-2}$ times functions of $\lambda x_i^{-1/2}B_i$, and of terms $F(\lambda,\lambda B)$, namely the other potentials, the logarithms of the determinants of the covariances along the chain, the localised Gaussian families of the steps and the counterterms at the backgrounds; in all of them $F(\lambda,\zeta)$ is holomorphic near $(0,0)$.

This is the form of $\log(\mathsf D^\lambda/\mathsf D^0)$ stated in (i), the term without $B$ being $\log(c_n(\lambda)/c_n(0))$. With $F=\sum_{k,\alpha}f_{k\alpha}\lambda^k\zeta^\alpha$ the coefficient of $\lambda^m$ of $\lambda^{-2}F(\lambda,\lambda B)$ is $\sum_{k+|\alpha|=m+2}f_{k\alpha}B^\alpha$, a polynomial of order $\le m+2$, and that of $F(\lambda,\lambda B)$ one of order $\le m$; this gives the orders of $W_1$ and $W_2$, and the identity between holomorphic functions of $\lambda$ passes from $(0,1]$ to a disc about $0$. Odd functions have the mean $0$ in $\mathsf D^0dB$, the mass $1$ removes the mean of the order $2$ of $\mathsf D^\lambda/\mathsf D^0$, and the coefficient of $\lambda^2$ in $\lambda^{-4}\langle\mathsf O^\lambda_1;\mathsf O^\lambda_2\rangle$ is \eqref{C-350-5}; at $\lambda^0$ it is $\langle X_{1,0};X_{2,0}\rangle$.

By Proposition~\ref{C-350-7}(i), $\lambda^{-4}\sum_{j\le n}\mathsf n^\lambda_j=\mathsf h^{(n)}_4+\lambda^2\mathsf h^{(n)}_6+O(\lambda^4)$.

(ii) By Lemma~\ref{C-350-23}, on $\mathsf{Bx}$ the $X_{a,m}$ are at most $C_{\rm S}\mathsf N^8\|w_a\|_\infty$ and $W_1$, $W_2$ at most $2C_{\rm S}\mathsf N^7$; a cumulant of order $\ge2$ does not see the constant in $W_2$. By (iii) of that Lemma the variances of $\nu^\circ_i$ are at most $\frac4{3\gamma^{\rm tr}}$; $\mathsf{Bx}$ holds the cube of the half-sides $\mathfrak b(x_{n-1})/\sqrt3$ in the components, as $x_i\ge x_{n-1}$ for $i<n$; so $t_i^2\ge\frac14\gamma^{\rm tr}\mathfrak b(x_{n-1})^2\ge20$ and $\epsilon\le10N^3e^{-\gamma^{\rm tr}\mathfrak b(x_{n-1})^2/8}$. Lemma~\ref{C-350-10} at the seven cumulants, whose entries have the total order $k\le10$, and $\beta^3a_K^2\le\mathsf N^3$ give (ii).
\end{proof}

\begin{lemma}[The suprema on the boxes]\label{C-350-23}
In the setting of Proposition~\ref{C-350-11}(i) there is $C_{\rm S}\ge1$, determined by the constants, such that for $|\lambda|<1$ and $B\in\mathsf{Bx}$:
\begin{enumerate}
\item[(i)] $|\mathsf O^\lambda_a(B)|\le C_{\rm S}\mathsf N^8\|w_a\|_\infty$ for $a=1,2$;
\item[(ii)] $\mathsf D^\lambda(B)\ne0$ and $|\mathsf D^\lambda(B)|\le e^{C_{\rm S}\mathsf N^7}|\mathsf D^0(B)|$;
\item[(iii)] for $0\le i<n$, the height $0$ included, the density of the step of the height $i$ at $\lambda=0$, at the trivial configuration, is proportional to $\exp[-\frac12(B_i,(\mathsf M_i\otimes1)B_i)]$ on the box of the height: its Gaussian part is $\nu^\circ_i$, the centred Gaussian measure with the real symmetric precision $\mathsf M_i\otimes1\ge\frac34\gamma^{\rm tr}$, and its normalised restriction to the box is $\nu^{\rm r}_i$, a probability measure; at the height $0$, $\tau_0=1$ and $\mathsf M_0=\mathsf A^{\rm cv}_0$.
\end{enumerate}
Consequently, on $\mathsf{Bx}$, $|X_{a,m}|\le C_{\rm S}\mathsf N^8\|w_a\|_\infty$ for $m=0,1,2$, and $|W_1|,|W_2|\le2C_{\rm S}\mathsf N^7$.
\end{lemma}

\begin{proof}[Proof sketch]
(i) The first member of $\mathsf O^\lambda_a$ is $a_K^{-1}\sum_pw_a(p)a_p(\tilde U^\lambda_1)$, at the background $\tilde U^\lambda_1$ of the level $1$ of the chain of the proof of Proposition~\ref{C-350-11}(i), a configuration on the bare bonds whose values stay near $\mathbf 1$ in $\mathrm{SL}(2,\mathbb C)$ by the radii of the chain, with at most $3N^3$ plaquettes and $a_K^{-1}\le N$; the second, the source term of the height $0$, is a sum over at most $8N^3$ domains of members of the family of the source, each bounded on the box by the size bounds of the step at a complex field scale (Proposition~\ref{IV-71-37}) and by $\mathfrak b(\beta)^2\le C\beta$.

(ii) $\mathsf D^\lambda=\prod_{0\le i<n}\rho^\lambda_i/Z^\lambda_i$ along the chain at the complex field scale $\lambda$, a product over the at most $K\le N$ heights. Each $\rho^\lambda_i$ is the exponential of a finite sum, so $\rho^\lambda_i\ne0$; its exponent is a sum over at most $8N^3$ places of quadratic forms in the field, whose coefficients are holomorphic along the chain, and of potentials, and on the box each of them is bounded by the size bounds of the step at a complex field scale (Proposition~\ref{IV-71-37}) and by $\mathfrak b(x_i)^2\le C\beta$, the radii of the chain keeping the points of all the heights in the discs where these bounds hold.

By the integral identity of the step at a complex field scale, $Z^\lambda_i$ is $\det^{1/2}$ of the covariance of the step times the exponential of a sum of local terms of the step, each bounded by the size bounds of the step; so $Z^\lambda_i\ne0$, and $|Z^\lambda_i|$ is bounded below by a determinant bound for its Gaussian part. Hence $\mathsf D^\lambda\ne0$, and, with at most $CN^3$ terms a height, each at most a number determined by the constants times $\mathfrak b(\beta)^2\le C\beta\le C\mathsf N$, this gives the bound with $\mathsf N^7$. (iii) For a height $1\le i<n$, at $\lambda=0$ the chain is trivial and the potentials vanish, and the density of the step at the trivial triple is $\exp[-\frac{\tau_i}2(B,\mathbf C_\ast^{-1}B)-\frac{1-\tau_i}2\sum_\square(B,\hat{\mathsf G}_\square B)]$.

Here the whole covariance of the step at the trivial triple is $\mathbf C_\ast=(\mathsf A^{\rm cv}_i)^{-1}\otimes1$ (equation~\YMMref{(2297)}{836}, from Proposition~\YMMref{120.3}{819} at the place of the height $i$), and the sum over the cubes of the local gauge-fixing forms is $\hat{\mathsf G}\otimes1$ (equation~\YMMref{(828)}{395} at the trivial block field, with Proposition~\YMMref{120.3}{819}); so the density is $\exp[-\frac12(B,((\tau_i\mathsf A^{\rm cv}_i+(1-\tau_i)\hat{\mathsf G})\otimes1)B)]=\exp[-\frac12(B,(\mathsf M_i\otimes1)B)]$, as $\mathsf A^{\rm cv}_i=\mathsf B^{\rm cv}_i+\hat{\mathsf G}$. At the height $0$, $\tau_0=1$, because $\xi_0\equiv x_0$, and the gauge-fixing term drops out.

The provision of the step $1$ is the canonical provision of the step at $(x_0;\mathfrak a^{\rm abs}_0)$, the quadruple $(0,x_0,\mathfrak a^{\rm abs}_0,\mathfrak w)$ being a place by the setting (J) of Definition~\ref{def:weighted-history-setting} at $k=K-1$ (Proposition~\ref{C-350-6}(iv)); by Proposition~\YMMref{120.3}{819} and Proposition~\YMMref{120.6}{821} at $k=0$, its whole covariance at the trivial triple is the scalar operator with the real matrix $(\mathsf A^{\rm cv}_0)^{-1}$. So the density of the height $0$ is $\exp[-\frac12(B_0,(\mathsf A^{\rm cv}_0\otimes1)B_0)]$, and $\mathsf M_0=\mathsf A^{\rm cv}_0$. At every height $\mathsf B^{\rm cv}_i$ and $\hat{\mathsf G}$ are real symmetric matrices $\ge0$ (equation~\YMMref{(2231)}{819}), $\mathsf A^{\rm cv}_i\ge\gamma^{\rm tr}$ (Proposition~\YMMref{120.3}{819}) and $\frac34\le\tau_i\le\frac43$ (Proposition~\ref{C-350-6}(ii)); so $\mathsf M_i$ is real and symmetric, $\mathsf M_i\ge\min\{\tau_i,1\}\mathsf A^{\rm cv}_i\ge\frac34\gamma^{\rm tr}$, the Gaussian part of the density is $\nu^\circ_i$, and $\nu^{\rm r}_i$ is a probability measure.

The consequences: Cauchy's estimate on $|\lambda|<1$ for $\lambda\mapsto\mathsf o_a(\lambda B)$, whose coefficient of $\lambda^{m+2}$ is $X_{a,m}(B)$, and Lemma~\ref{C-350-22} for $\lambda\mapsto\mathsf D^\lambda(B)/\mathsf D^0(B)$.
\end{proof}

\begin{proposition}[The composed chart; the lattice number of a flat top]\label{C-350-12}
Let $N=\ell L^K$, $\ell=2L^m$, and $1\le t\le K$ with $2L^{m+K-t+1}\ge C_LM$; the levels are those of the lattice $(\ell,K)$, $\mathcal V_j$ is the space of the configurations of the level $j$, and $\mathsf Y$ the space of the $y=(\mathcal A_j)_{j<t}$, $\mathcal A_j$ a real $\mathfrak g$-valued $1$-form on the non-pivot bonds of the level $j$ (Section~\YMMref{4}{16} and the pivot definition preceding Lemma~\YMMref{40.1}{297}).
\begin{enumerate}
\item[(i)] There are an open neighbourhood $\mathsf O_t$ of $\mathbf 1$ in $\mathcal V_t$, $\rho_2\in(0,1]$ and $C_{\rm H}>0$ such that, with $\mathsf Y_{\rho_2}=\{y\in\mathsf Y:|\mathcal A_j|<\rho_2,\ j<t\}$: the chain $W_t=W$, $W_j=V_{W_{j+1}}[\mathcal A_j]$ of the charts of Section~\ref{sec:sf-chart-interface} is defined on $\mathsf O_t\times\mathsf Y_{\rho_2}$, with $\mathcal M(W_j)=W_{j+1}$; $\Psi(W,y)=W_0$ is a real analytic diffeomorphism of $\mathsf O_t\times\mathsf Y_{\rho_2}$ onto an open set $\mathsf U\ni\mathbf 1$; and there is a real continuous $\theta$, real analytic near $(\mathbf 1,0)$, with $\int H\,dU=C_{\rm H}\int_{\mathsf O_t}dW\int_{\mathsf Y_{\rho_2}}dy\,e^{-\theta}H\circ\Psi$ for every bounded Borel $H$ that vanishes off $\mathsf U$.
\item[(ii)] For $\kappa\in(0,\infty)^t$ let $F_\kappa=A+\sum_{j<t}\kappa_jG_j(\mathcal M^j\cdot)$, $G_j$ the gauge-fixing function \YMMref{(54)}{27} at the level $j$, and $\phi_\kappa=F_\kappa\circ\Psi(\mathbf 1,\cdot)$. Then $\phi_\kappa\ge0$, $\phi_\kappa(0)=0$, $\nabla\phi_\kappa(0)=0$, and the Hessian of $\phi_\kappa$ at $0$ is block diagonal for the decomposition of $\mathsf Y$ by the levels and positive definite.
\item[(iii)] The numbers $\mathsf g_{m'}[N,t](p,q)=\gamma^{\rm L}_{m'}[\phi_\kappa,\theta(\mathbf 1,\cdot),a_p\circ\Psi(\mathbf 1,\cdot),a_q\circ\Psi(\mathbf 1,\cdot)]$, $m'=0,2$, do not depend on $\kappa$, and they are the same in every adapted chart. An \emph{adapted chart} of $(\mathcal V_0,dU,\operatorname{Lg}\mathcal M^t)$, $\operatorname{Lg}$ being the Lie algebra logarithm of the intrinsic logarithm on $G\setminus\{-1\}$, is $(\mathsf E',\Psi',\theta',C',\varrho')$ with: $\mathsf E'$ an open neighbourhood of $(0,0)$ in $\mathsf W_t\times\mathbb R^{\dim\mathsf Y}$, $\mathsf W_t$ the space of the $\mathfrak g$-valued $1$-forms $w$ on the bonds of the level $t$; $\Psi'$ a $C^\infty$ diffeomorphism of $\mathsf E'$ onto an open subset of the domain of $\operatorname{Lg}\mathcal M^t$, with $\Psi'(0,0)=\mathbf 1$ and $\operatorname{Lg}\mathcal M^t\Psi'(w,y)=w$; $\theta'$ a real continuous function on $\mathsf E'$, with $\theta'(0,\cdot)$ of class $C^\infty$; $C'>0$; $\varrho'>0$ a continuous function of $w$ alone; and $\int H\,dU=C'\int_{\mathsf E'}\varrho'(w)\,e^{-\theta'(w,y)}H(\Psi'(w,y))\,dw\,dy$ for every bounded Borel $H$ that vanishes off $\Psi'(\mathsf E')$.

In it the numbers are $\gamma^{\rm L}_{m'}[F_\kappa\circ\Psi'(0,\cdot),\theta'(0,\cdot),a_p\circ\Psi'(0,\cdot),a_q\circ\Psi'(0,\cdot)]$. They depend on $N$, $t$, $L$, $p$, $q$ alone.
\item[(iv)] One may take $\mathsf O_t=\{e^{iw}:|w|<\rho_2\}$ for every $\rho_2\in(0,\rho_\star]$, with a $\rho_\star\in(0,1]$. Then $W_j$, $V^{\rm bf}_j(W_{j+1})$, $G_j(W_j)$, $\theta$ and $a_p\circ\Psi$ are real analytic in $(w,y)$ on $\mathsf O_t\times\mathsf Y_{\rho_2}$, $W_j(\mathbf 1,0)=\mathbf 1$, $0<e^{-\theta}\le\prod_{j<t}(27L^9)^{\#\Gamma_j}$, and $C_{\rm H}=\prod_{j<t}(2\pi^2)^{-\#\mathfrak X^+_j}L^{-3\#\Gamma_j}$, where $\mathfrak X^+_j$ are the positively oriented non-pivot bonds of the level $j$ and $\Gamma_j$ the positively oriented bonds of the level $j+1$. One may also take $\mathsf O_t=\{W\in\mathcal V_t:\max_p|\partial W(p)-\mathbf 1|<\varepsilon_*\}$, $p$ running over the plaquettes of the level $t$, and $\rho_2=\rho_*$, with numbers $\varepsilon_*,\rho_*\in(0,1]$ determined by $L$ and the constants and independent of $t$, $K$ and $m$; (i)--(iii) hold for this choice as well.
\end{enumerate}
\end{proposition}

\begin{proof}[Proof sketch]
(i), (iv) One real number $\mathsf t_*>1$ serves at all levels: with it the constants, the lattice of each step and $\mathsf t_*$ are data under (Im) of Definition~\ref{def:sf-chart-data}. For such data Lemma~\ref{lem:sf-chart-integral} gives, on an open set of block fields $W$ and for $|\mathcal A|$ below the bound that belongs to $\mathsf t_*$: the averaged configuration $V^{\rm bf}_j(W)$ is defined and $\mathcal M(V^{\rm bf}_j(W))=W$; $\mathcal M(V_W[\mathcal A])=W$; and $\operatorname{Lg}(V_W[\mathcal A](b)V^{\rm bf}_j(W)(b)^{-1})=\mathcal A(b)$ on the non-pivot bonds. Lemma~\ref{lem:sf-chart-integral} gives the one-step integral formula, with the constant $(2\pi^2)^{-\#\mathfrak X^+_j}L^{-3\#\Gamma_j}$, a Jacobian factor in $(0,(27/8)^{\#\Gamma_j}]$ and a normalisation in $(0,(8L^9)^{\#\Gamma_j}]$.

Neighbourhoods of $\mathbf 1$ chosen downwards, by continuity in $(W,\mathcal A)$, keep the chain where these statements apply; by induction over the levels the integral formula holds for every Borel subset of the natural domain of the chain, with $C_{\rm H}$ the product of the constants. The map
\[
U\mapsto\Bigl(\mathcal M^tU,\ \bigl(\operatorname{Lg}(\mathcal M^jU(b)\,V^{\rm bf}_j(\mathcal M^{j+1}U)(b)^{-1})\bigr)_{j<t}\Bigr)
\]
is a real analytic left inverse of $\Psi$ with injective differential, and $\dim\mathcal V_t+\dim\mathsf Y=\dim\mathcal V_0$; the inverse function theorem gives (i) after shrinking, and (iv) with $W=e^{iw}$, $dW=\varrho_{\rm H}(w)\,dw$, $\varrho_{\rm H}>0$ the continuous density of the Haar measure of $\mathcal V_t$ in the variable $w$, on one set that carries (i)--(iii). For the second choice in (iv) one follows, level by level, a regularity number $\mathfrak r_j(W_j)$, which measures the plaquette variables of the averaged configuration of the level $j$ and is comparable with $\max_p|\partial W_j(p)-\mathbf 1|$: one step gives $\mathfrak r_j(W_j)\le L^{-2}\mathfrak r_{j+1}(W_{j+1})+C_{\mathfrak r}|\mathcal A_j|$, with $C_{\mathfrak r}$ determined by $L$ and the constants.

So a threshold $\mathfrak r_\flat$ for $\mathfrak r$, with $\varepsilon_*$ a fixed fraction of it and $\rho_*\le(1-L^{-2})\mathfrak r_\flat/C_{\mathfrak r}$, keeps every level in the sets where the statements of one step apply; nothing in this sees $t$, $K$ or $m$. (ii) $a_p\ge0$ and $G_j\ge0$; at $\mathbf 1$, $a_p(e^{iA})=\frac12|\mathrm dA(p)|^2+O(|A|^3)$, $G_j$ is half the sum of the squares of the linearised contour variables up to $O(|A|^3)$, and a fibre variable whose linearised curvature and contour variables vanish is $0$. At the linear order the composed chart is the chart of the free template: the bare curvature of the field of the level $j$ is $\mathsf T_j$ applied to it, with values in the slice $\mathcal C_{(j+1)}$ (Proposition~\YMMref{120.3}{819}), the slices are mutually orthogonal, and the linearised contour variables of the level $j$ are those of $\mathcal A_j$ alone; this makes the Hessian block diagonal.

(iii) By (i) and (iv), near $(0,0)$ the composed chart $(w,y)\mapsto\Psi(e^{iw},y)$, with $C_{\rm H}$, $\varrho_{\rm H}$ and $\theta(e^{iw},y)$, is an adapted chart: it is a real analytic diffeomorphism onto an open set, $\Psi(\mathbf 1,0)=\mathbf 1$ ($W_j(\mathbf 1,0)=\mathbf 1$), $\operatorname{Lg}\mathcal M^t\Psi(e^{iw},y)=w$, and $\theta$ is real analytic. Let $(\mathsf E',\Psi',\theta',C',\varrho')$ be an adapted chart, and $\mathsf E_1$ an open neighbourhood of $(0,0)$ that the composed chart maps into $\Psi'(\mathsf E')$. Then $\sigma:=\Psi'^{-1}\circ\Psi$ is a $C^\infty$ diffeomorphism of $\mathsf E_1$ onto an open set of the form $\sigma(w,y)=(w,s(w,y))$, both charts carrying $\operatorname{Lg}\mathcal M^t$ to the first variable, and $\sigma(0,0)=(0,0)$, as $\Psi(\mathbf 1,0)=\mathbf 1=\Psi'(0,0)$.

The two formulas for $dU$ give
\[
C_{\rm H}\,\varrho_{\rm H}(w)\,e^{-\theta(e^{iw},y)}=C'\varrho'(w)\,e^{-\theta'(\sigma(w,y))}\,|\det D_ys(w,y)|
\]
almost everywhere on $\mathsf E_1$, and, both sides being continuous, at every point of $\mathsf E_1$; in particular on the fibre $w=0$, a null set, where the numbers are read. So $y\mapsto s(0,y)$, a diffeomorphism near $0$ that fixes $0$, carries the phase $F_\kappa\circ\Psi(\mathbf 1,\cdot)$, the observables and the density $e^{-\theta(\mathbf 1,y)}dy$ into $F_\kappa\circ\Psi'(0,\cdot)$, the observables and $e^{-\theta'(0,y')}dy'$, up to the constant factor $C'\varrho'(0)/(C_{\rm H}\varrho_{\rm H}(0))$; the Taylor polynomials are taken at the minimum $0$ in both charts. The two coefficients of Laplace's method in either chart are those of the expansion of one quantity, the covariance of $a_p$ and $a_q$ under $e^{-\beta F_\kappa}$ on the fibre $\mathcal M^tU=\mathbf 1$ near $\mathbf 1$, and so they are the same in both.

The gauge transformations equal to $\mathbf 1$ on the top lattice preserve the fibre, $A$ and the $a_p$, and the normalised gauge-fixing factor (Definition~\YMMref{6.1}{27} and equations~\YMMref{(51)--(52)}{27}) integrates to one over their orbits at every strength (Lemma~\YMMref{5.3}{26} and equation~\YMMref{(55)}{28}); integrating over the orbits, the finest sites first, the covariance does not depend on $\kappa$ up to errors $e^{-c\beta}$, and neither do its coefficients. Near $(0,0)$ the composed chart is itself an adapted chart.
\end{proof}

\begin{proposition}[The identification of the order six]\label{C-350-13}
Under the assumptions of Proposition~\ref{C-350-11}, $\beta^3a_K^2\gamma^\circ_2=\mathsf T[N,n](w)$, the plaquettes being read in bare units on $\Lambda_N$; $\mathsf T[N,n](w)$ depends on $N$, $n$, $L$ and the lattice sources alone. The identity has no error term and holds at every size of the boxes.
\end{proposition}

\begin{proof}[Proof sketch]
By Proposition~\ref{C-350-6}(i), Proposition~\ref{C-350-12} applies at $t=n$ on the lattice of the datum. For real $0<\lambda\le1$ the exact formula of the steps (Proposition~\YMMref{123.1}{829}), read at the fields $\lambda B$, and the sum rules of the chain classes (Proposition~\YMMref{124.2}{834}, equation~\YMMref{(2289)}{834}) add the exponents of the steps of Proposition~\ref{C-350-11}(i) up to
\[
\mathsf D^\lambda(B)=c_n(\lambda)\exp\Bigl[-\lambda^{-2}\Bigl(x_0A(W_0)+\sum_{j<n}x_jG_j(W_j)\Bigr)-\theta\Bigr],
\]
where $c_n(\lambda)$ does not depend on $B$ and $(W_j)$ is the composed chart over the flat top $\mathbf 1$ of the level $n$ at the fields $\mathcal A_i=\lambda x_i^{-1/2}\check B_i$, $\check B_i=\operatorname{Ad}(u_{i+1})B_i$, with $u_{i+1}\in G$ a function of the fields of the heights $>i$ alone: the chain of Proposition~\ref{C-350-11} runs through configurations in another gauge than the composed chart. The observables become $a_K^{-1}\sum_pw_a(p)a_p(W_0)$ at the same fields. No $\tau_i$ and no counterterm $e_i$ is left, because
\[
\lambda^{-2}\xi_j(\lambda)L^j+\sum_{i\le j}L^ie_i=\lambda^{-2}x_0
\]
by the run equation; this is what $\xi_j$ is chosen for. So the density is the Wilson weight at $\beta\lambda^{-2}$ with the gauge-fixing strengths $\kappa_j=x_j/\beta$ on the fibre $\mathcal M^nU=\mathbf 1$. The triangular change $B\mapsto\check B$ preserves the volume and the boxes, which are Euclidean balls $|B_i(b)|<\mathfrak b(x_i)$ and are preserved by the rotations $\operatorname{Ad}(u)$; its linear part at $0$ is a constant rotation of the colours, so it preserves the Gaussian part $\bigotimes_{i<n}\nu^\circ_i$; and the two numbers of Laplace's method are kept by a diffeomorphism that carries phase, observables and density into each other.

So $\gamma^\circ_2$ is the second coefficient of Laplace's method \eqref{C-350-31} at the parameter $\beta\lambda^{-2}$, and $\beta^3a_K^2\gamma^\circ_2=\sum_{p,q}w_1(p)w_2(q)\gamma^{\rm L}_2$ in the chart of Proposition~\ref{C-350-12}, which does not depend on $\kappa$. The dilation $x\mapsto L^{K-n}x$ carries the lattice $(\ell,K)$ with its flat top at the level $n$ to the lattice $(\ell L^{K-n},n)$, with the same bare torus, plaquettes to plaquettes with their barycentres; so the number is $\mathsf g_2[N,n](p,q)$. The argument uses that the boxes are Euclidean balls; it would not apply to boxes that are cubes in the components.
\end{proof}

\begin{lemma}[An implicit function, Stokes's formula in a box, Wick's rule]\label{C-350-24}
\begin{enumerate}
\item[(i)] Let $\mathsf V$ be a finite-dimensional complex normed space, $0<\varsigma\le\frac1{120}$, and $F$ holomorphic from $\{\|y\|<\varsigma\}\times\{\lambda\in\mathfrak g^{\mathbb C}:|\lambda|<\varsigma\}$ to $\mathfrak g^{\mathbb C}$ with $|F|\le\frac1{40}$, $F(0,0)=0$ and $D_\lambda F(0,0)=1$. Then for $\|y\|<\varsigma^3$ there is exactly one $h(y)$ with $|h(y)|\le\varsigma^2$ and $F(y,h(y))=0$; $h$ is holomorphic, $|h(y)|\le\|y\|/(15\varsigma)$, and $h$ is $\mathfrak g$-valued at the real points if $F$ is.
\item[(ii)] Let $\mathsf B$ be the injective image in a bare torus of a box of $\mathbb Z^3$. For every closed bare path $\gamma$ in $\mathsf B$ there is an integer function $\mathsf k_\gamma$ on the plaquettes with all corners in $\mathsf B$, $|\mathsf k_\gamma|\le|\gamma|$ (the number of bonds of $\gamma$), with $a(\gamma)=\sum_p\mathsf k_\gamma(p)\,(\mathrm da)(p)$ for every $1$-cochain $a$.
\item[(iii)] Let $(\zeta_\iota)$ be a centred Gaussian family with covariance $C$, and $P_1,\dots,P_r$ homogeneous polynomials in it of orders $n_1,\dots,n_r$, $n=\sum n_a$, $\|P\|_1$ the sum of the moduli of the coefficients. For monomials the joint cumulant $\langle P_1;\dots;P_r\rangle$ is the sum, over the pairings of the factors that are connected (the graph on $\{1,\dots,r\}$ with an edge for every pair that joins two of the $P_a$ is connected), of the products of the covariances, and by multilinearity this extends to polynomials; the cumulant vanishes if $n$ is odd. With $\mathsf c\ge\sup|C|$ on the supports, $|\langle P_1;\dots;P_r\rangle|\le(n-1)!!\prod_a\|P_a\|_1\mathsf c^{n/2}$; if every connected pairing has a pair with $|C|\le\varepsilon$, the bound holds with $\mathsf c^{n/2-1}\varepsilon$ for $\mathsf c^{n/2}$; and for two covariances $C$, $C'$ bounded by $\mathsf c$ the two cumulants differ by at most $(n-1)!!\frac n2\prod_a\|P_a\|_1\mathsf c^{n/2-1}\sup|C-C'|$.
\end{enumerate}
\end{lemma}

\begin{proof}[Proof sketch]
(i) Put $g(y,\lambda)=\lambda-F(y,\lambda)$, so $|g|\le\frac1{40}+\varsigma\le\frac1{30}$, $g(0,0)=0$, $D_\lambda g(0,0)=0$. Cauchy's estimate and Schwarz's lemma give $|D_\lambda g(y,\lambda)|\le\frac2{15}\varsigma^{-2}\max\{\|y\|,|\lambda|\}$ for $\max\{\|y\|,|\lambda|\}\le\frac12\varsigma$, and $|g(y,0)|\le\|y\|/(30\varsigma)$. So for $\|y\|<\varsigma^3$ the map $\lambda\mapsto g(y,\lambda)$ sends $\{|\lambda|\le\varsigma^2\}$ into itself and contracts by $\frac2{15}$; its fixed point is $h(y)$, the uniform limit of holomorphic iterates, and $|h|\le|g(y,0)|+\frac2{15}|h|$. (ii) With the paths $P_z$ from a corner $o$ of the box to $z$ that run first in the direction $1$, then $2$, then $3$, the quantity $a(P_z)+a(e)-a(P_{z+e_\mu})$ of a bond $e=\langle z,z+e_\mu\rangle$ is the sum of $\pm(\mathrm da)$ over a strip of plaquettes of the box; along a closed path the terms $a(P_\cdot)$ telescope.

(iii) By Wick's theorem the moments are sums over all pairings; a pairing determines the partition of the factors into the components of its graph and is connected on each block, which identifies the sums over connected pairings with the joint cumulants. The bounds count the $(n-1)!!$ pairings and use $|\prod_{i\le m}c_i-\prod_{i\le m}c'_i|\le m\,\mathsf c^{m-1}\max|c_i-c'_i|$.
\end{proof}

\begin{proposition}[Two bare tori under one flat top]\label{C-350-14}
Let $N=2L^{\mathsf m}$, $N'=L^sN$ ($s\ge1$), $1\le t<\mathsf m$, $\pi:\Lambda_{N'}\to\Lambda_N$ the reduction, and
\[
\nu=N/L^t\ge\nu_{\rm W}=\max\Bigl\{192,\ \frac{48(3\log L+\log2)}{\delta_{\rm tp}L(L-1)},\ \frac{C_LM}L\Bigr\}.
\]
For plaquettes $p',q'$ of $\Lambda_{N'}$ whose base points lie within $N/8$ of each other (the base point of the plaquette $p_{\mu\nu}(x)$ being $x$, the distance being the sup distance of the bare torus, in bare units) and $m'=0,2$,
\[
|\mathsf g_{m'}[N,t](\pi p',\pi q')-\mathsf g_{m'}[N',t](p',q')|\le C'_{\rm W}N'^{237}e^{-\delta_{\rm tp}L\nu/48},
\]
with $C'_{\rm W}$ determined by $L$ and $C_{\rm tp}$. Hence, for real $w_1,w_2$ on the plaquettes of $\Lambda_N$ with supports in balls of radius $r\ge8$, all base points within $N/8$ of each other, and $\breve w_a$ equal to $w_a\circ\pi$ on one lift of a ball of radius $N/4$ that holds the supports and $0$ elsewhere,
\[
|\mathsf T[N,t](w)-\mathsf T[N',t](\breve w)|\le900r^6C'_{\rm W}N'^{237}e^{-\delta_{\rm tp}L\nu/48}\|w_1\|_\infty\|w_2\|_\infty .
\]
\end{proposition}

\begin{proof}[Proof sketch]
\emph{Signs.} Let $S\in\{N,N'\}$, and read the objects of the lattice $(S/L^t,t)$ in bare units on $\Lambda_S$; distances are those of the bare torus, $|\cdot|_\infty$ its sup distance, and $x(p)$ is the barycentre of a plaquette $p$, its base point being $x$ if $p=p_{\mu\nu}(x)$. For a bond $c$ of the top lattice $\mathbb T^{(t)}$, $B^t(c_\pm)$ are the blocks of order $t$ of its end points, its double block is their union, and its corridor ${\rm cor}(c)$ is the set of the bare bonds $b$ with $b_-\in B^t(c_-)$, $b_+\in B^t(c_+)$; $\mathsf 1_c$ is the bare $1$-cochain that is $1$ on ${\rm cor}(c)$, $-1$ on the reversed bonds and $0$ elsewhere. As $\nu\ge8$, every bare bond lies in one top block or in exactly one corridor, and no bond of the corridor of $c'$ has both end points in the double block of $c$ unless $c'\in\{c,\bar c\}$; the bonds of the double block of $c$ are those with both end points in it.

For a bare site $x\notin\mathbb T^{(t)}$, $j(x)=\max\{j:x\in\mathbb T^{(j)}\}<t$, its parent $\hat x\in\mathbb T^{(j(x)+1)}$ is the centre of the block of order $1$ that holds $x$, and ${\rm anc}(z)=\{z,\hat z,\hat{\hat z},\dots\}\setminus\mathbb T^{(t)}$; the sums over $x$ below run over the bare sites off $\mathbb T^{(t)}$. $\mathsf Q_k$ is the linear part of the $k$-fold block average, and $\xi_A(x)=(\mathsf Q_{j(x)}A)(\hat x,x)$ is the linearised contour variable of the site $x$. A $\mathfrak g$-valued $1$-form $y$ has the coordinates $y(\iota)$, $\iota=(b,\alpha)$, $b>0$ a bare bond and $\alpha$ a colour; $\iota+\gamma$ is the translate of $\iota$ by $\gamma$, and ${\rm dist}(\iota,\iota')=|b_--b'_-|_\infty$.

By Proposition~\ref{C-350-12}(iii) the numbers do not depend on $\kappa$; we take $\kappa\equiv1$.

\emph{A chart made for locality.} $\mathsf Q_k$ is a mean over bare paths: for every bond $c$ of $\mathbb T^{(k)}$ there is a probability measure on the bare paths from $c_-$ to $c_+$ in the double block of $c$, with at most $(4L)^k$ bonds, whose mean of $A(\gamma)$ is $(\mathsf Q_kA)(c)$; and $\mathsf Q_t\mathsf 1_c$ is $1$ at $c$ and $0$ at the other top bonds. $\mathcal M^t(U)(c)$ depends on $U$ on the double block of the top bond $c$ alone. So, for $y$ in $Y_S:=\ker\mathsf Q_t$ near $0$, the constraint $\mathcal M^t(e^{i\hat y})=\mathbf 1$ for $\hat y:=y+\sum_{c>0}h_c(y)\mathsf 1_c$ is solved one top bond at a time, with one element $h_c(y)\in\mathfrak g$ (three unknowns) for each: Lemma~\ref{C-350-24}(i) at $\varsigma=\varsigma_t:=(75L)^{-t}/(60L)$, applied to $(y,\lambda)\mapsto\operatorname{Lg}\mathcal M^t(e^{i(y+\lambda\mathsf 1_c)})(c)$, which is holomorphic and bounded by $\frac1{40}$ on the polydisc of the Lemma, vanishes at $(0,0)$ and has the derivative $1$ in $\lambda$ there, gives exactly one holomorphic $h_c$ with $|h_c(y)|\le\|y\|_\infty/(15\varsigma_t)$ for $\|y\|_\infty<\varsigma_t^3$, $\mathfrak g$-valued at real $y$; it reads $y$ on the double block of $c$ alone, $\hat y=y$ on the top blocks, and $\hat y=y+O(y^2)$ on $Y_S$.

The map $A\mapsto(\operatorname{Lg}\mathcal M^t(e^{iA}),\,A-\sum_{c>0}(\mathsf Q_tA)(c)\mathsf 1_c)$ is a real analytic diffeomorphism near $0$, as its differential at $0$ has the inverse $(w,y)\mapsto y+\sum_cw(c)\mathsf 1_c$; its inverse, followed by $A\mapsto e^{iA}$, is an adapted chart $\Psi'$ in the sense of Proposition~\ref{C-350-12}(iii), with $\varrho'=1$, a real analytic $\theta'$ and $\Psi'(0,y)=e^{i\hat y}$ for $y\in Y_S$; so both numbers may be computed in it. In it the action, the gauge terms, the Haar density (the exponential-coordinate Haar formula following equation~\YMMref{(911)}{420}) and the Jacobian, the product over the top bonds $c$ of the determinants of the derivatives of $\operatorname{Lg}\mathcal M^t(e^{i\cdot})(c)$ in the direction $\mathsf 1_c$, are sums of terms, each holomorphic and bounded by $1$ on the polydisc $\|y\|_\infty<\varsigma_t^3$ and depending on $y$ on at most $12L^{3t}$ bonds with end points in a set of diameter $\le2L^t$; for a plaquette $p$ these are the top blocks that hold a corner of $p$, within $L^t$ of its base point.

The reduction $\pi$ carries the lattices, blocks, contours, corridors and the block average of $\Lambda_{N'}$ to those of $\Lambda_N$ on every set of diameter $\le N/2$, so these are the same functions on both tori after the renaming of the variables by $\pi$.

\emph{The lines.} The gauge terms and the Haar terms of the chart read $y$ on the bonds, and not its curvature alone; so the lines of Wick's rule in this chart are the entries of the covariance $C_S$ of the chart variable $y\in Y_S$ under the centred Gaussian measure with the density proportional to $e^{-\frac12\mathsf H_S(y,y)}$, $\frac12\mathsf H_S(y,y)$ being the second Taylor polynomial at $0$ of the action with the gauge terms,
\[
\mathsf H_S(A,A')=(\mathrm dA,\mathrm dA')+\sum_x\langle\xi_A(x),\xi_{A'}(x)\rangle .
\]
Let $\mathcal X_S$ be the space of the $\mathfrak g$-valued functions on the bare sites that vanish on $\mathbb T^{(t)}$. Then $\mathrm d\mathcal X_S\subset Y_S$, as $\mathsf Q_t\mathrm d\chi=\mathrm d(\chi|_{\mathbb T^{(t)}})$, and $\chi\mapsto\xi_{\mathrm d\chi}=(\chi(x)-\chi(\hat x))_x$ is a bijection of $\mathcal X_S$ onto the functions on the sites off $\mathbb T^{(t)}$, with the inverse $\chi(x)=\sum_{x''\in{\rm anc}(x)}\xi(x'')$. So $Y_S$ is the direct sum of its axial part $Y^{\rm ax}_S=\{A\in Y_S:\xi_A=0\}$ and of $\mathrm d\mathcal X_S$, $y=y^{\rm ax}+\mathrm d\chi_y$, and
\[
\mathsf H_S(y,y)=\|\mathrm dy^{\rm ax}\|^2+\sum_x|\xi_y(x)|^2 ;
\]
$\mathrm d$ maps $Y^{\rm ax}_S$ one to one onto the curvatures whose plaquette averages of the order $t$ vanish, and $\mathsf H_S$ is positive definite on $Y_S$. In the Gaussian measure $y^{\rm ax}$ and the $\xi(x)$ are independent: the $\xi(x)$, one for each site off the top lattice, are independent with the variance $1$ in each colour, $(\mathrm d\chi)(b)=\sum_x\sigma(b,x)\xi(x)$ with $\sigma(b,x)=\mathbf 1[x\in{\rm anc}(b_+)]-\mathbf 1[x\in{\rm anc}(b_-)]$, and $\mathrm dy^{\rm ax}$ has the covariance $\Pi^{\rm fl}_t=\Pi^{\rm tp}_0-\Pi^{\rm tp}_t$ in each colour. On $Y^{\rm ax}_S$ the means of $y$ over two families of bare paths vanish: over the paths from the centre of the top block of a site $x$ to $x$, inside that block and with at most $(4L)^t$ bonds, whose mean of $y(\gamma)$ is $\sum_{x''\in{\rm anc}(x)}\xi_y(x'')=0$ (the contour averages of the levels composed from the top down), and over the paths of $\mathsf Q_t$ along a top bond, as $y\in\ker\mathsf Q_t$.

So $y(b)$ is the mean of the values of $y$ on closed bare paths of at most $4(4L)^t$ bonds in the box of $b$, the top block that holds $b$, or the double block of $c$ if $b\in{\rm cor}(c)$, and by Lemma~\ref{C-350-24}(ii), the box being the injective image of a box of $\mathbb Z^3$, $y(b)=\sum_p\tau(b,p)\,(\mathrm dy)(p)$ with real Stokes coefficients, $|\tau(b,p)|\le4(4L)^t$, $\tau(b,p)=0$ unless all the corners of $p$ lie in the box of $b$. Hence $C_S((b,\alpha),(b',\alpha'))=\delta_{\alpha\alpha'}[C^{\rm ax}_S+C^{\rm g}_S](b,b')$,
\[
C^{\rm ax}_S(b,b')=\sum_{p,p'}\tau(b,p)\tau(b',p')\,\Pi^{\rm fl}_t(p,p'),\qquad C^{\rm g}_S(b,b')=\sum_x\sigma(b,x)\sigma(b',x),
\]
and $C^{\rm g}_S(b,b')=0$ unless an end point of $b$ and one of $b'$ lie in one top block. A box holds at most $6L^{3t}$ plaquettes, $|\Pi^{\rm fl}_t|\le1$ (the orthogonal projections in~\YMMref{(2226)}{818}), and $\sigma(b,\cdot)\ne0$ at no more than $2t$ sites; so $\sup|C_S|\le\mathsf c_\tau^2+2t$, $\mathsf c_\tau=24L^{3t}(4L)^t$. By Proposition~\YMMref{120.6}{821} on the lattice $(S/L^t,t)$ at $s=0$, uniformly in the volume,
\[
|\Pi^{\rm fl}_t(p,q)|\le2L^3C_{\rm tp}L^{-3t}e^{-\delta_{\rm tp}L|x(p)-x(q)|_\infty/L^t}\qquad\text{if }\delta_{\rm tp}L(L-1)|x(p)-x(q)|_\infty\ge(3\log L+\log2)L^t ,
\]
the barycentres being taken in bare units. The barycentre of a plaquette in the box of $b$ lies within $2L^t$ of $b_-$; so if ${\rm dist}(\iota,\iota')\ge R\ge8L^t$ and $\delta_{\rm tp}L(L-1)(R-4L^t)\ge(3\log L+\log2)L^t$, then $C^{\rm g}_S(b,b')=0$ and
\[
|C_S(\iota,\iota')|\le\mathsf c_\tau^2\,2L^{3-3t}C_{\rm tp}\exp\bigl[-\delta_{\rm tp}L(R-4L^t)/L^t\bigr].
\]
At $R=N/24$ these conditions hold, and the exponent is at least $\delta_{\rm tp}L\nu/48$: $R=\frac\nu{24}L^t\ge8L^t$ and $R-4L^t\ge\frac\nu{48}L^t$ as $\nu\ge192$, and $\delta_{\rm tp}L(L-1)\frac\nu{48}\ge3\log L+\log2$ by the second member of $\nu_{\rm W}$. Finally the covariance on $\Lambda_N$ is the periodisation of that on $\Lambda_{N'}$, exactly: with $\Gamma=N\mathbb Z^3/N'\mathbb Z^3$, $|\Gamma|=(N'/N)^3$,
\[
C_N(\pi\iota,\pi\iota')=\sum_{\gamma\in\Gamma}C_{N'}(\iota,\iota'+\gamma).
\]
For $\pi^*$ commutes with $\mathrm d$, the $\mathsf Q_k$ and $\xi$, and $\pi$ is $|\Gamma|$ to one on the plaquettes and on the sites, so that $\mathsf H_{N'}(\pi^*A,\pi^*A')=|\Gamma|\,\mathsf H_N(A,A')$ and $\pi^*Y_N=Y_{N'}\cap\pi^*\mathcal E_N$, $\mathcal E_N$ the $1$-forms on $\Lambda_N$; $u=\sum_\gamma C_{N'}(\cdot,\iota'+\gamma)$ is the one solution in $Y_{N'}$ of $\mathsf H_{N'}(u,v)=\sum_\gamma v(\iota'+\gamma)$ for all $v\in Y_{N'}$, a problem invariant under $\Gamma$, so $u=\pi^*u_0$ with $u_0\in Y_N$, and the test functions $v=\pi^*v_0$ give $\mathsf H_N(u_0,v_0)=v_0(\pi\iota')$ for all $v_0\in Y_N$, that is $u_0=C_N(\cdot,\pi\iota')$.

\emph{Near and far.} By Laplace's method in the chart, each of the seven cumulants of $\mathsf g_{m'}$ is, by multilinearity, a sum over at most $64N'^6$ tuples of places of cumulants of homogeneous Taylor polynomials of the terms, of total order $\le10$, in the Gaussian family with the covariance $C_S$; Lemma~\ref{C-350-24}(iii) writes each as a sum over connected pairings of products of lines. A place is near if the set on which its term reads $y$ lies within $N/4$ of the base point of $p'$, and far otherwise; $\pi$ maps the near places of $\Lambda_{N'}$ one to one onto those of $\Lambda_N$, and $p'$, $q'$ are near, as $L^t\le N/192$. A tuple of near places gives the same polynomials on both tori, renamed by $\pi$, so that its cumulant on $\Lambda_N$ is that of the polynomials of $\Lambda_{N'}$ in the covariance $C_N(\pi\iota,\pi\iota')$; its indices lie within $N/2$ of each other, so that by the periodisation the two covariances differ by $\sum_{\gamma\ne0}C_{N'}(\iota,\iota'+\gamma)$, a sum of fewer than $|\Gamma|$ lines of length $\ge N/2$.

If a tuple has a far place, every connected pairing of it contains a line of length $\ge N/24$, on either torus: a path without repetitions in the graph of the pairing leads from the place $p'$ to the far place in at most three edges, the polynomials of one place depend on $y$ on a set of diameter $\le2L^t\le\frac2{192}N$, and so the lengths of the at most three lines of the path add up to more than $\frac N4-7L^t>\frac N8$. By the bound of the lines at $R=N/24$, such a line is at most $\varepsilon_C=2L^{3-3t}C_{\rm tp}\mathsf c_\tau^2e^{-\delta_{\rm tp}L\nu/48}$, and $\sup|C_S|\le2\mathsf c_\tau^2$; the bounds of Lemma~\ref{C-350-24}(iii) then bound the far tuples on each torus by a multiple of $\varepsilon_C$ and the difference at a near tuple by a multiple of $|\Gamma|\varepsilon_C$, with Cauchy's estimate on the polydisc for the coefficients of the polynomials, at most $(108L^{3t}\varsigma_t^{-3})^k$ in the sum of their moduli at the order $k$.

The powers of $L^t\le N'$, the factor $|\Gamma|\le N'^3$ and the number of the tuples are collected in $N'^{237}$, with $C'_{\rm W}$ determined by $L$ and $C_{\rm tp}$. A ball of radius $r\ge8$ holds at most $30r^3$ barycentres of plaquettes, which gives the second display.
\end{proof}

\begin{proposition}[The anchor]\label{C-350-15}
Let $\ell=2L^m\ge N_{\mathsf H}$ with $\ell L\ge C_LM$, and $K\ge1$. For all plaquettes $p,q$ of $\Lambda_{\ell L^K}$ and all plaquettes $\tilde p,\tilde q$ of $\mathbb Z^3$ over them whose base points are at sup-distance $\le\frac18\ell L^K$,
\[
|\mathsf g_2[\ell L^K,K](p,q)-\mathsf B^{\rm lat}_1(\tilde p,\tilde q)|\le C_\star\bigl[\ell^{29/2}KL^{-K/2}+\ell^{18}L^{-2K}+\ell^{19}K^2L^{-5K}\bigr]\le3C_\star\ell^{19}KL^{-K/2},
\]
with $C_\star$ determined by the constants and independent of the side; if $\ell^3\le L^K$, the left member is at most $3C_\star\ell^{29/2}KL^{-K/2}$.
\end{proposition}

\begin{proof}[Proof sketch]
Write $\mathsf N=\ell L^K$. Distances are those of the bare torus $\Lambda_{\mathsf N}$, in bare units, $|\cdot|$ its Euclidean distance and $|\cdot|_\infty$ its sup distance, $\langle z\rangle=\max\{1,|z|\}$; $x_p$ is the base point of the plaquette $p$, $\langle b-b'\rangle:=\langle b_--b'_-\rangle$, $\langle x_p-b\rangle:=\langle x_p-b_-\rangle$; and $o$ is a fixed site of the top lattice $\mathbb T^{(K)}$.

\emph{The covariant gauge below the flat top.} By Proposition~\ref{C-350-12}(iii) and its proof, $\mathsf g_2[\mathsf N,K](p,q)$ is the coefficient of $\beta^{-3}$ of one gauge-invariant quantity, the covariance of $a_p$ and $a_q$ under $e^{-\beta A}$ on the fibre $\mathcal M^KU=\mathbf 1$ near $\mathbf 1$; so every normalised gauge-fixing factor may be used to compute it. We take the covariant one with the gauge parameter $1$ of Definition~\textup{\ref{IV-73-15}}, with its Faddeev--Popov determinant as in Lemma~\textup{\ref{IV-73-12}}, for the whole bare gauge group modulo the constant transformations, and widen the flat top to its orbit under the gauge transformations of the top lattice: the top field runs over the pure gauges $w(\mu')(c)=\operatorname{Lg}(e^{i\mu'(c_-)}e^{-i\mu'(c_+)})$, $\mu'$ a $\mathfrak g$-valued function on $\mathbb T^{(K)}$ with $\mu'(o)=0$.

The constraint $\operatorname{Lg}\mathcal M^K(e^{i\mathsf a})=w(\mu')$ is linearised in straight frames: the linear part of $\operatorname{Lg}\mathcal M^K(e^{i\cdot})$ is $\mathbf Q^s_K-\mathrm d\vartheta_K$, where the straight averages $\mathbf Q^s_K$ of the free template (Section~\YMMref{120}{817}) map the bare $1$-forms to the $1$-forms of the top lattice, with $\mathrm d\mathbf Q^s_K=Q^{\rm pl}_K\mathrm d$ (Lemma~\YMMref{120.1}{818}), $Q^{\rm pl}_K$ the plaquette averages of the level $K$, and the exact difference $\mathrm d\vartheta_K$ on the top lattice goes into $\mu'$. So the linearised constraint is $\mathbf Q^s_K\mathsf a+\mathrm d\mu'=0$.

Laplace's method under the constraint, with the harmonic right inverse $\omega\mapsto(\mathsf a_\omega,\mu'_\omega)$ of the linearised constraint, which parametrises the constraint set up to the order three as a graph over the kernel
\[
\hat E=\{(\mathsf a,\mu'):\ \mathbf Q^s_K\mathsf a+\mathrm d\mu'=0,\ \mu'(o)=0\},
\]
gives the exact splitting $\mathsf g_2[\mathsf N,K]=\mathsf B^{\hat E}_1+\sum_{i\le11}\mathsf t_i$. Here $\mathsf B^{\hat E}_1$ collects the terms without the inverse: it is the formula \textup{(\ref{IV-73-24})} of $\mathsf B_1^{\rm lat}$ at the gauge parameter $1$, with the same vertices and measure terms, its kernels $\mathsf C,\mathsf E,\Pi$ being those of the centred Gaussian field $(\mathsf a,\mu')$ on $\hat E\otimes\mathfrak g$ with the form $\frac12\|\mathrm d\mathsf a\|^2+\frac12\|\mathrm d^*\mathsf a\|^2$, diagonal in the colours, and its kernels $\mathsf g,\mathsf D$ being formed with the Green's function $\mathsf G^{\rm tor}$ of $\mathrm d^*\mathrm d$ on the torus off the constants.

The $\mathsf t_i$ hold the inverse: $\mathsf t_1,\dots,\mathsf t_5$ come from the observables, $\mathsf t_6$, $\mathsf t_7$ from the cubic vertex, $\mathsf t_8$ from the quadratic form at the correction, $\mathsf t_9$, $\mathsf t_{10}$ from the traces of the Jacobian of the constraint, and $\mathsf t_{11}$ from the measure of $\mu'$; each is a cumulant of polynomials in the Gaussian field, in which the quadratic and cubic terms of the block average enter through the response $\omega\mapsto\mathsf a_\omega$ of the field to a source $\omega$ on the top lattice.

\emph{The lines.} Let $\mathsf C^{\rm tor}$ be the inverse of $\mathrm d^*\mathrm d+\mathrm d\mathrm d^*$ off the constant $1$-forms and $0$ on them, and $\mathsf S_0:=\mathsf C^{\rm tor}\mathrm d^*$ on the exact curvatures: $\mathsf S_0F$ is the primitive of $F$ in the range of $\mathrm d^*$, and $\mathrm d\mathsf S_0=1$. The $\mathsf a$-members of $\hat E$ are the $\mathsf a=\mathsf S_0F+\mathrm d\lambda$ with $F$ a curvature whose plaquette averages of the level $K$ vanish; in the Gaussian measure $F$ and $\lambda$ are independent, with the covariances $\Pi^{\rm tp}_0-\Pi^{\rm tp}_K$ and $(\mathsf G^{\rm tor})^2$. So the covariance of $\mathsf a$ is $C_{\hat E}=\mathsf C^{\rm tor}-\mathsf S_0\Pi^{\rm tp}_K\mathsf S_0^{\mathsf T}$, with $\mathrm dC_{\hat E}\,\mathrm d^*=\Pi:=\Pi^{\rm tp}_0-\Pi^{\rm tp}_K$; and the response is $\mathsf a_\omega=\mathsf S_0\mathsf H^{\rm cv}_K\omega$ plus a constant $1$-form, with $\mathrm d\mathsf a_\omega=\mathsf H^{\rm cv}_K\omega$, the curvature response of the level $K$, and $\mathrm d^*\mathsf a_\omega=0$.

The kernels $\mathsf C^{\rm tor}$, $\mathrm d\mathsf C^{\rm tor}$ and $\Pi^{\rm tp}_0=\mathrm d\mathsf C^{\rm tor}\mathrm d^*$ are Fourier sums $\mathsf N^{-3}\sum_{k\ne0}e^{ik\cdot z}\sigma(k)$ over the non-zero modes of the torus, $z$ the difference of the base points, with the symbols $\sigma$ of the lines of $\mathbb Z^3$ (Lemma~\textup{\ref{IV-73-51}}) at the gauge parameter $1$, of the classes $\mathcal S^{-2}$, $\mathcal S^{-1}$, $\mathcal S^0$; by Lemma~\textup{\ref{IV-73-75}}(a) they are at most $c_{\rm T}\langle z\rangle^{-1}$, $c_{\rm T}\langle z\rangle^{-2}$, $c_{\rm T}\langle z\rangle^{-3}$, with one number $c_{\rm T}$ for all sides. With $|\Pi^{\rm tp}_K(p,q)|\le CL^{-3K}=C\ell^3\mathsf N^{-3}$ uniformly in the volume (Proposition~\YMMref{120.5}{821}) and $\sum_z\langle z\rangle^{-2}\le14\mathsf N$ over the torus, the corrections $\mathsf S_0\Pi^{\rm tp}_K\mathsf S_0^{\mathsf T}$, $\Pi^{\rm tp}_K\mathsf S_0^{\mathsf T}$ and $\Pi^{\rm tp}_K$ are at most $c\,\ell^3\mathsf N^{-1}$, $c\,\ell^3\mathsf N^{-2}$ and $c\,\ell^3\mathsf N^{-3}$, every distance of the torus is below $\mathsf N$, and $|\Pi|\le1$, $\Pi$ being the kernel of an orthogonal projection (the orthogonal projections in~\YMMref{(2226)}{818}).

So, with $\mathsf c_0=c\,\ell^3$, the lines obey
\[
|C_{\hat E}(b,b')|\le\mathsf c_0\langle b-b'\rangle^{-1},\qquad|\mathrm dC_{\hat E}(p,b)|\le\mathsf c_0\langle x_p-b\rangle^{-2},\qquad|\Pi(p,q)|\le\min\{1,\mathsf c_0\langle x_p-x_q\rangle^{-3}\},
\]
and the kernels $\mathsf g,\mathsf D$ formed with $\mathsf G^{\rm tor}$ obey the majorants of \textup{(\ref{IV-73-56})} in the same way. The response has $|\mathsf a_\omega|\le C\ell L^{-K}$ and $|\mathrm d\mathsf a_\omega|\le CL^{-2K}$ in the sup norm of $\omega$: the latter from the row sums $\le CL^{-2K}$ of the response kernel at every admissible side (equation~\YMMref{(2234)}{820}, whose row sum follows by summing its exponential kernel bound), the former from these and the row sums $\le c\,\mathsf N$ of $\mathsf S_0$.

\emph{The Gaussian member against $\mathbb Z^3$.} Let $\mathcal L_{\mathsf N}$ be a Fourier sum over the non-zero modes of the torus of side $\mathsf N$ with a symbol of the class $\mathcal S^s$, $s>-3$, and $\mathcal L_\infty$ the kernel of $\mathbb Z^3$ with the same symbol. Then $\sum_{v\in\{0,1\}^3}\mathcal L_{2\mathsf N}(z+\mathsf Nv)=\mathcal L_{\mathsf N}(z)$ exactly, as $\sum_ve^{ik\cdot\mathsf Nv}$ is $8$ on the modes of the smaller torus and $0$ on the others; for $|z|_\infty\le\frac34\mathsf N$ the seven images with $v\ne0$ are at most $C\mathsf N^{-3-s}$ each, by Lemma~\textup{\ref{IV-73-75}}(a), and telescoping over the sides $2^i\mathsf N$, with Lemma~\textup{\ref{IV-73-75}}(b), gives $|\mathcal L_{\mathsf N}(z)-\mathcal L_\infty(z)|\le C\mathsf N^{-3-s}$.

With the three corrections from $\Pi^{\rm tp}_K$, the difference at $|z|_\infty\le\frac34\mathsf N$ between a line of $\mathsf B^{\hat E}_1$ and the line of $\mathbb Z^3$ at the same base points is at most $C\mathsf N^{-1/2}$ times the majorant $\mathsf c_0\langle z\rangle^{-n}$ of the line with its exponent $n$ lowered by $\frac12$. The two-loop sums of \textup{(\ref{IV-73-24})}, written out over monomials, connected pairings and places by Lemmas~\textup{\ref{IV-73-21}} and~\textup{\ref{IV-73-53}}, converge absolutely with these majorants, uniformly in the side, also with one line weakened by $\frac12$, as in Lemma~\textup{\ref{IV-73-55}}(ii), the sums over a torus being bounded by those over $\mathbb Z^3$ at the nearest images of the arguments: the sum is at most $C\mathsf c_0^5$.

The terms all of whose places lie within $\frac14\mathsf N$ of $x_p$ in the sup distance correspond one to one on the torus and on $\mathbb Z^3$, with the same monomials, pairings and coefficients, and together they differ by at most $C\mathsf N^{-1/2}\mathsf c_0^5$. Every other term has a place that is the end of a line whose other end is $p$ or $q$, of length $\ge\frac1{16}\mathsf N$, as $|x_p-x_q|_\infty\le\frac18\mathsf N$; so the sum of the other terms is at most $(16/\mathsf N)^{1/2}C\mathsf c_0^5$ on each side. Hence $|\mathsf B^{\hat E}_1-\mathsf B^{\rm lat}_1|\le C\ell^{15}(\ell L^K)^{-1/2}$, $\mathsf B^{\rm lat}_1$ being taken at the gauge parameter $1$, as in \textup{(\ref{IV-73-25})}.

\emph{The other terms.} Each $\mathsf t_i$ is at most $C\ell^{q_i}$ times one of the rates $KL^{-K}$, $K^2L^{-2K}$, $L^{-3K/2}$, $KL^{-K/2}$, $L^{-K/2}$, $L^{-5K/2}$, $KL^{-3K}$, $L^{-2K}$, $K^2L^{-5K}$, none worse than $KL^{-K/2}$, and $\sum_i|\mathsf t_i|\le C[\ell^{13}KL^{-K/2}+\ell^{18}L^{-2K}+\ell^{19}K^2L^{-5K}]$. The side enters only through $\mathsf c_0$, through paths of at most $\frac32\ell$ bonds on the top lattice, through one sum of $\langle\cdot\rangle^{-2}$ over the torus, through the number of plaquettes and through the dimension $9\ell^3$ of the top fields; no inverse on the top lattice is bounded by a norm, and so $C_\star$ does not depend on the side. For $\ell^3\le L^K$ the last two members are below the first.
\end{proof}

\begin{proposition}[One lattice is enough]\label{C-350-16}
Let $0<g\le\gamma$ and $\ell\in\mathfrak L(g)$. There are integers $K_0,K_\flat$ and numbers $C_\diamond<\infty$, $d_{\rm K}>0$, depending on $g$, $\ell$ and the constants, such that, with $K_\sharp(d)=\max\{K_0,K_\flat,\lceil3600\log_L(\sqrt3M/d)\rceil\}$: for $0<d\le1$, $20d\le\ell$, $\vartheta\in(0,1)$, $(f,h)\in\mathfrak T_{T_\ell}(d;\vartheta)$, $b\in[\gamma^{-2},L^2(g^{-2}+\bar B)]$ and $K'\ge K\ge K_\sharp(d)$,
\[
|S^T_{2;K',\ell,L^{K'}b}-S^T_{2;K,\ell,L^Kb}|(f,h)\le C_\diamond\vartheta^{-2}d^{178}\eta,
\]
also with $S_2^{(\ell,b)}$ for the first member; and for $0<d\le d_{\rm K}$, $s(d)<K_\sharp(d)$ and $(\sqrt3M/d)^{3600}\le L^{K_\sharp(d)}\le d^{-3601}$.
\end{proposition}

\begin{proof}[Proof sketch]
Corollary~\YMMref{117.1}{812}, equation~\YMMref{(2207)}{812} gives, for $K'\ge K\ge K_1(\mathsf d_\infty)$ and the labels of the interval $[\gamma^{-2},L^2(g^{-2}+\bar B)]$,
\[
\begin{multlined}
|S^T_{2;K'}-S^T_{2;K}|(f,h)\le\mathsf C'\ell^9K^{\mathsf p_2}L^{-K/10}\|f\|_\sharp\|h\|_\sharp,\qquad\|f\|_\sharp=\max\{\|f\|_\infty,40M\operatorname{Lip}f\} \\
\le40M\|f\|_\infty/(\vartheta d),
\end{multlined}
\]
with $\mathsf C'$ and $\mathsf p_2=6p_0+\frac12$ the constants of that statement; the limit $K'\to\infty$ exists along the full sequence. Here $\mathsf d_\infty\ge d/\sqrt3$ is the sup distance of the supports, and $K_1(\mathsf d_\infty)$ (the constants and thresholds in Theorem~\YMMref{AR}{810}) is the least integer $\ge K_0$ from which on, with $u=\lceil K/1800\rceil$, the two threshold inequalities at the end of the proof of Theorem~\YMMref{AR}{810} hold: $\mathsf d_\infty L^u\ge6M$, and $c_9\ell^3L^{3u}e^{-\mathsf d_\infty L^u/M}\le\frac12\mathsf C\ell^9K^{\mathsf p_2}L^{-K/10}$, with the numbers $c_9$, $\mathsf C$ of the two threshold inequalities at the end of the proof of Theorem~\YMMref{AR}{810} at $n=2$: $\mathsf C$ is the constant \YMMref{(2206)}{811}, and $c_9\ell^3L^{3u}$ is the factor of $e^{-\mathsf d_\infty L^u/M}$ in the left member of the second condition.

Let $K\ge K_\sharp(d)$ and $y=K/1800$. The second condition follows, in logarithms, from $\mathsf d_\infty L^u/M\ge c_{10}+(3u+\frac K{10})\log L$, $c_{10}=\max\{0,\log(2c_9/\mathsf C)\}$, and this from the two inequalities $L^{y/2}\ge c_{10}+3\log L+200y\log L$, which defines $K_\flat=\lceil1800y_0\rceil$ through the least $y_0$ from which on it holds, and $L^{y/2}\ge\sqrt3M/d$, which is $K\ge3600\log_L(\sqrt3M/d)$; the first condition follows from the same product. So $K_1(\mathsf d_\infty)\le K_\sharp(d)$, and $K^{\mathsf p_2}L^{-K/10}\le cL^{-K/20}\le c(d/\sqrt3M)^{180}$, which gives the factor $C_\diamond\vartheta^{-2}d^{178}$. By the minimality of $s(d)$, $L^{s(d)}d<LC_{W,0}A_0^2(\log L^{s(d)}+\log b_+)^{2p_0}$, so that $L^{s(d)}\le d^{-2}$ for small $d$; and $(\sqrt3M/d)^{3600}\le L^{K_\sharp(d)}<L(\sqrt3M/d)^{3600}\le d^{-3601}$ for small $d$.
\end{proof}

\begin{proposition}[The agreement of states]\label{C-350-21}
Let $0<g\le\gamma^{\rm sd}_W$, $(K,\beta)$ a datum at $g$ with $K\ge3$ and label $b$, $\vartheta\in(0,1)$, and $0<d\le\frac1{180}$ with $a_K\le9d$. A \emph{state at $g$ at $(K,\beta)$} is the periodically extended torus measure $\tilde\mu_{K,\ell',\beta}$ of a side $\ell'\in\mathfrak L(g)$, or a periodic Gibbs state at $(a_K,\beta)$; $S^T_2[\tilde\mu_{K,\ell',\beta}](\tilde f,\tilde h)$ is $S^T_{2;K,\ell',\beta}$ at the periodisations of $\tilde f,\tilde h$ on $T_{\ell'}$. Put $\mathsf d_\star(g,\vartheta)=4\sqrt3M\vartheta^{-1}(K^\vee_{\rm L}/K^\vee_{\rm Q})^2g^{-2}$, with the pair $K^\vee_{\rm L},K^\vee_{\rm Q}$ of equation~\YMMref{(2166)}{800}. There are $\mathsf Q^{\rm T}_2(g)\ge1$ and $C^{\rm T}_2(g)<\infty$, depending on $g$, on the fixed radius $\varrho$ and the common sourced choice of Section~\YMMref{82}{665} specified at the beginning of this section, and on no $K$, $\beta$, side, state, $d$, $\vartheta$, pair or position, with the following property.

Let $(\tilde f,\tilde h)$ be real Lipschitz functions on $\mathbb R^3$, neither $0$, each vanishing off a closed ball of radius $d$, with disjoint sets $\{\tilde f\ne0\}$, $\{\tilde h\ne0\}$, both vanishing off one ball of radius $9d$, and with $\operatorname{Lip}\tilde f\le\|\tilde f\|_\infty/(\vartheta d)$, $\operatorname{Lip}\tilde h\le\|\tilde h\|_\infty/(\vartheta d)$; no sign and no lower bound of the integrals is asked. Then for every state $\mathfrak s$ at $g$ at $(K,\beta)$
\[
|S^T_2[\mathfrak s]-S^T_2[\tilde\mu_{K,\ell_g,\beta}]|(\tilde f,\tilde h)\le(30d)^3\cdot2\,\mathsf Q^{\rm T}_2(g)^2\,\mathsf A_\star(d)^{-2}\eta,
\]
with a number $\mathsf A_\star(d)$ of the order $(b\vartheta d)^{1/2}$ formed with the datum; and for $d\le\mathsf d_\star(g,\vartheta)$, in particular for the pairs of $\mathfrak T_{\mathbb R^3}(d;\vartheta)$,
\begin{equation}\label{C-350-17}
\bigl|S^T_2[\mathfrak s]-S^T_2[\tilde\mu_{K,\ell_g,\beta}]\bigr|(\tilde f,\tilde h)\le C^{\rm T}_2(g)\vartheta^{-1}d^2b^{-2}\eta .
\end{equation}
Any two states at $g$ at $(K,\beta)$ differ at such pairs by at most twice this.
\end{proposition}

\begin{proof}[Proof sketch]
Let $\ell'=\ell_g=2L^m$, $\mu'=\tilde\mu_{K,\ell',\beta}$ and $\mathsf l=\ell'/2=L^m$. Let $x_1$ be the centre of the ball of radius $9d$ off which $\tilde f$ and $\tilde h$ vanish, and let $o\in a_K\mathbb Z^3$ be a site with $|o_\mu-(x_1)_\mu|\le\frac12a_K$ for $\mu=1,2,3$. Consider the board of the least torus: the eight cells $\square_{\mathbf k}=o+\mathsf l\mathbf k+[-\frac{\mathsf l}2,\frac{\mathsf l}2]^3$, $\mathbf k\in\{0,1\}^3$, with the reflections $x_\mu\mapsto2o_\mu+\mathsf l-x_\mu$ in the planes $x_\mu=o_\mu\pm\frac{\mathsf l}2$, which carry one cell into another. As $o\in a_K\mathbb Z^3$ and $\mathsf l$ is a multiple of $a_K$, these reflections are maps of the lattice, and they carry plaquettes to plaquettes.

The source families used below are those of
Proposition~\ref{prop:levelwise-source-density}, at the common
sourced choice fixed in this section. Its numerical source
conditions are the ones retained in
Section~\ref{sec:levelwise-first-creation}.
The radius condition for Proposition~\ref{prop:first-creation-envelopes}
holds at the present coupling: $g\le\gamma_W^{\rm sd}\le
\gamma_W^{\rm ref}$ and equation~\YMMref{(2326)}{845} give
$\varepsilon(x_t)\le c_A\alpha_0/(8MB_3)$ at every level.
Also $\mathsf l\ge L^2M$ by the least-side definition,
so both entries in the minimum defining $d_*(t)+1$ in
\eqref{eq:first-creation-radius} are at least $4$.
The calibrated rarity hypotheses are those of the same real
history, as in Proposition~\ref{prop:ward-fixed-volume-rarity}.
Thus Corollary~\ref{cor:packed-source-coefficients} applies
without a further restriction on this datum.

As $1+\frac12a_K<\frac12\mathsf l$, the ball $\mathbb B(x_1,1)$ lies in the interior of $\square_0$; the signed sums of (3) below vanish off it. Its images under the reflections are eight balls of radius $1$, one in each cell, with the centres $o_\mu+\mathsf lk_\mu+(-1)^{k_\mu}(x_1-o)_\mu$ ($\mu=1,2,3$); any two of these centres differ, in a coordinate $\mu$ in which their cells differ, by $\pm(\mathsf l-2(x_1-o)_\mu)$ modulo $2\mathsf l$, so that they lie at a sup distance at least $\mathsf l-2|(x_1-o)_\mu|\ge\mathsf l-a_K$ on $T_{\ell'}$.

(1)

\emph{One cell against the least torus.} For a bounded
complex Borel function $F$ of the inner bonds of one cell,
the reflected copies, conjugated on odd reflection parity,
satisfy
\[
 \|F\|_{\rm ch}^8=
       \mu'\Bigl(\prod_{\mathbf k\in\{0,1\}^3}F^{[\mathbf k]}\Bigr)
       \ge0.
\]
For $\ell\in\mathfrak L(g)$,
Corollary~\ref{IV-18-18} gives
\[
 |\tilde\mu_{K,\ell,\beta}(F)|
 \le\|F\|_{{\rm ch},\ell}
 \le\|F\|_{\rm ch}\exp[\mathsf l^3(p_{\ell'}-p_\ell)].
\]
Lemma~\ref{lem:canonical-source-system} bounds
$p_{\ell'}-p_\ell$ by
$C_{\rm gap}+C_{\rm S}+C_g\le2C_{\rm S}+C_g$.
Thus one may take
$\mathsf Q_0=\exp[\mathsf l^3(2C_{\rm S}+C_g)]\ge1$.
For a periodic Gibbs state, Proposition~\YMMref{91.9}{705}
gives
\[
 |\mathfrak s(F)|\le
 \liminf_{r\to\infty}\|F\|_{{\rm ch},\ell_gL^r}
 \le\mathsf Q_0\|F\|_{\rm ch}.
\]
The constants are those of Lemma~\ref{lem:canonical-source-system},
with $C_{\rm S}\ge C_{\rm gap}$ chosen there. In particular
$\mathsf Q_0$ is independent of the cutoff and the larger side.

This is the one place where Theorems~\YMMref{AJ}{673} and~\YMMref{AD}{607} enter, and the place through which $\mathsf Q^{\rm T}_2(g)$ and $C^{\rm T}_2(g)$ depend on the radius and on the choice of the longer list.

(2)

\emph{Only linking cumulants are left.} With $X_a=\mathcal O(F_a)$, $\mathcal P_a=X_a-\mu'(X_a)$ and $\mathcal P_{12}=\mathcal P_1\mathcal P_2-S^T_2[\mu'](F_1,F_2)$ one has $S^T_2[\mathfrak s](F_1,F_2)-S^T_2[\mu'](F_1,F_2)=\mathfrak s(\mathcal P_{12})-\mathfrak s(\mathcal P_1)\mathfrak s(\mathcal P_2)$. Expanding $\mu'(\prod_{\mathbf k}\mathcal P_{12}^{[\mathbf k]})$ in joint cumulants of the sixteen copies $X_{a,\mathbf k}$, the blocks that stay in one cell cancel exactly, because $\mathcal P_{12}$ is centred by the covariance of $\mu'$ itself; what is left is a sum, over the at most $\mathsf{Bl}_{16}$ partitions into linking blocks (blocks that meet two cells), of products of linking cumulants, $\mathsf{Bl}_n$ being the Bell numbers.

So if every linking cumulant of the copies is at most $C_{\rm LK}\prod u_a$, then $|S^T_2[\mathfrak s]-S^T_2[\mu']|(F_1,F_2)\le[\mathsf Q_0\mathsf{Bl}_{16}^{1/8}+\mathsf Q_0^2\mathsf{Bl}_8^{1/4}]C_{\rm LK}u_1u_2\le2\mathsf Q_1^2u_1u_2$, $\mathsf Q_1=\mathsf Q_0\mathsf{Bl}_{16}^{1/8}C_{\rm LK}$.

(3)

\emph{The packing.} Apply (2) to the signed sums $F^\epsilon_a=\sum_v\epsilon_v\tau_vf_a$, $f_1=\tilde f$, $f_2=\tilde h$, over a packing of translates $\mathbb B(x_1+v,9d)$, $v\in a_K\mathbb Z^3$, in the ball $\mathbb B(x_1,1)$, with at least $(30d)^{-3}$ members; the mean over the signs and the invariance of every state under the lattice translations give $|S^T_2[\mathfrak s]-S^T_2[\mu']|(\tilde f,\tilde h)\le(30d)^3\cdot2\mathsf Q_1^2u_1u_2$.

(4)

\emph{The linking cumulants under $\mu'$.} Their generating function is $Z[W_z]/Z$ on $T_{\ell'}$, with the source $W_z=\sum_{\mathbf k,a}z_{\mathbf k,a}F^\epsilon_{a,\mathbf k}\circ x$ read on the torus, $F^\epsilon_{a,\mathbf k}$ the copy in the cell $\mathbf k$. By the property of the level-wise class recalled in the proof of Proposition~\ref{C-350-8}, with $b^\pm_K=x_K\pm\bar B$ in the place of $b^\vee$, $b^\wedge$, the source lies in the class with norm at most $\max|z_{\mathbf k,a}|u_a$, $u_a=\|f_a\|_\infty/\mathsf A_\star(d)$. Write $\log(Z[W_z]/Z)=\log\Omega_8(z)+[\mathcal E(W_z)-\sum_{\mathbf k}\mathcal E(w^{(\mathbf k)})]+\sum_{\mathbf k}\mathcal E(w^{(\mathbf k)})$, $\Omega_8=Z[W_z]e^{-\mathcal E(W_z)}/Z$, $w^{(\mathbf k)}$ the part of $W_z$ in the cell $\mathbf k$.

A linking derivative kills the last sum. Apply
Corollary~\ref{cor:packed-source-coefficients} to the eight
reflected balls and their sixteen variables. Its hypotheses on
the bumps hold because reflections preserve their amplitude,
Lipschitz bounds and disjoint nonzero sets. The same amplitude
calculation in Lemma~\ref{lem:packed-levelwise-source}, now with
$b_K^-=x_K-\bar B$ and $b_K^+=x_K+\bar B$, applies: these
bound $a_{K-j}x_j$ below and above. Thus the square-free
derivatives of $\Omega_8$ at zero are at most
$\mathsf Q_{\rm D}\prod u_a$, with
$\mathsf Q_{\rm D}=Q_{8,16}(g)$.

We estimate the bracket in the generating function by using
the local energies of Proposition~\ref{prop:levelwise-source-density}:
\[
 e_j(X;w)=\sum_{z\in X\cap\mathbb T^{(j)}}
             \Delta F_j[w](X,\mathbf1,z).
\]
They depend only on the source plaquettes in $X$ and vanish
at zero. The counting calculation in
\eqref{eq:running-source-energy-bound}, applied to
\eqref{eq:levelwise-family-bound}, gives
\[
 |e_j(X;w)|\le16M^3E_{\rm s}\|w\|_{\rm lw}
                                      e^{-\kappa d_j(X)}.
\]
Here $E_{\rm s}$ is the fixed family ceiling of
Section~\ref{sec:levelwise-first-creation}.
If $X$ meets source plaquettes in at most one reflected
ball, its contribution to
$\mathcal E(W_z)-\sum_{\mathbf k}\mathcal E(w^{(\mathbf k)})$
is zero: locality identifies the one possible nonzero term,
and all other terms vanish. Only domains meeting two balls
remain.

The centres of the reflected balls are at sup distance at
least $\mathsf l-a_K$. Their source barycentres are therefore
at distance at least $\mathsf l-a_K-2$, and a corner is
within $a_K/2$ of its barycentre. Thus a domain that contributes
contains two points at distance at least $\mathsf l-3$,
since $K\ge3$ and $L\ge13$ imply $2a_K\le1$.
Lemma~\ref{lem:near-domain-geometry} gives
\[
 d_j(X)\ge\frac{\mathsf l-3}{MD_j}-2
            \ge12L^{K-j}-2,\qquad D_j=L^{j-K}.
\]
For a wrapped domain the bound follows instead from
$M(d_j(X)+2)\ge\ell'/(2D_j)=\mathsf l/D_j$, which is
stronger. In the last inequality we used
$\mathsf l\ge L^2M$ and $M\ge5$. The least-side formula
\YMMref{(1528)}{615} gives this side bound: its maximum
includes $L^{\lceil\log_L C_{\rm s}\rceil}\ge L$.
In particular the displayed lower bound is positive also at
$j=K$.

Put $r=\max_{\mathbf k,a}|z_{\mathbf k,a}|u_a<1$.
The full source and each cell source have level-wise norm
at most $r$, so the contribution of one domain is at most
$9\cdot16M^3E_{\rm s}r e^{-\kappa d_j(X)}$.
A ball of radius $1$ meets source plaquettes whose corners
lie in a cube of side $2+a_K\le3$. There are at most
$(3L^{K-j}/M+2)^3\le27L^{3(K-j)}$ level-$j$ cubes containing
these corners. The eight balls therefore require at most
$216L^{3(K-j)}$ such cubes. Apply
Lemma~\ref{lem:source-support-sum} with $t=j$, $q=0$ and
$D=12L^{K-j}-2$. Its rate condition holds because
$\delta_{\rm w}\kappa>\kappa_{\mathcal D}+2$ and
$0<\delta_{\rm w}<1/2$ in equation~\YMMref{(877)}{412}.
The extra factor $\#(X\cap\mathbb T^{(j)})$ in that lemma
is at least one. Summing first in $X$ and then in $j$ gives
\[
 \left|\mathcal E(W_z)
              -\sum_{\mathbf k}\mathcal E(w^{(\mathbf k)})\right|
 \le\mathsf C_\Delta r,\qquad
 \mathsf C_\Delta=
 248832M^6E_{\rm s}C_{\mathrm E,1}
                  \sum_{n\ge0}L^{3n}e^{2-12L^n}.
\]
The series converges, since the ratio of successive terms is
$L^3e^{-12(L-1)L^n}\to0$. Cauchy's formula on the source
polydisc bounds each nonempty square-free derivative of this
bracket by $\mathsf C_\Delta\prod u_a$.

So $C_{\rm LK}=\mathsf{cu}_{16}\mathsf Q_{\rm D}^{16}+\mathsf C_\Delta$, with a combinatorial number $\mathsf{cu}_{16}$, and $\mathsf Q^{\rm T}_2(g):=\mathsf Q_1$ gives the first display.

(5) For $d\le\mathsf d_\star(g,\vartheta)$ the second member of the minimum that defines $\mathsf A_\star(d)$ is the smaller one, so that $\mathsf A_\star(d)^{-2}=(K^\vee_{\rm L}/\mu_\star)^2(b^+_K/(b^-_K)^2)\,4\sqrt3M/(\vartheta d)$, and $b^+_K/(b^-_K)^2\le b^{-2}(X_g+\bar B)^3/(g^{-2}-\bar B)^2$. With $30^3\cdot2=54000$ this is \eqref{C-350-17} with $C^{\rm T}_2(g)=54000[\mathsf Q^{\rm T}_2(g)(K^\vee_{\rm L}/\mu_\star)(X_g+\bar B)^{3/2}(g^{-2}-\bar B)^{-1}(4\sqrt3M)^{1/2}]^2$. The last assertion is the triangle inequality through $\mu'$. Of this proposition the proof of clause (c) uses \eqref{C-350-17} alone, for one torus measure against the least torus.
\end{proof}

\begin{proposition}[One datum on two sides]\label{C-350-25}
Let $0<g\le\gamma_{\rm s}$, $(K,\beta)$ a datum at $g$ with $K\ge3$ and label $b$, $8a_K\le d\le L^{-2}$, $\vartheta\in(0,1)$, $s\ge0$, $(f,h)\in\mathfrak T_{T_{\ell_g}}(d;\vartheta)$, and $(\breve f,\breve h)$ the pair on $T_{\ell_gL^s}$ with the same lift; the datum, with its one run, is taken on both tori. Put $\mathcal y=\log(10LMb/d)$.
\begin{enumerate}
\item[(i)] For every level $1\le j\le K$, $\mathsf n_j[\ell_g](f,h)-\mathsf n_j[\ell_gL^s](\breve f,\breve h)=\mathsf R_j[\ell_g]-\mathsf R_j[\ell_gL^s]$, where $\mathsf R_j[\ell]$ is the part of the sum that defines $\mathsf n_j$ on $T_\ell$ over the domains $X$ that meet both supports and are not the image of a domain of the infinite lattice lying in a cuboid of sides $<\frac12\ell_g$; each such domain has linear size $\ge\ell_g/(2MD_j)-2$.
\item[(ii)] On every side of $\mathfrak L(g)$, at every level $j$ with $D_j>Ld/8$,
\[
|\mathsf n_j(f,h)|\le C'_{\mathsf n}\,\mathcal y^{3p_0}\Bigl(\frac d{D_j}\Bigr)^3\bigl[b^{-3/2}D_j^{-1/2}+b^{-1}d^3D_j^{-4}\bigr]\eta .
\]
\item[(iii)] $|\mathsf J_2[\ell_g](f,h)-\mathsf J_2[\ell_gL^s](\breve f,\breve h)|\le C_{\rm V}\,\mathcal y^{3p_0}\,b^{-1}d^3\,\eta$.
\end{enumerate}
$C'_{\mathsf n}$ and $C_{\rm V}$ are determined by the constants; the bounds depend on no $s$, $K$, $\vartheta$ or pair beyond $\eta$, and no bound on $L^K$ and no witness depth is asked.
\end{proposition}

\begin{proof}[Proof sketch]
If a source vanishes, all the terms vanish. (i) A domain of the level $j$ of the infinite lattice that lies in a cuboid of sides $<\frac12\ell_g$ and meets the lifted supports projects onto exactly one domain of each of the two lattices $(\ell_g,K)$ and $(\ell_gL^s,K)$, with their levels; these do not wrap, and it is a near lift of each. Every domain of either lattice that meets both supports and is not of this kind either wraps the torus or lies in no such cuboid, and in both cases its linear size is $\ge\ell_g/(2MD_j)-2$. These are the following facts on the localisation domains (Lemma~\ref{lem:near-domain-geometry}). The projection of a domain of the infinite lattice is a domain of the lattice $(\ell,K)$; a domain of the lattice $(\ell,K)$ does not wrap if and only if it is the projection of a domain of the infinite lattice that lies in a cuboid of lengths less than half the side; such domains are its near lifts, and they are translates of one another by periods and pairwise disjoint; the projection keeps the linear size $d_j(X)$, the number of edges of a least edge tree of $X$ in the units of the level $j$ (Section~\YMMref{72}{607} (the definitions of $\mathcal D_j$ and $d_j$ at its beginning)); every domain of the infinite lattice lies in a cuboid of lengths $\le MD_j(d_j(X)+2)$; and on $T_\ell$ a domain with $M(d_j(X)+2)<\ell/(2D_j)$ does not wrap.

A cuboid of lengths $<\frac12\ell_g$ has lengths less than half the side of either lattice. On a near lift the two pairs of lattice sources agree: the lifted supports lie in one ball of radius $6d$, and their translates by non-zero periods are more than $\frac12\ell_g$ away. On a domain that does not wrap the local terms of the weighted progression are then the same number on the two tori. This is proved by induction over the levels: the regular progression without sources is one family for all the lattices $(\ell,K)$ of the level (Theorem~\YMMref{Y}{438}, including its common-companion assertion); one step of the progression is natural under the chart of the domain from one lattice to the other, the two lattices having one $K$: bare arrays of the two lattices that correspond under the chart (the common-input condition at the beginning of Section~\YMMref{52.5}{425}) give local terms of the next level that are defined together and are equal (equation~\YMMref{(928)}{426}, including its separately equal three summands); and the bare arrays of the two lattices correspond under this chart, their members of the lower levels by the induction over the levels and their weight members because the sources agree on the plaquettes of the domain.

So the terms of these domains cancel in the difference, which is (i).

(ii) The first jets of the coarse levels gain the factor $(d/D_j)^3$. Here $\dot E^{(j)}_a$ and $\ddot E^{(j)}$ are the derivatives $\partial_{\mu_a}E^{(j)}_\mu$ and $\partial_{\mu_1}\partial_{\mu_2}E^{(j)}_\mu$ at $\mu=0$ of the local terms of the weighted progression, and $\|\cdot\|$ is the size \YMMref{(759)}{375} of a family of the level $j$ at the rate $\kappa^+=(1+4\beta_{\rm s})\kappa$, the rate at which the part of the local terms that depends on the sources decays in the linear size of the domain (Proposition~\YMMref{122.1}{826}); $\kappa$ is the decay rate and $\beta_{\rm s}\in(0,\frac15]$ the enlargement fraction among the constants (Definition~\ref{def:weighted-progression-data}), and $\|\cdot\|_{j,\kappa}\le\|\cdot\|$ is the size at the rate $\kappa$.

The members of $\dot E^{(i)}_a$ vanish on the domains that hold no plaquette of $\operatorname{supp}w_a$, and those of $\ddot E^{(i)}$ on the domains that do not hold a plaquette of each support (Proposition~\textup{\YMMref{122.1}{826}}(c)). The cubes of the level $i$ are the cubes of side $MD_i$ of the partition $\pi_i$ of the torus whose wall-connected unions are the domains of the level $i$ (the level-$i$ localization partition of Section~\YMMref{45}{351}); for a set $\mathcal S$ of plaquettes, $\mathsf N_i(\mathcal S)$ is the number of the cubes of the level $i$ in which a corner of a plaquette of $\mathcal S$ lies.

\emph{The level $j_1$.} Let $\mathsf r$ be the least sup distance between a corner of a plaquette of $\operatorname{supp}w_1$ and a corner of a plaquette of $\operatorname{supp}w_2$. The barycentres of the plaquettes of $\operatorname{supp}w_1$ and of $\operatorname{supp}w_2$ lie in $\operatorname{supp}f$ and $\operatorname{supp}h$, which lie in balls of radius $d$ with centres at most $8d$ apart and are at a Euclidean distance at least $d$, and a corner lies at the sup distance $\frac12a_K\le\frac1{16}d$ from the barycentre; so $\frac25d\le\mathsf r<11d$. Let $j_1$ be the least level $j\ge1$ with $\mathsf r<3MD_j$. As $3MD_{K-1}=3M/L>11d$, it exists and $j_1\le K-1$; and $\mathsf r\ge3MD_j$ at every level $1\le j<j_1$.

As a consequence of this definition, $D_{j_1}$ lies in the interval $(\frac{2d}{15M},\frac{Ld}8]$: $D_{j_1}>\mathsf r/3M\ge\frac{2d}{15M}$; and $D_{j_1}\le\frac{Ld}8$, since $D_{j_1}=LD_{j_1-1}\le L\mathsf r/3M<\frac{11Ld}{3M}\le\frac{Ld}8$ if $j_1\ge2$, and $D_{j_1}=La_K\le\frac{Ld}8$ if $j_1=1$. The level $j_1$ is formed on each of the two tori; the bounds below are the same on both. All the corners of the plaquettes of $\operatorname{supp}w_a$ lie in a cube of side $\frac{17}8d$, so that $\mathsf N_i(\operatorname{supp}w_a)\le(\frac{17d}{8MD_i}+2)^3$; as $d/MD_{j_1}<\frac{15}2$, this is less than $18^3$ for $i\ge j_1$, and less than $18^3L^{3(j_1-i)}$ for $i\le j_1$.

So the bare array of the first derivative in the direction $w_a$ sits at few places at the levels from $j_1$ on. For such an array the size of the family that one step forms is bounded with the weights $L^{-4(j-i)}$ times the number of the places, where the general size has $L^{3(j-i)}$ places times $L^{-4(j-i)}$: if the counterterms of a bare array are $0$, and its members $G^{(i)}$, $i\le j$, and its weight member $H^{(j)}$ vanish on the domains that hold no plaquette of a set $\mathcal S$ of plaquettes, then the family that the rule of the step $j+1$ forms of it has a size at most
\[
c_\sharp\Bigl[\sum_{i\le j}L^{-4(j-i)}\,\mathsf N_i(\mathcal S)\,\|G^{(i)}\|_{i,\kappa}+\mathsf N_j(\mathcal S)\,\|H^{(j)}\|_{j,\kappa}\Bigr],
\]
$H^{(j)}$ being the weight member and $c_\sharp$ a number determined by the constants. The proof is that of the bound of the bracket localisation of a bare array (the proof of Proposition~\YMMref{51.16}{412}, especially the rooted-domain count and equation~\YMMref{(881)}{413}), with the count of the places changed: as the counterterms vanish, the members enter the rule as they stand; a member of the level $i$ enters the rule of the height $j$ either on a domain whose linear size is at least $\max\{2L^{j-i-1}-2,0\}$, which costs $e^{-\max\{2L^{j-i-1}-2,0\}}\le L^{4-4(j-i)}$, or through its octet function, which lies within $\mathsf C_{\rm oc}\|G^{(i)}\|_{i,\kappa}(d_i(X)+2)^2L^{4-4(j-i)}$ of a constant (Proposition~\YMMref{51.14}{408}, at fixed radii), where $\mathsf C_{\rm oc}$ is the larger of $1$ and its constant $C_{\rm oc}$; the weight member enters at the level $j$ alone; and the sum over the domains that hold a plaquette of $\mathcal S$ is bounded by Lemma~\ref{lem:source-support-sum}, which gives the factor $\mathsf N_i(\mathcal S)$ in the place of the $L^{3(j-i)}$ places of the general bound.

Two bare arrays are of this kind. The bare array of the first derivative in the direction $w_a$ at the height $j$ has the members $\dot E^{(i)}_a$, $i\le j$, the counterterms $0$, since the counterterms do not depend on the sources (Proposition~\textup{\YMMref{122.1}{826}}(e)), and as weight member that of the source $L^{K-j}w_a$, whose members sit at the domains $X(p)$ of the plaquettes $p\in\operatorname{supp}w_a$, each of which holds its plaquette; it is of this kind with $\mathcal S=\operatorname{supp}w_a$. The bare array of the second derivative has the members $\ddot E^{(i)}$, $i\le j$, the counterterms $0$ and the weight member $0$; it is of this kind with $\mathcal S=\operatorname{supp}w_1$.

The step depends on its bare array through the size of this family: by the Lipschitz bound and the holomorphy of the step in the family that the rule forms (Theorem~\YMMref{V}{416}, including equation~\YMMref{(899)}{416}), the first derivative of the step at a bare array of the progression in the direction of an array is at most $\theta_j$ times the size of the family formed of that array, and the mixed second derivative in the directions of two arrays at most $2\theta_j/N^{[j]}_\star$ times the product of the two sizes, by Cauchy's estimate; here $\theta_j=\theta^{[j]}=K_{\mathbf M}(1+\beta_{\rm s})\delta(x_j)/\alpha_0$ and $N^{[j]}_\star$ are the numbers of \eqref{eq:weighted-progression-numbers}, $K_{\mathbf M}\ge1$ is the step constant and $\alpha_0$ the first fixed radius in Definition~\ref{def:weighted-progression-data}.

Every running coupling of the datum lies above the threshold $x_{\rm J}$ of the definitions preceding \eqref{eq:weighted-history-threshold} (Lemma~\YMMref{126.2}{845}, as in the proof of Proposition~\ref{C-350-6}), and the conditions \eqref{eq:weighted-history-threshold} that define $x_{\rm J}$ hold at every $x\ge x_{\rm J}$, as $\delta$ decreases there; so $C^\flat_0\mathsf S_{\rm A}\theta_j\le1$ and $\frac43C^\flat_1L\mathcal T(3)\theta_j\le1$ at every height $j<K$, with the numbers $C^\flat_0$, $C^\flat_1$ of Section~\ref{sec:weighted-progression-interface}, $\mathsf S_{\rm A}$ of \eqref{eq:weighted-running-numbers} and the function $\mathcal T(u)=\sum_{n\ge1}(n+1)^{2p_0+1}e^{-4\beta_{\rm s}\kappa[L^{n-1}u-2]_+}$, written $\mathsf T$ in \eqref{eq:weighted-mixed-numbers}, all determined by the constants.

\emph{The first jets.} Fix $a$; let $\dot{\mathsf B}_j=\sum_{i\le j}L^{-4(j-i)}\mathsf N_i(\operatorname{supp}w_a)\|\dot E^{(i)}_a\|$ be the weighted sum of the sizes of the first jets of the levels up to $j$, and $\mathsf v_j$ the size of the weight member of the level $j$. For $j_1\le j\le K-1$ the bound just stated, with $\mathsf N_j(\operatorname{supp}w_a)<18^3$, gives the family formed of the bare array of the first derivative at the height $j$ a size at most $c_\sharp[\dot{\mathsf B}_j+18^3\mathsf v_j]$; the first derivative of the step gives $\|\dot E^{(j+1)}_a\|$ at most $\theta_j$ times this size; and $\dot{\mathsf B}_{j+1}=L^{-4}\dot{\mathsf B}_j+\mathsf N_{j+1}(\operatorname{supp}w_a)\|\dot E^{(j+1)}_a\|$.

The power $d^3D_j^{-4}$ comes from two inputs. The first is the inhomogeneous term: the weight member of the level $j$ is the action term $\eta_j^{-1}\sum_pL^{K-j}w_a(p)\,a_p$, $\eta_j=L^{-j}$ being the bare spacing in the units of the level $j$; as $\operatorname{supp}w_a$ has at most $3(\frac{17}8)^3(d/a_K)^3<29(d/a_K)^3$ plaquettes, as $|a_p|<\frac12\alpha_0^2\eta_j^4$ at the fields of the radii $(\alpha_0,\alpha_1)$, and as the domains $X(p)$ have the linear size $0$, its size is $\mathsf v_j\le\frac{29}2\alpha_0^2d^3D_j^{-4}|w_a|_\infty$. The second is the start: by the general bound of the first jets, $\|\dot E^{(i)}_a\|\le C^\flat_0\mu_{i-1,a}$ with $\mu_{i,a}=12M^3\alpha_0^2\theta_i|w_a|_\infty/D_i$ (Proposition~\YMMref{122.1}{826}, equation~\YMMref{(2258)}{827}, \eqref{eq:weighted-running-numbers}), by $\mu_{i-1,a}\le(\frac43)^{1/2}L^{n/2}(1+n)^{p_0}\mu_{j_1-1,a}$ for $n=j_1-i\ge0$ (the running-coupling comparison following equation~\YMMref{(2258)}{827}) and by $\mathsf N_i(\operatorname{supp}w_a)<18^3L^{3n}$, the powers of $L$ leave the convergent series $\sum_n(1+n)^{p_0}L^{-n/2}$, and $\dot{\mathsf B}_{j_1}\le C\mu_{j_1-1,a}\le C'\alpha_0^2|w_a|_\infty/D_{j_1}\le C'(\frac L8)^3\alpha_0^2|w_a|_\infty d^3D_{j_1}^{-4}$, as $\theta_{j_1-1}\le(C^\flat_0\mathsf S_{\rm A})^{-1}$ and $D_{j_1}\le\frac{Ld}8$.

With $\hat{\mathsf v}=\frac{29}2\alpha_0^2d^3D_{j_1}^{-4}|w_a|_\infty$, so that $L^{4(j-j_1)}\mathsf v_j\le\hat{\mathsf v}$ for $j\ge j_1$, the three relations give
\[
L^{4(j+1-j_1)}\dot{\mathsf B}_{j+1}+18^3\hat{\mathsf v}\le\bigl(1+L^4\cdot18^3c_\sharp\theta_j\bigr)\bigl(L^{4(j-j_1)}\dot{\mathsf B}_j+18^3\hat{\mathsf v}\bigr)\qquad(j_1\le j\le K-2).
\]
So it is the quantity $L^{4(j-j_1)}\dot{\mathsf B}_j+18^3\hat{\mathsf v}$, and not the weighted sum itself, that grows from the level $j_1$ on by at most the factor $\prod_i(1+L^4\cdot18^3c_\sharp\theta_i)\le\exp[L^4\cdot18^3c_\sharp\sum_i\theta_i]$; and $\sum_{i<K}\theta_i\le\mathsf S_{\rm B}\theta_{K-1}\le\mathsf S_{\rm B}/(C^\flat_0\mathsf S_{\rm A})$ (the running-coupling comparison following equation~\YMMref{(2258)}{827}, with the number $\mathsf S_{\rm B}$ of \eqref{eq:weighted-running-numbers}). Hence $\dot{\mathsf B}_j+18^3\mathsf v_j\le L^{-4(j-j_1)}\bigl(L^{4(j-j_1)}\dot{\mathsf B}_j+18^3\hat{\mathsf v}\bigr)\le Cd^3D_{j_1}^{-4}L^{-4(j-j_1)}\alpha_0^2|w_a|_\infty=Cd^3D_j^{-4}\alpha_0^2|w_a|_\infty$; the family formed of the bare array of the first derivative at the height $j$ has a size at most $C_{\rm F}\alpha_0^2|w_a|_\infty d^3D_j^{-4}$, and the recursion of the first jets closes as $\|\dot E^{(j+1)}_a\|\le\theta_jC_{\rm F}\alpha_0^2|w_a|_\infty d^3D_j^{-4}$ for $j_1\le j\le K-1$, with a number $C_{\rm F}$ determined by the constants.

\emph{The second jets.} At the levels up to $j_1$ the general bound holds: $\|\ddot E^{(i)}\|\le C^\flat_1\varpi_{i-1}$ for $1\le i\le j_1$, where $\varpi_i=c_\nu\delta(x_i)^2|w_1|_\infty|w_2|_\infty D_i^{-2}$ is the number written $\nu_i$ in \eqref{eq:weighted-mixed-numbers} and $c_\nu$ is determined by the constants. This is the induction of Proposition~\YMMref{122.2}{827}, carried up to the level $j_1$: its condition $\frac43C^\flat_1L\theta_j\mathcal T(\mathsf r/MD_j)\le1$ holds at the levels $1\le j<j_1$, as $\mathsf r/MD_j\ge3$ there and $\mathcal T$ decreases (the monotonicity in \eqref{eq:weighted-mixed-numbers}). Let $\mathsf B_j=\sum_{i\le j}L^{-4(j-i)}\mathsf N_i(\operatorname{supp}w_1)\|\ddot E^{(i)}\|_{i,\kappa}$ collect the second jets of the levels up to $j$, at the rate $\kappa$.

By the two derivatives of the step, at the arrays of the first and of the second derivative, the second jets obey $\|\ddot E^{(j+1)}\|\le\nu^\sharp_j+\theta_jc_\sharp\mathsf B_j$, where $\nu^\sharp_j\le C\delta(x_j)^2d^6D_j^{-8}|w_1|_\infty|w_2|_\infty$ for $j\ge j_1$ comes from the product of the two first jets; and $\mathsf B_{j+1}=L^{-4}\mathsf B_j+\mathsf N_{j+1}(\operatorname{supp}w_1)\|\ddot E^{(j+1)}\|_{j+1,\kappa}$. The recursion has the same form as that of the first jets. The start: at a level $i<j_1$ a domain that holds plaquettes of both supports has a linear size at least $\mathsf r/MD_i-2$, so that $\|\ddot E^{(i)}\|_{i,\kappa}\le e^{-4\beta_{\rm s}\kappa[\mathsf r/MD_i-2]_+}\|\ddot E^{(i)}\|$ (equation~\YMMref{(2259)}{827}), with $\mathsf r/MD_i\ge3L^{j_1-1-i}$; with $\varpi_{i-1}\le\frac43L^n(1+n)^{2p_0}\varpi_{j_1-1}$ and $\mathsf N_i(\operatorname{supp}w_1)<18^3L^{3n}$, $n=j_1-i$, the powers of $L$ cancel and the series $\sum_{n\ge1}(1+n)^{2p_0}e^{-4\beta_{\rm s}\kappa[3L^{n-1}-2]_+}$ converges, so that $\mathsf B_{j_1}\le C\varpi_{j_1-1}\le C'(\log x_{j_1-1})^{2p_0}|w_1|_\infty|w_2|_\infty/(bD_{j_1})$, by $x_{j_1-1}\ge\frac34b/D_{j_1-1}$.

The quantity $L^{4(j-j_1)}\mathsf B_j$ grows from one level to the next by at most the factor $1+L^4\cdot18^3c_\sharp\theta_j$, beside the inhomogeneous terms $L^{4(j+1-j_1)}18^3\nu^\sharp_j$, whose sum over $j\ge j_1$ is at most 
\[
C(\log x_{j_1-1})^{2p_0}b^{-1}|w_1|_\infty|w_2|_\infty d^6D_{j_1}^{-7}\le C'(\log x_{j_1-1})^{2p_0}|w_1|_\infty|w_2|_\infty/(bD_{j_1})
,\]
as $\delta(x_j)^2\le\frac43A_1^2(\log x_j)^{2p_0}D_j/b$, $\sum_{j\ge j_1}D_j^{-3}\le D_{j_1}^{-3}/(1-L^{-3})$ and $d/D_{j_1}<\frac{15}2M$. So, as for the first jets, 
\[
\mathsf B_j\le C(\log x_{j_1-1})^{2p_0}b^{-1}|w_1|_\infty|w_2|_\infty d^3D_j^{-4}
\]
 for $j\ge j_1$, by $D_{j_1}\le\frac{Ld}8$; and $\|\ddot E^{(j+1)}\|\le\nu^\sharp_j+\theta_jc_\sharp\mathsf B_j$, with $\theta_j\le C(\log x_j)^{p_0}(D_j/b)^{1/2}$, gives 
\[
\|\ddot E^{(j)}\|\le C\mathcal y^{3p_0}[b^{-3/2}d^3D_j^{-7/2}+b^{-1}d^6D_j^{-7}]|w_1|_\infty|w_2|_\infty
\]
 for $j>j_1$, as $\log x_j\le\log x_{j_1-1}<\mathcal y$ for $j\ge j_1-1$, by $x_{j_1-1}\le\frac{12}{11}b/D_{j_1-1}$ (proof of Proposition~\ref{C-350-6}(ii)) and $D_{j_1-1}=D_{j_1}/L>\frac{2d}{15LM}$.

A level with $D_j>\frac{Ld}8\ge D_{j_1}$ lies above $j_1$, and the count of the domains that meet a support (Lemma~\ref{lem:source-support-sum}, at most $18^3$ cubes of the level $j$ holding a corner of a plaquette of $\operatorname{supp}w_1$) gives (ii).

(iii) By (i) the difference is at most $\sum_j(|\mathsf R_j[\ell_g]|+|\mathsf R_j[\ell_gL^s]|)$. On each torus the local second jets decay like $e^{-\kappa^+d_j(X)}$ in the linear size $d_j(X)$, with the rate $\kappa^+=(1+4\beta_{\rm s})\kappa$ of (ii) (Proposition~\YMMref{122.1}{826}); and the count of the domains of linear size at least $R$ that meet a support, $\sum_X\#(X\cap\mathbb T^{(j)})\,e^{-\kappa^+d_j(X)}\le c\,\nu\,e^{-R}$ over the domains $X$ with $d_j(X)\ge R$ that meet the $\nu$ cubes of the level $j$ in which a corner of a plaquette of the support lies (Lemma~\ref{lem:source-support-sum}, as $\kappa^+$ exceeds the number $\kappa_{\mathcal D}+2$ that it asks), taken at $R=A/D_j-2$, gives $|\mathsf R_j|\le C\|\ddot E^{(j)}\|e^{2-A/D_j}$, $A=\ell_g/(2M)\ge L$.

For the levels above $j_1$, (ii) and $\sum_{n\ge0}L^{7n}e^{-AL^n}\le\frac85e^{-A}$ give $C\mathcal y^{3p_0}b^{-1}d^3\eta$. For the levels $j\le j_1$ the second jets grow at most like $L^n(1+n)^{2p_0}$, $n=j_1-j$, by their general bound $C^\flat_1\varpi_{j-1}\le\frac43C^\flat_1L^n(1+n)^{2p_0}\varpi_{j_1-1}$ up to the level $j_1$ in the proof of (ii), against $e^{-A/D_j}\le e^{-8/d}e^{-8L^2(L^n-1)}$ and $e^{-8/d}\le\frac3{512}d^4$; this gives again at most $C\mathcal y^{2p_0}b^{-1}d^3\eta$. Double for the two tori.
\end{proof}

\begin{lemma}[Two sides at the order six]\label{C-350-26}
Let $0<g\le\gamma_{\rm s}$, $(K,\beta)$ a datum at $g$, $\ell\le\breve\ell$ in $\mathfrak L(g)$, $d\le L^{-1}$, $s(d)<K$, $20d\le\ell$, $(f,h)\in\mathfrak T_{T_\ell}(d;\vartheta)$, and $(\breve f,\breve h)$ the pair on $T_{\breve\ell}$ with the same lift. Then
\[
\begin{multlined}
|\mathsf h_6[\ell](f,h)-\mathsf h_6[\breve\ell](\breve f,\breve h)|\le|\mathsf J_2[\ell](f,h)-\mathsf J_2[\breve\ell](\breve f,\breve h)|+\Bigl[\bigl(2C_{\rm E}b^{-1}d\,l^{8p_0}+2\bar\varepsilon_4(g,d)\bigr)b^{-3}d \\
+3C_{\rm tm}\vartheta^{-1}b^{-2}\frac{a_K}d\Bigr]\eta .
\end{multlined}
\]
\end{lemma}

\begin{proof}[Proof sketch]
On each side, $\mathsf h_6=\mathsf J_2-\mathsf h_4-[\mathsf J_2-\mathsf h_4-\mathsf h_6]$. The bracket is at most $C_{\rm E}b^{-1}dl^{8p_0}\cdot b^{-3}d\,\eta$ on each side (Proposition~\ref{C-350-7}(ii)), and $\mathsf h_4$ is $c_2b^{-2}\mathcal K_6$ up to $[\bar\varepsilon_4b^{-3}d+C_{\rm lat}b^{-2}a_K/d]\eta$ on each side (Proposition~\ref{C-350-9}, with one $\bar\varepsilon_4$ for all sides and data at $g$); $\mathcal K_6$ and $\eta$ are the same for the two pairs. Subtract.
\end{proof}

\begin{proposition}[Two lattices with one bare torus]\label{C-350-27}
Let $0<g\le\gamma_{\rm s}$; $D_0=(K_0,\beta_0)$ a datum at $g$ with label $b$; $d\le L^{-1}$, $s(d)<K_0$; $(f,h)\in\mathfrak T_{T_{\ell_g}}(d;\vartheta)$, $81d\le\ell_g$, $r=dL^{K_0}$; $b_*\in I(g)$; and integers $K'\ge K_0$ and $2\le n_0\le n_1$ with $n_0\le K_0$, $n_1\le K'$, $L^{n_0}\ge r$. Put $\breve\ell=\ell_gL^{K'-K_0}$, $N'=\ell_gL^{K'}=\breve\ell L^{K_0}$, $D'=(K',L^{K'}b_*)$, $l'=\max\{1,\log(2b_*L^{K'}/r)\}$, $\mathsf N'=N'+\max\{\beta_0,L^{K'}b_*\}$, and let $(\breve f,\breve h)$ be the pair on $T_{\breve\ell}$ with the same lift. Let the two conditions of Proposition~\ref{C-350-11}(ii) hold at the three places $(D_0,\breve\ell;n_0)$, $(D',\ell_g;n_0)$ and $(D',\ell_g;n_1)$, and let $x_\flat$ denote the running coupling $x_{n-1}$ of the place; and let $\sigma\ge0$ be a number with $|\mathsf J_2[\ell_g](f,h)-\mathsf J_2[\breve\ell](\breve f,\breve h)|\le\sigma b^{-3}d\eta$ at $D_0$. Then, with $\tilde w$ the lattice sources of the lift on $\mathbb Z^3$,
\[
\begin{multlined}
\bigl|\mathsf H^{(K_0)}[D_0,\ell_g]-\mathsf T^{\rm lat}(\tilde w)\bigr|\le\Bigl\{\bigl[\sigma+2C_{\rm E}b^{-1}dl^{8p_0}+2\bar\varepsilon_4(g,d)+C_{\rm R}(l^{6p_0}+l'^{6p_0})(r/L^{n_0})^{3/2}\bigr]r \\
+\frac{3C_{\rm tm}b}{\vartheta d}+\sum_{\rm three\ \relax places}\mathcal U(x_\flat,\mathsf N')+2700C_\star(\ell_gL^{K'-n_1})^{19}n_1L^{-n_1/2}r^6\Bigr\}\eta .
\end{multlined}
\]
\end{proposition}

\begin{proof}[Proof sketch]
If $\tilde w_1$ or $\tilde w_2$ vanishes, both numbers are $0$. Consider three places: (1) the datum $D_0$ on $T_{\ell_g}$ with $(f,h)$; (2) the datum $D_0$ on $T_{\breve\ell}$ with $(\breve f,\breve h)$, where $\breve\ell\in\mathfrak L(g)$; (3) the datum $D'$ on $T_{\ell_g}$ with $f'=\breve f\circ\mathsf D^{-1}$, $h'=\breve h\circ\mathsf D^{-1}$, $\mathsf D:T_{\breve\ell}\to T_{\ell_g}$ the dilation $x\mapsto L^{-(K'-K_0)}x$. $D'$ is a datum at $g$, as $b_*\in I(g)$; $f',h'$ are real and bounded, with the norms of $f,h$, with supports in balls of radius $d'=rL^{-K'}$ at a distance in $[d',6d']$, and $8a_{K'}\le d'$ as $r\ge8$; so Proposition~\ref{C-350-6} applies at the three places, as it asks no depth.

$\mathsf D$ carries the bare lattice of $(\breve\ell,K_0)$ onto that of $(\ell_g,K')$, plaquettes to plaquettes and barycentres to barycentres: in bare units the sources of (2) and of (3) are one pair of functions on the plaquettes of $\Lambda_{N'}$, the periodisation of $\tilde w$, and $r$ is the same at the three places. The chain has seven links. 1. $\mathsf H^{(K_0)}[D_0,\ell_g]\to\mathsf H^{(K_0)}[D_0,\breve\ell]$: Lemma~\ref{C-350-26} with $\sigma$, multiplied by $\beta_0^3a_{K_0}^2=b^3/a_{K_0}$. 2. $\to\mathsf H^{(n_0)}[D_0,\breve\ell]$: Proposition~\ref{C-350-7}(iii) at (2). 3. $\to\mathsf T[N',n_0]$: Proposition~\ref{C-350-11}(i), (ii), with Lemma~\ref{C-350-23}, and Proposition~\ref{C-350-13} at (2), as for the link X of Step~3 below. 4.

$\to\mathsf H^{(n_0)}[D',\ell_g]$: the same at (3); the number $\mathsf T[N',n_0]$ is one, as it depends on $N'$, $n_0$, $L$ and the lattice sources alone, and the two runs that meet in it leave no trace, by Proposition~\ref{C-350-12}(iii). 5. $\to\mathsf H^{(n_1)}[D',\ell_g]$: Proposition~\ref{C-350-7}(iii) at (3), where $l$ is $l'$ and $d'/D'_{n_0}=r/L^{n_0}$, $D'_{n_0}=a_{K'-n_0}$ being the spacing of the level $n_0$ of the lattice $(\ell_g,K')$. 6. $\to\mathsf T[N',n_1]$: Proposition~\ref{C-350-11}(i), (ii), with Lemma~\ref{C-350-23}, and Proposition~\ref{C-350-13} at (3). 7. $\to\mathsf T^{\rm lat}(\tilde w)$: Proposition~\ref{C-350-15} on the lattice with $\ell'=N'/L^{n_1}=\ell_gL^{K'-n_1}$ top sites and $n_1$ levels, where $\ell'\ge\ell_g\ge N_{\mathsf H}$, $\ell'L\ge C_LM$, and all base points lie within $10r+1$ of each other, $10r+1\le\frac18N'$ as $81d\le\ell_g$ and $r\ge8$; at most $900r^6$ pairs.

In link~1, $b^3a_{K_0}^{-1}\cdot b^{-3}d=r$ and $b^3a_{K_0}^{-1}\cdot3C_{\rm tm}\vartheta^{-1}b^{-2}a_{K_0}/d=3C_{\rm tm}b/(\vartheta d)$. Add the errors of the seven links.
\end{proof}

\begin{lemma}[The splines of Definition~\ref{C-350-2}]\label{C-350-29}
Let $q,q'\ge1$ be odd.
\begin{enumerate}
\item[(i)] $N_3(t)=\frac34-t^2$ for $|t|\le\frac12$, $N_3(t)=\frac12(\frac32-|t|)^2$ for $\frac12\le|t|\le\frac32$, and $N_3(t)=0$ beyond; $\Phi$ is of class $C^1$, $0\le\Phi\le(\frac34)^3$ and $\{\Phi\ne0\}=(-\frac32,\frac32)^3$, so $|x|<2.6$ on $\operatorname{supp}\Phi$; $\int\Phi=1$, $\int\Phi x_\mu=0$ and $\int\Phi|x|^2=\frac34$.
\item[(ii)] $\Phi(x/q)=\sum_jc^q_j\Phi(x-j)$, with $c^q_j\ge0$ and $c^q_j=0$ unless $|j_\mu|\le\frac32(q-1)$ for all $\mu$; $\sum_jc^q_j=q^3$, $\sum_jc^q_jj=0$, $\sum_jc^q_j|j|^2=\frac34q^3(q^2-1)$; hence $q^{-1}\sum_jc^q_jP(\frac{y+j}q)=P(y)$ for all $y$.
\item[(iii)] For $q\ge3$, $\Upsilon^{(q)}$ is of class $C^1$ with support in $\{7\le|y|\le11q\}$, and the distance of the supports of $\Phi$ and $\Upsilon^{(q)}$ exceeds $4.4$; $0\le\Upsilon^{(q)}\le64q^2$, and $\Upsilon^{(q)}=63.25$ on $|y|=8$.
\item[(iv)] $e(u)=|u|^2$ for $|u|\le4.4$, $e(u)=0$ for $|u|\ge10.6$, $0\le e(u)\le|u|^2$ otherwise, and $\int\Phi(x)\Upsilon^{(q)}(x+u)\,dx=q^2e(u/q)-e(u)$.
\item[(v)] $\mathsf k_6\in[\frac{4\pi}{10.6},\frac{4\pi}{4.4}]$, and for $q\ge3$, on $\mathbb R^3$: $\mathcal K_5(\Phi,\Upsilon^{(q)})=4\pi\log q$ and $\mathcal K_6(\Phi,\Upsilon^{(q)})=(1-q^{-1})\mathsf k_6$.
\item[(vi)] $\Upsilon^{(qq')}(y)=\Upsilon^{(q)}(y)+q^{-1}\sum_jc^q_j\Upsilon^{(q')}(\frac{y+j}q)$.
\end{enumerate}
\end{lemma}

\begin{proof}[Proof sketch]
(i) $N_3$ is the density of the sum of three independent variables uniform on $[-\frac12,\frac12)$: integral $1$, mean $0$, variance $\frac14$. (ii) The intervals $[i-\frac12,i+\frac12)$, $i\in J_q$, tile $[-\frac q2,\frac q2)$, so $\mathbf 1_{[-1/2,1/2)}(\cdot/q)=\sum_{i\in J_q}\mathbf 1_{[-1/2,1/2)}(\cdot-i)$, and the threefold convolution gives $q^2N_3(x/q)=\sum_{i_1,i_2,i_3\in J_q}N_3(x-i_1-i_2-i_3)$; the moments of the $c^q_j$ are those of the sum of three variables uniform on $J_q$, and the identity for $P$ follows by expanding the square. (iii) If $|y|\le4.4q+2.6$ and $c^q_j\ne0$, then $|y+j|\le7q$, so all the $\Xi(\frac{y+j}q)$ are $P(\frac{y+j}q)$, and by (ii) $\Upsilon^{(q)}(y)=P(y)(1-\zeta(|y|))$, which is $0$ for $|y|\le7$ and $\ge0$ beyond.

If $|y|\ge4.4q+2.6$, then $\Xi(y)=0$ and $|y+j|/q>1.8$, where $\Xi\ge0$. If $|y|>10.6q-2.6$ all the terms vanish. The upper bound follows from $\Xi\le64-\frac34$. (iv) For $|u|\le4.4$, $|x+u|<7$ on $\operatorname{supp}\Phi$, and (i) gives $e(u)=\frac34+|u|^2-\frac34$; the last formula is (ii) after the change of variables $x\mapsto x-j$. (v) With $\varepsilon(t)=t^{-2}\int_{S^2}e(t\omega)\,d\omega$, which is $4\pi$ on $(0,4.4]$ and $0$ on $[10.6,\infty)$, polar coordinates and (iv) give $\mathcal K_p(\Phi,\Upsilon^{(q)})=\int_0^\infty t^{4-p}[\varepsilon(t/q)-\varepsilon(t)]\,dt$. At $p=5$ the integral over $t\ge\delta$ is $\int_{\delta/q}^\delta\varepsilon(s)\frac{ds}s=4\pi\log q$ for small $\delta$; at $p=6$ the substitution $t=qs$ gives $(1-q^{-1})\mathsf k_6$.

(vi) $(i,i')\mapsto i+qi'$ is a bijection of $J_q\times J_{q'}$ onto $J_{qq'}$, so $c^{qq'}_J=\sum_{j+qj'=J}c^q_jc^{q'}_{j'}$.
\end{proof}

\subsection{Proof of Theorem~\ref{C-350-3}}

\emph{Step 1: the front and the expansion.} We take $\gamma_{\rm e}=\min\{\gamma_{\rm s},\gamma_{\rm C}\}$. The number $d_{\rm e}$ is the least of $L^{-13}$, $\frac1{200}$, $d_{\rm A}(g)$, $d_{\rm K}$ of Proposition~\ref{C-350-16} and of the thresholds of Steps~2, 4 and~8: the smallness of $d$ needed for $2\le n_0<n_1$ and $12r+2\le N_0/8$, for the two conditions of Proposition~\ref{C-350-11}(ii) at the places of the chain, and for the three inequalities on $d$ stated at the head of Step~8. Let $D=(K,\beta)$ be a datum at $g$ with $s(d)<K$ and $L^K\le d^{-P}$, and $(f,h)\in\mathfrak T_{T_\ell}(d;\vartheta)$; then $K\ge3$, $8a_K\le d$ and $r\ge L^{s(d)}d\ge C_W(s(d))\ge(\log\frac1d)^2\ge8$. If $w_1$ or $w_2$ vanishes, all the jets vanish and only the last link of Step~3 is used. Propositions~\ref{C-350-8}, \ref{C-350-7}(ii) and~\ref{C-350-9}, the last two at every translate $(\tau_vf,\tau_vh)$, $v\in\mathsf{Pk}$, give, as $L^+_b$ depends on distances alone,
\begin{equation}\label{C-350-18}
|S^T_2-L_b^+|(f,h)\le\Bigl[\sum_{i=1}^3E_i\,b^{-3}d+C_{\rm lat}b^{-2}\frac{a_K}d\Bigr]\eta+b^{-3}a_K\max_{v\in\mathsf{Pk}}\bigl|\mathsf H^{(K)}[D]-\varkappa\mathcal K_5/a_K\bigr|(\tau_vf,\tau_vh),
\end{equation}
with $E_1=C_{\rm A}(g)\vartheta^{-1}bd$, $E_2=C_{\rm E}b^{-1}dl^{8p_0}$ and $E_3=\bar\varepsilon_4(g,d)$; here $\mathsf h_6=b^{-3}a_K\mathsf H^{(K)}[D]$, and $b^{-3}a_K\cdot r\eta=b^{-3}d\eta$. It remains to bound $|\mathsf H^{(K)}[D]-\varkappa\mathcal K_5/a_K|$ in multiples of $r\eta$ at every position; we write $(f,h)$ for the translate.

\emph{Step 2: the levels and a second datum.} Put
\[
K'=\max\{K,\lceil24\log_Lr\rceil\},\quad N_0=\ell L^K,\quad N'=\ell L^{K'},\quad n_1=K'-\lceil K'/78\rceil,
\]
so that $r^{24}\le L^{K'}\le\max\{d^{-P},Lr^{24}\}\le Ld^{-24P}$. If $K'=K$ (Case A) let $D'=D$. If $K'>K$ (Case B), then $r^{24}>L^K=r/d$, that is $r>d^{-1/23}$; fix $b_*\in I(g)$ once, let $D'=(K',L^{K'}b_*)$, a datum at $g$ (Corollary~\YMMref{74.2}{616}), and let $n_0$ be the level with $D_{n_0}=a_{K-n_0}\in[d^{1/2},Ld^{1/2})$. Then $L^{n_0}\in[rd^{-1/2},Lrd^{-1/2})$; $2\le n_0\le K$ for $d\le L^{-2}$, as $L^{n_0}\ge8L$ and $D_{n_0}<1$; and $n_0<n_1$ for small $d$, as $L^{n_0}<Lrd^{-1/2}<Lr^{12.5}$ and $L^{n_1}\ge L^{77K'/78-1}\ge r^{23.6}/L$. In bare units the base points of the plaquettes of the two supports lie within $12r+2$ of each other, and $12r+2\le N_0/8=\ell r/(8d)$ for small $d$.

The lifts $\breve w_a$ of Proposition~\ref{C-350-14} on $\Lambda_{N'}$ are the lattice sources, at $D'$, of real bounded functions on $T_\ell$ with supports in balls of radius $d'=ra_{K'}$ at a distance $\ge d'$, $8a_{K'}\le d'$; Proposition~\ref{C-350-6} applies to them at $D'$, as it puts no condition on the depth, the signs or the Lipschitz constants. The unit $r\eta$ is the same at $D$ and at $D'$, since $d'L^{K'}=r$.

\emph{Step 3: the chain.} In Case B,
\begin{align}
\mathsf H^{(K)}[D]&\xrightarrow{\rm R}\mathsf H^{(n_0)}[D]\xrightarrow{\rm X}\mathsf T[N_0,n_0](w)\xrightarrow{\rm W}\mathsf T[N',n_0](\breve w)\xrightarrow{\rm X}\mathsf H^{(n_0)}[D']\nonumber\\
&\xrightarrow{\rm R}\mathsf H^{(n_1)}[D']\xrightarrow{\rm X}\mathsf T[N',n_1](\breve w)\xrightarrow{\rm A}a_K^{-1}\mathsf P^{\rm lat}_1(\tilde f,\tilde h;a_K)\xrightarrow{\rm T}\varkappa\mathcal K_5(f,h)/a_K .\label{C-350-19}
\end{align}
Here R is Proposition~\ref{C-350-7}(iii). X is Proposition~\ref{C-350-11}(i), (ii), with Lemma~\ref{C-350-23}, and Proposition~\ref{C-350-13}: $\mathsf H^{(n)}[D]=\beta^3a_K^2\mathsf h^{(n)}_6=\beta^3a_K^2\gamma^{\rm r}_2$ by Proposition~\ref{C-350-11}(i), the cumulants being taken in $\mathsf D^0dB=\bigotimes_{0\le i<n}\nu^{\rm r}_i$; $\beta^3a_K^2\gamma^\circ_2=\mathsf T[N,n](w)$ by Proposition~\ref{C-350-13}; and $\beta^3a_K^2|\gamma^{\rm r}_2-\gamma^\circ_2|\le\mathcal U(x_{n-1},\mathsf N)\eta$ by Proposition~\ref{C-350-11}(ii) under the two conditions there, its proof reading the bounds of Lemma~\ref{C-350-23} on the boxes and the precisions $\mathsf M_i\ge\frac34\gamma^{\rm tr}$ of the Gaussian measures of the heights.

So $|\mathsf H^{(n)}[D]-\mathsf T[N,n](w)|\le\mathcal U(x_{n-1},\mathsf N)\eta$ under the two conditions; the two middle members $\mathsf T[N',n_0](\breve w)$ are one number, as it depends on no datum. W is Proposition~\ref{C-350-14} at $t=n_0$, $\nu=\ell/D_{n_0}>\ell d^{-1/2}/L\ge\nu_{\rm W}$, with the base points within $12r+2\le N_0/8$. A is Proposition~\ref{C-350-15} on the lattice with $\ell'=N'/L^{n_1}=\ell L^{\lceil K'/78\rceil}$ top sites and $n_1$ levels: $\ell'\ge\ell\ge N_{\mathsf H}$ and $\ell'L\ge C_LM$ by Proposition~\ref{C-350-6}(i), the base points lie within $12r+2\le\frac18N'-1$, and the sum runs over at most $900r^6$ pairs; with $\frac{19}{78}+\frac1{156}-\frac12=-\frac14$ and $r^5\le L^{5K'/24}$ the error is $E_{10}$ below, and $a_K\mathsf T^{\rm lat}(\tilde w)=\mathsf P^{\rm lat}_1(\tilde f,\tilde h;a_K)$ for the lift $(\tilde f,\tilde h)\in\mathfrak T_{\mathbb R^3}(d;\vartheta)$.

T is Corollary~\textup{\ref{IV-73-73}} for the lift: $|\mathsf P_1^{\rm lat}(\tilde f,\tilde h;a)-\kappa^{\rm lat}\mathcal K_5(f,h)|\le C_{\rm T}\vartheta^{-1}a\lg(d/a)\eta$ for $0<a\le d$, with $\kappa^{\rm lat}=\varkappa=c_1c_2$ (Theorems~\textup{\ref{IV-73-98}}, \textup{\ref{IV-73-100}}); only here enters $c_1=\kappa^{\rm lat}/c_2$, which depends on no $L$, $g$ or label. The errors of the links, in multiples of $r\eta$, are
\begin{align}
E_4&=C_{\rm R}l^{6p_0}d^{3/4}, & E_5&=\mathcal U(x_{n_0-1},\mathsf N_0)/r, & E_6&=900r^5C'_{\rm W}N'^{237}e^{-\delta_{\rm tp}\ell d^{-1/2}/48},\nonumber\\
E_7&=\mathcal U(x'_{n_0-1},\mathsf N')/r, & E_8&=C_{\rm R}l'^{6p_0}d^{3/4}, & E_9&=\mathcal U(x'_{n_1-1},\mathsf N')/r,\nonumber\\
E_{10}&=2700C_\star\ell^{19}L^{20}K'L^{-K'/24}, & E_{11}&=C_{\rm T}\vartheta^{-1}\lg(r)/r,\label{C-350-20}
\end{align}
where $x$, $x'$ are the running couplings of $D$, $D'$; $\mathsf N_0=N_0+\beta\le(\ell+b_+)d^{-P}$ and $\mathsf N'=N'+L^{K'}b_*\le(\ell+b_+)Lr^{24}$; by Proposition~\ref{C-350-6}(ii), $x_{n_0-1}\ge\frac34bL/D_{n_0}\ge\frac34d^{-1/2}$, $x'_{n_0-1}\ge\frac34r^{23}d^{1/2}\ge\frac34\max\{r^{11.5},d^{-1/2}\}$ and $x'_{n_1-1}\ge L^{K'/78}$; $l'\le\log(2b_+)+K'\log L$; and $(r/L^{n_0})^{3/2}\le d^{3/4}$. In Case A the chain is $\mathsf H^{(K)}[D]\to\mathsf H^{(n_1)}[D]$ by R, with the error $C_{\rm R}l^{6p_0}L^{\frac32(1-0.94K)}$, as $r/L^{n_1}\le L^{1-0.94K}$ and $l\le\log(2b_+)+K\log L$, followed by the last three links with $D'=D$.

\emph{Step 4: clause (a).} At each place of X the two conditions of Proposition~\ref{C-350-11}(ii) hold once $\frac18\gamma^{\rm tr}A_1^2(\log x_\flat)^{2p_0}\ge10+\log(10240\mathsf N^3)$, $x_\flat$ being the lower bound of the running coupling at that place in Step~3; and $\log\mathsf N\le C_1+C_2\log x_\flat$ with $C_1=\log((\ell+b_+)L)+P$ and $C_2=2P$ at the first place ($\log\frac1d\le2\log x_\flat+0.6$), $C_2=3$ at the second ($\log x_\flat\ge11\log r$ for $r\ge2$), $C_2=78$ at the third ($\log x_\flat\ge\frac{K'}{78}\log L$). As $2p_0\ge2$,
\[
\mathcal U(x_\flat,\mathsf N)\le C_{\rm X}\exp\bigl[36C_1+36C_2\log x_\flat-\tfrac18\gamma^{\rm tr}A_1^2(\log x_\flat)^{2p_0}\bigr]
\]
tends to $0$ faster than every power of $x_\flat^{-1}$, and $x_\flat\to\infty$ uniformly as $d\to0$: $x_\flat\ge\frac34d^{-1/2}$ at the first two places and $x_\flat\ge L^{s(d)/78}$ at the third. Further $E_6\le900C'_{\rm W}(\ell L)^{237}d^{-5P-5688P}e^{-\delta_{\rm tp}\ell d^{-1/2}/48}$; $E_4,E_8\le C_{\rm R}[\log(2b_+L)+24P\log\frac1d]^{6p_0}d^{3/4}$; $E_{10}$ and the error of Case A tend to $0$ as $K'\ge K>s(d)\to\infty$; and $E_{11}\le C_{\rm T}\vartheta^{-1}\sup_{\rho\ge(\log1/d)^2}\lg(\rho)/\rho$. As $b\in[g^{-2},b_+]$, each $E_i$ is bounded by a function of $(g,\vartheta,P,d)$ alone that tends to $0$ with $d$. With $\bar\varepsilon_{\rm a}(d)$ the supremum over $d''\le d$ of the sum of these bounds, \eqref{C-350-18} gives (a).

\emph{Step 5: clause (b) and the rate.} Let $b\in I(g)$; every $(K,L^Kb)$ is a datum at $g$, and $I(g)\subset[\gamma^{-2},L^2(g^{-2}+\bar B)]$. For $s(d)<K\le K_\sharp(d)$, (a) applies, as $L^K\le d^{-3601}\le d^{-P}$ (Proposition~\ref{C-350-16}). For $K>K_\sharp=K_\sharp(d)$, and for the limit, Proposition~\ref{C-350-16} and (a) at $K_\sharp$ give, with $a_{K_\sharp}\le(d/\sqrt3M)^{3600}$,
\[
|S^T_{2;K,\ell,L^Kb}-L_b^+|(f,h)\le\bigl[\bar\varepsilon_{\rm a}(d)+C_\diamond\vartheta^{-2}b_+^3d^{177}+C_{\rm lat}b_+d^{3598}\bigr]b^{-3}d\,\eta ,
\]
so that $\bar\varepsilon(d)=\bar\varepsilon_{\rm a}(d)+\sup_{d''\le d}[C_\diamond\vartheta^{-2}b_+^3d''^{177}+C_{\rm lat}b_+d''^{3598}]$ serves in (a) and (b). For the limit we apply (a) at the single cutoff $K_\sharp$, where $r\ge d^{-3599}$ and Case B holds: $E_5$, $E_6$, $E_7$, $E_9$ are below every power of $d$; $E_{10},E_{11}\le Cd^{3598}\log\frac1d$; $E_1+E_2\le Cd(\log\frac1d)^{8p_0}$; $E_4+E_8\le Cd^{3/4}(\log\frac1d)^{6p_0}$; and $E_3\le C_4(g)d^{1/2}$ (Proposition~\ref{C-350-9}). So the error of the limit is at most $c_{1/2}d^{1/2}b^{-3}d\,\eta$; its slowest member is that of the order four.

\emph{Step 6: clause (c).} Let $g_{\rm a}=L^{1/4}g$ and
$g_{\rm b}=L^{-1/4}g$. For $g''\in\{g,g_{\rm a},g_{\rm b}\}$,
$g''\le\min\{\gamma_{\rm e},\gamma^{\rm sd}_W\}$, and
$\mathsf x:=g^{-2}\ge L^{1/2}(\gamma^{\rm sd}_W)^{-2}
\ge4L^{1/2}\bar B\ge5\bar B$. All these reference couplings use
the single assignment fixed at the beginning of the section.
The real constants are its unchanged restriction; hence the canonical
flow and every sourced construction are evaluated at the same values.
This is the common-choice convention of Section~\YMMref{82}{665},
with its radius conditions imposed uniformly for $x\ge\gamma^{-2}$.

So there is one flow of the coupling: at every coupling the running couplings of a datum, which depend on $(K,\beta)$ alone, the data, the histories of the sourced chain and the torus measures are the same objects in the real restriction of the common assignment and in its sourced construction, as fixed in Section~\YMMref{82}{665}; and $C_{W,0}$ and $P$ are fixed before the couplings. Hence (a) and (b), which hold at this choice, hold at the three couplings $g''$, and these clauses and Proposition~\ref{C-350-21} speak of the same objects. The labels at $g$ lie in $[\mathsf x,L(\mathsf x+\bar B)]$, and $I(g)$ leaves out $[\mathsf x,\mathsf x+\bar B]$ and $(L(\mathsf x+B_\sharp)-\bar B,L(\mathsf x+\bar B)]$.

As $\mathsf x\ge5\bar B$ and $L\ge13$,
\[
\begin{multlined}
L^{-1/2}\mathsf x+\bar B<\mathsf x,\quad\mathsf x+\bar B\le L^{1/2}\mathsf x+LB_\sharp-\bar B,\quad L^{1/2}\mathsf x+\bar B\le L(\mathsf x+B_\sharp)-\bar B,\quad L(\mathsf x+\bar B) \\
\le L(L^{1/2}\mathsf x+B_\sharp)-\bar B,
\end{multlined}
\]
which place the first interval in $I(g_{\rm a})=(L^{-1/2}\mathsf x+\bar B,L^{1/2}\mathsf x+LB_\sharp-\bar B]$ and the second in $I(g_{\rm b})=(L^{1/2}\mathsf x+\bar B,L^{3/2}\mathsf x+LB_\sharp-\bar B]$. Let $d'_{\rm e}$ be the least, over the three couplings $g''$, of $\frac1{180}$, $L^{-2}$, $d_{\rm e}(g'')$, $\mathsf d_\star(g'',\vartheta)$ and $(LC_w(g''))^{-2}$, where $C_w(g'')=\sup_{s\ge0}C_W(s)L^{-s/2}$ is formed at $g''$; for $d\le(LC_w(g''))^{-2}$ the witness depth at $g''$ has $L^{s(d)}\le d^{-3}$, since $L^{s-1}d<C_W(s-1)\le C_wL^{(s-1)/2}$ for $s=s(d)\ge1$. Let $(K,\beta)$ be a datum at $g$ with $K>s(d)$. If $L^K\le d^{-P}$, (a) applies; if $b\in I(g)$, (b) applies.

Otherwise $b\in I(g_\bullet)$ with $g_\bullet\in\{g_{\rm a},g_{\rm b}\}$, and $(K,\beta)=(K,L^Kb)$ is a datum at $g_\bullet$ (Corollary~\YMMref{74.2}{616}), with $K\ge3$, $a_K\le d$ and $K>s(d)$ at $g_\bullet$, since $L^K>d^{-P}\ge d^{-3}$. Clause (b) at $g_\bullet$ gives the law on $T_{\ell_{g_\bullet}}$ for the pair with the same lift as $(f,h)$, which has the same $\eta$, $\mathcal K_5$ and $\mathcal K_6$. Of the sides $\ell_g$, $\ell_{g_\bullet}$, both of the form $2L^m$, let $g_\circ\in\{g,g_\bullet\}$ be the coupling of the smaller; the larger lies in $\mathfrak L(g_\circ)$, so both torus measures are states at $g_\circ$ at the datum $(K,\beta)$, one of them on the least torus, and \eqref{C-350-17} at $g_\circ$ bounds their difference at the lift by
\[
C^{\rm T}_2(g_\circ)\vartheta^{-1}d^2b^{-2}\eta\le C^{\rm T}_2(g_\circ)\vartheta^{-1}b_+d\cdot b^{-3}d\,\eta .
\]
So (c) holds with $\bar\varepsilon'(d)=\max_{g''}\bar\varepsilon_{g''}(d)+\max_{g''}C^{\rm T}_2(g'')\vartheta^{-1}b_+d$, $\bar\varepsilon_{g''}$ being the function of (a) and (b) at $g''$.

\emph{Step 7: clause (d).} (1) By Lemma~\ref{C-350-29}(ii), $\Phi(x/q)=\sum_jc^q_j\Phi(x-j)$; so every bilinear form $\mathsf S$ invariant under the translations by $\rho\mathbb Z^3$ satisfies
\[
(q\rho)^{-1}\mathsf S(\Phi_{q\rho},\Xi_{q\rho})-\rho^{-1}\mathsf S(\Phi_\rho,\Xi_\rho)=\rho^{-1}\mathsf S(\Phi_\rho,\Upsilon^{(q)}_\rho).
\]
At a cutoff $K\ge k_0+n$ the covariance $S^T_{2;K,\ell,L^Kb}$ is bilinear and invariant under the translations by $a_K\mathbb Z^3\supset\varrho_k\mathbb Z^3$, $\varrho_k=L^{-k}r_0$, whatever the supports (Setting and notation). Taken at $\rho=\varrho_k$, $q=L$, summed over $1\le k\le n$, and compared with the identity at $\rho=r_1$, $q=L^n$, it gives
\[
r_1^{-1}S(\Phi_{r_1},\Upsilon^{(L^n)}_{r_1})=\sum_{k=1}^n\varrho_k^{-1}S(\Phi_{\varrho_k},\Upsilon^{(L)}_{\varrho_k});
\]
these pairs have disjoint supports (Lemma~\ref{C-350-29}(iii)), which lie in balls of radius $11r_0<\ell/4$, and $K\to\infty$ is taken member by member, by the existence and the bilinearity of the limit. (2) By Lemma~\ref{C-350-29}(v), $\mathcal K_5(\Phi,\Upsilon^{(q)})=4\pi\log q$ and $\mathcal K_6(\Phi,\Upsilon^{(q)})=(1-q^{-1})\mathsf k_6$; with $\mathcal K_q(f_\rho,h_\rho)=\rho^{6-q}\mathcal K_q(f,h)$ and $4\pi c_1c_2=\frac6{\pi^2}$, the subtracted members of $\Delta_n$ are $\sum_{k\le n}\varrho_k^{-1}L^+_b(\Phi_{\varrho_k},\Upsilon^{(L)}_{\varrho_k})$, so that
\[
\Delta_n=\sum_{k=1}^n\delta_k,\qquad\delta_k(b)=\varrho_k^{-1}\bigl[S_2^{(\ell,b)}-L^+_b\bigr](\Phi_{\varrho_k},\Upsilon^{(L)}_{\varrho_k}).
\]
(3) The tents $\chi_z(x)=\prod_\mu\max\{0,1-|2x_\mu-z_\mu|\}$, $z\in\mathbb Z^3$, form a partition of unity and split $\Phi$ and $\Upsilon^{(L)}$, which is $\ge0$ (Lemma~\ref{C-350-29}(iii)), into finitely many nonnegative Lipschitz pieces. Each pair of pieces $(\chi_z\Phi,\chi_{z'}\Upsilon^{(L)})$, neither $0$, lies in $\mathfrak T_{\mathbb R^3}(d_{zz'};\vartheta_0)$ with $d_{zz'}=\max\{1,\frac16\operatorname{dist}\}\in[1,2L]$, the distance of the two supports lying in $(4.4,11L+2.6)$, and with $\vartheta_0\in(0,\frac12]$ depending on $L$ alone; scaled by $\varrho_k$ and read on $T_\ell$, it lies in $\mathfrak T_{T_\ell}(\varrho_kd_{zz'};\vartheta_0)$, since all the points of its two supports lie within $22L\varrho_k\le22r_0<\ell/2$ of each other and $20\varrho_kd_{zz'}\le40r_0<\ell$. Clause (b) at $(g,\vartheta_0,P_\sharp)$, with its $d^0_{\rm e}$ and $c^0_{1/2}$, gives for $2L\varrho_k\le d^0_{\rm e}$
\[
|\delta_k(b)|\le C_\Psi c^0_{1/2}(2L\varrho_k)^{1/2}b^{-3},\qquad C_\Psi=\sum_{z,z'}d_{zz'}\|\chi_z\Phi\|_\infty\|\chi_{z'}\Upsilon^{(L)}\|_\infty,
\]
$C_\Psi$ depending on $L$ alone. (4) For the finitely many $k$ with $2L\varrho_k>d^0_{\rm e}$, $\delta_k$ is continuous on the compact set $\overline{I(g)}\subset(0,\infty)$, by the continuity of the limit in $b$, so it is bounded there. (5) As $b^{-3}<g^6$ on $I(g)$, $\sum_k\sup_{I(g)}|\delta_k|<\infty$; so $\Delta_n\to\Delta_\infty=\sum_k\delta_k$ uniformly on $I(g)$, $\Delta_\infty$ is bounded and continuous, $\sup_{I(g)}|\Delta_n|\le\sum_k\sup_{I(g)}|\delta_k|$, and for $n$ beyond the coarse scales
\[
\sup_{I(g)}|\Delta_n-\Delta_\infty|\le C_\Psi g^6c^0_{1/2}(2Lr_0)^{1/2}\frac{L^{-(n+1)/2}}{1-L^{-1/2}} .
\]
Of the rate in (b) the argument uses only that the tolerances at the scales $2L\varrho_k$ are summable over $k$; a tolerance of the order $1/\log(1/d)$ would not give it.

\emph{Step 8: clause (e).} Here the levels are cut at the scale $d^{1/5}$, and the bare torus is changed through a change of the side at one datum. Let $\Lambda=\log\frac1d$, $l_+=\log(2b_+)+\Lambda$, $l'_+=\log(2b_+L)+23(P_\sharp-1)\Lambda$, $A_\ell=2700C_\star\ell_g^{19}L^{39/2}$,
\[
\mathsf N_{\max}(d)=(\ell_g+b_+)Ld^{-24(P_\sharp-1)},\qquad E_{\rm X}(d)=C_{\rm X}\mathsf N_{\max}(d)^{36}\exp\bigl[-\tfrac18\gamma^{\rm tr}A_1^2(\Lambda/78)^{2p_0}\bigr],
\]
and recall that the number $d_{\rm e}\le\frac1{200}$ of Step~1 is so small that for $0<d\le d_{\rm e}$: $\gamma^{\rm tr}A_1^2(\Lambda/78)^{2p_0}\ge80$, $10240\,\mathsf N_{\max}(d)^3\exp[-\frac18\gamma^{\rm tr}A_1^2(\Lambda/78)^{2p_0}]\le1$, and
\[
\epsilon_{\rm r}(d):=g^2d\bigl[A_\ell\bigl(1+24(P_\sharp-1)\Lambda/\log L\bigr)+C_{\rm T}\vartheta^{-1}\bigl(1+(P_\sharp-1)\Lambda\bigr)\bigr]\le1 .
\]
Let $0<d\le d_{\rm e}$, $(f,h)\in\mathfrak T_{T_\ell}(d;\vartheta)$, and let $(K,\beta)$ be a datum at $g$ with $s(d)<K\le K_\sharp(d)$; for the first assertion, $\beta=L^Kb$ with $b\in I(g)$. Write $K_0=K$, $D_0=(K_0,\beta)$ and $r=dL^{K_0}$; then $L^{K_0}\le d^{-P_\sharp}$ (Proposition~\ref{C-350-16}), $K_0\ge3$, $8a_{K_0}\le d$ and $r\ge L\Lambda^2$ (Proposition~\ref{C-350-9}). As in Step~1,
\[
\Bigl|S^T_2-c_2b^{-2}\mathcal K_6-\frac1{\#\mathsf{Pk}}\sum_v\mathsf h_6(\tau_vf,\tau_vh)\Bigr|\le\Bigl[(C_{\rm A}\vartheta^{-1}bd+C_{\rm E}b^{-1}dl^{8p_0}+\bar\varepsilon_4)\,b^{-3}d+C_{\rm lat}b^{-2}\frac{a_{K_0}}d\Bigr]\eta,
\]
and, at every position, by the last link T,
\[
|\mathsf h_6-\varkappa b^{-3}\mathcal K_5|\le b^{-3}a_{K_0}\,|\mathsf H^{(K_0)}[D_0,\ell_g]-\mathsf T^{\rm lat}(\tilde w)|+b^{-3}C_{\rm T}\vartheta^{-1}a_{K_0}\lg(r)\,\eta .
\]
Let $n_0$ be the level with $D_{n_0}\in[d^{1/5},Ld^{1/5})$; then $4\le n_0\le K_0-1$ and $r/L^{n_0}\le d^{4/5}$.

\emph{Case A}, $r^{24}\le L^{K_0}$. With $e=\lceil K_0/78\rceil$ and $n_1=K_0-e$: $2\le n_1\le K_0-1$, $(r/L^{n_1})^{3/2}\le d^{6/5}$, $L^e\ge d^{-1/78}$ and $r^6\le L^{K_0/4}$. Propositions~\ref{C-350-7}(iii) from $K_0$ to $n_1$, \ref{C-350-11}(i), (ii) and~\ref{C-350-13} at $n_1$, and~\ref{C-350-15} on the lattice with $\ell_gL^e$ top sites and $n_1$ levels give $|\mathsf H^{(K_0)}-\mathsf T^{\rm lat}|\le\{C_{\rm R}l^{6p_0}d^{6/5}r+E_{\rm X}(d)+A_\ell K_0\}\eta$, since $19e-\frac12n_1<\frac{39}2-\frac{K_0}4$ gives $(\ell_gL^e)^{19}n_1L^{-n_1/2}r^6\le\ell_g^{19}L^{39/2}K_0$.

\emph{Case B}, $r^{24}>L^{K_0}$. With $K'=\lceil24\log_Lr\rceil$, $e=\lceil K'/78\rceil$ and $n_1=K'-e$: $K'>K_0$, $r^{23}>1/d$, $r^{24}\le L^{K'}<Lr^{24}\le Ld^{-24(P_\sharp-1)}$, $n_0<n_1\le K'-1$, $L^e\ge d^{-1/78}$ and $K'\le1+24(P_\sharp-1)\Lambda/\log L$. Proposition~\ref{C-350-27}, with $b_*=b$ if $b\in I(g)$ and a fixed $b_*\in I(g)$ otherwise, and with $\sigma=C_{\rm V}\mathcal y^{3p_0}b^2d^2$ from Proposition~\ref{C-350-25}(iii) at $D_0$ and $s=K'-K_0$, gives
\[
|\mathsf H^{(K_0)}-\mathsf T^{\rm lat}|\le\Bigl\{\bigl[\sigma+2C_{\rm E}b^{-1}dl^{8p_0}+2\bar\varepsilon_4+C_{\rm R}(l^{6p_0}+l'^{6p_0})d^{6/5}\bigr]r+\frac{3C_{\rm tm}b}{\vartheta d}+3E_{\rm X}(d)+A_\ell K'\Bigr\}\eta,
\]
with $l'\le\log(2b_+Lr^{23})\le l'_+$.

\emph{The places of Proposition~\ref{C-350-11}(ii).} At every one of them $\mathsf N\le\mathsf N_{\max}(d)$, and the running coupling $x_\flat=x_{n-1}$ of the place exceeds $d^{-1/78}$: $x_\flat\ge\frac34bL^{K_0-n_0+1}>d^{-1/5}$ at the level $n_0$, on either side, and $x_\flat\ge\frac34b_*L^{e+1}\ge L^e$ at the level $n_1$, as $\frac34b\ge1$. So the two conditions hold by the choice of $d_{\rm e}$, and every $\mathcal U$ is at most $E_{\rm X}(d)$, which is smaller than every power of $d$, as $2p_0\ge2$.

\emph{Collecting.} Multiply by $b^{-3}a_{K_0}$: $b^{-3}a_{K_0}r=b^{-3}d$; $b^{-3}a_{K_0}\cdot3E_{\rm X}\le\frac38E_{\rm X}b^{-3}d$, as $r\ge8$; and $b^{-3}a_{K_0}\cdot3C_{\rm tm}b/(\vartheta d)=3C_{\rm tm}\vartheta^{-1}b^{-2}a_{K_0}/d$. The two members that grow with $r$ go into the lattice term: as $\lg r\le1+(P_\sharp-1)\Lambda$ and $b\ge g^{-2}$,
\[
b^{-3}a_{K_0}\bigl[A_\ell K'+C_{\rm T}\vartheta^{-1}\lg r\bigr]=\frac db\bigl[A_\ell K'+C_{\rm T}\vartheta^{-1}\lg r\bigr]\,b^{-2}\frac{a_{K_0}}d\le\epsilon_{\rm r}(d)\,b^{-2}\frac{a_{K_0}}d\le b^{-2}\frac{a_{K_0}}d ,
\]
with $K_0$ for $K'$ in Case A. So for $s(d)<K\le K_\sharp(d)$
\[
|S^T_2-L_b^+|(f,h)\le\Bigl[\bar\varepsilon_\flat(d)\,b^{-3}d+\bigl(\tfrac32+3\bigr)C_{\rm tm}\vartheta^{-1}b^{-2}\frac{a_K}d+b^{-2}\frac{a_K}d\Bigr]\eta=\Bigl[\bar\varepsilon_\flat(d)\,b^{-3}d+C^\flat_{\rm lat}b^{-2}\frac{a_K}d\Bigr]\eta,
\]
where $C_{\rm E}b^{-1}dl^{8p_0}$ and $\bar\varepsilon_4$ enter once through the front and twice through the first link of Proposition~\ref{C-350-27}, and
\[
\begin{multlined}
\bar\varepsilon_\flat(d)=C_{\rm A}\vartheta^{-1}b_+d+3C_{\rm E}g^2d\,l_+^{8p_0}+3\bar\varepsilon_4(g,d)+C_{\rm V}[\log(10LMb_+/d)]^{3p_0}b_+^2d^2+C_{\rm R}\bigl[l_+^{6p_0} \\
+l_+'^{\,6p_0}\bigr]d^{6/5}+\tfrac38E_{\rm X}(d).
\end{multlined}
\]
This is the second assertion of (e). For $b\in I(g)$ and $K>K_\sharp(d)$, and for the limit, Proposition~\ref{C-350-16} at $K_\sharp(d)$ adds $\epsilon_\diamond(d)b^{-3}d\,\eta$, $\epsilon_\diamond(d)=C_\diamond\vartheta^{-2}b_+^3d^{177}+C^\flat_{\rm lat}b_+d^{3598}$, and the lattice term may be dropped, as in Step~5. Finally every member is at most a multiple of $d^{1/2}$: $d\,l_+^{8p_0}\le Cd^{1/2}$; $\bar\varepsilon_4(g,d)\le C_4(g)d^{1/2}$; the member of the two sides is $O(d^2\Lambda^{3p_0})$; $d^{6/5}l_+'^{\,6p_0}\le Cd$; $E_{\rm X}$ and $\epsilon_\diamond$ are smaller. So $\sup_{d''\le d}[\bar\varepsilon_\flat(d'')+\epsilon_\diamond(d'')]\le c_{\rm r}d^{1/2}$, with $c_{\rm r}$ depending on $g$, $\vartheta$ and the constants; this is (e). The slowest member is again that of the order four.

\emph{Step 9: clause (f).} Let $g$, $g_{\rm a}$ and $g_{\rm b}$ be as in Step~6, and let $d''_{\rm e}$ be the least, over the three couplings $g''$, of the number $d_{\rm e}(g'')$ of Step~1 at $P=P_\sharp$, $\frac1{180}$, $L^{-2}$, $\mathsf d_\star(g'',\vartheta)$ and $(LC_w(g''))^{-2}$. Let $(K,\beta)$ be a datum at $g$ with $K>s(d)$ and label $b$. If $K\le K_\sharp(d)$, the second assertion of (e) at $g$ applies. If $K>K_\sharp(d)$ and $b\in I(g)$, the first assertion of (e) at $g$ applies. Otherwise $b\in I(g_\bullet)$ with $g_\bullet\in\{g_{\rm a},g_{\rm b}\}$ (Step~6), $(K,\beta)$ is a datum at $g_\bullet$, and $K>s(d)$ at $g_\bullet$, since $L^K>L^{K_\sharp(d)}\ge(\sqrt3M/d)^{3600}>d^{-3}$; the first assertion of (e) at $g_\bullet$ gives the law with the rate on $T_{\ell_{g_\bullet}}$ at the pair with the same lift, and \eqref{C-350-17}, as in Step~6, bounds the difference of the two torus measures at the lift by $C^{\rm T}_2(g_\circ)\vartheta^{-1}b_+d\cdot b^{-3}d\,\eta$.

So (f) holds with $c'_{\rm r}=\max_{g''}c_{\rm r}(g'')+\max_{g''}C^{\rm T}_2(g'')\vartheta^{-1}b_+(d''_{\rm e})^{1/2}$.

\subsection{Proof of Corollary~\ref{C-350-28}}

(i) Let $\mathfrak s$ be a state at $g$ at a datum $(K,\beta)$ with $K>s(d)$; as $d\le L^{-1}$, $K\ge3$ and $a_K<d/13<9d$. By \eqref{C-350-17},
\[
|S^T_2[\mathfrak s]-S^T_2[\tilde\mu_{K,\ell_g,\beta}]|(\tilde f,\tilde h)\le C^{\rm T}_2(g)\vartheta^{-1}d^2b^{-2}\eta\le C^{\rm T}_2(g)\vartheta^{-1}b_+d\cdot b^{-3}d\,\eta,
\]
and $S^T_2[\tilde\mu_{K,\ell_g,\beta}](\tilde f,\tilde h)=S^T_{2;K,\ell_g,\beta}(f,h)$ at the periodisation $(f,h)\in\mathfrak T_{T_{\ell_g}}(d;\vartheta)$, to which clauses (c) and (f) of Theorem~\ref{C-350-3} apply; $L_b^+(f,h)=L_b^+(\tilde f,\tilde h)$. The torus measure of a side of $\mathfrak L(g)$, periodically extended, and every periodic Gibbs state at $(a_K,\beta)$ are states at $g$ at the datum. (ii) For $b\in I(g)$ every $(K,L^Kb)$ is a datum at $g$. Apply (i) to the torus measure of the side $\ell'$ at the lift of $(f,h)$, and let $K\to\infty$: the lattice term tends to $0$ and $S^T_{2;K,\ell',L^Kb}(f,h)\to S_2^{(\ell',b)}(f,h)$. (iii) The labels at $g$ lie in $[g^{-2},b_+]$, and so does their limit. The witness depth depends on no cutoff, so $K_j>s(d)$ for large $j$; apply (i) at $(K_j,\beta_j)$ and let $j\to\infty$: $a_{K_j}\to0$ and $L^+_{b_j}\to L^+_b$.

\emph{The labels.} Let the states $\mathfrak s_j$ and the set of pairs be as in the second assertion of (iii), and let $b$, $b'$ be two limit points of the labels, along two subsequences. Along each of them the first assertion applies, and at a pair $F=(\tilde f,\tilde h)$ of the set that lies in $\mathfrak T_{\mathbb R^3}(d;\vartheta)$, $d\le d_\infty$, the limit $S(F)$ of $S^T_2[\mathfrak s_j](F)$ is a limit point along both; so $|L_b^+(F)-L_{b'}^+(F)|\le\bar\varepsilon_\infty(d)(b^{-3}+b'^{-3})d\,\eta$. On the class, $\mathcal K_6(F)\ge(10d)^{-6}\int\tilde f\int\tilde h\ge10^{-6}\vartheta^2\eta$, and $\mathcal K_5(F)\le d^{-5}\int\tilde f\int\tilde h\le(\frac{4\pi}3)^2d\,\eta$, as $|x-y|\ge d$ on the supports and each support lies in a ball of radius $d$. Hence
\[
10^{-6}\vartheta^2c_2\,\bigl|b^{-2}-b'^{-2}\bigr|\le\Bigl[\bar\varepsilon_\infty(d)\bigl(b^{-3}+b'^{-3}\bigr)+\bigl(\tfrac{4\pi}3\bigr)^2c_1c_2\,\bigl|b^{-3}-b'^{-3}\bigr|\Bigr]d ,
\]
and $d\to0$ along the scales of the set gives $b=b'$. The labels lie in the compact interval $[g^{-2},b_+]$ and have one limit point, so they converge.

\emph{Limits and components.} Let $(K_j,\beta_j,\nu_j)$ be a convergent Gibbs sequence at $g$ and $\sigma$ its limit (Definition~\ref{C-365-5}(2)). Then $S^\sigma_2(X)=\lim_jS^T_2[\nu_j](X)$ at every separated tuple $X\in\mathfrak S_2(\mathbb R^3)$ (Theorem~\YMMref{AP}{780}). A pair of $\mathfrak T_{\mathbb R^3}(d;\vartheta)$ whose entries are of class $C^1$ is a separated tuple: its entries lie in $C^1_c(\mathbb R^3)$ and have disjoint supports (Section~\ref{sec:state-limit-conventions} and the test-pair convention preceding equation~\YMMref{(2343)}{849}); and for every $d>0$ and $\vartheta\in(0,1)$ the class $\mathfrak T_{\mathbb R^3}(d;\vartheta)$ holds a pair of $C^\infty$ functions (Lemma~\ref{lem:ward-test-pairs}(c)).

So the second assertion of (iii), with the states $\nu_j$ and the $C^1$ pairs in the class, shows that the labels converge, to a number $b$. For a component $\sigma_\xi$ of a canonical decomposition of $\sigma$, Theorems~\YMMref{AL}{746} and~\YMMref{AP}{780} show that $\sigma_\xi$ is a limit in the frame of the sequence along all $j$: $\sigma_\xi(X)=\lim_j\rho_j(X^{\nu_j}_{K_j})$ for every $X\in\mathfrak S(\mathbb R^3)$, $X^{\nu_j}_{K_j}$ being the product of the centred observables $\mathcal O(f_q)-\nu_j(\mathcal O(f_q))$ of the entries of $X$, with the approximants $\rho_j(F)=\nu_j(\mathbf 1_{E_j}F)/\nu_j(E_j)$, where $E_j$ is a Borel set invariant under the translations by $a_{K_j}\mathbb Z^3$ with $\nu_j(E_j)>0$; and $\rho_j$ is a periodic Gibbs state at $(a_{K_j},\beta_j)$ (Theorem~\YMMref{AP}{780}).

The convergence on every separated tuple also gives convergence of all connected functions. Indeed, put $\sigma_j(X)=\rho_j(X^{\nu_j}_{K_j})$. The inverse moment--cumulant identity in the proof of Proposition~\YMMref{92.1}{710} gives, for $n\ge2$,
\[
 S_n^T[\rho_j](X)
 =\sum_{\pi\in\operatorname{Part}(n)}
   (-1)^{|\pi|-1}(|\pi|-1)!\prod_{B\in\pi}\sigma_j(X_B).
\]
Centering changes only the first cumulant. Every subtuple $X_B$ remains separated, so each factor tends to $\sigma_\xi(X_B)$. Passing to the limit in this finite sum proves $S^{\sigma_\xi}_n(X)=\lim_jS_n^T[\rho_j](X)$ for every $n\ge2$ and $X\in\mathfrak S_n(\mathbb R^3)$. For $\sigma$ itself $\rho_j=\nu_j$. In both cases $S^{\sigma'}_2$, $\sigma'\in\{\sigma,\sigma_\xi\}$, is at a $C^1$ pair in the class the limit of $S^T_2[\rho_j]$ along states at $g$ at the data $(K_j,\beta_j)$, whose labels converge to $b$, and the first assertion of (iii) gives the law.

\emph{Volume-first limits.} A defining sequence of states of a volume-first limit $S$ at $g$ is a Gibbs sequence at $g$; it converges, and its limit is the centred functional $\sigma_S$ of $S$ (Proposition~\YMMref{92.1}{710} and its moment--cumulant identity); the ergodic components of $S$ are the components of the canonical decompositions of $\sigma_S$ (Section~\ref{sec:state-limit-conventions} and Definition~\ref{C-365-5}(2)). So the labels along every defining sequence converge, and the law holds for $S$ and its ergodic components at the pairs of class $C^1$. Two defining sequences give one label: the one function $S_2=S^{\sigma_S}_2$ obeys the law with both limits at the $C^1$ pairs in the class at every small $d$, and the inequality displayed above gives their equality.

This is the unique label of Theorem~\YMMref{AW}{852}. For $b\in I(g)$ the pairs $(K,L^Kb)$, $K\in\mathbb N$, are data at $g$ (Corollary~\YMMref{74.2}{616}), and Proposition~\YMMref{92.1}{710}, applied to torus-limit states at these data, gives a volume-first limit at $g$ with a defining sequence of states $(K_j,\beta_j,\nu_j)$ with $a_{K_j}\beta_j=b$ for every $j$; its label is $b$.

(iv) $\mathcal K_5\ge10^{-5}\vartheta^2d\,\eta$ on the class; divide by $c_2b^{-3}\mathcal K_5$. In the case (i) the lattice term $C_{\rm lat}b^{-2}(a_K/d)\,\eta$, divided by $c_2b^{-3}\cdot10^{-5}\vartheta^2d\,\eta$, gives the member $10^5c_2^{-1}\vartheta^{-2}C_{\rm lat}\,b\,a_K/d^2$, which is small only if $a_K$ is small against $d^2/b$. In the cases (ii) and (iii), $|S-c_2b^{-2}\mathcal K_6|\le[(\frac{4\pi}3)^2c_1c_2+\bar\varepsilon_\infty(d)]\,b^{-3}d\,\eta$ and $\mathcal K_6\ge10^{-6}\vartheta^2\eta$, so that $S/(c_2\mathcal K_6)\to b^{-2}$ as $d\to0$ along the pairs at which the law holds: each of these limits determines its label, and then the quotient determines the number $\frac8{3\pi}$. A law of this form with a tolerance that tends to $0$ with $d$ thus fixes both its label and its coefficient.

\begin{remark}\label{C-350-33}
At a finite cutoff the tolerance of Theorem~\ref{C-350-3} is $\bar\varepsilon(d)$, or $c_{\rm r}d^{1/2}$, together with the lattice term $C_{\rm lat}b^{-2}a_K/d$, which is larger than the correction $c_1c_2b^{-3}\mathcal K_5$, of the order $b^{-3}d$, unless $a_K$ is small against $d^2/b$; the number $c_1$ is read off at such cutoffs or in the continuum limit. Clause (d) is weaker than clause (b): a mean over scales does not see an oscillation in the scale. The arguments use one finite lattice or one torus at a time, at scales below the least side; the agreement of states uses besides one uniformity in the volume, which comes from reflection positivity and from the bounds of Theorems~\YMMref{AJ}{673} and~\YMMref{AD}{607} on the partition functions of the tori at the source $0$, and no property of the correlations at large distances.
\end{remark}


\section{Natural scaling and failure of dilation covariance}\label{C-351-1}

We now express the two-term law in the natural parameter $d/b$.
The resulting tolerance is uniform over the admitted reference couplings.
The law holds at every positive label on a fixed torus and, on
$\mathbb R^3$, in every volume-first limit at the stated coupling ceiling
and in each of its canonical components. The nonzero second term excludes
dilation covariance of the two-point function.

We give two passages to this conclusion: a direct transfer of the law
between labels and an alternative using blocking factors with
incommensurable logarithms. We also retain the leading law for signed
test functions, with its sharp linear rate. These conclusions do not by
themselves establish interaction. Proposition~\ref{C-351-24} constructs a
quadratic Gaussian model with the same two-term law; its three-point
behavior distinguishes it from the torus limits considered here.

\subsection{Parameters and test pairs}

We retain the construction setting of Section~\ref{C-350-1}, and
write $b$ and $d$ for the label $\mathsf b$ and scale $\mathsf d$.
A \emph{choice} $\mathfrak c$ consists of a mollifier, a radius
$\varrho>0$, the common real and sourced numerical assignment of
Proposition~\YMMref{69.1}{595} and Section~\YMMref{82}{665} at source
radius $2\varrho$, and the canonical auxiliary data for these
constructions. Every constant in the real restriction keeps its value.
The choice fixes an odd $L\ge13$, $B_\sharp\ge0$,
$\bar B=B_\sharp L/(L-1)$, $M$, and $0<\gamma\le1/2$.
The short-distance ceilings $\gamma_W$ of
Theorem~\YMMref{AV}{851} and $\gamma_W^{\rm sd}$ of
the choices in equations~\YMMref{(2326)--(2327)}{845} are both at most $\gamma$.
The ceiling $\gamma_{\rm e}\in(0,\gamma]$, depending only on the
constants, is fixed at the beginning of the subsection on the least
torus, and $\Gamma:=\min\{\gamma_{\rm e},\gamma_W^{\rm sd}\}$.

We write $a_K=L^{-K}$, $X_g=L(g^{-2}+B_\sharp)$, $I(g)=(g^{-2}+\bar B,X_g-\bar B]$ for the interval of labels at $g$ (Section~\YMMref{74}{614}) and $I^\circ(g)$ for its interior, $\mathfrak L(g)=\{2L^m:m\ge m_{\mathfrak L}(g)\}$ for the set of the sides at $g$ (equation~\YMMref{(1528)}{615}) and $\ell_g$ for its least member, and $T_\ell=\mathbb R^3/\ell\mathbb Z^3$. A triple $(g,b,\ell)$ is \emph{admissible} if $0<g\le\gamma_W^{\rm sd}$, $b\in I^\circ(g)$ and $\ell\in\mathfrak L(g)$. Since $LB_\sharp=(L-1)\bar B$, we have $I^\circ(g)=(g^{-2}+\bar B,\,Lg^{-2}+(L-2)\bar B)$, an interval of length $(L-1)g^{-2}+(L-3)\bar B>0$.

We use the lattices, first-passage data, labels and torus cumulants of
Section~\ref{C-350-1}, with the same fixed choice of construction data.
The continuum torus cumulants $S_n^{(\ell,b)}$ exist for every $b>0$ on
Lipschitz tuples with no common point, by Theorem~\YMMref{A}{5}.
A volume-first limit $S$ at $g$ is defined by a sequence $(K_j,\beta_j,\nu_j)$ as in Section~\ref{sec:state-limit-conventions}; for $g\le\gamma_W$ it has one label $\mathsf b(S)=\lim_ja_{K_j}\beta_j$. $(\mathsf X,\lambda,(\sigma_\xi))$ denotes a canonical decomposition of its functional $\sigma_S$, with the probability measure $\lambda$ on $\mathsf X$; $S_2$ and $S_2^\xi$ are the two-point functions of $\sigma_S$ and $\sigma_\xi$, and $\mathsf N\subset\mathsf X$ is the set of the $\xi$ for which $S_2^\xi$ is not dilation covariant.

The class $\mathfrak T_{\mathbb E}(d;\vartheta)$ is the test-pair class
of Section~\ref{C-350-1}. $\mathfrak T^1(d;\vartheta)$ is the set of its pairs on $\mathbb R^3$ of class $C^1$. Over $\mathbb R^3$ the pairs of real $C^1$ functions with compact disjoint supports are called \emph{separated}; the pairs of $\mathfrak T^1(d;\vartheta)$ are separated. On a torus we admit real Lipschitz entries with disjoint supports. Every $\mathfrak T^1(d;\vartheta)$, $d>0$, $\vartheta\in(0,1)$, contains a pair (Lemma~\ref{lem:ward-test-pairs}(c)); for $\vartheta\le3\sqrt3/8$ one may take $f(u)=(1-|u-x|^2/d^2)_+^2$, $h(u)=(1-|u-y|^2/d^2)_+^2$ with $3d\le|x-y|\le8d$. For a pair we write $|\Psi|=\|f\|_\infty\|h\|_\infty$, $\Psi_r=(f(\cdot/r),h(\cdot/r))$ for $r>0$, $\Psi^{\langle\ell\rangle}$ for the periodisation, and
\[
\mathcal K_q(\Psi)=\iint f(x)h(y)\,|x-y|^{-q}\,dx\,dy ,
\]
on $T_\ell$ with the Euclidean distance of $T_\ell$; further $T^{(\ell,b)}(\Psi)=S_2^{(\ell,b)}(\Psi^{\langle\ell\rangle})$, $c_2=\frac9{16\pi^2}$, $c_1=\frac8{3\pi}$, $\varkappa=c_1c_2=\frac3{2\pi^3}$ and $L_b^+=c_2b^{-2}[\mathcal K_6+c_1b^{-1}\mathcal K_5]$. By Lemma~\ref{lem:ward-test-pairs} and equation~\YMMref{(2247)}{824} and~\ref{lem:ward-test-pairs}(b): $\Psi_r\in\mathfrak T^1(rd;\vartheta)$ for $\Psi\in\mathfrak T^1(d;\vartheta)$, with the same $\vartheta$, $|\Psi_r|=|\Psi|$ and $\mathcal K_q(\Psi_r)=r^{6-q}\mathcal K_q(\Psi)$; translations keep the classes, the suprema and the $\mathcal K_q$; for $d<\ell/20$ the periodisation maps $\mathfrak T_{\mathbb R^3}(d;\vartheta)$ into $\mathfrak T_{T_\ell}(d;\vartheta)$, keeps the suprema, and $\mathcal K_q(\Psi^{\langle\ell\rangle})=\mathcal K_q(\Psi)$ for every real $q>0$. As $d\le|x-y|\le10d$ on the supports of a pair of the class, $(10d)^{-q}\int f\int h\le\mathcal K_q(\Psi)\le d^{-q}|\Psi|(\frac{4\pi}3d^3)^2$, that is,
\begin{equation}\label{C-351-2}
10^{-5}\vartheta^2d\,|\Psi|\le\mathcal K_5(\Psi)\le\bigl(\tfrac{4\pi}3\bigr)^2d\,|\Psi|,\qquad10^{-6}\vartheta^2|\Psi|\le\mathcal K_6(\Psi)\le\bigl(\tfrac{4\pi}3\bigr)^2|\Psi|,\qquad\varkappa\bigl(\tfrac{4\pi}3\bigr)^2=\tfrac8{3\pi};
\end{equation}
note $10^{-5}\varkappa=4.837\ldots\cdot10^{-7}$.

A functional $\sigma$ on separated tuples over $\mathbb R^3$ \emph{obeys the law} (L) \emph{with} $(b,c,\varepsilon,d_7,\vartheta)$, where $\varepsilon$ is non-decreasing on $(0,d_7]$ and $\varepsilon(d)\to0$ as $d\to0$, if
\[
|S_2^\sigma(\Psi)-c_2b^{-2}[\mathcal K_6+cb^{-1}\mathcal K_5](\Psi)|\le\varepsilon(d)\,b^{-3}d\,|\Psi|\qquad(0<d\le d_7,\ \Psi\in\mathfrak T^1(d;\vartheta)).
\]
$S_2^\sigma$ is \emph{dilation covariant} if there is a real $\varpi$ with $S_2^\sigma(\Phi_r)=r^\varpi S_2^\sigma(\Phi)$ for all $r>0$ and all separated pairs $\Phi$. For a function $\mathsf S$ on pairs, $t>1$ and $r>0$ we call
\[
\mathcal N^{\mathsf S}_t(\Psi;r)=\mathsf S(\Psi_{t^2r})-(1+t^6)\,\mathsf S(\Psi_{tr})+t^6\,\mathsf S(\Psi_r)
\]
the $t$-scale number; at $t=2$ it is the \emph{three-scale number} $\mathsf S(\Psi_{4r})-65\,\mathsf S(\Psi_{2r})+64\,\mathsf S(\Psi_r)$. The weights annul constants and multiples of $r^6$, and give the function $r\mapsto r$ the factor $t^2-(1+t^6)t+t^6=-t(t-1)(t^5-1)$, which is $-62$ at $t=2$. For a triple $\tau=(g,b,\ell)$ and a pair $\Psi$ we put $E_\tau=2C_2^{\rm T}(g)\vartheta^{-1}b^2$, with the number $C_2^{\rm T}$ of Proposition~\ref{C-351-7}, and $Y_\tau(\delta)=b^2\,T^{(\ell,b)}(\Psi_{\delta b})$.

\subsection{The law at natural scales}

\begin{theorem}[The components]\label{C-351-3}
Let $\mathfrak c$ be a choice and $\vartheta\in(0,1)$. There are $\delta_{\rm c}>0$ and a non-decreasing function $\varepsilon^{\rm nat}$ on $(0,\delta_{\rm c}]$ with $\varepsilon^{\rm nat}(\delta)\le C\delta^{1/2}$, where $\delta_{\rm c}$, $\varepsilon^{\rm nat}$ and $C$ depend only on $\mathfrak c$ and $\vartheta$, such that for every $0<g\le\gamma_W$, every volume-first limit $S$ at $g$ with $b=\mathsf b(S)$, every canonical decomposition $(\mathsf X,\lambda,(\sigma_\xi))$ of $\sigma_S$, and for $\mathsf S=S_2$ and $\mathsf S=S_2^\xi$, $\xi\in\mathsf X$, the following holds.
\begin{enumerate}
\item[(1)] $|\mathsf S(\Psi)-L_b^+(\Psi)|\le\varepsilon^{\rm nat}(d/b)\,b^{-3}d\,|\Psi|$ for $0<d\le\delta_{\rm c}b$ and $\Psi\in\mathfrak T^1(d;\vartheta)$. So $\sigma_S$ and every $\sigma_\xi$ obey (L) with $(b,\frac8{3\pi},\varepsilon^S,d^S,\vartheta)$, where $\varepsilon^S(d)=\varepsilon^{\rm nat}(g^2d)$ and $d^S=\delta_{\rm c}g^{-2}$ depend only on $\mathfrak c$, $g$ and $\vartheta$. They obey (L) also with a tolerance formed at any dilated label: if $z$ is an integer with $L^zb\ge2\Gamma^{-2}$ and $g'=(2/(L^zb))^{1/2}$, then (L) holds with $(b,\frac8{3\pi},\varepsilon',d',\vartheta)$, where $\varepsilon'(d)=\varepsilon^L_{g',\vartheta}(L^zd)+2C_2^{\rm T}(g')\vartheta^{-1}L^{2z}b\,d$ and $d'=L^{-z}\min\{d^L,d_\flat\}(g',\vartheta)$, with the objects of Theorem~\ref{C-351-4}~(a) and Proposition~\ref{C-351-7}.
\item[(2)] Let $d_1>0$, $\Psi\in\mathfrak T^1(d_1;\vartheta)$ and $B_S=\varkappa b^{-3}\mathcal K_5(\Psi)$, so that $B_S\ge10^{-5}\varkappa\vartheta^2b^{-3}d_1|\Psi|>0$; let $t>1$. Then
\begin{align}
&|\mathsf S(\Psi_r)-c_2b^{-2}\mathcal K_6(\Psi)-B_Sr|\le\varepsilon^{\rm nat}(rd_1/b)\,b^{-3}d_1|\Psi|\,r&&(rd_1\le\delta_{\rm c}b),\notag\\
&|\mathcal N^{\mathsf S}_t(\Psi;r)+t(t-1)(t^5-1)B_Sr|\le(t^2+t+t^6+t^7)\,\varepsilon^{\rm nat}(t^2rd_1/b)\,b^{-3}d_1|\Psi|\,r&&(t^2rd_1\le\delta_{\rm c}b).\notag
\end{align}
So $\mathsf S(\Psi_r)=c_2b^{-2}\mathcal K_6(\Psi)+B_Sr+O(r^{3/2})$; $|r^{-1}[\mathsf S(\Psi_{tr})-\mathsf S(\Psi_r)]-(t-1)B_S|\le(t+1)\,\varepsilon^{\rm nat}(trd_1/b)\,b^{-3}d_1|\Psi|$, so that $r^{-1}[\mathsf S(\Psi_{tr})-\mathsf S(\Psi_r)]\to(t-1)B_S$; $r^{-1}\mathcal N^{\mathsf S}_t(\Psi;r)\to-t(t-1)(t^5-1)B_S$, and so do the means $n^{-1}\sum_{j<n}(t^{-j}r_0)^{-1}\mathcal N^{\mathsf S}_t(\Psi;t^{-j}r_0)$ as $n\to\infty$, for every $r_0$ with $t^2r_0d_1\le\delta_{\rm c}b$; and $\mathcal N^{\mathsf S}_t(\Psi;r)<0$ as soon as $(t^2+t+t^6+t^7)\,\varepsilon^{\rm nat}(t^2rd_1/b)<t(t-1)(t^5-1)\cdot10^{-5}\varkappa\vartheta^2$. At $t=2$: the three-scale number is $-62B_Sr+O(r^{3/2})$, within $198\,\varepsilon^{\rm nat}(4rd_1/b)\,b^{-3}d_1|\Psi|\,r$ of $-62B_Sr$, and negative as soon as $198\,\varepsilon^{\rm nat}(4rd_1/b)<62\cdot10^{-5}\varkappa\vartheta^2$, that is, for all small $r$.
\item[(3)] Neither $S_2$ nor any $S_2^\xi$ is dilation covariant: $\mathsf N=\mathsf X$ and $\lambda(\mathsf N)=1$.
\end{enumerate}
\end{theorem}

\begin{theorem}[The tori]\label{C-351-4}
Let $\mathfrak c$ be a choice and $\vartheta\in(0,1)$; let $d_{\rm e}$, $\bar\varepsilon=\bar\varepsilon_{g,\vartheta}$, $\bar\varepsilon_\infty=\bar\varepsilon_{\infty;g,\vartheta}$ and $C_{\rm cut}$ be the objects of Proposition~\ref{C-351-14} and the depth $s(d)$ that of the subsection on the least torus, and $d_\flat$, $C_2^{\rm T}$ those of Proposition~\ref{C-351-7}.
\begin{enumerate}
\item[(a)]

\emph{The least torus.} For $0<g\le\gamma_{\rm e}$ there are $d^L(g,\vartheta)>0$ and a non-decreasing function $\varepsilon^L_{g,\vartheta}$ on $(0,d^L(g,\vartheta)]$ with $\varepsilon^L_{g,\vartheta}(d)\le C(g,\vartheta)d^{1/2}$, all depending only on $\mathfrak c$, $g$ and $\vartheta$, such that
\[
|T^{(\ell_g,b)}(\Psi)-L_b^+(\Psi)|\le\varepsilon^L_{g,\vartheta}(d)\,b^{-3}d\,|\Psi|\qquad\bigl(b\in I(g),\ 0<d\le d^L(g,\vartheta),\ \Psi\in\mathfrak T_{\mathbb R^3}(d;\vartheta)\bigr).
\]
\item[(b)]

\emph{Every torus of the family; every state at a finite cutoff.} Let $0<g\le\Gamma$, $b\in I(g)$, $\ell\in\mathfrak L(g)$, $0<d\le\min\{d_{\rm e},\ell_g/20,d_\flat\}(g,\vartheta)$ and $\Psi\in\mathfrak T_{\mathbb R^3}(d;\vartheta)$. Then
\[
|T^{(\ell,b)}(\Psi)-L_b^+(\Psi)|\le\bigl[\bar\varepsilon_\infty(d)+2C_2^{\rm T}(g)\vartheta^{-1}bd\bigr]\,b^{-3}d\,|\Psi|,
\]
and the bracket is at most $C(g,\vartheta)(1+b)d^{1/2}$. Moreover, for every cutoff $K>s(d)$ and every state $\mathfrak s$ at the datum $(K,L^Kb)$ in the sense of Proposition~\ref{C-351-7},
\[
|S^T_2[\mathfrak s](\Psi)-L_b^+(\Psi)|\le\Bigl[\bigl(\bar\varepsilon(d)+C_2^{\rm T}(g)\vartheta^{-1}bd\bigr)\,b^{-3}d+C_{\rm cut}\,b^{-2}\frac{a_K}d\Bigr]|\Psi| ,
\]
and the same holds with $C_{\rm r}d^{1/2}$ in the place of $\bar\varepsilon(d)$ and $C'_{\rm cut}$ in the place of $C_{\rm cut}$, with the numbers of Proposition~\ref{C-351-14}~(c).
\item[(c)]

\emph{Every side, every label.} Let $\ell=2L^m$, $m\ge1$, $b>0$, $s\ge0$ an integer and $0<g''\le\Gamma$ a coupling for which $(g'',L^sb,L^s\ell)$ is admissible; this holds for all large $s$ with $g''=(2/(L^sb))^{1/2}$. Then for $0<d\le L^{-s}\min\{d_{\rm e},\ell_{g''}/20,d_\flat\}(g'',\vartheta)$ and $\Psi\in\mathfrak T^1(d;\vartheta)$
\[
|T^{(\ell,b)}(\Psi)-L_b^+(\Psi)|\le\bigl[\bar\varepsilon_\infty(L^sd)+2C_2^{\rm T}(g'')\vartheta^{-1}L^{2s}bd\bigr]\,b^{-3}d\,|\Psi|,
\]
with $\bar\varepsilon_\infty$ formed at $(g'',\vartheta)$; the bracket is at most $Cd^{1/2}$, with $C$ depending on $\mathfrak c$, $g''$, $s$, $b$ and $\vartheta$.
\item[(d)]

\emph{Exact dilation by powers of $L$.} $T^{(\ell,b)}(\Psi)=L^{2s}\,T^{(L^s\ell,L^sb)}(\Psi_{L^s})$ for $\ell=2L^m$, $b>0$, integers $s\ge0$ and pairs $\Psi$ of real Lipschitz functions with compact supports whose periodisations to $T_\ell$ have disjoint supports. In particular, for $0<g\le\gamma_{\rm e}$, $b\in I(g)$ and every integer $z\ge0$, the law of~(a) holds on the side $L^z\ell_g$ at the label $L^zb$ with no comparison of states:
\[
|T^{(L^z\ell_g,L^zb)}(\Psi)-L^+_{L^zb}(\Psi)|\le\varepsilon^L_{g,\vartheta}(L^{-z}d)\,(L^zb)^{-3}d\,|\Psi|\qquad\bigl(0<d\le L^zd^L(g,\vartheta),\ \Psi\in\mathfrak T_{\mathbb R^3}(d;\vartheta)\bigr).
\]
\item[(e)]

\emph{One coefficient.} If at one side $2L^m$ and one label the inequality of~(c) holds with a number $c$ and with a number $c'$ in the place of $\frac8{3\pi}$, each with a tolerance that tends to $0$, then $c=c'$.
\item[(f)]

\emph{Compact sets of labels.} At a fixed side, one tolerance, one range of $d$ and one constant $C$ serve in~(c) for all the labels of a compact subset of $(0,\infty)$.
\end{enumerate}
Let now $0<g\le\gamma_{\rm e}$, $b\in I(g)$, $T=T^{(\ell_g,b)}$, $\varepsilon^L=\varepsilon^L_{g,\vartheta}$, $d^L=d^L(g,\vartheta)$, $d_1>0$, $\Psi\in\mathfrak T^1(d_1;\vartheta)$, $\mathsf A=c_2b^{-2}\mathcal K_6(\Psi)$, $\mathsf B=\varkappa b^{-3}\mathcal K_5(\Psi)$ and $\mathsf u=b^{-3}d_1|\Psi|$. Then $10^{-5}\varkappa\vartheta^2\mathsf u\le\mathsf B\le\frac8{3\pi}\mathsf u$, and for $0<rd_1\le d^L$
\begin{equation}\label{C-351-5}
|T(\Psi_r)-\mathsf A-\mathsf Br|\le\varepsilon^L(rd_1)\,\mathsf u\,r\le C(g,\vartheta)\,d_1^{1/2}\,\mathsf u\,r^{3/2}.
\end{equation}
\begin{enumerate}
\item[(i)] For $t>1$ and $0<trd_1\le d^L$: $|r^{-1}[T(\Psi_{tr})-T(\Psi_r)]-(t-1)\mathsf B|\le(t+1)\,\varepsilon^L(trd_1)\,\mathsf u\to0$; for $0<t^2rd_1\le d^L$: $|r^{-1}\mathcal N^T_t(\Psi;r)+t(t-1)(t^5-1)\mathsf B|\le(t^2+t+t^6+t^7)\,\varepsilon^L(t^2rd_1)\,\mathsf u\to0$. At $t=2$: $|r^{-1}[T(\Psi_{2r})-T(\Psi_r)]-\mathsf B|\le3\varepsilon^L(2rd_1)\mathsf u$.
\item[(ii)] $r^{-2}|T(\Psi_r)-\mathsf A|\ge r^{-1}\mathsf u\,[10^{-5}\varkappa\vartheta^2-\varepsilon^L(rd_1)]\to\infty$ as $r\to0$.
\item[(iii)] For all $d_n\to0$ and all $\Psi_n\in\mathfrak T^1(d_n;\vartheta)$: $|T(\Psi_n)-c_2b^{-2}\mathcal K_6(\Psi_n)|/(d_n^2|\Psi_n|)\ge b^{-3}[10^{-5}\varkappa\vartheta^2-\varepsilon^L(d_n)]/d_n\to\infty$.
\item[(iv)] Let $g\le\Gamma$ and $b\in I^\circ(g)$. For every $0<C<\infty$, every $d\le d^L$ with $\varepsilon^L(d)\le\frac12\cdot10^{-5}\varkappa\vartheta^2$ and $d/b<\frac12\cdot10^{-5}\varkappa\vartheta^3/C$, and every $\Psi\in\mathfrak T^1(d;\vartheta)$: $b^2|T(\Psi)-c_2b^{-2}\mathcal K_6(\Psi)|>C\vartheta^{-1}(d/b)^2|\Psi|$.
\end{enumerate}
\end{theorem}

\begin{corollary}[The leading term with the rate $r$, for signed functions]\label{C-351-22}
Let $\mathfrak c$ be a choice, $\ell=2L^m$ with $m\ge1$, $b>0$, $p\in T_\ell$ with a lift $\tilde p$, let $f_0,h_0$ be real Lipschitz functions on $\mathbb R^3$ with compact disjoint supports, and put $f_r=(f_0((\cdot-\tilde p)/r))^{\langle\ell\rangle}$ and $h_r=(h_0((\cdot-\tilde p)/r))^{\langle\ell\rangle}$ for $r>0$.
\begin{enumerate}
\item[(a)] There are $C,r_*>0$, independent of $p$ and the same for all the labels of a compact subset of $(0,\infty)$, with
\[
|S_2^{(\ell,b)}(f_r,h_r)-c_2b^{-2}\mathcal K_6(f_0,h_0)|\le Cr\qquad(0<r\le r_*).
\]
In particular, for $\theta_\mu x=x-2x_\mu e_\mu$ and every real Lipschitz function $\psi_0$ with compact support in $\{x_\mu>0\}$, the pair $(f_0,h_0)=(\psi_0\circ\theta_\mu,\psi_0)$ gives $S_2^{(\ell,b)}(f_r,h_r)=c_2b^{-2}\mathcal K_6(\psi_0\circ\theta_\mu,\psi_0)+O(r)$.
\item[(b)] The rate $r$ is attained: if $\psi_0\ge0$ is of class $C^1$, not $0$, with support in the ball $\mathbb B(e_\mu,\frac12)$, then
\[
S_2^{(\ell,b)}(f_r,h_r)-c_2b^{-2}\mathcal K_6(\psi_0\circ\theta_\mu,\psi_0)=\varkappa\,b^{-3}\mathcal K_5(\psi_0\circ\theta_\mu,\psi_0)\,r+O(r^{3/2}),\qquad\mathcal K_5(\psi_0\circ\theta_\mu,\psi_0)>0 ;
\]
so no bound $o(r)$ holds in~(a).
\item[(c)] For $0<b_1<b_2$, with the constant $C$ of~(a) for the labels of $[b_1,b_2]$ and $0<r\le r_*$,
\[
\Bigl|\int_{b_1}^{b_2}\bigl[b^2S_3^{(\ell,b)}(1,f_r,h_r)-2b\,S_2^{(\ell,b)}(f_r,h_r)\bigr]\,db\Bigr|\le(b_1^2+b_2^2)\,C\,r ,
\]
where $1$ is the constant function on $T_\ell$.
\end{enumerate}
\end{corollary}

\subsection{Construction inputs}

The proof rests on Theorems~\YMMref{D}{8}, \YMMref{AJ}{673} and~\YMMref{AD}{607} among the construction results, with its conditions on the constants, through the statements cited below. Among them are Section~\YMMref{74}{614}, \ref{sec:setup-first-passage-arithmetic}, \YMMref{(1528)}{615} and~\YMMref{82}{665}, the real-sequence conditions and smallness floor~\YMMref{(972)}{438} preceding Theorem~\YMMref{Y}{438}, the bounds~\YMMref{(1531)}{615} on labels, Corollary~\YMMref{74.2}{616} and Lemma~\ref{lem:setup-label-coverage} on couplings and labels, which we use in the form of Lemma~\ref{C-351-25} below, the exact dilation by powers of $L$~\YMMref{(2215)}{814}, Theorem~\textup{\YMMref{AS}{813}} on the limits at a fixed label (their existence, the rate of removal of the cutoff, the bounds~\YMMref{(2214)}{814} uniform in the cutoff and in the labels of a compact set, the continuity in the label and the convergence along converging labels), Proposition~\YMMref{118.3}{814} on the derivatives in the label, the uniform cutoff estimate of Theorem~\YMMref{AR}{810} on the removal of the cutoff, the weighted progression and its source energies (Definition~\ref{def:weighted-progression-data} and Corollary~\ref{cor:weighted-canonical-data}), its setting (J) for the canonical objects (Definition~\ref{def:weighted-history-setting} and equation~\eqref{eq:weighted-history-threshold} and Proposition~\ref{prop:weighted-canonical-history}), the places with a lower place of a step at a complex field scale (Definition~\ref{IV-71-12} and Proposition~\ref{IV-71-50}), the sourced renormalisation-group chain and its source energies (Proposition~\ref{prop:running-source-chain} and equations~\eqref{eq:running-source-energies}--\eqref{eq:running-source-energy-linear}), the histories, the sum over the level of the first creation of a large-field region and the rarity of large fields (Section~\YMMref{123}{828}, equation~\eqref{eq:ward-first-creation} and Proposition~\ref{prop:ward-fixed-volume-rarity}), the reflection positivity of the limits on the tori (proved in Proposition~\ref{C-351-24}(c) from Theorem~\YMMref{A}{5} and finite-cutoff positivity, for every side $2L^m$, every $b>0$ and $C^1$ entries at every coordinate mirror) and that of the Wilson measures on the tori at the site and bond mirrors (Proposition~\YMMref{91.4}{700}), which, with the existence of the limits along the full sequence of the cutoffs, gives it for real Lipschitz entries at the mirrors $\{x_\mu=c\}$, $c\in\bigcup_{K\in\mathbb N}\frac12L^{-K}\mathbb Z$, of the lattices of the cutoffs, the block averages (Lemma~\ref{lem:sf-chart-integral}), the template sum of the free lattice curvature (Proposition~\YMMref{121.2}{824} and Lemma~\ref{lem:ward-fluctuation-template}), Laplace's method (Lemma~\textup{\ref{IV-73-26}}) and the coefficients of lattice perturbation theory (Definition~\textup{\ref{IV-73-22}}, Corollary~\textup{\ref{IV-73-73}}, Theorems~\textup{\ref{IV-73-98}} and~\textup{\ref{IV-73-100}}), the chessboard estimate (Corollary~\ref{IV-18-18}), the representation of the components by conditioned periodic Gibbs states (Theorem~\textup{\YMMref{AP}{780}}), and, for $g\le\gamma_W$, the leading law of Theorem~\YMMref{AW}{852} for $\sigma_S$ and, using the conditioned approximants in Theorem~\YMMref{AP}{780}, for its components: for $0<d\le\mathsf d_0(g,\vartheta)$ and the $C^1$ pairs of $\mathfrak T_{\mathbb R^3}(d;\vartheta)$, $\mathsf S(\Psi)=c_2b^{-2}\mathcal K_6(\Psi)(1+\eta)$ with $|\eta|\le\eta_{g,\vartheta}(d)\le\frac13$ and $\eta_{g,\vartheta}(d)\to0$ as $d\to0$.

Two further statements are proved below from these: the agreement of all states at one datum at two points, which uses no decay of correlations (Proposition~\ref{C-351-7}), and the next-order law on the least torus (Proposition~\ref{C-351-14}).

\begin{lemma}[Couplings and labels]\label{C-351-25}
\begin{enumerate}
\item[(a)] $\gamma^{-2}>\bar B$. For every real $b'>2\bar B$, in particular for every $b'\ge2\gamma^{-2}$, the coupling $g'=(2/b')^{1/2}$ satisfies $b'\in I^\circ(g')$; if $b'\ge2\gamma^{-2}$, then $g'\le\gamma$.
\item[(b)] Let $0<g'\le\gamma$, and let $(K_j,\beta_j)$ be pairs of integers $K_j\ge1$ and reals $\beta_j>0$ with $a_{K_j}\beta_j\to b'\in I^\circ(g')$. Then $(K_j,\beta_j)$ is a first-passage datum at $g'$ for all large $j$.
\end{enumerate}
\end{lemma}

\begin{proof}[Proof sketch]
(a) Corollary~\YMMref{74.2}{616} and Lemma~\ref{lem:setup-label-coverage} at the label $2\gamma^{-2}$, where the coupling is $\gamma$, puts $2\gamma^{-2}$ in $I(\gamma)=(\gamma^{-2}+\bar B,X_\gamma-\bar B]$; so $\gamma^{-2}>\bar B$. With $g'^{-2}=b'/2$ the formula for $I^\circ$ in the setting gives $I^\circ(g')=(\frac12b'+\bar B,\,\frac12Lb'+(L-2)\bar B)$. Its upper end exceeds $b'$, as $L\ge3$, and its lower end is less than $b'$ if and only if $b'>2\bar B$, which holds for $b'\ge2\gamma^{-2}$; and then $g'\le(2/(2\gamma^{-2}))^{1/2}=\gamma$. (b) As $I^\circ(g')$ is open, $a_{K_j}\beta_j\in I(g')$ for all large $j$, and Corollary~\YMMref{74.2}{616} at the label $a_{K_j}\beta_j$ says that $(K_j,L^{K_j}a_{K_j}\beta_j)=(K_j,\beta_j)$ is a first-passage datum at $g'$.
\end{proof}

\subsection{Agreement of states at two points}

\begin{proposition}[Agreement of two states at two points]\label{C-351-7}
Let $\mathfrak c$ be a choice, $0<g'\le\gamma_W^{\rm sd}$, $(K,\beta)$ a first-passage datum at $g'$ with $K\ge3$ and label $b=a_K\beta$, $\vartheta\in(0,1)$, $0<d\le d_\flat(g',\vartheta)$ and $a_K\le9d$. A \emph{periodic Gibbs state} of the Wilson action at $\beta$ on $a_K\mathbb Z^3$ is a Gibbs state at $(a_K,\beta)$ that is invariant under the translations of a subgroup of finite index of $a_K\mathbb Z^3$ (Section~\YMMref{91}{697}); these states form the set $\mathcal G^{\rm per}_{K,\beta}$. For every $K$ and every $\beta>0$ this set is the set $\mathcal G^{\rm inv}_{K,\beta}$ of the periodic Gibbs states that are invariant under all the translations of $a_K\mathbb Z^3$ (Proposition~\YMMref{103.4}{758}; Section~\ref{sec:state-limit-conventions}), and it contains the states of the tori, the weak limit points of the periodic extensions of the Wilson measures on the tori $T_{2L^m}$, $m\to\infty$ (Proposition~\YMMref{91.2}{699}; Section~\ref{sec:state-limit-conventions}).

A \emph{state at the datum} is the periodic extension to $a_K\mathbb Z^3$ of the Wilson measure on $T_\ell$ at $(K,\beta)$, for some $\ell\in\mathfrak L(g')$, or a periodic Gibbs state of the Wilson action at $\beta$ on $a_K\mathbb Z^3$: a state at $g'$ at $(K,\beta)$ in the sense of Section~\ref{sec:state-limit-conventions}. Every state at the datum is invariant under the translations of $a_K\mathbb Z^3$, and so are its cumulants: $S^T_n[\mathfrak s](\tau_vf_1,\dots,\tau_vf_n)=S^T_n[\mathfrak s](f_1,\dots,f_n)$ for all $v\in a_K\mathbb Z^3$ (Section~\ref{sec:state-limit-conventions}). Let $\mu'$ be the periodic extension of the Wilson measure on the least torus $T_{\ell_{g'}}$ at $(K,\beta)$.
\begin{enumerate}
\item[(a)] Every state $\mathfrak s$ at the datum satisfies $|S^T_2[\mathfrak s](\Psi)-S^T_2[\mu'](\Psi)|\le C_2^{\rm T}(g')\vartheta^{-1}d^2b^{-2}|\Psi|$ for all $\Psi\in\mathfrak T_{\mathbb R^3}(d;\vartheta)$; hence any two states $\mathfrak s,\mathfrak s'$ at the datum have
\[
|S^T_2[\mathfrak s](\Psi)-S^T_2[\mathfrak s'](\Psi)|\le2C_2^{\rm T}(g')\,\vartheta^{-1}d^2(a_K\beta)^{-2}|\Psi| .
\]
\item[(b)] Here $d_\flat(g',\vartheta)=\min\{\frac1{180},d_\star(g',\vartheta)\}$ with $d_\star(g',\vartheta)=c_{\rm a}\vartheta^{-1}g'^{-2}$, and
\[
C_2^{\rm T}(g')=C_{\rm a}\bigl[Q_0(g')\,Q_1(g')\bigr]^2\,\frac{(X_{g'}+\bar B)^3}{(g'^{-2}-\bar B)^2},\qquad Q_0(g')=\exp\bigl[\tfrac18\ell_{g'}^3(2C_{\rm S}+C_{g'})\bigr],
\]
where $C_{\rm S}$ and $C_{g'}$ are the numbers of the bounds of Theorems~\YMMref{AJ}{673} and~\YMMref{AD}{607} on the partition functions of the tori, $Q_1(g')\ge1$ is the bound on the linking cumulants of step~(4) of the proof, and $c_{\rm a},C_{\rm a}>0$ depend only on the constants. So $d_\star>0$ and $C_2^{\rm T}(g')<\infty$ depend only on the radius, the constants and $g'$ (and on $\vartheta$, for $d_\star$), and on no $K$, $\beta$, side, state, pair or position.
\end{enumerate}
\end{proposition}

\begin{proof}[Proof sketch]
Let $\ell'=\ell_{g'}$, let $x_1$ be a centre of a ball of radius $d$ that holds ${\rm supp}f$, and consider on $a_K\mathbb Z^3$ the board of cubic cells of side $\ell'/2$ one of which, $\square_0$, is centred within $\frac12a_K$ of $x_1$; on $T_{\ell'}$ there are eight cells, the mirror images of $\square_0$ under lattice reflections.

(1)

\emph{One cell against the least torus.} For a bounded function $F$ of the configuration that depends only on the bonds in the interior of $\square_0$ put $\|F\|_{\rm ch}^8=\mu'(\prod_{\mathbf k\in\{0,1\}^3}F^{[\mathbf k]})\ge0$, the $F^{[\mathbf k]}$ being the mirror copies of $F$ in the eight cells, conjugated when the number of reflections is odd. This is the chessboard seminorm $\|F\|_{{\rm ch},\ell'}$ of the side $\ell'$. Then $|\mathfrak s(F)|\le Q_0\|F\|_{\rm ch}$ for every state at the datum. For the torus of a side $\ell\ge\ell'$ of $\mathfrak L(g')$ the chessboard estimate at the mirrors of the cells gives $|\mathfrak s(F)|\le\|F\|_{{\rm ch},\ell}$ (clause~(c) of Corollary~\ref{IV-18-18}), and the comparison of two sides gives $\|F\|_{{\rm ch},\ell}\le\|F\|_{{\rm ch},\ell'}\,e^{(\ell'/2)^3(p_{\ell'}-p_\ell)}$ with the pressures $p_\ell=\ell^{-3}\log Z_{K,\ell,\beta}$ (clause~(e) there); Theorems~\YMMref{AJ}{673} and~\YMMref{AD}{607} at the source $0$ bound $p_{\ell'}-p_\ell$ by $2C_{\rm S}+C_{g'}$, so that $e^{(\ell'/2)^3(p_{\ell'}-p_\ell)}\le Q_0$.

A periodic Gibbs state $\mathfrak s$ is dominated through tori of growing sides: by the domination of the means of a periodic Gibbs state by the chessboard seminorms of the tori (Proposition~\YMMref{91.9}{705}, for the board above, the one cell $\square_0$ and a sequence of sides $\ell_k\in\mathfrak L(g')$ with $\ell_k\to\infty$), $|\mathfrak s(F)|\le\liminf_k\|F\|_{{\rm ch},\ell_k}$, and every $\|F\|_{{\rm ch},\ell_k}$ is at most $Q_0\|F\|_{\rm ch}$ by the comparison of two sides.

(2)

\emph{Normalised moments; only linking cumulants remain.} For bounded real $F_1,F_2$ localised in $\square_0$ let $X_a=\mathcal O_K(F_a)$, $M'(z)=\mu'(e^{z_1X_1+z_2X_2})$, $M_{\mathfrak s}(z)=\mathfrak s(e^{z_1X_1+z_2X_2})$, and $c_I=\partial_I(1/M')(0)$, $I\subset\{1,2\}$, which are real. For $\emptyset\ne B\subset\{1,2\}$ the function $P_B=\sum_{I\subset B}c_{B\setminus I}\prod_{a\in I}X_a$ has $\mathfrak s(P_B)=\partial_B(M_{\mathfrak s}/M')(0)$. Its mirror copies are the same polynomials in the mirror images $X_{a,\mathbf k}=\mathcal O_K(F_a\circ\hat{\mathsf r}^{-1}_{\mathbf k})$, and as $\mu'$ is invariant under the lattice isometries,
\[
\|P_B\|_{\rm ch}^8=\mu'\Bigl(\prod_{\mathbf k}P_B^{[\mathbf k]}\Bigr)=\sum_\pi\prod_{\gamma\in\pi}\bigl\langle X_{a,\mathbf k};\,(\mathbf k,a)\in\gamma\bigr\rangle_{\mu'},
\]
the sum running over the partitions $\pi$ of $\{0,1\}^3\times B$ into \emph{linking} blocks, blocks that hold images in two different cells: all cumulants within one cell cancel. So if every linking cumulant is at most $Q_1\prod_{(\mathbf k,a)\in\gamma}u_a$, then $|\partial_B(M_{\mathfrak s}/M')(0)|\le Q_0\mathsf{Bl}_{16}^{1/8}Q_1^{|B|/2}\prod_{a\in B}u_a$, $\mathsf{Bl}_N$ the number of partitions of a set of $N$ elements, and the two cumulants $\partial_1\partial_2\log M_{\mathfrak s}(0)$, $\partial_1\partial_2\log M'(0)$ differ by $|\partial_1\partial_2\log(M_{\mathfrak s}/M')(0)|\le2(Q_0\mathsf{Bl}_{16}^{1/8}Q_1)^2u_1u_2$.

(3)

\emph{Packing and signs.} Let $\iota=a_K\lceil18d/a_K\rceil$ and $\mathsf{Pk}=\{v\in\iota\mathbb Z^3:|v_\mu|\le0.575-9d\}$. The balls $\mathbb B(x_1+v,9d)$, $v\in\mathsf{Pk}$, are disjoint and lie in $\mathbb B(x_1,1)\subset\square_0$, both supports of $\Psi$ lie in $\mathbb B(x_1,9d)$, and $\#\mathsf{Pk}\ge(0.3/9d)^3=(30d)^{-3}$. For signs $\epsilon=(\epsilon_v)$ apply~(2) to $F_1=\sum_v\epsilon_v\tau_vf$, $F_2=\sum_v\epsilon_v\tau_vh$. The difference $\mathsf D=S^T_2[\mathfrak s]-S^T_2[\mu']$ is bilinear, its mean over the signs at $(F_1,F_2)$ is $\sum_v\mathsf D(\tau_vf,\tau_vh)$, and both cumulants are invariant under the translations of $a_K\mathbb Z^3$ (Section~\ref{sec:state-limit-conventions}); so $|\mathsf D(f,h)|\le(30d)^3\sup_\epsilon|\mathsf D(F_1,F_2)|$.

(4)

\emph{The linking cumulants on the least torus.} For a complex source $w$ on the plaquettes of $T_{\ell'}$ let $Z[w]=\int\rho_0\,e^{\mathcal O(w)}\,dU$ be the partition function of $T_{\ell'}$ with the source $w$, $\rho_0$ the Wilson density at $(K,\beta)$ and $\mathcal O(w)(U)=L^K\sum_pw(p)\,a_p(U)$, and $Z=Z[0]$. The generating function of the sixteen mirrored, packed and signed observables is $Z[W_z]/Z$, with the source $W_z=\sum_{\mathbf k,a}z_{\mathbf k,a}(F_a\circ\hat{\mathsf r}_{\mathbf k}^{-1})^{\langle\ell'\rangle}\circ x$; it vanishes off the plaquettes whose barycentres lie in the eight mirror images of $\mathbb B(x_1,1)$, closed balls of radius $1$ whose centres are at pairwise sup distances $\ge\frac12\ell'-a_K$. Write
\[
\log\frac{Z[W]}Z=\log\frac{Z[W]e^{-\mathcal E(W)}}Z+\Bigl[\mathcal E(W)-\sum_{\mathbf k}\mathcal E(w^{(\mathbf k)})\Bigr]+\sum_{\mathbf k}\mathcal E(w^{(\mathbf k)}),
\]
$\mathcal E$ the source energy of $T_{\ell'}$, the sum over the levels and localisation domains of the source energies \eqref{eq:running-source-energies} of the sourced renormalisation-group chain \eqref{eq:running-source-recursion}, and $w^{(\mathbf k)}$ the part of $W$ in the cell $\mathbf k$.

(4a) \emph{The setting on $T_{\ell'}$.} Let $\ell'=2L^m$, $m=m_{\mathfrak L}(g')$. By~\YMMref{(1528)}{615}, $L^m=LM\max\{R(L(g'^{-2}+\bar B)),L^{\lceil\log_LC_{\rm s}\rceil}\}$, and the values of the cube-size function $R$ are powers $L^{\mathsf u}$ with $\mathsf u\ge1$; so the cells have the side $\mathsf l:=\frac12\ell'\ge L^2M$. As $g'\le\gamma_W^{\rm sd}$, the least side is small over the run: $MR(x_K)\le L^m$ (Lemma~\YMMref{126.3}{846}). As $K\ge3$, Lemma~\YMMref{126.2}{845}, whose head asks $K\ge3$, gives $x_j\ge x_{\rm J}$ for $0\le j\le K$ and makes the cuts of the run, on the lattice $(\ell',K)$, frames of the construction; so the setting (J) holds with the last height $K-1$ (Proposition~\ref{prop:weighted-canonical-history}), and the canonical sourced chain of $T_{\ell'}$ is formed, its conditions holding by Corollary~\ref{cor:canonical-running-source}, at the sources of a class of small sources, an open neighbourhood of $0$.

By the running estimate~\YMMref{(1530)}{615} at the index $K$ and the endpoint bounds~\YMMref{(1529)}{615}, $b_-\le a_{K-j}x_j\le b_+$ for $0\le j\le K$, with $b_\pm=x_K\pm\bar B$; and $g'^{-2}-\bar B<b_-$, $b_+\le X_{g'}+\bar B$. The places, the classes of the sum rule, the image sets and the chart integrals of the steps, the backgrounds and the minimal orbits that the setting (J) names (Lemma~\ref{lem:weighted-passing-fields} and Proposition~\ref{prop:weighted-history-step}) are the objects that the sourced chain names (Proposition~\ref{prop:running-source-density} and its one-step proof): both use the same definitions (among them Definition~\ref{def:sf-chart-data} and equation~\eqref{eq:sf-chart-integral} of the chart integral and Definition~\ref{def:sf-chart-data} of the image set), at the same lattice $(\ell',K)$, the same height $j$ and the same coupling $\mathsf t=x_j$.

(4b)

\emph{The level-wise norm; the families and the energies.} For a source $w$ let $v^w_j=-L^{K-j}w/x_j$ be its weight at the level $j$. The \emph{level-wise norm} $\|w\|_{\rm lw}$ is the largest, over $0\le j<K$, of $|v^w_j|_\infty/\sigma_{\rm a}$ and of $\operatorname{sw}_j(v^w_j)/(\mu_{\rm a}x_j^{-1/2})$, where $\operatorname{sw}_j(v)$ is the largest difference of the values of $v$ at two plaquettes whose barycentres are at a sup distance $\le Ma_{K-j}$, and $\sigma_{\rm a}=\sigma_{\rm src}$, $\mu_{\rm a}x_j^{-1/2}=o_{\rm src}(x_j)$ are the radii of the class of the weights of \eqref{eq:source-weight-room}, $\sigma_{\rm a}$ and $\mu_{\rm a}$ depending only on the constants.

It is a norm; its open unit ball $\mathsf{Src}^{\rm lw}$ is convex, consists of the sources whose weights lie at every level in the class of the weights at which Proposition~\ref{prop:source-entry} is stated, and contains the class of small sources of~(4a); and $|w|_\infty\le\sigma_{\rm a}b_+\|w\|_{\rm lw}$, as $|v^w_j|_\infty=|w|_\infty/(a_{K-j}x_j)$. Proposition~\ref{prop:running-source-chain} runs the induction on $\mathsf{Src}^{\rm lw}$; consequently the sourced families of all levels are formed there, holomorphic in the source, bounded by $E_0$, a constant of the canonical list, in the norm of the localised terms at the rate $\kappa+1$, and the family on a domain reads the source on the plaquettes of that domain alone.

So the source energies $e_j(X;w)$ of the domains $X$ of the level $j$ are holomorphic on $\mathsf{Src}^{\rm lw}$, vanish at $w=0$ and read $w$ on the plaquettes of $X$ alone (equation~\eqref{eq:running-source-energies} and its proof), and, by~\eqref{eq:running-source-energy-bound} at the bound $2E_0$ of the differences of the families and by Schwarz's lemma for $\zeta\mapsto e_j(X;\zeta w/\|w\|_{\rm lw})$ on the unit disc,
\[
|e_j(X;w)|\le16M^3E_0\,\|w\|_{\rm lw}\,e^{-\kappa d_j(X)},
\]
$d_j(X)$ the linear size of $X$ in the units of the level $j$. Put $\mathcal E_t(w)=\sum_{j\le t}\sum_{X\in\mathcal D_j}e_j(X;w)$, $\mathcal D_j$ the localisation domains of the level $j$; so $\mathcal E=\mathcal E_K$, and these functions are holomorphic on $\mathsf{Src}^{\rm lw}$ and vanish at $0$.

(4c)

\emph{The term without large fields, with a source.} For $0\le t\le K$ let $h_t$ be the history without large-field regions up to the level $t$ and $\tau_{h_t}[\cdot]$ its term (Section~\YMMref{123}{828}), a complex Borel measure on the configurations $\mathcal V_t$ of the level $t$, and $\mathsf m_t=\tau_{h_t}[\rho_0]$: a positive measure of mass at most $Z$ (equation~\eqref{eq:ward-first-creation}), equal to $\rho_0\,dU$ at $t=0$ and carried, for $t\ge1$, by the open set $\mathcal V^\chi_t$ of the configurations that pass the density tests of the level (Proposition~\YMMref{123.1}{829}). For $w\in\mathsf{Src}^{\rm lw}$ and $0\le t\le K$,
\[
\tau_{h_t}\bigl[\rho_0e^{\mathcal O(w)}\bigr]=e^{\mathcal E_t(w)}\,e^{\mathsf R_t(w,\cdot)}\,\mathsf m_t ,
\]
an identity of complex Borel measures on $\mathcal V_t$, with $\mathsf R_0(w,\cdot)=\mathcal O(w)$ and, for $t\ge1$ and $W\in\mathcal V^\chi_t$,
\[
\mathsf R_t(w,W)=\mathcal O(w)(U)+S[\mathfrak f_t[w]](U)-S[\mathfrak f_t[0]](U),
\]
$U$ an element of the minimal orbit over $W$ and $S[\mathfrak f_t[w]]$ the action of the sourced bare array of the level $t$; this value does not depend on the element of the orbit, the domain of definition of the action being a condition on the configuration alone, and $\mathsf R_t(w,\cdot)$ is a bounded Borel function on $\mathcal V^\chi_t$. For the small sources of~(4a) and $1\le t\le K$, the density of $\tau_{h_t}[\rho_0e^{\mathcal O(w)}]$ is
\[
\exp\Bigl[\mathsf C_t(w)-A^{\eta_t}\bigl(x^w_t,U_t(W)\bigr)+S[\mathfrak f_t[w]]\bigl(U_t(W)\bigr)\Bigr]\quad\text{on }\mathcal V^\chi_t,\ \text{and $0$ off it},\qquad\mathsf C_t(w)=\sum_{1\le j\le t}\mathsf c_j+\mathcal E_t(w),
\]
with numbers $\mathsf c_j$ formed with the source $0$ alone, the weight $x^w_t=x_t-L^{K-t}w$ of the weighted action $A^{\eta_t}$, $\eta_t=L^{-t}$, and the background $U_t(W)$ over $W$. This follows by induction over $t$. At $t=0$ the term is $\rho_0e^{\mathcal O(w)}dU$, and $\rho_0e^{\mathcal O(w)}=e^{-A^{\eta_0}(x^w_0,\cdot)}$ (equation~\eqref{eq:running-weighted-action}). From $t$ to $t+1$: the term of the history is carried by one step, $\tau_{h_{t+1}}[\cdot]=\mathbf T(\mathsf g_{t+1}\tau_{h_t}[\cdot])$ with the factor $\mathsf g_{t+1}\ge0$ of the step, and a measure whose density on $\mathcal V^\chi_t$ is $e^{\mathsf e}$ times a bounded Borel function $F$ is carried to the density $e^{\mathsf e}\,\mathsf I[F]$ on $\mathcal V^\chi_{t+1}$, and $0$ off it, $\mathsf I[F]$ being the chart integral of the step (Proposition~\ref{prop:weighted-history-step}); here, for $t\ge1$, $\mathsf e=\mathsf C_t(w)$ and $F$ is the function $F^w_t$ in the proof of Proposition~\ref{prop:running-source-density} at $k=t$, which on $\mathcal V^\chi_t$ is $\exp[-A^{\eta_t}(x^w_t,U_t(V))+S[\mathfrak f_t[w]](U_t(V))]$ (the definition preceding \eqref{eq:running-source-one-step}), and for $t=0$, $\mathsf e=0$ and $F=\rho_0e^{\mathcal O(w)}$, whose chart integral is that of $F^w_0$ on the class of the sum rule (the height-zero case in that proof); $\mathcal V^\chi_{t+1}$ lies in the class of the sum rule of the step and in the set on which the chart integral is formed (Lemma~\ref{lem:weighted-passing-fields}); and there the one-step identity of the sourced chain \eqref{eq:running-source-one-step} gives
\[
\mathsf I[F^w_t](W)=\exp\Bigl[\mathsf c_{t+1}+\sum_{X\in\mathcal D_{t+1}}e_{t+1}(X;w)-A^{\eta_{t+1}}\bigl(x^w_{t+1},U_{t+1}(W)\bigr)+S[\mathfrak f_{t+1}[w]]\bigl(U_{t+1}(W)\bigr)\Bigr].
\]
The weighted action is linear in the weight, $A^{\eta_t}(x^w_t,U)=A^{\eta_t}(x_t,U)-\mathcal O(w)(U)$; so the quotient of the density by the same density at $w=0$, which is the density of $\mathsf m_t$, is $e^{\mathcal E_t(w)+\mathsf R_t(w,\cdot)}$ on $\mathcal V^\chi_t$: the identity holds for the small sources. For $w\in\mathsf{Src}^{\rm lw}$, $w\ne0$, and a bounded Borel function $\psi$ on $\mathcal V_t$, the function $\zeta\mapsto\int\psi\,d\tau_{h_t}[\rho_0e^{\zeta\mathcal O(w)}]$ is entire, and $\zeta\mapsto e^{\mathcal E_t(\zeta w)}\int\psi\,e^{\mathsf R_t(\zeta w,\cdot)}\,d\mathsf m_t$ is holomorphic on $\{|\zeta|\,\|w\|_{\rm lw}<1\}$, by~(4b); they agree for small $\zeta$, hence at $\zeta=1$.

(4d)

\emph{The remainders.} Let $w\in\mathsf{Src}^{\rm lw}$ vanish off the plaquettes whose barycentres lie in $N_{\rm b}$ closed balls of radius $1$. Then
\[
\begin{multlined}
|\mathsf R_K(w,\cdot)|\le\|w\|_{\rm lw}N_{\rm b}\,\mathsf g_\emptyset\ \relax \ \relax \text{on }\mathcal V^\chi_K,\qquad|\mathsf R_{n-1}(w,\cdot)|+|\mathcal E(w)-\mathcal E_{n-1}(w)| \\
\le\|w\|_{\rm lw}N_{\rm b}\,C_{\rm cnt}\,L^{4(K-n+1)}\ \relax \ \relax \text{on }\mathcal V^\chi_{n-1}\quad(1\le n\le K),
\end{multlined}
\]
on all of $\mathcal V_0$ for $n=1$, with numbers $\mathsf g_\emptyset$ and $C_{\rm cnt}\ge1$ that depend only on $g'$ and the constants. Indeed, let $W\in\mathcal V^\chi_s$. The background over $W$ has $0\le a_p(U_s(W))<2B_3^2\varepsilon_s^2L^{-4s}$ at every bare plaquette~\eqref{eq:weighted-background-plaquettes}, $\varepsilon_s=\varepsilon(x_s)\le\varepsilon(g'^{-2})$ being the threshold of the small fields at the coupling $x_s$ and $B_3$ a constant of the canonical list; at most $81N_{\rm b}L^{3K}$ plaquettes have their barycentres in the balls; and $|w|_\infty\le\sigma_{\rm a}b_+\|w\|_{\rm lw}$. So $|\mathcal O(w)(U_s(W))|\le162N_{\rm b}B_3^2\varepsilon_s^2\sigma_{\rm a}b_+L^{4(K-s)}\|w\|_{\rm lw}$, and at $s=0$, where $0\le a_p\le2$, $|\mathcal O(w)|\le162N_{\rm b}\sigma_{\rm a}b_+L^{4K}\|w\|_{\rm lw}$.

The difference of the two actions is the sum, over the levels $j\le s$, the domains $X$ of the level $j$ that meet the support of $w$ and their points $z$, of the differences between the values at $U$ and at $\mathbf 1$ of the local terms $(E^{(j)}[w]-E^{(j)}[0])(X,\cdot,z)$, each of a norm at most $2E_0\|w\|_{\rm lw}e^{-(\kappa+1)d_j(X)}$ by~(4b). Every element of the minimal orbit over $W$ is regular of the level $s$ at the radius $2B_3\varepsilon_s$ (Lemma~\ref{lem:weighted-passing-fields}); at such a configuration a local term of the level $j$ moves from its value at $\mathbf 1$ by at most its norm times $C(\varepsilon_s/\alpha_0)^2L^{-4(s-j)}(1+d_j(X))^2$ on the domains that are small against the radius, and by at most twice its norm on the other ones, whose linear sizes are large (equation~\YMMref{(2263)}{828} and its proof), $\alpha_0$ being the fixed radius of the localised terms; the bound in Lemma~\ref{lem:source-support-sum}, with at most $27N_{\rm b}L^{3(K-s)}$ cubes of the level $s$ holding a corner of a plaquette of the support, sums these differences to at most $CN_{\rm b}L^{3(K-s)}\|w\|_{\rm lw}$.

The energies of the levels $>s$ are bounded in the same way, by~(4b) and~(q3), by $CN_{\rm b}L^{3(K-s)}\|w\|_{\rm lw}$. With $b_+\le X_{g'}+\bar B$ this gives the two bounds.

(4e) \emph{The coefficients of $Z[w]e^{-\mathcal E(w)}/Z$.} By \eqref{eq:ward-first-creation} at level $K$, applied to the bounded gauge invariant function $e^{\mathcal O(w)}$, with the canonical history measures of Section~\YMMref{123}{828},
\[
Z[w]=Z_\emptyset[w]+\sum_{n=1}^KZ_n[w],\qquad Z_\emptyset[w]=\tau_{h_K}\bigl[\rho_0e^{\mathcal O(w)}\bigr](\mathcal V_K),\qquad Z_n[w]=\int\mathsf g^{\rm w}_n\,d\tau_{h_{n-1}}\bigl[\rho_0e^{\mathcal O(w)}\bigr],
\]
where $\mathsf g^{\rm w}_n\ge0$ is the factor of the step $n$ at which a large-field region is first created, and $\mathsf w_n=Z^{-1}\int\mathsf g^{\rm w}_n\,d\mathsf m_{n-1}$ is the weight of that event; each term is an entire function of $w$. By~(4c),
\[
Z_\emptyset[w]e^{-\mathcal E(w)}=\int e^{\mathsf R_K(w,\cdot)}\,d\mathsf m_K,\qquad Z_n[w]e^{-\mathcal E(w)}=e^{-(\mathcal E-\mathcal E_{n-1})(w)}\int\mathsf g^{\rm w}_n\,e^{\mathsf R_{n-1}(w,\cdot)}\,d\mathsf m_{n-1},
\]
so that, by~(4d) and as $\mathsf m_K$ has a mass at most $Z$,
\[
\bigl|Z_\emptyset[w]e^{-\mathcal E(w)}\bigr|\le e^{\|w\|_{\rm lw}N_{\rm b}\mathsf g_\emptyset}\,Z,\qquad\bigl|Z_n[w]e^{-\mathcal E(w)}\bigr|\le\exp\bigl[\|w\|_{\rm lw}N_{\rm b}C_{\rm cnt}L^{4(K-n+1)}\bigr]\,\mathsf w_n\,Z .
\]
At $w=W_z$ take $N_{\rm b}=8$ and the polydisc $\mathbb P=\{|z_{\mathbf k,a}|<r_a\}$, $r_a=A_\star(d)/\|f_a\|_\infty$, on which $\|W_z\|_{\rm lw}\le\max_{\mathbf k,a}|z_{\mathbf k,a}|/r_a<1$ by~(4g). The function $\Phi_{\rm D}(z)=Z[W_z]e^{-\mathcal E(W_z)}/Z$ is holomorphic on $\mathbb P$, $Z[w]$ being entire, and $\Phi_{\rm D}(0)=1$; its term without large fields is at most $e^{8\mathsf g_\emptyset}$ on the whole of $\mathbb P$, and its term of the level $n$ at most $e\,\mathsf w_n$ on the polydisc smaller by the factor $8C_{\rm cnt}L^{4(K-n+1)}$. Cauchy's estimate on each polydisc gives, for every set $J$ of at most sixteen of the variables,
\[
|\partial_J\Phi_{\rm D}(0)|\le\Bigl[e^{8\mathsf g_\emptyset}+e\,(8C_{\rm cnt}L^4)^{16}\sum_{n=1}^KL^{64(K-n)}\mathsf w_n\Bigr]\prod_{i\in J}r_i^{-1}.
\]
By \eqref{eq:ward-first-creation}, $\mathsf w_n\le\widetilde\nu_K\{Z_n\ne\emptyset\}$. Equation~\eqref{eq:ward-weighted-tail}, at $q=64$ and $t=K$, therefore bounds $\sum_nL^{64(K-n)}\mathsf w_n$ by
\[
 4M^{-3}(\ell')^3e^{\mathsf B(X_{g'})(\ell')^3}
          \mathsf C_{\rm ag}(65)e^{-c'p_0(g_K)^2/2}.
\]
The assumptions hold by the canonical choice of the fixed history thresholds in Proposition~\ref{prop:ward-fixed-volume-rarity}; the ceiling $g'\le\gamma_W^{\rm sd}$ preserves them. Since $\ell'$ depends only on $g'$ and the constants, the square bracket is bounded by a number $Q_{\rm D}(g')\ge1$ with the same dependence.

(4f)

\emph{The linking cumulants.} The last sum of the decomposition is a sum of functions each of which reads the variables of one cell; it is annulled by every linking derivative. In the bracket only the domains $X$ that meet the supports of two parts $w^{(\mathbf k)}$ contribute, the source energies being local and vanishing at $0$. Such a domain of the level $j$ holds corners of plaquettes whose barycentres lie in two different balls, and two such corners are at a sup distance $\ge\mathsf l-3$, while two points of a domain of the level $j$ are at a sup distance at most $M(d_j(X)+2)a_{K-j}$; so $d_j(X)\ge(\mathsf l-3)L^{K-j}/M-2\ge166L^{K-j}$, as $\mathsf l\ge L^2M$ and $L\ge13$. The bound of~(4b) at the nine sources $W$, $w^{(\mathbf k)}$ and the bound in Lemma~\ref{lem:source-support-sum} give
\[
\Bigl|\mathcal E(W)-\sum_{\mathbf k}\mathcal E(w^{(\mathbf k)})\Bigr|\le C_\Delta\max_{\mathbf k}\bigl\{\|W\|_{\rm lw},\|w^{(\mathbf k)}\|_{\rm lw}\bigr\},
\]
this exponential decay of the source energies in the size of their domains summing over the levels to a number $C_\Delta$ depending only on the constants. On $\mathbb P$ the bracket is holomorphic and at most $C_\Delta$, so that Cauchy's estimate bounds its squarefree coefficients by $C_\Delta\prod_ir_i^{-1}$. A derivative $\partial_J$ at $0$ of the logarithm of a function equal to $1$ at $0$ is the sum, over the partitions $\pi$ of $J$, of $(-1)^{|\pi|-1}(|\pi|-1)!$ times the product over the blocks of its derivatives; so the squarefree coefficients of $\log\Phi_{\rm D}$ are at most $\mathsf{cu}_{|J|}Q_{\rm D}^{|J|}\prod_{i\in J}r_i^{-1}$, with $\mathsf{cu}_N=\sum_\pi(|\pi|-1)!$ over the partitions of $N$ elements.

(4g)

\emph{The packed bumps.} For a source made of bumps with pairwise disjoint sets of non-vanishing, of height at most $A$ and Lipschitz constant at most $A/(\vartheta d)$, the supremum of the weight of the level $j$ is at most $A/b_-$, and its oscillation at most $(A/b_-)\min\{2,2\sqrt3Ma_{K-j}/(\vartheta d)\}$; as $a^{-1/2}\min\{2,ca\}\le(2c)^{1/2}$ for all $a>0$, its level-wise norm is at most $A/A_\star(d)$ at every level, with
\[
A_\star(d)=\min\Bigl\{\sigma_{\rm a}b_-,\ \mu_{\rm a}\frac{b_-}{b_+^{1/2}}\Bigl(\frac{\vartheta d}{4\sqrt3M}\Bigr)^{1/2}\Bigr\}\ge c\,(g'^{-2}-\bar B)(X_{g'}+\bar B)^{-3/2}(\vartheta d)^{1/2}\,b\qquad(d\le d_\star(g',\vartheta)),
\]
$d_\star$ being the scale below which the second member of the minimum is the smaller one. The bumps of $W_z$ are the mirror images of the $z_{\mathbf k,a}\epsilon_v\tau_vf_a$; their sets of non-vanishing are pairwise disjoint, their heights are at most $A=\max_{\mathbf k,a}|z_{\mathbf k,a}|\,\|f_a\|_\infty$ and their Lipschitz constants at most $A/(\vartheta d)$, as ${\rm Lip}f_a\le\|f_a\|_\infty/(\vartheta d)$; so $\|W_z\|_{\rm lw}\le\max_{\mathbf k,a}|z_{\mathbf k,a}|/r_a$.
So, by~(4f), every linking cumulant is at most $(\mathsf{cu}_{16}Q_{\rm D}^{16}+C_\Delta)\prod_{(\mathbf k,a)\in\gamma}r_a^{-1}$: the linking cumulants obey the bound of~(2) with $u_a=\|f_a\|_\infty/A_\star(d)$ and $Q_1=\mathsf{cu}_{16}Q_{\rm D}^{16}+C_\Delta$.

(5)

\emph{Conclusion.} By (1)--(4), 
\[
|S^T_2[\mathfrak s](\Psi)-S^T_2[\mu'](\Psi)|\le(30d)^3\,2(Q_0\mathsf{Bl}_{16}^{1/8}Q_1)^2A_\star(d)^{-2}|\Psi|\le C_2^{\rm T}(g')\vartheta^{-1}d^2b^{-2}|\Psi|
\]
 for $d\le d_\flat$, which is~(a) with the constants of~(b); two states differ by at most twice this.
\end{proof}

\begin{remark}
The proof uses only the following properties of the pair: real Lipschitz functions, neither $0$, each vanishing off a ball of radius $d$, whose sets of non-vanishing are disjoint and lie in one ball of radius $9d$, with ${\rm Lip}f_a\le\|f_a\|_\infty/(\vartheta d)$. So Proposition~\ref{C-351-7} holds also for such pairs of either sign. With $n$ entries in the place of two, every pair of which lies in $\mathfrak T_{\mathbb R^3}(d;\vartheta)$, the same argument gives a difference of at most $2C_n^{\rm T}(g')\vartheta^{-n/2}d^{3-n/2}b^{-n}\prod_a\|f_a\|_\infty$, which is small against the natural size $b^{-n}\prod_a\|f_a\|_\infty$ for $n\le5$; we use only $n=2$.
\end{remark}

\subsection{The coefficient on the least torus}

Here the constants are an admissible assignment of the canonical list. The statements used in Lemmas~\ref{C-351-8}--\ref{C-351-12} are made for the data (Wp1)--(Wp6) of a weighted progression (Definition~\ref{def:weighted-progression-data}), in its setting (J) (Definition~\ref{def:weighted-history-setting}) and, where a step is continued to a complex field scale, at a place with a lower place (Definition~\ref{IV-71-12}); we call this the setting of the weighted progression (the statements used in Lemma~\ref{C-351-13}, the constants and thresholds in Theorem~\YMMref{AR}{810}, Theorem~\YMMref{AR}{810}, equation~\YMMref{(2203)}{810}, Corollary~\YMMref{117.1}{812}, Lemma~\ref{lem:cutoff-admissible-labels} and Corollary~\YMMref{118.4}{815}, are made in the setting (Th1), (Th2), whose datum is the choice of the constants).

For the canonical objects at every admissible assignment of the canonical list it holds: the data are given by Corollary~\ref{cor:weighted-canonical-data}; the setting (J) holds by Proposition~\ref{prop:weighted-canonical-history} as soon as the couplings are at least the threshold $x_{\rm J}$ of the definitions preceding \eqref{eq:weighted-history-threshold}, a number depending only on the constants, and at a first-passage datum at $g$ this is the case whenever $g^{-2}\ge x_{\rm J}$, all the couplings of the run exceeding $g^{-2}$ (Lemma~\ref{C-351-8}); and the places with a lower place are those of Proposition~\ref{IV-71-50}. Besides the constants the setting involves two auxiliary real numbers: the lower place $\mathsf t_\flat\ge\gamma^{-2}$, and a room $\mathsf q>0$, such that a bound on the sizes of the bare arrays formed with the constants alone, increased by $2\mathsf q$, does not exceed the size ceiling at the lower place.

For the canonical objects $\mathsf q=1$, and $\mathsf t_\flat$ is a number depending only on the constants. No further condition on the constants enters. $\gamma_{\rm e}\in(0,\gamma]$ is the largest number under which the setting holds and all the conditions on $g$ of these statements hold (each has the form $g\le\gamma'$ or $\log g^{-2}\ge Y$ with numbers $\gamma'$, $Y$ depending only on the constants), together with $g^{-2}\ge4\bar B$; for $0<g\le\gamma_{\rm e}$, $d_{\rm e}(g,\vartheta)\in(0,L^{-13}]$ is the largest number below which their conditions on $d$ and those of the proof of Proposition~\ref{C-351-14} hold (the arithmetic of the levels there uses $d\le L^{-13}$, and the removal of the cutoff there uses $d\le d_{\rm K}$, the number of Lemma~\ref{C-351-13} at $\ell=\ell_g$).

The depth of the scale $d$ at $g$ is the witness depth $s(d)=\min\{s\ge0:L^sd\ge C_W(s)\}$, $C_W(s)=\mathsf c_W(\log[L^{s+1}(g^{-2}+\bar B)])^{2p_0}$, with a number $\mathsf c_W\ge1$ depending only on the constants (equation~\YMMref{(2321)}{843}); so $C_W(s(d))\le L^{s(d)}d<LC_W(s(d))$, $C_W(s(d))\ge(\log\frac1d)^2$, $s(d)\to\infty$ as $d\to0$, and $K>s(d)$ entails $d/a_K\ge L(\log\frac1d)^2$, hence $K\ge3$ and $8a_K\le d$. We write $D_j=a_{K-j}$ (the spacing of the level $j$), $\mathsf l=\max\{1,\log(2b/d)\}$, $\mathfrak b(x)=A_1(\log x)^{p_0}$ for the field cut-off, $A_1$ and the integer $p_0\ge1$ being constants of the canonical list, $\|u\|_\sharp=\max\{\|u\|_\infty,40M\,{\rm Lip}\,u\}$ and $\mathbb Z^3_N=\mathbb Z^3/N\mathbb Z^3$; $C$ denotes numbers that depend only on the constants, and on $g$ and $\vartheta$ where this is shown.

\begin{lemma}[The run]\label{C-351-8}
Let $0<g\le\gamma_{\rm e}$ and $(K,\beta)$ a first-passage datum at $g$ with $K\ge3$ and label $b$. Then $|x_j-L^{-(j-i)}x_i|\le\bar B$ for $0\le i\le j\le K$, $g^{-2}<x_K\le X_g$, and $\frac34b/D_j\le x_j\le2b/D_j$ for $0\le j\le K$. On every $\ell\in\mathfrak L(g)$ the setting of the weighted progression holds with the sources $f\circ x$, $h\circ x$ for all real bounded $f,h$ on $T_\ell$ supported in closed balls of radius $d$ at a distance $\ge d$, with $8a_K\le d$, neither $f\circ x$ nor $h\circ x$ vanishing identically; no other condition on $d$ enters. If $f\circ x$ or $h\circ x$ vanishes identically, then so does $(\tau_vf)\circ x$, respectively $(\tau_vh)\circ x$, for every $v\in a_K\mathbb Z^3$, and $S^T_{2;K,\ell,\beta}(f,h)$, the jets $\mathsf n_j(f,h)$ and $\mathsf n^{\mathsf z}_j(f,h)$ of Lemma~\ref{C-351-10} and every sum $\sum_{p,q}(f\circ x)(p)\,(h\circ x)(q)\,\mathsf k(p,q)$ against a kernel $\mathsf k$ on the plaquettes vanish.

So, for $b_*\in I(g)$, $K'\ge3$ and real functions $w_1,w_2$ on the plaquettes of $\mathbb Z^3_{\ell L^{K'}}$, neither identically $0$, supported in balls of radius $r\ge8$ at a distance $\ge r$, $w_a=f'_a\circ x$ for such a pair $(f'_1,f'_2)$ at the scale $ra_{K'}$ and the first-passage datum $(K',L^{K'}b_*)$.
\end{lemma}

\begin{proof}[Proof sketch]
The first two assertions are the running and endpoint estimates~\YMMref{(1529)--(1530)}{615}. As $L^{-j}x_0=b/D_j$ and $x_j\ge x_K>g^{-2}\ge4\bar B$: $b/D_j\le x_j+\bar B\le\frac43x_j$, and $x_j\le b/D_j+\bar B\le2b/D_j$ because $b/D_j\ge b>g^{-2}\ge\bar B$. The data of the weighted progression are those of the canonical objects (Corollary~\ref{cor:weighted-canonical-data}) at the calibrations $x_0,\dots,x_{K-1}$ of the run, with the lower place $\mathsf t_\flat$ and the room $\mathsf q=1$; the lattice $(\ell,K)$, with the run, is admissible over every cut of the run, for every $\ell\in\mathfrak L(g)$ (Proposition~\YMMref{74.1}{615}); and the setting (J) holds, as $x_j\ge x_K>g^{-2}\ge x_{\rm J}$ (Proposition~\ref{prop:weighted-canonical-history}).

The other conditions of the setting are lower bounds on $x_{K-1}$ and on $\log x_K$ by numbers formed with the constants and with $\mathsf t_\flat$; they hold as $x_{K-1}>X_g\ge Lg^{-2}$ and $x_K>g^{-2}$ once $g$ is small, and they enter $\gamma_{\rm e}$. Of the sources the setting reads only that they are real and bounded, supported in balls of radius $d$ at a distance $\ge d$, with $8a_K\le d$, and not identically $0$; a translation by $v\in a_K\mathbb Z^3$ permutes the barycentres of the plaquettes. If $f\circ x\equiv0$, then $\mathcal O_K(f)=0$, the source $\mu_1f\circ x+\mu_2h\circ x$ does not read $\mu_1$, so that the source energies of the weighted progression do not read it at any field scale, and a sum of the values $(f\circ x)(p)$ against anything vanishes.

For the last sentence let $f'_a$ take the values of $w_a$ at the barycentres of the plaquettes of spacing $a_{K'}$ and vanish elsewhere: they are real and bounded, supported in balls of radius $ra_{K'}$ at a distance $\ge ra_{K'}$, with $8a_{K'}\le ra_{K'}$, and no condition on the depth of the scale $ra_{K'}$ enters; $(K',L^{K'}b_*)$ is a first-passage datum since $b_*\in I(g)$.
\end{proof}

\begin{lemma}[The chart of the iterated block averages]\label{C-351-9}
Let $\ell=2L^m$, $1\le n\le K$ and $2L^{m+K-n+1}\ge C_LM$, $C_L$ a constant of the canonical list. For a block field $W$ of the level $n$ and fluctuation fields $y=(\mathcal A_j)_{j<n}$, $\mathcal A_j$ an $\mathfrak{su}(2)$-valued $1$-form on the non-pivot bonds of the level $j$, let $W_n=W$, $W_j=V_{W_{j+1}}[\mathcal A_j]$ (Section~\ref{sec:sf-chart-interface}), $\Upsilon(W,y)=W_0$, and $\theta_{\rm ch}=-\sum_{j<n}[\mathsf J_{V_j}[\mathcal A_j]+\log z^{(j)}_{V_j}]$, with the Jacobian term $\mathsf J$ and the normalisation $z^{(j)}$ of Lemma~\ref{lem:sf-chart-integral}, $V_j=V^{\rm bf}_j(W_{j+1})$ being the averaged background over the block field $W_{j+1}$ (Section~\ref{sec:sf-chart-interface}).

There are an open neighbourhood $\mathsf O_n$ of $\mathbf 1$ among the block fields of the level $n$ and a number $\rho_2\in(0,1]$, both defined by bounds that depend only on the constants and not on $n$, $K$ or $m$ (they are written out in the proof), such that, with $\Omega=\{|\mathcal A_j|<\rho_2,\ j<n\}$ and
\[
C_{\rm H}=\prod_{j<n}(2\pi^2)^{-\#\mathfrak X^+_j}L^{-3\#\mathfrak B^+_{j+1}},
\]
$\mathfrak X^+_j$ the positively oriented non-pivot bonds of the level $j$ and $\mathfrak B^+_{j+1}$ the positively oriented bonds of the level $j+1$:
\begin{enumerate}
\item[(i)] on $\mathsf O_n\times\Omega$ the chain and $\theta_{\rm ch}$ are formed, $W_{j+1}$ is the block average of $W_j$, and $\theta_{\rm ch}$ is real analytic;
\item[(ii)] $\Upsilon$ is a real analytic diffeomorphism of $\mathsf O_n\times\Omega$ onto an open set $\mathcal U\ni\mathbf 1$;
\item[(iii)] $\int H\,dU=C_{\rm H}\int_{\mathsf O_n}dW\int_\Omega dy\,e^{-\theta_{\rm ch}}H(\Upsilon(W,y))$ for every bounded Borel function $H$ that vanishes off $\mathcal U$, and the same holds with every open subset of $\mathsf O_n\times\Omega$ and its image.
\end{enumerate}
\end{lemma}

\begin{proof}[Proof sketch]
\emph{The numbers and the good sets.} For a configuration $V$ of a level let $\varpi(V)=\max_p|\partial V(p)-\mathbf 1|$, the maximum over the plaquettes of that level, and $\varpi_c(V)$ the same maximum over the plaquettes of the double block of a bond $c$ of the next level; $\mathcal M$ is the block average and $\operatorname{Lg}$ the $\mathfrak{su}(2)$-valued logarithm. With the numbers $\varpi_\chi$, $\varpi_1=2\varpi_\chi$, $r_\chi$, $C_\chi$ of the chart of the fibre (Section~\ref{sec:sf-chart-interface}), which depend only on $L$, the small-curvature radius $c^{\rm SO}_2$ of Proposition~\YMMref{14.1}{61} and the response constant $K_{\rm B}$ of Lemma~\ref{lem:background-curvature-response}, and the constants $B_3$, $\bar\varepsilon$, $\varepsilon_{\rm v}$, $r_{\rm B}$ of the canonical list, put
\[
\begin{multlined}
\mu_\flat=\min\Bigl\{c^{\rm SO}_2,\ \relax \tfrac12\bar\varepsilon,\ \relax \tfrac14L^2\varpi_\chi,\ \relax \tfrac18\varepsilon_{\rm v}\Bigr\},\qquad\rho_2=\min\Bigl\{\frac{r_\chi}{C_\chi},\ \relax \frac{r_{\rm B}}{4C_\chi},\ \relax \frac{(1-L^{-2})\,\mu_\flat}{4K_{\rm B}C_\chi},\ \relax \frac{\varepsilon_{\rm v}}{64C_\chi}\Bigr\},\qquad\mathsf O_n \\
=\Bigl\{W:\ \relax \varpi(W)<\frac{\mu_\flat}{B_3}\Bigr\};
\end{multlined}
\]
they read neither $n$ nor $K$ nor $m$, and $\rho_2\le r_\chi/C_\chi\le1$. For $1\le i\le n$ let $\mathsf{Gd}_i$ be the set of the block fields $W$ of the level $i$ over which the background exists and has a regularity number $\mu_i(W)<\mu_\flat$ (Definition~\ref{def:sf-chart-data}; the number is the same at all the elements of the minimal orbit, by unitary gauge invariance), and $\mathsf{Gd}_0=\{V:\varpi(V)<\frac14\varepsilon_{\rm v}\}$.

(1)

\emph{The averaged background over a good block field.} $\mathsf{Gd}_i$ is open and $\{\varpi<\mu_\flat/B_3\}\subset\mathsf{Gd}_i\subset\{\varpi<2\mu_\flat\}$ for $1\le i\le n$: over a block field with $\varpi<\mu_\flat/B_3$ the background exists and has a regularity number $<\mu_\flat$ (Theorem~\YMMref{E}{10} in the setting of Section~\YMMref{52.7}{432}), a regularity number $<\mu_\flat$ gives $\varpi<2\mu_\flat\le\bar\varepsilon$ (Lemma~\ref{lem:background-iterated-average}), and the representative of the background is a real analytic function of the block field (Theorem~\YMMref{E}{10} and Lemma~\ref{lem:background-curvature-response}); so $\mathbf 1\in\mathsf O_n\subset\mathsf{Gd}_n$. Let $0\le j<n$ and $W\in\mathsf{Gd}_{j+1}$. Then $V_j(W)=V^{\rm bf}_j(W)$ is defined and $\mathcal MV_j(W)=W$: it is the average $\mathcal M^jU$ of an element $U$ of the minimal orbit over $W$, carried to the axial gauge by a gauge transformation (the definition of $p(W)$ in Definition~\ref{def:sf-chart-data}), which conjugates the plaquette variables; so Lemma~\ref{lem:background-iterated-average} gives $\varpi(V_j(W))\le2L^{-2}\mu_{j+1}(W)<\frac12\varpi_\chi$.

The map $W\mapsto V_j(W)$ is real analytic on $\mathsf{Gd}_{j+1}$: $\mathcal M^jU^\circ(W)$, with the representative $U^\circ(W)$ of the background, is real analytic by Theorem~\YMMref{E}{10} and Lemma~\ref{lem:background-iterated-average}, on the good-set and actual-average domains just established, and the gauge transformation to the axial gauge is $v(x)=(\mathcal M^jU^\circ(W))(y,x)$, $x$ in the block of $y$, an averaged contour variable, a real analytic function of the configuration where it is formed. And $V_j(\mathbf 1)=\mathbf 1$: $\mathbf 1$ is flat, lies in the minimal orbit over $\mathbf 1$, and every path product and its average are $\mathbf 1$; uniqueness of the minimal orbit in Theorem~\YMMref{E}{10} identifies its axial representative with $\mathbf 1$.

(2)

\emph{One step.} Let $\mathsf O$ be an open subset of $\mathsf{Gd}_{j+1}$, $0<\rho\le r_\chi/C_\chi$, $G(\rho)=\{u\in SU(2):|\operatorname{Lg}u|<\rho\}$, and
\[
\mathcal N=\bigl\{V:\ \mathcal MV\in\mathsf O,\ \ V(\mathsf e)\,V_j(\mathcal MV)(\mathsf e)^{-1}\in G(\rho)\ \text{at every non-pivot bond }\mathsf e\bigr\}.
\]
At $W'\in\mathsf O$ the configuration $V_j(W')$, with $\varpi_c(V_j(W'))\le\varpi_\chi$, is a reference configuration of the chart of the fibre, and Proposition~\YMMref{53.2}{441}, with empty exterior, gives for $|\mathcal A|<\rho$: $V_j(W')$ is real and admissible, with invertible pivot coefficients; the chart is formed at $\mathcal A$, $V_{W'}[\mathcal A]=e^{iB'}V_j(W')$ with a real $B'$, $\sup|B'|\le2C_\chi|\mathcal A|$; $\mathcal MV_{W'}[\mathcal A]=W'$; $\operatorname{Lg}(V_{W'}[\mathcal A](\mathsf e)V_j(W')(\mathsf e)^{-1})=\mathcal A(\mathsf e)$ at the non-pivot bonds; $\varpi_c(V_{W'}[\mathcal A])<\varpi_1$; and $\mathsf J_{V_j}[\mathcal A]$, $z^{(j)}_{V_j}$ are formed, real and continuous (the continuity statement and the local bounds in its proof).

By the converse in that proposition, $\Phi_j(W',\mathcal A)=V_{W'}[\mathcal A]$ is a bijection of $\mathsf O\times\{|\mathcal A|<\rho\}$ onto the open set $\mathcal N$, with the inverse $\Xi_j(V)=(\mathcal MV,(\operatorname{Lg}(V(\mathsf e)V_j(\mathcal MV)(\mathsf e)^{-1}))_{\mathsf e})$. The map $\Xi_j$ is real analytic on $\mathcal N$: $\mathcal M$ is real analytic at configurations whose plaquettes in the double blocks are small, $V_j$ is by~(1), and $\operatorname{Lg}$ is the logarithmic series where $|u-\mathbf 1|<1$. Its differential is injective: if the derivatives of $\mathcal MV$ and of the second member of $\Xi_j(V)$ vanish along a tangent vector $\dot a$, then $\dot a=0$ at the non-pivot bonds, and then the derivative of $(\mathcal MV)(c)$ at a bond $c$ of the level $j+1$ is $\mathsf n_c\,\dot a(b_0(c))$, with the invertible kernel $\mathsf n_c$ at the pivot bond $b_0(c)$ (the pivot differential preceding equation~\YMMref{(911)}{420}), so that $\dot a=0$.

The dimensions agree, the positively oriented bonds of the level $j$ being the non-pivot ones and one pivot bond for each positively oriented bond of the level $j+1$. By the inverse function theorem $\Xi_j$ is a local real analytic diffeomorphism; being injective, it has the real analytic inverse $\Phi_j$, a real analytic diffeomorphism of $\mathsf O\times\{|\mathcal A|<\rho\}$ onto $\mathcal N$. The integral identity of the chart of the fibre (equation~\YMMref{(986)}{442}, with $e^{\mathsf z_c(V)}=L^{-3}z^{(j)}_V(c)$ because $\mathsf n_c=L\mathsf m_c$ acts on a three-dimensional pivot space) gives, for every Borel function $H\ge0$ on the configurations of the level $j$,
\[
\begin{multlined}
\int_{\mathcal N}H\,dV=c_j\int_{\mathsf O}dW'\,z^{(j)}_{V_j}\int_{|\mathcal A|<\rho}e^{\mathsf J_{V_j}[\mathcal A]}\,H\bigl(V_{W'}[\mathcal A]\bigr)\,d\mathcal A,\qquad V_j=V_j(W'),\quad c_j \\
=(2\pi^2)^{-\#\mathfrak X^+_j}L^{-3\#\mathfrak B^+_{j+1}},
\end{multlined}
\]
and the exponent has the closed form
\[
\mathsf J_{V_j}[\mathcal A]+\log z^{(j)}_{V_j}=\sum_{\mathsf e\in\mathfrak X^+_j}\log\sigma_1\bigl(|\mathcal A(\mathsf e)|^2\bigr)+\sum_c\mathsf z_c\bigl(V_{W'}[\mathcal A]\bigr)+3\,\#\mathfrak B^+_{j+1}\log L,\qquad\sigma_1(s^2)=\Bigl(\frac{\sin s}s\Bigr)^2,
\]
where $\mathsf z_c=-\log\det\mathsf n_c$ is a real analytic function of the configuration on $\{\varpi_c<\varpi_1\}$.

(3)

\emph{The step keeps the good sets, with constants alone.} Let $W'\in\mathsf{Gd}_{j+1}$, $|\mathcal A|<\rho_2$ and $W''=V_{W'}[\mathcal A]$. For $j\ge1$: $V_j(W')=\mathcal M^jU_+$ with an element $U_+$ of the minimal orbit over $W'$, which lies in the minimal orbit over $V_j(W')$ (Lemma~\ref{lem:sf-class-window}(2)), so that $\mu_j(V_j(W'))=L^{-2}\mu_{j+1}(W')$ (the level scaling in Definition~\ref{def:sf-chart-data}); and $W''=e^{iB'}V_j(W')$ with $\sup|B'|\le2C_\chi|\mathcal A|<\frac12r_{\rm B}$ moves the representative of the background so that its regularity number moves by at most $2K_{\rm B}\sup|B'|$ (Lemma~\ref{lem:background-curvature-response}): the background over $W''$ exists and
\[
\mu_j(W'')\le L^{-2}\mu_{j+1}(W')+4K_{\rm B}C_\chi|\mathcal A|<L^{-2}\mu_\flat+(1-L^{-2})\mu_\flat=\mu_\flat .
\]
For $j=0$: every bond moves by $|e^{iB'(\mathsf e)}-\mathbf 1|\le|B'(\mathsf e)|$, and a plaquette variable by at most the sum over its four bonds, so that $\varpi(W'')\le\varpi(V_0(W'))+8C_\chi|\mathcal A|\le2L^{-2}\mu_1(W')+8C_\chi|\mathcal A|<\frac14\varepsilon_{\rm v}$. So the step maps $\mathsf{Gd}_{j+1}\times\{|\mathcal A|<\rho_2\}$ into $\mathsf{Gd}_j$.

(4)

\emph{Composition.} Put $\mathcal U_n=\mathsf O_n$ and, downwards, $\mathcal U_j=\Phi_j(\mathcal U_{j+1}\times\{|\mathcal A_j|<\rho_2\})$. By~(2) at $\mathsf O=\mathcal U_{j+1}$, which lies in $\mathsf{Gd}_{j+1}$ by~(3), every $\mathcal U_j$ is open and lies in $\mathsf{Gd}_j$, and $\Phi_j\times{\rm id}$ is a real analytic diffeomorphism. $\Upsilon$ is their composite: a real analytic diffeomorphism of $\mathsf O_n\times\Omega$ onto the open set $\mathcal U=\mathcal U_0$, injective on all of $\mathsf O_n\times\Omega$, with the inverse $U\mapsto(\mathcal M^nU,(\operatorname{Lg}(\mathcal M^jU(\mathsf e)\,V_j(\mathcal M^{j+1}U)(\mathsf e)^{-1}))_{\mathsf e,j})$, and $\Upsilon(\mathbf 1,0)=\mathbf 1$ by~(1). Along the chain the backgrounds exist (at level $0$ the background is $U_0(W_0)=W_0$, and $|J(W_0)|\le4\varpi(W_0)<\varepsilon_{\rm v}$, as in the height-zero part of the proof of Proposition~\YMMref{52.3}{433}), $\mathcal MW_j=W_{j+1}$ and $\theta_{\rm ch}$ is formed; this is~(i) and~(ii), and by the closed form of~(2)
\[
\theta_{\rm ch}=-\sum_{j<n}\Bigl[\sum_{\mathsf e\in\mathfrak X^+_j}\log\sigma_1\bigl(|\mathcal A_j(\mathsf e)|^2\bigr)+\sum_c\mathsf z_c(W_j)\Bigr]-3\log L\sum_{j<n}\#\mathfrak B^+_{j+1},
\]
a real analytic function on $\mathsf O_n\times\Omega$. The identities of the steps, applied downwards with Tonelli's theorem, give~(iii) for Borel $H\ge0$ that vanish off $\mathcal U$, with $C_{\rm H}=\prod_jc_j$, and for bounded complex Borel $H$ by linearity. As $\Upsilon$ is injective on all of $\mathsf O_n\times\Omega$, a function that vanishes off the image of an open subset $E$ gives an integrand that vanishes off $E$: this is~(iii) for $E$. The density is fixed by~(iii): if a continuous $\theta'$ satisfied~(iii) as well, the function $H$ equal to the sign of $e^{-\theta'}-e^{-\theta_{\rm ch}}$, carried to $\mathcal U$ by $\Upsilon$, would give $\int|e^{-\theta'}-e^{-\theta_{\rm ch}}|\,dW\,dy=0$, and a continuous function with this property vanishes.
\end{proof}

\begin{lemma}[The jets of the weighted progression]\label{C-351-10}
Let $0<g\le\gamma_{\rm e}$, $(K,\beta)$ and $b$ as in Lemma~\ref{C-351-8}, $\ell\in\mathfrak L(g)$, and $f,h$ as there, with $d<1$; $\Psi=(f,h)$. $\mathsf n_j(f,h)$ is the mixed second derivative at $0$ in $(\mu_1,\mu_2)$ of the source energy of the level $j$ of the weighted progression with the source $\mu_1f\circ x+\mu_2h\circ x$, summed over the localisation domains and points of the level $j$; $\mathsf J_2=\sum_{j\le K}\mathsf n_j$. When all fluctuation fields are multiplied by a complex field scale $\mathsf z$ and each $x_i$ is replaced by $L^{-i}\beta+\mathsf z^2(x_i-L^{-i}\beta)$, $\mathsf n_{j+1}$ becomes an even function $\mathsf n^{\mathsf z}_{j+1}=\sum_m\mathsf n_{j+1,[2m]}\mathsf z^{2m}$, holomorphic for $|\mathsf z|<R_j$ with $R_j^{-2}\le C(\log x_j)^{2p_0}D_j/b$, and $\mathsf n^1_{j+1}=\mathsf n_{j+1}$; put $\mathsf h^{(n)}_{2m}=\sum_{j<n}\mathsf n_{j+1,[2m]}$ and $\mathsf h_{2m}=\mathsf h^{(K)}_{2m}$. Then:
\begin{enumerate}
\item[(i)] $\mathsf n_{j,[0]}=\mathsf n_{j,[2]}=0$, and $|\mathsf J_2-\mathsf h_4-\mathsf h_6|\le Cb^{-4}d^2\mathsf l^{8p_0}|\Psi|$;
\item[(ii)] for $0\le n\le K-1$ with $D_n\ge d$: $|\mathsf h_6-\mathsf h_6^{(n)}|\le Cb^{-3}\mathsf l^{6p_0}d\,(d/D_n)^{3/2}|\Psi|$;
\item[(iii)] for $(f,h)\in\mathfrak T_{T_\ell}(d;\vartheta)$, $20d\le\ell$, $d\le L^{-1}$ and $K>s(d)$: $|\mathsf h_4-c_2b^{-2}\mathcal K_6(f,h)|\le[\bar\varepsilon_4(g,d)\,b^{-3}d+\frac32C_{\rm tm}\vartheta^{-1}b^{-2}a_K/d]|\Psi|$, with $\bar\varepsilon_4(g,d)\le C(g)d^{1/2}$ and the constant $C_{\rm tm}$ of Proposition~\YMMref{121.2}{824} and Lemma~\ref{lem:ward-fluctuation-template};
\item[(iv)] for $\ell=\ell_g$, $(f,h)\in\mathfrak T_{T_\ell}(d;\vartheta)$, $d\le d_{\rm e}(g,\vartheta)$ and $s(d)\le K$: there is a set $\mathsf{Pk}=\mathsf{Pk}(a_K,d)\subset a_K\mathbb Z^3$, depending on $a_K$ and $d$ alone, with $|v|<1$ on it and $\#\mathsf{Pk}\ge(30d)^{-3}$, such that
\[
\Bigl|S^T_{2;K,\ell,\beta}(f,h)-\frac1{\#\mathsf{Pk}}\sum_{v\in\mathsf{Pk}}\mathsf J_2(\tau_vf,\tau_vh)\Bigr|\le C_{\rm A}(g)\,\vartheta^{-1}d^2b^{-2}|\Psi|,
\]
$\tau_v$ the translation by $v$ and $C_{\rm A}(g)$ a number depending only on $g$ and the constants.
\end{enumerate}
\end{lemma}

\begin{proof}[Proof sketch]
(i), (ii). By Lemma~\ref{C-351-8}, if $f\circ x$ or $h\circ x$ vanishes identically, all the jets vanish; let neither vanish.

\emph{The steps at one field scale.} When the fluctuation fields of all the steps are multiplied by one complex field scale $\mathsf z$ and the coefficients of the action are moved as in the statement, the step of the level $j+1$ is a step at the place $x_j$ with the lower place $\mathsf t_\flat$ (Definition~\ref{IV-71-12}; Proposition~\ref{IV-71-50} for the canonical objects), the deformed activities of all the earlier levels standing in its bare array: the room $\mathsf q$ bounds their size uniformly in $\mathsf z$, so that one induction over the levels forms all the steps at once. So the activities of the level $j+1$ are pointed families, holomorphic and even in $\mathsf z$ on $|\mathsf z|<R_j$, $R_j=\delta(\mathsf t_\flat)/(2\delta(x_j))$ with $\delta(\mathsf t)=\mathsf t^{-1/2}\mathfrak b(\mathsf t)$, equal at $\mathsf z=1$ to the activities of the weighted progression, and local in the weights; and $R_j^{-2}=4\delta(x_j)^2/\delta(\mathsf t_\flat)^2\le C(\log x_j)^{2p_0}D_j/b$, as $x_j\ge\frac34b/D_j$ (Lemma~\ref{C-351-8}).

One step at the field scale $\mathsf z$, as a map of its bare array (the activities of the earlier levels and the weight of the level), has first and second derivatives at most $C\,q(|\mathsf z|/R_j)$ times the sizes of the increments, $q(t)=t^2/(1-t^2)$, by Cauchy's estimate, its order $0$ in $\mathsf z$ reading no bare array (Proposition~\ref{IV-71-37} with Proposition~\ref{IV-71-46}); the chain rule gives the derivatives in the two parameters $\mu_1,\mu_2$ of the activities of the level $j+1$ from those of the earlier levels and from the weight.

\emph{The derivatives in the sources.} The weight of the source $a$ at the level $i$ is $L^{K-i}w_a$, $w_1=f\circ x$, $w_2=h\circ x$, of size $\|f_a\|_\infty/D_i$ at a plaquette, $f_1=f$, $f_2=h$. The weight member of the bare array of the level $i$ is $L^i$ times its sum against the $a_p$ over the plaquettes of a domain, $a_p$ being at most of the order $L^{-4i}$ on the analyticity core; a domain of one cube holds at most of the order $L^{3i}$ plaquettes, and at most $35(d/a_K)^3$ plaquettes carry $w_a$; as $L^{-3i}(d/a_K)^3=(d/D_i)^3$, the size of this member is at most $C(\|f_a\|_\infty/D_i)\psi_i$, $\psi_i=\min\{1,(d/D_i)^3\}$. An activity of an earlier level $i-n$ that vanishes off the domains meeting a set of sites enters the bare array of the level $i$ with a size at most $C[\varpi_n+(n+2)L^{-4n}\min\{L^{3n},N_{\rm cb}\}]$ times its norm, $\varpi_n=L^{3n}e^{-\max\{2L^{n-1}-2,0\}}$ and $N_{\rm cb}$ the number of the cubes of the level $i-n$ that hold such a site; for the supports of the sources $N_{\rm cb}\le77\max\{1,(d/D_{i-n})^3\}$.

Along the run $\delta(x_{i-n})\le C(1+n)^{p_0}L^{-n/2}\delta(x_i)$, so $R_{i-n}^{-2}\le C(1+n)^{2p_0}L^{-n}R_i^{-2}$. By induction over $i\le j$ one obtains, for $|\mathsf z|\le\eta_0R_j$, $\eta_0\in(0,\frac12]$ a number depending only on the constants: the first derivative in $\mu_a$ of the activity of the level $i+1$ has a norm at most $C|\mathsf z|^2R_i^{-2}(\|f_a\|_\infty/D_i)\psi_i$, and the mixed second derivative a norm at most $C|\mathsf z|^2R_i^{-2}(\|f\|_\infty/D_i)(\|h\|_\infty/D_i)\chi_i$, $\chi_i=\min\{1,(d/D_i)^{5/2}\}$; here $\psi_i^2\le\chi_i$ bounds the product of two first derivatives, and the second derivatives of the earlier levels, carried up with the sizes above, sum to at most $C\chi_i$, the levels below the scale $d$ through the decay of the activities across the distance of the two supports, those above it through the factor $L^{-n/2}$ that is left when the factor $L^{-4n}$ of the size is set against the growth $L^{n}$ of the radii and weights and $L^{5n/2}$ of the profile.

\emph{The jet.} $\mathsf n^{\mathsf z}_{j+1}$ is the sum, over the domains of the level $j+1$ and their points, of the mixed second derivative at the trivial configuration. Only the domains that meet both supports contribute, the activities being local in the weights; the cubes of the level $j+1$ that hold a corner of a plaquette of a support number at most $77\max\{1,(d/D_{j+1})^3\}$; two corners of plaquettes of the two supports are at a sup distance $\ge d/4$, so that such a domain has a linear size $\ge[d/(4MD_{j+1})-2]_+$ in the units of the level $j+1$, and the decay of the localised terms in the size of their domains, summed by the bound in Lemma~\ref{lem:source-support-sum}, gives the factor $e^{-[d/(4MD_{j+1})-2]_+}$. Hence, on $|\mathsf z|\le\eta_0R_j$,
\[
\begin{multlined}
|\mathsf n^{\mathsf z}_{j+1}|\le C|\mathsf z|^2R_j^{-2}D_j^{-2}\chi_j\,\Xi_{j+1}|\Psi|,\qquad\chi_j=\min\{1,(d/D_j)^{5/2}\},\qquad\Xi_{j+1}=\max\{1, \\
(d/D_{j+1})^3\}\,e^{-[d/(4MD_{j+1})-2]_+};
\end{multlined}
\]
$\chi_j$ records that a source of width $d$ is seen at the spacing $D_j$ through its mean, and $\Xi_{j+1}$ counts the cubes of the level $j+1$ that meet the supports and the decay of the activities across the distance of the two supports; the bound is local, and it reads no $K$, side, position or pair.

\emph{The orders.} At $\mathsf z=0$ the steps add nothing that reads the sources, so the order $0$ in $\mathsf z$ vanishes. The order $\mathsf z^2$ comes from the step of the level $j+1$ alone, the earlier levels being taken at the field scale $0$, where they read no source; it is the covariance, in the measure of the step, of the parts linear in the fluctuation field of the two source potentials at the trivial configuration. These linear forms are invariant under the adjoint action of the constant gauge transformations, which leave the trivial configuration and the gauge invariant terms unchanged, and the mean of ${\rm Ad}(u)$ over $u\in SU(2)$ is $0$ on $\mathfrak{su}(2)$; so they vanish, and the order $2$ vanishes. Cauchy's estimate on $|\mathsf z|=\eta_0R_j$ gives
\[
|\mathsf n_{j+1,[2m]}|\le C\eta_0^{-2m}R_j^{-2m}\,D_j^{-2}\chi_j\,\Xi_{j+1}|\Psi|\le Cb^{-m}(\log x_j)^{2mp_0}D_j^{m-2}\chi_j\,\Xi_{j+1}|\Psi| ,
\]
and the same with $m=4$ for the sum of the orders $\ge8$ at $\mathsf z=1$, as $\eta_0^2R_j^2\ge2$ for $g\le\gamma_{\rm e}$ (geometric summation of the even coefficient bounds in Proposition~\ref{IV-71-37}(cx8)).

\emph{The sums over the levels.} For $D_j\ge d$: $\Xi_{j+1}=1$, $D_j^{m-2}\chi_j=d^{m-2}(d/D_j)^{9/2-m}$ is geometric in $j$ with the ratio $L^{-(9/2-m)}\le L^{-1/2}$ for $m\le4$, and $\log x_j\le\log(2b/D_j)\le\mathsf l$. For $D_j<d$: $D_j^{m-2}\chi_j\le d^{m-2}$, $\log x_j\le\mathsf l\,(1+(r+1)\log L)$ at the $r$-th such level counted from above, and $\Xi_{j+1}\le\sup_{u\in[L^{r-1},L^r]}\max\{1,u^3\}e^{-[u/4M-2]_+}$, which is summable against every power of $r$. Summing over all the levels gives (i), with $|\mathsf h_{2m}|\le Cb^{-m}d^{m-2}\mathsf l^{2mp_0}|\Psi|$ for $m=2,3$; summing over the levels $j\ge n$ alone, where $D_j\ge D_n\ge d$, gives $\sum_{j\ge n}|\mathsf n_{j+1,[6]}|\le Cb^{-3}\mathsf l^{6p_0}d\,(d/D_n)^{3/2}|\Psi|$, which is (ii).

(iii). If $f\circ x$ or $h\circ x$ vanishes identically, $\mathsf h_4=0$ and the template sum below vanishes, so that its bound alone gives the assertion; let neither vanish.

\emph{At a real field scale.} The coefficients $\xi_i(\mathsf z)=L^{-i}\beta+\mathsf z^2(x_i-L^{-i}\beta)$ of the action are chosen so that, the counterterms $e_{i'}$ of the progression being kept, $\mathsf z^{-2}\xi_i(\mathsf z)L^i+\sum_{1\le i'\le i}L^{i'}e_{i'}=\mathsf z^{-2}\beta$, by the run equation $L^ix_i+\sum_{1\le i'\le i}L^{i'}e_{i'}=\beta$; and the sum rules of the steps are stated for every field of the box. So at a real field scale $0<\mathsf z\le1$ the chain of the weighted progression is the chain of the chart integrals of the Wilson weight at $\beta/\mathsf z^2$, read at the fields $\mathsf z\zeta$, $\zeta=(\zeta_i)_i$ the fluctuation fields of the levels in their boxes $|\zeta_i|<\mathfrak b(x_i)$, and $\sum_{2\le j\le K}\mathsf n^{\mathsf z}_j$ is, up to a term of the first level, the covariance of the first source derivatives of the exponent of the level $1$, taken at the background that the chain builds at the level $1$ from these fields, in a complex measure of mass $1$ on the product of the boxes, the product over the levels $1\le i<K$ of normalised measures on the box of the level $i$, each of which depends on the fields of the higher levels.

As $\mathsf z\to0$ at a fixed lattice, the backgrounds tend to the trivial configuration uniformly on the boxes, all the potentials vanish, and the coefficient of the action of the level $i$ tends to $L^{-i}x_0$, so that this measure tends in total variation to the product $\mathbb Q$ of the measures $\nu^{\rm r}_i$ below; dividing by $\mathsf z^4$,
\[
\sum_{2\le j\le K}\mathsf n_{j,[4]}=\langle q_1;q_2\rangle_{\mathbb Q}+\mathsf e_1,\qquad q_a=\tfrac12a_K^{-1}\sum_pw_a(p)\,|\Theta_p|^2,
\]
$\Theta_p$ the linear curvature of the plaquette $p$ in the fields $\zeta_i$, $\mathbb Q=\bigotimes_{1\le i<K}\nu^{\rm r}_i$, $\nu^{\rm r}_i$ the centred Gaussian measure of the quadratic part of the step of the level $i$ at the coefficient $L^{-i}x_0$ of the action, restricted to the box of the level $i$ and normalised, and $\mathsf e_1$ a term of the first level. With $\mathsf n_{1,[4]}$ and the slice of the level $1$ below, $\mathsf e_1$ costs at most $C\,b^{-2}\mathsf l^{4p_0}e^{-c_\flat d/a_K}|\Psi|$, $c_\flat=\min\{1/(8ML),\delta_{\rm tp}/(2\sqrt3)\}$, $\delta_{\rm tp}$ the rate of Propositions~\YMMref{120.5}{821} and~\YMMref{120.6}{821}, as the supports are $d/a_K$ bare spacings apart.

\emph{The Gaussian values; the running couplings drop out.} Under the Gaussian measure of the level $i$ at the coefficient $L^{-i}x_0$ the curvature of the fluctuation field of the level $i$ has the covariance $x_0^{-1}\mathcal C_{(i+1)}$, $\mathcal C_{(i+1)}$ the slice of the level $i+1$ of the free lattice curvature (Proposition~\YMMref{120.3}{819}): the coupling $x_i$ enters the scaled field and the coefficient of the action in inverse proportions and drops out. The slices $\mathcal C_{(i+1)}$, $1\le i<K$, add up to $C^{\rm fl}_{K,0}-\mathcal C_{(1)}$, where $C^{\rm fl}_{K,0}=\sum_{1\le j\le K}\mathcal C_{(j)}$ is the fluctuation kernel of depth $0$ of Section~\YMMref{120}{817} (Propositions~\YMMref{120.5}{821} and~\YMMref{120.6}{821}), the kernel of the orthogonal projection of the plaquette $2$-cochains, in each colour, onto the curvatures of the fields whose $K$-fold linear block average vanishes, that is, onto $\mathsf{Fl}_K=\mathrm d(\ker\mathbf Q^s_K)$, $\mathbf Q^s_K$ the straight average of Section~\YMMref{120}{817}; it has no dimension and $|C^{\rm fl}_{K,0}|\le1$; the slice $\mathcal C_{(1)}$ of the level $1$ goes into $\mathsf e_1$.

The variables of different levels are independent and every $\nu^{\rm r}_i$ is even; by Wick's rule, with three colours, the Gaussian value of $\langle q_1;q_2\rangle$ is, up to that slice,
\[
\tfrac14a_K^{-2}\cdot2\cdot3\,x_0^{-2}\sum_{p,q}w_1(p)\,w_2(q)\,C^{\rm fl}_{K,0}(p,q)^2=\tfrac32\,b^{-2}\,\mathcal K[C^{\rm fl}_{K,0}](f,h),\qquad b=a_Kx_0,
\]
with the template sum $\mathcal K[\mathsf k](f,h)=\sum_{p,q}f(x(p))\,h(x(q))\,\mathsf k(p,q)^2$; no further power of $a_K$ enters.

\emph{The boxes.} Every joint cumulant of order $2$ or $4$ of the components of the field of the level $i$ under $\nu^{\rm r}_i$ differs from the Gaussian one by at most $C_{\rm bx}e^{-\frac12\gamma_2\mathfrak b(x_i)^2}$, uniformly in the volume and in the cutoff (Lemma~\ref{lem:ward-box-cumulants}, the components entering as insertions), with numbers $C_{\rm bx}$, $\gamma_2$ depending only on the constants; the odd cumulants vanish. The map from the fields of the level $i$ to the curvature has row sums at most $CL^{-2i}$, and at most $35^2(d/a_K)^6$ pairs of plaquettes have $w_1(p)w_2(q)\ne0$; with $L^{-3i}(d/a_K)^3=(d/D_i)^3$ the boxes change $\langle q_1;q_2\rangle$ by at most $Cb^{-2}[\mathsf E^{(3)}+(\mathsf E^{(3)})^2+\mathsf E^{(6)}]|\Psi|$, $\mathsf E^{(q)}=\sum_{1\le i<K}u_i^q\epsilon_i$, $u_i=d/D_i$, $\epsilon_i=e^{-\frac12\gamma_2\mathfrak b(x_i)^2}\le\phi_{\rm bx}(\frac34b/D_i)$, where $\phi_{\rm bx}(s)=e^{-c_{\rm bx}(\log s)^2}$, $c_{\rm bx}=\frac12\gamma_2A_1^2$, since $\mathfrak b(x_i)^2\ge A_1^2(\log x_i)^2$ and $x_i\ge\frac34b/D_i$ (Lemma~\ref{C-351-8}).

For $s,s'\ge1$, $\phi_{\rm bx}$ decreases, $\phi_{\rm bx}(ss')\le\phi_{\rm bx}(s)\phi_{\rm bx}(s')$ and $s^q\phi_{\rm bx}(s)\le e^{q^2/(4c_{\rm bx})}$. Cut the levels at $D_i\approx d^{1/2}$: the levels with $D_i<d$ give at most $C\phi_{\rm bx}(\frac34b)\phi_{\rm bx}(d^{-1/2})$, their $u_i$ being spaced by the factor $L$; those with $d\le D_i\le d^{1/2}$ have $u_i\le1$, $\sum u_i^3\le2$ and $\epsilon_i\le\phi_{\rm bx}(\frac34b)\phi_{\rm bx}(d^{-1/2})$; those with $D_i>d^{1/2}$ have $\sum u_i^3\le2d^{3/2}$ and $\epsilon_i\le\phi_{\rm bx}(\frac34b)$. As $\phi_{\rm bx}(d^{-1/2})\le Cd^{3/2}$ and $\phi_{\rm bx}(\frac34b)\le C/b$, $\mathsf E^{(3)},\mathsf E^{(6)}\le Cb^{-1}d^{3/2}$, and also $\le C$; so the boxes cost at most $Cb^{-3}d^{3/2}|\Psi|$, with $C$ depending only on the constants.

The bound is uniform in $K$ and in the side through the locality of the cumulant bounds; no count of the restrictions enters.

\emph{The template sum.} Lemma~\ref{lem:ward-fluctuation-template} at the depth $0$, where $a_0=1\ge d$:
\[
\Bigl|\mathcal K[C^{\rm fl}_{K,0}](f,h)-\frac3{8\pi^2}\mathcal K_6(f,h)\Bigr|\le C_{\rm tm}\Bigl[\frac{a_K}{\vartheta d}+\Bigl(\frac d\ell\Bigr)^3+d^3\Bigr]|\Psi|,
\]
and $\frac32\cdot\frac3{8\pi^2}=c_2$. Collecting,
\[
|\mathsf h_4-c_2b^{-2}\mathcal K_6(f,h)|\le b^{-2}\Bigl\{\tfrac32C_{\rm tm}\Bigl[\frac{a_K}{\vartheta d}+\Bigl(\frac d\ell\Bigr)^3+d^3\Bigr]+C\,\frac db\,d^{1/2}+C\,\mathsf l^{4p_0}e^{-c_\flat d/a_K}\Bigr\}|\Psi| .
\]
As $\ell\ge1$, $d/a_K\ge L(\log\frac1d)^2$ for $K>s(d)$, and $g^{-2}<b\le L(g^{-2}+\bar B)$, the members other than $\frac32C_{\rm tm}b^{-2}a_K/(\vartheta d)$ are at most $\bar\varepsilon_4(g,d)\,b^{-3}d\,|\Psi|$ with
\[
\bar\varepsilon_4(g,d)=\sup_{g^{-2}<b\le L(g^{-2}+\bar B)}\Bigl\{3C_{\rm tm}\,b\,d^2+C\,d^{1/2}+C\,\frac bd\,\mathsf l^{4p_0}e^{-c_\flat L(\log\frac1d)^2}\Bigr\}\le C(g)\,d^{1/2},
\]
$e^{-c(\log\frac1d)^2}$ being below every power of $d$.

(iv). Let $Z[w]$ be the generating function of $T_\ell$, the partition function with the source $w$, and $\mathcal E(w)$ its source energy, the sum over the levels and localisation domains of the source energies \eqref{eq:running-source-energies} of the sourced chain on $T_\ell$ \eqref{eq:running-source-recursion}. On the least torus the setting and the identity of~(4a)--(4c) in the proof of Proposition~\ref{C-351-7} hold, the conditions on $g$ of the statements read there entering $\gamma_{\rm e}$: the term without large-field regions in the sum over the level of the first creation of a large-field region (equation~\eqref{eq:ward-first-creation}, with the histories of Section~\YMMref{123}{828}) is $e^{\mathcal E(w)}\int e^{\mathsf R_K(w,\cdot)}\,d\mathsf m_K$ for the sources of level-wise norm $<1$, with $|\mathsf R_K(w,\cdot)|\le\|w\|_{\rm lw}\mathsf g_\emptyset$ for a source carried by one ball of radius $1$, and with a measure $\mathsf m_K$ of mass at most $Z[0]$ (steps~(4c), (4d) there).

The weighted progression (Definition~\ref{def:weighted-progression-data}) and the sourced chain are made level by level by the same renormalisation step, with the same couplings, the same auxiliary data and the same lattice, and the source enters both only through the weights $x_j-L^{K-j}w$ of the plaquettes. At the trivial configuration every localised term is taken at the trivial pair, and $a_p(\mathbf 1)=0$, so that there the relative exponent of the weighted progression after the level $K$ is the sum over the levels and localisation domains of its source energies; following the two constructions level by level, this sum is $\mathcal E(w)$ for all sources near $0$, with no condition on the depth.

So the mixed second derivative of $\mathcal E$ at $(\tau_vf,\tau_vh)$ is $\mathsf J_2(\tau_vf,\tau_vh)$, and $\log(Z[w]/Z[0])-\mathcal E(w)$ is holomorphic near $0$ with the mixed second derivative $\mathsf R=S^T_2-\mathsf J_2$. Take $\iota=a_K\lceil18d/a_K\rceil$ and $\mathsf{Pk}=\{v\in\iota\mathbb Z^3:|v_\mu|\le0.575-9d\}$; the balls of radius $9d$ about the $x_1+v$, $x_1$ a centre of a ball of radius $d$ that holds ${\rm supp}f$, are disjoint and lie in $\mathbb B(x_1,1)$. For signs $\epsilon_v$ the source $z_1\sum_v\epsilon_v\tau_vf+z_2\sum_v\epsilon_v\tau_vh$ has level-wise norm $<1$ for $|z_1|\|f\|_\infty,|z_2|\|h\|_\infty<A_\star$, with $A_\star\ge c(g)(\vartheta d)^{1/2}b$, and on this polydisc the squarefree coefficients of $Z[w]e^{-\mathcal E(w)}/Z[0]$ in $(z_1,z_2)$ are at most $Q_{\rm D}(g)\prod_a\|f_a\|_\infty/A_\star$, with the number $Q_{\rm D}(g)$ of~(4e) in the proof of Proposition~\ref{C-351-7} formed with one ball and two variables: split at the level of the first creation of a large field, bound the part in which no field is large, and apply Cauchy's estimate on a polydisc for each level, as in~(4e) there.

So $|\mathsf R(\sum_v\epsilon_v\tau_vf,\sum_v\epsilon_v\tau_vh)|\le2Q_{\rm D}^2A_\star^{-2}|\Psi|$; its mean over the signs is $\sum_v\mathsf R(\tau_vf,\tau_vh)$; $S^T_2$ is invariant under the translations of $a_K\mathbb Z^3$; divide by $\#\mathsf{Pk}\ge(30d)^{-3}$.
\end{proof}

For $N=2L^{\mathsf M}$, $1\le n<\mathsf M$ with $2L^{\mathsf M-n+1}\ge C_LM$, and plaquettes $p,q$ of $\mathbb Z^3_N$, let $\mathsf g_2[N,n](p,q)$ be the coefficient of $\beta^{-3}$ in the expansion by Laplace's method (Lemma~\textup{\ref{IV-73-26}}), as $\beta\to\infty$, of the covariance of $a_p$ and $a_q$ on the fibre $\{\mathcal M^nU=\mathbf 1\}$ of $\mathbb Z^3_N$ under $\exp[-\beta(A+\sum_{j<n}\mathsf k_jG_j\circ\mathcal M^j)]$, written in the chart $\Upsilon(\mathbf 1,\cdot)$ of Lemma~\ref{C-351-9}; here $A=\sum_pa_p$ is the Wilson action, $G_j$ the gauge-fixing function~\YMMref{(54)}{27} of the level $j$, and $\mathsf k\in(0,\infty)^n$ are strengths. This is the number $\mathsf g_2[N,n]$ of
Proposition~\ref{C-350-12}(iii): both definitions take the coefficient of
the covariance on the same fibre, and that proposition proves independence
from the adapted chart and from the gauge-fixing strengths. In particular,
it depends only on the bare torus, the number of blocking levels, $L$ and
the plaquettes, with no running coupling or cutoff remaining. $\mathsf B_1^{\rm lat}$ and $\mathsf P_1^{\rm lat}(f,h;a)=a\sum_{\tilde p,\tilde q}f(ax(\tilde p))h(ax(\tilde q))\mathsf B_1^{\rm lat}(\tilde p,\tilde q)$, sums over the plaquettes of $\mathbb Z^3$, are the coefficients of Definition~\textup{\ref{IV-73-22}}; Corollary~\textup{\ref{IV-73-73}} and Theorems~\textup{\ref{IV-73-98}}, \textup{\ref{IV-73-100}}~(a) give, for $(f,h)\in\mathfrak T_{\mathbb R^3}(d;\vartheta)$, $0<a\le d$ and $\lg(\cdot)=1+|\log(\cdot)|$,
\begin{equation}\label{C-351-11}
|\mathsf P_1^{\rm lat}(f,h;a)-\varkappa\,\mathcal K_5(f,h)|\le C_\varkappa\vartheta^{-1}a\lg(d/a)\,\|f\|_\infty\|h\|_\infty,\qquad\varkappa=\tfrac3{2\pi^3}=c_1c_2,
\end{equation}
with an absolute number $C_\varkappa$.

\begin{lemma}[Lattice perturbation theory]\label{C-351-12}
\begin{enumerate}
\item[(i)] In the setting of Lemma~\ref{C-351-10} let $N=\ell L^K$, $w_1=f\circ x$ and $w_2=h\circ x$ on $\mathbb Z^3_N$, $2\le n\le K$, 
\[
\mathsf H^{(n)}=\beta^3a_K^2\mathsf h_6^{(n)}
\]
 and 
\[
\mathsf T[N,n](w)=\sum_{p,q}w_1(p)w_2(q)\mathsf g_2[N,n](p,q).
\]
 If $\gamma^{\rm tr}\mathfrak b(x_{n-1})^2\ge80$ and 
\[
10240\,N^3e^{-\gamma^{\rm tr}\mathfrak b(x_{n-1})^2/8}\le1,
\]
 $\gamma^{\rm tr}>0$ a number depending only on the constants, then
\[
|\mathsf H^{(n)}-\mathsf T[N,n](w)|\le C(N+\beta)^{36}e^{-\gamma^{\rm tr}\mathfrak b(x_{n-1})^2/8}\,\|w_1\|_\infty\|w_2\|_\infty .
\]
\item[(ii)] Let $N=2L^{\mathsf M}$, $N'=L^sN$ with $s\ge1$, $1\le n<\mathsf M$ with $N/L^n\ge\max\{192,\,48(3\log L+\log2)/(\delta_{\rm tp}L(L-1)),\,C_LM/L\}$, $\delta_{\rm tp}>0$ the rate of Propositions~\YMMref{120.5}{821} and~\YMMref{120.6}{821}, and $\pi$ the reduction $\mathbb Z^3_{N'}\to\mathbb Z^3_N$. For plaquettes $p',q'$ of $\mathbb Z^3_{N'}$ whose base points have a distance $\le N/8$: $|\mathsf g_2[N,n](\pi p',\pi q')-\mathsf g_2[N',n](p',q')|\le CN'^{237}e^{-\delta_{\rm tp}N/(48L^{n-1})}$.
\item[(iii)] Let $\ell'=2L^{m'}\ge N_{\mathsf H}$, the side threshold fixed in Section~\YMMref{52.6}{426} and retained for the kernel bounds in Section~\YMMref{120.3}{820}, $\ell'L\ge C_LM$ and $K\ge1$. For plaquettes $p,q$ of $\mathbb Z^3_{\ell'L^K}$ and plaquettes $\tilde p,\tilde q$ of $\mathbb Z^3$ over them with $|x_{\tilde p}-x_{\tilde q}|_\infty\le\ell'L^K/8$: $|\mathsf g_2[\ell'L^K,K](p,q)-\mathsf B_1^{\rm lat}(\tilde p,\tilde q)|\le3C_\star\ell'^{19}KL^{-K/2}$, with a number $C_\star$ independent of the side.
\end{enumerate}
\end{lemma}

\begin{proof}[Proof sketch]
(i) In the chart of Lemma~\ref{C-351-9} at the trivial block field of the level $n$, an exact identity, made one level at a time along matched chains of fields, shows that at a real field scale $\mathsf z$ the sum of the jets of the levels $\le n$ is the covariance of two observables in the probability measure $\propto\exp[-\mathsf z^{-2}\mathcal Q(\mathsf zg\zeta)-\mathcal Q'(\mathsf zg\zeta)]d\zeta$ on the product of the boxes, $g\zeta=(x_i^{-1/2}\zeta_i)_i$. Here $\mathcal Q\ge0$ has a positive definite Hessian, block diagonal over the levels; $\mathcal Q$, $\mathcal Q'$ and the observables are real analytic near $0$ and contain no $\mathsf z$, no counterterm and no box; and, up to colour rotations at each level, which keep the boxes and the Lebesgue measure, they are $\beta$ times the action with the gauge fixing of the strengths $\mathsf k_j=x_j/\beta$, the function $\theta_{\rm ch}$, and $\mathcal O(w_a)\circ\Upsilon$.

The chart is read there on the product of the boxes $|\mathcal A_i|<x_i^{-1/2}\mathfrak b(x_i)$ of the fields $\mathcal A_i=x_i^{-1/2}\zeta_i$, and these boxes lie in $\Omega$ once $\delta(x_i)=x_i^{-1/2}\mathfrak b(x_i)<\rho_2$; as $\rho_2$ depends only on the constants (it is written out in the proof of Lemma~\ref{C-351-9}), and as $\delta$ decreases on $[e^{2p_0},\infty)$ while all the couplings of the run exceed the inverse square of the coupling $g$ of the datum (Lemma~\ref{C-351-8}), this is a condition on the coupling $g$ alone, and it enters $\gamma_{\rm e}$. So $\mathsf h^{(n)}_6$ is the second Laplace coefficient computed with the Gaussian measures restricted to the boxes.

Without the boxes, $\beta^3a_K^2$ times this coefficient is $\mathsf T[N,n](w)$, by the intrinsic fibre-coefficient identity of Proposition~\ref{C-350-12}(iii) and because the second Laplace coefficient of $c\,\mathcal Q$ scales like $c^{-3}$. For these matched chains, Proposition~\ref{C-350-11}(ii) bounds the change on restricting the Gaussian cumulants to the boxes: the total polynomial order is at most $10$, there are at most $N^3$ places, and its Gaussian precision and normalisation bounds give the stated $(N+\beta)^{36}$ times exponential error.
(ii) Apply Proposition~\ref{C-350-14} with $t=n$ and
$m'=2$. Its lower bound on $N/L^n$ is the one in this statement, and
$\delta_{\rm tp}L(N/L^n)/48=\delta_{\rm tp}N/(48L^{n-1})$.
The coefficients are the same intrinsic fibre coefficients by
Proposition~\ref{C-350-12}(iii), so its estimate gives the assertion.

(iii) This is the coarser bound already stated in
Proposition~\ref{C-350-15}, at the side $\ell'$. The side threshold, room
condition and separation of the lifted plaquettes are exactly those of
that proposition. Its detailed proof keeps both the torus-to-lattice
kernel comparison and the correction from the flat top.
\end{proof}

\begin{lemma}[Removal of the cutoff at a fixed label]\label{C-351-13}
Let $0<g\le\gamma_{\rm e}$, $\ell\in\mathfrak L(g)$ and $\vartheta\in(0,1)$. Let $\mathsf J_g=[\gamma^{-2},L(X_g+\bar B)]$, let $\mathsf C'=\mathsf C'_2(\mathsf J_g)$ be the number of the definition in Corollary~\YMMref{117.1}{812} at $n=2$ on the side $\ell$, $\mathsf p_2=6p_0+\frac12$ the exponent~the exponent defined before Theorem~\YMMref{AR}{810} at $n=2$, $c_\diamond=\sup_{k\ge1}k^{\mathsf p_2}L^{-k/20}<\infty$, and
\[
\epsilon_\diamond(d)=(X_g+\bar B)^3\,\mathsf C'\,\ell^9\,c_\diamond\Bigl(\frac{40M}\vartheta\Bigr)^2d^{87}.
\]
There are $d_{\rm K}\in(0,\frac1{200}]$ and integers $K_\sharp(d)>s(d)$ with $d^{-1800}<L^{K_\sharp(d)}\le d^{-1801}$ for $0<d\le d_{\rm K}$, such that the following holds at every label $b\in I(g)$, at which every pair $(K,L^Kb)$, $K\in\mathbb N$, is a first-passage datum at $g$, for $0<d\le d_{\rm K}$ and $(f,h)\in\mathfrak T_{T_\ell}(d;\vartheta)$.
\begin{enumerate}
\item[(a)] For every $K\ge K_\sharp(d)$
\[
|S^T_{2;K,\ell,L^Kb}(f,h)-S^T_{2;K_\sharp(d),\ell,L^{K_\sharp(d)}b}(f,h)|\le\epsilon_\diamond(d)\,b^{-3}d\,\|f\|_\infty\|h\|_\infty ,
\]
and the same with $S_2^{(\ell,b)}(f,h)$ in the place of the first term.
\item[(b)] If, for numbers $\epsilon_{\rm h},C_{\rm h}\ge0$, the cutoff $K_\sharp(d)$ has $|S^T_{2;K_\sharp(d),\ell,L^{K_\sharp(d)}b}-L_b^+|(f,h)\le[\epsilon_{\rm h}\,b^{-3}d+C_{\rm h}\,b^{-2}a_{K_\sharp(d)}/d]\,\|f\|_\infty\|h\|_\infty$, then for every $K\ge K_\sharp(d)$
\[
|S^T_{2;K,\ell,L^Kb}-L_b^+|(f,h)\le\bigl[\epsilon_{\rm h}+\epsilon_\diamond(d)+C_{\rm h}(X_g+\bar B)\,d^{1798}\bigr]\,b^{-3}d\,\|f\|_\infty\|h\|_\infty ,
\]
with no term in $a_K$, and the same with $S_2^{(\ell,b)}$ in the place of $S^T_{2;K,\ell,L^Kb}$.
\end{enumerate}
\end{lemma}

\begin{proof}[Proof sketch]
\emph{The rate.} the constants and thresholds in Theorem~\YMMref{AR}{810} and Theorem~\YMMref{AR}{810}, equation~\YMMref{(2203)}{810} are made in the setting (Th1), (Th2), whose datum is the choice of the constants and whose numbers are formed from it. The interval $\mathsf J_g$ contains $I(g)$, and it is an interval of labels under (Tor) on the side $\ell$ (Lemma~\ref{lem:cutoff-admissible-labels}). Let $\mathsf C=\mathsf C_2(\mathsf J_g)$, $K_0=K_0(\ell,2,\mathsf J_g)$ and $K_1(\cdot)=K_1(\ell,2,\mathsf J_g,\cdot)$ be the other numbers of the constants and thresholds in Theorem~\YMMref{AR}{810} at $n=2$; so $\mathsf C'=\mathsf C\sum_{j\ge0}(1+j)^{\mathsf p_2}L^{-j/10}$. For real Lipschitz functions $f,h$ on $T_\ell$, neither of them $0$, whose supports have the spread $\mathsf d_\infty(f,h)\ge d'>0$ in the sup distance, the uniform cutoff estimate of Theorem~\YMMref{AR}{810} gives, at every label $b\in\mathsf J_g$ and every integer $K\ge K_1(d')$, the inequality~\YMMref{(2203)}{810} at $n=2$,
\[
|S^T_{2;K+1,\ell,L^{K+1}b}(f,h)-S^T_{2;K,\ell,L^Kb}(f,h)|\le\mathsf C\,\ell^9K^{\mathsf p_2}L^{-K/10}\,\|f\|_\sharp\|h\|_\sharp ,
\]
and its sum over the cutoffs, with $(K+j)^{\mathsf p_2}\le K^{\mathsf p_2}(1+j)^{\mathsf p_2}$, is the inequality~\YMMref{(2207)}{812} of Corollary~\YMMref{117.1}{812}:
\[
|S^T_{2;K',\ell,L^{K'}b}(f,h)-S^T_{2;K,\ell,L^Kb}(f,h)|\le\mathsf C'\,\ell^9K^{\mathsf p_2}L^{-K/10}\,\|f\|_\sharp\|h\|_\sharp\qquad(K'\ge K\ge K_1(d')).
\]
The supports of a pair $(f,h)\in\mathfrak T_{T_\ell}(d;\vartheta)$ are at the Euclidean distance $\ge d$, so at the sup distance $\ge d/\sqrt3$, both distances of $T_\ell$ being minima over the lifts and $|v|_\infty\ge|v|/\sqrt3$ on $\mathbb R^3$; so $\mathsf d_\infty(f,h)\ge d/\sqrt3$ (the definition preceding Lemma~\YMMref{114.6}{805}), and $d'=d/\sqrt3$ serves; put $K_\diamond(d)=K_1(d/\sqrt3)$. Moreover $\|f\|_\sharp\le40M\|f\|_\infty/(\vartheta d)$, as ${\rm Lip}\,f\le\|f\|_\infty/(\vartheta d)$ and $40M\ge\vartheta d$, and the same for $h$.

\emph{The threshold.} By~the two threshold inequalities at the end of the proof of Theorem~\YMMref{AR}{810} at $n=2$, $K_1(d')$ is the least integer $k\ge K_0$ such that every integer $K\ge k$ has
\[
d'L^u\ge6M\qquad\text{and}\qquad A_\diamond\,N_{[u]}\,e^{-d'L^u/M}\le\tfrac12\mathsf C\,\ell^9K^{\mathsf p_2}L^{-K/10},
\]
with the depth $u=\lceil K/1800\rceil$ of~the depth choice preceding Theorem~\YMMref{AR}{810} at $n=2$, $N_{[u]}=\ell^3L^{3u}$, and $A_\diamond=8e^3E_0C_{\rm E,0}\varrho^{-2}$ the factor in front of $N_{[u]}$ in~the two threshold inequalities at the end of the proof of Theorem~\YMMref{AR}{810}, $\varrho$ being the radius of the class of sources~in Section~\ref{sec:cutoff-calibration} that stands in the formulas of the constants and thresholds in Theorem~\YMMref{AR}{810}, not the radius of the choice. The first condition at $K=K_\diamond(d)$, with $u<K/1800+1$, gives $L^{K_\diamond(d)}>(6\sqrt3M/(Ld))^{1800}$; as $6\sqrt3M>L$ (the cube size is $M=L^{m'}$ with $m'\ge5$), $a_{K_\diamond(d)}<d^{1800}$. Conversely let $\Lambda_\diamond(d)=4000+366\log(\sqrt3M/d)$, and let $K$ be an integer whose number $v=L^u$ has $v\ge(\sqrt3M/d)\Lambda_\diamond(d)$.

The first condition holds at $d'=d/\sqrt3$, as $\Lambda_\diamond(d)\ge6$. As $K\le1800u$, $L^{-K/10}\ge v^{-180}$; with $K^{\mathsf p_2}\ge1$ and $\ell\ge1$ the second condition follows from $A_\diamond v^{183}e^{-d'v/M}\le\frac12\mathsf C$, that is, from $\varphi_\diamond(v)=d'v/M-183\log v\ge\log(2A_\diamond/\mathsf C)$, and $2A_\diamond<\mathsf C$: at $n=2$, as $\log(L\cdot L(X_g+\bar B))\ge1$ and $\bar{\mathsf C}_{\mathsf M}\ge\varrho^{-1}$ in~the majorant definitions in the proof of Theorem~\YMMref{AR}{810}, the number $\mathsf C$ of~\YMMref{(2206)}{811} is at least $4\,(16\varrho^{-1})^2\,2^3(640/\pi)^2(1+8E_0C_{\rm E,0})>2A_\diamond$. So $\varphi_\diamond(v)\ge0$ suffices. The function $\varphi_\diamond$ increases for $v\ge183M/d'$, and at $v=(M/d')\Lambda_\diamond(d)=(\sqrt3M/d)\Lambda_\diamond(d)$, since $183\log\lambda<\frac12\lambda$ for $\lambda\ge4000$, $\varphi_\diamond>\frac12\Lambda_\diamond(d)-183\log(\sqrt3M/d)=2000$; so $\varphi_\diamond(v)>0$ for all these $v$.

Every $K\ge1800\log_L(\sqrt3M\Lambda_\diamond(d)/d)$ has such a $v$, and $L^{\lceil x\rceil}<L\cdot L^x$. So
\[
\Bigl(\frac{6\sqrt3M}{Ld}\Bigr)^{1800}<L^{K_\diamond(d)}\le\max\Bigl\{L^{K_0},\ L\Bigl(\frac{\sqrt3M\,\Lambda_\diamond(d)}d\Bigr)^{1800}\Bigr\}.
\]

\emph{(a).} For $K\ge K_\diamond(d)$: $K^{\mathsf p_2}L^{-K/10}\le c_\diamond L^{-K/20}\le c_\diamond\,a_{K_\diamond(d)}^{1/20}<c_\diamond d^{90}$. So for $b\in I(g)\subset\mathsf J_g$, $K'\ge K\ge K_\diamond(d)$ and $(f,h)\in\mathfrak T_{T_\ell}(d;\vartheta)$
\[
|S^T_{2;K',\ell,L^{K'}b}(f,h)-S^T_{2;K,\ell,L^Kb}(f,h)|\le\mathsf C'\ell^9c_\diamond\Bigl(\frac{40M}\vartheta\Bigr)^2d^{88}\,\|f\|_\infty\|h\|_\infty\le\epsilon_\diamond(d)\,b^{-3}d\,\|f\|_\infty\|h\|_\infty ,
\]
as $d^{88}=b^{-3}d\cdot b^3d^{87}$ and $b\le X_g+\bar B$ on $I(g)$. Put $K_\sharp(d)=\max\{K_\diamond(d),s(d)+1\}$; then $L^{K_\sharp(d)}\ge L^{K_\diamond(d)}>d^{-1800}$. From $L^{s(d)}d<LC_W(s(d))=L\mathsf c_W(\log L^{s(d)}+\log(X_g+\bar B))^{2p_0}$ one gets $L^{s(d)}\le d^{-2}$, so $L^{s(d)+1}\le d^{-3}$, for small $d$; and, $K_0$ being fixed with $g$ and $\ell$, $L^{K_0}\le L(\sqrt3M\Lambda_\diamond(d)/d)^{1800}\le d^{-1801}$ for small $d$. Let $d_{\rm K}\in(0,\frac1{200}]$ be a number below which these hold; then $L^{K_\sharp(d)}\le d^{-1801}$ for $0<d\le d_{\rm K}$. The last display, with $K\ge K_\sharp(d)$ in the place of $K'$ and $K_\sharp(d)$ in the place of $K$, is~(a) at the finite cutoffs. At the labels of $I(g)$ the limit $S_2^{(\ell,b)}(f,h)$ exists along the full sequence of the cutoffs (Corollary~\YMMref{118.4}{815}), and a bound that holds at all large $K$ with one right member holds for the limit.

\emph{(b).} For $b\in I(g)$ the pairs $(K,L^Kb)$, $K\in\mathbb N$, are the first-passage data at $g$ with the label $b$ (Corollary~\YMMref{118.4}{815}). The main term $L_b^+(f,h)$ reads the label and the pair, not the cutoff. By the triangle inequality, (a) and the inequality assumed at $K_\sharp(d)$, the left member at $K\ge K_\sharp(d)$, and at the limit, is at most $[\epsilon_{\rm h}b^{-3}d+\epsilon_\diamond(d)b^{-3}d+C_{\rm h}b^{-2}a_{K_\sharp(d)}/d]\|f\|_\infty\|h\|_\infty$, and $b^{-2}a_{K_\sharp(d)}/d=b^{-3}d\cdot b\,a_{K_\sharp(d)}d^{-2}<b^{-3}d\cdot(X_g+\bar B)\,d^{1798}$, as $a_{K_\sharp(d)}\le a_{K_\diamond(d)}<d^{1800}$.
\end{proof}

\begin{proposition}[The law on the least torus]\label{C-351-14}
Let $0<g\le\gamma_{\rm e}$, $\vartheta\in(0,1)$, $\ell=\ell_g$ and $C_{\rm cut}=\frac32C_{\rm tm}\vartheta^{-1}$. There is a non-decreasing function $\bar\varepsilon=\bar\varepsilon_{g,\vartheta}$ on $(0,d_{\rm e}(g,\vartheta)]$, $\bar\varepsilon(d)\to0$, depending on $g$, $\vartheta$ and the constants and on no $K$, $\beta$, label, pair or position, such that for $0<d\le d_{\rm e}(g,\vartheta)$ and $\Psi=(f,h)\in\mathfrak T_{T_\ell}(d;\vartheta)$:
\begin{enumerate}
\item[(a)] at every first-passage datum $(K,\beta)$ at $g$ with $s(d)<K$ and $L^K\le d^{-1801}$, whatever its label $b=a_K\beta$,
\[
|S^T_{2;K,\ell,\beta}-L_b^+|(f,h)\le\bigl[\bar\varepsilon(d)\,b^{-3}d+C_{\rm cut}\,b^{-2}a_K/d\bigr]|\Psi| ;
\]
\item[(b)] for $b\in I(g)$ the same holds at $(K,L^Kb)$ for every $K>s(d)$, and $|S_2^{(\ell,b)}-L_b^+|(f,h)\le\bar\varepsilon_\infty(d)\,b^{-3}d\,|\Psi|$ with a non-decreasing $\bar\varepsilon_\infty=\bar\varepsilon_{\infty;g,\vartheta}\le\bar\varepsilon$ and $\bar\varepsilon_\infty(d)\le C(g,\vartheta)d^{1/2}$;
\item[(c)] with $C'_{\rm cut}=\frac92C_{\rm tm}\vartheta^{-1}+1$ in the place of $C_{\rm cut}$, the tolerance has the rate $d^{1/2}$ at every cutoff: there is $C_{\rm r}(g,\vartheta)<\infty$ with $|S^T_{2;K,\ell,L^Kb}-L_b^+|(f,h)\le[C_{\rm r}\,b^{-3}d^{3/2}+C'_{\rm cut}\,b^{-2}a_K/d]|\Psi|$ for all $b\in I(g)$ and $K>s(d)$, and $|S_2^{(\ell,b)}-L_b^+|(f,h)\le C_{\rm r}\,b^{-3}d^{3/2}|\Psi|$.
\end{enumerate}
\end{proposition}

\begin{proof}[Proof sketch]
\emph{Before the chain.} If $f\circ x$ or $h\circ x$ vanishes identically, then $S^T_{2;K,\ell,\beta}(f,h)$, all the jets and every lattice sum of the sources vanish (Lemma~\ref{C-351-8}), and Lemma~\ref{C-351-10}~(iii) holds by the template sum alone; in the chain below every arrow but the last then reads $0=0$, and the last is~\eqref{C-351-11} as it stands, with $\mathsf P_1^{\rm lat}=0$. Let $\mathsf r=d/a_K$, so that $\mathsf r\ge L(\log\frac1d)^2$. Translations keep the class and the $\mathcal K_q$, so Lemma~\ref{C-351-10}~(iv), (i), (iii) at the pairs $(\tau_vf,\tau_vh)$ give
\[
\begin{multlined}
|S^T_{2;K,\ell,\beta}-L_b^+|(f,h)\le\Bigl[C(g,\vartheta)\bigl(d^{1/2}+d(\log\tfrac1d)^{8p_0}\bigr)b^{-3}d+C_{\rm cut}\,b^{-2}\frac{a_K}d\Bigr]|\Psi|+\max_{v\in\mathsf{Pk}}|\mathsf h_6 \\
-\varkappa b^{-3}\mathcal K_5|(\tau_vf,\tau_vh),
\end{multlined}
\]
the packing mean costing $C_{\rm A}(g)\vartheta^{-1}bd$ and the orders $\ge8$ costing $Cb^{-1}d\,\mathsf l^{8p_0}$ in the unit $b^{-3}d|\Psi|$. As $\mathsf h_6=b^{-3}a_K\mathsf H^{(K)}$ and $b^{-3}a_K\mathsf r=b^{-3}d$, it remains to show $|\mathsf H^{(K)}-\varkappa\mathcal K_5/a_K|\le\eta(d)\,\mathsf r\,|\Psi|$ with $\eta(d)\to0$ uniformly; all the costs below are in the unit $\mathsf r|\Psi|$.

\emph{The levels.} Let $K'=\max\{K,\lceil24\log_L\mathsf r\rceil\}$, so that $L^{K'}\ge\mathsf r^{24}$ and $K'\le C\log\frac1d$; let $n_1=K'-\lceil K'/78\rceil$, $N_0=\ell L^K$ and $N'=\ell L^{K'}$. If $K'>K$, let $n_0$ be the level with $D_{n_0}\in[d^{1/2},Ld^{1/2})$, $b_*\in I(g)$ fixed, $D'=(K',L^{K'}b_*)$, a first-passage datum at $g$, and $\breve w_a$ the sources lifted to one copy in $\mathbb Z^3_{N'}$ of a ball of radius $N_0/4$ that holds the supports; by Lemma~\ref{C-351-8} the lifted sources are a pair at the datum $D'$, with no condition on its depth. The chain is
\[
\mathsf H^{(K)}\to\mathsf H^{(n_0)}\to\mathsf T[N_0,n_0]\to\mathsf T[N',n_0]\to\mathsf H^{(n_0)}[D']\to\mathsf H^{(n_1)}[D']\to\mathsf T[N',n_1]\to\mathsf T^{\rm lat}\to\varkappa\mathcal K_5/a_K ,
\]
with $\mathsf T^{\rm lat}=\sum\tilde w_1(\tilde p)\tilde w_2(\tilde q)\mathsf B_1^{\rm lat}(\tilde p,\tilde q)$ for the sources of the lifts $\tilde f,\tilde h$ to $\mathbb R^3$. By Lemma~\ref{C-351-10}~(ii) at $n=n_0$, the first arrow costs at most $C\mathsf l^{6p_0}(d/D_{n_0})^{3/2}\le C\mathsf l^{6p_0}d^{3/4}$. By Lemma~\ref{C-351-12}~(i), with $x_{n_0-1}\ge\frac34b/D_{n_0-1}\ge\frac34d^{-1/2}$ (Lemma~\ref{C-351-8}) and $N_0+\beta\le(\ell+X_g+\bar B)d^{-1801}$, as $b\le L(g^{-2}+\bar B)=X_g+\bar B$, the second costs at most $e^{-c(\log\frac1d)^{2p_0}}$ times a power of $1/d$. By Lemma~\ref{C-351-12}~(ii), applied pair by pair to the at most $900\mathsf r^6$ pairs of plaquettes of the supports, with $N_0/L^{n_0}=\ell/D_{n_0}\ge\ell d^{-1/2}/L$, the third costs at most $900\mathsf r^5CN'^{237}e^{-\delta_{\rm tp}\ell d^{-1/2}/48}$.

By Lemma~\ref{C-351-12}~(i) at $(D',n_0)$, the number $\mathsf T[N',n_0]$ being free of $\beta$, and with $x'_{n_0-1}\ge\frac34\mathsf r^{11}$, the fourth costs less than every power of $d$. By Lemma~\ref{C-351-10}~(ii) at $D'$ the fifth costs at most $C\mathsf l'^{6p_0}d^{3/4}$, $\mathsf l'\le\log(2X_g+2\bar B)+K'\log L$. By Lemma~\ref{C-351-12}~(i) at $(D',n_1)$, with $x'_{n_1-1}\ge L^{K'/78}$, the sixth costs less than every power of $d$. By Lemma~\ref{C-351-12}~(iii) on the torus of the side $N'/L^{n_1}=\ell L^{\lceil K'/78\rceil}$ with $n_1$ levels, over at most $900\mathsf r^6$ pairs, the seventh costs at most $C\ell^{19}K'\mathsf r^5L^{-K'/4}\le CK'/\mathsf r$, since $\mathsf r^{24}\le L^{K'}$.

By~\eqref{C-351-11} at $a=a_K$, as $a_K\mathsf T^{\rm lat}=\mathsf P_1^{\rm lat}(\tilde f,\tilde h;a_K)$, the eighth costs at most $C_\varkappa\vartheta^{-1}\lg(\mathsf r)/\mathsf r$. If $K'=K$, that is $\mathsf r^{24}\le L^K$, there is no second datum and no second torus: the chain is $\mathsf H^{(K)}\to\mathsf H^{(n_1)}\to\mathsf T[N_0,n_1]\to\mathsf T^{\rm lat}\to\varkappa\mathcal K_5/a_K$, its first arrow by Lemma~\ref{C-351-10}~(ii) at $n=n_1$, with the cost $C\mathsf l^{6p_0}(\mathsf r/L^{n_1})^{3/2}$, $\mathsf r/L^{n_1}\le L^{1-0.94K}$ and $\mathsf l\le\log(2X_g+2\bar B)+K\log L$, and the other three as the last three above with $D'=D$ and $N'=N_0$.

\emph{The conditions and the limits.} At the three places where Lemma~\ref{C-351-12}~(i) is used its two conditions hold for $d\le d_{\rm e}$, because $\frac18\gamma^{\rm tr}\mathfrak b(X)^2=\frac18\gamma^{\rm tr}A_1^2(\log X)^{2p_0}$ with $2p_0\ge2$ outgrows $10+\log(10240N^3)$, the logarithm of $N$ being at most a constant plus a fixed multiple of $\log X$ at the lower bounds $X$ of the three places. Every cost tends to $0$ with $d$, uniformly ($\mathsf r\to\infty$, $K'\le C\log\frac1d$), and depends only on $g$, $\vartheta$ and the constants; $\bar\varepsilon(d)$ is the supremum over $d''\le d$ of their sum. This is~(a).

\emph{Removal of the cutoff.} For $b\in I(g)$ the pairs $(K,L^Kb)$ are first-passage data at $g$, and for $s(d)<K\le K_\sharp(d)$, (a) applies, as $L^K\le d^{-1801}$. At $K_\sharp(d)$ the inequality of (a) is the one that Lemma~\ref{C-351-13}~(b) asks there, with $\epsilon_{\rm h}=\bar\varepsilon(d)$ and $C_{\rm h}=C_{\rm cut}$; so for $K\ge K_\sharp(d)$ and for the limit the inequality holds with the tolerance $\bar\varepsilon(d)+\epsilon_\diamond(d)+C_{\rm cut}(X_g+\bar B)d^{1798}$ and without the second member, and $\bar\varepsilon+\sup_{d''\le d}[\epsilon_\diamond(d'')+C_{\rm cut}(X_g+\bar B)d''^{1798}]$ serves in (a) and (b). At $K_\sharp$ one has $\mathsf r\ge d^{-3}$, so the second, third, fourth and sixth costs are below every power of $d$, $K'/\mathsf r$ and $\lg(\mathsf r)/\mathsf r$ are at most $Cd^2$, the packing mean and the orders $\ge8$ give $Cd(\log\frac1d)^{8p_0}$, the order six $C\mathsf l^{6p_0}d^{3/4}$, and the slowest member is the order four, $\bar\varepsilon_4(g,d)\le C(g)d^{1/2}$: this is the rate of $\bar\varepsilon_\infty$.

\emph{The rate at every cutoff.} Run the same chain with the cut at $D_{n_0}\in[d^{1/5},Ld^{1/5})$; the arithmetic of the levels goes through as before for $d\le L^{-13}$. The first and fifth arrows now cost at most $C\mathsf l^{6p_0}d^{6/5}$ and $C\mathsf l'^{6p_0}d^{6/5}$, the third $e^{-c\ell d^{-1/5}}$ times a power of $1/d$, and the second, fourth and sixth stay below every power of $d$. The seventh and eighth members, which near the smallest $\mathsf r$ decay only like a power of $\log\frac1d$, contribute at most $C(K'+\lg\mathsf r)\mathsf r^{-1}\,b^{-3}d|\Psi|=C(K'+\lg\mathsf r)(d/b)\,b^{-2}(a_K/d)|\Psi|\le\eta(d)\,b^{-2}(a_K/d)|\Psi|$, with $\eta(d)=Cg^2d\log\frac1d\le1$ for $d\le d_{\rm e}$; they go into the second member, whose constant becomes at most $C_{\rm cut}+1\le C'_{\rm cut}$.

All the remaining members are $O(d\,(\log\frac1d)^{8p_0})$ but the order four, which is at most $C(g)d^{1/2}$. Beyond $K_\sharp(d)$, Lemma~\ref{C-351-13}~(b), with $C_{\rm h}=C'_{\rm cut}$, removes the second member at the price $\epsilon_\diamond(d)+C'_{\rm cut}(X_g+\bar B)d^{1798}$, at most a constant times $d^{1/2}$. This is~(c).
\end{proof}

The coefficient $\frac8{3\pi}$ comes from $\varkappa=c_1c_2$ in~\eqref{C-351-11} alone; in particular it does not depend on $L$, $g$ or the label.

\subsection{Transfer to continuum components}

\begin{lemma}[From a limit at a fixed label to the components]\label{C-351-15}
Let $\mathfrak c$ be a choice, $0<g\le\gamma_W$, $S$ a volume-first limit at $g$, $b=\mathsf b(S)$, $\mathsf S$ the function $S_2$ or a function $S_2^\xi$ of a canonical decomposition of $\sigma_S$, $\vartheta\in(0,1)$, $z\in\mathbb Z$, and $\tau'=(g',b',\ell')$ an admissible triple with $b'=L^zb$. For $L^zd\le d_\flat(g',\vartheta)$ and $\Psi\in\mathfrak T^1(d;\vartheta)$,
\[
|b^2\mathsf S(\Psi)-b'^2T^{(\ell',b')}(\Psi_{L^z})|\le E_{\tau'}\,(d/b)^2|\Psi| .
\]
\end{lemma}

\begin{proof}[Proof sketch]
Let $(K_j,\beta_j,\nu_j)$ define $S$; then $a_{K_j}\beta_j\to b$. $\mathsf S$ is the limit of $S_2^T[\rho_j]$ for periodic Gibbs states $\rho_j$ at $(K_j,\beta_j)$, the class of Proposition~\ref{C-351-7}, which is $\mathcal G^{\rm per}_{K_j,\beta_j}=\mathcal G^{\rm inv}_{K_j,\beta_j}$ and holds the states of the tori: $\rho_j=\nu_j$ for $S_2$ (Section~\ref{sec:state-limit-conventions}), and for $S_2^\xi$ the conditioned states of Theorems~\YMMref{AL}{746} and~\YMMref{AP}{780}. By Lemma~\ref{C-351-17}~(a) with $s=L^z$ the dilates $\rho_j^{(z)}$ are such states at $(K_j-z,\beta_j)$, and $S_2^T[\rho_j](\Phi)=L^{2z}S_2^T[\rho_j^{(z)}](\Phi_{L^z})$. As $a_{K_j-z}\beta_j=L^za_{K_j}\beta_j\to b'\in I^\circ(g')$, the pairs $(K_j-z,\beta_j)$ are first-passage data at $g'$ for large $j$ (Lemma~\ref{C-351-25}~(b)). Proposition~\ref{C-351-7} there, the $\rho_j^{(z)}$ being states at the datum, at the scale $L^zd\le\frac1{180}<\ell'/20$ and with $a_{K_j-z}\le9L^zd$ for large $j$, gives
\[
|S_2^T[\rho_j^{(z)}](\Psi_{L^z})-S^T_{2;K_j-z,\ell',\beta_j}((\Psi_{L^z})^{\langle\ell'\rangle})|\le2C_2^{\rm T}(g')\vartheta^{-1}(L^zd)^2(a_{K_j-z}\beta_j)^{-2}|\Psi| .
\]
Let $j\to\infty$, using the convergence along converging labels (Theorem~\textup{\YMMref{AS}{813}}), and multiply by $b'^2=L^{2z}b^2$; then $2C_2^{\rm T}(g')\vartheta^{-1}L^{2z}d^2=E_{\tau'}(d/b)^2$.
\end{proof}

\begin{lemma}[A second term excludes covariance]\label{C-351-16}
Let a functional $\sigma$ on separated tuples over $\mathbb R^3$ obey (L) with $(b,c,\varepsilon,d_7,\vartheta)$, $c\ne0$; let $d_1>0$, $\Psi\in\mathfrak T^1(d_1;\vartheta)$, $\mathsf B=c\,c_2b^{-3}\mathcal K_5(\Psi)$, $\mathsf u=b^{-3}d_1|\Psi|$ and $t>1$. Then
\[
\begin{multlined}
|S^\sigma_2(\Psi_r)-c_2b^{-2}\mathcal K_6(\Psi)-\mathsf Br|\le\varepsilon(rd_1)\,\mathsf u\,r\quad(rd_1\le d_7),\qquad|\mathcal N^{S_2^\sigma}_t(\Psi;r)+t(t-1)(t^5-1)\mathsf Br| \\
\le(t^2+t+t^6+t^7)\,\varepsilon(t^2rd_1)\,\mathsf u\,r\quad(t^2rd_1\le d_7);
\end{multlined}
\]
$\mathcal N^{S_2^\sigma}_t(\Psi;r)\ne0$, with the sign of $-c$, as soon as $(t^2+t+t^6+t^7)\varepsilon(t^2rd_1)<t(t-1)(t^5-1)\cdot10^{-5}|c|c_2\vartheta^2$; and $S_2^\sigma$ is not dilation covariant.
\end{lemma}

\begin{proof}[Proof sketch]
(L) at $\Psi_r\in\mathfrak T^1(rd_1;\vartheta)$, with $|\Psi_r|=|\Psi|$, $\mathcal K_6(\Psi_r)=\mathcal K_6(\Psi)$ and $\mathcal K_5(\Psi_r)=r\mathcal K_5(\Psi)$, gives the first bound. In $\mathcal N_t$ the weights $1,-(1+t^6),t^6$ at $t^2r,tr,r$ annul the constant $c_2b^{-2}\mathcal K_6(\Psi)$, give $\mathsf Br$ the factor $t^2-(1+t^6)t+t^6=-t(t-1)(t^5-1)$, and give the remainders a sum at most $(t^2+t+t^6+t^7)\varepsilon(t^2rd_1)\mathsf ur$, $\varepsilon$ being non-decreasing; at $t=2$ these are $-62$ and $198$. By~\eqref{C-351-2}, $|\mathsf B|\ge10^{-5}|c|c_2\vartheta^2\mathsf u$, which gives the sign. If $S_2^\sigma(\Phi_r)=r^\varpi S_2^\sigma(\Phi)$ for all $r>0$ and all separated $\Phi$, take $\Phi=\Psi$: by the first bound the left side tends to $c_2b^{-2}\mathcal K_6(\Psi)\in(0,\infty)$ as $r\to0$, so $S_2^\sigma(\Psi)\ne0$ and $r^\varpi$ has a finite non-zero limit, whence $\varpi=0$; then $S_2^\sigma(\Psi_r)=S_2^\sigma(\Psi)=c_2b^{-2}\mathcal K_6(\Psi)$ for all $r$, and the first bound gives $|\mathsf B|\le\varepsilon(rd_1)\mathsf u\to0$, against $|\mathsf B|\ge10^{-5}|c|c_2\vartheta^2\mathsf u>0$.
\end{proof}

\subsection{An alternative transfer using two blocking factors}

\begin{lemma}[Three facts on the lattices]\label{C-351-17}
\begin{enumerate}
\item[(a)] For $a,s,\beta>0$ the map $x\mapsto sx$ carries the Gibbs states of the Wilson action at $\beta$ on $a\mathbb Z^3$, and the periodic ones, onto those on $sa\mathbb Z^3$. For the image $D_s\nu$ of $\nu$, $S^T_n[\nu](f_1,\dots,f_n)=s^nS^T_n[D_s\nu](f_1(\cdot/s),\dots,f_n(\cdot/s))$; so $(sa\beta)^2S^T_2[D_s\nu](\Psi_{sr})=(a\beta)^2S^T_2[\nu](\Psi_r)$.
\item[(b)] For $\ell=2L^m$, a compact set $J'\subset(0,\infty)$ and $d_*>0$ there is $C$ with $|S^T_{2;K,\ell,L^Kb}(u,w)|\le C\|u\|_\sharp\|w\|_\sharp$ for all $K$, all $b\in J'$ and all real Lipschitz $u,w$ on $T_\ell$ at a sup-distance $\ge d_*$. Hence, if $K_j\to\infty$, $a_{K_j}\beta_j\to b>0$ and $r_j\to r>0$, then $S^T_{2;K_j,\ell,\beta_j}((\Psi_{r_j})^{\langle\ell\rangle})\to T^{(\ell,b)}(\Psi_r)$ for every separated pair $\Psi$ whose dilate $\Psi_r$ has both supports in one ball of radius $<\ell/16$; for $\Psi\in\mathfrak T^1(d_1;\vartheta)$ it suffices that $80rd_1<\ell$.
\item[(c)] If $\log L_1/\log L_2\notin\mathbb Q$ and $\sigma_0>0$, there are integers $K_j,K'_j\to\infty$ with $L_1^{K_j}L_2^{-K'_j}\to\sigma_0$.
\end{enumerate}
\end{lemma}

\begin{proof}[Proof sketch]
(a) The action and the Haar measure depend on the lattice as a graph, and $a^{-1}\sum_pf(x(p))a_p=s\,(sa)^{-1}\sum_pf(sx(p)/s)\,a_p$, the plaquette $p$ of $a\mathbb Z^3$ going to the plaquette $sp$ of $sa\mathbb Z^3$; cumulants are multilinear. (b) The bound is Theorem~\textup{\YMMref{AS}{813}} and its uniform bound~\YMMref{(2214)}{814}. For $C^1$ entries with compact supports $\|f(\cdot/r_j)-f(\cdot/r)\|_\sharp\to0$, and the dilated supports stay at a distance bounded below; so, by bilinearity, the pair may be fixed at $\Psi_r$, and the convergence along converging labels, Theorem~\textup{\YMMref{AS}{813}}, applies. A pair of $\mathfrak T^1(d_1;\vartheta)$ dilated by $r$ lies in two balls of radius $rd_1$ whose centres are at most $8rd_1$ apart, so in one ball of radius $5rd_1$. (c) The multiples $K\log L_1$ are dense modulo $\log L_2$ (Kronecker's theorem).
\end{proof}

\begin{lemma}[A second choice]\label{C-351-18}
Let $\varrho>0$, values of the constants of $\mathcal T^{\rm S}[2\varrho]$ with the blocking factor $L$, and an odd integer $Q\ge1$ be given. There is $j_0$ such that for every $j\ge j_0$ there is an admissible assignment of $\mathcal T^{\rm S}[2\varrho]$ in which the constants before $L$ keep their values and the blocking factor is $\tilde L=(LQ)^j+2$. $\tilde L$ is odd and prime to $L$ and to $Q$; $\log L/\log\tilde L\notin\mathbb Q$, and $\log Q/\log\tilde L\notin\mathbb Q$ if $Q>1$. With the radius and the mollifier of a choice $\mathfrak c$ it is a \emph{second choice} $\tilde{\mathfrak c}$ of $\mathfrak c$; at $Q=L$, $\tilde L=L^{2j}+2$.
\end{lemma}

\begin{proof}[Proof sketch]
By the ordered choices in the proofs of Proposition~\YMMref{69.1}{595} and Lemma~\YMMref{82.1}{665}, for an admissible assignment $\sigma_0$ of the constants before $L$ there is an odd integer $L_\uparrow(\sigma_0)$ such that $(\sigma_0,L')$ is admissible for every odd $L'\ge L_\uparrow(\sigma_0)$: the conditions attached to the blocking factor are oddness and floors formed with the earlier constants; and every admissible assignment up to a place of the order of the constants is the truncation of a full admissible assignment. For large $j$, $\tilde L\ge L_\uparrow$. $\tilde L$ is odd; a prime that divides $\tilde L$ and $L$ (or $Q$) divides $2$, which is impossible. If $\log L/\log\tilde L=p/q$ with integers $p,q\ge1$, then $L^q=\tilde L^p$ with coprime $L,\tilde L>1$, which is impossible; the same for $Q>1$.
\end{proof}

\begin{lemma}[Two blocking factors]\label{C-351-19}
Let $\mathfrak c_1,\mathfrak c_2$ be choices with $\log L_1/\log L_2\notin\mathbb Q$, $\tau_i=(g_i,b_i,\ell_i)$ an admissible triple of $\mathfrak c_i$ (for $i=1$, $b_1\in I(g_1)$ suffices), $\vartheta\in(0,1)$, $d_1>0$, $\Psi\in\mathfrak T^1(d_1;\vartheta)$, and $\delta>0$ with $\delta b_1d_1\le d^{(1)}_\flat(g_1,\vartheta)$ and $\delta b_2d_1<d^{(2)}_\flat(g_2,\vartheta)$, the thresholds of Proposition~\ref{C-351-7} at $\mathfrak c_1$ and $\mathfrak c_2$. With $T$, $E$ and $Y$ formed at the respective choice,
\[
|Y_{\tau_1}(\delta)-Y_{\tau_2}(\delta)|\le(E_{\tau_1}+E_{\tau_2})\,d_1^2|\Psi|\,\delta^2 .
\]
\end{lemma}

\begin{proof}[Proof sketch]
Lemma~\ref{C-351-17}~(c) gives $K_j,K'_j\to\infty$ with $s_j=L_2^{K'_j}L_1^{-K_j}\to b_1/b_2$. Put $\beta_j=L_1^{K_j}b_1$: $(K_j,\beta_j)$ is a first-passage datum of $\mathfrak c_1$ at $g_1$ with the label $b_1$, and $b_{2,j}=L_2^{-K'_j}\beta_j=b_1/s_j\to b_2\in I^\circ(g_2)$, so that $(K'_j,\beta_j)$ is a first-passage datum of $\mathfrak c_2$ at $g_2$ for large $j$ (Lemma~\ref{C-351-25}~(b) at $\mathfrak c_2$). Both data have the bare coupling $\beta_j$ on the lattice $\mathbb Z^3$ in lattice units. Let $\rho_j$ be a periodic Gibbs state at $\beta_j$ there, and $\rho_j^{(1)}$, $\rho_j^{(2)}$ its images on $L_1^{-K_j}\mathbb Z^3$ and $L_2^{-K'_j}\mathbb Z^3$; then $\rho_j^{(1)}=D_{s_j}\rho_j^{(2)}$, and by Lemma~\ref{C-351-17}~(a)
\[
b_1^2S^T_2[\rho^{(1)}_j](\Psi_{\delta b_1})=b_{2,j}^2S^T_2[\rho^{(2)}_j](\Psi_{\delta b_{2,j}}) .
\]
Proposition~\ref{C-351-7} at $\mathfrak c_1$ and at $\mathfrak c_2$ compares $\rho_j^{(1)}$ and $\rho_j^{(2)}$ with the Wilson measures on $T_{\ell_1}$ and $T_{\ell_2}$; after multiplication by $b_1^2$ and $b_{2,j}^2$ the errors are $E_{\tau_1}d_1^2|\Psi|\delta^2$ and, in the limit, $E_{\tau_2}d_1^2|\Psi|\delta^2$. As $j\to\infty$ the torus terms tend to $Y_{\tau_1}(\delta)$, by the limit at a fixed label along the full sequence, and to $Y_{\tau_2}(\delta)$, by Lemma~\ref{C-351-17}~(b), since $80\,\delta b_2d_1<1<\ell_2$.
\end{proof}

\begin{corollary}[Two labels; all real dilations]\label{C-351-23}
Let $\mathfrak c$ be a choice, $\tilde{\mathfrak c}$ a second choice of $\mathfrak c$ (Lemma~\ref{C-351-18} with $Q=L$) with an admissible triple $\tau_2=(g_2,b_2,\ell_2)$ and the threshold $\tilde d_\flat$, and $\vartheta\in(0,1)$.
\begin{enumerate}
\item[(i)] For admissible triples $\tau=(g,b,\ell)$ and $\tau'=(g',b',\ell')$ of $\mathfrak c$, $d_1>0$, $\Psi\in\mathfrak T^1(d_1;\vartheta)$ and $\delta>0$ with $\delta bd_1\le d_\flat(g,\vartheta)$, $\delta b'd_1\le d_\flat(g',\vartheta)$ and $\delta b_2d_1<\tilde d_\flat(g_2,\vartheta)$,
\[
|Y_\tau(\delta)-Y_{\tau'}(\delta)|\le(E_\tau+E_{\tau'}+2E_{\tau_2})\,d_1^2|\Psi|\,\delta^2 .
\]
Equivalently, with $s=b'/b$, for $\Phi\in\mathfrak T^1(d;\vartheta)$ with $d\le d_\flat(g,\vartheta)$, $sd\le d_\flat(g',\vartheta)$ and $(b_2/b)d<\tilde d_\flat(g_2,\vartheta)$: $|b'^2T^{(\ell',b')}(\Phi_s)-b^2T^{(\ell,b)}(\Phi)|\le(E_\tau+E_{\tau'}+2E_{\tau_2})(d/b)^2|\Phi|$. So, up to errors $O((d/b)^2)$, the limits on the tori at the labels $b$ and $b'$ are carried into each other by the real dilation by $b'/b$; for $b'/b$ a power of $L$ and $\ell'/\ell=b'/b$ this holds exactly, by Theorem~\ref{C-351-4}~(d).
\item[(ii)] If at one admissible triple $\tau_*=(g_*,b_*,\ell_*)$ of $\mathfrak c$ a function $T_*=T^{(\ell_*,b_*)}$ satisfies $|T_*(\Phi)-c_2b_*^{-2}[\mathcal K_6+cb_*^{-1}\mathcal K_5](\Phi)|\le\varepsilon(d)b_*^{-3}d|\Phi|$ for $0<d\le d_7$ and $\Phi\in\mathfrak T^1(d;\vartheta)$, with a real $c$ and a non-decreasing $\varepsilon\to0$, then at every admissible triple $\tau=(g,b,\ell)$ of $\mathfrak c$ the same inequality holds with the same $c$, with $b$ in the place of $b_*$ and with the tolerance $\varepsilon(b_*d/b)+(E_\tau+E_{\tau_*}+2E_{\tau_2})b^{-1}d$, for $d$ below the least of $(b/b_*)d_7$, $(b/b_*)d_\flat(g_*,\vartheta)$, $d_\flat(g,\vartheta)$ and $\frac12(b/b_2)\tilde d_\flat(g_2,\vartheta)$.

In particular the law of Theorem~\ref{C-351-4}~(a) at one admissible triple $(g^\circ,b^\circ,\ell_{g^\circ})$ with $g^\circ\le\Gamma$ gives the law with the same $\frac8{3\pi}$ at every admissible triple of $\mathfrak c$, on every side of $\mathfrak L(g)$, with the tolerance $\varepsilon^L_{g^\circ,\vartheta}(b^\circ d/b)+Cb^{-1}d$.
\item[(iii)] Let $S$ and $S'$ be volume-first limits at couplings $g$ and $g'$ with $0<g,g'\le\gamma_W$, with the labels $b=\mathsf b(S)$ and $b'=\mathsf b(S')$; let $\mathsf S$ be $S_2$ or a function $S_2^\xi$ of a canonical decomposition of $\sigma_S$, and $\mathsf S'$ be $S'_2$ or a function $S'^{\,\eta}_2$ of a canonical decomposition of $\sigma_{S'}$; let $z,z'\in\mathbb Z$, and let $\tau=(g_\tau,L^zb,\ell_\tau)$ and $\tau'=(g_{\tau'},L^{z'}b',\ell_{\tau'})$ be admissible triples of $\mathfrak c$. For $d_1>0$, $\Psi\in\mathfrak T^1(d_1;\vartheta)$ and $\delta>0$ with $\delta L^zbd_1\le d_\flat(g_\tau,\vartheta)$, $\delta L^{z'}b'd_1\le d_\flat(g_{\tau'},\vartheta)$ and $\delta b_2d_1<\tilde d_\flat(g_2,\vartheta)$,
\[
|b^2\mathsf S(\Psi_{\delta b})-b'^2\mathsf S'(\Psi_{\delta b'})|\le(2E_\tau+2E_{\tau'}+2E_{\tau_2})\,d_1^2|\Psi|\,\delta^2 .
\]
Equivalently, with $s=b'/b$, for $\Phi\in\mathfrak T^1(d;\vartheta)$ with $L^zd\le d_\flat(g_\tau,\vartheta)$, $L^{z'}sd\le d_\flat(g_{\tau'},\vartheta)$ and $(b_2/b)d<\tilde d_\flat(g_2,\vartheta)$: $|b'^2\mathsf S'(\Phi_s)-b^2\mathsf S(\Phi)|\le(2E_\tau+2E_{\tau'}+2E_{\tau_2})(d/b)^2|\Phi|$. So, up to errors $O((d/b)^2)$, the two-point functions of the volume-first limits and of all their components at the couplings $g\le\gamma_W$ are carried into each other by the real dilation by the quotient of their labels. The law of Theorem~\ref{C-351-3}~(1) at the two labels gives the same comparison with the error $2\varepsilon^{\rm nat}(d/b)(d/b)|\Phi|$, which is only $O((d/b)^{3/2})|\Phi|$.
\end{enumerate}
\end{corollary}

\begin{proof}[Proof sketch]
(i) Lemma~\ref{C-351-19} for the pair $(\mathfrak c,\tilde{\mathfrak c})$, whose blocking factors have incommensurable logarithms, at $(\tau,\tau_2)$ and at $(\tau',\tau_2)$, and the triangle inequality. For the second form put $\Phi=\Psi_{\delta b}$ and $d=\delta bd_1$; then $\Psi_{\delta b'}=\Phi_{b'/b}$ and $d_1^2\delta^2=(d/b)^2$. (ii) Let $\Phi\in\mathfrak T^1(d;\vartheta)$ and $\Phi_*=\Phi_{b_*/b}\in\mathfrak T^1(b_*d/b;\vartheta)$. By~(i), $b^2T^{(\ell,b)}(\Phi)$ is within $(E_\tau+E_{\tau_*}+2E_{\tau_2})(d/b)^2|\Phi|$ of $b_*^2T_*(\Phi_*)$, and by the assumption of~(ii) $b_*^2T_*(\Phi_*)$ is within $\varepsilon(b_*d/b)(d/b)|\Phi|$ of $c_2[\mathcal K_6(\Phi_*)+cb_*^{-1}\mathcal K_5(\Phi_*)]=c_2[\mathcal K_6(\Phi)+cb^{-1}\mathcal K_5(\Phi)]$; divide by $b^2$.

(iii) Lemma~\ref{C-351-15}, for $S$ with $z$ and $\tau$ at the pair $\Psi_{\delta b}\in\mathfrak T^1(\delta bd_1;\vartheta)$, and for $S'$ with $z'$ and $\tau'$ at the pair $\Psi_{\delta b'}$, gives $|b^2\mathsf S(\Psi_{\delta b})-Y_\tau(\delta)|\le E_\tau d_1^2|\Psi|\delta^2$ and $|b'^2\mathsf S'(\Psi_{\delta b'})-Y_{\tau'}(\delta)|\le E_{\tau'}d_1^2|\Psi|\delta^2$, as $(\Psi_{\delta b})_{L^z}=\Psi_{\delta L^zb}$ and $(d/b)^2=d_1^2\delta^2$ at $d=\delta bd_1$; (i) at $(\tau,\tau')$ and the triangle inequality give the bound, and the second form follows as in~(i). By Theorem~\ref{C-351-3}~(1) at $\Phi$ and at $\Phi_s\in\mathfrak T^1(sd;\vartheta)$, $b^2\mathsf S(\Phi)$ and $b'^2\mathsf S'(\Phi_s)$ lie within $\varepsilon^{\rm nat}(d/b)(d/b)|\Phi|$ of the one number $c_2[\mathcal K_6(\Phi)+\frac8{3\pi}b^{-1}\mathcal K_5(\Phi)]$, as $sd/b'=d/b$ and $b'^{-1}\mathcal K_5(\Phi_s)=b^{-1}\mathcal K_5(\Phi)$.
\end{proof}

\begin{proposition}[Non-covariance from the law at one triple]\label{C-351-20}
Let $\mathfrak c_*$ and $\mathfrak c^\circ$ be choices, equal or not, $\vartheta\in(0,1)$, and $\tau^\circ=(g^\circ,b^\circ,\ell^\circ)$ an admissible triple of $\mathfrak c^\circ$ such that, with $T^\circ=T^{(\ell^\circ,b^\circ)}$ formed at $\mathfrak c^\circ$, there are $d_n\to0$ and $\Psi_n\in\mathfrak T^1(d_n;\vartheta)$ with $|T^\circ(\Psi_n)-c_2b^{\circ-2}\mathcal K_6(\Psi_n)|/(d_n^2|\Psi_n|)\to\infty$. Then for every $0<g\le\gamma_W$ of $\mathfrak c_*$, every volume-first limit $S$ at $g$ and every canonical decomposition $(\mathsf X,\lambda,(\sigma_\xi))$ of $\sigma_S$:
\begin{enumerate}
\item[(i)] neither $S_2$ nor any $S_2^\xi$ is dilation covariant: $\mathsf N=\mathsf X$ and $\lambda(\mathsf N)=1$;
\item[(ii)] there is no $d_*>0$ such that the three-scale number $S_2(\Psi_{4r})-65S_2(\Psi_{2r})+64S_2(\Psi_r)$ vanishes for all $d,r$ with $4rd\le d_*$ and all $\Psi\in\mathfrak T^1(d;\vartheta)$.
\end{enumerate}
\end{proposition}

\begin{proof}[Proof sketch]
Let $b=\mathsf b(S)$ and $\mathsf S=S_2$ or $S_2^\xi$. Both assertions follow once a bound
\begin{equation}\label{C-351-21}
|\mathsf S(\Phi)-c_2b^{-2}\mathcal K_6(\Phi)|\le C\,d^2|\Phi|\qquad(0<d\le d_0,\ \Phi\in\mathfrak T^1(d;\vartheta))
\end{equation}
is excluded for every $C$ and $d_0>0$. If $\mathsf S$ is dilation covariant with the exponent $\varpi$, the leading laws~\YMMref{(2353)}{852} and~\YMMref{(2356)}{854} gives $\mathsf S(\Phi_r)\to c_2b^{-2}\mathcal K_6(\Phi)\in(0,\infty)$ as $r\to0$, so $\varpi=0$ and $\mathsf S=c_2b^{-2}\mathcal K_6$ on the pairs of all the classes $\mathfrak T^1(d;\vartheta)$: \eqref{C-351-21} with $C=0$. If the three-scale number vanishes as in~(ii), we may take $d_*\le4\mathsf d_0(g,\vartheta)$; for $\Phi\in\mathfrak T^1(d';\vartheta)$ with $d_*/8<d'\le d_*/4$ the numbers $u_k=S_2(\Phi_{2^{-k}})$ satisfy $u_k-65u_{k+1}+64u_{k+2}=0$, so $u_k=\alpha+\alpha'64^{-k}$; by the leading law $\alpha=\lim_ku_k=c_2b^{-2}\mathcal K_6(\Phi)$ and $|\alpha'|=|u_0-\alpha|\le\frac13c_2b^{-2}\mathcal K_6(\Phi)\le C|\Phi|$.

Every pair of $\mathfrak T^1(d;\vartheta)$ with $d\le d_*/4$ is $\Phi_{2^{-k}}$ for such a $\Phi$ and some $k\ge0$, so $|S_2(\Phi_{2^{-k}})-c_2b^{-2}\mathcal K_6(\Phi_{2^{-k}})|\le C|\Phi|\,2^{-6k}\le C(8/d_*)^6d^6|\Phi|$: \eqref{C-351-21} with $d^6$ in the place of $d^2$.

Suppose now~\eqref{C-351-21}. Take $z$ with $b'=L^zb\ge2(\gamma_W^{\rm sd})^{-2}$ and $g'=(2/b')^{1/2}$; then $b'\in I^\circ(g')$ by Lemma~\ref{C-351-25}~(a), as $b'\ge2(\gamma_W^{\rm sd})^{-2}\ge2\gamma^{-2}$, and $\tau'=(g',b',\ell_{g'})$ is an admissible triple of $\mathfrak c_*$. For $\Psi\in\mathfrak T^1(d_1;\vartheta)$, \eqref{C-351-21} at $\Phi=\Psi_{\delta b}$ and Lemma~\ref{C-351-15} give
\[
|Y_{\tau'}(\delta)-c_2\mathcal K_6(\Psi)|\le(Cb^4+E_{\tau'})\,d_1^2|\Psi|\,\delta^2
\]
for small $\delta$. Let $\tilde{\mathfrak c}$ be a second choice of $\mathfrak c_*$ with $Q$ the blocking factor of $\mathfrak c^\circ$ (Lemma~\ref{C-351-18}), and $\tau_2$ an admissible triple of $\tilde{\mathfrak c}$. Lemma~\ref{C-351-19}, for the pair $(\mathfrak c_*,\tilde{\mathfrak c})$ at $(\tau',\tau_2)$ and for the pair $(\mathfrak c^\circ,\tilde{\mathfrak c})$ at $(\tau^\circ,\tau_2)$, gives for small $\delta$
\[
|Y_{\tau^\circ}(\delta)-c_2\mathcal K_6(\Psi)|\le C'd_1^2|\Psi|\,\delta^2,\qquad C'=Cb^4+2E_{\tau'}+2E_{\tau_2}+E_{\tau^\circ}.
\]
At $d_1=1$, $\delta=d_n/b^\circ$ and $\Psi=(\Psi_n)_{1/d_n}$ this reads $|T^\circ(\Psi_n)-c_2b^{\circ-2}\mathcal K_6(\Psi_n)|\le C'b^{\circ-4}d_n^2|\Psi_n|$, against the assumption on $T^\circ$. As $\lambda$ is a probability measure, $\lambda(\mathsf N)=\lambda(\mathsf X)=1$.
\end{proof}

\subsection{A Gaussian model satisfying the same two-term law}

\begin{proposition}[A quadratic functional of a Gaussian field]\label{C-351-24}
Let $\ell>0$, $\alpha>0$, $\upsilon_0>0$, $\upsilon=\upsilon_0/\ell$ and $b>0$. Let $\varphi_1,\varphi_2,\varphi_3$ be independent centred Gaussian fields on $T_\ell$ with the covariance operator $\hat G_b(D)$, $\hat G_b(k)=(b|k|^2+\alpha|k|+\upsilon)^{-1}$ for $k\in\frac{2\pi}\ell\mathbb Z^3$, and let $\mathcal C$ be the kernel of $\hat G_b(D)$. A \emph{mollifier} is a real radial function $\jmath\in C^\infty_c(\mathbb B(0,1))$ with $\int\jmath=1$. For a mollifier $\jmath$ and $0<\varepsilon_{\rm m}<\frac\ell2$ put $\jmath_{\varepsilon_{\rm m}}=\varepsilon_{\rm m}^{-3}\jmath(\cdot/\varepsilon_{\rm m})$ and $\varphi_{a,\varepsilon_{\rm m}}=\varphi_a*\jmath_{\varepsilon_{\rm m}}$, and for a real Lipschitz function $f$ on $T_\ell$ let
\[
\mathcal O_{\varepsilon_{\rm m}}(f)=\sum_a\int f\,\bigl[-\tfrac12{:}|\nabla\varphi_{a,\varepsilon_{\rm m}}|^2{:}-{:}\varphi_{a,\varepsilon_{\rm m}}\Delta\varphi_{a,\varepsilon_{\rm m}}{:}\bigr]
\]
be the observable of the mollified field. For $f\ne0$ the variance of $\mathcal O_{\varepsilon_{\rm m}}(f)$ tends to $\infty$ as $\varepsilon_{\rm m}\to0$, so that there is no $L^2$ limit observable. For $n\ge2$ and real Lipschitz functions $f_1,\dots,f_n$ on $T_\ell$ whose supports have no common point we put
\[
\mathcal S^{(\ell,b)}_n(f_1,\dots,f_n):=\lim_{\varepsilon_{\rm m}\to0}\bigl\langle\mathcal O_{\varepsilon_{\rm m}}(f_1);\dots;\mathcal O_{\varepsilon_{\rm m}}(f_n)\bigr\rangle ,
\]
the limit of the joint cumulants of the observables of the mollified field; $\mathsf d_\infty$ is the spread of such a tuple, as in Theorem~\textup{\YMMref{AS}{813}}.
\begin{enumerate}
\item[(a)] $(bu^2+\alpha u+\upsilon)^{-1}=\int_0^\infty\rho_{\rm sp}(m)\,(u^2+m^2)^{-1}dm$ for $u\ge0$, with $\rho_{\rm sp}(m)=\frac{2\alpha}\pi\,m^2\bigl((\upsilon-bm^2)^2+\alpha^2m^2\bigr)^{-1}>0$ and $\int_0^\infty\rho_{\rm sp}=\frac1b$.
\item[(b)] The limit exists, depends on no mollifier, and is symmetric and multilinear, as long as the tuples keep supports without a common point; at tuples with pairwise disjoint supports it is the Wick integral of the kernels of the vertices and of $\mathcal C$. For every $n\ge2$, every compact subset of $(0,\infty)$ and every $d_0>0$ there is $c$ with $|\mathcal S^{(\ell,b)}_n(f_1,\dots,f_n)|\le c\prod_i\|f_i\|_\sharp$ for all labels $b$ of that subset and all such tuples with $\mathsf d_\infty\ge d_0$. At pairs $(f,h)$ of real Lipschitz functions on $T_\ell$ with disjoint supports, $\mathcal S_2^{(\ell,b)}(f,h)=\iint f(x)h(y)\,s_2(x-y)\,dx\,dy$ with a kernel $s_2$ smooth off $0$ and
\[
s_2(x)=c_2b^{-2}\Bigl[|x|^{-6}+\frac{4\alpha}{3\pi}\,b^{-1}|x|^{-5}\Bigr]+O\bigl(|x|^{-4}(1+\log^+(\ell/|x|))\bigr)\qquad(x\to0).
\]
So at every label $b$, for $(f,h)\in\mathfrak T_{T_\ell}(d;\vartheta)$ and small $d$, 
\[
|\mathcal S_2^{(\ell,b)}(f,h)-c_2b^{-2}[\mathcal K_6+\frac{4\alpha}{3\pi}b^{-1}\mathcal K_5](f,h)|\le C(\ell,\alpha,\upsilon_0,b)\,d^2(1+\log(\ell/d))\,\|f\|_\infty\|h\|_\infty
\]
: the two-term law with the coefficient $\frac{4\alpha}{3\pi}$, which is $\frac8{3\pi}$ at $\alpha=2$, and with a remainder $O(d^2\log\frac1d)$.
\item[(c)] Let $\{x_\mu=s\}$ be a coordinate mirror of $T_\ell$, $s$ any real number, $\theta$ the reflection in it and $\Sigma_+=\{s<x_\mu<s+\frac\ell2\}$ one of its open sides. Then $(w\circ\theta,\mathcal Cw)\ge0$ for every real function $w$ carried by $\Sigma_+$; hence the Gaussian measure is reflection positive. For a tuple $\mathcal F=(f_i)_{i\in I}$, $I$ a finite set, let $\mathfrak m(\mathcal F)=\sum_\pi\prod_{J\in\pi}\mathcal S^{(\ell,b)}_{|J|}((f_i)_{i\in J})$, the sum running over the partitions $\pi$ of $I$ into blocks of at least two elements, and $\mathfrak m(\varnothing)=1$; let $\theta\mathcal F=(f_i\circ\theta)_{i\in I}$, and let $\theta\mathcal F\cdot\mathcal F'$ be the tuple of the entries of $\theta\mathcal F$ and of $\mathcal F'$. Then
\[
\sum_{i,i'}\bar z_iz_{i'}\,\mathfrak m(\theta\mathcal F_i\cdot\mathcal F_{i'})\ge0
\]
for every finite family of complex numbers $z_i$ and of tuples $\mathcal F_i$, the empty tuple among them, of real Lipschitz functions with pairwise disjoint supports in compact subsets of $\Sigma_+$; the same holds for entries that overlap on $\Sigma_+$, as long as every cumulant that occurs stands at a tuple whose supports have no common point. The limits on the tori have the same property at the mirrors of the lattices of the cutoffs: for $\ell=2L^m$, every $b>0$, $\mu\in\{1,2,3\}$ and every $c\in\bigcup_{K\in\mathbb N}\frac12L^{-K}\mathbb Z$, with $\theta$ the reflection in $\{x_\mu=c\}$ and $\Sigma_+=\{c<x_\mu<c+\frac\ell2\}$, the inequality above holds with the $S^{(\ell,b)}_n$ in the place of the $\mathcal S^{(\ell,b)}_n$, for every finite family of complex numbers $z_i$ and of tuples $\mathcal F_i$, the empty tuple among them, of real Lipschitz functions with pairwise disjoint supports in compact subsets of $\Sigma_+$.

For entries of class $C^1$ the same inequality holds at every mirror $\{x_\mu=c\}$, $c$ any real number (the moving-mirror limit argument in the proof below, using Theorem~\YMMref{A}{5}). As the model is reflection positive at every coordinate mirror, it is so at all the mirrors at which the limits are. On the tuples whose supports have no common point, the family $(\mathcal S^{(\ell,b)}_n)_n$ is invariant under the isometries of $T_\ell$, and
\[
\mathcal S_n^{(r\ell,rb)}\bigl(f_1(\cdot/r),\dots,f_n(\cdot/r)\bigr)=r^{-n}\mathcal S_n^{(\ell,b)}(f_1,\dots,f_n)\qquad(r>0),
\]
the number $\upsilon_0$ being kept.
\item[(d)] For every such tuple $\underline f=(f_1,\dots,f_n)$ the map $b\mapsto\mathcal S_n^{(\ell,b)}(\underline f)$ extends holomorphically to the half plane ${\rm Re}\,b>0$, and
\[
\partial_b\mathcal S^{(\ell,b)}_n(\underline f)=-\mathcal S^{(\ell,b)}_{n+1}(1,\underline f),\qquad(-\partial_b)^{\mathsf j}\mathcal S^{(\ell,b)}_n(\underline f)=\mathcal S^{(\ell,b)}_{n+\mathsf j}(1,\dots,1,\underline f)\quad(\mathsf j\ge0),
\]
with $\mathsf j$ entries equal to the constant function $1$ on $T_\ell$: the rule that Proposition~\YMMref{118.3}{814}, equation~\YMMref{(2217)}{814} gives for the limits on the tori. Consequently, for $0<b_1<b_2$ and every pair $(f,h)$ of real Lipschitz functions on $T_\ell$ with disjoint supports,
\[
\int_{b_1}^{b_2}\bigl[b^2\mathcal S_3^{(\ell,b)}(1,f,h)-2b\,\mathcal S_2^{(\ell,b)}(f,h)\bigr]\,db=b_1^2\mathcal S_2^{(\ell,b_1)}(f,h)-b_2^2\mathcal S_2^{(\ell,b_2)}(f,h),
\]
the identity on which Corollary~\ref{C-351-22}~(c) rests.
\item[(e)] Let $\ell\in\{2L,2L^2,2L^3,\dots\}$, $b>0$ and $n\ge2$. The limits on the tori have the following property: if $\underline f^{(j)}$, $j\in\mathbb N$, are tuples of $n$ real Lipschitz functions on $T_\ell$, none of them $0$, whose supports have no common point, with spreads bounded below by a positive number, norms $\|f^{(j)}_i\|_\sharp$ bounded in $j$, and $\prod_i\|f^{(j)}_i\|_\infty\to0$, then $S^{(\ell,b)}_n(\underline f^{(j)})\to0$. The model fails this at $n=3$ and at every label: if $\psi\in C^\infty(T_\ell)$ is real and $h\ge0$ is Lipschitz, neither of them $0$, with ${\rm supp}\,\psi\cap{\rm supp}\,h=\varnothing$, and $f_{\mathbf q}=|\mathbf q|^{-1}\psi\cos(\mathbf q\cdot x)$ for $\mathbf q\in\frac{2\pi}\ell\mathbb Z^3\setminus\{0\}$, then the triples $(f_{\mathbf q},f_{\mathbf q},h)$ satisfy these conditions as $|\mathbf q|\to\infty$, and
\[
\mathcal S^{(\ell,b)}_3(f_{\mathbf q},f_{\mathbf q},h)\to-\frac3{2b}\iint\psi(x)^2h(y)\,\mathsf p(x-y)\,dx\,dy<0\qquad(|\mathbf q|\to\infty),
\]
with $\mathsf p=|\nabla\mathcal C|^2+2\,\mathcal C\Delta\mathcal C$, which is positive off $0$.
\end{enumerate}
\end{proposition}

\begin{proof}[Proof sketch]
(a) With $P(\zeta)=b\zeta^2+\alpha\zeta+\upsilon$, whose roots have negative real parts, $\rho_{\rm sp}(m)=\frac1{\pi i}\,m\,[P(-im)^{-1}-P(im)^{-1}]$, and closing the contour of $\int_{\mathbb R}m\,dm/((m^2+u^2)P(im))$ in the lower half plane, where the only pole is $m=-iu$, gives $1/P(u)$; letting $u\to\infty$ in $u^2/P(u)$ gives the mass $\frac1b$.
(b)

\emph{The limit and the bound.} For smooth $u,v$ the vertex $\mathsf V_f(u,v)=\int f\,\nabla u\cdot\nabla v+\int\nabla f\cdot\nabla(uv)$ equals $\int f[-\nabla u\cdot\nabla v-u\Delta v-v\Delta u]$, so the observable of the mollified field is, in each colour, the Wick ordered half of $\mathsf V_f$ at that field. The operator $V_f$ with $(u,V_fv)=\mathsf V_f(u,v)$ maps $H^1$ to $H^{-1}$ with a norm at most $\|f\|_\infty+{\rm Lip}\,f\le c\|f\|_\sharp$, and $\hat G_b(D)$ maps $H^{-1}$ to $H^1$, uniformly for $b$ in a compact subset of ${\rm Re}\,b>0$. The joint cumulant of half squares of a Gaussian vector is half the sum, over the cyclic orders of the entries, of the traces of the products of the forms against the covariance; so, the three colours giving the factor $3$, the joint cumulant of the $\mathcal O_{\varepsilon_{\rm m}}(f_i)$ is $\frac32$ times the sum, over the $(n-1)!$ cyclic orders of $f_1,\dots,f_n$, of the traces of the products of the operators $V_{f_i}|\hat\jmath(\varepsilon_{\rm m}D)|^2\hat G_b(D)$ in these orders, $\hat\jmath$ the Fourier transform of $\jmath$.

If the supports have no common point, split every propagator into a part of range less than $\mathsf d_\infty/n$ and a smooth remainder. The product of the short parts vanishes: the vertices are local, and a closed chain of $n$ short propagators through the $n$ supports would give points of all the supports at mutual distances less than the spread. Every other product contains a smooth propagator, a trace class operator, among bounded ones. So the products are trace class, with trace norms at most $c\prod_i\|f_i\|_\sharp$ for $b$ in a compact subset of ${\rm Re}\,b>0$ and $\mathsf d_\infty\ge d_0$, and their traces are continuous under the convergence of the propagators as $\varepsilon_{\rm m}\to0$.

This gives the limit, its independence of the mollifier, its symmetry and multilinearity, and the bound; at pairwise disjoint supports the traces are the integrals of $\prod_if_i(x_i)$ against the products of the kernels off the diagonal, the Wick integrals.

\emph{The variance.} At $n=2$ and $f_1=f_2=f$ the same trace formula is the variance of $\mathcal O_{\varepsilon_{\rm m}}(f)$, and it is a sum of terms $\ge0$. In the orthonormal basis $e_k=\ell^{-3/2}e^{ik\cdot x}$, $k\in\frac{2\pi}\ell\mathbb Z^3$, the matrix of $V_f$ has the entries $\int\bar e_p\,V_fe_q=\ell^{-3}v(p,q)\hat f(p-q)$, with $v(p,q)=|p|^2+|q|^2-p\cdot q=p\cdot q+|p-q|^2$ and $\hat f(k)=\int f(x)e^{-ik\cdot x}dx$: the first member of $\mathsf V_f$ gives $p\cdot q$, and the second, after one integration by parts, $|p-q|^2$. As $f$ is real, the entries at $(p,q)$ and at $(q,p)$ are conjugate; as $\hat\jmath$ is real, the mollified field has the covariance $\hat\jmath(\varepsilon_{\rm m}D)^2\hat G_b(D)$, and
\[
{\rm Var}\,\mathcal O_{\varepsilon_{\rm m}}(f)=\frac32\,\ell^{-6}\sum_{p,q\in\frac{2\pi}\ell\mathbb Z^3}\hat\jmath(\varepsilon_{\rm m}p)^2\,\hat\jmath(\varepsilon_{\rm m}q)^2\,\hat G_b(p)\,\hat G_b(q)\,v(p,q)^2\,|\hat f(p-q)|^2 ,
\]
every term of which is $\ge0$. Let $f\ne0$, and let $k_0$ have $\hat f(k_0)\ne0$. Without the mollifier, the terms with $p=q+k_0$ tend to $b^{-2}|\hat f(k_0)|^2>0$ as $|q|\to\infty$, since $v(q+k_0,q)=|q|^2+q\cdot k_0+|k_0|^2$ and $|q|^2\hat G_b(q)\to b^{-1}$; so their series diverges. As $\hat\jmath(\varepsilon_{\rm m}k)\to\hat\jmath(0)=\int\jmath=1$ at every $k$ when $\varepsilon_{\rm m}\to0$, Fatou's lemma for sums of non-negative terms gives ${\rm Var}\,\mathcal O_{\varepsilon_{\rm m}}(f)\to\infty$.

\emph{The kernel.} By (a) the covariance kernel is $\mathcal C=\int_0^\infty\rho_{\rm sp}(m)\mathcal C_m\,dm$, 
\[
\mathcal C_m(x)=\sum_{j\in\mathbb Z^3}e^{-m|x+\ell j|}/(4\pi|x+\ell j|)>0,
\]
 and $\mathcal H=\Delta\mathcal C=\int\rho_{\rm sp}(m)m^2\mathcal C_m\,dm$ off $0$. Since $m^2\rho_{\rm sp}(m)=\frac{2\alpha}{\pi b^2}+O((1+m^2)^{-1})$, the Laplace transform $\hat\rho_{\rm sp}$ of $\rho_{\rm sp}$ has $\hat\rho_{\rm sp}(0)=\frac1b$ and $\hat\rho_{\rm sp}''(|x|)=\frac{2\alpha}{\pi b^2}|x|^{-1}+O(1)$, so that, with $\omega_0=\frac1{4\pi b}$ and $\omega_1=\frac\alpha{2\pi^2b^2}$,
\begin{align*}
\mathcal C&=\frac{\omega_0}{|x|}+O(\lg(\ell/|x|)),\\
\nabla\mathcal C&=-\omega_0\frac{\hat x}{|x|^2}+O\Bigl(\frac{\lg(\ell/|x|)}{|x|}\Bigr),\\
\nabla\nabla^{\mathsf T}\mathcal C&=\omega_0\frac{3\hat x\hat x^{\mathsf T}-1}{|x|^3}
+\omega_1\frac{1-2\hat x\hat x^{\mathsf T}}{|x|^2}+O(|x|^{-1}),\\
\mathcal H&=\frac{\omega_1}{|x|^2}+O(|x|^{-1}),\\
\nabla\mathcal H&=-\frac{2\omega_1\hat x}{|x|^3}+O(|x|^{-2}),\\
\Delta\mathcal H&=\frac{2\omega_1}{|x|^4}+O(|x|^{-2}),
\end{align*}
with $\hat x=x/|x|$ and $\lg(\ell/|x|)=1+\log(\ell/|x|)$, as $|x|<\ell$ on $T_\ell$, the images under $\ell\mathbb Z^3$ being smooth near $0$. The observable is half the quadratic form of the vertex in each colour, so Wick's rule, with three colours, gives 
\[
s_2=\frac32[{\rm tr}(\nabla\nabla^{\mathsf T}\mathcal C)^2+4\nabla\mathcal C\cdot\nabla\mathcal H+2\mathcal H^2+2\mathcal C\Delta\mathcal H]
\]
 off $0$. The leading terms give $\frac32\cdot6\,\omega_0^2|x|^{-6}=c_2b^{-2}|x|^{-6}$; the terms of the order $|x|^{-5}$ are $-8\omega_0\omega_1$, $+8\omega_0\omega_1$ and $+4\omega_0\omega_1$, in all 
\[
\frac32\cdot4\,\omega_0\omega_1|x|^{-5}=\frac{3\alpha}{4\pi^3}b^{-3}|x|^{-5}=c_2\frac{4\alpha}{3\pi}b^{-3}|x|^{-5};
\]
the rest is $O(\lg(\ell/|x|)\,|x|^{-4})$. Integrating the remainder against a pair of the class, on whose supports $d\le|x-y|\le10d$, gives the bound.

(c) Write $x=(x_\mu,x_\perp)$, and let $w$ be carried by $\Sigma_+$. A Fourier series in $x_\perp$, with the frequency $k_\perp$ and $m_\perp=(|k_\perp|^2+m^2)^{1/2}$, writes $(w\circ\theta,\mathcal C_mw)$ as a sum over $k_\perp$ of double integrals, in the distances $\mathsf v,\mathsf v'\in(0,\frac\ell2)$ of two points of $\Sigma_+$ to the mirror, against the kernel $\cosh m_\perp(\frac\ell2-\mathsf v-\mathsf v')/(2m_\perp\sinh\frac{m_\perp\ell}2)$, and $\cosh m_\perp(\frac\ell2-\mathsf v-\mathsf v')=\cosh m_\perp(\frac\ell4-\mathsf v)\cosh m_\perp(\frac\ell4-\mathsf v')+\sinh m_\perp(\frac\ell4-\mathsf v)\sinh m_\perp(\frac\ell4-\mathsf v')$ makes each of them a sum of two squares; integrating against $\rho_{\rm sp}\ge0$ gives $(w\circ\theta,\mathcal Cw)\ge0$, and with it the reflection positivity of the Gaussian measure, as for every centred Gaussian measure with a reflection positive covariance.

For $\varepsilon_{\rm m}$ less than the distance of the supports from the mirror, the observables of the mollified field of functions carried by compact subsets of $\Sigma_+$ are functions of the field on $\Sigma_+$, and, the mollifier being radial, the reflection carries $\mathcal O_{\varepsilon_{\rm m}}(f)$ to $\mathcal O_{\varepsilon_{\rm m}}(f\circ\theta)$. The moments of centred variables are the sums, over the partitions without blocks of one element, of the products of their cumulants; so $\sum\bar z_iz_{i'}\mathfrak m_{\varepsilon_{\rm m}}(\theta\mathcal F_i\cdot\mathcal F_{i'})$, with $\mathfrak m_{\varepsilon_{\rm m}}$ formed with the cumulants of the mollified field, is the expectation of $\overline{\Theta F}\,F$ for $F=\sum_{i'}z_{i'}\prod_{f\in\mathcal F_{i'}}\mathcal O_{\varepsilon_{\rm m}}(f)$, $\Theta$ the reflection of the field, and it is $\ge0$.

As $\varepsilon_{\rm m}\to0$, each cumulant that occurs stands at a tuple whose supports have no common point and converges by~(b). For the limits on the tori, let $2c\in a_{K_0}\mathbb Z$. For every $K\ge K_0$ the plane $\{x_\mu=c\}$ is a site or a bond mirror of the lattice of the spacing $a_K$, at which the Wilson measure on $T_\ell$ at $\beta=L^Kb$ is reflection positive (Proposition~\YMMref{91.4}{700}). For large $K$ the function $F_K=\sum_{i'}z_{i'}\prod_{f\in\mathcal F_{i'}}(\mathcal O_K(f)-\langle\mathcal O_K(f)\rangle)$ reads only bonds in $\Sigma_+$, the measure is invariant under the reflection, and $\langle\overline{\Theta F_K}\,F_K\rangle\ge0$ is the inequality with the cumulants $S^T_{n;K,\ell,L^Kb}$ in the place of the $\mathcal S^{(\ell,b)}_n$.

At the tuples that occur, whose supports are pairwise disjoint,
the cumulants converge along the full sequence of cutoffs by
Theorem~\YMMref{A}{5}; hence the finite positive matrix converges
to the asserted matrix. For a real mirror position $c$, choose
$c_j\in\bigcup_K\frac12L^{-K}\mathbb Z$ with $c_j\to c$.
The fixed compact positive-side supports remain in the positive
half-torus for all large $j$, and their reflected supports remain in
disjoint compact neighbourhoods. For $C^1$ entries the translated
reflections converge in the sharp norm. The multilinear bound in
Theorem~\YMMref{AS}{813}, with the all-side dilation of
Corollary~\YMMref{119.2}{816}, therefore passes each cumulant,
and hence each centred moment, to the limit. Positivity is closed
under this finite matrix limit. This is the moving-mirror argument
of Proposition~\YMMref{92.4}{713}, here on the two half-tori.
It proves the assertion at every real $c$, for every $b>0$ and every
side $2L^m$ allowed by Theorem~\YMMref{A}{5}.

The invariance is that of $\hat G_b$ and of the mollifiers, which are radial; the dilation identity follows from $\hat G_{rb}(k/r)$ on $T_{r\ell}$ being $r\hat G_b(k)$ on $T_\ell$, with $\upsilon=\upsilon_0/\ell$: the field on $T_{r\ell}$ has the law of $x\mapsto r^{-1}\varphi(x/r)$, the mollifier at the scale $\varepsilon_{\rm m}$ goes to that at the scale $r\varepsilon_{\rm m}$, and each observable, with its two derivatives, scales by $r^{-1}$; the identity passes to the limits.

(d) The trace formula of~(b) defines $\mathcal S^{(\ell,b)}_n$ for ${\rm Re}\,b>0$, and its terms are holomorphic in $b$, locally uniformly, as $\hat G_b$ is; so $\mathcal S^{(\ell,b)}_n(\underline f)$ is holomorphic there. The vertex of the constant function $1$ is $V_1=-\Delta$, and $\partial_b\hat G_b=-|k|^2\hat G_b^2$, that is, $\partial_b\hat G_b(D)=-\hat G_b(D)V_1\hat G_b(D)$: the derivative of a propagator inserts the vertex of $1$ between its ends. Summing over the $n$ propagators of each product and over the cyclic orders of $f_1,\dots,f_n$ gives the cyclic orders of $(1,f_1,\dots,f_n)$; hence $\partial_b\mathcal S_n^{(\ell,b)}(\underline f)=-\mathcal S_{n+1}^{(\ell,b)}(1,\underline f)$, and the higher derivatives follow by induction, the tuples with entries $1$ having supports without a common point.

At the entry $1$ the observable of the mollified field is $\frac12\sum_a\int{:}|\nabla\varphi_{a,\varepsilon_{\rm m}}|^2{:}$, since $\int\varphi\Delta\varphi=-\int|\nabla\varphi|^2$ on the torus. The integral identity is the integral over $[b_1,b_2]$ of $\partial_b[-b^2\mathcal S_2^{(\ell,b)}(f,h)]=b^2\mathcal S_3^{(\ell,b)}(1,f,h)-2b\,\mathcal S_2^{(\ell,b)}(f,h)$.

(e) For the limits, let $d_0>0$ bound the spreads from below. By Theorem~\textup{\YMMref{AS}{813}}, at the interval $\{b\}$ and with $d_0$ in the place of $\mathsf d$, for all large $K$ the term of index $K$ of the sequence of the cutoffs differs from $S^{(\ell,b)}_n(\underline f^{(j)})$ by at most $\prod_i\|f^{(j)}_i\|_\sharp$ times a number that tends to $0$ as $K\to\infty$ and does not depend on $j$; and, as $|\mathcal O_K(f)|\le6\ell^3L^{4K}\|f\|_\infty$, the term of index $K$ is at most a number depending only on $n$ times $\prod_i(6\ell^3L^{4K}\|f^{(j)}_i\|_\infty)$, as in the first member of~\YMMref{(2214)}{814}. Let $j\to\infty$ and then $K\to\infty$.

For the model, $\|f_{\mathbf q}\|_\infty\le|\mathbf q|^{-1}\|\psi\|_\infty$, ${\rm Lip}\,f_{\mathbf q}$ is bounded in $\mathbf q$, and the spread of $(f_{\mathbf q},f_{\mathbf q},h)$ is at least the distance of ${\rm supp}\,\psi$ and ${\rm supp}\,h$ in the norm $|\cdot|_\infty$; so the conditions hold. By~(b), $\mathcal S^{(\ell,b)}_3(f_{\mathbf q},f_{\mathbf q},h)$ is three times the trace of $V_{f_{\mathbf q}}\hat G_b(D)V_{f_{\mathbf q}}\hat G_b(D)V_h\hat G_b(D)$, the two cyclic orders giving equal traces. Now $V_{f_{\mathbf q}}u=-f_{\mathbf q}\Delta u-\nabla\cdot(u\nabla f_{\mathbf q})$; the first member tends to $0$, and the second, of the order $1$ in $H^{-1}$, carries the frequencies $\pm\mathbf q+O(1)$, at which $\hat G_b(k)kk^{\mathsf T}\to b^{-1}\hat{\mathbf q}\hat{\mathbf q}^{\mathsf T}$, $\hat{\mathbf q}=\mathbf q/|\mathbf q|$.

So $V_{f_{\mathbf q}}\hat G_b(D)V_{f_{\mathbf q}}$ tends to the multiplication by $(2b)^{-1}\psi^2$, the mean of $b^{-1}\psi^2\sin^2(\mathbf q\cdot x)$; as the supports of $\psi$ and $h$ are disjoint, $\hat G_b(D)V_h\hat G_b(D)$ has a smooth kernel near ${\rm supp}\,\psi$, and the traces tend to $\frac3{2b}{\rm Tr}[\psi^2\hat G_b(D)V_h\hat G_b(D)]=-\frac3{2b}\iint\psi(x)^2h(y)\,\mathsf p(x-y)\,dx\,dy$, since $\mathsf V_h(u,u)=-\int h(|\nabla u|^2+2u\Delta u)$ at $u=\mathcal C(x-\cdot)$. Finally $\mathsf p=|\nabla\mathcal C|^2+2\mathcal C\mathcal H>0$ off $0$, as $\mathcal C>0$ and $\mathcal H>0$ there, by the representation of $\mathcal C$ above.
\end{proof}

\subsection{Proof of Theorem~\ref{C-351-4}}

\emph{Step 1: the choice.} At $\pi^+(\varsigma^{\rm S})$ every constant keeps its value (the unchanged restriction of Section~\YMMref{82}{665}); so $L$, $I(g)$, $\mathfrak L(g)$, $\ell_g$, the first-passage data and the functions $S^T_{2;K,\ell,\beta}$ formed there are those of $\mathfrak c$, and Proposition~\ref{C-351-14} speaks of $\mathfrak c$.

\emph{Step 2: (a).} Put $d^L=\min\{d_{\rm e},\ell_g/21\}$. For $d\le d^L$ and $\Psi\in\mathfrak T_{\mathbb R^3}(d;\vartheta)$ the periodisation $\Psi^{\langle\ell_g\rangle}$ lies in $\mathfrak T_{T_{\ell_g}}(d;\vartheta)$ and has the suprema and the kernels $\mathcal K_q$ of $\Psi$, and Proposition~\ref{C-351-14}~(b) gives
\[
|T^{(\ell_g,b)}(\Psi)-L_b^+(\Psi)|\le\bar\varepsilon_\infty(d)\,b^{-3}d\,|\Psi|\qquad(b\in I(g)).
\]
Let $\varepsilon^L_{g,\vartheta}(d)$ be the supremum, over $0<d''\le d$, $b\in I(g)$ and $\Psi\in\mathfrak T_{\mathbb R^3}(d'';\vartheta)$, of $|T^{(\ell_g,b)}(\Psi)-L_b^+(\Psi)|/(b^{-3}d''|\Psi|)$: the least non-decreasing tolerance; it is at most $\bar\varepsilon_\infty\le C(g,\vartheta)d^{1/2}$.

\emph{Step 3: the consequences (i)--(iv).} (a) at $\Psi_r\in\mathfrak T^1(rd_1;\vartheta)$, with $\mathcal K_6(\Psi_r)=\mathcal K_6(\Psi)$, $\mathcal K_5(\Psi_r)=r\mathcal K_5(\Psi)$ and $|\Psi_r|=|\Psi|$, gives~\eqref{C-351-5}: a term of the order $r$ against an error $O(r^{3/2})$ and the threshold $r^2$ of a dilation covariant function; the bounds on $\mathsf B$ are~\eqref{C-351-2}. Write $T(\Psi_r)=\mathsf A+\mathsf Br+\rho(r)$ with $|\rho(r)|\le\varepsilon^L(rd_1)\mathsf ur$; then $T(\Psi_{tr})-T(\Psi_r)=(t-1)\mathsf Br+\rho(tr)-\rho(r)$ with $|\rho(tr)-\rho(r)|\le(t+1)\varepsilon^L(trd_1)\mathsf ur$, and $\mathcal N^T_t$ is treated as in Lemma~\ref{C-351-16}: (i). For $\Phi\in\mathfrak T^1(d;\vartheta)$ and $d\le d^L$, by~\eqref{C-351-2},
\[
|T(\Phi)-c_2b^{-2}\mathcal K_6(\Phi)|\ge\varkappa b^{-3}\mathcal K_5(\Phi)-\varepsilon^L(d)\,b^{-3}d|\Phi|\ge[10^{-5}\varkappa\vartheta^2-\varepsilon^L(d)]\,b^{-3}d\,|\Phi| ;
\]
at $\Phi=\Psi_r$, $d=rd_1$, divided by $r^2$, this is (ii), and at $\Phi=\Psi_n$, divided by $d_n^2|\Psi_n|$, it is (iii). Under $\varepsilon^L(d)\le\frac12\cdot10^{-5}\varkappa\vartheta^2$ the same line gives $b^2|T(\Psi)-c_2b^{-2}\mathcal K_6(\Psi)|\ge\frac12\cdot10^{-5}\varkappa\vartheta^2b^{-1}d|\Psi|$, which exceeds $C\vartheta^{-1}(d/b)^2|\Psi|$ if and only if $d/b<\frac12\cdot10^{-5}\varkappa\vartheta^3/C$; such $d$ exist, as $\varepsilon^L\to0$: (iv).

\emph{Step 4: (b).} As $d_\flat\le\frac1{180}<\ell_g/21$, Step~2 applies. At the data $(K,L^Kb)$ with $a_K\le9d$ the periodic extensions of the Wilson measures on $T_\ell$ and $T_{\ell_g}$ are two states at the datum; Proposition~\ref{C-351-7} and $K\to\infty$, along the full sequence, give
\[
|T^{(\ell,b)}(\Psi)-T^{(\ell_g,b)}(\Psi)|\le2C_2^{\rm T}(g)\vartheta^{-1}d^2b^{-2}|\Psi|=2C_2^{\rm T}(g)\vartheta^{-1}bd\cdot b^{-3}d\,|\Psi| ,
\]
and Proposition~\ref{C-351-14}~(b) on $T_{\ell_g}$ gives the first assertion; $\bar\varepsilon_\infty(d)\le C(g,\vartheta)d^{1/2}$ and $d<1$ bound the bracket. At a cutoff $K>s(d)$ one has $K\ge3$ and $8a_K\le d$, so $(K,L^Kb)$ is a first-passage datum at $g$ at which Proposition~\ref{C-351-7}~(a) compares every state with the periodic extension $\mu'$ of the Wilson measure on $T_{\ell_g}$; and $S^T_2[\mu'](\Psi)=S^T_{2;K,\ell_g,L^Kb}(\Psi^{\langle\ell_g\rangle})$, to which Proposition~\ref{C-351-14}~(b), respectively~(c), applies.

\emph{Step 5: (d).} This is~\YMMref{(2215)}{814} at $n=2$: the dilation by $L^s$ carries $T_\ell$ with the spacing $a_K$ onto $T_{L^s\ell}$ with the spacing $a_{K-s}$, keeps $\beta=L^Kb=L^{K-s}(L^sb)$, and $\mathcal O_K(f)=L^s\mathcal O_{K-s}(f(\cdot/L^s))$; so $S^T_{2;K,\ell,L^Kb}(\Psi^{\langle\ell\rangle})=L^{2s}S^T_{2;K-s,L^s\ell,L^{K-s}L^sb}((\Psi_{L^s})^{\langle L^s\ell\rangle})$, and we let $K\to\infty$. For the second sentence let $\Psi\in\mathfrak T_{\mathbb R^3}(d;\vartheta)$ with $d\le L^zd^L$; then $\Psi_{L^{-z}}\in\mathfrak T_{\mathbb R^3}(L^{-z}d;\vartheta)$, $T^{(L^z\ell_g,L^zb)}(\Psi)=L^{-2z}T^{(\ell_g,b)}(\Psi_{L^{-z}})$, $L^{-2z}L_b^+(\Psi_{L^{-z}})=L^+_{L^zb}(\Psi)$, and $L^{-2z}b^{-3}L^{-z}d=(L^zb)^{-3}d$; apply~(a) at $\Psi_{L^{-z}}$.

\emph{Step 6: (c).} For large $s$ the coupling $g''=(2/(L^sb))^{1/2}$ is $\le\Gamma$, $L^sb\in I^\circ(g'')$ by Lemma~\ref{C-351-25}~(a), as $L^sb\ge2\Gamma^{-2}\ge2\gamma^{-2}$ for large $s$, and $L^s\ell\in\mathfrak L(g'')$ (equation~\YMMref{(1528)}{615}), as $m_{\mathfrak L}(g'')$ grows only like a multiple of $\log\log g''^{-2}$ while $m+s$ grows like $s$. Apply~(b) at $(g'',L^sb,L^s\ell)$ to $\Psi_{L^s}\in\mathfrak T^1(L^sd;\vartheta)$, and~(d), with
\[
\mathcal K_6(\Psi_{L^s})=\mathcal K_6(\Psi),\qquad(L^sb)^{-1}\mathcal K_5(\Psi_{L^s})=b^{-1}\mathcal K_5(\Psi),\qquad L^{2s}(L^sb)^{-3}L^sd=b^{-3}d ;
\]
the bracket of~(b) at $(g'',L^sb,L^sd)$ is $\bar\varepsilon_\infty(L^sd)+2C_2^{\rm T}(g'')\vartheta^{-1}L^{2s}bd\le C(g'',\vartheta)(1+L^sb)(L^sd)^{1/2}$.

\emph{Step 7: (e), (f).} Two inequalities as in~(e), at a pair $\Psi\in\mathfrak T^1(d;\vartheta)$, give $|c-c'|\,c_2b^{-3}\mathcal K_5(\Psi)\le(\varepsilon+\varepsilon')(d)\,b^{-3}d|\Psi|$, and $\mathcal K_5(\Psi)\ge10^{-5}\vartheta^2d|\Psi|$; so $|c-c'|\,c_2\,10^{-5}\vartheta^2\le(\varepsilon+\varepsilon')(d)\to0$. For a compact set $J$ of labels one $s$ serves for all $b\in J$, and the compact set $L^sJ$ is covered by finitely many open intervals $I^\circ(g_i)$, $g_i\le\Gamma$, $L^s\ell\in\mathfrak L(g_i)$, since each point $b''$ of $L^sJ$ lies in $I^\circ((2/b'')^{1/2})$ (Lemma~\ref{C-351-25}~(a)); the least of the ranges and the largest of the tolerances and constants of the $g_i$ serve for all $b\in J$. \qed

\subsection{Proof of Theorem~\ref{C-351-3}}

\emph{Step 1: a period of labels.} Let $\gamma_o=\min\{\gamma_{\rm e},\gamma_W,\gamma_W^{\rm sd}\}$, $y=\gamma_o^{-2}$, $\Lambda=2y$, $\hat g_1=\gamma_o$ and $\hat g_2=\gamma_o/\sqrt2$. As $y\ge\gamma_{\rm e}^{-2}\ge4\bar B$ and $y\ge4$, we have $y>\bar B$, so that
\[
I^\circ(\hat g_1)=(y+\bar B,\,Ly+(L-2)\bar B)\supset[2y,Ly),\qquad I^\circ(\hat g_2)=(2y+\bar B,\,2Ly+(L-2)\bar B)\supset[3y,2Ly),
\]
and, as $L\ge3$, $[\Lambda,L\Lambda)\subset I^\circ(\hat g_1)\cup I^\circ(\hat g_2)$. Put $\delta_{\rm c}=(L\Lambda)^{-1}\min_i\min\{d^L,d_\flat\}(\hat g_i,\vartheta)$ and, for $0<\delta\le\delta_{\rm c}$,
\[
\varepsilon^{\rm nat}(\delta)=\max_{i=1,2}\varepsilon^L_{\hat g_i,\vartheta}(L\Lambda\delta)+2\max_{i=1,2}C_2^{\rm T}(\hat g_i)\,\vartheta^{-1}L^2\Lambda^2\delta :
\]
non-decreasing, depending only on $\mathfrak c$ and $\vartheta$, and $\varepsilon^{\rm nat}(\delta)\le C\delta^{1/2}$ by Theorem~\ref{C-351-4}~(a).

\emph{Step 2: a dilated label.} For $S$ with the label $b$ there is exactly one $z\in\mathbb Z$ with $b'=L^zb\in[\Lambda,L\Lambda)$; take $i$ with $b'\in I^\circ(\hat g_i)$. As $\hat g_i\le\gamma_W^{\rm sd}$, the triple $\tau'=(\hat g_i,b',\ell_{\hat g_i})$ is admissible, and $E_{\tau'}=2C_2^{\rm T}(\hat g_i)\vartheta^{-1}b'^2\le2C_2^{\rm T}(\hat g_i)\vartheta^{-1}L^2\Lambda^2$.

\emph{Step 3: (1).} Let $d\le\delta_{\rm c}b$, $\Psi\in\mathfrak T^1(d;\vartheta)$ and $d'=L^zd=b'd/b\le L\Lambda\delta_{\rm c}$. Theorem~\ref{C-351-4}~(a) at $(\hat g_i,b')$, where $\hat g_i\le\gamma_{\rm e}$ and $b'\in I(\hat g_i)$, for the pair $\Psi_{L^z}\in\mathfrak T^1(d';\vartheta)$, with $\mathcal K_6(\Psi_{L^z})=\mathcal K_6(\Psi)$, $b'^{-1}\mathcal K_5(\Psi_{L^z})=b^{-1}\mathcal K_5(\Psi)$ and $b'^{-1}d'=b^{-1}d$, together with Lemma~\ref{C-351-15}, gives
\[
\bigl|b^2\mathsf S(\Psi)-c_2[\mathcal K_6+\tfrac8{3\pi}b^{-1}\mathcal K_5](\Psi)\bigr|\le\Bigl[\varepsilon^L_{\hat g_i,\vartheta}\Bigl(L\Lambda\frac db\Bigr)+2C_2^{\rm T}(\hat g_i)\vartheta^{-1}L^2\Lambda^2\frac db\Bigr]\frac db\,|\Psi|\le\varepsilon^{\rm nat}(d/b)\,b^{-1}d\,|\Psi| ;
\]
divide by $b^2$. The labels of first-passage data at $g$ exceed $g^{-2}$ by~\YMMref{(1531)}{615}, so $b\ge g^{-2}$, $d/b\le g^2d$, and $d\le d^S$ implies $d\le\delta_{\rm c}b$: the law with $\varepsilon^S$ and $d^S$. For the form at a dilated label, let $L^zb\ge2\Gamma^{-2}$ and $g'=(2/(L^zb))^{1/2}$; then $g'\le\Gamma$ and $L^zb\in I^\circ(g')$ by Lemma~\ref{C-351-25}~(a), as $2\Gamma^{-2}\ge2\gamma^{-2}$, so $\tau'=(g',L^zb,\ell_{g'})$ is admissible; the same computation, with Theorem~\ref{C-351-4}~(a) at $(g',L^zb)$ and Lemma~\ref{C-351-15}, gives the tolerance $\varepsilon^L_{g',\vartheta}(L^zd)+E_{\tau'}b^{-1}d$ for $L^zd\le\min\{d^L,d_\flat\}(g',\vartheta)$, and $E_{\tau'}b^{-1}=2C_2^{\rm T}(g')\vartheta^{-1}L^{2z}b$.

\emph{Step 4: (2).} Lemma~\ref{C-351-16} with $c=\frac8{3\pi}$, $\varepsilon(d)=\varepsilon^{\rm nat}(d/b)$ and $d_7=\delta_{\rm c}b$; there $\mathsf B=cc_2b^{-3}\mathcal K_5(\Psi)=B_S$, and the lower bound on $B_S$ is~\eqref{C-351-2}. The first bound at $tr$ and at $r$ gives the bound on $r^{-1}[\mathsf S(\Psi_{tr})-\mathsf S(\Psi_r)]$; $\varepsilon^{\rm nat}(\delta)\le C\delta^{1/2}$ turns the remainders $o(r)$ into $O(r^{3/2})$; the means converge because their terms do; and at $t=2$ the numbers $t(t-1)(t^5-1)$ and $t^2+t+t^6+t^7$ are $62$ and $198$.

\emph{Step 5: (3).} By Lemma~\ref{C-351-16}, $S_2$ and every $S_2^\xi$ fail to be dilation covariant; so $\mathsf N=\mathsf X$, and $\lambda(\mathsf N)=1$ as $\lambda$ is a probability measure.

\emph{Step 6: (3) from the law at a single triple of any choice.} Let $b^\circ\in I^\circ(\Gamma)$. The triple $(\Gamma,b^\circ,\ell_\Gamma)$ is admissible, and clause~(iii) of Theorem~\ref{C-351-4} holds there, as $\Gamma\le\gamma_{\rm e}$. Proposition~\ref{C-351-20} with $\mathfrak c_*=\mathfrak c^\circ=\mathfrak c$ gives~(3) again, and moreover the failure of the three-scale number of $S_2$ to vanish at pairs and $r$ as small as is wished. This second proof uses the law at one triple and one $\vartheta$ only, and no law at the labels of the limits; the triple may belong to any choice. \qed

\begin{remark}[The constants in the second proof of Theorem~\ref{C-351-3}~(3)]\label{C-351-6}
The departure of clause~(iv) of Theorem~\ref{C-351-4} dominates explicitly the errors of the comparisons by which the proof of Proposition~\ref{C-351-20} passes from a limit to the least torus, when the label of the limit is first dilated into one period $[\Lambda_W,L\Lambda_W)$. Let $g\le\Gamma$ and $b\in I^\circ(g)$; let $\tilde{\mathfrak c}$ be a second choice of $\mathfrak c$ (Lemma~\ref{C-351-18} with $Q=L$), with its number $\tilde C_2^{\rm T}$ and an admissible triple $\tau_2=(g_2,b_2,\ell_2)$; put $\Lambda_W=2(\gamma_W^{\rm sd})^{-2}$, $g'_1=\gamma_W^{\rm sd}$ and $g'_2=\gamma_W^{\rm sd}/\sqrt2$. Since $(\gamma_W^{\rm sd})^{-2}\ge\gamma^{-2}>\bar B$ (Lemma~\ref{C-351-25}~(a)), the computation of Step~1 above gives $[\Lambda_W,L\Lambda_W)\subset I^\circ(g'_1)\cup I^\circ(g'_2)$; every $b'>0$ has exactly one dilate $L^zb'$ in $[\Lambda_W,L\Lambda_W)$. The number
\[
c_{\max}=\max_{i=1,2}\ \sup\bigl\{2C_2^{\rm T}(g'_i)\,b'^2:\ b'\in[\Lambda_W,L\Lambda_W)\cap I^\circ(g'_i)\bigr\}\le2L^2\Lambda_W^2\max_iC_2^{\rm T}(g'_i)
\]
is finite, and so is $\mathsf C_{\rm pair}=2C_2^{\rm T}(g)b^2+4\tilde C_2^{\rm T}(g_2)b_2^2+2c_{\max}$. For every admissible triple $\tau'=(g'_i,b',\ell')$ with $b'\in[\Lambda_W,L\Lambda_W)\cap I^\circ(g'_i)$,
\[
\vartheta^{-1}\mathsf C_{\rm pair}\ge2E_{\tau'}+2E_{\tau_2}+E_{(g,b,\ell_g)},
\]
the sum of the constants of the three comparisons by which the proof of Proposition~\ref{C-351-20}, with $\tau'$ of this kind and $\tau^\circ=(g,b,\ell_g)$, passes from a limit at the label $b'$, through the second choice, to $T^{(\ell_g,b)}$. Clause~(iv) of Theorem~\ref{C-351-4}, at $C=\mathsf C_{\rm pair}$, says that for small $d/b$ the departure $b^2|T^{(\ell_g,b)}(\Psi)-c_2b^{-2}\mathcal K_6(\Psi)|$ exceeds $(2E_{\tau'}+2E_{\tau_2}+E_{(g,b,\ell_g)})(d/b)^2|\Psi|$ at every pair $\Psi\in\mathfrak T^1(d;\vartheta)$.
\end{remark}

\subsection{Proof of Corollary~\ref{C-351-22}}

\emph{Step 1: pieces.} Let $\delta_0={\rm dist}({\rm supp}f_0,{\rm supp}h_0)>0$ and $\rho_0=\delta_0/4$, and let $\phi_i,\chi_k\ge0$, finitely many, be Lipschitz functions supported in closed balls of radius $\rho_0$ that lie within $\delta_0/4$ of ${\rm supp}f_0$, respectively of ${\rm supp}h_0$, with $\sum_i\phi_i=1$ on ${\rm supp}f_0$ and $\sum_k\chi_k=1$ on ${\rm supp}h_0$. With $M_0=\|f_0\|_\infty+1$, $f_0=\sum_i[(f_0+M_0)\phi_i-M_0\phi_i]$ is a signed sum of non-negative Lipschitz pieces, and likewise $h_0$; we drop the pieces that vanish. A piece $F$ of $f_0$ and a piece $H$ of $h_0$ lie in balls of radius $\rho_0$ at the distance $\delta'={\rm dist}({\rm supp}F,{\rm supp}H)\ge\delta_0/2=2\rho_0$; so $(F,H)\in\mathfrak T_{\mathbb R^3}(\delta'/2;\vartheta_0)$ for one $\vartheta_0\in(0,1)$ that serves for all the finitely many pairs of pieces.

\emph{Step 2: one pair of pieces.} Let $s$ and $g''$ be as in Step~6 of the proof of Theorem~\ref{C-351-4}. Clauses~(b) at $(g'',L^sb,L^s\ell)$ and~(d), which hold on all of $\mathfrak T_{\mathbb R^3}$, applied to the translate by $\tilde p$ of $(F,H)_r\in\mathfrak T_{\mathbb R^3}(r\delta'/2;\vartheta_0)$, with $\mathcal K_6((F,H)_r)=\mathcal K_6(F,H)$, $\mathcal K_5((F,H)_r)=r\,\mathcal K_5(F,H)$ and~\eqref{C-351-2}, give for small $r$, with a tolerance $\varepsilon$ independent of the position,
\[
\bigl|T^{(\ell,b)}((F,H)_r(\cdot-\tilde p))-c_2b^{-2}\mathcal K_6(F,H)\bigr|\le\Bigl[\frac8{3\pi}+\varepsilon(r\delta'/2)\Bigr]\,b^{-3}r\,\frac{\delta'}2\,\|F\|_\infty\|H\|_\infty .
\]

\emph{Step 3: (a).} $S_2^{(\ell,b)}$ is bilinear on the functions supported in two disjoint closed sets, and so is $\mathcal K_6$; summing over the pairs of pieces, with their signs, gives $|S_2^{(\ell,b)}(f_r,h_r)-c_2b^{-2}\mathcal K_6(f_0,h_0)|\le Cr$, with $C$ independent of $p$ and, by the covering argument of Step~7 of the proof of Theorem~\ref{C-351-4}, the same for the labels of a compact set. For $(\psi_0\circ\theta_\mu,\psi_0)$ the supports are compact subsets of $\{x_\mu<0\}$ and $\{x_\mu>0\}$, so at a positive distance.

\emph{Step 4: (b).} The supports of $\psi_0\circ\theta_\mu$ and $\psi_0$ lie in $\mathbb B(-e_\mu,\frac12)$ and $\mathbb B(e_\mu,\frac12)$, at a distance in $[1,3]$; so $(\psi_0\circ\theta_\mu,\psi_0)\in\mathfrak T^1(\frac12;\vartheta_1)$ for some $\vartheta_1\in(0,1)$, and its translate by $\tilde p$ dilated by $r$ lies in $\mathfrak T^1(\frac r2;\vartheta_1)$. Clause~(c) of Theorem~\ref{C-351-4}, whose bracket is at most $C(r/2)^{1/2}$, gives the expansion, and $\mathcal K_5(\psi_0\circ\theta_\mu,\psi_0)>0$ by~\eqref{C-351-2}.

\emph{Step 5: (c).} At a finite cutoff $\beta=L^Kb$, so that $\partial_b=L^K\partial_\beta$, and $\partial_\beta$ of a cumulant of observables under $e^{-\beta A}$ is minus the cumulant with the action $A=\sum_pa_p=a_K\mathcal O_K(1)$ added; hence $\partial_bS^T_{2;K,\ell,L^Kb}(f,h)=-S^T_{3;K,\ell,L^Kb}(1,f,h)$, which is~the finite-cutoff differentiation identity in the proof of Proposition~\YMMref{118.3}{814} for the first derivative. By Proposition~\YMMref{118.3}{814}, equation~\YMMref{(2217)}{814}, applied to the pair $(f_r,h_r)$, whose supports are disjoint, $b\mapsto S_2^{(\ell,b)}(f_r,h_r)$ is infinitely differentiable on $(0,\infty)$ and $\partial_bS_2^{(\ell,b)}(f_r,h_r)=-S_3^{(\ell,b)}(1,f_r,h_r)$, the derivatives at the finite cutoffs converging to it uniformly on $[b_1,b_2]$; so the integrand is $-\partial_b[b^2S_2^{(\ell,b)}(f_r,h_r)]$, and the integral equals
\[
b_1^2\bigl[S_2^{(\ell,b_1)}(f_r,h_r)-c_2b_1^{-2}\mathcal K_6(f_0,h_0)\bigr]-b_2^2\bigl[S_2^{(\ell,b_2)}(f_r,h_r)-c_2b_2^{-2}\mathcal K_6(f_0,h_0)\bigr],
\]
each bracket being at most $Cr$ by~(a). \qed

\subsection{Scope of the conclusions}

The theorems concern short distances at a fixed coupling; on $\mathbb R^3$ they hold for $0<g\le\gamma_W$. Not being dilation covariant is not being interacting: three free photon fields together with an independent generalised free field of two-point kernel $\varkappa b^{-3}|x-y|^{-5}$, whose spectral density $\varkappa b^{-3}\frac{2\pi}3m^3$ is positive, obey (L) with $c=\frac8{3\pi}$ and remainder $0$ and have $\mathsf N=\mathsf X$, and the Gaussian model of Proposition~\ref{C-351-24} has the same two-term law on a torus. The numbers $\gamma_{\rm e}$, $d_{\rm e}$, $d^L$, $\delta_{\rm c}$ and the tolerances are given by existence.

\section{Compact gauge groups}\label{C-368-1}

For a general compact gauge group, the lattice geometry of the construction is unchanged, but the colour algebra
and the matrix chart are not. At the trivial field every linear operator is a scalar lattice operator tensored with
the colour space. This observation determines the leading two-point coefficient. At one loop the dependence on the
representation is carried by two Casimir numbers. Away from the trivial field, the construction also requires
uniform estimates for propagators, critical configurations and localised Gaussian integrals.

We first state the local calculations and a conditional extension of the multiscale construction. The hypotheses
for that extension are given explicitly below. They separate the local group calculations from the estimates at
general regular backgrounds. The finite-box calculation of the next perturbative coefficient is independent of this
extension. For products acting on separate blocks, the measures factor, and the corresponding conclusions follow
directly from the $SU(2)$ construction. For quotients, the preliminary-gauge
and quantitative local assertions of Proposition~\ref{C-368-27}(iii),(iv)
retain their separately stated analytic hypotheses.

\subsection{Notation and hypotheses}

\begin{definition}[The pair]\label{C-368-2}
Let $N\ge2$ and let $G\subset U(N)$ be nontrivial, compact and connected; the \emph{pair} is $G$ with its representation on
$\mathbb C^N$. Let $\mathfrak g=\{Y=Y^*:e^{itY}\in G\text{ for all }t\in\mathbb R\}$, $n=\dim\mathfrak g$,
$\mathfrak g^c=\mathfrak g+i\mathfrak g$, and $G^c$ the group generated by $\exp\mathfrak g^c$; $\langle
X,Y\rangle=\frac1N\mathrm{tr}\,XY$ with the norm $\|X\|=(\frac1N\mathrm{tr}\,X^*X)^{1/2}$, and $|\cdot|$ the
operator norm; $P$ the projection of $\mathbb C^{N\times N}$ onto $\mathfrak g^c$ along
$\{M:\mathrm{tr}\,MZ=0\text{ for all }Z\in\mathfrak g^c\}$; $(t_a)_{a\le n}$ an orthonormal basis of $\mathfrak g$
and $\Omega=\sum_at_a^2$; $\kappa_G=\max_{0\ne Y\in\mathfrak g}|Y|/\|Y\|$ and $\tau_G=\sup_{\mathbf 1\ne V\in
G}|V-\mathbf 1|/\|V-\mathbf 1\|$;

$G^\pi=\exp(i\{Y\in\mathfrak g:|Y|<\pi\})$; $r_G$ the largest $r$ with $\{u\in G:|u-\mathbf 1|<r\}\subset G^\pi$,
and $\rho_G=2\arcsin\frac{r_G}2\in(0,\pi]$, the largest number such that every $u\in G$ whose eigen-angles are
below $\rho_G$ in modulus lies in $G^\pi$; $z_G=\min\{1,\mathrm{dist}(\mathbf 1,Z(G)\setminus\{\mathbf 1\})\}$, and
$z_G=1$ if $Z(G)=\{\mathbf 1\}$.

With $\log$ the power series on $|X-\mathbf 1|<1$, let $D^c=\{X\in G^c:|X-\mathbf 1|<1,\ \log X\in\mathfrak g^c\}$
and $D=D^c\cap G$; further $X^c=(X^*)^{-1}$, $m(X)=\max\{|X-\mathbf 1|,|X^{-1}-\mathbf 1|\}$,
$c(X)=\frac12(X+X^{-1})$ and $H_X=\frac1{2i}(X-X^{-1})$; $|P|$ is the norm of $P$ for $|\cdot|$, and
$\varsigma_G=\sup\{|PH_V|/|V-\mathbf 1|:\mathbf 1\ne V\in G\}$, $\varsigma^D_G=\sup\{|PH_V|/|V-\mathbf 1|:\mathbf
1\ne V\in D\}$, $\kappa'=\kappa'(G)=4\sup\{|(1-P)H_V|/|V-\mathbf 1|^3:\mathbf 1\ne V\in D\}$, which
Lemma~\ref{C-368-8}(ii) bounds, and $j_G=\max\{1+\kappa'z_G^2,3\varsigma^D_G\}$ are numbers of the pair, $(1-P)H_V$
being of third order in $V-\mathbf 1$.

The action density is $a_p=a(\partial U(p))$ with $a(X)=1-\frac1{2N}\mathrm{tr}(X+X^{-1})$, which is
$1-\frac1N\mathrm{Re}\,\mathrm{tr}\,X$ on $G$; $F(p)=\xi^{-2}PH_{\partial U(p)}$ and $J=D^{\mathrm T}F$, with
$\partial U$, $D$ and $\xi$ of Definition~\YMMref{9.1}{37}; the \emph{Casimir density} is
$k_p=n-\frac1N\mathrm{Re}\,\mathrm{tr}(\Omega\,\partial U(p))$.

For simple $\mathfrak g$: $C_A=\mathrm{tr}(\mathrm{ad}\,H)^2$, the trace taken on $\mathfrak g^c$, and
$C_\rho=\frac1N\mathrm{tr}(\Omega H^2)$ for $H\in\mathfrak g$ with $\|H\|=1$, numbers that do not depend on $H$
(Lemma~\ref{C-368-11}); $\varphi_G=\sum_a\mathrm{tr}(\mathrm{ad}\,t_a)^2=nC_A$;
$\gamma_G=\frac12C_\rho-\frac3{16}C_A$; $c_2(G)=\frac{3n}{16\pi^2}$ and $c_1(G)=\frac{\varphi_G}{3\pi n}$.

For semisimple $\mathfrak g=\bigoplus_{i\le r}\mathfrak g_i$ with simple ideals $\mathfrak g_i$: $n_i=\dim\mathfrak
g_i$, $P_i$ the orthogonal projection onto $\mathfrak g_i$, $\kappa_i=\max_{0\ne Y\in\mathfrak g_i}|Y|/\|Y\|$, and
$C^{(i)}_A$, $\theta_i$ are formed as $C_A$, $C_\rho$ with $H\in\mathfrak g_i$.
\end{definition}

If $\mathfrak g$ is semisimple, then $G\subset SU(N)$: $\det$ restricted to $G$ is a character, and its
differential $\mathrm{tr}$ vanishes on $[\mathfrak g,\mathfrak g]=\mathfrak g$. For every $G$, $|P|\le N^{1/2}$:
$\mathfrak g^c$ being closed under adjoints, $P$ is the orthogonal projection for $\mathrm{tr}\,X^*Y$, so that
$|PZ|\le(\mathrm{tr}\,Z^*Z)^{1/2}\le N^{1/2}|Z|$. For $SU(N)$: $n=N^2-1$, $\kappa_G^2=\tau_G^2=N-1$,
$r_G=2\sin\frac\pi N$, $\rho_G=\frac{2\pi}N$, $z_G=\min\{1,r_G\}$, which is $1$ exactly for $N\le6$,
$\Omega=n\mathbf 1$, $|P|\le2$, $\varsigma_G\le2-\frac2N$ (Remark~\ref{C-368-30}), $\kappa'\le2-\pi/\sqrt3$
(Lemma~\ref{C-368-8}(ii)), $\varsigma^D_G=1$ and $j_G=3$ (on $D$, where $|V-\mathbf 1|<1$,
$|PH_V|\le|H_V|+|(1-P)H_V|\le|V-\mathbf 1|-\frac18|V-\mathbf 1|^3+\frac14\kappa'|V-\mathbf 1|^3<|V-\mathbf 1|$, and
the quotient tends to $1$ at $\mathbf 1$), $(C_A,C_\rho)=(2N^2,N^2-1)$, $\gamma_G=\frac{N^2-4}8$ and
$c_1(G)=\frac{2N^2}{3\pi}$; at $SU(2)$ this is $(C_A,C_\rho)=(8,3)$, $\gamma_G=0$, $c_1=\frac8{3\pi}$, and at
$SU(3)$ it is $(n,C_A,C_\rho,\gamma_G)=(8,18,8,\frac58)$, $r_G=\sqrt3$, $z_G=1$, $\kappa_G^2=2$. The blocking
factor $L$, the cutoff $K$, the levels, $\eta=L^{-k}$ and $a_K=L^{-K}$ are those of Section~\YMMref{3}{13};
$B_\sharp$ bounds the increments $L^{-1}x_{k-1}-x_k$ of the run, $\bar B=B_\sharp\frac L{L-1}$, $\mathsf G_\sharp$
bounds the coefficient, and $x_k\ge\gamma^{-2}$, as in Theorem~\YMMref{D}{8}; the reference coupling and first-passage notation are those of
Proposition~\YMMref{74.1}{615}: $0<g\le\gamma$,
$X_g=L(g^{-2}+B_\sharp)$ and
$I(g)=(g^{-2}+\bar B,X_g-\bar B]$, with the admitted sides
determined by the geometric room condition there. The short-distance
ceiling $\gamma_W\le\gamma$ and the scale and cutoff restrictions
are those of Theorem~\YMMref{AV}{851}, transferred under
Definition~\ref{C-368-3} and the construction hypotheses below.
Set $\mathcal K_q(f,h)=\iint f(x)h(y)|x-y|^{-q}\,dx\,dy$.

The averaging radius of the $SU(2)$ construction is $\omega_M=\frac1{200L^2}$. A \emph{number of $L$ and the pair}
is a number that depends on $L$ and on $G$ with its representation, and on nothing else.

\begin{definition}[Regular classes at the local scale]
\label{def:group-regular-classes}
For a separated sequence $\Omega_0\supset\cdots\supset\Omega_k$,
put $\eta=L^{-k}$. The local index of a bare bond or plaquette
$\mathfrak e$ meeting $\Omega_0$ is the largest $j\le k$ for which
one of its corners lies in $\Omega_j$, and
$\mathsf L(\mathfrak e)=L^{j-k}$. A plaquette is inner if all
its corners lie in $\Omega_0$; a bond is full if its four
incident plaquettes are inner. These are the scales of
Definition~\YMMref{9.3}{40}.

Both $\mathcal U'(r)=\mathfrak U'_k(\mathbf B,r)$ and
$\mathcal U(r)=\mathfrak U_k(\mathbf B,r)$ consist of real
$G$-valued configurations on the inner bare bonds with
\[
 |\partial U(p)-\mathbf1|<r\eta^2\mathsf L(p)^{-2}
 \qquad(p\text{ inner}).
\]
The additional condition, at every full bond $b$, is respectively
\[
 |D_U^{\mathrm T}\partial U(b)|<r\eta^2\mathsf L(b)^{-3},
 \qquad |J(U)(b)|<r\mathsf L(b)^{-3}.
\]
Here the divergence and current use bare spacing $\eta$, and
$J$ is the projected current of Definition~\ref{C-368-2}.
The definitions themselves impose no inclusion between the two classes.
\end{definition}

\begin{lemma}[Transpose bound and depth zero]
\label{lem:group-transpose-depth-zero}
For a complex linear map $T$ of $\mathfrak g^c$, let $T^{\mathrm T}$
be its transpose for $\langle X,Y\rangle=N^{-1}\mathrm{tr}(XY)$.
With operator norms induced by $|\cdot|$,
\[
 |T^{\mathrm T}|\le N|T|.
\]
At depth zero the determining set consists of all inner bare bonds,
and its constraint fibre over a real datum $V$ is $\{V\}$.
If each inner plaquette of $V$ has curvature less than
$\varepsilon_1$, then
\[
 |J(V)(b)|<4\varsigma_G\varepsilon_1
 \quad(b\text{ full}),\qquad
 V\in\mathcal U(4\varsigma_G\varepsilon_1).
\]
In particular, if $B_3\ge4\varsigma_G$ and
$\varepsilon_0\ge B_3\varepsilon_1$, this singleton is the
set of critical configurations in the regularity class of radius
$\varepsilon_0$.
\end{lemma}

\begin{proof}
In the real orthonormal basis $(t_a)$, the trace pairing is the
ordinary bilinear coordinate pairing and $\|X\|^2=\sum_a|x_a|^2$.
Singular values give $\|X\|\le|X|\le\sqrt N\|X\|$.
The Hilbert operator norm is unchanged by transposition, since
the transpose is the complex conjugate of the adjoint. Applying
the two norm comparisons, first to the output and then to the
input, gives
\[
 |T^{\mathrm T}X|\le\sqrt N\|T^{\mathrm T}\|_{\mathrm{op}}\|X\|
 =\sqrt N\|T\|_{\mathrm{op}}\|X\|
 \le N|T|\,|X|.
\]
At depth zero there are no inward bonds and the averaging map
is the identity. Thus the fibre is a singleton and every curve
in it is constant. In bare units the current at a full bond is
a sum of four transported terms $\pm PH_{\partial V(p)}$.
Each has norm less than $\varsigma_G\varepsilon_1$; unitary
transport preserves this norm. This proves the current bound.
Since $G$ has positive dimension, the ratio defining $\varsigma_G$
tends to one along $e^{itY}$ for any nonzero $Y\in\mathfrak g$;
hence $\varsigma_G\ge1$, and the curvature bound also has the
claimed radius. The final assertion follows by increasing that
radius to $\varepsilon_0$.
\end{proof}

\begin{lemma}[The averaging row constant]
\label{lem:group-averaging-row-constant}
Suppose that the linearized average at a positive-level row $c$
has the counting kernel
\[
 (\hat QA)(c)=\sum_{b\in\mathsf P_c}k(c,b)\xi_cA(b),\qquad
 |k(c,b)|\le C_{\rm av}(\xi/\xi_c)^3,
\]
where $\xi$ is the bare spacing and $C_{\rm av}\ge1$. Use the actual double blocks and
neighbourhoods of the averaging map: each $\mathsf P_c$ has at most
$6(\xi_c/\xi)^3$ positive bonds; each bare bond lies in at most
eighteen such blocks; and the scale ratios used in either row or
column sum are at most $L$. Then the two linear averaging bounds
in equation~\YMMref{(1045)}{466} hold with
\[
 C_Q=18NL^3 C_{\rm av}.
\]
At level zero the average is the identity. The choice
$C_{\rm av}=3$ therefore gives $C_Q=54NL^3$, including
$108L^3$ at $SU(2)$. For a general pair this numerical choice
requires the displayed kernel estimate with $C_{\rm av}=3$
on the admitted averaging domain; it is a quantitative averaging
input in the transferred construction.
\end{lemma}

\begin{proof}
For a row over $y$, the row sum is bounded by
$6LC_{\rm av}\max_{N_L(y)}|A|$, with the dimensionless field norms
of equation~\YMMref{(1045)}{466}. The transpose uses the bare weight
$\xi^3$ and the row weight $\xi_c^{-1}$. The latter cancels the
factor $\xi_c$ in the counting kernel, giving
\[
 (\hat Q^{\mathrm T}\varphi)(b)
      =\xi^{-3}\sum_{c:b\in\mathsf P_c}k(c,b)^{\mathrm T}\varphi(c).
\]
If $b$ is over $y$, Lemma~\ref{lem:group-transpose-depth-zero}
and the eighteen-block count give
\[
 \xi_y^3|\hat Q^{\mathrm T}\varphi(b)|
 \le\sum_{c:b\in\mathsf P_c}
       NC_{\rm av}(\xi_y/\xi_c)^3|\varphi(c)|
 \le18NL^3C_{\rm av}\max_{N_L(y)}|\varphi|.
\]
The same constant dominates the row bound and the identity rows.
This is the row/column count following
\YMMref{(1045)}{466}, with the transpose factor $N$ retained.
\end{proof}

\begin{definition}[Transfer to $G$]\label{C-368-3}
The transfer to $G$ of a definition or a statement of the $SU(2)$ construction is made by kind.
\begin{enumerate}
\item[(1)] What stands for one object is replaced alike at all places: $N\times N$ matrices, $\frac1N\mathrm{tr}$,
$G$, $G^c$, $\mathfrak g$, $\mathfrak g^c$ for $2\times2$ matrices, $\frac12\mathrm{tr}$, $SU(2)$, $SL(2,\mathbb
C)$ and their algebras; colour indices run to $n$; $a_p$, $F$ and $J$ are those of Definition~\ref{C-368-2}, and
$a(X)$ stands for every trace of a group-valued argument at complex fields, the gauge-fixing term included;
$P[\frac12\{c(X),\theta\}]$ stands where $SU(2)$ has $\frac12(\mathrm{tr}\,X)\theta$; $|\cdot|$ stands in suprema,
radii, series, domains and tests, $\|\cdot\|$ and $\langle\cdot,\cdot\rangle$ in Hilbert sums, pairings and
quadratic forms; the cut-off of the bond shift (the cutoff preceding Lemma~\YMMref{53.3}{444}), $\vartheta_{\rm co}(|Y|/\delta_k)$, is
$\vartheta_{\rm co}(\frac43+\lambda(\nu(Y)/\delta_k-\frac43))$, with $\nu(Y)=(\frac1N\mathrm{tr}\,Y^{2p})^{1/2p}$,
$p$ an integer with $w=N^{1/2p}<\frac54$, and $\lambda=\frac w{5-4w}$, and $\lambda\|\vartheta'_{\rm co}\|_\infty$
stands for $\|\vartheta'_{\rm co}\|_\infty$ in every condition and bound; as $\nu\le|\cdot|\le w\nu$, it is smooth
and invariant, $0$ on $|Y|\le\frac43\delta_k$ and $1$ on $|Y|\ge\frac53\delta_k$ (Lemma~\ref{C-368-10}(ii)); and
``a constant depending only on $L$ and the pair'' stands for ``a number of $L$''.
\item[(2)] A numeral that counts colours, compares $|\cdot|$ with $\|\cdot\|$, comes from a transpose or comes from
the projection $P$ takes its value for $G$: $n$, $\frac{3n}2$, $3n$ for $3$, $\frac92$, $9$; the mean over $n$ colours
replaces the mean over three colours in the quadratic coefficient of a local term; $N$ replaces the $2$ of $|T^{\mathrm
T}|\le2|T|$ (Lemma~\ref{lem:group-transpose-depth-zero}); $\sqrt N$ for $\sqrt2$; a factor $N^{1/2}$ where an inverse bounded in
$\ell^2$ is read in $\sup|\cdot|$; $4\varsigma_G$ for the $4$ at depth $0$ (Lemma~\ref{lem:group-transpose-depth-zero}), which is at
most $8(1-\frac1N)$ for $SU(N)$; $3|P|$ for the $3$ of the bound $3\alpha_0$ of the current at an inner bond that
is not full (Proposition~\YMMref{30.1}{236}), which is at most $6$ for $SU(N)$, and, under $\alpha_0\le z_G$, $j_G$ for
the $3$ of that lemma, which is $3$ for $SU(N)$ (Definition~\ref{C-368-2}, Remark~\ref{C-368-30},
Lemma~\ref{C-368-8}(ii)); $6\hat\varsigma_G$ for the $12$ of the bound $\xi^3|J_W(U)(b)|\le12a$ of the $SU(2)$
construction on the current $J_W$ of a region $W$ at a bond $b$, at $G^c$-valued fields whose plaquette variables
at $b$ lie within $a$ of $\mathbf 1$ and whose bond variables there have condition number $|u||u^{-1}|\le2$,
$\hat\varsigma_G=\sup|PH_g|/|g-\mathbf 1|$ over the plaquette variables $g\ne\mathbf 1$ of such fields, which is at
most $2$ at $SU(2)$ and at most $\frac{17}2|P|$ for semisimple $\mathfrak g$ (Remark~\ref{C-368-30}); the factor
$|P|$ in the kernel and Weitzenb\"ock bounds of the Hessian in $|\cdot|$ (Lemma~\ref{C-368-14}); for the averaging row constant of Lemma~\ref{lem:group-averaging-row-constant},
$C_Q=18NL^3C_{\rm av}$, which is $54NL^3$ when the normalized kernel bound
$C_{\rm av}=3$ is supplied, and $108L^3$ at $SU(2)$; the constants
formed from $6LC$ and $12LC$ move with this value; $z<(\frac54)^nL^{3n}$ for the bound $z<2L^9$ of Lemma~\YMMref{65.3}{581}; and an inequality
that absorbs a colour factor, such as $\frac34\le1$ in the determinant counting argument of Lemma~\ref{lem:colour-determinant-count}, must hold anew with
the new numbers.
\item[(3)] The following statements of the $SU(2)$ construction are false for $G$ as worded, and their transfers
are the forms named, each with the statement of this section that bears it. The inclusion $\mathcal
U'(r)\subset\mathcal U(r)$ of the two regular classes $\mathcal U'(r)=\mathfrak U'_k(\mathbf B,r)$ and $\mathcal
U(r)=\mathfrak U_k(\mathbf B,r)$ of Definition~\ref{def:group-regular-classes} (the $SU(2)$ inclusions are Lemma~\YMMref{9.4}{40} and equation~\YMMref{(465)}{237}), the first
bounding $D^{\mathrm T}_U\partial U$ and the second the current, becomes $\mathcal U'(r)\subset\mathcal
U(\kappa_Gr)$, and for $r\le z_G$ the sharper $\mathcal U'(r)\subset\mathcal U((1+\kappa'r^2)r)$, the form that
Lemma~\ref{C-368-8}(ii) states for $SU(N)$ (Remark~\ref{C-368-30}). The second inclusion of the two classes in the
$SU(2)$ construction, $\mathcal U(r)\subset\mathcal U'(r+r^2)$, becomes $\mathcal U(r)\subset\mathcal
U'(r+r^2+\kappa'r^3)$ for every $G$ and $r\le z_G$, which at $SU(2)$, where $\kappa'=0$ and $z_G=1$, is the
inclusion of the $SU(2)$ construction for $r\le1$: the plaquette variables of a field of $\mathcal U(r)$ or of
$\mathcal U'(r)$ lie within $z_G$ of $\mathbf 1$, so in $D$ (Lemma~\ref{C-368-6}(iii)), where
$|(1-P)H_V|\le\frac14\kappa'|V-\mathbf 1|^3$ (Definition~\ref{C-368-2}), and the two bounds $\eta|(D^{\mathrm
T}\partial U)(b)|\le\rho^2+\eta^3|J(b)|+\kappa'\rho^3$ and $\eta^3|J(b)|\le\eta|(D^{\mathrm T}\partial
U)(b)|+\kappa'\rho^3$ at a bond $b$ whose plaquette variables lie in $D$ within $\rho\le z_G$ of $\mathbf 1$, which
Lemma~\ref{C-368-8}(ii) states for every $G$, give the two inclusions at each full bond, as in the proof of that
lemma for $SU(N)$.

The $SU(2)$ null-set exception $G\setminus G^\pi=\{-\mathbf1\}$ becomes the indicator of $G^\pi$ in the integrand
(Lemma~\ref{C-368-10}(i)); and where the $SU(2)$ construction uses $|\log j|<1$ on the ball of radius $\frac38$,
the bound $|\log j_G(Y)|\le0.173\,n|Y|^2$ stands, its constant entering a maximum (Lemma~\ref{C-368-10}(i)).
``$|X-\mathbf 1|<1$'' becomes ``$X\in D^c$'' wherever a series logarithm is to lie in $\mathfrak g^c$
(Lemma~\YMMref{5.1}{23}; Lemma~\ref{C-368-6}(iii)). The mean of a family of diameter below $\frac1{40}$
(Proposition~\YMMref{4.3}{18}) takes the proviso that the quotients lie in $D^c$, which holds if the diameter is
below $z_G$ (Lemma~\ref{C-368-6}(iv)). The second assertion of Lemma~\YMMref{14.2}{62} takes the proviso that
every eigen-angle of each $W_x$ is below $\rho_G$ in modulus, or that $W_x$ is a product of exponentials
$e^{i\gamma_j}$, $\gamma_j\in\mathfrak g^c$, $\sum_j|\gamma_j|<\ln2$ (Lemma~\ref{C-368-7}(ii)). ``$u\ne-\mathbf
1$'' becomes ``$u\in G^\pi$'' (Lemmas~\ref{C-368-6}(ii), \ref{C-368-7}(ii)). $|X-\mathbf 1|$ becomes $m(X)$ in
conditions on complex fields that are to be kept by $X\mapsto X^c$ and wherever $|X^{-1}-\mathbf 1|=|X-\mathbf
1|=|X^*-\mathbf 1|$ was used; ``the conjugate point lies in the domain if and only if the point does''
(Lemma~\ref{lem:two-by-two-conjugation}) and ``the function $z_c$ of the elimination factor is formed at $(U^*)^{-1}$ if and only
if at $U$'' take the two-sided guard $m(X)$; and the symmetric part $Q^s$ of a complex domain is replaced by its
real-connected part $Q^r$, or moved into one, the conditions of Lemma~\ref{C-368-15}(iii) on the radii standing
then among the conditions on the constants of Definition~\ref{C-368-31}(K) (Lemmas~\ref{C-368-24}(ii),
\ref{C-368-15}(iii), \ref{C-368-10}(iv)).

The averaging radius $\omega_M$ (equation~\YMMref{(30)}{19}) becomes $\omega^G=\min\{\omega_M,z_G\}$, or the
floor $200L^2z_G\ge1$ is laid once (Lemma~\ref{C-368-7}(i)).

The radius $\frac12$ that is to keep $\det T$ off $(-\infty,0]$ (the three-dimensional determinant argument in the proof of Proposition~\YMMref{50.2}{386}) becomes
$r_n=\sin\frac\pi{2n}$, with the numbers of Lemma~\ref{C-368-10}(iv) in the chains in the proof of Proposition~\YMMref{50.2}{386} and in the elimination-factor estimates of Lemma~\ref{C-368-10}(iii), (iv). The full constant $C_{\rm mob}$ in the local determinant bound takes the factor $\max\{1,n2^{-11}\}$, as shown in Lemma~\ref{lem:colour-determinant-count}; its trace-dependent part is evaluated with $n$ colours.

The comb gauge (Lemma~\YMMref{14.4}{64}) takes the guard $\varpi<z_G$ (Lemma~\ref{C-368-8}(iii)), and
Lemma~\ref{lem:group-contour-gauge} the guard $\varpi_y(W)<z_G$, under which its contour quotients lie in $D^c$ and their means
in $G$ (Lemma~\ref{C-368-6}(iii), (iv)). The quantitative cube-gauge input (Lemma~\ref{C-368-8}(iii)) takes constants $c$, $B$ of
the pair (Lemma~\ref{C-368-8}(iii)). The orthogonal field maps in the covariance argument following Proposition~\YMMref{47.4}{365} are required here to keep
$|\cdot|$, as $\mathrm{Ad}(G)$ does (Proposition~\ref{C-368-17}). The membership in $\mathcal U'(2s)$,
$s=2\alpha_0+3\alpha_2$, of the curvature/current conversion in Lemma~\ref{C-368-8}(ii), with both of its quantitative hypotheses, takes the guard $\alpha_0+3\alpha_2<z_G$, under which the
plaquette variables there, within $\alpha_0+3\alpha_2$ of $\mathbf 1$, lie in $D$ (Lemma~\ref{C-368-6}(iii)); the
bound of Lemma~\ref{C-368-8}(ii) at a full bond $b$, with $\rho=(\alpha_0+3\alpha_2)\eta^2\mathsf L(b)^{-2}$ and
the bound $s\mathsf L(b)^{-3}$ of the current there, then gives for every $G$ the membership in $\mathcal
U'(s+s^2+\kappa's^3)$, which lies in $\mathcal U'(2s)$ as soon as $s(1+\kappa's)\le1$, so at $SU(N)$ for
$s\le\frac13$, $\kappa'$ being at most $2-\pi/\sqrt3$ there.

The bound $|\mathrm{Bl}(\mathsf g,\mathsf f)|\le\|\mathsf g\|_2\|\mathsf f\|_2$ that the $SU(2)$ construction draws
by polarisation from $|\mathrm{Qf}(\mathsf a)|\le\|\mathsf a\|_2^2$, $\mathrm{Qf}$ a quadratic form on $\mathfrak
g^c$-valued functions $\mathsf a$ on bonds, $\mathrm{Bl}$ its polar form, $\|\mathsf a\|_2^2=\sum_b|\mathsf
a(b)|^2$, at real $\mathsf g$, $\mathsf f$, takes the factor $\kappa_G^2$, and keeps the factor $1$ where
$|X+iY|^2\le|X|^2+|Y|^2$ for all $X,Y\in\mathfrak g$ (Remark~\ref{C-368-30}, Lemma~\ref{C-368-6}(i)).

For $G\subset SU(N)$ the energy of block averages takes $\frac1{4\tau_G^2}$ and the threshold
$e_*=e_L/(2\tau_G^2)$, or $N-1$ in the place of $\tau_G^2$, and the lower bounds of the action on large-field
regions take the factor $\tau_G^{-2}$ (Lemma~\ref{C-368-9}).

The assembly of the bond maps (Lemma~\YMMref{53.3}{444}) takes $r_G$ for $2$, with the condition $\mathsf c_Lc_4\le2Lr_G$ (Proposition~\ref{C-368-16}(iv)).

The regularity radius $\varepsilon_{\rm v}$ of the critical configurations (Definition~\YMMref{9.5}{41}) takes
the ceilings $\varepsilon_{\rm v}\le a_G$, with a fixed admitted positive orbit radius $a_G$ from Definition~\ref{C-368-12},
$\varepsilon_{\rm v}\le\varepsilon_G$, $\varepsilon_G$ the number of the uniqueness clause of
Definition~\ref{C-368-31}(V), and $\varepsilon_{\rm v}\le L^2r_G$ (Proposition~\ref{C-368-18}(f)).

The $\pi$ of the ceiling $\alpha_1<\pi L^j$, $1\le j\le k$, under which the local terms are continuous on the
analyticity cores (Lemma~\ref{lem:full-core-continuity}), which a proof of the chain from level $k$ to level $k+1$ quotes,
becomes $\rho_G$, the bound on the eigen-angles under which an element of $G$ lies in $G^\pi$
(Definition~\ref{C-368-2}); the ceiling reads $\alpha_1<\rho_GL^j$ and is met by the row $\alpha_1<\rho_GL$ of the
table (Proposition~\ref{C-368-17}).
\end{enumerate}
These rules define the transferred statement. They do not, by themselves, prove it. When no correction in (3) applies, the convention uses (1) and (2) alone.
\end{definition}

\begin{remark}[Where the literal transfer fails]\label{C-368-30}
A rule that replaced the signs alike and kept every numeral would not transfer the proofs, as the following
instances show. Each names the statement of the $SU(2)$ construction that it concerns, the least $N$ at which the
literal transfer fails for $SU(N)$, or the least $n$, and the statement of this section that proves the form that
holds.

\emph{Numerals.} The $2$ of $|T^{\mathrm T}|\le2|T|$ (Lemma~\ref{lem:group-transpose-depth-zero}) is $N$, sharply at $SU(N)$,
$\kappa^2=N-1$ being attained by general kernels (Proposition~\ref{C-368-18}(g)). The $4$ of
Lemma~\ref{lem:group-transpose-depth-zero} at $SU(2)$, without its general-group factor $\varsigma_G$, fails from $N=6$: let $\Omega_0$ hold a bare bond $b$ with its four
plaquettes, and let the datum be $X_N=\mathrm{diag}(e^{ia},e^{-ia},\dots,e^{-ia})\in SU(N)$, $a=\frac{2\pi}{N-2}$,
at $b$ and $\mathbf 1$ at the other bonds; with $\varpi=|X_N-\mathbf 1|=2\sin\frac\pi{N-2}$ its plaquette variables
are within $\varpi$ of $\mathbf 1$, while its current at $b$ is $4PH_{X_N}$ up to sign, of norm $4r_N\varpi$,
$r_N=(2-\frac2N)\cos\frac\pi{N-2}\ge1.178$ for $N\ge6$; so for $\varepsilon_1\in(\varpi,r_N\varpi]$ the datum is
$\varepsilon_1$-regular and does not lie in $\mathfrak U(\mathbf B;4\varepsilon_1)$. The current at $b$ is a sum of
four transported terms $\pm PH_{\partial V(p)}$, so the lemma holds for every $G$ with $4\varsigma_G$ for $4$. For
$SU(N)$ this is at most $8(1-\frac1N)$: for $V\in SU(N)$ with eigen-angles $\theta_j$ the eigenvalues of $PH_V$ are
$\sin\theta_j-\frac1N\sum_l\sin\theta_l$, so that $|PH_V|\le(2-\frac2N)\max_j|\sin\theta_j|\le(2-\frac2N)|V-\mathbf
1|$ (Proposition~\ref{C-368-18}(g)). The value $8(1-\frac1N)$ is that of $SU(N)$. For the spin representation of
$Spin(2m+1)$ on $\mathbb C^{2^m}$, simple, simply connected and with $z_G=1$, let $(h_a)_{a\le m}$ be the
orthonormal basis of a Cartan subalgebra that acts on a basis of weight vectors by the diagonal matrices of signs;
$P$ commutes with the adjoint action of the maximal torus, so it maps the diagonal matrices to the Cartan
subalgebra, and at $V=\exp(i\sum_ax_ah_a)$
\[PH_V=\sum_a\sin x_a\prod_{b\ne a}\cos x_b\;h_a,\qquad|PH_V|=\sum_a|\sin x_a|\prod_{b\ne a}|\cos x_b|,\]
every sign vector being a weight. At $x_a=x$, $\tan^2x=\frac1{m-1}$, this is $\sqrt m(1-\frac1m)^{(m-1)/2}$, which
exceeds $4$ from $m=43$; as $|V-\mathbf 1|\le2$, $\varsigma_G>2$ there. The inequality $\frac34\le1$ by which the
proof of Lemma~\ref{lem:colour-determinant-count} absorbs a colour factor reads $\frac n4\le1$, false from $n=5$, and
$(\frac54)^3<2$ behind the bound $2L^9$ of Lemma~\YMMref{65.3}{581} fails from $n=4$
(Proposition~\ref{C-368-18}(d)). The $12$ of the bound $\xi^3|J_W(U)(b)|\le12a$ on the current of a region at
$G^c$-valued fields whose bond variables have condition number at most $2$ (Definition~\ref{C-368-3}(2)) fails for
every $N\ge5$: on the thirteen bonds of the four plaquettes at $b$ a field with values in $SL(N,\mathbb R)$ of
condition number at most $2$, products of plane rotations and of diagonal matrices
$\mathrm{diag}(d_1,d_2,d_3,d_3,d_3,1,\dots,1)$ with rational entries, and $\mathbf 1$ at the other bonds, has its
four plaquette variables at $b$ within $a=\frac{103}{100}$ of $\mathbf 1$ and $\xi^3|J_W(U)(b)|>12.98>12a$. The
current at $b$ is $\xi^{-3}$ times a sum over the two directions $\nu$ of the plaquettes at $b$, each term being
$\pm PH_g$ at the plaquette variable $g$ of one plaquette of direction $\nu$ and $\mp u^{-1}PH_{g'}u$ at that of
the other, $u$ a bond variable, with $|u^{-1}Xu|\le|u||u^{-1}||X|\le2|X|$; so the bound holds for every semisimple
$\mathfrak g$ with $2(1+2)\hat\varsigma_G=6\hat\varsigma_G$ for $12$. At $SU(2)$, $g-g^{-1}=2(g-\mathbf
1)-\mathrm{tr}(g-\mathbf 1)\mathbf 1$ gives $|H_g|\le2|g-\mathbf 1|$, so that $\hat\varsigma_G\le2$ and
$6\hat\varsigma_G\le12$; for semisimple $\mathfrak g$, $G^c\subset SL(N,\mathbb C)$ (Lemma~\ref{C-368-6}(iii)), so
a bond variable $u$ of condition number at most $2$ has $|u|\le2$ and $|u^{-1}|\le2$, a plaquette variable $g$ has
$|g^{-1}|\le16$, and $|PH_g|\le\frac12|P|(|g-\mathbf 1|+|g^{-1}||g-\mathbf 1|)\le\frac{17}2|P||g-\mathbf 1|$.

\emph{One matrix and the logarithm.} The identity $X+X^{-1}=(\mathrm{tr}\,X)\mathbf 1$ of \YMMref{(70)}{38}, which
(1) transfers to $X+X^{-1}=\frac2N(\mathrm{tr}\,X)\mathbf 1$, holds for every $X\in G$ if and only if $X+X^{-1}$ is
scalar for every $X\in G$, that is, if and only if $G$ is conjugate to $\{u\otimes\mathbf 1_k:u\in SU(2)\}$ (for
$n\ge2$; Lemma~\ref{C-368-6}(i)); with it fail at $SU(3)$ the identity $a_p=\frac12|\partial U(p)-\mathbf 1|^2$,
the membership $H_X\in\mathfrak g^c$, $|Y|=\|Y\|$ on $\mathfrak g$ and $|X^{-1}-\mathbf 1|=|X-\mathbf 1|$ on $G^c$
(at $X=\mathrm{diag}(\sqrt2,\sqrt2,\frac12)$); $P$, the two norms, $m(X)$ and $G^\pi$ take their places
(Lemmas~\ref{C-368-6}(i), \ref{C-368-24}(ii)). ``$|X-\mathbf 1|<1$ gives $\log X\in\mathfrak g^c$''
(Lemma~\YMMref{5.1}{23}) fails for $N\ge7$, at $X=e^{2\pi i/N}\mathbf 1$; ``$X\in D^c$'' takes its place
(Lemma~\ref{C-368-6}(iii)). ``A family in $G$ of diameter less than $\frac1{40}$ has its mean in $G$''
(Proposition~\YMMref{4.3}{18}) fails for $N\ge252$: the mean of $\mathbf 1$ and $e^{2\pi i/N}\mathbf 1$ is $e^{\pi
i/N}\mathbf 1$, of determinant $-1$; the quotients in $D^c$ take its place (Lemma~\ref{C-368-6}(iv)).

The averaging radius $\omega_M$ of equation~\YMMref{(30)}{19} fails for $N\ge212372$ at $L=13$; $\omega^G$, or
the floor $200L^2z_G\ge1$, takes its place (Lemma~\ref{C-368-7}(i)). ``$|\mathbf 1-T|\le\frac12$ keeps $\det T$ off
$(-\infty,0]$'' (the three-dimensional determinant argument in the proof of Proposition~\YMMref{50.2}{386}) fails for $n\ge6$; the radius $r_n$ takes its place
(Lemma~\ref{C-368-10}(iv)).

\emph{Fields.} The inclusion $\mathcal U'(r)\subset\mathcal U(r)$ of the two regular classes
(Lemma~\YMMref{9.4}{40}) fails for $N\ge3$ at every small $r$ (Lemma~\ref{C-368-8}(ii)). For every $G$ it holds at
the radius $\kappa_Gr$: the current is $J(U)=\eta^{-2}PT$ with $T=\frac1{2i}(S-S^*)$ and $S=D^{\mathrm T}_U\partial
U$, the transports commuting with the adjoint, and $P$ is the orthogonal projection onto $\mathfrak g$ on hermitian
matrices, so that $|PT|\le\kappa_G\|PT\|\le\kappa_G|T|\le\kappa_G|S|$; and for $r\le z_G$ it holds at the radius
$(1+\kappa'r^2)r$, by the second bound at a bond of Lemma~\ref{C-368-8}(ii), the plaquette variables lying in $D$
(Definition~\ref{C-368-3}(3)). The bound $3\alpha_0$ of the current at an inner bond that is not full
(Proposition~\YMMref{30.1}{236}) fails for $N\ge8$ (Lemma~\ref{C-368-8}(ii)). It holds for every $G$ with
$3|P|\alpha_0$, as $|PH_V|\le|P||H_V|\le|P||V-\mathbf 1|$ at unitary $V$, which is at most $6\alpha_0$ at $SU(N)$;
and under $\alpha_0\le z_G$, where the plaquette variables lie in $D$, the bound of that lemma holds with
$j_G\alpha_0$: at a full bond by the inclusion at the radius $(1+\kappa'\alpha_0^2)\alpha_0$, and at an inner bond
$b$ that is not full, where $\mathsf L(b)^3$ times the current is a sum of at most three transported terms $\pm
PH_{\partial U(p)}$, with $3\varsigma^D_G\alpha_0$. The $3$ fails under $\alpha_0\le z_G$ at the pair of $SU(2)$ on
$\mathbb C^4\oplus(\mathbb C^2)^{100}$, spin $\frac32$ and a hundred copies of spin $\frac12$, which is simply
connected, with simple Lie algebra and $z_G=1$: with $J_3$ the image of $\frac12\sigma_3$, $\mathrm{tr}\,J_3^2=55$,
and $w=e^{itJ_3}$ one has $PH_w=c(t)J_3$, $c(t)=\mathrm{tr}(J_3\sin
tJ_3)/\mathrm{tr}\,J_3^2=\frac1{55}(3\sin\frac{3t}2+101\sin\frac t2)$, so that
$|PH_w|=\frac32c(t)=\frac32t-0.1034\,t^3+O(t^5)$ exceeds $|w-\mathbf
1|=2\sin\frac{3t}4=\frac32t-0.1406\,t^3+O(t^5)$ for small $t>0$; with $\Omega_0$ not the whole torus, the field
$U_0$ equal to $w$ at an inner bond $b$ in the outermost layer of $\Omega_0$ along a face, three of whose four
plaquettes are inner, and to $\mathbf 1$ at the other bonds lies in $\mathfrak U'_k(\mathbf B,\alpha_0)$ for
$\alpha_0>|w-\mathbf 1|$, and $\mathsf L(b)^3|J_{\Omega_0}(U_0)(b)|=3|PH_w|$, above $3\alpha_0$ for
$\alpha_0<|PH_w|$.

The comb gauge of a box of Lemma~\YMMref{14.4}{64} needs $\varpi<z_G$ for $N\ge7$, and the form of that clause with
some gauge of the box in the place of the comb, that $2s\varpi<2$ alone gives a gauge in which $W^h=e^{ia^0}$ with
$a^0$ in $\mathfrak g$ and $\sup|a^0|\le\pi s\varpi$, fails for $N\ge14$; the gauge on a cube of
Theorem~\YMMref{F}{199} has the quantitative general-group counterpart assumed in Lemma~\ref{C-368-8}(iii). That $G\setminus
G^\pi$ is a null set, as it is for $SU(2)$, fails for $N\ge3$; the indicator of $G^\pi$ takes its place
(Lemma~\ref{C-368-10}(i)). The cut-off $\vartheta_{\rm co}(|Y|/\delta_k)$ of the cutoff preceding Lemma~\YMMref{53.3}{444} is not
differentiable for rank at least $2$; the cut-off $\vartheta_{\rm co}(\frac43+\lambda(\nu(Y)/\delta_k-\frac43))$ of
Definition~\ref{C-368-3}(1) takes its place, with $\lambda\|\vartheta'_{\rm co}\|_\infty$ for $\|\vartheta'_{\rm
co}\|_\infty$; at $SU(3)$, $p=3$ and $\lambda=6.1194$ (Lemma~\ref{C-368-10}(ii)). The orthogonal colour maps in the covariance argument following Proposition~\YMMref{47.4}{365} keep the ball of $|\cdot|$ for $SU(2)$, and general rotations of $\mathfrak g$ do not
keep it for rank at least $2$; the rotations are to keep $|\cdot|$, as $\mathrm{Ad}(G)$ does, and covariance under
$\mathrm{Ad}(G)$ is what the proofs of the induction use (Proposition~\ref{C-368-17}).

The bound $|\mathrm{Bl}(\mathsf g,\mathsf f)|\le\|\mathsf g\|_2\|\mathsf f\|_2$ that the $SU(2)$ construction draws
by polarisation from $|\mathrm{Qf}(\mathsf a)|\le\|\mathsf a\|_2^2$ (Definition~\ref{C-368-3}(3)) fails at $SU(3)$,
of rank $2$: on two bonds $b_1$, $b_2$ let $x$, $y$, $z$ be the diagonal entries of $\mathsf a(b_1)$,
$\zeta_\pm=\frac12(\mathsf a(b_2)_{11}-\mathsf a(b_2)_{22}\pm i\mathsf a(b_2)_{12}\pm i\mathsf a(b_2)_{21})$,
$\psi_1=\frac12(e^{-i\pi/4}\zeta_++e^{i\pi/4}\zeta_-)$, $\psi_2=\frac1{2i}(e^{i\pi/4}\zeta_--e^{-i\pi/4}\zeta_+)$
and $\mathrm{Qf}(\mathsf a)=\frac12(x-z)\psi_1+\frac34\psi_1^2+\frac{3\sqrt6}4y\psi_2-\frac18\psi_2^2$; then
$|\mathrm{Qf}(\mathsf a)|\le|\mathsf a(b_1)|^2+|\mathsf a(b_2)|^2$ for all $\mathfrak{sl}(3,\mathbb C)$-valued
$\mathsf a$, while the real $\mathsf g=(\frac35\mathrm{diag}(1,-1,0),\frac45\sigma_3\oplus0)$ and $\mathsf
f=(\frac35\mathrm{diag}(0,1,-1),\frac45\sigma_1\oplus0)$ at $(b_1,b_2)$, $\sigma_1$ and $\sigma_3$ the Pauli
matrices, have $\|\mathsf g\|_2=\|\mathsf f\|_2=1$ and $\mathrm{Bl}(\mathsf g,\mathsf
f)=\frac1{25}(7+3\sqrt2+9\sqrt3)>1.07$. For every $G$ the bound holds with $\kappa_G^2$: as
$\|\cdot\|\le|\cdot|\le\kappa_G\|\cdot\|$ on $\mathfrak g$ (Lemma~\ref{C-368-6}(i)) and $\|\cdot\|$ obeys the
parallelogram law, $4|\mathrm{Bl}(\mathsf g,\mathsf f)|\le\|\mathsf g+\mathsf f\|_2^2+\|\mathsf g-\mathsf
f\|_2^2\le2\kappa_G^2(\|\mathsf g\|_2^2+\|\mathsf f\|_2^2)$, and $\mathsf g\mapsto t\mathsf g$, $\mathsf f\mapsto
t^{-1}\mathsf f$ gives $|\mathrm{Bl}(\mathsf g,\mathsf f)|\le\kappa_G^2\|\mathsf g\|_2\|\mathsf f\|_2$, which is
the bound of the $SU(2)$ construction at $SU(2)$; and where $|X+iY|^2\le|X|^2+|Y|^2$ for all $X,Y\in\mathfrak g$ it
holds with $1$, by $4i\,\mathrm{Bl}(\mathsf g,\mathsf f)=\mathrm{Qf}(\mathsf g+i\mathsf f)-\mathrm{Qf}(\mathsf
g-i\mathsf f)$ and $|X-iY|=|X+iY|$ for hermitian $X$, $Y$.

\emph{Complex fields.} The bare trace of the gauge-fixing term (the definition of $\Psi^{\rm G}$ in Section~\YMMref{51}{389}, p.~393) is not real at real
fields for $N\ge3$; $a(X)$ takes its place (Proposition~\ref{C-368-22}). ``The conjugate point lies in the domain
if and only if the point does'', which for $SU(2)$ rests on the isometry $\Theta$ of Lemma~\ref{lem:two-by-two-conjugation}, fails
for $N\ge3$; the two-sided guard $m(X)$ takes its place (Lemma~\ref{C-368-24}(ii), (iv)).

For $SU(N)$ with $N\le6$ the forms with $D^c$ coincide with the statements of the $SU(2)$ construction, and the
proviso on means of Lemma~\ref{C-368-6}(iv) is void for $N\le251$.
\end{remark}


\begin{definition}[Construction hypotheses]\label{C-368-31}
Let $G$ be simply connected with $\mathfrak g$ simple, $L\ge13$ odd and $200L^2z_G\ge1$; constants are constants
depending only on $L$ and the pair, uniform in the volume. We say that the pair satisfies the construction
hypotheses if the following five assertions hold. These hypotheses concern the complete multiscale construction;
they are not consequences of the matrix identities alone.
\begin{enumerate}
\item[(P)] \emph{Propagators and covariances.} Assume the uniform
linear estimates at regular backgrounds required below. Their $SU(2)$
forms are Theorems~\YMMref{G}{204}, \YMMref{I}{225}
and~\YMMref{L}{298}, and Proposition~\YMMref{31.1}{252}:
the constrained inverse and its local expansion, the real regional
inverse, the fluctuation covariance, and the current multiplier.
For $G$ these are assumptions on the natural field classes and on
every regional geometry used by the construction; the cited reduced-row
theorems alone do not assert their extension to other determining sets.
The quantitative local gauges in Lemma~\ref{C-368-8}(iii), and the
local current, difference, commutator, action-remainder and layer estimates
used to construct these inverses, are also required wherever used.
The required inverse bounds include: for a separated sequence of regions with $M$ above a floor and fields of the natural class with small
parameters, the operator $\mathsf K$ of the random-walk expansion is invertible, $\mathsf K^{-1}=\sum_wE_w$ over
finite sequences $w$ of pieces, local and holomorphic, with
\[|E_w\iota g|_y\le8B_4\,q_M^{|w|}\,\sup e^{-\delta_2\sum_id(y_i,y_{i+1})}|g|,\qquad q_M=C_qM^{-1}e^{-\delta_2M/(8L)},\]
the rate $\delta_2$ being that of the $SU(2)$ construction and the floor on $M$ a constant depending only on $L$
and the pair; and for $G$-valued $U$ in the class $\mathfrak U'_k(\mathbf B,\alpha)$ of Definition~\ref{def:group-regular-classes},
$\alpha$ small, $\Delta_U+Q'^*\mathsf aQ'$ is positive definite, with an inverse $G'$ that obeys $|G'f|_y\le B_{\rm
ps}\xi_y^2e^{-\delta_{\rm ps}d(y,y')}\sup|f|$, and the like for its covariant derivative, for $(Q'G'^2Q'^*)^{-1}$
and for the gauge-fixed operator.
\item[(V)] \emph{Critical configurations off the trivial field.}
Assume existence of a regular critical orbit over every admitted regular
datum, uniqueness of that orbit at one fixed positive regularity radius,
improvement to a constant times the datum radius, and holomorphic
response with exponential decay, uniformly in depth and volume.
The $SU(2)$ reduced-data models are Theorem~\YMMref{E}{10},
Proposition~\YMMref{37.1}{278}, Theorems~\YMMref{K}{269}
and~\YMMref{M}{315}, and Lemma~\YMMref{51.6}{398}.
The precise existence requirement is stated in Section~\ref{C-368-32}.
Restrictions, joined data and localised data must remain in the admitted
classes; their critical backgrounds must satisfy the comparison in
Proposition~\ref{C-368-16}(iii). These are assumptions for $G$,
including the uniform near-trivial background, tower and analyticity-core bounds of Proposition~\ref{C-368-16}(ii), with the response constant $C_\partial=12\kappa_6+4L$, on the
objects of a step at the fields over a regular base (the global and the localised background of the step and the
minimal orbits of its three data), and on the towers off the trivial field. Among them, the uniqueness of the
regular critical orbit near the trivial field and its normalized representatives
as stated in Proposition~\ref{cg-near-trivial-uniqueness}, for the full
regional geometry of Definition~\ref{def:cg-uniform-slice} and every
reversal-invariant set $\mathfrak F$ of removed inward bonds. The radii and
bounds there are uniform in depth, sub-domain, cutoff and volume.
\item[(C)] \emph{Conjugation.} With $a(X)$ for every trace of a group-valued argument at complex fields, the
gauge-fixing term included, and with the domains written with $m(X)$ and with $D^c$ at the logarithms of the
averages, every map between sets of complex fields that the induction uses commutes with $p\mapsto p^c$ on the
real-connected parts of its domains (Lemma~\ref{C-368-15}) and maps them into the real-connected parts of its
targets. The covariance assertion follows from Proposition~\ref{C-368-22}(i)
for these maps wherever they are formed; the source and target domain memberships
required to apply it remain among \textup{(V)}, \textup{(II)}, and \textup{(K)}.

\item[(II)] \emph{The second part.} The transfers to $G$ of the statements of the second part of the $SU(2)$
construction that the induction cites hold, in the forms of Proposition~\ref{C-368-22}(ii), every proposition that
the proof of Theorem~\YMMref{D}{8} quotes among them, and the objects of the definitions that these statements
read, formed for $G$, have the properties that the $SU(2)$ construction proves for them; so do the statements of
the second part that use the averaging map off the small fields, the pivot bond or the measure of $G\setminus
G^\pi$, the uniform box-gauge bound of Lemma~\ref{C-368-8}(iv), and the large-field chain of Lemma~\ref{C-368-9}(ii); and at every citation in the proofs named in
Proposition~\ref{C-368-17} the assumptions of the cited statement, in its transferred form, hold at the data of the
citation. Among them: the pivot bond (Proposition~\YMMref{14.11}{75}) holds for $G$ in all its clauses with the
constants of equation~\YMMref{(128)}{75}, no numeral changed; and the localised Gaussian family
(Proposition~\YMMref{50.2}{386}) is a pointed family of its level on the analyticity cores of $G$, holomorphic on balls
of one radius, with the contraction bound of its expansion.
\item[(K)] \emph{The constants.} Every condition on the constants read for $G$, the places where $z_G$ enters
included, passes its test in Proposition~\ref{C-368-18}(b): it holds on a large set, and at each two-sided constant
the set of the earlier choices at which its window meets all its ends together is large. The conditions read for
$G$ are every condition of the $SU(2)$ construction transferred to $G$, every condition of
Proposition~\ref{C-368-18}(c)--(g), the conditions $\varepsilon_{\rm v}\le\varepsilon_G$ and $c_{11}\le
r_G/(2C_{\rm pg})$ of the uniqueness clause of (V), $c_{11}$ the radius of Definition~\ref{def:cg-uniform-slice} and $C_{\rm
pg}$ the admitted preliminary-gauge constant of Definition~\ref{def:cg-uniform-slice}, the ceiling $\mathsf c_Lc_4\le2Lr_G$ of
Proposition~\ref{C-368-16}(iv), the row $\alpha_1<\rho_GL$ of Proposition~\ref{C-368-17}, and, at the memberships
reached from memberships of restrictions had apart, the conditions of Lemma~\ref{C-368-15}(iii): the moved first
radius $b'_0=b_0+2e^{4\xi b_1}b_1^2$, for the regular pairs $\alpha_0+2\vartheta^2e^{4\vartheta\alpha_1}\alpha_1^2$
with $\vartheta\ge1$ the largest ratio of the spacings around a plaquette, at most the first radius $B_0$ at which
the statement read there is had, and for the analyticity cores besides $b_1\le B_1$, $K(b'_0+b_1^2)\le B_0$ and
$b'_0+b_1\le a^\sharp$, with $K\ge16L^2$ and $a^\sharp$ the constant and the radius of that lemma, which are
positive ceilings on the second radius where $b_0<B_0$, $Kb_0<B_0$ and $b_0<a^\sharp$. Each number of $G$ that
stands in such a condition has a place before $\gamma$ in the order of the table and is a function of the constants
before that place and of the pair, which reads no radius.
\end{enumerate}
The near-trivial uniqueness clause of (V) is reduced in Proposition~\ref{cg-near-trivial-uniqueness} to the transferred uniform slice and contraction estimates; their uniform nonlinear bounds are assumed there.
The orbit and energy hypothesis of Definition~\ref{C-368-12}, with a fixed admitted positive radius $a_G$, is also required wherever its ceiling is used in (K). For quotients, the sector statements are those of
Proposition~\ref{C-368-27}.
\end{definition}

\subsection{Construction and local coefficients}

\begin{theorem}[Conditional construction and local coefficients]\label{C-368-4}
Let $G\subset U(N)$ be compact and connected with $\mathfrak g$ simple, and let $L\ge13$ be odd with
$200L^2z_G\ge1$ (for $SU(N)$ at $L=13$ this is $N\le212371$). Every constant below is a positive number of $L$ and
the pair, uniform in the volume. The construction conclusions in (a), in the second half of (b), in the existence
assertion of (d), and in the second half of (e) assume (P)--(K) of Definition~\ref{C-368-31}. The fixed-lattice
identities, one-loop coefficients, finite-box expansion, and assertions about an already convergent sequence do not
require these construction hypotheses.
\begin{enumerate}
\item[(a)] Let $G$ be simply connected and satisfy the construction hypotheses. Then Theorem~\YMMref{D}{8}
holds for $G$ at every choice of the constants of the $SU(2)$ construction that meets the conditions on the
constants read for $G$ (Definition~\ref{C-368-31}(K)): its conditions transferred to $G$ and the conditions of
Proposition~\ref{C-368-18}, among them $\varepsilon_{\rm v}\le a_G$, $\varepsilon_{\rm v}\le\varepsilon_G$ and
$\varepsilon_{\rm v}\le L^2r_G$ on the regularity radius $\varepsilon_{\rm v}$ of the critical configurations
(Definition~\YMMref{9.5}{41}), $a_G$ a fixed positive radius admitted by the hypothesis of Definition~\ref{C-368-12} and $\varepsilon_G$
the number of the uniqueness clause of Definition~\ref{C-368-31}(V), and the ceiling $\mathsf c_Lc_4\le2Lr_G$ of
Proposition~\ref{C-368-16}(iv); such choices exist, in the order of the $SU(2)$ construction with $\gamma$ last
(Definition~\ref{C-368-31}(K) and Proposition~\ref{C-368-18}(b)). The run is the one sequence $x_0=\beta$,
$x_k=L^{-1}x_{k-1}-b_k$, $b_k$ the mean of the $n$ quadratic colour coefficients of the local terms of second order of step
$k$: it is real, does not depend on the torus, and $|b_k|\le\mathsf G_\sharp\le B_\sharp$.
\item[(b)] Let $A(v,V)=\sum_pv(p)a_p(V)$ for a complex weight $v$ of the plaquettes, and $\mathcal
O(w)=a_K^{-1}A(w,\cdot)$ for a source $w$. For every compact $G\ne\{\mathbf 1\}$ in $U(N)$: the chessboard norm of
a source is $\|e^{\mathcal O(w)}\|_{\rm ch}^{\mathsf N}=Z_\beta[w^{\rm ch}]/Z_\beta\ge0$
(equation~\YMMref{(1759)}{682}, transferred to $G$); a source $w$ with $\mathrm{Re}\,w\le(1-c_{\rm pr})a_K\beta$,
$c_{\rm pr}\in(0,1]$, leaves every plaquette at least the fraction $c_{\rm pr}$ of its price,
$|e^{-A(\beta-a_K^{-1}w,V)}|\le e^{-c_{\rm pr}\beta\sum_pa_p(V)}$, as $a_p\ge0$ on $G$; and on a torus with
$|\Lambda|$ sites the Wilson measure with one coupling $s\ge1$ at all plaquettes, which at $s=c_{\rm pr}\beta$
dominates the weight with the source by the price, has
\[\langle a_p\rangle_s\le\frac Cs+\frac{n\log s}{|\Lambda|\,s}\]
at every plaquette, $C$ a number of the pair, while for plaquette couplings $v(p)$ with $s\le v(p)\le C's$,
$C'\ge1$, the mean $\frac1{3|\Lambda|}\sum_p\langle a_p\rangle_v$ over the plaquettes of the torus is at most
$(C+n\log C')/s+n\log s/(|\Lambda|s)$, and at couplings invariant under the lattice symmetries every plaquette has
this mean; for $c_{\rm pr}\beta\ge1$ the couplings $\beta-a_K^{-1}\mathrm{Re}\,w(p)$ of a source with
$0\le\mathrm{Re}\,w\le(1-c_{\rm pr})a_K\beta$ are of this kind, with $s=c_{\rm pr}\beta$ and $C'=1/c_{\rm pr}$. On
the tori of the continuum theorem at its first-passage data, which have at least $2L^K$ sites a side and couplings
at most $L^{K+1}(g^{-2}+\bar B)$, $n|\Lambda|^{-1}\log s$ is below a number of $g$, the pair and the constants. Let
$G$ be simply connected and satisfy the construction hypotheses. Then the sourced-density and stability assertions of Theorems~\YMMref{AJ}{673} and~\YMMref{AD}{607} hold for $G$ with
the run of (a): the effective densities with a source have at every level the structure that the $SU(2)$
construction proves, and the stability holds through the whole induction with complex plaquette weights, uniformly
in the volume; writing $\widetilde Z(w)=e^{-E_{(K)}-\mathcal E_K(w)}Z_{K,\ell,\beta}[w]$, with $\mathcal E_K$ defined by \eqref{eq:running-source-energies}, gives
\[\int e^{-A(\beta-a_K^{-1}w,V)}\,dV=\exp\Bigl(E_{(K)}+\sum_{j=1}^K\ \sum_{X\in\mathcal D_j}e_j(X;w)\Bigr)\,\widetilde Z(w),\]
$E_{(K)}$ the vacuum constant at the source $0$, $e_j(X;w)$ the source energy of level $j$ on the domain $X$ of
level $j$ \eqref{eq:running-source-energies} and $\widetilde Z(w)$ the normalised partition function with the source,
with the bounds of Theorem~\YMMref{AD}{607} on them. With it the connected functions of the action density obey
the bound per ball of Proposition~\textup{\YMMref{89.4}{693}}: for every reference coupling $0<g\le\gamma$, every order
$\mathsf m\ge2$, every $R>0$ and every $\delta>0$, at every first-passage datum $(K,\beta)$ at $g$, on every side
$\ell$ admitted at $g$ with $\ell\ge16\mathsf m(R+1)$ and for real Lipschitz functions $f_1,\dots,f_{\mathsf m}$ on
the torus, each supported in a ball of radius $R$, the supports pairwise at a distance at least $\delta$,
$|\langle\mathcal O(f_1);\dots;\mathcal O(f_{\mathsf m})\rangle|$ is at most the constant $C^\infty_{\mathsf
m}(\delta;R)$ of that proposition times the product of the norms of the $f_i$ in that proposition,
$C^\infty_{\mathsf m}(\delta;R)$ being formed from $g$, $\mathsf m$, $\delta$, $R$ and the constants of
Theorem~\YMMref{AD}{607} for $G$ alone, free of $K$, $\beta$ and $\ell$: the connected functions are derivatives
at $0$ of the logarithm of the normalised partition function with sources, and Cauchy's inequality on the polydiscs
of sources on which the stability holds gives the bound, as in Proposition~\textup{\YMMref{89.4}{693}}.
\item[(c)] For $k\ge1$ and on every torus, the one-loop functional $\Psi^G_k$ of Theorem~\YMMref{O}{325} for $G$
is tempered (Section~\YMMref{43.1}{324}), with the coefficient
\[b[\Psi^G_k]=\frac{C_A}8\,b[\Psi_k]+\Bigl(\frac{C_\rho}3-\frac{C_A}8\Bigr)\bigl(L^{-k}-L^{-4k}\bigr),\]
$\Psi_k$ that of $SU(2)$; Theorems~\YMMref{O}{325} and~\ref{III-12-105} hold for $G$ with constants of $L$,
$C_A$ and $C_\rho$; and near the trivial block field the quadratic coefficient of the logarithm of the Gaussian
normaliser of the step to level $l$ is the one-loop increment (Theorem~\YMMref{O}{325})
\[b^0_l(G)=\frac{C_A}8\,b^0_l+\Bigl(\frac{C_\rho}3-\frac{C_A}8\Bigr)(L^3-1)L^{-4l}.\]
\item[(d)] Let the run have real increments that depend on no torus and are bounded by $B_\sharp$, as in (a). Let
$0<g\le\gamma$, $X_g=L(g^{-2}+B_\sharp)$, $K(\beta,g)$ the least $k$ with $x_k\le X_g$, and let $K_j\to\infty$, let
$l_j\to\infty$ be sides admitted at $g$, and $K(\beta_j,g)=K_j$. For $\mathsf m\ge2$ and real $C^1$ functions $f_i$
on $\mathbb R^3$ of compact pairwise disjoint supports let $S^{(j)}_{\mathsf m}(f_1,\dots,f_{\mathsf
m})=\langle\mathcal O(f_1);\dots;\mathcal O(f_{\mathsf m})\rangle$ on the torus of side $l_j$ at $(K_j,\beta_j)$,
$\mathcal O(f)=a_{K_j}^{-1}\sum_pf(x(p))a_p$, $f$ periodised. Then along a subsequence $a_{K_j}\beta_j$ converges
in $[g^{-2}+B_\sharp-\bar B/L,\,L(g^{-2}+\bar B)]$. For simply connected $G$ satisfying the construction
hypotheses, with the run of (a), the $S^{(j)}_{\mathsf m}$ obey the bound per ball of (b) uniformly in $j$, which
holds at every first-passage datum and on every side admitted at $g$ and at least $16\mathsf m(R+1)$ with one
constant for given $\mathsf m$, $R$ and $\delta$, each of the compact pairwise disjoint supports of the $f_i$ lying
in a ball of some radius $R$ and the supports lying at mutual distances at least some $\delta>0$; so along a
further subsequence every $S^{(j)}_{\mathsf m}$ converges.
For simply connected $G$ under all the construction hypotheses
\textup{(P)}, \textup{(V)}, \textup{(C)}, \textup{(II)} and
\textup{(K)}, the following additional conclusions are transferred
with the same parameter restrictions. On a fixed torus the
fixed-label and moving-label cutoff limits have the properties of
Theorem~\YMMref{A}{5}. Growing-torus and periodic-Gibbs sequences
have the subsequential compactness, multilinear bounds and inherited
coordinate reflection positivity of Propositions~\YMMref{92.1}{710},
\YMMref{92.3}{712} and~\YMMref{92.4}{713}. For periodic-Gibbs
limits the canonical component and reconstruction conclusions are
those of Theorems~\YMMref{AL}{746} and~\YMMref{AP}{780}:
each component is a conditioned periodic-Gibbs limit, is Euclidean
invariant and reflection positive, clusters, and has a unique vacuum;
the undecomposed limit has the same reconstruction, with zero-energy
space $L^2(\mathsf X,\lambda)$. For a general growing-torus limit,
full rotation invariance retains the Ward criterion in \textup{(g)}.
The short-distance conclusions retain the restrictions in
\textup{(e)} and use $c_2(G)$ in place of $c_2$.
For
every compact connected $G$ and along every sequence of tori on which all $S^{(j)}_{\mathsf m}$ converge, each
limit $S_{\mathsf m}$ is $\mathsf m$-linear, real and symmetric, obeys $|S_{\mathsf m}(f_1,\dots,f_{\mathsf m})|\le
C(A_1,\dots,A_{\mathsf m})\prod_i\|f_i\|_{C^1}$ for supports in fixed disjoint compact sets $A_i$, with $C$
depending on the sequence, is invariant under the translations of $\mathbb R^3$ and the $48$ signed permutations of
the axes, and is reflection positive at every coordinate mirror. On one torus of side $l=2L^m$ with $K_j\to\infty$
the same holds with the translations of the torus, the signed permutations and the pairs of mirrors
$\{y_\mu=s\}\cup\{y_\mu=s+\frac l2\}$.
\item[(e)] The Gaussian measure defined by the quadratic action at $\mathbf 1$ is the product of $n$ identical
one-colour measures. Its leading two-point coefficient is
\[
 c_2(G)=\frac{3n}{16\pi^2}.
\]
If $G$ is simply connected and satisfies all the construction
hypotheses, the finite-cutoff estimate of Theorem~\YMMref{AV}{851}
and its continuum version, Theorem~\YMMref{AW}{852}, transfer to
its tori and their limits:
\[
 \langle\mathcal O(f);\mathcal O(h)\rangle
 =c_2(G)\mathsf b^{-2}\mathcal K_6(f,h)(1+\varepsilon'),
 \qquad |\varepsilon'|\le\bar\varepsilon_g(\mathsf d),
\]
At finite cutoff this means $0<g\le\gamma_W$,
$(K,\beta)\in\mathfrak d(g)$,
$0<\vartheta<1$, $0<\mathsf d\le d_0(g,\vartheta)$,
$K>s(\mathsf d)$, $\ell\in\mathfrak L(g)$, and a test pair
$(f,h)\in\mathfrak T_{T_\ell}(\mathsf d;\vartheta)$;
then $\mathsf b=a_K\beta$ and the tolerance is the transferred
$\bar\varepsilon_{g,\vartheta}(\mathsf d)\le1/3$ of
Theorem~\YMMref{AV}{851}. In a continuum limit the pairs are $C^1$,
the cutoffs tend to infinity, and $\mathsf b$ is the limiting label
of Theorem~\YMMref{AW}{852}. For block sums of $SU(2)$, the corresponding statement is Theorem~\ref{C-368-5}(iii).
\item[(f)] For every compact connected $G$, the fixed-box expansion is that of Lemma~\ref{C-368-20}. Its
coefficient of $\beta^{-3}$ has an infinite-box limit; after summing over orientations, its leading part at
separation $\mathsf r$ is $\varphi_G\mathsf r^{-5}/(16\pi^3)$. Relative to the coefficient in (e), its ratio is
\[
 c_1(G)=\frac{\varphi_G}{3\pi n}.
\]
The expansion and the limit are iterated, with the box fixed during the expansion in $\beta^{-1}$.
\item[(g)] At every cutoff, on every torus and for every $\beta>0$, the
finite-cutoff rotational Ward identity used in
Section~\ref{C-360-1} holds for $G$ with $\frac12k_p$ in the place
of $\frac32a_p$; and $k_p=na_p$ if and only if $\Omega=n\mathbf 1$, as for an irreducible representation or a sum
of copies of one. Let $(K_j,l_j,\beta_j)$ be a sequence of tori with $l_j\to\infty$ along which all connected
functions converge, $S_{\mathsf m}$ their limits as in (d), and let $(\omega,\mathsf m,\mathcal R,f)$ be an
admitted quadruple (Section~\YMMref{112}{781}, before equation~\YMMref{(2115)}{784}): $\omega\in\mathfrak{so}(3)$, $\mathsf m\ge2$, $\mathcal R>0$,
and $f=(f_1,\dots,f_{\mathsf m})$ smooth functions with compact pairwise disjoint supports in
$\{|y|<\frac12\mathcal R\}$. For the $j$ with $2\mathcal R+2a_{K_j}<\frac12l_j$ the bulk term and the flux term of
the identity at the cut-off radius $\mathcal R$ are defined, and their sum has a limit, which equals
\[\sum_iS_{\mathsf m}(f_1,\dots,L_\omega f_i,\dots,f_{\mathsf m})=-\frac d{dt}S_{\mathsf m}(e^{t\omega}f)\Big|_{t=0},\]
$L_\omega$ the rotation field of $\omega$, $t\mapsto S_{\mathsf m}(e^{t\omega}f)$ being differentiable, and which
does not depend on $\mathcal R$. The limit of the connected functions is invariant under $O(3)$ at all $C^1$ tuples
of (d) if and only if this limit of the sum vanishes for one $\omega_0\ne0$ and every admitted quadruple
$(\omega_0,\mathsf m,\mathcal R,f)$, and if and only if it vanishes at every admitted quadruple: the conjugates of
$\omega_0$ under the signed permutations span $\mathfrak{so}(3)$, and the invariance passes from smooth tuples to
$C^1$ tuples by mollifying, with the bound of (d). It is the limit of the flux term alone wherever the bulk term
tends to $0$. If the limit is invariant, the rotations and reflections about the third axis act on the Hilbert
space of the reflection positive form of the limit by a strongly continuous unitary representation $\mathsf U$ of
$O(2)$ that fixes the vacuum vector, and $\mathsf U(R)[\tau_{te_3}F]=[\tau_{te_3}((R\oplus1)F)]$ for $t\ge0$ on the
classes of the functions $F$ supported in $\{y_3>0\}$, $\tau_{te_3}$ being the translation by $te_3$.
\item[(h)] Let instead $\mathfrak g=\bigoplus_{i\le r}\mathfrak g_i$ be semisimple, $r\ge2$.
\begin{enumerate}
\item[(h1)] The mean of the $n$ quadratic colour coefficients of a local term of second order is
$\mathsf{cf}=\sum_i\frac{n_i}n\mathsf{cf}_i$, where $\mathsf{cf}_i$ is its coefficient on the $i$-th ideal; moreover,
$|\mathsf{cf}_i|\le\kappa_i^2C_{\rm am}r_0^{-2}\|E\|=(\kappa_i/\kappa_G)^2\mathsf G_\sharp\|E\|/E_0$ for the terms
$E$ of Theorem~\YMMref{D}{8}, $C_{\rm am}$ the constant of the coefficient bound~\YMMref{(940)}{428}, $r_0=a^{\rm
us}(\alpha_0,\alpha_1)$ the radius of the $SU(2)$ construction and $\mathsf G_\sharp=\kappa_G^2C_{\rm
am}r_0^{-2}E_0$ (Proposition~\ref{C-368-17}), $C_{\rm am}$ and $r_0$ reading constants of the $SU(2)$ construction
that come before them in its order. If $|e^{(i)}_k|\le\Gamma$, the sequences $x^{(i)}_0=\beta$,
$x^{(i)}_k=L^{-1}x^{(i)}_{k-1}-e^{(i)}_k$ and their mean $x_k=\sum_i\frac{n_i}nx^{(i)}_k$ satisfy
$x_k=L^{-1}x_{k-1}-\sum_i\frac{n_i}ne^{(i)}_k$ and $|x^{(i)}_k-x_k|\le2(1-\frac{n_i}n)\Gamma\frac
L{L-1}(1-L^{-k})$.
\item[(h2)] The one-loop theorem holds in block form at every semisimple pair, the ideals being tied by the
representation or not: on $\mathfrak g_i$,
$b^{(i)}[\Psi^G_k]=\frac{C^{(i)}_A}8b[\Psi_k]+(\theta_i-\frac38C^{(i)}_A)t_k$, $t_k=\frac13(L^{-k}-L^{-4k})$,
$b^{(i)}[\Psi^G_k]$ being the coefficient on $\mathfrak g_i$ of the one-loop functional $\Psi^G_k$ of the pair and
$\Psi_k$ that of $SU(2)$, with one set of sequences for all ideals; here $\theta_i=\frac1N\mathrm{tr}(\Omega H^2)$
for a unit $H\in\mathfrak g_i$, $\Omega$ being summed over a basis of all of $\mathfrak g$, so that
$\theta_i=\sum_jn_j\tau_{ij}$ with numbers $\tau_{ij}$ of the pair (Proposition~\ref{C-368-13}(iv)). Equivalently,
$b^{(i)}[\Psi^G_k]$ is the coefficient of the simple pair of $\mathfrak g_i$ plus $(\sum_{j\ne i}n_j\tau_{ij})$
times the coefficient $t_k$ of one scalar form $\Upsilon_k$ of the lattice, which is tempered. For a unit
$H\in\mathfrak g_i$ the numbers $\mathrm{tr}(\mathrm{ad}\,H)^2$ and $\frac1N\mathrm{tr}(\Omega H^2)$ that
Definition~\ref{C-368-2} forms for simple $\mathfrak g$ are $C^{(i)}_A$ and $\theta_i$; they depend on the ideal
wherever these pairs are not all equal, as at $SU(2)\times SU(3)$ on $\mathbb C^2\oplus\mathbb C^3$, where
$C^{(i)}_A$ is $20$ and $30$ (Lemma~\ref{C-368-11}(i)), and the coefficient is then given ideal by ideal, as above.
The $r$ coefficients are bounded ideal by ideal,
$|b^{(i)}[\Psi^G_k]|\le\frac{C^{(i)}_A}8\sup_k|b[\Psi_k]|+|\frac{\theta_i}3-\frac{C^{(i)}_A}8|L^{-1}$, and so by
the largest of these $r$ bounds (Proposition~\ref{C-368-13}(iv)).
\item[(h3)] The action of $G$ has \emph{property (T)} at a choice of the constants if at that choice, for the tied
action, the backgrounds have open slices as the setting preceding Definition~\YMMref{52.1}{415} asks, the Gaussian structure
of Theorem~\YMMref{U}{378} and the normaliser of Proposition~\YMMref{50.2}{386} are formed with the properties
that the $SU(2)$ construction proves for them, Proposition~\YMMref{51.16}{412} (the localisation) and
Theorem~\YMMref{V}{416} (one step, the counterterm numbers being free) hold in their transferred forms, and
Theorem~\YMMref{X}{427} holds at the simple pairs $(G_i,\Lambda)$ of the ideals on $\mathbb C^N$, with the fixed
weight $\Lambda=\frac1N\mathbf 1$ of the action of $G$, which lies in the commutant (Proposition~\ref{C-368-17}).
If the action of $G$ has property (T) at a choice of the constants, then one run suffices, the mean
$x_k=L^{-1}x_{k-1}-\sum_i\frac{n_i}ne^{(i)}_k$, $e^{(i)}_k$ the coefficient on $\mathfrak g_i$ of the local terms
of second order of step $k$: the sequences $x^{(i)}_k$ of (h1) stay within $2(1-\frac{n_i}n)\mathsf G_\sharp\frac
L{L-1}(1-L^{-k})$ of $x_k$, by (h1) with $\Gamma=\mathsf G_\sharp$, as $|e^{(i)}_k|\le(\kappa_i/\kappa_G)^2\mathsf
G_\sharp\le\mathsf G_\sharp$, and all clauses of the $SU(2)$ construction on the increments hold, the increment
theorem, the regular progression, the sequence and the bounds on polydiscs among them. At the least tied pair, the additional propagator assertions require the regular-domain and relative-smallness hypotheses stated in Proposition~\ref{C-368-26}(v); under them the
kinematic and interaction parts of the first step, the Hessian lemma and the propagators at every level hold
(Proposition~\ref{C-368-26}(v)).
\item[(h4)] If $G=G_1\times\dots\times G_r$ acts on $\mathbb C^{N_1}\oplus\dots\oplus\mathbb C^{N_r}$ by blocks,
the Wilson measure at $\beta$ is the product of those of the $G_i$ at $\beta N_i/N$, cumulants add with the weights
$(N_i/N)^{\mathsf m}$, and for simply connected $G_i$ with simple Lie algebras and $200L^2z_{G_i}\ge1$, each
satisfying the construction hypotheses, clauses (a)--(g) hold factor by factor at these couplings, with precisely
the qualifications stated there.
\end{enumerate}
\end{enumerate}
\end{theorem}

\begin{remark}[Leading law and perturbative coefficients]\label{C-368-23}
\emph{The constants.} The Hessian of $a$ at $\mathbf 1$ is $\langle\cdot,\cdot\rangle$, so to leading order the
field consists of $n$ independent copies of the free field with this form, colour by colour, each contributing
$\frac3{16\pi^2}$ as at $SU(2)$, where $3\cdot\frac3{16\pi^2}=\frac9{16\pi^2}$; this is
$c_2(G)=\frac{3n}{16\pi^2}$, the constant of the law of leading order. For $SU(2)^r$ the product $c_2(G)\,\mathsf
b^{-2}$ is the sum $\sum_i\lambda_i^2\frac9{16\pi^2}(\lambda_i\mathsf b)^{-2}=r\frac9{16\pi^2}\mathsf b^{-2}$ that
Step~3 of the proof of Theorem~\ref{C-368-5} forms from the laws of the factors. The next term is governed by the
coefficient of $\beta^{-3}$ of $\langle a_p;a_q\rangle$; by Lemma~\ref{C-368-20} its leading part, summed over
orientations, is $\frac{\varphi_G}{16\pi^3}\mathsf r^{-5}$ at distance $\mathsf r$, and
$c_1(G)=\frac{\varphi_G}{16\pi^3}\big/c_2(G)=\frac{\varphi_G}{3\pi n}$ is the ratio of the next perturbative
coefficient to the leading coefficient. At $SU(2)$, $\varphi=24$, and this is
$\frac3{2\pi^3}\big/\frac9{16\pi^2}=\frac8{3\pi}$, the coefficient $c_1$ of Theorem~\textup{\ref{IV-73-100}}.

\emph{The law of leading order.} Let $G$ be simply connected with $\mathfrak g$ simple, $L\ge13$ odd and
$200L^2z_G\ge1$, and assume the construction hypotheses. On the tori of the continuum theorem and in their limits
along tori, under the finite-cutoff restrictions of Theorem~\YMMref{AV}{851}
and the continuum restrictions of Theorem~\YMMref{AW}{852},
as stated in Theorem~\ref{C-368-4}(e),
\[\langle\mathcal O(f);\mathcal O(h)\rangle=c_2(G)\,\mathsf b^{-2}\mathcal K_6(f,h)\,(1+\varepsilon'),\qquad|\varepsilon'|\le\bar\varepsilon_g(\mathsf d),\]
$\mathsf b=a_K\beta$ the label and $\bar\varepsilon_g$ the number of that law, which tends to $0$ with $\mathsf d$;
this is the law of leading order of Theorem~\ref{C-368-4}(e). For $SU(2)^r$ the law holds with one label
(Theorem~\ref{C-368-5}(iii)). For simply connected $G$ with $\mathfrak g$ simple its proof is the proof of the law
of Theorems~\YMMref{AV}{851} and~\YMMref{AW}{852}, with the colour sum running to $n$. That proof rests on Theorems~\YMMref{D}{8}, \YMMref{AJ}{673} and~\YMMref{AD}{607} and on the
count of colours: Theorem~\ref{C-368-4}(a) gives Theorem~\YMMref{D}{8} for $G$ along the one run; Theorems~\YMMref{AJ}{673} and~\YMMref{AD}{607}
for $G$ with that run (Theorem~\ref{C-368-4}(b)) gives the stability with sources; the Gaussian measure that the
quadratic part of the action at the trivial field defines is the product of $n$ copies of that of one colour of
$SU(2)$, with the same normalisation (Theorem~\ref{C-368-4}(e)); and the remainders are bounded by the same terms
of the effective action as for $SU(2)$. The passage to limits along tori is that of Theorem~\ref{C-368-4}(d). The
law concerns the tori and their limits along tori.

The fixed-box expansion in Lemma~\ref{C-368-20} has an error constant depending on the box. Its large-box limit is
taken after the expansion at fixed box. Thus the calculation of $c_1(G)$ in this lemma is an iterated perturbative
limit.
\end{remark}

\begin{corollary}[$SU(3)$]\label{C-368-28}
Let $G=SU(3)$ and $L\ge13$ odd. Then $z_G=1$, so that no floor on $L$ is needed and the forms with $D^c$ are the
statements of the $SU(2)$ construction; no number of the averaging layer depends on the group; $r_G=\sqrt3$,
$\kappa_G^2=\tau_G^2=2$, $n=8$; the least admissible $p$ of the cut-off is $3$, with $\lambda=6.1194$; and the
extra conditions of Lemma~\ref{C-368-7}(v) and Proposition~\ref{C-368-16}(iv) are void. The local coefficients in
Theorem~\ref{C-368-4} have the following values. Its construction conclusions hold subject to
Definition~\ref{C-368-31}:
\[b[\Psi^G_k]=\frac94\,b[\Psi_k]+\frac5{12}\bigl(L^{-k}-L^{-4k}\bigr),\qquad c_2(G)=\frac3{2\pi^2},\qquad c_1(G)=\frac6\pi .\]
Besides, $\mathfrak{su}(3)$ has a basis, orthonormal for $\frac12\mathrm{tr}\,XY$, whose eight members have the
eigenvalues $1,-1,0$, as the Pauli matrices, with the eigenvalues $1,-1$, are such a basis of $\mathfrak{su}(2)$.
\end{corollary}

\begin{proof}[Sketch of proof]
The numbers are those of $SU(N)$ at $N=3$ (Definition~\ref{C-368-2}); $\frac{C_A}8=\frac94$,
$\frac{C_\rho}3-\frac{C_A}8=\frac5{12}$, $\frac{3\cdot8}{16\pi^2}=\frac3{2\pi^2}$, $\frac{C_A}{3\pi}=\frac6\pi$.
The cut-off asks $3^{1/2p}<\frac54$, first met at $p=3$, and $\lambda=w/(5-4w)$ at $w=3^{1/6}$. With $r_G=\sqrt3$
the condition $\varpi_L<r_G$ and the condition $\mathsf c_Lc_4\le2Lr_G$ are met by the ceilings of the $SU(2)$
construction. For the basis: the rotation $R$ by $\frac{2\pi}5$ about $e_3$ acts on the real symmetric traceless
matrices, an orthogonal sum $V_0\oplus V_1\oplus V_2$ of the multiples of $\mathrm{diag}(-1,-1,2)$, of the span of
$e_1e_3^{\mathrm T}+e_3e_1^{\mathrm T}$ and $e_2e_3^{\mathrm T}+e_3e_2^{\mathrm T}$, and of the span of
$\mathrm{diag}(1,-1,0)$ and $e_1e_2^{\mathrm T}+e_2e_1^{\mathrm T}$, as the identity, the rotation by
$\frac{2\pi}5$ and the rotation by $\frac{4\pi}5$. So for $X$ of unit norm with squared components $(c_0,c_1,c_2)$,
$\frac12\mathrm{tr}(R^jXR^{-j}R^kXR^{-k})=c_0+c_1\cos\frac{2\pi(j-k)}5+c_2\cos\frac{4\pi(j-k)}5$, which vanishes
for $j\not\equiv k$ exactly when $(c_0,c_1,c_2)=(\frac15,\frac25,\frac25)$. For $X=ab^{\mathrm T}+ba^{\mathrm T}$
with orthonormal $a,b\in\mathbb R^3$, which has the eigenvalues $1,-1,0$, one finds $c_0=3a_3^2b_3^2$ and
$c_2=(1-a_3^2)(1-b_3^2)$; the two conditions say that $a_3^2$, $b_3^2$ are the roots $\frac{5\pm\sqrt{10}}{15}$ of
$15x^2-10x+1=0$, whose sum $\frac23$ leaves $\frac13$ for the square of the third component of $a\times b$, so such
$a,b$ exist. The five $R^kXR^{-k}$ and the three imaginary antisymmetric Gell-Mann matrices, which are orthogonal
to all real symmetric matrices, have the eigenvalues $1,-1,0$ and are pairwise orthogonal.
\end{proof}

\begin{theorem}[Sums of copies of $SU(2)$]\label{C-368-5}
Let $G=SU(2)^r$ act on $\mathbb C^2\oplus\dots\oplus\mathbb C^2$ by blocks,
$a_p=\sum_i\lambda_i(1-\frac12\mathrm{tr}\,\partial U^{(i)}(p))$ with $\lambda_i>0$, $\sum_i\lambda_i=1$ (a block
repeated $m_i$ times in $\mathbb C^N$ gives $\lambda_i=2m_i/N$), one coupling $\beta$, $n=3r$, $\lambda_\pm$ the
largest and the least $\lambda_i$, and let the constants of the $SU(2)$ construction, $\mathsf G_\sharp$ among them
with its value for $SU(2)$ (Proposition~\ref{C-368-17}), obey besides $\lambda_-B_\sharp\ge\lambda_+\mathsf
G_\sharp$.
\begin{enumerate}
\item[(i)] On every finite lattice and for $\beta>0$ the Wilson measure is the product of those of $SU(2)$ at
$\lambda_i\beta$, and $\langle\mathcal O(f_1);\dots;\mathcal O(f_{\mathsf m})\rangle_\beta=\sum_i\lambda_i^{\mathsf
m}\langle\mathcal O(f_1);\dots;\mathcal O(f_{\mathsf m})\rangle^{SU(2)}_{\lambda_i\beta}$.
\item[(ii)] For $\lambda_-\beta\ge\gamma^{-2}$ the mean run $x_0=\beta$,
$x_k=L^{-1}x_{k-1}-\frac1r\sum_i\lambda_i^{-1}c^{(i)}_k$, with $c^{(i)}_k$ the coefficient of step $k$ of $SU(2)$
along the couplings $\lambda_ix_0,\dots,\lambda_ix_{k-1}$, is defined and real while
$\lambda_-x_{k-1}\ge\gamma^{-2}$. Along it Theorem~\YMMref{D}{8} holds in every factor, the step of the factor
$i$ at level $j$ being the step of $SU(2)$ at $\lambda_ix_j$ with its own thresholds and histories; and
$x^{(i)}_0=\beta$, $x^{(i)}_k=L^{-1}x^{(i)}_{k-1}-\lambda_i^{-1}c^{(i)}_k$ have the mean $x_k$ and
$|x^{(i)}_k-x_k|\le2\lambda_-^{-1}\mathsf G_\sharp\frac L{L-1}$.
\item[(iii)] Let $\lambda_-g^{-2}\ge\gamma_W^{-2}+\bar B$, $K$ the least $k$ with $x_k\le L(g^{-2}+\frac{\mathsf
G_\sharp}r\sum_i\lambda_i^{-1})$, $\mathsf b=a_K\beta$, and let the hypotheses of the law of two points of
Theorem~\YMMref{AV}{851} on the scale $\mathsf d$, the cutoff
and the test pair $(f,h)$ hold separately at the $2r$
reference couplings $g_{i,1}$, $g_{i,2}$ with $g_{i,1}^{-2}=\lambda_ig^{-2}-\bar B$,
$g_{i,2}^{-2}=L\lambda_ig^{-2}-3\bar B$, and let the tori have sides admitted at each of the $2r$ couplings
$g_{i,s}$, that is, sides $2L^m$ with $m$ at least the largest of the $2r$ least exponents admitted at them, an
unbounded set. Then
\[\langle\mathcal O(f);\mathcal O(h)\rangle=\frac{3n}{16\pi^2}\,\mathsf b^{-2}\mathcal
K_6(f,h)\,(1+\varepsilon'),\qquad|\varepsilon'|\le\max_{i,s}\bar\varepsilon_{g_{i,s}}(\mathsf
d)+\tfrac{r-1}{12}\varepsilon^\times(\mathsf d),\]
uniformly in $K$ and in $\beta$ in its window, on these tori, in their limits along such tori, and in every Gibbs
state of $e^{-\beta\sum_pa_p}$ whose marginals on the factors are periodic; $\bar\varepsilon_g$ is the number of
that law, and $\varepsilon^\times(\mathsf d)$, a number of $g$, the $\lambda_i$, the parameter of the test pairs
and the constants, is at most a constant times $\mathsf d$ and vanishes on the tori and in product states.\qedhere
\end{enumerate}
\end{theorem}


\subsection{Inputs from the $SU(2)$ construction}

We start from Theorems~\YMMref{D}{8}, \YMMref{AJ}{673} and~\YMMref{AD}{607} for $SU(2)$ and from their proofs: the inductive
step, Proposition~\YMMref{69.3}{603} and Theorem~\YMMref{AC}{604}; the flow bound~\YMMref{(1468)}{597} and the intrinsic coupling construction in Section~\YMMref{52.10}{439}; the large-field chain, Proposition~\YMMref{44.4}{347}, Proposition~\YMMref{44.5}{348} and Lemma~\YMMref{65.2}{580} and the local-deposit argument in the proof of Proposition~\YMMref{65.4}{582}; and the verification of the conditions on the constants,
the ordered choice and its numerical verification in Proposition~\YMMref{69.1}{595}. From the first part we use the mean and the block
average (equation~\YMMref{(25)}{18} and~\YMMref{(37)--(39)}{20}, Proposition~\YMMref{4.3}{18}), the propagators in a background
(the precise input (P) of Definition~\ref{C-368-31}), the critical
configurations and their localisation (input (V) there and
Proposition~\ref{C-368-16}(iii)), the bond map (Lemma~\YMMref{53.3}{444}), the elimination factor
(Lemma~\ref{C-368-10}(iii)), the one-loop Theorem~\YMMref{O}{325}, with the determinant representation and telescoping argument in Section~\YMMref{43.3}{329} and Lemma~\ref{lem:gaussian-flow-bound}, and its limit, Theorem~\ref{III-12-105}; and the continuum
theorems with their proofs. Of $G$ we use the closed subgroup theorem, the spectral theorem, Schur's lemma, the
Haar measure, the Peter--Weyl theorem, and for groups other than $SU(N)$ maximal tori and roots.


\subsection{Matrix charts, averages and gauges}

\begin{lemma}[Norms, the domain $D^c$ and means]\label{C-368-6}
\begin{enumerate}
\item[(i)] $\|Y\|\le|Y|\le\kappa_G\|Y\|$ on $\mathfrak g$, with $\kappa_G^2\le N-1$ for semisimple $\mathfrak g$;
$a(V)=\frac12\|V-\mathbf 1\|^2$ on $U(N)$; on $SU(N)$, $\frac1{2(N-1)}|V-\mathbf 1|^2\le a(V)\le\frac12|V-\mathbf
1|^2$ and $|V-\mathbf 1|\le(N-1)^{1/2}\|V-\mathbf 1\|$, all three sharp. So $a(V)\ge\frac12\tau_G^{-2}|V-\mathbf
1|^2$ on $G$, with $\tau_G^2\le N-1$ for $G\subset SU(N)$. For $n\ge2$, $X+X^{-1}$ is scalar for every $X\in G$ if
and only if $G$ is conjugate to $\{u\otimes\mathbf 1_k:u\in SU(2)\}$.
\item[(ii)] $Y\mapsto e^{iY}$ maps $\{Y\in\mathfrak g:|Y|<\pi\}$ bijectively onto $G^\pi$, which is the component
of $\mathbf 1$ in $\{u\in G:-1\notin\mathrm{spec}\,u\}$. If $A_j\in\mathfrak g$ and $\sum_j|A_j|<\pi$, then
$e^{iA_1}\cdots e^{iA_m}=e^{iA}$ with $A\in\mathfrak g$, $|A|\le\sum_j|A_j|$. For hermitian $Y$, $Y'$ with
$|Y|,|Y'|\le t<\pi$: $|Y-Y'|\le\frac t{\sin t}|e^{iY}-e^{iY'}|$. For a matrix $X$ with $s=|X-\mathbf 1|<1$:
$|X^{-1}-\mathbf 1|\le\frac s{1-s}$; and $|X^{-1}|\le|\mathrm{Ad}\,X|$ for $X\in SL(N,\mathbb C)$, $\mathrm{Ad}\,X$
being read on all $N\times N$ matrices, where $|\mathrm{Ad}\,X|=|X|\,|X^{-1}|$; with the norm of $\mathrm{Ad}\,X$
on $\mathfrak g^c$ alone in its place the inequality fails for some pairs, as the proof of Lemma~\ref{C-368-14}
shows.
\item[(iii)] For compact connected $G\subset U(N)$ the statements $\mathfrak g^c=\{Z:\mathrm{tr}\,Z=0\}$,
$G^c=\{X:\det X=1\}$ and $G=SU(N)$ are equivalent; for semisimple $\mathfrak g$, $\mathrm{tr}\,Z=0$ on $\mathfrak
g^c$ and $\det X=1$ on $G^c$, and either inclusion the other way holds only at $SU(N)$. For every compact connected
$G$: if $Z_1,\dots,Z_q\in\mathfrak g^c$ and $\sum_j|Z_j|<\ln2$, then $|e^{Z_1}\cdots e^{Z_q}-\mathbf 1|<1$ and
$\log(e^{Z_1}\cdots e^{Z_q})\in\mathfrak g^c$, and it lies in $i\mathfrak g$ if the $Z_j$ do. Let $\mathfrak g$ be
semisimple. Then $Z(G^c)=Z(G)$ is a finite set of unitary matrices; for $X\in G^c$ with $|X-\mathbf 1|<1$,
$z(X)=\exp((1-P)\log X)$ lies in $Z(G)$, $|z(X)-\mathbf 1|\le|X-\mathbf 1|$, $\log X=\log z(X)+P\log X$, and $z$ is
locally constant. So $D^c=\{z(X)=\mathbf 1\}$ is the connected component of $\mathbf 1$ in $\{X\in G^c:|X-\mathbf
1|<1\}$ and the trace on $G^c$ of an open set of matrices; it contains $\{X\in G^c:|X-\mathbf 1|<z_G\}$, and no
larger ball if $z_G<1$; it is kept by the powers $X^t=e^{t\log X}$, $0\le t\le1$, by continuous paths in the unit
ball that begin in it, by products of elements $X_j$ of $D^c$ with $\prod_j(1+|X_j-\mathbf 1|)<2$, by the products
$e^{Z_1}\cdots e^{Z_q}$ with $Z_j\in\mathfrak g^c$, $\sum_j|Z_j|<\ln2$, by conjugation with $G$, by inverses in the
unit ball and by adjoints; $\{X\in G^c:|\log X|<2\arcsin\frac{z_G}2\}\subset D^c$; and $D=\exp(i\{Y\in\mathfrak
g:|Y|<\frac\pi3\})=G^\pi\cap\{|X-\mathbf 1|<1\}$. Let $0<c_j\le1$. Within the domain cut out in stages by
$|X_j-\mathbf 1|<c_j$ for continuous $G^c$-valued maps, retain these radius conditions and add $X_j\in D^c$ at each stage. The resulting domain is open and closed in the original domain,
and equals it if $z_G=1$.
\item[(iv)] Let $V_1,\dots,V_q$ be invertible, $d=\max_{i,j}|V_iV_j^{-1}-\mathbf 1|\le\frac1{40}$. Then

\[
\Phi(Y)=Y+\frac1q\sum_i\log(V_iV_1^{-1}e^{-Y})
\]
 is a contraction with factor $\frac25$ of $\{|Y|\le\frac18\}$; its
fixed point has $|Y_*|\le2d$, and $\mathrm{av}=e^{Y_*}V_1$ satisfies $\sum_i\log(V_i\mathrm{av}^{-1})=0$ and
$|V_i\mathrm{av}^{-1}-\mathbf 1|\le4d$, is holomorphic in $(V_i)$, and is the only unitary mean of a unitary
family. Let $\mathfrak g$ be semisimple, let the $V_i$ lie in $G^c$, and let the quotients $V_iV_1^{-1}$ lie in
$D^c$, which holds if $d<z_G$. Then $Y_*\in\mathfrak g^c$, $\mathrm{av}\in G^c$ and $V_i\mathrm{av}^{-1}\in D^c$
for all $i$; and $\mathrm{av}\in G$ for a family in $G$. Further, 
\[
\mathrm{av}(W\mathbf VW')=W\,\mathrm{av}(\mathbf
V)\,W'
\]
 for $W,W'\in G$. For semisimple $\mathfrak g$, if $z_G\le\frac1{40}$ and $z\in Z(G)$ has $|z-\mathbf
1|=z_G$, the family $(\mathbf 1,z)$ has no mean in $G$.
\end{enumerate}
\end{lemma}

\begin{proof}[Sketch of proof]
(i) For $V$ with the eigen-angles $\theta_j$, $|V-\mathbf 1|=\max_j2|\sin\frac{\theta_j}2|$ and
$a(V)=\frac2N\sum_j\sin^2\frac{\theta_j}2=\frac12\|V-\mathbf 1\|^2$. On $SU(N)$, $\sum_j\theta_j\in2\pi\mathbb Z$,
so $|\sin\frac{\theta_1}2|\le\sum_{j\ge2}|\sin\frac{\theta_j}2|$, and Cauchy's inequality gives
$\sin^2\frac{\theta_1}2\le(N-1)\sum_{j\ge2}\sin^2\frac{\theta_j}2$, which is $|V-\mathbf 1|^2\le(N-1)\|V-\mathbf
1\|^2$; the same argument with the eigenvalues of a traceless hermitian $Y$ gives $\kappa_G^2\le N-1$. The
constants are approached along $\exp(it\,\mathrm{diag}(N-1,-1,\dots,-1))$, $t\to0$, and $\frac12$ is attained at
$e^{2\pi i/N}\mathbf 1$. If $X+X^{-1}$ is scalar on $G$, then $2\cos(tY)$ is scalar for $Y\in\mathfrak g$ and all
$t$, so $Y^2=\|Y\|^2\mathbf 1$ and orthogonal elements of $\mathfrak g$ anticommute; for orthonormal
$X,Y\in\mathfrak g$ the element $Z=iXY=\frac i2[X,Y]$ lies in $\mathfrak g$, and an element $W$ orthogonal to $X$,
$Y$, $Z$ would anticommute with $X$ and $Y$, hence commute with $Z$, and anticommute with $Z$, which is absurd as
$Z^2=W^2=\mathbf 1$; so $\mathfrak g$ is spanned by $X,Y,Z$, these generate a quaternion algebra, and $\mathbb C^N$
is a multiple of its two-dimensional representation.

(ii) The path $s\mapsto e^{isA_1}\cdots e^{isA_m}$ has speed at most $\sum_j|A_j|$, so its eigen-angles stay below
$\pi$, and its logarithm, continuous along it and in $\mathfrak g$ near $\mathbf 1$ by the closed subgroup theorem,
stays in $\mathfrak g$; the Lipschitz bound is the spectral calculus of $Y\mapsto e^{iY}$; the bound on $X^{-1}$ is
the series $\sum_j(\mathbf 1-X)^j$; and for $X\in SL(N,\mathbb C)$ with singular values
$\sigma_1\ge\dots\ge\sigma_N$, $\prod\sigma_j=1$, one has
$|X^{-1}|=\sigma_N^{-1}\le\sigma_1\sigma_N^{-1}=|X|\,|X^{-1}|=|\mathrm{Ad}\,X|$, the norm of $\mathrm{Ad}\,X$ on
all $N\times N$ matrices being attained at $vw^*$ for unit vectors $v$, $w$ with $|Xv|=|X|$ and
$|(X^{-1})^*w|=|X^{-1}|$.

(iii) $G^c=G\exp(\mathfrak g)=\{ue^H:u\in G,\ H\in\mathfrak g\}$, $e^H$ being positive (a polar decomposition got
by integrating $e^{tZ}$ inside the set $\{ue^H\}$), and $\det e^Z=e^{\mathrm{tr}\,Z}$; for semisimple $\mathfrak g$
every element of $\mathfrak g^c$ is a sum of brackets, so its trace on every invariant subspace vanishes. The rule
for products holds for every compact connected $G$: along $X_t=e^{tZ_1}\cdots e^{tZ_q}$, $0\le t\le1$,
$|X_t-\mathbf 1|\le e^{t\sum_j|Z_j|}-1<1$, and $X_t'X_t^{-1}=\sum_j\mathrm{Ad}(e^{tZ_1}\cdots e^{tZ_{j-1}})Z_j$
lies in $\mathfrak g^c$, as $\mathrm{Ad}(G^c)$ keeps $\mathfrak g^c$; $W_t=\log X_t$ has
$X_t'X_t^{-1}=\Xi(\mathrm{ad}\,W_t)W_t'$ with $\Xi(w)=(e^w-1)/w$, and $\Xi(\mathrm{ad}\,W_t)$ is invertible, the
eigenvalues of $\mathrm{ad}\,W_t$ being differences of logarithms of eigenvalues of $X_t$, with imaginary parts in
$(-\pi,\pi)$; so $W_t$ solves $W'=\Xi(\mathrm{ad}\,W)^{-1}(X'X^{-1})$ from $W_0=0$. For a real Lie algebra
$\mathfrak h$ of matrices that holds $W$, $\Xi(\mathrm{ad}\,W)$, a power series in $\mathrm{ad}\,W$ with real
coefficients, maps $\mathfrak h$ into itself and, being invertible, onto itself; so the equation keeps $\mathfrak
h=\mathfrak g^c$, and $\mathfrak h=i\mathfrak g$ if the $Z_j$ lie there, and $W_t$ stays in $\mathfrak h$. With
$W=\log X$ the eigenvalues of $\mathrm{ad}\,W$ have imaginary parts in $(-\pi,\pi)$, so $\mathrm{ad}\,W$ is a
polynomial in $e^{\mathrm{ad}\,W}=\mathrm{Ad}(X)$ and keeps $\mathfrak g^c$; hence $(1-P)W$ commutes with
$\mathfrak g^c$, it is a scalar on each isotypic component, the mean there of the eigenvalues of $W$, and
$z(X)=Xe^{-PW}\in G^c$ is central. A mean $c$ of numbers $\log\lambda_j$ with $|\lambda_j-1|<r\le1$ has
$|e^c-1|<r$, the set $\{|e^{x+iy}-1|<r,\ |y|<\frac\pi2\}$ being convex; this gives $|z(X)-\mathbf 1|\le|X-\mathbf
1|$, and $z$, continuous with values in the finite set $Z(G)$, is locally constant. The products are joined to
$\mathbf 1$ inside the unit ball by $t\mapsto e^{tZ_1}\cdots e^{tZ_q}$ and $t\mapsto X_1^t\cdots X_q^t$; the inner
ball comes from $|\log X|\ge\theta(z(X))$, $\theta(V)$ the largest modulus of the argument of an eigenvalue of $V$.

(iv) $|\log(\mathbf 1+E)|\le-\log(1-|E|)$ and the bound on the derivative of $\log$ for $|E|\le\frac1{10}$ give the
contraction; if the quotients lie in $D^c$, so does $V_iV_1^{-1}e^{-Y}$ for $Y\in\mathfrak g^c$, $|Y|\le\frac18$,
by the rule for products of (iii), so $\Phi$ keeps $\mathfrak g^c$, $Y_*\in\mathfrak g^c$ and
$V_i\mathrm{av}^{-1}=V_iV_1^{-1}e^{-Y_*}\in D^c$; and $\mathrm{av}=e^{Y_*}V_1\in G^c$. For unitary $V_i$, $V'$ with
$\sum_i\log(V_iV'^{-1})=0$ the spectral theorem gives that every angle of $V_jV'^{-1}$ is at most
$\frac{2\pi}{3\sqrt3}d$, whence uniqueness. At $(\mathbf 1,z)$, $\mathrm{av}=e^{\frac12\log z}$, which is not in
$G^c$ (for $SU(N)$ its determinant is $-1$).
\end{proof}

\begin{lemma}[A bond in the contour gauge]\label{lem:group-contour-gauge}
Let $W$ be a real $G$-valued field on the bonds of a block $B(y)$,
and let $\varpi_y(W)$ be the largest plaquette distance from
$\mathbf1$ in that block. Suppose
\[
 \tfrac34(L-1)^2\varpi_y(W)<\omega_{\rm c},
 \qquad \varpi_y(W)<z_G,
\]
where $\omega_{\rm c}:=1/40$, the contour-mean radius admitted by
Proposition~\YMMref{4.3}{18}. Then the contour means $W(y,x)$
exist in $G$, and every bond $b$ inside the block satisfies
\[
 |W(y,b_-)W(b)W(y,b_+)^{-1}-\mathbf1|
       \le2\kappa^{\rm SO}_1\varpi_y(W).
\]
Here $\kappa^{\rm SO}_1=\pi(L-1)^2/(2\sqrt3)$.
\end{lemma}

\begin{proof}
Two coordinate-order paths from $y$ to $x$ differ by at most
$\frac34(L-1)^2$ plaquettes. Telescoping their transports gives
the displayed diameter bound for the contour family. Its quotients
are products of conjugated plaquette variables. The guard
$\varpi_y(W)<z_G$ and Lemma~\ref{C-368-6}(iii) put their logarithms
in $\mathfrak g$; Lemma~\ref{C-368-6}(iv) therefore gives the mean
in $G$. Apply the inverse-exponential inequality of
Lemma~\ref{C-368-6}(ii) to the intrinsic logarithms of these
quotients, then the zero-sum-logarithm argument of
Lemma~\YMMref{14.5}{64}. This gives
\[
 |W(\Gamma)W(y,x)^{-1}-\mathbf1|
       \le\kappa^{\rm SO}_1\varpi_y(W)
\]
for every coordinate-order path $\Gamma$ from $y$ to $x$.
Reverse $b$ if necessary so that its step moves away from $y$
in its coordinate direction. Choose that direction last in the
coordinate order. The resulting path $\Gamma'$ from $y$ to $b_+$
ends with $b$, and deleting this last bond gives a coordinate-order
path $\Gamma$ to $b_-$. Thus
\[
 W(y,b_-)W(b)W(y,b_+)^{-1}
 = [W(\Gamma)W(y,b_-)^{-1}]^{-1}
       [W(\Gamma')W(y,b_+)^{-1}].
\]
The two factors lie in $G$ and have distance at most
$\kappa^{\rm SO}_1\varpi_y(W)$ from $\mathbf1$.
The triangle inequality and unitary invariance prove the claim.
\end{proof}


\begin{lemma}[Block averages]\label{C-368-7}
Let $L\ge13$ be odd and let $c$ be a coarse bond.
\begin{enumerate}
\item[(i)] Let $U$ be $G$-valued with $|\partial U(p)-\mathbf 1|<\min\{\omega_M,r_G\}$ at the plaquettes of the
double block of $c$. Then the contour quotients lie in $D$ within $\frac1{160}$ of $\mathbf 1$, each contour family
(Section~\YMMref{4}{16} (coordinate-order contours)) has exactly one mean in $G$, within $\frac1{40}$, the $W_x$ of
equations~\YMMref{(37)--(39)}{20} lie in $D$ within $\frac1{14}$, and $\mathcal M(U)(c)=\exp(L^{-3}\sum_x\log
W_x)\,U[c]\in G$ is real-analytic in $U$. This holds for every compact connected $G$; with $\omega_M$ in the place
of $\min\{\omega_M,r_G\}$ it fails for every $G$ with $r_G<\omega_M$. Let now the centre of $\mathfrak g$ be $0$.
Then $Z(G)$ is finite and $\min\{1,r_G\}=z_G$, so that the statement holds with $\omega_M$ in the place of
$\min\{\omega_M,r_G\}$ if and only if $\omega_M\le z_G$. At $G^c$-valued fields the same formula is holomorphic in
the matrix entries where the contour quotients and the $W_x$ lie in $D^c$; without this condition it leaves $G^c$
for every such $G$ with $z_G<1$.
\item[(ii)] The averaging map of equations~\YMMref{(37)--(39)}{20} transferred to $G$, its three cases read with ``$u\in
G^\pi$'' and $\mathrm{Log}\,e^{iY}=iY$ for $Y\in\mathfrak g$, $|Y|<\pi$, is $G$-valued and Borel on all fields,
depends on $U$ on the double block of $c$ alone, is compatible with the reversal of bonds, and $\mathcal
M(U^h)(c)=h(c_-)\mathcal M(U)(c)h(c_+)^{-1}$ with the same case; so $\mathcal M^q(U^h)=(\mathcal M^qU)^h$. As a map
it needs no condition on $z_G$. For $SU(N)$ with $\omega^G=\min\{\omega_M,2\sin\frac\pi N\}$ the averaging map and
the lemmas of the $SU(2)$ construction on it, Lemma~\YMMref{14.2}{62} among them, hold with their constants but for
one clause: the second assertion of Lemma~\YMMref{14.2}{62}, that $L^{-3}\sum_x\log W_x\in i\mathfrak g$ and
$\mathcal M(W)(c)\in G$ for every $G$-valued $W$ of its domain, holds for every $G$ with $\rho_G\ge\frac\pi3$, that
is $r_G\ge1$, and fails where $\rho_G<\frac\pi3$, that is $r_G<1$, so for $SU(N)$ exactly when $N\ge7$. It holds
for every $G$ with the proviso that every eigen-angle of each $W_x$ is below $\rho_G$ in modulus, or that $W_x$ is
a product of exponentials $e^{i\gamma_j}$, $\gamma_j\in\mathfrak g^c$, with $\sum_j|\gamma_j|<\ln2$; the proviso is
met under the averaging radius.
\item[(iii)] If $|\partial U(p)-\mathbf 1|\le\alpha_0L^{-2k}$ at all bare plaquettes and $\alpha_0\le
c(G)=\min\{3.9\cdot10^{-9}L^{-2},\frac14r_G\}$, then for $0\le q\le k$ the plaquette variables of $\mathcal M^qU$
are within $2\alpha_0L^{2(q-k)}$ of $\mathbf 1$, $\mathcal M^kU(c)U(\Gamma_c)^{-1}$ is a product of exponentials
$e^{iA_j}$, $A_j\in\mathfrak g$, of total size at most $12\alpha_0$, so that it lies in $G^\pi$, and $|\mathcal
M^kU(c)-U(\Gamma_c)|\le12\alpha_0$, $\Gamma_c$ the straight bare path from $c_-$ to $c_+$. The same holds for every
$G$-valued averaging map that is the formula of (i) on small fields, and in a local form, for fields given beneath
the double block of $c$ alone.
\item[(iv)] Let $G=\tilde G/\mathsf Z$ with $\tilde G$ simply connected. On a torus a field as in (iii) has a
sector $\zeta(U)\in\mathsf Z^3$, the class of the $\mathsf Z$-valued cocycle of the central parts of the plaquette
variables of a lift of $U$ to $\tilde G$, and $\mathcal M^kU$ lies in the sector of $U$. More generally, a coarse
field $V$ with $|V(c)-U(\Gamma_c)|\le\epsilon$ over a fine field with plaquettes within $\kappa$ of $\mathbf 1$
lies in the sector of $U$ as soon as $\mathcal L^2\kappa+4\epsilon<r_G$, $\mathcal L$ the ratio of the spacings; in
angles the threshold is $\pi$, and it is sharp for $SO(3)$. Fields whose block averages are all $\mathbf 1$ lie in
the trivial sector, and the block averages alone fix the sector.
\item[(v)] Let $1\le j\le K$, let $\mathsf Y$ be the two end points of a bond of level $j$ and $\mathsf
Y_i=B^{j-i}(\mathsf Y)$, $0\le i\le j$, the sites of level $i$ in their blocks of order $j-i$; let $U_0$ be
$G$-valued on the bare bonds with both end points in $\mathsf Y_0$, with $|\partial U_0(p)-\mathbf 1|\le\alpha
L^{-2j}$ at the bare plaquettes with all corners in $\mathsf Y_0$, where $0<\alpha\le c^{\rm g}_2$, and let $A_0$
be $\mathfrak g$-valued on these bonds with $L^j\sup|A_0|<c_4$; here $c^{\rm g}_2\le\frac14\min\{\omega_M,r_G\}$
and the radius $c_4\le\frac14r_L$ of the averages over a box, $r_L$ the radius of the one step below, are numbers
of $L$ and $r_G$, so small that $\sum_lC_{\rm P}\varpi_l\le\ln2$ and $\sum_lC\lambda_l\le\ln2$ in the proof. Then,
uniformly in the number $j$ of steps, for $0\le i\le j$ the averages $\mathcal M^iU_0$ and $\mathcal
M^i(e^{iA_0}U_0)$ are formed on the bonds with both end points in $\mathsf Y_i$, by the formula of (i) at every
step, the plaquette variables of $\mathcal M^iU_0$ lie within $2\alpha L^{-2(j-i)}$ of $\mathbf 1$, and there are
$\mathfrak g$-valued $A_i$ and $f_i:\mathsf Y_i\to G$, $f_0=\mathbf 1$, with
\[\mathcal M^i(e^{iA_0}U_0)(c)=f_i(c_-)\,e^{iA_i(c)}\,(\mathcal M^iU_0)(c)\,f_i(c_+)^{-1},\qquad|A_i|\le2L^i\sup|A_0|,\]
$f_{i+1}(y)=f_i(y)e^{i\sigma_i(y)}$ with $\sigma_i(y)\in\mathfrak g$, $|\sigma_i(y)|\le\kappa_6\sup|A_i|$, and
$f_i\in G^\pi$ with $|\frac1i\log f_i|\le6L^i\sup|A_0|$, so that the bound $|f_i-\mathbf 1|\le C_vL^i\sup|A_0|$ of
the averages over a box holds with $C_v=6$. One step is governed by the numbers $r_L=\frac1{1600L^2}$,
$\varpi_L=\frac12\min\{\omega_M,r_G\}$ and $\kappa_6=2L$: if $W_0$ is $G$-valued on the double block of a bond $c$
with plaquette variables within $\varpi_L$ of $\mathbf 1$, and $a$ is $\mathfrak g$-valued there with
$\sup|a|<r_L$, then the contour variables of $W=e^{ia}W_0$ and of $W_0$ are formed, $|W(y,x)W_0(y,x)^{-1}-\mathbf
1|<1$, and its logarithm lies in $i\mathfrak g$ with $|\frac1i\log(W(y,x)W_0(y,x)^{-1})|\le\kappa_6\sup|a|$. The
condition $\varpi_L<r_G$ cannot be spared. At complex arguments there is a number $a_{\rm c}>0$ of $L$ and the pair
such that, at every real background $U_0$ beneath a bond $c$ of level $q$ whose bare plaquette variables lie within
$a_{\rm c}L^{-2q}$ of $\mathbf 1$, and for $\mathfrak g^c$-valued $a$ with $L^q\sup|a|<a_{\rm c}$, $\mathcal
M^q(e^{ia}U_0)(c)=e^{iQ}\,\mathcal M^qU_0(c)$ with $Q\in\mathfrak g^c$ holomorphic in $a$, in $\mathfrak g$ for
real $a$, and $|Q|\le167\,L^q\sup|a|$, uniformly in $q$; and $|\frac1i\log\mathcal M^q(e^{ia})|\le9.2\,L^q\sup|a|$,
uniformly in $q$ and $N$. At the trivial background, for $\xi=L^{-k}$ and a $\mathfrak g^c$-valued form $f$ beneath
the double block of $c$ with $\sup|f|<4\cdot10^{-7}$, the averages $\mathcal M^j(e^{i\xi f})$, $0\le j\le k$, are
$e^{ia_j}$ with $\mathfrak g^c$-valued $a_j$, holomorphic in $f$, real for real $f$, and
$\sup|a_j|\le8L^{j-k}\sup|f|$. For $SU(3)$ no number changes.
\end{enumerate}
\end{lemma}

\begin{proof}[Sketch of proof]
(i) A contour quotient is a product of at most $\frac98L^2$ conjugated plaquette variables, those of the area
between two shortest paths in a block, so it lies within $(1+\omega_M)^{9L^2/8}-1\le e^{9/1600}-1\le\frac1{160}$ of
$\mathbf 1$. A plaquette variable within $r_G$ of $\mathbf 1$ is $e^{iF}$ with $F\in\mathfrak g$,
$|F|\le\frac\pi2|\partial U(p)-\mathbf 1|$, and so is each of its conjugates by $G$; so the quotient is $e^{iA}$
with $A\in\mathfrak g$, $|A|\le\frac{9\pi}{16}L^2\omega_M<\frac1{100}$, and lies in $D$ (Lemma~\ref{C-368-6}(ii)).
Lemma~\ref{C-368-6}(iv) gives exactly one unitary mean, within $4d\le\frac1{40}$, and its map $\Phi$ keeps
$i\mathfrak g$, each $V_iV_1^{-1}e^{-Y}$ with $Y\in i\mathfrak g$, $|Y|\le\frac18$, being a product of two
exponentials of elements of $i\mathfrak g$ of total size below $\ln2$ (Lemma~\ref{C-368-6}(iii)); so the fixed
point lies in $i\mathfrak g$ and the mean in $G$. $P_x=U(\Gamma)U([x,x'])U(\Gamma')^{-1}U[c]^{-1}$ takes
$\frac{25}8L^2$ plaquettes and lies within $e^{1/64}-1\le\frac1{60}$; then $|W_x-\mathbf
1|\le\frac{6541}{96000}<\frac1{14}$, and $W_x$, the product of $P_x$ with conjugates of the quotients of two
contour variables by their means, is again $e^{iB}$ with $B\in\mathfrak g$ small, so it lies in $D$. Holomorphy
holds on the open sets of matrices cut out by the conditions with $D^c$, on which $\mathrm{av}$ and $\log$ are
holomorphic. If $r_G<\omega_M$, one bond set to an element $u_*\in G\setminus G^\pi$ with $|u_*-\mathbf 1|=r_G$
gives a contour family without a mean in $G$, as in (v). Let the centre of $\mathfrak g$ be $0$. A central
$z\ne\mathbf 1$ lies outside $G^\pi$, since $z=e^{iY}$ with $Y\in\mathfrak g$, $|Y|<\pi$, would give
$e^{i\,\mathrm{ad}\,Y}=1$ with the eigenvalues of $\mathrm{ad}\,Y$ below $2\pi$ in modulus, so $\mathrm{ad}\,Y=0$
and $Y=0$; and an element $u\in G\setminus G^\pi$ with $|u-\mathbf 1|<1$ has the central factor $z(u)\ne\mathbf 1$,
with $|z(u)-\mathbf 1|\le|u-\mathbf 1|$ (Lemma~\ref{C-368-6}(iii)); so $\min\{1,r_G\}=z_G$. If $z_G<\omega_M$, one
bond set to a central $z$ with $|z-\mathbf 1|=z_G$ gives the contour family $(z,\mathbf 1)$, which has no mean in
$G$ (Lemma~\ref{C-368-6}(iv)).

(ii) The three cases are read with $G^\pi$, on which $\mathrm{Log}$ inverts the chart of Lemma~\ref{C-368-6}(ii);
each is covariant, as $\mathrm{av}(h[V]h')=h\,\mathrm{av}([V])h'$ and
$\mathrm{Log}(huh^{-1})=h(\mathrm{Log}\,u)h^{-1}$; the sets on which the cases apply are open, closed and Borel,
and the map is continuous on each. The domain of Lemma~\YMMref{14.2}{62} asks $|W_x-\mathbf 1|<1$, so the
eigen-angles of $W_x$ are below $\frac\pi3$ in modulus; if $\rho_G\ge\frac\pi3$, then $\log W_x\in i\mathfrak g$ by
the definition of $\rho_G$ and Lemma~\ref{C-368-6}(ii), and under the proviso the same holds by
Lemma~\ref{C-368-6}(ii), (iii). If $\rho_G<\frac\pi3$, there is $V\in G$ whose eigen-angles are at most $\rho_G$ in
modulus and whose hermitian logarithm $Y$ does not lie in $\mathfrak g$, for $SU(N)$ the element $e^{2\pi
i/N}\mathbf 1$; with $V$ on one bond of the corridor of $c$ other than the pivot bond and $\mathbf 1$ at the
others, $L^{-3}\sum_x\log W_x=iL^{-2}Y\notin i\mathfrak g$.

(iii) We never take the logarithm of an element that is merely near $\mathbf 1$: every object of a step is written
as a product of exponentials $e^{iA_j}$, $A_j\in\mathfrak g$ small, and $\sum_j|A_j|$ is added up
(Lemma~\ref{C-368-6}(ii)). One step gives $\mathcal M(W)(c)=e^{iY}W[c]$ with $Y\in\mathfrak g$,
$|Y|\le\kappa_3\max_p|\partial W(p)-\mathbf 1|$, $\kappa_3<6L^2$; the plaquette variables of an average are within
$L^2\varpi+C_0\varpi^2$ of $\mathbf 1$ if those beneath are within $\varpi$, $C_0$ a number of $L$; and the
corrections sum over the levels as geometric series, with the factor $\frac{2\kappa_3}{L(L-1)}<11.99$; the number
$c(G)$ reads of the pair only $r_G$.

(iv) If $\sum_j|X_j|<\pi$ in the norm of $\mathbb C^N$, then $\prod_je^{i\tilde X_j}$ is the canonical lift of
$\prod_je^{iX_j}$; by (iii) the plaquette variables of every average are such products of small exponentials of
those beneath, so the central parts of the plaquette variables of the lift multiply under blocking. A field with
trivial block averages is a gauge transform of $e^{ia}$, $4\sup|a|<\pi$ (Proposition~\ref{C-368-27}), whose lift
$e^{i\tilde a}$ has central parts $\mathbf 1$. Four bonds turned by $\frac\pi4+\delta$ about one axis around one
plaquette of $SO(3)$ give a field near the trivial one in a twisted sector.

(v) One step, relative to the background. On the double block of $c$ a tree gauge, the comb of
Lemma~\YMMref{14.4}{64}, writes $W_0=e^{iA_0}$ and $W=e^{ia}W_0=e^{iA}$ with $\mathfrak g$-valued $A_0$, $A$ of size
at most $2\pi L\varpi+\sup|a|$ and $\sup|A-A_0|\le1.01\sup|a|$ (Lemma~\ref{C-368-6}(ii)). In this gauge every
contour variable, every $W_x$ and the average are exponentials $e^{iX}$ of elements $X\in\mathfrak g$ formed from
$A$ by means and by logarithms of products of small exponentials, and the Lipschitz bound of the remainder of
Campbell and Hausdorff, $|F(X)-F(X')|\le4S\sum_j|X_j-X_j'|$ for $F(X)=\log(e^{X_1}\cdots e^{X_r})-\sum_jX_j$ and
$S=\sum_j\max\{|X_j|,|X_j'|\}\le\frac14$, compares each of them with its linear part, a straight average with real
weights. Let $\xi(y,x)=\frac1i\log(W(y,x)W_0(y,x)^{-1})$, so that $|\xi|\le\kappa_6\sup|a|$, and let
$\sigma(y)=L^{-3}\sum_{x\in B(y)}\xi(y,x)$. Then $\mathcal M(W)(c)=e^{i\sigma(c_-)}e^{ia'(c)}\mathcal
M(W_0)(c)e^{-i\sigma(c_+)}$ with $a'(c)\in\mathfrak g$: in the linear part the means of $A-A_0$ over the contours
from $c_\pm$ cancel against $\sigma(c_\pm)$, and what is left is $L^{-3}\sum_{x\in B(c_-)}(A-A_0)([x,x'])$, a sum
over the $L$ bonds of each $[x,x']$, whence the exact factor $L$; the rest is bounded by the Lipschitz bound, so
that $|a'(c)|\le L\sup|a|\,(1+C\lambda)$ with $\lambda$ a number of $L$ times $\varpi+\sup|a|$. In the same way the
plaquette variables of the average lie within $L^2\varpi(1+C_{\rm P}\varpi)$ of $\mathbf 1$ if those beneath lie
within $\varpi$, $C_{\rm P}$ a number of $L$. Over the levels the plaquette bounds $\varpi_i=2\alpha L^{-2(j-i)}$
and the numbers $\lambda_i$, a number of $L$ times $\varpi_i+2c_4L^{i-j}$, decrease geometrically in $L^{-(j-i)}$,
with $\sum_{l<j}\varpi_l\le\frac{2c^{\rm g}_2}{L^2-1}$ and $\sum_{l<j}2c_4L^{l-j}\le\frac{2c_4}{L-1}$. At level $i$
the one step applies, as $\varpi_i\le2c^{\rm g}_2\le\varpi_L$ and
$\sup|A_i|\le2L^i\sup|A_0|<2c_4L^{i-j}\le\frac12r_L$; and the two conditions on $c^{\rm g}_2$ and $c_4$,
$\sum_lC_{\rm P}\varpi_l\le\ln2$ and $\sum_lC\lambda_l\le\ln2$, keep the products $\prod_{l<i}(1+C_{\rm
P}\varpi_l)$ and $\prod_{l<i}(1+C\lambda_l)$ at most $2$, as $1+x\le e^x$; this closes the induction on $i$ and
gives $|A_i|\le2L^i\sup|A_0|$ and the plaquettes within $2\alpha L^{-2(j-i)}$. Further $f_{i+1}=f_ie^{i\sigma_i}$
with $\sum_{l<i}|\sigma_l|\le\kappa_6\sum_{l<i}2L^l\sup|A_0|\le6L^i\sup|A_0|<\pi$, so $f_i\in G^\pi$
(Lemma~\ref{C-368-6}(ii)), and $|f_i-\mathbf 1|\le|\frac1i\log f_i|$, as $|e^{i\theta}-1|\le|\theta|$ for real
$\theta$. At complex arguments the same recursion runs with $\mathfrak g^c$-valued exponents: the bound of Campbell
and Hausdorff holds for complex matrices, the products $e^{Z_1}\cdots e^{Z_q}$ with $Z_k\in\mathfrak g^c$,
$\sum_k|Z_k|<\ln2$, have their logarithms in $\mathfrak g^c$ (Lemma~\ref{C-368-6}(iii)), and the averages,
composites of exponentials, products, means and logarithms, are holomorphic in $a$; this gives $Q\in\mathfrak g^c$
and the bounds with $167$, $9.2$ and $8$, uniformly in $q$. If $\varpi_L\ge r_G$, one bond set to $u_*\in
G\setminus G^\pi$ with $|u_*-\mathbf 1|=r_G$ gives two contours without a mean in $G$.
\end{proof}

\begin{lemma}[Regular fields and gauges]\label{C-368-8}
\begin{enumerate}
\item[(i)] For every $G$ the gradient of $\xi^{-1}\sum_pa_p$ for the pairing of Definition~\YMMref{9.1}{37} is
$J=D^{\mathrm T}F$, at $G$- and at $G^c$-valued fields.
\item[(ii)] Let $\kappa'$ be the number of the pair of Definition~\ref{C-368-2}. Then
$4|(1-P)H_V|\le\kappa'|V-\mathbf 1|^3$ for $V\in D$ with $|V-\mathbf 1|\le z_G$;
$|(1-P)H_V|\le(\frac\pi3-\frac{\sqrt3}2)|1-P|\,|V-\mathbf 1|^3$ for every $V\in D$, $|1-P|$ the norm of $1-P$ for
$|\cdot|$, so that $\kappa'\le4(\frac\pi3-\frac{\sqrt3}2)|1-P|$; at $SU(2)$, $\kappa'=0$; and at $SU(N)$, for every
$N$, $4|(1-P)H_V|\le(2-\pi/\sqrt3)|V-\mathbf 1|^3$ for every $V\in D$, so that $\kappa'\le2-\pi/\sqrt3$,
$|(1-P)H_V|\le\frac15|V-\mathbf 1|^3$ and $|PH_V|\le(1+\frac15|V-\mathbf 1|^2)|V-\mathbf 1|$ on $D$. For every $G$,
if the plaquette variables at the four plaquettes of a bond $b$ lie in $D$ within $\rho\le z_G$ of $\mathbf 1$,
then
\[\eta|(D^{\mathrm T}\partial U)(b)|\le\rho^2+\eta^3|J(b)|+\kappa'\rho^3,\qquad\eta^3|J(b)|\le\eta|(D^{\mathrm T}\partial U)(b)|+\kappa'\rho^3 .\]
For every $G$ and every $r$, $\mathcal U'(r)\subset\mathcal U(\kappa_Gr)$, the classes being those of Definition~\ref{def:group-regular-classes}; for every $G$ and $r\le z_G$, $\mathcal U(r)\subset\mathcal
U'(r+r^2+\kappa'r^3)$; Moreover, let $\alpha_0,\alpha_2>0$ and let $E$ be a real field whose plaquette variables obey
$|\partial E(p)-\mathbf1|\le(\alpha_0+3\alpha_2)\eta^2\mathsf L(p)^{-2}$
and whose current at every full bond obeys
$|J(E)(b)|\le s\mathsf L(b)^{-3}$, where
$s=2\alpha_0+3\alpha_2$ and $\alpha_0+3\alpha_2<z_G$.
Then $E\in\mathcal U'(s+s^2+\kappa's^3)$, and hence
$E\in\mathcal U'(2s)$ if $s(1+\kappa's)\le1$.
For $SU(2)$ the two hypotheses are supplied by the action expansion
under equations~\YMMref{(466)}{238}--\YMMref{(467)}{238}.
Their general-group counterparts are quantitative inputs in (P),
wherever this conversion is applied to a solved field. The bound $\mathsf L(b)^3|J_{\Omega_0}(U_0)(b)|\le3\alpha_0$
of Proposition~\YMMref{30.1}{236} at an inner bond that is not full holds for every $G$ with $3|P|\alpha_0$ in the place
of $3\alpha_0$. Let now $G=SU(N)$ and $r\le z_G$. Then $\mathcal U'(r)\subset\mathcal U((1+\kappa'r^2)r)$; at
$SU(2)$, where $\kappa_G=1$, $\kappa'=0$ and $z_G=1$, the inclusions are the two inclusions of the $SU(2)$
construction. Under the same curvature and current hypotheses, the membership in
$\mathcal U'(2s)$ holds for $s\le\frac13$; and the bound at an inner bond that is not full holds with $6\alpha_0$, with $3\alpha_0$ exactly
when $N\le7$, and for every $N$ with $3\alpha_0$ if $\alpha_0\le2\sin\frac\pi N$, so if $\alpha_0\le z_G$. If
$|\partial U(p)-\mathbf 1|\le\rho\le1$, $\rho<z_G$, at the four plaquettes of a bond $b$, then $J(U)(b)\in\mathfrak
g$ and
\[|(D^{\mathrm T}\partial U)(b)-i\eta^2J(U)(b)|\le\eta^{-1}\rho^2\bigl(1+\tfrac35\rho^2\bigr)^{1/2}\le\tfrac{13}{10}\eta^{-1}\rho^2 .\]
The inclusion $\mathcal U'(r)\subset\mathcal U(r)$ is false for $N\ge3$ at every small $r$, and the guard $\rho<z_G$ cannot be dropped.
\item[(iii)] The following quantitative cube-gauge assertion is an
additional local analytic input when it is used for a general closed
$G\subset U(N)$ of positive dimension. It requires constants $c,B>0$
depending only on the pair: if $Q$ is a cube
of $l^3$ blocks at level $k$, the concentric cube $\tilde Q$ of side $l+4$ is shorter than half the torus,
$l\alpha\le c$, and $|\partial U(p)-\mathbf 1|$ and the covariant divergence of $\partial U$ are below
$\alpha\eta^2$ on $\tilde Q$, then on $\tilde Q$ the field is gauge equivalent to $e^{i\eta A}$ with $|A|,|\nabla
A|,|d^*dA|,|\Delta A|<Bl\alpha$ on $Q$.
For $SU(2)$ all four bounds hold. The first three are
Theorem~\YMMref{F}{199}; the Laplacian bound follows from the
same construction, as proved below, after enlarging $B$.

The comb gauge of a box of side $s$ (Lemma~\YMMref{14.4}{64}) gives $|W^h(b)-\mathbf 1|\le\min\{2s,|x-y|_1\}\varpi$,
and $W^h=e^{ia^0}$ with $\sup|a^0|\le\pi s\varpi$ if besides $2s\varpi<2$ also $\varpi<z_G$; at $SU(N)$ this guard
bites from $N=7$ on. The form of that clause with some gauge of the box in the place of the comb, that $2s\varpi<2$
alone gives a gauge in which $W^h=e^{ia^0}$ with $a^0$ in $\mathfrak{su}(N)$ and $\sup|a^0|\le\pi s\varpi$, holds
for $SU(N)$ exactly when $N\le13$. The quantitative regular-cube assertion for $SU(N)$ uses the input
just stated; the comb estimate alone does not supply its derivative bounds.
\item[(iv)] For $SU(2)$ the following box-gauge estimate is Proposition~\ref{prop:uniform-box-gauge}. For other $G$ with centre-free Lie algebra it is an additional analytic input, required in (II): there is $C$ such that for every $s$ and $0<\kappa\le C^{-2}$
every $G$-valued field on $\{0,\dots,s\}^3$ with $|\partial U(p)-\mathbf 1|\le\kappa$ has a gauge transform within
$C\sqrt\kappa$ of $\mathbf 1$ at all bonds. The general-group assertion requires a uniform quantitative filling estimate. For a
closed $G$ that is not connected this holds if and only if it holds for the identity component, with
$C(G)=\max\{C(G_0),(2/r_G)^{1/2}\}$. It fails for every $G$ whose Lie algebra has a centre, and on periodic
lattices for $SO(3)$.
\end{enumerate}
\end{lemma}

\begin{proof}[Sketch of proof]
(i) For $\theta\in\mathfrak g^c$ and every $X\in G^c$, near $\mathbf 1$ or not, $\frac
d{dt}a(e^{it\theta}X)|_{t=0}=\frac1N\mathrm{tr}(\theta H_X)=\frac1N\mathrm{tr}(\theta
PH_X)=\langle\theta,PH_X\rangle$, since $(1-P)M$ is trace-orthogonal to $\mathfrak g^c$ by the definition of $P$
(Definition~\ref{C-368-2}). A change $e^{it\theta}$ of one bond changes the plaquette variables of its four
plaquettes by conjugates of $e^{\pm it\theta}$ under transports, and $P\,\mathrm{Ad}(g)=\mathrm{Ad}(g)P$ for $g\in
G^c$, as $\mathrm{Ad}(g)$ keeps $\mathfrak g^c$ and the trace form; summing over the four plaquettes, with the
factors $\xi$ of Definition~\YMMref{9.1}{37}, gives $J=D^{\mathrm T}F$ with $F=\xi^{-2}PH_{\partial U}$.

(ii) For $V=e^{iY}\in D$, $Y\in\mathfrak g$, $t=|Y|<\frac\pi3$: $H_V=\sin Y$ and $(1-P)Y=0$, so
$(1-P)H_V=(1-P)(\sin Y-Y)$, of norm at most $|1-P|(t-\sin t)\le(\frac\pi3-\frac{\sqrt3}2)|1-P|\,|V-\mathbf 1|^3$,
as $|V-\mathbf 1|=2\sin\frac t2$ and $(t-\sin t)/(2\sin\frac t2)^3$ increases on $(0,\frac\pi3]$. For $V\in G$ with
$0<|V-\mathbf 1|<z_G$, $V\in D$ (Lemma~\ref{C-368-6}(iii)), so that $\kappa'$ is at most four times the supremum of
$|(1-P)H_V|/|V-\mathbf 1|^3$ over $D\setminus\{\mathbf 1\}$; and at $V=e^{iY}\in D$ with $|V-\mathbf 1|=z_G$ the
bound $4|(1-P)H_V|\le\kappa'|V-\mathbf 1|^3$ is the limit of those at $e^{isY}$, $s\uparrow1$, which lie in $G$
within $z_G$ of $\mathbf 1$. At $SU(2)$, $H_V$ lies in $\mathfrak g^c$. At $SU(N)$, $|1-P|=1$ and
$(1-P)H_V=\frac1N\mathrm{Im}\,\mathrm{tr}\,V\cdot\mathbf 1=\frac1N\sum_j(\sin\theta_j-\theta_j)\mathbf 1$, the
eigen-angles $\theta_j$ of $V$ being at most $t<\frac\pi3$ in modulus and summing to $0$. With $f(x)=x-\sin x$,
which is odd, convex on $[0,t]$ and concave on $[-t,0]$, the sum is extreme with a share $p\le\frac12$ of the
angles at $t$ and the rest equal, at $-\frac{pt}{1-p}$: $\frac1N|\sum_jf(\theta_j)|\le\max_{0\le p\le1/2}\psi(p)$,
$\psi(p)=pf(t)-(1-p)f(\frac{pt}{1-p})$. The share $p=\frac13$, which is extreme to third order in $t$, gives
$\psi(\frac13)=\frac23\sin\frac t2\,(1-\cos\frac t2)$, that is $4\psi(\frac13)=\frac{|V-\mathbf
1|^3}{3(1+\cos(t/2))}$, at most $0.1787\,|V-\mathbf 1|^3$. The supremum of $4\psi(p)/|V-\mathbf 1|^3$ over $p$ and
$t\le\frac\pi3$ is below $0.179$, so that $4|(1-P)H_V|\le0.179\,|V-\mathbf 1|^3$, below $(2-\pi/\sqrt3)|V-\mathbf
1|^3$. The bound with $2-\pi/\sqrt3$ comes from the tangent of $\psi$ at $p=\frac13$: $\psi$ is concave,
$\psi''(p)=-t^2(1-p)^{-3}\sin\frac{pt}{1-p}$; with $s=\frac t2$, $\psi'(\frac13)=3s\cos s-\sin s-\sin2s\le0$, as
$2\sin s+\tan s\ge3s$; so $\psi\le\psi(\frac13)-\frac13\psi'(\frac13)=\sin s-s\cos s$ on $[0,\frac12]$, and $4(\sin
s-s\cos s)/(2\sin s)^3$ increases in $s$, its derivative having the sign of $x(2+\cos x)-3\sin x\ge0$ at $x=2s$, to
$2-\pi/\sqrt3$ at $s=\frac\pi6$. At a bond $b$, $S=\eta(D^{\mathrm T}\partial U)(b)$ is the sum of the four
plaquette variables transported to $b$, two with the sign $+$ and two with the sign $-$, and $S=R+iT$ with $R$, $T$
the like signed sums of the transported $c(\partial U(p))$ and $H_{\partial U(p)}$, both hermitian at real fields,
and $PT=\eta^3J(U)(b)$, $P$ commuting with the transports. The four constants $\mathbf 1$ cancel and each
$c(\partial U(p))-\mathbf 1$ lies between $-\frac12\rho^2$ and $0$, so $|R|\le\rho^2$; and
$|(1-P)T|\le\kappa'\rho^3$, the four plaquette variables lying in $D$ within $\rho\le z_G$ of $\mathbf 1$. Hence
$|S|\le|R|+|PT|+|(1-P)T|\le\rho^2+\eta^3|J(b)|+\kappa'\rho^3$, and
$\eta^3|J(b)|=|PT|\le|T|+|(1-P)T|\le|S|+\kappa'\rho^3$. At $G$-valued fields, for every $G$, $T$ is hermitian and
$P$ commutes with taking adjoints, $\mathfrak g^c$ being closed under them, so $PT\in\mathfrak g$ and
$|PT|\le\kappa_G\|PT\|\le\kappa_G\|T\|\le\kappa_G|T|\le\kappa_G|S|$, $P$ being orthogonal for
$\langle\cdot,\cdot\rangle$ (Definition~\ref{C-368-2}, Lemma~\ref{C-368-6}(i)); at a full bond of a field of
$\mathcal U'(r)$ this gives $|J(U)(b)|<\kappa_Gr\,\mathsf L(b)^{-3}$, and as $\kappa_G\ge1$, $\mathcal
U'(r)\subset\mathcal U(\kappa_Gr)$. For every $G$ and $r\le z_G$ the plaquette variables of the two classes lie
within $z_G$ of $\mathbf 1$, so in $D$ (Lemma~\ref{C-368-6}(iii)); at a full bond $b$ they are within
$\rho=r\eta^2\mathsf L(b)^{-2}$ of $\mathbf 1$ at the four plaquettes $p$ of $b$, as $\mathsf L(p)\ge\mathsf L(b)$,
each $p$ containing the end points of $b$ (Definition~\ref{def:group-regular-classes}); with $\eta\le\mathsf L(b)$ the first bound
gives $|(D^{\mathrm T}\partial U)(b)|<(r+r^2\frac\eta{\mathsf L(b)}+\kappa'r^3\frac{\eta^3}{\mathsf
L(b)^3})\eta^2\mathsf L(b)^{-3}\le(r+r^2+\kappa'r^3)\eta^2\mathsf L(b)^{-3}$ on $\mathcal U(r)$, which is the
second inclusion, and the second bound gives $|J(U)(b)|<(1+\kappa'r^2)r\,\mathsf L(b)^{-3}$ on $\mathcal U'(r)$,
which is the inclusion $\mathcal U'(r)\subset\mathcal U((1+\kappa'r^2)r)$. Under the stated curvature and current hypotheses the
plaquette variables at a full bond $b$ lie within
$\rho=(\alpha_0+3\alpha_2)\eta^2\mathsf L(b)^{-2}$ of $\mathbf 1$, so in $D$, and the current there is at most
$s\,\mathsf L(b)^{-3}$; the first bound gives $|(D^{\mathrm T}\partial U)(b)|\le(s+s^2+\kappa's^3)\eta^2\mathsf
L(b)^{-3}$, as $\alpha_0+3\alpha_2\le s$ and $\eta\le\mathsf L(b)$. As $\eta^3J=PT$ and $|PT|\le|P|\,|T|$, the
bound $3\alpha_0$ at an inner bond that is not full holds for every $G$ with $3|P|\alpha_0$, its proof bounding $T$
as at $SU(2)$, where $PT=T$; at $SU(N)$, $|P|\le2$. At $SU(N)$, $(1-P)T$ is a real multiple of $\mathbf 1$, so
$|R+i(1-P)T|^2\le|R|^2+|(1-P)T|^2\le\rho^4(1+\kappa'^2\rho^2)$, and $(D^{\mathrm T}\partial
U)(b)-i\eta^2J(U)(b)=\eta^{-1}(R+i(1-P)T)$ gives the display, as $\kappa'^2\le\frac35$. One excited bond with all
other variables $\mathbf 1$ refutes $\mathcal U'(r)\subset\mathcal U(r)$; the central plaquette variables $z\mathbf
1,\mathbf 1,\mathbf 1,z\mathbf 1$, $z=e^{2\pi i/N}$, have $J=0$ and the left member $2\eta^{-1}z_G$. So each later
proof that uses the bound of the regular class goes through with $\kappa_G\alpha_0$, and at $SU(N)$ under
$\alpha_0\le z_G$ with $(1+\kappa')\alpha_0$, each threshold on $\alpha_0$ read as a ceiling being divided by
$\kappa_G$, or by $1+\kappa'$.

(iii) For $SU(2)$, use the gauge constructed in the proof of
Theorem~\YMMref{F}{199}. We record the ordinary Laplacian bound
as well. In its box notation, put $\epsilon=s\alpha$,
$s=l+4\le5l$, and let $d(x)$ be the depth of a site in the box.
The smoothed section defines a quaternion-valued one-form $\omega$
and a smooth map $\Theta$ with
\[
 A=\eta^{-1}\Theta(\mathbf1+\eta\omega),\qquad
 |\omega|\le20\epsilon,\qquad
 |\nabla^\eta\omega|\le C_M\epsilon\quad(d\ge1).
\]
Its equation and the interior estimate in that proof give
\[
 |\Delta^\eta\omega_\mu(x)|
 \le55\epsilon+120\epsilon
       \max_{x'\sim x}\eta^{-1}|\omega_\mu(x')-\omega_\mu(x)|
 \le(55+C_M)\epsilon\quad(d(x)\ge1+\eta),
\]
provided $120\epsilon\le1$. Choose $c_\Theta$ to bound the
first two derivatives of $\Theta$ on
$|q-\mathbf1|\le1/2$. Taylor's formula at each of the six
neighbours, with the constant terms cancelling and the linear terms combining into the Laplacian of $\omega$, yields
\[
 |\Delta^\eta A_\mu(x)|
 \le c_\Theta|\Delta^\eta\omega_\mu(x)|
    +6c_\Theta\eta
       \max_{x'\sim x}
       \bigl(\eta^{-1}|\omega_\mu(x')-\omega_\mu(x)|\bigr)^2
 \le c_\Theta(55+C_M)\epsilon
       +6c_\Theta\eta C_M^2\epsilon^2.
\]
All arguments of $\Theta$ lie in this ball by
$20\eta\epsilon\le1/2$. Since $\eta\le1/13$ and $\epsilon$
has a fixed upper bound, the last expression is at most
$C\epsilon$. Bonds with an endpoint in $Q$ have base depth
at least $2-\eta\ge3/2>1+\eta$; hence the bound is at most
$5Cl\alpha$ there. Enlarging $B$ proves the fourth estimate
without a further hypothesis for $SU(2)$.

For general $G$, the quantitative assertion also requires the
corresponding derivative bounds for the retracted smoothed section;
it remains the local input stated in (iii).
For the elementary comb estimate: under
$\varpi<z_G$ the bond variables of the comb are products of at most $2s$ conjugated plaquette variables, with
logarithms in $\mathfrak g$ of size at most $\pi s\varpi$ (Lemma~\ref{C-368-6}(ii)); at $SU(N)$, $N\le13$, a box
with $z_G\le\varpi$ and $2s\varpi<2$ has $s=1$, $N\ge7$, or $s=2$, $N=13$, since
$2\sin\frac\pi{13}<\frac12<2\sin\frac\pi{12}$ and $\frac13<2\sin\frac\pi{18}$. For $s=1$, $7\le N\le13$ and
$\varpi<1$, plaquette variables within $\varpi$ of $\mathbf 1$ at the six faces of $\{0,1\}^3$ give $U^v=e^{iA}$
with $|A|\le6\arcsin\frac\varpi2\le\pi\varpi$, even near a central element; at $N=13$ and $\varpi=z_G$ this is
$\frac{6\pi}{13}=\frac\pi2\cdot\frac{12}{13}$, the least that a gauge can give. For $s=2$, $N=13$ and
$\varpi\in[2\sin\frac\pi{13},\frac12)$, a gauge of $\{0,1,2\}^3$ has $|A(b)|\le6\theta\le2\pi\varpi$ at its $54$
bonds, $\theta=2\arcsin\frac\varpi2$ being the largest eigen-angle of a plaquette variable. For $N\ge14$ one
plaquette variable $e^{2\pi i/N}\mathbf 1$, at the distance $2\sin\frac\pi N<1$, refutes the clause at $s=1$:
$e^{iA_1}\cdots e^{iA_4}=e^{\pm2\pi i/N}\mathbf 1$ with $A_j\in\mathfrak{su}(N)$ forces
$\sum_j|A_j|\ge2\pi\frac{N-1}N$, so $\max_j|A_j|\ge\frac\pi2\frac{N-1}N>2\pi\sin\frac\pi N$.

(iv) The $SU(2)$ proof uses a quantitative filling of boundary maps in the quaternionic sphere. For general simply connected $G$, avoidance of the singular set
$\{u:\dim\ker(\mathrm{Ad}\,u-1)\ge\mathrm{rk}\,G+2\}$ alone does not give a logarithm or extension with a uniform Lipschitz bound.
Such a quantitative filling bound, together with its lifting to quotients, is needed to deduce the asserted estimate; it is retained as part of the analytic input in (II).
For $G=U(1)$ the box estimate fails at constant field strength: let $U$ have the
plaquette angle $\kappa$ at every plaquette of one plane direction and $0$ at the others, as $U(\langle
x,x+e_2\rangle)=e^{i\kappa x_1}$ and $\mathbf 1$ at the other bonds does. Its plaquette variables are within
$\kappa$ of $\mathbf 1$. In a gauge with all bonds within $C\sqrt\kappa$ of $\mathbf 1$, $\kappa\le\frac14C^{-2}$,
the angles $a(b)$ of the bonds obey $|a(b)|\le\frac\pi2C\sqrt\kappa$, so $da=\kappa$ exactly at those plaquettes,
and summing over a square of side $l$ in that plane gives $\kappa l^2\le4l\cdot\frac\pi2C\sqrt\kappa$, which is
false for $l>2\pi C\kappa^{-1/2}$; so no $C$ serves for all $s$, and the same field strength in a central direction
gives the failure whenever $\mathfrak g$ has a centre. With the lemma fail the vanishing of invariant linear forms
(the constant-gauge argument following equation~\YMMref{(943)}{429}) and further statements, among them a step of the proof of
Proposition~\YMMref{42.1}{321}, so the theorems ask $\mathfrak g$ semisimple. A flat $SO(3)$-field in a twisted
sector of a torus is no gauge transform of $\mathbf 1$.
\end{proof}

\begin{lemma}[Energy of block averages and large fields]\label{C-368-9}
\begin{enumerate}
\item[(i)] Let $G\subset SU(N)$ be compact, $L\ge13$ odd, $\mathcal M$ a $G$-valued averaging map that equals the
formula of Lemma~\ref{C-368-7}(i) where $|\partial U(p)-\mathbf 1|\le\varpi'(L,r_G)\le\frac12r_G$, and
$e_*=e_L/(2\tau_G^2)$, where $e_L>0$ is a constant depending only on $L$ and the pair with
$(2e_L/L)^{1/2}\le\varpi'(L,r_G)$ and $(N-1)^{1/2}C_{\rm en}(L)(2e_L)^{1/2}\le\frac12(L^{1/2}-1)\ln2$, $C_{\rm
en}(L)$ being the number of $L$ in the cost of one step of the proof. For all $K$, tori, $1\le t\le K$, sets $S$ of
plaquettes of level $t$ and fields $U$, with $\mathcal N(S)$ the bare plaquettes beneath $S$, of all three pairs of
directions (the cones defined at the beginning of Section~\YMMref{15}{78}),
\[L^t\sum_{\mathcal N(S)}a_p(U)\ge\min\Bigl\{\frac1{4\tau_G^2}\sum_{p'\in S}|\partial\mathcal M^tU(p')-\mathbf 1|^2,\ e_*\Bigr\};\]
as $\tau_G^2\le N-1$, the same holds with $N-1$ in the place of $\tau_G^2$. No number above $\frac1{2\kappa_G^2}$
serves for $\frac1{4\tau_G^2}$, and $\frac14$ fails where $\kappa_G^2>2$, so at $SU(N)$ for $N\ge4$. At $SU(3)$,
$\frac14$ serves for every set $S$ in which no site is a corner of two plaquettes of $S$ of different pairs of
directions, with a threshold $e(L)>0$ in the place of $e_*$; and $\frac1{4(1+\delta)}$ serves for every $S$ and
every $\delta>0$, with a threshold of $L$ and $\delta$.
\item[(ii)] Let $\mathfrak g$ be semisimple, so that $G\subset SU(N)$ and (i) applies, $L\ge13$ odd, and let the
averaging radius $\min\{\omega_M,r_G\}$ stand for $\omega_M$. Assume the transferred witness and local-term
estimates in (II). Then the large-field chain holds for $G$ with the following prices. Its first link is
Proposition~\YMMref{44.4}{347} and Proposition~\YMMref{44.5}{348}: a failed plaquette test or a failed approximate
fluctuation test of a test cube (equations~\YMMref{(636)--(637)}{340}; conditions in $|\cdot|$ on plaquette variables
and on averaged contour variables) forces a plaquette variable or an averaged contour variable farther than the
numbers of the $SU(2)$ construction from $\mathbf 1$, or an averaged contour variable that is not formed.
The fluctuation-test implication is used with
$V_t=\mathcal M(V_{t-1})$ on the inner level-$t$ bonds of the tested region,
as required by Proposition~\YMMref{44.5}{348}; the successive block fields
satisfy this local averaging identity. Its second link is Lemma~\YMMref{65.2}{580}, on the witnesses of the cells created at one step, with ``$\notin
G^\pi$'' for ``$=-\mathbf 1$''; The local-deposit argument in the proof of Proposition~\YMMref{65.4}{582} uses its alternatives for a failed plaquette or fluctuation test, a failed density test, and a large or undefined stored bond. Its
third link is the local-deposit argument in the proof of Proposition~\YMMref{65.4}{582}, the lower bounds of the amount of action near a created cell, which hold
with the expressions of the $SU(2)$ construction, taken in constants of the pair, times $\tau_G^{-2}$. Here an energy pair $(c,e)$ means that the action beneath one level-$t$
plaquette $p'$ obeys
\[
 L^t\sum_{p\in\mathcal N(p')}a_p(U)
 \ge\min\{c|\partial\mathcal M^tU(p')-\mathbf1|^2,e\}.
\]
Taking $S=\{p'\}$ in (i) gives the pair
$(\frac1{4\tau_G^2},e_*)$; and the deposit constant
\[
 c_\star=5^{-3}\min\{\kappa_{\rm W},\kappa_{\rm P},\kappa_{\rm Q},\kappa_{\rm F}\}
\]
of
equation~\YMMref{(1429)}{582}, where
\[
 \kappa_{\rm W}=\tfrac12,\qquad
 \kappa_{\rm P}=\tfrac12c_{\rm P},\qquad
 \kappa_{\rm Q}=\tfrac12c_{\rm Q}(A_1/A_0)^2,\qquad
 \kappa_{\rm F}=\tfrac14\gamma_{\rm kp}\varkappa_{\rm w}^2(A_1/A_0)^2,
\]
formed from the four prices per created cell, so that each created cell of level $t$
gives the small factor $e^{-c_\star p_0(g_t)^2}$, is at least $\tau_G^{-2}$ times its expression in the $SU(2)$
construction, taken in constants of the pair.
\end{enumerate}
\end{lemma}

\begin{proof}[Sketch of proof]
(i) Sizes are followed in $|\cdot|$ and energies in $\|\cdot\|$, where $a=\frac12\|V-\mathbf 1\|^2$ exactly, and at
one plaquette $|V-\mathbf 1|^2\le2\tau_G^2a(V)$ (Definition~\ref{C-368-2}, Lemma~\ref{C-368-6}(i)). Let $E_j$ be
the energy $\sum a_p(\mathcal M^jU)$ over the plaquettes of level $j$ beneath $S$, of all three pairs of
directions, and $\varpi_j$ the largest $|\partial\mathcal M^jU(p)-\mathbf 1|$ among them. Where
$\varpi_j\le\varpi'(L,r_G)$ the averages are given by the formula of Lemma~\ref{C-368-7}(i), and one step costs,
between square roots of energies, the factor $L^{1/2}(1+(N-1)^{1/2}C_{\rm en}(L)\varpi_j)$, the exact $L^{1/2}$
carrying the comparison of two cutoffs. Let $L^tE_0<e_*$, the other case being plain. By induction on $j$,
$E_j\le2L^jE_0$: if this holds up to $j<t$, then $\varpi_i^2\le2\tau_G^2E_i\le4\tau_G^2L^iE_0<2e_LL^{-(t-i)}$ for
$i\le j$, so that $\varpi_i\le(2e_L/L)^{1/2}\le\varpi'(L,r_G)$, and the product of the factors up to $j$ is at most
$\exp\bigl((N-1)^{1/2}C_{\rm en}(L)(2e_L)^{1/2}\sum_{m\ge1}L^{-m/2}\bigr)\le\sqrt2$, which gives
$E_{j+1}\le2L^{j+1}E_0$. At $j=t$, $\sum_{p'\in S}|\partial\mathcal M^tU(p')-\mathbf
1|^2\le2\tau_G^2E_t\le4\tau_G^2L^tE_0$, which is the inequality. The constant $\frac1{2\kappa_G^2}$ is approached
by abelian fields $e^{i\theta Y}$ of small constant field strength in one plane along a direction $Y\in\mathfrak g$
with $|Y|=\kappa_G\|Y\|$: the averages of such a field are exact, $L^t$ times the bare energy beneath a coarse
plaquette, carried by $L^{3t}$ parallel bare plaquettes of angle $\varphi$, is $\frac12L^{4t}\varphi^2\|Y\|^2$ to
leading order, and the coarse plaquette variable is within $L^{2t}\varphi|Y|$ of $\mathbf 1$. At $SU(3)$ let $S$
have one pair of directions, and let $E_\parallel$ and $E_\perp$ be the sums of $a_p(U)$ over the bare plaquettes
beneath $S$ parallel and transverse to that pair. Below a threshold of $L$, $\sum_{p'\in S}a_{p'}(\mathcal M^tU)\le
L^t(E_\parallel+w_\perp E_\perp)$ with a weight $w_\perp\le C(L)\bigl(L^t\sum_{\mathcal
N(S)}a_p(U)\bigr)^{1/2}\le\frac12$. For $W=e^{i\epsilon A}$ with constant field strengths $F_{\mu\nu}=dA$ in
$\mathfrak g$, the defect of one step at one coarse plaquette of the pair $12$ contains the cubic term
$-\epsilon^3\frac{L^2-1}{12}L^4\langle F_{12},i[F_{13},F_{23}]\rangle$, which has either sign; it stands at every
coarse plaquette, coming from the variance $\frac{L^2-1}{12}$ of the transverse coordinate of the bare plaquettes
beneath, and it is paid by a small part of the energy of the transverse plaquettes beneath, whose weights stay
below $1$ over the $t$ steps, the errors summing as a geometric series in $L^{-1/2}$. If no site is a corner of two
plaquettes of $S$ of different pairs of directions, the cones of members of $S$ of different pairs are disjoint, so
the transverse energy spent for one pair is claimed by no plaquette of $S$, and $\sum_{p'\in S}a_{p'}(\mathcal
M^tU)\le L^t\sum_{\mathcal N(S)}a_p(U)$. For general $S$ the plaquettes beneath members of two pairs add at most
$2C(L)\bigl(L^t\sum_{\mathcal N(S)}a_p(U)\bigr)^{1/2}$ times $L^t\sum_{\mathcal N(S)}a_p(U)$, which below a
threshold of $L$ and $\delta$ is at most $\delta$ times it. With $a(V)\ge\frac14|V-\mathbf 1|^2$ on $SU(3)$ this
gives the two forms.

(ii) A test is a condition in $|\cdot|$ on plaquette variables and on averaged contour variables. A test passed at
all plaquettes of a double block puts the field under the averaging radius $\min\{\omega_M,r_G\}$, where every
contour family has exactly one mean, within $4d$ (Lemmas~\ref{C-368-6}(iv), \ref{C-368-7}(i)); by contraposition a
failed test forces a plaquette variable or an averaged contour variable farther than the numbers of the $SU(2)$
construction from $\mathbf 1$, or an averaged contour variable that is not formed: this is the first link, with the
numerals of the $SU(2)$ construction, the straightening of Lemma~\ref{C-368-7}(iii) holding for every closed $G$.
The geometric and running-coupling statements used in the local-deposit argument of Proposition~\YMMref{65.4}{582}, besides Lemma~\YMMref{65.2}{580} involve cubes,
regions and runs of couplings alone, read no group and hold for $G$ as written. The witness and local-deposit
estimates of Lemma~\YMMref{65.2}{580} and the local-deposit argument in the proof of Proposition~\YMMref{65.4}{582} are the second-part inputs assumed in (II);
their inputs ask that the terms of each level be bounded in norm by $E_0$ and the increments of the run by
$B_\sharp$, separately, and the bounds of a step add these sizes; and a family of local terms splits off from the
effective action of a large-field region, with a contribution bounded by a small factor times its norm times a
factor that decays exponentially in the size of the region.

The group enters the prices alone, where the local-deposit argument in the proof of Proposition~\YMMref{65.4}{582} turns the witnesses of
Lemma~\YMMref{65.2}{580} into lower bounds of the action: at one plaquette through
$a(V)\ge\frac12\tau_G^{-2}|V-\mathbf 1|^2$ (Lemma~\ref{C-368-6}(i)), and at averaged variables through (i), whose
pair $(\frac1{4\tau_G^2},\frac{e_L}{2\tau_G^2})$ takes the place of the pair $(\frac14,\frac{e_L}2)$ of the $SU(2)$
construction. So every price of the local-deposit argument in the proof of Proposition~\YMMref{65.4}{582} is the expression of the $SU(2)$ construction, in constants
of the pair, times $\tau_G^{-2}$, each of the four prices per created cell is at least $\tau_G^{-2}$ times its
expression, and so is $c_\star$. For simply connected $G$ with $\mathfrak g$ simple this chain is among the
statements of Definition~\ref{C-368-31}(II).
\end{proof}

\begin{lemma}[The Haar measure in the chart, the cut-off and determinants]\label{C-368-10}
\begin{enumerate}
\item[(i)] For $u_0\in G$ and Borel $F\ge0$, $\int_{G^\pi
u_0}F\,du=\mathrm{vol}(G)^{-1}\int_{|Y|<\pi}F(e^{iY}u_0)\,j_G(Y)\,dY$,
$j_G=\prod_\alpha(\frac{\sin(\alpha/2)}{\alpha/2})^2$ over the positive eigenvalues $\alpha$ of $\mathrm{ad}\,Y$;
$|\log j_G(Y)|\le(n-\mathrm{rk}\,G)(-\log\frac{\sin|Y|}{|Y|})$ on $|Y|<\pi$, and $|\log j_G(Y)|\le0.173\,n|Y|^2$ on
$|Y|\le1$; the quadratic part of $\log j_G$ is $-\frac1{24}\sum_iC^{(i)}_A\|P_iY\|^2$. The identities of the
$SU(2)$ construction that integrate in the chart hold for $G$ with the indicator of $e^{iA}e^{iH}\in G^\pi$ in the
integrand, $A,H\in\mathfrak g$ being the two exponents of the product that such an identity integrates in the
chart. $G\setminus G^\pi$ is a null set if and only if $r_G=2$; otherwise it has inner points.
\item[(ii)] Let $\vartheta_{\rm co}$ be the profile of the cut-off of the bond shift (the cutoff preceding Lemma~\YMMref{53.3}{444}),
smooth, $0$ on $(-\infty,\frac43]$ and $1$ on $[\frac53,\infty)$, with the absolute constant $\|\vartheta_{\rm
co}'\|_\infty\ge3$, and $\delta\in(0,1]$ its scale. If $w=N^{1/2p}<\frac54$ and $\lambda=\frac w{5-4w}$, then, with
$\vartheta_G(s):=\vartheta_{\rm co}(\frac43+\lambda(s-\frac43))$, the function $\vartheta_G(\nu(Y)/\delta)$ is
smooth and invariant, $0$ on $|Y|\le\frac43\delta$ and $1$ on $|Y|\ge\frac53\delta$, and
$\|\vartheta_G'\|_\infty=\lambda\|\vartheta_{\rm co}'\|_\infty\ge3$; the cut-off $\vartheta_{\rm co}(|Y|/\delta)$
itself is not differentiable for rank at least $2$.
\item[(iii)] If $T_j=L^3\mathrm{Ad}(g_j)\,\mathsf m\,\mathrm{Ad}(g'_j)^{-1}$, $g_j,g'_j\in G$, $|T_j-1|\le r<1$
($j=1,2$), then $\mathsf m$ is invertible, $|\mathrm{tr}\log T_j|\le-n\log(1-r)$ and $\mathrm{tr}\log
T_1=\mathrm{tr}\log T_2$. For every compact connected $G$ the elimination factor $\mathsf m_c(V)$ of
Lemma~\ref{C-368-10}(iii) is formed, as a complex linear map of $\mathfrak g^c$ into itself, at every $V$ of the
domain of the complex average at $c$: the $G^c$-valued $V$ on $\mathsf P_c$ whose contour families have diameter
below $\frac1{40}$ and whose contour quotients and variables $W_x$ lie in $D^c$; at $G$-valued $V$ it maps
$\mathfrak g$ into $\mathfrak g$. For $\mathfrak g^c$-valued $a$ on $\mathsf P_c$ with $\sup|a|<\frac34r_*$,
\[|\mathsf m_c(e^{ia})-L^{-3}|\le L^{-1}C_S\sup|a|,\]
where $0<r_*\le\frac1{600L}$, $r_1=\frac{r_*}{40}$, $m_1=\max\{\frac1{r_1},\frac4{r_1^2}\}$ and
$C_S\ge\frac{8m_1}{3r_*}$: numbers of $L$, free of $k$, $N$ and $G$.
\item[(iv)] Let $n\ge2$. If $T$ acts on an $n$-dimensional space and $|1-T|\le r_n:=\sin\frac\pi{2n}$, then $\det
T\notin(-\infty,0]$, $\mathrm{Log}\det T=\mathrm{tr}\log T$ and $|\mathrm{Log}\det T|\le-n\log(1-r_n)=:\rho_n$,
with $\frac\pi2<\rho_n<\frac52$; and $|\mathrm{Log}\det(1+E)|\le c_n|E|\le\rho_n$ for $|E|\le r_n$, where
$c_n:=\rho_n/r_n=-\frac n{r_n}\log(1-r_n)$, which is below $5$ for $n\le3$ and above $5$ for $n\ge4$.

The radius $\frac12$ keeps $\det T$ off $(-\infty,0]$ if and only if $n\le5$. The following two forms of the elimination estimate use this radius; their three-dimensional counterparts are the proof of Proposition~\YMMref{50.2}{386} and the argument preceding equation~\YMMref{(1393)}{572}. Let $K:=LC_SC''_{\mathrm{ia}}$, with $C_S$, $r_*$ as in (iii) and
$0<c_\sharp\le4\cdot10^{-7}$ and $C''_{\mathrm{ia}}\ge8$, using the last, trivial-background estimate of Lemma~\ref{C-368-7}(v),
read at the trivial background. In the chain of the localised Gaussian family, whose radius $a_{\mathsf N}$ is taken with $0<a_{\mathsf
N}\le\min\{\frac12,\frac{Lc_\sharp}8,\frac{Lr_*}{8C''_{\mathrm{ia}}},\frac{r_n}{4K}\}$, the normalised factor $T$
stays within $r_n$ of $1$ and the function $z_c$ of the elimination factor obeys $|z_c-2n\log L|\le\rho_n$, for
every $G$. The corresponding linear estimate is $|z_c-2n\log L|\le C_za$ on fields of size $a$ for every
$C_z\ge\max\{1,4c_nK\}$ and every $a$ with
\[0<a\le\min\Bigl\{\frac12,\ \frac{Lc_\sharp}8,\ \frac{3Lr_*}{32C''_{\mathrm{ia}}},\ \frac{r_n}{4K},\ \frac{L\varrho^{\log}_G}4\Bigr\},\]
$\varrho^{\log}_G\in(0,\log2]$ being the radius of the logarithm on $G^c$, the largest $\varrho\le\log2$ such that
$X\in G^c$, $|X-\mathbf 1|<1$ and $|\log X|<\varrho$ give $\log X\in\mathfrak g^c$; $\varrho^{\log}_G$ is positive,
$G^c=G\exp(\mathfrak g)$ being a closed subgroup of $GL(N,\mathbb C)$ with the Lie algebra $\mathfrak g^c$, it is
at least $\min\{\log2,2\arcsin\frac{z_G}2\}$ for semisimple $\mathfrak g$ (Lemma~\ref{C-368-6}(iii)), and at
$SU(N)$ and at $U(N)$ the last member of the minimum is implied by the others. So all these constants are explicit
in $L$ and $n$ at $SU(N)$ and $U(N)$, and for the other $G$ the radius of the second chain reads the pair through
$\varrho^{\log}_G$.
\end{enumerate}
\end{lemma}

\begin{proof}[Sketch of proof]
(i) For $SU(2)$, write $U=\cos r\,\mathbf1+i\sin r\,\omega$,
where $0\le r\le\pi$ and $\omega\in\mathfrak{su}(2)$ has norm one.
This identifies $SU(2)$ with the unit sphere in
$\mathbb R\mathbf1+i\mathfrak{su}(2)$. Left and right multiplication
preserve its Euclidean norm, so its normalized surface measure is
Haar measure. In these coordinates it is
$(2\pi^2)^{-1}\sin^2r\,dr\,d\omega$, whereas $dY=r^2dr\,d\omega$
for $Y=r\omega$. Therefore
\[
 \int_{SU(2)}F(U)\,dU
 =\frac1{2\pi^2}\int_{|Y|<\pi}
       F(e^{iY})\left(\frac{\sin|Y|}{|Y|}\right)^2dY.
\]
The omitted point $-\mathbf1$ has surface measure zero.
For general $G$, the differential of $\exp(iY)$, after translation
back to the identity, is $i\int_0^1 e^{-it\,\mathrm{ad}Y}\,dt$.
The paired eigenvalues $\alpha,-\alpha$ of $\mathrm{ad}Y$ contribute
$(\sin(\alpha/2)/(\alpha/2))^2$ to its real Jacobian; a zero eigenvalue
contributes one. The principal logarithm is unique on $|Y|<\pi$,
so change of variables on $G^\pi u_0$ gives the formula in (i),
with the normalization $\mathrm{vol}(G)^{-1}$. The range of the chart
is the one in Lemma~\ref{C-368-6}(ii);
the eigenvalues of $\mathrm{ad}\,Y$ are at most $2|Y|$ in modulus and at least $\mathrm{rk}\,G$ of them vanish, and
on $|Y|\le1$ the bound $0.173\,n|Y|^2$ follows because $x^{-2}(-\log\frac{\sin x}x)$ increases and is $0.1726$ at
$x=1$. For the nonnegative function $F$ in (i), the same formula with $\int_G$ on the left holds whenever $F$ vanishes Haar-almost everywhere off $G^\pi u_0$. Conversely, if the integral over $G^\pi u_0$ is finite, equality requires this vanishing. The $SU(2)$ construction bounds the one-bond functions $\mathsf{ls}_b$ defined before equation~\YMMref{(1381)}{569},
the logarithms $-\log\sigma_1$ of the density at the arguments $g_k\mathsf a(b)$, by $1$ on their domain
$\{g_k|\mathsf a(b)|<1\}$ (the Haar-term calculation in the proof of Proposition~\YMMref{62.1}{569}), and evaluates them with room at arguments of modulus below
$\frac38$ (the same calculation: $g_k|\mathsf a+\mathsf e|<1/8+2\rho_{\rm bd}\le3/8$). For $G$ the logarithm of $j_G$ is formed and holomorphic on
$\{Y\in\mathfrak g^c:|Y|<\pi\}$, so the domain stays, and with $\lambda(x)=-\log\frac{\sin x}x$ its modulus is at
most $(n-\mathrm{rk}\,G)\lambda(1)$, a number of the pair, on $\{Y\in\mathfrak g^c:|Y|<1\}$. On $\mathfrak g$ at
$SU(N)$, $\lambda$ being convex and increasing, the supremum of $|\log j_G|$ on a ball $|Y|\le r$ is taken with the
eigenvalues of $Y$ at $\pm r$ (and one at $0$ for odd $N$): it is at most $\frac{N^2}2\lambda(r)$, with equality
for even $N$, so about $0.0863N^2$ at $r=1$ and $0.0118N^2$ at $r=\frac38$, and at $r=\frac38$ it is below $1$
exactly when $N\le9$. The bound enters Proposition~\YMMref{53.15}{463}, the transfer of these one-bond functions,
linearly, as the size of the function transferred in equation~\YMMref{(1039)}{463}; so a constant of the pair takes the place
of $1$ there, the constants that use this transfer bound absorb it in a maximum, and neither the domain nor the
radius $\frac38$ moves.

(ii) $\nu$ is smooth off $0$, invariant, and $\nu\le|\cdot|\le w\nu$; so $\nu(Y)/\delta\le\frac43$ on
$|Y|\le\frac43\delta$, and on $|Y|\ge\frac53\delta$ the argument of $\vartheta_{\rm co}$ is at least
$\frac43+\lambda(\frac5{3w}-\frac43)=\frac53$.

(iii) $T_2$ is conjugate to $T_1W^{-1}$ with $W=\mathrm{Ad}(h)$; $\mathrm{tr}\log X-\mathrm{tr}\log(XW^{-1})$ lies
in $2\pi i\mathbb Z$, is constant on the convex set $\{|X-1|\le r,|X-W|\le r\}$, and vanishes at $\frac12(1+W)$,
the eigen-angles of $\mathrm{Ad}(h)$ being symmetric. The elimination factor is read off the relative average
$B\mapsto\frac1i\log[\mathcal M(e^{iB}V)(c)\mathcal M(V)(c)^{-1}]$: $L\,\mathsf m_c(V)X$ is the derivative at $s=0$
of its value at the field equal to $sX$ on the pivot bond and to $0$ on the other bonds of $\mathsf P_c$, so that
$L\,\mathsf m_c(V)$ is the coefficient of the pivot bond in its derivative at $B=0$. For $V$ in the domain and $B$
small, $e^{iB}V$ stays in it, its quotients and its $W_x$ staying in $D^c$ by the rule for products of
Lemma~\ref{C-368-6}(iii), so the relative average is formed by the series, holomorphic and $\mathfrak g^c$-valued,
$\mathfrak g$-valued at $G$-valued $V$ and real $B$, and $\mathsf m_c(V)$ maps $\mathfrak g^c$ into itself. At
$V=e^{ia}$ with $\sup|a|<\frac34r_*$ and $\sup|B|<r_1$ the relative average is at most $0.55$ and holomorphic in
$(a,B)$; so, by Cauchy's estimate in $B$, $|L\,\mathsf m_c(e^{ia})|\le\frac1{r_1}$, and $L\,\mathsf m_c(e^{ia})$ is
holomorphic in $a$. At $V=\mathbf 1$ the relative average is $\frac1i\log\mathcal M(e^{iB})(c)$, whose derivative
at $B=0$ is the linear average with real weights, which gives the pivot bond the weight $L^{-2}$; so $L\,\mathsf
m_c(\mathbf 1)=L^{-2}$, that is $\mathsf m_c(\mathbf 1)=L^{-3}$. Finally $h(a)=L\,\mathsf m_c(e^{ia})-L^{-2}$ is
holomorphic on $\sup|a|<\frac34r_*$, vanishes at $0$ and has $|h|\le\frac1{r_1}+L^{-2}\le2m_1$, so Schwarz's lemma
gives $|h(a)|\le\frac{8m_1}{3r_*}\sup|a|$; as $\mathsf m_c(e^{ia})-L^{-3}=L^{-1}h(a)$, this is the bound for every
$C_S\ge\frac{8m_1}{3r_*}$.

(iv) Each eigenvalue of $T$ lies within $r_n$ of $1$ and so has an argument of modulus at most $\arcsin
r_n=\frac\pi{2n}$; the arguments add up to at most $\frac\pi2$; at $\cos\frac\pi n\,e^{i\pi/n}1$ the radius
$\frac12$ fails for $n\ge6$. As $-\log(1-x)$ is convex on $[0,r_n]$,
$|\mathrm{Log}\det(1+E)|\le-n\log(1-|E|)\le\frac{|E|}{r_n}\rho_n=c_n|E|$; and $c_2<3.48$, $c_3=6\log2<4.16$,
$c_n>5$ for $n\ge4$. In the chains the radii are bounded from above only and the bounds from below only. What the
chains need at $G$ is $4Ka\le r_n$, which keeps $T$ within $r_n$ of $1$ through $|1-L^3\mathsf m_c|\le
L^2C_SC''_{\mathrm{ia}}\sup|\varphi|$ at potentials with $\sup|\varphi|<\frac{4a}L$, and $C_z\ge4c_nK$; for $n\ge4$
the radii of the $SU(2)$ construction, which give $4Ka\le\frac12$, fall by the factor $2r_n$, and its $C_z\ge20K$
rises by the factor $\frac{c_n}5$, which gives these two inequalities, the two of the $SU(2)$ construction holding
besides as $r_n<\frac12$ and $c_n>5$; for $n=2,3$ the constants of the $SU(2)$ construction serve, as $r_n\ge\frac12$
and $c_n<5$. The member $\frac{L\varrho^{\log}_G}4$ of the second radius keeps $\frac{4a}L\le\varrho^{\log}_G$, so
that the logarithms of that chain lie in $\mathfrak g^c$. The one input ``$z_c$ is formed at $(U^*)^{-1}$ if and
only if at $U$'' fails from $SU(3)$ on (at $\mathrm{diag}(1.9,1.9,\frac1{3.61})$); the chain needs only conjugate
values where both are formed (Lemma~\ref{C-368-24}), and with the two-sided guard $m(X)$ the clause holds for every
$G$.
\end{proof}

\subsection{Scalar kernels and the one-loop calculation}

\begin{lemma}[Colour at the trivial field]\label{C-368-11}
\begin{enumerate}
\item[(i)] Let $\mathfrak g$ be semisimple. Every $\mathrm{Ad}(G)$-invariant complex bilinear form on $\mathfrak
g^c$ is $q=\sum_iq_i\langle P_i\cdot,P_i\cdot\rangle$, $q_i=\frac1{n_i}\sum_aq(t^{(i)}_a,t^{(i)}_a)$ over an
orthonormal basis of $\mathfrak g_i$; every invariant form is a multiple of $\langle\cdot,\cdot\rangle$ if and only
if $\mathfrak g$ is simple, and then $C_A$ and $C_\rho$ do not depend on $H$; invariant linear forms vanish, and a
vector fixed by $\mathrm{Ad}(G)$ is $0$. For semisimple $\mathfrak g$ and a unit $H$,
$C_A(H)=\mathrm{tr}(\mathrm{ad}\,H)^2=\sum_iC^{(i)}_A\|P_iH\|^2$, which does not depend on $H$ if and only if all
$C^{(i)}_A$ agree: for $SU(2)\times SU(3)$ on $\mathbb C^2\oplus\mathbb C^3$ the $C^{(i)}_A$ are $20$ and $30$, for
$SU(2)\times SU(2)$ on $\mathbb C^2\oplus\mathbb C^2$ both are $16$. For simple $\mathfrak g$ and $z>0$ the Hessian
at $U=\mathbf 1$ of $\ln\det(D^{\mathrm T}D+z)-\frac12\ln\det(A''+DD^{\mathrm T}+z)$, $A''$ the Hessian of the
action, depends on the pair only through $C_A$ and $C_\rho$, exactly, on every finite lattice.
\item[(ii)] At $U=\mathbf 1$ every linear operator of a step (averaging operators, Laplacians, gauge-fixing
projections, propagators, the fluctuation operator) is $K\otimes1_{\mathfrak g^c}$, $K$ a real matrix formed from
the lattice alone; so the linear average is $Q_k\otimes1$, the Hessian is $d^*d\otimes1$, the projection $R$ is
$R_{\rm real}\otimes1$ and the propagator is $\mathsf M^{-1}\otimes1$, $\mathsf M=d^*d+dR_{\rm real}d^*+Q_k^*Q_k$.
For such operators: $K\mapsto K\otimes1_V$ respects sums, products, inverses and adjoints; for pairwise disjoint
sets $P_j$ that cover the support of $K(x',\cdot)$ and numbers $w_j\ge0$, $|(K_Vf)(x')|\le\sum_jw_j\sup_{P_j}|f|$
holds for one normed space $V$ if and only if $\sum_{x\in P_j}|K(x',x)|\le w_j$, so for one $V$ if and only if for
all; the best constants of $\ell^2$ bounds and of lower bounds for inner product norms do not depend on $V$; and
$|l_0(f)|\le\Phi(|l_1(f)|,\dots,|l_m(f)|)$, with real linear forms $l_j$ and non-decreasing $\Phi$, holds for
functions with values in one normed space if and only if it holds for real functions. So the bounds of the $SU(2)$
construction at the trivial field hold for $G$ with their constants, each in the norms to which its kind passes:
the bounds in suprema, by row sums and by non-decreasing functions of real linear forms in $|\cdot|$ and in
$\|\cdot\|$ alike; the $\ell^2$ bounds, the quadratic forms and the lower bounds, the lower bound of the
fluctuation operator among them, with its number of $L$ alone, in the inner-product norm $\|\cdot\|$, the sign
$|X|^2$ of the $SU(2)$ construction for $X\in\mathfrak g$ being read as $\langle X,X\rangle$. A bound of a sum of
squares of $|\cdot|$ retains its constant where it follows from Schur's test on $|K|$ or from convexity. Otherwise the comparison
$\|X\|\le|X|\le N^{1/2}\|X\|$ costs at most $N^{1/2}$ in an unsquared $\ell^2$ norm estimate and at most $N$ in its squared form.
\end{enumerate}
\end{lemma}

\begin{proof}[Sketch of proof]
(i) is Schur's lemma for $\mathrm{Ad}(G)$ on each simple ideal; the ideals are pairwise orthogonal for every
invariant form, and invariant linear forms vanish because $[\mathfrak g,\mathfrak g]=\mathfrak g$. For
$\mathfrak{su}(m)$, $\mathrm{tr}(\mathrm{ad}\,H)^2=2m\,\mathrm{tr}\,H^2$, and $\|H\|=1$ in $\frac15\mathrm{tr}$
means $\mathrm{tr}\,H^2=5$; this gives $20$ and $30$. The statement on the Hessian is the computation of
Proposition~\ref{C-368-13}, made on one finite lattice.

(ii) The averages of $\mathbf 1$ are $\mathbf 1$; by the chain rule the first-order part of the average at $\mathbf
1$ is the linear average, the straight average plus the gradient of the block mean of the contour sums, and that of
the straightened averages is the straight average, the contour terms cancelling; all have real weights formed from
the lattice alone; the second variation of the action at $\mathbf 1$ is $\sum_p\|dA(p)\|^2$, colour by colour; and
a holomorphic map that is $\mathfrak g$-valued at real points is $\mathfrak g^c$-valued. The transfer of bounds is
linear algebra, and a norming functional reduces a normed space to $\mathbb R$; for the row sums, $f$ equal to the
sign of $K(x',\cdot)$ times a unit vector on one $P_j$ gives the converse. In an orthonormal basis $(t_a)$ of
$\mathfrak g$, $(F,(X\otimes1)F)=\sum_a\langle F^a,XF^a\rangle$ for real $F$, so quadratic forms and their lower
bounds pass between inner-product norms; $|\cdot|$ on $\mathfrak g$ is no inner-product norm for rank at least $2$,
and read with the operator norm a clause of this kind can fail: at $SU(3)$ the floor of the $SU(2)$ construction
for a datum $X\delta_b$, $X=\mathrm{diag}(2,-1,-1)$, gives $4\|X\|^2=8$, against $4|X|^2=16$. Where a sharp linear
bound for $G$ is wanted, it is read off a nonlinear bound of the $SU(2)$ construction along $f=tg\sigma_3$,
$t\to0$, with Cauchy's inequality for the rest.
\end{proof}

The uniqueness argument needs a uniform gauge slice and a contraction
estimate in that slice. We state the required input explicitly, including
the extra rows at the interfaces. The result remains conditional for a
general group. The reduced $SU(2)$ problem is treated in
Propositions~\YMMref{30.9}{248} and~\YMMref{37.3}{284}; those statements
do not assert the additional interface-row estimates below.

\begin{definition}[Uniform slice hypotheses]\label{def:cg-uniform-slice}
Fix the pair, an odd $L\ge13$ with $200L^2z_G\ge1$, and the separated
regional geometry, with
\[
 \mathfrak M\ge M_{\rm ps},\qquad R_1\ge L,\qquad
 R_1\mathfrak M\ge w_{\rm S},\qquad
 \#\mathbb T^{(k)}\text{ in each direction}\ge6LM_{\rm ps}.
\]
The separation condition is that of $\mathfrak r_1$ for the determining
set $\mathbf B$. Permit every reversal-invariant set $\mathfrak F$ of
removed inward bonds. Lengths are level-$k$ lengths:
$\eta=L^{-k}$ and $\xi_j=L^j\eta$. All constants in the following
hypotheses depend only on $L$ and the pair, uniformly in this geometry,
$\mathfrak F$, depth, cutoff and volume.

The reference is $U_0=\mathbf1$. Let $a_j[A]$ and $v_j[A]$ be the
centered averaging logarithms and the accumulated straightening factors,
so that, on a bond $b$ of level $j$,
\begin{equation}\label{eq:cg-centered-factorization}
 \mathcal M^j(e^{i\eta A})(b)
 =v_j[A](b_-)\,\mathsf x_A(b)\,v_j[A](b_+)^{-1},
 \qquad \mathsf x_A(b)=e^{ia_j[A](b)}.
\end{equation}
The value at a reversed bond is its inverse. At the outer endpoint $z$
of an outward bond the factors satisfy
$v_j[A](z)=v_{j-1}[A](z)\,\mathsf v(a_{j-1}[A])(z)$.
The last factor depends only on interior bonds in the preceding block;
write $\kappa_6\ge1$ for its uniform small-field bound.

There is one row for each positively oriented interior or outward bond.
At a seam node $x$, order the retained inward bonds ending there as
$b_1(x),\ldots,b_m(x)$ and add the rows $(b_i(x),b_1(x))$, $2\le i\le m$.
Their scale is $\xi_j$ if the bonds have level $j$. Let $\mathfrak W$
be the space of $\mathfrak g$-valued row functions, with norm $\sup|\cdot|$
and inner product
\[
 (w,w')_1=\sum_\iota\xi_\iota^{-1}\langle w(\iota),w'(\iota)\rangle.
\]
Here $\langle\cdot,\cdot\rangle$ is the inner product of the pair;
it is not replaced by the operator norm. Fine-field adjoints use the
$\eta^3$-weighted inner product. The nonlinear row map is
\begin{equation}\label{eq:cg-natural-rows}
 \begin{aligned}
 \mathcal Q^\natural(A)(c)&=a_j[A](c)
       &&(c\text{ interior or outward}),\\
 \mathcal Q^\natural(A)(b_i,b_1)
       &=\frac1i\log\{\mathsf x_A(b_i)\mathsf x_A(b_1)^{-1}\}
       &&(\text{seam rows}).
 \end{aligned}
\end{equation}
For $\sup|B|\le c_{11}/2$, the products defining the datum in
\eqref{eq:cg-natural-datum} lie on the intrinsic branch and obey
\eqref{eq:cg-natural-datum-bound}. All logarithms below are on
the intrinsic branch through $\mathbf1$.
The hypotheses include their existence and analyticity on a neighbourhood
of the closed balls in which they are used.

\begin{enumerate}
\item[(S)] There are $c_{11}>0$, $\mathsf B_1\ge1$ and
$C_\sharp\ge1+\kappa_6$. A regular real configuration $U$ with
regularity parameter $\alpha_0$ and datum $e^{iB}$, where
$\alpha_1=\sup|e^{iB}-\mathbf1|$ and $\alpha_0+\alpha_1\le c_{11}$,
has a chart and gauge satisfying
\begin{equation}\label{eq:cg-uniform-slice}
 U^g=e^{i\eta A},\qquad g=v_j[A]\text{ on }\Lambda_j,
 \qquad \mathsf RD^*A=0,\qquad
 \|A\|_{X,\mathbf B}\le\tfrac14\mathsf B_1(\alpha_0+C_\sharp\alpha_1).
\end{equation}
At every retained inward bond $b$ ending at a seam node $x$,
\begin{equation}\label{eq:cg-seam-gauge}
 \mathsf x_A(b)\,v_j[A](x)^{-1}g(x)=e^{iB_j(b)}.
\end{equation}
The preliminary gauge used to obtain this chart bounds the discrepancy
from $\mathbf1$ at every level by $C_{\rm pg}(\alpha_0+\alpha_1)$;
its radius is included in $c_{11}$, and
$c_{11}\le r_G/(2C_{\rm pg})$. The regular class here bounds curvature
by $\alpha_0\eta^2\mathsf L(p)^{-2}$ and the current on full bonds
by $\alpha_0\mathsf L(b)^{-3}$.

\item[(N)] There are positive regularity ceilings $c_8,c_9$, a chart
radius $r_\circ$ and constants $K_1\ge4$, $K_7\ge1$. Set
$K_5=2(K_7+K_1)$ and $K_8=K_1(16+9K_1^3)$.
For regularity parameters at most $\min\{c_8,c_9\}$, the following
operators and estimates hold on $\|A\|_{X,\mathbf B}\le r_\circ$.
Let $F(A)$ be the real action in the chart, $F'(0)=0$, and write
\[
 \Delta_{\mathsf R}=F''(0)+D\mathsf RD^*,\qquad
 F'(A)+D\mathsf RD^*A=\Delta_{\mathsf R}A+\mathcal V(A).
\]
There are real linear maps $Q,H,\widetilde G$ and an isometric row
map $T$ preserving both the row supremum norm and its inner product, with
\begin{equation}\label{eq:cg-constrained-identities}
 QH=1,\quad Q\widetilde G=0,\quad\widetilde GQ^*=0,
 \quad\widetilde G\Delta_{\mathsf R}=1-HQ,
 \quad\ker\widetilde G=\operatorname{ran}Q^*.
\end{equation}
Define
\[
 \mathcal C(A)=T^{-1}\mathcal Q^\natural(A)-QA,
 \qquad K_A=D\mathcal Q^\natural(A),\qquad
 S_A=1+H^*\mathcal C'(A)^*.
\]
The map $S_A$ is invertible, with $\|S_A^{-1}\|_{\infty\to\infty}\le8/7$.
The differential of the actual constraint, after removing the gauge
factors, is $K_A^\flat=\Xi'_A K_A$, with $\Xi'_A$ invertible. In
particular a configuration $e^{i\eta A}$ is critical for its actual
averaging constraint precisely when $F'(A)$ annihilates $\ker K_A$.
For $B^\natural\in\mathfrak W$, put
\begin{equation}\label{eq:cg-contraction-map}
 \begin{aligned}
 w(A)&=S_A^{-1}H^*(F'(A)+D\mathsf RD^*A),\\
 N(A)&=-\mathcal V(A)+\mathcal C'(A)^*w(A),\\
 \Psi_{B^\natural}(A)&=\widetilde G N(A)
                +H\{T^{-1}B^\natural-\mathcal C(A)\}.
 \end{aligned}
\end{equation}
Whenever $\beta=\sup|B^\natural|$, $K_5\beta\le r\le r_\circ$ and
$\|A\|,\|A_1\|,\|A_2\|\le r$, the uniform bounds are
\begin{equation}\label{eq:cg-contraction-bounds}
 \|\Psi_{B^\natural}(A)\|\le\tfrac12K_5\beta+\tfrac12r,
 \qquad
 \|\Psi_{B^\natural}(A_1)-\Psi_{B^\natural}(A_2)\|
       \le\tfrac12\|A_1-A_2\|.
\end{equation}
The norms in this display are $\|\cdot\|_{X,\mathbf B}$.
The radii are also chosen so that the chart and row transformations
just described remain on their intrinsic logarithm branches.
\end{enumerate}
These are uniform analytic hypotheses. Colour-scalar linear operators
at the trivial reference do not by themselves prove either (S) or the
nonlinear estimates in (N).
\end{definition}

We describe explicitly which row datum corresponds to $e^{iB}$. At an
interior bond put $B^\natural(c)=B_j(c)$. For an outward bond with outer
endpoint $z$, put $v_z[B]=\mathsf v(B_{j-1})(z)$ and
\begin{equation}\label{eq:cg-natural-datum}
 \begin{aligned}
 B^\natural(c)&=\frac1i\log(e^{iB_j(c)}v_z[B])&&(z=c_+),\\
 B^\natural(c)&=\frac1i\log(v_z[B]^{-1}e^{iB_j(c)})&&(z=c_-),\\
 B^\natural(b_i,b_1)&=\frac1i\log(e^{iB_j(b_i)}e^{-iB_j(b_1)})
       &&(\text{seam rows}).
 \end{aligned}
\end{equation}
The uniform small-field bounds included in the hypotheses give
\begin{equation}\label{eq:cg-natural-datum-bound}
 \sup|B^\natural|\le(1+\kappa_6)\sup|B|\le C_\sharp\sup|B|.
\end{equation}
Indeed the outward products have logarithm norm at most
$|B_j(c)|+\kappa_6\sup|B|$; the seam products have norm at most
$2\sup|B|$, and $\kappa_6\ge1$. These products are taken within the
specified branch, as are the one-step factors defining $v_z[B]$.

\begin{lemma}\label{lem:cg-slice-system}
Under Definition~\ref{def:cg-uniform-slice}, the row derivative $K_A$
is onto for $\|A\|_{X,\mathbf B}\le r_\circ$. For
$K_5\sup|B^\natural|\le r_\circ$ the system
\begin{equation}\label{eq:cg-slice-system}
 F'(A)+D\mathsf RD^*A=K_A^*\omega,\qquad
 \mathcal Q^\natural(A)=B^\natural
\end{equation}
has exactly one solution with $\|A\|_{X,\mathbf B}\le r_\circ$.
The multiplier is unique without any further bound on its size.
\end{lemma}

\begin{proof}
Write $\widehat K_A=T^{-1}K_A=Q+\mathcal C'(A)$ and
$\widehat\omega=T^{-1}\omega$. Since $H^*Q^*=1$,
$H^*\widehat K_A^*=S_A$. The invertibility of $S_A$ implies the
injectivity of $K_A^*$, and hence the surjectivity of $K_A$ in these
finite-dimensional real inner-product spaces.

Set $\mathcal G(A)=F'(A)+D\mathsf RD^*A$.
The first equation in \eqref{eq:cg-slice-system} implies
$S_A\widehat\omega=H^*\mathcal G(A)$, so
$\widehat\omega=w(A)$. With this value of the multiplier the two
equations become
\[
 \Delta_{\mathsf R}A-N(A)=Q^*w(A),\qquad
 QA=T^{-1}B^\natural-\mathcal C(A).
\]
Apply $\widetilde G$ and use \eqref{eq:cg-constrained-identities}:
these equations imply $A=\Psi_{B^\natural}(A)$.
Conversely, a fixed point has the second equation after applying $Q$.
Applying $\widetilde G$ shows that
$\Delta_{\mathsf R}A-N(A)$ belongs to $\operatorname{ran}Q^*$.
Its coefficient is $w(A)$, because applying $H^*$ gives
$H^*\mathcal G(A)-H^*\mathcal C'(A)^*w(A)=w(A)$.
Thus the fixed-point equation is equivalent to the system.

By \eqref{eq:cg-contraction-bounds}, $\Psi_{B^\natural}$ maps the
closed radius-$r_\circ$ ball into itself and contracts distances by
$1/2$. This ball is complete, so it has one fixed point. The formula
$\omega=Tw(A)$ gives the only multiplier. The same argument on the
radius-$K_5\beta$ ball, or directly the fixed-point bound when
$\beta=0$, gives $\|A\|_{X,\mathbf B}\le K_5\beta$.
\end{proof}

\begin{proposition}[Uniqueness from uniform slice estimates]\label{cg-near-trivial-uniqueness}
Let $G$ be compact and connected, let $L\ge13$ be odd, and assume
$200L^2z_G\ge1$ and the uniform hypotheses of
Definition~\ref{def:cg-uniform-slice}. There are
$\varepsilon_G,\varrho_G>0$, depending only on $L$ and the pair,
with the following property in every geometry of that definition.
For any $\mathfrak g$-valued datum $B$ on the positively oriented
bonds of $\mathbf B_{\mathfrak F}$ with $\sup|B|<\varrho_G$, any
two real critical configurations in
$C(\mathbf B_{\mathfrak F},e^{iB})$ satisfying
\[
 |\partial U(p)-\mathbf1|<\varepsilon_G\eta^2\mathsf L(p)^{-2},
 \qquad |J(U)(b)|<\varepsilon_G\mathsf L(b)^{-3}
\]
at inner plaquettes and full bonds, respectively, belong to the
same residual gauge orbit. Thus $U'=U^u$, where $u=\mathbf1$
at every node of $\mathbf B_{\mathfrak F}$. Each configuration has
a representative $U^g=e^{i\eta A}$ with $g=v_j[A]$ on the tops,
$\mathsf RD^*A=0$ and
\[
 \|A\|_{X,\mathbf B}
       \le\tfrac14\mathsf B_1(\varepsilon_G+C_\sharp\sup|B|).
\]
The two radii are uniform in depth, sub-domain, cutoff and volume.
\end{proposition}

\begin{proof}
Take the constants of Definition~\ref{def:cg-uniform-slice} and put
\[
 \varepsilon_G=\min\{c_{11}/2,c_8,c_9,
                         r_\circ/(K_5L\mathsf B_1)\},\qquad
 \varrho_G=\min\{c_{11}/2,
                         r_\circ/(K_5L\mathsf B_1C_\sharp)\}.
\]
For two configurations $U_a,U_b$ as in the statement set
$\alpha_0=\varepsilon_G$ and
$\alpha_1=\sup|e^{iB}-\mathbf1|\le\sup|B|<\varrho_G$.
Hypothesis (S) gives $A_*,g_*$, $*=a,b$, with
\[
 U_*^{g_*}=e^{i\eta A_*},\qquad
 \mathsf RD^*A_*=0,\qquad
 \|A_*\|_{X,\mathbf B}\le
 \tfrac14\mathsf B_1(\alpha_0+C_\sharp\alpha_1)
 \le\frac{r_\circ}{2K_5L}.
\]
This gives the asserted representatives.

Both have the row datum \eqref{eq:cg-natural-datum}. To check this,
at an interior bond the two endpoint gauges in
\eqref{eq:cg-centered-factorization} cancel, giving
$e^{ia_j[A_*](c)}=e^{iB_j(c)}$ and hence $a_j[A_*](c)=B_j(c)$.
At an outward bond the uncancelled factor is
$\mathsf v(a_{j-1}[A_*])(z)$; it equals $v_z[B]$, since it reads
only the interior bonds of the preceding block. Moving this factor
to the other side gives the first two formulas in
\eqref{eq:cg-natural-datum}. Finally,
\eqref{eq:cg-seam-gauge} at $b_i$ and $b_1$ gives
\[
 \mathsf x_{A_*}(b_i)\mathsf x_{A_*}(b_1)^{-1}
       =e^{iB_j(b_i)}e^{-iB_j(b_1)}.
\]
Taking the stipulated intrinsic logarithms proves
$\mathcal Q^\natural(A_a)=B^\natural=\mathcal Q^\natural(A_b)$.
By \eqref{eq:cg-natural-datum-bound},
$K_5\sup|B^\natural|\le K_5C_\sharp\sup|B|
<r_\circ/(L\mathsf B_1)\le r_\circ$.

Criticality is preserved by gauge transformations. The criticality
clause of (N), together with the surjectivity in
Lemma~\ref{lem:cg-slice-system}, gives a multiplier $\omega_*$
with $F'(A_*)=K_{A_*}^*\omega_*$. Since $\mathsf RD^*A_*=0$,
both pairs solve \eqref{eq:cg-slice-system} for the same datum.
Lemma~\ref{lem:cg-slice-system} therefore gives $A_a=A_b$.
Consequently $U_b=U_a^{g_b^{-1}g_a}$. The two gauges agree at every
top by \eqref{eq:cg-uniform-slice}. At a seam node, one retained
inward bond in \eqref{eq:cg-seam-gauge} determines $g_*(x)$ from
$A_*$ and $B$, so they agree there as well. Every node of
$\mathbf B_{\mathfrak F}$ is a top or such a seam node. Thus
$g_b^{-1}g_a$ is a residual gauge transformation.
\end{proof}


\begin{definition}[Uniform orbit and energy hypothesis]\label{C-368-12}
Let $G\subset U(N)$ be compact and connected and $L\ge13$ odd. Let $\mathcal M$ be a $G$-valued averaging map that,
at every coarse bond $c$, is the formula of Lemma~\ref{C-368-7}(i) at every field whose plaquettes on the double
block of $c$ are within $\omega_G=\min\{\omega_M,r_G\}$ of $\mathbf 1$, the radius below which
Lemma~\ref{C-368-7}(i) forms that formula in $G$ for every compact connected $G$, also where $\mathfrak g$ has a
centre and $z_G=0$; where the centre of $\mathfrak g$ is $0$, $\omega_G$ is the radius
$\omega^G=\min\{\omega_M,z_G\}$ of Definition~\ref{C-368-3}(3). The radial reduction below reads $\mathcal M$ at such fields alone.
Let $\Lambda_0,\dots,\Lambda_t$ be the lattices of the tower, $\Lambda_0$ the bare one and $\Lambda_t$ the block
lattice. In this definition and its reductions, in exception to Definition~\ref{C-368-3}(1), for functions $f$ on the bonds
or on the plaquettes of one lattice $\|f\|^2=\sum|f(\cdot)|^2$, $\|f\|_1=\sum|f(\cdot)|$ and
$|f|_\infty=\max|f(\cdot)|$, the values being measured in the operator norm; the energy norm of a $\mathfrak
g$-valued $1$-form $\delta$ on $\Lambda_0$ is $\|\delta\|_E^2=\|d\delta\|^2+L^{-2t}\|\delta\|^2$, in bare units;
and the two matrix norms are compared only through $\|X\|\le|X|\le N^{1/2}\|X\|$. The uniform orbit and energy hypothesis is that there are
$a_G,\varepsilon_*\in(0,1]$, constants depending only on $L$ and the pair that read $\omega_G$, such that for every
$t\ge1$, on every torus whose lattice $\Lambda_t$ has at least $12$ sites a side, all regions being the whole
torus:
\begin{enumerate}
\item[(i)] a configuration $U$ with $\mathcal M^tU=\mathbf 1$ that is critical for the action on $\{\mathcal
M^tU=\mathbf 1\}$ and has $|\partial U(p)-\mathbf 1|<\epsilon L^{-2t}$ at all bare plaquettes, $\epsilon\le a_G$,
is flat;
\item[(ii)] Independently of this hypothesis, flat fields $U$ with $\mathcal M^tU=\mathbf 1$ are the $\mathbf 1^u$ with $u=\mathbf 1$ on $\Lambda_t$; under (i), therefore,
the small critical configurations over $\mathbf 1$ with these plaquettes are exactly this orbit;
\item[(iii)] for each $t$, each torus and each $\epsilon'\le a_G$ there is a neighbourhood of $\mathbf 1$,
depending on them, at each point $V$ of which two critical configurations of $\{\mathcal M^tU=V\}$ with $|\partial
U(p)-\mathbf 1|<\epsilon'L^{-2t}$ lie in one orbit of the gauge functions equal to $\mathbf 1$ on $\Lambda_t$,
which depends holomorphically on $V$;
\item[(iv)] if $U$ is critical for the action on $\{\mathcal M^tU'=V\}$, $V=\mathcal M^tU$, and $|\partial U(p)-\mathbf 1|\le\varepsilon_*L^{-2t}$, then
\[\sum_p|\partial U(p)-\mathbf 1|^2\le4N\varepsilon_*L^{-t}\sum_{p'}\Bigl|\frac1i\log\partial V(p')\Bigr|,\]
the sums running over the plaquettes of $\Lambda_0$ and of $\Lambda_t$; in particular over a flat block field the
critical configurations with these plaquettes are flat, and over a block field near $\mathbf 1$ the regular
critical configurations form one gauge orbit.
\end{enumerate}
This hypothesis requires $a_G$ and $\varepsilon_*$ to be independent of $t$ and the torus. If all the uniform analytic estimates in the reduction below hold, its radial argument permits
$a_G=\min\{\varepsilon_5/2,[4N(68C_sC_\alpha+132C_R)]^{-1}\}$ with the numbers $\varepsilon_5$, $C_s$, $C_\alpha$,
$C_R$ of those estimates. Write $a_*(G)$ for the supremum of the admitted orbit radii, taking it to be $0$ if no positive radius is admitted.
The construction uses a fixed admitted positive $a_G$, not an assumed maximiser. Positivity of this radius for general $G$, and a lower bound uniform in $N$ for $Na_*(SU(N))$, are not established by the matrix estimates alone.
\end{definition}

\begin{proof}[Flat fields and conditional reductions]
For (i), assume a uniform gauge $g$, equal to $\mathbf 1$ on $\Lambda_t$, in which the potentials of $U$ and of all its
averages are small, $(\mathcal M^iU)^g=e^{i\alpha_i}$ with $|\alpha_i|_\infty\le C_s\epsilon L^{-(t-i)}$ and
$\|\alpha_i\|\le C_\alpha L^{t-i/2}\mathcal E^{1/2}$, $\mathcal E=\sum_p|\partial U(p)-\mathbf 1|^2$. Let $F(a')$
be the action at $e^{ia'}$, $a=\alpha_0$, and $T$ the derivative at $a$ of the $t$-fold average read in the chart.
The required first-variation estimate is $dF(a)\ge(\frac5{6N}-68C_sC_\alpha\epsilon)\mathcal E$, from
$\langle\theta,H_{e^{i\theta}}\rangle\ge\frac56\|e^{i\theta}-\mathbf 1\|^2$ for $|\theta|\le1$. Along the ray
$s\mapsto e^{isa}$ an abelian field would stay in the constraint set; for $G$ it leaves it at second order, and we
correct it. Assume also that the derivative has the form $T\delta=d\lambda_\delta+K\delta$ with $K=(Q'+E_{t-1})\cdots(Q'+E_0)$,
$\|Q'\rho\|\le L^{-1/2}\|\rho\|$, $\|E_i\rho\|\le C(|\alpha_i|_\infty+|\alpha_{i+1}|_\infty)\|\rho\|$; the
interpolation $H_0$ of the straight average, composed with the inverse of its average, has $Q'^tH_0=1$ and
$\|H_0\psi\|_E\le33L^{-t/2}\|\psi\|$. The number $\varepsilon_5$ is the last of a falling chain
$\varepsilon_1\ge\varepsilon_2\ge\dots\ge\varepsilon_5$ in $(0,1]$; existence of these positive thresholds uniformly in depth and torus is part of the analytic input:
below $\varepsilon_1$ all the averages $\mathcal M^iU$ are formed by the formula, their plaquettes staying within
$2\epsilon L^{-2(t-i)}$ of $\mathbf 1$, below $\omega_G$, and $\|\partial\mathcal M^iU-\mathbf
1\|\le2L^{i/2}\mathcal E^{1/2}$; below $\varepsilon_2$ the gauge $g$ exists with the two bounds above; below
$\varepsilon_3$ the first variation obeys the bound above; below $\varepsilon_4$ the derivative $T$ has the
structure above, and the radial direction the form below; and below $\varepsilon_5$, besides,
$\|(K-Q'^t)H_0\|\le\frac12$. For $\epsilon\le\varepsilon_5$, $KH_0=1+\Delta$ with $\|\Delta\|\le\frac12$, so $K$
has the right inverse $H=H_0(1+\Delta)^{-1}$ with $\|H\psi\|_E\le66L^{-t/2}\|\psi\|$, and $T$ is onto; and
$Ta=d\lambda_*+R$ with $\|R\|\le C_R\epsilon L^{t/2}\mathcal E^{1/2}$. The direction
$a-\delta_{\lambda_*}a-HR+\delta_{\lambda'}a$, $\delta_\lambda$ the infinitesimal gauge transformations and
$\lambda'=\lambda_{HR}$, lies in $\ker T$, and $dF$ vanishes on gauge directions; so the critical equation gives
\[\Bigl(\frac5{6N}-68C_sC_\alpha\epsilon\Bigr)\mathcal E\le dF(a)=dF(HR)\le2\mathcal
E^{1/2}\|HR\|_E\le132\,\mathcal E^{1/2}L^{-t/2}\|R\|\le132\,C_R\epsilon\,\mathcal E,\]
and $\mathcal E=0$ for $\epsilon\le a_G$. This proves the radial reduction under the displayed uniform estimates, not their existence. (ii) A flat field has trivial transports around closed paths in boxes, and its averages are the
transports along straight paths. This gives the flat-field assertion without the uniform analytic input:
lifting to $\mathbb Z^3$ and using the comb gauge on growing cubes gives
$U=\mathbf 1^u$, and the averages $\mathbf 1$ force $u=\mathbf 1$ on $\Lambda_t$. (iii) Assuming the gauge-fixed constrained Hessian is invertible at the flat orbit, near $\mathbf 1$ the
implicit function theorem for the system ``gradient of the action orthogonal to the constraint, constraint $=V$'',
on the orthogonal complement of the gauge directions, gives a local branch and local uniqueness. Exclusion of other critical branches throughout the stated plaquette domain is part of (iii).
(iv) The terminal-field energy bound is an additional analytic input. Its descending construction would require the following estimates. Let $V_j=\mathcal M^jU$,
$\phi_{V}(p)=|\partial V(p)-\mathbf 1|$, $\Phi_V(p)=\frac1i\log\partial V(p)$, $D_j=\max_p\phi_{V_j}(p)$, and $D_V$
the derivative of the plaquette variables along a tangent vector in the left chart. The required hierarchical axial-gauge descent
solves, one step down, $\mathcal M'(V_j)X_j=X_{j+1}$ with $\|X_j\|\le
L(1+CD_j)\|X_{j+1}\|+C\|\phi_{V_j}\|$, starting from $X_t=0$, and the defect $E_j=D_{V_j}X_j-\Phi_{V_j}$ obeys
$\|E_j\|_1\le L\|E_{j+1}\|_1+C\|\phi_{V_j}\|(\|\phi_{V_j}\|+\|X_{j+1}\|)$; $X_0$ is tangent to the constraint set,
so $0=\sum_p\langle\Phi_U(p)+E_0(p),\sin\Phi_U(p)\rangle$, and $\langle\Phi,\sin\Phi\rangle\ge\frac1{2N}|\Phi|^2$
does not by itself provide the uniform aggregate error bound needed for (iv); that bound is part of the terminal-field input. At $SU(N)$ a central element $e^{2\pi i/N}\mathbf 1$ on one bond off the
block axes gives a nonflat configuration with zero current and plaquettes at distance $2\sin(\pi/N)$. If the chosen extension of $\mathcal M$ averages this field to $\mathbf 1$, it obstructs an orbit radius above $2L^2\sin(\pi/N)$
at the first level. Agreement with the small-field averaging formula alone does not impose this behaviour on the extension.
\end{proof}

\begin{proposition}[One loop]\label{C-368-13}
Let $L\ge13$ be odd and $k\ge1$.
\begin{enumerate}
\item[(i)] Let $\mathfrak g$ be simple. On every torus $\Psi^G_k$ is tempered: real-analytic near $\mathbf 1$, zero
at $\mathbf 1$, invariant under gauge transformations and lattice isometries, with the symbol
$\delta_{aa'}\hat\pi^G_{\mu\nu}$, $\hat\pi^G=\frac{C_A}8\hat\pi^{\Psi_k}+\gamma_G\hat\pi^{\mathfrak B_k}$,
holomorphic and bounded on a strip; here $\mathfrak B_k=(\mathfrak B_{k,s})_s$ is the scalar functional of the
tadpole, formed for $SU(2)$ on the tower of height $k$ and side $s$ (Section~\YMMref{42.1}{318}),
\[\mathfrak B_{k,s}(W)=\sum_pw_{k,s}(p)\,a_p\bigl(c_k(W)\bigr)\qquad(W\text{ near }\mathbf 1\text{ on }\Lambda_k),\]
$c_k$ the critical branch (Theorem~\YMMref{N}{318}) and $w_{k,s}(p)\in[0,1]$ the diagonal of the orthogonal
projection of the plaquette functions of $\Lambda_0$ onto $d(\ker\mathsf Q^{[k]})$, $\mathsf Q^{[k]}$ the $k$-fold
linear average, which reads $L$, $k$, $s$ and no group; $\mathfrak B_0=0$, and with the weights $1$ in the place of
$w_{k,s}$ it would be $\mathcal S_k=A\circ c_k$. The functional $\mathfrak B_k$ is tempered
(Section~\YMMref{43.1}{324}), with $b[\mathfrak B_k]=j_k:=\frac23(L^{-k}-L^{-4k})$; the coefficient of $\Psi^G_k$
is that of Theorem~\ref{C-368-4}(c), and the increments $b[\Psi^G_l]-L^{-1}b[\Psi^G_{l-1}]$ are the coefficients of
tempered functionals and telescope as in Section~\YMMref{43.3}{329}. Theorem~\YMMref{O}{325} holds for $G$
with the remainder constant $\frac{C_A}8C_{\rm rem}(L)+\frac23L^{-4}|\gamma_G|\delta_0^2$, $C_{\rm rem}(L)$ and
$\delta_0$ those of $SU(2)$, the least that the argument gives, and
\[\sup_k|b[\Psi^G_k]|\le\frac{C_A}8\sup_k|b[\Psi_k]|+\Bigl|\frac{C_\rho}3-\frac{C_A}8\Bigr|L^{-1};\]
so the condition $B_\sharp\ge(1+L^{-1})c_2+1$ of the $SU(2)$ construction, $c_2$ that of
Theorem~\YMMref{O}{325}, and Lemma~\ref{lem:gaussian-flow-bound} stand with this number, which depends on $L$, $C_A$,
$C_\rho$ alone. For $SU(N)$ the two factors are $\frac{N^2}4$ and $\frac{N^2-4}{12}$.
\item[(ii)] Let $\mathfrak g$ be simple. Theorem~\ref{III-12-105} holds for $G$ in every clause: the numbers $b^0_l(G)$ converge geometrically,
\[\Bigl|b^0_l(G)-\frac{C_A}8b^0_\infty\Bigr|\le\frac{C_A}8C_{\rm geo}L^{-l/6}+\frac23|\gamma_G|L^{3-4l},\]
the functionals of that theorem for $G$ are $\frac{C_A}8$ times those of $SU(2)$ plus $\gamma_G$ times one scalar
functional bounded by $5C_h^2L^{-k}$, $C_{\rm re}(G)\le\frac{C_A}8C_{\rm re}+10|\gamma_G|C_h^2$, and $C_{\rm
inc}(G)\le17(2C_{\rm rem}(G)C_{\rm re}(G))^{1/2}\delta_0^{-2}$ with $C_{\rm rem}(G)=\frac{C_A}8C_{\rm
rem}+5|\gamma_G|C_h^2/L$ in this clause, a larger number than the remainder constant of (i).
\item[(iii)] Let $G$ be compact and connected, $l\ge1$ and, in this clause, $k=l-1\ge0$. On the tower of height $l$
and side $s$, formed for $G$, let $p_j(W)$ be the critical point near $\mathbf 1$ of $\mathcal
S_{j-1}+\mathrm{GF}_j$ on the fibre $\{\mathcal M(V)=W\}$, $c_l=p_1\circ\dots\circ p_l$ the critical branch, with
$\mathcal M^lc_l(W)=W$ and $c_l(\mathbf 1)=\mathbf 1$, and $\mathcal S_l=A\circ c_l$
(Theorem~\YMMref{N}{318}); let $\mathsf Z_l(W)$ be the normaliser of the first part formed for $G$, the
Gaussian integral of the Hessian of $\mathcal S_{l-1}+\mathrm{GF}_l$ on the fibre at $p_l(W)$
(Theorem~\YMMref{N}{318}), whose logarithm has the quadratic coefficient $b^0_l(G)$ for simple $\mathfrak g$
(Theorem~\YMMref{O}{325}, Theorem~\ref{C-368-4}(c)); and let $\mathsf Z^{(k)}[\mathbf P]=\prod_c|\det\mathsf
m_c(V)|^{-1}\int dB\,e^{-\frac12(B,\mathsf A^\sharp(\mathbf P)B)_1}$, $V=\mathcal M^kU$, be the Gaussian integral
of the step of the second part at a triple $\mathbf P=(U,0,J')$ (Proposition~\YMMref{50.2}{386}, with $\dim\mathfrak g$ real coordinates per positive nonpivot bond), with the
fluctuation operator $\mathsf A^\sharp$ of the step, whose value at the plain triples over $W$ is the Gaussian
normaliser $Z^{(k)}(W)=Z_k[P_W]$ preceding equation~\YMMref{(957)}{435}. Near $W=\mathbf 1$ it is formed at the plain triple
$\mathbf P_W=(c_l(W),0,L^3J(c_l(W)))$ of the branch, is positive, is the same along the gauge orbit of this specified branch, and
identification with every regular critical configuration over $W$ requires the orbit hypothesis of Definition~\ref{C-368-12}, or (V) in its regular domain. At the specified branch one has
\[\mathsf Z^{(k)}[\mathbf P_W]=L^{nn_l}L^{-kd_l/2}\mathsf Z_l(W),\qquad n_l=3s^3,\quad d_l=2n(L^3-1)s^3;\]
so the quadratic coefficient of $\log Z^{(l-1)}$ is $b^0_l(G)$ for simple $\mathfrak g$, and for semisimple
$\mathfrak g$ it is, on each ideal $\mathfrak g_i$,
$b^{0,(i)}_l=\frac{C^{(i)}_A}8b^0_l+(\theta_i-\frac38C^{(i)}_A)\frac13(L^3-1)L^{-4l}$. The localised Gaussian
family (Proposition~\YMMref{50.2}{386}), formed for $G$, has its members formed and holomorphic near the trivial triple,
and there the same Hessian as $\log Z^{(k)}$; its coefficient there is one number for all data: for simple
$\mathfrak g$
\[\frac{C_A}8b^0_{k+1}+\gamma_G(j_{k+1}-L^{-1}j_k)=b^0_{k+1}(G),\]
with $j_k=\frac23(L^{-k}-L^{-4k})$ and $b^0_{k+1}$ the increment of $SU(2)$, and for semisimple $\mathfrak g$, on
$\mathfrak g_i$, $\frac{C^{(i)}_A}8b^0_{k+1}+\gamma^{(i)}(j_{k+1}-L^{-1}j_k)=b^{0,(i)}_{k+1}$,
$\gamma^{(i)}=\frac12\theta_i-\frac3{16}C^{(i)}_A$. That this family is a pointed family of its level on the
analyticity cores of $G$, holomorphic on balls of one radius, with the contraction bound of its expansion, is part
of (II) of Definition~\ref{C-368-31} (Proposition~\ref{C-368-22}(iii)), and therefore is asserted here under
hypothesis (II) for simply connected $G$ with simple $\mathfrak g$ and $200L^2z_G\ge1$.
\item[(iv)] Let $\mathfrak g=\bigoplus_i\mathfrak g_i$ be semisimple, $\Omega_i$ the sum of $t_a^2$ over an
orthonormal basis of $\mathfrak g_i$, $\tau_{ij}=\frac1N\mathrm{tr}(t^2s^2)$ for unit $t\in\mathfrak g_i$,
$s\in\mathfrak g_j$, $j\ne i$, and $\tau_{ii}=\frac1{n_iN}\mathrm{tr}(\Omega_it^2)$ for unit $t\in\mathfrak g_i$,
numbers of the pair; so $\theta_i=\frac1N\mathrm{tr}(\Omega H^2)=\sum_jn_j\tau_{ij}$ for a unit $H\in\mathfrak
g_i$, and $n_i\tau_{ii}$ is the number $C_\rho$ of the simple pair of $\mathfrak g_i$, the subgroup with Lie
algebra $\mathfrak g_i$ on $\mathbb C^N$. With one set of sequences and functions for all ideals and all groups,
the one-loop coefficients on $\mathfrak g_i$ are $\frac{C^{(i)}_A}8$ times those of $SU(2)$ plus
$(\theta_i-\frac38C^{(i)}_A)$ times one scalar sequence: at every semisimple pair, the ideals being tied by the
representation or not,
\[b^{(i)}[\Psi^G_k]=\frac{C^{(i)}_A}8b[\Psi_k]+\Bigl(\theta_i-\frac38C^{(i)}_A\Bigr)t_k,\qquad t_k=\tfrac13(L^{-k}-L^{-4k}),\]
$b^{(i)}[\Psi^G_k]$ being the coefficient on $\mathfrak g_i$ of the one-loop functional $\Psi^G_k$ of the pair and
$\Psi_k$ that of $SU(2)$. The same holds with weights, for an action $\sum_mw_m$ times the Wilson actions in
representations $\rho_m$, $w_m>0$, with $C^{(i)}_A(w)$, $\theta_i(w)$ formed with the weighted form $\langle
X,Y\rangle_w=\sum_mw_mN_m^{-1}\mathrm{tr}\,\rho_m(X)\rho_m(Y)$ in the place of $\langle\cdot,\cdot\rangle$, a
rational weight being a multiplicity; the limits of the coefficients, $\frac{C^{(i)}_A(w)}8b^0_\infty$, read the
weights through this normalisation alone, and the limit as a form on $\mathfrak g$,
$\frac{b^0_\infty}8\mathrm{tr}(\mathrm{ad}\,X)^2$, does not read them. Equivalently, for a unit vector $t$ of
$\mathfrak g_i$, with $\Pi^F=F''(\mathbf 1)$ the polarisation of a functional $F$ (Section~\YMMref{43.1}{324}),
\[\Pi^{\Psi^G_k}[tf,tf']=\Pi^{\Psi^{(i)}_k}[tf,tf']+\Bigl(\sum_{j\ne
i}n_j\tau_{ij}\Bigr)\Upsilon_k[f,f'],\qquad0\le\Upsilon_k[f,f]\le\tfrac12\mathsf h_k(f,f),\]
$\Psi^{(i)}_k$ that of the simple pair of $\mathfrak g_i$, $\mathsf h_k=\mathcal S_k''(\mathbf 1)$ the Hessian at
$\mathbf 1$, for one colour, of $\mathcal S_k=A\circ c_k$ (Theorem~\YMMref{N}{318}), and
$\Upsilon_k[f,f']=\frac12\sum_p\mathsf v_k(p)(d\mathsf c_kf)(p)(d\mathsf c_kf')(p)$ one scalar form of the lattice
that reads no group, $\mathsf v_k(p)\in[0,1]$ the diagonal of the orthogonal projection onto $d(\ker\mathsf
Q^{[k]})$, a variance of the scalar Gaussian theory on $\ker\mathsf Q^{[k]}$, and $\mathsf c_k$ the linear part of
the chain at $\mathbf 1$. The form $\Upsilon_k=\frac12\Pi^{\mathfrak B_k}$ is the polarisation of the tempered
functional $\frac12\mathfrak B_k$ of (i), with $b[\Upsilon_k]=t_k\in[0,\frac13L^{-k}]$; so
$b^{(i)}[\Psi^G_k]=b[\Psi^{(i)}_k]+(\sum_{j\ne i}n_j\tau_{ij})\,t_k$. For block sums $\tau_{ij}=0$ for $j\ne i$.
The form of the coefficients on $\mathfrak g$ is $\sum_ib^{(i)}[\Psi^G_k]\langle P_i\cdot,P_i\cdot\rangle$, with
$r$ numbers, which are bounded ideal by ideal,
\[|b^{(i)}[\Psi^G_k]|\le\frac{C^{(i)}_A}8\sup_k|b[\Psi_k]|+\Bigl|\frac{\theta_i}3-\frac{C^{(i)}_A}8\Bigr|L^{-1},\]
and so by the largest of these $r$ bounds. For a unit $H\in\mathfrak g_i$ the numbers
$\mathrm{tr}(\mathrm{ad}\,H)^2$ and $\frac1N\mathrm{tr}(\Omega H^2)$ that Definition~\ref{C-368-2} forms for simple
$\mathfrak g$ are $C^{(i)}_A$ and $\theta_i$. The pairs $(C^{(i)}_A,\theta_i)$ may differ between the ideals, as at
$SU(2)\times SU(3)$ on $\mathbb C^2\oplus\mathbb C^3$, where $C^{(i)}_A$ is $20$ and $30$
(Lemma~\ref{C-368-11}(i)), the coefficient being then given ideal by ideal, as above, or agree, as $(16,6)$ on both
ideals of $SU(2)\times SU(2)$ on $\mathbb C^2\oplus\mathbb C^2$ and $(8,6)$ on both on $\mathbb C^2\otimes\mathbb
C^2$; and
$b^{(1)}[\Psi^G_k]-b^{(2)}[\Psi^G_k]=\frac18(C^{(1)}_A-C^{(2)}_A)(b[\Psi_k]-3t_k)+(\theta_1-\theta_2)t_k$.
\end{enumerate}
\end{proposition}

\begin{proof}[Sketch of proof]
(i) By invariance the background may be taken, to second order, along one unit $H\in\mathfrak g$, and the
fluctuations split into the eigenspaces of $\mathrm{ad}\,H$. Every object of a step but the action is built from
products, logarithms and means of group elements, which commute with homomorphisms of Lie algebras; with a
parameter whose cube vanishes this shows that the sector of eigenvalue $c$ of $\mathrm{ad}\,H$ in the field $\theta
H$ is the charged sector of $SU(2)$ in the field $\frac c2\theta t_3$, $\mathrm{ad}\,t_3$ having the eigenvalues
$0,\pm2$. Summing over the sectors, the group enters through $\sum c^2=C_A$. The action adds one term:
$\cos(fH)=\mathbf 1-\frac12f^2H^2+O(f^4)$ gives the tadpole $-\frac12d^{\mathrm T}(d\theta)^2d\otimes M_H$,
$M_HY=\frac12P\{H^2,Y\}$, $\mathrm{tr}\,M_H=C_\rho$, whose propagator is the orthogonal projection onto $d(\ker
Q^{[k]})$, $Q^{[k]}$ the $k$-fold linear average; its trace is $2s^3(\eta^{-3}-1)$, two thirds of the plaquettes
less the block lattice. Eliminating the group-free second differential,
\[\Xi^{G\prime\prime}(\mathbf 1)(\theta H,\theta H)=\frac{C_A}8\,\Xi''(\mathbf 1)(\theta t_3,\theta
t_3)+\frac12\Bigl(C_\rho-\frac38C_A\Bigr)\sum_pw_k(p)(d\theta)(p)^2,\]
$\Xi$ the logarithm of the normaliser and $w_k$ the diagonal of the projection; the count of the trace gives the
formula for $b[\Psi^G_k]$. The functional $\mathfrak B_k$, formed here, is tempered by a step of its own. It is
real-analytic near $\mathbf 1$ and $0$ at $\mathbf 1$; it is gauge invariant, as $c_k(W^v)=c_k(W)^{\tilde v}$ for
$v$ near $\mathbf 1$ or constant and $a_p$ is a class function; and it is invariant under the isometries of the
torus that map $\Lambda_k$ onto itself, which keep $d(\ker\mathsf Q^{[k]})$ and so $w_{k,s}$. On every torus the
orthogonal projection onto $d(\ker\mathsf Q^{[k]})$ is the matrix of the periodic operator $\mathcal P=d(\mathsf
X^{-1}-\mathsf X^{-1}Q^{\mathrm T}C^{-1}Q\mathsf X^{-1})d^{\mathrm T}$, $\mathsf X=d^{\mathrm T}d+d\varrho
d^{\mathrm T}+Q^{\mathrm T}Q$, $\varrho=1-R^{\mathrm T}R$, $C=Q\mathsf X^{-1}Q^{\mathrm T}$, formed with the
operators $Q$, $R$ of the response of the $SU(2)$ construction (Section~\YMMref{43.2}{327}), which is real, acts trivially on the colour and has an exponentially decaying
kernel: its quadratic form at a plaquette function $g$ is the maximum of $2(dB,g)-|dB|^2-|\varrho d^{\mathrm
T}B|^2$ over $B\in\ker Q$, and a gauge direction $d\lambda$, $\lambda\in\ker R$, removes the last term, which
leaves $|\mathsf P_{\mathcal C}g|^2$, $\mathsf P_{\mathcal C}$ the projection onto $\mathcal C=d(\ker
Q)=d(\ker\mathsf Q^{[k]})$. So $w_{k,s}(p)-w_k(p)=\sum_{m\ne0}\mathcal P(p,p+sm)$ tends to $0$ exponentially in
$s$, $w_k$ the diagonal at the infinite side, and the symbol $\hat\pi^{\mathfrak
B_k}_{\mu\nu}(q)=\eta^4\sum_{p\bmod\mathbb Z^3}w_k(p)\,(\hat f(q)e_\mu)(p)\,(\hat f(-q)e_\nu)(p)$, $\hat f=\hat
d\hat h$ with the symbol $\hat h$ of the response (Section~\YMMref{43.2}{327}), is holomorphic and bounded on the
strip of width $2\delta_0$. The trace $\sum_pw_{k,s}(p)=\dim d(\ker\mathsf Q^{[k]})=2s^3(\eta^{-3}-1)$, a third on
each orientation, gives $b[\mathfrak B_k]=\eta^4\cdot\frac23(\eta^{-3}-1)=j_k$. With this the theorems of the
$SU(2)$ construction on tempered functionals apply to $\mathfrak B_k$: the coefficient (Section~\YMMref{43.1}{324});
the composition with one step of the branch (Section~\YMMref{43.3}{329}), $b[\mathfrak B_{l-1}\circ
p_l]=L^{-1}b[\mathfrak B_{l-1}]$, hence the increments and their telescoping sum (Section~\YMMref{43.3}{329});
and the remainder, the term of $\mathfrak B_k$ adding at most $\frac23|\gamma_G|L^{-4}\delta_0^2$ to the constant.
Temperedness, the telescoping sum, the remainder and the uniform bound of $\Psi^G_k$ follow from the theorems for
$SU(2)$ applied to the two terms, with $0\le L^{-k}-L^{-4k}\le L^{-1}$. The gauge-fixing term is a trace too; it
enters at its zero, the branch being in the axial gauge, with the Hessian $|\mathsf L_UB|^2$, $\mathsf L_U$ being
built from products, logarithms and means, so it brings no second tadpole.

A second proof walks the chain of the $SU(2)$ construction for $G$: at the trivial configuration all operators are
scalar times $1_{\mathfrak g^c}$ (Lemma~\ref{C-368-11}(ii)), so the numbers $\delta_0$, $C_0$, $C_h$ pass from
three colours to $n$.

(ii) Along one colour direction every block of the limit theorem is a scalar kernel tensor a polynomial in
$\mathrm{ad}\,H$, but for the tadpole, and a $\log\det$ sees $\mathrm{tr}(\mathrm{ad}\,H)^2$ and $\mathrm{tr}\,M_H$
alone; the two scalar contributions in the proof of Theorem~\ref{III-12-105} have the same lattice kernels.
Their estimates add, with factors $C_A/8$ and $|\gamma_G|$, giving the bounds in (ii).

(iii) For the Gaussian identity, work at the specified local critical branch. Identification of all regular critical configurations with this branch requires
Definition~\ref{C-368-12}, or (V) on its regular domain. The difference $\Theta=A^\eta-L^k\mathcal S_k\circ\mathcal M^k$ of the action
$A^\eta=\eta^{-1}A=L^kA$, written in the units of level $k$, and the effective action of the averaged field is
gauge invariant, vanishes on the branch and is stationary there. Its
second-order expansion identifies the quadratic form at this branch with the Hessian of $L^k\mathcal S_k$ plus the
gauge-fixing term on the fibre $\{\mathcal M=W\}$, a Schur complement. At $W=\mathbf 1$ the pivot coefficients are
$L^{-3}\,\mathrm{id}$ and the fluctuation operator plus $Q^{\mathrm T}Q$ is positive definite, so everything is
formed near $\mathbf 1$; the Gaussian integral at the point of the branch gives the powers of $L$, with $nn_l$ for
$3n_l$ and $d_l=2n(L^3-1)s^3$; $\mathrm{Log}\det$ is taken at the radius $r_n$ (Lemma~\ref{C-368-10}(iv)); and the
value is the same on the orbit, the mean being covariant under all of $G\times G$ and the pivot determinants being
those of their restrictions to $\mathfrak g$. The step of the $SU(2)$ construction that $|1-T|\le\frac12$ keeps
$\det T$ off the negative axis (the three-dimensional determinant argument in the proof of Proposition~\YMMref{50.2}{386}) fails from $n=6$ and is replaced by the radius $r_n$.

The localised Gaussian family is formed for $G$ near the trivial triple from the operators $\mathsf A^\sharp_Y$ of
the regions $Y$, of whose $\mathrm{tr}\log$ its members are the differences over regions; at complex fields near
$\mathbf 1$ these operators are holomorphic germs, built from means, averages, logarithms, the action and its
derivatives, and so are the members. The coefficient of the family reads of it the terms of second order at the
trivial triple alone, in the directions $f\,t_a$, $f$ scalar; these read the lattice, $L$ and the pair and none of
the data, so the coefficient is one number for all data. Along an abelian direction $\theta H$ every operator of
which $\mathsf A^\sharp_Y$ is made is $o_0\otimes1+o_1\otimes\mathrm{ad}(\theta H)+o_2\otimes(\mathrm{ad}(\theta
H))^2+o_\tau\otimes M_{\theta H}+O(\theta^3)$ with matrices $o$ that read no group, as in the proof of (i); so the
term of second order for $G$ is $\frac{C_A}8$ times that of $SU(2)$ plus $\gamma_G$ times that of one gauge
invariant family of $SU(2)$, formed with the Wilson action with the numerical weights of the tadpole on the
plaquettes, to which the statements of the second part at $SU(2)$ apply as they stand. The same splitting holds for
$\log Z^{(k)}$ by (i) and the identity above; and at $SU(2)$ the coefficient of the localised family is that of
$\log Z^{(k)}$ (equation~\YMMref{(790)}{386} identifies their totals up to a constant, and Theorem~\YMMref{X}{427} identifies the coefficient of such totals), for both families. This gives the Hessian of the family at the
trivial triple and its coefficient, $\frac{C_A}8b^0_{k+1}+\gamma_G(j_{k+1}-L^{-1}j_k)$, the scalar family having
the coefficient $j_{k+1}-L^{-1}j_k=b[\mathfrak B_{k+1}]-L^{-1}b[\mathfrak B_k]$ by (i); this is the germ of second
order of the family at the trivial triple, which the proof of Proposition~\ref{C-368-22}(iii) computes in the same
way, by the sectors of (i); on an ideal $C^{(i)}_A$, $\gamma^{(i)}$ take the places of $C_A$, $\gamma_G$, by the
splitting of (iv).

(iv) A basis split along the ideals. For weights the form $\langle\cdot,\cdot\rangle_w$, the Hessian of the
weighted action at $\mathbf 1$, takes the place of $\langle\cdot,\cdot\rangle$, and $C^{(i)}_A$, $\theta_i$,
$\gamma^{(i)}=\frac12\theta_i-\frac3{16}C^{(i)}_A$ become $C^{(i)}_A(w)$, $\theta_i(w)$, $\gamma^{(i)}(w)$; a
rational weight is a multiplicity of a representation, which brings the weighted action to an unweighted one in a
larger representation, at a coupling multiplied by a number, and both members of the display are analytic in the
weights, the rationals being dense: so it holds at all positive weights. The term with $\gamma^{(i)}(w)$ drops in
the limit, $t_k\to0$, and $C^{(i)}_A(w)\langle P_iX,P_iX\rangle_w=\mathrm{tr}(\mathrm{ad}\,P_iX)^2$, whose sum over
$i$ is $\mathrm{tr}(\mathrm{ad}\,X)^2$. For $G_i$-valued $W$ near $\mathbf 1$ the chain of the pair is that of the
simple pair of $\mathfrak g_i$; the linear averages and the second variation split over the orthogonal sum
$\bigoplus_l\mathfrak g_l\otimes\mathbb R^{\mathsf b_0}$, $\mathbb R^{\mathsf b_0}$ the real functions on the bonds
of $\Lambda_0$, the ideals $j\ne i$ seeing the plaquettes of $U$ only through scalar weights; so the normaliser
factorises into that of the simple pair times $n_j$-th powers of scalar Gaussian integrals, whose second derivative
is the display. These scalar integrals are those of the tadpole of (i), whose propagator is the projection onto
$d(\ker\mathsf Q^{[k]})$ with the diagonal $\mathsf v_k=w_k$; so $\Upsilon_k=\frac12\Pi^{\mathfrak B_k}$, which is
real, symmetric, vanishes on closed forms, is invariant under the isometries and is tempered by (i), with the
coefficient $\frac12j_k=t_k$, $j_k=\frac23(L^{-k}-L^{-4k})$ being that of $\mathfrak B_k$ by (i), (iii) and
Theorem~\ref{C-368-4}(c); $0\le\mathsf v_k\le1$ gives $0\le\Upsilon_k\le\frac12\mathsf h_k$, with $b[\mathsf
h_k]=L^{-k}$, the composite $A\circ p_1\circ\dots\circ p_k$ having the coefficient $L^{-k}$. Adding the coefficient
$\frac{C^{(i)}_A}8b[\Psi_k]+(n_i\tau_{ii}-\frac38C^{(i)}_A)t_k$ of the simple pair and $(\sum_{j\ne
i}n_j\tau_{ij})t_k$ gives the formula with $\theta_i$. The bound ideal by ideal is that of (i) with $C^{(i)}_A$,
$\theta_i$ for $C_A$, $C_\rho$, as $0\le t_k\le\frac13L^{-1}$. At $SU(2)\times SU(2)$ a unit $H$ of an ideal has a
factor $h\in\mathfrak{su}(2)$ and $C^{(i)}_A=\mathrm{tr}(\mathrm{ad}\,h)^2=4\,\mathrm{tr}\,h^2$, which is $16$ on
$\mathbb C^2\oplus\mathbb C^2$, where $\mathrm{tr}\,h^2=4$, and $8$ on $\mathbb C^2\otimes\mathbb C^2$, where
$\mathrm{tr}\,h^2=2$; in both $\Omega=6\cdot\mathbf 1$, so $\theta_i=6$ on both ideals.
\end{proof}

\subsection{Hessians and conjugation at complex fields}

\begin{lemma}[The Hessian at complex fields]\label{C-368-14}
Let $U$ be $SL(N,\mathbb C)$-valued, $p$ a plaquette with its bonds $b_1,\dots,b_4$ numbered as in the $SU(2)$
construction, $U_i=U(b_i)$ and $Y_p=\partial U(p)=U_1U_2U_3^{-1}U_4^{-1}$; for an $\mathfrak{sl}(N,\mathbb
C)$-valued form $A$ let $a_1,\dots,a_4$ be its values on the bonds of $p$ transported to the base point along the
plaquette in its order,
\[a_1=A(b_1),\quad a_2=\mathrm{Ad}(U_1)A(b_2),\quad a_3=\mathrm{Ad}(U_1U_2U_3^{-1})A(b_3),\quad a_4=\mathrm{Ad}(Y_p)A(b_4),\]
so that $\partial(e^{it\xi A}U)(p)=e^{it\xi a_1}e^{it\xi a_2}e^{-it\xi a_3}e^{-it\xi a_4}Y_p$; with the signs
$\epsilon=(1,1,-1,-1)$ let $S_p=\sum_i\epsilon_ia_i$. Here and below $|\mathrm{Ad}\,g|=|g|\,|g^{-1}|$ is the norm
of $X\mapsto gXg^{-1}$ on all $N\times N$ matrices. With no condition on $U$, the Hessian of the action of $SU(N)$
is
\[(A,\Delta(U)A)=\xi^3\sum_p\Bigl\{\tfrac1N\mathrm{tr}\bigl(\xi^{-1}S_p\,c(Y_p)\,\xi^{-1}S_p\bigr)+i\sum_{i<j}\epsilon_i\epsilon_j\langle[a_i,a_j],F(p)\rangle\Bigr\}.\]
Let $f^\circ(b)=\xi^{-2}\max_{p\ni b}|Y_p-\mathbf 1|$, and $f(b)$ the same maximum over the plaquettes with a
corner within sup distance $\xi$ of $b_-$. If $|\mathrm{Ad}(U(b))^{\pm1}|\le2$ at all bonds and $|Y_p-\mathbf
1|\le\frac12$, the kernel $\Delta(U)(b,b')$ vanishes unless $b$ and $b'$ lie in a common plaquette,
$\xi^2|\Delta(U)(b,b')|\le259$, and on forms on a region $\mathsf R$ interior to the ambient one
\[
\begin{multlined}
\Delta_{\mathsf R}(U)+DD^{\mathrm T}=\Delta^\circ+Z(U),\qquad|(Z(U)A)(b)|\le1346f^\circ(b)\max|A|\le3456f(b)\max|A|, \\
Z(\mathbf 1)=0,
\end{multlined}
\]
$\Delta_{\mathsf R}$ the compression of $\Delta(U)$ to the forms on $\mathsf R$ and $\Delta^\circ$ the covariant
Dirichlet Laplacian, $(\Delta^\circ\phi)(x)=\xi^{-2}[6\phi(x)-\sum_\nu(\mathrm{Ad}_\nu(x)\phi(x+\xi
e_\nu)+\mathrm{Ad}_\nu(x-\xi e_\nu)^{-1}\phi(x-\xi e_\nu))]$, $\mathrm{Ad}_\nu(x)=\mathrm{Ad}(U(x,x+\xi e_\nu))$.
Under the guard $|Y_p-\mathbf 1|\le1$ of the $SU(2)$ construction the kernel is at most $180\xi^{-2}$ and
$|(Z(U)A)(b)|\le868f(b)\max|A|$; so the numerals $1536$ and $3456$ of the $SU(2)$ construction serve in all its
clauses. If $U$ is $SU(N)$-valued with $|Y_p-\mathbf 1|\le\sqrt2$, then $-12\sup
f^\circ\,(A,A)\le(A,\Delta(U)A)\le(16\xi^{-2}+12\sup f^\circ)(A,A)$ for $\mathfrak{su}(N)$-valued $A$. The
constants are absolute, the same for every $N\ge2$. Let $\mathfrak g$ be semisimple, so that $G\subset SU(N)$ and
$G^c\subset SL(N,\mathbb C)$, and let $U$ be $G^c$-valued. Then the Hessian of the action of $G$ is this form
restricted to $\mathfrak g^c$-valued forms, and its operator and its Weitzenb\"ock term are those of $SU(N)$
followed by $P$, the orthogonal projection onto $\mathfrak g^c$ for $\mathrm{tr}\,X^*Y$: the bounds on the
quadratic form and the bounds in the normalised trace norm stay, and the kernel and Weitzenb\"ock bounds in
$|\cdot|$ hold with the factor $|P|\le N^{1/2}$, the norm of $P$ for $|\cdot|$, the condition
$|\mathrm{Ad}(U(b))^{\pm1}|\le2$ being read on all $N\times N$ matrices. Read with $\mathrm{Ad}$ on $\mathfrak g^c$
alone, the clauses under $|Y_p-\mathbf 1|\le\frac12$ stand, with $|Y_p^{-1}|\le2$ by the Neumann series, and under
the guard $|Y_p-\mathbf 1|\le1$ the inverse obeys $|Y_p^{-1}|\le|Y_p|^{N-1}\le2^{N-1}$, a number of the pair, in
the place of $16$.
\end{lemma}

\begin{proof}[Sketch of proof]
The identity is the second derivative of $a(e^{it\xi a_1}\cdots)$ written with $c(Y_p)$ and $H_{Y_p}$; in operator
form the bilinear maps $N_{ij}$, such as $X\mapsto\frac12(XY_p+Y_p^{-1}X)$, have norm at most
$\frac12(|Y_p|+|Y_p^{-1}|)$ and differ from the identity by at most $\frac12(|Y_p-\mathbf 1|+|Y_p^{-1}-\mathbf
1|)$.

The inverse of a plaquette variable is the one new quantity, and $|X^{-1}|\le|\mathrm{Ad}\,X|$ on $SL(N,\mathbb
C)$, $\mathrm{Ad}$ acting on all $N\times N$ matrices (Lemma~\ref{C-368-6}(ii)), bounds it with no $N$:
$|\mathrm{Ad}(U(b))^{\pm1}|\le2$ at four bonds gives $|Y_p^{-1}|\le|\mathrm{Ad}\,Y_p|\le16$; on $\mathfrak g^c$
alone the inequality $|X^{-1}|\le|\mathrm{Ad}\,X|$ fails, as at the representation of spin $2$ of $SU(2)$ at
$X=e^{(\log2)J_3}$, where $|X^{-1}|=4$ and the norm of $\mathrm{Ad}\,X$ on $\mathfrak g^c$ is $2$.

Before the projection $P$ the count of the kernel is the same in $|\cdot|$ and in $\|\cdot\|$:
$|N_{ij}|\le\frac12(\frac32+2)=\frac74$, the transports give $|R_i^{-1}||R_i|=1$ at the first position of a bond in
a plaquette and at most $(3\cdot2)^2=36$ at the third, $R_3=\mathrm{Ad}(Y_pU_4)$, and the four plaquettes of a bond
hold it at most at the positions $1,3,1,3$, as for a bond of the first direction: so
$\xi^2|\Delta(U)(b,b')|\le\frac74\cdot74=129.5$ before $P$; $P$ has the norm $1$ for $\|\cdot\|$ and $2-\frac2N<2$
for $|\cdot|$, whence $259$ in $|\cdot|$, with no passage between the two norms; the formula of Weitzenb\"ock, with
the covariant Laplacian, which is the operator of $U$ itself and not the free one, gives $1346$, and $f^\circ\le
f$.

Under the guard $|Y_p-\mathbf 1|\le1$ the finer count, in which $\mathrm{Ad}(Y_p)$ drops out of every block, gives $180$ and $868$.

With the eigen-angles of a unitary $Y_p$ the bounds on the form follow as in the $SU(2)$ construction. For
semisimple $\mathfrak g$ the action of $G$ is the restriction of that of $SU(N)$ to $G^c$-valued fields,
$\mathrm{Ad}(U(b))$, $DD^{\mathrm T}$ and the covariant Laplacian keep $\mathfrak g^c$-valued forms, and
$\langle[a_i,a_j],F(p)\rangle=\langle[a_i,a_j],PF(p)\rangle$, $\langle A,PB\rangle=\langle A,B\rangle$ for
$A\in\mathfrak g^c$; so the form is the restriction, and the operator and $Z$ are $P$ times those of $SU(N)$. As
$P$ is orthogonal for $\mathrm{tr}\,X^*Y$, the form bounds and the bounds in the trace norm stay; in $|\cdot|$,
$|PX|\le(\mathrm{tr}(PX)^*PX)^{1/2}\le N^{1/2}|X|$.
\end{proof}

\begin{lemma}[Conjugation]\label{C-368-24}
\begin{enumerate}
\item[(i)] $X\mapsto X^c=(X^*)^{-1}$ is an antiholomorphic involutive automorphism of $GL(N,\mathbb C)$ whose fixed
points are the unitary matrices; it keeps $G^c$; $X^c-\mathbf 1=(X^{-1}-\mathbf 1)^*$; $\partial(U^c)(p)=\partial
U(p)^c$, $H_{X^c}=H_X^*$, and $F$, $J$ and covariant differences are covariant likewise; $a(X^c)=\overline{a(X)}$.
\item[(ii)] For a set $Q$ of points, $Q^s=Q\cap Q^c$ is the largest subset kept by the conjugation; it holds the
real points, commutes with intersections and is open with $Q$; if $Q=\{|X_\alpha(p)-\mathbf 1|<r_\alpha\}$ with
covariant $X_\alpha$ of group type, then $Q^s=\{m(X_\alpha)<r_\alpha\}$. $m(X^{-1})=m(X^*)=m(X^c)=m(gXg^{-1})=m(X)$
for unitary $g$, and $x\le m(X)\le\frac x{1-x}$ for $x=|X-\mathbf 1|<1$. If $|Y-\mathbf 1|\le\varphi$ on a set kept
by the conjugation, $Y$ covariant and $\varphi$ symmetric, then $m(Y)\le\varphi$. If $m(X)<1$, then $\log
X^{-1}=-\log X$ and $\frac1i\log X^c=(\frac1i\log X)^*$; and $\mathrm{av}(V^c)=\mathrm{av}(V)^c$.
\item[(iii)] The chain of objects of a step from the fields to the fluctuation operator commutes with the
conjugation wherever both sides are formed, $A^\sharp(p^c)=\theta A^\sharp(p)\theta^{-1}$ with $\theta Z=Z^*$, on
the domain of the $SU(2)$ construction written with $m(X)<1$ and with $X\in D^c$ at the logarithms of the averages,
the latter condition being void for $SU(N)$ with $N\le6$. On this domain the symmetric action $a(X)$ suffices while
the proviso on means of Lemma~\ref{C-368-6}(iv), the contour quotients in $D^c$, is void, that is for
$z_G>\frac1{40}$, at $SU(N)$ exactly for $N\le251$; from $N=252$ on the quotients in $D^c$ are added to the domain.
This applies in particular to the intrinsic cube response of
Proposition~\YMMref{54.2}{469}, with the uniqueness ball and radii of
Proposition~\YMMref{54.1}{468}; to the unparameterized response obtained
there by setting $s=\mathbf1$ and $\varkappa=0$; and to the potential
and axial gauge $(A,h)$ of Theorem~\YMMref{E}{10}, constructed in
Section~\YMMref{38}{288}. Each is taken on its stated domain.
For another group the same covariance assertion applies wherever the
corresponding construction and uniqueness bounds are supplied under
\textup{(V)}; it does not supply those bounds.
\item[(iv)] For $SU(N)$, $N\ge3$, the moduli are not kept: with $H=\mathrm{diag}(1,\dots,1,-(N-1))$ and $s>0$,
$e^{-sH}$ is the conjugate of $e^{sH}$ and $|e^{sH}-\mathbf
1|=\max\{e^s-1,1-e^{-(N-1)s}\}<e^{(N-1)s}-1=|e^{-sH}-\mathbf 1|$, at $SU(3)$ with $H=\mathrm{diag}(1,1,-2)$; so
neither the analyticity cores of the $SU(2)$ construction nor the moduli of the quantities of its towers are kept
by $X\mapsto X^c$ unless the conditions are written with $m(X)$.
\end{enumerate}
\end{lemma}

\begin{proof}[Sketch of proof]
(i) is direct, $G^c$ being closed under adjoints and inverses; all of it says that $X\mapsto(X,X^{-1})$ carries the
conjugation to the involution $\Theta(E,F)=(F^*,E^*)$ of the ring $M_N\times M_N^{\rm op}$, of which $m$, $a$ and
$H_X$ are a norm, a trace and the odd part. (ii) The adjoint keeps the operator norm, so
$m(X^c)=\max\{|X^{-1}-\mathbf 1|,|X-\mathbf 1|\}$; the bound by $\frac x{1-x}$ is Lemma~\ref{C-368-6}(ii); for a
covariant $Y$ the bound at the point and at the conjugate point gives both members of $m$; the logarithms follow
from the series, and the mean from its uniqueness (Lemma~\ref{C-368-6}(iv)). (iii) Products, inverses, means, logarithms, the action $a$ and
its derivatives commute with the stated conjugation by (i), (ii).
For the regional response retain the two equations
\YMMref{(1051)}{468}, with the regional substitutions
\YMMref{(1052)}{469}. In particular the identity in $S-I$ remains
untruncated. The current is extended by zero, and the nonlinear maps
and every surviving walk term use only the pair in the region.
The proof of Proposition~\YMMref{54.2}{469} identifies the regional
first-difference norm on the response carrier with that for any
continuation. Thus conjugation carries both equations and their
uniqueness ball to the equations and ball at the conjugate data:
\[
 \max\{\|T\|,|\lambda-H^{\mathfrak k}(s)^{\mathrm T}
          \mathbf J^0|_\infty\}
 \le C_{\mathrm D,\mathrm s}\varepsilon_{\mathrm D,\mathrm s}
                       e^{-8\varkappa}.
\]
The parameter and datum restrictions remain
$|s|\le e^\varkappa$ and
$|B'|_\infty,\gamma_0\le
\varepsilon_{\mathrm D,\mathrm s}e^{-16\varkappa}$.
Here the field and multiplier are conjugated by $Z\mapsto Z^*$,
which preserves $|Z|$ and intertwines the covariant differences;
the group-valued domain restrictions use $m$. Uniqueness gives the
claimed response identity. The same argument at $s=\mathbf1$ uses
the unparameterized bounds, with $\varkappa=0$.

For the canonical pair, the inverse datum twist
\YMMref{(580)}{289} commutes with conjugation. Its constrained response
is the one of Theorem~\YMMref{K}{269}, and the inverse formula for
its multiplier in Section~\YMMref{38}{288} consists of covariant
operations. Finally the pillar recursion \YMMref{(584)}{290}
constructs $h$ from the response by products, inverses and means.
These operations give $A(B^*)=A(B)^*$ and $h(B^*)=h(B)^c$ at the
real reference of that construction. Every radius used here is the radius of the cited
construction, restricted further by this lemma's two-sided guards.
For $G$, the same calculation applies to the supplied counterparts
under \textup{(V)}. (iv) is a computation with the eigenvalues.
\end{proof}

\begin{lemma}[Reflection on the real-connected part]\label{C-368-15}
Let $Q$ be a set of points $p=(\varrho,v)$, $\varrho$ real parameters and $v$ complex fields and numbers; let
$p\mapsto p^c$ act by $X\mapsto X^c$ on group-valued members, by $Z\mapsto Z^*$ on $\mathfrak g^c$-valued ones and
by complex conjugation on numbers; $p$ is real if $p^c=p$. Let $Q^s=\{p\in Q:p^c\in Q\}$, and let $Q^r$, the
\emph{real-connected part}, be the set of $p\in Q^s$ whose connected component in the slice $(Q^s)_\varrho$ holds a
real point.
\begin{enumerate}
\item[(i)] $Q^r$ is kept by the conjugation, holds the real points of $Q$, is monotone in $Q$, has open slices if
$Q$ has, and is the largest subset of $Q^s$ every slice component of which holds a real point; polydiscs and
products pass; $Q^r=Q^s$ if the slices of $Q^s$ are star-shaped about real points, or if the slices of $Q$ are
convex and have no group-valued member, the conjugation being then an antilinear involution; the convex clause asks
for this last condition, as at $G=U(2)$, where the segment from $\mathrm{diag}(2,3)$ to
$\mathrm{diag}(\frac12,\frac13)$ has conjugate ends and keeps the distance $0.08$ from $U(2)$, its open
$0.05$-neighbourhood $Q$ in $|\cdot|$, which is convex, lies in $GL(2,\mathbb C)$ and holds no unitary matrix, has
$Q^s\ne\emptyset$ and $Q^r=\emptyset$; and $(\bigcap_iQ_i)^r\subset\bigcap_iQ_i^r$, possibly strictly.
\item[(ii)] Let the slices of $Q_1^r$ be open and let $\Phi$ be holomorphic in $v$ there, with real members
constant on slice components. Then $\Phi(p^c)=\Phi(p)^c$ on $Q_1^r$ if and only if $\Phi(p)$ is real for every real
$p\in Q_1$; and then $\Phi(Q_1^r)\subset Q_2$ implies $\Phi(Q_1^r)\subset Q_2^r$. Composites of such maps and of
continuous covariant maps are covariant; a composite of maps built from products, inverses, means and logarithms
and of continuations on convex subsets of a vector space kept by the conjugation (slices without group-valued
member) is covariant on all of $Q^s$, and so are the towers wherever they exist at both points. If the slices of
$Q$ are open, the function equal to $i$ on $Q^s\setminus Q^r$ and to $0$ on $Q^r$ is holomorphic, real at real
points, and not covariant at any point of $Q^s\setminus Q^r$.
\item[(iii)] Let $\mathcal Q(b_0,b_1,c_0)$ be the set of the $SU(2)$ construction of the triples $(U,A',\mathbf
J)$, $U$ real with $|\partial U(p)-\mathbf 1|<b_0\xi^2$, $A'$ with $|A'|$ and its covariant differences below
$b_1$, $|\partial(e^{i\xi A'}U)(p)-\mathbf 1|<b_0\xi^2$ at the plaquettes and $|\mathbf J|<c_0$; its symmetric part
has $m(\partial(e^{i\xi A'}U)(p))$ in the place of the one-sided distance. Then
\[\mathcal Q(b_0,b_1,c_0)^s\subset\mathcal Q\bigl(b_0+2e^{4\xi b_1}b_1^2,b_1,c_0\bigr)^r,\]
the ray $(U,tA',t\mathbf J)$, $0\le t\le1$, joining each point to the real point $(U,0,0)$ inside the larger set.
For the classes of regular pairs the same holds with $\alpha_0+2\vartheta^2e^{4\vartheta\alpha_1}\alpha_1^2$ in the
place of $\alpha_0$, $\vartheta\ge1$ the largest ratio of the spacings around a plaquette, and the like holds for
the natural classes. For the analyticity cores, $\mathcal U^\circ(b_0,b_1,c_0)^s\subset\mathcal
U^\circ(B_0,B_1,c_0)^r$ whenever $b'_0=b_0+2e^{4\xi b_1}b_1^2\le B_0$, $b_1\le B_1$, $K(b'_0+b_1^2)\le B_0$ and
$b'_0+b_1\le a^\sharp$, where $K\ge16L^2$ and $a^\sharp$ are the constant and the radius of the inclusion of the
simple sets in the cores, in its symmetric form, assumed for $G$ through Proposition~\ref{C-368-16}(ii) and the
uniqueness clause of Definition~\ref{C-368-31}(V). With $Q^r$ in the place of $Q^s$ no number moves but at a
membership reached from memberships of restrictions had apart: such a membership is reached in $Q^s$, and the
inclusions above carry it into a real-connected part, at the radii named there. The symmetric parts of the domains
of the $SU(2)$ construction are cut out by conditions $m(X)<r$ and commute with intersections
(Lemma~\ref{C-368-24}(ii)); the real-connected parts pass to intersections as the inclusion
$(\bigcap_iQ_i)^r\subset\bigcap_iQ_i^r$ of (i).
\end{enumerate}
\end{lemma}

\begin{proof}[Sketch of proof]
(i) is set topology. In the instance at $U(2)$ the distance in $|\cdot|$ of a positive diagonal matrix from $U(2)$
is the largest distance of its entries from $1$, and from the singular matrices its least entry, singular values
moving by at most the norm of a difference; on the segment, $\mathrm{diag}(2-\frac32t,3-\frac83t)$ with $0\le
t\le1$, the first is least at $t=\frac{18}{25}$, where it is $0.08$, and the second is at least $\frac13$. (ii)
$p\mapsto\Phi(p^c)^c$ is holomorphic in $v$, as a composite of two antiholomorphic maps, and agrees with $\Phi$ at
the real points, which form in each slice a totally real submanifold of maximal dimension; by the identity theorem
the two agree on every component that holds a real point. No formula of $\Phi$ is used, so sums of group elements,
complex constants and maps defined as the one solution in a ball are covered. Conversely covariance at a real point
says that its image is real. (iii) Along the ray $tA'$ a product of four exponentials differs from its chord by
\[|X_t-(1-t)X_0-tX_1|\le8\,t(1-t)\,\xi^2d^2e^{4\xi b_1},\]
where $d$ is the maximum operator norm of the four transported parameter increments $A'$ in this plaquette product, so $d<b_1$; the same estimate holds for $X_t^{-1}$. Thus
$m(X_t)\le(1-t)m(X_0)+tm(X_1)+8t(1-t)\xi^2d^2e^{4\xi b_1}$, and with $X_0=\partial U(p)$, $m(X_0)=|\partial
U(p)-\mathbf 1|<b_0\xi^2$, $m(X_1)<b_0\xi^2$ at a point of the symmetric part and $d<b_1$ this gives
$m(X_t)<b_0\xi^2+2e^{4\xi b_1}b_1^2\xi^2$, as $8t(1-t)\le2$, the bounds on $tA'$ and $t\mathbf J$ being smaller
than those on $A'$, $\mathbf J$; so the ray from $(U,A',\mathbf J)$ to the real point $(U,0,0)$ stays in the larger
set. For the regular pairs the ray is the same, the factor $\vartheta$ entering through the spacings of the
plaquettes. For the cores, $\mathcal U^\circ(b_0,b_1,c_0)^s\subset\mathcal Q(b_0,b_1,c_0)^s$, and the ray lies in
$\mathcal Q(b'_0,b_1,c_0)^s$, which lies in $\mathcal U^\circ(B_0,B_1,c_0)^s$ by the inclusion of the simple sets
under the conditions named, and it ends at a real point.
\end{proof}


\begin{proposition}[Fields off the trivial field]\label{C-368-16}
Let $\mathfrak g$ be simple, $L\ge13$ odd, $200L^2z_G\ge1$; constants are constants depending only on $L$ and the
pair, uniform in the volume. In (i), the completed propagator construction assumes (P). The uniform critical-background, tower and core assertions in (ii), as well as (iii)--(v), assume (V).
The explicit Hessian identities and bounds, the averaging formulas on their guarded domains, and holomorphic germs at each fixed finite tower are local statements.
A fixed finite-dimensional germ does not assert a radius or a nonlinear majorant uniform in depth and volume.
\begin{enumerate}
\item[(i)] Let $G=SU(N)$, $N\ge2$ arbitrary. The Hessian is
given by Lemma~\ref{C-368-14}, the current by
Lemma~\ref{C-368-8}(i), and the trivial-field operators by
Lemma~\ref{C-368-11}. Their gauge covariance follows from their
transport formulas. The remaining quantitative inputs---the local
action remainders, difference and commutator estimates in the natural
norms, their variable-layer forms, the local gauges, and the
order-zero inverse at rate $\frac14\log(37/36)$---are included in (P).
For $SU(2)$ the local nonlinear estimates are supplied by
Propositions~\YMMref{30.1}{236} and~\YMMref{33.1}{259};
the completed inverse is Theorem~\YMMref{G}{204}.
The following Wilson-Hessian bounds are proved directly below. In the Wilson part $\Delta^W_U=D^{\mathrm
T}D+\Delta'(U)$ the operator $\Delta'(U)$ is local, symmetric and real on real forms, with row sums at most
$18\alpha\mathsf L_{\mathbf B}(b)^{-2}$ in $|\cdot|$ and $20\alpha\mathsf L_{\mathbf B}(b)^{-2}$ in $\|\cdot\|$ for
$U$ in the class $\mathfrak U_k(\mathbf B,\alpha)$ of Definition~\ref{def:group-regular-classes} with $\alpha\le\frac38$, and, for $U$ in
that class at every $\alpha$, $(A,\Delta^W_UA)\ge\eta^3\sum\|DA\|^2-32\alpha\eta^3\sum_b\mathsf L_{\mathbf
B}(b)^{-2}\|A(b)\|^2$, with no $N$, $\mathsf L_{\mathbf B}$ being the local block side of the determining set
$\mathbf B$ (Definition~\YMMref{9.3}{40}): $\mathsf L_{\mathbf B}(\mathfrak e)=L^{j-k}$, $j$ the largest index
for which the bond or plaquette $\mathfrak e$ has an end point or a corner in $\Omega_j$, so that $\eta^2\mathsf
L_{\mathbf B}(\mathfrak e)^{-2}=L^{-2j}$; and $|\tau^*-U(\Gamma)|\le\kappa_7\alpha_0$ for the transports of the
chains, with the transport-comparison constant $\kappa_7(L)$
in the proof of Lemma~\YMMref{30.3}{240}. The same path and mean
comparisons apply on the guarded domains of Lemma~\ref{C-368-7}.
The completed propagator and covariance estimates are hypothesis (P) of
Definition~\ref{C-368-31}; bounds of them in $|\cdot|$ give the same bounds in the normalised trace norm, with the
factors $(N-1)^{1/2}$ and $N^{1/2}$ between the two (Lemma~\ref{C-368-6}(i)). For the other simple $G$ of the
proposition Lemma~\ref{C-368-14} gives the Hessian among these inputs, with the factor $|P|$ in its kernel and
Weitzenb\"ock bounds.
\item[(ii)] For every $G$ of the proposition, averages exist, with logarithms in $\mathfrak g^c$, at every regular
real background. Over every real $e^{iB}$ with $\sup|B|$ small there is a regular critical configuration with a
holomorphic continuation $W(B)$ on a polydisc whose radius depends on no scale, and $|\partial W-\mathbf 1|\le
C\sup|B|\,\eta^2\mathsf L_{\mathbf B}^{-2}$, $|J(W)|\le C\sup|B|\,\mathsf L_{\mathbf B}^{-3}$ at the plaquettes and
the bonds, $\mathbf B$ the determining set and $\mathsf L_{\mathbf B}$ its local block side, as in (i); the
flatness $L^{-2q}$ of towers of depth $q$, with plaquettes within $cL^{-2q}$ of $\mathbf 1$ at every truncated depth, is included in this hypothesis. The linearised tower kernel is independently the same real scalar kernel
for all $G$; control of the nonlinear remainder at this scale is required in (V). Under the uniform estimates of
Proposition~\ref{cg-near-trivial-uniqueness}, small fields lie in the analyticity cores of $G$ uniformly in the
scale, the sub-domain, the cutoff and the volume.
\item[(iii)] The localisation input in (V) is the following
counterpart of Proposition~\YMMref{53.5}{447}. Use exactly the
separated joined sequence, the one or two disjoint wide collars,
the regular original and joined data, and the torus-size hypothesis
of Section~\YMMref{53.4}{445}. For arbitrary elements $U,U'$ of
the two minimal orbits, require gauges $w,w'$ and a real odd field $a$
such that $(U|_\Omega)^w=e^{ia}(U')^{w'}$, and, outside the wide collars,
\[
 |a(b)|\le\mathsf C_{\rm lc}\varepsilon_{\rm r}
       L^{-j(b)}e^{-\delta_{\rm lc}D(b)/2},\qquad
 w'(y)=1,\quad |w(y)-1|\le\mathsf C_{\rm lc}\varepsilon_{\rm r}
       e^{-\delta_{\rm lc}D(y)/2}
\]
at the nodes $y$. Here $j(b)$ is the largest level containing an
endpoint, and $D$ is the weighted distance to the endpoints of the
joined tower bonds, with $D=+\infty$ when there are none. The constants
are independent of depth and volume. Retain all five entries of the
smallness ceiling in equation~\YMMref{(1000)}{446}, evaluated at the
admitted constants of the pair, as well as its geometric floors and
response hypotheses. This is the comparison used in (iv); no
localisation estimate for a different joining geometry is inferred.
\item[(iv)] Let the localisation of (iii) hold for $G$, and let the objects of a step be formed for $G$ at the
fields over a regular base, which is part of (V): the global and the localised background of the step are defined
there; the minimal orbits of the three data of the step exist and have elements whose $k$-fold averages are in the
axial gauge; their elements $U'$ have $|\partial U'(p)-\mathbf 1|<2B_3\varepsilon_tL^{-2j}$ at the inner bare
plaquettes $p$, $j$ the largest index for which $p$ has a corner in $\Omega'_j$; and the data lie in the regular
sets of radius $2\varepsilon_t$. For simply connected $G$ they follow under (V), by (iii). The zone-bond conditions use the following explicit ceilings. The background constants are those admitted in (V), and $c_{\rm Q}$, $\mathsf c_{\rm s}$ and $\mathsf c_{\rm fl}$ have the formulas in equation~\YMMref{(1026)}{457}, evaluated at the constants of the pair. Put
\begin{gather*}
 \epsilon_{\rm Q}=\min\{\varepsilon_\flat,\bar\varepsilon/2,
 c^{\rm SO}_2/(2B_3),\omega_{\rm p}/(2+4B_3L^{-2}),1/520\},\\
 \epsilon_{\rm fl}=\min\{\epsilon_{\rm Q},(1+\mathsf c_{\rm s})^{-1}\},\qquad
 \mathsf c_{\rm H}=\pi\kappa_{\mathsf c}\sqrt{2L}
                  \max\{2\bar C_{\rm reg},4B_3\},\\
 \epsilon_{\rm H}=\min\{\epsilon_{\rm fl},c^{\rm SO}_2/(2B_3),
       \pi/[\pi(1+\mathsf c_{\rm s})+\mathsf c_{\rm H}]\},\\
 \epsilon_{\rm cu}=\min\{\epsilon_{\rm lc}/2,
             c^{\rm g}_2/(2B_3),c_4/(4\mathsf C_{\rm lc})\},\qquad
 \mathsf c_L=4+2C_v+4\kappa_6,\\
 \epsilon_{\rm xp}=\min\{\epsilon_{\rm H},\epsilon_{\rm cu},
 \epsilon_{\rm lc}/2,c^{\rm g}_2/(2L^2B_3),
 c_4/(4L\mathsf C_{\rm lc}),
 (2\mathsf c_L\mathsf C_{\rm lc})^{-1},(2+\mathsf c_{\rm s})^{-1}\}.
\end{gather*}
Here $\bar C_{\rm reg}$ is the coupling-independent regularity bound, $\varepsilon_\flat$ is the background radius, $\kappa_{\mathsf c}$ is the contour-comparison constant, and $c^{\rm g}_2,c_4,C_v$ are the curvature, perturbation and gauge-factor constants of the centered averaging estimate in Lemma~\ref{C-368-7}(v). Assume $\varepsilon_t\le\epsilon_{\rm xp}$; $A_1/A_0\le c_{\rm Q}$; $\gamma\le\gamma^G_{\rm xp}$, where $\gamma^G_{\rm xp}=e^{-y^G_{\rm xp}/2}$ and $y^G_{\rm xp}$ is formed by equation~\YMMref{(1016)}{452} with $\|\vartheta_G'\|_\infty=\lambda\|\vartheta_{\rm co}'\|_\infty$ of
Lemma~\ref{C-368-10}(ii) in the place of $\|\vartheta_{\rm co}'\|_\infty$ in $\Upsilon_1$, the one formula of the
chain that changes; and $C_\chi\mathsf c_{\rm fl}\delta_k\le r_\chi\le1\le C_\chi$, with the chart constants $r_\chi,C_\chi$ of Section~\ref{sec:sf-chart-interface}, whose transferred bounds are among the construction hypotheses. Here $\varepsilon_t$ and $\delta_k$ are the thresholds of the levels $t$
and $k=t-1$ of the step, $g_k$ the running coupling and $\check R_t$ the cube size (equation~\YMMref{(634)}{339},
\YMMref{(634)}{339}), $\mathfrak b_k=\delta_k/g_k$ is the radius of the scaled fluctuation field
(equation~\YMMref{(634)}{339}), $\mathsf c_{\rm fl}$ is the constant of the fluctuation range
(equation~\YMMref{(1026)}{457}), $\mathsf C_{\rm lc}$, $\delta_{\rm lc}$ are those of the localisation, and $\mathsf
c_L=4+2C_v+4\kappa_6$ as above, with the radius $c_4$ and the constants $C_v$, $\kappa_6$ of the averages
over a box and of the contour variables (Lemma~\ref{C-368-7}(v), with the one-step contour estimate of Lemma~\YMMref{14.7}{66}), which for $G$ are
constants depending only on $L$ and the pair with $|f_i-\mathbf 1|\le C_vL^i\sup|A_0|$ whenever $L^j\sup|A_0|<c_4$,
in the notation of Lemma~\ref{C-368-7}(v). Put $\varsigma=2\mathsf C_{\rm lc}\varepsilon_t\exp(-\frac14\delta_{\rm
lc}M\check R_t)$; the fifth member of $\epsilon_{\rm xp}$, $4L\mathsf C_{\rm lc}\varepsilon_t\le c_4$, gives
$\varsigma<\frac{c_4}{2L}$. Let finally $\mathsf c_Lc_4\le2Lr_G$, which is void at $SU(3)$. Then at each bond $b$
of the cut-off zone the quotient $E_b$ of the global and the localised background has $|E_b-\mathbf 1|\le\mathsf
c_L\varsigma<r_G$; $H_b=\frac1i\log E_b\in\mathfrak g$, the shift of \YMMref{(989)}{444}, has
$|H_b|\le\frac\pi2\mathsf c_L\varsigma$ and $8\|\vartheta_G'\|_\infty|H_b|=8\lambda\|\vartheta_{\rm
co}'\|_\infty|H_b|\le\delta_k$; with $\theta(W)=\vartheta_G(\nu(Y(W))/\delta_k)$ where $-1$ is no eigenvalue of
$W$, $Y(W)$ the principal hermitian logarithm, and $\theta(W)=1$ elsewhere, the bond map $W\mapsto
We^{-i\theta(W)H_b}e^{iH_b}$ of Lemma~\YMMref{53.3}{444} is a diffeomorphism of $G$ onto itself, its inverse at $W$
being the fixed point of the contraction $W'\mapsto We^{-iH_b}e^{i\theta(W')H_b}$ of factor at most $\frac\pi{16}$,
with Jacobian in $[\frac78,\frac98]$; and as $g_k\mathsf c_{\rm fl}\mathfrak b_k+|H_b|=\mathsf c_{\rm
fl}\delta_k+|H_b|\le1+\frac1{24}<\pi$, the substitution $W\mapsto We^{-i\theta(W)H_b}$, $W=e^{ig_kY(b)}$, is a
homeomorphism between open sets each holding $\{g_k|Y(b)|+|H_b|<\pi\}$, outside of which the integrands vanish.
\item[(v)] The domains of analyticity are defined for $G$, are gauge stable, and are the same on all large tori and
on $\mathbb R^3$; the totals do not depend on the chosen background; and at the trivial real factor the
continuations and the bounds of the towers are those of (ii). The continuations and the towers off the trivial
field are those of (V), which are assumed in (V) for the simply connected construction.
\end{enumerate}
\end{proposition}

\begin{proof}[Sketch of proof]
(i) The local identities follow from the current and Hessian formulas.
The quantitative inverse and local-remainder bounds specified in (P)
are assumptions here. The following calculation proves the asserted
Wilson-Hessian bounds without using those assumptions.

The group enters them through the Hessian; through the commutators, $|[X,Y]|\le2|X||Y|$; through sums, products and
Neumann series of decaying kernels in $|\cdot|$, which are sums over the lattice and keep the constants of the
$SU(2)$ construction for every $N$; through the transpose of a kernel, which takes $N$ for the $2$ of $|T^{\mathrm
T}|\le2|T|$, and the bounds of Schur's kind, whose numeral $3$ becomes $n$ (Definition~\ref{C-368-3}(2),
Proposition~\ref{C-368-18}(g)); through the table at the trivial field (Lemma~\ref{C-368-11}); and through the
small-field theorem in the gauge of the class (Lemma~\ref{C-368-8}(iii)).

The Wilson part $\Delta'(U)=\Delta(U)-D^{\mathrm T}D$ is read off the identity of Lemma~\ref{C-368-14}, with
$\xi=\eta$: at a unitary $Y_p$ its terms are $\xi^{-2}\frac1N\mathrm{tr}\bigl(S_p(c(Y_p)-\mathbf 1)S_p\bigr)$ and
$i\sum_{i<j}\epsilon_i\epsilon_j\langle[a_i,a_j],F(p)\rangle$, where $|c(Y_p)-\mathbf 1|=\frac12x_p^2$ with
$x_p=|Y_p-\mathbf 1|$, and $\langle[a_i,a_j],F(p)\rangle=\xi^{-2}\langle[a_i,a_j],H_{Y_p}\rangle$ with
$|H_{Y_p}|\le x_p$, as $\langle A,PB\rangle=\langle A,B\rangle$ for $A\in\mathfrak g^c$. A bond lies in four
plaquettes and a plaquette has four bonds. So, with $x$ the largest $x_p$ at the plaquettes of $b$, which is
$\xi^2f^\circ(b)$, the row sums of $\Delta'(U)$ at $b$ are at most $(12+8x)f^\circ(b)$ in $\|\cdot\|$ and
$(12+16x)f^\circ(b)$ in $|\cdot|$, where $|P|\le2$ at $SU(N)$; in the class $\mathfrak U_k(\mathbf B,\alpha)$ one
has $f^\circ(b)\le\alpha\mathsf L_{\mathbf B}(b)^{-2}$ and $x\le\alpha$, and for $\alpha\le\frac38$ this gives
$18\alpha\mathsf L_{\mathbf B}(b)^{-2}$ in $|\cdot|$ and $15\alpha\mathsf L_{\mathbf B}(b)^{-2}\le20\alpha\mathsf
L_{\mathbf B}(b)^{-2}$ in $\|\cdot\|$. As $\Delta(U)$ and $D^{\mathrm T}D$ are Hessians of real functions at
$G$-valued $U$, $\Delta'(U)$ is real on real forms. With $\|S_p\|^2\le4\sum_i\|a_i\|^2$ the same count gives, bond
by bond, $(A,\Delta'(U)A)\ge-\eta^3\sum_b(12+8x)f^\circ(b)\|A(b)\|^2$, $x$ taken at $b$, and $12+8x\le28\le32$ at
every unitary $Y_p$, where $x\le2$: this is the lower bound. The transport-comparison constant $\kappa_7(L)$ in the proof of Lemma~\YMMref{30.3}{240}
reads no $N$ (Proposition~\ref{C-368-18}(g)).

(ii) The averages: Lemmas~\ref{C-368-6}(iii) and~\ref{C-368-7} with the rule $D^c$.

At a fixed finite tower the implicit function theorem gives a local critical germ over $e^{iB}$. The uniform radius and nonlinear majorants asserted in (ii) are part of (V), not consequences of this finite-dimensional theorem. The
linearisation is formed from the operators of Lemma~\ref{C-368-11}(ii), with Lemma~\ref{C-368-14}; a
holomorphic map that is $\mathfrak g$-valued at real points is $\mathfrak g^c$-valued. On a tower of depth $q$,
truncated at a level $q'\le q$, the towers ask at this configuration the flatness $L^{-2q}$, which a bound of first
order in $\sup|B|$ gives only up to the factor $L^{q-q'}$. The linear term has a kernel
$\hat D$ that is real, scalar and the same for all $G$: for a real scalar $1$-form $\psi$, $|\hat D[\psi]|\le K_{\rm
lin}\sup_p\mathsf L_-(p)^2|d\psi(p)|$, where $\mathsf L_-(p)$ is the least of the local scales $L^{\lambda(x)-k}$
at the corners $x$ of $p$, $\lambda(x)$ being the level of the sequence of regions at $x$ and $\eta=L^{-k}$, and
the constant $K_{\rm lin}$ of the $SU(2)$ construction is read off its nonlinear bound of the towers
along $\psi$ times a fixed unit of $\mathfrak{su}(2)$, as the factor tends to $0$ (Lemma~\ref{C-368-11}(ii));
Cauchy's estimate controls higher orders only on an admitted polydisc with a uniform majorant. The uniform nonlinear tower-flatness estimate itself remains part of (V). To lie in the analyticity core a small field asks, on every
sub-domain, for both tower rules and at every depth, that its averages exist with logarithms in $\mathfrak g^c$ and
that over each real datum near $\mathbf 1$ the regular critical configurations form one gauge orbit whose
representative continues holomorphically with the tower bounds. These uniform assertions are the near-trivial input in (V);
the slice and contraction reduction of Proposition~\ref{cg-near-trivial-uniqueness} gives the one-orbit conclusion when its uniform hypotheses hold.
The local matrix and colour-scalar linear estimates do not establish those hypotheses by themselves.

(iii) The $SU(2)$ statement, including the choice of one common
triple for both the bond and node estimates, is proved in
Section~\YMMref{53.4}{445}. Its general-group counterpart with the
same domains is the explicit localisation input of (V). The
logarithms used to align the collar data are formed from products
of small exponentials under Lemma~\ref{C-368-6}(ii).

(iv) At a bond $b$ of the cut-off zone $E_b$ is a product of seven factors, and their chords add: the localisation
of (iii) bounds the relative field of the two backgrounds beneath $b$ by $L^{-k}\varsigma$ and gives two factors at
the ends within $\varsigma$ of $\mathbf 1$; Lemma~\ref{C-368-7}(v), as $L^k$ times that field is below
$\varsigma<c_4$, gives the averaged exponent, at most $2\varsigma$, the two factors $f_k$ at the ends, within
$C_v\varsigma$ of $\mathbf 1$, and the two exponents of the contour variables at the ends, at most
$2\kappa_6\varsigma$ each; so $|E_b-\mathbf 1|\le(4+2C_v+4\kappa_6)\varsigma=\mathsf c_L\varsigma$. The fifth
member of $\epsilon_{\rm xp}$ gives $\varsigma<\frac{c_4}{2L}$, so $\mathsf c_Lc_4\le2Lr_G$ gives $|E_b-\mathbf
1|<r_G$, $E_b\in G^\pi$ and $|H_b|\le\frac\pi2|E_b-\mathbf 1|$; the bound
$8\|\vartheta_G'\|_\infty|H_b|\le\delta_k$ comes from $A_1/A_0\le c_{\rm Q}$ and $\gamma\le\gamma^G_{\rm xp}$, and
$\mathsf c_{\rm fl}\delta_k\le1$ from $C_\chi\mathsf c_{\rm fl}\delta_k\le r_\chi\le1$; the assembly is that of the
$SU(2)$ construction with $r_G$ for $2$.

A further argument gives $E_b\in G^\pi$ from the block constraints. On the double block of the coarse bond $c$ of
$b$ let $V$ and $V'$ be the $k$-fold averages of the global and of the localised background. They have the same
contour variables $X_\pm$ at the two end blocks, both being in the axial gauge at the sites of the region of level
$k+1$ and both being given by the datum at the others, as an element of a minimal orbit satisfies its block
constraints and the joined datum is the datum at the common bonds; and they have the same average $X_0$ at $c$. So
$V(b)=X_-^{-1}e^{i\beta}X_0X_+$ and $V'(b)=X_-^{-1}e^{i\beta'}X_0X_+$ with $\beta,\beta'\in\mathfrak g$,
$|\beta|+|\beta'|<\pi$, and $E_b=V(b)V'(b)^{-1}=X_-^{-1}e^{i\beta}e^{-i\beta'}X_-$, which lies in $G^\pi$ by
Lemma~\ref{C-368-6}(ii), $G^\pi$ being kept by conjugation with $G$. The statement that $G\setminus G^\pi$ is a
null set (the $SU(2)$ case of Lemma~\ref{C-368-10}(i) and Proposition~\YMMref{53.4}{445}) is false for $SU(N)$, $N\ge3$, at every
$H_b\ne0$, and it is not needed: the substitution enters only integrands that vanish outside the chart
(the chart-supported substitution in the proof of Proposition~\YMMref{53.4}{445}).

(v) The proofs of the $SU(2)$ construction stand; ``trace $0$'' becomes $\mathrm{Ad}(g)\mathfrak g^c=\mathfrak
g^c$, and $F=\xi^{-2}PH$ with $P\,\mathrm{Ad}(g)=\mathrm{Ad}(g)P$; a holomorphic map that is $\mathfrak g$-valued
at real points is $\mathfrak g^c$-valued; and $G^c=G\exp(\mathfrak g)$ is closed.
\end{proof}

\subsection{Dependence on regular backgrounds}\label{C-368-32}

The construction hypotheses have different roles. The propagator estimates control inverses at general regular
fields. The variational estimates then construct and localise the constrained critical orbit. The second-part
hypotheses provide the holomorphic Gaussian families and the one-step estimates on their full domains. Finally, the
constant conditions must be compatible in the prescribed order. The following reductions identify the group
dependence in these requirements.

\emph{Propagators.}
Proposition~\ref{C-368-16}(i) proves the local Hessian and current
identities and separates them from the quantitative inputs in (P).
For other simple $G$ the Hessian bounds include the factor $|P|$.
Changing the colour norm, taking a transpose, and summing components
may introduce the factors $\kappa_G$, $N^{1/2}$ and $N$ in
Definition~\ref{C-368-3}(2). The constants of an entire inverse or
tree expansion must be chosen from its admitted quantitative bounds;
these norm comparisons alone do not certify a particular numerical
choice. H\"older and $\ell^2$ seminorms of $\mathfrak g^c$-valued
functions compare with those of real functions by the factors $N^{1/2}$ and $(2n)^{1/2}$, since
$\|\cdot\|\le|\cdot|\le N^{1/2}\|\cdot\|$ and such a function has $2n$ real components; the estimate required in
(P) keeps the rate $\delta_2$, and the floor on $M$ moves with $N$ through the natural bound $B_{\rm nat}$, of the
order $N^{1/2}$. The covariance numbers, thresholds, rates and cell sizes are formed from those of the expansion.

\emph{Variational estimates.}
The variational estimates in (V) require the background estimates of (P) with the changes of
Definition~\ref{C-368-3}, among which these concern the variational part: a proof that uses the current bound of
the regular class runs with the class $\mathcal U(r)$, into which $\mathcal U'(r)$ falls at the radius $\kappa_Gr$,
and for $r\le z_G$ at $(1+\kappa'r^2)r$, thresholds on $\alpha_0$ read as ceilings being divided by $\kappa_G$, or
by $1+\kappa'$ (Remark~\ref{C-368-30}, Lemma~\ref{C-368-8}(ii)); the bound $3\alpha_0$ of Proposition~\YMMref{30.1}{236}
becomes $3|P|\alpha_0$, at most $6\alpha_0$ at $SU(N)$, and $j_G\alpha_0$ under $\alpha_0\le z_G$, which is
$3\alpha_0$ at $SU(N)$; the second inclusion becomes $\mathcal U(r)\subset\mathcal U'(r+r^2+\kappa'r^3)$ for every
$G$ and $r\le z_G$, and the membership in $\mathcal U'(2s)$, $s=2\alpha_0+3\alpha_2$, of the curvature/current conversion in Lemma~\ref{C-368-8}(ii), with both of its quantitative hypotheses,
holds for every $G$ under $\alpha_0+3\alpha_2<z_G$ and $s(1+\kappa's)\le1$, so for $SU(N)$ under
$\alpha_0+3\alpha_2<z_G$ and $s\le\frac13$, both by the bound at a bond of Lemma~\ref{C-368-8}(ii), for the
membership at each full bond $b$ with $\rho=(\alpha_0+3\alpha_2)\eta^2\mathsf L(b)^{-2}$
(Definition~\ref{C-368-3}(3)); Lemma~\ref{lem:group-contour-gauge}, whose averaged contour variables are means of contour
families (Proposition~\YMMref{4.3}{18}), takes the guard $\varpi_y(W)<z_G$, under which the contour quotients,
products of at most $\frac98L^2$ conjugated plaquette variables within $\varpi_y(W)$ of $\mathbf 1$, lie in $D^c$
by the rule for products, so that the means lie in $G$ (Lemma~\ref{C-368-6}(iii), (iv)), and the comb gauge of
Lemma~\YMMref{14.4}{64} takes the guard $\varpi<z_G$ (Lemma~\ref{C-368-8}(iii)); and wherever the towers use
$|X^c-\mathbf 1|=|X-\mathbf 1|$, an equality of moduli at the plaquettes alone, they take the form with $m(X)$,
with $|\mathrm{Ad}(g)-1|\le m(g)(2+m(g))$, and a logarithm is guarded by $X\in D^c$.


The existence assertion required in (V) has the following uniform content: there must be a power $\mathfrak r(G)$ of $L$, a number
$B_3(G)\ge1$ and, for every power $\rho\ge\mathfrak r(G)$ of $L$, a number $a_1(\rho,G)>0$, numbers of $L$, $\rho$
and the pair with $\kappa_GB_3(G)a_1(\rho,G)\le a_0(G)$, $a_0(G)>0$ the fixed uniqueness radius required in (V), such that in every setting of the $SU(2)$ construction at the separation $\rho$,
for every reversal-invariant set $\mathfrak F$ of inward bonds and over every real datum $V$ on $\mathbf
B_{\mathfrak F}$ whose plaquette variables, and on a sequence of regions whose holonomies at the tested plaquettes,
lie within $\varepsilon_1\le a_1(\rho,G)$ of $\mathbf 1$, the set $C(\mathbf B_{\mathfrak F},V)$ holds a critical
configuration in $\mathcal U'(B_3(G)\varepsilon_1)$. Together with uniqueness of the regular critical orbit this
assumption gives one nonempty regular critical orbit, contained in
$\mathcal U'(B_3(G)\varepsilon_1)$; near the trivial datum it is
the conditional assertion of Proposition~\ref{C-368-16}(ii). Its route is the corresponding $SU(2)$ proof, the existence of a critical
configuration in the class and the improvement of the regularity radius (the $SU(2)$ model is Proposition~\YMMref{37.1}{278}) being run for $G$
with the bounds that the variational part takes from the earlier sections at a regular reference configuration
$U_0\ne\mathbf 1$: those of the linear operators at $U_0$, which rest on the propagators of (P) at regular
backgrounds, and those of the straightened averages, the action and the gauge transformations in the chart, the
rows and the gauge functional at $U_0$.


\emph{Conjugation.}
Proposition~\ref{C-368-22}(i) gives (C) for the maps listed there and for their composites, finite sums and
continuations on convex subsets of a vector space kept by the conjugation (slices without group-valued member,
Lemma~\ref{C-368-15}(i), (ii)). The maps of the induction are of these kinds: each is a chain of factors, each
factor being built from products, inverses, means, logarithms in $D^c$, the action $a$ and its derivatives, which
commute with the conjugation by Lemma~\ref{C-368-24}(i), (ii), or one of the three implicit systems of
Lemma~\ref{C-368-24}(iii), or a tower, or a continuation on a convex subset of a vector space kept by the
conjugation, or a map holomorphic on the open slices of the real-connected part of its domain that takes real
fields to real values, every trace of a group-valued argument at complex fields being $a(X)$
(Definition~\ref{C-368-3}(1)); and the domains of the $SU(2)$ construction pass into real-connected parts by
Lemma~\ref{C-368-15}(iii). A map that is holomorphic on the open slices of the real-connected part of its domain,
its real members being constant on slice components, is of them as soon as it takes real fields to real values, and
it then maps the real-connected part of its domain into that of its target as soon as it maps it into the target
(Lemma~\ref{C-368-15}(ii)).


\emph{The pivot bond and Gaussian coefficient.}
For the pivot bond, the unknown element $x\in G$ at the pivot bond enters the quotient $R=W[c]G_1[c]^{-1}$ through
a homeomorphism of $G$ onto itself; in each of the three cases of the averaging map a solution has the eigen-angles
of $R$ below $\frac1{272L^2}$ in modulus, so that it lies among the fields with plaquette variables within
$\omega_M$ of $\mathbf 1$ on the double block, where the average is the formula of Lemma~\ref{C-368-7}(i) and the
other two cases have no solution; every element that occurs is a product of conjugated plaquette variables within
$\omega_M\le z_G\le r_G$ of $\mathbf 1$, so that its logarithm lies in $i\mathfrak g$ (Lemma~\ref{C-368-6}(ii),
(iii)); and in the chart the equation reads $r=L^2(t-\bar s-E(r))$ for $r=\log R$, with $t$ and $\bar s$ averages
over the block of logarithms that do not read $x$ and with $E$ holomorphic, the right member being a contraction of
$\{r\in i\mathfrak g:|r|\le\frac1{50}\}$ into itself with constant below $\frac1{121L^2}$; its fixed point gives
the one solution and the bound at the pivot bond, with the numerals of the $SU(2)$ construction. The large-field
chain is Lemma~\ref{C-368-9}(ii). The localised Gaussian family is formed from the fluctuation operators of (P) by
the formulas of the $SU(2)$ construction, and at the trivial triple its Hessian is that of the normaliser
(Proposition~\ref{C-368-13}(iii)).


\emph{Choice of constants.}
Proposition~\ref{C-368-18} gives the shape of the table and the criterion: one test for each condition and one
joint test for each two-sided constant. A floor on a constant free upward and a positive ceiling on a constant free
downward are admissible. A condition whose last constant is determined must be checked from its formula. The
conclusion that all the conditions pass these tests is hypothesis (K).


The conditions of Lemma~\ref{C-368-15}(iii) at the memberships reached from memberships of restrictions had apart
are positive ceilings on the second radius where the first radius lies below $B_0$, and for the cores $Kb_0<B_0$
and $b_0<a^\sharp$; the ceiling $\mathsf c_Lc_4\le2Lr_G$ is a positive ceiling on $c_4$, and the row
$\alpha_1<\rho_GL$ one on $\alpha_1$ (Propositions~\ref{C-368-16}(iv), \ref{C-368-17}).

The comparison of the two norms, the regularity window and the central-element guards are specified in
Proposition~\ref{C-368-18}(c), (e), (f). The numbers of the pair must occur before $\gamma$ in the order of choice;
they may depend only on the constants preceding them and on the pair.


\subsection{Multiscale induction and the choice of constants}

\begin{proposition}[Conjugation and the inputs of the induction]\label{C-368-22}
The conjugation assertions apply to the maps wherever they are formed. The global second-part assertions and the pointed-family assertion in (iii) assume (II) of Definition~\ref{C-368-31}.
Let $\mathfrak g$ be simple, $L\ge13$ odd, $200L^2z_G\ge1$; constants are constants depending only on $L$ and the pair, uniform in the volume.
\begin{enumerate}
\item[(i)] Let $a(X)$ stand for every trace of a group-valued argument at complex fields, the gauge-fixing term
included, and let the domains be written with $m(X)$ and with $D^c$ at the logarithms of the averages. Then
plaquettes, path products, gauge transforms, covariant differences, $F$, $J$, means, averages and their iterates,
the chain to the fluctuation operator and the three implicit systems of Lemma~\ref{C-368-24}(iii), the towers, and
the composites and finite sums of these maps and their continuations on convex subsets of a vector space kept by
the conjugation (slices without group-valued member) commute with $p\mapsto p^c$ on the real-connected parts of
their domains and map them into the real-connected parts of their targets. A map holomorphic on the open slices of
$Q_1^r$, with real members constant on slice components, commutes with $p\mapsto p^c$ there if and only if it takes
the real points of $Q_1$ to real points, and then it maps $Q_1^r$ into $Q_2^r$ as soon as it maps $Q_1^r$ into
$Q_2$ (Lemma~\ref{C-368-15}(ii)).
\item[(ii)] In the second-part statements assumed in (II), the clauses that are false for $G$ as worded are
required in the following corrected forms, among them these: $m(X)$ where $|X^{-1}-\mathbf 1|=|X-\mathbf
1|=|X^*-\mathbf 1|$ was used, which is false at $SU(3)$; the radius $r_n$ where $|1-T|\le\frac12$ was to keep $\det
T$ off $(-\infty,0]$ (the three-dimensional determinant argument in the proof of Proposition~\YMMref{50.2}{386}), which is false from $n=6$; the factor $\max\{1,n2^{-11}\}$ in
the bound \eqref{eq:colour-determinant-bound} (Proposition~\ref{C-368-18}(d)); the bound $z<(\frac54)^nL^{3n}$ in the place of
$z<2L^9$ in Lemma~\YMMref{65.3}{581}; orthogonal colour maps that keep $|\cdot|$ in the covariance argument following Proposition~\YMMref{47.4}{365}; the indicator of $G^\pi$ where $G\setminus G^\pi$ is used as a null set, and
integrands that vanish outside the chart in the substitution of Proposition~\YMMref{53.4}{445}; $a(X)$ for the
bare trace of the gauge-fixing term of the definition of $\Psi^{\rm G}$ in Section~\YMMref{51}{389}, p.~393; $E^{\rm fam}\ge e_n$ in the bound for the local Gaussian constants, as derived in Proposition~\ref{C-368-18}(d); and the two-sided guard $m(X)$ in the elimination-factor estimate of Lemma~\ref{C-368-10}(iv). Of the clauses in which a proof of the second part works with one colour, those at the
trivial field and in the averaging layer hold for $G$ with the constants of the $SU(2)$ construction; those that
concern the background near trivial block data, and the radius $r_0=a^{\rm us}(\alpha_0,\alpha_1)$ in $\mathsf
G_\sharp$ (Proposition~\ref{C-368-17}), hold as soon as on every truncated tower the background exists, is one
orbit, is holomorphic on a polydisc of uniform radius and has there the bounds of Proposition~\ref{C-368-16}(ii),
whose uniform conclusions here are assumed in (V), including uniqueness of the critical orbit;
and one more ceiling, $\varepsilon_{\rm v}\le a_G$, enters on the regularity radius $\varepsilon_{\rm v}$ of the
critical configurations (Definition~\YMMref{9.5}{41}), with a fixed admitted positive radius $a_G$ from
Definition~\ref{C-368-12} (Proposition~\ref{C-368-18}(f)).
\item[(iii)] Near $W=\mathbf 1$ the localised Gaussian family of the second part (Proposition~\YMMref{50.2}{386}) has at
the trivial triple the Hessian of the Gaussian normaliser of Proposition~\ref{C-368-13}(iii); its quadratic
coefficient there is one number for all data, $b^0_l(G)$. Under (II), for simply connected $G$ the family is a
pointed family of its level on the analyticity cores of $G$, holomorphic on balls of one radius, with the
contraction bound of its expansion (Definition~\ref{C-368-31}(II)).
\end{enumerate}
\end{proposition}

\begin{proof}[Sketch of proof]
(i) A map is covariant on the real-connected part if it is a chain of factors each of which is either built from
products, inverses, means, logarithms in $D^c$, the action and its derivatives, and so covariant by
Lemma~\ref{C-368-24}, or holomorphic on open slices and real at real fields, and so covariant by
Lemma~\ref{C-368-15}(ii), or a continuation on a convex subset of a vector space kept by the conjugation (a slice
without group-valued member); a second way is to check that each of its formulas is real and holds at the doubled
pair $X\mapsto\mathrm{diag}(X,(X^{\mathrm T})^{-1})$, for which $|\mathrm{diag}(X,(X^{\mathrm T})^{-1})-\mathbf
1|=m(X)$ and the action, the norm and the form restrict to those of $G$. The maps listed in (i) are of these kinds:
plaquettes, path products, gauge transforms, covariant differences, $F$, $J$, means, averages and their iterates
and the chain to the fluctuation operator are built from products, inverses, means, logarithms in $D^c$, the action
and its derivatives (Lemma~\ref{C-368-24}(i)--(iii)); the three implicit systems are covariant by
Lemma~\ref{C-368-24}(iii); the towers and the continuations on convex subsets of a vector space kept by the
conjugation by Lemma~\ref{C-368-15}(ii); and composites and finite sums of covariant maps are covariant.
Lemma~\ref{C-368-15}(iii) moves the domains of the $SU(2)$ construction into real-connected parts at the cost named
there. Two statements of the $SU(2)$ construction take another form for $G$: the gauge-fixing term of
the definition of $\Psi^{\rm G}$ in Section~\YMMref{51}{389}, p.~393, set with a bare trace at complex fields, whose imaginary part $\frac
tN\mathrm{tr}\sin(tH)$ at $SU(3)$ does not vanish at a real field, takes $a(X)$, which keeps its quadratic part and
its bounds; and the statement that the conjugate point lies in the domain if and only if the point does, which for
$SU(2)$ rests on the isometry $\Theta$ of Lemma~\ref{lem:two-by-two-conjugation} and is false beyond $SU(2)$, takes the two-sided
guard $m(X)$.

(ii) The forms that hold are those of Lemmas~\ref{C-368-24}(ii), \ref{C-368-10}(iv) and
Proposition~\ref{C-368-18}(d) for the first three, and those of Definition~\ref{C-368-3}(2), (3),
Remark~\ref{C-368-30}, (i) above, Propositions~\ref{C-368-16}(iv), \ref{C-368-18}(d) and Lemma~\ref{C-368-10}(iv)
for the others. At the trivial pair every operator is scalar times $1_{\mathfrak g^c}$ (Lemma~\ref{C-368-11}(ii)),
with the lower bound of the fluctuation operator and only $3$ replaced by $n$, and the averaging layer is
Lemma~\ref{C-368-7}; near trivial block data the clauses need the existence, one orbit, holomorphy on a polydisc of
uniform radius and the bounds of the background on every truncated tower, of which Proposition~\ref{C-368-16}(ii)
gives these conclusions under (V). The ceiling uses a fixed admitted orbit radius from Definition~\ref{C-368-12};
the central-field obstruction described there applies only when the chosen averaging extension sends that field to the trivial block datum.
It does not establish existence or a sharp value of a positive orbit radius for every averaging extension.

(iii) At the trivial triple the germ of second order of the localised family is computed as that of the normaliser
in Proposition~\ref{C-368-13}(i): by invariance the background is taken, to second order, along one unit
$H\in\mathfrak g$; the fluctuations split into the eigenspaces of $\mathrm{ad}\,H$, each of which is a charged
sector of $SU(2)$; and the action adds the one tadpole, with $\mathrm{tr}\,M_H=C_\rho$. So the coefficient of the
family there is $\frac{C_A}8$ times that of the localised family of $SU(2)$, which is one number for all data and
equals the coefficient of $\log Z^{(k)}$ (equation~\YMMref{(790)}{386} identifies their totals up to a constant, and Theorem~\YMMref{X}{427} identifies the coefficient of such totals), the increment $b^0_{k+1}$, plus
$\gamma_G(j_{k+1}-L^{-1}j_k)=(\frac{C_\rho}3-\frac{C_A}8)(L^3-1)L^{-4(k+1)}$, $j_k=\frac23(L^{-k}-L^{-4k})$; this
is $b^0_{k+1}(G)$, the coefficient of the normaliser (Theorem~\ref{C-368-4}(c), Proposition~\ref{C-368-13}(iii)),
and it reads no datum.
\end{proof}

\begin{proposition}[Induction under the construction hypotheses]\label{C-368-17}
For a simple Lie algebra, assume (P)--(K) of Definition~\ref{C-368-31}. For a semisimple Lie algebra, assume the
transferred one-step statements and their applicability at the data of each citation below. For tied ideals, the
construction conclusion has property (T) of Theorem~\ref{C-368-4}(h3) as its hypothesis. The coefficient estimates
below apply whenever the indicated terms and their analyticity bounds are given.
Let $\mathfrak g$ be semisimple. Each of the following proofs of the $SU(2)$ construction derives its transferred
statement from the transferred statements that it cites, the hypotheses of each citation, those that the forms (3)
of Definition~\ref{C-368-3} add to a cited statement among them, being met at its data, as in
Definition~\ref{C-368-31}(II): the regular progression in Theorem~\YMMref{Y}{438}; the actual one-step estimates of Theorem~\YMMref{W}{423}, their quadratic coefficient in Theorem~\YMMref{X}{427}, and the canonical coupling construction in Section~\YMMref{52.10}{439}, with $\mathsf G_\sharp$ as its coefficient bound; the historical step in Proposition~\YMMref{69.3}{603}, its iteration in Theorem~\YMMref{AC}{604} and Corollary~\YMMref{69.4}{604}; and the history bounds of Section~\YMMref{70}{604}, which give Theorem~\YMMref{D}{8}. In the historical step the proof explicitly checks the coordinate changes and tests, the separate on-shell identities for the old exponent and retained factors, the conditional Gaussian and cluster identity, the normalizer, the kept and rejected forms, the strip sum, and the canonical identification of the new exponents. The positive change-of-variables theorem used there is Theorem~\YMMref{AB}{566}; the boundary and retained-factor bounds are those of Section~\YMMref{62}{567}. Their transferred hypotheses, including their domains and measurability assertions, remain required here. Of the group these proofs use only that $\mathrm{Ad}(G)$ keeps
$\langle\cdot,\cdot\rangle$ and $|\cdot|$, that invariant linear forms vanish, that $|t_a|\le\kappa_G$, and that
$a(wuw^{-1})=a(u)$. The part of second order of the effective terms enters them as one real number, the mean
coefficient $\mathsf{cf}$, the mean of its $n$ quadratic colour coefficients: it is defined, bounded, additive, real and holomorphic in the parameters,
never a multiple of the Hessian of the action, with $|\mathsf{cf}|\le\mathsf G_\sharp=\kappa_G^2C_{\rm
am}r_0^{-2}E_0\le B_\sharp$, where $C_{\rm am}$ is the constant of the coefficient bound~\YMMref{(940)}{428} and $r_0=a^{\rm
us}(\alpha_0,\alpha_1)$ the radius of the $SU(2)$ construction; $C_{\rm am}$ and $r_0$ read constants of the
$SU(2)$ construction that come before them in its order, $\mathsf G_\sharp$ reads besides the pair, through
$\kappa_G$, and $E_0$, and none of them reads $\gamma$, $B_\sharp$, a level, the cutoff, $\beta$ or a field. On an
ideal $\mathfrak g_i$, $|\mathsf{cf}_i|\le\kappa_i^2C_{\rm am}r_0^{-2}E_0=(\kappa_i/\kappa_G)^2\mathsf G_\sharp$.
The numbers $\frac92$ and $3$ of Theorem~\YMMref{D}{8} become $\frac{3n}2$ and $n$; the $\pi$ of the ceiling
$\alpha_1<\pi L^j$, $1\le j\le k$, under which the local terms are continuous on the analyticity cores, which a
proof of the chain from level $k$ to level $k+1$ quotes, becomes $\rho_G$ of Definition~\ref{C-368-2}, the bound on
the eigen-angles under which an element of $G$ lies in $G^\pi$, the domain of the chart, so that the ceiling is
$\alpha_1<\rho_GL$, one more row of the table of Proposition~\ref{C-368-18}(a), a positive ceiling on $\alpha_1$,
which holds of itself for $\rho_G>\frac1{13}$, as $\alpha_1\le1$ and $L\ge13$; as $\rho_G\le\pi$, this row gives
$\alpha_1<\pi L^j$, under which the exponential of $\mathfrak g^c$ is a local diffeomorphism at every $Z_0$ with
$|Z_0|<\pi$, the eigenvalues of $\mathrm{ad}\,Z_0$ being at most $2|Z_0|<2\pi$ in modulus for matrices of every
size; the $\frac32$ in the size of an array stays if the analyticity condition is read with $m(\partial U)$, since
$|a(u)|\le\frac12|u-\mathbf 1||u^{-1}-\mathbf 1|$; and the transferred statements of these proofs differ from those
of the $SU(2)$ construction only in these numbers and in ``of $L$ and the pair'' for ``of $L$'', the sign $\mathsf
G_\sharp$ keeping its places with its value for $G$. The large-field chain of the $SU(2)$ construction,
Proposition~\YMMref{44.4}{347}, Proposition~\YMMref{44.5}{348} and Lemma~\YMMref{65.2}{580} and the local-deposit argument in the proof of Proposition~\YMMref{65.4}{582}, is
transferred by Lemma~\ref{C-368-9}(ii), its statements taking the prices of the $SU(2)$ construction times
$\tau_G^{-2}$ and the averaging radius $\min\{\omega_M,r_G\}$. The same holds for an action
$a(u)=1-\frac12\mathrm{tr}\,\Lambda(u+u^{-1})$ with a weight $\Lambda>0$ in the commutant of $G$,
$\mathrm{tr}\,\Lambda=1$, the form $\mathrm{tr}(\Lambda XY)$ and $\kappa_\Lambda=\max|Y|/\mathrm{tr}(\Lambda
Y^2)^{1/2}$ in the places of $\langle\cdot,\cdot\rangle$ and $\kappa_G$.
\end{proposition}

\begin{proof}[Sketch of proof]
The construction hypotheses provide the cited one-step estimates and their applicability on the transferred
domains. There are five additional guards to retain. Logarithms of group variables are taken in $D^c$; a logarithm
in the real group is taken in $G^\pi$; complex domains use $m(X)$; determinant logarithms use the radius $r_n$; and
local terms on the analyticity cores satisfy $\alpha_1<\rho_GL^j$. Lemmas~\ref{C-368-6}, \ref{C-368-10},
\ref{C-368-24} and Proposition~\ref{C-368-16}(iv) supply the first four at their stated data. The row
$\alpha_1<\rho_GL$ supplies the last because $L^j\ge L$. Compatibility of these rows is part of (K).

The colour algebra in the induction consists of invariance of the trace form and of the operator norm, vanishing of
invariant linear forms, $|t_a|\le\kappa_G$, and conjugation invariance of $a$. The comparison of the two volume
exponents in Theorem~\YMMref{D}{8} becomes $n\le3n/2$. With the one-step bounds and domain conditions assumed,
the induction proceeds on the same histories.

Terms of first order in the fluctuation field are bounded by Schwarz's lemma: if $h$ is holomorphic on
$|\lambda|<\rho$ with $\rho>1$, $h(0)=0$ and $|h|\le c$, then $|h(1)|\le c/\rho$, nothing being asked of $h'(0)$. The damping
$L^{-(k-j)}$ of old levels at $U=\mathbf 1$ comes from the vanishing of invariant linear forms. In the nine proofs
of the chain from level $k$ to level $k+1$ the inputs ask $\|E^{(j)}\|\le E_0$ and $|b_j|\le B_\sharp$ separately,
and the bounds of a step add these sizes. The increments enter only as the real numbers $e_j=b_j$ with $|b_j|\le
B_\sharp$, no relation between them and the entries of a term or their derivatives is used, and the one
stationarity used is that of the one action. The bounds of the coefficient: by Cauchy's estimate on the ball of
radius $r_0$ the part $q$ of second order of a term $E$ obeys $|q(X,X)|\le C_{\rm am}r_0^{-2}\|E\||X|^2$, as in the
$SU(2)$ construction, where the orthonormal basis has $|t_a|=1$; with $|t_a|\le\kappa_G$, and $|t_a|\le\kappa_i$ on
$\mathfrak g_i$, this gives $|\mathsf{cf}_i|\le\kappa_i^2C_{\rm am}r_0^{-2}E_0$ and
$|\mathsf{cf}|\le\sum_i\frac{n_i}n\kappa_i^2C_{\rm am}r_0^{-2}E_0\le\mathsf G_\sharp$. The factor $\kappa_i^2$
against $C_{\rm am}$ cannot be dropped: for $SU(2)^r$ on blocks, $r\ge2$, one has $\kappa_i^2=r$, and a term that
reads one factor has on its ideal $r$ times its coefficient for $SU(2)$.
\end{proof}

The colour factor in the determinant estimate is most easily checked before
it is absorbed into the constant. The following calculation records the
precise contraction and closed-chain bounds it requires. Their availability
for the transferred operators remains part of the construction hypotheses.

\begin{lemma}[The colour factor in a local determinant bound]
\label{lem:colour-determinant-count}
Let $n\ge1$, $c>0$, $0<q<1$, $Q>1$ and $\delta_1>0$. Put
\[
\begin{aligned}
 \delta&=\frac{\delta_1}{128}
        \frac{\log(1/q)}{\log Q+\log(1/q)},\\
 C_{\rm a}&=n\left(|\log c|+\log\frac1{1-q}\right),\\
 C_{\rm d}&=\frac Q{Q-1}+\frac1{1-q},\qquad
 C_{\rm mob}=2^7C_{\rm a}+C_{\rm d}.
\end{aligned}
\]
Assume $\delta M\ge9\log2$. Let $X$ be a finite set of coarse
cubes, $d=d_t(X)\ge0$, and let $\emptyset\ne S\subset X$ satisfy
$d_t(S)=0$ and $\#X\le8(d+1)$. For each $S\subset Y\subset X$
let $T_Y$ be an operator on bond fields with $n$ colour coordinates,
with $\|T_Y\|\le q$, and with a common bond $b$ in its domain. Define
\[
 f_b(Y)=n\log c-\sum_{\nu\ge1}\frac1\nu
                  \operatorname{tr}T_Y^\nu(b,b),\qquad
 E(X,b)=\sum_{S\subset Y\subset X}(-1)^{\#X-\#Y}f_b(Y).
\]
The trace is over the colour coordinates. When $d>0$, suppose the
closed-chain sums
\[
 \Sigma_\nu=\sum_{S\subset Y\subset X}(-1)^{\#X-\#Y}
                         \operatorname{tr}T_Y^\nu(b,b)
\]
satisfy, with $\mathcal N=2^{\#X-\#S}$,
\begin{equation}\label{eq:colour-determinant-majorants}
 |\Sigma_\nu|\le n\mathcal N q^\nu,
 \qquad
 |\Sigma_\nu|\le n\mathcal N Q^\nu
                      e^{-\delta_1M(d+1)/64}.
\end{equation}
Then
\begin{equation}\label{eq:colour-determinant-bound}
 |E(X,b)|\le2^{-d}\max\{1,n2^{-11}\}C_{\rm mob}e^{-\delta Md}.
\end{equation}
\end{lemma}

\begin{proof}
The contraction bound makes each series absolutely convergent and gives
$|f_b(Y)|\le C_{\rm a}$. There are at most $2^{8(d+1)-1}$ sets
$Y$. If $d=0$, their sum is bounded by $2^7C_{\rm a}\le C_{\rm mob}$.
For $d>0$, $X\ne S$, so the alternating sum of the constant $n\log c$
vanishes. Thus $E(X,b)=-\sum_{\nu\ge1}\Sigma_\nu/\nu$.

Put
\[
 \theta=\frac{\delta_1}{64(\log Q+\log(1/q))},\qquad
 \nu_*=\lfloor\theta M(d+1)\rfloor.
\]
Use the second majorant up to $\nu_*$ and the first afterwards.
Since $\theta\log(1/q)=2\delta$ and
$\theta\log Q-\delta_1/64=-2\delta$, the two geometric sums give
\[
 |E(X,b)|\le n\mathcal N
       \left(\frac Q{Q-1}+\frac1{1-q}\right)e^{-2\delta M(d+1)}.
\]
One exponential absorbs the region count:
\[
 |E(X,b)|\le\frac n4\,2^{-d}C_{\rm d}e^{-\delta M(d+1)}
       \le n2^{-11}\,2^{-d}C_{\rm d}e^{-\delta Md}.
\]
Here the second inequality retains the further factor
$e^{-\delta M}\le2^{-9}$. Together with the case $d=0$ and
$C_{\rm d}\le C_{\rm mob}$, this proves the claim.
\end{proof}


\begin{proposition}[The constants]\label{C-368-18}
Let $L\ge13$ be odd and $200L^2z_G\ge1$. The constants are chosen in the order of the $SU(2)$ construction,
$\gamma$ last, each chosen constant ranging, given the earlier ones, over the range that the $SU(2)$ construction
declares for it; each is free upward, free downward, two-sided (its admissible values, given the earlier ones, form
a window) or determined by a formula from earlier ones. The table of the constants of the $SU(2)$ construction
carries besides a radius $\rho>0$, which is not chosen and stands in the rows on $\gamma$ alone. A set of choices
of all the constants is \emph{large} if it is large at the empty choice, where, from the last constant down to the
first, a set is large at a choice of the constants before the $i$-th one if the values of the $i$-th constant at
which it is large at the longer choice hold, for some $t$ and some $c>0$, the part $[t,\infty)$ of its range when
this constant is free upward, the interval $(0,c]$ when it is free downward, all its admissible values when it is
two-sided, and its value when it is determined; at a choice of all the constants a set is large if it holds that
choice. A set that reads $\gamma$ or $\rho$ is large if so is the set of the choices of the constants before
$\gamma$ at which, for every $\rho>0$, some interval $(0,c]$ consists of admissible values of $\gamma$ at which it
holds.
\begin{enumerate}
\item[(a)] For $SU(N)$, every $N\ge2$, the table of the constants of the $SU(2)$ construction, formed with the
numbers of $G$ as data (any positive values of these numbers in the ranges that the $SU(2)$ construction states for
them serve for its shape), with the changes of Definition~\ref{C-368-3} and with those conditions of (c)--(g) that
are a floor on a constant free upward or a positive ceiling on a constant free downward laid as rows, has
admissible shape, the shape under which, at every admissible choice of the earlier constants, the admissible values
of a constant free upward hold a half-line of its range, those of a constant free downward an interval $(0,c]$, and
the window of a two-sided constant is not empty; $\gamma$ is last, and choices exist at every value of the radius
$\rho$, the admissible choices of the constants before $\gamma$ being the same for all $\rho$; a further condition
read for $G$ is laid on it as a row when it is large. For a semisimple pair with a weight in the commutant, the
table of the second part of the $SU(2)$ construction, its chosen constants, their order, kinds and ranges kept, is,
as the table of the $SU(2)$ construction is, of admissible shape and without cap, with $\gamma$ last, and choices
exist, given that the numbers of the first part of the pair exist with the ranges that the $SU(2)$ construction
states for them. In both, a floor on a constant free upward and a positive ceiling on a constant free downward
never obstruct; $L$ is free upward, with floors alone, and a floor on $L$ may be laid once at the root, every other
constant keeping the same admissible set. This is the shape of the table; that every condition read for $G$ is
large, so that all of them are rows of a table of admissible shape and hold at each of its choices, is (K).
\item[(b)] Large sets form a filter: the intersection of two sets of choices is large if and only if both are.
Finitely many conditions are each large if and only if there is a table obtained from that of (a) by floors on free upward constants,
positive ceilings on free downward constants, and further constants determined by formulas from earlier ones,
at every choice of which all the conditions hold. These additions retain every original two-sided window whole;
one floor or ceiling on each free constant suffices. Separately, further ends on a two-sided constant are jointly feasible
at an earlier choice exactly when its original window meets all those ends together. Requiring a large set of such earlier choices
is a compatibility test for the restricted windows, not largeness of the ends as subsets of the original full table.
Hypothesis (K) requires both the largeness conditions and the joint-window compatibility tests used by the construction.
\item[(c)] For $SU(N)$ the two-sided conditions that compare $|\cdot|$ with $\|\cdot\|$ hold together for every
$N\ge2$ at the ratio $\frac1{N-1}=\kappa_G^{-2}$ of the squares of the two norms, with the kernel bounds of (P)
read in the normalised trace norm, which the bounds in $|\cdot|$ give with the factors $(N-1)^{1/2}$ and $N^{1/2}$
between the two norms (Proposition~\ref{C-368-16}(i)). The two-sided condition on the regularity radius
$\varepsilon_{\rm v}$, which comes from the inclusions $\mathcal U'(r)\subset\mathcal U((1+\kappa'r^2)r)$ and
$\mathcal U(r)\subset\mathcal U'(r+r^2+\kappa'r^3)$, $r\le z_G$, of Lemma~\ref{C-368-8}(ii), is not empty while
$\frac1{400L^2}<2\sin\frac\pi N$, a constant that precedes $\varepsilon_{\rm v}$ in the order being multiplied by
$1+\kappa'$, with the numbers of the critical configuration of (V), and one floor on the separation constant $R_1$ makes its ends compatible; for the two-sided conditions that a central element near $\mathbf
1$ enters when $z_G<1$ see (e).
\item[(d)] The chains that use $\mathrm{Log}\det$ (the proof of Proposition~\YMMref{50.2}{386} and the elimination-factor estimates of Lemma~\ref{C-368-10}(iii), (iv))
hold at the radius $r_n$, with $4Ka\le r_n$ and $C_z\ge4c_nK$, $c_n=\rho_n/r_n$, the pair $4Ka\le\frac12$,
$C_z\ge20K$ of the $SU(2)$ construction holding besides, and with the further member of the radius of the second
chain that Lemma~\ref{C-368-10}(iv) names, which keeps its logarithms in $\mathfrak g^c$; their constants are
constants depending only on $L$ and the pair, explicit in $L$ and $n$ at $SU(N)$, that obey the conditions in which
they enter: there the radii are bounded from above only and the bounds from below only, so the ordered choice
closes again provided every subsequent use of the normalizer and elimination constants requires only these upper bounds on radii and lower bounds on estimate constants. This compatibility with the later inequalities remains part of (K). Write $E^{\rm fam}$ for the local Gaussian bound constant, the counterpart of $E_b$ in Proposition~\YMMref{50.2}{386}; here $E_b$ already denotes the quotient of Proposition~\ref{C-368-16}(iv). Require
\[
 E^{\rm fam}\ge e_n:=\frac{2(2n\log L+\rho_n)}{L^3+1}.
\]
Indeed, Lemma~\ref{C-368-10}(iv) then gives
\[
 |z_c|\le2n\log L+\rho_n
          \le\tfrac12(L^3+1)E^{\rm fam},
\]
which is the bound needed for the local Gaussian constant at an elimination term. The chosen constants must satisfy this guard; it is not inferred merely from a dimension-independent lower bound on $E^{\rm fam}$. The inequality $\frac34\le1$ by which the proof of Lemma~\ref{lem:colour-determinant-count} absorbs a
colour factor reads $\frac n4\le1$, which fails from $n=5$; a factor $2^{-9}$ of the same estimate, which that
proof does not use, restores it for $n\le2048$, and for every $n$ the precise bound is \eqref{eq:colour-determinant-bound}. Here $C_{\rm mob}$ is the full constant, including the zero-size term, with its trace contribution evaluated in dimension $n$. $(\frac54)^3<2$ becomes $(\frac54)^n$, which fails from $n=4$, and
the bound $z<2L^9$ of Lemma~\YMMref{65.3}{581} becomes $z<(\frac54)^nL^{3n}$. The devices by which the
$SU(2)$ construction fixes the branch of a logarithm or of a root or power of a determinant ask something of the
operator and not of the value of its determinant, and serve in every dimension, but for those that read the value:
the principal logarithm of the elimination factor (Lemma~\ref{C-368-10}(iii)); the power series of the logarithm
about $1$ in the Jacobian of the exponential map (Lemma~\ref{C-368-10}(i)); the power series of the logarithm about
$1$ at the Jacobian $1-\mathsf{jc}$ of the bond shift with its cut-off (the cutoff preceding Lemma~\YMMref{53.3}{444}), $\mathsf{jc}$
being the derivative of the cut-off in the direction of the shift, so that $1-\mathsf{jc}$ is the determinant of a
perturbation of rank one of a map of determinant $1$, whose value is computed outright and in which no count of the
dimension enters; and the power series of the logarithm about $1$ at a quotient $\det(\mathsf T(0)^{-1}\mathsf
T(\zeta))$ of the determinants of a holomorphic family $\mathsf T(\zeta)$ of operators on $\mathfrak g^c$-valued
fields, formed near $\zeta=0$, where the quotient is near $1$ by continuity. Among the devices that ask something
of the operator is the number $2n\log L-\mathrm{tr}\log T$ of the elimination, $T=L^3\mathrm{Ad}(g)\,\mathsf
m\,\mathrm{Ad}(g')^{-1}$ with $g,g'\in G$ as in Lemma~\ref{C-368-10}(iii): it asks $|T-1|<1$ at some $g,g'$, and it
does not depend on their choice among those with $|T-1|<1$, for every $n$ (Lemma~\ref{C-368-10}(iii)). Where a
count of the dimension, $n=3$ against $\pi$ or $2\pi$ or $n-\mathrm{rk}\,G=2$ against $1$, fixes a branch, as in
the determinant argument in Proposition~\YMMref{50.2}{386}, the elimination estimates of Lemma~\ref{C-368-10}(iii), (iv), the auxiliary-frame independence proved in Lemma~\ref{C-368-10}(iii) and the bound for the Gaussian constants just proved,
the radius $r_n$ and the numbers of Lemma~\ref{C-368-10}(iv) take its place, or, where it compares $\mathrm{tr}\log
T$ at two choices of $g,g'$, the equality of Lemma~\ref{C-368-10}(iii); and where the sign
$(-1)^3=(-1)^{\dim\mathfrak g}$ stands twice in one identity, in the adjoint determinant calculation, it cancels for
every $n$. That the branch is fixed in one of these ways at every place where the $SU(2)$ construction forms a
logarithm or a root or power of a determinant over $\mathfrak g^c$ is part of (II) and (K).
\item[(e)] If $z_G<1$, the conditions that fall to $z_G$ include a ceiling on $\gamma$, a floor on $L$ and the
radii that the $SU(2)$ construction defines as minima with the averaging radius. For the rows that fall to $z_G$
the floor on $L$ is $L^2\ge\frac1{200z_G}$ together with the row $\alpha_0\le z_G^2/(4C^2)$, $C$ of
Lemma~\ref{C-368-8}(iv), or $L>\frac1{z_G}$; the floor $L>2z_G^{-2}$, which implies $L>\frac1{z_G}$ and
$200L^2z_G\ge1$, carries with a ceiling on $\gamma$ all the relations in $N$ of the table. The numbers
$\frac1{200}$, $\frac12$, $\frac13$, the radius $1$ and the like do not fall, for there ``$X\in D^c$'' is supplied
by a product or a path; the mean and the block average keep their numerals under $\omega_M\le z_G$, which the floor
$200L^2z_G\ge1$ gives; so do the charts, the gluing of gauges and the gauges at fields whose plaquette variables
lie in $D$ under the averaging radius, while the gauges that Definition~\ref{C-368-3}(3) names take there the forms
named: the comb gauge of Lemma~\YMMref{14.4}{64} takes the guard $\varpi<z_G$ and Lemma~\ref{lem:group-contour-gauge} the guard
$\varpi_y(W)<z_G$, which bite at $SU(N)$ from $N=7$, and the quantitative cube-gauge input of Lemma~\ref{C-368-8}(iii) takes
constants $c$, $B$ of the pair; the box clause of Lemma~\YMMref{14.4}{64} read with some gauge of the box in the
place of the comb holds at $SU(N)$ exactly for $N\le13$ (Lemma~\ref{C-368-8}(iii)); an element known by a chord
$|V-V'|$ alone takes $z_G$ in its radius; and the statements that place a logarithm in $\mathfrak g^c$ by
``$|X-\mathbf 1|<1$'' are false for $z_G<1$ as worded and hold with ``$X\in D^c$''. Among the conditions that the
centre enters, a condition can bound $N$ only through an inequality with no free constant, $b<2\sin\frac\pi N$,
where a proof concludes from ``$X$ is within $b$ of $\mathbf 1$'' that $\log X$ is traceless; read with $D^c$, such
a conclusion asks instead for a product or a path, while a ball of radius $z_G$ in the place of $D^c$ would let the
number $\frac1{14}$ of Lemma~\ref{C-368-7}(i) cap $N$ at $87$. The other inequalities with no free constant that
bound $n$ or $N$ take the forms named at their places: $\frac n4\le1$, $(\frac54)^n<2$ and $|\mathbf
1-T|\le\frac12$ in (d), $|\log j|<1$ on the ball of radius $\frac38$ in Lemma~\ref{C-368-10}(i), and the ratio of
the two norms in (c). A field whose bond variables lie within $r$ of $\mathbf 1$, $(1+r)^4<2$, is $\zeta e^{iA}$
with $\zeta(b)=z(U(b))\in Z(G)$ and $\mathfrak g$-valued $A$, and $\partial\zeta(p)=z(\partial U(p))$; so $\zeta$
is flat exactly when every plaquette variable lies in $D^c$, and a flat $\zeta$ is a central gauge transform of
$\mathbf 1$ on a box and is fixed by its three holonomies up to such a transform on a torus. A gauge function $h$
on a connected set with $|h(x)-\mathbf 1|<1$ and $h(x)h(y)^{-1}\in D^c$ at neighbouring sites $x$, $y$, as when $h$
carries a field with bond variables in $D$ to another such field, is a central constant times $e^{i\psi}$, $\psi$
being $\mathfrak g$-valued. The sets of fields cut out by plaquette tests on a block, a double block or a tower are
connected to the field $\mathbf 1$ inside themselves, and a field that passes the plaquette test $t$ on a box is
joined to $\mathbf 1$ inside the set cut out by the test $t+4.49t^2$; and at these places no two-sided condition
becomes empty. That ``$X\in D^c$'' holds at every place where a logarithm is taken, the places being closed under
the statements that inherit a radius, is part of (K).
\item[(f)] For every $G$, also at $z_G=1$, the conditions include $\varepsilon_{\rm v}\le L^2r_G$ and
$\varepsilon_{\rm v}\le a_G$, where a positive orbit radius $a_G$ admitted by Definition~\ref{C-368-12} is fixed as part of the analytic input.
No unconditional positivity or two-sided $SU(N)$ estimate for the supremum $a_*(G)$ is used. At $SU(N)$, $N\ge3$, the central-field obstruction
at $\varepsilon_{\rm v}>2L^2\sin(\pi/N)$ applies if the chosen averaging extension sends the displayed central field to $\mathbf 1$:
the nonflat zero-current field of Definition~\ref{C-368-12} is then admitted at the first level. The conditions also include
$\varepsilon_{\rm v}\le\varepsilon_G$ and $c_{11}\le
r_G/(2C_{\rm pg})$ of the uniqueness clause of Definition~\ref{C-368-31}(V), $c_{11}$ the radius of
Definition~\ref{def:cg-uniform-slice} and $C_{\rm pg}$ the admitted preliminary-gauge constant of Definition~\ref{def:cg-uniform-slice}; the
ceiling $\mathsf c_Lc_4\le2Lr_G$ of Proposition~\ref{C-368-16}(iv), a positive ceiling on $c_4$, void at $SU(3)$;
the row $\alpha_1<\rho_GL$ of Proposition~\ref{C-368-17}, a positive ceiling on $\alpha_1$; and, at the memberships
reached from memberships of restrictions had apart, the conditions of Lemma~\ref{C-368-15}(iii): the moved first
radius $b'_0=b_0+2e^{4\xi b_1}b_1^2$, for regular pairs $\alpha_0+2\vartheta^2e^{4\vartheta\alpha_1}\alpha_1^2$, at
most the first radius $B_0$ at which the statement read there is had, and for the analyticity cores besides $b_1\le
B_1$, $K(b'_0+b_1^2)\le B_0$ and $b'_0+b_1\le a^\sharp$, with $K\ge16L^2$ and $a^\sharp$ of that lemma; where
$b_0<B_0$, and for the cores $Kb_0<B_0$ and $b_0<a^\sharp$, these are positive ceilings on the second radius.
\item[(g)] The constants in the matrix product, logarithm and
contour comparisons of Lemmas~\ref{C-368-6}(ii),
\ref{C-368-7}(iii) and~\ref{C-368-14} have the stated dependence
on the pair. The $SU(2)$ contour coefficients are those of
Lemmas~\YMMref{14.5}{64} and~\YMMref{14.6}{65}; the accumulated
transport coefficient is constructed in the proof of
Lemma~\YMMref{30.3}{240}. Exact numerical choices for a general-group
Hessian or nonlinear-remainder bound must be supplied as part of
(P). In particular, this paragraph does not deduce the values
$12,1536,3456$, the action radius and bound $1/128,808$, or the
gauge radius and bound $1/64,93$ from matrix norm comparison alone. Of the
constants of the first part of the $SU(2)$ construction these move with $N$: the averaging radius; the $2$ of
$|T^{\mathrm T}|\le2|T|$ (Lemma~\ref{lem:group-transpose-depth-zero}), which becomes $N$, sharply; a numeral $3$ that counts
colours in a bound of Schur's kind of the opening lemmas of the propagator construction
(Proposition~\ref{C-368-16}(i)), which becomes $n$ (Definition~\ref{C-368-3}(2)); the unmodified $SU(2)$ coefficient $4$ in Lemma~\ref{lem:group-transpose-depth-zero}
at depth $0$, which becomes $4\varsigma_G$, at most $8(1-\frac1N)$ for $SU(N)$; the averaging constant $C_Q=18NL^3C_{\rm av}$ of
Lemma~\ref{lem:group-averaging-row-constant}, with its admitted kernel
bound. When $C_{\rm av}=3$ is supplied, its value is $54NL^3$, compared
with $108L^3$ at $SU(2)$; the constants formed from $6LC$ and $12LC$
then move by $N/2$ and their reciprocal thresholds by $2/N$; and one constant, which takes the factor
$\kappa_G$, for $SU(N)$ and radii $r\le z_G$ the factor $1+\kappa'$, the thresholds on $\alpha_0$ being divided by
it. Two bounds of the $SU(2)$ construction take new numbers besides (Definition~\ref{C-368-3}(2), (3),
Remark~\ref{C-368-30}): the $12$ of the bound $\xi^3|J_W(U)(b)|\le12a$ on the current of a region at complex fields
whose bond variables have condition number at most $2$ becomes $6\hat\varsigma_G$, the $12$ failing at $SU(N)$ from
$N=5$; and the factor $1$ of the bound $|\mathrm{Bl}(\mathsf g,\mathsf f)|\le\|\mathsf g\|_2\|\mathsf f\|_2$ drawn
by polarisation becomes $\kappa_G^2$, and stays $1$ where $|X+iY|^2\le|X|^2+|Y|^2$ on $\mathfrak g$. Further
numbers move in Definition~\ref{C-368-3}(2) and Lemmas~\ref{C-368-8}(ii), (iii), \ref{C-368-9}(i)
and~\ref{C-368-10}(i), (ii). The constants of the propagators and the covariances, those of the critical
configurations and those of the sections, which rest on both, are set by the statements of
Proposition~\ref{C-368-16} and of (P), (V); and all relations among the constants follow from the form of the
displays. A condition of the table need not stay true when a threshold that stands in a conclusion is lowered: the
floor $\delta_{2,\mathrm b}R_1\mathfrak M\ge64(3\log L+\log2)$ in the required separation conditions of (K) fails when the rate
$\delta_{2,\mathrm b}$ is lowered with $R_1$ and $\mathfrak M$ fixed, already at $SU(2)$; the ordered choice of (a)
uses no such monotonicity.
\end{enumerate}
\end{proposition}

\begin{proof}[Sketch of proof]
(a), (b) The table of (a) follows the ordered choice of Proposition~\YMMref{69.1}{595}, with the numbers of $G$ as data and the compatibility required in Definition~\ref{C-368-31}(K): the constants of
$SU(2)$ are kept, the value for $G$ stands behind each function symbol, and new letters enter only where a formula
needs them; a floor on a constant free upward or a positive ceiling on one free downward keeps the shape admissible
and leaves the admissible sets of the other constants as they are, and a condition whose last constant is
determined is read off its formula. The recursion that defines large sets commutes with intersections, so the
intersection of two sets is large exactly when both are; a large set holds, at each choice of the earlier constants
at which it is large, a half-line, an interval $(0,c]$ or the whole window of the next constant, so one floor or
ceiling on each free constant, laid in the order of the choice, makes every admissible choice lie in it, and a set
in which every choice of such a table lies is large, provided the original two-sided windows have been retained whole. Further ends on a two-sided constant
instead restrict its window. Their joint feasibility at an earlier choice means exactly that the original window meets their intersection;
this does not make the ends large in the original table. The radius stands in rows on $\gamma$ alone, so the tables at two radii differ in rows on $\gamma$ alone
and have the same admissible choices before $\gamma$; a set that reads $\gamma$ or $\rho$ and is large is met by
laying its set of earlier choices as rows before $\gamma$, which read no $\rho$, and then at each $\rho$ a ceiling
on $\gamma$. (c) On $\mathfrak{su}(N)$, $|\cdot|\le(N-1)^{1/2}\|\cdot\|$ (Lemma~\ref{C-368-6}(i)), that is
$\|\cdot\|^2\ge\frac1{N-1}|\cdot|^2$, and the conditions that compare the norms hold together at this ratio of the
squares. The two-sided condition on $\varepsilon_{\rm v}$ is that of the inclusions of Lemma~\ref{C-368-8}(ii),
valid for $r\le z_G$: its window, which uses the two inclusions, is not empty with the numbers of the critical
configuration of (V), and its ends are made compatible by one floor on $R_1$. (d) Lemma~\ref{C-368-10}(iv): the
chains need $4Ka\le r_n$, which keeps the normalised factor within $r_n$ of $1$, and $C_z\ge4c_nK$; for $n\ge4$ the
radii of the $SU(2)$ construction fall by the factor $2r_n$ and its constant $C_z$ rises by the factor
$\frac{c_n}5$, which gives both, the pair $4Ka\le\frac12$, $C_z\ge20K$ of the $SU(2)$ construction holding besides
as $r_n<\frac12$ and $c_n>5$; for $n=2,3$ the constants of the $SU(2)$ construction serve, as $r_n\ge\frac12$ and
$c_n<5$; and the further member of the radius of the second chain keeps the logarithms of that chain in $\mathfrak
g^c$. The other devices, a trace series, an integral along a segment, a Gaussian integral, positivity along a path,
uniqueness by continuity and the number $2n\log L-\mathrm{tr}\log T$ of the elimination, ask their conditions of
the operator or of each of its eigenvalues by itself, and so read no dimension. The last asks $|T-1|<1$ at one
choice of $g,g'$; at two choices $T_1$, $T_2$ with $|T_j-1|\le r<1$, $T_2$ is conjugate to $T_1W^{-1}$,
$W=\mathrm{Ad}(h)$ with $h\in G$, and $\mathrm{tr}\log X-\mathrm{tr}\log(XW^{-1})$ lies in $2\pi i\mathbb Z$, is
constant on the convex set $\{|X-1|\le r,\ |X-W|\le r\}$ and vanishes at $\frac12(1+W)$, the eigen-angles of $W$
being symmetric, so that the number is the same at both, for every $n$ (Lemma~\ref{C-368-10}(iii)). Of the devices
that read the value, the series at the Jacobian of the bond shift and the series at the quotient of determinants
read it without a count: $\det(A+u\otimes v)=\det A\,(1+v(A^{-1}u))$ makes the Jacobian of the bond shift, a map of
determinant $1$ plus the rank-one term that the cut-off brings, the scalar $1-\mathsf{jc}$, and the bound
$8\|\vartheta_G'\|_\infty|H_b|\le\delta_k$ on the shift keeps $|\mathsf{jc}|$ below $1$, the Jacobian lying in
$[\frac78,\frac98]$ (Proposition~\ref{C-368-16}(iv)); and $\det(\mathsf T(0)^{-1}\mathsf T(\zeta))$ is $1$ at
$\zeta=0$, so within $1$ of $1$ near $\zeta=0$. (e) The central factor $z(X)$ is a local homomorphism, invariant
under conjugation in $G^c$, and $|\log X|\ge\theta(z(X))$; so the inequality by which a product is shown to lie
near $\mathbf 1$ shows its membership in $D^c$, and $\{|\log X|<2\arcsin\frac{z_G}2\}\subset D^c$ exactly
(Lemma~\ref{C-368-6}(iii)); a radius known only through a chord meets the central elements at distance $z_G$.
Around a plaquette whose bond variables lie within $r$ of $\mathbf 1$, $(1+r)^4<2$, the partial products stay in
the unit ball, so $z(\partial U(p))=\partial\zeta(p)$; and a gauge function with $h(x)h(y)^{-1}\in D^c$ at
neighbours has $z(h(x))=z(h(y))$ there, so $z\circ h$ is constant on a connected set. (f) Use the admitted orbit radius of Definition~\ref{C-368-12}; for an averaging extension that sends its central field to $\mathbf 1$: at
$SU(N)$ the field with $e^{2\pi i/N}\mathbf 1$ on one bond off the block axes is critical, without current,
averages to $\mathbf 1$ and is no pure gauge, and its plaquettes are at distance $2\sin\frac\pi N$ from $\mathbf
1$, so it is admitted at the first level as soon as $2\sin\frac\pi N<\varepsilon_{\rm v}L^{-2}$; the further
conditions of (f) are those that Definition~\ref{C-368-31}(V), Propositions~\ref{C-368-16}(iv), \ref{C-368-17} and
Lemma~\ref{C-368-15}(iii) set. (g) Each constant is given by the statement that sets it, in which $N$ enters only
as named; $|T^{\mathrm T}|\le N|T|$ is sharp, $\kappa^2=N-1$ being attained by general kernels; $4\varsigma_G$,
$6\hat\varsigma_G$ and $\kappa_G^2$ are those of Remark~\ref{C-368-30}, where $8(1-\frac1N)$ is the bound at
$SU(N)$ and $4\varsigma_G$ exceeds $8$ at the spin representation of $Spin(2m+1)$ from $m=43$. The displayed separation floor bounds from below a product in which $\delta_{2,\mathrm b}$ is a factor, so it is lost when
$\delta_{2,\mathrm b}$ alone is lowered; in the ordered choice each constant is chosen after those it depends on.
\end{proof}

\begin{proposition}[Reflection positivity and Gibbs states]\label{C-368-19}
Let $G\ne\{\mathbf 1\}$ be a closed subgroup of $U(N)$ and $\beta\ge0$, and consider the Wilson measure of a periodic lattice.
\begin{enumerate}
\item[(i)] With $\hat g=\mathrm{diag}(g,\bar g)$,
\[e^{-\beta a(gg'^{-1})}=e^{-\beta}\sum_{\mathsf n\ge0}\frac{(\beta/2N)^{\mathsf n}}{\mathsf n!}\sum_{I,J}\hat g_{IJ}\,\overline{\hat g'_{IJ}},\]
absolutely and uniformly convergent, with positive coefficients, where $I=(i_1,\dots,i_{\mathsf n})$ and
$J=(j_1,\dots,j_{\mathsf n})$ run over $\{1,\dots,2N\}^{\mathsf n}$ and $\hat g_{IJ}=\prod_k\hat g_{i_kj_k}$, the
empty product being $1$; so the measure is reflection positive at every mirror $v\cdot x=c$, $v=e_\mu$ with
$2c\in\mathbb Z$ or $v=e_\mu\pm e_\nu$ with $c\in\mathbb Z$, whose period lattice, with periods at least $3$,
splits along $v$, for every bounded Borel function of the bonds on one side, gauge invariant or not; and for
$\beta>0$ every $\lambda_\pi(\beta)=d_\pi^{-1}\int e^{-\beta a}\chi_\pi$ obeys $\lambda_\pi(\beta)\ge
e^{-\beta}\frac{(\beta/2N)^{\mathsf n(\pi)}}{\mathsf n(\pi)!\,d_\pi}>0$, $\mathsf n(\pi)$ being an $\mathsf n$ for
which $\pi$ occurs in the $\mathsf n$-th tensor power of $\mathbb C^N\oplus\overline{\mathbb C^N}$. At a mirror
through the midpoints of bonds reflection positivity holds if and only if the crossing operator $\mathsf K$ is
positive; $\mathsf K$ is the average over gauge transformations times the convolution with the weight on every
bond, with the eigenvalues $\prod_bc_{\pi_b}$, $c_\pi=d_\pi^{-1}\int k\,\chi_\pi$ for a real even class weight $k$
in the place of $e^{-\beta a}$, so that $c_\pi=\lambda_\pi(\beta)$ at $k=e^{-\beta a}$; two site mirrors ask
nothing of the weight, a bond mirror asks these products to be non-negative, and a diagonal pair asks every
$c_\pi\ge0$.
\item[(ii)] The two chessboard estimates of the $SU(2)$ construction hold, with mirror copies conjugated at odd reflection parity in the slab products and in $\mathrm{Fold}$:

\[
|\mu(\prod_{j\le2M}f_j^{[j]})|\le\prod_j\mu(\prod_if_j^{[i]})^{1/2M}
\]
 for slabs and 
\[
|\langle\prod_\kappa
F_\kappa\rangle_\beta|\le\prod_\kappa\langle\mathrm{Fold}\,F_\kappa\rangle_\beta^{1/\mathsf N}
\]
 for cubic cells,
and so do the statements of the $SU(2)$ construction that follow them: for sources $\|e^{\mathcal O(w)}\|_{\rm
ch}^{\mathsf N}=Z_\beta[w^{\rm ch}]/Z_\beta\ge0$ (equation~\YMMref{(1759)}{682}); subadditivity,

\[
Z_{\underline{\mathsf n}}[\mathrm{Fold}\,F]\le Z_{\underline{\mathsf n}'}[\mathrm{Fold}\,F]^{\prod_\mu\mathsf
n_\mu/\mathsf n'_\mu}
\]
 for boards with even numbers $\mathsf n'_\mu\le\mathsf n_\mu$ of cells along the axes and
one function $F$ of a cell on both (Proposition~\YMMref{86.5}{682}); and the volume (Corollary~\YMMref{86.6}{683}): on
tori of sides $\ell\ge\ell'\ge2\ell_{\rm c}$, $\ell_{\rm c}$ the side of the cells, with $p_\ell=\ell^{-3}\log
Z_\beta$ on the torus of side $\ell$, $p_\ell\le p_{\ell'}\le0$ and $\|F\|_{{\rm ch},\ell}\le\|F\|_{{\rm
ch},\ell'}e^{\ell_{\rm c}^3(p_{\ell'}-p_\ell)}$. The comparison of the Gibbs kernel of a box with the torus measure
(Proposition~\YMMref{91.8}{704}) holds for every such $G$, as it reads of the group only $0\le a_p\le2$. The Wilson
measure on a torus and on the periodic configurations of the infinite lattice are one measure. Given $\mathsf
K\ge0$ in the three directions, reflection positivity and the two chessboard estimates hold for every real bounded
even class function in the place of $\beta a$, and so do the statements that follow them, with $p_\ell\le
p_{\ell'}$, and $p_{\ell'}\le0$ for a non-negative function.
\item[(iii)] For representations $\rho_i$ whose sum is the inclusion and all $\beta_i>0$, every periodic Gibbs
state of $\prod_p\exp[-\sum_i\beta_i(1-\frac1{N_i}\mathrm{Re}\,\mathrm{tr}\,\rho_i(\partial U(p)))]$ is invariant
under the reflections in bond mirrors and under the lattice isometries, and is reflection positive at the nine
families of mirrors; for one block this holds also at $\beta=0$; and $\nu(\mathrm{Im}\,\chi_\pi(\partial U(p)))=0$
at every plaquette. The evenness $k(g^{-1})=k(g)$ of the weight cannot be dropped.
\item[(iv)] Let $G$ be compact, connected and semisimple, $w(V)=P[\frac1{2i}(V-V^*)]$, and $W$ the holonomy of a
plaquette $p\ni b$ based at $b$. Then $\sum_a\frac d{ds}\big|_{s=0}\langle t_a,w(e^{ist_a}W)\rangle=n-k_p$. With
the isotypic decomposition $\mathbb C^N=\bigoplus_sW_s\otimes\mathbb C^{k_s}$, $d_s=\dim W_s$, $\Omega$ acting on
$W_s$ by $\omega_s$ and $w_s=1-\mathrm{Re}\,\chi_s/d_s$: $a_p=\sum_s\frac{d_sk_s}Nw_s(\partial U(p))$ and
$k_p=\sum_s\omega_s\frac{d_sk_s}Nw_s(\partial U(p))$; so $k_p$ is a multiple of $a_p$ if and only if the $\omega_s$
of the non-trivial constituents agree, and then $k_p=\frac{nN}{N-k_0}a_p$, $k_0$ the multiplicity of the trivial
constituent, on which $w_s=0$; under this condition, $k_p=na_p$ if no constituent is trivial, as for an irreducible
representation and for sums of copies of one, and $k_p=\frac92a_p$ at $\{\mathrm{diag}(u,1):u\in SU(2)\}\subset
U(3)$. For $SU(2)$ on $\mathbb C^2\oplus\mathbb C^3$, with $a_p=1-\frac15\mathrm{Re}\,\mathrm{tr}\,\partial U(p)$,
no multiple of $a_p$ serves for $k_p$. With the weight $\Lambda=\mathrm{diag}(0.26,0.26,0.16,0.16,0.16)$ in the
commutant, at first order of strong coupling, the integrated identity holds with the Casimir density, both sides
being $0.206$, and fails with $\frac n2a_p$ in the place of $\frac12k_p$, the side that carries it becoming
$0.2796$ against $0.206$. The same holds for an action with a weight $\Lambda$ in the commutant: the bonds move by
the vector part $w$, the action is differentiated by the gradient part, which is $P[N\Lambda(V-V^*)/2i]$ for the
pairing $\langle\cdot,\cdot\rangle$, $w$ itself at $\Lambda=\frac1N\mathbf 1$, and $P[\Lambda(V-V^*)/2i]$ for the
pairing $\mathrm{tr}\,XY$, and the flux carries the same Casimir density.
\end{enumerate}
\end{proposition}

\begin{proof}[Sketch of proof]
(i) We expand $e^{\frac\beta{2N}\mathrm{tr}(\hat g\hat g'^*)}$; every term is $f(g)\overline{f(g')}$ with a
positive coefficient, which gives reflection positivity, and one term bounds $\lambda_\pi$; the square root of the
weight for the convolution, $k(gh^{-1})=\int_Gq(gx^{-1})q(xh^{-1})\,dx$ with $q$ a real continuous even class
function whose coefficients are the $\lambda_\pi^{1/2}$, makes each plaquette cut by a mirror a square, $\partial
U(p)=g_p(U)g_p(\theta U)^{-1}$; at a bond mirror the partition function is a trace of powers of $\mathsf
K^{1/2}\mathsf S\mathsf K^{1/2}$, $\mathsf K$ holding the crossing plaquettes alone and $\mathsf S$, the slab
operator, the plaquettes inside a layer.

(ii) The partition function with insertions is the trace of a product in which the crossing operator alternates
with the slab operators, which carry the plaquettes inside a layer and the insertions; the eigenvalues of the
crossing operator are products of the $\lambda_\pi\ge0$, so $\mathsf K\ge0$, and H\"older's inequality for traces,
$|\mathrm{Tr}(A_{2M}\cdots A_1)|\le\prod_j\|A_j\|_{2M}$, gives the estimates; where the proofs of the chessboard
estimates of the $SU(2)$ construction (Proposition~\YMMref{86.3}{681}) use $\mathrm{tr}\,X^{-1}=\mathrm{tr}\,X$, false at
$i\in U(1)$, they need only $a(X^{-1})=a(X)=a(hXh^{-1})$, that is, that flipping a plaquette keeps its action.
Reflections and translations keep action, observable and measure; folded integrals on two tori are traces of powers
of one positive operator, so that $\mathrm{Tr}\,B^{\mathsf p}\le(\mathrm{Tr}\,B^{\mathsf p'})^{\mathsf p/\mathsf
p'}$ for $B\ge0$ and integers $\mathsf p\ge\mathsf p'\ge1$ gives subadditivity and the comparison of the volumes
for every weight with $\mathsf K\ge0$, and $Z_\beta\le1$, that is $p_{\ell'}\le0$, where the weight is at most $1$;
the proof of Proposition~\YMMref{91.8}{704} reads of the group only $0\le a_p\le2$, which holds on every $G\subset
U(N)$; and the cell estimate also follows from the slab estimate in three directions, with no trace formula.

(iii) $\log(Z_\Gamma[t\mathbb A]/Z_\Gamma[0])\ge t|\mathbb A|\nu(A)$ minus a surface term; the trace formula with
insertions has modulus at most $\prod\|X\|$; at $X=e^{\pm\varepsilon\mathsf D}$, $\mathsf D$ real antisymmetric,
$\nu(A_{\mathsf D})=0$, so neighbouring layers are exchangeable, and two layers fix the chain of layers because the
kernel is injective on measures, all $\lambda_\pi$ being positive. For the nine families, reflection positivity on
adapted tori, Schwarz's inequality with a longest run and a separation argument identify every periodic Gibbs state
with a reflection positive one.

(iv) The derivative of $\langle t_a,w(e^{ist_a}W)\rangle$ at $0$ is $\frac1N\mathrm{Re}\,\mathrm{tr}(t_a^2W)$;
summing gives $\frac1N\mathrm{Re}\,\mathrm{tr}(\Omega W)=n-k_p$. The decomposition follows from Schur's lemma. At
the trivial constituent $w_s=0$; a constituent $s$ and its conjugate $s^*$ have $w_s=w_{s^*}$ and
$\omega_s=\omega_{s^*}$; and the $w_s$ of the non-trivial constituents, one for each pair $\{s,s^*\}$, are linearly
independent functions on $G$, as $1$ and the real parts of their characters are, by the orthogonality of
characters. So $k_p=ca_p$ if and only if $\omega_s=c$ at every non-trivial constituent, and
$\frac1N\mathrm{tr}\,\Omega=\sum_s\omega_s\frac{d_sk_s}N=n$ gives $c=\frac{nN}{N-k_0}$. For the action
$1-\frac12\mathrm{tr}\,\Lambda(u+u^{-1})$ with $\Lambda$ in the commutant, $\frac
d{ds}\big|_{s=0}\bigl(1-\frac12\mathrm{tr}\,\Lambda(e^{isX}W+W^{-1}e^{-isX})\bigr)=\mathrm{tr}(X\Lambda
H_W)=\langle X,P[N\Lambda H_W]\rangle$ for $X\in\mathfrak g$, as $\Lambda$ commutes with $W$.
\end{proof}

\begin{lemma}[The coefficient of $\beta^{-3}$ in a box]\label{C-368-20}
Let $G\subset U(N)$ be compact and connected, $n\ge1$, semisimple or not, and consider the box $\{0,\dots,s\}^3$
with trivial Dirichlet data, $\beta\ge1$. Then $|\langle
a_p;a_q\rangle-\beta^{-2}B_0^{[s]}(p,q)-\beta^{-3}B_1^{[s]}(p,q)|\le C\beta^{-4}$, with $C$ depending on the box,
as in Proposition~\textup{\ref{IV-73-27}} for $SU(2)$, and
\[B_1^{[s]}(G)=\frac{\varphi_G}{24}\,B_1^{[s]}(SU(2))+c_G\,Y^{[s]},\qquad c_G=\frac1N\mathrm{tr}\,\Omega^2-\frac38\varphi_G,\]
$\varphi_G=\sum_a\mathrm{tr}(\mathrm{ad}\,t_a)^2$ and $Y^{[s]}$ the coefficient for $U(1)$, $0\le Y^{[s]}(p,q)\le
C\langle x_q-x_p\rangle^{-6}$. As $s\to\infty$ the coefficient converges, the points being recentred at the centre
of the box, and summed over orientations its part at distance $\mathsf r$ is $\frac{\varphi_G}{16\pi^3}\mathsf
r^{-5}$ up to $C_G\mathsf r^{-6}(1+|\log\mathsf r|)$.
\end{lemma}

\begin{proof}[Sketch of proof]
Laplace's method about the gauge orbit of $\mathbf 1$ in the box gives the expansion. The Wick pairings that decay
like the fifth inverse power fall into six classes, and in each the structure constants enter through
$\sum_{abc}f_{abc}^2=\varphi_G$: the numeral $24$ of $SU(2)$ becomes $\varphi_G$ in five classes and $12$ becomes
$\frac12\varphi_G$ in one. The representation enters through $\Omega$ only in terms that decay like the sixth
power. The limit and the leading part are those of $SU(2)$ (Theorem~\textup{\ref{IV-73-71}}), whose constant
$\frac3{2\pi^3}$ is multiplied by $\frac{\varphi_G}{24}$.
\end{proof}

\begin{lemma}[Coefficients on the simple ideals]\label{C-368-21}
Let $\mathfrak g=\bigoplus_{i\le r}\mathfrak g_i$ be semisimple. The mean of the $n$ quadratic colour coefficients of a local
term of second order is $\sum_i\frac{n_i}n\mathsf{cf}_i$, $\mathsf{cf}_i$ its coefficient on $\mathfrak g_i$, and
$|\mathsf{cf}_i|\le\kappa_i^2C_{\rm am}r_0^{-2}E_0=(\kappa_i/\kappa_G)^2\mathsf G_\sharp$ for the terms of
Theorem~\YMMref{D}{8}, with $C_{\rm am}$, $r_0$ and $\mathsf G_\sharp$ of Proposition~\ref{C-368-17}. If
$|e^{(i)}_k|\le\Gamma$, the sequences $x^{(i)}_0=\beta$, $x^{(i)}_k=L^{-1}x^{(i)}_{k-1}-e^{(i)}_k$ and
$x_k=\sum_i\frac{n_i}nx^{(i)}_k$ satisfy $x_k=L^{-1}x_{k-1}-\sum_i\frac{n_i}ne^{(i)}_k$ and
\[|x^{(i)}_k-x_k|\le2\Bigl(1-\frac{n_i}n\Bigr)\Gamma\frac L{L-1}(1-L^{-k}).\]
\end{lemma}

\begin{proof}[Sketch of proof]
The first claim is the splitting of an orthonormal basis along the ideals; the bound is that of the $SU(2)$
construction with $|t_a|\le\kappa_i$ on $\mathfrak g_i$ (Proposition~\ref{C-368-17}). Further
$x^{(i)}_k-L^{-k}\beta=-\sum_{s\le k}L^{-(k-s)}e^{(i)}_s$, so $|x^{(i)}_k-x^{(m)}_k|\le2\Gamma\frac
L{L-1}(1-L^{-k})$, and $x^{(i)}_k-x_k=\sum_{m\ne i}\frac{n_m}n(x^{(i)}_k-x^{(m)}_k)$.
\end{proof}

\begin{proposition}[Semisimple Lie algebras and tied ideals]\label{C-368-26}
Let $\mathfrak g=\bigoplus_{i\le r}\mathfrak g_i$ be semisimple, $r\ge2$, and $G_i$ the subgroup with Lie algebra $\mathfrak g_i$.
\begin{enumerate}
\item[(i)] The averages $E_j(M)=\int_{G_j}hMh^{-1}\,dh$ and their products over $j\notin I$ are conditional
expectations; so the thresholds of the $SU(2)$ construction pass down to the ideals with the constant $1$, and back
up with the factor $r$.
\item[(ii)] For $u_i=e^{iZ_i}\in G_i$, $a(u_1\cdots u_r)=\sum_ia(u_i)-t(u)$ with the tie
$t=\frac14\sum_{i<j}\tau_{ij}\langle Z_i,Z_i\rangle\langle Z_j,Z_j\rangle+O(Z^6)$, $\tau_{ij}$ of
Proposition~\ref{C-368-13}(iv); the tie vanishes exactly for block sums. In the Hessian of the action at $u=u_iv$,
$v=\prod_{j\ne i}u_j$, the block of the ideal $\mathfrak g_i$ is $\frac1N\mathrm{Re}\,\mathrm{tr}(Sc(u)T)$,
$S,T\in\mathfrak g_i$, with $c(u)=c(u_i)c(v)-H_{u_i}H_v$: it is the Hessian of the simple pair of $\mathfrak g_i$
with the weight $c(v)/N$, which is hermitian, lies in the commutant of $G_i$, is positive near $\mathbf 1$ and
reads the field $v$ of the other factors, less the term $\frac1N\mathrm{Re}\,\mathrm{tr}(H_vSH_{u_i}T)$, at most
$|S||T||u_i-\mathbf 1||v-\mathbf 1|$. This term vanishes when on each constituent $W_s\otimes M_s$ of $\mathbb
C^N$, $W_s$ irreducible under $G_i$ and $M_s$ its space of multiplicity under the other factors, one of $W_s$ and
$M_s$ is self-conjugate, so at every ideal $SU(2)$ and at $SU(2)\times SU(3)$ on $\mathbb C^2\otimes\mathbb C^3$;
at $SU(3)\times SU(3)$ on $\mathbb C^3\otimes\mathbb C^3$ it is not zero in general. The gradient on $\mathfrak
g_i$ carries the like term $P_i[c(u_i)H_v]$ beside $P_i[c(v)H_{u_i}]$. Between two ideals stands one mixed block
$\frac1N\mathrm{Re}\,\mathrm{tr}(SS'(c(u)-\mathbf 1))$, $S\in\mathfrak g_i$, $S'\in\mathfrak g_j$, $i\ne j$, at
most $|S||S'||u_i-\mathbf 1||u_j-\mathbf 1|$, which vanishes at $u=\mathbf 1$ and need not vanish at a background
(at $SU(2)\times SU(2)$ on $\mathbb C^2\otimes\mathbb C^2$, with $S=\sigma_3\otimes\mathbf 1$, $S'=\mathbf
1\otimes\sigma_3$ and $u=e^{it\sigma_3}\otimes e^{it\sigma_3}$, it is $-\sin^2t$). The block averages on small
fields and the exponential map are products over the factors, the density $j_G=\prod_ij_{G_i}\circ P_i$ is a
product, and the quadratic part of the gauge-fixing term is a sum over the ideals; the gauge-fixing term $a(X)$
itself is $\sum_ia(X_i)-t(X)$, with the tie of fourth order; and the domain $\{|Y|<\pi\}$ of the chart, $G^\pi$
with its indicator and the cut-off $\nu$ are compared with those of the factors through (i), with the constant $1$
down to the ideals and the factor $r$ back up.
\item[(iii)] Let $F_i=P_iF$ be the part of the field strength on $\mathfrak g_i$, and $Q_i=\sum_p\frac12\langle
P_iH_{\partial U(p)},P_iH_{\partial U(p)}\rangle$, the sum running over the bare plaquettes. By the transferred
statements of the step on its lists, its potentials and their edge parts (Proposition~\ref{C-368-17}), the leftover
$\sum_i(\mathsf{cf}_i-\mathsf{cf})\frac12\langle F_i,F_i\rangle$ of a level splits off at the level of the
potentials, with size at most
$\frac74\alpha_0^2\max_i|\mathsf{cf}_i-\mathsf{cf}|\le\frac72(1-\min_i\frac{n_i}n)\alpha_0^2\max_i|\mathsf{cf}_i|$,
its edge parts $8M^3$ times this; its terms of first order are of third order in the field strength and are bounded
by Cauchy's estimate whether or not they vanish; in dimension three the part of second order of an earlier level is
bounded and is not taken out; the leftover of level $j$ is $L^j\sum_i(e^{(i)}_j-b_j)Q_i$, $e^{(i)}_j$ the
coefficient on $\mathfrak g_i$ and $b_j$ their mean; and the leftovers of all levels add up to
$-L^k\sum_i(x^{(i)}_k-x_k)Q_i$, which is bounded uniformly in $k$, as $L^kQ_i$ is of the order $\alpha_0^2$ on a
block of level $k$ and $|x^{(i)}_k-x_k|$ is bounded by Lemma~\ref{C-368-21}.
\item[(iv)] If the action of $G$ has property~(T) of Theorem~\ref{C-368-4}(h3) at a choice of the constants, then
at that choice one run, the mean, suffices: the $r$ couplings $x^{(i)}_k$ stay within $2(1-\frac{n_i}n)\mathsf
G_\sharp\frac L{L-1}$ of it (Lemma~\ref{C-368-21} with $\Gamma=\mathsf G_\sharp$, as
$|e^{(i)}_k|\le(\kappa_i/\kappa_G)^2\mathsf G_\sharp\le\mathsf G_\sharp$, $\kappa_i\le\kappa_G$ since $\mathfrak
g_i\subset\mathfrak g$), and all clauses on the increments hold, the increment theorem, the regular progression,
the sequence and the bounds on polydiscs among them; a format with one exponent, one integrand, one total and a
carried term is the format of one run.
\item[(v)] At the least tied pair, $SU(2)\times SU(2)$ acting on $\mathbb C^4\oplus\mathbb C^2$ by
$(u_1,u_2)\mapsto(u_1\otimes u_2)\oplus u_2$ with every weight, that is, with the action
$\lambda_1(1-\frac14\mathrm{Re}\,\mathrm{tr}(u_1\otimes
u_2))+(1-\lambda_1)(1-\frac12\mathrm{Re}\,\mathrm{tr}\,u_2)$ for every $\lambda_1\in(0,1)$, one has
$1-a=(1-a^{(2)})(1-\lambda_1a^{(1)})$, $a^{(i)}=1-\frac12\mathrm{tr}\,u_i$; so the gradient and the Hessian of the
action are the weighted ones of the $SU(2)$ construction, with the plaquette weights $\lambda_1(1-a^{(2)}_p)$ on
the first factor and $1-\lambda_1a^{(1)}_p$ on the second, plus the block $-\lambda_1(\nabla a^{(1)}\otimes\nabla
a^{(2)}+\nabla a^{(2)}\otimes\nabla a^{(1)})$, which annihilates infinitesimal gauge directions. The fluctuation operator of the first
step is $A^\sharp(p_1)\oplus A^\sharp(p_2)+C^{\mathrm T}TC$, where $p_1$, $p_2$ are the points of the two factors,
each with its plaquette weights, $A^\sharp$ is the fluctuation operator of the $SU(2)$ construction, $C$ is the
linear map from the fluctuation fields to the forms on the bonds through which the $SU(2)$ construction forms
$A^\sharp$ from the Hessian of the action, and $T$, the Hessian of the action less the direct sum of the Hessians
of the two factors, is symmetric, of the range of one plaquette, vanishes at $\mathbf 1$ and at fields with values
in one factor, and obeys $|(TA)(b)|\le36300\max_{p\ni b}|\partial U(p)-\mathbf 1|^2\,|A|_{(b)}$, $|A|_{(b)}$ the
largest $|A|$ on the bonds of the plaquettes that hold $b$. For the propagator conclusion, assume at each level a regular domain on which the diagonal weighted operator $A_0=A^\sharp(p_1)\oplus A^\sharp(p_2)$ has the required inverse and decay bounds, and
$\|A_0^{-1}C^{\mathrm T}TC\|<1$ in the corresponding forcing and decay norm, uniformly in the level. Under these additional analytic inputs the Neumann inverse exists. The kinematic, interaction and Hessian conclusions require the corresponding regular-domain estimates; the algebraic identities alone do not establish them or property (T).
\end{enumerate}
\end{proposition}

\begin{proof}[Sketch of proof]
(i) $E_j$ is the average of the conjugation over a compact group, so it is a positive unital projection onto the
commutant of $G_j$, a conditional expectation, and does not increase $|\cdot|$, also at complex arguments; the
$E_j$ commute, and their product over $j\notin I$ is the average over the product of these $G_j$, a conditional
expectation onto its commutant. On $\mathfrak g^c$, $E_j$ is $0$ on $\mathfrak g_j^c$, a vector of the simple ideal
$\mathfrak g_j$ fixed by $\mathrm{Ad}(G_j)$ being $0$ (Lemma~\ref{C-368-11}(i)), and the identity on the other
ideals, which commute with $G_j$; so $\prod_{j\ne i}E_j=P_i$ there, whence $|P_iX|\le|X|$ and
$|X|\le\sum_i|P_iX|\le r\max_i|P_iX|$.

(ii) For the commuting $Z_i$, $1-\frac1N\mathrm{Re}\,\mathrm{tr}(u_1\cdots
u_r)=1-\frac1N\mathrm{tr}\cos(\sum_iZ_i)$. Its terms of second order are those of the $a(u_i)$, as
$\mathrm{tr}\,Z_iZ_j=0$ for $i\ne j$; of its terms of fourth order those in which some $Z_j$ stands to the first
power vanish, the trace of an element of a simple ideal being $0$ on every invariant subspace, and the cross terms
that remain are $-\frac1{4N}\sum_{i<j}\mathrm{tr}(Z_i^2Z_j^2)$, where
$\frac1N\mathrm{tr}(Z_i^2Z_j^2)=\tau_{ij}\langle Z_i,Z_i\rangle\langle Z_j,Z_j\rangle$, an invariant quadratic form
on a simple ideal being a multiple of $\langle\cdot,\cdot\rangle$. They vanish exactly when $Z_i^2Z_j^2=0$ for all
$i\ne j$, a product of two commuting positive matrices with trace $0$ being $0$, that is, when the $G_i$ act on
separate blocks. At $u=u_iv$, $v=\prod_{j\ne i}u_j$, the elements $u_i$ and $v$ commute, so that
$c(u_iv)=c(u_i)c(v)-H_{u_i}H_v$ and $H_{u_iv}=H_{u_i}c(v)+c(u_i)H_v$, and $c(v)$ and $H_v$ commute with $G_i$ and
with $\mathfrak g_i$. So the block $\frac1N\mathrm{Re}\,\mathrm{tr}(Sc(u)T)$ of $\mathfrak g_i$, the second
derivative along $S,T\in\mathfrak g_i$, is $\frac1N\mathrm{Re}\,\mathrm{tr}(c(v)\,Sc(u_i)T)$, the Hessian of the
simple pair of $\mathfrak g_i$ with the weight $c(v)/N$, less $\frac1N\mathrm{Re}\,\mathrm{tr}(H_vSH_{u_i}T)$,
whose modulus is at most $|S||T||H_{u_i}||H_v|\le|S||T||u_i-\mathbf 1||v-\mathbf 1|$, as $|H_w|\le|w-\mathbf 1|$
for unitary $w$; and the gradient on $\mathfrak g_i$ is $P_iH_u=P_i[c(v)H_{u_i}]+P_i[c(u_i)H_v]$. On a constituent
$W_s\otimes M_s$ the elements $S$, $T$ and $H_{u_i}$ act on $W_s$ and $H_v$ on $M_s$, with
$\mathrm{tr}_{M_s}H_v=\mathrm{Im}\,\chi_{M_s}(v)$, and $\mathrm{Re}\,\mathrm{tr}(SHT)=\mathrm{tr}(\frac12\{S,T\}H)$
for hermitian $S$, $H$, $T$; so the odd term is
$\frac1N\sum_s\mathrm{tr}_{W_s}(\frac12\{S,T\}H_{u_i})\,\mathrm{Im}\,\chi_{M_s}(v)$. An antiunitary $J$ that
commutes with the group on a self-conjugate side has $JSJ^{-1}=-S$ for the elements $S$ of its Lie algebra and
$JH_wJ^{-1}=-H_w$ for its elements $w$, and $\mathrm{tr}(JXJ^{-1})=\overline{\mathrm{tr}\,X}$; so the factor of
that side, a real number, equals its negative and vanishes. Every representation of $SU(2)$ is self-conjugate, and
at $SU(2)\times SU(3)$ on $\mathbb C^2\otimes\mathbb C^3$ the side of the other factor of the ideal of $SU(3)$ is
the $\mathbb C^2$ of $SU(2)$; at $SU(3)\times SU(3)$ on $\mathbb C^3\otimes\mathbb C^3$, with $S=T=D$,
$Z_i=\epsilon D$ and $v=e^{i\epsilon'D}$, $D=\mathrm{diag}(1,1,-2)$, the two factors are
$\epsilon\,\mathrm{tr}\,D^3+O(\epsilon^3)$ and $-\frac{\epsilon'^3}6\mathrm{tr}\,D^3+O(\epsilon'^5)$, with
$\mathrm{tr}\,D^3=-6$. Between two ideals the second derivative along $S\in\mathfrak g_i$, $S'\in\mathfrak g_j$ is
$\frac1N\mathrm{Re}\,\mathrm{tr}(SS'u)=\frac1N\mathrm{Re}\,\mathrm{tr}(SS'(c(u)-\mathbf 1))$, as
$\mathrm{tr}\,SS'=0$; on a constituent of the factor $i$, the factor $j$ and the rest, $\mathrm{tr}(SS'u)$ is
$\mathrm{tr}(S(u_i-\mathbf 1))\,\mathrm{tr}(S'(u_j-\mathbf 1))$ times the trace of the rest, $S$ and $S'$ having
trace $0$ there, which gives the bound; and the example is a computation. The block averages on small fields and
the exponential map are products over the factors, the mean of a family of products of commuting factors being the
product of their means and the logarithm of a product of small commuting factors the sum of their logarithms;
$j_G(Y)=\prod_ij_{G_i}(P_iY)$, the eigenvalues of $\mathrm{ad}\,Y$ being those of the $\mathrm{ad}\,P_iY$ on the
$\mathfrak g_i$; the gauge-fixing term is a trace $a(X)$, to which the expansion above applies at complex fields,
so that its quadratic part is a sum over the ideals; and the chart, $G^\pi$ with its indicator and $\nu$ are
compared with those of the factors by the bounds of (i).

(iii) The leftover is a sum over the ideals of the difference of their coefficient and the mean times their part
$\frac12\langle F_i,F_i\rangle$ of the Yang--Mills term. By the transferred statements of the step on its lists,
its potentials and their edge parts (Proposition~\ref{C-368-17}) it splits off as a named term of the potentials of
each level, as the $SU(2)$ construction splits off a term of the potentials, with the size
$\frac74\alpha_0^2\max_i|\mathsf{cf}_i-\mathsf{cf}|$ and its edge parts $8M^3$ times this; and
$\mathsf{cf}_i-\mathsf{cf}=\sum_m\frac{n_m}n(\mathsf{cf}_i-\mathsf{cf}_m)$ gives
$|\mathsf{cf}_i-\mathsf{cf}|\le2(1-\frac{n_i}n)\max_m|\mathsf{cf}_m|$, whence the second bound. Its terms of first
order in the fluctuation field are of third order in the field strength, and Cauchy's estimate bounds them, as
Schwarz's lemma bounds such terms in the proof of Proposition~\ref{C-368-17}, whether or not they vanish. In
dimension three the part of second order of an earlier level is bounded, as in the $SU(2)$ construction, and stays
in the potentials. At level $j$ the leftover is $L^j\sum_i(e^{(i)}_j-b_j)Q_i$; as $x^{(i)}_k-x_k=-\sum_{j\le
k}L^{-(k-j)}(e^{(i)}_j-b_j)$, the leftovers of the levels up to $k$ add up to $-L^k\sum_i(x^{(i)}_k-x_k)Q_i$, which
is bounded uniformly in $k$, $L^kQ_i$ being of the order $\alpha_0^2$ on a block of level $k$ and $|x^{(i)}_k-x_k|$
bounded by Lemma~\ref{C-368-21}.

(iv) The proof of the $SU(2)$ construction meets the colour through one theorem on a whole class of families of
terms, Theorem~\YMMref{X}{427}, so a term of the tied action may be restricted to an ideal, and the leftover of
(iii) is carried as a term; with property (T) the remaining steps are those of the $SU(2)$ construction, and
Lemma~\ref{C-368-21} with $\Gamma=\mathsf G_\sharp$ gives the distance of the couplings, as
$|e^{(i)}_k|=|\mathsf{cf}_i|\le(\kappa_i/\kappa_G)^2\mathsf G_\sharp\le\mathsf G_\sharp$
(Proposition~\ref{C-368-17}), $\kappa_i\le\kappa_G$ since $\mathfrak g_i\subset\mathfrak g$. A format with one
exponent, one integrand and one total, the leftover of (iii) being the carried term, is the format of the $SU(2)$
construction for one run.

(v) At this pair $1-a=(1-a^{(2)})(1-\lambda_1a^{(1)})$, $a^{(1)}$ and $a^{(2)}$ being the plaquette functions of
the two factors, so that $1-a$ is a product of two factors each of which reads one group; hence the gradient and
the Hessian of the action are the weighted ones of the $SU(2)$ construction, one weight on each ideal,
$\lambda_1(1-a^{(2)}_p)$ and $1-\lambda_1a^{(1)}_p$, plus the block $-\lambda_1(\nabla a^{(1)}\otimes\nabla
a^{(2)}+\nabla a^{(2)}\otimes\nabla a^{(1)})$, which annihilates infinitesimal gauge directions, the $a^{(i)}_p$ being gauge invariant.
The fluctuation operator of the first step is the direct sum $A^\sharp(p_1)\oplus A^\sharp(p_2)$ plus $C^{\mathrm
T}TC$, $T$ being this block: it is symmetric, of the range of one plaquette and of second order in the plaquettes,
each gradient being of first order in its plaquette variable, and it vanishes at $\mathbf 1$ and at fields with
values in one factor, where the gradient of the factor whose field is $\mathbf 1$ vanishes. If the regular-domain inverse hypotheses of (v) hold,
$(A_0+C^{\mathrm T}TC)^{-1}=\sum_{j\ge0}(-A_0^{-1}C^{\mathrm T}TC)^jA_0^{-1}$ in the stated norm. A pointwise second-order bound on $T$ does not establish this relative-smallness condition
after the maps $C$, scale changes and gauge fixing, nor does it supply the all-level regular-domain estimates. These remain analytic inputs for the additional assertions of (v).
\end{proof}

\begin{lemma}[Block sums]\label{C-368-29}
Let $G=G_1\times\dots\times G_r$ act on $\mathbb C^{N_1}\oplus\dots\oplus\mathbb C^{N_r}$ by blocks, $G_i\ne\{\mathbf 1\}$ closed.
\begin{enumerate}
\item[(i)] $a(u)=\sum_i\frac{N_i}Na_i(u_i)$ and the Haar measure is a product; so on every finite lattice the
Wilson measure at $\beta$ is the product of those of the $G_i$ at $\beta N_i/N$, the fields of the factors are
independent, and $\langle\mathcal O(f_1);\dots;\mathcal O(f_{\mathsf m})\rangle_\beta=\sum_i(\frac{N_i}N)^{\mathsf
m}\langle\mathcal O^{(i)}(f_1);\dots;\mathcal O^{(i)}(f_{\mathsf m})\rangle_{\beta N_i/N}$; with a weight
$\Lambda=\bigoplus\Lambda_i$, $\frac{N_i}N$ is replaced by $\mathrm{tr}\,\Lambda_i$.
\item[(ii)] Let $\nu$ be a Gibbs state of the product weight. Then $\nu(\prod_iF_i(U^{(i)}))=\nu(\prod_i\bar
F_i(U^{(i)}))$ for bounded Borel functions $F_i$, $\bar F_i$ any representative of the conditional expectation of
$F_i$ given the tail field $\mathcal T$ in the marginal $\nu_i$: the factors are independent given the tails. For a
Gibbs state $\nu_1$ of one factor and $\mathsf m\ge1$ there is exactly one probability on $\mathsf m$ copies whose
value at $F_1\otimes\dots\otimes F_{\mathsf m}$ is $\nu_1(\prod_a\mathbb E_{\nu_1}[F_a|\mathcal T])$; it is a Gibbs
state of the product weight of $\mathsf m$ copies at equal couplings, all its marginals are $\nu_1$, and it is
periodic if $\nu_1$ is. Let the marginals $\nu_i$ of $\nu$ be periodic for the factors $i$ of a set $I'$, and for
$i\in I'$ let $F^{(i)}_{\mathbf k}\in\mathcal A(\square^{(i)}_{\mathbf k})$, $\mathbf k$ in a finite set $\mathcal
K^{(i)}_0$, be functions of the cells of a board of $\mathbb L_K$, each factor with its own board and its own
sequence of sides $\ell^{(i)}_j\to\infty$. Then the chessboard domination of the $SU(2)$ construction
(Proposition~\YMMref{91.9}{705}) holds for their product:
\[\Bigl|\nu\Bigl(\prod_{i\in I'}\prod_{\mathbf k}F^{(i)}_{\mathbf k}(U^{(i)})\Bigr)\Bigr|\le\prod_{i\in
I'}\liminf_{j\to\infty}\prod_{\mathbf k}\bigl\|F^{(i)}_{\mathbf k}\bigr\|^{(i)}_{{\rm ch},\ell^{(i)}_j},\]
$\|\cdot\|^{(i)}_{{\rm ch},\ell}$ being the chessboard seminorm of the Wilson measure of $G_i$ at its coupling.
\item[(iii)] Let $G=SU(2)^r$ in $\mathbb C^2\oplus\dots\oplus\mathbb C^2$, $0<g_i\le\gamma$, $(K;(\beta_i))$ a
first-passage datum at $(g_i)$ with the labels $\mathsf b_i=a_K\beta_i$, $\nu$ a Gibbs state whose marginals
$\nu_i$, $\nu_j$ are periodic, $i\ne j$, $\vartheta\in(0,1)$, $a_K\le\mathsf d\le\frac1{20}$, and $(f,h)$ a test
pair of scale $\mathsf d$ and parameter $\vartheta$ as defined in
Section~\ref{C-350-1}. Then
\[|\langle\mathcal O^{(i)}(f);\mathcal O^{(j)}(h)\rangle_\nu|\le\frac14\varepsilon^\times_{ij}(\mathsf
d)\,\frac3{16\pi^2}(\mathsf b_i\mathsf b_j)^{-1}\mathcal K_6(f,h),\]
with $\varepsilon^\times_{ij}(\mathsf d)$ at most a number of $g_i$, $g_j$, $\vartheta$ and the constants times
$\mathsf d$, reading none of $K$, $\beta_i$, $\nu$, $f$, $h$. In every Gibbs state of the product weight whose
marginals are periodic in all the factors $i_1,\dots,i_{\mathsf m}$ of a mixed cumulant $\langle\mathcal
O^{(i_1)}(f_1);\dots;\mathcal O^{(i_{\mathsf m})}(f_{\mathsf m})\rangle$, $\mathsf m\ge2$, the $i_q$ not all equal,
this cumulant obeys the per-ball and factorial growth estimates obtained by the normalizer argument in the proof of Proposition~\YMMref{91.13}{709}, with constants depending on the factor ceilings and the radii of the balls containing the supports. No separation of the supports is required: the mixed derivative of every single-factor normalizing exponent is zero, as checked below.\qedhere
\end{enumerate}
\end{lemma}

\begin{proof}[Sketch of proof]
(i) is the product structure; cumulants of sums of independent families add, and $\mathcal O=\sum_i\frac{N_i}N\mathcal O^{(i)}$ gives the powers.

(ii) The Gibbs kernels of the product weight are products of those of the factors, and the marginal $\nu_i$ is a
Gibbs state of the factor $i$. For finite sets of bonds $B_1\subset B_2\subset\cdots$ that exhaust the lattice,
$\nu(\prod_iF_i)=\nu(\prod_i\gamma^{(i)}_{B_l}F_i)$, $\gamma^{(i)}_{B_l}$ the Gibbs kernel of the factor $i$;
$\gamma^{(i)}_{B_l}F_i$ is the conditional expectation of $F_i$ in $\nu_i$ given the bonds off $B_l$, these
$\sigma$-fields decrease to $\mathcal T$, and the downward martingale theorem gives the convergence to $\bar F_i$
off a $\nu_i$-null set; dominated convergence gives the first statement. The space of configurations being compact
and metrisable, $\nu_1$ has a regular conditional probability $\pi$ given $\mathcal T$, and the $\mathsf m$-copy
state is $\int\nu_1(d\omega)\,\pi(\omega,\cdot)^{\otimes\mathsf m}$: its values at products fix it; it obeys the
Gibbs equations because $\mathcal T$ lies in the $\sigma$-field of the bonds off every finite set; and a
translation that keeps $\nu_1$ keeps $\mathcal T$ and so keeps it. The domination is the proof of
Proposition~\YMMref{91.9}{705} run with the factors, first for a periodic Gibbs state of the product weight and
functions that are products over the factors: the comparison of the Gibbs kernel of a box with the torus measure
(Proposition~\YMMref{91.8}{704}) holds factor by factor, by Fubini, the exponents adding, and reads of a group only
$0\le a_p\le2$; the master inequality (the proof of Proposition~\YMMref{91.9}{705}) holds with the diagonal translations, the geometry
of boards (the proof of Proposition~\YMMref{91.9}{705}) reading no group; and after the expansion of the product of the translates each
term is a product over the factors of functions of one factor alone, whose mean under the product of the torus
measures is the product of the means, each bounded by the chessboard estimate of its factor
(Proposition~\ref{C-368-19}(ii)). Periodic marginals suffice: with $\bar F^{(i)}$ the conditional expectation given
$\mathcal T$ of the product $F^{(i)}$ of the functions of the factor $i$, the first statement and H\"older's
inequality give $|\nu(\prod_iF^{(i)})|\le\prod_{i\in I'}\nu_i(|\bar F^{(i)}|^{\mathsf m})^{1/\mathsf m}$, $\mathsf
m$ the least even integer at least $\#I'$; and $\nu_i(|\bar F^{(i)}|^{\mathsf m})$ is the value of the $\mathsf
m$-copy state of $\nu_i$, a periodic Gibbs state of the product weight of $\mathsf m$ copies, at the product of
$\frac{\mathsf m}2$ copies of $F^{(i)}$ and $\frac{\mathsf m}2$ copies of its conjugate, which the domination for
the product weight bounds, the conjugate having the seminorm of the function.

(iii) By (ii) the mixed covariance is the covariance of the conditional expectations given the tails, so that
$|\langle\mathcal O^{(i)}(f);\mathcal O^{(j)}(h)\rangle_\nu|^2\le\mathsf V_{\nu_i}(f)\,\mathsf V_{\nu_j}(h)$,
$\mathsf V_{\nu_i}(f)$ the variance of $\mathbb E_{\nu_i}[\mathcal O(f)|\mathcal T]$ in $\nu_i$. Let $\mathsf
D(f_1,f_2)$ be the covariance of $\mathbb E_{\nu_i}[\mathcal O(f_1)|\mathcal T]$ and $\mathbb E_{\nu_i}[\mathcal
O(f_2)|\mathcal T]$ in $\nu_i$, a real bilinear form; it is the mixed covariance of the two copies in the two-copy
state of $\nu_i$. \emph{The unit scale.} Let $F_1$, $F_2$ vanish off $\mathbb B(0,1)$. The systems of the $SU(2)$
construction give at the first-passage datum $(K,\beta_i)$ the normalisers (Nr1) of the one-entry tuples on the
board of side $L$ (equation~\YMMref{(1844)}{708}, Proposition~\YMMref{91.12}{708}\,(1)); they pass to the two-copy
state, on the product of the polydiscs of the two copies, by the domination of (ii), as in the transfer principle
(Proposition~\YMMref{91.11}{707}); at this mixed word the derivative of the sum of the two normalising exponents
vanishes, and Cauchy's estimate gives $|\mathsf D(F_1,F_2)|\le C'_V(g_i)\|F_1\|_\sharp\|F_2\|_\sharp$, with
$C'_V(g)$ a number of $g$ and the constants and $\|F\|_\sharp=\max\{\sup|F|,40M\,\mathrm{Lip}\,F\}$. \emph{The
factor of the radius.} $\nu_i$ is invariant under all translations of $a_K\mathbb Z^3$
(Proposition~\ref{C-368-19}(iii)), so $\mathsf D(\tau_vf_1,\tau_vf_2)=\mathsf D(f_1,f_2)$ for $v\in a_K\mathbb
Z^3$. Let $f_1$, $f_2$ vanish off a ball $\mathbb B(x_0,\mathsf r)$, $a_K\le\mathsf r\le\frac1{20}$. There is a
packing of at least $(0.3/\mathsf r)^3$ translates of this ball in $\mathbb B(0,1)$ by vectors of $a_K\mathbb Z^3$
(the packing argument in the proof of Proposition~\YMMref{127.1}{848}); the signed sums $\sum_v\epsilon_v\tau_vf_l$ over the packing, $\epsilon_v=\pm1$, have
the norms $\|f_l\|_\sharp$, and the mean over the signs of $\mathsf D$ at these sums is the sum over the packing of
the $\mathsf D(\tau_vf_1,\tau_vf_2)$, the cross terms cancelling (the packing argument in the proof of Proposition~\YMMref{127.1}{848}). So $|\mathsf
D(f_1,f_2)|\le C'_V(g_i)(\mathsf r/0.3)^3\|f_1\|_\sharp\|f_2\|_\sharp$, and $\mathsf V_{\nu_i}(f)\le
C'_V(g_i)(\mathsf r/0.3)^3\|f\|_\sharp^2$. \emph{The test pairs.} The two functions of a test pair of scale
$\mathsf d$ and parameter $\vartheta$ vanish off balls of radius $\mathsf d$, so $\mathsf r=\mathsf d$; as
$\vartheta<1$, $\mathsf d\le\frac1{20}$ and $M\ge1$, $\|f\|_\sharp\le\frac{40M}{\vartheta\mathsf d}\|f\|_\infty$,
and the same for $h$; $\|f\|_\infty\|h\|_\infty\le10^6\vartheta^{-2}\mathcal K_6(f,h)$ and $\mathsf b_i\le
L(g_i^{-2}+\bar B)$; this gives the bound, with $\varepsilon^\times_{ij}(\mathsf d)$ a number of $g_i$, $g_j$,
$\vartheta$ and the constants times $\mathsf d$. \emph{The higher cumulants.} For each factor of the word, the
graded normaliser that the systems give to the sub-tuple of its entries (equation~\eqref{eq:overlap-graded-normalizer}), with no
separation of the supports, passes to the state by the domination of (ii), on the product of the polydiscs of the
factors; at a mixed word the derivative of the sum of the exponents vanishes, each exponent reading the entries of
its own factor alone, so that Cauchy's estimate bounds the cumulant itself, and the per-ball bound and the growth
bound follow from the placement and tent-decomposition calculations in the proof of Proposition~\YMMref{91.13}{709}. The separation-dependent contribution there is the mixed derivative of the normalizing exponent; here that contribution vanishes because every exponent depends on the variables of only one factor. The remaining Cauchy and counting estimates require no separation.
\end{proof}

\begin{lemma}[Change of units and one-step normalization]\label{lem:cg-level-normalization}
\begin{enumerate}
\item[(i)] Remove the last level of a regional sequence, and denote its
length convention and determining set by $\eta^-=L\eta$ and
$\mathbf B^-$. For a real form $A^-$ put $A_0=LA^-$. Then
\[
 \eta A_0=\eta^-A^-,\qquad
 \|A^-\|_{X,\mathbf B^-}\le\|A_0\|_{X,\mathbf B}
                         \le L^2\|A^-\|_{X,\mathbf B^-}.
\]
This holds for every closed matrix group and every background.
\item[(ii)] Let $k\ge1$ and suppose all regions through level $k$ are the whole torus.
Let $U_0,U'$ be regular backgrounds with parameter $\alpha_0$, and
let $(A_0,g_0)$ be a real chart pair,
$U'^{g_0}=e^{i\eta A_0}U_0$, normalized on the preceding top:
$g_0=v_{k-1}[A_0]$ on the level-$k$ sites.
Assume the gauge-correction estimate of
Lemma~\YMMref{30.4}{242} and the slice estimate of
Lemma~\YMMref{30.10}{250} hold at these backgrounds, with their
admitted radii and constants, and assume the one-step factor bound
\[
 \sup_{\Lambda_k}|\mathsf v_{k-1}[A_0]-\mathbf1|
                 \le2\kappa_6L^{-1}\|A_0\|_{X,\mathbf B}.
\]
Here the bounds may be the transferred bounds when a transfer is
explicitly supplied. In particular this statement does not assert an
unconditional transfer to an arbitrary group.
If $\alpha_0$ and $\|A_0\|_{X,\mathbf B}$ are below the thresholds
specified in the proof, and
\[
 \sup_c|\mathcal M^kU'(c)\mathcal M^kU_0(c)^{-1}-\mathbf1|
       \le\gamma,
\]
there is a new chart pair $(A,g)$ with
\[
 \mathsf RD^*A=0,\qquad g=v_k[A]\text{ on }\Lambda_k,
 \qquad \sup_c|a_k[A](c)|\le\tfrac\pi2\gamma,
\]
and
\[
 |A|_{(-1)}\le k_1\|A_0\|_{X,\mathbf B},\qquad
 \|A\|_{X,\mathbf B}+|D^*DA|_{(-3)}
                  \le C_{\rm sl}(\alpha_0+\gamma).
\]
The constants and the thresholds depend only on those of the stated
inputs and on $L$.
\end{enumerate}
\end{lemma}

\begin{proof}
(i) If $j(b)$ is the level of a bare bond, its weighted field values
in the two conventions are
\[
 L^{j(b)-k}L|A^-(b)|,
 \qquad L^{\min\{j(b),k-1\}-(k-1)}|A^-(b)|.
\]
They agree unless $j(b)=k$, when their ratio is $L$. A covariant
difference gains two factors of $L$, one from the form and one from
the inverse bare spacing. Its weighted values therefore agree below
level $k$ and have ratio $L^2$ at level $k$. The background transports
are identical. Taking the maxima gives both inequalities. In
particular, on the constant sequence this computation is independent
of the sector.

(ii) Write $\alpha_3=\|A_0\|_{X,\mathbf B}$ and prescribe
$r=\mathsf v_{k-1}[A_0]$ on the top sites. Let $a_{\rm gc}$ be the
radius admitted by gauge correction and $a_{\rm sl}$ the field radius
admitted by the slice estimate, decreased to stay in the averaging
chart. Set
\[
 k_1=1+2C_{\rm gc}K_3(3+2\kappa_6L^{-1}).
\]
Require $\alpha_0$ below both input regularity ceilings and require
\[
 \alpha_3\le a_{\rm gc},\qquad
 2\kappa_6L^{-1}\alpha_3\le a_{\rm gc},\qquad
 k_1\alpha_3\le a_{\rm sl}.
\]
Gauge correction gives $w=e^{i\lambda}$ with
$\mathsf RD^*A_0^w=0$ and induced top factor $r$. Since
$|D^*A_0|_{(-2)}\le3\alpha_3$, its bounds give
\[
 \|\lambda\|_Y\le K_3(3+2\kappa_6L^{-1})\alpha_3,
 \qquad |A_0^w|_{(-1)}\le k_1\alpha_3.
\]
Put $A=A_0^w$ and $g=wg_0$.
The initial residual factor at a level-$k$ site is
$v_k[A_0]^{-1}g_0=\mathsf v_{k-1}[A_0]^{-1}$.
Under the correction it is multiplied on the left by the prescribed
$r$, so it becomes $\mathbf1$. Thus $g=v_k[A]$ on the top.

Every retained row on the constant sequence is an interior row at
level $k$. The normalized factorization of the actual average gives
\[
 e^{ia_k[A](c)}\mathcal M^kU_0(c)=\mathcal M^kU'(c).
\]
For a unitary matrix on the intrinsic branch,
$|\log V|\le(\pi/2)|V-\mathbf1|$, by the scalar spectral formula;
hence $\sup|a_k[A]|\le\pi\gamma/2$. Gauge invariance preserves
regularity of $e^{i\eta A}U_0=U'^g$. The slice estimate now applies
with row bound $\pi\gamma/2$, proving the two final bounds after
increasing $C_{\rm sl}$. For the $SU(2)$ constants of
Lemma~\YMMref{30.10}{250}, one may take
$C_{\rm sl}=20\pi K_1$. This proof uses only the stated local bounds
and group identities, so the same calculation applies whenever those
bounds have been supplied in another sector.
\end{proof}

\begin{proposition}[Quotients]\label{C-368-27}
Let $G=\tilde G/\mathsf Z$ with $\tilde G$ simply connected and $\mathsf Z$ finite and central, and $L\ge13$ odd.
For $g=e^{iY}\in G^\pi$, $Y\in\mathfrak g$, $|Y|<\pi$, the \emph{canonical lift} $s(g)\in\tilde G$ is $e^{i\tilde
Y}$, $\tilde Y$ the element of the Lie algebra of $\tilde G$ that the covering carries to $Y$; for a lift $\tilde
U$ to $\tilde G$ of a field $U$ whose plaquette variables lie in $G^\pi$, the \emph{central part} of $\tilde U$ at
a plaquette $p$ is $z_{\tilde U}(p)=\partial\tilde U(p)\,s(\partial U(p))^{-1}\in\mathsf Z$.
\begin{enumerate}
\item[(i)] The sector of Lemma~\ref{C-368-7}(iv) is defined, gauge invariant and locally constant on the fields
with small plaquettes, and kept by blocking. Let a separated sequence of regions have the whole
torus at the bottom. Write the finite sequence as
$\Omega_0,\ldots,\Omega_k$, with $\Omega_{-1}=\Omega_0$
and $\Omega_{k+1}=\varnothing$. In bare units its members $\Omega_q$ are unions
of cubes of side $\mathfrak M L^q$, with $\mathfrak M$ and $R_1$
powers of $L$, $R_1\ge L$, $R_1\mathfrak M\ge L^2$, and
$\operatorname{dist}_\infty(\Omega_q^c,\Omega_{q+1})>
R_1\mathfrak M L^q$. Let $U$ be a field with $|\partial U(p)-\mathbf 1|<\varepsilon L^{-2\mathsf j(p)}$ at every
bare plaquette, $\mathsf j(p)$ the level of the sequence at $p$, where $L^2\varepsilon\le c(G)$ of
Lemma~\ref{C-368-7}(iii), and let $V$ be its block datum. For a coordinate plane $P$ through a site of the top
lattice, the \emph{plane number} of $U$ at $P$ is the product of the central parts $z_{\tilde U}(p)$ over the bare
plaquettes $p$ of $P$, for a lift $\tilde U$ of $U$; the \emph{cells} $S$ of $P$ are defined as follows. Write $Q_r(p)$ for the square of side
$L^r$ in $P$ containing the bare plaquette $p$, with corners on
the level-$r$ lattice. Its cell is $Q_q(p)$ for the largest
$q\le k$ at which a corner lies in $\Omega_q$. These cells
tile $P$. A side of a level-$q$ cell having an end in
$\Omega_q$ is one determining bond of level $q$; a side with
both ends outside $\Omega_q$ consists of $L$ determining
bonds of level $q-1$. Call these bonds the elements of the
cell. Thus each cell has at most $4L$ elements, and
$\partial_SV$ is their ordered product, with the signs
of a walk around the boundary of $S$.
Then every $\partial_SV$ lies in $G^\pi$; for any lifts $\hat V$ of the values of $V$, $z_{\hat
V}(S):=\partial_S\hat V\,s(\partial_SV)^{-1}$ lies in $\mathsf Z$; and the plane number of $U$ at $P$ is
$\prod_Sz_{\hat V}(S)$ over the cells $S$ of $P$. So the plane numbers depend on $V$ alone, and two such fields
over one datum lie in one sector; over $V=e^{iB}$ with $4L\sup|B|<\pi$ the sector is trivial.
\item[(ii)] Let $G\subset U(N)$ be closed and let the averaging
map agree with the formula of Lemma~\ref{C-368-7}(i) on small
fields. Let the sequence have the whole torus at the bottom
and satisfy the following stronger separation conditions in
bare units: $\Omega_q$ is a union of cubes of side $L^{q+1}$,
and $\operatorname{dist}_\infty(\Omega_q^c,\Omega_{q+1})>
L^{q+2}$. Empty and full members are allowed.
Then a field $U$ with $|\partial U(p)-\mathbf 1|<\varepsilon
L^{-2\mathsf j(p)}$ at every bare plaquette, $\mathsf j(p)$ the level of the sequence at $p$, and block datum
$\mathbf 1$, $\varepsilon\le\varepsilon_2=\min\{\varepsilon_2(L),\frac1{106}r_G\}$, has a gauge transform, by a
gauge function equal to $\mathbf 1$ at the tops and within $C\varepsilon$ of $\mathbf 1$ at the nodes, of the form
$e^{ia}$ with $|a(b)|\le C\varepsilon L^{-\mathsf j(b)}$ and $\sum_bL^{-2\mathsf j(b)}|a(b)|^2\le C\sum_p|\partial
U(p)-\mathbf 1|^2$, $C$ a number of $L$; and a flat field with block datum $\mathbf 1$ is a pure gauge. So these
statements on the whole torus hold as worded; their hypothesis, block datum $\mathbf 1$, selects the trivial
sector.
\item[(iii)] Suppose that all regions are the whole torus. For every
closed $G\subset U(N)$, the preliminary gauge, the change of level
(Lemma~\ref{lem:cg-level-normalization}(i)) and the section of a datum
(Lemma~\YMMref{36.3}{274}) hold in every sector. More precisely,
two configurations in $\mathcal U'(\alpha_0)$ whose terminal data
differ by at most $\alpha_1$ admit a gauge equal to $\mathbf1$ at
the top sites that identifies their averaged contour variables and
obeys
\[
 \bigl|\mathcal M^n(U^{\hat g})(b)
             \mathcal M^nU_0(b)^{-1}-\mathbf1\bigr|
 \le C_{\rm pg}(\alpha_0+\alpha_1)
\]
at every level. Here $C_{\rm pg}$ depends only on $L$, and the
ceiling on $\alpha_0$ is $c(G)$ of Lemma~\ref{C-368-7}(iii).
\item[(iv)] Let $G=SO(3)\subset U(3)$, all regions being the whole torus, and let $U_0$ be a background of the regular class $\mathfrak U'_k(\mathbf B,\alpha_0)$ of the $SU(2)$ construction, in any sector. There are two numbers of $L$, a ceiling on $\alpha_0$ and a floor $6LM$ on the number of sites of the top lattice along each axis, $M\ge M_{\rm ps}$ a power of $L$, in the place of the floor $6LM_{\rm ps}$ of the $SU(2)$ construction (Theorem~\YMMref{I}{225}), such that under them the operator $\Delta^{\rm W}_{U_0}+D\mathsf RD^*+\mathsf Q^*\mathsf Q$ of the $SU(2)$ construction is formed on the whole torus and is positive definite, and its inverse $\mathsf G_0$, the operator $\mathsf H_0=\mathsf G_0\mathsf Q^*(\mathsf Q\mathsf G_0\mathsf Q^*)^{-1}$ and $\mathsf G_0-\mathsf H_0\mathsf Q\mathsf G_0$ obey the bounds of the $SU(2)$ construction with numbers of $L$ in the places of its constants.

The clauses of the $SU(2)$ construction that read the background on pieces of bounded size hold for $SO(3)$ in every sector. A bound of such a clause that is of degree $d$ in sizes of forms and of degree $e$ in chords of values in the group, chords at most $\alpha$, $\alpha\in(0,1]$ the parameter of the regular class $\mathfrak U'_k(\mathbf B,\alpha)$ of the background, keeps its number $C$ of the $SU(2)$ construction where $d+e\ge2$, where $d=1$ and $e=0$ at real forms, and, for a linear map whose kernel is real, in bounds on the norms of its kernel and on their sums along rows, and takes $(1+0.04\alpha)C$ where $d=0$ and $e=1$, as the constant of the gauge on a cube in the improvement of the regularity radius does, at most $\sqrt2C$ in the other bounds with $d=1$ and $e=0$ at complex forms, and $2C$ where $d=e=0$, while a bound on the norm from suprema to suprema of an operator on complex forms takes a factor at most $\sqrt2$.

The constant $C_{4,\mathrm{ec}}=808$ of the $SU(2)$ normalization in the bounds of the action of Proposition~\YMMref{30.1}{236} gives way to an absolute number, which serves at real and at complex forms and at every scale of that clause, a scale being a number $\xi\ge\eta$ at which the plaquette variables at the bond meet the condition of the clause. The comparison of the two regular classes at complex fields holds with $r+5r^2$ in the place of $r+4r^2$ at the fields $e^{i\eta A}U_0$ with $\eta\sup|A|\le\frac18$.

For an averaging map that is covariant at every field and the formula of Lemma~\ref{C-368-7}(i) only below a smaller radius, the same holds with three ceilings, each the least of the number of the $SU(2)$ construction and a fixed multiple of that radius $\varpi$: on $\alpha_0$, the ceiling $L^{-2}c^{\rm g}_2$ of the global chart with $\frac14L^{-2}\varpi$; on the size $|A|_{(-1)}$ of the potential $A$ of a chart pair $(A,g)$ of the global chart, the ceiling $L^{-1}c_4$, $c_4$ the number of the local chart estimates, with $\frac1{32}L^{-1}\varpi$; and on the radius of the one-step chart maps, the maps $a\mapsto\xi_\star(W_0,a)$, $\frac1i\log\mathsf v_\star(W_0,a)$ and $\mathcal Q^1_\star(W_0,a)$ of the one-step chart, the chart radius $r_L$ with $\frac18\varpi$. The backgrounds of these maps then lie within $\frac12\varpi$ of $\mathbf 1$, and the fields of a step within $\varpi$, in the way in which Lemma~\ref{C-368-7}(v) keeps them within $\min\{\omega_M,r_G\}$.

With these ceilings, the gauge correction (Lemma~\YMMref{30.4}{242}) holds in every sector with constants $K_3$ and $C_{\rm gc}$ that are numbers of $L$; and one step (Lemma~\ref{lem:cg-level-normalization}(ii)) and one step of the ascent (Proposition~\YMMref{37.2}{282}) hold in every sector, with constants of $L$ in the places of the numerals of the $SU(2)$ construction, the improvement of the regularity radius (Proposition~\YMMref{37.1}{278}) holding with free constants. In particular, over every datum whose plaquette variables lie within $\varepsilon'$ of $\mathbf 1$ and which lies within $\varepsilon'$ of the averages of a configuration of the regular class one level down, $\varepsilon'$ below a number of $L$, there is a critical configuration of the regular class one level up. The sector does enter objects on the whole torus: at a flat $SO(3)$-field the covariant difference on $\mathfrak g$-valued site functions has a non-zero kernel in the trivial sector and none in a twisted one.
\end{enumerate}
\end{proposition}

\begin{proof}[Sketch of proof]
(i) The squares $Q_r(p)$ are nested, so the largest-level rule
assigns the same cell to every bare plaquette of that cell.
A level-$q$ cell has a corner in $\Omega_q$ and no corner in
$\Omega_{q+1}$. The block structure then puts no level-$q$
site of the cell in $\Omega_{q+1}$. Separation places the
cell and the double blocks of its boundary elements in
$\Omega_{q-1}$: their distance from that corner is at most
$\frac32L^q<L^{q+1}$. On a side with no end in $\Omega_q$,
the level-$(q-1)$ sites are outside $\Omega_q$, so the
$L$ subdivisions are indeed determining bonds.

An element of level $r$ has an end in $\Omega_r$ and neither
end in $\Omega_{r+1}$. The cell across it has level at least
$r$ and at most $r+1$: a cell of level $q'\ge r+2$ would
put an endpoint within $L^{q'}$ of $\Omega_{q'}$, hence in
$\Omega_{q'-1}\subset\Omega_{r+1}$, a contradiction.
At level $r$ the element is a whole side; at level $r+1$
it is one of the $L$ subdivisions of a side. Consequently
each element belongs to exactly two cell boundaries, with
opposite orientations.

If $\sum_j|X_j|<\pi$ in the norm of $\mathbb C^N$, $X_j\in\mathfrak g$, then $\prod_je^{i\tilde X_j}$ is the
canonical lift of $\prod_je^{iX_j}$, and a product of small exponentials that lies in $\mathsf Z$ is $\mathbf 1$.
By Lemma~\ref{C-368-7}(iii) in its local form, under $L^2\varepsilon\le c(G)$, the value of the datum at a bond is
the transport of $U$ along the straight bare path of the bond times a product of exponentials of elements of
$\mathfrak g$ of small total size; the paths of the sides of a cell $S$ compose its boundary, so that $\partial_SV$
is the transport of $U$ around $S$, a product of conjugated plaquette variables of $U$ in $S$, times such a
product, of total size below $\pi$. So $\partial_SV\in G^\pi$, and, the values of $V$ being lifted along the lifted
paths, $\partial_S\hat V=\prod_{p\subset S}z_{\tilde U}(p)\cdot s(\partial_SV)$. Another choice of the lifts of the
values of $V$ changes the $z_{\hat V}(S)$ by central factors that cancel in the product over the cells of $P$, each
bond of the datum in $P$ bounding two cells with opposite orientations; and the product over the cells of the
central parts of their plaquettes is the plane number. The plane numbers, with the cube numbers, which are $\mathbf
1$ for small plaquettes, fix the sector. Over $V=e^{iB}$ the product around a cell has at most $4L$ factors $e^{\pm
iB(e)}$, and their lifts $e^{\pm i\tilde B(e)}$ give $z_{\hat V}(S)=\mathbf 1$.

(ii) We construct the gauge by descending through the layers. The
interpolation used in the descent is needed both for the pointwise
bound and for the energy estimate; setting every contour variable
equal to the identity would not give the latter estimate.

Write $n$ for the last level, put $\Omega_{-1}=T_\ell$ and
$\Omega_{n+1}=\emptyset$, and set
\[
 j(x)=\max\{q:x\in\Omega_q\},\qquad
 j(e)=\max_{x\text{ a corner of }e}j(x).
\]
Here $e$ may be a bond or a plaquette of any level. Let
$\Lambda_q=\mathbb T^{(q)}\cap(\Omega_q\setminus\Omega_{q+1})$;
these are the tops. The datum prescribes the averages on the
level-$q$ bonds with an endpoint in $\Lambda_q$, with any of the
inward bonds omitted as permitted by the given bond set. In
particular it prescribes all such bonds with no endpoint in
$\Omega_{q+1}$. A node is an endpoint of a retained bond.
All bond and plaquette sums below use positive orientations.

We first record the geometric consequences of the stated separation.
The blocks $B^q(y)$, $y\in\Lambda_q$, partition the bare sites,
and $j=q$ throughout each such block. Indeed, membership in
$\Omega_q$, a union of cubes of side $L^{q+1}$, is constant on
each level-$q$ block; assign a site to the block of the largest
level whose region contains it. Furthermore,
\begin{equation}\label{eq:layer-gauge-nearby}
 |x-x'|_\infty\le L^{j(x)+1}
 \quad\Longrightarrow\quad j(x')\ge j(x)-1.
\end{equation}
For $j(x)\ge1$ this follows from the distance between
$\Omega_{j(x)}$ and $\Omega_{j(x)-1}^c$; for $j(x)=0$ it is
automatic. A node which is not a top is the inner endpoint of
a retained inward bond whose other endpoint is a top.

For level $i$, let $\mathcal B_i$ and $\mathcal P_i$ be the
bonds and plaquettes with all corners in $\Omega_{i-1}$, and put
\[
 V_i=\mathcal M^iU|_{\mathcal B_i},\qquad
 u_i(p)=|\partial V_i(p)-\mathbf1|.
\]
The bare cone $Y_i(p)$ of a plaquette is the union of the
level-$i$ blocks of its four corners. The double block beneath
a bond of $\mathcal B_{i+1}$ lies in $\Omega_i$, so each
averaging step below is local to the indicated sets.

\emph{The two estimates for an averaging step.}
We use constants depending only on $L$. If $v$ is the largest
plaquette distance on the four blocks of a coarse plaquette $P$,
then, for sufficiently small $v$,
\begin{equation}\label{eq:layer-gauge-curvature-step}
 u_{i+1}(P)\le (\mathsf W_i u_i)(P)+C_0v^2,
 \qquad
 (\mathsf W_i f)(P)=L^{-3}
       \sum_{x\in B(P_-)}\sum_{p\in S_x(P)}f(p).
\end{equation}
Here $P_-$ is the base point of $P$, and $S_x(P)$ is the
$L\times L$ square of fine plaquettes parallel to $P$ with
base point $x$. The kernel is nonnegative and has row sum
$L^2$ and column sum $L^{-1}$; consequently
$\|\mathsf W_i\|_{2\to2}\le L^{1/2}$.

For completeness, the matrix calculation behind this estimate is
the one-step calculation of Lemma~\YMMref{14.9}{68}.
Put the field in the comb gauge on the four blocks. Each resulting
bond is a product of at most $C_L$ conjugated plaquette variables,
so its logarithm has size at most $C_Lv$. Differentiating the
mean formula at the identity gives
\[
 (Qa)(c)=L^{-3}\sum_{x\in B(c_-)}
 \bigl[\bar a(c_-,x)+a([x,x'])-\bar a(c_+,x')\bigr],
 \qquad x'=x+c_+-c_- ,
\]
where $\bar a(y,x)$ is the arithmetic mean of the line sums
along the coordinate-order contours. On going around $P$, the
contour terms at its four corners cancel. The remaining sum is
$L^{-3}\sum_x\sum_{p\in S_x(P)} da(p)$. The product-exponential
estimate and the quadratic Taylor remainder of the mean give
$C_Lv^2$ for the discarded terms. This proves
\eqref{eq:layer-gauge-curvature-step}. All estimates are matrix
norm estimates. Each logarithm is intrinsic to $G$: the comb
products and the contour quotients are products of small
exponentials in $i\mathfrak g$, and Lemmas~\ref{C-368-6}(ii),
(iv) and~\ref{C-368-7}(i) apply. Thus the argument uses no
$SU(2)$ identity. Counting $L^3$ translated squares, each with
$L^2$ plaquettes, gives the row sum; a fixed fine plaquette
has total multiplicity $L^2$ in these squares as the coarse
plaquette varies, giving the column sum $L^{-1}$.

These facts imply the two multiscale bounds
\begin{align}
 u_i(p)&\le2\varepsilon L^{2(i-k)}
       &&\text{if }Y_i(p)\subset\Omega_k,
                    \quad k\ge\max\{i-1,0\},\label{eq:layer-gauge-curvature-sup}\\
 \sum_{p\in\mathcal P}u_i(p)^2
       &\le2L^i\sum_{p_0\subset\bigcup_{p\in\mathcal P}Y_i(p)}
                         u_0(p_0)^2
       &&(\mathcal P\subset\mathcal P_i).
       \label{eq:layer-gauge-curvature-energy}
\end{align}
The bare plaquettes on the right are counted once. To see the
first bound, take maxima $P_r$ on the nested cones beneath $p$.
Equation~\eqref{eq:layer-gauge-curvature-step} gives
$P_{r+1}\le L^2P_r+C_0P_r^2$. Write
$P_r=\varepsilon L^{2(r-k)}p_r$. Starting with $p_0\le1$,
the induction $p_r\le2$ closes because
\[
 p_r\le\exp\left(2C_0\varepsilon L^{-2}
                         \sum_{t<r}L^{2(t-k)}\right)<2
\]
when $\varepsilon$ is below a number of $L$.
For the second bound, restrict at each level to the cones of
$\mathcal P$. A fine plaquette belongs to the four-block set
of at most twelve coarse plaquettes. Hence the square root of
the plaquette energy increases at one step by at most
\[
 L^{1/2}\bigl(1+4C_0L^{-1/2}v_r\bigr),
 \qquad v_r\le2\varepsilon L^{2(r-i+1)}\quad(r<i).
\]
Since $\sum_{r<i}v_r<4\varepsilon$, the product of the
parentheses is at most $\sqrt2$, after decreasing the same
ceiling on $\varepsilon$. This proves
\eqref{eq:layer-gauge-curvature-energy}.

\emph{Interpolation and its error.}
Measure a fine spacing as one, so the coarse lattice is
$L\mathbb Z^3$. For $b=\langle x,x+e_\mu\rangle$, write
$x=LY+r$, $0\le r_\nu<L$, and put
\[
 (I\alpha)(b)=L^{-1}\sum_{s_\mu=0,\ s_\nu\in\{0,1\}}
       \prod_{\nu\ne\mu}
        (1-r_\nu/L)^{1-s_\nu}(r_\nu/L)^{s_\nu}
       \alpha\bigl(\langle L(Y+s),L(Y+s)+Le_\mu\rangle\bigr).
\]
This is periodic and independent of the chosen coarse origin.
Its weights are nonnegative and sum to one. In particular
\[
 |I\alpha|\le I|\alpha|\le L^{-1}\sup|\alpha|,
 \qquad \|I\alpha\|_2\le L^{1/2}\|\alpha\|_2.
\]
For the second inequality apply convexity of the square: the
total weight of the fine bonds reading any fixed coarse bond
is $L^3$, which cancels two powers of the prefactor $L^{-1}$.
The line sum of $I\alpha$ along a straight coarse bond is
$\alpha(c)$. Directly taking transverse differences also gives
$dI\alpha=L^{-2}$ times a convex combination of the two adjacent
parallel values of $d\alpha$.

Suppose $s=\sup|\alpha|\le r_*/2$ on the coarse bonds within
one coarse spacing of a double block, and $v$ bounds the
plaquette distances of $e^{i\alpha}$ there. Choose
$r_*=r_*(L)>0$ small enough that the logarithm of the mean at
the identity is holomorphic and bounded by one on the complex
ball of radius $r_*$, and $r_*\le1/(200L)$. Such a choice
follows from the last, trivial-background estimate of
Lemma~\ref{C-368-7}(v). For $\beta=I\alpha$, the preceding
identities and the elementary product estimates give
\begin{equation}\label{eq:layer-gauge-interpolation}
 \begin{split}
 \sup|\beta|&\le L^{-1}s,\\
 \varpi_c(e^{i\beta})&\le L^{-2}(v+51s^2),\\
 |\mathcal M(e^{i\beta})(c)e^{-i\alpha(c)}-\mathbf1|
       &\le v+(35+2r_*^{-2})s^2.
 \end{split}
\end{equation}
Here is the error calculation. On a coarse plaquette,
$|d\alpha|\le v+35s^2$ by expanding its four exponentials.
The curl identity for $I$ and the same expansion on a fine
plaquette add at most $16L^{-2}s^2$, giving the second line.
The logarithm of the mean differs from $Q\beta$ by at most
$2r_*^{-2}s^2$, by Cauchy's estimate. Moving each of the $L$
straight steps past a coordinate-order contour uses at most
$L(L-1)$ plaquettes. Thus
$|Q\beta(c)-\alpha(c)|\le v+35s^2$, proving the last line.

We also use the following rigidity estimate on a double block.
If $W,W_0$ have identical averaged contour variables on its
two blocks, then
\[
 |W(b)W_0(b)^{-1}-\mathbf1|
 \le|\mathcal M(W)(c)\mathcal M(W_0)(c)^{-1}-\mathbf1|
         +C_L\max\{\varpi_c(W),\varpi_c(W_0)\}.
\]
Inside either block, choose a coordinate-order path whose last
step is the bond. The two quotients of its transports by the
endpoint contour means have size at most $C_L\varpi$, by the
path comparison in Lemma~\ref{C-368-7}(i); applying this to
both fields gives the bound. For a bond between the blocks, insert
the two endpoint contour means; move its straight crossing
to the central coarse path and then replace that path by its
mean. Each replacement costs $C_L$ times the largest
plaquette distance, and the endpoint means of the two fields
cancel. This is precisely the double-block reconstruction
calculation of Lemma~\YMMref{44.3}{347}. It proves the displayed
estimate with a constant of $L$ for the present $G$ as well.
In addition, the quotients in this calculation are products
of conjugated small plaquette exponentials and the mean
corrections. Their intrinsic logarithms therefore lie in
$\mathfrak g$ when the total exponent size is less than $\pi$,
by Lemma~\ref{C-368-6}(ii). For a closed group which is not
connected these products lie in its identity component; the
unrestricted transports themselves need not do so.

\emph{The descending gauge.}
Set $g=\mathbf1$ at the tops. The other bare sites form the
disjoint sets
\[
 N_i=(\mathbb T^{(i)}\setminus\mathbb T^{(i+1)})
           \cap\Omega_{i+1},\qquad i=n-1,\ldots,0.
\]
When $g$ is known on $\mathbb T^{(i+1)}$, set
\[
 W_{i+1}=V_{i+1}^{,g},\quad
 a_{i+1}=\frac1{\mathrm i}\log W_{i+1},\quad
 \beta_i=Ia_{i+1},\quad W_{0,i}=e^{i\beta_i},
\]
and, for $x\in B(y)\setminus\{y\}$,
$y\in\mathbb T^{(i+1)}\cap\Omega_{i+1}$, define
\[
 g(x)=W_{0,i}(y,x)^{-1}g(y)V_i(y,x).
\]
The expressions $W(y,x)$ are averaged contour variables. This
choice makes $W_i=V_i^g$ and $W_{0,i}$ have the same contour
variables on each such block; gauge covariance gives
$\mathcal M(W_i)=W_{i+1}$.

A bond of $\mathcal B_i$ with no endpoint in $\Omega_{i+1}$
has both endpoints at tops, and $W_i=\mathbf1$ there. If one
endpoint lies in $\Omega_i$, this follows from the trivial
datum. Otherwise its double block consists of bonds of the
level-$i-1$ datum, all equal to $\mathbf1$, so its average is
also $\mathbf1$. In particular $W_n=\mathbf1$ and $a_n=0$.
This observation also treats the exterior block of a double
block crossing the boundary of $\Omega_{i+1}$: both its
contour variables and the interpolated reference are trivial.

Here are the estimates that justify every step of the construction.
For a bond $b$ of level $i$ with an endpoint in $\Omega_{i+1}$,
choose a coarse double block containing it, with one centre
$y\in\Omega_{i+1}$, and put $k=j(y)-1\ge i$. Enlarge the
bare double block by $2L^{i+1}$. Equation~\eqref{eq:layer-gauge-nearby}
shows that $k\le j\le k+2$ throughout this enlargement.
Let $v_i(b)$ be the maximum of $u_i$ on the double block and
of $u_{i+1}$ within one coarse spacing of it, and let $s_{i+1}(b)$
be the maximum of $|a_{i+1}|$ on the coarse bonds in that
neighbourhood. Inductively assume
\[
 |a_q(c)|\le C_A\varepsilon L^{q-j(c)}\quad(q>i).
\]
Equation~\eqref{eq:layer-gauge-curvature-sup} gives
\[
 v_i(b)\le2\varepsilon L^{2(i+1-k)},\qquad
 s_{i+1}(b)\le C_A\varepsilon L^{i+1-k},\qquad
 s_{i+1}(b)^2\le\varepsilon L^{2(i+1-k)}
\]
if $C_A^2\varepsilon\le1$. Hence the reference plaquettes
are at distance at most $53\varepsilon$ from $\mathbf1$,
by \eqref{eq:layer-gauge-interpolation}.
Rigidity and that equation give
\begin{equation}\label{eq:layer-gauge-defect}
 W_i(b)=e^{iz_i(b)}e^{i\beta_i(b)},\qquad
 |z_i(b)|\le C_z\bigl(v_i(b)+s_{i+1}(b)^2\bigr).
\end{equation}
The constant $C_z=C_z(L)\ge1$ may be enlarged to control the
total size of the small exponentials just described. Thus the
logarithm defining $z_i$ is in $\mathfrak g$, independently
of the size of $r_G$. On bonds with no endpoint in
$\Omega_{i+1}$ put $z_i=v_i=s_{i+1}=0$.

Choose
\[
 N_2=3(8L)^3,\qquad
 C_A\ge\max\{4C_zL^8,375C_z^2N_2L\},
\]
and choose $\varepsilon_2(L)>0$ below the ceilings used in
\eqref{eq:layer-gauge-curvature-sup}--\eqref{eq:layer-gauge-curvature-energy}
and such that
\[
 C_A^2\varepsilon_2(L)\le1,\qquad
 C_AL\varepsilon_2(L)\le r_*/2,\qquad
 53\varepsilon_2(L)\le\omega_M/2,\qquad
 2C_zC_AN_2^{1/2}L^{3/2}\varepsilon_2(L)\le\log2.
\]
We may decrease it further by constants depending only on $L$
to keep the contour means in their stated domains. The
additional ceiling $\varepsilon\le r_G/106$ gives
$53\varepsilon\le r_G/2$. In particular every interpolated
reference has the small plaquettes required by
Lemma~\ref{C-368-7}(i). The actual intermediate fields do as
well: $C_A\ge4L^8$ implies
$\varepsilon_2(L)\le L^{-16}/16<3.9\cdot10^{-9}L^{-2}$,
so $\varepsilon\le c(G)$. The local version of
Lemma~\ref{C-368-7}(iii) on each required double block gives
intermediate plaquette distances at most $2\varepsilon$.
Thus each chosen averaging map agrees with the small-field
formula throughout this construction.

The small exponential product in \eqref{eq:layer-gauge-defect}
and Lemma~\ref{C-368-6}(ii) give
$|a_i|\le I|a_{i+1}|+|z_i|$. Iterating, using positivity of
$I$ and $a_n=0$, gives
\[
 |a_i|\le\sum_{r\ge i}I^{r-i}|z_r|.
\]
For a fixed bond $b$ of level $i$, put $J=j(b)$. A nonzero
term at level $r$ can influence $b$ only if $r\le J$;
its endpoints lie within $2L^r$ of those of $b$, and its
layer index is at least $J-1$, by
\eqref{eq:layer-gauge-nearby}. The local bound above implies
$|z_r(b')|\le3C_zL^6\varepsilon L^{2r-2j(b')}$.
Since each interpolation costs $L^{-1}$ in supremum norm,
\[
 |a_i(b)|\le3C_zL^8\varepsilon L^{i-2J}
                    \sum_{r=i}^JL^r
 \le4C_zL^8\varepsilon L^{i-J}
 \le C_A\varepsilon L^{i-J}.
\]
This closes the induction and justifies all logarithms in it.
At $i=0$ it proves the claimed pointwise bound. At a node
which is not a top, the retained inward bond of level $q$
joining it to a top has $V_q=\mathbf1$. Its transformed value
is $g(z)$ or $g(z)^{-1}$, and the same bound gives
$|g(z)-\mathbf1|\le C_A\varepsilon$.

\emph{The energy estimate.}
Keep the term quadratic in the coarse potential in
\eqref{eq:layer-gauge-defect}. With all functions extended by
zero off their indicated bond sets, it gives
\[
 |a_r|\le T_r|a_{r+1}|+\zeta_r,\qquad
 \zeta_r=C_zv_r,\qquad T_r=I+B_r,
\]
where, on an active bond $b$,
\[
 (B_rf)(b)=C_zC_A\varepsilon L^{r+3-j(b)}
                    \max_{\mathcal N_b}f.
\]
The set $\mathcal N_b$ is the coarse neighbourhood used in
the definition of $s_{r+1}(b)$; on an inactive bond $B_rf=0$.
The maps $T_r$ are monotone and subadditive on nonnegative
functions. Split $\zeta_r$ according to its layer:
$\zeta_r^J=\mathbf1_{\{j=J\}}\zeta_r$, which is zero
unless $r<J$, and put
\[
 f^J=\sum_{r<J}T_0\cdots T_{r-1}\zeta_r^J.
 \qquad\text{Then}\qquad |a|\le\sum_Jf^J.
\]
Each step spreads support by at most $3L^{r+1}$. Therefore
all the descendants of $\zeta_r^J$ are within $4L^{J-1}$
of a site of layer $J$, and their layer indices lie in
$[J-1,J+1]$, by \eqref{eq:layer-gauge-nearby}.

A coarse bond belongs to $\mathcal N_b$ for at most $N_2$
fine bonds: their base points lie in a cube of side
$6L^{r+1}$, with three bond directions at each site. Thus,
on such a layer-supported descendant,
\[
 \|T_tf\|_2\le
 \bigl(L^{1/2}+C_zC_AN_2^{1/2}\varepsilon L^{t+4-J}\bigr)
 \|f\|_2.
\]
Since $t<r<J$, the sum of $L^{t+7/2-J}$ over $t<r$ is
less than $2L^{3/2}$. The chosen smallness ceiling therefore
gives
\[
 \|T_0\cdots T_{r-1}\zeta_r^J\|_2
          \le2L^{r/2}\|\zeta_r^J\|_2.
\]
Let
\[
 E^J=\sum_{p:\,|j(p)-J|\le2}u_0(p)^2.
\]
Every plaquette entering $v_r(b)$ for $j(b)=J$ has its
bare cone in these layers. The same multiplicity bound
$N_2$, now for plaquettes, and
\eqref{eq:layer-gauge-curvature-energy} give
\[
 \|\zeta_r^J\|_2^2
       \le4C_z^2N_2L^{r+1}E^J,
 \qquad \sum_JE^J\le5\sum_pu_0(p)^2.
\]
Consequently
\[
 \|f^J\|_2\le
 4C_zN_2^{1/2}L^{1/2}(E^J)^{1/2}\sum_{r<J}L^r
 \le5C_zN_2^{1/2}L^{J-1/2}(E^J)^{1/2}.
\]
At a bare bond at most three $f^J$ are nonzero, those with
$|J-j(b)|\le1$. Cauchy--Schwarz now yields
\[
 \begin{split}
 \sum_bL^{-2j(b)}|a(b)|^2
 &\le3\sum_JL^{-2(J-1)}\|f^J\|_2^2\\
 &\le75C_z^2N_2L\sum_JE^J
 \le375C_z^2N_2L\sum_p|\partial U(p)-\mathbf1|^2.
 \end{split}
\]
This is the asserted energy bound, with $C=C_A$.

Finally suppose the field is flat. In each pillar let $t(x)$
be transport from its top to $x$. Flatness makes this independent
of the path inside the pillar. On a determining bond the straight
transport equals its block average, hence is $\mathbf1$.
For a bare bond within one pillar, its transport and the two
pillar paths form a contractible loop. For a bond joining
pillars of equal level, close the loop by their determining
bond. For adjacent pillars of levels $q$ and $q+1$, first
pass through the level-$(q+1)$ double block, using its outer
determining bond and the level-$q$ determining bonds in the
outer block. All closing transports are $\mathbf1$; the
separation and block nesting place the loop inside that
double block. These three cases give
$t(b_-)U(b)t(b_+)^{-1}=\mathbf1$ on every bare bond. The
same closing path at a retained inward bond gives $t=\mathbf1$
at its inner node, as at every top. Thus $U=\mathbf1^{t^{-1}}$,
with the residual gauge equal to $\mathbf1$ at all nodes.
In the quotient case the constructed potential also satisfies
$4\sup|a|\le4C_A\varepsilon\le4/C_A<\pi$; lifting its
small exponentials therefore has trivial central plaquette
parts. No sector condition was assumed in the construction;
the trivial datum enforces it.

(iii) We give the preliminary gauge argument, including its estimate.
Write $W_n=\mathcal M^nU$ and $W_{0,n}=\mathcal M^nU_0$.
The local straightening estimate of Lemma~\ref{C-368-7}(iii)
keeps every mean used below in its small-field domain. Its contour
means lie in $G$ and are gauge covariant, by
Lemmas~\ref{C-368-6}(iv) and~\ref{C-368-7}(i).
Set $g_k=\mathbf1$ on the top lattice and descend through the
blocks. If $x\in B(y)$ at level $n$, set
\[
 g_n(x)=W_{0,n}(y,x)^{-1}g_{n+1}(y)W_n(y,x),
\]
where $W_n(y,x)$ is the averaged contour variable from $y$ to $x$.
It equals $\mathbf1$ at $x=y$, so $g_n$ extends $g_{n+1}$.
Thus $\hat g=g_0$ is a gauge, equal to $\mathbf1$ at the top
sites. Gauge covariance gives
\[
 (\mathcal M^nU^{\hat g})(y,x)=W_{0,n}(y,x).
\]
These identities also determine $\hat g$ uniquely, by downward
induction.

For the quantitative assertion, put
\[
 \delta_n(b)=|\mathcal M^n(U^{\hat g})(b)W_{0,n}(b)^{-1}
                  -\mathbf1|.
\]
At the terminal level $\delta_k\le\alpha_1$. Within a single
block, the equality of the contour variables and the contour
straightening bound give $\delta_n\le C_L\alpha_0$. Across
two adjacent blocks, with coarse bond $c$, the same comparison
gives
\[
 \delta_n(b)\le\delta_{n+1}(c)
             +C_L\alpha_0L^{2(n-k)}.
\]
Indeed, a path product changes by at most the sum of the changes
of its factors, since the other factors are unitary. Replacing
the paths by their means costs at most a constant of $L$ times
the largest plaquette discrepancy in the two blocks; the two
end-point contour discrepancies vanish by the identities above.
The straightening estimate bounds this plaquette discrepancy by
$2\alpha_0L^{2(n-k)}$. Iterating until the two end points
enter one block or reach the top and summing
$\sum_{j\ge1}L^{-2j}=(L^2-1)^{-1}$ proves the asserted bound.
This argument uses only multiplication, inversion, unitarity and
the mean estimates, so applies to every closed $G$. Change of
level is a rescaling; the section of a datum is constructed from
its group values by the same contour identities. Neither uses a
lift or a choice of sector.

(iv) \emph{What the operators read.} The operators of the $SU(2)$ construction read the background through $\mathrm{Ad}(U_0(b))$ at the bonds and through the plaquette functions, on pieces of bounded size. Unrolled to the covering lattice the field is periodic; on a box, which is simply connected, every sector is trivial, and the field has a lift $\tilde U$ to $SU(2)$ with $\partial\tilde U(p)=s(\partial U_0(p))$, which lies in the regular class of $SU(2)$ at the parameter $\frac\alpha2(1+0.04\alpha)$, $\alpha\le1$ the parameter of $U_0$: at the plaquettes by the ratio $\frac1{2\cos(\theta/4)}$ of the chords of $s(g)$ and of $g$, $\theta$ the angle of $g$, which is at most $\frac12(1+0.04\alpha)$ for chords at most $\alpha$; and at the bonds, as the scalar part of the divergence of $\partial\tilde U$ at a bond $b$ is at most $0.3\alpha^2\eta^3$ and the norm of its part $i\eta^2J(\tilde U)(b)$ at most half that of the divergence of $\partial U_0$ at $b$ plus $0.3\alpha^3\eta^5$, which keeps the divergence of $\partial\tilde U$ below $\frac\alpha2(1+0.002\alpha^2)\eta^2$ at $\eta\le\frac1{13}$. At the lift the operators of $SO(3)$ are those of $SU(2)$ conjugated by the differential of the covering, which doubles the operator norm on $\mathfrak{su}(2)$, the pairings changing by the one factor $\frac83$; and the action of the pair is $a_3=a^\flat(1-\frac3{16}a^\flat)$, with $a^\flat=\frac83\,a_2\circ s$ and $a_2$ the action of $SU(2)$.

\emph{The operator on the whole torus.} The boxes are taken as the first members of separated sequences of larger tori, at growing radii; there the propagator statements of the $SU(2)$ construction for $SU(2)$ (Theorem~\YMMref{I}{225} and its local inverse construction), read at the lift, give the operators with the Hessian $\Delta^{{\rm W},\flat}$ of $\sum_pa^\flat$, positive and with inverses that decay exponentially, uniformly in the box.

Their limits are kernels on the covering lattice, invariant under the periods of the torus, and their periodisations are, on the whole torus and in every sector, the operators with $\Delta^{{\rm W},\flat}_{U_0}$: positive definite, with inverses that decay exponentially. Since the distance $d_T$ on the torus is the least distance $d$ of two lifts, the sum over the periods keeps with its constant a bound on the sums along rows weighted by $e^{\delta d}$, and it carries a pointwise bound $Be^{-\delta d}$ of a kernel, $d$ in the spacing of the top lattice, to $9Be^{-\delta d_T}$ once $LM\delta\ge2$, which the floor $6LM$ gives for $M$ large. This periodisation factor must be at least $8$ on a torus of even side, where the point antipodal to a site has eight images at one distance from that site.

The Hessian of $a_3$ differs from $\Delta^{{\rm W},\flat}_{U_0}$ by a local symmetric operator $\mathsf E$, $|(\mathsf EA)(b)|\le96\,\alpha_0^2\eta^2\max|A|$, the maximum over the bonds of the four plaquettes at $b$, of the relative order $|\partial U_0(p)-\mathbf 1|^2$; $\mathsf G_0=(1-\mathsf G^\flat_0\mathsf E)^{-1}\mathsf G^\flat_0$, $\mathsf G^\flat_0$ the inverse with $\Delta^{{\rm W},\flat}_{U_0}$, by a Neumann series in norms weighted by the decay, for $\alpha_0$ below the ceiling. Positivity holds along $\Delta^{{\rm W},\flat}_{U_0}-t\mathsf E$, $0\le t\le1$, on which the operator stays invertible.

The bounds take numbers of $L$ in the places of the constants of the $SU(2)$ construction. With $q<1$ the norm of $\mathsf G^\flat_0\mathsf E$ for the sums along rows weighted by the decay, kept below $1$ by the ceiling on $\alpha_0$, a bound $B$ of these sums for $\mathsf G^\flat_0$ becomes $B/(1-q)$ for $\mathsf G_0$, and its pointwise bound $9Be^{-\delta d_T}$, with the factor $9$ of the sum over the periods, becomes $9B(1-q)^{-1}e^{-\delta d_T}$ at the same rate, and the bounds of $\mathsf H_0$ and of $\mathsf G_0-\mathsf H_0\mathsf Q\mathsf G_0$, formed from those of $\mathsf G_0$, take numbers of $L$ in the same way.

\emph{The local clauses.} The same lifts, on boxes shorter than half the torus: the clause of the $SU(2)$ construction for $SU(2)$ is read at the lift and carried back by the covering. Its differential doubles $|\cdot|$ on $\mathfrak{su}(2)$ and multiplies it on $\mathfrak{su}(2)^c=\mathfrak{sl}(2,\mathbb C)$ by a factor between $\sqrt2$ and $2$, as $|z\cdot\sigma|^2=|z|^2+2|x\times y|$ and $|2z\cdot T|=2|z|$ for $z=x+iy\in\mathbb C^3$, $\sigma$ the Pauli matrices and $T$ the generators of $\mathfrak{so}(3)$ to which it takes $\frac12\sigma$. The chords at most $\alpha$ of values in the group are halved at the lift up to the factor $1+0.04\alpha$, as above. The way back multiplies the size of a form by at most $2$ and a chord by $2\cos(\theta/4)\le2$, keeps the norms of operators on real forms, and keeps on complex forms the norms of a real kernel and their sums along rows, a real $3\times3$ matrix having one norm on $\mathbb R^3$ and on $\mathbb C^3$.

Thus a bound $C\prod_i\sigma_i^{d_i}$ of a form or a chord, of degree $d$ in sizes of forms and $e$ in such chords, comes back with $2^{1-d-e}(1+0.04\alpha)^eC$ at real forms and with $2^{1-d/2-e}(1+0.04\alpha)^eC$ at complex ones: with its constant $C$ at $d+e\ge2$, and at $d=1$, $e=0$ for real forms; with $(1+0.04\alpha)C$ at $d=0$, $e=1$, the gauge on a cube among these, of degree one in $\alpha$ and none in forms, whose bound $Bl\alpha$ comes back as $2Bl\cdot\frac\alpha2(1+0.04\alpha)$; with $\sqrt2C$ at $d=1$, $e=0$ for complex forms, and with $C$ there for a linear map whose kernel is real, in bounds on the norms of its kernel or on their sums along rows; and with $2C$ at $d=e=0$. A bound on the norm from suprema to suprema of an operator on complex forms comes back with a factor at most $\sqrt2$, and a clause of $SU(2)$ on the forms below $r$ serves for the complex forms below $\sqrt2r$ and the real ones below $2r$. The bounds of the action are taken anew for $a_3$, whence the absolute number in the place of $C_{4,\mathrm{ec}}$ and the form with $r+5r^2$. The three ceilings keep the fields of a step below the radius under which the averaging map is the formula.

\emph{The three lemmas.} The gauge correction, one step and the ascent are the proofs of the $SU(2)$ construction read at such backgrounds with the operator and these clauses, and so are the statements that these proofs cite and those that these cite in turn. Among them the estimate in the slice (Lemma~\YMMref{30.10}{250}) and the solution of the system (Proposition~\YMMref{30.9}{248}) are statements on the whole torus, whose proofs read the background only through the estimates on $\mathsf G_0$ and $\mathsf H_0$ and through clauses on pieces of bounded size, so that they too hold in every sector. The improvement of the regularity radius is the proof of Proposition~\YMMref{37.1}{278} read so with all its numerals free.

The constants are fixed in this order: $K_3$, $C_{\rm gc}$ and the constant of the gauge on a cube, at most $1.04$ times its value in the $SU(2)$ construction, as numbers of $L$; then the numbers that the $SU(2)$ construction forms from them by its formulas, those of Proposition~\YMMref{37.1}{278} among them; and the ceilings on $\alpha_0$ and on $\varepsilon'$ last, so that the inequalities by which the improvement and the ascent bring the radius back into the regular class, hold with the constants as moved, each being a ceiling on a constant chosen after the numbers that it reads; the rate $\delta$ of the decay stays that of the $SU(2)$ construction, and the floor $6LM$ is laid after it.

\emph{Flat fields.} At a flat field the equation $D_U\lambda=0$ says $\lambda(b_-)=\mathrm{Ad}(U(b))\lambda(b_+)$ at all bonds, and a twist leaves no invariant vector.
\end{proof}

\subsection{Proof of Theorem~\ref{C-368-4}}

For the construction conclusions, assume (P)--(K) throughout. The local calculations in the proof are also valid
without these assumptions, as specified in the statement.

\emph{Step 1: the objects of a step.} The following definitions of the first part of the $SU(2)$ construction,
transferred to $G$ by Definition~\ref{C-368-3}, define objects with the properties that the $SU(2)$ construction
proves: the logarithm on $D^c$, the mean, the block average on small fields and on all fields and its iterates
(Lemmas~\ref{C-368-6}, \ref{C-368-7}); the gradient (Lemma~\ref{C-368-8}(i)); the regular classes, with the
inclusions $\mathfrak U'_k(\mathbf B,r)\subset\mathfrak U_k(\mathbf B,\kappa_Gr)$ for every $G$
(Remark~\ref{C-368-30}) and $\mathfrak U_k(\mathbf B,r)\subset\mathfrak U'_k(\mathbf B,r+r^2+\kappa'r^3)$ for every
$G$ and $r\le z_G$ (Definition~\ref{C-368-3}(3), by the bound at a bond of Lemma~\ref{C-368-8}(ii)), and for
$SU(N)$ and $r\le z_G$ the sharper $\mathfrak U'_k(\mathbf B,r)\subset\mathfrak U_k(\mathbf B,(1+\kappa'r^2)r)$
(Lemma~\ref{C-368-8}(ii)); the gauges (Lemma~\ref{C-368-8}(iii), (iv)); the chart with the indicator of $G^\pi$,
the Haar density and the smooth cut-off (Lemma~\ref{C-368-10}(i), (ii)); the elimination factor of
Lemma~\ref{C-368-10}(iii), a complex linear map of $\mathfrak g^c$ into itself at every $V$ of the domain of the
complex average at $c$, the $G^c$-valued $V$ on $\mathsf P_c$ whose contour families have diameter below
$\frac1{40}$ and whose contour quotients and variables $W_x$ lie in $D^c$, with its branch of $\mathrm{tr}\log$ and
$|\mathsf m_c(e^{ia})-L^{-3}|\le L^{-1}C_S\sup|a|$ for $\mathfrak g^c$-valued $a$ with $\sup|a|<\frac34r_*$
(Lemma~\ref{C-368-10}(iii)); $\mathrm{Log}\det$ at the radius $r_n$ (Lemma~\ref{C-368-10}(iv)); and the domains of
analyticity with their structural facts (Proposition~\ref{C-368-16}(v)). The propagators, the covariances, the
critical configurations and the towers off the trivial field are the objects of (P) and (V) of
Definition~\ref{C-368-31}. Every logarithm is taken of an element of $D^c$, the averages being those of
Lemma~\ref{C-368-7}(i) with $W_x\in D$; integrals over $G\setminus G^\pi$, which has inner points wherever $r_G<2$,
so for $SU(N)$ with $N\ge3$ (Lemma~\ref{C-368-10}(i)), carry the indicator of $G^\pi$; and the uses of the
averaging map off the small fields, of the pivot bond and of the measure of $G\setminus G^\pi$ in the second part
are those of (II) in Definition~\ref{C-368-31}. Among them the pivot bond (Proposition~\YMMref{14.11}{75}) holds for
$G$ in all its clauses with all its constants: the quotients that it reads are products of conjugated plaquette
variables within $\omega_M\le z_G\le r_G$ of $\mathbf 1$, so in $G^\pi$, the averaged contour variables have one
mean each (Lemma~\ref{C-368-6}(iv)), and the equation for the variable at the pivot bond is solved in the chart of
$G^\pi$ by a contraction, the nonlinear part being bounded by Cauchy's estimate.

\emph{Step 2: the trivial field; clause (c).} By Lemma~\ref{C-368-11} every linear operator of a step at $\mathbf
1$ is $K\otimes1_{\mathfrak g^c}$, $K$ a real matrix formed from the lattice alone, and every bound of the $SU(2)$
construction there holds for $G$ with its constant, each in the norms to which its kind passes
(Lemma~\ref{C-368-11}(ii)). The bounds in suprema, by row sums and by non-decreasing functions of real linear forms
hold in $|\cdot|$ and in $\|\cdot\|$ alike:
\[|l_0(f)|\le\Phi(|l_1(f)|,\dots,|l_m(f)|)\ \text{for all}\ f:\Lambda\to\mathfrak g^c\iff\text{the same for all}\ f:\Lambda\to\mathbb R .\]
The $\ell^2$ bounds, the quadratic forms and the lower bounds, the lower bound of the fluctuation operator among
them, hold in $\|\cdot\|$, the sign $|X|^2$ of the $SU(2)$ construction being read as $\langle X,X\rangle$ for
$X\in\mathfrak g$; sums of squares in $|\cdot|$ retain their constants under Schur's test on $|K|$ or convexity. Otherwise
the comparison $\|X\|\le|X|\le N^{1/2}\|X\|$ costs $N^{1/2}$ for an unsquared $\ell^2$ estimate and $N$ for its squared form.
Identification of all small critical configurations with the near-trivial branch requires the orbit input of Definition~\ref{C-368-12}; the local coefficient calculation in
Proposition~\ref{C-368-13}(i)--(iii) gives (c).

\emph{Step 3: fields off the trivial field.} Lemma~\ref{C-368-14} gives the Hessian at $SL(N,\mathbb C)$-valued
fields with the absolute numerals $12$, $1536$, $3456$ of the $SU(2)$ construction, the inverse of a plaquette
variable being bounded by $|X^{-1}|\le|\mathrm{Ad}\,X|$, $\mathrm{Ad}$ acting on all $N\times N$ matrices
(Lemma~\ref{C-368-6}(ii)). For another $G$, that is, for semisimple $\mathfrak g$, so that $G\subset SU(N)$, and
$G^c$-valued $U$, the Hessian of the action of $G$ is the restriction of this form to $\mathfrak g^c$-valued forms;
its two-sided bound and its bounds in the normalised trace norm stay, and the kernel and Weitzenb\"ock bounds in
$|\cdot|$ carry the norm $|P|\le N^{1/2}$ of the projection of the pair, the condition
$|\mathrm{Ad}(U(b))^{\pm1}|\le2$ being read on all $N\times N$ matrices. Proposition~\ref{C-368-16}(i) gives the
inputs of the propagator construction for $SU(N)$, and, for the other simple $G$, the Hessian among them with the
factor $|P|$; (ii) the averages and the critical configuration near the trivial field; (iii) the localisation,
assuming (V); and (iv) the bond maps, with $|E_b-\mathbf 1|<r_G$ at every bond of the cut-off zone under the
ceiling $\mathsf c_Lc_4\le2Lr_G$ of Proposition~\ref{C-368-16}(iv). This is a positive ceiling on the constant
$c_4$, which is free downward, so that it never obstructs the ordered choice (Proposition~\ref{C-368-18}(a)); it
stands among the conditions on the constants read for $G$ under which clause (a) holds. With (P) and (V) of
Definition~\ref{C-368-31} these are the propagators, the covariances, the minimisers over several scales with their
towers, the localisation and the bond maps of the $SU(2)$ construction, transferred to $G$.

\emph{Step 4: conjugation.} Each symmetric domain $Q^s$ of the $SU(2)$ construction is read as its real-connected
part $Q^r$ (Lemma~\ref{C-368-15}): no number moves, but at the memberships reached from memberships of restrictions
had apart. The analyticity of the maps is had on the open domains cut out by conditions $m(X)<r$ and by $X\in D^c$
at the logarithms of the averages, and the clauses on the conjugation are stated on the real-connected parts; $Q^r$
is the part of $Q^s$ whose slice components hold real points, and it passes to intersections as the inclusion
$(\bigcap_iQ_i)^r\subset\bigcap_iQ_i^r$. A map of the induction is a composite of maps built from products,
inverses, means, logarithms in $D^c$, the action and its derivatives, which commute with $p\mapsto p^c$ by
Lemma~\ref{C-368-24}, of maps holomorphic on open slices and real at real points, and of continuations on convex
subsets of a vector space kept by the conjugation (slices without group-valued member); by Lemma~\ref{C-368-15}(ii)
and Proposition~\ref{C-368-22}(i) such a map commutes with $p\mapsto p^c$ on $Q_1^r$ and maps $Q_1^r$ into $Q_2^r$
as soon as it maps $Q_1^r$ into $Q_2$; that every map of the induction is of this kind is (C) of
Definition~\ref{C-368-31}. A membership that is reached from memberships of restrictions had apart is reached in
$Q^s$, and there Lemma~\ref{C-368-15}(iii) carries it into a real-connected part: $\mathcal
Q(b_0,b_1,c_0)^s\subset\mathcal Q(b_0+2e^{4\xi b_1}b_1^2,b_1,c_0)^r$, the like for regular pairs and the natural
classes, and for the analyticity cores $\mathcal U^\circ(b_0,b_1,c_0)^s\subset\mathcal U^\circ(B_0,B_1,c_0)^r$
under the four conditions of Lemma~\ref{C-368-15}(iii), $b'_0=b_0+2e^{4\xi b_1}b_1^2\le B_0$, $b_1\le B_1$,
$K(b'_0+b_1^2)\le B_0$ and $b'_0+b_1\le a^\sharp$, where $K\ge16L^2$ and $a^\sharp$ are the constant and the radius
of the inclusion of the simple sets in the cores. The moved first radius, which is to be at most the first radius
$B_0$ at which the statement read at such a membership is had, and these conditions stand among the conditions on
the constants read for $G$ (Proposition~\ref{C-368-18}): where the first radius of the membership lies below $B_0$,
and for the cores $Kb_0<B_0$ and $b_0<a^\sharp$, they are positive ceilings on the second radius, which is free
downward.

\emph{Step 5: large fields; sectors.} Lemma~\ref{C-368-9} gives the energy of block averages and the large-field
chain, the prices being those of the $SU(2)$ construction times $\tau_G^{-2}$, by $a(V)=\frac12\|V-\mathbf
1\|^2\ge\frac12\tau_G^{-2}|V-\mathbf 1|^2$. Lemma~\ref{C-368-9}(i) holds with the pair $(\frac1{4\tau_G^2},e_*)$,
$e_*=e_L/(2\tau_G^2)$: at one plaquette its proof reads the pair through $|V-\mathbf 1|^2\le2\tau_G^2a(V)$ alone;
and $\tau_G^2\le N-1$. The chain has three links, those of Lemma~\ref{C-368-9}(ii). The first link is
Proposition~\YMMref{44.4}{347} and Proposition~\YMMref{44.5}{348}, on the failed tests. A failed small-field test forces a
plaquette variable or an averaged contour variable farther than the constants of the $SU(2)$ construction from
$\mathbf 1$: by contraposition, a test passed at all plaquettes of a double block puts the field under the
averaging radius $\min\{\omega_M,r_G\}$, where the contour families have one mean within $4d$
(Lemmas~\ref{C-368-6}(iv), \ref{C-368-7}(i)), the straightening of Lemma~\ref{C-368-7}(iii) holding for every
closed $G$. The second link is Lemma~\YMMref{65.2}{580}, on the witnesses of the cells created at one step,
with ``$\notin G^\pi$'' for ``$=-\mathbf 1$''. The third link is the local-deposit argument in the proof of Proposition~\YMMref{65.4}{582}, the lower bounds of the
amount of action near a created cell, which cites clauses of the second. Beside these, the proof of the third link
cites statements on cubes, regions and runs of couplings alone, which read no group and hold for $G$ as written.
The second and the third link hold for $G$ by Lemma~\ref{C-368-9}(ii), by which Proposition~\ref{C-368-17}
transfers the chain and which (II) of Definition~\ref{C-368-31} includes: their proofs derive them, by the device
of the proof of Proposition~\ref{C-368-17}, from the transferred statements that they cite, their inputs asking
$\|E^{(j)}\|\le E_0$ and $|b_j|\le B_\sharp$ separately; and the prices follow from these three links and
Lemma~\ref{C-368-9}(i): the group enters them where the third link turns the witnesses of the second into lower
bounds of the action, at one plaquette through $a(V)\ge\frac12\tau_G^{-2}|V-\mathbf 1|^2$ and at averaged variables
through the pair $(\frac1{4\tau_G^2},e_*)$ of Lemma~\ref{C-368-9}(i), so that the clauses of the third link hold
with the prices of the $SU(2)$ construction divided by $\tau_G^2$. As $G$ is simply connected, the fields on a
torus with small plaquettes form one sector, and the statements on the whole torus that the induction uses hold as
for $SU(2)$; $SU(N)$ is simply connected, and for quotients Proposition~\ref{C-368-27} gives what holds in each
sector.

\emph{Step 6: clause (a).} Let $G$ be simply connected. Assume the five construction hypotheses of
Definition~\ref{C-368-31}. By Steps 1--5, Proposition~\ref{C-368-22} and (II), the statements cited in the proofs
of the induction hold for $G$, with $m(X)$ where $|X^{-1}-\mathbf 1|=|X-\mathbf 1|=|X^*-\mathbf 1|$ was used, with
the radius $r_n$ where $|1-T|\le\frac12$ was to keep $\det T$ off $(-\infty,0]$, with $\max\{1,n2^{-11}\}$ in the
bound \eqref{eq:colour-determinant-bound}, with the ceilings $\varepsilon_{\rm v}\le a_G$, $\varepsilon_{\rm
v}\le\varepsilon_G$ and $\varepsilon_{\rm v}\le L^2r_G$ on the regularity radius (Proposition~\ref{C-368-22}(ii),
Proposition~\ref{C-368-18}(f), Definition~\ref{C-368-31}(V)), the member $r_G/(2C_{\rm pg})$ of the radius $c_{11}$
of Definition~\ref{def:cg-uniform-slice}, the ceiling $\mathsf c_Lc_4\le2Lr_G$ of Step 3 and the conditions of Step 4 on the
radii at the memberships reached from memberships of restrictions had apart among the conditions on the constants
(Definition~\ref{C-368-31}(K)), and with the colour clauses near trivial block data from
Proposition~\ref{C-368-16}(ii) and (V). The ceiling $b_1\le\pi L^j$ under which the local terms are continuous on
the analyticity cores (Lemma~\ref{lem:full-core-continuity}), which a proof of the chain from
level $k$ to level $k+1$ quotes at the fixed radii as $\alpha_1<\pi L^j$, $1\le j\le k$, takes $\rho_G$ of
Definition~\ref{C-368-2} in the place of $\pi$ (Proposition~\ref{C-368-17}): $\alpha_1<\rho_GL^j$, which the row
$\alpha_1<\rho_GL$ of the table gives, as $L^j\ge L$. This row is a positive ceiling on $\alpha_1$ and holds of
itself for $\rho_G>\frac1{13}$, as $\alpha_1\le1$ and $L\ge13$, so for $SU(N)$ with $N\le81$, where
$\rho_G=\frac{2\pi}N$. As $\rho_G\le\pi$, it gives $\alpha_1<\pi L^j$, under which the exponential of $\mathfrak
g^c$ is a local diffeomorphism at the $Z_0$ with $|Z_0|<\pi$ that the continuity on the cores reads, for matrices
of every size, the eigenvalues of $\mathrm{ad}\,Z_0$ being at most $2|Z_0|<2\pi$ in modulus. The maps of the
induction commute with the conjugation by (C). By Proposition~\ref{C-368-17}, the hypotheses of each citation,
those that the forms (3) of Definition~\ref{C-368-3} add among them, being met at its data by (II), the proofs of
the induction then give Theorem~\YMMref{D}{8} for $G$, the mean coefficient being the one number of second
order that they read, so that, with $b_k$ the mean of the $n$ quadratic colour coefficients of step $k$, the run is
\[x_k=L^{-1}x_{k-1}-b_k,\qquad|b_k|\le\mathsf G_\sharp\le B_\sharp,\]
real and free of the torus as for $SU(2)$. By (K) and Proposition~\ref{C-368-18}(b) the conditions on the constants
read for $G$, those of Proposition~\ref{C-368-18}(c)--(g) included, hold together on a large set, so choices exist,
in the order of the $SU(2)$ construction with $\gamma$ last.

\emph{Step 7: clause (b).} The chessboard norm of a source is Proposition~\ref{C-368-19}(ii), for every compact
$G\ne\{\mathbf 1\}$. With $A(v,V)=\sum_pv(p)a_p(V)$ for a complex weight $v$ of the plaquettes and $\mathcal
O(w)=a_K^{-1}A(w,\cdot)$ for a source $w$, the weight with the source is $e^{-A(\beta-a_K^{-1}w,V)}$, of modulus
$e^{-\sum_p(\beta-a_K^{-1}\mathrm{Re}\,w(p))a_p(V)}$, with $a_p\ge0$ on $G$; so $\mathrm{Re}\,w\le(1-c_{\rm
pr})a_K\beta$, $c_{\rm pr}\in(0,1]$, gives $|e^{-A(\beta-a_K^{-1}w,V)}|\le e^{-c_{\rm pr}\beta\sum_pa_p(V)}$: every
plaquette keeps at least the fraction $c_{\rm pr}$ of its price, and the Wilson measure at the one coupling
$s=c_{\rm pr}\beta$ dominates the weight with the source.

The bound on the mean plaquette reads of the group its dimension $n$ and that $a$ vanishes on $G$ at $\mathbf 1$
alone, through $\frac12\tau_G^{-2}|V-\mathbf 1|^2\le a(V)\le\frac12|V-\mathbf 1|^2$ (Lemma~\ref{C-368-6}(i)) and
through the Haar measure of balls, $cr^n\le\mathrm{Haar}\{|V-\mathbf 1|<r\}\le C''r^n$, the first for $r\le1$, $c$
and $C''$ numbers of the pair, which the chart of Lemma~\ref{C-368-10}(i) gives, its density being at most $1$ and
bounded below on $|Y|\le\frac\pi3$. With the first of the two bounds on $a$ they give $\int_Ge^{-\frac
s2a(W)}\,dW\le C_1s^{-n/2}$ for $s>0$, $C_1$ a number of the pair; in the chart, with
$a(e^{iY})\ge\frac2{\pi^2}\|Y\|^2$ for $|Y|<\pi$ and $a\ge\frac12\tau_G^{-2}r_G^2$ off $G^\pi$,
$\int_Ge^{-ya(u)}\,du$ lies between two numbers of the pair times $y^{-n/2}$, the lower bound for $y\ge1$ and the
upper one for $y>0$. Let the couplings of the plaquettes of a torus with $|\Lambda|$ sites satisfy $s\le v(p)\le
C's$, $s\ge1$, $C'\ge1$, let $Z(v)=\int e^{-A(v,V)}\,dV$, and let $\langle\cdot\rangle_v$ be the Wilson measure at
the couplings $v$. As $\log Z(tv)$ is convex in $t$ with the derivative $-\sum_pv(p)\langle a_p\rangle_{tv}$, this
weighted sum does not increase with $t$, its derivative being minus a variance, and $\frac s2\sum_p\langle
a_p\rangle_v\le\frac12\sum_pv(p)\langle a_p\rangle_v\le\log Z(\frac v2)-\log Z(v)$.

From below: in the gauge in which a maximal tree carries $\mathbf 1$ the $2|\Lambda|+1$ bonds off the tree are
independent Haar variables, and restricting them to within $(C's)^{-1/2}$ of $\mathbf 1$ gives
$a_p\le\frac12(4(C's)^{-1/2})^2$ by the second bound on $a$, so $v(p)a_p\le8$ at each of the $3|\Lambda|$
plaquettes, and $Z(v)\ge(c(C's)^{-n/2})^{2|\Lambda|+1}e^{-24|\Lambda|}$.

From above: of all the bonds, $2|\Lambda|-2$ can be ordered and matched to distinct plaquettes so that each lies in
its own plaquette and in none of the plaquettes of the earlier ones; dropping the other plaquettes, as $a_p\ge0$,
and integrating these bonds one by one, the last first, by the invariance of the Haar measure, gives $Z(\frac
v2)\le(C_1s^{-n/2})^{2|\Lambda|-2}$. So the mean $\frac1{3|\Lambda|}\sum_p\langle a_p\rangle_v$ is at most
$(C+n\log C')/s+n\log s/(|\Lambda|s)$, $C$ a number of the pair. At couplings invariant under the lattice
symmetries every plaquette has this mean, so at the one coupling $s$, where $C'=1$, $\langle a_p\rangle_s\le
C/s+n\log s/(|\Lambda|s)$ at every plaquette. The term $n\log s/(|\Lambda|s)$ comes from the three holonomies of
the torus, $2|\Lambda|+1$ bonds off the tree against $2|\Lambda|-2$ matched bonds; on the tori of the continuum
theorem, with at least $2L^{K+1}$ sites a side, at couplings at most $L^{K+1}(g^{-2}+\bar B)$, that of its
first-passage data, $n|\Lambda|^{-1}\log s$ is below a number of $g$ and the constants. At equal couplings
equipartition gives the mean plaquette $\frac n{3s}(1+o(1))$ on large tori as $s\to\infty$; so the number $C$ in
the bound reads the pair, and it is at least $\frac n3$.

Let $G$ be simply connected and assume the construction hypotheses. The sourced-density and stability assertions of Theorems~\YMMref{AJ}{673} and~\YMMref{AD}{607} hold for $G$ with
the run of (a), by the corresponding $SU(2)$ proof run for $G$: a source enters as a weight of the plaquette
function $a_p$, which for simple $\mathfrak g$ is the same class function as in the action, so the steps of the
induction of (a) run for the weighted action, whose gradient and Hessian are the weighted ones, through complex
plaquette weights near $\beta$ and uniformly in the volume, the price of a plaquette keeping the prices of the
large fields; the effective densities with a source have at every level the structure that the $SU(2)$ construction
proves, and writing $\widetilde Z(w)=e^{-E_{(K)}-\mathcal E_K(w)}Z_{K,\ell,\beta}[w]$, with $\mathcal E_K$ defined by \eqref{eq:running-source-energies}, gives the factorisation of the partition
function with a source that (b) displays, with the bounds of Theorem~\YMMref{AD}{607} on its factors. With it the
bound per ball of (b) is the passage of the $SU(2)$ construction from the stability to
Proposition~\textup{\YMMref{89.4}{693}}, run for $G$, at every reference coupling $0<g\le\gamma$, at every first-passage
datum $(K,\beta)$ at $g$ and on every side $\ell$ admitted at $g$ with $\ell\ge16\mathsf m(R+1)$, for supports in
balls of radius $R$ pairwise at a distance at least $\delta$: the tuple of test functions is placed on a board by a
lattice translation (Lemma~\YMMref{89.3}{692}); the stability gives on each cell a normaliser without zeros, with its
bound per cell (Proposition~\YMMref{89.1}{690}, Proposition~\YMMref{89.1}{690}, Proposition~\YMMref{89.2}{692}); the
expectation with sources is dominated by the chessboard norms of the occupied cells, through subadditivity and the
comparison of two volumes (equation~\YMMref{(1759)}{682}, Proposition~\YMMref{86.5}{682}, Corollary~\YMMref{86.6}{683},
which hold for $G$ by Proposition~\ref{C-368-19}(ii)), the price of a plaquette and the bound on the mean plaquette
entering as in the $SU(2)$ construction; and the graded Cauchy inequality on the polydisc of sources
(Lemmas~\YMMref{87.1}{684} and~\YMMref{87.2}{684}) gives the bound per ball, uniformly in the torus, with the constant
$C^\infty_{\mathsf m}(\delta;R)$ of Proposition~\textup{\YMMref{89.4}{693}}, which these sections form from $g$, $\mathsf
m$, $\delta$, $R$ and the constants of Theorem~\YMMref{AD}{607} for $G$ alone, free of $K$, $\beta$ and $\ell$.
These sections of the $SU(2)$ construction read the group through the statements named alone.

\emph{Step 8: clause (d).} The increments of the run are real, depend on no torus and are bounded by $B_\sharp$, by
(a); so $|x_j-L^{-(j-i)}x_i|\le\bar B(1-L^{-(j-i)})$, for $\beta>X_g$ there is exactly one first passage $K$, and
\[g^{-2}<x_K\le X_g,\qquad\frac\beta{L(g^{-2}+\bar B)}\le L^K<\beta g^2,\qquad g^{-2}+B_\sharp-\frac{\bar B}L<a_K\beta\le L(g^{-2}+\bar B);\]
the window $I(g)$ has the length $(L-1)g^{-2}+B_\sharp\frac{L(L-3)}{L-1}>0$, and $L^KI(g)$ consists of
first-passage couplings; so $a_{K_j}\beta_j$ has a convergent subsequence in the closed interval. For simply
connected $G$ satisfying the construction hypotheses, by (b) the bound per ball holds uniformly in $j$, at the
first-passage data $(K_j,\beta_j)$ and on the sides $l_j$ admitted at $g$, from the $j$ with $l_j\ge16\mathsf
m(R+1)$ on: $|S^{(j)}_{\mathsf m}(f)|\le C^\infty_{\mathsf m}(\delta;R)\prod_i\|f_i\|$, in the norms of Lipschitz
functions of Proposition~\textup{\YMMref{89.4}{693}}, each at most a constant times $\|f_i\|_{C^1}$, for supports in balls of
radius $R$ pairwise at a distance at least $\delta$, $C^\infty_{\mathsf m}(\delta;R)$ being formed from $g$,
$\mathsf m$, $\delta$, $R$ and the constants of Theorem~\YMMref{AD}{607} for $G$ alone; on one torus of side $l$
the bound $C_{\mathsf m}(\delta)e^{\mathsf mc_Vl^3}\prod_i\|f_i\|_{C^1}$ holds (Proposition~\YMMref{88.3}{688});
with it, a diagonal subsequence over a countable dense set of tuples and $\mathsf m$-linearity give convergence.
The properties of the limits need convergence alone and hold for every compact connected $G$: the bound by uniform
boundedness for $\mathsf m$-linear forms on Banach spaces, which is why $C$ depends on the sequence; the
translations and the signed permutations from the lattice symmetries, which keep the measure, the action and the
observables; reflection positivity from Proposition~\ref{C-368-19}(i); and on one torus its translations and the
pairs of parallel mirrors at the distance $\frac l2$. For the additional conclusions assume throughout that $G$ is
simply connected and that \textup{(P)}, \textup{(V)}, \textup{(C)},
\textup{(II)} and \textup{(K)} hold. The transferred construction
in (a), with the sourced estimates in (b), supplies the inputs of
the fixed-torus comparison in Theorem~\YMMref{A}{5}; its proof
uses these estimates at the unchanged label and admitted sides.
The bounds in (b) give the compactness and distributional limits
of Propositions~\YMMref{92.1}{710}--\YMMref{92.4}{713} by the
same diagonal extraction and passage of finite positive matrices
used above. For Gibbs limits, Proposition~\ref{C-368-19}(iii)
supplies periodic Gibbs states, their lattice symmetries and all the
required reflection families. Conditioning on a translation-invariant
event preserves the Gibbs specification by the conditional-expectation
argument of Proposition~\YMMref{106.1}{764}, which uses no special
identity of the group. Together with the transferred
cumulant and one-point bounds, these are exactly the hypotheses
used for the component construction in Theorem~\YMMref{AL}{746}
and its application in Theorem~\YMMref{AP}{780}. The
reflection-Hilbert-space and reconstruction arguments concern
the resulting scalar functionals and use no further identity of
$SU(2)$. Thus they give the Gibbs conclusions stated in (d).
For growing-torus limits this argument supplies coordinate
reflection positivity; their full rotational invariance is the
criterion proved in Step~10. The colour coefficient in the
short-distance law is computed in Step~9.

\emph{Step 9: the leading law and the perturbative ratio.}
The quadratic part of the action at $\mathbf 1$ is $\frac12\sum_p\|dA(p)\|^2$, colour by colour
(Lemma~\ref{C-368-11}(ii)). In an orthonormal basis of $\mathfrak g$, its Gaussian measure is therefore the product
of $n$ copies of the one-colour Gaussian measure for $SU(2)$, with the same normalisation. Each colour contributes
$3/(16\pi^2)$, and hence
\[
 c_2(G)=\frac{3n}{16\pi^2}.
\]
Under the construction hypotheses, the estimates in (a) and (b) give the effective-action remainders used in the
finite-cutoff leading law of Theorem~\YMMref{AV}{851}. Its colour sum now runs to $n$, and the label and continuum passage use Theorem~\YMMref{AW}{852}. The passage to limits is the one in (d).
For $SU(2)^r$ the product argument is given in Theorem~\ref{C-368-5}(iii).

For (f), Lemma~\ref{C-368-20} gives the coefficient of $\beta^{-3}$ first at fixed box and then at infinite volume.
Its leading part at large separation is $\varphi_G\mathsf r^{-5}/(16\pi^3)$, after summing over orientations.
Dividing this coefficient by the leading one gives
\[
 c_1(G)=\frac{\varphi_G}{16\pi^3}\Big/\frac{3n}{16\pi^2}
       =\frac{\varphi_G}{3\pi n}.
\]
For $SU(2)$, $\varphi_G=24$, so that $c_1=8/(3\pi)$.

\emph{Step 10: clause (g).} Proposition~\ref{C-368-19}(iv) replaces the formula of the $SU(2)$ construction for
each generator. The rest of its proof of the Ward identity at finite cutoff (Proposition~\YMMref{112.1}{782}) is an
integration by parts on $G$ at every cutoff, $\int v_{\mathcal X}\Phi\,dV=-\int\Phi\,\mathrm{div}\,v_{\mathcal
X}\,dV$ for a field of directions $\mathcal X$, and it is unchanged but for the divergence of the rotation field,
which carries $\frac12k_p$ where the $SU(2)$ construction has $\frac32a_p$. For $\omega\in\mathfrak{so}(3)$ let
$(L_\omega f)(y)=\sum_\rho(\omega y)_\rho\,\partial_\rho f(y)$, let $e^{t\omega}f$ be the tuple of the functions
$f_i(e^{-t\omega}\,\cdot\,)$, and let $\nabla^{\rm s}_\omega$ be the symmetric lattice difference along $\omega y$,
$(\nabla^{\rm s}_\omega\phi)(y)=\sum_\rho(\omega y)_\rho(2a)^{-1}[\phi(y+ae_\rho)-\phi(y-ae_\rho)]$. A quadruple
$(\omega,\mathsf m,\mathcal R,f)$ is admitted (Section~\YMMref{112}{781}, before equation~\YMMref{(2115)}{784}) if $\omega\in\mathfrak{so}(3)$,
$\mathsf m\ge2$, $\mathcal R>0$ and the $f_i$ are smooth, with compact pairwise disjoint supports in
$\{|y|<\frac12\mathcal R\}$. Let $(K_j,l_j,\beta_j)$ be a sequence of tori as in (d), $K_j\to\infty$, with
$l_j\to\infty$. For the $j$ with $2\mathcal R+2a_{K_j}<\frac12l_j$ the ball $\{|y|<\frac53\mathcal R+a_{K_j}\}$,
which carries the cut-off of radius $\mathcal R$ and its shell, has a diameter below $l_j$ and lies in the torus
without wrapping, so the bulk term and the flux term of the identity at the cut-off radius $\mathcal R$ are
defined; $j_1(\mathcal R,f)$ is the least index from which on $2\mathcal R+2a_{K_j}<\frac12l_j$,
$2\sqrt3\,a_{K_j}\le\frac12\mathcal R$ and $a_{K_j}\le\frac14\min\{1,\mathrm{sep}(f)\}$, $\mathrm{sep}(f)$ the
least distance of two of the supports. For $j\ge j_1(\mathcal R,f)$ the identity at the cutoff $K_j$, with the
radial cut-off $\chi_{\mathcal R}(y)=\chi_0(|y|/\mathcal R)$ (choose $\chi_0$ smooth, equal to $1$ on $[0,4/3]$ and zero on $[5/3,\infty)$, as permitted in Section~\YMMref{112}{781}), which equals $1$ on
$\{|y|\le\frac43\mathcal R\}$ and so covers the $f_i$, reads
\[\sum_iS^{(j)}_{\mathsf m}(f_1,\dots,\nabla^{\rm s}_\omega f_i,\dots,f_{\mathsf m})=\mathfrak
U^{(j)}_{\omega,\mathcal R}(f)+\mathrm{Fl}^{(j)}_{\omega,\mathcal R}(f),\]
the bulk term $\mathfrak U^{(j)}_{\omega,\mathcal R}(f)$ collecting the spin and orbital insertions, and the flux
$\mathrm{Fl}^{(j)}_{\omega,\mathcal R}(f)$ being the connected function of the $\mathcal O(f_i)$ with the shell
flux, a sum over the plaquettes $p$ with $\frac43\mathcal R-a<|x(p)|<\frac53\mathcal R+a$ that carries $\frac12k_p$
(Proposition~\YMMref{112.1}{782}, under~\YMMref{(2116)}{784}). In the limit, along such a sequence on which all connected functions converge,
\[\lim_j\bigl[\mathfrak U^{(j)}_{\omega,\mathcal R}(f)+\mathrm{Fl}^{(j)}_{\omega,\mathcal
R}(f)\bigr]=\sum_iS_{\mathsf m}(f_1,\dots,L_\omega f_i,\dots,f_{\mathsf m})=-\frac d{dt}S_{\mathsf
m}(e^{t\omega}f)\Big|_{t=0},\]
independently of $\mathcal R$ among the radii that admit $f$: the bound of (d) (Step~8), on the spaces of $C^1$
functions with supports in fixed disjoint compact neighbourhoods of the supports, with $\|\nabla^{\rm s}_\omega
f_i-L_\omega f_i\|_\sharp\le c\,a_{K_j}^2$ and $\|t^{-1}(e^{t\omega}f_i-f_i)+L_\omega f_i\|_\sharp\le c|t|$, $c$ a
number of $\omega$ and $f_i$, gives the limit and the differentiability of $t\mapsto S_{\mathsf m}(e^{t\omega}f)$
(Proposition~\YMMref{112.3}{784}). The conjugates of one $\omega_0\ne0$ under the signed permutations span
$\mathfrak{so}(3)$, and the one-parameter group of $\omega_0$ with the cubic group generates $O(3)$
(Lemma~\YMMref{112.4}{785} and the proof of Theorem~\YMMref{AQ}{785}), which gives the equivalence at the admitted quadruples (Theorem~\YMMref{AQ}{785}); the
invariance under $O(3)$ passes to the $C^1$ tuples of (d) by mollifying, with the bound of (d): the mollified
entries and their gradients converge uniformly, with supports in fixed disjoint compact sets, on whose $C^1$
functions that bound is continuous. The limit of $\mathrm{Fl}^{(j)}_{\omega,\mathcal R}(f)$ alone is this limit
exactly when $\mathfrak U^{(j)}_{\omega,\mathcal R}(f)\to0$. When the limit is invariant, the rotations about the
third axis and the reflections in planes through it keep the half-space $\{y_3>0\}$, commute with the reflection in
$\{y_3=0\}$ and keep the reflection positive form, so $\mathsf U(R)[F]=[(R\oplus1)F]$, on the classes of the
functions $F$ supported in $\{y_3>0\}$, is well defined modulo the null space, unitary and fixes the vacuum vector;
as $R\oplus1$ commutes with the translation by $te_3$, $\mathsf U(R)[\tau_{te_3}F]=[\tau_{te_3}((R\oplus1)F)]$ for
$t\ge0$. Strong continuity comes from the continuity of the $S_{\mathsf m}$.

\emph{Step 11: clause (h).} (h1) is Lemma~\ref{C-368-21} with the bound of Proposition~\ref{C-368-17}; (h2) is
Proposition~\ref{C-368-13}(iv); (h3) is Proposition~\ref{C-368-26}(iv), with Lemma~\ref{C-368-21} at
$\Gamma=\mathsf G_\sharp$, as $|e^{(i)}_k|\le(\kappa_i/\kappa_G)^2\mathsf G_\sharp\le\mathsf G_\sharp$
(Proposition~\ref{C-368-17}), which gives
\[|x^{(i)}_k-x_k|\le2\Bigl(1-\frac{n_i}n\Bigr)\mathsf G_\sharp\frac L{L-1}(1-L^{-k}),\]
and so the same with $\kappa_G^2\mathsf G_\sharp$ in the place of $\mathsf G_\sharp$, as $\kappa_G\ge1$; and (h4)
is Lemma~\ref{C-368-29}(i) with clauses (a)--(g) for each factor, in the forms stated there. $\square$


\subsection{Proof of Theorem~\ref{C-368-5}}

\emph{Step 1: clause (i).} $a(u)=\sum_i\lambda_ia_i(u_i)$ and the Haar measure is a product, so the fields of the
factors are independent (Lemma~\ref{C-368-29}(i)); cumulants of sums of independent families add, and $\mathcal
O=\sum_i\lambda_i\mathcal O^{(i)}$ gives the factor $\lambda_i^{\mathsf m}$.

\emph{Step 2: clause (ii).} Proposition~\YMMref{52.4}{435} proves covariance of the
pointed totals for the actual arrays admitted by Theorem~\YMMref{W}{423}.
Its proof uses the earlier array through the sum in
equation~\YMMref{(963)}{436}: earlier pointed totals, constants,
and multiples of the plaquette action. The constants and the action
are invariant, independently of the counterterm coefficients.
Thus the covariance induction applies to these arrays whenever
the earlier pointed totals are covariant. So symmetry passes from level to level
along every good sequence, $|L^{-1}x_{j-1}-x_j|\le B_\sharp$ and $x_j\ge\gamma^{-2}$, and along it the coefficient
of $SU(2)$ is defined, real and at most $\mathsf G_\sharp$, and Theorem~\YMMref{D}{8} holds whatever its
quadratic leftovers. The sequences $\lambda_ix_j$ are good: their increments are
$\lambda_i\frac1r\sum_m\lambda_m^{-1}c^{(m)}_j$, at most $\lambda_+\lambda_-^{-1}\mathsf G_\sharp\le B_\sharp$, and
$\lambda_ix_j\ge\lambda_-x_j\ge\gamma^{-2}$. Then
\[|x^{(i)}_k-x_k|\le\mathsf G_\sharp\Bigl(\lambda_i^{-1}+\frac1r\sum_m\lambda_m^{-1}\Bigr)\frac
L{L-1}(1-L^{-k})\le2\lambda_-^{-1}\mathsf G_\sharp\frac L{L-1}.\]

\emph{Step 3: clause (iii).} The first-passage arithmetic of Proposition~\YMMref{74.1}{615} and Corollary~\YMMref{74.2}{616} applies to the mean run, whose increments are real and bounded by $B_\sharp$. The
windows $I(g_{i,1})=(\lambda_ig^{-2},L\lambda_ig^{-2}-2\bar B]$ and $I(g_{i,2})=(L\lambda_ig^{-2}-2\bar
B,L^2\lambda_ig^{-2}-2(L+1)\bar B]$ are adjacent, and $\lambda_ia_K\beta$ lies in one of them, $I(g_{i,s(i)})$, so
$(K,\lambda_i\beta)$ is a first-passage datum of the factor $i$ at the reference coupling $g_{i,s(i)}$, $s(i)$
being read off $(K,\lambda_i\beta)$. The law of the continuum theorem for the factor $i$ is cited at $g_{i,s(i)}$,
with the hypotheses on the scale, the cutoff and the test pair that (iii) lays at the $2r$ couplings, and on tori
whose side is admitted at $g_{i,s(i)}$, which the hypothesis that the sides are admitted at all $2r$ couplings
$g_{i,s}$ provides; these laws and (i) give
\[\sum_i\lambda_i^2\,\frac9{16\pi^2}(\lambda_i\mathsf b)^{-2}\mathcal K_6=\frac{3n}{16\pi^2}\,\mathsf b^{-2}\mathcal K_6 .\]
In a Gibbs state with periodic marginals the factors are independent given the tails, the diagonal entries are the
laws of the factors in their periodic marginals, and the mixed entries are bounded by Lemma~\ref{C-368-29}(iii),
whose range $a_K\le\mathsf d\le\frac1{20}$ holds, as the hypotheses of the law on the scale and the cutoff give
$a_K<\mathsf d\le\frac1{200}$; as $\sum_{i\ne
j}\lambda_i\lambda_j\frac14\varepsilon^\times\frac3{16\pi^2}(\lambda_i\lambda_j\mathsf
b^2)^{-1}=\frac{r-1}{12}\varepsilon^\times\cdot\frac{9r}{16\pi^2}\mathsf b^{-2}$, this gives the error term
$\frac{r-1}{12}\varepsilon^\times(\mathsf d)$, the divergent parts never mixing the factors. $\square$


\subsection{Scope of the conclusions}

The full construction for a simply connected simple group is conditional on Definition~\ref{C-368-31}. Its two
principal outputs are the source stability in Theorem~\ref{C-368-4}(b) and the resulting bound per ball. The latter
yields subsequential limits. In contrast, the algebraic coefficients and the finite-box perturbative expansion
require no global construction hypothesis. The error constant in that expansion depends on the box, so its stated
limit is iterated.

For any compact connected group, convergence of the connected functions along a sequence of tori gives the
translation symmetry, cubic symmetry and reflection positivity stated in Theorem~\ref{C-368-4}(d). Its rotational
Ward identity gives the criterion in (g): full $O(3)$ invariance is equivalent to the vanishing of the indicated
bulk-plus-flux limit. No vanishing assertion is included in that criterion.

Products acting on separate blocks are governed by exact factorisation. For tied simple ideals, the one-loop
calculation is block diagonal at the trivial field, while the nonlinear action couples the ideals. The assertion
that the mean run suffices has property (T) as its explicit hypothesis. Quotients require the sector distinctions
in Proposition~\ref{C-368-27}. The regular-field construction here assumes a semisimple Lie algebra: a nonzero
centre obstructs the box gauge in Lemma~\ref{C-368-8}(iv), and invariant linear forms then need not vanish.

\section{Wilson loops on a fixed torus}\label{C-372-1}

We consider three limits of lattice Wilson loops. Background-smoothed loops converge along almost every block-field thread. Sharp polygonal loops have renormalised expectation limits in the stated perimeter range. The unrenormalised trace has a semicircle limit, and division by its vanishing mean produces divergent variance.

These are different limiting statements. The sharp-loop theorem concerns expectations on a fixed torus; its stronger mixed-function conclusions have their own perimeter and separation restrictions. The proof retains the ordered-product estimate, the inserted-line renormalisation step and the summability restriction on the normalised errors.

\subsection{Notation and hypotheses}

\emph{Constants.} Use a common real numerical choice as in Proposition~\YMMref{69.1}{595}. We use the odd blocking factor $L\ge13$; $M=L^{m'}$ with $m'\ge5$; $r\in\mathbb N$; $\gamma\in(0,\frac12]$ with $\log\gamma^{-2}\ge\log4>1$; $A_0>0$, $\mathfrak p_0\in\mathbb N$ and the threshold $\varepsilon(t)=A_0t^{-1/2}(\log t)^{\mathfrak p_0}$, $t>1$, of \YMMref{(634)}{339}, which decreases for $\log t>2\mathfrak p_0$; $B_3\ge1$; $c_\rho=\max\{2B_3,\,C_HA_1/((1-L^{-3/2})A_0)\}$ with the constants $C_H$, $A_1$ of this formula (Section~\ref{sec:cutoff-calibration}), so that $2\le2B_3\le c_\rho<\infty$; $c'\in(0,\frac1{250})$, $\mathsf C^\Sigma_{\rm maj}\ge0$ and $c_{\rm low}\ge24$.

The number $t_\natural\ge\gamma^{-2}$ is the least real number from which on the three smallness conditions (cut1)--(cut3) of Section~\ref{sec:cutoff-calibration} on $\varepsilon(t)$ hold; they hold at every $t\ge t_\natural$, and there $2c_\rho\varepsilon(t)\le\varepsilon_{\mathcal W}\alpha_0\le\bar\varepsilon\le1$, $\bar\varepsilon$ the radius of Theorem~\YMMref{E}{10} in the setting of Section~\YMMref{52.7}{432} and $\varepsilon_{\mathcal W}$, $\alpha_0$ the numbers of (cut1), formed from the constants alone (equation~\YMMref{(2134)}{790}). $\Theta_\star\ge t_\natural$ is the least real number from which on the five conditions of the comparison of two consecutive cutoffs (Proposition~\YMMref{113.9}{795}) hold; both exist (the threshold choices in Sections~\YMMref{113}{786} and~\YMMref{115}{805}, with Corollary~\ref{cor:canonical-running-source}), and $2c_\rho\varepsilon(t)<1/200L^2$ for $t\ge\Theta_\star$.

The lowered choice of the constants (Definition~\ref{def:rarity-choices}: $A_1$ multiplied by $\lambda'_\natural$, $B_\sharp$ not smaller, $A_0$ and the constants fixed before $A_1$ kept) gives the number $c'_\star$, which is $c_\star$ of \eqref{eq:rarity-constants} at that choice, and the numbers $\mathsf Y^\natural_1,\dots,\mathsf Y^\natural_4$ of \eqref{eq:rarity-seven-conditions} formed from the two choices; we put $c^\natural=c'_\star/(4\cdot49^3)>0$ and $Y^\natural=\max_{i\le4}\mathsf Y^\natural_i$. All these numbers depend on the constants alone.

\emph{The one-torus setting}: a label $\mathsf b\ge\gamma^{-2}$ and a side $\ell=2L^m$, $m\in\mathbb N$, with $L^m\ge M(\log\mathsf b)^r$, equivalently $M\,Q(\mathsf b)\le L^m$ (Lemma~\ref{C-372-7}).

\emph{Fields and laws.} $T_\ell=(\mathbb R/\ell\mathbb Z)^3$, $a_d=L^{-d}$ and $\Lambda_d=a_d\mathbb Z^3/\ell\mathbb Z^3$ (Section~\YMMref{3}{13}). $\mathcal V_d$ is the set of the fields on the bonds of $\Lambda_d$ with values in $G=\mathrm{SU}(2)$, $V(\bar e)=V(e)^{-1}$, with the gauge transformations $V^v(e)=v(e_-)V(e)v(e_+)^{-1}$ (Section~\YMMref{3}{13}); $V(\Gamma)$ is the ordered product along a path $\Gamma$, $\partial V(p)$ the ordered product along the boundary of a plaquette $p$, and $\varpi(V)=\max_p|\partial V(p)-\mathbf 1|$, in the operator norm . $\mathcal M:\mathcal V_{d+1}\to\mathcal V_d$ is the block averaging map (equations~\YMMref{(37)--(39)}{20}); it is Borel and gauge covariant, and one map serves all cutoffs.

The Wilson measure $\langle\cdot\rangle_K$ on $\mathcal V_K$ has the density $Z^{-1}e^{-\beta\sum_pa_p}$ with respect to the Haar measure $dU$, where $a_p=1-\frac12\operatorname{tr}\partial U(p)$ and $\beta=L^K\mathsf b$. $\mu^K_d$ is the law of $\mathcal M^{K-d}U$ under $\langle\cdot\rangle_K$, and $\mu_d=\lim_K\mu^K_d$ in total variation along the full sequence of cutoffs (Corollary~\textup{\ref{IV-62-66}}). $\mathcal X=\prod_{d\ge0}\mathcal V_d$ carries the product topology and the coordinates $\pi_d$; $\mathsf{Thr}=\{V_\bullet:V_d=\mathcal M(V_{d+1})\ \forall d\}$ is the set of the threads; $\mathcal F_d=\sigma(\pi_0,\dots,\pi_d)$; $(\theta_vV_\bullet)_d=V_d^{\,v|\Lambda_d}$ for $v:\bigcup_d\Lambda_d\to G$.

Finally $\mu$ denotes the Borel probability measure on $\mathcal X$ whose image under $V_\bullet\mapsto(V_0,\dots,V_D)$ is, for every $D$, the image of $\mu_D$ under $V\mapsto(\mathcal M^DV,\dots,\mathcal MV,V)$; it exists and is unique (Proposition~\ref{C-372-8}(a)).

\emph{Depths.} We write $x_{[d]}=L^d\mathsf b$ and $\varepsilon_{[d]}=\varepsilon(x_{[d]})$, with brackets, so as to tell these depth quantities from the run quantities $x_t$, $\varepsilon_t$ of the frame below. In $\mathcal V_d$ let $I_d=\{\varpi<2\varepsilon_{[d]}\}$ and $J_d=\{\varpi<c_\rho\varepsilon_{[d]}/B_3\}$; $I_d\subset J_d$ since $c_\rho\ge2B_3$. Put
\[R_d=\frac{2\ell^3L^{3d}}{M^3}\,e^{-c'A_0^2(\log x_{[d]})^{2\mathfrak p_0}+(\mathsf C^\Sigma_{\rm maj}+2+c_{\rm low}+\frac92\log\mathsf b)\ell^3},\qquad R^\natural_d=\ell^3L^{3d}e^{-c^\natural A_0^2(\log x_{[d]})^{2\mathfrak p_0}}.\]
A depth $d$ is admitted if $x_{[d]}\ge\Theta_\star$. $d_\natural$ and $d_\Theta$ are the least integers $d\ge1$ with $x_{[d]}\ge t_\natural$, resp.\ $x_{[d]}\ge\Theta_\star$, and $s_\natural$ is the least integer $s\ge0$ with $\log x_{[s]}\ge Y^\natural$; they depend on $\mathsf b$ and the constants, not on $\ell$, and $d_\Theta\ge d_\natural$. $\tau_0(V_\bullet)$ is the least $D\ge d_\natural$ with $V_d\in I_d$ for all $d\ge D$ ($\infty$ if there is none); the starting depth is $\tau=\max\{\tau_0,d_\Theta\}$, and $\Omega=\mathsf{Thr}\cap\{\tau<\infty\}$.

\emph{Loops and their smoothing.} A lattice loop $C$ of $\Lambda_{d'}$ with $N$ bonds has the length $|C|=Na_{d'}$; for $j\ge d'$, $C^{\langle j\rangle}$ arises by replacing $j-d'$ times every bond by the straight path of $L$ bonds of the next lattice, and $W_C(V)=\frac12\operatorname{tr}V(C^{\langle j\rangle})$ for $V\in\mathcal V_j$. For $W\in\mathcal V_d$ and $K\ge d$ the backgrounds over $W$ form the set $\mathcal C_K(W)$ of the $U\in\mathcal V_K$ with $\mathcal M^{K-d}U=W$ that are critical for $\sum_pa_p$ on this constraint set and regular at the radius $\varepsilon_{\rm v}$ (Definitions~\YMMref{9.3}{40} and~\YMMref{9.5}{41}); $\mathcal G^d_K$ is the group of the gauge transformations of $\Lambda_K$ that equal $\mathbf 1$ on $\Lambda_d$, and $w^K_C(W)=W_C(U)$ for $U\in\mathcal C_K(W)$.

For $X\in\mathcal V_s$ let $H^K_s(X)=\mathcal M^{K-s-1}U_+\in\mathcal V_{s+1}$, $U_+$ an element of $\mathcal C_K(X)$ for which the average $\mathcal M^{K-s-1}U_+$ is in the axial gauge on $\Lambda_s$ (Section~\ref{sec:sf-chart-interface}); $H^\infty_s=\lim_KH^K_s$ and $H^\bullet_{d,n}=H^\bullet_{d+n-1}\circ\dots\circ H^\bullet_d$. The smoothed loop is $W^{[d]}_C(W)=\lim_nW_C(H^\infty_{d,n}(W))$; $F^{[d]}_C$ is $W^{[d]}_C$ on $J_d$ and $0$ off $J_d$; $Y_{d,C}=F^{[d]}_C\circ\pi_d$ and $Y^K_{d,C}=(\mathbf 1_{J_d}w^K_C)\circ\pi_d$. A curve is a continuous $\hat\gamma:[0,1]\to\mathbb R^3$ of finite length $|\gamma|$ with $\hat\gamma(1)-\hat\gamma(0)\in\ell\mathbb Z^3$; a shadow $\Gamma^{(j)}$ of it on $\Lambda_j$ cuts it into at most $|\gamma|L^j+2$ pieces of sup-diameter $\le a_j$, takes sites of $\Lambda_j$ within $\frac12a_j$ of the cut points and joins them by shortest paths of $\Lambda_j$; $W^{[d]}_\gamma(W)=\lim_nW_{\Gamma^{(d+n)}}(H^\infty_{d,n}(W))$, $w^K_\gamma(W)=W_{\Gamma^{(K)}}(U)$ for $U\in\mathcal C_K(W)$, and $F_\gamma$ is $W^{[d]}_\gamma$ on $J_d$ and $0$ off it.

\emph{Polygons.} A polygon is a closed self-avoiding lattice path with finitely many axis-parallel pieces (maximal straight parts) and corners in some $\Lambda_{i_0}$, lying in a ball of radius $<\ell/4$. A family $\mathsf C=(C_1,\dots,C_{\#\mathsf C})$ consists of polygons at pairwise positive distances; $|\mathsf C|=\sum_i|C_i|$ and $W_{\mathsf C}=\prod_iW_{C_i}$. For a polygon or a family, $n_{\mathsf C}$ is the number of pieces, $s_{\mathsf C}$ the least length of a piece, $\delta_{\mathsf C}$ the least distance of two different parallel pieces of one polygon, $\mathfrak s_{\mathsf C}=n_{\mathsf C}^2+|\mathsf C|^2/\delta_{\mathsf C}^2$, $\mathfrak d_{\mathsf C}=\min\{s_{\mathsf C},\delta_{\mathsf C}\}$, $c_{\mathsf C}\in(0,1]$ a chord-arc constant ($|x-y|\ge c_{\mathsf C}\operatorname{arc}(x,y)$ for $x,y$ on one $C_i$, $\operatorname{arc}$ the shorter arc), $\delta'_{\mathsf C}$ the least sup-distance of two different polygons ($\infty$ for one polygon), and $\rho_{\mathsf C}>0$ the least $|x-y|_\infty$ over the pairs of points that lie on two different polygons, or on one polygon with both arcs between them of length $\ge\frac12s_{\mathsf C}$.

We put $k=\frac3{4\pi}$, $\kappa=k|\mathsf C|/\mathsf b$ and $\mathcal W_K(\mathsf C)=a_K^{-\kappa}\langle W_{\mathsf C}\rangle_K$; for a square $n_{\mathsf C}=4$ and one may take $c_{\mathsf C}=\frac12$, numbers that dilations do not move.

\emph{Numbers of the free field.} By this we mean numbers formed from the Gaussian field alone, in which no measure of the interacting theory enters. For loops $C$, $D$ with corners in $\Lambda_K$, $\mathcal I_a(C,D)=a\sum_\mu\sum_{y,y'}j^C_\mu(y)\,g_N(\frac{y-y'}a)\,j^D_\mu(y')$ is the lattice Coulomb energy at the spacing $a=a_K$, $N=\ell/a$, $j^C$ the current of $C$ and $g_N$ the Green function of the lattice Laplacian on $\mathbb Z^3/N\mathbb Z^3$ without its zero mode. $G_\ell$ is the Green function of $-\Delta$ on $T_\ell$ without its zero mode, $\mathcal I_\ell(C,D)=\oint_C\oint_DG_\ell(x-y)\,dx\cdot dy$ for loops at positive distance, $\mathcal J_\ell(C)$ the limit as $u\to0$ of the part $|x-y|>u$ of $\oint_C\oint_CG_\ell(x-y)\,dx\cdot dy$ less $\frac{|C|}{2\pi}\log\frac1u$, and $\mathcal J(C)$ the same with $\frac1{4\pi|x|}$ in the place of $G_\ell$; ${\rm Lg}(s)=(1+\log(e+s))^{3/2}$. With Euler's constant $\gamma_E$, $S(z)=\frac{2\sin(z/2)}z$, $p=q+2\pi n$, $\mathsf N=\sum_{n\in\mathbb Z^2}\frac{S(p_2)S(p_3)}{p_2^2+p_3^2}$ and $\mathsf D$ the same with $S^2$ in the place of $S$:
\[
\begin{multlined}
c_{\rm lat}=\frac{5\log2+2\gamma_E}{4\pi}=0.36766\ldots,\quad c_r=\frac1{4\pi^2}\int_{[-\pi,\pi]^2}\Bigl[\frac{\mathsf N^2}{\mathsf D}-\frac1{4\sin^2\frac{q_2}2+4\sin^2\frac{q_3}2}\Bigr]dq_2\,dq_3= \\
-0.03149\ldots,\quad c_Q=c_{\rm lat}+c_r=0.33617\ldots
\end{multlined}
\]

\emph{Action density; mirrors.} $\mathcal O(f)=a_K^{-1}\sum_pf(x(p))a_p$ for real Lipschitz $f$, $x(p)$ the centre of $p$, and $\mathcal W_K(\mathsf C;f_1,\dots,f_n)=a_K^{-\kappa}\langle W_{\mathsf C}\prod_q(\mathcal O(f_q)-\langle\mathcal O(f_q)\rangle_K)\rangle_K$. $S^{\sf T}_{m,K}(f_1,\dots,f_m)$ is the joint cumulant of $\mathcal O(f_1),\dots,\mathcal O(f_m)$ under $\langle\cdot\rangle_K$, the truncated $m$-point function of the action density; for pairwise disjoint supports its limit $S_m(f_1,\dots,f_m)$ exists along the full sequence \YMMref{(2212)}{813}. A coordinate mirror is $\theta x=x-2(x_\mu-c)e_\mu$ with $c\in\frac12a_{i_0}\mathbb Z$; it fixes the planes $x_\mu=c$ and $x_\mu=c+\frac\ell2$, it is a site mirror at the cutoff $K$ if $c\in a_K\mathbb Z$ and a bond mirror otherwise, and $S^+_\theta=\{c\le x_\mu\le c+\frac\ell2\}$. It acts on functions of the field by $(\theta F)(U)=F(\theta U)$ and on test functions by $(\theta f)(x)=f(\theta x)$, so that $\theta W_C=W_{\theta C}$ and $\theta\mathcal O(f)=\mathcal O(\theta f)$.

\emph{Three free photon fields.} $X_C$ is the exact bare $2$-cochain with $d^{\sf T}X_C=j_C$; in lattice units $a_K(X_C,X_D)=\mathcal I_{a_K}(C,D)$. The three free photon fields at the label $\mathsf b$ are the curvatures $F^a=\beta^{-1/2}\Pi\xi^a$, $a=1,2,3$, $\xi^a$ independent standard Gaussian variables on the bare plaquettes and $\Pi$ the orthogonal projection onto the exact $2$-cochains; $\phi^a_C=(F^a,X_C)$ is the flux of $F^a$ through $C$, so that $\mathbb E\,\phi^a_C\phi^{a'}_D=\delta_{aa'}\mathcal I_{a_K}(C,D)/\mathsf b$; their loop is $W^{\rm fr}_C=\prod_{a=1}^3\cos\phi^a_C$, their action density $\mathcal O^{\rm fr}(f)=a_K^{-1}\sum_pf(x(p))\,\frac12\sum_a(F^a(p))^2$, and their block curvatures $F^{(d)}$ are the curvatures of the block averages at the depth $d$, linear images of $F$.

\emph{The frame of the renormalisation group} (Theorem~\ref{C-372-4} and Theorem~\ref{C-372-3}(5)). A reference coupling $g\le\gamma$ with $g^{-2}\ge4\bar B$, $\log g^{-2}\ge\max\{6\mathfrak p_0,A_0^{-1/\mathfrak p_0},(1+\log LM)^{1/2}\}$ and $g(\log g^{-2})^{\mathfrak p_0}\le c_{\rm th}$, $c_{\rm th}>0$ formed from the constants (true for small $g$); the flow bound $B_\sharp$ of the run and $\bar B=\frac L{L-1}B_\sharp$; $\mathsf I(g)=(g^{-2}+\bar B,L(g^{-2}+B_\sharp)-\bar B]$ (Section~\YMMref{74}{614}); a side $\ell\in\mathfrak L(g)$, the set of the sides at $g$ (equation~\YMMref{(1528)}{615}); a compact $\mathsf J\subset\mathsf I(g)$, $\mathsf b\in\mathsf J$, $\mathsf b_-=\min\mathsf J$, $\mathsf b_+=\max\mathsf J$.

The run is $x_0=\beta$, $x_t=x_{t-1}/L-u_t$ with $|u_t|\le B_\sharp$, so that $|x_t-L^{K-t}\mathsf b|\le\bar B$ and $\frac23x_{[d]}\le x_{K-d}\le x_{[d]}+\bar B$; $g_t=x_t^{-1/2}$, $\varepsilon_t=\varepsilon(x_t)$, $\mathfrak b_t=A_1(\log x_t)^{\mathfrak p_0}$ and $\delta_t=g_t\mathfrak b_t$. The level $t$ is the lattice $\Lambda_{K-t}$, with the tame set $\mathsf O_t=\{\varpi<2\varepsilon_t\}$. Theorem~\YMMref{D}{8}, the inductive description of the effective densities, with the actual-history conclusion of Corollary~\YMMref{69.4}{604}, provides the following objects. Finite positive Borel measures $m_t$ on $\mathcal V_{K-t}$, carried by $\mathsf O_t$ for $t\ge1$, with $m_0=e^{-\beta\sum_pa_p}dU$; bounded Borel functions $\mathsf h_t,\mathsf h^{\rm w}_t\ge0$ on $\mathcal V_{K-t+1}$, the weights of the step $t$ without and with a new large field, with $\int G\,dm_t=\int(G\circ\mathcal M)\,\mathsf h_t\,dm_{t-1}$; for every bounded Borel $F$ on $\mathcal V_K$ quotients $\Phi^{(t)}[F]$ with $\Phi^{(0)}[F]=F$, $|\Phi^{(t)}[F]|\le\sup|F|$ and $\int G\,\Phi^{(t)}[F]\,dm_t=\int(G\circ\mathcal M)\,\Phi^{(t-1)}[F]\,\mathsf h_t\,dm_{t-1}$ for bounded Borel $G$, out of which functions of coarser block fields pull, $\Phi^{(t)}[F\cdot(G\circ\mathcal M^{t'})]=\Phi^{(t)}[F]\cdot(G\circ\mathcal M^{t'-t})$ for $t\le t'$; and, for gauge invariant $F$ and every last level $0\le\bar t\le K$,
\begin{equation}\label{C-372-2}
\begin{multlined}
Z\langle F\rangle_K=\int\Phi^{(\bar t)}[F]\,dm_{\bar t}+\sum_{t\le\bar t}\int\Phi^{(t-1)}[F]\,\mathsf h^{\rm w}_t\,dm_{t-1},\qquad\mathsf w_t:=\frac1Z\int\mathsf h^{\rm w}_t\,dm_{t-1} \\
\le\mathsf V_\ell C_{\rm r}(p)L^{-p(K-t)}\ \relax \ \relax (p\ge0),
\end{multlined}
\end{equation}
the last inequality being the \emph{rarity bound}; it gives $\sum_{t'\le t}L^{p(K-t')}\mathsf w_{t'}\le2\mathsf V_\ell C_{\rm r}(p+1)L^{-(K-t)}$. Here $C_{\rm r}(p)$ is formed from the constants and $p$, and $\mathsf V_\ell$, free of $K$, is at most $2M^{-3}\ell^3e^{B\ell^3}$ on every side of $\mathfrak L(g)$, with $B$ formed from the constants and $g$, and equals $1$ on the least sides of $\mathfrak L(g)$ when $g$ is small. We take for $\mathfrak{su}(2)$ the real space of the Hermitian traceless $2\times2$ matrices with the scalar product $\langle X,Y\rangle=\frac12\operatorname{tr}XY$, whose norm $|X|$ is the operator norm on it; $t_1,t_2,t_3$ is an orthonormal basis, with $t_at_b+t_bt_a=2\delta_{ab}\mathbf 1$ and $\sum_at_a^2=3\cdot\mathbf 1$, the components of $X$ are $X^a=\langle t_a,X\rangle$, a complex $2\times2$ matrix is written $x=x^0\mathbf 1+\sum_ax^a\,it_a$, and ${\rm R}(h)\in SO(3)$, ${\rm R}(h)_{ab}=\langle t_a,ht_bh^{-1}\rangle$, is the rotation of $h\in\mathrm{SU}(2)$.

For $W\in\mathsf O_{t+1}$ there are a chart $B\mapsto V_W(B)$, from the real $\mathfrak{su}(2)$-valued $1$-forms $B$ on a set of bonds of $\Lambda_{K-t}$ (the fluctuation field of the step) with $|B|<\mathfrak b_t$ at every bond into $\mathcal V_{K-t}$, with $\mathcal MV_W(B)=W$ and $V_W(B)(b)=e^{ig_tB(b)}V_W(0)(b)$ at the fluctuation bonds, and the kernel $\nu_W$ of one step, a complex measure of mass $1$ on $\{|B|<\mathfrak b_t\}$, with the \emph{disintegration}
\[\int H(\mathcal MV)\,G(V)\,\mathsf h_{t+1}(V)\,m_t(dV)=\int H(W)\Bigl[\int G(V_W(B))\,\nu_W(dB)\Bigr]m_{t+1}(dW)\]
for bounded Borel $G$, $H$; so $\Phi^{(t+1)}[F](W)=\int\Phi^{(t)}[F](V_W(B))\,\nu_W(dB)$ for $m_{t+1}$-almost every $W$. On $\{|B|<\mathfrak b_t\}$ the kernel is $\nu_W=\Xi_W^{-1}e^{\mathsf E_W}\mu_{C_W}$: the Gaussian measure $\mu_{C_W}$ of the covariance $C_W$ of the step at the background triple of $W$, which is formed from a background $U$ over $W$ with $\mathcal M^tU$ in the axial gauge and from the current of $U$ (the triple defined before equation~\YMMref{(827)}{395} and Section~\ref{sec:normalized-step-measure}), times $e^{\mathsf E_W}$, $\mathsf E_W$ the sum of two local potentials, the interaction and the part of the bare array of couplings, each with a norm with quadratic growth at most $c\,\delta_t$, $c$ formed from the constants, and $\Xi_W$ the number that makes the mass $1$ (\eqref{eq:normalized-step-measure} and Proposition~\ref{prop:normalized-step-kernel}). $W^{(t)}_{\mathsf C}(V)=W_{\mathsf C}(U)$ for a background $U$ over $V\in\mathcal V_{K-t}$, and $U_+$ denotes a background over $W\in\mathsf O_{t+1}$. $\Pi_n$ is the orthogonal projection of the bare $2$-cochains onto the linearised curvatures of the backgrounds of $n$ block steps ($\Pi_0$: onto the exact ones), $\mathcal C_{(n)}=\Pi_{n-1}-\Pi_n$, and
\[\lambda_n=\prod_i\exp\Bigl[-\frac3{2\mathsf b}\,a_K(X_{C_i},\mathcal C_{(n)}X_{C_i})\Bigr]\in(0,1],\qquad\Xi_t=\prod_{s\le t}\lambda_s,\qquad v_{K,d}(C)=a_K(X_C,\Pi_{K-d}X_C).\]

\subsection{Main results}

\begin{theorem}[Smoothed loops along almost every thread]\label{C-372-3}
Assume the one-torus setting.
\begin{enumerate}
\item[(1)] \emph{The starting depth.} $\tau$ is a Borel function on $\mathcal X$ with values in $\{d_\Theta,d_\Theta+1,\dots\}\cup\{\infty\}$; for every integer $D\ge d_\Theta$, $\{\tau\le D\}=\bigcap_{d\ge D}\pi_d^{-1}(I_d)$, a countable intersection of open sets; $\Omega$ is Borel and $\mu(\Omega)=1$. For $D\ge d_\Theta$: $\mu(\tau>D)\le\sum_{d\ge D}R_d$, which is at most $\frac L{L-1}R_D$ if $c'A_0^2\log x_{[D]}\ge2$; both bounds exceed $1$ as long as $c'A_0^2(\log x_{[D]})^{2\mathfrak p_0}\le26\ell^3$. For $D\ge\max\{d_\Theta,s_\natural\}$: $\mu(\tau>D)\le\sum_{d\ge D}R^\natural_d$, which is at most $\frac L{L-1}R^\natural_D$ if $c^\natural A_0^2\log x_{[D]}\ge2$; and $R^\natural_D\le\eta\in(0,1]$ as soon as
\[D\log L\ge\max\Bigl\{\frac6{c^\natural A_0^2},\Bigl[\frac{2(3\log\ell+\log\frac1\eta)}{c^\natural A_0^2}\Bigr]^{1/2}\Bigr\}.\]
So $\tau$ is at most of the order $(\log\ell)^{1/2}$ on events of probability at least $1-\frac L{L-1}\eta$.
\item[(2)] \emph{Loops from the starting depth on.} If $K>d$ and $\varpi(W)<\bar\varepsilon$, $\bar\varepsilon$ the radius of Theorem~\YMMref{E}{10} in the setting of Section~\YMMref{52.7}{432}, then $\mathcal C_K(W)$ is one non-empty orbit of $\mathcal G^d_K$; and $\{\varpi<2c_\rho\varepsilon_{[d]}\}\subset\{\varpi<\bar\varepsilon\}$ if $x_{[d]}\ge t_\natural$. Hence for $d\ge\tau_0(V_\bullet)$ and $K\ge d+1$ the backgrounds over $V_d$ form one non-empty orbit, and $w^K_C(V_d)$ is one number for every closed path $C$ of $\Lambda_{d'}$, $d'\le K$. If $d\ge\tau(V_\bullet)$, then $d$ is admitted, $V_d\in I_d\subset J_d$, and for every lattice loop $C$ of $\Lambda_{d'}$, $d'\ge d$,
\[
\begin{multlined}
\lim_{K\to\infty}w^K_C(V_d)=W^{[d]}_C(V_d)=\lim_{n\to\infty}W_C(H^\infty_{d,n}(V_d)),\qquad|w^K_C(V_d)-W^{[d]}_C(V_d)| \\
\le|C|L^d\mathsf c_{\rm f}L^{-(K-d)/3}\quad(K\ge K_2),
\end{multlined}
\]
with the numbers $\mathsf c_{\rm f}$, $K_2$ of Proposition~\ref{C-372-9}. For a curve the same holds with $w^K_\gamma$, $W^{[d]}_\gamma$ and the bound $(|\gamma|L^d+2)\mathsf c_{\rm g}L^{-(K-d)/3}$, and $W^{[d]}_\gamma$ does not depend on the shadows. At these $V_d$ the smoothed loop, as a function on $J_d$, is continuous and gauge invariant with values in $[-1,1]$, and $0\le1-W^{[d]}_C(V_d)\le\frac12(2B_3\varpi(V_d)L^{2d}s_1s_2)^2$ if $C$ bounds a coordinate rectangle with sides $s_1$, $s_2$. $\Omega$ depends on none of $d$, $C$, $\gamma$, $K$. The statements at finite cutoffs hold at every depth $d\ge\tau_0$, the limits as $K\to\infty$ at every depth $d\ge\tau$. No step of (2) and (3) uses the thread property $V_d=\mathcal M(V_{d+1})$: they hold at every point of $\{\tau<\infty\}\subset\mathcal X$.
\item[(3)] \emph{Random variables.} Let $d$ be admitted and $C$ a lattice loop of $\Lambda_{d'}$, $d'\ge d$. Then $Y_{d,C}$ is Borel on $\mathcal X$, $[-1,1]$-valued, $\sigma(\pi_d)$-measurable, $\theta_v$-invariant, and continuous on the open set $\pi_d^{-1}(J_d)$, where it is $W^{[d]}_C(V_d)$; for $K\ge\max\{2,d+1,d'\}$, $Y^K_{d,C}$ is Borel and $\sup_{\mathcal X}|Y^K_{d,C}-Y_{d,C}|\to0$; $\{\tau\le d\}$ is Borel, $\theta_v$-invariant and contained in $\pi_d^{-1}(I_d)$, and $Y_{d,C}=\lim_Kw^K_C(V_d)$ on it. With $Y_{d,C}:=Y_{d,C^{\langle d\rangle}}$ for $d>d'$, $(Y_{d,C})_{d\ge d_\Theta}$ is an $(\mathcal F_d)$-adapted sequence with values in $[-1,1]$ on $(\mathcal X,\mu)$ whose members of index $\ge\tau$ are, on $\Omega$, the smoothed loops along $C$. The same holds for rectifiable closed curves, with one shadow fixed on each $\Lambda_j$, with $F_\gamma\circ\pi_d$ and $w^K_\gamma$, $K\ge\max\{2,d+1\}$.
\item[(4)] \emph{Means.} For admitted $d$, 
\[
\int Y_{d,C}\,d\mu=\int F^{[d]}_C\,d\mu_d=\lim_K\langle F^{[d]}_C(\mathcal M^{K-d}U)\rangle_K=\lim_K\langle(\mathbf 1_{J_d}w^K_C)(\mathcal M^{K-d}U)\rangle_K,
\]
with the rates of Proposition~\ref{C-372-9}(d); the same holds for curves. For $d\ge d_\Theta$, 
\[
|\int Y_{d,C}\,d\mu-\int\mathbf 1\{\tau\le d\}Y_{d,C}\,d\mu|\le\sum_{d'\ge d}R_{d'},
\]
which is at most $\frac L{L-1}R_d$ if $c'A_0^2\log x_{[d]}\ge2$; if $G_d$ is Borel on $\mathcal V_d$ with $|G_d|\le1$ and $G_d=W^{[d]}_C$ on $J_d$, then $G_d(V_d)=Y_{d,C}$ for all $d\ge\tau$ and $\int|G_d(V_d)-Y_{d,C}|\,d\mu\le R_d$. For admitted depths $d_1,\dots,d_n$ and loops $C_i$ of depths $\ge d_i$, $\langle\prod_iF^{[d_i]}_{C_i}(\mathcal M^{K-d_i}U)\rangle_K\to\int\prod_iY_{d_i,C_i}\,d\mu$.
\item[(5)] \emph{Locality.} In the frame, let $d$ be admitted and at least the depth of the corners of $C$, $M_1a_d\le\frac{\delta_1}8$, $K>d$ and $2\varepsilon_{K-d}\le\epsilon_{\rm lc}$; let $\Omega'$ be a slab of the depth $d$ with faces in $M_1a_d(\mathbb Z+\frac12)$, or the intersection of two such slabs, with ${\rm dist}(C,T_\ell\setminus\Omega')\ge\frac{\delta_1}2$, and let $Y_{K,d}(V)=W_C(U')$, $U'$ the localised background formed from the bonds of $V$ in $\Omega'$ ($M_1$, $\delta_{\rm lc}$, $\mathsf C_{\rm lc}$, $\epsilon_{\rm lc}$ are the constants of the localisation of backgrounds, Proposition~\YMMref{53.5}{447} and Section~\YMMref{53.4}{445}). Then for every $V\in\mathsf O_{K-d}$
\[|w^K_C(V)-Y_{K,d}(V)|\le2\mathsf C_{\rm lc}\hat\varepsilon_d|C|L^de^{-\delta_{\rm lc}\delta_1L^d/8},\qquad\hat\varepsilon_d=A_0(\tfrac23x_{[d]})^{-1/2}\log^{\mathfrak p_0}(x_{[d]}+\bar B),\]
a bound free of $K$ that falls faster than every power of $a_d$; so two fields of $J_d$ that lie in $\mathsf O_{K-d}$ for all large $K$ and agree on the bonds of $\Omega'$ have smoothed loops within twice this bound.
\item[(6)] \emph{Every label.} Let $\mathsf b>0$, $\ell=2L^m$ with $m\in\mathbb N$, and let $s\ge0$ be an integer with $L^s\mathsf b\ge\gamma^{-2}$ and $M\,Q(L^s\mathsf b)\le L^{m+s}$ (Lemma~\ref{C-372-7}). Then (1)--(4) hold at $(\ell,\mathsf b)$ for the depths $d\ge s$, the objects of the depth $d$ being those of $(L^s\ell,L^s\mathsf b)$ at the depth $d-s$, carried back by $x\mapsto L^{-s}x$, with $R_{d-s}[L^s\ell,L^s\mathsf b]$ in the place of $R_d$ and with the depths $d_\natural$, $d_\Theta$, $s_\natural$ of $(L^s\ell,L^s\mathsf b)$ increased by $s$; the backgrounds so obtained are backgrounds of the lattices of side $\ell$. Indeed, $x\mapsto L^{-s}x$ maps $\Lambda_{d-s}$ of side $L^s\ell$ onto $\Lambda_d$ of side $\ell$, one graph with $2L^{m+d}$ sites in each direction, and the cutoff $K-s$ onto the cutoff $K$, with the same number $K-d$ of block steps above the depth $d$; it carries the Wilson measure at $(K-s,L^s\ell,L^s\mathsf b)$ onto the one at $(K,\ell,\mathsf b)$, as $\beta=L^{K-s}(L^s\mathsf b)$.

The block averaging map with its averaging radius $1/200L^2$ (equations~\YMMref{(37)--(39)}{20}); the constraint sets, the action $\sum_pa_p$, criticality and the regularity at the radius $\varepsilon_{\rm v}$, measured in the units of the level, with the bare spacing $L^{-(K-d)}$ (Definitions~\YMMref{9.3}{40} and~\YMMref{9.5}{41}); the residual gauge groups and the minimal orbit (Section~\ref{sec:analytic-core-complements}); the axial gauge of the averaged backgrounds (Section~\ref{sec:sf-chart-interface}); the radius $\bar\varepsilon$ and the lattices of Theorem~\YMMref{E}{10} in the setting of Section~\YMMref{52.7}{432}, whose lattice hypotheses are imposed through the number of the sites in each direction at the levels used; and the plaquette size $\varpi$: all are formed from these graphs, from $K-d$ and from the constants, and hold neither the side, nor the label, nor the spacing.

And $x_{[d]}=L^d\mathsf b=L^{d-s}(L^s\mathsf b)$, so that $\varepsilon_{[d]}$, $I_d$ and $J_d$ at $(\ell,\mathsf b)$ are the sets of $(L^s\ell,L^s\mathsf b)$ at the depth $d-s$. So $\mathcal C_K(W)$, $H^K_d$, $w^K_C$ and $W^{[d]}_C$ at $(\ell,\mathsf b)$ are the carried objects.
\end{enumerate}
\end{theorem}

\begin{theorem}[The renormalised sharp loop]\label{C-372-4}
Work in the frame; $\mathsf C$ is a family of polygons and $\kappa=k|\mathsf C|/\mathsf b$.
\begin{enumerate}
\item[(1)] \emph{Ceiling.} If $\kappa<\frac12$, that is $|\mathsf C|<\frac{2\pi}3\mathsf b=2.094\ldots\mathsf b$, then $\sup_{K>s_0}|\mathcal W_K(\mathsf C)|<\infty$; precisely, for $\vartheta\in(\kappa,\frac12)$ and $K>d\ge s_0$,
\begin{equation}\label{C-372-5}
\Bigl|\mathcal W_K(\mathsf C)-a_K^{-\kappa}\Xi_{K-d}\frac1Z\int W^{(K-d)}_{\mathsf C}\,dm_{K-d}\Bigr|\le\mathsf E'_d:=\epsilon_1L^{-(\vartheta-\kappa)(d+1)}+(c_\Xi L^\kappa+\epsilon_1)\,2\mathsf V_\ell C_{\rm r}(\kappa+1)L^{-d},
\end{equation}
where $\epsilon_1=c_\Xi C_\eta/(1-L^{\kappa-\vartheta})$, $C_\eta$ and $s_0$ are the numbers of Proposition~\ref{C-372-10}(a) at the exponent $\vartheta$, and $c_\Xi$ is the number of Proposition~\ref{C-372-11}(c). So $\sup_{K>s_0}|\mathcal W_K(\mathsf C)|\le c_\Xi L^{\kappa s_0}+\mathsf E'_{s_0}$, and $|\Phi^{(t)}[W_{\mathsf C}]|\le(c_\Xi+\epsilon_1)L^{-\kappa t}$ $m_t$-almost everywhere for $t\le K-s_0$. On tori with $\ell\le1.48\,\mathsf b$ every coordinate rectangle has $\kappa<\frac12$. If $\kappa\ge\frac12$, then $|\langle W_{\mathsf C}\rangle_K|\le c_\vartheta a_K^\vartheta$ for every $\vartheta<\frac12$.
\item[(2)] \emph{The limit.} If $\kappa<\frac12$, then $\mathcal W(\mathsf C)=\lim_K\mathcal W_K(\mathsf C)$ exists along the full sequence of cutoffs. Let $F^{[d]}_{\mathsf C}=\prod_iW^{[d]}_{C_i}$ on $J_d$ and $0$ off $J_d$, and
\[\mathsf l_d=\lim_Ka_K^{-\kappa}\Xi_{K-d}=a_d^{-\kappa}\exp\Bigl[-\frac3{2\mathsf b}\bigl((c_{\rm lat}-c_Q)|\mathsf C|+O_{\mathsf C,\ell}(a_d(1+\log^+(1/a_d)))\bigr)\Bigr].\]
Then for all $d\ge\max\{s_0,d_0\}$, $d_0$ the depth of Proposition~\ref{C-372-9}(e),
\[\Bigl|\mathcal W(\mathsf C)-\mathsf l_d\int F^{[d]}_{\mathsf C}\,d\mu_d\Bigr|\le\mathsf E_d:=\mathsf E'_d+2c_\Xi\mathsf V_\ell C_{\rm r}(\kappa+1)L^{-d}+2c_\Xi L^{\kappa d}R_d\to0,\]
and so $|a_d^{-\kappa}\int F^{[d]}_{\mathsf C}\,d\mu_d-\mathcal W(\mathsf C)\,e^{\frac3{2\mathsf b}(c_{\rm lat}-c_Q)|\mathsf C|}|\le c\,(\mathsf E_d+a_d(1+\log^+(1/a_d))/\mathsf b)$, whether or not $\mathcal W(\mathsf C)$ vanishes; here $c_{\rm lat}-c_Q=0.03149\ldots$ and $\frac32(c_{\rm lat}-c_Q)=0.047235\ldots$. One power of the spacing renormalises; no further divergent counterterm is needed. The ball on which $F^{[d]}_{\mathsf C}$ is formed may be taken at the label $\mathsf b$ or at the label $\mathsf b-\bar B$: the two means differ by at most $R_d$. If $k|\mathsf C|/\mathsf b_-<\frac12$, then $\mathcal W_K\to\mathcal W$ uniformly on $\mathsf J$, $\mathcal W$ is continuous on $\mathsf J$, and $\mathcal W_K$ formed at labels $\mathsf b_K\in\mathsf J$ with $\mathsf b_K\to\mathsf b$, $\mathsf b_K$ in the exponent too, tends to $\mathcal W$ at $\mathsf b$.
\item[(3)] \emph{Every torus, every label.} For every integer $m$, $\ell=2L^m$, every $\mathsf b>0$ and every family with $\kappa<\frac12$ the limit $\mathcal W^{(\ell,\mathsf b)}(\mathsf C)$ exists along the full sequence, and $\mathcal W^{(L^z\ell,L^z\mathsf b)}(L^z\mathsf C)=L^{-z\kappa}\mathcal W^{(\ell,\mathsf b)}(\mathsf C)$ for every integer $z$.
\item[(4)] \emph{Positivity; perimeter law.} Wherever they exist, the limits have the properties of limits listed in Proposition~\ref{C-372-13}. For every coordinate rectangle $C$ with $\kappa<\frac12$ there are $0<c\le c_2$ with $c\le\mathcal W_K(C)\le c_2$ for all large $K$, uniformly in $\mathsf b\in\mathsf J$ if $k|C|/\mathsf b_-<\frac12$; so $\mathcal W(C)>0$. For rectangles $C_i=n_1a_i\times n_2a_i$ with corners in $\Lambda_i$, $n_1$, $n_2$ fixed, $d_i=\lceil(\frac43-\frac u2)i\rceil$ with $u\in(0,\frac13)$, and a fixed $\vartheta\in(0,\frac12)$, let $s_0(i)=\max\{i,s'_0\}$ and $C_\eta(i)=C(\mathsf J)L^{\vartheta i}$ be the numbers of the one-step bound linear in $|C_i|/\mathsf b$ (Proposition~\ref{C-372-10}(b)) for $C_i$, $C(\mathsf J)$ depending on the constants, $\vartheta$, $\mathsf J$, $n_1$ and $n_2$, and let $i$ be so large that $(n_1^2+n_2^2)a_i^2<\ell^2/4$ (so that $C_i$ is a polygon), $\kappa(C_i)\le k|C_i|/\mathsf b_-<\vartheta$, $s_0(i)\le d_i$ and $\theta_{d_i}\le1$, with the number $\theta_{d_i}$ below, which tends to $0$ as $i\to\infty$ (so that the factor $1-\frac12\theta_{d_i}^2$ is at least $\frac12$). Then
\[e^{-\frac3{2\mathsf b_-}\eta_i}(1-\tfrac12\theta_{d_i}^2)(1-2\mathsf V_\ell C_{\rm r}(1)L^{-d_i})-\mathsf E'_{d_i}\le\mathcal W_K(C_i)\le c_\Xi L^{\kappa d_i}+\mathsf E'_{d_i}\qquad(K>d_i,\ \mathsf b\in\mathsf J),\]
where $\mathsf E'_{d_i}$ is the number of \eqref{C-372-5} formed with $C_\eta(i)$ and $s_0(i)$ in the place of the numbers of Proposition~\ref{C-372-10}(a), $\eta_i=\sup_{K\ge i}|c_{\rm lat}|C_i|+\mathcal J_\ell(C_i)+E_{a_K}(C_i)|$ with $E_a$ the remainder of the lattice Coulomb energy (Proposition~\ref{C-372-11}(a)), and $\theta_{d_i}=2B_3\sup_{K>d_i,\mathsf b\in\mathsf J}\varepsilon_{K-d_i}L^{2d_i}n_1n_2a_i^2$ (Proposition~\ref{C-372-12}(d)); both ends tend to $1$: $\sup\{|\mathcal W_K(C_i)-1|:\mathsf b\in\mathsf J,\ K>d_i\}\to0$ as $i\to\infty$. For the squares $Q_i$ of side $a_i$ with corners in $\Lambda_i$ there are $i_2$ and $n_Q$ with $\mathcal W_K(Q_i)\ge\frac18$ for all $\mathsf b\in\mathsf J$, $i\ge i_2$ and $K>i+n_Q$. On sides with $\mathsf V_\ell>1$ these small rectangles must be so small that $\mathsf V_\ell L^{-i}$ is small. If $s\times T$ and $s\times2T$ are rectangles with $\kappa<\frac12$, then $\liminf_K\mathcal W_K(s\times nT)\ge\mathcal W(s\times2T)[\mathcal W(s\times2T)/\mathcal W(s\times T)]^{n-2}$ for every $n\ge2$ with $nT<\ell/2$.
\item[(5)] \emph{Block fields; a static source.} Let $\kappa<\frac12$ and $H$ be a bounded Borel gauge invariant function on $\mathcal V_{s_1}$. For $K>d'\ge\max\{s_0,s_1\}$,
\[\Bigl|a_K^{-\kappa}\langle W_{\mathsf C}H(\mathcal M^{K-s_1}U)\rangle_K-a_K^{-\kappa}\Xi_{K-d'}\frac1Z\int W^{(K-d')}_{\mathsf C}\,(H\circ\mathcal M^{d'-s_1})\,dm_{K-d'}\Bigr|\le\sup|H|\,\mathsf E'_{d'};\]
so $\lim_Ka_K^{-\kappa}\langle W_{\mathsf C}H(\mathcal M^{K-s_1}U)\rangle_K$ exists along the full sequence and equals $\lim_d\mathsf l_d\int F^{[d]}_{\mathsf C}\,(H\circ\mathcal M^{d-s_1})\,d\mu_d$. For a rectangle $C$ the quotients $\langle W_CH(\mathcal M^{K-s_1}U)\rangle_K/\langle W_C\rangle_K$ converge; their limit $\omega(H)$ is a linear functional on the bounded gauge invariant Borel functions of all block fields, real on the real ones, normalised, $\omega(1)=1$, consistent between the depths, $\omega(H\circ\mathcal M)=\omega(H)$, and bounded at each depth: $|\omega(H)|\le\sup|H|\,(c_\Xi L^{\kappa d'}+\mathsf E'_{d'})/\mathcal W(C)$ with $d'=\max\{s_0,s_1\}$, $\mathcal W(C)>0$ by (4); $c^{-1}a_d^\kappa\le\int Y_{d,C}\,d\mu\le c\,a_d^\kappa$ for large $d$, and $\int H(V_{s_1})Y_{d,C}\,d\mu/\int Y_{d,C}\,d\mu$ tends to $\omega(H)$ as $d\to\infty$. For $H$ bounded Borel but not gauge invariant all of this holds with the gauge average $\bar H(V)=\int H(V^v)\,dv$ (Haar measure on the gauge transformations of $\Lambda_{s_1}$) in the place of $H$, the Wilson measure, $\mu$ and $Y_{d,C}$ being gauge invariant.
\item[(6)] \emph{Mixed functions.} Let a coordinate mirror $\theta$ part $\mathsf C$, at distance $\ge\delta_1>0$ from its planes, from the supports of real Lipschitz $f_1,\dots,f_n$, at distance $\ge\delta>0$ from the planes, and let $2\kappa<\frac12$. Then $Q_K(z)=a_K^{-\kappa}\langle W_{\mathsf C}e^{\mathcal O(w_z)}\rangle_K/\langle e^{\mathcal O(w_z)}\rangle_K$, $w_z=\sum_qz_qf_q$, obeys $|Q_K(z)|\le\sqrt6\,e^{\Delta(\delta)/2}\mathcal W_K(\mathsf C,\theta\mathsf C)^{1/2}$ for $K>s_0$ on a polydisc free of $K$, and converges locally uniformly there to a holomorphic $Q$; the supports need not be disjoint. If they are pairwise disjoint, then $\lim_K\mathcal W_K(\mathsf C;f_1,\dots,f_n)$ exists along the full sequence and equals $\sum_{A\subset\{1,\dots,n\}}\partial_AQ(0)\sum_\pi\prod_{\beta\in\pi}S_{|\beta|}(f_\beta)$, $\pi$ running over the partitions of the complement of $A$ into blocks of at least two elements, $f_\beta=(f_q)_{q\in\beta}$.

The same holds for two orthogonal coordinate mirrors $\theta_1$, $\theta_2$ if $4\kappa<\frac12$, $\mathsf C$ lies on one side of both at distance $\ge\delta_1$ from their planes, and the $f_q$ fall into two groups, the first with supports on the other side of $\theta_1$, the second with supports on the side of $\mathsf C$ of $\theta_1$ and on the other side of $\theta_2$, the two groups at sup-distance $\ge2\delta$ and all supports at distance $\ge\delta$ from the planes. For $n=1$ it holds, if $2\kappa<\frac12$ and $\mathsf C$ has extent $<\frac\ell2$ in each coordinate direction, for every real Lipschitz $f$ whose support misses the least closed coordinate box that holds $\mathsf C$. Where no such mirror parts them, but $N$ equal coordinate boxes between equally spaced walls in $\frac12a_i\mathbb Z$, for an index $i$, hold each polygon in one box, the number of the boxes being even in every direction that the walls cut, so that both fixed planes $c$ and $c+\frac\ell2$ of every wall mirror are walls (on $\ell=2L^m$, $2L^j$ or $4L^j$ slabs in each cut direction), each $f_q$ is supported in one box, the supports lie off the walls and no box meets both a polygon and a support, one has $\langle W_{\mathsf C^{\rm f}_\iota}\rangle_K\ge0$ and $|Q_K(z)|\le M_S\prod_\iota\mathcal W_K(\mathsf C^{\rm f}_\iota)^{1/N}$ for $\sum_q|z_q|\|f_q\|_\sharp<\rho_D/N$, a polydisc free of $K$, with $\|f\|_\sharp$, $\rho_D$ and $M_S$ of Proposition~\ref{C-372-14}, $\mathsf C^{\rm f}_\iota$ being the family of the images under the wall reflections of the polygons $\mathsf C_\iota$ of the box $\iota$, which has exactly one image in every box, with $\kappa(\mathsf C^{\rm f}_\iota)=N\kappa(\mathsf C_\iota)$; so, if $N\kappa(\mathsf C_\iota)<\frac12$ for every box $\iota$, $\sup_K|\partial_AQ_K(0)|<\infty$ for every $A\subset\{1,\dots,n\}$, and this holds whatever the supports off the walls, a function whose support misses the walls being the sum of pieces each supported in one box; and $\sup_K|\mathcal W_K(\mathsf C;f_1,\dots,f_n)|<\infty$ when the supports are pairwise disjoint, the truncated functions of the action density being bounded in $K$ there \YMMref{(2212)}{813}; for $n=1$, $\mathcal W_K(\mathsf C;f_1)=\partial_1Q_K(0)$, with no further condition.

The limits obey the law of dilations $\mathcal W^{(L^z\ell,L^z\mathsf b)}(L^z\mathsf C;f^{(z)})=L^{-z\kappa-zn}\mathcal W^{(\ell,\mathsf b)}(\mathsf C;f)$, $f^{(z)}_q=f_q(L^{-z}\cdot)$, and the positive form of Proposition~\ref{C-372-13}(b) with mixed entries: let $\theta'$ be the coordinate mirror $\theta'x=x-2(x_\mu-\hat c)e_\mu$ with a place $\hat c\in\frac12a_{i_0}\mathbb Z$, let families $G_1,\dots,G_s$ lie in $\{\hat c<x_\mu<\hat c+\frac\ell4\}$ and tuples $f_1,\dots,f_s$ of real Lipschitz functions, each tuple with pairwise disjoint supports, have their supports in $\{\hat c+\frac\ell4<x_\mu<\hat c+\frac\ell2\}$, all at positive distance from the planes $x_\mu\in\hat c+\frac\ell4\mathbb Z$, and let $2[\kappa(G_i)+\kappa(G_j)]<\frac12$ for all $i$, $j$; then each entry $\mathcal W(\theta'G_i,G_j;\theta'f_i,f_j)$ is a limit of the kind above, the coordinate mirror with the place $\hat c+\frac\ell4$, which lies in $\frac12a_{i_0}\mathbb Z$ as $\frac\ell4=\frac12L^m$, parting the family $(\theta'G_i,G_j)$ from the supports of $\theta'f_i$ and $f_j$, and $\sum_{i,j}\bar c_ic_j\mathcal W(\theta'G_i,G_j;\theta'f_i,f_j)\ge0$ for all complex $c_1,\dots,c_s$.
\item[(7)] \emph{The action density beside a loop.} Let $C$ be a polygon with $\psi=\langle W_C\rangle_K>0$ (every coordinate rectangle has it), $f$ bounded and real, and $\theta$ a site mirror with $C$ in one open half torus and every plaquette $p$ with $f(x(p))\ne0$ with its bonds in the other closed half $S$. Write $\langle X;Y\rangle_K=\langle XY\rangle_K-\langle X\rangle_K\langle Y\rangle_K$. At every cutoff, for every bounded real Borel function $H_0$ of the bonds of $S$ and $H=\mathcal O(f)-H_0$: $\langle W_C;W_{\theta C}\rangle_K\ge0$, $\langle\theta H;H\rangle_K\ge0$, and
\[\Bigl|\frac{\langle W_C;\mathcal O(f)\rangle_K}\psi-\frac{\langle W_C;H_0\rangle_K}\psi\Bigr|\le\Bigl[\frac{\langle W_CW_{\theta C}\rangle_K}{\psi^2}-1\Bigr]^{1/2}\langle\theta H;H\rangle_K^{1/2};\]
at $H_0=0$ the left side is $|\langle W_C\mathcal O(f)\rangle_K/\psi-\langle\mathcal O(f)\rangle_K|$. Let moreover, in the frame, $g\le\gamma_W$, the ceiling of equation~\YMMref{(2346)}{850}; let $\vartheta_{\rm sh}\in(0,1)$ be a shape parameter, $\mathsf d\in(0,\mathsf d_0]$ with the threshold $\mathsf d_0=\mathsf d_0(g,\vartheta_{\rm sh})\in(0,\frac1{200}]$ of equation~\YMMref{(2350)}{851}, and $K-s(\mathsf d)\ge1$, $s(\mathsf d)$ the witness depth at $g$ of equation~\YMMref{(2321)}{843}; and let $(\theta f,f)$ be a test pair of the class $\mathfrak T_{T_\ell}(\mathsf d;\vartheta_{\rm sh})$ of Lemma~\ref{lem:ward-test-pairs}, on a torus with $20\mathsf d\le\ell$: $f$ is non-negative, Lipschitz and not $\equiv0$, the supports of $\theta f$ and $f$ lie in closed balls of radius $\mathsf d$ and are at a distance in $[\mathsf d,6\mathsf d]$, $\int f\ge\vartheta_{\rm sh}\|f\|_\infty\mathsf d^3$ and ${\rm Lip}\,f\le\|f\|_\infty/\vartheta_{\rm sh}\mathsf d$.

As $(K,L^K\mathsf b)$ is a first-passage datum at $g$ with the label $\mathsf b\in\mathsf I(g)$ (Corollary~\YMMref{74.2}{616}), on the side $\ell\in\mathfrak L(g)$ of the frame, the law of the truncated two-point function of the action density, Theorem~\textup{\YMMref{AV}{851}}, gives $\langle\theta\mathcal O(f);\mathcal O(f)\rangle_K=\mathsf c_2\mathsf b^{-2}\int\!\!\int f(\theta x)f(y)|x-y|^{-6}\,dx\,dy\,(1+\epsilon')$, with $\mathsf c_2=\frac9{16\pi^2}$ the constant of \YMMref{C}{7}, the double integral the free kernel \YMMref{(2240)}{822} of the test pair, and $|\epsilon'|\le\epsilon_{\rm tp}:=\bar\varepsilon_{g,\vartheta_{\rm sh}}(\mathsf d)$, the tolerance of equation~\YMMref{(2349)}{851}, which is at most $\frac13$ in that theorem; with the bound above and $\mathsf c_2^{1/2}=\frac3{4\pi}$ this gives
\[\Bigl|\frac{\langle W_C\mathcal O(f)\rangle_K}\psi-\langle\mathcal O(f)\rangle_K\Bigr|\le\frac3{4\pi\mathsf b}\Bigl[(1+\epsilon_{\rm tp})\int\!\!\int\frac{|f(\theta x)f(y)|}{|x-y|^6}\,dx\,dy\Bigr]^{1/2}\Bigl[\frac{\langle W_CW_{\theta C}\rangle_K}{\psi^2}-1\Bigr]^{1/2};\]
and, as the free kernel of a test pair is at most $(\frac{4\pi}3)^2\|f\|_\infty^2$ (equation~\YMMref{(2247)}{824}; two balls of radius $\mathsf d$ at a distance $\ge\mathsf d$), the same two-point estimate gives on this whole class $\langle\theta\mathcal O(f);\mathcal O(f)\rangle_K\le\frac43(\frac{4\pi}3)^2\mathsf c_2\mathsf b^{-2}\|f\|_\infty^2=\frac43\mathsf b^{-2}\|f\|_\infty^2$, so that, with no condition beyond this class,
\[\Bigl|\frac{\langle W_C\mathcal O(f)\rangle_K}\psi-\langle\mathcal O(f)\rangle_K\Bigr|\le\frac2{\sqrt3}\,\mathsf b^{-1}\|f\|_\infty\Bigl[\frac{\langle W_CW_{\theta C}\rangle_K}{\psi^2}-1\Bigr]^{1/2}.\]
Besides, at every cutoff,
\[
\begin{multlined}
\frac{\langle W_C\mathcal O(f)\rangle_K}\psi-\langle\mathcal O(f)\rangle_K=-\frac1{a_K}\sum_pf(x(p))\langle W_{\partial p}\rangle_K\Bigl[\frac{\langle W_CW_{\partial p}\rangle_K}{\psi\langle W_{\partial p}\rangle_K}-1\Bigr], \\
\frac d{d\mathsf b}\log\langle W_C\rangle_K=-\Bigl[\frac{\langle W_C\mathcal O(1)\rangle_K}\psi-\langle\mathcal O(1)\rangle_K\Bigr].
\end{multlined}
\]
For a rectangle with $2\kappa<\frac12$ at positive distance from the planes of $\theta$ and Lipschitz $f$ with support at positive distance from them, $\langle W_C;\mathcal O(f)\rangle_K/\psi=\mathcal W_K(C;f)/\mathcal W_K(C)$, $\langle W_CW_{\theta C}\rangle_K/\psi^2=\mathcal W_K(C,\theta C)/\mathcal W_K(C)^2$ and $\langle\theta\mathcal O(f);\mathcal O(f)\rangle_K=S^{\sf T}_{2,K}(\theta f,f)$ converge as $K\to\infty$.
\item[(8)] \emph{Short distances.} Let $0<\vartheta<\frac12$ and $p\ge1$. (a) For rectangles $C_i=n_1a_i\times n_2a_i$ with corners in $\Lambda_i$, $n_1$, $n_2$ fixed, $\varrho_i=\max\{n_1,n_2\}a_i$, $\mathsf n=\max\{L,n_1,n_2\}$, all large $i$ and all $K>i$ with $\Delta_i\le\frac12$,
\[\Bigl|\log\mathcal W_K(C_i)+\frac3{2\mathsf b}\bigl[c_{\rm lat}|C_i|+\mathcal J_\ell(C_i)+E_{a_K}(C_i)\bigr]\Bigr|\le\kappa\log L+c\frac{|C_i|}{\mathsf b}+2\Delta_i,\]
\[\Delta_i=4B_3^2A_0^2\mathsf n^3\frac{\varrho_i}{\mathsf b}\Bigl(1+\log\frac{\mathsf n\mathsf b}{\varrho_i}\Bigr)^{2\mathfrak p_0}+C_1\frac{|C_i|}{\mathsf b}\Bigl(\frac{\varrho_i}{\mathsf b}\Bigr)^{\vartheta/2}+C_3\mathsf V_\ell C_{\rm r}(p)\varrho_i^{p-1/2},\]
$E_a$ the remainder of Proposition~\ref{C-372-11}(a) and $c$, $C_1$, $C_3$ depending on the constants, $\vartheta$, $n_1$, $n_2$ and $\mathsf J$: an error of the order of the main term, which gives the two-sided bound of (4) at these rectangles. (b) For rectangles $C$ with sides in $[c_0\varrho,10\varrho]$ and corners in $\Lambda_i$, $i$ the integer with $a_i\le\frac12c_0\varrho<La_i$, put $a_*=\varrho(\mathsf b/\varrho)^{\vartheta/(3+\vartheta)}$, a length coarser than the loop. At every level $t$ with $a_{K-t}\le a_*$ the mean over the block field of the defect of one step, $\bar{\mathsf D}_{t+1}$ of Proposition~\ref{C-372-12}(b) for $C$, has modulus at most $C_*\frac\varrho{\mathsf b}(\frac{a_{K-t}}{\mathsf b})^\vartheta$, $C_*$ depending on the constants, $c_0$, $\vartheta$ and $\mathsf J$ alone: at the levels finer than the corners, $K-t>s_0=\max\{i,s'_0\}$, by Proposition~\ref{C-372-10}(b), and at the levels with $a_i\le a_{K-t}\le a_*$, about $\frac\vartheta{3+\vartheta}\log_L\frac{\mathsf b}\varrho$ of them, where the loop passes a number of the cubes of the level that is bounded in terms of $L$ and $c_0$, by the argument of one step with this bounded number of segments. Hence, as $\varrho\to0$, with a constant depending on the constants, $c_0$, $\vartheta$, $\mathsf J$ and $\ell$,
\[\log\mathcal W(C)=-\frac3{2\mathsf b}\bigl[\mathcal J(C)+c_{\rm lat}|C|\bigr]+O\Bigl(\Bigl(\frac\varrho{\mathsf b}\Bigr)^{1+\frac{3\vartheta}{3+\vartheta}}\Bigl(1+\log\frac{\mathsf b}\varrho\Bigr)^{2\mathfrak p_0}+\frac{\varrho^4}{\mathsf b\,\ell^3}\Bigr);\]
the exponent $1+\frac{3\vartheta}{3+\vartheta}$ reaches $\frac{10}7$ as $\vartheta\to\frac12$, the length $a_*$ balancing the step errors, $\frac\varrho{\mathsf b}(\frac{a_*}{\mathsf b})^\vartheta$, against the last level, $\frac{\varrho^4}{\mathsf ba_*^3}$, and the remainder, which holds $\log(\mathsf b/\varrho)$ and not $\log\varrho$, is invariant under dilations.
\end{enumerate}
\end{theorem}

\begin{theorem}[Without the renormalisation group]\label{C-372-6}
Let $C$ be a coordinate rectangle on a torus of side $\ell$, $\mathsf b>0$, $a$ the spacing, $\beta=\mathsf b/a\ge e^3$, $\Lambda=\max\{1,|C|/\mathsf b\}$, $\rho=\lfloor(\frac{\beta}{\Lambda\log^3\beta})^{1/4}\rfloor\ge2$, the sides of $C$ at least $(4\rho+1)a$, and the points within $(\rho+1)a$ of $C$ in a cube of side $\le\frac\ell2$. Then
\[|\langle W_C\rangle|\le6e^{3.1\mathsf b/\ell}\,1.58^{|C|/\mathsf b}\Bigl[\Lambda\frac a{\mathsf b}\log^3\frac{\mathsf b}a\Bigr]^{k|C|/4\mathsf b}.\]
With a dyadic $\rho=2^S$ instead, sides at least $(4\rho+1)a$ and $k|C|/\mathsf b<3$, the ceiling of the tube with the profiles of the scales $2,4,\dots,2^S$ reads
\[|\langle W_C\rangle|\le e^{\delta_0+2.2144k|C|/\mathsf b}\rho^{-k|C|/\mathsf b}\Bigl[1+\frac{72\,\bar a\,\rho^3\log^2\rho}{2^{3-k|C|/\mathsf b}-1}\Bigr]+35\,\bar a,\qquad\delta_0=\frac{16k\rho\log\rho+k^2(|C|/\mathsf b)\log^2\rho}\beta,\]
with the plaquette mean $\bar a=\langle a_p\rangle\le\beta^{-1}\log\mathfrak M_\nu$ (Lemma~\ref{C-372-17}(f)), the condition $k|C|/\mathsf b<3$ keeping the denominator positive; at $\rho\in(\frac12\beta^{1/3},\beta^{1/3}]$, where $\bar a\rho^3\log^2\rho$ is at most a multiple of $\log^3\beta$, and $\beta\ge\beta_2$, a threshold depending on $k|C|/\mathsf b$ alone (the term $192\frac{\mathsf b}\ell$ of $\log\mathfrak M_\nu$ is at most $384\mathsf b/|C|$, as $|C|<2\ell$): $|\langle W_C\rangle|\le C(\frac{k|C|}{\mathsf b})(\frac a{\mathsf b})^{|C|/4\pi\mathsf b}\log^3\frac{\mathsf b}a$, as $\rho^{-k|C|/\mathsf b}\le2^{k|C|/\mathsf b}(\frac a{\mathsf b})^{k|C|/3\mathsf b}$ and $\frac k3=\frac1{4\pi}$. The first bound holds in the same form, with a constant depending also on the number of pieces, for polygons whose pieces have at least $4\rho+1$ steps and whose pieces without a common corner are at sup-distance at least $(2\rho+1)a$, for families of such polygons in which also any two pieces of different polygons are at sup-distance at least $(2\rho+1)a$, also with a further factor $G$, $|G|\le1$, a Borel function of the bonds farther than $\rho a$ from the polygons; and for the characters $\chi_j$ of spin $j$, with $\kappa_j=\frac43j(j+1)\le20\log\beta$ and $A=6e^{3.1\mathsf b/\ell}$:
\[|\langle\chi_j(U(C))\rangle|\le(2j+1)\bigl(A\,1.58^{|C|/\mathsf b}\bigr)^{\kappa_j}e^{\kappa_j^2/100}\Bigl[\Lambda\frac a{\mathsf b}\log^3\frac{\mathsf b}a\Bigr]^{\kappa_jk|C|/4\mathsf b}.\]
These bounds hold with $3$ in the place of $6e^{3.1\mathsf b/\ell}$ for every periodic Gibbs state of the Wilson action at $\beta$ on $\mathbb Z^3$ that is invariant under the isometries of $\mathbb Z^3$ and reflection positive in the lattice planes. So, as $a\to0$ with $\mathsf b$ in a compact set, uniformly in $\ell\ge\ell_0$, for polygons: $\langle h(W_C)\rangle\to\sigma(h)=\frac2\pi\int_{-1}^1h(w)\sqrt{1-w^2}\,dw$ for continuous $h$, $\langle W_C^2\rangle\to\frac14$, and for polygons $C_i$ at pairwise positive distances, continuous $h_i$ and Borel $G$, $|G|\le1$, depending only on the bonds at distance $\ge\delta>0$ from them, $\langle\prod_ih_i(W_{C_i})\,G\rangle-\prod_i\sigma(h_i)\langle G\rangle\to0$.

For a coordinate rectangle, whose mean $\langle W_C\rangle$ is positive (Proposition~\ref{C-372-13}(c)), the normalised sharp loop has the variance $\mathrm{Var}(W_C/\langle W_C\rangle)=(1+\langle\chi_1(U(C))\rangle)/(4\langle W_C\rangle^2)-1$, which tends to infinity as $a\to0$, at least like a positive constant times $[\Lambda\frac a{\mathsf b}\log^3\frac{\mathsf b}a]^{-k|C|/2\mathsf b}$; and no multiple $c_aW_C$ of the sharp loop has both a limit law and a mean $c_a\langle W_C\rangle$ that stays away from $0$. In the other direction, $\langle W_C\rangle\ge\langle W_{\partial p}\rangle^{A_C/a^2}\ge\exp[-\frac32\frac{A_C}{a\mathsf b}(192\frac{\mathsf b}\ell+56.6+12\log\beta)]$ for a rectangle of area $A_C$ and sides $<\ell/2$, the second inequality when $192\frac{\mathsf b}\ell+56.6+12\log\beta\le\frac\beta2$; for the Gibbs states above and rectangles with even numbers of steps, $\langle W_C\rangle\ge\exp[-\frac{A_C}{a\mathsf b}(43+9\log\frac{\mathsf b}a)]$ when $28.3+6\log\beta\le\frac\beta2$.

Without the renormalisation group the floor falls like the exponential of the area in units of $a\mathsf b$.
\end{theorem}

\subsection{Construction estimates}

Theorem~\YMMref{D}{8} with Corollary~\YMMref{69.4}{604} on the actual admissible histories enters through the objects of the frame and through several of the following statements, cited by their keys, which we use: the limit laws of the block fields, with one block map for all cutoffs, with their rate and with the two depths of their proof (Corollary~\textup{\ref{IV-62-66}}, \textup{(\ref{IV-62-68})}, Section~\ref{sec:cutoff-calibration}), the numbers $\mathsf C_{\rm bf}$, $K_{\rm bf}$ of the rate being formed at an interval of labels under the condition (Tor) (Definitions~\ref{sec:cutoff-calibration} and~\textup{\ref{IV-62-62}}(5)); the bound on the history without large-field regions (Proposition~\textup{\YMMref{115.2}{806}}, \YMMref{(2187)}{806}); the rarity of large-field regions (Proposition~\ref{prop:ward-fixed-volume-rarity}); the bound at one place at the fixed and at the lowered constants (Lemma~\ref{lem:rarity-chessboard}, Definition~\ref{def:rarity-choices}); the three smallness conditions that define $t_\natural$ and the facts on them (Section~\ref{sec:cutoff-calibration}, Section~\ref{sec:cutoff-calibration}, Section~\ref{sec:cutoff-calibration}); the comparison of two consecutive cutoffs (Proposition~\YMMref{113.9}{795}, \YMMref{(2154)}{795}) with the competitors, the lower bound on their excess, the chart of the fibre at the pivot bonds and the numbers of the comparison (Lemma~\YMMref{113.8}{794} and Proposition~\YMMref{113.5}{792}, the chart bounds in equation~\YMMref{(2138)}{791}, the global hypotheses preceding equation~\YMMref{(2132)}{790} and the constants in and before equation~\YMMref{(2133)}{790}), the averaged backgrounds (Proposition~\YMMref{52.3}{433}, the axial-normalization argument in the proof of Proposition~\YMMref{52.3}{433}), and the response of the background to a change of its datum with the averages of a moved field (Theorem~\YMMref{E}{10} and Lemma~\YMMref{51.6}{398}, with the thresholds of Section~\ref{sec:background-step-complements}, Lemma~\ref{C-368-7}(v)), which give the comparison of the averaged backgrounds of two consecutive cutoffs over one block field and along the tower, and the Lipschitz bound of the tower in the block field (Lemma~\ref{C-372-31}); the existence of the background over a small field as one orbit (Theorem~\YMMref{E}{10} in the setting of Section~\YMMref{52.7}{432}, Section~\ref{sec:analytic-core-complements}); the localisation of backgrounds (Proposition~\YMMref{53.5}{447}, Section~\YMMref{53.4}{445}, including equations~\YMMref{(998)}{446}--\YMMref{(1000)}{446}), with the localised background of a region (equation~\YMMref{(77)}{42}) and the arguments that give its hypotheses for a slab of $M_1$-cubes (the geometric verification in Step~5 of the proof of Lemma~\ref{C-372-17}, using Proposition~\YMMref{44.2}{345}); the covariance of the block averages and of the backgrounds under the lattice isometries (Section~\ref{sec:setup-lattice-symmetries}, Lemma~\ref{lem:setup-lattice-covariance}); for one step, the triples, the analyticity cores, the chart at the trivial pair and the ball about the trivial triple in the core (Section~\ref{sec:analytic-core-complements}, Definition~\YMMref{49.1}{374}, and equation~\YMMref{(933)}{427}), the background triple and its regularity (the triple defined before equation~\YMMref{(827)}{395} and Section~\ref{sec:normalized-step-measure}), the kernel of one step (\eqref{eq:normalized-step-measure} and Proposition~\ref{prop:normalized-step-kernel}), the provision of the step on the whole core (Definition~\YMMref{52.1}{415}, Theorem~\YMMref{W}{423}), the shift of Proposition~\YMMref{51.5}{396} with their brackets (equations~\YMMref{(806)}{392}--\YMMref{(809)}{392} and Lemma~\ref{lem:analytic-rooted-bracket}), the chart and the pivot values at $t=0$ (equation~\YMMref{(804)}{391} and the bounds and covariance proved with it), the potential on a cube of any side (Theorem~\YMMref{F}{199}), the cluster expansion of the step with up to four potential
families (Theorem~\YMMref{R}{361}, Corollary~\YMMref{47.3}{364}
and Proposition~\YMMref{47.4}{365}, with the following covariance
argument), applied to the actual Gaussian structure and the two
potentials of Proposition~\ref{prop:normalized-step-kernel}, the constrained variational equation for critical curves at
the trivial field (the proof of Lemma~\YMMref{124.1}{832}, through
its gauge normalization), and its curvature identification
(equation~\YMMref{(2230)}{819}), the curvature response (the curvature-response construction preceding Lemma~\YMMref{120.2}{819}, the curvature-response calculation preceding Lemma~\YMMref{120.2}{819}), the slice of the free Coulomb energy at the trivial triple (Proposition~\YMMref{120.3}{819}, Proposition~\YMMref{120.3}{819}) and the linearised shift system (Lemma~\YMMref{124.3}{835}, whose proof, with the two changes named in the proof of Proposition~\ref{C-372-10}(a), runs for the objects of the frame on every side of $\mathfrak L(g)$); and the limits of the truncated functions of the action density \YMMref{(2212)}{813}.

Theorems~\YMMref{AJ}{673} and~\YMMref{AD}{607}, the bound for a history with a source on a polydisc uniformly in $K$, is used in this section only in Proposition~\ref{C-372-14}, hence in Theorem~\ref{C-372-4}(6) and in the convergence of the three quantities of Theorem~\ref{C-372-4}(7); the explicit ceiling of Theorem~\ref{C-372-4}(7) uses besides the law of the truncated two-point function of the action density, Theorem~\textup{\YMMref{AV}{851}}, with its test pairs (Lemma~\ref{lem:ward-test-pairs}), the free kernel \YMMref{(2240)}{822} and its bound on a test pair (equation~\YMMref{(2247)}{824}), the constant $\mathsf c_2$ \YMMref{C}{7}, its ceiling $\gamma_W$ equation~\YMMref{(2346)}{850}, its threshold $\mathsf d_0$ equation~\YMMref{(2350)}{851}, its witness depth $s(\mathsf d)$ in equation~\YMMref{(2321)}{843} and its tolerance equation~\YMMref{(2349)}{851}, and the first-passage data (Corollary~\YMMref{74.2}{616}).

The facts on the constants (Proposition~\YMMref{69.1}{595}, the calibration in Section~\ref{sec:cutoff-calibration}, and Corollary~\ref{cor:canonical-running-source}) make fixed constants satisfy the conditions of the setting. Besides these we use standard facts: reflection positivity of the Wilson measure on the torus at site and bond mirrors, the chessboard estimate, Kolmogorov's extension theorem, Vitali's theorem, Schwarz's lemma and Cauchy's estimates. In Theorem~\ref{C-372-3}(1)--(5) and in Propositions~\ref{C-372-8} and~\ref{C-372-9}(a)--(d) one torus and one label are fixed, and the limit is $K\to\infty$, with the limit $n\to\infty$ of the tower inside it; in the frame the numbers of the rate of Proposition~\ref{C-372-9}(d) serve all the labels of $\mathsf J$.

Theorem~\ref{C-372-4} takes besides the limit $d\to\infty$ of means after $K\to\infty$ ((2), (5)), lets the label move in a compact set $\mathsf J$ and along labels $\mathsf b_K\to\mathsf b$ ((2)), lets rectangles shrink, $i\to\infty$ and $\varrho\to0$ ((4), (8)), and moves the torus and the label together by dilations ((3), as Theorem~\ref{C-372-3}(6) does). Theorem~\ref{C-372-6} takes $a\to0$ uniformly in $\ell\ge\ell_0$ with $\mathsf b$ in a compact set, and holds also for periodic Gibbs states on $\mathbb Z^3$. In Theorems~\ref{C-372-3} and~\ref{C-372-4} the side $\ell$ of the torus is fixed, or moves together with the label, by dilations.

\subsection{Auxiliary estimates}

\begin{lemma}\label{C-372-7}
Let $L\ge2$, $m'\in\mathbb Z_{\ge0}$, $r,m\in\mathbb N$, $M=L^{m'}$, $\mathsf b$ real with $\log\mathsf b>1$, and $Q(\mathsf b)=\min\{L^j:j\in\mathbb Z_{\ge0},\ L^j\ge(\log\mathsf b)^r\}$. Then $M\,Q(\mathsf b)\le L^m$ if and only if $L^m\ge M(\log\mathsf b)^r$. The hypothesis $\log\mathsf b>1$ is needed: at $L=3$, $m'=2$, $m=1$, $r=2$, $\mathsf b=1.2$ the second condition holds and the first fails. In the setting it holds since $\gamma\le\frac12$.
\end{lemma}

\begin{proof}[Sketch of proof]
$Q\ge(\log\mathsf b)^r$ gives $M(\log\mathsf b)^r\le MQ\le L^m$. Conversely $L^{m-m'}\ge(\log\mathsf b)^r>1$, so $j=m-m'$ is a positive integer, $L^j$ lies in the set whose least element is $Q$, and $Q\le L^{m-m'}$. In the example $(\log1.2)^2=0.033\ldots$, so $Q=1$ and $MQ=9>3=L^m\ge9\cdot0.033$.
\end{proof}

\begin{proposition}[Laws on the threads; small plaquettes]\label{C-372-8}
In the one-torus setting:
(a) there is exactly one Borel probability measure $\mu$ on $\mathcal X$ whose image under $V_\bullet\mapsto(V_0,\dots,V_D)$ is, for every $D$, the image of $\mu_D$ under $V\mapsto(\mathcal M^DV,\dots,\mathcal MV,V)$; $\mu(\mathsf{Thr})=1$, the image of $\mu$ under $\pi_d$ is $\mu_d$, $\mu$ is $\theta_v$-invariant, and the law of $(\mathcal M^KU,\dots,\mathcal M^{K-D}U)$ under $\langle\cdot\rangle_K$ tends in total variation to the image of $\mu$ under $(\pi_0,\dots,\pi_D)$;
(b) each $I_d$ is open and gauge invariant; for $d\ge d_\natural$ and $K>d$, $\mu^K_d(\mathcal V_d\setminus I_d)\le R_d$ and $\mu_d(\mathcal V_d\setminus I_d)\le R_d$; for $D\ge d_\natural$, $\mu(\tau_0>D)\le\sum_{d\ge D}R_d$, which is at most $\frac L{L-1}R_D$ if $c'A_0^2\log x_{[D]}\ge2$; and $R_d$ tends to $0$ faster than every $L^{-pd}$;
(c) for $d\ge s_\natural$, $K>d$ and every set $S$ of sites of $\Lambda_d$, the $\mu^K_d$- and the $\mu_d$-probability that some plaquette with a corner in $S$ has $|\partial V(p)-\mathbf 1|\ge2\varepsilon_{[d]}$ are at most $\#S\,e^{-c^\natural A_0^2(\log x_{[d]})^{2\mathfrak p_0}}$;
(d) with the constants $C_5,\mathsf C_{\rm maj}\ge0$ of Theorem~\YMMref{D}{8}, $s_d=\frac14\min\{\bar\varepsilon,B_3^{-1}\varepsilon_{[d]}\}$,
\[
\begin{multlined}
\mathsf p^{[d]}=\exp\bigl[-(c_{\rm low}+\tfrac92\log x_{[d]}+\mathsf C_{\rm maj}+2)(\ell L^d)^3\bigr],\qquad\mathsf q^{[d]}=\exp\Bigl[-\Bigl(C_5+\mathsf C_{\rm maj}+2+3\log\frac7{s_d^3} \\
+24B_3^2s_d^2x_{[d]}\Bigr)(\ell L^d)^3\Bigr],
\end{multlined}
\]
one has at every admitted depth $d$ and for $K\ge\max\{2,d+1\}$: $\mu^K_d(J_d)\ge\mu^K_d(I_d)\ge\mathsf q^{[d]}>0$, and $\mu^K_d(I_d)\ge\max\{1-R_d,\mathsf p^{[d]},\mathsf q^{[d]}\}$ if $d\ge1$; the same holds for $\mu_d$. The floors hold no $K$ and fall like $e^{-c(\ell L^d)^3}$.
\end{proposition}

\begin{proof}[Sketch of proof]
(a) By one block map, $\mu_{d+1}\circ\mathcal M^{-1}=\mu_d$ (Corollary~\textup{\ref{IV-62-66}}(c)), so the images of the $\mu_D$ form a projective family of Borel probabilities on the finite products of the compact metric spaces $\mathcal V_d$, which has exactly one extension to $\mathcal X$ (Kolmogorov's theorem). Each marginal is carried by the Borel set $\{V_d=\mathcal MV_{d+1},\ d<D\}$, whence $\mu(\mathsf{Thr})=1$; gauge invariance of the $\mu_d$ and covariance of $\mathcal M$ give the invariance; the joint law at the cutoff $K$ is the image of the single law $\mu^K_D$ under $V\mapsto(\mathcal M^DV,\dots,V)$, and $\mu^K_D\to\mu_D$ in total variation. (b) $\varpi$ is continuous and $\varpi(V^v)=\varpi(V)$.

Here Theorem~\YMMref{D}{8} is applied along the line of couplings $x_{[d']}=L^{d'}\mathsf b$, $0\le d'\le K$ (the run with all $u_t=0$, at the label $\mathsf b$), with the inner ball $I_{d'}=\{\varpi<2\varepsilon_{[d']}\}$ at the depth $d'$: the partition function is a sum over the histories of large-field regions, and the small-field chain of this line (Proposition~\textup{\YMMref{115.2}{806}}), which at the depths $d'>d$ cuts the fields at the radii $c_\rho\varepsilon_{[d']}$, measured on their backgrounds (the definitions preceding Lemma~\YMMref{113.10}{796}), is a bounded Borel weight $\mathfrak w\ge0$ on the bare fields with three properties: the image of $\mathfrak w\,e^{-\beta\sum_pa_p}dU$ under $\mathcal M^{K-d}$ is carried by $I_d$ (clause (a) there); $\int F(1-\mathfrak w)e^{-\beta\sum_pa_p}dU\ge0$ for every gauge invariant bounded Borel $F\ge0$ (clause (b), which places $\mathfrak w$ between the history without large-field regions and the full density); and $\int(1-\mathfrak w)e^{-\beta\sum_pa_p}dU\le R_dZ$, $R_d$ being \YMMref{(2187)}{806} at the depth $d$ and the label $\mathsf b$ (clause (c)).

As $I_d$ is gauge invariant, $\mu^K_d(\mathcal V_d\setminus I_d)=Z^{-1}\int\mathbf 1_{\mathcal V_d\setminus I_d}(\mathcal M^{K-d}U)(1-\mathfrak w)e^{-\beta\sum_pa_p}dU\le R_d$, uniformly in $K$. This chain of the line is another object than the measures $m_t$ of the frame, which follow the run, with $|x_{K-d}-x_{[d]}|\le\bar B$ and the thresholds $2\varepsilon(x_t)$: the tame set $\mathsf O_{K-d}$ that carries $m_{K-d}$ is placed by Proposition~\ref{C-372-9}(e) in the larger ball $J'_d$. The bound for $\mu^K_d$ passes to $\mu_d$ in total variation. For $D\ge d_\natural$, $\{\tau_0>D\}=\bigcup_{d\ge D}\pi_d^{-1}(\mathcal V_d\setminus I_d)$, and the image of $\mu$ under $\pi_d$ is $\mu_d$ by (a); countable subadditivity gives $\mu(\tau_0>D)\le\sum_{d\ge D}R_d$.

As $(y+\log L)^{2\mathfrak p_0}-y^{2\mathfrak p_0}\ge2y\log L$ for $y\ge1$, and $\log x_{[d]}\ge\log\mathsf b>1$, $R_{d+1}/R_d\le L^{3-2c'A_0^2\log x_{[d]}}$, which is at most $L^{-1}$ for all $d\ge D$ if $c'A_0^2\log x_{[D]}\ge2$; then $\sum_{d\ge D}R_d\le\frac L{L-1}R_D$. $R_d$ falls faster than every $L^{-pd}$ because $(\log x_{[d]})^{2\mathfrak p_0}\ge(d\log L)^2$. (c) One witness function of a large plaquette at one place has a mean at most $e^{-c^\natural A_0^2(\log x_{[d]})^{2\mathfrak p_0}}$, free of the volume: Lemma~\ref{lem:rarity-chessboard}, which applies Theorem~\YMMref{D}{8} at the fixed constants and at the lowered choice and frees the bound of the volume by the chessboard estimate for the reflection positive Wilson measure.

Sum over the corners in $S$; the bound passes to $\mu_d$ in total variation. (d) For $\mathsf p^{[d]}$: Theorem~\YMMref{D}{8}, along the line of couplings of (b), bounds, at the level of $\Lambda_d$, the mass of the history without large-field regions below by $e^{E}e^{-(c_{\rm low}+\frac92\log x_{[d]})(\ell L^d)^3}$ and the partition function above by $e^{E}e^{(\mathsf C_{\rm maj}+2)(\ell L^d)^3}$, with one number $E$, and that history is carried by fields whose block field of the depth $d$ lies in the inner ball $I_d$ of the line. For $\mathsf q^{[d]}$: Theorem~\YMMref{D}{8} along the line bounds the density of the block field of the depth $d$ below, as a measure, by $e^{E-C_5(\ell L^d)^3}$ times the small-field indicator times the Boltzmann factor of the background over $V$; on the ball of the fields with $|V(b)-\mathbf 1|<s_d$ at every bond, whose plaquettes are $<4s_d\le\varepsilon_{[d]}$ so that it lies in $I_d$, the integral of this density is at least $\exp[-(3\log\frac7{s_d^3}+24B_3^2s_d^2x_{[d]})(\ell L^d)^3]$; with the upper bound on the partition function this gives $\mathsf q^{[d]}$.

The floors hold no $K$; let $K\to\infty$.
\end{proof}

\begin{lemma}[Two consecutive cutoffs over one block field; the Lipschitz bound of the tower]\label{C-372-31}
In the one-torus setting put $\varepsilon_{\rm r}(s)=2c_\rho\varepsilon_{[s]}$; for two fields $V$, $V'$ on one lattice put ${\rm dist}(V,V')=\max_b\theta(V(b)V'(b)^{-1})$, $\theta(X)\in[0,\pi]$ the angle of $X\in G$, $\cos\theta(X)=\frac12\operatorname{tr}X$; for a real $1$-form $B$ on the bonds of a lattice (its values $B(b)$ hermitian and traceless, $B(\bar b)=-B(b)$) put $\sup|B|=\max_b|B(b)|$ and $(e^{iB}X')(b)=e^{iB(b)}X'(b)$ at the positively oriented bonds $b$. With the chart constant $C_{\rm db}=8L^3$ and the constant $\gamma_{\rm r}\in(0,1]$ of the lower bound on the excess of a competitor (the global hypotheses preceding equation~\YMMref{(2132)}{790}, Proposition~\YMMref{113.5}{792}), the constant $\mathsf K_u$ of the comparison of the background actions at two consecutive cutoffs (equation~\YMMref{(2155)}{795}, \YMMref{(2154)}{795}), the numbers $\varepsilon_{\rm cp}$, $\varrho_{\rm no}=\min\{r_{\rm B},\varrho^0_{\rm R}\}$ and $\mathsf S_{\rm e}$ of the definitions preceding and in equation~\YMMref{(2133)}{790}, the radius $c_4\in(0,1]$ and the positive curvature ceiling $c'_2=c^{\rm g}_2$ of Lemma~\ref{C-368-7}(v), decreasing $c_4$ to $1$ if necessary, and the constant $B^{\rm e}_5\ge8$ of the response of the background to a change of its datum (Section~\ref{sec:background-step-complements}), we put
\[
\begin{multlined}
\mathsf c_{\rm A}=16\,C_{\rm db}\,c_\rho\,(\mathsf K_u/\gamma_{\rm r})^{1/2},\qquad a_s(K)=\mathsf c_{\rm A}\,\varepsilon_{[s]}\,(\ell L^s)^{3/2}L^{-(K-s+1)/2},\qquad r_s(K) \\
=\sum_{K'\ge K}a_s(K')=\frac{a_s(K)}{1-L^{-1/2}}=r_s(s)L^{-(K-s)/2},
\end{multlined}
\]
\[\varepsilon_{\rm L}=\min\{\varepsilon_{\rm cp},\,c'_2/B_3\},\qquad\beta_{\rm L}=\min\Bigl\{1,\ \tfrac12\varrho_{\rm no},\ \frac{c_4}{2B^{\rm e}_5\mathsf S_{\rm e}}\Bigr\},\qquad\mathsf c_9=\frac{2L}{L^{1/2}-1}\,B^{\rm e}_5\mathsf S_{\rm e};\]
$\mathsf c_{\rm A}$, $\varepsilon_{\rm L}$, $\beta_{\rm L}$ and $\mathsf c_9$ depend on the constants alone, $a_s(\cdot)$ decreases, and $\varepsilon_{\rm r}(s)\le\varepsilon_{\rm L}$ at every admitted depth $s$, by the conditions of the comparison of two consecutive cutoffs that define $\Theta_\star$.
(a) \emph{One depth.} Let $s$ be admitted, $K\ge\max\{2,s+1\}$ and $X\in\mathcal V_s$ with $\varpi(X)<\varepsilon_{\rm r}(s)$. Then $H^K_s(X)$ and $H^{K+1}_s(X)$ are defined; each is in the axial gauge on $\Lambda_s$, has the average $X$ and has plaquettes at most $2B_3L^{-2}\varpi(X)$; and
\[{\rm dist}\bigl(H^{K+1}_s(X),H^K_s(X)\bigr)\le a_s(K).\]
Hence $H^\infty_s(X)=\lim_KH^K_s(X)$ exists, ${\rm dist}(H^\infty_s(X),H^K_s(X))\le r_s(K)$ uniformly on $\{\varpi<\varepsilon_{\rm r}(s)\}$, $H^\infty_s$ is continuous there, $\mathcal MH^\infty_s(X)=X$, $\varpi(H^\infty_s(X))\le2B_3L^{-2}\varpi(X)$, and $H^\infty_s(X^u)=H^\infty_s(X)^{u\circ\pi}$ for every gauge transformation $u$ of $\Lambda_s$, $\pi$ mapping each site of $\Lambda_{s+1}$ to the site of $\Lambda_s$ whose block holds it.
(b) \emph{Along the tower.} Let $d$ be admitted, $K\ge\max\{2,d+1\}$, $W\in J_d$, and $Y_j=H^{K+1}_{d,j}(W)$, $0\le j\le K+1-d$, the tower over $W$ at the cutoff $K+1$. Then for $0\le j<K-d$
\[{\rm dist}\bigl(Y_{j+1},H^K_{d+j}(Y_j)\bigr)\le L^{-j/2}a_d(K),\]
whereas (a) at the depth $d+j$ gives the number $a_{d+j}(K)=(\varepsilon_{[d+j]}/\varepsilon_{[d]})L^{2j}a_d(K)$, up to $L^{2j}a_d(K)$.
(c) \emph{The Lipschitz bound in the block field.} Let $K\ge2$, $d'<K$, $X'\in\mathcal V_{d'}$ with $\varpi(X')<\varepsilon_{\rm L}$, $B$ a real $1$-form on the bonds of $\Lambda_{d'}$ with $\sup|B|\le\beta_{\rm L}$, and $X=e^{iB}X'$. Then $\mathcal C_K(X')$ and $\mathcal C_K(X)$ are non-empty, each one orbit of $\mathcal G^{d'}_K$, and for $U'\in\mathcal C_K(X')$, $U\in\mathcal C_K(X)$, every depth $d''$ with $d'\le d''\le K$ and every closed path $\Gamma$ of $\Lambda_{d''}$ with $N$ bonds
\[\Bigl|\tfrac12\operatorname{tr}(\mathcal M^{K-d''}U)(\Gamma)-\tfrac12\operatorname{tr}(\mathcal M^{K-d''}U')(\Gamma)\Bigr|\le2B^{\rm e}_5\mathsf S_{\rm e}\,N\,L^{-(d''-d')}\sup|B|;\]
the constant holds neither $K$, nor $d'$, nor the number $d''-d'$ of the depths between.
(d) \emph{Two consecutive cutoffs: the loops at the depth $d+n$.} Let $d$ be admitted, $n\ge1$, $K\ge\max\{2,d+n+1\}$ with $a_d(K)\le\beta_{\rm L}$, $W\in J_d$, $\Gamma$ a closed path of $\Lambda_{d+n}$ with $N$ bonds and $W_\Gamma(V)=\frac12\operatorname{tr}V(\Gamma)$. Then
\[
\begin{multlined}
\bigl|W_\Gamma(H^{K+1}_{d,n}W)-W_\Gamma(H^K_{d,n}W)\bigr|\le\mathsf c_9\,N\,L^{-n/2}a_d(K),\qquad\bigl|W_\Gamma(H^\infty_{d,n}W)-W_\Gamma(H^K_{d,n}W)\bigr| \\
\le\mathsf c_9\,N\,L^{-n/2}r_d(K),
\end{multlined}
\]
$H^\infty_{d,n}W$ being the limit of the towers $H^K_{d,n}W$ (Proposition~\ref{C-372-9}(a)).
\end{lemma}

\begin{proof}[Sketch of proof]
(a) Let $t=K-s$ and $\tilde U$ a background over $X$ at the cutoff $K+1$ (it exists by Theorem~\YMMref{E}{10} in the setting of Section~\YMMref{52.7}{432}, as $\varepsilon_{\rm r}(s)\le\bar\varepsilon$). Its average $U'=\mathcal M\tilde U$, a field of the bare lattice of the cutoff $K$ with $\mathcal M^tU'=X$, is a competitor there (Lemma~\YMMref{113.8}{794}): its action, in the units of the level $t$, exceeds the action of the backgrounds over $X$ at the cutoff $K$ by at most $2\mathsf K_u\varepsilon_{\rm r}(s)^2L^{-t}(\ell L^s)^3$, the comparison of the background actions at the cutoffs $K$ and $K+1$ (equation~\YMMref{(2155)}{795}) and the energy of the background (equation~\YMMref{(2149)}{794}) bounding the two parts of this excess.

The excess is at least $\frac18\gamma_{\rm r}\sum_{k<t}L^{t-k}\|\mathcal A_k\|^2$ (Proposition~\YMMref{113.5}{792}), where $\|\cdot\|$ is the $\ell^2$ norm over the bonds and $\mathcal A_k$ the real $1$-form by which the average $\mathcal M^kU'$, in the gauge of its averaged contour variables, differs from the averaged background over $\mathcal M^{k+1}U'$; the term $k=t-1$ alone gives $\|\mathcal A_{t-1}\|\le4(\mathsf K_u/\gamma_{\rm r})^{1/2}\varepsilon_{\rm r}(s)(\ell L^s)^{3/2}L^{-(t+1)/2}$. In that gauge $\mathcal M^{t-1}U'$ is the averaged background over $X$ at the cutoff $K+1$, that is $H^{K+1}_s(X)$ (the axial-normalization argument in the proof of Proposition~\YMMref{52.3}{433}), and the chart of the fibre over $X$ writes it as $e^{iB'}H^K_s(X)$ with $\sup|B'|\le2C_{\rm db}\sup|\mathcal A_{t-1}|$ (equation~\YMMref{(2138)}{791}); with $\varepsilon_{\rm r}(s)=2c_\rho\varepsilon_{[s]}$ this is $a_s(K)$.

The axial gauge, the average and the plaquettes are equation~\YMMref{(953)}{433}. Summed over the cutoffs, a geometric series of ratio $L^{-1/2}$, the bound makes $(H^K_s(X))_K$ a Cauchy sequence in a compact space, uniformly in $X$; each $H^K_s$ is continuous (Proposition~\YMMref{52.3}{433}) and covariant, a bare gauge transformation equal to $u\circ\pi$ on $\Lambda_{s+1}$ carrying the elements of the minimal orbit over $X$ whose averages are in the axial gauge to those over $X^u$; $\mathcal M$ is continuous on the fields with plaquettes $<1/200L^2$, which holds the fields $H^K_s(X)$; and the limit inherits all this. (b) The same argument with $\tilde U=Y_{t+1}$, $t=K-d$, the bare field at the end of the tower of the cutoff $K+1$: it is a background over $W$ and over every $Y_j$, with $\mathcal M^{t+1-j}\tilde U=Y_j$ (the minimal orbit over $H^{K+1}_{d+j}(Y_j)$ lies in the one over $Y_j$, Lemma~\ref{lem:sf-class-window}(2)).

The competitor $U'=\mathcal M\tilde U=Y_t$ has the averages $\mathcal M^kU'=Y_{t-k}$, and its excess is bounded once, at the depth $d$, by $2\mathsf K_u\varepsilon_{\rm r}(d)^2L^{-t}(\ell L^d)^3$. At $k=t-1-j$ the term of the lower bound has the weight $L^{t-k}=L^{j+1}$, so $\|\mathcal A_k\|\le4(\mathsf K_u/\gamma_{\rm r})^{1/2}\varepsilon_{\rm r}(d)(\ell L^d)^{3/2}L^{-(t+1)/2}L^{-j/2}$; at the point $Y_j$ the chart of the fibre and the axial-normalization argument in the proof of Proposition~\YMMref{52.3}{433} give $Y_{j+1}=e^{iB'}H^K_{d+j}(Y_j)$ with $\sup|B'|\le2C_{\rm db}\sup|\mathcal A_k|\le L^{-j/2}a_d(K)$. The single bound on the excess at the depth $d$ pays for all the depths of the tower together. (c) The response of the background to a change of its datum (Theorem~\YMMref{E}{10} and Lemma~\YMMref{51.6}{398}, with the thresholds of Section~\ref{sec:background-step-complements}, for the $n=K-d'$ averaging steps on the whole torus, the datum $X'$ and the radius $\varepsilon_{\rm L}$): $\mathcal C_K(X')$ and $\mathcal C_K(e^{iB}X')$ are single orbits of $\mathcal G^{d'}_K$, and a background over $X$ is, up to a bare gauge transformation, $e^{iL^{-n}\mathcal H(B)}U^\circ$ with $U^\circ\in\mathcal C_K(X')$ and $|\mathcal H(B)(b)|\le\sigma:=B^{\rm e}_5\mathsf S_{\rm e}\sup|B|$ at every bare bond $b$, the exponentially decaying weights of the response summing to at most $\mathsf S_{\rm e}$.

The averages of every order $i\le n$ of a field so moved are, up to gauge, $e^{ia_i}\mathcal M^iU^\circ$ with $|a_i|\le2L^{-(n-i)}\sigma$ (Lemma~\ref{C-368-7}(v)). Along a closed path a gauge transformation conjugates the holonomy; the product bound $|\prod_iA_i-\prod_iA'_i|\le\sum_i|A_iA_i'^{-1}-\mathbf 1|$ on $G$, $|e^{iY}-\mathbf 1|\le|Y|$ and $|\operatorname{tr}Y|\le2|Y|$ give the bound at $i=K-d''$. (d) A telescope over the depths $d+j$, $0\le j<n$, along the tower $Y_j=H^{K+1}_{d,j}(W)$ of the finer cutoff. For a field $Y$ of a depth between $d$ and $d+n$ and $K'\in\{K,K+1\}$ let $\mathsf u^{K'}(Y)$ be the loop along $\Gamma$ of the average down to $\Lambda_{d+n}$ of a background over $Y$ at the cutoff $K'$: one number on the orbit, with $\mathsf u^{K'}(Y)=\mathsf u^{K'}(H^{K'}_s(Y))$ at the depth $s$ of $Y$ (Lemma~\ref{lem:sf-class-window}(2)) and $\mathsf u^{K'}(Y)=W_\Gamma(Y)$ at the depth $d+n$.

With $\Delta_j=\mathsf u^{K+1}(Y_j)-\mathsf u^K(Y_j)$ one has $\Delta_n=0$, $\Delta_0$ the left side of the first bound, and $\Delta_j-\Delta_{j+1}=\mathsf u^K(Y_{j+1})-\mathsf u^K(H^K_{d+j}(Y_j))$. By (b), $H^K_{d+j}(Y_j)=e^{iB}Y_{j+1}$ with $\sup|B|\le L^{-j/2}a_d(K)\le\beta_{\rm L}<\pi$, and $\varpi(Y_{j+1})\le2B_3\varpi(W)L^{-2(j+1)}<\varepsilon_{\rm r}(d+j+1)\le\varepsilon_{\rm L}$; so (c) at $d'=d+j+1$ and $d''=d+n$, $n-j-1$ depths apart, gives $|\Delta_j-\Delta_{j+1}|\le2B^{\rm e}_5\mathsf S_{\rm e}NL^{-(n-j-1)}L^{-j/2}a_d(K)$. The roles of the two fields are so chosen because (c) asks a bound on the plaquettes of $X'$ alone, which the tower of the finer cutoff has. The sum over $j$ is a geometric series of the ratio $L^{1/2}$, $\sum_{j<n}L^{j/2}\le L^{n/2}/(L^{1/2}-1)$, which gives $\mathsf c_9NL^{-n/2}a_d(K)$.

The conditions persist for $K'\ge K$, as $a_d$ decreases, and the sum over $K'\ge K$ with $\sum_{K'\ge K}a_d(K')=r_d(K)$ gives the second bound.
\end{proof}

\begin{proposition}[Smoothed loops over a fixed field]\label{C-372-9}
In the one-torus setting let $d$ be admitted, $W\in J_d$ and $C$ a lattice loop of $\Lambda_{d'}$, $d'\ge d$. Put $\varepsilon_{\rm r}=2c_\rho\varepsilon_{[d]}$; let $\mathsf c_{\rm A}>0$ be the constant with which the comparison of two consecutive cutoffs (Proposition~\YMMref{113.9}{795}) bounds the distance of the averaged backgrounds of two consecutive cutoffs over one block field (Lemma~\ref{C-372-31}(a)), and $\mathsf c_9,\beta_{\rm L}>0$ the constant and the radius of the Lipschitz bound of the tower in the block field (Lemma~\ref{C-372-31}(c), (d)), the numbers formed in Lemma~\ref{C-372-31}; all three depend on the constants alone. Put $a_s(K)=\mathsf c_{\rm A}\varepsilon_{[s]}(\ell L^s)^{3/2}L^{-(K-s+1)/2}$ and $r_d(K)=\sum_{K'\ge K}a_d(K')=r_d(d)L^{-(K-d)/2}$; let $K_2$ be the least $K\ge2$ with $\lfloor\frac{K-d}3\rfloor\ge\max\{1,d'-d\}$ and $a_d(K)\le\beta_{\rm L}$; and put $\mathsf c_{\rm f}=\frac{28L\varepsilon_{\rm r}}{1-L^{-1}}+\mathsf c_9r_d(d)$ and $\mathsf c_{\rm g}=\frac{1544L^2\varepsilon_{\rm r}}{1-L^{-1}}+6\mathsf c_9r_d(d)$, numbers that depend on $d$, $\ell$, $\mathsf b$ and the constants ($K_2$ also on $d'-d$).
(a) $V_n=H^\infty_{d,n}(W)$ satisfies $\mathcal MV_{n+1}=V_n$, $\varpi(V_n)\le2B_3\varpi(W)L^{-2n}$ and, for $n\ge d'-d$, $|W_C(V_n)-W_C(V_{n+1})|\le28B_3|C|L^d\varpi(W)L^{-n}$. So $W^{[d]}_C$ exists uniformly on $J_d$, with $|W^{[d]}_C-W_C(V_n)|\le\frac{14}{1-L^{-1}}|C|L^d\varepsilon_{\rm r}L^{-n}$ there; it is continuous and gauge invariant on $J_d$ with values in $[-1,1]$; $F^{[d]}_C$ is Borel; and $0\le1-W^{[d]}_C(W)\le\frac12(2B_3\varpi(W)L^{2d}s_1s_2)^2$ for a coordinate rectangle with sides $s_1$, $s_2$.
(b) For $K\ge\max\{2,d+1,d'\}$, $\mathcal C_K(W)$ is one non-empty orbit of $\mathcal G^d_K$; $w^K_C$ is continuous and gauge invariant on $J_d$; $\sup_{J_d}|w^K_C-W^{[d]}_C|\to0$, and this supremum is at most $|C|L^d\mathsf c_{\rm f}L^{-(K-d)/3}$ for $K\ge K_2$.
(c) For curves and any shadows, $|W_{\Gamma^{(d+n)}}(V_n)-W_{\Gamma^{(d+n+1)}}(V_{n+1})|\le772(|\gamma|L^{d+n+1}+2)L^{-2n}\varepsilon_{\rm r}$; so $W^{[d]}_\gamma$ exists uniformly on $J_d$, does not depend on the shadows, equals $W^{[d]}_C$ along a lattice loop $C$, and is continuous and gauge invariant, and $F_\gamma$ is Borel. $w^K_\gamma(W)$ is one number on $\mathcal C_K(W)$; two shadows on the bare lattice change it by at most $1496(|\gamma|L^K+2)L^{-2(K-d)}\varepsilon_{\rm r}$; and $\sup_{J_d}|w^K_\gamma-W^{[d]}_\gamma|\le(|\gamma|L^d+2)\mathsf c_{\rm g}L^{-(K-d)/3}$ for $K\ge K_2$, taken at $d'=d$.
(d) For $K\ge K_{\rm bf}$ with $\lceil K/1260\rceil\ge d$,
\[
\begin{multlined}
\Bigl|\langle F^{[d]}_C(\mathcal M^{K-d}U)\rangle_K-\int F^{[d]}_C\,d\mu_d\Bigr|\le\mathfrak r_K:=\mathsf C_{\rm bf}S_L\ell^3K^{2\mathfrak p_0+1/2}L^{-K/10},\qquad S_L=\sum_{j\ge0}(1 \\
+j)^{2\mathfrak p_0+1/2}L^{-j/10},
\end{multlined}
\]
a bound free of $C$ and $d$, where $\mathsf C_{\rm bf}$ and $K_{\rm bf}$ are the numbers of Definition~\textup{\ref{IV-62-62}}(5), formed at an interval of labels under the condition (Tor) on the side $\ell$ (Section~\ref{sec:cutoff-calibration}) that holds $\mathsf b$, and depend on the label only through that interval (Corollary~\textup{\ref{IV-62-66}}(a)); in the frame $[\mathsf b_-,\mathsf b_+]$ is such an interval, and it gives one pair of numbers for all $\mathsf b\in\mathsf J$, with the range $K\ge K_{\rm bf}$, $\lceil K/1260\rceil\ge d$; for $K\ge K_2$ besides, $|\langle(\mathbf 1_{J_d}w^K_C)(\mathcal M^{K-d}U)\rangle_K-\int F^{[d]}_C\,d\mu_d|\le|C|L^d\mathsf c_{\rm f}L^{-(K-d)/3}+\mathfrak r_K$. The same holds for curves, with $F_\gamma$, $w^K_\gamma$ and $(|\gamma|L^d+2)\mathsf c_{\rm g}$.
(e) In the frame put $\mathsf b'=\mathsf b-\bar B$ and $J'_d=\{\varpi<c_\rho\varepsilon(L^d\mathsf b')/B_3\}$, and let $d\ge d_0$, the least integer $d_0\ge1$ with $L^{d_0}\gamma^{-2}\ge\Theta_\star$ and $\log(L^{d_0}\gamma^{-2})>2\mathfrak p_0$. Then the one-torus setting holds at the labels $\mathsf b$ and $\mathsf b'$ and $d$ is admitted at both; $I_d\subset J_d\subset J'_d$ and $\mathsf O_{K-d}\subset J'_d$ for every $K>d$, so that $m_{K-d}$ gives no mass to $\mathcal V_d\setminus J'_d$, and $\mu^K_d(\mathcal V_d\setminus J'_d)\le R_d$; over every $V\in J'_d$ and every $K>d$ the backgrounds form one non-empty orbit, $w^K_C$ is Borel, and (a), (b) hold on $J'_d$, also for families: $\sup_{J'_d}|w^K_{\mathsf C}-\prod_iW^{[d]}_{C_i}|\to0$.
\end{proposition}

\begin{proof}[Sketch of proof]
(a) Over $W$ the backgrounds of the finer cutoffs are built one block step at a time, $H^K_{d,n}(W)$, and all of them are, up to gauge, block averages $\mathcal M^{K-d-n}U$ of one bare background $U$, whose plaquettes are $<B_3\varpi(W)(a_K/a_d)^2$ (Theorem~\YMMref{E}{10} in the setting of Section~\YMMref{52.7}{432}); so the factor $2B_3$ is paid once, not at every step: $\varpi(\mathcal M^{K-d-n}U)\le2B_3\varpi(W)L^{-2n}<\varepsilon_{\rm r}L^{-2n}\le L^{-3n/2}\,2c_\rho\varepsilon_{[d+n]}$. Induction on $n$, uniform continuity on the compact set $\{\varpi\le L^{-3n/2}2c_\rho\varepsilon_{[d+n]}\}$ and Lemma~\ref{C-372-31}(a) give $H^K_{d,n}\to H^\infty_{d,n}$ uniformly on $J_d$; the bound passes to $V_n$, which lies in $\{\varpi<2c_\rho\varepsilon_{[d+n]}\}$, so that $H^\infty_{d+n}(V_n)$ is defined.

The ball $\{\varpi<2c_\rho\varepsilon_{[d]}\}$, larger than $J_d$ by the factor $2B_3\ge2$, is the domain of the single step $H^\infty_d$; the whole tower needs room for the factor $2B_3$ that the backgrounds pay, which is why the proposition asks $W\in J_d$, and why the smaller ball $I_d\subset J_d$, and not the larger one, enters the starting depth. One refinement costs $|W_\Gamma(\mathcal MV)-W_\Gamma(V)|\le14L^2N\varpi(V)$ for a loop of $N$ bonds if $\varpi(V)<1/200L^2$ (in a gauge fixed along the blocks the average of a nearly flat field moves each bond by $O(L^2\varpi(V))$), which holds as $2c_\rho\varepsilon_{[d]}<1/200L^2$; with $N=|C|L^{d+n}$ the number of bonds grows like $L^n$ while the plaquettes fall like $L^{-2n}$: a geometric series, whose tail with $\varpi(W)<\varepsilon_{\rm r}/2B_3$ is the stated bound.

For a rectangle, $a_p=\frac12|\partial U(p)-\mathbf 1|^2$ and the holonomy of $\partial R$ is an ordered product of conjugated plaquette variables, so $1-W_{\partial R}\le(\sum_{p\subset R}\sqrt{a_p})^2$ (Proposition~\ref{C-372-13}(c)), and the bare plaquettes of the background are $<2B_3\varpi(W)(a_K/a_d)^2$.
(b) The tower ends in the non-empty set $\mathcal C_K(W)$, one orbit by Theorem~\YMMref{E}{10} in the setting of Section~\YMMref{52.7}{432} and Section~\ref{sec:analytic-core-complements}. The comparison of two consecutive cutoffs (Proposition~\YMMref{113.9}{795}) gives ${\rm dist}(H^{K+1}_s(X),H^K_s(X))\le a_s(K)$ for $\varpi(X)<2c_\rho\varepsilon_{[s]}$ (Lemma~\ref{C-372-31}(a)). At a depth $d+n$ we split into three terms. $|w^K_C(W)-W_C(\mathcal M^{K-d-n}U)|$ and $|W^{[d]}_C(W)-W_C(V_n)|$ are at most $\frac{14}{1-L^{-1}}|C|L^d\varepsilon_{\rm r}L^{-n}$, uniformly in $K$, by the series of (a); and $|W_C(H^K_{d,n}W)-W_C(V_n)|\le|C|L^{d+n}\sup{\rm dist}(H^K_{d,n},H^\infty_{d,n})\to0$ at fixed $n$. Choose $n$, then $K$. For the rate we use the Lipschitz bound of the tower in the block field (Lemma~\ref{C-372-31}(d)): for $n\ge1$, $K\ge\max\{2,d+n+1\}$ with $a_d(K)\le\beta_{\rm L}$, $W\in J_d$ and every closed path $\Gamma$ of $\Lambda_{d+n}$ with $N$ bonds, $|W_\Gamma(H^{K+1}_{d,n}W)-W_\Gamma(H^K_{d,n}W)|\le\mathsf c_9NL^{-n/2}a_d(K)$.

It is a telescope over the depths $d+j$, $0\le j<n$, along the tower $Y_j=H^{K+1}_{d,j}(W)$ of the finer cutoff: at the depth $d+j$ the one-depth smoothings $H^{K+1}_{d+j}$ and $H^K_{d+j}$ of the two cutoffs, applied to $Y_j$, give fields of the depth $d+j+1$ at distance at most $L^{-j/2}a_d(K)$ (Lemma~\ref{C-372-31}(b): the excess of the action of the averaged background of the finer cutoff is bounded once, at the depth $d$, and the term of the depth $d+j$ in the lower bound on the excess carries the weight $L^{j+1}$; the bound $a_{d+j}(K)$ of one depth, applied at $Y_j$, is up to $L^{2j}$ times larger), and the loop along $\Gamma$ of the average down to $\Lambda_{d+n}$ of the background over such a field, $n-j-1$ depths below it, moves by at most $2B^{\rm e}_5\mathsf S_{\rm e}NL^{-(n-j-1)}$ times that distance, a constant free of the cutoff and of the number of the depths between (Lemma~\ref{C-372-31}(c)); the sum over $j$ is a geometric series of the ratio $L^{1/2}$, which gives $\mathsf c_9$.

Summed over $K'\ge K$, with $\sum_{K'\ge K}a_d(K')=r_d(K)$, at $\Gamma=C^{\langle d+n\rangle}$ it bounds the middle term by $\mathsf c_9|C|L^dL^{n/2}r_d(K)$; at $n=\lfloor\frac{K-d}3\rfloor$ both $L^{-n}$ and $L^{n/2}L^{-(K-d)/2}$ are at most $L\cdot L^{-(K-d)/3}$, which gives $\mathsf c_{\rm f}$. Continuity and gauge invariance of $w^K_C$ on $J_d$ come with the tower.
(c) Two shadows of consecutive lattices, evaluated at $V_{n+1}$, differ by a band of $O(|\gamma|L^{d+n+1}+2)$ plaquettes of size $O(L^{-2n}\varepsilon_{\rm r})$; two shadows on the bare lattice, evaluated at a background over $W$, differ by a band of $O(|\gamma|L^K+2)$ bare plaquettes of size $O(L^{-2(K-d)}\varepsilon_{\rm r})$; then argue as in (a), (b).
(d) By Corollary~\textup{\ref{IV-62-66}} and \textup{(\ref{IV-62-68})}, two consecutive cutoffs differ on a gauge invariant bounded Borel function $F$ of the block field of a depth at most $\lceil K/1260\rceil$ ($1260=180(3N+4)$ at $N=1$ in Section~\ref{sec:cutoff-calibration}) by at most $\mathsf C_{\rm bf}\ell^3K^{2\mathfrak p_0+1/2}L^{-K/10}\sup|F|$; the comparison there is made through the small-field chains of the two cutoffs at that depth, and its numbers, formed at an interval of labels under (Tor), serve every label of the interval (Corollary~\textup{\ref{IV-62-66}}(a)). In the frame $[\mathsf b_-,\mathsf b_+]$ is such an interval on the side $\ell$: $\mathsf b_->g^{-2}\ge\gamma^{-2}$, and $M\,Q(\mathsf b_+)\le L^m$, with $Q$ of Lemma~\ref{C-372-7}, holds on the sides of $\mathfrak L(g)$ at every label of $\mathsf I(g)$, as in (e). Sum over $K'\ge K$ with $F=F^{[d]}_C$; the second mean adds the rate of (b).
(e) $\gamma^{-2}\le g^{-2}<\mathsf b'<\mathsf b<L(g^{-2}+\bar B)$, so the sides at $g$ satisfy the condition on the side at both labels, and $L^d\mathsf b'\ge L^d\gamma^{-2}\ge\Theta_\star$; the constraint sets, the action and the regularity radius hold no label, so (a), (b) apply at the label $\mathsf b'$ on its ball $J'_d$. $x_{K-d}\ge L^d\mathsf b-\bar B\ge L^d\mathsf b'$, $\varepsilon$ decreases for $\log t>2\mathfrak p_0$ and $c_\rho\ge2B_3$; so $2\varepsilon(x_{K-d})\le2\varepsilon(L^d\mathsf b')\le c_\rho\varepsilon(L^d\mathsf b')/B_3$, that is $\mathsf O_{K-d}\subset J'_d$, and $J_d\subset J'_d$ likewise. $\mu^K_d(\mathcal V_d\setminus J'_d)\le\mu^K_d(\mathcal V_d\setminus I_d)\le R_d$ by Proposition~\ref{C-372-8}(b). For families $|\prod_ix_i-\prod_iy_i|\le\sum_i|x_i-y_i|$ on $[-1,1]$.
\end{proof}

The new estimate is one renormalisation step with lines inserted. Its soft part rests on three facts about ordered products of matrices under a measure.

\begin{lemma}[Ordered products under a measure]\label{C-372-16}
(a) Let $\mathfrak A$ be a unital normed algebra, $I$ a finite totally ordered set, and for $j\in I$ let $\mathsf{At}_j$ be a finite set of atoms $\alpha$, each with an element $\tau_\alpha\in\mathfrak A$, $|\tau_\alpha|\le1$. A block is a non-empty set of atoms of pairwise different indices, a packing $\pi$ a set of blocks with pairwise disjoint sets of indices, and $\tau^\to(\pi)$ the product of the $\tau_\alpha$, $\alpha$ in the blocks of $\pi$, in the order of $I$. For complex weights $w$ on the blocks put $Z[w]=\sum_\pi\prod_{\beta\in\pi}w(\beta)\,\tau^\to(\pi)$ and $|w|_1=\sum_\beta|w(\beta)|$. Then $|Z[w]|\le e^{|w|_1}$ and $|Z[w_1+w_2]-Z[w_1]|\le|w_2|_1e^{|w_1|_1+|w_2|_1}$. If bounded functions $f_\alpha$ on a space with a complex measure $\nu$ of mass $1$ and finite total variation have $\int e^{\sum_\alpha z_\alpha f_\alpha}d\nu=e^{F(z)}$ with $F$ holomorphic near $0$ and $F(0)=0$, then
\[\int\prod^{\to}_{j\in I}\Bigl(\mathbf 1+\sum_{\alpha\in\mathsf{At}_j}f_\alpha\tau_\alpha\Bigr)d\nu=Z[\kappa],\qquad\kappa(\beta)=\Bigl(\prod_{\alpha\in\beta}\partial_{z_\alpha}\Bigr)F(0),\]
the joint cumulants of the $f_\alpha$, $\alpha\in\beta$.
\emph{Generators.} Here and below $\mathfrak{su}(2)$ is the real space of the Hermitian traceless $2\times2$ matrices, as the chart $e^{ig_tB(b)}V_W(0)(b)$ of the frame asks, with the scalar product $\langle X,Y\rangle=\frac12\operatorname{tr}XY$, and $t_1,t_2,t_3$ is an orthonormal basis of it: $t_at_b+t_bt_a=2\delta_{ab}\mathbf 1$, $t_a^2=\mathbf 1$, $|t_a|=1$ in the operator norm and $\sum_at_a^2=3\cdot\mathbf 1$ (the Pauli matrices up to a rotation). The components of $X\in\mathfrak{su}(2)$ are $X^a=\langle t_a,X\rangle$, so that $X=\sum_aX^at_a$ and $|X|=(\sum_a(X^a)^2)^{1/2}$; this is used for the variables $y_j$, the fluctuation field $B$, the test forms $\varphi$ and the rotations $\Upsilon$ below, and $(\varphi,B)=\sum_b\sum_a\varphi^a(b)B^a(b)$. A complex $2\times2$ matrix is written $x=x^0\mathbf 1+\sum_ax^a\,it_a$, with $x^0=\frac12\operatorname{tr}x$, $x^a=-\frac i2\operatorname{tr}(xt_a)$ and $|x^0|,|x^a|\le|x|$; as $(\sum_ax^at_a)^2=\sum_a(x^a)^2\mathbf 1$, $\det x=(x^0)^2+\sum_a(x^a)^2$. For $h\in\mathrm{SU}(2)$, ${\rm R}(h)\in SO(3)$ is the rotation with $ht_bh^{-1}=\sum_a{\rm R}(h)_{ab}t_a$, ${\rm R}(h)_{ab}=\langle t_a,ht_bh^{-1}\rangle$.
(b) Let $\mathfrak A=M_2(\mathbb C)^{\otimes r}$ with the operator norm and $I=I_1<\dots<I_r$ a union of intervals; let $c_j$ ($j\in I$) and $p_e$ ($e=\{j<j'\}$ in a set of pairs, each inside one $I_i$) be complex numbers, and $w_{\rm g}$ the weight that is $c_j$ at one block $\{\alpha^0_j\}$ with $\tau=\mathbf 1$, $p_e$ at the three blocks $\{\alpha^a_j,\alpha^a_{j'}\}$, $a=1,2,3$, with $\tau=it_a$ in the $i$-th factor, and $0$ elsewhere. Two pairs cross if their end points interlace. With $\mathfrak T=\sum_j|c_j|+3\sum_e|p_e|$, $\mathfrak X=\sum9|p_e||p_{e'}|$ over the crossing pairs of pairs, and $\mathfrak O=\sum|\hat w(\beta)||\hat w(\beta')|$ over the ordered pairs (equal ones included) of overlapping sets $\beta,\beta'$ among the $\{j\}$ and the $e$, $\hat w=c_j$ resp.\ $-3p_e$:
\[\Bigl|Z[w_{\rm g}]-\exp\Bigl(\sum_jc_j-3\sum_ep_e\Bigr)\mathbf 1\Bigr|\le(2\mathfrak X+\tfrac12\mathfrak O)\,e^{\mathfrak T}.\]
(c) Let $J=J_1\sqcup\dots\sqcup J_r$ be finite cyclically ordered sets (the lines), $g_j\in\mathrm{SU}(2)$, and $w(y)=\prod_i\frac12\operatorname{tr}\prod_{j\in J_i}(e^{i\sum_ay^a_jt_a}g_j)$ for $y=(y_j)\in(\mathbb R^3)^J$; every derivative of $w$ of every order has modulus at most $1$. Let $\varrho$ be a complex Borel measure of mass $1$ on $\{|y_j|\le\bar y\ \forall j\}$ whose logarithmic Laplace transform $\ell(z)=\log\int e^{(z,y)}d\varrho$ is holomorphic on a polydisc of radius $r_1>1$, with the cumulants $\kappa_\gamma=\partial^\gamma\ell(0)$, and let $s_{jj'}=s_{j'j}$ be real numbers, for the pairs $\{j,j'\}$ of one line and for $j=j'$, with $S=\sum_js_{jj}+2\sum_{\{j,j'\}}s_{jj'}\ge0$. For a pair of one line let $j\to j'$ be the arc from $j$ forward to $j'$, $G_{j\to j'}=g_jg_{j+1}\cdots g_{j'-1}$, ${\rm R}(h)\in SO(3)$ the rotation with $ht_bh^{-1}=\sum_a{\rm R}(h)_{ab}t_a$, $\kappa^{(2)}_{jj'}=(\kappa_{e_{ja}+e_{j'b}})_{a,b}$, $d_{jj'}=\kappa^{(2)}_{jj'}-s_{jj'}{\rm R}(G_{j\to j'})$, $d_{jj}=\kappa^{(2)}_{jj}-s_{jj}\mathbf 1$, and $|m|_1=\sum_{a,b}|m^{ab}|$. Then
\[
\begin{multlined}
\Bigl|\int w\,d\varrho-e^{-\frac32S}w(0)\Bigr|\le|\varrho|\,(E_1+E_2+E_3),\qquad E_1=\sum_\alpha|\kappa_{e_\alpha}|,\quad E_2=\sup_{|y_j|\le\bar y}\mathsf e(y),\quad E_3 \\
=\sum_{|\gamma|\ge3}\frac{|\gamma|}{\gamma!}|\kappa_\gamma|,
\end{multlined}
\]
$|\varrho|$ the total variation, with
\[
\begin{multlined}
\mathsf e(y)=\sum_{\{j,j'\}}\Bigl[|d_{jj'}|_1+6|s_{jj'}|\sum_{i\in(j\to j')}|y_i|+|\kappa^{(2)}_{jj'}|_1(|y_j|+|y_{j'}|+|y_j||y_{j'}|)\Bigr]+\sum_j\bigl[\tfrac12|d_{jj}|_1 \\
+|\kappa^{(2)}_{jj}|_1|y_j|\bigr]+\sum_{\rm two\ \relax lines}|\kappa^{(2)}_{jj'}|_1,
\end{multlined}
\]
the first sum over the pairs of one line, $i$ over the indices strictly between $j$ and $j'$ on the arc. If $\ell=\sum_X\ell_X$, $|\ell_X|\le A_X$ on the polydisc, $\ell_X$ depending on at most $N_X$ of the variables, and $r_1\ge2$, then $E_3\le3\sum_XA_X(N_X/r_1)^3e^{2(N_X+3)/r_1}$.
\end{lemma}

\begin{proof}[Sketch of proof]
(a) $|\tau^\to(\pi)|\le1$, and a packing with $n$ blocks arises from $n!$ of the $n$-tuples of blocks, so $\sum_\pi\prod|w(\beta)|\le e^{|w|_1}$; expanding $\prod(w_1+w_2)(\beta)$ and isolating one block of $w_2$ gives the second bound. For the identity, multiply out: the left side is $\sum_{S\subset I}\sum_{(\alpha_j)_{j\in S}}\int\prod_{j\in S}f_{\alpha_j}d\nu\prod^\to\tau_{\alpha_j}$; the moment is $\prod\partial_{z_{\alpha_j}}e^F$ at $0$, and by induction from $\partial e^F=(\partial F)e^F$ the derivatives of $e^F$ in pairwise different variables are, at $0$, sums over the partitions of the variables of products of derivatives of $F$; a choice of $S$, of the atoms and of a partition is a packing. (b) A packing in the support of $w_{\rm g}$ is a set of singletons, a set $\mathsf m$ of disjoint pairs and a colouring of $\mathsf m$.

If no two pairs of $\mathsf m$ cross, the sum over the colours is $(-3)^{\#\mathsf m}\mathbf 1$: take a shortest pair; between its two factors $\tau_a$ stand only factors $\mathbf 1$, and $\sum_a\tau_a\tau_a=-\sum_at_a^2=-3$; induct. The packings with a crossing contribute at most $2\mathfrak Xe^{\mathfrak T}$; and $\exp\sum\hat w-Z_{\rm nc}$, $Z_{\rm nc}$ the sum over the packings with the values $(-3)^{\#\mathsf m}$, is a sum over tuples of sets with an overlap, at most $\frac12\mathfrak Oe^{\mathfrak T}$. (c) Let $f(\tau)=\int w(\tau y)\,d\varrho$. The cumulant formula $\int y_\alpha G\,d\varrho=\sum_\beta\frac1{\beta!}\kappa_{\beta+e_\alpha}\int\partial^\beta G\,d\varrho$, for entire $G$ with bounded derivatives at $0$ (from $\partial_\alpha Z=Z\,\partial_\alpha\ell$ for $Z=e^\ell$, coefficient by coefficient, then by Taylor polynomials and dominated convergence), gives $f'(\tau)=\sum_{\alpha,\beta}\frac{\tau^{|\beta|}}{\beta!}\kappa_{\beta+e_\alpha}\int(\partial^\beta\partial_\alpha w)(\tau y)\,d\varrho$.

The terms $|\beta|=0$ are at most $|\varrho|E_1$, those with $|\beta|\ge2$ at most $|\varrho|E_3$. For $|\beta|=1$ the matrix identity $\sum c^{ab}_{jj'}\partial_{y^a_j}\partial_{y^b_{j'}}w=-3Sw+\mathcal R$, $|\mathcal R|\le2\mathsf e(y)$, at $c=\kappa^{(2)}$ holds: by cyclicity the trace of the line of a pair is $\frac12\operatorname{tr}(u_jHu_{j'}H')$ with $u_j=e^{i\sum_ay^a_jt_a}$; moving $e^{iy\cdot t}$ across $t_a$ costs $|y|$; $\sum_{a,b}{\rm R}(G)_{ab}t_aHt_b=\sum_b(Gt_bG^{-1})Ht_b=3H+\rho$ with $|\rho|\le6|H-G|\le6\sum_{i\in(j\to j')}|y_i|$ (a telescoping product of unitaries; $\sum_bt_b^2=3$); on the diagonal $\partial_a\partial_bu=-\delta_{ab}u$ up to $2|y|$. So $f'=-3\tau Sf+\mathcal E$ with $|\mathcal E|\le|\varrho|(E_1+2\tau E_2+E_3)$, and $f(1)=e^{-\frac32S}f(0)+\int_0^1e^{-\frac32S(1-\tau^2)}\mathcal E\,d\tau$, $S\ge0$.

The bound on $E_3$: Cauchy's estimates for each $\ell_X$ and $\sum_mm\binom{m+N-1}mx^m=Nx(1-x)^{-N-1}$, $x=1/r_1\le\frac12$.
\end{proof}

\begin{proposition}[One step with lines inserted]\label{C-372-10}
In the frame let $\mathsf C$ be a family with corners in $\Lambda_{i_0}$.
(a) There are $C_\eta<\infty$ and an integer $s_0\ge i_0$, depending on the constants, $g$, $\mathsf J$ and on $\mathsf C$ only through $\#\mathsf C$, $|\mathsf C|$, $s_{\mathsf C}$, $\rho_{\mathsf C}$ and $i_0$, not on $K$, $\mathsf b\in\mathsf J$, $\ell$, $t$, $W$, such that for $\mathsf b\in\mathsf J$, $K>s_0$, $0\le t<K-s_0$ and every $W\in\mathsf O_{t+1}$
\[\Bigl|\int W^{(t)}_{\mathsf C}(V_W(B))\,\nu_W(dB)-\lambda_{t+1}W^{(t+1)}_{\mathsf C}(W)\Bigr|\le C_\eta L^{-(K-t)/4};\]
more precisely, with $s=K-t$ (in (a) and its proof), $\Lambda_{\mathsf C}=\frac{3|\mathsf C|}{2M\mathsf b_-}$, $\mathsf g_s=(\frac3{2\mathsf b_-})^{1/2}L^{-s/2}$ and $\mathsf y_s=\log(2\mathsf b_+)+s\log L$, the left side is at most
\[C'(1+\#\mathsf C+\Lambda_{\mathsf C})e^{C''\Lambda_{\mathsf C}}\bigl[\mathsf y_s^{2\mathfrak p_0}(\mathsf g_s+L^{-s/2})+e^{-\rho_{\mathsf C}L^s/LM}\bigr]+\tfrac32\,\mathsf c\,C_X\mathsf b_-^{-2}|\mathsf C|L^{-s},\qquad\mathsf c=\tfrac54\bar B,\]
with $C'$, $C''$, $C_X$ depending on the constants alone; so for every $\vartheta<\frac12$ the same holds with $L^{-\vartheta(K-t)}$ and another $C_\eta$. For one polygon the same bound holds for the holonomy matrix, $|\int{\rm Hol}_C(e^{i\eta_t\mathbf A(B)}U_+)\,\nu_W(dB)-\lambda_{t+1}{\rm Hol}_C(U_+)|\le C_\eta L^{-(K-t)/4}$ in the operator norm, from a fixed base point, $e^{i\eta_t\mathbf A(B)}U_+$ being the configuration of the shift formula of step~1 of the proof (at $t=0$, where $\eta_0=1$, the chart $e^{iB'}U_+$ itself). No positivity of $\nu_W$ is used.
(b) If $\nu_W\ge0$, which is the case for $m_{t+1}$-almost every $W$, and $K-t>s_0$, the left side is at most
\[C_\sharp\frac{|\mathsf C|}{\mathsf b}\Bigl[(1+\mathfrak b_t)^2(\varepsilon_{t+1}+\delta_t+L^{-(K-t)})+e^{-\mathsf D_t}+\frac{\bar B}{\mathsf b}L^{-(K-t)}\Bigr],\qquad\mathsf D_t=\frac{\delta'_{\mathsf C}L^{K-t}}{LM}-2,\]
where $C_\sharp$ depends only on the constants, $n_{\mathsf C}$ and $c_{\mathsf C}$, and $s_0$ is the least integer $\ge i_0$ such that every $x\ge\frac23\mathsf b_-L^{s_0-1}$ satisfies explicit smallness conditions formed from the constants and $n_{\mathsf C}$, so that $s_0=\max\{i_0,s'_0\}$ with $s'_0$ formed from the constants, $\mathsf b_-$ and $n_{\mathsf C}$ alone; with the Euclidean distance in the place of the sup-distance in $\delta'_{\mathsf C}$, $\mathsf D_t$ would carry a further factor $1/\sqrt3$. At every $W\in\mathsf O_{t+1}$ a bound of the same form holds with the total variation of $\nu_W$ as a further factor. The right side is at most $C_\eta L^{-\vartheta(K-t)}$ for every $\vartheta<\frac12$, with $C_\eta$ depending on the constants, $\vartheta$, $\mathsf J$, $|\mathsf C|$, $n_{\mathsf C}$, $c_{\mathsf C}$, $\delta'_{\mathsf C}$, and on $g$ through $\mathsf J$ alone. For the squares $Q_i$ of side $a_i$ with corners in $\Lambda_i$ one has $s_0=\max\{i,s'_0\}=i+O(1)$, so $s_0=i$ for $i\ge s'_0$, and the bound $C(\mathsf J)\frac{a_i}{\mathsf b}L^{-\vartheta(K-t)}\le C(\mathsf J)L^{-\vartheta(K-t-i)}$, which depends on the scale of the step only through $a_{K-t}/a_i$.
(c) (The drift along a small square.) Let $C$ be a square of side $a_s$ in a coordinate plane with corners in $\Lambda_s$, let $\sigma_1,\dots,\sigma_N$ be its $N=4L^{K-t-s}$ bonds of $\Lambda_{K-t}$, and let $\Upsilon_i(B)=\frac1i\log[{\rm hol}_{\sigma_i}(e^{i\eta_t\mathbf A(B)}U_+)\,{\rm hol}_{\sigma_i}(U_+)^{-1}]\in\mathfrak{su}(2)$ be the rotation of $\sigma_i$ in the shift formula of step~1 of the proof of (a) (at $t=0$ the $\sigma_i$ are bare bonds, $\Upsilon_i=B'(\sigma_i)$ at a positively oriented $\sigma_i$ and $\Upsilon_i=-\operatorname{Ad}(U_+(\sigma_i))B'(\bar\sigma_i)$ at a reversed one). For every $\alpha<\frac12$ there are $c_1(\alpha)$ and $s_A$, depending on the constants, $\mathsf J$, $\alpha$ and a lower bound of $\ell$, not on $g$, $K$, $t$, $s$, $\mathsf b$, $W$, such that for $s\ge s_A$, $K>s$, $0\le t<K-s$ and every $W\in\mathsf O_{t+1}$
\[\sum_{i,a}\Bigl|\int\Upsilon^a_i\,d\nu_W\Bigr|\le c_1(\alpha)L^{-s}L^{-\alpha(K-t)}.\]
These are the rotations of the shift formula; the loop, being a trace, does not see the difference to the rotations of the background over $V_W(B)$ in any other gauge.
(d) (The mean of the fluctuation field.) For every $W\in\mathsf O_{t+1}$ and every real $1$-form $\varphi$ on the fluctuation bonds, $|\int(\varphi,B)\,\nu_W(dB)|\le C_4\delta_t(\varepsilon_{t+1}+a_{K-t}/\ell)\sum_b\max_a|\varphi^a(b)|$, $C_4$ depending on the constants: the mean of the fluctuation field is of the order of the coupling with its logarithm, $\delta_t$, times the bound $\varepsilon_{t+1}$ of the curvature of a tame $W$, the term $a_{K-t}/\ell$ coming from the domains that do not fit into the torus; on tame $W$ it is so of the order of the coupling squared, up to logarithms. The order $g_t^2$ times the curvature, up to a logarithm, is that of the mean of one rotation $\Upsilon_i$ of (c), which holds the further factor $g_t$ of $g_tB$.
(e) (Levels at least as coarse as the corners.) Let $0<\vartheta<\frac12$, $c_0>0$ and $\varrho>0$, let $C$ be a rectangle with sides in $[c_0\varrho,10\varrho]$ and corners in $\Lambda_i$, $i$ the integer with $a_i\le\frac12c_0\varrho<La_i$, and put $a_*=\varrho(\mathsf b/\varrho)^{\vartheta/(3+\vartheta)}$, a length coarser than the loop for small $\varrho/\mathsf b$. There is $C_*$, depending on the constants, $c_0$, $\vartheta$ and $\mathsf J$ alone, such that at every cutoff $K\ge i$, at every level $0\le t<K$ with $a_i\le a_{K-t}\le a_*$ and at every $W\in\mathsf O_{t+1}$ with $\nu_W\ge0$, hence for $m_{t+1}$-almost every $W$,
\[\Bigl|\int W^{(t)}_C(V_W(B))\,\nu_W(dB)-\lambda_{t+1}W^{(t+1)}_C(W)\Bigr|\le C_*\frac\varrho{\mathsf b}\Bigl(\frac{a_{K-t}}{\mathsf b}\Bigr)^\vartheta;\]
so the mean of the defect over the block field, $\bar{\mathsf D}_{t+1}$ of Proposition~\ref{C-372-12}(b) for $C$, has modulus at most $C_*\frac\varrho{\mathsf b}(\frac{a_{K-t}}{\mathsf b})^\vartheta$, as $Z^{-1}m_{t+1}$ has mass at most $1$ (Proposition~\ref{C-372-12}(c)). At the levels finer than the corners, $K-t>s_0=\max\{i,s'_0\}$, (b) gives the same bound with another $C_*$ of the same kind: there $|C|\le40\varrho$, $n_C=4$, the chord-arc constant of $C$ is bounded below in terms of $c_0$, and the bracket of (b), whose distance term vanishes for one polygon ($\delta'_C=\infty$), is at most a multiple of $(a_{K-t}/\mathsf b)^\vartheta$ formed from the constants, $\vartheta$ and $\mathsf J$. The levels $i<K-t\le s'_0$ occur only when $i<s'_0$, and then $\varrho>2a_{s'_0}/c_0$ and $a_{K-t}\ge a_{s'_0}$, so that there the bound $2$ of the defect at $\nu_W\ge0$ is at most $C_*\frac\varrho{\mathsf b}(\frac{a_{K-t}}{\mathsf b})^\vartheta$ for $C_*$ of the same kind. So the bound of one step holds at every level with $a_{K-t}\le a_*$; this is the bound of one step in Theorem~\ref{C-372-4}(8)(b).
\end{proposition}

\begin{proof}[Sketch of proof]
(a) \emph{0. The objects of the step.} A \emph{triple} on a domain $X$ is $P=(U,A',\mathbf J)$, $U$ an $\mathrm{SU}(2)$-valued configuration and $A'$, $\mathbf J$ complex $1$-forms on the bare bonds of $X$; it enters through its pair $\mathbf p(P)=(e^{i\eta_{t+1}A'}U,\mathbf J)$, and a gauge transformation $u$ acts by $P^u=(U^u,\operatorname{Ad}(u)A',\operatorname{Ad}(u)\mathbf J)$ (Section~\ref{sec:analytic-core-complements}). The \emph{background triple} of $W\in\mathsf O_{t+1}$ is $P_W=(U_+,0,J^{(t+1)}(U_+))$, $U_+$ a background over $W$ whose average $\mathcal M^tU_+$ is in the axial gauge and $J^{(t+1)}(U_+)$ its current (the background triple defined before equation~\YMMref{(827)}{395}, Section~\ref{sec:normalized-step-measure}); by Theorem~\YMMref{E}{10} in the setting of Section~\YMMref{52.7}{432} at the radius $2\varepsilon_{t+1}$, $|\partial U_+(p)-\mathbf 1|$ and $|(D^{\sf T}_{U_+}\partial U_+)(b)|$ are below $\alpha\eta_{t+1}^2$ at all bare plaquettes and bonds, $\alpha:=2B_3\varepsilon_{t+1}$, and $|J^{(t+1)}(U_+)|<\alpha$ (Section~\ref{sec:normalized-step-measure} at the level $t+1$).

The \emph{trivial triple} is $(\mathbf 1,0,0)$; the constant gauge transformations fix it. For a pair $v=(A,\mathbf J)$ of $1$-forms on $X$, $P[v]=(\mathbf 1,A,\mathbf J)$, and $\|v\|_X$ is the largest of $|A|$, $|\nabla A|$, $|\mathbf J|$ on the bonds of $X$ (Section~\ref{sec:analytic-core-complements}). The \emph{analyticity core} $\mathcal U^\circ(X)$ is a set of complex triples on $X$ given by bounds whose radii are formed from the constants alone (Definition~\YMMref{49.1}{374}); it is carried onto itself by $P\mapsto P^u$, kept by restriction to smaller domains, holds $P_W|_X$ for tame $W$ (Proposition~\ref{prop:normalized-step-kernel}), and holds the ball of the $P[v]$, $\|v\|_X<r_\alpha$, about the trivial triple, $r_\alpha$ formed from the constants (equation~\YMMref{(933)}{427}).

Three cited facts enter. (i) \emph{The kernel.} On $\{|B|<\mathfrak b_t\}$, $\nu_W=\Xi_W^{-1}e^{\mathsf E_W}\mu_{C_W}$: the Gaussian measure $\mu_{C_W}$ of the covariance $C_W$ of the step at $P_W$, times the exponential of two local potentials, the interaction and the part of the bare array of couplings, each with a norm with quadratic growth at most $c\,\delta_t$ and below the admitted norm $w_\star$ of the cluster expansion, normalised to the mass $1$ (\eqref{eq:normalized-step-measure} and Proposition~\ref{prop:normalized-step-kernel}). (ii) \emph{The provision on the core.} At every triple $P$ of $\mathcal U^\circ(X)$, complex ones included, the Gaussian structure and both potentials of the step are given, holomorphic in $(A',\mathbf J)$ at a fixed $U$, functions of the pair, gauge covariant, local (their restrictions to a smaller domain are the objects formed there) and with the same bounds as at $P_W$ (Definition~\YMMref{52.1}{415} and Theorem~\YMMref{W}{423}); so the cluster expansion for data with up to four potentials (Theorem~\YMMref{R}{361} for the logarithm and bound; Corollary~\YMMref{47.3}{364} for an added potential and the comparison with the Gaussian part; Proposition~\YMMref{47.4}{365} for holomorphy and locality, and the covariance argument following it) bounds every cluster term on the whole core, and not only at $P_W$. (iii) \emph{A potential on a cube.} If $|\partial U-\mathbf 1|$ and $|D^{\sf T}_U\partial U|$ are below $\alpha\eta^2$ near a cube $Q$ of integer side $\mathsf l$ in level units, with $\mathsf l\alpha\le c_*$ and $\mathsf l+4$ less than half the side of the torus, then $U^{u^{-1}}=e^{i\eta A}$ on $Q$ for a gauge transformation $u$ and a real $1$-form $A$ with $|A|,|\nabla A|<B_*\mathsf l\alpha$, $B_*$ and $c_*$ formed from the constants and free of $\mathsf l$ (Theorem~\YMMref{F}{199}).
\emph{1. The exact formula.} For $t\ge1$ the shift of Proposition~\YMMref{51.5}{396}, taken at the background triple and the $1$-form $g_tB$, gives a $1$-form $\mathbf A(B)$ on the bare bonds such that $e^{i\eta_t\mathbf A(B)}U_+$ is a gauge transform of a background over $V_W(B)$ (equation~\YMMref{(827)}{395}); the loop being gauge invariant, $W^{(t)}_{\mathsf C}(V_W(B))=W_{\mathsf C}(e^{i\eta_t\mathbf A(B)}U_+)$ exactly, $\eta_t=a_K/a_{K-t}$. $\mathbf A$ is holomorphic in $B$ and in the triple, at most $\mathsf c_{\mathfrak h}\sup|g_tB|$, local through the decoupling parameters, gauge covariant and real at real data (the shift properties proved in Proposition~\YMMref{51.5}{396}). At $t=0$ there is no shift system and no decoupling parameter: $\eta_0=1$, $W^{(0)}_{\mathsf C}=W_{\mathsf C}$, and the chart itself gives $V_W(B)=e^{iB'}U_+$, $U_+=V_W(0)$, with $B'=g_0B$ on the fluctuation bonds and, at the pivot bond $b_0(c)$ of each block bond $c$, $B'(b_0(c))=\beta_c(P_W,g_0B)$, a pivot value that reads the double block of $c$ alone, vanishes at $B=0$, is at most $2C_P\sup|g_0B|$, is holomorphic in the triple and in $B$, covariant at its first end point and $\mathfrak{su}(2)$-valued at real data (equation~\YMMref{(804)}{391} and the bounds and covariance proved with it); we put $\mathbf A(B):=B'$.

In both cases $U':=e^{i\eta_t\mathbf A(B)}U_+$ is the configuration of the shift formula. \emph{2. Segments.} The trace does not see where a closed path begins: each $C_i$ is cut at one site, its base point, and run from there. For $s\ge i_0+m'+1$, which $s_0$ ensures, the corners are centres of cubes of side $M$ of the level $t$, and the cubes that $\mathsf C$ passes are all different; we cut $\mathsf C$ into $N=|\mathsf C|/(Ma_{K-t})$ segments, one for each such cube, the segment $j$ being the steps of the path that start in the $j$-th cube. The segments of $C_i$ form an interval $I_i$ of indices, in the order of the path from the base point, and $I=I_1<\dots<I_r$. In the frames $R_x$ carried along the loop by $U_+$, $R_x$ the product of the $U_+(b')$ over the steps of $C_i$ from the base point to $x$, the holonomy of $U'$ from the base point is $\prod^\to_{j\in I_i}(\mathbf 1+k_j)\cdot{\rm Hol}_{C_i}(U_+)$, where $\mathbf 1+k_j$ is the ordered product over the steps $b$ of the segment $j$ of $R_be^{i\eta_t\sigma_b\mathbf A(b^+)}R_b^{-1}$, $b^+$ the positively oriented bond of the step $b$, $\sigma_b=\pm1$ its orientation and $R_b$ the frame at the first end point of $b^+$ (a step against the orientation gives $U'(b)=U_+(b)e^{-i\eta_t\mathbf A(b^+)}$); $\det(\mathbf 1+k_j)=1$.

For $t\ge1$ the brackets in the decoupling parameters of the shift system (equations~\YMMref{(806)}{392}--\YMMref{(809)}{392} and Lemma~\ref{lem:analytic-rooted-bracket}) make each $k_j$ a sum of local terms over the domains of the level $t$ that hold the cube of $j$, each with a factor $e^{-c\,d}$ in its size $d$, $c>1$ formed from the constants; at $t=0$, $k_j$ is one local term, on the cube of the segment with its six neighbours. The linear part of $k_j$ in $B$, which has no component along $\mathbf 1$ (its exponents have trace $0$), and its remainder have, in the norm with quadratic growth of the cluster expansion, the sizes $\mathsf n_\ell\le c\,g_t$ and $\mathsf n_{\rm r}\le c\,g_t^2$, $c$ formed from the constants, and not $c\,\delta_t$, which is what keeps every sum below of the order $|\mathsf C|/\mathsf b$; a sum of them over any set of segments has the same norms up to the factors $3$ and $4$.

In components, $k_j=\sum_a\ell^a_j\,it_a+\sum_{\mu=0}^3r^\mu_j\tau_\mu$, $\tau_0=\mathbf 1$, $\tau_a=it_a$, with the linear parts $\ell^a_j$ and the remainders $r^\mu_j$; so the atoms of $j$ are three linear ones, with $\tau=it_a$, and four remainder ones, with $\tau=\mathbf 1$ and $\tau=it_a$, and $k^0_j=r^0_j$, $k^a_j=\ell^a_j+r^a_j$. \emph{3. Cumulants.} The linear parts and the remainders are entered as two further potentials into the cluster expansion of the step (Theorem~\YMMref{R}{361}, with four potential families,
and Proposition~\YMMref{47.4}{365} for holomorphy and locality;
the two families already in $\nu_W$ are those of
Proposition~\ref{prop:normalized-step-kernel}), the linear parts as potentials given at every real $B$, the remainders as potentials given below the threshold $\mathfrak b_t$: $\log\int\exp(\sum_\alpha z_\alpha f_\alpha)\,d\nu_W=\sum_XF_X(z)$, the sum over the domains $X$ of the level $t+1$, with $F_X$ holomorphic, $F_X(0)=0$, $|F_X|\le c\,e^{-\lambda d(X)}$, $c$ and $\lambda>2$ constants of the cluster expansion ($\lambda-1$ exceeding the growth constant of the number of domains) and $d(X)$ the linear size of $X$, and $F_X$ depending on the variables of the segments in $X$ alone.

The polydisc is the one on which the inserted potentials stay below the admitted norm: its radii are $\rho_\ell=w_\star/(3\mathsf n_\ell)$, of the order $1/g_t$, in the variables of the linear atoms and $\rho_{\rm r}=w_\star/(4\mathsf n_{\rm r})$, of the order $1/g_t^2$, in those of the remainder atoms, the inverses of the sizes of step~2; they do not depend on $N$. Cauchy's estimates give all joint cumulants $\kappa(\beta)=\partial_\beta\sum_XF_X(0)$, with tree decay: a block $\beta$ with $n_\ell$ linear and $n_{\rm r}$ remainder atoms has $|\kappa(\beta)|\le\sum_Xc\,\rho_\ell^{-n_\ell}\rho_{\rm r}^{-n_{\rm r}}e^{-\lambda d(X)}$, of the order $g_t^{n_\ell+2n_{\rm r}}$, the sum over the $X$ that hold the segments of $\beta$.

We write $E(X)$ for the $X$-term of this expansion with one or two insertions, that is a first or a second derivative of $F_X$ at $z=0$, regarded as a function of the triple at which the step is formed. \emph{4. The ordered product.} By Lemma~\ref{C-372-16}(a), with $\mathfrak A=M_2(\mathbb C)^{\otimes r}$, one tensor factor for each polygon, the integral of the ordered product $\bigotimes_i\prod^\to_{j\in I_i}(\mathbf 1+k_j)$ is $Z[\kappa]$; the functional $T\mapsto2^{-r}\operatorname{tr}(T\cdot\bigotimes_i{\rm Hol}_{C_i}(U_+))$ has norm at most $1$, sends $\mathbf 1$ to $W^{(t+1)}_{\mathsf C}(W)$ and the integral of the ordered product to $\int W^{(t)}_{\mathsf C}(V_W(B))\,\nu_W(dB)$, so that the left side of (a) is at most $|Z[\kappa]-\lambda_{t+1}\mathbf 1|$.

For all $j\le j'$ in $I$ let $\delta_{cc'}p_{jj'}$ be the second cumulant of the linear parts $\ell^c_j$, $\ell^{c'}_{j'}$ formed at the trivial triple with the frames $\mathbf 1$, under the measure of the step there; it is a multiple of $\delta_{cc'}$ by step~5, it depends neither on $W$ nor on the colour, and $p_{j'j}:=p_{jj'}$. Let $c_j:=\kappa(\{(j;{\rm r},0)\})=\int k^0_j\,d\nu_W$, the first cumulant of the central remainder atom. A pair $j<j'$ of one polygon is \emph{near} if $j'-j\le N_\sharp:=\lfloor s_{\mathsf C}/(2Ma_{K-t})\rfloor$; the arc from the segment $j$ to the end of the segment $j'$ then has length at most $\frac12s_{\mathsf C}+Ma_{K-t}\le s_{\mathsf C}$ ($s_0$ ensures $N_\sharp\ge1$), so that it holds at most one corner, and it does not pass the base point.

It is \emph{straddling} if it is not near and $N_i-(j'-j)\le N_\sharp$, $N_i=\#I_i$ (near across the base point), and \emph{far} otherwise; pairs on two polygons are far. A far pair is at sup-distance at least $\rho_{\mathsf C}/a_{K-t}$ in level units, so that the domains that hold it have $d(X)\ge\rho_{\mathsf C}/(LMa_{K-t})-2$. We write $\kappa=w_{\rm g}+w_{\rm b}$, $w_{\rm g}$ the weight of Lemma~\ref{C-372-16}(b) with these $c_j$ and with $p_e=p_{jj'}$ at the near pairs $e=\{j<j'\}$, and $w_{\rm b}=\kappa-w_{\rm g}$; $w_{\rm b}$ holds the coloured singletons, the colour defects $\kappa(\{(j;\ell,a),(j';\ell,a')\})-\delta_{aa'}p_{jj'}$ of the near pairs, the other blocks with a remainder atom (all but the central singletons), the blocks of three or more atoms, and the far and the straddling pairs.

With $\mathsf g_s\ge g_t$, $N\mathsf g_s^2=\Lambda_{\mathsf C}$, and the bounds of steps~3 and~5, $|w_{\rm b}|_1\le c(1+\Lambda_{\mathsf C})\mathsf y_s^{2\mathfrak p_0}(\mathsf g_s+L^{-s/2})+c\Lambda_{\mathsf C}e^{2-\rho_{\mathsf C}L^s/LM}+c\,\#\mathsf C\,\mathsf g_s^2$, the far pairs giving the term with $\rho_{\mathsf C}$ and the straddling ones, of the order $g_t^2$ for each polygon, the term with $\#\mathsf C$. Further $\mathfrak T\le c\Lambda_{\mathsf C}$, and $\mathfrak X,\mathfrak O\le c\Lambda_{\mathsf C}\mathsf g_s^2$: a near pair that crosses a near pair $\{j_1<j_2\}$ has an end point among the $j_2-j_1$ indices between them, and $\sum_{j'}|p_{jj'}|$ and $\sum_{j'}|j'-j||p_{jj'}|$, the latter over the near pairs, are at most $c\,g_t^2$.

Lemma~\ref{C-372-16}(a), (b) give $Z[\kappa]=\exp(\sum_jc_j-3\sum_ep_e)\mathbf 1$ up to these errors. \emph{5. Symmetry.} What is not small by size, the first cumulants of the colour components and the colour defects of the near pairs and of the second cumulants $\langle\ell^a_j;\ell^{a'}_j\rangle$ of one segment, is small by symmetry. By fact (ii) each cluster term $E(X)$ is defined at every triple of $\mathcal U^\circ(X)$, holomorphic there in the sense of the slices, and bounded there as at $P_W$. Let $X$ be small, $d(X)\le\frac12s\log L$, and let $Q_X$ be the cube of side $\mathsf l_X=2M(d(X)+2)+1$ (level-$(t+1)$ units) about a cube of $X$; it holds $X$ and the arc of every near pair whose segments lie in $X$.

For $s\ge s_0$ the conditions of fact (iii) hold on $Q_X$ for every small $X$ ($\mathsf l_X\alpha\le c\,Ms\log L\,\varepsilon_{t+1}$ is small and $\mathsf l_X+4$ is far below the side of the torus), and fact (iii) at the $\alpha$ of step~0 gives $u$ and a real $A$ with $U_+^u=e^{i\eta_{t+1}A}$ on $Q_X$ and $|A|,|\nabla A|<B_*\mathsf l_X\alpha$; with $v_X:=(A,\operatorname{Ad}(u)J^{(t+1)}(U_+))$ on the bonds of $X$, $\|v_X\|_X<B'_*\mathsf l_X\alpha$, $B'_*=\max\{B_*,1\}$, and $(P_W|_X)^u$ and $P[v_X]$ have one pair. So the triple lies, up to the gauge transformation $u$, on the disc $\zeta\mapsto P[\zeta v_X]$, $|\zeta|<\rho_X$, about the trivial triple, inside the core by equation~\YMMref{(933)}{427}; its radius $\rho_X=\min\{r_\flat,1/(3\mathsf l_X)\}/(B'_*\mathsf l_X\alpha)$, $r_\flat$ formed from the constants, has $1/\rho_X\le cM^2\varepsilon_{t+1}(d(X)+2)^2$, $c$ formed from the constants, keeps the frames along the arcs in $Q_X$ within a factor $e$ of $\mathbf 1$, and is at least $2$ for large $s$.

The locality, the covariance and the dependence on the pair alone of the cluster terms (Proposition~\YMMref{47.4}{365} and the covariance argument following it, applied with fact (ii)) turn $E(X)$ at $P_W$, formed on the torus, into $E(X)$ at $P[v_X]$, formed on $X$, the frames being changed by one constant rotation, so that the colours are turned by one $O_X\in SO(3)$. At $\zeta=0$ the constant gauge rotations fix the trivial triple: a colour vector that they turn, the first derivative of a term in one coloured atom, vanishes there, and a bilinear form in the colours that commutes with all of them is a multiple of $\delta_{ab}$, $\mathbb C^3$ being irreducible under $SO(3)$. Schwarz' lemma for $\zeta\mapsto E(X)[P[\zeta v_X]]-E(X)[P[0]]$ on the disc gives the factor $1/\rho_X\le cM^2\varepsilon_{t+1}(d(X)+2)^2$, paid by the decay $e^{-\lambda d(X)}$; the domains with $d(X)>\frac12s\log L$ cost $e^{-\frac12s\log L}=L^{-s/2}$.

The size that this factor multiplies is the size of step~3, except for the first cumulant of a linear part, whose size is smaller: the Gaussian part of a cluster term, the same term formed with the linear insertions alone under the Gaussian measure of the step without threshold, is even under $B\mapsto-B$ (apply the covariance argument following Proposition~\YMMref{47.4}{365} to the reflection $B\mapsto-B$), so that its derivative in one linear atom vanishes at $z=0$; the first cumulant of a linear part is therefore the derivative of the difference of the term and its Gaussian part, which is at most $c\,\delta_t\,e^{-\lambda d(X)}$ (Corollary~\YMMref{47.3}{364}, equation~\YMMref{(724)}{364}): it has the amplitude $c\,\delta_t\mathsf n_\ell$, that of the non-Gaussian potentials.

With $\alpha\le c\,\mathsf g_s\mathsf y_s^{\mathfrak p_0}$, $\delta_t\le A_1\mathsf g_s\mathsf y_s^{\mathfrak p_0}$ and the sums over the $N$ segments, these are the bounds of step~4. \emph{6. Singletons.} $\det(\mathbf 1+k_j)=1$ and $\det x=(x^0)^2+\sum_a(x^a)^2$ give $(1+k^0_j)^2+\sum_a(k^a_j)^2=1$, that is $k^0_j=-\frac12\sum_a(k^a_j)^2-\frac12(k^0_j)^2$; integrated with $\int f^2\,d\nu_W=\langle f;f\rangle+(\int f\,d\nu_W)^2$, $\langle f;f\rangle$ the second cumulant, and with $\langle\ell^a_j;\ell^a_j\rangle=p_{jj}$ plus a diagonal defect, this gives $c_j=-\frac32p_{jj}$ up to errors of the size of step~4 (the diagonal defects, the mixed second cumulants with a remainder, the squares of first cumulants, $c_j^2$).

As $\sum_{j,j'\in I_i}p_{jj'}=\sum_{j\in I_i}p_{jj}+2\sum_{j<j'}p_{jj'}$,
\[\sum_jc_j-3\sum_ep_e=-\frac32\sum_i\sum_{j,j'\in I_i}p_{jj'}+\sum_j\bigl(c_j+\tfrac32p_{jj}\bigr)+3\!\!\sum_{j<j'\text{ in one }I_i,\text{ not near}}\!\!p_{jj'},\]
and the last two sums are of the size of the errors of step~4. \emph{7. The flat value.} At the trivial triple with the frames $\mathbf 1$ the linear part of a segment is linear in $B$, and the sum over the segments of $C_i$ of its colour component $a$ is $g_t(j_{C_i},\eta_t(\dot{\mathbf A}B)^a)=g_t(X_{C_i},\eta_t\,{\rm d}(\dot{\mathbf A}B)^a)$, $\dot{\mathbf A}$ the derivative at $0$ of $\mathcal A\mapsto\mathbf A$ formed at the trivial triple (at $t=0$ the derivative $\mathsf C_{\mathbf 1}=\mathsf C_{\mathbf 1,\mathbb R}\otimes1$ of the chart), $j_C$ the current of $C$ and ${\rm d}^{\sf T}X_C=j_C$. The linear part of the shift system at the trivial triple produces, up to one sign, the curvature response of the step: on every side of $\mathfrak L(g)$ and for $1\le t<K$ there is $\mathsf{sg}_t\in\{1,-1\}$ with $\eta_t\,{\rm d}(\dot{\mathbf A}\mathcal A)=\mathsf{sg}_t(\mathsf T^{\rm cv}_t\otimes1)\mathcal A$ for every $1$-form $\mathcal A$ on the fluctuation bonds, $\mathsf T^{\rm cv}_t=\mathsf H^{\rm cv}_t\mathsf C_{\mathbf 1,\mathbb R}$, $\mathsf H^{\rm cv}_t$ the curvature response of the level $t$ (the curvature-response construction preceding Lemma~\YMMref{120.2}{819}).

The proof follows that of Lemma~\YMMref{124.3}{835}, which is stated for data with two real sources and on large sides only; two changes make it run for the objects of the frame on every side of $\mathfrak L(g)$. First, the block field $\mathbf 1$ lies in $\mathsf O_{t+1}$, so that it has the properties that the kernel lemma of fact (i) (Proposition~\ref{prop:normalized-step-kernel}) gives to every tame block field and on which the hypotheses of (F7) rest; the configuration $\mathbf 1$ is a background over it, its average $\mathcal M^t\mathbf 1=\mathbf 1$ is in the axial gauge and its current vanishes, so that its background triple is the trivial triple; and the value of the shift on the curve below is taken from (F7) at the block field $\mathbf 1$.

Second, we use only the constrained variational equation in
the proof of Lemma~\YMMref{124.1}{832} and the curvature identity
\YMMref{(2230)}{819}. These calculations hold on every side;
the large-side identification with the full Landau inverse is
not needed. Constant gauge rotations fix the trivial triple,
so covariance makes the shift derivative commute with every
$\operatorname{Ad}(g_0)$, $g_0\in\mathrm{SU}(2)$.
It is therefore one real kernel times the identity on colours.
For real $\mathcal A$, (F7) makes
$U_\tau=e^{i\eta_t\mathbf A(\tau\mathcal A)}$ a curve of
configurations critical over their own averages. Write
$\dot A=\dot{\mathbf A}\mathcal A$ and
$\mathcal M^tU_\tau=\mathbf1+i\tau\dot V+o(\tau)$.
The cited variational calculation gives a real fine-site function
$\theta$ and a multiplier $\omega$ such that
\[
 h=\dot A+\nabla\theta,\qquad
 D_{\mathbf1}^{\mathsf T}D_{\mathbf1}h=Q^{\mathsf T}\omega,
 \qquad Qh=\dot V.
\]
Indeed, differentiation of the action on the moving constraint
kernel makes its flat Hessian pairing with $\dot A$ vanish on
the kernel of the linearized average. Factoring through that average gives the
multiplier; its value on gradients is zero, and the gauge
normalization restores the straight average $Qh=\dot V$.
Equation~\YMMref{(2230)}{819} now gives
\[
 \eta_t\,{\rm d}\dot A
 =\eta_t^2D_{\mathbf1}h
 =(\mathsf H_t^{\rm cv}\otimes1)\dot V.
\]
The chart over $\mathbf1$ identifies $\dot V$, up to a constant
rotation $\operatorname{Ad}(g_*)$, with
$\mathsf C_{\mathbf1}\mathcal A$ minus a gradient.
The curvature response annihilates this gradient. On a single
colour the remaining rotation gives a sign; linearity in
$\mathcal A$ makes it the same sign for all $\mathcal A$,
as in the last paragraph of the proof of
Lemma~\YMMref{124.3}{835}.

At $t=0$, $\mathsf H^{\rm cv}_0={\rm d}$ (the curvature-response calculation preceding Lemma~\YMMref{120.2}{819}) and the sign is $+$. At the trivial triple the covariance of the step is $\mathbf C_\ast(P_{\mathbf 1})=(\mathsf A^{\rm cv}_t)^{-1}\otimes1$, and $\mathsf T^{\rm cv}_t(\mathsf A^{\rm cv}_t)^{-1}(\mathsf T^{\rm cv}_t)^{\sf T}=\eta_t\,\mathcal C_{(t+1)}$ (Proposition~\YMMref{120.3}{819}, Proposition~\YMMref{120.3}{819}); so the Gaussian variance of the sum of the linear parts along $C_i$ is $g_t^2\eta_t(X_{C_i},\mathcal C_{(t+1)}X_{C_i})$ in each colour, the slice $(X_{C_i},\mathcal C_{(t+1)}X_{C_i})$ of the free Coulomb energy, and the sign drops out. The true pair sum differs from the Gaussian one, as in step~5, by the difference of the terms and their Gaussian parts: $\sum_{j,j'\in I_i}p_{jj'}=g_t^2\eta_t(X_{C_i},\mathcal C_{(t+1)}X_{C_i})$ up to $cN_i\delta_tg_t^2$, $N_i$ the number of segments of $C_i$.

With step~6, $\sum_jc_j-3\sum_ep_e=-\frac32S$ up to the errors of step~4, where 
\[
S:=g_t^2\eta_t\sum_i(X_{C_i},\mathcal C_{(t+1)}X_{C_i})\ge0,
\]
 $\mathcal C_{(t+1)}$ being an orthogonal projection: three colours and $\frac12\operatorname{tr}$ give $\frac32$. Between different polygons all pairs are far, and 
\[
|a_K(X_{C_i},\mathcal C_{(t+1)}X_{C_{i'}})|\le C_X|\mathsf C|e^{2-\rho_{\mathsf C}/LMa_{K-t}}
\]
 for $i\ne i'$; and $0\le a_K(X_C,\mathcal C_{(t+1)}X_C)\le C_X|C|$, $C_X$ formed from the constants. \emph{8. The label.} $g_t^2\eta_t=a_K/(a_{K-t}x_t)$ and $|1/(a_{K-t}x_t)-1/\mathsf b|\le\mathsf c\,\mathsf b^{-2}L^{-(K-t)}$, $\mathsf c=\frac54\bar B$, turn $e^{-\frac32S}$ into $\lambda_{t+1}$ up to $\frac32\mathsf cC_X\mathsf b_-^{-2}|\mathsf C|L^{-s}$, both exponents being $\le0$.

Steps~4 to~8 together give the bound, with $C'$, $C''$ formed from the constants; $\sup_s\mathsf y_s^{2\mathfrak p_0}L^{-s/4}<\infty$ and $\sup_sL^{s/4}e^{-\rho_{\mathsf C}L^s/LM}<\infty$ give $C_\eta$. The matrix form is step~4 without the trace, for one polygon from its base point.

(b) The variables are now the rotations $\Upsilon_\sigma(B)=\frac1i\log[(e^{i\eta_t\mathbf A(B)}U_+)(\sigma)U_+(\sigma)^{-1}]$ of the $|\mathsf C|/a_{K-t}$ bonds $\sigma$ of $\Lambda_{K-t}$ along $\mathsf C$, oriented as $\mathsf C$ runs (at $t=0$ the bare bonds of $\mathsf C$, with $\Upsilon_\sigma=B'(\sigma)$ at a positively oriented $\sigma$ and $-\operatorname{Ad}(U_+(\sigma))B'(\bar\sigma)$ at a reversed one), $|\Upsilon_\sigma|\le c\delta_t$, so that $W^{(t)}_{\mathsf C}(V_W(B))=w(\Upsilon(B))$, $W^{(t+1)}_{\mathsf C}(W)=w(0)$, with $w$ as in Lemma~\ref{C-372-16}(c), $g_\sigma=U_+(\sigma)$. By the brackets each $\Upsilon_\sigma$ is a sum of local terms, gauge covariant at $\sigma_-$ ($\Upsilon_\sigma\mapsto\operatorname{Ad}(u(\sigma_-))\Upsilon_\sigma$); their linear parts and remainders enter the cluster expansion as in (a), with $\log\nu_W(e^{(z,\Upsilon)})=\sum_X\ell_X(z)$ on a polydisc of radius $r_1\ge c/(n_{\mathsf C}g_t)$, one variable $z_\sigma\in\mathbb C^3$ for the whole of $\Upsilon_\sigma$.

Lemma~\ref{C-372-16}(c), with $\varrho$ the image of $\nu_W\ge0$ under $B\mapsto\Upsilon(B)$, and with $s_{\sigma\sigma'}\delta_{aa'}$ the covariance of the linear parts of $\Upsilon^a_\sigma$, $\Upsilon^{a'}_{\sigma'}$ formed at the trivial triple under the Gaussian measure of the step there (real, symmetric, free of $W$; a multiple of $\delta_{aa'}$ by the constant rotations), gives the bound with $E_1\le C_1\mathsf lg_t^2(\mathfrak b_t+\mathfrak b_t^2)\epsilon^\circ_t$, $E_2\le C_2\mathsf lg_t^2[\epsilon^\circ_t+(1+\mathfrak b_t)^2\delta_t+e^{-\mathsf D_t}]$ and $E_3\le C_3\mathsf lg_t^3$, $\mathsf l=|\mathsf C|/a_{K-t}$, where $\epsilon^\circ_t\le c(\varepsilon_{t+1}+L^{-(K-t)})$ is the price of the symmetry argument: the first cumulants and the defects of the second cumulants from $s_{\sigma\sigma'}{\rm R}(U_+(\sigma_-\to\sigma'_-))$ are small by the symmetry of (a), step 5, the transport running along the shorter arc, whose length the chord-arc constant bounds; the pairs on two polygons have $d(X)\ge\mathsf D_t$. $\sum_{\sigma,\sigma'\in C_i}s_{\sigma\sigma'}=g_t^2\eta_t(X_{C_i},\mathcal C_{(t+1)}X_{C_i})\ge0$ as in (a), step 7, and $\mathsf lg_t^2\le\frac32|\mathsf C|/\mathsf b$; the label costs $\frac94C_2\bar B|\mathsf C|\mathsf b^{-2}L^{-(K-t)}$.

For a complex $\nu_W$ of mass $1$ Lemma~\ref{C-372-16}(c) holds with the total variation as a factor. Positivity: by the disintegration with $m_t$, $m_{t+1}$, $\mathsf h_{t+1}\ge0$, $\int H(W)[\int G(V_W(B))\nu_W(dB)]m_{t+1}(dW)\ge0$ for all bounded Borel $H,G\ge0$; testing a countable set of continuous $G\ge0$ dense among them, the image of $\nu_W$ under $B\mapsto V_W(B)$ is a positive measure of mass $1$ for $m_{t+1}$-almost every $W$; as $V_W(B)(b)=e^{ig_tB(b)}V_W(0)(b)$ at the fluctuation bonds and $|g_tB(b)|<\delta_t<\log2$, the chart is injective with a Borel inverse on $\{|B|<\mathfrak b_t\}$, so $\nu_W$ itself is $\ge0$ for almost every $W$. The exponent $\frac12$ is the order of the error: five of its members are of the order $\frac{|\mathsf C|}{\mathsf b}x_t^{-1/2}$ up to logarithms, two of them of first order in the curvature of $W$.

For $Q_i$: $n_{Q_i}=4$ and $c_{Q_i}=\frac12$ are free of $i$, there is no distance term, and $|Q_i|/\mathsf b=4a_i/\mathsf b$.

(c), (d) A covariant insertion based at a site $x_0$ is a family of local terms of the step that turns under gauge transformations as a vector at $x_0$. For such an insertion each term $D(X;P)$ of the first cumulant, formed at the triple $P$ on a domain $X$ (the derivative at $0$ of the cluster term of $X$ with the insertion multiplied by $z$), is holomorphic in $P$ on the analyticity core and bounded there by $\mathsf A(\mathbb W)e^{-\lambda d(X)}$ (fact (ii) in step~0 of the proof of (a)), is a function of the pair, turns as a vector under constant rotations, and so vanishes at the trivial triple; a tame background on $X$ is gauge equivalent to a potential of norm $\le c\,M(d(X)+2)\varepsilon_{t+1}$, $c$ formed from the constants (fact (iii) on the cube of side $2M(d(X)+2)+1$ about $X$), as long as this cube fits into the torus and $c\,M(d(X)+2)\varepsilon_{t+1}$ stays below a number formed from the constants, the other domains costing $a_{K-t}/\ell$, resp.\ $\varepsilon_{t+1}$, with the factor $d(X)+2$; Schwarz' lemma gives $\sum_a|\int\mathbb W^{\rm tot}_a\,d\nu_W|\le c\,\mathsf A(\mathbb W)\varsigma_t$, $\varsigma_t=c\varepsilon_{t+1}+6LMa_{K-t}/\ell$, where $\mathbb W^{\rm tot}_a$ is the sum of the local terms of the insertion and $\mathsf A(\mathbb W)$ its amplitude in the cluster expansion.

For the fluctuation field, an odd insertion linear in $B$, the amplitude is that of the non-Gaussian potentials, $\lesssim\delta_t$, because the Gaussian part is even: this is (d). For the rotations $\Upsilon_i$ the linear part is an odd insertion of norm $\le cg_t$ and the remainder an insertion of norm $\le cg_t^2$, so each bond costs $cg_t^2(\log x_t)^{\mathfrak p_0}\varsigma_t$; and $Ng_t^2\le c\,a_s/\mathsf b$ for the $N=4L^{K-t-s}$ bonds, which gives (c). At $t=0$ the rotations are those of the chart, step~1 of the proof of (a), and the same argument runs.

(e) As $a_{K-t}\ge a_i>c_0\varrho/2L$ and $|C|\le40\varrho$, $|C|/a_{K-t}<80L/c_0$, and $C$ lies in a cube of side less than $20L/c_0$ in the units of the level. If $\varrho/\mathsf b$ is at least a number $\eta_0>0$ formed from the constants, $c_0$, $\vartheta$ and $\mathsf J$, the bound $2$ of the defect at $\nu_W\ge0$ is at most $C_*\frac\varrho{\mathsf b}(\frac{c_0\varrho}{2L\mathsf b})^\vartheta\le C_*\frac\varrho{\mathsf b}(\frac{a_{K-t}}{\mathsf b})^\vartheta$ for $C_*\ge2(2L/c_0)^\vartheta\eta_0^{-1-\vartheta}$; so let $\varrho/\mathsf b<\eta_0$, and with it $a_{K-t}/\mathsf b\le a_*/\mathsf b=(\varrho/\mathsf b)^{3/(3+\vartheta)}$, $g_t$, $\delta_t$ and $\varepsilon_{t+1}$ be small, and $K-t$ large.

The proof of (b) runs with the pieces of $C$ in the cubes of the level in the place of the bonds of the level. \emph{The exact formula.} Step~1 of the proof of (a) holds for every bare loop, the corners entering nowhere: $W^{(t)}_C(V_W(B))=W_C(U')$, $U'=e^{i\eta_t\mathbf A(B)}U_+$, at every level. \emph{Pieces.} From a base point, $C$ is cut into the maximal runs $C_1,\dots,C_n$ of consecutive steps that start in one of the half-open cubes of side $a_{K-t}$ about the sites of $\Lambda_{K-t}$, the cubes of the level that $C$ passes; each of the four sides of $C$ gives at most its length over $a_{K-t}$ plus $2$ of them, so that $n<80L/c_0+8$, a number formed from $L$ and $c_0$. With $g_j=U_+(C_j)$, the holonomy of $U_+$ along the piece, and $\Upsilon_j(B)=\frac1i\log[U'(C_j)\,U_+(C_j)^{-1}]\in\mathfrak{su}(2)$, one has $|\Upsilon_j|\le c\,\delta_t|C_j|/a_{K-t}$, $c$ formed from the constants (the shift is at most $\mathsf c_{\mathfrak h}\sup|g_tB|$, at $t=0$ at most $2C_P\sup|g_0B|$, at each of the $|C_j|/a_K$ bare steps, and $\eta_t|C_j|/a_K=|C_j|/a_{K-t}$), and $W^{(t)}_C(V_W(B))=w(\Upsilon(B))$, $W^{(t+1)}_C(W)=w(0)$, with $w$ as in Lemma~\ref{C-372-16}(c) for one line.

\emph{Insertions.} By the brackets each $\Upsilon_j$ is a sum of local terms of the level, gauge covariant at the first end point of $C_j$, with a linear part in $B$ of norm at most $c\,g_t|C_j|/a_{K-t}$ and a remainder of norm at most $c\,g_t^2\frac{|C_j|}{a_{K-t}}(1+\frac{|C_j|}{a_{K-t}})$ in the norm of the cluster expansion; each piece is one insertion, of norm of the order $g_t$ times its length over $a_{K-t}$. Entered into the cluster expansion of the step as in step~3 of the proof of (a), they give $\log\nu_W(e^{(z,\Upsilon)})=\sum_X\ell_X(z)$, one variable $z_j\in\mathbb C^3$ for the whole of $\Upsilon_j$, on a polydisc of radius $r_1\ge c\,a_{K-t}/(g_t|C|)\ge c\,c_0/(Lg_t)\ge2$, and Cauchy's estimates give a cumulant of order $m$ the factor $(c\,g_t|C|/a_{K-t})^m$ and the decay of its domains.

\emph{Symmetry.} Let $s_{jj'}\delta_{aa'}$ be the covariance of the linear parts of $\Upsilon^a_j$, $\Upsilon^{a'}_{j'}$ formed at the trivial triple under the Gaussian measure of the step there (real, symmetric, free of $W$; a multiple of $\delta_{aa'}$ by the constant rotations). At the trivial triple the first cumulants, colour vectors, vanish, and the pair cumulants are multiples of $\delta_{aa'}$, which differ from $s_{jj'}\delta_{aa'}$ by parts of the relative size $c\,\delta_t$, from the non-Gaussian potentials, and $c\,g_t$, from the remainders (step~5 of the proof of (a)). For a domain $X$ that holds pieces of $C$, with $d(X)\le\frac12(K-t)\log L$, fact (iii) of step~0 of the proof of (a), on a cube of side at most $2M(d(X)+2)+1+40L/c_0$ that holds $X$ and $C$, where its conditions hold for $\varrho/\mathsf b<\eta_0$, puts the background triple, up to a gauge transformation, on a disc about the trivial triple whose radius $\rho_X$ has $1/\rho_X\le c\,\varepsilon_{t+1}(d(X)+2)^2$, $c$ formed from the constants, $L$ and $c_0$, the frames along $C$ staying within a factor $e$ of $\mathbf 1$; Schwarz' lemma on this disc bounds the terms of $X$ of the first cumulants, and of the pair cumulants less their values at the trivial triple turned by the frames, by $1/\rho_X$ times their size, paid by the decay $e^{-\lambda d(X)}$, and the domains with $d(X)>\frac12(K-t)\log L$ cost $L^{-(K-t)/2}$.

So the first cumulants and the defects $d_{jj'}$ of the pair cumulants from $s_{jj'}{\rm R}(G_{j\to j'})$ carry, relative to their size, the factor $c\,(\varepsilon_{t+1}+L^{-(K-t)/2}+\delta_t+g_t)$, and the first cumulant of a linear part carries besides the amplitude $c\,\delta_t$ of the non-Gaussian potentials, its Gaussian part being even. \emph{The pair term.} The argument of step~7 of the proof of (a) holds for every bare loop: at the trivial triple with the frames $\mathbf 1$ the sum of the linear parts of the $\Upsilon_j$ is $g_t\eta_t(j_C,\dot{\mathbf A}B)$, so that $S:=\sum_{j,j'}s_{jj'}=g_t^2\eta_t(X_C,\mathcal C_{(t+1)}X_C)\ge0$, the slice of the free Coulomb energy at the step, at the depths coarser than the corners as at the finer ones.

By Proposition~\ref{C-372-11}(c), (d), which hold at every depth, $0\le a_K(X_C,\mathcal C_{(t+1)}X_C)=v_{K,K-t}(C)-v_{K,K-t-1}(C)\le C^{\rm bg}A_C^2/a_{K-t}^3\le c\,|C|$, with the area $A_C\le100\varrho^2$ of $C$, $|C|\ge4c_0\varrho$ and $c$ formed from $L$ and $c_0$; so step~8 of the proof of (a) turns $e^{-\frac32S}$ into $\lambda_{t+1}$ up to $c\,|C|\mathsf b^{-2}a_{K-t}$. \emph{The bound.} Lemma~\ref{C-372-16}(c), its measure being the image of $\nu_W\ge0$ under $B\mapsto\Upsilon(B)$, a probability measure on $\{|y_j|\le c\,\delta_t|C|/a_{K-t}\}$, bounds the defect by $E_1+E_2+E_3$ and this term of the label, the sum over two lines being empty. With $g_t^2\le\frac32a_{K-t}/\mathsf b$, $\delta_t=g_t\mathfrak b_t$, $\varepsilon_{t+1}$ and $L^{-(K-t)/2}$ of the order $(a_{K-t}/\mathsf b)^{1/2}$ up to logarithms, and the sizes above, each of $E_1$, $E_2$, $E_3$ is at most $\frac{|C|}{\mathsf b}(\frac{a_{K-t}}{\mathsf b})^{1/2}y_t^{\mathfrak p}$, $y_t=\log x_t$, $\mathfrak p$ formed from $\mathfrak p_0$, times a polynomial in $|C|/a_{K-t}$ with coefficients formed from the constants and $\mathsf J$, and the term of the label is at most $c\frac{|C|}{\mathsf b}\frac{a_{K-t}}{\mathsf b}$.

As $|C|/a_{K-t}<80L/c_0$, $|C|\le40\varrho$ and $(\frac{a_{K-t}}{\mathsf b})^{1/2}y_t^{\mathfrak p}\le c\,(\frac{a_{K-t}}{\mathsf b})^\vartheta$ for $\vartheta<\frac12$, $c$ formed from the constants, $\vartheta$ and $\mathsf J$, the defect is at most $C_*\frac\varrho{\mathsf b}(\frac{a_{K-t}}{\mathsf b})^\vartheta$. The mean over the block field follows, $|\bar{\mathsf D}_{t+1}|$ being at most the bound of the defect at $m_{t+1}$-almost every $W$ times $Z^{-1}m_{t+1}(\mathcal V_{K-t-1})\le1$. The levels finer than the corners are treated in the statement of (e) itself.
\end{proof}

\begin{proposition}[Numbers of the free field]\label{C-372-11}
Let $L\ge3$ be odd, $\ell=2L^m$, and $C$ a polygon with corners in $\Lambda_i$; no measure and no $\mathrm{SU}(2)$ enter.
(a) For $K\ge i$ and $a=a_K$: $\mathcal I_a(C,C)=\frac{|C|}{2\pi}\log\frac1a+c_{\rm lat}|C|+\mathcal J_\ell(C)+E_a(C)$ with $|E_a(C)|\le c[n_Ca(1+\log\frac{|C|}a)+a^2|C|^2{\rm Lg}(\ell/a)(\delta_C^{-3}+\ell^{-3})]$, $c$ absolute; pieces in perpendicular directions do not interact in the sum; $|\mathcal J_\ell(C)-\mathcal J(C)|\le c|C|^2{\rm diam}(C)^2\ell^{-3}$; for a rectangle $s\times t$,
\[\mathcal J(C)=\frac1\pi\Bigl[s\log s+t\log t-s\operatorname{asinh}\frac st-t\operatorname{asinh}\frac ts+2\sqrt{s^2+t^2}-2(s+t)\Bigr].\]
For polygons $C$, $D$ at distance $\delta>0$: $|\mathcal I_a(C,D)-\mathcal I_\ell(C,D)|\le ca^2|C||D|\,{\rm Lg}(\ell/a)(\delta^{-3}+\ell^{-3})$.
(b) For $0\le d\le K$ and every closed path $C'$ of $\Lambda_K$, $v_{K,d}(C')$ is the largest value of $2\mathcal I_{a_K}(j_{C'},Q^{\sf T}\iota)-\mathcal I_{a_K}(Q^{\sf T}\iota,Q^{\sf T}\iota)$ over the closed currents $\iota$ of $\Lambda_d$ without zero mode, $Q$ the linearisation of $\mathcal M^{K-d}$ at the trivial field (the mean over a block of the straight bare paths): the energy of the best smeared current. For the polygon $C$ and $i\le d\le K$,
\[v_{K,d}(C)=\frac{|C|}{2\pi}\log\frac1{a_d}+(c_{\rm lat}+c_{r,K-d})|C|+\mathcal J_\ell(C)+{\rm Rem}_{K,d}(C),\]
$|{\rm Rem}_{K,d}(C)|\le|E_{a_d}(C)|+c[n_C^2a_d+n_Ca_d(1+\log\frac{|C|}{a_d})+a_d|C|^2/\delta_C^2]$, a bound free of $K$. Here the numbers $c_{r,j}$, $j\ge0$ an integer, are formed as $c_r$ in the setting, with $S_{(j)}(z)=\frac{\sin(z/2)}{L^j\sin(z/2L^j)}$ in the place of $S(z)$, $E_{(j)}(p_\nu)=4L^{2j}\sin^2\frac{p_\nu}{2L^j}$ in the place of $p_\nu^2$ in the denominators, and the sums over the $n\in\mathbb Z^2$ with $|n_2|,|n_3|\le\frac{L^j-1}2$ in the places of the sums over $\mathbb Z^2$, in $\mathsf N$ and in $\mathsf D$; at $j=K-d$, $L^j=a_d/a_K$ is the ratio of the two spacings. The number $c_{r,j}$ depends on $L^j$ alone, $c_{r,0}=0$ (then $S_{(0)}=1$ and only $n=0$ enters), $|c_{r,j}|\le c_3$, and $c_{r,j}\to c_r$ as $j\to\infty$. The limit $v_d=\lim_Kv_{K,d}$ exists and obeys the same remainder bound:
\[\Bigl|v_d-\frac{|C|}{2\pi}\log\frac1{a_d}-c_Q|C|-\mathcal J_\ell(C)\Bigr|\le |E_{a_d}(C)|+c\Bigl[n_C^2a_d+n_Ca_d\Bigl(1+\log\frac{|C|}{a_d}\Bigr)+\frac{a_d|C|^2}{\delta_C^2}\Bigr].\]
For a fixed polygon the right side is $O_{C,\ell}(a_d(1+\log^+(1/a_d)))$ as $d\to\infty$. The constants $c$ and $c_3$ are free of the side; $c_Q$ is a number of the plane, free of $L$, $\ell$ and $C$.
(c) $\prod_{n\le K-d}\lambda_n=\exp[-\frac3{2\mathsf b}(\mathcal I_{a_K}(C,C)-v_{K,d}(C))]$ for all $0\le d\le K$. Hence for $i\le d\le K$
\[a_K^{-\kappa}\prod_{n\le K-d}\lambda_n=L^{\kappa d}\exp\Bigl[\frac3{2\mathsf b}\bigl(c_{r,K-d}|C|+{\rm Rem}_{K,d}(C)-E_{a_K}(C)\bigr)\Bigr],\]
and one step, $n=K-d$ with $d\ge i$, gives $\lambda_n=L^{-\kappa}\exp[-\frac3{2\mathsf b}((c_{r,n-1}-c_{r,n})|C|+{\rm Rem}_{K,d+1}-{\rm Rem}_{K,d})]$: each step takes away $\kappa\log L$ and a multiple of the length; the steps sum to $\sum_{n\ge1}(c_{r,n-1}-c_{r,n})=-c_r=c_{\rm lat}-c_Q$, every term depending on $L$ and the sum not. For families, $c_\Xi^{-1}L^{\kappa d}\le a_K^{-\kappa}\Xi_{K-d}\le c_\Xi L^{\kappa d}$ for $i_0\le d\le K$, with $c_\Xi$ free of $d$ and $K$ and uniform for $\mathsf b$ in a compact set, and $\mathsf l_d=\lim_Ka_K^{-\kappa}\Xi_{K-d}=a_d^{-\kappa}\exp[-\frac3{2\mathsf b}((c_{\rm lat}-c_Q)|\mathsf C|+O_{\mathsf C,\ell}(a_d(1+\log^+(1/a_d))))]$, so $a_d^\kappa\mathsf l_d\to e^{-0.047235\ldots|\mathsf C|/\mathsf b}$.
(d) For a rectangle $C$ of area $A$ with corners in $\Lambda_K$, in any position, and $0\le d\le K$: $0\le v_{K,d}(C)\le C^{\rm bg}A^2/a_d^3$, $C^{\rm bg}$ absolute.
\end{proposition}

\begin{proof}[Sketch of proof]
(a) Fourier analysis of $g_N$: a component of the current couples only to the same component, so perpendicular pieces do not meet; the sum of $g_N$ along an axis gives the logarithm and the constant $c_{\rm lat}$, a computation with the lattice Green function on a line; the torus gives $\mathcal J_\ell$; distant pieces are compared with the continuum with an error $O(a^2)$, the corners with an error $O(a)$ each. The formula for the rectangle is the double integral of $\frac1{4\pi|x-y|}$ over its four sides, the parallel pieces interacting with the sign $-1$ and the perpendicular ones not at all. (b) The range of $\Pi_{K-d}$ is spanned by the curvatures of the bare fields $\Delta^{-1}Q^{\sf T}\iota$, $\iota$ a closed current of $\Lambda_d$ (the Lagrange multiplier of the linearised background problem); as $(X,\Pi X)$ is the largest value of $2(X,Y)-(Y,Y)$ over the range and $(X_{C'},{\rm d}A)=(j_{C'},A)$, $v_{K,d}$ is the stated maximum, at every depth and for every closed bare path.

Since the corners of the polygon lie in $\Lambda_d$ for $d\ge i$, the Fourier transform of the sharp loop at the momenta $(q+2\pi n)/a_d$ is the transform $\hat\jmath(q)$ of its chain on $\Lambda_d$ times the form factor of one bond; the quadratic form is diagonal in $q$ and in the direction of the current, and the maximum is solved mode by mode. Without the constraint that $\iota$ be closed one gets $a_dN^{-3}\sum_{q\ne0}\sum_\mu|\hat\jmath_\mu(q)|^2\mathsf m^{(K-d)}_\mu(q)$, $N=\ell/a_d$, where $\mathsf m^{(j)}_\mu=(\mathsf N^{(j)}_\mu)^2/\mathsf D^{(j)}_\mu$ is the propagator that the block averages see: with $p=q+2\pi n$, $s_{(j)}(p)=\prod_\nu S_{(j)}(p_\nu)$ and $\Delta_{(j)}(p)=\sum_\nu E_{(j)}(p_\nu)$, $\mathsf N^{(j)}_\mu$ sums, over the aliases ($n$ in the $L^{3j}$ classes $|n_\nu|\le\frac{L^j-1}2$ at a finite cutoff, $j=K-d$, and in $\mathbb Z^3$ at $K=\infty$, where $S_{(\infty)}=S$ and $E_{(\infty)}(z)=z^2$), the bare propagator $1/\Delta_{(j)}(p)$ times the form factors $S_{(j)}(p_\mu)$ of the sharp bond and $S_{(j)}(p_\mu)s_{(j)}(p)$ of the smeared bond, and $\mathsf D^{(j)}_\mu$ the same with the smeared bond twice; at $q_\mu=0$ only $n_\mu=0$ survives, and they reduce to the sums $\mathsf N$, $\mathsf D$ of $c_{r,j}$ (of $c_r$ at $K=\infty$).

The constraint costs a non-negative term at most $cn_C^2a_d$: with $\partial_\mu(q)=e^{iq_\mu}-1$ and $\langle z,z'\rangle=\sum_\mu\bar z_\mu z'_\mu$, $\langle\partial(q),\hat\jmath(q)\rangle=0$ for a closed $C$ (the sum over the steps telescopes), $|\bar\partial_\mu\hat\jmath_\mu|$ is at most twice the number of pieces in the direction $\mu$, $\mathsf N^{(j)}_\mu/\mathsf D^{(j)}_\mu$ differs from $1/s_{(j)}(q)$, the inverse form factor of the block, which does not depend on $\mu$, by $O(|q|^2)$, and $\sum_\mu|\partial_\mu|^2/\mathsf D^{(j)}_\mu\ge c|q|^4$. With $\Delta_1(q)=\sum_\nu4\sin^2\frac{q_\nu}2$, the symbol $\mathsf m^{(j)}_\mu-1/\Delta_1$ is bounded with derivatives at most $c|q|^{-|\alpha|}$, $c$ free of $j$; its kernel falls like $|x|^{-2}$; along one straight piece of $n$ bonds its sum gives $c_{r,j}$ per unit length (the sum along an axis is the integral over the plane $q_\mu=0$, whence the two-dimensional formula) up to $O(a_d(1+\log n))$, and parallel pieces at distance $\ge\delta_C$ interact through it by $O(a_d|C|^2/\delta_C^2)$.

With (a) at the spacing $a_d$ this gives the law at every finite cutoff, with a remainder free of $K$; $c_{r,j}\to c_r$ by dominated convergence, and $|c_{r,j}|\le c_3$ by the bound on the symbol. (c) $\sum_{n\le K-d}\mathcal C_{(n)}=\Pi_0-\Pi_{K-d}$ at every $d$, and $a_K(X_C,\Pi_0X_C)=\mathcal I_{a_K}(C,C)$; insert (a) and (b): $\frac3{2\mathsf b}\frac{|C|}{2\pi}\log\frac{a_d}{a_K}=\kappa\log\frac{a_d}{a_K}$; one step is $a_K(X_C,\mathcal C_{(K-d)}X_C)=v_{K,d+1}-v_{K,d}$.

For families the factors multiply, and one may take $c_\Xi=\prod_i\exp[\frac3{2\mathsf b}(c_3|C_i|+2\mathsf r_i)]$, $\mathsf r_i$ the supremum over $d\ge i_0$ of the bound on $|{\rm Rem}_{K,d}(C_i)|$ in (b), which also bounds $|E_{a_K}(C_i)|$ for $K\ge d$. (d) $v_{K,d}\ge0$, $\Pi_{K-d}$ being an orthogonal projection. It vanishes on the $2$-cochains $Y$ with ${\rm d}^{\sf T}Y=0$, as its range consists of exact cochains; so $v_{K,d}(C)=a_K(X,\Pi_{K-d}X)$ for every $2$-cochain $X$ with ${\rm d}^{\sf T}X=j_C$, in particular for $X$ the indicator of the bare plaquettes of the rectangle, with $\|X\|_1=\sum_p|X(p)|=A/a_K^2$. The maximum of (b) holds at every depth $0\le d\le K$ for this bare current, wherever the corners lie (corners in $\Lambda_d$ are needed only to factor $\hat\jmath$ in the proof of (b)).

Drop the constraint that $\iota$ be closed: the maximum can only grow. With $\mathfrak q=a_d/a_K=L^{K-d}$, the bare momenta $\varkappa=(q+2\pi n)/\mathfrak q$, $q\in(2\pi/N)\mathbb Z^3$, $|n_\nu|\le\frac{\mathfrak q-1}2$, and $p=\mathfrak q\varkappa$, the transform of $Q^{\sf T}\iota$ at $\varkappa$ is $\hat\iota_\mu(q)\varphi_\mu(\varkappa)$, $|\varphi_\mu(\varkappa)|=\mathfrak q|S_{(K-d)}(p_\mu)s_{(K-d)}(p)|$ the form factor of the smeared bond, and $\Delta_1(\varkappa)=\mathfrak q^{-2}\Delta_{(K-d)}(p)$; the maximum over all $\iota$, taken mode by mode, gives
\[v_{K,d}(C)\le a_K(N\mathfrak q)^{-3}\sum_{q\ne0}\sum_\mu\frac{|\mathcal B_\mu(q)|^2}{\mathfrak q^4\mathsf D^{(K-d)}_\mu(q)},\qquad\mathcal B_\mu(q)=\sum_n\frac{\hat\jmath_\mu(\varkappa)\,\overline{\varphi_\mu(\varkappa)}}{\Delta_1(\varkappa)},\]
$\hat\jmath$ the transform of the bare current $j_C$. Three bounds hold uniformly in $\mathfrak q$: $|\hat\jmath_\mu(\varkappa)|\le\Delta_1(\varkappa)^{1/2}\|X\|_1$, as $j_C={\rm d}^{\sf T}X$ and $\sum_\nu|e^{i\varkappa_\nu}-1|^2=\Delta_1(\varkappa)$; $\sum_n|S_{(K-d)}(p_\mu)s_{(K-d)}(p)|\Delta_{(K-d)}(p)^{-1/2}\le c/|q|$, the term $n=0$ being at most $(0.3|q|)^{-1}$ and, for $n\ne0$, each one-dimensional factor at most $51/\max\{1,|n_\nu|\}$ and $\Delta_{(K-d)}(p)\ge0.2|n|^2$, so that the terms $n\ne0$ add up to at most $c\sum_n\prod_\nu\max\{1,|n_\nu|\}^{-4/3}$; and $\mathsf D^{(K-d)}_\mu(q)\ge0.3^8|q|^{-2}$ ($S_{(j)}\ge0.3$ on $[-\frac32\pi,\frac32\pi]$ and $E_{(j)}(z)\le z^2$). So $|\mathcal B_\mu(q)|\le c\,\mathfrak q^2\|X\|_1/|q|$, each of the $N^3$ modes contributes at most $c\|X\|_1^2$, and $v_{K,d}(C)\le c\,a_K\mathfrak q^{-3}\|X\|_1^2=c\,A^2/a_d^3$, as $a_K\mathfrak q=a_d$: $C^{\rm bg}$ is absolute, for every odd $L\ge3$, every $0\le d\le K$ and every position. Equivalently, the kernel of $\Pi_{K-d}$ on the bare plaquettes is at most $c\,(a_K/a_d)^3$ in modulus: at one plaquette this is the bound on the diagonal, and Schwarz' inequality gives the rest.
\end{proof}

\begin{proposition}[The tower]\label{C-372-12}
In the frame let $F=W_{\mathsf C}$ and let $0\le\bar t\le K$ be the last level of the tower.
(a) Suppose $|\int W^{(t)}_{\mathsf C}(V_W(B))\,\nu_W(dB)-\lambda_{t+1}W^{(t+1)}_{\mathsf C}(W)|\le\eta_{t+1}$ for $m_{t+1}$-almost every $W$ and $0\le t<\bar t$, and put $\mathsf e_t=\sum_{s\le t}\Xi_{s-1}\eta_s$. Then $|\Phi^{(t)}[F]-\Xi_tW^{(t)}_{\mathsf C}|\le\mathsf e_t$ $m_t$-almost everywhere for $t\le\bar t$, and
\[\Bigl|\langle F\rangle_K-\Xi_{\bar t}\frac1Z\int W^{(\bar t)}_{\mathsf C}\,dm_{\bar t}\Bigr|\le\mathsf e_{\bar t}+\sum_{t\le\bar t}(\Xi_{t-1}+\mathsf e_{t-1})\mathsf w_t;\]
the same holds with a further factor $H(\mathcal M^{K-s_1}U)$ in $\langle F\rangle_K$ and $H\circ\mathcal M^{K-\bar t-s_1}$ in the integral, for $H$ bounded Borel gauge invariant on $\mathcal V_{s_1}$ and $K-\bar t\ge s_1$, with $\sup|H|$ multiplying the right side. No positivity of $\nu_W$ is used.
(b) Without any hypothesis, let $\mathsf A_t=Z^{-1}\int W^{(t)}_{\mathsf C}\,dm_t$ and $\bar{\mathsf D}_t=Z^{-1}\int[\int W^{(t-1)}_{\mathsf C}(V_W(B))\,\nu_W(dB)-\lambda_tW^{(t)}_{\mathsf C}(W)]\,m_t(dW)$, the mean of the defect of the step over the block field. Then $\mathsf A_{t-1}=\lambda_t\mathsf A_t+\bar{\mathsf D}_t+\omega_t$ with $|\omega_t|\le\mathsf w_t$, and so
\[\langle W_{\mathsf C}\rangle_K=\Xi_{\bar t}\mathsf A_{\bar t}+\sum_{t\le\bar t}\Xi_{t-1}(\bar{\mathsf D}_t+\omega_t):\]
as an identity between numbers, only the means of the defects over the block field enter.
(c) On gauge invariant Borel sets $\mu^K_{K-\bar t}-Z^{-1}m_{\bar t}$ is a positive measure of mass $\sum_{t\le\bar t}\mathsf w_t$, and $1-2\mathsf V_\ell C_{\rm r}(1)L^{-(K-\bar t)}\le Z^{-1}m_{\bar t}(\mathcal V_{K-\bar t})\le1$.
(d) For a coordinate rectangle $C$ with sides $s_1$, $s_2$: $1-\frac12\theta_d^2\le W^{(K-d)}_C\le1$ on $\mathsf O_{K-d}$, with $\theta_d=2B_3\varepsilon_{K-d}L^{2d}s_1s_2$.
\end{proposition}

\begin{proof}[Sketch of proof]
(a) For $m_{t+1}$-almost every $W$,
\[
\begin{multlined}
\Phi^{(t+1)}[F](W)-\Xi_{t+1}W^{(t+1)}_{\mathsf C}(W)=\int[\Phi^{(t)}[F]-\Xi_tW^{(t)}_{\mathsf C}](V_W(B))\,\nu_W(dB) \\
+\Xi_t\Bigl[\int W^{(t)}_{\mathsf C}(V_W(B))\,\nu_W(dB)-\lambda_{t+1}W^{(t+1)}_{\mathsf C}(W)\Bigr].
\end{multlined}
\]
The second member is at most $\Xi_t\eta_{t+1}$ almost everywhere by hypothesis. The first, $\tilde\Psi(W):=\int\Psi_t(V_W(B))\,\nu_W(dB)$ with $\Psi_t:=\Phi^{(t)}[F]-\Xi_tW^{(t)}_{\mathsf C}$, is at most $\mathsf e_t$ in modulus for $m_{t+1}$-almost every $W$ whenever $|\Psi_t|\le\mathsf e_t$ $m_t$-almost everywhere; this duality step reads only $m_t,m_{t+1},\mathsf h_{t+1}\ge0$. The disintegration at $G=\Psi_t$, a bounded Borel function, and the transport identity $\int G'\,dm_{t+1}=\int(G'\circ\mathcal M)\,\mathsf h_{t+1}\,dm_t$ at $G'=|H|$ give, for every bounded Borel $H$,
\[\Bigl|\int H\tilde\Psi\,dm_{t+1}\Bigr|=\Bigl|\int(H\circ\mathcal M)\,\Psi_t\,\mathsf h_{t+1}\,dm_t\Bigr|\le\mathsf e_t\int|H\circ\mathcal M|\,\mathsf h_{t+1}\,dm_t=\mathsf e_t\int|H|\,dm_{t+1};\]
at $H=\overline{{\rm sgn}\,\tilde\Psi}\,\mathbf 1\{|\tilde\Psi|>\mathsf e_t+\epsilon\}$ this says $(\mathsf e_t+\epsilon)\,m_{t+1}(|\tilde\Psi|>\mathsf e_t+\epsilon)\le\mathsf e_t\,m_{t+1}(|\tilde\Psi|>\mathsf e_t+\epsilon)$, so the set has measure $0$, for every $\epsilon>0$. No law of $V_W(B)$, no null set of $m_t$ under $\nu_W$ and no positivity of $\nu_W$ enter. Induction on $t$, with $\mathsf e_{t+1}=\mathsf e_t+\Xi_t\eta_{t+1}$; then \eqref{C-372-2} with $|W^{(t)}_{\mathsf C}|\le1$ and $|\Phi^{(t-1)}[F]|\le\Xi_{t-1}+\mathsf e_{t-1}$. With $H$: $\Phi^{(t)}[W_{\mathsf C}\,H\circ\mathcal M^{K-s_1}]=\Phi^{(t)}[W_{\mathsf C}]\,H\circ\mathcal M^{K-t-s_1}$ for $K-t\ge s_1$. (b) For gauge invariant bounded Borel $G$ on $\mathcal V_{K-t+1}$, \eqref{C-372-2} at $F=G\circ\mathcal M^{t-1}$ for $\bar t=t-1$ and $\bar t=t$, with $\Phi^{(t-1)}[G\circ\mathcal M^{t-1}]=G$ (functions pull out and $\Phi^{(t-1)}[1]=1$), gives $\int G\,dm_{t-1}=\int G\,(\mathsf h_t+\mathsf h^{\rm w}_t)\,dm_{t-1}$.

At $G=W^{(t-1)}_{\mathsf C}$ the disintegration with $H=1$ turns $Z^{-1}\int W^{(t-1)}_{\mathsf C}\mathsf h_t\,dm_{t-1}$ into $\lambda_t\mathsf A_t+\bar{\mathsf D}_t$, and $\omega_t=Z^{-1}\int W^{(t-1)}_{\mathsf C}\mathsf h^{\rm w}_t\,dm_{t-1}$; $\mathsf A_0=\langle W_{\mathsf C}\rangle_K$. (c) \eqref{C-372-2} at $F=\mathbf 1_E\circ\mathcal M^{\bar t}$, $E$ gauge invariant: $Z\mu^K_{K-\bar t}(E)-m_{\bar t}(E)=\sum_{t\le\bar t}\int(\mathbf 1_E\circ\mathcal M^{\bar t-t+1})\mathsf h^{\rm w}_t\,dm_{t-1}\ge0$, of total mass $Z\sum_{t\le\bar t}\mathsf w_t$; the rarity bound at $p=0$ in its summed form. (d) Over a tame field the background has bare plaquettes $<2B_3\varepsilon_{K-d}(a_K/a_d)^2$ (Theorem~\YMMref{E}{10} in the setting of Section~\YMMref{52.7}{432}), and $1-W_{\partial R}\le(\sum_{p\subset R}\sqrt{a_p})^2$ over the $s_1s_2/a_K^2$ plaquettes of the rectangle (Proposition~\ref{C-372-13}(c)).
\end{proof}

\begin{proposition}[Identities and inequalities at every cutoff]\label{C-372-13}
For every $\beta>0$, every torus and every cutoff:
(a) $\mathcal W^{(L^z\ell,L^z\mathsf b)}_{K-z}(L^z\mathsf C;f^{(z)})=L^{-z\kappa-zn}\mathcal W^{(\ell,\mathsf b)}_K(\mathsf C;f_1,\dots,f_n)$, $f^{(z)}_q=f_q(L^{-z}\cdot)$;
(b) if families $G_1,\dots,G_s$ lie in the interior of $S^+_\theta$ for a coordinate mirror $\theta$ and the supports of the tuples $f_i$ lie in $S^+_\theta$ (and do not meet the fixed planes if $\theta$ is a bond mirror), then $\sum_{i,j}\bar c_ic_j\mathcal W_K(\theta G_i,G_j;\theta f_i,f_j)\ge0$ for all complex $c_i$; so $\mathcal W_K(G_1,\theta G_2)^2\le\mathcal W_K(G_1,\theta G_1)\,\mathcal W_K(G_2,\theta G_2)$ for real entries; if $\theta$ cuts a loop $C_{12}$ into two halves and $C_{11}$, $C_{22}$ are the loops made of each half and its mirror image, then $\mathcal W_K(C_{12})^2\le\mathcal W_K(C_{11})\,\mathcal W_K(C_{22})$;
(c) the plaquette mean $\langle W_{\partial p}\rangle_K$ is $>0$; for bare rectangles $p\times q$ with $p,q<\ell/2a_K$, $w(p,q)=\langle W_{p\times q}\rangle_K\ge w(1,1)^{pq}>0$, and $\log w$ is convex in $p$ and in $q$; and $0\le1-W_{\partial R}\le(\sum_{p\subset R}\sqrt{a_p})^2$ for every field and every rectangle $R$.
For the rectangles $s\times T$ of one coordinate plane, $s$ and $T$ lengths in $\bigcup_ia_i\mathbb Z$, put $\mathcal W(s\times T):=\lim_K\mathcal W_K(s\times T)=\lim_Ka_K^{-\kappa}\langle W_{s\times T}\rangle_K$, $\kappa=\kappa(s\times T)$, wherever the limit exists. A \emph{range} of the limits is a set of such rectangles on which they exist and are finite, closed downwards (with $s\times T$ it holds every $s'\times T'$ with $s'\le s$, $T'\le T$) and open upwards (with $s\times T$ it holds a rectangle with both sides larger); in the frame, and by dilation on every torus and at every label, the polygons $s\times T$ with $\kappa<\frac12$ form one (Theorem~\ref{C-372-4}(1)--(3)). Every limit inherits (a), (b), $\mathcal W\ge0$ and the midpoint inequality $\mathcal W(s\times\frac{T'+T''}2)^2\le\mathcal W(s\times T')\,\mathcal W(s\times T'')$, and the same in $s$, whenever the three rectangles lie in a range, the powers of $a_K$ matching because $2|s\times\frac{T'+T''}2|=|s\times T'|+|s\times T''|$; a limit that is $>0$ at one rectangle of a range is $>0$ at all rectangles of that range. Strict positivity of a limit is not inherited; it comes from the floor of Theorem~\ref{C-372-4}(4).
\end{proposition}

\begin{proof}[Sketch of proof]
(a) The Wilson measure at $(K,\ell,\mathsf b)$ is the one at $(K-z,L^z\ell,L^z\mathsf b)$ relabelled, as $\beta=L^K\mathsf b=L^{K-z}(L^z\mathsf b)$: lengths are multiplied by $L^z$, $\kappa$ is unchanged, $a_K^{-\kappa}$ gains $L^{-z\kappa}$ and each $\mathcal O(f)$ the factor $L^{-z}$. (b) Reflection positivity of the Wilson measure at site and bond mirrors, for functions of the bonds of one closed (resp.\ open) half torus, and Schwarz' inequality: each $a_K^{-\kappa(G_i)}W_{G_i}\prod_q(\mathcal O(f_{iq})-\langle\mathcal O(f_{iq})\rangle_K)$ lives on one side of $\theta$, and the powers of $a_K$ split because lengths add. For the cut loop, $W_{C_{12}}=\frac12\operatorname{tr}(U(\gamma_1)U(\gamma_2))$ with the two halves $\gamma_1$, $\gamma_2$, and reflection positivity for matrix-valued functions with Schwarz' inequality gives the bound; $|C_{11}|+|C_{22}|=2|C_{12}|$. (c) $s\mapsto\log\int e^{s\sum_p\frac12\operatorname{tr}\partial U(p)}dU$ has the derivative $0$ at $s=0$ (the Haar mean of a non-trivial character vanishes) and is strictly convex, so its derivative at $s=\beta$, the number of plaquettes times the plaquette mean, is $>0$.

Reflection positivity in the two directions of the plane of the rectangle gives $w(p,q)^2\le w(p,q-1)\,w(p,q+1)$ and the same in $p$; with $w(p,0)=1$ the ratios $w(p,q)/w(p,q-1)$ increase, so $w(p,q)\ge w(p,1)^q\ge w(1,1)^{pq}$. For $g\in\mathrm{SU}(2)$, $1-\frac12\operatorname{tr}g=\frac12|g-\mathbf 1|^2$; the holonomy of $\partial R$ is an ordered product of conjugated plaquette variables, and $|\prod_ig_i-\mathbf 1|\le\sum_i|g_i-\mathbf 1|$. Limits keep $\ge0$, the positive forms and the inequalities; the midpoint inequality of $\mathcal W_K$ is the convexity of $\log w$ in each side, multiplied by $a_K^{-2\kappa(s\times\frac{T'+T''}2)}=a_K^{-\kappa(s\times T')-\kappa(s\times T'')}$, at every cutoff at which the three rectangles have their corners in $\Lambda_K$.

Spreading in a range $\mathsf R$, the limits being finite there (Theorem~\ref{C-372-4}(1)): if $\mathcal W(s\times T_0)>0$ and $s\times(T_0\pm\epsilon)\in\mathsf R$, the product $\mathcal W(s\times(T_0-\epsilon))\,\mathcal W(s\times(T_0+\epsilon))\ge\mathcal W(s\times T_0)^2>0$ of two finite numbers $\ge0$ makes both positive; going down from the lowest positive point with an upper point of $\mathsf R$ fixed, the distance to it doubles at each step, so that every $T$ below $T_0$ is reached, and upwards in the same way with a lower point; then the same in $s$; and two rectangles of $\mathsf R$ are joined through the one with the smaller sides, which lies in $\mathsf R$. The values of $\epsilon$ are taken in $\bigcup_ia_i\mathbb Z$.
\end{proof}

\begin{proposition}[A source off the loop]\label{C-372-14}
In the frame let $f_1,\dots,f_n$ be real Lipschitz, $\|f\|_\sharp=\max\{\sup|f|,40M\,{\rm Lip}f\}$, and $z$ in the polydisc $\sum_q|z_q|\|f_q\|_\sharp<\rho_D$, where $\rho_D=2\varrho_S/(2e^{(C_S+C_g)\ell^3}+2)$ is formed from the radius $\varrho_S$ and the constants $C_S$, $C_g$ of Theorems~\YMMref{AJ}{673} and~\YMMref{AD}{607}; $\rho_D$ depends on neither $K$ nor $\mathsf b\in\mathsf I(g)$. Write $w_z=\sum_qz_qf_q$. The radius estimates below are stated for a family with $\max_q\|f_q\|_\sharp>0$. If every $f_q$ vanishes, the source is zero, its derivatives of positive order vanish, and the assertions reduce directly to the zero-source statements.
(a) $\langle e^{\mathcal O(w_z)}\rangle_K=e^{\mathcal E(w_z)}h_K(z)$ with $\frac12\le|h_K(z)|\le\frac32$, whatever the supports, $\mathcal E(w)$ a sum of local energies of the source; the energies of sources $v_i=\sum_qz^{(i)}_qf^{(i)}_q$, $i=1,\dots,s$, $f^{(i)}_q$ real Lipschitz, with $\sum_i\sum_q|z^{(i)}_q|\|f^{(i)}_q\|_\sharp<\rho_D$ and supports at pairwise sup-distance $\ge2\delta$, add up to an error: $|\mathcal E(\sum_iv_i)-\sum_i\mathcal E(v_i)|\le\Delta_s(\delta)=c(s+1)(\ell/2\delta)^3$, $c$ formed from the constants of Theorems~\YMMref{AJ}{673} and~\YMMref{AD}{607}, free of $K$; so, with $\mathsf e(v)=|\langle e^{\mathcal O(v)}\rangle_K|$, $\frac12(\frac23)^se^{-\Delta_s(\delta)}\prod_i\mathsf e(v_i)\le\mathsf e(\sum_iv_i)\le\frac32\,2^se^{\Delta_s(\delta)}\prod_i\mathsf e(v_i)$; we write $\Delta(\delta)=\Delta_2(\delta)$. In particular $\langle e^{\mathcal O(w_z)}\rangle_K\ne0$ on the polydisc, for all $K$.
(b) The objects of the frame exist with the weight $e^{\mathcal O(w)}$. For complex bounded $w$ and bounded Borel $F$ on $\mathcal V_K$ let $m^w_t[F]=\Phi^{(t)}[Fe^{\mathcal O(w)}]\,m_t$, a complex measure on $\mathcal V_{K-t}$, and $m^w_t=m^w_t[1]$. Then $\int G\,dm^w_t[F]=\int(G\circ\mathcal M)\,\mathsf h_t\,dm^w_{t-1}[F]$ for bounded Borel $G$, and, for gauge invariant $F$ and $G$ and $0\le\bar t\le K$, the sum over the level of the first creation of a large field
\[
\begin{aligned}
Z\bigl\langle F\,G(\mathcal M^{\bar t}U)\,e^{\mathcal O(w)}\bigr\rangle_K
&=\int G\,dm^w_{\bar t}[F]\\
&\quad+\sum_{t\le\bar t}\int(G\circ\mathcal M^{\bar t-t+1})\,\mathsf h^{\rm w}_t\,dm^w_{t-1}[F].
\end{aligned}
\]
At real $w$ only, $\Phi^{(t)}[e^{\mathcal O(w)}]>0$ $m_t$-almost everywhere, and the quotients 
\[
\Phi^{(t),w}[F]=\Phi^{(t)}[Fe^{\mathcal O(w)}]/\Phi^{(t)}[e^{\mathcal O(w)}],
\]
 with $|\Phi^{(t),w}[F]|\le\sup|F|$, and the kernel 
\[
\nu^w_W(dB)=\Phi^{(t)}[e^{\mathcal O(w)}](V_W(B))\,\nu_W(dB)/\Phi^{(t+1)}[e^{\mathcal O(w)}](W),
\]
a probability measure for $m_{t+1}$-almost every $W$, obey \eqref{C-372-2} and the disintegration for $(m^w,\nu^w)$, with $Z\langle e^{\mathcal O(w)}\rangle_K$ in the place of $Z$; $m^w_t$ is then a positive measure with the null sets of $m_t$, and on gauge invariant sets $\mu^{K,w}_{K-\bar t}-m^w_{\bar t}/(Z\langle e^{\mathcal O(w)}\rangle_K)$ is a positive measure of mass $\sum_{t\le\bar t}\mathsf w^w_t$, $\mathsf w^w_t=(Z\langle e^{\mathcal O(w)}\rangle_K)^{-1}\int\mathsf h^{\rm w}_t\,dm^w_{t-1}$, $\mu^{K,w}_d$ being the law of $\mathcal M^{K-d}U$ under $e^{\mathcal O(w)}\langle\cdot\rangle_K/\langle e^{\mathcal O(w)}\rangle_K$.

The rarity bound holds with a source, at complex $z$ for the integral of $\mathsf h^{\rm w}_t$: 
\[
|Z\langle e^{\mathcal O(w)}\rangle_K|^{-1}|\int\mathsf h^{\rm w}_t\,dm^w_{t-1}|\le\mathsf V_SC^S_{\rm r}(p)L^{-p(K-t)},
\]
with $\mathsf V_S$, $C^S_{\rm r}(p)$ formed from the constants of Theorems~\YMMref{D}{8}, \YMMref{AJ}{673} and~\YMMref{AD}{607}, $p$ and $\ell$; at real $z$ this bounds $\mathsf w^w_t$, and $\sum_{t\le\bar t}L^{\kappa(K-t)}\mathsf w^w_t\le2\mathsf V_SC^S_{\rm r}(\kappa+1)L^{-(K-\bar t)}$.
(c) For bounded Borel gauge invariant $F$ on $\mathcal V_d$, 
\[
\mu^{K,w_z}_d(F)=\langle F(\mathcal M^{K-d}U)e^{\mathcal O(w_z)}\rangle_K/\langle e^{\mathcal O(w_z)}\rangle_K
\]
 obeys $|\mu^{K,w_z}_d(F)|\le B_d\sup|F|$, with $B_d$ formed from the constants of Theorems~\YMMref{D}{8}, \YMMref{AJ}{673} and~\YMMref{AD}{607}, $d$, $\mathsf b$ and $\ell$, and free of $K$. The limit $\mu^{w_z}_d(F)=\lim_K\mu^{K,w_z}_d(F)$ exists for $z$ in the polydisc $|z_q|<\rho_*:=\rho_D/(n\max_q\|f_q\|_\sharp)$, which lies in the polydisc of the head and is free of $K$, uniformly on its half $|z_q|\le\frac12\rho_*$ and in $\sup|F|\le1$; it is holomorphic in $z$, a probability measure at real $z$, and $\mu^w_{d+1}(\mathcal M^{-1}E)=\mu^w_d(E)$. No condition on the supports is asked; so the derivatives at $z=0$, the mixed functions of $F(\mathcal M^{K-d}U)$ and of the $\mathcal O(f_q)$, converge too.
(d) If a coordinate mirror $\theta$ parts $\mathsf C$, at distance $\ge\delta_1$ from its planes, from the supports, at distance $\ge\delta$, and $z$ lies in half the polydisc, then for every bounded function $X$ of the bonds on the side of $\mathsf C$
\[|\langle Xe^{\mathcal O(w_z)}\rangle_K|\le M_0\|X\|_\theta\,|\langle e^{\mathcal O(w_z)}\rangle_K|,\qquad\|X\|_\theta^2=\langle X\,\overline{\theta X}\rangle_K\ge0,\qquad M_0=\sqrt6\,e^{\Delta(\delta)/2};\]
so $|Q_K(z)|\le M_0\,\mathcal W_K(\mathsf C,\theta\mathsf C)^{1/2}$. For two orthogonal mirrors, the sources in two groups as in Theorem~\ref{C-372-4}(6), and $z$ in a quarter of the polydisc, $\sum_q|z_q|\|f_q\|_\sharp<\rho_D/4$ (for one mirror, half the polydisc is $\sum_q|z_q|\|f_q\|_\sharp<\rho_D/2$), 
\[
|\langle Xe^{\mathcal O(w_z)}\rangle_K|^4\le M_0'^4\langle X\,\overline{\theta_1X}\,\overline{\theta_2X}\,\theta_1\theta_2X\rangle_K\,|\langle e^{\mathcal O(w_z)}\rangle_K|^4
\]
 for $X$ on the side of $\mathsf C$ of both mirrors, $M_0'^4=24\cdot36\cdot(\frac92)^4e^{\Delta_4(\delta)+6\Delta_2(\delta)}$, the mean on the right holding no source and being $\ge0$. Where equally spaced coordinate walls part the polygons from the supports as in Theorem~\ref{C-372-4}(6), with an even number of boxes in each direction that the walls cut, so that both fixed planes $c$ and $c+\frac\ell2$ of every wall mirror are walls ($N$ boxes in all; on $\ell=2L^m$, with walls in $\frac12a_i\mathbb Z$, the number of slabs in a cut direction divides $4L^{m+i}$, and it is taken even, $2L^j$ or $4L^j$), and with each $f_q$ supported in one box: $\langle W_{\mathsf C^{\rm f}_\iota}\rangle_K\ge0$, $\kappa(\mathsf C^{\rm f}_\iota)=N\kappa(\mathsf C_\iota)$, and
\[
\begin{multlined}
|Q_K(z)|\le M_S\prod_\iota\mathcal W_K(\mathsf C^{\rm f}_\iota)^{1/N}\quad\text{for }\sum_q|z_q|\|f_q\|_\sharp<\rho_D/N,\qquad M_S \\
=2(\tfrac32)^{n_{\rm b}}e^{\Delta_{n_{\rm b}}(\delta)}\bigl[\tfrac32\,2^Ne^{\Delta_N(\delta)}\bigr]^{n_{\rm b}/N},
\end{multlined}
\]
$n_{\rm b}$ the number of boxes that meet the supports and $\delta$ the distance of the supports to the walls; a function whose support misses the walls is the sum of such pieces, one for each box.
(e) $\mathcal W_K(\mathsf C;f_1,\dots,f_n)=\sum_{A}\partial_AQ_K(0)\,\partial_{A^c}E_K(0)$, $E_K(z)=\langle\exp\sum_qz_q(\mathcal O(f_q)-\langle\mathcal O(f_q)\rangle_K)\rangle_K$, and $\partial_{A^c}E_K(0)=\sum_\pi\prod_{\beta\in\pi}S^{\sf T}_{|\beta|,K}(f_\beta)$, $\pi$ over the partitions of $A^c$ into blocks of at least two elements: only the derivatives of $Q_K$ in pairwise different variables enter.
(f) Let $0<\vartheta'\le1$, $\bar t(K)=K-\lceil\vartheta'K\rceil$, $\eta=(24\ell^3\sum_q\sup|f_q|)^{-1}$, and let $q_{K,\bar t}(z)=a_K^{-\kappa}\int\Phi^{(\bar t)}[W_{\mathsf C}e^{\mathcal O(w_z)}]\,dm_{\bar t}\big/\int\Phi^{(\bar t)}[e^{\mathcal O(w_z)}]\,dm_{\bar t}$ be the same quotient formed with the measure $m_{\bar t}$ without large fields. For every multi-index $\alpha$, with $p=(4|\alpha|+\kappa+1)/\vartheta'$, and all $K$ with $2\mathsf V_\ell C_{\rm r}(p)L^{-p\vartheta'K}\le0.16$: $|\partial^\alpha Q_K(0)-\partial^\alpha q_{K,\bar t(K)}(0)|\le14\,\alpha!\,\eta^{-|\alpha|}\mathsf V_\ell C_{\rm r}(p)L^{-K}$: the Taylor coefficients of $Q_K$ at $0$ are those of the chain without large fields.
\end{proposition}

\begin{proof}[Sketch of proof]
(a) Theorems~\YMMref{AJ}{673} and~\YMMref{AD}{607} for the Wilson measure with the source $w_z$ writes $\langle e^{\mathcal O(w_z)}\rangle_K$ as $e^{\mathcal E(w_z)}h$, $\mathcal E(w)=\sum_{j\le K}\sum_Xe_j(X;w)$ a sum of local terms $e_j(X;w)$ over the levels $j$ and the domains $X$ of the level $j$, which read $w$ near $X$ alone and vanish at $w=0$, with $|e_j(X;w)|\le C_Se^{-\kappa_Sd_j(X)}$, $\kappa_S$ the decay number of Theorems~\YMMref{AJ}{673} and~\YMMref{AD}{607} and $d_j(X)$ the size of $X$ in cubes of side $M$ of the level $j$; and $h$ holomorphic with $h(0)=1$ and $|h|\le e^{(C_S+C_g)\ell^3}$. With $\|z\|=\sum_q|z_q|\|f_q\|_\sharp$, Schwarz' lemma for $\zeta\mapsto h(w_{\zeta z})-1$ on $|\zeta|<2\varrho_S/\|z\|$ gives $|h-1|\le\frac12$ on the polydisc.

In $\mathcal E(\sum_iv_i)-\sum_i\mathcal E(v_i)$ a local term that sees one support alone cancels; one that sees two supports at sup-distance $\ge2\delta$ belongs to a domain $X$ with $d_j(X)\ge2\delta L^{K-j}/M-4$ and is at most $(s+1)C_Se^{-\kappa_Sd_j(X)}$; the sum over such domains and levels is at most $c(s+1)(\ell/2\delta)^3$. The bounds on $\mathsf e(\sum_iv_i)$ follow, as $\frac12e^{{\rm Re}\,\mathcal E(v)}\le\mathsf e(v)\le\frac32e^{{\rm Re}\,\mathcal E(v)}$ for every source $v$ of the polydisc. (b) A history is linear in the bare density, so $\Phi^{(t)}[Fe^{\mathcal O(w)}]$ is the density of the history without large fields with the source, at complex $w$ as well. The transport identity is that of the frame at $Fe^{\mathcal O(w)}$, and the sum over the first creation is \eqref{C-372-2} at the gauge invariant function $F\,(G\circ\mathcal M^{\bar t})\,e^{\mathcal O(w)}$, the factor $G\circ\mathcal M^{\bar t}$ pulling out of $\Phi^{(\bar t)}$ and of $\Phi^{(t-1)}$ as $G$ and as $G\circ\mathcal M^{\bar t-t+1}$.

The quotients and the kernel are formed at real $w$: there $e^{\mathcal O(w)}\ge e^{-\sup|\mathcal O(w)|}>0$ gives, at every level $t$, $\Phi^{(t)}[e^{\mathcal O(w)}]\ge e^{-\sup|\mathcal O(w)|}$ off one $m_t$-null set $\mathcal N_t$, so $m^w_t$ has the null sets of $m_t$. The kernel $\nu^w_W$ is $\nu_W$ times the density $B\mapsto\Phi^{(t)}[e^{\mathcal O(w)}](V_W(B))/\Phi^{(t+1)}[e^{\mathcal O(w)}](W)$, and it is a probability measure for $m_{t+1}$-almost every $W$: $\nu_W\ge0$ for $m_{t+1}$-almost every $W$ (Proposition~\ref{C-372-10}(b)); the denominator is $\ge e^{-\sup|\mathcal O(w)|}$ for $W\notin\mathcal N_{t+1}$; the numerator is $\ge e^{-\sup|\mathcal O(w)|}$ off $V_W^{-1}\mathcal N_t$, and $\nu_W(V_W^{-1}\mathcal N_t)=0$ for $m_{t+1}$-almost every $W$, the one set $\mathcal N_t$ being fixed first and the exceptional set of $W$ depending on it: the disintegration at $G=\mathbf 1_{\mathcal N_t}$ gives $\int H(W)\,\nu_W(V_W^{-1}\mathcal N_t)\,m_{t+1}(dW)=\int(H\circ\mathcal M)\,\mathbf 1_{\mathcal N_t}\,\mathsf h_{t+1}\,dm_t=0$ for every bounded Borel $H$; and the mass is $1$ by the disintegration at $F=e^{\mathcal O(w)}$, $\Phi^{(t+1)}[e^{\mathcal O(w)}](W)=\int\Phi^{(t)}[e^{\mathcal O(w)}](V_W(B))\,\nu_W(dB)$ for $m_{t+1}$-almost every $W$.

The disintegration for $(m^w,\nu^w)$ is that of the frame at $G\,\Phi^{(t)}[e^{\mathcal O(w)}]$. The wild mass is a sum over the histories of the level $K$ in which a large field is created at the step $t$; Theorems~\YMMref{AJ}{673} and~\YMMref{AD}{607} bounds every history with a source as Theorem~\YMMref{D}{8} bounds it without one, with the constants of Theorems~\YMMref{AJ}{673} and~\YMMref{AD}{607}, and (a) bounds the denominator below; the rarity bound then follows as in the square-completion estimate \eqref{eq:ward-age-square} and the geometric-series calculation in the proof of Proposition~\ref{prop:ward-fixed-volume-rarity}, with the constants of Theorems~\YMMref{AJ}{673} and~\YMMref{AD}{607}. (c) Theorems~\YMMref{AJ}{673} and~\YMMref{AD}{607} at one depth bounds $|\langle F(\mathcal M^{K-d}U)e^{\mathcal O(w_z)}\rangle_K|$ by $\sup|F|$ times a number free of $K$ times $|\langle e^{\mathcal O(w_z)}\rangle_K|$: this is $B_d$.

For the limit let $\sup|F|\le1$ and $G_K(z)=\mu^{K,w_z}_d(F)$, holomorphic and bounded by $B_d$ on the polydisc $|z_q|<\rho_*$. We use an interpolation inequality: if $G$ is holomorphic on the polydisc of radius $\rho$ in $n$ variables, $|G|\le B$ there and $|G|\le\sigma$ on the real cube $[-r,r]^n$, $0<r\le\rho/2$, then its Taylor coefficients at $0$ obey $|c_\alpha(G)|\le C_{n,j}r^{-|\alpha|}(\sigma+2B(r/\rho)^{j+1})$ for $|\alpha|\le j$ (the Taylor polynomial of order $j$ is within $2B(r/\rho)^{j+1}$ of $G$ on the cube, by Cauchy's estimates on the complex lines through $0$, and on the polynomials of order at most $j$ in $n$ variables the largest coefficient and the supremum on a cube are two norms).

It suffices that every coefficient $c_\alpha(G_K)$ converges, uniformly in $F$: Cauchy's estimate $|c_\alpha(G_K)|\le B_d\rho_*^{-|\alpha|}$ makes the tails of the Taylor series small on $|z_q|\le\frac12\rho_*$, uniformly in $K$ and $F$. For $\alpha=0$ this is Corollary~\textup{\ref{IV-62-66}}. Fix $j=|\alpha|\ge1$. The proof of the limit laws of the block fields compares two consecutive cutoffs through the small-field chains at the depth $u=\lceil\phi_1K\rceil$, $\phi_1=1/(180(3N'+4))$, $N'$ the number of the test functions (Section~\ref{sec:cutoff-calibration}); we prolong the tuple $f$ by copies of $f_1$ to $N'=\max\{n,j^2\}$ functions, the new variables being $0$, so that $u\ge d$ for large $K$. Let $G^{\rm sf}_K$ be the quotient of $F$ formed with the weight of the history without large-field regions down to the depth $u$, and $D_K=G_K-G^{\rm sf}_K$; then $c_\alpha(G_{K+1})-c_\alpha(G_K)=c_\alpha(D_{K+1})-c_\alpha(D_K)+c_\alpha(G^{\rm sf}_{K+1}-G^{\rm sf}_K)$, and the interpolation inequality is used twice, on two polydiscs. (i) On the polydisc of radius $\rho'=(24n\ell^3L^{4(K+1)}\max_q\sup|f_q|)^{-1}$, on which $|\mathcal O(w_z)|\le\frac14$ at the cutoffs $K$ and $K+1$ (there are $3(\ell/a_K)^3$ plaquettes and $0\le a_p\le2$), $|G_K|,|G^{\rm sf}_K|<1.8$ and $|D_K|\le3.6$; at its real points the measure without large-field regions is below the full one with a mass defect of the order of the rarity at the depth $u$, so that $|D_K|\le C\exp[3K\log L-cK^2]$, below every power of $L^K$; with $r=L^{-(4j+1)K}\rho'$ the coefficients of $D_K$ are at most $C(L^{(4j^2+5j+3)K}e^{-cK^2}+L^{-K})$, summable in $K$. (ii) On the polydisc on which the proof of the limit laws bounds the small-field chains, of radii $\rho''$ with $1/\rho''\le cK^{2\mathfrak p_0}L^{3u}$, so that $\rho''^{-j}\le(cK^{2\mathfrak p_0})^jL^{3j+K/(180j)}$ by $3\phi_1j\le1/(180j)$, which is where $N'\ge j^2$ enters, $|G^{\rm sf}_K|\le1.8$; at its real points that proof goes through with the source, the contact terms of the source standing in a factor that drops out of the quotient, and gives $|G^{\rm sf}_{K+1}-G^{\rm sf}_K|\le cK^{2\mathfrak p_0+1/2}L^{-K/10}$; with $r=L^{-K/(90j)}\rho''$ the coefficients are at most a power of $K$ times $L^{K(-\frac1{10}+\frac1{90}+\frac1{180j})}+L^{K(-\frac1{90j}+\frac1{180j})}$, summable in $K$, both exponents being negative.

So every $c_\alpha(G_K)$ converges with a summable rate, uniformly in $\sup|F|\le1$, and $G_K$ converges uniformly on $|z_q|\le\frac12\rho_*$; the limit is holomorphic, a complex measure of mass $1$ in $F$ (a limit, uniform in $\sup|F|\le1$, of such measures), a probability measure at real $z$, where the $\mu^{K,w}_d$ are, and $\mathcal M$-compatible, as $\mu^{K,w}_{d+1}(\mathcal M^{-1}E)=\mu^{K,w}_d(E)$ at every $K$. (d) Reflection positivity at $\theta$ and Schwarz' inequality give $|\langle XH\rangle_K|^2\le\|X\|_\theta^2\langle\overline{\theta H}H\rangle_K$ for $H$ on the other side; at $H=e^{\mathcal O(w)}$, $\overline{\theta H}H=e^{\mathcal O(w+w^\theta)}$ with $w^\theta=\overline{\theta w}=\sum\bar z_q\theta f_q$.

The mean, not the energy, is invariant: $\langle e^{\mathcal O(w^\theta)}\rangle_K=\overline{\langle e^{\mathcal O(w)}\rangle_K}$, as $\theta$ is a symmetry of the Wilson measure and the $a_p$ are real. By (a) for $w$, $w^\theta$ and $w+w^\theta$, which lie in the polydisc of the head as $z$ lies in its half and whose supports lie $2\delta$ apart, $\langle e^{\mathcal O(w+w^\theta)}\rangle_K\le\frac32e^{\Delta}e^{{\rm Re}\,\mathcal E(w)+{\rm Re}\,\mathcal E(w^\theta)}$, and $e^{{\rm Re}\,\mathcal E(v)}\le2|\langle e^{\mathcal O(v)}\rangle_K|$ at $v=w$ and $v=w^\theta$; so $\langle e^{\mathcal O(w+w^\theta)}\rangle_K\le6e^{\Delta}|\langle e^{\mathcal O(w)}\rangle_K|^2$. At $X=W_{\mathsf C}$, $\|W_{\mathsf C}\|_\theta^2=\langle W_{\mathsf C}W_{\theta\mathsf C}\rangle_K=a_K^{2\kappa}\mathcal W_K(\mathsf C,\theta\mathsf C)$.

Two mirrors: Schwarz at $\theta_1$ with $Xe^{\mathcal O(w_2)}$ on the side of $\mathsf C$, then at $\theta_2$; the sources that arise hold four images of the second group and two of the first, pairwise $2\delta$ apart, and together they stay in the polydisc of (a), $\sum_q|z_q|\|f_q\|_\sharp<\rho_D$ with each $z_q$ counted as often as its images, as soon as $z$ lies in a quarter of it; then (a) bounds the source factors. Walls: in a cut direction with an even number $n$ of slabs the reflections in the walls map walls to walls, both fixed planes of each being walls, and generate the $n/2$ translations by even multiples of the width and $n/2$ reflections, a group of $n$ elements that acts simply transitively on the slabs; so each box has exactly one image in every box, an image has a parity (the number of reflections modulo $2$), and the folds below are defined. (With an odd number of slabs the translation by one width lies in the group, and so does the reflection of a box in its own mid-plane, which is no wall; the evenness excludes this.) The chessboard estimate $|\langle\prod_\iota F_\iota\rangle_K|\le\prod_\iota\langle{\rm Fold}\,F_\iota\rangle_K^{1/N}$ for functions $F_\iota$ of the bonds of the boxes, ${\rm Fold}\,F_\iota$ the product of the $N$ images of $F_\iota$ under the group (conjugated at the odd ones), each $\langle{\rm Fold}\,F_\iota\rangle_K$ being $\ge0$, is applied with $F_\iota=W_{\mathsf C_\iota}$ for a box with polygons, $F_\iota=e^{\mathcal O(w_\iota)}$ for a box with a support, $w_\iota$ the part of $w$ in it, and $F_\iota=1$ otherwise. ${\rm Fold}\,W_{\mathsf C_\iota}=W_{\mathsf C^{\rm f}_\iota}$, with $\kappa(\mathsf C^{\rm f}_\iota)=N\kappa(\mathsf C_\iota)$, so that the powers of $a_K$ split. ${\rm Fold}\,e^{\mathcal O(w_\iota)}=e^{\mathcal O(\sum_gw^g_\iota)}$, the $N$ images $w^g_\iota$ lying in different boxes, pairwise $2\delta$ apart, each with $\mathsf e(w^g_\iota)=\mathsf e(w_\iota)$ and together in the polydisc of (a) for $\sum_q|z_q|\|f_q\|_\sharp<\rho_D/N$; by (a) at $s=N$, $\langle{\rm Fold}\,e^{\mathcal O(w_\iota)}\rangle_K^{1/N}\le[\frac32\,2^Ne^{\Delta_N(\delta)}]^{1/N}\mathsf e(w_\iota)$.

In the denominator, (a) at $s=n_{\rm b}$ gives $|\langle e^{\mathcal O(w)}\rangle_K|\ge\frac12(\frac23)^{n_{\rm b}}e^{-\Delta_{n_{\rm b}}(\delta)}\prod_\iota\mathsf e(w_\iota)$; the factors $\mathsf e(w_\iota)$ cancel and $M_S$ remains. (e) Leibniz' rule for $\partial_{\{1,\dots,n\}}[Q_K(z)E_K(z)]$ at $0$; the derivatives of $E_K$ at $0$ in pairwise different variables are sums over the partitions of products of joint cumulants, the blocks of one element vanishing because the $\mathcal O(f_q)$ are centred. (f) \eqref{C-372-2} at $F=W_{\mathsf C}e^{\mathcal O(w_z)}$ and at $F=e^{\mathcal O(w_z)}$, on the polydisc $|z_q|\le\eta a_K^4$, where $|\mathcal O(w_z)|\le\frac14$ (there are $3(\ell/a_K)^3$ plaquettes and $0\le a_p\le2$), together with the rarity bound at $p$, bounds the difference of the two quotients by $c\mathsf V_\ell C_{\rm r}(p)a_K^{-\kappa}L^{-p(K-\bar t)}$; Cauchy's estimate on this polydisc costs $\alpha!(\eta a_K^4)^{-|\alpha|}$, and $p(K-\bar t)\ge(4|\alpha|+\kappa+1)K$.
\end{proof}

\begin{proposition}[Three free photon fields]\label{C-372-15}
(a) Uniformly in $K$, $\mathbb P(\max_p|F^{(d)}(p)|\ge2\varepsilon_{[d]})\le18\ell^3L^{3d}e^{-\frac23A_0^2(\log x_{[d]})^{2\mathfrak p_0}}$, and the probability that some plaquette with a corner in a set $S$ of sites of $\Lambda_d$ has $|F^{(d)}(p)|\ge2\varepsilon_{[d]}$ is at most $72\,\#S\,e^{-\frac23A_0^2(\log x_{[d]})^{2\mathfrak p_0}}$. These are bounds at one depth, uniform in $K$: the first has the form of $R^\natural_d$, with $18$ and $\frac23$ in the place of $1$ and $c^\natural$, the second the local form of Proposition~\ref{C-372-8}(c), with $72$ and $\frac23$ in the place of $1$ and $c^\natural$; summed over the depths $d\ge D$ (a union bound) they bound the probability that some depth $d\ge D$ has a plaquette with $|F^{(d)}(p)|\ge2\varepsilon_{[d]}$. There the smoothing is linear and needs no ball.
(b) For a family $\mathsf C=(C_1,\dots,C_r)$, at every length,
\[a_K^{-\kappa}\,\mathbb E\prod_iW^{\rm fr}_{C_i}\to\prod_ie^{-\frac3{2\mathsf b}[c_{\rm lat}|C_i|+\mathcal J_\ell(C_i)]}\Bigl[2^{-r}\sum_{\epsilon\in\{\pm1\}^r}e^{-\frac1{\mathsf b}\sum_{i<j}\epsilon_i\epsilon_j\mathcal I_\ell(C_i,C_j)}\Bigr]^3>0;\]
at $r=2$ the bracket is $\cosh(\mathcal I_\ell(C_1,C_2)/\mathsf b)$.
(c) For one loop the step is exact: the conditional mean of $W^{{\rm fr},(t)}_C$, the loop of the conditional mean of the field given the block field of the level $t$, given the block field of the level $t+1$ is $\lambda_{t+1}W^{{\rm fr},(t+1)}_C$; for a family the defect is at most $(e^{3x}-1)\lambda_{t+1}$, $x=\frac1{\mathsf b}\sum_{i<j}|a_K(X_{C_i},\mathcal C_{(t+1)}X_{C_j})|$, which for polygons at fixed distances falls faster than every power of $a_{K-t}$; so for families the factor without cross terms is the right one. The exponent $\kappa$, the perimeter constant $c_{\rm lat}|C|+\mathcal J_\ell(C)$ and the factor $\lambda_{t+1}$ of one step are those of three free photon fields.
(d) At every cutoff $K$, for every polygon $C$ with corners in $\Lambda_K$ and every bounded real $f$,
\[\frac{\mathbb E[W^{\rm fr}_C\mathcal O^{\rm fr}(f)]}{\mathbb E\,W^{\rm fr}_C}-\mathbb E\,\mathcal O^{\rm fr}(f)=-\frac3{2\mathsf b^2}\,a_K^{-1}\sum_pf(x(p))\,\mathcal I_{a_K}(C,\partial p)^2.\]
If $C$ has its corners in a fixed $\Lambda_i$ and $f$ is Lipschitz with support at positive distance from $C$, then as $K\to\infty$
\[
\begin{multlined}
-\frac3{2\mathsf b^2}\,a_K^{-1}\sum_pf(x(p))\,\mathcal I_{a_K}(C,\partial p)^2\ \relax \longrightarrow\ \relax -\frac3{2\mathsf b^2}\int_{T_\ell}f(x)\,|B_{C,\ell}(x)|^2\,dx,\qquad B_{C,\ell}(x) \\
={\rm curl}_x\oint_CG_\ell(x-y)\,dy,
\end{multlined}
\]
the field of $C$ on the torus; in the ball of radius $<\ell/4$ that holds the polygon, $|B_{C,\ell}(x)-B_C(x)|\le c\,|C|\,{\rm diam}(C)\,\ell^{-3}$, $c$ absolute, with $B_C(x)=\oint_C\frac{dy\times(x-y)}{4\pi|x-y|^3}$ the Biot--Savart field of $C$ in $\mathbb R^3$. For a bounded $f$ that is not Lipschitz the sums over the plaquette centres need not tend to the integral: an $f$ that vanishes at the plaquette centres of all cutoffs, a countable set, gives $0$ at every cutoff.
\end{proposition}

\begin{proof}[Sketch of proof]
(a) $F^{(d)}$ is a linear image of $F$; each component at a plaquette of $\Lambda_d$ is a centred Gaussian variable of variance at most $1/x_{[d]}$ (the smeared self-energy of the boundary of one plaquette of $\Lambda_d$, divided by $\beta$). $|F^{(d)}(p)|\ge2\varepsilon$ forces a component $\ge2\varepsilon/\sqrt3$, of probability at most $2e^{-\frac23\varepsilon^2x_{[d]}}$, and $\frac23\varepsilon_{[d]}^2x_{[d]}=\frac23A_0^2(\log x_{[d]})^{2\mathfrak p_0}$; a union bound over the $3(\ell L^d)^3$ plaquettes and the three components gives $18$, over the $12$ plaquettes at a site $72$. (b) $\mathbb E\prod_i\cos\phi^a_{C_i}=2^{-r}\sum_\epsilon\mathbb E\,e^{i\sum_i\epsilon_i\phi^a_{C_i}}=2^{-r}\sum_\epsilon\exp[-\frac1{2\mathsf b}\sum_{i,j}\epsilon_i\epsilon_j\mathcal I_{a_K}(C_i,C_j)]$; three independent colours; Proposition~\ref{C-372-11}(a) for the diagonal terms and for the distant pairs. (c) Given the block field of the level $t+1$, the curvature is its conditional mean, the curvature $\Pi_{t+1}F$ of the linear background, plus an independent Gaussian fluctuation; between the levels $t$ and $t+1$ the fluctuation of the flux through $C$ has the variance $\frac{a_K}{\mathsf b}(X_C,\mathcal C_{(t+1)}X_C)$ per colour, and $\mathbb E[\cos(\phi+\xi)\mid\phi]=e^{-\frac12{\rm Var}\,\xi}\cos\phi$; three colours give $\lambda_{t+1}$.

For a family the cross covariances $\frac{a_K}{\mathsf b}(X_{C_i},\mathcal C_{(t+1)}X_{C_j})$ enter as factors $e^{\mp\frac{a_K}{\mathsf b}(X_{C_i},\mathcal C_{(t+1)}X_{C_j})}$, at most $e^{x}$ for each colour. (d) For jointly Gaussian $\phi$, $F$: $\mathbb E[\cos\phi\,F^2]/\mathbb E\cos\phi-\mathbb E F^2=-{\rm Cov}(\phi,F)^2$; ${\rm Cov}(\phi^a_C,F^a(p))=\beta^{-1}(X_C,X_{\partial p})=\mathcal I_{a_K}(C,\partial p)/\mathsf b$; $\frac12$ and three colours give $\frac32$. The identity holds for every bounded $f$, as nothing in it is a limit. For the limit let the corners of $C$ lie in a fixed $\Lambda_i$: away from $C$, uniformly on the support of $f$, $\mathcal I_{a_K}(C,\partial p)=a_K^2(B_{C,\ell}(x(p))\cdot n_p+O(a_K{\rm Lg}(\ell/a_K)))$, the mutual energy of $C$ with a small loop being the flux through it of the field of $C$ formed with the Green function of the torus: Proposition~\ref{C-372-11}(a) for the two polygons $C$ and $\partial p$ at positive distance compares the lattice energy with $\mathcal I_\ell(C,\partial p)=\oint_{\partial p}\oint_CG_\ell(x-y)\,dy\cdot dx$ up to $O(a_K^3{\rm Lg}(\ell/a_K))$, and Stokes' formula makes $\mathcal I_\ell(C,\partial p)$ the flux of $B_{C,\ell}$ through $p$; summing $a_K^3f(x(p))(B_{C,\ell}(x(p))\cdot n_p)^2$ over the three orientations of the plaquettes gives a Riemann sum of $f|B_{C,\ell}|^2$, which converges for Lipschitz $f$.

Finally $G_\ell(x)=\frac1{4\pi|x|}+H_\ell(x)$ with $H_\ell(x)=\ell^{-1}H_1(x/\ell)$ smooth on $|x|<\ell$; as $\oint_Cdy=0$, 
\[
{\rm curl}\oint_CH_\ell(x-y)\,dy=\oint_C[\nabla H_\ell(x-y)-\nabla H_\ell(x-y_0)]\times dy,
\]
of modulus at most 
\[
|C|\,{\rm diam}(C)\sup_{|u|\le\ell/2}|\nabla^2H_\ell(u)|=|C|\,{\rm diam}(C)\,\ell^{-3}\sup_{|u|\le1/2}|\nabla^2H_1(u)|.
\]

\end{proof}

The proof of Theorem~\ref{C-372-6} uses the following tools; we state them in lattice units, with spacing $1$ and $\beta=\mathsf b/a$.

\begin{lemma}[Tools of the tube]\label{C-372-17}
Let $\nu$ be the Wilson measure at $\beta$ on a torus $(\mathbb Z/N\mathbb Z)^3$, $N$ even, or a periodic Gibbs state as in Theorem~\ref{C-372-6}. For a step $c=\langle x_-,x_+\rangle$ of a closed path let $c(y)$, $y\in\mathbb Z^2$, be its translates in the lattice plane through $x_-$ orthogonal to $c$; for an integer $\rho\ge2$ let $T_\rho(c)=\{c(y):|y|<\rho\}$ be the slab and $K_\rho(c)=\{x_-+y+se:|y|_\infty\le\rho,\ 0\le s\le1\}$ the box of $c$, $e$ its direction; let $\phi_\rho(y)=\min\{1,\log^+(\rho/|y|)/\log\rho\}$, $\phi_\rho(0)=1$, so that $\phi_\rho=0$ at $|y|\ge\rho$, $E_\rho=(2\pi\log\rho+7.63)/\log^2\rho$, and $c_j=j(j+1)$.
(a) (Admissible steps.) A set $I$ of steps of closed paths $C_1,\dots,C_r$ is admissible at $\rho$ if no other step of the family runs through a bond of $T_\rho(c_i)$, $i\in I$, and no bond of $T_\rho(c_{i'})$, $i'\ne i$, lies in $K_\rho(c_i)$. For closed paths without self-intersection with $s$ straight pieces in all, each of at least $4\rho+1$ steps, whose pieces without a common corner are at sup-distance at least $2\rho+1$, the steps $2\rho\le j\le l'-1-2\rho$ of each piece of $l'$ steps form an admissible set of $n-4\rho s$ steps, $n$ the number of steps; for a rectangle, $n-16\rho$.
(b) (The profile.) $\sum_{\langle x,y\rangle}(\phi_\rho(x)-\phi_\rho(y))^2\le E_\rho$, the sum over the bonds of $\mathbb Z^2$, and $(\log\rho-1.2144)/2\pi\le1/E_\rho\le\log\rho/2\pi$.
(c) (An information inequality, three colours at once.) Let $\varrho>0$ be a $C^2$ density on $\mathrm{SU}(2)^Y$ with $0\in Y$, $\phi:Y\to[0,\infty)$ with $\phi_0=1$, frames $P_y\in\mathrm{SU}(2)$ with $P_0=\mathbf 1$, and for a unit vector $\mathbf n=(\mathbf n_1,\mathbf n_2,\mathbf n_3)\in\mathbb R^3$ the moves $v_y\mapsto\exp[is\phi_yP_y^{-1}(\frac12\sum_a\mathbf n_a\sigma_a)P_y]v_y$, $\sigma_1,\sigma_2,\sigma_3$ the Pauli matrices, with the derivative $D_{\mathbf n}$ at $s=0$; let $\bar I=\sup_{|\mathbf n|=1}\mathbb E[-D_{\mathbf n}^2\log\varrho]$. Then $\|\mathbb E\,D^j(v_0)\|^2\le\bar I/(\bar I+c_j)$ in the operator norm for every spin $j>0$, where $D^j$ is the representation of spin $j$. The spin-zero representation has norm $1$.
(d) (The second variation.) For these moves on the bonds $c(y)$ of a slab, with any frames, a bond $e'=\langle x,y\rangle$ of the plane, $p(e')$ the plaquette spanned by $e'$ and $c$, and $H_{e'}=P_yU(e')^{-1}P_x^{-1}$: $D_{\mathbf n}^2a_{p(e')}\le\frac14(\phi_x-\phi_y)^2+\phi_x\phi_y(3a(H_{e'})+a_{p(e')})$, $a(g)=1-\frac12\operatorname{tr}g$.
(e) (Ceilings multiply over cells.) Let $B=\bigsqcup_lB_l$ be a set of bonds made of cells $B_l$, no plaquette having bonds in two cells, and $C=\gamma_1\delta_1\cdots\gamma_s\delta_s$ a closed path with $\gamma_l$ through no bond of $B\setminus B_l$ and $\delta_l$ through no bond of $B$. If $Q_l$ is the conditional mean under $\nu$ of $D^j(U(\gamma_l))$ given the bonds off $B_l$, then $\nu(\chi_j(U(C))\,G)=\nu(\operatorname{tr}[\prod_lQ_lD^j(U(\delta_l))]\,G)$ for every bounded Borel $G$ of the bonds off $B$, and $|\nu(\chi_j(U(C))\,G)|\le(2j+1)\sup|G|\,\nu(\prod_l\|Q_l\|)$.
(f) (One plaquette.) With $\mathfrak M_\nu=\sup_p\nu(e^{\beta a_p})$: $\nu(a_p)\le\beta^{-1}\log\mathfrak M_\nu$ and $\nu(a_p\ge s)\le\mathfrak M_\nu e^{-\beta s}$; for $\beta\ge1$, $\log\mathfrak M_\nu\le192\frac{\mathsf b}\ell+4(14.15+3\log\beta)$ on the torus of side $\ell$, and $\log\mathfrak M_\nu\le28.3+6\log\beta$ for the Gibbs states.
(g) (The ceiling of a tube.) Let $\beta>0$, $\rho\ge2$, $C_l$ closed paths with spins $j_l$, $I$ admissible at $\rho$, $c_i=c_{j_l}$ for a step of $C_l$, $G$ a bounded Borel function of the bonds off $B=\bigcup_{i\in I}T_\rho(c_i)$, $\mathcal G=\sum_{i\in I}\frac{2c_i}{\beta E_\rho}(1-\frac{2c_i}{\beta E_\rho})$ and $\omega_\rho=40\rho^4\sum_{i\in I}c_i/(\beta E_\rho)^2$. If $\omega_\rho\le1$, then $|\nu(G\prod_l\chi_{j_l}(U(C_l)))|\le\sup|G|\prod_l(2j_l+1)\,e^{-\mathcal G}\mathfrak M_\nu^{\omega_\rho}$; and $2c_j/E_\rho\ge\frac43c_jk(\log\rho-1.2144)$: the number per scale is $k$.
(h) (The ceiling with dyadic profiles.) Let $\rho=2^S\ge2$ and $C$ a coordinate rectangle of $n$ steps with sides of at least $4\rho+1$ steps and, on the torus, with the points within $\rho+1$ of $C$ in a cube of side at most $\frac N2$; let $W_C=\frac12\operatorname{tr}U(C)$, $p=kn/\beta<3$, $\delta_0=(16k\rho\log\rho+pk\log^2\rho)/\beta$ and $\bar a=\nu(a_p)$, the plaquette mean. Then, with the profiles $\phi_r$ of the scales $r=2,4,\dots,2^S$,
\[|\nu(W_C)|\le e^{\delta_0+2.2144p}\rho^{-p}\Bigl[1+\frac{72\,\bar a\rho^3\log^2\rho}{2^{3-p}-1}\Bigr]+35\,\bar a;\]
the plaquette mean enters linearly, and $p<3$ keeps the denominator positive.
\end{lemma}

\begin{proof}[Sketch of proof]
(a) A step at distance $\ge2\rho$ from the ends of its piece has a slab that meets no other step, and the slabs of two steps of pieces without a common corner at sup-distance $\ge2\rho+1$ do not reach each other's boxes. (b) The sum is a lattice form of $\int|\nabla\phi_\rho|^2=2\pi\int_1^\rho(r\log\rho)^{-2}r\,dr=2\pi/\log\rho$, with explicit corrections at the centre and at the rim. (c) The moves are one-parameter groups of maps of $\mathrm{SU}(2)^Y$ that keep the Haar measure; under the move of $v_0$, $D^j(v_0)$ is multiplied by $e^{isJ_{\mathbf n}}$, $J_{\mathbf n}$ the generator of spin $j$, the image of $\frac12\sum_a\mathbf n_a\sigma_a$, and $\sum_{\mathbf n}J_{\mathbf n}^2=c_j$ over an orthonormal basis of directions.

With $M=\mathbb E\,D^j(v_0)$ and a unit vector $u$, integration by parts gives $iJ_{\mathbf n}Mu=-\mathbb E[(D^j(v_0)-M)u\,D_{\mathbf n}\log\varrho]$. Put $s_a=D_{e_a}\log\varrho$ and $s=(s_1,s_2,s_3)$. Since the supremum defining $\bar I$ runs over every unit direction, integration by parts gives $\mathbb E(ss^{\sf T})\le\bar I\mathbf 1$. The vector form of Schwarz's inequality therefore gives $\sum_a|\mathbb E[Xs_a]|^2\le\bar I\mathbb E|X|^2$ for $X=(D^j(v_0)-M)u$. Using $\sum_a|J_{e_a}Mu|^2=c_j|Mu|^2$ and $\mathbb E|X|^2=1-|Mu|^2$ yields $c_j|Mu|^2\le\bar I(1-|Mu|^2)$. (d) Expansion of $a_{p(e')}$ in $s$: the abelian part gives $\frac14(\phi_x-\phi_y)^2$, and the commutator terms are bounded by the defect $a(H_{e'})$ of the frames and by $a_{p(e')}$. (e) The cells interact with the rest only through plaquettes with bonds in one cell and off $B$, so conditioning cell by cell replaces $D^j(U(\gamma_l))$ by $Q_l$, the $\delta_l$ and $G$ being functions of the bonds off $B$; $|\operatorname{tr}\prod A_i|\le(2j+1)\prod\|A_i\|$. (f) Jensen's and Chebyshev's inequalities; the chessboard estimate bounds $\nu(e^{\beta a_p})$ by a power of a ratio of partition functions, and $\int\exp[-\beta\sum_{p\in Q}a_p]\,dU\ge(\mathsf K\beta^3)^{-|Q|}$, $\mathsf K=1.39\cdot10^6$, for a box $Q$ of sites ($\log\mathsf K<14.15$); the torus adds $192\mathsf b/\ell$. (g) In each slab of $I$ we use comb frames: in the plane through $x_-$ orthogonal to $c_i$ a coordinate axis through $x_-$ is the base, the lines of the plane orthogonal to it are the teeth, and $P_y$ is the transport from $x_-$ along the base and then along the tooth to $x_-+y$, a product over bonds of the plane, which lie off $B$.

Then $H_{e'}=\mathbf 1$ on the base and on the teeth, and at a rung $e'$, a bond parallel to the base at the height $t_{e'}\ne0$, $H_{e'}$ is conjugate to the holonomy of the boundary of the strip $S_{e'}$ of the $|t_{e'}|$ plaquettes between the rung and the base, so that $a(H_{e'})\le|t_{e'}|\sum_{q\in S_{e'}}a_q$ (Proposition~\ref{C-372-13}(c) and Schwarz' inequality). Write $\nu_B(F)$ for the conditional mean of $F$ under $\nu$ given the bonds off $B$. As no plaquette has bonds in two cells, given the bonds off $B$ the cells are independent, and the conditional law of the bonds of the slab of $c_i$ given the bonds off it is their conditional law given the bonds off $B$. The bonds of the slab enter the action only through the plaquettes $p(e')$, $e'$ a bond of the plane, so that $-D_{\mathbf n}^2$ of the logarithm of this conditional density is $\beta\sum_{e'}D_{\mathbf n}^2a_{p(e')}$, with $\phi$ taken as $0$ off the slab, whose bonds alone move; by (b), (d) its conditional mean is $\bar I\le\frac\beta4(E_\rho+X_i)$, with
\[X_i=4\sum_{e'}\phi_x\phi_y\bigl[3a(H_{e'})+\nu_B(a_{p(e')})\bigr],\]
a function of the bonds off $B$. It is at most $\sum_qw^{(i)}_q\,\nu_B(a_q)$, a sum over plaquettes with weights $w^{(i)}_q\ge0$: $4\phi_x\phi_y$ at the plaquette $p(e')$ of every bond $e'$ of the disc, and $12|t_{e'}|\phi_x\phi_y$ at each plaquette of the strip of every rung $e'$, where $\nu_B(a_q)=a_q$, the bonds of the strip lying in the box $K_\rho(c_i)$ and so off $B$. The total weight $\sum_qw^{(i)}_q$ is at most $20\rho^4$: the strips give about $12\int_{|y|<\rho}y_2^2\,dy=3\pi\rho^4$, the plaquettes $p(e')$ a multiple of $\rho^2$. (c) gives $\|Q_i\|^2\le(1+4c_i/(\beta(E_\rho+X_i)))^{-1}$, so, by $-\frac12\log(1+z)\le-\frac z2+\frac{z^2}4$ and $\frac1{E+X}\ge\frac1E-\frac X{E^2}$, $\log\|Q_i\|\le-\frac{2c_i}{\beta E_\rho}(1-\frac{2c_i}{\beta E_\rho})+\frac{2c_iX_i}{\beta E_\rho^2}$; (e) multiplies these over the disjoint cells.

Put $w_q=\sum_i2c_iw^{(i)}_q/(\beta E_\rho^2)$. Jensen's inequality for the conditional mean $\nu_B$ gives $\exp\sum_i\frac{2c_iX_i}{\beta E_\rho^2}\le\nu_B\bigl(\prod_qe^{w_qa_q}\bigr)$; as $\sum_qw_q/\beta\le40\rho^4\sum_ic_i/(\beta E_\rho)^2=\omega_\rho\le1$, H\"older's inequality with the exponents $\beta/w_q$ and (f), with $\mathfrak M_\nu\ge1$, give $\nu(\prod_qe^{w_qa_q})\le\prod_q\nu(e^{\beta a_q})^{w_q/\beta}\le\mathfrak M_\nu^{\omega_\rho}$. The last inequality is (b) with $\frac43k=\frac1\pi$. (h) Keep the cells $T_\rho(c_i)$ of the $n-16\rho$ steps of the admissible set of (a) (spin $\frac12$, $c_{1/2}=\frac34$). On each cell and for every dyadic $2\le r\le\rho$ the moves with the profile $\phi_r$, which vanishes at $|y|\ge r$, rotate bonds of the cell alone, and (b)--(e), as in the proof of (g), give at every field $\prod_i\|Q_i\|\le\min_re^{-\mathcal G_r+Y_r}$, with $\mathcal G_r$ the number $\mathcal G$ of (g) formed with $E_r$, $Y_r=\sum_i2c_{1/2}X^{(r)}_i/(\beta E_r^2)$ and $X^{(r)}_i$ the $X_i$ of the proof of (g) formed with $\phi_r$; by (b), $e^{-\mathcal G_r}\le e^{\delta_0+1.2144p}r^{-p}$, $\delta_0$ paying for the $16\rho$ omitted steps and for the second-order term of $\mathcal G_r$, so that $e^{-\mathcal G_r+Y_r}\le e^{\delta_0+2.2144p}r^{-p}$ where $Y_r\le p$.

The sharper estimate of the non-abelian member of the second variation bounds the mean of $Y_r$ in the place of an exponential moment: in comb frames $a(H_{e'})$ at a rung $e'$ of height $t_{e'}$ is $1-W$ of the boundary of the strip $S_{e'}$ of $1\times|t_{e'}|$ plaquettes, whose mean is $1-w(1,|t_{e'}|)\le1-w(1,1)^{|t_{e'}|}\le|t_{e'}|\,\bar a$ by Proposition~\ref{C-372-13}(c) (for the Gibbs states the strips of even height by the reflections in the planes of sites, those of odd height by $\sqrt{a(gh)}\le\sqrt{a(g)}+\sqrt{a(h)}$), so that $\nu(X^{(r)}_i)\le c\,r^3\bar a$ and $\nu(Y_r)\le c\,p\,r^3\log^2r\,\bar a$: the plaquette mean enters linearly, and this bound of a mean takes the place of the condition $\omega_\rho\le1$ and of the weight $40\rho^4$ of (g).

Let $r_*$ be the largest dyadic $r\le\rho$ with $Y_r\le p$; for $r<\rho$ the event $\{r_*=r\}$ lies in $\{Y_{2r}>p\}$, of probability at most $\nu(Y_{2r})/p$ by Markov's inequality, and on the event that $Y_r>p$ at every $r$ the bound $1$ serves. The sum over the dyadic $r<\rho$ of $r^{-p}(2r)^3\log^2(2r)$ is geometric and led by its last term exactly when $p<3$, which gives the display.
\end{proof}

\subsection{Proof of Theorem~\ref{C-372-3}}

\emph{Step 1: (1).} For $D\ge d_\Theta\ge d_\natural$ (as $\Theta_\star\ge t_\natural$), $\{\tau\le D\}=\{\tau_0\le D\}=\bigcap_{d\ge D}\pi_d^{-1}(I_d)$: if $\tau_0\le D$, then $V_d\in I_d$ for all $d\ge\tau_0$, so for all $d\ge D$; if $V_d\in I_d$ for all $d\ge D$, then $D$ is one of the integers of which $\tau_0$ is the least. Each member is open ($I_d$ is open, $\pi_d$ continuous), $\{\tau\le D\}=\emptyset$ for $D<d_\Theta$, and $\{\tau=\infty\}=\{\tau_0=\infty\}$; a function into a countable set whose sublevel sets are Borel is Borel. As $\pi_d(\mu)=\mu_d$, Proposition~\ref{C-372-8}(b) and countable subadditivity give the tail by the sum, and the same clause of Proposition~\ref{C-372-8}(b) the bound $\frac L{L-1}R_D$ when $c'A_0^2\log x_{[D]}\ge2$, as then $R_{d+1}\le L^{-1}R_d$ for all $d\ge D$; $\sum_dR_d<\infty$ and $\mu(\mathsf{Thr})=1$ give $\mu(\Omega)=1$, and for this the bound $\mu_d(\mathcal V_d\setminus I_d)\le R_d$ with $R_d\to0$ faster than every power of $L^{-d}$ suffices.

If $c'A_0^2(\log x_{[D]})^{2\mathfrak p_0}\le26\ell^3$, the exponent in $R_D$ is at least $(\mathsf C^\Sigma_{\rm maj}+2+c_{\rm low}-26)\ell^3\ge0$ by $c_{\rm low}\ge24$, and $2\ell^3L^{3D}M^{-3}\ge16$ as $\ell\ge2M$; so $R_D>1$, and the two bounds, $\sum_{d\ge D}R_d$ and $\frac L{L-1}R_D$, both at least $R_D$, exceed $1$. Next, $\{\tau>D\}=\bigcup_{d\ge D}\pi_d^{-1}(\mathcal V_d\setminus I_d)$, and $\mathcal V_d\setminus I_d$ is the event of Proposition~\ref{C-372-8}(c) for the set of all $\ell^3L^{3d}$ sites, which gives the bound by $\sum_{d\ge D}R^\natural_d$. As $(y+\log L)^{2\mathfrak p_0}-y^{2\mathfrak p_0}\ge2y\log L$ for $y\ge1$, $R^\natural_{d+1}/R^\natural_d\le L^{3-2c^\natural A_0^2\log x_{[d]}}$, which is at most $L^{-1}$ for $d\ge D$ if $c^\natural A_0^2\log x_{[D]}\ge2$.

With $y=D\log L\le\log x_{[D]}$, $\log R^\natural_D\le3\log\ell+3y-c^\natural A_0^2y^2$; the first member of the condition gives $3y\le\frac12c^\natural A_0^2y^2$ and $c^\natural A_0^2\log x_{[D]}\ge6$, the second $\frac12c^\natural A_0^2y^2\ge3\log\ell+\log\frac1\eta$; so $R^\natural_D\le\eta$.

\emph{Step 2: (2).} Theorem~\YMMref{E}{10} in the setting of Section~\YMMref{52.7}{432} and Section~\ref{sec:analytic-core-complements} give the orbit for $\varpi(W)<\bar\varepsilon$; the inclusion $\{\varpi<2c_\rho\varepsilon_{[d]}\}\subset\{\varpi<\bar\varepsilon\}$ for $x_{[d]}\ge t_\natural$ is the condition (cut1) at $t=x_{[d]}$, which gives $2c_\rho\varepsilon(t)\le\varepsilon_{\mathcal W}\alpha_0$, with $\varepsilon_{\mathcal W}\alpha_0\le\bar\varepsilon$ (equation~\YMMref{(2134)}{790}). For $d\ge\tau_0\ge d_\natural$, $x_{[d]}\ge t_\natural$ and $V_d\in I_d\subset\{\varpi<2c_\rho\varepsilon_{[d]}\}$. If $d\ge\tau(V_\bullet)$, then $d\ge d_\Theta$ is admitted and $d\ge\tau_0$, so $V_d\in I_d\subset J_d$. Each clause of Proposition~\ref{C-372-9} asks of $W$ only that it lie in $J_d$ at an admitted depth, not that it be a block average or a member of a thread; we substitute $W:=V_d$, and Proposition~\ref{C-372-9}(a)--(c) give the limits, the rate, the curves, the continuity, the gauge invariance and the bound for rectangles. The thread property enters no step, and $\Omega$ is defined without $d$, $C$, $\gamma$, $K$.

\emph{Step 3: (3).} $J_d$ is open and $W^{[d]}_C$ is continuous and gauge invariant on it; the inverse image of a Borel set under $F^{[d]}_C$ is a Borel subset of $J_d$, joined with the closed set $\mathcal V_d\setminus J_d$ if it contains $0$; $\pi_d$ is continuous; so $Y_{d,C}$ is Borel and $\sigma(\pi_d)$-measurable, $\sigma(\pi_d)\subset\mathcal F_d$, and its restriction to the Borel set $\mathsf{Thr}$ is a Borel function on the threads. $\sup_{\mathcal X}|Y^K_{d,C}-Y_{d,C}|\le\sup_{J_d}|w^K_C-W^{[d]}_C|\to0$ by Proposition~\ref{C-372-9}(b). $\theta_v$ is a bijection, with inverse $\theta_{v^{-1}}$, that maps every $\pi_j^{-1}(I_j)$ onto itself, and $Y_{d,C}(\theta_vV_\bullet)=F^{[d]}_C(V_d^{\,v|\Lambda_d})=Y_{d,C}(V_\bullet)$. $(C^{\langle d\rangle})^{\langle j\rangle}=C^{\langle j\rangle}$, so $W_{C^{\langle d\rangle}}=W_C$ on $\mathcal V_j$, and every $d\ge d_\Theta$ is admitted. For curves Proposition~\ref{C-372-9}(c) replaces (a), (b); the shadows are fixed lattice by lattice because $w^K_\gamma$ depends on the shadow on the bare lattice, and a curve has no refinement.

\emph{Step 4: (4).} $\pi_d(\mu)=\mu_d$ gives the first equality, Proposition~\ref{C-372-9}(b), (d) the others and the rates. $|Y_{d,C}|\le1$ and Step~1 give $|\int Y_{d,C}\,d\mu-\int\mathbf 1\{\tau\le d\}Y_{d,C}\,d\mu|\le\mu(\tau>d)\le\sum_{d'\ge d}R_{d'}$, at most $\frac L{L-1}R_d$ if $c'A_0^2\log x_{[d]}\ge2$ (Step~1). $G_d-F^{[d]}_C$ vanishes on $J_d$ and has modulus at most $1$ off it, and $\mu_d(\mathcal V_d\setminus J_d)\le\mu_d(\mathcal V_d\setminus I_d)\le R_d$; for $d\ge\tau$, $V_d\in J_d$ and so $G_d(V_d)=Y_{d,C}$. For the joint means, with $D=\max_id_i$: the law of $(V_0,\dots,V_D)$ is the image of $\mu_D$ under $V\mapsto(\mathcal M^DV,\dots,V)$, $\prod_iF^{[d_i]}_{C_i}\circ\mathcal M^{D-d_i}$ is one bounded Borel function on $\mathcal V_D$, $\mathcal M^{K-d_i}=\mathcal M^{D-d_i}\mathcal M^{K-D}$, and $\mu^K_D\to\mu_D$ in total variation; no measure on an infinite product is needed here. The means depend on the single laws $\mu_d$; what $\mu$ adds to them are the events that involve infinitely many depths, $\{\tau\le D\}$ and $\Omega$.

\emph{Step 5: (5).} We use the localisation of backgrounds (Proposition~\YMMref{53.5}{447}) at $n=K-d$, for the constant sequence $\Omega_0=\dots=\Omega_{K-d}=T_\ell$, whose datum is $V$ on the bonds of $\Lambda_d$ and whose minimal orbit is the set of the backgrounds over $V$; for the region $\Omega'$, a region for the partition into the $M_1$-cubes of the depth $d$, with the one part $\mathsf Y_1=T_\ell\setminus\Omega'$ of its complement and the index $l_1=K-d$; at $\varepsilon_{\rm r}=2\varepsilon_{K-d}\le\epsilon_{\rm lc}$; and at $V\in\mathsf O_{K-d}$. We verify its hypotheses for this region of $M_1$-cubes, using the joined sequence defined in equation~\YMMref{(999)}{446}: the joined sequence is the $\mathfrak r_2$-tower of $\Omega'$, whose minimal orbit over the datum induced by $V$ defines the localised background $U'$ (equation~\YMMref{(77)}{42}); it is separated; its tower bonds are the bonds of its members of index $<K-d$, with both end points in the collar $\mathsf Y_1^{\sim2}\setminus\mathsf Y_1$, which lies in the wide collar $\mathsf Y_1^{\sim3}\setminus\mathsf Y_1$ (Section~\YMMref{53.4}{445}), the points of $\Omega'$ within three layers of $M_1$-cubes of the depth $d$ of $T_\ell\setminus\Omega'$; the core $\Omega'^{\sim-2}$ lies in its last member, so that $j(b)=K-d$ at the bonds of the core; and both data are regular at the radius $2\varepsilon_{K-d}$, as $V\in\mathsf O_{K-d}$ (Proposition~\YMMref{44.2}{345}).

So, after gauge transformations, the background over $V$ and $U'$ differ by factors $e^{i\mathsf a(b)}$ at the bare bonds $b$ with no end point in the wide collar, with $|\mathsf a(b)|\le\mathsf C_{\rm lc}\,2\varepsilon_{K-d}L^{d-K}e^{-\delta_{\rm lc}D(b)/2}$, $D(b)$ the distance of $b$ from the tower bonds in the weighted distance of the tower (Section~\YMMref{53.4}{445}). The bare bonds of $C$ are such bonds of the core, as ${\rm dist}(C,T_\ell\setminus\Omega')\ge\frac{\delta_1}2\ge4M_1a_d$. The tower-bond endpoints lie within the two-layer collar: a row below the last level has both endpoints outside the next tower member, which contains the core $\Omega'^{\sim-2}$. Moreover every local spacing in the joined sequence is at most $a_d$. For any connected path $\Gamma$ joining two points, the layer pieces cover $\Gamma$, so its weighted length is at least $a_d^{-1}\operatorname{length}(\Gamma)$, and hence at least $a_d^{-1}$ times the periodic sup-distance between those points. Taking the infimum over paths proves the same bound for weighted distance; the sup-distance to the complement of a slab, or of the intersection of two slabs, is the distance; so on $C$, $D(b)\ge{\rm dist}(C,T_\ell\setminus\Omega')/a_d-2M_1\ge\frac{\delta_1}2L^d-2M_1\ge\frac{\delta_1}4L^d$, as $M_1a_d\le\frac{\delta_1}8$.

The loop is gauge invariant, $W_C=\frac12\operatorname{tr}$, and $|\operatorname{tr}\prod_iA_i-\operatorname{tr}\prod_iB_i|\le2\sum_i|A_i-B_i|$ for unitary matrices; summing $|e^{i\mathsf a(b)}-\mathbf 1|\le|\mathsf a(b)|$ over the $|C|/a_K$ bare bonds of $C$ gives $|w^K_C(V)-Y_{K,d}(V)|\le\frac{|C|}{a_K}\,2\mathsf C_{\rm lc}\varepsilon_{K-d}L^{d-K}e^{-\delta_{\rm lc}\delta_1L^d/8}=2\mathsf C_{\rm lc}\varepsilon_{K-d}|C|L^de^{-\delta_{\rm lc}\delta_1L^d/8}$, and $\varepsilon_{K-d}\le\hat\varepsilon_d$ because $\frac23x_{[d]}\le x_{K-d}\le x_{[d]}+\bar B$. Two fields that agree on the bonds of $\Omega'$ have the same $Y_{K,d}$; let $K\to\infty$ with Proposition~\ref{C-372-9}(b), (e).

\emph{Step 6: (6).} The Wilson measure at $(K,\ell,\mathsf b)$ is the one at $(K-s,L^s\ell,L^s\mathsf b)$ carried by $x\mapsto L^sx$, since $\beta=L^K\mathsf b=L^{K-s}(L^s\mathsf b)$, and $\Lambda_d$ goes to $\Lambda_{d-s}$; so the limit laws, the measure on the threads from the depth $s$ on and Proposition~\ref{C-372-8} hold at every $\mathsf b>0$ with $R_{d-s}[L^s\ell,L^s\mathsf b]$. At $(L^s\ell,L^s\mathsf b)$ the one-torus setting holds (Lemma~\ref{C-372-7}), and (1)--(4) hold there at the depth $d-s$, with the depths $d_\natural$, $d_\Theta$, $s_\natural$ of $(L^s\ell,L^s\mathsf b)$; increased by $s$, these depths take the places of $d_\natural$, $d_\Theta$, $s_\natural$ in (1)--(4) at $(\ell,\mathsf b)$, the depth $d-s$ of $(L^s\ell,L^s\mathsf b)$ being the depth $d$ of $(\ell,\mathsf b)$.

The dilation maps $\Lambda_j$ of $T_\ell$ onto $\Lambda_{j-s}$ of $T_{L^s\ell}$, one graph with the same $\ell L^j=2L^{m+j}$ sites in each direction, bond onto bond, block onto block and plaquette onto plaquette, and the cutoff $K$ onto the cutoff $K-s$, with the same number $K-d$ of block steps above the depth $d$; it leaves $\beta$, $x_{[d]}=L^d\mathsf b=L^{d-s}(L^s\mathsf b)$, $\varepsilon_{[d]}$, the balls $I_d$, $J_d$ and the length in units of the lattice, $|C|L^d$, unchanged. Every object in the definition of the smoothed loop is formed from these graphs, from $K-d$ and from the constants, in the units of the lattices, and holds neither the side, nor the label, nor the spacing. Item by item: the block averaging map with its averaging radius $1/200L^2$ (equations~\YMMref{(37)--(39)}{20}); the constraint sets $\{\mathcal M^{K-d}U=W\}$, the action $\sum_pa_p$, criticality for it on them and the regularity at the radius $\varepsilon_{\rm v}$, a constant, all measured in the units of the level, with the bare spacing $L^{-(K-d)}$ (Definitions~\YMMref{9.3}{40} and~\YMMref{9.5}{41}); the residual gauge groups $\mathcal G^d_K$ and the minimal orbit (Section~\ref{sec:analytic-core-complements}); the axial gauge of the averaged backgrounds in the definition of $H^K_s$ (Section~\ref{sec:sf-chart-interface}); the radius $\bar\varepsilon$ and the lattices of Theorem~\YMMref{E}{10} in the setting of Section~\YMMref{52.7}{432}, whose lattice hypotheses are imposed through the number of the sites in each direction at the levels used, so that the backgrounds over a field with $\varpi<\bar\varepsilon$ form one non-empty orbit at $(\ell,\mathsf b)$ and at $(L^s\ell,L^s\mathsf b)$ alike; and the plaquette size $\varpi$, a maximum over the plaquettes of one graph .

So the transported backgrounds over $V_d$ are the backgrounds over $V_d$ on the lattices of side $\ell$, the transported $\mathcal C_K(W)$, $H^K_d$, $H^\infty_{d,n}$, $w^K_C$, smoothed loops, balls and variables are those of $(\ell,\mathsf b)$ at the depth $d$, and (1)--(4) at $(L^s\ell,L^s\mathsf b)$ and the depth $d-s$ are (1)--(4) at $(\ell,\mathsf b)$ and the depth $d$.

\subsection{Proof of Theorem~\ref{C-372-4}}

\emph{Step 1: the ceiling.} Fix $\vartheta\in(\kappa,\frac12)$. Proposition~\ref{C-372-10}(a) gives the hypothesis of Proposition~\ref{C-372-12}(a) with $\eta_s=C_\eta L^{-\vartheta(K-s+1)}$ up to the last level $\bar t=K-d$, $d\ge s_0$. By Proposition~\ref{C-372-11}(c), $a_K^{-\kappa}\Xi_{s-1}\le c_\Xi L^{\kappa(K-s+1)}$, the depth $K-s+1\ge d$ being at least $i_0$; so
\[a_K^{-\kappa}\mathsf e_t\le c_\Xi C_\eta\sum_{j\ge K-t+1}L^{-(\vartheta-\kappa)j}=\epsilon_1L^{-(\vartheta-\kappa)(K-t+1)}:\]
the additive errors of the steps, weighed by $\Xi_t\approx L^{-\kappa t}$, sum after division by $a_K^\kappa$ exactly when $\vartheta>\kappa$; what ends at $\kappa=\vartheta$ is the additive error, not the tower. The rarity bound at $p=\kappa+1$ gives $\sum_{t\le\bar t}L^{\kappa(K-t)}\mathsf w_t\le2\mathsf V_\ell C_{\rm r}(\kappa+1)L^{-d}$, and $a_K^{-\kappa}(\Xi_{t-1}+\mathsf e_{t-1})\le(c_\Xi L^\kappa+\epsilon_1)L^{\kappa(K-t)}$; Proposition~\ref{C-372-12}(a) gives \eqref{C-372-5}. At $d=s_0$, $|W^{(\bar t)}_{\mathsf C}|\le1$ and $Z^{-1}m_{\bar t}(\mathcal V_{K-\bar t})\le1$ give the supremum, and $|\Phi^{(t)}[W_{\mathsf C}]|\le\Xi_t+\mathsf e_t\le(c_\Xi+\epsilon_1)L^{-\kappa t}$. A rectangle in a ball of radius $<\ell/4$ has $|C|<\sqrt2\,\ell$, which is $<\frac{2\pi}3\mathsf b$ if $\ell\le1.48\,\mathsf b$. If $\kappa\ge\frac12>\vartheta$, the same tower with $\Xi_{s-1}\le c\,L^{-\kappa(s-1)}$ gives $\mathsf e_{\bar t}\le c\,a_K^\vartheta$ and $\Xi_{\bar t}\le c\,a_K^\kappa$, and the rarity terms are smaller.

\emph{Step 2: the limit.} Let $d\ge\max\{s_0,d_0\}$. By Proposition~\ref{C-372-9}(e), $m_{K-d}$ is carried by $\mathsf O_{K-d}\subset J'_d$, where $W^{(K-d)}_{\mathsf C}=w^K_{\mathsf C}\to F':=\prod_iW^{[d]}_{C_i}$ uniformly. By Proposition~\ref{C-372-12}(c), $Z^{-1}m_{K-d}$ and $\mu^K_d$ differ on gauge invariant sets by a positive measure of mass $\sum_{t\le K-d}\mathsf w_t\le2\mathsf V_\ell C_{\rm r}(\kappa+1)L^{-(\kappa+1)d}$; $\mu^K_d\to\mu_d$ in total variation; $a_K^{-\kappa}\Xi_{K-d}\to\mathsf l_d\le c_\Xi L^{\kappa d}$; and $\mathbf 1_{J'_d}F'$ differs from $F^{[d]}_{\mathsf C}$ only off $J_d$, a set of $\mu_d$-measure at most $R_d$, which is also the difference of the two means. With \eqref{C-372-5}, $\limsup_K\mathcal W_K-\liminf_K\mathcal W_K\le2\mathsf E_d\to0$, so the limit exists, and $|\mathcal W(\mathsf C)-\mathsf l_d\int F^{[d]}_{\mathsf C}\,d\mu_d|\le\mathsf E_d$.

The value of $\mathsf l_d$ (Proposition~\ref{C-372-11}(c)) gives the asymptotics of $a_d^{-\kappa}\int F^{[d]}_{\mathsf C}\,d\mu_d$, its remainder giving $c\,a_d(1+\log^+(1/a_d))/\mathsf b$ for a fixed polygon, which tends to zero. Uniformity on $\mathsf J$: the numbers of Proposition~\ref{C-372-10} depend on $\mathsf J$ alone; $c_\Xi$ is uniform on $\mathsf J$; $a_K^{-\kappa}\prod_n\lambda_n\to\mathsf l_d$ locally uniformly in $\mathsf b$; the interval $[\mathsf b_-,\mathsf b_+]$ is an interval of labels under (Tor) on the side $\ell$ (Section~\ref{sec:cutoff-calibration}: $\mathsf b_-\ge\gamma^{-2}$, and $M\,Q(\mathsf b_+)\le L^m$, with $Q$ of Lemma~\ref{C-372-7}, which is the function $R$ of that definition, holds on the sides of $\mathfrak L(g)$ at every label of $\mathsf I(g)$, as in the proof of Proposition~\ref{C-372-9}(e)), so the numbers $\mathsf C_{\rm bf}$, $K_{\rm bf}$ of Definition~\textup{\ref{IV-62-62}}(5) formed at this interval depend on the label only through it, and Corollary~\textup{\ref{IV-62-66}}(a) bounds the difference of the laws of the block field of the depth $d$ at two consecutive cutoffs $K'$, $K'+1$, on the gauge invariant sets, by $\mathsf C_{\rm bf}\ell^3K'^{2\mathfrak p_0+1/2}L^{-K'/10}$ at every $\mathsf b\in\mathsf J$ and every $K'\ge K_{\rm bf}$ with $\lceil K'/1260\rceil\ge d$; summed over $K'\ge K$, with Corollary~\textup{\ref{IV-62-66}}(b), (c) for the limit and for all Borel sets, $\sup_E|\mu^K_d(E)-\mu_d(E)|\le\sum_{K'\ge K}\mathsf C_{\rm bf}\ell^3K'^{2\mathfrak p_0+1/2}L^{-K'/10}$ for all $\mathsf b\in\mathsf J$ and all $K\ge K_{\rm bf}$ with $\lceil K/1260\rceil\ge d$; $R_d$ and the rate of Proposition~\ref{C-372-9}(b) are explicit and continuous in $\mathsf b$; the balls $J'_d$ may move with $\mathsf b$.

So Steps~1 and~2 hold uniformly on $\mathsf J$, $\mathcal W_K\to\mathcal W$ uniformly, each $\mathcal W_K$ is continuous in $\mathsf b$, and the same bounds hold at labels $\mathsf b_K\in\mathsf J$.

\emph{Step 3: every torus and label.} Given $m\in\mathbb Z$ and $\mathsf b>0$, for large $s$ the number $g_s=(L^s\mathsf b-\bar B-1)^{-1/2}$ is a reference coupling below the thresholds of the frame, $L^s\mathsf b\in\mathsf I(g_s)$, and $2L^{m+s}\in\mathfrak L(g_s)$; the exponent of Proposition~\ref{C-372-10} is the same for all $g$, $\kappa$ is invariant under dilations, and Proposition~\ref{C-372-13}(a) gives $\mathcal W^{(\ell,\mathsf b)}_K(\mathsf C)=L^{s\kappa}\mathcal W^{(L^s\ell,L^s\mathsf b)}_{K-s}(L^s\mathsf C)$; so the limit exists and obeys the law of dilations.

\emph{Step 4: the floor.} For $C_i=n_1a_i\times n_2a_i$, Proposition~\ref{C-372-10}(b) applies with $n_{C_i}=4$ and $c_{C_i}$ fixed, both invariant under dilations, with $s_0(i)=\max\{i,s'_0\}=i+O(1)$ and no distance term; its bound is at most $C(\mathsf J)\frac{a_i}{\mathsf b}L^{-\vartheta(K-t)}\le C(\mathsf J)L^{-\vartheta(K-t-i)}$, an error that depends on the scale of the step only through $a_{K-t}/a_i$; it holds for $m_{t+1}$-almost every $W$, which is all that Proposition~\ref{C-372-12}(a) asks, so Step~1 runs with $C_\eta(i)=C(\mathsf J)L^{\vartheta i}$ and $s_0(i)$ in the place of the numbers of Proposition~\ref{C-372-10}(a). Let $i$ be so large that $(n_1^2+n_2^2)a_i^2<\ell^2/4$, so that $C_i$ lies in a ball of radius $<\ell/4$ and is a polygon, that $\kappa(C_i)\le k|C_i|/\mathsf b_-<\vartheta$ for all $\mathsf b\in\mathsf J$, that $s_0(i)\le d_i$, which holds for large $i$ as $d_i\ge\frac76i$, and that $\theta_{d_i}\le1$, which holds for large $i$ as $\theta_{d_i}\to0$ (below).

We run Proposition~\ref{C-372-12} down to the depth $d_i=\lceil(\frac43-\frac u2)i\rceil$. There 
\[
\theta_{d_i}=2B_3\sup_{K>d_i,\mathsf J}\varepsilon_{K-d_i}L^{2d_i}n_1n_2a_i^2\le c\,L^{\frac32d_i-2i}(1+d_i)^{\mathfrak p_0}\to0,
\]
as $\frac32d_i\le(2-\frac34u)i+\frac32$ and $\varepsilon_{K-d}\le A_0(2a_d/\mathsf b_-)^{1/2}(\log(\mathsf b_-/2a_d))^{\mathfrak p_0}$ for large $d$. The numbers of the free field of $C_i$ tend to $0$: $\eta_i\to0$ by Proposition~\ref{C-372-11}(a), so $e^{-\frac3{2\mathsf b_-}\eta_i}\le a_K^{-\kappa}\Xi_{K-d_i}$ (as $v_{K,d}\ge0$), $c_\Xi\to1$ and $L^{\kappa(d_i+1)}\to1$ ($\kappa d_i\le c\,a_ii$); and $\mathsf E'_{d_i}\to0$, as $\epsilon_1\le2C(\mathsf J)L^{\vartheta i}/(1-L^{-\vartheta/2})$ and $d_i\ge(\frac43-\frac u2)i$.

Proposition~\ref{C-372-12}(a), (c), (d) give the display of (4): by (d) and $\theta_{d_i}\le1$, $W^{(K-d_i)}_{C_i}\ge1-\frac12\theta_{d_i}^2\ge\frac12$ on $\mathsf O_{K-d_i}$, which carries $m_{K-d_i}$, so 
\[
Z^{-1}\int W^{(K-d_i)}_{C_i}\,dm_{K-d_i}\ge(1-\frac12\theta_{d_i}^2)Z^{-1}m_{K-d_i}(\mathcal V_{d_i})\ge(1-\frac12\theta_{d_i}^2)(1-2\mathsf V_\ell C_{\rm r}(1)L^{-d_i})
\]
 by (c), whatever the sign of the last factor, and this integral, being $\ge0$, is multiplied by $a_K^{-\kappa}\Xi_{K-d_i}\ge e^{-\frac3{2\mathsf b_-}\eta_i}$; both ends are $1+o(1)$ uniformly in $K>d_i$ and $\mathsf b\in\mathsf J$. For the squares $Q_i$ we use the mean form, Proposition~\ref{C-372-12}(b), at one $i$: by Proposition~\ref{C-372-10}(b) the mean defects are at most $C(\mathsf J)L^{-\vartheta(K-t-i)}$.

With $d=i+n_Q$, choose $n_Q$ and then $i$ so large that $\kappa(Q_i)\le\frac\vartheta2$, the free factor $a_K^{-\kappa}\Xi_{K-d}$ lies within a factor $\sqrt2$ of $L^{\kappa d}$, $L^{\kappa(d+1)}\le\sqrt2$, the summed step errors $2C(\mathsf J)L^{-\vartheta(n_Q+1)}/(1-L^{-\vartheta/2})$ are at most $\frac1{32}$, $2B_3L^{2n_Q}\sup_{K>d}\varepsilon_{K-d}\le1$ (so that the loop at the last level is at least $\frac12$ on tame fields), the mass of the measure without large fields is at least $\frac78$, and the rarity terms are at most $\frac1{32}$; then $\mathcal W_K(Q_i)\ge\frac1{\sqrt2}\cdot\frac12\cdot\frac78-\frac1{16}>\frac18$. For other rectangles we spread with Proposition~\ref{C-372-13}(b): $\rho_K(s,T)=\mathcal W_K(s\times T)$ has $\rho_K(s,\frac{T'+T''}2)^2\le\rho_K(s,T')\rho_K(s,T'')$, the powers of $a_K$ matching because lengths add; so $T\mapsto\log\rho_K(s,T)$ is convex, $\rho_K(s,nT)\ge\rho_K(s,2T)[\rho_K(s,2T)/\rho_K(s,T)]^{n-2}$, and a finite $\limsup$ (Step~1) on a set of rectangles closed downwards together with a positive $\liminf$ at one of them gives a positive $\liminf$ at all rectangles of the range, going down by $\log\rho(T)\ge2\log\rho(T')-\log\rho(2T'-T)$ and up by convexity, in each side.

\emph{Step 5: block fields; the static source.} $H\circ\mathcal M^{K-s_1}$ pulls out of the quotients, $\Phi^{(t)}[W_{\mathsf C}\,H\circ\mathcal M^{K-s_1}]=\Phi^{(t)}[W_{\mathsf C}]\,H\circ\mathcal M^{K-t-s_1}$ for $K-t\ge s_1$; so Steps~1 and~2 run with the factor $H$, $\sup|H|$ multiplying the errors (Proposition~\ref{C-372-12}(a)), and $Y_{d,C}$-means are those of $F^{[d]}_C$. For a rectangle we divide by $\mathcal W(C)>0$ (Step~4): the limit quotients $\omega(H)=\lim_K\langle W_CH(\mathcal M^{K-s_1}U)\rangle_K/\langle W_C\rangle_K$ form a real linear functional with $\omega(1)=1$ on the bounded gauge invariant Borel functions of all block fields, consistent between depths, $\omega(H\circ\mathcal M)=\omega(H)$, because $H\circ\mathcal M^{d'-s_1}$ is the same function seen at a finer depth; at the depth $s_1$ it is bounded, $|\omega(H)|\le\sup|H|\,(c_\Xi L^{\kappa d'}+\mathsf E'_{d'})/\mathcal W(C)$ with $d'=\max\{s_0,s_1\}$, by the display of (5) with $|W^{(K-d')}_C|\le1$, $Z^{-1}m_{K-d'}(\mathcal V_{d'})\le1$ and $a_K^{-\kappa}\Xi_{K-d'}\le c_\Xi L^{\kappa d'}$.

On the threads $H(V_{s_1})=H(\mathcal M^{d-s_1}V_d)$ $\mu$-almost surely, so $\int H(V_{s_1})Y_{d,C}\,d\mu=\int F^{[d]}_C\,(H\circ\mathcal M^{d-s_1})\,d\mu_d$, and Step~2, with $\mathsf l_d\asymp a_d^{-\kappa}$, gives the limits in $d$ and the two-sided bound for $\int Y_{d,C}\,d\mu$. For $H$ not gauge invariant: for a gauge transformation $v$ of $\Lambda_K$, $\mathcal M^{K-s_1}(U^v)=(\mathcal M^{K-s_1}U)^{v|\Lambda_{s_1}}$; averaging $v$ over the Haar measure and using the invariance of the Wilson measure, of $W_{\mathsf C}$, of $\mu$ and of $Y_{d,C}$ replaces $H$ by $\bar H$.

\emph{Step 6: mixed functions.} By Proposition~\ref{C-372-14}(d) at $X=W_{\mathsf C}$, $|Q_K(z)|\le M_0\mathcal W_K(\mathsf C,\theta\mathsf C)^{1/2}$ on half the polydisc; $(\mathsf C,\theta\mathsf C)$ is a family (its polygons lie $2\delta_1$ apart) with $\kappa(\mathsf C,\theta\mathsf C)=2\kappa<\frac12$, so Step~1 bounds the right side uniformly in $K>s_0$. For the limit at real $z$ we take the source off the loop by Proposition~\ref{C-372-14}(d) applied to the difference of the loop and a function of the block field on the side of $\mathsf C$; no step with a source is used. Let $\bar t=K-d$. As $c\in\frac12a_{i_0}\mathbb Z$, $\theta$ is a lattice isometry of $\mathsf{Iso}_{i_0}$ (Section~\ref{sec:setup-lattice-symmetries}); let $d_1:=i_0+\log_LM_1$, so that from the depth $d_1$ on $\theta$ maps $\Lambda_d$ and the partition into the $M_1$-cubes of the depth $d$, cubes of side $M_1a_d$, onto themselves (Lemma~\ref{lem:setup-lattice-covariance}).

Let $d$ be at least $\max\{s_0(\mathsf C),s_0(\mathsf C,\theta\mathsf C),d_0,d_1\}$, $s_0(\cdot)$ the number of Proposition~\ref{C-372-10}(a) at the family named, and so large that $M_1a_d\le\frac{\delta_1}8$ and $2\hat\varepsilon_d\le\epsilon_{\rm lc}$, so that $2\varepsilon_{K-d}\le\epsilon_{\rm lc}$ for every $K>d$ and Theorem~\ref{C-372-3}(5) applies at $d$; let $K>d$. Then $\theta$ commutes with the block averages, $\mathcal M^j(\theta U)=\theta\,\mathcal M^jU$ for $j\le\bar t$, and carries the constraint sets, criticality, regularity and the backgrounds over $V$ to those over $\theta V$ (Lemma~\ref{lem:setup-lattice-covariance}); $\varpi$ is $\theta$-invariant, so $\mathsf O_{\bar t}$ is; and the lattice Green function is $\theta$-invariant and the linearised block averages are $\theta$-covariant (Lemma~\ref{lem:setup-lattice-covariance}), so the projections $\Pi_n$ commute with $\theta$ and the lattice Coulomb energies are $\theta$-invariant.

Hence $W^{(\bar t)}_{\mathsf C}(\theta^{-1}V)=W^{(\bar t)}_{\theta\mathsf C}(V)$ on $\mathsf O_{\bar t}$ and $\Xi_{\bar t}(\theta\mathsf C)=\Xi_{\bar t}(\mathsf C)$. Let $\mathcal Y_{K,d}(V)$ be the loop of $\mathsf C$ at the background localised to a slab $\Omega'$ on the side of $\mathsf C$ as in Theorem~\ref{C-372-3}(5); then $Y=\Xi_{\bar t}(\mathsf C)\mathcal Y_{K,d}(\mathcal M^{\bar t}U)$ is a function of the bonds on the side of $\mathsf C$, and $\sup_{\mathsf O_{\bar t}}|W^{(\bar t)}_{\mathsf C}-\mathcal Y_{K,d}|\le\epsilon^{\rm loc}_d$, the bound of Theorem~\ref{C-372-3}(5). Now $\|W_{\mathsf C}-Y\|_\theta^2=\langle W_{\mathsf C}W_{\theta\mathsf C}\rangle_K-2\langle W_{\mathsf C}\,\theta Y\rangle_K+\langle Y\,\theta Y\rangle_K$, and three towers without source evaluate the three means: Proposition~\ref{C-372-12}(a) for the family $(\mathsf C,\theta\mathsf C)$, whose loop at the background is $W^{(\bar t)}_{\mathsf C}W^{(\bar t)}_{\theta\mathsf C}$ and whose factor is $\Xi_{\bar t}^2$; Proposition~\ref{C-372-12}(a) for $\mathsf C$ with the gauge invariant factor $\theta\mathcal Y_{K,d}$ of the block field; and Proposition~\ref{C-372-12}(c) for the function $\mathcal Y_{K,d}\,\theta\mathcal Y_{K,d}$ of the block field.

The main terms add up to $\Xi_{\bar t}^2Z^{-1}\int[W^{(\bar t)}_{\mathsf C}(W^{(\bar t)}_{\theta\mathsf C}-\theta\mathcal Y_{K,d})-(W^{(\bar t)}_{\mathsf C}-\mathcal Y_{K,d})\,\theta\mathcal Y_{K,d}]\,dm_{\bar t}$, of modulus at most $2\Xi_{\bar t}^2\epsilon^{\rm loc}_d$. With $a_K^{-\kappa}\Xi_{\bar t}\le c_\Xi L^{\kappa d}$ and the rarity bound at $p=2\kappa+1$,
\[a_K^{-2\kappa}\|W_{\mathsf C}-Y\|_\theta^2\le\sigma_d:=c_\Xi^2L^{2\kappa d}\bigl[2\epsilon^{\rm loc}_d+2\mathsf V_\ell C_{\rm r}(2\kappa+1)L^{-(2\kappa+1)d}\bigr]+\mathsf E'_d(\mathsf C,\theta\mathsf C)+2c_\Xi L^{\kappa d}\mathsf E'_d(\mathsf C),\]
free of $K$, and $\sigma_d\to0$ because $2\kappa<\vartheta<\frac12$ and $\epsilon^{\rm loc}_d$ falls faster than every power of $a_d$. At real $z$, Proposition~\ref{C-372-14}(d) at $X=W_{\mathsf C}-Y$ gives $|Q_K(z)-a_K^{-\kappa}\Xi_{\bar t}\,\mu^{K,w_z}_d(\mathcal Y_{K,d})|\le M_0\sigma_d^{1/2}$. By Proposition~\ref{C-372-14}(b) the probability $\mu^{K,w}_d$ gives at most $\sum_{t\le\bar t}\mathsf w^w_t\le2\mathsf V_SC^S_{\rm r}(\kappa+1)L^{-(\kappa+1)d}$ to the complement of $\mathsf O_{\bar t}\cap J'_d$, since the sourced measure without large fields has the null sets of $m_{\bar t}$ and is carried by $\mathsf O_{\bar t}\subset J'_d$; on $\mathsf O_{\bar t}\cap J'_d$, $|\mathcal Y_{K,d}-F'_d|\le\epsilon^{\rm loc}_d+\sup_{J'_d}|w^K_{\mathsf C}-F'_d|$ with $F'_d=\mathbf 1_{J'_d}\prod_iW^{[d]}_{C_i}$ (Proposition~\ref{C-372-9}(e)). With $a_K^{-\kappa}\Xi_{\bar t}\to\mathsf l_d$ and $\mu^{K,w}_d\to\mu^w_d$ (Proposition~\ref{C-372-14}(c)),
\[\limsup_K\bigl|Q_K(z)-\mathsf l_d\,\mu^{w_z}_d(F'_d)\bigr|\le M_0\sigma_d^{1/2}+c_\Xi L^{\kappa d}\epsilon^{\rm loc}_d+4c_\Xi\mathsf V_SC^S_{\rm r}(\kappa+1)L^{-d}\to0,\]
so $\lim_KQ_K(z)=\lim_d\mathsf l_d\,\mu^{w_z}_d(F'_d)$ exists at every real $z$ of half the polydisc. A bounded family of holomorphic functions that converges on a real cube converges locally uniformly to a holomorphic limit (Vitali's theorem); so $Q_K\to Q$, the derivatives converge, and Proposition~\ref{C-372-14}(e) with the limits $S_m$ of \YMMref{(2212)}{813}, for pairwise disjoint supports, gives the limit of $\mathcal W_K(\mathsf C;f)$. For two orthogonal mirrors the same argument runs with the four images of $\mathsf C$ and Proposition~\ref{C-372-14}(d) twice, on the quarter $\sum_q|z_q|\|f_q\|_\sharp<\rho_D/4$ of the polydisc, the towers at the subfamilies of $(\mathsf C,\theta_1\mathsf C,\theta_2\mathsf C,\theta_1\theta_2\mathsf C)$, of $\kappa$-value up to $4\kappa<\frac12$, and a region $\Omega'$ that is the intersection of two slabs; the limit then exists at the real $z$ with $|z_q|<\rho_D/(4n\max_q\|f_q\|_\sharp)$, a real cube in that quarter, and Vitali's theorem applies there.

For $n=1$, $\mathcal W_K(\mathsf C;f)$ is linear in $f$; as each coordinate extent of $\mathsf C$ is $<\frac\ell2$, every point outside the least closed coordinate box that holds $\mathsf C$ lies, with a neighbourhood, beyond a coordinate mirror that misses the box, with $c\in\frac12a_i\mathbb Z$ for some $i$; a Lipschitz partition of unity subordinate to finitely many such mirrors cuts $f$ into pieces, each of which the case of one mirror covers. With walls, the number of boxes is even in every direction that the walls cut and the walls lie in $\frac12a_i\mathbb Z$, so that both fixed planes $c$ and $c+\frac\ell2$ of every wall mirror are walls and each box has exactly one image in every box; and each $f_q$ is supported in one box, a function whose support misses the walls being the sum of its pieces in the boxes, with the derivatives in its variable the sums of those in the variables of its pieces.

Then Proposition~\ref{C-372-14}(d) on the polydisc $\sum_q|z_q|\|f_q\|_\sharp<\rho_D/N$, Step~1 for each family $\mathsf C^{\rm f}_\iota$, of $\kappa$-value $N\kappa(\mathsf C_\iota)<\frac12$, and Cauchy's estimates bound the derivatives $\partial_AQ_K(0)$ uniformly in $K$, for all supports off the walls; by Proposition~\ref{C-372-14}(e) this bounds $\mathcal W_K(\mathsf C;f_1,\dots,f_n)$ uniformly in $K$ for $n=1$, where $\mathcal W_K(\mathsf C;f_1)=\partial_1Q_K(0)$, and for pairwise disjoint supports, where the cumulants $S^{\sf T}_{|\beta|,K}(f_\beta)$ of the blocks of at least two elements converge \YMMref{(2212)}{813} and so are bounded in $K$. The law of dilations and the positive form with mixed entries hold at every cutoff (Proposition~\ref{C-372-13}(a), (b)) and pass to the limits; for the latter the plane $x_\mu=\hat c+\frac\ell4$, $\hat c\in\frac12a_{i_0}\mathbb Z$ the place of the mirror $\theta'$ of (6), a fixed plane of the coordinate mirror with the place $\hat c+\frac\ell4$, which lies in $\frac12a_{i_0}\mathbb Z$ as $\frac\ell4=\frac12L^m$, parts polygons and supports, so that every entry converges.

\emph{Step 7: beside the loop.} Let $X=W_C-\psi$, a function of the bonds on the side of $C$, and $Y=H-\langle H\rangle_K$, a real function of the bonds of $S$. Reflection positivity at the site mirror and Schwarz' inequality (as in Proposition~\ref{C-372-13}(b)) give $|\langle XY\rangle_K|^2\le\langle X\,\theta X\rangle_K\langle\theta Y\,Y\rangle_K$, with $\langle X\,\theta X\rangle_K=\langle W_CW_{\theta C}\rangle_K-\psi^2\ge0$ by the $\theta$-invariance of the measure, $\langle\theta Y\,Y\rangle_K=\langle\theta H;H\rangle_K\ge0$, and $\langle XY\rangle_K=\langle W_C;\mathcal O(f)\rangle_K-\langle W_C;H_0\rangle_K$; divide by $\psi$. The explicit ceiling uses the law of the truncated two-point function of the action density, Theorem~\textup{\YMMref{AV}{851}}.

Let, as in (7), $g\le\gamma_W$ in the frame, $\gamma_W$ the ceiling of equation~\YMMref{(2346)}{850}; let $\vartheta_{\rm sh}\in(0,1)$ be a shape parameter, and $\mathsf d\in(0,\mathsf d_0(g,\vartheta_{\rm sh})]$ with $K-s(\mathsf d)\ge1$, $\mathsf d_0(g,\vartheta_{\rm sh})$ the threshold of equation~\YMMref{(2350)}{851} and $s(\mathsf d)$ the witness depth of equation~\YMMref{(2321)}{843} at $g$. As $\mathsf b\in\mathsf I(g)$, in the frame $(K,L^K\mathsf b)$ is a first-passage datum at $g$ with the label $\mathsf b$ (Corollary~\YMMref{74.2}{616}), and the side $\ell$ of the frame lies in $\mathfrak L(g)$. Then for every test pair $(h_1,h_2)\in\mathfrak T_{T_\ell}(\mathsf d;\vartheta_{\rm sh})$ on a torus $T_\ell$ with $20\mathsf d\le\ell$ (Lemma~\ref{lem:ward-test-pairs}: non-negative Lipschitz functions, neither $0$, with supports in closed balls of radius $\mathsf d$ at a distance between $\mathsf d$ and $6\mathsf d$, $\int h_j\ge\vartheta_{\rm sh}\|h_j\|_\infty\mathsf d^3$ and ${\rm Lip}\,h_j\le\|h_j\|_\infty/(\vartheta_{\rm sh}\mathsf d)$)
\[
\begin{multlined}
S^{\sf T}_{2,K}(h_1,h_2)=a_K^{-2}\sum_{p,q}h_1(x(p))h_2(x(q))\langle W_{\partial p};W_{\partial q}\rangle_K=\frac9{16\pi^2}\mathsf b^{-2}\int\!\!\int h_1(x)h_2(y)|x \\
-y|^{-6}\,dx\,dy\,(1+\epsilon'),\qquad|\epsilon'|\le\epsilon_{\rm tp},
\end{multlined}
\]
with $\frac9{16\pi^2}$ the constant $\mathsf c_2$ of \YMMref{C}{7}, the double integral the free kernel \YMMref{(2240)}{822}, and $\epsilon_{\rm tp}:=\bar\varepsilon_{g,\vartheta_{\rm sh}}(\mathsf d)$, the tolerance equation~\YMMref{(2349)}{851}, one function for all these data, non-decreasing, tending to $0$ with $\mathsf d$ and at most $\frac13$ on $(0,\mathsf d_0(g,\vartheta_{\rm sh})]$ (Theorem~\textup{\YMMref{AV}{851}}); so $|\epsilon'|\le\frac13$ and the bound $(\frac{4\pi}3)^2\|h_1\|_\infty\|h_2\|_\infty$ of the free kernel of a test pair, two balls of radius $\mathsf d$ at a distance $\ge\mathsf d$ (equation~\YMMref{(2247)}{824}), give $S^{\sf T}_{2,K}(h_1,h_2)\le\frac43\,(\frac{4\pi}3)^2\,\frac9{16\pi^2}\,\mathsf b^{-2}\|h_1\|_\infty\|h_2\|_\infty$, the bound of Theorem~\textup{\YMMref{AV}{851}}.

As $\mathcal O(f)=a_K^{-1}\sum_pf(x(p))(1-W_{\partial p})$, $\langle\theta\mathcal O(f);\mathcal O(f)\rangle_K=a_K^{-2}\sum_{p,q}f(\theta x(p))f(x(q))\langle W_{\partial p};W_{\partial q}\rangle_K=S^{\sf T}_{2,K}(\theta f,f)$, and $\sqrt{9/16\pi^2}=\frac3{4\pi}$: for a test pair $(\theta f,f)$ the inequality at $H_0=0$ gives the bound of (7) with the factor $1+\epsilon_{\rm tp}$. On the whole class of the law, $\frac43\,(\frac{4\pi}3)^2\,\frac9{16\pi^2}=\frac43$ gives $\langle\theta\mathcal O(f);\mathcal O(f)\rangle_K\le\frac43\mathsf b^{-2}\|f\|_\infty^2$, and the inequality at $H_0=0$ gives the ceiling $\frac2{\sqrt3}\mathsf b^{-1}\|f\|_\infty[\langle W_CW_{\theta C}\rangle_K/\psi^2-1]^{1/2}$ with no further condition.

The identity follows from the same expression of $\mathcal O(f)$: 
\[
\langle W_C\mathcal O(f)\rangle_K/\psi-\langle\mathcal O(f)\rangle_K=a_K^{-1}\sum_pf(x(p))[\langle W_{\partial p}\rangle_K-\langle W_CW_{\partial p}\rangle_K/\psi],
\]
and $\langle W_{\partial p}\rangle_K>0$ (Proposition~\ref{C-372-13}(c)). The derivative is $\frac d{d\mathsf b}=L^K\frac d{d\beta}$ applied to $\log\langle W_C\rangle_K$. For a rectangle with $2\kappa<\frac12$ the three quantities converge by (6) at $n=1$ with (4), by (2) at the family $(C,\theta C)$, and by \YMMref{(2212)}{813}, the supports of $\theta f$ and $f$ being disjoint.

\emph{Step 8: short distances.} (a) We run the tower in its mean form, Proposition~\ref{C-372-12}(b), down to the depth $i$ of the corners, that is to the level $t_i=K-i$; when $\max\{n_1,n_2\}<L$ this is the level $t_\varrho$ with $a_{K-t_\varrho}\le\varrho_i<La_{K-t_\varrho}$. The number $s_0$ of Proposition~\ref{C-372-10}(b) for $C_i$ is $\max\{i,s'_0\}$, $s'_0$ formed from the constants, $\mathsf b_-$ and $n_{C_i}=4$; so $s_0=i$ for all large $i$, and every step that the tower uses, at a level $t<t_i$, has $K-t>i=s_0$. There Proposition~\ref{C-372-10}(b) bounds the steps with a constant free of the size of the loop; weighed by the free factors, whose ratios between the levels are bounded by Proposition~\ref{C-372-11}(c), they contribute $C_1\frac{|C_i|}{\mathsf b}(\frac{\varrho_i}{\mathsf b})^{\vartheta/2}$, and the rarity terms contribute $C_3\mathsf V_\ell C_{\rm r}(p)\varrho_i^{p-1/2}$.

At the depth $i$ a tame background has $1-W^{(t_i)}_{C_i}\le\frac12\theta^2=2B_3^2n_1^2n_2^2\varepsilon_{K-i}^2\le4B_3^2A_0^2\mathsf n^3\frac{\varrho_i}{\mathsf b}(1+\log\frac{\mathsf n\mathsf b}{\varrho_i})^{2\mathfrak p_0}$ (Proposition~\ref{C-372-12}(d), with $x_{K-i}\ge\frac23\mathsf b/a_i$, $n_1^2n_2^2a_i\le\mathsf n^3\varrho_i$ and $\mathsf b/a_i\le\mathsf n\mathsf b/\varrho_i$). The free factor down to that level is $a_K^{-\kappa}\Xi_{t_i}=\exp[-\frac3{2\mathsf b}(c_{\rm lat}|C_i|+\mathcal J_\ell(C_i)+E_{a_K}(C_i)-v_{K,i}(C_i))]$ (Proposition~\ref{C-372-11}(a), (c)), and $0\le\frac3{2\mathsf b}v_{K,i}(C_i)\le\kappa\log^+\frac{\varrho_i}{a_i}+c\frac{|C_i|}{\mathsf b}\le\kappa\log L+c\frac{|C_i|}{\mathsf b}$, $c$ depending on $n_1$, $n_2$: at the depth of the corners the logarithm of $v$ (Proposition~\ref{C-372-11}(b)) and that of $\mathcal J_\ell$ cancel up to $\frac{|C_i|}{2\pi}\log\max\{n_1,n_2\}$ and a multiple of $|C_i|$, and $\kappa\log\max\{n_1,n_2\}\le\kappa\log L+k\frac{|C_i|}{\mathsf b}\log\frac{\mathsf n}L$; $|\log(1-z)|\le2|z|$ for $|z|\le\frac12$. (b) Let $a_*:=\varrho(\mathsf b/\varrho)^{\vartheta/(3+\vartheta)}$ and let $d_*$ be the least depth with $a_{d_*}\le a_*$, so that $a_*/L<a_{d_*}\le a_*$; for small $\varrho$, $a_{d_*}>a_i$, and the depth $d_*$ is coarser than the corners.

At all the levels with $a_{K-t}\le a_*$ the defect of one step is at most $C_*\frac\varrho{\mathsf b}(\frac{a_{K-t}}{\mathsf b})^\vartheta$ at $m_{t+1}$-almost every $W$, hence in the mean, $|\bar{\mathsf D}_{t+1}|\le C_*\frac\varrho{\mathsf b}(\frac{a_{K-t}}{\mathsf b})^\vartheta$ with $\bar{\mathsf D}_{t+1}$ of Proposition~\ref{C-372-12}(b), $C_*$ depending on the constants, $c_0$, $\vartheta$ and $\mathsf J$ alone: this is the bound of one step of (8)(b), shown at the end of this step. We stop the tower, in its mean form, Proposition~\ref{C-372-12}(b), at the depth $d_*$. Proposition~\ref{C-372-11}(c), valid at every depth, and (d) give the factor $e^{-\frac3{2\mathsf b}(\mathcal I_{a_K}(C,C)-v_{K,d_*})}$ with $0\le v_{K,d_*}(C)\le C^{\rm bg}A_C^2/a_{d_*}^3$, $A_C$ the area of $C$; on tame fields $0\le1-W^{(K-d_*)}_C\le4B_3^2A_0^2\frac{A_C^2}{\mathsf ba_{d_*}^3}(1+\log\frac{\mathsf b}{a_{d_*}})^{2\mathfrak p_0}$ (Proposition~\ref{C-372-12}(d)).

As $A_C\le100\varrho^2$ and $a_*/L<a_{d_*}\le a_*$, these are $O(\varrho^4/a_*^3)$ and $O(\frac{\varrho^4}{\mathsf ba_*^3}(1+\log\frac{\mathsf b}{a_*})^{2\mathfrak p_0})$, with constants that hold $10^4L^3$ (and $\log L$ in the logarithm); the step errors add up to $O(\frac\varrho{\mathsf b}(\frac{a_*}{\mathsf b})^\vartheta)$; the terms of the large fields, $\sum_{t\le K-d_*}\Xi_{t-1}\omega_t$ with $|\omega_t|\le\mathsf w_t$, are by the rarity bound at $p=2$ at most $c\,\mathsf V_\ell C_{\rm r}(2)\,a_{d_*}^2\le c\,\mathsf V_\ell C_{\rm r}(2)\,a_*^2$ times the free factor, and $a_*^2=\mathsf b^2(\frac\varrho{\mathsf b})^{2-2\alpha}$ with $\alpha=\frac\vartheta{3+\vartheta}<\frac17$, below the remainder as $2-2\alpha>1+3\alpha$; and $\frac\varrho{\mathsf b}(\frac{a_*}{\mathsf b})^\vartheta=\frac{\varrho^4}{a_*^3\mathsf b}=(\frac\varrho{\mathsf b})^{1+\frac{3\vartheta}{3+\vartheta}}$: the choice of $a_*$ balances the two.

Proposition~\ref{C-372-11}(a) and $|\mathcal J_\ell-\mathcal J|\le c\varrho^4\ell^{-3}$ finish; $\log(\mathsf b/\varrho)$, and not $\log\varrho$, is the logarithm that dilations leave invariant. The bound of one step at the levels with $a_{K-t}\le a_*$. At the levels finer than the lattice of the corners, $K-t>s_0=\max\{i,s'_0\}$, which is $i$ for small $\varrho$, it is Proposition~\ref{C-372-10}(b): its constant depends on the constants, $n_C=4$ and the chord-arc constant of $C$, which is bounded below in terms of $c_0$, its factor $|C|/\mathsf b$ is at most $40\varrho/\mathsf b$, and its bracket, without a distance term for one polygon, is at most $c\,(a_{K-t}/\mathsf b)^\vartheta$ with $c$ depending on the constants, $\vartheta$ and $\mathsf J$.

At the remaining levels, $a_i\le a_{K-t}\le a_*$, at least as coarse as the lattice of the corners, $a_{K-t}\ge a_i>c_0\varrho/2L$ and $|C|\le40\varrho$, so that the loop passes a number of cubes of the level that is bounded in terms of $L$ and $c_0$, and $|C|/a_{K-t}<80L/c_0$. There the proof of Proposition~\ref{C-372-10}(b) runs with the segments $C_j$ of $C$, one for each cube of the level that the loop passes, in the place of the bonds of the level, down to one segment where the cells are larger than the loop. The shift of Proposition~\YMMref{51.5}{396} gives $W^{(t)}_C(V_W(B))=W_C(e^{i\eta_t\mathbf A(B)}U_+)$ exactly for every bare loop (step~1 of the proof of Proposition~\ref{C-372-10}(a)); so $W^{(t)}_C(V_W(B))=w(\Upsilon(B))$ and $W^{(t+1)}_C(W)=w(0)$, with $w$ as in Lemma~\ref{C-372-16}(c), $g_j=U_+(C_j)$ the holonomy of $U_+$ along the segment, and the rotations $\Upsilon_j(B)=\frac1i\log[(e^{i\eta_t\mathbf A(B)}U_+)(C_j)\,U_+(C_j)^{-1}]$ of the segments.

Each segment is one insertion into the four-family Gaussian
cluster expansion of Theorem~\YMMref{R}{361}, applied to the
two potentials in $\nu_W$ as in Step~3 of the proof of
Proposition~\ref{C-372-10}(a), with a linear part of norm of the order $g_t|C_j|/a_{K-t}$ and a remainder of the order $(g_t|C_j|/a_{K-t})^2$, so that $\log\nu_W(e^{(z,\Upsilon)})$ is a sum of local terms on a polydisc of radius of the order $a_{K-t}/(g_t|C|)$. The first cumulants, and the defects of the pair cumulants from their values $s_{jj'}{\rm R}(G_{j\to j'})$ at the trivial triple, are colour vectors and colour defects that vanish at the flat background, and Schwarz' lemma in the background gives them the factor $\varepsilon_{t+1}$, as in step~5 of the proof of Proposition~\ref{C-372-10}(a); the pair values at the trivial triple sum to the slice $g_t^2\eta_t(X_C,\mathcal C_{(t+1)}X_C)$ of the free Coulomb energy, as in its step~7, at depths coarser than the corners as well (Proposition~\ref{C-372-11}(c) holds at every depth); and $\nu_W\ge0$ for $m_{t+1}$-almost every $W$.

Lemma~\ref{C-372-16}(c) then bounds the defect of the step by $E_1+E_2+E_3$ and the term of the label, at most $c\,|C|\mathsf b^{-2}a_{K-t}$; each of $E_1$, $E_2$, $E_3$ is at most $\frac{|C|}{\mathsf b}(\frac{a_{K-t}}{\mathsf b})^{1/2}y_t^{\mathfrak p}$, $y_t=\log x_t$, $\mathfrak p$ formed from $\mathfrak p_0$, times a polynomial in $|C|/a_{K-t}$ with coefficients formed from the constants, since a cumulant of order $m$ carries the factor $(g_t|C|/a_{K-t})^m$, the first cumulants besides the amplitude $\delta_t=g_t\mathfrak b_t$ of the non-Gaussian potentials, the Gaussian part being even, the variables are at most $c\,\delta_t|C|/a_{K-t}$, $g_t^2$ is of the order $a_{K-t}/\mathsf b$, and the first cumulants and the defects of the pair cumulants carry the factor $\varepsilon_{t+1}$, of the order $(a_{K-t}/\mathsf b)^{1/2}y_t^{\mathfrak p_0}$.

As $|C|/a_{K-t}<80L/c_0$, $|C|\le40\varrho$ and $(\frac a{\mathsf b})^{1/2}y^{\mathfrak p}\le c\,(\frac a{\mathsf b})^\vartheta$ for $\vartheta<\frac12$, the bound is at most $C_*\frac\varrho{\mathsf b}(\frac{a_{K-t}}{\mathsf b})^\vartheta$. No level coarser than $a_*$ is used; at the levels with $a_{K-t}^3\ge\varrho^2\mathsf b$, coarser than $a_*$, the step bound holds by itself, since there $1-\lambda_{t+1}$ and $1-W^{(t)}_C$ on tame fields are at most of the order $\varrho^4/(\mathsf ba_{K-t}^3)$ up to logarithms, below $C_*\frac\varrho{\mathsf b}(\frac{a_{K-t}}{\mathsf b})^\vartheta$.

\subsection{Proof of Theorem~\ref{C-372-6}}

We work in lattice units, $\beta=\mathsf b/a$, $n=|C|/a$ steps. By Lemma~\ref{C-372-17}(a) the steps of $C$ at distance $\ge2\rho$ from the corners form an admissible set $I$ of $n-16\rho$ steps (of $n-4\rho s$ steps for polygons with $s$ pieces), whose slabs lie in disjoint cells. Lemma~\ref{C-372-17}(g) at spin $\frac12$, $c_i=\frac34$, gives $|\langle W_C\rangle|\le e^{-\mathcal G}\mathfrak M_\nu^{\omega_\rho}$: across each admissible step the bonds of the slab that cross the disc of radius $\rho$ about it are rotated together with the profile $\phi_\rho$, whose Dirichlet energy is about $2\pi/\log\rho$, at an action cost $\approx2\pi\beta/\log\rho$; the information inequality in three colours gives the factor $e^{-k\log\rho/\beta}$ per slab, up to the second-order term of $\mathcal G$ and the constant $1.2144$, and the bounds multiply exactly over the disjoint cells; with $n$ slabs, $e^{-kn\log\rho/\beta}=\rho^{-k|C|/\mathsf b}$.

The second variation (Lemma~\ref{C-372-17}(d)) limits $\rho$: its second member, summed over the disc, asks $\omega_\rho=40\rho^4\sum_ic_i/(\beta E_\rho)^2\le1$, $E_\rho\sim2\pi/\log\rho$; with $\rho=\lfloor(\frac{\beta}{\Lambda\log^3\beta})^{1/4}\rfloor$ one has $\omega_\rho\le c/\log\beta$, so that $\mathfrak M_\nu^{\omega_\rho}$, with Lemma~\ref{C-372-17}(f), the corrections of $\mathcal G$ and the $16\rho$ omitted steps give the prefactor $6e^{3.1\mathsf b/\ell}1.58^{|C|/\mathsf b}$, and $\rho^{-1}\le2(\Lambda\log^3\beta/\beta)^{1/4}$ gives the power $[\Lambda\frac a{\mathsf b}\log^3\frac{\mathsf b}a]^{k|C|/4\mathsf b}$, the factor $2^{k|C|/\mathsf b}$ going into the prefactor. For the second bound let $\rho=2^S\in(\frac12\beta^{1/3},\beta^{1/3}]$ and $p=k|C|/\mathsf b=kn/\beta<3$, and keep the cells $T_\rho(c_i)$ of the $n-16\rho$ admissible steps.

On each cell, for every dyadic $2\le r\le\rho$, the moves with the profile $\phi_r$, taken as $0$ at $|y|\ge r$, rotate bonds of the cell alone, and Lemma~\ref{C-372-17}(b)--(e), as in the proof of (g), give at every field $\prod_i\|Q_i\|\le\min_re^{-\mathcal G_r+Y_r}$, with $\mathcal G_r=\sum_i\frac{2c_i}{\beta E_r}(1-\frac{2c_i}{\beta E_r})$, $Y_r=\sum_i2c_iX^{(r)}_i/(\beta E_r^2)$ and $X^{(r)}_i$ the $X_i$ of that proof formed with $\phi_r$; by (b), $e^{-\mathcal G_r}\le e^{\delta_0+1.2144p}r^{-p}$, where $\delta_0=(16k\rho\log\rho+pk\log^2\rho)/\beta$ pays for the $16\rho$ omitted steps and the second-order term of $\mathcal G_r$. The sharper estimate bounds the mean of $Y_r$ instead of an exponential moment: in comb frames $a(H_{e'})$ at a rung $e'$ of height $t_{e'}$ is $1-W$ of the boundary of the strip $S_{e'}$ of $1\times|t_{e'}|$ plaquettes, whose mean is $1-w(1,|t_{e'}|)\le1-w(1,1)^{|t_{e'}|}\le|t_{e'}|\,\nu(a_p)$ by Proposition~\ref{C-372-13}(c), in the place of the bound $|t_{e'}|\sum_{q\in S_{e'}}a_q$ at every field; so $\nu(X^{(r)}_i)\le c\,r^3\nu(a_p)$, $r^3$ in the place of $r^4$, and $\nu(Y_r)\le c\,p\,r^3\log^2r\,\bar a$ with $\bar a=\nu(a_p)\le\beta^{-1}\log\mathfrak M_\nu$ (Lemma~\ref{C-372-17}(f)): the plaquette mean enters linearly, and this bound of a mean takes the place of the condition $\omega_\rho\le1$ and of the weight $40\rho^4$.

Let $r_*$ be the largest dyadic $r\le\rho$ with $Y_r\le p$; on $\{r_*=r\}$, $r<\rho$, $Y_{2r}>p$, an event of probability at most $\nu(Y_{2r})/p$ by Markov's inequality, and on the event that $Y_r>p$ at every $r$ the bound $1$ serves. So, as Lemma~\ref{C-372-17}(h) states,
\[|\langle W_C\rangle|\le e^{\delta_0+2.2144p}\rho^{-p}\Bigl[1+\frac{72\,\bar a\rho^3\log^2\rho}{2^{3-p}-1}\Bigr]+35\,\bar a,\]
the sum over the dyadic $r<\rho$ of $r^{-p}(2r)^3\log^2(2r)$ being geometric and led by its last term exactly when $p<3$. With $\rho^3\le\beta$, so that $\bar a\rho^3\le\log\mathfrak M_\nu$, with $\rho^{-p}\le2^p\beta^{-p/3}$, $\bar a\le\beta^{-p/3}\log\mathfrak M_\nu$ and $\log\mathfrak M_\nu\le192\frac{\mathsf b}\ell+4(14.15+3\log\beta)$, where $192\frac{\mathsf b}\ell\le\frac{384k}p$ as $|C|<2\ell$, this is at most $C(\frac{k|C|}{\mathsf b})(\frac a{\mathsf b})^{|C|/4\pi\mathsf b}\log^3\frac{\mathsf b}a$ for $\beta\ge\beta_2$, the threshold, depending on $p$ alone, from which on $n\ge16\rho+4$ and $\delta_0\le1$: $\frac p3=\frac{|C|}{4\pi\mathsf b}$, and the third logarithm is $\log^2\rho$ times $\log\mathfrak M_\nu$.

For the Gibbs states the strips of even height are treated by the reflections in the planes of sites, those of odd height by $\sqrt{a(gh)}\le\sqrt{a(g)}+\sqrt{a(h)}$. For the spin $j$, $c_j=j(j+1)$ replaces $\frac34$, that is $k\mapsto\frac43c_jk=\frac{j(j+1)}\pi$, and $\kappa_j\le20\log\beta$ keeps $\omega_\rho$ admissible. For a family of such polygons whose pieces, those of different polygons included, are at sup-distance at least $(2\rho+1)a$ when they have no common corner, Lemma~\ref{C-372-17}(a) applies to the whole family and gives one admissible set of $n-4\rho s$ steps, $n$ and $s$ the numbers of all the steps and of all the pieces, and Lemma~\ref{C-372-17}(e), (g) bound $\langle\prod_iW_{C_i}\,G\rangle$ in the same way, a factor $G$ of the bonds farther than $\rho a$ from the polygons being a function of the bonds off $B$; a fixed family at positive distances has this property as soon as $(2\rho+1)a$ is below these distances, that is for all small $a$, which is what the limits as $a\to0$ use.

The proof uses of $\nu$ only the conditional laws of finite sets of bonds, reflection positivity and invariance (through Lemma~\ref{C-372-17}(f), and for the second bound through Proposition~\ref{C-372-13}(c) for the strips); for the Gibbs states the bound of (f) gives $3$ in the place of $6e^{3.1\mathsf b/\ell}$. So all characters of the holonomy of non-zero spin have means tending to $0$: its class tends to the Haar law of $\mathrm{SU}(2)$, whose image under $\frac12\operatorname{tr}$ has the density $\frac2\pi\sqrt{1-w^2}$ (the $\chi_j$ are the Chebyshev polynomials of the second kind in $W_C$, orthogonal for $\sigma$, and polynomials are dense among the continuous $h$); and $\langle W_C^2\rangle=\frac14(1+\langle\chi_1(U(C))\rangle)\to\frac14$, from $W_C^2=\frac14\chi_{1/2}(U(C))^2=\frac14(1+\chi_1(U(C)))$.

For a rectangle $\langle W_C\rangle>0$ (Proposition~\ref{C-372-13}(c)), so $\mathrm{Var}(W_C/\langle W_C\rangle)=\langle W_C^2\rangle/\langle W_C\rangle^2-1=(1+\langle\chi_1(U(C))\rangle)/(4\langle W_C\rangle^2)-1$; as $\langle\chi_1(U(C))\rangle\to0$ and the first bound gives $\langle W_C\rangle^{-2}\ge(6e^{3.1\mathsf b/\ell}1.58^{|C|/\mathsf b})^{-2}[\Lambda\frac a{\mathsf b}\log^3\frac{\mathsf b}a]^{-k|C|/2\mathsf b}$, the variance tends to infinity at least at this rate. If $c_aW_C$ had a limit law while $|c_a\langle W_C\rangle|\ge c>0$, then $|c_a|\ge c/|\langle W_C\rangle|\to\infty$; as the law of $W_C$ tends to the semicircle law, which gives no mass to $0$, $\langle\mathbf 1\{|c_aW_C|>R\}\rangle\to1$ for every $R$, and the laws of $c_aW_C$ are not tight: a contradiction.

The rotations touch only bonds within $\rho a\to0$ of the polygons, act on several polygons at once and leave $G$ unchanged, so $\langle\prod_i\chi_{j_i}(U(C_i))\,G\rangle\to0$ unless all $j_i=0$, which is the asymptotic independence. The floor: $\langle W_C\rangle\ge\langle W_{\partial p}\rangle^{A_C/a^2}$ by Proposition~\ref{C-372-13}(c), and $\langle W_{\partial p}\rangle=1-\langle a_p\rangle\ge1-\beta^{-1}\log\mathfrak M_\nu\ge e^{-\frac32\beta^{-1}\log\mathfrak M_\nu}$ when $\beta^{-1}\log\mathfrak M_\nu\le\frac12$ (Lemma~\ref{C-372-17}(f)), with $\frac{A_C}{a^2}\beta^{-1}=\frac{A_C}{a\mathsf b}$; for the Gibbs states the same argument, with reflections in planes of sites and even numbers of steps, and $\frac32(28.3+6\log\beta)\le43+9\log\beta$.

\subsection{Scope of the conclusions}

At a fixed depth the variable $Y_{d,C}$ of Theorem~\ref{C-372-3} is the smoothed loop on $\{V_d\in J_d\}$ and $0$ off it, and Proposition~\ref{C-372-8}(d) bounds the probability of $\{V_d\in J_d\}$ below; the starting depth is random and by its form no stopping time, since $\{\tau\le D\}$ involves all the coordinates of index $\ge D$; the loop of the depth $d$ uses $H^\infty_{d,n}(V_d)$ and not the finer members of the thread, so that $W^{[d]}_C(V_d)$ and $W^{[d+1]}_C(V_{d+1})$ are two different functions on the threads, and the limit in the depth is a limit of means and of quotients of means (Theorem~\ref{C-372-4}(2), (5)); for the sharp loop itself no multiple has both a limit law and a mean that stays away from $0$, by the non-tightness consequence of the semicircle law (Theorem~\ref{C-372-6}).

The bounds of the convergence in $n$ and in $K$ hold the factor $L^d\varepsilon_{[d]}=A_0\mathsf b^{-1/2}L^{d/2}(\log x_{[d]})^{\mathfrak p_0}$, which grows with $d$. On the side $\ell$ depend the tails of $\tau$, which hold $\ell^3$ or $e^{C\ell^3}$; the floors $\mathsf p^{[d]}$, $\mathsf q^{[d]}$; the speed of $K\to\infty$, as $a_s(K)$ holds $(\ell L^s)^{3/2}$, $\mathfrak r_K$ holds $\ell^3$ and $K_{\rm bf}$ depends on $\ell$; and $\mathsf V_\ell$; the event $\{V_d\in J_d\}$ asks all the $3(\ell L^d)^3$ plaquettes. Free of $\ell$ are the bounds in $n$ on the tower, on curves and on small rectangles (Proposition~\ref{C-372-9}(a), the first bound of (c)), the locality bound of Theorem~\ref{C-372-3}(5), the numbers $c_{\rm lat}$, $c_Q$, $c$ and $c_3$ (Proposition~\ref{C-372-11}), and, for $\ell\ge\ell_0$, all the ceilings of Theorem~\ref{C-372-6}.

Theorem~\ref{C-372-4} concerns $|\mathsf C|<\frac{2\pi}3\mathsf b$ and is a perimeter law, two-sided for rectangles. Three free photon fields at the label $\mathsf b$ have the same features: their block curvatures obey at each depth, uniformly in $K$, bounds of the same form as those of Proposition~\ref{C-372-8}(c) and of $R^\natural_d$ (Proposition~\ref{C-372-15}(a)), and their sharp loops have the same exponent $\kappa$, the same perimeter constant $c_{\rm lat}|C|+\mathcal J_\ell(C)$ and the same factor $\lambda_{t+1}$ of one step as the loops of $\mathrm{SU}(2)$ in Theorem~\ref{C-372-4} and Proposition~\ref{C-372-10} (Proposition~\ref{C-372-15}(b), (c)). Theorem~\ref{C-372-4}(7) is an upper bound: for three free photon fields the response of the action density is $-\frac3{2\mathsf b^2}\int f|B_{C,\ell}|^2\,dx$, with the field $B_{C,\ell}$ of $C$ on the torus, within $c|C|\,{\rm diam}(C)\ell^{-3}$ of its Biot--Savart field $B_C$ in the ball that holds the polygon (Proposition~\ref{C-372-15}(d)).

The quotients of Theorem~\ref{C-372-4}(5) formed for three free photon fields, $\lim_K\mathbb E[W^{\rm fr}_C\,H]/\mathbb E\,W^{\rm fr}_C$ for $H$ a bounded function of a block field, take negative values at functions $H\ge0$: by Proposition~\ref{C-372-15}(c) the conditional mean of $W^{\rm fr}_C$ given the block field of the level $t$ is $\Xi_tW^{{\rm fr},(t)}_C$, and $W^{{\rm fr},(t)}_C$ is a product of three cosines of independent centred Gaussian variables of the variance $v_{K,K-t}(C)/\mathsf b$, which grows like $\frac{|C|}{2\pi\mathsf b}\log\frac1{a_{K-t}}$ (Proposition~\ref{C-372-11}(b)); so $W^{{\rm fr},(t)}_C$ is negative with positive probability, and the quotient at the indicator of $\{W^{{\rm fr},(t)}_C<0\}$ is negative.

Alternatively, the levels $\varrho<a_{K-t}\le a_*$ of the proof of Theorem~\ref{C-372-4}(8)(b), whose cells are larger than the loop, can be traded for a law of one point at the depth of the size of the loop: $1-W^{(t)}_C$ is, on tame fields, a quadratic form in the curvature, and the law follows from a bound on $|Z^{-1}\int(1-W^{(t)}_C)\,dm_t-\frac3{2\mathsf b}v_{K,K-t}(C)|$ at that level, the mean of one bounded function of the block field against its free value.

\section{Centre flux and twisted partition functions}\label{C-366-1}

Centre twists admit a transfer-operator description in terms of electric-flux sectors. We establish the crossing and Fourier relations, the positivity properties of admissible mixtures and the continuum limits that follow from the construction estimates.

A single twisted class need not be reflection positive in every direction. The additional hypotheses for the flux asymptotics are stated in the results. We also distinguish a gap over a fixed small section from a gap of the infinite-volume continuum theory. In the conditional area-law regime the quotient $Z^-/Z$ tends to $1$; the small quantity is its defect.

The continuum centre-flux theorem below also assumes \textbf{(Hess)} and \textbf{(Flux-window)} in the clauses specified there. The former concerns the full fluctuation Hamiltonian on a doubled domain; the latter is a quantitative bound on a twisted window. Neither follows from the Gaussian quadratic estimate alone.

\subsection{Notation and hypotheses}

\emph{The lattice.} Throughout, $G=\mathrm{SU}(2)$, the directions are $\mu\in\{1,2,3\}$, and $\mathrm{tr}$ is the trace of $2\times2$ matrices, not halved. A box is a triple $N=(N_1,N_2,N_3)$ of integers $\ge3$ with its periodic lattice; $|N|=N_1N_2N_3$, a cube has equal sides, and $N^{\times3}=(N,N,N)$. $\partial U(p)$ is the plaquette variable of a plaquette $p$, $t_p=\mathrm{tr}\,\partial U(p)$, and $dU$ is the product of the Haar probability measures on the bonds. A sign field $\gamma$ on the bonds acts by $(\gamma U)(b)=\gamma(b)U(b)$, and $d\gamma(p)$ is the product of $\gamma$ over the four bonds of $p$. For a sign field $\eta$ on the plaquettes, $\beta\ge0$ and a bounded Borel function $F$ of the field,
\begin{equation}\label{C-366-2}
A_\eta(U)=\sum_p\bigl(1-\tfrac12\eta(p)\,\mathrm{tr}\,\partial U(p)\bigr),\qquad Z_\beta[\eta;F]=\int F\,e^{-\beta A_\eta(U)}\,dU,
\end{equation}
and $A=A_\eta$ at $\eta\equiv1$.

\emph{Classes.} A class is $n=(n_{12},n_{13},n_{23})\in\mathbb Z_2^3$; a plane $\mu\nu$ with $n_{\mu\nu}=1$ is a twisted plane of $n$, and $n_{\mu\nu}$ also names the class whose one twisted plane is $\mu\nu$. A column of the plane $\mu\nu$ is the set of the plaquettes of that plane over one point of it; $\eta^n$ is $-1$ exactly on one column of each twisted plane; $Z[n;F]=Z_\beta[\eta^n;F]$, $Z[n]=Z[n;1]$, $r[n]=Z[n]/Z[0]$, $F^n=-\log r[n]$, $\mu_n=Z[n;\cdot\,]/Z[n]$ and $\langle\cdot\rangle=\mu_0$. On a cube, $r_i$ is $r[n]$ for a class with $i$ twisted planes; it depends on $i$ alone. A function of the field is centre-blind if it is invariant under $U\mapsto\gamma U$ for every sign field $\gamma$ on the bonds. A sign field $s$ on the plaquettes is closed if its product over the faces of each unit cube is $1$; its product over a lattice plane parallel to $\mu\nu$ is then $(-1)^{n_{\mu\nu}}$ for one class $n(s)$, and $[n]\ni\eta^n$ is the set of the closed $s$ with $n(s)=n$. A bounded function $G(U,s)$ of a field and a closed sign field, Borel in $U$, is covariant if $G(\gamma U,s\,d\gamma)=G(U,s)$ for all sign fields $\gamma$ on the bonds; for it $Z_\beta[s;G(\cdot,s)]$ is the same number $\mathsf M_n[G]$ for all $s\in[n]$, and $\mu_n(G)=\mathsf M_n[G]/Z[n]$.

\emph{Transfer operators.} With $\mu$ as the time and $\nu<\rho$ the other two directions we write $n=(k,m)_\mu$, with the electric part $k=(n_{\mu\nu},n_{\mu\rho})$ and the magnetic part $m=n_{\nu\rho}$. The section is the periodic lattice $N_\nu\times N_\rho$, and a layer is the field on one section. On $L^2$ of the fields of the section, $T_m=D_m\mathsf KD_m$, where $\mathsf K$ is the crossing operator defined in Section~\YMMref{85}{677} by integrating the time-like bonds between two layers and $D_m^2=e^{-\beta A^m_\Sigma}$ is the Boltzmann weight of the section $\Sigma$, with the plaquette whose corner is a fixed site $\xi_*$ of the section reversed if $m=1$; a section with $m=1$ is called twisted. For a site $\xi$ of the section and $j\in\{\nu,\rho\}$, $(\xi,j)$ is the bond from $\xi$ to $\xi+e_j$, and the plaquette with the corner $\xi$ has the bonds $(\xi,\nu)$, $(\xi+e_\nu,\rho)$, $(\xi+e_\rho,\nu)$, $(\xi,\rho)$. The flip $C_j$ reverses the bonds $(\xi,j)$ with $\xi_j=(\xi_*)_j$, the bonds of direction $j$ that leave the line of the section through $\xi_*$ normal to $j$; $C^k=C_\nu^{k_1}C_\rho^{k_2}$ is composition with a map $c^k$ of the fields; $P_e=\frac14\sum_k(-1)^{e\cdot k}C^k$ for $e\in\mathbb Z_2^2$; and $Z_{e,m}=\mathrm{Tr}(P_eT_m^{N_\mu})$, written $Z^{[\mu]}_{e,m}$ where the time is to be named.

\emph{Numbers of a box.} $F_0=\log Z[0]$, $P=P(\beta)=\inf F_0/|N|$ over all boxes, the excess $\delta=F_0-|N|P$, and $\delta^n=\log Z[n]-|N|P=\log r[n]+\delta$. Over a section and with the height $M$, $\mathsf G(M)=\log(\mathrm{Tr}\,T_0^M/\|T_0\|^M)$ and $\mathsf d(M)=2F_0(M)-F_0(2M)$, the section being fixed. For integers $l\ge3$, $D_q(l)=q^3F_0(l^{\times3})-F_0((ql)^{\times3})$, $D_2(l)=8F_0(l,l,l)-F_0(2l,2l,2l)$ and $\mathsf e(l)=2F_0(l,2l,2l)-F_0(2l,2l,2l)$; over a section $(l_1,l_2)$, $\varepsilon_{\rm v}(l_1,l_2)=\log(\|T_0\|/\|T_1\|)$, and $\varepsilon_{\rm v}(l)=\varepsilon_{\rm v}(l,l)$. For a loop $C$, $W_C=\frac12\mathrm{tr}\,U(C)$; $W_{R\times T}$ is $W_C$ for the boundary of a rectangle $R\times T$ of a coordinate plane.

\emph{Gibbs states.} For a box $N$, $\mu_N$ is the measure $\langle\cdot\rangle=\mu_0$ of the periodic box $N$ (the index of $\mu_N$ is a box, that of $\mu_n$ a class); where it stands beside a measure on the fields of the infinite lattice, it is extended periodically. A Gibbs state at $\beta$ on the infinite lattice $a\mathbb Z^3$, $a>0$, is, as in Section~\YMMref{91}{697}, a Borel probability measure $\nu$ on its fields with $\nu(\gamma_BF)=\nu(F)$ for every finite set $B$ of bonds and every continuous function $F$ of finitely many bonds, $\gamma_B$ the Gibbs kernel of $B$ with the weight $\exp[-\beta\sum_p(1-\frac12\mathrm{tr}\,\partial U(p))]$, the sum over the plaquettes $p$ that hold a bond of $B$ (equation~\YMMref{(1818)}{698}); it is periodic if it is invariant under the translations of a subgroup of finite index of $a\mathbb Z^3$ (Section~\YMMref{91}{697}), and every periodic Gibbs state is invariant under all isometries of the lattice (Proposition~\YMMref{103.4}{758}). In Theorem~\ref{C-366-3}(10), (11) the lattice is $\mathbb Z^3$.

\emph{Mixtures and boards.} A weight $w\colon\mathbb Z_2^3\to[0,\infty)$, $w\ne0$, is admissible if 
\[
\hat w(e,m)=\sum_k(-1)^{e\cdot k}w((k,m)_\mu)\ge0
\]
 for all $\mu,e,m$; then $Z_w=\sum_nw(n)Z[n;\cdot\,]$, $|Z_w|=Z_w[1]$ and $\mu_w=Z_w/|Z_w|$. A board is a box with $N_\mu=\mathsf c\,\mathsf n_\mu$, $\mathsf c\ge3$ odd and $\mathsf n_\mu\ge2$ even, cut into $\mathsf N=\prod_\mu\mathsf n_\mu$ cells by mirrors through bonds, the columns placed where mirrors cross. For a bounded complex Borel function $F$ of one cell, $\mathrm{Fold}\,F$ is the product over the cells of the mirror copies of $F$, the copy in a cell being complex conjugated if the cell is odd relative to the given one, that is, reached from it by an odd number of reflections in the mirrors (equation~\YMMref{(1749)}{680}, \YMMref{(1749)}{680}, \YMMref{(1749)}{680}); for real $F$ the copies are the mirror images, and for complex $F$ the conjugation makes $Z[0;\mathrm{Fold}\,F]\ge0$ (with the plain mirror images the constant $e^{i\pi/8}$ on a board of $8$ cells would have the fold $-1$). 
\[
\|F\|_{\rm ch}^{\mathsf N}=Z[0;\mathrm{Fold}\,F]/Z[0]
\]
 is the periodic chessboard seminorm (Proposition~\YMMref{86.3}{681}), and $\|F\|_w^{\mathsf N}=Z_w[\mathrm{Fold}\,F]/|Z_w|$.

\emph{Cutoffs.} $L$, an odd integer with $L\ge13$, and $a_K=L^{-K}$ are the blocking factor and the spacings; the label is $b>0$ and $\beta=L^Kb$. A box $s$ of the cutoff $K$ has sides $s_\mu\in a_K\mathbb N$, $s_\mu\ge3a_K$; an index $K$ on a number above denotes that number at $N_\mu=s_\mu/a_K$ and $\beta=L^Kb$, as in $Z^z_K(s;b)=Z_K[z](s;b)$, $r_K[z](s;b)$ and $\delta^0_K$; with $p_\infty(K,b)=a_K^{-3}P(L^Kb)$ we put $\delta^n_K(s;b)=\log Z_K[n](s;b)-|s|\,p_\infty(K,b)$. A length is $L$-adic if it lies in $\mathbb D=\bigcup_Ka_K\mathbb N$. For $K,m\in\mathbb N$ the cube $\ell=2L^m$ is the lattice $\Lambda_{K,\ell}$ of the torus $T_\ell$ (Section~\YMMref{3}{13}); for other $m$ it is reached by the dilation $Z_K[n](s;b)=Z_{K-j}[n](L^js;L^jb)$. $Z_{K,\ell}$ is its periodic partition function, $\Delta_K[\ell,b]=\log Z_{K+1,\ell}-\log Z_{K,\ell}$, and on cubes $t=\ell/b$.

A box $s$ is mixed for the class $z$ if all $s_\mu\in L^{\mathbb Z}\cup2L^{\mathbb Z}$ and each twisted plane of $z$ has a side in $L^{\mathbb Z}$. We write $\rho^z(N,\beta)$ for $r[z]$ on the box $N$ at $\beta$, and, for a fixed class $z$, $\phi_j(t)=\rho^z((2L^j)^{\times3},2L^j/t)$. The constants $\gamma$, $M=L^{m'}$, $r_\star$ and $\mathfrak p_0$ are those of Proposition~\YMMref{69.1}{595}, with $r_\star=r$ and $2\mathfrak p_0>3r_\star+2$. Put $R(b)=\min\{L^a:a\in\mathbb N_0,\ L^a\ge(\log b)^{r_\star}\}$, $\ell_\flat(b)=2MR(b)$ and $\ell_\sharp(b)=L\ell_\flat(b)$, and set
\[
 m_0(b)=1+m'+\max\{\log_LR(b+\bar B),\lceil\log_LC_{\rm s}\rceil\},
 \qquad b\ge\gamma^{-2}.
\]
Here $C_{\rm s}$ is the torus-floor constant, chosen at least $C_L$ as in the paragraph following Proposition~\YMMref{74.5}{617}. With $\sigma$ and $\mathfrak b$ of Theorem~\ref{C-366-3}(3) at $\beta=L^Kb$, $\sigma_K(l)=a_K^{-2}\sigma(l/a_K)$ and $\mathfrak b_K(t,\lambda;l)$ is $\mathfrak b$ on the box $(t,\lambda,l)/a_K$.

\emph{Levels and block fields.} Over a section $l=(l_1,l_2)$ and with a time $\mu$, the levels are $E^{K,m,e}_i=a_K^{-1}\log(\|T_0\|/\lambda_i)$, $\lambda_1\ge\lambda_2\ge\dots>0$ the eigenvalues of $P_eT_m$, each repeated according to its multiplicity, and $E^{K,m,e}_i=\infty$ beyond them. $\mathcal M$ is the averaging map; on the cube $\ell$ at the cutoff $K$ the block field of depth $d\le K$ is $\mathcal M^{K-d}U$, a field on the block lattice of spacing $a_d$, and $\mu^K_d$ is its law under the periodic measure. A block line is a line of a block lattice.

\subsection{Principal results}

The untwisted partition function supplies the reference for the
centre-flux ratios. We record the stability estimate at a fixed
physical label, including the room needed to compare two tori.

\begin{lemma}[Stability at a fixed label]\label{lem:pure-scaling-stability}
Use the common construction constants, with $C_{\rm s}\ge C_L$
as in the torus-floor choice following Proposition~\YMMref{74.5}{617},
and let $b\ge\gamma^{-2}$, $K\ge1$ and $m_i\ge m_0(b)$ for
$i=1,2$. Put $\ell_i=2L^{m_i}$ and $\theta_j=L^{K-j}b$.
The vacuum constants $E_i=E_{(K)}[((\ell_i,K),\theta_\bullet)]$
of the real progression are real, and
\begin{equation}\label{eq:pure-scaling-stability}
 \begin{gathered}
 e^{-[c_{\rm low}+\frac92\log b]\ell_i^3}
 \le e^{-E_i}Z_{K,\ell_i,L^Kb}
 \le e^{(C_{\rm maj}+2)\ell_i^3},\\
 \left|E_1\ell_1^{-3}-E_2\ell_2^{-3}\right|
 \le2C_{\rm tt}C'_{\rm tt}E_0
       e^{-\kappa\min\{\ell_1,\ell_2\}/(5M)}.
 \end{gathered}
\end{equation}
Here $C_{\rm tt}=4096\sum_{q\ge0}e^{-q}$ and $C'_{\rm tt}=2$;
the rate and size constants are those of Proposition~\YMMref{74.5}{617}.
In particular, the first logarithmic error is bounded by
$(C_{\rm lab}+\frac92\log b)\ell_i^3$, where
$C_{\rm lab}=\max\{c_{\rm low},C_{\rm maj}+2\}$.
\end{lemma}

\begin{proof}
The increments of the sequence are
$L^{-1}\theta_{j-1}-\theta_j=0$, and
$\theta_j\ge b\ge\gamma^{-2}$. Thus it is an admissible
prescribed sequence in equation~\YMMref{(1524)}{614}, with last
radius $R(b)$. The definition of $m_0$ gives
\[
 L^{m_i}\ge LM R(b+\bar B)\ge LM R(b),\qquad
 \ell_i=2L^{m_i}\ge2LM C_{\rm s}\ge C_LM.
\]
Both tori are therefore admitted, including the stronger room
required for their comparison. Apply Proposition~\YMMref{74.4}{617}
at $k=K$, with $u_0=L^Kb$ and $u_K=b$, and
Proposition~\YMMref{74.5}{617} to this same sequence on the two
tori. These give the two bounds and the reality of the vacuum
constants. The vacuum terms are those of the canonical progression
of the prescribed sequence; its bare counterterms are its zero
increments. Finally $b\ge\gamma^{-2}\ge4$ implies $\log b>0$,
so the definition of $C_{\rm lab}$ gives the symmetric logarithmic
bound.
\end{proof}

\begin{theorem}[One lattice]\label{C-366-3}
Let $N$ be a box, $\beta\ge0$, and $\mu$ a direction, taken as the time, $\nu<\rho$ the other two.
\begin{enumerate}
\item \emph{(Transfer operator; order; momentum; crossing.)} $T_m\ge0$ is Hilbert--Schmidt with a continuous symmetric kernel $>0$; $C_\nu,C_\rho$ are commuting unitary involutions that commute with $T_m$. For every class $Z[(k,m)_\mu]=\mathrm{Tr}(T_m^{N_\mu}C^k)$, and $Z_{e,m}=\frac14\sum_k(-1)^{e\cdot k}Z[(k,m)_\mu]=\sum_i\lambda_i(P_eT_m)^{N_\mu}\ge0$ with eigenvalues $\lambda_i(P_eT_m)$ that do not depend on $N_\mu$; $Z_{e,m}<Z_{00,m}$ for $e\ne0$, and $Z_{e,m}>0$ if $\beta>0$. $\|T_m\|$ is a simple eigenvalue with an eigenfunction $>0$ in the range of $P_{00}$, and $\|T_1\|\le\|T_0\|$. Moreover
\[ e^{-2\beta(n_{23}N_1+n_{13}N_2+n_{12}N_3)}Z[0]\le Z[n]\le Z[(0,m)_\mu]\le Z[0], \]
the last two inequalities strict if $\beta>0$ and the classes differ; so $r[n]\le r[n']$ for every one-plane class $n'$ whose plane is twisted in $n$. Over a twisted section, with the reversed plaquette at the corner $\xi_*$ and the lines of $C_\nu$, $C_\rho$ through $\xi_*$, as in the setting, let $\mathcal T_\nu$ be the composition with the map $V\mapsto\mathsf S_\nu(\gamma_\nu V)$ of the fields, where $(\mathsf S_\nu V)_{(\xi,j)}=V_{(\xi+e_\nu,j)}$ is the step by one site in the direction $\nu$ and $\gamma_\nu$ reverses the one bond $(\xi_*+e_\nu,\rho)$, which the reversed plaquette shares with its neighbour in the direction $\nu$; and let $\mathcal T_\rho$ be formed likewise, with $\nu$ and $\rho$ exchanged and the bond $(\xi_*+e_\rho,\nu)$.

On all of $L^2$, $\mathcal T_\nu$ and $\mathcal T_\rho$ are unitaries that commute with $T_1$ and map the gauge invariant functions onto themselves, and $\mathcal T_\nu^{N_\nu}=C_\rho$, $\mathcal T_\rho^{N_\rho}=C_\nu$; so $\mathcal T_\nu$ commutes with $C_\rho$, and $\mathcal T_\rho$ with $C_\nu$. The products $\mathcal T_\nu\mathcal T_\rho$ and $\mathcal T_\rho\mathcal T_\nu$ differ by the reversal of the four bonds at the site $\xi_*+e_\nu+e_\rho$, and $\mathcal T_\nu C_\nu$ and $C_\nu\mathcal T_\nu$, as $\mathcal T_\rho C_\rho$ and $C_\rho\mathcal T_\rho$, by the gauge transformation that is $-1$ at the sites of one line of the section and $1$ elsewhere: gauge transformations with the values $\pm1$, which act as the identity on the gauge invariant functions.

On the gauge invariant functions, off which $T_1$ vanishes ($T_1=\Pi T_1\Pi$, $\Pi$ the gauge average) and which for $\beta>0$ form the closure of the range of $T_1$, the five operators $T_1$, $\mathcal T_\nu$, $\mathcal T_\rho$, $C_\nu$, $C_\rho$ commute pairwise. So the range of $P_e$, $e=(e_\nu,e_\rho)$, in the gauge invariant functions, on which $C_\nu=(-1)^{e_\nu}$ and $C_\rho=(-1)^{e_\rho}$, has an orthonormal basis of joint eigenvectors of $T_1$, $\mathcal T_\nu$ and $\mathcal T_\rho$, and the joint eigenvalues $(e^{i\theta_\nu},e^{i\theta_\rho})$ of $(\mathcal T_\nu,\mathcal T_\rho)$ there have $\theta_\nu\in\frac{2\pi}{N_\nu}(\mathbb Z+\frac{e_\rho}2)$ and $\theta_\rho\in\frac{2\pi}{N_\rho}(\mathbb Z+\frac{e_\nu}2)$; and $Z[(k,1)_\mu]=\mathrm{Tr}(T_1^{N_\mu}\mathcal T_\nu^{k_2N_\nu}\mathcal T_\rho^{k_1N_\rho})$.

For the times $1$ and $2$, with $e=(e_1,e_2)$ dual to $(n_{12},n_{13})$ and $m=n_{23}$ for the time $1$, and $(e_1,e_2')$ dual to $(n_{12},n_{23})$ and $m'=n_{13}$ for the time $2$, the crossing relation
\[ Z^{[1]}_{e_1e_2,m}=\tfrac12\sum_{m',e_2'\in\{0,1\}}(-1)^{e_2m'+e_2'm}\,Z^{[2]}_{e_1e_2',m'} \]
holds, and likewise for every pair of times: the flux of the common plane is kept, and the pairs $(e_2,m)$ and $(m',e_2')$ are exchanged by the Fourier transform of $\mathbb Z_2^2$. Hence, for every time, $Z_{e_1e_2,m}\le\min\{Z_{e_10,0},Z_{0e_2,0}\}$, $\|T_mP_{e_1e_2}\|\le\min\{\|T_0P_{e_10}\|,\|T_0P_{0e_2}\|\}$ and $\|\frac12(1-C_\nu)T_0\|=\|T_0P_{10}\|$: the least level of odd electric flux along $\nu$ has no flux along $\rho$. For $\beta>0$ and a fixed section, $\varepsilon_{\rm mag}=\log(\|T_0\|/\|T_1\|)=-\lim_{N_\mu\to\infty}N_\mu^{-1}\log r[n_{\nu\rho}]\in[0,\infty)$ and $\varepsilon_{\rm el}=\log(\|T_0\|/\|T_0P_{10}\|)=-\lim_{N_\mu\to\infty}N_\mu^{-1}\log\frac{Z[0]-Z[n_{\mu\nu}]}{Z[0]+Z[n_{\mu\nu}]}\in(0,\infty)$.
\item \emph{(Two lengths; the excess.)} Each of $\frac12(Z[0]\pm Z[n_{12}])$ is a sum of $N_1$-th powers of numbers $\ge0$ that do not depend on $N_1$, and a sum of $N_2$-th powers of numbers $\ge0$ that do not depend on $N_2$. $F_0/|N|$ does not increase in any side and tends to $P\ge-6\beta$ as all sides tend to infinity. Over every section and for every time, $0\le\mathsf G\le\delta$, $\mathsf G$ does not increase in the height, and $\mathsf G(M)\le\mathsf d(M)=2\delta(M)-\delta(2M)=2\mathsf G(M)-\mathsf G(2M)\le2\mathsf G(M)$. For every integer $q\ge2$, $qF_0(N)-F_0(N^{[\mu]})=q\delta(N)-\delta(N^{[\mu]})\ge0$, $N^{[\mu]}$ the box with $N_\mu$ multiplied by $q$. And $\delta(N')/|N'|\le\min\{1,\sum_\mu N_\mu/N'_\mu\}\,\delta(N)/|N|$ if $N'_\mu\ge N_\mu$ for all $\mu$.
\item \emph{(Electric flux tube.)} Let $\beta>0$ and $l\ge3$. Over the sections of sides $\lambda$ and $l$ let $E_0=-\log\|T_0\|$ and $E_{\rm e}=-\log\|\frac12(1-C')T_0\|=-\log\|T_0P_{10}\|$, $C'$ the flip of the bonds parallel to the side $\lambda$, and $\varepsilon=E_{\rm e}-E_0$; then $0\le E_0\le E_{\rm e}<\infty$. For $f\in\{E_0,E_{\rm e}\}$ and $\lambda,\lambda'\ge3$: $2f(\lambda)\ge f(\lambda-j)+f(\lambda+j)$ for $\lambda-j\ge3$, $f(\lambda+\lambda')\ge f(\lambda)+f(\lambda')$, $f(\lambda)/\lambda$ does not decrease and tends to a real number $s_f$, $s_f\le[f(2\lambda)-f(\lambda)]/\lambda$, and $s_f\lambda-f(\lambda)\ge0$ is convex and does not increase. With $s_0=s_{E_0}$ and $R(\lambda)=s_0\lambda-E_0(\lambda)$, the tension $\sigma(l)=s_{E_{\rm e}}-s_0=\lim\varepsilon(\lambda)/\lambda\in[0,\infty)$ obeys
\[ \frac{\varepsilon(\lambda)-R(\lambda)}\lambda\le\sigma(l)\le\frac{\varepsilon(2\lambda)-\varepsilon(\lambda)+E_0(2\lambda)-2E_0(\lambda)}\lambda\qquad(\lambda\ge3). \]
On the boxes $(t,\lambda,l)$ let $r=r[n_{12}]$, $q=\frac12(Z[0]-Z[n_{12}])$, $p=\frac12(Z[0]+Z[n_{12}])$, $\Phi=F_0+t\lambda s_0$ and $\mathfrak b=\frac12(1-r)e^{\Phi}=q\,e^{t\lambda s_0}$. Then $-s_0=\inf_{t',\lambda'\ge3}F_0(t',\lambda',l)/t'\lambda'\ge lP$, $0\le\Phi\le\delta$, $\Phi/t\lambda$ does not increase in $t$ nor in $\lambda$, $p\ge e^{-t\lambda s_0}$ and so $\mathfrak b\le e^\Phi-1$; and for $t\ge t_0\ge3$, $\lambda\ge\lambda_0\ge3$
\[ \frac{1-r}{1+r}=\frac qp\le\mathfrak b(t,\lambda)\le\mathfrak b(t_0,\lambda_0)^{t\lambda/t_0\lambda_0}\le\bigl(e^{\Phi(t_0,\lambda_0)}-1\bigr)^{t\lambda/t_0\lambda_0}, \]
so that $\Phi(t_0,\lambda_0)<\log2$, a number of periodic boxes, is enough for an area law, and $r(t,\lambda,l)\ge(1-x)/(1+x)$ with $x=\mathfrak b(t_0,\lambda_0)^{t\lambda/t_0\lambda_0}$. Further $-(t\lambda)^{-1}\log\mathfrak b(t,\lambda)$ does not decrease in $t$ nor in $\lambda$, and
\[ \sigma(l)=\sup_{t,\lambda}\frac{\log(1/\mathfrak b(t,\lambda))}{t\lambda}=\lim_{t,\lambda\to\infty}\frac{\log(1/\mathfrak b(t,\lambda))}{t\lambda}; \]
so $\sigma(l)>0$ if and only if $\mathfrak b<1$ on one box. For $t>t_0$, $\lambda>\lambda_0$: $-\log\mathfrak b(t,\lambda)\ge(t-t_0)[(\lambda-\lambda_0)\sigma(l)-R(\lambda_0)]-(\lambda/\lambda_0)\Phi(t_0,\lambda_0)$. Finally $\varrho=q/Z[0]$ has $\varrho(2t)\le\varrho(t)^2e^{2F_0(t)-F_0(2t)}$, and the same in $\lambda$.
\item \emph{(Ceilings of the tension; Casimir terms; slabs.)} Let $\beta>0$ and $l\ge3$, with the signs of (3). (a) For integer heights $3\le h_0<h_1$ and every length $\lambda\ge3$,
\[ (h_1-h_0)\,\varepsilon(\lambda)\le\log\bigl(e^{\mathsf G(h_0;\lambda,l)}-1\bigr)+\log\frac2{1-r(h_1,\lambda,l)}, \]
$\mathsf G(h_0;\lambda,l)$ being $\mathsf G$ at the height $h_0$ over the section $(\lambda,l)$; and for $3\le\lambda_0<\lambda_1$, with $\kappa'=\log(2/(1-r(h_1,\lambda_1,l)))$,
\[ \sigma(l)\le\frac1{h_1-h_0}\Bigl\{\frac{\Phi(h_0,\lambda_0)}{\lambda_0}+\frac{\kappa'+\Phi(h_1,\lambda_0)}{\lambda_1-\lambda_0}\Bigr\}+\frac{\Phi(h_1,\lambda_0)}{\lambda_0h_1}, \]
and the same with $\delta(\cdot,\cdot,l)$ in the place of $\Phi$. (b) Over a transverse side $b'\ge3$, for all integers $a\ge3$, $k'\ge2$: $R(a)\le\mathsf G(a;b',k'a)/((k'-1)a)$, $R$ formed over the sections $(a,b')$; and $\delta(M,M,M)\le\mathsf G(M;M,M)+\mathsf G(M;M,2M)+\frac12\mathsf G(M;2M,2M)+\frac14\mathsf G(2M;2M,4M)$ for every $M\ge3$. (c) For all integers $t,\lambda_1,l\ge3$,
\[ \sigma(l)\ge\frac1{\lambda_1}\Bigl\{\frac1t\Bigl[\log\frac2{1-r(t,\lambda_1,l)}-\log\frac{Z[0](t,\lambda_1,l)^2}{Z[0](2t,\lambda_1,l)}\Bigr]-\frac1{\lambda_1}\log\frac{Z[0](\lambda_1,2\lambda_1,l)^2}{Z[0](2\lambda_1,2\lambda_1,l)}\Bigr\}: \]
a positive tension is one inequality between five partition functions. (d) On the slab $\{0,\dots,A'\}\times\{0,\dots,B'\}\times\mathbb Z_{N'}$, free in the first two directions and periodic in the third, $N'\ge3$, with the Wilson action on all its plaquettes, $|\langle\mathrm{tr}\,U(C)\rangle|\le2(\tanh\beta N')^{RT}$ for the boundary $C$ of every rectangle of $R\times T$ plaquettes in a plane normal to the third direction, at every $\beta\ge0$. (e) On a box $(t,\lambda,l)$ let the $t\lambda$ plaquettes of the periodic $t\times\lambda$ lattice of the twisted plane be parted into sets $D$, each connected through common edges; call a bond inner to $D$ if it is parallel to the twisted plane and lies over an edge both of whose plaquettes lie in $D$, or normal to it and lies over a site all four of whose plaquettes lie in $D$. Then $q/p=\int\prod_D\rho_D\,d\mu$, where $\mu$ is a probability measure on the configurations of the bonds that are inner to no $D$, and $\rho_D=(Z_D^+-Z_D^-)/(Z_D^++Z_D^-)\in(-1,1)$, $Z_D^\pm$ the integrals over the bonds inner to $D$ of the weight of the plaquettes parallel to the twisted plane over $D$ and of the plaquettes normal to it over the inner edges of $D$, without and with one reversed column over $D$; $Z_D^-$ does not depend on the place of that column.
\item \emph{(How floors and ceilings travel.)} The time $\mu$ and the section are fixed, the boxes differ in $N_\mu$ alone, and all numbers of layers are $\ge3$. (a) For $k\ne0$ and $m\in\{0,1\}$, 
\[
Z[(k,m)_\mu](N_\mu)^2\ge2Z[(0,m)_\mu](2N_\mu)-Z[(0,m)_\mu](N_\mu)^2;
\]
so with $\mathcal R=Z[0](N_\mu)^2/Z[0](2N_\mu)$, $r[(k,0)_\mu]^2\ge2/\mathcal R-1$, and $1/\mathcal R\le\frac14\sum_kr[(k,0)_\mu]^2$. With $Z'=Z[(0,m)_\mu]$ and $q_0=Z_{00,m}/Z'\ge\|T_m\|^{N_\mu}/\mathrm{Tr}\,T_m^{N_\mu}$: $Z[(k,m)_\mu]/Z'\ge2q_0-1$ for every $k$; $\sum_{k\ne0}Z[(k,m)_\mu]/Z'=4q_0-1$, so that $\sum_{k\ne0}r[(k,0)_\mu]\ge4e^{-\mathsf G(N_\mu)}-1\ge4e^{-\delta}-1$; and 
\[
Z[((1,1),m)_\mu]\ge Z[((1,0),m)_\mu]+Z[((0,1),m)_\mu]-Z'.
\]
(b) \emph{A ceiling moves inside a twisted plane.} Let $n=(k,0)_\mu$, $k\ne0$, and $r(N_\mu)=r[n]$. For $N_\mu-j\ge3$,
\[ (1-r(N_\mu-j))(1-r(N_\mu+j))\ge(1-r(N_\mu))^2e^{2\delta(N_\mu)-\delta(N_\mu-j)-\delta(N_\mu+j)}. \]
(c) \emph{A floor moves along its own time.} For $\emptyset\ne S\subset\mathbb Z_2^2$, $P_S=\sum_{e\in S}P_e$, $f(N')=\mathrm{Tr}(P_ST_m^{N'})/\mathrm{Tr}\,T_0^{N'}$, $\lambda=\|P_ST_m\|$ and $\Theta=e^{\mathsf G}$: 
\[
(\lambda/\|T_0\|)^{N'}/\Theta(N')\le f(N')\le(\lambda/\|T_0\|)^{N'-N''}\Theta(N'')
\]
 for $3\le N''<N'$; so for $3\le M_0<M_1$ and $M_2\ge3$, 
\[
f(M_2)\ge\Theta(M_2)^{-1}[f(M_1)/\Theta(M_0)]^{M_2/(M_1-M_0)}.
\]
In particular, if $r[n_{\nu\rho}](M_1)\ge\varrho>0$, then $\varepsilon_{\rm v}\le[\mathsf G(M_0)+\log(1/\varrho)]/(M_1-M_0)$ and $r[n_{\nu\rho}](M)\ge e^{-\mathsf G(M)}(\varrho e^{-\mathsf G(M_0)})^{M/(M_1-M_0)}$ for every $M\ge3$. (d) \emph{Floors move inside a twisted plane.} With $n$ and $r$ as in (b): for $3\le N_1<N_\mu<N_2$ and $\vartheta=(N_2-N_\mu)/(N_2-N_1)$,
\[ 1-r(N_\mu)\le(1-r(N_1))^\vartheta(1-r(N_2))^{1-\vartheta}e^{\vartheta\delta(N_1)+(1-\vartheta)\delta(N_2)-\delta(N_\mu)}; \]
for $N_\mu\ge N_0\ge3$, 
\[
\frac12(1-r(N_\mu))\le[\frac12(1-r(N_0))e^{\mathsf G(N_0)}]^{N_\mu/N_0};
\]
and for every class $z$ one of whose twisted planes holds $\mu$, all integers $n'\ge1$, $m_1,m_2\ge2$, $m=m_1+m_2$, and every odd $k'\ge1$,
\[ F^z(k'n'+m)\le k'F^z(n'+m)+\tfrac{k'-1}2[\delta(2m_1)+\delta(2m_2)]+\delta(k'n'+m)-k'\delta(n'+m), \]
the argument being the number of layers in the time $\mu$. (e) \emph{Cubes.} For odd $L'\ge3$, $N'=2L'^j$ with $j\ge2$, every class $z\ne0$ and every $\beta$, on the cubes,
\[ F^z(L'N')\le(2L'+1)^3F^z(N')+C_{L'}\,\delta(N'/L'),\qquad C_{L'}=34L'^5+53L'^4+27L'^3+4L'^2, \]
$\delta(N'/L')$ the excess of the cube of side $N'/L'$.
\item \emph{(The overlap; odd functions.)} Let $N_\mu=2\tau$, $k\ne0$, $\Psi$ the kernel of $T_m^\tau$, $\mu_2=\Psi^2dV\,dW/Z[(0,m)_\mu]$ the law of two layers half a period apart in the class $(0,m)_\mu$, $\langle\cdot\rangle_2$ the mean under $\mu_2$, $g(V,W)=(V,c^kW)$, $r'=Z[(k,m)_\mu]/Z[(0,m)_\mu]$, $\mathsf m=(\Psi^2-\Psi^2\circ g)/(\Psi^2+\Psi^2\circ g)$, $p_1$ the law of one layer and $E_c=\{\Psi\circ g\ge c\Psi\}$. Then $r'=\langle\Psi\circ g/\Psi\rangle_2=\int\sqrt{d\mu_2\,d(g_*\mu_2)}=\langle(1-\mathsf m^2)^{1/2}\rangle_2$; $1-\mathrm{TV}\le r'\le(1-\mathrm{TV}^2)^{1/2}$, $\mathrm{TV}$ the total variation distance of $\mu_2$ and $g_*\mu_2$; $c\,\mu_2(E_c)\le r'\le c+\mu_2(E_c)^{1/2}$ for every $c>0$; $r'\le2\mu_2(E)^{1/2}$ for every Borel $E$ with $gE=E^c$; $1-r'\le\sup\langle\Phi\rangle_2^2/\langle\Phi^2\rangle_2=\langle\mathsf m^2\rangle_2\le1-r'^2$, the supremum over the bounded Borel $\Phi$ with $\Phi\circ g=-\Phi$ and $\langle\Phi^2\rangle_2>0$; $r'\ge1-2\,\mathrm{TV}(\mu_2,p_1\otimes p_1)$; if $\mu_2\ge c_0\,p_1\otimes p_1$ on a Borel set $\mathcal G$, then $r'\ge c_0\,(p_1\otimes p_1)(\mathcal G\cap g\mathcal G)$; $-\log r'\le\frac12\mathrm{KL}(\mu_2|g_*\mu_2)=\langle\log(\Psi/\Psi\circ g)\rangle_2<\infty$; and for every bimeasurable bijection $\Phi$ of the fields that keeps $dU$ and the layer $0$ and maps the layer $\tau$ by $c^k$, with $\mu'$ the law of the class $(0,m)_\mu$ and $A'=A_{\eta^{(0,m)_\mu}}$,
\[ r'\ge\int\sqrt{d\mu'\,d(\Phi_*\mu')}\ge\exp\bigl[-\tfrac\beta2\mu'(A'\circ\Phi^{-1}-A')\bigr],\qquad\tfrac12\mathrm{KL}(\mu_2|g_*\mu_2)\le\tfrac\beta2\mu'(A'\circ\Phi^{-1}-A'). \]
Odd functions: let $n=(k,0)_\mu$ have one twisted plane, $N_\mu=2k'\vartheta$ with integers $\vartheta\ge2$, $k'\ge2$, let $D,D'$ be bounded complex Borel functions of the bonds inside one slab of $\vartheta$ layers, $D'$ odd under $c^k$ applied to all layers of the slab, and $\mathsf t_jD'$ the translate of $D'$ by $j\vartheta$ layers, $1\le j\le2k'-1$. Then
\[ |\langle D\,\mathsf t_jD'\rangle|\le\sup|D|\,\sup|D'|\sum_{\varsigma=\pm1}\Bigl(\frac{1+\varsigma r[n]}2\Bigr)^{\frac{2k'-j-1}{2k'}}\Bigl(\frac{1-\varsigma r[n]}2\Bigr)^{\frac{j-1}{2k'}}, \]
and at $j=k'$, $|\langle D\,\mathsf t_{k'}D'\rangle|\le2\sup|D|\,\sup|D'|\,[\frac14(1-r[n]^2)]^{\frac12-\frac1{2k'}}$.
\item \emph{(Reflection positivity.)} Let $N_\mu\ge4$. On centre-blind functions $Z[n;\cdot\,]$ does not depend on the places of the columns, and for $\beta>0$ it is reflection positive at the mirrors normal to $\mu$ if and only if $\mu$ lies in no twisted plane of $n$; when $\mu$ lies in a twisted plane, the form is negative at a centre-blind gauge invariant function of one slab, a product of even class functions of the plaquettes through $\mu$ in one crossing. For centre-blind $F$ of a slab and each $(e,m)$, $\sum_k(-1)^{e\cdot k}Z[(k,m)_\mu;\overline{\theta F}F]\ge0$. For each of the eight sets $\mathfrak s$ of coordinate planes, $\sum\{Z[n;\cdot\,]:n_{\mu\nu}=0\text{ for }\mu\nu\notin\mathfrak s\}$ is reflection positive on centre-blind functions at the mirrors normal to every direction $\lambda$ with $N_\lambda\ge4$; of the single classes, $n=0$ alone is so in all three directions; and $Z[0]\le\sum_nZ[n]\le8Z[0]$.

If $N_\mu$ is even, $2\le\vartheta\le N_\mu/2$, $F$ is a bounded Borel function of a slab of $\vartheta$ layers, and the columns of the planes through $\mu$ lie in a crossing outside the slab, then $|Z[(k,m)_\mu;F]|^2\le Z[(0,m)_\mu]\,Z[(0,m)_\mu;\overline{\theta F}F]$, $\theta$ the reflection that leaves $\frac12N_\mu-\vartheta+1$ crossings on each side, and $Z[(k,m)_\mu;\mathbf 1_E]\le Z[(0,m)_\mu]\,\mu_{(0,m)_\mu}(E)^{1/2}$ for every event $E$ of the slab: what is rare without electric flux is rare with it.
\item \emph{(Chessboard bounds.)} On a board, for bounded complex Borel functions $F_\kappa$ of the cells and every class, $|Z[n;\prod_\kappa F_\kappa]|\le Z[0]\prod_\kappa\|F_\kappa\|_{\rm ch}$, and the same with $\eta^n$ replaced by any sign field carried by the plaquettes in the mirrors. For admissible $w$: $Z_w$ is reflection positive at the mirrors of the three directions; $w(0)Z[0]\le|Z_w|=Z[0]\sum_nw(n)r[n]\le8w(0)Z[0]$; $|Z_w[\prod_\kappa F_\kappa]|\le|Z_w|\prod_\kappa\|F_\kappa\|_w$; $\|F\|_w^{\mathsf N}\le\frac{\sum_nw(n)}{\sum_nw(n)r[n]}\|F\|_{\rm ch}^{\mathsf N}\le8\|F\|_{\rm ch}^{\mathsf N}$, and for $F\ge0$, $\|F\|_w^{\mathsf N}\ge\frac{w(0)}{\sum_nw(n)r[n]}\|F\|_{\rm ch}^{\mathsf N}\ge\frac18\|F\|_{\rm ch}^{\mathsf N}$; so $|\mu_w(\prod_\kappa F_\kappa)|\le8\prod_\kappa\|F_\kappa\|_{\rm ch}$.

All this holds for covariant functions $G_\kappa(U,s)$ of the cells, with $F_\kappa=G_\kappa(\cdot,1)$, every $s\in[n]$ and the columns anywhere. The constants are attained: the functions of the cells that cancel the weight of rods of plaquettes through the centres of the cells give equality in the first bound and in the bound of $\|\cdot\|_w$ at $w\equiv1$; and with the fold taken in a class $n\ne0$ itself the estimate fails for $\beta>0$: there are positive functions of the cells with $Z[n;\prod_\kappa F_\kappa]=Z[0]$ and $\prod_\kappa Z[n;\mathrm{Fold}\,F_\kappa]^{1/\mathsf N}=Z[n]<Z[0]$. In a class: for real $q\ge1$, $\mu_n(\prod_\kappa|G_\kappa|)\le r[n]^{-1/q}\prod_\kappa\||F_\kappa|^q\|_{\rm ch}^{1/q}$; for a covariant event $E$ of one cell, $\mu_n(E)\le r[n]^{-1}\|\mathbf 1_E\|_{\rm ch}$; and the law of the inner bonds of one cell under $\mu_n$ has a density $h_n$ against their law under $\langle\cdot\rangle$ with $\int h_n\log h_n\le\mathsf N\log(1/r[n])$, so that the two laws differ in total variation by at most $(\frac{\mathsf N}2\log\frac1{r[n]})^{1/2}$.

Without a board: let $N_\mu=2h$, $h\ge2$, and let $G$ be covariant and depend only on the bonds and signs of the $h$ layers on one side of a mirror normal to $\mu$, $F=G(\cdot,1)$. If $\mu$ lies in every twisted plane of $n$, then $r[n]|\mu_n(G)|\le\langle\overline{\theta F}F\rangle^{1/2}\le\langle|F|^2\rangle^{1/2}$, the first inequality attained, and
\[ |\mu_n(G)-\langle F\rangle|\le(r[n]^{-2}-1)^{1/2}(\langle|F|^2\rangle-|\langle F\rangle|^2)^{1/2}, \]
which is at most $\frac12(r[n]^{-2}-1)^{1/2}$ for events; the condition on $\mu$ cannot be dropped, since an event of a half torus may be determined by the flux through a plane inside it. If besides $N_\nu=2h'$, $h'\ge2$, and $G$ depends only on a quarter of the torus cut by mirrors normal to $\mu$ and to $\nu$, then $r[n]|\mu_n(G)|\le\langle|F|^4\rangle^{1/4}$ for every $n$, and for every $w\ge0$, $w\ne0$, and every covariant event $E$ of the quarter, $\mathsf M_w[E]/|Z_w|\le\frac{\sum_nw(n)}{\sum_nw(n)r[n]}\langle\mathbf 1_E\rangle^{1/4}$, where $\mathsf M_w=\sum_nw(n)\mathsf M_n$. Pairs: if $N_\mu=2h$, $h\ge2$, $w\ge0$ has $\hat w\ge0$ at the time $\mu$, and $\mathbb M_w[G]=\sum_sw(n(s))Z_\beta[s;G(\cdot,s)]$ over the closed $s$, then $\mathbb M_w[G\,\overline{\theta G}]\ge0$ for every bounded $G(U,s)$, Borel in $U$ and not supposed covariant, that depends only on the bonds and plaquette signs of the $h$ layers on one side of a plane midway between two layers normal to $\mu$, $\theta$ the reflection in it.

The infinite lattice: let $\beta>0$, $K,m\in\mathbb N$, $\ell=2L^m$, let $\mathrm{pr}$ be the projection of $a_K\mathbb Z^3$ onto $\Lambda_{K,\ell}$, on sites, bonds and plaquettes, and $\mathrm{per}_\ell U=U\circ\mathrm{pr}$, on the bonds, the periodic extension. For a class $n$ and $s\in[n]$ the sign field $s\circ\mathrm{pr}$ on the plaquettes of $a_K\mathbb Z^3$ is closed; closed sign fields on the plaquettes of $a_K\mathbb Z^3$ being exact, $s\circ\mathrm{pr}=d\chi$ for a sign field $\chi$ on the bonds of $a_K\mathbb Z^3$, unique up to a factor $d\epsilon$, $\epsilon$ a sign field on the sites and $d\epsilon$ at a bond the product of $\epsilon$ at its two ends; and the plaquette variables of the twisted extension $\chi\,\mathrm{per}_\ell U$ are $\partial(\chi\,\mathrm{per}_\ell U)(p)=s_{\mathrm{pr}\,p}\,\partial U(\mathrm{pr}\,p)$.

The law $\tilde\mu_n$ of the class $n$ on $a_K\mathbb Z^3$ is the image of $e^{-\beta A_s}dU/Z[n]$ under $U\mapsto\chi\,\mathrm{per}_\ell U$; on the functions invariant under $U\mapsto(d\epsilon)U$ for all $\epsilon$ it depends neither on $s\in[n]$ nor on $\chi$ (a change $s\mapsto s\,d\gamma$ is undone by $U\mapsto\gamma U$, with $\chi(\gamma\circ\mathrm{pr})$ for $\chi$, and $\chi$ is fixed up to $d\epsilon$), and $\tilde\mu_0$ is the periodic measure extended periodically. For admissible $w$ let $\tilde\mu_w=\sum_nw(n)r[n]\,\tilde\mu_n/\sum_nw(n)r[n]$; for a board of $a_K\mathbb Z^3$ with cells of side $\mathsf l\le\ell/2$ and a function $F$ of a cell let $\|F\|_{w,\ell}=\|F\circ\mathrm{per}_\ell\|_w$, on the board of $\Lambda_{K,\ell}$ under it, and let $\|F\|_{{\rm ch},\ell}$, the seminorm of Corollary~\ref{IV-18-18}, be the same number at $w=\delta_0$.

Then on the local functions invariant under all $U\mapsto(d\epsilon)U$ the law $\tilde\mu_w$ satisfies Corollary~\ref{IV-18-18}(a)--(d) with $\|\cdot\|_{w,\ell}$ for $\|\cdot\|_{{\rm ch},\ell}$ and $Z_w$ for the periodic partition function, and clause (e) of it with the pressure $p^w_\ell=\ell^{-3}\log(|Z_w|/\underline w)$ for $p_\ell$, where $|Z_w|$ is formed on $\Lambda_{K,\ell}$ and $\underline w$ is the least positive value of $\hat w(e,m)$ over all $\mu,e,m$ ($\underline w=4$ at $w\equiv1$, $\underline w=1$ at $w=\delta_0$): for $\ell\ge\ell'\ge2\mathsf l$, $p^w_\ell\le p^w_{\ell'}$ and $\|F\|_{w,\ell}\le\|F\|_{w,\ell'}\,e^{\mathsf l^3(p^w_{\ell'}-p^w_\ell)}$; at $w\equiv1$ the same holds with $\ell^{-3}\log|Z_w|$ for $p^w_\ell$. (The division by $\underline w$ makes $e^{\mathsf l^3p^w_\ell}$, the root of order $(\ell/\mathsf l)^3$ of $|Z_w|/\underline w$, non-increasing in $\ell$: with the time $\mu$, a fold integral of $Z_w$ over $\mathsf n_\mu$ cells in the direction $\mu$ is $\sum_{e,m}\hat w(e,m)\,\mathrm{Tr}\,B_{e,m}^{\mathsf n_\mu/2}$ with positive operators $B_{e,m}$ that do not depend on $\mathsf n_\mu$; for $q=\mathsf n_\mu/\mathsf n'_\mu\ge1$, $\mathrm{Tr}\,B^{q\mathsf n'_\mu/2}\le(\mathrm{Tr}\,B^{\mathsf n'_\mu/2})^q$, $\omega\le\omega^q$ for $\omega=\hat w(e,m)/\underline w\in\{0\}\cup[1,\infty)$, and $\sum_ix_i^q\le(\sum_ix_i)^q$ for $x_i\ge0$; the three directions are taken in turn.

At $w=t\delta_0$ the undivided pressure is $p_\ell+\ell^{-3}\log t$, and for small $t$ it is larger at $\ell$ than at $\ell'<\ell$.) Further $r[n]\,|\tilde\mu_n(\prod_\kappa F_\kappa)|\le\prod_\kappa\|F_\kappa\|_{{\rm ch},\ell}$ for such functions $F_\kappa$ of the cells of one period box. For all local functions $F_\kappa$ of the cells, invariant or not, with the plain extension and the columns of every class on the edges of the board of $\Lambda_{K,\ell}$, so that the functional $\mu_w$ names the board, $|\mu_w(\prod_\kappa F_\kappa\circ\mathrm{per}_\ell)|\le\prod_\kappa\|F_\kappa\|_{w,\ell}$.
\item \emph{(The density of the mixture.)} Let $\beta>0$, $\mathsf w_{\rm mix}=\sum_se^{-\beta A_s}$ over the closed sign fields $s$, and $\hat{\mathsf w}_{\rm mix}=2^{2-2|N|}\mathsf w_{\rm mix}$, $|N|$ the number of sites. For every class and every bounded centre-blind Borel $F$, $\int F\sum_{s\in[n]}e^{-\beta A_s}dU=2^{2|N|-2}Z[n;F]$, whatever the places of the columns; so $\hat{\mathsf w}_{\rm mix}$ is the density of $\sum_nZ[n;\cdot\,]$ on the centre-blind functions. It is the restriction of an entire function of the matrix entries, centre-blind, gauge invariant and invariant under the isometries of the lattice, and $\hat{\mathsf w}_{\rm mix}\le8$. With $\bar A=\sum_p(1-\frac12|t_p|)$, $z_p=e^{-\beta|t_p|}$ and $\mathsf M(U)$ the set of the unit cubes on whose faces the signs of the $t_p$ have the product $-1$, $\mathsf w_{\rm mix}=e^{-\beta\bar A}\sum_\Gamma\prod_{p\in\Gamma}z_p$, the sum over the sets $\Gamma$ of plaquettes whose boundary, the set of the cubes with an odd number of faces in $\Gamma$, is $\mathsf M(U)$. $\log\mathsf w_{\rm mix}$ is no sum of functions of the four bonds of single plaquettes.

If the angles $\theta_p=\arccos\frac12|t_p|$ of the faces of a unit cube add up to less than $\pi$, no $t_p$ there vanishes and the product of their signs is $1$, and $\pi$ cannot be replaced by a larger number; so on $\mathfrak S=\{\theta_p<\frac\pi6\text{ for all }p\}$ the signs of the $t_p$ form a closed sign field $s(U)$, and $\mathfrak S$ is the disjoint union of the eight open, non-empty, centre-blind, gauge invariant sets $\mathfrak S_n=\{U\in\mathfrak S:n(s(U))=n\}$. Let now all sides be $\ge4$. On $\mathfrak S$, for $\beta\ge4$,
\[ \log\bigl(e^{\beta A_{s(U)}}\mathsf w_{\rm mix}(U)\bigr)=\sum_XE(X),\qquad|E(X)|\le\frac{|X|}{198}\prod_{p\in X}(400z_p), \]
the sum over the connected sets $X$ of plaquettes that are unions of vortices, where in this item two plaquettes are neighbours if they are faces of one unit cube, that is, dual bonds with a common dual site, and a set is connected if it is connected for this relation (so the consecutive plaquettes of a straight column, opposite faces of one cube, are neighbours; the sets $D$ of (4)(e) are connected through common edges instead), a vortex being a non-empty connected set of plaquettes that every unit cube meets in an even number of faces (it has at least $4$ plaquettes, and at least $\min_\mu N_\mu$ if the sign field that is $-1$ on it is not exact), with $E(X)=\log(1+\prod_{p\in X}z_p)$ for a vortex $X$ that holds no smaller one: thin vortices are local terms. For $\vartheta\le\frac\pi6$ with $\zeta=100e^{-2\beta\cos\vartheta}\le\frac12$ and every centre-blind $F\ge0$ vanishing off $\mathfrak S_n\cap\{\theta_p<\vartheta\ \forall p\}$, $(1-\epsilon)\sum_{n'}Z[n';F]\le Z[n;F]\le\sum_{n'}Z[n';F]$ with $\epsilon=\frac{3|N|}{50}\zeta^{\min_\mu N_\mu}$.
\item \emph{(Area laws from one inequality.)} Let $\beta>0$ and $D_2(l_0)\le\frac1{20}$ at an integer $l_0\ge3$. Then $D_2(l_0)>0$; put $m_*=l_0^{-1}\log\frac1{18D_2(l_0)}$, let $l_1$ be the least $2^il_0$ with $m_*l_1\ge3$, and $\sigma_*=m_*/l_1$. Then $\sigma_*=m_*/l_0$ if $m_*l_0\ge3$ and $\sigma_*>\frac16m_*^2$ otherwise; $D_2(l)\le\frac1{18}e^{-m_*l}$ for $l\in\{2^il_0\}$; $l\,\varepsilon_{\rm v}(l)\le e^{-\sigma_*l^2}$ for $l\in\{2^jl_1\}$; for every integer $h\ge4l_1$ and all $t,\lambda\ge(l_1h)^{1/2}$, on the box $(t,\lambda,h)$ whose first two sides span the twisted plane,
\[ -\log\bigl[\tfrac12(1-r)\bigr]\ge\sigma_*t\lambda-\frac{4t\lambda}{l_1h}\Bigl(\frac12\log\frac{4h}{l_1}-0.66\Bigr), \]
in particular, on every cube of side $N\ge4l_1$, $-\log[\frac12(1-r_1)]\ge\sigma_*N^2-\frac{4N}{l_1}(\frac12\log\frac{4N}{l_1}-0.66)$; over the section $(N,N)$, $N\ge4l_1$, the norm $\Lambda_-(N)$ of $T_0$ on the functions odd under one of the two flips has $\log(\|T_0\|/\Lambda_-(N))\ge\sigma_*N-\frac2{l_1}\log\frac{4N}{l_1}$; and on the even cubes of side $N\ge\max\{2R,T+1,4l_1\}$, $|\mu_{N^{\times3}}(W_{R\times T})|\le\exp[-\sigma_*RT+\frac TN\log8+\frac{2RT}{l_1N}\log\frac{4N}{l_1}]$. There is exactly one periodic Gibbs state $\nu$, the measures of the periodic cubes converge to it, and $|\nu(W_C)|\le e^{-\sigma_*|A(C)|}$ for every simple closed path $C$ of a coordinate plane that each line parallel to one fixed axis of the plane, through midpoints of bonds, meets at most twice, $|A(C)|$ the number of plaquettes enclosed; such a $C$ is made of $T$ rows of plaquettes along the axis, and if they lie in a segment of length $w'$ along it, $|\mu_{N^{\times3}}(W_C)|\le8^{T/N}e^{1.13D_2(N/2)}[\frac12(1-r_1(N))]^{|A(C)|/N^2}$ on the cubes $N\in\{2^jl_0\}$ with $N\ge\max\{4l_1,2l_0,2T,2w'\}$. If $\mathsf e(l)\le0.069$ at an integer $l\ge3$, then $D_2(8l)\le0.04736$, so that all of this holds at $l_0=8l$; and on every cube of side $N\ge2l$, $\delta(N^{\times3})\le0.13736\,\theta^{N/2l}$ and $1-r[n]\le0.4121\,\theta^{N/2l}$ for every $n\ne0$, $\theta=12.9\,\mathsf e(l)\le0.8901$.
\item \emph{(Rectangular loops.)} Let $\beta>0$, $\nu$ a periodic Gibbs state, $\mathsf m(R,T)=\nu(W_{R\times T})$ and $\mathsf w=\mathsf m(1,1)$. Then $\mathsf m$ depends on $R,T$ alone, $\mathsf m(R,T)=\mathsf m(T,R)$, $\mathsf m(R,T)^2\le\mathsf m(R,T-1)\,\mathsf m(R,T+1)$ with $\mathsf m(R,0)=1$, $\mathsf w\ge\frac{\beta}{4(\beta+1)}$, $\mathsf m(R,T)\ge\mathsf w^{RT}>0$, and $\mathsf m(R,T)\le(8\beta/(8\beta+3))^{2(R+T)-4}$ for $R,T\ge2$. Hence $V(R)=\lim_T-T^{-1}\log\mathsf m(R,T)=\inf_T(-T^{-1}\log\mathsf m(R,T))$ exists, is concave and non-decreasing on $\{0,1,2,\dots\}$ with $V(0)=0$; $\sigma_{\rm W}=\lim V(R)/R=\inf_RV(R)/R$; $\mathsf m(R,T)\le e^{-TV(R)}\le e^{-\sigma_{\rm W}RT}$; $V(1)\ge\log(1+\frac3{8\beta})$ and $V(R)\ge2\log(1+\frac3{8\beta})$ for $R\ge2$, while $V(1)<2\log(1+\frac3{8\beta})$ for $\beta<\frac1{24}$; and $0\le\sigma_{\rm W}\le V(1)\le-\log\mathsf w\le\log\frac{4(\beta+1)}\beta$.
\item \emph{(Purity floor.)} Let $N_3=2h$, $n=(k,0)_3$ with $k\ne0$, and $F\in[0,1]$ a Borel function of the fields on the closed half slab of the layers $0,\dots,h$ that is invariant under $c^k$ applied to all its layers; $F$ may be such a function of the block fields of the half slab. On the box $(N_1,N_2,2ih)$ lay $F$ into each of the $2i$ half slabs, mirrored at the odd places; let $\mathcal Z_i$ be the periodic integral of the product, $w_F=\mathcal Z_1/Z[0]$. If $\mathcal Z_1>0$, put $\mathfrak q_i=\mathcal Z_i/\mathcal Z_1^i$. Then, for $i\ge2$,
\[ r[n]\ge w_F\,\frac{\mathfrak q_i-2^{1-i}}{1-2^{1-i}},\qquad r[n]\ge w_F\bigl((2\mathfrak q_2-1)_+\bigr)^{1/2},\qquad r[n]\ge w_F\Bigl(\frac{2\|\mathsf P\|}{\mathrm{Tr}\,\mathsf P}-1\Bigr), \]
where $\mathsf P=\mathsf Y^*\mathsf Y$, $\mathsf Y$ the operator of $F$ on the half slab, so that $\mathcal Z_i=\mathrm{Tr}\,\mathsf P^i$. If $\mathcal Z_1=0$, positivity gives $\mathsf P=0$ and $\mathcal Z_i=0$ for every $i$; no purity quotient is asserted in this case.
\item \emph{(Small cubes, exactly.)} On every cube $N^{\times3}$, $N\ge3$, with $r_1\ge r_2,r_3$ its flux quotients and every time: $|Z_{e,0}/Z[0]-\frac14|\le\frac14(2r_1+r_2)\le\frac34r_1$ for all $e$; $\delta(N^{\times3})\ge\mathsf G(N)\ge\log\frac4{1+2r_1+r_2}\ge\log\frac4{1+3r_1}$; $D_q(N)\ge q^2(q-1)\,\mathsf G(N)$ for $q\ge2$; $Z[0](N^{\times3})^2/Z[0](2N,N,N)\ge4/(1+2r_1^2+r_2^2)\ge4/(1+3r_1^2)$; if $N$ is even, $N\ge6$ and $r_1\le\frac13$, then $\mathsf e(N/2)\ge\log2$; and for $\beta>0$ the bracket of (3) on the cube has $\mathfrak b(N,N)\ge\frac12(1-r_1)e^{\mathsf G(N)}\ge2(1-r_1)/(1+3r_1)$. So on a cube with $r_1\le\frac15$: $\delta\ge\log\frac52$, $\mathfrak b\ge1$, and, when $N$ is even with $N\ge6$, $\mathsf e(N/2)\le0.069$ fails.
\item \emph{(A floor and the bulk by construction.)} Let $\lambda(\beta)=\frac92\log\beta+31.4$. (a) Let $\mu\nu$ be any coordinate plane and $\rho$ the third direction. If $\beta\ge1$ and $h'\ge1$ is an integer with $2h'\le\min\{N_\mu,N_\nu\}$, then
\[ \log r[n_{\mu\nu}]\ge-N_\rho\bigl[32\beta/(h'-\tfrac12)^2+2432\,h'\lambda(\beta)\bigr]; \]
if $\min\{N_\mu,N_\nu\}\ge4+2(\beta/\lambda(\beta))^{1/3}$, then $\log r[n_{\mu\nu}]\ge-N_\rho[2464\,\beta^{1/3}\lambda(\beta)^{2/3}+4864\,\lambda(\beta)]$; and if $\min_\lambda N_\lambda\ge4h'+2$, then for every $n\ne0$, $\log r[n]\ge-[32\beta/(h'-\frac12)^2+2432\,h'\lambda(\beta)]\sum N_\rho$, the sum over the twisted planes of $n$, $\rho$ normal to the plane. (b) With $c_\infty$ of Proposition~\ref{C-366-13} and $P^{\rm G}(\beta)=\log\frac2{\pi\beta^3}-\frac32c_\infty$: $P(\beta)-P^{\rm G}(\beta)\le2.2\cdot10^5(\log\beta)\,\beta^{-1/6}$ for $\beta\ge e^{148}$.
\end{enumerate}
\end{theorem}

\paragraph{Additional hypotheses for the centre-flux limits.}
The assertions below concerning the restriction of the small-field chain
to a mixed box use the following hypothesis.

\smallskip
\noindent\textbf{(Hess).}
For every mixed box, every embedding in a periodic covering torus used in
Lemma~\ref{C-366-7}, every level of its small-field chain, and every admitted
real background, remove the factors independent of the fluctuation field.
The remaining real fluctuation integral has the form
\(\int_{D_r}e^{-H(B)}\,dB\), where
\(D_r=\{B:|B_y|<r\text{ for every fluctuation site }y\}\).
Here \(H\) includes the interaction, Jacobian, and all other
field-dependent factors of this integral, and is invariant under the
covering group. Its restriction to the fixed fields, with the action
divided by the order of the group and with Lebesgue measure at orbit
representatives, is the fluctuation integral of the mixed box.
There are no additional boundary factors or tests in either integral;
alternatively, any such factors have a common invariant convex
regularization for which the following bounds hold uniformly and which
converges to these integrals.

Let $c_->0$, $c_+>0$, $\mu>0$ and $g_0\ge0$. The function \(H\) extends to a real smooth function on \(D_{2r}\), and,
on that whole set,
\[
 \nabla^2H\ge c_-,\qquad
 \sup_y\sum_{y'}\sup_{D_{2r}}
   |\nabla_y\nabla_{y'}H|e^{\mu d(y,y')}\le c_+,\qquad
 \sup_y|\nabla_yH(0)|\le g_0.
\]
The constants can be chosen uniformly over the cutoffs, backgrounds,
embeddings and labels in a compact subset of \((0,\infty)\), with the
dependence on the level, radius and box controlled so that the right
side of Proposition~\ref{C-366-8} is at most
\[
 C|s|L^{3(K-i)}\exp[-c s_{\min}L^{K-i}]
\]
at level \(i\). The same realization and bounds hold for the intermediate
small-field integrals and boundary regularizations used in the
comparison of two consecutive cutoffs.

\smallskip
\noindent\textbf{(Flux-window).}
For the flip-invariant window \(F\) defined in Step~11 of the proof of
Theorem~\ref{C-366-4}, there are a fixed sufficiently large \(m_0\),
absolute constants $C_2>0$, $\epsilon_F\ge0$, and $t_0>0$ such that, on every cube
\(\ell=2L^m\), at every cutoff, for
\(t=\ell/b\le t_0\) and the one-plane class \(n=n_{13}\),
\[
 Z_K[0;F]\le C_2tZ_K[0],\qquad Z_K[n;1-F]\le\epsilon_F t Z_K[0].
\]
The window is read at the block layer next to a mirror at the cutoffs
\(K\ge K_\sharp\), and at the bare layer specified in that step below
\(K_\sharp\); the flux columns are placed away from these layers. The
same estimates hold for its reflected window. All constants are
independent of \(K\) and of the representative obtained by the dilations
used in that step.

\smallskip
The hypotheses are additional assumptions. The lower bound
for the Gaussian quadratic form (Proposition~\YMMref{55.2}{484},
equation~\YMMref{(1099)}{484}) does not control the full Hessian
of the interacting density required by \textbf{(Hess)}, and the
flip of the two plane means does not itself prove
\textbf{(Flux-window)}. The proof of the mixed-box and spectral-limit
claims is conditional on these inputs. The fixed-cutoff transfer
identities and the confinement implications under the stated box
inequality do not require them.

\begin{theorem}[Uniformly in the cutoff, and the limit]\label{C-366-4}
Let the constants be fixed as in Proposition~\YMMref{69.1}{595}, and $b>0$.
The mixed-box floor in \textup{(2)(a)} and the two-cutoff comparison in \textup{(4)} assume \textbf{(Hess)}. This hypothesis is also imposed on their consequences in \textup{(2)(b,c)}, the continuum and full-sequence spectral assertions of \textup{(3)}, \textup{(5)(b,c)}, the continuum-function and zero-free conclusions of \textup{(6)}, the full-sequence limits in \textup{(7)}, the uniform-cutoff and limiting assertions of \textup{(8)}, and \textup{(9)}. In \textup{(10)}, the large-$N$ conclusion of \textup{(a)} and the normalization by positive limiting class masses in \textup{(d)} retain this hypothesis unless that positivity is assumed separately.

The ceiling in \textup{(5)(a)} and its consequence \textup{(5)(d)} assume \textbf{(Flux-window)}; so do the strict sector-survival and $v_z<1$ conclusions of \textup{(3)} and their uses. The additional \textbf{(Window estimates)} in \textup{(10)(c)} and the separate hypotheses \textup{(R)}, \textup{(W)}, \textup{(W')} and \textup{(M)} remain as stated. Fixed-cutoff identities, reflection positivity and estimates whose hypotheses do not invoke these comparisons are unaffected.

\begin{enumerate}
\item \emph{(Window.)} For every box $s$ of the cutoff $K$ and $n\ne0$, $e^{-2bL^{2K}(n_{23}s_1+n_{13}s_2+n_{12}s_3)}\le r_K[n](s;b)<1$, and $r_K[n]\le r_K[n']$ for every one-plane class $n'$ contained in $n$; $\delta^0_K\ge0$, $\delta^n_K=\log r_K[n]+\delta^0_K\le\delta^0_K$, $P(L^Kb)\ge-6L^Kb$, and the identities and inequalities of Theorem~\ref{C-366-3}(2) hold for $\delta^0_K$ in physical units. For every compact set $\mathsf J\subset(0,\infty)$ of labels and every $s_{\min}>0$ there are $j$ and $C$, depending on $\mathsf J$, $s_{\min}$ and the constants, with $\delta^0_K(s;b)\le2C|s|$ for all $b\in\mathsf J$ and $K>j$ on every box whose sides are $\ge s_{\min}$, and $\mathrm{Tr}\,T_m^{t/a_K}/\|T_0\|^{t/a_K}\le e^{2Ctl_1l_2}$ over the sections $(l_1,l_2)$ for $t,l_1,l_2\ge s_{\min}$, $m=0,1$, $b\in\mathsf J$ and $K>j$: with $\mathsf J\subset[b_-,b_+]$ one may take $j$ least with $L^jb_-\ge\gamma^{-2}$ and $L^js_{\min}\ge2L^{m_0(L^jb_+)}$, and $C=L^{3j}C_*(L^jb_+)$, where $C_*(b)$ is the constant of the comparison $|\log Z_K[0](s;b)-\varepsilon_K|s||\le C_*(b)|s|$, with numbers $\varepsilon_K(b)$, on all boxes with sides $\ge2L^{m_0(b)}$ at $b\ge\gamma^{-2}$, which Lemma~\ref{lem:pure-scaling-stability} and the monotone density of $F_0$ give; $C_*(b)$ and $m_0(b)$ do not decrease in $b$.

On the cubes, $0\le\delta^0_K((2L^m)^{\times3};b)\le(2L^m)^3[c_R+\frac92\log b]$ for $b\ge\gamma^{-2}$, $m\ge m_0(b)$ and $K\ge1$, with $c_R$ formed from the constants of Lemma~\ref{lem:pure-scaling-stability}, and $\delta^0_K(\ell'^{\times3};b)\le C_f\ell'^3\log b$ at $\ell'\in\{\ell_\flat(b)/L,\ell_\flat(b)\}$ for $b\ge b_\sharp$ of (4) and $K\ge4$, $C_f$ formed from the constants; and the supremum of $\delta^0$ over all lattices and couplings with $N_\mu/\beta$ in a compact subset of $(0,\infty)$ is finite. On every box with $L$-adic sides $\sup_K\delta^n_K<\infty$, and $\inf_Kr_K[n]>0$ if and only if $\inf_K\delta^n_K>-\infty$.
\item \emph{(Floor.)} (a) If $s$ is mixed for $z$, there is $C(s,b)<\infty$, depending on $s$, $b$ and the constants and bounded for $b$ in compact sets, with $\log Z^z_K(s;b)-|s|\,p_\infty(K,b)\ge-C(s,b)$ at every $K$ at which $s$ is a box of the cutoff $K$; hence $\inf_Kr_K[z](s;b)>0$, uniformly for $b$ in compact sets. (For $b$ in a compact set $\mathsf J$, (4) for $z$ and for the class $0$, for which $s$ is mixed too, gives $|\log r_{K+1}[z]-\log r_K[z]|\le2CL^{-K/10}$ for $K\ge K_0$, with the $C$, $K_0$ of (4) at $\mathsf J$; the first sentence of (1) bounds $r_K[z]$ from below, uniformly on $\mathsf J$, at the finitely many $K<K_0$; and $\delta^0_K\ge0$.) (b) For a class with one twisted plane the same holds on every box with $L$-adic sides, uniformly for $b$ in compact sets, and $\inf\{\rho^z(N,\beta):N_\mu\ge3,\ N_\mu/\beta\in I\ \forall\mu\}>0$ for every compact interval $I\subset(0,\infty)$. (c) For every class the same holds on every box with sides in $2L^{\mathbb Z}$, the cubes $\ell=2L^m$ among them, uniformly for $b$ in compact sets; $c(I)=\inf\{\rho^z(N^{\times3},\beta):N\ge3,\ N/\beta\in I\}>0$ for every compact interval $I\subset(0,\infty)$, so that $\sup_{j\ge1}\sup_{t\in I}(-\log\phi_j(t))\le\log(1/c(I))$; and the infimum of $\rho^z$ over the lattices with $N/\beta$ in a compact subset of $\{v_z>0\}$ of (3) is positive.

On these boxes the weights $r_K[n]/\sum_{n'}r_K[n']$ of the mixture converge, to numbers in $[v_n/8,\frac12]$ for $n\ne0$.
\item \emph{(Limit.)} For every class there is a continuous $v_z\colon(0,\infty)^3\to[0,1]$ with $\rho^z(N_j,\beta_j)\to v_z(\tau)$ along every sequence with $\beta_j\to\infty$ and $N_j/\beta_j\to\tau$; in particular $r_K[z](s;b)\to V_z(s;b)=v_z(s/b)$ along the full sequence of the cutoffs on every box with $L$-adic sides, uniformly for $b$ in compact sets. In each side, at all real values of the two others, $v_z$ is real-analytic on $(0,\infty)$, and $|V_z(s;b)-V_z(s';b)|\le\frac4{s_-}e^{\Gamma}|s-s'|_1$ for $s,s'\in[s_-,s_+]^3$, where $\Gamma$ is the supremum of $\delta^0$ over the lattices with $N_\mu/\beta\in[\tau_-/2,2\tau_+]$ for all $\mu$, $[\tau_-,\tau_+]$ an interval that holds $[s_-/b,s_+/b]$ in its interior; $\Gamma$ is finite by (1). $v_z>0$ on the boxes of (2), $\{v_z>0\}$ is open, $v_z>0$ everywhere for one twisted plane, and $v_z<1$ on every box with $L$-adic sides.

For each time and all $e$, $m$, the sector weight $\frac14\sum_k(-1)^{e\cdot k}v_{(k,m)_\mu}\ge0$ is $>0$ on all of $(0,\infty)^3$ or $=0$ on all of it, and so is $v_z$ for a class with one twisted plane; so a positive lower limit in the cutoff on one box gives one on all boxes, uniform for the sides divided by the label in a compact set; the limits obey the order of Theorem~\ref{C-366-3}(1), $\frac14\sum_k(-1)^{e\cdot k}v_{(k,m)_\mu}\ge0$, and $v_{((1,1),m)_\mu}\ge v_{((1,0),m)_\mu}+v_{((0,1),m)_\mu}-v_{(0,m)_\mu}$. For each time, each section $l\in\mathbb D^2$ and all $m,e,i$, $E^{K,m,e}_i\to E^{m,e}_i(l;b)\in[0,\infty]$ along the full sequence; the limits extend to real sections, are continuous in $b$, and $E^{m,e}_i(\lambda'l;\lambda'b)=E^{m,e}_i(l;b)/\lambda'$ for real $\lambda'>0$.

Let $H_{l,m}$ be the diagonal operator whose eigenvalues are the finite limits, the eigenvectors labelled by $(e,i)$, and $C^k=(-1)^{e\cdot k}$ on the label $e$. Then $H_{l,m}\ge0$ has finite multiplicities and no finite point of accumulation; $\mathrm{Tr}\,e^{-tH_{l,1}}\le\mathrm{Tr}\,e^{-tH_{l,0}}<\infty$ for $t>0$, and for every $L$-adic height $h_0$ the number of levels $\le\Lambda$ is at most $\exp[h_0\Lambda+2Ch_0l_1l_2]$, $C$ that of (1) at $\mathsf J=\{b\}$ and $s_{\min}=\min\{h_0,l_1,l_2\}$; $0$ is an eigenvalue of $H_{l,0}$ with label $e=00$; the space of $H_{l,1}$ and each sector $e$ of $H_{l,0}$ are not $\{0\}$; with $H_{l,m}$ at the label $b$,
\[ v_{(k,m)_\mu}\bigl((t,l)/b\bigr)=\mathrm{Tr}\bigl(e^{-tH_{l,m}}C^k\bigr)/\mathrm{Tr}\,e^{-tH_{l,0}}, \]
and, for every time $\mu$, the sector weights $w^{[\mu]}_{e,m}=\frac14\sum_k(-1)^{e\cdot k}v_{(k,m)_\mu}$, functions on $(0,\infty)^3$, are by this formula $w^{[\mu]}_{e,m}(s/b)=\mathrm{Tr}(P_ee^{-s_\mu H_{l,m}})/\mathrm{Tr}\,e^{-s_\mu H_{l,0}}$ on the boxes $s$ with $L$-adic sides, $l$ the section of $s$ normal to $\mu$ and $P_e$ the projection on the label $e$, each trace divided by the periodic trace of its own time and section; at every point of $(0,\infty)^3$ they obey the crossing relation of Theorem~\ref{C-366-3}(1) with $w^{[\mu]}_{e,m}$ in the place of $Z^{[\mu]}_{e,m}$, $w^{[1]}_{e_1e_2,m}=\frac12\sum_{m',e_2'}(-1)^{e_2m'+e_2'm}w^{[2]}_{e_1e_2',m'}$, and likewise for every pair of times, since on each lattice the relation of Theorem~\ref{C-366-3}(1) divided by $Z[0]$ is this relation between the numbers $\frac14\sum_k(-1)^{e\cdot k}r[(k,m)_\mu]$; $E_{\rm mag}(l)=\inf\mathrm{spec}\,H_{l,1}=\lim_Ka_K^{-1}\log(\|T_0\|/\|T_1\|)$ is finite, and $E_{\rm el}(l)=\min\{E^{0,e}_1:e_1=1\}=-\lim_{t\to\infty}t^{-1}\log\frac{1-v}{1+v}$, $v=v_{((1,0),0)_\mu}((t,l)/b)$, is finite.

For each time, $L$-adic section $l$, direction $\nu$ of it and integer $q\ge2$, $a_K^{-1}[q\log\|T_0(l)\|-\log\|T_0(l^{[\nu]})\|]$ converges to a finite number $\ge0$, $l^{[\nu]}$ the section with $l_\nu$ multiplied by $q$. $\delta^0_K$ converges on every box with $L$-adic sides. On the cubes $\ell=2L^m$, $D(t)=\lim_K[L^3\log Z_{K,\ell}-\log Z_{K,L\ell}]$ exists, depends on $t=\ell/b$ alone, is continuous, and $D(t)-(L^3-1)\log(1/t)$ has a finite limit as $t\to0$; $\delta(t)=\lim_K\delta^0_K(\ell^{\times3};b)=\sum_{r\ge0}L^{-3(r+1)}D(L^rt)$, locally uniformly, and $\lim_K\delta^n_K=\log v_n(t,t,t)+\delta(t)$ is finite for every $n$. The limits do not depend on the choice of the constants: they are the same for the blocking factor of every other admissible choice.

By (6) and Theorem~\textup{\YMMref{AD}{607}}, $t\mapsto v_z(t,t,t)$ and $D$ are real-analytic on $(0,\infty)$.
\item \emph{(Two consecutive cutoffs.)} If $s$ is mixed for $z$, $\ell\in2L^{\mathbb Z}$ and $\mathsf J\subset(0,\infty)$ is compact, there are $C,K_0$ with
\[ \Bigl|\log Z^z_{K+1}(s;b)-\log Z^z_K(s;b)-\frac{|s|}{\ell^3}\Delta_K[\ell,b]\Bigr|\le CL^{-K/10}\qquad(b\in\mathsf J,\ K\ge K_0); \]
so for two classes $z,z'$ for which $s$ is mixed, $|\log(Z^z_K/Z^{z'}_K)(s;b)-\log(V_z/V_{z'})(s;b)|\le C'L^{-K/10}$, the limits lying in $(0,\infty)$. On boxes with sides in $2L^{\mathbb Z}$, in every class, the left member tends to $0$. Uniformly in the label: there are numbers $b_\sharp,c_1,C_V,C_G,j_0$, formed from the constants, such that with $\varpi=3+2\mathfrak p_0/r_\star$ and $c_2=20\varpi$, for $b\ge b_\sharp$, every $s$ mixed for $z$ with sides in $L^{\mathbb N}\cup2L^{\mathbb N}$, all $\ge\ell_\flat(b)$, $S=\max_\mu s_\mu$, and every $K\ge c_2\log_LS$ with $((\lceil K/1260\rceil-j_0)\log L+\log b)^{2\mathfrak p_0}\ge C_V|s|\log b$, the left member at $\ell=\ell_\flat(b)$ is at most $c_1S^\varpi L^{-K/10}$; the last condition is void where $|s|\le(\log b)^{2\mathfrak p_0-1}/C_V$, so that for every $0\le q'<2\mathfrak p_0-3r_\star-1$, all large $b$ and $s_\mu=\frac12\ell_\sharp(b)L^{i_\mu}$ with $L^{i_1+i_2+i_3}\le(\log b)^{q'}$ only $K\ge c_2\log_LS$ is asked. And $0\le\frac{|s|}{\ell_\flat^3}\log Z_{K,\ell_\flat}-\log Z^z_K(s;b)\le C_G|s|\log b$ for $b\ge b_\sharp$, these boxes and $K\ge j_0+2$.
\item \emph{(Small cubes.)} (a) There are an absolute constant $C_\square$ and $t_0>0$ depending on the constants with $r_K[n]\le C_\square t$ on every cube $\ell=2L^m$ with $t=\ell/b\le t_0$, at every $K$, for every $n\ne0$; so $-\log r_K[n]\ge\log(1/t)-\log C_\square$, and no floor is uniform as $t\to0$. (b) As $t\to0$, $v_z(t,t,t)/t\to\mathscr G_z$, the number \eqref{C-366-14}, whose numerical evaluation gives $0.10868$, $0.08251$, $0.07574$ for one, two, three twisted planes; so $C^{-1}t\le v_z(t,t,t)\le Ct$ for $t\le t_0$; and for $z,z'\ne0$, with $\vartheta$ and $\theta^a$ of Proposition~\ref{C-366-13},
\[ \lim_{t\to0}\frac{v_z(t,t,t)}{v_{z'}(t,t,t)}=\exp\Bigl\{\frac12\int_0^\infty\sum_a\Bigl[\prod_\mu\vartheta(u,\theta^a_\mu(z))-\prod_\mu\vartheta(u,\theta^a_\mu(z'))\Bigr]\frac{du}u\Bigr\}, \]
which for two and three twisted planes against one is $0.75926$ and $0.69695$. (c) On the odd cubes $s=\frac12\ell_\sharp(b)$, $b\ge b_\sharp$, at every cutoff, $|\log r_K[n](s^{\times3};b)-\log(t\,C^n(N)/N)|\le\omega_1(t)$, $N=s/a_K$, $t=s/b$, $C^n(N)=\lim_{\beta\to\infty}\beta\rho^n(N^{\times3},\beta)$, with $\omega_1(t)\to0$ as $t\to0$. (d) There is $t_1>0$, formed from the constants, such that on every cube $2L^m$ with $t\le t_1$, at every cutoff, $r_1\le\frac15$; there, by Theorem~\ref{C-366-3}(13), $\delta^0_K\ge\log\frac52$, $\mathfrak b_K\ge1$ on the cube and $\mathsf e\le0.069$ fails at the half side; and $\liminf_{t\to0}\inf_K\delta^0_K\ge\log4$, the four electric sectors being equally filled in the limit.
\item \emph{(The label.)} Let a reference coupling $g\le\gamma$, a side $\ell\in\mathfrak L(g)$ (equation~\YMMref{(1528)}{615}) and a label $b_0\in I(g)$ (Corollary~\YMMref{74.2}{616}) be given, with the radius and the constants of Theorem~\textup{\YMMref{AD}{607}}. There are $R_0,R_0'>0$ and $\mathsf M_{\rm S}<\infty$, formed from these data and the constants and not depending on $K$, such that for every $K$ the periodic partition function of the cube $\ell$ has no zero at the complex labels $b_0-c$, $|c|<R_0$, and $c\mapsto r_K[n](\ell;b_0-c)$ is holomorphic in $|c|<R_0$ with $|r_K[n]|\le27\mathsf M_{\rm S}$ on $|c|\le R_0/8$ for one or two twisted planes, and holomorphic in $|c|<R_0'$ with $|r_K[n]|\le12600\mathsf M_{\rm S}$ on $|c|\le R_0'/15$ for three; so $\sup_K|\partial_b^\nu r_K[n](\ell;b_0)|\le27\mathsf M_{\rm S}\,\nu!\,(8/R_0)^\nu$ for one or two planes, $\le12600\mathsf M_{\rm S}\,\nu!\,(15/R_0')^\nu$ for three.

For every $t'>0$ there are $j'$, $\varrho>0$ and $\mathsf M$ such that for $j\ge j'$ the function $\phi_j$ extends holomorphically to $\{|t-t'|<\varrho\}$ with $|\phi_j|\le\mathsf M$ there; hence $\phi_j\to v_n(t,t,t)$ with all derivatives, locally uniformly on a complex neighbourhood of $(0,\infty)$. There is an open $\Omega\subset\mathbb C$ containing $(0,\infty)$ such that for every compact $Q\subset\Omega$ there is $j_Q$ with $Z[n]\ne0$ on the cube of $2L^j$ sites at $\beta=2L^j/t$ for $t\in Q$, $j\ge j_Q$. With $m_0=\inf_Kr_K[n](\ell;b_0)>0$, $\int\exp[-L^K(b_0-c)A_{\eta^n}]dU\ne0$ for all complex $|c|<R_1$ and all $K$, where $R_1=\min\{R_0/16,m_0R_0/864\mathsf M_{\rm S}\}$ for one or two twisted planes and $R_1=\min\{R_0'/30,m_0R_0'/756000\mathsf M_{\rm S}\}$ for three.

For labels in $[b_1,b_2]\subset I(g)$, with $R_*=R_0(b_1)$, $R_*'=R_0'(b_1)$, $\mathsf m$ the infimum of $r_K[n](\ell;b)$ over $K$ and $b\in[b_1,b_2]$, $\mathcal O=a_K^{-1}A$ under $\mu_0$ and $\mathcal O^n=a_K^{-1}A_{\eta^n}$ under $\mu_n$: at every $K$, for one or two twisted planes, the means differ by at most $V_1=216\mathsf M_{\rm S}/(R_*\mathsf m)$ and the variances by at most $3456\mathsf M_{\rm S}/(R_*^2\mathsf m)+V_1^2$; for three, by $V_1=189000\mathsf M_{\rm S}/(R_*'\mathsf m)$ and $5670000\mathsf M_{\rm S}/(R_*'^2\mathsf m)+V_1^2$.
\item \emph{(Block pairs.)} Extend the averaging map $\mathcal M$ to pairs $(U,s)$, $s$ closed, by $\mathcal M_s(U)(c)=\gamma[c]\,\mathcal M(\gamma U)(c)$, where $\gamma$ is a sign field on the bonds with $d\gamma=s$ on the two blocks at the ends of the block bond $c$ and $\gamma[c]$ is its product over the $L$ bonds of the straight path of $c$; the block sign field is the product of $s$ over the $L^2$ plaquettes of a block plaquette, and $s^{(d)}$ is its $(K-d)$-th iterate. Let $w\ge0$, $w\not\equiv0$, $\ell=2L^m$ with $m\ge1$, $K\ge2$ and $0\le d\le K$. Under $\sum_nw(n)\sum_{s\in[n]}e^{-\beta A_s}dU$ the law of the block pairs $(\mathcal M^{K-d}_sU,s^{(d)})$ of depth $d$ is reflection positive at the mirrors midway between neighbouring layers of the block lattice of depth $d$ normal to each direction in which $\hat w\ge0$; for admissible $w$ in all three directions; and so is every weak limit point as $K\to\infty$ of the laws of the towers $((\mathcal M^{K-d}_sU,s^{(d)}))_{1\le d\le K}$. At every cutoff the block sign field $s^{(d)}$ of a closed $s$ is closed and $n(s^{(d)})=n(s)$, its product over a plane of the block lattice being the product of $s$ over the plane of the lattice that holds it: the class is a function of the block pairs. Let $b\ge\gamma^{-2}$ and $MR(b)\le L^m$, let $d_\natural$ be the least $d\in\mathbb N$ with $L^db\ge\mathsf t_\natural$, and for $d\ge d_\natural$ let
\[ \mathsf R(d)=2M^{-3}(\ell L^d)^3\exp\bigl[-c'A_0^2(\log L^db)^{2\mathfrak p_0}\bigr]\exp\bigl[(C+\tfrac92\log b)\,\ell^3\bigr]\le C(\ell,b)L^{3d}e^{-cd^2}, \]
where $\mathsf t_\natural$, $c'$, $A_0$, $C$, $c$ are formed from the constants and $C(\ell,b)$ from $\ell$, $b$ and the constants; $\mathsf R(d)$ does not depend on $K$. For integers $d\ge d_\natural$ and $K\ge d+1$ let $\hat E_{K,d}$ be the event of the pairs $(U,s)$ that for some depth $d'\in[d,K]$ and some block plaquette $P$ of depth $d'$ the matrix $s^{(d')}_P\,\partial(\mathcal M^{K-d'}_sU)(P)$ is farther than $1/400L^2$ from $\mathbf 1$: a union of covariant local events, one for each block plaquette, which at $s=1$ is the event that some plaquette of the block field $\mathcal M^{K-d'}U$ of some depth $d'\in[d,K]$ is farther than $1/400L^2$ from $\mathbf 1$. Its periodic probability is at most $\mathsf R(d)$; off it, at each depth $d'\in[d,K]$, the signs of the traces of the block plaquettes of the block field $\mathcal M^{K-d'}_sU$ form the block sign field $s^{(d')}$, whose class is $n(s)$, so that off $\hat E_{K,d}$ the class is a function of each of these block fields alone; and $r_K[n]\,\mu_n(\hat E_{K,d})\le16\mathsf R(d)^{1/4}$, while under the normalised mixture of weight $w$ with $w(0)>0$ the probability of $\hat E_{K,d}$ is at most $\frac{\sum_nw(n)}{w(0)}16\mathsf R(d)^{1/4}$.

Other pairs $(\ell,b)$ are brought into this range by the dilation by a power $L^j$ with $L^jb\ge\gamma^{-2}$ and $MR(L^jb)\le L^{m+j}$; then the same holds at the depths $d\ge j+d_\natural(L^jb)$, a threshold that depends on $\ell$ and $b$. Under these conditions, in every weak limit point, with $r_n=\lim_Kr_K[n]$ and $\mathbb P^n$ the limit law of the towers in the class $n$: the class is, up to null sets, a function of the tower of block fields, and the classes with $r_n>0$ are carried by disjoint sets of towers; each class is locally dominated by the periodic limit law, that is, for $d\ge1$ and every covariant Borel event $B$ of the block pairs of depth $d$ that depends only on a quarter of the torus cut by block mirrors normal to two directions, $r_n\mathbb P^n(B)\le\mu_d(B_1)^{1/4}$, where $B_1=\{V:(V,1)\in B\}$ and $\mu_d$ is the periodic limit law of the block field of depth $d$, and $r_n\mathbb P^n(B)\le\mu_d(B_1)^{1/2}$ if $B$ depends only on a half of the torus whose normal lies in every twisted plane of $n$; and in every class with $r_n>0$ the towers are consistent under $\mathcal M_s$: with $(V_d,s_d)$ the pair of depth $d$ of a tower, almost surely $(V_d,s_d)=(\mathcal M_{s_{d+1}}V_{d+1},\ \text{the block sign field of }s_{d+1})$ for all $d\ge1$.
\item \emph{(Tension of a slab.)} For $\ell=2L^j$ and $K+j\ge2$, with $r_1$ on the cube $\ell$ and $\kappa_K=\log(2/(1-r_1))$,
\[ \Bigl(1-\frac1L\Bigr)\ell\,E_{{\rm el},K}(\ell,\ell)\le L^2\delta^0_K\bigl((\ell/L)^{\times3};b\bigr)+\kappa_K,\qquad\sigma_K(\ell)\,\ell^2\le\Bigl(\frac L{L-1}\Bigr)^2\Bigl[L^3\,\delta^0_K\bigl((\ell/L)^{\times3};b\bigr)+\kappa_K\Bigr], \]
$E_{{\rm el},K}=a_K^{-1}\varepsilon_{\rm el}$ over the section $(\ell,\ell)$, and $(L/(L-1))^2\le1.1737$ for $L\ge13$. Both right members are bounded in $K$; so $\sup_K\sigma_K(\ell)<\infty$; by Theorem~\ref{C-366-3}(13) the lower limits in $K$ of $\delta^0_K(\ell^{\times3};b)$, of $\mathsf G$ over the section of the cube and of every $D_q(\ell/a_K)$ are positive; and in every limit the least level of the sector $e=10$ of $H_{(\ell,\ell),0}$ is at most $\frac L{(L-1)\ell}[L^2\sup_K\delta^0_K((\ell/L)^{\times3};b)+\log\frac2{1-v_1(t,t,t)}]$, $t=\ell/b$. For $l\in\mathbb D$, $\sigma_K(l)\ge-(t\lambda)^{-1}\log\mathfrak b_K(t,\lambda;l)$ at every $L$-adic $(t,\lambda)$; for $t>t_0$, $\lambda>\lambda_0$ in $\mathbb D$ and $s'=2L^i\le\min\{t_0,\lambda_0,l\}$,
\[ -\log\mathfrak b_K(t,\lambda;l)\ge(t-t_0)\Bigl[(\lambda-\lambda_0)\sigma_K(l)-\frac{2\lambda_0l}{s'^3}\bar\delta(s')\Bigr]-\frac{t_0\lambda l}{s'^3}\bar\delta(s'),\qquad\bar\delta(s')=\sup_K\delta^0_K(s'^{\times3};b), \]
at all $K$ at which the boxes and the cube are formed; so $\liminf_K\sigma_K(l)>0$ if and only if $\mathfrak b_K(t,\lambda;l)\le c<1$ at one $L$-adic $(t,\lambda)$ for all large $K$.
\item \emph{(Small twisted section.)} There is $c_0>0$, depending on the constants, such that for $l\in L^{\mathbb Z}$ with $t=l/b\le c_0$, over the section $(l,l)$ with the time $1$ at every cutoff with $l/a_K\in L^{\mathbb N}$, the following hold. In this item $\lambda_0\ge\lambda_1\ge\lambda_2\ge\dots$ are all the eigenvalues of $T_1$, the four sectors together, each repeated according to its multiplicity, and the levels of $T_1$ are the numbers $a_K^{-1}\log(\lambda_0/\lambda_i)$, $i\ge0$; as $\lambda_0=\|T_1\|$, they are the levels $E^{K,1,e}_j$ of the setting, the four sectors together, each diminished by $a_K^{-1}\log(\|T_0\|/\|T_1\|)$, and by (3) they converge along the full sequence to the eigenvalues of $H_{(l,l),1}$, each diminished by $E_{\rm mag}(l,l)=\inf\mathrm{spec}\,H_{(l,l),1}$, which is the least eigenvalue of the sector $00$. (a) The top eigenvalue $\lambda_0$ of $T_1$ is simple, and the gap $a_K^{-1}\log(\lambda_0/\lambda_1)$, $\lambda_1$ the next eigenvalue, lies between $1.2142/l$ and $\pi L/((L-1)l)\le3.4034/l$; for $t\le c_0'$, $c_0'\le c_0$ formed from the constants, it lies between $3.0117/l$ and $3.19/l$ at every such cutoff, and at all large cutoffs every level of $T_1$ but the least, in all four sectors, is at least $1.2519/l$, with $\sum e^{-lE}\le e^{0.2515}-1<0.2860$ over these levels $E$. (b) With $E(e)$ the least eigenvalue of $H_{(l,l),1}$ in the sector $e$ minus its least eigenvalue $E_{\rm mag}(l,l)$, $lE(e)\to\pi|e|$ as $t\to0$: $\pi,\pi,\sqrt2\pi$ for $e=10,01,11$; and $l$ times the second eigenvalue of $H_{(l,l),1}$ in the sector $00$, counted with multiplicity, minus $E_{\rm mag}(l,l)$ tends to $2\pi$ as $t\to0$: in the normalisation of this item, $l$ times the limit in the cutoff of the least level of $T_1$ in the sector $00$ above the vacuum tends to $2\pi$. (c) Let $t\le c_0'$, and let the cutoffs be those with $l/a_K=L^J$, $J\ge2$.

A plaquette function is a bounded Borel function of the numbers $1-\frac12\eta(p)\,\mathrm{tr}\,\partial U(p)$ of finitely many plaquettes $p$, $\eta$ a closed sign field of the class; it is a covariant function. Let $\omega_K(G)=\lim_{h\to\infty}\mu_{n_{23}}(G)$ on the boxes $(h,l/a_K,l/a_K)$, for plaquette functions $G$ of finitely many layers: the state of the twisted cylinder $C_l=\mathbb R\times(\mathbb R/l\mathbb Z)^2$ at the cutoff $K$, invariant under the translations of its lattice and under the reflection $\theta_K$ in the plane at the time $a_K/2$, midway between two layers. For $f\in{\rm Lip}_c(C_l)$ let $\mathcal O_K(f)=a_K^{-1}\sum_pf(x(p))(1-\frac12\eta(p)\,\mathrm{tr}\,\partial U(p))$, $x(p)$ the barycentre of $p$, $e_K(f)=a_K^{-1}\sum_pf(x(p))\bar a_{K,p}$, $\bar a_{K,p}$ the mean of $1-\frac12\mathrm{tr}\,\partial U(p')$ in the periodic measure of the cube $2l/L$ at the cutoff $K$ for a plaquette $p'$ of the orientation of $p$, $\mathcal O^-_K(f)=\mathcal O_K(f)-e_K(f)$, and $\|f\|_\sharp=\max\{\sup|f|,40M\,{\rm Lip}\,f\}$.

For $n'\ge1$ and $d,R>0$ there are $C_{n'}(d,R)$ and $K_{n'}$, formed from $n'$, $d$, $R$, $l$, $t$ and the constants, with $|\omega_K(\prod_{q\le n'}\mathcal O^-_K(f_q))|\le C_{n'}(d,R)\prod_q\|f_q\|_\sharp$ for $K\ge K_{n'}$ whenever the supports of $f_1,\dots,f_{n'}$ have mutual distances $\ge d$ and each lies in a ball of radius $R$. Every sequence of cutoffs tending to infinity has a subsequence along which the levels of $T_1$ converge, to $E_0\le E_1\le\dots$, the eigenvalues of $H_{(l,l),1}$ diminished by $E_{\rm mag}(l,l)$, and $S_{n'}(f_1,\dots,f_{n'})=\lim\omega_K(\prod_q\mathcal O^-_K(f_q))$ exists for all $n'$ and all tuples of functions of class $C^1$ with compact and disjoint supports; each $S_{n'}$ is multilinear, symmetric, invariant under the translations of $C_l$ and the reflection of the time, bounded as above, and reflection positive.

For every such limit, the Hilbert space $\mathcal H$ formed from the tuples supported at positive times, with the vacuum $\Omega$, carries a contraction semigroup $e^{-sH}$ with $H\ge0$ of pure point spectrum contained in $\{E_i\}$; $E_0=0$ and $3.0117/l\le E_1\le3.19/l$; $\ker H=\mathbb C\Omega$ and $\mathrm{spec}\,H\subset\{0\}\cup[3.0/l,\infty)$; for tuples $F,G$ supported at positive times and $s\ge0$,
\[ |S(\theta F,\tau_sG)-S(F)S(G)|\le e^{-3.0\,s/l}\,\bigl(S(\theta F,F)-S(F)^2\bigr)^{1/2}\bigl(S(\theta G,G)-S(G)^2\bigr)^{1/2}, \]
$\theta$ the reflection of the time and $\tau_s$ the translation by $s$; with the law of (4) used also on the boxes $(L^3l,l,l)$, $3.0$ may be replaced by $6.1$ in the last two assertions, the action density being blind to the electric sectors other than $00$. Finally $\mathcal H\ne\mathbb C\Omega$: with $l=L^m$, let $j(t)$ be the least integer for which the reference coupling $g=(L^{j(t)}/2t)^{-1/2}$ is at most the ceiling $\gamma_W$ of the reference couplings of Theorem~\textup{\YMMref{AV}{851}} and $2L^{j(t)-1}$ is a side at $g$ (equation~\YMMref{(1528)}{615}), an integer formed from $t$ and the constants; there is a number $d_*>0$, formed from $t$ and the constants, such that the bump $f^\circ(y)=(1-|y-(2d^\circ,\xi)|^2/d^{\circ2})_+^2$ of the radius $d^\circ=L^{m-j(t)}d_*$, $\xi$ a point of the section, has $S_2(\theta f^\circ,f^\circ)-S_1(f^\circ)^2\ge\frac18\,10^{-6}c_{\rm sd}b^{-2}>0$, $b$ the label and $c_{\rm sd}=9/16\pi^2$; so $d^\circ$ is $l$ times a number formed from $t$ and the constants.
\item \emph{(Even cubes.)} Let $n\ne0$, $q_K(b)=r_K[n](\ell_\sharp(b)^{\times3};b)$, $K_1(b)=\lfloor\frac1{68}\log_Lb\rfloor$ and $C^n(N)=\lim_{\beta\to\infty}\beta\rho^n(N^{\times3},\beta)$. There is $t_{\rm w}>0$ such that for $t\le t_{\rm w}$ there is a least integer $\mu(t)$ with $b'(t)=2L^{\mu(t)}/t\ge b_\sharp$ and $\ell_\sharp(b'(t))=2L^{\mu(t)}$; then $\log\frac1t\le\log b'(t)\le2\log\frac1t$ and $\phi_j(t)=q_{j-\mu(t)}(b'(t))$. (a) $|\log\rho^n(N^{\times3},N/t)-\log(tC^n(N)/N)|\le\omega(t)\to0$ ($t\to0$) for $N\in2L^{\mathbb N}$ with $N\ge N(t)$, and for $N\in2L^{\mathbb N}$ with $N\le\ell_\sharp(b'(t))\,b'(t)^{1/68}$. (b) Assume \textbf{(R)}: for some $\kappa>0$, $|\log q_{K+1}-\log q_K|\le c_1\ell_\sharp^\varpi L^{-\kappa K}$ for $b\ge b_\sharp$, $K\ge c_2\log_L\ell_\sharp$, as (4) gives on mixed boxes.

Then $\sup_{b'\ge b}\sum_{K\ge K_1(b')}|\log q_{K+1}-\log q_K|(b')\to0$ as $b\to\infty$; the bound of (a) holds for all $N\in2L^{\mathbb N}$ at $t\le t_0$, so $ct\le r_K[n]\le Ct$ on the small cubes $2L^m$ at every cutoff; and $\sum_j|\phi_{j+1}-\phi_j|<\infty$ on $(0,t_{\rm w}]$. For the lower bound $ct\le r_K[n]$ and for the bounded variation the one-sided form $\log q_{K+1}\ge\log q_K-c_1\ell_\sharp^\varpi L^{-\kappa K}$ of (R) is enough. (c) Assume \textbf{(Window estimates)} of Definition~\ref{C-366-15}, as well as \textbf{(W)}: for one $\Lambda<\lambda_\flat$ and some $c,\eta,t_*>0$, $\varsigma_K\ge e^{-\Lambda}\varsigma_{K_\sharp}-e^{-c(\log b)^{1+\eta}}$ for $K\ge K_\sharp$, $t\le t_*$ (Lemma~\ref{C-366-12} and Definition~\ref{C-366-15}, with $F_t$ or $F'_t$).

Then $r_1(N,N/t)\ge c't$ for small $t$ and all $N\in2L^{\mathbb N}$. $\Lambda\le0.55L-1.1$ is admissible for $F_t$, and $\Lambda=\frac L2$, a loss of $\frac12$ for each cube of the box $(\ell,\ell,L\ell)$, for $F'_t$, as $L\ge7$. (W) may be replaced by the same inequality for $F'_t$ on the box $(\ell,\ell,i\ell)$ at one $2\le i\le L$, with $L$ replaced by $i$ in the definition of $\varsigma_K$ and $\Lambda<\log(0.42\cdot1.8712^i)$; or by \textbf{(W$'$)}: for some $\gamma'>0$, $\|\mathsf P_K\|\ge\frac12e^{\gamma'}\mathrm{Tr}\,\mathsf P_K$ for $K\ge K_\sharp$ and $t\le t_*$, $\mathsf P_K$ the operator $\mathsf P$ of Theorem~\ref{C-366-3}(12) for $F_t$ on the cube $\ell$ at the cutoff $K$; under (W$'$), $r_1(N,N/t)\ge\frac12\gamma'c_\flat t$ for small $t$ and all $N\in2L^{\mathbb N}$.

With the window $F''_t$ of Definition~\ref{C-366-15} each of these gives the same for $r_2$. (d) On a cube $2L^m$ with $b\ge\gamma^{-2}$ and $MR(b)\le L^m$ the weak limit points of the laws of (7) are, in every class $n$ with $\lim_Kr_K[n]>0$, carried by towers consistent under $\mathcal M_s$, by (7). Assume \textbf{(M)}: $|\mu^{K+1}_{\rm mix}(H\circ\mathcal M^{K+1-d}_s)-\mu^K_{\rm mix}(H\circ\mathcal M^{K-d}_s)|\le o_K\sup|H|$, $\sum_Ko_K<\infty$, for $\mu_{\rm mix}=\mathbb M_w/\mathbb M_w[1]$ at $w\equiv1$ and all bounded Borel $H$ of the block pairs of depth $d$ with $H(\gamma V,s\,d\gamma)=H(V,s)$. Then in every class the laws of the block pairs under $\mu_n$ converge in total variation.
\end{enumerate}
\end{theorem}

\begin{theorem}[Area laws uniformly in the cutoff, under one inequality]\label{C-366-5} The final small-side exclusion uses \textbf{(Hess)} and \textbf{(Flux-window)}; the fixed-cutoff confinement implications use only their stated box hypothesis.
Let the constants be fixed as in Proposition~\YMMref{69.1}{595} and $b>0$, and assume the hypothesis

\noindent\textbf{(H)} there are $s_0>0$ and $K_0$ such that for every $K\ge K_0$ the number $l_K=s_0/2a_K$ is an integer $\ge3$ and $\mathsf e(l_K)\le0.069$ at $\beta=L^Kb$.

\noindent At each cutoff $K\ge K_0$, (H) is one inequality between the partition functions of two periodic boxes, $\mathsf e(l_K)=2F_0(l_K,2l_K,2l_K)-F_0(2l_K,2l_K,2l_K)\le0.069$; the theorem is conditional on this hypothesis. Then for every $K\ge K_0$, at $\beta=L^Kb$: Theorem~\ref{C-366-3}(10) applies at $l_0=8l_K$, with $l_0a_K=4s_0$, $D_2(l_0)\le0.047355$, $m_*l_0\ge0.1597$, $4s_0\le l_1a_K\le150.4\,s_0$ and $\sigma_*a_K^{-2}\ge2.65\cdot10^{-4}s_0^{-2}$; on every cube of physical side $\ell\ge602\,s_0$, $-\log[\frac12(1-r_1)]\ge2.65\cdot10^{-4}\ell^2/s_0^2-(\ell/2s_0)\log(\ell/s_0)$; the periodic Gibbs state $\nu_K$ of the lattice $a_K\mathbb Z^3$ (Section~\YMMref{91}{697}) is unique and $|\nu_K(W_{R\times T})|\le\exp[-2.65\cdot10^{-4}(Ra_K)(Ta_K)/s_0^2]$ for every rectangle; on every cube of physical side $\ell\ge s_0$, $1-r_K[n]\le0.4121\cdot0.8901^{\ell/s_0}$ for every $n\ne0$; and $\sigma_K(\ell)\ge1.9\,\ell^{-2}$ for every $\ell\in\mathbb D$ with $\ell\ge s_0$ and $\ell/a_K\in\mathbb N$. All these bounds are uniform in $K\ge K_0$. Under \textbf{(Hess)} and \textbf{(Flux-window)}, with $t_1$ of Theorem~\ref{C-366-4}(5), (H) fails if $s_0\le t_1b$: it can hold only at physical sides of the order of the label or more.
\end{theorem}

\subsection{Construction estimates}

We use Theorem~\YMMref{D}{8}, the inductive description of the effective densities, with its constants and their conditions (Proposition~\YMMref{69.1}{595} and Corollary~\YMMref{69.4}{604}), and the lattices and tori $T_\ell$ (Section~\YMMref{3}{13}); through Theorem~\YMMref{D}{8} we use the small-field chain on $T_\ell$ with its small-field tests and its control of large fields (Proposition~\ref{prop:ward-fixed-volume-rarity}), the comparison of two consecutive cutoffs on $T_\ell$ (Section~\YMMref{116}{807}), reflection positivity on tori, the statements on blocks and on the averaging map $\mathcal M$ (its covariance under lattice isometries and sign changes), and the convergence in total variation of the laws of the block fields on $T_\ell$ (Corollary~\textup{\ref{IV-62-66}}).

We use the stability bound on periodic cubes (Lemma~\ref{lem:pure-scaling-stability}): for $b\ge\gamma^{-2}$ and $m_\iota\ge m_0(b)$, $\log Z_{K,\ell_\iota}$ at $\ell_\iota=2L^{m_\iota}$ lies within a constant times $\ell_\iota^3$ of a vacuum energy $E_\iota$, and $|E_1/\ell_1^3-E_2/\ell_2^3|$ is bounded, no constant depending on $K$ or $m_\iota$. We use Theorem~\YMMref{AD}{607}, the local source-energy and stability theorem, with the windows of labels (Corollary~\YMMref{74.2}{616}), the sides (equation~\YMMref{(1528)}{615}), the extents (\eqref{eq:analytic-domain-diameter}) and the bound on the number of localization domains (Lemma~\ref{lem:analytic-domain-geometry}); and, of the constants, that the blocking factor $L$ is an odd integer with $L\ge13$, and that the ordered choice in the proof of Proposition~\YMMref{69.1}{595}, with the further floors of Section~\YMMref{82}{665}, permits every sufficiently large odd $L$ with the preceding constants fixed, so that both $L$ and $L+2$ serve.

For Theorem~\ref{C-366-4}(9)(c) we use besides, with the radius and the constants of Theorems~\YMMref{AJ}{673} and~\YMMref{AD}{607}, the bounds on the joint cumulants of the action density on periodic tori, uniform in the cutoff (Lemma~\ref{lem:dilated-torus-cumulant-bound}), and the short-distance law of its two-point function (Theorem~\textup{\YMMref{AV}{851}}), whose leading term is $c_{\rm sd}b^{-2}$ times the free kernel, $c_{\rm sd}=9/16\pi^2$ (Theorem~\YMMref{C}{7} and equation~\YMMref{(2240)}{822}), with the witness depth (equation~\YMMref{(2321)}{843}); and, for its last assertion, the normalisers at the first-passage data (Proposition~\YMMref{91.12}{708}\,(1)), the transfer principle for one law and for two laws under one normaliser (Proposition~\YMMref{91.11}{707}), the first facts on states (Section~\ref{sec:state-limit-conventions}), the device of sign averaging over a packing with its count (the packing and sign-averaging argument in the proof of Proposition~\YMMref{127.1}{848}), and the first-passage data at the labels of the windows (Corollary~\YMMref{74.2}{616} and Lemma~\ref{lem:setup-label-coverage}).

For Proposition~\ref{C-366-19} we use besides the lemma on rearranged paths (Lemma~\YMMref{4.1}{17}). On one lattice we use the exact statements among the construction results: the periodic lattice and its means, the crossing operator, slabs and slab operators, boards and the periodic chessboard estimate (Proposition~\YMMref{86.3}{681}), the chessboard theorem on the infinite lattice (Corollary~\ref{IV-18-18}), and the statements on periodic Gibbs states (Proposition~\YMMref{91.2}{699} and Proposition~\YMMref{91.9}{705}). A Gibbs state at $\beta$ on the infinite lattice $a\mathbb Z^3$ is, as in Section~\YMMref{91}{697}, a Borel probability measure $\nu$ on its fields with $\nu(\gamma_BF)=\nu(F)$ for every finite set $B$ of bonds and every continuous function $F$ of finitely many bonds, $\gamma_B$ the conditional kernel of $B$ with the weight $\exp[-\beta\sum_p(1-\frac12\mathrm{tr}\,\partial U(p))]$, the sum over the plaquettes $p$ that hold a bond of $B$ (equation~\YMMref{(1818)}{698}); it is periodic if it is invariant under the translations of a subgroup of finite index of $a\mathbb Z^3$ (Section~\YMMref{91}{697}), and every periodic Gibbs state is invariant under all isometries of the lattice (Proposition~\YMMref{103.4}{758}).

For a box $N$, $\mu_N$ is the measure $\langle\cdot\rangle=\mu_0$ of the periodic box $N$; where it stands beside a measure on the fields of the infinite lattice, it is extended periodically. Theorem~\ref{C-366-5} rests, besides, on its hypothesis (H): at each cutoff $K\ge K_0$, the one inequality $\mathsf e(l_K)\le0.069$ between the partition functions of the two periodic boxes $(l_K,2l_K,2l_K)$ and $(2l_K,2l_K,2l_K)$.

\subsection{Estimates and proofs}

\begin{lemma}[Power sums; lists of levels]\label{C-366-6}
Let $T\ge0$ be compact, $Q$ a projection commuting with $T$, $z(N)=\mathrm{Tr}(QT^N)<\infty$ for integers $N\ge N_*$, and $\lambda=\|QT\|$. Then $z(N)^2\le z(N-j)z(N+j)$ for $N-j\ge N_*$; $z(N)^{1/N}$ does not increase and tends to $\lambda$; $z(N+N')\le\lambda^{N'}z(N)\le z(N)z(N')$; and, if $\lambda>0$, $z(N)/\lambda^N\ge1$ does not increase and is at most $(z(N)^q/z(qN))^{1/(q-1)}$ for every integer $q\ge2$. (a) If $(E_i)$, $(E'_i)$ are non-decreasing sequences in $[0,\infty]$ with $f(t)=\sum_ie^{-tE_i}$ and $f'(t)=\sum_ie^{-tE'_i}$ finite for all $t>0$, and $f(t_r)=f'(t_r)$ along one sequence $t_r\to\infty$, then $E_i=E'_i$ for all $i$. (b) If non-decreasing lists $(E^K_i)_i$ in $[0,\infty]$ have $\sum_ie^{-t_0E^K_i}\le B(t_0)$ for all large $K$ and every $t_0$ of a set with infimum $0$, then $E^K_i\ge t_0^{-1}\log(i/B(t_0))$; every infinite set of indices $K$ holds an infinite subset along which all $E^K_i$ converge in $[0,\infty]$; and along it $\sum_ie^{-tE^K_i}\to\sum_ie^{-tE_i}<\infty$ for every $t>0$, locally uniformly on $(0,\infty)$, the limit list having finitely many members in every bounded set.
\end{lemma}
\begin{proof}[Proof sketch]
If $QT=0$, all the unnormalised power sums vanish. Otherwise $\lambda>0$. With $z(N)=\sum_i\lambda_i^N$, $\lambda_i$ the eigenvalues of $QT$, the first assertions are the Cauchy--Schwarz inequality, the monotonicity of $\ell^p$ norms in $p$, and $\sum_i\lambda_i^{N+N'}\le\lambda^{N'}\sum_i\lambda_i^N$, $\lambda^{N'}\le z(N')$; with $\mu_i=\lambda_i/\lambda$, $x(N)=\sum_i\mu_i^N\ge1$ does not increase and $x(qN)\le x(N)$, so $x(N)^q/x(qN)\ge x(N)^{q-1}$, the powers of $\lambda$ cancelling in $z(N)^q/z(qN)$. For (a), with $f_i(s)=\sum_{p\ge i}e^{-sE_p}$, $e^{-sE_i}\le f_i(s)\le e^{-(s-s_1)E_i}f_i(s_1)$ for $s\ge s_1$, so $E_i=-\lim_rt_r^{-1}\log f_i(t_r)$; by induction $f_i$ and $f'_i$ agree along $(t_r)$, and so do $E_i$ and $E'_i$. Without the finiteness of $f$ at all $t>0$ the conclusion fails: $E_i=1$ and $E'_i=2$ for all $i$. For (b), $i\,e^{-t_0E^K_i}\le\sum_{p\le i}e^{-t_0E^K_p}\le B(t_0)$; $[0,\infty]^{\mathbb N}$ is compact and metrisable; and for $t\ge t_1>t_0$ the terms are bounded by $\min\{1,(B(t_0)/i)^{t_1/t_0}\}$, summable in $i$, so dominated convergence applies.
\end{proof}

\begin{lemma}[The path inequality]\label{C-366-17}
Let $X$ be a continuous, strictly positive, symmetric kernel on a compact metric space with a Borel probability measure, $u,v$ continuous and strictly positive, and $a\ge1$ odd. Then $\langle u,X^av\rangle\ge\langle u,Xv\rangle^a(\|u\|_2\|v\|_2)^{1-a}$. In particular, if $t^{(h)}$ is the kernel of $T_m^h$, $g$ a flip $c^k$ of the fields of a section, $\Delta(h)=\int t^{(h)}(V,gV)\,dV$ and $\mathsf Z(h)=\int t^{(h)}(V,V)\,dV$, then for integers $n'\ge1$, $m_1,m_2\ge1$, $m=m_1+m_2$, and odd $k'$,
\[ t^{(k'n'+m)}(V,gW)\ge\frac{t^{(n'+m)}(V,gW)^{k'}}{[t^{(2m_1)}(V,V)\,t^{(2m_2)}(W,W)]^{(k'-1)/2}},\qquad\Delta(k'n'+m)\ge\frac{\Delta(n'+m)^{k'}}{[\mathsf Z(2m_1)\mathsf Z(2m_2)]^{(k'-1)/2}}. \]
\end{lemma}
\begin{proof}[Proof sketch]
Let $\|u\|_2=\|v\|_2=1$, $c=\langle u,Xv\rangle$, and $\nu(x,y)=u(x)X(x,y)v(y)/c$, a probability density with the marginals $\nu_1=u\,Xv/c$ and $\nu_2=v\,Xu/c$ (symmetry). On $(v_0,\dots,v_a)$ the chain $\nu(v_0,v_1)\prod_{i=1}^{a-1}q_i$, with $q_i=\nu(v_{i+1},v_i)/\nu_2(v_i)$ for odd $i$ and $\nu(v_i,v_{i+1})/\nu_1(v_i)$ for even $i$, has the $i$-th marginal $\nu_1$ for even $i$ and $\nu_2$ for odd $i$; as $a$ is odd, $v_a$ carries an odd index. Writing $\langle u,X^av\rangle$ as the mean of a quotient under this chain, Jensen's inequality and Gibbs' inequality give $\log\langle u,X^av\rangle\ge a\log c+\frac{a-1}2[\mathrm{KL}(\nu_1|u^2)+\mathrm{KL}(\nu_2|v^2)]\ge a\log c$. For the second part, $X(V,W)=t^{(n')}(V,gW)$ is symmetric, since $g$ commutes with $T_m$ and $g^2=1$, and $X^{a}(V,W)=t^{(an')}(V,g^aW)$; with $u=t^{(m_1)}(\cdot,V)$, $v=t^{(m_2)}(\cdot,W)$ and odd $a$, $\langle u,X^av\rangle=t^{(an'+m)}(V,gW)$ and $\|u\|_2^2=t^{(2m_1)}(V,V)$; apply the first part at $a=k'$ and $a=1$, put $W=V$, and integrate with H\"older's inequality (exponents $k'$, $k'/(k'-1)$) and the Cauchy--Schwarz inequality.
\end{proof}

\begin{lemma}[Doubling]\label{C-366-9}
Let $\beta\ge0$, and for a box $N$ and a direction $\mu$ let $\mathsf d^{(\mu)}(N)=2F_0(N)-F_0(N^{(\mu)})$, $N^{(\mu)}$ the box with $N_\mu$ doubled. Then $\mathsf d^{(\mu)}\ge0$, and $\mathsf d^{(\mu)}(N^{(\mu)})\le\psi(\mathsf d^{(\mu)}(N))$ whenever $\mathsf d^{(\mu)}(N)<\log2$, with $\psi(x)=2\log(1+\varrho_x^2)$, $\varrho_x=(e^x-1)/(1+\sqrt{1-(e^x-1)^2})$. Further $D_2(N)\ge4\mathsf G(N)>0$ for $\beta>0$, $\mathsf G$ over the section $(N,N)$ at the height $N$, and $\mathsf e(l)\le D_2(l)$; if $\mathsf e(l)\le\frac1{20}$, then $D_2(2l)\le17.72\,\mathsf e(l)^2$; if $\mathsf e(l)\le0.069$, then $\mathsf e(2l)\le12.9\,\mathsf e(l)^2$, $D_2(2l)\le20\,\mathsf e(l)^2$, and $\delta(M,M,M)\le0.01717\,\theta^{M/l}$ for $M=2^jl$, $j\ge1$, $\theta=12.9\,\mathsf e(l)$. On every cube of side $N$ and for every $n\ne0$, $1-r[n]\le3(1-e^{-\mathsf G(N)})$, $\mathsf G$ over the section of the cube.
\end{lemma}
\begin{proof}[Proof sketch]
With the time $\mu$, $Z[0]=\|T_0\|^{N_\mu}(1+\mathsf r)$, $\mathsf r=\sum_{i\ge1}(\lambda_i/\lambda_0)^{N_\mu}$, and the doubled side has $\mathsf r'\le\mathsf r^2$; so $e^{\mathsf d^{(\mu)}}=(1+\mathsf r)^2/(1+\mathsf r')\ge1+2\mathsf r/(1+\mathsf r^2)$, which gives $\mathsf r\le\varrho_x$, and the defect at the doubled side is at most $2\log(1+\mathsf r^2)$. By Theorem~\ref{C-366-3}(2), $D_2(N)$ is a sum of three defects of doubling, each at least $\mathsf d\ge\mathsf G$ of its section, the first of them $4\mathsf d^{(1)}(N,N,N)$. By the symmetry of the axes $\mathsf e(l)=\mathsf d^{(1)}(l,2l,2l)=\mathsf d^{(2)}(2l,l,2l)=\mathsf d^{(3)}(2l,2l,l)$; doubling once in each direction, the bulk cancelling in sums of defects, gives $\mathsf a=\mathsf d^{(2)}(2l,2l,2l)\le\psi(\mathsf e)$, $\mathsf b=\mathsf d^{(3)}(2l,4l,2l)\le\psi(2\mathsf e+\mathsf a)$, $\mathsf d^{(1)}(l,4l,4l)\le4\mathsf e+2\mathsf a+\mathsf b$, $\mathsf e(2l)\le\psi(4\mathsf e+2\mathsf a+\mathsf b)$, and
\[ D_2(2l)=4\mathsf a+2\mathsf b+\mathsf e(2l),\qquad\delta(M,M,M)=\sum_{s\ge0}8^{-s-1}D_2(2^sM); \]
with $\kappa(y_0)=\varrho_{y_0}/y_0$, as $\varrho_y/y$ increases and $\log(1+u)\le u$, $\psi(y)\le2\kappa(y_0)^2y^2$ on $[0,y_0]$, $y_0<\log2$, which gives the constants of the lemma. On a cube, $\sum_kr[(k,0)_\mu]=4Z_{00,0}/Z[0]\ge4e^{-\mathsf G}$; with $x_i=1-r_i$ this is $2x_1+x_2\le4(1-e^{-\mathsf G})$, and the order of Theorem~\ref{C-366-3}(1) with $Z_{11,0},Z_{11,1}\ge0$ gives $x_1\le x_2\le2x_1$ and $x_3\le2x_2-x_1$.
\end{proof}

\begin{lemma}[The fourth power; loops]\label{C-366-10}
Let $\beta>0$. On a box $(t,t',h)$ whose first two sides span the plane of the class $n$ put $\zeta=\frac12(1-r[n])$, and over the section $(t,t')$, with the third direction as the time, $\mathsf r_h=\mathrm{Tr}\,T_0^h/\|T_0\|^h-1$ and $\mathsf d(h;t,t')=2F_0(t,t',h)-F_0(t,t',2h)$. For integers $t,t',h\ge3$ and $u\ge t$, $u'\ge t'$:
\[
\begin{multlined}
\zeta(u,u';h)\le\bigl[\zeta(t,t';h)\,e^{\delta(t,t',h)}\bigr]^{uu'/tt'}e^{-\delta(u,u',h)},\qquad2\zeta\le h\varepsilon_{\rm v}(t,t')+\mathsf r_h,\qquad h\varepsilon_{\rm v}(t,t')\le\mathsf d(h; \\
t,t')-\log(1-2\zeta(t,t';2h)),
\end{multlined}
\]
and $\zeta(t,t';h)\le1-e^{-\mathsf G(t;t',h)}\le\delta(t,t',h)$. If $\bar z=[(h\varepsilon_{\rm v}(t,t')+\frac12\mathsf r_{2h})e^{\delta(t,t',2h)}]^4<\frac12$, then
\[ h\,\varepsilon_{\rm v}(2t,2t')\le\mathsf d(h;2t,2t')-\log(1-2\bar z). \]
If $t'$ is even and $C$ is a rectangle parallel to that plane with sides $R\le t-1$ along $t$ and $T\le t'-1$ along $t'$, then $|\langle W_C\rangle|\le(4\langle\mathsf P\mathsf P'\rangle)^{T/t'}$ and $0\le\langle\mathsf P\mathsf P'\rangle\le(1-\zeta)^{1-\varrho}\zeta^\varrho+\zeta^{1-\varrho}(1-\zeta)^\varrho\le2\zeta^{\min\{\varrho,1-\varrho\}}$, $\varrho=R/t$, $\mathsf P,\mathsf P'$ the half traces of the two straight closed lines through the sides of length $T$; the second chain holds for all $t'\ge3$.
\end{lemma}
\begin{proof}[Proof sketch]
Along the third side $Z[0]=\|T_0\|^h(1+\mathsf r_h)$ and $Z[n]=\mathrm{Tr}\,T_1^h\ge\|T_1\|^h$, so $2\zeta\le1-e^{-h\varepsilon_{\rm v}}/(1+\mathsf r_h)\le h\varepsilon_{\rm v}+\mathsf r_h$. By Theorem~\ref{C-366-3}(2), $q=\frac12(Z[0]-Z[n])$ is a power sum in each length of the plane, so $q(u,u')\le q(t,t')^{uu'/tt'}$ with the volumes taken out by $\delta$: this is the first display, and doubling both short sides raises $\zeta$ to the fourth power. Back along the third side, $\|T_1\|^h\ge\mathrm{Tr}\,T_1^{2h}/\mathrm{Tr}\,T_1^h\ge(1-2\zeta(2h))Z[0](2h)/Z[0](h)$ and $\|T_0\|^h\le Z[0](h)$ give the third inequality; with the time along $t$, $1-\zeta\ge Z_{00,0}/Z[0]\ge e^{-\mathsf G}$. The fourth power follows by chaining these. For the loop, $Z\langle W_C\rangle=\mathrm{Tr}(T_0^{t'-T}\Gamma^*T_{\rm or}^T\Gamma)$ with the transfer operator $T_{\rm or}$ with two static charges, acting on functions with values in $2\times2$ matrices; H\"older's inequality for the double series over the two eigenbases gives the first bound, and $\langle\mathsf P\mathsf P'\rangle$, with the time along $t$, is a sum of the weights $1-\zeta$ and $\zeta$ of the two flux sectors at the distances $R$ and $t-R$.
\end{proof}

\begin{lemma}[Lipschitz bounds in the sides]\label{C-366-11}
Let $\beta>0$, $N$ a box, $\lambda$ a direction with $N_\lambda\ge5$, $z$ a class, $N^\flat$ the box $N$ with $\lceil N_\lambda/2\rceil$ in place of $N_\lambda$ and $N^\natural$ the box $N$ with $N_\lambda+1+\lfloor N_\lambda/2\rfloor$ in place of $N_\lambda$. Then
\[ |\rho^z(N+e_\lambda,\beta)-\rho^z(N,\beta)|\le\frac4{N_\lambda}e^{\delta(N^\flat)},\qquad|\delta(N+e_\lambda)-\delta(N)|\le\frac3{N_\lambda}\max\{\delta(N^\flat),\delta(N^\natural)\}. \]
\end{lemma}
\begin{proof}[Proof sketch]
With $\lambda$ as the time, $\rho^z=\sum_e(-1)^{e\cdot k}\mathcal W(e;m)/\sum_e\mathcal W(e;0)$, $\mathcal W(e;m)=\sum_i\vartheta_i^{N_\lambda}$ with the eigenvalues $\vartheta_i\in[0,1]$ of $P_eT_m/\|T_0\|$, and the denominator is $\ge1$. One more layer changes each term by $\vartheta^{N}(1-\vartheta)\le\vartheta^{\lceil N/2\rceil}/(\lfloor N/2\rfloor+1)$, and $\sum\vartheta_i^{\lceil N/2\rceil}$ over all sectors is at most $e^{\mathsf G}\le e^{\delta(N^\flat)}$ for each $m$, by $\mathrm{Tr}\,T_1^M\le\mathrm{Tr}\,T_0^M$. For $\delta$, $n\mapsto\log\sum_ia_i^n$ is convex and $\delta\ge0$: the increment at $N_\lambda$ lies between the mean increments over the heights $\lceil N_\lambda/2\rceil\dots N_\lambda$ and $N_\lambda+1\dots N_\lambda+1+\lfloor N_\lambda/2\rfloor$, and $\lfloor N_\lambda/2\rfloor\ge N_\lambda/3$.
\end{proof}

\begin{lemma}[Periodic Gibbs states]\label{C-366-16}
Let $\beta>0$. (a) For every periodic Gibbs state $\nu$, every bounded Borel function $F$ of the bonds in a cube of side $q$ and every cube of side $l\in2q\mathbb N$, $|\nu(F)-\mu_{l^{\times3}}(F)|\le\|F\|_\infty(2\delta(l^{\times3}))^{1/2}$. (b) Every weak limit point of the measures of the periodic cubes, extended periodically, is a periodic Gibbs state. (c) So if $\delta(l^{\times3})\to0$ as $l\to\infty$, there is exactly one periodic Gibbs state and the measures of the periodic cubes converge to it; this is so if $D_2(l_0)\le\frac1{20}$ at one $l_0\ge3$ or $\mathsf e(l)\le0.069$ at one $l\ge3$, and then $|\nu(F)-\mu_{l^{\times3}}(F)|\le0.126\|F\|_\infty e^{-m_*l/2}$ for $l\in\{2^il_0\}\cap2q\mathbb N$ in the first case.
\end{lemma}
\begin{proof}[Proof sketch]
(a) Let $A$ be a real continuous function of the bonds in the cube of side $q$. For $n\in2q\mathbb N$ let $\mathbb A_n$ be the sum of the mirror images of $A$ over the $(n/q)^3$ cells of side $q$ of the cube $n$, an array periodic under $2q$, and $F_t(n)=\log Z[0;e^{t\mathbb A_n}](n^{\times3})$. By reflection positivity at the mirrors between the cells, $F_t$ is a sum of powers in each side in $2q\mathbb N$, so $F_t(n)/n^3$ does not increase (Lemma~\ref{C-366-6}) and tends to its infimum $\pi(t)$; $\pi(0)=P$. The convex form of the domination of periodic Gibbs states by the chessboard seminorms (Proposition~\YMMref{91.9}{705}; its proof, the Gibbs equations of $\nu$ on a large box and the chessboard estimate on the periodic cube, holds for cells of every side) gives the tangent property $\pi(t)-\pi(0)\ge t\nu(A)/q^3$.

On the cube $l$, with $N_c=(l/q)^3$ cells, the $N_c$ mirror images of $A$ have one mean $\mu_{l^{\times3}}(A)$ under the measure $\mu_{l^{\times3}}$ of the periodic cube, and $|\mathbb A_l|\le N_c\|A\|_\infty$; Hoeffding's lemma gives $F_t(l)\le F_0(l)+tN_c\mu_{l^{\times3}}(A)+\frac12t^2N_c^2\|A\|_\infty^2$. As $l^3\pi(t)\le F_t(l)$ and $F_0(l)=l^3P+\delta(l^{\times3})$, with $s=tN_c>0$ this is $\nu(A)-\mu_{l^{\times3}}(A)\le\delta/s+s\|A\|_\infty^2/2$, whose infimum over $s$ is $\|A\|_\infty(2\delta)^{1/2}$; the same holds for $-A$, and bounded Borel functions follow by the regularity of the two measures. (b) is Proposition~\YMMref{91.2}{699}. (c) Two periodic Gibbs states agree on $F$ up to $2\|F\|_\infty(2\delta(l^{\times3}))^{1/2}$ for all large $l$; with (b) the periodic cubes converge.

Under $D_2(l_0)\le\frac1{20}$, Lemma~\ref{C-366-9} and Theorem~\ref{C-366-3}(2) give $\delta(N^{\times3})\le0.14108\,(N/l)^3D_2(l)\le8\cdot0.14108\cdot\frac1{18}e^{-m_*N/2}$ with $l\le N<2l$ in $\{2^il_0\}$; under $\mathsf e(l)\le0.069$, $\delta(N^{\times3})\le0.13736\,\theta^{N/2l}$. At $l\in\{2^il_0\}$, $(2\cdot0.14108\cdot\frac1{18}e^{-m_*l})^{1/2}\le0.126\,e^{-m_*l/2}$.
\end{proof}

\begin{lemma}[Chessboard bounds for every number of cells; the cylinder over a twisted section]\label{C-366-21}
Let $\beta\ge0$, $\lambda'\ge3$ odd and $\mathsf q\ge2$ even, and let $Q$ be the periodic cube of side $\mathsf q\lambda'$, a board with $\mathsf q^3$ cells of side $\lambda'$, with the fold and the chessboard seminorm (equation~\YMMref{(1749)}{680}, Proposition~\YMMref{86.3}{681}): for a bounded complex Borel function $\phi$ of one cell, $\mathrm{Fold}\,\phi$ is the product over the cells of $Q$ of the mirror copies of $\phi$, the copy in a cell being complex conjugated if the cell is odd relative to the given one, $Z_Q[0;\mathrm{Fold}\,\phi]\ge0$, and $\|\phi\|_{{\rm ch};Q}=(Z_Q[0;\mathrm{Fold}\,\phi]/Z_Q[0])^{1/\mathsf q^3}$; for real $\phi$ these are the fold and the seminorm of the setting. A plaquette function of a cell is a bounded complex Borel function $\phi$ of the numbers $1-\frac12\mathrm{tr}\,\partial U(p)$ of the plaquettes $p$ of the cell; in a class it is taken with $1-\frac12\eta(p)\,\mathrm{tr}\,\partial U(p)$ in the place of these numbers, and $\phi^{(u)}$ is its copy in the cell $u$.
\begin{enumerate}
\item[(a)] On every box with $k_\mu\lambda'$ sites in the direction $\mu$, $k_\mu\ge\mathsf q$ integers, odd ones admitted, cut into cells of side $\lambda'$, and in every class $n$,
\[ \Bigl|Z\Bigl[n;\prod_u\phi_u^{(u)}\Bigr]\Bigr|\le\prod_uZ_Q[0;\mathrm{Fold}\,\phi_u]^{1/\mathsf q^3};\qquad\text{in particular}\quad Z[n]\le Z_Q[0]^{|N|/|Q|}. \]
\item[(b)] Let $\beta>0$ and $N=k\lambda'$ with an integer $k\ge\mathsf q$; let $\omega=\lim_{h\to\infty}\mu_{n_{23}}$ on the boxes $(h,N,N)$, the state of the twisted cylinder $\mathbb Z\times(\mathbb Z/N)^2$ on the plaquette functions of finitely many layers, $\lambda_0>\lambda_1\ge\lambda_2\ge\dots$ the eigenvalues of $T_1$ over the twisted section $(N,N)$ with the time $1$, and $\Xi=Z_Q[0]^{k^2/\mathsf q^3}\lambda_0^{-\lambda'}$. For plaquette functions $\phi_u$ of the cells $u\in\mathbb Z\times(\mathbb Z/k)^2$, all but finitely many equal to $1$,
\[ \Bigl|\omega\Bigl(\prod_u\phi_u^{(u)}\Bigr)\Bigr|\le\Xi^{J'}\prod_u\|\phi_u\|_{{\rm ch};Q}, \]
$J'$ the number of the slabs of cells in time that carry a function other than $1$. Moreover $1\le\Xi\le\exp[\Phi'/(k_0-\mathsf q)]$ for every integer $k_0>\mathsf q$, where $\Phi'=\frac{k_0k^2}{\mathsf q^3}\log Z_Q[0]-\log Z_{(k_0\lambda',N,N)}[n_{23}]$, and $\sum_i(\lambda_i/\lambda_0)^{j\lambda'}\le\Xi^j$ for every integer $j\ge\mathsf q$. With $N'=\mathsf q\lambda'$ the side of $Q$, $k^2/\mathsf q^3=N^2\lambda'/N'^3$, so that $\log\Xi=\lambda'[(N^2/N'^3)\log Z_Q[0]-\log\lambda_0]$: at a fixed partner $Q$ and a fixed section $(N,N)$, $\log\Xi$ is linear in the number $\lambda'$ of the sites of a side of a cell.
\end{enumerate}
\end{lemma}
\begin{proof}[Proof sketch]
(a) Without a twist, cut the box in a direction $\mu$ into its $k_\mu\ge\mathsf q$ slabs of cells. By the trace formula the integral is the trace of the product of the slab operators; H\"older's inequality for traces with $\mathsf q$ factors, the operator norm being at most the norm of order $\mathsf q$, bounds its modulus by the product of the norms of order $\mathsf q$ of the slab operators. The adjoint of a slab operator $Y$ is the operator of the mirror copy of the slab carrying the complex conjugate function; so the $\mathsf q$-th power of the norm of order $\mathsf q$, $\mathrm{Tr}\,(Y^*Y)^{\mathsf q/2}$, is the integral, over the box with $\mathsf q$ slabs in the direction $\mu$, of the mirror copies of the function of the slab, complex conjugated at every other place.

The three directions in turn give the bound, the copies composing to the folds on $Q$, each of which is thus an integral $\ge0$. In a class $n$ place the columns inside cells; on a reversed plaquette $e^{-\beta(1+\frac12\mathrm{tr}\,u)}=e^{-\beta(1-\frac12\mathrm{tr}\,u)}e^{-\beta\,\mathrm{tr}\,u}$, so the twist is a family of functions of cells and the periodic bound applies. After the fold the reversed plaquettes of the $\mathsf q^3$ copies form, $\mathsf q$ being even, the boundary $d\gamma$ of a sign field $\gamma$ on the bonds of $Q$, and $U\mapsto\gamma U$ removes them: the folds are those of the periodic class. (b) By the trace formula with the time $1$, $\omega(\prod_u\phi_u^{(u)})$ is $\lambda_0^{-w}$ times the mean, in the top eigenvector of $T_1$, of the product of the operators of the slabs of cells, $w$ the number of their layers; a slab with the function $1$ has the operator $T_1^{\lambda'}$, of norm $\lambda_0^{\lambda'}$.

The norm of the operator of an occupied slab is at most its norm of order $\mathsf q$, whose $\mathsf q$-th power is, as in (a), a twisted integral over the box $(\mathsf q\lambda',N,N)$ with the mirror copies of the functions of the slab, complex conjugated at every other place in time, bounded by (a), the fold of a mirror copy being the fold of the function itself (Proposition~\YMMref{86.4}{682}). At $\phi\equiv1$ this is $\lambda_0^{\lambda'}\le Z_Q[0]^{k^2/\mathsf q^3}$, that is $\Xi\ge1$. By (a), $\mathrm{Tr}\,T_1^{j\lambda'}=Z_{(j\lambda',N,N)}[n_{23}]\le Z_Q[0]^{jk^2/\mathsf q^3}$ for $j\ge\mathsf q$, which is the last assertion, and with $\mathrm{Tr}\,T_1^{k_0\lambda'}\le\lambda_0^{(k_0-\mathsf q)\lambda'}\mathrm{Tr}\,T_1^{\mathsf q\lambda'}$ it gives $\Xi^{k_0-\mathsf q}\le e^{\Phi'}$.

The linearity of $\log\Xi$ in $\lambda'$ is the identity $k^2/\mathsf q^3=N^2\lambda'/N'^3$, $\lambda_0$ and $Z_Q[0]$ reading $N$, $N'$ and $\beta$ alone.
\end{proof}

\begin{lemma}[A box with flux as a sector of a periodic torus]\label{C-366-7}
Let $\ell=2L^m$ with $m\ge1$, $K\ge2$, $s_\mu\in\{L^j,2L^j:0\le j\le m\}$, $n_\mu=\ell/s_\mu$, and $w_{\mu\nu}\in\{0,1\}$ for $\mu\ne\nu$ with $w_{\mu\nu}=1$ only if $n_\mu$ is even; let $z_{\mu\nu}=w_{\mu\nu}+w_{\nu\mu}$ mod $2$, $\mathfrak T=\{\sum_\mu k_\mu s_\mu e_\mu\}\subset\mathbb Z^3/\ell\mathbb Z^3$ and $n=|\mathfrak T|=\ell^3/|s|$. On the lattice $L^{-j}\mathbb Z^3$, $j\ge0$, which covers the lattice of $T_\ell$ of spacing $L^{-j}$, let $n_\nu(x)=\lfloor[x]_\nu/s_\nu\rfloor$, $[x]$ the nearest point of $\mathbb Z^3$, $\varepsilon_\nu(x)=(-1)^{n_\nu(x)}$ and $\varsigma_\nu(\langle x,x'\rangle)=\varepsilon_\nu(x)\varepsilon_\nu(x')$; $\varsigma_\nu$ is invariant under the shifts by $s_\mu e_\mu$, $\mu=1,2,3$, and so it is a sign field on the bonds of the lattice of $T_\ell$; and for $v=\sum_\mu k_\mu s_\mu e_\mu$
\[ (\gamma_vV)(b)=\chi_v(b)\,V(b-v),\qquad\chi_v=\prod_\mu\prod_{\nu\ne\mu}\varsigma_\nu^{\,w_{\mu\nu}k_\mu}. \]
A twist triple of the class $z$ is $\Omega\in Q_8^3$, $Q_8=\{\pm\mathbf 1,\pm i\sigma_1,\pm i\sigma_2,\pm i\sigma_3\}$, with $\Omega_\mu\Omega_\nu=(-1)^{z_{\mu\nu}}\Omega_\nu\Omega_\mu$ for $\mu\ne\nu$ (for the class whose one twisted plane is $12$, $\Omega=(i\sigma_1,i\sigma_2,\mathbf 1)$); as every element of $Q_8$ is its own inverse up to the sign, the conjugations by $\Omega_1,\Omega_2,\Omega_3$ commute. The twisted fields of a box $N$ of the class $z$ are the fields $\tilde U$ on the bonds of $\mathbb Z^3$ with $\tilde U(b+N_\mu e_\mu)=\Omega_\mu\tilde U(b)\Omega_\mu^{-1}$, $\mu=1,2,3$, with the Haar measure $d\tilde U$ at the bonds that begin in one period box and the action $A'(\tilde U)$, the sum of $1-\frac12\mathrm{tr}\,\partial\tilde U(p)$ over the $3|N|$ plaquettes whose least corner lies in that box; neither depends on the period box. Then $\int e^{-\beta A'(\tilde U)}d\tilde U=Z[z]$ on the box $N$ for every $\beta\ge0$; $v\mapsto\gamma_v$ is an action of $\mathfrak T$ on configurations that keeps the Haar measure, and the underlying translations act freely on sites, bonds and plaquettes; every operation of the small-field chain on $T_\ell$ (averaging, gauge fixing, the minimisers about which it expands and the backgrounds, the charts and their Jacobians, the small-field tests and the cuts) commutes with each $\gamma_v$; and, with $dU^{\mathfrak T}$ the Haar measure at representatives of the bonds modulo $\mathfrak T$,
\[ \int_{\{U:\ \gamma_vU=U\ \forall v\}}e^{-\beta A(U)/n}\,dU^{\mathfrak T}=Z^z_K(s;b)\qquad\text{on }\Lambda_{K,\ell}. \]
Every box that is mixed for $z$ is of this form, for some $m$ and $w$, after a dilation by a power of $L$. The construction holds on every periodic torus with multiplicities $n_\mu$: for a box $s$ of the lattice $\mathbb Z^3$, integers $n_\mu\ge1$ and $w_{\mu\nu}\in\{0,1\}$ with $w_{\mu\nu}=1$ only where $n_\mu$ is even, and $z$ formed from $w$ as above, the same $\varsigma_\nu$, $\chi_v$ and $\gamma_v$, formed on $\mathbb Z^3$ with $n_\nu(x)=\lfloor x_\nu/s_\nu\rfloor$, give on the periodic torus of the sides $n_\mu s_\mu$ an action on configurations of $\mathfrak T=\{\sum_\mu k_\mu s_\mu e_\mu\}$, of order $n=\prod_\mu n_\mu$, that keeps the Haar measure, with free underlying translations on sites, bonds and plaquettes, and $\int_{\{U:\ \gamma_vU=U\ \forall v\}}e^{-\beta A(U)/n}\,dU^{\mathfrak T}=Z[z]$ on the box $s$ for every $\beta\ge0$.
\end{lemma}
\begin{proof}[Proof sketch]
$\varsigma_\nu=-1$ exactly on the bonds that cross the walls $x_\nu\in s_\nu\mathbb Z-\frac12$; as $L$ is odd, $\varepsilon_\nu$ is constant on the blocks of every order and the walls lie between blocks; $n_\nu(x+s_\mu e_\mu)=n_\nu(x)+\delta_{\mu\nu}$, so the two ends of a bond change sign together under these shifts. $k_\mu$ is given modulo $n_\mu$, and $\chi_v$ reads it modulo $2$ only where $w_{\mu\nu}=1$, that is, where $n_\mu$ is even; so $\chi_v$ is well defined, $\chi_{v+v'}(b)=\chi_v(b)\chi_{v'}(b-v)$ makes $v\mapsto\gamma_v$ an action, and $\mathfrak T$ acts freely on the sites, the bonds and the plaquettes, its least displacement being the least side $s_\mu<\ell$. So $\gamma_v$ is the translation by $v$ followed by the multiplication by $\chi_v$, a closed central sign field which on the lifts is $\varepsilon_v(x)\varepsilon_v(x')$, $\varepsilon_v=\prod_\mu\prod_{\nu\ne\mu}\varepsilon_\nu^{w_{\mu\nu}k_\mu}=\pm1$ constant on the blocks of every order: near every block it is a gauge transformation by $\pm\mathbf 1$.

Every operation of the chain commutes with the translations (its covariance under the lattice isometries), and with the multiplication by $\chi=\chi_v$ for the following reasons. (1) The contours that define $V(y,x)$, $x$ in the block of $y$, lie inside that block, where $\chi=1$: $V(y,x)$ is kept, and with it the gauge factors and the axial gauge, which read the family $(V(y,x))$ alone; the straight path $[x,x']$ over a block bond $c$ takes the sign $\chi(c)$, which stands twice in each $W_x$ and each $Y_{x,\Gamma,\Gamma'}$ of the definition of $\mathcal M$; so the same case of that definition applies at $c$, and $\mathcal M(\chi V)=\chi\,\mathcal M(V)$. (2) $\chi$ keeps every plaquette variable and, near every bond, is a gauge transformation by $\pm\mathbf 1$, which keeps the current; with (1) iterated, the critical regular configurations over $V$ go onto those over $\chi V$: the minimisers are covariant, and the cuts and the inner balls, which read the plaquette variables of $V$ and of its minimiser, are invariant. (3) The backgrounds are covariant, $V^{\rm bf}(\chi W)=\chi V^{\rm bf}(W)$, defined and stably defined together. (4) In the argument $V(b)V^{\rm bf}(\mathcal MV)(b)^{-1}$ of a small-field test the sign $\chi(b)^2=1$ cancels: the tests are invariant; and the charts are covariant, $V_{\gamma_vW}(\gamma_vB)=\gamma_vV_W(B)$, $\gamma_v$ acting on the fluctuation fields $B$ by the translation alone. (5) The Jacobians of the charts are invariant, $\varrho_{c+v}(\gamma_vV)=\varrho_c(V)$.

Hence every density of the chain is invariant under $\mathfrak T$, and the chain of $T_\ell$ restricted to the fixed fields is a chain for them.

For the class, put $U(b)=\tilde U(b)\Omega_\mu$ at the bonds $b$ of the direction $\mu$ that begin in the period box $\prod_\mu[0,N_\mu)$ and end outside it, and $U(b)=\tilde U(b)$ at the other bonds that begin in it: a bijection of the twisted fields onto the periodic fields of the box $N$ that keeps the Haar measure. In a plaquette the factors $\Omega_\mu^{\pm1}$ cancel, except in the plaquettes of a plane $\mu\nu$ whose least corner has $x_\mu=N_\mu-1$ and $x_\nu=N_\nu-1$, where they leave the commutator $\Omega_\nu^{-1}\Omega_\mu^{-1}\Omega_\nu\Omega_\mu=(-1)^{z_{\mu\nu}}$; so $\mathrm{tr}\,\partial\tilde U(p)=\eta^z(p)\,\mathrm{tr}\,\partial U(p)$, $\eta^z$ being $-1$ exactly on one column of each twisted plane, and $\int e^{-\beta A'}d\tilde U=Z[z]$. With a twist triple $\Omega$ of $z$ and $g(x)=(-1)^{\sum_{\nu<\mu}w_{\mu\nu}n_\mu(x)n_\nu(x)}\Omega_1^{n_1(x)}\Omega_2^{n_2(x)}\Omega_3^{n_3(x)}$ on the lifts,
\[ g(x+s_\mu e_\mu)=\textstyle\prod_{\nu\ne\mu}\varepsilon_\nu(x)^{w_{\mu\nu}}\,g(x)\,\Omega_\mu ; \]
so the gauge transformation by $g^{-1}$, constant on all blocks, maps the fixed fields bijectively onto the fields with $\tilde V(b+s_\mu e_\mu)=\Omega_\mu\tilde V(b)\Omega_\mu^{-1}$, the twisted fields of the class $z$ of the box $s$ measured in the spacing of the lattice, and $dU^{\mathfrak T}$ onto $d\tilde U$, the action of a fixed field being $n$ times that of its image. On $\Lambda_{K,\ell}$, where that box is $s/a_K$ and $\beta=L^Kb$, the identity $\int e^{-\beta A'}d\tilde U=Z[z]$ gives the display; and the flat field $\tilde V\equiv\mathbf 1$ gives a flat fixed field. $n_\mu$ is even exactly when $s_\mu\in L^{\mathbb Z}$, and $z_{\mu\nu}=1$ needs $w_{\mu\nu}=1$ or $w_{\nu\mu}=1$, that is, a side of the plane in $L^{\mathbb Z}$.

On a periodic torus with the sides $n_\mu s_\mu$ the same computations, on the lattice $\mathbb Z^3$ with $n_\nu(x)=\lfloor x_\nu/s_\nu\rfloor$, read the box $s$ and the multiplicities alone: $n_\nu(x+s_\mu e_\mu)=n_\nu(x)+\delta_{\mu\nu}$, so that $\varsigma_\nu$ is a sign field on that torus; $\chi_v$ is well defined, as $w_{\mu\nu}=1$ only where $n_\mu$ is even, and $\mathfrak T$ acts freely on sites, bonds and plaquettes; the identity for $g$ gives the bijection of the fixed fields onto the twisted fields of the class $z$ of the box $s$, the action of a fixed field being $n$ times that of its image; and $\int e^{-\beta A'}d\tilde U=Z[z]$ gives the identity of the integrals. The statements on the blocks and on the chain are those of $T_\ell$.
\end{proof}

\begin{proposition}[The root of an invariant integral]\label{C-366-8}
Let $\mathfrak Y$ be a finite set with a pseudo-distance $\mathrm d$, $\mathcal A=(\mathbb R^3)^{\mathfrak Y}$, and $\mathfrak T$ a group of order $n$ acting freely on $\mathfrak Y$, keeping $\mathrm d$, and on $\mathcal A$ by $(\gamma x)_{\gamma y}=\varrho(\gamma)x_y$, $\varrho(\gamma)\in\mathrm O(3)$; $\sigma=\min\{\mathrm d(y,\gamma y):\gamma\ne1\}$ when $n>1$. For $n=1$ the two integrals below agree exactly; assume $n>1$ for the estimate. Let $c_->0$, $c_+>0$, $\mu>0$, $g_0\ge0$, $r>0$, $\rho=\frac32r$, $D_a=\{|x_y|<a\ \forall y\}$, and $H\in C^\infty(D_{2r})$ real and $\mathfrak T$-invariant with $\nabla^2H\ge c_-$ on $D_{2r}$, $\sup_y\sum_{y'}\sup_{D_{2r}}|\nabla_y\nabla_{y'}H|\,e^{\mu\mathrm d(y,y')}\le c_+$, and $|\nabla_yH(0)|\le g_0$. With $\alpha=\mu c_-/2c_+$ and $S_\nu=\sup_y\sum_{y'}e^{-\nu\mathrm d(y,y')}$,
\[ \Bigl|\log\int_{D_r\cap\mathcal A^{\mathfrak T}}e^{-H/n}d^{\mathfrak T}x-\frac1n\log\int_{D_r}e^{-H}dx\Bigr|\le\frac{\#\mathfrak Y}n\Bigl[\frac{4\rho c_+}{c_-}\bigl(g_0+\tfrac52\rho c_+\bigr)S_{\alpha/2}e^{-\alpha\sigma/4}+\rho^2c_+e^{-\mu\sigma/2}\Bigr], \]
$\mathcal A^{\mathfrak T}$ the fixed points and $d^{\mathfrak T}x$ Lebesgue measure at representatives of the orbits.
\end{proposition}
\begin{proof}[Proof sketch]
Interpolate $H_\tau=\frac{c_-}2|x|^2+\tau(H-\frac{c_-}2|x|^2)$; at $\tau=0$ both integrals factor over the sites and agree. The $\tau$-derivatives of the two free energies are sums of means of the same local pieces $\hat u_y(x)=x_y\cdot\int_0^1\nabla_yH(sx)\,ds-\frac{c_-}2|x_y|^2$, under the full measure and under the product of $n$ copies of the quotient measure. Fix $y_0$ and representatives containing the ball of radius $\sigma/2$ about $y_0$; the gradients of the two Hamiltonians differ only through interactions that leave this ball, and in the norm weighted by $\mathsf W=\mathrm{diag}(e^{-\alpha\mathrm d(y_0,y)})$, in which the Hessian stays $\ge\frac12c_-$ by Schur's test,
\[ \sup|\mathsf W\nabla(h_1-h_2)|\le2\rho c_+S_{\alpha/2}e^{-\alpha\sigma/4}. \]
Two Langevin dynamics for $h_1$ and $h_2$ driven by one Brownian path stay within $2/c_-$ times this in that norm, which bounds the difference of the means of $\hat u_{y_0}$; the walls of $D_r$ are softened by a convex penalty and removed at the end, and the choice of representatives costs the second term. Integrate over $\tau$ and sum over the $\#\mathfrak Y/n$ representatives.
\end{proof}

\begin{lemma}[Block functions at two cutoffs]\label{C-366-12}
(a) For $b\ge\gamma^{-2}$, a cube $\ell=2L^m$ with $MR(b)\le L^m$, $K\ge2$ at which the conditions (k0)--(k4) and (k6) of Corollary~\textup{\ref{IV-62-66}}(a) hold (Section~\ref{sec:cutoff-calibration} at $n=1$, whose two depths are $u=\lceil K/1260\rceil$ and $u'=\lceil K/9\rceil$, the latter written $s$ there), $0\le d\le u$ and every gauge invariant Borel function $F$ with values in $[0,1]$ of the block field of depth $d$: $\mu^{K+1}_d(F)\le e^{m_K}\mu^K_d(F)+\epsilon_K$ and $\mu^K_d(F)\le e^{m_K}\mu^{K+1}_d(F)+\epsilon_K$, where, with the numbers $\epsilon_{\rm fg}$, $\epsilon_{\rm cl}$, $\epsilon_{\rm sh}$ of \YMMref{(2192)--(2194)}{808} and $\mathsf R_{t_1}$ of Section~\ref{sec:cutoff-calibration},
\[ m_K=4\epsilon_{\rm fg}+2\epsilon_{\rm cl}+2\epsilon_{\rm sh}+2\mathsf R_{t_1},\qquad\epsilon_K=\tfrac43\epsilon_{\rm sh}+\mathsf R_{t_1}. \]
These are numbers of $K$, $\ell$, $b$ and the constants, free of $F$ and $d$: $\epsilon_{\rm fg}$ is the error of forgetting the start of the small-field flow, $\epsilon_{\rm cl}$ that of the classical start, $\epsilon_{\rm sh}$ the relative mass where the small-field tests of the two cutoffs differ, and $\mathsf R_{t_1}$ the relative mass of the large fields of the periodic cube at the depth $u$, of the form of $\mathsf R(u)$ of Theorem~\ref{C-366-4}(7). Up to factors formed from the constants, $\epsilon_{\rm fg}$ is $\ell^3L^{4u-u'}$, $\epsilon_{\rm cl}$ is $(\log L^Kb)^{2\mathfrak p_0}\ell^3K^{1/2}L^{3u'-(K-u')/2}$, $\epsilon_{\rm sh}$ is $(u'-u)\ell^3L^{3u'}e^{-c_{\rm sh}(\log L^ub)^{2\mathfrak p_0}}$ and $\mathsf R_{t_1}$ is $\ell^3L^{3u}e^{-c'A_0^2(\log L^ub)^{2\mathfrak p_0}}e^{(C+\frac92\log b)\ell^3}$, with numbers $c_{\rm sh},c',A_0,C>0$ formed from the constants. (b) There are $A,c>0$ and $b_0\ge\gamma^{-2}$, formed from the constants, such that for $b\ge b_0$, $\ell=2L^m$ with $MR(b)\le L^m$ and $\ell\le2ML^2(\log b)^{r_\star}$, $K_\sharp=\lceil A\log_L\log b\rceil$, every $K\ge K_\sharp$ and every such $F$ of the block field of depth one,
\[ e^{-E}\mu^{K_\sharp}_1(F)-\epsilon\le\mu^K_1(F)\le e^E\bigl(\mu^{K_\sharp}_1(F)+\epsilon\bigr),\qquad E\le1/\log b,\quad\epsilon\le e^{-c(\log b)^{2\mathfrak p_0}}. \]
As $2\mathfrak p_0>3r_\star+2>1$ and $t=\ell/b\ge2MR(b)/b\ge2M(\log b)^{r_\star}/b$, for every $k>0$ there is $b_k$, formed from $k$ and the constants, with $\epsilon\le t^k$ for $b\ge b_k$: on these cubes $\epsilon=o(t^k)$ as $b\to\infty$, for every $k$.
\end{lemma}
\begin{proof}[Proof sketch]
(a) The comparison of two consecutive cutoffs (Section~\YMMref{116}{807}) restarts the small-field flows of the two lattices from the classical action at the depth $u'$ and compares them at the depth $u$. For $F$ as in (a), $\mu^K_d(F)$ is the periodic mean of $F(\mathcal M^{K-d}U)$, and $F\circ\mathcal M^{u-d}$ is a gauge invariant function of the block field of depth $u$. Split the periodic measure into the small-field measure of the chain at the depth $u$, carried by the inner ball, and the rest, whose relative mass is at most $\mathsf R_{t_1}$ by the control of large fields (Proposition~\ref{prop:ward-fixed-volume-rarity}); the gauge factors integrate to at most $1$ against gauge invariant functions. The small-field masses at the cutoffs $K$ and $K+1$ are compared pointwise at the fields of depth $u'$: each density differs from that of its classical run by a factor $e^{\pm\epsilon_{\rm fg}}$, the flow forgetting its start; the two classical starts differ by a factor $e^{\pm\epsilon_{\rm cl}}$; and the two systems of small-field tests differ on a set whose relative mass is at most $\epsilon_{\rm sh}$. Putting the factors and the masses together for $F$ and for $F=1$, which compares the two normalisations, gives (a) with the stated $m_K$ and $\epsilon_K$.

(b) Let $\ell\le2ML^2(\log b)^{r_\star}$, so that $\ell^3\le C(\log b)^{3r_\star}$. As $4u-u'\le-\frac K{10}+4$ and $3u'-\frac12(K-u')\le-\frac K9+4$, $\sum_{K\ge K_\sharp}(\epsilon_{\rm fg}+\epsilon_{\rm cl})\le C(\log b)^{2\mathfrak p_0+3r_\star}K_\sharp^{2\mathfrak p_0+1}L^{-K_\sharp/10}$, which is at most $1/(8\log b)$ for $K_\sharp=\lceil A\log_L\log b\rceil$ with $A$ large and $b$ large. In $\epsilon_{\rm sh}$ and $\mathsf R_{t_1}$, $(u\log L+\log b)^{2\mathfrak p_0}\ge(\log b)^{2\mathfrak p_0}+2\mathfrak p_0(\log b)^{2\mathfrak p_0-1}u\log L$, as $2\mathfrak p_0\ge1$; with $u\ge K/1260$ the second term gives a factor $L^{-K}$ for large $b$, which beats $L^{3u'}\le L^{K/3+3}$ and makes the sum over $K$ converge, and the first leaves $e^{-c(\log b)^{2\mathfrak p_0}}$.

The volume factor of $\mathsf R_{t_1}$ is $e^{(C+\frac92\log b)\ell^3}\le e^{C(\log b)^{3r_\star+1}}$, beaten by $e^{-\frac12c'A_0^2(\log b)^{2\mathfrak p_0}}$ because $2\mathfrak p_0>3r_\star+1$: this is why $\ell$ is bounded above. So $\sum_{K\ge K_\sharp}(\epsilon_{\rm sh}+\mathsf R_{t_1})\le e^{-2c(\log b)^{2\mathfrak p_0}}$ for large $b$, $E:=\sum_{K\ge K_\sharp}m_K\le1/\log b$ and $\epsilon:=\sum_{K\ge K_\sharp}\epsilon_K\le e^{-c(\log b)^{2\mathfrak p_0}}$. The same bounds give the conditions (k0)--(k4) and (k6) at every $K\ge K_\sharp$ for large $b$ ((k0) asks $K\ge10$, (k1) $b\ge\Theta_\star$, and (k4) a power $\ell^3L^{-5K/9+4}$ that is small). Now iterate (a) at $d=1$: $\mu^K_1\le e^{\sum_jm_j}\mu^{K_\sharp}_1+\sum_j\epsilon_je^{\sum_{i>j}m_i}\le e^E(\mu^{K_\sharp}_1+\epsilon)$, and $\mu^{j+1}_1\ge e^{-m_j}(\mu^j_1-\epsilon_j)$ gives $\mu^K_1\ge e^{-E}\mu^{K_\sharp}_1-\epsilon$, the sums running over $K_\sharp\le j<K$.

Finally $t\ge2M(\log b)^{r_\star}/b\ge1/b$ for large $b$, and $\epsilon\le b^{-k}$ once $c(\log b)^{2\mathfrak p_0-1}\ge k$, which holds for large $b$ as $2\mathfrak p_0>1$.
\end{proof}

\begin{proposition}[Rigidity]\label{C-366-20}
Let $q\ge2$ be an integer, $\tau\in\mathbb D^3$, $\mathfrak S=\prod_\mu\tau_\mu q^{\mathbb N_0}$, and $s^{[\mu]}$ the box $s$ with $s_\mu$ multiplied by $q$. For $b\in(0,\infty)$ let $\mathscr S_b$ be the class of the sequences $(K_j,b_j)_{j\ge1}$ of cutoffs and labels with $K_j\to\infty$ and $b_j\to b$; it is closed under subsequences, and the full sequence of the cutoffs at the fixed label $b$, $K_j=j$ and $b_j=b$, belongs to it. Along a sequence of $\mathscr S_b$ a number written with the index $K$ at the label $b$, as $r_K[n](s;b)$ or $E^{K,m,e}_i$, is taken at the cutoff $K_j$ and the label $b_j$, that is at $\beta=L^{K_j}b_j$, and it converges if it converges as $j\to\infty$. Along every sequence of $\mathscr S_b$ the following are equivalent: \textbf{(R1)} for every $s\in\mathfrak S$, $r_K[n](s;b)$ converges for $n\ne0$, and $q\log Z_K[0](s;b)-\log Z_K[0](s^{[\mu]};b)$ converges for each $\mu$; \textbf{(R2)} for every $s\in\mathfrak S$ and every $n$, $\delta^n_K(s;b)$ converges in $[-\infty,\infty)$; \textbf{(R3)} for each time, every $L$-adic section and all $m,e,i$, $E^{K,m,e}_i$ converges in $[0,\infty]$; \textbf{(R4)} (R1) with every integer $q'\ge2$ in the place of $q$, and (R2), hold on every box with $L$-adic sides.

If they hold along a sequence of $\mathscr S_b$, operators $H_{l,m}$ are formed from the limits along it as in Theorem~\ref{C-366-4}(3), and for each time, $L$-adic section $l$, direction $\nu$ of it and integer $q'\ge2$, the difference of the norms $a_K^{-1}[q'\log\|T_0(l)\|-\log\|T_0(l^{[\nu]})\|]$ converges along it to a finite number $\ge0$, $l^{[\nu]}$ the section with $l_\nu$ multiplied by $q'$. If they hold along every sequence of $\mathscr S_b$, then all levels, $r_K[z]$ on every $L$-adic box, $\delta^0_K$ and the differences of the norms converge along each of these sequences, to limits that are the same for all of them and so depend on $b$ alone; and if this is so for every $b$ of an open set $I\subset(0,\infty)$, these limits are continuous in $b\in I$ and are reached uniformly for the labels in every compact subset of $I$, the levels in a metric of $[0,\infty]$.
\end{proposition}
\begin{proof}[Proof sketch]
Fix a sequence of $\mathscr S_b$. Its labels and $b$ form a compact set $\mathsf J\subset(0,\infty)$, at which the window of Theorem~\ref{C-366-4}(1) holds with one constant $C$ at all large cutoffs, uniformly along the sequence; the window is the one input of the argument that reads the label. (R2) gives (R1): $\delta^0_K$ is bounded (Theorem~\ref{C-366-4}(1)), $r_K[n]=e^{\delta^n_K-\delta^0_K}$, and $q\log Z_K[0](s)-\log Z_K[0](s^{[\mu]})=q\delta^0_K(s)-\delta^0_K(s^{[\mu]})$. Assume (R1). Over a fixed section and with $\Theta_K(t)=\mathrm{Tr}\,T_0^{t/a_K}/\|T_0\|^{t/a_K}\in[1,B]$ (Theorem~\ref{C-366-4}(1)), the powers of $\|T_0\|$ cancel in $x^K_p=q\log Z_K[0](q^pt_0,l)-\log Z_K[0](q^{p+1}t_0,l)=q\log\Theta_K(q^pt_0)-\log\Theta_K(q^{p+1}t_0)\in[-\log B,q\log B]$, and
\[ \log\Theta_K(t_0)=\sum_{p\le J}q^{-p-1}x^K_p+O(q^{-J}) \]
uniformly in $K$; so $\Theta_K$ converges at the heights $t_0q^p$, and with it the eight sector traces $\Theta_K\cdot\frac14\sum_k(-1)^{e\cdot k}r_K[(k,m)]$. Two limits of the levels along subsequences (Lemma~\ref{C-366-6}(b)) have Laplace transforms that agree along $t_0q^p\to\infty$, hence are equal (Lemma~\ref{C-366-6}(a), applied to the sector traces, whose coefficients are positive): all levels over the section converge. To turn the box, $q\log Z_K[0](s_0,l)-\log Z_K[0](s_0,l')=s_0\kappa_K+q\log\Theta_K(s_0;l)-\log\Theta_K(s_0;l')$ with $\kappa_K=a_K^{-1}[q\log\|T_0(l)\|-\log\|T_0(l')\|]$, which does not depend on $s_0$: its convergence at one height gives it at all; three turns of the box reach every $L$-adic box, which gives (R3).

From (R3) the quotients and the differences of (R1) converge at every $L$-adic box and every $q'$, and $\delta^0_K(s)/|s|$, telescoped over the boxes $q^ps$ with a tail $\le6Cq^{-J}$ by the last inequality of Theorem~\ref{C-366-3}(2) and the window, converges: (R4). $\kappa_K\ge0$ because $\log\|T_0(\Sigma)\|/|\Sigma|=\lim_tF_0(t,\Sigma)/(t|\Sigma|)$ does not increase in a side of $\Sigma$. Let now (R1)--(R4) hold along every sequence of $\mathscr S_b$. Two of these sequences interlace into a third one, which belongs to $\mathscr S_b$; along it the limits exist, so they are the same along the two. Let this be so for every $b\in I$, and $b_i\to b$ in $I$: choose $K_i\to\infty$ with the values at $(K_i,b_i)$ within $1/i$ of the limits at the label $b_i$; the sequence $(K_i,b_i)$ belongs to $\mathscr S_b$, so the values along it, and with them the limits at the labels $b_i$, tend to the limits at $b$, and the limits are continuous on $I$.

If the convergence were not uniform on a compact set $\mathsf J'\subset I$, there would be $K_j\to\infty$ and $b_j\in\mathsf J'$, $b_j\to b\in\mathsf J'$, with the values at $(K_j,b_j)$ at a distance at least $\epsilon>0$ from the limits at the labels $b_j$; these limits tend to those at $b$, by the continuity, and so do the values along $(K_j,b_j)$, a sequence of $\mathscr S_b$, so that the distance tends to $0$.
\end{proof}

\begin{lemma}[One-sided recursions]\label{C-366-18}
Let $0\le x_K\le1$ with $x_{K+1}\ge e^{-\varepsilon_K}x_K-x'_K$, $\varepsilon_K,x'_K\ge0$, $\sum_K(\varepsilon_K+x'_K)=D'<\infty$. Then $\sum_K(x_K-x_{K+1})_+\le D'$, $\sum_K|x_{K+1}-x_K|\le1+2D'$, and $x_K$ converges. By Theorem~\ref{C-366-4}(6), if $\lim_j\phi_j(t)$ exists for every $t$ of a set with a limit point in $(0,\infty)$, it exists at every $t>0$, locally uniformly with all derivatives.
\end{lemma}
\begin{proof}[Proof sketch]
$x_K-x_{K+1}\le(1-e^{-\varepsilon_K})x_K+x'_K\le\varepsilon_K+x'_K$; the positive and the negative parts of the increments differ by at most $1$ in their sums. The second assertion is Vitali's theorem for the normal family of Theorem~\ref{C-366-4}(6).
\end{proof}

\begin{proposition}[Laplace asymptotics on one lattice]\label{C-366-13}
For a box $N$ and $\varphi\in[-\pi,\pi]^3$ let $S_\varphi(N)=\sum\log\hat k^2$ over the momenta $k_\mu\in(2\pi\mathbb Z+\varphi_\mu)/N_\mu$ mod $2\pi$ with $\hat k^2=\sum_\mu(2-2\cos k_\mu)\ne0$, $c_\infty=(2\pi)^{-3}\int_{[-\pi,\pi]^3}\log\hat k^2\,dk$, and $\vartheta(u,\varphi)=\sum_{n'\in\mathbb Z}e^{-u(2\pi n'+\varphi)^2}$. $Z^z_\beta(N)$ is $Z[z]$ on the box $N$ at $\beta$, so that $\rho^z(N,\beta)=Z^z_\beta(N)/Z^0_\beta(N)$. For $z\ne0$ let $\theta^a(z)\in\{0,\pi\}^3$, $a=1,2,3$, be the phases of the three colours under a twist triple $\Omega$ of the class $z$ (Lemma~\ref{C-366-7}) ($\theta^a_\mu=\pi$ where the $a$-th diagonal entry of the adjoint action of $\Omega_\mu$ is $-1$; for one twisted plane $12$ they are $(0,\pi,0)$, $(\pi,0,0)$, $(\pi,\pi,0)$, and none is $0$), and $\mathscr D^z(N)=\exp[-\frac12\sum_aS_{\theta^a(z)}(N)]$. There are absolute constants $C_T,C_P,C_S$ with the following properties.
\begin{enumerate}
\item[(a)] For every box, every class $z\ne0$ and $\beta\ge e$ with $\beta/(\log\beta)^3\ge C_T\bar N^{36}$, $\bar N=\max_\mu N_\mu$: 
\[
\bigl|\log Z^z_\beta(N)-|N|\log\frac2{\pi\beta^3}-\log\mathscr D^z(N)\bigr|\le C_T\bar N^{18}(\log\beta)^{3/2}\beta^{-1/2}.
\]

\item[(b)] For every $N\ge3$ and $\beta\ge e$ with $\beta^{1/4}\ge C_P(N\log\beta)^{C_P}$: 
\[
C_P^{-1}\beta/N\le Z^0_\beta(N^{\times3})(\pi\beta^3/2)^{N^3}e^{\frac32N^3c_\infty}\le C_P\beta/N.
\]
So for every $c\in(0,1]$ there are $\beta_c$ and $C_c$ with $C_c^{-1}N/\beta\le Z^z_\beta(N^{\times3})/Z^0_\beta(N^{\times3})\le C_cN/\beta$ for every $N\ge3$, $z\ne0$ and $\beta\ge\max\{\beta_c,\exp(N^c)\}$.
\item[(c)] For every cube $N\ge2$ and $\varphi\in[-\pi,\pi]^3\setminus\{0\}$, $|S_\varphi(N^{\times3})-N^3c_\infty-\log|\varphi|^2|\le C_S$ and $|S_0(N^{\times3})-N^3c_\infty-2\log N|\le C_S$. The limits $f(\varphi)=\lim_N[S_\varphi(N^{\times3})-N^3c_\infty]$, $\varphi\ne0$, and $f_0=\lim_N[S_0(N^{\times3})-N^3c_\infty-2\log N]$ exist; $f(\varphi)=\log|\varphi|^2+h(\varphi)$ and $f_0=h(0)$ with the continuous bounded function
\[ h(\varphi)=\int_0^\infty\Bigl[(4\pi u)^{-3/2}-\prod_\mu\vartheta(u,\varphi_\mu)+e^{-u|\varphi|^2}-e^{-u}\Bigr]\frac{du}u, \]
even in each $\varphi_\mu$; and $\lim_N[S_\theta(N^{\times3})-S_{\theta'}(N^{\times3})]=\int_0^\infty[\prod_\mu\vartheta(u,\theta'_\mu)-\prod_\mu\vartheta(u,\theta_\mu)]\frac{du}u$ for $\theta,\theta'\in\{0,\pi\}^3\setminus\{0\}$.
\item[(d)] For $\lambda,R>0$, $\int_{|c_\mu|\le R}e^{-\lambda\sum_{\mu<\nu}|c_\mu\times c_\nu|^2}d^9c\le48\pi^3R\lambda^{-2}$, and the integral is $\ge\pi^3e^{-5}R\lambda^{-2}$ if $\lambda R^4\ge4$.
\item[(e)] With
\begin{equation}\label{C-366-14}
\mathscr G_z=\frac{8\pi^2\exp\bigl[\frac12f_0-\frac12\sum_af(\theta^a(z))\bigr]}{\int_{[-\pi,\pi]^3}e^{-f(\varphi)}d\varphi},
\end{equation}
$\frac\beta NZ^z_\beta(N^{\times3})/Z^0_\beta(N^{\times3})\to\mathscr G_z$ as $N\to\infty$ with $\beta\ge\exp(N^c)$, for every $c\in(0,1]$; so $C^z(N)/N\to\mathscr G_z$, $C^z(N)=\lim_{\beta\to\infty}\beta\rho^z(N^{\times3},\beta)$ the number of (g).
\item[(f)] (Free values.) For $n\ne0$ put $d_n(N)=-\frac12\sum_a[S_{\theta^a(n)}(N)-|N|c_\infty]$, so that $\log\mathscr D^n(N)=d_n(N)-\frac32|N|c_\infty$. On the boxes $N=N's$ of a fixed integer shape $s$, $d_n(N's)=\log D_n(s)+d^{(1)}_n(s)N'^{-2}+O(N'^{-4}(\log N')^{7/2})$ with coefficients in closed form; on cubes $\log D_n=-\frac12\sum_af(\theta^a(n))$ is $-0.127453$, $-0.402863$, $-0.488493$ and $d^{(1)}_n$ is $0.476777$, $-0.755058$, $-0.834736$ for one, two, three twisted planes; and $\sum_{n\ne0}\log D_n(s)=-3\log2$ on every shape. With Theorem~\ref{C-366-3}(14)(b), in the range of (a) with $\beta\ge e^{148}$,
\[ \delta^n(N)\ge d_n(N)-C_T\bar N^{18}(\log\beta)^{3/2}\beta^{-1/2}-2.2\cdot10^5\,|N|(\log\beta)\beta^{-1/6}. \]
\item[(g)] (A fixed cube.) On every cube $N\ge3$, as $\beta\to\infty$, $\beta^{3N^3}Z^z_\beta(N^{\times3})$ tends to $\mathcal Z^z(N)=(\frac2\pi)^{N^3}\mathscr D^z(N^{\times3})$ for $z\ne0$, and $\beta^{3N^3-1}Z^0_\beta(N^{\times3})$ tends to
\[ \mathcal Z(N)=\Bigl(\frac2\pi\Bigr)^{N^3}\frac{e^{-\frac12S_0(N^{\times3})}}{8\pi^2}\int_{[-\pi,\pi]^3}e^{-S_\varphi(N^{\times3})}\,d\varphi\in(0,\infty); \]
so for every $z\ne0$, $\beta\rho^z(N^{\times3},\beta)$ tends to $C^z(N)=\mathcal Z^z(N)/\mathcal Z(N)\in(0,\infty)$. There are absolute constants $c_3$, $x_0>0$ such that for $\beta\ge3$
\[ \Bigl|\frac{\beta\rho^z(N^{\times3},\beta)}{C^z(N)}-1\Bigr|\le3c_3\Bigl(\frac{N^{66}\log\beta}\beta\Bigr)^{1/8}\qquad\text{whenever}\quad c_3\Bigl(\frac{N^{66}\log\beta}\beta\Bigr)^{1/8}\le x_0 . \]
\end{enumerate}
\end{proposition}
\begin{proof}[Proof sketch]
(a) In a class $z\ne0$ the minima of the twisted action form exactly two orbits under the gauge group, with stabiliser $\pm1$ and without zero modes. The Gaussian integral about them is the determinant of the lattice Laplacian with momenta shifted by the colour phases, which is $\mathscr D^z$; the cubic and quartic remainders, cut at fields of size $(\log\beta/\beta)^{1/2}$, give the error, and $(2\pi)^3/(2\pi^2)^2=2/\pi$ is the Haar normalisation. (b) In the periodic class the nine constant modes $c\in(\mathbb R^3)^3$ are kept as coordinates; by
\[ 1-\tfrac12\mathrm{tr}(ABA^{-1}B^{-1})=2\sin^2\alpha\sin^2\alpha'|\mathsf a\times\mathsf a'|^2\qquad(A=e^{i\alpha\mathsf a\cdot\sigma},\ B=e^{i\alpha'\mathsf a'\cdot\sigma},\ |\mathsf a|=|\mathsf a'|=1) \]
their action is quartic, between $\frac{32}{\pi^4N}\sum_{\mu<\nu}|c_\mu\times c_\nu|^2$ and $\frac2N\sum_{\mu<\nu}|c_\mu\times c_\nu|^2$ for $|c_\mu|\le\frac\pi2N$, and (d) gives one power of $\beta/N$ more than in the twisted classes; the linear term of the fluctuation vanishes, and the determinant of the remaining modes is bounded by (c). The two-sided bound of the quotient follows from (a) and (b) with $\log\beta\ge N^c$. (c) $S_\varphi-N^3c_\infty$ is a Riemann sum minus its integral; the theta-function representation $\log x=\int_0^\infty(e^{-u}-e^{-ux})\frac{du}u$, applied to every momentum, with Poisson summation in each direction, gives the bounds, the limits and the formula for $h$. (e) Over the moduli of the flat periodic fields, the orbit volume, the three directions $N^{3/2}d^3\theta$ and the determinants $e^{-\frac12S_0-S_\varphi}$ give
\[ Z^0_\beta(N^{\times3})\sim\Bigl(\frac2{\pi\beta^3}\Bigr)^{N^3}e^{-\frac32N^3c_\infty}\frac\beta N\,\frac{e^{-f_0/2}}{8\pi^2}\int_{[-\pi,\pi]^3}e^{-f}\,d\varphi \]
in the range of the limit, and with (a) the quotient tends to $\mathscr G_z$; the numerical evaluation of \eqref{C-366-14} uses the integral representation of $f$ and $f_0$. (f) The expansion is Poisson summation with the lattice factor $\hat k^2$ expanded to the fourth order. The eight phase sets $\{0,\pi\}^3$ of the box $N$ together give the momenta of the box $2N$: $\sum_{\varphi\in\{0,\pi\}^3}S_\varphi(N)=S_0(2N)$; so $\sum_{\varphi\ne0}[S_\varphi(N)-|N|c_\infty]=[S_0(2N)-8|N|c_\infty-2\log(2N)]-[S_0(N)-|N|c_\infty-2\log N]+2\log2$ tends to $2\log2$, and each non-zero $\varphi$ occurs three times among the colour phases of the seven classes, whence $\sum_{n\ne0}\log D_n=-\frac32\cdot2\log2$.

The last inequality is (a) with $P\le P^{\rm G}+2.2\cdot10^5(\log\beta)\beta^{-1/6}$. (g) Put $y=N^{66}\beta^{-1}\log\beta$. The twisted half is (a) at the fixed cube: for $y\le1$ and $\beta\ge3$, $N^{36}(\log\beta)^3/\beta=y^{6/11}(\log\beta)^{27/11}\beta^{-5/11}\le5.4\,y^{6/11}$, the maximum of $(\log\beta)^{27/11}\beta^{-5/11}$ being $(5.4/e)^{27/11}<5.4$, at $\log\beta=5.4$; so for $y\le y_0$, $y_0\in(0,1]$ absolute, the range of (a) holds and $|\log(\beta^{3N^3}Z^z_\beta(N^{\times3})/\mathcal Z^z(N))|\le C_T(5.4\,y^{6/11})^{1/2}\le C_T\sqrt{5.4}\,y^{1/8}\le1$, which gives the limit $\mathcal Z^z(N)$ and a relative error at most $2C_T\sqrt{5.4}\,y^{1/8}$. In the periodic class the minima of the action are the flat fields; at the flat field with the commuting holonomies $e^{i\phi_\mu\mathsf a\cdot\sigma}$ the Gaussian integral over the modes transverse to the flat fields carries the factor $e^{-\frac12S_0-S_\theta}$, $\theta=2\phi$, from one neutral colour and two charged ones of the phases $\pm\theta$, and the integral over the moduli of the flat fields, in the variable $\theta\in[-\pi,\pi]^3$, with the orbit volume and the quartic action of the constant modes of (b), gives the limit $\mathcal Z(N)$.

Its rate comes in four steps; $A$ is the action, $|M|^2=\frac12\mathrm{tr}\,M^*M$ for $2\times2$ matrices, and $a_\beta=4N^3\beta^{-1}\log\beta$. \emph{The split.} A gauge transformation sets to $\mathbf 1$ the bonds of the spanning tree formed by the bonds $\langle x,x+e_3\rangle$ with $x_3\le N-2$, $\langle x,x+e_2\rangle$ with $x_3=0$, $x_2\le N-2$, and $\langle x,x+e_1\rangle$ with $x_2=x_3=0$, $x_1\le N-2$, the sites being $x\in\{0,\dots,N-1\}^3$; in integrals of gauge invariant functions this removes the gauge group exactly, without a Jacobian. Then the variable $h_\mu$ of the bond $\langle(N-1)e_\mu,Ne_\mu\rangle$ is the holonomy along the axis $\mu$, and the field is bounded by its curvature, through chains of plaquettes along the tree: the variable of every bond of the direction $\mu$ outside the tree lies within $2NA^{1/2}$ of $h_\mu$ if the bond starts at $x_\mu=N-1$, and of $\mathbf 1$ otherwise, and $|h_\mu h_\nu h_\mu^{-1}h_\nu^{-1}-\mathbf 1|\le N(2A)^{1/2}$.

Let $\mathsf s_\mu=1-(\frac12\mathrm{tr}\,h_\mu)^2$, which is $\sin^2\frac{\theta_\mu}2$ on the flat field of the moduli $\theta$; for $0<r\le\frac14$ let $\psi_\mu=\min\{1,\max\{0,\frac13(\mathsf s_\mu/r^2-1)\}\}$, $\chi_1=\psi_1$, $\chi_2=(1-\psi_1)\psi_2$, $\chi_3=(1-\psi_1)(1-\psi_2)\psi_3$ and $\chi_{\rm s}=\prod_\mu(1-\psi_\mu)$, so that $\chi_1+\chi_2+\chi_3+\chi_{\rm s}=1$. $\chi_i$ vanishes unless $\mathsf s_i\ge r^2$: this is the regular part, at a distance at least $2r$ from $\theta=0$; $\chi_{\rm s}$ vanishes unless $\mathsf s_\mu<4r^2$ for every $\mu$: this is the part within a distance of the order $r$ of $\theta=0$, the image of the eight central flat fields $h_\mu=\pm\mathbf 1$.

The fields with $A\ge a_\beta$ carry at most $e^{-\beta a_\beta}=\beta^{-4N^3}$, which is at most $(C/\beta)^{N^3}\beta^{1-3N^3}\mathcal Z(N)$, as $\mathcal Z(N)\ge c^{N^3}$ with an absolute $c>0$. \emph{The regular part.} On the support of $\chi_1$ write $h_1=e^{i\phi_1\mathsf a\cdot\sigma}$, $|\mathsf a|=1$, $h_\nu=e^{i\alpha_\nu\mathsf a\cdot\sigma}e^{iy_\nu\cdot\sigma}$ with $y_\nu\perp\mathsf a$, $|y_\nu|<\frac\pi2$, for $\nu=2,3$, and $e^{iX(b)\cdot\sigma}$, $|X(b)|<\pi$, for the variable of every other bond $b$ outside the tree divided by $h_\mu$ or by $\mathbf 1$ as above; the flat fields are $y_2=y_3=0$, $X=0$, with the coordinates $\phi_1\mathsf a$, $\alpha_2$, $\alpha_3$ along them, and $\nu=(y_2,y_3,X)$ are the $6N^3-2$ coordinates transverse to them.

As $|h_1h_\nu h_1^{-1}h_\nu^{-1}-\mathbf 1|=2|\sin\phi_1|\sin|y_\nu|$ and $|\sin\phi_1|=\mathsf s_1^{1/2}\ge r$, the bounds by the curvature give $\|\nu\|^2\le\Lambda_rA$ with $\Lambda_r=2\pi^2N^5+\frac{\pi^2}4N^2r^{-2}$; the third derivatives of the plaquette terms along the rays through the flat field give $|A-Q|\le43N^2\|\nu\|^3$, $Q$ the quadratic form of the transverse modes at the flat field, and so $Q\ge\|\nu\|^2/\Lambda_r$. On $\{A<a_\beta\}$ then $\|\nu\|\le\varrho_r=(\Lambda_ra_\beta)^{1/2}$ and $|A-Q|\le\epsilon_rQ$, $\epsilon_r=43N^2\Lambda_r\varrho_r$: there the integral of $\chi_1e^{-\beta A}$ lies between the Gaussian integrals of $(1+\epsilon_r)Q$ and of $(1-\epsilon_r)Q$, which are those of $Q$ times $(1\pm\epsilon_r)^{1-3N^3}$, while over $\|\nu\|\le\varrho_r$ the Haar density changes by a factor within $e^{\pm1.2\varrho_r}$ and $\chi_1$ by at most $r^{-2}\varrho_r^2$, and the Gaussian measures give the outside of $\{A<a_\beta\}$ a relative mass at most $\beta^{-N^3}$ for $\log\beta\ge3$.

Let $\mathcal I_1$ be the Gaussian integral of $Q$ against $\chi_1$ over the flat fields, which is the limit of $\beta^{3N^3-1}\int\chi_1e^{-\beta A}dU$ as $\beta\to\infty$ at fixed $r$, and $\mathcal I_2$, $\mathcal I_3$ the same for $\chi_2$, $\chi_3$; the same Gaussian integral without $\chi_i$ over $\{\mathsf s_i\ge r^2\}$ is at most $\mathcal Z(N)$ by the limit, and so
\[ \Bigl|\beta^{3N^3-1}\int_{\{A<a_\beta\}}\chi_ie^{-\beta A}dU-\mathcal I_i\Bigr|\le\bigl[e^{6N^3\epsilon_r+1.2\varrho_r}-1+\beta^{-N^3}+2r^{-2}\varrho_r^2\bigr]\mathcal Z(N), \]
where, if $\Lambda_r\le3N^2r^{-2}$, $6N^3\epsilon_r\le2682\,N^{19/2}r^{-3}(\beta^{-1}\log\beta)^{1/2}$ and $r^{-2}\varrho_r^2\le12N^5r^{-4}\beta^{-1}\log\beta$: the relative error of Laplace's method on the regular part is $N^{19/2}r^{-3}(\beta^{-1}\log\beta)^{1/2}$ up to an absolute factor. \emph{The part about $\theta=0$.} Where $\chi_{\rm s}\ne0$, one of the eight changes of sign of the bonds of some directions $\mu$ that start at $x_\mu=N-1$, which keep the tree, the action and $\chi_{\rm s}$ and change the signs of these $h_\mu$, makes $\mathrm{tr}\,h_\mu>0$ for every $\mu$, and then $|h_\mu-\mathbf 1|\le\arcsin2r\le\pi r$; if besides $A<a_\beta$, the bounds by the curvature give $\sum_b|U(b)-\mathbf 1|^2\le\varrho_{\rm s}^2=6\pi^2N^2r^2+24N^5a_\beta$; let $\varrho_{\rm s}\le\frac12$.

A gauge transformation that maximises $\sum_b\frac12\mathrm{tr}\,U(b)$ (Landau's gauge) keeps this bound and writes the field as the transform, by a gauge transformation $u$ with $u(0)=\mathbf 1$, of the field $(1-|Y(b)|^2)^{1/2}\mathbf 1+iY(b)\cdot\sigma$ with $\sum_\mu[Y(\langle x,x+e_\mu\rangle)-Y(\langle x-e_\mu,x\rangle)]=0$ at every site; $Y=c+w$, $c=(c_\mu)$ the nine constant modes, with $|c_\mu|\le\varrho_{\rm s}N^{-3/2}$, and $w$ the $6N^3-6$ others. By Kirchhoff's matrix-tree theorem the Jacobian of $(u,Y)\mapsto U$ is at most $(e^{S_0(N^{\times3})}/N^3)^{3/2}(1+\epsilon_J)^{3N^3}$, $\epsilon_J=7N^2\varrho_{\rm s}$. The action of a plaquette is at least half the square of the vector part of the difference of the products along its two halves, which at constant fields is exactly $-2c_\mu\times c_\nu$; so $A\ge\frac12(1-\epsilon_A)\|d_1w\|^2+N^3\sum_{\mu<\nu}|c_\mu\times c_\nu|^2$, $\|d_1w\|^2$ the sum over the plaquettes of the squares of the circulations of $w$, $\epsilon_A=4\mathsf kN\varrho_{\rm s}+3\mathsf k^2\varrho_{\rm s}^2$ with an absolute $\mathsf k$.

The area formula, the Gaussian integral over $w$, whose determinant is $e^{6S_0(N^{\times3})}$, and (d) at $\lambda=\beta N^3$ and $R=\varrho_{\rm s}N^{-3/2}$, which gives the factor $48\pi^3R\lambda^{-2}$, together give, the powers of $2\pi^2$, of $2\pi$ and of $e^{S_0(N^{\times3})}$ cancelling against those of $\mathcal Z(N)$,
\[ \beta^{3N^3-1}\int_{\{A<a_\beta\}}\chi_{\rm s}e^{-\beta A}dU\le\mathsf E_{\rm s}\,\mathcal Z(N),\qquad\mathsf E_{\rm s}=\frac{192}{\mathcal J_N}N^{-1/2}\varrho_{\rm s}\exp\bigl[21N^5\varrho_{\rm s}+\varrho_{\rm s}^2+6N^3\epsilon_A\bigr], \]
$\mathcal J_N=N^{-2}\int_{[-\pi,\pi]^3}e^{S_0(N^{\times3})-S_\varphi(N^{\times3})}d\varphi\ge e^{-2C_S}\int_{[-\pi,\pi]^3}|\varphi|^{-2}d\varphi$ by (c). So, where $N^5\varrho_{\rm s}$ is bounded and $r^2\ge N^3a_\beta$, the mass about $\theta=0$ is of the order $N^{1/2}r$ against $\beta^{1-3N^3}\mathcal Z(N)$, the one power of $r$ being the one power of $R$ in (d). \emph{The choice of $r$.} As $r$ decreases, the error $N^{19/2}r^{-3}(\beta^{-1}\log\beta)^{1/2}$ of the regular part grows and the mass $N^{1/2}r$ about $\theta=0$ shrinks; the two are equal at $r^4=N^9(\beta^{-1}\log\beta)^{1/2}$, that is at $r=(N^{18}\beta^{-1}\log\beta)^{1/8}=N^{-6}y^{1/8}$. The root of order $8$ is this balance between the third power of $1/r$ beside $(\beta^{-1}\log\beta)^{1/2}$ and the one power of $r$ from (d); and at this $r$, $N^5\varrho_{\rm s}$, of the order $N^6r=y^{1/8}$, stays bounded, which is why the power in $y$ is $66=8\cdot6+18$.

Indeed, with $r=N^{-6}y^{1/8}$, $\beta^{-1}\log\beta=N^{-66}y$ and $y\le1$: $\Lambda_r\le3N^2r^{-2}=3N^{14}y^{-1/4}$, $\varrho_r^2\le12N^{-49}y^{3/4}$, $\epsilon_r\le447N^{-17/2}y^{1/8}$, $6N^3\epsilon_r\le2682N^{-11/2}y^{1/8}\le7y^{1/8}$, $r^{-2}\varrho_r^2\le12N^{-37}y^{1/2}$, $\varrho_{\rm s}^2\le60N^{-10}y^{1/4}$ and $N^5\varrho_{\rm s}\le8y^{1/8}$; so the exponent in $\mathsf E_{\rm s}$ is at most an absolute multiple of $y^{1/8}$, and $\mathsf E_{\rm s}\le CN^{-11/2}y^{1/8}$. As $\beta\to\infty$ at fixed $r$, $\beta^{3N^3-1}\int\chi_{\rm s}e^{-\beta A}dU$ tends to $\mathcal I_{\rm s}=\mathcal Z(N)-\mathcal I_1-\mathcal I_2-\mathcal I_3$, by the limit $\mathcal Z(N)$, and $\mathcal I_{\rm s}\le\mathsf E_{\rm s}\mathcal Z(N)$, $\mathsf E_{\rm s}$ being increasing in $a_\beta$.

Summing over the four parts and the fields with $A\ge a_\beta$,
\[ \Bigl|\frac{\beta^{3N^3-1}Z^0_\beta(N^{\times3})}{\mathcal Z(N)}-1\Bigr|\le3\bigl[e^{6N^3\epsilon_r+1.2\varrho_r}-1+\beta^{-N^3}+2r^{-2}\varrho_r^2\bigr]+2\mathsf E_{\rm s}+\Bigl(\frac C\beta\Bigr)^{N^3}\le c_3y^{1/8} \]
for $y\le y_0$, with an absolute $c_3$, any larger number serving as well, and $y_0$ taken smaller if need be, so that $\epsilon_r,\epsilon_A\le\frac12$, $\varrho_r<1$, $\varrho_{\rm s}\le\frac12$, $r\le\frac14$ and $\beta^{-N^3}+(C/\beta)^{N^3}\le y^{1/8}$. So $C^z(N)=\mathcal Z^z(N)/\mathcal Z(N)=8\pi^2e^{\frac12S_0(N^{\times3})}\mathscr D^z(N^{\times3})/\int_{[-\pi,\pi]^3}e^{-S_\varphi(N^{\times3})}d\varphi$, and with (c) and dominated convergence $C^z(N)/N\to\mathscr G_z$. Take $c_3\ge2C_T\sqrt{5.4}$ and $x_0\le\min\{\frac13,c_3y_0^{1/8}\}$: where the second right member $x=c_3y^{1/8}$ is at most $x_0$, $y\le y_0$, so that the twisted relative error is at most $x$ as well; and a quotient of two numbers within $x\le\frac13$ of $1$ lies within $3x$ of $1$.
\end{proof}

\begin{definition}[Quantitative window estimates]\label{C-366-15}
We use \textbf{(Window estimates)} for the following additional hypothesis on the three explicit windows constructed below. There are absolute constants $0<c_\flat\le C_\flat$, a number $x_\flat(L)>0$, $t_*>0$, and for $0<t\le t_*$ Borel functions $F_t,F'_t$ with values in $[0,1]$ of the block field of depth one on a half slab of normal $3$, invariant under the flip $c^k$ of the class $n_{13}=((1,0),0)_3$, and $F''_t$ invariant under the flip of the class $((1,1),0)_3$, such that on every cube of $N=\mathsf nL^j$ sites, $\mathsf n\in2L^{\mathbb N}$, $j\ge1$, at $\beta=N/t\ge3$, where they are read on the block field $\mathcal M^jU$ of $\mathsf n$ sites a side (at the cutoff $K$ on the cube $\ell$, $\mathsf n=L\ell$ and $j=K-1$, and $\mathcal M^{K-1}U$ is the block field of depth one), the numbers of Theorem~\ref{C-366-3}(12) obey: for $F_t$, if $N^{66}\beta^{-1}\log\beta\le x_\flat(L)$, $c_\flat t\le w_F\le C_\flat t$ and $\log(2^{L-1}\mathfrak q_L)\ge\lambda_\flat=0.5516L-0.7162$; for $F'_t$, if $N^{68}\beta^{-1}\log\beta\le x_\flat(L)$, $w_F\ge c_\flat t$, $\mathfrak q_i\ge0.84\cdot0.9356^i$ for $i\le L$, and $\log(2^{L-1}\mathfrak q_L)\ge\lambda_\flat=0.6266L-0.873$; and for $F''_t$ in the same range, $w_F\ge c_\flat t$ and $\mathfrak q_i\ge0.84\cdot0.944^i$ for $i\le L$.

On a cube $\ell=2L^m$ with $MR(b)\le L^m\le ML^2(\log b)^{r_\star}$, $w_K$ and $\mathfrak q_L(K)$ are these numbers at the cutoff $K$, and $\varsigma_K=w_K^L\mathfrak q_L(K)$ is an integral over the box $(\ell,\ell,L\ell)$ divided by $Z_{K,\ell}^L$.
\end{definition}
\begin{proof}[Construction and elementary properties of the windows]
\emph{The windows.} On the cube of $N=\mathsf nL^j$ sites with the time $3$ let $h=N/2$; the half slab holds the layers $0\le x_3\le h$, and the block layers inside it are those of the block lattice of $\mathcal M^jU$ at the heights $c=L^jy_3$, $1\le y_3\le\frac{\mathsf n}2-1$. In the block layer at the height $c$, for a block site $y$ of it let $\bar h_1(y)$, $\bar h_2(y)$ be the products of the block field along the straight closed block lines from $y$ in the directions $1$ and $2$, and let $\bar s_\pm(c)$, $\bar s_1(c)$, $\bar s_2(c)$ be the means over the $\mathsf n^2$ sites $y$ of the layer of $\frac12\mathrm{tr}[\bar h_1(y)\bar h_2(y)^{\pm1}]$, $\frac12\mathrm{tr}\,\bar h_1(y)$ and $\frac12\mathrm{tr}\,\bar h_2(y)$: $\bar s_\pm$ are the means of the half traces of the closed block lines that wind once in the direction $1$ and once, forwards or backwards, in the direction $2$. With $c_1=L^j$, the block layer next to the lower end of the half slab, and $c_2=h-L^j$, the one next to its upper end, $F_t=\Phi_t\circ\mathcal M^j$, $F'_t=\Phi'_t\circ\mathcal M^j$, $F''_t=\Phi''_t\circ\mathcal M^j$ with
\[ \Phi_t=\mathbf 1\{\bar s_2(c_1)>\tfrac12\}\,e^{-\frac4t(\bar s_+(c_1)^2+\bar s_-(c_1)^2)},\qquad\Phi'_t=\prod_{c\in\{c_1,c_2\}}\mathbf 1\{\bar s_2(c)>\tfrac12\}\,e^{-\frac2t(\bar s_+(c)^2+\bar s_-(c)^2)}, \]
\[ \Phi''_t=\prod_{c\in\{c_1,c_2\}}\mathbf 1\{\bar s_+(c)<-\tfrac12\}\,e^{-\frac4t(\bar s_1(c)^2+\bar s_2(c)^2)}: \]
one pin for $F_t$, two pins for $F'_t$ and $F''_t$. They are Borel, with values in $[0,1]$, and gauge invariant. The flip of the class $n_{13}$, taken across a line in a wall of the blocks of every order ($L$ is odd), passes the averaging map and becomes the change of sign of the block bonds of the direction $1$ that cross that line: $\bar h_1\mapsto-\bar h_1$ and $\bar h_2$ is kept, so $\bar s_\pm$ change sign, $\bar s_2$ is kept, and $F_t$, $F'_t$ are invariant; the flip of $((1,1),0)_3$ changes the signs of $\bar h_1$ and of $\bar h_2$, so $\bar s_1$, $\bar s_2$ change sign, $\bar s_+$ is kept, and $F''_t$ is invariant. On a flat field with the holonomies $e^{i\varphi_\mu\mathsf a\cdot\sigma}$, $|\mathsf a|=1$, the means are $\bar s_\pm=\cos(\varphi_1\pm\varphi_2)$ and $\bar s_\mu=\cos\varphi_\mu$.

The flip of $n_{13}$ acts by $\varphi_1\mapsto\varphi_1+\pi$; up to $\varphi\mapsto-\varphi$ its fixed points are $\varphi_1=\frac\pi2$, $\varphi_2\in\{0,\pi\}$, and inside $\{\bar s_2>\frac12\}$ the common zeros of $\bar s_+$ and $\bar s_-$ are the one fixed point $(\frac\pi2,0)$, at which they vanish transversally and independently: at $(\frac\pi2+a,b)$, $\bar s_+^2+\bar s_-^2=2(a^2+b^2)+O(|(a,b)|^4)$. Thus both transverse angle variables have scale $t^{1/2}$ in this flat parametrisation. Likewise the flip of $((1,1),0)_3$ acts by $\varphi_\mu\mapsto\varphi_\mu+\pi$, $\mu=1,2$; the indicator $\{\bar s_+<-\frac12\}$ singles out its fixed point $(\frac\pi2,\frac\pi2)$, about which $\bar s_1^2+\bar s_2^2$ cuts both angles at the width $t^{1/2}$: at $(\frac\pi2+a,\frac\pi2+b)$, $\bar s_1^2+\bar s_2^2=\sin^2a+\sin^2b=a^2+b^2+O(|(a,b)|^4)$.

The numbers of Theorem~\ref{C-366-3}(12) are $w_F=\mathcal Z_1/Z[0]$ and $\mathfrak q_i=\mathcal Z_i/\mathcal Z_1^i$, where $\mathcal Z_i$ is the periodic integral over the box $(N,N,iN)$ of the product of the $2i$ mirrored copies of the window. The integral identity in the definition follows from the slab trace formula. The bounds called \textbf{(Window estimates)} are additional uniform estimates on these inserted integrals.



The uninserted periodic Laplace estimate of Proposition~\ref{C-366-13} does not establish these bounds for the moving pins. In particular, \textbf{(Window estimates)} includes the stated weight bounds, numerical flat ratios and growing-lattice ranges.
\end{proof}

\begin{proposition}[Winding block loops on small tori]\label{C-366-19}
Let $\ell=2L^m$, $b\ge\gamma^{-2}$ with $MR(b)\le L^m$, and $d\ge1$ an integer with $L^db\ge\mathsf t_\natural$, $\delta_d=\exp[7(\ell L^d)^2\varepsilon(L^db)]-1\le\frac1{20}$ and $\mathsf R(d)\le\frac1{40}$, where $\mathsf t_\natural$, $\varepsilon(x)=A_0x^{-1/2}(\log x)^{\mathfrak p_0}$ and $\mathsf R(d)$ (of Theorem~\ref{C-366-4}(7)) are formed from the constants. For a site $y$ of the block lattice of depth $d$ let $\gamma_1,\gamma_2,\gamma_{12}$ be the closed block lines through $y$ that wind once in the direction $1$, once in the direction $2$, and once in $1$ and then once in $2$, $W_\gamma$ the half trace of the block field along $\gamma$, $\mathsf t$ the translation by $\frac\ell2e_3$, and $\mathsf m_\gamma=\langle W_\gamma\cdot\mathsf tW_\gamma\rangle$ at the cutoff $K>d$. Then $\mathsf m_{\gamma_2}=\mathsf m_{\gamma_1}$ and
\[ f^K=4\mathsf m_{\gamma_1}+\mathsf m_{\gamma_{12}}\ge(1-\tfrac14\delta_d^2)(1-\mathsf R(d))-5(\delta_d+2\mathsf R(d))\ge0.4743, \]
while $|f^K|\le10[\frac14(1-r_1^2)]^{\frac12-\frac1{\ell L^d}}$ with $r_1$ on the cube $\ell$; so $r_1\le1-2\,(0.045)^{2/(1-2L^{-d}/\ell)}\le0.99596$ at every cutoff $K>d$, and there $\delta^0_K\ge\mathsf G\ge0.003034$ by Theorem~\ref{C-366-3}(13). These conditions hold at $\ell=2L^m$, $b_m=\exp[L^{(m-m')/r_\star}]$, $d=1$ for all large $m$, $M=L^{m'}$, where $\ell/b_m\to0$. For every $b>0$, $f^K(b)$ converges as $K\to\infty$ to a number $f(b)$, and, by Theorems~\YMMref{AJ}{673} and~\YMMref{AD}{607}, $f$ is real-analytic on $(0,\infty)$.
\end{proposition}
\begin{proof}[Proof sketch]
$W_{\gamma_1}$ and $W_{\gamma_{12}}$ are real, of modulus $\le1$, depend only on the bonds of the slab of $L^{K-d}$ layers that holds the block layer of $y$, and are odd under the flip $c^k$ of the class $n_{13}$ applied to all layers of that slab; Theorem~\ref{C-366-3}(6), with $2k'=\ell L^d$ slabs of that thickness and $j=k'$, gives the upper bound. By the small-field chain and the rarity of its large-field regions on the periodic torus (Proposition~\ref{prop:ward-fixed-volume-rarity}), off an event of periodic probability at most $\mathsf R(d)$ every plaquette of the block field $V$ of depth $d$ is within $2\varepsilon$ of $\mathbf 1$, $\varepsilon=\varepsilon(L^db)$, in the operator norm $|\cdot|$; let $\mathsf B$ be the event that this holds, and $\mathsf n=\ell L^d$ the number of block bonds of a straight closed block line.

On $\mathsf B$ the lemma on rearranged paths (Lemma~\YMMref{4.1}{17}), read on the block lattice of depth $d$, where a step sequence may wind about the torus, applies to every pair of paths $\Gamma,\Gamma'$ from one site whose step sequences, of a length $q$, are rearrangements of one another: their transports obey $|V(\Gamma')V(\Gamma)^{-1}-\mathbf 1|\le(1+2\varepsilon)^{q(q-1)/2}-1\le e^{\varepsilon q^2}-1$. (i) Let $\mathsf H_i=V(\gamma_i)$, the transport along $\gamma_i$ from $y$, so that $W_{\gamma_1}=\frac12\mathrm{tr}\,\mathsf H_1$, $W_{\gamma_2}=\frac12\mathrm{tr}\,\mathsf H_2$ and $W_{\gamma_{12}}=\frac12\mathrm{tr}\,\mathsf H_1\mathsf H_2$. The sequences $(e_1^{\mathsf n},e_2^{\mathsf n})$ and $(e_2^{\mathsf n},e_1^{\mathsf n})$, of the length $q=2\mathsf n$, have the transports $\mathsf H_1\mathsf H_2$ and $\mathsf H_2\mathsf H_1$, so that $|\mathsf H_2\mathsf H_1\mathsf H_2^{-1}\mathsf H_1^{-1}-\mathbf 1|\le e^{4\mathsf n^2\varepsilon}-1\le\delta_d$.

On $\mathrm{SU}(2)$, $|g-\mathbf 1|^2=2-\mathrm{tr}\,g$, and $\mathrm{tr}(ABA^{-1}B^{-1})=(\mathrm{tr}\,A)^2+(\mathrm{tr}\,B)^2+(\mathrm{tr}\,AB)^2-\mathrm{tr}\,A\,\mathrm{tr}\,B\,\mathrm{tr}\,AB-2$ (by $X+X^{-1}=(\mathrm{tr}\,X)\mathbf 1$); so the half traces $x_1=W_{\gamma_1}$, $x_2=W_{\gamma_2}$, $x_{12}=W_{\gamma_{12}}$ obey $x_1^2+x_2^2+x_{12}^2-2x_1x_2x_{12}\ge1-\frac14\delta_d^2$, and, as $|x_{12}|\le1$ and $2|x_1x_2|\le x_1^2+x_2^2$,
\[ 2W_{\gamma_1}^2+2W_{\gamma_2}^2+W_{\gamma_{12}}^2\ge1-\tfrac14\delta_d^2\qquad\text{on }\mathsf B . \]
In the case $\delta_d=0$, that of two commuting transports $\mathsf H_\mu=e^{i\theta_\mu\mathsf a\cdot\sigma}$ with one axis $\mathsf a$, this inequality is the identity
\[ 2\cos^2\theta_1+2\cos^2\theta_2+\cos^2(\theta_1+\theta_2)=1+(\cos\theta_1+\cos\theta_2\cos(\theta_1+\theta_2))^2+\cos^2\theta_2\sin^2(\theta_1+\theta_2)\ge1 . \]
(ii) Let $\gamma$ be one of the three lines, $\mathfrak s$ its step sequence, of the length $\mathsf n$ or $2\mathsf n$, $\Gamma_3$ the straight block path of $\frac{\mathsf n}2$ steps in the direction $3$ from $y$ to $\mathsf ty$, and $\mathsf v$ the transport along it. The sequences $(\mathfrak s,e_3^{\mathsf n/2})$ and $(e_3^{\mathsf n/2},\mathfrak s)$, of a length $q\le\frac52\mathsf n$, have the paths $\gamma$ followed by $\Gamma_3$, and $\Gamma_3$ followed by $\mathsf t\gamma$; so $|\mathsf v\,V(\mathsf t\gamma)\,\mathsf v^{-1}V(\gamma)^{-1}-\mathbf 1|\le e^{6.25\,\mathsf n^2\varepsilon}-1\le\delta_d$, which is the place of the numeral $7$ in $\delta_d$, and, with $|\frac12\mathrm{tr}\,X|\le|X|$, $|VV'^{-1}-\mathbf 1|=|V-V'|$ for unitary $V,V'$ and the invariance of the trace under conjugation, $|W_{\mathsf t\gamma}-W_\gamma|\le\delta_d$ on $\mathsf B$. (iii) As $\mathsf tW_\gamma=W_{\mathsf t\gamma}$ and $|W_\gamma|\le1$, $W_\gamma\,\mathsf tW_\gamma\ge W_\gamma^2-|W_\gamma-W_{\mathsf t\gamma}|$, and $|W_\gamma-W_{\mathsf t\gamma}|$ is at most $\delta_d$ on $\mathsf B$ and at most $2$ off it; so $\mathsf m_\gamma\ge\langle W_\gamma^2\rangle-\delta_d-2\mathsf R(d)$. With the weights $2,2,1$, the inequality of (i) on $\mathsf B$, $W_\gamma^2\ge0$ off $\mathsf B$ and the periodic probability of $\mathsf B$ at least $1-\mathsf R(d)$,
\[ f^K=2\mathsf m_{\gamma_1}+2\mathsf m_{\gamma_2}+\mathsf m_{\gamma_{12}}\ge(1-\tfrac14\delta_d^2)(1-\mathsf R(d))-5(\delta_d+2\mathsf R(d)), \]
which decreases in $\delta_d$ and in $\mathsf R(d)$ and equals $\frac{1599}{1600}\cdot\frac{39}{40}-\frac12=0.47439\ldots$ at $\delta_d=\frac1{20}$, $\mathsf R(d)=\frac1{40}$; $\mathsf m_{\gamma_2}=\mathsf m_{\gamma_1}$ is the symmetry of the axes, the exchange of the directions $1$ and $2$ about $y$ keeping the periodic law of the block field and commuting with $\mathsf t$. The limit in $K$ is the convergence in total variation of the laws of the block fields on $T_\ell$ (Theorem~\YMMref{D}{8}); its analyticity is Theorems~\YMMref{AJ}{673} and~\YMMref{AD}{607}, $f^K$ being at every label the mean of one bounded Borel function of the block field.
\end{proof}

\begin{proof}[Proof of Theorems~\ref{C-366-3}, \ref{C-366-4} and~\ref{C-366-5}]
In this proof A$(k)$ and B$(k)$ denote part $(k)$ of Theorems~\ref{C-366-3} and~\ref{C-366-4}. Theorem~\ref{C-366-3} is proved on one lattice. In Theorems~\ref{C-366-4} and~\ref{C-366-5}, the first sentence of B(1), the inequalities of B(8), the reflection positivity of B(7) at a fixed cutoff and Theorem~\ref{C-366-5}, under its hypothesis (H), at each cutoff follow from Theorem~\ref{C-366-3} and the properties of the averaging map; the rest of B(1) uses Lemma~\ref{lem:pure-scaling-stability}; everything else uses Theorem~\YMMref{D}{8}, and B(6), B(9)(c) and the statements on analyticity use Theorems~\YMMref{AJ}{673} and~\YMMref{AD}{607} as well, B(9)(c) with the radius and the constants of Theorems~\YMMref{AJ}{673} and~\YMMref{AD}{607}, through Lemma~\ref{lem:dilated-torus-cumulant-bound} and Theorem~\textup{\YMMref{AV}{851}} and, for its last assertion, through the normalisers at the first-passage data (Proposition~\YMMref{91.12}{708}\,(1)), the transfer principle for one law and for two laws under one normaliser (Proposition~\YMMref{91.11}{707}), the first facts on states (Section~\ref{sec:state-limit-conventions}), the device of sign averaging over a packing with its count (the packing and sign-averaging argument in the proof of Proposition~\YMMref{127.1}{848}) and the windows of first passage (Corollary~\YMMref{74.2}{616}).

\emph{Step 1: the trace and the order (A(1)).} For a sign field $\gamma$ on the bonds, $-1$ being central and Haar measure invariant, $Z_\beta[\eta\,d\gamma;F]=Z_\beta[\eta;F\circ\gamma]$: columns move, two columns over one plane cancel, and on centre-blind $F$ the places of the columns do not matter. Cut the box along $\mu$ into layers $V^\tau$ and crossings $S^\tau$, with the columns through $\mu$ in the crossing $\tau_*$:
\[ A_{\eta^n}(U)=\sum_\tau A^m_\Sigma(V^\tau)+\sum_{\tau\ne\tau_*}A_\times(V^\tau,S^\tau,V^{\tau+1})+A_\times(V^{\tau_*},S^{\tau_*},c^kV^{\tau_*+1}), \]
$A_\times$ the action of the plaquettes of one crossing; integrating the $S^\tau$ gives the trace formula. As $\mathrm{tr}(vw^{-1})=\sum_{ij}v_{ij}\bar w_{ij}$, $\prod_be^{-\beta(1-\frac12\mathrm{tr}\,V_bW_b^{-1})}$ is a series of squares of monomials with coefficients $\ge0$, so $\mathsf K=\Pi\mathsf K_0\Pi\ge0$, $\Pi$ the gauge average, strictly on gauge invariant functions if $\beta>0$ (Stone--Weierstrass). A continuous kernel $>0$ gives a simple top with a positive eigenfunction, which $C^k$ fixes. $Z[(0,m)_\mu]-Z[(k,m)_\mu]=2\sum_{e\cdot k=1}Z_{e,m}\ge0$ removes the planes through $\mu$, and the same with the time in the remaining twisted plane removes it; dually $Z_{00,m}-Z_{e,m}=\frac12\sum_{e\cdot k=1}Z[(k,m)_\mu]>0$ for $e\ne0$, every $Z[n]$ being positive; $A_{\eta^n}\le A+2\#\{\eta^n=-1\}$ gives the lower bound. $Z_{e,m}>0$: the product of the traces of straight closed lines of the directions $\nu,\rho$ with the exponents $e$ lies in the range of $P_e$, and Jensen's inequality for its spectral measure bounds $\mathrm{Tr}(P_eT_m^N)$ below. $\|T_1\|\le\|T_0\|$ follows from $\mathrm{Tr}\,T_1^N\le\mathrm{Tr}\,T_0^N$ by Lemma~\ref{C-366-6}.

A step in the direction $\nu$ moves the reversed plaquette of the section by one site, and reversing the bond that it shares with its image moves it back: $\mathcal T_\nu$ keeps $dV$ and $A^1_\Sigma$ and commutes with $\mathsf K$, so with $T_1$, and it maps the gauge invariant functions onto themselves, the step and the sign changes carrying gauge transformations to gauge transformations. For a sign field $\gamma$ on the bonds, $\mathsf S_\nu(\gamma V)=\gamma'\,\mathsf S_\nu V$ with $\gamma'(\xi,j)=\gamma(\xi+e_\nu,j)$; so the map of $\mathcal T_\nu$ is the step followed by the reversal of the bond $(\xi_*,\rho)$, and by induction $\mathcal T_\nu^q$ is the composition with the step by $q$ sites followed by the reversal of the $q$ bonds $(\xi_*-ie_\nu,\rho)$, $0\le i<q$, of the direction $\rho$ on the line of the section through $\xi_*$ normal to $\rho$.

At $q=N_\nu$ the step is the identity and these are the bonds that $C_\rho$ reverses: $\mathcal T_\nu^{N_\nu}=C_\rho$, and likewise $\mathcal T_\rho^{N_\rho}=C_\nu$. In the same way $\mathcal T_\nu\mathcal T_\rho$ and $\mathcal T_\rho\mathcal T_\nu$ differ by the reversal of the four bonds at the site $\xi_*+e_\nu+e_\rho$, the gauge transformation that is $-1$ at that site. Let $c_\nu^{(s)}$ reverse the bonds $(\xi,\nu)$ with $\xi_\nu=s$, so that $C_\nu$ is the composition with $c_\nu^{(s)}$ at $s=(\xi_*)_\nu$; then $\mathsf S_\nu c_\nu^{(s)}=c_\nu^{(s-1)}\mathsf S_\nu$, and $c_\nu^{(s-1)}c_\nu^{(s)}$ is the gauge transformation that is $-1$ at the sites of the line $\xi_\nu=s$ and $1$ elsewhere: $\mathcal T_\nu C_\nu$ and $C_\nu\mathcal T_\nu$ differ by it, and $\mathcal T_\rho C_\rho$ and $C_\rho\mathcal T_\rho$ likewise.

Gauge transformations act as the identity on the gauge invariant functions, off which $T_1=D_1\mathsf KD_1$ vanishes; for $\beta>0$ these are the closure of the range of $T_1$, as $\mathsf K$ is strictly positive on them and $D_1$ is invertible and gauge invariant. On them $T_1$, $\mathcal T_\nu$, $\mathcal T_\rho$, $C_\nu$, $C_\rho$ commute pairwise, and on the range of $P_e$ there, where $C_\nu=(-1)^{e_\nu}$ and $C_\rho=(-1)^{e_\rho}$, $\mathcal T_\nu^{N_\nu}=(-1)^{e_\rho}$ and $\mathcal T_\rho^{N_\rho}=(-1)^{e_\nu}$: the commuting unitaries $\mathcal T_\nu$, $\mathcal T_\rho$ have finitely many eigenvalues there, $N_\nu$-th roots of $(-1)^{e_\rho}$ and $N_\rho$-th roots of $(-1)^{e_\nu}$, which gives the momenta of A(1); their joint eigenspaces are invariant under the compact self-adjoint $T_1$, and an orthonormal eigenbasis of $T_1$ in each of them is a joint eigenbasis.

Finally $C^k=C_\nu^{k_1}C_\rho^{k_2}=\mathcal T_\nu^{k_2N_\nu}\mathcal T_\rho^{k_1N_\rho}$ gives the trace formula with the momenta. For the crossing relation, expand $Z^{[1]}_{e,m}=\frac14\sum_k(-1)^{e\cdot k}Z[(k,m)_1]$ and $Z[n]=\sum_{e'}(-1)^{e'\cdot k'}Z^{[2]}_{e',m'}$, $k'=(n_{12},n_{23})$, $m'=n_{13}$: the sum over $n_{12}$ forces the first entries to agree, and what is left is the display. All $Z^{[2]}\ge0$, so its modulus is at most its value at $e_2=0$, $m=0$, which is $Z^{[1]}_{e_10,0}$; the time $3$ in the place of $2$ gives $Z_{0e_2,0}$; $N_\mu$-th roots give the norms. $\varepsilon_{\rm mag}$ and $\varepsilon_{\rm el}$ are Lemma~\ref{C-366-6} for the sector traces, the top of $T_0$ being simple with its eigenfunction in the range of $P_{00}$.

\emph{Step 2: lengths (A(2)--(4)).} The class $n_{12}$ is electric for the times $1$ and $2$, so $\frac12(Z[0]\pm Z[n_{12}])=\mathrm{Tr}(\frac12(1\pm C)T_0^{N_\mu})$ with either direction of the plane as the time, and Lemma~\ref{C-366-6} applies in each length. $N_\mu^{-1}F_0$ decreases to $\log\|T_0\|\ge|\Sigma|P$, $P\ge-6\beta$ because $A\le6|N|$, and $\delta(t,\Sigma)=t[\log\|T_0\|-|\Sigma|P]+\mathsf G(t)$ gives $\mathsf G\le\delta$, the identities for $\mathsf d$ and, the bulk cancelling, $qF_0(N)-F_0(N^{[\mu]})=q\delta(N)-\delta(N^{[\mu]})$; $F_0/|N|$ minus its limit in one side is $\mathsf G/|N|$, and the last inequality of A(2) follows side after side, $\mathsf G$ not increasing in its height.

For A(3), $q(t,\lambda)$ is a power sum in $t$ and in $\lambda$, and $p\ge\|T_0\|^t\ge e^{-t\lambda s_0}$, the top eigenfunction lying in the range of $\frac12(1+C)$ and $E_0(\lambda)\le s_0\lambda$; so $\mathfrak b=q\,e^{t\lambda s_0}\le\mathfrak b(t_0,\lambda_0)^{t\lambda/t_0\lambda_0}$, $\mathfrak b\le Z[0]e^{t\lambda s_0}-1=e^\Phi-1$ and $q/p\le\mathfrak b$; $-t^{-1}\log$ and $t\to\infty$ carry the log-convexity and the submultiplicativity in $\lambda$ to $E_0$ and $E_{\rm e}$, and the monotone increments give $s_f$ and the bounds on $s_f$ and on $\sigma$. The roots of the power sums in both lengths give the monotonicity of $-(t\lambda)^{-1}\log\mathfrak b$ and its limit $s_{E_{\rm e}}-s_0=\sigma$; the lower bound on $-\log\mathfrak b(t,\lambda)$ follows from the convexity in $t$ continued from $t_0$, with the Casimir term $R(\lambda_0)\ge R(\lambda)$; and $q(2t)\le q(t)^2$ gives the bound on $\varrho$.

For A(4)(a), $g(h)=\log(q(h)/\|T_0\|^h)$ is convex in the height $h$ and $g(h)/h\to-\varepsilon(\lambda)$; its chord from $h_0$ to $h_1$, continued beyond $h_1$, lies below it, so $(h_1-h_0)\varepsilon\le g(h_0)-g(h_1)$, with $g(h_0)\le\log(e^{\mathsf G(h_0)}-1)$, the top eigenvalue lying in the even part, and $-g(h_1)\le\log(2/(1-r(h_1)))$; (a) in the two lengths, with A(3), gives the bound on $\sigma$. For (b), $\Lambda(a,b')=\log\|T_0\|$ over the section $(a,b')$ is symmetric in its two arguments, $F_0(a,b',c)=c\Lambda(a,b')+\mathsf G(c;a,b')$ with any side as the time, $\Lambda_1(b')=\lim_a\Lambda(a,b')/a$, and $\varphi(a)=\Lambda(a,b')-a\Lambda_1(b')=R(a)\ge0$ is convex and does not increase; writing $F_0(a,b',k'a)$ with the time along $a$ gives $\varphi(a)\le\varphi(a)/k'+\mathsf G(a;b',k'a)/(k'a)$, which is the bound; the same loop in the second transverse side and $\delta(M^{\times3})=\mathsf G(M;M,M)+M\varphi(M;M)+M^2[\Lambda_1(M)-MP]$ give the bound on $\delta(M,M,M)$. (c) is the left bound on $\sigma$ in A(3) with $q(t)\ge e^{-tE_{\rm e}}$, $Z[0](2t)\le\|T_0\|^tZ[0](t)$ and (b) at $k'=2$.

For (d), let $\mathsf p$ run over the plaquettes of the planar rectangle and let $\gamma_{\mathsf p}=-1$ exactly on the planar edges crossed by the straight dual path from $\mathsf p$ to the free boundary, lifted to the bonds of the directions $1,2$ at all heights: $d\gamma_{\mathsf p}=-1$ exactly on the stack of the $N'$ plaquettes over $\mathsf p$, the plaquettes through the third direction keep their weight, and $\mathrm{tr}\,U(C)$ takes the product of the signs of the enclosed $\mathsf p$. Averaging over all products of these changes, with $h_{\mathsf p}=\frac\beta2\sum_{\rm stack}\mathrm{tr}\,\partial U$, $|h_{\mathsf p}|\le\beta N'$, gives $\langle\mathrm{tr}\,U(C)\rangle=\int\mathrm{tr}\,U(C)\prod_{\rm inside}\tanh h_{\mathsf p}\,d\mu'$ with a probability measure $\mu'$. (e) is the same averaging over the reversal of one column inside each container, at fixed walls.

\emph{Step 3: how bounds travel; the overlap (A(5), A(6)).} A(5)(a): with $B_\pm=\frac12(1\pm C^k)T_m^N\ge0$, $B_+B_-=0$, $\mathrm{Tr}\,T_m^{2N}=\sum_\pm\mathrm{Tr}\,B_\pm^2\le\sum_\pm(\mathrm{Tr}\,B_\pm)^2=\frac12(x^2+y^2)$, $x=\mathrm{Tr}\,T_m^N$, $y=\mathrm{Tr}(T_m^NC^k)$. With $z_e=Z_{e,m}\ge0$ and $z_{00}\ge\|T_m\|^N$: $Z[(k,m)]=\sum_e(-1)^{e\cdot k}z_e\ge Z'-2(Z'-z_{00})$, $\sum_kZ[(k,m)]=4z_{00}$, $\mathrm{Tr}\,T_m^{2N}\le\sum_ez_e^2=\frac14\sum_kZ[(k,m)]^2$, and $Z[((1,0),m)]+Z[((0,1),m)]-Z'=Z[((1,1),m)]-4z_{11}$. (b): $q(N)=\frac12(Z[0]-Z[n])(N)=\mathrm{Tr}(\frac12(1-C^k)T_0^N)$ is log-convex in $N$ (Lemma~\ref{C-366-6}), $q=\frac12Z[0](1-r)$ with $0\le1-r\le1$, and the volumes cancel in $Z[0](N)^2/Z[0](N-j)Z[0](N+j)=e^{2\delta(N)-\delta(N-j)-\delta(N+j)}$. (c): $\mathrm{Tr}(P_ST_m^{N'})\ge\lambda^{N'}$, and $\mathrm{Tr}(P_ST_m^{N'})\le\lambda^{N'-N''}\mathrm{Tr}(P_ST_m^{N''})\le\lambda^{N'-N''}\mathrm{Tr}\,T_0^{N''}$ by Lemma~\ref{C-366-6} and A(1), with $\mathrm{Tr}\,T_0^{N'}=\|T_0\|^{N'}\Theta(N')$. (d): the first inequality is the log-convexity of $q$ between two heights, the second the monotonicity of $q(N)^{1/N}$ (Lemma~\ref{C-366-6}) with $Z[0](N)\ge\|T_0\|^N$, and the third is Lemma~\ref{C-366-17} for the class $z$ with the time $\mu$, divided by the periodic numbers, the volumes cancelling: $k'(n'+m)-(k'-1)m-(k'n'+m)=0$; the denominator $\mathsf Z=Z[(0,m')]\le Z[0]$ stands the right way. (e): put $k'=2L'+1$, $n'=L'^{j-1}(L'-1)$, $m=L'^{j-1}(L'+1)$ and $m_1=m_2=m/2$, so that $n'+m=N'$ and $k'n'+m=L'N'$, and enlarge the three sides of the cube $N'$ one after the other.

In a direction that lies in a twisted plane of $z$ the third inequality of (d) applies; for a class with one twisted plane, in the normal of that plane (c) applies at $M_1=N'$, $M_0=N'/L'$, with the exponent $L'^2/(L'-1)\le2L'+1$. A(1) passes between the three expressions, and all excesses of the auxiliary boxes are bounded against the cube $N'/L'$ by the last inequality of A(2). A(6): $T_m^\tau C^k$ has the kernel $\Psi\circ g$, so $Z[(k,m)_\mu]=\iint\Psi\,\Psi\circ g$ by A(1), and $\mu_2$ is the joint law of the layers $0$ and $\tau$; $w=2\Psi\,\Psi\circ g/(\Psi^2+\Psi^2\circ g)$ is $g$-invariant, which gives the three expressions of $r'$. $\int\sqrt{f_1f_2}\ge\int\min(f_1,f_2)=1-\mathrm{TV}$ and the Cauchy--Schwarz inequality give the bracket; on $E_c$ the integrand $\Psi\circ g/\Psi$ is $\ge c$, off it $<c$, and $\langle(\Psi\circ g/\Psi)^2\rangle_2=1$; if $gE=E^c$, $r'=2\int_E\Psi\,\Psi\circ g/Z'$; for odd $\Phi$, $\langle\Phi\rangle_2=E_{\bar\nu}[\Phi\mathsf m]$, $\bar\nu=\frac12(\mu_2+g_*\mu_2)$, and the Cauchy--Schwarz inequality with $1-\sqrt{1-x}\le x$ gives the next bracket; $g_*(p_1\otimes p_1)=p_1\otimes p_1$ gives the bound by $\mathrm{TV}(\mu_2,p_1\otimes p_1)$; the bound on $\mathcal G\cap g\mathcal G$ is pointwise; Jensen's inequality for $-\log$ gives the entropy bound; and an affinity does not decrease under a map, whereas relative entropy does not increase, while $\Phi$ intertwines the projection on the two layers with $g$.

For the odd functions, cut the period into $2k'$ slabs of $\vartheta$ layers: $Z[0]=\mathrm{Tr}\,Y^{2k'}$ and $Z[n]=\mathrm{Tr}(\mathsf ZY^{2k'})$ with the transfer operator $Y$ over $\vartheta$ layers and the flip $\mathsf Z$ in one layer, $\mathsf ZY=Y\mathsf Z$; the projections $\frac12(1\pm\mathsf Z)$ carry the weights $\frac12(1\pm r[n])$, the operator of $D'$ maps each of their ranges into the other, and H\"older's inequality for traces over the $j-1$ and the $2k'-j-1$ slabs between the two functions gives the bound.

\emph{Step 4: reflection positivity, chessboard bounds, the mixture (A(7)--(9)).} For a slab with $\alpha,\alpha'$ crossings on its two sides and the columns through $\mu$ on the side of $\alpha'$, the trace formula gives
\[ Z[(k,m)_\mu;\overline{\theta F}F]=\sum_e(-1)^{e\cdot k}\bigl\|T_m^{\alpha/2}\hat F\,T_m^{\alpha'/2}P_e\bigr\|_2^2 , \]
$\hat F$ the slab operator; this is $\ge0$ at $k=0$, $\sum_k(-1)^{e\cdot k}$ of it is $4\|T_m^{\alpha/2}\hat FT_m^{\alpha'/2}P_e\|_2^2\ge0$, and summed over the classes allowed by $\mathfrak s$ the signs add up to $\prod_j(1+(-1)^{e_j})\ge0$. For $k\ne0$, $\beta>0$, write $e^{\frac\beta2\mathrm{tr}\,u}=\cosh\cdot(1+\tanh)$ at $\frac\beta2\mathrm{tr}\,u$, and let one plaquette of the flipped row of one crossing carry $\chi_J/\cosh$, $\chi_J$ the character of spin $J$, and the others a factor concentrated at $\mathrm{tr}\,u=0$, through which integer spins alone pass; the part even under the flip falls like $c^J/(2J)!$, while the odd part is a Fourier coefficient of a function with the pole of $\tanh(\beta\cos\theta)$ and is $\ge e^{-cJ}$ for infinitely many $J$: the form is negative. $\sum_nZ[n]\le8Z[0]$ is A(1).

For the last assertion of A(7), $Z[(k,m)_\mu;F]=\mathrm{Tr}(\hat FT_m^{N_\mu-\vartheta+1}C^k)$ with one slab, and the Cauchy--Schwarz inequality for traces. On a board a reversed plaquette in a mirror is a unitary flip commuting with $\mathsf K^{1/2}$, so $Z_\beta[\eta;\prod F_\kappa]=\mathrm{Tr}(C_r\mathfrak Y_r\cdots C_1\mathfrak Y_1)$, and H\"older's inequality for traces does not see unitaries; the adjoint of a slab operator is the operator of the mirror copy of the slab that carries the complex conjugate function, so the norms of the slab operators are fold integrals with the copies complex conjugated in the cells odd relative to the given one (equation~\YMMref{(1749)}{680}), which are $\ge0$.

Put each transverse electric-flip cut on a wall of the board. A cell function reads only bonds with both endpoints in that cell; the flip changes only bonds crossing the wall, so it fixes every inserted cell function, including a complex one. Applying the flip on all layers of a slab preserves its action and Haar measure, hence its inserted slab operator commutes with the flip. The crossing operator and its square root commute as well. Thus each $\mathfrak Y_i$, and each $\mathfrak Y_i^*\mathfrak Y_i$, commutes with the electric-sector projections. In $Z_w$ the flips add up to $\mathsf D^w_m=\sum_e\hat w(e,m)P_e\ge0$, which commutes with all $\mathfrak Y_i$: $Z_w[\prod F_\kappa]=\sum_m\mathrm{Tr}(\mathsf D^w_m\mathfrak Y_r\cdots\mathfrak Y_1)$, and $|Z_w|=Z[0]\sum w(n)r[n]$ lies between $w(0)Z[0]$ and $8w(0)Z[0]$ by A(1). Comparing $Z_w[\mathrm{Fold}\,F]$ with $\sum_nw(n)Z[n;\mathrm{Fold}\,F]$ and with its term $n=0$ gives the two bounds of $\|\cdot\|_w$. A covariant function of a cell is $G(U,s)=F(\gamma U)$ with $d\gamma=s$ on the cell and $F$ invariant under the gauge transformations with values $\pm1$, and $\mathsf M_n[\prod G_\kappa]=Z[n;\prod F_\kappa]$ with the columns in the mirrors. The functions of the cells that cancel the weight of the rods of plaquettes through the centres of the cells make $Z[n;\mathrm{Fold}\,F]=Z[0;\mathrm{Fold}\,F]$ for every $n$, which gives equality; the functions that remove the reversed columns from the weight give the failure of the fold in a class, since $Z[n]<Z[0]$.

In a class: H\"older's inequality with the first bound for $|G|^q$; the density bound by the first bound for $h_n^{\mathsf N/(\mathsf N-1)}$ and the entropy inequality, and Pinsker's inequality. Without a board, one mirror and the Cauchy--Schwarz inequality for $\mathsf M_{(k,m)}[\overline{\theta G}G']=\mathrm{Tr}(Y_m[G]^*C^kY_m[G'])$ give $r[n]|\mu_n(G)|\le\langle\overline{\theta F}F\rangle^{1/2}$ when $n=(k,0)_\mu$, that is, when $\mu$ lies in every twisted plane of $n$; the same for $G-\langle F\rangle$ gives the two-sided bound; a second mirror gives the quarter; and the event $\{\prod_{\text{plane}}s=-1\}$ of a plane normal to $\mu$ inside the half has periodic probability $0$ and probability $1$ in a class twisted in that plane.

For pairs, summing over the signs inside each half gives $\mathbb M_{\delta_{(k,m)_\mu}}[\overline{\theta G}G']=\mathrm{Tr}(Y_m[G]^*C^kY_m[G'])$, and $\sum_kw((k,m)_\mu)C^k=\sum_e\hat w(e,m)P_e\ge0$. On $a_K\mathbb Z^3$ a closed sign field is $d\chi$, $\chi$ on the bonds, closed sign fields being exact on every box of sites: a class is the image of its measure under $U\mapsto\chi\,\mathrm{per}_\ell U$, the plain Wilson measure of a field periodic up to gauge transformations with values $\pm1$, and its law on the invariant local functions depends neither on $s$ nor on $\chi$; the proof of Corollary~\ref{IV-18-18} carries over, and for other functions the columns are put on the edges of the board.

Its clause (e) holds with $p^w_\ell$: with the time $\mu$ and $\mathsf n_\mu$ cells in the direction $\mu$, $Z_w[\mathrm{Fold}\,F]=\sum_{e,m}\hat w(e,m)\,\mathrm{Tr}\,(P_eB_mP_e)^{\mathsf n_\mu/2}$, where $B_m=\mathfrak Y_m^*\mathfrak Y_m\ge0$ is the operator of a slab of cells and its mirror copy and commutes with $P_e$, the same operators for every number of cells in the direction $\mu$. For $\mathsf n_\mu=q\mathsf n'_\mu$, $\mathrm{Tr}\,B^{q\mathsf n'_\mu/2}\le(\mathrm{Tr}\,B^{\mathsf n'_\mu/2})^q$ for $B\ge0$, $\omega\le\omega^q$ for $\omega=\hat w(e,m)/\underline w\in\{0\}\cup[1,\infty)$, and $\sum_ix_i^q\le(\sum_ix_i)^q$ for $x_i\ge0$; so $Z_w[\mathrm{Fold}\,F]/\underline w$ over $\mathsf n_\mu$ cells is at most the $q$-th power of the same number over $\mathsf n'_\mu$ cells, and its root of the order of the number of cells does not increase when the board grows in one direction.

The three directions in turn carry this from the board of $\Lambda_{K,\ell'}$ to that of $\Lambda_{K,\ell}$; at $F=1$ it is $p^w_\ell\le p^w_{\ell'}$, and the inequality at $F$ and at $F=1$ gives $\|F\|_{w,\ell}\le\|F\|_{w,\ell'}e^{\mathsf l^3(p^w_{\ell'}-p^w_\ell)}$. At $w\equiv1$, $\hat w(e,m)\in\{0,4\}$, and the same holds without the division by $\underline w$. A(9): for a closed $s=\eta^nd\gamma$, $\gamma$ is fixed up to the closed sign fields on the bonds and $e^{-\beta A_s}(U)=e^{-\beta A_{\eta^n}}(\gamma U)$; integrating against a centre-blind $F$ gives the first identity, $2^{2|N|-2}$ being the number of the closed sign fields of one class. The mean of $\sum_ne^{-\beta A_{\eta^n}}$ over all sign changes $\gamma$ is $\hat{\mathsf w}_{\rm mix}$, which is therefore $\le8$, and centre-blind, gauge invariant and invariant under the isometries.

Writing $e^{-\beta A_s}=e^{-\beta\bar A}\prod_{p:\,s_p\ne{\rm sign}\,t_p}z_p$, the sum over the closed $s$ becomes the sum over the sets $\Gamma=\{p:s_p\ne{\rm sign}\,t_p\}$ with boundary $\mathsf M(U)$. At the constant field $U^0(x,\mu)=\cos\alpha+\sin\alpha\,i\sigma_\mu$, $\sin^2\alpha=2^{-1/2}$, all $t_p$ vanish, and a sixth mixed derivative of $\log\mathsf w_{\rm mix}$ along a hexagon of edges of one cube, two of which lie in no common plaquette, has the modulus $\beta^6(2\sqrt2-2)^{1/2}=0.91018\,\beta^6$, while a sum of functions of single plaquettes has none. On one cube the plaquette variables of the faces, carried to a common corner, multiply to $\mathbf 1$; if the angles add up to less than $\pi$, the product of the signs of the traces is $1$, and a cube with six angles $\frac\pi6$ and the product $-1$ shows that $\pi$ is sharp.

On $\mathfrak S$ the sign field $s(U)$ is locally constant and closed, $\mathsf w_{\rm mix}=e^{-\beta A_{s(U)}}\Xi(z)$ with $\Xi(z)=\sum_\Gamma\prod_{p\in\Gamma}z_p$ over the closed sets $\Gamma$, and the expansion of $\log\Xi$ in connected sets at $|z_p|\le\frac1{200}$ gives the $E(X)$, which vanish unless $X$ is connected and a union of vortices, connected for the relation of A(9): two plaquettes are neighbours if they are faces of one unit cube, that is, dual bonds with a common dual site. For this relation the faces of one cube are pairwise neighbours, so each cube meets one connected component of a closed set only, the components of a closed set are closed, the closed sets are the unions of vortices pairwise apart, and $\Xi$ is a gas of vortices; for a vortex $X$ that holds no smaller one the only clusters with the support $X$ are the repetitions of $X$, whose terms add up to $\log(1+\prod_{p\in X}z_p)$.

The last bound is the smallness, on $\mathfrak S_n$, of the weights $\sum_{s\in[n']}e^{-\beta A_s}/\mathsf w_{\rm mix}$ of the other classes, whose sets $\Gamma$ wind about the torus and so have at least $\min_\mu N_\mu$ plaquettes.

\emph{Step 5: one inequality; loops (A(10), A(11)).} Under the hypothesis $D_2(l_0)\le\frac1{20}$ of A(10), one inequality between the partition functions of the periodic cubes $l_0$ and $2l_0$, Lemma~\ref{C-366-9} gives $18D_2(2^il_0)\le(18D_2(l_0))^{2^i}$, hence $\delta(s,u,w)\le0.14108\,(suw/l^3)D_2(l)$ for $s,u,w\ge l\in\{2^il_0\}$ and $\mathsf r_h(l,l)\le\frac1{18}e^{-m_*h}$. Lemma~\ref{C-366-10} then gives, for $\bar\varepsilon\ge\varepsilon_{\rm v}(l)$ and an integer $h\ge2l$ with $e^{-m_*h}\le(h\bar\varepsilon)^4$, $h\bar\varepsilon\le\frac1{10}$ and $2(h/l)D_2(l)\le\frac14$,
\[ 2l\,\varepsilon_{\rm v}(2l)\le4.84\,(h/l)^3(l\bar\varepsilon)^4, \]
whose iteration from $l_1$ gives $l\varepsilon_{\rm v}(l)\le e^{-\sigma_*l^2}$; the bound $\sigma_*\ge\min\{\frac16m_*^2,m_*/l_0\}$ is the definition of $l_1$. With $2\zeta\le h\varepsilon_{\rm v}+\mathsf r_h$ and the doubling of the short sides (Lemma~\ref{C-366-10}) this is the bound on the boxes $(t,\lambda,h)$ and on the cubes; the bound on $\Lambda_-(N)$ is the same chain for the norm of $T_0$ on the odd functions, which is $\|T_0P_{10}\|$ by A(1). The loop bounds are Lemma~\ref{C-366-10} on the even cubes $N$, the factor $8^{T/N}$ tending to $1$ as $N\to\infty$; loops made of rows go the same way, row by row, with the norms of the transfer operators with static charges, the excess entering through $e^{1.13D_2(N/2)}$.

Uniqueness and convergence are Lemma~\ref{C-366-16}. Under the hypothesis $\mathsf e(l)\le0.069$, Lemma~\ref{C-366-9} gives $\theta\le0.8901$, $\mathsf e(4l)\le\theta^4/12.9$ and $D_2(8l)\le20\,\mathsf e(4l)^2\le20\theta^8/166.41\le0.047355$; for a cube $N\ge2l$ and the largest $M\in\{2^jl\}$, $j\ge1$, with $M\le N$, $\delta(N^{\times3})\le(N/M)^3\delta(M^{\times3})\le8\cdot0.01717\,\theta^{M/l}\le0.13736\,\theta^{N/2l}$; and $1-r[n]\le3(1-e^{-\mathsf G})\le3\delta$ (Lemma~\ref{C-366-9}), with $3\cdot0.13736\le0.4121$. A(11): $\nu$ is reflection positive at site mirrors and at bond mirrors, the exact statements on periodic Gibbs states and translation invariance comparing it with the periodic tori, where A(7) holds at $k=0$.

For a rectangle of height $T$, the paths from a point of the mirror that run $k$ bonds up, $R$ across and $k$ down, with $k=h+1,h$ at a site mirror for $T=2h+1$ and $k=h,h-1$ at a bond mirror for $T=2h$, form half loops whose Gram matrix has the entries $2\mathsf m(R,T+1)$, $2\mathsf m(R,T)$, $2\mathsf m(R,T-1)$; its positivity is the log-convexity. So $\mathsf m(R,T)/\mathsf m(R,T-1)$ does not decrease in $T$, $\mathsf m(R,T)\ge\mathsf m(R,1)^T\ge\mathsf w^{RT}$ by the symmetry, $V$ exists, is concave in $R$ by the log-convexity in $R$, and $V(R)/R$ does not increase; the bounds on $\mathsf w$, on $\mathsf m$ for $R,T\ge2$ and on $V(1)$, $V(R)$ come from integrating the bonds of the loop in their plaquettes.

\emph{Step 6: purity; small cubes; constructions (A(12)--(14)).} A(12): $\mathsf P=\mathsf Y^*\mathsf Y\ge0$ commutes with $C^k$; with $B_\pm=\frac12(1\pm C^k)\mathsf P\ge0$ and $d=\mathrm{Tr}(C^k\mathsf P)/\mathrm{Tr}\,\mathsf P\in[0,1]$,
\[ \mathrm{Tr}\,\mathsf P^i=\sum_\pm\mathrm{Tr}\,B_\pm^i\le\sum_\pm(\mathrm{Tr}\,B_\pm)^i=(\mathrm{Tr}\,\mathsf P)^i\sum_\pm\bigl(\tfrac{1\pm d}2\bigr)^i; \]
convexity in $d$ gives $d\ge(\mathfrak q_i-2^{1-i})/(1-2^{1-i})$, at $i=2$ the display gives $d^2\ge2\mathfrak q_2-1$, and as $i\to\infty$ it gives $\|\mathsf P\|\le\frac12(1+d)\,\mathrm{Tr}\,\mathsf P$; and $r[n]\ge Z[n;F\otimes\theta F]/Z[0]=d\,w_F$. For block fields, the averaging map commutes with the flip in a wall of the blocks of all orders, $L$ being odd. A(13): with the time $\mu$, $Z_{e,0}/Z[0]=\frac14\sum_k(-1)^{e\cdot k}r[(k,0)_\mu]$, and the three classes with $k\ne0$ have the quotients $r_1,r_1,r_2$: the first bound; $e^{-\mathsf G}=\|T_0\|^N/Z[0]\le Z_{00,0}/Z[0]$; $D_q(N)$ is a sum of defects of $q$-fold enlargement, the first $q^2[qF_0(N^{\times3})-F_0(qN,N,N)]\ge q^2(q-1)\mathsf G(N)$; the bound on $Z[0](N^{\times3})^2/Z[0](2N,N,N)$ is A(5)(a); if $r_1\le\frac13$, each sector trace at the height $N$ is at most $\frac12Z[0]$, and with $\mathrm{Tr}(P_eT_0^N)\le(\mathrm{Tr}\,P_eT_0^{N/2})^2$ this gives $(\mathrm{Tr}\,T_0^{N/2})^2\ge2\,\mathrm{Tr}\,T_0^N$, that is $\mathsf e(N/2)\ge\log2$; and $\Phi\ge F_0-N\log\|T_0\|=\mathsf G(N)$ on the cube gives the bound on $\mathfrak b$.

A(14)(a): there are two Haar-preserving, gauge covariant bijections $\Phi_\pm$ of the fields, moving only the bonds of the directions $\mu,\nu$ inside a tube of radius $h'$ about the reversed column, which spread the flux of the column over the cross-section of the tube; pointwise, $\frac12[A_\eta(\Phi_+U)+A_\eta(\Phi_-U)]-A(U)$ is at most a sum over the layers of the squares of the spread curvature, of order $(h'-\frac12)^{-2}$ in total, and of the actions of ladders of plaquettes of length of order $h'$ weighted by the angles of rotation. As $\Phi_\pm$ keep $dU$, $Z[\eta]=\int e^{-\beta A_\eta(\Phi_\pm U)}\,dU$ for both signs; the arithmetic mean of the two integrands is at least their geometric mean, and Jensen's inequality in the periodic measure gives $\log r\ge-\beta\langle\frac12[A_\eta(\Phi_+U)+A_\eta(\Phi_-U)]-A(U)\rangle$.

For $\beta\ge1$ the mean actions $\bar a_{12},\bar a_{13},\bar a_{23}$ of the plaquettes of the three orientations in the periodic measure have $\bar a_{12}+\bar a_{13}+\bar a_{23}\le\lambda(\beta)/\beta$, by $\log Z[0]\ge-|N|[\frac92\log\beta+30]$ and the convexity of $\log Z[0]$ in $\beta$; and the mean action of a ladder of $1\times k$ plaquettes of one orientation is $k$ times the mean of that orientation. The second form is the choice of $h'$ of the order of $(\beta/\lambda(\beta))^{1/3}$, and for several twisted planes the tubes are disjoint. (b): $\log Z[0]\ge-|N|[\frac92\log\beta+30]$ by an elementary estimate; the chessboard estimate for $e^{\frac\beta2(1-\frac12t_p)}$ in one cell bounds its periodic mean by $\beta^{121.5}e^{810}$, so up to a factor $2$ in the partition function of the cube all plaquettes lie within $(130\pi^2\log\beta/\beta)^{1/2}$ of $\mathbf 1$ once $\log\beta\ge147.6$; cut the cube into eight free blocks, fix the gauge on a comb in each, and bound the action from below by a Gaussian whose determinant is the product over the inner Neumann momenta, the blocks having a side of the order $\beta^{1/6}$.

\emph{Step 7: the window (B(1)).} The first sentence of B(1) is A(1) and A(2) at $N_\mu=s_\mu/a_K$, $\beta=L^Kb$, with $|s|p_\infty=|N|P$; $\delta^0_K\ge0$ is the definition of $P$. On the cubes $2L^m$, $m\ge m_0(b)$, $b\ge\gamma^{-2}$, Lemma~\ref{lem:pure-scaling-stability} places $\log Z_K[0]$ within $C_*(b)\ell^3$ of a vacuum energy whose density varies boundedly from cube to cube, no constant depending on $K$ or $m$; the monotone density squeezes a box with sides $\ge2L^{m_0(b)}$ between two cubes, so $|\log Z_K[0](s)-\epsilon_K|s||\le C_*(b)|s|$ and $p_\infty\ge\epsilon_K-C_*(b)$, whence $\delta^0_K\le2C_*(b)|s|$ and, on the cubes, the explicit bound. For a compact set $\mathsf J\subset[b_-,b_+]$ of labels and $s_{\min}>0$, all labels of $\mathsf J$ and all sides $\ge s_{\min}$ are brought into this range at once by
\[ Z_K[n](s;b)=Z_{K-j}[n](L^js;L^jb), \]
with the least $j$ for which 
\[
L^jb_-\ge\gamma^{-2}
\]
 and 
\[
L^js_{\min}\ge2L^{m_0(L^jb_+)},
\]
which exists because 
\[
L^{-j}2L^{m_0(L^jb_+)}\to0
\]
: for $b\in\mathsf J$, $L^jb\ge\gamma^{-2}$ and, $m_0$ not decreasing, $L^js_{\min}\ge2L^{m_0(L^jb)}$; and, $C_*$ not decreasing, $C_*(L^jb)\le C_*(L^jb_+)$, so that $C=L^{3j}C_*(L^jb_+)$ serves for all $b\in\mathsf J$ and $K>j$. The bound on 
\[
\mathrm{Tr}\,T_m^{t/a_K}/\|T_0\|^{t/a_K}
\]
 follows from $\mathrm{Tr}\,T_1^M\le\mathrm{Tr}\,T_0^M$ and $\mathsf G\le\delta$. A lattice with $N_\mu/\beta$ in a compact set and $\beta$ large is a box of a cutoff $K$ at a label in $[1,L)$ with sides in a compact set, where these bounds hold uniformly, by the above at $\mathsf J=[1,L]$; for bounded $\beta$, $\delta^0\le6\beta|N|$. On a box with $L$-adic sides the remaining cutoffs are finitely many, each with a finite $\delta^0_K$ and $r_K[n]>0$, and $\delta^n_K-\log r_K[n]=\delta^0_K$ is bounded.

\emph{Step 8: floor and two cutoffs on mixed boxes (B(2)(a), B(4)).} After a dilation, $s$ is a centre sector of some $T_\ell$ (Lemma~\ref{C-366-7}), and the chain of $T_\ell$ restricted to the fixed fields is a chain for $Z^z_K(s;b)$. At the level $i$ the fluctuation integral over the fixed fields is, by Proposition~\ref{C-366-8} in the coordinates (using \textbf{(Hess)} for the full fluctuation Hamiltonian, and with least displacement $s_{\min}L^{K-i}$ in the units of that level), the $n$-th root of that of the cover up to $\epsilon_i\le C|s|L^{3(K-i)}e^{-cs_{\min}L^{K-i}}$. With $E(d)=\sum_{i\le K-d}\epsilon_i$, bounded independently of $K$, $\Sigma$ the small-field mass of the cover at the depth $d$ and $C_0$ formed from the constants,
\[ \log Z^z_K(s;b)\ge\frac{\log\Sigma}n-E(d)-|s|L^{3d}\bigl[C_0+\tfrac92\log(L^db)\bigr],\qquad\Sigma\ge\tfrac12Z_{K,\ell},\qquad\frac{|s|}{\ell^3}\log Z_{K,\ell}\ge|s|\,p_\infty(K,b), \]
the second when the large-field part of the cover is at most $\frac12$, the third by A(2); finitely many $K$ remain, at which the integrand is positive. Against the small periodic cube $\ell_\flat(b)$, box and cube being two sectors of one torus, the same chain gives the bound with $C_G$: the cover drops out. For two cutoffs, the small-field masses at $K$ and $K+1$ are compared pointwise at fixed fields as on $T_\ell$, the sector taking $|s|/\ell^3$ times the error of the cover; large fields are rare in the box, since an event of a slab passes by the Cauchy--Schwarz inequality at mirrors, with the power $\frac18$, to the small periodic cube $\ell_\flat(b)$, where Proposition~\ref{prop:ward-fixed-volume-rarity} controls it; box and cube are two sectors of one torus, whose partition function cancels.

The rate $L^{-K/10}$ and the uniformity in the label are those of the cover, with the arithmetic of the constants, which makes the condition on $K$ void for $|s|\le(\log b)^{2\mathfrak p_0-1}/C_V$. For two classes the terms $\frac{|s|}{\ell^3}\Delta_K$ cancel in the difference of the two laws, and the increments are summable. So, for a compact set $\mathsf J$ of labels, B(4) at $\mathsf J$ for $z$ and for the class $0$, for which $s$ is mixed too, gives $|\log r_{K+1}[z]-\log r_K[z]|\le2CL^{-K/10}$ for $K\ge K_0(\mathsf J)$; the first sentence of B(1) bounds $r_K[z]$ from below, uniformly on $\mathsf J$, at the finitely many $K<K_0$ at which $s$ is a box of the cutoff $K$; and as $\delta^0_K\ge0$, $\log Z^z_K(s;b)-|s|\,p_\infty(K,b)=\log r_K[z]+\delta^0_K\ge\log r_K[z]$: the floor of B(2)(a) holds with $C(s,b)$ bounded on $\mathsf J$, and $\inf_Kr_K[z]>0$ uniformly on $\mathsf J$.

On the boxes with sides in $2L^{\mathbb Z}$ the left member is the increment of $\log r_K[z]$, which tends to $0$ by B(2)(c) and B(3), plus the left member of the periodic class, for which every such box is mixed.

\emph{Step 9: compactness, no zero, the full sequence (B(2)(b), B(3)).} By B(1) at a compact set $\mathsf J$ of labels that holds the $b_j$, $i\,e^{-t_0E^K_i}\le e^{2Ct_0l_1l_2}$: Lemma~\ref{C-366-6}(b) applies, so every sequence $(K_j,b_j)$, $b_j\to b$, has a subsequence along which all levels converge, the sector traces divided by $\|T_0\|^{t/a_K}$ tend to $\mathcal W(e;m)(t;l)=\sum_ie^{-tE^{m,e}_i}$, holomorphic in $\mathrm{Re}\,t>0$, and $v_{(k,m)_\mu}=\sum_e(-1)^{e\cdot k}\mathcal W(e;m)/\mathcal D$, $\mathcal D=\sum_e\mathcal W(e;0)\ge1$, in each of the three expressions. The order of A(1) and A(5)(a) pass to the limit, and so does the crossing relation, for the sector weights $w^{[\mu]}_{e,m}=\frac14\sum_k(-1)^{e\cdot k}v_{(k,m)_\mu}=\mathcal W(e;m)/\mathcal D$, each the sector trace divided by the periodic trace of its own time and section: on each lattice the relation of A(1) divided by $Z[0]$ is the same relation between the numbers $Z^{[\mu]}_{e,m}/Z[0]=\frac14\sum_k(-1)^{e\cdot k}r[(k,m)_\mu]$, which converge to the $w^{[\mu]}_{e,m}$. A weight $w=\mathcal W(e;m)/\mathcal D$ vanishes at one height if and only if at all; it is real-analytic in a side $x$ of the section, and if it vanishes at $x_0$ then, for $q=L^p$ and $\mathsf c_q\ge\mathcal D(t/q;l_x)$ from B(1),
\[ w_t(x)\le\mathcal W(t;l_x)\le\mathcal W(t/q;l_x)^q\le\bigl[\mathsf c_q\,w_{t/q}(x)\bigr]^q=O(|x-x_0|^q), \]
so $w\equiv0$ in that side, and, side after side, on all boxes: this is the alternative of B(3), which passes to the limits of Step~10, and for a class with one twisted plane, whose quotient is the sum of its four sector weights with the time normal to its plane, it holds for $v_z$. With the floor on the box $(1,1,1)$, which is mixed for every class, this shows that the space of $H_{l,1}$ is not $\{0\}$ and, with the floor along the normal of A(5)(c) and compactness, B(2)(b); with B(5), which makes each electric weight on a small cube at least $\frac14(1-2C_\square t)$, each electric sector is not $\{0\}$, so $v_z<1$. On boxes with sides in $L^{\mathbb Z}$, B(4) makes the increments of $\log r_K[z]$ and of $L\log Z_K[0](s)-\log Z_K[0](s^{\mu+})$, $s^{\mu+}$ the box with $s_\mu$ multiplied by $L$, summable uniformly on compacts of labels, the terms $\Delta_K$ cancelling in the second: this is (R1) of Proposition~\ref{C-366-20} on the ray of the boxes with sides in $L^{\mathbb Z}$, and Proposition~\ref{C-366-20} gives the convergence of all levels along the full sequence, of $r_K[z]$ on every $L$-adic box, of $\delta^0_K$ and of the differences of the norms.

The same holds along every sequence $(K_j,b_j)$ of cutoffs and labels with $K_j\to\infty$ and $b_j\to b\in(0,\infty)$, in the place of the full sequence at the fixed label $b$; these sequences form a class that is closed under subsequences. (R1) holds along each of them: each $r_K[z](s;\cdot)$ and each difference $L\log Z_K[0](s)-\log Z_K[0](s^{\mu+})$ is continuous in the label at a fixed $K$, and their increments in $K$ are summable uniformly for the labels in a compact set, so that they converge uniformly on compacts of labels to limits that are continuous in the label, and so converge along $(K_j,b_j)$. The proof of Proposition~\ref{C-366-20} then runs along such a sequence as it stands: its only input, the window, is uniform for the labels in a compact set (B(1): the supremum of $\delta^0$ over the lattices with $N_\mu/\beta$ in a compact set is finite), and Lemma~\ref{C-366-6}(a), (b) and the turning of the box do not read the label.

So all levels, $r_K[z]$ on every $L$-adic box, $\delta^0_K$ and the differences of the norms converge along every such sequence. Two such sequences with one limit label $b$ interlace into a third one, along which the limits exist: the limits depend on $b$ alone. They are continuous in $b$: for $b_i\to b$ a choice of $K_i\to\infty$ with the values at $(K_i,b_i)$ within $1/i$ of the limits at $b_i$ gives a sequence of the class. And the convergence is uniform for the labels in a compact set, since otherwise a compact set of labels would hold a sequence $(K_j,b_j)$ with $b_j\to b$ whose values stay at a fixed distance from the limits at the labels $b_j$, which tend to those at $b$, against the convergence along it. $E_{\rm mag}$ and $E_{\rm el}$ are Lemma~\ref{C-366-6} for the limit lists, the space of $H_{l,1}$ not being $\{0\}$.

On the cubes, $L^3\log Z_{K,\ell}-\log Z_{K,L\ell}$ is a difference of periodic numbers of two cubes in which the terms $\Delta_K$ of B(4) cancel, so $D(t)$ exists; it depends on $t$ alone by the argument of Step~10; $\delta^0_K(\ell^{\times3})$, telescoped over the cubes $L^r\ell$, gives $\delta(t)=\sum_rL^{-3(r+1)}D(L^rt)$; and the behaviour of $D$ as $t\to0$ is the asymptotics of Proposition~\ref{C-366-13}(b), (e) for the periodic cubes, the bulk cancelling, carried to every cutoff by B(4) as in Step~11. By Theorems~\YMMref{AJ}{673} and~\YMMref{AD}{607}, through Theorem~\textup{\YMMref{AD}{607}}, the periodic quotients are normal families in the label, and B(6) gives the analyticity.

\emph{Step 10: every lattice; two blocking factors (B(2)(c), B(3)).} By Lemma~\ref{C-366-11} and B(1) the quotients are equicontinuous in the sides for $N_\mu/\beta$ in a compact set; with Step~9, in its form for the sequences $(K_j,b_j)$ with a moving label $b_j=L^{-K_j}\beta_j\to b$, this gives a continuous $V^L_z(s;b)$, the limit along every sequence with $L^{-K_j}N_j\to s$, $L^{-K_j}\beta_j\to b$, with the Lipschitz bound of B(3); in each side, at real values of the others, it is $\sum_e(-1)^{e\cdot k}\mathcal W(e;m)/\mathcal D$, a quotient of Laplace transforms, and real-analytic. The constants may be chosen with the blocking factors $L_1=L$ and $L_2=L+2$, and $\log L_1/\log L_2\notin\mathbb Q$; so for $b_1,b_2>0$ there are $K_j,M_j\to\infty$ with $K_j\log L_1-M_j\log L_2\to\log(b_2/b_1)$, and $\beta_j=L_1^{K_j}b_1$, $N_j\approx\tau\beta_j$ is a sequence of both scales, of the label $b_1$ at the scale $L_1$ and of the moving label $L_2^{-M_j}\beta_j\to b_2$ at the scale $L_2$, to which Step~9 applies in its form for sequences:
\[ V^{L_1}_z(\tau b_1;b_1)=V^{L_2}_z(\tau b_2;b_2), \]
so $V_z(s;b)=v_z(s/b)$, every sequence with $\beta_j\to\infty$, $N_j/\beta_j\to\tau$ converges ($\beta_j=L^{K_j}b_j$, $b_j\in[1,L)$), and the limits do not depend on the choice of the constants. For $s\in(2L^{\mathbb Z})^3$, $V_z(s;b)=V_z(\frac12s;\frac12b)>0$, the box $\frac12s$ being mixed for every class; positivity of the limit and of each $r_K$ is the floor, uniform for the labels in a compact set, since the convergence is uniform there (Step~9), the limits are continuous and positive, and the first sentence of B(1) bounds the finitely many remaining cutoffs; and an infimum $0$ over the cubes with $N/\beta\in I$, or over the lattices with $N/\beta$ in a compact subset of $\{v_z>0\}$, would give, along $\beta_j\to\infty$, a point of that set with $v_z=0$. The same for the levels gives their extension to real sections and their homogeneity. The weights of the mixture converge with the $r_K[n]$, and $1\le\sum_nv_n\le8$ with $v_0=1\ge v_n$, so that their limits lie in $[v_n/8,\frac12]$ for $n\ne0$.

\emph{Step 11: small cubes (B(5)).} (a), the ceiling. By A(1), $r_2,r_3\le r_1$, and by the symmetry of the axes it is enough to bound $r[n]$ for $n=n_{13}=((1,0),0)_3$, whose twisted plane holds the time $3$, on the cubes $\ell=2L^m$. A dilation by a power of $L$ keeps $t$ and the lattices and brings $(\ell,b)$ into the range $b\ge b_0$, $MR(b)\le L^m$, $\ell\le2ML^2(\log b)^{r_\star}$ of Lemma~\ref{C-366-12}(b), where small $t$ is large $b$. In a block layer normal to $3$ of the block field of depth one, with $\bar h_1(y)$, $\bar h_2(y)$ the products of the block field along the straight closed block lines from $y$ in the directions $1$ and $2$ (as in Definition~\ref{C-366-15}), let $\bar s_1$ and $\bar s_+$ be the means over the sites $y$ of the layer of $\frac12\mathrm{tr}\,\bar h_1(y)$ and of $\frac12\mathrm{tr}[\bar h_1(y)\bar h_2(y)]$. The flip of $n$ changes the signs of both; on a flat field with the holonomies $e^{i\varphi_\mu\mathsf a\cdot\sigma}$ they are $\cos\varphi_1$ and $\cos(\varphi_1+\varphi_2)$, which vanish together exactly at the fixed points $\varphi_1=\pm\frac\pi2$, $\varphi_2\in\{0,\pi\}$ of the flip $\varphi_1\mapsto\varphi_1+\pi$, taken up to $\varphi\mapsto-\varphi$. Let
\[ F=\mathbf 1\bigl\{\max(|\bar s_1|,|\bar s_+|)<m_0t^{1/2}\bigr\}, \]
with an absolute $m_0$, be read at the block layer next to a block mirror normal to $3$ at the cutoffs $K\ge K_\sharp$ of Lemma~\ref{C-366-12}(b), and at the bare layer next to a mirror between two layers below $K_\sharp$; $F$ is gauge invariant, Borel with values in $[0,1]$ and flip-blind, and its mirror image $\theta F$ is $F$ read at the layer on the other side of the mirror. By \textbf{(Flux-window)}, with absolute $C_2$ and $\epsilon_F$, at every cutoff and with the columns of $n$ away from the layers of $F$ and $\theta F$: \textbf{(F1)} $\mu(F)\le C_2t$, and \textbf{(F2)} $Z[n;1-F]\le\epsilon_F\,t\,Z[0]$.

Both estimates are part of \textbf{(Flux-window)}. In particular, (F1) is a bound for a shrinking event of the block means, and does not follow from the uninserted periodic Laplace estimate of Proposition~\ref{C-366-13}. The sign reversal of the plane means under the twist likewise does not by itself establish (F2).

Now
\[ r[n]=\frac{Z[n;\theta F\,F]}{Z[0]}+\frac{Z[n;1-\theta F\,F]}{Z[0]} . \]
The first term is bounded through the overlap representation of the class with flux, in the form that A(6) takes for a function of a slab and its mirror image, the trace formula of Step~4 with $\alpha,\alpha'$ as there: $Z[n;\theta F\,F]=\sum_e(-1)^{e\cdot k}\|T_0^{\alpha/2}\hat F\,T_0^{\alpha'/2}P_e\|_2^2\le Z[0;\theta F\,F]\le Z[0;F]=\mu(F)\,Z[0]$. In the second, $1-\theta F\,F\le(1-F)+(1-\theta F)$, and $Z[n;1-\theta F]=Z[n;1-F]$ by the reflection, the columns being moved away from the layers of $F$ (Step~1). So $r_1\le(C_2+2\epsilon_F)\,t$ at every cutoff, with the absolute $C_\square=C_2+2\epsilon_F$; $t_0$ and both bounds are as specified in (Flux-window). (b) and (c): on the odd cube $s=\frac12\ell_\sharp(b)=LMR(b)$, $b\ge b_\sharp$, $S=s$, B(4) holds uniformly in the label from $K_0=\lceil c_2\log_LS\rceil$ on, $c_2=20\varpi$: its second condition is void on this cube, since $|s|\le L^6M^3(\log b)^{3r_\star}\le(\log b)^{2\mathfrak p_0-1}/C_V$ for large $b$, as $2\mathfrak p_0>3r_\star+1$, so that the low cutoff of B(4) is $K_0$.

At $K_0$ the lattice has $N_0=SL^{K_0}\in[S^{1+c_2},LS^{1+c_2}]$ sites a side, so $N_0\le C(\log b)^{r_\star(1+c_2)}$, at $\beta_0=bL^{K_0}\ge b$; this lattice and those of the cutoffs $K<K_0$, which are smaller at $\beta\ge b$, lie in the range of Proposition~\ref{C-366-13} for $c<1/(r_\star(1+c_2))$ and large $b$. From $K_0$ on, B(4), uniformly in the label, carries the lattice value to every cutoff with a relative error $o(1)$ as $b\to\infty$, the sum of $c_1S^\varpi L^{-K/10}$ over $K\ge K_0$ being at most a multiple of $S^{\varpi-c_2/10}=S^{-\varpi}$; this is (c), the values $t=LMR(b)/b$ filling $(0,t_0]$; $v_z$ depends on $t$ alone (B(3)), which gives $C^{-1}t\le v_z\le Ct$ and, with Proposition~\ref{C-366-13}(e), $v_z(t,t,t)/t\to\mathscr G_z$.

The ratios are Proposition~\ref{C-366-13}(c), (e) for two classes, the periodic denominators cancelling. (d): with the ceiling, $r_1\le\frac15$ on the cubes $2L^m$ with $t\le t_1=\min\{t_0,1/5C_\square\}$ at every cutoff, and A(13) applies; $\delta^0_K\ge\mathsf G\ge\log(4/(1+3r_1))\ge\log(4/(1+3C_\square t))$.

\emph{Step 12: the label (B(6)).} A source is a complex function $\mathsf h$ on $T_\ell$ in the ball of the norm of Theorem~\textup{\YMMref{AD}{607}} (the source $w$ there); for a class $n$, with $x(p)$ the barycentre of the plaquette $p$, put
\[ Z_K[n;\mathsf h]=\int\exp\Bigl[-\sum_pL^K\bigl(b_0-\mathsf h(x(p))\bigr)\bigl(1-\tfrac12\eta^n(p)\,\mathrm{tr}\,\partial U(p)\bigr)\Bigr]dU ; \]
at $n=0$ this is the left member of the identity of Theorem~\textup{\YMMref{AD}{607}} at $\beta=L^Kb_0$, and at a constant source $\mathsf h\equiv c$ it is $Z_K[n]$ at the label $b_0-c$. Let $n$ have one or two twisted planes, $\mu$ a direction that lies in each of them, and $\theta$ the reflection $x_\mu\mapsto-x_\mu$, with the mirror planes $x_\mu\in\{0,\frac\ell2\}$, or $x_\mu\mapsto\frac\ell2-x_\mu$, with the mirror planes $x_\mu\in\{\frac\ell4,\frac{3\ell}4\}$; the two mirror planes cut $T_\ell$ into the closed halves $H_-$ and $H_+=\theta H_-$. Let $\mathsf h$ be symmetric under $\theta$ and real within $\frac\ell{16}$ of the two mirror planes, and let its conjugated fold $\mathsf h^\flat$ be $\mathsf h$ on $H_-$ and $\bar{\mathsf h}$ on $H_+$.

With $\mu$ as the time and the mirror layers $x_\mu\in\{0,\frac\ell2\}$, $Z_K[n;\mathsf h]=\iint\Psi_-(V,V')\,\Psi_+(V,c^kV')\,dV\,dV'$, where $V,V'$ are the fields of the two mirror layers, on whose plaquettes the weights $L^K(b_0-\mathsf h(x(p)))$ are real and positive and give the square roots of their factors to each half, $\Psi_\mp$ are the integrals over the bonds inside $H_\mp$, and $c^k$ is the flip of the electric part of $n$, whose magnetic part vanishes because $\mu$ lies in every twisted plane. The reflection $\theta$ keeps the Haar measure and the traces of the plaquettes, so the symmetry of $\mathsf h$ gives $\Psi_+=\Psi_-$, while $\overline{\Psi_-}$ is the integral over $H_-$ at the source $\bar{\mathsf h}$, that is, by $\theta$, over $H_+$ at $\bar{\mathsf h}$; the Cauchy--Schwarz inequality in the fields of the mirror layers, $c^k$ keeping their Haar measure, gives
\[ |Z_K[n;\mathsf h]|\le\|\Psi_-\|_2^2=Z_K[0;\mathsf h^\flat] . \]
At mirrors midway between two layers the same holds, the crossing kernel at the mirror being positive semidefinite at real positive weights (Step~1).
For three twisted planes the same with the time $3$, which lies in the planes $13$ and $23$ and whose section is twisted, gives $|Z_K[n;\mathsf h]|\le Z_K[n_{12};\mathsf h^\flat]$, the right member being $\ge0$, and a second fold with the time $1$, which lies in the plane $12$, at mirrors normal to the direction $1$ at which the source is symmetric and real as well, gives $Z_K[n_{12};\mathsf h^\flat]\le Z_K[0;\mathsf h^{\flat\flat}]$, $\mathsf h^{\flat\flat}$ being $\mathsf h$ or $\bar{\mathsf h}$ according to the product of the signs of the two halves. The identity of Theorem~\textup{\YMMref{AD}{607}} writes $Z_K[0;\mathsf h^\flat]/|Z_K[0;\mathsf h]|$ as the exponential of the sum of the differences of the real parts of the local energies $e_j(X;\cdot)$ at $\mathsf h^\flat$ and at $\mathsf h$, times the quotient of the moduli of the two remainders $\widetilde Z$.

An energy at $\bar{\mathsf h}$ is the conjugate of the energy at $\mathsf h$, and it reads the source near its domain; so the real parts agree on every localization domain that does not meet both $H_-\cap\{\mathrm{Im}\,\mathsf h\ne0\}$ and $H_+\cap\{\mathrm{Im}\,\mathsf h\ne0\}$, two sets at a distance $\ge\frac\ell8$. A domain of the level $j$ that meets both has an extent of at least $\ell L^{K-j}/8M-4$ (\eqref{eq:analytic-domain-diameter}); the number of the localization domains (Lemma~\ref{lem:analytic-domain-geometry}) and the decay of the energies in the extent bound their sum by a number that does not depend on $K$, and with the bounds on $\widetilde Z$, which also give $Z_K[0;\mathsf h]\ne0$ for small sources,
\[ |Z_K[n;\mathsf h]|\le Z_K[0;\mathsf h^\flat]\le\mathsf M_{\rm S}|Z_K[0;\mathsf h]|,\qquad Z_K[0;\mathsf h]\ne0 , \]
and the same with $\mathsf h^{\flat\flat}$ for three twisted planes. For the constant source $b_0\mapsto b_0-c$, written as $c\chi_1+c\chi_2$ with a partition of unity $\chi_1+\chi_2=1$ in $x_\mu$, $\chi_1$ vanishing within $\frac\ell{16}$ of the planes $x_\mu\in\{0,\frac\ell2\}$ and $\chi_2$ within $\frac\ell{16}$ of the planes $x_\mu\in\{\frac\ell4,\frac{3\ell}4\}$, the quotient $r_K=Z_K[n;\cdot\,]/Z_K[0;\cdot\,]$ at the source $\zeta_1\chi_1+\zeta_2\chi_2$, in which one of the two parts may be complex while the other is real, is bounded by $\mathsf M_{\rm S}$ on the two crosses of a bidisc in $(\zeta_1,\zeta_2)$, and the following lemma bounds it on a smaller bidisc. If $F$ is holomorphic near $\bar D_R^2$ and $|F|\le\mathsf M$ wherever one variable lies in $\bar D_R$ and the other in $[-R,R]$, then $|F|\le27\,\mathsf M$ on $\bar D_{R/8}^2$; if $F$ is holomorphic near $\bar D_R^4$ and $|F|\le\mathsf M$ wherever one variable lies in $\bar D_R$ and the three others in $[-R,R]$, then $|F|\le12600\,\mathsf M$ on $\bar D_{R/15}^4$ (this form serves for three twisted planes).

For the proof put $\rho_\circ=2+\sqrt3$ and $\mathsf j(u)=\frac R4(u+u^{-1})$, which maps the annulus $\rho_\circ^{-1}\le|u|\le\rho_\circ$ into $\bar D_R$ and the unit circle onto $[-\frac R2,\frac R2]$. Moving the circle of one variable of the Cauchy integral to $|u|=\rho_\circ^{\pm1}$, the other variables staying on the unit circle, where $\mathsf j$ is real, bounds the Laurent coefficients of $F\circ\mathsf j$ by $\mathsf M\rho_\circ^{-\max_i|k_i|}$; summing them on the annuli $1.3^{-1}\le|u_i|\le1.3$, whose images hold $\bar D_{R/8}$, gives at most $26.2\,\mathsf M$, and on $1.15^{-1}\le|u_i|\le1.15$, whose images hold $\bar D_{R/15}$, at most $12563\,\mathsf M$ in four variables. This bounds $r_K$ on a disc whose radius does not depend on $K$, and Cauchy's estimate gives the derivatives.

In $t$ this is the normal family of the $\phi_j$, and Vitali's theorem gives the analyticity of the limit. With $m_0=\inf_Kr_K(b_0)>0$, Cauchy's estimate keeps $r_K$ away from $0$ on a smaller disc, which, $Z[0]$ having no zero there, gives the zero-free discs and $\Omega$; the first two derivatives of $\log r_K$ in the label are differences of means and of variances of the total action, bounded by Cauchy's estimate.

\emph{Step 13: block pairs (B(7)).} A block mirror lies midway between two layers, $L$ being odd; for $G$ bounded Borel in the block pairs on one side, $(U,s)\mapsto G(\mathcal M^{K-d}_sU,s^{(d)})$ depends on the bonds and signs on that side, and the reflection commutes with $\mathcal M_s$: the clause on pairs in A(8) gives the positivity, which weak limits keep; at each cutoff $\mathcal M^{K-d}_s=\mathcal M_s\circ\mathcal M^{K-d-1}_s$ makes the tower consistent. Under the conditions of B(7), the small-field chain among the construction results on the periodic torus, with the rarity of its large-field regions (Proposition~\ref{prop:ward-fixed-volume-rarity}), bounds the periodic probability of a rough block plaquette of some depth $d'\in[d,K]$ by $\mathsf R(d)$, uniformly in $K$; the dilation carries this to the other pairs $(\ell,b)$, the depth $d$ becoming $d-j$.

With $\gamma$ a chart of $s$ near a block plaquette $P$, the sign of $P$ in the block sign field times $\partial(\mathcal M^k_sU)(P)$ is $\partial(\mathcal M^k(\gamma U))(P)$: the event of a rough plaquette is covariant and local, and off it the signs of the block traces are the block sign field. In $\mathbb Z/2\lambda$, $\lambda$ odd, every arc of at most $\frac12(\lambda+1)$ points lies in one of four arcs of $\lambda$ points, so each local event lies in one of sixteen quarters of the torus cut by mirrors in two directions; the quarter bound of A(8) carries the union of the events of one quarter into the class with the root $\frac14$, which gives $16\mathsf R(d)^{1/4}$, and the bound for the mixtures.

In a weak limit the class is then, up to null sets, the class of the block sign fields of the tower. For the other statements on limit points pass to a subsequence along which the laws of the towers converge in every class. The quarter bound of A(8) at each cutoff, $r_K[n]\mu_n(B)\le\mu^K_d(B_1)^{1/4}$, and the convergence in total variation of $\mu^K_d$ to $\mu_d$ (through Theorem~\YMMref{D}{8}) give $r_n\mathbb P^n(B)\le\mu_d(B_1)^{1/4}$ for open $B$ by the portmanteau theorem, and for Borel $B$ by outer regularity, the group of the sign changes being finite; the half bound of A(8) gives the exponent $\frac12$. For the consistency of the limit towers the identity at each cutoff is joined to Lusin's theorem, $\mathcal M_s$ being a Borel map and not continuous.

The relation between the block sign fields is a closed condition and passes to the limit. For the fields fix a class with $r_n>0$, a depth $d\ge1$, a block bond $c$ of depth $d$ and $\epsilon>0$; $\mathcal M(\cdot)(c)$ is a Borel function of the block field of depth $d+1$ near $c$, and by Lusin's theorem it is continuous on a compact set $\mathcal K$ of these local fields of $\mu_{d+1}$-measure at least $1-\epsilon$. The towers with $V_d(c)=(\mathcal M_{s_{d+1}}V_{d+1})(c)$ whose local field near $c$, read in a chart of $s_{d+1}$, lies in $\mathcal K$ form a closed set; at each cutoff it holds the tower off a covariant local event whose probability under $\mu_n$ is at most $r_K[n]^{-1}\mu^K_{d+1}(\mathcal K^c)^{1/4}$ by the quarter bound.

By the portmanteau theorem and the convergence of $\mu^K_{d+1}$ in total variation, $\mathbb P^n$ gives this closed set a probability at least $1-\epsilon^{1/4}/r_n$, and $\epsilon\to0$ gives the consistency at $c$, almost surely; there are finitely many $c$ and countably many $d$.

\emph{Step 14: tension of a slab (B(8)).} On one lattice, A(4)(a) at the heights and lengths $N/L$ and $N$, $N=\ell/a_K$, with $\mathsf G,\Phi\le\delta$ and the last inequality of A(2), $\delta(N/L,N,N)\le L^2\delta((N/L)^{\times3})$ and $\delta(N/L,N/L,N)\le L\delta((N/L)^{\times3})$, gives the two inequalities in lattice units, the three terms of the bound on $\sigma$ adding up to $(L/(L-1))^2L^3\delta((N/L)^{\times3})$ exactly; in physical units a product of two lattice lengths carries $L^{2K}$. The right members are bounded in $K$: $\delta^0_K$ by B(1), and $\kappa_K\to\log(2/(1-v_1(t,t,t)))<\infty$ by B(3), since $v_1<1$ on the cubes with $L$-adic sides. In the limit the least level of the sector $e=10$ is $\varepsilon_{\rm el}$ (A(1)) in physical units. The lower bound on $\sigma_K$ is A(3) in physical units, and the bound on $-\log\mathfrak b_K$ is the last bound of A(3) with $R_K(\lambda_0)\le\mathsf G_K(\lambda_0;l,2\lambda_0)/\lambda_0\le\delta^0_K(\lambda_0,l,2\lambda_0)/\lambda_0\le(2\lambda_0l/s'^3)\bar\delta(s')$ by A(4)(b) and A(2), and $\Phi_K(t_0,\lambda_0)\le(t_0\lambda_0l/s'^3)\bar\delta(s')$; the equivalence follows by choosing first $\lambda$, then $t$, large.

\emph{Step 15: the small twisted section (B(9)).} The levels are read in the normalisation of B(9): they are the numbers $a_K^{-1}\log(\lambda_0/\lambda_i)$, $\lambda_0=\|T_1\|$, the four sectors together, so that the least of them, the vacuum, is $E_0=0$; they are the levels of the setting diminished by $a_K^{-1}\log(\|T_0\|/\|T_1\|)$, and their limits are the eigenvalues of $H_{(l,l),1}$ diminished by $E_{\rm mag}(l,l)$. Over the twisted section $(N,N)$, $N=l/a_K$, at the time $1$, $\mathrm{Tr}\,T_1^h=Z[n_{23}](h,N,N)$; with $x(h)=\sum_{i\ge1}(\lambda_i/\lambda_0)^h$, in
\[ L\log\mathrm{Tr}\,T_1^h-\log\mathrm{Tr}\,T_1^{Lh}=L\log(1+x(h))-\log(1+x(Lh)) \]
the bulk and all that is proportional to the time cancel. The chain that bounds it has these links. (i) A dilation by a power of $L$ keeps $t$ and the lattices and takes $(l,b)$ to $(l',b')$ with $l'=\frac12\ell_\sharp(b')=LMR(b')\ge\ell_\flat(b')$, the values $LMR(b')/b'$ filling $(0,t_0]$ (Step~11); small $t$ is large $b'$. (ii) The odd boxes $(l',l',l')$, $(Ll',l',l')$, $(L^2l',l',l')$ and, for each fixed $j\ge1$, $(L^jl',l',l')$ have their sides in $L^{\mathbb N}$ and are mixed for $n_{23}$; on them B(4) holds uniformly in the label from its low cutoff $\lceil c_2\log_LS\rceil$ on, its second condition being void for large $b'$, since $|s|\le L^{j+6}M^3(\log b')^{3r_\star}\le(\log b')^{2\mathfrak p_0-1}/C_V$ as $2\mathfrak p_0>3r_\star+1$.

In the left member of the display at $h=N$, which in the dilated units is $L\log Z[n_{23}](l'^{\times3})-\log Z[n_{23}](Ll',l',l')$, and at $h=LN$, the terms $\frac{|s|}{\ell^3}\Delta_K$ of B(4) cancel, so that B(4) carries its value at the low cutoff to every cutoff with an error summable in $K$. (iii) At and below the low cutoff the lattices lie in the range of the twisted clause of Proposition~\ref{C-366-13}(a), whose volume terms cancel in the display and whose Gaussian term gives the free value: with $\theta^a$ the phases of the three colours $a$ of $n_{23}$ and the momenta $p\in((2\pi\mathbb Z+\theta^a_2)/N)\times((2\pi\mathbb Z+\theta^a_3)/N)$ mod $2\pi$, none of them $0$, $\prod_{j<h}(2\cosh\omega-2\cos\frac{2\pi j}h)=2\cosh h\omega-2$ gives $\log\mathscr D^{n_{23}}(h,N,N)=-hE_{\rm free}+\log(1+x_{\rm free}(h))$, $E_{\rm free}=\frac12\sum_{a,p}\omega(p)$, with
\[ 1+x_{\rm free}(h)=\prod_{a,p}\bigl(1-e^{-h\omega(p)}\bigr)^{-1},\qquad2\cosh\omega(p)-2=\hat p^2 , \]
and the free value of the display is that of $x_{\rm free}$. (iv) With $L\ge13$ of the setting, the free value at $h=N$ is at most $0.26(L-1)$: at $N=L=13$ it is $3.0917$ against $0.26(L-1)=3.12$, the margin $0.028$; at $h=LN$ it is about $4Le^{-\pi L}$. (v) $c_0$, formed from the constants, is chosen so that for $t\le c_0$ the errors of Laplace's method and of B(4) stay below that margin at $h=N$, and are small enough against the free value at $h=LN$ that $x(LN)\ge0.29693\,e^{-\pi L}$ follows. Then at every cutoff the display at $h=N$ is at most $0.26(L-1)$; as $x$ does not increase in $h$, $(L-1)\log(1+x(N))\le0.26(L-1)$, so $x(N)\le e^{0.26}-1=0.29693$ and $N\log(\lambda_0/\lambda_1)\ge\log(1/0.29693)\ge1.2142$; at $h=LN$, with this bound on $x(N)$ and $x(LN)\le(\lambda_1/\lambda_0)^{(L-1)N}x(N)$, $N\log(\lambda_0/\lambda_1)\le\pi L/(L-1)$.

The sharper bounds use the same comparison at smaller $t$ with the free values over the twisted section, whose least excitation is $\omega_N/N$ with $3.1264\le\omega_N<\pi$ for $N\ge13$. Likewise B(4) on the odd boxes $(L^jl',l',l')$, $j\ge1$, carries the quotients $Z[(k,1)_1]/Z[(0,1)_1]$ to their free values; the sector traces $\frac14\sum_k(-1)^{e\cdot k}Z[(k,1)_1]$ at the heights $L^jl$, which are $L^jl'$ in the dilated units, fix the levels (Lemma~\ref{C-366-6}); in the normalisation of B(9) the least free level of the sector $e\ne00$ is $\pi|e|/l$, its momenta lying in $\frac{2\pi}l(\mathbb Z+\frac{e_\rho}2)\times\frac{2\pi}l(\mathbb Z+\frac{e_\nu}2)$ by A(1), while in the sector $00$, whose least level is the vacuum $E_0=0$, the least positive free level, the second level of $00$ minus the vacuum, is $2\pi/l$.

For (c), the state of the cylinder has $\omega_K(\overline{\theta_KF}\,\tau_{n'}G)=\langle\psi(G),(T_1/\lambda_0)^{n'}\psi(F)\rangle$ for plaquette functions of finitely many layers, with vectors $\psi(F)$ formed from the slab operators and the top eigenvector of $T_1$. Lemma~\ref{C-366-21}(b) with cells of side $l/L$, that is $\lambda'=L^{J-1}$ sites and $k=L$ cells a side of the section, and with $\mathsf q=2$, compares the cylinder with the periodic cube $Q$ of side $2l/L$, at the cost of one factor $\Xi$ for each slab of cells in time that carries a function; at $k_0=L$, $\Xi\le\exp[\Phi'/(L-2)]$, and $\Phi'=\frac{L^3}8\delta^0_K(Q)-\delta^{n_{23}}_K(l^{\times3})$, the bulk cancelling, is bounded uniformly in $K$ by the window of B(1) on the cube $2l/L$ and the floor of B(2) on the odd cube $l$.

A partition of unity of each $f_q$ into pieces supported inside cells of translates of the board cuts a tuple into finitely many tuples placed in cells, at a cost formed from $n'$, $R$ and $l$; for a placed tuple the product of the centred observables of one cell is a polynomial in the numbers $1-\frac12\mathrm{tr}\,\partial U(p)$ of its plaquettes, and its seminorm on $Q$ is a centred moment of the observables of the eight mirror copies in the periodic measure of $Q$, the means $e_K$ being those of $Q$; the bounds on the joint cumulants of the action density on that periodic cube (Lemma~\ref{lem:dilated-torus-cumulant-bound}) then give $C_{n'}(d,R)$. Compactness (Lemma~\ref{C-366-6}(b) for the levels, a diagonal choice for the $S_{n'}$) gives the limit points; reflection positivity is A(7) at the mirrors normal to the time, which lies in no twisted plane of $n_{23}$.

In every limit, $S(\theta F,\tau_sG)=\sum_ie^{-sE_i}c_i(F,G)$ with a positive semidefinite matrix of coefficients, which defines an isometry of $\mathcal H$ into the $\ell^2$ space of the limit levels carrying $e^{-sH}$ to multiplication by $e^{-sE_i}$; $E_0=0$ is simple and belongs to $\Omega$, and the gap gives the spectral bound and the clustering. The action density is invariant under the flips $C^k$, so its vectors lie in the sector $00$, whose least positive level exceeds $6.1/l$ at small $t$ when the law of (4) is used on the boxes $(L^3l,l,l)$. Finally the lower bound on $S_2(\theta f^\circ,f^\circ)-S_1(f^\circ)^2$. Let $j(t)$ be the least integer for which $g=(L^{j(t)}/2t)^{-1/2}$ is at most the ceiling $\gamma_W$ of the reference couplings of Theorem~\textup{\YMMref{AV}{851}} and $2L^{j(t)-1}$ is a side at $g$ (equation~\YMMref{(1528)}{615}); it is formed from $t$ and the constants.

Put $\hat L=L^{j(t)-m}$ and read everything in the units dilated by $\hat L$, which name the same lattice and the same $\beta$: the section has the side $\hat l=\hat Ll=L^{j(t)}$, the label is $\hat b=\hat Lb=L^{j(t)}/t=2g^{-2}$, the cutoff is $\hat K=K-j(t)+m$, the bump has the radius $d_*=\hat Ld^\circ$, and the partner of the cylinder is the periodic cube of side $2\hat l/L$, which is the cube $Q$ of side $2l/L$ above; by Corollary~\YMMref{74.2}{616} and Lemma~\ref{lem:setup-label-coverage} every $(\hat K,L^{\hat K}\hat b)$, $\hat K\ge1$, is a first-passage datum at $g$ with the label $\hat b$. (1) The state of the cylinder defines a law $\tilde\omega$ on the configurations of $a_{\hat K}\mathbb Z^3$, carried by the fields that are periodic up to the twist in the two directions of the section, with $\tilde\omega(G)=\omega_K(G^{\rm per})$ for the continuous plaquette functions $G$ of finitely many plaquettes, $G^{\rm per}$ reading each plaquette at its class modulo the section; on these functions $\tilde\omega$ is invariant under the translations of $a_{\hat K}\mathbb Z^3$. (2) Let $F_1,F_2$ be Lipschitz functions supported in one ball of radius $2$ whose centre lies at the distance $\ge3$ from the walls of a board of side $L$, and $\mathsf u=\zeta_1F_1+\zeta_2F_2$, $\zeta\in\mathbb C^2$.

Lemma~\ref{C-366-21}(b), with cells of side $L$, that is of $L^{\hat K+1}$ sites a side, one slab of cells in time that carries a function, and the same partner $Q$, gives $|\tilde\omega(e^{\mathcal O_{\hat K}(\mathsf u)})|\le\hat\Xi\,\|e^{\mathcal O_{\hat K}(\mathsf u)}\|_{{\rm ch};Q}$, where $\hat\Xi=\Xi^{L^{2-j(t)}}\le\Xi$, because at a fixed partner $\log\Xi$ is linear in the number of sites of a side of a cell; so $\hat\Xi$ is bounded uniformly in $K$ by a number $\Xi_*$ formed from $t$ and the constants. (3) By the normalisers among the construction results at the first-passage data (Proposition~\YMMref{91.12}{708}\,(1)) there is a function $\mathcal E$, holomorphic in $\zeta$ on the domain $\mathbb D_\varrho(F_1,F_2)$ among the construction results and with $\mathcal E(0)=0$, with $\|e^{\mathcal O_{\hat K}(\mathsf u)}\|_{{\rm ch};\ell''}\le|e^{\mathcal E(\zeta)}|\,\mathsf Q_1(L)$ for all sides $\ell''$ at $g$, the side $2\hat l/L$ among them; so $|e^{-\mathcal E(\zeta)}\mathfrak s(e^{\mathcal O_{\hat K}(\mathsf u)})|\le Q'=\Xi_*\mathsf Q_1(L)$ for $\mathfrak s=\tilde\omega$ and, by the transfer principle (Proposition~\YMMref{91.11}{707}), for the periodic state $\tilde\mu$ of $Q$, read as a law on $a_{\hat K}\mathbb Z^3$. (4) The transfer between two laws under one normaliser (Proposition~\YMMref{91.11}{707}) gives $|S^T_2[\tilde\omega](F_1,F_2)-S^T_2[\tilde\mu](F_1,F_2)|\le C_D\|F_1\|_\sharp\|F_2\|_\sharp$, $C_D=2\log2\,(2Q'+2)^2(2/\varrho)^2$, $\varrho$ the radius of that domain, with no separation of the supports. (5) Put
\[
 D(f,h)=S^T_2[\tilde\omega](f,h)-S^T_2[\tilde\mu](f,h).
\]
This is a bilinear form invariant under lattice translations, by (1)
and the translation invariance of the periodic law. We apply the
packing argument in the proof of Proposition~\YMMref{127.1}{848};
only these two properties and the unit-ball bound in (4) enter that
argument. In particular, it does not require the cylinder law to be a
periodic Gibbs state.

Let $r=5\mathsf d$, with $a_{\hat K}\le r\le1/20$, and suppose that
$f_1,f_2$ vanish outside one ball of radius $r$. Set
$\delta=a_{\hat K}\lceil2r/a_{\hat K}\rceil$, so that
$2r\le\delta<3r$. Translates of this ball whose centres have all
coordinates of modulus at most $0.575-r$ lie in the unit ball and
have disjoint interiors when their translation vectors lie in
$\delta\mathbb Z^3$. Each coordinate interval contains at least
$0.3/r$ points of the required translate of $\delta\mathbb Z$.
We thus obtain a set $P$ of lattice translations with
$N=\#P\ge(0.3/r)^3$. For independent signs $\epsilon_v=\pm1$ put
\[
 F_{i,\epsilon}=\sum_{v\in P}\epsilon_v\tau_vf_i,\qquad i=1,2.
\]
Their supremum and Lipschitz norms do not exceed those of $f_i$.
For the Lipschitz assertion, each function vanishes on the boundary
of its supporting ball; on a segment joining different balls,
bound the endpoint values by the distances to the respective exit
and entry points. The sum of these distances is no greater than
the length of the segment. Consequently
$\|F_{i,\epsilon}\|_\sharp\le\|f_i\|_\sharp$.
Bilinearity and sign averaging give
\[
 \mathbb E D(F_{1,\epsilon},F_{2,\epsilon})
 =\sum_{v\in P}D(\tau_vf_1,\tau_vf_2)=N D(f_1,f_2).
\]
Applying (4) and dividing by $N$ yields
$|D(f_1,f_2)|\le C_D(r/0.3)^3\|f_1\|_\sharp\|f_2\|_\sharp$.
If $\operatorname{Lip}f_i\le2\|f_i\|_\infty/\mathsf d$, then
$\|f_i\|_\sharp\le80M\|f_i\|_\infty/\mathsf d$, and hence
\[
 |S^T_2[\tilde\omega](f_1,f_2)-S^T_2[\tilde\mu](f_1,f_2)|
 \le C'_2\mathsf d\,\|f_1\|_\infty\|f_2\|_\infty,
 \qquad C'_2=C_D(50/3)^3(80M)^2.
\] (6) Let $f_+$ be the bump of radius $d_*$ about the point $x_+=(2d_*,\hat L\xi)$, the dilate of $f^\circ$, and $f_-$ its mirror image in the plane at the time $a_{\hat K}/2$, the bump about the mirror point $x_-$ of $x_+$.

When $a_{\hat K}\le d_*/4$, $3d_*\le|x_+-x_-|\le8d_*$, both bumps vanish off one ball of radius $5d_*$ and have Lipschitz constants $\le2/d_*$, and $(f_-,f_+)$ is a pair of the law of two points at $\vartheta=\frac12$. For the state $\tilde\mu$ of the torus of side $2\hat l/L$, a state at $g$, Theorem~\textup{\YMMref{AV}{851}} gives, for $d_*\le\mathsf d_0(g,\frac12)$ and $s(d_*)<\hat K$, $s(\cdot)$ the witness depth (equation~\YMMref{(2321)}{843}), $S^T_2[\tilde\mu](f_-,f_+)=c_{\rm sd}\hat b^{-2}\mathcal K_{\rm free}(1+\varepsilon')$ with $|\varepsilon'|\le\frac13$, where the free kernel of the pair has $\mathcal K_{\rm free}=\mathcal K_{\rm free}(f_-,f_+)\ge\frac14\,10^{-6}$. (7) With $d_*=\min\{\mathsf d_0(g,\frac12),\,c_{\rm sd}/(24\cdot10^6C_2'\hat b^2)\}$, a number formed from $t$ and the constants, $C_2'd_*\le\frac16c_{\rm sd}\hat b^{-2}\cdot\frac14\,10^{-6}\le\frac16c_{\rm sd}\hat b^{-2}\mathcal K_{\rm free}$, and so, by (5) at $\mathsf d=d_*$,
\[ S^T_2[\tilde\omega](f_-,f_+)\ge\bigl(1-\tfrac13-\tfrac16\bigr)c_{\rm sd}\hat b^{-2}\mathcal K_{\rm free}\ge\tfrac18\,10^{-6}c_{\rm sd}\hat b^{-2} \]
at every cutoff with $\hat K>s(d_*)$ and $a_{\hat K}\le d_*/4$. Back in the original units, $\mathcal O_{\hat K}(f^{[\hat L]})=\hat L^{-1}\mathcal O_K(f)$ for $f^{[\hat L]}(x)=f(x/\hat L)$, the barycentres being multiplied by $\hat L$ and the factor $a^{-1}$ divided by it; so the covariance of $\mathcal O_K(\theta_Kf^\circ)$ and $\mathcal O_K(f^\circ)$ under $\omega_K$ is $\hat L^2$ times this one, and $\hat L^2\hat b^{-2}=b^{-2}$: the bound $\frac18\,10^{-6}c_{\rm sd}b^{-2}$ holds at all large cutoffs. A covariance does not see the centring by $e_K$, and $\theta_Kf^\circ$ is the translate of $\theta f^\circ$ by $a_K$ in the time; so the bound holds for $S_2(\theta f^\circ,f^\circ)-S_1(f^\circ)^2$ in every limit.

\emph{Step 16: even cubes (B(10)).} The reduction is arithmetic in $\ell_\sharp$ and $R$: for $t\le t_{\rm w}$ the label $b'(t)$ exists, and the cube $2L^j$ at $\beta=2L^j/t$ is the cube $\ell_\sharp(b')L^{j-\mu(t)}$ at $\beta=L^{j-\mu(t)}b'$. (a) follows for $N\ge N(t)$ from B(3), B(5) and $C^n(N)/N\to\mathscr G_n$ (Proposition~\ref{C-366-13}(e)); and up to $K_1(b')$ the cube $\ell_\sharp(b')$ has $N\le\ell_\sharp b'^{1/68}$ sites a side at $\beta=L^Kb'$, so that $N^{66}\log\beta/\beta\le2\ell_\sharp^{66}b'^{-3/68}\log b'$, which tends to $0$ as $b'\to\infty$, and Proposition~\ref{C-366-13}(g) gives the lattice value. Under (R), $\sum_{K\ge K_1(b)}c_1\ell_\sharp^\varpi L^{-\kappa K}\le C(\log b)^{r_\star\varpi}b^{-\kappa/68}\to0$, also under $\sup_{b'\ge b}$; summing from $K_1(b')$ gives the bound of (a) at every $N\in2L^{\mathbb N}$, with its lower bound on $C^n(N)/N$ the floor $ct$, and the bounded variation, $\phi_{j+1}-\phi_j$ being a difference of two consecutive cutoffs on the cubes $\ell_\sharp(b'(t))L^{j-\mu(t)}$; with the one-sided form, $\log q_K\ge\log q_{K_1(b')}-o(1)$ for $K\ge K_1(b')$ gives $ct\le r_K$, and Lemma~\ref{C-366-18} gives the bounded variation.

Under (W) and (Window estimates), at $K_\sharp$ the lattice lies in the range of Definition~\ref{C-366-15}: $w_{K_\sharp}\ge c_\flat t$ and $\varsigma_{K_\sharp}\ge2^{1-L}e^{\lambda_\flat}(c_\flat t)^L$, beside which $e^{-c(\log b)^{1+\eta}}$ is $o(1)$; Lemma~\ref{C-366-12}(b) puts $w_K$ within $e^{\pm E}$ of $w_{K_\sharp}$ up to $\epsilon$, which is $o(t)$, indeed $o(t^k)$ for every $k$, as Lemma~\ref{C-366-12}(b) states ($m_K$ and $\epsilon_K$ being formed from the numbers \YMMref{(2192)--(2194)}{808} and $\mathsf R_{t_1}$ of Section~\ref{sec:cutoff-calibration}), and so negligible beside $w_{K_\sharp}\ge c_\flat t$; so $\log(2^{L-1}\mathfrak q_L(K))\ge\lambda_\flat-\Lambda-o(1)>0$, and A(12) gives
\[ r_K[n_{13}]\ge w_K\,\frac{\mathfrak q_L(K)-2^{1-L}}{1-2^{1-L}}\ge c't ; \]
smaller lattices lie in the window. On the box $(\ell,\ell,i\ell)$ the same holds with $2^{i-1}\mathfrak q_i\ge0.42\cdot1.8712^i$ for $F'_t$. Under (W$'$), A(12) gives $r_K[n_{13}]\ge w_K(2\|\mathsf P_K\|/\mathrm{Tr}\,\mathsf P_K-1)\ge w_K(e^{\gamma'}-1)$, with $w_K\ge\frac12c_\flat t$ for small $t$ by Lemma~\ref{C-366-12}(b), its $\epsilon$ being $o(t)$. With $F''_t$ the same gives the quotient of the class $((1,1),0)_3$, which on the cube is $r_2$. $0.55L-1.1<0.5516L-0.7162$, and $\frac L2<0.6266L-0.873$, as $L\ge13$. In (d), the consistency is that of B(7); as $\{n(s)=n\}$ is invariant under $s\mapsto s\,d\gamma$, (M) makes the unnormalised laws of each class converge, and B(2)(c) normalises them.

\emph{Step 17: Theorem~\ref{C-366-5}.} Under the hypothesis (H), at each $K\ge K_0$, Lemma~\ref{C-366-9} at $l=l_K$ gives $\theta\le0.8901$, $\mathsf e(4l_K)\le\theta^4/12.9$ and $D_2(8l_K)\le20\,\mathsf e(4l_K)^2\le0.047355$, so A(10) applies at $l_0=8l_K$; $m_*l_0\ge\log\frac1{18\cdot0.047355}\ge0.1597$, and if $m_*l_0<3$ then $l_1<6/m_*\le37.6\,l_0$ and $\sigma_*>m_*^2/6\ge0.00425\,l_0^{-2}$, while if $m_*l_0\ge3$ then $\sigma_*=m_*/l_0\ge3l_0^{-2}$; with $l_0a_K=4s_0$ this is the stated bound in physical length. The cube and loop bounds are A(10) in physical length, with $4N/l_1\le\ell/s_0$ and $N\ge4l_1$ for $\ell\ge602\,s_0$; the uniqueness of $\nu_K$ is Lemma~\ref{C-366-16}. On a cube of side $N\ge2l_K$, A(10) gives $\Phi\le\delta\le0.13736\,\theta^{N/2l_K}$ and $1-r[n]\le0.4121\,\theta^{N/2l_K}$; so $\mathfrak b\le e^{0.13736}-1<0.1473$ on that cube, and A(3) gives
\[ \sigma(N)\ge N^{-2}\log(1/0.1473)\ge1.9\,N^{-2},\qquad\text{that is}\qquad\sigma_K(\ell)\ge1.9\,\ell^{-2}. \]
For the final exclusion assume (Hess) and (Flux-window), and let $s_0\le t_1b$. On the cube of physical side $s_0=2l_Ka_K$, $r_1\to v_1(s_0/b,s_0/b,s_0/b)\le C_\square t_1\le\frac15$ by B(3) and B(5), so $r_1\le\frac13$ for all large $K$, and A(13) gives $\mathsf e(l_K)\ge\log2>0.069$ there, against the hypothesis (H).
\end{proof}

\subsection{Scope}

Theorem~\ref{C-366-3} is a statement on one lattice, for every box and every $\beta\ge0$, in lattice units; besides the periodic boxes it treats the infinite lattice at one $\beta$, in the last assertion of~(8), in~(10), under its hypothesis, one inequality between the partition functions of two periodic cubes, with its periodic Gibbs state and its loops, and in~(11), for every periodic Gibbs state, and, in~(4)(d), a slab with free boundaries in two directions. Theorem~\ref{C-366-4} concerns the periodic boxes of a fixed physical size, at a fixed label or at labels in a compact set, at every cutoff and as $K\to\infty$, the Hamiltonians $H_{l,m}$ over fixed sections, and, in~(9)(c), the cylinder over one small twisted section.

Theorem~\ref{C-366-5} concerns the infinite lattice $a_K\mathbb Z^3$ at each cutoff $K\ge K_0$: under the hypothesis~(H), the inequality $\mathsf e(l_K)\le0.069$ between the partition functions of the two periodic boxes $(l_K,2l_K,2l_K)$ and $(2l_K)^{\times3}$, it gives the uniqueness of the periodic Gibbs state $\nu_K$, an area law for the loops of every rectangle with a physical tension at least $2.65\cdot10^{-4}s_0^{-2}$, and bounds on the cubes of every physical side $\ge s_0$, resp.\ $\ge602\,s_0$, all uniform in $K$; it is a statement of the kind of confinement at each cutoff, conditional on the hypothesis~(H). Under (Hess) and (Flux-window), the final exclusion further shows that (H) can hold only at physical sides $s_0$ of the order of the label or more.

The vacuum, the gap and the clustering of Theorem~\ref{C-366-4}(9) belong to one small twisted section, the uniqueness of $\nu_K$ is that of a Gibbs state of the lattice of one cutoff, and the tension $\sigma_K(l)$ of Theorems~\ref{C-366-4}(8) and~\ref{C-366-5} is that of a slab of fixed width $l$. Theorem~\ref{C-366-4}(10)(b)--(d) and Theorem~\ref{C-366-5} are implications from (R), (W), (W$'$), (M) and (H), with (Window estimates) additionally required in (10)(c) and (Hess) and (Flux-window) for the final small-side exclusion; and Theorem~\ref{C-366-3}(7), (8) give reflection positivity for the eight sums over sets of coordinate planes and for the admissible mixtures. A box of the class $z$ is a centre sector of a periodic torus, that is, its fields of the class $z$ are the fields of the torus fixed by a group generated by the shifts by the sides of the box, each followed by the multiplication by a sign field $\chi$ on the bonds with $d\chi\equiv1$, if and only if $n_\mu n_\nu$ is even in every twisted plane $\mu\nu$ of $z$, $n_\mu$ being the side of the torus divided by that of the box.

Indeed, if it is such a sector, a fixed field is determined by its values at representatives of the bonds, which are free, its plaquette variables are periodic under the shifts, and on one copy of the box they are those, $\eta^z\partial U'$, of a field $U'$ of the box; for values in the diagonal subgroup of $\mathrm{SU}(2)$, which holds $-\mathbf 1$, the product of the plaquette variables over a plane $\mu\nu$ of the torus is $\mathbf 1$, each bond entering it twice in opposite senses, and it is $(-1)^{z_{\mu\nu}n_\mu n_\nu}$, the product over the $n_\mu n_\nu$ copies of a plane of the box; the converse is the construction of Lemma~\ref{C-366-7}, read on that torus, which asks $w_{\mu\nu}=1$ only where $n_\mu$ is even.

On the tori $T_{\ell'}$, $\ell'=2L^{m+i}$ with $i\ge0$, the multiplicities of the cube $2L^m$ are $L^i$, odd, and a dilation by a power of $L$ does not change them: the cube $2L^m$ with flux is a centre sector of no torus $T_{\ell'}$, and no box of a class with a twisted plane whose two sides lie in $2L^{\mathbb Z}$ is a centre sector of a torus $T_{\ell'}$. The cube $2L^m$ of the class $z$ is a centre sector of the periodic tori with one doubled side in each twisted plane of $z$: of $(4L^m,2L^m,2L^m)$, the doubled side lying in the twisted plane, for one twisted plane, and in both twisted planes for two; of $(4L^m,4L^m,2L^m)$ or of the cube $4L^m$ for three. The chain and the two-cutoff comparison are stated on the tori $T_\ell$, and Lemma~\ref{C-366-12} on periodic cubes.

For the pairs $(U,s)$ of the mixture $\mu_{\rm mix}$ of~(M), let a chain be run with the averaging map $\mathcal M_s$ of Theorem~\ref{C-366-4}(7) and with the small-field tests applied to the signed plaquette variables, $s(p)\,\partial U(p)$ on the lattice and $s^{(d)}(P)\,\partial(\mathcal M^{K-d}_sU)(P)$ on the block lattice of depth $d$. These variables, and so the small-field and the large-field regions of this chain, do not change under $(U,s)\mapsto(\gamma U,s\,d\gamma)$, $\gamma$ a sign field on the bonds: it multiplies $\partial U(p)$ and $s(p)$ by $d\gamma(p)$, and, at each depth, $\mathcal M^{K-d}_sU$ by the sign field $\bar\gamma$ on the block bonds, $\bar\gamma(c)$ the product of $\gamma$ over the straight path of the block bond $c$, and $s^{(d)}$ by $d\bar\gamma$.

When the small-field region of this chain does not wind about the torus, the closed sign field $s$ is $d\gamma$ on it for a sign field $\gamma$ on the bonds, and the change of variables $(U,s)\mapsto(\gamma U,s\,d\gamma)$, under which $A_{s\,d\gamma}(\gamma U)=A_s(U)$ and which the functions $H$ of~(M) do not see, makes the integrand on that region the periodic one; $s\,d\gamma$ is then a closed sign field of the class $n(s)$ that is $1$ on that region, and so, for $n(s)\ne0$, $-1$ on some plaquette of the large-field region of this chain (the probability estimate for such large-field events is Proposition~\ref{prop:ward-fixed-volume-rarity}).

\section{Latent heat and one-point constants}\label{lh:sec}

Let $P=\lim_{N\to\infty}N^{-3}\log Z_N$ be the lattice pressure and $J(\beta)=P^{\prime}(\beta+)-P^{\prime}(\beta-)$ its derivative jump. We prove $J(\beta)=O(\beta^{-4-c})$ for every $c<1/180$. The proof combines a reflection-positive lower bound, comparison with finite tori and an oscillation estimate for two incommensurable blocking factors.

The cutoff threshold depends on the size of the box. We retain this dependence: the result is a decay estimate for the jump and its continuum one-point consequences. The two-period lemma and the improved discarded-mass estimate are the main quantitative steps.

\subsection{Notation and hypotheses}

\emph{One lattice.} For a plaquette $p$ of $\mathbb Z^3$ with barycentre $x(p)$ and holonomy $\partial U(p)$ we put $a_p=1-\frac12\operatorname{tr}\partial U(p)\in[0,2]$, and $\mathcal O_1(f)=\sum_pf(x(p))a_p$ for bounded $f$ of compact support. For an integer $N\ge3$, $S_N$ is the sum of the $a_p$ over the $3N^3$ plaquettes of the periodic cube with $N^3$ sites, $Z_N(\beta)=\int e^{-\beta S_N}dU$ with the Haar probability measure of $\mathrm{SU}(2)$ on every bond, $\mu_{N,\beta}$ is the Wilson measure of the cube, read as a law on the configurations of $\mathbb Z^3$ that is carried by the $N\mathbb Z^3$-periodic configurations (the periodic Wilson law in equation~\YMMref{(1817)}{698}, with $\Gamma=N\mathbb Z^3$), and $P_N=N^{-3}\log Z_N$, so $P=\lim_NP_N$.

For $x\in\mathbb Z^3$ let $c_x=\frac12\sum a_p$, the sum over the six faces of the unit cube $x+[0,1]^3$. A Gibbs state of Wilson's action at the spacing $a>0$ and the coupling $\beta$ is a Borel probability measure $\nu$ on the configurations of $a\mathbb Z^3$ such that $\nu(\gamma_BF)=\nu(F)$ for every finite set $B$ of bonds and every continuous function $F$ of finitely many bonds; here the Gibbs kernel $\gamma_B$ is the normalised average over the configurations on $B$, with the Haar measure of $\mathrm{SU}(2)$ on every bond of $B$ and the weight $e^{-\beta\sum_pa_p}$, the sum over the plaquettes $p$ that contain a bond of $B$, the configuration off $B$ being kept fixed (Section~\YMMref{91}{697}).

A Gibbs state is periodic if it is invariant under the translations by a subgroup of finite index of $a\mathbb Z^3$ (Section~\YMMref{91}{697}); by Proposition~\YMMref{103.4}{758} every periodic Gibbs state is invariant under all isometries of $a\mathbb Z^3$, so the periodic Gibbs states are exactly the Gibbs states that are invariant under $a\mathbb Z^3$. $\mathcal G(\beta)$ is the set of all periodic Gibbs states on $\mathbb Z^3$ at $\beta$, that is, at the spacing $1$ (the set $\mathcal G^{\rm per}_{1,\beta}$ of Section~\YMMref{91}{697}), and $e_\nu=\nu(a_p)$ for $\nu\in\mathcal G(\beta)$; this is one number for all $p$ (Lemma~\ref{lh:lat}). For a law $\mathfrak s$ that is invariant under $\mathbb Z^3$ we write $\langle f;h\rangle_{\mathfrak s}=\operatorname{Cov}_{\mathfrak s}(\mathcal O_1(f),\mathcal O_1(h))$ and $\Gamma_{\mathfrak s}(z)=\operatorname{Cov}_{\mathfrak s}(c_0,c_z)$.

\emph{Tents and the reference.} $\Lambda_\eta(u)=(1-|u|/\eta)_+$ for real $u$ and $\eta>0$. For an integer $\eta\ge1$ and $X\in\mathbb Z^3$ the tent is $\varphi^\eta_X(z)=\prod_{i\le3}\Lambda_\eta(z_i-\frac12-\eta X_i)$, the hull is $h^\eta=\sum_{|X|_\infty\le2}\varphi^\eta_X$, the sum of the $5^3$ tents nearest to $\varphi^\eta_0$, and $\varphi^\eta_{\rm c}(x)=\prod_i\Lambda_\eta(x_i)$ for $x\in\mathbb Z^3$. For a law $\mathfrak s$ invariant under $\mathbb Z^3$ we put $s^\eta_{\mathfrak s}(X)=\langle\varphi^\eta_0;\varphi^\eta_X\rangle_{\mathfrak s}$, and the near part of the specific heat is $n_\eta(\mathfrak s)=\eta^{-3}\langle\varphi^\eta_0;h^\eta\rangle_{\mathfrak s}$. For $r_1>0$ the reference is
\[
\hat n(\beta)=\liminf_{N\to\infty}n_{\lceil r_1\beta\rceil}(\mu_{N,\beta}),
\]
the lower limit along all integers $N\ge3$; it is formed from Wilson's action and $r_1$ alone. $P''(d\beta)$ is the second derivative of the convex function $P$, a positive measure, and $\mathbb M$ is the image under $s=\log\beta$ of the signed measure $\beta^4[P''(d\beta)-\hat n(\beta)d\beta]$. Thus $\mathbb M(\{s\})=e^{4s}J(e^s)$, and
\[
|\mathbb M|([s,s+\delta])=\int_{[\beta,e^\delta\beta]}x^4\,|P''(dx)-\hat n(x)\,dx|\qquad(\beta=e^s).
\]

\emph{Settings.} A setting consists of a radius $\varrho>0$, a common choice $\varsigma^{\rm S}$ of the real and source construction constants and auxiliary data in Section~\YMMref{82}{665} at the radius required by Theorems~\YMMref{AJ}{673} and~\YMMref{AD}{607}, with \eqref{eq:canonical-source-quarter-room}, and the associated Wilson measures and source normalization. Write $\mathcal T^{\rm S}[2\varrho]$ for this ordered choice and $\varsigma=\pi^+(\varsigma^{\rm S})$ for its restriction to the real construction, denoted by $\mathcal T^{\rm F+}$. The numerical constants retain their values: the real hypotheses are those of Proposition~\YMMref{69.1}{595}, Proposition~\YMMref{69.3}{603} and Corollary~\YMMref{69.4}{604}; the additional source floors are imposed in their stated order. The coupling ceiling $\gamma\in(0,\frac12]$ is chosen last.

Among the constants chosen in Proposition~\YMMref{69.1}{595} are the odd block size $L\ge13$, the cell size $M$, a power of $L$, the flow bound $B_\sharp$ with $\bar B=B_\sharp L/(L-1)$, the exponent $\mathfrak p_0\in\mathbb N$ of the large-field threshold, the exponent $r$ of the cube sizes, with which the cube size at a coupling $x>1$ is $\mathcal R(x)=\min\{L^c:c\in\mathbb Z_{\ge0},\ L^c\ge(\log x)^r\}$ (equation~\YMMref{(634)}{339}; below it is read at couplings $x\ge\gamma^{-2}$ alone), and the amplitudes $A_1$ and $A_0$. A lowered choice $\varsigma'$ of a choice $\varsigma$ of $\mathcal T^{\rm F+}$ (Definition~\ref{def:rarity-choices}) keeps the values of the constants before $A_1$ and of $A_0$, has $B_\sharp'\ge B_\sharp$, and has $A_1'=\lambda'_\natural A_1$ with $\lambda'_\natural=\lambda_\natural/(2+30L)$, where $\lambda_\natural\in(0,\frac1{8\pi}]$ is a number of $\varsigma$; so the small-field thresholds $\delta_j$ of the two choices satisfy $\delta'_j=\lambda'_\natural\delta_j$.

For the ordered choices constructed in Lemma~\ref{lem:rarity-lowered-choice}, lowered choices exist: decreasing $A_1$ preserves the eight inequalities at $A_0$, and the later constants are reselected in the order proved there. Definition~\ref{def:rarity-choices} attaches to the pair $(\varsigma,\varsigma')$ a floor $\mathsf Y^\natural$ for the logarithm of the last coupling and the rate $c_2^\natural=\frac12\min\{\frac18c_1^\natural,\lambda_\natural^2A_1^2/2A_0^2\}$, with a number $c_1^\natural>0$ of the pair. We put $a_K=L^{-K}$, $\mathbb L_K=a_K\mathbb Z^3$, $\lambda=\log L$, $\ell=2L^m$ and $T_\ell=\mathbb R^3/\ell\mathbb Z^3$; the label of a pair $(K,\beta)$ is $b=a_K\beta$, and
\[
a_K^{-4}J(\beta)=b^{-4}\beta^4J(\beta).
\]

\emph{The physical scale.} For the plaquettes $p$ of $\mathbb L_K$ we put $\mathcal O_K(f)=a_K^{-1}\sum_pf(x(p))a_p$ and $r_K(f)=\frac13a_K^3\sum_pf(x(p))$, a Riemann sum of $\int f$; and $\|f\|_\sharp=\max\{\|f\|_\infty,40M\operatorname{Lip}f\}$. $\mathcal G_{K,\beta}$ is the set of all periodic Gibbs states on $\mathbb L_K$ at $\beta$, that is, at the spacing $a_K$ (the set $\mathcal G^{\rm per}_{K,\beta}$ of Section~\YMMref{91}{697}), a convex set (Propositions~\YMMref{91.1}{698} and~\YMMref{103.4}{758}), and its extreme points are the pure phases (Proposition~\YMMref{106.1}{764}). $S^T_n[\nu](f_1,\dots,f_n)$ is the joint cumulant of $\mathcal O_K(f_1),\dots,\mathcal O_K(f_n)$ in a law $\nu$ on $\mathbb L_K$; for real Lipschitz $f_i$ on $T_\ell$, $S^T_{n;K,\ell,\beta}(f_1,\dots,f_n)$ is the joint cumulant under the Wilson measure of $T_\ell$ at the spacing $a_K$, that is, of the periodic cube with $\ell L^K$ sites a side. $f^{\langle\ell\rangle}=\sum_{v\in\ell\mathbb Z^3}f(\cdot-v)$ is the periodisation. The spread of $(f_1,\dots,f_n)$ is the least sup-diameter of a set of points taken one from each support; for $n=2$ it is the sup-distance of the two supports. On $T_\ell$ one has $|\mathcal O_K(f)|\le6\ell^3L^{4K}\|f\|_\infty$, so that
\[
|S^T_{n;K,\ell,\beta}(f_1,\dots,f_n)|\le\mathsf{op}_n\prod_{i\le n}6\ell^3L^{4K}\|f_i\|_\infty ,
\]
with a combinatorial number $\mathsf{op}_n$ of $n$; we call this the crude bound.

\emph{Labels and windows.} For $0<g\le\gamma$ let $X_g=L(g^{-2}+B_\sharp)$ and $I(g)=(g^{-2}+\bar B,X_g-\bar B]$ (Section~\YMMref{74}{614}), and let $\mathfrak d(g)$ be the canonical family at $g$ (Section~\ref{sec:state-limit-conventions}), the set of the first-passage data at $g$ (Corollary~\YMMref{74.2}{616}): the pairs $(K,\beta)$, $K\in\mathbb N$, for which $K$ is the least index $k$ at which the running coupling $x_k(\beta)$, the member of index $k$ of the sequence of the increments from $x_0(\beta)=\beta$, is defined, real and at most $X_g$; in words, the running coupling arrives at the scale of $g$ at the last step. We use two facts on $\mathfrak d(g)$. By Section~\ref{sec:state-limit-conventions}, the labels $a_K\beta$ of the pairs of $\mathfrak d(g)$ lie in the compact interval $I^{\rm lab}(g)=[g^{-2}+B_\sharp\frac{L-2}{L-1},\,L(g^{-2}+\bar B)]$ of Section~\ref{sec:state-limit-conventions}, whose lower end is at least $g^{-2}$; so $a_K\beta\ge g^{-2}$ on $\mathfrak d(g)$. By Section~\ref{sec:state-limit-conventions}, $(K,L^Kb)\in\mathfrak d(g)$, with the label $b$, for every $b\in I(g)$ and $K\in\mathbb N$. A window is an interval of labels; in a setting we use the window
\[
\mathsf I=[b_-,b_+],\qquad b_-=\gamma^{-2}+\bar B+\tfrac12,\qquad b_+=L(\gamma^{-2}+3\bar B+1)-2\bar B,
\]
and $g_2=(\gamma^{-2}+2\bar B+1)^{-1/2}$. $m_0\ge1$ is the least integer such that the sides $2L^m$, $m\ge m_0$, are admitted at $\gamma$ and at $g_2$ (the sides that the section of Section~\YMMref{74}{614} attaches to a reference coupling), and $\ell_0=2L^{m_0}$. With $\varepsilon=L^{-1}$, the tent $\mathsf f(y)=\prod_i\Lambda_\varepsilon(y_i-\frac12)$ and its hull $\mathsf h=\sum_{|X|_\infty\le2}\mathsf f(\cdot-\varepsilon X)$, the far part is
\[
y_{m,K}(\upsilon)=b^5\varepsilon^{-3}S^T_{2;K,\ell,L^Kb}\bigl(\mathsf f^{\langle\ell\rangle},1-\mathsf h^{\langle\ell\rangle}\bigr)\big|_{b=e^\upsilon},\qquad\ell=2L^m .
\]
When two settings $(\mathrm S_1)$, $(\mathrm S_2)$ are used, an index $i$ marks the objects of $(\mathrm S_i)$: $L_i$, $\lambda_i=\log L_i$, $M_i$, $\mathsf I_i=[b_{-,i},b_{+,i}]$, $m_{0,i}$, $\ell_{0,i}$, $\varepsilon_i$, $y^{(i)}_{m,K}$, and so on.

\emph{Tuples and functionals.} A separated tuple is $X=(f_1,\dots,f_n)$ with real $f_i\in C^1_c(\mathbb R^3)$ and pairwise disjoint supports; the empty tuple is one of them, $(f)$ is the tuple with the one entry $f$, $|X|=n$ is the number of entries, $X_A=(f_i)_{i\in A}$ for $A\subset\{1,\dots,n\}$, $(\tau_vf)(x)=f(x-v)$, $\tau_vX=(\tau_vf_i)_i$, and $X\cdot Y$ is the joined tuple, formed when its entries have pairwise disjoint supports. A symmetric functional is a map $\sigma$ on the separated tuples that is multilinear and symmetric in the entries, with $\sigma$ of the empty tuple equal to $1$. For a law $\nu$ on $\mathbb L_K$ we put
\[
X^\nu_K=\prod_{i\le n}\bigl[\mathcal O_K(f_i)-\nu(\mathcal O_K(f_i))\bigr],\qquad\sigma_\nu(X)=\nu(X^\nu_K).
\]
A Gibbs sequence at $g$ is $(K_j,\beta_j,\nu_j)_j$ with $(K_j,\beta_j)\in\mathfrak d(g)$, $K_j\to\infty$ and $\nu_j\in\mathcal G_j=\mathcal G_{K_j,\beta_j}$. A limit in its frame is $\sigma'(X)=\lim_j\rho_j(X^{\nu_j}_{K_j})$, for all separated $X$, with approximants $\rho_j\in\mathcal G_j$. The sequence converges if $\sigma(X)=\lim_j\sigma_{\nu_j}(X)$ exists for every separated $X$; then $\sigma$ is its limit, and $S^\sigma_2(f,h)=\sigma((f,h))-\sigma((f))\sigma((h))$. For separated $X^1,\dots,X^q$ the iterated plateau of a symmetric functional $\sigma$ is the iterated limit
\[
\lim_{|v_q|\to\infty}\cdots\lim_{|v_2|\to\infty}\sigma\bigl(X^1\cdot\tau_{v_2}X^2\cdots\tau_{v_q}X^q\bigr),
\]
the translations receding one after the other, $v_2$ first; at $q=1$ it is $\sigma(X^1)$, and at $q=2$ it is $\lim_{|v|\to\infty}\sigma(X^1\cdot\tau_vX^2)$. A canonical decomposition $(\mathsf X,\varpi,(\sigma_\xi)_{\xi\in\mathsf X})$ of $\sigma$ consists of (Can1) a compact metrisable space $\mathsf X$, (Can2) a Borel probability measure $\varpi$ of full support on it, (Can3) a map $\xi\mapsto\sigma_\xi$ into the symmetric functionals that is continuous for the pointwise convergence and injective, such that (Can4) $\sigma$ has iterated plateaux and $\int\prod_{i\le q}\sigma_\xi(X^i)\,d\varpi$ is the iterated plateau of $(X^1,\dots,X^q)$, for all $q\ge1$ and all separated $X^1,\dots,X^q$; at $q=1$, $\sigma=\int\sigma_\xi\,d\varpi$.

The $\sigma_\xi$ are the components. A symmetric functional has one vacuum if it has a canonical decomposition whose space is a point. By Theorem~\textup{\YMMref{AP}{780}}, the limit $\sigma$ of a convergent Gibbs sequence at $g$ has a canonical decomposition, unique up to isomorphism; every component $\sigma_\xi$ is a limit in the frame along the full sequence, with the approximants $\nu_j(\cdot\,|B_j)$, $B_j$ invariant with $\nu_j(B_j)>0$; $\sigma_\xi((f))=\mathfrak c_\xi\int f$ with a number $\mathfrak c_\xi$; $|\sigma_\xi(X\cdot\tau_vY)-\sigma_\xi(X)\sigma_\xi(Y)|\le C^g(X,Y)/|v|$ for $|v|\ge r(X,Y)$, with the numbers $C^g$, $r$ of the paragraph on one cutoff below; and $\#\mathsf X=\dim\ker H_\sigma$, both finite and equal or both infinite, where $H_\sigma\ge0$ is the Hamiltonian of the reconstruction of $\sigma$ in any of the six coordinate directions and $H_\sigma\Omega_\sigma=0$ for its vacuum vector $\Omega_\sigma$.

So such a limit has one vacuum if and only if $\ker H_\sigma=\mathbb C\Omega_\sigma$. We put $\kappa=\int\mathfrak c_\xi^2\,d\varpi$.

\emph{Limits along sequences of pairs.} A sequence of pairs is $\mathfrak p=(K_j,\beta_j)_j$ with $K_j\in\mathbb N$, $\beta_j>0$ and $K_j\to\infty$; it is at the label $b>0$ if $a_{K_j}\beta_j\to b$. $\mathsf K_{\mathfrak p}$ is the set of the functionals $\sigma$ for which there are $\nu_j\in\mathcal G_{K_j,\beta_j}$ with $\sigma(X)=\lim_j\sigma_{\nu_j}(X)$ for every separated $X$, along the full sequence; it carries the topology of pointwise convergence. $\mathsf K^\bullet_{\mathfrak p}$ is the set of its members with one vacuum. ${\rm Lim}[b]$ is the union of the sets $\mathsf K_{\mathfrak p}$ over the sequences of pairs at the label $b$, and ${\rm Lim}^\bullet[b]$ the set of its members with one vacuum; for a set $\mathfrak q$ of pairs, ${\rm Lim}_{\mathfrak q}[b]$ and ${\rm Lim}^\bullet_{\mathfrak q}[b]$ are the same unions over the sequences all of whose pairs lie in $\mathfrak q$. A sequence of pairs $\mathfrak p$ is saturated if $\mathsf K_{\mathfrak p}\ne\emptyset$ and $\mathsf K_{\mathfrak p'}=\mathsf K_{\mathfrak p}$ for every subsequence $\mathfrak p'$. For a Borel probability measure $\mu$ on a compact metrisable set of symmetric functionals, $\operatorname{bar}(\mu)=\int\psi\,\mu(d\psi)$ is its barycentre.

\emph{Dilations; the reconstruction at every label.} For $z\in\mathbb Z$ and a separated tuple $X=(f_1,\dots,f_n)$ put $X^{[z]}=(f_i(L^z\cdot))_i$, again a separated tuple, and for a symmetric functional $\sigma$ put $\sigma^{[z]}(X)=L^{-z|X|}\sigma(X^{[z]})$. The dilation $x\mapsto L^zx$ carries a law $\nu$ on $\mathbb L_K$ to a law $\nu^{(z)}$ on $\mathbb L_{K-z}$, maps $\mathcal G_{K,\beta}$ onto $\mathcal G_{K-z,\beta}$ (Lemma~\ref{lem:lattice-law-dilation}) and turns $\mathcal O_{K-z}(f)$ into $L^{-z}\mathcal O_K(f(L^z\cdot))$ (Lemma~\ref{lem:lattice-law-dilation}, applied to the function $f(L^z\cdot)$: the dilated function $f^{[z]}$ of that Proposition is $f(L^{-z}\cdot)$, whereas the tuple $X^{[z]}$ above has the entries $f_i(L^z\cdot)$); so $\sigma_{\nu^{(z)}}=(\sigma_\nu)^{[z]}$, and the pair $(K-z,\beta)$ has the label $L^za_K\beta$.

Hence, if $\mathfrak p=(K_j,\beta_j)_j$ is a sequence of pairs at the label $b$ and $\sigma\in\mathsf K_{\mathfrak p}$, then $\sigma^{[z]}\in\mathsf K_{\mathfrak p'}$ for the sequence of pairs $\mathfrak p'=(K_j-z,\beta_j)_{j\ge j_0}$ at the label $L^zb$, where $K_j>z$ for $j\ge j_0$. As $\sigma^{[z]}(X\cdot\tau_vY)=L^{-z(|X|+|Y|)}\sigma(X^{[z]}\cdot\tau_{L^{-z}v}Y^{[z]})$, the conditions (Can1)--(Can4) pass through the dilation, the receding translations by $v$ becoming translations by $L^{-z}v$: $(\mathsf X,\varpi,(\sigma_\xi)_{\xi\in\mathsf X})$ is a canonical decomposition of $\sigma$ if and only if $(\mathsf X,\varpi,(\sigma_\xi^{[z]})_{\xi\in\mathsf X})$ is one of $\sigma^{[z]}$, with the same space $\mathsf X$.

For every $b>0$ there are an integer $z\ge0$ and $0<g\le\gamma$ with $g^{-2}+\bar B<L^zb<X_g-\bar B$, because $L^zb>\gamma^{-2}+\bar B$ for large $z$ and these open intervals overlap one another as $g$ decreases and cover $(\gamma^{-2}+\bar B,\infty)$; we fix such a pair $(z,g)$ for $b$, with $z=0$ if $b$ itself lies in one of these intervals. Then $L^za_{K_j}\beta_j\in I(g)$ for large $j$, so $(K_j-z,\beta_j)\in\mathfrak d(g)$ by Section~\ref{sec:state-limit-conventions}, and $\sigma^{[z]}$ is the limit of a convergent Gibbs sequence at $g$ over these pairs, to which Theorem~\textup{\YMMref{AP}{780}} applies. For $\sigma\in\mathsf K_{\mathfrak p}$ we put $H_\sigma=L^zH_{\sigma^{[z]}}$ and $\Omega_\sigma=\Omega_{\sigma^{[z]}}$, on the space of the reconstruction of $\sigma^{[z]}$, in each of the six coordinate directions, which the dilation keeps, as it commutes with the coordinate mirrors: the time $t$ of $\sigma^{[z]}$ is the time $L^{-z}t$ of $\sigma$, and the kernel is that of $H_{\sigma^{[z]}}$.

As the canonical decompositions of $\sigma$ and of $\sigma^{[z]}$ correspond to one another with the same space $\mathsf X$, which is unique up to homeomorphism (Theorem~\textup{\YMMref{AP}{780}} at $\sigma^{[z]}$), $\dim\ker H_\sigma=\#\mathsf X$, both finite and equal or both infinite, and $\ker H_\sigma=\mathbb C\Omega_\sigma$ if and only if $\sigma$ has one vacuum; these statements read $\mathsf X$ alone, so they do not depend on the choice of $(z,g)$, and at $z=0$ they are those of Theorem~\textup{\YMMref{AP}{780}} for $\sigma$ itself.

\emph{One cutoff.} $\mathcal I$ is the $\sigma$-field of the Borel sets invariant under the translations by $a_K\mathbb Z^3$. By Proposition~\YMMref{106.3}{765}, for $(K,\beta)\in\mathfrak d(g)$ the set $\mathcal G_{K,\beta}$ is a class at the cutoff $K$ that obeys the bound system at $g$: it is not empty; its members are invariant under the isometries of $\mathbb L_K$ and reflection positive at every lattice mirror; with $\nu$ it contains $h\nu/\nu(h)$ for every bounded $\mathcal I$-measurable $h\ge0$ with $\nu(h)>0$, so in particular the conditioned states $\nu(\cdot\,|B)$, $B\in\mathcal I$, $\nu(B)>0$; and it obeys the bound system $\mathsf C^g=(\hat C^g_1,(\hat C^g_n)_{n\ge2})$ at $g$ (equation~\YMMref{(2030)}{765}; Section~\YMMref{100}{743}), which has two members: for $n\ge2$ the cumulants of separated tuples of type $(\mathsf d;R)$, that is, with supports at mutual distances $\ge\mathsf d$ and each in a ball of radius $R$, satisfy $|S^T_n[\rho](f_1,\dots,f_n)|\le\hat C^g_n(\mathsf d;R)\prod_i\|f_i\|_\sharp$ for $\rho\in\mathcal G_{K,\beta}$; and the one-point functions of two members of the class satisfy $|\rho(\mathcal O_K(f))-\nu(\mathcal O_K(f))|\le\hat C^g_1(R)\|f\|_\sharp$ for $\rho,\nu\in\mathcal G_{K,\beta}$ and every $f\in C^1_c(\mathbb R^3)$ with support in a closed ball of radius $R$.

These numbers are free of $K$, $\beta$ and the states. For such a class the plateau theorem, Proposition~\YMMref{100.1}{744}, provides for separated $X$ and $\bar R\in a_K\mathbb N$ numbers $\delta(X;\bar R)>0$, formed from $X$ and $\bar R$, and $\mathsf N(X;\bar R)$, formed from $X$, $\bar R$ and both members of the bound system at $g$, and a constant $\mathfrak k_0$ (equations~\YMMref{(1976)--(1980)}{743}); none of them depends on the cutoff, the coupling or the state. Its clustering inequality says: if $\nu$ belongs to the class, $X,Y$ are separated with supports in $(-\bar R,\bar R)^3$, $a_K\le2\min\{\delta(X;\bar R),\delta(Y;\bar R)\}$, $v\in a_K\mathbb Z^3$ and $|v|\ge4\sqrt3\bar R$, then $X\cdot\tau_vY$ is formed and
\[
\bigl|\sigma_\nu(X\cdot\tau_vY)-\nu\bigl(E_\nu[X^\nu_K|\mathcal I]\,E_\nu[Y^\nu_K|\mathcal I]\bigr)\bigr|\le\mathfrak k_0\,[\mathsf N(X;\bar R)\mathsf N(Y;\bar R)]^{1/2}\,\bar R/|v| .
\]
With it comes the statement on iterated plateaux (Proposition~\YMMref{100.2}{745}, with the conditional means defined immediately before it): the limit $\sigma$ of a convergent Gibbs sequence at $g$ has iterated plateaux, and the plateau of $(X^1,\dots,X^q)$ is $\lim_j\nu_j\bigl(\prod_{i\le q}E_{\nu_j}[(X^i)^{\nu_j}_{K_j}|\mathcal I]\bigr)$. We put $r(X,Y)=4\sqrt3\bar R$ and $C^g(X,Y)=\mathfrak k_0[\mathsf N(X;\bar R)\mathsf N(Y;\bar R)]^{1/2}\bar R$ where $\bar R(X,Y)$ is the least $L^k$, $k\in\mathbb Z$, for which $\operatorname{supp}X\cup\operatorname{supp}Y\subset(-L^k/2,L^k/2)^3$, and is $1$ when both supports are empty; these are the numbers of the clustering of components above, and they depend on $X$, $Y$, the bounds at $g$, $L$ and $M$ alone.

\subsection{Principal results}

\begin{theorem}\label{lh:main}
\begin{enumerate}
\item[(A)] \emph{The limit.}
\begin{enumerate}
\item[(1)] $\beta^4J(\beta)\to0$ as $\beta\to\infty$ through the reals; equivalently, $\beta^4\sup\{|\nu(a_p)-\nu'(a_p)|:\nu,\nu'\in\mathcal G(\beta)\}\to0$.
\item[(2)] Let $(\mathrm S_1)$, $(\mathrm S_2)$ be settings at one radius with $L_2=L_1+2$ (such pairs exist at every radius), and let $r_1>0$ with $13r_1b_{+,i}\le1$, $i=1,2$. There is $C<\infty$, formed from the two settings and $r_1$, such that $\limsup_{s\to\infty}|\mathbb M|([s,s+\delta])\le C\delta$ for every $\delta>0$.
\item[(3)] Let the two settings of (2) have one common value of $\mathfrak p_0$; for every radius and every integer $a\ge1$ there are such pairs with $\mathfrak p_0\ge a$. There are $C_{\rm A},s_{\rm A}<\infty$, formed from the two settings and $r_1$, such that $\beta^4J(\beta)\le C_{\rm A}(\log\beta)^{-\mathfrak p_0/3}$ for $\beta\ge3$, and $|\mathbb M|([s,s+\delta])\le C_{\rm A}\delta$ for $s\ge s_{\rm A}$ and $C_{\rm A}s^{-\mathfrak p_0/3}\le\delta\le1$. Hence for every real $q_0$ a pair of settings with $\mathfrak p_0\ge3q_0$ gives $\sup_{\beta\ge3}(\log\beta)^{q_0}\beta^4J(\beta)<\infty$, and then, in every setting and for every $0<g\le\gamma$, $a_K^{-4}J(\beta)\le C'K^{-q_0}$ for all $(K,\beta)\in\mathfrak d(g)$, with a number $C'$ free of $K$ and $\beta$.
\item[(4)] In either of the two settings of (2) put, for $K\in\mathbb N$ and $b\in\mathsf I$,
\[
u_K(b)=a_K^{-3}\bigl[P(L^Kb)-P_{N_0}(L^Kb)\bigr],\qquad N_0=\ell_0L^K,
\]
the free energy per physical volume of the infinite lattice minus that of the torus $T_{\ell_0}$. Then $-3C_{\rm v}\le u_K\le0$ on $\mathsf I$, with the number $C_{\rm v}$ of Proposition~\ref{lh:blk}, and there are $C_{\rm sc},C_u<\infty$, formed from the two settings and $r_1$, such that $u_K+\frac12C_{\rm sc}b^2$ is convex on $(b_-,b_+)$ for every $K$. Every sequence of cutoffs has a subsequence along which $u_K$ converges uniformly on the compact subsets of $(b_-,b_+)$; every such limit $u$ is continuously differentiable, $|u'(b)-u'(b')|\le C_u|b-b'|$, $-3C_{\rm v}\le u\le0$, $u+\frac12C_{\rm sc}b^2$ is convex; and along the subsequence the one-sided derivatives $u_K'(b\pm)$ converge to $u'(b)$ uniformly on the compact subsets. For $\nu\in\mathcal G_{K,L^Kb}$ and bounded $f$ of compact support,
\[
\nu(\mathcal O_K(f))-\tilde\mu_{K,\ell_0,L^Kb}(\mathcal O_K(f))=-\mathsf u\,r_K(f)\qquad\text{with a number }\mathsf u\in[u_K'(b-),u_K'(b+)],
\]
where $\tilde\mu_{K,\ell_0,\beta}$ is the Wilson measure of $T_{\ell_0}$ at the spacing $a_K$: along the subsequence the energy density of the periodic Gibbs states, counted from that of the torus $T_{\ell_0}$, converges to $-u'(b)$, uniformly in the state and in the label on compact sets.
\end{enumerate}
\item[(B)] \emph{The power.} Let $(\mathrm S_1)$, $(\mathrm S_2)$ and $r_1$ be as in (A)(2), and let $\varsigma'_i$ be a lowered choice of $\varsigma_i$, $i=1,2$.
\begin{enumerate}
\item[(1)] For every real $A>90$ there are $C_{\rm h},C_{\rm r},\beta_1<\infty$, given by \eqref{lh:cst}, which depend only on $A$, $r_1$, the two settings and the pairs $(\varsigma_i,\varsigma'_i)$, such that for $\beta\ge\beta_1$, $s=\log\beta$ and $0\le\delta\le1$
\[
|\mathbb M|([s,s+\delta])\le C_{\rm h}\delta+C_{\rm r}\beta^{-1/2A},\qquad\beta^4J(\beta)\le C_{\rm r}\beta^{-1/2A}.
\]
At $A=91$ the power is $\beta^{-1/182}$.
\item[(2)] Consequently $c_*(c)=\sup_{\beta\ge1}\beta^{4+c}J(\beta)$ is finite for every real $c<\frac1{180}$; it is a number of Wilson's action alone.
\item[(3)] Consequently $\sup_{\beta\ge3}(\log\beta)^q\beta^4J(\beta)<\infty$ for every real $q$, at this one pair of settings.
\item[(4)] The two bounds of (1) hold, with other constants formed from the same objects, also when only $\varsigma'_1$ is given; then Theorem~\YMMref{D}{8} is used at the three choices $\varsigma_1$, $\varsigma'_1$, $\varsigma_2$.
\item[(5)] Let $A>90$ and $C_{\rm h}$, $C_{\rm r}$, $\beta_1$ be as in (1). For $i=1,2$, all $m\ge m_{0,i}$, all $K\in\mathbb N$ with $L_i^Kb_{-,i}\ge\beta_1$ and all $G\in C^2(\mathbb R)$ with support in $\log\mathsf I_i$,
\[
\Bigl|\int G(\upsilon)\,y^{(i)}_{m,K}(\upsilon)\,d\upsilon\Bigr|\le(C_{\rm h}+C_{\rm r}+C_N)\|G\|_1+(C_{\rm h}+C_{\rm r})(L_i^Kb_{-,i})^{-1/2A}\|G'\|_1+c_{{\rm P},i}L_i^{-m}\mathsf Q(G),
\]
with $\mathsf Q(G)=\|G''\|_1+7\|G'\|_1+12\|G\|_1$, $c_{{\rm P},i}$ of Proposition~\ref{lh:blk}(c) and $C_N=\max_iC^{\rm ref}_i$ of Proposition~\ref{lh:blk}(d): in the mean over the labels the far part is bounded on all sides, beyond a cutoff that does not depend on the side.
\end{enumerate}
\item[(C)] \emph{The physical scale.} In any setting let $0<g\le\gamma$, $(K,\beta)\in\mathfrak d(g)$, $\rho,\nu\in\mathcal G_{K,\beta}$ and $0\le c<\frac1{180}$; $e_\rho=\rho(a_p)$ is one number for all plaquettes $p$ of $\mathbb L_K$. For every plaquette $p$ of $\mathbb L_K$ and every bounded $f$ with support in $[-R,R]^3$,
\[
3a_K^{-4}|\rho(a_p)-\nu(a_p)|\le a_K^{-4}J(\beta)\le g^{8+2c}c_*(c)L^{-cK},
\]
\[
|\rho(\mathcal O_K(f))-\nu(\mathcal O_K(f))|=3a_K^{-4}|e_\rho-e_\nu|\,|r_K(f)|\le a_K^{-4}J(\beta)(2R+a_K)^3\|f\|_\infty .
\]
\item[(D)] \emph{One-point functions.} In any setting let $0<g\le\gamma$ and let $(K_j,\beta_j,\nu_j)_j$ be a Gibbs sequence at $g$.
\begin{enumerate}
\item[(1)] For $(K,\beta)\in\mathfrak d(g)$, $a_K^{-4}J(\beta)\le g^8\sup_{\beta'\ge L^Kg^{-2}}\beta'^4J(\beta')$, which tends to $0$ as $K\to\infty$. So $\rho_j(\mathcal O_{K_j}(f))-\nu_j(\mathcal O_{K_j}(f))\to0$ for all $\rho_j\in\mathcal G_j$ and all bounded $f$ of compact support.
\item[(2)] Every limit $\sigma'$ in the frame of the sequence has $\sigma'((f))=0$ for all $f\in C^1_c(\mathbb R^3)$.
\item[(3)] If the sequence converges, with the limit $\sigma$, then $\mathfrak c_\xi=0$ for every canonical decomposition of $\sigma$ and every $\xi$, so $\kappa=0$; every component $\sigma_\xi$ is the limit of the Gibbs sequence $(K_j,\beta_j,\nu_j(\cdot\,|B_j))_j$, along the full sequence, with $B_j\in\mathcal I$, $\nu_j(B_j)>0$, so that Theorem~\textup{\YMMref{AP}{780}} applies to it, with its reconstruction, and it has one vacuum. For a component of such a limit the three properties $\mathfrak c_\xi=0$, ``$\sigma_\xi$ is the limit of a convergent Gibbs sequence at $g$ with the pairs $(K_j,\beta_j)$'' and ``$\sigma_\xi$ is such a limit and has one vacuum'' are equivalent, by an argument that does not use (A).
\item[(4)] In (3), $|S^\sigma_2(f,\tau_vh)|\le C^g((f),(h))/|v|$ for $f,h\in C^1_c(\mathbb R^3)$ and $|v|\ge r((f),(h))$.
\end{enumerate}
\item[(E)] \emph{The rate on a torus.} Let $\varsigma$ be a choice of $\mathcal T^{\rm F+}$, $\varsigma'$ a lowered choice of it, $n\ge2$ an integer, $0<b'\le b''$ and $\mathsf d>0$. There are $k^{\rm rt}\ge0$ and $C^{\rm rt},C^{\rm pol}<\infty$, depending only on the values of the constants at $\varsigma$, on $c_2^\natural$, $\mathsf Y^\natural$, $n$, $b'$, $b''$ and $\mathsf d$, with the following property. Let $m\in\mathbb N$, $\ell=2L^m$, $b\in[b',b'']$, and let $f_1,\dots,f_n$ be real Lipschitz functions on $T_\ell$, none of them $0$, of spread at least $\mathsf d$; put $S_K=S^T_{n;K,\ell,L^Kb}(f_1,\dots,f_n)$, $\Pi=\prod_i\|f_i\|_\sharp$ and $\mathsf p_n=2\mathfrak p_0(n+1)+\frac12$. Then:
\begin{enumerate}
\item[(a)] $|S_{K'}-S_K|\le C^{\rm rt}\ell^{3(n+1)}K^{\mathsf p_n}L^{-K/10}\Pi$ for all integers $K'\ge K\ge k^{\rm rt}+54m$; so $S_\infty=\lim_KS_K$ exists, and $|S_\infty-S_K|$ obeys the same bound;
\item[(b)] $|S_K|\le C^{\rm pol}\ell^{219n}\Pi$ for all $K\in\mathbb N$, and $|S_\infty|\le C^{\rm pol}\ell^{219n}\Pi$;
\item[(c)] for every real $A>30(n+1)$ and all integers $K'\ge K\ge k^{\rm rt}+Am$,
\[
|S_{K'}-S_K|\le2^{3(n+1)}C^{\rm rt}K^{\mathsf p_n}L^{-(\frac1{10}-\frac{3(n+1)}A)K}\Pi,
\]
a bound that does not depend on the side and tends to $0$;
\item[(d)] at $n=2$, $[b',b'']=[2\gamma^{-2},L^2\gamma^{-2}]$ and $\mathsf d=\frac14$, with $k^{\rm rt}$, $C^{\rm rt}$, $C^{\rm pol}$ of these data and
\[
C=\max\bigl\{11,\ 54/\lambda+k^{\rm rt}/\log2,\ C^{\rm rt}(110\mathsf p_2/e\lambda)^{\mathsf p_2},\ C^{\rm pol}\bigr\}:
\]
for all $m\in\mathbb N$, $b\in[2\gamma^{-2},L^2\gamma^{-2}]$, real Lipschitz $f_1,f_2$ on $T_\ell$ whose supports are at sup-distance at least $\frac14$, and all integers $K\ge C\log\ell$, $|S_K-S_\infty|\le C\ell^{438}L^{-K/C}\|f_1\|_\sharp\|f_2\|_\sharp$ and $|S_\infty|\le C\ell^{438}\|f_1\|_\sharp\|f_2\|_\sharp$.
\end{enumerate}
\item[(F)] \emph{The plateau at one cutoff.} In a setting let $0<g\le\gamma$, $(K,\beta)\in\mathfrak d(g)$ and $\nu\in\mathcal G_{K,\beta}$.
\begin{enumerate}
\item[(1)] There is one bounded $\mathcal I$-measurable function $\mathsf e$ with $\mathsf e=E_\nu[a_p|\mathcal I]$ $\nu$-almost surely for every plaquette $p$ of $\mathbb L_K$; and $-3\mathsf e\in[P'(\beta-),P'(\beta+)]$ $\nu$-almost surely.
\item[(2)] Let $f,h\in C^1_c(\mathbb R^3)$, $\bar R\in a_K\mathbb N$ with 
\[
\operatorname{supp}f\cup\operatorname{supp}h\subset(-\bar R,\bar R)^3
\]
 and 
\[
a_K\le2\min\{\delta((f);\bar R),\delta((h);\bar R)\},
\]
and $v\in a_K\mathbb Z^3$ with $|v|\ge4\sqrt3\bar R$. Then
\[
\bigl|S^T_2[\nu](f,\tau_vh)-9a_K^{-8}r_K(f)r_K(h)\operatorname{Var}_\nu\mathsf e\bigr|\le\mathfrak k_0[\mathsf N((f);\bar R)\mathsf N((h);\bar R)]^{1/2}\bar R/|v|,
\]
\[
0\le9a_K^{-8}\operatorname{Var}_\nu\mathsf e\le\tfrac14\bigl(a_K^{-4}J(\beta)\bigr)^2\le\tfrac14\bigl(g^{8+2c}c_*(c)L^{-cK}\bigr)^2\qquad(0\le c<\tfrac1{180}).
\]
\item[(3)] There are $\nu_\pm\in\mathcal G_{K,\beta}$ with $-3e_{\nu_\pm}=P'(\beta\pm)$. For any two such states and $\nu=\frac12(\nu_++\nu_-)$, $9a_K^{-8}\operatorname{Var}_\nu\mathsf e=\frac14(a_K^{-4}J(\beta))^2$. If $\nu$ is a pure phase, or $-3e_\nu\in\{P'(\beta-),P'(\beta+)\}$, then $\operatorname{Var}_\nu\mathsf e=0$.
\item[(4)] Let $X,Y$ be separated tuples, and $\bar R$, $v$ as in (2) with $\delta(X;\bar R)$, $\delta(Y;\bar R)$ in the place of $\delta((f);\bar R)$, $\delta((h);\bar R)$. Then
\[
\bigl|\sigma_\nu(X\cdot\tau_vY)-\sigma_\nu(X)\sigma_\nu(Y)-\operatorname{Cov}_\nu\bigl(E_\nu[X^\nu_K|\mathcal I],E_\nu[Y^\nu_K|\mathcal I]\bigr)\bigr|\le\mathfrak k_0[\mathsf N(X;\bar R)\mathsf N(Y;\bar R)]^{1/2}\bar R/|v|,
\]
and the modulus of the covariance is at most $\frac14\operatorname{osc}(X^\nu_K)\operatorname{osc}(Y^\nu_K)$, where $\operatorname{osc}(F)=\sup_\rho\rho(F)-\inf_\rho\rho(F)$, $\rho$ running through $\mathcal G_{K,\beta}$. For every real continuous local function $F$ of the configuration, $\operatorname{osc}(F)\le D^+\bar p^F(0)-D^-\bar p^F(0)$, the kink at $t=0$ of the pressure $\bar p^F(t)$ perturbed by $tF$ (Lemma~\ref{lh:lat}(b)); at $X=(f)$, $\operatorname{osc}((f)^\nu_K)=a_K^{-4}J(\beta)|r_K(f)|$.
\end{enumerate}
\item[(G)] \emph{The limits at a label.} In a setting let $b>0$, let $\mathfrak p=(K_j,\beta_j)_j$ be a sequence of pairs at the label $b$, $\mathsf K=\mathsf K_{\mathfrak p}$ and $\mathsf K^\bullet=\mathsf K^\bullet_{\mathfrak p}$; for $\sigma\in\mathsf K$, $H_\sigma$ and $\Omega_\sigma$ are those of the paragraph on dilations.
\begin{enumerate}
\item[(1)] $\mathsf K$ is convex, compact and metrisable; the empty set is not excluded.
\item[(2)] Let $\sigma=\lim_j\sigma_{\nu_j}\in\mathsf K$. For every canonical decomposition of $\sigma$ and every $\xi$ there are $B_j\in\mathcal I$ with $\nu_j(B_j)>0$ such that $\sigma_\xi=\lim_j\sigma_{\nu_j(\cdot|B_j)}$ along the full sequence, and $\sigma_\xi$ has one vacuum; so $\sigma_\xi\in\mathsf K^\bullet$. The set $\mathsf K^\bullet$ is compact.
\item[(3)] $\operatorname{bar}:\mathrm{Prob}(\mathsf K^\bullet)\to\mathsf K$ is an affine homeomorphism onto $\mathsf K$. For $\sigma\in\mathsf K$ and every canonical decomposition of $\sigma$, the map $\xi\mapsto\sigma_\xi$ is a homeomorphism of $\mathsf X$ onto the support of the one measure $\mu_\sigma$ with $\operatorname{bar}\mu_\sigma=\sigma$, and it carries $\varpi$ to $\mu_\sigma$; and $\dim\ker H_\sigma=\#\operatorname{supp}\mu_\sigma$, both finite and equal or both infinite.
\item[(4)] For $\sigma\in\mathsf K$ these say the same: $\sigma$ has one vacuum; $\ker H_\sigma=\mathbb C\Omega_\sigma$ in one, and then in each, of the six directions; $\sigma$ is an extreme point of $\mathsf K$; $\sigma(X\cdot\tau_vX)\to\sigma(X)^2$ as $|v|\to\infty$, for every separated $X$; $|\sigma(X\cdot\tau_vY)-\sigma(X)\sigma(Y)|\le c_{X,Y}/|v|$ for $|v|\ge r_{X,Y}$, with numbers $c_{X,Y}$, $r_{X,Y}$ free of $\sigma$ (in a window they are $C^g(X,Y)$, $r(X,Y)$); $\sigma$ is the limit, along the full sequence, of $\sigma_{\omega_j}$ with pure phases $\omega_j\in\mathcal G_{K_j,\beta_j}$; $\operatorname{Var}_{\nu_j}(E_{\nu_j}[X^{\nu_j}_{K_j}|\mathcal I])\to0$ for every separated $X$, for one, and then for every, choice of states $\nu_j\in\mathcal G_{K_j,\beta_j}$ with $\sigma_{\nu_j}\to\sigma$.

For a subsequence $\mathfrak p'$ of $\mathfrak p$, $\mathsf K$ is a closed face of $\mathsf K_{\mathfrak p'}$; if $\tilde{\mathfrak p}$ is a second sequence of pairs at the label $b$ and $\mathfrak p''$ interlaces $\mathfrak p$ and $\tilde{\mathfrak p}$, then $\mathsf K_{\mathfrak p''}=\mathsf K\cap\mathsf K_{\tilde{\mathfrak p}}$. If $\mathsf K\ne\emptyset$, these say the same: $\mathsf K$ is a point; $\mathsf K^\bullet$ is a point; $\mathsf K=\mathsf K^\bullet$; $\operatorname{Cov}_{\nu_j}(E_{\nu_j}[X^{\nu_j}_{K_j}|\mathcal I],E_{\nu_j}[Y^{\nu_j}_{K_j}|\mathcal I])\to0$ for all separated $X,Y$ and all choices of $\nu_j\in\mathcal G_{K_j,\beta_j}$ for which $\sigma_{\nu_j}$ converges. At single fields this covariance tends to $0$ in every case:
\[
\bigl|\operatorname{Cov}_{\nu_j}\bigl(E_{\nu_j}[(f)^{\nu_j}_{K_j}|\mathcal I],E_{\nu_j}[(h)^{\nu_j}_{K_j}|\mathcal I]\bigr)\bigr|\le\tfrac14\bigl(a_{K_j}^{-4}J(\beta_j)\bigr)^2|r_{K_j}(f)r_{K_j}(h)|\to0 .
\]
\item[(5)] ${\rm Lim}[b]$ has exactly one element if and only if ${\rm Lim}^\bullet[b]$ has exactly one element. The same holds for ${\rm Lim}_{\mathfrak q}[b]$ and ${\rm Lim}^\bullet_{\mathfrak q}[b]$ for every set $\mathfrak q$ of pairs that contains the pairs of a sequence of pairs at the label $b$; in particular for $\mathfrak q=\{(K,L^Kb):K\in\mathbb N\}$. If $g^{-2}+\bar B<b<X_g-\bar B$ for some $0<g\le\gamma$, a third equivalent statement is that every Gibbs sequence at $g$ with $a_{K_j}\beta_j\to b$ converges; and then the one limit has one vacuum.
\item[(6)] ${\rm Lim}[b]$ and ${\rm Lim}^\bullet[b]$ are non-empty, compact and metrisable. Every sequence of pairs at the label $b$ has a saturated subsequence, and every subsequence of a saturated sequence is saturated, with the same set of limits; ${\rm Lim}[b]$ is the union of the sets $\mathsf K_{\mathfrak p}$, and ${\rm Lim}^\bullet[b]$ the union of their sets of extreme points, over the saturated sequences $\mathfrak p$ at the label $b$. Every member of ${\rm Lim}[b]$ is the barycentre of exactly one Borel probability measure on ${\rm Lim}^\bullet[b]$. The map $\operatorname{bar}$ is an affine homeomorphism of $\mathrm{Prob}({\rm Lim}^\bullet[b])$ onto the closed convex hull $\overline{\rm co}\,{\rm Lim}[b]$ of ${\rm Lim}[b]$ for the pointwise convergence, and for $\sigma\in{\rm Lim}[b]$ these say the same: $\sigma$ has one vacuum; $\sigma$ is an extreme point of $\overline{\rm co}\,{\rm Lim}[b]$; $\sigma$ is an extreme point of ${\rm Lim}[b]$.

The following four statements are equivalent, and they hold if ${\rm Lim}[b]$ has exactly one element: ${\rm Lim}[b]$ is convex; ${\rm Lim}[b]=\mathsf K_{\mathfrak q}$ for one sequence of pairs $\mathfrak q$ at the label $b$; ${\rm Lim}[b]=\operatorname{bar}(\mathrm{Prob}({\rm Lim}^\bullet[b]))$; every finite subset of ${\rm Lim}^\bullet[b]$ lies in $\mathsf K^\bullet_{\mathfrak p}$ for one sequence of pairs $\mathfrak p$ at the label $b$.
\item[(7)] Let $0<g\le\gamma_W$, where $\gamma_W$ is the ceiling of the law of two points, the short-distance law for the two-point function (equation~\YMMref{(2346)}{850}); $\gamma_W<\gamma$ and $\gamma_W^{-2}\ge L^2\gamma^{-4}$. Let $\mathfrak p$ be a sequence of pairs in $\mathfrak d(g)$; no label is supposed. If $\mathsf K_{\mathfrak p}\ne\emptyset$, then $a_{K_j}\beta_j$ converges, to a label $b$ in the interval $I^{\rm lab}(g)$ of Section~\ref{sec:state-limit-conventions}, and for every $\sigma\in\mathsf K_{\mathfrak p}$, every $\vartheta\in(0,1)$ and every pair $(f,h)$ of the set $\mathfrak T^1_{\mathbb R^3}(\mathsf d_0(g,\vartheta);\vartheta)$, which is not empty,
\[
b^{-2}=\lim_{t\to0}\frac{S^\sigma_2(f_t,h_t)}{\mathsf c_2\,\mathcal K_{\rm free}(f,h)},\qquad f_t(x)=f(x/t),\qquad h_t(x)=h(x/t).
\]
Here $\mathfrak T^1_{\mathbb R^3}(\mathsf d;\vartheta)$ is the set of the test pairs of class $C^1$ at the scale $\mathsf d$ with the shape parameter $\vartheta$ (Lemma~\ref{lem:ward-test-pairs} and~the test-pair convention preceding equation~\YMMref{(2343)}{849}): pairs of non-negative $C^1$ functions $f,h$, neither of them $0$, with supports in balls of radius $\mathsf d$ at a distance between $\mathsf d$ and $6\mathsf d$, $\int f\ge\vartheta\|f\|_\infty\mathsf d^3$ and $\operatorname{Lip}f\le\|f\|_\infty/\vartheta\mathsf d$, and the same for $h$; $\mathsf d_0(g,\vartheta)>0$ is the threshold of equations~\YMMref{(2349)--(2350)}{851}; and $\mathsf c_2=\frac9{16\pi^2}$ and $\mathcal K_{\rm free}(f,h)=\iint f(x)h(y)|x-y|^{-6}\,dx\,dy>0$ are the constant and the free kernel of Theorem~\YMMref{C}{7} and equation~\YMMref{(2240)}{822}. So the sets ${\rm Lim}_{\mathfrak d(g)}[b]$, $b>0$, are pairwise disjoint and their union is the set of the limits of all convergent Gibbs sequences at $g$; a sequence of pairs in $\mathfrak d(g)$ that interlaces two sequences with different labels has $\mathsf K_{\mathfrak p}=\emptyset$; and a convex combination $t\sigma+(1-t)\sigma'$, $0<t<1$, of limits at two different labels is the limit of no Gibbs sequence at $g$.
\end{enumerate}
\end{enumerate}
\end{theorem}

The second theorem has a hypothesis on two periodic boxes. Let $Z(n_1,n_2,n_3)$ be the partition function of Wilson's action at the coupling $\beta$ on the periodic lattice with $n_1,n_2,n_3$ sites in the three directions; it does not depend on the spacing. In a setting let $b>0$, $\ell_*=2L^{m_*}$ with $m_*\in\mathbb Z$, $K_*$ an integer with $K_*+m_*\ge1$, and $0<x_0\le0.069$. The box inequality $\mathrm{BI}[b,\ell_*,K_*;x_0]$ says: for every integer $K\ge K_*$, with $l_K=\ell_*/2a_K=L^{K+m_*}$,
\[
\mathcal E(l_K)=2\log Z(l_K,2l_K,2l_K)-\log Z(2l_K,2l_K,2l_K)\le x_0\qquad\text{at }\beta=L^Kb .
\]

\begin{theorem}\label{lh:box}
In a setting suppose $\mathrm{BI}[b,\ell_*,K_*;x_0]$, and let $\underline K=\max\{K_*,1\}$. Then:
\begin{enumerate}
\item[(1)] for every $K\ge\underline K$ the set $\mathcal G_{K,L^Kb}$ has exactly one element $\nu_K$, which is a pure phase, and $J(L^Kb)=0$;
\item[(2)] for every $n\ge2$ and every tuple $X$ of $n$ real Lipschitz functions on $\mathbb R^3$ with pairwise disjoint compact supports, $\lim_KS^T_n[\nu_K](X)$ exists along all $K\ge\underline K$ and equals $\lim_{m'\to\infty}\lim_{K\to\infty}S^T_{n;K,\ell',L^Kb}(X^{\langle\ell'\rangle})$, $\ell'=2L^{m'}$;
\item[(3)] for $\mathfrak p=(K,L^Kb)_{K\ge\underline K}$ the set $\mathsf K_{\mathfrak p}$ is one point $\sigma=\lim_K\sigma_{\nu_K}$, and $\sigma$ has one vacuum.
\end{enumerate}
For every $z\in\mathbb Z$, $\mathrm{BI}[L^zb,L^z\ell_*,K_*-z;x_0]$ is the same statement as $\mathrm{BI}[b,\ell_*,K_*;x_0]$, so (1)--(3) hold at the label $L^zb$ as well.
\end{theorem}

\subsection{Construction estimates}

The proof starts from Theorem~\YMMref{D}{8}, with the common numerical choice of Proposition~\YMMref{69.1}{595} and its actual-history application in Corollary~\YMMref{69.4}{604}, and from what is derived from it on the tori: the discarding of the histories with a large-field region (Proposition~\textup{\YMMref{115.2}{806}}), with its weight (the definitions preceding Proposition~\YMMref{115.2}{806}) and with the statement on two nearby positive functionals and their cumulants (Lemma~\ref{lem:positive-functional-comparison}), by which the discarding holds at every bound on the lost mass that is at most $\frac12$; the comparison of two consecutive cutoffs (Theorem~\YMMref{AR}{810}); and the convergence on each torus (Theorem~\textup{\YMMref{AS}{813}}). The last says, at every choice: let $\ell=2L^m$, $n\ge2$, let $f_1,\dots,f_n$ be real Lipschitz functions on $T_\ell$, none of them $0$, whose supports have no common point, let $0<b'\le b''$, let $\mathsf d>0$ be at most the spread of $(f_1,\dots,f_n)$, and let $s\ge0$ be an integer with $L^sb'\ge\gamma^{-2}$ and $M\mathcal R(L^sb'')\le L^{m+s}$; such integers exist (Lemma~\YMMref{118.2}{813}, equation~\YMMref{(2210)}{813}), and $s=0$ is one of them when $b'\ge\gamma^{-2}$ and $M\mathcal R(b'')\le L^m$. Then for all $b\in[b',b'']$ and all integers $K'\ge K\ge K_{1,s}(\mathsf d)+s$
\[
\bigl|S^T_{n;K',\ell,L^{K'}b}(f_1,\dots,f_n)-S^T_{n;K,\ell,L^Kb}(f_1,\dots,f_n)\bigr|\le L^{sn}C'_s\,(L^s\ell)^{3(n+1)}(K-s)^{\mathsf p_n}L^{-(K-s)/10}\prod_i\|f_i\|_\sharp ,
\]
so that $S^T_{n;K,\ell,L^Kb}(f_1,\dots,f_n)$ converges as $K\to\infty$, uniformly in $b\in[b',b'']$. Here $C'_s=\mathsf C'_n([L^sb',L^sb''])$ and $K_{1,s}(\mathsf d)=K_1(L^s\ell,n,[L^sb',L^sb''],L^s\mathsf d)$ are numbers of the constants and thresholds in Theorem~\YMMref{AR}{810}, which attaches to a side $\ell$, to $n$, to an interval $\mathsf J$ of labels with $\min\mathsf J\ge\gamma^{-2}$ and $M\mathcal R(\max\mathsf J)\le L^m$, and to $\mathsf d>0$ the thresholds $K_0(\ell,n,\mathsf J)\le K_1(\ell,n,\mathsf J,\mathsf d)$ and the rate constant $\mathsf C'_n(\mathsf J)$. At $s=0$ the bound reads $\mathsf C'_n([b',b''])\,\ell^{3(n+1)}K^{\mathsf p_n}L^{-K/10}\prod_i\|f_i\|_\sharp$ for $K'\ge K\ge K_1(\ell,n,[b',b''],\mathsf d)$; for other sides and windows it carries the dilation exponent $s$.

For a choice and a lowered choice, the setting (Th) gives an admissible lattice by Lemma~\YMMref{118.1}{812}. Section~\ref{sec:local-large-field-rarity} supplies the common lattice and background definitions, the explicit threshold, the witnesses for each creation, the gauge integration, and the chessboard estimate. Its Proposition~\ref{prop:local-large-field-rarity} gives the volume-uniform probability bound. The history weights and masses are those of Theorem~\YMMref{D}{8}.

From Theorems~\YMMref{AJ}{673} and~\YMMref{AD}{607}, it uses the zero-source stability bounds (ii), (iii) of Theorem~\textup{\YMMref{AD}{607}}; the comparison of two laws in one ball of radius $1$ at a first-passage datum, the common-normalizer comparison in the proof of Proposition~\YMMref{91.13}{709}, whose hypothesis (Nr1) Proposition~\YMMref{91.12}{708} supplies at every pair of $\mathfrak d(g)$; the bound at separated supports in every periodic Gibbs state, Proposition~\textup{\YMMref{91.13}{709}}; and the family of bounds at $g$ (Proposition~\YMMref{106.3}{765}, with equation~\YMMref{(2030)}{765}); for the far part on all sides also the bound for connected functions on the tori at first-passage data, Proposition~\YMMref{89.5}{694}.

For the limits it uses Theorem~\textup{\YMMref{AP}{780}}, Proposition~\YMMref{100.1}{744} with equations~\YMMref{(1976)--(1980)}{743}, the statement on iterated plateaux (Proposition~\YMMref{100.2}{745}), and Proposition~\textup{\YMMref{106.1}{764}} (pure phases drawn from a regular conditional law given $\mathcal I$); for (G)(7) the law of two points: its test pairs, its constant and its free kernel (Theorem~\YMMref{C}{7} and equation~\YMMref{(2240)}{822}, \ref{lem:ward-test-pairs} and~the convention in equation~\YMMref{(2343)}{849}), with non-empty sets of test pairs (Lemma~\ref{lem:ward-test-pairs}(c)), its ceiling and its threshold (equations~\YMMref{(2346)}{850} and~\YMMref{(2349)--(2350)}{851}), the law in the limits through periodic Gibbs states, with the convergence of the labels (Proposition~\YMMref{127.8}{854}), and the law in every component (Corollary~\textup{\YMMref{127.11}{854}}); for the box inequality the domination theorem (Proposition~\YMMref{91.9}{705}), the seminorm as a mean over a period box (equation~\eqref{eq:period-box-seminorm}), the chessboard estimate on the tori of $\mathbb L_K$ and the comparison of two sides (Corollary~\ref{IV-18-18}(b), (c), (e)), a coarse and a fine board (the proof of Lemma~\YMMref{104.1}{759}), the functions of a cell (the plaquette-support argument in Section~\ref{sec:periodic-chessboard}), the chessboard estimate (Proposition~\YMMref{86.3}{681}), and the standing choice of an odd blocking factor $L\ge13$.

The choices used here are supplied by the ordered proof of Proposition~\YMMref{69.1}{595} and its sourced extension in Section~\YMMref{82}{665}. The exponents are fixed before the block size; every sufficiently large odd block size is then admitted, and the remaining constants are chosen in the stated order. Besides, we use the first-passage data at $g$ and the canonical family $\mathfrak d(g)$ (Corollary~\YMMref{74.2}{616} and~\ref{sec:state-limit-conventions}) with the two facts on its labels (Section~\ref{sec:state-limit-conventions}), the cube size $\mathcal R$ (equation~\YMMref{(634)}{339}), the dilation of the tori, of their partition functions and connected functions and of the labels (Lemma~\YMMref{118.2}{813}), the dilation of the lattices $\mathbb L_K$, of the laws on them, of their observables and of their periodic Gibbs states (Lemma~\ref{lem:lattice-law-dilation}), the regions of a history, with $Z_t=\Lambda_t^c$ defined by equation~\YMMref{(644)}{341}, and statements on one lattice that use neither Theorem~\YMMref{D}{8} nor Theorems~\YMMref{AJ}{673} and~\YMMref{AD}{607}: the Gibbs states, the periodic Gibbs states and the pure phases (Section~\YMMref{91}{697} and Proposition~\YMMref{106.1}{764}); periodic laws and their bond representatives (equation~\YMMref{(1817)}{698} and the construction immediately preceding it); the lower tangent inequality for periodic Gibbs states (Lemma~\YMMref{91.6}{702}); weak limit points of torus measures are periodic Gibbs states (Proposition~\YMMref{91.2}{699}), and weak limits of Gibbs states at couplings $\beta_k\to\beta$ are Gibbs states at $\beta$ (Proposition~\YMMref{91.1}{698}, with the family of the functions $-a_p$); every periodic Gibbs state is invariant under the isometries of the lattice and reflection positive at every lattice mirror (Proposition~\YMMref{103.4}{758}), and the periodic Gibbs states at one coupling form a convex set (Propositions~\YMMref{91.1}{698} and~\YMMref{103.4}{758}) whose members are integrals of its ergodic members (Proposition~\YMMref{106.2}{764}); reflection positivity of the torus measures at the coordinate mirrors (Proposition~\YMMref{91.4}{700}, with $\Gamma=a(N_1\mathbb Z\times N_2\mathbb Z\times N_3\mathbb Z)$ and height period $N_\mu a$ in direction $e_\mu$); the reconstruction of the transfer operator with its spectral measure (Lemma~\YMMref{99.1}{737}); the trace formula for boxes of slabs (Proposition~\YMMref{85.2}{679}, with the spectral argument in the proof of Proposition~\YMMref{86.5}{682} and the slab operators of equation~\YMMref{(1744)}{679}); and the crossing operator on the Hilbert space of one layer (Section~\YMMref{85}{677}), which is positive and of trace class (Proposition~\YMMref{85.1}{678}, with the rank-one expansion in its proof).

Lemmas~\ref{lh:lat}, \ref{lh:rp}, \ref{lh:cond}, Propositions~\ref{lh:per}, \ref{lh:slope}, \ref{lh:many}, Lemma~\ref{lh:rec}(a), (c) and Proposition~\ref{lh:onepair}(a)--(e) use neither Theorem~\YMMref{D}{8} nor Theorems~\YMMref{AJ}{673} and~\YMMref{AD}{607}, and Proposition~\ref{lh:onepair}(a)--(e) does not use the supply of choices either. Clause (E) uses Theorem~\YMMref{D}{8} alone, at a choice and at a lowered choice of it. Clause (A) uses Theorems~\YMMref{D}{8}, \YMMref{AJ}{673} and~\YMMref{AD}{607} at the choices of its two settings, Theorem~\YMMref{D}{8} through the convergence on each torus; it uses no lowered choice and no statement on the rarity of a large field. Clause (B), in (1)--(3) and (5) and with the constants \eqref{lh:cst}, uses Theorems~\YMMref{AJ}{673} and~\YMMref{AD}{607} at the two choices $\varsigma^{\rm S}_1$, $\varsigma^{\rm S}_2$ of its settings and Theorem~\YMMref{D}{8} at the four choices $\varsigma_1$, $\varsigma_2$, $\varsigma'_1$, $\varsigma'_2$, at each pair $(\varsigma_i,\varsigma'_i)$ through (E).

With other constants, (B)(4) uses Theorems~\YMMref{AJ}{673} and~\YMMref{AD}{607} at the same two choices and Theorem~\YMMref{D}{8} at the three choices $\varsigma_1$, $\varsigma'_1$, $\varsigma_2$ alone: at $(\varsigma_1,\varsigma'_1)$ through (E), and at $\varsigma_2$ through the convergence on each torus, by Proposition~\ref{lh:slope}. In (A) and (B) Theorems~\YMMref{AJ}{673} and~\YMMref{AD}{607} enter only through the stability bounds at zero source (Theorem~\textup{\YMMref{AD}{607}}), the comparison of two laws in one unit ball (the common-normalizer comparison in the proof of Proposition~\YMMref{91.13}{709}, with Proposition~\YMMref{91.12}{708}) and the bound at separated supports in every periodic Gibbs state (Proposition~\textup{\YMMref{91.13}{709}}).

Clause (C) uses, besides (B)(2), only that a setting forms $\mathfrak d(g)$; its setting need not be one of the settings of (B). Clauses (D), (F) and (G) use Theorems~\YMMref{D}{8}, \YMMref{AJ}{673} and~\YMMref{AD}{607} at the choice of the setting, together with (A) where the decay of $a_K^{-4}J(\beta)$ to $0$ enters and with (B) where its rate enters; the bound by $\frac14(a_K^{-4}J(\beta))^2$ in (F) uses the choice of the setting alone. In Theorem~\ref{lh:box}, clause (1) uses neither Theorem~\YMMref{D}{8} nor Theorems~\YMMref{AJ}{673} and~\YMMref{AD}{607} nor the supply of choices; clause (2) uses besides Theorem~\YMMref{D}{8} at the choice of the setting, through the convergence on each torus in Proposition~\ref{lh:onepair}(f); and clause (3), one vacuum, uses besides Theorem~\textup{\YMMref{AP}{780}} and the family of bounds at $g$ (Proposition~\YMMref{106.3}{765}) through (G)(4), that is, Theorems~\YMMref{D}{8}, \YMMref{AJ}{673} and~\YMMref{AD}{607} at the choice of the setting.

It does not use (A): by (1), $J(L^Kb)=0$, so that $a_K^{-4}J(L^Kb)=0$ along the sequence.

\subsection{Auxiliary results}

\begin{lemma}[One lattice]\label{lh:lat}
\begin{enumerate}
\item[(a)] For $\beta\ge0$ the limit $P(\beta)=\lim_NP_N(\beta)$ exists along all integers $N\ge3$. Each $P_N$ is smooth and convex, $-6\beta\le P_N\le0$, $P_N'=-N^{-3}\mu_{N,\beta}(S_N)\in[-6,0]$ and $P_N''=N^{-3}\operatorname{Var}_{\mu_{N,\beta}}(S_N)$. $P$ is finite and convex, its difference quotients lie in $[-6,0]$, and $0\le J\le6$.
\item[(b)] Let $\beta>0$ and $\nu\in\mathcal G(\beta)$. Then $\nu(a_p)$ is one number $e_\nu$ for all $p$, and $P(\beta+t)-P(\beta)\ge-3e_\nu t$ for $t>-\beta$; so $-3e_\nu\in[P'(\beta-),P'(\beta+)]$, and $3|e_\rho-e_\nu|\le J(\beta)$ for $\rho,\nu\in\mathcal G(\beta)$. There are $\nu_\pm\in\mathcal G(\beta)$ with $-3e_{\nu_\pm}=P'(\beta\pm)$; hence $J(\beta)=3\max\{|e_\rho-e_\nu|:\rho,\nu\in\mathcal G(\beta)\}$. More generally, let $F$ be a real continuous function of the bonds in a cube of side $k\in\mathbb N$, let $F_N$ be the sum of its translates over a period cell of the cube with $N^3$ sites, and $p^F_N(t)=N^{-3}\log\mu_{N,\beta}(e^{tF_N})$, $\bar p^F=\limsup_Np^F_N$. These functions are convex, $|p^F_N(t)|\le|t|\|F\|_\infty$, and for every $\nu\in\mathcal G(\beta)$, $N\ge3$ and real $t$
\[
p^F_N(t)\ge t\,\nu(F)-\bigl[144\beta+12(2k-1)^3|t|\,\|F\|_\infty\bigr]/N ;
\]
so $D^-\bar p^F(0)\le\nu(F)\le D^+\bar p^F(0)$. The same holds on $\mathbb L_K$ for $\mathcal G_{K,\beta}$.
\item[(c)] For laws $\mathfrak s$ invariant under $\mathbb Z^3$ and integers $\eta\ge1$: $0\le c_x\le6$, $|\Gamma_{\mathfrak s}|\le9$, $|s^\eta_{\mathfrak s}|\le9\eta^6$ and $|n_\eta(\mathfrak s)|\le1125\eta^3$; on the periodic configurations $S_N=\sum_xc_x$ over a period cell, and $P_N''=\sum_z\Gamma_{\mu_{N,\beta}}(z)$ over a period cell. If $\eta$ divides $N$ and $N\ge7\eta$, then, with $\mathcal D=\{X\in\mathbb Z^3:-\frac N{2\eta}\le X_i<\frac N{2\eta}\}$,
\begin{equation}\label{lh:split}
P_N''(\beta)=n_\eta(\mu_{N,\beta})+\eta^{-3}\sum_{X\in\mathcal D,\ |X|_\infty\ge3}s^\eta_{\mu_{N,\beta}}(X).
\end{equation}
\end{enumerate}
\end{lemma}

\begin{proof}[Proof sketch]
(a) Removing $n$ plaquette terms changes $\log Z$ by at most $2\beta n$, and a bond that no remaining term reads integrates to $1$. A torus and the free box of the same side differ by at most $6N^2$ plaquettes, so their logarithms differ by at most $12\beta N^2$. A free box of side $N'=qN+r'$, $0\le r'<N$, holds $q^3$ disjoint free boxes of side $N$ and at most $9N'^2N+6N'^3/N$ further plaquettes; so the free energies per site $p^{\rm f}_N$ of free boxes satisfy $|p^{\rm f}_{N'}-(qN/N')^3p^{\rm f}_N|\le2\beta(9N/N'+6/N)$, hence $\limsup_{N'}|p^{\rm f}_{N'}-p^{\rm f}_N|\le12\beta/N$, and $(P_N)$ is a Cauchy sequence. As $0\le a_p\le2$, $e^{-6\beta N^3}\le Z_N\le1$; the derivatives are those of the logarithm of a Laplace transform on a compact space; convexity and the bounds on the difference quotients pass to the limit.
(b) Periodic Gibbs states are invariant under the lattice isometries, which act transitively on the plaquettes. The lower tangent inequality (Lemma~\YMMref{91.6}{702}), applied on the cube with $N^3$ sites, whose boundary layer holds at most $6N^2$ sites, gives $\log\mu_{N,\beta}(e^{tF_N})\ge tN^3\nu(F)-144\beta N^2-12(2k-1)^3|t|\|F\|_\infty N^2$; divide by $N^3$. Convexity follows from H\"older's inequality, and the bound by $|t|\|F\|_\infty$ from $|F_N|\le N^3\|F\|_\infty$; letting $N\to\infty$ gives $\bar p^F(t)\ge t\nu(F)$, which is the statement on the one-sided derivatives. At $F=-(a_{p_{12}(0)}+a_{p_{13}(0)}+a_{p_{23}(0)})$, the three plaquettes at the site $0$, we have $k=1$, $\|F\|_\infty\le6$, $F_N=-S_N$ and $\mu_{N,\beta}(e^{tF_N})=Z_N(\beta+t)/Z_N(\beta)$, so $p^F_N(t)=P_N(\beta+t)-P_N(\beta)$ and $\nu(F)=-3e_\nu$:
\[
\log Z_N(\beta+t)-\log Z_N(\beta)\ge-3e_\nu tN^3-144\beta N^2-72|t|N^2 ,
\]
and $-3e_\nu$ is a subgradient of $P$ at $\beta$. For the ends take points $\beta_k\downarrow\beta$ (or $\beta_k\uparrow\beta$) at which $P$ is differentiable and $\nu_k\in\mathcal G(\beta_k)$ invariant under $\mathbb Z^3$, weak limit points of torus measures (Proposition~\YMMref{91.2}{699}): $-3e_{\nu_k}=P'(\beta_k)\to P'(\beta\pm)$, and weak limit points of $(\nu_k)$ are Gibbs states at $\beta$ (Proposition~\YMMref{91.1}{698}) invariant under $\mathbb Z^3$, so periodic Gibbs states at $\beta$. On $\mathbb L_K$ use the dilation by $L^K$, which maps $\mathcal G_{K,\beta}$ onto $\mathcal G(\beta)$ (Lemma~\ref{lem:lattice-law-dilation} at $z=K$).
(c) $0\le c_x\le6$ since $0\le a_p\le2$, and a variable with values in $[0,6]$ has variance at most $9$. Next, $\sum_j\Lambda_\eta(u-\eta j)=1$, $\sum_{x\in\mathbb Z}\Lambda_\eta(x)=\eta$, and $\sum_{x\in\mathbb Z}\Lambda_\eta(x)\Lambda_1(u-x)=\Lambda_\eta(u)$, the left member being the piecewise affine interpolation of $\Lambda_\eta$ at the integers; so $\mathcal O_1(\varphi^\eta_X)=\sum_x\varphi^\eta_{\rm c}(x-\eta X)c_x$, which gives $|s^\eta_{\mathfrak s}|\le9\eta^6$ and, with the $125$ points $|X|_\infty\le2$, $|n_\eta|\le1125\eta^3$. Every plaquette is a face of exactly two unit cubes, so $S_N=\sum_xc_x$ and $P_N''=N^{-3}\operatorname{Var}(S_N)=\sum_z\Gamma(z)$. On the periodic configurations $\sum_{X\in\mathcal D}\mathcal O_1(\varphi^\eta_X)=S_N$, $\mathcal D$ being a set of representatives of $\mathbb Z^3$ modulo $N/\eta$; by invariance $\operatorname{Cov}(\mathcal O_1(\varphi^\eta_0),S_N)=\eta^3\operatorname{Cov}(c_0,S_N)=\eta^3P_N''$. Split $\mathcal D$ into $|X|_\infty\le2$, which gives the hull (all these points lie in $\mathcal D$ as $N/2\eta\ge3.5$), and the rest.
\end{proof}

\begin{lemma}[Reflection positivity: the specific heat is at least its near part]\label{lh:rp}
\begin{enumerate}
\item[(a)] Let $\beta\ge0$, $\nu\in\mathcal G(\beta)$ and $\eta\ge1$ an integer. There is a finite positive Borel measure $\mathsf m$ on $[0,1]\times S^1\times S^1$ with
\[
s^\eta_\nu(n+2,w)=\int t^{\eta n}\zeta_2^{\eta w_2}\zeta_3^{\eta w_3}\,d\mathsf m(t,\zeta_2,\zeta_3)\qquad(n\ge0,\ w\in\mathbb Z^2);
\]
so $s^\eta_\nu(3e_1)\ge0$, and $\sum_{|X|_\infty\ge3}e^{-\epsilon|X|_1}s^\eta_\nu(X)\ge-870\,s^\eta_\nu(3e_1)$ for every $\epsilon>0$.
\item[(b)] Let $I\subset(0,\infty)$ be an open interval, $\eta\ge1$ an integer and $q$ a Borel function, bounded on the compact subsets of $I$, with $q(\beta)\le n_\eta(\nu)-870\eta^{-3}s^\eta_\nu(3e_1)$ for all $\beta\in I$ and $\nu\in\mathcal G(\beta)$. Then $P''(d\beta)\ge q(\beta)\,d\beta$ on $I$.
\item[(c)] Let $\mathfrak s$ be invariant under $\mathbb Z^3$, $\eta,\eta'\ge1$ integers, $\bar\eta=\max\{\eta,\eta'\}$, $w=\varphi^\eta_{\rm c}/\eta^3$, $w'=\varphi^{\eta'}_{\rm c}/\eta'^3$ (non-negative, each of sum $1$), and $\Theta(y)=\prod_i\min\{1,(5-|y_i-\frac12|/\bar\eta)_+\}$. Then
\begin{equation}\label{lh:tent}
n_\eta-n_{\eta'}=\sum_v\Bigl[w'(v)\frac{\langle\varphi^\eta_0;\Theta-\Theta(\cdot+v)\rangle}{\eta^3}-w(v)\frac{\langle\varphi^{\eta'}_0;\Theta-\Theta(\cdot+v)\rangle}{\eta'^3}\Bigr]-\frac{\langle\varphi^\eta_0;\Theta-h^\eta\rangle}{\eta^3}+\frac{\langle\varphi^{\eta'}_0;\Theta-h^{\eta'}\rangle}{\eta'^3},
\end{equation}
all brackets in $\mathfrak s$; in each bracket the second entry vanishes within the sup-distance $\min\{\eta,\eta'\}$ of the support of the first, and off $\{|y-(\frac12,\frac12,\frac12)|_\infty\le6\bar\eta\}$.
\item[(d)] Let $N\ge3$, $\beta\ge0$, $\delta\ge0$, and let $\omega_1,\omega_2,\omega_3$ be real, even, summable functions on $\mathbb Z$ with non-negative Fourier transforms and $\sum_j\omega_i(jN)\le1+\delta$. Then $P_N''(\beta)\ge(1+\delta)^{-3}\sum_{z\in\mathbb Z^3}\omega_1(z_1)\omega_2(z_2)\omega_3(z_3)\Gamma_{\mu_{N,\beta}}(z)$.
\end{enumerate}
\end{lemma}

\begin{proof}[Proof sketch]
(d) Put $\omega_i'(t)=(1+\delta)^{-1}\sum_j\omega_i(t+jN)$ on $\mathbb Z/N$; its Fourier coefficients are non-negative, so $|\omega_i'|\le\omega_i'(0)\le1$. Reflection positivity of the torus measure at the coordinate mirrors makes every slice sum $t\mapsto\sum_{u}e^{2\pi i(k_1u_1+k_2u_2)/N}\Gamma(t,u)$, the sum over the two other coordinates, non-negative: for $1\le t\le N-1$ the mirror $\{x_3=\frac{t+1}2\}$ carries the cubes of the layer $t$ to those of the layer $0$, and at $t=0$ the sum is a square. Writing $1-\omega_1'\omega_2'\omega_3'=(1-\omega_3')+\omega_3'(1-\omega_2')+\omega_3'\omega_2'(1-\omega_1')$ and expanding the factors $\omega'$ in characters with non-negative coefficients, each term gives a non-negative sum against $\Gamma$; so $\sum_z\Gamma(z)\ge\sum_z\omega'_1\omega'_2\omega'_3\Gamma(z)$, and the left member is $P_N''$ by Lemma~\ref{lh:lat}(c).
(a) $\nu$ is invariant and reflection positive at the mirrors $\{x_1=\frac12\}$ and $\{x_1=1\}$. The reconstruction of the transfer operator $T$, $0\le T\le1$, with the step $e_1$, and of the tangential translations, with their joint spectral measure, gives $\mathsf m=\langle[F],\mathfrak E(\cdot)[F]\rangle$ for $F=\mathcal O_1(\varphi^\eta_{e_1})-\nu(\mathcal O_1(\varphi^\eta_{e_1}))$, which reads only bonds in $\{x_1\ge1\}$ and whose mirror image is the centred $\mathcal O_1(\varphi^\eta_{-e_1})$; this is the representation, and $s^\eta_\nu(3e_1)=\int t^\eta d\mathsf m\ge0$. Distribute the $X$ with $|X|_\infty\ge3$ over the six cones about the coordinate half-axes, with the weights $\kappa(k,w)=1$ inside a cone, $\frac12$ on a face and $\frac13$ on an edge, which sum to one.

By the invariance under the permutations and reflections of the axes the left member becomes $6e^{-2\epsilon}\int\Sigma\,d\mathsf m$, where, with $x=e^{-\epsilon}t^\eta$, $e^{iq_i}=\zeta_i^\eta$ and $\Psi_x(w)=\sum_{k\ge3}x^{k-2}\kappa(k,w)$, $\Sigma=\sum_we^{-\epsilon(|w_2|+|w_3|)}e^{iq\cdot w}\Psi_x(w)$. For $x\le\frac12$, $|\Sigma|\le\sum_w\Psi_x(w)=x\sum_{i\ge0}x^i(4(i+3)^2+\frac13)\le145x$. For $x\ge\frac12$, $x^2\Psi_x=\Phi_x-\mathsf R_x$ with $\Phi_x(w)=x^{|w|_\infty}\bigl(\frac{1+x}{2(1-x)}-\frac16\mathbf 1_{\{|w_2|=|w_3|\}}\bigr)$ and a remainder $\mathsf R_x\ge0$ supported in $|w|_\infty\le2$ of total mass $\frac13+\frac{13}3x+\frac{49}3x^2$. Two elementary Fourier inequalities for the transforms $\hat E_x$ of $x^{|w|_\infty}$ on $\mathbb Z^2$ and $\hat D_x$ of its restriction to the two diagonals, namely $\hat D_x\ge-1$ and $\hat E_x\ge\frac{1-x}{2(1+x)}(\hat D_x+1)$, give $\hat\Phi_x=\frac{1+x}{2(1-x)}\hat E_x-\frac16\hat D_x\ge\frac1{12}\hat D_x+\frac14\ge\frac16$; the transform of $e^{-\epsilon(|w_2|+|w_3|)}\Phi_x$ is then a convolution of non-negative functions, so $\Sigma\ge-x^{-2}\sum_w\mathsf R_x(w)\ge-53x$.

In both cases $\Sigma\ge-145x$, so the left member is at least $-870e^{-3\epsilon}\int t^\eta d\mathsf m\ge-870\,s^\eta_\nu(3e_1)$; and $6\cdot145=870$.
(b) The window $\tilde w_1(z)=\eta^{-1}\sum_Xe^{-\epsilon|X|}\sum_x\Lambda_\eta(x)\Lambda_\eta(x+z-\eta X)$ is even, has values in $[0,1]$ and the transform $\eta^{-1}\sum_ue^{-\epsilon|u|}e^{iu\eta\theta}\,|\sum_x\Lambda_\eta(x)e^{-ix\theta}|^2\ge0$, and $\delta_N=\sum_{j\ne0}\tilde w_1(jN)\to0$. By (d) with $\omega_i=\tilde w_1$, $P_N''(\beta)\ge(1+\delta_N)^{-3}\eta^{-3}\sum_Xe^{-\epsilon|X|_1}s^\eta_{\mu_{N,\beta}}(X)$. Integrate against $0\le G\in C^\infty_c(I)$, let $N\to\infty$ (dominated convergence for $\int G''P_N$, Fatou's lemma for the right member, and weak limit points of $\mu_{N,\beta}$ lie in $\mathcal G(\beta)$), use (a) to bound the terms with $|X|_\infty\ge3$, and let $\epsilon\to0$.
(c) $\sum_v\varphi^{\eta'}_{\rm c}(v)\varphi^\eta_0(\cdot-v)=\sum_{v,x}\varphi^{\eta'}_{\rm c}(v)\varphi^\eta_{\rm c}(x)\varphi^1_0(\cdot-v-x)$ is symmetric in $(\eta,\eta')$; expand $\eta'^3\langle\varphi^\eta_0;\Theta\rangle$ with it, exchange $\eta$ and $\eta'$, divide by $\eta^3\eta'^3$, and write $h=\Theta-(\Theta-h)$. The vectors $v$ with a weight $\ne0$ have $|v|_\infty<\bar\eta$, so $\Theta-\Theta(\cdot+v)$ vanishes where $|y-(\frac12,\frac12,\frac12)|_\infty\le3\bar\eta$, and $\Theta-h^\eta$ vanishes where this distance is at most $2\eta$.
\end{proof}

\begin{proposition}[For data: two incommensurable periods]\label{lh:per}
Let $\beta_0>0$, let $P$ be convex on $(\beta_0,\infty)$, $\hat n$ Borel and bounded on the compact subsets of $(\beta_0,\infty)$, and $C_{\rm low}\ge0$ with
\begin{equation}\label{lh:low}
P''(d\beta)\ge\bigl[\hat n(\beta)-C_{\rm low}\beta^{-5}\bigr]d\beta\qquad\text{on }(\beta_0,\infty);
\end{equation}
let $\mathbb M$ be the image under $s=\log\beta$ of $\beta^4[P''(d\beta)-\hat n\,d\beta]$, and $J_P(\beta)=P'(\beta+)-P'(\beta-)$. For $i=1,2$ let $\lambda_i>0$ with $\lambda_1/\lambda_2\notin\mathbb Q$, $\alpha_i\in\mathbb R$, $\varkappa_i>0$, $\mathsf J_i=[\alpha_i,\alpha_i+\lambda_i+2\varkappa_i]$, integers $k_i$ with $\alpha_i+k_i\lambda_i>\log\beta_0$, and integers $m^0_i\ge0$; for integers $m\ge m^0_i$ and $K\ge k_i$ let $P^{(i)}_{m,K}$ be a $C^2$ function and $n^{(i)}_{m,K}$ a bounded Borel function on $W^{(i)}_K=\exp(\mathsf J_i+K\lambda_i)$, with numbers $\epsilon^{(i)}_m,C_N\ge0$ such that on $W^{(i)}_K$
\begin{equation}\label{lh:data}
\beta^3|P-P^{(i)}_{m,K}|\le\epsilon^{(i)}_m,\qquad\beta^5|n^{(i)}_{m,K}-\hat n|\le C_N ;
\end{equation}
put $y^{(i)}_{m,K}(\upsilon)=\beta^5[(P^{(i)}_{m,K})''-n^{(i)}_{m,K}](\beta)$ at $\beta=e^{\upsilon+K\lambda_i}$, $\upsilon\in\mathsf J_i$, and $\varkappa=\min\{\varkappa_1,\varkappa_2,1\}$.
\begin{enumerate}
\item[(a)] Let $\sup_{K\ge k^0}\sup_{\mathsf J_1}|y^{(1)}_{m^0,K}|\le Y^0$ for some $m^0\ge m^0_1$ and some integer $k^0\ge k_1$; let $m_i\ge m^0_i$, integers $\bar K_i\ge k_i$ with $\bar K_1\ge k^0$, and $\omega_i\ge0$ satisfy $\sup_{\mathsf J_i}|y^{(i)}_{m_i,K}-y^{(i)}_{m_i,K'}|\le\omega_i$ for $K,K'\ge\bar K_i$. Let $\epsilon=\max_i\epsilon^{(i)}_{m_i}\le\varkappa^2/4$,
\[
\mathsf M_1=(Y^0+C_N+C_{\rm low})\varkappa_1+\epsilon^{(1)}_{m^0}\bigl(15/\varkappa_1+14+12\varkappa_1\bigr),\qquad\mathsf A_0=7\mathsf M_1/\varkappa_1+C_{\rm low}+4C_N+\omega_1+\omega_2 .
\]
Then, with $\mathbb M^+=\mathbb M+C_{\rm low}ds$, for $s\ge s_*=\max_i(\alpha_i+2\varkappa_i+\bar K_i\lambda_i)$: $\mathbb M^+([s,s+a])\le\mathsf A_0a+(\mathsf A_0+134)\epsilon^{1/2}$ for $0\le a\le\varkappa/2$; and for $0\le\delta\le1$
\[
e^{4s}J_P(e^s)\le(\mathsf A_0+134)\epsilon^{1/2},\qquad|\mathbb M|([s,s+\delta])\le(2\mathsf A_0+C_{\rm low}+134)\delta+(\mathsf A_0+134)\epsilon^{1/2}.
\]
Hence, if such thresholds $\bar K_i$, of whatever size, exist for all $m_i\ge m^0_i$ with one pair $\omega_1,\omega_2$, and $\inf_m\epsilon^{(i)}_m=0$, then $e^{4s}J_P(e^s)\to0$ and $\limsup_{s\to\infty}|\mathbb M|([s,s+\delta])\le(2\mathsf A_0+C_{\rm low}+134)\delta$.
\item[(b)] Let $\epsilon^{(i)}_m=c_{{\rm P},i}e^{-\lambda_im}$ with $c_{{\rm P},i}>0$, and let there be $\mathsf a_i\ge0$, $A_i>0$, $\omega_i\ge0$ with $\sup_{\mathsf J_i}|y^{(i)}_{m,K}-y^{(i)}_{m,K'}|\le\omega_i$ for all $m\ge m^0_i$ and all integers $K,K'\ge\bar K_i(m)=\max\{k_i,\lceil\mathsf a_i+A_im\rceil\}$. Then the hypothesis of (a) on $Y^0$ holds at $m^0=m^0_1$, $k^0=\bar K_1(m^0_1)$, with $Y^0=\sup_{\mathsf J_1}|y^{(1)}_{m^0_1,k^0}|+\omega_1$. Let $\mathsf A_0$ be that of (a) at these numbers, $A=\max_iA_i$,
\[
c_{3,i}=|\alpha_i|+2\varkappa_i+\lambda_i(\max\{k_i,0\}+\mathsf a_i+1),\quad c_6=\max_ic_{{\rm P},i}e^{\lambda_i+c_{3,i}/A_i},\quad T_i(m)=\alpha_i+2\varkappa_i+\lambda_i\bar K_i(m),
\]
$s_\sharp=\max\{0,T_1(m^0_1),T_2(m^0_2),A\log^+(4c_6/\varkappa^2)\}$, $C_{\rm h}=2\mathsf A_0+C_{\rm low}+134$ and $C_{\rm r}=(\mathsf A_0+134)c_6^{1/2}$. Then for $s\ge s_\sharp$, $\beta=e^s$ and $0\le\delta\le1$: $\beta^4J_P(\beta)\le C_{\rm r}\beta^{-1/2A}$ and $|\mathbb M|([s,s+\delta])\le C_{\rm h}\delta+C_{\rm r}\beta^{-1/2A}$.
\item[(c)] (The cut.) Let the hypothesis of (a) on $Y^0$ hold, and let there be one integer $K_*\ge\max\{k_1,k_2,k^0\}$ and numbers $\omega_i$ such that $\sup_{\mathsf J_i}|y^{(i)}_{m,K}-y^{(i)}_{m,K'}|\le\omega_i$ for $K,K'\ge K_*$ and every $m\ge m^0_i$, and $\inf_m\epsilon^{(i)}_m=0$, $i=1,2$. Then $J_P(e^s)=0$ for $s\ge s_*=\max_i(\alpha_i+2\varkappa_i+K_*\lambda_i)$, and on $[s_*,\infty)$ the measure $\mathbb M$ has a density with values in $[-C_{\rm low},\mathsf A_0-C_{\rm low}]$, $\mathsf A_0$ being that of (a).
\item[(d)] (The converse.) Let $\mathbb M$ be any signed Radon measure on $(s_0,\infty)$, $s_2>s_0$, $\epsilon_1:[s_2,\infty)\to(0,1]$ non-increasing, and $C_{\rm h},C_{\rm r}\ge0$ with $|\mathbb M|([s,s+\delta])\le C_{\rm h}\delta+C_{\rm r}\epsilon_1(s)$ for $s\ge s_2$ and $0\le\delta\le1$. Then for $s_3\ge s_2$ and $G\in C^1(\mathbb R)$ with compact support in $[s_3,\infty)$, $|\langle\mathbb M,G\rangle|\le(C_{\rm h}+C_{\rm r})(\|G\|_1+\epsilon_1(s_3)\|G'\|_1)$; so $|\langle\mathbb M,G(\cdot-t)\rangle-\langle\mathbb M,G\rangle|\le2(C_{\rm h}+C_{\rm r})(\|G\|_1+\epsilon_1(s_3)\|G'\|_1)$ for every $t\ge0$.
\item[(e)] (The exponent is attained.) For all $\lambda_1,\lambda_2>0$ with $\lambda_1/\lambda_2\notin\mathbb Q$, all $\alpha_i\in\mathbb R$, $\varkappa_i>0$, $\mathsf a_i\ge0$ and every $A>0$ there are data as in (b), with $\beta_0=1$, $m^0_i=0$, $A_1=A_2=A$, $C_N=0$, $C_{\rm low}=7$, $\omega_1=\omega_2=0$ and $c_{{\rm P},i}=\frac14e^{(1+|\alpha_i|)/A}$, such that $\beta_k^4J_P(\beta_k)=\beta_k^{-1/2A}$ at $\beta_k=e^{2k}$, $k\in\mathbb N$.
\end{enumerate}
\end{proposition}

\begin{proof}[Proof sketch]
$\mathbb M^+$ is a positive measure by \eqref{lh:low}; it has the same atoms as $\mathbb M$, $\mathbb M(\{s\})=e^{4s}J_P(e^s)$, and $\mathbb M\le\mathbb M^+$, $|\mathbb M|\le\mathbb M^++C_{\rm low}ds$. For $G\in C^2$ supported in $\mathsf J_i+K\lambda_i$ put $g(\beta)=\beta^4G(\log\beta)$, so that $g''=\beta^2(G''+7G'+12G)(\log\beta)$. Then $\langle\mathbb M,G\rangle=\int g''(P-P^{(i)}_{m,K})d\beta+\int g\,(n^{(i)}_{m,K}-\hat n)d\beta+\int g\beta^{-5}y^{(i)}_{m,K}(\log\beta-K\lambda_i)d\beta$, by two integrations by parts against $P^{(i)}_{m,K}$; with \eqref{lh:data} this is the identity
\begin{equation}\label{lh:ident}
\begin{multlined}
\langle\mathbb M,G\rangle=\int G(x)\,y^{(i)}_{m,K}(x-K\lambda_i)\,dx+r,\qquad|r|\le\epsilon^{(i)}_m\mathsf Q(G)+C_N\|G\|_1,\quad\mathsf Q(G)=\|G''\|_1 \\
+7\|G'\|_1+12\|G\|_1 .
\end{multlined}
\end{equation}
Let $G$ be supported in $[-\varkappa,\varkappa]$ and $A_G(s)=\langle\mathbb M,G(\cdot-s)\rangle$. For $s\ge\alpha_i+\varkappa_i+\bar K_i\lambda_i$ write $s=s_0+K\lambda_i$ with $s_0\in[\alpha_i+\varkappa_i,\alpha_i+\varkappa_i+\lambda_i)$, so $K\ge\bar K_i$ and $G(\cdot-s)$ is supported in $\mathsf J_i+K\lambda_i$; a shift of $s$ by a multiple of $\lambda_i$ changes $K$ alone, and \eqref{lh:ident} at the two cutoffs, with the oscillation bound, gives $|A_G(s)-A_G(s')|\le2\epsilon^{(i)}_{m_i}\mathsf Q(G)+(2C_N+\omega_i)\|G\|_1$ whenever $s'-s\in\lambda_i\mathbb Z$, whatever the number of periods. As $\lambda_1/\lambda_2\notin\mathbb Q$, the set $\{k\lambda_1-k'\lambda_2:k,k'\ge0\}\cap[0,\infty)$ is dense in $[0,\infty)$ (by the pigeon-hole principle), and $A_G$ is continuous; moving first forward by $k\lambda_1$ and then back by $k'\lambda_2$ we get, for $s\ge s_*$ and every $t\ge0$,
\[
|A_G(s+t)-A_G(s)|\le2\bigl(\epsilon^{(1)}_{m_1}+\epsilon^{(2)}_{m_2}\bigr)\mathsf Q(G)+(4C_N+\omega_1+\omega_2)\|G\|_1 ;
\]
only the irrationality of $\lambda_1/\lambda_2$ is used. At the exponent $m^0$, \eqref{lh:ident} with $Y^0$ gives $\mathbb M^+([s,s+\frac12\varkappa_1])\le\mathsf M_1$ for $s\ge\alpha_1+\varkappa_1+k^0\lambda_1$. Now let $G=\Phi_{a,\flat}$ be a plateau function: the indicator of $[-\frac\flat2,a+\frac\flat2]$ smoothed with the kernel $\frac{15}{16}(1-y^2)_+^2$ scaled to the width $\frac\flat2$, which is $C^2$, has values in $[0,1]$, equals $1$ on $[0,a]$, is supported in $[-\flat,a+\flat]$ and has $\mathsf Q(G)\le\frac{15}{2\flat}+14+12(a+\flat)$. Averaging the last display over $t\in[0,\varkappa_1]$ and using $\mathbb M\le\mathbb M^+$ bounds $A_G(s)$ by $7\mathsf M_1\|G\|_1/\varkappa_1$ plus the error; hence, for $0\le a\le\varkappa/2$ and $\flat=\epsilon^{1/2}$,
\[
\mathbb M^+([s,s+a])\le\mathsf A_0(a+\flat)+4\epsilon\bigl(\tfrac{15}{2\flat}+26\bigr)\le\mathsf A_0a+(\mathsf A_0+134)\epsilon^{1/2}.
\]
At $a=0$ this bounds the atom; covering $[s,s+\delta]$ by at most $2\delta/\varkappa+1$ intervals of length $\le\varkappa/2$ gives the bound for $|\mathbb M|$.
(b) $y^{(1)}_{m^0_1,k^0}$ is bounded on $\mathsf J_1$, and the oscillation bound at $m^0_1$ bounds $y^{(1)}_{m^0_1,K}$, $K\ge k^0$, by it plus $\omega_1$. At a given $s\ge s_\sharp$ let $m_i$ be the largest $m\ge m^0_i$ with $T_i(m)\le s$; as $\bar K_i(m)\le\max\{k_i,0\}+\mathsf a_i+A_im+1$, $T_i(m_i+1)>s$ gives $\lambda_im_i>(s-c_{3,i})/A_i-\lambda_i$, so $\epsilon\le c_6e^{-s/A}\le\varkappa^2/4$, and (a) applies with $\bar K_i=\bar K_i(m_i)$ and $s_*\le s$; $\mathsf A_0$ does not depend on $m_1,m_2$.
(c) $\mathsf A_0$ and $s_*$ do not depend on $m_1,m_2$, and the lower bound of $\epsilon$ over the admitted pairs is $0$; so $\mathbb M^+([s,s+a])\le\mathsf A_0a$ for $s\ge s_*$: there are no atoms, and by outer regularity $\mathbb M^+(B)\le\mathsf A_0|B|$ for Borel sets $B\subset[s_*,\infty)$.
(d) Cut $[s_3,\infty)$ into intervals $I_j$ of length $\delta=\epsilon_1(s_3)$; then $|\mathbb M|(I_j)\le(C_{\rm h}+C_{\rm r})\delta$, and $\delta\sup_{I_j}|G|\le\int_{I_j}|G|+\delta\int_{I_j}|G'|$; sum over $j$. The translate $G(\cdot-t)$ has the same norms and its support lies in $[s_3,\infty)$.
(e) Let $s_k=2k$, $\theta_k=e^{-s_k/2A}$, $\Phi(x)=-\frac12(1-|x|)_+^2$, so that $\Phi''=2\delta_0-\mathbf 1_{(-1,1)}$, and $\varphi=\sum_k\frac12\theta_k^2\Phi((\cdot-s_k)/\theta_k)$. Put $P(\beta)=\beta^{-3}[\varphi(\log\beta)+\frac7{12}]$ and $\hat n(\beta)=7\beta^{-5}$. For continuous $\phi$ the image of $\beta^4\,d^2(\beta^{-3}\phi(\log\beta))$ is the distribution $\phi''-7\phi'+12\phi$; so the image of $\beta^4P''(d\beta)$ is $\sum_k\theta_k\delta_{s_k}+(\rho+7)dx$ with a density $\rho\ge-\frac12-\frac72-3=-7$: $P$ is convex, $\mathbb M=\sum_k\theta_k\delta_{s_k}+\rho\,dx$, \eqref{lh:low} holds with $C_{\rm low}=7$, and $\beta_k^4J_P(\beta_k)=\theta_k=\beta_k^{-1/2A}$. As approximants take $n^{(i)}_{m,K}=\hat n$ and $P^{(i)}_{m,K}=\beta^{-3}[\varphi^{[i,m]}(\log\beta)+\frac7{12}]$, where $\varphi^{[i,m]}$ keeps, mollified to within $\epsilon^{(i)}_m$, only the bumps with $s_k+1\le\alpha_i+\lambda_i\bar K_i(m)$: the bumps left out have $|\varphi_k|\le\frac14\theta_k^2\le\epsilon^{(i)}_m$, and $y^{(i)}_{m,K}=0$ on $\mathsf J_i$ for $K\ge\bar K_i(m)$.
\end{proof}

\begin{remark}[One period does not suffice, and the lower bound cannot be dropped]\label{lh:onep}
(1) Let $\lambda>0$, $0<\theta<\min\{1,\lambda/2\}$, $\Phi$ as in the proof of (e), $\varphi(x)=\sum_{k\in\mathbb Z}\frac12\theta^2\Phi((x-k\lambda)/\theta)$, $P(\beta)=\beta^{-3}[\varphi(\log\beta)+\frac7{12}]$ and $\hat n=7\beta^{-5}$. With $\lambda$-periodic smooth $\varphi^{[m]}$, $\|\varphi-\varphi^{[m]}\|_\infty\le e^{-\lambda m}$, the approximants $\beta^{-3}[\varphi^{[m]}(\log\beta)+\frac7{12}]$ and $n_{m,K}=\hat n$ satisfy \eqref{lh:low} with $C_{\rm low}=7$, \eqref{lh:data} with $C_N=0$ and $\epsilon_m=e^{-\lambda m}$, the bound on $Y^0$ at every $m^0$, and the oscillation bound with $\omega=0$ at every cutoff, a threshold free of $m$, since $y_{m,K}$ does not depend on $K$; yet $\beta^4J_P(\beta)=\theta$ at every $\beta=e^{k\lambda}$. So the hypotheses at one period, or at two commensurable periods $p\lambda$, $q\lambda$ (with $\varphi^{[pm]}$, $\varphi^{[qm]}$ in the place of $\varphi^{[m]}$, so that $\epsilon^{(i)}_m=e^{-p\lambda m}$, $e^{-q\lambda m}$), give neither (a) nor (b) nor (c). (2) The lower bound \eqref{lh:low} cannot be dropped from (a), nor from (b), at any two periods. Let $\beta_j=4^j$ for the integers $j\ge j_0\ge1$, numbers $w_j$ with $2\beta_j^{-2}\le w_j\le\frac14\beta_j$ and $w_j/\beta_j\to0$, $c_j=\frac16w_j\beta_j^{-4}$, $\tau(x)=-(1-|x|)_+^3$, $\beta_0=1$,
\[
P(\beta)=-\log\beta+\sum_{j\ge j_0}c_j\,\tau\Bigl(\frac{\beta-\beta_j}{w_j}\Bigr)\qquad\text{and}\qquad\hat n(\beta)=\beta^{-2}.
\]
The $j$-th term vanishes off $[\beta_j-w_j,\beta_j+w_j]\subset[\frac34\beta_j,\frac54\beta_j]$, and these intervals are disjoint. The function $\tau$ is $C^2$ off $0$, $\tau'(0\pm)=\pm3$, and $\tau''(x)=-6(1-|x|)$ for $0<|x|<1$; so on the $j$-th interval, off $\beta_j$, $P''\ge\beta^{-2}-6c_jw_j^{-2}=\beta^{-2}-\beta_j^{-4}w_j^{-1}\ge(\frac{16}{25}-\frac12)\beta_j^{-2}>0$, while at $\beta_j$ the slope of $P$ jumps up by $6c_j/w_j=\beta_j^{-4}$. So $P$ is convex on $(\beta_0,\infty)$, and $\beta_j^4J_P(\beta_j)=1$ for every $j$. Moreover $\beta^3c_j|\tau((\beta-\beta_j)/w_j)|\le(\frac54\beta_j)^3c_j<\frac12w_j/\beta_j$. For given $\lambda_i$, $\alpha_i$, $\varkappa_i$, $k_i$ as in Proposition~\ref{lh:per}, put $m^0_i=0$ and $n^{(i)}_{m,K}=\hat n$, so that $C_N=0$, and, for all $K$,
\[
P^{(i)}_{m,K}(\beta)=-\log\beta+\sum_{j_0\le j<j(m)}c_j\,\tau_m\Bigl(\frac{\beta-\beta_j}{w_j}\Bigr),
\]
with integers $j(m)\to\infty$ and $C^2$ functions $\tau_m$ with support in $[-1,1]$ and $\sup|\tau_m-\tau|\le\epsilon_m:=\sup_{j\ge j(m)}w_j/\beta_j$. Then \eqref{lh:data} holds with $\epsilon^{(i)}_m=\epsilon_m$, and $\epsilon_m\to0$. On $W^{(i)}_K$ the function $y^{(i)}_{m,K}$ vanishes as soon as $e^{\alpha_i+K\lambda_i}>\frac54\beta_{j(m)-1}$, because there $(P^{(i)}_{m,K})''=\beta^{-2}=\hat n$; so the oscillation bound holds with $\omega_i=0$ beyond a finite threshold at each $m$, and the bound on $Y^0$ holds at every $m^0$, as the cutoffs below the threshold are finitely many. Thus every condition of (a) and of its last sentence holds except \eqref{lh:low}: near $\log\beta_j$ the measure $\mathbb M$ has the density $\beta^5c_jw_j^{-2}\tau''((\beta-\beta_j)/w_j)$, which tends to $-\beta_j/w_j$ at $\log\beta_j$, so \eqref{lh:low} fails for every $C_{\rm low}$; and the conclusion of (a) fails, as $e^{4s}J_P(e^s)=1$ at every $s=\log\beta_j$.

Let now $w_j=\min\{\frac14\beta_j,\beta_j^{1-1/A}\}$ with $A>\frac13$ and $j_0$ so large that $w_j\ge2\beta_j^{-2}$ for $j\ge j_0$, let $\mathsf a_i\ge0$ and $\bar K_i(m)=\max\{k_i,\lceil\mathsf a_i+Am\rceil\}$, and at the period $\lambda_i$ keep in $P^{(i)}_{m,K}$ the terms with $2\beta_j\le e^{\alpha_i+\lambda_i\bar K_i(m)}$, with $\tau_m$ within $c_{{\rm P},i}e^{-\lambda_im}$ of $\tau$, where $c_{{\rm P},i}=2^{1/A}e^{|\alpha_i|/A}$. A term left out has $\beta_j>\frac12e^{\alpha_i+\lambda_i\bar K_i(m)}\ge\frac12e^{-|\alpha_i|}e^{\lambda_iAm}$, so that $w_j/\beta_j\le\beta_j^{-1/A}<c_{{\rm P},i}e^{-\lambda_im}$; and $y^{(i)}_{m,K}=0$ on $\mathsf J_i$ for $K\ge\bar K_i(m)$. So every condition of (b) holds, with $A_1=A_2=A$, $\omega_1=\omega_2=0$, $C_N=0$ and $\epsilon^{(i)}_m=c_{{\rm P},i}e^{-\lambda_im}$, except \eqref{lh:low}, and the conclusion of (b) fails.

In the form in which the proofs of (a) and (b) use the data, namely a measure with bounded variation on unit intervals whose smoothed translates are nearly invariant, a lower bound is needed as well: $\mathbb M=\sum_j(\delta_{s_j}-\delta_{s_j+\eta_j})$, with $s_{j+1}-s_j\ge1$, $0<\eta_j<1$ and $\eta_j\to0$, has $|\mathbb M|(I)\le4$ on every interval $I$ of length $1$ and $\langle\mathbb M,G(\cdot-s)\rangle\to0$ as $s\to\infty$ for every $G\in C^1_c(\mathbb R)$, but atoms $1$ at the points $s_j$. Its atoms $-1$ at the points $s_j+\eta_j$ set it apart from the measures $\mathbb M$ of data, whose atoms $e^{4s}J_P(e^s)$ are non-negative; in the instance above the density $-\beta_j/w_j$ takes their place. (3) By (d), a bound $|\mathbb M|([s,s+\delta])\le C\delta+Ce^{-cs}$ gives the uniform translation bound of Proposition~\ref{lh:per}(d), with its persistent $L^1$ term and an exponentially small derivative correction; the proofs of (a) and (b) turn near invariance at such a rate under two incommensurable periods, with \eqref{lh:low}, into the same bound at the rate $e^{-cs/2}$; and by (e) this loss of one half is real for data.
\end{remark}

\begin{proposition}[For data: only one slope in the exponent]\label{lh:slope}
Let the data of Proposition~\ref{lh:per} be given, with $\epsilon^{(i)}_m=c_{{\rm P},i}e^{-\lambda_im}$, $c_{{\rm P},i}>0$. Suppose:
\begin{enumerate}
\item[(i)] there are $\mathsf a_1\ge0$, $A_1>0$, $\omega_1\ge0$ with $\sup_{\mathsf J_1}|y^{(1)}_{m,K}-y^{(1)}_{m,K'}|\le\omega_1$ for all $m\ge m^0_1$ and all integers $K,K'\ge\max\{k_1,\mathsf a_1+A_1m\}$, and $\sup_{K\ge k_1}\sup_{\mathsf J_1}|y^{(1)}_{m^0_1,K}|\le Y^0$;
\item[(ii)] there is $\omega_2\ge0$ such that for every $m$ in an unbounded set of integers $\ge m^0_2$ there is an integer $\bar K_2(m)\ge k_2$ with $\sup_{\mathsf J_2}|y^{(2)}_{m,K}-y^{(2)}_{m,K'}|\le\omega_2$ for $K,K'\ge\bar K_2(m)$.
\end{enumerate}
Put $\mathsf M_0=\varkappa_1(Y^0+C_N+C_{\rm low})+c_{{\rm P},1}(4/\varkappa_1+14+12\varkappa_1)$, $\mathsf H_*=\omega_1+\omega_2+4C_N+120+16\mathsf M_0/\varkappa_1+12\mathsf M_0/\lambda_2$ and $C_\vartheta=\max\{1,c_{{\rm P},1}^{1/2}\}\exp[(\alpha_1+\varkappa_1+\lambda_1(\mathsf a_1+1))/2A_1+\lambda_1/2]$. There is $s_\sharp$ such that for $s\ge s_\sharp$, $\beta=e^s$ and $\delta\ge0$
\[
|\mathbb M|([s,s+\delta])\le(2\mathsf H_*+C_{\rm low})\delta+2\mathsf H_*C_\vartheta\beta^{-1/2A_1},\qquad\beta^4J_P(\beta)\le\mathsf H_*C_\vartheta\beta^{-1/2A_1}.
\]
Only the slope $A_1$ of the first period enters the exponent; of the second period only one bound $\omega_2$, some finite threshold at each $m$, and $\epsilon^{(2)}_m\to0$ are used.
\end{proposition}

\begin{proof}[Proof sketch]
\emph{The identity for tents.} For $\varphi_{c,\vartheta}(x)=(1-|x-c|/\vartheta)_+$ with $[c-\vartheta,c+\vartheta]\subset\mathsf J_i+K\lambda_i$ the identity \eqref{lh:ident} holds with $\mathsf Q(G)$ replaced by $4/\vartheta+14+12\vartheta$: with $Q(s)=e^{3s}P(e^s)$ and $Q_{m,K}$ formed in the same way from $P^{(i)}_{m,K}$, the image of $\beta^4P''(d\beta)$ is the distribution $Q''-7Q'+12Q$, $D=Q-Q_{m,K}$ has $|D|\le\epsilon^{(i)}_m$, and $\int\varphi_{c,\vartheta}\,d(D''-7D'+12D)=[D(c-\vartheta)-2D(c)+D(c+\vartheta)]/\vartheta+7\int\varphi'_{c,\vartheta}D+12\int\varphi_{c,\vartheta}D$.
\emph{Step 1 (mass).} With $c_1=\alpha_1+k_1\lambda_1+\varkappa_1$, every $c\ge c_1$ is $\upsilon_0+K\lambda_1$ with $K\ge k_1$ and $\upsilon_0\in[\alpha_1+\varkappa_1,\alpha_1+\varkappa_1+\lambda_1]$; the identity at $m^0_1$, $\vartheta=\varkappa_1$, with (i), gives $\int\varphi_{c,\varkappa_1}d\mathbb M^+\le\mathsf M_0$, so $\mathbb M^+([c-\frac{\varkappa_1}2,c+\frac{\varkappa_1}2])\le2\mathsf M_0$, and an interval of length $l$ in $[c_1,\infty)$ has $\mathbb M^+$-mass at most $2\mathsf M_0(l/\varkappa_1+1)$.
\emph{Step 2 (one period).} For an admitted $m$, $0<\vartheta\le\min\{\varkappa_i,1\}$ with $\epsilon^{(i)}_m\le\vartheta^2$, $\upsilon_0\in[\alpha_i+\varkappa_i,\alpha_i+\varkappa_i+\lambda_i]$ and cutoffs $K,K'$ beyond the threshold of $m$, the identity at the two cutoffs gives $|\int\varphi_{\upsilon_0+K\lambda_i,\vartheta}d\mathbb M^+-\int\varphi_{\upsilon_0+K'\lambda_i,\vartheta}d\mathbb M^+|\le t_i\vartheta$ with $t_i=\omega_i+2C_N+60$.
\emph{Step 3 (two periods).} Let $m\ge m^0_1$, $\vartheta\le\min\{\varkappa_1,\varkappa_2,1,\lambda_2/6\}$ with $\vartheta^2\ge c_{{\rm P},1}e^{-\lambda_1m}$, $c=\upsilon_0+K\lambda_1$ with $K\ge\max\{k_1,\mathsf a_1+A_1m\}$, and $\mathsf H=\vartheta^{-1}\int\varphi_{c,\vartheta}d\mathbb M^+$. Take $m_2$ in the set of (ii) with $\epsilon^{(2)}_{m_2}\le\vartheta^2$, an integer $K_2\ge\bar K_2(m_2)$ with $S=\alpha_2+\varkappa_2+K_2\lambda_2\ge c_1+1$, $n=\lfloor\lambda_2/3\vartheta\rfloor$, and the intervals $I_r=(S+3\vartheta(r-1),S+3\vartheta(r-1)+\vartheta)\subset[S,S+\lambda_2)$, $r\le n$. As $\lambda_1/\lambda_2\notin\mathbb Q$, the points $c+k\lambda_1$, $k\ge k_\flat$, are dense modulo $\lambda_2$ for every $k_\flat$; so there are integers $k_r\ge0$ and $j_r\ge1$ with $p_r=c+k_r\lambda_1-j_r\lambda_2\in I_r$. Step~2 at $i=1$, with the cutoffs $K$ and $K+k_r$, both beyond the threshold however large $k_r$ is, and then at $i=2$, with the cutoffs $K_2+j_r$ and $K_2$, gives $\int\varphi_{p_r,\vartheta}d\mathbb M^+\ge\vartheta(\mathsf H-t_1-t_2)$. The $p_r$ are at least $2\vartheta$ apart, so $\sum_r\varphi_{p_r,\vartheta}\le1$, and by Step~1, $n\vartheta(\mathsf H-t_1-t_2)\le\mathbb M^+([S-\vartheta,S+\lambda_2+\vartheta])\le2\mathsf M_0(\frac{4\lambda_2}{3\varkappa_1}+1)$; with $n\vartheta\ge\lambda_2/6$ this gives $\mathsf H\le\mathsf H_*$.
\emph{Step 4.} For large $s$ let $m(s)=\lfloor((s-\alpha_1-\varkappa_1)/\lambda_1-\mathsf a_1-1)/A_1\rfloor$ and $\vartheta(s)=\max\{1,c_{{\rm P},1}^{1/2}\}e^{-\lambda_1m(s)/2}\le C_\vartheta e^{-s/2A_1}$; choose $s_\sharp$ such that for $s\ge s_\sharp$, $m(s)\ge m^0_1$, $(s-\alpha_1-\varkappa_1)/\lambda_1\ge k_1+1$ and $\vartheta(s)\le\min\{\varkappa_1,\varkappa_2,1,\lambda_2/6\}$. Then every $c\ge s$ is admitted in Step~3 with $m(s)$ and $\vartheta(s)$. So the atom at $s$ is at most $\int\varphi_{s,\vartheta}d\mathbb M^+\le\mathsf H_*\vartheta$; $[s,s+\delta]$ is covered by at most $\delta/\vartheta+1$ intervals $[c-\frac\vartheta2,c+\frac\vartheta2]$, $c\ge s$, each of $\mathbb M^+$-mass at most $2\mathsf H_*\vartheta$; and $|\mathbb M|\le\mathbb M^++C_{\rm low}ds$.
\end{proof}

\begin{proposition}[One block size]\label{lh:blk}
In a setting:
\begin{enumerate}
\item[(a)] $(K,L^Kb)\in\mathfrak d(\gamma)\cup\mathfrak d(g_2)$ for $b\in\mathsf I$ and $K\in\mathbb N$; $b_+-Lb_-=(2L-2)\bar B+\frac L2>0$, so $\bigcup_KL^K\mathsf I=[Lb_-,\infty)$.
\item[(b)] Let $K,m\in\mathbb N$, $\ell=2L^m$, $N=2L^{m+K}$, $\beta>0$ and $b=a_K\beta$. The partition function of $T_\ell$ at the spacing $a_K$ is $Z_N(\beta)$, so its free energy per physical volume is $a_K^{-3}P_N(\beta)$; the dilation by $L^K$ maps $\mathcal G_{K,\beta}$ onto $\mathcal G(\beta)$ and the Wilson measure of $T_\ell$ to $\mu_{N,\beta}$; for a law $\nu$ on $\mathbb L_K$ invariant under $a_K\mathbb Z^3$, with image $\nu^{(K)}$, and an integer $\eta\ge1$, $\beta^5n_\eta(\nu^{(K)})=b^5(a_K\eta)^{-3}S^T_2[\nu](\mathsf f_\eta,\mathsf h_\eta)$, where $\mathsf f_\eta$ is the tent of half width $a_K\eta$ about a point of $a_K(\mathbb Z+\frac12)^3$ and $\mathsf h_\eta$ its hull; and, with $\eta_K=L^{K-1}$,
\begin{equation}\label{lh:far}
\beta^5P_N''(\beta)=\beta^5n_{\eta_K}(\mu_{N,\beta})+y_{m,K}(\log b).
\end{equation}
\item[(c)] For $b\in\mathsf I$, $K\in\mathbb N$, $\beta=L^Kb$, $m\ge m_0$ and $N=2L^{m+K}$,
\[
0\le a_K^{-3}[P_N(\beta)-P(\beta)]\le3\ell_0C_{\rm v}/\ell,\qquad\text{so}\qquad0\le\beta^3[P_N(\beta)-P(\beta)]\le c_{\rm P}L^{-m},
\]
with $c_{\rm P}=\frac32b_+^3\ell_0C_{\rm v}$ and $C_{\rm v}=2C_{\rm S}+\max\{C_\gamma,C_{g_2}\}$, where $C_{\rm S}$ is the constant in Theorem~\textup{\YMMref{AD}{607}}(ii), (iii) and $C_g=C_{\rm S}(1+\log g^{-2})$ is its lower-bound constant at fixed $g$; $C_{\rm S}$ does not depend on $\varrho$, $\gamma$, $g$, $K$, $\beta$, $\ell$, and $C_g$ not on $K$, $\beta$, $\ell$.
\item[(d)] Let $r_1>0$ with $13r_1b_+\le1$. Then $\hat n$ is Borel, $|\hat n(\beta)|\le1125\lceil r_1\beta\rceil^3$, and $\hat n(\beta)=n_{\lceil r_1\beta\rceil}(\nu_\beta)$ for some $\nu_\beta\in\mathcal G(\beta)$. With the constant $C_{\rm nr}$ of the comparison of two laws in one ball of radius $1$ (the common-normalizer comparison in the proof of Proposition~\YMMref{91.13}{709}) and the constants $C_{\rm sp}(\mathsf d)$ of the bound at separated supports in every periodic Gibbs state (Proposition~\textup{\YMMref{91.13}{709}}), both taken at $g\in\{\gamma,g_2\}$ as in the proof, $\mathsf d_1=\min\{r_1b_-,\varepsilon\}$ and
\[
C^{\rm ref}=b_+^5\mathsf d_1^{-3}(40\sqrt3M/\mathsf d_1)^2\bigl[C_{\rm nr}+8C_{\rm sp}(\mathsf d_1)+C_{\rm sp}(r_1b_-)\bigr],
\]
for $b\in\mathsf I$, $K\in\mathbb N$, $\beta=L^Kb$, $m\ge m_0$ and $N=2L^{m+K}$:
\[
|\beta^5P_N''(\beta)-y_{m,K}(\log b)-\beta^5\hat n(\beta)|\le C^{\rm ref};
\]
$\beta^5|n_{\lceil r_1\beta\rceil}(\nu)-\hat n(\beta)|\le C^{\rm ref}$ and $0\le\beta^5\lceil r_1\beta\rceil^{-3}s^{\lceil r_1\beta\rceil}_\nu(3e_1)\le C^{\rm ref}$ for every $\nu\in\mathcal G(\beta)$; and $P''(d\beta)\ge[\hat n(\beta)-871C^{\rm ref}\beta^{-5}]d\beta$ on $(Lb_-,\infty)$.
\item[(e)] For all $K\in\mathbb N$, $m\ge m_0$ and $\upsilon\in\log\mathsf I$, $|y_{m,K}(\upsilon)|\le b_+^5L^3C_3(\varepsilon)(40\sqrt3ML)^2\ell^3$, where $C_3(\mathsf d)=(40\sqrt3M+1)C_{\rm tt}(\mathsf d)$ and $C_{\rm tt}(\mathsf d)$ is the larger, over $g\in\{\gamma,g_2\}$, of the numbers $A(\mathsf d)B(1)^2\,2^{\varkappa_\star}$ of Proposition~\YMMref{89.5}{694} at $g$: the constant of the bound for connected two-point functions on the tori at first-passage data, for two functions each supported in a closed sup-ball of radius $1$, at distance at least $\mathsf d$; it does not depend on $K$, $\beta$ or the side. This is better than the bound by $C^{\rm pol}\ell^{438}$ that (E)(b) gives.
\end{enumerate}
\end{proposition}

\begin{proof}[Proof sketch]
(a) With $y_1=\gamma^{-2}+\bar B\ge4+\bar B$ and $L\ge13$, $(L-1)y_1\ge4\bar B+1$, which says $y_2=g_2^{-2}+\bar B=y_1+2\bar B+1\le Ly_1-2\bar B$; so $I(\gamma)\cup I(g_2)=(y_1,Ly_2-2\bar B]\supset\mathsf I$, as $b_+=Ly_2-2\bar B$. And $b_+-Lb_-=L(y_2-y_1)-2\bar B-\frac L2=(2L-2)\bar B+\frac L2$, so consecutive intervals $L^K\mathsf I$ overlap.
(b) The dilation by $L^K$ carries bonds, plaquettes and the Haar measures of $\mathbb L_K$ to those of $\mathbb Z^3$ and $\ell\mathbb Z^3$ to $N\mathbb Z^3$, and $\mathcal O_K(f)$ becomes $L^K\mathcal O_1(f(L^{-K}\cdot))$; so $S^T_2[\nu](f,h)=L^{2K}\langle f(L^{-K}\cdot);h(L^{-K}\cdot)\rangle_{\nu^{(K)}}$, and $\beta^5\eta^{-3}L^{-2K}=b^5(a_K\eta)^{-3}$. As $L^K$ is odd, $\frac12\in a_K(\mathbb Z+\frac12)$, and $\varepsilon=a_K\eta_K$. At $\eta=\eta_K$, which divides $N$ with $N/\eta_K=2L^{m+1}\ge7$, \eqref{lh:split} becomes \eqref{lh:far}: the rescaled tents are the translates $\mathsf f(\cdot-\varepsilon X)$, their sum over the period cell is $1$ on $T_\ell$, so the far sum is $1-\mathsf h^{\langle\ell\rangle}$.
(c) By (a), $(K,\beta)\in\mathfrak d(g)$ for some $g\in\{\gamma,g_2\}$. Let $p(j)$, $j\in\mathbb N^3$, be the free energy per physical volume of the periodic box with $j_\mu\ell_0/a_K$ sites in the direction $\mu$, at the spacing $a_K$ and the coupling $\beta$. With $Z_{K,\ell,\beta}=Z_{\ell L^K}(\beta)$, the partition function of $T_\ell$ at the spacing $a_K$, (b) gives $p(k,k,k)=(k\ell_0)^{-3}\log Z_{K,k\ell_0,\beta}=a_K^{-3}P_{k\ell_0L^K}(\beta)$ for every power $k$ of $L$. In the first direction such a box consists of $2j_1$ slabs, and the trace formula (Proposition~\YMMref{85.2}{679}) writes its partition function as $\operatorname{Tr}B^{j_1}$ with one positive trace-class operator $B$ for all $j_1$, formed at $j_2,j_3$.

With the top eigenvalue $\|B\|$, $u_k=\frac1k\log\operatorname{Tr}B^k$ satisfies $0\le k(u_k-\log\|B\|)\le u_1-\log\|B\|$ and does not increase in $k$; hence $p$ does not increase in $j_1$ and $0\le p(k',j_2,j_3)-p(k,j_2,j_3)\le\frac1{k'}[p(1,j_2,j_3)-\lim_np(n,j_2,j_3)]$ for $k'\le k$. The same argument in the second and in the third direction shows that $p$ does not increase in any of its three arguments, with the same bound in each direction. The bracket is at most $C_{\rm v}$, for all $j_2,j_3$. By the monotonicity in each argument, $p(1,j_2,j_3)\le p(1,1,1)$, and $p(n,j_2,j_3)\ge p(L^i,L^i,L^i)$ whenever $L^i\ge\max\{n,j_2,j_3\}$; so the bracket is at most $\sup_{i\ge0}[p(1,1,1)-p(L^i,L^i,L^i)]$, where $p(1,1,1)$ and $p(L^i,L^i,L^i)$ belong to the tori with the sides $\ell_0=2L^{m_0}$ and $L^i\ell_0=2L^{m_0+i}$, both admitted at $\gamma$ and at $g_2$.

To $(K,\beta)$ and to each side $\ell$ admitted at $g$, Theorem~\textup{\YMMref{AD}{607}} attaches a real number $E_{K,\ell,\beta}$, and its stability bounds say: by the zero-source factorisation and (ii), $-C_g\le\ell^{-3}\log Z_{K,\ell,\beta}-\ell^{-3}E_{K,\ell,\beta}\le C_{\rm S}$, that is, $\ell^{-3}\log Z_{K,\ell,\beta}$ lies within $C_{\rm S}$ above and $C_g$ below the energy per volume $\ell^{-3}E_{K,\ell,\beta}$; by (iii) the energies per volume at two such sides differ by at most $C_{\rm S}$; and by the sentence following (iii), $C_{\rm S}$ is free of $\varrho,\gamma,g,K,\beta,\ell$. At the two tori this gives $p(1,1,1)-p(L^i,L^i,L^i)\le C_{\rm S}+C_{\rm S}+C_g\le C_{\rm v}$ for every $i$; and the brackets of the directions $2$ and $3$ obey the same bound.

Changing one direction at a time, from $(k',k',k')$ through $(k,k',k')$ and $(k,k,k')$ to $(k,k,k)$, with $k'=\ell/\ell_0=L^{m-m_0}$ and a power $k\ge k'$ of $L$, each step changes $p$ by at most $C_{\rm v}/k'$, so $0\le p(k',k',k')-p(k,k,k)\le3C_{\rm v}/k'$, that is, $0\le a_K^{-3}[P_N-P_{k\ell_0L^K}]\le3\ell_0C_{\rm v}/\ell$, as $N=\ell L^K$; let $k\to\infty$ through the powers of $L$ (Lemma~\ref{lh:lat}(a)), and multiply by $b^3=\beta^3a_K^3$.
(d) Theorems~\YMMref{AJ}{673} and~\YMMref{AD}{607} enter through two bounds at first-passage data, with constants free of $K$, $\beta$, $m$ and the state; by (a) the pair $(K,\beta)$ lies in $\mathfrak d(g)$ for some $g\in\{\gamma,g_2\}$. By the common-normalizer comparison in the proof of Proposition~\YMMref{91.13}{709} at $n=2$, whose hypothesis (Nr1) Proposition~\YMMref{91.12}{708} supplies at every pair of $\mathfrak d(g)$, two laws, each a torus measure on an admitted side $2L^m$, $m\ge m_0$, or a periodic Gibbs state at $(K,\beta)$, satisfy $|S^T_2[\mathfrak s](F_1,F_2)-S^T_2[\mathfrak s'](F_1,F_2)|\le C_{\rm nr}\|F_1\|_\sharp\|F_2\|_\sharp$ for $F_1,F_2$ supported in one closed ball of radius $1$ about a lattice point; here $C_{\rm nr}$ is the larger, over $g\in\{\gamma,g_2\}$, of the numbers $8\log2\,(2\mathsf Q_1(L)+2)^2\varrho^{-2}$, with $\mathsf Q_1(L)$ of equation~\eqref{eq:overlap-bound-constants} at $g$.

By Proposition~\textup{\YMMref{91.13}{709}} at $n=2$ and $R=2$, in every periodic Gibbs state $|S^T_2[\nu](F_1,F_2)|\le C_{\rm sp}(\mathsf d)\|F_1\|_\sharp\|F_2\|_\sharp$ for supports in balls of radius $2$ at distance at least $\mathsf d$; here $C_{\rm sp}(\mathsf d)$ is the larger, over $g\in\{\gamma,g_2\}$, of the numbers $C^{\rm per}_2(\mathsf d;2)$ of equation~\eqref{eq:overlap-bound-constants} at $g$. The physical half widths $a_K\lceil r_1\beta\rceil\in[r_1b,r_1b+a_K]$ and $a_K\eta_K=\varepsilon$ lie in $[\mathsf d_1,\frac2{13}]$, so all tents and hulls involved, and the function $\Theta$ of \eqref{lh:tent} rescaled, have $\|\cdot\|_\sharp\le40\sqrt3M/\mathsf d_1$ and live in a ball of radius $1$ or $2$ about one lattice point.

On each interval on which $\lceil r_1\beta\rceil$ is constant, $\beta\mapsto n_{\lceil r_1\beta\rceil}(\mu_{N,\beta})$ is continuous; so $\hat n$ is Borel, and it is attained along a subsequence $\mu_{N_k,\beta}\to\nu_\beta\in\mathcal G(\beta)$, the near part being weakly continuous. Write
\[
n_{\eta_K}(\mu_{N,\beta})-\hat n(\beta)=\bigl[n_{\eta_K}(\mu_{N,\beta})-n_{\eta_K}(\nu_\beta)\bigr]+\bigl[n_{\eta_K}(\nu_\beta)-n_{\lceil r_1\beta\rceil}(\nu_\beta)\bigr]:
\]
the first bracket compares two states in one unit ball, and contributes at most $b_+^5\mathsf d_1^{-3}(40\sqrt3M/\mathsf d_1)^2C_{\rm nr}$ after multiplication by $\beta^5$; this is the one place where $m$ enters, and the bound does not see it. By \eqref{lh:tent} the second bracket is a sum of four means, with weights summing to $1$, of covariances in the one state $\nu_\beta$ between a tent and a function $H$ with $\|H\|_\sharp\le2\cdot40\sqrt3M/\mathsf d_1$ which vanishes within the distance $\mathsf d_1$ of the support of the tent; after multiplication by $\beta^5$ each mean is at most $2b_+^5\mathsf d_1^{-3}(40\sqrt3M/\mathsf d_1)^2C_{\rm sp}(\mathsf d_1)$. With \eqref{lh:far} this gives the first inequality. In the same way $\beta^5|n_{\lceil r_1\beta\rceil}(\nu)-\hat n|$ compares two states in a unit ball, and $\beta^5\lceil r_1\beta\rceil^{-3}s_\nu(3e_1)$ is a covariance at the distance $a_K\lceil r_1\beta\rceil\ge r_1b_-$, non-negative by Lemma~\ref{lh:rp}(a). Finally, every $\beta>Lb_-$ is $L^Kb$ with $b\in\mathsf I$ by (a); on each open interval of $(Lb_-,\infty)$ on which $\lceil r_1\beta\rceil$ is constant, Lemma~\ref{lh:rp}(b) applies with $q=\hat n-(1+870)C^{\rm ref}\beta^{-5}$; the end points of these intervals are countably many, and there $P''$ has atoms $\ge0$ while $q\,d\beta$ has none.
(e) Let $\chi_k$, $k\in(\mathbb Z/\ell\mathbb Z)^3$, be the periodised products $\prod_i(1-|y_i-k_i|)_+$: they sum to $1$ on $T_\ell$, each is supported in a closed sup-ball of radius $1$, $\operatorname{Lip}\chi_k\le\sqrt3$, and so $\|\chi_kw\|_\sharp\le(40\sqrt3M+1)\|w\|_\sharp$ and $\operatorname{supp}(\chi_kw)\subset\operatorname{supp}w$. By bilinearity $S^T_2(u,w)=\sum_kS^T_2(u,\chi_kw)$, a sum of $\ell^3$ terms, each at most $C_{\rm tt}(\mathsf d)(40\sqrt3M+1)\|u\|_\sharp\|w\|_\sharp$ when $u$ is supported in a sup-ball of radius $1$ and the supports of $u$ and $w$ are at distance at least $\mathsf d$. Apply this to $u=\mathsf f^{\langle\ell\rangle}$, $w=1-\mathsf h^{\langle\ell\rangle}$, $\mathsf d=\varepsilon$, with $\|u\|_\sharp,\|w\|_\sharp\le40\sqrt3ML$ and $b^5\varepsilon^{-3}\le b_+^5L^3$.
\end{proof}

\begin{lemma}[Two block sizes]\label{lh:two}
For every $\varrho>0$ and every integer $a\ge1$ there are two settings at the radius $\varrho$ with the block sizes $L_1$ and $L_2=L_1+2$ and one common value $\mathfrak p_0\ge a$. For every odd $L_1\ge3$, $\log L_1/\log(L_1+2)\notin\mathbb Q$.
\end{lemma}

\begin{proof}[Proof sketch]
Use the ordered construction in the proof of Proposition~\YMMref{69.1}{595}: choose $r\ge2$, then an integer $\mathfrak p_0\ge\max\{a,6r+1\}$ and $q\ge41\mathfrak p_0+1$, followed by the reserve fractions and the other constants preceding $L$. The proof admits every sufficiently large odd $L$ with these values fixed. The additional floors in Section~\YMMref{82}{665} have the same property. Choose an odd $L_1$ above their common threshold and complete the prescribed construction separately at $L_1$ and $L_1+2$; with $\varrho$ and the associated Wilson measures and source normalization, each is a setting. If $\log L_1/\log L_2=u/v$ with positive integers, then $L_1^v=L_2^u$, and a prime factor of $L_1$ divides $L_2-L_1=2$, which is impossible for odd $L_1$. Every statement that we use is made for every radius and every choice, and no object of one choice is compared with an object of the other; the two settings meet in $P$ and $\hat n$ alone.
\end{proof}

\begin{lemma}[The lost mass on tori, without an exponential of the volume]\label{lh:lost}
Let $\varsigma'$ be a lowered choice of a choice $\varsigma$ of $\mathcal T^{\rm F+}$, and let the setting (Th) of Section~\ref{sec:cutoff-calibration} hold at $\varsigma$ and at a label $b\ge\gamma^{-2}$ with $\log b\ge\mathsf Y^\natural$: $K\ge2$, $\ell=2L^m$ with $M\mathcal R(b)\le L^m$, where $\mathcal R(x)=\min\{L^c:c\in\mathbb Z_{\ge0},\ L^c\ge(\log x)^r\}$ (equation~\YMMref{(634)}{339}), the line of couplings $x_j=L^{K-j}b$, and a top level. Let $Z=\sum_h\mathsf t(h)$ be the sum of the masses $\mathsf t(h)$ of the histories (Theorem~\YMMref{D}{8}), $Z_t(h)=\Lambda_t(h)^c$ the large-field region defined by equation~\YMMref{(644)}{341}, and $N_t=\ell^3L^{3(K-t)}$. Then for $1\le t\le K$
\[
\sum_{h:\,Z_t(h)\ne\emptyset}\mathsf t(h)\le\mathsf R_tZ,\qquad\mathsf R_t=\tfrac L{L-1}M^{-3}N_t\,e^{-c_2^\natural A_0^2(\log x_t)^{2\mathfrak p_0}}.
\]
The side enters $\mathsf R_t$ only through the number $N_t$ of the sites of the level $t$; and $\mathsf R_t$ is smaller by the factor $25000$ than the number $2\cdot24^3M^{-3}N_te^{-c_2^\natural A_0^2(\log x_t)^{2\mathfrak p_0}}$ that the rarity of a large field at a point, Proposition~\ref{prop:local-large-field-rarity}, gives.
\end{lemma}

\begin{proof}
The line $x_j=L^{K-j}b$ is admissible under the stated side and coupling
conditions, by Lemma~\YMMref{118.1}{812}. Definition~\ref{def:rarity-choices}
and Lemma~\ref{lem:rarity-lowered-choice} put it in the common setting
of Proposition~\ref{prop:local-large-field-rarity}. That proposition
bounds the probability of a prescribed creation at level $t'$ by
$e^{-c_2^\natural p_0(g_{t'})^2}$.

If $Z_t\ne\emptyset$, Proposition~\YMMref{67.2}{588} gives a creation
at some level $t'\le t$. There are at most
$M^{-3}N_{t'}=M^{-3}N_tL^{3(t-t')}$ cells of that level, since
$\check R_{t'}\ge1$. Write $a=t-t'$. On this line,
$y_{t'}=y_t+a\log L$, and convexity gives
\[
 y_{t'}^{2\mathfrak p_0}-y_t^{2\mathfrak p_0}
 \ge2\mathfrak p_0y_t^{2\mathfrak p_0-1}a\log L.
\]
Condition (Y.7) of \eqref{eq:rarity-seven-conditions} therefore bounds
the contribution of level $t-a$ by $L^{-a}$ times the contribution of
level $t$. The union bound and the geometric series yield
\[
 \sum_{h:Z_t(h)\ne\emptyset}\mathsf t(h)
 \le M^{-3}N_te^{-c_2^\natural p_0(g_t)^2}Z
                 \sum_{a=0}^{t-1}L^{-a}
 \le\frac L{L-1}M^{-3}N_te^{-c_2^\natural p_0(g_t)^2}Z.
\]
Finally $L/(L-1)\le13/12<2\cdot24^3/25000$.
\end{proof}

\begin{lemma}[The threshold on a torus is of exact order $\ell^{3/2\mathfrak p_0}$]\label{lh:thr}
In a setting let $K_0(\ell)=K_0(\ell,2,\mathsf I)$ and $K_1(\ell)=K_1(\ell,2,\mathsf I,\varepsilon)$ be the thresholds of the convergence on the torus $T_\ell$ (Theorem~\textup{\YMMref{AS}{813}} at $s=0$, with the constants and thresholds in Theorem~\YMMref{AR}{810}) at $n=2$, for the labels of $\mathsf I$ and the spread $\varepsilon$, as they are set there. There are $c_1,c_2<\infty$ such that for all $m\ge m_0$ with $M\mathcal R(b_+)\le L^m$, $\ell=2L^m$,
\[
\frac{1800}\lambda\Bigl[\Bigl(\frac{\Xi\ell^3}{c'A_0^2}\Bigr)^{1/2\mathfrak p_0}-\log(Lb_-)\Bigr]\le K_0(\ell)\le K_1(\ell)\le c_1+c_2\ell^{3/2\mathfrak p_0},
\]
where $\Xi>0$ is the number in the factor $e^{\Xi\ell^3}$ that the bound on the lost mass, with which these thresholds are formed, carries, and $c'>0$ the number in its large-field weight $e^{-c'A_0^2(\log x)^{2\mathfrak p_0}}$.
\end{lemma}

\begin{proof}[Proof sketch]
The conditions that define the thresholds read the side only through $N_{[d]}=\ell^3L^{3d}$ at depths $d\le K/9+1$ and through the factor $e^{\Xi\ell^3}$. At $n=2$ the depth $u$ of the comparison of the cutoffs satisfies $K/1800\le u\le K/1800+1$. One condition asks that the bound on the lost mass, $2M^{-3}\ell^3L^{3u}\exp[\Xi\ell^3-c'A_0^2(\log(L^ub_-))^{2\mathfrak p_0}]$, be at most $\frac12$; as $2M^{-3}\ell^3\ge16$ and $\log(L^ub_-)\le K\lambda/1800+\log(Lb_-)$, it fails unless $(K\lambda/1800+\log(Lb_-))^{2\mathfrak p_0}\ge\Xi\ell^3/c'A_0^2$, which is the lower bound. For the upper bound, each condition says that a quantity $c\,\ell^3L^{-cK}$, or $c\exp[c\log K+c\log\ell+cK+\Xi\ell^3-c''K^{2\mathfrak p_0}]$ with $c''>0$ formed from $A_0$, $A_1$ and the constants that enter the definition of the thresholds, is at most a fixed number. For $K\ge c_2\ell^{3/2\mathfrak p_0}$ we have $\Xi\ell^3\le\frac12c''K^{2\mathfrak p_0}$ and $\log\ell\le\frac{2\mathfrak p_0}{3c_2}K$, and all conditions hold once also $K\ge c_1$, because $2\mathfrak p_0\ge2$.
\end{proof}

\begin{lemma}[One cutoff: conditional expectations]\label{lh:cond}
Let $K\in\mathbb N$, and let $\mathcal G$ be a set of laws on $\mathbb L_K$, each invariant under the isometries of $\mathbb L_K$, such that $h\nu/\nu(h)\in\mathcal G$ for $\nu\in\mathcal G$ and every bounded $\mathcal I$-measurable $h\ge0$ with $\nu(h)>0$. For $\rho\in\mathcal G$, $e_\rho=\rho(a_p)$ is one number for all plaquettes $p$; put $e_+=\sup_{\rho\in\mathcal G}e_\rho$, $e_-=\inf_{\rho\in\mathcal G}e_\rho$, and let $\nu\in\mathcal G$.
\begin{enumerate}
\item[(a)] There is one bounded $\mathcal I$-measurable $\mathsf e$ with $\mathsf e=E_\nu[a_p|\mathcal I]$ $\nu$-almost surely for every plaquette $p$; $e_-\le\mathsf e\le e_+$ $\nu$-almost surely, and $\operatorname{Var}_\nu\mathsf e\le\frac14(e_+-e_-)^2$.
\item[(b)] For bounded $f$ of compact support, $E_\nu[\mathcal O_K(f)-\nu(\mathcal O_K(f))\,|\,\mathcal I]=3a_K^{-4}r_K(f)(\mathsf e-e_\nu)$ $\nu$-almost surely.
\item[(c)] For every bounded real Borel function $F$, $E_\nu[F|\mathcal I]$ takes $\nu$-almost surely values in $[\inf_{\rho\in\mathcal G}\rho(F),\sup_{\rho\in\mathcal G}\rho(F)]$.
\item[(d)] If $\nu$ is ergodic, or $e_\nu\in\{e_-,e_+\}$, then $\mathsf e$ is $\nu$-almost surely constant.
\end{enumerate}
\end{lemma}

\begin{proof}[Proof sketch]
Let $k$ be bounded and $\mathcal I$-measurable, $c=1+\sup|k|$ and $h=k+c\ge1$. The laws $h\nu/\nu(h)$ and $\nu$ lie in $\mathcal G$, so they are invariant, and the isometries act transitively on the plaquettes; so $\nu(ha_p)=\nu(ha_{p'})$ and $\nu(ca_p)=\nu(ca_{p'})$ for all $p,p'$, and subtracting, $\nu(ka_p)=\nu(ka_{p'})$: all $a_p$ have one conditional expectation given $\mathcal I$. (c): if $\{E_\nu[F|\mathcal I]>\sup_\rho\rho(F)+\epsilon\}$ had positive mass, conditioning on it would give a member of $\mathcal G$ in which the mean of $F$ exceeds $\sup_\rho\rho(F)$; likewise below. At $F=a_p$ this gives the range of $\mathsf e$, and a variable with values in an interval of length $e_+-e_-$ has variance at most $\frac14(e_+-e_-)^2$. (b): linearity, and $a_K^{-1}\sum_pf(x(p))=3a_K^{-4}r_K(f)$. (d): if $\nu$ is ergodic, $\mathcal I$ is trivial modulo $\nu$; if $\nu(\mathsf e)=e_\nu=e_+$, the function $\mathsf e\le e_+$ has the mean $e_+$, so it is constant; likewise at $e_-$.
\end{proof}

\begin{proposition}[One cutoff: the plateau for several fields]\label{lh:many}
Let $\mathcal G$ be a class at the cutoff $K$ (Section~\YMMref{100}{743}) that obeys the bound system $\mathsf C^g=(\hat C^g_1,(\hat C^g_n)_{n\ge2})$ at $g$, with the properties listed in the paragraph on one cutoff: for $n\ge2$ and every tuple $X=(f_1,\dots,f_n)$ of type $(\mathsf d;R)$, $|S^T_n[\rho](X)|\le\hat C^g_n(\mathsf d;R)\prod_i\|f_i\|_\sharp$ for $\rho\in\mathcal G$, and $|\rho(\mathcal O_K(f))-\rho'(\mathcal O_K(f))|\le\hat C^g_1(R)\|f\|_\sharp$ for $\rho,\rho'\in\mathcal G$ and $f\in C^1_c(\mathbb R^3)$ with support in a closed ball of radius $R$; the numbers $\delta(X;\bar R)$ and $\mathsf N(X;\bar R)$ are those of the plateau theorem, the second formed with this bound system. Let $\nu\in\mathcal G$. Let $X,Y$ be separated, $\bar R\in a_K\mathbb N$ with all supports in $(-\bar R,\bar R)^3$, $a_K\le2\min\{\delta(X;\bar R),\delta(Y;\bar R)\}$, $v\in a_K\mathbb Z^3$ and $|v|\ge4\sqrt3\bar R$. Then
\[
\bigl|\sigma_\nu(X\cdot\tau_vY)-\sigma_\nu(X)\sigma_\nu(Y)-\operatorname{Cov}_\nu\bigl(E_\nu[X^\nu_K|\mathcal I],E_\nu[Y^\nu_K|\mathcal I]\bigr)\bigr|\le\mathfrak k_0[\mathsf N(X;\bar R)\mathsf N(Y;\bar R)]^{1/2}\bar R/|v|,
\]
\[
\bigl|\operatorname{Cov}_\nu\bigl(E_\nu[X^\nu_K|\mathcal I],E_\nu[Y^\nu_K|\mathcal I]\bigr)\bigr|\le\tfrac14\operatorname{osc}_{\mathcal G}(X^\nu_K)\operatorname{osc}_{\mathcal G}(Y^\nu_K),\qquad\operatorname{osc}_{\mathcal G}(F)=\sup_{\rho\in\mathcal G}\rho(F)-\inf_{\rho\in\mathcal G}\rho(F);
\]
and $\operatorname{osc}_{\mathcal G}((f)^\nu_K)=3a_K^{-4}|r_K(f)|(e_+-e_-)$.
\end{proposition}

\begin{proof}[Proof sketch]
Clause (b) of the plateau theorem at $\nu$ compares $\sigma_\nu(X\cdot\tau_vY)$ with 
\[
\nu(E_\nu[X^\nu_K|\mathcal I]E_\nu[Y^\nu_K|\mathcal I]),
\]
which is $\sigma_\nu(X)\sigma_\nu(Y)$ plus the covariance, as $\nu(E_\nu[A|\mathcal I])=\nu(A)$. By Lemma~\ref{lh:cond}(c) each conditional expectation lies almost surely in an interval of length $\operatorname{osc}_{\mathcal G}$, so its variance is at most $\frac14\operatorname{osc}_{\mathcal G}^2$; then the Cauchy--Schwarz inequality. At $X=(f)$, $\rho((f)^\nu_K)=3a_K^{-4}r_K(f)(e_\rho-e_\nu)$.
\end{proof}

\begin{lemma}[Recentring and clustering]\label{lh:rec}
\begin{enumerate}
\item[(a)] For laws $\rho,\nu,\nu'$ on $\mathbb L_K$ and $X=(f_1,\dots,f_n)$,
\begin{equation}\label{lh:rece}
\begin{aligned}
\rho(X^{\nu'}_K)={}&\sum_{A\subset\{1,\dots,n\}}\ \prod_{i\in A}
\bigl[\nu(\mathcal O_K(f_i))-\nu'(\mathcal O_K(f_i))\bigr]\\
&\qquad\rho\bigl((X_{A^c})^\nu_K\bigr).
\end{aligned}
\end{equation}
\item[(b)] Let $0<g\le\gamma$ and $B^g_n(\mathsf d;R)=\sum_\pi\prod_{A\in\pi}\hat C^g_{|A|}(\mathsf d;R)$, the sum over the partitions of $\{1,\dots,n\}$ into blocks of at least two elements ($B^g_0=1$, $B^g_1=0$). The set $\mathcal E^g$ of the symmetric functionals $\Phi$ with $|\Phi(X)|\le B^g_n(\mathsf d;R)\prod_i\|f_i\|_\sharp$ for every $X$ of type $(\mathsf d;R)$ is convex, and compact and metrisable for pointwise convergence; it contains $\sigma_\nu$ for $(K,\beta)\in\mathfrak d(g)$ and $\nu\in\mathcal G_{K,\beta}$.
\item[(c)] Let $\mathsf T$ be a compact metrisable set of symmetric functionals that are real-valued on real tuples, with pointwise convergence, such that for all separated $X,Y$ there are $r_{X,Y},c_{X,Y}<\infty$ with $\sup_{\psi\in\mathsf T}|\psi(X\cdot\tau_vY)-\psi(X)\psi(Y)|\le c_{X,Y}/|v|$ for $|v|\ge r_{X,Y}$. Then for every Borel probability measure $\mu$ on $\mathsf T$, $(\operatorname{supp}\mu,\mu,(\psi)_{\psi\in\operatorname{supp}\mu})$ is a canonical decomposition of $\operatorname{bar}(\mu)$, and $\mu\mapsto\operatorname{bar}(\mu)$ is affine, continuous and injective.
\item[(d)] Let $0<g\le\gamma$, $(K,\beta)\in\mathfrak d(g)$, $\nu^1,\nu^2\in\mathcal G_{K,\beta}$, $0\le t\le1$, $\nu=t\nu^1+(1-t)\nu^2$, $X$ of $n$ entries and of type $(\mathsf d;R)$, and $\eta=a_K^{-4}J(\beta)(2R+1)^3$. Then
\[
\bigl|\sigma_\nu(X)-t\sigma_{\nu^1}(X)-(1-t)\sigma_{\nu^2}(X)\bigr|\le\sum_{k=1}^n\binom nk\eta^kB^g_{n-k}(\mathsf d;R)\prod_i\|f_i\|_\sharp ,
\]
and at $n=2$ exactly 
\[
\sigma_\nu((f,h))=t\sigma_{\nu^1}((f,h))+(1-t)\sigma_{\nu^2}((f,h))+9t(1-t)a_K^{-8}r_K(f)r_K(h)(e_{\nu^1}-e_{\nu^2})^2.
\]
With (C), $\eta\le(2R+1)^3g^{8+2c}c_*(c)L^{-cK}$ for $0\le c<\frac1{180}$: the defect of convexity of $\nu\mapsto\sigma_\nu$ at one cutoff falls like a power of the spacing.
\end{enumerate}
\end{lemma}

\begin{proof}[Proof sketch]
(a) Write $\mathcal O_K(f_i)-\nu'(\mathcal O_K(f_i))=[\mathcal O_K(f_i)-\nu(\mathcal O_K(f_i))]+[\nu-\nu'](\mathcal O_K(f_i))$ and multiply out. (b) A centred moment $\sigma_\nu(X)$ is the sum, over the partitions without singletons, of the products of the cumulants of the blocks, which the family of bounds at $g$ bounds; the defining conditions are inequalities and equalities at finitely many tuples, so they pass to convex combinations and to pointwise limits, and Tychonoff's theorem gives compactness. Approximating each entry of a tuple in $\|\cdot\|_\sharp$ by functions from a countable family supported in finite unions of rational cubes, and using multilinearity with the bound, one finds a countable set of tuples on which two members of $\mathcal E^g$ agree only if they are equal, and on which pointwise convergence implies convergence at every tuple; so $\mathcal E^g$ is metrisable. (c) By induction on $q$, the uniform clustering shows that $\int\prod_{i\le q}\psi(X^i)\,d\mu$ is the iterated plateau of $\operatorname{bar}(\mu)$ at $(X^1,\dots,X^q)$: this is (Can4); the other conditions are plain.

So $\operatorname{bar}(\mu)$ determines the integrals of the products of the functions $\psi\mapsto\psi(X)$, which are real-valued and separate the points of $\mathsf T$; their products and constants form a real unital algebra, so by the Stone--Weierstrass theorem it determines $\mu$. The barycentre is affine in $\mu$, and continuous since the integrands are continuous. (d) $\sigma_\nu(X)=t\nu^1(X^\nu_K)+(1-t)\nu^2(X^\nu_K)$, and by \eqref{lh:rece} with $\rho=\nu^i$, $\nu'=\nu$ and the centring $\nu^i$, $\nu^i(X^\nu_K)-\sigma_{\nu^i}(X)$ is a sum over $A\ne\emptyset$ of products of $|A|$ one-point differences, each at most $a_K^{-4}J(\beta)(2R+a_K)^3\|f\|_\infty\le\eta\|f\|_\sharp$ by (C), times $\sigma_{\nu^i}(X_{A^c})$, which is at most $B^g_{n-|A|}(\mathsf d;R)\prod_{A^c}\|f_i\|_\sharp$ by (b).

At $n=2$, $\operatorname{Cov}_\nu(F,H)=t\operatorname{Cov}_{\nu^1}(F,H)+(1-t)\operatorname{Cov}_{\nu^2}(F,H)+t(1-t)[\nu^1-\nu^2](F)\,[\nu^1-\nu^2](H)$ for $F=\mathcal O_K(f)$, $H=\mathcal O_K(h)$, and $[\nu^1-\nu^2](\mathcal O_K(f))=3a_K^{-4}r_K(f)(e_{\nu^1}-e_{\nu^2})$.
\end{proof}

\begin{remark}[The one-point differences and the convexity of the limits]\label{lh:arc}
In the proof of (G)(1) the vanishing of the one-point differences, (D)(1), makes the limits of the mixtures the mixtures of the limits; the identity at $n=2$ in (d) shows, for data, what can happen without it. Let, along a sequence of pairs, the class at the $j$-th pair be the segment $\{s\nu^+_j+(1-s)\nu^-_j:0\le s\le1\}$ between two ergodic laws $\nu^\pm_j$ whose one-point functions differ by an amount that tends to $c\int f$ with $c\ne0$, and let $\sigma_{\nu^\pm_j}\to\sigma^\pm$. Then the mixtures $\nu^s_j=s\nu^+_j+(1-s)\nu^-_j$ have
\[
\sigma_{\nu^s_j}((f,h))\to\sigma_s((f,h))=s\sigma^+((f,h))+(1-s)\sigma^-((f,h))+s(1-s)c^2\int f\int h .
\]
If moreover $(f,h)\mapsto\sigma^+((f,h))-\sigma^-((f,h))$ is not a multiple of $(f,h)\mapsto\int f\int h$, the limits $\sigma_s$, $0\le s\le1$, form an arc and not a segment, and the set of the limits along the sequence, which is $\{\sigma_s:0\le s\le1\}$, does not contain the midpoint of $\sigma^+$ and $\sigma^-$: it is not convex. For the model, (D)(1) excludes such sequences among the periodic Gibbs states at first-passage data.
\end{remark}

\begin{proposition}[What a rate on the tori gives]\label{lh:pow}
Let $(\mathrm S_1)$, $(\mathrm S_2)$ and $r_1$ be as in (A)(2); no lowered choice is given. Suppose that for $i=1,2$ there are numbers $C_i\ge1$, $c_{1,i}\ge0$, integers $m_{1,i}\ge m_{0,i}$ and, for $m\ge m_{1,i}$, $\ell=2L_i^m$ and $b\in\mathsf I_i$, real numbers $S^{(\ell,b)}$ such that for all integers $K\ge C_i\log\ell$
\[
\bigl|S^T_{2;K,\ell,L_i^Kb}(f_1,f_2)-S^{(\ell,b)}\bigr|\le C_i\ell^{c_{1,i}}L_i^{-K/C_i}\|f_1\|_\sharp\|f_2\|_\sharp\qquad\text{at the one pair }f_1=\mathsf f^{\langle\ell\rangle},\ f_2=1-\mathsf h^{\langle\ell\rangle}.
\]
Then the oscillation hypothesis of Proposition~\ref{lh:per}(b) holds for $m\ge m_{1,i}$, with $\mathsf a_i=C_i\log2$, $A_i=C_i\max\{\lambda_i,c_{1,i}\}$ and $\omega_i=2^{1+c_{1,i}}\cdot4800b_{+,i}^5M_i^2L_i^5C_i$, and the two bounds of (B)(1) hold with the power $\beta^{-1/2A}$, $A=\max_iA_i$, beyond some $\beta_1$. More generally, any oscillation bound for the far parts with thresholds $\mathsf a_i+A_im$ on all sides $m\ge m_{1,i}$ gives $\beta^4J(\beta)\le C\beta^{-1/2\max_iA_i}$ for large $\beta$.
\end{proposition}

\begin{proof}[Proof sketch]
Index $i$ dropped. For $m\ge m_1$ and integers $K,K'\ge\mathsf a+Am$ we have $K\ge C\log2+C\lambda m=C\log\ell$, the same for $K'$, and $\ell^{c_1}L^{-K/C}=2^{c_1}L^{c_1m-K/C}\le2^{c_1}$, since $K/C\ge c_1m$. So the two torus functions differ by at most $2C2^{c_1}\|f_1\|_\sharp\|f_2\|_\sharp$; with $b^5\varepsilon^{-3}\le b_+^5L^3$ and $\|f_1\|_\sharp,\|f_2\|_\sharp\le40\sqrt3ML$, $\sup_{\log\mathsf I}|y_{m,K}-y_{m,K'}|\le\omega$. Enter in Proposition~\ref{lh:per}(b) as in Step~1 of the proof of Theorem~\ref{lh:main} below, with $m^0_i=m_{1,i}$; the second sentence is the same argument with the given thresholds.
\end{proof}

\begin{proposition}[A threshold that does not grow with the side gives $J=0$ on a half line]\label{lh:cut}
Let $(\mathrm S_1)$, $(\mathrm S_2)$ and $r_1$ be as in (A)(2). Suppose that there are an integer $K_*$ and numbers $\omega_i$ such that $\sup_{\log\mathsf I_i}|y^{(i)}_{m,K}-y^{(i)}_{m,K'}|\le\omega_i$ for all integers $K,K'\ge K_*$ and all $m\ge m_{0,i}$, $i=1,2$. Then there is $\beta_*$ with $J(\beta)=0$ for $\beta\ge\beta_*$, and on $(\beta_*,\infty)$ the measure $\beta^4[P''(d\beta)-\hat n\,d\beta]$ has a bounded density. In terms of one lattice: with $Y_N(\beta)=\beta^5[P_N''(\beta)-\hat n(\beta)]$, Proposition~\ref{lh:blk}(d) puts $Y_N(\beta)$ within $C^{\rm ref}$ of $y_{m,K}(\log b)$ for $N=2L^{m+K}$, $\beta=L^Kb$, $b\in\mathsf I$, $m\ge m_0$. Given lowered choices, the oscillation bound of Step~5 of the proof of Theorem~\ref{lh:main} gives $|Y_{L^kN}(L^k\beta)-Y_N(\beta)|\le\omega_A+2C^{\rm ref}$ for all $k\ge0$ at these $N$, $\beta$ with $N/\beta\le2b_+^{-1-1/A}L^{-k^{\rm rt}/A}\beta^{1/A}$; the hypothesis here asks such a bound at all of them.
\end{proposition}

\begin{proof}[Proof sketch]
Enter in Proposition~\ref{lh:per}(c) with the data of Step~1 of the proof of Theorem~\ref{lh:main} and with $\max\{K_*,k_1,k_2\}$ in the place of $K_*$: $\epsilon^{(i)}_m=c_{{\rm P},i}L_i^{-m}$ has the lower bound $0$, and the hypothesis on $Y^0$ holds at $m^0=m_{0,1}$ with $k^0=\max\{K_*,k_1\}$, since $|y^{(1)}_{m_{0,1},K}|\le|y^{(1)}_{m_{0,1},k^0}|+\omega_1$ for $K\ge k^0$. Since $P$ is the pressure, $J_P=J$. For the last sentence: $L^kN=2L^{m+K+k}$ and $L^k\beta=L^{K+k}b$, and $K\ge k^{\rm rt}+Am$ reads $N/\beta\le2b^{-1-1/A}L^{-k^{\rm rt}/A}\beta^{1/A}$.
\end{proof}

\begin{proposition}[The inequality of two boxes at one pair, and along a label]\label{lh:onepair}
Let $a>0$ and $\beta>0$. For the periodic boxes $\underline n=(n_1,n_2,n_3)$, $n_\mu\ge3$, of the lattice $a\mathbb Z^3$ at the coupling $\beta$ let $F(\underline n)=\log Z(\underline n)$, $|\underline n|=n_1n_2n_3$, $\Delta(\underline n)=F(\underline n)-|\underline n|P(\beta)$, write $F(N^3)$ and $\Delta(N^3)$ for the cube of side $N$, and $\mathcal E(l)=2F(l,2l,2l)-F(2l,2l,2l)$.
\begin{enumerate}
\item[(a)] $\Delta(\underline n)\ge0$; $\Delta(\underline n')/|\underline n'|\le\Delta(\underline n)/|\underline n|$ for $\underline n'\ge\underline n$ in each component; and $\mathcal E(N)\le8\Delta(N^3)-\Delta((2N)^3)\le8\Delta(N^3)$ for $N\ge3$.
\item[(b)] If $\mathcal E(l)\le x_0$ for one integer $l\ge3$ and one number $0<x_0\le0.069$, then, with $\theta=12.9x_0\le0.8901$: $\mathcal E(2l)\le12.9\,\mathcal E(l)^2$; $\Delta(N'^3)<0.13736\,\theta^{N'/2l}$ for every $N'\ge2l$; and $\mathcal E(l')\le0.743\,\theta^{l'/2l}$ for every $l'\ge l$.
\item[(c)] If $\mathcal E(l)\le0.069$ for one integer $l\ge3$, the lattice $a\mathbb Z^3$ has exactly one periodic Gibbs state at $\beta$; it is a pure phase, and $J(\beta)=0$.
\item[(d)] If the lattice has two periodic Gibbs states at $\beta$, then $\mathcal E(l)>0.069$ for every $l\ge3$, $\Delta(N^3)>0.0086$ for every $N\ge3$, and $\liminf_{j\to\infty}\Delta((2L^j)^3)\ge\log2$ for every odd integer $L\ge13$.
\item[(e)] At every $\beta>0$, either $\mathcal E(l)>0.069$ for every $l\ge3$, or $\mathcal E(l')\to0$ exponentially in $l'$.
\item[(f)] In a setting suppose $\mathrm{BI}[b,\ell_*,K_*;x_0]$, let $\underline K=\max\{K_*,1\}$, $\theta=12.9x_0$, and let $\nu_K$ be the one periodic Gibbs state of (c) at $a=a_K$, $\beta=L^Kb$. Let $n\ge2$, and let $X=(f_1,\dots,f_n)$ be real Lipschitz functions, none of them $0$, with pairwise disjoint supports in $(-\frac{\ell''}4,\frac{\ell''}4)^3$ for a side $\ell''=2L^{m''}\ge\ell_*$, $m''\in\mathbb N$. There is $C(X)<\infty$, formed from $X$, $\ell''$ and $b$ and from no cutoff and no further side, such that for all $K\ge\underline K$ and all sides $\ell'=2L^{m'}\ge\ell''$
\[
\bigl|S^T_n[\nu_K](X)-S^T_{n;K,\ell',L^Kb}(X^{\langle\ell'\rangle})\bigr|\le C(X)\,\theta^{\ell'/4\ell_*}.
\]
\end{enumerate}
\end{proposition}

\begin{proof}[Proof sketch]
(a) Fix a direction of a box as the time and the two other sides, $n_\nu\ge3$, the section $\Sigma$, and write $F(t,\Sigma)$ and $Z(t,\Sigma)$ for the box with $t\ge3$ layers. On the Hilbert space of the configurations of the bonds of one layer (Section~\YMMref{85}{677}) let $\mathsf C_\Sigma$ be the crossing operator of Section~\YMMref{85}{677}, written $\mathsf K$ there, and $\mathsf D_\Sigma$ the multiplication by $e^{-\beta A_\Sigma}$, the weight of the plaquettes of one layer. $\mathsf C_\Sigma$ is positive (Proposition~\YMMref{85.1}{678}), and it is of trace class: the series the rank-one expansion in the proof of Proposition~\YMMref{85.1}{678} in the proof of that clause writes it as $\sum_\alpha c_\alpha|\chi_\alpha\rangle\langle\chi_\alpha|$ with $\|\chi_\alpha\|\le1$ and $c_\alpha\ge0$, $\sum_\alpha c_\alpha<\infty$ the rank-one expansion in the proof of Proposition~\YMMref{85.1}{678}.

So the transfer operator of one layer, 
\[
T_\Sigma:=\mathsf C_\Sigma^{1/2}\mathsf D_\Sigma\mathsf C_\Sigma^{1/2},
\]
is positive, as 
\[
\langle\psi,T_\Sigma\psi\rangle=\int e^{-\beta A_\Sigma}|\mathsf C_\Sigma^{1/2}\psi|^2\,dV,
\]
and of trace class; like $\mathsf C_\Sigma$ and $\mathsf D_\Sigma$ it does not depend on the number of layers (Proposition~\YMMref{85.2}{679}). Integrating the variables of the bonds between two layers turns the weight of a slab into a chain of kernels: the slab operator of the function $1$ on $\vartheta\ge2$ layers is $\mathsf D_\Sigma(\mathsf C_\Sigma\mathsf D_\Sigma)^{\vartheta-1}$, so that 
\[
\mathfrak Y[1]=\mathsf C_\Sigma^{1/2}\mathsf D_\Sigma(\mathsf C_\Sigma\mathsf D_\Sigma)^{\vartheta-1}\mathsf C_\Sigma^{1/2}=T_\Sigma^\vartheta
\]
 (equation~\YMMref{(1744)}{679}).

The trace formula, Proposition~\YMMref{85.2}{679} with the function $1$ on two slabs of $2$ and $t-2$ layers, gives $Z(t,\Sigma)=\operatorname{Tr}T_\Sigma^t$ for $t\ge4$; at $t=3$ the integral over the three layers is $\operatorname{Tr}(\mathsf C_\Sigma\mathsf D_\Sigma)^3$, by the Hilbert--Schmidt kernel identity $\operatorname{Tr}(AB)=\iint A(y,x)B(x,y)\,dx\,dy$ used in the proof of Proposition~\YMMref{85.2}{679}, applied to the operators $(\mathsf C_\Sigma\mathsf D_\Sigma)^2$ and $\mathsf C_\Sigma\mathsf D_\Sigma$, which have bounded kernels, and this is $\operatorname{Tr}T_\Sigma^3$ by the cyclicity of the trace. With the eigenvalues $\mu_i\ge0$ of the positive compact operator $T_\Sigma$, counted with their multiplicities, the trace of a power is the sum of the powers of the eigenvalues, by computing the trace in an orthonormal eigenbasis, and
\[
e^{F(t,\Sigma)}=Z(t,\Sigma)=\operatorname{Tr}T_\Sigma^t=\sum_i\mu_i^t\qquad\text{for every integer }t\ge3,
\]
with numbers $\mu_i\ge0$ that do not depend on $t$, odd $t$ and $t=3$ included. By the spectral argument in the proof of Proposition~\YMMref{86.5}{682}, $(\operatorname{Tr}T_\Sigma^t)^{1/t}$ does not increase in $t$; so $F/|\underline n|$ does not increase when a side grows, and as $F(N^3)/N^3\to P$ (Lemma~\ref{lh:lat}(a)), $P=\inf F/|\underline n|$, and $\Delta\ge0$ with the monotonicity per site. In particular $F(N,2N,2N)\le4F(N^3)$, so $\mathcal E(N)\le8F(N^3)-F((2N)^3)=8\Delta(N^3)-\Delta((2N)^3)$.
(b) \emph{One direction.} Fix a direction of a box as the time and the two other sides, the section $\Sigma$, and write $F(t,\Sigma)$ for the box with $t\ge3$ layers. By (a), with $\lambda_0>0$ the largest of the numbers $\mu_i$ and $\mathsf y_j\in[0,1]$ the quotients $\mu_i/\lambda_0$ of the others, $F(t,\Sigma)=t\log\lambda_0+\log(1+\mathsf r_t)$ with $\mathsf r_t=\sum_j\mathsf y_j^{\,t}$. So $\mathsf r_t$ does not increase in $t$, $\mathsf r_{2t}\le\mathsf r_t^2$, and
\[
\mathsf d(t):=2F(t,\Sigma)-F(2t,\Sigma)=2\log(1+\mathsf r_t)-\log(1+\mathsf r_{2t})\ \ge\ \log(1+\mathsf r_t)\ \ge0 .
\]
Let $\mathsf d(t)\le x<\log2$, $\eta=e^x-1<1$ and $\rho(\eta)=\eta/(1+\sqrt{1-\eta^2})$. Then $\mathsf r_t\le\eta$, and $1+2\mathsf r_t/(1+\mathsf r_t^2)\le(1+\mathsf r_t)^2/(1+\mathsf r_{2t})=e^{\mathsf d(t)}\le1+\eta$; as $r\mapsto2r/(1+r^2)$ increases on $[0,1]$ and takes the value $\eta$ at $r=\rho(\eta)$, we get $\mathsf r_t\le\rho(\eta)$, and so
\[
\mathsf d(2t)\le2\log(1+\mathsf r_{2t})\le2\log(1+\mathsf r_t^2)\le\psi(x):=2\log\bigl(1+\rho(e^x-1)^2\bigr).
\]
The function $\psi$ increases on $[0,\log2)$, and, as $(e^y-1)/y$ increases, $\psi(y)\le2\kappa(y_0)^2y^2$ for $0\le y\le y_0<\log2$, with $\kappa(y_0)=\frac{e^{y_0}-1}{y_0}\bigl(1+\sqrt{1-(e^{y_0}-1)^2}\bigr)^{-1}$; and $\kappa(0.069)\le0.51832$, $\kappa(0.14056)\le0.53995$, $\kappa(0.29264)\le0.59869$.

\emph{Doubling.} For a box $\underline n$ and a direction $\mu$ let $\mathsf d^{(\mu)}(\underline n)$ be $2F(\underline n)$ minus $F$ of the box with $n_\mu$ doubled. By the above, read in the direction $\mu$, $\mathsf d^{(\mu)}\ge0$, and $\mathsf d^{(\mu)}$ of the box with $n_\mu$ doubled is at most $\psi(\mathsf d^{(\mu)}(\underline n))$ when $\mathsf d^{(\mu)}(\underline n)<\log2$. Let $\epsilon=\mathcal E(l)\le0.069$. As $F$ does not depend on the order of the three sides, $\epsilon=\mathsf d^{(1)}(l,2l,2l)=\mathsf d^{(2)}(2l,l,2l)=\mathsf d^{(3)}(2l,2l,l)$. Then:
\begin{enumerate}
\item[(i)] $\mathsf a:=\mathsf d^{(2)}(2l,2l,2l)\le\psi(\epsilon)\le2\kappa(0.069)^2\epsilon^2\le0.53732\,\epsilon^2\le0.03708\,\epsilon$;
\item[(ii)] since $F(2l,4l,l)\le2F(2l,2l,l)$ and $F(2l,4l,2l)=2F(2l,2l,2l)-\mathsf a$, we have $\mathsf d^{(3)}(2l,4l,l)\le2\,\mathsf d^{(3)}(2l,2l,l)+\mathsf a=2\epsilon+\mathsf a\le2.03708\,\epsilon\le0.14056$, so $\mathsf b:=\mathsf d^{(3)}(2l,4l,2l)\le\psi(2\epsilon+\mathsf a)\le2.41970\,\epsilon^2\le0.16696\,\epsilon$;
\item[(iii)] since $F(l,4l,4l)\le4F(l,2l,2l)$ and $F(2l,4l,4l)=2F(2l,4l,2l)-\mathsf b=4F(2l,2l,2l)-2\mathsf a-\mathsf b$, we have $\mathsf c:=\mathsf d^{(1)}(l,4l,4l)\le4\epsilon+2\mathsf a+\mathsf b\le4.24112\,\epsilon\le0.29264<\log2$;
\item[(iv)] $\mathcal E(2l)=\mathsf d^{(1)}(2l,4l,4l)\le\psi(\mathsf c)\le2\cdot0.59869^2\cdot4.24112^2\,\epsilon^2\le12.8943\,\epsilon^2$.
\end{enumerate}
So $\mathcal E(2l)\le12.9\,\mathcal E(l)^2$ when $\mathcal E(l)\le0.069$. As $12.9\cdot0.069=0.8901<1$, induction shows that $\epsilon_k=\mathcal E(2^kl)$ does not increase in $k$ and that $12.9\,\epsilon_{k+1}\le(12.9\,\epsilon_k)^2$; so $12.9\,\epsilon_k\le(12.9\,\epsilon_0)^{2^k}\le\theta^{2^k}$, that is, $\mathcal E(2^kl)\le x_k:=\theta^{2^k}/12.9\le x_0\le0.069$.

\emph{Cubes.} Put $D(N)=8F(N^3)-F((2N)^3)=8\Delta(N^3)-\Delta((2N)^3)$. At $n=2^kl$, with $\mathsf a$ and $\mathsf b$ formed at $n$ in the place of $l$, the three brackets of
\[
\begin{multlined}
D(2n)=\bigl[8F((2n)^3)-4F(2n,4n,2n)\bigr]+\bigl[4F(2n,4n,2n)-2F(2n,4n,4n)\bigr]+\bigl[2F(2n,4n, \\
4n)-F((4n)^3)\bigr]
\end{multlined}
\]
are $4\mathsf a$, $2\mathsf b$ and $\mathcal E(2n)$; so $D(2n)\le(2.14928+4.83940+12.8943)\,\epsilon_k^2\le20\,\epsilon_k^2$. By (a) the density $F/|\underline n|$ decreases to $P$ along the cubes of the sides $2^sM$, $s\to\infty$, so that $\Delta(M^3)=\sum_{s\ge0}8^{-(s+1)}D(2^sM)$. At $M=2^jl$, $j\ge1$, every $D(2^sM)$ is at most $20\,\epsilon_{s+j-1}^2\le20\,\epsilon_{j-1}^2$, and
\[
\Delta(M^3)\le\tfrac{20}7\,\epsilon_{j-1}^2\le\tfrac{20}{7\cdot12.9^2}\,\theta^{2^j}\le0.017170\,\theta^{M/l}.
\]
In particular $\Delta((2n)^3)\le\frac{20}7\,\mathcal E(n)^2\le\frac{20}7x_k^2$ at $n=2^kl$. With $\mathcal E(N)\le8\Delta(N^3)$ from (a) and $\frac{20}7\cdot64<182.9$ this gives $\Delta((2N)^3)\le183\,\Delta(N^3)^2$ when $\Delta(N^3)\le0.008625$, for then $\mathcal E(N)\le0.069$; if moreover $\Delta(N^3)>0$, the hypothesis of (b) holds at $l=N$ with $x_0=8\Delta(N^3)$, so with $\theta=103.2\,\Delta(N^3)$, as $12.9\cdot8=103.2$. For $N'\ge2l$ let $j\ge1$ with $M=2^jl\le N'<2M$; by (a), $\Delta(N'^3)\le(N'/M)^3\Delta(M^3)<8\cdot0.017170\,\theta^{M/l}\le0.13736\,\theta^{N'/2l}$, since $M/l>N'/2l$ and $0<\theta<1$.

\emph{All sides $l'\ge l$.} At $n=2^kl$, the identity $\mathcal E(n)=2\Delta(n,2n,2n)-\Delta((2n)^3)$ and the two bounds $\mathcal E(n)\le x_k$, $\Delta((2n)^3)\le\frac{20}7x_k^2$ give $\Delta(n,2n,2n)\le\frac{x_k}2+\frac{10}7x_k^2$. Let $2^kl\le l'<2^{k+1}l$. As $\Delta\ge0$, and by (a) applied to the boxes $(2^kl,2^{k+1}l,2^{k+1}l)\le(l',2l',2l')$, whose volumes have the ratio $(l'/2^kl)^3<8$,
\[
\begin{multlined}
\mathcal E(l')\le2\Delta(l',2l',2l')\le16\,\Delta(2^kl,2^{k+1}l,2^{k+1}l)\le8x_k+\tfrac{160}7x_k^2\le\Bigl(\frac8{12.9}+\frac{160\cdot0.8901}{7\cdot12.9^2}\Bigr)\theta^{2^k} \\
\le0.743\,\theta^{2^k},
\end{multlined}
\]
using $x_k^2=\theta^{2^{k+1}}/12.9^2\le0.8901\,\theta^{2^k}/12.9^2$; and $\theta^{2^k}\le\theta^{l'/2l}$, since $2^k>l'/2l$.
(c) Let $L$ be an odd integer $\ge13$, an allowed blocking factor, and $K\in\mathbb N$. Wilson's action reads the lattice only as a graph, so the dilation $x\mapsto(a_K/a)x$ carries the periodic boxes and the periodic Gibbs states of $a\mathbb Z^3$ at $\beta$ onto those of $\mathbb L_K$, with the same partition functions, and we take $a=a_K$. The cube with $(2L^j)^3$ sites, $j>K$, is then the torus $T_\ell$, $\ell=2L^{j-K}$, at the spacing $a_K$; we put $p_\ell=\ell^{-3}\log Z_{K,\ell,\beta}$. The domination theorem (Proposition~\YMMref{91.9}{705}) bounds the number $n$ of pairwise mutually singular periodic Gibbs states by $\liminf_je^{\Delta((2L^j)^3)}$. Indeed, let $\Psi$ be a function of the cell $[-\frac\ell4,\frac\ell4]^3$ of the board of $\mathbb L_K$ with the side $\frac\ell2$, half the side of $T_\ell$, and the offset $0$, and let $\nu$ be a periodic Gibbs state. Proposition~\YMMref{91.9}{705} gives $|\nu(\Psi)|\le\liminf_{\ell'\to\infty}\|\Psi\|_{{\rm ch},\ell'}$, and the comparison of two sides, Corollary~\ref{IV-18-18}(e), gives $\|\Psi\|_{{\rm ch},\ell'}\le\|\Psi\|_{{\rm ch},\ell}\,e^{(\ell/2)^3(p_\ell-p_{\ell'})}$ for $\ell'\ge\ell$, where $(\ell/2)^3(p_\ell-p_{\ell'})\to\frac18\Delta$, $\Delta$ the excess of the cube with $\ell/a_K$ sites a side, as $p_{\ell'}\to a_K^{-3}P$; so
\[
|\nu(\Psi)|\le e^{\Delta/8}\,\|\Psi\|_{{\rm ch},\ell},
\]
and $\|\Psi\|_{{\rm ch},\ell}^8$ is the mean, in the torus measure $\tilde\mu$ of $T_\ell$, of the product of the eight mirror copies of $\Psi$ (equation~\eqref{eq:period-box-seminorm}). Mutually singular laws $\nu_1,\dots,\nu_n$ have pairwise disjoint carriers, which local events approximate: for $\epsilon>0$ there are pairwise disjoint local events $A_i$ with $\nu_i(A_i)\ge1-n\epsilon$, and for large $\ell$ they are events of the bonds of the cell. At their indicators the last bound reads $\nu_i(A_i)\le e^{\Delta/8}\tilde\mu(E_i)^{1/8}$, where $E_i$ is the intersection of the eight mirror images of $A_i$; the $E_i$ are pairwise disjoint, so $n(1-n\epsilon)^8\le e^{\Delta}$. By (b) the right member tends to $1$ as $\ell\to\infty$; as $\epsilon$ is arbitrary, $n\le1$. Every periodic Gibbs state is an integral of ergodic periodic Gibbs states, and two different ergodic ones are mutually singular; so there is one, it is extreme, and $J(\beta)=0$ by Lemma~\ref{lh:lat}(b).
(d) The first statement is (c); the second follows from (a), as $8\cdot0.008625=0.069$; the third from the bound of (c) with $n=2$, at every odd $L\ge13$, two different ergodic periodic Gibbs states being mutually singular.
(e) is (b) and (c).
(f) Let $G$ be a product of centred $\mathcal O_K(f_i)$, $i\in I\subset\{1,\dots,n\}$, centred at the means on $T_{\ell''}$. Its chessboard seminorm $\|G\|_{{\rm ch},\ell''}$, on the board of $T_{\ell''}$ whose cell holds the supports, has as its eighth power a centred moment, on the fixed torus $T_{\ell''}$, of $8|I|$ observables with pairwise disjoint supports, the mirror copies of the $\mathcal O_K(f_i)$; it is a polynomial in cumulants of orders $\ge2$, each of which converges as $K\to\infty$ by Theorem~\textup{\YMMref{AS}{813}}; so it is bounded in $K$. Let $\mu'$ be the measure of $T_{\ell'}$ at the spacing $a_K$, read on $\mathbb L_K$, $Q'=e^{\Delta'/8}$ with $\Delta'$ the excess of the cube with $\ell'/a_K$ sites a side, $Q''=e^{\Delta''/8}$ likewise at $\ell''$, $p_\ell=\ell^{-3}\log Z_{K,\ell,L^Kb}$ as in the proof of (c), and $\mathsf g=Q''\|G\|_{{\rm ch},\ell''}$. Then
\begin{equation}\label{lh:cmp}
|\nu_K(G)-\mu'(G)|\le10.7\,[Q'(Q'-1)]^{1/2}\,\mathsf g\qquad\text{if }\Delta'\le1.81 .
\end{equation}
Four facts enter. (i) For a real function $\Psi$ of the cell $[-\frac{\ell'}4,\frac{\ell'}4]^3$ of the board of $\mathbb L_K$ with the side $\frac{\ell'}2$ and the offset $0$, the coarse board, $|\nu_K(\Psi)|\le Q'\|\Psi\|_{{\rm ch},\ell'}$, by Proposition~\YMMref{91.9}{705} and Corollary~\ref{IV-18-18}(e), as in the proof of (c); $G$ is such a function (the plaquette-support argument in Section~\ref{sec:periodic-chessboard}). (ii) The board with the side $\frac{\ell''}2$ and the offset $0$, the fine board, has the cell $[-\frac{\ell''}4,\frac{\ell''}4]^3$, which holds the supports. As $\ell'/\ell''=L^{m'-m''}$ is odd, the coarse board is coarser than the fine one, and each of its mirror maps carries the cells of the fine board onto cells of the fine board, with the same parity (the proof of Lemma~\YMMref{104.1}{759}).

So the mirror copies $G^{[k]}$ of $G$ in the eight cells $k$ of the coarse board of $T_{\ell'}$ are mirror copies of $G$ in eight different cells of a period box of the fine board, and the chessboard estimate on the fine board of $T_{\ell'}$ (Corollary~\ref{IV-18-18}(b), (c), which rests on Proposition~\YMMref{86.3}{681}) gives $|\mu'(\prod_{k\in S}G^{[k]})|\le\|G\|_{\rm f}^{\,|S|}$ for every set $S$ of these cells, $\|G\|_{\rm f}$ being the seminorm of $G$ at the side $\ell'$ on the fine board. (iii) By Corollary~\ref{IV-18-18}(e) on the fine board, at the sides $\ell'\ge\ell''$, and as (a) gives $\Delta\ge0$ at the cube with $\ell'/a_K$ sites a side, that is, $p_{\ell'}\ge a_K^{-3}P$: $\|G\|_{\rm f}\le\|G\|_{{\rm ch},\ell''}\,e^{(\ell''/2)^3(p_{\ell''}-p_{\ell'})}\le\mathsf g$. (iv) $\mu'$ is invariant under the isometries of $T_{\ell'}$, so $\mu'(G^{[k]})=\mu'(G)$, and $|\mu'(G)|\le\|G\|_{\rm f}\le\mathsf g$.

Now let $t$ be real, $\Psi=1+tG$, $x=|t|\mathsf g$ and $R(x)=(1+x)^8-1-8x$. By equation~\eqref{eq:period-box-seminorm} on the coarse board, (ii) and (iv),
\[
\|\Psi\|_{{\rm ch},\ell'}^8=\mu'\Bigl(\prod_k\bigl(1+tG^{[k]}\bigr)\Bigr)=1+8t\mu'(G)+\sum_{|S|\ge2}t^{|S|}\mu'\Bigl(\prod_{k\in S}G^{[k]}\Bigr)\le1+8t\mu'(G)+R(x),
\]
so $\|\Psi\|_{{\rm ch},\ell'}\le1+t\mu'(G)+\frac18R(x)$, as $(1+y)^{1/8}\le1+\frac y8$ for $y\ge-1$; and $1+t\nu_K(G)=\nu_K(\Psi)\le Q'\|\Psi\|_{{\rm ch},\ell'}$ by (i). Hence $t(\nu_K(G)-\mu'(G))\le(Q'-1)(1+x)+\frac18Q'R(x)$ for both signs of $t$, with $Q'\ge1$ by (a). At $\mathsf g=0$ this gives $\nu_K(G)=\mu'(G)$; at $\mathsf g>0$
\[
|\nu_K(G)-\mu'(G)|\le\mathsf g\Bigl[\frac{(Q'-1)(1+x)}x+\frac{Q'R(x)}{8x}\Bigr]\qquad\text{for every }x>0 .
\]
For $Q'>1$ and $x^2=(Q'-1)/Q'$ the bracket is $[Q'(Q'-1)]^{1/2}\bigl(1+x+R(x)/8x^2\bigr)$, and $1+x+R(x)/8x^2\le10.7$ for $x\le0.45$, that is, for $\Delta'\le1.81$; for $Q'=1$ let $x\to0$. This is \eqref{lh:cmp}. By (b) at $l=l_K=\ell_*/2a_K$, $\Delta'<0.13736\,\theta^{\ell'/\ell_*}$ and $\Delta''<0.13736$; so $Q''<1.02$, $[Q'(Q'-1)]^{1/2}\le0.14\,\theta^{\ell'/2\ell_*}$, and \eqref{lh:cmp} gives $|\nu_K(G)-\mu'(G)|\le2\,\theta^{\ell'/2\ell_*}\|G\|_{{\rm ch},\ell''}$. A cumulant of order $n\ge2$ is one polynomial in the means $\mathfrak s(G_I)$, $I\subset\{1,\dots,n\}$, for both laws $\mathfrak s\in\{\nu_K,\mu'\}$; by (iv) and the last bound these means are bounded in $K$ and $\ell'$, and as $\theta^{\ell'/2\ell_*}\le\theta^{\ell'/4\ell_*}$ this gives the display.
\end{proof}

\subsection{Proof of Theorem~\ref{lh:main}}

\emph{Step 1: (A)(1) and (A)(2).} Take two settings at one radius with $L_2=L_1+2$ (Lemma~\ref{lh:two}) and $r_1>0$ with $13r_1b_{+,i}\le1$. We enter in Proposition~\ref{lh:per}: $\beta_0=\max_iL_ib_{-,i}$; $C_{\rm low}=871C^{\rm ref}_1$; $\lambda_i=\log L_i$, so $\lambda_1/\lambda_2\notin\mathbb Q$ by Lemma~\ref{lh:two}; $\alpha_i=\log b_{-,i}$ and $\varkappa_i=\frac12[\log(b_{+,i}/b_{-,i})-\lambda_i]>0$, so that $\mathsf J_i=\log\mathsf I_i$; $k_i$ the least integer $\ge1$ with $L_i^{k_i}b_{-,i}>\beta_0$; $m^0_i=m_{0,i}$; $P^{(i)}_{m,K}=P_N$ and $n^{(i)}_{m,K}=n_{\eta_K}(\mu_{N,\beta})$ with $N=2L_i^{m+K}$ and $\eta_K=L_i^{K-1}$. By Proposition~\ref{lh:blk}(b)--(d), \eqref{lh:data} holds with $\epsilon^{(i)}_m=c_{{\rm P},i}L_i^{-m}$ and $C_N=\max_iC^{\rm ref}_i$, \eqref{lh:low} holds on $(\beta_0,\infty)$, and by \eqref{lh:far} the function $y^{(i)}_{m,K}$ of Proposition~\ref{lh:per} is the far part of the setting $(\mathrm S_i)$. By the convergence on each torus (Theorem~\textup{\YMMref{AS}{813}} at $\varsigma_i$; the supports of $\mathsf f^{\langle\ell\rangle}$ and $1-\mathsf h^{\langle\ell\rangle}$ have no common point), for each $m$ the far part converges as $K\to\infty$, uniformly on $\mathsf J_i$; so there are non-decreasing integers $K^{(i)}(m)\ge k_i$, the maxima over $m_{0,i}\le m'\le m$ of the thresholds at the single sides, with
\begin{equation}\label{lh:one}
\sup_{\mathsf J_i}|y^{(i)}_{m,K}-y^{(i)}_{m,K'}|\le1\qquad(K,K'\ge K^{(i)}(m)).
\end{equation}
At $m^0=m_{0,1}$ and $k^0=K^{(1)}(m^0)$ this gives $Y^0=\sup_{\mathsf J_1}|y^{(1)}_{m^0,k^0}|+1<\infty$, the far part at one cutoff being bounded on $\mathsf J_1$. Proposition~\ref{lh:per}(a) with $\omega_i=1$, large $m_i$ and $\bar K_i=K^{(i)}(m_i)$ gives, with $\mathsf A_0$ independent of $m_1,m_2$ and for $0\le\delta\le1$,
\[
\begin{multlined}
\limsup_{s\to\infty}e^{4s}J(e^s)\le(\mathsf A_0+134)\max_i(c_{{\rm P},i}L_i^{-m_i})^{1/2},\quad\limsup_{s\to\infty}|\mathbb M|([s,s+\delta])\le(2\mathsf A_0+C_{\rm low} \\
+134)\delta+(\mathsf A_0+134)\max_i(c_{{\rm P},i}L_i^{-m_i})^{1/2}.
\end{multlined}
\]
Letting $m_1,m_2\to\infty$ gives (A)(1), and (A)(2) with $C=2\mathsf A_0+C_{\rm low}+134$, first for $\delta\le1$ and then for all $\delta>0$ by cutting $[s,s+\delta]$ into intervals of length at most $1$. The equivalent form of (A)(1) is Lemma~\ref{lh:lat}(b).

\emph{Step 2: (A)(3).} Let the two settings have one value of $\mathfrak p_0$; such pairs exist with $\mathfrak p_0\ge a$ by Lemma~\ref{lh:two}. The pair $(\mathsf f^{\langle\ell\rangle},1-\mathsf h^{\langle\ell\rangle})$ has the spread $\varepsilon$ and norms $\|\cdot\|_\sharp\le40\sqrt3ML$ (values in $[0,1]$, Lipschitz constants at most $\sqrt3L$), and $b^5\varepsilon^{-3}\le b_+^5L^3$. By the rate of the convergence on each torus and Lemma~\ref{lh:thr}, in each setting and for the sides with $M\mathcal R(b_+)\le L^m$,
\[
\sup_{\log\mathsf I}|y_{m,K}-y_{m,K'}|\le2b_+^5L^3(40\sqrt3ML)^2C'\ell^9K^{\mathsf p_2}L^{-K/10}\le1\qquad\text{for }K,K'\ge c\bigl(1+\ell^{3/2\mathfrak p_0}\bigr),
\]
with the rate constant $C'=\mathsf C'_2(\mathsf I)$ of Theorem~\textup{\YMMref{AS}{813}} at $n=2$ and $s=0$, which is admitted as $b_-\ge\gamma^{-2}$ and $M\mathcal R(b_+)\le L^m$, and a number $c$ independent of $m$, because $\log\ell\le\frac{2\mathfrak p_0}3\ell^{3/2\mathfrak p_0}$. We may therefore take $K^{(i)}(m)\le c(1+(2L_i^m)^{3/2\mathfrak p_0})$ in Step~1. At a large $s$ let $m_i$ be the largest integer with $\alpha_i+2\varkappa_i+\lambda_iK^{(i)}(m_i)\le s$; then $s<c'(1+L_i^{3(m_i+1)/2\mathfrak p_0})$, so $L_i^{-m_i}\le c''s^{-2\mathfrak p_0/3}$, and Proposition~\ref{lh:per}(a) gives $e^{4s}J(e^s)\le Cs^{-\mathfrak p_0/3}$ and $|\mathbb M|([s,s+\delta])\le(2\mathsf A_0+C_{\rm low}+134)\delta+Cs^{-\mathfrak p_0/3}$ for large $s$; the second is at most $C_{\rm A}\delta$ when $C_{\rm A}\ge\max\{2\mathsf A_0+C_{\rm low}+135,C\}$ and $C_{\rm A}s^{-\mathfrak p_0/3}\le\delta$. On bounded sets of $\beta\ge3$ use $J\le6$. For a real $q_0$ choose $\mathfrak p_0\ge3q_0$; as $\log\beta\ge1$ for $\beta\ge3$, $(\log\beta)^{-\mathfrak p_0/3}\le(\log\beta)^{-q_0}$. In any setting, for $(K,\beta)\in\mathfrak d(g)$ we have $b=a_K\beta\ge g^{-2}\ge4$, so $\log\beta\ge K\lambda$ and $a_K^{-4}J(\beta)=b^{-4}\beta^4J(\beta)\le g^8C(K\lambda)^{-q_0}$.

\emph{Step 3: (A)(4).} In one of the two settings let $b\in\mathsf I$, $a=a_K$, $\beta=L^Kb$. By Proposition~\ref{lh:blk}(c) at $m=m_0$, $0\le a^{-3}[P_{N_0}(\beta)-P(\beta)]\le3C_{\rm v}$, which is $-3C_{\rm v}\le u_K\le0$. $u_K$ is the difference of a convex and a smooth function of $b$; $u_K'(b\pm)=a^{-4}[P'(\beta\pm)-P_{N_0}'(\beta)]$, and as measures on $(b_-,b_+)$, by \eqref{lh:far} at $m=m_0$,
\[
u_K''(db)=b^{-4}\mathbb M_K(db)+\vartheta_K(b)\,db,\qquad\vartheta_K=-b^{-5}\bigl[\beta^5\bigl(n_{\eta_K}(\mu_{N_0,\beta})-\hat n(\beta)\bigr)+y_{m_0,K}(\log b)\bigr],
\]
where $\mathbb M_K$ is the image of $\mathbb M$ under $s\mapsto e^{s-K\lambda}$. The far parts at $m_0$ are bounded at each cutoff and converge by \eqref{lh:one}, so $Y^*_0=\sup_{K\ge1}\sup_{\log\mathsf I}|y_{m_0,K}|<\infty$, and by Proposition~\ref{lh:blk}(d), $|\vartheta_K|\le b_-^{-5}(C^{\rm ref}+Y_0^*)$. By the last assertion of Proposition~\ref{lh:blk}(d), in the setting at hand, $P''(d\beta)\ge[\hat n(\beta)-871C^{\rm ref}\beta^{-5}]d\beta$ on $(Lb_-,\infty)$, which contains $L^K(b_-,b_+)$ for every $K\ge1$; so $\mathbb M\ge-871C^{\rm ref}ds$ on $(\log(Lb_-),\infty)$, and, as $ds=db/b$, $b^{-4}\mathbb M_K(db)\ge-871C^{\rm ref}b_-^{-5}db$ on $(b_-,b_+)$. So $u_K''\ge-C_{\rm sc}\,db$ with $C_{\rm sc}=b_-^{-5}(871C^{\rm ref}+C^{\rm ref}+Y_0^*)$, for every $K\ge1$, the numbers $C^{\rm ref}$, $b_-$ and $Y_0^*$ being those of the setting at hand.

As $\log(b+h)-\log b\le h/b_-$, (A)(2) gives $\limsup_K\sup\{|u_K''|([b,b+h]):b_-<b,\ b+h<b_+\}\le C_uh$ for every $h>0$, with $C_u=b_-^{-5}(C+C^{\rm ref}+Y_0^*)$, the upper limit of (A)(2) being over the real $s\to\infty$ and so uniform in $b$. Let $I'=[\alpha,\alpha']\subset(b_-,b_+)$. By the last bound and a covering, $\sup_{K\ge K_1}|u_K''|(I')<\infty$ for some $K_1$; the mean of $u_K'$ over $I'$ is $(u_K(\alpha')-u_K(\alpha))/|I'|$, of modulus at most $3C_{\rm v}/|I'|$, and $u_K'(\cdot\pm)$ differs from it by at most $|u_K''|(I')$; so the $u_K$, $K\ge K_1$, are bounded and have one Lipschitz constant on $I'$. The theorem of Arzel\`a and Ascoli and a diagonal sequence give the subsequence.

For $b,b+h\in I'$, $|u_K'(b\pm)-h^{-1}(u_K(b+h)-u_K(b))|\le|u_K''|([b,b+h])$; so along the subsequence $\limsup_{K,K'}\sup_{I'}|u_K'(\cdot\pm)-u_{K'}'(\cdot\pm)|\le2C_uh$ for every $h>0$: the one-sided derivatives converge uniformly on $I'$ to a function $w$, $u(b')-u(b)=\int_b^{b'}w$, and $|w(b')-w(b)|\le\limsup_K|u_K''|([b,b'])\le C_u|b'-b|$; so $u'=w$. The other properties pass to uniform limits. Last, by Lemma~\ref{lh:lat}(b) on $\mathbb L_K$, $\nu(\mathcal O_K(f))=3a^{-4}e_\nu r_K(f)$ with $-3e_\nu\in[P'(\beta-),P'(\beta+)]$; the torus measure $\tilde\mu=\tilde\mu_{K,\ell_0,\beta}$ is invariant under the isometries of $\mathbb L_K$, so $\tilde\mu(\mathcal O_K(f))=3a^{-4}e_*r_K(f)$ with $-3e_*=P_{N_0}'(\beta)$; subtract.

\emph{Step 4: (E).} (i) Let first $b'\ge\max\{\gamma^{-2},e^{\mathsf Y^\natural}\}$ and $M\mathcal R(b'')\le L^m$. Discarding (Proposition~\textup{\YMMref{115.2}{806}}, its lost-mass and cumulant bounds) holds with any number $\epsilon_{\rm R}$ that bounds the lost mass, $\sum_{h:Z_t(h)\ne\emptyset}\mathsf t(h)\le\epsilon_{\rm R}Z$, in the place of the bound on the lost mass that its statement takes from the definitions preceding Proposition~\YMMref{115.2}{806}: if $\epsilon_{\rm R}\le\frac12$, then the generating function $T$ of the kept histories is entire, $T(0)\ge\frac12Z$, and
\[
|S^T_{n;K,\ell,\beta}(\underline f)-\partial_z^n\log T(0)|\le4\,\mathsf{op}_n\epsilon_{\rm R}\prod_i6\ell^3L^{4K}\|f_i\|_\infty .
\]
Indeed, with the bare density $\rho_0$ and the weight $\mathfrak w^{\rm sf}_t\ge0$ of the definitions preceding Proposition~\YMMref{115.2}{806}, $\mathsf\Lambda(F)=\int F\rho_0\,dU$ and $\mathsf\Lambda'(F)=\int F\,\mathfrak w^{\rm sf}_t\rho_0\,dU$ are positive functionals on the gauge-invariant bounded Borel functions, $\mathsf\Lambda-\mathsf\Lambda'$ is positive by Proposition~\textup{\YMMref{115.2}{806}}, its positive-domination inequality, $\mathsf\Lambda(1)=Z$, and $T(z)=\mathsf\Lambda'(e^{\sum_iz_i\mathcal O_K(f_i)})$. The proof of equation~\YMMref{(2190)}{806} first bounds the lost mass by the sum of the histories with $Z_t(h)\ne\emptyset$, before estimating that sum. It therefore gives $\mathsf\Lambda(1)-\mathsf\Lambda'(1)\le\epsilon_{\rm R}\mathsf\Lambda(1)$; and the statement on two nearby positive functionals and their cumulants (Lemma~\ref{lem:positive-functional-comparison}, with $|\mathcal O_K(f_i)|\le6\ell^3L^{4K}\|f_i\|_\infty$) gives $T(0)=\mathsf\Lambda'(1)\ge\frac12Z$ and the display, the cumulants of $\mathsf\Lambda$ and of $\mathsf\Lambda'$ being $S^T_{n;K,\ell,\beta}(\underline f)$ and $\partial^n_z\log T(0)$.

These are the last two steps of the proof of Proposition~\textup{\YMMref{115.2}{806}}, made with $\epsilon_{\rm R}$ in the place of that bound. The proof of Theorem~\YMMref{AR}{810} (the comparison of the cutoffs $K$ and $K+1$, with its three lemmas) uses of the bound $\epsilon_{\rm R}$ only these two facts, that it bounds the lost mass and that it is at most $\frac12$. Its error terms are of three kinds: terms exponentially small in $K$ times $N_{[d]}=\ell^3L^{3d}$ at depths $d\le K/9+1$; the term $8\,\mathsf{op}_n\epsilon_{\rm R}\prod_i6\ell^3L^{4(K+1)}\|f_i\|_\infty$; and, for $n\ge2$ and spread $\mathsf d$, a contact term at most $c\,\ell^3L^{3u}e^{-\mathsf dL^u/M}$, with the depth $u=\lceil K/180(3n+4)\rceil$.

The bound on the lost mass with which the thresholds of Theorem~\textup{\YMMref{AS}{813}} are formed carries a factor $e^{{\rm const}\,\ell^3}$. We take instead $\epsilon_{\rm R}=\mathsf R_{t_1}$ of Lemma~\ref{lh:lost} at the level $t_1=K-u$: since $\log(L^ub)\ge u\lambda\ge K\lambda/180(3n+4)$ and $2\mathfrak p_0\ge2$, $\mathsf R_{t_1}\le c\,\ell^3L^{3u}e^{-c''K^{2\mathfrak p_0}}$, which beats the factor $\ell^3$ and the products of the crude bound. The most demanding of the remaining conditions is $\ell^3L^{-K/18}\le c$. Hence there is an integer $k_\flat$, free of $m$, $\ell$ and $K$, such that for $K\ge k_\flat+54m$ all conditions hold and
\[
|S_{K+1}-S_K|\le C_n\ell^{3(n+1)}K^{\mathsf p_n}L^{-K/10}\Pi .
\]
Summing over the cutoffs, with $C_n\sum_{k\ge K}k^{\mathsf p_n}L^{-k/10}\le C_n'K^{\mathsf p_n}L^{-K/10}$, which defines $C'_n$, gives (a) with $k_\flat$ and $C'_n$ in this range. For $K\le k_\flat+54m$ the crude bound and $L^{4n(k_\flat+54m)}\le L^{4nk_\flat}\ell^{216n}$ give $|S_K|\le\mathsf{op}_n6^nL^{4nk_\flat}\ell^{219n}\Pi$, and for larger $K$ (a) compares $S_K$ with $S_{k_\flat+54m}$; this gives (b) with $C^{\rm pol}_n=\mathsf{op}_n6^nL^{4nk_\flat}+C'_n\sup_kk^{\mathsf p_n}L^{-k/10}$, as $3(n+1)\le219n$. (ii) Let $\jmath\ge0$ be the least integer with $L^\jmath\ge\max\{\gamma^{-2}/b',e^{\mathsf Y^\natural}/b',b''\}$ and $L^{\jmath/2}\ge M(4r/e)^r$. For $K>\jmath$, Lemma~\YMMref{118.2}{813} gives, with $f^{(\jmath)}=f(L^{-\jmath}\cdot)$,
\[
S^T_{n;K,\ell,L^Kb}(\underline f)=L^{\jmath n}S^T_{n;K-\jmath,L^\jmath\ell,L^{K-\jmath}(L^\jmath b)}(\underline f^{(\jmath)}):
\]
both members are cumulants of the same variables on the same periodic cube with $\ell L^K$ sites a side, since $\mathcal O_K(f)=L^\jmath\mathcal O_{K-\jmath}(f^{(\jmath)})$ under the dilation by $L^\jmath$. The dilation keeps the supremum, divides the Lipschitz constant by $L^\jmath$ and multiplies the spread by $L^\jmath$, so $\|f^{(\jmath)}\|_\sharp\le\|f\|_\sharp$. The window $L^\jmath[b',b'']$ lies above both floors; and as $\mathcal R(x)<L(\log x)^r$, $\log(L^\jmath b'')\le2\jmath\lambda$ and $(2\jmath\lambda)^r\le(4r/e)^rL^{\jmath/2}$, the sides $2L^{m+\jmath}$, $m\ge1$, satisfy $M\mathcal R(L^\jmath b'')\le L^{m+\jmath}$. So (i) applies on the dilated data, with numbers $k_{\flat,\jmath}$, $C'_{n,\jmath}$, $C^{\rm pol}_{n,\jmath}$; as $n+3(n+1)+\frac1{10}\le4n+4$ and $K-\jmath\ge k_{\flat,\jmath}+54(m+\jmath)$ for $K\ge k_{\flat,\jmath}+55\jmath+54m$, this gives (a) with $k^{\rm rt}=k_{\flat,\jmath}+55\jmath$ and $C^{\rm rt}=L^{\jmath(4n+4)}C'_{n,\jmath}$, and $|S_K|\le L^{220n\jmath}C^{\rm pol}_{n,\jmath}\ell^{219n}\Pi$ for $K>\jmath$.

For $1\le K\le\jmath$ the crude bound gives $|S_K|\le\mathsf{op}_n6^nL^{4n\jmath}\ell^{3n}\Pi$; so (b) holds with $C^{\rm pol}=\max\{L^{220n\jmath}C^{\rm pol}_{n,\jmath},\mathsf{op}_n6^nL^{4n\jmath}\}$. (c): $A>30(n+1)\ge90>54$, so (a) applies for $K\ge k^{\rm rt}+Am$, and $m\le K/A$ gives $\ell^{3(n+1)}\le2^{3(n+1)}L^{3(n+1)K/A}$. (d): $K\ge C\log\ell$ with $C\ge54/\lambda+k^{\rm rt}/\log2$ gives $K\ge54m+k^{\rm rt}$, as $\log\ell\ge\log2$ and $\log\ell=\log2+m\lambda$; with $C\ge11$, $\frac1{10}-\frac1C\ge\frac1{110}$ and $K^{\mathsf p_2}L^{-K/110}\le(110\mathsf p_2/e\lambda)^{\mathsf p_2}$, so (a) gives the first bound, as $\ell^9\le\ell^{438}$; (b) gives the second.

\emph{Step 5: the far part beyond $k^{\rm rt}+Am$.} In a setting apply (E)(a) at $n=2$ to $\varsigma$ and a lowered choice $\varsigma'$, with $[b',b'']=\mathsf I$, $\mathsf d=\varepsilon$ and the pair of Step~2. With $(40\sqrt3ML)^2=4800M^2L^2$, for all $m\in\mathbb N$, $\upsilon\in\log\mathsf I$ and $K'\ge K\ge k^{\rm rt}+54m$,
\begin{equation}\label{lh:osc}
|y_{m,K'}(\upsilon)-y_{m,K}(\upsilon)|\le c_{\rm y}\ell^9K^{\mathsf p_2}L^{-K/10},\qquad c_{\rm y}=4800b_+^5M^2L^5C^{\rm rt}.
\end{equation}
For real $A>90$ and $K,K'\ge k^{\rm rt}+Am$ we have $\ell^9=2^9L^{9m}\le2^9L^{9K/A}$ and $\frac1{10}>\frac9A$; so
\[
\sup_{\log\mathsf I}|y_{m,K}-y_{m,K'}|\le\omega_A=2^9c_{\rm y}\sup_{k\ge1}k^{\mathsf p_2}L^{-(\frac1{10}-\frac9A)k}<\infty .
\]
This is the oscillation hypothesis of Proposition~\ref{lh:per}(b) with the slope $A$; and $\frac1{10}>\frac9A$ holds exactly for $A>90$. With it, $\sup_K\sup_{\log\mathsf I}|y_{m,K}|<\infty$ at each $m$, the finitely many cutoffs below the threshold being bounded one by one. (B) uses (E) here and nowhere else.

\emph{Step 6: (B).} Enter in Proposition~\ref{lh:per}(b) with the data of Step~1, $\mathsf a_i=k^{\rm rt}_i$, $A_i=A$ and $\omega_i=\omega_{A,i}$ of Step~5 at $(\varsigma_i,\varsigma'_i)$. At each single $(m,K)$, by \eqref{lh:far} and Lemma~\ref{lh:lat}(c), $|y_{m,K}|\le\beta^5(P_N''+|n_{\eta_K}(\mu_{N,\beta})|)\le(L^Kb_+)^5(9N^3+1125\eta_K^3)$. This gives (B)(1) with the constants
\begin{align}
&\begin{gathered}
k^0=\max\{k_1,\lceil k^{\rm rt}_1+Am_{0,1}\rceil\},\qquad\\
 Y^0\le(L_1^{k^0}b_{+,1})^5\bigl(72L_1^{3(m_{0,1}+k^0)}+1125L_1^{3(k^0-1)}\bigr)+\omega_{A,1},
\end{gathered}\label{lh:cst}\\
&\begin{gathered}
\mathsf M_1=(Y^0+C_N+C_{\rm low})\varkappa_1+c_{{\rm P},1}L_1^{-m_{0,1}}\bigl(15/\varkappa_1+14+12\varkappa_1\bigr),\qquad\\
\mathsf A_0=7\mathsf M_1/\varkappa_1+C_{\rm low}+4C_N+\omega_{A,1}+\omega_{A,2},
\end{gathered}\notag\\
&\begin{gathered}
c_{3,i}=|\log b_{-,i}|+2\varkappa_i+\lambda_i(k_i+k^{\rm rt}_i+1),\qquad\\
 c_6=\max_ic_{{\rm P},i}e^{\lambda_i+c_{3,i}/A},\qquad\\
 T_i(m)=\log b_{-,i}+2\varkappa_i+\lambda_i\max\{k_i,\lceil k^{\rm rt}_i+Am\rceil\},
\end{gathered}\notag\\
&\begin{gathered}
C_{\rm h}=2\mathsf A_0+C_{\rm low}+134,\qquad\\
 C_{\rm r}=(\mathsf A_0+134)c_6^{1/2},\qquad\\
\log\beta_1=\max\{0,T_1(m_{0,1}),T_2(m_{0,2}),A\log^+(4c_6/\varkappa^2)\},
\end{gathered}\notag
\end{align}
since $P$ is the pressure, $J_P=J$, and $|\mathbb M|([s,s+\delta])$ is the integral of (B)(1). (2): for $0<c<\frac1{180}$ take $A=\frac1{2c}>90$; then $\beta^{4+c}J(\beta)\le C_{\rm r}$ for $\beta\ge\beta_1$, and $\beta^{4+c}J\le6\max\{1,\beta_1\}^{4+c}$ on $[1,\beta_1]$; for $c\le0$, $\beta^{4+c}J\le\beta^{4+1/182}J$ on $[1,\infty)$. $c_*(c)$ is defined by $P$ alone. (3): $(\log\beta)^q\beta^{-1/182}$ is bounded on $[3,\infty)$. (4): enter in Proposition~\ref{lh:slope} with the data of Step~1. At $i=1$, Step~5 at $(\varsigma_1,\varsigma'_1)$ gives (i) with $\mathsf a_1=k^{\rm rt}_1$, $A_1=A$ and $\omega_1=\omega_{A,1}$, and $\sup_{K\ge k_1}\sup|y^{(1)}_{m_{0,1},K}|$ is finite, the cutoffs below the threshold being finitely many; at $i=2$, \eqref{lh:one} at $\varsigma_2$ gives (ii) with $\omega_2=1$ for every $m\ge m_{0,2}$, and $\epsilon^{(2)}_m=c_{{\rm P},2}L_2^{-m}\to0$.

Theorem~\YMMref{D}{8} is used at $\varsigma_1$ and $\varsigma'_1$ through (E), and at $\varsigma_2$ through the convergence on each torus. (5): by \eqref{lh:ident} with the data of Step~1 and $G_K=G(\cdot-K\lambda_i)$, which is supported in $\mathsf J_i+K\lambda_i\subset[s_3,\infty)$, $s_3=\log(L_i^Kb_{-,i})\ge\log\beta_1$ (and $K\ge k_i$, as $\beta_1>\beta_0$), $\int G\,y^{(i)}_{m,K}=\langle\mathbb M,G_K\rangle-r$ with $|r|\le c_{{\rm P},i}L_i^{-m}\mathsf Q(G)+C_N\|G\|_1$; by (1), the hypothesis of Proposition~\ref{lh:per}(d) holds with $s_2=\log\beta_1$ and $\epsilon_1(s)=e^{-s/2A}$, and it bounds $|\langle\mathbb M,G_K\rangle|$ by $(C_{\rm h}+C_{\rm r})(\|G\|_1+(L_i^Kb_{-,i})^{-1/2A}\|G'\|_1)$.

\emph{Step 7: (C).} By the dilation, $\rho^{(K)},\nu^{(K)}\in\mathcal G(\beta)$, and $\rho(a_p)=e_{\rho^{(K)}}$; Lemma~\ref{lh:lat}(b) gives $3|e_\rho-e_\nu|\le J(\beta)$, and invariance gives $\nu(\mathcal O_K(f))=3a_K^{-4}e_\nu r_K(f)$. For each of the three orientations the barycentres form a translate of $a_K\mathbb Z^3$, of which at most $(2R/a_K+1)^3$ points lie in $[-R,R]^3$; so $|r_K(f)|\le(2R+a_K)^3\|f\|_\infty$. On $\mathfrak d(g)$, $b=a_K\beta\ge g^{-2}\ge4$, so $\beta\ge1$, and, as $4+c>0$,
\[
a_K^{-4}J(\beta)=b^{-4}\beta^4J(\beta)\le c_*(c)\,b^{-4}\beta^{-c}=c_*(c)\,b^{-4-c}L^{-cK}\le g^{8+2c}c_*(c)L^{-cK}.
\]
The setting is used only to form $\mathfrak d(g)$; $c_*(c)$ is a number of Wilson's action.

\emph{Step 8: (D).} (1) On $\mathfrak d(g)$, $b\ge g^{-2}$ and $\beta=L^Kb\ge L^Kg^{-2}$, so $a_K^{-4}J(\beta)=b^{-4}\beta^4J(\beta)\le g^8\sup_{\beta'\ge L^Kg^{-2}}\beta'^4J(\beta')$, which tends to $0$ by (A)(1); and by (C) with $a_K\le1$ the difference of the one-point functions is at most $a_K^{-4}J(\beta)(2R+1)^3\|f\|_\infty$. (2) At $X=(f)$ a limit in the frame is $\lim_j[\rho_j(\mathcal O_{K_j}(f))-\nu_j(\mathcal O_{K_j}(f))]=0$. (3) Each component is a limit in the frame along the full sequence, with the approximants $\rho_j=\nu_j(\cdot|B_j)$ (Theorem~\textup{\YMMref{AP}{780}}); by (2), $0=\sigma_\xi((f))=\mathfrak c_\xi\int f$; take $\int f=1$. As the one-point values of the component vanish, \eqref{lh:rece} with $\rho=\nu'=\rho_j$ and $\nu=\nu_j$ gives $\rho_j(X^{\rho_j}_{K_j})-\rho_j(X^{\nu_j}_{K_j})\to0$, every term with $A\ne\emptyset$ having a factor that tends to $0$ and a factor that converges; so the Gibbs sequence $(K_j,\beta_j,\rho_j)$ converges, with the limit $\sigma_\xi$.

The component clusters with the numbers $C^g$, $r$ (Theorem~\textup{\YMMref{AP}{780}}); so for its own canonical decomposition $(\mathsf X',\varpi',(\sigma'_\zeta))$, (Can4) at $q=2$ gives $\int\sigma'_\zeta(X)^2d\varpi'=\lim_{|v|\to\infty}\sigma_\xi(X\cdot\tau_vX)=\sigma_\xi(X)^2=(\int\sigma'_\zeta(X)d\varpi')^2$; so $\zeta\mapsto\sigma'_\zeta(X)$ is constant on the support of $\varpi'$, which is $\mathsf X'$, and the injectivity in (Can3) makes $\mathsf X'$ a point. This argument used of (A) only that $\mathfrak c_\xi=0$; conversely, a limit of a convergent Gibbs sequence $(K_j,\beta_j,\rho'_j)$ has the one-point values $\lim_j\rho'_j((f)^{\rho'_j}_{K_j})=0$, so $\mathfrak c_\xi=0$; this is the equivalence. (4) The one-point values vanish; with (Can4) at $q=1$ and the clustering of the components,
\[
S^\sigma_2(f,\tau_vh)=\sigma((f)\cdot\tau_v(h))=\int\sigma_\xi((f)\cdot\tau_v(h))\,d\varpi,\qquad|\sigma_\xi((f)\cdot\tau_v(h))|\le C^g((f),(h))/|v| .
\]

\emph{Step 9: (F).} (1) By Proposition~\YMMref{106.3}{765}, $\mathcal G_{K,\beta}$ satisfies the hypotheses of Lemma~\ref{lh:cond}, and by Lemma~\ref{lh:lat}(b) on $\mathbb L_K$, $-3e_+=P'(\beta-)$, $-3e_-=P'(\beta+)$ and $e_+-e_-=\frac13J(\beta)$. (2) The clustering inequality of the plateau theorem at $X=(f)$, $Y=(h)$ compares $\sigma_\nu((f)\cdot\tau_v(h))=S^T_2[\nu](f,\tau_vh)$ with $\nu(E_\nu[(f)^\nu_K|\mathcal I]\,E_\nu[(h)^\nu_K|\mathcal I])$, which is $9a_K^{-8}r_K(f)r_K(h)\operatorname{Var}_\nu\mathsf e$ by Lemma~\ref{lh:cond}(b), since $\nu(\mathsf e)=e_\nu$; and $\operatorname{Var}_\nu\mathsf e\le\frac14(J/3)^2$; the last bound is (C). (3) $\nu_\pm$ exist by Lemma~\ref{lh:lat}(b) on $\mathbb L_K$. By Lemma~\ref{lh:cond}(d), $\mathsf e$ is almost surely constant in $\nu_\pm$, in every $\nu$ with $-3e_\nu$ at an end of the interval, and in a pure phase, which is ergodic. For $\nu=\frac12(\nu_++\nu_-)$, which lies in the convex set $\mathcal G_{K,\beta}$, and bounded $F,F'$,
\[
\operatorname{Cov}_\nu(F,F')=\tfrac12\operatorname{Cov}_{\nu_+}(F,F')+\tfrac12\operatorname{Cov}_{\nu_-}(F,F')+\tfrac14\bigl(\nu_+(F)-\nu_-(F)\bigr)\bigl(\nu_+(F')-\nu_-(F')\bigr).
\]
Take $F=\mathcal O_K(f)$ and $F'=\mathcal O_K(\tau_vf)$, $v\in a_K\mathbb Z^3$, with $0\le f\in C^1_c((-\frac12,\frac12)^3)$ and $r_K(f)>0$, for which (2) applies at $\bar R=1$. Then $\nu_+(F)-\nu_-(F)=3a_K^{-4}r_K(f)(e_{\nu_+}-e_{\nu_-})=-a_K^{-4}J(\beta)r_K(f)$, the same for $F'$, as $r_K(\tau_vf)=r_K(f)$; (2) in the three states, with $|v|\to\infty$, gives $9a_K^{-8}r_K(f)^2\operatorname{Var}_\nu\mathsf e=\frac14(a_K^{-4}J(\beta)r_K(f))^2$. (4) Proposition~\ref{lh:many} with $\mathcal G=\mathcal G_{K,\beta}$; $X^\nu_K$ is a real continuous local function, a polynomial in finitely many $a_p$, and by Lemma~\ref{lh:lat}(b) every $\rho\in\mathcal G_{K,\beta}$ has $D^-\bar p^F(0)\le\rho(F)\le D^+\bar p^F(0)$ for such $F$; at $X=(f)$, $3a_K^{-4}|r_K(f)|(e_+-e_-)=a_K^{-4}J(\beta)|r_K(f)|$.

\emph{Step 10: (G) in a window.} Let $g^{-2}+\bar B<b<X_g-\bar B$ for some $0<g\le\gamma$. From some index on $a_{K_j}\beta_j\in I(g)$, so $(K_j,\beta_j)\in\mathfrak d(g)$; we drop the first members, which changes no set of limits. Let $\mathsf S_j=\{\sigma_\nu:\nu\in\mathcal G_{K_j,\beta_j}\}\subset\mathcal E^g$ and fix a metric of $\mathcal E^g$.

(1) $\mathsf K=\{\sigma\in\mathcal E^g:\operatorname{dist}(\sigma,\mathsf S_j)\to0\}$, a closed subset of the compact metrisable space $\mathcal E^g$ (Lemma~\ref{lh:rec}(b)). Convexity: if $\sigma^i=\lim_j\sigma_{\nu^i_j}$ and $\nu_j=t\nu^1_j+(1-t)\nu^2_j\in\mathcal G_{K_j,\beta_j}$, then
\[
\sigma_{\nu_j}(X)=t\,\nu^1_j(X^{\nu_j}_{K_j})+(1-t)\,\nu^2_j(X^{\nu_j}_{K_j}),
\]
and by \eqref{lh:rece} with the centrings $\nu^i_j$ and $\nu_j$ each bracket is at most $a_{K_j}^{-4}J(\beta_j)(2R+1)^3\|f\|_\infty\to0$ by (D)(1), while the moments $\nu^i_j((X_A)^{\nu^i_j}_{K_j})$ converge; so $\sigma_{\nu_j}\to t\sigma^1+(1-t)\sigma^2$.

(2) $(K_j,\beta_j,\nu_j)_j$ is a convergent Gibbs sequence at $g$; by (D)(3) each component is $\lim_j\sigma_{\nu_j(\cdot|B_j)}$ and has one vacuum. Components cluster at the rate $C^g(X,Y)/|v|$. A member of $\mathsf K$ that clusters, in the sense $\sigma(X\cdot\tau_vX)\to\sigma(X)^2$ for all $X$, has one vacuum, by (Can4) at $q=2$ as in Step~8; conversely, a member with one vacuum is its own component and clusters at the rate $C^g/|v|$. As $C^g$ and $r$ do not depend on the member, pointwise limits of members of $\mathsf K^\bullet$ cluster: $\mathsf K^\bullet$ is closed in $\mathsf K$, so compact.

(3) Barycentres of measures on $\mathsf K^\bullet$ lie in $\mathsf K$: for finite supports by (1), and in general by Riemann sums, the uniform continuity of $\psi\mapsto\psi(X)$ on the compact $\mathsf K^\bullet$, and the closedness of $\mathsf K$. Onto: for $\sigma\in\mathsf K$ the image of $\varpi$ under $\xi\mapsto\sigma_\xi$ is a measure on $\mathsf K^\bullet$ by (2), with the barycentre $\sigma$ by (Can4) at $q=1$. Injective: Lemma~\ref{lh:rec}(c) with $\mathsf T=\mathsf K^\bullet$. A continuous bijection of the compact $\mathrm{Prob}(\mathsf K^\bullet)$ onto the Hausdorff space $\mathsf K$ is a homeomorphism; by injectivity the image of $\varpi$ is $\mu_\sigma$; $\xi\mapsto\sigma_\xi$ is continuous and injective on the compact $\mathsf X$, and its image is the support of $\mu_\sigma$, $\varpi$ having full support. $\dim\ker H_\sigma=\#\mathsf X=\#\operatorname{supp}\mu_\sigma$ by Theorem~\textup{\YMMref{AP}{780}}.

(4) One vacuum is $\ker H_\sigma=\mathbb C\Omega_\sigma$ by $\#\mathsf X=\dim\ker H_\sigma$, in each of the six directions. By (3), the extreme points of $\mathsf K$ are the barycentres of the extreme points of $\mathrm{Prob}(\mathsf K^\bullet)$, the point masses: extreme points are the members with one vacuum. The two forms of clustering are (2). \emph{Pure phases.} If $\sigma=\lim_j\sigma_{\omega_j}$ with pure phases $\omega_j$, each $\omega_j$ is ergodic, so $E_{\omega_j}[\cdot\,|\mathcal I]$ is almost surely constant, and by the statement on iterated plateaux $\int\sigma_\xi(X)^2d\varpi-\sigma(X)^2=\lim_j\operatorname{Var}_{\omega_j}(E_{\omega_j}[X^{\omega_j}_{K_j}|\mathcal I])=0$ for every $X$; so $\xi\mapsto\sigma_\xi(X)$ is constant on $\mathsf X$, and $\mathsf X$ is a point.

Conversely, let $\sigma\in\mathsf K^\bullet$ be the limit of $\sigma_{\nu_j}$. By the variance statement below, $E_{\nu_j}[X^{\nu_j}_{K_j}|\mathcal I]\to\sigma(X)$ in $L^2(\nu_j)$ for every $X$. Let $(\nu_j^U)_U$ be a regular conditional law of $\nu_j$ given $\mathcal I$; by Proposition~\textup{\YMMref{106.1}{764}} its members are, off a $\nu_j$-null set, pure phases in $\mathcal G_{K_j,\beta_j}$, and $\nu_j^U(X^{\nu_j}_{K_j})=E_{\nu_j}[X^{\nu_j}_{K_j}|\mathcal I](U)$ for $\nu_j$-almost every $U$. Take a countable family $X_1,X_2,\dots$ of tuples on which pointwise convergence in $\mathcal E^g$ implies convergence at every tuple (Lemma~\ref{lh:rec}(b)). By Chebyshev's inequality there are, for each $j$, a point $U_j$ outside the null sets and numbers $\epsilon_j\to0$, $k_j\to\infty$ with $|\nu_j^{U_j}((X_k)^{\nu_j}_{K_j})-\sigma(X_k)|\le\epsilon_j$ for $k\le k_j$.

The pure phases $\omega_j=\nu_j^{U_j}$ then have $\omega_j((X_k)^{\nu_j}_{K_j})\to\sigma(X_k)$ for every $k$; by \eqref{lh:rece} and (D)(1), $\sigma_{\omega_j}(X_k)\to\sigma(X_k)$; and as $\sigma_{\omega_j}\in\mathcal E^g$, $\sigma_{\omega_j}\to\sigma$. \emph{Variances.} By the statement on iterated plateaux, $\operatorname{Var}_{\nu_j}(E_{\nu_j}[X^{\nu_j}_{K_j}|\mathcal I])\to\int\sigma_\xi(X)^2d\varpi-\sigma(X)^2$, which does not depend on the choice of the $\nu_j$, and which vanishes for every $X$ if and only if $\mathsf X$ is a point. \emph{Faces.} For a subsequence $\mathfrak p'$, $\mathsf K\subset\mathsf K_{\mathfrak p'}$. If $\sigma\in\mathsf K$ is $t\sigma_1+(1-t)\sigma_2$ with $\sigma_i\in\mathsf K_{\mathfrak p'}$, $0<t<1$, the uniqueness of the representing measure on ${\rm Lim}^\bullet[b]$ (Lemma~\ref{lh:rec}(c), with the compact set ${\rm Lim}^\bullet[b]$ of (6)) gives $\mu_\sigma=t\mu_{\sigma_1}+(1-t)\mu_{\sigma_2}$; so the $\mu_{\sigma_i}$ live on $\operatorname{supp}\mu_\sigma\subset\mathsf K^\bullet$, and $\sigma_i\in\mathsf K$ by (3).

For the interlacing $\mathfrak p''$, $\mathsf K_{\mathfrak p''}\subset\mathsf K\cap\mathsf K_{\tilde{\mathfrak p}}$, as $\mathfrak p$ and $\tilde{\mathfrak p}$ are subsequences of $\mathfrak p''$; conversely, two Gibbs sequences that converge to $\sigma$ along $\mathfrak p$ and along $\tilde{\mathfrak p}$ interlace to one that converges to $\sigma$ along $\mathfrak p''$. \emph{One point.} Let $\mathsf K\ne\emptyset$. If $\mathsf K$ is a point, $\mathsf K^\bullet$ is not empty, as it holds the components of that point, and lies in $\mathsf K$; if $\mathsf K^\bullet$ is a point, so is $\mathrm{Prob}(\mathsf K^\bullet)$ and hence $\mathsf K$; if $\mathsf K$ is a point, $\mathsf K=\mathsf K^\bullet$; if $\mathsf K=\mathsf K^\bullet$ had two members $\sigma_1\ne\sigma_2$, then $\frac12(\sigma_1+\sigma_2)\in\mathsf K=\mathsf K^\bullet$ would be an extreme point.

By the statement on iterated plateaux, $\operatorname{Cov}_{\nu_j}(E_{\nu_j}[X^{\nu_j}_{K_j}|\mathcal I],E_{\nu_j}[Y^{\nu_j}_{K_j}|\mathcal I])\to\int\sigma_\xi(X)\sigma_\xi(Y)d\varpi-\sigma(X)\sigma(Y)$, which vanishes for all $X,Y$ if and only if $\mathsf X$ is a point; for all convergent choices this says $\mathsf K=\mathsf K^\bullet$. At single fields, Lemma~\ref{lh:cond}(b) and (F)(2) give the displayed bound, with $|r_K(f)|\le(2R+1)^3\|f\|_\infty$, and it tends to $0$ by (D)(1).

(6) The pairs $(K,L^Kb)$, $K\in\mathbb N$, lie in $\mathfrak d(g)$; as $\mathcal E^g$ is compact, a subsequence of $(\sigma_{\nu_K})$ converges, so ${\rm Lim}[b]\ne\emptyset$. ${\rm Lim}[b]=\bigcap_k\overline{\bigcup\{\mathsf S_{K,\beta}:K\ge k,\ |a_K\beta-b|\le\frac1k\}}$, with $\mathsf S_{K,\beta}=\{\sigma_\nu:\nu\in\mathcal G_{K,\beta}\}$, is compact, and ${\rm Lim}^\bullet[b]$ is closed in it by the uniform clustering. A saturated subsequence: with a countable base $(U_r)$ of $\mathcal E^g$, pass to nested subsequences along which, for each $r$, $\mathsf S_j\cap U_r$ is either non-empty for all members or empty for all members, and take the diagonal; for it and for each of its subsequences a functional is a limit if and only if each $U_r$ that holds it meets $\mathsf S_j$ for all large $j$, so all these sets of limits agree, and they are not empty by compactness. Every $\sigma\in\mathsf K_{\mathfrak p}$ lies in $\mathsf K_{\mathfrak p'}$ for a saturated subsequence $\mathfrak p'$, and the extreme points are the members with one vacuum by (4); this gives the two unions. Uniqueness of the representing measure on ${\rm Lim}^\bullet[b]$: Lemma~\ref{lh:rec}(c) with $\mathsf T={\rm Lim}^\bullet[b]$; existence: (3).

\emph{The hull.} ${\rm Lim}^\bullet[b]$ is compact and its members cluster uniformly, with the numbers $C^g$, $r$; so by Lemma~\ref{lh:rec}(c) the map $\operatorname{bar}$ is affine, continuous and injective on the compact convex set $\mathrm{Prob}({\rm Lim}^\bullet[b])$, hence an affine homeomorphism onto its image, a compact convex set. The image contains ${\rm Lim}[b]$, by the existence just shown, and it is the closure of the image of the finitely supported measures, which is the convex hull of ${\rm Lim}^\bullet[b]\subset{\rm Lim}[b]$; so it is $\overline{\rm co}\,{\rm Lim}[b]$. Its extreme points are the images of the point masses, that is, the members of ${\rm Lim}^\bullet[b]$. Hence a member of ${\rm Lim}[b]$ with one vacuum is an extreme point of $\overline{\rm co}\,{\rm Lim}[b]$, and so of ${\rm Lim}[b]$; and a member $\sigma$ that is an extreme point of ${\rm Lim}[b]$ is an extreme point of every $\mathsf K_{\mathfrak p}$ that contains it, so it has one vacuum by (4).

\emph{Convexity.} If ${\rm Lim}[b]=\mathsf K_{\mathfrak q}$, it is convex by (1). If ${\rm Lim}[b]$ is convex, it is compact and convex, so it equals $\overline{\rm co}\,{\rm Lim}[b]=\operatorname{bar}(\mathrm{Prob}({\rm Lim}^\bullet[b]))$. If this equality holds and $\psi_1,\dots,\psi_n\in{\rm Lim}^\bullet[b]$, then $\frac1n\sum_i\psi_i$ lies in some $\mathsf K_{\mathfrak p}$; by (3) it is the barycentre of a measure on $\mathsf K^\bullet_{\mathfrak p}$, which by the uniqueness on ${\rm Lim}^\bullet[b]$ is $\frac1n\sum_i\delta_{\psi_i}$, so that $\psi_1,\dots,\psi_n\in\mathsf K^\bullet_{\mathfrak p}$. If every finite subset of ${\rm Lim}^\bullet[b]$ lies in one $\mathsf K^\bullet_{\mathfrak p}$, let $\psi_1,\psi_2,\dots$ be dense in ${\rm Lim}^\bullet[b]$ and $\psi_1,\dots,\psi_n\in\mathsf K^\bullet_{\mathfrak p_n}$; for each $n$ there is a pair $(K'_n,\beta'_n)$ of $\mathfrak p_n$ with $K'_n\ge n$ and $|a_{K'_n}\beta'_n-b|\le\frac1n$ which carries, for each $i\le n$, a state $\nu$ with $\operatorname{dist}(\sigma_\nu,\psi_i)\le\frac1n$.

Along $\mathfrak q=(K'_n,\beta'_n)_n$, a sequence of pairs at the label $b$, every $\psi_i$ is a limit; so the closed set $\mathsf K_{\mathfrak q}$ contains ${\rm Lim}^\bullet[b]$, and by (3) for $\mathfrak q$ it contains $\operatorname{bar}(\mathrm{Prob}({\rm Lim}^\bullet[b]))\supset{\rm Lim}[b]\supset\mathsf K_{\mathfrak q}$. Finally, a set with one element is convex.

(5) Every $\sigma\in{\rm Lim}[b]$ lies in some $\mathsf K_{\mathfrak q}$; its components lie in $\mathsf K^\bullet_{\mathfrak q}\subset{\rm Lim}^\bullet[b]$, and $\sigma=\int\sigma_\xi\,d\varpi$. If ${\rm Lim}^\bullet[b]=\{\tau\}$, all components are $\tau$ and $\sigma=\tau$; if ${\rm Lim}[b]=\{\sigma\}$, then ${\rm Lim}^\bullet[b]$ is not empty by (6) and lies in $\{\sigma\}$. The same lines serve for ${\rm Lim}_{\mathfrak q}[b]$, the components of a limit along pairs of $\mathfrak q$ being limits along the same pairs, and ${\rm Lim}_{\mathfrak q}[b]$ being non-empty by (6) applied to a sequence with pairs in $\mathfrak q$. For the third statement: a Gibbs sequence at $g$ with the label $b$ has convergent subsequences (Theorem~\textup{\YMMref{AP}{780}}), whose limits lie in ${\rm Lim}[b]$; if this is one point, every subsequence has a subsequence that converges to it, so the sequence converges; conversely, interlacing two convergent Gibbs sequences with the label $b$ gives one, which converges, so their limits agree. The one limit then has one vacuum, as ${\rm Lim}^\bullet[b]\ne\emptyset$.

The clauses (1)--(4) use of (A) only that $a_{K_j}^{-4}J(\beta_j)\to0$ along $\mathfrak p$. Along $\mathfrak p=(K,L^Kb)_K$ this holds without (A) for every $b$ outside the countable set $\bigcup_KL^{-K}\{\beta:J(\beta)>0\}$, as the convex function $P$ is differentiable off a countable set; at such labels these clauses use Theorems~\YMMref{D}{8}, \YMMref{AJ}{673} and~\YMMref{AD}{607} at the choice of the setting alone.

\emph{Step 11: (G) at every label.} For $z\in\mathbb Z$ put, as in the paragraph on dilations, $X^{[z]}=(f_i(L^z\cdot))_i$ and $\sigma^{[z]}(X)=L^{-z|X|}\sigma(X^{[z]})$. The dilation $x\mapsto L^zx$ carries a law $\nu$ on $\mathbb L_K$ to a law $\nu^{(z)}$ on $\mathbb L_{K-z}$, maps $\mathcal G_{K,\beta}$ onto $\mathcal G_{K-z,\beta}$ and conditioned states to conditioned states, as it carries the invariant events to the invariant events (Lemma~\ref{lem:lattice-law-dilation}), and turns $\mathcal O_{K-z}(f)$ into $L^{-z}\mathcal O_K(f(L^z\cdot))$ (Lemma~\ref{lem:lattice-law-dilation}, applied to the function $f(L^z\cdot)$); so
\[
\sigma_{\nu^{(z)}}=(\sigma_\nu)^{[z]},\qquad a_{K-z}^{-4}J(\beta)=L^{-4z}a_K^{-4}J(\beta).
\]
So $\sigma\mapsto\sigma^{[z]}$ is a linear homeomorphism that carries $\mathsf K_{\mathfrak p}$, ${\rm Lim}[b]$, ${\rm Lim}^\bullet[b]$ onto the same sets at the label $L^zb$ for the pairs $(K_j-z,\beta_j)$, and canonical decompositions to canonical decompositions, component by component, with the same space $\mathsf X$. It carries the reconstruction as well. Let $(z,g)$ be a pair for $b$, that is, $z\in\mathbb Z$ and $0<g\le\gamma$ with $g^{-2}+\bar B<L^zb<X_g-\bar B$; the pair fixed for $b$ in the paragraph on dilations is one, with $z\ge0$. For $\sigma\in\mathsf K_{\mathfrak p}$, $\sigma^{[z]}$ is then, as in that paragraph, the limit of a convergent Gibbs sequence at $g$ over the pairs $(K_j-z,\beta_j)$ of large index, and Theorem~\textup{\YMMref{AP}{780}} gives its reconstruction, with the Hamiltonian $H_{\sigma^{[z]}}$ and the vacuum vector $\Omega_{\sigma^{[z]}}$, in each of the six coordinate directions.

For a coordinate direction $u$, the reflection $\theta$ in the plane $\{x\cdot u=0\}$ commutes with the dilation, so that $\sigma^{[z]}(\theta X\cdot Y)=L^{-z(|X|+|Y|)}\sigma(\theta X^{[z]}\cdot Y^{[z]})$ for the tuples $X$, $Y$ in the half space $\{x\cdot u>0\}$; hence the form $\langle[X],[Y]\rangle_\sigma=\sigma(\theta X\cdot Y)$ is positive together with that of $\sigma^{[z]}$, and $W:[X]\mapsto L^{-z|X|}[X^{[z]}]$ is a unitary map of the space of the reconstruction of $\sigma^{[z]}$ onto the space that this form builds from $\sigma$, with $W[\emptyset]=[\emptyset]$. As $(\tau_vX)^{[z]}=\tau_{L^{-z}v}X^{[z]}$, $W$ carries the translations by $tu$ to the translations by $L^{-z}tu$; so, with $\hat H_\sigma$ the generator of the semigroup $e^{-t\hat H_\sigma}[X]=[\tau_{tu}X]$, $t\ge0$, in the space that the form builds from $\sigma$, $\hat\Omega_\sigma=[\emptyset]$ the class of the empty tuple there, and $e^{-tH_{\sigma^{[z]}}}[X]=[\tau_{tu}X]$, $\Omega_{\sigma^{[z]}}=[\emptyset]$ in the reconstruction of $\sigma^{[z]}$,
\[
H_{\sigma^{[z]}}=L^{-z}\,W^*\hat H_\sigma W,\qquad W\Omega_{\sigma^{[z]}}=\hat\Omega_\sigma,\qquad\ker\hat H_\sigma=W\ker H_{\sigma^{[z]}},
\]
in each of the six directions: the kernels and the vacua correspond. A linear homeomorphism carries convex sets, closed sets, extreme points, faces and one-point sets to the same, and the dilation carries the clustering numbers to dilated ones. For every $b>0$ there are pairs for $b$, already with $z\ge0$, because $L^zb>\gamma^{-2}+\bar B$ for large $z$ and the open intervals $(g^{-2}+\bar B,X_g-\bar B)$ overlap one another as $g$ decreases and cover $(\gamma^{-2}+\bar B,\infty)$. So $W$ carries the operator $L^zH_{\sigma^{[z]}}$ and the vector $\Omega_{\sigma^{[z]}}$, which are $H_\sigma$ and $\Omega_\sigma$ of the paragraph on dilations at the pair $(z,g)$ fixed for $b$, to $\hat H_\sigma$ and $\hat\Omega_\sigma$, the generator of the translations and the class of the empty tuple in the space that the form builds from $\sigma$.

These read no window, so another pair $(z,g)$ gives a unitarily equivalent operator with the corresponding vector, and $\dim\ker H_\sigma=\dim\ker H_{\sigma^{[z]}}=\#\mathsf X$, both finite and equal or both infinite. Step~10 at $(L^zb,(K_j-z,\beta_j))$ gives (1)--(6) at $(b,\mathfrak p)$. In particular the kernel clauses of (3) and (4) pass to $(b,\mathfrak p)$, with $H_\sigma=L^zH_{\sigma^{[z]}}$ and $\Omega_\sigma=\Omega_{\sigma^{[z]}}$ of the paragraph on dilations at the pair $(z,g)$ fixed for $b$: the dilation carries $\mu_\sigma$ to $\mu_{\sigma^{[z]}}$, so that $\dim\ker H_\sigma=\dim\ker H_{\sigma^{[z]}}=\#\mathsf X=\#\operatorname{supp}\mu_\sigma$, and $\ker H_\sigma=\mathbb C\Omega_\sigma$ in one, and then in each, of the six directions if and only if $\sigma$ has one vacuum.

\emph{Step 12: (G)(7).} Let $0<g\le\gamma_W$ and let $\mathfrak p$ be a sequence of pairs in $\mathfrak d(g)$ with $\mathsf K_{\mathfrak p}\ne\emptyset$; $\sigma\in\mathsf K_{\mathfrak p}$ is the limit of a convergent Gibbs sequence over $\mathfrak p$. By Proposition~\YMMref{127.8}{854}, applied with $\rho_j=\nu_j$ (see also Corollary~\textup{\YMMref{127.11}{854}}), along a convergent Gibbs sequence at $g$ the labels $a_{K_j}\beta_j$ converge, to a number $b\in I^{\rm lab}(g)$, and the two-point function of the limit, $S^\sigma_2=\lim_jS^T_2[\nu_j]$, has the displayed limit with this $b$, for every $\vartheta\in(0,1)$ and every pair of $\mathfrak T^1_{\mathbb R^3}(\mathsf d_0(g,\vartheta);\vartheta)$; this set is not empty by Lemma~\ref{lem:ward-test-pairs}(c).

The right member of the formula reads $\sigma$ alone; so the sets ${\rm Lim}_{\mathfrak d(g)}[b]$ are pairwise disjoint, and every convergent Gibbs sequence at $g$ is over a sequence of pairs with a non-empty set of limits, which has a label. An interlacing of two sequences with different labels has no label, so its set of limits is empty. Suppose that $\sigma''=t\sigma+(1-t)\sigma'$, $0<t<1$, with $\sigma\in{\rm Lim}_{\mathfrak d(g)}[b]$, $\sigma'\in{\rm Lim}_{\mathfrak d(g)}[b']$, were the limit of a Gibbs sequence over $\mathfrak p''$; then $\mathfrak p''$ has a label $b''$. The sets of the limits with one vacuum at the labels $b$ and $b'$, along sequences in $\mathfrak d(g)$, are compact, by the diagonal argument of Step~10 (6), and their members cluster uniformly, with the numbers $C^g$, $r$; with the representing measures $\mu_\sigma$, $\mu_{\sigma'}$ of (G)(6), carried by the components of $\sigma$ and of $\sigma'$, Lemma~\ref{lh:rec}(c) makes $(\operatorname{supp}\mu'',\mu'',(\psi))$, $\mu''=t\mu_\sigma+(1-t)\mu_{\sigma'}$, a canonical decomposition of $\sigma''$.

By Corollary~\textup{\YMMref{127.11}{854}}, the formula in every component of every canonical decomposition, each component of $\sigma''$ has the coefficient $b''^{-2}$, whereas the members of $\operatorname{supp}\mu_\sigma$, which are the components of $\sigma$ by (G)(3), have $b^{-2}$, and those of $\operatorname{supp}\mu_{\sigma'}$ have $b'^{-2}$; so $b=b''=b'$. \qed

\begin{proof}[Proof of Theorem~\ref{lh:box}]
(1) is Proposition~\ref{lh:onepair}(c) at $a=a_K$, $\beta=L^Kb$ and $l=l_K=L^{K+m_*}\ge L$, as $K+m_*\ge K_*+m_*\ge1$; that $P$ is differentiable at $\beta$ is $J(\beta)=0$. (2) Let first the entries be none $0$ and their supports lie in $(-\frac{\ell''}4,\frac{\ell''}4)^3$ for a side $\ell''=2L^{m''}\ge\ell_*$, $m''\in\mathbb N$; an entry $0$ makes every cumulant $0$, and every compact set lies in such a cube. The torus cumulants converge as $K\to\infty$ (Theorem~\textup{\YMMref{AS}{813}}), so Proposition~\ref{lh:onepair}(f) gives $\limsup_KS^T_n[\nu_K](X)-\liminf_KS^T_n[\nu_K](X)\le2C(X)\theta^{\ell'/4\ell_*}$ for every side $\ell'\ge\ell''$; let $\ell'\to\infty$. The same bound shows that the limit is within $C(X)\theta^{\ell'/4\ell_*}$ of $\lim_KS^T_{n;K,\ell',L^Kb}(X^{\langle\ell'\rangle})$. (3) Since $\mathcal G_{K,L^Kb}=\{\nu_K\}$, the set $\mathsf K_{\mathfrak p}$ has at most one point, the limit of $\sigma_{\nu_K}$ if it exists. For a separated $X$, $\sigma_{\nu_K}(X)$ is a centred moment, so
\[
\sigma_{\nu_K}(X)=\sum_\pi\prod_{A\in\pi}S^T_{|A|}[\nu_K](X_A),
\]
the sum over the partitions of $\{1,\dots,|X|\}$ into blocks of at least two elements; each factor converges by (2). So $\mathsf K_{\mathfrak p}=\{\sigma\}$; the point of a one-point convex set is extreme, so $\sigma$ has one vacuum by Theorem~\ref{lh:main}(G)(4). The last sentence: $l_K=\ell_*/2a_K$ and $\beta=L^Kb$ do not change when $(b,\ell_*,K)$ is replaced by $(L^zb,L^z\ell_*,K-z)$, and $K\ge K_*$ says $K-z\ge K_*-z$.
\end{proof}

\subsection{Scope of the conclusions}

The bounds (A)--(C) are upper bounds; they do not show that $J$ vanishes at any finite $\beta$. The numbers $54$, $90$ and $\frac1{180}$ are those of sufficient conditions in the estimates: $54=3\cdot18$ comes from the condition $\ell^3L^{-K/18}\le c$ of Step~4 and stands in the threshold $k^{\rm rt}+54m$ of (E), which the slope $A>90$ of Step~5 covers in (B); $90=9\cdot10$ comes from the comparison of $\ell^9$ with $L^{-K/10}$ in Step~5; and $\frac1{180}=\sup_{A>90}\frac1{2A}$ comes from Proposition~\ref{lh:per}(b) in Step~6. Proposition~\ref{lh:per}(e) is a statement on data and on the map from the slope to the exponent: for every slope $A>0$ there are data as in (b), with $A_1=A_2=A$, whose atoms are $\beta_k^{-1/2A}$, so that the loss of one half is real for data (Remark~\ref{lh:onep}(3)); and the data of (e) at $A=90$ satisfy the hypothesis of (b) at every slope $A'\ge90$, a larger threshold being a weaker hypothesis, and have the atoms $\beta_k^{-1/180}$.

By Proposition~\ref{lh:pow} the exponent follows the slope that a rate on the tori gives, and by Proposition~\ref{lh:cut} a threshold that does not grow with the side gives $J=0$ on a half line. Theorem~\ref{lh:box} assumes the box inequality and speaks of the couplings $L^Kb$ of one label. Clause (G)(5) is an equivalence and asserts neither of its two sides; clause (G)(6) asserts none of its four equivalent statements on the convexity of ${\rm Lim}[b]$; and (G)(1)--(4) allow $\mathsf K_{\mathfrak p}=\emptyset$ for the given sequence, whereas (G)(6) gives non-empty sets along saturated subsequences. Clause (E) concerns $n\ge2$ and tuples of positive spread, with a threshold $k^{\rm rt}+54m$ that grows with the side, and (A)(4) concerns limit points along subsequences of the cutoffs.

\section{Cube means and infrared bounds}\label{C-365-1}

We average the action density over cubes with a boundary collar. Chessboard estimates turn control of the source pressure into bounds on correlations of distinct cube means. Comparing two blocking factors with incommensurable logarithms excludes concentration of the pressure response.

The resulting continuum bounds give an $O(r^{-2})$ radial two-point kernel and at most linear spectral mass near zero. They exclude an atom at zero without giving a positive gap. The growing-torus conclusions use only the positivity and symmetry established for those limits. The stronger exponential conclusions explicitly assume the box criterion.

\subsection{Notation and hypotheses}

Lengths are measured on the physical scale, on which the bare lattice spacing is $a_K=L^{-K}$; thus $\ell$, $M$, $\mathsf l$, $T$, $q$ and $r$ are lengths on this scale. A floor is a lower threshold, a window a compact set of labels; a quantity is free of another if it does not depend on it. $|\cdot|$ is the Euclidean norm, $|\cdot|_\infty$ the sup-norm and $\operatorname{dist}_\infty$ its distance (on a torus, the distance of the torus); $[t]_+=\max\{t,0\}$; $\operatorname{int}I$ and $|I|$ are the interior and the length of an interval $I$; $\operatorname{Lip}$ is the Lipschitz constant for the Euclidean distance.

\begin{definition}[Frames; the numbers at a reference coupling]\label{C-365-2}
\begin{enumerate}
\item[(1)] A \emph{frame} consists of a radius $\varrho>0$, a common real and sourced choice of constants obtained by the ordered construction in the proof of Proposition~\YMMref{69.1}{595} and Section~\YMMref{82}{665}, and the corresponding Wilson measures and source normalization of Theorems~\YMMref{AJ}{673} and~\YMMref{AD}{607}. The choice retains the source-room condition~\eqref{eq:canonical-source-quarter-room} and the array radius~\eqref{eq:array-borel-radius}. A frame determines an odd integer $L\ge13$; $M=L^{m'}$ with an integer $m'$; numbers $B_\sharp\ge0$ and $\bar B=B_\sharp L/(L-1)$; $\gamma\in(0,\frac12]$; and the upper stability constant $C_S$ and decay rate $\kappa>\log9216+2$ of Theorem~\YMMref{AD}{607}. Choose $C_{\rm st}$ and $C_S\ge C_{\rm fp}^{\rm st}$ as in \eqref{eq:setup-pressure-constants}. Lemma~\ref{lem:setup-pressure-stability} gives the lower stability bound with $C_g=C_{\rm st}+\frac92\log X_g$ and the estimate $C_g\le C_S(1+\log g^{-2})$. We put $S(\kappa)=2^{12}(1-9216\,e^{-\kappa})^{-1}$, the number of the count of domains.
\item[(2)] At a reference coupling $0<g\le\gamma$ we put
\[X_g=L(g^{-2}+B_\sharp),\quad I(g)=(g^{-2}+\bar B,\,X_g-\bar B],\quad C_g=C_{\rm st}+\tfrac92\log X_g,\quad B'=2C_S+C_g,\]
and $I^{\rm lab}(g)=[g^{-2}+B_\sharp\frac{L-2}{L-1},\,L(g^{-2}+\bar B)]$. A first-passage datum at $g$ is a pair $(K,\beta)$ of a cutoff and a bare coupling such that $K$ is the first index at which the running inverse coupling started at $\beta$ is $\le X_g$, as in Proposition~\YMMref{74.1}{615}; $\mathfrak d(g)$ is the set of these pairs. $\mathfrak L(g)=\{2L^m:m\ge m_{\mathfrak L}(g)\}$ is the set of the sides at $g$ of equation~\YMMref{(1528)}{615}, $m_{\mathfrak L}(g)$ being the integer \YMMref{(1528)}{615} there, which the frame and $g$ fix. The set $\mathfrak L(g)$ is closed upward among the sides $2L^m$, and by Lemma~\ref{lem:setup-label-coverage} and equation~\YMMref{(1533)}{615}, for $0<g'\le g\le\gamma$, $m_{\mathfrak L}(g)\le m_{\mathfrak L}(g')$, so that $\mathfrak L(g')\subset\mathfrak L(g)$. These are the sides of which Theorem~\textup{\YMMref{AD}{607}} and the transfer (M7) speak.
\item[(3)] A second frame, whose objects wear a hat, is a \emph{partner} of the first if $\log L/\log\hat L\notin\mathbb Q$.
\end{enumerate}
\end{definition}

\begin{definition}[The lattice, the pressures, the sources]\label{C-365-3}
\begin{enumerate}
\item[(1)] For an integer $K\ge0$, $a_K=L^{-K}$ and $\mathbb L_K=a_K\mathbb Z^3$; for $\ell=2L^m$, $T_\ell=(\mathbb R/\ell\mathbb Z)^3$ carries the lattice of spacing $a_K$, with $\mathsf N=\ell/a_K$ sites a side. For a plaquette $p$, $a_p=1-\frac12\operatorname{tr}\partial U(p)\in[0,2]$, and $x(p)$ is its barycentre. For a function $f$ we put $\mathcal O(f)=\mathcal O_K(f)=a_K^{-1}\sum_pf(x(p))a_p$ and $\operatorname{Rie}_K(f)=\frac13a_K^3\sum_pf(x(p))$, the sums over the plaquettes of $\mathbb L_K$ or of the torus.
\item[(2)] $Z_{K,\ell,\beta}[w]=\int e^{-\beta\sum_pa_p+\mathcal O(w)}\,dU$, with the Haar measure of mass $1$; $Z_{K,\ell,\beta}=Z_{K,\ell,\beta}[0]$ defines the law $\langle\cdot\rangle_{K,\ell,\beta}$. $S^T_{n;K,\ell,\beta}(f_1,\dots,f_n)$ and $S^T_n[\nu](f_1,\dots,f_n)$ are the joint cumulants of $\mathcal O(f_1),\dots,\mathcal O(f_n)$ under the law of the torus and under a law $\nu$; $\operatorname{Cov}$ and $\operatorname{Var}$ are covariance and variance. $\mathcal G^{\rm per}_{K,\beta}$ is the set of the periodic Gibbs states on $\mathbb L_K$ at $\beta$, in the sense of the Definitions of kernels and Gibbs states (Section~\YMMref{91}{697}).
\item[(3)] The label of $(K,\beta)$ is $b=a_K\beta$. We put $p_{K,\ell}(b)=\ell^{-3}\log Z_{K,\ell,L^Kb}$ and $p_{K,\infty}(b)=\inf_\ell p_{K,\ell}(b)$ over the sides $\ell=2L^m$. For a reference side $\ell'$: $D_{K,\ell}=p_{K,\ell}-p_{K,\ell'}$, $D_K=p_{K,\infty}-p_{K,\ell'}$, and $\bar\eta(\ell)=3\ell'B'/\ell$. $P$ is the pressure of the Wilson model on $\mathbb Z^3$ (Lemma~\ref{C-365-16}) and $J(\beta)=P'(\beta+)-P'(\beta-)$. For a function $D$ on an interval such that $D+\frac12Cy^2$ is convex, $D''$ is the signed measure $(D+\frac12Cy^2)''-C\,dy$; $\Delta^2_sD(b)=D(b+s)+D(b-s)-2D(b)$.
\item[(4)] A \emph{source} is a complex Lipschitz function $w$ on $T_\ell$ with $\|w\|_\sharp=\max\{\sup|w|,40M\operatorname{Lip}w\}<2\varrho$; this class lies in the class of the sources of Theorem~\textup{\YMMref{AD}{607}}. A real number $x$ in the place of a source is the constant function $x\cdot1$. Since $\mathcal O(x\cdot1)=L^Kx\sum_pa_p$, we have the \emph{shift of the label}
\[Z_{K,\ell,L^Kb}[w+x]=Z_{K,\ell,L^K(b-x)}[w].\]
\item[(5)] For $(K,\beta)\in\mathfrak d(g)$ and $\ell\in\mathfrak L(g)$, Theorem~\textup{\YMMref{AD}{607}} provides the source energies: a number $E_\ell$ and local terms $e_j(X;w)$, $1\le j\le K$, $X$ a localisation domain of level $j$ with linear size $d_j(X)$. We put $\mathcal E_\ell=\sum_{j,X}e_j(X;\cdot)$ and $\widetilde Z_\ell(w)=e^{-E_\ell-\mathcal E_\ell(w)}Z_{K,\ell,\beta}[w]$; the index $\ell$ is dropped where one torus is fixed. $\mathsf G$ is the finite group generated by the translations of $T_\ell$ by $a_K\mathbb Z^3$ and the reflections $x_\mu\mapsto-x_\mu$, acting on sources by $(\mathsf rw)(y)=w(\mathsf r^{-1}y)$; it contains the reflections at all mirrors of all boards. We put $\mathcal E^{\rm av}(w)=|\mathsf G|^{-1}\sum_{\mathsf r\in\mathsf G}\mathcal E(\mathsf rw)$ and $\widetilde Z^{\rm av}=e^{-E-\mathcal E^{\rm av}}Z$.
\end{enumerate}
\end{definition}

\begin{definition}[Boards, profiles, cube means]\label{C-365-4}
\begin{enumerate}
\item[(1)] $\mathsf l$ and $\mathsf l_1$ are the least powers of $L$ with $\mathsf l\ge32M$ and $\mathsf l_1\ge8M+16$; $T_0=640\sqrt3M$. The width of the collar is a real number $T\ge\max\{T_0,\mathsf l\}$; the side of a board that carries profiles and cube means is a power $q=L^i\ge2T$; we write $q_i=L^i$.
\item[(2)] For a power $\mathsf j=L^i$ with $\mathsf j\ge La_K$, a \emph{board} of side $\mathsf j$ with offset $o\in a_K\mathbb Z^3$ is the partition of $\mathbb L_K$, or of $T_\ell$ when $\ell\ge2\mathsf j$, into the cubes $\square_c$ of side $\mathsf j$ centred at the points $o+\mathsf jc$, $c\in\mathbb Z^3$ (on $T_\ell$, $c$ modulo $\ell/\mathsf j$), the \emph{cells}; $c$ is the \emph{index} of the cell. Since $\mathsf j/a_K$ is odd, the centres are sites and the \emph{mirrors}, the planes half way between neighbouring centres, pass through midpoints of bonds. On $T_\ell$ there are $N=(\ell/\mathsf j)^3$ cells, an even number in each direction, and every cell is even or odd, by the parity of $c_1+c_2+c_3$. The mirror copy $F^{[c']}$ of a function $F$ of the cell $c$ in the cell $c'$ is its image under the reflections at the mirrors that carry $c$ to $c'$, complex conjugated if $c$ and $c'$ have different parities; $\|F\|_{\rm ch}=\langle\prod_{c'}F^{[c']}\rangle^{1/N}$ is the chessboard seminorm (equation~\YMMref{(1756)}{681}). Profiles and cube means are formed on boards of the sides $q$ of (1); boards of the sides $\mathsf l_1$ and $\mathsf l$ enter the proofs as well.
\item[(3)] For a cell $\square_c$ of a board of side $q$, $\mathfrak P_T(\square_c)$ is the set of the real $C^1$ functions $f$ with $0\le f\le1$, $f=1$ at sup-distance $\ge T/2$ from $\partial\square_c$, $f=0$ off $\square_c$ and at sup-distance $\le T/8$ from $\partial\square_c$, and $\operatorname{Lip}f\le16\sqrt3/T$; then $\|f\|_\sharp=1$. Its members are the \emph{profiles} of the cell.
\item[(4)] The \emph{cube mean} of the cell $c$ with the profile $f$ is $x_c=q^{-3}\mathcal O(f)$.
\item[(5)] $\mathsf c=\max\{1,40\sqrt3M(16/T+4/\mathsf l)\}\in[1,1+5\sqrt3]$; $C_w=96\mathsf c(C_S+C_g)/\varrho$, and $C_w^-$, $C_w^+$ are the same number with $1$, respectively $1+5\sqrt3$, in the place of $\mathsf c$; $C_2=72e^4C_w^2$; $C_I'=16(C_S+C_g)/(9\varrho^2)$; $x_q=1/(3C_wTq^2)$ and $I_q(b)=[b-x_q,b+x_q]$.
\item[(6)] On $T_\ell$, $\ell\ge2q$, for a profile $f$ of a cell $c$: $F_e$ and $F_o$ are the sums of the mirror copies of $f$ over the cells of the parity of $c$ and over the others, $F=F_e+F_o$, and $1-F$ are the \emph{walls}. Copies in different cells have supports at sup-distance $\ge T/4$, and $\|F\|_\sharp=\|F_e-F_o\|_\sharp=1$. For complex $z$, $w_z=zF_e+\bar zF_o$ is the source of the source formula for the source $zf$ of the cell, and $\Lambda(z)=\log\|e^{z\mathcal O(f)}\|_{\rm ch}=N^{-1}\log(Z[w_z]/Z)$.
\item[(7)] The model profile is $f_0=\prod_\mu\varphi_0(x_\mu)$, with $\varphi_0$ even, $C^1$, $0\le\varphi_0\le1$, non-increasing in $|s|$, $\varphi_0=1$ for $|s|\le q/2-3T/8$, $\varphi_0=0$ for $|s|\ge q/2-T/4$, $|\varphi_0'|\le16/T$; then $f_0(\cdot-y)\in\mathfrak P_T(\square_c)$ whenever $|y-y_c|_\infty\le T/8$, $y_c$ the centre of $\square_c$. With $f_0$ formed at the side $q_i$: $w_i=q_i^{-3}f_0(\cdot-\frac12q_ie_3)$ and $\theta:x_3\mapsto-x_3$; for $\mu=1,2,3$, $w_i^{(\mu)}=q_i^{-3}f_0(\cdot-\frac12q_ie_\mu)$ and $\theta_\mu:x_\mu\mapsto-x_\mu$.
\end{enumerate}
\end{definition}

\begin{definition}[Limits]\label{C-365-5}
\begin{enumerate}
\item[(1)] $\mathfrak S_n(O)$ is the set of the $n$-tuples of real functions of $C^1_c(O)$ with pairwise disjoint supports, and $\mathfrak S(O)=\bigcup_n\mathfrak S_n(O)$, the separated tuples of Section~\ref{sec:state-limit-conventions}. $X\cdot Y$ is the concatenation of tuples, $\tau_a$ the translation by $a$, $(h)$ the tuple with the one entry $h$; translations and reflections act entry by entry. All assertions below on limits concern such tuples.
\item[(2)] A \emph{Gibbs sequence} at $g$ is a sequence $(K_j,\beta_j,\nu_j)$ with $(K_j,\beta_j)\in\mathfrak d(g)$, $K_j\to\infty$ and $\nu_j\in\mathcal G^{\rm per}_{K_j,\beta_j}$; its labels are $b_j=a_{K_j}\beta_j$. It \emph{converges} if $\sigma(X)=\lim_j\nu_j\bigl(\prod_i(\mathcal O(f_i)-\nu_j(\mathcal O(f_i)))\bigr)$ exists for every $X=(f_1,\dots,f_n)\in\mathfrak S(\mathbb R^3)$. By the theorem on their limits (Theorem~\textup{\YMMref{AP}{780}}), $\sigma$ is invariant under $\mathbb R^3\rtimes\mathrm O(3)$ and reflection positive at every mirror; its connected functions $S^\sigma_n$ are the limits of the cumulants $S^T_n[\nu_j]$; $\sigma((f))=0$; and in a canonical decomposition $(\mathsf X,\lambda,(\sigma_\xi))$ of $\sigma$, in the sense of Theorem~\YMMref{AL}{746}, applied to the classes
of periodic Gibbs states in Theorem~\YMMref{AP}{780}, every component is a limit along the full sequence, $\sigma_\xi(X)=\lim_j\rho_j\bigl(\prod_i(\mathcal O(f_i)-\nu_j(\mathcal O(f_i)))\bigr)$ with states $\rho_j\in\mathcal G^{\rm per}_{K_j,\beta_j}$, is invariant and reflection positive likewise, and has $\sigma_\xi((f))=\mathfrak c_\xi\int f$ with a number $\mathfrak c_\xi$, its one-point constant. Both obey the bound tables (Proposition~\YMMref{106.3}{765}) of Theorem~\textup{\YMMref{AP}{780}}.
\item[(3)] For $\sigma'\in\{\sigma,\sigma_\xi\}$ there is a real-analytic function $\mathbf s\ge0$ on $(0,\infty)$ such that
\[
 \sigma'(f,h)=\iint f(x)h(y)\mathbf s(|x-y|)\,dx\,dy,
 \qquad (f,h)\in\mathfrak S_2(\mathbb R^3).
\]
Indeed, Theorem~\YMMref{AP}{780} supplies the reconstruction data, and Proposition~\YMMref{108.3}{771} gives their real-analytic two-point kernel off the diagonal. Translation and rotation invariance identify this kernel with $\mathbf s(|x-y|)$. At separation $re_3$, its point-operator formula is
\[
 \mathbf s(r)
 =\langle\Omega,\mathsf B(e_3,\tfrac r2e_3)
                    \mathsf B(\tfrac r2e_3,e_3)\Omega\rangle
 =\|\mathsf B(\tfrac r2e_3,e_3)\Omega\|^2\ge0.
\]
Here the first equality follows by joining the two inner semigroup intervals and using that the vacuum is fixed by translations; the second is the adjoint identity for the point operators. Holomorphy of these operators for $\operatorname{Re}r>0$ proves the stated real analyticity.
\item[(4)] The reflection Hilbert space and translation construction in Proposition~\YMMref{94.1}{720} gives, for such a $\sigma'$ and the direction $e_3$, a Hilbert space $\mathcal H(\sigma',e_3)$ with vectors $[X]$, $X\in\mathfrak S(\mathbb E_+)$, $\mathbb E_+=\{x_3>0\}$, $\langle[X],[Y]\rangle=\sigma'(\theta X\cdot Y)$; a self-adjoint $H\ge0$ and unitary operators $\mathcal U(v)$, $v\perp e_3$, commuting and strongly continuous, with $[\tau_{te_3+v}X]=e^{-tH}\mathcal U(v)[X]$; and the joint spectral measure $\mathcal P$ of $(H,P_1,P_2)$, carried by $[0,\infty)\times\mathbb R^2$, with $e^{-tH}\mathcal U(v)=\int e^{-tE+iv\cdot\mathsf q}\,d\mathcal P(E,\mathsf q)$. For real $h\in C^1_c(\mathbb E_+)$ we put $\psi_h=[(h)]$ and $\mu_h=\langle\psi_h,\mathcal P(\cdot)\psi_h\rangle$; $Q_r=[0,r]\times[-r,r]^2$.
\item[(5)] A continuum limit $S$ at $g$ \emph{along growing tori} has a defining sequence $(K_j,\ell_j,\beta_j)$ with $(K_j,\beta_j)\in\mathfrak d(g)$, $\ell_j\in\mathfrak L(g)$, $K_j\to\infty$ and $\ell_j\to\infty$, no relation between $K_j$ and $\ell_j$ being asked, and $S_n(X)=\lim_jS^T_{n;K_j,\ell_j,\beta_j}(X^{\langle\ell_j\rangle})$ for $n\ge2$, $X^{\langle\ell\rangle}$ the periodisation; $\sigma_S$ is its centred functional, formed from the $S_n$ with $S_1=0$. By Propositions~\YMMref{92.1}{710}, \YMMref{92.3}{712}
and~\YMMref{92.4}{713}, $\sigma_S$ obeys a bound table with one-point constant $0$, is invariant under the translations of $\mathbb R^3$ and the cubic group, and is reflection positive at every coordinate mirror; the reconstruction applies to it in each of the directions $e_1,e_2,e_3$.
\end{enumerate}
\end{definition}

\begin{lemma}[Every frame has a partner]\label{C-365-6}
For every frame there is a partner frame with the same radius and the same values of the constants that precede $L$ in the list $\mathcal T^{\rm S}[2\varrho]$: one may take $\hat L=L^k+2$ for every large integer $k$. For every radius there are also partner frames whose block sizes are $L_1$ and $L_1+2$. The cited construction theorems apply to each of these common choices, with its own subsequent constants and coupling ceiling.
\end{lemma}

\begin{proof}[Sketch of proof]
Keep the exponents, reserve fractions and other constants preceding $L$ in the ordered construction that defines the frame. The proof of Proposition~\YMMref{69.1}{595} admits every sufficiently large odd value of the new block size with these preceding values fixed. The further floors of Section~\YMMref{82}{665} preserve this property. Thus $\hat L=L^k+2$ is admitted for every sufficiently large $k$. Complete the same construction at $\hat L$, choosing all subsequent constants anew and retaining the source-room and array-radius conditions at their respective choices; choose the coupling ceiling last. This gives a second frame with the same radius and the same pre-$L$ data. Since $L$ is odd, $L^k+2$ is odd and $\gcd(L^k+2,L)=\gcd(2,L)=1$; if $L^p=\hat L^{p'}$ with positive integers $p,p'$, a prime factor of $L$ would divide $\hat L$, so
\[\log L/\log\hat L\notin\mathbb Q.\]
For an odd $L_1$ above the threshold of admission, $L_1$ and $L_1+2$ are both admitted, and a common prime factor would divide $2$. Theorem~\YMMref{AJ}{673} and the source-energy bounds of Theorem~\YMMref{AD}{607} apply separately at both frames. The argument does not require equality of constants chosen after $L$.
\end{proof}

\begin{definition}[The Setting; the constants; the floors]\label{C-365-7}
\begin{enumerate}
\item[(1)] \emph{Frames.} Two frames, the second a partner of the first (Lemma~\ref{C-365-6}); a hat marks the objects of the second.
\item[(2)] \emph{Windows and reference sides.} In the first frame: $0<g\le\gamma$; a nonempty compact set $W\subset I(g)$; a number $g_1\in(0,g]$ with $W\subset\operatorname{int}I(g_1)$, which exists by Lemma~\ref{C-365-29}; a number $\varepsilon>0$ with $W_1=[\min W-3\varepsilon,\max W+3\varepsilon]\subset I(g_1)$; the interval $W'=[\min W-2\varepsilon,\max W+2\varepsilon]$; and a reference side $\ell'\in\mathfrak L(g_1)$ with $\ell'\ge2\mathsf l_1$. In the second frame: $0<\hat g\le\hat\gamma$; a compact interval $\hat W=[\hat w_-,\hat w_+]\subset\hat I(\hat g)$ of positive length; and a reference side $\hat\ell'\in\hat{\mathfrak L}(\hat g)$ with $\hat\ell'\ge2\hat{\mathsf l}_1$. All numbers of the first frame are formed at $g_1$ (so $C_g$ stands for $C_{g_1}$) with the reference side $\ell'$, and those of the second at $\hat g$ with $\hat\ell'$.
\item[(3)] \emph{The constants of the argument.}
\begin{itemize}
\item $C_\times(y)=2592\,[65y^3+2]\,C_SS(\kappa)$ for $y>0$, and $M_\times=1348488\,C_SS(\kappa)$.
\item $\mathsf Q_1$ is the constant of the transfer (Subsection~\ref{C-365-13}, item (M7)); it depends on $L$ and on $g_1$ and on none of $K$, $\beta$, $\ell$.
\item $\mathsf c_\chi=40\sqrt3ML$; $\mathsf Q_2=\exp[(C_S+C_g)\mathsf l_1^3+C_\times(1/8)]$; $\mathsf Q_*=2\max\{\mathsf Q_1,\mathsf Q_2\}+2$; $C_{\rm lg}=54L^3C_SS(\kappa)$; $C'_{\rm lg}=27(10m'+75)C_SS(\kappa)$;
\[h_1=\frac{\varrho}{250e\,\mathsf c_\chi\mathsf Q_*},\qquad\mathsf A_1=2C_{\rm lg}+\frac{L^3}{125}\bigl(4\log2+2C'_{\rm lg}\bigr).\]
\item $r_2=\varrho/(e^{(C_S+C_g)\ell'^3}+1)$; $C_\omega=2\log2/(\ell'^3r_2^2)$; $C_2'=\mathsf c^2\varrho^{-2}[1152\,C_\times(7/2)\,\mathsf l^{-3}+65536\,(C_S+C_g)]$; $r_1=\min\{\varrho/(256\mathsf c),\,r_2/2,\,h_1/2\}$;
\[C_I=C_I'+C_2'+C_\omega+2\mathsf A_1/h_1^2+M_\times T^{-3}\varrho^{-2},\qquad C_0=2C_I'+C_\omega+2\mathsf A_1/h_1^2;\]
$r_1^-$ and $C_I^+$ are $r_1$ and $C_I$ with $1+5\sqrt3$ in the place of $\mathsf c$ and $T_0$ in the place of $T$.
\item With $w_\pm$ the ends of $W_1$: $\lambda_-=\hat w_-/w_+$, $\lambda_+=\hat w_+/w_-$, and 
\[
\delta_{\max}=\min\{\varepsilon/2,\,|\hat W|/(48\lambda_+)\}.
\]

\item $k_7$ is the least integer $\ge1$ with $\varsigma_7=L^{-k_7}\le2/(3\sqrt3\max\{1,\lambda_+\})$; the \emph{tent} is $\chi(y)=\prod_\mu(1-|y_\mu|/\varsigma_7)_+$ and the \emph{plateau} $\eta(y)=\prod_\mu\min\{1,(3-|y_\mu|/\varsigma_7)_+\}$; $\mathsf c_7=40\sqrt3M/\varsigma_7$; $c_T=2\log2\,(2\mathsf Q_1+2)^2(2/\varrho)^2$; $c_N=\varsigma_7^{-3}c_T\mathsf c_7^2$; $\hat c_N=\hat L^3\hat c_T(40\sqrt3\hat M\hat L)^2$; $\hat{\mathsf c}_8=\max\{1,40\sqrt3\hat M/(\lambda_-\varsigma_7)\}$; 
\[
c_F=\sup_{K\ge k_7}\sup_{b\in W_1}|S^T_{2;K,\ell',L^Kb}(\chi,1-\eta)|;
\]
$\hat c_F$ is the bound that the convergence on fixed tori (item (M8)) gives in the second frame at the side $\hat\ell'$ for $n=2$, the interval $\hat W$, the spread $\lambda_-\varsigma_7$ and norms $\le\hat{\mathsf c}_8$; both are finite by (M8);
\[c_2=\varsigma_7^{-3}\bigl[c_T\mathsf c_7^2+\lambda_+^2\hat c_T\hat{\mathsf c}_8^2+c_F+\lambda_+^2\hat c_F\bigr],\qquad\hat{\mathfrak m}=16\hat B'/|\hat W|+\hat C_0|\hat W|,\]
\[C_\flat=\lambda_+^4\bigl[24\lambda_+\hat{\mathfrak m}/|\hat W|+10\lambda_+\hat c_N+1\bigr]+6c_N+3(C_0+c_2)+1.\]
\item $\mu_K=(D_K+\frac12C_0y^2)''$ on $\operatorname{int}I(g_1)$ (a positive measure by Proposition~\ref{C-365-21}). For $0<\delta\le\delta_{\max}$, $K_1(\delta)$ is the least integer $K_1\ge1$ such that $\mu_K(I)\le C_\flat(|I|+\delta)$ for all $K\ge K_1$ and all closed intervals or points $I\subset W'$; such integers exist by Proposition~\ref{C-365-28}, and $K_1(\delta)$ does not increase in $\delta$.
\end{itemize}
\item[(4)] \emph{The floors.}
\[q_0=\max\Bigl\{\frac{4C_\flat+C_I^+}{9(C_w^-)^2T_0^2},\ \bigl(3C_w^-T_0\min\{r_1^-,\varepsilon/2\}\bigr)^{-1/2}\Bigr\},\qquad\delta^\circ(q)=\min\Bigl\{\delta_{\max},\ \frac{3C_w^-T_0}{2C_\flat q}\Bigr\},\]
\[K_0(q)=K_1(\delta^\circ(q)),\qquad\ell_0(q)=\max\Bigl\{2q,\ \ell',\ \frac{12\ell'B'}{C_\flat\varepsilon^2},\ \frac{16\,C_\flat\ell'B'q^2}{3(C_wT)^2}\Bigr\}.\]
\item[(5)] \emph{Dependence.} All these numbers are formed from the two frames, $g_1$, $\hat g$, $W$, $\varepsilon$, $\hat W$, $\ell'$, $\hat\ell'$, and those that contain $\mathsf c$ or $T$ (namely $\mathsf c$, $C_w$, $C_2$, $C_2'$, $C_I$, $r_1$, $x_q$ and $\ell_0(\cdot)$) also from $T$; none depends on $K$, $\ell$, $b$, a profile, an offset or a state. $q_0$, $C_0$, $C_\flat$, $\delta_{\max}$, $K_1(\cdot)$ and $K_0(\cdot)$ do not depend on $T$, and $K_0(q)$ and $\ell_0(q)$ do not depend on each other. $K_0(q)$ does not decrease in $q$.
\end{enumerate}
\end{definition}

\subsection{Principal results}

\begin{theorem}[Cube means in every periodic Gibbs state, on large tori and in every continuum limit]\label{C-365-8}
Let Theorem~\YMMref{D}{8}, Theorems~\YMMref{AJ}{673} and~\YMMref{AD}{607} and the statements (M1)--(M12) among the construction results listed in Subsection~\ref{C-365-13} be granted at each of the two frames of a Setting (Definition~\ref{C-365-7}); by Lemma~\ref{C-365-6} every frame is the first frame of a Setting. Let $T\ge\max\{T_0,\mathsf l\}$. Then for every $q=L^i\ge\max\{2T,q_0\}$ the following hold.
\begin{enumerate}
\item[(a)] \emph{Periodic Gibbs states.} Let $K\ge K_0(q)$, $b\in W$ and $\nu,\nu'\in\mathcal G^{\rm per}_{K,L^Kb}$. For every board of $\mathbb L_K$ of side $q$, with any offset, and profiles of $\mathfrak P_T$ in distinct cells $c,c',c_1,\dots,c_n$ of it:
\[
\begin{multlined}
|\operatorname{Cov}_\nu(x_c,x_{c'})|\le C_2(T/q)^2,\qquad|\nu(x_c)-\nu'(x_c)|\le12e^2C_wT/q,\qquad\Bigl|\nu\Bigl(\prod_{i=1}^n\bigl(x_{c_i} \\
-\nu(x_{c_i})\bigr)\Bigr)\Bigr|\le\bigl(12e^2C_wT/q\bigr)^n.
\end{multlined}
\]
\item[(b)] \emph{Tori.} Let $K\ge K_0(q)$, $b\in W$ and $\ell\in\mathfrak L(g)$ with $\ell\ge\ell_0(q)$. For the law $\langle\cdot\rangle_{K,\ell,L^Kb}$ and every board of $T_\ell$ of side $q$ the first and the third bound of (a) hold.
\item[(c)] \emph{Limits of Gibbs sequences.} Let $(K_j,\beta_j,\nu_j)$ be a convergent Gibbs sequence at $g$ whose labels lie in $W$ from some index on (they need not converge), $\sigma$ its limit, $(\mathsf X,\lambda,(\sigma_\xi))$ a canonical decomposition of $\sigma$, $\sigma'=\sigma$ or $\sigma'=\sigma_\xi$ for any $\xi$, $\mathbf s$ the radial kernel of $\sigma'$ and $\mathcal H(\sigma',e_3)$ its reconstruction. Let $j_0$ be an index with $L^{j_0}\ge\max\{2T,q_0\}$, $r_0=\sqrt6L^{j_0}$ and $C_{\mathbf s}=6L^2(1-3T/(4L^{j_0}))^{-6}C_2$. Then:
\begin{enumerate}
\item[1.] $\mathfrak c_\xi=0$ for every $\xi$.
\item[2.] $S^{\sigma'}_2(\theta w_i,w_i)\le C_2(T/q_i)^2$ for all $i\ge j_0$.
\item[3.] There is exactly one positive Borel measure $\rho$ on $[0,\infty)$ with $\mathbf s(r)=\int e^{-mr}(4\pi r)^{-1}\,d\rho(m)$ for all $r>0$.
\item[4.] $\mathbf s(r)\le C_{\mathbf s}T^2/r^2$ for $r\ge r_0$; in particular $r\mathbf s(r)\to0$ as $r\to\infty$.
\item[5.] $|S^{\sigma'}_2(f,\tau_ah)|\le C_{\mathbf s}T^2\|f\|_1\|h\|_1(|a|-R_f-R_h)^{-2}$ for all $f,h\in C^1_c(\mathbb R^3)$ with supports in the balls of radii $R_f$, $R_h$ about $0$ and all $a$ with $|a|\ge R_f+R_h+r_0$.
\item[6.] $\rho([0,m])\le4\pi eC_{\mathbf s}T^2m$ for $0<m\le1/r_0$.
\item[7.] $\rho(\{0\})=0$.
\item[8.] For real $h\in C^1_c(\mathbb E_+)$: $\|e^{-tH}\psi_h\|^2\le C_{\mathbf s}T^2\|h\|_1^2/(4t^2)$ for $t\ge r_0/2$; $\psi_h$ is orthogonal to $\ker H$; $\int_0^\infty\sigma'(\theta h,\tau_{te_3}h)\,dt=\|H^{-1/2}\psi_h\|^2<\infty$; and $\langle\psi_h,1_{[0,\epsilon]}(H)\psi_h\rangle\le\frac{e^2}4C_{\mathbf s}T^2\|h\|_1^2\epsilon^2$ for $0<\epsilon\le2/r_0$.
\item[9.] Let $X=(f_1,\dots,f_n)\in\mathfrak S_n(\mathbb R^3)$ have its supports in the ball $B(0,R_X)$ and $h\in C^1_c(\mathbb R^3)$ its support in $B(0,R_h)$; let $\mathsf N(X)^2=\sup_u\sigma'(X\cdot\theta_uX)$, the supremum over the unit vectors $u$, $\theta_u$ the reflection in the plane $\{u\cdot x=R_X+1\}$; it is finite. Then for $|a|\ge R_X+R_h+2+r_0$
\[|\sigma'(X\cdot\tau_ah)|\le\frac{\mathsf N(X)\,\|h\|_1\,C_{\mathbf s}^{1/2}T}{2(|a|-R_X-R_h-2)},\qquad\bigl|S^{\sigma'}_{n+1}(f_1,\dots,f_n,\tau_ah)\bigr|\le\frac{C(X)\,\|h\|_1\,T}{|a|-R_X-R_h-2},\]
with a number $C(X)$ formed from $C_{\mathbf s}$ and from the values of $\sigma'$ at the subtuples of $X$ and at their concatenations with their reflections: a connected function with one argument moved away tends to $0$ at the rate $|a|^{-1}$, in every direction.
\item[10.] The converse and the measure that attains the bounds (Proposition~\ref{C-365-30}~(d), (e)): for the limit of every convergent Gibbs sequence and every component of it, with the constant $\kappa'$ and the measure $\rho$ of Proposition~\ref{C-365-30}, a bound $\sigma'(\theta w_i,w_i)\le C(T/q_i)^2$ for all large $i$, a bound $\mathbf s(r)=O(r^{-2})$ as $r\to\infty$, the pair $\kappa'=0$, $\rho([0,m])=O(m)$ as $m\to0$, and a bound $\mu_h(Q_r)\le A_\mu\|h\|_1^2r^2$ for small $r$ are equivalent; the measure $\rho=c_\rho\,dm$, whose kernel is $\mathbf s(r)=c_\rho/(4\pi r^2)$, satisfies them and attains 4 and 6 in order, so that these conclusions do not exclude it; and with no hypothesis on the labels the cube means of $\sigma'$ obey the bound of the order $1/q$ of Proposition~\ref{C-365-30}~(d), and $q$ times the quantity that it bounds tends to $0$ if and only if $\rho(\{0\})=0$.
\end{enumerate}
\item[(d)] \emph{Growing tori.} Let $S$ be a continuum limit at $g$ along growing tori, with no relation between $K_j$ and $\ell_j$, whose defining sequence has its labels in $W$ from some index on, and let $j_0$ be as in (c). Then:
\begin{enumerate}
\item[1.] $\sigma_S(\theta_\mu w_i^{(\mu)},w_i^{(\mu)})\le C_2(T/q_i)^2$ for all $i\ge j_0$ and $\mu=1,2,3$.
\item[2.] For real $h\in C^1_c(\mathbb E_+)$, in $\mathcal H(\sigma_S,e_3)$: 
\[
\mu_h(Q_r)\le12.46\,(1-3T/(4q_{j_0}))^{-6}L^2C_2\|h\|_1^2T^2r^2
\]
 for $0<r\le1/q_{j_0}$, and $\mathcal P(\{0\})\psi_h=0$.
\item[3.] For 
\[
X=(f_1,\dots,f_n)\in\mathfrak S_n(\mathbb R^3)
\]
 and $h\in C^1_c(\mathbb R^3)$: $\sigma_S(X\cdot\tau_ah)\to0$ and $S_{n+1}(f_1,\dots,f_n,\tau_ah)\to0$ as $\min_\mu|a_\mu|\to\infty$.
\end{enumerate}
\item[(e)] \emph{Every label and every reference coupling.}
\begin{enumerate}
\item[1.] For every $b_\infty>0$ there are an integer $s\ge0$ and $g_2\in(0,\gamma]$ with $L^sb_\infty\in\operatorname{int}I(g_2)$.
\item[2.] Let $s\ge0$ be an integer, $0<g_2\le\gamma$, $W^{(2)}\subset I(g_2)$ compact, and form the numbers of Definition~\ref{C-365-7} for a Setting built on $(g_2,W^{(2)})$ in the place of $(g,W)$, with the same second frame, and with the width $L^sT$. If $L^sb\in W^{(2)}$, $L^sT\ge\max\{T_0,\mathsf l\}$, $L^sq\ge\max\{2L^sT,q_0\}$ and $K-s\ge K_0(L^sq)$, then at $(K,L^Kb)$ the three bounds of (a) hold multiplied by $L^{8s}$, $L^{4s}$ and $L^{4ns}$ respectively, $T/q$ standing as it is; if besides $L^s\ell\in\mathfrak L(g_2)$ and $L^s\ell\ge\ell_0(L^sq)$, the bounds of (b) hold on $T_\ell$ multiplied likewise.
\item[3.] For every compact set $\mathsf B\subset(0,\infty)$ there are a constant $C_{\mathsf B}$ and floors $q_{\mathsf B}$, $K_{\mathsf B}(\cdot)$, $\ell_{\mathsf B}(\cdot)$, none of them depending on $T$, such that
\[|\operatorname{Cov}_\nu(x_c,x_{c'})|\le C_{\mathsf B}(T/q)^2,\qquad|\nu(x_c)-\nu'(x_c)|\le(2C_{\mathsf B})^{1/2}T/q\]
for all $b\in\mathsf B$, $T\ge\max\{T_0,\mathsf l\}$, $q=L^i\ge\max\{2T,q_{\mathsf B}\}$, $K\ge K_{\mathsf B}(q)$, $\nu,\nu'\in\mathcal G^{\rm per}_{K,L^Kb}$ and cells and profiles as in (a); and the covariance bound alone holds for the law $\langle\cdot\rangle_{K,\ell,L^Kb}$ on boards of $T_\ell$ for all sides $\ell=2L^m\ge\ell_{\mathsf B}(q)$.
\item[4.] At every $g'\in(0,\gamma]$, clauses (c) and (d) hold for every convergent Gibbs sequence at $g'$ and for every continuum limit at $g'$ along growing tori, with no condition on the labels, with $C_{\mathsf B}$ for $\mathsf B=I^{\rm lab}(g')$ in the place of $C_2$ and with $q_{\mathsf B}$ in the place of $q_0$.
\end{enumerate}
\end{enumerate}
\end{theorem}

\begin{theorem}[The curvature of the pressure at large cutoffs; mean actions; the second bound of differences of means]\label{C-365-9}
Assume the construction inputs listed in Theorem~\ref{C-365-8}, and let a Setting be given.
\begin{enumerate}
\item[(a)] \emph{No pile of specific heat.} For $0<\delta\le\delta_{\max}$ and $K\ge K_1(\delta)$: $\mu_K(I)\le C_\flat(|I|+\delta)$ for every closed interval or point $I\subset W'$; $|D_K''|(I)\le(C_\flat+C_0)(|I|+\delta)$; and
\[-C_0s^2\le\Delta^2_sD_K(b)\le C_\flat s^2+C_\flat\delta s\qquad\text{whenever }[b-s,b+s]\subset W'.\]
With $\delta_0(K)=\inf\{\delta\in(0,\delta_{\max}]:K_1(\delta)\le K\}$, which tends to $0$ as $K\to\infty$, these bounds hold with $\delta_0(K)$ in the place of $\delta$. For $\delta'>0$ let $\mathsf p=D_K*\phi_{\delta'}$, $\phi_{\delta'}(t)=\delta'^{-2}(\delta'-|t|)_+$; at every $b$ with $[b-\delta',b+\delta']\subset W'$: $-C_0\le\mathsf p''(b)\le C_\flat(1+\delta_0(K)/\delta')$ and $-C_0\delta'^2/12\le\mathsf p(b)-D_K(b)\le C_\flat(\delta'^2/12+\delta_0(K)\delta'/6)$; at $\delta'=\delta_0(K)>0$: $-C_0\le\mathsf p''\le2C_\flat$ and $|\mathsf p-D_K|\le\frac14\max\{C_\flat,C_0/3\}\delta_0(K)^2$.
\item[(b)] \emph{The atom at one label.} For $0<\delta\le\delta_{\max}$, $b\in W'$ and $K\ge K_1(\delta)$:
\[\mu_K(\{b\})=p'_{K,\infty}(b+)-p'_{K,\infty}(b-)=L^{4K}J(L^Kb)\le C_\flat\delta .\]
And $\beta^4J(\beta)\to0$ as $\beta\to\infty$ through the reals.
\item[(c)] \emph{Mean actions.} For every integer $K\ge0$, every $b>0$ and $\nu\in\mathcal G^{\rm per}_{K,L^Kb}$, the number $\nu(a_p)=\mathsf a(\nu)$ is the same for all plaquettes $p$; with $e(\nu)=3a_K^{-4}\mathsf a(\nu)$ we have $\nu(\mathcal O(f))=\operatorname{Rie}_K(f)\,e(\nu)$ for every bounded $f$ with compact support, and $p_{K,\infty}(b')\ge p_{K,\infty}(b)-(b'-b)e(\nu)$ for all $b'\ge0$. There are $\nu_\pm\in\mathcal G^{\rm per}_{K,L^Kb}$ with $-e(\nu_\pm)=p'_{K,\infty}(b\pm)$. Hence
\[\sup\{|e(\rho)-e(\nu)|:\rho,\nu\in\mathcal G^{\rm per}_{K,L^Kb}\}=J_K(b):=p'_{K,\infty}(b+)-p'_{K,\infty}(b-)=L^{4K}J(L^Kb),\]
and $J_K(b)=0$ if and only if all periodic Gibbs states at $(K,L^Kb)$ have the same mean action.
\item[(d)] \emph{The second bound of differences of means.} For every integer $K\ge0$, $b>0$, $\nu,\nu'\in\mathcal G^{\rm per}_{K,L^Kb}$, every board of $\mathbb L_K$ of side $q$ and every $f$ with $0\le f\le1$ that vanishes off a cube of side $\le q-a_K$ (in particular every profile, every $q$ and every cell): $|\nu(x_c)-\nu'(x_c)|\le J_K(b)$. In the Setting, $J_K(b)=\mu_K(\{b\})\le C_\flat\delta$ for $b\in W'$ and $K\ge K_1(\delta)$. So in Theorem~\ref{C-365-8}~(a) the second bound may be read as $\min\{12e^2C_wT/q,\ \mu_K(\{b\})\}$, and $\sup_{b\in W}|\nu(x_c)-\nu'(x_c)|$ tends to $0$ as $K\to\infty$ uniformly in $q$, in the cells and in the profiles.
\item[(e)] \emph{The plateau.} In any frame let $0<g\le\gamma$, $(K,\beta)\in\mathfrak d(g)$, $b=a_K\beta$, $\nu\in\mathcal G^{\rm per}_{K,\beta}$, and let $\mathcal I$ be the $\sigma$-field of the Borel sets of configurations fixed by all translations by $a_K\mathbb Z^3$. There is one bounded real $\mathcal I$-measurable function $\mathsf e$ with $\nu(ka_p)=\nu(k\mathsf e)$ for every plaquette $p$ and every bounded $\mathcal I$-measurable $k$. For $f,h\in C^1_c(\mathbb R^3)$, $\bar R\in a_K\mathbb N$ with $\operatorname{supp}f\cup\operatorname{supp}h\subset(-\bar R,\bar R)^3$ and $a_K\le2\min\{\delta((f);\bar R),\delta((h);\bar R)\}$, and $v\in a_K\mathbb Z^3$ with $|v|\ge\mathfrak k_1\bar R$:
\[\Bigl|S^T_2[\nu](f,\tau_vh)-9a_K^{-8}\operatorname{Rie}_K(f)\operatorname{Rie}_K(h)\operatorname{Var}_\nu(\mathsf e)\Bigr|\le\mathfrak k_0\,\mathsf N((f);\bar R)^{1/2}\,\mathsf N((h);\bar R)^{1/2}\,\frac{\bar R}{|v|},\]
where $\mathfrak k_0$, $\mathfrak k_1=4\sqrt3$, $\delta(\cdot;\bar R)$ and $\mathsf N(\cdot;\bar R)$ are the numbers of the plateau statement (Proposition~\YMMref{100.1}{744} and equations~\YMMref{(1976)--(1980)}{743}); and $0\le9a_K^{-8}\operatorname{Var}_\nu(\mathsf e)\le\frac14J_K(b)^2$, with equality at $\nu=\frac12(\nu_++\nu_-)$, while $\operatorname{Var}_\nu(\mathsf e)=0$ if $\nu$ is ergodic or $\nu=\nu_\pm$.
\item[(f)] \emph{A power of the spacing, under a rate in the cutoff.} Let two frames at one radius be given whose block sizes are $L_1$ and $L_2=L_1+2$, which exist by Lemma~\ref{C-365-6}; an index $i\in\{1,2\}$ marks the objects of the frame $i$. Let Theorems~\YMMref{AJ}{673} and~\YMMref{AD}{607} be granted at both frames, and Theorem~\YMMref{D}{8} at the choices of its constants that the two frames determine and at a lowered choice of each (Definition~\ref{def:rarity-choices}). In the frame $i$ let $\mathcal W_i=[\gamma_i^{-2}+\bar B_i+\frac12,\,L_i(\gamma_i^{-2}+3\bar B_i+1)-2\bar B_i]$, let $\chi^{(i)}(y)=\prod_\mu(1-L_i|y_\mu-\frac12|)_+$ and $\eta^{(i)}(y)=\prod_\mu\min\{1,(3-L_i|y_\mu-\frac12|)_+\}$ be the tent and the plateau of half-width $1/L_i$ about the point $(\frac12,\frac12,\frac12)$, and for integers $m,K\ge1$, $\ell=2L_i^m$ and $b>0$ let
\[\mathsf y^{(i)}_{m,K}(b)=b^5L_i^3\,S^T_{2;K,\ell,L_i^Kb}\bigl((\chi^{(i)})^{\langle\ell\rangle},1-(\eta^{(i)})^{\langle\ell\rangle}\bigr)\]
be the far part of the specific heat of $T_\ell$. Let the far part settle in the cutoff at a rate on every torus: there are numbers $\mathsf k_i\ge0$ and $\mathsf c_{{\rm y},i}<\infty$, free of $m$, $K$ and $b$, such that for all integers $m\ge1$, $\ell=2L_i^m$, all $b\in\mathcal W_i$ and all integers $K'\ge K\ge\mathsf k_i+54m$
\[\bigl|\mathsf y^{(i)}_{m,K'}(b)-\mathsf y^{(i)}_{m,K}(b)\bigr|\le\mathsf c_{{\rm y},i}\,\ell^9K^{\mathsf p_{2,i}}L_i^{-K/10},\qquad\mathsf p_{2,i}=6\mathfrak p_{0,i}+\tfrac12,\]
$\mathfrak p_{0,i}$ the exponent of the theorem on two consecutive cutoffs (Theorem~\YMMref{AR}{810}) in the frame $i$. Then for every real $A>90$ there are numbers $\mathsf C_{\rm r}$ and $\beta_1$ with $\beta^4J(\beta)\le\mathsf C_{\rm r}\beta^{-1/(2A)}$ for all $\beta\ge\beta_1$. Hence for every $0\le\alpha<1/180$ the number $\mathsf c_*(\alpha)=\sup_{\beta\ge1}\beta^{4+\alpha}J(\beta)$, a number of the Wilson model on $\mathbb Z^3$, is finite; and in every frame, for $0<g\le\gamma$, $(K,\beta)\in\mathfrak d(g)$ and $\rho,\nu\in\mathcal G^{\rm per}_{K,\beta}$:
\[3a_K^{-4}|\rho(a_p)-\nu(a_p)|\le a_K^{-4}J(\beta)\le g^{8+2\alpha}\,\mathsf c_*(\alpha)\,L^{-\alpha K};\]
$|\rho(\mathcal O_K(f))-\nu(\mathcal O_K(f))|\le a_K^{-4}J(\beta)(2R+a_K)^3\|f\|_\infty$ for every bounded $f$ with support in $[-R,R]^3$; and $|\rho(x_c)-\nu(x_c)|\le g^{8+2\alpha}\mathsf c_*(\alpha)L^{-\alpha K}$ for every cube mean, free of $q$ and of $T$, with no floor in the cutoff.
\end{enumerate}
\end{theorem}

The hypothesis K1 is the following inequality between the free energies of two periodic boxes. Its data are a label $\mathsf b>0$, a side $\ell_\circ=2L^{m_\circ}$ with $m_\circ\in\mathbb Z$, an integer $K_\circ$ with $K_\circ+m_\circ\ge1$ and a number $0<x_\circ\le0.069$; with $Z(n_1,n_2,n_3;\beta)$ the partition function of Wilson's action on the periodic lattice with $n_1,n_2,n_3$ sites a side, $l_K=\ell_\circ/(2a_K)=L^{K+m_\circ}$ and $\beta_K=L^K\mathsf b$, the hypothesis $\mathrm{K1}[\mathsf b,\ell_\circ,K_\circ;x_\circ]$ requires
\[2\log Z(l_K,2l_K,2l_K;\beta_K)-\log Z(2l_K,2l_K,2l_K;\beta_K)\le x_\circ\qquad\text{for every integer }K\ge K_\circ .\]
We put $\theta_\circ=12.9\,x_\circ\le0.8901$ and $m_*=2\log(1/\theta_\circ)/\ell_\circ\ge0.2328/\ell_\circ$.

\begin{theorem}[The two-point function of the action density under the hypothesis K1 at one label]\label{C-365-10}
Assume the construction inputs listed in Theorem~\ref{C-365-8}; let $0<g\le\gamma$, $\mathsf b\in I(g)$, and assume $\mathrm{K1}[\mathsf b,\ell_\circ,K_\circ;x_\circ]$.
\begin{enumerate}
\item[(a)] \emph{The gap} (from (M13) of Subsection~\ref{C-365-13}). For each $K\ge\max\{K_\circ,1\}$ there is exactly one periodic Gibbs state $\nu_K$ at $(K,L^K\mathsf b)$; the Gibbs sequence $(K,L^K\mathsf b,\nu_K)_{K\ge\max\{K_\circ,1\}}$ converges; its limit $\sigma$ is its own single component, has a unique vacuum $\Omega$ and is invariant under $\mathbb R^3\rtimes\mathrm O(3)$; the limits at this label in every order of cutoff and volume, that is, the limits of $S^T_{n;K,\ell,L^K\mathsf b}(X^{\langle\ell\rangle})$ for $n\ge2$ and $X\in\mathfrak S_n(\mathbb R^3)$, at the couplings $L^K\mathsf b$ themselves and with $\ell$ running through the sides $2L^m$, as $K$ and $\ell$ tend to infinity jointly in any relation, or $\ell$ first and then $K$, or $K$ first and then $\ell$, are the connected functions $S^\sigma_n(X)$; in particular every continuum limit along growing tori at this label, that is, with a defining sequence in which $\beta_j=L^{K_j}\mathsf b$ from an index on, has the centred functional $\sigma$; the spectrum of $H$ in $\mathcal H(\sigma,e_3)$ lies in $\{0\}\cup[m_*,\infty)$, and for $X,Y\in\mathfrak S(\mathbb E_+)$ and $t\ge0$
\[\bigl|\sigma(\theta X\cdot\tau_{te_3}Y)-\sigma(X)\sigma(Y)\bigr|\le e^{-m_*t}\,\|(1-P_\Omega)[X]\|\,\|(1-P_\Omega)[Y]\|,\]
$P_\Omega$ the projection on $\mathbb C\Omega$, and likewise in every direction.
\item[(b)] \emph{The kernel and the spectral measure.} Let $\mathbf s$ and $\rho$ be the radial kernel of $\sigma$ and its measure of Theorem~\ref{C-365-8}~(c)\,3, which applies to $\sigma$ by Theorem~\ref{C-365-8}~(e)\,4. Then $\mathbf s(r)\le e^{-m_*(r-r')}\mathbf s(r')$ for $0<r'\le r$; $\rho([0,m_*))=0$; and $\mathbf s(r)\le(r'/r)\,e^{-m_*(r-r')}\,\mathbf s(r')$ for $0<r'\le r$. So $\mathbf s(r)=O(r^{-3})$ as $r\to\infty$ and $\rho([0,m])=O(m^2)$ as $m\to0$.
\item[(c)] \emph{Cube means in the limit.} For $T\ge\max\{T_0,\mathsf l\}$, $q=L^i\ge2T$, a board of $\mathbb R^3$ of side $q$ (a partition of $\mathbb R^3$ into the cubes $\square_c$ of side $q$ centred at the points $o+qc$, $c\in\mathbb Z^3$, with an offset $o\in\mathbb R^3$) and profiles $f_c\in\mathfrak P_T(\square_c)$, $f_{c'}\in\mathfrak P_T(\square_{c'})$ of distinct cells:
\[\bigl|S^\sigma_2(q^{-3}f_c,q^{-3}f_{c'})\bigr|\le\chi_\circ\,q^{-3},\qquad\chi_\circ=4\pi\int_{T/4}^\infty r^2\mathbf s(r)\,dr<\infty;\]
if the closures of the two cells have no common point, that is $|c-c'|_\infty\ge2$, the supports of the two profiles are at distance $\ge q+T/4$, and the left member is at most $\mathbf s(q)$, which is exponentially small in $q$ by (b).
\end{enumerate}
\end{theorem}

\begin{proposition}[Two fixed tori: means, and the convergence of the difference of the pressures]\label{C-365-11}
Assume the construction inputs listed in Theorem~\ref{C-365-8}; let one frame and $0<g\le\gamma$ be given.
\begin{enumerate}
\item[(a)] \emph{An exact identity.} For $K\ge1$, a side $\lambda=2L^m$, $b\ge0$ and a bounded $f$ with compact support on $\mathbb R^3$: $\langle\mathcal O(f^{\langle\lambda\rangle})\rangle_{K,\lambda,L^Kb}=-\operatorname{Rie}_K(f)\,p'_{K,\lambda}(b)$. For two sides $\ell>\ell'$,
\[A_K(f;b):=\langle\mathcal O(f^{\langle\ell\rangle})\rangle_{K,\ell,L^Kb}-\langle\mathcal O(f^{\langle\ell'\rangle})\rangle_{K,\ell',L^Kb}=-\operatorname{Rie}_K(f)\,D'_{K,\ell}(b),\]
where $D'_{K,\ell}(b)=3a_K^{-4}\bigl[\langle a_p\rangle_{K,\ell',L^Kb}-\langle a_p\rangle_{K,\ell,L^Kb}\bigr]=-L^3A_K(\chi_0;b)$, $\chi_0(y)=\prod_\mu(1-L|y_\mu|)_+$.
\item[(b)] \emph{Bounded, uniformly in the cutoff and in both sides.} For $b\in I(g)$, $K\ge1$, $\ell,\ell'\in\mathfrak L(g)$ and $f$ with support in a ball of radius $\le2$: $|A_K(f;b)|\le2\log2\,(2\mathsf Q_1+2)\varrho^{-1}\|f\|_\sharp$; so $|D'_{K,\ell}(b)|\le2\log2\,(2\mathsf Q_1+2)\varrho^{-1}L^3\mathsf c_\chi$, free of $K$, $b$, $\ell$, $\ell'$. For a cube mean whose cell lies in a ball of radius $\le\ell'/8$, its expectations on the two tori differ by $-q^{-3}\operatorname{Rie}_K(f)D'_{K,\ell}(b)$, and $0\le q^{-3}\operatorname{Rie}_K(f)\le1$.
\item[(c)] \emph{Convergence of the difference, with a rate.} Let $\ell'<\ell$ be two sides of $\mathfrak L(g)$, $\mathsf J\subset I(g)$ a compact interval and $U_K=p_{K,\ell}-p_{K,\ell'}$. For $b\in\mathsf J$ and $K\ge\max\{K_{\rm bf}(\ell,\mathsf J),K_{\rm bf}(\ell',\mathsf J)\}$:
\[
\begin{multlined}
|U_{K+1}(b)-U_K(b)|\le2\,\mathsf C_{\rm bf}(\mathsf J)\,K^{\mathsf p_0}L^{-K/10}+2\mathsf v_K,\qquad\mathsf v_K=2C_{\rm tt}E_0\sum_{s\ge u}L^{3s}\exp\Bigl[ \\
-\frac{\kappa\ell'L^s}{5M}\Bigr],\quad u=\lceil K/1260\rceil,
\end{multlined}
\]
with $\mathsf p_0=2\mathfrak p_0+\frac12$; here the number $\mathsf C_{\rm bf}(\mathsf J)$, formed from the constants of the frame and the ends of $\mathsf J$ and free of the sides, the exponent $\mathsf p_0$ and the thresholds $K_{\rm bf}(\cdot,\mathsf J)$ are those of the theorem on two consecutive cutoffs (Theorem~\YMMref{AR}{810}, Definition~\textup{\ref{IV-62-62}}~(5), Corollary~\textup{\ref{IV-62-66}}~(a)), and $C_{\rm tt}$, $E_0$ those of the comparison of the vacuum energies of two tori (Lemma~\ref{lem:vacuum-prefix}). So $U_K$ converges as $K\to\infty$, uniformly on $\mathsf J$.
\item[(d)] \emph{Convergence of the derivative.} Let moreover $\ell'\ge2\mathsf l_1$, and put $r_3=\min\{h_1,\varrho/(e^{(C_S+C_g)\ell^3}+1)\}$ and $M_3=\mathsf A_1+(\ell^{-3}+\ell'^{-3})\log2$, both free of $K$ and of $b$. Then $|\partial_b^nU_K(b)|\le n!\,M_3r_3^{-n}$ for all $b\in I(g)$, $K\ge1$ and $n\ge1$. Let $\mathsf W=[b_-,b_+]$ with $\mathsf J=[b_--r_3,b_++r_3]\subset I(g)$, and $\mathsf o_K=4\mathsf C_{\rm bf}(\mathsf J)K^{\mathsf p_0}L^{-K/10}+4\mathsf v_K$. For $K\ge\max\{K_{\rm bf}(\ell,\mathsf J),K_{\rm bf}(\ell',\mathsf J)\}$
\[\sup_{\mathsf W}|U'_{K+1}-U'_K|\le\frac{4\mathsf o_K}{r_3}\Bigl[1+\log_2^+\frac{2M_3}{\mathsf o_K}\Bigr].\]
Hence $\sum_KK\sup_{\mathsf W}|U'_{K+1}-U'_K|<\infty$, $U'_K$ converges uniformly on $\mathsf W$, and $\sup_K\sup_{\mathsf W}|U'_K|<\infty$; and for every real Lipschitz $f\ne0$ with support in a ball of radius $\le\ell'/8$, the difference between the two tori of the derivatives at zero source of the contact terms that the theorem on two consecutive cutoffs splits off at the cutoffs $K$ and $K+1$ is summable in $K$, uniformly on $\mathsf W$.
\end{enumerate}
\end{proposition}

\begin{proposition}[Variances, the cost of a source on the walls, and covariances]\label{C-365-12}
Assume the construction inputs listed in Theorem~\ref{C-365-8}; let one frame, $0<g\le\gamma$, $T\ge\max\{T_0,\mathsf l\}$ and $q=L^i\ge2T$ be given.
\begin{enumerate}
\item[(a)] \emph{Variances from below.} Let $b\in I(g)$, $K\ge1$, a board of $\mathbb L_K$ of side $q$ with any offset, a profile $f$ of one of its cells $c$, and $X_\kappa=\mathcal O(f^{[\kappa]})$ for the cells $\kappa$. For $\ell\in\mathfrak L(g)$, $\ell\ge2q$, with the board of $T_\ell$ of the same offset, let $\mathsf e_\ell=N^{-1}\partial_s^2\mathcal E^{\rm av}_\ell(s(F_e-F_o))|_{s=0}$ and $\bar{\mathsf e}=\limsup_{\ell\to\infty}\mathsf e_\ell$, which is finite. Then for every $\nu\in\mathcal G^{\rm per}_{K,L^Kb}$ and real numbers $y_\kappa$, finitely many of them non-zero,
\[\operatorname{Var}_\nu\Bigl(\sum_\kappa y_\kappa X_\kappa\Bigr)\ge(\bar{\mathsf e}-2C_I'q^3)\sum_\kappa y_\kappa^2 .\]
Consequently $\mathsf V_\nu(\kappa):=\operatorname{Var}_\nu(X_\kappa)-\bar{\mathsf e}\ge-2C_I'q^3$, and $|\operatorname{Cov}_\nu(X_\kappa,X_{\kappa'})|\le\frac12[\mathsf V_\nu(\kappa)+\mathsf V_\nu(\kappa')]+2C_I'q^3$ for $\kappa\ne\kappa'$.
\item[(b)] \emph{The cost of a source on the walls controls the cube means.} Let $b\in I(g)$, $K\ge1$, $\ell\in\mathfrak L(g)$ with $\ell\ge2q$, a board of $T_\ell$ of side $q$, a profile $f$ of a cell $c$, and for $0<r\le\varrho/2$
\[\mathsf S_f(r)=N^{-1}\log\frac{\widetilde Z^{\rm av}(rF)\,\widetilde Z^{\rm av}(-rF)}{\widetilde Z^{\rm av}(0)^2}.\]
Then $\mathsf S_f(r)\ge-2C_I'q^3r^2$, and there is a real number $m$ with
\[\|x_c-m\|_{\rm ch}\le(rq^3)^{-1}\exp\bigl\{\tfrac12\mathsf S_f(r)+(2C_I'+M_\times T^{-3}\varrho^{-2})q^3r^2\bigr\}.\]
So for profiles $f$, $f'$ of distinct cells $c$, $c'$: $|\operatorname{Cov}_{K,\ell,L^Kb}(x_c,x_{c'})|\le2(rq^3)^{-2}\exp\{\frac12\mathsf S_f(r)+\frac12\mathsf S_{f'}(r)+2(2C_I'+M_\times T^{-3}\varrho^{-2})q^3r^2\}$; in particular, if $\mathsf S_f(r),\mathsf S_{f'}(r)\le\mathsf C$ at $r=\hat rq^{-3+p/2}\le\varrho/2$ with $p\in[2,3]$ and $0<\hat r\le1$, then $|\operatorname{Cov}_{K,\ell,L^Kb}(x_c,x_{c'})|\le2\hat r^{-2}e^{\mathsf C+2(2C_I'+M_\times T^{-3}\varrho^{-2})\hat r^2}q^{-p}$.
\item[(c)] \emph{The cost from above.} In a Setting, with the numbers at $g_1$: for $b\in I(g_1)$, $K\ge1$, $\ell\in\mathfrak L(g_1)$ with $\ell\ge\max\{2q,\ell'\}$, and $0<r\le\min\{r_1,b\}$,
\[\mathsf S_f(r)\le q^3\Delta^2_rD_{K,\ell}(b)+6C_wTq^2r+2C_Iq^3r^2 .\]
In particular, for $b\in W$, $\ell\ge\max\{2q,\ell',6\ell'B'q^3\}$, $0<\hat r\le1$ with $\hat rq^{-3/2}\le\min\{r_1,\varepsilon/2\}$, and $K\ge K_1(\min\{\delta_{\max},q^{-3/2}\})$:
\[-2C_I'\hat r^2\le\mathsf S_f(\hat rq^{-3/2})\le6C_wT\hat r\,q^{1/2}+2C_I\hat r^2+2C_\flat+1 .\]
\item[(d)] \emph{A power of $q$ passes to the kernel.} Let $0<g'\le\gamma$, $p\in[2,3]$, and let a convergent Gibbs sequence at $g'$ satisfy the hypothesis of Lemma~\ref{C-365-32} with numbers $C_{\rm LH}$, $c_{\rm LH}$, $q_{\rm LH}$ and, besides, for every $q=L^i\ge\max\{2T,q_{\rm LH}\}$ and $j\ge j(q)$, $|\operatorname{Cov}_\nu(x_c,x_{c'})|\le C_pq^{-p}$ for all $\nu\in\mathcal G^{\rm per}_{K_j,\beta_j}$ and profiles in two distinct cells of one board of side $q$. Then for $\sigma$ and every $\sigma_\xi$, with $L^{j_0}\ge\max\{2T,q_{\rm LH}\}$ and $C_p'=(\sqrt6L)^p(1-3T/(4L^{j_0}))^{-6}C_p$: $\mathbf s(r)\le C_p'r^{-p}$ for $r\ge\sqrt6L^{j_0}$, and $\rho([0,m])\le4\pi eC_p'm^{p-1}$ for $0<m\le1/(\sqrt6L^{j_0})$.
\item[(e)] \emph{One frame and a growth bound of the specific heat.} Let $W\subset\operatorname{int}I(g)$ be a compact interval, $\ell'\in\mathfrak L(g)$ with $\ell'\ge2\mathsf l_1$, $C_{\rm FS}=3\ell'B'$, and suppose that there are $Y_0$ and integers $K_Y(\ell)$ with $D''_{K,\ell}\le Y_0\ell^{1/3}$ on $W$ for all $\ell\in\mathfrak L(g)$ with $\ell\ge\ell'$ and all $K\ge K_Y(\ell)$. Let $q\ge q_1=\max\{2T,C_I/(9C_w^2T^2),(3C_wTr_1)^{-1/2}\}$ with $6C_{\rm FS}q^3\ge\max\{2q,\ell',2L^{m_{\mathfrak L}(g)}\}$, let $\ell_q$ be the least side of $\mathfrak L(g)$ with $\ell_q\ge6C_{\rm FS}q^3$, $K\ge K_Y(\ell_q)$ and $[b-x_q,b+x_q]\subset W$. Then for all $\nu,\nu'\in\mathcal G^{\rm per}_{K,L^Kb}$ and profiles of $\mathfrak P_T$ in distinct cells $c,c',c_1,\dots,c_n$ of one board of side $q$: $|\operatorname{Cov}_\nu(x_c,x_{c'})|\le2e^4\delta^2$, $|\nu(x_c)-\nu'(x_c)|\le2e^2\delta$ and $|\nu(\prod_i(x_{c_i}-\nu(x_{c_i})))|\le(2e^2\delta)^n$, with
\[\delta\le\frac Tq\Bigl[4C_w+\frac{(6LC_{\rm FS})^{1/3}Y_0}{3C_wT^2}\Bigr]+C_0x_q .\]
\item[(f)] \emph{The cost bounds the variances in every periodic Gibbs state.} Let $b\in I(g)$, $K\ge1$, a board of $\mathbb L_K$ of side $q$, a profile $f$ of one of its cells $c$, $X=\mathcal O(f)$, and sides $\ell_k\in\mathfrak L(g)$ with $\ell_k\ge2q$ and $\ell_k\to\infty$ be given, the boards of $T_{\ell_k}$ having the same offset; let $\bar{\mathsf e}$ be the number of (a), $\mathsf S_{f,\ell_k}(r)$ the cost of (b) on $T_{\ell_k}$, $c_X=M_\times q^3T^{-3}$ the constant of the cross terms between the cells, and $c_3=3C_I'q^3+4c_X/(9\varrho^2)$. Then for every $\nu\in\mathcal G^{\rm per}_{K,L^Kb}$ and $0<r\le\varrho/2$
\[\operatorname{Var}_\nu(X)\le\bar{\mathsf e}+2c_X\varrho^{-2}+2r^{-2}\exp\Bigl[\limsup_{k\to\infty}\mathsf S_{f,\ell_k}(r)+c_3r^2\Bigr].\]
If, with $\mathsf V_\nu$ as in (a), $\mathsf V_\nu(\kappa)\le C_{\rm var}q^{6-p}$ and $\mathsf V_\nu(\kappa')\le C_{\rm var}q^{6-p}$ at two distinct cells $\kappa,\kappa'$, with $p\in[2,3]$, then $|\operatorname{Cov}_\nu(q^{-3}X_\kappa,q^{-3}X_{\kappa'})|\le(C_{\rm var}+2C_I')q^{-p}$; at $p=3$ this is the bound $(C_{\rm var}+2C_I')q^{-3}$.
\item[(g)] \emph{A one-sided bound on the walls bounds the cost.} In the situation of (b) let
\[\Upsilon_f(x)=N^{-1}\bigl[\log\widetilde Z^{\rm av}(xF)-\log\widetilde Z^{\rm av}(x)\bigr],\qquad\upsilon_f(x)=N^{-1}\partial_t\log\widetilde Z^{\rm av}\bigl(xF+t(1-F)\bigr)\big|_{t=0},\]
the second being the derivative at the source $xF$ in the direction of the walls. Then $\Upsilon_f(x)\le-x\,\upsilon_f(x)+C_I'q^3x^2$ for real $|x|\le\varrho/2$. In a Setting, with the numbers at $g_1$, let $b\in W$, $\ell\in\mathfrak L(g_1)$ with $\ell\ge\max\{2q,\ell'\}$, $0<\delta\le\delta_{\max}$ and $K\ge K_1(\delta)$, and let $\Upsilon_f(x)\le\mathsf s_1x+C_\Upsilon$ for all real $|x|\le\bar x$, with real numbers $\mathsf s_1$, $C_\Upsilon$ and $\bar x>0$. Then for $0<r\le\min\{\bar x,r_1,\varepsilon/2\}$
\[\mathsf S_f(r)\le2C_\Upsilon+q^3\bigl[C_\flat r^2+C_\flat\delta r+2\bar\eta(\ell)\bigr]+2(C_\omega+2\mathsf A_1/h_1^2)q^3r^2 .\]
\end{enumerate}
\end{proposition}

\subsection{Construction estimates}\label{C-365-13}

The proofs use Theorem~\YMMref{D}{8} and Theorems~\YMMref{AJ}{673} and~\YMMref{AD}{607} among the construction results and the following statements (M1)--(M12) among the construction results, each taken at both frames of the Setting. Theorems~\YMMref{D}{8}, \YMMref{AJ}{673} and~\YMMref{AD}{607} are granted together with these statements, several of which rest on them: among them, the convergence on fixed tori (M8) and the theorem on two consecutive cutoffs in (M12) rest on Theorem~\YMMref{D}{8}, and the theorems on limits (M10) and Proposition~\YMMref{106.3}{765} in (M9) on Theorems~\YMMref{D}{8}, \YMMref{AJ}{673} and~\YMMref{AD}{607}.

Statements among the construction results that are not cited by number are cited by the names they carry there. Theorem~\ref{C-365-10} uses besides (M13), and clause (f) of Theorem~\ref{C-365-9} uses (M14). (M13) is proved in Section~\ref{C-361-1} on the hypothesis K1, under that hypothesis. (M14) is proved in Section~\ref{lh:sec} on the latent heat of the Wilson model: its power is derived there from Theorems~\YMMref{AJ}{673} and~\YMMref{AD}{607} at the two frames of (M14), Theorem~\YMMref{D}{8} at their choices and at a lowered choice of each, and the rate at which the far part of the specific heat settles in the cutoff on every torus, the rate that clause (f) of Theorem~\ref{C-365-9} assumes and that (M14) carries among its hypotheses: for $i=1,2$, numbers $\mathsf k_i\ge0$ and $\mathsf c_{{\rm y},i}<\infty$, free of $m$, $K$ and $b$, with $|\mathsf y^{(i)}_{m,K'}(b)-\mathsf y^{(i)}_{m,K}(b)|\le\mathsf c_{{\rm y},i}\,\ell^9K^{\mathsf p_{2,i}}L_i^{-K/10}$ for all integers $m\ge1$, $\ell=2L_i^m$, all $b\in\mathcal W_i$ and all integers $K'\ge K\ge\mathsf k_i+54m$.

The list is followed by the four statements (i)--(iv) that the proof of Proposition~\ref{C-365-11}~(c) derives from (M12).
\begin{itemize}
\item[(M1)] \emph{Source energies.}
Theorem~\YMMref{AD}{607}, Proposition~\YMMref{83.2}{675}, and
Proposition~\YMMref{74.4}{617} give the following at
$(K,\beta)\in\mathfrak d(g)$ and $\ell\in\mathfrak L(g)$.
Each $e_j(X;\cdot)$ is holomorphic along finite families
$z\mapsto w_0+\sum_qz_qf_q$, vanishes at zero, depends on the source
only on $X$ with one layer of its level cubes, obeys conjugation
reality, and has modulus at most $C_Se^{-\kappa d_j(X)}$.
Moreover,
\[
 |\widetilde Z_\ell(w)|\le e^{C_S\ell^3},\qquad
 \widetilde Z_\ell(0)\ge e^{-C_g\ell^3},\qquad
 |E_\ell\ell^{-3}-E_{\ell'}\ell'^{-3}|\le C_S.
\]
The vacuum splitting is equation~\YMMref{(1737)}{675}.
The translation law in Theorem~\YMMref{AJ}{673}, summed in
equation~\YMMref{(1736)}{675}, makes $\mathcal E_\ell$ invariant
under translations by $M\mathbb Z^3$.
\item[(M2)] \emph{The stability functional.}
Theorem~\YMMref{AJ}{673} gives the finite history decomposition,
its source-independent reference measures and evaluation domains,
and one null set outside which both exponents are defined for all
sources. Its equation~\YMMref{(1728)}{673} is the density
$\mathsf{Stay}_h^{\rm T}\chi_K(\Omega_K)
 e^{\boldsymbol\Phi_h^w+A_{K,h}^w}$.
The local, boundary and retained terms are holomorphic in the
source on these domains and obey conjugation reality.
The real construction supplies the nonnegative Borel test factors
and the reference measure; at a real source, the resulting
history measure is positive because it is the pushforward of the
positive bare density.
The bare action is linear in its weight:
\[
 \rho_0^w(U)=e^{-A(\beta-L^Kw,U)}
    =e^{-\beta\sum_pa_p+\mathcal O(w)},\qquad
 A(c,U)=\sum_pc(x(p))a_p(U),\qquad Z[w]=\int\rho_0^w\,dU.
\]
The source-independent coordinate maps and the measure identity
are those of Theorem~\YMMref{AI}{664} and Lemma~\YMMref{82.5}{670}.
Proposition~\YMMref{83.2}{675} identifies $E_\ell=E_{(K)}$,
the source energies in equation~\YMMref{(1736)}{675}, and
$\widetilde Z_\ell=\widetilde Z_K$ in
equation~\YMMref{(1737)}{675}. Its
equation~\YMMref{(1738)}{675} bounds the sum of the normalized
total variations of the history measures. This stronger estimate,
not merely a bound on $|\widetilde Z|$, is used below.
\item[(M3)] \emph{The count of domains.} For a domain $X$ of level $j$ and depth $s=K-j$, and sets $Y,Y'$: (a) if $X$ with one layer of cubes meets $Y$ and $Y'$, then $d_j(X)\ge L^s\operatorname{dist}_\infty(Y,Y')/M-4$; (b) for $\kappa'>\log9216$ and real $D'$, the sum of $e^{-\kappa'd_j(X)}$ over the $X$ that with one layer meet $Y$ and have $d_j(X)\ge D'$ is $\le27N_s(Y)S(\kappa')e^{-(\kappa'-\log9216)[D']_+}$; (c) $N_s(Y)\le(\ell L^s/M)^3$, and $N_s(Y)\le(2RL^s/M+2)^3$ if $Y$ lies in a sup-ball of radius $R$, $N_s(Y)$ being the number of the cubes of side $ML^{-s}$ of the standard partition that meet $Y$.
\item[(M4)] \emph{The isometries.} The translations by $a_K\mathbb Z^3$ and the reflections carry the Wilson weight and the Haar measure into themselves ($\operatorname{tr}g^{-1}=\operatorname{tr}g$ on $\mathrm{SU}(2)$), and $\mathcal O(\mathsf rw)=\mathcal O(w)\circ\mathsf r^{-1}$; so $Z[\mathsf rw]=Z[w]$ and $\|\mathsf rw\|_\sharp=\|w\|_\sharp$.
\item[(M5)] \emph{The passage and labels.}
For $b\in I(g)$ and $K\ge1$, the pair $(K,L^Kb)$ belongs to
$\mathfrak d(g)$, by Corollary~\YMMref{74.2}{616}.
Proposition~\YMMref{74.1}{615} gives the label bounds, in particular
$a_K\beta\ge g^{-2}$; their fixed compact interval is
$I^{\rm lab}(g)$, equation~\YMMref{1854}{710}.
\emph{The sides.} The family $\mathfrak L(g)$ is closed upward,
and $\mathfrak L(g')\subset\mathfrak L(g)$ for $0<g'\le g$,
as follows from its lower bound on the side.
\item[(M6)] \emph{The chessboard estimate}, at every $\beta\ge0$, on boards of every side $\mathsf j=L^i\le\ell/2$ (Definition~\ref{C-365-4}~(2)): $\|cF\|_{\rm ch}=|c|\,\|F\|_{\rm ch}$, $\|1\|_{\rm ch}=1$, $\|F\|_{\rm ch}\le\sup|F|$, the triangle inequality, and $|\langle\prod_sF_s\rangle|\le\prod_s\|F_s\|_{\rm ch}$ for functions $F_s$ of different cells. \emph{The source formula}: for a complex source $w$ that vanishes outside the interior of a cell, $\|e^{\mathcal O(w)}\|_{\rm ch}^N=Z_\beta[w^{\rm ch}]/Z_\beta\ge0$, $w^{\rm ch}$ the sum of the mirror copies of $w$ (conjugated in the cells of the other parity). \emph{The domination}: for $K\in\mathbb N$, $\beta>0$, $\nu\in\mathcal G^{\rm per}_{K,\beta}$, a board of $\mathbb L_K$ of any side $\ge La_K$ and any offset, functions $F_{\mathbf k}$ of finitely many different cells and every sequence of sides $\ell_k=2L^{m_k}\to\infty$: $|\nu(\prod F_{\mathbf k})|\le\liminf_k\prod\|F_{\mathbf k}\|_{{\rm ch},\ell_k}$.
\item[(M7)] \emph{The transfer.} There is $\mathsf Q_1$, depending on $L$ and on $g$ and on none of $K$, $\beta$, $\ell$, such that for $(K,\beta)\in\mathfrak d(g)$, all $n\ge1$, all $f_1,\dots,f_n\in\operatorname{Lip}_c(\mathbb R^3)$ with supports in one ball of radius $2$, in any mutual position, and any two of the laws $\nu\in\mathcal G^{\rm per}_{K,\beta}$ and $\langle\cdot\rangle_{K,\ell,\beta}$, $\ell\in\mathfrak L(g)$ (the latter carried to $\mathbb L_K$, so that its cumulants at $f$ are $S^T_{n;K,\ell,\beta}(f^{\langle\ell\rangle})$), the cumulants of order $n$ differ by at most $2\log2\,(2\mathsf Q_1+2)^n(n/\varrho)^n\prod_q\|f_q\|_\sharp$.
\item[(M8)] \emph{The convergence on fixed tori.} For every side $\ell=2L^m$, $n\ge2$, real Lipschitz $f_1,\dots,f_n$ on $T_\ell$, none of them $0$, whose supports have no common point and the spread $\ge\mathsf d>0$ (for $n=2$: the sup-distance of the supports), and every compact interval $\mathsf J\subset(0,\infty)$: $S^T_{n;K,\ell,L^Kb}(f_1,\dots,f_n)$ converges as $K\to\infty$ uniformly in $b\in\mathsf J$, and $\sup_K\sup_{b\in\mathsf J}|S^T_{n;K,\ell,L^Kb}(f)|\le\mathsf c'(\ell,n,\mathsf J,\mathsf d)\prod_i\|f_i\|_\sharp$, the right member reading of the functions the spread and the norms alone.
\item[(M9)] \emph{Periodic Gibbs states.} The Definitions of kernels and Gibbs states: the Gibbs kernel $\gamma_B$ of a finite set $B$ of bonds is formed with the exponent $H_B=-\beta\sum_{p\in\mathrm{Pl}[B]}a_p$, $\mathrm{Pl}[B]$ the plaquettes containing a bond of $B$; $\nu$ is a Gibbs state if $\nu(\gamma_BF)=\nu(F)$ for all $B$ and bounded Borel $F$, and periodic if it is invariant under a subgroup of finite index of the lattice translations. There is at least one periodic Gibbs state at every $(a,\beta)$, and weak limits of the laws of the tori are such states; every periodic Gibbs state is invariant under the translations and the cubic group of its lattice; weak limits of periodic Gibbs states at couplings $\beta_k\to\beta$ that are invariant under the lattice translations are periodic Gibbs states at $\beta$; the tangent inequality for periodic Gibbs states; the exact dilation facts: for real $s>0$ the dilation $\mathrm{dil}_s:x\mapsto sx$ carries the bonds, the plaquettes, the functions $a_p$, the product Haar measures and the glued configurations of $a\mathbb Z^3$ to those of $sa\mathbb Z^3$, with $\mathcal O_a(f)(U)=s\,\mathcal O_{sa}(f(\cdot/s))(\mathrm{dil}_sU)$, and for $s$ a power of $L$ it maps the periodic Gibbs states onto the periodic Gibbs states; the classes at a cutoff and the bound system at $g$ (equation~\YMMref{(2030)}{765}; Section~\YMMref{100}{743}), whose two members bound the cumulants of separated tuples in the members of a class and the differences of the one-point functions of two members, free of the cutoff, the coupling and the state; Proposition~\YMMref{106.3}{765}: for $(K,\beta)\in\mathfrak d(g)$ the set $\mathcal G^{\rm per}_{K,\beta}$ is a class at the cutoff $K$ that obeys the bound system at $g$, and with $\nu$ it contains $k\nu/\nu(k)$ for every bounded $\mathcal I$-measurable $k\ge0$ with $\nu(k)>0$; and the plateau statement for such a class (Proposition~\YMMref{100.1}{744}, with the numbers of equations~\YMMref{(1976)--(1980)}{743}), with its statements on conditional expectations given the invariant $\sigma$-field $\mathcal I$.
\item[(M10)] \emph{Limits.}
We use the Gibbs limits and their component laws in
Theorems~\YMMref{AP}{780} and~\YMMref{AL}{746}, and the compactness,
translation and cubic invariance, and coordinate reflection
positivity of growing-torus limits in
Propositions~\YMMref{92.1}{710}, \YMMref{92.3}{712}
and~\YMMref{92.4}{713}.
The radial kernel and the reconstruction, including their bounds
and continuity, are those stated in Definition~\ref{C-365-5}.
\item[(M11)] \emph{The ordered choice of constants.} Use the order in the proof of Proposition~\YMMref{69.1}{595}, with the source floors of Section~\YMMref{82}{665} inserted at their stated places. Every initial segment selected by this procedure extends to a complete choice: the next free constant has either a finite lower bound, a positive upper bound or the nonempty interval specified at that step, and each determined constant is then fixed by its formula. Continue through the finite list and choose $\gamma$ last. With the constants before $L$ fixed, all sufficiently large odd $L$ are allowed; the remaining constants are subsequently chosen in that same order.
\item[(M12)] For Proposition~\ref{C-365-11} alone, the following, on one torus $T_\lambda$, $\lambda=2L^m$, at one label (a hat marks the objects of the classical starts). The theorem among the construction results on two consecutive cutoffs (Theorem~\YMMref{AR}{810}), with its read-out of the one-point functions, for derivatives in the source and read here at the order $1$, up to the derivative at the source $0$ of the difference of its contact terms (Section~\YMMref{116}{807}), and with the bound of its error, summable in $K$ (its clause~(3)). The objects of its comparison (Section~\YMMref{116}{807}) for the lattices of the cutoffs $K$ and $K+1$, marked $\mathrm A$ and $\mathrm B$, with the small numbers $\epsilon_{\rm fg}$, $\epsilon_{\rm cl}$, $\epsilon_{\rm sh}$, $\epsilon_3=4\epsilon_{\rm cl}+6\epsilon_{\rm sh}$ and the depth $u=\lceil K/1260\rceil$ of Section~\ref{sec:cutoff-calibration} at $n=1$, where $\epsilon_{\rm cl},\epsilon_{\rm sh}\le\frac1{20}$ under the conditions of the theorem: on the inner ball $\mathsf{In}$, a set of block fields of the depth $u$ that contains the trivial block field $\mathbf 1$, the measures of the chains and of the classical starts are a constant times the exponential of an exponent times the Haar measure, and their masses, integrals and constants satisfy $T^\bullet=e^{c^\bullet}I^\bullet$, $\hat T^\bullet=e^{\hat c^\bullet}\hat I^\bullet$ and are holomorphic in the source (Proposition~\YMMref{114.2}{801} and the proof of Proposition~\YMMref{116.3}{809}); the measures of the classical starts are integrals over the block fields of a larger depth, against the weight of the run (Section~\ref{sec:block-field-laws}), and at real sources the functions so integrated and the measures are non-negative (the proof of Proposition~\YMMref{116.1}{808}).

The data of its read-out (Section~\YMMref{116}{807}): $|I^\bullet/\hat I^\bullet-1|\le2\epsilon_{\rm fg}\le\frac14$ and a bound of the constants of the classical starts (the proof of Proposition~\YMMref{116.3}{809}), and at real sources $\hat T^{\mathrm A},\hat T^{\mathrm B}>0$ with $|\hat T^{\mathrm B}/\hat T^{\mathrm A}-1|\le\epsilon_3\le\frac12$ (Proposition~\YMMref{116.1}{808}). The discarding bound $0\le Z-T^\bullet(0)\le\mathsf R_{t_1}Z$, $Z$ the partition function of the torus at the cutoff of $\bullet$ and $\mathsf R_{t_1}$ the discarding number, one number for both cutoffs (Proposition~\textup{\YMMref{115.2}{806}}, the proof of Theorem~\YMMref{AR}{810}), with $\mathsf R_{t_1}\le\frac12$ under the conditions of the theorem.

The number $\mathfrak o_K=8\mathsf R_{t_1}+12\epsilon_{\rm fg}+8\epsilon_{\rm cl}+12\epsilon_{\rm sh}$ and, for an interval $\mathsf J$ of labels of the theorem on the side $\lambda$ (Section~\ref{sec:cutoff-calibration}; every compact interval $\mathsf J\subset I(g)$ is one on every side of $\mathfrak L(g)$, as $\min\mathsf J>g^{-2}+\bar B\ge\gamma^{-2}$, $\max\mathsf J\le X_g-\bar B\le L(g^{-2}+\bar B)$ and $M\,R(\max\mathsf J)\le M\,R(L(g^{-2}+\bar B))<L^{m_{\mathfrak L}(g)}$ by equation~\YMMref{(1528)}{615}, the cube-size function $R$ being non-decreasing), its number $\mathsf C_{\rm bf}(\mathsf J)$, formed from the constants and the ends of $\mathsf J$, its exponent $\mathfrak p_0$ and its thresholds $K_{\rm bf}(\lambda,\mathsf J)$: for every label of $\mathsf J$ and every $K\ge K_{\rm bf}(\lambda,\mathsf J)$ the conditions of the theorem hold and $\mathfrak o_K\le\mathsf C_{\rm bf}(\mathsf J)\lambda^3K^{\mathsf p_0}L^{-K/10}$, $\mathsf p_0=2\mathfrak p_0+\frac12$ (Definition~\textup{\ref{IV-62-62}}~(4), (5), Corollary~\textup{\ref{IV-62-66}}~(a)); with these conditions the setting (Th) among the construction results holds for the run of each of the two cutoffs on the torus (Section~\ref{sec:cutoff-calibration}).

Its comparison of the two classical runs: their masses and their densities in $L^1$ (Proposition~\textup{\YMMref{116.1}{808}} and its proof); the pointwise quotient of the two classical starts at a real source, a factor in $[e^{-\epsilon_{\rm cl}},e^{\epsilon_{\rm cl}}]$ where both weights are positive, the two weights being equal there (Proposition~\textup{\YMMref{116.1}{808}}); and the bound, by $\epsilon_{\rm sh}$ times its mass, of the part of one classical run where the weight of the other vanishes (the positive-kernel argument in the proof of Proposition~\YMMref{116.1}{808}), which bounds the tails of a classical run by multiples of its measures as an order of measures, against every bounded Borel function of the block field. The exponent and the constant of a step of a run (Proposition~\YMMref{114.2}{801}), and the action of a bare array, in which every member enters less its value at the trivial pair (equation~\YMMref{(874)}{412}); the one value and the continuity of the weighted action of the background (the gauge-invariance argument preceding equation~\YMMref{(2133)}{790} and the local analytic-background argument in the proof of Proposition~\YMMref{52.3}{433}), the continuous choice of the background near a regular block field (Theorem~\YMMref{E}{10} and its real-analytic continuation), and the one value of the action of the array on the minimal orbit (Propositions~\YMMref{114.2}{801} and~\YMMref{113.13}{799}); a local term is a holomorphic, so continuous, function of the pair that it reads, on an open set of pairs (Lemma~\ref{lem:analytic-open-pairs}), and the pair of a configuration on a domain is a polynomial map of its matrix entries (Section~\ref{sec:analytic-core-complements}, \YMMref{(70)}{38}, \YMMref{(69)}{38}).

The line $x_j=L^{K-j}b$, $0\le j\le K$, of a label, whose increments vanish, is an admissible finite sequence with last index $K$, which with the integer $m$ of a side $2L^m$ forms admissible data, and its cut $(x_0,\dots,x_{K-1})$ with the lattice $(\lambda,K)$ and these couplings forms a frame of level $K-1$ among the construction results (the proof of Lemma~\YMMref{118.1}{812}). The family of the level $j$ of the chain at the source $0$ is the family of the level $j$ of the canonical progression of the cut $(x_0,\dots,x_{K-1})$ of the line $x_j=L^{K-j}b$ of the label (the proof of Lemma~\YMMref{118.1}{812}); at that frame the members of the levels $j\le K$ of this progression are those of the progression of the line, and they are regular families on every lattice of their level, with their companions (Theorem~\YMMref{Y}{438}).

The exact step identity \eqref{eq:vacuum-step-normalization} identifies the logarithm of the Gaussian normalizing factor with the extensive term and the sum of its local Gaussian constants. Adding the values $E^{(j+1)}(X,\mathbf1,z)$ of the regular family gives the vacuum number $\mathsf E^{(j)}(T_\ell,T_\ell)$ of that step. Their sum through step $n$ is $E_{(n)}$, as in \eqref{eq:vacuum-prefix-definition}; these definitions agree with equations~\YMMref{(1406)--(1407)}{574}. Lemma~\ref{lem:vacuum-prefix} proves the exact identity and the comparison needed below.

For the comparison on two tori we use the same admissible sequence and its regular families, with their common companions, on both sides. The real vacuum families and the rate conditions are those stated in Lemma~\ref{lem:vacuum-prefix}. The pair of sides is admitted once $LMR(x_K)\le L^{m_i}$ for $i=1,2$: the retained room condition \eqref{eq:running-source-history-room} gives $C_L\le2LR(\gamma^{-2})\le2LR(x_K)$, so $\ell_i=2L^{m_i}\ge2LMR(x_K)\ge C_LM$. For the line $x_j=L^{K-j}b$ of a label $b\in\mathsf J$, the last radius is $R(b)$ (equation~\YMMref{(634)}{339}), and on the sides of $\mathfrak L(g)$ the required room is $LMR(b)\le LMR(L(g^{-2}+\bar B))\le L^{m_{\mathfrak L}(g)}\le L^{m_i}$, by equation~\YMMref{(1528)}{615}.
\item[(M13)] For Theorem~\ref{C-365-10} alone: the consequences of the hypothesis K1 proved in Section~\ref{C-361-1} on one inequality between the partition functions of two periodic boxes. Let $0<g\le\gamma$, $\mathsf b\in I(g)$, let $\mathrm{K1}[\mathsf b,\ell_\circ,K_\circ;x_\circ]$ hold, and let $K\ge\max\{K_\circ,1\}$ and $\beta=L^K\mathsf b$. (i) For every side $\ell=2L^m\ge\ell_\circ$: $0\le\log Z_{K,\ell,\beta}-\ell^3p_{K,\infty}(\mathsf b)<0.13736\,\theta_\circ^{\ell/\ell_\circ}$, the bound being the same at every such cutoff. (ii) There is exactly one periodic Gibbs state $\nu_K$ at $(K,\beta)$. (iii) Over every cross-section of physical sides $\ge\ell_\circ$ the largest eigenvalue $\lambda_0$ of the transfer operator defined in Section~\ref{C-361-1} is simple, and the next one, $\lambda_1$, has $-a_K^{-1}\log(\lambda_1/\lambda_0)>m_*$. (iv) For $n\ge2$, a side $\ell_1=2L^{m_1}\ge\ell_\circ$ and $X=(f_1,\dots,f_n)$ of real Lipschitz functions, none of them $0$, with pairwise disjoint compact supports in $(-\ell_1/4,\ell_1/4)^3$, there is a number $C(X)$ with $|S^T_n[\nu_K](X)-S^T_{n;K,\ell_2,\beta}(X^{\langle\ell_2\rangle})|\le C(X)\,\theta_\circ^{\ell_2/(2\ell_\circ)}$ for all $K\ge\max\{K_\circ,1\}$ and all sides $\ell_2=2L^{m_2}\ge\ell_1$. (v) The Gibbs sequence $(K,L^K\mathsf b,\nu_K)_{K\ge\max\{K_\circ,1\}}$ converges; its limit $\sigma$ is its own single component; $\ker H=\mathbb C\Omega$ in $\mathcal H(\sigma,e_3)$; and the spectrum of $H$ lies in $\{0\}\cup[m_*,\infty)$.
\item[(M14)] For clause (f) of Theorem~\ref{C-365-9} alone: the power of the jump of the slope of the pressure that Section~\ref{lh:sec} on the latent heat of the $\mathrm{SU}(2)$ Wilson model on $\mathbb Z^3$ proves for two frames at one radius with the block sizes $L_1$ and $L_2=L_1+2$, under the following hypotheses. Theorems~\YMMref{AJ}{673} and~\YMMref{AD}{607} are granted at both frames, and Theorem~\YMMref{D}{8} at their choices and at a lowered choice of each. In each frame $i\in\{1,2\}$ the far part of the specific heat settles in the cutoff at a rate on every torus: with the window $\mathcal W_i$, the tent $\chi^{(i)}$ and the plateau $\eta^{(i)}$ of half-width $1/L_i$ about $(\frac12,\frac12,\frac12)$ of clause (f), and the far part $\mathsf y^{(i)}_{m,K}(b)=b^5L_i^3\,S^T_{2;K,\ell,L_i^Kb}\bigl((\chi^{(i)})^{\langle\ell\rangle},1-(\eta^{(i)})^{\langle\ell\rangle}\bigr)$ of the specific heat of $T_\ell$, $\ell=2L_i^m$, there are numbers $\mathsf k_i\ge0$ and $\mathsf c_{{\rm y},i}<\infty$, free of $m$, $K$ and $b$, such that for all integers $m\ge1$, all $b\in\mathcal W_i$ and all integers $K'\ge K\ge\mathsf k_i+54m$
\[\bigl|\mathsf y^{(i)}_{m,K'}(b)-\mathsf y^{(i)}_{m,K}(b)\bigr|\le\mathsf c_{{\rm y},i}\,\ell^9K^{\mathsf p_{2,i}}L_i^{-K/10},\qquad\mathsf p_{2,i}=6\mathfrak p_{0,i}+\tfrac12,\]
$\mathfrak p_{0,i}$ the exponent of the theorem on two consecutive cutoffs (Theorem~\YMMref{AR}{810}) in the frame $i$. Under these hypotheses, for every real $A>90$ there are numbers $\mathsf C_{\rm r}$ and $\beta_1$ with $\beta^4J(\beta)\le\mathsf C_{\rm r}\beta^{-1/(2A)}$ for all $\beta\ge\beta_1$.
\end{itemize}
The proof of Proposition~\ref{C-365-11}~(c) derives from (M12), in its Step~3 and on one torus $T_\lambda$, four statements: (i) the read-out at the order $0$; (ii) the comparison of the two classical starts as measures on the inner ball; (iii) the cancelling of the constants of the two classical starts, from the value $0$ and the continuity of their exponents at the trivial configuration; and (iv) the identity of the sum of the constants of the chain at the source $0$ with a vacuum energy of the comparison of two tori, a real number.
We use the following standard results: Cauchy's formula and estimates, Schwarz's lemma, the maximum principle, Poisson's formula, the inequalities of Jensen, H\"older and Cauchy--Schwarz, the theorems of Tonelli and Fubini, Morera's theorem, Kronecker's theorem (the multiples of an irrational number are dense modulo $1$), Helly's selection theorem, Bernstein's theorem on completely monotone functions and his inequality for the derivatives of polynomials, the spectral theorem, Stolz's theorem, and H\"older's inequality in the Schatten classes, which in the one case used is proved in Proposition~\ref{C-365-20}.

\subsection{Estimates and proofs}

Up to the window inequality the argument runs in one frame. The energies averaged over the isometries and their cross terms (Lemma~\ref{C-365-14}) give the energies of a constant source on two sides on a complex disc (Lemma~\ref{C-365-15}); finite size (Lemma~\ref{C-365-16}) is exact; the stability functional (Lemma~\ref{C-365-17}) gives the imaginary direction (Proposition~\ref{C-365-18}), which with the cross terms and the chessboard estimate gives the walls (Proposition~\ref{C-365-19}); these four give the main inequality and semi-convexity (Proposition~\ref{C-365-21}); with Cauchy's estimate in the chessboard seminorm (Lemma~\ref{C-365-22}) the main inequality gives the reduction on tori (Proposition~\ref{C-365-23}), and with finite size and the domination the window inequality in every periodic Gibbs state (Proposition~\ref{C-365-24}). The bound against piles of specific heat (Proposition~\ref{C-365-28}) uses semi-convexity and finite size in both frames, the transfer, the convergence on fixed tori and dilations by real factors (Lemmas~\ref{C-365-25}--\ref{C-365-27}). Clauses (a), (b) of Theorem~\ref{C-365-8} follow from the reduction and this bound; clauses (c), (d) from (a), (b) and the theorems on limits through Proposition~\ref{C-365-30} and Lemmas~\ref{C-365-31}--\ref{C-365-33}; clause (e) through Lemma~\ref{C-365-29}.

\begin{lemma}[The energies averaged over the isometries; cross terms]\label{C-365-14}
Let $(K,\beta)\in\mathfrak d(g)$ and $\ell\in\mathfrak L(g)$.
\begin{enumerate}
\item[(A1)] $\mathcal E^{\rm av}$ is holomorphic along finite families, $\mathcal E^{\rm av}(0)=0$, $\mathcal E^{\rm av}(\bar w)=\overline{\mathcal E^{\rm av}(w)}$, $\mathcal E^{\rm av}(\mathsf rw)=\mathcal E^{\rm av}(w)$ for $\mathsf r\in\mathsf G$, $\mathcal E^{\rm av}(z\cdot1)=\mathcal E(z\cdot1)$ for complex $|z|<2\varrho$, and $|\widetilde Z^{\rm av}(w)|=\prod_{\mathsf r\in\mathsf G}|\widetilde Z(\mathsf rw)|^{1/|\mathsf G|}\le e^{C_S\ell^3}$.
\item[(A2)] Let $\mathcal F\in\{\mathcal E,\mathcal E^{\rm av}\}$, and let $u,w_1,\dots,w_n$ be complex Lipschitz functions on $T_\ell$ such that the support of $w_r$ is covered by $n_r$ closed sup-balls of radius $R_0$, the supports of different $w_r$ are at sup-distance $\ge\Delta\ge4M$, and $\|u\|_\sharp$, $\|u+w_r\|_\sharp$, $\|u+\sum_rw_r\|_\sharp<2\varrho$. With $\mathcal F'(w)=\mathcal F(u+w)-\mathcal F(u)$, the cross term $\mathrm{cr}=\mathcal F'(\sum_rw_r)-\sum_r\mathcal F'(w_r)$ satisfies $|\mathrm{cr}|\le C_\times(R_0/\Delta)\sum_rn_r$.
\item[(A3)] If the $w_r$ depend linearly on $\zeta\in\mathbb C^k$ and (A2) applies for all $\zeta$ with $\max_i|\zeta_i|<\rho_3$, then $|\mathrm{cr}(\zeta)|\le C_\times(R_0/\Delta)\bigl(\sum_rn_r\bigr)(\max_i|\zeta_i|/\rho_3)^2$.
\end{enumerate}
These bounds are free of $K$; what grows with $K$ is only the bound of $\mathcal E$ itself.
\end{lemma}

\begin{proof}[Sketch of proof]
(A1) By (M4), $Z[w]=Z[\mathsf rw]=e^{E+\mathcal E(\mathsf rw)}\widetilde Z(\mathsf rw)$ for every $\mathsf r\in\mathsf G$; the geometric mean of the moduli over $\mathsf G$ gives $|Z[w]|=e^{E+\operatorname{Re}\mathcal E^{\rm av}(w)}\prod_{\mathsf r}|\widetilde Z(\mathsf rw)|^{1/|\mathsf G|}$, which is the formula for $|\widetilde Z^{\rm av}(w)|$, bounded by (M1). Holomorphy, the value at $0$ and the conjugation pass from $\mathcal E$ to its mean; the invariance holds because $\mathsf G$ is a group; constants are fixed by every $\mathsf r$. (A2) It suffices to take $\mathcal F=\mathcal E$, since the isometries preserve the hypotheses; then $\mathrm{cr}=\sum_{j,X}\mathsf b(X)$ with
\[\mathsf b(X)=e_j(X;u+\textstyle\sum_rw_r)-e_j(X;u)-\sum_r\bigl[e_j(X;u+w_r)-e_j(X;u)\bigr].\]
By locality $\mathsf b(X)=0$ unless $X$ with one layer of cubes meets at least two of the supports; then $\mathsf b(X)$ has at most $3\#R(X)$ terms that do not cancel, $R(X)$ the set of the supports met, each $\le C_Se^{-\kappa d_j(X)}$, and $d_j(X)\ge y_s-4$ with $y_s=L^s\Delta/M\ge4L^s$ by (M3)(a). The count of domains at $Y$ one of the $n_r$ balls, with $\kappa-\log9216>2$, bounds the sum over such $X$ by $n_r\cdot27S(\kappa)(2R_0L^s/M+2)^3e^{-2(y_s-4)}$; with $(2R_0L^s/M+2)^3\le32[(R_0/\Delta)^3y_s^3+1]$, $\sum_{s\ge0}y_s^3e^{-2(y_s-4)}\le65$ and $\sum_se^{-2(y_s-4)}\le2$ one obtains $|\mathrm{cr}|\le3\cdot27\cdot32\,[65(R_0/\Delta)^3+2]C_SS(\kappa)\sum_rn_r$, and $3\cdot27\cdot32=2592$. (A3) $\mathrm{cr}(\zeta)$ is holomorphic on the polydisc, vanishes at $0$ with its first derivatives, since $D\mathcal F'(0)$ is linear; the maximum principle for $t\mapsto\mathrm{cr}(t\zeta^0)/t^2$ gives the claim.
\end{proof}

\begin{lemma}[Two sides at a constant source, on a complex disc]\label{C-365-15}
Let $(K,\beta)\in\mathfrak d(g)$, $K\ge1$, and let $\ell_1,\ell_2\in\mathfrak L(g)$ with $\ell_1,\ell_2\ge2\mathsf l_1$; let $h_1$, $\mathsf A_1$ be the numbers of Definition~\ref{C-365-7} formed at $g$. The function
\[d(z)=\ell_1^{-3}\mathcal E_{\ell_1}(z\cdot1)-\ell_2^{-3}\mathcal E_{\ell_2}(z\cdot1)\]
is holomorphic on $|z|<2\varrho$, real on the real axis, $d(0)=0$, and $|d(z)|\le\mathsf A_1$ for all complex $|z|\le h_1$. Hence, with $\mu'=d'(0)\in\mathbb R$, $|d(x)-\mu'x|\le2\mathsf A_1x^2/h_1^2$ for real $|x|\le h_1/2$ and $|d''(0)|\le2\mathsf A_1/h_1^2$. The numbers $h_1$, $\mathsf A_1$ are free of $K$, $\beta$, $\ell_1$ and $\ell_2$.
\end{lemma}

\begin{proof}[Sketch of proof]
Write $z\cdot1=\sum_uz\chi_u$ with the tents $\chi_u$ of half-width $\varsigma=1/L$ about the points $u\in\varsigma\mathbb Z^3$; $\|\chi_u\|_\sharp\le\mathsf c_\chi$. Call a domain of depth $s$ large if $d_j(X)\ge A_s-4$, $A_s=L^{s-1-m'}$, and small otherwise; by (M3)(a) a small domain cannot meet two sets at sup-distance $\ge\varsigma$. The large part has a density bounded on the whole disc, torus by torus: $\ell^{-3}|\mathcal E^{\rm lg}_\ell(z\cdot1)|\le27C_SS(\kappa)L^3\sum_sA_s^3e^{-2[A_s-4]_+}\le C_{\rm lg}$. For the small part, the invariance of $\mathcal E^{\rm av}$ under the translations by $\varsigma\mathbb Z^3$ and the chain rule give
\[\ell^{-3}\partial_z^n\mathcal E^{\rm sm}_\ell(z\cdot1)\big|_0=L^3\sum_{\rm near}D^n\mathcal E^{\rm sm,av}_\ell[\chi_0,\chi_{v_2},\dots,\chi_{v_n}],\]
the sum over the at most $125^{n-1}$ near tuples, $|v_q-v_{q'}|_\infty\le2\varsigma$, the same index set on both tori (here $\ell_1,\ell_2\ge2\mathsf l_1$ is used: the $125$ points are distinct points of each torus). At a near tuple $D^n\mathcal E^{\rm av}$ is a joint cumulant minus a derivative of $\log G_\ell$, $G_\ell(\zeta)=\widetilde Z^{\rm av}_\ell(w)/\widetilde Z_\ell(0)$, $w=\sum\zeta_q\chi_{v_q}$. The chessboard estimate on the board of side $\mathsf l_1$ (an even number $\ge2$ of cells in each direction), with the source formula and the cross terms of Lemma~\ref{C-365-14} at $R_0/\Delta\le1/8$, spends the extensive bounds of (M1) on one cell and gives $|G_\ell|\le\mathsf Q_2$; Schwarz's lemma and Cauchy's estimate give $|\partial^n\log G_\ell(0)|\le\log2\,(n\mathsf c_\chi(2\mathsf Q_2+2)/\varrho)^n$; the large part at a near tuple is $\le C'_{\rm lg}(n\mathsf c_\chi/\varrho)^n$ by Cauchy; and the cumulants on the two tori differ by the transfer (M7) by at most $2\log2\,(2\mathsf Q_1+2)^n(n\mathsf c_\chi/\varrho)^n$. Altogether, with $n^n\le e^nn!$,
\[\frac{|\partial^n_zd^{\rm sm}(0)|}{n!}\le\frac{L^3}{125}\bigl(4\log2+2C'_{\rm lg}\bigr)\Bigl(\frac{125e\,\mathsf c_\chi\mathsf Q_*}{\varrho}\Bigr)^n=\frac{L^3}{125}\bigl(4\log2+2C'_{\rm lg}\bigr)(2h_1)^{-n},\]
which sums on $|z|\le h_1$; with the large part, $|d|\le\mathsf A_1$. The forms at real $x$ follow from Cauchy's estimate on $|z|=h_1$; $d$ is real on the real axis by the conjugation in (M1).
\end{proof}

\begin{lemma}[Finite size, exactly]\label{C-365-16}
\begin{enumerate}
\item[(a)] For even integers $n_1,n_2,n_3\ge2$ and $\beta\ge0$ let $Z(n_1,n_2,n_3)$ be the partition function of Wilson's action on the periodic box $\mathbb Z^3/(n_1\mathbb Z\times n_2\mathbb Z\times n_3\mathbb Z)$ and $\mathsf g=\log Z/(n_1n_2n_3)$. Then $\mathsf g$ does not increase in each argument; $P(\beta)=\inf\mathsf g=\lim_n\mathsf g(n,n,n)$, the limit along the even $n$, is convex in $\beta$ and $P\ge-6\beta$; and for even $\mathsf N\ge e_0\ge2$
\[0\le\mathsf g(\mathsf N,\mathsf N,\mathsf N)-P\le3(e_0/\mathsf N)\bigl[\mathsf g(e_0,e_0,e_0)-P\bigr].\]
\item[(b)] For every block size $L$, every $K\ge0$, $b\ge0$ and sides $\ell\ge\ell'$ of the form $2L^m$: $p_{K,\infty}(b)=\lim_\ell p_{K,\ell}(b)=L^{3K}P(L^Kb)$ with the one function $P$ of (a), $D_K\le D_{K,\ell}\le0$, and $0\le p_{K,\ell}-p_{K,\infty}\le3(\ell'/\ell)(p_{K,\ell'}-p_{K,\infty})$.
\item[(c)] In a frame, for $b\in I(g)$, $K\ge1$ and $\ell\ge\ell'$ in $\mathfrak L(g)$: $-B'\le D_K(b)\le D_{K,\ell}(b)\le0$ and $0\le D_{K,\ell}(b)-D_K(b)\le\bar\eta(\ell)$.
\end{enumerate}
Parts (a) and (b) are exact statements on one lattice, at every $\beta\ge0$.
\end{lemma}

\begin{proof}[Sketch of proof]
Fix even $n_2,n_3$ and let $K_n(\phi,\phi')$ be the slab kernel: the integral, over the bonds of the slab $\{0\le x_1\le n\}$ that lie in neither boundary plane, of $\exp[-\beta(\sum'a_p+\frac12\sum_{x_1=0}a_p+\frac12\sum_{x_1=n}a_p)]$, the bonds of the two planes being fixed to $\phi$, $\phi'$. It is continuous with values in $(0,1]$, symmetric (the reflection $\vartheta:x_1\mapsto n-x_1$ maps the slab onto itself, preserves the Haar measure and carries $a_p$ to $a_{\vartheta p}$, because $\operatorname{tr}g^{-1}=\operatorname{tr}g$ on $\mathrm{SU}(2)$) and multiplicative, $K_{n+m}(\phi,\phi'')=\int K_n(\phi,\phi')K_m(\phi',\phi'')\,d\phi'$ (Fubini; the plaquettes of the middle plane receive $\frac12+\frac12$). With $z(n)=Z(n,n_2,n_3)=\int K_n(\phi,\phi)\,d\phi$ one has $z(n+m)=\langle K_n,K_m\rangle_{L^2}$ and $\|K_n\|_2^2=z(2n)$, so the Cauchy--Schwarz inequality gives
\[z(n+m)^2\le z(2n)z(2m),\qquad z(4k)\le z(2k)^2 .\]
The first, at $n=e/2-1$, $m=e/2+1$, makes $\varphi(e)=\log z(e)\le0$ convex on the even $e$; its slopes increase to a finite $\Lambda_\infty\le0$, and $\varphi(e)/e\to\Lambda_\infty$ by Stolz. The excess $\mathrm{exc}(e)=\varphi(e)-e\Lambda_\infty$ does not increase, and it is $\ge0$ because $\varphi(2^ke)/2^ke$ does not increase by the second inequality; so $\varphi(e)/e=\Lambda_\infty+\mathrm{exc}(e)/e$ does not increase. This in each direction gives the monotonicity and $P=\inf\mathsf g=\lim\mathsf g(n,n,n)$, convex as a limit of convex functions and $\ge-6\beta$ since $a_p\le2$ and there are three plaquettes to a site. For the rate, $\mathrm{exc}(\mathsf N)\le\mathrm{exc}(e_0)\le e_0n_2n_3[\mathsf g(e_0,e_0,e_0)-P]$ gives $\mathsf g(\mathsf N,n_2,n_3)\le\mathsf g(n,n_2,n_3)+e_0[\mathsf g(e_0,e_0,e_0)-P]/\mathsf N$ for all even $n$; applied in the three directions in turn, and letting $n\to\infty$, this is (a).

For (b), $T_\ell$ at the cutoff $K$ is the box with $\mathsf N=\ell L^K$ sites a side, $p_{K,\ell}(b)=L^{3K}\mathsf g(\mathsf N,\mathsf N,\mathsf N)$ at $\beta=L^Kb$, and $e_0/\mathsf N=\ell'/\ell$ for $e_0=\ell'L^K$; the limit $P$ names no block size. For (c), by the passage and (M1), $p_{K,\ell}=\ell^{-3}[E_\ell+\log\widetilde Z_\ell(0)]\in\ell^{-3}E_\ell+[-C_g,C_S]$, and $|\ell^{-3}E_\ell-\ell'^{-3}E_{\ell'}|\le C_S$ gives $D_{K,\ell}\ge-B'$ for all $\ell\in\mathfrak L(g)$, hence $D_K\ge-B'$; enter $-D_K\le B'$ in (b).
\end{proof}

\begin{lemma}[The third member of the stability chain as a functional of the source]\label{C-365-17}
Let $(K,\beta)\in\mathfrak d(g)$ and $\ell\in\mathfrak L(g)$. There are a positive $\sigma$-finite measure space $(\Xi,\mathfrak M)$ and functions $\Psi_\varpi$, $\varpi\in\Xi$, on the class of the sources, neither $\Xi$ nor $\mathfrak M$ depending on the source, such that with $h_\varpi=\operatorname{Re}\Psi_\varpi$:
\begin{enumerate}
\item[(T1)] for every source $w$, $\varpi\mapsto e^{\Psi_\varpi(w)}$ is measurable for the completion of $\mathfrak M$, $\widetilde Z(w)=\int e^{\Psi_\varpi(w)}\,d\mathfrak M(\varpi)$, and
\[|\widetilde Z(w)|\le\mathsf T(w):=\int e^{h_\varpi(w)}\,d\mathfrak M(\varpi)\le e^{C_S\ell^3};\]
\item[(T2)] $\mathsf T(u)=\widetilde Z(u)>0$ at every real source $u$;
\item[(T3)] for every $\varpi$ and every finite family $f_1,\dots,f_n$, the map $z\mapsto e^{\Psi_\varpi(\sum_qz_qf_q)}$ is holomorphic and without zeros on the convex open set $\{z:\|\sum_qz_qf_q\|_\sharp<2\varrho\}$; so $z\mapsto h_\varpi(\sum_qz_qf_q)$ is the real part of a holomorphic function there;
\item[(T4)] for real $u$, $v$ there is an $\mathfrak M$-null set, depending on the pair, off which $h_\varpi(u+\bar\zeta v)=h_\varpi(u+\zeta v)$ for all $\zeta$ with $\|u+\zeta v\|_\sharp<2\varrho$; and $(\varpi,z)\mapsto h_\varpi(w(z))$ is jointly measurable.
\end{enumerate}
The functional $\mathsf T$ is the normalized total-variation sum in equation~\YMMref{(1738)}{675}, at $k=K$.
\end{lemma}

\begin{proof}[Sketch of proof]
Take $k=K$. By Theorem~\YMMref{AJ}{673}, there are finitely
many histories and a common source-independent null set $N_h$
for each history. Outside $N_h$, let $\mathsf D_w$ be its density:
it is zero where $\mathsf{Stay}^{\rm T}_h\chi_K(\Omega_K)=0$
and otherwise equals
$\mathsf{Stay}_h\exp(\boldsymbol\Phi_h^w+A_{K,h}^w)$.
Put $\Xi_h=\{\mathsf D_0\ne0\}$.
This set is the same for all sources, since the tests are unchanged
and an exponential has no zero. On it put
$d\mathfrak M_h=\mathsf{Stay}_h\,d\mathcal M_h$.
The reference measure is a finite product of Haar and Lebesgue
measures, hence is $\sigma$-finite; multiplication by the
nonnegative finite Borel test factor preserves this property.
Neither this measure nor $\Xi_h$ depends on the source.

Histories with $\mathfrak M_h(\Xi_h)=0$ contribute zero at
every source and may be omitted. Take the disjoint union of the
remaining $\Xi_h$, with measure $\mathfrak M_h$ on each piece, and set
\[
 \Psi_{(h,\phi)}(w)=\boldsymbol\Phi_h^w(\phi)
                   +A_{K,h}^w(\phi)-E_{(K)}^w.
\]
The total variation of $\mathsf D_w\mathcal M_h$ is
$|\mathsf D_w|\mathcal M_h$. Thus the integral of
$e^{\operatorname{Re}\Psi}$ is exactly the normalized
total-variation sum in equation~\YMMref{(1738)}{675}, which proves
the two inequalities in (T1). Without moduli, the history
decomposition in Theorem~\YMMref{AJ}{673} and the vacuum splitting
in equation~\YMMref{(1737)}{675} give
$\int e^{\Psi(w)}d\mathfrak M=e^{-E_{(K)}^w}Z[w]
=\widetilde Z(w)$.

For a real source, the bare density is positive, so each
history measure is positive; the vacuum constant is real by
conjugation. This proves (T2). To prove (T3), fix a history point
and a source-independent evaluation representative, as supplied by
Theorem~\YMMref{AJ}{673}. Every argument in its exponent is fixed
as the source varies. The active action is linear in $w$, and the
regular, boundary and retained terms are finite sums of the
holomorphic families in that theorem. The whole-torus vacuum
constant has the holomorphic finite-sum splitting in
Proposition~\YMMref{83.2}{675}.

The one member that no clause covers, the constant $c_h(w)=\exp(-\mathcal E^w_{K,\mathcal S(h)})$ of the label, which does not depend on $\phi$, is holomorphic as a quotient: on measurable sets $B_n\uparrow\Xi_h$ of finite measure on which the rest $e^{\Psi^\circ_\varpi}$ of the exponential is bounded near $z_0$, the map $z\mapsto\int F_01_{B_n}\,d\hat\tau^{{\rm co}\prime}_h[\rho^{w(z)}_0]$ is holomorphic by the bare-integral formula and holomorphy assertion of Proposition~\YMMref{73.5}{611}, and it equals $e^{E^{w(z)}_{(K)}}c_h(w(z))I_n(z)$ with $I_n(z)=\int_{B_n}F_0e^{\Psi^\circ_\varpi(w(z))}d\mathfrak M$ holomorphic by Morera and Fubini, and $I_n(z_0)=\int_{B_n}|e^{\Psi^\circ_\varpi(w(z_0))}|\,d\mathfrak M>0$ for all large $n$ when $F_0$ is the conjugate phase of the integrand at $z_0$, because the integrand has no zero, $B_n\uparrow\Xi_h$ and $\mathfrak M(\Xi_h)>0$.

So $e^{\Psi_\varpi(w(z))}$ is holomorphic without zeros on a convex set and has a holomorphic logarithm, whose real part is $h_\varpi(w(z))$. (T4) follows from conjugation of the history measures in Theorem~\YMMref{AJ}{673} and of the vacuum splitting in Proposition~\YMMref{83.2}{675}, first at the points $\zeta$ with rational coordinates and then by continuity; joint measurability from measurability in $\varpi$ and continuity in $z$.
\end{proof}

\begin{proposition}[The imaginary direction; a floor; Harnack; H\"older over several directions]\label{C-365-18}
Let $(K,\beta)\in\mathfrak d(g)$ and $\ell\in\mathfrak L(g)$, and let $u,v,v_1,\dots,v_n$ be real Lipschitz functions on $T_\ell$; $\|\cdot\|=\|\cdot\|_\sharp$.
\begin{enumerate}
\item[(P1)] If $\|u\|+\|v\|<2\varrho$: $|\widetilde Z(u+iv)|\le\mathsf T(u+iv)\le\widetilde Z(u)^{1-\vartheta}e^{\vartheta C_S\ell^3}$ with $\vartheta=2\|v\|^2/((2\varrho-\|u\|)^2+\|v\|^2)$.
\item[(P2)] If $\|u\|+\|v\|<2\varrho$: $\widetilde Z(u+v)\le\widetilde Z(u)^{1-\vartheta_1}e^{\vartheta_1C_S\ell^3}$ with $\vartheta_1=2\|v\|/(2\varrho-\|u\|+\|v\|)$.
\item[(P3)] If $\rho>0$, $\|u\|+\rho\sum_i\|v_i\|<2\varrho$ and $0\le nt\le\rho/4$: $\widetilde Z(u+t\sum_iv_i)\le\prod_i\widetilde Z(u+ntv_i)^{(1-\vartheta'')/n}e^{\vartheta''C_S\ell^3}$ with $\vartheta''=8n^2t^2/\rho^2$.
\end{enumerate}
For $\|u\|\le\varrho/2$: $\log\widetilde Z(u)\ge-(2C_g+C_S)\ell^3$ (the floor); if moreover $\|u\|+\|v\|<2\varrho$, then
\[
\begin{multlined}
\log|\widetilde Z(u+iv)|\le\log\widetilde Z(u)+C_I'\|v\|^2\ell^3\quad\text{(IM)},\qquad\log\widetilde Z(u+v)-\log\widetilde Z(u)\le\tfrac83(C_S \\
+C_g)\varrho^{-1}\|v\|\ell^3\quad\text{(the Harnack bound)},
\end{multlined}
\]
and $\partial_t^2\log\widetilde Z(u+tv)|_{t=0}\ge-2C_I'\|v\|^2\ell^3$ (semi-convexity). All of this holds for $\widetilde Z^{\rm av}$ in the place of $\widetilde Z$. $C_I'$ is free of $K$, $\beta$, $\ell$ and of the sources.
\end{proposition}

\begin{proof}[Sketch of proof]
For the Poisson kernel $\mathsf P_\rho(\zeta_0;\phi)=(\rho^2-|\zeta_0|^2)/|\rho e^{i\phi}-\zeta_0|^2$ one has (i) $\frac12[\mathsf P_\rho(iy;\phi)+\mathsf P_\rho(iy;-\phi)]\ge1-2y^2/(\rho^2+y^2)$ for $0\le y<\rho$; (ii) $\mathsf P_\rho(s;\phi)\ge1-2s/(\rho+s)$ for $0\le s<\rho$; (iii) $\prod_{i\le n}\mathsf P_\rho(t;\phi_i)\ge(1-8n^2t^2/\rho^2)\,n^{-1}\sum_i\mathsf P_\rho(nt;\phi_i)$ for $0\le nt\le\rho/4$; (i), (ii) are elementary, and (iii) follows from $\log\mathsf P_\rho(t;\phi)\ge\log(1-\tau^2)+2\tau\cos\phi-\tau^2/(1-\tau)$ and $n^{-1}\sum_i\mathsf P_\rho(nt;\phi_i)\le(1+\mathsf y)[1+\frac{16}3n^2\tau^2]$, $\tau=t/\rho$, $\mathsf y=2\tau\sum_i\cos\phi_i$. Thus each kernel is $(1-\vartheta)$ times the reference kernel plus $\vartheta R$ with $R\ge0$ of mean $1$. For (P1) take $1<\rho<(2\varrho-\|u\|)/\|v\|$; by (T3), (T4) of Lemma~\ref{C-365-17}, off a null set $G(\zeta)=h_\varpi(u+\zeta v)$ is harmonic near $|\zeta|\le\rho$ and $G(\bar\zeta)=G(\zeta)$. Poisson's formula at $\zeta=i$, made symmetric, gives $G(i)=(1-\vartheta)G(0)+\vartheta\int RG(\rho e^{i\phi})\,dm$; Jensen's inequality, integration in $\varpi$, H\"older's inequality with the exponents $1/(1-\vartheta)$, $1/\vartheta$ and Tonelli give
\[\mathsf T(u+iv)\le\mathsf T(u)^{1-\vartheta}\Bigl(\int R\,\mathsf T(u+\rho e^{i\phi}v)\,dm\Bigr)^{\vartheta}\le\widetilde Z(u)^{1-\vartheta}e^{\vartheta C_S\ell^3}\]
by (T1), (T2); let $\rho$ increase to its bound. (P2) is the same at $\zeta=1$ with (ii), and (P3) the same in $n$ variables with (iii), H\"older's inequality now with $n$ exponents $n/(1-\vartheta'')$ and one $1/\vartheta''$. The floor is (P2) at the base $u$ and the direction $-u$, where $\vartheta_1=\|u\|/\varrho\le\frac12$, with $\widetilde Z(0)\ge e^{-C_g\ell^3}$ from (M1). (IM): $\vartheta\le8\|v\|^2/(9\varrho^2)$ and $C_S\ell^3-\log\widetilde Z(u)\le2(C_S+C_g)\ell^3$ by the floor; the Harnack bound likewise with $\vartheta_1\le4\|v\|/(3\varrho)$. Semi-convexity: $\mathsf h(\zeta)=\log\widetilde Z(u+\zeta v)$ is holomorphic near $0$ and real on the real axis, and $\operatorname{Re}\mathsf h(iy)-\mathsf h(0)=-\frac12\mathsf h''(0)y^2+O(y^4)\le C_I'\|v\|^2y^2\ell^3$ by (IM). For $\widetilde Z^{\rm av}$ each inequality is one between products of powers of values of $\mathsf T$ whose exponents read norms alone; write it at the images under $\mathsf r\in\mathsf G$ and take the geometric mean, using (A1) of Lemma~\ref{C-365-14} and equality at real sources.
\end{proof}

\begin{proposition}[The walls]\label{C-365-19}
Let $b\in I(g)$, $K\ge1$, $T\ge\max\{T_0,\mathsf l\}$, $q=L^i\ge2T$, $\ell\in\mathfrak L(g)$ with $\ell\ge2q$, a board of $T_\ell$ of side $q$ with any offset, and $f\in\mathfrak P_T$ a profile of one of its cells, with $F$ as in Definition~\ref{C-365-4}. For real $x$ with $|x|\le\varrho/(256\mathsf c)$ and $x\le b$:
\[\log\widetilde Z^{\rm av}(xF)-\log\widetilde Z^{\rm av}(x)\le N\bigl[3C_wTq^2|x|+C_2'q^3x^2\bigr].\]
\end{proposition}

\begin{proof}[Sketch of proof]
By the passage, $(K,L^Kb)\in\mathfrak d(g)$. Fix $o_0\in a_K\mathbb Z$ and take eight boards of side $\mathsf l$: the board $\iota\in\{0,1\}^3$ has the offset $o_\iota$ with $(o_\iota)_\mu=o_0+\iota_\mu(\mathsf l+a_K)/2$, a point of $a_K\mathbb Z^3$ as $\mathsf l/a_K$ is odd. Take the partition of unity $\Theta_\iota=\prod_\mu\mathsf t_{\iota_\mu}(y_\mu)$, $\mathsf t_1=1-\mathsf t_0$, $\mathsf t_0$ the $\mathsf l$-periodic trapezoid equal to $1$ within $\mathsf l/8$ of $o_0+\mathsf l\mathbb Z$ and to $0$ beyond $3\mathsf l/8$, with $\operatorname{Lip}\Theta_\iota\le4\sqrt3/\mathsf l$; in a cell of the board $\iota$ the support of $\Theta_\iota$ lies in a sup-ball of radius $R_0\le7\mathsf l/16$ at sup-distance $\ge\mathsf l/16$ from the walls of that board.

The walls $\mathsf g=1-F$ split as $\mathsf g\Theta_\iota=\sum_rh_r$ into pieces of the cells of the board $\iota$, with $\|h_r\|_\sharp\le\mathsf c$ and mutual sup-distances $\ge\Delta=\mathsf l/8\ge4M$; since $\mathsf g\ne0$ only within $T/2$ of a mirror of the board of side $q$, at most $n_\iota\le9NTq^2/\mathsf l^3$ pieces are non-zero. H\"older's inequality (P3) at $u=x$, $n=8$, $v_\iota=-\operatorname{sgn}(x)\mathsf g\Theta_\iota$, $t=|x|$, $\rho=\varrho/(8\mathsf c)$, with the floor, costs $\vartheta''(2C_S+2C_g)\ell^3$, $\vartheta''=32768\mathsf c^2x^2/\varrho^2$ (this is the term in $x^2$, paid on the whole torus), and leaves the eight numbers $\mathsf B_\iota=\log[\widetilde Z^{\rm av}(x+z\mathsf g\Theta_\iota)/\widetilde Z^{\rm av}(x)]$, $z=-8x$.

Write $Z_b[w]=Z_{K,\ell,L^Kb}[w]$. On the board $\iota$, the shift of the label, $Z_b[x+w]/Z_b[x]=Z_{K,\ell,\beta'}[w]/Z_{K,\ell,\beta'}$ with $\beta'=L^K(b-x)\ge0$, and the chessboard estimate at $\beta'$ give $Z_b[x+z\mathsf g\Theta_\iota]/Z_b[x]\le\prod_r\|e^{z\mathcal O(h_r)}\|_{\rm ch}$, and by the source formula $\|e^{z\mathcal O(h_r)}\|_{\rm ch}^{N_f}=Z_b[x+zh_r^{\rm per}]/Z_b[x]$, where $h_r^{\rm per}$, the sum of the $N_f=(\ell/\mathsf l)^3$ mirror copies of $h_r$, is a dense array. The energies, whose bounds grow with $K$, cancel exactly up to cross terms. Put $\mathcal E'(w)=\mathcal E^{\rm av}(x+w)-\mathcal E^{\rm av}(x)$, the function $\mathcal F'$ of Lemma~\ref{C-365-14}~(A2) for $\mathcal F=\mathcal E^{\rm av}$ at the base source $u=x$; then $Z_b[x+w]/Z_b[x]=e^{\mathcal E'(w)}\widetilde Z^{\rm av}(x+w)/\widetilde Z^{\rm av}(x)$, and so, the sums being over the $r$ with $h_r\ne0$,
\[\mathsf B_\iota\le-\mathcal E'(z\mathsf g\Theta_\iota)+\sum_rN_f^{-1}\Bigl[\mathcal E'(zh_r^{\rm per})+\log\frac{\widetilde Z^{\rm av}(x+zh_r^{\rm per})}{\widetilde Z^{\rm av}(x)}\Bigr].\]
The mirror copies of $h_r$ are its images $\mathsf rh_r$ under reflections $\mathsf r\in\mathsf G$ at mirrors of the board $\iota$, which fix the constant $x$, so that $\mathcal E^{\rm av}(x+z\mathsf rh_r)=\mathcal E^{\rm av}(\mathsf r(x+zh_r))=\mathcal E^{\rm av}(x+zh_r)$ by (A1). By Lemma~\ref{C-365-14}~(A3), applied twice with the base source $u=x$, $\mathcal E'(z\mathsf g\Theta_\iota)=\sum_r\mathcal E'(zh_r)+\mathrm{cr}_\iota$ and $\mathcal E'(zh_r^{\rm per})=N_f\mathcal E'(zh_r)+\mathrm{cr}_r$; so the terms in $\mathcal E'$ of the bound of $\mathsf B_\iota$ are $-\mathrm{cr}_\iota+\sum_rN_f^{-1}\mathrm{cr}_r$, at most $2C_\times(7/2)n_\iota|z|^2\mathsf c^2/\varrho^2$. The dense arrays are bounded by the Harnack bound, $N_f^{-1}\log[\widetilde Z^{\rm av}(x+zh_r^{\rm per})/\widetilde Z^{\rm av}(x)]\le4\mathsf c(C_S+C_g)\varrho^{-1}\mathsf l^3|z|$. Summing,
\[\mathsf B_\iota\le n_\iota\Bigl[\frac{32\mathsf c(C_S+C_g)}{\varrho}\mathsf l^3|x|+\frac{128\,C_\times(7/2)\mathsf c^2x^2}{\varrho^2}\Bigr]\le N\Bigl[3C_wTq^2|x|+\frac{1152\,C_\times(7/2)\mathsf c^2}{\varrho^2\mathsf l^3}Tq^2x^2\Bigr],\]
since $72\cdot4\mathsf c(C_S+C_g)/\varrho=3C_w$; with $Tq^2\le q^3$ and the cost of H\"older's inequality this is the claim.
\end{proof}

\begin{proposition}[Slabs; the walls without a term in $x^2$ for products]\label{C-365-20}
\begin{enumerate}
\item[(a)] \emph{The slab inequality (exact).} On a periodic box $(\mathbb Z/\mathsf N\mathbb Z)^3$ let $c(p)$ be any real weights on the plaquettes and $Z(c)=\int e^{-\sum_pc(p)a_p}\,dU$. Let $\mathsf N=2k\,t$ with integers $k,t\ge1$, cut the box in one direction by the $2k$ planes of sites $x_\mu=n_0+jt$ into the closed slabs $\mathsf S_j$, and let $c_j^{\rm arr}$ be the weights that agree with $c$ on $\mathsf S_j$ and are symmetric under the reflections at all $2k$ planes. Then $Z(c)\le\prod_jZ(c_j^{\rm arr})^{1/(2k)}$; no sign of $c$ is supposed. For $c(p)=\beta-a^{-1}w(x(p))$ with $w$ real this reads $Z[w]\le\prod_jZ[w_j^{\rm arr}]^{1/(2k)}$, $w_j^{\rm arr}$ the function that agrees with $w$ on $\mathsf S_j$ and is even at all the planes.
\item[(b)] \emph{The walls for products.} Let $(K,\beta)\in\mathfrak d(g)$, $\ell\in\mathfrak L(g)$, $T\ge T_0$, $q=L^i\ge2T$, $\ell\ge2q$, a board of $T_\ell$ of side $q$ with mirrors $y_\mu\in s_\mu+q\mathbb Z$, and $T'$ the least power of $L$ that is $\ge2T$. Let $\varphi_\mu:\mathbb R/\ell\mathbb Z\to[0,1]$ have $\operatorname{Lip}\varphi_\mu\le16/T$ and $\varphi_\mu=1$ at distance $\ge T/2$ from $s_\mu+q\mathbb Z$, and $F(y)=\prod_\mu\varphi_\mu(y_\mu)$; let $x_0$, $x$ be real with $|x_0|+|x|\le\varrho/2$, and $C_V=16L(C_S+C_g)/\varrho+3C_\times(1)/(4\varrho T_0^3)$. Then
\[\log\widetilde Z^{\rm av}(x_0+xF)-\log\widetilde Z^{\rm av}(x_0+x)\le C_V|x|\,Tq^2N;\]
if every $\varphi_\mu$ has the period $q$ and is even at $s_\mu$ (as for the mirror copies of the centred $f_0$), the same holds for $\widetilde Z$ itself. No condition $x\le b$ is asked.
\end{enumerate}
\end{proposition}

\begin{proof}[Sketch of proof]
(a) Cut in the direction $\mu=1$. The kernel $K_j(\phi,\phi')$ of the slab $\mathsf S_j$, with half weights on its two faces and the bonds of the faces fixed to $\phi,\phi'$, is real, continuous and positive on a compact space, so a Hilbert--Schmidt operator, and Fubini gives $Z(c)=\operatorname{Tr}(K_0K_1\cdots K_{2k-1})$. The reflection at the mid-plane of a slab shows that $K_j^T$ is the kernel of the slab with mirrored weights, so $\operatorname{Tr}((K_jK_j^T)^k)=Z(c_j^{\rm arr})$; H\"older's inequality in the Schatten classes, $|\operatorname{Tr}(A_1\cdots A_{2k})|\le\prod_i\|A_i\|_{2k}$, is the claim. In the case used it is the classical chessboard argument: $\mathsf F(C_1,\dots,C_{2k})=\operatorname{Tr}(C_1C_2^TC_3C_4^T\cdots C_{2k-1}C_{2k}^T)$ is multilinear, invariant under a cyclic shift by one place with all letters transposed, and satisfies $|\mathsf F(C)|^2\le\mathsf F(C_1,\dots,C_k,C_k,\dots,C_1)\,\mathsf F(C_{2k},\dots,C_{k+1},C_{k+1},\dots,C_{2k})$ by Cauchy--Schwarz; after the normalisation $\mathsf F(C_i,\dots,C_i)=1$, among the words that maximise $|\mathsf F|$ one with the longest run of equal letters must have a run of length $\ge k$, whence the maximum is $\le1$. (b) Three steps, one direction at a time: $u_0=x_0+x$ and $u_\mu=x_0+x\prod_{\nu\le\mu}\varphi_\nu(y_\nu)$, so that $u_3=x_0+xF$.

In the direction $\mu$ the planes $y_\mu\in s_\mu+T'/2+T'\mathbb Z$ are planes of sites, because $q/a_K$ and $T'/a_K$ are odd, and they cut $T_\ell$ into the $\ell/T'$ closed slabs $\mathsf S_j=\{|y_\mu-s_\mu-jT'|\le T'/2\}$, $j\in\mathbb Z/(\ell/T')\mathbb Z$, an even number of them; as $q/T'$ is an integer, the mirrors are the mid-planes of the $\ell/q$ \emph{wall slabs}, those with $jT'\in q\mathbb Z$. The function $u_{\mu-1}$ does not read $y_\mu$, and $u_\mu-u_{\mu-1}=-x\prod_{\nu<\mu}\varphi_\nu(y_\nu)\,(1-\varphi_\mu(y_\mu))$ vanishes at distance $\ge T/2$ from the mirrors, hence off the wall slabs and near their faces, as $T'/2\ge T$. For a wall slab $\mathsf S_j$ put
\[w_j=(u_\mu-u_{\mu-1})\mathbf 1_{\mathsf S_j}=-x\prod_{\nu<\mu}\varphi_\nu(y_\nu)\,\bigl(1-\varphi_\mu(y_\mu)\bigr)\mathbf 1_{\mathsf S_j},\]
so that $u_\mu=u_{\mu-1}+\sum_{\text{wall slabs}}w_j$, and let $w_j^{\rm arr}$, the sum of the $\ell/T'$ mirror images of $w_j$ at the planes, one in each slab, be the even continuation of $w_j$ from $\mathsf S_j$. The even continuation of $u_\mu$ from a slab is $u_{\mu-1}$ for a slab that is not a wall slab and $u_{\mu-1}+w_j^{\rm arr}$ for a wall slab $\mathsf S_j$; so (a) at the real source $u_\mu$ gives
\[\log\frac{Z[u_\mu]}{Z[u_{\mu-1}]}\le\frac{T'}{\ell}\sum_{\text{wall slabs}}\log\frac{Z[u_{\mu-1}+w_j^{\rm arr}]}{Z[u_{\mu-1}]}.\]
The normalisers cancel exactly. Let $\mathcal F=\mathcal E^{\rm av}$ for the first assertion and $\mathcal F=\mathcal E$ in the periodic even case, $\widetilde Z^{\mathcal F}=e^{-E-\mathcal F}Z$ (that is, $\widetilde Z^{\rm av}$, respectively $\widetilde Z$), and $\mathcal F'(w)=\mathcal F(u_{\mu-1}+w)-\mathcal F(u_{\mu-1})$ as in Lemma~\ref{C-365-14}~(A2) at the base $u_{\mu-1}$, so that $Z[u_{\mu-1}+w]/Z[u_{\mu-1}]=e^{\mathcal F'(w)}\widetilde Z^{\mathcal F}(u_{\mu-1}+w)/\widetilde Z^{\mathcal F}(u_{\mu-1})$. The pieces of $w_j^{\rm arr}$ are images $\mathsf rw_j$ under reflections $\mathsf r\in\mathsf G$ at the planes, which fix $u_{\mu-1}$, so $\mathcal F'(\mathsf rw_j)=\mathcal F'(w_j)$ for $\mathcal F=\mathcal E^{\rm av}$ by (A1); in the periodic even case $1-\varphi_\mu$ is even at the mid-plane of $\mathsf S_j$, the pieces are the translates of $w_j$ by $T'\mathbb Ze_\mu$, which lies in $M\mathbb Z^3$ as $T'$ and $M$ are powers of $L$ with $T'>M$, and the same holds for $\mathcal E$ itself by (M1).

Hence $\mathcal F'(w_j^{\rm arr})=(\ell/T')\mathcal F'(w_j)+\mathrm{cr}_j$ and $\mathcal F'(\sum_jw_j)=\sum_j\mathcal F'(w_j)+\mathrm{cr}$, the sum over the wall slabs, and in $\log[\widetilde Z^{\mathcal F}(u_\mu)/\widetilde Z^{\mathcal F}(u_{\mu-1})]=\log(Z[u_\mu]/Z[u_{\mu-1}])-\mathcal F'(\sum_jw_j)$ the terms $\mathcal F'(w_j)$ cancel, leaving the cross terms. These are bounded by Lemma~\ref{C-365-14}~(A3), with $x$ replaced by $\zeta$, $|\zeta|<\varrho$, in the $w_j$ alone: each support lies within $T/2$ of the mid-plane of its slab and is covered by $(\ell/T')^2$ closed sup-balls of radius $R_0=T'/2$, and different supports are at sup-distance $\ge\Delta=T'-T\ge T'/2\ge4M$, so that $R_0/\Delta\le1$.

The Harnack bound at the base $u_{\mu-1}$, with $\|u_{\mu-1}\|_\sharp\le|x_0|+|x|\le\varrho/2$, and in the direction $w_j^{\rm arr}$, with $\|w_j^{\rm arr}\|_\sharp\le|x|$ as $40M\cdot16\sqrt3/T\le1$, costs $\frac83(C_S+C_g)\varrho^{-1}|x|\ell^3$ for each wall array, and $\#\{\text{wall slabs}\}\,T'/\ell=T'/q$; a slab needs no partition of unity across it, so (P3) and its term in $x^2$ do not enter. Summing the three steps with $2T\le T'<2LT$ and $|x|\le\varrho/2$ gives the constant $C_V$.
\end{proof}

\begin{proposition}[The main inequality; semi-convexity]\label{C-365-21}
Let $\ell'\in\mathfrak L(g)$ with $\ell'\ge2\mathsf l_1$ be a reference side, $T\ge\max\{T_0,\mathsf l\}$, $q=L^i\ge2T$, $\ell\in\mathfrak L(g)$ with $\ell\ge\max\{2q,\ell'\}$, $K\ge1$ and $b\in I(g)$, all numbers of Definition~\ref{C-365-7} being formed at $g$ with $\ell'$. Put $\omega(z)=\ell'^{-3}\log[\widetilde Z_{\ell'}(z)/\widetilde Z_{\ell'}(0)]$, let $\mu'=d'(0)$ for the $d$ of Lemma~\ref{C-365-15} at $(\ell_1,\ell_2)=(\ell,\ell')$, and for a profile $f$ let $\epsilon(z)=\mathcal E^{\rm av}_\ell(zf)+q^3[\omega'(0)-\mu']z$, a function holomorphic on $|z|<2\varrho$ with $\epsilon(0)=0$ and $\epsilon'(0)$ real.
\begin{enumerate}
\item[(a)] For every board of $T_\ell$ of side $q$, every offset and every profile $f\in\mathfrak P_T$, for $|z|\le r_1$ with $\operatorname{Re}z\le b$:
\[\Lambda(z)-\operatorname{Re}\epsilon(z)\le q^3\bigl[D_{K,\ell}(b-\operatorname{Re}z)-D_{K,\ell}(b)\bigr]+3C_wTq^2|\operatorname{Re}z|+C_Iq^3|z|^2 .\]
\item[(b)] $D''_{K,\ell}(b)\ge-C_0$ at every $b\in I(g)$. So $D_{K,\ell}+\frac12C_0y^2$ and $D_K+\frac12C_0y^2$ are convex on $I(g)$; $\mu_K=(D_K+\frac12C_0y^2)''$ and $\mu_{K,\ell}=(D_{K,\ell}+\frac12C_0y^2)''$ are positive Radon measures on $\operatorname{int}I(g)$; and for both, $\mu([y_1,y_2])\le2B'/t+C_0(y_2-y_1+t)$ whenever $[y_1-t,y_2+t]\subset I(g)$.
\end{enumerate}
\end{proposition}

\begin{proof}[Sketch of proof]
By the passage $(K,L^Kb)\in\mathfrak d(g)$, so Lemmas~\ref{C-365-14}, \ref{C-365-15}, \ref{C-365-17} and Propositions~\ref{C-365-18}, \ref{C-365-19} hold on both tori. The shift of the label gives on each torus, for real $|x|<2\varrho$, $\log\widetilde Z(x)-\log\widetilde Z(0)=\ell^3[p_{K,\ell}(b-x)-p_{K,\ell}(b)]-\mathcal E(x\cdot1)$; dividing by the volumes and subtracting,
\[\ell^{-3}\bigl[\log\widetilde Z_\ell(x)-\log\widetilde Z_\ell(0)\bigr]=\bigl[D_{K,\ell}(b-x)-D_{K,\ell}(b)\bigr]+\omega(x)-d(x).\]
$H(z)=\widetilde Z_{\ell'}(z)/\widetilde Z_{\ell'}(0)$ is holomorphic on $|z|<2\varrho$, $H(0)=1$ and $|H|\le e^{(C_S+C_g)\ell'^3}$ by (M1); Schwarz's lemma gives $|H-1|\le\frac12$ on $|z|\le r_2$, so $|\omega(x)-\omega'(0)x|\le C_\omega x^2$ for $|x|\le r_2/2$ and $|\omega''(0)|\le C_\omega$. For (a), with $z=x+i\tau$, the source formula gives $N\Lambda(z)=\operatorname{Re}\mathcal E^{\rm av}(w_z)+\log|\widetilde Z^{\rm av}(w_z)|-\log\widetilde Z(0)$. Then (i) (IM) for $\widetilde Z^{\rm av}$ at $u=xF$, $v=\tau(F_e-F_o)$ removes $\tau$ at the cost $C_I'\tau^2\ell^3$; (ii) the walls (Proposition~\ref{C-365-19}) replace $xF$ by the constant $x$ at the cost $N[3C_wTq^2|x|+C_2'q^3x^2]$; (iii) the identity, the bound of $\omega$ and Lemma~\ref{C-365-15} turn $\log\widetilde Z(x)-\log\widetilde Z(0)$ into at most $Nq^3\{[D_{K,\ell}(b-x)-D_{K,\ell}(b)]+[\omega'(0)-\mu']x+(C_\omega+2\mathsf A_1/h_1^2)x^2\}$; (iv) the cross terms between the cells, Lemma~\ref{C-365-14}~(A3) with $R_0=q/2$, $\Delta=T/4$, $\rho_3=\varrho$ and $C_\times(2q/T)\le M_\times(q/T)^3$, give $\operatorname{Re}\mathcal E^{\rm av}(w_z)\le N\operatorname{Re}\mathcal E^{\rm av}(zf)+NM_\times T^{-3}\varrho^{-2}q^3|z|^2$.

Adding and dividing by $N$ gives (a). For (b), semi-convexity in Proposition~\ref{C-365-18} at $u=0$, $v=1$ gives $\mathsf h''(0)\ge-2C_I'\ell^3$ for $\mathsf h(x)=\log\widetilde Z_\ell(x)$; differentiating the identity twice at $x=0$,
\[D''_{K,\ell}(b)=\ell^{-3}\mathsf h''(0)-\omega''(0)+d''(0)\ge-2C_I'-C_\omega-2\mathsf A_1/h_1^2=-C_0 .\]
$D_{K,\ell}$ is real-analytic in the label, and $D_K$ is its pointwise limit by Lemma~\ref{C-365-16}. For $\Phi=D+\frac12C_0y^2$ with $D\in\{D_K,D_{K,\ell}\}$, $\mu([y_1,y_2])=\Phi'(y_2+)-\Phi'(y_1-)$ is bounded by difference quotients over the length $t$, and $D\in[-B',0]$ by Lemma~\ref{C-365-16}~(c).
\end{proof}

\begin{lemma}[Cauchy's estimate in the chessboard seminorm; moments from the seminorm]\label{C-365-22}
On a board of $T_\ell$, at any $\beta\ge0$:
\begin{enumerate}
\item[(a)] If $X=\mathcal O(f)$ is a real function of one cell and $\tilde\epsilon$ is holomorphic near the disc $|z|\le r$ with $\tilde\epsilon(0)=0$, then $\|X-\tilde\epsilon'(0)\|_{\rm ch}\le r^{-1}\max_{|z|=r}\exp[\Lambda(z)-\operatorname{Re}\tilde\epsilon(z)]$.
\item[(b)] For functions $X_1,\dots,X_n$ of $n$ different cells and complex numbers $m_i$: $|\langle\prod_i(X_i-m_i)\rangle|\le\prod_i\|X_i-m_i\|_{\rm ch}$; $|\langle X_i\rangle-m_i|\le\|X_i-m_i\|_{\rm ch}$; $|\operatorname{Cov}(X_1,X_2)|\le2\|X_1-m_1\|_{\rm ch}\|X_2-m_2\|_{\rm ch}$; and the centred moments satisfy $|\langle\prod_i(X_i-\langle X_i\rangle)\rangle|\le\prod_i2\|X_i-m_i\|_{\rm ch}$.
\end{enumerate}
\end{lemma}

\begin{proof}[Sketch of proof]
(a) Pointwise on the configurations, $X-\tilde\epsilon'(0)=(2\pi i)^{-1}\oint_{|z|=r}e^{zX-\tilde\epsilon(z)}z^{-2}\,dz$; the Riemann sums converge uniformly since $X$ is bounded; the triangle inequality and the homogeneity of the seminorm, with $\|F\|_{\rm ch}\le\sup|F|$ for the remainder, give $\|X-\tilde\epsilon'(0)\|_{\rm ch}\le r^{-1}\max_{|z|=r}\|e^{zX-\tilde\epsilon(z)}\|_{\rm ch}$, and $\|e^{zX-\tilde\epsilon(z)}\|_{\rm ch}=e^{-\operatorname{Re}\tilde\epsilon(z)}(Z[w_z]/Z)^{1/N}$ by the source formula. (b) The first inequality is the chessboard estimate; the second is the first with one factor and $\|1\|_{\rm ch}=1$; the third follows from $\operatorname{Cov}(X_1,X_2)=\langle(X_1-m_1)(X_2-m_2)\rangle-(\langle X_1\rangle-m_1)(\langle X_2\rangle-m_2)$; and $\|X_i-\langle X_i\rangle\|_{\rm ch}\le\|X_i-m_i\|_{\rm ch}+|\langle X_i\rangle-m_i|\le2\|X_i-m_i\|_{\rm ch}$.
\end{proof}

\begin{proposition}[From a bound on increments of the pressure to the cube means, on a torus]\label{C-365-23}
In the setting of Proposition~\ref{C-365-21}, let $\mathsf s$, $\delta_2$, $\tilde C$, $\eta_+$, $x_0$ be real numbers with $\delta_2\ge0$, $\tilde C\ge0$, $\eta_+\ge0$ and $0<x_0\le b$, such that
\[D_{K,\ell}(b-x)-D_{K,\ell}(b)\le\mathsf sx+\delta_2|x|+\tilde Cx^2+\eta_+\qquad\text{for all real }|x|\le x_0 .\]
Put $\delta=3C_wT/q+\delta_2$, and let $q\ge\max\{(\tilde C+C_I)/(9C_w^2T^2),\,(3C_wT\min\{r_1,x_0\})^{-1/2}\}$. Then on every board of $T_\ell$ of side $q$, $\|x_c-m'\|_{\rm ch}\le e^{2+q^3\eta_+}\delta$ with the real number $m'=q^{-3}\epsilon'(0)+\mathsf s$, which depends on $\ell$ and on the profile. Hence, for profiles of $\mathfrak P_T$ in distinct cells,
\[|\operatorname{Cov}_{K,\ell,L^Kb}(x_c,x_{c'})|\le2e^{4+2q^3\eta_+}\delta^2,\qquad\Bigl|\bigl\langle\textstyle\prod_{i=1}^n(x_{c_i}-\langle x_{c_i}\rangle)\bigr\rangle_{K,\ell,L^Kb}\Bigr|\le\bigl(2e^{2+q^3\eta_+}\delta\bigr)^n .\]
\end{proposition}

\begin{proof}[Sketch of proof]
Put $\tilde\epsilon(z)=\epsilon(z)+q^3\mathsf sz$. For $|z|=r\le\min\{r_1,x_0\}$, so that $\operatorname{Re}z\le x_0\le b$, the main inequality and the hypothesis at $x=\operatorname{Re}z$ give, since $\delta_2\ge0$ and $\tilde C\ge0$ allow $\delta_2|x|\le\delta_2r$ and $\tilde Cx^2\le\tilde Cr^2$,
\[\Lambda(z)-\operatorname{Re}\tilde\epsilon(z)\le q^3\delta r+q^3(\tilde C+C_I)r^2+q^3\eta_+ .\]
Take $r=1/(q^3\delta)$; since $\delta\ge3C_wT/q$, $r\le1/(3C_wTq^2)\le\min\{r_1,x_0\}$ by the second floor; $q^3\delta r=1$; and $q^3(\tilde C+C_I)r^2=(\tilde C+C_I)/(q^3\delta^2)\le(\tilde C+C_I)/(9C_w^2T^2q)\le1$ by the first. Lemma~\ref{C-365-22}~(a) gives $\|\mathcal O(f)-\tilde\epsilon'(0)\|_{\rm ch}\le q^3\delta e^{2+q^3\eta_+}$; divide by $q^3$; for a profile of another cell the same holds with its own $\epsilon$ and $m'$; Lemma~\ref{C-365-22}~(b) gives the rest.
\end{proof}

\begin{proposition}[The window inequality, in one frame, at every cutoff]\label{C-365-24}
In the setting of Proposition~\ref{C-365-21}, let $b$ be a label with $I_q(b)\subset\operatorname{int}I(g)$, $q\ge q_1=\max\{2T,\,C_I/(9C_w^2T^2),\,(3C_wTr_1)^{-1/2}\}$, $K\ge1$, and
\[\delta_K(q,b)=3C_wT/q+\mu_K(I_q(b)).\]
\begin{enumerate}
\item[(a)] For all $\nu,\nu'\in\mathcal G^{\rm per}_{K,L^Kb}$, every board of $\mathbb L_K$ of side $q$ with any offset and profiles of $\mathfrak P_T$ in distinct cells $c,c',c_1,\dots,c_n$:
\[
\begin{multlined}
|\operatorname{Cov}_\nu(x_c,x_{c'})|\le2e^4\delta_K(q,b)^2,\qquad|\nu(x_c)-\nu'(x_c)|\le2e^2\delta_K(q,b),\qquad\Bigl|\nu\Bigl(\prod_i(x_{c_i} \\
-\nu(x_{c_i}))\Bigr)\Bigr|\le(2e^2\delta_K(q,b))^n .
\end{multlined}
\]
\item[(b)] The same holds with $\delta=3C_wT/q+C_0x_q+x_q^{-1}\Delta^2_{x_q}D_K(b)$, which is $\le\delta_K(q,b)$; and for every $\ell\in\mathfrak L(g)$ with $\ell\ge\ell'$,
\[
\begin{multlined}
x_q^{-1}\Delta^2_{x_q}D_K(b)\le x_q^{-1}\Delta^2_{x_q}D_{K,\ell}(b)+6C_wC_{\rm FS}Tq^2/\ell\le x_q\sup_{|y-b|\le x_q}D''_{K,\ell}(y) \\
+6C_wC_{\rm FS}Tq^2/\ell,\qquad C_{\rm FS}=3\ell'B':
\end{multlined}
\]
the bounds in every periodic Gibbs state are governed by the pressures of the two tori $T_\ell$, $T_{\ell'}$ at a side $\ell$ of the order $q^3$.
\end{enumerate}
\end{proposition}

\begin{proof}[Sketch of proof]
(a) \emph{A window gives the hypothesis of Proposition~\ref{C-365-23}.} $\Phi=D_K+\frac12C_0y^2$ is convex on $\operatorname{int}I(g)\supset I_q(b)$; with $\sigma_0=\Phi'(b-)$, for $0<x\le x_q$, $\Phi(b+x)-\Phi(b)\le\sigma_0x+x\mu_K([b,b+x))$ and $\Phi(b-x)-\Phi(b)\le-\sigma_0x+x\mu_K((b-x,b))$; since $D_K(b-x)-D_K(b)=\Phi(b-x)-\Phi(b)+C_0bx-\frac12C_0x^2$, the function $D_K$ has the hypothesis with $\mathsf s=C_0b-\sigma_0$, $\delta_2=\mu_K(I_q(b))\ge0$, $\tilde C=0$, $\eta_+=0$, $x_0=x_q$. \emph{From the tori to the states.} By Lemma~\ref{C-365-16}, $D_{K,\ell}=D_K+\mathsf w_\ell$ with $0\le\mathsf w_\ell\le\bar\eta(\ell)$, so $D_{K,\ell}$ has it with the tolerance $\eta_+=\bar\eta(\ell)$, and Proposition~\ref{C-365-23} gives $\|x_c-m'_\ell\|_{{\rm ch},\ell}\le e^{2+q^3\bar\eta(\ell)}\delta_K(q,b)$ for all $\ell\in\mathfrak L(g)$, $\ell\ge\max\{2q,\ell'\}$. The real numbers $m'_\ell$ are bounded at fixed $K$, since $|m'_\ell-\langle x_c\rangle_\ell|\le e^3\delta_K$ for large $\ell$ and $|x_c|$ is bounded. For finitely many profiles choose $\ell_k\to\infty$ with $m'_{\ell_k,i}\to m_{*,i}$; then $\limsup_k\|x_{c_i}-m_{*,i}\|_{{\rm ch},\ell_k}\le e^2\delta_K$, and the domination gives
\[\Bigl|\nu\Bigl(\prod_i(x_{c_i}-m_{*,i})\Bigr)\Bigr|\le(e^2\delta_K)^n\qquad\text{for every }\nu\in\mathcal G^{\rm per}_{K,L^Kb};\]
at $n=1$, $|\nu(x_c)-m_*|\le e^2\delta_K$ for $\nu$ and for $\nu'$, one common centre serving both; expanding as in Lemma~\ref{C-365-22}~(b) gives the three bounds. (b) In Proposition~\ref{C-365-23} the hypothesis enters only at $|x|\le r=1/(q^3\delta)\le x_q$. On $[b-r,b+r]$ the convex $\Phi$ lies below its chord, so $\Phi(b-x)-\Phi(b)\le-\mathsf s_rx+\frac12\Delta^2_r\Phi(b)$ with the slope $\mathsf s_r$ of the chord; and $\Delta^2_r\Phi(b)/r=r^{-1}\int_0^r\mu_K((b-u,b+u))\,du$ does not decrease in $r$, so $\frac12\Delta^2_r\Phi(b)\le r\,x_q^{-1}\Delta^2_{x_q}\Phi(b)$, which is $r(C_0x_q+x_q^{-1}\Delta^2_{x_q}D_K(b))$ and at most $r\mu_K(I_q(b))$. The comparison with the torus follows from $D_K\le D_{K,\ell}\le D_K+\bar\eta(\ell)$, $2\bar\eta(\ell)/x_q=6C_wC_{\rm FS}Tq^2/\ell$, and $\Delta^2_sD\le s^2\sup D''$ for the real-analytic $D_{K,\ell}$.
\end{proof}

\begin{lemma}[A sum rule; near and far parts; a pile recurs at all large cutoffs]\label{C-365-25}
Let one frame, $0<g\le\gamma$, a compact interval $W\subset I(g)$ of positive length and $\ell'\in\mathfrak L(g)$ with $\ell'\ge2\mathsf l_1$ be given, let $k_7\ge1$ be an integer with $\varsigma_7=L^{-k_7}\le2/(3\sqrt3)$, and let $\chi$, $\eta$, $\mathsf c_7$, $c_T$, $c_N$ be formed with this $\varsigma_7$ as in Definition~\ref{C-365-7}.
\begin{enumerate}
\item[(i)] \emph{Exact.} For $K\ge k_7$, every side $\ell>6\varsigma_7$ and $b\ge0$:
\[
\begin{multlined}
p''_{K,\ell}(b)=\varsigma_7^{-3}S^T_{2;K,\ell,L^Kb}(\chi,1)=V^{\rm near}_{K,\ell}(b)+V^{\rm far}_{K,\ell}(b),\qquad V^{\rm near}=\varsigma_7^{-3}S^T_2(\chi,\eta),\quad V^{\rm far} \\
=\varsigma_7^{-3}S^T_2(\chi,1-\eta).
\end{multlined}
\]
\item[(ii)] $|V^{\rm near}_{K,\ell}-V^{\rm near}_{K,\ell'}|\le c_N$ on $I(g)$ for all $K\ge k_7$ and $\ell\in\mathfrak L(g)$; and for each side $\ell$, $\omega_\ell(K,K'')=\sup_W|V^{\rm far}_{K,\ell}-V^{\rm far}_{K'',\ell}|\to0$ as $K,K''\to\infty$.
\item[(iii)] For every closed interval $I$ and $\delta>0$ with $I^{+\delta}=[\min I-\delta,\max I+\delta]\subset\operatorname{int}W$, all $K,K''\ge k_7$ and all $\ell\in\mathfrak L(g)$ with $\ell\ge\ell'$:
\[\mu_K(I)\le\mu_{K''}(I^{+\delta})+(|I|+2\delta)\bigl[2c_N+\omega_\ell(K,K'')+\omega_{\ell'}(K,K'')\bigr]+\frac{24\ell'B'}{\delta\ell}.\]
\end{enumerate}
\end{lemma}

\begin{proof}[Sketch of proof]
(i) 
\[
p''_{K,\ell}=\ell^{-3}L^{2K}\operatorname{Var}(\sum_pa_p)=a_K^{-5}\mathsf N^{-3}\operatorname{Var}(\sum_pa_p),
\]
while $S^T_2(\chi,1)=a_K^{-2}\sum_p\chi(x(p))\operatorname{Cov}(a_p,\sum_{p'}a_{p'})$. By translation invariance the covariance depends on the orientation of $p$ alone, the three orientations summing to $\mathsf N^{-3}\operatorname{Var}(\sum_pa_p)$, and for each orientation $\sum_p\chi(x(p))=(\varsigma_7/a_K)^3$, the sums of tents of the integer half-width $\varsigma_7/a_K$; the splitting is $1=\eta+(1-\eta)$. (ii) The supports of $\chi$ and $\eta$ lie in a ball of radius $3\sqrt3\varsigma_7\le2$, and $\|\chi\|_\sharp,\|\eta\|_\sharp,\|1-\eta\|_\sharp\le\mathsf c_7$, so the transfer (M7) at $n=2$ between the two tori, at the first-passage data given by the passage, bounds the near parts; the supports of $\chi$ and $1-\eta$ are at sup-distance $\varsigma_7$, so the convergence on fixed tori (M8) gives the far parts. (iii) Take a $C^2$ function $1_I\le\eta_I\le1_{I^{+\delta}}$ with $\|\eta_I''\|_1\le8/\delta$, the indicator of $[\min I-\delta/2,\max I+\delta/2]$ smoothed by the kernel $\frac{15}{16}(1-t^2)_+^2$ at the scale $\delta/2$. Since $\mu_K\ge0$, $\mu_K(I)\le\int\eta_I\,d\mu_K$ and $\int\eta_I\,d\mu_{K''}\le\mu_{K''}(I^{+\delta})$, and
\[\int\eta_I\,d(\mu_K-\mu_{K''})=\int\eta_I''\bigl[(D_K-D_{K,\ell})-(D_{K''}-D_{K'',\ell})\bigr]+\int\eta_I\,\bigl(D''_{K,\ell}-D''_{K'',\ell}\bigr),\]
where the bracket lies in $[-\bar\eta(\ell),\bar\eta(\ell)]$ by Lemma~\ref{C-365-16}, and by (i) $D''_{K,\ell}=(V^{\rm near}_{K,\ell}-V^{\rm near}_{K,\ell'})+V^{\rm far}_{K,\ell}-V^{\rm far}_{K,\ell'}$, so that $|D''_{K,\ell}-D''_{K'',\ell}|\le2c_N+\omega_\ell+\omega_{\ell'}$ on $W$. In words: a concentration of $\mu_K$ at one large cutoff recurs at all large cutoffs.
\end{proof}

\begin{lemma}[Dilations by real factors]\label{C-365-26}
For real $s>0$, $a>0$ and $\beta>0$ the image under $\mathrm{dil}_s$ maps the periodic Gibbs states at $(a,\beta)$ onto the periodic Gibbs states at $(sa,\beta)$, and $S^T_n[\nu](f_1,\dots,f_n)=s^n\,S^T_n[(\mathrm{dil}_s)_*\nu](f_1(\cdot/s),\dots,f_n(\cdot/s))$.
\end{lemma}

\begin{proof}[Sketch of proof]
The exponent $H_B=-\beta\sum_{p\in\mathrm{Pl}[B]}a_p$ of a Gibbs kernel contains no spacing. By the exact dilation facts (M9), $\mathrm{dil}_s$ carries $\mathrm{Pl}[B]$ onto $\mathrm{Pl}[\mathrm{dil}_sB]$, $a_{\mathrm{dil}_sp}(\mathrm{dil}_sU)=a_p(U)$, the product Haar measure on the bonds of $B$ to that on the bonds of $\mathrm{dil}_sB$, and the configuration that is $V$ on $B$ and $U$ off $B$ to the one that is $\mathrm{dil}_sV$ on $\mathrm{dil}_sB$ and $\mathrm{dil}_sU$ off it; hence
\[\mathrm{dil}_s(\gamma_BF)=\gamma_{\mathrm{dil}_sB}(\mathrm{dil}_sF),\]
and $((\mathrm{dil}_s)_*\nu)(\gamma_{\mathrm{dil}_sB}\mathrm{dil}_sF)=\nu(\gamma_BF)=\nu(F)$; as $B\mapsto\mathrm{dil}_sB$ and $F\mapsto\mathrm{dil}_sF$ are onto, the image is a Gibbs state. Invariance under the translations by a subgroup $\Gamma$ of finite index passes to $s\Gamma$, of the same index; $\mathrm{dil}_{1/s}$ is the inverse. The formula for the cumulants follows from $\mathcal O_a(f)(U)=s\,\mathcal O_{sa}(f(\cdot/s))(\mathrm{dil}_sU)$.
\end{proof}

\begin{lemma}[One pressure for both frames; the reference tori compared]\label{C-365-27}
In a Setting, let $K\ge k_7$, $\hat K\ge1$, $\lambda=L^K/\hat L^{\hat K}$ and $J$ a closed interval with $J\subset W_1$ and $\lambda J\subset\hat W$, so that $\lambda\in[\lambda_-,\lambda_+]$. Then $p''_{K,\infty}(J)=\lambda^4\hat p''_{\hat K,\infty}(\lambda J)$ and
\[D_K''(J)\le\lambda^4\hat D''_{\hat K}(\lambda J)+c_2|J| .\]
\end{lemma}

\begin{proof}[Sketch of proof]
By Lemma~\ref{C-365-16}~(b), $p_{K,\infty}(b)=L^{3K}P(L^Kb)$ and $\hat p_{\hat K,\infty}(\hat b)=\hat L^{3\hat K}P(\hat L^{\hat K}\hat b)$ with one convex $P$; for $J=[y_1,y_2]$, $p''_{K,\infty}(J)=L^{4K}[P'(L^Ky_2+)-P'(L^Ky_1-)]=\lambda^4\hat p''_{\hat K,\infty}(\lambda J)$, since $\hat L^{\hat K}\lambda=L^K$. With $\mathsf r_1=p''_{K,\ell'}$, $\mathsf r_2=\hat p''_{\hat K,\hat\ell'}$ this gives $D''_K(J)-\lambda^4\hat D''_{\hat K}(\lambda J)=\int_J[\lambda^5\mathsf r_2(\lambda b)-\mathsf r_1(b)]\,db$, and both $\mathsf r_1(b)$ and $\lambda^5\mathsf r_2(\lambda b)$ are $a_K^{-5}(\text{box})^{-3}$ times a variance of the action at the one coupling $\beta=L^Kb=\hat L^{\hat K}\lambda b$. In the units of the second frame the tent becomes $\chi(\cdot/\lambda)$, $\hat{\mathcal O}=\lambda^{-1}\mathcal O$ and $S_2=\lambda^2\hat S_2$; the sum rule of Lemma~\ref{C-365-25} holds on both reference boxes, as $6\lambda\varsigma_7<\hat\ell'$.

The far parts are bounded by $c_F$ and $\hat c_F$, the bound of (M8) reading of the functions only the spread $\lambda\varsigma_7\ge\lambda_-\varsigma_7$ and the norms $\le\hat{\mathsf c}_8$. For the near parts take $\nu\in\mathcal G^{\rm per}$ at $(a_K,\beta)$, which exists by (M9), and its image $\hat\nu$ under the dilation by the real factor $\lambda$, a periodic Gibbs state at $(\hat a_{\hat K},\beta)$ by Lemma~\ref{C-365-26}, with $S_2[\nu](\chi,\eta)=\lambda^2\hat S_2[\hat\nu](\chi(\cdot/\lambda),\eta(\cdot/\lambda))$; by the passage in each frame both pairs are first-passage data, and the transfer in each frame between the reference torus and the state ($3\sqrt3\lambda\varsigma_7\le2$) gives
\[|\lambda^5\mathsf r_2(\lambda b)-\mathsf r_1(b)|\le\varsigma_7^{-3}\bigl[c_T\mathsf c_7^2+\lambda_+^2\hat c_T\hat{\mathsf c}_8^2+c_F+\lambda_+^2\hat c_F\bigr]=c_2 .\qedhere\]
\end{proof}

\begin{proposition}[No pile of specific heat]\label{C-365-28}
In a Setting, with $(g_1,W_1,W')$ and $\operatorname{dist}(W',\partial W_1)=\varepsilon$: for every $0<\delta\le\delta_{\max}$ there is an integer $K^\circ$ with
\[\mu_K(I)\le C_\flat(|I|+\delta)\qquad\text{for all }K\ge K^\circ\text{ and all closed intervals or points }I\subset W'.\]
So $K_1(\delta)$ of Definition~\ref{C-365-7} exists, and it does not increase in $\delta$. Of the convergence on fixed tori (item (M8)) the proof reads two clauses, both at $n=2$, for pairs with disjoint supports and without a rate: (1) the convergence in the cutoff, uniformly in the label, of the cumulant $S^T_2$ at the pair $(\chi,1-\eta)$ on the tori $T_\ell$, $\ell\in\mathfrak L(g_1)$, $\ell\ge\ell'$, over $W_1$, and at the pair of the tent and the plateau of half-width $1/\hat L$ on the tori $T_{\hat\ell}$, $\hat\ell\in\hat{\mathfrak L}(\hat g)$, $\hat\ell\ge\hat\ell'$, over $\hat W$; (2) the bound clause, which gives $c_F$, for the one pair $(\chi,1-\eta)$ on $T_{\ell'}$ over $W_1$, and $\hat c_F$ on $T_{\hat\ell'}$, uniformly in $\hat K\ge1$, in $\hat b\in\hat W$ and in the pairs $(\chi(\cdot/\lambda),1-\eta(\cdot/\lambda))$, $\lambda\in[\lambda_-,\lambda_+]$, whose spread $\lambda\varsigma_7\ge\lambda_-\varsigma_7$ and whose norms $\le\hat{\mathsf c}_8$ are all that the clause reads of them. Thus on $T_{\hat\ell'}$ two kinds of pairs are read: the pair of half-width $1/\hat L$ through the convergence, and the family of the dilated pairs through the bound.
\end{proposition}

\begin{proof}[Sketch of proof]
\emph{Thresholds.} In the first frame Lemma~\ref{C-365-25} is used with $g_1$, $W_1$, $\ell'$ and $k_7$ of Definition~\ref{C-365-7}. Take $\ell(\delta)\in\mathfrak L(g_1)$, $\ell(\delta)\ge\ell'$, with $24\ell'B'/(\delta\ell(\delta))\le\delta/2$, and an integer $K^\circ\ge k_7$ with $\omega_{\ell(\delta)}(K,K'')+\omega_{\ell'}(K,K'')\le1/6$ for all $K,K''\ge K^\circ$ (Lemma~\ref{C-365-25}~(ii)). In the second frame Lemma~\ref{C-365-25} is used with $\hat g$, $\hat W$, $\hat\ell'$ and $\hat k_7=1$, so that the tent and the plateau have the half-width $1/\hat L\le2/(3\sqrt3)$ there and the constant of the near parts is $\hat c_N$. Put $\hat\delta=\lambda_+\delta$; take $\hat\ell\in\hat{\mathfrak L}(\hat g)$, $\hat\ell\ge\hat\ell'$, with
\[\frac{24\hat\ell'\hat B'}{\hat\delta\hat\ell}\le\frac\delta2,\]
and an integer $\hat K_*\ge1$ with $\hat\omega_{\hat\ell}(\hat K,\hat K'')+\hat\omega_{\hat\ell'}(\hat K,\hat K'')\le1/(10\lambda_+)$ for all $\hat K,\hat K''\ge\hat K_*$. \emph{One interval.} Let $I=[y-\delta/2,y+\delta/2]$, $y\in W'$, $J=I^{+\delta}\subset\operatorname{int}W_1$, $K\ge K^\circ$. Let $R=\lfloor|\hat W|/(12\lambda_+\delta)\rfloor\ge|\hat W|/(24\lambda_+\delta)$ and $t_r=\hat w_-+|\hat W|/4+(r-\frac12)6\lambda_+\delta$, $r=1,\dots,R$. As $\log L/\log\hat L$ is irrational, Kronecker's theorem gives integers $K''_r\ge K^\circ$ and $\hat K_r\ge\hat K_*$ with
\[\bigl|\lambda_ry-t_r\bigr|<\lambda_+\delta/2,\qquad\lambda_r=L^{K''_r}/\hat L^{\hat K_r};\]
then $\lambda_r\in[\lambda_-,\lambda_+]$, and the intervals $\mathsf I_r=(\lambda_rJ)^{+\hat\delta}$ are pairwise disjoint and lie in the middle half of $\hat W$. \emph{The chain.} Lemma~\ref{C-365-25}~(iii) in the first frame, for $I$ with the margin $\delta$ and the side $\ell(\delta)$, from $K$ to $K''_r$: $\mu_K(I)\le\mu_{K''_r}(J)+3\delta(2c_N+\frac16)+\frac\delta2=\mu_{K''_r}(J)+(6c_N+1)\delta$. Lemma~\ref{C-365-27} between the frames, with $\hat D''\le\hat\mu$ (as $\hat C_0\ge0$), $\hat\mu\ge0$ and $\lambda_r\le\lambda_+$: $\mu_{K''_r}(J)=D''_{K''_r}(J)+3C_0\delta\le\lambda_+^4\hat\mu_{\hat K_r}(\lambda_rJ)+3(c_2+C_0)\delta$. Lemma~\ref{C-365-25}~(iii) in the second frame, for the interval $\lambda_rJ$, of length $3\lambda_r\delta\le3\lambda_+\delta$, with the margin $\hat\delta$ and the side $\hat\ell$, from $\hat K_r$ to the one cutoff $\hat K_*$, whose added terms are at most $5\lambda_+\delta\,(2\hat c_N+\frac1{10\lambda_+})+\frac\delta2=10\lambda_+\hat c_N\delta+\delta$: $\hat\mu_{\hat K_r}(\lambda_rJ)\le\hat\mu_{\hat K_*}(\mathsf I_r)+10\lambda_+\hat c_N\delta+\delta$. Together:
\[\mu_K(I)\le\lambda_+^4\bigl[\hat\mu_{\hat K_*}(\mathsf I_r)+10\lambda_+\hat c_N\delta+\delta\bigr]+\bigl[6c_N+3(c_2+C_0)+1\bigr]\delta\qquad\text{for every }r.\]
Take the mean over $r$: the $\mathsf I_r$ are pairwise disjoint and lie in the middle half of $\hat W$, whose mass for $\hat\mu_{\hat K_*}$ is at most $8\hat B'/|\hat W|+\frac34\hat C_0|\hat W|\le\hat{\mathfrak m}$ by the mass bound of Proposition~\ref{C-365-21}~(b) in the second frame at $t=|\hat W|/4$; so $R^{-1}\sum_r\hat\mu_{\hat K_*}(\mathsf I_r)\le\hat{\mathfrak m}/R\le24\lambda_+\hat{\mathfrak m}\delta/|\hat W|$. This is $\mu_K(I)\le C_\flat\delta$. \emph{Any closed $I\subset W'$} is covered by at most $|I|/\delta+1$ such intervals. The property at $\delta$ implies it at every $\delta'>\delta$, so $K_1$ does not increase.
\end{proof}

\begin{lemma}[The labels; the exact dilation]\label{C-365-29}
\begin{enumerate}
\item[(i)] If $W\subset I(g)$ is nonempty and compact, then $W\subset\operatorname{int}I(g_1)$ for every $g_1<g$ with $g_1^{-2}+\bar B<\min W$, and such $g_1$ exist.
\item[(ii)] Every $b>\gamma^{-2}+\bar B$ lies in $\operatorname{int}I(g_2)$ for some $g_2\le\gamma$; hence for every $b_\infty>0$ there are an integer $s\ge0$ and $g_2\le\gamma$ with $L^sb_\infty\in\operatorname{int}I(g_2)$.
\item[(iii)] For integers $0\le s\le K$ the dilation $x\mapsto L^sx$ carries $\mathbb L_K$ onto $\mathbb L_{K-s}$, $\mathcal G^{\rm per}_{K,\beta}$ onto $\mathcal G^{\rm per}_{K-s,\beta}$, $T_\ell$ at the cutoff $K$ with its law onto $T_{L^s\ell}$ at the cutoff $K-s$, boards of side $q$ onto boards of side $L^sq$, $\mathfrak P_T(\square_c)$ onto $\mathfrak P_{L^sT}(L^s\square_c)$ by $\tilde f=f(\cdot/L^s)$, and the label $b$ into $L^sb$; and $x_c=L^{4s}\tilde x_c$.
\end{enumerate}
\end{lemma}

\begin{proof}[Sketch of proof]
(i) $I(g)$ is open on the left and $W$ is compact, so $\min W>g^{-2}+\bar B$; for $g_1^{-2}\in(g^{-2},\min W-\bar B)$ one has $g_1<g$, $g_1^{-2}+\bar B<\min W$ and $X_{g_1}-\bar B>X_g-\bar B\ge\max W$. (ii) Put $g_2^{-2}=b-\bar B-\epsilon$ with $\epsilon=\min\{1,b-\bar B-\gamma^{-2}\}$; then $g_2\le\gamma$ and $g_2^{-2}+\bar B<b$, and since $L(\bar B-B_\sharp)=\bar B$,
\[X_{g_2}-\bar B>b\iff(L-1)b>2\bar B+L\epsilon;\]
as $b>4+\bar B$ (because $\gamma\le\frac12$), $(L-1)b>4(L-1)+(L-1)\bar B\ge L+2\bar B\ge L\epsilon+2\bar B$. For $b_\infty>0$ take $s$ with $L^sb_\infty>\gamma^{-2}+\bar B$. (iii) The dilation is a graph isomorphism of the lattices, the Wilson weight reads the graph and $\beta$ alone, so the laws of the tori correspond, and the periodic Gibbs states correspond by (M9); $\beta=L^Kb=L^{K-s}(L^sb)$; $\tilde{\mathcal O}_{K-s}(\tilde f)=a_{K-s}^{-1}\sum_pf(x(p))a_p=L^{-s}\mathcal O_K(f)$, hence $\tilde x_c=(L^sq)^{-3}\tilde{\mathcal O}(\tilde f)=L^{-4s}x_c$; the collars and the Lipschitz constant scale with $L^s$, and offsets in $a_K\mathbb Z^3$ go to offsets in $a_{K-s}\mathbb Z^3$.
\end{proof}

\begin{proposition}[The K\"all\'en--Lehmann form]\label{C-365-30}
Let $\sigma'$ be the limit $\sigma$ of a convergent Gibbs sequence or a component $\sigma_\xi$ of it, and $\mathbf s$ its radial kernel.
\begin{enumerate}
\item[(a)] $\mathbf s$ is completely monotone on $(0,\infty)$; so $\mathbf s\ge0$ does not increase, and $\kappa'=\lim_{r\to\infty}\mathbf s(r)\in[0,\infty)$ exists.
\item[(b)] $r\mapsto r(\mathbf s(r)-\kappa')$ is completely monotone; there is exactly one positive Borel measure $\rho$ on $[0,\infty)$ with $\mathbf s(r)=\kappa'+\int e^{-mr}(4\pi r)^{-1}\,d\rho(m)$.
\item[(c)] $\rho(\{0\})=\lim_{r\to\infty}4\pi r(\mathbf s(r)-\kappa')$, and $\rho([0,m])\le4\pi e\,m^{-1}(\mathbf s(1/m)-\kappa')$ for $m>0$.
\item[(d)] \emph{Cube means, with no further hypothesis.} Let $T\ge\max\{T_0,\mathsf l\}$, $q=L^i\ge2T$, a board of $\mathbb R^3$ of side $q$ (a partition of $\mathbb R^3$ into the cubes $\square_c$ of side $q$ centred at the points $o+qc$, $c\in\mathbb Z^3$, with an offset $o\in\mathbb R^3$; $c$ is the index of the cell) and profiles $f_c\in\mathfrak P_T(\square_c)$, $f_{c'}\in\mathfrak P_T(\square_{c'})$ of two distinct cells be given. For every $0<r_0\le T/4$
\[0\le q^{-6}\Bigl[\sigma'(f_c,f_{c'})-\kappa'\int f_c\int f_{c'}\Bigr]\le24\pi\,r_0\bigl(\mathbf s(r_0)-\kappa'\bigr)/q,\]
and if the two cells share a face, the middle member is $\ge\rho(\{0\})(1-T/q)^6/(4\sqrt6\,\pi q)$. So $q$ times the middle member tends to $0$ as $q\to\infty$, uniformly in the pairs of distinct cells and in the profiles, if and only if $\rho(\{0\})=0$.
\item[(e)] \emph{The converse; the measure that attains the bounds.} The following are equivalent: (i) there are $T\ge\max\{T_0,\mathsf l\}$, $C$ and $j_0$ with $L^{j_0}\ge2T$ and $\sigma'(\theta w_i,w_i)\le C(T/q_i)^2$ for all $i\ge j_0$, the $w_i$ of Definition~\ref{C-365-4}~(7) being formed with this $T$; (ii) $\mathbf s(r)=O(r^{-2})$ as $r\to\infty$; (iii) $\kappa'=0$ and $\rho([0,m])=O(m)$ as $m\to0$; (iv) there are $A_\mu$ and $r_*>0$ with $\mu_h(Q_r)\le A_\mu\|h\|_1^2r^2$ for every real $h\in C^1_c(\mathbb E_+)$ and $0<r\le r_*$, $\mu_h$ formed in $\mathcal H(\sigma',e_3)$. Under them $\rho(\{0\})=0$, so that $\rho\ge\mathsf Z\delta_0$ with $\mathsf Z>0$ is excluded, and so is $\rho([0,m])\ge c_\rho m^\vartheta$ for all small $m$ with $c_\rho>0$ and $\vartheta<1$; the measure $\rho=c_\rho\,dm$ on $[0,\infty)$, $c_\rho>0$, whose kernel is $\mathbf s(r)=c_\rho/(4\pi r^2)$, satisfies (ii) and (iii) and attains both bounds in order.
\end{enumerate}
\end{proposition}

\begin{proof}[Sketch of proof]
We use Bernstein's theorem (a function on $(0,\infty)$ is completely monotone if and only if it is $\int e^{-tE}d\mu(E)$ with a unique positive $\mu$), that $-c'$ is completely monotone with $c$, and that a pointwise finite limit of completely monotone functions is completely monotone (Helly's theorem applied to the measures $e^{-t_0E}d\mu_n$). (a) By (M10), $\sigma'$ is invariant under $\mathbb R^3\rtimes\mathrm O(3)$ and reflection positive, so the reconstruction applies: for $f_\epsilon(x)=\eta_\epsilon(x_3)\mathsf g(x_1,x_2)$ with $0\le\eta_\epsilon\in C^1_c((s-\epsilon,s+\epsilon))$ of mass $1$ and real $\mathsf g\in C^1_c(\mathbb R^2)$, $c_\epsilon(t)=\sigma'(\theta f_\epsilon,\tau_{te_3}f_\epsilon)=\int e^{-tE}d\mu_{f_\epsilon}$ is completely monotone, and by the kernel it tends to $\mathsf P_{\mathsf g}(2s+t)$, $\mathsf P_{\mathsf g}(u)=\iint\mathsf g(y)\mathsf g(y')\mathbf s((u^2+|y-y'|^2)^{1/2})\,dy\,dy'$; letting $\mathsf g$ tend to a point mass gives (a). (b) For a radial $\mathsf g\ne0$ put $G=\mathsf g*\mathsf g$, with radial profile $\hat G$, and $\mathsf g_R=R^{-1}\mathsf g(\cdot/R)$, so that $\mathsf P_R(u)=2\pi\int_0^\infty\hat G(\varrho'/R)\mathbf s((u^2+\varrho'^2)^{1/2})\varrho'd\varrho'$; with $\partial_u\mathbf s((u^2+\varrho'^2)^{1/2})=(u/\varrho')\partial_{\varrho'}\mathbf s((u^2+\varrho'^2)^{1/2})$ one integration by parts gives
\[-\mathsf P_R'(u)=2\pi u\Bigl[\hat G(0)\mathbf s(u)+\int_0^\infty\hat G'(v)\,\mathbf s((u^2+R^2v^2)^{1/2})\,dv\Bigr]\ \longrightarrow\ 2\pi\hat G(0)\,u\,(\mathbf s(u)-\kappa')\qquad(R\to\infty),\]
and $-\mathsf P_R'$ is completely monotone; Bernstein's theorem gives $\rho$, normalised so that $4\pi r(\mathbf s(r)-\kappa')=\int e^{-mr}d\rho(m)$. (c) Monotone convergence, and $\int e^{-mr}d\rho(m)\ge e^{-1}\rho([0,1/r])$. (d) As $\mathbf s\ge\kappa'$ and $f_c,f_{c'}\ge0$, the middle member is $q^{-6}\iint f_c(x)f_{c'}(y)(\mathbf s(|x-y|)-\kappa')\,dx\,dy\ge0$. The supports are at distance $\ge T/4\ge r_0$, where $\mathbf s(|x-y|)-\kappa'\le r_0(\mathbf s(r_0)-\kappa')/|x-y|$, since $r(\mathbf s(r)-\kappa')$ does not increase by (b); and with $0\le f_c,f_{c'}\le1$ on the cells, $u=(y-x)/q$ and the tent $\chi_{\mathbf d}(u)=\prod_\mu(1-|u_\mu-d_\mu|)_+$ about the difference $\mathbf d\ne0$ of the indices of the two cells, $q^{-6}\iint_{\square_c\times\square_{c'}}|x-y|^{-1}dx\,dy=q^{-1}\int\chi_{\mathbf d}(u)|u|^{-1}du$, which is $\le q^{-1}$ if $|\mathbf d|_\infty\ge2$ and $\le q^{-1}\int_{|u|\le2\sqrt3}|u|^{-1}du=24\pi/q$ otherwise.

For cells that share a face, $4\pi r(\mathbf s(r)-\kappa')=\int e^{-mr}d\rho\ge\rho(\{0\})$, the profiles are $1$ on the concentric cubes of side $q-T$, and two points of these cubes are at distance $\le\sqrt6q$. If $\rho(\{0\})=0$, then $r(\mathbf s(r)-\kappa')\le\epsilon_1$ for $r\ge r_{\epsilon_1}$, and the pairs of points of $\square_c\times\square_{c'}$ at distance $<r_{\epsilon_1}$ have measure $\le8\pi q^2r_{\epsilon_1}^4$; so $q$ times the middle member is $\le24\pi\epsilon_1+8\pi(\mathbf s(T/4)-\kappa')r_{\epsilon_1}^4q^{-3}$, for every $\epsilon_1>0$. (e) (i) gives (ii): $\theta w_i,w_i\ge0$ have disjoint supports in a box of diameter $\sqrt6q_i$, $\int w_i\ge(1-3T/(4q_i))^3$, and $\mathbf s$ does not increase, so $C(T/q_i)^2\ge\sigma'(\theta w_i,w_i)\ge\mathbf s(\sqrt6q_i)(1-3T/(4q_i))^6$; for $\sqrt6q_i\le r<\sqrt6q_{i+1}$, $\mathbf s(r)\le\mathbf s(\sqrt6q_i)$ and $q_i^{-2}<6L^2/r^2$. (ii) gives (i): if $\mathbf s(r)\le A_{\mathbf s}r^{-2}$ for $r\ge r_0$, take $T\ge\max\{T_0,\mathsf l,4r_0\}$; the supports of $\theta w_i$ and $w_i$ are at distance $\ge T/4\ge r_0$, so $\sigma'(\theta w_i,w_i)\le A_{\mathbf s}q_i^{-2}\int\chi_{e_3}(u)|u|^{-2}du\le3\pi A_{\mathbf s}q_i^{-2}$, as $\chi_{e_3}(u)\le u_3$ on its support, which lies in $\{|u|\le\sqrt6,\ u_3>0\}$. (ii) gives (iii) by (c). (iii) gives (ii): if $\kappa'=0$ and $\rho([0,m])\le C_\rho m$ for $0<m\le m_0$, then $e^{-mr}=e^{-m_0r}+r\int_m^{m_0}e^{-ur}du$ and Fubini give $\int_{[0,m_0]}e^{-mr}d\rho(m)\le C_\rho m_0e^{-m_0r}+C_\rho/r$, and for $r\ge1$ the rest of $4\pi r\mathbf s(r)$ is $\le e^{-m_0(r-1)}4\pi\mathbf s(1)$. (iii) and (iv): for real $h\in C^1_c(\mathbb E_+)$ the kernel, the form (b) and the formula
\[\frac{e^{-m\sqrt{\tau^2+|w|^2}}}{4\pi\sqrt{\tau^2+|w|^2}}=\int_{\mathbb R^2}e^{i\mathsf q\cdot w}\,\frac{e^{-\tau\omega_m(\mathsf q)}}{2\omega_m(\mathsf q)}\,\frac{d^2\mathsf q}{(2\pi)^2},\qquad\omega_m(\mathsf q)=(m^2+|\mathsf q|^2)^{1/2},\ \tau>0,\ w\in\mathbb R^2,\]
give for $t\ge0$ and $v\perp e_3$
\[\langle\psi_h,e^{-tH}\mathcal U(v)\psi_h\rangle=\sigma'(\theta h,\tau_{te_3+v}h)=\kappa'\Bigl(\int h\Bigr)^2+\int d\rho(m)\int e^{-t\omega_m(\mathsf q)-iv\cdot\mathsf q}\,\frac{|\hat h_m(\mathsf q)|^2}{2\omega_m(\mathsf q)}\,\frac{d^2\mathsf q}{(2\pi)^2},\]
with $\hat h_m(\mathsf q)=\int h(y)e^{-\omega_m(\mathsf q)y_3+i\mathsf q\cdot y_\perp}dy$; as a finite measure on $[0,\infty)\times\mathbb R^2$ is determined by its transform, $\mu_h$ is $\kappa'(\int h)^2\delta_{(0,0)}$ plus the images under $\mathsf q\mapsto(\omega_m(\mathsf q),-\mathsf q)$ of the measures $|\hat h_m|^2(2\omega_m)^{-1}d^2\mathsf q/(2\pi)^2$, integrated in $d\rho(m)$. As $|\hat h_m|\le\|h\|_1$ and $\int_{\omega_m(\mathsf q)\le r}(2\omega_m(\mathsf q))^{-1}d^2\mathsf q/(2\pi)^2=(r-m)_+/(4\pi)$, this gives $\mu_h(Q_r)\le\kappa'\|h\|_1^2+(4\pi)^{-1}r\,\rho([0,r])\|h\|_1^2$, so (iii) gives (iv). Conversely, take $0\le h\ne0$ with support in $\{y_3\le R_h,\ |y_\perp|\le R_h\}$ and $r\le\min\{r_*,1/(2R_h)\}$: the points $(\omega_m(\mathsf q),-\mathsf q)$ with $m\le r/2$ and $|\mathsf q|\le r/2$ lie in $Q_r$, and there $|\hat h_m(\mathsf q)|\ge e^{-1}\cos(\frac12)\|h\|_1$; so $A_\mu\|h\|_1^2r^2\ge\mu_h(Q_r)\ge\kappa'\|h\|_1^2+e^{-2}\cos^2(\frac12)(32\pi)^{-1}r\,\rho([0,r/2])\|h\|_1^2$, which gives (iii). Under (iii), $\rho(\{0\})\le\rho([0,m])\to0$ as $m\to0$; and $\int_0^\infty e^{-mr}c_\rho\,dm/(4\pi r)=c_\rho/(4\pi r^2)$.
\end{proof}

\begin{lemma}[Three mirrors; the spectral form]\label{C-365-31}
\begin{enumerate}
\item[(a)] Let $\sigma'$ be a real, symmetric, normalized functional on $\mathfrak S(\mathbb R^3)$, multilinear on products of test spaces with disjoint closures, whose one-point function is $\sigma'(f)=c\int f$ and whose connected functions obey a bound table as in Section~\YMMref{90}{695}. Assume that $\sigma'$ is invariant under the translations of $\mathbb R^3$ and reflection positive at the three coordinate mirrors through $0$; let $X,Y\in\mathfrak S(\mathbb R^3)$ have their supports in $\{|x|_\infty<R\}$ and $G(a)=\sigma'(X\cdot\tau_aY)$. Then $G(a)$ converges as $\min_\mu a_\mu\to\infty$, jointly in the three coordinates, to $\langle[\theta X_0],\mathcal P(\{0\})[Y_0]\rangle$ in $\mathcal H(\sigma',e_3)$, with $X_0=\tau_{-(R+\frac12)e_3}X$ and $Y_0=\tau_{(R+\frac12)e_3}Y$; likewise in each of the eight orthants, $X$ and $Y$ being exchanged in the directions in which $a_\mu\to-\infty$.
\item[(b)] Let $\sigma'$ obey a bound table with one-point constant $0$, be invariant under the translations of $\mathbb R^3$ and reflection positive at $(e_3,0)$, and let $\sigma'(\theta w_i,w_i)\le C_{\rm LH}(T/q_i)^2$ for $i\ge j_0$, with $q_{j_0}\ge2T$. Then for real $h\in C^1_c(\mathbb E_+)$
\[\mu_h(Q_r)\le12.46\,(1-3T/(4q_{j_0}))^{-6}L^2C_{\rm LH}\|h\|_1^2T^2r^2\quad(0<r\le1/q_{j_0}),\qquad\mathcal P(\{0\})\psi_h=0 .\]
\end{enumerate}
Neither the cubic group nor $\mathrm O(3)$ nor a kernel is used.
\end{lemma}

\begin{proof}[Sketch of proof]
(a) In the picture $e_3$, for $t=a_3-2R-1\ge0$, translation invariance and the reconstruction give $G(a)=\langle[\theta X_0],e^{-tH}\mathcal U(a_\perp)[Y_0]\rangle=\int e^{-tE+ia_\perp\cdot\mathsf q}\,dm_3$, $m_3=\langle[\theta X_0],\mathcal P(\cdot)[Y_0]\rangle$, so $\sup_{a_\perp}|G(a)-G_3(a_\perp)|\le\varepsilon_3(t)=\int_{E>0}e^{-tE}d|m_3|\downarrow0$, with $G_3(a_\perp)$ the part at $E=0$; the same in the pictures $e_1$, $e_2$. Changing one coordinate at a time between two points with all coordinates $\ge A$,
\[|G(a)-G(a')|\le2\sum_\mu\varepsilon_\mu(A-2R-1)\to0,\]
so $G$ converges; its limit is the limit of the Ces\`aro means $A^{-3}\int_{[A,2A]^3}G$, whose three brackets in the picture $e_3$ tend to $1_{E=0}$, $1_{\mathsf q_1=0}$, $1_{\mathsf q_2=0}$, giving $m_3(\{(0,0,0)\})$. In the other orthants write $G(a)=\sigma'(\tau_{-a}X\cdot Y)$. (b) Let $w=w_i$, $q=q_i$, $i\ge j_0$; $w\ge0$ has its support in $\{T/4\le x_3\le q-T/4\}$. The convolution $w*h=\int w(y)\tau_yh\,dy=\int h(y)\tau_yw\,dy$ is read twice, the Riemann sums converging in $\|\cdot\|_\sharp$ and $h\mapsto\psi_h$ being linear and bounded by the bound table:
\[
\begin{multlined}
\psi_{w*h}=\mathsf G_w(H,P)\psi_h,\quad\mathsf G_w(E,\mathsf q)=\int w(y)e^{-y_3E+iy_\perp\cdot\mathsf q}\,dy,\qquad\psi_{w*h} \\
=\int h(y)e^{-y_3H}\mathcal U(y_\perp)\psi_w\,dy .
\end{multlined}
\]
The operators $e^{-tH}\mathcal U(v)$ are contractions, so 
\[
\int|\mathsf G_w|^2d\mu_h=\|\psi_{w*h}\|^2\le\|h\|_1^2\sigma'(\theta w,w)\le\|h\|_1^2C_{\rm LH}(T/q)^2.
\]
On $Q_{1/q}$, as $\varphi_0$ is even, 
\[
\mathsf G_w=q^{-3}\int\varphi_0(y_3-\frac q2)e^{-y_3E}dy_3\prod_{k=1,2}\int\varphi_0(y_k)\cos(y_k\mathsf q_k)dy_k\ge e^{-1}\cos^2(\tfrac12)(1-3T/(4q))^3;
\]
and $e^2/\cos^4(\frac12)\approx12.4577$. For $1/(Lq_i)<r\le1/q_i$, $\mu_h(Q_r)\le\mu_h(Q_{1/q_i})$ and $q_i^{-2}<L^2r^2$. Finally $\|\mathcal P(\{0\})\psi_h\|^2=\mu_h(\{0\})\le\mu_h(Q_r)\to0$.
\end{proof}

\begin{lemma}[Limits of Gibbs sequences, with free constants]\label{C-365-32}
Let a frame, $0<g'\le\gamma$, $T\ge\max\{T_0,\mathsf l\}$ and numbers $C_{\rm LH},c_{\rm LH},q_{\rm LH}>0$ be given, and let $(K_j,\beta_j,\nu_j)$ be a convergent Gibbs sequence at $g'$ with the following property: for every $q=L^i\ge\max\{2T,q_{\rm LH}\}$ there is $j(q)$ such that for all $j\ge j(q)$, all $\nu,\nu'\in\mathcal G^{\rm per}_{K_j,\beta_j}$, every board of $\mathbb L_{K_j}$ of side $q$ with any offset and profiles of $\mathfrak P_T$ in two distinct cells $c,c'$: $|\operatorname{Cov}_\nu(x_c,x_{c'})|\le C_{\rm LH}(T/q)^2$ and $|\nu(x_c)-\nu'(x_c)|\le c_{\rm LH}T/q$. Then the conclusions 1--9 of Theorem~\ref{C-365-8}~(c) hold for $\sigma$ and every $\sigma_\xi$, with $C_{\rm LH}$ in the place of $C_2$, $L^{j_0}\ge\max\{2T,q_{\rm LH}\}$, $r_0=\sqrt6L^{j_0}$ and $C_{\mathbf s}=6L^2(1-3T/(4L^{j_0}))^{-6}C_{\rm LH}$.
\end{lemma}

\begin{proof}[Sketch of proof]
\emph{The pair of cells.} Fix $q=q_i\ge\max\{2T,q_{\rm LH}\}$. On the board of $\mathbb L_{K_j}$ of side $q$ with the offset $(0,0,\frac12q+\frac12a_{K_j})$, which lies in $a_{K_j}\mathbb Z^3$ because $q/a_{K_j}$ is odd, the two cells with the centres $(0,0,\pm\frac12q+\frac12a_{K_j})$ have a common face, and $f_0(\cdot\mp\frac12qe_3)$ is moved from them by $\frac12a_{K_j}\le T/8$ for large $j$; so, $\varphi_0$ being even, $\mathcal O(w_i)$ and $\mathcal O(\theta w_i)$ are cube means $x_c$, $x_{c'}$ of two distinct cells. \emph{The lattice bounds pass to $\sigma$ and to every $\sigma_\xi$}, because they name the pair $(K_j,\beta_j)$ alone and so hold for the $\nu_j$ and for the approximants $\rho_j$ of $\sigma_\xi$:
\[\mathfrak c_\xi\int w_i=\sigma_\xi((w_i))=\lim_j(\rho_j-\nu_j)(x_c),\qquad|(\rho_j-\nu_j)(x_c)|\le c_{\rm LH}T/q,\qquad\int w_i\ge(1-\tfrac{3T}{4q})^3\ge(\tfrac58)^3,\]
for every $q$ above the floor, so $\mathfrak c_\xi=0$; this is shown before the kernel is used, so that $S^{\sigma'}_2(f,h)=\sigma'(f,h)$. Further $S^\sigma_2(\theta w_i,w_i)=\lim_j\operatorname{Cov}_{\nu_j}(x_{c'},x_c)$ and $\sigma_\xi(\theta w_i,w_i)=\lim_j[\operatorname{Cov}_{\rho_j}(x_{c'},x_c)+(\rho_j-\nu_j)(x_{c'})(\rho_j-\nu_j)(x_c)]$, whose second term tends to $0$; this is conclusion 2. \emph{From two cells with a common face to the kernel.} By Proposition~\ref{C-365-30}, $\mathbf s\ge0$ does not increase; $\theta w_i$, $w_i\ge0$ have disjoint supports in a box of diameter $\sqrt6q$, so $C_{\rm LH}(T/q_i)^2\ge\sigma'(\theta w_i,w_i)\ge\mathbf s(\sqrt6q_i)(1-3T/(4q_i))^6$; for $r\ge r_0$ take $i\ge j_0$ with $\sqrt6q_i\le r<\sqrt6Lq_i$, so $q_i^{-2}<6L^2/r^2$ and $\mathbf s(r)\le C_{\mathbf s}T^2/r^2$.

Hence $\kappa'=0$, conclusion 3 is Proposition~\ref{C-365-30}~(b), conclusion 5 follows from $|x-y|\ge|a|-R_f-R_h\ge r_0$ on the supports, and 6, 7 from Proposition~\ref{C-365-30}~(c). \emph{One field.} For real $h\in C^1_c(\mathbb E_+)$, $\|e^{-tH}\psi_h\|^2=\sigma'(\theta h,\tau_{2te_3}h)\le\|h\|_1^2\mathbf s(2t)$, the supports being at distance $\ge2t$; so $e^{-tH}\psi_h\to0$, which kills the projection on $\ker H$; $\int_0^\infty\langle\psi_h,e^{-tH}\psi_h\rangle\,dt=\|H^{-1/2}\psi_h\|^2$ is finite by the bound at large $t$; and $e^{-2}\langle\psi_h,1_{[0,\epsilon]}(H)\psi_h\rangle\le\|e^{-H/\epsilon}\psi_h\|^2$. For conclusion 9 rotate $a$ to the $e_3$ axis by $\mathrm O(3)$ and put $X_0=\tau_{-(R_X+1)e_3}X$, $h_0=\tau_{(R_h+1)e_3}h$, $t=|a|-R_X-R_h-2$:
\[\sigma'(X\cdot\tau_ah)=\langle[\theta X_0],e^{-tH}\psi_{h_0}\rangle,\qquad|\sigma'(X\cdot\tau_ah)|\le\|[\theta X_0]\|\,\|e^{-tH}\psi_{h_0}\|\le\mathsf N(X)\|h\|_1C_{\mathbf s}^{1/2}T/(2t),\]
where $\|[\theta X_0]\|^2=\sigma'(X\cdot\theta_{e_3}X)\le\mathsf N(X)^2$, finite by the bound tables since the supports of $X\cdot\theta_uX$ lie in one ball and those of $X$ and $\theta_uX$ are at distance $\ge2$;
in the moment--cumulant formula for $S^{\sigma'}_{n+1}(f_1,\dots,f_n,\tau_ah)$ exactly one factor holds $\tau_ah$, and it is $\sigma'(X_J\cdot\tau_ah)$ for a subtuple $X_J$ or $\sigma'((\tau_ah))=0$.
\end{proof}

\begin{lemma}[Growing tori, with free constants]\label{C-365-33}
Let a frame, $0<g'\le\gamma$, $T\ge\max\{T_0,\mathsf l\}$, numbers $C_{\rm LH},q_{\rm LH}>0$ and $j_0$ with $L^{j_0}\ge\max\{2T,q_{\rm LH}\}$ be given, and let $S$ be a continuum limit at $g'$ along growing tori with a defining sequence $(K_j,\ell_j,\beta_j)$ that has the following property: for every $q=L^i\ge\max\{2T,q_{\rm LH}\}$ there is $j(q)$ such that for all $j\ge j(q)$, $\ell_j\ge2q$ and $|\operatorname{Cov}_{K_j,\ell_j,\beta_j}(x_c,x_{c'})|\le C_{\rm LH}(T/q)^2$ for profiles of $\mathfrak P_T$ in two distinct cells of every board of $T_{\ell_j}$ of side $q$. Then the conclusions 1--3 of Theorem~\ref{C-365-8}~(d) hold with $C_{\rm LH}$ in the place of $C_2$.
\end{lemma}

\begin{proof}[Sketch of proof]
By (M10), $\sigma_S$ obeys a bound table with one-point constant $0$, is invariant under the translations and the cubic group, and is reflection positive at the coordinate mirrors, so Lemma~\ref{C-365-31} applies at $u=e_1,e_2,e_3$. For $q=q_i$ and large $j$, as in the proof of Lemma~\ref{C-365-32} on the board of $T_{\ell_j}$ of side $q$ with the offset $(0,0,\frac12q+\frac12a_{K_j})$ ($2q$ divides $\ell_j$, and the supports lie in $(-\ell_j/4,\ell_j/4)^3$), $\theta w_i$ and $w_i$ are $q^{-3}$ times profiles of two distinct cells, so
\[\sigma_S(\theta w_i,w_i)=S_2(\theta w_i,w_i)=\lim_jS^T_{2;K_j,\ell_j,\beta_j}(\theta w_i,w_i)\le C_{\rm LH}(T/q_i)^2;\]
the cubic group carries this to the mirrors $(e_1,0)$, $(e_2,0)$. The spectral clauses are Lemma~\ref{C-365-31}~(b). For conclusion 3, Lemma~\ref{C-365-31}~(a) at $Y=(h)$ gives the limit $\langle[\theta X_0],\mathcal P(\{0\})\psi_{h_0}\rangle=0$, $h_0=\tau_{(R+\frac12)e_3}h\in C^1_c(\mathbb E_+)$, and $0$ likewise in the other orthants; in the moment--cumulant formula exactly one factor holds $\tau_ah$, and it is $\sigma_S(X_J\cdot\tau_ah)\to0$ or $\sigma_S((\tau_ah))=0$.
\end{proof}

\begin{proof}[Proof of Theorem~\ref{C-365-8}]
\emph{Step 1: the Setting is consistent.} $g_1$ exists by Lemma~\ref{C-365-29}~(i), hence $\varepsilon$ exists; Proposition~\ref{C-365-28} is applied to $(g_1,W_1,W')$, with $\operatorname{dist}(W',\partial W_1)=\varepsilon$. Since $C_w^-\le C_w$, $r_1^-\le r_1$, $C_I\le C_I^+$ and $T\ge T_0$, the floor $q\ge q_0$ gives
\[\text{(f1)}\quad q\ge\frac{4C_\flat+C_I}{9C_w^2T^2},\qquad\text{(f2)}\quad x_q=\frac1{3C_wTq^2}\le\min\{r_1,\varepsilon/2\},\]
so $I_q(b)\subset W'\subset\operatorname{int}I(g_1)$ for $b\in W$, and $\delta=\delta^\circ(q)\le\delta_{\max}$ has $C_\flat\delta\le\frac32C_w^-T_0/q\le\frac32C_wT/q$. The number $C_\flat$ is formed from $\lambda_+$, $\hat{\mathfrak m}$, $\hat c_N$, $c_N$, $c_2$ and $C_0$, none of which depends on $T$.

\emph{Step 2: clause (a).} Let $b\in W$ and $K\ge K_0(q)=K_1(\delta)$. Proposition~\ref{C-365-28} gives $\mu_K(I_q(b))\le C_\flat(2x_q+\delta)$, and $2C_\flat x_q\le\frac32C_wT/q$ is (f1); so
\[\delta_K(q,b)=3C_wT/q+\mu_K(I_q(b))\le6C_wT/q .\]
Proposition~\ref{C-365-24}~(a) at $g_1$ applies, as $q\ge q_1$ by $q\ge2T$, (f1) and (f2): $2e^4\cdot36=72e^4$ and $2e^2\cdot6=12e^2$. The clause names the states at $(K,L^Kb)$, which do not depend on the reference coupling.

\emph{Step 3: clause (b).} Let $\ell\in\mathfrak L(g)$, $\ell\ge\ell_0(q)\ge\ell'$; as $\mathfrak L(g_1)$ is closed upward, $\ell\in\mathfrak L(g_1)$. $\Psi=D_K+\frac12C_0y^2$ and $\Phi=D_{K,\ell}+\frac12C_0y^2$ are convex on $W_1$ and $0\le\Phi-\Psi\le\bar\eta=\bar\eta(\ell)$. With $\sigma_0=\Psi'(b-)$, $t=(\bar\eta/C_\flat)^{1/2}\le\varepsilon/2$ (this is $\ell\ge12\ell'B'/(C_\flat\varepsilon^2)$) and $0<x\le\varepsilon/2$, the slopes of a convex function do not decrease:
\[\frac{\Phi(b+x)-\Phi(b)}x\le\frac{\Phi(b+x+t)-\Phi(b+x)}t\le\frac{\Psi(b+x+t)-\Psi(b+x)+\bar\eta}t\le\sigma_0+C_\flat(x+t+\delta)+\frac{\bar\eta}t,\]
by Proposition~\ref{C-365-28} on $[b,b+x+t]\subset W'$, and the mirror bound holds on the left; $C_\flat t+\bar\eta/t=2(C_\flat\bar\eta)^{1/2}$. So $D_{K,\ell}$ has the hypothesis of Proposition~\ref{C-365-23} with $\mathsf s=C_0b-\sigma_0$, $\delta_2=2(C_\flat\bar\eta)^{1/2}+C_\flat\delta\ge0$, $\tilde C=C_\flat\ge0$, $\eta_+=0$, $x_0=\varepsilon/2<b$. Its floors are (f1), (f2); $2(3C_\flat\ell'B'/\ell)^{1/2}\le\frac32C_wT/q$ is $\ell\ge\frac{16}3C_\flat\ell'B'q^2/(C_wT)^2$; so $\delta_2\le3C_wT/q$, $\delta\le6C_wT/q$, and the bounds follow as in Step 2.

\emph{Step 4: clause (c).} Take $g'=g$ in Lemma~\ref{C-365-32}. For $q=L^i\ge\max\{2T,q_0\}$ and large $j$ we have $K_j\ge K_0(q)$ and $b_j\in W$; clause (a), which names the pair $(K_j,\beta_j)$ alone, gives the hypothesis of Lemma~\ref{C-365-32} with
\[(C_{\rm LH},c_{\rm LH},q_{\rm LH})=(C_2,\ 12e^2C_w,\ q_0),\]
and its conclusions are clause (c).

\emph{Step 5: clause (d).} Take $g'=g$ in Lemma~\ref{C-365-33}. For $q=L^i\ge\max\{2T,q_0\}$ and large $j$: $K_j\ge K_0(q)$, $b_j\in W$, and $\ell_j\in\mathfrak L(g)$, $\ell_j\ge\ell_0(q)\ge2q$. The two floors do not depend on each other, so no relation between $K_j$ and $\ell_j$ is used, and clause (b) gives the hypothesis of Lemma~\ref{C-365-33} with $(C_{\rm LH},q_{\rm LH})=(C_2,q_0)$; its conclusions are clause (d).

\emph{Step 6: clauses (e)\,1 and (e)\,2.} (e)\,1 is Lemma~\ref{C-365-29}~(ii). For (e)\,2, Lemma~\ref{C-365-29}~(iii) carries the states at $(K,L^Kb)$ and their boards of side $q$ to states at $(K-s,L^{K-s}\cdot L^sb)$ and boards of side $L^sq$, with
\[\operatorname{Cov}_\nu(x_c,x_{c'})=L^{8s}\operatorname{Cov}_{\tilde\nu}(\tilde x_c,\tilde x_{c'}),\qquad\nu(x_c)-\nu'(x_c)=L^{4s}(\tilde\nu-\tilde\nu')(\tilde x_c),\]
and $n$-cell moments carrying $L^{4ns}$; $(L^sT)/(L^sq)=T/q$; apply (a), (b) in the Setting built on $(g_2,W^{(2)})$ at $(L^sT,L^sq,K-s,L^s\ell)$.

\emph{Step 7: clause (e)\,3.} Each $b_\infty\in\mathsf B$ has, by (e)\,1, a pair $(s,g_2)$ and a compact interval $W^{(2)}\subset\operatorname{int}I(g_2)$ holding $L^sb_\infty$ in its interior; finitely many of the open sets $L^{-s}\operatorname{int}W^{(2)}$ cover $\mathsf B$, say $(s_i,g^{(i)},W^{(i)})$, $i\le n_{\mathsf B}$, each with its Setting in the frame $L$ and one second frame for all. Put
\[C_{\mathsf B}=\max_iL^{8s_i}72e^4\bigl(C_w^{+,(i)}\bigr)^2,\quad q_{\mathsf B}=\max_iL^{-s_i}q_0^{(i)},\quad K_{\mathsf B}(q)=\max_i\bigl[s_i+K_0^{(i)}(L^{s_i}q)\bigr],\]
and let $\ell_{\mathsf B}(q)$ be the maximum over $i$ of $L^{-s_i}$ times $\ell_0^{(i)}(L^{s_i}q)$ formed with $C_w^{-,(i)}T_0$ in the place of $C_wT$, which bounds $\ell_0^{(i)}(L^{s_i}q)$ at the width $L^{s_i}T$ because $(L^{s_i}q)^2/(C_wL^{s_i}T)^2=q^2/(C_wT)^2$. None of these depends on $T$. A side $L^{s_i}\ell=2L^{m+s_i}\ge\ell'^{(i)}$ lies in $\mathfrak L(g_1^{(i)})\subset\mathfrak L(g^{(i)})$ by (M5), so (b) at $g^{(i)}$ speaks of it; $T\ge\max\{T_0,\mathsf l\}$ gives $L^{s_i}T\ge\max\{T_0,\mathsf l\}$; and $(2C_{\mathsf B})^{1/2}\ge L^{4s_i}12e^2C_w^{(i)}$. Step 6 gives (e)\,3.

\emph{Step 8: clause (e)\,4.} At $g'\le\gamma$ the labels of every Gibbs sequence, and of every defining sequence of a limit along growing tori, lie in the compact set $\mathsf B=I^{\rm lab}(g')$ by (M5). For $q=L^i\ge\max\{2T,q_{\mathsf B}\}$ and large $j$, $K_j\ge K_{\mathsf B}(q)$ and $\ell_j\ge\ell_{\mathsf B}(q)\ge2q$, so (e)\,3 gives the hypotheses of Lemmas~\ref{C-365-32} and~\ref{C-365-33} with
\[(C_{\rm LH},c_{\rm LH},q_{\rm LH})=\bigl(C_{\mathsf B},(2C_{\mathsf B})^{1/2},q_{\mathsf B}\bigr),\]
and these lemmas, in which the theorems on limits are taken at $g'$ itself, give (c) and (d) at $g'$. The lemmas with free constants are needed here: no single window, dilated or not, holds $I^{\rm lab}(g')$, whose ends have a quotient $\ge L$, while the quotient of the ends of $I(g)$ is at most $L$ with the lower end open; so (e)\,4 does not follow from (c), (d) applied to the sequence taken whole, and the only lattice input of (c), (d) is a bound expressed through $(K_j,\beta_j)$ and $\ell_j$.
\end{proof}

\begin{lemma}[Mean actions are slopes of the pressure; both ends are taken; the plateau]\label{C-365-34}
Let $a>0$ and $\beta>0$.
\begin{enumerate}
\item[(i)] For $\nu\in\mathcal G^{\rm per}_{a,\beta}$ the number $\mathsf a(\nu)=\nu(a_p)$ is the same for all plaquettes $p$, and $P(\beta+t)-P(\beta)\ge-3\mathsf a(\nu)t$ for all $t>-\beta$; so $-3\mathsf a(\nu)\in[P'(\beta-),P'(\beta+)]$.
\item[(ii)] There are $\nu_\pm\in\mathcal G^{\rm per}_{a,\beta}$ with $-3\mathsf a(\nu_\pm)=P'(\beta\pm)$; so $3\sup\{|\mathsf a(\rho)-\mathsf a(\nu)|:\rho,\nu\in\mathcal G^{\rm per}_{a,\beta}\}=J(\beta)$.
\item[(iii)] At $a=a_K$, for $(K,\beta)\in\mathfrak d(g)$ and $\nu\in\mathcal G^{\rm per}_{K,\beta}$, the function $\mathsf e$ of Theorem~\ref{C-365-9}~(e) exists, takes its values $\nu$-almost surely in $[\inf_\rho\mathsf a(\rho),\sup_\rho\mathsf a(\rho)]$, an interval of length $J(\beta)/3$, and $3a_K^{-4}\operatorname{Rie}_K(f)(\mathsf e-\mathsf a(\nu))$ is a conditional expectation of $\mathcal O_K(f)-\nu(\mathcal O_K(f))$ given $\mathcal I$, for every bounded $f$ with compact support.
\end{enumerate}
\end{lemma}

\begin{proof}[Sketch of proof]
(i) The lattice isometries act transitively on the plaquettes, carry $a_p$ to $a_{\mathsf rp}$, and leave every periodic Gibbs state invariant (M9); so $\nu(a_p)$ is one number. The tangent inequality (M9) for the function $A=-(a_{p_{12}(0)}+a_{p_{13}(0)}+a_{p_{23}(0)})$ of the three plaquettes at one site, whose array over the periodic box of $N$ sites a side is $-\sum_pa_p$, reads
\[\log Z_N(\beta+t)-\log Z_N(\beta)\ge-3\mathsf a(\nu)\,tN^3-144\,\beta N^2-72\,|t|N^2,\]
where $Z_N$ is the partition function of that box, which does not depend on the spacing; divide by $N^3$ and let $N\to\infty$ (Lemma~\ref{C-365-16}). (ii) Take points $\beta_k\downarrow\beta$ at which $P$ is differentiable and $\nu_k\in\mathcal G^{\rm per}_{a,\beta_k}$, which exist by (M9); by (i), $-3\mathsf a(\nu_k)=P'(\beta_k)$; a subsequence converges weakly, its limit is a periodic Gibbs state at $\beta$ by the closedness in (M9), and $a_p$ being continuous, $-3\mathsf a(\nu_+)=\lim_kP'(\beta_k)=P'(\beta+)$; likewise $\beta_k\uparrow\beta$ gives $\nu_-$. (iii) By Proposition~\YMMref{106.3}{765} (M9), at a first-passage datum the set $\mathcal G^{\rm per}_{K,\beta}$ is a class at the cutoff $K$ that obeys the bound system at $g$, and with $\nu$ it contains $k\nu/\nu(k)$ for every bounded $\mathcal I$-measurable $k\ge0$ with $\nu(k)>0$; with the invariance used in (i) this gives $\nu(ka_p)=\nu(ka_{p_0})$ for all plaquettes $p$, so one $\mathcal I$-measurable $\mathsf e$ serves all plaquettes, and the statements on conditional expectations given $\mathcal I$ (M9) place its values between the infimum and the supremum of the $\mathsf a(\rho)$, an interval of length $J(\beta)/3$ by (ii); linearity gives the conditional expectation of $\mathcal O_K(f)$, since $a_K^{-1}\sum_pf(x(p))=3a_K^{-4}\operatorname{Rie}_K(f)$.
\end{proof}

\begin{proof}[Proof of Theorem~\ref{C-365-9}]
\emph{Step 1: clause (a).} The first bound is Proposition~\ref{C-365-28}; $|D_K''|(I)\le\mu_K(I)+C_0|I|$. For the convex $\Phi=D_K+\frac12C_0y^2$,
\[\Delta^2_s\Phi(b)=\int(s-|y-b|)_+\,d\mu_K(y)=\int_0^s\mu_K((b-u,b+u))\,du\le C_\flat(s^2+\delta s),\]
and $\Delta^2_sD_K=\Delta^2_s\Phi-C_0s^2\ge-C_0s^2$ because $\mu_K\ge0$. Since $K_1$ does not increase, the bounds hold for every $\delta>\delta_0(K)$, hence at $\delta_0(K)$, and $\delta_0(K)\le\delta$ once $K\ge K_1(\delta)$. For the tent $\phi_{\delta'}$, $\mathsf p''(b)=\int\phi_{\delta'}(b-y)\,d\mu_K(y)-C_0$, and the layer-cake formula with the bound on the open intervals $(b-\delta'(1-\delta't),b+\delta'(1-\delta't))$ gives $\int\phi_{\delta'}(b-y)\,d\mu_K(y)\le C_\flat(1+\delta_0/\delta')$; and $\mathsf p(b)-D_K(b)=\int_0^{\delta'}\phi_{\delta'}(t)\Delta^2_tD_K(b)\,dt$ with $\int_0^{\delta'}\phi_{\delta'}(t)t^2\,dt=\delta'^2/12$ and $\int_0^{\delta'}\phi_{\delta'}(t)t\,dt=\delta'/6$.

\emph{Step 2: clause (b).} $p_{K,\ell'}$ is real-analytic, so the atom of $\mu_K=(D_K+\frac12C_0y^2)''$ at $b$ is the jump of $p'_{K,\infty}$ there, and $p_{K,\infty}(b)=L^{3K}P(L^Kb)$ by Lemma~\ref{C-365-16}:
\[\mu_K(\{b\})=p'_{K,\infty}(b+)-p'_{K,\infty}(b-)=L^{4K}J(L^Kb)\le C_\flat\delta .\]
For the limit of $\beta^4J(\beta)$: by Lemma~\ref{C-365-29}~(ii) every point of $[1,L]$ has a dilate $L^{s}b_\infty$ in the interior of a window $W^{(\iota)}\subset\operatorname{int}I(g^{(\iota)})$; finitely many of the sets $L^{-s_\iota}\operatorname{int}W^{(\iota)}$ cover $[1,L]$, each window with its Setting. Every $\beta\ge L^{1+\max s_\iota}$ is $L^{K'}y$ with $y\in W^{(\iota)}$ for some $\iota$ and an integer $K'$ that tends to $\infty$ with $\beta$; so $\beta^4J(\beta)=y^4\mu^{(\iota)}_{K'}(\{y\})\le(\max W^{(\iota)})^4C_\flat^{(\iota)}\delta$ once $K'\ge K_1^{(\iota)}(\delta)$, for every $\delta$.

\emph{Step 3: clause (c).} Lemma~\ref{C-365-34}~(i), (ii) at $a=a_K$, $\beta=L^Kb$. With $e(\nu)=3a_K^{-4}\mathsf a(\nu)$: $\nu(\mathcal O(f))=a_K^{-1}\mathsf a(\nu)\sum_pf(x(p))=\operatorname{Rie}_K(f)e(\nu)$, and the tangent inequality in the label reads
\[p_{K,\infty}(b')-p_{K,\infty}(b)=L^{3K}\bigl[P(L^Kb')-P(L^Kb)\bigr]\ge-3\mathsf a(\nu)L^{4K}(b'-b)=-(b'-b)e(\nu),\]
since $p'_{K,\infty}(b\pm)=L^{4K}P'(L^Kb\pm)$; the end points are taken by $\nu_\pm$.

\emph{Step 4: clause (d).} By (c), $\nu(x_c)-\nu'(x_c)=q^{-3}\operatorname{Rie}_K(f)(e(\nu)-e(\nu'))$, and for $0\le f\le1$ vanishing off a cube of side $\le q-a_K$, each of the three orientations contributes at most $(q/a_K)^3$ plaquettes, so $0\le q^{-3}\operatorname{Rie}_K(f)\le1$:
\[|\nu(x_c)-\nu'(x_c)|\le|e(\nu)-e(\nu')|\le J_K(b)=\mu_K(\{b\}) .\]
Uniformity in $b\in W\subset W'$ follows from (b) at every $\delta$.

\emph{Step 5: clause (e).} Proposition~\YMMref{100.1}{744}, applied to $\nu$ and the tuples $(f)$, $(h)$, bounds the difference between $S^T_2[\nu](f,\tau_vh)$ and $\nu(E_\nu[\mathcal O(f)-\nu(\mathcal O(f))\,|\,\mathcal I]\,E_\nu[\mathcal O(h)-\nu(\mathcal O(h))\,|\,\mathcal I])$ by the right member; by Lemma~\ref{C-365-34}~(iii) the latter is $9a_K^{-8}\operatorname{Rie}_K(f)\operatorname{Rie}_K(h)\operatorname{Var}_\nu(\mathsf e)$, and a random variable with values in an interval of length $J(\beta)/3$ has variance $\le\frac14(J(\beta)/3)^2$, so $9a_K^{-8}\operatorname{Var}_\nu(\mathsf e)\le\frac14(a_K^{-4}J(\beta))^2=\frac14J_K(b)^2$. As $-3\mathsf a(\nu_\pm)=P'(\beta\pm)$ and $P'(\beta+)\ge P'(\beta-)$, $\mathsf a(\nu_+)=\inf_\rho\mathsf a(\rho)$ and $\mathsf a(\nu_-)=\sup_\rho\mathsf a(\rho)$; the mean of $\mathsf e$ in a state $\nu$ is $\nu(a_p)=\mathsf a(\nu)$. So at $\nu_+$ the function $\mathsf e$ is $\ge$ its mean, and at $\nu_-$ it is $\le$ its mean; in both cases the mean is an end of the interval, so $\mathsf e$ is almost surely constant. An ergodic $\nu$ makes the $\mathcal I$-measurable $\mathsf e$ constant. At $\nu=\frac12(\nu_++\nu_-)$, a periodic Gibbs state,
\[\operatorname{Cov}_\nu(F,F')=\tfrac12\operatorname{Cov}_{\nu_+}(F,F')+\tfrac12\operatorname{Cov}_{\nu_-}(F,F')+\tfrac14\bigl(\nu_+(F)-\nu_-(F)\bigr)\bigl(\nu_+(F')-\nu_-(F')\bigr);\]
at $F=\mathcal O(f)$, $F'=\mathcal O(\tau_vf)$ with $\operatorname{Rie}_K(f)>0$ the covariances in $\nu_\pm$ are $O(1/|v|)$ by the plateau at $\nu_\pm$, $\nu_+(F)-\nu_-(F)=-a_K^{-4}J(\beta)\operatorname{Rie}_K(f)$, and $|v|\to\infty$ gives the equality.

\emph{Step 6: clause (f).} (M14) is taken with what the clause supposes, which is what (M14) asks: two frames at one radius with the block sizes $L_1$ and $L_2=L_1+2$, Theorems~\YMMref{AJ}{673} and~\YMMref{AD}{607} granted at both, Theorem~\YMMref{D}{8} at the choices of its constants that they determine and at a lowered choice of each, and, in each frame $i$, the rate of the far part, that is, numbers $\mathsf k_i\ge0$ and $\mathsf c_{{\rm y},i}<\infty$, free of $m$, $K$ and $b$, with $|\mathsf y^{(i)}_{m,K'}(b)-\mathsf y^{(i)}_{m,K}(b)|\le\mathsf c_{{\rm y},i}\,\ell^9K^{\mathsf p_{2,i}}L_i^{-K/10}$ for all integers $m\ge1$, $\ell=2L_i^m$, all $b\in\mathcal W_i$ and all integers $K'\ge K\ge\mathsf k_i+54m$.

The power of (M14) is derived, in Section~\ref{lh:sec} on the latent heat of the $\mathrm{SU}(2)$ Wilson model on $\mathbb Z^3$, from these data, the rate among them; so (M14) gives for every real $A>90$ numbers $\mathsf C_{\rm r}$ and $\beta_1$ with $\beta^4J(\beta)\le\mathsf C_{\rm r}\beta^{-1/(2A)}$ for $\beta\ge\beta_1$. That proof, measured on the physical scale, compares the specific heats of the Wilson model at the block sizes $L_1$ and $L_2=L_1+2$, whose logarithms are rationally independent (Lemma~\ref{C-365-6}), through the exact finite-size bounds, a common near part and the far parts $\mathsf y^{(i)}_{m,K}$, as in Proposition~\ref{C-365-28} made quantitative. The rate enters at the cutoffs $K\ge\mathsf k_i+Am$: there $\ell^9=2^9L_i^{9m}\le2^9L_i^{9K/A}$, so that for $K'\ge K$
\[\sup_{b\in\mathcal W_i}\bigl|\mathsf y^{(i)}_{m,K'}(b)-\mathsf y^{(i)}_{m,K}(b)\bigr|\le2^9\mathsf c_{{\rm y},i}\sup_{k\ge1}k^{\mathsf p_{2,i}}L_i^{-(\frac1{10}-\frac9A)k},\]
a number free of $m$ and $K$, finite as $A>90$. With $A=1/(2\alpha)$ for $0<\alpha<1/180$, and $J\le6$ (the difference quotients of $P$ lie in $[-6,0]$, since $0\le\sum_pa_p\le6$ per site), $\mathsf c_*(\alpha)<\infty$; for $\alpha=0$ use $\alpha=1/182$. On $\mathfrak d(g)$ the label is $b=a_K\beta\ge g^{-2}$ by (M5), so $\beta\ge1$ and
\[a_K^{-4}J(\beta)=b^{-4}\beta^4J(\beta)\le b^{-4}\mathsf c_*(\alpha)\beta^{-\alpha}=\mathsf c_*(\alpha)\,b^{-4-\alpha}L^{-\alpha K}\le g^{8+2\alpha}\mathsf c_*(\alpha)L^{-\alpha K}.\]
By (c) at $b$, $3a_K^{-4}|\rho(a_p)-\nu(a_p)|=|e(\rho)-e(\nu)|\le a_K^{-4}J(\beta)$; for $f$ with support in $[-R,R]^3$ each orientation has at most $(2R/a_K+1)^3$ barycentres there, which gives the bound for $\mathcal O_K(f)$; and (d) gives the bound for the cube means.
\end{proof}

\begin{proof}[Proof of Theorem~\ref{C-365-10}]
\emph{Step 1: clause (a) and the limit at the label.} By the passage, $(K,L^K\mathsf b)\in\mathfrak d(g)$ for $K\ge1$, so $(K,L^K\mathsf b,\nu_K)_{K\ge\max\{K_\circ,1\}}$ is a Gibbs sequence at $g$, with the constant label $\mathsf b$. By (M13)\,(ii) and (v), $\nu_K$ is the one periodic Gibbs state at $(K,L^K\mathsf b)$, the sequence converges, its limit $\sigma$ is its own single component with $\ker H=\mathbb C\Omega$, and the spectrum of $H$ lies in $\{0\}\cup[m_*,\infty)$; $\sigma$ is invariant under $\mathbb R^3\rtimes\mathrm O(3)$, as every limit of a convergent Gibbs sequence (M10). We use two terms. A \emph{limit at the label $\mathsf b$} is a limit, for $n\ge2$ and $X\in\mathfrak S_n(\mathbb R^3)$, of $S^T_{n;K,\ell,L^K\mathsf b}(X^{\langle\ell\rangle})$ as $K$ and $\ell=2L^m$ tend to infinity, together or in either order, the bare coupling being $L^K\mathsf b$ at every cutoff $K$.

A \emph{continuum limit along growing tori at the label $\mathsf b$} is a continuum limit at $g$ along growing tori (Definition~\ref{C-365-5}~(5)) whose defining sequence $(K_j,\ell_j,\beta_j)$ has $\beta_j=L^{K_j}\mathsf b$, that is, all its labels equal to $\mathsf b$, from an index on. Let $X=(f_1,\dots,f_n)\in\mathfrak S_n(\mathbb R^3)$, $n\ge2$ (an entry $0$ makes both sides $0$), and take a side $\ell_1=2L^{m_1}\ge\ell_\circ$ with the supports in $(-\ell_1/4,\ell_1/4)^3$. By (M13)\,(iv), for all $K\ge\max\{K_\circ,1\}$ and all sides $\ell\ge\ell_1$,
\[\bigl|S^T_{n;K,\ell,L^K\mathsf b}(X^{\langle\ell\rangle})-S^T_n[\nu_K](X)\bigr|\le C(X)\,\theta_\circ^{\ell/(2\ell_\circ)},\]
a bound free of $K$, while $S^T_n[\nu_K](X)\to S^\sigma_n(X)$ as $K\to\infty$, the Gibbs sequence converging. So $S^T_{n;K,\ell,L^K\mathsf b}(X^{\langle\ell\rangle})\to S^\sigma_n(X)$ when $K$ and $\ell$ tend to infinity together; in the order $\ell$ first the inner limit is $S^T_n[\nu_K](X)$; in the order $K$ first the inner limit exists at every fixed side $\ell\ge\ell_1$ by (M8), the supports of the periodised functions being disjoint, and it lies within $C(X)\theta_\circ^{\ell/(2\ell_\circ)}$ of $S^\sigma_n(X)$. Along a defining sequence with $\beta_j=L^{K_j}\mathsf b$ from an index on, $K_j\to\infty$ and $\ell_j\to\infty$, the same bound gives $S_n(X)=\lim_jS^T_{n;K_j,\ell_j,\beta_j}(X^{\langle\ell_j\rangle})=S^\sigma_n(X)$ for $n\ge2$; as $\sigma_S$ and $\sigma$ are formed from their connected functions with the one-point functions $0$ ($S_1=0$, and $\sigma((f))=0$ by Definition~\ref{C-365-5}~(2)), $\sigma_S=\sigma$.

So every limit at the label $\mathsf b$, in any order of cutoff and volume, and every continuum limit along growing tori at the label $\mathsf b$ is $\sigma$. The class $[\varnothing]$ of the empty tuple has norm $\sigma(\varnothing)^{1/2}=1$ and is fixed by every $e^{-tH}$, so it spans $\ker H$; take $\Omega=[\varnothing]$. For $X,Y\in\mathfrak S(\mathbb E_+)$, $\langle[X],\Omega\rangle=\sigma(\theta X)=\sigma(X)$ by the invariance, and $P_\Omega$ commutes with $H$, so
\[\sigma(\theta X\cdot\tau_{te_3}Y)-\sigma(X)\sigma(Y)=\bigl\langle(1-P_\Omega)[X],e^{-tH}(1-P_\Omega)[Y]\bigr\rangle,\]
and $\|e^{-tH}(1-P_\Omega)\|\le e^{-m_*t}$ by the spectral theorem; the rotations carry this to every direction. By Theorem~\ref{C-365-8}~(e)\,4 at $g$, clause (c) of that theorem applies to $\sigma$: $\sigma((f))=0$, the kernel $\mathbf s\ge0$ does not increase, and
\[4\pi r\,\mathbf s(r)=\int e^{-mr}\,d\rho(m),\qquad\rho\ge0,\qquad\rho(\{0\})=0 .\]

\emph{Step 2: the gap passes to the kernel.} For real $f\in C^1_c(\mathbb E_+)$ the vector $\psi_f$ is orthogonal to $\Omega$, as $\langle\Omega,\psi_f\rangle=\sigma((f))=0$; so by the gap $\sigma(\theta f,\tau_{te_3}f)=\langle\psi_f,e^{-tH}\psi_f\rangle\le e^{-m_*t}\|\psi_f\|^2$. With $f_\epsilon(x)=\eta_\epsilon(x_3)\mathsf g(x_1,x_2)$ as in the proof of Proposition~\ref{C-365-30}, this reads $c_\epsilon(t)\le e^{-m_*t}c_\epsilon(0)$; letting $\epsilon\to0$ and then $\mathsf g$ tend to a point mass,
\[\mathbf s(2s+t)\le e^{-m_*t}\,\mathbf s(2s)\qquad(s>0,\ t\ge0),\]
that is $\mathbf s(r)\le e^{-m_*(r-r')}\mathbf s(r')$ for $0<r'\le r$.

\emph{Step 3: no spectral mass below the gap.} Fix $r'>0$. For $0\le m<m_*$ and $r\ge r'$, $\int e^{-m''r}d\rho(m'')\ge e^{-mr}\rho([0,m])$, so
\[\rho([0,m])\le e^{mr}\,4\pi r\,\mathbf s(r)\le4\pi r\,e^{-(m_*-m)r}e^{m_*r'}\mathbf s(r')\to0\qquad(r\to\infty),\]
and $\rho([0,m_*))=0$.

\emph{Step 4: the sharper decay.} As $\rho$ is carried by $[m_*,\infty)$, for $0<r'\le r$
\[4\pi r\,\mathbf s(r)=\int_{[m_*,\infty)}e^{-m''r}\,d\rho(m'')\le e^{-m_*(r-r')}\int e^{-m''r'}\,d\rho(m'')=e^{-m_*(r-r')}\,4\pi r'\mathbf s(r');\]
so $\mathbf s(r)=O(r^{-3})$, and $\rho([0,m])=0$ for $m<m_*$ gives $\rho([0,m])=O(m^2)$.

\emph{Step 5: cube means in the limit.} Profiles of distinct cells have supports at distance $\ge T/4$; as $\sigma((f))=0$, $S^\sigma_2=\sigma$ on pairs, and with $0\le f_c,f_{c'}\le1$ and $\int f_c\le q^3$,
\[\bigl|S^\sigma_2(q^{-3}f_c,q^{-3}f_{c'})\bigr|\le q^{-6}\int f_c(x)\int_{|y-x|\ge T/4}\mathbf s(|x-y|)\,dy\,dx\le q^{-3}\,4\pi\int_{T/4}^\infty r^2\mathbf s(r)\,dr,\]
which is finite by Step 4 and the continuity of $\mathbf s$ on $[T/4,\infty)$. The cells $\square_c$ and $\square_{c'}$ are centred at $o+qc$ and $o+qc'$, $c,c'\in\mathbb Z^3$, as in clause (c); let $\mathbf d=c'-c$ be the difference of their indices, as in the proof of Proposition~\ref{C-365-30}~(d). If the two cells have no common point, that is $|\mathbf d|_\infty\ge2$, then the cells are at sup-distance $q(|\mathbf d|_\infty-1)\ge q$, each support lies at sup-distance $\ge T/8$ from the boundary of its cell, so the supports are at distance $\ge q+T/4$, and as $\mathbf s$ does not increase the left member is at most $q^{-6}\|f_c\|_1\|f_{c'}\|_1\mathbf s(q)\le\mathbf s(q)$, which by Step 4 is at most $(r'/q)\,e^{-m_*(q-r')}\mathbf s(r')$ for $0<r'\le q$. For the $26$ cells with $|\mathbf d|_\infty=1$, which share a face, an edge or a corner with the first, the bound is $\chi_\circ q^{-3}$ of the display.
\end{proof}

\begin{lemma}[A derivative from the values on a diameter]\label{C-365-35}
Let $\mathsf u$ be holomorphic on $|z|<r$ with $|\mathsf u|\le M$ there, and $|\mathsf u(x)|\le o$ for real $|x|\le r/2$, where $o>0$. Then $|\mathsf u'(0)|\le(4o/r)[1+\log_2^+(M/o)]$.
\end{lemma}

\begin{proof}[Sketch of proof]
Write $\mathsf u(z)=\sum_mc_mz^m$ with $|c_m|\le Mr^{-m}$. If $o\ge M$, $|\mathsf u'(0)|\le M/r\le o/r$. Otherwise let $N=\lceil\log_2(M/o)\rceil\ge1$ and $Q(x)=\sum_{m\le N}c_mx^m$; for real $|x|\le r/2$, $|\mathsf u(x)-Q(x)|\le M\sum_{m>N}2^{-m}\le o$, so $|Q|\le2o$ there. Bernstein's inequality for a real polynomial $Q_1(t)=\sum_{k\le N}c'_kt^k$ with $|Q_1|\le1$ on $[-1,1]$ gives $|Q_1'(0)|\le N$; applied to $Q_1(t)=\operatorname{Re}[e^{i\varphi}Q(rt/2)]/(2o)$ with the phase $\varphi$ that makes $e^{i\varphi}Q'(0)\ge0$, it gives $|\mathsf u'(0)|=|Q'(0)|\le4oN/r$.
\end{proof}

The comparison of two cutoffs stops before the terminal scale. We
therefore need the vacuum constant of the steps already performed,
with the stopping index retained in the two-torus estimate.

\begin{lemma}[Vacuum constants before the terminal scale]
\label{lem:vacuum-prefix}
Fix the common construction choice and the real regular progression
of Theorem~\YMMref{Y}{438} for one admissible sequence
$x_0,\ldots,x_K$. Let $T_{\ell_1}$ and $T_{\ell_2}$ be admitted
tori for this sequence, where $\ell_i=2L^{m_i}$, and assume
\[
 LM R(x_K)\le L^{m_i},\qquad \ell_i\ge C_LM,\qquad
 \kappa\ge2\log9216+5,\qquad \delta_NM\ge\kappa,
 \qquad E_0\ge2E_N.
\]
Here $C_L$ is the room constant of Proposition~\YMMref{74.5}{617},
and $E_N,\delta_N$ are the normalizer constants of
equations~\YMMref{(1399)--(1401)}{573}.
At step $j+1$ put
\[
 \mathsf{vl}^{(j+1)}(X,z)=E^{(j+1)}(X,\mathbf1,z),
 \qquad \mathsf{vn}^{(j+1)}(X,z)=\mathsf n^{(j+1)}(X,z).
\]
The second family is the local Gaussian constant in
equation~\YMMref{(1399)}{573}. Both families are real and bounded
by $E_0e^{-\kappa d_{j+1}(X)}$. With the extensive term
$\mathsf N^{(j)}$ of equations~\YMMref{(1402), (1405)}{574}, define
\begin{equation}\label{eq:vacuum-prefix-definition}
\begin{split}
 \mathsf E^{(j)}(T_\ell,T_\ell)
   &=\mathsf N^{(j)}(T_\ell,T_\ell)
     +\sum_X\sum_{z\in X}
       \bigl(\mathsf{vl}^{(j+1)}(X,z)
                    +\mathsf{vn}^{(j+1)}(X,z)\bigr),\\
 E_{(n)}[T_\ell]&=\sum_{j=0}^{n-1}
                      \mathsf E^{(j)}(T_\ell,T_\ell),
                    \qquad 0\le n\le K.
\end{split}
\end{equation}
The domains and sites in the first sum have level $j+1$.
These are the constants in equations~\YMMref{(1406)--(1407)}{574}:
on a whole torus there is no boundary layer to remove from the
set of base points.

The normalizing factor of each real step satisfies
\begin{equation}\label{eq:vacuum-step-normalization}
 \log\bigl(\mathsf{pf}_j(x_j)\mathsf Z^{(j)}_{\mathbf1}\bigr)
   =\mathsf N^{(j)}(T_\ell,T_\ell)
            +\sum_X\sum_{z\in X}\mathsf{vn}^{(j+1)}(X,z).
\end{equation}
Consequently the constant accumulated by the zero-source chain
through step $n$ is $E_{(n)}[T_\ell]$.

Put $\ell_* =\min(\ell_1,\ell_2)$ and
$C_{\rm tt}=4096\sum_{q\ge0}e^{-q}$, as in
Proposition~\YMMref{74.5}{617}. Then
\begin{equation}\label{eq:vacuum-prefix-comparison}
 \left|\frac{E_{(n)}[T_{\ell_1}]}{\ell_1^3}
       -\frac{E_{(n)}[T_{\ell_2}]}{\ell_2^3}\right|
 \le2C_{\rm tt}E_0
       \sum_{s=K-n}^{K-1}L^{3s}e^{-\kappa\ell_*L^s/(5M)}.
\end{equation}
The sum is empty when $n=0$. For every integer $u\ge0$,
\begin{equation}\label{eq:vacuum-prefix-tail}
 \sum_{s\ge u}L^{3s}e^{-\alpha L^s}
       \le2L^{3u}e^{-\alpha L^u},
       \qquad \alpha=\frac{\kappa\ell_*}{5M}>4.
\end{equation}
\end{lemma}

\begin{proof}
We first identify the constant of one step. Write
$N_j=\#\mathbb T^{(j)}$, $D=N_j-N_{j+1}$ and
$\nu=\frac32\log(2\pi)$. There are $3D$ positive nonpivot
bonds and $3N_{j+1}$ pivot bonds. At the trivial triple the
fluctuation matrix $\mathsf A_{\mathbf1}$ is positive on a real
space of dimension $9D$, and the elimination map at every pivot
is $L^{-3}\mathrm{id}$ on the three-dimensional colour space.
The real Gaussian integral and these pivot determinants give
\[
 \log\mathsf Z^{(j)}_{\mathbf1}
   =27N_{j+1}\log L+3D\nu
                        -\tfrac12\log\det\mathsf A_{\mathbf1}.
\]
The Haar prefactor of the step, at $t=x_j$, is
\[
 \log\mathsf{pf}_j(t)
 =-D\log\mathsf z(t,\varepsilon(t))
   +3D\bigl(\log\sigma_{\rm H}-\tfrac32\log t\bigr)
   -9N_{j+1}\log L.
\]
Since $D=(L^3-1)N_{j+1}$, the definition of
$\mathsf N^{(j)}$ yields
\[
 \log\bigl(\mathsf{pf}_j(t)\mathsf Z^{(j)}_{\mathbf1}\bigr)
 =\mathsf N^{(j)}(T_\ell,T_\ell)
       +18N_{j+1}\log L-\tfrac12\log\det\mathsf A_{\mathbf1}.
\]
It remains to identify the last two terms with the local Gaussian
constants. The logarithm series defining the local normalizer in
the proof of Proposition~\YMMref{50.2}{386} has
$T=1-c_N^{-1}\mathsf A_{\mathbf1}$, with $\|T\|<1$.
Adding its diagonal terms over the nonpivot bonds gives
\[
 \sum_b f_b(T_\ell)
 =9D\log c_N-\sum_{q\ge1}\frac1q\operatorname{Tr}T^q
 =\log\det\mathsf A_{\mathbf1}.
\]
The series is absolutely convergent. The last equality follows
by applying $-\sum_{q\ge1}(1-v)^q/q=\log v$ to each positive
eigenvalue $v$ of $c_N^{-1}\mathsf A_{\mathbf1}$, which lies
in $(0,2)$. Summing the local finite differences of $f_b$ recovers
$f_b(T_\ell)$, and the tents sum to one for each bond. Thus the
nonpivot constants sum to $-\frac12\log\det\mathsf A_{\mathbf1}$.
Each of the $3N_{j+1}$ pivot bonds contributes $6\log L$, by
equation~\YMMref{(1399)}{573}. This proves
\eqref{eq:vacuum-step-normalization}. The chain identity in
Proposition~\YMMref{114.2}{801} adds to this logarithm the values
$E^{(j+1)}(X,\mathbf1,z)$ of the new local terms. Summing over
the steps proves the assertion about its accumulated constant.

For the comparison, we retain the level in the proof of
Proposition~\YMMref{74.5}{617}. At level $j$, the localization
cubes have side $ML^{j-K}$. A domain is common if it has a lift
contained in an axis-parallel box of side less than $\ell_*/2$.
The regular families have the same value on these lifts by
Theorem~\YMMref{Y}{438}. The Gaussian constants have the same
property: the local logarithm series, pivot bond and tent are
transported by the same chart, as in the proof following
equation~\YMMref{(1400)}{573}. Thus both families define one
function over the common lifts.

The sum of this function over a period of either torus, divided
by its volume, has the same value. Indeed, the larger side is
an integer multiple of the smaller one; a period of the larger
torus contains that multiple cubed of the lifts of each site
of the smaller torus. The extensive term also cancels, since
\[
 \ell_i^{-3}\mathsf N^{(j-1)}(T_{\ell_i},T_{\ell_i})
 =L^{3(K-j)}\bigl[3(L^3-1)
       (\mathsf{co}^{\rm b}_{j-1}+\mathsf{co}^{\rm i}_{j-1})
                         +\mathsf{co}^{\rm s}_{j-1}\bigr]
\]
depends only on the common sequence. Therefore, writing
$\mathsf S_i^{(j),\rm nc}$ for the sum of the two local families
over the non-common domains on $T_{\ell_i}$, we have
\[
 \frac{E_{(n)}[T_{\ell_1}]}{\ell_1^3}
 -\frac{E_{(n)}[T_{\ell_2}]}{\ell_2^3}
 =\sum_{j=1}^n
    \left(\ell_1^{-3}\mathsf S_1^{(j),\rm nc}
                     -\ell_2^{-3}\mathsf S_2^{(j),\rm nc}\right).
\]

The lifting argument in that proof gives
$d_j(X)\ge\ell_*L^{K-j}/(2M)-2$ for each non-common domain.
For a fixed site $z$, its domain-counting estimate is
\[
 \sum_{\substack{X\ni z\,;\ d_j(X)\ge D_0}}e^{-\kappa d_j(X)}
 \le2048\sum_{q\ge D_0}e^{-(\kappa-\log9216)q}
 \le\frac{C_{\rm tt}}2 e^{-\kappa D_0/2}.
\]
The hypotheses give $t_j=\ell_*L^{K-j}/M\ge2L^2\ge338$,
so $\frac12(t_j/2-2)\ge t_j/5$. Multiplying the last bound by
$2E_0$ and by the number $\ell_i^3L^{3(K-j)}$ of level-$j$
sites gives
\[
 \ell_i^{-3}|\mathsf S_i^{(j),\rm nc}|
   \le C_{\rm tt}E_0L^{3(K-j)}
                    e^{-\kappa\ell_*L^{K-j}/(5M)}.
\]
Sum only over $1\le j\le n$ and put $s=K-j$ to obtain
\eqref{eq:vacuum-prefix-comparison}. Finally the ratio of
successive terms in \eqref{eq:vacuum-prefix-tail} is
\[
 \exp\bigl(3\log L-\alpha L^s(L-1)\bigr)\le\tfrac12,
\]
because $L\ge13$, $\alpha\ge4$ and $\log L\le L-1$.
Summing this geometric majorant from $s=u$ proves the tail bound.
\end{proof}


\begin{proof}[Proof of Proposition~\ref{C-365-11}]
\emph{Step 1: the identity.} The translations by $a_K\mathbb Z^3$ and the permutations of the axes act on $T_\lambda$, preserve the weight and the Haar measure and are transitive on the plaquettes, so $\langle a_p\rangle_{K,\lambda,L^Kb}=\mathsf a_\lambda$ is one number; there are $3(\lambda/a_K)^3$ plaquettes, and $p'_{K,\lambda}(b)=-\lambda^{-3}L^K\langle\sum_pa_p\rangle=-3a_K^{-4}\mathsf a_\lambda$. As each plaquette of $\mathbb L_K$ has exactly one image on the torus, $\sum_{p\subset T_\lambda}f^{\langle\lambda\rangle}(x(p))=\sum_{p\subset\mathbb L_K}f(x(p))$, so
\[\langle\mathcal O(f^{\langle\lambda\rangle})\rangle_{K,\lambda,L^Kb}=a_K^{-1}\mathsf a_\lambda\sum_pf(x(p))=-\operatorname{Rie}_K(f)\,p'_{K,\lambda}(b);\]
subtract the two sides; for the tent $\chi_0$ of half-width $1/L$ about a site, $\operatorname{Rie}_K(\chi_0)=L^{-3}$.

\emph{Step 2: the bound.} By the passage both laws are at a first-passage datum, and the transfer (M7) at $n=1$, for $f$ with support in a ball of radius $2$, gives $|A_K(f;b)|\le2\log2\,(2\mathsf Q_1+2)\varrho^{-1}\|f\|_\sharp$; at $f=\chi_0$, $\|\chi_0\|_\sharp\le\mathsf c_\chi$. For a cube mean, $0\le q^{-3}\operatorname{Rie}_K(f)\le1$ as in the proof of Theorem~\ref{C-365-9}~(d).

\emph{Step 3: the read-out at the order $0$ on one torus.} Let $\lambda\in\{\ell,\ell'\}$, $b\in\mathsf J$ and $K\ge K_{\rm bf}(\lambda,\mathsf J)$. The theorem on two consecutive cutoffs (Theorem~\YMMref{AR}{810}) compares on $T_\lambda$, at the label $b$, the run $\mathrm A$ at the cutoff $K^{\mathrm A}=K$ with the run $\mathrm B$ at the cutoff $K^{\mathrm B}=K+1$; the mark $\bullet$ stands for $\mathrm A$ or $\mathrm B$. Its objects are taken at the order $1$, the tuple being any one real Lipschitz function on $T_\lambda$ that is not $0$, and at the source $0$, where no object reads that function (Section~\ref{sec:cutoff-calibration} and Section~\YMMref{116}{807}); its conditions hold at every such $K$ (Corollary~\textup{\ref{IV-62-66}}~(a)).

The objects are: the depth $u=\lceil K/1260\rceil$ and the level $t_1^\bullet=K^\bullet-u$ of $\bullet$ at this depth; the space of the configurations at the depth $u$, one space for both runs, with its Haar measure $dV$ and its inner ball $\mathsf{In}$, an open set that holds the trivial configuration $\mathbf 1$; the chain of $\bullet$, with its exponent $\Phi^\bullet$ on $\mathsf{In}$, the sum $c^\bullet(0)$ of the constants split off along it at the steps $1,\dots,t_1^\bullet$, the integral $I^\bullet=\int_{\mathsf{In}}e^{\Phi^\bullet}dV$ and the mass $T^\bullet=e^{c^\bullet(0)}I^\bullet$; and the classical start of $\bullet$ at the depth $\lceil K/9\rceil$, with its exponent $\hat\Phi^\bullet$ on $\mathsf{In}$, the sum $\hat c^\bullet(0)$ of its constants, $\hat I^\bullet=\int_{\mathsf{In}}e^{\hat\Phi^\bullet}dV$, its measure $\hat\nu^\bullet=e^{\hat c^\bullet(0)}e^{\hat\Phi^\bullet}dV$ on $\mathsf{In}$ and its mass $\hat T^\bullet=\hat\nu^\bullet(\mathsf{In})=e^{\hat c^\bullet(0)}\hat I^\bullet$ (the form of the theorem, the proof of Proposition~\YMMref{116.1}{808}).

The small numbers $\epsilon_{\rm fg}\le\frac18$, $\epsilon_{\rm cl}\le\frac1{20}$, $\epsilon_{\rm sh}\le\frac1{20}$ and $\mathsf R_{t_1}\le\frac12$ are those of Section~\ref{sec:cutoff-calibration} on $T_\lambda$, and $\epsilon_3=4\epsilon_{\rm cl}+6\epsilon_{\rm sh}$. We write $\mathsf R_K[T_\lambda]$ for the number $\mathfrak o_K$ of (M12) on the torus $T_\lambda$ (Definition~\textup{\ref{IV-62-62}}~(4)); with its bound of Corollary~\textup{\ref{IV-62-66}}~(a),
\[\mathsf R_K[T_\lambda]=\mathfrak o_K=8\mathsf R_{t_1}+12\epsilon_{\rm fg}+8\epsilon_{\rm cl}+12\epsilon_{\rm sh}\le\mathsf C_{\rm bf}\lambda^3K^{\mathsf p_0}L^{-K/10}.\]
The links (i)--(iv) below are the four statements of Subsection~\ref{C-365-13} that the proof of clause (c) uses, each derived here from the cited construction estimates. The links (i)--(iii) give
\[\bigl|\log Z_{K+1,\lambda,L^{K+1}b}-\log Z_{K,\lambda,L^Kb}-\operatorname{Re}(c^{\mathrm B}-c^{\mathrm A})(0)\bigr|\le\mathsf R_K[T_\lambda]\le\mathsf C_{\rm bf}\lambda^3K^{\mathsf p_0}L^{-K/10},\]
so that $\mathsf R_K[T_\lambda]$ bounds the error of the read-out at the order $0$ on the torus $T_\lambda$; the link (iv) identifies $c^\bullet(0)$ for Step~4.

(i) \emph{The read-out at the order $0$}, from the discarding estimate, the form and the data of the read-out at $z=0$. With $Z^\bullet=Z_{K^\bullet,\lambda,L^{K^\bullet}b}$, the discarding estimate (Proposition~\textup{\YMMref{115.2}{806}}, at the level $t_1^\bullet$) gives $(1-\mathsf R_{t_1})Z^\bullet\le T^\bullet\le Z^\bullet$. The data of the read-out at $z=0$ (Section~\YMMref{116}{807}; the proofs of Propositions~\YMMref{116.1}{808} and~\YMMref{116.3}{809}) give $|I^\bullet/\hat I^\bullet-1|\le2\epsilon_{\rm fg}$, and $\hat T^{\mathrm A}>0$, $\hat T^{\mathrm B}>0$ with $|\hat T^{\mathrm B}/\hat T^{\mathrm A}-1|\le\epsilon_3\le\frac12$. As $T^\bullet=e^{(c^\bullet-\hat c^\bullet)(0)}\,(I^\bullet/\hat I^\bullet)\,\hat T^\bullet$ with $T^\bullet>0$ and $\hat T^\bullet>0$, moduli and logarithms, with $-\log(1-t)\le2t\log2$ for $0\le t\le\frac12$ and $\log\frac{1+y}{1-y}\le\frac{32}{15}y$ for $0\le y\le\frac14$, give
\[
\begin{multlined}
\Bigl|\log\frac{Z^{\mathrm B}}{Z^{\mathrm A}}-\operatorname{Re}\bigl[(c^{\mathrm B}-\hat c^{\mathrm B})-(c^{\mathrm A}-\hat c^{\mathrm A})\bigr](0)\Bigr|\le1.39\,\mathsf R_{t_1}+4.27\,\epsilon_{\rm fg}+1.39\,\epsilon_3=1.39\,\mathsf R_{t_1}+4.27\,\epsilon_{\rm fg} \\
+5.56\,\epsilon_{\rm cl}+8.34\,\epsilon_{\rm sh}.
\end{multlined}
\]

(ii) \emph{The two classical starts are comparable as measures on the inner ball}, from the proof of the comparison of their masses. Let $Y$ run through the configurations at the depth $\lceil K/9\rceil$, with the Haar measure $dY$, $Y_{t_1}$ being the block average of $Y$ at the depth $u$, and let $\mathsf k^\bullet\ge0$ and $\hat F^\bullet\ge0$ be the weights and the starts of Section~\YMMref{116}{807} at the source $0$, so that $\int_{\mathsf{In}}H\,d\hat\nu^\bullet=\int H(Y_{t_1})\,\mathsf k^\bullet\hat F^\bullet\,dY$ for every bounded Borel function $H$ on the configurations at the depth $u$ (Section~\ref{sec:block-field-laws}). The positive-kernel argument in the proof of Proposition~\YMMref{116.1}{808} runs with a bounded Borel function $H\ge0$ of the block field at the depth $u$ as a further factor, since its kernel inequalities hold against every bounded nonnegative Borel function; it gives
\[
\begin{gathered}
\int H(Y_{t_1})\,\mathsf k^\bullet\,\mathbf 1\{\mathsf k^\circ=0\}\,\hat F^\bullet\,dY\le\epsilon_{\rm sh}\int_{\mathsf{In}}H\,d\hat\nu^\bullet\qquad\\
\bigl(\{\bullet,\circ\}=\{\mathrm A,\mathrm B\}\bigr).
\end{gathered}
\]
Where both weights are positive they are equal and 
\[
\hat F^{\mathrm B}/\hat F^{\mathrm A}\in[e^{-\epsilon_{\rm cl}},e^{\epsilon_{\rm cl}}]
\]
 (Proposition~\textup{\YMMref{116.1}{808}}). Dropping from $\int_{\mathsf{In}}H\,d\hat\nu^{\mathrm B}$ the part where $\mathsf k^{\mathrm A}=0$, comparing the starts where both weights are positive, and using the display for $(\bullet,\circ)=(\mathrm A,\mathrm B)$, and then the same with $\mathrm A$ and $\mathrm B$ exchanged, one obtains for every bounded Borel $H\ge0$
\[
\begin{gathered}
e^{-\epsilon_{\rm meas}}\int_{\mathsf{In}}H\,d\hat\nu^{\mathrm A}\le\int_{\mathsf{In}}H\,d\hat\nu^{\mathrm B}\le e^{\epsilon_{\rm meas}}\int_{\mathsf{In}}H\,d\hat\nu^{\mathrm A},\qquad\\
\epsilon_{\rm meas}=\epsilon_{\rm cl}-\log(1-\epsilon_{\rm sh})\le\epsilon_{\rm cl}+1.06\,\epsilon_{\rm sh}.
\end{gathered}
\]
At the real source $0$ the density $\mathsf D^\bullet=e^{\hat c^\bullet(0)}e^{\hat\Phi^\bullet}$ of $\hat\nu^\bullet$ is positive at every point of $\mathsf{In}$ (the positive classical density in the proof of Proposition~\YMMref{116.1}{808}), so $\log\mathsf D^\bullet=\operatorname{Re}\hat c^\bullet(0)+\operatorname{Re}\hat\Phi^\bullet$; taking for $H$ the indicator functions of the sets where $\mathsf D^{\mathrm B}<e^{-\epsilon_{\rm meas}}\mathsf D^{\mathrm A}$ or $\mathsf D^{\mathrm A}<e^{-\epsilon_{\rm meas}}\mathsf D^{\mathrm B}$, the display gives
\[\bigl|[\operatorname{Re}\hat c^{\mathrm B}(0)+\operatorname{Re}\hat\Phi^{\mathrm B}]-[\operatorname{Re}\hat c^{\mathrm A}(0)+\operatorname{Re}\hat\Phi^{\mathrm A}]\bigr|\le\epsilon_{\rm meas}\qquad\text{Haar-almost everywhere on }\mathsf{In}.\]

(iii) \emph{The constants of the two classical starts cancel}, from the value $0$ and the continuity of the exponent of the classical start at the trivial configuration. At the source $0$ the exponent $\hat\Phi^\bullet(V)$ of the classical start of $\bullet$ at the level $t_1^\bullet$ is the action of its bare array at a background $U(V)$ over $V$ less the Wilson action at that background weighted by the coupling of the level (Proposition~\YMMref{114.2}{801}). It vanishes at $V=\mathbf 1$ and is continuous there. Indeed the trivial configuration of the bare lattice has the action $0$ and the block average $\mathbf 1$, so it lies on the minimal orbit over $\mathbf 1$, on which both terms take one value (the gauge-invariance argument preceding equation~\YMMref{(2133)}{790}, Propositions~\YMMref{114.2}{801} and~\YMMref{113.13}{799}); the weighted action vanishes at it, and so does the action of the bare array, every member of which enters less its value at the trivial pair (equation~\YMMref{(874)}{412}).

The actual analytic-background argument in the proof of Proposition~\YMMref{52.3}{433} supplies a continuous representative $\Upsilon(V)$ near $\mathbf 1$. The weighted action is the finite sum $\eta_t^{-1}\sum_pc(p)a_p(\Upsilon(V))$ of continuous plaquette functions; gauge invariance identifies this sum with the action of the background. Thus it is continuous at $\mathbf 1$. The action of the array is continuous there as well: its members are local terms, holomorphic in the pair that they read (Lemma~\ref{lem:analytic-open-pairs}), the pair being a polynomial map of the configuration (Section~\ref{sec:analytic-core-complements}, \YMMref{(70)}{38}, \YMMref{(69)}{38}); so the action of the array is continuous at $\mathbf 1$ along the continuous choice of Theorem~\YMMref{E}{10} and its real-analytic continuation. So for every $\delta>0$ there is an open set $\mathcal N\subset\mathsf{In}$ with $\mathbf 1\in\mathcal N$ on which $|\operatorname{Re}\hat\Phi^{\mathrm A}|\le\delta$ and $|\operatorname{Re}\hat\Phi^{\mathrm B}|\le\delta$.

The comparison (ii) holds Haar-almost everywhere, so the argument reads the open set $\mathcal N$, which has positive Haar measure: at a point of $\mathcal N$ off the null set of (ii), $|\operatorname{Re}\hat c^{\mathrm B}(0)-\operatorname{Re}\hat c^{\mathrm A}(0)|\le\epsilon_{\rm meas}+2\delta$. As $\delta>0$ is arbitrary, $|\operatorname{Re}(\hat c^{\mathrm B}-\hat c^{\mathrm A})(0)|\le\epsilon_{\rm meas}$, and with (i)
\[\bigl|\log Z^{\mathrm B}-\log Z^{\mathrm A}-\operatorname{Re}(c^{\mathrm B}-c^{\mathrm A})(0)\bigr|\le1.39\,\mathsf R_{t_1}+4.27\,\epsilon_{\rm fg}+6.56\,\epsilon_{\rm cl}+9.40\,\epsilon_{\rm sh}\le\mathsf R_K[T_\lambda],\]
which is the display above.

(iv) \emph{The constants at zero source.} We claim that
\[
 c^\bullet(0)=E_{(n)}[T_\lambda],\qquad
 n=t_1^\bullet=K^\bullet-u,
\]
where the right side is the vacuum constant of the first $n$ steps of the actual progression at cutoff $K^\bullet$. The hypotheses at $K\ge\max\{K_{\rm bf}(\ell,\mathsf J),K_{\rm bf}(\ell',\mathsf J)\}$ give the admitted setting of Section~\ref{sec:cutoff-calibration} on both tori.

The couplings $x_j^\bullet=L^{K^\bullet-j}b$ have zero increments and form an admissible finite sequence. With either side $2L^m$ they give admitted data, by the proof of Lemma~\YMMref{118.1}{812}. The stronger room required in Lemma~\ref{lem:vacuum-prefix} follows from
$LMR(b)\le LMR(L(g^{-2}+\bar B))\le L^{m_{\mathfrak L}(g)}\le L^m$; its proof above also gives $2L^m\ge C_LM$.

At zero source, the family of level $j$ in the chain is the family of the canonical progression of the cut $(x_0^\bullet,\ldots,x_{K^\bullet-1}^\bullet)$, by the same proof of Lemma~\YMMref{118.1}{812}. Theorem~\YMMref{Y}{438} identifies these with the first $K^\bullet$ families of the progression of the full line, with their common companions on all admitted tori. They are real at the trivial configuration and satisfy the bound in Lemma~\ref{lem:vacuum-prefix}; the common choice retains $\delta_NM\ge\kappa$, $E_0\ge2E_N$ and $\kappa\ge2\log9216+5$.

By Proposition~\YMMref{114.2}{801}, step $i$ adds the logarithm of its normalizing factor and the sum of the local values $E^{(i)}(X,\mathbf1,z)$. Equation~\eqref{eq:vacuum-step-normalization} identifies the logarithm with $\mathsf N^{(i-1)}$ and the sum of $\mathsf{vn}^{(i)}$. Since a whole torus has no boundary, every site in each domain is a base point. Thus the constant added at step $i$ is exactly $\mathsf E^{(i-1)}(T_\lambda,T_\lambda)$. Summation over $1\le i\le n$ proves the claimed identity. In particular these accumulated constants are real.

\emph{Step 4: two tori.} Equation~\eqref{eq:vacuum-prefix-comparison} gives, per volume, \[
 \left|E_{(n)}[T_\ell]\ell^{-3}-E_{(n)}[T_{\ell'}]\ell'^{-3}\right|
 \le2C_{\rm tt}E_0\sum_{s=K-n}^{K-1}L^{3s}
       \exp[-\kappa\ell'L^s/(5M)];
\] divide Step 3 by the volumes and subtract the two tori:
\[|U_{K+1}(b)-U_K(b)|\le\ell^{-3}\mathsf R_K[T_\ell]+\ell'^{-3}\mathsf R_K[T_{\ell'}]+2\mathsf v_K\le2\mathsf C_{\rm bf}K^{\mathsf p_0}L^{-K/10}+2\mathsf v_K .\]
The right member is summable in $K$, so $U_K$ converges uniformly on $\mathsf J$: clause (c).

\emph{Step 5: a holomorphic extension free of the cutoff.} For $\lambda\in\{\ell,\ell'\}$, $\phi_\lambda(z)=\widetilde Z_\lambda(z)/\widetilde Z_\lambda(0)$ is holomorphic on $|z|<2\varrho$ with $\phi_\lambda(0)=1$ and $|\phi_\lambda|\le e^{(C_S+C_g)\lambda^3}\le e^{(C_S+C_g)\ell^3}$ by (M1); by Schwarz's lemma $|\phi_\lambda-1|\le\frac12$, so $\log\phi_\lambda$ is holomorphic with modulus $\le\log2$, on $|z|<\varrho/(e^{(C_S+C_g)\ell^3}+1)$. By the shift of the label, $\lambda^{-3}\log\phi_\lambda(x)=p_{K,\lambda}(b-x)-p_{K,\lambda}(b)-\lambda^{-3}\mathcal E_\lambda(x\cdot1)$ for real $x$, the logarithm being real there by the conjugation in (M1); so, with $d$ of Lemma~\ref{C-365-15},
\[\mathsf u_{K,b}(z)=\ell^{-3}\log\phi_\ell(z)-\ell'^{-3}\log\phi_{\ell'}(z)+d(z)\]
is holomorphic on $|z|<r_3$, bounded there by $M_3$, and equals $U_K(b-x)-U_K(b)$ for real $|x|<r_3$; Cauchy's estimate gives $|\partial_b^nU_K(b)|\le n!M_3r_3^{-n}$.

\emph{Step 6: the derivative.} For $b\in\mathsf W$ the function $\mathsf h=\mathsf u_{K+1,b}-\mathsf u_{K,b}$ is holomorphic on $|z|<r_3$ with $|\mathsf h|\le2M_3$, and for real $|x|\le r_3/2$ both $b$ and $b-x$ lie in $\mathsf J$, so by Step 4, with the number $\mathsf o_K=4\mathsf C_{\rm bf}(\mathsf J)K^{\mathsf p_0}L^{-K/10}+4\mathsf v_K$ of (d),
\[|\mathsf h(x)|=\bigl|(U_{K+1}-U_K)(b-x)-(U_{K+1}-U_K)(b)\bigr|\le\mathsf o_K .\]
Lemma~\ref{C-365-35} gives $|U'_{K+1}(b)-U'_K(b)|=|\mathsf h'(0)|\le(4\mathsf o_K/r_3)[1+\log_2^+(2M_3/\mathsf o_K)]$. As $\mathsf o_K$ decays like $K^{\mathsf p_0}L^{-K/10}$ and its logarithm grows like $K$, the series $\sum_KK\sup_{\mathsf W}|U'_{K+1}-U'_K|$ converges; this gives the uniform convergence of $U'_K$ and the bound of $\sup_K\sup_{\mathsf W}|U'_K|$.

\emph{Step 7: the contact terms.} The read-out of the theorem on two consecutive cutoffs at the order $1$ in the source (Theorem~\YMMref{AR}{810}, equation~\YMMref{(2202)}{810}, at $n=1$), on each torus with $f^{\langle\lambda\rangle}$, gives $S^T_{1;K+1}(f)-S^T_{1;K}(f)$ up to the derivative at zero source of the difference of the contact terms at the cutoffs $K+1$ and $K$ and an error that is summable in $K$; subtracting the two tori, by Step 1, $A_{K+1}(f)-A_K(f)=-\operatorname{Rie}_{K+1}(f)(U'_{K+1}-U'_K)(b)-[\operatorname{Rie}_{K+1}(f)-\operatorname{Rie}_K(f)]U'_K(b)$, with $|\operatorname{Rie}_{K+1}(f)-\operatorname{Rie}_K(f)|\le\sqrt3a_K\operatorname{Lip}(f)(2R+2)^3$ for $R$ the radius of a ball about the origin that holds the support of $f$, and $|U'_K|$ bounded; so the difference between the two tori of the derivatives of the contact terms is summable in $K$, uniformly on $\mathsf W$.
\end{proof}

\begin{proof}[Proof of Proposition~\ref{C-365-12}]
\emph{Step 1: characteristic functions from the imaginary axis.} In (a), by the passage $(K,L^Kb)\in\mathfrak d(g)$. For real $y$ and $t$ the function $e^{ityX_\kappa}$ of the cell $\kappa$ has, by the source formula, $\|e^{ityX_\kappa}\|_{{\rm ch},\ell}^N=Z[ity(F_e-F_o)]/Z$, a number $\ge0$, the same for both parities of $\kappa$; with $v=y(F_e-F_o)$, $\|v\|_\sharp=|y|$, and (IM) for $\widetilde Z^{\rm av}$ at $u=0$,
\[\|e^{ityX_\kappa}\|_{{\rm ch},\ell}\le\exp\bigl[N^{-1}\operatorname{Re}\mathcal E^{\rm av}_\ell(itv)+C_I't^2y^2q^3\bigr].\]
The Taylor coefficients of $s\mapsto N^{-1}\mathcal E^{\rm av}_\ell(s(F_e-F_o))$ are real by (A1) of Lemma~\ref{C-365-14}, and bounded uniformly in $\ell$ at fixed $K$, by the count of domains, which bounds the energies of the domains that meet one cell; so $N^{-1}\operatorname{Re}\mathcal E^{\rm av}_\ell(itv)=-\frac12t^2y^2\mathsf e_\ell+O(t^4)$ uniformly in $\ell$, and $\bar{\mathsf e}$ is finite.

\emph{Step 2: the variance from below.} Choose sides $\ell_k\to\infty$ with $\mathsf e_{\ell_k}\to\bar{\mathsf e}$. The domination, for the functions $e^{ity_\kappa X_\kappa}$ of different cells, gives
\[\bigl|\nu\bigl(e^{it\sum_\kappa y_\kappa X_\kappa}\bigr)\bigr|\le\exp\Bigl[-\tfrac12t^2(\bar{\mathsf e}-2C_I'q^3)\sum_\kappa y_\kappa^2+O(t^4)\Bigr],\]
while $|\nu(e^{itY})|\ge\operatorname{Re}\nu(e^{it(Y-\nu(Y))})\ge1-\frac12t^2\operatorname{Var}_\nu(Y)$; letting $t\to0$ gives (a).

\emph{Step 3: covariances by variances.} Apply (a) to $X_\kappa\pm X_{\kappa'}$: 
\[
\operatorname{Var}_\nu(X_\kappa)+\operatorname{Var}_\nu(X_{\kappa'})\pm2\operatorname{Cov}_\nu(X_\kappa,X_{\kappa'})\ge2(\bar{\mathsf e}-2C_I'q^3),
\]
which is the second assertion of (a).

\emph{Step 4: Cauchy's estimate with the chord.} In (b), for $z=x+i\tau$ with $|z|=r\le\varrho/2$, the source formula gives 
\[
N\Lambda(z)=\operatorname{Re}\mathcal E^{\rm av}(w_z)+\log|\widetilde Z^{\rm av}(w_z)|-\log\widetilde Z(0).
\]
(IM) at $u=xF$, $v=\tau(F_e-F_o)$ costs $C_I'\tau^2\ell^3$. By semi-convexity, $\mathsf g(t)=\log\widetilde Z^{\rm av}(tF)$ has $\mathsf g''\ge-2C_I'\ell^3$ on $[-r,r]$, so $\mathsf g+C_I'\ell^3t^2$ lies below its chord:
\[\mathsf g(x)\le\tfrac12[\mathsf g(r)+\mathsf g(-r)]+C_I'\ell^3r^2+\mathsf m_rx,\qquad\mathsf m_r=[\mathsf g(r)-\mathsf g(-r)]/(2r).\]
The cross terms between the cells give $\operatorname{Re}\mathcal E^{\rm av}(w_z)\le N\operatorname{Re}\mathcal E^{\rm av}(zf)+NM_\times T^{-3}\varrho^{-2}q^3|z|^2$, as in Proposition~\ref{C-365-21}. With $\tilde\epsilon(z)=\mathcal E^{\rm av}(zf)+N^{-1}\mathsf m_rz$, holomorphic with $\tilde\epsilon(0)=0$ and $\tilde\epsilon'(0)$ real, $\Lambda(z)-\operatorname{Re}\tilde\epsilon(z)\le\frac12\mathsf S_f(r)+(2C_I'+M_\times T^{-3}\varrho^{-2})q^3r^2$ on $|z|=r$, and Lemma~\ref{C-365-22}~(a) gives the bound, with $m=q^{-3}\tilde\epsilon'(0)$; Lemma~\ref{C-365-22}~(b) gives the covariance. At $r=\hat rq^{-3+p/2}$ one has $(rq^3)^{-2}=\hat r^{-2}q^{-p}$ and $q^3r^2=\hat r^2q^{p-3}\le\hat r^2$.

\emph{Step 5: the cost from below.} By semi-convexity at the bases $\pm tF$, $|t|\le r\le\varrho/2$,
\[\mathsf g(r)+\mathsf g(-r)-2\mathsf g(0)=\int_0^r(r-t)\bigl[\mathsf g''(t)+\mathsf g''(-t)\bigr]dt\ge-2C_I'\ell^3r^2,\]
and $\ell^3/N=q^3$.

\emph{Step 6: the cost from above.} In (c), the walls (Proposition~\ref{C-365-19}) at $x=\pm r$ ($r\le b$, $r\le\varrho/(256\mathsf c)$), the identity in the proof of Proposition~\ref{C-365-21}, the bounds $\omega(r)+\omega(-r)\le2C_\omega r^2$ and $|d(r)+d(-r)|\le4\mathsf A_1r^2/h_1^2$ give
\[\mathsf S_f(r)\le q^3\Delta^2_rD_{K,\ell}(b)+6C_wTq^2r+2q^3r^2\bigl[C_2'+C_\omega+2\mathsf A_1/h_1^2\bigr]\le q^3\Delta^2_rD_{K,\ell}(b)+6C_wTq^2r+2C_Iq^3r^2 .\]
At $r=\hat rq^{-3/2}$: $q^3\Delta^2_rD_{K,\ell}(b)\le q^3[\Delta^2_rD_K(b)+2\bar\eta(\ell)]$ by Lemma~\ref{C-365-16}, $q^3\Delta^2_rD_K(b)\le C_\flat\hat r^2+C_\flat\hat rq^{3/2}\delta\le2C_\flat$ by Theorem~\ref{C-365-9}~(a) with $\delta\le q^{-3/2}$ and $[b\pm r]\subset W'$, and $2q^3\bar\eta(\ell)=6\ell'B'q^3/\ell\le1$; the lower bound is Step 5.

\emph{Step 7: a power of $q$.} In (d) the proof of Lemma~\ref{C-365-32} applies with $C_pq_i^{-p}$ in the place of $C_{\rm LH}(T/q_i)^2$: $C_pq_i^{-p}\ge\sigma'(\theta w_i,w_i)\ge\mathbf s(\sqrt6q_i)(1-3T/(4q_i))^6$, and for $\sqrt6q_i\le r<\sqrt6Lq_i$, $q_i^{-p}<(\sqrt6L)^pr^{-p}$; then $\rho([0,m])\le4\pi e\,m^{-1}\mathbf s(1/m)$ by Proposition~\ref{C-365-30}.

\emph{Step 8: one frame and a growth bound.} In (e), Proposition~\ref{C-365-24}~(b) at the side $\ell_q$, which lies below $6LC_{\rm FS}q^3$, gives
\[\delta\le\frac{3C_wT}q+C_0x_q+x_qY_0\ell_q^{1/3}+\frac{6C_wC_{\rm FS}Tq^2}{\ell_q}\le\frac{3C_wT}q+C_0x_q+\frac{Y_0(6LC_{\rm FS})^{1/3}q}{3C_wTq^2}+\frac{C_wT}q,\]
which is the display; the window inequality gives the three bounds. Only one frame enters.

\emph{Step 9: the variances from above.} In (f), by the passage $(K,L^Kb)\in\mathfrak d(g)$. Let $\nu$ and $0<r\le\varrho/2$ be given, and on $T_{\ell_k}$, with $N_k$ cells, let $\tilde\epsilon_k(z)=\mathcal E^{\rm av}_{\ell_k}(zf)+N_k^{-1}\mathsf m_{r,k}z$, $\mathsf m_{r,k}$ the slope of the chord of Step~4. For $z=x+i\tau$ with $|z|=r$, the source formula, (IM) at $u=xF$, $v=\tau(F_e-F_o)$, the chord of Step~4, and the cross terms between the cells (Lemma~\ref{C-365-14}~(A3) with $R_0=q/2$, $\Delta=T/4$, $\rho_3=3\varrho/2$, so that $C_\times(2q/T)\le c_X$) give
\[\log\|e^{zX}\|_{{\rm ch},\ell_k}-\operatorname{Re}\tilde\epsilon_k(z)\le\tfrac12\mathsf S_{f,\ell_k}(r)+2C_I'q^3r^2+\tfrac49c_X\varrho^{-2}r^2\le\mathsf S_{f,\ell_k}(r)+c_3r^2,\]
the last step by $\mathsf S_{f,\ell_k}(r)\ge-2C_I'q^3r^2$ (Step~5). The functions $\tilde\epsilon_k$ are holomorphic on $|z|<2\varrho$, real on the real axis and $0$ at $0$, and they are bounded uniformly in $k$, by the count of domains for the domains that meet one cell and by (M1) and the floor for the slopes; along a subsequence they converge, locally uniformly, to such a function $\tilde\epsilon$. The domination gives $|\nu(e^{zX})|\le\liminf_k\|e^{zX}\|_{{\rm ch},\ell_k}$, so $G(z)=\nu(e^{zX})e^{-\tilde\epsilon(z)}$ is holomorphic on $|z|<2\varrho$, $G(0)=1$, and $|G|\le\exp[\limsup_k\mathsf S_{f,\ell_k}(r)+c_3r^2]$ on $|z|=r$. The coefficient of $z^2$ in $G$ is $\frac12[\nu((X-\tilde\epsilon'(0))^2)-\tilde\epsilon''(0)]$, and Cauchy's estimate gives
\[\operatorname{Var}_\nu(X)\le\nu\bigl((X-\tilde\epsilon'(0))^2\bigr)\le\tilde\epsilon''(0)+2r^{-2}\exp\Bigl[\limsup_k\mathsf S_{f,\ell_k}(r)+c_3r^2\Bigr].\]
Finally $\tilde\epsilon''(0)$ is the limit of $\partial_s^2\mathcal E^{\rm av}_{\ell_k}(sf)|_{s=0}$ along the subsequence. By (A1) of Lemma~\ref{C-365-14} and the cross terms at $\rho_3=\varrho$, $N^{-1}\mathcal E^{\rm av}_\ell(s(F_e-F_o))=\frac12[\mathcal E^{\rm av}_\ell(sf)+\mathcal E^{\rm av}_\ell(-sf)]+N^{-1}\mathrm{cr}(s)$ with $|\mathrm{cr}(s)|\le Nc_X|s|^2/\varrho^2$ for complex $|s|<\varrho$, so Cauchy's estimate gives $|\mathrm{cr}''(0)|\le2Nc_X/\varrho^2$, hence $\partial_s^2\mathcal E^{\rm av}_\ell(sf)|_{s=0}\le\mathsf e_\ell+2c_X\varrho^{-2}$ and $\tilde\epsilon''(0)\le\bar{\mathsf e}+2c_X\varrho^{-2}$. The second assertion of (f) is (a): $|\operatorname{Cov}_\nu(X_\kappa,X_{\kappa'})|\le C_{\rm var}q^{6-p}+2C_I'q^3\le(C_{\rm var}+2C_I')q^{6-p}$, as $p\le3$; divide by $q^6$.

\emph{Step 10: the walls and the cost.} In (g), for $t$ between $0$ and $x$ the source $xF+t(1-F)=t+(x-t)F$ has the norm $\le|x|\le\varrho/2$, and $\|1-F\|_\sharp\le1$; so semi-convexity gives $\partial_t^2\log\widetilde Z^{\rm av}(xF+t(1-F))\ge-2C_I'\ell^3$, and at $t=x$
\[\log\widetilde Z^{\rm av}(x)\ge\log\widetilde Z^{\rm av}(xF)+xN\upsilon_f(x)-C_I'\ell^3x^2,\]
which divided by $N$ is the first assertion. Next, by (A1) of Lemma~\ref{C-365-14} at constant sources,
\[\mathsf S_f(r)=\Upsilon_f(r)+\Upsilon_f(-r)+N^{-1}\bigl[\log\widetilde Z(r)+\log\widetilde Z(-r)-2\log\widetilde Z(0)\bigr];\]
the first two terms are $\le2C_\Upsilon$; by the identity in the proof of Proposition~\ref{C-365-21} and the bounds of $\omega$ and of $d$ there, the bracket is $q^3[\Delta^2_rD_{K,\ell}(b)+\omega(r)+\omega(-r)-d(r)-d(-r)]\le q^3[\Delta^2_rD_{K,\ell}(b)+2(C_\omega+2\mathsf A_1/h_1^2)r^2]$, as $r\le r_1\le\min\{r_2/2,h_1/2\}$; and $\Delta^2_rD_{K,\ell}(b)\le\Delta^2_rD_K(b)+2\bar\eta(\ell)$ by Lemma~\ref{C-365-16}, while $\Delta^2_rD_K(b)\le C_\flat r^2+C_\flat\delta r$ by Theorem~\ref{C-365-9}~(a), since $[b-r,b+r]\subset W'$.
\end{proof}

\section{Local nets and duality}\label{C-363-1}

We construct the local net of the action density from the continuum limits and their components. The lattice quarter turn determines the modular data of a wedge. Passing its bounded operators to the limit gives wedge duality and the affiliation of the closed field operators with their local algebras.

We distinguish the original net from its dual completion: the latter is Haag dual, while equality with the original net is an additional assertion. Energy bounds, Gaussian examples and the lattice ring calculation clarify the scope of this distinction. The fixed-depth block-loop arguments use Section~\ref{C-372-1}.

\subsection{Notation and hypotheses}

\paragraph{Constants.} We fix the radius $\varrho>0$ and the constants $(\sigma^-,\gamma)$ of the list $\mathcal T^{\rm S}[2\varrho]$ of the common sourced construction in Section~\YMMref{82}{665}, among them the integers $L,M,\mathfrak p_0$ and the number $B_\sharp\ge0$ that bounds the increments of the flow of the running coupling (see the objects below), and a coupling $0<g\le\gamma$; $\bar B=B_\sharp L/(L-1)$ is formed from $\sigma^-$ alone, with no radius, $\gamma$ or $g$, and $\bar B\le\frac{13}{24}\gamma^{-2}<\gamma^{-2}$ (indeed, $B_\sharp\gamma^2\le\frac12$ and $L\ge13$ give $\bar B\gamma^2\le L/[2(L-1)]\le13/24$). $C^g_{\rm pt}$ is the constant of the point bound at $g$ of equation~\YMMref{(2097)}{779} among the construction results, and $\bar{\mathfrak c}(\mathsf C^g)$, with the bound system $\mathsf C^g$ at $g$ of equation~\YMMref{(2030)}{765}, is the bound of the one-point constants of the components of equation~\YMMref{(1993)}{749}. The short-distance law of the two-point function among the construction results, the law of two points of Theorem~\textup{\YMMref{AV}{851}} as it holds in every component (Corollary~\textup{\YMMref{127.11}{854}}), holds at the couplings $0<g\le\gamma_{\rm sd}$, where $\gamma_{\rm sd}=\gamma_W\in(0,\gamma]$ is its ceiling equation~\YMMref{(2346)}{850}, a number formed from $(\sigma^-,\gamma)$ alone. For $0<g\le\gamma_{\rm sd}$ we put
$$c_{\rm sd}=c_{\rm sd}(g)=4.6\cdot10^{-9}L^{-2}\bigl(g^{-2}+\bar B\bigr)^{-2},\qquad r_{\rm sd}=r_{\rm sd}(g)=2.2\,\mathsf d_0\bigl(g,\tfrac1{10}\bigr),$$
with $\bar B=B_\sharp L/(L-1)$ and the threshold $\mathsf d_0$ of the law, equation~\YMMref{(2350)}{851}: two numbers that depend on $g$ and on no sequence, limit, decomposition or $\xi$, and $c_{\rm sd}(g)\le4.6\cdot10^{-9}g^4$. Then for every component $\sigma'_\xi$ of a limit of periodic Gibbs states at $g$ the kernel $\mathsf s$ of its two-point function, $\sigma'_\xi((f,h))=\iint f(x)h(y)\,\mathsf s(|x-y|)\,dx\,dy$ for $f,h$ of disjoint supports, is $\ge0$ and non-increasing, and $\mathsf s(r)\ge c_{\rm sd}r^{-6}$ for $0<r\le r_{\rm sd}$. Indeed, with the point vectors $\Psi_1$ of the reconstruction recalled below and $0<\epsilon<r/2$, $\mathsf s(r)=\|e^{-(r/2-\epsilon)H}\Psi_1(\epsilon e_3)\|^2$, and the truncated kernel, $\mathsf s$ minus the square of the one-point constant, is $\|e^{-(r/2-\epsilon)H}(1-|\Omega\rangle\langle\Omega|)\Psi_1(\epsilon e_3)\|^2$, the vacuum of a component being unique; both are $\ge0$ and non-increasing in $r$.

For $0<d\le\mathsf d_0(g,\frac1{10})$ let $f\ge0$ be of class $C^1$, equal to $1$ on the ball of radius $0.3d$ about $2de_3$ and to $0$ off the ball of radius $0.9d$ about it, with $\sup f=1$ and ${\rm Lip}\,f<3/d$: then $(\theta f,f)$ is a test pair at the scale $d$ with the shape parameter $\frac1{10}$ (Lemma~\ref{lem:ward-test-pairs}), its two supports lie at distances between $2.2d$ and $5.8d$, and $\int f<3.06\,d^3$. The lower bound of Corollary~\textup{\YMMref{127.11}{854}} makes the truncated two-point function at this pair at least $\frac23\,10^{-8}\mathsf c_2\mathsf b^{-2}$, with $\mathsf c_2=9/16\pi^2$ \YMMref{C}{7} and the label $\mathsf b\le L(g^{-2}+\bar B)$ (equation~\YMMref{(1854)}{710}); as the truncated kernel is non-increasing, its value at $2.2d$ is at least $\frac23\,10^{-8}\mathsf c_2\mathsf b^{-2}(3.06\,d^3)^{-2}$, and $\frac23\,10^{-8}\,\mathsf c_2\,(2.2)^6(3.06)^{-2}>4.6\cdot10^{-9}$ gives the bound at $r=2.2d$, for the truncated kernel and so for $\mathsf s$.

\paragraph{Space-time.} In $\mathbb R^3\ni x=(x_1,x_2,x_3)=(\vec x,x_3)$ the coordinate $x_3$ is Euclidean time and $\mathbb E_+=\{x_3>0\}$; $\theta,\theta_1$ are the reflections at $x_3=0$ and at $x_1=0$, $\tau_v$ is the translation by $v$, $\mathsf d(x_1,x_2,x_3)=(x_3,x_2,x_1)$, and $R_\tau x=(x_1\cos\tau-x_3\sin\tau,\,x_2,\,x_1\sin\tau+x_3\cos\tau)$, also on $\mathbb C^3$, with $R=R_{\pi/2}=\theta_1\mathsf d$; $\mathsf I=\{x_1>0,x_3>0\}$ and $\mathsf{II}=R\,\mathsf I=\{x_1<0,x_3>0\}$. An isometry $\mathsf r$ acts on functions by $\mathsf rh=h\circ\mathsf r^{-1}$ and on tuples entry by entry. In $\mathbb R^{1+2}\ni x=(x^0,x^1,x^2)$, $t=x^0$: $V_+=\{x^0>|\vec x|\}$, $W=\{x^1>|x^0|\}$, $W'=-W$, $\mathsf j(x)=(-x^0,-x^1,x^2)$, and $B(s)(t,x^1,x^2)=(t\cosh s+x^1\sinh s,\,t\sinh s+x^1\cosh s,\,x^2)$ are the boosts that keep $W$. A wedge is a region $\mathfrak gW$, $\mathfrak g$ in the Poincar\'e group; a double cone is a Poincar\'e image of $\{|x^0|+|\vec x|<r\}$, upright if it is a translate of one; $O^\times$ is the union of the open sets that are spacelike separated from $O$; $\mathcal F(O)=C^\infty_c(O;\mathbb C)$, $\mathcal F_{\mathbb R}(O)$ its real elements, $\mathcal S=\mathcal S(\mathbb R^{1+2})$, $\hat f$ the Fourier transform, and $f^{(a,\Lambda)}=f\circ(a,\Lambda)^{-1}$.

\paragraph{Data (Rec) and the reconstruction.} A separated tuple is a tuple of real $C^1_c$ functions on $\mathbb R^3$ with pairwise disjoint supports; $F,F'$ denote finite complex combinations of separated tuples, over a set if all supports lie in it, and $\bar F$ has the conjugate coefficients. Data (Rec) with a triple $(C_\wp,\wp,c_0)$ are the symmetric functionals $\sigma$ on separated tuples, $\sigma(\varnothing)=1$, with the following conditions (Section~\YMMref{107}{765}): (Rec1) $\sigma$ is invariant under the isometries of $\mathbb R^3$; (Rec2) $\sigma$ is reflection positive at $x_3=0$, the classes $[F]$ of the combinations over $\mathbb E_+$ span the Hilbert space $\mathcal H$, $\Omega=[\varnothing]$, and $e^{-tH}[F]=[\tau_{te_3}F]$ with $H\ge0$; (Rec3) the joint spectrum of $H$ and the momenta lies in $\{E\ge|q|\}$; (Rec4) $|\sigma(\theta\bar F\cdot(h)\cdot F')|\le C_\wp\eta^{-\wp}\|h\|_{L^1}\|[\tau_{-\eta e_3}F]\|\|[\tau_{-\eta e_3}F']\|$ for $0<\eta\le1$, $F,F'$ over $\{x_3>\eta\}$ and $h\in C^1_c(\{|x_3|<c_0\eta\})$. (Rec5) is the uniqueness of the vacuum, $\ker H=\mathbb C\Omega$.

The reconstruction in Sections~\YMMref{109}{772} and~\YMMref{110}{776} gives $N=1+H$, $\mathcal D_\infty=\bigcap_kD(H^k)$, a representation $U$ of the Poincar\'e group, and the field $\phi(f)$, $f\in\mathcal S$, on $\mathcal D_\infty$, with $\phi(f)^\dagger=\phi(\bar f)$ and $[H,\phi(f)]=i\phi(\partial_tf)$ on $\mathcal D_\infty$, and with Lorentz covariance and locality on the span $\mathcal D_0$ of the vectors $\phi(f_1)\cdots\phi(f_n)\Omega$ ($\mathcal D_0(O)$: all $f_k\in\mathcal F(O)$); for $Y=\phi(f_1)\cdots\phi(f_n)$ we write $Y^\dagger=\phi(\bar f_n)\cdots\phi(\bar f_1)$. It gives the field $\mathsf B(s,s')$ at the point $0$ between the Euclidean times $s,s'>0$, with $e^{-tH}\mathsf B(s,s')e^{-t'H}=\mathsf B(s+t,s'+t')$, $\mathsf B(s,s')^*=\mathsf B(s',s)$ and
\begin{equation}\label{C-363-2}
\|\mathsf B(s,s')\|\le C_\wp\min\{s,s',1\}^{-\wp},\qquad e^{-sH}\phi(f)e^{-sH}=\int f(x)\,U(x)\mathsf B(s,s)U(x)^{-1}\,dx,
\end{equation}
the insertion operators of Lemma~\YMMref{107.1}{766}: for $0<\eta\le1$ and $h\in L^1(\{|x_3|<c_0\eta\})$, $\mathsf B^{\rm in}_\eta(h)$ is the bounded operator, linear in $h$ and of norm $\le C_\wp\eta^{-\wp}\|h\|_{L^1}$, with $\langle[\tau_{-\eta e_3}F],\mathsf B^{\rm in}_\eta(h)[\tau_{-\eta e_3}F']\rangle=\sigma(\theta\bar F\cdot(h)\cdot F')$ for $F,F'$ over $\{x_3>\eta\}$ when $h\in C^1_c(\{|x_3|<c_0\eta\})$, and by the insertion formula \YMMref{(2039)}{767} $\langle\psi,\mathsf B^{\rm in}_\eta(h)\psi'\rangle=\int h(x)\langle\psi,U(\vec x)\mathsf B(\eta+x_3,\eta-x_3)U(\vec x)^{-1}\psi'\rangle\,dx$; the form $q_0$ of the field at the point $0$ and the time $0$, $q_0(e^{-sH}\psi',e^{-s'H}\psi)=\langle\psi',\mathsf B(s,s')\psi\rangle$, and point vectors $\Psi_n(z)$, holomorphic on the tube $\mathsf{Tb}^{(n)}=\{z\in(\mathbb C^3)^n:d(z_1)>0,\ d(z_{k+1}-z_k)>0\}$, $d(z)={\rm Re}\,z_3-|{\rm Im}(z_1,z_2)|$, of norm $\le(C_\wp\eta^{-\wp})^n$ where these numbers are $\ge2\eta$, $\eta\le1$, with $[G]=\int\prod_kh_k(x_k)\Psi_n(x)\,dx$ for ordered tuples $G=(h_k)$ in $\mathbb E_+$ whose supports lie in balls of radii below $c_0$ times the gaps of the heights.

With the continued translations $\mathcal T(\mathfrak w)$, $\mathcal T(\eta e_3)=e^{-\eta H}$, and $\iota(t,\vec x)=(\vec x,-it)$, Lemma~\YMMref{109.2}{773}, which links the field operators with the point vectors, reads
$$\text{(Br)}\qquad\mathcal T(\mathfrak w_1)\phi(f_1)\cdots\mathcal T(\mathfrak w_n)\phi(f_n)\Omega=\int\prod_kf_k(x_k)\,\Psi_n\bigl((\iota x_k+\textstyle\sum_{i\le k}\mathfrak w_i)_k\bigr)\,dx\qquad(d(\mathfrak w_i)>0),$$
and the products of fields are the limits of these vectors as the $\mathfrak w_i$ tend to $0$.

\paragraph{The net and the conditions.} $P_n=\mathbf 1_{[n,n+1)}(H)$, $b_{mn}(f)=\|P_m\phi(f)P_n\|$, $\bar\phi(f)$ is the closure of $\phi(f)$, $\mathcal P$ is the joint spectral measure of the translations, $U(B(s))=e^{isK}$ and $\Delta^s=e^{-2\pi sK}$. For open $O$
$$\mathcal C(O)=\{Q\in\mathcal B(\mathcal H):\langle\phi(\bar f)\psi,Q\psi'\rangle=\langle\psi,Q\phi(f)\psi'\rangle\ \ \forall f\in\mathcal F(O),\ \psi,\psi'\in\mathcal D_0\},\qquad\mathcal A(O)=\mathcal C(O)';$$
$\mathcal L(O)=\mathcal A(O)\cap\mathcal C(O^\times)$; $\mathcal H_0$ is the closure of $\mathcal C(W)\Omega$; for bounded open convex $O$, $\mathcal A^d(O)=\bigcap\{\mathcal A(\mathfrak gW):\mathfrak gW\supset O\}$, and for other open $G$, $\mathcal A^d(G)=\bigvee\{\mathcal A^d(O):O\subset G\ \text{a double cone}\}$. $\mathfrak P$ is the unital $*$-algebra generated linearly by $\mathcal F(\mathbb R^{1+2})$, with $f^*=\bar f$ and $gh=hg$ for real spacelike separated $g,h$; $\pi$ is its representation by the fields on $\mathcal D_0$ and $\mathcal W(f_1\cdots f_n)=\langle\Omega,\phi(f_1)\cdots\phi(f_n)\Omega\rangle$ its Wightman functional; for a subalgebra $\mathfrak A$, $\mathfrak A^+$ is the set of the finite sums $\sum Z_i^*Z_i$, $Z_i\in\mathfrak A$; $\mathfrak P(W)$ is generated by $1$ and $\mathcal F(W)$.
\begin{itemize}
\item[(W)] $\mathcal C(O_1)'\subset\mathcal C(O_2)$ for all non-empty open spacelike separated $O_1,O_2$.
\item[(CP)] For every upright double cone $K$, real $f\in\mathcal F(K)$ and $Y_0,\dots,Y_r$ in the algebra $\mathfrak A$ generated by $1$, $f$ and $\mathcal F(K^\times)$: if $\sum_n\lambda^nY_n\in\mathfrak A^+$ for all real $\lambda$, then $\mathcal W(\sum_nf^nY_n)\ge0$.
\item[(QCP$_f$)] For real $f\in\mathcal F(K)$, $K$ an upright double cone with $\bar K\subset W'$, and $\eta=\phi(f)\Omega$: if $Y_0,Y_1,Y_2\in\mathfrak P(W)$ and $Y_0+\lambda Y_1+\lambda^2Y_2\in\mathfrak P(W)^+$ for all real $\lambda$, then $\langle\Omega,Y_0\Omega\rangle+\langle\Omega,Y_1\eta\rangle+\langle\eta,Y_2\eta\rangle\ge0$. This is the second order of (CP) at $f$; its first order is the same statement with $Y_2=0$.
\item[(LB$_{u,c}$)] $|\langle\psi,\phi(u)\psi\rangle|\le c\langle\psi,N\psi\rangle$ on $\mathcal D_0$; (LB$^\pm_{u,c}$) are its halves $\pm\langle\psi,\phi(u)\psi\rangle\le c\langle\psi,N\psi\rangle$; (LB$_u$) holds if (LB$_{u,c}$) holds for some $c$, and $c_*(u)\in[0,\infty]$ is the least such $c$.
\item[(QA$_f$)] $\sum_n\|\phi(f)^n\Omega\|^{-1/n}=\infty$.
\end{itemize}

\paragraph{The objects.} For a state $\nu$ on $L^{-K}\mathbb Z^3$ or on a torus and a tuple $X=(f_i)$ let $(X)^\nu_K=\prod_i(\mathcal O_K(f_i)-\nu(\mathcal O_K(f_i)))$, extended linearly to combinations; $\theta$ and the bar act on lattice functions by reflection and conjugation. A pair $(K,\beta)$ is admitted at $g$ if it lies in the canonical family $\mathfrak d(g)$ of equation~\YMMref{(1843)}{708} (Section~\ref{sec:state-limit-conventions}), that is, if it is a first-passage datum at $g$ (Corollary~\YMMref{74.2}{616}), formed at the fixed $(\sigma^-,\gamma)$: along the flow $x_0=\beta$, $x_k=L^{-1}x_{k-1}-b_k$ of the inverse squared running coupling, whose increments obey $|b_k|\le B_\sharp$, $K$ is the least index with $x_K\le X_g=L(g^{-2}+B_\sharp)$.

Then $x_0>x_1>\dots>x_K\ge\gamma^{-2}$, and the label $L^{-K}\beta$ lies in the outer interval of the labels $I^{\rm lab}(g)=[g^{-2}+B_\sharp-L^{-1}\bar B,\,L(g^{-2}+\bar B)]$ of (equation~\YMMref{(1854)}{710}) (Corollary~\YMMref{74.2}{616}). A periodic Gibbs state at $(K,\beta)$ is a Gibbs state of the Wilson weight $e^{-\beta\sum_pa_p}$ on $L^{-K}\mathbb Z^3$ that is invariant under the translations of a subgroup of finite index of $L^{-K}\mathbb Z^3$; a torus state at $(K,\beta)$ is a weak limit point, as $m\to\infty$, of the Wilson measures of the tori $L^{-K}\mathbb Z^3/2L^m\mathbb Z^3$ (Section~\YMMref{91}{697}), and it is a periodic Gibbs state (Proposition~\YMMref{91.2}{699}). A volume-first limit at $g$ (Section~\ref{sec:state-limit-conventions}) has a defining sequence $(K_j,\beta_j,\nu_j)$: $K_j\to\infty$, pairs $(K_j,\beta_j)$ admitted at $g$, and torus states $\nu_j$ at $(K_j,\beta_j)$ whose connected functions converge; its centred functional is $\sigma(X)=\lim_j\nu_j((X)_j)$, $(X)_j=(X)^{\nu_j}_{K_j}$.

A limit of periodic Gibbs states at $g$ is $\sigma'(X)=\lim_j\rho_j((X)^{\rho_j}_{K_j})$ for periodic Gibbs states $\rho_j$ of the Wilson weight at admitted pairs $(K_j,\beta_j)$, $K_j\to\infty$. In general $(\cdot)_j$ is the product at $K_j$ centred by the $j$-th state of the sequence at hand: $(X)_j=(X)^{\nu_j}_{K_j}$ for the centred functional, $(X)_j=(X)^{\rho_j}_{K_j}$ for a limit of periodic Gibbs states $\rho_j$, also where the periodic Gibbs states $\rho_{j,\xi}$ that approximate a component integrate, the $\rho_j$ giving the centring, and, in Steps~2--7 of the proof of Theorem~\ref{C-363-4}, the product on the torus $\Lambda_j$ centred by the Wilson measure $\mu_j$ there. A canonical decomposition $(\mathsf X,\lambda,(\sigma_\xi))$ of $\sigma$ has $\mathsf X$ compact metrisable, $\lambda$ a probability measure of full support, $\xi\mapsto\sigma_\xi(X)$ continuous, $\xi\mapsto\sigma_\xi$ injective, and iterated cluster limits $\lim_{|v_q|\to\infty}\cdots\lim_{|v_1|\to\infty}\sigma(X^0\cdot\tau_{v_1}X^1\cdots\tau_{v_q}X^q)=\int\prod_i\sigma_\xi(X^i)\,\lambda(d\xi)$; the $\sigma_\xi$ are its components, and objects of $\sigma_\xi$ bear the index $\xi$.

\paragraph{Lattice, norms, slabs.} On $\Lambda=a\mathbb Z^3/\ell\mathbb Z^3$, $\ell/a$ even and $\ge4$: $G=\mathrm{SU}(2)$ with Haar measure, $t_p=\frac12{\rm tr}\,\partial U(p)$, $a_p=1-t_p$, the Wilson measure $\mu_\beta\propto e^{-\beta\sum_pa_p}dU$, and $\mathcal O(f)=a^{-1}\sum_pf(x(p))a_p$, $x(p)$ the barycentre of $p$ ($\mathcal O_K$ at $a=L^{-K}$). $\|h\|_\sharp=\max\{\sup|h|,40M\,{\rm Lip}\,h\}$ and $\|h\|_{\sharp,\mathsf l}=\max\{\sup|h|,40M\mathsf l\,{\rm Lip}\,h\}$; on $\mathbb R^2$, $|h|_{C^1}=\sup|h|+\sup|\nabla h|$; $\|f\|_\star=\int dt\sum_{\kappa\in\mathbb Z^2}\sup_{\vec x\in\kappa+[-1,1]^2}(|f|+|\nabla_{\vec x}f|)(t,\vec x)$. A board $(i,o)$ of $L^{-K}\mathbb Z^3$, $i\le K-1$ and $o$ a site of $L^{-K}\mathbb Z^3$, is its partition into the cells $\square_{\mathbf k}=o+\mathsf l\mathbf k+[-\frac{\mathsf l}2,\frac{\mathsf l}2]^3$, $\mathbf k\in\mathbb Z^3$, cubes of side $\mathsf l=L^{-i}$ centred at the sites $o+\mathsf l\mathbf k$; the middle of $\square_{\mathbf k}$ is $\{|x-o-\mathsf l\mathbf k|_\infty\le\frac{\mathsf l}4\}$; for odd $L$ the walls of the cells lie midway between sites and the middle planes of the cells pass through sites; and $\|\cdot\|_{\rm ch}$ is the seminorm of the chessboard estimate.

For $h\in L^1(\mathbb R^2)$, $\mathsf B_h(s,s')=\int h(\vec x)U(\vec x)\mathsf B(s,s')U(\vec x)^{-1}d^2x$ is the field at time $0$ smeared in space; $T_{mn}(h)=P_me^{sH}\mathsf B_h(s,s')e^{s'H}P_n$ does not depend on $s,s'$, and $\mathsf T(h)=\sum_{m,n}T_{mn}(h)$ is the form of the field at time $0$ on the vectors of bounded energy. With $\rho_{\mathsf l}(\varsigma)=\frac5{\mathsf l}\cos^2\frac{5\pi\varsigma}{2\mathsf l}$ on $|\varsigma|\le\frac{\mathsf l}5$ and $0$ elsewhere, the plate is $\mathsf I_{\mathsf l}(h)=\int\rho_{\mathsf l}(\varsigma)\mathsf B_h(\frac{\mathsf l}2+\varsigma,\frac{\mathsf l}2-\varsigma)\,d\varsigma$.

\subsection{Principal results}

\begin{theorem}[The net of data (Rec)]\label{C-363-3}
Let $\sigma$ be data (Rec) with the triple $(C_\wp,\wp,c_0)$; uniqueness of the vacuum is not assumed. Every statement holds also with $W$ and $W'$ exchanged.
\begin{enumerate}
\item[(a)] \emph{Frame.} For every $n$ the spaces $N^n\mathcal D_0$ and $N^n\mathcal D_0(O)$, $O\ne\emptyset$ open, are dense, and $\phi(f)$ on $\mathcal D_0$ and on $\mathcal D_\infty$ have one closure $\bar\phi(f)$. $\mathcal C(O)=\{Q:Q\bar\phi(f)\subset\bar\phi(f)Q,\ f\in\mathcal F(O)\}$ is a von Neumann algebra, $U(\mathfrak g)\mathcal C(O)U(\mathfrak g)^{-1}=\mathcal C(\mathfrak gO)$, and $\mathcal A$ is isotonous and covariant. For $O\ne\emptyset$, $\Omega$ is separating for $\mathcal C(O)$ and cyclic for $\mathcal A(O)$; $\bar\phi(f)$, $f\in\mathcal F(O)$, is affiliated with $\mathcal A(O)$, and $\mathcal A$ is the least net with this property; $\mathcal C(\bigcup_iO_i)=\bigcap_i\mathcal C(O_i)$ and $\mathcal A(\bigcup_iO_i)=\bigvee_i\mathcal A(O_i)$. $\mathcal L$ is isotonous, covariant and local; (W) holds if and only if $\mathcal L=\mathcal A$, and if and only if it holds for spacelike separated upright double cones.
\item[(b)] \emph{The order $\wp$.} For $f\in\mathcal S$ and $m,n\in\mathbb N_0$:
(b1) $b_{mn}(f)\le eC_\wp\|f\|_{L^1}(m+n+2)^\wp$;
(b2) $(|m-n|-1)^kb_{mn}(f)\le b_{mn}(\partial_t^kf)$ if $|m-n|\ge2$, $k\in\mathbb N$;
(b3) for $\epsilon\in(0,1]$, $k_1=\max\{0,\lceil(\wp-1)/\epsilon\rceil\}$ and $|m-n|>(m+n+2)^\epsilon$: $b_{mn}(f)\le3^{k_1}eC_\wp\|\partial_t^{k_1}f\|_{L^1}(m+n+2)$, so that the blocks are linearly bounded off every band $|m-n|\le(m+n+2)^\epsilon$, $\epsilon>0$;
(b4) $b_{mn}(h*f)\le\|h\|_{L^1}b_{mn}(f)$ for $h\in L^1(\mathbb R^{1+2})$ with $h*f\in\mathcal S$;
(b5) with $\mathsf s_0(q)=1+2\cdot2^q+2\cdot3^q$ and $\mathsf s(q,k)=\sum_{l\ge1}(l+3)^ql^{-k}$, for real $j\ge0$ and $\psi,\psi'\in\mathcal D_\infty$:
$$\|N^j\phi(f)\psi\|\le2^je^2C_\wp\bigl[\mathsf s_0(\wp+j)\|f\|_{L^1}+2\mathsf s(\wp+j,k_j)\|\partial_t^{k_j}f\|_{L^1}\bigr]\|N^{\wp+j}\psi\|,\qquad k_j=\lfloor\wp+j\rfloor+2,$$
$$|\langle\psi',\phi(f)\psi\rangle|\le e^2C_\wp\bigl[\mathsf s_0(\tfrac\wp2)\|f\|_{L^1}+2\mathsf s(\tfrac\wp2,k_\flat)\|\partial_t^{k_\flat}f\|_{L^1}\bigr]\|N^{\wp/2}\psi'\|\,\|N^{\wp/2}\psi\|,\qquad k_\flat=\lfloor\tfrac\wp2\rfloor+2,$$
and $D(H^\wp)\subset D(\bar\phi(f))$; at $\wp=4$ the two brackets are bounded above by $195\|f\|_{L^1}+537.82\|\partial_t^6f\|_{L^1}$ and $27\|f\|_{L^1}+37.20\|\partial_t^4f\|_{L^1}$.
\item[(c)] \emph{Wedge relation.} There is exactly one antiunitary $J$ with $J\Omega=\Omega$, $J\mathcal D_0=\mathcal D_0$ and $J\phi(f)\psi=\phi(\bar f\circ\mathsf j)J\psi$ for $\psi\in\mathcal D_0$; $J^2=1$, $JU(a)J=U(\mathsf ja)$, $JU(B(s))J=U(B(s))$ and $J\mathcal C(O)J=\mathcal C(\mathsf jO)$. For $f_1,\dots,f_n\in\mathcal F(W)$ and $Y=\phi(f_1)\cdots\phi(f_n)$: $Y\Omega\in D(\Delta^{1/2})$ and $J\Delta^{1/2}Y\Omega=Y^\dagger\Omega$; every boost-invariant dense subspace of $D(\Delta^{1/2})$ is a core of $\Delta^{1/2}$, and $\mathcal D_0(W)$ is one. For $Q\in\mathcal C(W)$: $JQ\Omega\in D(\Delta^{1/2})$ and $\Delta^{1/2}JQ\Omega=Q^*\Omega$. Moreover $\mathcal C(W')\subset\mathcal C(W)'$ on all of $\mathcal H$; hence $O\mapsto\mathcal C(O^\times)$ is isotonous, covariant and local on double cones, $\mathcal C(O^\times)\subset\mathcal A^d(O)$, and (W) holds if and only if $\mathcal C(O^\times)=\mathcal A^d(O)$ for all double cones $O$.
\item[(d)] \emph{What locality says.} For every $\epsilon>0$ there is a real $f_*\in C^\infty_c(\mathbb R^{1+2})$, $f_*\ge0$, $\int f_*=1$, even in each coordinate, with support in $[-\epsilon,\epsilon]^3$ and $\hat f_*(p)>0$ for all real $p$; it may be chosen with $\|\partial_t^kf_*\|_{L^1}\le A_0R_0^k(k!)^2$ for all $k$, $A_0=\frac{e^3\pi}2$, $R_0=\frac{3e\pi^2}{2\epsilon}$. The following are equivalent: (W); $\mathcal C(W)\Omega$ is dense; $\mathcal C(W)'=\mathcal C(W')$; $J\mathcal A(W)J\subset\mathcal C(W)$; (CP); $\phi(f)\Omega\in\mathcal H_0$ for all $f\in\mathcal F(W)$; $\phi(f)\Omega\in\mathcal H_0$ for all $f$ in a set $\mathfrak F\subset\mathcal S$ whose Fourier transforms have no common real zero (for instance $\mathfrak F=\{f_s\}$, one real $f_s$ with $\hat f_s$ free of real zeros); $\inf_{s\in\mathbb R}{\rm dist}(U(B(s))\phi(f_s)\Omega,\mathcal H_0)=0$ for one such $f_s$; for one real $f_0\ge0$, $f_0\ne0$, in $C^\infty_c$, every Lorentz transformation $\Lambda$ has a translation $a_\Lambda$ with $\phi(f_0^{(a_\Lambda,\Lambda)})\Omega\in\mathcal H_0$; (QCP$_{f_s}$) for one such $f_s$ with support in an upright double cone $K$, $\bar K\subset W'$.

The first order of (QCP$_f$) holds for all data (Rec), and (QCP$_f$) gives operators $Q_n\in\mathcal C(W)$ with $\|Q_n\|\le n$ and $Q_n\Omega\to\phi(f)\Omega$. $\mathsf D(f)={\rm dist}(\phi(f)\Omega,\mathcal H_0)$ is a seminorm on $\mathcal S$, invariant under translations and under the boosts $B(s)$, with $\mathsf D(h*f)\le\|h\|_{L^1}\mathsf D(f)$; (W) holds if and only if $\mathsf D\equiv0$. Under (W), $\mathcal A(W)=\mathcal C(W')$, $\Omega$ is cyclic and separating for it, its modular group is $\Delta^{it}=U(B(-2\pi t))$ and its modular conjugation is $J$. (CP) implies, for real $f,g_1,g_2$ with pairwise spacelike separated compact supports,
$$|\mathcal W(fg_1g_2)|^4\le\mathcal W(f^2g_1^2)\,\mathcal W(f^2g_2^2)\,\mathcal W(g_1^2g_2^2),$$
and $\mathcal W(p(f_1,f_2))\ge0$ for every real polynomial $p\ge0$ on $\mathbb R^2$ and real spacelike separated $f_1,f_2$; Motzkin's polynomial $x^4y^2+x^2y^4-3x^2y^2+1$, which is no sum of squares, is among them.
\item[(e)] \emph{Dual net.} Assume (W), and let $O,O_1,O_2$ be bounded, open and convex. Then $\mathcal A^d(O)=\mathcal C(O^\times)=\mathcal A(O^\times)'\supset\mathcal A(O)$ is a von Neumann algebra; $\mathcal A^d$ is isotonous and covariant, $J\mathcal A^d(O)J=\mathcal A^d(\mathsf jO)$, $\Omega$ is cyclic and separating for $\mathcal A^d(O)$, and $\bar\phi(f)$, $f\in\mathcal F(O)$, is affiliated with it; $\mathcal A^d(O_1)\subset\mathcal A^d(O_2)'$ for spacelike separated $O_1,O_2$; $\mathcal A^d(O)'=\mathcal A(O^\times)=\mathcal A^d(O^\times)=\bigvee\{\mathcal A(\tilde W'):\tilde W\supset O\text{ a wedge}\}$; every family $O\mapsto\mathfrak F(O)$ of bounded operators with $\mathfrak F(O)\subset\mathcal A(O_1)'$ for all double cones $O_1\subset O^\times$ has $\mathfrak F(O)\subset\mathcal A^d(O)$; $\mathcal A^d(\tilde W)=\mathcal A(\tilde W)$ for wedges $\tilde W$, so $(\mathcal A^d)^d=\mathcal A^d$; and $\mathcal A(O)=\mathcal A^d(O)$ if and only if $\mathcal C(O^\times)\subset\mathcal C(O)'$. $\mathcal A^d$ is outer regular: for a double cone $O=(a+V_+)\cap(b-V_+)$ and $O_\epsilon=(a-\epsilon e_0+V_+)\cap(b+\epsilon e_0-V_+)$, $\mathcal A^d(O)=\bigcap_{\epsilon>0}\mathcal A^d(O_\epsilon)$. Without (W): $O\mapsto\bigcap_{\mathfrak gW\supset O}\mathcal A(\mathfrak gW)$ is local if and only if (W) holds.
\item[(f)] \emph{Sufficient conditions.} Let $k_0=\max\{1,\lceil\wp\rceil\}$. If $f$ is real and $|\langle\psi,\phi(\partial_tf)\psi\rangle|\le c\langle\psi,N\psi\rangle$ on $\mathcal D_0$, then $\phi(f)$ is essentially self-adjoint on $\mathcal D_0$; if besides $\sup_{m,n}b_{mn}(\partial_t^jf)/(m+n+2)<\infty$ for $1\le j\le k_0+3$, the resolvents of $\bar\phi(f)$ keep $D(H^k)$, $k\le k_0$, off a strip about $\mathbb R$, and all of them lie in $\mathcal C(O)$ for every open $O$ spacelike separated from ${\rm supp}f$. Each of the following implies (W) and (CP).
(f1) There are real $f_i\in C^\infty_c$, ${\rm supp}f_i\subset\{|x|\le\delta_i\}$, $\delta_i\to0$, $\int f_i=1$, $\sup_i\|f_i\|_{L^1}<\infty$, and numbers $c_i,\epsilon_i>0$, nothing uniform in $i$, with $b_{mn}(\partial_tf_i)\le c_i(m+n+2)(1+\log(m+n+2))$ for $|m-n|\le(m+n+2)^{\epsilon_i}$.
(f2) There is a real $f_0\in C^\infty_c$, $f_0\ge0$, $f_0\ne0$, even in $t$, with $\langle\psi,\phi(\partial_tf_0)\psi\rangle\le c\langle\psi,N\psi\rangle$ on $\mathcal D_0$.
(f3) (QA$_{f_0}$) holds for one real $f_0\in C^\infty_c$, $f_0\ge0$, $f_0\ne0$.
(f4) $\wp\le1$, or $\|\mathsf B(s,s)\|\le Cs^{-1}(1+\log s^{-1})$ for $0<s\le1$; then every $\phi(f)$, $f\in\mathcal S$ real, is essentially self-adjoint on $\mathcal D_0$, and the closures commute strongly at spacelike separation.
(f5) There are bounded operators $\mathsf B^e_h(s,s')$, $s,s'>0$, real linear in $h\in C^1_c(\mathbb R^2)$, with $e^{-tH}\mathsf B^e_h(s,s')e^{-t'H}=\mathsf B^e_h(s+t,s'+t')$, $\mathsf B^e_h(s,s')^*=\mathsf B^e_h(s',s)$ and $\|\int\rho_e(\varsigma/\eta)\mathsf B^e_h(\eta+\varsigma,\eta-\varsigma)d\varsigma\|\le C_e|h|_{C^1}$ for all $\eta=\frac1{4(k+2)}$, $k\in\mathbb N_0$, for one measurable $\rho_e\ge0$ on $[-\frac12,\frac12]$ with $\bar\rho_e=\int\rho_e>0$; and with the time average $\mathsf e(f)=\int q_{f_t}(e^{-itH}\cdot,e^{-itH}\cdot)\,dt$ of the sharp-time forms $q_h$ of $\mathsf B^e$ (Lemma~\ref{C-363-20}), the boosts $B_j(s)$ in the direction $e_j$, one $s>0$ and real $\mathsf c,\mathsf c'$:
$$\langle\psi,\phi(f)\psi\rangle=\mathsf c\Bigl[(1+2\coth^2s)\,\mathsf e(f)(\psi,\psi)-\frac1{2\sinh^2s}\sum_{j=1,2}\sum_\pm\mathsf e\bigl(f\circ B_j(\pm s)\bigr)(\psi_{j,\pm},\psi_{j,\pm})\Bigr]+\mathsf c'\int f\,\|\psi\|^2$$
for real $f\in C^\infty_c$ and $\psi\in\mathcal D_0$, $\psi_{j,\pm}=U(B_j(\pm s))^{-1}\psi$. Then (LB$_f$) holds for every real $f\in C^\infty_c$.
If the bound of (f2) holds at $f_0=f_*$, the function of (d), then (f1) holds with the functions $f_i=f_**\mu_i$, $\mu_i\in\mathcal S$; this route to (W) does not use the wedge relation (c).
\item[(g)] \emph{Mixtures and limits.} Let $(\mathsf X,\lambda,(\sigma_\xi))$ be a canonical decomposition of $\sigma$ with all $\sigma_\xi$ data (Rec) with one triple. For real $u\in\mathcal S$ and $c\ge0$, (LB$_{u,c}$), and each of (LB$^\pm_{u,c}$), holds for $\sigma$ if and only if it holds for every $\sigma_\xi$; so do (CP) and (W); and every inequality $\mathcal W(a)\ge0$, $a\in\mathfrak P$ fixed, passes from all $\sigma_\xi$ to $\sigma$. If data (Rec) $\sigma_j$ with one triple converge to $\sigma$ on separated tuples, then $\mathcal W^{\sigma_j}(f_1\cdots f_n)\to\mathcal W^\sigma(f_1\cdots f_n)$ for all $f_k\in\mathcal S$: the real-time functions converge; every inequality $\mathcal W(a)\ge0$, $a\in\mathfrak P$, passes to $\sigma$, and (LB$^\pm_{u,c}$), (CP) and (W) pass from all $\sigma_j$ to $\sigma$.
\item[(h)] \emph{Two bounds towards a linear one.} (h1) Let $\mathsf l,r>0$, $Q\ge1$, $c_2\ge0$, $\mathfrak m\in\mathbb R$, $h'\in C^\infty_c(\mathbb R^2)$ real, $\rho_1\ge0$ even, of class $C^1$, carried by $(-\frac14,\frac14)$, $\int\rho_1=1$, and let $\hat D(x)$, $x\in[-r,r]$, be bounded operators with: (i) $x\mapsto\hat D(x)$ is the restriction of a function holomorphic in norm on a complex neighbourhood of $[-r,r]$; (ii) $\hat D(x)=\hat D(x)^*\ge0$ and $\|\hat D(x)\|\le Q$; (iii) $\hat D(0)=e^{-\mathsf lH}$; (iv) $\hat D'(0)=\mathsf I-\mathfrak me^{-\mathsf lH}$ with $\mathsf I=\int\mathsf l^{-1}\rho_1(\varsigma/\mathsf l)\mathsf B_{h'}(\frac{\mathsf l}2+\varsigma,\frac{\mathsf l}2-\varsigma)\,d\varsigma$; (v) for every $\Lambda<\infty$ and every $\psi$ in $\mathcal H_\Lambda=\mathbf 1_{[0,\Lambda]}(H)\mathcal H$: $\limsup_{\epsilon\downarrow0}\sup_{|x|\le r}\frac{d^2}{dx^2}\langle\psi,-\log(\hat D(x)+\epsilon)\psi\rangle\le c_2\langle\psi,N\psi\rangle$. Then for $\psi\in\mathcal D_\infty$
$$\bigl|\langle\psi,\phi(g_{\mathsf l}\otimes h')\psi\rangle-\mathfrak m\|\psi\|^2\bigr|\le\frac{\mathsf l\langle\psi,H\psi\rangle+\log(Q+1)\|\psi\|^2}r+\frac r2\,c_2\langle\psi,N\psi\rangle,$$
where $\hat g_{\mathsf l}(E)=\mathsf G(\mathsf lE)$, $\mathsf G(y)=\int\rho_1(s)\cosh(sy)\,ds\,\frac{y/2}{\sinh(y/2)}$. No injectivity of $\hat D(x)$ is assumed, and under (i)--(iv) the slope is the real-time average of the field: $\lim_{\epsilon\downarrow0}\frac d{dx}\langle\psi',\log(\hat D(x)+\epsilon)\psi\rangle|_{x=0}=\langle\psi',\phi(g_{\mathsf l}\otimes h')\psi\rangle-\mathfrak m\langle\psi',\psi\rangle$ for $\psi,\psi'\in\mathcal H_\Lambda$. If such bounds hold at all $\mathsf l=L^{-i}$ with constants $\le C\mathsf l^{-N_0}$, then (LB$_{k\otimes h'}$) holds for every real $k\in\mathcal S(\mathbb R)$. (h2) Let $\wp<6$. For a unit timelike vector $n\in V_+$, $\mu\ge1$, Borel sets $\Delta,\Delta'\subset\{p\in\bar V_+:n\cdot p\le\mu\}$ and $u\in\mathcal S$: $\|\mathcal P(\Delta')\phi(u)\mathcal P(\Delta)\|\le e^2C_\wp\|u\|_{L^1}\mu^\wp$. For $\mu\ge1$, $u\in\mathcal S$ and vectors $\psi,\psi'$ of bounded energy in the spectral subspace of $\{p\cdot p\le\mu^2\}$:
$$|\langle\psi',\phi(u)\psi\rangle|\le C_{\rm L}\bigl(\|u\|_{L^1}+\|\mathcal L^4u\|_{L^1}\bigr)\mu^\wp\|\psi'\|\|\psi\|,\qquad\mathcal L=-\partial_t^2-\partial_1^2-\partial_2^2,$$
with $C_{\rm L}=100e^2C_\wp\max\{18^\wp,6^\wp c\,\Sigma_\wp\}$, $\Sigma_\wp=\sum_{r\ge1}(r+1)r^{-8}(2r+3)^\wp$ and a numeral $c$, formed from $C_\wp,\wp$ alone; for the following two consequences assume additionally $\wp>0$. Then for all $m,n$, with $\mathcal P_<$ the spectral projection of $\{p\cdot p\le(m+n+2)^{2/\wp}\}$: $\|\mathcal P_<P_m\phi(u)P_n\mathcal P_<\|\le C_{\rm L}(\|u\|_{L^1}+\|\mathcal L^4u\|_{L^1})(m+n+2)$. With $\Pi$ the spectral projection of $\{p:p\cdot p\le(1+p^0)^{2/\wp}\}$, $\|\Pi e^{-sH}\phi(u)e^{-sH}\Pi\|\le c(u)s^{-1}$ for $0<s\le1$, $c(u)$ formed from $C_{\rm L}$ and the norms $\|\partial_t^ju\|_{L^1}+\|\mathcal L^4\partial_t^ju\|_{L^1}$, $j\le2$. For data (Rec) of any order $\wp$ and real $f\in\mathcal S$, energy and invariant mass together bound the field: $\pm\phi(f)\le c(f)(1+H+M^2)^{\wp/2}$ as forms on the vectors of bounded energy, $M^2=\int(E^2-|q|^2)\,d\mathcal P$ the operator of the invariant mass squared, $c(f)$ formed from $C_\wp$, $\wp$ and finitely many of the norms $\|\partial_\nu^kf\|_{L^1}$, $\nu=0,1,2$.
\end{enumerate}
\end{theorem}

\begin{theorem}[Locality of the continuum limits]\label{C-363-4}
Let $0<g\le\gamma$, let $\sigma$ be the centred functional of a volume-first limit at $g$, and let $\tau$ be $\sigma$ or a component of a canonical decomposition of $\sigma$; every such limit, decomposition and component is admitted. On $\mathcal H(\tau)$:
\begin{enumerate}
\item[(a)] $\mathcal C_\tau(W')\Omega$ is dense, and (W) holds: $O\mapsto\mathcal A_\tau(O)$ is an isotonous, covariant and local net of von Neumann algebras, $\mathcal L_\tau=\mathcal A_\tau$, and $\mathcal H_0=\mathcal H(\tau)$.
\item[(b)] $\mathcal A_\tau(W)=\mathcal C_\tau(W')=\mathcal C_\tau(W)'$ and $\mathcal A_\tau(W)'=\mathcal A_\tau(W')$; $\Omega$ is cyclic and separating for $\mathcal A_\tau(W)$, with modular group $\Delta^{it}=U(B(-2\pi t))$ and modular conjugation $J$; in particular $\mathcal A_\tau(W)\ne\mathbb C1$ unless $\mathcal H(\tau)=\mathbb C\Omega$.
\item[(c)] For open $O$ and $f\in\mathcal F(O)$, $\bar\phi(f)$ is affiliated with $\mathcal A_\tau(O)$.
\item[(d)] (CP) holds, with the two inequalities of Theorem~\ref{C-363-3}(d).
\item[(e)] The dual net $\mathcal A^d_\tau(O)=\mathcal C_\tau(O^\times)=\mathcal A_\tau(O^\times)'$ has all the properties of Theorem~\ref{C-363-3}(e): it is local, Haag dual, the largest family local relative to $\mathcal A_\tau$, and its own dual net.
\item[(f)] Let $\tau=\sigma$, let $\Lambda_j$ be the tori of its truncated functions and $\mathcal M_{\Lambda_j}$ the algebras of all bounded operators on the bonds of the right halves of their time-zero planes (Proposition~\ref{C-363-14}). Every weak limit $Q$, along a subsequence, of operators $Q_j\in\mathcal M_{\Lambda_j}$ lies in $\mathcal C_\sigma(W')$; the lattice modular flows converge to the boosts, $\Delta_{\Lambda_j}^{it}\psi_j\Rightarrow U(B(-2\pi t))\psi$ in the sense of Lemma~\ref{C-363-8}; for real $u\in C^\infty_c(\mathsf I)$ and $\epsilon>0$ there is $T^{(\epsilon)}_u\in\mathcal C_\sigma(W')$, a weak limit of operators of $\mathcal M_{\Lambda_j}$, with $\|T^{(\epsilon)}_u\|\le(2\sqrt\epsilon)^{-1}$ and $\|T^{(\epsilon)}_u\Omega-\phi(f_u)\Omega\|\le\frac{\sqrt\epsilon}2\pi e^{5\pi^2/8}\|[\theta_1u\cdot u]\|$, $f_u$ the function of Lemma~\ref{C-363-9}(c); and $\mathcal A_\sigma(W)$ is the von Neumann algebra generated by the $U(\mathfrak g)QU(\mathfrak g)^{-1}$, $\mathfrak gW\subset W$, with $Q$ or $Q^*$ such a weak limit.
\item[(g)] For $\tau=\sigma$ and a canonical decomposition: $[G]\mapsto([G]_\xi)_\xi$ is a unitary of $\mathcal H(\sigma)$ onto $\int^\oplus\mathcal H(\sigma_\xi)\,\lambda(d\xi)$, by which we identify the two, $[G]=\int^\oplus[G]_\xi$; the semigroup, the translations, the insertions, the operators $\mathsf B(s,s')$ and the fields decompose; the continuous functions on $\mathsf X$ act by multiplication (the operators of Lemma~\ref{C-363-13}(i)), generate the diagonal algebra and lie in the centre of $\mathcal A_\sigma(W)$; and $\mathcal A_\sigma(W)=\int^\oplus\mathcal A_{\sigma_\xi}(W)\,\lambda(d\xi)$: for $\lambda$-almost every $\xi$, among them every $\xi$ with $\lambda(\{\xi\})>0$, the fibre at $\xi$ of the algebra of (f) is the wedge algebra of the component $\sigma_\xi$.
\item[(h)] \emph{A supply.} Let $\sigma'$ be a limit of periodic Gibbs states $\rho_j$ at $g$, $c\in\bigcup_KL^{-K}\mathbb Z$, and $B_j$ Borel functions, $|B_j|\le\mathsf b$, of finitely many bonds of $L^{-K_j}\mathbb Z^3$ in $\{x_3=0,\ x_1\le c\}$. Then $|\rho_j(\overline{\theta(F')_j}B_j(F)_j)|\le\mathsf b\,\rho_j(\overline{\theta(F')_j}(F')_j)^{1/2}\rho_j(\overline{\theta(F)_j}(F)_j)^{1/2}$ for $F,F'$ over $\mathbb E_+$ and large $j$; along a subsequence $\rho_j(\overline{\theta(F')_j}B_j(F)_j)\to\langle[F'],Q[F]\rangle$ for all such $F,F'$, with $\|Q\|\le\mathsf b$; and every such $Q$ lies in $\mathcal C_{\sigma'}(W+ce_1)$. The same holds for every component $\sigma'_\xi$ of a canonical decomposition of $\sigma'$, with $\rho_j$ replaced, as the states that integrate, by the periodic Gibbs states $\rho_{j,\xi}$ that approximate the component (Theorem~\textup{\YMMref{AP}{780}} among the construction results, recalled in the results used), the $\rho_j$ still giving the centring.

By contrast, on the space of a component an operator $e^{-sH}V_1e^{-sH}$, $s>0$, $V_1$ bounded, lies in no $\mathcal C(O)$, $O\ne\emptyset$ open, unless it is $0$ or the space is $\mathbb C\Omega$: weak limits of the operators $[G]\mapsto[D_j\cdot\tau_{te_3}G]$ of lattice functions $D_j$ placed strictly between two times, in $\{s<x_3<t-s\}$ with $s,t\in\bigcup_KL^{-K}\mathbb Z$, whose inner operators $[G]\mapsto[\tau_{-se_3}D_j\cdot\tau_{(t-2s)e_3}G]$ are bounded uniformly in $j$, on a nontrivial component space are nonlocal unless they vanish (Proposition~\ref{C-363-18}(a)).
\end{enumerate}
\end{theorem}

\begin{theorem}[Energy bounds of the field]\label{C-363-5}
Let $0<g\le\gamma$ and let $\sigma'$ be a limit of periodic Gibbs states at $g$ or a component of a canonical decomposition of one. Let $f\in\mathcal S$, $m,n\in\mathbb N_0$, $\psi,\psi'\in\mathcal D_\infty$, and let $c_V,\bar\Phi,c'_V,\Phi'$ be the numbers of Lemma~\ref{C-363-16}; none of them depends on the cutoff, the sequence, $\sigma'$ or $\xi$.
\begin{enumerate}
\item[(a)] \emph{Order $3$.} With $\mathsf C_3=\frac{4\sqrt8}\varrho L^4\,40M(40+8L^2)e^{\frac83c_V+\frac12\bar\Phi}$ and $\mathsf c_k=27L^3\mathsf C_3\|\partial_t^kf\|_\star$: $b_{mn}(f)\le\mathsf c_0(m+n+2)^3$ for all $m,n$, and $b_{mn}(f)\le\mathsf c_k(m+n+2)^3(|m-n|-1)^{-k}$ if $|m-n|\ge2$, $k\in\mathbb N$; $\|\phi(f)\psi\|\le8(71\mathsf c_0+6.3\mathsf c_5)\|N^3\psi\|$, and $\|N^a\phi(f)\psi\|\le c_a(f)\|N^{a+3}\psi\|$ for real $a$, so that $D(N^3)\subset D(\bar\phi(f))$; $|\langle\psi',\phi(f)\psi\rangle|\le8(17.1\mathsf c_0+5.7\mathsf c_3)\|N^{3/2}\psi'\|\|N^{3/2}\psi\|$; and $\|e^{-sH}\phi(f)e^{-sH}\|\le c(f)s^{-3}$ for $0<s\le1$. If $f$ is real and $\|\partial_t^kf\|_\star\le A^{k+1}(k!)^{\mathsf s}$ for all $k\ge0$, with $\mathsf s\ge1$ (such $f\ge0$, $f\ne0$, of compact support exist for every $\mathsf s>1$), then $\|\phi(f)^n\Omega\|\le C^n(n!)^{3\mathsf s}$ for all $n$, with $C=C(A,\mathsf s,L^3\mathsf C_3)$.
\item[(b)] \emph{Order $2$.} With $\mathsf C_2=\frac{4(\sqrt8+1)}\varrho L^4\,40M(40+8L^2)e^{\frac83c'_V+\frac14\Phi'}$: $b_{mn}(f)\le27L^2\mathsf C_2\|f\|_\star(m+n+2)^2$, and $b_{mn}(f)\le27L^2\mathsf C_2\|\partial_t^kf\|_\star(m+n+2)^2(|m-n|-1)^{-k}$ if $|m-n|\ge2$; $\|\phi(f)\psi\|\le c(f)\|N^2\psi\|$, $|\langle\psi',\phi(f)\psi\rangle|\le c'(f)\|N\psi'\|\|N\psi\|$ and $\|e^{-sH}\phi(f)e^{-sH}\|\le c''(f)s^{-2}$ for $0<s\le1$; and for $f$ as at the end of (a), $\|\phi(f)^n\Omega\|\le C^n(n!)^{2\mathsf s}$ with $C=C(A,\mathsf s,L^2\mathsf C_2)$.
\item[(c)] \emph{At a point; at one time.} Here Theorems~\YMMref{AJ}{673} and~\YMMref{AD}{607} are used at every radius and for every admitted choice of the constants $(\sigma^-,\gamma)$, at all couplings at once. With $q=\mathfrak p_0+\frac16$ and $C$ formed from $\sigma^-,\gamma,\varrho,g$: $\|\mathsf B(s,s')\|\le C\mu^{-7/2}(1+\log\mu^{-1})^q$, $\mu=\min\{s,s',\frac12\}$; so for every $\epsilon>0$, $\sigma'$ is data (Rec) with a triple $(C_\epsilon,\frac72+\epsilon,\frac12)$, and Theorem~\ref{C-363-3} applies at the order $\frac72+\epsilon$. For real $k\in C^1_c(\mathbb R^2)$ carried by a square of side $2$: $\|\mathsf B_k(s,s')\|\le C|k|_{C^1}\mu^{-5/2}(\log\frac e\mu)^q$, $\mu=\min\{s,s',1\}$; so $b_{mn}(f)\le c(f)(m+n+2)^{5/2}(1+\log(m+n+2))^q$ for all $m,n$.
\item[(d)] \emph{Sharpness.} Let $0<g\le\gamma_{\rm sd}$ and let $\sigma'$ be a component. For continuous $k\ge0$, $k\ne0$, of compact support, $\liminf_{u\downarrow0}u^2\|\mathsf B_k(u,u)\|\ge(\pi c_{\rm sd}/32)^{1/2}\|k\|_{L^2}$; so the field at one time, the form $\mathsf T(k)$ on the vectors of bounded energy, has no bound $\pm\mathsf T(k)\le cN$. $\|\mathsf B(s,s)\|\ge\|\mathsf B(s,s)\Omega\|\ge\frac18c_{\rm sd}^{1/2}s^{-3}$ for $0<2s\le r_{\rm sd}$: the infimum of admissible pure-power exponents in the point bound lies in $[3,\frac72]$. For $0\ne h\ge0$ in $C^1_c(\mathbb R^2)$, $\|\mathsf I_{\mathsf l}(h)\Omega\|\ge c(h)\mathsf l^{-2}$ for small $\mathsf l$; so no bound $\|\mathsf I_{\mathsf l}(h)\|\le C\mathsf l^{-p}|h|_{C^1}$ holds with $p<2$, and the exponent $2$ of the plates in Lemma~\ref{C-363-16}(c) is exact. For $0\ne k\ge0$ in $C^1_c(\mathbb R^2)$, $\epsilon\in(0,1]$ and $q'<\frac32\epsilon$:
$$\sup\bigl\{\|T_{mn}(k)\|(m+n+2)^{-q'}:\ m,n\ge0,\ |m-n|\le(m+n+2)^\epsilon\bigr\}=\infty;$$
in particular no bound $\|T_{mn}(k)\|\le c(m+n+2)$ holds on a band with $\epsilon>\frac23$.
\item[(e)] \emph{Slab families.} Let $\sigma'$ be a component or the centred functional of a volume-first limit, $\mathsf l=L^{-i}$, $i\ge1$, $\vec y^*\in L^{-i-2}\mathbb Z^2$, $h'\ne0$ real and carried by $\{|\vec x-\vec y^*|_\infty\le\frac{\mathsf l}4\}$, $\rho_1$ as in Theorem~\ref{C-363-3}(h1), and $f(x)=\mathsf l^{-1}\rho_1(x_3/\mathsf l)h'(\vec x)$. The limits $\mathbb D(z)$, along a subsequence, of the lattice operators $e^{z\mathcal O(f)-\mathcal E(z)}$ between reflection-positive vectors, $\mathcal E$ the normaliser of the one-cell bound (Proposition~\YMMref{89.1}{690}), are holomorphic on $|z|<\varrho/\|f\|_\sharp$ with $\|\mathbb D(z)\|\le\bar{\mathsf Q}$, a number free of $\mathsf l,f,\xi$, $\mathbb D(0)=e^{-\mathsf lH}$ and $\mathbb D'(0)$ the insertion of $f$ plus a multiple of $e^{-\mathsf lH}$, and $\hat D(x)=\mathbb D(x)/\langle\Omega,\mathbb D(x)\Omega\rangle$, $|x|\le\varrho/(2(\bar{\mathsf Q}+1)\|f\|_\sharp)$, has (i)--(iv) of Theorem~\ref{C-363-3}(h1) with $Q=2\bar{\mathsf Q}$ and $\mathfrak m=q_0(\Omega,\Omega)\int h'$.\qedhere
\end{enumerate}
\end{theorem}

\begin{theorem}[Examples: Wick powers of free fields; free photons]\label{C-363-6}
\begin{enumerate}
\item[(a)] \emph{Gaussian families.} Let $(\mathsf M,\mu)$ be a $\sigma$-finite measure space with $L^2(\mathsf M,\mu)$ separable, $\omega>0$ and $\vec p$ measurable on it with $|\vec p|\le\omega$, $\kappa_\alpha$, $\alpha$ in a finite set $\mathsf A$, measurable, with signs $\varepsilon_\alpha$, such that $\|\omega^{-1/2}\kappa_\alpha e^{-s\omega}\|\le c_\kappa s^{-\delta}$ for $0<s\le1$, $\delta\ge0$, and $G=(G_{\alpha\beta})$ real and continuous on $\mathbb R^3\setminus\{0\}$ with $G_{\alpha\beta}(-x)=G_{\beta\alpha}(x)$, such that the vectors $e^\alpha_x=\kappa_\alpha e^{-x_3\omega-i\vec x\cdot\vec p}$, $x\in\mathbb E_+$, have the inner products $\langle e^\alpha_x,e^\beta_y\rangle=\varepsilon_\alpha G_{\alpha\beta}(\theta x-y)$.

Let $\mathrm O(3)$ act on $\mathbb R^{\mathsf A}$ by a homomorphism $\rho_{\mathsf A}$ into the orthogonal matrices, with $G(\mathsf Rx)=\rho_{\mathsf A}(\mathsf R)G(x)\rho_{\mathsf A}(\mathsf R)^{\rm T}$ for $\mathsf R\in\mathrm O(3)$ and $\rho_{\mathsf A}(\theta)={\rm diag}(\varepsilon_\alpha)$: the reflection acts on the components by the signs (for photons $\mathsf A$ indexes the three derivatives, times the colours, and $\rho_{\mathsf A}(\mathsf R)=\mathsf R\otimes1$; for a scalar family $\rho_{\mathsf A}$ is trivial). Let $\mathsf P(t)=\sum_{\underline\alpha}\mathsf p_{\underline\alpha}t_{\alpha_1}\cdots t_{\alpha_k}\ne0$ be real and $k$-homogeneous, $k\ge1$, with a symmetric real tensor $\mathsf p$, and invariant, $\mathsf P(\rho_{\mathsf A}(\mathsf R)t)=\mathsf P(t)$ for $\mathsf R\in\mathrm O(3)$; $|\mathsf p|_1=\sum|\mathsf p_{\underline\alpha}|$.

The Schwinger functions of ${:}\mathsf P(\Phi){:}$ for the Gaussian field $\Phi=(\Phi_\alpha)$ of the covariance $G$, that is, the sums over the complete pairings of $n$ vertices with $k$ legs each, no leg paired within its vertex, of the products of the $G_{\alpha\beta}(x_i-x_l)$ weighted by $\mathsf p$, are data (Rec) with $C_\wp=2^\wp2^k|\mathsf p|_1(2^\delta c_\kappa)^k(k/e)^{k/2}$, $\wp=k(\delta+\frac12)$, $c_0=\frac12$, and a unique vacuum. A unitary $V$ carries $\mathcal H$ onto a closed subspace of the Fock space over $L^2(\mathsf M,\mu)$ that reduces $d\Gamma(\omega)$, $d\Gamma(\vec p)$ and the Fock operators of the field at a point, with $V\Omega=\Omega_{\rm F}$ and $VHV^{-1}=d\Gamma(\omega)$ there, and $\langle\chi,V\phi(f)\psi\rangle=\Phi_f(\chi,V\psi)$ for every $\chi\in\mathcal D_\infty(d\Gamma(\omega))$ and $\psi\in\mathcal D_\infty$, $\Phi_f$ the form of the smeared Wick polynomial.
\item[(b)] \emph{Growth.} If $k\ge3$ and for every $u\in C^\infty_c\setminus\{0\}$ the function 
\[
w_u(m_1,\dots,m_k)=\hat u(\sum_i\omega(m_i),\sum_i\vec p(m_i))\sum_{\underline\alpha}\mathsf p_{\underline\alpha}\prod_i\kappa_{\alpha_i}(m_i)
\]
 lies in $L^2(\mathsf M^k)\setminus\{0\}$, then $\|\phi(u)^n\Omega\|\ge c_u^n((kn)!)^{1/2}$ for real $u\in C^\infty_c\setminus\{0\}$; so (QA$_u$) fails, no bound $|\langle\psi,\phi(u)\psi\rangle|\le c\langle\psi,N\psi\rangle$ holds at any real $u\ne0$ in $C^\infty_c$, no bound $b_{mn}(u)\le\beta(m+n+2)^q$ with $q<\frac k2$ holds on any band $|m-n|\le(m+n+2)^\epsilon$, and no triple of the datum has $\wp<\frac k2$.
\item[(c)] \emph{Wick powers.} For the free field $\chi$ of mass $\mathsf m>0$ and $\mathsf P=t^k$, (a) holds with every $\delta\in(0,\frac12]$ and $c_\kappa=(8\pi\delta)^{-1/2}(2\mathsf m)^{-\delta}$, so ${:}\chi^k{:}$ is data (Rec) of every order $\wp>\frac k2$; for $k\ge2$, $w_u\ne0$ as soon as $\hat u$ does not vanish identically on $\{E>(|\vec k|^2+k^2\mathsf m^2)^{1/2}\}$, which holds for every $u\in C^\infty_c\setminus\{0\}$, for every $u\in\mathcal S\setminus\{0\}$ that vanishes off $\bar W$, and for $u=\mathsf g\otimes h'$ with $\hat{\mathsf g}$ free of real zeros and $h'\in C^\infty_c(\mathbb R^2)\setminus\{0\}$; so (b) holds for $k\ge3$. ${:}\chi^6{:}$ is data (Rec) with $(93312(\pi e)^{-3}\mathsf m^{-1},4,\frac12)$, a unique vacuum, ${\rm spec}\,H\subset\{0\}\cup[2\mathsf m,\infty)$ and $S_2(x)=720(4\pi|x|)^{-6}e^{-6\mathsf m|x|}$. (W) and (CP) hold for every ${:}\chi^k{:}$, $k\ge1$, $\mathsf m>0$; for every real homogeneous Wick polynomial in the first derivatives $\partial_\alpha\varphi^a$ of $\mathsf n$ massless free fields $\varphi=(\varphi^a)_{a\le\mathsf n}$ that is invariant under $\mathrm O(3)$ and $\mathrm O(\mathsf n)$, for instance ${:}(\partial\varphi\cdot\partial\varphi)^2{:}$ with $\wp=6$; and for the Wick cube of the generalised free field with the kernel $r^{-2}$ ($\wp=3$, $S_2=6r^{-6}$, no gap).

So neither (LB$_u$) nor (QA$_u$), at any $u$, is necessary for (W) among data (Rec) with $\wp=4$, $c_0=\frac12$, a unique vacuum, a gap and $S_2\sim r^{-6}$; and for every $\wp>\frac32$ there are data (Rec) of order $\wp$ without (LB$_u$) at any $u\ne0$.
\item[(d)] \emph{Free photons.} Let $\varphi$ be the free massless field on $\mathbb R^{1+2}$ and $\Theta={:}\partial_\lambda\varphi\,\partial^\lambda\varphi{:}$. The Schwinger functions of ${:}\partial_\alpha\Phi\,\partial_\alpha\Phi{:}$ for the Euclidean massless free field $\Phi$ on $\mathbb R^3$ (the pairings of $-\partial_a\partial_b(4\pi|x|)^{-1}$) are data (Rec) with $(24/\pi e,3,\frac12)$, a unique vacuum and $S_2=12/16\pi^2r^6$; their space is the whole even Fock space, with $V\phi(f)V^{-1}=-\Theta(f)$. For real $f\in\mathcal S$: $\pm\Theta(f)\le16\pi|f|_8\,N$, $|f|_8$ the Schwartz seminorm in equation~\YMMref{(1873)}{716} with $|\hat f(k)|\le16\pi^2|f|_8(1+|k|^2)^{-2}$; $\|\Theta(f)\psi\|\le c(f)\|N\psi\|$; $\|\Theta(f_1)\cdots\Theta(f_n)\Omega\|\le2^nn!\prod_kc(f_k)$; and $\Theta(f)$ is essentially self-adjoint on every core of $H$; (W) and (CP) hold. $\|\mathsf B(s,s)\|\ge\|\mathsf B(s,s)\Omega\|=\frac{\sqrt3}{16\pi}s^{-3}$, so no triple has $\wp<3$, while on its own space the order is exactly $1$: for real $f\ne0$, $\alpha<1$ and every $c$ some $\psi\in\mathcal D_\infty$ has $|\langle\psi,\phi(f)\psi\rangle|>c\langle\psi,N^\alpha\psi\rangle$. At one time, with $\Theta_k=\int k(\vec x)\Theta(0,\vec x)d^2x$, $k\in\mathcal S(\mathbb R^2)$ real, and $G=m+n+2$:
$$\|P_m\Theta_kP_n\|\le(1+\sqrt3)\|k\|_\infty G+2\mathsf n_k(|m-n|+2)\sqrt G,\qquad\mathsf n_k=\tfrac1{2\pi}\bigl[10.4\|k\|_2^2+1.6\|k\|_2\|\nabla k\|_2\bigr]^{1/2},$$
and the second term is attained: for $k\ne0$, every integer $\omega\ge\max\{3,7\xi_0(k)+1\}$ and every $j\in\mathbb N$ there are $m,n$ with $|n-m-\omega|\le1$, $2j(\omega-\frac12)+\omega-1<G\le(j+1)(2\omega+3)$ and $\|P_n\Theta_kP_m\|\ge3\cdot10^{-3}\|k\|_2\,\omega\sqrt G-\kappa_k(\omega-2)G$, where $\kappa_k(R)=(2\pi)^{-2}\int_{|\xi|>R}|\tilde k|$ and $\xi_0(k)$ is the least $R$ with $\int_{|\xi|\le R}|\tilde k|^2\ge\frac12\|\tilde k\|_2^2$. So for real $k\ne0$ in $C^\infty_c$ the band bound $\|P_me^{-H/G}\Theta_ke^{-H/G}P_n\|\le CG$ on $|m-n|\le G^\epsilon$ holds if and only if $\epsilon\le\frac12$.
\item[(e)] \emph{$\mathsf n$ colours.} Let $\Theta=\sum_{a\le\mathsf n}{:}\partial_\lambda\varphi^a\partial^\lambda\varphi^a{:}$, $\mathsf n\ge1$. The closed span of $\Omega$ and the $\Theta(f_1)\cdots\Theta(f_n)\Omega$ ($f_k\in\mathcal S$, or $f_k\in C^\infty_c$ real) is the space of all $\mathrm O(\mathsf n)$-invariant vectors, and it is the closure of the Euclidean vectors; with $\mathrm{SO}(3)$ in the place of $\mathrm O(3)$ at $\mathsf n=3$ this is false. $\|\mathsf B(s,s)\Omega\|=\frac{\sqrt{3\mathsf n}}{16\pi}s^{-3}$. For real $f\ne0$, $\alpha<1$, every $c$ and either sign there is $\psi\in\mathcal D_0$ with $\pm\langle\psi,\phi(f)\psi\rangle>c\langle\psi,N^\alpha\psi\rangle$, while $\pm\phi(f)\le\mathsf n\,16\pi|f|_8N$: on the invariant space the order is exactly $1$, and $\sup_{m,n}b_{mn}(f)(1+\max\{m,n\})^{-p}=\infty$ for every $p<1$.
\item[(f)] \emph{The net of the action density of free photons.} Let $\mathsf n\ge1$ and $O$ be a double cone. For a family $\Phi$ of Wightman fields on the Fock space of $\mathsf n$ free fields $\varphi=(\varphi^a)_{a\le\mathsf n}$ of a mass $\mathsf m\ge0$, among them $\Theta=\sum_a{:}\partial_\lambda\varphi^a\partial^\lambda\varphi^a{:}$, $S={:}\varphi\cdot\varphi{:}=\sum_a{:}\varphi^a\varphi^a{:}$ and the multiplet $({:}\partial_\mu\varphi\cdot\partial_\nu\varphi{:})_{\mu\nu}$: $\mathcal P_\Phi(O)$ is the algebra of the polynomials in the $\Phi(f)$, $f\in\mathcal F(O)$, $\mathcal D_0^\Phi=\mathcal P_\Phi(\mathbb R^{1+2})\Omega$ and $\mathcal H_\Phi$ its closure, $\mathcal C_\Phi(O)$ the set of the bounded $Q$ on $\mathcal H_\Phi$ with $\langle\Phi(\bar f)\psi,Q\psi'\rangle=\langle\psi,Q\Phi(f)\psi'\rangle$ for $f\in\mathcal F(O)$ and $\psi,\psi'\in\mathcal D_0^\Phi$, $\mathcal A_\Phi(O)=\mathcal C_\Phi(O)'$, and $\mathcal A^d_\Phi(O)=\bigcap\{\mathcal A_\Phi(\tilde W):\tilde W\supset O\text{ a wedge}\}$ its dual net.

At the mass $0$ the unitary $V$ of (a), for the data of (d), (e), carries $\mathcal D_0$, $\mathcal C(O)$, $\mathcal A(O)$ and $\mathcal A^d(O)$ of the datum onto $\mathcal D_0^\Theta$, $\mathcal C_\Theta(O)$, $\mathcal A_\Theta(O)$ and $\mathcal A^d_\Theta(O)$, so that Theorem~\ref{C-363-3} applies to $\Theta$; there $\mathcal H_\Theta$, $\mathcal H_S$ and the space of the multiplet are all the space of the $\mathrm O(\mathsf n)$-invariant vectors, and $S$ is no datum (Rec), its point operators being unbounded in the infrared: its nets are those of the definition alone. For every real $f\in\mathcal F$, the massless square $S(f)$ is essentially self-adjoint on its polynomial cyclic domain $\mathcal D_0^S$. At a mass $\mathsf m>0$, $\Theta$ and $S$ are data (Rec) by (a), and the unitary $V$ of (a) carries the nets of each onto those of the definition.

At the mass $0$: $\mathcal A_\Theta(O)$ is the $\mathrm O(\mathsf n)$-invariant part of the net of the field strengths, restricted to the invariant vectors (for $\mathsf n=1$ it is generated by the $\cos\bar\varphi(h)$, $h\in\mathcal F_{\mathbb R}(O)$, $\int h=0$, where $\bar\varphi(h)$ is the closure of the field smeared with the real $h\in C^\infty_c(O)$ and $\int h$ is the integral over space-time); the multiplet ${:}\partial_\mu\varphi\cdot\partial_\nu\varphi{:}$ has the same net, and for $\mathsf n\in\{1,3\}$ the closures of the $T_{\mu\nu}(f)$, $T_{\mu\nu}={:}\partial_\mu\varphi\cdot\partial_\nu\varphi{:}-\frac12\eta_{\mu\nu}\Theta$ with $\eta={\rm diag}(1,-1,-1)$, $T_{00}(f)$ among them, $f\in\mathcal F_{\mathbb R}(O)$, are self-adjoint and affiliated with $\mathcal A_\Theta(O)$; (W) holds. $\mathcal A_\Theta(O)\ne\mathcal A^d_\Theta(O)$: for real $h$ carried by $O$ with $\int h\ne0$ the mean over $\mathrm O(\mathsf n)$ of $e^{i\bar\varphi^1(h)}$ lies in $\mathcal C_\Theta(O^\times)=\mathcal A^d_\Theta(O)$ and not in $\mathcal A_\Theta(O)$; for $\mathsf n=1$, $\cos(\mathsf s\bar\pi(g))$, $\bar\pi(g)$ the closure of the momentum at time $0$ smeared with a real $g\in C^\infty_c(\mathbb R^2)$ that is $1$ near the closed base of $O$, lies in $\mathcal C_\Theta(O)$ and does not commute with $\cos\bar\varphi(h)$ when $\mathsf s\int h\notin2\pi\mathbb Z$.

For $\mathsf n=1$ and closed subgroups $\Lambda\subset\mathbb R$ the even parts of the nets $\{e^{i\bar\varphi(f)}:f\in\mathcal F_{\mathbb R}(O),\int f\in\Lambda\}''$ are pairwise different local nets between $\mathcal A_\Theta$ and $\mathcal A^d_\Theta$ with the same wedge algebras. At a mass $\mathsf m>0$, where $\varphi^a(h)=\varphi^a(-\mathsf m^{-2}\Box h)$ makes every smeared field one of a test function of integral $0$ and $\Theta=\frac12(\Box+2\mathsf m^2)S$, and where, as at the mass $0$ and by the same steps, which hold at every $\mathsf m\ge0$, the commutators of the $\Theta(f)$, $f\in\mathcal F(O)$, reach the colour sums $\sum_a\varphi^a(h)\varphi^a(h')$, whose commutators have the same structure constants at every mass, for real $h,h'\in\mathcal F(O)$ of integral $0$, hence for all real $h,h'\in\mathcal F(O)$ (every $Q\in\mathcal C_\Theta(O)$ commutes weakly with them), among them the squares $\sum_a\varphi^a(g)^2$, whose closures generate $\mathcal A_S(O)$, the nets of $\Theta$ and of $S$ are one, $\mathcal A_\Theta=\mathcal A_S$ on $\mathcal H_\Theta=\mathcal H_S$, for every $\mathsf n\ge1$, and the net of $\Theta$, as that of $S$, is its own dual net.

The net of $S={:}\varphi\cdot\varphi{:}$ is its own dual net at every mass; at the mass $0$, where $\Theta(f)=\frac12S(\Box f)$, $\mathcal A_\Theta(O)\subsetneq\mathcal A_S(O)=\mathcal A^d_S(O)=\mathcal A^d_\Theta(O)$, and the two nets have the same wedge algebras: on one space, with one vacuum and one family of wedge algebras, the field $S$ is Haag dual and the field $\frac12\Box S$ is not.
\end{enumerate}
\end{theorem}

\begin{theorem}[Lattice algebras at one time; the defect $\mathbb Z_2$ and its Euclidean half]\label{C-363-7}
(a) Let $G$ be a compact group, $(\Sigma,\mathsf B)$ a finite graph, $\mathcal H_{\rm ph}\subset L^2(G^{\mathsf B})$ the gauge-invariant subspace, and $\mathsf A(Y)$, $Y\subset\mathsf B$, the algebra of all gauge-invariant bounded operators of the bonds of $Y$ on $\mathcal H_{\rm ph}$. Then $\mathsf A(Y)$ is a von Neumann algebra, $\mathsf A(Y)'=\mathsf A(\mathsf B\setminus Y)$ for every $Y$, $\mathsf A(\mathsf B)=\mathcal B(\mathcal H_{\rm ph})$, and $\mathsf A(Y_1\cup Y_2)=\mathsf A(Y_1)\vee\mathsf A(Y_2)$ if $Y_1,Y_2$ have no common site.

In (b)--(i), $G=\mathrm{SU}(2)$. On a layer $\mathbb Z/N_1\times\mathbb Z/N_2$, $N_i\ge4$, let, for a set $S$ of plaquettes, $\mathsf B_S$ be their bonds, $\mathsf A_\bullet(S)$ the von Neumann algebra generated by the traces $M_{t_p}$, $p\in S$, and the spin projections $\Pi^j_b$, $b\in\mathsf B_S$; $Z^1(S)$ the sign fields $c$ on $\mathsf B_S$ with $\prod_{b\in\partial p}c_b=1$ for $p\in S$; $B^1$ the coboundaries of the whole layer, $c_b=z(b_-)z(b_+)$ for a sign field $z$ on all its sites, and $B^1(S)\subset Z^1(S)$ the coboundaries on $\mathsf B_S$, $c_b=z(b_-)z(b_+)$, $b\in\mathsf B_S$, for a sign field $z$ on the sites of $\mathsf B_S$; $H^1(S)=Z^1(S)/B^1(S)$; $\mathsf V(c)$ the unitary of $V_b\mapsto c_bV_b$, a sign field on a set of bonds being continued by $1$ to the other bonds; and for a dual path $\gamma$, $\mathsf V_\gamma=\mathsf V(c)$ with $c=-1$ exactly on the bonds that $\gamma$ crosses. $S$ is admissible if it has at least three plaquettes, any two of them are joined by a chain of plaquettes of $S$ each sharing a bond with the next, and the plaquettes of $S$ at each site are so joined through bonds at that site.

For every closed path $C'$ and spin $j$: $\mathsf V(c)M_{\chi_j(U(C'))}=c(C')^{2j}M_{\chi_j(U(C'))}\mathsf V(c)$, $\mathsf V(c)\mathsf V(c')=\mathsf V(cc')$; $\mathsf V(c)=1$ on $\mathcal H_{\rm ph}$ for $c\in B^1$, $\mathsf V(c)$ being then the gauge transformation by the central field $z$; and for $c\in B^1(S)$, with $z$ continued by $1$ to the other sites, $\mathsf V(c)=\mathsf V(c^\circ)$ on $\mathcal H_{\rm ph}$, $c^\circ$ the restriction of $(z(b_-)z(b_+))_b$ to the bonds off $\mathsf B_S$, so that $\mathsf V(c)|_{\mathcal H_{\rm ph}}$, a function of spin projections of bonds off $\mathsf B_S$, lies in $\mathsf A(\mathsf B\setminus\mathsf B_S)=\mathsf A(\mathsf B_S)'$ and ${\rm Ad}\,\mathsf V(c)$ is trivial on $\mathsf A(\mathsf B_S)$; on a periodic layer the plaquettes on which a sign field is not closed are even in number: there is no single point operator.
\begin{enumerate}
\item[(b)] For admissible $S$: $\mathsf A_\bullet(S)=\mathsf A(\mathsf B_S)\cap\{\mathsf V(c):c\in Z^1(S)\}'$ and $\mathsf A_\bullet(S)'=\mathsf A(\mathsf B\setminus\mathsf B_S)\vee\{\mathsf V(c):c\in Z^1(S)\}$; the $\mathsf V(c)$, $c\in Z^1(S)$, lie in the centre of $\mathsf A_\bullet(S)$; ${\rm Ad}\,\mathsf V(c)$ on $\mathsf A(\mathsf B_S)$ depends on the class of $c$ in $H^1(S)=Z^1(S)/B^1(S)$ alone, so that one $c$ of each class serves in both identities, and $\mathsf A(\mathsf B_S)=\bigoplus_\kappa\mathsf A(\mathsf B_S)_\kappa$ over the characters $\kappa$ of $H^1(S)$, where $\mathsf A(\mathsf B_S)_\kappa=\{X\in\mathsf A(\mathsf B_S):\mathsf V(c)X\mathsf V(c)^*=\kappa([c])X\text{ for all }c\in Z^1(S)\}$ is a character subspace. This is a linear grading; every summand $\ne0$, the summand $\kappa=1$ being $\mathsf A_\bullet(S)$; if every $\mathbb Z_2$-cycle of $\mathsf B_S$ is a sum of boundaries $\partial p$, $p\in S$, then $\mathsf A_\bullet(S)=\mathsf A(\mathsf B_S)$. $\mathsf A_\bullet(S)$ is the join of the $\mathsf A(\mathsf B_T)$ over the chains $T\subset S$ of three plaquettes, and also of the $\mathsf A(\partial p)$, $p\in S$. The traces and spin projections of one plaquette, or of two adjacent plaquettes, do not generate all gauge-invariant operators of their bonds; the traces alone generate an abelian algebra, while $\mathsf A(\mathsf B_S)$ is not abelian.
\item[(c)] \emph{The ring.} Let $a_i<b_i<b_i'<a_i'$ be integers with $a_i'-a_i\le N_i-1$, $Q=[a_1,a_1']\times[a_2,a_2']$, $D=[b_1,b_1']\times[b_2,b_2']$, $R$ the plaquettes of $Q$ off $D$, $C=\partial D$, $W_C=M_{\frac12{\rm tr}\,U(C)}$, $u=M_{{\rm sign}\,{\rm tr}\,U(C)}$, $\gamma$ a straight dual path from a plaquette of $D$ to a plaquette outside $Q$, $c_*=-1$ on the bonds of $\mathsf B_R$ that $\gamma$ crosses, $v=\mathsf V(c_*)$, and $D_1,\dots,D_4$ the four rectangles $[a_1,a_1']\times[a_2,b_2]$, $[b_1',a_1']\times[a_2,a_2']$, $[a_1,a_1']\times[b_2',a_2']$, $[a_1,b_1]\times[a_2,a_2']$. Then $R$ is admissible, $H^1(R)=\{1,[c_*]\}\cong\mathbb Z_2$, and on $\mathcal H_{\rm ph}$
\[
\begin{gathered}
\mathsf A(\mathsf B\setminus\mathsf B_R)'=\mathsf A(\mathsf B_R)=\mathsf A_\bullet(R)\vee\{W_C\}\\
=\mathsf A_\bullet(R)\oplus\mathsf A_\bullet(R)u,\qquad\\
\mathsf A_\bullet(R)'=\mathsf A(\mathsf B\setminus\mathsf B_R)\vee\{\mathsf V_\gamma\}\\
=\mathsf A(\mathsf B\setminus\mathsf B_R)\oplus\mathsf A(\mathsf B\setminus\mathsf B_R)\mathsf V_\gamma,
\end{gathered}
\]
with the normal conditional expectation $X\mapsto\frac12(X+vXv)$ of $\mathsf A(\mathsf B_R)$ onto $\mathsf A_\bullet(R)$; $vW_Cv=-W_C$, $W_C\notin\mathsf A_\bullet(R)$, $\mathsf V_\gamma\notin\mathsf A(\mathsf B\setminus\mathsf B_R)$, $\mathsf A(\mathsf B\setminus\mathsf B_R)=\mathsf A_\bullet(R)'\cap\{u\}'$, $u\mathsf V_\gamma=-\mathsf V_\gamma u$; each inclusion has the index two. $v$ is central in $\mathsf A_\bullet(R)$, and for a closed path $C'$ in $\mathsf B_R$, $M_{\chi_j(U(C'))}\in\mathsf A_\bullet(R)$ if and only if $j$ is an integer or $c_*(C')=1$: the adjoint loop around the hole lies in $\mathsf A_\bullet(R)$, the fundamental one does not. So over a ring duality and additivity exclude each other: $Y\mapsto\mathsf A(Y)$ is dual and not additive at $\mathsf B_R$, $S\mapsto\mathsf A_\bullet(S)$ is additive and not dual at $R$, and no assignment of von Neumann algebras that equals $\mathsf A$ on the closed and on the open rectangles is both.

At every spacing $\mathsf A_\bullet(R)=\bigvee_{i\le4}\mathsf A(\mathsf B_{D_i})$; if no corner $D_i\cap D_{i+1}$ is one plaquette, the multiplications by flip-invariant gauge-invariant functions of the bonds of $R$ are generated by those of the $D_i$, and $M_{\chi_1(U(C))}=\sum_aM_{(T_1)_{a_4a_1}}M_{(T_2)_{a_1a_2}}M_{(T_3)_{a_2a_3}}M_{(T_4)_{a_3a_4}}$ with gauge-invariant functions $(T_i)_{ab}$ of the bonds of $D_i$ of modulus $\le1$.
\item[(d)] \emph{Two discs; the whole layer.} Let $D_1,D_2$ be rectangles without a common site, with a row and a column of plaquettes that meet neither, $S$ the plaquettes off $D_1\cup D_2$, $\mathsf B^\circ_{D_i}$ the bonds inside $D_i$, $\gamma$ a dual path from $D_1$ to $D_2$ whose other plaquettes lie in $S$ off that row and column, $\mathsf V_1,\mathsf V_2$ the two fluxes of the torus, and $\mathsf N=\mathsf A(\mathsf B^\circ_{D_1})\vee\mathsf A(\mathsf B^\circ_{D_2})$. Then $S$ is admissible, $H^1(S)=\mathbb Z_2^3$, $\mathsf A(\mathsf B_S)'=\mathsf N$, $\mathsf A_\bullet(S)'=\mathsf N\vee\{\mathsf V_\gamma,\mathsf V_1,\mathsf V_2\}$, $\mathsf V_\gamma\notin\mathsf N\vee\{\mathsf V_1,\mathsf V_2\}$ and $\mathsf V_i\notin\mathsf N\vee\{\mathsf V_\gamma\}$. On the whole layer $\mathsf A_\bullet(\mathsf P)'={\rm span}\{1,\mathsf V_1,\mathsf V_2,\mathsf V_1\mathsf V_2\}$, $\mathcal H_{\rm ph}$ is the sum of four flux sectors, none $0$, and $\mathsf A_\bullet(\mathsf P)$ is the sum of all bounded operators of the sectors; one disc has no defect but the fluxes: $\mathsf A_\bullet(\mathsf P\setminus D)'=\mathsf A(\mathsf B^\circ_D)\vee\{\mathsf V_1,\mathsf V_2\}$.
\item[(e)] \emph{Three space dimensions.} On a finite part of $\mathbb Z^3$, or on a periodic one with sides $\ge4$, at one time: for all the plaquettes $S$ of a box $[0,n_1]\times[0,n_2]\times[0,n_3]$, $\bigvee_{p\in S}\mathsf A(\partial p)=\mathsf A(\mathsf B_S)$; for the plaquettes $S$ of the eight unit cubes of $[0,3]^2\times[0,1]$ off $(1,2)^2$, $\bigvee_{p\in S}\mathsf A(\partial p)=\mathsf A(\mathsf B_S)\cap\{\mathsf V(c_*)\}'$ and $(\bigvee_p\mathsf A(\partial p))'=\mathsf A(\mathsf B\setminus\mathsf B_S)\vee\{\mathsf V(c_*)\}$, $c_*=-1$ on the four bonds of $[1,2]\times[2,3]\times[0,1]$ parallel to $e_1$, with $\mathsf V(c_*)W_C=-W_C\mathsf V(c_*)$ for the fundamental loop $C$ around the ring; for a dual surface $\mathfrak s$ that meets the ring in that section, $\mathsf V(c_{\mathfrak s})=\mathsf V(c_*)\,\mathsf V(c_{\mathfrak s}|_{\mathsf B\setminus\mathsf B_S})$ with the second factor in $\mathsf A(\mathsf B\setminus\mathsf B_S)$: the defect $\mathbb Z_2$ is carried by the flip across a surface whose rim is a closed dual line in the complement that links the ring, a closed vortex line.

For a set $S$ of plaquettes let $Z^1(S)$, $B^1(S)$ be formed as on a layer, and call $H^1(S)=Z^1(S)/B^1(S)$ its defect; call an order $p_1,\dots,p_n$ of $S$ a shelling if for each $k\ge2$, with $X_{k-1}$ the set of the bonds of $p_1,\dots,p_{k-1}$, either all bonds of $p_k$ lie in $X_{k-1}$, or $p_k$ has exactly one bond in common with some $p_j$, $j<k$, the set $\partial p_j\cup(\partial p_k\cap X_{k-1})$ is connected, and every site of a bond of $p_k$ off $X_{k-1}$ that is a site of $X_{k-1}$ lies on that set. If $S$ has a shelling, $\bigvee_{p\in S}\mathsf A(\partial p)=\mathsf A(\mathsf B_S)$; a box has one. If $S=S_1\cup S_2$, both with shellings, $I=\mathsf B_{S_1}\cap\mathsf B_{S_2}$ has $r$ connected components, each with a cycle, and every site common to $\mathsf B_{S_1}$ and $\mathsf B_{S_2}$ is a site of $I$, then $\bigvee_{p\in S}\mathsf A(\partial p)=\mathsf A(\mathsf B_S)\cap\{\mathsf V(c):c\in Z^1(S)\}'$ and the defect is $\mathbb Z_2^{r-1}$; the solid ring above is such a union, with $r=2$.
\item[(f)] \emph{The sign rule.} On every periodic lattice in three dimensions and at every $\beta$, for sign fields $\eta$ on the plaquettes and $s$ on the bonds, $(ds)_p=\prod_{b\in\partial p}s_b$, bounded Borel $\Phi$, closed paths $C_i$ and spins $j_i$:
\[
\begin{gathered}
Z[\eta\,ds;\Phi;C_1,\dots,C_r]=\prod_is(C_i)^{2j_i}\,Z[\eta;\Phi;C_1,\dots,C_r],\qquad\\
 Z[\eta;\Phi;C]=\int\Phi((\eta_pt_p)_p)\prod_i\chi_{j_i}(U(C_i))\,e^{-\beta\sum_p(1-\eta_pt_p)}dU.
\end{gathered}
\]
Carrying a dual string once across a fundamental loop changes the sign, across an adjoint loop it does not; the rule holds for bounded fat odd loops $F$, $F(sU)=s(C)F(U)$ for $s$ closed on the plaquettes of a tube around $C$; and every limit point of renormalised mixed functions of loops, pairs of monopoles and twisted action densities has it.
\item[(g)] \emph{Block loops.} For the constants of Theorem~\YMMref{D}{8}, $L$ odd, a label $\mathsf b\ge\gamma^{-2}$, a torus of side $\ell=2L^m$ with $MR(\mathsf b)\le L^m$ ($R$ the function of Theorem~\YMMref{D}{8}), $\beta=L^K\mathsf b$ on $\Lambda_{K,\ell}=L^{-K}\mathbb Z^3/\ell\mathbb Z^3$, $\langle\cdot\rangle_K$ the Wilson measure there and $\mathcal M^{K-d}U$ the block field on $\Lambda_{d,\ell}$ with law $\mu^K_d$: the block map is covariant under the centre, $\mathcal M^n(sU)(c)=s[c]\,\mathcal M^n(U)(c)$ for every $U$ if $(ds)_p=1$ at every plaquette with all corners in the double block of $c$, $s[c]$ the product of $s$ along the straight path from $c_-$ to $c_+$; for every closed path $C$ of $\Lambda_{d,\ell}$ and bounded Borel $h$, $\langle h(W_C(\mathcal M^{K-d}U))\rangle_K$ converges along the full sequence of cutoffs, consecutive terms differing by $\le C\ell^3K^{\mathfrak p_0}L^{-K/10}\sup|h|$ for $K\ge K_{\rm bf}$ with $\lceil K/1260\rceil\ge d$, $K_{\rm bf}$ the threshold of the rate $\mathfrak o_K$ of Theorem~\YMMref{D}{8} (Corollary~\textup{\ref{IV-62-66}} among the construction results), and so do all mixed moments of loops of block fields; the flips $\bar\xi$, $d\bar\xi=1$, of the torus are symmetries of the limit laws $\mu_d$, so that odd and even functions are orthogonal in $L^2(\mu_d)$; and a closed vortex line that bounds passes as a flip: for a finite union $\mathsf S$ of closed dual plaquettes of $\Lambda_{d_0,\ell}$, $d\ge d_0$, $\eta_K$ the sign field that is $-1$ exactly on the plaquettes of $\Lambda_{K,\ell}$ pierced by $\partial\mathsf S$, and $\bar\xi_c=(-1)^{\#(c\cap\mathsf S)}$: $Z_K[\eta_K;F(\mathcal M^{K-d}\cdot)]/Z_K[1;1]=\int F(\bar\xi V)\,d\mu_d^K(V)\to\int F(\bar\xi V)\,d\mu_d(V)$ for bounded Borel $F$ of the block bonds whose blocks miss $\partial\mathsf S$, the equality holding at every $K$ and the error being $\le\sup|F|\sum_{K'\ge K}\mathfrak o_{K'}$ for $K\ge K_{\rm bf}$ and $\lceil K/1260\rceil\ge d$, $\mathfrak o_K$ the rate of Theorem~\YMMref{D}{8}; for $F=W_CF'$, with $F'$ a bounded Borel function of such block bonds and $F'(\bar\xi V)=F'(V)$, $W_C(\bar\xi V)=(-1)^{\#(C\cap\mathsf S)}W_C(V)$ gives the factor $(-1)^{\#(C\cap\mathsf S)}$, the linking number modulo $2$ of $C$ and $\partial\mathsf S$.
\item[(h)] \emph{Sharp loops; flips.} In (g) let $C$ be a lattice polygon without self-intersection that does not wind and has a neighbourhood in a cube of side $\le\ell/2$, $h$ bounded on $[-1,1]$ and continuous off a Lebesgue null set, $\bar h=\frac2\pi\int h(w)\sqrt{1-w^2}\,dw$, and $F_K$ Borel functions, $|F_K|\le1$, of the bonds of $\Lambda_{K,\ell}$ at a distance $\ge\mathsf d>0$ from $C$. Then $\langle F_Kh(W_C)\rangle_K-\bar h\langle F_K\rangle_K\to0$, $W_C=\frac12{\rm tr}\,U(C)$. For continuous $h$ this is Lemma~\ref{C-363-41}, Theorem~\ref{C-372-6} of Section~\ref{C-372-1} read at one polygon (see the results used): every character of non-zero spin of the holonomy of $C$, times a function of modulus $\le1$ of the bonds at a fixed distance from $C$, has a mean of modulus at most a constant times a positive power of $L^{-K}\log^3L^K$ \eqref{C-363-42}, through a rotation of the bonds of a tube of radius $\rho L^{-K}$ about $C$, $\rho$ of the order of $\beta^{1/4}$ up to logarithms, so that the class of the holonomy tends to the Haar law, whose image under $\frac12{\rm tr}$ has the density $\frac2\pi\sqrt{1-w^2}$; a bounded $h$ continuous off a null set lies between continuous functions $h_-\le h\le h_+$ with $\bar h_+-\bar h_-$ as small as we please.

So, with $C$ in a site mirror, the odd unitary $u_K$ of (c) and every character $\chi_j(U(C))$, $j\ge\frac12$, the adjoint one included, tend weakly to $0$ between the vectors of such functions in the space of reflection positivity, while $u_K^2=1$ and $\langle W_C^2\rangle_K\to\frac14$. For a rectangular path $\gamma$ whose sides have at least $3$ and at most $N-2$ bonds, $N=\ell L^K$ the number of sites on a side of the torus, $\gamma'$ its bonds but the two at each corner (no two of them lie in one plaquette, which a side of $N-1$ bonds would allow), $n'=\#\gamma'$, and bounded Borel $F$ of the other bonds: $|\langle\frac12{\rm tr}\,U(\gamma)F\rangle_K|\le\mathsf r(4\beta)^{n'}\sup|F|$, $\mathsf r(x)=I_2(x)/I_1(x)$ the ratio of modified Bessel functions; the bound tends to $\exp(-3|\gamma|/8\mathsf b)$.

The flips fade on bounded vectors: let $\omega_K$ be the state on the layer of $\Lambda_{K,\ell}$ given by the transfer operator of Section~\YMMref{85}{677} (with its trace formula, Proposition~\YMMref{85.2}{679}), $\zeta$ a dual path of the layer with an end plaquette $p_0$, $Q_0=M_{\mathbf 1\{a_{p_0}<1\}}$, and $X,Y$ bounded operators commuting with $Q_0$; then
$$|\omega_K(X\mathsf V_\zeta Y)|\le2\|X\|\|Y\|\langle a_p\rangle_K^{1/2},\qquad\langle a_p\rangle_K\le(21+3\log\beta)/\beta\quad(\beta\ge1).$$
\item[(i)] \emph{The Euclidean $\mathbb Z_2$.} In (g) let $\mathcal R$ be a solid ring, a union of closed cubes of side $L^{-d}$ centred at sites of $\Lambda_{d,\ell}$, and $\bar\xi$ a flip, the sign change of the block variables crossing a dual surface $\mathsf S$ whose rim misses $\mathcal R$; $\mathsf h_D$ is Haar measure on the block fields at the depth $D$ and $\mathsf{In}_{[D]}=\{V:\max_p|\partial V(p)-\mathbf 1|<2\varepsilon(x_{[D]})\}$, $|\cdot|$ the operator norm, the set of the small fields of Theorem~\YMMref{D}{8} at the depth $D$, its inner ball (the common-depth definitions in Section~\YMMref{116}{807}, the definitions preceding Lemma~\YMMref{113.10}{796}), with $x_{[D]}=L^D\mathsf b$ (Section~\ref{sec:block-field-laws}) and $\varepsilon(\mathsf t)=A_0\mathsf t^{-1/2}(\log\mathsf t)^{\mathfrak p_0}$ \YMMref{(634)}{339}, $A_0$ a constant of the list of Theorem~\YMMref{D}{8}.

There are $D>d$ and numbers $\kappa_D,\bar\kappa_D>0$, $\mathsf R_{[D]}\to0$, free of $K$, with $\mu_D^K(E)\ge\kappa_D\,\mathsf h_D(E\cap\mathsf{In}_{[D]})$ and $\mu_D^K(E)\le\bar\kappa_D\,\mathsf h_D(E\cap\mathsf{In}_{[D]})+\mathsf R_{[D]}$ for all Borel $E$ and $\mu^K_D(\mathsf{In}_{[D]})\ge1-\mathsf R_{[D]}$, for all $K\ge D+1$ and for the limit law $\mu_D$. A thick vortex is a small field: for every $D\ge d+2$ there is a block field $V^1$ at the depth $D$ with values in an abelian subgroup, equal to $\bar\xi_c$ on the bonds of $\mathcal R$, whose plaquette angles are $\le16\pi n_{\mathsf S}L^{-2(D-d)}$, $n_{\mathsf S}$ the number of dual bonds in $\partial\mathsf S$. Hence one positive measure $\theta_0$ of mass $\varkappa_D>0$, which depends on no cutoff and no surface, lies below the law of the block fields of $\mathcal R$ and below all its flips, for all $K\ge D+1$ and in the limit; so the law and each flip differ by at most $1-\varkappa_D$ in total variation, every odd function of the block fields of $\mathcal R$ with values in $[-1,1]$ has its mean in $[\varkappa_D-1,1-\varkappa_D]$, and for a closed path $C$ in $\mathcal R$ with $\#(C\cap\mathsf S)$ odd the fundamental block loop $W_C$ and its sign keep an $L^1$-distance $\ge\mathsf m_C=\int|W_C|\,d\theta_0>0$ from every flip-invariant function, at every $K\ge D+1$ and in the limit.

The limit laws have densities for Haar measure by Lemma~\ref{C-363-26}. Conversely, let $\mathcal R$ be a rectangular solid ring of $\Lambda_{d,\ell}$: with $N=\ell L^d$ sites on a side of the torus, integers $a_i<b_i<b_i'<a_i'$ $(i=1,2)$ and $t_1\le t_2$, in the units of $\Lambda_{d,\ell}$, $\mathcal R$ is the union of the four bricks $[a_1,a_1']\times[a_2,b_2]\times[t_1,t_2]$, $[b_1',a_1']\times[a_2,a_2']\times[t_1,t_2]$, $[a_1,a_1']\times[b_2',a_2']\times[t_1,t_2]$ and $[a_1,b_1]\times[a_2,a_2']\times[t_1,t_2]$ of sites, where $t_2>t_1$ or no corner of two consecutive bricks is one plaquette, and where $b_i'-b_i\ge2$, $a_i'-a_i\le N-2$ and $t_2-t_1\le N-2$, so that the hole holds a site and the four bricks hold all the bonds between the sites of $\mathcal R$; let $\alpha$ be its flip, the sign change of the block variables of the bonds $\langle x,x+e_1\rangle$ of $\mathcal R$ with $x_1=b_1$ and $x_2\le b_2$, which cross one section of the ring, $\mu=\rho\,dV$ the limit law of its block fields, with a gauge-invariant density $\rho$, $u={\rm sign}\,W_C$ for the loop $C$ around the hole, $\mathfrak L$ the classes of all bounded gauge-invariant Borel functions of these block fields and $\mathfrak E$ those of the flip-invariant ones.

Then every $F\in\mathfrak L$ is $F=\Phi_0+\Phi_1u$ $\mu$-almost everywhere with $\Phi_0,\Phi_1\in\mathfrak E$ measurable for the four bricks (for the cubes of three sites if the bricks are three sites thick), $\sup|\Phi_i|\le\sup|F|$, and the same holds for the block fields of all coarser depths in $\mathcal R$; $\mathfrak E=\bigvee_i\mathfrak L({\rm brick}_i)$ and $\mathfrak L=\mathfrak E+\mathfrak Eu$; the adjoint loop is $\chi_1(V(C))={\rm tr}(T_1T_2T_3T_4)$ with entries $(T_i)_{ab}$ functions of the $i$-th brick of modulus $\le1$. On the closed subspace $\mathfrak h$ of $L^2(\mu)$ of the classes of the gauge-invariant functions, on which $\mathfrak L$ and $\mathfrak E$ act by multiplication, the commutant of $\mathfrak E$ is $\mathfrak L+\mathfrak Lv$, of type ${\rm I}_1\oplus{\rm I}_2$, where $(vF)(V)=e(V)(\rho(\alpha V)/\rho(V))^{1/2}F(\alpha V)$ and $e$ is the indicator of $\{\rho>0\}\cap\alpha\{\rho>0\}$: $v=v^*$, $v^2=e$, $vu=-uv$, and $\int e\,d\mu\ge\varkappa_D$, so that $\mu$ and its image under $\alpha$ are not mutually singular and $v\ne0$: $L$ being odd, $\mathcal R$ is also a union of closed cubes of side $L^{-(d+1)}$ centred at sites of $\Lambda_{d+1,\ell}$, and $\alpha$ agrees on the bonds of $\mathcal R$, up to a central gauge transformation, with the sign field into which the block map from the depth $d+1$ to the depth $d$ turns the flip of a flat sheet $\mathsf S_1$ of $\Lambda_{d+1,\ell}$ between two straight closed dual lines in the direction $e_3$ that lie over the open unit squares about a site of the hole and about a site of the same row outside $[a_1,a_1']\times[a_2,a_2']$, so that its rim misses $\mathcal R$; here $D$ and $\varkappa_D$ are those of the first half at the depth $d+1$.

And $v$ is the limit of the self-adjoint unitaries $v_K$ of the closed vortex line at the cutoffs: for all bounded Borel functions $F,F'$ of the block fields of $\mathcal R$, with $\mu^K$ their law at the cutoff $K$, $|\langle F,v_KF'\rangle_{L^2(\mu^K)}-\langle F,vF'\rangle_{L^2(\mu)}|\le2\sqrt2\sup|F|\sup|F'|(\sum_{K'\ge K}\mathfrak o_{K'})^{1/2}$ for $K\ge K_{\rm bf}$ and $\lceil K/1260\rceil\ge d$. If besides the hole is at least three bonds wide in both directions, the ring has positive height and it lies in an open ball of radius $<\ell/4$, then the law and its flip have the same null sets, $e=1$, $v$ is a self-adjoint unitary, and $\mathfrak L=\mathfrak E\oplus\mathfrak Eu$ is a grading by $\mathbb Z_2$.

For a band $\mathbb Z/N\times I_2\times I_3$ that winds around the torus, $I_2,I_3$ intervals of $3$ to $N-1$ sites, the defect is exactly $\mathbb Z_2$: $\mathfrak L=\mathfrak E\oplus\mathfrak Eu_1$, orthogonally, $\mathfrak Eu_1\ne0$, $u_1$ the sign of the straight loop that winds, $\mathfrak E$ is generated by the cubes of the band, and the flip is a unitary $v$ with $\mathfrak E'=\mathfrak L\vee\{v\}$ on the space of the gauge-invariant classes for the law of the block fields of the band.
\end{enumerate}
\end{theorem}

\subsection{Construction estimates}

We use the reconstruction from data (Rec) of Section~\YMMref{107}{765} and Sections~\YMMref{109}{772} and~\YMMref{110}{776}, as recalled in the setting: the lemma on the smeared insertion (Lemma~\YMMref{107.1}{766}) with the insertion formula \YMMref{(2039)}{767}, the point bound \eqref{C-363-2} (Proposition~\YMMref{107.2}{767} and Lemma~\YMMref{109.2}{773}), the uniqueness of the family of point operators (Proposition~\YMMref{107.2}{767}), the lemma on the energy bound (Lemma~\YMMref{107.3}{768}) and the definition of the energy form $q_0$ (Lemma~\YMMref{107.3}{768}), the field operators of Proposition~\YMMref{109.1}{773}, the boundary distributions and polynomial cyclicity proved in Section~\YMMref{109}{772}, and the spectrum condition in Theorem~\YMMref{AO}{778}, the lemma (Br) that links the field operators with the point vectors (Lemma~\YMMref{109.2}{773}), the propositions on Lorentz covariance and on locality (Propositions~\YMMref{110.1}{776} and~\YMMref{110.2}{777}), and the domain argument preceding Proposition~\YMMref{109.1}{773}; Lemma~\ref{lem:setup-real-identity}, by which a holomorphic function on a connected open subset of $\mathbb C^N$ that vanishes on a non-empty open set of real points vanishes; Proposition~\YMMref{97.2}{731}, which turns an insertion bound for sources in small balls with a constant $C_{\rm I}$ into the point bound of order $4$ with the constant $C_{\rm pt}=c_{\rm pt}C_{\rm I}$, $c_{\rm pt}=320\cdot8^4\,\mathsf{mv}_3(40L)^4M$ the coefficient of $C_{\rm I}$ in~\YMMref{(1919)}{729} at $c_{\rm I}=\frac1{40L}$, with $\mathsf{mv}_3=m_3$ as defined immediately after that equation, and the insertion bound for sources in small balls of one cell, Proposition~\YMMref{105.1}{762}, with its proof, which the proof of Theorem~\ref{C-363-5}(c) uses up to its last step together with the insertion bound on one lattice for a pair that carries systems (Lemma~\YMMref{104.4}{760}, with its normalizer~\YMMref{(2017)}{760}), the systems of localised source energies of a pair and their constants (Lemma~\ref{lem:canonical-source-system}, with \eqref{eq:setup-pressure-constants}), the window \YMMref{(1531)}{615} of the labels (Proposition~\YMMref{74.1}{615}), the two source-radius conditions of Proposition~\YMMref{82.2}{668}, at the same radius-independent numerical prefix, and the dilation among the construction results (Lemma~\ref{lem:lattice-law-dilation}); and Theorem~\textup{\YMMref{AP}{780}} among the construction results on the limits of periodic Gibbs states and their canonical decompositions, which follows from Theorem~\YMMref{D}{8}, at the constants of Proposition~\YMMref{69.1}{595} and Theorems~\YMMref{AJ}{673} and~\YMMref{AD}{607}.

By that theorem every limit of periodic Gibbs states at $g$, volume-first ones included, is data (Rec) with $(16C^g_{\rm pt},4,\frac12)$ and has a canonical decomposition, every component of which is data (Rec) with $(16C^g_{\rm pt}+\bar{\mathfrak c}(\mathsf C^g),4,\frac12)$ and a unique vacuum and is approximated by lattice states: if $\sigma'(X)=\lim_j\rho_j((X)^{\rho_j}_{K_j})$, every component is $\sigma'_\xi(X)=\lim_j\rho_{j,\xi}((X)^{\rho_j}_{K_j})$ along the full sequence, with the periodic Gibbs states $\rho_{j,\xi}(F)=\rho_j(\mathbf 1_{E_j}F)/\rho_j(E_j)$ at the same couplings, $E_j$ events with $\rho_j(E_j)>0$ (Theorem~\textup{\YMMref{AP}{780}}); these are the lattice approximants of the component; the numbers depend on no sequence, limit, decomposition or $\xi$; and a volume-first $\sigma$ has $S^\sigma_1=0$, its truncated functions $S^\sigma_n$, $n\ge2$, being the limits of the joint cumulants $S^T_{n;j}$ of the $\mathcal O_{K_j}(f)$ under the Wilson measures $\mu_j$ on tori $\Lambda_j=L^{-K_j}\mathbb Z^3/\ell_j\mathbb Z^3$, $\ell_j=2L^{m_j}\to\infty$.

We use besides the reflection positivity of periodic Gibbs states at the nine families of lattice mirrors (Proposition~\textup{\YMMref{103.4}{758}}), the chessboard estimate (Proposition~\YMMref{86.3}{681}, with equation~\YMMref{(1759)}{682}) and the multiple reflections (Proposition~\YMMref{104.2}{759}), the one-cell bound with a source (the cell bound, Proposition~\YMMref{89.1}{690}), the systems of localised source energies (Section~\YMMref{88}{685}) and the one-lattice source recursion of Lemma~\ref{C-363-23}, with its identification of the source energies and its proof on smaller cores, together with the modification of multiple reflections stated and proved in Lemma~\ref{C-363-16}(e), and, for Theorem~\ref{C-363-5}(c), Theorems~\YMMref{AJ}{673} and~\YMMref{AD}{607} at every radius and every admitted choice of the constants $(\sigma^-,\gamma)$; from Theorem~\YMMref{D}{8}, that on one torus the laws of the block fields at a fixed depth converge, at a rate (Corollary~\textup{\ref{IV-62-66}}), with a floor, an upper bound and a tail on the small fields, and a tail at a fraction of the threshold, read off the small-field chain (the setting of the chain, Section~\ref{sec:cutoff-calibration}; Proposition~\textup{\YMMref{115.2}{806}} with the definition preceding Proposition~\YMMref{115.2}{806}, \YMMref{(2187)}{806} and \YMMref{(2188)--(2190)}{806}; Propositions~\YMMref{114.2}{801} and~\YMMref{113.13}{799}; Lemma~\YMMref{113.10}{796} with the definitions preceding Lemma~\YMMref{113.10}{796}, \YMMref{(2159)}{796} and \YMMref{(2161)}{796}; equation~\YMMref{(2180)}{803} with \YMMref{(2179)}{803}; Theorem~\YMMref{E}{10} and the curvature comparison in the proof of Lemma~\YMMref{113.10}{796}; the condition $A_0\ge4B_3C_HA_1$ in the ordered choice, written explicitly in the proof of Lemma~\ref{lem:rarity-lowered-choice}; and \YMMref{(634)}{339}, (Section~\ref{sec:block-field-laws}), the common-depth definitions in Section~\YMMref{116}{807}); the block-averaging map (equations~\YMMref{(37)--(39)}{20}) with its locality, its real-analyticity on the small fields and its linear part at the unit (Proposition~\YMMref{4.5}{20}, \YMMref{4.5}{20} and~\YMMref{4.6}{22}); the transfer operator of a periodic lattice in one direction and its trace formula (Section~\YMMref{85}{677}, Proposition~\YMMref{85.2}{679}); and the short-distance law of the two-point function recalled in the setting (Theorem~\textup{\YMMref{AV}{851}}, Corollary~\textup{\YMMref{127.11}{854}}).

The theorem on sharp loops on tori on which Theorem~\ref{C-363-7}(h) rests is Theorem~\ref{C-372-6} of Section~\ref{C-372-1}, in its limit, as the spacing tends to $0$, for polygons with a further factor $G$, read at one polygon with $G=F_K$, together with the bound for the characters behind that limit; it is stated with a sketch of its proof as Lemma~\ref{C-363-41} below: at every label $\mathsf b>0$, with $\beta=L^K\mathsf b$ and $\langle\cdot\rangle_K$ the Wilson measure on $\Lambda_{K,\ell}$, for a lattice polygon $C$ without self-intersection that does not wind and has a neighbourhood in a cube of side $\le\ell/2$, and continuous $h$ on $[-1,1]$, $\langle F_Kh(\frac12{\rm tr}\,U(C))\rangle_K-\bar h\langle F_K\rangle_K\to0$ uniformly in Borel functions $|F_K|\le1$ of the bonds at a fixed positive distance from $C$, $\bar h=\frac2\pi\int h(w)\sqrt{1-w^2}\,dw$; its proof reads the Wilson weight on one torus, the Schwarz inequality at site mirrors and Haar measure alone.

The textbook theorems used are the spectral theorem, Stone's theorem on strongly continuous unitary groups, polar decomposition, the bicommutant theorem, von Neumann's theorem on $1+T^*T$, modular theory, reduction theory (direct integrals of separable Hilbert spaces and of von Neumann algebras over a compact metrisable space with a Borel measure: the decomposable operators are the commutant of the diagonal algebra, a von Neumann algebra between the diagonal algebra and its commutant is the direct integral of a measurable field of von Neumann algebras, whose fibres are fixed up to a null set by a countable generating family, and the commutant of a direct integral is the direct integral of the commutants), H\"older's inequality in Schatten classes, Jensen's inequality, Gronwall's lemma, the three-lines theorem, the theorems of Stone--Weierstrass, Montel, Vitali and Liouville, the extension theorem of M.~Riesz, the theorems of Peter--Weyl, Clebsch--Gordan and Burnside, the theorem on quasi-analytic vectors, the duality theorem for the free field, the first fundamental theorem for orthogonal groups, Fubini's theorem, the implicit function theorem, and the fact that the zeros of a real-analytic function that does not vanish identically on a connected open set form a null set.

\subsection{Auxiliary estimates}

\begin{lemma}[Frames]\label{C-363-8}
Let $\iota:\mathcal V\to\mathcal H$ be linear with dense range in a separable Hilbert space, and $\iota_j:\mathcal V\to\mathcal H_j$ linear, $\iota_jG$ defined for $j\ge j(G)$, with $\langle\iota_jG',\iota_jG\rangle\to\langle\iota G',\iota G\rangle$. Write $\psi_j\Rightarrow\psi$ if for every $\delta>0$ some $G$ has $\|\psi-\iota G\|<\delta$ and $\limsup_j\|\psi_j-\iota_jG\|<\delta$; $\psi_j\rightharpoonup\psi$ if $\sup_j\|\psi_j\|<\infty$ and $\langle\iota_jG,\psi_j\rangle\to\langle\iota G,\psi\rangle$ for all $G$; and, for bounded operators, $Q_j\rightharpoonup Q$ if $\sup_j\|Q_j\|<\infty$ and $\langle\iota_jG',Q_j\iota_jG\rangle\to\langle\iota G',Q\iota G\rangle$ for all $G,G'$. Then: (i) $\iota_jG\Rightarrow\iota G$; $\psi_j\Rightarrow\psi$ and $\psi'_j\rightharpoonup\psi'$ give $\langle\psi_j,\psi'_j\rangle\to\langle\psi,\psi'\rangle$; $\Rightarrow$ implies $\rightharpoonup$ and $\|\psi_j\|\to\|\psi\|$; and $\psi_j\rightharpoonup\psi$ gives $\|\psi\|\le\liminf_j\|\psi_j\|$. (ii) Bounded sequences of vectors and of operators have subsequences that converge in the sense of $\rightharpoonup$, the limits with $\|Q\|\le\liminf_j\|Q_j\|$. (iii) $Q_j\rightharpoonup Q$ and $\psi_j\Rightarrow\psi$ give $Q_j\psi_j\rightharpoonup Q\psi$, and $Q_j^*\rightharpoonup Q^*$. (iv) If $\pi_j$ are projections with $\|(1-\pi_j)\iota_jG\|\to0$ for all $G$, then $(\pi_j\mathcal H_j,\pi_j\iota_j)$ has the same properties, and for $\psi_j=\pi_j\psi_j$ the two meanings of $\Rightarrow$ and of $\rightharpoonup$ agree. (v) Let $\mathsf P_j$ and $\mathsf A$ be self-adjoint, $\mathcal V_0\subset\mathcal V$ a subspace, $\mathsf r:\mathcal V_0\to\mathcal V$ linear, with $\iota\mathcal V_0\subset D(\mathsf A)$, $(\mathsf A\pm i)\iota\mathcal V_0$ dense, $\mathsf A\iota G=\iota\mathsf rG$ and $\mathsf P_j\iota_jG=\iota_j\mathsf rG$ for $G\in\mathcal V_0$ and large $j$.

Then $\Phi(\mathsf P_j)\psi_j\Rightarrow\Phi(\mathsf A)\psi$ for every $\Phi\in C_0(\mathbb R)$ and all $\psi_j\Rightarrow\psi$.
\end{lemma}
\begin{proof}[Proof sketch]
(i) The triangle inequality and the convergence of the Gram matrices, as in $|\langle\psi_j,\psi'_j\rangle-\langle\psi,\psi'\rangle|\le\|\psi_j-\iota_jG\|\sup_j\|\psi'_j\|+|\langle\iota_jG,\psi'_j\rangle-\langle\iota G,\psi'\rangle|+\|\iota G-\psi\|\|\psi'\|$. (ii) On a countable $\mathbb Q(i)$-subspace with dense image the forms converge along a diagonal subsequence; their limits are bounded by $\sup_j\|Q_j\|\,\|\iota G'\|\|\iota G\|$, Riesz's theorem gives $Q$, and the convergence of the Gram matrices passes to all of $\mathcal V$. (iii), (iv) The triangle inequality. (v) Call a bounded sequence $(\mathsf X_j)$ with a bounded $\mathsf X$ good if $\psi_j\Rightarrow\psi$ implies $\mathsf X_j\psi_j\Rightarrow\mathsf X\psi$; sums, products and limits uniform in $j$ of good sequences are good. Given $\delta$, choose $G\in\mathcal V_0$ with $\|\psi-(\mathsf A\pm i)\iota G\|<\delta$; as $(\mathsf P_j\pm i)\iota_jG=\iota_j(\mathsf rG\pm iG)$,
$$\limsup_j\|(\mathsf P_j\pm i)^{-1}\psi_j-\iota_jG\|\le\limsup_j\|\psi_j-\iota_j(\mathsf rG\pm iG)\|<\delta ,$$
so the resolvents are good; the polynomials without constant term in $(x+i)^{-1},(x-i)^{-1}$ are dense in $C_0(\mathbb R)$, and $\|\Phi(\mathsf P_j)\|\le\sup|\Phi|$.
\end{proof}

\begin{lemma}[Rotations and boosts of point vectors]\label{C-363-9}
Let $\sigma$ be data (Rec), $\tilde B(\zeta)=\iota B(\zeta)\iota^{-1}$ on $\mathbb C^3$, applied to each point of a $z\in(\mathbb C^3)^n$, so that $\tilde B(i\tau)=R_\tau$; for $0<\delta\le\frac\pi2$ let $\mathsf I_\delta=\{x_1>0,\ 0<x_3<x_1\tan\delta\}$, $\mathsf I_{\pi/2}=\mathsf I$, and
$$\mathsf I_\delta^{(n)}=\{x\in(\mathbb R^3)^n:\ x_1\in\mathsf I_\delta,\ x_{k+1}-x_k\in\mathsf I_\delta\ (1\le k<n)\},$$
an open convex cone.
(a) $\tilde B(s)\mathsf{Tb}^{(n)}=\mathsf{Tb}^{(n)}$ and $U(B(s))\Psi_n(z)=\Psi_n(\tilde B(s)z)$ for real $s$; if $R_\tau z\in\mathsf{Tb}^{(n)}$ for $0\le\tau\le\tau_0$, then $\Psi_n(z)\in D(e^{-\tau_0K})$ and $e^{-\tau K}\Psi_n(z)=\Psi_n(R_\tau z)$ for $0\le\tau\le\tau_0$; and for $x\in\mathsf I_\delta^{(n)}$, $R_\tau x\in\mathsf{Tb}^{(n)}$ for $0\le\tau\le\pi-\delta$.
(b) A small tuple in $\mathsf I_\delta$ is a tuple $G=(h_1,\dots,h_n)$ of real $h_k\in C^1_c(\mathbb B(x^\circ_k,\epsilon))$, $\mathbb B(y,\epsilon)$ the open ball of $\mathbb R^3$, with centres $x^\circ=(x^\circ_k)_k$ such that $\prod_k\mathbb B(x^\circ_k,\epsilon)\subset\mathsf I_\delta^{(n)}$ and $\epsilon<c_0\eta_{x^\circ}$, where
$$\eta_{x^\circ}=\min\Bigl\{1,\ \min_{0\le\tau\le\pi-\delta}(R_\tau x^\circ_1)_3,\ \tfrac12\min_{k<n}\ \min_{0\le\tau\le\pi-\delta}\bigl(R_\tau(x^\circ_{k+1}-x^\circ_k)\bigr)_3\Bigr\}>0 ;$$
For $n=1$ the member involving $\min_{k<n}$ is omitted from this minimum; the small tuples in $\mathsf I$ are those in $\mathsf I_{\pi/2}$. For a small tuple $G$ in $\mathsf I_\delta$ and $0\le\tau\le\pi-\delta$: $[G]\in D(e^{-\tau K})$ and $e^{-\tau K}[G]=[R_\tau G]$; in particular $\Delta^{1/4}[G]=[RG]$ for the small tuples in $\mathsf I$. The span $\mathcal V_0$ of $\varnothing$ and of the small tuples in $\mathsf I$ gives a core $[\mathcal V_0]$ of $\Delta^{1/4}$, with $(\Delta^{1/4}\pm i)[\mathcal V_0]$ dense.
(c) For real $u\in C^\infty_c(\mathsf I)$ let $f_u(B(s)(0,\rho,x_2))=\int_0^{\pi/2}u(\rho\cos\alpha,x_2,\rho\sin\alpha)\,e^{-(s+i(\pi/4-\alpha))^2}d\alpha$ on $W$, and $f_u=0$ off $W$. Then $f_u\in\mathcal S$, ${\rm supp}f_u\subset W\cap\{r_1\le((x^1)^2-(x^0)^2)^{1/2}\le r_2\}$ for some $0<r_1<r_2$,
\begin{equation}\label{C-363-10}
\mathsf F(K)[u]=\phi(f_u)\Omega,\qquad\mathsf F(k)=\sqrt\pi\,e^{-k^2/4+\pi k/4},
\end{equation}
and for every $p\in\mathbb R^3$ some $u$ has $\hat f_u(p)\ne0$.
\end{lemma}
\begin{proof}[Proof sketch]
(a) The analytic Wightman functions are invariant under the complexified Lorentz transformations that keep the tube, and the inner products $\langle\Psi_m(w),\Psi_n(z)\rangle$ are such functions; so $\Omega\mapsto\Omega$, $\Psi_n(z)\mapsto\Psi_n(\tilde B(s)z)$ extends to an isometry. By (Br) at $\mathfrak w_k=\eta e_3$, then at $\mathfrak w_k=\eta\tilde B(s)e_3$, and $\eta\downarrow0$, it carries $\phi(f_1)\cdots\phi(f_n)\Omega$ to $\phi(f_1\circ B(-s))\cdots\phi(f_n\circ B(-s))\Omega=U(B(s))\phi(f_1)\cdots\phi(f_n)\Omega$; $\mathcal D_0$ is dense, so it is $U(B(s))$. The set of $\zeta$ with $\tilde B(\zeta)z\in\mathsf{Tb}^{(n)}$ is invariant under real shifts and holds $i[0,\tau_0]$, so it holds a strip; there $\zeta\mapsto\Psi_n(\tilde B(\zeta)z)$ is holomorphic in norm and equals $e^{i\zeta K}\Psi_n(z)$ on $\mathbb R$; comparison with $\int_{[-\mathsf M,\mathsf M]}e^{i\zeta k}dm^K_{\chi,\Psi_n(z)}$, where $m^K_{\chi,\psi}=\langle\chi,\mathcal P_K(\cdot)\psi\rangle$, $\mathcal P_K$ the spectral measure of $K$, and $m^K_\psi=m^K_{\psi,\psi}$, on bounded parts of the spectrum of $K$ at $\zeta=i\tau$ gives $\int e^{-2\tau k}dm^K_{\Psi_n(z)}\le\|\Psi_n(R_\tau z)\|^2$ and the identity.

A point $y\in\mathsf I_\delta$ is $(\rho\cos\alpha,y_2,\rho\sin\alpha)$ with $\rho>0$, $0<\alpha<\delta$, and $(R_\tau y)_3=\rho\sin(\alpha+\tau)>0$ for $0\le\tau\le\pi-\delta$; at real points $d(x)=x_3$, and $R_\tau$ is real and linear; so for $x\in\mathsf I_\delta^{(n)}$ the points $y=x_1$ and $y=x_{k+1}-x_k$ give $R_\tau x\in\mathsf{Tb}^{(n)}$. (b) Let $G$ be a small tuple in $\mathsf I_\delta$ with the centres $x^\circ$, and $0\le\tau\le\pi-\delta$. The balls $\mathbb B(R_\tau x^\circ_k,\epsilon)$ that carry the entries of $R_\tau G$ have a first height $\ge\eta_{x^\circ}$ and heights that grow by $\ge2\eta_{x^\circ}$, and $\epsilon<c_0\eta_{x^\circ}$; so $[R_\tau G]=\int\prod_kh_k(x_k)\,\Psi_n(R_\tau x)\,dx$, and on $\prod_k\mathbb B(x^\circ_k,\epsilon)\subset\mathsf I_\delta^{(n)}$, $\Psi_n(R_\tau x)=e^{-\tau K}\Psi_n(x)$ by (a); both integrands are continuous in norm and $e^{-\tau K}$ is closed, so $[G]\in D(e^{-\tau K})$ and $e^{-\tau K}[G]=[R_\tau G]$; at $\delta=\tau=\frac\pi2$ this is $\Delta^{1/4}[G]=[RG]$; and $\Delta^{1/4}\Omega=\Omega$.

Let $\chi\perp(\Delta^{1/4}+1)[\mathcal V_0]$; then $\chi\perp\Omega$. Let $0<\delta<\frac\pi2$ and $G'$ a small tuple in $\mathsf I_\delta$ with the centres $x^\circ$. For $0\le\beta\le\frac\pi2-\delta$ the tuple $R_\beta G'$ is a small tuple in $\mathsf I$ with the centres $R_\beta x^\circ$: $R_\beta\mathsf I_\delta\subset\mathsf I$, so that $R_\beta\mathsf I_\delta^{(n)}\subset\mathsf I_{\pi/2}^{(n)}$, and the heights of $R_\tau R_\beta x^\circ$, $0\le\tau\le\frac\pi2$, are those of $R_{\beta+\tau}x^\circ$ with $\beta+\tau\in[0,\pi-\delta]$, so that $\eta_{R_\beta x^\circ}\ge\eta_{x^\circ}>\epsilon/c_0$. The function $\mathsf f(z)=\int e^{-zk}dm^K_{\chi,[G']}(k)$ is continuous on $0\le{\rm Re}\,z\le\pi-\delta$, as $\int e^{-{\rm Re}\,z\,k}d|m^K_{\chi,[G']}|\le\|\chi\|(\|[G']\|+\|e^{-(\pi-\delta)K}[G']\|)$, holomorphic inside, and $\mathsf f(\beta)=\langle\chi,e^{-\beta K}[G']\rangle=\langle\chi,[R_\beta G']\rangle$.

As $\chi\perp(\Delta^{1/4}+1)[R_\beta G']=[R_\beta G']+[R_{\beta+\pi/2}G']$, $\mathsf f(\beta)+\mathsf f(\beta+\frac\pi2)=0$ for $0\le\beta\le\frac\pi2-\delta$, hence on the strip $0<{\rm Re}\,z<\frac\pi2-\delta$; at $z=\beta_0+it$ the finite measure $e^{-\beta_0k}(1+e^{-\pi k/2})m^K_{\chi,[G']}$ has Fourier transform $0$, so it vanishes, and $\chi\perp[G']$. Small tuples in $\mathsf I_\delta$, with arbitrary entries, are had about every point of the open set $\mathsf I_\delta^{(n)}$, and $\langle\chi,\Psi_n(\cdot)\rangle$ is continuous; so it vanishes on $\mathsf I_\delta^{(n)}$, an open set of real points of the tube, hence on the tube (Lemma~\ref{lem:setup-real-identity}), and $\chi\perp\mathcal D_0$ by (Br): $\chi=0$.

The closure of $\Delta^{1/4}$ on $[\mathcal V_0]$ is symmetric, $\ge0$, lies in $\Delta^{1/4}$, and the range $(\Delta^{1/4}+1)[\mathcal V_0]$ is dense; so it is self-adjoint, and it is $\Delta^{1/4}$. (c) $f_u$ is smooth with all Schwartz seminorms finite because $|\partial^a_sg(s+i\tau)|\le c_a(1+|s|)^ae^{-s^2+\tau^2}$, $g(s)=e^{-s^2}$. As $\int g(s+i\beta')e^{isk}ds=\sqrt\pi e^{-k^2/4+\beta'k}$, (a) used backwards gives for $y=(0,\rho,x_2)$ and $0<\gamma\le\alpha$
$$\mathsf F(K)\Psi_1(\iota B(i\alpha)y)=\int g\bigl(s+i(\gamma+\tfrac\pi4-\alpha)\bigr)\,\Psi_1(\iota B(s+i\gamma)y)\,ds ;$$
integrate against $u$ in polar coordinates of the $(x_1,x_3)$-plane, pair with $e^{-\eta H}\chi$, and let $\gamma\downarrow0$ by dominated convergence, the norms being $\le C_\wp\min\{1,\eta\}^{-\wp}$ because $d(\iota B(i\gamma)x+\eta e_3)\ge\eta$ on $W$; (Br) at $n=1$ gives $\langle e^{-\eta H}\chi,\mathsf F(K)[u]\rangle=\langle e^{-\eta H}\chi,\phi(f_u)\Omega\rangle$, and $e^{-\eta H}\mathcal H$ is dense. If $\hat f_u(p)=0$ for all $u$, the entire function $\beta\mapsto\int g(s+i\beta)e^{ip\cdot B(s)y}ds$ would vanish, so $g*e^{ip\cdot B(\cdot)y}=0$, which is absurd as $\hat g$ has no zero.
\end{proof}

\begin{proposition}[Operators between two Euclidean times]\label{C-363-18}
Let $\sigma$ be data (Rec) with a unique vacuum. (a) If $O\ne\emptyset$ is open, $Q\in\mathcal C(O)$, $s>0$ and $Q\Omega,Q^*\Omega\in D(e^{sH})$, then $Q\in\mathbb C1$. So a bounded operator $e^{-sH}V_1e^{-sH}$, $V_1$ bounded and $s>0$, lies in no $\mathcal C(O)$ unless it is $0$ or $H$ is bounded, that is, unless $\mathcal H=\mathbb C\Omega$; this covers $\mathsf B(s,s')$, every product $\mathcal T(\mathfrak z)\mathsf B(s,s')\mathcal T(\mathfrak z')\cdots$, and every weak limit $V$, along a subsequence and in the sense of Lemma~\ref{C-363-8}, of lattice operators $V_j:[G]\mapsto[D_j\cdot\tau_{te_3}G]$ on the spaces of reflection positivity at $x_3=0$ of lattice states $\rho_j$ that converge to $\sigma$ and are reflection positive at the site planes orthogonal to $e_3$, as periodic Gibbs states are (Proposition~\textup{\YMMref{103.4}{758}}), the lattice approximants of a component among them, with $s,t\in\bigcup_KL^{-K}\mathbb Z$ and $D_j$ placed in $\{s<x_3<t-s\}$, whose inner operators $V_{1,j}:[G]\mapsto[\tau_{-se_3}D_j\cdot\tau_{(t-2s)e_3}G]$ are bounded uniformly in $j$ on the spaces of reflection positivity at $x_3=0$ of the shifted states $\rho_j\circ\tau_{se_3}$, $(\rho_j\circ\tau_{se_3})(F)=\rho_j(\tau_{se_3}F)$, which are reflection positive there because $\rho_j$ is so at the site plane $x_3=s$. (b) For every open $O\ne\emptyset$ the von Neumann algebra generated by the $e^{-tH}$ and the $U(x)\mathsf B(s,s)U(x)^{-1}$, $t,s>0$, $x\in O$, is $\mathcal B(\mathcal H)$.
\end{proposition}
\begin{proof}[Proof sketch]
(a) Take $O_1\ne\emptyset$ open and $\delta>0$ with $O_1+te_0\subset O$ for $|t|<\delta$, and a product $X$ of fields of $O_1$. Weak commutation gives $\langle Q^*\Omega,e^{itH}X\Omega\rangle=\langle X^\dagger\Omega,e^{-itH}Q\Omega\rangle$ for $|t|<\delta$. The left member continues to a bounded holomorphic function on ${\rm Im}\,\zeta>-s$, the right member to one on ${\rm Im}\,\zeta<s$; they agree on an interval, so they form one bounded entire function, a constant (Liouville). So the spectral measure of $H$ for $(Q^*\Omega,X\Omega)$ sits at $0$; $\mathcal D_0(O_1)$ is dense (Theorem~\ref{C-363-3}(a)), so $Q^*\Omega\in\ker H=\mathbb C\Omega$, and $Q^*-\bar c\in\mathcal C(O)$ annihilates $\Omega$, which separates $\mathcal C(O)$. The operators named are of the form $V=e^{-sH}V_1e^{-sH}$ with a bounded $V_1$: $\mathsf B(s,s')$ and the products by the semigroup law; for a weak limit $V$ of the $V_j$, $se_3$ and $te_3$ are lattice vectors for large $j$, and $\tau_{se_3}\theta\tau_{se_3}=\theta$ gives, for every state $\rho_j$,
$$\rho_j(\theta\bar G'\cdot D_j\cdot\tau_{te_3}G)=(\rho_j\circ\tau_{se_3})\bigl(\theta\tau_{se_3}\bar G'\cdot\tau_{-se_3}D_j\cdot\tau_{(t-2s)e_3}\tau_{se_3}G\bigr),$$
that is $\langle[G'],V_j[G]\rangle=\langle[\tau_{se_3}G'],V_{1,j}[\tau_{se_3}G]\rangle$, the right member in the space of $\rho_j\circ\tau_{se_3}$; the shifted states converge to $\sigma\circ\tau_{se_3}=\sigma$, as $\sigma$ is invariant under the translations, so that their spaces, with the classes $[G]$, are frames for $\mathcal H(\sigma)$ as in Lemma~\ref{C-363-8}; along a further subsequence $V_{1,j}\rightharpoonup V_1$ in these frames, $V_1$ bounded (Lemma~\ref{C-363-8}(ii)), and $[\tau_{se_3}G]=e^{-sH}[G]$ gives $\langle[G'],V[G]\rangle=\langle e^{-sH}[G'],V_1e^{-sH}[G]\rangle$ for all $G,G'$. Then $V\Omega,V^*\Omega\in e^{-sH}\mathcal H\subset D(e^{sH})$, and $V=c\ne0$ would make $e^{sH}$ bounded. If $H$ is bounded, the joint spectrum of the translations, which is invariant under the Lorentz transformations and lies in $\{E\ge|q|\}$, is bounded, so it is $\{0\}$; then $H=0$ and $\mathcal H=\ker H=\mathbb C\Omega$. (b) An operator $Y$ commuting with the generators commutes with $\lim_te^{-tH}=|\Omega\rangle\langle\Omega|$, so $Y\Omega=y\Omega$ and $Y=y$ on the vectors $e^{-t_0H}\mathsf b(x_1)e^{-t_1H}\cdots\mathsf b(x_n)\Omega$, $\mathsf b(x)=U(x)\mathsf B(s,s)U(x)^{-1}$, $x_k\in O$; a vector orthogonal to all of them is, by analyticity in the differences of the points and in the times, orthogonal to all point vectors, hence to $\mathcal D_0$ by (Br).
\end{proof}

\begin{lemma}[A second reflection]\label{C-363-11}
Let $\sigma$ be data (Rec), $Q$ bounded and $\tilde q\in\mathcal H$ with $\langle[G_1],Q[G_2]\rangle=\langle\tilde q,[R^{-1}(\theta G_1\cdot G_2)]\rangle$ for all tuples $G_1,G_2$ over $\mathsf{II}$. Then $Q\in\mathcal C(W')$.
\end{lemma}
\begin{proof}[Proof sketch]
For $\mathfrak z\in\mathbb C^3$ put $\mathfrak z^\flat=R^{-1}\mathfrak z=(\mathfrak z_3,\mathfrak z_2,-\mathfrak z_1)$ and $\mathfrak z^\natural=(-\overline{\mathfrak z_3},\overline{\mathfrak z_2},-\overline{\mathfrak z_1})$, so that $\mathfrak z^\natural=R^{-1}\theta\mathfrak z$ for real $\mathfrak z$; for a point $x$ of $\mathsf{II}$ the third coordinate of $x^\flat$ and of $x^\natural$ is $-x_1>0$. (1) Let $w\in\mathsf{II}^m$ and $z\in\mathsf{II}^n$ be Euclidean points with growing heights, $w_{1,3}<\dots<w_{m,3}$ and $z_{1,3}<\dots<z_{n,3}$, whose $N=m+n$ numbers $-w_{k,1}$, $-z_{k,1}$ are pairwise different. Then the identity of the lemma holds for the kernels: $\langle\Psi_m(w),Q\Psi_n(z)\rangle=\langle\tilde q,\Psi_N(\mathsf s(w,z))\rangle$, $\mathsf s(w,z)$ the points $w_k^\natural$, $z_k^\flat$ in the order of their third coordinates.

Indeed, integrated against real entries $g_k$, $h_k$ carried by small balls about the $w_k$, $z_k$ in $\mathsf{II}$, the two members are the two members of that identity at $G_1=(g_k)$, $G_2=(h_k)$, whose classes, and the class of $R^{-1}(\theta G_1\cdot G_2)$ after a permutation of its entries, are integrals of point vectors; and the kernels are continuous. (2) Let $0<r_1<\dots<r_N$ be these numbers, which are the heights of the points of $\mathsf s(w,z)$, and $y_i$ the second coordinates, and continue in the third coordinates alone: a ket point is $(-r_i,y_i,\omega_i)$, a bra point $(-r_i,y_i,-\overline{\omega_i})$, so that $\mathsf s(w,z)$ has the points $\mathsf s_i=(\omega_i,y_i,r_i)$. The left member is holomorphic in $\omega$ on the convex open set $D_1$ where the real parts of the $\omega_i$ of the ket points increase from $0$ in the order of the ket entries and those of the bra points decrease from $0$ in the order of the bra entries; the right member is holomorphic on the convex open set $D_2=\{|{\rm Im}\,\omega_1|<r_1,\ |{\rm Im}(\omega_{i+1}-\omega_i)|<r_{i+1}-r_i\}$, as $d(\mathsf s_1)=r_1-|{\rm Im}\,\omega_1|$ and $d(\mathsf s_{i+1}-\mathsf s_i)=r_{i+1}-r_i-|{\rm Im}(\omega_{i+1}-\omega_i)|$, whatever the real parts are. $D_1\cap\mathbb R^N\subset D_2$, and the two members agree there by (1); by Lemma~\ref{lem:setup-real-identity} they agree on $D_1\cap D_2$. (3) Let $p\in W'$ and $\mathsf r>0$ with $\overline{\mathbb B(p,\mathsf r)}\subset W'$, and $\mathbb B_k=\mathbb B(p-8k\mathsf re_1,\mathsf r)\subset W'$, $k\in\mathbb N_0$, balls of $\mathbb R^{1+2}$, $e_1$ the unit vector of the $x^1$-axis.

For $x\in\mathbb B_k$ and $x'\in\mathbb B_{k'}$, $k<k'$: $-x^1>|x^0|$ and $(-x'^1)-(-x^1)\ge6\mathsf r>2\mathsf r\ge|x'^0-x^0|$. Take $N$ functions $f_i\in\mathcal F(\mathbb B_{k_i})$, one in each of $N$ balls of pairwise different indices, and let $Y$ be a product of some of the $\phi(f_i)$ and $X$ a product of the others. By (Br) at $\mathfrak w_a=\eta e_3$ the vector $X_\eta$ obtained from $X\Omega$ by putting $e^{-\eta H}$ before each field is the integral of $\Psi_n$ at the points $\iota x+a\eta e_3=(x^1,x^2,a\eta-ix^0)$, $a$ counting the fields of $X$ from the left, and the like holds for $Y_\eta$; so $r_i=-x^1_i$, $y_i=x^2_i$, ${\rm Im}\,\omega_i=-x^0_i$ and ${\rm Re}\,\omega_i=\pm a\eta$, and $\omega\in D_1\cap D_2$ by the inequalities just written.

By (2) and as $\eta\downarrow0$, 
\[
\langle Y\Omega,QX\Omega\rangle=\int\prod\bar f\prod f\,\mathsf L(x)\,dx
\]
 with the one kernel 
\[
\mathsf L(x)=\langle\tilde q,\Psi_N((-ix^0_i,x^2_i,-x^1_i)_i)\rangle,
\]
whichever of the fields stand in the bra $Y$ or in the ket $X$. (4) Hence 
\[
\langle\phi(\bar h)\psi,Q\psi'\rangle=\langle\psi,Q\phi(h)\psi'\rangle
\]
 for $h\in\mathcal F(\mathbb B_0)$, $\psi$ in the span of $\Omega$ and of the products of fields, one in each of distinct balls $\mathbb B_k$ of odd index, applied to $\Omega$, and $\psi'$ in the like span with even indices $k\ge2$. (5) These spans stay dense after every power of $1+H$ (analyticity in the translations), and the polynomial bound of the field carries the relation to $\mathcal D_0$: $Q$ commutes weakly with every $\phi(h)$, $h\in\mathcal F(\mathbb B(p,\mathsf r))$.

Every $p\in W'$ and every $\mathsf r>0$ with $\overline{\mathbb B(p,\mathsf r)}\subset W'$ are admitted, these balls cover $W'$, and a partition of unity on the support of $h\in\mathcal F(W')$ gives $Q\in\mathcal C(W')$.
\end{proof}

\begin{lemma}[Two mirrors]\label{C-363-17}
Let $\sigma$ be data (Rec) and $Q$ a bounded operator. For combinations $F,F'$ over $\mathsf Q_L=\{x_1<0,x_3>0\}$ and $h=h_++h_-$, $h_\pm\in C^1_c(\{x_1>0,\ \pm x_3>0\})$, put $\ell^Q_{F',F}(h)=\langle[F'],Q[F\cdot h_+]\rangle+\langle[F'\cdot\theta\bar h_-],Q[F]\rangle$. If for all $F,F'$ there is $C<\infty$ with $|\ell^Q_{F',F}(h)|\le C\,\sigma(\theta_1\bar h\cdot h)^{1/2}$ for all such $h$, then $Q\in\mathcal C(W)$.
\end{lemma}
\begin{proof}[Proof sketch]
$\sigma(\theta_1\bar h\cdot h)=\|[\mathsf dh]\|^2$ is the norm of reflection positivity at the mirror $x_1=0$. Step 1: by Riesz, $\ell(h)=\langle\Psi,[\mathsf dh]\rangle$ with some $\Psi\in\mathcal H$; with $v(y)=\mathsf B(y,1)\Omega$ and $[k]=\int k(x)U(\vec x)v(x_3)dx$ for $k\in C^1_c(\mathbb E_+)$, $\ell(h)=\int h(x)\Phi(x_1,x_2;x_3)dx$, $\Phi(x_1,x_2;\zeta)=\langle\Psi,U((\zeta,x_2))v(x_1)\rangle$, and by the spectrum condition $\Phi$ continues holomorphically to $|{\rm Im}\,\zeta|<x_1$: its domain holds all real $\zeta$ of both signs. Step 2: with $F,F'$ over $\{x_3>\tau_0\}$, $\psi=[\tau_{-\tau_0e_3}F]$, $\psi'=[\tau_{-\tau_0e_3}F']$ and $\mathsf B(\zeta,\tau_0-\zeta)=e^{-i{\rm Im}\zeta H}\mathsf B({\rm Re}\,\zeta,\tau_0-{\rm Re}\,\zeta)e^{i{\rm Im}\zeta H}$, the first mirror gives the two functions $\mathsf G_+(\vec x;\zeta)=\langle Q^*[F'],U(\vec x)\mathsf B(\zeta,\tau_0-\zeta)U(\vec x)^{-1}\psi\rangle$, $0<{\rm Re}\,\zeta<\tau_0$, and $\mathsf G_-(\vec x;\zeta)=\langle\psi',U(\vec x)\mathsf B(\tau_0+\zeta,-\zeta)U(\vec x)^{-1}Q[F]\rangle$, $-\tau_0<{\rm Re}\,\zeta<0$; they are $\Phi$ on real segments, so on $0<\pm{\rm Re}\,\zeta<\tau_0$, $|{\rm Im}\,\zeta|<x_1$.

Step 3: for $f\in\mathcal F(W)$, $|t|\le x_1-\delta$ on ${\rm supp}f$; by \eqref{C-363-2}, 
\[
\int f\,\mathsf G_+(\vec x;s-it)=\langle e^{-sH}Q^*[F'],\phi(f)e^{-(\tau_0-s)H}\psi\rangle\to\langle[F'],Q\phi(f)[F]\rangle
\]
 and $\int f\,\mathsf G_-(\vec x;-s-it)\to\langle\phi(\bar f)[F'],Q[F]\rangle$ as $s\downarrow0$, and both tend to $\int f\,\Phi(x_1,x_2;-it)$. Step 4: the classes over $\mathsf Q_L$ are dense after every power of $N$ (a translation along $e_1$ carries every class into them, and a vector orthogonal to them gives a function of the translation that vanishes on a half-line and, as $|q_1|\le E$ on the spectrum and every class is $e^{-\varepsilon H}$ of a class, continues holomorphically to a strip, hence vanishes), and the polynomial bound of the field carries the relation to $\mathcal D_0$.
\end{proof}

\begin{lemma}[Central positivity]\label{C-363-12}
(a) Let $\mathfrak A$ be a unital $*$-algebra, $\omega$ a state on it with GNS data $(\pi,\mathcal D,\Omega)$, and $X_0=X_0^*$ central. The following are equivalent: (CP) $\omega(\sum X_0^nY_n)\ge0$ whenever $\sum\lambda^nY_n\in\mathfrak A^+$ for all real $\lambda$; (SA) $\pi(X_0)$ has a self-adjoint extension on a larger Hilbert space whose bounded Borel functions $F$ satisfy $\langle\pi(Y^*)\psi,F\psi'\rangle=\langle\psi,F\pi(Y)\psi'\rangle$ for $Y\in\mathfrak A$, $\psi,\psi'\in\mathcal D$; (CP$^\sharp$) $\omega(a)\ge0$ whenever $(1+X_0^2)^ka\in\mathfrak A^+$ for some $k\ge0$; (Ext) with $\mathscr Q$ the polynomials $\sum_n\lambda^nY_n$ in a real variable $\lambda$ with coefficients $Y_n\in\mathfrak A$, ${\rm ev}:\mathscr Q\to\mathfrak A$ the $*$-homomorphism $\sum_n\lambda^nY_n\mapsto\sum_nX_0^nY_n$, $\mathfrak r=\mathbb C[\lambda,(1+\lambda^2)^{-1}]$ and $\mathscr Q_{\mathfrak r}$ the $*$-algebra of the functions $\sum_kg_k(\lambda)Y_k$, $g_k\in\mathfrak r$, $Y_k\in\mathfrak A$: $\omega\circ{\rm ev}$ has a linear extension $\tilde\omega$ to $\mathscr Q_{\mathfrak r}$ with $\tilde\omega(F^*F)\ge0$ for all $F\in\mathscr Q_{\mathfrak r}$. (b) A closed symmetric operator $T$ affiliated with a von Neumann algebra $\mathfrak M$ has a self-adjoint extension on $\mathcal H\otimes\mathbb C^2$ affiliated with $\mathfrak M\otimes M_2(\mathbb C)$. (c) Two closed symmetric operators affiliated with commuting von Neumann algebras have strongly commuting self-adjoint extensions on $\mathcal H\otimes\mathbb C^2\otimes\mathbb C^2$. (d) For a linear functional $L$ on $\mathbb C[x,y]$ with $L(|q|^2)\ge0$: $L(p)\ge0$ for all real $p\ge0$ on $\mathbb R^2$ if and only if $L(a)\ge0$ for all real $a$ such that $(1+x^2)^ka$ is a sum of squares for some $k$; Motzkin's polynomial $M$ has the level $k=1$: $(1+x^2)M=(x^2y^2-1)^2+(xy(x^2-1))^2+(x(y^2-1))^2$.
\end{lemma}
\begin{proof}[Proof sketch]
(a) (SA)$\Rightarrow$(CP), (CP$^\sharp$): with $E$ the spectral measure of the extension let $\mu_Y=\langle\Omega,E(\cdot)\pi(Y)\Omega\rangle$, $Y\in\mathfrak A$, a complex measure, linear in $Y$. By (SA), $\mu_{Z^*Z}=\langle\pi(Z)\Omega,E(\cdot)\pi(Z)\Omega\rangle\ge0$, so $\mu_b\ge0$ for $b\in\mathfrak A^+$; and $d\mu_{X_0^nY}=\lambda^nd\mu_Y$, as the extension continues $\pi(X_0)$ on $\pi(\mathfrak A)\Omega$. If $(1+X_0^2)^ka=b\in\mathfrak A^+$, then $(1+\lambda^2)^kd\mu_a=d\mu_b\ge0$, so $\mu_a\ge0$ and $\omega(a)=\mu_a(\mathbb R)\ge0$: (CP$^\sharp$). If $\sum_n\lambda^nY_n\in\mathfrak A^+$ for all real $\lambda$, then for each real $\lambda_0$ the measure $\sum_n\lambda_0^n\mu_{Y_n}$ is positive, so the densities $\rho_n$ of the $\mu_{Y_n}$ with respect to the sum of their variations satisfy $\sum_n\lambda_0^n\rho_n\ge0$ for all rational $\lambda_0$ off one null set, hence $\sum_n\lambda^n\rho_n(\lambda)\ge0$, and $\omega(\sum X_0^nY_n)=\int\sum_n\lambda^n\rho_n(\lambda)\ge0$: (CP). (CP)$\Rightarrow$(Ext): the extension theorem of M.~Riesz, for the real space of the hermitian $F\in\mathscr Q_{\mathfrak r}$, its cone $\mathsf K$ of the $F$ with $F(\lambda)\in\mathfrak A^+$ for every real $\lambda$, which holds every $F^*F$, and the subspace of the hermitian polynomials, on whose part in $\mathsf K$ the functional $\omega\circ{\rm ev}$ is $\ge0$ by (CP), gives an extension $\tilde\omega\ge0$ on $\mathsf K$, continued complex-linearly to $\mathscr Q_{\mathfrak r}$, because every hermitian $F$ is majorised by a hermitian polynomial: $F$ is a finite sum of terms $gS$, $g\in\mathfrak r$ real, $S=S^*\in\mathfrak A$ (as $gY+\bar gY^*=({\rm Re}\,g)(Y+Y^*)+({\rm Im}\,g)\,i(Y-Y^*)$), and for a real polynomial $q\ge|g|$
$$\tfrac12q(S^2+1)-gS=\tfrac14(q-g)(S+1)^2+\tfrac14(q+g)(S-1)^2\in\mathsf K ;$$
for $|g|\le C(1+\lambda^2)^r$ one may take $q=C(1+\lambda^2)^r$, so that the extension holds at each fixed $r\in\mathbb N_0$, from the polynomials $\sum_{n\le2r}\lambda^nY_n$ to the functions whose coefficients are $O((1+\lambda^2)^r)$. (CP$^\sharp$)$\Rightarrow$(Ext): the same theorem, with the cone of the finite sums of the $F^*F$, $F\in\mathscr Q_{\mathfrak r}$, in the place of $\mathsf K$. If a polynomial $p$ is such a sum $\sum_iF_i^*F_i$, then the $G_i=(1+\lambda^2)^kF_i$ are polynomials for some $k$, $(1+\lambda^2)^{2k}p=\sum_iG_i^*G_i$, and the $*$-homomorphism ${\rm ev}$ gives $(1+X_0^2)^{2k}{\rm ev}(p)\in\mathfrak A^+$, so $\omega({\rm ev}\,p)\ge0$ by (CP$^\sharp$). The majorisation: for $g=p_0/(1+\lambda^2)^k$, $p_0$ a real polynomial, $q=\frac12(1+p_0^2)\ge|g|$, and $q\mp g=n_\mp/(1+\lambda^2)^k$ with real polynomials $n_\mp\ge0$ on $\mathbb R$, which are $|m_\mp|^2$ for complex polynomials $m_\mp$, a polynomial that is $\ge0$ on $\mathbb R$ being a sum of two squares; so $q\mp g=|u_\mp|^2$ with $u_\mp=m_\mp(\lambda-i)^k(1+\lambda^2)^{-k}\in\mathfrak r$, and the identity above reads $\frac12q(S^2+1)-gS=(\frac12u_-(S+1))^*(\frac12u_-(S+1))+(\frac12u_+(S-1))^*(\frac12u_+(S-1))$. (Ext)$\Rightarrow$(SA): $\tilde\omega(F^*)=\overline{\tilde\omega(F)}$ by polarisation; in the GNS space of $\tilde\omega$, which holds that of $\omega$ isometrically, multiplication by $\lambda$ is symmetric with $(\lambda\pm i)\mathscr Q_{\mathfrak r}=\mathscr Q_{\mathfrak r}$, as $(\lambda\pm i)^{-1}\in\mathfrak r$, so its closure is self-adjoint and extends $\pi(X_0)$; the multiplications by the $Y\in\mathfrak A$ commute with its resolvents, and the bounded $C$ with $\langle\pi(Y^*)\psi,C\psi'\rangle=\langle\psi,C\pi(Y)\psi'\rangle$ for $Y\in\mathfrak A$, $\psi,\psi'\in\mathcal D$, form a weakly closed space that holds the polynomials in the two resolvents, hence, by the bicommutant theorem, all bounded Borel functions of the closure. (b) $T\oplus(-T)$ has the deficiency spaces $\mathcal N_+\oplus\mathcal N_-$ and $\mathcal N_-\oplus\mathcal N_+$, which the flip $(x,y)\mapsto(y,x)\in1\otimes M_2$ exchanges isometrically; the extension that this flip defines commutes with the unitaries of $\mathfrak M'\otimes1$. (c) Apply (b) in two tensor factors. (d) This is (CP)$\Leftrightarrow$(CP$^\sharp$) of (a) for $\mathfrak A=\mathbb C[x,y]$, $X_0=x$ and $\omega=L/L(1)$ (if $L(1)=0$, then $L=0$ by the Schwarz inequality), where $\mathfrak A^+$ is the set of the sums of squares of real polynomials, as $|q|^2=({\rm Re}\,q)^2+({\rm Im}\,q)^2$: there (CP) says $L(p)\ge0$ for every real $p\ge0$ on $\mathbb R^2$, since $\sum_n\lambda^nY_n\in\mathfrak A^+$ for all real $\lambda$ gives $\sum_nx^nY_n\ge0$ on $\mathbb R^2$ (put $\lambda=x$), and a real $p=\sum_nx^np_n(y)\ge0$ has $\sum_n\lambda^np_n\in\mathfrak A^+$ for every real $\lambda$, a polynomial in one variable that is $\ge0$ being a sum of two squares; and if $(1+x^2)^ka$ is a sum of squares, $a$ is real.

The identity is checked by expanding, the three squares summing to $x^6y^2+x^4y^4-2x^4y^2+x^2y^4-3x^2y^2+x^2+1$.
\end{proof}

\begin{lemma}[Decomposition and limits]\label{C-363-13}
Let $(\mathsf X,\lambda,(\sigma_\xi))$ be a canonical decomposition of data (Rec) $\sigma$, all $\sigma_\xi$ data (Rec) with one triple, and $\mathcal W_\xi$ their Wightman functionals.
(i) There is exactly one unital $*$-representation $g\mapsto M_g$ of $C(\mathsf X)$ on $\mathcal H(\sigma)$ with $\langle[F'],M_g[F]\rangle=\int g(\xi)\sigma_\xi(\theta\bar F'\cdot F)\lambda(d\xi)$; $\|M_g\|\le\sup|g|$, and $M_g$ commutes with $e^{-tH}$, the $U(x)$, the insertions and the $\mathsf B(s,s')$, keeps every $D(H^a)$, and commutes with every $\phi(f)$ on $\mathcal D_\infty$; so $M_g\in\mathcal C(O)$ for every $O$.
(ii) For $u\in L^1(\mathbb R^{1+2})$ let $\phi_u(\psi',\psi)=\int u(x)\,q_0(U(x)^{-1}\psi',U(x)^{-1}\psi)\,dx$ on the vectors $\psi=e^{-\eta H}\chi$, $\psi'=e^{-\eta H}\chi'$, $\eta>0$, among them the classes $[F]$ and the $M_g[F]$, $F$ over $\mathbb E_+$: the form of the field smeared with $u$, with $|\phi_u(\psi',\psi)|\le\|u\|_{L^1}C_\wp\min\{\eta,1\}^{-\wp}\|\chi'\|\|\chi\|$ and $\phi_u(\psi',\psi)=\langle\psi',\phi(u)\psi\rangle$ for $u\in\mathcal S$, by \eqref{C-363-2}; $\phi^\xi_u$ is the same form for $\sigma_\xi$. For combinations $F,F'$ over $\mathbb E_+$ and $u\in L^1(\mathbb R^{1+2})$ the functions $\xi\mapsto\langle[F']_\xi,[F]_\xi\rangle$, $\langle[F']_\xi,H_\xi[F]_\xi\rangle$ and $\phi^\xi_u([F']_\xi,[F]_\xi)$ are continuous, and $\phi_u([F'],M_g[F])=\int g\,\phi^\xi_u([F']_\xi,[F]_\xi)\,d\lambda$; the like for $1$ and for $H$.
(iii) For $a\in\mathfrak P$, $\xi\mapsto\mathcal W_\xi(a)$ is continuous and $\langle\Omega,M_g\pi(a)\Omega\rangle=\int g\,\mathcal W_\xi(a)\,d\lambda$; in particular $\mathcal W_\sigma=\int\mathcal W_\xi\,d\lambda$. For $a,b_1,\dots,b_q\in\mathfrak P$ and spatial translations $v_i$, $\lim_{|v_q|\to\infty}\cdots\lim_{|v_1|\to\infty}\mathcal W_\sigma(a\,\tau_{v_1}b_1\cdots\tau_{v_q}b_q)=\int\mathcal W_\xi(a)\prod_i\mathcal W_\xi(b_i)\,\lambda(d\xi)$, each limit existing in turn.
(iv) Let $a\in\mathfrak P$, $a^*=a$, and suppose that for all $b_1,\dots,b_r\in\mathfrak P$ there is $R$ such that $\mathcal W_\sigma(b^*ab)\ge0$ for every $b$ obtained from a complex polynomial $\mathsf p$ in $r$ variables by putting, at each occurrence of its $i$-th variable, a mean $\frac1{\mathsf n}\sum_{\nu=1}^{\mathsf n}\tau_{v_\nu}b_i$ over spatial translations $v_\nu=(0,\vec v_\nu)$ of length $\ge R$, all the translations that occur in $b$ being pairwise different. Then $\mathcal W_\xi(a)\ge0$ for every $\xi$.
(v) If data (Rec) $\sigma_j$ with one triple converge to $\sigma$ on separated tuples, the forms of the fields and of the energy on the classes of combinations converge, and $\mathcal W^{\sigma_j}(f_1\cdots f_n)\to\mathcal W^\sigma(f_1\cdots f_n)$ for all $f_k\in\mathcal S$.
\end{lemma}
\begin{proof}[Proof sketch]
(i) On $C(\mathsf X)\odot\{F\}$ the form $\int\bar g'g\langle[F']_\xi,[F]_\xi\rangle d\lambda$ is positive; multiplication by $g$ acts on its Hilbert space $\mathfrak K$, and $[F]\mapsto1\otimes F$ is an isometry $\mathsf W$ of $\mathcal H(\sigma)$ into $\mathfrak K$. For a tuple $Y$ over $\mathbb E_+$ and far horizontal translations $v$ the vectors $[F\cdot\tau_vY]$ tend weakly to $\mathsf W^*\mathsf m(g_Y)\mathsf W[F]$, $g_Y(\xi)=\sigma_\xi(Y)$, and the iterated cluster limits give $\|\mathsf W^*\mathsf m(g_Y)\mathsf W[F]\|=\|\mathsf m(g_Y)\mathsf W[F]\|$, so $\mathsf m(g_Y)$ keeps the range of $\mathsf W$. The $g_Y$ separate the points of $\mathsf X$ ($\xi\mapsto\sigma_\xi$ is injective); by Stone--Weierstrass every $\mathsf m(g)$ keeps the range, which is then all of $\mathfrak K$, and $M_g=\mathsf W^{-1}\mathsf m(g)\mathsf W$.

It commutes with $e^{-tH}$ because each $e^{-tH_\xi}$ is symmetric, with the translations, and with the insertions because the vector $\mathsf B^{\rm in}_\eta(h)[G]$ is a class given by one formula for $\sigma$ and for all $\sigma_\xi$; hence with all $\mathsf B(s,s')$ and with the fields. (ii) For an insertion smeared in space, $\mathsf G_{\mathsf k}(z)=\int\mathsf k(\vec x)\langle e^{-(\frac\eta2+\bar z)H}\chi',U(\vec x)\mathsf B(\frac\eta2,\frac\eta2)U(\vec x)^{-1}e^{-(\frac\eta2-z)H}\chi\rangle d^2x$ is holomorphic on the strip $|{\rm Re}\,z|<\frac\eta2$ and bounded by $\|\mathsf k\|_{L^1}C_\wp(\eta/2)^{-\wp}\|\chi'\|\|\chi\|$, the same numbers for all $\xi$; at real $z$ it gives the Euclidean insertion, at $z=-it$ the real-time field.

Montel's theorem and the identity theorem carry the continuity in $\xi$ of the Euclidean values to the imaginary axis; the same with the weight $M_g$. (iii) In $n$ variables: 
\[
\mathcal W^{(\varepsilon)}_n(f_1,\dots,f_n)=\langle\Omega,\phi(f_1)e^{-\varepsilon H}\cdots e^{-\varepsilon H}\phi(f_n)\Omega\rangle
\]
 is a value on the imaginary axes of a function holomorphic and bounded on a polystrip, uniformly in $\xi$, and $|\mathcal W_n-\mathcal W_n^{(\varepsilon)}|\le\varepsilon\,\mathsf C(f_1,\dots,f_n)$ with $\mathsf C$ formed from the triple and the $\|\partial_t^kf_i\|_{L^1}$; Montel and the identity theorem give the continuity. For the cluster limits, the analytic Wightman functions are bounded by $(C_\wp\eta^{-\wp})^n$ uniformly in $\xi$ and in horizontal translations, and the rate
\[
\begin{gathered}
\bigl|\mathcal W(f_1\cdots f_n)-\langle\Omega,e^{-\eta H}\phi(f_1)\cdots e^{-\eta H}\phi(f_n)\Omega\rangle\bigr|\\
\le c\int_0^\eta(1+|\log s|)\,ds,\qquad\\
 c=2^{\mathsf m-1}(2^\wp C_\wp)^n\max_{l\le\mathsf m+1}\Bigl\|\mathsf D^l\prod_kf_k\Bigr\|_{L^1},
\end{gathered}
\]
$\mathsf D=\sum_kk\,\partial/\partial x^0_k$, $\mathsf m$ an integer $\ge\max\{2,n\wp\}$ (by (Br) at $\mathfrak w_k=\eta e_3$ the integrand is holomorphic in the $x^0_k+ik\eta$, and $\mathsf m+1$ integrations by parts in the real times, with the bound of the point vectors at $\frac\eta2$, bound the derivatives in $\eta$), with $c$ formed from the triples ($C_\wp$, $\wp$ the larger of the numbers of the triple of $\sigma$ and of the one triple of the $\sigma_\xi$), $n$ and norms that translations of the $f_k$ do not change, so free of $\xi$ and of the translations, exchanges the cluster limits with $\eta\downarrow0$. (iv) Expand $b^*ab$ in the means and let the translations tend to infinity one after the other, a factor whose translation goes off commuting in $\mathfrak P$ with the others once it is spacelike to them.

By (iii) the terms in which no translation of $b^*$ is one of $b$ tend to $\int\mathcal W_\xi(a)$ times the product of the $\mathcal W_\xi(b_i)$ and $\mathcal W_\xi(b_i^*)=\overline{\mathcal W_\xi(b_i)}$ that they carry, and the others, a share $O(\frac1{\mathsf n})$ of the terms, are bounded uniformly, $\xi\mapsto\mathcal W_\xi(c')$ being continuous on the compact $\mathsf X$ for each $c'\in\mathfrak P$. So $0\le\int\mathcal W_\xi(a)(|\mathsf q(\xi)|^2+O(\frac1{\mathsf n}))\,d\lambda$, $\mathsf q=\mathsf p(\mathcal W_\xi(b_1),\dots,\mathcal W_\xi(b_r))$, and $\int\mathcal W_\xi(a)|\mathsf q|^2d\lambda\ge0$. These $\mathsf q$ form a unital $*$-algebra of continuous functions that separates the points of $\mathsf X$, as $\xi\mapsto\sigma_\xi$ is injective and $\mathcal W_\xi$ fixes $\sigma_\xi$ (its analytic continuation to Euclidean points gives the Schwinger functions of $\sigma_\xi$); it is dense in $C(\mathsf X)$ (Stone--Weierstrass). $\mathcal W_\xi(a)$ is real and continuous and $\lambda$ has full support: $\mathcal W_\xi(a)\ge0$ for every $\xi$. (v) The proofs of (ii) and (iii) use of the family only one triple and the convergence on separated tuples.
\end{proof}

\begin{lemma}[Bump functions with transforms free of zeros; division]\label{C-363-19}
(a) For $\epsilon_k>0$ with $\sum_k\epsilon_k=\epsilon<\infty$ let $\mathsf u$ be the infinite convolution of the probability densities $\frac{e^{-x/\epsilon_k}}{\epsilon_k(1-e^{-1})}$ on $[0,\epsilon_k]$, and $w=\mathsf u*\mathsf u(-\,\cdot)$. Then $w\in C^\infty_c(\mathbb R)$, $w\ge0$, $w$ is even, $\int w=1$, ${\rm supp}\,w\subset[-\epsilon,\epsilon]$, $\hat w$ is entire, $\hat w(s)=|\hat{\mathsf u}(s)|^2>0$ on $\mathbb R$, and $\hat w$ has no zero in $|{\rm Im}\,\zeta|<1/\epsilon$; $f_*=w\otimes w\otimes w$ is the function of Theorem~\ref{C-363-3}(d). At $\epsilon_k=\frac{6\epsilon}{\pi^2k^2}$: $\|w^{(k)}\|_{L^1}\le A_0R_0^k(k!)^2$, $A_0=\frac{e^3\pi}2$, $R_0=\frac{3e\pi^2}{2\epsilon}$. (b) For $\delta>0$, $\gamma>1$ and $a_k=\delta k^{-\gamma}/\sum_l l^{-\gamma}$ let $g_{\delta,\gamma}$ be the infinite convolution of the uniform probability densities on $[-\frac{a_k}2,\frac{a_k}2]$: $g_{\delta,\gamma}\in C^\infty_c(\mathbb R)$ is $\ge0$, even, $\int g_{\delta,\gamma}=1$, ${\rm supp}\,g_{\delta,\gamma}\subset[-\frac\delta2,\frac\delta2]$.

If $\epsilon_k=\epsilon2^{-k}$ and $\gamma>1$, or $\epsilon_k=6\epsilon/\pi^2k^2$ and $1<\gamma<2$, then $g_{\delta,\gamma}=w*\mu$ with a real even $\mu\in\mathcal S(\mathbb R)$, whatever the sizes of $\delta$ and $\epsilon$; so $g_{\delta,\gamma}^{\otimes3}=f_**\mu^{\otimes3}$. (c) For data (Rec), real $u\in\mathcal S$ and real $\mu\in L^1(\mathbb R^{1+2})$ with $u*\mu\in\mathcal S$: $\langle\psi,\phi(u*\mu)\psi'\rangle=\int\mu(a)\langle U(a)^{-1}\psi,\phi(u)U(a)^{-1}\psi'\rangle\,da$ on $\mathcal D_\infty$, so $c_*(u*\mu)\le\|\mu\|_{L^1}c_*(u)$ and $b_{mn}(u*\mu)\le\|\mu\|_{L^1}b_{mn}(u)$; and $c_*(u\circ\Lambda^{-1})\le\bigl((\Lambda e_0)^0+|\overrightarrow{\Lambda e_0}|\bigr)c_*(u)$ for Lorentz transformations $\Lambda$. (d) For $f_0\in C^\infty_c$, $f_0\ge0$, $f_0\ne0$, every real $p$ has a Lorentz transformation $\Lambda$ with $\hat f_0(\Lambda^{\rm T}p)\ne0$.
\end{lemma}
\begin{proof}[Proof sketch]
(a) $\hat{\mathsf u}(\zeta)=\prod_k\frac{1-e^{-(1-i\zeta\epsilon_k)}}{(1-e^{-1})(1-i\zeta\epsilon_k)}$, a convergent product since $\sum\epsilon_k<\infty$; no factor vanishes for $|{\rm Im}\,\zeta|<1/\epsilon_k$, as $|e^{-(1-i\zeta\epsilon_k)}|<1$ there. The decay of $\hat{\mathsf u}$ faster than every power gives smoothness; on $|{\rm Im}\,\zeta|\le\frac1{4\epsilon}$, $|\hat w(s+iq)|\ge e^{-2\epsilon|q|}\Pi(s)^{-2}$, $\Pi(s)=\prod_k(1+\frac{16}9s^2\epsilon_k^2)^{1/2}$, with $\Pi(s)\le e^{1/6}(1+y^2)^{(1+\log_2^+y)/2}$, $y=\frac43\epsilon|s|$, at the dyadic choice and $\Pi(s)\le\exp(2.6(8\epsilon|s|/\pi^2)^{1/2})$ at the other; there also $\hat w(s)\le\exp(2-2(2\epsilon|s|/e\pi^2)^{1/2})$, which gives the Gevrey bound. (b) $\hat g_{\delta,\gamma}(\zeta)=\prod_k\frac{\sin(a_k\zeta/2)}{a_k\zeta/2}$ has $|\hat g_{\delta,\gamma}(s+iq)|\le\exp(\frac\delta2|q|+1-(\delta|s|/2e\zeta_\gamma)^{1/\gamma})$, $\zeta_\gamma=\sum k^{-\gamma}$: it falls faster than $|\hat w|$ falls, and $\hat w$ has no zero on a strip; so $\hat g_{\delta,\gamma}/\hat w$ is holomorphic on the strip and its restriction to $\mathbb R$ is a Schwartz function. (c) $\phi(u*\mu)=\int\mu(a)U(a)\phi(u)U(a)^{-1}da$, and $U(a)$ commutes with $N$ and with the $P_m$; for a Lorentz transformation $U(\Lambda)HU(\Lambda)^{-1}=\int(\Lambda^0{}_0E+\Lambda^j{}_0q_j)d\mathcal P\le((\Lambda e_0)^0+|\overrightarrow{\Lambda e_0}|)H$ on the spectrum $\{E\ge|q|\}$. (d) Rotate $p$ to $p'$ with $p'_2=0$; $s\mapsto\hat f_0(B(s)^{\rm T}p')$ is entire, and at $s=i\frac\pi2$ it is $\int e^{p'_1x^0+p'_0x^1}f_0(x)dx>0$; so it does not vanish at some real $s$.
\end{proof}

\begin{lemma}[Sharp loops on a torus]\label{C-363-41}
Let $\mathsf b>0$, $\ell=2L^m$, $a_K=L^{-K}$, $\beta=L^K\mathsf b$, and let $\langle\cdot\rangle_K$ be the Wilson measure $\mu_\beta$ on $\Lambda_{K,\ell}=a_K\mathbb Z^3/\ell\mathbb Z^3$. For $j\in\frac12\mathbb N$ let $D^j$ be the irreducible representation of $\mathrm{SU}(2)$ of dimension $2j+1$ and $\chi_j$ its character. Let $C$ be a closed path of $\Lambda_{K_0,\ell}$ for some $K_0$, read as a path of every $\Lambda_{K,\ell}$, $K\ge K_0$, that has no self-intersection, does not wind (the vectors of its bonds add up to $0$ in $\mathbb R^3$) and has a neighbourhood in a cube of side $\le\ell/2$; $|C|$ is its length. Let $\mathsf d>0$ and let $F_K$ be Borel functions, $|F_K|\le1$, of the bonds of $\Lambda_{K,\ell}$ at a distance $\ge\mathsf d$ from $C$. Then for every $j\ge\frac12$ there are $K^{\rm sl}_j$, formed from $j$, $C$, $\mathsf b$, $\ell$ and $\mathsf d$, and $A^{\rm sl}_j$, formed from $j$, $C$, $\mathsf b$ and $\ell$, with
\begin{equation}\label{C-363-42}
\bigl|\langle F_K\,\chi_j(U(C))\rangle_K\bigr|\le A^{\rm sl}_j\,\bigl(a_K\log^3a_K^{-1}\bigr)^{j(j+1)|C|/4\pi\mathsf b}\qquad(K\ge K^{\rm sl}_j).
\end{equation}
Hence, for every continuous $h$ on $[-1,1]$, $\langle F_Kh(W_C)\rangle_K-\bar h\langle F_K\rangle_K\to0$ uniformly in such $F_K$, where $W_C=\frac12{\rm tr}\,U(C)$ and $\bar h=\frac2\pi\int_{-1}^1h(w)\sqrt{1-w^2}\,dw$.
\end{lemma}
\begin{proof}[Proof sketch]
We count in units of $a_K$: $C$ has $n=|C|/a_K$ steps and $\mathsf p_C$ sides; $\mathsf c_j=j(j+1)$, $\Lambda_C=\max\{1,|C|/\mathsf b\}$ and $\rho=\lfloor(\beta/\Lambda_C\log^3\beta)^{1/4}\rfloor$. \emph{Slabs.} For a step $c$ of $C$ with the first point $x_-$ and the vector $\mathsf e$ let $\Pi(c)$ be the lattice plane through $x_-$ orthogonal to $c$, $\mathsf T_\rho(c)$, the slab of $c$, the set of the translates $c(y)$ of $c$ by the vectors $y$ of $\Pi(c)-x_-$ with $|y|<\rho$, and $\{x_-+y+s'\mathsf e:|y|_\infty\le\rho,\ 0\le s'\le1\}$ its box. Let $K\ge K^{\rm sl}_j$, where $K^{\rm sl}_j$ is so large that for $K\ge K^{\rm sl}_j$: $\ell/a_K\ge10$, $\beta\ge e^3$, $\rho\ge2$, $\frac43\mathsf c_j\le20\log\beta$, every side of $C$ has at least $4\rho+1$ steps, two sides without a common corner are at a sup-distance $\ge2\rho+1$, the points at a sup-distance $\le\rho+1$ from $C$ lie in a cube of side $\le\ell/2a_K$, and $\rho a_K<\mathsf d$.

The steps $c_i$ of $C$ at least $2\rho$ steps away from its corners form a set $I$ of $n_I\ge n-4\rho\mathsf p_C$ steps; no step of $C$ but $c_i$ runs through a bond of $\mathsf T_\rho(c_i)$, the box of $c_i$ holds no bond of another slab of $I$, and no plaquette contains bonds of two of these slabs. So, given the bonds off $B=\bigcup_{i\in I}\mathsf T_\rho(c_i)$, the slabs are independent under the Wilson measure, and $F_K$ and the paths $\delta_i$ of $C$ from $c_i$ to the next step of $I$ read no bond of $B$; with $Q_i$ the conditional mean of $D^j(U(c_i))$, $\langle F_K\chi_j(U(C))\rangle_K=\langle F_K\,{\rm tr}\prod_iQ_iD^j(U(\delta_i))\rangle_K$, and $|\langle F_K\chi_j(U(C))\rangle_K|\le(2j+1)\langle\prod_i\|Q_i\|\rangle_K$.
\emph{One slab.} Given the outer bonds, the variables $v_y=U(c(y))$, $|y|<\rho$, of the slab of $c=c_i$ have a density $\varrho_\circ\propto\exp(-\beta\sum_ea_{p(e)})$ for Haar measure, $p(e)$ the plaquette spanned by a bond $e$ of $\Pi(c)$ and the direction of $c$. Move $v_y\mapsto\exp(is\varphi_\rho(y)\,\mathsf n_y\cdot\frac\sigma2)v_y$, with the logarithmic weight $\varphi_\rho(y)=\min\{1,\log(\rho/|y|)/\log\rho\}$, $\varphi_\rho(0)=1$, $\varphi_\rho(y)=0$ for $|y|\ge\rho$, and the colour $\mathsf n_y\cdot\sigma=P_y^{-1}(\mathsf t\cdot\sigma)P_y$, $\mathsf t$ a unit vector of $\mathbb R^3$, $P_y=U(\pi_y)$ for the path $\pi_y$ of $\Pi(c)$ from $x_-$ along the first axis and then parallel to the second axis to $x_-+y$.

These moves are left translations and keep Haar measure; integration by parts along them and the Schwarz inequality give the information inequality for the three colours at once, $\|Q_i\|^2\le(1+\mathsf c_j/\bar{\mathcal I})^{-1}$, $\bar{\mathcal I}$ the largest, over $\mathsf t$, conditional mean of $-\partial_s^2\log\varrho_\circ$ at $s=0$. Along a move, for a bond $e$ of $\Pi(c)$ from $x_-+y$ to $x_-+y'$, $\partial_s^2a_{p(e)}\le\frac14(\varphi_\rho(y)-\varphi_\rho(y'))^2$, plus $\varphi_\rho(y)\varphi_\rho(y')(3a(H_e)+a_{p(e)})$, $a(g)=1-\frac12{\rm tr}\,g$, if $e$ is a rung of the comb of the $\pi_y$, a bond parallel to the first axis at a distance $t_e\ge1$ from it; there $H_e=P_{y'}U(e)^{-1}P_y^{-1}$ is a conjugate of the holonomy around the strip $S_e$ of the $t_e$ plaquettes of $\Pi(c)$ between $e$ and the first axis, and $a(H_e)\le t_e\sum_{q\in S_e}a_q$ (Stokes for a strip).

With $\sum(\varphi_\rho(y)-\varphi_\rho(y'))^2\le\bar E_\rho=(2\pi\log\rho+7.63)/\log^2\rho$, the sum over the bonds $\langle y,y'\rangle$ of $\mathbb Z^2$: $\bar{\mathcal I}\le\frac\beta4(\bar E_\rho+X_i)$, $X_i=\sum_e\varphi_\rho(y)\varphi_\rho(y')(12a(H_e)+4\,\mathbb E_i(a_{p(e)}))$ over the rungs, $\mathbb E_i$ the conditional mean, and
$$\log\|Q_i\|\le-\frac{2\mathsf c_j}{\beta\bar E_\rho}\Bigl(1-\frac{2\mathsf c_j}{\beta\bar E_\rho}\Bigr)+\frac{2\mathsf c_j}{\beta\bar E_\rho^2}\,X_i .$$
\emph{The mean.} Jensen's inequality for the conditional means, the independence of the slabs and H\"older's inequality bound $\langle\exp\sum_i\frac{2\mathsf c_j}{\beta\bar E_\rho^2}X_i\rangle_K$ by $\mathsf M_\beta^\varpi$, $\mathsf M_\beta=\sup_p\langle e^{\beta a_p}\rangle_K$, $\varpi=40\rho^4\mathsf c_jn_I/(\beta\bar E_\rho)^2\le1$. And $\log\mathsf M_\beta\le192\mathsf b/\ell+4(14.15+3\log\beta)$: the Schwarz inequality at site mirrors, used $3\mathsf k$ times, gives $\mathsf M_\beta\le\langle\prod_qe^{\beta a_q}\rangle_K^{\mathsf w^{-3}}$, the product over the $\mathsf w^3$ plaquettes $q$ parallel to one plaquette in $[0,\mathsf w]^2\times\{0,2,\dots,2\mathsf w-2\}$, $\mathsf w=2^{\mathsf k}\in(\frac\ell{4a_K},\frac\ell{2a_K}]$; and the equation of the Wilson measure for the bonds of the box $[0,\mathsf w]^2\times[0,2\mathsf w-2]$ bounds the conditional mean of that product by $e^{2\beta|\mathsf P_\partial|}$, $|\mathsf P_\partial|\le24\mathsf w^2$ the number of the plaquettes that stick out of the box, divided by the partition function of the box with free boundary, which is $\ge(1.4\cdot10^6\beta^3)^{-(\mathsf w+1)^2(2\mathsf w-1)}$ (a spanning tree of the box, and a ball of radius $(4\beta)^{-1/2}$ about $\mathbf 1$ for the other bonds).

As $2\mathsf c_j/\bar E_\rho\ge\frac{\mathsf c_j}\pi(\log\rho-1.2144)$, $n_I/\beta\ge|C|/\mathsf b-4\rho\mathsf p_C/\beta$, $\rho\ge\frac12(\beta/\Lambda_C\log^3\beta)^{1/4}$ and $\varpi\log\mathsf M_\beta$ stays bounded, the product of the $\|Q_i\|$ gives \eqref{C-363-42}.
\emph{The law.} $\chi_{l/2}(g)=\mathsf U_l(\frac12{\rm tr}\,g)$ with the Chebyshev polynomials $\mathsf U_l$ of the second kind, $\int\chi_j(g)\,dg=0$ for $j>0$, and the law of $\frac12{\rm tr}\,g$ under Haar measure is $\frac2\pi\sqrt{1-w^2}\,dw$. For a polynomial $h=\sum_l\hat h_l\mathsf U_l$, $\langle F_Kh(W_C)\rangle_K-\bar h\langle F_K\rangle_K=\sum_{l\ge1}\hat h_l\langle F_K\chi_{l/2}(U(C))\rangle_K\to0$ by \eqref{C-363-42}, uniformly in $F_K$; the polynomials are dense in $C[-1,1]$, and $|W_C|\le1$.
\end{proof}

\begin{lemma}[From slabs to linear blocks]\label{C-363-20}
Let $H\ge0$ be self-adjoint, $V$ a real vector space with a seminorm $|\cdot|_V$, and $\mathsf B_h(s,s')$, $h\in V$, $s,s'>0$, bounded operators, linear in $h$, with $e^{-tH}\mathsf B_h(s,s')e^{-t'H}=\mathsf B_h(s+t,s'+t')$, $\mathsf B_h(s,s')^*=\mathsf B_h(s',s)$, and $\|\int\rho_e(\varsigma/\eta)\mathsf B_h(\eta+\varsigma,\eta-\varsigma)d\varsigma\|\le C_e|h|_V$ for $\eta=\frac1{4(k+2)}$, $k\in\mathbb N_0$, where $\rho_e\ge0$ is measurable on $[-c_1,c_1]$, $c_1\le\frac12$, $\bar\rho_e=\int\rho_e\in(0,\infty)$. Then: (a) there is exactly one family of hermitian forms $q_h$ on the vectors of bounded energy with $q_h(e^{-sH}\chi,e^{-s'H}\chi')=\langle\chi,\mathsf B_h(s,s')\chi'\rangle$; (b) $|q_h(P_m\psi,P_n\psi')|\le C_\flat|h|_V(m+n+2)\|P_m\psi\|\|P_n\psi'\|$ with $C_\flat=8C_e/\bar\rho_e$; (c) for $t\mapsto f_t\in V$ three times continuously differentiable for $|\cdot|_V$ and of compact support, $\nu_j(f)=\int|\partial_t^jf_t|_Vdt$ and $\nu^*(f)=\sum_{j\le3}\binom3j2^{3-j}\nu_j(f)$, the time average $\mathsf e(f)(\psi,\psi')=\int q_{f_t}(e^{-itH}\psi,e^{-itH}\psi')dt$ has $|\mathsf e(f)(\psi,\psi)|\le c_e(f)\langle\psi,(1+H)\psi\rangle$, $c_e(f)=2C_\flat[\nu_0(f)+2(\zeta(2)+\zeta(3))\nu^*(f)]$, and extends to the form domain of $H$ with the same bound; (d) if, for data (Rec), the forms $\mathsf e$ of such a family with $V=C^1_c(\mathbb R^2)$ and the field satisfy the identity of Theorem~\ref{C-363-3}(f5), then (LB$_{f,c(f)}$) holds with $c(f)=|\mathsf c|[(1+2\coth^2s)c_e(f)+\frac{e^s}{2\sinh^2s}\sum_{j,\pm}c_e(f\circ B_j(\pm s))]+|\mathsf c'||\int f|$.

For the free massless field of Theorem~\ref{C-363-6}(d), with $T_{\mu\nu}={:}\partial_\mu\varphi\,\partial_\nu\varphi{:}-\frac12\eta_{\mu\nu}{:}\partial_\lambda\varphi\,\partial^\lambda\varphi{:}$ and $\eta={\rm diag}(1,-1,-1)$, one has $T_{00}=\frac12{:}(\dot\varphi^2+|\nabla\varphi|^2){:}$ and $\Theta=-2\eta^{\mu\nu}T_{\mu\nu}$, so that the field of the datum is $-\Theta=2\eta^{\mu\nu}T_{\mu\nu}$; the identity holds with $\mathsf e$ the form of $T_{00}$, $\mathsf c=2$ and $\mathsf c'=0$, because every symmetric tensor satisfies $\eta^{\mu\nu}T_{\mu\nu}=(1+2\coth^2s)T_{00}-\frac1{2\sinh^2s}\sum_{j=1,2}\sum_\pm T_{uu}$, $u=\cosh s\,e_0\pm\sinh s\,e_j$.
\end{lemma}
\begin{proof}[Proof sketch]
(a) Two pairs of times are compared through a pair larger than both. (b) Let $X$ be the operator $P_n\mathcal H\to P_m\mathcal H$ of $q_h$, and write $H=m+\delta_m$ on $P_m\mathcal H$, $0\le\delta_m<1$. Then $P_m\mathsf B_h(\eta+\varsigma,\eta-\varsigma)P_n=e^{-(\eta+\varsigma)m-(\eta-\varsigma)n}e^{-(\eta+\varsigma)\delta_m}Xe^{-(\eta-\varsigma)\delta_n}$ and $\|e^{-\alpha\delta_m}Xe^{-\alpha'\delta_n}-X\|\le(\alpha+\alpha')\|X\|$; with $\kappa=\int\rho_e(\varsigma/\eta)e^{-(\eta+\varsigma)m-(\eta-\varsigma)n}d\varsigma$ the slab operator $\mathsf S_\eta$ has $\|P_m\mathsf S_\eta P_n-\kappa X\|\le2\eta\kappa\|X\|$, so $\kappa(1-2\eta)\|X\|\le C_e|h|_V$; at $\eta=\frac1{4(m+n+2)}$, $\kappa\ge\eta\bar\rho_ee^{-3/8}$, and $\frac{16}3e^{3/8}<8$. (c) The block $(m,n)$ of $\mathsf e(f)$ is $\int e^{it(m-n)}e^{it\delta_m}X_{mn}(f_t)e^{-it\delta_n}dt$; three integrations by parts for $m\ne n$ give the decay $|m-n|^{-3}\nu^*(f)$ times the bound (b); with $A+B\le2\sqrt{AB}(1+|A-B|)$, $A=m+1$, $B=n+1$, Schur's test gives the form bound. (d) In a boosted frame $U(\Lambda)HU(\Lambda)^{-1}\le e^sH$ on the spectrum, so each of the five terms is bounded by its constant times $e^s\langle\psi,N\psi\rangle$.

The free identity: $T_{uu}=\cosh^2s\,T_{00}\pm2\cosh s\sinh s\,T_{0j}+\sinh^2s\,T_{jj}$; the sum over the signs and over $j$ is $4\cosh^2s\,T_{00}+2\sinh^2s(T_{11}+T_{22})$, and the combination is $T_{00}-T_{11}-T_{22}=\eta^{\mu\nu}T_{\mu\nu}$. Since $\eta^{\mu\nu}\eta_{\mu\nu}=3$, $\eta^{\mu\nu}T_{\mu\nu}=\Theta-\frac32\Theta=-\frac12\Theta$, and $V\phi(f)V^{-1}=-\Theta(f)$ (Theorem~\ref{C-363-6}(d)) gives $\mathsf c=2$; by covariance $\mathsf e(f\circ B_j(\pm s))(\psi_{j,\pm},\psi_{j,\pm})$ is the form of $T_{uu}(f)$ at $\psi$.
\end{proof}

\begin{lemma}[Locality through a dilation]\label{C-363-21}
(a) Let $\mathcal R_0$ be a $*$-algebra of bounded operators on a Hilbert space $\mathcal F$, $\Omega_{\rm F}$ cyclic for it, $P$ a projection with $P\Omega_{\rm F}=\Omega_{\rm F}$, and $\mathcal K=P\mathcal F$. Every bounded $T$ on $\mathcal K$ that commutes with all $PAP|_{\mathcal K}$, $A\in\mathcal R_0$, is $PT_1P|_{\mathcal K}$ with $T_1\in\mathcal R_0'$, $\|T_1\|\le\|T\|$. (b) For data (Rec) and non-empty open $O_1,O_2$ the following say the same: $\mathcal C(O_1)'\subset\mathcal C(O_2)$; there are a Hilbert space $\mathcal F$, an isometry $V:\mathcal H\to\mathcal F$ and a $*$-algebra $\mathcal R_0\subset\mathcal B(\mathcal F)$ with $\mathcal R_0V\Omega$ dense, $V^*\mathcal R_0V\subset\mathcal C(O_1)$ and $V^*\mathcal R_0'V\subset\mathcal C(O_2)$.
\end{lemma}
\begin{proof}[Proof sketch]
(a) With $T$ also $T^*T$ commutes with $X=PA^*AP|_{\mathcal K}\ge0$, so $\|AT\Omega_{\rm F}\|^2=\langle T\Omega_{\rm F},XT\Omega_{\rm F}\rangle=\langle X^{1/2}\Omega_{\rm F},T^*TX^{1/2}\Omega_{\rm F}\rangle\le\|T\|^2\|A\Omega_{\rm F}\|^2$; so $T_1A\Omega_{\rm F}=AT\Omega_{\rm F}$ defines a bounded $T_1$, and $T_1\in\mathcal R_0'$. For $\xi\in\mathcal K$ and $A_i\Omega_{\rm F}\to\xi$: $PT_1A_i\Omega_{\rm F}=(PA_iP)T\Omega_{\rm F}=T(PA_iP)\Omega_{\rm F}\to T\xi$, so $PT_1P|_{\mathcal K}=T$. (b) If the dilation exists and $T\in\mathcal C(O_1)'$, then $VTV^*$ commutes on $V\mathcal H$ with the compressions of $\mathcal R_0$, so by (a) $T=V^*T_1V$ with $T_1\in\mathcal R_0'$, and $T\in\mathcal C(O_2)$. Conversely take $\mathcal F=\mathcal H$, $V=1$, $\mathcal R_0=\mathcal C(O_1)$; it holds $\mathcal A(O_2)=\mathcal C(O_2)'\subset\mathcal C(O_1)$, and $\mathcal A(O_2)\Omega$ is dense.
\end{proof}

\begin{lemma}[Weyl operators and Wick forms]\label{C-363-22}
Let $\mathfrak h=L^2(\mathbb R^2)$, $\omega=(p^2+\mathsf m^2)^{1/2}$ with $\mathsf m>0$, $\mathcal F$ the Fock space over $\mathfrak h$, $H_{\rm F}=d\Gamma(\omega)$, $\kappa^x=(2\pi)^{-1}(2\omega)^{-1/2}e^{i(t\omega-\vec x\cdot p)}$, $u_h=\int h(x)\kappa^xdx$, $A(u)=a^*(u)+a(u)$, $\mathsf W[u]=e^{i\overline{A(u)}}$, $\mathcal R_0(O)$ the span of the $\mathsf W[u_h]$, $h\in\mathcal F_{\mathbb R}(O)$, and $\Phi_f$ the form of the smeared Wick power ${:}\chi^k{:}(f)$ of the free field $\chi(h)=A(u_h)$. (a) For $u\in\bigcap_rD(\omega^r)$ and every $n$: $\mathsf W[u]D(H_{\rm F}^n)=D(H_{\rm F}^n)$ and $(1+H_{\rm F})^n\mathsf W[u]=\mathsf W[u](1+H_{\rm F}+A(i\omega u)+\langle u,\omega u\rangle)^n$ there. (b) $\Phi_f(\zeta,\zeta')=\sum_j\binom kj\int f(x)\langle a(\kappa^x)^j\zeta,a(\kappa^x)^{k-j}\zeta'\rangle dx$ on $\mathcal D_\infty(H_{\rm F})$, with $|\Phi_f(\zeta,\zeta')|\le2^k(4\pi\mathsf m)^{-k/2}\|f\|_{L^1}\|(1+H_{\rm F})^k\zeta\|\|(1+H_{\rm F})^k\zeta'\|$; it is the form of an operator $\Phi(f)$ on $\mathcal D_\infty(H_{\rm F})$. (c) For real $h\in C^\infty_c$ and $x$ spacelike to ${\rm supp}\,h$, $\langle\kappa^x,u_h\rangle=\int h(y)\frac{e^{-\mathsf m\mathsf d_{xy}}}{4\pi\mathsf d_{xy}}dy$, $\mathsf d_{xy}=(|\vec x-\vec y|^2-(x^0-y^0)^2)^{1/2}$: a real number; so Weyl operators of spacelike separated regions commute, and $\Phi_f(\mathsf W[u_h]\zeta,\mathsf W[u_h]\zeta')=\Phi_f(\zeta,\zeta')$ if ${\rm supp}f$ and ${\rm supp}\,h$ are spacelike separated. (d) For every non-empty open $O$ the vectors $\mathsf W[u_h]\Omega_{\rm F}$, $h\in\mathcal F_{\mathbb R}(O)$, are total in $\mathcal F$. (e) For $f\in\mathcal F(O)$ and $B\in\mathcal R_0(O)'$: $\langle\Phi(\bar f)\zeta,B\zeta'\rangle=\langle\zeta,B\Phi(f)\zeta'\rangle$ for $\zeta,\zeta'\in\mathcal D_\infty(H_{\rm F})$. (f) \emph{Gaussian families.} Let $(\kappa_\alpha,\varepsilon_\alpha,\mathsf P)$, $\alpha\in\mathsf A$, be a Gaussian family of Theorem~\ref{C-363-6}(a) on $(\mathsf M,\mu)$, and now let $\mathfrak h=L^2(\mathsf M,\mu)$, $\mathcal F$ the Fock space over it, $H_{\rm F}=d\Gamma(\omega)$, $\kappa^x_\alpha=e^{i(t\omega-\vec x\cdot\vec p)}\kappa_\alpha$, and $\Phi_f$ the form of the smeared Wick polynomial ${:}\mathsf P(\Phi){:}(f)$.

Put $\varsigma_\alpha=1$ if $\varepsilon_\alpha=+1$ and $\varsigma_\alpha=i$ if $\varepsilon_\alpha=-1$, $u^\alpha_h=\varsigma_\alpha\hat h(\omega,\vec p)\kappa_\alpha$ for $h\in\mathcal S$, let $\mathfrak k(O)$ be the real span of the $u^\alpha_h$, $h\in\mathcal F_{\mathbb R}(O)$, $\alpha\in\mathsf A$, and let $\mathcal R_0(O)$ now be the span of the $\mathsf W[u]$, $u\in\mathfrak k(O)$. Suppose \textup{(f1)} $\mathfrak k(O)+i\,\mathfrak k(O)$ is dense in $\mathfrak h$ for every non-empty open $O$, and \textup{(f2)} $\overline{\langle\kappa^x_\alpha,u\rangle}=\varepsilon_\alpha\langle\kappa^x_\alpha,u\rangle$ for $\alpha\in\mathsf A$, $u\in\mathfrak k(O)$ and $x$ spacelike to $O$. Then (a)--(e) hold with the $u\in\mathfrak k(O)$ in the place of the $u_h$, $h\in\mathcal F_{\mathbb R}(O)$, and with three changes.

In (a), $u\in\bigcap_rD(\omega^r)$ with $\omega^{-1/2}u\in\mathfrak h$, which holds for every $u\in\mathfrak k(O)$. In (b), the form is
$$\Phi_f(\zeta,\zeta')=\sum_j\binom kj\sum_{\underline\alpha}\mathsf p_{\underline\alpha}\prod_{i>j}\varepsilon_{\alpha_i}\int f(x)\Bigl\langle\prod_{i\le j}a(\kappa^x_{\alpha_i})\zeta,\prod_{i>j}a(\kappa^x_{\alpha_i})\zeta'\Bigr\rangle dx ,$$
and with the weight $\varpi=\omega^{1/2}(1+\omega)^{r'}$, $r'>\delta$, the number $\mathsf n_*=\max_\alpha\|\varpi^{-1}\kappa_\alpha\|<\infty$ and the order $L_*=\lceil k(r'+\frac12)\rceil$ it obeys $|\Phi_f(\zeta,\zeta')|\le2^k|\mathsf p|_1\mathsf n_*^k\|f\|_{L^1}\|(1+H_{\rm F})^{L_*}\zeta\|\|(1+H_{\rm F})^{L_*}\zeta'\|$, the operator $\Phi(f)$ being bounded by a multiple of $\|(1+H_{\rm F})^{2L_*}\cdot\|$. In (c), the invariance holds for $u\in\mathfrak k(O)$ and ${\rm supp}f$ spacelike to $O$, and Weyl operators of $\mathfrak k(O_1)$, $\mathfrak k(O_2)$ commute for spacelike separated $O_1,O_2$. The hypotheses \textup{(f1)}, \textup{(f2)} hold for the first derivatives $\partial_\alpha\varphi^a$ of $\mathsf n$ massless free fields, the free photons with $\mathsf n$ colours of Theorem~\ref{C-363-6}(c)--(e), and for the generalised free fields of Theorem~\ref{C-363-6}(c), the free fields of a mixture $d\rho(\mathsf m)$ of masses with $\varepsilon=+$, the kernel $r^{-2}$ being the mixture $d\rho=4\pi\,d\mathsf m$.
\end{lemma}
\begin{proof}[Proof sketch]
(a) On exponential vectors $\mathsf W[u]{\rm Exp}(v)=e^{-\frac12\|u\|^2+i\langle u,v\rangle}{\rm Exp}(v+iu)$, whence $H_{\rm F}\mathsf W[u]=\mathsf W[u](H_{\rm F}+A(i\omega u)+\langle u,\omega u\rangle)$ there; exponential vectors are dense in every $D(H_{\rm F}^n)$, $H_{\rm F}$ is closed, and an induction gives the domains. (b) Wick's theorem, with $\|a(\kappa^x)^j\zeta\|\le\|\omega^{-1}\kappa^x\|^j\|H_{\rm F}^j\zeta\|$ and $\|\omega^{-1}\kappa^x\|^2=1/4\pi\mathsf m$; the operator by Riesz after the commutator identity $\Phi_f(H_{\rm F}\zeta,\zeta')-\Phi_f(\zeta,H_{\rm F}\zeta')=i\Phi_{\partial_tf}(\zeta,\zeta')$. (c) The Euclidean two-point function $e^{-\mathsf m r}/4\pi r$ continued to real time is real at spacelike separation; on exponential vectors the conjugation by $\mathsf W[u_h]$ shifts each annihilator $a(\kappa^x)$ by $i\langle\kappa^x,u_h\rangle$ in the ket and conjugately in the bra, and the reality of the shift makes the two changes cancel. (d) $\{u_h:h\in\mathcal F(O)\}$, $h$ complex, is dense in $\mathfrak h$ by the reflection principle and uniqueness of Fourier transforms: for $g\perp u_h$, $h\in\mathcal F(O)$, and $\eta$ carried by a small ball, $x\mapsto\langle g,u_{\eta(\cdot-x)}\rangle=\int\bar g\kappa\hat\eta(\omega,p)e^{i(t\omega-\vec x\cdot p)}d^2p$ vanishes near a point and continues boundedly into every timelike direction, so it vanishes, and its Fourier transform at $t=0$ gives $\bar g\kappa\hat\eta=0$; along an approximate unit $\hat\eta\to1$, and $\kappa>0$, so $g=0$.

The vectors of (d) are $\mathsf W[u_h]\Omega_{\rm F}=e^{-\frac12\|u_h\|^2}{\rm Exp}(iu_h)$ with real $h$. Let $\Psi$ be orthogonal to all of them. For real $h_1,\dots,h_n\in\mathcal F_{\mathbb R}(O)$ the function $z\mapsto\langle\Psi,{\rm Exp}(\sum_jz_ju_{h_j})\rangle$ is entire on $\mathbb C^n$ and vanishes on $i\mathbb R^n$, since $\sum_jiy_ju_{h_j}=iu_{\sum_jy_jh_j}$; so all its Taylor coefficients at $0$ vanish: $\Psi\perp\Omega_{\rm F}$ and $\Psi\perp{\rm Sym}(u_{h_1}\otimes\cdots\otimes u_{h_n})$. This last product is complex-linear in each $u_{h_j}$, and $h\mapsto u_h$ is complex-linear; so $\Psi$ is orthogonal to it for all $h_j\in\mathcal F(O)$, and these symmetric tensors are total in each $n$-particle space by the density; hence $\Psi=0$. (e) The Wick power of one smeared field is a real polynomial in it: ${:}A(u)^k{:}=\mathsf h_k(A(u))$ with $\mathsf h_0=1$, $\mathsf h_1=\lambda$, $\mathsf h_{k+1}=\lambda\mathsf h_k-k\|u\|^2\mathsf h_{k-1}$; for $u=u_{\delta_\epsilon(\cdot-x)}$, $\delta_\epsilon$ an approximate unit and $x\in{\rm supp}f$, the unitary group of $\overline{A(u)}$ lies in $\mathcal R_0(O)$, so $B$ commutes weakly with ${:}A(u)^k{:}$ on $\mathcal D_\infty(H_{\rm F})$; integrate against $f$ and let $\epsilon\downarrow0$, the regularised Wick powers being bounded by $c(f)\|(1+H_{\rm F})^{2k}\cdot\|$ uniformly in $\epsilon$ and converging weakly. (f) The changes, clause by clause. (a): $\|a(w)\zeta\|\le\|\omega^{-1/2}w\|\|H_{\rm F}^{1/2}\zeta\|$ and $\|a^*(w)\zeta\|^2=\|a(w)\zeta\|^2+\|w\|^2\|\zeta\|^2$ on each $n$-particle space take the place of the bound by the number of particles, which uses $\mathsf m>0$; so $A(u)$ is defined on $D(H_{\rm F}^{1/2})$ for $u,\omega^{-1/2}u\in\mathfrak h$, it is essentially self-adjoint on the vectors of finitely many particles in $\mathcal D_\infty(H_{\rm F})$, and the proof of (a) runs as it stands.

Every $u^\alpha_h$ lies in all $D(\omega^r)$ and has $\omega^{-1/2}u^\alpha_h\in\mathfrak h$, because $\hat h$ decays faster than every power and $\varpi^{-1}\kappa_\alpha\in\mathfrak h$. (b): in the bound of the annihilators the weight $\varpi$ takes the place of $\omega$, 
\[
\|a(g_1)\cdots a(g_m)\zeta\|\le\prod_i\|\varpi^{-1}g_i\|\|(1+H_{\rm F})^{m(r'+\frac12)}\zeta\|,
\]
and 
\[
\|\varpi^{-1}\kappa_\alpha\|^2=\Gamma(2r')^{-1}\int_0^\infty s^{2r'-1}e^{-s}\|\omega^{-1/2}\kappa_\alpha e^{-s\omega/2}\|^2ds<\infty
\]
 by the bound $c_\kappa s^{-\delta}$ of Theorem~\ref{C-363-6}(a) and $r'>\delta$; the form is hermitian, $\overline{\Phi_{\bar f}(\zeta',\zeta)}=\Phi_f(\zeta,\zeta')$, because $\mathsf P(\varepsilon t)=\mathsf P(t)$; the operator comes as before, from the commutator with $H_{\rm F}$. (c): on exponential vectors 
\[
\Phi_f({\rm Exp}\,v,{\rm Exp}\,v')=e^{\langle v,v'\rangle}\int f(x)\,\mathsf P\bigl(\bar z(x)+\varepsilon z'(x)\bigr)dx,
\]
$z_\alpha=\langle\kappa^x_\alpha,v\rangle$, $z'_\alpha=\langle\kappa^x_\alpha,v'\rangle$; under $v\to v+iu$, $v'\to v'+iu$ the argument of $\mathsf P$ changes by $-i\bar\zeta_\alpha+i\varepsilon_\alpha\zeta_\alpha$, $\zeta_\alpha=\langle\kappa^x_\alpha,u\rangle$, which \textup{(f2)} makes $0$; and $\langle u^\alpha_{h'},u^\beta_h\rangle=\bar\varsigma_\alpha\int h'(x)\langle\kappa^x_\alpha,u^\beta_h\rangle dx$ is real by \textup{(f2)}, so the Weyl operators commute. (d): the argument of (d) with the $u_j\in\mathfrak k(O)$ in the place of the $u_{h_j}$, and \textup{(f1)} in the place of the density of the $u_h$. (e): with $g_\alpha=\rho_\epsilon\kappa^x_\alpha$, $\rho_\epsilon=\hat\delta_\epsilon(\omega,\vec p)$, the regularised fields $\Phi_\alpha=a^*(g_\alpha)+\varepsilon_\alpha a(g_\alpha)$ are $\bar\varsigma_\alpha A(\varsigma_\alpha g_\alpha)$, as $\bar\varsigma_\alpha^2=\varepsilon_\alpha$, and $\varsigma_\alpha g_\alpha=u^\alpha_{\delta_\epsilon(\cdot-x)}\in\mathfrak k(O)$ for $x\in{\rm supp}f$ and small $\epsilon$, so the unitary groups of the closures of the $A(\varsigma_\alpha g_\alpha)$ lie in $\mathcal R_0(O)$.

The recursion 
\[
{\Phi_{\alpha_0}}{:}\Phi_{\underline\alpha}{:}={:}\Phi_{(\alpha_0,\underline\alpha)}{:}+\sum_i\varepsilon_{\alpha_0}\langle g_{\alpha_0},g_{\alpha_i}\rangle{:}\Phi_{\underline\alpha\setminus i}{:}
\]
 makes every Wick monomial a polynomial with complex coefficients in several non-commuting $A(\varsigma_\alpha g_\alpha)$; $B$ passes through one factor at a time, so $\langle\zeta,BZ\zeta'\rangle=\langle Z^\dagger\zeta,B\zeta'\rangle$ for such a polynomial $Z$ and $Z^\dagger$ the polynomial with the conjugate coefficients and the products reversed; and $Z^\dagger=Z$ on $\mathcal D_\infty(H_{\rm F})$ for the regularised Wick polynomial $Z=\sum_{\underline\alpha}\mathsf p_{\underline\alpha}{:}\Phi_{\underline\alpha}{:}$, by the hermiticity of (b) at $B=1$.

Then integrate against $f$ and let $\epsilon\downarrow0$ as before, with $\varpi$ for $\omega$ in the weights and $L_*$ for $k$ in the norms. \emph{The families.} For the photons, $\kappa_0=(2\pi)^{-1}(2|p|)^{-1/2}$, $\kappa_j=-ip_j\kappa_0$, $\kappa_3=-\omega\kappa_0$ and $\varepsilon=(+,+,-)$ (the proof of Theorem~\ref{C-363-6}(d)), so $u^j_h=\hat h'(\omega,p)\kappa_0$ with $h'=-\partial_jh$, and $u^3_h=\hat h'(\omega,p)\kappa_0$ with $h'=\partial_th$. For real $h'\in C^\infty_c$ the function $F_{h'}(x)=\langle\kappa^x_0,\hat h'(\omega,p)\kappa_0\rangle$ is smooth and, by the computation of (c) at $\mathsf m=0$, real on the open set of the $x$ spacelike to ${\rm supp}\,h'$; there $\langle\kappa^x_j,\hat h'\kappa_0\rangle=\partial_jF_{h'}(x)$ is real and $\langle\kappa^x_3,\hat h'\kappa_0\rangle=-i\partial_tF_{h'}(x)$ is imaginary: \textup{(f2)}. $\mathfrak k(O)+i\,\mathfrak k(O)$ holds the vectors $\hat h(\omega,p)\,\omega\kappa_0$, $h\in\mathcal F(O)$, and the reflection argument of (d) with $\omega\kappa_0>0$ in the place of $\kappa$, $\omega-\vec v\cdot p\ge(1-|\vec v|)|p|>0$, gives \textup{(f1)}.

For $\mathsf n$ colours $\mathfrak h$ is the sum of $\mathsf n$ copies, $\langle\kappa^x_{(a,\alpha)},u^{(b,\beta)}_h\rangle=\delta_{ab}\langle\kappa^x_\alpha,u^\beta_h\rangle$, and both hypotheses hold copy by copy. For a mixture of masses, $\mathsf M=\mathbb R^2\times(0,\infty)$, $d\mu=d^2p\,d\rho(\mathsf m)$, $\omega=(p^2+\mathsf m^2)^{1/2}$ and one field: $\langle\kappa^x,u_h\rangle=\int h(y)\int\frac{e^{-\mathsf m\mathsf d_{xy}}}{4\pi\mathsf d_{xy}}d\rho(\mathsf m)\,dy$ is real at spacelike separation, by (c) and a $\rho$-integrable majorant, and the reflection argument, applied to the image of the finite measure $\bar g\kappa\hat\eta(\omega,\vec p)\,\mu$ under $(\vec p,\mathsf m)\mapsto(\omega,\vec p)$, a homeomorphism onto $\{E>|\vec p|\}$, gives the density; only the weight changes.
\end{proof}

\begin{proposition}[Four squares]\label{C-363-14}
Let $\Lambda=a\mathbb Z^3/\ell\mathbb Z^3$, $\ell/a$ even and $\ge4$, $\beta\ge0$, $\mu=\mu_\beta$ with the normalisation $Z$. The planes $\mathsf P_\mu=\{x_\mu\in\{0,\frac\ell2\}\}$, $\mu=1,3$, pass through sites and cut the torus into $\mathsf Q_1=\{0<x_1,x_3<\frac\ell2\}$, $\mathsf Q_2=\theta_1\mathsf Q_1$, $\mathsf Q_3=\theta\theta_1\mathsf Q_1$, $\mathsf Q_4=\theta\mathsf Q_1$; giving each plaquette in a plane half of its weight in each of the two adjacent closed quadrants, $e^{-\beta\sum_pa_p}=w_1w_2w_3w_4$, $w_i$ depending only on the bonds of $\bar{\mathsf Q}_i$. Let $\mathsf B_c$ be the bonds on the four corner lines $\mathsf P_1\cap\mathsf P_3$; $\mathsf B_R$ the other bonds of $\mathsf P_3\cap\{0\le x_1\le\frac\ell2\}$, in both planes $x_3\in\{0,\frac\ell2\}$; $\mathsf B_L=\theta_1\mathsf B_R$; $\mathsf B_U=\mathsf d\mathsf B_R$, the bonds of $\mathsf P_1\cap\{0\le x_3\le\frac\ell2\}$ off the corner lines; $\mathsf X=G^{\mathsf B_R}$, $\mathsf C=G^{\mathsf B_c}$, $\kappa$ the map of $\mathsf C$ induced by $\mathsf d$, $\mathcal K=L^2(\mathsf X)$, and $\mathcal H_\Lambda=L^2(\mathsf C\times\mathsf X\times\mathsf X)$, regarded as the space of the fields $c\mapsto\hat\xi_c$ of Hilbert--Schmidt operators on $\mathcal K$ whose left index is the $\theta_1$-image $x_L$ and whose right index is $x_R$.

For a bounded Borel function $F$ of the bonds of $\bar{\mathsf Q}_1$ let $\mathsf A^F_c$ be the Hilbert--Schmidt operator on $\mathcal K$ with the kernel $Z^{-1/4}\int Fw_1$, integrated over the bonds of $\bar{\mathsf Q}_1$ in no plane, between the configuration $x$ of $\mathsf B_R$ and the configuration $x'$ of $\mathsf B_U$ carried to $\mathsf B_R$ by $\mathsf d$, at the corner configuration $c$; $Z_c=\mathsf A^1_c$. For a bounded Borel $F$ of the bonds of $\{0\le x_3\le\frac\ell2\}$ let $\iota_\Lambda F=Z^{-1/2}\int Fw_1w_2$ over the bonds of that half off $\mathsf B_c\cup\mathsf B_R\cup\mathsf B_L$, $\Omega_\Lambda=\iota_\Lambda1$; and $\mathcal M_\Lambda=\{\hat\xi_c\mapsto\hat\xi_cq_c\}$, $\mathcal M'_\Lambda=\{\hat\xi_c\mapsto b_c\hat\xi_c\}$ with bounded measurable fields $q_c,b_c$ of operators of $\mathcal K$.

Then:
\begin{enumerate}
\item[(L1)] $\langle\iota_\Lambda F',\iota_\Lambda F\rangle=\mu(\theta\bar F'\cdot F)$; $J_\Lambda\hat\xi_c=\hat\xi^*_c$ has $J_\Lambda\iota_\Lambda F=\iota_\Lambda(\theta_1\bar F)$ and $J_\Lambda\mathcal M_\Lambda J_\Lambda=\mathcal M'_\Lambda$; $\mathcal M_\Lambda$ and $\mathcal M'_\Lambda$ are von Neumann algebras, each the commutant of the other.
\item[(L2)] For $F_1$ of $\bar{\mathsf Q}_1$ and $F_2$ of $\bar{\mathsf Q}_2$: $(\iota_\Lambda(F_2F_1))^\wedge_c=\mathsf A^{R^{-1}F_2}_{\kappa c}\mathsf A^{F_1}_c$; $(\mathsf A^F_c)^*=\mathsf A^{\mathsf d\bar F}_{\kappa c}$, $Z^*_c=Z_{\kappa c}$, and $\hat\Omega_c=Z_c^*Z_c=\rho_c^{1/2}$ with $\rho_c$ of trace class, $\int{\rm Tr}\,\rho_c\,dc=1$.
\item[(L3)] For $Q\in\mathcal M_\Lambda$ there is $\tilde q\in\mathcal H_\Lambda$, $\|\tilde q\|\le\|Q\|$, with $\langle\iota_\Lambda F',Q\iota_\Lambda F\rangle=\langle\tilde q,\iota_\Lambda(R^{-1}(\theta\bar F'\cdot F))\rangle$ for $F,F'$ of the bonds of $\bar{\mathsf Q}_2$.
\item[(L4)] Let $e_c$ be the projection on $(\ker Z_c)^\perp$, $e_R:\hat\xi_c\mapsto\hat\xi_ce_c$ in $\mathcal M_\Lambda$, $e_L=J_\Lambda e_RJ_\Lambda$, $\pi=e_Le_R$, and $e^{-2\pi sK_\Lambda}\hat\xi_c=\rho_c^{-s}\hat\xi_c\rho_c^s$ on $\pi\mathcal H_\Lambda$: $\Delta_\Lambda=e^{-2\pi K_\Lambda}$ is the modular operator of $(\mathcal M_\Lambda,\Omega_\Lambda)$. There is a self-adjoint involution $\Sigma$ of $\pi\mathcal H_\Lambda$ commuting with $K_\Lambda$ such that $\mathsf P_\Lambda=\Sigma e^{-\frac\pi2K_\Lambda}$ is self-adjoint, $|\mathsf P_\Lambda|=e^{-\frac\pi2K_\Lambda}$, and $\pi\iota_\Lambda F\in D(\mathsf P_\Lambda)$, $\mathsf P_\Lambda\pi\iota_\Lambda F=\pi\iota_\Lambda(RF)$ for $F$ of $\bar{\mathsf Q}_1$: the quarter turn is the fourth root of the modular operator, up to a sign.
\item[(L5)] For such $F$ and $\xi=\mathsf F(K_\Lambda)\pi\iota_\Lambda F$ there is a closed, densely defined $T$ commuting with the unitaries of $\mathcal M'_\Lambda$, with $T\Omega_\Lambda=\xi$ and $\|T^*T\Omega_\Lambda\|\le\pi e^{5\pi^2/8}\|\iota_\Lambda(\theta_1\bar F\cdot F)\|=\pi e^{5\pi^2/8}\mu(\theta\theta_1F\cdot\theta\bar F\cdot\theta_1\bar F\cdot F)^{1/2}$; so $T^{(\epsilon)}=T(1+\epsilon T^*T)^{-1}\in\mathcal M_\Lambda$, $\|T^{(\epsilon)}\|\le(2\sqrt\epsilon)^{-1}$, and $\|T^{(\epsilon)}\Omega_\Lambda-\xi\|\le\frac{\sqrt\epsilon}2\|T^*T\Omega_\Lambda\|$.
\end{enumerate}
No coupling enters, and no injectivity of $\rho_c$ is assumed: gauge invariance destroys it.
\end{proposition}
\begin{proof}[Proof sketch]
The maps $\theta,\theta_1,\mathsf d$ keep $\Lambda$, the Haar measure and the Wilson weight ($a_{\mathsf rp}(\mathsf rU)=a_p(U)$, as ${\rm tr}\,g^{-1}={\rm tr}\,g$ on $\mathrm{SU}(2)$). (L1) $\mu(\theta\bar F'\cdot F)=Z^{-1}\int\overline{F'(\theta U)}F(U)w_1w_2(U)w_1w_2(\theta U)dU$; at fixed $(c,x_L,x_R)$, which $\theta$ fixes, the integrand is a function of the upper bonds times one of the lower bonds, and the two integrals are $Z^{1/2}\iota_\Lambda F$ and $Z^{1/2}\overline{\iota_\Lambda F'}$; $\theta_1$ exchanges $w_1,w_2$ and $x_L,x_R$; $\mathcal H_\Lambda=L^2(\mathsf C)\otimes\mathcal K\otimes\mathcal K$ and $\mathcal M_\Lambda=L^\infty(\mathsf C)\otimes1\otimes\mathcal B(\mathcal K)$. (L2) At $F=F_2F_1$ the integral over the bonds of $\mathsf Q_1\cup\mathsf Q_2$ factorises at fixed configuration of $\mathsf B_U$; the substitution $U\to\theta_1U$ in the factor of $\mathsf Q_2$ and the diagonal reflection $\mathsf d$, which gives $(\mathsf A^F_c)^{\rm T}=\mathsf A^{\mathsf dF}_{\kappa c}$, yield the product formula; the kernel of $Z_c$ is real. (L3) With $E=R^{-1}F$, $E'=R^{-1}F'$ of $\bar{\mathsf Q}_1$: $\langle\iota_\Lambda F',Q\iota_\Lambda F\rangle=\int{\rm Tr}((Z_cq_c^*Z_c^*)^*(\iota_\Lambda(\theta_1\bar E'\cdot E))^\wedge_{\kappa c})dc$ and $\theta_1\bar E'=R^{-1}\theta\bar F'$; put $\hat{\tilde q}_{\kappa c}=Z_cq_c^*Z_c^*$; H\"older gives $\|Z_cq_c^*Z_c^*\|_2\le\|q_c\|\|Z_c\|_4^2$, and $\int\|Z_c\|_4^4dc=\|\Omega_\Lambda\|^2=1$. (L4) With the polar decomposition $Z_c=W_c\rho_c^{1/4}$, $W_{\kappa c}=W_c^*$; put $(\Sigma\xi)^\wedge_c=W_c^*\hat\xi_{\kappa c}W_c$, a symmetric involution of $\pi\mathcal H_\Lambda$ commuting with the modular flow; then $(\pi\iota_\Lambda F)^\wedge_c=\rho_c^{1/4}\mathsf c_c$ with $\mathsf c_c=W_c^*\mathsf A^F_ce_c=e_c\mathsf c_ce_c$, while $(\Sigma\pi\iota_\Lambda RF)^\wedge_c=\mathsf c_c\rho_c^{1/4}=(e^{-\frac\pi2K_\Lambda}\pi\iota_\Lambda F)^\wedge_c$. (L5) The Schur multipliers $h_\gamma(\omega)=\sqrt\pi e^{-\omega^2/16\pi^2+\gamma\omega}$ of the eigenvalue differences $\omega$ of $\log\rho_c$ have the norm $\sqrt\pi e^{4\pi^2\gamma^2}$ on every Schatten class; $\mathsf F(\omega/2\pi)=h_{1/8}(\omega)$, so $\hat\xi_c=\rho_c^{1/4}h_{1/8}(\mathsf c_c)$, and with $\hat\xi^\sharp_c=\rho_c^{1/4}h_{1/8}(\mathsf c^*_c)$, $\hat\zeta_c=h_{3/8}(\mathsf c_c)h_{1/8}(\mathsf c_c^*)$,
$$\rho_c^{1/2}\hat\xi_c^*=\hat\xi^\sharp_c\rho_c^{1/2},\qquad\hat\xi_c\hat\xi^*_c=\hat\zeta_c\rho_c^{1/2}.$$
The first makes $Q'\Omega_\Lambda+\eta\mapsto Q'\xi$, $Q'\in\mathcal M'_\Lambda$, $\eta\in(1-e_R)\mathcal H_\Lambda$, closable, with a closure $T$ that commutes with the unitaries of $\mathcal M'_\Lambda$; the second gives $T^*\xi=\zeta$ with $\|\hat\zeta_c\|_2\le\pi e^{5\pi^2/8}\|\mathsf A^F_c\|_4^2$ and $\|\mathsf A^F_c\|_4^4=\|(\mathsf A^F_c)^*\mathsf A^F_c\|_2^2=\|(\iota_\Lambda(\theta_1\bar F\cdot F))^\wedge_c\|_2^2$; (L1) turns $\|\iota_\Lambda(\theta_1\bar F\cdot F)\|$ into the four-point function. Von Neumann's theorem and the bicommutant theorem give $T^{(\epsilon)}\in\mathcal M_\Lambda$, and $T^{(\epsilon)}\Omega_\Lambda-\xi=-\epsilon T^{(\epsilon)}T^*T\Omega_\Lambda$.
\end{proof}

\paragraph{The one-lattice source recursion.}
We need more than the bound on the source energies on their original
source ball: the two-cell estimate uses their continuation to larger
source discs at the earlier levels. We give the induction that provides
this continuation, retaining the dependence on the background radius.
The local maps are those of Theorem~\YMMref{W}{423},
Proposition~\YMMref{51.16}{412}, Lemma~\YMMref{114.1}{800}
and Proposition~\YMMref{77.2}{632}.

Here is the notation used for those maps. On a domain $Y$ at level $j$,
$P$ is its background triple, $B$ is its fluctuation field, and
$|B|<\mathsf C_{\mathcal A}\mathfrak b(x_j)$ is the field polydisc.
For a family $V$ vanishing at $B=0$, its norm at a fixed background is
\[
 \|V\|_{{\rm f},P}
 =\sup_{Y,B}\frac{e^{\lambda_Vd_j(Y)}|V(Y;P,B)|}
                         {1+q_Y(B)},\qquad
 q_Y(B)=M^{-3}\sum_b|B(b)|^2.
\]
The norm $\mathsf n^{\rm f}(V)$ also takes the supremum over the
specified background cores. For a pointed family $E$ of level $j$,
$\|E\|_{j,\kappa+1}$ is the supremum of
$e^{(\kappa+1)d_j(X)}|E(X,P,z)|$ over its domains, background cores
and marked sites. The local old-potential map is denoted by
$\mathbb V'[\mathfrak f]$, the pointed integration map by $\mathbf F$,
and the source-independent normalizer by $\mathbf N$.

For the admitted run $x_0>\cdots>x_K$ put
$b_i=L^{-1}x_{i-1}-x_i$, so that $|b_i|\le B_\sharp$, and
$v_j^w=-L^{K-j}w/x_j$. The recursion starts with the empty array and
is
\[
 \mathfrak f_j[w]=((E^{(i)}[w],b_i)_{1\le i\le j},\emptyset),\qquad
 E^{(j+1)}[w]
 =\mathbf N^{(j+1)}+
   \mathbf F\bigl[\mathbb V'[\mathfrak f_j[w]]
                            -\mathbb W^{\rm src}[v_j^w]\bigr].
\]
The direct source family is equation~\YMMref{2165}{800}:
\[
 \mathbb W^{\rm src}[v](Y;P,B)
 =x_j\mathbb V_v^{\rm act}(Y;P,x_j^{-1/2}B)
 -x_j\sum_bv(p_b)D_A\mathbb V_1^{\rm act}(Y;P,0)
                                      [x_j^{-1/2}B(b)\delta_b].
\]
Here $p_b$ is the fixed reference plaquette in the cube containing
$b_-$, and $\mathbb V_v^{\rm act}$ is the localized weighted Wilson
action. At level zero it is the finite plaquette sum with the pivot
completion. All these maps use the fixed auxiliary data. Changing a
core radius restricts their domains; it does not change their formulas.

This recursion gives the regular terms used by the source energies.
Indeed, its negative direct source family is the one-lattice escort
of Proposition~\YMMref{78.6}{642}: their weighted action and reference
subtraction are identical, with opposite signs. On the history with
no large-field region every bond is an integration bond, so the
pointed integration map is the same. Thus the regular recursion is
equation~\YMMref{1683}{655}, and equation~\YMMref{1736}{675} gives
\[
 e_i(X;w)=\sum_{z\in X\cap\mathbb T^{(i)}}
       \bigl[E^{(i)}[w](X,\mathbf1,z)
                       -E^{(i)}[0](X,\mathbf1,z)\bigr].
\]
In particular, the continuation constructed below continues the
source energies in use, not a second choice of their factorization.

\begin{lemma}[The source recursion on smaller cores]\label{C-363-23}
Let $(K,\beta)$ be admitted at $g$, with run
$x_0>\cdots>x_K\ge\gamma^{-2}$. Use the fixed one-lattice data above
and the reference torus $T_{\ell'}$, where $\ell'$ is the least side
in $\mathfrak L(g)$; in particular $\ell'\ge2L^7$.
Write $\bar B=B_\sharp L/(L-1)$ and
$\delta(x)=x^{-1/2}\mathfrak b(x)$,
$\mathfrak b(x)=A_1(\log x)^{\mathfrak p_0}$.
The fixed core radii and size constants obey
\[
 \begin{gathered}
 \alpha_0+\alpha_1\le a^{\rm S}_\star,\qquad
 c_{\rm sh}M\alpha_0<\alpha_1,\qquad
 \alpha_1^2\le c_\natural\alpha_0,\qquad
 s_\gamma:=\frac{(1+\beta_{\rm s})\delta(\gamma^{-2})}{\alpha_0}
                                           \le s^{\rm S}_\star,\\
 E_0\ge2E_{\mathsf N},\qquad
 N_\star^{\rm s}\ge\frac L{L-1}
                       (E_0+\tfrac32\alpha_0^2B_\sharp),\\
 K_{\mathbf M}\bigl(\delta(\gamma^{-2})+
       s_\gamma N_\star^{\rm s}
       +e^{-\gamma_2\mathfrak b(\gamma^{-2})^2/2}\bigr)\le E_0/4,
 \qquad 8K_\Lambda s_\gamma N_\star^{\rm s}\le w_\star,
 \qquad 8K_Ps_\gamma\le w_\star .
 \end{gathered}
\]
These are the absolute-core and source-recursion conditions, with
their size bounds as in equations~\YMMref{2169}{801}
and~\YMMref{2170}{801}. Fix uniform bounds
$K_{\rm Q}^\vee,K_{\rm L}^\vee$ for the quadratic and linear source
coefficients, let
\[
 \mu_\star=\min\{w_\star/32,E_0/(16K_{\mathbf M}^{\rm f})\},
 \qquad \sigma_{\rm src}=\mu_\star/K_{\rm Q}^\vee,
 \qquad K^{\rm d}\le K_{\mathbf M}^{\rm f},
\]
and retain $4\beta_{\rm s}\kappa\ge1$ and
$\log\gamma^{-2}>2\mathfrak p_0$.
Then the following statements hold.

\begin{enumerate}
\item[(i)] On cores of radius $\bar a_0\le1$, at every level
$j\ge0$, the part of the direct source linear in $B$ satisfies
\[
 |\mathbb W^{\rm lin}[v](Y;P,B)|
 \le\bar a_0K_{\rm L}^\vee x_j^{1/2}
       \operatorname{sw}_j(v)(1+q_Y(B))e^{-\lambda_Vd_j(Y)}.
\]
The constants may be fixed uniformly over all levels and all the
smaller cores in (ii). Consequently the one-step estimates apply
when the oscillation allowance is divided by $\bar a_0$.

\item[(ii)] Fix $1\le k\le K$, and put
$\theta_k=\delta(x_{k-1})/\delta(\gamma^{-2})$.
For every $\theta\in[\theta_k,1]$ the recursion through level $k$
is defined on the cores with radii
$(\theta\alpha_0,\theta\alpha_1)$ for sources satisfying
\[
 |w|_\infty<\varrho_0^\star:=\sigma_{\rm src}(x_K-\bar B),
 \qquad \frac{40}{\sqrt3}M\operatorname{Lip}_\infty w
 <\varrho_1^{[k]}(\theta)
 :=\frac{40}{\sqrt3}\frac{\mu_\star}{K_{\rm L}^\vee}
                    \frac{x_{k-1}^{1/2}}{\theta\alpha_0}.
\]
It is jointly holomorphic in the source and the open background
slices, local in the source, and its pointed families satisfy
$\|E^{(i)}[w]\|_{i,\kappa+1}\le E_0$ for $1\le i\le k$.
At $\theta=1$ the estimate without the factor $\alpha_0$ in (i)
also gives the same conclusions with second source radius
$(40/\sqrt3)(\mu_\star/K_{\rm L}^\vee)x_{k-1}^{1/2}$.

\item[(iii)] The source class at $\theta=1$ contains the original
class with radii $2\varrho,2\varrho$. On common source classes,
the families constructed at $\theta$ and at $1$ agree on the
smaller cores.

\item[(iv)] The pointed integration at one step is a holomorphic
function of its fourth-family entry. On the admitted norm ball it
is Lipschitz, with bound $K^{\rm d}e^{-\lambda_{\rm out}d}$.

\item[(v)] At the trivial background the linear source potential
vanishes. The old potential vanishes to second order in the
fluctuation field. If $N$ is its weighted array size and $s$ its
absolute-core factor, then
\[
 |\mathbb V'(Y;\mathbf1,B)|
 \le K_\Lambda sNe^{-\lambda_Vd_j(Y)}
       \left(\frac{\sup|B|}
                    {\mathsf C_{\mathcal A}\mathfrak b(x_j)}\right)^2.
\]
\end{enumerate}

For the consequence used below, let $f\ne0$ be a real Lipschitz
function vanishing off the middle of a cell of side $\mathsf l$,
put $r=\varrho/\|f\|_\sharp$, and let $v'^+,v'^-$ be its even and
odd mirror sums on $T_{\ell'}$. For $1\le k\le K$, put $s=K-k$ and
\[
 \begin{aligned}
 \lambda_s&=(1+(s+1)\log L)^{\mathfrak p_0},\\
 \epsilon_s&=\max\left\{\frac{\mathsf l}{160M},
   \min\left\{\left(\frac{24}{11}L^{-s-1}\right)^{1/2},
                   \frac{24\alpha_0}{11}L^{-s-1}\lambda_s\right\}\right\},\\
 \hat\epsilon_s&=\max\left\{\frac{\mathsf l}{160M},
                                  \frac{24\alpha_0}{11}L^{-s-1}\right\}.
 \end{aligned}
\]
If $X$ is a level-$k$ domain that sees two cells, then
$G_{k,X}(z,z')=e_k(X;zv'^++z'v'^-)$ continues holomorphically to a
bidisc of radius $R_k\ge4r/\epsilon_s$, with
\[
 |G_{k,X}(z,z')|\le2M^3E_0e^{-\kappa d_k(X)},\qquad
 |\partial_z\partial_{z'}G_{k,X}(0,0)|
 \le\frac{M^3E_0c_2}{32}\frac{\hat\epsilon_s^2}{r^2}
                                     e^{-\kappa d_k(X)}.
\]
Here $c_2$ depends only on $M,\mathsf C_{\mathcal A},A_1,
\mathfrak p_0,\gamma$.
Using the last assertion of (ii) instead gives a bidisc of radius
at least $4r/\epsilon_s^\circ$, where
\[
 \epsilon_s^\circ=\max\left\{\frac{\mathsf l}{160M},
                         \left(\frac{24}{11}L^{-s-1}\right)^{1/2}\right\}.
\]
\end{lemma}

\begin{proof}
We first retain the core-radius factor in the direct source
calculation. At positive level write $a=\bar a_0$ and
$\varepsilon=\min\{\varepsilon_{\mathfrak h},
 c_{\rm os}a/(2c_{\mathfrak h})\}$. The proof of
Lemma~\YMMref{114.1}{800} gives, with $N=\#_{\pi_j}Y$,
\[
 e^{r'd_j(Y)}N\sup_{|A|<\varepsilon}
    |\mathbb V_v^{\rm act}(Y;P,A)|
       \le672M^3a^2|v|_Y\qquad(0\le r'\le\bar\kappa-1).
\]
For its linear Taylor term Cauchy's formula divides this bound
by $\varepsilon$. Retain $a^2/\varepsilon=a(a/\varepsilon)$,
instead of replacing it by $a/\varepsilon$. The bond sum obeys
$\sum_b|B(b)|\le(3/2)M^3N(1+q_Y(B))$.
The path from a plaquette to its reference plaquette has at most
$d_j(Y)+2$ links, and $e^{-d}(d+2)\le2$. These estimates give (i).
The quotient $a/\varepsilon$ is bounded by its value at the
original radius, so the constant is uniform on smaller cores.
At level zero the potential is zero unless $Y$ is the domain of
one plaquette. On the polydisc scaled by $a$, each shifted
plaquette angle is bounded by a fixed multiple of $a$. Its Wilson
potential is therefore bounded by a fixed multiple of $a^2$;
Cauchy's estimate for its linear term loses one factor $a$.
The pivot calculation in Proposition~\YMMref{51.4}{393} gives (i)
at this level as well, with one fixed $K_{\rm L}^\vee$.

We next verify that the smaller cores admit every required step.
The first three absolute-core conditions become
\[
 \theta(\alpha_0+\alpha_1)\le a^{\rm S}_\star,\qquad
 c_{\rm sh}M\theta\alpha_0<\theta\alpha_1,\qquad
 \theta^2\alpha_1^2\le c_\natural\theta\alpha_0.
\]
They follow from $0<\theta\le1$. For $j<k$, monotonicity of
$\delta$ gives
\[
 \delta(x_j)\le\delta(x_{k-1})
       =\theta_k\delta(\gamma^{-2})
       \le\theta\delta(\gamma^{-2}),\qquad
 s_j^\theta:=\frac{(1+\beta_{\rm s})\delta(x_j)}{\theta\alpha_0}
       \le s_\gamma.
\]
This is the fourth absolute-core condition. It also preserves the
field-room inequalities: the allowed Taylor radius at
$\theta\alpha_0$ is at least $\theta$ times the original one.
The fixed auxiliary maps restrict to these cores.

For the relative weight of a source in (ii),
Proposition~\YMMref{73.3}{609} gives
\[
 x_j\ge L^{K-j}(x_K-\bar B),\qquad
 |v_j^w|\le\frac{|w|}{x_K-\bar B}<\sigma_{\rm src},\qquad
 \operatorname{sw}_j(v_j^w)\le x_j^{-1}M\operatorname{Lip}_\infty w.
\]
Hence
\[
 \theta\alpha_0x_j^{1/2}\operatorname{sw}_j(v_j^w)
 <\frac{\mu_\star}{K_{\rm L}^\vee}
                   (x_{k-1}/x_j)^{1/2}
 \le\frac{\mu_\star}{K_{\rm L}^\vee}.
\]
The quadratic and higher source terms cost at most
$2K_{\rm Q}^\vee|v_j^w|$, by the same half-radius Taylor estimate
as in Lemma~\YMMref{114.1}{800}. With (i), the full source norm is
therefore less than $3\mu_\star$.

We can now carry out the induction, including its counterterms.
If the earlier pointed families have norm at most $E_0$, their
weighted array has size
\[
 N_j=\sum_{i=1}^jL^{-(j-i)}
       \bigl(\|E^{(i)}[w]\|_{i,\kappa}
                      +\tfrac32\theta^2\alpha_0^2|b_i|\bigr)
 \le\frac L{L-1}(E_0+\tfrac32\alpha_0^2B_\sharp)
 \le N_\star^{\rm s}.
\]
The old-potential rule of Proposition~\YMMref{51.16}{412} thus has
norm at most $K_\Lambda s_j^\theta N_j\le w_\star/8$.
Adding the source gives at most
$w_\star/8+3w_\star/32=7w_\star/32<w_\star/4$.
The internal potentials have total size at most
$3K_Ps_j^\theta\le3w_\star/8$.
Theorem~\YMMref{W}{423} bounds the unsourced output by $3E_0/4$.
Proposition~\YMMref{77.2}{632} bounds the added source contribution
by $3K^{\rm d}\mu_\star\le3E_0/16$.
The new pointed family therefore has norm at most $15E_0/16$.
At the first step the old array is empty, so the same argument
starts the induction. This is the proof of
Proposition~\YMMref{114.2}{801} with the displayed counterterms
and relative weights retained. Holomorphy follows from the
parameter version of the pointed expansion and the linearity of
the old-potential rule. Locality follows inductively because a
term over $X$ reads only earlier terms and source values inside
$X$. At $\theta=1$, omit the factor $\alpha_0$ from (i) and repeat
the same argument with the smaller source radius stated in (ii).

The source-radius conditions give, for $x\ge\gamma^{-2}$,
\[
 4\varrho\le\sigma_{\rm src}(x-\bar B),\qquad
 4K_{\rm L}^\vee\varrho x^{-1/2}\le40\mu_\star/\sqrt3.
\]
At $x=x_K$ and $x=x_{k-1}$ these imply
$\varrho_0^\star\ge4\varrho$ and
$\varrho_1^{[k]}(1)\ge4\varrho$.
The recursion contains no core radius except through its domain;
induction on the level therefore identifies the constructions
on their common domains. This proves (iii).
Assertion (iv) is the holomorphic parameter statement of
Lemma~\YMMref{77.1}{631} and the entry estimate of
Proposition~\YMMref{77.2}{632}, at the norm reserves just checked.

For (v), the old potential is zero at $B=0$. Its differential
there is invariant under all constant gauge transformations.
Averaging a traceless matrix over these transformations gives zero,
so the differential vanishes. Schwarz's lemma on the field
polydisc of radius $\mathsf C_{\mathcal A}\mathfrak b(x_j)$ gives
the stated square. The same invariance makes the linear source
term zero at the trivial background.

We turn to the two-variable consequence. The two mirror sums have
disjoint supports in cell interiors. Thus their source belongs to
the class in (ii) at $\theta_k$ whenever $|z|,|z'|<R_k$, with
\[
 R_k=\min\left\{\frac{\varrho_0^\star}{\sup|f|},
                 \frac{\varrho_1^{[k]}(\theta_k)}{\|f\|_\sharp}\right\}.
\]
The first denominator is controlled by
$\sup|f|\le\mathsf l\operatorname{Lip}f/4
          =\mathsf l\|f\|_\sharp/(160M)$.
For the second, the run and
$\bar B\le(13/24)\gamma^{-2}$ give
\[
 x_{k-1}\ge L^{s+1}(x_K-\bar B)
          \ge(11/24)L^{s+1}\gamma^{-2},\qquad
 \frac{40}{\sqrt3}\frac{\mu_\star}{K_{\rm L}^\vee}
                                      \ge4\varrho\gamma.
\]
Using first $\theta_k\alpha_0\le1$ and then the formula for
$\theta_k$ gives respectively
\[
 \frac{\gamma x_{k-1}^{1/2}}{\theta_k\alpha_0}
 \ge\left(\frac{11}{24}L^{s+1}\right)^{1/2},\qquad
 \frac{\gamma x_{k-1}^{1/2}}{\theta_k\alpha_0}
 =\frac{\gamma^2x_{k-1}}{\alpha_0}
       \frac{\mathfrak b(\gamma^{-2})}{\mathfrak b(x_{k-1})}
 \ge\frac{11}{24\alpha_0}L^{s+1}\lambda_s^{-1}.
\]
For the last inequality use the decrease of
$x^{-1}(\log x)^{\mathfrak p_0}$ and
$\log((11/24)L^{s+1}\gamma^{-2})/\log\gamma^{-2}
 \le1+(s+1)\log L$.
These are exactly the bounds $R_k\ge4r/\epsilon_s$.
At $\theta=1$, using the last class in (ii) gives
$R_k\ge4r/\epsilon_s^\circ$.
Summing the pointed terms at the trivial background, and
subtracting their values at source zero, gives
$|G_{k,X}|\le2M^3E_0e^{-\kappa d_k(X)}$.

It remains to improve the mixed coefficient. Put
$x=x_{k-1}$, $b_x=\mathfrak b(x)$ and $R=R_k$.
Use the fourth-family norm at the trivial background, with its
actual denominator $1+q_Y(B)$. Assertion (v) gives
\[
 \|V_{\rm old}(z,z')\|_{\rm f}\le C\mathsf A b_x^{-2},\qquad
 \|(V_{\rm old})_z\|_{\rm f}\le C\mathsf A b_x^{-2}R^{-1},\qquad
 \|(V_{\rm old})_{zz'}\|_{\rm f}
                              \le C\mathsf A b_x^{-2}R^{-2},
\]
where $\mathsf A$ is the earlier-family majorant from the induction;
$C$ here depends on $M$ and $\mathsf C_{\mathcal A}$.
The direct source is linear in $(z,z')$, so the last entry is
\[
 V(z,z')=V_{\rm old}(z,z')+zW_++z'W_-.
\]
Its part linear in $B$ vanishes at the trivial background, and the
quadratic and higher direct-source estimates give
$\|W_\pm\|_{\rm f}\le C\mu_\star\mathsf l/r$.
Here one uses $L^{s+1}/x\le(x_K-\bar B)^{-1}$,
$4\varrho\le\sigma_{\rm src}(x_K-\bar B)$ and
$\sup|f|\le\mathsf l\varrho/(160Mr)$.
The direct source has zero mixed entry derivative; its first
derivatives do not receive a factor $b_x^{-2}$.

Let $\mathfrak F_X$ denote the one-entry map with its
source-independent terms omitted. At $V_0=V(0,0)$ the mixed
coefficient is
\[
 D\mathfrak F_X(V_0)[(V_{\rm old})_{zz'}]
 +D^2\mathfrak F_X(V_0)
       [(V_{\rm old})_z+W_+,(V_{\rm old})_{z'}+W_-].
\]
The bounds $\mathsf A\le w_\star/8$,
$\mu_\star\le w_\star/32$,
$K^{\rm d}\mathsf A\le13E_0/72$ and
$K^{\rm d}\mu_\star\le E_0/16$ leave a fixed norm reserve.
The Lipschitz bound in (iv) and Cauchy's estimate in that reserve
bound the first and second derivatives by
$CK^{\rm d}e^{-\lambda_{\rm out}d_k(X)}$ and
$CK^{\rm d}w_\star^{-1}e^{-\lambda_{\rm out}d_k(X)}$.
Bounding the cross products by the two squares therefore gives
\[
 |\partial_z\partial_{z'}G_{k,X}(0,0)|
 \le CM^3E_0\bigl[b_x^{-2}R^{-2}+\mathsf l^2/r^2\bigr]
                                             e^{-\kappa d_k(X)}.
\]
Finally $\lambda_s/b_x$ is bounded in terms of
$A_1,\mathfrak p_0,\gamma$, since
$x\ge(11/24)L^{s+1}\gamma^{-2}$ and
$\log\gamma^{-2}>2\mathfrak p_0$.
Together with $R\ge4r/\epsilon_s$, this bounds the bracket by
$C\hat\epsilon_s^2/r^2$. Choosing $c_2$ as in the statement
proves the mixed-coefficient estimate.
\end{proof}
\begin{lemma}[Cell bounds]\label{C-363-16}
Let $(K,\beta)$ be admitted at $g$, $(i,o)$ a board of $L^{-K}\mathbb Z^3$ with cells of side $\mathsf l=L^{-i}$, $1\le i\le K-1$, $f$ real Lipschitz, $f\ne0$, vanishing off the middle of one cell, $r=\varrho/\|f\|_\sharp$ and $t=|z|/r$. Let $c_V=E_++C_{\rm low}+C_{\rm gap}$ be the constants of the systems of localised source energies that the pairs admitted at $g$ carry on the sides of the family $\mathfrak L(g)$ of Lemma~\ref{C-363-23} (Theorem~\YMMref{AD}{607}), and $\bar\Phi=\sup_{0<\mathsf l\le1}\Phi_1(\frac{\mathsf l}4,\frac{\mathsf l}4;\mathsf l)$ the supremum in $\mathsf l$ of the second-order constant of the one-cell bound, $\Phi_1$ the function of equation~\YMMref{(1786)}{690} at the constants of these systems. There are numbers $c'_V$, formed from $(\sigma^-,\gamma)$ and $g$, and $\Phi'$, $\Phi^{(1)}$, formed from $(\sigma^-,\gamma)$, none of them depending on the cutoff, $\mathsf l$, the cell, $f$ or the volume, and holomorphic normalisers $\mathcal E^\flat,\mathcal E^\natural$, formed on the reference torus $T_{\ell'}$ of Lemma~\ref{C-363-23} from its system, the board, the cell and $f$, and free of the side of the torus on which the seminorm is taken, such that on every torus $T_\ell$, $\ell\in\mathfrak L(g)$, $\ell\ge2\mathsf l$:
\begin{enumerate}
\item[(a)] $\|e^{\mathcal O_K(zf)}\|_{\rm ch}\le|e^{\mathcal E^\flat(z)}|\exp[c_V\mathsf l^3+\bar\Phi\,t^2/(1-t)]$ for $|z|<r$;
\item[(a$'$)] $\|e^{\mathcal O_K(zf)}\|_{\rm ch}\le|e^{\mathcal E^\natural(z)}|\exp[c_V\mathsf l^3+\Phi^{(1)}\mathsf l\,t^2]$ for $|z|\le r$, using the last source class in Lemma~\ref{C-363-23}(ii);
\item[(b)] $\|e^{\mathcal O_K(zf)}\|_{\rm ch}\le|e^{\mathcal E^\natural(z)}|\exp[c'_V\mathsf l^3+\Phi'(|z|\|f\|_{\sharp,\mathsf l}/\varrho)^2]$ for $|z|\le r$;
\item[(c)] for every $\sigma'$ as in Theorem~\ref{C-363-5}, $i\ge1$ and real $h\in C^1_c(\mathbb R^2)$ carried by a square of side $2$: $\|\mathsf I_{\mathsf l}(h)\|\le\mathsf C_3\mathsf l^{-3}|h|_{C^1}$ and $\|\mathsf I_{\mathsf l}(h)\|\le\mathsf C_2\mathsf l^{-2}|h|_{C^1}$, with the numbers of Theorem~\ref{C-363-5}(a), (b); from (a$'$), $\|\mathsf I_{\mathsf l}(h)\|\le C\mathsf l^{-5/2}|h|_{C^1}$;
\item[(d)] for $\sigma'$ a limit of periodic Gibbs states $\nu_j$, $i\ge1$, a board $(i,o)$ with $o\in L^{-i-2}\mathbb Z^3$ and $o_3=0$, a finite set $\mathcal K\subset\mathbb Z^2\times\{0\}$, so that the cells $\square_{\mathbf k}$, $\mathbf k\in\mathcal K$, lie in the layer $\{|x_3|\le\frac{\mathsf l}2\}$, real $f_{\mathbf k}\in C^1_c(\mathbb R^3)$, $f_{\mathbf k}\ne0$, vanishing off the middles of the cells $\square_{\mathbf k}$, $r_{\mathbf k}=\varrho/\|f_{\mathbf k}\|_\sharp$, $t_{\mathbf k}=|z_{\mathbf k}|/r_{\mathbf k}$, and $D_j[z]=\prod_{\mathbf k}\exp(z_{\mathbf k}\mathcal O_{K_j}(f_{\mathbf k})-\mathcal E^\flat_{\mathbf k,j}(z_{\mathbf k}))$: along a subsequence $\nu_j(\overline{\theta(\tau_{\frac{\mathsf l}2e_3}G')_j}D_j[z](\tau_{\frac{\mathsf l}2e_3}G)_j)\to\langle[G'],\mathbb D[z][G]\rangle$ locally uniformly, with bounded operators $\mathbb D[z]$, $|z_{\mathbf k}|<r_{\mathbf k}$, holomorphic in norm, $\|\mathbb D[z]\|\le\exp[c_V\#\mathcal K\mathsf l^3+\bar\Phi\sum_{\mathbf k}t_{\mathbf k}^2/(1-t_{\mathbf k})]$, $\mathbb D[0]=e^{-\mathsf lH}$ and $\partial_{z_{\mathbf k}}\mathbb D[0]=\mathsf B^{\rm in}_{\mathsf l/2}(f_{\mathbf k})+c_{\mathbf k}e^{-\mathsf lH}$; for one cell and $f$ even under the reflection at the middle plane of the cell, $\mathbb D(x)$ is at real $x$ a number times a positive operator, and $\mathbb D(x)/\langle\Omega,\mathbb D(x)\Omega\rangle\ge0$ where $\langle\Omega,\mathbb D(x)\Omega\rangle\ne0$.

For a component $\sigma'_\xi$, with its lattice approximants $\rho_{j,\xi}$ (recalled in the results used) as the states that integrate, the family is had for one cell: bounded operators $\mathbb D(z)$ on $\mathcal H(\sigma'_\xi)$, limits along a subsequence, holomorphic on $|z|<\varrho/\|f\|_\sharp$, with $\|\mathbb D(z)\|\le\bar{\mathsf Q}$, $\bar{\mathsf Q}$ the number of Theorem~\ref{C-363-5}(e), formed from the cell bound in Proposition~\YMMref{89.1}{690}, as in the proof of Theorem~\ref{C-363-5}(e), and free of $\mathsf l$, $f$ and $\xi$, with $\mathbb D(0)=e^{-\mathsf lH}$, $\mathbb D'(0)$ the insertion of $f$ plus a multiple of $e^{-\mathsf lH}$, and with the same positivity at real $x$ for even $f$; and $\|\mathsf B^{{\rm in},\xi}_\eta(h)\|\le\|\mathsf B^{\rm in}_\eta(h)\|$ for every insertion.
\item[(e)] \emph{Reflections for many cells of one layer.} Let $\nu$ be a law on $L^{-K}\mathbb Z^3$ invariant under lattice isometries and reflection positive at every coordinate bond mirror. Fix a sequence $\ell_\bullet$ of sides tending to $\infty$ and $\mathfrak C\in\{\mathcal F_{\rm loc},\mathcal C_{\rm loc}\}$, the bounded Borel or continuous local functions, respectively. Assume also chessboard domination: on every board and for every finite family of cell functions $F_{\mathbf k}\in\mathfrak C$,
\[
 \left|\nu\left(\prod_{\mathbf k}F_{\mathbf k}\right)\right|
 \le \liminf_{m\to\infty}
             \prod_{\mathbf k}\|F_{\mathbf k}\|_{{\rm ch},\ell_m}.
\]
These are the three properties needed below, as every periodic Gibbs state of the Wilson weight has them (Proposition~\textup{\YMMref{103.4}{758}}, Proposition~\YMMref{91.9}{705}); let $\{x_3=c\}$ be a wall plane of the board, $\mathcal K\subset\mathbb Z^3$ finite and non-empty with $\square_{\mathbf k}\subset\{c\le x_3\le c+\mathsf l\}$ for $\mathbf k\in\mathcal K$, $D_{\mathbf k}\in\mathcal A(\square_{\mathbf k})\cap\mathfrak C$ functions of the bonds of the cells, and $\mathsf N\ge0$ with $\prod_{\mathbf k\in\mathcal K}\|D_{\mathbf k}\|_{{\rm ch},\ell_m}\le\mathsf N$ for all large $m$. Then for all $F,F'\in\mathfrak C(\mathbb H^+_c)$, functions of the bonds of $\mathbb H^+_c=\{x_3\ge c\}$,
$$\Bigl|\nu\Bigl(\overline{\theta_cF'}\cdot\prod_{\mathbf k\in\mathcal K}D_{\mathbf k}\cdot\tau_{\mathsf le_3}F\Bigr)\Bigr|\le\mathsf N\,\|F'\|_{\nu,c}\|F\|_{\nu,c},\qquad\|F\|_{\nu,c}^2=\nu\bigl(\overline{\theta_cF}\,F\bigr).$$
For one cell the proof of Proposition~\YMMref{104.2}{759} uses exactly these three properties and gives this estimate.
\end{enumerate}
\end{lemma}
\begin{proof}[Proof sketch]
(a) By the root identity of equation~\YMMref{(1759)}{682}, the sorting identity and the two counting estimates in the proof of the cell bound, Proposition~\YMMref{89.1}{690} among the construction results, the source energies that see two cells, $\mathscr T(z,z')=(\mathsf l/\ell')^3\sum_j\sum_{X}e_j(X;zv'^++z'v'^-)$, the sum over the levels $0\le j\le K$ and the domains $X$ of the level $j$ on $T_{\ell'}$ that see at least two cells, with $\ell'$, $T_{\ell'}$, $e_j$ and the mirror sums $v'^\pm$ of Lemma~\ref{C-363-23}, are holomorphic on $|z|+|z'|<2r$ and bounded by $\bar\Phi$; their part linear in $z$ goes into the normaliser, $\mathcal E^\flat(z)=\mathcal E^\circ(z)+\mathsf az$, $\mathcal E^\circ$ the normaliser of the cell bound (Proposition~\YMMref{89.1}{690}), which is the normaliser $\mathcal E$ of Theorem~\ref{C-363-5}(e), and Cauchy's estimate bounds the rest of ${\rm Re}\,\mathscr T(z,\bar z)$ by $\bar\Phi\sum_{k\ge2}t^k$. $\mathscr T$, $\mathcal E^\circ$, $\mathsf a$ and the normalisers are formed on $T_{\ell'}$ and read no side $\ell$; the root identity carries the bound from $T_{\ell'}$ to every $T_\ell$, $\ell\in\mathfrak L(g)$, $\ell\ge2\mathsf l$, the difference of the two energy densities being at most $C_{\rm gap}$. (a$'$), (b) The normaliser $\mathcal E^\natural=\mathcal E^\circ+\mathscr T(z,0)+\overline{\mathscr T(0,\bar z)}$ absorbs all unmixed terms, and only the mixed part $\mathscr T^{\rm mx}$ counts.

The two counting estimates for domains that see two cells, in the proof of Proposition~\YMMref{89.1}{690}, give at the level $k=K-s$, with $u_s=\mathsf lL^s/M$, $(\mathsf l/\ell')^3\sum_Xe^{-\kappa d_k(X)}\le S(\kappa)\,m_s$, the sum over the domains $X$ of the level $k$ that see at least two cells and on which the mixed part of $G_{k,X}$ is not $0$, with $m_s\le\min\{u_s^3,27(\frac{u_s}2+2)^3e^{-\delta_e[u_s/2-4]_+}\}$; here $\delta_e=\kappa-\kappa_{\mathcal D}>0$ is the rate that the count of the domains leaves (Section~\YMMref{89}{688}, at $\kappa_e=\kappa$), and $\sum_sm_s\le\sigma_\Phi$, $\sigma_\Phi$ the number, formed from $L$, $M$ and $\kappa$, that the sum of the two counting bounds in that proof gives at $R=D=\frac{\mathsf l}4$.

Cauchy's estimate on the bidiscs of Lemma~\ref{C-363-23}, summed over the levels with these counts, gives $|\mathscr T^{\rm mx}(z,\bar z)|\le\Phi^{(1)}\mathsf l\,t^2$ from the radius $4r/\epsilon^\circ_s$, using the final assertion of Lemma~\ref{C-363-23}(ii), whose source radius grows with the depth, and $|\mathscr T^{\rm mx}(z,\bar z)-\tau_{11}|z|^2|\le\Phi^{(3)}\mathsf l^3\lambda(\mathsf l)^4t^3$, $\lambda(\mathsf l)=(1+\log\frac{LM}{\mathsf l})^{\mathfrak p_0}$, from the radius $4r/\epsilon_s$, that is, from (i)--(iii); the coefficient $\tau_{11}=\partial_z\partial_{z'}\mathscr T(0,0)$ obeys $|\tau_{11}|\le\Phi^{(4)}\mathsf l^2r^{-2}$ by the last bound of Lemma~\ref{C-363-23}; and $\mathsf l\,t=|z|\|f\|_{\sharp,\mathsf l}/\varrho$ for $f$ in the middle of a cell, where $\sup|f|\le\frac{\mathsf l}4{\rm Lip}f$.

So $\Phi'=\Phi^{(4)}+\Phi^{(3)}\sup_{\mathsf l\le1}\mathsf l\lambda(\mathsf l)^4$. (c) Tents of side $\mathsf l/L^2$ cut the plate function $\rho_{\mathsf l}(x_3)h(\vec x)$ into pieces $f_{\mu,k}$, sorted into $L^4$ classes $\mu$ so that the pieces of one class lie in the middles of distinct cells of one layer; $\|f_{\mu,k}\|_\sharp\le F=40M(40+8L^2)\mathsf l^{-2}|h|_{C^1}$ and at most $n\le8\mathsf l^{-2}$ pieces in a class. For one class the reflections for many cells of one layer, (e), bound the expectation of the product of the exponentials between reflected vectors by the product of their chessboard seminorms; with (a), Cauchy's estimate at the radius $R=\varrho/2S_\mu$, $S_\mu^2=\sum_k\|f_{\mu,k}\|_\sharp^2$, gives the insertion bound $\frac4\varrho e^{c_Vn\mathsf l^3+\bar\Phi/2}S_\mu$: $n$ cells of a layer cost $\sqrt n$, not $n$; summing the classes gives $\mathsf C_3\mathsf l^{-3}$.

With (b), Cauchy at $R=\min\{\varrho/2\mathsf lS_\mu,\min_k\varrho/\|f_{\mu,k}\|_\sharp\}$, where the quadratic terms add up to $\le\frac14\Phi'$ and $R^{-1}\le(2\mathsf lS_\mu+F)/\varrho\le2(\sqrt8+1)F/\varrho$, gives $\mathsf C_2\mathsf l^{-2}$. With (a$'$), Cauchy at $R=\min\{\varrho/(S_\mu(\Phi^{(1)}\mathsf l)^{1/2}),\min_k\varrho/\|f_{\mu,k}\|_\sharp\}$, where the quadratic terms add up to $\le1$ and $R^{-1}\le((8\Phi^{(1)})^{1/2}\mathsf l^{-1/2}+1)F/\varrho$, gives the exponent $\frac52$. The cutoff is removed with bounds uniform in $K$, through the lemma on the smeared insertion, Lemma~\YMMref{107.1}{766}. (d) The reflections (e) with (a) for each cell, Montel's theorem in $\#\mathcal K$ variables and a diagonal choice over a countable set of combinations with dense classes give the limit; at $z=0$ the semigroup, and the derivative is the insertion plus a number times $e^{-\mathsf lH}$.

For real $x$ and $f$ even under the reflection at the middle plane, which passes through sites, $e^{x\mathcal O(f)}$ is a function of the mirror bonds, positive, times a product of a function of one side and its mirror image, so reflection positivity makes the lattice operators the numbers $e^{-\mathcal E^\flat_j(x)}$ times positive operators; as $\mathcal E^\flat_j=\mathcal E^\circ_j+\mathsf a_jz$ with complex $\mathsf a_j$, these numbers need not be positive, but along a further subsequence their phases converge, so $\mathbb D(x)$ is a number times a positive operator, and the division by $\langle\Omega,\mathbb D(x)\Omega\rangle$ removes the number. For a component the family of one cell is formed in the same way with its lattice approximants, the one-cell bound among the construction results giving the bound $\bar{\mathsf Q}$, free of $\mathsf l$, $f$ and $\xi$, as in the proof of Theorem~\ref{C-363-5}(e): the approximants $\rho_{j,\xi}$ are periodic Gibbs states at the couplings of the sequence (Theorem~\textup{\YMMref{AP}{780}}), so they have the three conditions of (e) and the bounds of the cell bound, and their moments, centred by the $\rho_j$, converge to those of $\sigma'_\xi$ along the full sequence; so the reflections, Montel's theorem, the diagonal choice and the positivity at the middle plane run on $\mathcal H(\sigma'_\xi)$ as on $\mathcal H(\sigma')$.

For a component, 
\[
\|M_g\mathsf B^{\rm in}_\eta(h)[G]\|^2=\int g^2\|[X]_\xi\|^2d\lambda
\]
 with $X$ the tuple of the insertion (Lemma~\ref{C-363-13}(i)), and full support gives the bound at every $\xi$. (e) The multiple-reflection proof of Proposition~\YMMref{104.2}{759}, with $D=\prod_{\mathbf k\in\mathcal K}D_{\mathbf k}$ in the place of its function of one cell. All the cells lie in the one layer $c\le x_3\le c+\mathsf l$, and the mirror $\theta'$ at its middle plane keeps every cell; so $D$ is a function of $\mathfrak C(\mathbb H^+_c)$, 
\[
D^\sharp=\overline{\theta'D}=\prod_{\mathbf k}\overline{\theta'D_{\mathbf k}},
\]
and the operator $\hat DX=D\cdot\tau_{\mathsf le_3}X$ has the adjoint $\widehat{D^\sharp}$ for the form of the mirror, which, with the convexity of the moments $\mathsf s_n=\langle F,(\widehat{D^\sharp}\hat D)^nF\rangle_{\nu,c}$, holds as in that proof.

The moment $\mathsf s_n$ is the mean of a product over $2n$ layers of the translates of the $D_{\mathbf k}$ and of the $D^\sharp_{\mathbf k}$, each with the chessboard seminorm of $D_{\mathbf k}$. The coarser board is chosen so that one of its cells holds the bonds of $F$ and all the cells $\square_{\mathbf k}$, $\mathbf k\in\mathcal K$; the chessboard domination on it, with the estimate from fine to coarse cells (Lemma~\YMMref{104.1}{759}), the function $1$ standing in the other fine cells, gives $\mathsf s_n\le(\sup|F|)^2\mathsf N^{2n}$; the $n$-th root as $n\to\infty$ and the Schwarz inequality of the form give the claim.
\end{proof}

\begin{lemma}[Gluing]\label{C-363-24}
Let $G=\mathrm{SU}(2)$ and $(\Sigma,\mathsf B)$ a finite graph. (a) Let $Y_1,Y_2\subset\mathsf B$, $I=Y_1\cap Y_2$, $\mathsf R_2=Y_2\setminus I$, every site of $Y_1$ on $\mathsf R_2$ a site of $I$, and every connected component of $I$ with a cycle. Then the gauge-invariant operators of $Y_1$ and of $Y_2$ generate those gauge-invariant operators of $Y_1\cup Y_2$ that commute with the gauge transformations of $\mathsf R_2$ by central elements that are constant on each component of $I$ and $1$ off $I$; for connected $I$, all of them. For two plaquettes with exactly one common bond (seven bonds, six sites), the gauge-invariant operators of the two plaquettes generate all those of their union. (b) Let connected subgraphs $\Gamma_i$ (boxes) cover the graph, such that at every site $x$ any two boxes at $x$ are joined by a chain of boxes at $x$ in which consecutive boxes meet in a component at $x$ with two independent cycles.

Let $Z$ be the sign fields $c$ with $c(C')=1$ for every closed path $C'$ in one box, $B$ the coboundaries and $\mathsf H=Z/B$. On an open set $X_0$ of full Haar measure, invariant under gauge transformations and sign fields: for $V\in X_0$, a configuration $V'$ looks as $V$ in every box up to the gauge transformations of the box if and only if $V'=g\cdot(cV)$ with a gauge transformation $g$ and $c\in Z$; the Borel sets invariant under gauge transformations and under $Z$ agree on $X_0$ with the sets of the $\sigma$-algebra $\mathcal S_{\rm add}$ that the boxes generate; and for every character $\kappa$ of $\mathsf H$ there is a closed path $C_\kappa$ with $c(C_\kappa)=\kappa([c])$, such that every bounded gauge-invariant Borel $F$ is $\sum_\kappa\Phi_\kappa s_\kappa$ on a set of full Haar measure in $\mathcal S_{\rm add}$, $s_\kappa={\rm sign}\,{\rm tr}\,V(C_\kappa)$, $\Phi_\kappa$ $\mathcal S_{\rm add}$-measurable with $\sup|\Phi_\kappa|\le\sup|F|$. (c) For a cyclic chain of boxes $\Gamma_1,\dots,\Gamma_n$ with components $\Delta_i\subset\Gamma_i\cap\Gamma_{i+1}$ that hold two independent cycles, frames $\mathfrak f_i(V)\in\mathrm{SO}(3)$ formed from two holonomies of $\Delta_i$ at a site $y_i$, and paths $C_i$ in $\Gamma_i$ from $y_{i-1}$ to $y_i$: the entries of $T_i=\mathfrak f_{i-1}^{\rm T}{\rm Ad}(V(C_i))\mathfrak f_i$ are invariant under gauge transformations and under $Z$, measurable for $\Gamma_i$ and of modulus $\le1$, and $\chi_1(V(C_1\cdots C_n))={\rm tr}(T_1\cdots T_n)$ on $X_0$.
\end{lemma}
\begin{proof}[Proof sketch]
(a) Write $\mathfrak F(Y)$ for the bounded operators of the bonds of $Y$, acting on $L^2(G^Y)$ and tensored with $1$, $\Gamma_Y(g)$ for the gauge transformation by $g$ on the bonds of $Y$ alone, $\Gamma$ for $\Gamma_{\mathsf B}$, and $\Sigma_*=\Sigma(Y_1)\cap\Sigma(\mathsf R_2)$ for the common sites, which lie in $\Sigma(I)$ by hypothesis. \emph{Fullness.} A spanning tree of a connected set $J$ of bonds with $n\ge1$ independent cycles fixes the gauge at all but one site, so the gauge action $\Gamma_J$ on $L^2(G^J)$ is induced from the conjugation action of $G$ on $L^2(G^n)$; by Peter--Weyl every irreducible representation of the gauge group of the sites of $J$ with trivial total central character occurs in it, since every irreducible representation of $\mathrm{SO}(3)$ occurs in the conjugation action of $\mathrm{SU}(2)$ on $L^2(\mathrm{SU}(2))$: $J$ is full.

\emph{Contained in.} $\mathfrak F(Y_1)$ commutes with $\Gamma_{\mathsf R_2}(\zeta)$; on $\mathfrak F(Y_2)$ the adjoint action of $\Gamma_{\mathsf R_2}(\zeta)$ is that of $\Gamma(\zeta)$, because $\Gamma_I(\zeta)=1$, both ends of a bond of $I$ lying in one component, on which $\zeta$ is a constant central element; so the invariant operators of $Y_2$ commute with $\Gamma_{\mathsf R_2}(\zeta)$. \emph{Contains,} in three steps. (1) The bonds of $Y_1\cup Y_2$ are those of $Y_1$ and of $\mathsf R_2$, so its operators are weak limits of sums of products $AC$, $A\in\mathfrak F(Y_1)$, $C\in\mathfrak F(\mathsf R_2)$; averaging over the gauge group and over the $\Gamma_{\mathsf R_2}(\zeta)$, and writing $C$ as a weak limit of sums of members of multiplets, the right member is the weakly closed span of the sums $\sum_kA_kD_k^*$, where $(A_k)$ in $\mathfrak F(Y_1)$ and $(D_k)$ in $\mathfrak F(\mathsf R_2)$ are multiplets of one irreducible representation $\sigma$ of $G^{\Sigma_*}$, invariant under the gauge transformations at the other sites, $\sigma$ of trivial central character on the sites of each component of $I$, the mean over the $\zeta$ removing the other $\sigma$. (2) The common sites lie in $\Sigma(I)$, so the gauge group of $\Sigma_*$ acts on the bonds of $I$ (here the hypothesis on the common sites is read), and each component of $I$ is full; so there are multiplets $\chi^{(\mathsf r)}$ of $\sigma$ in $\mathfrak F(I)$, invariant at the sites of $I$ off $\Sigma_*$, with $\sum_{\mathsf r}\chi_k^{(\mathsf r)*}\chi_{k'}^{(\mathsf r)}=\delta_{kk'}$, the partial sums bounded by $1$: a multiplet of $\sigma$ in $\mathfrak F(I)$ is an intertwiner from $V_\sigma\otimes L^2(G^I)$ to $L^2(G^I)$, every type in $\sigma\otimes\Gamma_I$ has trivial central character on each component of $I$ and so occurs in $\Gamma_I$, and partial isometries between the copies of each type give the $\chi^{(\mathsf r)}$. (3) Then $\sum_kA_kD_k^*=\sum_{\mathsf r}a^{(\mathsf r)}b^{(\mathsf r)}$ with $a^{(\mathsf r)}=\sum_kA_k\chi_k^{(\mathsf r)*}$ and $b^{(\mathsf r)}=\sum_k\chi_k^{(\mathsf r)}D_k^*$; $a^{(\mathsf r)}$ is gauge invariant of $Y_1$, by the unitarity of $\sigma$, and $b^{(\mathsf r)}$ is gauge invariant of $Y_2$, because the sites of $I$ off $\Sigma_*$ are no sites of $\mathsf R_2$ and the sites of $\mathsf R_2$ off $\Sigma_*$ no sites of $I$, $\Sigma(I)$ lying in $\Sigma(Y_1)$.

For connected $I$ every $\sigma$ that occurs is trivial on the constant central elements of $\Sigma_*$, a constant central gauge transformation of all the sites of $\mathsf R_2$ acting trivially on $\mathfrak F(\mathsf R_2)$; so the condition of commutation is void. \emph{The two plaquettes.} On the space of the gauge-invariant vectors of a periodic layer with $N_1,N_2\ge4$ the invariant operators of two plaquettes with one common bond generate those of their seven bonds, by the argument with the full algebras of single plaquettes as cells that the last paragraph of the proof of Theorem~\ref{C-363-7}(b), on two plaquettes with one common bond, gives on $L^2(G^7)$ itself (spin labels, Lemma~\ref{C-363-27}, one plaquette made large, a monotone path, and the identity for the cross term at the ends of the common bond), every closed sign field of the seven bonds being a coboundary.

This is a statement on the abstract graph of the seven bonds: if the complement of a set $Y$ of bonds of a finite graph is connected, holds a cycle and meets every site, then $X\mapsto XP$, $P$ the projection on the gauge-invariant vectors, is injective on the invariant operators of $Y$, because such an operator is a sum of blocks $1\otimes\mathcal B(\mathsf K_\tau)$ over the types $\tau$ of $\Gamma_Y$, all of trivial total central character, and the block of $\tau$ is killed by $P$ only if the conjugate type does not occur in the gauge action on the complement, which, full and meeting every site, holds it. In a periodic layer with $N_1,N_2\ge4$ the complement of two plaquettes with a common bond is connected, holds cycles and meets each of their six sites. (b) Two elements of $\mathrm{SU}(2)$ that do not commute have the centraliser $\{\pm1\}$: in a component with two independent cycles the two holonomies at a site fix the gauge transformation there up to the centre, and the chains carry this along; the quotients by the compact gauge groups are standard Borel spaces, and the sign of the trace of a loop that detects a character separates the classes. (c) Under a gauge transformation the frames turn by ${\rm Ad}(g(y_i))$, which keeps scalar and vector products; sign fields do not change ${\rm Ad}$; and $\chi_1(V(C))={\rm tr}\,{\rm Ad}(V(C))$.
\end{proof}

\begin{lemma}[Sharp loops and flips on one lattice]\label{C-363-25}
On a torus $\Lambda_{K,\ell}$ at the coupling $\beta$, with $N=\ell L^K$ sites on a side and $\langle\cdot\rangle$ the Wilson measure: (a) for a rectangular path $\gamma$ with sides of at least three and at most $N-2$ bonds, $\gamma'$ its bonds but the two at each corner, $n'=\#\gamma'$, and bounded Borel $F$ of the other bonds, $|\langle\frac12{\rm tr}\,U(\gamma)F\rangle|\le(I_2(4\beta)/I_1(4\beta))^{n'}\sup|F|$; (b) on a cubic torus at $\beta\ge1$, $\langle a_p\rangle\le(21+3\log\beta)/\beta$ for every plaquette; (c) for a state $\omega$, a bounded $V$ and a projection $Q_0$ with $Q_0VQ_0=0$, and bounded $X,Y$ commuting with $Q_0$: $|\omega(XVY)|\le2\|X\|\|V\|\|Y\|\omega(1-Q_0)^{1/2}$; for the flip $\mathsf V_\zeta$ of a dual path $\zeta$ of a layer with an end plaquette $p_0$ and $Q_0=M_{\mathbf 1\{a_{p_0}<1\}}$, $Q_0\mathsf V_\zeta Q_0=0$ and $\omega(1-Q_0)\le\omega(M_{a_{p_0}})$.
\end{lemma}
\begin{proof}[Proof sketch]
(a) Two bonds of $\gamma'$ lie in no common plaquette: collinear neighbours share none; bonds of two adjacent sides could share one only at a corner, and the corners are removed; and two parallel bonds lie in a common plaquette only at the distance of one bond, while two opposite sides of $\gamma$ are at least three bonds apart in the plane of $\gamma$ and, as each side has at most $N-2$ bonds, at least two bonds apart also across the seam of the torus. Given all other bonds, the weight factorises over the bonds of $\gamma'$, and $\frac12{\rm tr}\,U(\gamma)$ is linear in each of them. For a bond $b$ the sum $A_b$ of the four staples of its plaquettes is $a\,g$ with $0\le a\le4$ and $g\in\mathrm{SU}(2)$, and by the character expansion $\int Ue^{\frac\beta2{\rm tr}(UA_b)}dU/\int e^{\frac\beta2{\rm tr}(UA_b)}dU=\frac{I_2(\beta a)}{I_1(\beta a)}g^{-1}$, a ratio that increases with its argument; the product over $\gamma'$ gives the bound. (b) $\log Z(\beta)$ is convex with $Z(0)=1$, so $\langle\sum_pa_p\rangle_\beta\le-\log Z(\beta)/\beta$; restricting the integral to $|U_b-1|\le\epsilon$ at all bonds, where $a_p=\frac12|\partial U(p)-1|^2\le8\epsilon^2$ and the Haar measure of $\{|U-1|\le\epsilon\}$ is $\ge8\epsilon^3/3\pi^3$, and using that a cubic torus has as many bonds as plaquettes, $\langle a_p\rangle\le[\log(3\pi^3/8)+3\log\epsilon^{-1}]/\beta+8\epsilon^2$; at $\epsilon^2=1/8\beta$ this is $\le(7+\frac32\log\beta)/\beta$. (c) $XVY=X(1-Q_0)VY+XQ_0V(1-Q_0)Y$; the Schwarz inequality for $\omega$, with $X(1-Q_0)X^*\le\|X\|^2(1-Q_0)$ and $Y^*(1-Q_0)Y\le\|Y\|^2(1-Q_0)$, bounds each term by $\|X\|\|V\|\|Y\|\omega(1-Q_0)^{1/2}$. $\mathsf V_\zeta$ reverses $t_{p_0}$, the end plaquette having exactly one crossed bond, so it carries $\{a_{p_0}<1\}$ onto $\{a_{p_0}>1\}$; and $\mathbf 1\{a_{p_0}\ge1\}\le a_{p_0}$.
\end{proof}

\begin{lemma}[Haar-null sets under the block map]\label{C-363-26}
Let $0\le k<K$ and let $\mathcal M$ be the block-averaging map of one step from the configurations of the fine lattice $\mathbb T^{(k)}$ to those of the coarse lattice $\mathbb T^{(k+1)}$ (equations~\YMMref{(37)--(39)}{20}). For every Haar-null Borel set $N$ of coarse configurations, $\mathcal M^{-1}(N)$ is a Haar-null Borel set of fine configurations; so the same holds for every iterate $\mathcal M^n$. Hence, in the setting of Theorem~\ref{C-363-7}(g), the laws $\mu^K_d$ of the block fields, $K\ge\max\{d,1\}$, are absolutely continuous for the Haar measure $\mathsf h_d$, and so are the limit laws $\mu_d$.
\end{lemma}
\begin{proof}
The positively oriented fine bonds are the disjoint union of the bonds with both ends in one block and of the sets $X_c$, $c$ a positively oriented coarse bond, where $X_c$ is made of the $L^2$ fine bonds that lie over $c$, parallel to $c$, one at each offset across $c$: each coarse bond has $L^2$ fine bonds of its own. $\mathcal M(U)(c)$ depends only on the values of $U$ on the bonds inside the two blocks of $c$ and on $X_c$ (Proposition~\YMMref{4.5}{20}). At fixed values $\iota$ of the bonds inside the blocks, $\mathcal M(\iota,\cdot)$ is thus a product over $c$ of Borel maps $\Psi^\iota_c$ of the bonds of $X_c$, and by Fubini's theorem it suffices that, for $\iota$ off a null set, every $\Psi^\iota_c$ pulls null sets back to null sets, a property that passes to products.

A contour variable is the ordered product of the $U$ along a path; the intrinsic logarithm is the inverse of $\exp$ on $\{|X|<\pi\}$, real-analytic on $\mathrm{SU}(2)\setminus\{-\mathbf 1\}$; $U[c]$ is the contour variable of the straight path from $c_-$ to $c_+$, and $x'=x+(c_+-c_-)$ for $x\in B(c_-)$. Equations~\YMMref{(37)--(39)}{20} give $\mathcal M(U)(c)$ by three cases, with disjoint domains that exhaust all configurations (Proposition~\YMMref{4.5}{20}). \emph{First case:} on the open set $\mathcal S_c$ of the configurations whose plaquettes with all corners in the double block of $c$ lie within $\omega_M=\frac1{200L^2}$ of $\mathbf 1$ (the small fields of the double block, a condition on its bonds alone), $\mathcal M(U)(c)=\exp(L^{-3}\sum_{x\in B(c_-)}\log W_x)U[c]$ with $W_x=U(c_-,x)U([x,x'])U(x',c_+)U[c]^{-1}$, the $U(y,x)$ being the group means of the contour variables along the staircase paths of a block from its centre $y$ to $x$, and $\log$ the series; the first case applies exactly on $\mathcal S_c$.

\emph{Second case:} off $\mathcal S_c$, where none of the closed contour variables 
\[
Y_{x,\Gamma,\Gamma'}=U(\Gamma)U([x,x'])U(\Gamma')^{-1}U[c]^{-1},
\]
$x\in B(c_-)$, $\Gamma$ and $\Gamma'$ staircase paths from $c_-$ to $x$ and from $c_+$ to $x'$, equals $-\mathbf 1$, 
\[
\mathcal M(U)(c)=\exp\bigl(L^{-3}\sum_x{\rm mean}_{\Gamma,\Gamma'}\log Y_{x,\Gamma,\Gamma'}\bigr)U[c]
\]
 with the intrinsic logarithm. \emph{Third case:} at all other configurations $\mathcal M(U)(c)=U[c]$. A contour variable $Y_{x,\Gamma,\Gamma'}$ with $x$ on the line through $c_-$ parallel to $c$ runs forth and back on that line and is identically $\mathbf 1$; every other one passes once through the bond of $X_c$ on $[x,x']$ and once through the bond of $X_c$ on $[c_-,c_+]$, two different bonds, all its other bonds lying inside the blocks; so the set where it equals $-\mathbf 1$ is the graph of one bond of $X_c$ over the other bonds, a closed real-analytic submanifold of codimension $3$.

The finite union $\mathcal N_c$ of these graphs is closed and null, and its section at every value $\iota$ of the bonds inside the blocks is null; its complement $\hat{\mathcal O}$ is open, connected, as the complement of finitely many closed submanifolds of codimension $3$ of a connected manifold, and holds $\mathbf 1$. The domain of the second case lies in $\hat{\mathcal O}$, that of the third in $\mathcal N_c$. In the formula of the second case every logarithm is taken of one of these closed contour variables, and no group mean occurs: products of bond variables, the intrinsic logarithm on $\mathrm{SU}(2)\setminus\{-\mathbf 1\}$ and the exponential make it a real-analytic map $F_c$ on $\hat{\mathcal O}$, jointly in the bonds inside the blocks and in $X_c$.

Its differential in the variables of $X_c$ at $\mathbf 1$: the factors $U(\Gamma)$, $U(\Gamma')$ do not read $X_c$, the differential of the logarithm at $\mathbf 1$ is the identity, so $dY_{x,\Gamma,\Gamma'}=dU([x,x'])-dU[c]$ there, and the differential of $F_c$ is $L^{-3}\sum_xdU([x,x'])=L^{-2}\sum_{b\in X_c}du_b$, each bond of $X_c$ lying on $L$ of the $L^3$ straight paths $[x,x']$: the mean over the $L^2$ bonds, the value that Lemma~\YMMref{4.6}{22} gives for the formula of the first case. The Gram determinant of this differential is a real-analytic function on the connected set $\hat{\mathcal O}$ with the value $L^{-6}$ at $\mathbf 1$; so its zero set $\mathcal Z_c$ is null, and off it $F_c$ is a submersion in the variables of $X_c$.

We now prove that the first branch is a submersion at every point of
$\mathcal S_c$. This avoids any assumption about its connected components.
Let $\nu$ be the direction of $c$. Index the $L^2$ lines parallel to $c$
in its double block by their transverse coordinates $z$. Each line has
one crossing bond $b_z\in X_c$; write $u_z=U(b_z)$ and let $z_0$ index
the line through $c_-$. For $x\in B(c_-)$ on the line $z$, factor the
straight transport as
\[
 U([x,x'])=P_xu_zQ_x,\qquad U[c]=g=P_0u_{z_0}Q_0.
\]
The factors $P_x,Q_x$ use only bonds inside the blocks. The contour means
also use only those bonds. Thus, with
\[
 a_x=U(c_-,x)P_x,\qquad
 b_x=Q_xU(c_+,x')^{-1},
\]
the first-branch variables are $W_x=a_xu_zb_xg^{-1}$, and $a_x,b_x$
stay fixed when the crossing variables vary.

Let $t_z$ be the point of the line $z$ whose $\nu$-coordinate is that of
$c_-$, so $t_{z_0}=c_-$ and $a_{c_-}=P_0$. Given
$\xi\in\mathfrak g$, vary all crossing bonds together by
\[
 u_z(\epsilon)=e^{i\epsilon X_z}u_z,
 \qquad X_z=\operatorname{Ad}(a_{t_z})^{-1}\xi.
\]
The norm of this tangent vector in the Euclidean product of the crossing
variables is $L|\xi|$, and $g(\epsilon)=e^{i\epsilon\xi}g$. Moreover,
\[
 W_x(\epsilon)
 =e^{i\epsilon(\xi+\rho_x)}W_xe^{-i\epsilon\xi},
 \qquad
 \xi+\rho_x=\operatorname{Ad}(a_xa_{t_z}^{-1})\xi.
\]
The error $\rho_x$ is small because the contour mean is close to every
contour in its family. Indeed,
\[
 a_xa_{t_z}^{-1}
 =U(c_-,x)U([x,t_z])U(c_-,t_z)^{-1}=h_1h_2.
\]
To obtain the last factorization, choose the contour from $c_-$ to $x$
which traverses the $\nu$ direction last. It is a contour
$\Gamma_\perp$ from $c_-$ to $t_z$ followed by $[t_z,x]$. Then
\[
 h_1=U(c_-,x)U(\Gamma)^{-1},\qquad
 h_2=U(\Gamma_\perp)U(c_-,t_z)^{-1}.
\]
Lemma~\YMMref{4.4}{19} gives $|h_i-\mathbf1|\le1/40$ and
$|W_x-\mathbf1|<1/14$. Unitarity and
$|\operatorname{Ad}(h)\xi-\xi|\le2|h-\mathbf1||\xi|$ therefore give
\[
 |h_1h_2-\mathbf1|\le\frac1{20},\qquad
 |\rho_x|\le\frac1{10}|\xi|.
\]
Write $W_x=e^{iw_x}$ with $w_x\in\mathfrak g$ and
$|w_x|<0.0715$. The differential of the logarithm gives
\[
 \left.\frac d{d\epsilon}\log W_x(\epsilon)\right|_{\epsilon=0}
 =i\eta_x,\qquad
 |\eta_x|\le\left(1.001\cdot\frac1{10}+2|w_x|\right)|\xi|
 <\frac14|\xi|.
\]
Here the first term bounds the logarithmic differential of the left
perturbation $\rho_x$, and the second is the commutator from conjugating
$w_x$ by $e^{i\epsilon\xi}$.

Put $\mathsf E=L^{-3}\sum_xw_x$ and
$\bar\eta=L^{-3}\sum_x\eta_x$. The formula for the average now yields
\[
 \left.-i\frac d{d\epsilon}\mathcal M(U(\epsilon))(c)
                 \mathcal M(U)(c)^{-1}\right|_{\epsilon=0}
 =\mathsf T_U\xi
 :=\mathsf f(\operatorname{ad}(i\mathsf E))\bar\eta
                 +\operatorname{Ad}(e^{i\mathsf E})\xi,
 \qquad \mathsf f(s)=\frac{e^s-1}{s}.
\]
The operator $\mathsf f(\operatorname{ad}(i\mathsf E))$ is the average,
for $0\le t\le1$, of the isometries
$\operatorname{Ad}(e^{it\mathsf E})$. Consequently
$|\mathsf T_U\xi|\ge3|\xi|/4$. It is an invertible map of
$\mathfrak g$, and
\[
 \zeta\longmapsto
 \bigl(\operatorname{Ad}(a_{t_z})^{-1}\mathsf T_U^{-1}\zeta\bigr)_z
\]
is a right inverse of the crossing-variable differential, of norm at
most $4L/3$.

A submersion pulls Haar-null sets back to null sets, since in local
coordinates it is a projection and Haar measure has a smooth positive
density. The first branch is a submersion everywhere on $\mathcal S_c$.
For the second branch, the exceptional set $\mathcal Z_c$ found above is
null; the third branch is contained in $\mathcal N_c$. Fubini's theorem
therefore gives, for almost every fixed internal configuration $\iota$,
the null-set property for $\Psi_c^\iota$. There are finitely many coarse
bonds, so these properties hold simultaneously outside one exceptional
set of $\iota$ of measure zero. The product map has the same property,
again by Fubini.

So $\mathcal M^{-1}(N)$ is null. Then $\mu^K_d(N)=Z_K^{-1}\int\mathbf 1_N(\mathcal M^{K-d}U)\,e^{-\beta\sum_pa_p}dU=0$ for a null $N$, and $\mu_d(N)=\lim_K\mu^K_d(N)=0$, the laws converging uniformly on all Borel sets (Corollary~\textup{\ref{IV-62-66}}).
\end{proof}

\begin{lemma}[One vertex]\label{C-363-27}
Let $V_{s_1},V_{s_2},V_{s_3}$ be irreducible representations of $\mathrm{SU}(2)$, $J_i=(J^a_i)_{a\le3}$ the spins acting in the $i$-th factor of $V_{s_1}\otimes V_{s_2}\otimes V_{s_3}$, $[J^a_i,J^b_i]=i\epsilon_{abc}J^c_i$, and $t_{ik}=\sum_aJ^a_iJ^a_k$. For every spin $J$ the algebra generated by $1$, $t_{12}$ and $t_{23}$ acts irreducibly on the multiplicity space $\mathsf M^J={\rm Hom}_{\mathrm{SU}(2)}(V_J,V_{s_1}\otimes V_{s_2}\otimes V_{s_3})$: it is the algebra of all operators of $\mathsf M^J$.
\end{lemma}
\begin{proof}[Proof sketch]
The enveloping algebra $\mathfrak U$ of three copies of the Lie algebra of $\mathrm{SU}(2)$ maps onto the algebra of all operators of the irreducible module $V_{s_1}\otimes V_{s_2}\otimes V_{s_3}$ (Burnside), equivariantly for the diagonal adjoint action, which is locally finite; averaging over the group shows that the invariant part $\mathfrak U^{\rm inv}$ maps onto the commutant of the diagonal action, the sum of the algebras of all operators of the spaces $\mathsf M^J$. By symmetrisation the graded algebra of $\mathfrak U^{\rm inv}$ is the algebra of the polynomial invariants of three vectors of $\mathbb C^3$ under $\mathrm{SO}(3)$, which is generated by the scalar products and the determinant (first fundamental theorem); by induction along the filtration of $\mathfrak U$ by the order, $\mathfrak U^{\rm inv}$ is generated by the $J_i^2$, the $t_{ik}$ and $T_3=\sum\epsilon_{abc}J^a_1J^b_2J^c_3$, and $[t_{12},t_{23}]=\sum J^a_1J^c_3[J^a_2,J^c_2]=-iT_3$. On the isotypic part $V_J\otimes\mathsf M^J$ the $J_i^2=s_i(s_i+1)$ are numbers and, the square of the total spin being $J(J+1)$, $t_{13}=\frac12\bigl(J(J+1)-\sum_is_i(s_i+1)\bigr)-t_{12}-t_{23}$. So $1$, $t_{12}$, $t_{23}$ generate the image of $\mathfrak U^{\rm inv}$ in the operators of $\mathsf M^J$, which is all of them.
\end{proof}

\subsection{Proof of Theorem~\ref{C-363-3}}

\emph{(a)} If $\eta\perp N^n\mathcal D_0$, then $\langle\eta,U(x)N^n\chi\rangle=0$ for all $x$ and $\chi\in\mathcal D_0$, as $U(x)$ keeps $\mathcal D_0$ and commutes with $N$; so the spectral measure of the pair $(\eta,N^n\chi)$ vanishes and $\langle\eta,\chi\rangle=\int(1+E)^{-n}dm_{\eta,N^n\chi}=0$. If $\psi\perp N^n\mathcal D_0(O)$, the function $(a_k)_{k\le q}\mapsto\langle\psi,N^nU(a_1)\phi(f_1)\cdots U(a_q)\phi(f_q)\Omega\rangle$, $f_k\in\mathcal F(O_1)$, with arbitrary $q\ge1$ with $O_1+\{|a|<\delta\}\subset O$, vanishes for small $a_k$ and is the boundary value of a function holomorphic in $a_k+iV_+$ (spectrum condition); by the reflection principle it vanishes for all $a_k$, and a partition of unity gives $\psi\perp N^n\mathcal D_0$.

By the bound (b5) and the core $\mathcal D_0$ of $N^n$, $n\ge\wp$, the closures from $\mathcal D_0$ and from $\mathcal D_\infty$ agree, $\langle\phi(\bar f)\chi,\eta\rangle=\langle\chi,\bar\phi(f)\eta\rangle$ for $\chi\in\mathcal D_\infty$, $\eta\in D(\bar\phi(f))$, and the relation that defines $\mathcal C(O)$ passes from $\mathcal D_0$ to $\mathcal D_\infty$. A smoothing in time turns it into $Q\bar\phi(f)\subset\bar\phi(f)Q$. Let $Q\in\mathcal C(O)$, $f\in\mathcal F(O)$, $\delta>0$ with ${\rm supp}f+(-\delta,\delta)e_0\subset O$, $g\in C^\infty_c((-\delta,\delta))$ and $Q(g)=\int g(t)e^{itH}Qe^{-itH}dt$. As $e^{-itH}\phi(f)=\phi(f(\cdot+te_0))e^{-itH}$ on $\mathcal D_\infty$ and $f(\cdot+te_0)\in\mathcal F(O)$ for $|t|<\delta$, $Q(g)$ has the same relation for this $f$ on $\mathcal D_\infty$; and $\langle N^n\chi,Q(g)\psi\rangle=\sum_{k\le n}\binom nki^k\langle\chi,Q(g^{(k)})N^{n-k}\psi\rangle$ for $\chi,\psi\in\mathcal D_\infty$ gives $Q(g)\mathcal D_\infty\subset\mathcal D_\infty$.

With $\phi(f)^\dagger=\phi(\bar f)$: $\phi(f)Q(g)=Q(g)\phi(f)$ on $\mathcal D_\infty$; as $\bar\phi(f)$ is closed, $Q(g)\bar\phi(f)\subset\bar\phi(f)Q(g)$; for an approximate unit $g_k$, $Q(g_k)\to Q$ strongly, and closedness once more gives $Q\bar\phi(f)\subset\bar\phi(f)Q$. Conversely a bounded $Q$ with $Q\bar\phi(f)\subset\bar\phi(f)Q$ for all $f\in\mathcal F(O)$ lies in $\mathcal C(O)$, by the relation between $\phi(\bar f)$ and $\bar\phi(f)$ above. So $\mathcal C(O)$ is closed under products; it holds $1$, is closed under adjoints and weakly closed: a von Neumann algebra; covariance is $U(\mathfrak g)\phi(f)U(\mathfrak g)^{-1}=\phi(f^{\mathfrak g})$ on $\mathcal D_0$. If $Q\in\mathcal C(O)$ and $Q\Omega=0$, then $\langle\psi,QX\Omega\rangle=\langle X^\dagger\psi,Q\Omega\rangle=0$ for products $X$ of fields of $O$, so $Q=0$.

The projection onto the closure of $\mathcal A(O)\Omega$ lies in $\mathcal A(O)'=\mathcal C(O)$, commutes with every $\bar\phi(f)$, $f\in\mathcal F(O)$, and fixes $\Omega$, so that closure holds $\mathcal D_0(O)$. The unitaries of $\mathcal A(O)'=\mathcal C(O)$ commute with $\bar\phi(f)$; and a net $\mathfrak N$ with this affiliation has $\mathfrak N(O)'\subset\mathcal C(O)$, so $\mathfrak N(O)\supset\mathcal A(O)$. A partition of unity writes $f\in\mathcal F(\bigcup O_i)$ as a finite sum of $f_i\in\mathcal F(O_i)$: $\mathcal C(\bigcup O_i)=\bigcap\mathcal C(O_i)$, and the commutants give $\mathcal A(\bigcup O_i)=\bigvee\mathcal A(O_i)$. For $O_2\subset O_1^\times$, $\mathcal L(O_1)\subset\mathcal C(O_1^\times)\subset\mathcal C(O_2)$ commutes with $\mathcal A(O_2)\supset\mathcal L(O_2)$. (W) says $\mathcal A(O_1)\subset\mathcal C(O_2)$ for all $O_2\subset O_1^\times$, that is $\mathcal A(O_1)\subset\mathcal C(O_1^\times)$ by additivity, that is $\mathcal L=\mathcal A$; and an open set is the union of the upright double cones in it, so additivity reduces (W) to such cones.

\emph{(b)} (b1) $P_m\phi(f)P_n=P_me^{sH}(e^{-sH}\phi(f)e^{-sH})e^{sH}P_n$ with $\|P_me^{sH}\|\le e^{s(m+1)}$; at $s=(m+n+2)^{-1}$, \eqref{C-363-2} gives $b_{mn}(f)\le e\,C_\wp\|f\|_{L^1}(m+n+2)^\wp$. (b2) With $X=P_m\phi(f)P_n$ and $[H,\phi(f)]=i\phi(\partial_tf)$,
$$(m-n)X=iP_m\phi(\partial_tf)P_n-(H-m-\tfrac12)X+X(H-n-\tfrac12),\qquad\|(H-m-\tfrac12)P_m\|\le\tfrac12,$$
so $(|m-n|-1)b_{mn}(f)\le b_{mn}(\partial_tf)$; iterate. (b3) If $\wp\le1$, then $k_1=0$ and the claim follows directly from (b1). Suppose $\wp>1$, so $k_1\ge1$. For $|m-n|>(m+n+2)^\epsilon\ge2^\epsilon$, $|m-n|\ge2$ and $|m-n|-1\ge\frac13|m-n|$, so (b2) at $k_1$ and (b1) at $\partial_t^{k_1}f$ give $b_{mn}(f)\le3^{k_1}|m-n|^{-k_1}eC_\wp\|\partial_t^{k_1}f\|_{L^1}(m+n+2)^\wp$, and $|m-n|^{k_1}>(m+n+2)^{\epsilon k_1}\ge(m+n+2)^{\wp-1}$. (b4) $\phi(h*f)=\int h(a)U(a)\phi(f)U(a)^{-1}da$, and $U(a)$ commutes with $P_m$. (b5) At $s=(1+\max\{m,n\})^{-1}$ the argument of (b1) gives $b_{mn}(f)\le e^2C_\wp\|f\|_{L^1}(1+\max\{m,n\})^\wp$, and with $\|(H-m)P_m\|\le1$ the argument of (b2) gives $(|m-n|-2)^kb_{mn}(f)\le b_{mn}(\partial_t^kf)$ for $|m-n|\ge3$. With $(1+\max\{m,n\})^2\le(1+m)(1+n)(1+|m-n|)$: $b_{mn}(f)\le e^2C_\wp(1+m)^{\wp/2}(1+n)^{\wp/2}\gamma(m-n)$, $\gamma(d)=(1+|d|)^{\wp/2}\|f\|_{L^1}$ for $|d|\le2$ and $\gamma(d)=(1+|d|)^{\wp/2}(|d|-2)^{-k_\flat}\|\partial_t^{k_\flat}f\|_{L^1}$ for $|d|\ge3$; $\sum_d\gamma(d)$ is the bracket of the form bound, and Schur's test with $\sum_a(1+a)^{2q}\|P_a\psi\|^2\le\|N^q\psi\|^2$ gives it. For the operator bound, $(2+a)^j\le2^j(1+b)^j(1+|a-b|)^j$ and $(1+\max\{a,b\})^\wp\le(1+b)^\wp(1+|a-b|)^\wp$, and the second form of Schur's test. At $\wp=4$: $\mathsf s_0(4)=195$, $2\mathsf s(4,6)<537.82$, $\mathsf s_0(2)=27$, $2\mathsf s(2,4)<37.20$. $\mathcal D_\infty$ is a core of $N^\wp$.

\emph{(c)} Put $J\phi(f_1)\cdots\phi(f_n)\Omega=\phi(\bar f_1\circ\mathsf j)\cdots\phi(\bar f_n\circ\mathsf j)\Omega$; it is antiunitary because $\mathcal W(F_1\circ\mathsf j\otimes\cdots\otimes F_N\circ\mathsf j)=\mathcal W(F_N\otimes\cdots\otimes F_1)$, which at Euclidean points is the invariance of $\sigma$ under the rotation $R_\pi$ by $\pi$ in the $(x_1,x_3)$-plane, the image of $\mathsf j=B(i\pi)$, which reverses the order of the Euclidean times, continued by Lemma~\ref{lem:setup-real-identity}; uniqueness holds because $\mathcal D_0$ is dense. Call $X=\phi(f_1)\cdots\phi(f_n)$, $f_k\in C^\infty_c(\mathbb R^{1+2})$, an ordered product in $W$ if ${\rm supp}f_1\times\cdots\times{\rm supp}f_n$ lies in $W^{(n)}=\{(x_k)_k:\ x_1\in W,\ x_{k+1}-x_k\in W\ (1\le k<n)\}$ (an order in real time along the wedge, not the order of the Euclidean heights of the setting), and put $X^{\mathsf j}=\phi(f_1\circ\mathsf j)\cdots\phi(f_n\circ\mathsf j)$, in the same order, and $X^\dagger=\phi(\bar f_n)\cdots\phi(\bar f_1)$.

As $W$ is an open convex cone, $x_k\in W$ and $x_l-x_k\in W$ for $k<l$ on the supports: all $f_k$ lie in $\mathcal F(W)$, and their supports are pairwise spacelike separated. With $\mathsf m(y)=y^1-|y^0|$, so that $W=\{\mathsf m>0\}$ and $\mathsf m(B(s)y)\ge e^{-|s|}\mathsf m(y)$, one has $d(\iota B(s+i\tau)y)=\sin\tau\,\mathsf m(B(s)y)\ge\sin\tau\,e^{-|s|}\mathsf m(y)$ for $y\in W$, $0\le\tau\le\pi$; read at $y=x_1$ and at $y=x_{k+1}-x_k$, this puts $(\iota B(\zeta)x_k)_k$ into $\mathsf{Tb}^{(n)}$ for $0<{\rm Im}\,\zeta<\pi$, and $(\iota B(\zeta)x_k+k\eta e_3)_k$ for $0\le{\rm Im}\,\zeta\le\pi$, $\eta>0$. So the function $\zeta\mapsto\langle\chi,U(B(\zeta))X\Omega\rangle$ continues through the point vectors, by (Br) and $\eta\downarrow0$ at the edges, to a function continuous on the closed strip $0\le{\rm Im}\,\zeta\le\pi$ and holomorphic inside, with the boundary value $\langle\chi,U(B(s))X^{\mathsf j}\Omega\rangle$ at ${\rm Im}\,\zeta=\pi$, $B(s+i\pi)=B(s)\mathsf j$; the comparison with the spectral integrals of $K$ on its bounded spectral sets gives $X\Omega\in D(\Delta^{1/2})$ and $\Delta^{1/2}X\Omega=X^{\mathsf j}\Omega$.

By the definition of $J$ ($\mathsf j^2=1$), $JX^{\mathsf j}\Omega=\phi(\bar f_1)\cdots\phi(\bar f_n)\Omega$, and this is $X^\dagger\Omega$, the order being reversed by locality on $\mathcal D_0$, as the supports are pairwise spacelike separated. A boost-invariant dense subspace $\mathcal D$ of $D(\Delta^{1/2})$ is a core: if $\langle\xi,\chi\rangle+\langle\Delta^{1/2}\xi,\Delta^{1/2}\chi\rangle=0$ for $\chi\in\mathcal D$, then the measure $(1+e^{-2\pi k})dm^K_{\xi,\chi}$ has Fourier transform $0$, and $\langle\xi,\chi\rangle=0$ follows. The span of $\Omega$ and of the $X\Omega$, $X$ ordered in $W$, is such a subspace: it is boost invariant, as $B(s)^{\times n}W^{(n)}=W^{(n)}$, and it is dense by the argument of (a): products of fields in the balls of radius $\frac14$ about the points $(0,2k,0)$ stay ordered in $W$ under small translations of each factor, the reflection principle carries the orthogonality to all translations, and a partition of unity to all products of fields.

Locality removes the ordering: for $Y=\phi(h_1)\cdots\phi(h_r)$, $h_k\in\mathcal F(W)$ in any order, and $\xi=X\Omega$, $X$ ordered in $W$, the fields of $(X^{\mathsf j})^\dagger=\phi(\bar f_n\circ\mathsf j)\cdots\phi(\bar f_1\circ\mathsf j)$ lie in $\mathcal F(W')$, so they commute on $\mathcal D_0$ with those of $Y$, and $(X^{\mathsf j})^\dagger\Omega=J\xi$; hence $\langle\Delta^{1/2}\xi,Y\Omega\rangle=\langle\Omega,(X^{\mathsf j})^\dagger Y\Omega\rangle=\langle Y^\dagger\Omega,(X^{\mathsf j})^\dagger\Omega\rangle=\langle Y^\dagger\Omega,J\xi\rangle=\langle\xi,JY^\dagger\Omega\rangle$, also at $\xi=\Omega$; on the core this gives $Y\Omega\in D(\Delta^{1/2})$ and $J\Delta^{1/2}Y\Omega=Y^\dagger\Omega$ for all $h_k\in\mathcal F(W)$.

For $Q\in\mathcal C(W)$ the weak commutation, used $n$ times, gives $\langle Q^*\Omega,X\Omega\rangle=\langle X^\dagger\Omega,Q\Omega\rangle=\langle J\Delta^{1/2}X\Omega,Q\Omega\rangle$ on the core, so $\Delta^{1/2}JQ\Omega=Q^*\Omega$. Opposite wedges: let $Q\in\mathcal C(W)$, $Q'\in\mathcal C(W')$ and $Q'_g=\int g(a)U(a)Q'U(a)^{-1}da$, $g\in C^\infty_c(W)$; as $W'\subset W'+a$ for $a\in W$, $Q'_g\in\mathcal C(W')$; it keeps $\mathcal D_\infty$ (smoothness in the time translation) and commutes there with the fields of $W'$. With $S=J\Delta^{1/2}$, $J$ carries the graph of $\Delta^{1/2}$ onto that of $\Delta^{-1/2}$, so $\mathcal D_0(W')$ is a core of $S^*=J\Delta^{-1/2}$, and every $R$ on $\mathcal D_\infty$ with an adjoint there that commutes with the fields of $W'$ has $SR\Omega=R^\dagger\Omega$; and $S^*Q\Omega=Q^*\Omega$.

For products $X,Y$ of fields of $W$ and $R=Y^\dagger Q'_gX$:
$$\langle Y\Omega,QQ'_gX\Omega\rangle=\langle S^*Q\Omega,R\Omega\rangle=\langle SR\Omega,Q\Omega\rangle=\langle Y\Omega,Q'_gQX\Omega\rangle ;$$
so $QQ'_g=Q'_gQ$, and $g\to\delta$ gives $QQ'=Q'Q$. For a wedge $\tilde W\supset O$, $\tilde W'\subset O^\times$, so $\mathcal C(O^\times)\subset\mathcal C(\tilde W')\subset\mathcal C(\tilde W)'=\mathcal A(\tilde W)$, whence $\mathcal C(O^\times)\subset\mathcal A^d(O)$; two spacelike separated double cones are parted by a wedge, which gives the locality; if $\mathcal C(O^\times)=\mathcal A^d(O)$ for all double cones, then $\mathcal A(O)\subset\mathcal C(O^\times)$ and (W) holds by (a), and the converse is (e).

\emph{(d)} The function $f_*$ is that of Lemma~\ref{C-363-19}(a). (W) implies density: $\mathcal A(W')\subset\mathcal C(W)$ and $\mathcal A(W')\Omega$ is dense. Density implies wedge duality: $\Omega$ is then cyclic and separating for $\mathcal N=\mathcal C(W)$; by (c), $Q\Omega\mapsto Q^*\Omega$ is the restriction of $\Delta^{1/2}J=J\Delta^{-1/2}$ ($JKJ=-K$), and $J\mathcal N\Omega$ is a boost-invariant dense subspace of $D(\Delta^{1/2})$, hence a core; so modular theory gives the conjugation $J$, the modular operator $\Delta^{-1}$, and $\mathcal N'=J\mathcal NJ=\mathcal C(W')$. Wedge duality implies (W): spacelike separated upright double cones lie in some $\mathfrak gW'$ and $\mathfrak gW$, and $\mathcal C(K_1)'\subset\mathcal C(\mathfrak gW')'=\mathcal C(\mathfrak gW)\subset\mathcal C(K_2)$. $J\mathcal A(W)J=\mathcal C(W')'$, so $J\mathcal A(W)J\subset\mathcal C(W)$ is (W) for the pair $(W',W)$, equivalent to (W) by the same chain with $W$ and $W'$ exchanged. (W) implies (CP): $\bar\phi(f)$, $f\in\mathcal F(K)$ real, is affiliated with $\mathcal A(K)\subset\mathcal C(K^\times)$; the extension of Lemma~\ref{C-363-12}(b) has its bounded Borel functions in $\mathcal A(K)\otimes M_2$, which commute weakly with the fields of $K^\times$; Lemma~\ref{C-363-12}(a) at $\mathfrak A$ generated by $1$, $f$ and $\mathcal F(K^\times)$, $X_0=f$. (CP) implies (QCP$_f$), as $W\subset K^\times$ for $\bar K\subset W'$. (QCP$_f$) gives $\phi(f)\Omega\in\mathcal H_0$.

Let $\mathsf L(Y_0+\lambda Y_1+\lambda^2Y_2)=\langle\Omega,Y_0\Omega\rangle+\langle\Omega,Y_1\eta\rangle+\langle\eta,Y_2\eta\rangle$ on the quadratic polynomials in $\lambda$ with coefficients in $\mathfrak P(W)$; $f$ commutes with $\mathfrak P(W)$ in $\mathfrak P$, as $\bar K$ and $W$ are spacelike separated. On a square $(Z_0+\lambda Z_1)^*(Z_0+\lambda Z_1)$, $\mathsf L$ is $\|Z_0\Omega+Z_1\eta\|^2\ge0$ for all data; (QCP$_f$) is the positivity of $\mathsf L$ on the larger cone of the quadratic polynomials whose values at all real $\lambda$ lie in $\mathfrak P(W)^+$, and the extension below is positive on this cone. Let $\mathscr C$ be the space of the functions $G=\sum_kg_k(\lambda)Y_k$, $Y_k\in\mathfrak P(W)$, $g_k$ complex Borel with $|g_k|\le C(1+\lambda^2)$, and $\mathscr C^+$ the cone of those whose values at all real $\lambda$ lie in $\mathfrak P(W)^+$.

Every hermitian element of $\mathscr C$ is a sum of terms $gS$, $g$ real, $S=S^*$, and $\frac12q(S^2+1)-gS=\frac14(q-g)(S+1)^2+\frac14(q+g)(S-1)^2\in\mathscr C^+$ for $q=C(1+\lambda^2)$, $|g|\le q$, with $\frac12q(S^2+1)$ a quadratic polynomial: the step of the proof of Lemma~\ref{C-363-12}(a) at $r=1$. The extension theorem of M.~Riesz extends $\mathsf L$ to a linear $\hat\omega$ on $\mathscr C$ with $\hat\omega\ge0$ on $\mathscr C^+$. For $G,G'$ with $|g_k|\le C(1+|\lambda|)$, $G^*G\in\mathscr C^+$, so $\langle[G],[G']\rangle=\hat\omega(G^*G')$ gives a Hilbert space; $V\pi(Y)\Omega=[Y]$ on the dense $\mathcal D_0(W)$ is an isometry $V$ of $\mathcal H$ into it; $\hat\pi(g)[G]=[gG]$, $g$ bounded Borel, is bounded by $\sup|g|$, as $(\sup|g|^2-|g|^2)G^*G\in\mathscr C^+$; and $Q_g=V^*\hat\pi(g)V$ commutes weakly with the fields of $W$ on $\mathcal D_0(W)$, so on $\mathcal D_0$ by (a) and (b5): $Q_g\in\mathcal C(W)$, $\|Q_g\|\le\sup|g|$.

With $g_n(\lambda)=\lambda\mathbf 1_{|\lambda|\le n}$, $Q_{g_n}\Omega\to\phi(f)\Omega$, because $\|[\mathbf 1_{|\lambda|>n}Y]\|^2\le n^{-2}\mathcal W(f^2Y^*Y)$ and $V^*[\lambda]=\phi(f)\Omega$. The first order holds for all data: if $Y_0+\lambda Y_1\in\mathfrak P(W)^+$ for all real $\lambda$, then $\langle\psi,\pi(Y_1)\psi\rangle=0$ for $\psi\in\mathcal D_0$, so $\pi(Y_1)=0$ there by polarisation, and $\mathcal W(Y_0+fY_1)=\langle\Omega,Y_0\Omega\rangle\ge0$. One vector decides: $\mathcal H_0$ is invariant under $\mathcal C(W)$, under the boosts, and under the translations $U(-a)$, $a\in\bar W$ (as $\mathcal C(W)\subset\mathcal C(W+a)$), hence, by the spectrum condition and the reflection principle along the light rays $(\pm1,1,0)$ and along $e_2$, under all translations; so its projection $P_0$ commutes with $\mathcal P$, and $\mathsf D(f)=\|(1-P_0)\phi(f)\Omega\|$ is a seminorm invariant under translations and boosts, with $\mathsf D(h*f)\le\|h\|_{L^1}\mathsf D(f)$.

For $\chi\perp\mathcal H_0$ the tempered distribution $\mathsf T(h)=\langle\chi,\phi(h)\Omega\rangle$ vanishes on all $f*k$, $f\in\mathfrak F$, $k\in\mathcal S$; if the $\hat f$, $f\in\mathfrak F$, have no common real zero, a partition of unity shows that $\mathsf T$ vanishes on all $h$ with $\hat h\in C^\infty_c$, so $\phi(h)\Omega\in\mathcal H_0$ for all $h\in\mathcal S$. Then for $\chi\in\mathcal D_\infty\cap\mathcal H_0^\perp$ and $h\in\mathcal F(W)$, $\langle Q\Omega,\phi(\bar h)\chi\rangle=\langle\bar\phi(h)Q\Omega,\chi\rangle=\langle Q\phi(h)\Omega,\chi\rangle=0$ for $Q\in\mathcal C(W)$; for $\psi\in\mathcal H_0\cap\mathcal D_\infty$ the vector $\chi=(1-P_0)\phi(h)\psi$ has $\|\chi\|^2=\langle\phi(\bar h)\chi,\psi\rangle=0$; by induction $\mathcal D_0(W)\subset\mathcal H_0$, and $\mathcal H_0$ is dense.

The infimum over the boosts is the distance itself, by the boost invariance of $\mathcal H_0$; for one $f_0\ge0$ the transforms of its Lorentz images have no common real zero by Lemma~\ref{C-363-19}(d), and translations do not change zero sets. So all the conditions are equivalent, and the modular data under (W) are those of the density step, applied to $\mathcal A(W)=\mathcal C(W')$. (CP) gives (W) also without the wedge relation: for spacelike separated upright double cones $K_1,K_2$, $Q\in\mathcal C(K_1)'$ and real $g\in\mathcal F(K_2)$, the extension of Lemma~\ref{C-363-12}(a) for $(K_2,g)$ gives bounded $Q_h=V^*h(\hat T)V\in\mathcal C(K_2^\times)\subset\mathcal C(K_1)$, which commute with $Q$, and $Q_{g_n}\psi\to\phi(g)\psi$ for $\psi$ in the domain of the GNS representation, which holds $\mathcal D_0(K_2^\times)$; so $Q$ commutes weakly with $\phi(g)$ there, and the polynomial bound carries this to $\mathcal D_0$.

The inequalities: under (W) the closures of $\phi(g_1),\phi(g_2)$ are affiliated with commuting algebras, and Lemma~\ref{C-363-12}(c) gives strongly commuting self-adjoint extensions $\hat A_1,\hat A_2$; with $\psi=\phi(f)\Omega$ and their joint spectral measure,
\[
\begin{gathered}
|\mathcal W(fg_1g_2)|=\Bigl|\int\lambda_1\lambda_2\,d\mu_{\psi,\Omega}\Bigr|\\
\le\Bigl(\int|\lambda_1\lambda_2|d\mu_\psi\Bigr)^{1/2}\Bigl(\int|\lambda_1\lambda_2|d\mu_\Omega\Bigr)^{1/2}\\
\le\bigl(\|\hat A_1\psi\|\|\hat A_2\psi\|\bigr)^{1/2}\|\hat A_1\hat A_2\Omega\|^{1/2},
\end{gathered}
\]
and locality on $\mathcal D_0$ identifies the squares of the three norms with $\mathcal W(f^2g_1^2)$, $\mathcal W(f^2g_2^2)$, $\mathcal W(g_1^2g_2^2)$; likewise $\mathcal W(p(f_1,f_2))=\int p\,d\mu_\Omega\ge0$, and Motzkin's polynomial is covered by Lemma~\ref{C-363-12}(d).

\emph{(e)} Under (W), $\mathcal A(\tilde W)=\mathcal C(\tilde W')$ for every wedge by (d) and covariance. Two spacelike separated open convex sets are parted by a wedge (Hahn--Banach), so $O^\times=\bigcup\{\tilde W':\tilde W\supset O\}$, and additivity gives $\mathcal A^d(O)=\bigcap\mathcal C(\tilde W')=\mathcal C(O^\times)$. Locality: a wedge parts $O_1$ from $O_2$. Haag duality: $\mathcal A^d(O)'=\mathcal A(O^\times)=\bigvee\mathcal A(O_3)\subset\bigvee\mathcal A^d(O_3)=\mathcal A^d(O^\times)\subset\mathcal A^d(O)'$, the joins over the double cones $O_3\subset O^\times$; and as $O^\times=\bigcup\{\tilde W':\tilde W\supset O\}$, (a) gives $\mathcal A(O^\times)=\bigvee\{\mathcal A(\tilde W'):\tilde W\supset O\}$.

Outer regularity: a point with a neighbourhood spacelike to $O$ is spacelike to the compact set $\bar O$, hence to $\bar O_\epsilon$ for small $\epsilon$, the compact sets $\bar O_\epsilon$ decreasing to $\bar O$; so $O^\times=\bigcup_{\epsilon>0}O_\epsilon^\times$, and (a) gives $\mathcal A^d(O)=\mathcal C(O^\times)=\bigcap_\epsilon\mathcal C(O_\epsilon^\times)=\bigcap_\epsilon\mathcal A^d(O_\epsilon)$. An operator commuting with all $\mathcal A(O_1)$, $O_1\subset O^\times$, commutes with $\mathcal A(O^\times)$, so it lies in $\mathcal A^d(O)$. For $O\subset\tilde W$: $\mathcal A(O)\subset\mathcal A^d(O)\subset\mathcal A(\tilde W)$, so the wedge algebras of $\mathcal A^d$ are those of $\mathcal A$, and $(\mathcal A^d)^d=\mathcal A^d$. $\mathcal A(O)=\mathcal A^d(O)$ if and only if $\mathcal C(O^\times)\subset\mathcal A(O)=\mathcal C(O)'$.

Covariance, $J$, cyclicity and affiliation follow from (a) and (c). Without (W): if the intersections are local, then $\mathcal A(O_1)\subset\mathcal A^d(O_1)\subset\mathcal A^d(O_2)'\subset\mathcal A(O_2)'=\mathcal C(O_2)$ for spacelike separated double cones.

\emph{(f)} Let $A_j=\phi(\partial_t^jf)$ and $R_\epsilon=(1+\epsilon N)^{-1}$. From $[H,A_0]=iA_1$: $R^r_\epsilon A_0\chi=A_0R^r_\epsilon\chi-i\epsilon\sum_{j=1}^rR^j_\epsilon A_1R^{r-j+1}_\epsilon\chi$. For $A_0^*\eta=\pm i\eta$ and $\eta_\epsilon=R^{k_0}_\epsilon\eta$, the imaginary part of $\langle\eta_\epsilon,\bar A_0\eta_\epsilon\rangle$ and the form bound of $A_1$ give $\|\eta_\epsilon\|^2\le c\sum_j\epsilon^{1/2}\|N^{1/2}R_\epsilon^{k_0-j+1}\eta_\epsilon\|\,\epsilon^{1/2}\|N^{1/2}R^j_\epsilon\eta\|\to0$, while $\eta_\epsilon\to\eta$; so $\bar A_0$ is self-adjoint. $N^k\phi(f)=\sum_j\binom kji^j\phi(\partial_t^jf)N^{k-j}$ gives $\bar A_0N^{-k}=N^{-k}(\bar A_0+K_k)$ with $K_k$ bounded under the block bounds, whence the resolvents keep $D(H^k)$ off a strip and, through cores, commute with $\bar\phi(h)$ for $h$ spacelike to $f$; Neumann series reach all $z\notin\mathbb R$; so the resolvents lie in $\mathcal C(O)$. (f1) With the logarithm the self-adjointness and the commutation come from unitary groups along which the exponent of the energy falls.

Put $\mathsf w(x)=x(1+\log x)$, and for real $h\in C^\infty_c$ let $v=f_i*h$. Off the band (b3) bounds the blocks of $\partial_tf_i$ linearly, so $b_{mn}(\partial_tf_i)\le c\,\mathsf w(m+n+2)$ for all $m,n$; by (b4), $\partial_t^jv=(\partial_tf_i)*\partial_t^{j-1}h$ has the same bound with $c\|\partial_t^{j-1}h\|_{L^1}$, $j\ge1$, and (b2) adds the decay $(|m-n|-1)^{-3}$ for $|m-n|\ge2$; as $\mathsf w(m+n+2)\le3.4(1+|m-n|)\,\mathsf w(m+1)^{1/2}\mathsf w(n+1)^{1/2}$, Schur's test with the weights $\mathsf w(n+1)$ gives $|\langle\psi',\phi(\partial_t^jv)\psi\rangle|\le c_j\|\mathsf w(N)^{1/2}\psi'\|\|\mathsf w(N)^{1/2}\psi\|$ on $\mathcal D_\infty$, $j\ge1$. By $\log x=\int_0^\infty[(1+t)^{-1}-(x+t)^{-1}]dt$, $[N+t,\phi(\partial_t^jv)]=i\phi(\partial_t^{j+1}v)$ and $\int_0^\infty\mathsf w(x)(x+t)^{-2}dt=1+\log x$, this turns into $|\langle\log N\,\psi,\phi(\partial_t^jv)\psi\rangle-\langle\phi(\partial_t^jv)\psi,\log N\,\psi\rangle|\le c_{j+1}\langle\psi,(1+\log N)\psi\rangle$, $j\ge0$, the left member being $|\mathfrak l^{(j)}(\psi,\psi)|$ for the hermitian form $\mathfrak l^{(j)}(\psi',\psi)=\int_0^\infty\langle(N+t)^{-1}\psi',\phi(\partial_t^{j+1}v)(N+t)^{-1}\psi\rangle dt$.

The three-lines theorem, on vectors of bounded energy, carries the bound to $\mathfrak l^{(j)}(N^\tau\psi',N^{-\tau}\psi)$ with the constant $\sum_{l\le\lceil|\tau|\rceil}\binom{\lceil|\tau|\rceil}lc_{j+1+l}$; as $\frac d{d\tau}\langle N^\tau\chi',AN^{-\tau}\chi\rangle=i\,\mathfrak l^{(0)}(N^\tau\chi',N^{-\tau}\chi)$ for $A=\phi(v)$, and a second derivative brings in $\mathfrak l^{(1)}$, integration over $|\tau|\le\frac\rho2$ at $\chi'=\chi$ gives
$$|\langle N^\rho\chi,A\chi\rangle-\langle\chi,AN^\rho\chi\rangle|\le\rho\,c_1\langle\chi,N^\rho(1+\log N)\chi\rangle+\rho^2c'(\rho)\langle\chi,N^\rho\chi\rangle\qquad(\chi\in\mathcal D_\infty,\ \rho\ge0),$$
$c'(\rho)=\sum_{l\le\lceil\rho/2\rceil}\binom{\lceil\rho/2\rceil}lc_{l+2}$. The bounded $A_s=e^{-sH}Ae^{-sH}$ satisfy the same with the same numbers. For $\psi\in\mathcal D_\infty$ and $y(\lambda)=\|N^{\rho(\lambda)/2}e^{i\lambda A_s}\psi\|^2$ with the falling exponent $\rho(\lambda)=\rho_0e^{-c_1|\lambda|}$, the term $\rho'\langle N^\rho\log N\rangle$ of $y'$ cancels the term with the logarithm: $y'\le(\rho_0c_1+\rho_0^2c'(\rho_0))\,y$ for $\lambda\ge0$, and Gronwall's lemma bounds $\|N^{\rho(\lambda)/2}e^{i\lambda A_s}\psi\|$ uniformly in $s$. As $s\downarrow0$, $\|(A_s-A_{s'})\chi\|\le c|s-s'|\|N^{q+1}\chi\|$ ($q\ge\wp$ an integer) makes the groups $e^{i\lambda A_s}$ converge strongly to a unitary group $V$ with $V(\lambda)D(N^r)\subset D(N^{re^{-c_1|\lambda|}})$ for $r\ge0$, so $V(\lambda)\mathcal D_\infty=\mathcal D_\infty$, and $\lambda\mapsto V(\lambda)\psi$ is continuously differentiable in every norm $\|N^a\cdot\|$ with the derivative $iAV(\lambda)\psi$; its generator is the closure of $A$ on $\mathcal D_\infty$ (for $\eta$ with $(A|_{\mathcal D_\infty})^*\eta=\pm i\eta$, $\lambda\mapsto\langle V(\lambda)\psi,\eta\rangle$ is bounded with the derivative $\pm$ itself, hence $0$): $\phi(v)$ is essentially self-adjoint on $\mathcal D_\infty$ and on $\mathcal D_0$.

The group loses moments of the energy and keeps $\mathcal D_\infty$. For $h'\in\mathcal S$ whose support is spacelike separated from ${\rm supp}\,v$, $\phi(v)\phi(h')=\phi(h')\phi(v)$ on $\mathcal D_\infty$ (locality on $\mathcal D_0$, the core $\mathcal D_0$ of $N^n$, (b5)), and $\frac d{d\lambda}\langle\phi(\bar h')V(\lambda)\psi',V(\lambda)\psi\rangle=0$ gives $V(\lambda)\bar\phi(h')\subset\bar\phi(h')V(\lambda)$: by (a) the bounded Borel functions of $\bar\phi(v)$ lie in $\mathcal C(O)$ for every open $O$ spacelike separated from ${\rm supp}\,v$. For real $h\in\mathcal F(O)$ the $f_i*h$ lie in $\mathcal F(O)$ for large $i$ and tend to $h$ in $\mathcal S$; so $\mathcal C(O)$ is the commutant of the groups of the $\bar\phi(f_i*h)$ with ${\rm supp}(f_i*h)\subset O$ (by (a) every $Q\in\mathcal C(O)$ commutes with them; conversely weak commutation with the $\phi(f_i*h)$ passes to $\phi(h)$ by (b5)), and for spacelike separated $O_1,O_2$ the groups of $O_1$ lie in $\mathcal C(O_2)$: $\mathcal C(O_1)'\subset\mathcal C(O_2)$, which is (W) without a wedge, and (CP) by (d). (f2) The time reversal $[F]\mapsto[\bar F]$ gives the other sign for $f_0$ even in $t$; $\partial_t^j(h*f_0)=(\partial_t^{j-1}h)*\partial_tf_0$ gives the hypotheses of the first part at every $h*f_0$; by Lemma~\ref{C-363-19}(d) the images of $f_0$ under rotations, boosts and translations into $W'$ have no common real zero of their transforms, and the resolvents of the corresponding operators lie in $\mathcal C(W)$ and produce from $\Omega$ a total set; (d) applies. (f3) For products $X$ of fields spacelike to $f_0$, $\|\phi(f_0)^nX\Omega\|^2=\langle\phi(f_0)^{2n}\Omega,X^\dagger X\Omega\rangle\le\|\phi(f_0)^{2n}\Omega\|\|X^\dagger X\Omega\|$, so these total vectors are quasi-analytic, $\bar\phi(f_0)$ is self-adjoint with resolvents commuting with the outside fields, and the argument of (f2) applies. (f4) At $s=(m+n+2)^{-1}$ the argument of (b1) gives $b_{mn}(f)\le e\,C\|f\|_{L^1}(m+n+2)(1+\log(m+n+2))$ for every $f\in\mathcal S$ (for $\wp\le1$ (b1) gives it without the logarithm), so for every $\partial_t^jf$: the hypothesis of (f1) holds for every approximate unit $f_i$, and the argument of (f1) runs with $v=f$ itself; so every real $\phi(f)$ is essentially self-adjoint on $\mathcal D_0$, and $e^{i\lambda\bar\phi(f)}\bar\phi(h)\subset\bar\phi(h)e^{i\lambda\bar\phi(f)}$ for real $h$ whose support is spacelike separated from that of $f$; $\bar\phi(h)$ being self-adjoint, the group of $\bar\phi(f)$ commutes with its spectral projections: the closures commute strongly. (f5) Lemma~\ref{C-363-20}(b)--(d) gives (LB$_f$) for every real $f\in C^\infty_c$, so (f2) at every $f_0$.

The last assertion of (f): if (f2) holds at $f_*$, time reversal gives the two-sided bound $c_*(\partial_tf_*)<\infty$, and with Lemma~\ref{C-363-19}(b), (c) the functions $f_i=g^{\otimes3}_{1/i,3/2}=f_**\mu_i^{\otimes3}$ have $c_*(\partial_tf_i)\le\|\mu_i\|^3_{L^1}c_*(\partial_tf_*)$, hence $b_{mn}(\partial_tf_i)\le\|\mu_i\|_{L^1}^3c_*(\partial_tf_*)(m+n+2)$ by polarisation; these functions have the properties that (f1) asks of the $f_i$.

\emph{(g)} By Lemma~\ref{C-363-13}(ii), $\phi_u([F],M_{g^2}[F])=\int g^2\phi^\xi_u\,d\lambda$ and $\langle[F],M_{g^2}N[F]\rangle=\int g^2\langle[F]_\xi,N_\xi[F]_\xi\rangle d\lambda$; at $g=1$ the bound passes from all $\sigma_\xi$ to $\sigma$; conversely, the bound for $\sigma$ at the vectors $M_g[F]$ gives $\int g^2(c\,\mathsf n-\mathsf a)d\lambda\ge0$ for all real $g$, with $\mathsf a(\xi)=\phi^\xi_u([F]_\xi,[F]_\xi)$ and $\mathsf n(\xi)=\langle[F]_\xi,N_\xi[F]_\xi\rangle$ continuous, and full support gives $\mathsf a\le c\,\mathsf n$ at every $\xi$; the classes $[F]_\xi$ form a core of $N_\xi$. The premise of (CP) speaks of $\mathfrak P$ alone, and $\mathcal W_\sigma=\int\mathcal W_\xi d\lambda$ (Lemma~\ref{C-363-13}(iii)): (CP) passes from all $\sigma_\xi$ to $\sigma$, and so does every fixed inequality $\mathcal W(a)\ge0$.

Conversely, for $a=\sum f^nY_n$ as in (CP), which is hermitian, and $b$ formed as in Lemma~\ref{C-363-13}(iv), a polynomial in means of translates of elements $b_i$ of $\mathfrak P$ over pairwise different spatial translations of length $\ge R$, with $R$ so large that the translated $b_i$ are spacelike to $K$, which puts $b$ in the algebra of (CP), $\sum\lambda^nb^*Y_nb\in\mathfrak A^+$ and $f$ is central, so $\mathcal W_\sigma(b^*ab)\ge0$, and Lemma~\ref{C-363-13}(iv) gives (CP) for every $\sigma_\xi$; (d) turns (CP) into (W). The limits: Lemma~\ref{C-363-13}(v); inequalities $\mathcal W(a)\ge0$ pass to the limit.

\emph{(h)} (h1) Fix $\Lambda<\infty$; on $\mathcal H_\Lambda$ let $q$ be the bounded operator of the form 
\[
\int h'(\vec x)q_0(U(-\vec x)\cdot,U(-\vec x)\cdot)d^2x;
\]
then 
\[
\langle\psi',\mathsf I\psi\rangle=\int\mathsf l^{-1}\rho_1(\varsigma/\mathsf l)\langle e^{-(\frac{\mathsf l}2+\varsigma)H}\psi',q\,e^{-(\frac{\mathsf l}2-\varsigma)H}\psi\rangle d\varsigma
\]
 and 
\[
\langle\psi',\phi(g_{\mathsf l}\otimes h')\psi\rangle=\int g_{\mathsf l}(t)\langle e^{-itH}\psi',q\,e^{-itH}\psi\rangle dt.
\]
An operator integral $\int\alpha_\tau(H)q\beta_\tau(H)d\mu(\tau)$ depends on $(\alpha,\beta,\mu)$ only through its symbol \[\int\alpha_\tau(E)\beta_\tau(E')d\mu(\tau)\] (approximate the functions by step functions on a fine partition of $[0,\Lambda]$). With $A(x)=\hat D(x)+\epsilon$ and $\log A=\int_0^\infty[(1+t)^{-1}-(t+A)^{-1}]dt$, $L_\epsilon(x)=-\log(\hat D(x)+\epsilon)$ is $C^2$ in norm and $L_\epsilon'(0)=-\int_0^\infty R^\epsilon_t\hat D'(0)R^\epsilon_tdt$, $R^\epsilon_t=(t+\epsilon+e^{-\mathsf lH})^{-1}$. As $\epsilon\downarrow0$ on $\mathcal H_\Lambda$, $-L_\epsilon'(0)$ tends to the operator integral with the symbol
$$e^{-\frac{\mathsf l}2(E+E')}\int\mathsf l^{-1}\rho_1(\varsigma/\mathsf l)e^{-\varsigma(E-E')}d\varsigma\int_0^\infty\frac{dt}{(t+e^{-\mathsf lE})(t+e^{-\mathsf lE'})}=\mathsf G(\mathsf l(E-E')),$$
minus $\mathfrak m$; and $\phi(g_{\mathsf l}\otimes h')$ has the same symbol: this is the slope. Taylor's formula for $u_\epsilon(x)=\langle\psi,L_\epsilon(x)\psi\rangle$ at $x=\pm r$, with $L_\epsilon\ge-\log(Q+1)$ and $L_\epsilon(0)\le\mathsf lH$, gives 
\[
r|u'_\epsilon(0)|\le\mathsf l\langle H\rangle_\psi+\log(Q+1)\|\psi\|^2+\frac{r^2}2\max\{\sup_{|x|\le r}u_\epsilon'',0\};
\]
let $\epsilon\downarrow0$ and use (v); extend from $\bigcup_\Lambda\mathcal H_\Lambda$ to $\mathcal D_\infty$ by continuity in a norm $\|N^a\cdot\|$. For every real $k\in\mathcal S(\mathbb R)$, a dyadic partition $\chi_i$ of the energy transfer and the functions $\mu_i$ with $\hat\mu_i=\hat k\chi_i/\mathsf G(\mathsf l_i\,\cdot)$, $\|\mu_i\|_{L^1}\le C_N\mathsf l_i^N$, write $k=\sum_i\mu_i*g_{\mathsf l_i}$ in $L^1$, and Lemma~\ref{C-363-19}(c) sums the bounds. (h2) By Lorentz covariance on $\mathcal D_\infty$ the bound of (b1) holds in every frame for vectors with $n\cdot p\le\mu$: boost the two spectral sets to the rest frame of $n$, where the argument of (b1) gives $e^2C_\wp\|u\|_{L^1}\mu^\wp$.

The set $\{p\cdot p\le\mu^2\}$ is cut into cells: $\{p\cdot p\le\mu^2,\ E<2\mu\}$ and, for $a\ge1$, the cells $\{p\cdot p\le\mu^2,\ 2^a\mu\le E<2^{a+1}\mu\}$ on which the direction of $\vec p$ lies in one of $\lceil2\pi2^a\rceil$ equal sectors, of angle about $2^{-a}$. The frame of such a cell has the rapidity $a\log2$ in the direction of the middle of its sector, and $n\cdot p<3\mu$ on the cell. For two cells at the distance $D$ in momentum space let $n$ be the midpoint of the hyperbolic segment between their frames, whose length is $\mathsf d$: then $n\cdot p\le3\mu e^{\mathsf d/2}$ on both cells, and $e^{\mathsf d}\le4(1+2D/\mu)^2$, so $n\cdot p\le6\mu+12D$ there. A test function whose Fourier transform vanishes for $|k|\ge\frac34D$ has a vanishing block between the two cells, the spectral projections being smeared translations; so, cutting $\hat u=\hat\kappa(\cdot/D)\hat u+(1-\hat\kappa(\cdot/D))\hat u$ with a fixed smooth $\hat\kappa$ equal to $1$ on $\{|k|\le\frac12\}$ and to $0$ on $\{|k|\ge\frac34\}$, the block is that of a function of $L^1$ norm $\le c\,D^{-8}\|\mathcal L^4u\|_{L^1}$.

For $r\mu\le D<(r+1)\mu$, $r\ge1$, the block is therefore $\le e^2C_\wp\,c\,(r\mu)^{-8}\|\mathcal L^4u\|_{L^1}(6\mu(2r+3))^\wp$, and for $D<\mu$ it is $\le e^2C_\wp\|u\|_{L^1}(18\mu)^\wp$. At most $100(r+1)$ cells lie at a distance $<(r+1)\mu$ from a given cell; Schur's test over the cells, with the weight $r+1$ of this count, gives the bound with the constant $C_{\rm L}$, the sum over $r$ converging for $\wp<6$, with no dependence on the bound of the energy. For these two consequences let $\wp>0$. On $P_m\mathcal H\cap\mathcal P_<\mathcal H$ the invariant mass is $\le(m+n+2)^{1/\wp}$, which gives the bound of the blocks. For the last statement, $\Pi P_m=\mathcal P_<\Pi P_m$ and $P_n\Pi=P_n\Pi\mathcal P_<$ with the $\mathcal P_<$ of the pair $(m,n)$, since $p\cdot p\le(1+p^0)^{2/\wp}\le(m+n+2)^{2/\wp}$ on both ranges; $\Pi$ commutes with $H$, so the argument of (b2) gives $(|m-n|-1)^j\|\Pi P_m\phi(u)P_n\Pi\|\le\|\Pi P_m\phi(\partial_t^ju)P_n\Pi\|$; the blocks of $\Pi e^{-sH}\phi(u)e^{-sH}\Pi$ are then at most $c\,e^{-s(m+n)}(m+n+2)\min\{1,(|m-n|-1)^{-2}\}$, and Schur's test with $\sup_{x\ge0}(x+2)e^{-sx}\le e/s$ gives the bound.

The bound by $1+H+M^2$: cut momentum space into unit cubes; for two cubes at the sup-distance $d$, $D=\sqrt3(d+1)$, and a point $p_0$ of energy $E_0$ and invariant mass $m_0$ in one of them, a timelike vector $a$ with $|a|=(\sqrt2\mathsf M_1)^{-1}$, $\mathsf M_1^2\le2D^2(1+E_0+m_0^2)$, has $a\cdot p\le1$ on both cubes, and the point bound in the frame of $a$ bounds the block of $\phi(f)$ between the two cubes by $e^2C_\wp\|f\|_{L^1}(2D)^\wp$ times the least value of $(1+p^0+p\cdot p)^{\wp/2}$ on either cube; derivatives in a coordinate along which the cubes are $d$ apart give the decay $(d-1)^{-k}$ as in (b2), and Schur's test with the weights $(1+p^0+p\cdot p)^{\wp/2}$ of the cubes gives the form bound.

\subsection{Proof of Theorem~\ref{C-363-4}}

\emph{Step 1.} By Theorem~\textup{\YMMref{AP}{780}} among the construction results on the limits of periodic Gibbs states, $\sigma$ is data (Rec) with $(16C^g_{\rm pt},4,\frac12)$ and $S^\sigma_1=0$, and its truncated functions $S^\sigma_n$, $n\ge2$, are the limits of the joint cumulants $S^T_{n;j}$ of the $\mathcal O_{K_j}(f)$ under $\mu_j=\mu_{\beta_j}$ on $\Lambda_j$, with $\ell_j/a_j=2L^{K_j+m_j}$ even, $a_j=L^{-K_j}$ the spacing of $\Lambda_j$; every component is data (Rec) with one triple.

\emph{Step 2 (frames).} Let $\mathcal H_j,\iota_{\Lambda_j},\mathcal M_j,\pi_j,K_{\Lambda_j},\mathsf P_j$ be the objects of Proposition~\ref{C-363-14} on $\Lambda_j$ at $\beta_j$, and, for tuples $G=(h_k)$ over $\mathbb E_+$, $(G)_j=\prod_k(\mathcal O_{K_j}(h_k)-\mu_j(\mathcal O_{K_j}(h_k)))$, the functions $h_k$ periodised on $\Lambda_j$, and $\iota_jG=\iota_{\Lambda_j}((G)_j)$, both extended linearly. In Steps 2--7 the sign $(\cdot)_j$ is this product on the torus $\Lambda_j$, centred by $\mu_j$; it is not the product $(X)^{\nu_j}_{K_j}$ on $L^{-K_j}\mathbb Z^3$ of the setting, and for a tuple $G$ over a quadrant it is, for large $j$, a real function of the bonds of the closed quadrant of $\Lambda_j$, with $\mathsf r(G)_j=(\mathsf rG)_j$ for $\mathsf r\in\{\theta,\theta_1,R\}$ ($\mathsf r\mathcal O(f)=\mathcal O(\mathsf rf)$, $\mu_j$ being invariant).

For large $j$ the supports lie in $(-\frac{\ell_j}4,\frac{\ell_j}4)^3$ and farther than $2a_j$ from the time-zero plane. For $X=\theta G'\cdot G=(f_k)$, (L1) and the invariance of the weight give $\langle\iota_jG',\iota_jG\rangle=\mu_j\prod_k(\mathcal O(f_k)-\mu_j\mathcal O(f_k))=\sum_\varpi\prod_{\mathsf b\in\varpi}S^T_{|\mathsf b|;j}(X_{\mathsf b})$, the sum over the partitions without singletons (a moment of centred variables, and a cumulant of order $\ge2$ does not see the centring); by Step 1 it tends to $\sum_\varpi\prod S^\sigma_{|\mathsf b|}(X_{\mathsf b})=\sigma(X)=\langle[G'],[G]\rangle$. So $(\mathcal H_j,\iota_j)$ is a frame for $(\mathcal H(\sigma),[\cdot])$, and Lemma~\ref{C-363-8} applies.

\emph{Step 3 (limits of half-plane operators lie in $\mathcal C(W')$).} Let $Q_j\in\mathcal M_j$, $Q_j\rightharpoonup Q$ along a subsequence; take the $\tilde q_j$ of (L3), $\|\tilde q_j\|\le\sup\|Q_j\|$, and a further subsequence with $\tilde q_j\rightharpoonup\tilde q$. For tuples $G_1,G_2$ over $\mathsf{II}$ and large $j$, the $(G_i)_j$ are real functions of the bonds of $\bar{\mathsf Q}_2$, and $R^{-1}(\theta(G_1)_j\cdot(G_2)_j)=(R^{-1}(\theta G_1\cdot G_2))_j$; the limit of (L3) is the hypothesis of Lemma~\ref{C-363-11}, so $Q\in\mathcal C(W')$.

\emph{Step 4 (the lattice modular flows converge to the boosts).} Every subsequence of the support projections $e_{R,j}\in\mathcal M_j$ has a further one with $e_{R,j}\rightharpoonup\mathsf E$; by Step 3, $\mathsf E\in\mathcal C(W')$, and by Lemma~\ref{C-363-8}(iii), $\mathsf E\Omega$ is the weak limit of $e_{R,j}\Omega_{\Lambda_j}=\Omega_{\Lambda_j}=\iota_j\varnothing$: $\mathsf E\Omega=\Omega$. $\Omega$ separates $\mathcal C(W')$, so $\mathsf E=1$ and $\|(1-e_{R,j})\iota_jG\|^2=\langle\iota_jG,(1-e_{R,j})\iota_jG\rangle\to0$ along the full sequence; by (L1), $\|(1-e_{L,j})\iota_jG\|=\|(1-e_{R,j})\iota_j(\theta_1\bar G)\|\to0$. So Lemma~\ref{C-363-8}(iv) applies to $\pi_j$, and by (L4) (for $G\in\mathcal V_0$ and large $j$, $(G)_j$ is a function of $\bar{\mathsf Q}_1$ with $R(G)_j=(RG)_j$) and Lemma~\ref{C-363-9}(b), Lemma~\ref{C-363-8}(v) applies to $\mathsf P_j$, $\Delta^{1/4}$, $\mathsf r=R$:
$$\Phi(\mathsf P_j)\psi_j\Rightarrow\Phi(\Delta^{1/4})\psi\quad(\Phi\in C_0(\mathbb R)),\qquad e^{-2\pi itK_{\Lambda_j}}\psi_j\Rightarrow U(B(-2\pi t))\psi,$$
the second from the even functions $\Phi(x)=|x|^{4it}\chi_{\mathsf M}(x)$, $\chi_{\mathsf M}$ a cut-off near $0$ and $\infty$, since $|\mathsf P_j|=e^{-\frac\pi2K_{\Lambda_j}}$ and $\Delta^{1/4}$ is injective.

\emph{Step 5 (the field on the vacuum is reached).} For real $u\in C^\infty_c(\mathsf I)$ let $\Phi(x)=\mathsf F(-\frac2\pi\log|x|)$, $\Phi(0)=0$: an even function of $C_0(\mathbb R)$ with $\Phi(\mathsf P_j)=\mathsf F(K_{\Lambda_j})$ and $\Phi(\Delta^{1/4})=\mathsf F(K)$. Step 4 and \eqref{C-363-10} give $\xi_j=\mathsf F(K_{\Lambda_j})\pi_j\iota_j(u)\Rightarrow\phi(f_u)\Omega$. (L5) with $F=(u)_j$ gives $T^{(\epsilon)}_j\in\mathcal M_j$; along a subsequence $T^{(\epsilon)}_j\rightharpoonup T^{(\epsilon)}_u$, which lies in $\mathcal C(W')$ by Step 3, and by Lemma~\ref{C-363-8}(i), (iii) and Step 2
\begin{equation}\label{C-363-15}
\|T^{(\epsilon)}_u\Omega-\phi(f_u)\Omega\|\le\tfrac{\sqrt\epsilon}2\,\pi e^{5\pi^2/8}\|[\theta_1u\cdot u]\|,
\end{equation}
one Euclidean four-point function at separated points. With $\epsilon\to0$, $\phi(f_u)\Omega$ lies in the closure of $\mathcal C(W')\Omega$.

\emph{Step 6 (density, locality).} By Lemma~\ref{C-363-9}(c) the transforms of the $f_u$ have no common real zero; Theorem~\ref{C-363-3}(d), with $W$ and $W'$ exchanged, gives the density of $\mathcal C(W')\Omega$, (W), (CP), the duality for wedges and the modular data; Theorem~\ref{C-363-3}(a), (e) give (a)--(e) for $\tau=\sigma$; $\mathcal A(W)=\mathbb C1$ would make $\mathcal H=\overline{\mathcal A(W)\Omega}=\mathbb C\Omega$.

\emph{Step 7 (the wedge algebra is made of lattice limits).} Let $\mathfrak M$ be the von Neumann algebra generated by the operators of (f) and $\mathcal N=\mathcal C(W')=\mathcal A(W)$. For $\mathfrak gW\subset W$ one has $\mathfrak gW'\supset W'$, so $U(\mathfrak g)\mathcal C(W')U(\mathfrak g)^{-1}=\mathcal C(\mathfrak gW')\subset\mathcal C(W')$, and $\mathfrak M\subset\mathcal N$ by Step 3. The closure of $\mathfrak M\Omega$ is invariant under $\mathfrak M$ and under the $U(a)$, $a\in W$, hence under all translations, and it holds every $\phi(f_u)\Omega$ (Step 5); the one-vector argument of Theorem~\ref{C-363-3}(d), with $\mathfrak M$ in the place of $\mathcal C(W')$, makes $\mathfrak M\Omega$ dense. $\mathfrak M\Omega$ and $\mathcal N\Omega$ are boost-invariant dense subspaces of $D(\Delta^{1/2})$ on which $Q\Omega\mapsto Q^*\Omega$ is the restriction of $J\Delta^{1/2}$ (Theorem~\ref{C-363-3}(c) for $W'$), hence cores of it; so the two algebras have the one modular conjugation $J$, and $\mathfrak M'=J\mathfrak MJ\subset J\mathcal NJ=\mathcal N'$: $\mathcal N\subset\mathfrak M$.

\emph{Step 8 (components).} All $\sigma_\xi$ are data (Rec) with one triple. For $a=\sum f^nY_n$ as in (CP), which is hermitian, and $b$ formed as in Lemma~\ref{C-363-13}(iv), a polynomial in means of translates of elements $b_i$ of $\mathfrak P$ over pairwise different spatial translations of length $\ge R$, with $R$ so large that the translated $b_i$ are spacelike to $K$ and so in the algebra of (CP), $\sum\lambda^nb^*Y_nb\in\mathfrak A^+$ for all real $\lambda$ and $f$ is central, so $\mathcal W_\sigma(b^*ab)\ge0$ by Step 6; Lemma~\ref{C-363-13}(iv) gives (CP) for every $\sigma_\xi$, and Theorem~\ref{C-363-3}(d), (a), (e) give (a)--(e) for $\tau=\sigma_\xi$. The same follows from the vectors $M_g\Omega$ of Lemma~\ref{C-363-13}(i): $M_g\in\mathcal C(O)$ for every $O$, and the extension of the proof of (W)$\Rightarrow$(CP) to the vector $M_g\Omega$ gives $\int g^2\mathcal W_\xi(a)\,d\lambda\ge0$ for all real $g\in C(\mathsf X)$.

\emph{Step 9 (fibres).} The operators $M_g$ of Lemma~\ref{C-363-13}(i) commute with the semigroup, the translations and the fields, so they lie in $\mathcal C_\sigma(O)$ for every $O$, hence in $\mathcal C_\sigma(W')\cap\mathcal C_\sigma(W)=\mathcal A_\sigma(W)\cap\mathcal A_\sigma(W)'$ (Step 6). The rest is reduction theory: on a separable Hilbert space that is a direct integral over $(\mathsf X,\lambda)$, $\mathsf X$ compact metrisable and $\lambda$ a Borel measure, the decomposable operators are the commutant $\mathcal Z'$ of the diagonal algebra $\mathcal Z$; a von Neumann algebra $\mathcal N$ with $\mathcal Z\subset\mathcal N\subset\mathcal Z'$ is $\int^\oplus\mathcal N_\xi\,d\lambda$ for a measurable field of von Neumann algebras, whose fibres are those of a countable generating family and are fixed up to a $\lambda$-null set; and $(\int^\oplus\mathcal N_\xi\,d\lambda)'=\int^\oplus\mathcal N_\xi'\,d\lambda$. (1) The classes depend continuously on the entries, with constants of one triple for $\sigma$ and for all $\xi$; so $\mathcal H(\sigma)$ is separable, and there are countably many combinations $F_n$ whose classes $[F_n]_\xi$ are total in every $\mathcal H(\sigma_\xi)$; and $\xi\mapsto\sigma_\xi(\theta\bar F'\cdot F)$ is continuous: the fields $\xi\mapsto[F_n]_\xi$ are a fundamental sequence.

The Hilbert space $\mathfrak K$ of the proof of Lemma~\ref{C-363-13}(i) is the direct integral $\int^\oplus\mathcal H(\sigma_\xi)\lambda(d\xi)$: $g\otimes F\mapsto(g(\xi)[F]_\xi)_\xi$ is isometric by the form of $\mathfrak K$, and onto, as a field orthogonal to all $g[F_n]$, $g\in C(\mathsf X)$, has $\langle[F_n]_\xi,x_\xi\rangle=0$ for almost every $\xi$ and all $n$. So $[F]\mapsto([F]_\xi)_\xi$ is unitary, $\mathcal H(\sigma)=\int^\oplus\mathcal H(\sigma_\xi)\lambda(d\xi)$, and $M_g$ is the multiplication by $g$; the diagonal algebra $\mathcal Z$ is the weak closure of the $M_g$. (2) The semigroup, the translations and the insertions decompose, being given by one formula for $\sigma$ and for all $\sigma_\xi$; so do the bounded $e^{-sH}\phi(f)e^{-sH}$, by \eqref{C-363-2}; as the $e^{-sH_\xi}$ are injective and the polynomial bound holds with one triple, a subsequence converging fibrewise gives $(\phi(f)\psi)_\xi=\phi_\xi(f)\psi_\xi$ for almost every $\xi$, $\psi\in\mathcal D_\infty$, and by induction the fibres of $\phi(f_1)\cdots\phi(f_n)\Omega$ are $\phi_\xi(f_1)\cdots\phi_\xi(f_n)\Omega_\xi$ almost everywhere. (3) $M_g\in\mathcal A_\sigma(W)'$ and weak closure give $\mathcal Z\subset\mathcal A_\sigma(W)'$, so $\mathcal A_\sigma(W)\subset\mathcal Z'$ consists of decomposable operators; $M_g\in\mathcal A_\sigma(W)$ gives $\mathcal Z\subset\mathcal A_\sigma(W)$: $\mathcal Z$ lies in the centre of $\mathcal A_\sigma(W)$, and $\mathcal A_\sigma(W)=\int^\oplus\mathfrak A_\xi\,d\lambda$ for a measurable field of von Neumann algebras $\mathfrak A_\xi$. (4) A decomposable bounded $Q$ lies in $\mathcal C_\sigma(W')$ if and only if $Q_\xi\in\mathcal C_{\sigma_\xi}(W')$ for almost every $\xi$, and the same with $W$ for $W'$: ``if'' by integration of the relation of the fibres, by (2); ``only if'' because $\mathcal C_\sigma(W')$ is an algebra that holds the $M_g$, so that $\int g\,[\langle\phi_\xi(\bar f)\psi_\xi,Q_\xi\psi'_\xi\rangle-\langle\psi_\xi,Q_\xi\phi_\xi(f)\psi'_\xi\rangle]\,d\lambda=0$ for all $g\in C(\mathsf X)$; the bracket vanishes almost everywhere, with one null set for countably many test functions $f\in\mathcal F(W')$ and products of fields, and the polynomial bound with one triple carries it to all. (5) A countable family that generates $\mathcal A_\sigma(W)=\mathcal C_\sigma(W')$ has, by (4), its fibres in the von Neumann algebra $\mathcal C_{\sigma_\xi}(W')$ for almost every $\xi$, and these fibres generate $\mathfrak A_\xi$: $\mathfrak A_\xi\subset\mathcal C_{\sigma_\xi}(W')$ almost everywhere. (6) $\mathcal A_\sigma(W)'=\mathcal C_\sigma(W)=\int^\oplus\mathfrak A_\xi'\,d\lambda$, and (4) for $W$, read on a countable generating family, gives $\mathfrak A_\xi'\subset\mathcal C_{\sigma_\xi}(W)$ almost everywhere; with wedge duality for each $\sigma_\xi$ (Step 8), $\mathfrak A_\xi\supset\mathcal C_{\sigma_\xi}(W)'=\mathcal C_{\sigma_\xi}(W')$.

So $\mathfrak A_\xi=\mathcal C_{\sigma_\xi}(W')=\mathcal A_{\sigma_\xi}(W)$ for $\lambda$-almost every $\xi$, every $\xi$ with $\lambda(\{\xi\})>0$ among them, as such a point is no null set, and $\mathcal A_\sigma(W)=\int^\oplus\mathcal A_{\sigma_\xi}(W)\,\lambda(d\xi)$: for $\lambda$-almost every $\xi$ the wedge algebra of the component $\sigma_\xi$ is the fibre of the algebra of (f), the fibres of a direct integral being fixed up to null sets.

\emph{Step 10 (the supply).} Here $(F)_j=(F)^{\rho_j}_{K_j}$ is the product at $K_j$ centred by $\rho_j$, as in (h), also for a component, where the states $\rho_{j,\xi}$ integrate and the $\rho_j$ give the centring. Reflection positivity of $\rho_j$ at the site mirror $x_3=0$: a bounded Borel function $B$ of bonds in the mirror is fixed by $\theta$, and with $\mathsf c=(\sup|B|^2-|B|^2)^{1/2}$, $\sup|B|^2\rho_j(\overline{\theta G}G)-\rho_j(\overline{\theta(BG)}BG)=\rho_j(\overline{\theta(\mathsf cG)}\mathsf cG)\ge0$; the Schwarz inequality gives the first bound. $\mathcal H(\sigma')$ is separable; a diagonal subsequence over a countable set of combinations with dense classes, and the first bound, give the limit form and $Q$, $\|Q\|\le\mathsf b$. For $c=0$ and combinations $F,F'$ over $\mathsf Q_L$ put $A_j=B_j\overline{\theta(F')_j}(F)_j$; then $\ell^Q_{F',F}(h)=\lim_j\rho_j(A_j(h)_j)$, and the Schwarz inequality for the reflection positivity at the site mirror $x_1=0$ gives
$$|\rho_j(A_j(h)_j)|^2\le\rho_j(A_j\,\theta_1\bar A_j)\,\rho_j(\overline{\theta_1(h)_j}(h)_j),$$
whose second factor tends to $\sigma'(\theta_1\bar h\cdot h)$; in the first, with $Y=F\cdot\theta_1\bar F$, $Y'=F'\cdot\theta_1\bar F'$ and the function $B_j\theta_1\bar B_j$ of the bonds of the mirror $x_3=0$, $A_j\theta_1\bar A_j=\overline{\theta(Y')_j}\,B_j\theta_1\bar B_j\,(Y)_j$, and the first mirror, used a second time, bounds it by $\mathsf b^2\|[Y']\|\|[Y]\|$ in the limit. So the hypothesis of Lemma~\ref{C-363-17} holds and $Q\in\mathcal C_{\sigma'}(W)$. For $c\ne0$, $ce_1$ is a lattice vector for large $j$; the translate $\tau_{-ce_1}B_j$ is of the case $c=0$, and by the invariance of $\rho_j$ its limit is $U(ce_1)^{-1}QU(ce_1)$. The proof uses of the states that integrate only reflection positivity at the site mirrors $x_3=0$, $x_1=0$ and invariance under the lattice translations, which the lattice approximants of a component have; so the same holds for every component.

Theorems~\YMMref{D}{8}, \YMMref{AJ}{673} and~\YMMref{AD}{607} enter only through Theorem~\textup{\YMMref{AP}{780}}.

\subsection{Proof of Theorem~\ref{C-363-5}}

Every bound below passes to the components through the last sentence of Lemma~\ref{C-363-16}(d), which carries the norm bound of every insertion from $\sigma'$ to each component; every component is data (Rec) with one triple.

\emph{(a)} Step 1: the blocks at one time. Let $h\in C^1_c(\mathbb R^2)$ be real and carried by a square of side $2$, $q=T_{mn}(h)$, $A=(H-m)P_m$, $B=(H-n)P_n$, so that $0\le A,B\le1$. As $P_m\mathsf B_h(s,s')P_n=e^{-sH}q\,e^{-s'H}$,
$$e^{\frac{\mathsf l}2(m+n)}P_m\mathsf I_{\mathsf l}(h)P_n=\int\rho_{\mathsf l}(\varsigma)e^{-\varsigma(m-n)}\,e^{-(\frac{\mathsf l}2+\varsigma)A}\,q\,e^{-(\frac{\mathsf l}2-\varsigma)B}\,d\varsigma ,$$
and $\|e^{-uA}qe^{-u'B}-q\|\le(u+u')\|q\|$ with $u+u'=\mathsf l$; since $\int\rho_{\mathsf l}(\varsigma)e^{-\varsigma(m-n)}d\varsigma=\int\rho_{\mathsf l}(\varsigma)\cosh(\varsigma(m-n))d\varsigma\ge1$, $\|q\|\le(1-\mathsf l)^{-1}e^{\frac{\mathsf l}2(m+n)}\|\mathsf I_{\mathsf l}(h)\|$. Take the least $i\ge1$ with $L^{-i}\le2/(m+n+2)$: then $e^{\frac{\mathsf l}2(m+n)}\le e$ and $\mathsf l^{-1}\le L(m+n+2)$, and Lemma~\ref{C-363-16}(c) gives
$$\|T_{mn}(h)\|\le\tfrac{3e}2L^3\mathsf C_3(m+n+2)^3|h|_{C^1}.$$
Step 2: the field. By \eqref{C-363-2}, $\langle\chi',P_m\phi(f)P_n\chi\rangle=\int\langle e^{-itH}\chi',T_{mn}(f_t)e^{-itH}\chi\rangle\,dt$ with $f_t=f(t,\cdot)$; a partition of unity on $\mathbb R^2$ subordinate to the squares $\kappa+[-1,1]^2$, $\kappa\in\mathbb Z^2$, splits the real and imaginary parts of $f_t$ into pieces of norm $|\cdot|_{C^1}\le3.3\sup_{\kappa+[-1,1]^2}(|f_t|+|\nabla f_t|)$, and $2\cdot3.3\cdot\frac{3e}2<27$ gives $b_{mn}(f)\le\mathsf c_0(m+n+2)^3$; Theorem~\ref{C-363-3}(b2), which holds for all data, gives the decay. Step 3: the operator and form bounds. $\|P_m\phi(f)\psi\|\le\sum_nb_{mn}(f)\|P_n\psi\|$, and $m+n+2\le2(1+n)(1+|m-n|)$; so $\|\phi(f)\psi\|\le\|\kappa\|_{\ell^1}\|N^3\psi\|$ by Young's inequality, $\kappa(d)=8(1+|d|)^3\min\{\mathsf c_0,\mathsf c_5(|d|-1)^{-5}\}$, and $\sum_{|d|\le2}(1+|d|)^3=71$, $\sum_{|d|\ge3}(1+|d|)^3(|d|-1)^{-5}<6.3$; with a weight $(1+m)^a\le(1+n)^a(1+|d|)^{|a|}$ and a larger $k$ the same gives the bound with $N^{a+3}$.

For the form, $(m+n+2)^2\le4(1+m)(1+n)(1+|d|)$ and Schur's test with $\kappa'(d)=8(1+|d|)^{3/2}\min\{\mathsf c_0,\mathsf c_3(|d|-1)^{-3}\}$, whose two sums are less than $17.1$ and $5.7$, respectively. The blocks of $e^{-sH}\phi(f)e^{-sH}$ are $\le e^{-s(m+n)}b_{mn}(f)$, and $\sup_{x\ge0}(x+2)^3e^{-sx}\le(3/s)^3e^{2s}$. Step 4: growth. For $|m-n|=d\ge2$, the bound of Step 2 at $\partial_t^kf$ with $\|\partial_t^kf\|_\star\le A^{k+1}(k!)^{\mathsf s}\le A^{k+1}k^{\mathsf sk}$, at $k=\lfloor((d-1)/e^{\mathsf s}A)^{1/\mathsf s}\rfloor$, gives
$$b_{mn}(f)\le\mathsf G\,(m+n+2)^3e^{-\gamma'|m-n|^{1/\mathsf s}}\qquad\text{for all }m,n,\qquad\gamma'=\mathsf se^{-1}A^{-1/\mathsf s},\ \mathsf G=27L^3\mathsf C_3Ae^{\mathsf s+\gamma'} .$$
As $\Omega=P_0\Omega$, $\|\phi(f)^n\Omega\|\le\sum\prod_{k\le n}b_{m_km_{k-1}}(f)$ over paths $0=m_0,m_1,\dots,m_n$; with $\mathsf D=\sum_k|m_k-m_{k-1}|\ge\max_km_k$, at most $2^n$ paths for each sequence of steps, and $\sum_k|d_k|^{1/\mathsf s}\ge\mathsf D^{1/\mathsf s}$,
$$\|\phi(f)^n\Omega\|\le(2\mathsf G)^n\sup_{\mathsf D\ge0}(2\mathsf D+2)^{3n}e^{-\frac{\gamma'}2\mathsf D^{1/\mathsf s}}\Bigl(\sum_{d\ge0}e^{-\frac{\gamma'}2d^{1/\mathsf s}}\Bigr)^n\le C^n(n!)^{3\mathsf s}.$$
Such $f$ exist: $w\otimes b$ with $b\in C^\infty_c(\mathbb R^2)$, $b\ge0$, and $w$ the infinite convolution of the uniform probability densities on $[0,i^{-\mathsf s}]$, $i\ge1$, for which $\|w^{(k)}\|_{L^1}\le\prod_{i\le k}2i^{\mathsf s}=2^k(k!)^{\mathsf s}$, because the $k$-th derivative falls on the first $k$ factors, each becoming a difference of two point masses divided by the width.

\emph{(b)} The same four steps with the second bound of Lemma~\ref{C-363-16}(c), $\|T_{mn}(h)\|\le\frac{3e}2L^2\mathsf C_2(m+n+2)^2|h|_{C^1}$, the exponent $2$ in the place of $3$, the weight $(1+n)^2$ in Step 3 for the operator bound and $[(1+m)(1+n)]^{1}$ for the form bound, $\sup_x(x+2)^2e^{-sx}\le(2/s)^2e^{2s}$, and $(n!)^{2\mathsf s}$ at the end.

\emph{(c)} \emph{The point.} Put
$\bar B=B_\sharp L/(L-1)$, where $B_\sharp\ge0$ is one of the fixed
constants $\sigma^-$. Thus $\bar B$ is independent of the source
radius, $\gamma$ and $g$, and $\bar B<\gamma^{-2}\le g^{-2}$.
For an integer $z\ge2$ set
\[
 v_z^{-2}=L^{z-1}(g^{-2}+B_\sharp-L^{-1}\bar B)
                       +L^{-1}\bar B-B_\sharp,
 \qquad G=g^{-2}+2\bar B.
\]
Then $L^{z-1}g^{-2}\le v_z^{-2}\le L^zG$. Define
\[
 \Lambda_z=1+\log(L^zG),\qquad
 \varrho_z=\frac12\min\{
   \sigma^{\rm S}(v_z^{-2}-\bar B),\,
   \sigma^{\rm S}_{\rm L}v_z^{-1}
                         (\log v_z^{-2})^{-\mathfrak p_0}\}.
\]
Here $\sigma^{\rm S}$ and $\sigma^{\rm S}_{\rm L}$ depend only on
$\sigma^-$. We check explicitly that this radius is allowed.
The radius-independent conditions remain valid when the coupling
ceiling is decreased from $\gamma$ to $v_z$.
The two source-radius conditions in
Proposition~\YMMref{82.2}{668}, with source radius $2\varrho_z$, are
\[
 2\varrho_z\le\sigma^{\rm S}(x-\bar B),\qquad
 2\varrho_z x^{-1/2}(\log x)^{\mathfrak p_0}
                              \le\sigma^{\rm S}_{\rm L}
                 \quad(x\ge v_z^{-2}).
\]
The first left margin is smallest at $x=v_z^{-2}$. The calibration
also gives $\log\gamma^{-2}>2\mathfrak p_0$. Therefore
$x^{-1/2}(\log x)^{\mathfrak p_0}$ is decreasing for
$x\ge v_z^{-2}\ge\gamma^{-2}$, and the second condition is likewise
strongest at that endpoint. Both conditions hold by the definition
of $\varrho_z$. This proves the required admissibility and also
identifies $\varrho_z$ as the largest radius permitted by these two
source conditions at the fixed ceiling $v_z$.

Let $(K,\beta)$ be admitted at $g$, $K>z$. Its label $\mathsf b=L^{-K}\beta$ lies in the window \YMMref{(1531)}{615} of Proposition~\YMMref{74.1}{615}, $g^{-2}+B_\sharp\frac{L-2}{L-1}<\mathsf b\le L(g^{-2}+\bar B)$; so the running coupling $g'$, $g'^{-2}=L^{-1}(L^z\mathsf b+\bar B)-B_\sharp$, at which the label $L^z\mathsf b$ of $(K-z,\beta)$ passes first, has $v_z^{-2}<g'^{-2}\le L^zG$ and $1+\log g'^{-2}\le\Lambda_z$. At the radius $\varrho_z$ and the choice $(\sigma^-,v_z)$ the dilated pair $(K-z,\beta)$ is a first-passage datum at $g'\le v_z$, and it carries systems of localised source energies with the constants $(2\varrho_z;C_{\rm S},\kappa;C_{\rm S},C_{g'};C_{\rm gap})$ (Lemma~\ref{lem:canonical-source-system}), where $C_{\rm S},\kappa,C_{\rm gap}$ are the numbers of the pairs admitted at $g$, formed from $\sigma^-$ alone, and $C_{g'}\le C_{\rm S}(1+\log g'^{-2})\le C_{\rm S}\Lambda_z$ \eqref{eq:setup-pressure-constants}.

As $v_z^{-2}-\bar B\ge\frac{12}{13}L^{z-1}g^{-2}$, $v_z\le gL^{-(z-1)/2}$, $\log v_z^{-2}<\Lambda_z$ and $g\le\frac12$,
$$\frac1{\varrho_z}\le2c_\sigma\,g\,L^{-(z-1)/2}\Lambda_z^{\mathfrak p_0},\qquad\Lambda_z=1+z\log L+\log G,\qquad G=g^{-2}+2\bar B,$$
$c_\sigma=\max\{\frac{13}{24\sigma^{\rm S}},\frac1{\sigma^{\rm S}_{\rm L}}\}$, formed from $\sigma^-$ alone: the radius of analyticity in the source grows like $1/g'$, the square root of the label $g'^{-2}$, up to logarithms. This is where Theorems~\YMMref{AJ}{673} and~\YMMref{AD}{607} are used at every radius and every admitted choice. Fix $0<\eta\le1$, a real $h\in C^1_c$ carried by a ball $B(y,\frac\eta{40L})$ with $|y_3|\le\frac\eta8$, and $F,F'$ over $\{x_3>\frac\eta4\}$. The cell: $\mathsf l=L^{-i}$ with $\frac\eta{5L}<\mathsf l\le\frac\eta5$ and a centre $y^*\in L^{-i-2}\mathbb Z^3$ within $\frac{\mathsf l}{2L^2}$ of $y$; then the ball lies in the middle of the cell, and the cell in $\{|x_3|<\frac\eta4\}$.

The dilation by $L^z$ among the construction results (Lemma~\ref{lem:lattice-law-dilation}) carries periodic Gibbs states, mirrors and boards to the like and $\mathcal O_K(h)$ to $\mathcal O_{K-z}(L^zh(L^{-z}\cdot))$; it carries the cell to one of side $L^z\mathsf l\le L^z\eta/5$, and $h$ to a source of norm $\max\{L^z\sup|h|,40M\,{\rm Lip}\,h\}$, which is $\le\|h\|_\sharp$ as long as $L^z\eta\le5$, since $h$ vanishes on the boundary of its ball, so that $\sup|h|\le\frac\eta{40L}{\rm Lip}\,h$. The insertion bound on one lattice for a pair that carries systems, Lemma~\YMMref{104.4}{760} among the construction results, read on the dilated lattice at the radius $\varrho_z$, has the constant $2\mathsf Q_{\rm mid}(L^z\mathsf l)/\varrho_z$ at the side of the one cell, with $\log\mathsf Q_{\rm mid}(\mathsf l')=(C_{\rm S}+C_{g'}+C_{\rm gap})\mathsf l'^3+\Phi_1(\frac{\mathsf l'}4,\frac{\mathsf l'}4;\mathsf l')$ (equation~\YMMref{(2017)}{760}) and $\Phi_1(\frac{\mathsf l'}4,\frac{\mathsf l'}4;\mathsf l')\le\bar\Phi$ for $0<\mathsf l'\le1$, where
\[
 \bar\Phi=C_{\rm S}S(\kappa)
 \left[\frac{27}{26}
       +27\{32(C''_\Sigma+\log_L10)+4C'_\Sigma\}\right]
\]
is the contribution of $\Phi_1$ in equation~\YMMref{(2018)}{760}; so $\mathsf Q_{\rm mid}(L^z\mathsf l)\le Q_\star=e^{2C_{\rm S}+C_{\rm gap}+\bar\Phi}$, free of $z$, as long as $(L^z\eta/5)^3\Lambda_z\le1$.

This is the proof of Proposition~\YMMref{105.1}{762} up to its last step, with the side of the one cell in the place of the supremum over the sides. With $z$ fixed and the cutoff removed, $|\sigma'(\theta\bar F\cdot(h)\cdot F')|\le\frac{2Q_\star}{\varrho_z}\|h\|_\sharp\|[\tau_{-\frac\eta4e_3}F]\|\|[\tau_{-\frac\eta4e_3}F']\|$ for every $z\ge2$ with $(L^z\eta/5)^3\Lambda_z\le1$. As $z\mapsto(L^z\eta/5)^3\Lambda_z$ increases, such a $z$ exists if and only if $\eta\le\eta_1=5L^{-2}\Lambda_2^{-1/3}$, a number $<L^{-1}$. At the largest such $z$: the condition gives $L^z\le5/\eta$, so $\Lambda_{z+1}=1+\log L+z\log L+\log G\le\mathsf L(\eta)$, $\mathsf L(\eta)=1+\log(5LG)+\log\frac1\eta$; its failure at $z+1$ gives $L^{-z}<\frac{L\eta}5\Lambda_{z+1}^{1/3}$; hence $L^{-(z-1)/2}<\frac L{\sqrt5}\eta^{1/2}\mathsf L(\eta)^{1/6}$, and the insertion bound for a source in a ball of radius $\eta/40L$ holds for $\eta\le\eta_1$ with the constant $A_1\eta^{1/2}\mathsf L(\eta)^q$, $A_1=\frac4{\sqrt5}LQ_\star c_\sigma g$, $q=\mathfrak p_0+\frac16$, formed from $\sigma^-$ and $g$ alone.

For $\eta_1<\eta\le1$ Proposition~\YMMref{105.1}{762} itself, for the undilated pairs, gives the constant $C_{\rm I}=2\bar{\mathsf Q}_{\rm mid}/\varrho$, and $C_{\rm I}\le C_{\rm I}\eta_1^{-1/2}\eta^{1/2}\mathsf L(\eta)^q$ as $\mathsf L(\eta)\ge1$. So for all $0<\eta\le1$ the bound holds with the running constant $\mathsf c_*(\eta)=A\,\eta^{1/2}\mathsf L(\eta)^{q}$, $A=\max\{A_1,C_{\rm I}\eta_1^{-1/2}\}$; the radius $\varrho$ enters $A$ through $C_{\rm I}$ alone. The data dilated by $L^k$, $\sigma'_k=\sigma'\circ\mathsf D_k^{-1}$, where $\mathsf D_k(f_1,\dots,f_n)=(L^kf_i(L^{-k}\cdot))_i$, extended linearly, commutes with $\theta$, the conjugation and the product of tuples, has $\mathsf D_k\tau_v=\tau_{L^kv}\mathsf D_k$ and keeps the norms of the classes, are again data with this property, with $\mathsf c_*(L^{-k}\cdot)$, as $\|L^{-k}h''(L^k\cdot)\|_\sharp\le\|h''\|_\sharp$ for $k\ge0$; and Proposition~\YMMref{97.2}{731} among the construction results, which turns an insertion bound for sources in small balls with a constant $C_{\rm I}$ into the point bound of order $4$ with the constant $c_{\rm pt}C_{\rm I}$, $c_{\rm pt}=320\cdot8^4\,\mathsf{mv}_3(40L)^4M$ the coefficient of $C_{\rm I}$ in equation~\YMMref{(1919)}{729} at $c_{\rm I}=\frac1{40L}$ (with $\mathsf{mv}_3=m_3$ in the notation following that equation), applied to $\sigma'_k$ at $\eta''=L^k\eta\in(L^{-1},1]$ and carried back, gives
\[
\begin{gathered}
|\sigma'(\theta\bar F\cdot(h)\cdot F')|\\
\le16\,c_{\rm pt}\,A\,s_q\,L^{1/2}\,\eta^{-7/2}\mathsf L(\eta)^q\,\|h\|_{L^1}\|[\tau_{-\eta e_3}F]\|\|[\tau_{-\eta e_3}F']\|,\qquad\\
 s_q=\sup_{u\ge0}e^{-u/2}(1+u)^q,
\end{gathered}
\]
for $F,F'$ over $\{x_3>\eta\}$ and $h$ carried by $\{|x_3|<\frac12\eta\}$; this is (Rec4) with $(C_\epsilon,\frac72+\epsilon,\frac12)$, $C_\epsilon=16c_{\rm pt}As_qL^{1/2}\sup_{\eta\le1}\eta^\epsilon\mathsf L(\eta)^q$, and the point bound with the logarithm follows by taking $\epsilon=q/\mathsf L(\mu)$.

\emph{One time.} A chain of dilates $y=0,1,\dots,Y$ of one pair bounds the coefficients of the mixed part of the source energies that see two cells. Let $(K,\beta)$ be admitted at $g$, with the label $\mathsf b=L^{-K}\beta\ge g^{-2}$, $(i,o)$ a board with cells of side $\mathsf l=L^{-i}$, $f\ne0$ a real Lipschitz source vanishing off the middle of one cell, $\ell'$ the side of a reference torus, $v'^\pm$ the sums of the mirror copies of $f$ in the even and in the odd cells of the board on the torus of side $\ell'$, and $\mathsf N=(\ell'/\mathsf l)^3$. The normalised partition function $F(z,z')=Z_{K,\ell',\beta}[zv'^++z'v'^-]/Z_{K,\ell',\beta}$, $Z_{K,\ell',\beta}[w]$ the partition function of the Wilson weight with the source $w$ on the torus $L^{-K}\mathbb Z^3/\ell'\mathbb Z^3$, is entire with $F(0,0)=1$, and under the dilation by $L^y$ it is the same number for the pair $(K-y,\beta)$ on the torus of side $L^y\ell'$ with the arrays $L^yv'^\pm(L^{-y}\cdot)$: one function for all $0\le y<K$.

For $y\ge1$ put $x_y=\frac12L^y\mathsf b$, $g_y=x_y^{-1/2}$, $2\varrho_y=\min\{\sigma^{\rm S}(x_y-\bar B),\sigma^{\rm S}_{\rm L}x_y^{1/2}(\log x_y)^{-\mathfrak p_0}\}$, $\lambda_y=2+\log x_y$, and $\varrho_0=\varrho$, $\lambda_0=2+\log g^{-2}$: at the radius $\varrho_y$ and the choice $(\sigma^-,g_y)$ the pair $(K-y,\beta)$ is a first-passage datum at $g_y$ and carries systems of localised source energies with the constants $(2\varrho_y;C_{\rm S},\kappa;C_{\rm S},C_{g_y};C_{\rm gap})$, $C_{\rm S}+C_{g_y}\le C_{\rm S}\lambda_y$, as in (c). The system of the $y$-th dilate gives a factorisation $F=e^{S_y}T_y$ on the bidisc of radius $r_y=\varrho_y/\|f\|_\sharp$ (the dilated source keeps the norm $\|f\|_\sharp$ while $L^y\mathsf l\le1$): $S_y$, with $S_y(0,0)=0$, is the sum of its source energies at the dilated arrays, and $T_y$ is the remainder normalised at $0$, with $|T_y|\le e^{\mathsf NV_y}$, $V_y=C_{\rm S}\lambda_y(L^y\mathsf l)^3$, by the stability bounds of the systems (Section~\YMMref{88}{685}).

If one function on a bidisc has two factorisations $e^{S_1}T_1=e^{S_2}T_2$ with $S_\iota(0,0)=0$, $T_\iota(0,0)=1$, $|T_\iota|\le e^{B_\iota}$, Jensen's inequality on the torus of the bidisc, $\int|\log|T_\iota||\le2B_\iota$ for each factor, bounds the Taylor coefficients of $S_1-S_2$ on a bidisc of radius $r$: $|(S_1)_{ab}-(S_2)_{ab}|\le4(B_1+B_2)r^{-(a+b)}$ for $(a,b)\ne(0,0)$, as they are Fourier coefficients of $2{\rm Re}(S_1-S_2)=2(\log|T_2|-\log|T_1|)$. For the mixed parts per cell, $\mathsf m_y=\mathsf N^{-1}S_y^{\rm mix}$, $S^{\rm mix}(z,z')=S(z,z')-S(z,0)-S(0,z')+S(0,0)$ the sum of the terms $z^az'^b$ with $a,b\ge1$, $\mathsf m_0$ being the mixed part of the proof of Lemma~\ref{C-363-16}, the dilates $y$ and $y+1$ give the chain
$$|\mathsf m_{0;ab}|\le|\mathsf m_{Y;ab}|+\sum_{y=0}^{Y-1}4(V_y+V_{y+1})\min\{r_y,r_{y+1}\}^{-(a+b)}\qquad(a,b\ge1),$$
and in the top dilate Cauchy's estimate on the bidisc of radius $r_Y$ with the two counting estimates in the proof of the cell bound (Proposition~\YMMref{89.1}{690}) gives $|\mathsf m_{Y;ab}|\le\bar\Phi r_Y^{-(a+b)}$. The radii have $\varrho_y\ge\underline\varrho_y=\frac12\sigma_\wedge x_y^{1/2}(\log x_y)^{-\mathfrak p_0}$, $\sigma_\wedge=\min\{\frac{11}{13}\sigma^{\rm S},\sigma^{\rm S}_{\rm L}\}$, and $\underline\varrho_y/\underline\varrho_{y'}\ge L^{-(y'-y)/2}$ for $1\le y\le y'$. As $V_y\le L^{-3(Y-y)}V_Y$ while $\underline\varrho_y^{-k}\le L^{(Y-y)k/2}\underline\varrho_Y^{-k}$, the terms $1\le y<Y$ of the chain form a geometric sum, at most $5.6L^{k/2}V_Y\underline\varrho_Y^{-k}\|f\|_\sharp^k$ for $k=a+b\le5$, and the term $y=0$ is at most $8L^3C_{\rm S}\bar\lambda_1\mathsf l^3\varrho_\wedge^{-k}\|f\|_\sharp^k$, with $\bar\lambda_1=2+\log(\frac12L^2(g^{-2}+\bar B))$ and $\varrho_\wedge=\min\{\varrho,\frac12\sigma_\wedge(\frac12Lg^{-2})^{1/2}(\log(\frac12Lg^{-2}))^{-\mathfrak p_0}\}$.

Take $Y=i-r(\mathsf l)$, $r(\mathsf l)$ the least $r\ge1$ with $L^{3r}\ge\bar\lambda(\mathsf l)=2+\log\frac{L(g^{-2}+\bar B)}{2\mathsf l}$: in the dilate $Y$ the cell has the side $L^{-r(\mathsf l)}$, of the order $\bar\lambda(\mathsf l)^{-1/3}$, so that $V_Y\le C_{\rm S}$ and $\underline\varrho_Y^{-2}\le8L\sigma_\wedge^{-2}g^2\,\mathsf l\,\bar\lambda(\mathsf l)^{2\mathfrak p_0+1/3}$; and $\mathsf lL^2\bar\lambda(\mathsf l)^{1/3}\le1$ gives $1\le Y\le i-1$. The reference tori are those of the sides $2L^m$, $m\ge m_*$, with $L^y$ times the side in the family of sides of the $y$-th dilate for all $y\le Y$, the threshold $m_*$ being formed from $\sigma^-$ and $g$ alone. The two factorisations bound the mixed parts themselves, 
\[
|S_1^{\rm mix}-S_2^{\rm mix}|(z,z')\le4(B_1+B_2)\frac{tt'}{(1-t)(1-t')},
\]
$t=|z|/r$, $t'=|z'|/r$, and $t\le\frac12$ on all the bidiscs of the chain for $|z|\|f\|_\sharp\le\frac12\varrho_\wedge$; this gives, for a source $f$ in the middle of a cell of side $\mathsf l$ with $\mathsf lL^2\bar\lambda(\mathsf l)^{1/3}\le1$ and $|z|\le\frac12\varrho_\wedge/\|f\|_\sharp$,
\[
\begin{gathered}
\|e^{\mathcal O_K(zf)}\|_{\rm ch}\le|e^{\mathcal E_f(z)}|\exp\bigl[c_V\mathsf l^3+\mathsf a(\mathsf l)|z|^2\|f\|_\sharp^2\bigr],\qquad\\
\mathsf a(\mathsf l)=\bigl[3\bar\Phi+22.4LC_{\rm S}\bigr]\frac{8Lg^2}{\sigma_\wedge^2}\,\mathsf l\,\bar\lambda(\mathsf l)^{2\mathfrak p_0+\frac13}\\
+32L^3C_{\rm S}\bar\lambda_1\mathsf l^3\varrho_\wedge^{-2},
\end{gathered}
\]
with $\mathcal E_f$ holomorphic on $|z|<\varrho/\|f\|_\sharp$, $\mathcal E_f(0)=0$, one function for all the reference sides, the normaliser of the one-cell bound that absorbs the unmixed terms. In the multiple reflections for the $n$ cells of one layer, Lemma~\ref{C-363-16}(e) (as in the proof of Lemma~\ref{C-363-16}(c)), Cauchy's estimate at the radius $\min\{\mathsf a(\mathsf l)^{-1/2}\mathcal N^{-1},\varrho_\wedge/2\mathcal N_\infty\}$, $\mathcal N^2=\sum_k\|f_k\|_\sharp^2$, $\mathcal N_\infty=\max_k\|f_k\|_\sharp$, bounds the insertion of the layer by $2e^{nc_V\mathsf l^3+1}[\mathsf a(\mathsf l)^{1/2}\mathcal N+2\mathcal N_\infty/\varrho_\wedge]$; for the pieces of $\rho_{\mathsf l}\otimes k$ this is $\le C\mathsf l^{-5/2}\bar\lambda(\mathsf l)^{\mathfrak p_0+1/6}|k|_{C^1}$, and the limit of the cutoffs gives $\|\mathsf I_{\mathsf l}(k)\|\le C\mathsf l^{-5/2}\bar\lambda(\mathsf l)^{\mathfrak p_0+1/6}|k|_{C^1}$. The operator at one time has no larger norm than the plate: on the vectors of bounded energy $\mathsf B_k(\frac{\mathsf l}2+\varsigma,\frac{\mathsf l}2-\varsigma)=e^{-\varsigma H}\mathsf B_k(\frac{\mathsf l}2,\frac{\mathsf l}2)e^{\varsigma H}$, so the plate is the operator integral of $\mathsf B_k(\frac{\mathsf l}2,\frac{\mathsf l}2)$ with the symbol $\tilde\rho(E-E')$,
$$\tilde\rho(x)=\int\rho_{\mathsf l}(\varsigma)\cosh(\varsigma x)\,d\varsigma=\frac{\sinh(ax)}{ax}\,\frac{\pi^2}{\pi^2+a^2x^2},\qquad a=\frac{\mathsf l}5,\qquad\frac1{\tilde\rho(x)}=\prod_{j\ge2}\Bigl(1+\frac{a^2x^2}{j^2\pi^2}\Bigr)^{-1},$$
and $1/\tilde\rho$ is the Fourier transform of a probability density $\mathsf g$, a product of Fourier transforms of Laplace densities; hence $\mathsf B_k(\frac{\mathsf l}2,\frac{\mathsf l}2)=\int\mathsf g(t)e^{itH}\mathsf I_{\mathsf l}(k)e^{-itH}dt$ and $\|\mathsf B_k(\frac{\mathsf l}2,\frac{\mathsf l}2)\|\le\|\mathsf I_{\mathsf l}(k)\|$. This gives the bound at $s=\frac{\mathsf l}2$, $\mathsf l=L^{-i}$, for $\mathsf lL^2\bar\lambda(\mathsf l)^{1/3}\le1$, Lemma~\ref{C-363-16}(c) gives it for the finitely many larger $\mathsf l=L^{-i}$, and the semigroup property gives it for all $s,s'$. The block bound follows from $\|P_m\phi(f)P_n\|\le e^{s(m+n+2)}\|e^{-sH}\phi(f)e^{-sH}\|\le e\int\|\mathsf B_{f_t}(s,s)\|dt$ at $s=1/(m+n+2)$ and the partition of unity of (a).

\emph{(d)} For a component, $\mathsf B_k(u,u)\Omega$ is the class of the insertion of $k$ at the height $u$, and
$$\|\mathsf B_k(u,u)\Omega\|^2=\iint k(\vec x)k(\vec y)\,\mathsf s\bigl((|\vec x-\vec y|^2+4u^2)^{1/2}\bigr)\,d^2x\,d^2y .$$
With $\mathsf s\ge0$, $\mathsf s(r)\ge c_{\rm sd}r^{-6}$ for $r\le r_{\rm sd}$, and $u^4(|z|^2+4u^2)^{-3}$, an approximate identity on $\mathbb R^2$ of the mass $2\pi u^4\int_0^\infty r(r^2+4u^2)^{-3}dr=\frac\pi{32}$, the liminf follows for continuous $k\ge0$ of compact support. A bound $\pm\mathsf T(k)\le cN$ gives $|\langle\psi',\mathsf T(k)\psi\rangle|\le c\|N^{1/2}\psi'\|\|N^{1/2}\psi\|$ by polarisation, so $\|\mathsf B_k(u,u)\|\le c\sup_E(1+E)e^{-2uE}\le c(1+(2eu)^{-1})$, against the liminf. $\|\mathsf B(s,s)\Omega\|^2=\mathsf s(2s)\ge c_{\rm sd}(2s)^{-6}$; with the upper bound in (c), the infimum of admissible pure-power exponents lies in $[3,\frac72]$. For the plates, $\mathsf I_{\mathsf l}(h)\Omega=\tilde\rho(H)\mathsf B_h(\frac{\mathsf l}2,\frac{\mathsf l}2)\Omega$ with $\tilde\rho(E)=\int\rho_{\mathsf l}(\varsigma)\cosh(\varsigma E)d\varsigma\ge1$, so $\|\mathsf I_{\mathsf l}(h)\Omega\|\ge(1-o(1))(\pi c_{\rm sd}/32)^{1/2}\|h\|_{L^2}(\mathsf l/2)^{-2}$.

For the blocks, put $w_D=T_{D0}(k)\Omega$, and for $0<\eta\le1$ let $\rho=\rho^{(\eta)}\in C^1_c((-\frac\eta4,\frac\eta4))$ be even and $\ge0$ with $\int\rho=1$, and $Y_\eta=\int\rho(\varsigma)\mathsf B_k(\eta+\varsigma,\eta-\varsigma)d\varsigma$, the insertion of $\rho\otimes k$ at the height $\eta$ (an admitted insertion, as $\rho\otimes k$ is carried by $\{|x_3|<\frac12\eta\}$). As $\rho$ is even, $\|P_mY_\eta P_n\|\le e^{-\eta(m+n)}\cosh(\frac\eta4|m-n|)\|T_{mn}(k)\|$, and $P_DY_\eta\Omega=e^{-\eta H}\tilde\rho_\eta(H)w_D$ with $\tilde\rho_\eta(E)=\int\rho(\varsigma)\cosh(\varsigma E)d\varsigma\ge1$. First, for every $D\ge0$ and every $n$ with $P_n\ne0$,
$$\sum_{j=-1}^1\|T_{n+D+j,n}(k)\|^2\ge\|w_D\|^2 :$$
for a unit vector $\psi\in P_n\mathcal H$ and its far translates $\psi_t=U(te_1)\psi$, the spectral measure of $H$ in $Y_\eta\psi_t$ tends from below on open sets to the convolution of those of $Y_\eta\Omega$ and of $e^{-2\eta H}\psi$, by clustering, which the unique vacuum of a component gives ($e^{-tH}\to|\Omega\rangle\langle\Omega|$); so the column $w_D$ of the vacuum reappears, up to the factors $e^{-\eta(\cdot)}$, in the blocks leaving every non-empty energy shell $P_n\mathcal H$: on the open set $(n+D-\frac12,n+D+2)$, with the two bounds above,
$$\sum_{j=-1}^1e^{-2\eta(2n+D+j)}\cosh^2\bigl(\tfrac\eta4(D+1)\bigr)\|T_{n+D+j,n}(k)\|^2\ge e^{-2\eta(D+1)-4\eta(n+1)}\|w_D\|^2,$$
and $\eta\downarrow0$ gives the inequality, the blocks and $w_D$ depending neither on $\eta$ nor on $\rho$. Secondly, $\|w_D\|$ is large for infinitely many $D$: by the two-point law as above, the heights $2\eta+\varsigma+\varsigma'$ of the two insertions being $\le\frac52\eta$, $\|Y_\eta\Omega\|^2\ge(1-o(1))\,c_{\rm sd}\frac\pi2(\frac25)^4\eta^{-4}\|k\|_{L^2}^2$ as $\eta\downarrow0$; while on $[D,D+1)$, $e^{-\eta E}\cosh(\varsigma E)\le e^{-(\eta-|\varsigma|)E}\le e^{-\frac34\eta D}$ for $|\varsigma|\le\frac\eta4$, so that $\|Y_\eta\Omega\|^2=\sum_D\|P_DY_\eta\Omega\|^2\le\sum_De^{-\frac32\eta D}\|w_D\|^2$; if $\|w_D\|^2\le(D+2)^{2q''}$ for all large $D$ with $q''<\frac32$, the right member would be $O(\eta^{-2q''-1})=o(\eta^{-4})$. Thirdly, every shell $P_n\mathcal H$ with $n\ge E_0$ is non-empty, because the support of the spectral measure of the translations holds a point $(E_0,q_0)$ with $E_0\ge1$ and is invariant under the boosts, along which the energy runs to infinity. Now let $q'<\frac32\epsilon$, choose $q''$ with $q'/\epsilon<q''<\frac32$, and for the infinitely many $D$ with $\|w_D\|^2>(D+2)^{2q''}$ put $n=\lceil(D+1)^{1/\epsilon}\rceil$: the three blocks $(n+D+j,n)$ lie in the band and $m+n+2\le6(D+1)^{1/\epsilon}$, so a bound $C(m+n+2)^{q'}$ in the band would give $(D+2)^{2q''}<3C^26^{2q'}(D+1)^{2q'/\epsilon}$, false for large $D$.

\emph{(e)} The multiple reflections for one cell (Proposition~\YMMref{104.2}{759}) and the one-cell bound (Proposition~\YMMref{89.1}{690}), with Montel's theorem and a diagonal choice over a countable set of combinations with dense classes, give $\mathbb D(z)$ on the whole disc with $\|\mathbb D(z)\|\le\bar{\mathsf Q}$, $\mathbb D(0)=e^{-\mathsf lH}$ and $\mathbb D'(0)=\mathsf I+c\,e^{-\mathsf lH}$, $\mathsf I$ the insertion of $f$, which is the slab operator of Theorem~\ref{C-363-3}(h1)(iv) by the insertion formula \YMMref{(2039)}{767}. $G(x)=\langle\Omega,\mathbb D(x)\Omega\rangle$ is holomorphic, $G(0)=1$, $|G|\le\bar{\mathsf Q}$; by the lemma of Schwarz $|G(x)-1|\le(1+\bar{\mathsf Q})|x|\|f\|_\sharp/\varrho\le\frac12$ on the smaller disc, so $\|\hat D(x)\|\le2\bar{\mathsf Q}$.

At real $x$ the lattice operators are $e^{-\mathcal E(x)}$ times positive operators, by reflection positivity at the middle plane of the cell, a site plane, $f$ being even under that reflection; the quotient $\hat D(x)$ removes the factor $e^{-\mathcal E(x)}$, so $\hat D(x)\ge0$. Finally $\hat D'(0)=\mathbb D'(0)-G'(0)e^{-\mathsf lH}=\mathsf I-\langle\Omega,\mathsf I\Omega\rangle e^{-\mathsf lH}$, and $\langle\Omega,\mathsf I\Omega\rangle=q_0(\Omega,\Omega)\int h'$ because $\langle\Omega,\mathsf B(s,s')\Omega\rangle=q_0(\Omega,\Omega)$ for all $s,s'$ and $\int\rho_1=1$. For a component the same holds with its lattice approximants, the periodic Gibbs states $\rho_{j,\xi}$ recalled in the results used.

\subsection{Proof of Theorem~\ref{C-363-6}}

\emph{(a)} Let $e^\alpha_x=\kappa_\alpha e^{-x_3\omega-i\vec x\cdot\vec p}$ for $x\in\mathbb E_+$, and for pairwise different points of $\mathbb E_+$
$$J_n(x)=\prod_i\mathsf P(\partial_{t_i})\Bigl[\exp\Bigl(\sum_{i<l}t_iG(x_i-x_l)t_l\Bigr){\rm Exp}\Bigl(\sum_it_i\cdot e_{x_i}\Bigr)\Bigr]_{t=0},\qquad J(G)=\int\prod_kh_k(x_k)J_n(x)\,dx ,$$
${\rm Exp}(u)=\sum_mu^{\otimes m}/\sqrt{m!}$. The covariance of $G$ and the invariance of $\mathsf P$ under $\mathrm O(3)$ are, as in the statement, those of the orthogonal representation $\rho_{\mathsf A}$ of $\mathrm O(3)$ on $\mathbb R^{\mathsf A}$: $G(\mathsf Rx)=\rho_{\mathsf A}(\mathsf R)G(x)\rho_{\mathsf A}(\mathsf R)^{\rm T}$, $\mathsf P(\rho_{\mathsf A}(\mathsf R)t)=\mathsf P(t)$, with $\rho_{\mathsf A}(\theta)={\rm diag}(\varepsilon_\alpha)$; and the tensor $\mathsf p$ is symmetric. So the invariance of $\mathsf P$ under $\theta$ is $\mathsf P(\varepsilon t)=\mathsf P(t)$, $(\varepsilon t)_\alpha=\varepsilon_\alpha t_\alpha$, that is $\mathsf p_{\underline\alpha}\prod_i\varepsilon_{\alpha_i}=\mathsf p_{\underline\alpha}$, and the covariance under $\theta$ is $\varepsilon G(x)\varepsilon=G(\theta x)$, $\varepsilon={\rm diag}(\varepsilon_\alpha)$.

As $\langle{\rm Exp}\,u,{\rm Exp}\,v\rangle=e^{\langle u,v\rangle}$ and $\langle e^\alpha_x,e^\beta_y\rangle=\varepsilon_\alpha G_{\alpha\beta}(\theta x-y)$, the substitution $s_i=\varepsilon\tilde s_i$, with $\mathsf P(\varepsilon t)=\mathsf P(t)$ in the derivatives and $\varepsilon G(x_i-x_l)\varepsilon=G(\theta x_i-\theta x_l)$ in the exponent, gives $\langle J_m(x),J_n(y)\rangle=S_{m+n}(\theta x,y)$, the Schwinger function of ${:}\mathsf P(\Phi){:}$; so $V[G]=J(G)$ is an isometry, (Rec1) follows from the covariance of $G$ and the invariance of $\mathsf P$, $e^{-tH_{\rm F}}J_n(x)=J_n(x+te_3)$ gives (Rec2), $|\sum\vec p(m_i)|\le\sum\omega(m_i)$ gives (Rec3), and $\ker d\Gamma(\omega)=\mathbb C\Omega_{\rm F}$ gives (Rec5).

The point field is the Fock operator
$$\mathsf B_{\rm F}(s,s')=\sum_j\binom kj\sum_{\underline\alpha}\mathsf p_{\underline\alpha}\prod_{i>j}\varepsilon_{\alpha_i}\,a^*(e^{\alpha_1}_{se_3})\cdots a^*(e^{\alpha_j}_{se_3})\,\Gamma(e^{-(s+s')\omega})\,a(e^{\alpha_{j+1}}_{s'e_3})\cdots a(e^{\alpha_k}_{s'e_3}),$$
with $\mathsf B_{\rm F}(s,s')J_n(y)=J_{n+1}(se_3,y+(s+s')e_3)$; from 
\[
\|a(u_1)\cdots a(u_m)e^{-cH_{\rm F}}\|\le\prod_i\|\omega^{-1/2}u_i\|(m/2ec)^{m/2}
\]
 and $\|\omega^{-1/2}e^\alpha_{ue_3}\|\le c_\kappa u^{-\delta}$ one gets $\|\mathsf B_{\rm F}(s,s')\|\le C_{\mathsf B}\min\{s,s',1\}^{-\wp}$, $C_{\mathsf B}=2^k|\mathsf p|_1(2^\delta c_\kappa)^k(k/e)^{k/2}$, hence (Rec4) with $C_\wp=2^\wp C_{\mathsf B}$, $c_0=\frac12$. The uniqueness of the family of point operators, Proposition~\YMMref{107.2}{767}, identifies $V^{-1}\mathsf B_{\rm F}V$ with $\mathsf B$, and so the field with the form of the Wick polynomial, which on vectors with finitely many particles and bounded energy is
$$\Phi_f(\chi,\chi')=\sum_j\binom kj\sum_{\underline\alpha}\mathsf p_{\underline\alpha}\prod_{i>j}\varepsilon_{\alpha_i}\int f(x)\Bigl\langle\prod_{i\le j}a(\kappa^x_{\alpha_i})\chi,\prod_{i>j}a(\kappa^x_{\alpha_i})\chi'\Bigr\rangle dx,\qquad\kappa^x_\alpha=e^{i(t\omega-\vec x\cdot\vec p)}\kappa_\alpha ;$$
the projection on $V\mathcal H$ commutes with $H_{\rm F}$, the translations and $\mathsf B_{\rm F}$, which gives the identity for every $\chi\in\mathcal D_\infty(d\Gamma(\omega))$.

\emph{(b)} The component with $kn$ particles of $V\phi(u)^n\Omega$ is $((kn)!)^{1/2}{\rm Sym}(w_u^{\otimes n})$, and no component has more particles: in $\Phi_u(\chi,\chi')$ the $j$-th term changes the number of particles by $2j-k$, the top term $j=k$ is the creation of $w_u$, by Fubini and $\int u(x)\prod_ie^{i(t\omega(m_i)-\vec x\cdot\vec p(m_i))}dx=\hat u(\sum\omega,\sum\vec p)$, and $w_u$ is a symmetric function, the tensor $\mathsf p$ being symmetric, so that ${\rm Sym}(w_u)=w_u$. The symmetric tensors $v^{\otimes k}$ are total among the symmetric ones, so some unit $v$ has $c_u=|\langle v^{\otimes k},w_u\rangle|>0$, and $\|\phi(u)^n\Omega\|\ge|\langle v^{\otimes kn},{\rm Sym}(w_u^{\otimes n})\rangle|((kn)!)^{1/2}=c_u^n((kn)!)^{1/2}$; $m!\ge(m/e)^m$ makes $\sum_n\|\phi(u)^n\Omega\|^{-1/n}$ converge for $k\ge3$.

Let a real $u_0\ne0$ in $C^\infty_c$ have $b_{mn}(u_0)\le\beta(m+n+2)^q$ on a band, $q<\frac k2$; outside the band Theorem~\ref{C-363-3}(b3) gives a linear bound, so $b_{mn}(u_0)\le\beta'(m+n+2)^{\max\{q,1\}}$ for all $m,n$; a linear form bound at $u_0$ gives $b_{mn}(u_0)\le c\sqrt{(m+2)(n+2)}$ as well. Choose $\mathsf s>1$ with $\max\{q,1\}\mathsf s<\frac k2$, and $u=h*u_0$ with $h=g\otimes g_2$, $g$ a bump of the Gevrey class $\mathsf s$ (as in the proof of Theorem~\ref{C-363-5}(a)) and $g_2\ge0$ in $C^\infty_c(\mathbb R^2)$: $u\in C^\infty_c$ is real and $\ne0$, $\hat u=\hat h\hat u_0$ being a product of entire functions $\ne0$. Theorem~\ref{C-363-3}(b2), (b4) and the choice of the number of derivatives as in the proof of Theorem~\ref{C-363-5}(a) give $b_{mn}(u)\le C e^{-\mathsf a|m-n|^{1/\mathsf s}}b_{mn}(u_0)$, and the sum over the paths from $\Omega$ gives $\|\phi(u)^n\Omega\|\le C^n(n!)^{\max\{q,1\}\mathsf s}$, against $c_u^n(n!)^{k/2}$.

The same with a triple of order $\wp<\frac k2$ in the place of the band bound.

\emph{(c)} For the free field, $\|\omega^{-1/2}\kappa e^{-s\omega}\|^2=\frac1{4\pi}\int_{2s\mathsf m}^\infty e^{-v}\frac{dv}v\le\frac1{8\pi\delta}(2s\mathsf m)^{-2\delta}$ for $0<\delta\le\frac12$. For $k\ge2$, the map $(p_1,\dots,p_k)\mapsto(\sum\omega(p_i),\sum p_i)$ has an image that holds $\mathcal V_k=\{E>(|\vec k|^2+k^2\mathsf m^2)^{1/2}\}$, and $w_u$ is continuous with $\kappa>0$; so $w_u\ne0$ as soon as $\hat u\ne0$ at a point of $\mathcal V_k$. For $u\in C^\infty_c\setminus\{0\}$, $\hat u$ is entire. For $u\in\mathcal S$ vanishing off $\bar W$ and every $k_2$, $(E,k_1)\mapsto\hat u(E,k_1,k_2)$ continues boundedly into the directions $\eta$ with $\eta_1<-|\eta_0|$, since $x^0\eta_0-x^1\eta_1\ge0$ on $\bar W$; if it vanished on the open set where the plane meets $\mathcal V_k$, it would vanish on the plane.

For $u=\mathsf g\otimes h'$ take $\vec k_*$ with $\hat h'(\vec k_*)\ne0$ and $E$ large. At $k=6$, $\delta=\frac16$: $\wp=4$ and $C_\wp=2^4\cdot2^6\cdot2c_\kappa^6\cdot(6/e)^3$ with $c_\kappa^6=(\frac{4\pi}3)^{-3}(2\mathsf m)^{-1}$, which is $93312(\pi e)^{-3}\mathsf m^{-1}$; the field changes the number of particles by an even number, so the space lies in the even Fock space, where $d\Gamma(\omega)\ge2\mathsf m$ off $\Omega_{\rm F}$; $S_2=6!\,G^6$, $G=e^{-\mathsf mr}/4\pi r$. \emph{Locality.} Let $O_1,O_2$ be spacelike separated, $P$ the projection on $V\mathcal H$ and $\mathcal R_0=\mathcal R_0(O_2)$. By Lemma~\ref{C-363-22}(a), (c), for real $h\in\mathcal F(O_2)$ the compression $V^{-1}P\mathsf W[u_h]V$ commutes weakly with the fields of $O_1$ on $\mathcal D_0$, so $P\mathcal R_0P|_{V\mathcal H}\subset V\mathcal C(O_1)V^{-1}$; by Lemma~\ref{C-363-22}(d) $\Omega_{\rm F}$ is cyclic for $\mathcal R_0$, and Lemma~\ref{C-363-21}(a) gives $(P\mathcal R_0P|)'=P\mathcal R_0'P|$; by Lemma~\ref{C-363-22}(e), $V^{-1}P\mathcal R_0'PV\subset\mathcal C(O_2)$.

So
$$\mathcal C(O_1)'\subset V^{-1}(P\mathcal R_0P|)'V=V^{-1}P\mathcal R_0'PV\subset\mathcal C(O_2),$$
which is (W), and (CP) follows by Theorem~\ref{C-363-3}(d). For the derivatives $\partial_\alpha\varphi^a$ of massless fields and for the generalised free fields the same proof runs with Lemma~\ref{C-363-22}(f), whose hypotheses \textup{(f1)}, \textup{(f2)} hold for them, in the place of Lemma~\ref{C-363-22}(a)--(e): $\mathcal R_0(O)$ is the span of the $\mathsf W[u]$, $u\in\mathfrak k(O)$, the real span of the $u^\alpha_h=\varsigma_\alpha\hat h(\omega,\vec p)\kappa_\alpha$, $h\in\mathcal F_{\mathbb R}(O)$, with the weight $\omega^{1/2}(1+\omega)^{r'}$, $r'>\delta$, in the place of $\omega$ in the bounds of Lemma~\ref{C-363-22}, the reality of $\langle\kappa^x,u_h\rangle$ replaced by the reality or pure imaginarity of $\langle\kappa^x_\alpha,u\rangle$, $u\in\mathfrak k(O)$, according to $\varepsilon_\alpha$ (the derivatives in space and in time of a kernel that is real at spacelike separation), and the density of the one-particle vectors of a region by the same reflection argument with $\omega\kappa_0$ in the place of $\kappa$.

For ${:}(\partial\varphi\cdot\partial\varphi)^2{:}$: $k=4$, $\delta=1$, $\wp=6$; for the kernel $r^{-2}$, the mixture of masses $4\pi\,d\mathsf m$ has $\|\omega^{-1/2}\kappa e^{-s\omega}\|^2=\frac1{2s}$, so $\delta=\frac12$ and $\wp=3$ at $k=3$, $S_2=3!\,r^{-6}$, and the spectrum reaches $0$. The last two assertions of (c) follow from (b), applied to ${:}\chi^6{:}$ and to ${:}\chi^3{:}$.

\emph{(d)} The photons are the Gaussian family of (a) with $\mathsf M=\mathbb R^2\setminus\{0\}$, $d\mu=d^2p$, $\omega=|p|$, $\mathsf P=\sum_\alpha t_\alpha^2$, $\kappa_0=(2\pi)^{-1}(2|p|)^{-1/2}$, $\kappa_j=-ip_j\kappa_0$, $\kappa_3=-\omega\kappa_0$, $\varepsilon=(+,+,-)$, $G_{\alpha\beta}=-\partial_\alpha\partial_\beta(4\pi|x|)^{-1}$; the form of (a) is ${:}|\nabla\varphi|^2-\dot\varphi^2{:}=-\Theta$. With $\|\omega^{-1/2}e^3_{ue_3}\|^2=(16\pi u^2)^{-1}$ and $\|\omega^{-1/2}e^j_{ue_3}\|^2=(32\pi u^2)^{-1}$ the three terms of $\mathsf B_{\rm F}$ are bounded by $(\pi e)^{-1}\min\{s,s'\}^{-3}$ each, so $C_\wp=2^3\cdot3/\pi e=24/\pi e$; $S_2=2\sum_{\alpha\beta}G_{\alpha\beta}^2=12/16\pi^2r^6$ and $\|\mathsf B(s,s)\Omega\|^2=S_2(2s)$.

The span of the $\Theta$-vectors is the even Fock space, as in (e) at $\mathsf n=1$. \emph{The linear bound.} Let $\tilde f(k)=\int f(x)e^{i(k^0t-\vec k\cdot\vec x)}dx$ (the transform $\hat f$ of the setting) and $\kappa_f(\vec k)=\frac1{2\pi^2}\sup_{k^0}|\tilde f(k^0,\vec k)|\le8|f|_8(1+|\vec k|^2)^{-2}$, $\theta(f)=\int\kappa_f\,d^2k\le8\pi|f|_8$. The field $\phi(f)=-\Theta(f)$, $f$ real, has a part with one creator and one annihilator, with the kernel $t_f(p,q)=2\kappa_0(p)\kappa_0(q)(\vec p\cdot\vec q-\omega_p\omega_q)\tilde f(\omega_p-\omega_q,\vec p-\vec q)$, of modulus $\le\sqrt{\omega_p\omega_q}\,\kappa_f(\vec p-\vec q)$, bounded by $\theta(f)H$ by Young's inequality; and the two pair parts $\mathsf Z^+(w_f)+\mathsf Z^+(w_f)^\dagger$, $\mathsf Z^+(c)=\iint c(p,q)a^*_pa^*_q\,dp\,dq$, with the kernel $w_f(p,q)=\sum_\alpha\kappa_\alpha(p)\kappa_\alpha(q)\tilde f(\omega_p+\omega_q,\vec p+\vec q)$, of modulus $\le\frac12\sqrt{\omega_p\omega_q}\,\kappa_f(\vec p+\vec q)$, so that $w_f/\sqrt{\omega_p\omega_q}$ has the operator norm $\le\frac12\theta(f)$.

For them a square is completed: with $\epsilon>0$ and $g_p=\pm w_f(p,\cdot)/\epsilon\omega_p$, $0\le\int\omega_p\|(a_p+a^*(g_p))\psi\|^2dp$ is $\langle\psi,H\psi\rangle$, plus the pair form at $\psi$ divided by $\pm\epsilon$, plus $\int\omega_p\|a(g_p)\psi\|^2dp\le\frac{\theta(f)^2}{4\epsilon^2}\langle\psi,H\psi\rangle$, plus $\int\omega_p\|g_p\|^2dp\,\|\psi\|^2\le\frac{\mathsf n(f)}{4\epsilon^2}\|\psi\|^2$, $\mathsf n(f)=\frac1{4\pi^4}\iint\omega_q|\tilde f(\omega_p+\omega_q,\vec p+\vec q)|^2\le32\pi^3|f|_8^2$; so $\pm(\mathsf Z^+(w_f)+\mathsf Z^+(w_f)^\dagger)\le(\epsilon+\frac{\theta(f)^2}{4\epsilon})H+\frac{\mathsf n(f)}{4\epsilon}$, and at $\epsilon=\frac\vartheta2$, $\vartheta\ge\theta(f)$: $\le\vartheta H+\mathsf n(f)/2\vartheta$.

At $\vartheta=8\pi|f|_8\ge\theta(f)$: $\pm\Theta(f)\le16\pi|f|_8H+2\pi^2|f|_8\le16\pi|f|_8N$. The same bound for $\Theta(\partial_tf)=-i[H,\Theta(f)]$ gives the operator bound and essential self-adjointness on every core of $H$; $\|\Theta(f)\psi\|\le c(f)(n+2)\|\psi\|$ on vectors with at most $n$ particles gives the growth; (W) and (CP) by (c). \emph{Order exactly one} on the even space: the squeezed vectors of (e) at $\mathsf n=1$. \emph{At one time.} For $m\le n$ write $\Theta_k=\Theta^{(0)}+\Theta^{(-)}+\Theta^{(+)}$ (number-preserving, annihilating and creating parts), with $\tilde k(\xi)=\int k(\vec x)e^{-i\xi\cdot\vec x}d^2x$, $\Theta^{(0)}=\iint(4\pi^2)^{-1}\sqrt{\omega_p\omega_q}(1-\hat p\cdot\hat q)\tilde k(p-q)\,a^*_pa_q$, $\Theta^{(-)}=\iint\Xi(p,q)\,a_pa_q$, $\Xi(p,q)=-(8\pi^2)^{-1}\sqrt{\omega_p\omega_q}(1-\hat p\cdot\hat q)\tilde k(-p-q)$, and $\Theta^{(+)}=(\Theta^{(-)})^\dagger$.

The annihilation and creation densities select energies: $\langle\chi,a_pa_q\psi\rangle=0$ unless $\omega_p+\omega_q\in(n-m-1,n-m+1)$ for $\chi\in P_m\mathcal F$, $\psi\in P_n\mathcal F$. The integral operators with the kernels $\tilde k(q\mp p)$ have norms $\le(2\pi)^2\|k\|_\infty$ (Plancherel), and with the factor $1-\hat p\cdot\hat q$ at most twice that; multiplication of a kernel by a window $\eta(\omega_p+\omega_q-\omega)$, $\eta$ the inverse Fourier transform of an integrable $\hat\eta$, multiplies its norm by at most $\frac1{2\pi}\|\hat\eta\|_{L^1}$, which is at most $\sqrt3$ (Schwarz) for $\eta_2=\mathbf 1_{[-\frac32,\frac32]}*\mathbf 1_{[-\frac12,\frac12]}$, a function equal to $1$ on $[-1,1]$ and carried by $[-2,2]$.

With $\omega=n-m$ and $\Xi_\eta(p,q)=\Xi(p,q)\eta_2(\omega_p+\omega_q-\omega)$, $P_m\Theta^{(-)}P_n=P_m\iint\Xi_\eta a_pa_qP_n$. So $\|P_m\Theta^{(0)}P_n\|\le2\|k\|_\infty\sqrt{(m+1)(n+1)}$ and, with $g_q=\overline{\Xi_\eta(\cdot,q)}\omega_q^{-1/2}$, 
\[
\mathsf h(\omega)^2=\iint|\Xi_\eta(p,q)|^2\omega_q^{-1}dp\,dq=\int\|g_q\|^2dq
\]
 and $\|a^*(g)\chi\|^2=\|a(g)\chi\|^2+\|g\|^2\|\chi\|^2$,
$$\|P_m\Theta^{(-)}P_n\|\le\sqrt{n+1}\,\bigl[\sqrt3\|k\|_\infty\sqrt{m+1}+\mathsf h(\omega)\bigr],\qquad\mathsf h(\omega)^2\le\frac{\omega+2}{(2\pi)^4}\int|\tilde k(\xi)|^2|\mathcal A_\xi|\,d\xi\le(\omega+2)^2\mathsf n_k^2,$$
$\mathcal A_\xi=\{p:|p|+|\xi-p|\in[\omega-2,\omega+2]\}$, of area $\le10.4(\omega+2)$ for $|\xi|\le\frac12(\omega+2)$ (confocal ellipses) and $\le\frac\pi2(\omega+2)|\xi|$ otherwise; $\Theta^{(+)}$ vanishes between these shells for $m<n$. This is the upper bound. For the lower bound let $\Xi_1=\Xi\,\eta_1(\omega_p+\omega_q-\omega)$, $\eta_1=(\frac12-|\cdot|)_+$, $L_1=\iint\Xi_1a_pa_q$; then $[L_1,L_1^*]\ge2\|\Xi_1\|_2^2$, and the vectors $\chi_j=(L_1^*)^j\Omega$ have $\|\chi_{j+1}\|^2\ge2(j+1)\|\Xi_1\|_2^2\|\chi_j\|^2$; as $L_1^*$ raises the energy by $\omega\pm\frac12$, one shell $P_m$ in the range $[j(\omega-\frac12),j(\omega+\frac12)]$ has $\|P_nL_1^*P_m\|\ge(\frac{j+1}3)^{1/2}\|\Xi_1\|_2$ for one of three $n$; the window costs the factor $\frac12$, and $\|\Xi_1\|_2^2\ge2.3\cdot10^{-5}\|k\|_2^2\omega^3$ for $\omega\ge\max\{3,7\xi_0(k)+1\}$, because on the ring $\mathcal A'_\xi=\{|p|+|\xi-p|\in[\omega-\frac14,\omega+\frac14]\}$, $|\xi|\le\xi_0$, where $|p|\ge3\xi_0$, the kernel has $|\Xi(p,\xi-p)|^2\ge\frac43|p|^2|\tilde k(\xi)|^2/64\pi^4$; the number-preserving part is at most $\kappa_k(\omega-2)G$ there, as only momentum transfers $|p-q|>\omega-2$ occur. The band statement follows by taking $j\approx(\omega+1)^{1/\epsilon}/\omega$.

\emph{(e)} \emph{The cyclic space.} The $\Theta$-vectors are $\mathrm O(\mathsf n)$-invariant. For the converse use the explicit bilocal calculation in steps (1)--(4) of (f) below, first as identities of Wick polynomials, before taking commutants. The timelike matrix $\mathsf H_{\mu\nu}=\partial_\mu\partial_\nu\Delta_{\rm c}$ is invertible. Point splitting gives its traceless symmetric contractions, and adjoining the already available trace $\Theta$ gives every local symmetric component. The two contractions $5\mathsf X+\mathsf H\mathsf t$ and $8\mathsf t$ in (3) then recover every bilocal component at timelike separated supports. The commutator identity and overlapping-cone chain in (4) give the colour sums $\sum_a\varphi^a(h)\varphi^a(h')$ for arbitrary real integral-zero tests in a connected double cone, including spacelike separated supports. These are limits of polynomials in the $\Theta(f)$ on the common finite-particle energy-smooth core, with the finite-particle bounds used in those limits. This argument applies to the polynomial vectors themselves and does not use the cyclic-space identification it proves. In particular, for real $g$ of integral zero in a double cone $O$, the sums of squares $\sum_a\varphi^a(g)^2$ are such limits.

One more commutator gives the squares: let $O_3\subset O$ be a smaller double cone that holds ${\rm supp}\,g$, so that $g=\partial_\mu\mathsf G^\mu$ with $\mathsf G\in C^\infty_c(O_3)^3$, take a ball $O_1\subset O$ spacelike to $O_3$, a ball $O_2\subset O$ in the common future of $O_1$ and $O_3$, and $g_1,g_3\in\mathcal F_{\mathbb R}(O_1)$ of integral $0$ with $\Delta_{\rm c}(g_1,g_3)\ne0$; then $\sum_a\varphi^a(g_i)\varphi^a(g)$, $i=1,3$, are such limits up to numbers, and, as $\Delta_{\rm c}(g,g_3)=\Delta_{\rm c}(g_1,g)=\Delta_{\rm c}(g,g)=0$,
$$\Bigl[\sum_a\varphi^a(g_1)\varphi^a(g),\sum_b\varphi^b(g_3)\varphi^b(g)\Bigr]=i\Delta_{\rm c}(g_1,g_3)\sum_a\varphi^a(g)^2 ;$$
by the first fundamental theorem for $\mathrm O(\mathsf n)$ the invariant polynomials in the $\varphi^a(g)$ are generated by the scalar products $\sum_a\varphi^a(g)\varphi^a(g')$, which polarisation reduces to squares; and the invariant polynomial vectors are dense in the invariant vectors. The Euclidean vectors give the same space by (a). With $\mathrm{SO}(3)$: $\sum\epsilon_{abc}a_a^*(u)a_b^*(v)a_c^*(w)\Omega\ne0$ is $\mathrm{SO}(3)$-invariant, has three particles and is orthogonal to the even space. $\|\mathsf B(s,s)\Omega\|^2=\mathsf nS_2(2s)$ with the $S_2$ of one colour. \emph{Order one.} For a bounded unit $u\in\mathfrak h$ of compact support put $A^*=\sum_aa^*_a(u)^2$; then $A(A^*)^k\Omega=2k(2k-2+\mathsf n)(A^*)^{k-1}\Omega$, and the singlet squeezed vectors $\chi_\lambda=(1-\lambda)^{\mathsf n/4}\exp(\frac{\sqrt\lambda e^{i\beta}}2A^*)\Omega$, $0\le\lambda<1$, are invariant unit vectors of $\mathcal D_\infty(H)$.

With the one-particle operator $T_f$ of the kernel $\frac{\omega_p\omega_q-\vec p\cdot\vec q}{4\pi^2\sqrt{\omega_p\omega_q}}\tilde f(\omega_p-\omega_q,\vec p-\vec q)$ (the kernel $-t_f$ of (d)) and $c_f(p,q)=\frac{\omega_p\omega_q-\vec p\cdot\vec q}{8\pi^2\sqrt{\omega_p\omega_q}}\tilde f(\omega_p+\omega_q,\vec p+\vec q)$ (the $w_f$ of (d)), the field of $\mathsf n$ colours is $\Theta(f)=d\Gamma(T_f\otimes1)-\mathsf Z^+(c_f)-\mathsf Z^+(c_{\bar f})^\dagger$, $\mathsf Z^+(c)=\sum_a\iint c(p,q)a^*_a(p)a^*_a(q)\,dp\,dq$, and $\phi(f)=-\Theta(f)$; so
$$\langle\chi_\lambda,\phi(f)\chi_\lambda\rangle=\frac{\mathsf n}{1-\lambda}\bigl[2\sqrt\lambda\,{\rm Re}(e^{i\beta}\mathsf w)-\lambda\mathsf t\bigr],\qquad\langle\chi_\lambda,H\chi_\lambda\rangle=\frac{\mathsf n\lambda\,\mathsf e}{1-\lambda},$$
$\mathsf t=\langle u,T_fu\rangle$, $\mathsf w=\langle c_f,u\otimes u\rangle$, $\mathsf e=\langle u,\omega u\rangle$. Some $u$ has $|\mathsf t|+2|\mathsf w|>0$: otherwise $\tilde f$ would vanish on $\{|k^0|\le|\vec k|\}$ and on $\pm\bar V_+$. With Jensen's inequality $\langle\psi,N^\alpha\psi\rangle\le\langle\psi,N\psi\rangle^\alpha$ the quotient grows like $(1-\lambda)^{\alpha-1}$. Either sign: the one-particle operator $T_f$ has a continuous kernel that vanishes on the diagonal, so it is $0$ or indefinite; choose $u$ with $\mathsf t$ of the sign needed, or with $\mathsf w\ne0$ if $T_f=0$, and $\beta$. These vectors are limits in every $\|N^l\cdot\|$ of vectors of $\mathcal D_0$. The upper bound adds the bounds of the $\mathsf n$ colours. If $b_{mn}(f)\le\beta_p(1+\max\{m,n\})^p$ with $p<1$, an interpolation of this bound with the decay of Theorem~\ref{C-363-3}(b2) and the crude bound $b_{mn}(f)\le c(2+\min\{m,n\})^{\wp}$, from Theorem~\ref{C-363-3}(b5) at $j=0$ and $b_{mn}(f)=b_{nm}(\bar f)$, gives $b_{mn}(f)\le C(2+m)^{p'/2}(2+n)^{p'/2}(1+|m-n|)^{-2}$ with $p<p'<1$, hence a form bound with $N^{p'}$, which the squeezed vectors exclude.

\emph{(f)} For a field $\Phi$ among $\Theta$, $S={:}\varphi\cdot\varphi{:}$, the multiplet $\Theta_{\mu\nu}={:}\partial_\mu\varphi\cdot\partial_\nu\varphi{:}$ and the $T_{\mu\nu}=\Theta_{\mu\nu}-\frac12\eta_{\mu\nu}\Theta$, all on the Fock space, $\mathcal D^\Phi_0$ is the span of $\Omega$ and of the $\Phi(f_1)\cdots\Phi(f_n)\Omega$, $\mathcal H_\Phi$ its closure, and $\mathcal C_\Phi(O)$, $\mathcal A_\Phi(O)=\mathcal C_\Phi(O)'$ and $\mathcal A^d_\Phi(O)=\bigcap\{\mathcal A_\Phi(\tilde W):\tilde W\supset O\}$ are formed as in the setting, with $\Phi$ and $\mathcal D^\Phi_0$ in the place of $\phi$ and $\mathcal D_0$. $\Theta$ at the mass $0$ (by (d), (e)), $\Theta$ at a mass $\mathsf m>0$ (by (a), with $\delta=1$, as $\|\omega^{-1/2}\kappa_\alpha e^{-s\omega}\|^2\le(16\pi s^2)^{-1}$) and $S$ at $\mathsf m>0$ (by (a), as in (c)) are data (Rec), and $V$ carries their nets onto these; so Theorem~\ref{C-363-3}(e) applies to them. For $S$ at the mass $0$, $\mathcal A^d_S$ is the net of this definition through wedges. $\mathfrak R(\mathsf K)$ is the von Neumann algebra of the Weyl operators of the vectors of a closed real subspace $\mathsf K$ of the one-particle space, and $\mathsf K'$ the symplectic complement.

Let first $\mathsf n=1$, $\mathsf m=0$. With $\mathsf H_{\kappa\lambda}=\partial_\kappa\partial_\lambda\Delta_{\rm c}$ and $\mathsf X_{\kappa\lambda}(x,y)={:}\partial_\kappa\varphi(x)\partial_\lambda\varphi(y){:}$, on vectors with finitely many particles and for timelike separated supports,
$$[\Theta(f_1),\Theta(f_2)]=-4i\iint f_1(x)f_2(y)\,\mathsf H^{\kappa\lambda}(x-y)\,\mathsf X_{\kappa\lambda}(x,y)\,dx\,dy+c\,1 ,$$
and $\Delta_{\rm c}(z)=\mp(2\pi)^{-1}(z\cdot z)^{-1/2}\ne0$ on $\pm V_+$, inside the cone in $1+2$ dimensions; on $V_+$, $-2\pi(z\cdot z)^{3/2}\mathsf H_{\kappa\lambda}(z)=3z_\kappa z_\lambda/(z\cdot z)-\eta_{\kappa\lambda}$, whose mixed form has the eigenvalues $2,-1,-1$, so that $\mathsf H$ is invertible at timelike points. For $\Theta$ the following identities are used on the common energy-smooth core. First time-smooth the bounded weakly commuting operator $Q$ as in the proof of Theorem~\ref{C-363-3}(a), working in a region with closure in $O$; the linear energy bounds and essential self-adjointness extend the polynomial-domain identities to this core. Products and commutators are taken there, and the displayed Wick and point-splitting limits pass with their finite-particle bounds. At the end let the smoothing tend to the identity. This is a statement about these cores and limits, not an unrestricted rule for weak commutants of unbounded operators. Let $Q\in\mathcal C_\Theta(O)$, $O$ open and connected. (1) The display gives the contractions $\mathsf X[\Psi\mathsf H]=\iint\Psi\,\mathsf H^{\kappa\lambda}\mathsf X_{\kappa\lambda}$, $\Psi$ scalar in $C^\infty_c$ of the timelike pairs of $O\times O$: one component of nine. (2) For $f\in\mathcal F(O)$, $\chi\in C^\infty_c(V_+)$ and small $\epsilon$, $\Psi(x,y)=-2\pi f(y)\epsilon^{-3}\chi(\frac{x-y}\epsilon)((x-y)\cdot(x-y))^{3/2}$ makes $\Psi\mathsf H$, as $3z_\kappa z_\lambda/(z\cdot z)-\eta_{\kappa\lambda}$ is homogeneous of degree $0$, the function $f(y)\epsilon^{-3}\zeta(\frac{x-y}\epsilon)$ with $\zeta(z)=\chi(z)(3z_\kappa z_\lambda/(z\cdot z)-\eta_{\kappa\lambda})$, and point splitting gives $3\mathsf N^{\kappa\nu}[\chi]\Theta_{\kappa\nu}(f)-(\int\chi)\Theta(f)$, $\mathsf N^{\kappa\nu}[\chi]=\int\chi\,z^\kappa z^\nu/(z\cdot z)$; these tensors span the traceless symmetric ones, and adjoining the trace $\Theta$, with which $Q$ already commutes, gives all symmetric components: $Q$ commutes weakly with all $\Theta_{\kappa\nu}(f)$, $f\in\mathcal F(O)$. (3) $[\Theta_{\kappa\lambda}(f_1),\Theta_{\mu\nu}(f_2)]=-i\iint f_1f_2\,\mathsf Y_{\kappa\lambda,\mu\nu}+c$, $\mathsf Y_{\kappa\lambda,\mu\nu}=\mathsf H_{\lambda\mu}\mathsf X_{\kappa\nu}+\mathsf H_{\lambda\nu}\mathsf X_{\kappa\mu}+\mathsf H_{\kappa\mu}\mathsf X_{\lambda\nu}+\mathsf H_{\kappa\nu}\mathsf X_{\lambda\mu}$, and with $\mathsf t=(\mathsf H^{-1})^{\lambda\mu}\mathsf X_{\lambda\mu}$: $(\mathsf H^{-1})^{\lambda\mu}\mathsf Y_{\kappa\lambda,\mu\nu}=5\mathsf X_{\kappa\nu}+\mathsf H_{\kappa\nu}\mathsf t$ and $(\mathsf H^{-1})^{\kappa\nu}(\mathsf H^{-1})^{\lambda\mu}\mathsf Y_{\kappa\lambda,\mu\nu}=8\,\mathsf t$; so $\iint\Psi^{\kappa\nu}\mathsf X_{\kappa\nu}=\iint\Upsilon^{\kappa\lambda\mu\nu}\mathsf Y_{\kappa\lambda,\mu\nu}$ with $\Upsilon^{\kappa\lambda\mu\nu}=\frac15[\Psi^{\kappa\nu}(\mathsf H^{-1})^{\lambda\mu}-\frac18(\Psi^{\rho\sigma}\mathsf H_{\rho\sigma})(\mathsf H^{-1})^{\kappa\nu}(\mathsf H^{-1})^{\lambda\mu}]$: here the invertibility of $\mathsf H$ solves for every bilocal field $\mathsf X[\Psi]$, $\Psi$ a matrix in $C^\infty_c$ of the timelike pairs. (4) For double cones $K,K'$ in $O$ write $K\approx K'$ if $Q$ commutes weakly with all $\varphi(h)\varphi(h')$, $h\in\mathcal F_{\mathbb R}(K)$, $h'\in\mathcal F_{\mathbb R}(K')$ of integral $0$; such $h$ are divergences $\partial_\kappa\mathsf k^\kappa$, so $\varphi(h)\varphi(h')$ is a bilocal field up to a number, and (3) gives $K\approx K'$ when all pairs of $K\times K'$ are timelike.

In
\[
\begin{gathered}
[\varphi(k)\varphi(k'),\varphi(l)\varphi(l')]\\
=i\Delta_{\rm c}(k',l)\varphi(k)\varphi(l')\\
+i\Delta_{\rm c}(k',l')\varphi(k)\varphi(l)\\
+i\Delta_{\rm c}(k,l)\varphi(l')\varphi(k')\\
+i\Delta_{\rm c}(k,l')\varphi(l)\varphi(k'),
\end{gathered}
\]
with $k,k',l,l'$ of integral $0$: $k=l$ in $K$ and $k',l'$ in $K'$ give $K\approx K$ when $K\approx K'$ (the left member and the first and fourth terms are known, the third has the factor $\Delta_{\rm c}(k,k)=0$, and $\Delta_{\rm c}$ does not vanish identically on the pairs of functions of integral $0$ in $K'$); $k$ in $K$, $k',l$ in $K''$ and $l'$ in $K'$ give $K\approx K'$ from $K\approx K''$, $K''\approx K''$, $K''\approx K'$ (the left member and the second, third and fourth terms are known). Every point of $O$ has a small double cone $K$ with $K\approx K$ (with a translate of $K$ along the time axis), the classes of the relation so induced on $O$ are open, and $O$ is connected: there is one class.

With a Lebesgue number of a cover of a compact connected set that holds the supports, every $h\in\mathcal F_{\mathbb R}(O)$ of integral $0$ is cut into pieces of integral $0$ in small double cones along a chain, $h=\sum_i(\chi_ih-\vartheta_ib_i)+\sum_i\vartheta_i(b_i-b_1)$, $\vartheta_i=\int\chi_ih$, $\int b_i=1$, $\chi_i$ a partition of unity and $b_i-b_1$ cut along the chain: $Q$ commutes weakly with all $\varphi(h)\varphi(h')$, $h,h'\in\mathcal F_{\mathbb R}(O)$ of integral $0$. (5) The cosine series converges on vectors with finitely many particles: $Q$ commutes with the $\cos\bar\varphi(h)$, $h\in\mathcal F_{\mathbb R}(O)$, $\int h=0$. Conversely, a bounded $Q$ on the even space that commutes with these cosines lies in $\mathcal C_\Theta(O)$: $t^{-2}(1-\cos t\bar\varphi(h))\psi\to\frac12\varphi(h)^2\psi$ on vectors with finitely many particles, polarisation gives the $\varphi(h)\varphi(h')$, with $h=\partial_\kappa\mathsf k^\kappa$ all bilocal fields smeared in $O\times O$, and point splitting the $\Theta_{\kappa\lambda}(f)$ and $\Theta(f)$, $f\in\mathcal F(O)$, whose polynomial domain lies in that of these limits.

So $\mathcal C_\Theta(O)=\{\cos\bar\varphi(h):\int h=0\}'$ on the even space, $\mathcal A_\Theta(O)$ is the even part of the algebra of the field strengths, $\mathfrak R(\mathsf K_0(O))$ restricted to the even space (with $\mathsf K_0(O)$ below), and the multiplet has the same weak commutants; the closures of the $T_{\kappa\lambda}(f)$ are self-adjoint by the linear bound and affiliated by this identification; (W) follows from the Weyl relations. For $\mathsf n$ colours the same proof gives the net, with the $\mathrm O(\mathsf n)$-invariant parts, the bounded functions of the operators $\sum_a(\sum_i\alpha_ix_{ia}+\beta_ip_{ia})^2$ on $L^2(\mathbb R^{n'\times\mathsf n})$, for every finite number $n'$ of modes, generating the commutant of the rotations of the colours; at $\mathsf n=3$ it gives the affiliation of the $T_{\kappa\lambda}(f)$ as well.

\emph{The defect.} For a region $O$ let $\mathsf K(O)$ and $\mathsf K_0(O)$ be the closed real spans of the one-particle vectors of the $\varphi(h)$, $h\in\mathcal F_{\mathbb R}(O)$, respectively with $\int h=0$; at one time, $\mathsf K[D]$ is the closed real span of the one-particle vectors $u_f+v_g$ of $\varphi(0,f)$, $\pi(0,g)$, $f,g\in C^\infty_c(D)$ real, and $\mathsf K_0[D]$ that with $\int f=0$; $D_\rho=\{|\vec x|<\rho\}$ and $D^e_\rho=\{|\vec x|>\rho\}$. For $O$ the double cone over the disc $D_\rho$ the Cauchy data at the time $0$ and the finite speed give $\mathsf K(O)=\mathsf K[D_\rho]$, $\mathsf K_0(O)=\mathsf K_0[D_\rho]$, $\mathsf K(O^\times)=\mathsf K[D^e_\rho]$; and outside the charges escape to infinity: for $b\in C^\infty_c(D^e_\rho)$ with $\int b=1$, $b_n=n^{-2}b(\cdot/n)$ has $\|u_{b_n}\|^2=n^{-1}\|u_b\|^2\to0$, so $u_f=\lim_nu_{f-(\int f)b_n}$ and $\mathsf K_0(O^\times)=\mathsf K(O^\times)=\mathsf K[D^e_\rho]$.

The symplectic complements: $\mathsf K[D^e_\rho]'=\mathsf K[D_\rho]$, because for $w=w_1+iw_2\in\mathsf K[D^e_\rho]'$, $w_1,w_2$ in the real subspace of the $w$ with $w(-\vec k)=\overline{w(\vec k)}$, the distributions $F,G$ with $\tilde F=w_1$, $\tilde G=w_2/\omega$ have ${\rm Im}\langle u_f+v_g,w\rangle=\frac12(\int fG-\int gF)$, so they vanish on $D^e_\rho$ and are carried by $\bar D_\rho$; a dilation by $\lambda<1$, uniformly bounded on $L^2(d^2k/\omega)$ and on $L^2(\omega\,d^2k)$ and strongly continuous as $\lambda\uparrow1$, and a mollifier make $w$ a limit of vectors $u_f+v_g$ with $f,g\in C^\infty_c(D_\rho)$; so $\mathsf K[D_\rho]'=\mathsf K[D^e_\rho]''=\mathsf K[D^e_\rho]$; and $\mathsf K_0[D_\rho]'=\mathsf K[D^e_\rho]+\mathbb Rv_{g_1}$, $g_1=1$ near $\bar D_\rho$, because $u_f+v_g\mapsto\int f=2{\rm Im}\langle u_f+v_g,v_{g_1}\rangle$ is a continuous functional on $\mathsf K[D_\rho]$ with the kernel $\mathsf K_0[D_\rho]$.

The duality theorem for the free field, $\mathfrak R(\mathsf K)'=\mathfrak R(\mathsf K')$, turns these into $\mathcal C_\Theta(O^\times)=\mathcal A_\Theta(O^\times)'$, the even part of $\mathfrak R(\mathsf K[D_\rho])$, and $\mathcal C_\Theta(O)=\mathcal A_\Theta(O)'$, the even part of $\mathfrak R(\mathsf K[D^e_\rho]+\mathbb Rv_{g_1})$; the Weyl relation $e^{i\bar\varphi(h)}e^{i\mathsf s\bar\pi(g_1)}=e^{-i\mathsf s\int h}e^{i\mathsf s\bar\pi(g_1)}e^{i\bar\varphi(h)}$ gives the non-commuting pair. For $\mathsf n\ge1$ colours, $h\in\mathcal F_{\mathbb R}(O)$ with $\mathsf c=\int h\ne0$ and $\vec n$ uniform on the sphere $S^{\mathsf n-1}$, the mean over $\mathrm O(\mathsf n)$ of $e^{i\bar\varphi^1(h)}$ is $X=\int e^{i\bar\varphi(h\vec n)}d\vec n$, $\varphi(h\vec n)=\sum_an_a\varphi^a(h)$, and its restriction to the invariant vectors lies in $\mathcal C_\Theta(O^\times)$ by the Weyl relations.

If its restriction to the invariant vectors lay in $\mathcal A_\Theta(O)$, then, $\Omega$ being separating for the algebra of the field of $O$, so that the restriction to the invariant vectors is injective on its invariant part, $X$ would lie in the algebra of the field strengths of $O$, which commutes with the shifts $e^{i\vec s\cdot\bar\pi(g_1)}$, $\bar\pi(g_1)=(\bar\pi^a(g_1))_a$; but conjugation by them multiplies $e^{i\bar\varphi(h\vec n)}$ by $e^{i\mathsf c\vec s\cdot\vec n}$, and the one-particle part of the difference on $\Omega$, $ie^{-\|u_h\|^2/2}u_h\otimes\int(e^{i\mathsf c\vec s\cdot\vec n}-1)\vec n\,d\vec n$, $u_h$ the one-particle vector of $\varphi(h)$, is $\ne0$ for small $\vec s\ne0$.

So the restriction of $X$ lies in $\mathcal A^d_\Theta(O)=\mathcal C_\Theta(O^\times)$ and outside $\mathcal A_\Theta(O)$. The even parts, restricted to the even Fock space, of the nets $\mathfrak F_\Lambda(O)=\{e^{i\bar\varphi(f)}:f\in\mathcal F_{\mathbb R}(O),\ \int f\in\Lambda\}''$ are local, lie between, are injective in $\Lambda$ (an $h$ with $\int h\in\Lambda\setminus\Lambda'$ separates), and have the wedge algebras of $\Theta$, because in a wedge the field strengths generate the field, the charges escaping to infinity inside the wedge: for $f\in\mathcal F_{\mathbb R}(W)$, $f_n=n^{-3}f(\cdot/n)\in\mathcal F_{\mathbb R}(W)$ has $\int f_n=\int f$ and $\|u_{f_n}\|^2=n^{-1}\|u_f\|^2\to0$, while $f-f_n$ has the integral $0$. At $\mathsf m>0$ the cosines of all $h\in\mathcal F_{\mathbb R}(O)$ are reached, as $\varphi(h)=\varphi(\partial_\kappa(-\mathsf m^{-2}\partial^\kappa h))$, and the one-particle duality, which holds with $\omega=(|\vec k|^2+\mathsf m^2)^{1/2}$, gives $\mathcal C_\Theta(O^\times)=\mathcal A_\Theta(O)$.

This is the case $\mathsf n=1$. For $\mathsf n$ colours at a mass $\mathsf m>0$ the steps (1)--(4) run with the colour sums in the place of the products of one field: the commutator of two $\Theta(f_i)$ is the colour sum of the bilocal fields, weighted by the matrix $\mathsf H$ of the second derivatives of the commutator function of the mass $\mathsf m$, which is invertible at short timelike separations, $\mathsf m(z\cdot z)^{1/2}<2$, and differs from that of the mass $0$ by terms less singular at $z=0$, which the point splitting of (2) does not see; the double cones of (4) are taken small; and the commutators of the colour sums $\sum_a\varphi^a(h)\varphi^a(h')$ have the structure constants of the display of (4) at every mass.

So every $Q\in\mathcal C_\Theta(O)$, $O$ open and connected, commutes weakly with the colour sums of all real $h,h'\in\mathcal F(O)$ of integral $0$, and, by $\varphi^a(h)=\varphi^a(\partial_\kappa(-\mathsf m^{-2}\partial^\kappa h))$, with those of all real $h,h'\in\mathcal F(O)$; among them are the squares $\sum_a\varphi^a(g)^2$, whose closures generate $\mathcal A_S(O)$, and point splitting gives the $S(f)$, $f\in\mathcal F(O)$: $\mathcal C_\Theta(O)\subset\mathcal C_S(O)$, so $\mathcal A_S(O)\subset\mathcal A_\Theta(O)$ on $\mathcal H_\Theta=\mathcal H_S$. $\Theta(f)=\frac12S((\Box+2\mathsf m^2)f)$ gives the converse. So $\mathcal A_\Theta=\mathcal A_S$ on every connected open region, wedges included, and Theorem~\ref{C-363-3}(e) for $\Theta$ with the duality theorem for the free field gives, as at the mass $0$ below, $\mathcal A^d_\Theta(O)=\mathcal C_\Theta(O^\times)$, the invariant part of the algebra of the field of $O$, which is $\mathcal A_S(O)=\mathcal A_\Theta(O)$: the net of $\Theta$, as that of $S$, is its own dual net, for every $\mathsf n\ge1$.

\emph{The massless square.}
We first establish the operator core for $S={:}\varphi\cdot\varphi{:}$.
This uses particle number, rather than the energy bound used for
$\Theta$. Work in the one-particle realization $L^2(\mathbb R^2,d^2p)$,
with $\omega_p=|p|$. The Wick expansion of $S(f)$ has a number-preserving
part and pair-creation and pair-annihilation parts. Apart from the fixed
Fourier normalization and the factor $2$ in the number-preserving part,
their kernels are
\[
 K_f(p,q)=\frac{\hat f(\omega_p-\omega_q,p-q)}
                    {\sqrt{\omega_p\omega_q}},\qquad
 B_f(p,q)=\frac{\hat f(\omega_p+\omega_q,p+q)}
                    {\sqrt{\omega_p\omega_q}}.
\]
For $f\in\mathcal S$, the first kernel defines a bounded operator and
the second is Hilbert--Schmidt. To see the first assertion, on
$|p|,|q|\ge1$ use Schur's test and
$|\hat f(\omega_p-\omega_q,p-q)|\le C_M(1+|p-q|)^{-M}$, with $M>2$.
The complementary part is Hilbert--Schmidt: $|p|^{-1}$ is locally
integrable in two dimensions, and the same decay controls the other
variable when its radius is large. For the pair kernel, decay in the
energy sum gives
\[
 \iint |B_f(p,q)|^2\,d^2p\,d^2q
 \le C_M\int_0^\infty\!\int_0^\infty
                 (1+r+s)^{-M}\,dr\,ds<\infty
 \qquad(M>2).
\]
The finitely many colours only change the constants.

Let $\mathcal F_{\le k}$ denote the subspace of vectors with at most
$k$ particles. The creation and annihilation formulas therefore give
\[
 S(f)\mathcal F_{\le k}\subset\mathcal F_{\le k+2},\qquad
 \|S(f)\psi\|\le c_f(k+2)\|\psi\|.
\]
These are particle-sector bounds; the operator $N=1+H$ retains its
earlier meaning. For real $f$ and $\psi\in\mathcal F_{\le k}$, iteration
gives
\[
 \|S(f)^r\psi\|\le
 c_f^r\prod_{j=1}^r(k+2j)\,\|\psi\|.
\]
The series $\sum_{r\ge0}t^r\|S(f)^r\psi\|/r!$ converges for
$2c_f|t|<1$. Every vector of $\mathcal D_0^S$ has finitely many particles,
and $S(f)$ preserves $\mathcal D_0^S$. This domain is dense in
$\mathcal H_S$ by its definition. The analytic-vector theorem, in the
form of \cite[Theorem~2, pp.~180--181]{Nussbaum1965}, thus proves
that $S(f)$ is essentially self-adjoint on $\mathcal D_0^S$, without
using any identification of its weak commutant or its net.

The identification of $\mathcal H_S$ with the invariant subspace is
also made before using weak commutants. The bilocal commutator
construction used in (e), with $S$ in place of $\Theta$, gives the
colour sums $\sum_a\varphi^a(g)\varphi^a(g')$ and their squares for all
real compactly supported tests, with no integral-zero restriction.
They are limits of polynomials in the $S(f)$ on vectors of finitely
many particles, with bounds uniform on each $\mathcal F_{\le k}$.
The third commutator and one further commutator give the squares;
polarization gives the colour sums. The first fundamental theorem for
$\mathrm O(\mathsf n)$ expresses every invariant field polynomial in
these colour sums. Applying their polynomial approximants to the vacuum
and using those finite-particle bounds shows that
$\mathcal H_S$ contains the invariant polynomial vectors. Those vectors
are dense in the invariant subspace, while every $S$-polynomial vector
is invariant. Hence the two spaces agree.

On the common finite-particle energy-smooth domain the symmetric Wick
operator extends its restriction to $\mathcal D_0^S$. Essential
self-adjointness identifies its closure with the closure just obtained.
If a bounded operator $Q$ commutes weakly with $S(f)$ on
$\mathcal D_0^S$, the adjoint-domain criterion first gives
$Q\mathcal D_0^S\subset D(\overline{S(f)})$ and
$\overline{S(f)}Q\psi=QS(f)\psi$ there. Approximation in the graph norm
then gives $Q\overline{S(f)}\subset\overline{S(f)}Q$; the same holds
for $Q^*$. Thus $Q$ commutes with the bounded functions of
$\overline{S(f)}$, and the Wick products and the limits described above can
be taken on the common core.

For $S$, the squares $\sum_a\varphi^a(g)^2$ are reached for every real
$g\in\mathcal F(O)$ by these commutators. Consequently
$\mathcal A_S(O)$ is the invariant part of the algebra of the field
of $O$, and
$\mathcal A_\Theta(O)\subsetneq\mathcal A_S(O)$ at $\mathsf m=0$.
The field strengths generate the field in every wedge, so
$\mathcal A_S(\tilde W)=\mathcal A_\Theta(\tilde W)$ and
$\mathcal A_S^d=\mathcal A_\Theta^d$. Finally,
Theorem~\ref{C-363-3}(e) for $\Theta$ and the duality theorem for the
free field give
$\mathcal A_\Theta^d(O)=\mathcal C_\Theta(O^\times)$, the invariant
part of the algebra of the field of $O$. This is $\mathcal A_S(O)$,
and therefore
$\mathcal A_S=\mathcal A_S^d=\mathcal A_\Theta^d$.

\subsection{Proof of Theorem~\ref{C-363-7}}

\emph{(a)} Let $\Gamma(g)$, $g\in G^\Sigma$, be the gauge transformations on $L^2(G^{\mathsf B})$, $\mathfrak F(Y)=\mathcal B(L^2(G^Y))\otimes1$, and $P=\int\Gamma(g)\,dg$ the projection on $\mathcal H_{\rm ph}$. $\mathfrak N=\mathfrak F(Y)\cap\Gamma(G^\Sigma)'$ is a von Neumann algebra and $P\in\Gamma(G^\Sigma)''\subset\mathfrak N'$, so $\mathsf A(Y)=\mathfrak NP$ is a von Neumann algebra with the commutant $P\mathfrak N'P|_{\mathcal H_{\rm ph}}$, and $\mathfrak N'=\mathfrak F(\mathsf B\setminus Y)\vee\Gamma(G^\Sigma)''$ is the weak closure of the span of the $X\Gamma(g)$, $X\in\mathfrak F(\mathsf B\setminus Y)$. As $P\Gamma(g)=P$, $PX\Gamma(g)P=PXP=PX^\natural P$, $X^\natural=\int\Gamma(g)X\Gamma(g)^{-1}dg\in\mathfrak F(\mathsf B\setminus Y)\cap\Gamma'$; so $\mathsf A(Y)'=\mathsf A(\mathsf B\setminus Y)$. At $Y=\varnothing$ this is $\mathsf A(\mathsf B)=\mathcal B(\mathcal H_{\rm ph})$. Without a common site the gauge groups of $Y_1$ and $Y_2$ are independent, $(X_1\otimes X_2)^\natural=X_1^\natural\otimes X_2^\natural$, and the products are weakly total.

\emph{The flips.} $V\mapsto cV$ keeps the Haar measure; $\chi_j(-g)=(-1)^{2j}\chi_j(g)$ gives $\mathsf V(c)M_{\chi_j(U(C'))}=c(C')^{2j}M_{\chi_j(U(C'))}\mathsf V(c)$. Two groups of coboundaries are told apart. If $c_b=z(b_-)z(b_+)$ on all the bonds of the layer, $z$ a sign field on all its sites, $\mathsf V(c)$ is the gauge transformation $\Gamma(z)$ by the central field $z$, which is $1$ on $\mathcal H_{\rm ph}$. For a set $S$ of plaquettes, the coboundaries $B^1(S)=\{dz|_{\mathsf B_S}\}$, $z$ a sign field on the sites of $\mathsf B_S$, form the subgroup of $Z^1(S)$ by which $H^1(S)=Z^1(S)/B^1(S)$ is formed. For $c=dz|_{\mathsf B_S}$ and $\tilde z$ the continuation of $z$ by $1$ to all sites, $\mathsf V(c)=\Gamma(\tilde z)\mathsf V(c^\circ)$, where $c^\circ=d\tilde z$ on $\mathsf B\setminus\mathsf B_S$ and $c^\circ=1$ on $\mathsf B_S$, the sign field $c^\circ$ of the statement; so on $\mathcal H_{\rm ph}$, $\mathsf V(c)=\mathsf V(c^\circ)$, a product of the parities $\sum_j(-1)^{2j}\Pi^j_b$ of bonds $b\notin\mathsf B_S$, which lies in $\mathsf A(\mathsf B\setminus\mathsf B_S)$; it is not $1$ when, say, $z=-1$ at one site $x$ alone and a plaquette $q\notin S$ has at $x$ one bond in $\mathsf B_S$ and one off it, for then it anticommutes with $M_{t_q}\ne0$.

As $\Gamma(\tilde z)$ commutes with every gauge-invariant operator and $\mathsf V(c^\circ)$ with every operator of the bonds of $\mathsf B_S$, ${\rm Ad}\,\mathsf V(c)$ is the identity on $\mathsf A(\mathsf B_S)$ for $c\in B^1(S)$; and on $\mathcal H_{\rm ph}$ the flips of one class of $H^1(S)$ differ by elements of $\mathsf A(\mathsf B\setminus\mathsf B_S)$, so that beside $\mathsf A(\mathsf B\setminus\mathsf B_S)$ one flip of each class is enough. On the bonds where $c=-1$, $\mathsf V(c)$ acts in the spin-$j$ sector of the bond as $(-1)^{2j}$, so it lies in the algebra of the spin projections of these bonds. On a periodic layer every bond lies in two plaquettes, so $\prod_p\prod_{b\in\partial p}c_b=\prod_bc_b^2=1$, and the plaquettes where $c$ is not closed are even in number.

\emph{(b)} Let $Y$ commute with the $M_{t_p}$, $p\in S$, and the $\Pi^j_b$, $b\in\mathsf B_S$, on $L^2(G^{\mathsf B_S})$. Then $Y$ keeps the polynomials in the matrix entries and commutes with the Laplacian $\Delta_b=-\sum_jj(j+1)\Pi^j_b$ of every bond; for functions $F,F'$ with $M_F,M_{F'}$ in the algebra generated,
$$[[\Delta_b,M_F],M_{F'}]=2M_{D_b(F,F')},\qquad D_b(F,F')=\sum_a(J^a_bF)(J^a_bF'),$$
$J^a_b$ the infinitesimal left translations of the bond $b$, so the real algebra $\mathfrak W$ of such functions is closed under every $D_b$. For two closed paths $\mathsf M=V_b\mathsf a$, $\mathsf N=V_b\mathsf c$ that run once through $b$, $D_b({\rm tr}\,\mathsf M,{\rm tr}\,\mathsf N)=\frac14[{\rm tr}\,\mathsf M^{-1}\mathsf N-{\rm tr}\,\mathsf M\mathsf N]$ and ${\rm tr}\,\mathsf M\,{\rm tr}\,\mathsf N={\rm tr}\,\mathsf M\mathsf N+{\rm tr}\,\mathsf M^{-1}\mathsf N$: paths may be fused at a common bond. Four fusions in every chain $T$ of three plaquettes, in a row or around a site, give the seven traces $\frac12{\rm tr}$ of $\alpha_i$, $\alpha_i\alpha_k$, $\alpha_1\alpha_2\alpha_3$, $\alpha_i$ the boundaries at one site. With $g=\mathsf s(g)+i\mathsf v(g)\cdot\sigma$,
$$\mathsf s(g_1g_2g_3)=\mathsf s_1\mathsf s_2\mathsf s_3-\mathsf s_1\mathsf v_2\cdot\mathsf v_3-\mathsf s_2\mathsf v_1\cdot\mathsf v_3-\mathsf s_3\mathsf v_1\cdot\mathsf v_2+\det(\mathsf v_1,\mathsf v_2,\mathsf v_3),$$
so the seven numbers are the Gram matrix and the determinant of the $\mathsf v(V(\alpha_i))$, which fix the triple up to a rotation; every closed path of $T$ is a word in the $\alpha_i$, and so two configurations of $T$ with the same seven numbers differ by a gauge transformation of $T$. Two chains that share two plaquettes fix the gauge transformation on their common sites up to the centre, as two non-commuting elements of $\mathrm{SU}(2)$ have the centraliser $\{\pm1\}$; along the chains of $S$ the configurations with the same functions in $\mathfrak W$ are those that differ by a gauge transformation and a closed sign field. Let $\mathcal K_S$ be the compact group of the pairs $(g,c)$, $g$ a gauge transformation of the sites of $\mathsf B_S$ and $c\in Z^1(S)$, acting by $((g,c)V)_b=c_bg(b_-)V_bg(b_+)^{-1}$, and $\lambda$ its representation on $L^2(G^{\mathsf B_S})$; $\lambda(\mathcal K_S)$ commutes with the $M_{t_p}$ and the $\Pi^j_b$.

Let $\delta_0=\prod|\mathsf v(V(\alpha_p))\times\mathsf v(V(\alpha_q))|^2$, the product over the pairs of adjacent plaquettes of $S$, $\alpha_p,\alpha_q$ their boundaries begun at an end of the common bond; it is an invariant polynomial, $\not\equiv0$ on the connected space $G^{\mathsf B_S}$, so its zero set is a null set. Off it the stabiliser of $V$ in $\mathcal K_S$ is $\{(z,dz)\}$, $z$ central, which acts trivially (an element of $\mathrm{SU}(2)$ that commutes with two elements that do not commute is $\pm1$), and the seven numbers of the chains separate the orbits; so by Stone--Weierstrass on the compact sets $\{|\delta_0|\ge\epsilon\}$, $M_F$ lies in the generated algebra for every continuous $\mathcal K_S$-invariant $F$, and $Y$ commutes with all of them.

The polynomials are the sum of the isotypic parts ${\rm Pol}_\pi$ of the irreducible representations $(\pi,\mathsf V_\pi)$ of $\mathcal K_S$ that are trivial on the stabiliser; at a point $V$ off the zeros of $\delta_0$ the functionals $v\mapsto m_v(V)$ of the equivariant linear maps $m:\mathsf V_\pi\to{\rm Pol}$ fill $\mathsf V_\pi^*$, the orbit of $V$ being $\mathcal K_S$ modulo the stabiliser (Peter--Weyl and Stone--Weierstrass on it). With $\mathsf n=\dim\pi$ such maps $m^{(i)}$ whose functionals at one point $V_1$ form a basis of $\mathsf V_\pi^*$, $\mathsf M=(m^{(i)}_{e_k})_{ik}$ in a basis $(e_k)$ of $\mathsf V_\pi$ and $\delta_\pi=\det\mathsf M$, Cramer's rule writes $|\delta_\pi|^2m_v$ as a combination of the $m^{(i)}_v$ with invariant polynomial coefficients, with which $Y$ commutes; so $Ym_{e_k}=\sum_lm_{e_l}\tau_{lk}$ for every equivariant $m$ on $\{\delta_\pi\ne0\}$, with $\tau=\mathsf M^{-1}(Ym^{(i)}_{e_k})_{ik}$.

The vector fields $\xi_{b,F}=\sum_a(J^a_bF)J^a_b=\frac12([\Delta_b,M_F]-M_{\Delta_bF})$, $F$ a real invariant polynomial, commute with $Y$ on the polynomials and with $\lambda(\mathcal K_S)$, and carry equivariant maps to equivariant maps; the identity for the equivariant map $\xi_{b,F}m$, set beside $\xi_{b,F}$ applied to the identity for $m$ ($\xi_{b,F}$ is a derivation), gives $\mathsf M\,\xi_{b,F}\tau=0$, so $\xi_{b,F}\tau=0$ on $\{\delta_\pi\ne0\}$. The $\xi_{b,F}(V)$ span the tangent space on a dense open set: at a bond $b$ of a plaquette $p$ of a chain $T\subset S$, with $\alpha$ the boundary of $p$ begun with $b$ and $w_1,w_2$ closed paths at $b_-$ through the two bonds $c_1,c_2$ that a spanning tree of $\mathsf B_T\setminus\{b\}$ leaves out, the traces ${\rm tr}\,V(\alpha)$, ${\rm tr}\,V(\alpha w_i)$ are invariant and their derivatives $(J^a_b\,\cdot\,)_a$ are $-\mathsf v(V(\alpha))$, $-\mathsf v(V(\alpha w_i))$, whose determinant is a polynomial that is not $0$ at a configuration with $\mathbf 1$ on the tree and $V(\alpha)=i\sigma_3$, $V(\alpha w_1)=i\sigma_1$, $V(\alpha w_2)=i\sigma_2$ (the bonds $b,c_1,c_2$ are free there).

So $d\tau=0$ on a dense open part of $\{\delta_\pi\ne0\}$, $\tau$ is constant, $\tau^\pi$, near $V_1$, and the matrix of polynomials $(Ym^{(i)}_{e_k})_{ik}-\mathsf M\tau^\pi$, which vanishes near $V_1$, vanishes identically: here analyticity enters. Hence $Ym_v=m_{\tau^\pi v}$ for all equivariant $m$, and with $y_\pi(\kappa)=\dim\pi\,{\rm tr}(\tau^\pi\pi(\kappa)^{-1})$, $YP_\pi=\lambda(y_\pi)\in\lambda(\mathcal K_S)''$, $P_\pi$ the projection on the isotypic part; summed over $\pi$, $Y\in\lambda(\mathcal K_S)''$. So the generated algebra is the commutant of the gauge transformations of $S$ and of the $\mathsf V(c)$, $c\in Z^1(S)$; restricted to $\mathcal H_{\rm ph}$ and combined with (a), this is the pair of identities for $\mathsf A_\bullet(S)$, and it exhibits $\mathsf A_\bullet(S)$ as the join of the $\mathsf A(\mathsf B_T)$ over the chains.

For admissible $S$ the join of the $\mathsf A(\partial p)$, $p\in S$, is $\mathsf A_\bullet(S)$: the generators $M_{t_p}$ and $\Pi^j_b$, $b\in\partial p$, lie in $\mathsf A(\partial p)$, and $\mathsf A(\partial p)\subset\mathsf A(\mathsf B_T)\subset\mathsf A_\bullet(S)$ for a chain $T\subset S$ of three plaquettes through $p$, which admissibility provides. The $\mathsf V(c)$, $c\in Z^1(S)$, are functions of spin projections of $\mathsf B_S$ and commute with $\mathsf A_\bullet(S)$: they are central. ${\rm Ad}\,\mathsf V(c)$ is the identity on $\mathsf A(\mathsf B_S)$ for $c\in B^1(S)$ (the flips above), so $H^1(S)$ acts, and the decomposition into its characters is the spectral decomposition of this action of a finite abelian group; the summand of a character $\kappa$ holds $M_{\chi_{1/2}(U(C_\kappa))}\ne0$ for a closed path $C_\kappa$ with $c(C_\kappa)=\kappa([c])$.

If every $\mathbb Z_2$-cycle is a sum of boundaries, every closed $c$ is a coboundary. For one plaquette, the cyclic shift of the four bonds times the projection on the spin $1$ at one corner and $0$ at the others commutes with $M_{t_p}$ and all $\Pi^j_b$ and moves the charge to the next corner; for two adjacent plaquettes use reflection in the perpendicular bisector of the common bond, composed with the projection on charge $1$ at one end and $0$ at all other vertices. On this projected block the zero charges at the outer vertices force equal bond spins along each of the two outer three-bond paths. Reflection reverses each path and preserves each individual bond-spin projection, as well as each plaquette trace, but moves the endpoint charge; neither lies in the von Neumann algebra generated by the gauge transformations and the closed flips, which is the commutant of all invariant operators, so the traces and spin projections generate less. The traces generate multiplications, and $M_{t_p}$ does not commute with $\Pi^j_b$, $b\in\partial p$.

\emph{Two plaquettes with one common bond: the full algebras generate.} Let $p,q$ share the bond $c_0$ with the ends $x,y$; their union has seven bonds and six sites, and at $x$, as at $y$, three bonds meet: $a\in\partial p$, $c_0$ and $a'\in\partial q$. Let $\mathfrak B$ be the algebra of all gauge-invariant bounded operators on $L^2(G^7)$, the functions of these bonds, and $\mathfrak N\subset\mathfrak B$ the von Neumann algebra generated by the gauge-invariant operators of the four bonds of $p$ and those of the four bonds of $q$; we show $\mathfrak N=\mathfrak B$, which, restricted to $\mathcal H_{\rm ph}$ by a normal homomorphism, gives $\mathsf A(\partial p)\vee\mathsf A(\partial q)=\mathsf A(\partial p\cup\partial q)$.

By Peter--Weyl, $L^2(G^7)=\bigoplus_{\mathbf j}\mathcal K(\mathbf j)$ over the spins $\mathbf j$ of the seven bonds, $\mathcal K(\mathbf j)=\bigotimes_z\mathcal V_z(\mathbf j)$, $\mathcal V_z(\mathbf j)$ the product of the spin-$j_b$ representations of the bonds at the site $z$, on which the gauge group at $z$ acts; $\mathcal V_z(\mathbf j)=\bigoplus_JV_J\otimes\mathsf M_z(\mathbf j,J)$, and a label $(\mathbf j,\mathbf J)$, $\mathbf J$ the charges at the six sites, occurs if the polygon inequalities hold at every site. So $\mathfrak B=\bigoplus_{\mathbf J}\mathcal B(\mathsf M_{\mathbf J})$, $\mathsf M_{\mathbf J}=\bigoplus_{\mathbf j}\bigotimes_z\mathsf M_z(\mathbf j,J_z)$. (i) \emph{Labels.} Let $Y\in\mathfrak B$ commute with $\mathfrak N$.

It commutes with the Casimir operators of the bonds, so it keeps every $\mathcal K(\mathbf j)$; being gauge invariant, it keeps the charges $\mathbf J$; and for a corner $(b_1,b_2)$ of $p$ or $q$ at $z$ the spectral projections of the Casimir operator of the gauge group at $z$ acting on the two bonds of that plaquette lie in $\mathfrak N$, so $Y$ commutes with $t^z_{b_1b_2}=J_{b_1,z}\cdot J_{b_2,z}$ on every $\mathcal K(\mathbf j)$. At the four sites with two bonds $\mathsf M_z$ is one-dimensional; at $x$ and at $y$, $t_{ac_0}$ and $t_{c_0a'}$ generate all operators of $\mathsf M_z(\mathbf j,J_z)$ by Lemma~\ref{C-363-27}. So $Y=\sum y(\mathbf j,\mathbf J)P(\mathbf j,\mathbf J)$, $P(\mathbf j,\mathbf J)$ the projection on the label. (ii) \emph{One plaquette made large.} If two labels with the same charges occur and agree off $\partial p$, $y$ takes one value at them.

Indeed $L^2(G^{\partial p})=\bigoplus_\pi V_\pi\otimes{\rm Mult}_\pi$ over the charges $\pi=(f_z)$ at the four corners of $p$, ${\rm Mult}_\pi$ with the basis $e_{\mathbf i}$ indexed by the spins $\mathbf i$ of $\partial p$ with $f_z\in i_{b_1}\otimes i_{b_2}$ at every corner (Clebsch--Gordan is free of multiplicities); the gauge-invariant operators of $\partial p$ are the operators $1\otimes\mathcal B({\rm Mult}_\pi)$, and $Y$, commuting with them, is the identity on every ${\rm Mult}_\pi$; so $y$ takes one value on the spins $\mathbf i$ of $\partial p$ of every $\pi$ that occurs with the given spins off $\partial p$ and the given charges; and with $N$ above all charges that occur and $e_b\in\{0,\frac12\}$ with $e_{b_1}+e_{b_2}\equiv f_z$ modulo $1$ at the four corners (they exist, the charges at the corners being fixed modulo $1$ by the charges $\mathbf J$ and the spins off $\partial p$, with an integer sum), the spins $(N+e_b)_{b\in\partial p}$ belong to every such $\pi$, which joins them all. (iii) \emph{A monotone path.} Raising the four spins of one plaquette by $\frac12$ keeps the charges and the polygon inequalities, two spins at each corner growing by $\frac12$.

By such raisings two labels with the same charges are brought to an integer difference, twice their difference modulo $2$ being a cycle and so a sum of $\partial p$ and $\partial q$; raised $2N$ times more on both plaquettes, $N$ large, they are joined through the label that takes the larger of the two spins on the six bonds that lie on one plaquette and the smaller on $c_0$, changing first the bonds of $p$, then those of $q$; all labels on the way occur. With (ii) at every change, $y$ depends on $\mathbf J$ alone. (iv) \emph{The charges.} The projections $E^z_J$ on the charge $J$ at $z$ lie in $\mathfrak N$. At the four outer sites they are spectral projections of Casimir operators of one plaquette.

At $x$, with $A=J_{a,x}$, $B=J_{c_0,x}$, $D=J_{a',x}$, the Casimir operator $(A+B+D)^2$ holds the cross term $t_{AD}$ of the two outer bonds, which belongs to neither plaquette; the identity of three commuting spins
$$(B^2+t_{AB})\,t_{AD}=A^2t_{BD}+t_{AB}t_{BD}-[t_{AB},t_{BD}]-[t_{AB},[t_{AB},t_{BD}]]$$
puts $\mathsf Ft_{AD}$ in $\mathfrak N$ on every block of spins, $\mathsf F=B^2+t_{AB}=\frac12((A+B)^2+B^2-A^2)$, and with it $\mathsf F(A+B+D)^2$. As $(A+B+D)^2$ commutes with all of $\mathfrak B$, and $\mathsf F$ is the number $\frac12(f(f+1)+s_c(s_c+1)-s_a(s_a+1))$ on the range where the spins of $a,c_0$ are $s_a,s_c$ and $(A+B)^2=f(f+1)$, a range given by projections of $\mathfrak N$, $(A+B+D)^2$ times such a range lies in $\mathfrak N$ where this number is not $0$, and with it $E^x_J$ times the range; the same holds where $f=0$, as $(A+B+D)^2=D^2$ there. The weakly closed ideal of the $n\in\mathfrak N$ with $nE^x_J\in\mathfrak N$ is $\mathfrak N\mathsf p$ with a central projection $\mathsf p$ of $\mathfrak N$ that holds these ranges; $1-\mathsf p$ commutes with $\mathfrak N$, so by (i)--(iii) it is the sum of the $P(\mathbf j,\mathbf J)$ over the labels of some charges, and if it were not $0$, each of its ranges of the kind above would have $f>0$ and $f(f+1)=s_a(s_a+1)-s_c(s_c+1)$; but raising $p$ by $\frac12$ keeps the charges and $f$ and turns such a range into one where the number is $\frac12(s_c-s_a)\ne0$, a range of $\mathsf p$: so $\mathsf p=1$. (v) $\mathfrak N$ holds the projections on the $\mathsf M_{\mathbf J}$, and an operator on $\mathsf M_{\mathbf J}$ that commutes with $\mathfrak N$ is a number by (i)--(iii); so $\mathfrak N$ cut down to $\mathsf M_{\mathbf J}$ is irreducible, and $\mathfrak N=\bigoplus_{\mathbf J}\mathcal B(\mathsf M_{\mathbf J})=\mathfrak B$.

This is the two-plaquette case of Lemma~\ref{C-363-24}(a).

\emph{(c)} $R$ deformation retracts to a circle around the hole: $H^1(R)=\mathbb Z_2$, generated by $[c_*]$, and $c_*(C)=-1$. $\mathsf A(\mathsf B\setminus\mathsf B_R)'=\mathsf A(\mathsf B_R)$ by (a). As $v$ represents the generator, (b) gives $\mathsf A(\mathsf B_R)=\mathsf A_\bullet(R)\oplus\mathsf A(\mathsf B_R)_-$, and $u={\rm sign}\,{\rm tr}\,U(C)$ is an odd self-adjoint unitary in $\{W_C\}''$, so $\mathsf A(\mathsf B_R)_-=\mathsf A_\bullet(R)u$; $\frac12(X+vXv)$ is the projection on the even part. $\mathsf A_\bullet(R)'=\mathsf A(\mathsf B\setminus\mathsf B_R)\vee\{v\}$ by (b), $v$ being one flip of the class that is not trivial and the flips of $B^1(R)$ acting on $\mathcal H_{\rm ph}$ through $\mathsf A(\mathsf B\setminus\mathsf B_R)$; and on $\mathcal H_{\rm ph}$, $v=\mathsf V_\gamma\mathsf V(c')$ with $c'$ carried by $\mathsf B\setminus\mathsf B_R$ and $\mathsf V(c')\in\mathsf A(\mathsf B\setminus\mathsf B_R)$; ${\rm Ad}\,u$ fixes $\mathsf A(\mathsf B\setminus\mathsf B_R)$ and changes the sign of $\mathsf V_\gamma$ ($c_\gamma(C)=-1$), which gives the second decomposition and $\mathsf A(\mathsf B\setminus\mathsf B_R)=\mathsf A_\bullet(R)'\cap\{u\}'$.

For a closed path $C'$, $vM_{\chi_j(U(C'))}v=c_*(C')^{2j}M_{\chi_j(U(C'))}$. If an assignment $\mathscr N$ equal to $\mathsf A$ on closed and open rectangles were additive and dual, then $\mathscr N(\mathsf B_R)=\bigvee_i\mathsf A(\mathsf B_{D_i})=\mathsf A_\bullet(R)$ and $\mathscr N(\mathsf B\setminus\mathsf B_R)=\mathsf A(\mathsf B_{\rm in})\vee\mathsf A(\mathsf B_{\rm out})=\mathsf A(\mathsf B\setminus\mathsf B_R)$, the two pieces of the complement having no common site; duality would give $\mathsf V_\gamma\in\mathsf A(\mathsf B\setminus\mathsf B_R)$. The four rectangles: every chain of three plaquettes of $R$ lies in one $D_i$ or in the union of two adjacent ones, and two adjacent rectangles glue by Lemma~\ref{C-363-24}(a), their overlap being a connected rectangle with a cycle; so $\mathsf A_\bullet(R)=\bigvee_i\mathsf A(\mathsf B_{D_i})$.

The commutative statement and the formula for the adjoint loop are Lemma~\ref{C-363-24}(b), (c) for the four rectangles as boxes, whose overlaps hold two independent cycles when no corner is one plaquette.

\emph{(d)} $S$ deformation retracts to the torus with two holes: $H^1(S)=\mathbb Z_2^3$, generated by the cut $\gamma$ and the two winding cuts. $\mathsf A(\mathsf B_S)'=\mathsf A(\mathsf B^\circ_{D_1}\cup\mathsf B^\circ_{D_2})=\mathsf N$ by (a), and (b) gives $\mathsf A_\bullet(S)'=\mathsf N\vee\{\mathsf V_\gamma,\mathsf V_1,\mathsf V_2\}$, the products of $\mathsf V_\gamma,\mathsf V_1,\mathsf V_2$ meeting every class of $H^1(S)$ and the flips of $B^1(S)$ acting on $\mathcal H_{\rm ph}$ through $\mathsf N$. The fundamental loop around $D_1$ lies in $\mathsf A(\mathsf B_S)$, commutes with $\mathsf N$, $\mathsf V_1$, $\mathsf V_2$ and anticommutes with $\mathsf V_\gamma$; the winding loops along the free row and column anticommute with one $\mathsf V_i$ each and commute with $\mathsf N$ and $\mathsf V_\gamma$: hence the two non-memberships. On the whole layer $H^1=\mathbb Z_2^2$ and $\mathsf A(\varnothing)=\mathbb C$; the four joint eigenspaces of $\mathsf V_1,\mathsf V_2$ are carried into one another by the fundamental winding loops, so none is $0$. Off one disc, $H^1$ is generated by the two fluxes, the loop around the disc being a boundary in homology modulo $2$ there.

\emph{(e)} A box has a shelling, an order of its plaquettes in which each new plaquette either has all its bonds already present or is glued to an earlier one along a single bond so that the overlap with what is present is connected, and every site of a new bond of $p_k$ that is already a site of $X_{k-1}$ lies on the boundary $\partial p_j$ of that earlier plaquette together with the overlap $\partial p_k\cap X_{k-1}$; Lemma~\ref{C-363-24}(a), the two-plaquette case included, carries the generation from step to step: at a step that brings new bonds it applies to $Y_1=X_{k-1}$ and $Y_2=\partial p_j\cup\partial p_k$, with $I=\partial p_j\cup(\partial p_k\cap X_{k-1})$ connected and holding the cycle $\partial p_j$, $\mathsf R_2=\partial p_k\setminus X_{k-1}$, and the site clause of the shelling as its hypothesis on the common sites, while the gauge-invariant operators of $Y_2$ are generated by those of $\partial p_j$ and of $\partial p_k$, by the two-plaquette case.

The solid ring is the union of two shellable halves, the cubes $[i_1,i_1+1]\times[i_2,i_2+1]\times[0,1]$ of the ring with $i_1\in\{0,1\}$ and those with $i_1\in\{1,2\}$, whose intersection has two components, the bonds of the two cubes with $i_1=1$, $i_2\in\{0,2\}$, each with a cycle; a site common to the two halves has $x_1\in\{1,2\}$, and the sites of the first half with $x_1=2$, as those of the second with $x_1=1$, are corners of these two cubes, so every site common to $\mathsf B_{S_1}$ and $\mathsf B_{S_2}$ is a site of $I$, the hypothesis of Lemma~\ref{C-363-24}(a) on the common sites, as in the general union of (e), where a site of $\mathsf B_{S_1}$ on a bond of $\mathsf B_{S_2}$ off $I$ is common to both; Lemma~\ref{C-363-24}(a) gives the join as the invariant operators that commute with the central gauge transformations constant on each component, that is with $\mathsf V(c_*)$, and duality gives the commutant; $\mathsf V(c_*)$ anticommutes with $W_C$ because $c_*(C)=-1$; for a dual surface meeting the ring in that section, $c_{\mathfrak s}c_*$ is carried by the complement.

With $r$ components the central gauge transformations constant on components, modulo the global ones, form $\mathbb Z_2^{r-1}$.

\emph{(f)} The substitution $U\mapsto sU$ keeps the product Haar measure, $t_p(sU)=(ds)_pt_p(U)$ and $\chi_j(U(C))(sU)=s(C)^{2j}\chi_j(U(C))(U)$. For a fat odd loop the same with $F(sU)=s(C)F(U)$. Both members may be multiplied by any non-zero numbers, which carries the rule to limit points.

\emph{(g)} The block map is formed from parallel transports along paths in the double block of $c$ and from averages of them, and it is covariant under gauge transformations; if $ds=1$ on the double block, $s$ is there the coboundary of a central field $z$, and $\mathcal M^n(sU)(c)=z(c_-)\mathcal M^n(U)(c)z(c_+)=s[c]\mathcal M^n(U)(c)$. The convergence of $\langle h(W_C(\mathcal M^{K-d}U))\rangle_K$, with the stated rate, is the convergence of the laws of the block fields at a fixed depth of Theorem~\YMMref{D}{8} (Corollary~\textup{\ref{IV-62-66}}), a bounded Borel function of the block field being tested. A flip $\bar\xi$ of the torus is the image of a closed sign field $\xi$ of the fine lattice, $d\xi=1$, under which the Wilson measure is invariant; by covariance the block-field laws are invariant under $\bar\xi$. For the bounding vortex line, $\eta_K=ds$ off the plaquettes pierced by $\partial\mathsf S$, $s=-1$ on the fine bonds that cross $\mathsf S$; the sign rule gives $Z_K[\eta_K;F\circ\mathcal M^{K-d}]=Z_K[1;F\circ\mathcal M^{K-d}\circ s]$ and covariance gives $\mathcal M^{K-d}(sU)=\bar\xi\mathcal M^{K-d}(U)$ on the blocks that miss $\partial\mathsf S$; the limit and its error come from Theorem~\YMMref{D}{8}; for $F=W_CF'$ the factor is $\bar\xi(C)=(-1)^{\#(C\cap\mathsf S)}$.

\emph{(h)} For continuous $h$ the first statement is Lemma~\ref{C-363-41} (Theorem~\ref{C-372-6} of Section~\ref{C-372-1}, read at one polygon), applied to the polygon $C$, a closed path of some $\Lambda_{K_0,\ell}$ taken in every $\Lambda_{K,\ell}$, $K\ge K_0$, with the functions $F_K$ of the bonds at the distance $\ge\mathsf d$ from $C$, uniformly in them; a bounded $h$ that is continuous off a null set lies between continuous functions $h_-\le h\le h_+$ with $\overline{h_+}-\overline{h_-}$ as small as we please, and $F_K$ is a combination of four functions with values in $[0,1]$. The law of $\frac12{\rm tr}\,g$ under Haar measure is $\frac2\pi\sqrt{1-w^2}\,dw$; so $\bar h=0$ for $h={\rm sign}$ and for the characters $\chi_j$, $j\ge\frac12$, written as functions of $w=\frac12{\rm tr}$, and $\bar h=\frac14$ for $h(w)=w^2$.

With $C$ in a site mirror and $F_K=\overline{\theta G'_K}G_K$, $|G_K|,|G'_K|\le1$ at a fixed distance, this is the fading between the vectors of the space of reflection positivity. The Bessel bound is Lemma~\ref{C-363-25}(a), and $\mathsf r(x)=1-\frac3{2x}+O(x^{-2})$, $n'\approx|\gamma|L^K$, $4\beta=4L^K\mathsf b$ give the limit $\exp(-3|\gamma|/8\mathsf b)$. The fading of the flips is Lemma~\ref{C-363-25}(b), (c), with $\omega_K(1-Q_0)\le\omega_K(M_{a_{p_0}})=\langle a_{p_0}\rangle_K$, the multiplication operators of one layer having the Wilson measure as their law in the state of the transfer operator.

\emph{(i)} \emph{Notation.} $x_{[D]}=L^D\mathsf b$ (Section~\ref{sec:block-field-laws}) is the coupling at the depth $D$; $\varepsilon(\mathsf t)=A_0\mathsf t^{-1/2}(\log\mathsf t)^{\mathfrak p_0}$ and $\delta(\mathsf t)=(A_1/A_0)\varepsilon(\mathsf t)$ \YMMref{(634)}{339}, with the constants $A_0,A_1>0$, $\mathfrak p_0\in\mathbb N$ of the list of Theorem~\YMMref{D}{8}; $\varpi(V)=\max_p|\partial V(p)-\mathbf 1|$, $|\cdot|$ the operator norm; and the small fields of the depth $D$ are those of the inner ball $\mathsf{In}_{[D]}=\{V:\varpi(V)<2\varepsilon(x_{[D]})\}$ of Theorem~\YMMref{D}{8} (the common-depth definitions in Section~\YMMref{116}{807}, the definitions preceding Lemma~\YMMref{113.10}{796}), an open gauge-invariant set (Lemma~\YMMref{113.10}{796}). $\mathsf R_{[D]}=2M^{-3}(\ell L^D)^3\exp[-c'A_0^2(\log x_{[D]})^{2\mathfrak p_0}]\exp[(\mathsf C^\Sigma_{\rm maj}+2+C_{\rm low}(\mathsf b))\ell^3]$ is the number \YMMref{(2187)}{806} at the level of the depth $D$; it reads $D$ and not $K$, and it tends to $0$ as $D\to\infty$, since $c'A_0^2>0$ and $2\mathfrak p_0\ge2$. An element $g$ of $G$ with the eigenvalues $e^{\pm i\phi}$, $\phi\in[0,\pi]$, has $|g-\mathbf 1|=2\sin\frac\phi2\le\phi$, so $|e^{iw\sigma_3}-\mathbf 1|\le|w|$; and replacing the four factors of $\partial(kV)(p)$ by those of $\partial k(p)$ one by one, $|\partial(kV)(p)-\partial k(p)|\le\sum_{b\in\partial p}|V_b-\mathbf 1|$ for $G$-valued fields $k$, $V$ and the bondwise product $kV$.

\emph{The floor.} Let $D$ be so large that $x_{[D]}\ge\mathsf t_\natural$ and $\mathsf R_{[D]}<1$, let $K\ge D+1$, and let $t=K-D$ be the level of the depth $D$. The setting of the small-field chain of Theorem~\YMMref{D}{8} asks of its top level only that it lie between $1$ and $K-1$ and carry a coupling $\ge\mathsf t_\natural$ (Section~\ref{sec:cutoff-calibration}); so it holds with the top level $t$, at a fixed depth. At the source $0$, Proposition~\textup{\YMMref{115.2}{806}} gives a bounded Borel weight $\mathfrak w^{\rm sf}_t\ge0$ defined before Proposition~\YMMref{115.2}{806} such that the image of $\mathfrak w^{\rm sf}_t\rho_0\,dU$ under $\mathcal M^t$ is the measure $\mathbf 1_{\mathsf{In}_{[D]}}\nu_t$ of the chain \YMMref{(2188)}{806}, $\int F\mathfrak w^{\rm sf}_t\rho_0\,dU\le\int F\rho_0\,dU$ for gauge-invariant bounded Borel $F\ge0$ \YMMref{(2189)}{806}, and $0\le Z-\int\mathfrak w^{\rm sf}_t\rho_0\,dU\le\mathsf R_{[D]}Z$ \YMMref{(2190)}{806}, $Z=\int\rho_0\,dU$.

So the measure of the chain lies below $Z$ times the law on gauge-invariant functions (histories with a large field at a finer level, which may end in $\mathsf{In}_{[D]}$, are left out of it), and what it leaves out has a mass $\le\mathsf R_{[D]}Z$. On $\mathsf{In}_{[D]}$, which lies in the field cut (Lemma~\YMMref{113.10}{796}), the measure of the chain has the density $p=e^{\mathsf c}e^{\Phi_t}$ for $\mathsf h_D$, $\mathsf c$ a complex number (Proposition~\YMMref{114.2}{801}), and there $|\Phi_t|\le\mathsf O_D:=(\ell L^D)^3[\frac32c_\rho^2A_0^2(\log x_{[D]})^{2\mathfrak p_0}+C_{\rm old}\frac L{L-1}E_0]$: the bound \YMMref{(2161)}{796} of the action on a field cut, at the radius $c_\rho\varepsilon(x_{[D]})$ and the weight $x_{[D]}$, gives the first term, because $\mathsf t\,\varepsilon(\mathsf t)^2=A_0^2(\log\mathsf t)^{2\mathfrak p_0}$, and the bound $\mathsf m^S_t$ \YMMref{(2179)}{803} of the added action (equation~\YMMref{(2180)}{803}) gives the second; $\mathsf O_D$ reads $D$ and not $K$.

A complex density of a positive measure is $\ge0$ almost everywhere; so $p=e^{{\rm Re}\,\mathsf c}e^{{\rm Re}\,\Phi_t}$ lies in $e^{{\rm Re}\,\mathsf c}[e^{-\mathsf O_D},e^{\mathsf O_D}]$, and its mass, which lies in $[(1-\mathsf R_{[D]})Z,Z]$, puts $e^{{\rm Re}\,\mathsf c}/Z$ between $(1-\mathsf R_{[D]})e^{-\mathsf O_D}/\mathsf h_D(\mathsf{In}_{[D]})$ and $e^{\mathsf O_D}/\mathsf h_D(\mathsf{In}_{[D]})$: the constant $\mathsf c$, the sum of the vacuum energies of the steps, drops out. For a gauge-invariant Borel $F$ with $0\le F\le1$ on the fields of the depth $D$, $F\circ\mathcal M^t$ is gauge invariant, the block map being covariant, and this gives
\[
\begin{gathered}
\kappa_D\int_{\mathsf{In}_{[D]}}F\,d\mathsf h_D\ \le\ \int F\,d\mu^K_D\ \le\ \bar\kappa_D\int_{\mathsf{In}_{[D]}}F\,d\mathsf h_D+\mathsf R_{[D]},\qquad\\
\kappa_D=\frac{(1-\mathsf R_{[D]})e^{-2\mathsf O_D}}{\mathsf h_D(\mathsf{In}_{[D]})},\quad\\
\bar\kappa_D=\frac{e^{2\mathsf O_D}}{\mathsf h_D(\mathsf{In}_{[D]})},
\end{gathered}
\]
the right inequality by \YMMref{(2189)}{806} at $1-F$ together with \YMMref{(2190)}{806}; at $F=\mathbf 1_{\mathsf{In}_{[D]}}$, $\mu^K_D(\mathsf{In}_{[D]})\ge1-\mathsf R_{[D]}$. For a Borel set $E$, $\mu^K_D(E)=\int\bar{\mathbf 1}_E\,d\mu^K_D$ with the mean $\bar{\mathbf 1}_E$ of $\mathbf 1_E$ over the gauge group, the law being gauge invariant, while $\mathsf{In}_{[D]}$ and $\mathsf h_D$ are gauge invariant; so the floor $\mu^K_D(E)\ge\kappa_D\mathsf h_D(E\cap\mathsf{In}_{[D]})$ and the upper bound $\mu^K_D(E)\le\bar\kappa_D\mathsf h_D(E\cap\mathsf{In}_{[D]})+\mathsf R_{[D]}$ hold for all Borel $E$, free of $K$, and for the limit law, the laws converging on every Borel set (Corollary~\textup{\ref{IV-62-66}}).

\emph{The thick vortex.} On the bonds $c$ of the depth $D$ let $n(c)=\#(c\cap\mathsf S)$, let $j_0$ be an integer cycle carried by the dual bonds of the rim with $j_0\equiv dn$ modulo $2$ (the rim, with an even number of its dual bonds at every dual site, is a union of closed dual paths, which are oriented), and $\mathsf r=\lfloor L^{D-d}/4\rfloor$. The translates $\tau_vj_0$ by the $(2\mathsf r+1)^3$ lattice vectors $v$ of a box of radius $\mathsf rL^{-D}$ differ from $j_0$ by coboundaries $d\beta_v$, $\beta_v$ carried within the distance $<L^{-d}$ from the rim: with the prism $(\mathsf{pr}_\lambda j)(x;\mu)=j(x;\lambda\mu)$, which takes functions on plaquettes to functions on bonds and functions on cubes to functions on plaquettes, $d\,\mathsf{pr}_\lambda j+\mathsf{pr}_\lambda dj=\tau_\lambda j-j$ and $dj_0=0$, summed along a path to $v$.

With $\varphi$ the mean of the $\pi\beta_v$, $\omega$ the mean of the $\pi\tau_vj_0$, $\theta=\pi n+\varphi$ and $V^1_c=e^{i\theta(c)\sigma_3}$: $d\theta=\omega+2\pi\mathsf k$ with $\mathsf k$ integer, so $\partial V^1(p)=e^{i\omega_p\sigma_3}$, and $|\omega_p|\le\pi n_{\mathsf S}(2\mathsf r+1)^{-2}\le16\pi n_{\mathsf S}L^{-2(D-d)}$, as a segment of the rim normal to $p$ meets at most $2\mathsf r+1$ of the translates of $p$; the rim misses $\mathcal R$, so it keeps the distance $\ge L^{-d}$ from it (the closed cubes of side $L^{-d}$ about sites that touch a dual bond of $\Lambda_{d,\ell}$ in the rim cover its neighbourhood of that radius, and none of them belongs to $\mathcal R$), and on the bonds of $\mathcal R$, $\varphi=0$ and $V^1=\bar\xi$.

Counted in $\Lambda_{d,\ell}$, $n_{\mathsf S}$ is at most the number $3(\ell L^d)^3$ of all its dual bonds; so for $D$ so large that $48\pi(\ell L^d)^3L^{-2(D-d)}\le\varepsilon(x_{[D]})$, which holds for large $D$ as $\varepsilon(x_{[D]})$ decays only like $L^{-D/2}$ up to logarithms, the thick vortex of every flip has $\varpi(V^1)\le\varepsilon(x_{[D]})$: it is a small field, and one $D$ serves all surfaces. \emph{The common part.} Let $\eta=\frac14\varepsilon(x_{[D]})$ and ${\rm vol}(\eta)$ the Haar measure of $\{g\in G:|g-\mathbf 1|<\eta\}$. If every bond of a field $W$ is within $\eta$ of $\mathbf 1$, then $\varpi(W)<4\eta<2\varepsilon(x_{[D]})$ and $\varpi(V^1W)<\varpi(V^1)+4\eta\le2\varepsilon(x_{[D]})$, by the bound of a product of four: $W$ and $V^1W$ (bondwise product) are small fields, and $(V^1W)|_{\mathcal R}=\bar\xi\,W|_{\mathcal R}$.

By the floor, the law $\mu_{\mathcal R}$ of the block fields of $\mathcal R$ at the depth $D$ dominates the measure $\kappa_D\mathsf g(w)\,dw$, $\mathsf g(w)$ the Haar measure of the continuations of $w$ to small fields; on the configurations $w$ of $\mathcal R$ within $\eta$ of $\mathbf 1$, $\mathsf g(w)$ and $\mathsf g(\bar\xi w)$ are $\ge c_\eta={\rm vol}(\eta)^{3(\ell L^D)^3-n_{\mathcal R}}$, $n_{\mathcal R}$ the number of the bonds of $\mathcal R$ at the depth $D$ (continue $w$ by bonds within $\eta$ of $\mathbf 1$, and $\bar\xi w$ by the same bonds times those of $V^1$); so $\theta_0=\kappa_Dc_\eta\mathbf 1_{\{|w_c-\mathbf 1|<\eta\ \forall c\}}\,dw$ lies below $\mu_{\mathcal R}$ and below $\bar\xi_*\mu_{\mathcal R}$, for every flip of a surface whose rim misses $\mathcal R$ and at every cutoff $K\ge D+1$, and in the limit; it reads no surface and no cutoff, and its mass $\varkappa_D=\kappa_D{\rm vol}(\eta)^{3(\ell L^D)^3}$ reads $D$ alone.

The block map from the depth $D$ to the depth $d$, which reads on the bonds of $\mathcal R$ the bonds of $\mathcal R$ alone and turns $\bar\xi$ at the depth $D$ into $\bar\xi$ at the depth $d$ by the first clause of (g), carries this down. Then $\mu_{\mathcal R}\wedge\bar\xi_*\mu_{\mathcal R}$ has mass $\ge\varkappa_D$, which is the bound in total variation; for odd $F$ with values in $[-1,1]$, $\int F\,d\mu_{\mathcal R}=\frac12\int F\,d(\mu_{\mathcal R}-\bar\xi_*\mu_{\mathcal R})$ lies in $[\varkappa_D-1,1-\varkappa_D]$; and for flip-invariant $\Phi$, since $\mu\ge\frac12(\theta_0+\bar\xi_*\theta_0)$ and $W_C\circ\bar\xi=-W_C$,
$$\int|W_C-\Phi|\,d\mu\ge\frac12\int\bigl(|W_C-\Phi|+|W_C+\Phi|\bigr)\,d\theta_0\ge\int|W_C|\,d\theta_0=\mathsf m_C>0,$$
$|W_C|$ being near $1$ on the configurations near $1$; the same for ${\rm sign}\,W_C$. \emph{Densities.} If $\mathsf h_d(E)=0$, the preimage of $E$ under the block map from the depth $D$ is Haar-null by Lemma~\ref{C-363-26}, and the upper bound found with the floor gives $\mu_d(E)\le\mathsf R_{[D]}$ for every large $D$, so $\mu_d(E)=0$. \emph{The ceiling.} On the open set $X_0$ of full Haar measure of Lemma~\ref{C-363-24}(b), which has full $\mu$-measure, the four bricks with their overlaps determine a configuration up to a gauge transformation and one flip; Lemma~\ref{C-363-24}(b) with $\mathsf H=\mathbb Z_2$, $C_\kappa=C$ for the non-trivial character, gives $F=\Phi_0+\Phi_1u$ with $\Phi_i$ measurable for the bricks; for bricks three sites thick the cubes generate the bricks by the same lemma; the block field of a coarser depth in $\mathcal R$ is a function of that of the depth $d$ in $\mathcal R$.

Lemma~\ref{C-363-24}(c) gives the adjoint loop. \emph{The commutant.} $\mathfrak L$ is maximal abelian on $\mathfrak h$, the space of the gauge-invariant classes of $L^2(\mu)$, the density $\rho$ being chosen gauge invariant; $\mathfrak E$ is its part fixed by the flip $\alpha$; on $Y=\{\rho>0\}\cap\alpha\{\rho>0\}$ the flip is implemented by $v$, a partial isometry with $v^2=e$ and $vM_Fv=M_{F\circ\alpha}e$, and on the rest of $\{\rho>0\}$ the flip moves no mass inside; so $\mathfrak E'=\mathfrak L+\mathfrak Lv$, isomorphic to $L^\infty(Y_1)\oplus L^\infty(Y\cap\{u=1\})\otimes M_2(\mathbb C)$, $Y_1=\{\rho>0\}\setminus Y$; $vu=-uv$, and $v\ne0$ because $\mu(Y)\ge\varkappa_D$: the measure $\mu\wedge\alpha_*\mu$, of the density $\min\{\rho,\rho\circ\alpha\}$, vanishes off $Y$ and lies above $\theta_0$.

For this last point $\alpha$ is realised as the flip of a surface whose rim misses $\mathcal R$, at the depth $d+1$. In the units of $\Lambda_{d,\ell}$ let $y_0=(b_1+1,b_2+1)$, a site of the hole ($b_i'-b_i\ge2$), and $y_1=(a_1'+1,b_2+1)$, a site beyond $Q$ ($a_1'-a_1\le N-2$); let $\mathsf S_1$ be the flat sheet of $\Lambda_{d+1,\ell}$ in the plane $x_2=b_2+1+\frac1{2L}$ from the straight closed dual line of $\Lambda_{d+1,\ell}$ in the direction $3$ over $y_0+\frac1{2L}(1,1)$ to the one over $y_1+\frac1{2L}(1,1)$, across the arm of the ring at $x_1=b_1'$. Each line lies over the open square of side $1$ about $y_0$ or $y_1$, which meets no closed cube of $\mathcal R$: the rim misses $\mathcal R$; and $C$ crosses $\mathsf S_1$ once. $\mathcal R$ is a union of closed cubes of side $L^{-(d+1)}$ about sites of $\Lambda_{d+1,\ell}$, $L$ being odd; so the first half of (i) and (g) hold for $\mathsf S_1$ at the depth $d+1$.

The block map from the depth $d+1$ to the depth $d$ reads on the bonds of $\mathcal R$ the bonds of $\mathcal R$ alone, and by the first clause of (g), no plaquette in $\mathcal R$ being pierced by the rim, it turns the flip of $\mathsf S_1$ into the sign field $(-1)^{\#(c\cap\mathsf S_1)}$ on the bonds $c$ of $\mathcal R$ at the depth $d$; this sign field is closed on the plaquettes of the bricks and odd on $C$, so it differs from $\alpha$ by a coboundary, which acts trivially on the gauge-invariant functions. Hence $\theta_0$, carried down from the depth $d+1$, lies below $\mu$ and below $\alpha_*\mu$, and its mass is $\varkappa_D$. This holds for every ring admitted, $D$ and $\varkappa_D$ being those of the first half of (i) at the depth $d+1$.

If besides $b_i'-b_i\ge3$ $(i=1,2)$ and $a_i'-a_i\le N-3$ for one $i$, say $i=1$, a sheet of $\Lambda_{d,\ell}$ itself serves as well, with $D$ and $\varkappa_D$ at the depth $d$: in the units of $\Lambda_{d,\ell}$, the flat strip $\{x_2=b_2+\frac32,\ b_1+\frac32\le x_1\le a_1'+\frac32\}$ between the straight closed dual lines in the direction $3$ over $(b_1+\frac32,b_2+\frac32)$ and over $(a_1'+\frac32,b_2+\frac32)$. A closed cube of side $1$ about a site meets such a line only if the site lies in one of the four columns of sites nearest to the line, and these columns hold no site of $\mathcal R$: about the first line they lie over $b_1+1,b_1+2$ and $b_2+1,b_2+2$, in the hole, as $b_i+2<b_i'$; about the second they lie over $a_1'+1,a_1'+2$ in the direction $1$, beyond $Q$, as $a_1'+2<a_1+N$.

So the four cubes about each line, at every height, are free, and the rim of the strip misses $\mathcal R$; $C$ crosses the strip once, on its side $x_1=b_1'$, as $b_1<b_1+\frac32<b_1'<a_1'+\frac32$ and $b_2<b_2+\frac32<b_2'$; the flip of the strip is closed on the plaquettes of the bricks and odd on $C$, so it differs from $\alpha$ by a coboundary, and the first half of (i) at the depth $d$ puts the measure $\theta_0$ of that depth below $\mu$ and below $\alpha_*\mu$. \emph{The vortex line.} At the cutoff $K$ the flip of the block fields of $\mathcal R$ is the insertion of the closed vortex line $\partial\mathsf S_1$: (g) at the depth $d+1$, for the functions of the block fields of $\mathcal R$ at the depth $d+1$ that factor through the block map, together with the covariance just used, gives $\mu^K\circ\alpha=x_K\mu^K$ with $x_K=\mathbb E_K[X_K\,|\,V]$, $X_K$ the ratio of the twisted to the untwisted weight, at every $K$; so $x_K\cdot x_K\circ\alpha=1$, $v_KF=x_K^{1/2}F\circ\alpha$ is a self-adjoint unitary, and, with $\rho_K$ the density of $\mu^K$ for Haar measure (Lemma~\ref{C-363-26}), $|\langle F,v_KF'\rangle-\langle F,vF'\rangle|\le\sup|F|\sup|F'|\int|\sqrt{\rho_K\,\rho_K\circ\alpha}-\sqrt{\rho\,\rho\circ\alpha}|\,d\mathsf h$, which the Schwarz inequality bounds by $2\sqrt2\sup|F|\sup|F'|$ times the square root of the total variation between $\mu^K$ and $\mu$, at most $\sum_{K'\ge K}\mathfrak o_{K'}$.

\emph{Wide rings.} Let moreover the hole be at least three bonds wide in both directions, and the ring be of positive height and lie in an open ball of radius $<\ell/4$. Then a sheet of $\Lambda_{d,\ell}$ itself realises $\alpha$: in the units of $\Lambda_{d,\ell}$, the flat sheet $\mathsf S$ in the plane $x_2=b_2+\frac32$ from the straight closed dual line in the direction $3$ over the centre $(b_1+\frac32,b_2+\frac32)$ of four cubes of the hole (here the width three is used) to the one over the centre of four cubes beyond the ring on the same row (for which the ball leaves room) has its rim off $\mathcal R$ and crosses $C$ once. Let $\xi^{(D)}$ be its flip at the depth $D$, $\pi_D$ the limit law of the block fields of $\mathcal R$ at the depth $D$, so that $\pi_d=\mu$, and $\Phi_D$ the block map from these fields to those at the depth $d$; by the first clause of (g), $\Phi_D\circ\xi^{(D)}=\bar\xi\circ\Phi_D$ with the flip $\bar\xi$ of $\mathsf S$ at the depth $d$, which differs from $\alpha$ by a coboundary; as $\mu$ is invariant under gauge transformations, we take $\alpha=\bar\xi$ here.

We show that $\mu(\alpha A)=0$ gives $\mu(A)=0$ for a Borel set $A$ of the block fields of $\mathcal R$ at the depth $d$.

(1) \emph{From the measure to the fields.} Let $\mathcal G_D$ be the set of the $V\in\mathsf{In}_{[D]}$ whose flipped restriction $\xi^{(D)}V|_{\mathcal R}$ continues to a field of $\mathsf{In}_{[D]}$. If $\mu(\alpha A)=0$, then $\mu(A)\le1-\mu_D(\mathcal G_D)+\mathsf R_{[D]}$. Indeed, with $\tilde A=\Phi_D^{-1}A$, $\pi_D(\xi^{(D)}\tilde A)=\mu(\alpha A)=0$; by the floor, $\int_{\xi^{(D)}\tilde A}q_D\,dw=0$, $q_D(w)$ the Haar measure of the continuations of $w$ into $\mathsf{In}_{[D]}$, and $q_D>0$ exactly on the projection $\mathsf{Pr}_D$ of the open set $\mathsf{In}_{[D]}$ to the bonds of $\mathcal R$; so $\tilde A\cap\xi^{(D)}\mathsf{Pr}_D$ is Haar-null. The fields of $\mathcal G_D$ whose restriction to $\mathcal R$ lies in $\tilde A$ have it in $\tilde A\cap\xi^{(D)}\mathsf{Pr}_D$; they form a Haar-null set, of $\mu_D$-measure $\le\mathsf R_{[D]}$ by the upper bound; the others with the restriction in $\tilde A$ lie off $\mathcal G_D$.

(2) \emph{Insertion of a straight line into a small field.} In units of $L^{-D}$ let $\mathsf l_P$ be the straight closed dual line in the direction $3$ over a dual point $P$ of the $(1,2)$-plane, $\varrho_1\ge1$ an integer with $2\varrho_1+4<\ell L^D$, $\mathcal T$ the sites whose horizontal part lies within the sup-distance $\varrho_1+\frac32$ of $P$ (the tube), $s$ a sign field with $(ds)_p=-1$ exactly at the plaquettes with corners in $\mathcal T$ that $\mathsf l_P$ pierces, and $V$ a field with $|\partial V(p)-\mathbf 1|\le\delta_0$ at the plaquettes with corners in $\mathcal T$. Then there is a field $V'$ with $V'_c=s_cV_c$ on every bond that has not both ends in $\mathcal T$, $|\partial V'(p)-\mathbf 1|\le6\varrho_1^{-2}+3\delta_0$ at the plaquettes with all corners in $\mathcal T$, and $\partial V'(p)=(ds)_p\partial V(p)$ at all the others.

Idea: take an axial gauge $\gamma$ of $\mathcal T$, $\mathbf 1$ on the spanning tree made of all the bonds in the direction $3$ at a corner $o$ of the unit plaquette centred at $P$ but one, one line in the direction $1$ through $o$ in each layer, and all the bonds in the direction $2$, the holonomy $\Upsilon$ around the torus along the bonds in the direction $3$ at $o$ being made diagonal by a constant gauge transformation; then, $x$ measured from $o$, a bond in the direction $1$ at $x$ is within $|x_2|\delta_0$ of $\mathbf 1$, and a bond in the direction $3$ at $x$ is $\mathsf Y_x\Upsilon_x$ with $\Upsilon_x\in\{\mathbf 1,\Upsilon\}$ and $\mathsf Y_x$ within $|x|_1\delta_0$ of $\mathbf 1$. For positively oriented bonds put $V'_c=s_c\gamma(c_-)^{-1}e^{-i\varsigma_c\sigma_3}\gamma(c_-)V_c$ on the bonds with both ends in $\mathcal T$, with the profile $\varsigma_c=(1-f_c)\vartheta_c$ on the horizontal bonds, $\vartheta_c$ half the signed angle that $c$ subtends at $P$, $f_c=\min\{1,(r_c^2-\frac14)/\varrho_1^2\}$, $r_c$ the sup-distance of the midpoint of $c$ from $P$, and $\varsigma_c=0$ on the vertical bonds; define the values on reversed bonds by inversion.

Then $d\varsigma=\pi$ at the pierced plaquettes, which takes away the sign of $ds$ there, and $|d\varsigma|\le6\varrho_1^{-2}$ at the other horizontal ones, as $d\vartheta=0$ there: three terms $\pm(f_{c'}-f_c)\vartheta_c$ remain, each $\le\frac{2r}{\varrho_1^2}\cdot\frac1r$, $r\ge1$ the sup-distance of the centre of the plaquette from $P$; this is the source of the $6$. The rotations $e^{-i\varsigma_c\sigma_3}$ commute with one another and with $\Upsilon$. With $o=(0,0)$ and $P=(\frac12,\frac12)$ in a horizontal layer, the half-angle bound is $|\vartheta_c|\,|c_--o|_1\le1$ for the positively oriented horizontal bonds. It is paired with the gauge error measured from the same $o$. For a horizontal plaquette the direction-$2$ bonds are $\mathbf1$, leaving one commutator with a direction-$1$ bond. For a vertical plaquette first factor out the original plaquette product; the remaining error is one commutator with $\mathsf Y_x$ in $\mathsf Y_x\Upsilon_x$, since the inserted rotations commute with $\Upsilon_x$. Its cost is at most $2\delta_0$ in either case. Adding the original plaquette error $\delta_0$ gives $3\delta_0$. At the boundary $f_c=1$, so the insertion matches the prescribed exterior field; no boundary value of the gauge at $o$ is required.

(3) \emph{The witness.} Let $\varrho_1=\frac12(L^{D-d}-5)$. At the depth $D$ the hole holds a square of $(2\varrho_1+4)^2$ sites about a dual point $P_1$ of the $(1,2)$-plane, and beyond the ring, on its row, there is a dual point $P_2$ whose square of that size misses $Q$ (the ball leaves room); the tubes $\mathcal T_1,\mathcal T_2$ over them are disjoint and hold no site of $\mathcal R$. Let $s=-1$ on the bonds in the direction $2$ that cross the flat strip between the vertical lines over $P_1$ and $P_2$, and $s=1$ elsewhere: $(ds)_p=-1$ exactly at the plaquettes that the two lines pierce. On the bonds of $\mathcal R$ at the depth $D$, $s$ and $\xi^{(D)}$ are closed on the plaquettes of the bricks and odd on the loop around the hole, so $s=\xi^{(D)}dz$ there, with signs $z$ on the sites of $\mathcal R$ (Lemma~\ref{C-363-24}(b) for the four bricks, whose group $Z/B$ is $\mathbb Z_2$ and is read on that loop).

Let $\theta'<\frac23$ and $D$ be so large that $6\varrho_1^{-2}+3\theta'\varepsilon(x_{[D]})<2\varepsilon(x_{[D]})$, as $\varrho_1^{-2}$ is of the order of $L^{-2(D-d)}$ and $\varepsilon(x_{[D]})$ of the order of $L^{-D/2}$ up to logarithms. If $\varpi(V)<\theta'\varepsilon(x_{[D]})$, $V$ is a small field, and the insertion in both tubes with $\delta_0=\theta'\varepsilon(x_{[D]})$ gives $V'$ with $V'=sV$ off the tubes, so on the bonds of $\mathcal R$, and $\varpi(V')<2\varepsilon(x_{[D]})$; the gauge transform of $V'$ by $z$, continued by $1$, is a small field equal to $\xi^{(D)}V$ on the bonds of $\mathcal R$. So $\{\varpi<\theta'\varepsilon(x_{[D]})\}\subset\mathcal G_D$.

(4) \emph{The tail at a fraction of the threshold.} For integers $n\ge0$ and $D'=D-n$ with $x_{[D']}\ge\mathsf t_\natural$, let $\Theta_n=4B_3L^{-3n/2}+2\frac{C_HA_1}{A_0}\sum_{j<n}L^{-3j/2}$, where $B_3\ge1$ and $C_H\ge1$ are the constants of the list of Theorem~\YMMref{D}{8} that enter the bound of the background over a small field and the contraction of the backgrounds from one level to the next (Theorem~\YMMref{E}{10} and the curvature comparison in the proof of Lemma~\YMMref{113.10}{796}). Then $\mu^K_D\{\varpi\ge\Theta_n\varepsilon(x_{[D]})\}\le\mathsf R_{[D']}$ for $K\ge D+1$, and in the limit. Run the chain with its top level at the depth $D'$. On the histories with $\mathfrak w^{\rm sf}\ne0$ the field at the depth $D'$ lies in the inner ball, and the field of every finer level lies in the field cut and is the field over the next coarser one moved by a fluctuation from the ball of the step; so the regularity numbers $\mathsf m$ of the backgrounds obey $\mathsf m<2B_3\varepsilon(x_{[D']})$ at the top (Theorem~\YMMref{E}{10}) and $\mathsf m_k\le L^{-2}\mathsf m_{k+1}+C_H\delta(x_k)$ from the level $k+1$ to the finer level $k$ (the proof of Lemma~\YMMref{113.10}{796}).

Summed down the $n$ levels, with $\varepsilon(x_{[D-j]})\le L^{j/2}\varepsilon(x_{[D]})$ and the same for $\delta$, this gives $\mathsf m<\frac12\Theta_n\varepsilon(x_{[D]})$ at the depth $D$, and a field whose background has the regularity number below $\mathsf m$ has its plaquettes below $2\mathsf m$ (as in the inclusions \YMMref{(2159)}{796}): $\varpi<\Theta_n\varepsilon(x_{[D]})$. The other histories have the mass $\le\mathsf R_{[D']}Z$ by \YMMref{(2189)}{806} at $1-F$ and \YMMref{(2190)}{806}, $F$ the indicator of $\{\varpi(\mathcal M^{K-D}U)\ge\Theta_n\varepsilon(x_{[D]})\}$; the limit follows.

(5) \emph{The end.} Let $n_0$ be the least integer with $L^{3n_0/2}\ge80B_3$; it is fixed by $L$ and $B_3$, and chosen so that $\Theta_{n_0}<\frac35$: the first term of $\Theta_{n_0}$ is $\le\frac1{20}$, and the condition $A_0\ge4B_3C_HA_1$, one of the conditions on $A_0$ written in the proof of Lemma~\ref{lem:rarity-lowered-choice}, with $B_3\ge1$ and $L\ge13$, makes the second $\le\frac12(1-13^{-3/2})^{-1}<0.52$. For $D$ so large that $x_{[D-n_0]}\ge\mathsf t_\natural$, $\mathsf R_{[D]}<1$ and $6\varrho_1^{-2}<\frac15\varepsilon(x_{[D]})$, (1) and (3) at $\theta'=\frac35$ give $\mu(A)\le\mu_D\{\varpi\ge\frac35\varepsilon(x_{[D]})\}+\mathsf R_{[D]}$, and (4) at $n=n_0$ bounds the first term by $\mathsf R_{[D-n_0]}$. So $\mathsf R_{[D-n_0]}$ bounds the mass of the histories that the chain with its top level at the depth $D-n_0$ leaves out, this chain controlling the plaquettes at the depth $D$, and $\mathsf R_{[D]}$ comes from the upper bound in (1). Both tend to $0$ as $D\to\infty$, and $\mu(A)$ reads no $D$: $\mu(A)=0$. With $\alpha A$ in the place of $A$ the converse holds; so the law and its flip have the same null sets, $e=1$, $v$ is a self-adjoint unitary, and with the ceiling, whose hypothesis on the corners the positive height meets, $\mathfrak L=\mathfrak E\oplus\mathfrak Eu$ is a grading by $\mathbb Z_2$.

\emph{A band that winds.} The flip $\bar\xi_1$ of a band $\mathbb Z/N\times I_2\times I_3$ is the restriction of a flip of the torus, a symmetry of the law: $v$ is unitary, odd and even functions are orthogonal, $\mathfrak Eu_1\ne0$ as $u_1^2=1$, and the cubes of the band generate $\mathfrak E$ by Lemma~\ref{C-363-24}(b) with the cubes as boxes; the only closed sign fields on the cubes of the band modulo coboundaries are $1$ and the winding one.

\subsection{Scope of the conclusions}

Theorem~\ref{C-363-4}(a)--(g) concern the centred functionals of volume-first limits and their components, for which the proof uses a periodic lattice; for other limits of periodic Gibbs states only (h) is asserted. The proofs of Theorems~\ref{C-363-3} and~\ref{C-363-4} use no energy bound beyond the polynomial one of the reconstruction, no gap, no uniqueness of the vacuum of $\sigma$ and no rate of clustering. Theorem~\ref{C-363-5}(d) bounds from below the norms of the field at one instant and of one Euclidean slab. Theorem~\ref{C-363-7}(g)--(i) are statements on the commutative algebras of the Euclidean block fields at a fixed depth on one torus.

\section{Continuum limits of the lattice Langevin dynamics}\label{C-367-1}

At a fixed physical label and on a fixed torus, the Wilson Langevin
processes have subsequential limits on threads of block fields. We prove
that each limit is a nonconstant stationary reversible Markov process on
the gauge-invariant $\sigma$-algebra. The lattice Dirichlet forms converge
in the strong and weak sense. Conditional Poincar\'e and logarithmic
Sobolev inequalities control the fine field; a connected coarse-scale
region then gives weak inequalities with defects tending to zero. These
estimates imply mixing and entropy relaxation. Their constants may
depend on the fixed torus, and the result does not identify the processes
obtained along different subsequences.

We also study observables regularised by the Wilson gradient flow. The
finite-cutoff flow identities and upper bounds use no probabilistic
assumption. The moment estimates use the admissibility conclusion
\textup{(Adm)} of Proposition~\ref{C-367-12}, concerning conditional layer means and colour covariances.
The identified continuum flow limits also assume \textup{(Flow approximation)}.
The nonlinear defect bounds are proved in Lemma~\ref{C-367-14}.
Lemma~\ref{lem:short-time-lower-propagation} proves lower propagation
on a regular time interval; Proposition~\ref{C-367-23}(c) applies it
on the interval of the short-time flow theorem. The extension beyond
that interval is the separate input \textup{(Lower propagation)}.
Neither this extension nor \textup{(Flow approximation)} is inferred
from \textup{(E)}.

The generator requires separate treatment. Some bounded flowed
observables have generator images whose norms grow exponentially in the
inverse square of the spacing. We therefore first insert condenser
cutoffs, pass to the continuum limit, and then remove the cutoffs.
The resulting generator identity is distributional.

\subsection{The lattice process and its limiting objects}

$G=\mathrm{SU}(2)$; $\mathfrak g$ is the space of hermitian traceless $2\times2$ matrices, $\langle X,Y\rangle=\frac12\mathrm{tr}\,XY$, $|X|=\langle X,X\rangle^{1/2}$, so that $X^2=|X|^2\mathbf 1$ and $e^{iX}\in G$; $\mathsf d$ is the geodesic distance of $G$, the unit sphere $S^3$, with $\mathsf d(e^{iA},\mathbf 1)=|A|$ for $|A|\le\pi$ and $|g-h|=2\sin\frac12\mathsf d(g,h)$. The numbers $L$, $a_d=L^{-d}$, $\ell=2L^m$, the torus $T_\ell=\mathbb R^3/\ell\mathbb Z^3$ of side $\ell$, the lattices $\Lambda_{d,\ell}=a_d\mathbb Z^3/\ell\mathbb Z^3\subset T_\ell$ (Section~\YMMref{3}{13}) and the averaging map $\mathcal M$ (equations~\YMMref{(37)--(39)}{20}) have the conventions just stated. The label $\mathsf b$ and setting (E) are those of Section~\ref{sec:cutoff-calibration}: the constants of the construction are an admissible choice $\varsigma=(\sigma^-,\gamma)$ of the canonical list of constants and conditions (Section~\ref{sec:cutoff-calibration}), subject to the conditions (Th1) and (Th2) of Section~\ref{sec:cutoff-calibration}, with $\gamma$ the fixed reference-coupling ceiling; the label satisfies $\mathsf b\ge\gamma^{-2}$ and $M\,R(\mathsf b)\le L^m$, with the cube size $M=L^{m'}$ and the function $R$ of Section~\ref{sec:cutoff-calibration}, $R(x)$ the least power $L^j$, $j\ge0$, that is not smaller than a fixed power of $\log x$; and $\beta=L^K\mathsf b$, so that $\mathsf b=a_K\beta$.

$\Omega_d$ is the set of configurations on $\Lambda_{d,\ell}$, $dV$ (on $\Omega_K$: $dU$) its Haar probability measure, $V^u(b)=u(b_-)V(b)u(b_+)^{-1}$ the gauge action; invariant means gauge invariant. $\partial V(p)$ is the plaquette variable, $a_p=1-\frac12\mathrm{tr}\,\partial V(p)$, $\varpi_p=|\partial V(p)-\mathbf 1|=(2a_p)^{1/2}$, $\varpi(V)=\max_p\varpi_p$, $A_K=\sum_pa_p$, $A^a=a_K^{-1}A_K$, and $\mathcal O(f)=a_K^{-1}\sum_pf(x(p))\,a_p$ for real $f$ on $T_\ell$, $x(p)$ the barycentre of $p$. On $\Omega_K$ we write $N=\ell/a_K$ for the number of sites along an axis. Tangent vectors are real $\mathfrak g$-valued $1$-forms $X$, with the curves $s\mapsto e^{isX}V$; $\langle X,\partial F(V)\rangle=\frac d{ds}F(e^{isX}V)|_{s=0}$, $\|X\|^2=\sum_b|X(b)|^2$; $\partial_bF$ is the value of the $1$-form $\partial F$ at the bond $b$, and $\partial^\alpha_bF=\langle\sigma_\alpha,\partial_bF\rangle=\frac d{ds}F(e^{is\sigma_\alpha\delta_b}V)|_{s=0}$, with the Pauli matrices $\sigma_1,\sigma_2,\sigma_3$, a basis of $\mathfrak g$ that is orthonormal for $\langle\cdot,\cdot\rangle$, and $\delta_b$ the function of the bonds that is $1$ at $b$ and $0$ at the others; $\Gamma_d(F)=a_d^{-1}|\partial F|^2$ on $\Omega_d$, and $\Gamma_d(F,F')$ is its polarisation.

The vectors $\mathrm d_V\lambda$, $(\mathrm d_V\lambda)(b)=\lambda(b_-)-\mathrm{Ad}(V(b))\lambda(b_+)$ for $\mathfrak g$-valued site functions $\lambda$, are the gauge directions at $V$; a tangent vector orthogonal to them is horizontal, and the gradient of an invariant function is horizontal. $\mathrm{osc}=\sup-\inf$.

\emph{Dynamics.} $\mu^K=Z_K^{-1}e^{-\beta A_K}dU$; $\Delta_K=a_K^{-1}\sum_{b,\alpha}(\partial_b^\alpha)^2$; $\mathcal L_K=\Delta_K-a_K^{-1}\beta\sum_b\langle\partial_bA_K,\partial_b\,\cdot\,\rangle$; $\mathcal E_K(F)=\int\Gamma_K(F)\,d\mu^K$ on the closure $D(\mathcal E_K)$ of $C^\infty(\Omega_K)$, with $\mathcal E_K(F)=\infty$ off $D(\mathcal E_K)$; $P^K_t=e^{t\mathcal L_K}$, and $(U_t)_{t\in\mathbb R}$ is the stationary process. $J(U)=a_K^{-3}\partial A_K(U)$ is the current (\YMMref{9.2}{38}), $(J,J)_a=a_K^{-3}|\partial A_K|^2$; $\Phi_s$ is the gradient flow $\frac d{ds}U(b)=-ia_KJ(U)(b)U(b)$, $U_s=\Phi_sU$, $v=-a_K^{-2}\sum_b\langle\partial_bA_K,\partial_b\,\cdot\,\rangle$ its vector field, and $\Lambda_\sigma(U)=\int_0^\sigma a_K^{-2}\varpi(U_s)\,ds$.

\emph{Threads and limits.} $\mathcal Y=\prod_{d\ge0}\Omega_d$ is compact and metrisable; a thread is a point $Y=(Y_d)$ of $\mathcal Y$ with $Y_d=\mathcal MY_{d+1}$ for all $d$; $\pi_{K,d}=\mathcal M^{K-d}$. The laws $\mu^K_d$ of $\pi_{K,d}U$ under $\mu^K$ converge in total variation to measures $\mu_d$, with $\mu_{d+1}\circ\mathcal M^{-1}=\mu_d$ (Corollary~\textup{\ref{IV-62-66}}), and $\mu=\mu_{\ell,\mathsf b}$ is the measure on threads with the marginals $\mu_d$. A cylinder function is $\Psi(Y_d)$, $\Psi$ bounded Borel on $\Omega_d$, and $\Psi_K=\Psi\circ\pi_{K,d}$; $L^2_{\rm inv}$ is the closure in $L^2(\mu)$ of the invariant cylinder functions, which generate the invariant $\sigma$-algebra, and $\Pi$ is the orthogonal projection onto $L^2_{\rm inv}$.

Along $K_j\to\infty$, $H_j\in L^2(\mu^{K_j})$ converge \emph{strongly} to $F\in L^2(\mu)$ if there are cylinder functions $\Psi_m\to F$ in $L^2(\mu)$ with $\lim_m\limsup_j\|H_j-\Psi_{m,K_j}\|_{L^2(\mu^{K_j})}=0$, and \emph{weakly} if $\sup_j\|H_j\|<\infty$ and $\int H_j\Psi_{K_j}d\mu^{K_j}\to\int F\Psi\,d\mu$ for every cylinder function $\Psi$. A \emph{limit} is a sequence $K_j\to\infty$ along which all means $\mathbb E\prod_i\Psi_i(\pi_{K_j,d}U_{t_i})$ of bounded Borel $\Psi_i$ converge, with its limit process $(Y_t)$ on threads; $(F,Q_tF')=\mathbb EF(Y_0)F'(Y_t)$, $\mathcal E_\star(F)=\sup_{t>0}t^{-1}(F,(1-Q_t)F)$, $A_\star$ is the operator of $\mathcal E_\star$, and the limit is frozen if $\mathcal E_\star$ vanishes on $L^2_{\rm inv}$, that is, if $Q_t=1$ on $L^2_{\rm inv}$ for all $t$ ($\mathcal E_\star$ is finite only on $L^2_{\rm inv}$, Lemma~\ref{C-367-6}(b)).

The forms converge \emph{in the strong and weak sense} if (M1) $\mathcal E_\star(F)\le\liminf_j\mathcal E_{K_j}(H_j)$ whenever $H_j\to F$ weakly, and (M2) every $F$ has $H_j\to F$ strongly with $\mathcal E_{K_j}(H_j)\to\mathcal E_\star(F)$.

\emph{The chain.} We recall the following objects of the small-field chain on one torus (Sections~\YMMref{113}{786}, \YMMref{114}{799}, \YMMref{115}{805}), which rests on the results on one step of Proposition~\ref{prop:source-entry}. $\delta(\mathsf t)=A_1\mathsf t^{-1/2}(\log\mathsf t)^{\mathfrak p_0}$, $\varepsilon=\frac{A_0}{A_1}\delta$, $\mathfrak b(\mathsf t)=\mathsf t^{1/2}\delta(\mathsf t)$, the threshold functions of equation~\YMMref{(634)}{339}, with the constants $A_0,A_1>0$ and $\mathfrak p_0\in\mathbb N$ of the construction; $\varepsilon_n=\varepsilon(L^n\mathsf b)$; $\sum_d\delta(L^d\mathsf b)<\infty$. For an integer $n\ge1$ with $L^n\mathsf b\ge\mathsf t_\natural$, $\mathsf t_\natural\ge\gamma^{-2}$ the fixed calibration of Section~\ref{sec:cutoff-calibration}, a number formed from the constants alone, and for $K>n$, put $t=K-n$, so that $1\le t\le K-1$, as required of a top level in Section~\ref{sec:cutoff-calibration}: level $k$ is depth $d=K-k$, $V_k=\mathcal M^kU$, $x_k=L^d\mathsf b$, so that $x_k^{-1}=a_d/\mathsf b$.

The cuts $\mathsf{In}_k\subset\mathsf{Cut}_k\subset\Omega_d$ of the definitions preceding Lemma~\YMMref{113.10}{796} are open and invariant (Lemma~\YMMref{113.10}{796}); $\mathsf{Cut}_k\subset\{\varpi<\theta_d\}$ with $\theta_d$ below; $\mathsf I_n=\mathsf{In}_t=\{\varpi<2\varepsilon_n\}\subset\Omega_n$, and $\mathsf D_n$ is the connected component of $\mathbf 1$ in $\mathsf I_n$. $\mathsf R_n$ is the number \YMMref{(2187)}{806} at the level $t=K-n$, which depends on $n$, $\ell$ and $\mathsf b$ but not on $K$; it bounds the mass of the large fields (Proposition~\textup{\YMMref{115.2}{806}}, lost-mass bound), and by \YMMref{(2187)}{806} $\mathsf R_n\le C(\ell,\mathsf b)\,L^{3n}\exp[-c'A_0^2(\log L^n\mathsf b)^{2\mathfrak p_0}]$ with the rate $c'>0$ of Section~\ref{sec:cutoff-calibration} and $2\mathfrak p_0\ge2$, so that it tends to $0$ faster than every exponential in $n$.

The chain measure is $\sigma^K=Z_K^{-1}\mathfrak w^{\rm sf}_te^{-\beta A_K}dU$, with the constant $Z_K$ of $\mu^K$ and the weight of the small-field chain of level $t$, the function defined before Proposition~\YMMref{115.2}{806},
\[
\mathfrak w^{\rm sf}_t(U)=\mathbf 1_{\mathsf{In}_t}(V_t)\prod_{k=0}^{t-1}\bigl[\mathbf 1_{\mathsf{Cut}_k}\,\mathfrak G_k[x_k]\,\chi^{\rm sf}_{x_k}\bigr](V_k),
\]
where $\mathfrak G_k[x_k]$ is the gauge-fixing factor of the axial gauge on the blocks of the level $k+1$ and $\chi^{\rm sf}_{x_k}$ the fluctuation test of the step $k+1$ (Section~\ref{sec:sf-chart-interface}); the weight is a bounded Borel function with values in $[0,\infty)$ (Proposition~\textup{\YMMref{115.2}{806}}, proof and the definitions preceding Proposition~\YMMref{115.2}{806}). So $\sigma^K$ is a positive measure of mass in $[1-\mathsf R_n,1]$ (Lemma~\ref{C-367-7}(i)), normalised by the constant of $\mu^K$ and not by its own mass; and as its weight carries gauge-fixing factors, $\sigma^K$ is not gauge invariant, and its comparisons with $\mu^K$ are made on invariant functions. $\sigma^K_d$ is its image at depth $d$, $\sigma^K_n$ the top law.

For $0\le k<t$ and $W\in\mathsf{Cut}_{k+1}$, a background $U_+$ over $W$ is any element of the minimal orbit over $W$ with $\mathcal M^kU_+$ in the axial gauge, and $V_W(B)$ is the chart of the fibre $\{\mathcal MV=W\}$ (Section~\ref{sec:sf-chart-interface}), defined on the convex box $\{|B(b)|<\mathfrak b(x_k)\}$ of the bonds that are not pivot bonds, with $V_W(B)(b)=e^{ix_k^{-1/2}B(b)}(\mathcal M^kU_+)(b)$ there; it carries the box into $\mathsf{Cut}_k$ (Lemma~\YMMref{113.10}{796}). The fixed radii $0<\alpha_0<\alpha_1$ are two constants of the construction (Definition~\YMMref{49.1}{374} and Lemma~\YMMref{82.1}{665}, Lemma~\YMMref{82.1}{665}), free of $n$, at which the analyticity cores of the chain are taken (Definition~\YMMref{49.1}{374}); $\mathsf{Cut}_{k+1}$ lies in the class $\mathcal W_{k+1}(x_k;\alpha_0,\alpha_1)$ of the sum rule (Lemma~\YMMref{113.10}{796}).

One fluctuation integral $p_W$ is set at every $W\in\mathsf{Cut}_{k+1}$ by the chart integral $\mathsf I$ of \eqref{eq:sf-chart-integral} at $\mathsf t=x_k$ (Definition~\ref{def:sf-chart-data} and equation~\eqref{eq:sf-chart-integral}): it is the measure on the box with
\[
\int h(V_W(B))\,p_W(dB)=\frac{\mathsf I[hF_k](W)}{\mathsf I[F_k](W)}\qquad\text{for every bounded Borel }h\text{ on }\Omega_d ,
\]
$F_k$ the integrand of the chain at the level $k$ and at the source $0$ (Proposition~\YMMref{114.2}{801}); that is, $p_W$ has on the box the density proportional to $\exp[-x_kG_W(B)+\mathsf J_W(B)]\,F_k(V_W(B))$, with the gauge-fixing form $G_W$ and the exponent $\mathsf J_W$ of the Jacobians of the chart (Section~\ref{sec:sf-chart-interface}). The denominator vanishes at no $W\in\mathsf{Cut}_{k+1}$: on the class of the sum rule $\mathsf I[F_k]=e^{\mathsf c_{k+1}(0)}F_{k+1}\ne0$, $\mathsf c_{k+1}(0)$ the constant \YMMref{(2173)}{801} of the step (Proposition~\YMMref{114.2}{801}). $p_W$ depends on $K$, $k$ and $W$, and not on the top level; that it is a probability measure with a positive continuous density on the box, and that its images under $V_W$ form a regular conditional law of $V_k$ given $V_{k+1}=W$ under $\sigma^K$, is Lemma~\ref{C-367-7}(ii).

The tower of $H$ is $h_0=H$, $h_{k+1}(W)=\int h_k(V_W(B))\,p_W(dB)$ at every $W\in\mathsf{Cut}_{k+1}$. For a finite measure $\nu$, $\mathrm{Ent}_\nu(f)=\inf_{c>0}\int\Xi(f,c)\,d\nu$, $\Xi(x,c)=x\log\frac xc-x+c\ge0$.

\emph{Constants.} $\mathsf o_n=[2C_{\rm nl}\frac L{L-1}E_0+\frac32c_\rho^2A_0^2(\log L^n\mathsf b)^{2\mathfrak p_0}](\ell L^n)^3$, with the number $C_{\rm nl}=C'_{\rm A}+2C_{\mathrm E,0}(C_{\rm lg}+L^{4l_0})$, with $C'_{\rm A}$, $l_0$, $C_{\rm lg}$ and $C_{\mathrm E,0}$ defined exactly as in equation~\YMMref{(2163)}{798} and its following sentence, using $K_{\rm A}=K^\sharp_{\rm tw}$ and $a_{\rm A}=a^\sharp$. The proof following that equation gives the bound on the non-local action, that is, the action of a bare array at the background of a field in a cut (equation~\YMMref{(2180)}{803}). We also use the size $E_0>0$ of Section~\ref{sec:step-input-conventions} and the constant $c_\rho$ of the field cuts Section~\ref{sec:cutoff-calibration} (and $C_H$ defined in Section~\ref{sec:regular-background-complements}); $\theta_d=2c_\rho\varepsilon(L^d\mathsf b)$; $\gamma_{\rm r}\in(0,1]$ and $c_{\rm r}\ge1$ are the numbers of Theorem~\YMMref{U}{378}, entered through Section~\ref{sec:step-input-conventions}, which depend on $L$ alone: the floor and the ceiling of a real reference form of the Gaussian integral of one step, from which the form of the step itself differs, in the $\ell^2$ operator norm, by at most the defect $\frac12\mathsf u_{\rm G}(LM)^{-3}\le\frac12\gamma_{\rm r}$ (the completeness (St4) of the Gaussian structure of the step, Definition~\YMMref{47.2}{359} and Theorem~\YMMref{W}{423}, with $\mathsf u_{\rm G}=4C_{\rm Str}$ of Theorem~\YMMref{U}{378} and the proof of Proposition~\YMMref{113.4}{791}); so the form of the step lies between $\frac12\gamma_{\rm r}$ and $c_{\rm r}+\frac12\gamma_{\rm r}$; $\Lambda_+=c_{\rm r}+\gamma_{\rm r}$; $C_Q$, $c_Q$, $a_Q$ are the constants defined in Lemma~\ref{lem:plain-average-remainder}, which depend only on $L$; for a $1$-form $X$ on $\Lambda_{K-j,\ell}$, $m^j_y(X)=\sum_{c'}e^{-\delta_0d(y,c'_-)}|X(c')|$, $\delta_0>0$ the rate of decay in the bounds of the response, are the averages that enter these bounds, and $\mathsf B$, $\mathsf S$ are the constants with $|\mathcal H_V(\zeta)(b)|\le\mathsf B\,m^j_y(\zeta)$ for $b_-$ in the block of $y$ and $\sum_ym^j_y(X)^2\le\mathsf S^2\|X\|^2$ (Section~\ref{sec:background-step-complements}, Lemma~\ref{lem:response-current-zero}); $\mathsf M=\max\{1,\mathsf B\mathsf S\}$.

$C_{\rm P}=4C_{\rm ch}/\gamma_{\rm r}$ (Lemma~\ref{C-367-7}); $\epsilon^\flat_d=[1+(C_\flat+C_{\rm cov})\delta(L^d\mathsf b)]^2-1$ and $\Pi^\flat_n=\prod_{d>n}(1+\epsilon^\flat_d)$, finite and tending to $1$ (Proposition~\ref{C-367-9}); $C^\sigma=3\mathsf M^2\Pi^\flat_{n_{\rm a}}$. Thresholds: $n_2$ is the least integer $\ge1$ with $L^{n_2}\mathsf b\ge\mathsf t_\natural$ such that $C_{\rm Hs}\delta(L^d\mathsf b)\le\frac14\gamma_{\rm r}$ and $\theta_d\le\omega_L$ for $d\ge n_2$, with $C_{\rm Hs}$ the constant of Lemma~\ref{C-367-7}(iv) and a number $\omega_L\le(3200L^3)^{-1}$ that depends only on $L$; $n_3\ge n_2$ is the least integer such that for $d\ge n_3$ the number $\theta_d$ lies below the radius of Section~\ref{sec:background-step-complements} and $2\theta_d\le\varepsilon_{\mathcal W}\alpha_0$, with the fixed radii $\alpha_0<\alpha_1$ of the chain and the number $\varepsilon_{\mathcal W}$ of Proposition~\YMMref{52.3}{433}, so that $\theta_d$ is at most the threshold $\varepsilon_{\mathcal W}\cdot\frac12\alpha_0$ of Proposition~\YMMref{52.3}{433} at the radii $(\frac12\alpha_0,\frac12\alpha_1)$; $n_{\rm a}$ is the least integer $\ge\max\{1,n_3\}$ such that for all $d\ge n_{\rm a}$: $L^d\mathsf b\ge\mathsf t_\natural$; $\theta_d\le\min\{a_Q,\varpi_{\rm db}\}$ (that $\theta_d$ lies below the radius of Section~\ref{sec:background-step-complements} holds already for $d\ge n_3$); $4\delta(L^d\mathsf b)\le r_2$; $4C_{\rm db}\delta(L^d\mathsf b)\le r_1$; and $C_z\delta(L^d\mathsf b)\le\min\{\frac12r_{\rm h},\frac12,(36+6C_{\rm pz}^2)^{-1/2}\}$.

Here $\varpi_{\rm db}=r_{\mathsf y}^2/(8\pi L^3)$ and $C_{\rm db}=8L^3$ are the numbers of Definition~\ref{def:sf-chart-data}, $r_{\rm h}=\min\{r_{\mathrm B},\varrho^0_{\mathrm R}\}$ is the radius of holomorphy of the response in Section~\ref{sec:background-step-complements}, $r_1=\min\{\frac{c_Q}{4L},\frac{c_Q}{32C_QL^4}\}$, $r_2=\frac{r_1}{8C_QL^3}$, $C_z=4C_QL^3$ and $C_{\rm pz}=24C_QL^3+128C_QL^4/c_Q$, so that these numbers depend only on $L$ and the constants of Section~\ref{sec:background-step-complements}; $n''$ is the least integer $\ge n_{\rm a}$ with $\mathsf R_e\le\frac12$ for all $e\ge n''$; $n_1$ is the least $n\ge2$ with $L^{n-1}\mathsf b\ge\mathsf t_\natural$ and $4\theta_d\le(48000L^4)^{-1}$ for $d\ge n$.
\paragraph{An additional estimate for gauges at several block levels.} \emph{(Multiscale gauge)} There are finite constants $R_{\rm S},C_{\rm S}\ge1$, independent of the odd integer $L\ge13$, the side $\ell$, the integers $n,J$ and the field, with the following property. If $V\in\Omega_n$, $J\ge1$, $n-J\ge1$, $L^J\ge R_{\rm S}$, $\varpi(\mathcal M^jV)<(200L^2)^{-1}$ for $0\le j<J$, and $\frac\pi2L^J\varpi(\mathcal M^JV)\le1$, there is a periodic site gauge $u$ such that \[\mathsf d(V^u(b),\mathbf1)\le44\pi L^2\sum_{j=0}^{J-1}\varpi(\mathcal M^jV)+C_{\rm S}L^{-J}\quad\hbox{at every bond }b.\] This is an additional quantitative hypothesis. Only (B1)(e), (B2)(d) and the explicit-top variants of (A4)(f) and (A5) below assume it; the assertions using $j_1$, $w_n$ and $C^{\rm t}_n$ retain their stated hypotheses.
With $C_{\rm BG}=2\lambda_{{\rm BG},3}$ of \eqref{eq:uniform-box-gauge-constants}, $C''=1+20\pi^2C_{\rm BG}^2$, $\kappa_0=(8\pi C_{\rm BG})^{-2}$, $\vartheta_\infty=C_HA_1/((1-L^{-3/2})A_0)$, which is $<1$ under (E), and $j_1$ the least $j\ge1$ with $L^{-3j/2}c_\rho\le\min\{\frac12(1-\vartheta_\infty),1/C''\}$: $n_{\rm c}$ is the least integer such that for $n\ge n_{\rm c}$: $n-j_1\ge1$, $L^{n-j_1}\mathsf b\ge\mathsf t_\natural$, $2\varepsilon_n\le C''\kappa_0$, $\ell L^n\ge\max\{\lambda_{{\rm BG},3},(C''/2\varepsilon_n)^{1/2}\}$; $\mathsf r_n=\mathsf R_n+\mathsf R_{n-j_1}$. For the assertions assuming \textup{(Multiscale gauge)}, let $\alpha_n=\frac18(2\arcsin\varepsilon_n)^{1/2}$ for $n\ge2$ (the indices $0$ and $1$ belong to the fixed radii), and let $J_n$ be the least integer $J\ge1$ with $L^J\ge\max\{R_{\rm S},8C_{\rm S}/\alpha_n\}$. Then $J_n\le n/4+c$, with $c$ independent of $n$ and $K$.

Under \textup{(Multiscale gauge)}, $n_{\rm e}$ is the least integer $\ge2$ such that for $n\ge n_{\rm e}$: $\varepsilon_n\le\frac1{51}$, $\log L^n\mathsf b>2\mathfrak p_0$, $n-J_n\ge1$, $L^{n-J_n}\mathsf b\ge\mathsf t_\natural$, $\theta_d<(200L^2)^{-1}$ for $d\ge n-J_n$, $44\pi L^2\sum_{j<J_n}\theta_{n-j}\le\frac18\alpha_n$, and $\frac\pi2L^{J_n}\theta_{n-J_n}\le1$. No constant of this section depends on $K$; $n_1,n_2,n_3,n_{\rm a},n_{\rm e},C^\sigma,\Pi^\flat_n$ depend only on $\mathsf b$ and the constants of the construction, and $\mathsf R_n$, $\mathsf o_n$, $n''$, $n_{\rm c}$ and the constants defined from them ($\mathsf r_n$, $C^{\rm t}_n$, $C^\flat_n$, $\mathsf C_n$, $\mathsf C'_n$) depend besides on $\ell$.

\emph{Flowed observables.} Fix $0<\kappa<\frac16$. For $s>0$ and $V\in\Omega_K$, $\mathsf m_s(V)$ is the least 
\[
\max\{s^{1/4+\kappa/2}\|A\|_\infty,s^{3/4+\kappa/2}\|\nabla A\|_\infty\}
\]
 over the real $1$-forms $A$ with $e^{ia_KA}$ gauge equivalent to $V$ ($\nabla$ forward difference quotients); 
\[
\mathsf O_K(C,\bar s)=\{U:\mathsf m_s(U_s)\le C\text{ for }0<s\le\bar s\}.
\]
The \emph{flowed observables} are the functions $H^\sigma_K=\varphi(\mathcal O(f_1),\dots,\mathcal O(f_r))\circ\Phi_\sigma$, $\sigma>0$, $f_i\in C^\infty(T_\ell)$ real, $\varphi\in C^\infty(\mathbb R^r)$ bounded with all its derivatives. For $0<\alpha<1$ let $\mathcal O^\alpha$ be the space of orbits of $\mathfrak g^3$-valued $C^{0,\alpha}$ connections $\mathfrak A$ on $T_\ell$ under $C^{1,\alpha}$ gauge transformations; a \emph{flow line} is a continuous map $(0,\infty)\to\mathcal O^\alpha$ that near every time is the orbit of a mild solution of $\partial_s\mathfrak A=-\mathrm d_{\mathfrak A}^*F_{\mathfrak A}+\mathrm d_{\mathfrak A}(\partial_\nu\mathfrak A_\nu)$, the Yang--Mills gradient flow up to an infinitesimal gauge transformation; $\mathcal O^0(f)(\mathfrak A)=\frac12\int f\sum_{\mu<\nu}|F_{\mu\nu}(\mathfrak A)|^2$.

With $P_s=e^{s\Delta_a}$ the heat semigroup of $\Lambda_{K,\ell}$, $a=a_K$, $|f|_{\eta;a}=\sup_{0<s\le1}s^{-\eta/2}\|P_sf\|_\infty$ and $X\wedge Y=\frac1{2i}[X,Y]$, the rough regularity quantity of a $1$-form $\mathbf B$ is
\[
\|\mathbf B\|_{\kappa;a}=|\mathbf B|_{-\frac12-\kappa;a}+\sup_{0<s\le1}s^{\frac34+\kappa}\max\bigl|(\tau_h\nabla_\lambda P_s\mathbf B_\nu)\wedge(\tau_{h'}P_s\mathbf B_{\nu'})\bigr|_{-\frac12;a},
\]
the maximum over directions $\lambda,\nu,\nu'$ and translations $\tau_h,\tau_{h'}$ in a fixed finite set containing every translation in the quadratic term of the gauge-fixed flow equation below; it bounds $s^{\frac14+\frac\kappa2}\|P_s\mathbf B\|_\infty$ and $\frac12s^{\frac34+\frac\kappa2}\|\nabla P_s\mathbf B\|_\infty$. The rough distance of $\mathbf B$, $\mathbf B'$ adds to $|\mathbf B-\mathbf B'|_{-\frac12-\kappa;a}$ the same supremum for the products $\nabla P_s(\mathbf B-\mathbf B')\wedge P_s\mathbf B$ and $\nabla P_s\mathbf B'\wedge P_s(\mathbf B-\mathbf B')$.

\subsection{Markov limits, functional inequalities and flowed observables}

\begin{theorem}[Nontrivial Markov limits and mixing]\label{C-367-2}
Assume (E).
\begin{enumerate}
\item[(A1)]

\emph{Motion.} For every limit, every integer $n\ge n''$ and every invariant $\Phi\in C^1(\Omega_n)$, with $|\mathsf I_n|$ the Haar measure of $\mathsf I_n$,
\begin{equation}\label{C-367-3}
\mathcal E_\star(\Phi(Y_n))\ge\frac{e^{-\mathsf o_n}}{2C^\sigma|\mathsf I_n|}\int_{\mathsf I_n}\Gamma_n(\Phi)\,dV .
\end{equation}
Hence $(\Psi,Q_t\Psi)<\|\Psi\|^2$ for $\Psi=a_p(Y_n)$ and every $t>0$: no limit is frozen.
\item[(A2)]

\emph{Compactness.} Along every sequence $K_j\to\infty$, invariant $H_j\in D(\mathcal E_{K_j})$ with $\sup_j(\sup|H_j|+\mathcal E_{K_j}(H_j))<\infty$ have a strongly convergent subsequence; its limit $G$ has $|G|\le\sup_j\sup|H_j|$.
\item[(A3)]

\emph{The Markov property and form convergence.} For every limit, $\mathcal E_\star$ is a closed, densely defined Markovian form on $L^2_{\rm inv}$, and $Q_t=e^{-tA_\star}$ is a strongly continuous symmetric Markov semigroup on $L^2_{\rm inv}$; $P^{K_j}_tH_j\to Q_tF$ strongly for $t>0$ whenever $H_j\to F$ strongly, in particular $P^{K_j}_t\Psi_{K_j}\to Q_t\Psi$ for every invariant cylinder function $\Psi$; (M1) and (M2) hold, and $(\alpha-\mathcal L_{K_j})^{-1}W_j\to(\alpha+A_\star)^{-1}W$ strongly for $\alpha>0$ and invariant $W_j\to W$ strongly; $(Y_t)$ is a stationary reversible Markov process on the invariant $\sigma$-algebra with the transition operators $Q_t$ and the invariant measure $\mu_{\ell,\mathsf b}$; and $\{G\in L^2_{\rm inv}:|G|\le1,\ \mathcal E_\star(G)\le1\}$ is compact.

If invariant $H_j\in C^\infty(\Omega_{K_j})$ converge strongly to $F$ and $\sup_j\|\mathcal L_{K_j}H_j\|_{L^2(\mu^{K_j})}<\infty$, then $F\in D(A_\star)$, $\mathcal L_{K_j}H_j\to-A_\star F$ weakly and $\mathcal E_{K_j}(H_j)\to\mathcal E_\star(F)$; under the assumptions of (C7), this hypothesis fails at the flowed observables of (C7)(b), (c), (e). If invariant $u_j\to u$ strongly with $\mathcal E_{K_j}(u_j)\to\mathcal E_\star(u)$, and invariant $v_j\to v$ weakly with $\sup_j\mathcal E_{K_j}(v_j)<\infty$, then $\mathcal E_{K_j}(u_j,v_j)\to\mathcal E_\star(u,v)$.
\item[(A4)]

\emph{Variance bounds and mixing.} Let $n\ge\max\{n'',n_{\rm c}\}$ and $\mathsf C_n=C^\sigma[\frac{C_{\rm P}}{L^2-1}\frac{a_n^2}{\mathsf b}+\frac12C^{\rm t}_n]$, with $C^{\rm t}_n$ of (B1). For all $K>n$: (a) $\mathrm{Var}_{\mu^K}(H)\le\mathsf C_n\mathcal E_K(H)+\mathsf r_n\mathrm{osc}(H)^2$ for bounded invariant $H\in D(\mathcal E_K)$; (b) $\|P^K_tH-\mu^K(H)\|^2_{L^2(\mu^K)}\le(\mathsf r_n+\frac14e^{-2t/\mathsf C_n})\mathrm{osc}(H)^2$ for bounded Borel invariant $H$ and $t\ge0$. For every limit: (c) $\mathrm{Var}_\mu(F)\le\mathsf C_n\mathcal E_\star(F)+\mathsf r_n\mathrm{osc}(F)^2$ for bounded $F\in L^2_{\rm inv}$, so $\mathcal E_\star(F)=0$ only for constant $F$, and $(F,(1-Q_t)F)>0$ for every non-constant bounded $F\in L^2_{\rm inv}$ and $t>0$; (d) $0\le\mathbb EF(Y_0)F(Y_t)-(\int F\,d\mu)^2\le\inf_n(\mathsf r_n+\frac14e^{-t/\mathsf C_n})\mathrm{osc}(F)^2$ for bounded $F\in L^2_{\rm inv}$, the infimum over the admitted $n$, which tends to $0$ as $t\to\infty$; and a probability measure $\rho\,\mu$, $\rho\in L^2_{\rm inv}$, with $Q_t\rho=\rho$ for one $t>0$ is $\mu$.

(e) If $F_K\to F$ strongly along the full sequence of cutoffs, $\sup_K\sup|F_K|\le M_F$, $F$ is not constant, and $H_K$ are invariant Lipschitz functions with $\|H_K-F_K\|_{L^2(\mu^K)}\to0$, then $\liminf_K\mathcal E_K(H_K)\ge\sup_n(\mathrm{Var}_\mu F-4\mathsf r_nM_F^2)/\mathsf C_n>0$; and for invariant Lipschitz $H_K$ with $\sup_K\mathrm{osc}(H_K)\le O$, $\liminf_K\mathcal E_K(H_K)\ge\sup_n[\liminf_K\mathrm{Var}_{\mu^K}(H_K)-\mathsf r_nO^2]/\mathsf C_n$. (f) Under \textup{(Multiscale gauge)}, all of (a)--(e) holds as well for $n\ge\max\{n'',n_{\rm e}\}$ with $\mathsf R_n<1$, with $C^\flat_n$ of (B1)(e) in the place of $C^{\rm t}_n$ and $\mathsf R_n+\mathsf R_{n-J_n}$ in the place of $\mathsf r_n$.
\item[(A5)]

\emph{Entropy relaxation.} Let $n\ge\max\{n_{\rm a},n_{\rm c}\}$, $\mathsf C'_n$ and $\varsigma^K$ as in (B2). For all $K>n$: (a) if $G\in C^1(\Omega_K)$ is real and invariant, $M_G=\sup|G|$ and $m_2=\int G^2d\mu^K>\mathsf r_nM_G^2$, then
\[
\mathrm{Ent}_{\mu^K}(G^2)\le\mathsf C'_n\mathcal E_K(G)+\mathsf r_nM_G^2\Bigl[1+\log\frac{M_G^2}{m_2-\mathsf r_nM_G^2}\Bigr];
\]
(b) if $\psi\in C^1(\Omega_K)$ is real and invariant with $\Gamma_K(\psi)\le l^2$ and $\bar\psi=\int\psi\,d\varsigma^K/\varsigma^K(\Omega_K)$, then $\mu^K(\psi-\bar\psi\ge u)\le e^{-u^2/(\mathsf C'_nl^2)}+\mathsf r_n$ for $u>0$; (c) if $f_0\ge0$ is bounded Borel invariant, $m_1=\int f_0d\mu^K>0$, $\bar M=\sup f_0$ and $\mathsf r_n\bar M\le\frac12m_1$, then for $\tau\ge0$
\[
\mathrm{Ent}_{\mu^K}(P^K_\tau f_0)\le e^{-4\tau/\mathsf C'_n}\mathrm{Ent}_{\mu^K}(f_0)+\mathsf r_n\bar M\Bigl[1+\log\frac{2\bar M}{m_1}\Bigr].
\]
For every limit, (a) holds for bounded $F\in D(\mathcal E_\star)$ with $\mu,\mathcal E_\star$ in the places of $\mu^K,\mathcal E_K$; and if $F\in L^2_{\rm inv}$, $0\le F\le\bar M$, $m_1=\int F\,d\mu>0$ and $\mathsf r_n\bar M<\frac12m_1$, then $\mathrm{Ent}_\mu(Q_\tau F)\le e^{-4\tau/\mathsf C'_n}\mathrm{Ent}_\mu(F)+\mathsf r_n\bar M[1+\log\frac{2\bar M}{m_1}]$; so $\mathrm{Ent}_\mu(Q_\tau F)\to0$ as $\tau\to\infty$, and $m_1^{-1}(Q_\tau F)\,\mu\to\mu$ in total variation on the invariant $\sigma$-algebra. Under \textup{(Multiscale gauge)}, the same holds for $n\ge\max\{n_{\rm a},n_{\rm e}\}$ with $\mathsf R_n<1$, with the explicit top of (B1)(e).
\end{enumerate}
\end{theorem}

\begin{theorem}[Functional inequalities uniform in the cutoff]\label{C-367-4}
Assume (E).
\begin{enumerate}
\item[(B1)]

\emph{The coarse-scale law.} For $n\ge n_{\rm c}$ there are an invariant Borel weight $w_n:\Omega_n\to[0,1]$ and a constant $C^{\rm t}_n<\infty$ such that for all $K>n$ and every bounded invariant locally Lipschitz $\phi$ on $\mathsf I_n$: (a) $w_n\le\mathbf 1_{\mathsf D_n}$ and $\sigma^K_n(w_n<1)\le\mathsf R_{n-j_1}$; (b) $\mathrm{Ent}_{w_n\sigma^K_n}(\phi^2)\le C^{\rm t}_n\int\Gamma_n(\phi)\,w_n\,d\sigma^K_n$; (c) $\int(\phi-c)^2d\sigma^K_n\le\frac12C^{\rm t}_n\int\Gamma_n(\phi)\,d\sigma^K_n+\mathsf R_{n-j_1}\mathrm{osc}(\phi)^2$, $c$ the mean of $\phi$ under $w_n\sigma^K_n$; (d) $\mathrm{Ent}_{\sigma^K_n}(\phi^2)\le C^{\rm t}_n\int\Gamma_n(\phi)\,d\sigma^K_n+\mathsf R_{n-j_1}\sup_{\mathsf I_n}\Xi(\phi^2,c')$, $c'$ the mean of $\phi^2$ under $w_n\sigma^K_n$.

Here $w_n=\min\{1,\rho/\rho_0\}$, $\rho(V)=\int\mathbf 1_{S'}(V^g)\,dg$ for a connected open $S'\subset\mathsf D_n$ that is a finite union of chart balls, $\rho_0=\min_{\mathcal K_1}\rho>0$ on a compact invariant set $\mathcal K_1\subset\mathsf D_n$, and $C^{\rm t}_n=a_n\,c(S')\,e^{\mathsf o_n}/\rho_0$, $c(S')$ the log-Sobolev constant of Haar measure on $S'$. (e) (An explicit top.) Assume in addition \textup{(Multiscale gauge)}. Let $n\ge n_{\rm e}$, $\mathsf R_n<1$, $N=\ell L^n$, $w^{\rm e}_n(V)=\int\mathbf 1\{\mathsf d(V^u(b),\mathbf 1)<\alpha_n\text{ for all }b\}\,du$ ($du$ Haar measure of the gauge group), $c_n=[(\frac12\alpha_n-\sin\frac12\alpha_n)/2\pi]^{N^3}$ and $\chi_n=\min\{1,w^{\rm e}_n/c_n\}$.

Under this additional hypothesis, (a)--(d) hold with $\chi_n\mathbf 1_{\mathsf I_n}$ in the place of $w_n$, $\mathsf R_{n-J_n}$ in the place of $\mathsf R_{n-j_1}$, and the constant
\[
\begin{multlined}
C^\flat_n=\frac{a_ne^{2\mathsf o_n}C_{\mathscr D}}{(1-\mathsf R_n)c_n},\quad C_{\mathscr D}=2\sqrt e\Bigl(\frac{\alpha_n}{\sin\alpha_n}\Bigr)^{6N^3}\Bigl(\frac{2\rho_{\max}}{r'}\Bigr)^{9N^3}\Bigl(\rho_{\max}+\frac{8\pi\rho_{\max}^2}{r'}\Bigr)^2,\quad\rho_{\max} \\
=\alpha_n\sqrt{3N^3},\ \relax r'=64\alpha_n^2,
\end{multlined}
\]
in the place of $C^{\rm t}_n$; the set $\{\chi_n>0\}\cap\mathsf I_n$ is open, invariant, path connected and contained in $\mathsf D_n$; and for bounded invariant $\Phi\in C^1(\Omega_n)$, $\mathrm{Var}_{\mu^K_n}(\Phi)\le\frac12C^\flat_n\int_{\mathsf I_n}\Gamma_n(\Phi)\,d\mu^K_n+(\mathsf R_n+\mathsf R_{n-J_n})\mathrm{osc}(\Phi)^2$, and the same with $\mu_n$ in the place of $\mu^K_n$.
\item[(B2)]

\emph{The chain with a connected coarse-scale region.} Let $n\ge\max\{n_{\rm a},n_{\rm c}\}$, $K>n$, $\varsigma^K=w_n(V_t)\,\sigma^K$ and
\[
\mathsf C'_n=\Bigl[\frac{2C_{\rm P}}{L^2-1}\frac{a_n^2}{\mathsf b}+C^{\rm t}_n\Bigr]\,3\mathsf M^2\,\Pi^\flat_n .
\]
(a) $\int f\,d\varsigma^K\le\int f\,d\sigma^K\le\int f\,d\mu^K\le\int f\,d\varsigma^K+\mathsf r_n\sup f$ for bounded Borel invariant $f\ge0$; (b) $\mathrm{Ent}_{\varsigma^K}(G^2)\le\mathsf C'_n\int\Gamma_K(G)\,d\varsigma^K$ for real invariant $G\in C^1(\Omega_K)$; (c) for $\nu\in\{\sigma^K,\mu^K\}$: $\mathrm{Ent}_\nu(G^2)\le\mathsf C'_n\int\Gamma_K(G)\,d\nu+\mathsf r_n\sup_U\Xi(G(U)^2,c')$, $c'$ the mean of $G^2$ under $\varsigma^K$. (d) Assume in addition \textup{(Multiscale gauge)}. For Lebesgue-almost every label $\mathsf b$, all $m$ and all $n\ge\max\{n_{\rm a},n_{\rm e}\}$ with $\mathsf R_n<1$, (a)--(c) hold with $\mathbf 1_{\mathsf D_n}(V_t)\,\sigma^K$ in the place of $\varsigma^K$, a finite constant $C^\sharp_n$ in the place of $C^{\rm t}_n$, and $\mathsf R_n+\mathsf R_{n-J_n}$ in the place of $\mathsf r_n$; for every $\mathsf b$ and $\eta>0$ they hold with the component of $\mathbf 1$ in $\{\varpi<\tilde r\}$, for some $\tilde r\in(\varepsilon_n,2\varepsilon_n)$ that does not depend on $K$, in the place of $\mathsf D_n$, and $\mathsf R_n+\mathsf R_{n-J_n}+\eta$ in the place of $\mathsf r_n$; with the explicit top of (B1)(e) they hold with $\chi_n(V_t)\sigma^K$, $C^\flat_n$ and $\mathsf R_n+\mathsf R_{n-J_n}$.

(e) (No weight.) For $n\ge n_{\rm a}$, $K>n$, real invariant $G\in C^1(\Omega_K)$, $f_t$ the top of the tower of $G^2$ and $\epsilon>0$:
\[
\begin{multlined}
\mathrm{Ent}_{\sigma^K}(G^2)\le\frac{6C_{\rm P}\mathsf M^2\Pi^\flat_n}{L^2-1}\frac{a_n^2}{\mathsf b}\int\Gamma_K(G)\,d\sigma^K+\mathrm{Ent}_{\sigma^K_n}(f_t),\qquad\int\Gamma_n\bigl((f_t+\epsilon)^{1/2}\bigr)d\sigma^K_n \\
\le3\mathsf M^2\Pi^\flat_n\int\Gamma_K(G)\,d\sigma^K .
\end{multlined}
\]
\item[(B3)]

\emph{The fine field conditioned on a block field.} Let $n\ge n_{\rm a}$, $K>n$, $1\le t'\le K-n$, $e=K-t'$, $W\in\mathsf{Cut}_{t'}$, and let $\mathsf P_W$ be the law of $V_0$ when $V_{t'}=W$ and $V_k=V_{V_{k+1}}(B_k)$, $B_k$ of law $p_{V_{k+1}}$. Then $\mathsf P_W$ is a regular conditional law of $U$ given $\mathcal M^{t'}U=W$ under $\sigma^K$ and under $\varsigma^K$, and for real invariant $G\in C^1(\Omega_K)$ with tower $(g_k)$ and every $W\in\mathsf{Cut}_{t'}$
\[
\begin{multlined}
\mathrm{Ent}_{\mathsf P_W}(G^2)\le C^{\rm fib}_e\!\int\Gamma_K(G)\,d\mathsf P_W,\qquad\Gamma_e(g_{t'})(W)\le3\mathsf M^2\Pi^\flat_e\!\int\Gamma_K(G)\,d\mathsf P_W,\qquad C^{\rm fib}_e \\
=\frac{6\mathsf M^2C_{\rm P}\Pi^\flat_e}{L^2-1}\frac{a_e^2}{\mathsf b},
\end{multlined}
\]
a constant that depends neither on $K$ nor on $\ell$. If $\psi\in C^1(\Omega_K)$ is real and invariant and $\Gamma_K(\psi)\le l^2$ $\mathsf P_W$-almost everywhere, then $\mathsf P_W(\psi-\psi_{t'}(W)\ge u)\le e^{-u^2/(C^{\rm fib}_el^2)}$, $\psi_{t'}$ the member of the tower of $\psi$; and if $\Gamma_K(\psi)\le l^2$ on all of $\Omega_K$, then, with $\psi_{t'}:=0$ off $\mathsf{Cut}_{t'}$, $\mu^K(|\psi-\psi_{t'}(\mathcal M^{t'}U)|\ge u)\le2e^{-u^2/(C^{\rm fib}_el^2)}+\mathsf R_n$.
\end{enumerate}
\end{theorem}

\paragraph{Admissible top depths.}
The assertions about flowed observables use the following condition, in
the notation of Proposition~\ref{C-367-12}.

A top depth $d_0$ is \emph{admissible with constant $C_{\rm adm}$} if
$d_0\ge n_{\rm a}$, $d_0\ge d_{\rm num}(\ell,\mathsf b)$,
$c_\rho\varepsilon(L^{d_0}\mathsf b)\le C_{\rm BG}^{-2}$ and
$4c_{\mathsf h}\delta(L^{d_0}\mathsf b)\le1/4$,
and, for every $K>d_0$ and $d_0<d\le K$, (M) and (Col) hold
with that same constant.
Here $d_{\rm num}$ is the explicit numerical depth in the proof below.
The mean and covariance-error estimates hold at every
$W\in\mathsf{Cut}_{K-d+1}$; the Gaussian layer-kernel estimates hold
almost surely in the chain gauge, at every pair of bare bonds.
\textup{(Adm)} Under (E), there are $d_{\rm adm}$ and $C_{\rm adm}$,
depending on $\ell$, $\mathsf b$ and the construction constants but not
on $K$ or the chosen top depth, such that every $d_0\ge d_{\rm adm}$
is admissible with $C_{\rm adm}$. This is proved in
Proposition~\ref{C-367-12} below.

The proof of Proposition~\ref{C-367-12} establishes \textup{(M)} and
\textup{(Col)} uniformly in the cutoff and top depth whenever
$d_0\ge d_{\rm adm}$. Thus the flowed estimates use \textup{(Adm)}
as a consequence of \textup{(E)}.
\paragraph{Approximation of the fixed-depth background flow.}
Fix $0<\alpha<1$. For $D$ sufficiently large put
$\mathcal K_D=\{V\in\Omega_D:\varpi(V)\le\theta_D\}$, and let
$U_{K,D}(V)$ be any member of the minimal orbit over $V$ on
$\Lambda_{K,\ell}$, $K>D$. The unflowed comparison in the proof of
Proposition~\ref{C-367-12}(b) defines
$w^D_C(V)=\lim_{K\to\infty}W_C(U_{K,D}(V))$ for each closed lattice path $C$.
We isolate the following remaining approximation input, denoted
\emph{(Flow approximation)}. For each fixed $D$, every sequence
$K_\nu\to\infty$, $V_\nu\to V$ in $\mathcal K_D$ and arbitrary minimal
backgrounds $U_\nu$ over $V_\nu$ has a subsequence whose positive-time
flows converge modulo periodic gauges, locally uniformly in time
in $\mathcal O^\alpha$, to an orbit-valued flow line. More explicitly,
on each compact positive time interval, periodic gauges give real
bond potentials $A_\nu(s)$ representing $\Phi_sU_\nu$ as
$e^{ia_{K_\nu}A_\nu(s)}$ and local representatives $\mathfrak A_s$ of
the limit for which
\[
 \sup_{\tau\le s\le T}\max_b
 |A_\nu(s,b)-\mathfrak A_{s,\mu(b)}(x_b)|\longrightarrow0,
 \qquad 0<\tau<T<\infty.
\]
Here $x_b$ is the bond midpoint and $\mu(b)$ its direction.
The limit has local mild representatives satisfying
\[
 \partial_s\mathfrak A=-\mathrm d_{\mathfrak A}^{*}F_{\mathfrak A}
       +\mathrm d_{\mathfrak A}(\partial_\nu\mathfrak A_\nu),
 \qquad \lim_{s\downarrow0}W_C(\mathfrak A_s)=w^D_C(V),
\]
and there is at most one such orbit-valued flow line with these
initial holonomies. The compactness, equation identification and
uniqueness in this assertion remain an additional input. The
positive-time derivative bounds used below are proved separately
by the single-scale gauge and the short-time regularisation in
the proof of Proposition~\ref{C-367-12}(c).
\begin{theorem}[Observables regularised by the gradient flow]\label{C-367-5}
Assume (E) and $0<\kappa<\frac16$; (C1) holds for all $K$, $m$, $\beta>0$, with $\mathsf b=a_K\beta$. (C2)(b), (C2)(c) and (C3)--(C7), which read the moments, the good sets or the field $\mathbf B_K$ of Proposition~\ref{C-367-12}(a), use the consequence (Adm) of (E) proved in Proposition~\ref{C-367-12}: every top depth from some depth on, a depth free of $K$, is admissible in the sense of Proposition~\ref{C-367-12}. (C1) and (C2)(a) are free of it. The assertions in (C3), and those in (C5)(d) involving the identified limit $\mathsf G_\sigma$, also assume \textup{(Flow approximation)}. The assertions in (C7) use the nonlinear defect estimates of Lemma~\ref{C-367-14}.
\begin{enumerate}
\item[(C1)] For $G\in C^\infty(\Omega_K)$, $H^\sigma=G\circ\Phi_\sigma$ and $\sigma\ge0$: $\mathcal L_KH^\sigma=(\mathsf b\,\partial_\sigma+\Delta_K)H^\sigma$. Hence $\mathsf b\,\partial_\sigma\langle H^\sigma\rangle_{\mu^K}=-\langle\Delta_KH^\sigma\rangle_{\mu^K}$; for $\chi\in D(\mathcal E_K)\cap L^\infty$, $\int\chi\,(\mathsf b\partial_\sigma+\Delta_K)H^\sigma\,d\mu^K=-\int\Gamma_K(\chi,H^\sigma)\,d\mu^K=-\int(1-\chi)\,\mathcal L_KH^\sigma\,d\mu^K$; for $t\ge0$, $P^K_tH^\sigma=H^{\sigma+\mathsf bt}+\int_0^tP^K_s\Delta_KH^{\sigma+\mathsf b(t-s)}\,ds$; and for the stationary process and $0\le r\le t$, $H^{\sigma+\mathsf b(t-r)}(U_r)-H^{\sigma+\mathsf bt}(U_0)-\int_0^r(\Delta_KH^{\sigma+\mathsf b(t-s)})(U_s)\,ds$ is a martingale with the bracket $2\int_0^r\Gamma_K(H^{\sigma+\mathsf b(t-s)})(U_s)\,ds$: in the frame that moves along $\sigma+\mathsf bt=\mathrm{const}$ the drift of the action is gone.
\item[(C2)] (a) For $\bar C\ge1$, $0<\bar s\le1$ there is $K_0$ such that for every flowed observable, $0<\sigma\le1$ and $K\ge K_0$: $\Gamma_K(H^\sigma_K)\le C_\Gamma$ and $|\partial_\sigma H^\sigma_K|\le C_\partial$ on $\mathsf O_K(\bar C,\bar s)$, with $C_\Gamma,C_\partial$ depending on $\bar C,\bar s,\ell,\kappa,\sigma$ and the norms of $\nabla\varphi$, $f_i$, $\nabla f_i$, bounded for $\sigma$ in compact subsets of $(0,1]$; for $\sigma\le\bar s$ this holds for every $K$. (b) For $0<\sigma\le1$, $1\le p<\infty$ and all $K$: $\|\partial_\sigma H^\sigma_K\|_{L^p(\mu^K)}\le C_p(\sigma)\sum_i\|\partial_i\varphi\|_\infty(2\|f_i\|_\infty+4\|\nabla f_i\|_\infty)$, $C_p(\sigma)$ bounded on compact subsets of $(0,1]$; so $\|H^\sigma_K-H^{\sigma'}_K\|_{L^p(\mu^K)}\le C|\sigma-\sigma'|$ locally uniformly in $\sigma,\sigma'\in(0,1]$ and in $K$, and, by (C1), $\sup_K\|\mathcal L_KH^\sigma_K\|_{L^p(\mu^K)}<\infty$ if and only if $\sup_K\|\Delta_KH^\sigma_K\|_{L^p(\mu^K)}<\infty$.

(c) For compact $I\subset(0,1]$, $1\le q<\infty$: $\sup_K\|\partial_\sigma H^{\sigma''}_K-\partial_\sigma H^{\sigma'}_K\|_{L^q(\mu^K)}\le C_I(q)|\sigma''-\sigma'|^{1/2}$ for $\sigma',\sigma''\in I$.
\item[(C3)] For every flowed observable and $\sigma>0$, $H^\sigma_K\to\mathsf G_\sigma$ strongly along the full sequence of cutoffs, where
\[
\mathsf G_\sigma(Y)=\varphi\bigl(\mathcal O^0(f_1)(\mathfrak F(Y)_\sigma),\dots,\mathcal O^0(f_r)(\mathfrak F(Y)_\sigma)\bigr),
\]
$\mathfrak F$ the map of Proposition~\ref{C-367-12}(b); the same holds for real Lipschitz $f_i$ and bounded Lipschitz $\varphi$, and $\sigma\mapsto\mathsf G_\sigma$ is locally Lipschitz from $(0,1]$ into $L^2(\mu)$. For $0<\sigma<1$ there is $\dot{\mathsf G}_\sigma\in L^2_{\rm inv}$ with $\int\Psi_K\,\partial_\sigma H^\sigma_K\,d\mu^K\to\int\Psi\,\dot{\mathsf G}_\sigma\,d\mu$ along the full sequence for every invariant cylinder function $\Psi$; $\sigma\mapsto\int\Psi\dot{\mathsf G}_\sigma\,d\mu$ is locally H\"older continuous of exponent $\frac12$ and is the derivative of the continuously differentiable function $\sigma\mapsto\int\Psi\mathsf G_\sigma\,d\mu$.
\item[(C4)]

\emph{Regularisation and condenser cutoffs.} There are $R_\flat\ge1$ and, for $R\ge R_\flat$, numbers $\bar s(R)\in(0,1]$, $\bar K(R)$, $\bar\varkappa(R)>0$ with $R^q\bar\varkappa(R)\to0$ for every $q$, and for $K\ge\bar K(R)$ invariant $\chi_{K,R}\in C^\infty(\Omega_K)$ with values in $[0,1]$ such that $\chi_{K,R}=1$ on $\{\mathsf G^*\le R\}$ and $\chi_{K,R}=0$ off $\{\mathsf G^*\le\frac52R\}\subset\mathsf O_K(8R,\bar s(R))$, $\mathsf G^*$ the functional of Step~8; $\mathcal E_K(\chi_{K,R})+\int(1-\chi_{K,R})^2d\mu^K\le\bar\varkappa(R)$; $\sup_{K\ge\bar K(R)}\int\Gamma_K(\chi_{K,R})^pd\mu^K<\infty$ for every $p<\infty$; and $\mathbf 1_{\{\chi_{K,R}>0\}}\le\chi_{K,R'}$ for $R'\ge\frac52R$, $K\ge\bar K(R')$.

Let $w_{K,R}$ be the potential of Lemma~\ref{C-367-15} for $\eta=1-\chi_{K,R}$, and $\chi^{\rm c}_{K,R}=1-w_{K,R}$, the condenser cutoff. Then $\chi^{\rm c}_{K,R}\in D(\mathcal E_K)$ is invariant, takes values in $[0,1]$, is $1$ where $\chi_{K,R}=1$ and $0$ where $\chi_{K,R}=0$, has $\Gamma_K(\chi^{\rm c}_{K,R},\cdot)=0$ off the shell $\{0<\chi_{K,R}<1\}$, $\mathcal E_K(\chi^{\rm c}_{K,R})+\int(1-\chi^{\rm c}_{K,R})^2d\mu^K\le\bar\varkappa(R)$, and $|\int\Gamma_K(\chi^{\rm c}_{K,R},\psi)\,d\mu^K|\le\mathrm{osc}(\psi)\,\bar\varkappa(R)$ for $\psi\in D(\mathcal E_K)\cap L^\infty$. A second family of smooth invariant cutoffs, defined without choosing a gauge, is given by Proposition~\ref{C-367-22}.
\item[(C5)]

\emph{Static invariance and the generator.} (a) For every flowed observable, $\sigma>0$, $R\ge R_\flat$ and $K\ge\bar K(R)$,
\[
\Bigl|\int\chi^{\rm c}_{K,R}\,(\mathsf b\,\partial_\sigma+\Delta_K)H^\sigma_K\,d\mu^K\Bigr|\le\mathrm{osc}(\varphi)\,\bar\varkappa(R);
\]
for $0<\sigma\le1$, with $C_\partial=\sup_K\|\partial_\sigma H^\sigma_K\|_{L^2(\mu^K)}$, $|\int\chi^{\rm c}_{K,R}\Delta_KH^\sigma_K\,d\mu^K+\mathsf b\,\partial_\sigma\langle H^\sigma_K\rangle_{\mu^K}|\le\mathrm{osc}(\varphi)\bar\varkappa(R)+\mathsf b\,C_\partial\,\bar\varkappa(R)^{1/2}$; and $|\int\chi^{\rm c}_{K,R}\mathcal L_KH\,d\mu^K|\le\mathrm{osc}(H)\bar\varkappa(R)$ for every $H\in C^2(\Omega_K)$. (b) With $q(u)=3u^2-2u^3$, $q(\chi^{\rm c}_{K,R})\in D(\mathcal E_K)$ is invariant, takes values in $[0,1]$, is $1$ where $\chi_{K,R}=1$ and $0$ where $\chi_{K,R}=0$, has $\mathcal E_K(q(\chi^{\rm c}_{K,R}))\le\frac94\bar\varkappa(R)$, and its generator is a function: $\int q(\chi^{\rm c}_{K,R})\,\mathcal L_KH\,d\mu^K=\int H\,\mathsf l_{K,R}\,d\mu^K$ for $H\in C^2(\Omega_K)$, $\|\mathsf l_{K,R}\|_{L^1(\mu^K)}\le\frac{15}2\bar\varkappa(R)$. (c) For $H,H'\in C^2(\Omega_K)$,
\[
\begin{multlined}
\Bigl|\int\chi^{\rm c}_{K,R}\,(H'\mathcal L_KH-H\mathcal L_KH')\,d\mu^K\Bigr|\le\mathrm{osc}(HH')\,\bar\varkappa(R) \\
+2\sup|H|\,\bar\varkappa(R)^{1/2}\Bigl(\int_{\{0<\chi_{K,R}<1\}}\Gamma_K(H')\,d\mu^K\Bigr)^{1/2}.
\end{multlined}
\]
(d) Let $(K_j)$ be a limit and $F\in D(A_\star)$ with a test sequence, that is, invariant $F_j\in D(\mathcal L_{K_j})$ with $\sup_j\|F_j\|_\infty<\infty$, $F_j\to F$ and $\mathcal L_{K_j}F_j\to-A_\star F$ strongly; test sequences exist for $F=\alpha(\alpha+A_\star)^{-1}W$, $\alpha>0$, whenever $W$ is the strong limit of uniformly bounded invariant $W_j$ (for instance an invariant cylinder function or a $\mathsf G_\sigma$). Then for $R\ge R_\flat$ with $\bar\varkappa(R)\le1$
\[
\limsup_{j\to\infty}\Bigl|\int\chi^{\rm c}_{K_j,R}F_j\,(\mathsf b\,\partial_\sigma+\Delta_{K_j})H^\sigma_{K_j}\,d\mu^{K_j}+\int(A_\star F)\,\mathsf G_\sigma\,d\mu\Bigr|\le c_F\,\bar\varkappa(R)^{1/2},
\]
$c_F=\sup|\varphi|\sup_j[2\|F_j\|_\infty+2\mathcal E_{K_j}(F_j)^{1/2}+\|\mathcal L_{K_j}F_j\|_{L^2}]$: the displayed tested error tends to zero as $R\to\infty$. For every $\tau\in C^\infty_c(0,1)$ the corresponding distributional statement is \[\begin{multlined}\lim_{R\to\infty}\limsup_{j\to\infty}\Bigl|\int_0^1\tau(\sigma)\!\int\chi^{\rm c}_{K_j,R}F_j\,\Delta_{K_j}H^\sigma_{K_j}\,d\mu^{K_j}\,d\sigma\\+\int_0^1\bigl[\tau(\sigma)(A_\star F,\mathsf G_\sigma)-\mathsf b\,\tau'(\sigma)(F,\mathsf G_\sigma)\bigr]d\sigma\Bigr|=0.\end{multlined}\] For two flowed observables $H_K,H'_K$ at $\sigma,\sigma'$ with limits $\mathsf G,\mathsf G'$ and $\alpha>0$, \[\begin{multlined}\lim_R\limsup_j\Bigl|\int\chi^{\rm c}_{K_j,R}[\alpha(\alpha-\mathcal L_{K_j})^{-1}H'_{K_j}]\,\mathcal L_{K_j}H_{K_j}\,d\mu^{K_j}\\ +(\alpha A_\star(\alpha+A_\star)^{-1}\mathsf G',\mathsf G)\Bigr|=0,\end{multlined}\] a limit that is symmetric in the two observables.

If $C_1=\liminf_R\limsup_j\int\chi^{\rm c}_{K_j,R}|\mathcal L_{K_j}H^\sigma_{K_j}|\,d\mu^{K_j}<\infty$, then $\mathsf G_\sigma\in D(\mathcal E_\star)$ and $\mathcal E_\star(\mathsf G_\sigma)\le\sup|\varphi|\,C_1$; if $C_2=\liminf_R\limsup_j\|\chi^{\rm c}_{K_j,R}\mathcal L_{K_j}H^\sigma_{K_j}\|_{L^2(\mu^{K_j})}<\infty$, then $\mathsf G_\sigma\in D(A_\star)$ and $\|A_\star\mathsf G_\sigma\|\le C_2$. Without a cutoff, the criterion of (A3) gives $\mathsf G_\sigma\in D(A_\star)$, with $\mathcal L_{K_j}H^\sigma_{K_j}\to-A_\star\mathsf G_\sigma$ weakly and $\mathcal E_{K_j}(H^\sigma_{K_j})\to\mathcal E_\star(\mathsf G_\sigma)$, whenever $\sup_K\|\mathcal L_KH^\sigma_K\|_{L^2(\mu^K)}<\infty$; this fails at the members of (C7)(b), (c), (e).
\item[(C6)]

\emph{Lower energy bounds.} There is $p_\ast>0$, depending only on the constants of the construction, and there are $\sigma_\ast,c_\ast>0$, $S_\ast<\infty$ such that for $0<\sigma\le\sigma_\ast$, every $\varphi\in C^1(\mathbb R)$ and every invariant Borel $\chi_K:\Omega_K\to[0,1]$ with $\mu^K(\chi_K<1)\le p_\ast/8$,
\[
\liminf_{K\to\infty}\int\chi_K\,\Gamma_K(\varphi(A^a)\circ\Phi_\sigma)\,d\mu^K\ \ge\ c_\ast\,\sigma^{-\frac52+\kappa}\inf_{[0,S_\ast\sigma^{-3/2-\kappa}]}\varphi'^2 .
\]
For every $\vartheta<1$ there is $\sigma_\vartheta>0$ such that the left side is positive for $0<\sigma\le\sigma_\vartheta$, every $\varphi$ with $\varphi'>0$ and every such $\chi_K$ with $\mu^K(\chi_K<1)\le\vartheta$; in particular $\liminf_K\mathcal E_K(\varphi(A^a)\circ\Phi_\sigma)>0$. There are $\sigma_{\ast\ast}>0$ and $0<c_\circ\le C_\circ$, $c_\circ$ depending only on the constants of the construction, with
\[
c_\circ\,\mathsf b^{-1}\ell^3\sigma^{-3/2}\le\liminf_K\langle A^a\circ\Phi_\sigma\rangle_{\mu^K}\le\limsup_K\langle A^a\circ\Phi_\sigma\rangle_{\mu^K}\le C_\circ\,\mathsf b^{-1}\ell^3\sigma^{-3/2}\qquad(0<\sigma\le\sigma_{\ast\ast});
\]
for every $\varrho\in(0,1)$ there are $c_\varrho,\sigma_\varrho>0$ such that for $0<\sigma\le\sigma_\varrho$ and all large $K$, with $\mu^K$-probability at least $1-\varrho$, the flowed action exceeds $c_\varrho\mathsf b^{-1}\ell^3\sigma^{-3/2}$ and $\Gamma_K(A^a\circ\Phi_\sigma)$ exceeds $c_\varrho\sigma^{-5/2+\kappa}$; and for $0<\sigma\le\sigma_\ast$ the flowed current does not vanish in probability: $\liminf_K\mu^K((J,J)_a(U_\sigma)\ge c_J(\sigma))\ge\frac14p_\ast$ for some $c_J(\sigma)>0$.
\item[(C7)]

\emph{The generator without a regularisation cutoff.} Let $\bar s,\bar C$ be given by Proposition~\ref{C-367-12}(a) for some $0<\epsilon<1$, and $0<\sigma\le\bar s$. There is $K_1$ such that for $K\ge K_1$, $a=a_K$: (a) $\int|\partial(A_K\circ\Phi_\sigma)|^q\,d\mu^K\ge\exp[(6.6\,q-24.64)\,\sigma/a^2]$ for $q\in[3.74,8]$, and $\int\Gamma_K(A^a\circ\Phi_\sigma)^2d\mu^K\ge e^{1.76\,\sigma/a^2}$; (b) $\|\mathcal L_K\arctan(A^a\circ\Phi_\sigma)\|_{L^2(\mu^K)}\ge e^{0.8\,\sigma/a^2}$; (c) for $f(x)=1-\cos(2\pi x_1/\ell)$, $E^{\rm a}_\sigma=\frac32\ell^3(2\bar C(\frac\sigma2)^{-\frac34-\frac\kappa2}+8\bar C^2(\frac\sigma2)^{-\frac12-\kappa})^2$, $M_\sigma=(\frac4\sigma+\frac{11400}{\ell^2})(\frac32\sigma E^{\rm a}_\sigma+1)+1$ and every $\varphi\in C^\infty(\mathbb R)$, bounded with all its derivatives, with $\inf_{[0,M_\sigma]}|\varphi'|>0$ ($K_1$ depending on $\varphi$): $\|\mathcal L_K\varphi(\mathcal O(f)\circ\Phi_\sigma)\|_{L^2(\mu^K)}\ge e^{0.8\,\sigma/a^2}$; (d) for $p\in[1.87,4]$: $\|\mathcal L_K\cos(A^a\circ\Phi_\sigma)\|_{L^p(\mu^K)}+\|\mathcal L_K\sin(A^a\circ\Phi_\sigma)\|_{L^p(\mu^K)}\ge e^{0.04\,\sigma/(pa^2)}$, and for every $p>1.8637$ the left side is unbounded in $K$; (e) for every real $f\in C^\infty(T_\ell)$, $f\not\equiv0$, with $K_1$ and a constant $C$ depending on $f$ but not on $K$: $\|\mathcal L_K\sin(\mathcal O(f)\circ\Phi_\sigma)\|_{L^2(\mu^K)}\ge\exp[\frac\sigma{5a^2}-\frac{10\,\mathsf b}a-C]$, and $\|\mathcal L_K\varphi(\mathcal O(f)\circ\Phi_\sigma)\|_{L^2(\mu^K)}\ge e^{\sigma/20a^2}$ for every non-constant periodic $\varphi\in C^\infty(\mathbb R)$, $K_1$ depending besides on $\varphi$; (f) $\sup_K\int e^{13\Lambda_\sigma}d\mu^K=\infty$.
\end{enumerate}
\end{theorem}

\subsection{Renormalisation inputs}

We use Theorems~\YMMref{D}{8}, \YMMref{AJ}{673} and~\YMMref{AD}{607} among the construction results through its results on one step (Proposition~\ref{prop:source-entry} at zero source) and on the small-field chain on one torus (Sections~\YMMref{113}{786}, \YMMref{114}{799}, \YMMref{115}{805}, Section~\YMMref{116}{807}): setting (E) (Section~\ref{sec:cutoff-calibration}), the field cuts (the definitions preceding Lemma~\YMMref{113.10}{796}, Lemma~\YMMref{113.10}{796}) and the gauge action (Section~\YMMref{3}{13}), the chart at the pivots (Section~\ref{sec:sf-chart-interface}), the iteration from a level with its exponents and the bound of the non-local action (Proposition~\YMMref{114.2}{801}, Proposition~\YMMref{114.2}{801}, equation~\YMMref{(2163)}{798}, equation~\YMMref{(2180)}{803}), the comparison of the chain with the bare density through the weight $\mathfrak w^{\rm sf}_t$ and the bound $\mathsf R_n$ (the definitions preceding Proposition~\YMMref{115.2}{806}, Proposition~\textup{\YMMref{115.2}{806}}), and the convergence in total variation of the laws of the block fields (Corollary~\textup{\ref{IV-62-66}}).

Further construction results are cited where they are used: the constants and thresholds of the construction (equation~\YMMref{(634)}{339}, Section~\ref{sec:step-input-conventions}, \ref{sec:cutoff-calibration}, \ref{sec:cutoff-calibration}, \ref{sec:cutoff-calibration}, Section~\ref{sec:cutoff-calibration}), with the old-term constants defined in equation~\YMMref{(2163)}{798} and its following sentence; the current (\YMMref{9.2}{38}); the lattices and the localisation domains with their linear sizes (Section~\YMMref{3}{13}, \YMMref{72}{607} (the definitions of $\mathcal D_j$ and $d_j$ at its beginning), with \YMMref{72}{607} (the definitions of $\mathcal D_j$ and $d_j$ at its beginning)); averages (equations~\YMMref{(37)--(39)}{20}, Proposition~\YMMref{51.1}{390} and Lemma~\ref{lem:plain-average-remainder}), and of the averaging map its gauge covariance (Proposition~\YMMref{4.5}{20}, \YMMref{(40)}{20}), its real-analyticity on the small fields and its measurability (Proposition~\YMMref{4.5}{20}), its linear part at the identity, which is the linear average \YMMref{(42)}{22} of equation~\YMMref{(42)}{22} (Lemma~\YMMref{4.6}{22}), with the $k$-fold composites $\mathbf Q^{\rm lin}_k$ of the linear average, and the bound for rearranged paths (Lemma~\YMMref{4.1}{17}); backgrounds over a block field, whose averages give back the block field (Definitions~\YMMref{9.3}{40} and~\YMMref{9.5}{41}, Section~\ref{sec:analytic-core-complements}); the fixed radii and the analyticity cores (Definition~\YMMref{49.1}{374} and Lemma~\YMMref{82.1}{665}); the bounds of the response (Lemma~\ref{lem:response-current-zero}, Section~\ref{sec:background-step-complements}) and its Landau gauge on the kernel of the site averages (Section~\ref{sec:response-step-conventions}, equation~\YMMref{(837)}{398}); the constant-frame row map and its response identity, with $C_{\mathcal Q}=34(L^3+2)$ and the counting-measure kernel bounds (Lemma~\ref{lem:constant-frame-row}); the double-block constants of Definition~\ref{def:sf-chart-data}; the regularity numbers and the classes of covariance (Definition~\ref{def:sf-chart-data}, Lemma~\ref{lem:sf-class-window}, and the curvature comparison in the proof of Lemma~\YMMref{113.10}{796}, with $C_H$ as fixed in Section~\ref{sec:regular-background-complements}); backgrounds and small plaquettes (Proposition~\YMMref{52.3}{433}, Proposition~\ref{prop:uniform-box-gauge}); the local potentials and their bounds (Proposition~\YMMref{51.4}{393}), with their sums in equations~\YMMref{(828)}{395}--\YMMref{(830)}{395}; the outer substitution (Proposition~\YMMref{51.10}{403}); the sum over polymers \eqref{eq:analytic-domain-tail}; Section~\ref{sec:step-input-conventions}; Theorem~\YMMref{W}{423} and equation~\YMMref{(925)}{424}; the Gaussian structure of a step and its numbers (Definition~\YMMref{47.2}{359}, Theorem~\YMMref{U}{378} and Section~\ref{sec:step-input-conventions}); the gauge-fixing factor, fluctuation test and chart integral of Section~\ref{sec:sf-chart-interface}, with the measure identity of Lemma~\ref{lem:sf-chart-integral}; the number $\varepsilon_{\mathcal W}$ Proposition~\YMMref{52.3}{433}; and the consequences of the setting (Section~\ref{sec:cutoff-calibration}).

These are the construction results used in this section. The admissibility conclusion \textup{(Adm)} follows from (E). The additional hypotheses \textup{(Multiscale gauge)}, \textup{(Flow approximation)} and \textup{(Lower propagation)} are imposed only on the assertions that invoke them. In particular, the continuum-flow identification in Proposition~\ref{C-367-12}(b) uses the compactness, equation identification, initial holonomies and uniqueness specified in \textup{(Flow approximation)}.

Within the section, the proofs of (C2)(b), (C2)(c) and (C3)--(C7) start besides from the moments of the gauge-fixed field $\mathbf B_K$ of Proposition~\ref{C-367-12}(a), gathered before Step~6, and from its proved conclusion (Adm) on the top depths. The moment bound under the chain measure from an admissible top depth rests on Lemma~\ref{C-367-7}(iv), the bounds of the response and of the heat kernel, the grid lemma, and three statements on one fluctuation integral, set in the proof of Proposition~\ref{C-367-12}(a): the bounds (M) of the conditional mean of a layer, the bounds (Col) of the part of its conditional covariance that is not scalar in colour, and the exponential bound (Q) of centred quadratic forms under the log-concave law of one integral.

A top depth $d_0$ is admissible when it meets the numerical and uniform-estimate conditions stated before Theorem~\ref{C-367-5}; under (Adm) the chains of all the large top depths give the moments under $\mu^K$. The statements made for all $K$, $m$ and $\beta>0$, namely Lemma~\ref{C-367-6}(a), (C1), the short-time theory of the gauge-fixed flow in Proposition~\ref{C-367-12}(a) (the flow on $[0,s_0(R)]$ with its four bounds, the Lipschitz bound in the rough distance and the bound of the flowed action), Lemmas~\ref{C-367-13}, \ref{C-367-21} and \ref{C-367-15}, Lemma~\ref{C-367-14}, Proposition~\ref{C-367-22}(a), (b) and Proposition~\ref{C-367-23}(a), (b) and the upper bound in (c), use no continuum-limit result. The lower bound in Proposition~\ref{C-367-23}(c) is the separate input \textup{(Lower propagation)}. The identified continuum flow also assumes \textup{(Flow approximation)}.

\subsection{The lattice dynamics and a common core}

\begin{lemma}[Lattice dynamics, subsequential limits and a common core]\label{C-367-6}
(a) For all $K$, $m$, $\beta>0$: $\int G\,\mathcal L_KF\,d\mu^K=-\int\Gamma_K(G,F)\,d\mu^K$ for $F,G\in C^\infty(\Omega_K)$; $\beta\,\partial_bA_K=\mathsf b\,a_K^2J(U)(b)$, so $\mathcal L_K=\Delta_K+\mathsf b\,v$ is the generator of $dA=-\mathsf b\,J(U)\,dt+\sqrt2\,d\mathcal W$, $dU(b)=ia_K\circ dA(b)\,U(b)$, with $\mathcal W$ the Brownian motion of the $1$-forms for the scalar product $a_K^3\sum_b\langle\cdot,\cdot\rangle$: no constant is tuned. $P^K_t$ is a symmetric Markov semigroup on $L^2(\mu^K)$ that commutes with gauge transformations; $(U_{-t})$ has the law of $(U_t)$; $\mathbb E|F(U_t)-F(U_0)|^2=2(F,(1-P^K_t)F)\le2t\,\mathcal E_K(F)$; for $F\in D(\mathcal L_K)$, $\mathbb E|F(U_t)-F(U_0)|^2=2t\,\mathcal E_K(F)-\mathbb E|\int_0^t\mathcal L_KF(U_s)\,ds|^2\ge2t\,\mathcal E_K(F)-t^2\|\mathcal L_KF\|^2$; and $\mathcal E_K(P^K_tF)\le\|F\|^2/2et$; $v(G\circ\Phi_\sigma)=\partial_\sigma(G\circ\Phi_\sigma)$.

With $\Pi_y$ the Haar average over the gauge transformations at a site $y$ and $\Pi_S=\prod_{y\in S}\Pi_y$: $\mathcal E_K(F)\ge(2a_K)^{-1}\|F-\Pi_yF\|^2$ and $\|P^K_t(F-\Pi_SF)\|\le e^{-t/2a_K}\|F-\Pi_SF\|$.

(b) Assume (E). Every sequence of cutoffs has a subsequence that is a limit. For every limit: $(Y_t)$ is stationary, reversible and positive under time reflection ($\mathbb E[\overline{\theta\mathsf F}\,\mathsf F]\ge0$ for finite sums $\mathsf F$ of products of cylinder functions at times $\ge0$, $\theta$ the reflection of the times), with one-time law $\mu$; its multi-time means are real-analytic functions of the strictly positive consecutive time differences; its finite-dimensional laws are invariant under gauge transformations chosen independently at different times; $Q_t$ is a symmetric positivity preserving contraction of $L^2(\mu)$, $Q_t1=1$, $Q_t=\Pi Q_t\Pi$, and $\|Q_tF\|^2\le(F,Q_{2t}F)$.

If $H_j\to F$ strongly, the spectral measures of $H_j$ for $-\mathcal L_{K_j}$, taken as measures on $[0,\infty]$, converge weakly to a measure $\nu_F$ of mass $\|F\|^2$ with $(F,Q_tF)=\int e^{-t\lambda}d\nu_F$ for $t>0$; $\mathcal E_\star(F)=\int\lambda\,d\nu_F$; $\mathcal E_\star$ is closed, Markovian and finite only on $L^2_{\rm inv}$; (M2) holds, and $\mathcal E_\star(F)\le\liminf_j\mathcal E_{K_j}(H_j)$ whenever $H_j\to F$ strongly. The following are equivalent: $\|Q_t\Psi\|^2=(\Psi,Q_{2t}\Psi)$ for invariant cylinder functions $\Psi$ and $t>0$; $Q_{t+s}=Q_tQ_s$; $P^{K_j}_t\Psi_{K_j}\to Q_t\Psi$ strongly for these $\Psi$; $(Y_t)$ is a Markov process on the invariant $\sigma$-algebra with the transition operators $Q_t$.

They imply (M1), and $P^{K_j}_tH_j\to Q_tF$ strongly whenever $H_j\to F$ strongly. In general $Q_t=J^*R_tJ$, where $J$ is an isometry of $L^2(\mu)$ into a Hilbert space $\mathcal K$ and $(R_t)$ is a semigroup of positive self-adjoint contractions of $\mathcal K$, and $(F,Q_{2t}F)-\|Q_tF\|^2=\|(1-JJ^*)R_tJF\|^2$; the four properties hold if and only if $R_tJL^2_{\rm inv}\subset JL^2_{\rm inv}$ for all $t$.

(c) Assume (E). For every integer $n\ge1$ with $L^n\mathsf b\ge\mathsf t_\natural$ and every $K>n$, $\mu^K(\varpi(\pi_{K,d}U)<\theta_d$ for all $n\le d\le K)\ge1-\mathsf R_n$, and $\mu(\varpi(Y_d)<\theta_d$ for all $d\ge n)\ge1-\mathsf R_n$. Fix $\varphi_{\rm cut}\in C^\infty([0,\infty);[0,1])$ with $\varphi_{\rm cut}=1$ on $[0,1]$, $\varphi_{\rm cut}=0$ on $[4,\infty)$ and $|\varphi_{\rm cut}'|\le1$, and put
\[
\chi_d=\prod_p\varphi_{\rm cut}\bigl(\varpi_p^2/(2\theta_d)^2\bigr)\quad\text{on }\Omega_d,\qquad\mathsf k_d=24L^{d/2}\theta_d^{-1}(\ell L^d)^3,
\]
the product over the plaquettes of $\Lambda_{d,\ell}$: $\chi_d$ is smooth and invariant, takes its values in $[0,1]$, is $1$ on $\{\varpi\le2\theta_d\}$ and $0$ on $\{\varpi\ge4\theta_d\}$, and $\Gamma_d(\chi_d)^{1/2}\le\mathsf k_d\mathbf 1\{\varpi\ge\theta_d\}$; the number $\mathsf k_d$ depends on $\ell$ and not on $K$, and it is at most a constant times $\ell^3L^{4d}$. The span $\mathcal C$ of the functions $F=\Phi(Y_n)\prod_{d\ge n}\chi_d(Y_d)$, $n\ge n_1$, $\Phi\in C^\infty(\Omega_n)$ invariant, is dense in $L^2_{\rm inv}$; the lattice forms $F_K=\Phi\circ\pi_{K,n}\prod_{n\le d\le K}\chi_d\circ\pi_{K,d}$ are smooth and invariant, $|F_K|\le\sup|\Phi|$, $F_K\to F$ strongly along the full sequence, and $\sup_K\int\Gamma_K(F_K)^pd\mu^K<\infty$ for every $p<\infty$; the cutoffs $\chi^K_n=\prod_{n\le d\le K}\chi_d\circ\pi_{K,d}$ have $\mu^K(\chi^K_n<1)\le\mathsf R_n$ and $\|\Gamma_K(\chi^K_n)^{1/2}\|_{L^{2p}(\mu^K)}\to0$ as $n\to\infty$, uniformly in $K$.

Consequently, for every limit, $\mathcal E_\star$ is densely defined, $\nu_F(\{\infty\})=\|F-\Pi F\|^2$ for $F\in L^2(\mu)$, $(F,Q_tF)\to\|\Pi F\|^2$ and $Q_t\to1$ strongly on $L^2_{\rm inv}$ as $t\to0$, and $\frac12\mathbb E|F(Y_t)-F(Y_0)|^2\le t\sup_K\mathcal E_K(F_K)$ for $F\in\mathcal C$.
\end{lemma}

\begin{proof}[Proof sketch]
(a) Haar measure is invariant under left translations in each factor, so 
\[
\int\partial^\alpha_b(FGe^{-\beta A_K})\,dU=0,
\]
which is the integration by parts; (\YMMref{9.2}{38}) and $\beta=\mathsf b/a_K$ give the drift. Writing a tangent vector as $X=a_KB$ with a $1$-form $B$ per unit length, $a_K\|X\|^2=(B,B)_{a_K}$, and $\Gamma_K(F)$ is the square of the gradient for $(\cdot,\cdot)_{a_K}$: this is the stochastic equation. The closure of $(\mathcal E_K,C^\infty)$ is a Dirichlet form, gauge transformations act unitarily and preserve it, so they commute with $P^K_t$; the two bounds are $1-e^{-t\lambda}\le t\lambda$ and $\lambda e^{-2t\lambda}\le1/2et$ under the spectral measures of the spectral theorem, and the identity is 
\[
2\int(1-e^{-t\lambda})\,d\nu=2t\int\lambda\,d\nu-2\int(t\lambda-1+e^{-t\lambda})\,d\nu,
\]
with 
\[
\mathbb E|\int_0^t\mathcal L_KF(U_s)\,ds|^2=2\int(t\lambda-1+e^{-t\lambda})\,d\nu\le t^2\|\mathcal L_KF\|^2
\]
 by stationarity, $\nu$ the spectral measure of $F$.

The generators $X_\alpha$ of the rotations at $y$ act on six bonds, $\sum_\alpha|X_\alpha F|^2\le6|\partial F|^2$, and the first non-zero eigenvalue $3$ of the Laplacian of the unit sphere $S^3=G$ gives $\int\sum_\alpha|X_\alpha F|^2d\mu^K\ge3\|F-\Pi_yF\|^2$. The commuting projections $\Pi_y$ split $L^2(\mu^K)$ into joint eigenspaces invariant under $P^K_t$, on which $\frac d{dt}\|P^K_tF\|^2=-2\mathcal E_K(P^K_tF)$. Since $v$ is the field of $\Phi$, $v(G\circ\Phi_\sigma)(U)=\frac d{ds}G(\Phi_\sigma\Phi_sU)|_{s=0}$.

(b) The laws of the block fields at rational times lie in a compact set of measures on the compact metrisable space $\mathcal Y^{\mathbb Q}$. The means $\langle\Psi_0,P^K_{s_1}\Psi_1\cdots P^K_{s_n}\Psi_n\rangle$ extend holomorphically and boundedly to $\mathrm{Re}\,s_i>0$, so normal families carry convergence from rational to all times and give the real analyticity; convergence in total variation of the laws of the block fields carries it from continuous to Borel functions. Stationarity, reversibility and positivity under reflection hold at every cutoff, where $\mathbb E[\overline{\theta\mathsf F}\mathsf F]$ is the square of a norm by the Markov property, and pass to the limit.

By (a) the non-invariant parts decay in a time $O(a_K)$; this gives the independent gauges and $Q_t=\Pi Q_t\Pi$. Since $1$ and the $e^{-t\lambda}$ span a dense algebra of $C([0,\infty])$, the spectral measures converge; $\frac1t(1-e^{-t\lambda})\uparrow\lambda$ gives $\mathcal E_\star(F)=\int\lambda\,d\nu_F$, lower semicontinuity and closedness, and $(F,(1-Q_t)F)=\frac12\mathbb E(F(Y_t)-F(Y_0))^2$ the Markov property of the form. $\|Q_tF\|^2\le(F,Q_{2t}F)$ is Cauchy--Schwarz at the cutoff; recovery sequences are $\Lambda(\Lambda-\mathcal L_{K_j})^{-1}G_j$ with $\Lambda\to\infty$ slowly. In the equivalences, weak convergence with convergence of the norms is strong convergence, and $\|P^K_t\Psi_K\|^2=(\Psi_K,P^K_{2t}\Psi_K)$; the strong convergence of $P^{K_j}_t\Psi_{K_j}$ gives (M1) through the quadratic forms $\langle H,(1-P^{K_j}_t)H\rangle\ge2\langle H,(1-P^{K_j}_t)G_j\rangle-\langle G_j,(1-P^{K_j}_t)G_j\rangle$, and the non-invariant parts are handled by (a).

For the last statement, the form $\mathbb E[\overline{\theta\mathsf F}\,\mathsf G]$ is hermitian and positive; $\mathcal K$ is its completion, $JF=[F(Y_0)]$, $R_s$ the shift of the times by $s$, self-adjoint because $\theta$ reverses the shift and contractive by repeated Cauchy--Schwarz; $(F,Q_tF')=\langle JF,R_tJF'\rangle$, and Pythagoras gives the identity.

(c) The weight defined before Proposition~\YMMref{115.2}{806} of the chain with top depth $n$ vanishes unless all $\pi_{K,d}U$, $n\le d\le K$, lie in their cuts, and by the comparison of the chain with the bare density (Proposition~\textup{\YMMref{115.2}{806}}, domination and lost-mass bounds) its integral against $e^{-\beta A_K}dU$ falls short of $Z_K$ by at most $\mathsf R_nZ_K$; so off a $\mu^K$-mass $\le\mathsf R_n$ all $\pi_{K,d}U$ lie in the cuts, and $\varpi(\pi_{K,d}U)<\theta_d$; for $\mu$ by total variation at finitely many depths (Corollary~\textup{\ref{IV-62-66}}) and a monotone limit. Four statements among the construction results on the averaging map $\mathcal M$ enter: its gauge covariance $\mathcal M(U^u)=\mathcal M(U)^{u|}$, $u|$ the restriction of a gauge transformation $u$ to the sites of the coarser lattice (Proposition~\YMMref{4.5}{20}, \YMMref{(40)}{20}); $\mathcal M$ is Borel, and on the small fields, where the plaquette variables of the double block of a coarse bond lie within $(200L^2)^{-1}$ of $\mathbf 1$, it is given at that bond by the first of the three formulas of its definition and is real-analytic (Proposition~\YMMref{4.5}{20}); its linear part at $\mathbf 1$ (Lemma~\YMMref{4.6}{22}, \YMMref{(42)}{22}) is the straight average along the bonds plus the difference of the averages along the contours at the two ends of the coarse bond, which is the coarse gradient of a site function; and the bound for rearranged paths (Lemma~\YMMref{4.1}{17}), which gives on the triple block of a coarse bond $c$, the union of the blocks of $c_-$, of $c_+$ and of the end point of the next bond $c^+$ in the direction of $c$, a gauge in which every bond variable is $e^{iA_b}$ with $A_b\in\mathfrak g$, $|A_b|\le15L^2\epsilon$, if all plaquette variables there lie within $\epsilon\le(200L^2)^{-1}$ of $\mathbf 1$.

With them the averaging map acts on gradients of invariant functions as follows: if $\varpi(\mathcal M^jU)\le\epsilon'_j\le(48000L^4)^{-1}$ for $j<k$, then $\mathcal M^k$ is real-analytic near $U$ and
\[
\Gamma_{d+k}(\Psi\circ\mathcal M^k)(U)\le\prod_{j<k}(1+C_D\epsilon'_j)^2\,\Gamma_d(\Psi)(\mathcal M^kU),\qquad C_D=C_D(L),
\]
for invariant $\Psi\in C^1(\Omega_d)$: the transpose of the differential of $\mathcal M$ maps horizontal vectors to horizontal vectors with norm at most $L^{-1/2}(1+C_D\epsilon)$, because the linear average is the straight average along the bonds, of norm $L^{-1/2}$ by Cauchy--Schwarz, plus the coarse gradient of the contour averages, which is a gauge direction of the block field and is orthogonal to the horizontal vectors, and the rest has norm at most $37L^2C_L\epsilon$, by the gauge covariance in the gauge of the triple block and the mean value inequality on the ball $\{|A_b|\le(3200L^2)^{-1}\}$, $C_L$ the largest norm there of the differentials of $A\mapsto\mathcal M(e^{iA})(c)$ and of the rest, so that $C_D=37L^{5/2}C_L$ depends on $L$ alone; and $a_{d+k}^{-1}L^{-k}=a_d^{-1}$.

The cutoffs: $|\partial(\varpi_p^2)|\le4\varpi_p$ by the computation of Lemma~\ref{C-367-13}(a) on $\Omega_d$, so each factor of $\chi_d$ has a gradient of norm at most $4\varpi_p/(2\theta_d)^2\le4/\theta_d$, which vanishes unless $2\theta_d\le\varpi_p\le4\theta_d$, and there are $3(\ell L^d)^3$ plaquettes; this gives $\Gamma_d(\chi_d)^{1/2}\le\mathsf k_d\mathbf 1\{\varpi\ge\theta_d\}$. $F_K$ is smooth: by induction from the depth $K$ down, the closure of the set where $\prod_{e<e'\le K}\chi_{e'}\circ\pi_{K,e'}\ne0$ lies in $\{\varpi(\pi_{K,e'}U)\le4\theta_{e'}$ for $e<e'\le K\}$, where $4\theta_{e'}\le(48000L^4)^{-1}$ as $e'>e\ge n\ge n_1$, so that $\pi_{K,e}$ is real-analytic near it; on that open neighbourhood the product with $\chi_e\circ\pi_{K,e}$ is smooth, it is $0$ on the open complement of the closure, and the two open sets cover $\Omega_K$; the same holds for the factor $\Phi\circ\pi_{K,n}$.

This is the role of the cutoffs: they keep every block field that enters $F_K$ where $\mathcal M$ is given by its real-analytic formula. By the product rule and the averaging bound at $\epsilon'_j=4\theta_{K-j}$, with $\prod_j(1+C_D\epsilon'_j)\le\mathsf c_*=\exp(4C_D\sum_{e\ge n_1}\theta_e)$,
\[
\Gamma_K(F_K)^{1/2}\le\mathsf c_*\Bigl[\sup\Gamma_n(\Phi)^{1/2}+\sup|\Phi|\sum_{n\le d\le K}\mathsf k_d\,\mathbf 1\{\varpi(\pi_{K,d}U)\ge\theta_d\}\Bigr],
\]
and the event in the sum has $\mu^K$-mass $\le\max\{\mathsf R_d,\mathsf R_{d-1}\}$, by the first statement of (c) at the depth $d$, or at the depth $d-1$ if $d=K$. Minkowski's inequality and the decay of $\mathsf R_d$ faster than every exponential, against $\mathsf k_d=O(L^{4d})$, bound every $L^p$ norm uniformly in $K$, and $\|\Gamma_K(\chi^K_n)^{1/2}\|_{L^{2p}(\mu^K)}\le\mathsf c_*\sum_{d\ge n}\mathsf k_d\max\{\mathsf R_d,\mathsf R_{d-1}\}^{1/2p}\to0$ as $n\to\infty$. Density: a bounded Borel invariant function of a small depth is $\mu$-almost surely a Borel function of a large depth, it is approximated in $L^2(\mu_d)$ by invariant polynomials, and the cutoffs change it by $\le\mathsf R_d^{1/2}$.

Then $\mathcal E_\star(F)\le\liminf_K\mathcal E_K(F_K)$, $\|Q_tF-F\|^2\le2(F,(1-Q_t)F)\le2t\sup_K\mathcal E_K(F_K)$, and $\nu_F(\{\infty\})=\|F\|^2-\lim_{t\to0}(F,Q_tF)$. Each property of $(Y_t)$ and of $\mathcal E_\star$ stated in (b) and (c) is shared by the frozen process, which keeps $Y_0$ up to gauge transformations of the threads drawn from the Haar measure independently at different times, so that $Q_t=\Pi$ for $t>0$ and $\mathcal E_\star=0$ on $L^2_{\rm inv}$: motion is the content of (A1).
\end{proof}

\begin{proposition}[The limiting process on the common core]\label{C-367-17}
Assume (E); let $\mathcal C$ and $F_K$ be as in Lemma~\ref{C-367-6}(c), and $T>0$.
\begin{enumerate}
\item[(a)] (Paths.) For $F\in\mathcal C$ and $1\le p<\infty$, $\mathbb E|F_K(U_t)-F_K(U_s)|^{2p}\le C_p(F)|t-s|^p$ uniformly in $K$; the processes $(F_K(U_t))_{t\le T}$ are tight in $C[0,T]$, and along every limit $F(Y_t)$ has a continuous modification, locally H\"older continuous of every exponent $<\frac12$. For $d\ge n_1$ there is $\mathsf p_d(T)$, which does not depend on $K$ and tends to $0$ faster than every exponential in $d$, such that for all $K>d$, with probability at least $1-\mathsf p_d(T)$, $\varpi(\pi_{K,e}U_t)<2\theta_e$ for all $d\le e\le K$ and all $t\in[0,T]$.
\item[(b)] (Energies.) For $F\in\mathcal C$ the limit $\mathcal E_\infty(F)=\lim_K\mathcal E_K(F_K)$ exists along the full sequence of cutoffs, and $\mathcal E_\infty(F)=\int\Gamma_\infty(F)\,d\mu$ with
\[
\Gamma_\infty(F)(Y)=\lim_{d\to\infty}a_d^{-1}\,\mathsf t[\tfrac23,\tfrac16;Y_d]\bigl(\partial F^{[d]}(Y_d)\bigr)\quad\mu\text{-almost surely and in }L^1(\mu),
\]
where $F^{[d]}$ is the lattice form of $F$ at depth $d$ and 
\[
\mathsf t[\alpha,\beta;V](B)=\alpha\|B\|^2+2\beta\sum_c\langle B(c),\mathrm{Ad}(V(c))B(c^+)\rangle,
\]
$c^+$ the bond that follows $c$ in its direction. $\mathcal E_\infty(F)$ depends only on the class of $F$ in $L^2(\mu)$, $\mathcal E_\star\le\mathcal E_\infty$ on $\mathcal C$ for every limit, and for $F=\Phi(Y_d)\prod_{e\ge d}\chi_e(Y_e)$ with $\mathsf R_d<1$
\[
\mathcal E_\infty(F)\ge\bigl(\tfrac13-\tfrac12\bar\eta_d\bigr)(1-\mathsf R_d)\,e^{-\mathsf o_d}\,a_d^{-1}|\mathsf I_d|^{-1}\int_{\mathsf I_d}|\partial\Phi|^2dV,\qquad\bar\eta_d\to0 .
\]
\item[(c)] (The laws of the block fields.) For an integer $d\ge1$ with $L^d\mathsf b\ge\mathsf t_\natural$ and every bounded Borel invariant $f\ge0$ on $\Omega_d$,
\[
\begin{multlined}
(1-\mathsf R_d)\,e^{-\mathsf o_d}|\mathsf I_d|^{-1}\!\int_{\mathsf I_d}f\,dV\le\int f(Y_d)\,\mathbf 1\{\varpi(Y_e)<\theta_e\text{ for all }e\ge d\}\,d\mu\le\int f\,d\mu_d \\
\le e^{\mathsf o_d}|\mathsf I_d|^{-1}\!\int_{\mathsf I_d}f\,dV+\mathsf R_d\sup f,
\end{multlined}
\]
and the same at every cutoff $K>d$, with $\mu^K_d$ in the place of $\mu_d$ and with the middle member $\int f(\pi_{K,d}U)\,\mathbf 1\{\varpi(\pi_{K,e}U)<\theta_e\text{ for }d\le e\le K\}\,d\mu^K$. So an invariant Haar-null set has $\mu_d$- and $\mu^K_d$-mass $\le\mathsf R_d$; and if $\mathsf R_d<1$, every non-empty open subset of $\mathsf I_d$ has positive $\mu_d$-mass and $\mu^K_d$-mass bounded below uniformly in $K$.
\item[(d)] (Square field and martingales.) Let $(K_j)$ be a limit and $F\in\mathcal C$. There is $\Gamma_\star(F)\in L^2_{\rm inv}$ with $(2t)^{-1}\mathbb E[\phi(Y_0)(F(Y_t)-F(Y_0))^2]\to\int\phi\,\Gamma_\star(F)\,d\mu$ as $t\to0$ for bounded invariant $\phi$; $0\le\Gamma_\star(F)\le\Gamma_\infty(F)$, $\int\Gamma_\star(F)\,d\mu=\mathcal E_\star(F)$, and $\mathcal E_\star(\psi(F))=\int\psi'(F)^2\Gamma_\star(F)\,d\mu$ for $\psi\in C^2(\mathbb R)$. Along a further subsequence the lattice martingales converge, and $F(Y_t)=F(Y_0)+M^F_t+D^F_t$, with a continuous martingale $M^F$ with $\langle M^F\rangle_t=2\int_0^t\Gamma_\infty(F)(Y_s)\,ds$, and a continuous process $D^F$ with stationary increments, $\mathbb E|D^F_t|^2=2[t\,\mathcal E_\infty(F)-(F,(1-Q_t)F)]$, $\mathbb E[D^F_t(F(Y_t)-F(Y_0))]=0$, and $\lim_{t\to0}t^{-1}\mathbb E|D^F_t|^2=2[\mathcal E_\infty(F)-\mathcal E_\star(F)]$; in reversed time, $F(Y_t)-F(Y_0)=\frac12M^F_t-\frac12(\hat M^F_T-\hat M^F_{T-t})$ with a martingale $\hat M^F$ of the law of $M^F$.

The following are equivalent: $\mathcal E_\star(F)=\mathcal E_\infty(F)$; $\Gamma_\star(F)=\Gamma_\infty(F)$ $\mu$-almost everywhere; $D^F$ has zero energy; $\lim_{\Lambda\to\infty}\limsup_j\int_{(\Lambda,\infty)}\lambda\,d\nu^{K_j}_F=0$, $\nu^K_F$ the spectral measure of $F_K$. And $F$ is frozen ($Q_tF=F$ for all $t$) if and only if $D^F=-M^F$.
\item[(e)] (Motion from fields with a divergence.) A vector field with a divergence at the cutoff $K$ is a $\zeta\in L^2(\mu^K;\mathfrak g^{\rm bonds})$ with $\delta_K\zeta\in L^2(\mu^K)$ such that $a_K^{-1}\int\langle\zeta,\partial H\rangle\,d\mu^K=-\int H\,\delta_K\zeta\,d\mu^K$ for $H\in C^\infty(\Omega_K)$. If $a_{K_j}^{-1}\int|\zeta_j|^2d\mu^{K_j}\le\mathsf A^2$, $\|\delta_{K_j}\zeta_j\|\le\mathsf D$ and $\delta_{K_j}\zeta_j\to g$ weakly, then for $F\in L^2(\mu)$ and $t>0$
\[
\mathcal E_\star(F)\ge\frac{(F,g)^2}{\mathsf A^2},\qquad(F,(1-Q_t)F)\ge\frac{t\,(F,g)^2}{\mathsf A^2+t\,\mathsf D^2};
\]
conversely, if $\mathcal E_\star(F)>0$, such fields with $(F,g)\ne0$ exist. A field that does not depend on the configuration pairs to zero with every invariant gradient: $\int\langle\partial F,z\rangle\,d\mu^K=0$ for constant $z$ and invariant $F$.
\item[(f)] (An example.) In a limit of reversible diffusions observed through a block variable whose law converges in total variation, the inequality $\mathcal E_\star\le\mathcal E_\infty$ can be strict: for the reversible diffusion on $\mathbb R^2$ with the weight $e^{-\frac\kappa2(x-w)^2-\frac12w^2}$, the form $\int(|\partial_xF|^2+s^2|\partial_wF|^2)$ with fixed $s>0$ and the block variable $x$, as $\kappa\to\infty$ the law of $x$ converges in total variation, the square field of $x$ is $1$, the limit is an Ornstein--Uhlenbeck process, Markov and moving, and $\mathcal E_\star(x)=s^2/(1+s^2)<1=\mathcal E_\infty(x)$.
\end{enumerate}
\end{proposition}

\begin{proof}[Proof sketch]
(a) At every cutoff, for $H\in C^\infty(\Omega_K)$ and $0\le t\le T$, $H(U_t)-H(U_0)=\frac12M_t-\frac12(\hat M_T-\hat M_{T-t})$ with forward and backward martingales of brackets $2\int\Gamma_K(H)$ (the drift drops out, and the square field alone enters); the martingale moment inequalities and stationarity give $\mathbb E|H(U_t)-H(U_s)|^{2p}\le c_p2^p|t-s|^p\int\Gamma_K(H)^pd\mu^K$, and $\mathbb P(\sup_{t\le T}|H(U_t)-H(U_0)|\ge\lambda)\le10T\lambda^{-2}\mathcal E_K(H)$. With the bounds of Lemma~\ref{C-367-6}(c) this is uniform in $K$, and the criterion of Kolmogorov and Chentsov gives tightness and H\"older continuity. For the small fields apply the maximal bound to $H=\prod_{d\le e\le K}\chi'_e\circ\pi_{K,e}$, built from a second family of cutoffs, $\chi'_e=\prod_p\varphi_{\rm cut}(\varpi_p^2/\theta_e^2)$ on $\Omega_e$ with the function $\varphi_{\rm cut}$ of Lemma~\ref{C-367-6}(c): $\chi'_e$ is $1$ where $\varpi\le\theta_e$ and $0$ where $\varpi\ge2\theta_e$, and, as there, each factor has a gradient of norm at most $4\varpi_p/\theta_e^2\le8/\theta_e$, so that $\Gamma_e(\chi'_e)^{1/2}\le\mathsf k_e\mathbf 1\{\varpi\ge\theta_e\}$ with the same numbers $\mathsf k_e=24L^{e/2}\theta_e^{-1}(\ell L^e)^3$, which thus bound $\Gamma_e^{1/2}$ of both families.

As in the proof of Lemma~\ref{C-367-6}(c), the energy of $H$ is at most $\mathsf c_*^2\Sigma_d^2$, $\Sigma_d=\sum_{e\ge d}\mathsf k_e\max\{\mathsf R_e,\mathsf R_{e-1}\}^{1/2}$, with $\mathsf c_*=\exp(4C_D\sum_{e\ge n_1}\theta_e)$; $H(U_0)=1$ off a mass $\le\mathsf R_d$, and $H(U_t)>0$ forces $\varpi(\pi_{K,e}U_t)<2\theta_e$ for $d\le e\le K$; so $\mathsf p_d(T)=\mathsf R_d+10T\mathsf c_*^2\Sigma_d^2$.

(b) On the small fields $F_K=F^{[d]}\circ\pi_{K,d}$, $\partial F_K=\mathcal D^{(k)}(U)^*\partial F^{[d]}$ with $\mathcal D^{(k)}$ the differential of $\mathcal M^k$, and the transpose of the differential carries the quadratic forms along: $L^k\mathsf t[\alpha,\beta;U](\mathcal D^{(k)*}B)=\mathsf t[\mathsf A_L^k(\alpha,\beta);V](B)+O(\sum_j\epsilon'_j)\|B\|^2$ for horizontal $B$ at $V=\mathcal M^kU$, where $\mathsf A_L$ is the linear map of the flat average, with the eigenvector $(\frac23,\frac16)$ of eigenvalue $1$ (the Gram matrix of the kernels that are a tent along the bond and a box across it) and $(\frac13,-\frac16)$ of eigenvalue $L^{-2}$.

So $\Gamma_K(F_K)$ is close in $L^1(\mu^K)$, uniformly in $K$, to $a_d^{-1}\mathsf t[\frac23,\frac16;\cdot](\partial F^{[d]})\circ\pi_{K,d}$, and total variation gives the limit along the full sequence; the same estimate between two depths along a thread gives $\Gamma_\infty$. A function of $\mathcal C$ that vanishes $\mu$-almost everywhere has a gradient that vanishes off an invariant Haar-null set, whose mass is controlled by (c); this gives the dependence on the class. $\mathcal E_\star\le\mathcal E_\infty$ is the lower semicontinuity of Lemma~\ref{C-367-6}(b). The lower bound: on the threads with $\varpi(Y_e)<2\theta_e$ for all $e\ge d$, where all cutoffs $\chi_e$ equal $1$ near $Y_e$ and $F^{[e]}=\Phi\circ\mathcal M^{e-d}$, the same estimate between the depths $d$ and $e\to\infty$ gives $|\Gamma_\infty(F)(Y)-a_d^{-1}\mathsf t[\frac23,\frac16;Y_d](\partial\Phi(Y_d))|\le\frac12\bar\eta_da_d^{-1}|\partial\Phi(Y_d)|^2$, with $\bar\eta_d=4C_{\rm T}\mathsf c_*^2\sum_{e>d}\theta_e$, $C_{\rm T}=C_{\rm T}(L)$ the constant of the term $O(\sum_j\epsilon'_j)$ above and $\mathsf c_*$ that of the proof of Lemma~\ref{C-367-6}(c), so $\Gamma_\infty(F)(Y)\ge(\frac13-\frac12\bar\eta_d)a_d^{-1}|\partial\Phi(Y_d)|^2$ as $\mathsf t[\frac23,\frac16;V]\ge\frac13\|\cdot\|^2$; elsewhere $\Gamma_\infty(F)\ge0$.

The event of the middle member of (c) lies in that set of threads, and (c) at the invariant function $f=|\partial\Phi|^2$ gives the bound.

(c) At the top depth $d$ the chain from level $0$ has the density $e^{\Sigma+\Phi_t}$ with respect to Haar measure on $\mathsf I_d$, $\Sigma$ a constant (Proposition~\textup{\YMMref{115.2}{806}}, image identity, Proposition~\YMMref{114.2}{801}), its non-classical part bounded by $C_{\rm nl}\frac L{L-1}E_0(\ell L^d)^3$ (equation~\YMMref{(2180)}{803}) and the classical part between $0$ and $\frac32c_\rho^2A_0^2(\log L^d\mathsf b)^{2\mathfrak p_0}(\ell L^d)^3$ (by \YMMref{(2161)}{796} at the weight $x_t$); its mass lies in $[1-\mathsf R_d,1]$; the comparison of the chain with the bare density (Proposition~\textup{\YMMref{115.2}{806}}, domination and lost-mass bounds) gives the sandwich for $\mu^K_d$, and total variation (Corollary~\textup{\ref{IV-62-66}}) for $\mu_d$.

The middle member: the weight of the chain with top depth $d$ vanishes unless $\pi_{K,e}U$ lies in its cut, so in $\{\varpi<\theta_e\}$, for $d\le e\le K$; the function $f(\pi_{K,d}U)\mathbf 1\{\varpi(\pi_{K,e}U)<\theta_e$ for $d\le e\le K\}$ is invariant, by the gauge covariance of $\mathcal M$; so, $\sigma^K$ being at most $\mu^K$ on invariant functions (Lemma~\ref{C-367-7}(i)), $\int f\,d\sigma^K_d=\int f(\pi_{K,d}U)\mathbf 1\{\cdots\}\,d\sigma^K\le\int f(\pi_{K,d}U)\mathbf 1\{\cdots\}\,d\mu^K$, and $\int f\,d\sigma^K_d\ge(1-\mathsf R_d)e^{-\mathsf o_d}|\mathsf I_d|^{-1}\int_{\mathsf I_d}f\,dV$ by the lower bound of the density of $\sigma^K_d$ (Lemma~\ref{C-367-7}(iii)).

For $\mu$: at a depth $D>d$ the event $\{\varpi(Y_e)<\theta_e$ for $d\le e\le D\}$ is an event of the block field of depth $D$, it contains the event at the cutoff $K>D$, and total variation at the depth $D$ carries the bound to $\mu$; then $D\to\infty$, a monotone limit. Positivity: for a non-empty open $O\subset\mathsf I_d$ and a smaller non-empty open $O'$ with closure in $O$, the gauge orbit of $O'$ is open, invariant and of positive Haar measure, and it is covered by finitely many gauge translates of $O$, which all have the mass of $O$.

(d) For bounded invariant $\phi$, the forward--backward decomposition gives $(2t)^{-1}\mathbb E[\phi(U_0)(F_K(U_t)-F_K(U_0))^2]\le\int\phi\,\Gamma_K(F_K)+c\,\langle\phi,(1-P_t)\phi\rangle^{1/2}$ for $\phi\ge0$; the limit in $K$ and then in $t$ gives $\Gamma_\star\le\Gamma_\infty$, and, with $\mathsf e_t(\phi;F)=(2t)^{-1}\mathbb E[\phi(Y_0)(F(Y_t)-F(Y_0))^2]$, the identity $2t\,\mathsf e_t(\phi;F)=2\langle\phi F,(1-Q_t)F\rangle-\langle\phi,(1-Q_t)F^2\rangle$ gives the existence. The chain rule follows from fourth moments of the increments, which are $O(t^2)$. The lattice processes $F_K(U_t)$, their martingales $M^{K,F}$ and brackets are tight in $C[0,T]$, limits of $L^4$-bounded martingales are martingales, and the brackets converge to $2\int\Gamma_\infty(F)(Y_s)ds$ by (b).

At the cutoff, with the spectral measure $\nu$ of $F_K$, $\mathbb E|\int_0^t\mathcal L_KF_K(U_s)ds|^2=2\int(t\lambda-1+e^{-t\lambda})d\nu=2[t\,\mathcal E_K(F_K)-(F_K,(1-P^K_t)F_K)]$ and the drift is orthogonal to the increment; these pass to the limit. The equivalences: $\mathcal E_\infty(F)-\mathcal E_\star(F)=\lim_\Lambda\lim_j\int(\lambda-\Lambda)_+d\nu^{K_j}_F$, the energy that escapes to infinite frequency, is the energy of the drift and the integral of $\Gamma_\infty-\Gamma_\star$.

(e) Let $H_j\to F$ strongly with $\mathcal E_{K_j}(H_j)\to\mathcal E_\star(F)$; then $|\langle H_j,\delta_{K_j}\zeta_j\rangle|\le\mathsf A\,\mathcal E_{K_j}(H_j)^{1/2}$, and the left side tends to $|(F,g)|$. At a positive time, $1-e^{-x}\ge\frac x{1+x}$ for $x\ge0$ (as $e^x\ge1+x$) and $\int\frac{t\lambda}{1+t\lambda}d\nu=\min_G\{\|H-G\|^2+t\,\mathcal E_K(G)\}$, which is bounded below, for every real $s$, by $2s\langle H,\delta_K\zeta\rangle-s^2(\|\delta_K\zeta\|^2+t^{-1}a_K^{-1}\int|\zeta|^2)$. Conversely, $\zeta_j=\partial(1-\mathcal L_{K_j})^{-1}H_j$ has $\delta_{K_j}\zeta_j=(1-\mathcal L_{K_j})^{-1}H_j-H_j\to(1+A_\star)^{-1}F-F$ by (A3), and $(F,(1+A_\star)^{-1}F-F)=-\int\frac\lambda{1+\lambda}d\nu_F\ne0$. For constant $z$, $\int\partial_bF\,d\mu^K$ is invariant under $\mathrm{Ad}(g)$ for all $g$, hence $0$.

(f) The operator of the diffusion on linear functions has the trace $\kappa(1+s^2)+s^2$ and the determinant $\kappa s^2$; its eigenvalues are $\lambda_+\sim\kappa(1+s^2)$ and $\lambda_-\to s^2/(1+s^2)$, and the spectral measure of $x$ puts a mass asymptotic to $\frac1{1+s^2}\lambda_+^{-1}$ at $\lambda_+$, so $\int\lambda\,d\nu^\kappa_x=1$ for all $\kappa$ while $\mathcal E_\star(x)=s^2/(1+s^2)$.
\end{proof}

\subsection{Conditional fluctuation laws and transport estimates}

\begin{lemma}[Conditional laws along the renormalisation chain]\label{C-367-7}
Assume (E), let $n\ge1$ be an integer with $L^n\mathsf b\ge\mathsf t_\natural$, and let $K>n$. (i) $1-\mathsf R_n\le\sigma^K(\Omega_K)\le1$, and $\int f\,d\sigma^K\le\int f\,d\mu^K\le\int f\,d\sigma^K+\mathsf R_n\sup f$ for bounded Borel invariant $f\ge0$. (ii) For $0\le k<t$, $d=K-k$, and every $W\in\mathsf{Cut}_{k+1}$, $p_W$ is a probability measure with a positive continuous density on the box; $\sigma^K_{d-1}$ is carried by $\mathsf{Cut}_{k+1}$, and
\[
\int\Bigl[\int f(V_W(B))\,p_W(dB)\Bigr]\sigma^K_{d-1}(dW)=\int f\,d\sigma^K_d\qquad\text{for every Borel }f\ge0\text{ on }\Omega_d ;
\]
so the images of the $p_W$ under the $V_W$ form a regular conditional law of $V_k$ given $V_{k+1}=W$ under $\sigma^K$, and under $\sigma^K$ the sequence $V_t,\dots,V_0$ is a Markov chain from coarse to fine whose kernel at every $W\in\mathsf{Cut}_{k+1}$ is the image under $V_W$ of $p_W$; the kernels carry bounded continuous functions on $\mathsf{Cut}_k$ to continuous functions on $\mathsf{Cut}_{k+1}$, that is, $W\mapsto\int h(V_W(B))\,p_W(dB)$ is continuous on $\mathsf{Cut}_{k+1}$ for every bounded continuous $h$ on $\mathsf{Cut}_k$; kernels and towers do not depend on the top level; $h_k(V_k)$ is the conditional mean of $H$ given $V_k$; the tower of an invariant $H$ is invariant at every point of $\mathsf{Cut}_k$, with $\inf H\le h_k\le\sup H$; and a density that is a function of the first member $V_t$ changes no kernel and no tower.

(iii) $\sigma^K_n=\mathsf g\,dV$ on $\mathsf I_n$, $\sup\mathsf g\le e^{\mathsf o_n}\inf\mathsf g$, and $(1-\mathsf R_n)e^{-\mathsf o_n}\le\mathsf g\,|\mathsf I_n|\le e^{\mathsf o_n}$ at every point. (iv) For $n\ge n_2$, $0\le k<t$, $W\in\mathsf{Cut}_{k+1}$: $p_W=e^{-\Psi_W}dB$ with $\Psi_W$ real-analytic and $\frac14\gamma_{\rm r}\le D^2\Psi_W\le\Lambda_+$; every tangent vector $X$ of the fibre $\{\mathcal MV=W\}$ at $V_W(B)$ has $\|X\|^2\le C_{\rm ch}\sum|X(b)|^2$, the sum over the bonds that are not pivot bonds; and for bounded $g\in C^1(\mathsf{Cut}_k)$
\[
\mathrm{Var}_{p_W}(g\circ V_W),\ \ \tfrac12\mathrm{Ent}_{p_W}\bigl((g\circ V_W)^2\bigr)\ \le\ C_{\rm P}\frac{a_d^2}{\mathsf b}\int\Gamma_d(g)(V_W(B))\,p_W(dB),
\]
with $C_{\rm P}=4C_{\rm ch}/\gamma_{\rm r}$, which depends only on the constants of the construction, and not on $K$, $k$, $W$, $\ell$ or $\mathsf b$.
\end{lemma}

\begin{proof}[Proof sketch]
(i) is Proposition~\textup{\YMMref{115.2}{806}}, domination and lost-mass bounds, the comparison of the small-field chain with the bare density, with the weight $\mathfrak w^{\rm sf}_t$ and the bound $\mathsf R_n$ of the definitions preceding Proposition~\YMMref{115.2}{806}; (ii): the weight defined before Proposition~\YMMref{115.2}{806} vanishes unless every $V_j$ lies in its cut, so $\sigma^K_{d-1}$ is carried by $\mathsf{Cut}_{k+1}$. Positivity at every point, by induction over the scales: with $\Sigma_k$ the sum of the constants of the steps $1,\dots,k$ at the source $0$, $e^{\Sigma_0}F_0=F_0$ is the bare density on $\mathsf{Cut}_0$ (the definition preceding Proposition~\YMMref{114.2}{801}), and if $e^{\Sigma_k}F_k>0$ on $\mathsf{Cut}_k$, then $e^{\Sigma_{k+1}}F_{k+1}=\mathsf I[e^{\Sigma_k}F_k]>0$ on $\mathsf{Cut}_{k+1}$ \YMMref{(2174)}{801}, because the prefactors of the chart integral are positive, $G_W$ and $\mathsf J_W$ are real (Lemma~\ref{lem:sf-chart-integral}) and the chart carries the box into $\mathsf{Cut}_k$ (Lemma~\YMMref{113.10}{796}); so the density of $p_W$, a positive multiple of $e^{-x_kG_W+\mathsf J_W}(e^{\Sigma_k}F_k)\circ V_W$, is positive at every point of the box, and it is continuous, as $(W,B)\mapsto V_W(B),G_W(B),\mathsf J_W(B)$ are continuous and the chart integral of a continuous function is continuous (Lemma~\ref{lem:sf-chart-integral}), which makes every $F_k$ continuous on $\mathsf{Cut}_k$ by the same induction.

The kernels carry bounded continuous functions to continuous functions, in the whole range $L^n\mathsf b\ge\mathsf t_\natural$: at a bounded continuous $h$ on $\mathsf{Cut}_k$ the kernel at $W$ is $\mathsf I[hF_k](W)/\mathsf I[F_k](W)$, the normalised chart integral \eqref{eq:sf-chart-integral} of the integrand of the chain; $\mathsf{Cut}_{k+1}$ lies in the set $\mathsf O^{\rm im}$ (Lemma~\YMMref{113.10}{796}), where the chart, its gauge-fixing form and the exponent of its Jacobians are continuous in $(W,B)$ (Lemma~\ref{lem:sf-chart-integral}), and $F_k$ is continuous on $\mathsf{Cut}_k$, into which the chart carries the box; so the chart integrals of the non-negative functions $(\sup|h|\pm h)F_k$ integrate over the fixed box functions that are continuous in $(W,B)$, Fatou's lemma on the box makes them lower semicontinuous in $W$, and as their sum $2\sup|h|\,\mathsf I[F_k]=2\sup|h|\,e^{\mathsf c_{k+1}(0)}F_{k+1}$ is continuous by the induction above, both are continuous; the denominator $\mathsf I[F_k]$ is continuous and positive.

The identity: the factors of the weight of the levels $>k$ are functions of $V_{k+1}$; the measure of the chain of the level $k$ is $\nu_k=e^{\Sigma_k}\mathbf 1_{\mathsf{Cut}_k}F_k\,dV$ the measure recursion in Proposition~\YMMref{114.2}{801}, \YMMref{(2175)}{802}; and by the identity of measures of the step (Lemma~\ref{lem:sf-chart-integral}) the image under $\mathcal M$ of $\mathfrak G_k[x_k]\chi^{\rm sf}_{x_k}\,h\,\nu_k$ has the density $e^{\Sigma_k}\mathsf I[h\mathbf 1_{\mathsf{Cut}_k}F_k]$ on the open set $\mathsf O^{\rm im}$, which contains $\mathcal W_{k+1}(x_k;\alpha_0,\alpha_1)\supset\mathsf{Cut}_{k+1}$ (Lemma~\YMMref{113.10}{796}), for every bounded Borel $h$ on $\Omega_d$, and on $\mathsf{Cut}_{k+1}$ this density is $e^{\Sigma_k}\mathsf I[hF_k]$, the chart carrying the box into $\mathsf{Cut}_k$; dividing by $e^{\Sigma_k}\mathsf I[F_k]$ gives the identity for bounded Borel $f\ge0$, and monotone convergence for every Borel $f\ge0$.

As $V_{k+2},\dots,V_t$ are functions of $V_{k+1}$, the identity at $f=h\cdot(g\circ\mathcal M)$ is the Markov property with the kernel $p_W$. The tower of an invariant $H$ is invariant at every point by the covariance of the chain under the gauge transformations that are constant on blocks; (iii) follows from the statements on the iteration from a level: by Proposition~\textup{\YMMref{115.2}{806}}, image identity and Proposition~\YMMref{114.2}{801}, at the top the density is $e^{\Sigma+\Phi_t}$ on $\mathsf I_n$, $\Sigma$ a constant, with a non-classical part bounded by $C_{\rm nl}\frac L{L-1}E_0(\ell L^n)^3$ (equation~\YMMref{(2180)}{803}) and a classical part between $0$ and $\frac32c_\rho^2A_0^2(\log L^n\mathsf b)^{2\mathfrak p_0}(\ell L^n)^3$ (by \YMMref{(2161)}{796} at the weight $x_t$), so it oscillates by at most $\mathsf o_n$, and its mass is that of (i).

(iv): at every $W\in\mathsf{Cut}_{k+1}$, which lies in the class $\mathcal W_{k+1}(x_k;\alpha_0,\alpha_1)$ of the sum rule (Lemma \YMMref{113.10}{796}), the sum identities in equations~\YMMref{(828)}{395}--\YMMref{(830)}{395} and their preceding Jacobian calculation and the step (Proposition~\ref{prop:source-entry}) give the density of $p_W$ the form $e^{\mathcal E(P_W;B)}$, normalised, of Proposition~\ref{C-367-8}(b). This form holds at every depth, at every $W$ of the class of the sum rule, where the restrictions of $P_W$ to the polymers lie in the analyticity cores and the sum rules hold, so that every term of $\mathcal E(P_W;\cdot)$ is formed and holomorphic in $B$ on $\{|B(b)|<C_A\mathfrak b(x_k)\}$ (Theorem~\YMMref{W}{423}, Proposition~\YMMref{51.4}{393} and its proof, the potential definitions in Section~\YMMref{51}{389}, and Section~\ref{sec:step-input-conventions}); the conditions of $n_3$ enter Proposition~\ref{C-367-8} only through the ball of complex triples of (a) and through the response.

So the form holds for every integer $n\ge1$ with $L^n\mathsf b\ge\mathsf t_\natural$, from $n_2$ on in particular, as Proposition~\ref{C-367-8}(b) states. The density being positive by (ii), $e^{\mathcal E(P_W;\cdot)}$ is a constant multiple of a positive function, so the imaginary part of $\mathcal E(P_W;\cdot)$ is constant on the connected box, and $\Psi_W=-\mathcal E(P_W;\cdot)+\log\int e^{\mathcal E(P_W;B')}dB'$, with the branch of the logarithm that makes it real, is real-analytic near the closed box. Its Hessian is $D^2\Psi_W=\mathsf A_W-D^2\sum_Yv_Y$, with the form $\mathsf A_W=\sum_Y\mathsf b_Y(P_W)+\sum_\square\hat{\mathsf G}_\square(P_W)$ of the Gaussian integral of one step at $P_W$, whose inverse is the whole covariance of the step at $P_W$.

It lies within $\frac12\mathsf u_{\rm G}(LM)^{-3}\le\frac12\gamma_{\rm r}$, in the $\ell^2$ operator norm, of a real reference form whose floor and ceiling are $\gamma_{\rm r}$ and $c_{\rm r}$ (the completeness (St4) of Definition~\YMMref{47.2}{359}, Theorem~\YMMref{W}{423}, the proof of Proposition~\YMMref{113.4}{791}), so $\frac12\gamma_{\rm r}\le\mathsf A_W\le c_{\rm r}+\frac12\gamma_{\rm r}$. The potentials are bounded by a constant times $\delta(x_k)\mathfrak b(x_k)^2e^{-\lambda d(Y)}$ on the whole product of the cores and the polydisc of radius $C_A\mathfrak b(x_k)$, so at $P_W$ for every $W$ of the class of the sum rule (Theorem~\YMMref{W}{423}, Proposition~\YMMref{51.4}{393} and its proof, as in the proof of Proposition~\ref{C-367-8}(d)), and by Cauchy's estimate their Hessians are at most $C_{\rm Hs}\delta(x_k)\le\frac14\gamma_{\rm r}$ for $n\ge n_2$; hence $\frac14\gamma_{\rm r}\le D^2\Psi_W\le c_{\rm r}+\frac34\gamma_{\rm r}\le\Lambda_+$.

The log-Sobolev inequality $\mathrm{Ent}_p(f^2)\le\frac2\kappa\int|\nabla f|^2dp$ of a probability measure $p=e^{-\Psi}dx$ on an open convex set with $D^2\Psi\ge\kappa$, at $\kappa=\frac14\gamma_{\rm r}$ on the convex box, gives $\mathrm{Ent}_{p_W}(F^2)\le\frac8{\gamma_{\rm r}}\int|\nabla_BF|^2dp_W$, and its linearisation the Poincar\'e inequality with $\frac4{\gamma_{\rm r}}$; since $x_k^{-1}=a_d/\mathsf b$, the chart gives $|\nabla_B(g\circ V_W)|^2\le C_{\rm ch}x_k^{-1}|\partial g|^2=C_{\rm ch}\frac{a_d^2}{\mathsf b}\Gamma_d(g)\circ V_W$.
\end{proof}

\begin{proposition}[Log-concavity of the fluctuation law]\label{C-367-8}
Assume (E). For $0\le k<K$, lengths at the level $k+1$ are measured in its units, in which the bare spacing is $\eta_{k+1}=L^{-(k+1)}$; $J^{(k+1)}(U)=\eta_{k+1}^{-3}\partial A_K(U)=L^{-3(K-k-1)}J(U)$ is the current (\YMMref{9.2}{38}) formed with this spacing; for a $1$-form $A$ and a configuration $U$, $(\nabla^U_\mu A_\nu)(x)=\eta_{k+1}^{-1}[\mathrm{Ad}(U(\langle x,x+\eta_{k+1}e_\mu\rangle))A_\nu(x+\eta_{k+1}e_\mu)-A_\nu(x)]$, and $|\nabla_UA|(b)=\max_\mu|(\nabla^U_\mu A_\nu)(x)|$ at the bond $b=\langle x,x+\eta_{k+1}e_\nu\rangle$. $\pi^M_k$ is the partition of the torus into the cubes of side $M$ in the units of level $k$ (the partition $\pi_k$ of Section~\YMMref{72}{607} (the definitions of $\mathcal D_j$ and $d_j$ at its beginning)); $\mathcal D_k$ is the set of its localisation domains, the unions of cubes of $\pi^M_k$ that are connected through their faces, which we call polymers; $d(Y)=d_k(Y)$ is the linear size \YMMref{72}{607} (the definitions of $\mathcal D_j$ and $d_j$ at its beginning) of a polymer $Y$, the least number of edges of a tree of edges of its cubes that meets all of them; $\mathrm{Bd}(Y)$ is the set of the bare bonds with an end point in $Y$; and for a function $f$ of triples and a pair $(a,\mathbf j)$ of complex $1$-forms, $D_{(a,\mathbf j)}f(P_W)=\frac d{dz}f(U_+,za,\mathbf J_++z\mathbf j)|_{z=0}$.

There are $r_{\rm c}\in(0,\alpha_0/256]$, determined by the fixed radii $\alpha_0<\alpha_1$ of the notation and the constants of Proposition~\YMMref{51.10}{403}, and $C_{\rm mx}<\infty$, depending only on the constants of the construction, such that for $K>n\ge n_3$, $0\le k<K-n$ ($k=0$ included), $W\in\mathsf{Cut}_{k+1}$, any background $U_+$ over $W$, $\mathbf J_+=J^{(k+1)}(U_+)$ and $P_W=(U_+,0,\mathbf J_+)$:
\begin{enumerate}
\item[(a)] for every polymer $Y\in\mathcal D_k$ and complex $1$-forms $A'$, $\mathbf j$ on $\mathrm{Bd}(Y)$ with $|A'|,|\nabla_{U_+}A'|,|\mathbf j|<r_{\rm c}$ there, the triple $(U_+,A',\mathbf J_++\mathbf j)$ on $\mathrm{Bd}(Y)$ lies in the analyticity core $\mathcal U^\circ_{k+1}(Y;\alpha_0,\alpha_1)$ (Definition~\YMMref{49.1}{374}, Proposition~\YMMref{51.10}{403}, first assertion);
\item[(b)] $p_W\propto e^{\mathcal E(P_W;B)}dB$, where
\[
\mathcal E(P;B)=\sum_{Y\in\mathcal D_k}\bigl[-\tfrac12(B,\mathsf b_Y(P)B)+v_Y(P;B)\bigr]-\tfrac12\sum_{\square\in\pi^M_k}(B,\hat{\mathsf G}_\square(P)B),
\]
$\mathsf b_Y=x_kD_B^2\Psi^{\rm A}(Y;P,0)$ with $\Psi^{\rm A}$ of equations~\YMMref{(815)}{393} and~\YMMref{(816)}{393}, $v_Y$ the potentials of Theorem~\YMMref{W}{423} and of the fourth family at zero source (Proposition~\ref{prop:source-entry}), $\hat{\mathsf G}_\square$ of Section~\ref{sec:step-input-conventions}; the quadratic form, too, is a sum of local terms; and $(A',\mathbf j,B)\mapsto\mathcal E$ is holomorphic on the product of the ball of (a) with $\{|B(b)|<C_A\mathfrak b(x_k)\}$, $C_A\ge3$, which contains the closed box. The formula for $p_W$, the locality of its quadratic form and the holomorphy of $B\mapsto\mathcal E(P_W;B)$ on $\{|B(b)|<C_A\mathfrak b(x_k)\}$ hold as well for every integer $n\ge1$ with $L^n\mathsf b\ge\mathsf t_\natural$, so from $n_2$ on, with $K>n$, $0\le k<K-n$, at every $W\in\mathsf{Cut}_{k+1}$ and any background $U_+$ over $W$: they need of $W$ only that it lies in the class $\mathcal W_{k+1}(x_k;\alpha_0,\alpha_1)$ of the sum rule (Lemma~\YMMref{113.10}{796}), at which the restrictions of $P_W$ to the polymers lie in the analyticity cores;
\item[(c)] for $W$ at a depth $\ge n_{\rm a}$, with $\mathcal H^+$ the response of $U_+$ at level $k+1$ (Section~\ref{sec:background-step-complements}), $\mathbf U_\xi=e^{iL^{-k-1}\mathcal H^+(\xi)}U_+$, $W'_\xi=\mathcal M^{k+1}\mathbf U_\xi$ and $P_\xi=(U_+,\mathcal H^+(\xi),J^{(k+1)}(\mathbf U_\xi))$, there is $r_0>0$, which may depend on $W$ and $K$, such that for real $\xi$ with $\sup|\xi|<r_0$: $W'_\xi\in\mathsf{Cut}_{k+1}$ is a gauge transform of $e^{i\xi}W$, and $\mathbf U_\xi$ lies in the minimal orbits over $W'_\xi$ and over $\mathcal M^k\mathbf U_\xi$; there is a chart $\check V_\xi(B)$ of the fibre over $W'_\xi$, real-analytic in $(\xi,B)$ near the closed box, $\check V_0=V_W$, with values in a compact invariant subset of $\mathsf{Cut}_k$; for every bounded Borel invariant $H$, $h_{k+1}(e^{i\xi}W)=\int h_k(\check V_\xi(B))\,\check p_\xi(dB)$, where $\check p_\xi$ has the density $e^{\mathcal E(P_\xi;B)}/\int e^{\mathcal E(P_\xi;B')}dB'$, real-analytic in $\xi$ uniformly on the closed box, and $\check p_0=p_W$; the frame, $\check p_\xi$ and $\psi_\xi=\frac d{ds}\mathcal E(P_{s\xi};\cdot)|_{s=0}$ do not depend on the choice of the axial element;
\item[(d)] for real $\zeta$ on the box and every pair $(a,\mathbf j)$ of complex bare $1$-forms, with $\mu_y(a,\mathbf j)$ the largest of $|a|$, $|\nabla_{U_+}a|$, $|\mathbf j|$ at the bare bonds with an end point in the block of $y\in\Lambda_{K-k-1,\ell}$,
\[
\bigl|\partial_\zeta D_{(a,\mathbf j)}\mathcal E(P_W;B)\bigr|\le C^\sharp_{\rm mx}(x_k)\,\mathfrak b(x_k)\,\|\zeta\|\Bigl(\sum_y\mu_y(a,\mathbf j)^2\Bigr)^{1/2},
\]
with $C^\sharp_{\rm mx}$ depending on $x_k$ only through $\delta(x_k)\le1$; and for every real $\xi$ on $\Lambda_{K-k-1,\ell}$, $\sup_B|\partial_\zeta\partial_s\mathcal E(P_{s\xi};B)|_{s=0}|\le C_{\rm mx}\mathfrak b(x_k)\|\zeta\|\|\xi\|$, $C_{\rm mx}=\mathsf B\mathsf S\sup_{\mathsf t\ge\mathsf t_\natural}C^\sharp_{\rm mx}(\mathsf t)$.
\end{enumerate}
\end{proposition}

\begin{proof}[Proof sketch]
(a) Proposition~\YMMref{52.3}{433} holds for free radii $b_0,b_1\in(0,1]$, with the threshold $\varepsilon_{\mathcal W}b_0$ for $\varpi(W)$; the threshold that enters $n_3$ is the one at the radii $b_0=\frac12\alpha_0$, $b_1=\frac12\alpha_1$, that is, $2\theta_d\le\varepsilon_{\mathcal W}\alpha_0$ for $d\ge n_3$. So every $W\in\mathsf{Cut}_{k+1}$ with $K-k-1\ge n\ge n_3$ has $\varpi(W)<\theta_{K-k-1}\le\frac12\varepsilon_{\mathcal W}\alpha_0$, and by the plain-core verification in its proof $P_W$ lies in the core $\mathcal U^\circ_{k+1}(Y;\frac12\alpha_0,\frac12\alpha_1)$ with half the radii. The outer substitution from half the radii (Proposition~\YMMref{51.10}{403} (first assertion) at the radii $\alpha_0,\alpha_1$) carries this triple, moved by $(A',\mathbf j)$ with $|A'|,|\nabla_{U_+}A'|,|\mathbf j|<r_{\rm c}$, into the core with the radii $\alpha_0,\alpha_1$: the halved radii leave the room $r_{\rm c}$.

(b) The density in the chart comes from the sum identities in equations~\YMMref{(828)}{395}--\YMMref{(830)}{395} and their preceding Jacobian calculation and the step (Proposition~\ref{prop:source-entry}); the quadratic form is a sum of local terms, obtained by subtracting polynomials that agree on an open set. The formula for $p_W$, the locality of its quadratic form and the holomorphy of $B\mapsto\mathcal E(P_W;B)$ on $\{|B(b)|<C_A\mathfrak b(x_k)\}$ need of $W$ only that it lies in the class $\mathcal W_{k+1}(x_k;\alpha_0,\alpha_1)$ of the sum rule (Lemma~\YMMref{113.10}{796}), where the restrictions of $P_W$ to the polymers lie in the analyticity cores and the sum rules hold, so that every term of $\mathcal E(P_W;\cdot)$ is formed and holomorphic in $B$ there (Theorem~\YMMref{W}{423}, Proposition~\YMMref{51.4}{393} and its proof); so they hold for every integer $n\ge1$ with $L^n\mathsf b\ge\mathsf t_\natural$, from $n_2$ on in particular.

Only the joint holomorphy in $(A',\mathbf j)$ on the ball of (a) uses (a), and with it the condition $2\theta_d\le\varepsilon_{\mathcal W}\alpha_0$ of $n_3$. (c) Moving $W$ along the response keeps the background in the minimal orbit, so no control of the axial gauge is needed. With the residual gauge $h_\xi$ such that the background over $W'_\xi$ is $\mathbf U_\xi^{h_\xi}$, $\check V_\xi(B)$ is $V_{W'_\xi}(\mathrm{Ad}(h_\xi)B)$ carried back by $h_\xi$; its values at the pivots follow from the contraction there (the contraction leading to \YMMref{(800)}{391}, with the inverse of the one-step pivot derivative bounded by $2L^2$ in \YMMref{(798)}{390}) and the implicit function theorem. Every term of $\mathcal E$ is invariant under $(P,B)\mapsto(P^u,\mathrm{Ad}(u)B)$ and depends on $P=(U,A',\mathbf J)$ only through $(e^{iL^{-k-1}A'}U,\mathbf J)$ (Theorem~\YMMref{W}{423}, Definition~\YMMref{52.1}{415}, Proposition~\YMMref{51.4}{393} and its proof, Section~\ref{sec:step-input-conventions}, and the pair and gauge covariance of Proposition~\ref{prop:source-entry}), and $\mathrm{Ad}(h_\xi)$ keeps $dB$ and the box: so the density depends on $P_\xi$ alone, and no derivative of $h_\xi$ enters.

The values lie in a compact subset because the regularity number of the backgrounds stays strictly below the cut. (d) Each local term is bounded on the whole product of the core and the polydisc (Proposition~\YMMref{51.4}{393} and its proof, Theorem~\YMMref{W}{423}), for instance $|v_Y|\le32C_V\delta(x_k)(d(Y)+1)C_A^2\mathfrak b(x_k)^2e^{-\lambda d(Y)}$ with constants $C_V$, $\lambda$ of the step; Cauchy's estimate on a bidisc, the Cauchy--Schwarz inequality and the sum over polymers \eqref{eq:analytic-domain-tail} give the first bound. At $s=0$, the derivative of $P_{s\xi}$ near $y$ is at most $\mathsf B\,m_y(\xi)$ by Lemma~\ref{lem:response-current-zero}, which also gives $\sum_ym_y(\xi)^2\le\mathsf S^2\|\xi\|^2$.
\end{proof}

\begin{proposition}[One-step transport of the Dirichlet energy]\label{C-367-9}
For a level $j$ and a $\mathfrak g$-valued function $w$ of the bare bonds let $[w]_j^2=\sum_y(L^j\max\{|w(b)|:b_-\in B^j(y)\})^2$, $B^j(y)$ the cube of $L^{3j}$ bare sites centred at $y\in\Lambda_{K-j,\ell}$; then $L^{j/2}[w]_j\le L^{(j+1)/2}[w]_{j+1}$, $\|w\|^2\le3L^j[w]_j^2$, $[w]_0\le\|w\|$, and $[\mathrm{Ad}(g)w]_j=[w]_j$. For $V\in\mathsf{Cut}_j$ at a depth $\ge n_{\rm a}$, $U$ in the minimal orbit over $V$ and a real $1$-form $X$ on $\Lambda_{K-j,\ell}$, let $N^\flat_V(X)=L^{j/2}\min_\lambda[L^{-j}D\mathcal H_U(0)X-\mathrm d_U\lambda]_j$ for $j\ge1$, $N^\flat_V(X)=\min_\lambda[X-\mathrm d_V\lambda]_0$ for $j=0$, and $\Gamma^\flat_{K-j}(g)(V)=L^K\sup_{X:\,N^\flat_V(X)>0}\langle X,\partial g(V)\rangle^2/N^\flat_V(X)^2$.

$N^\flat_V$ is a seminorm that does not depend on the choice of $U$ and vanishes exactly on the gauge directions; $N^\flat_V(X)\le L^{j/2}\mathsf M\|X\|$, so $\Gamma_d\le\mathsf M^2\Gamma^\flat_d$; $\Gamma^\flat_K(H)\le3\Gamma_K(H)$ for invariant $H$; and $\Gamma^\flat_d(g)$ is Borel for invariant $g\in C^1$. There are $n_{\rm a}$ as in the notation and $C_\flat$, depending only on $L$, $C_Q$, $\mathsf M$ and the constants of Section~\ref{sec:background-step-complements}, such that, with $C_{\rm cov}=\frac4{\gamma_{\rm r}}C_{\rm ch}^{1/2}\mathsf M\sqrt6C_QL^{-1/2}C_{\rm mx}$, for $K>n\ge n_{\rm a}$, $0\le k<K-n$, $W\in\mathsf{Cut}_{k+1}$:
\begin{enumerate}
\item[(a)] every real $\xi'$ on $\Lambda_{K-k-1,\ell}$ has a representative $\xi$ modulo gauge at $W$, depending on $W$, $U_+$, $\xi'$ and not on $B$, with $\|\xi\|\le\sqrt6C_QL^{-(k+1)/2}N^\flat_W(\xi')$ and $N^\flat_{V_W(B)}(\check{\mathsf d}_k(W,B)\xi)\le(1+C_\flat\delta(x_k))N^\flat_W(\xi')$ for all $B$ in the box, $\check{\mathsf d}_k(W,B)\xi=\frac d{ds}\check V_{s\xi}(B)|_{s=0}$; more precisely, for every bare velocity $w$ of backgrounds at $U_+$ of level $k+1$ and $\xi=D\mathcal M^{k+1}(U_+)w$, $N^\flat_{V_W(B)}(\check{\mathsf d}_k(W,B)\xi)\le(1+C_\flat\delta(x_k))L^{(k+1)/2}[w]_{k+1}$;
\item[(b)] for real invariant $H\in C^1(\Omega_K)$ the tower has $h_k\in C^1(\mathsf{Cut}_k)$, and for $W\in\mathsf{Cut}_{k+1}$ and real $\xi$, with $\psi_\xi$ of Proposition~\ref{C-367-8}(c),
\[
\langle\xi,\partial h_{k+1}(W)\rangle=\int\langle\check{\mathsf d}_k(W,B)\xi,\partial h_k(V_W(B))\rangle\,p_W(dB)+\mathrm{Cov}_{p_W}(h_k\circ V_W,\psi_\xi),
\]
\[
\Gamma^\flat_{d-1}(h_{k+1})(W)^{1/2}\le(1+C_\flat\delta(x_k))\int\Gamma^\flat_d(h_k)^{1/2}\circ V_W\,dp_W+C_{\rm cov}\,\delta(x_k)\Bigl(\int\Gamma^\flat_d(h_k)\circ V_W\,dp_W\Bigr)^{1/2},
\]
so $\Gamma^\flat_{d-1}(h_{k+1})(W)\le(1+\epsilon^\flat_d)\int\Gamma^\flat_d(h_k)\circ V_W\,dp_W$ at every $W$; the same holds for $\phi_k=f_k^{1/2}$, $(f_k)$ the tower of $G^2+\epsilon$, $G$ real invariant, $\epsilon>0$;
\item[(c)] for $0\le k\le K-n$,
\begin{equation}\label{C-367-10}
\int\Gamma_{K-k}(h_k)\,d\sigma^K_{K-k}\le3\mathsf M^2\Pi^\flat_n\int\Gamma_K(H)\,d\sigma^K ,
\end{equation}
and, pointwise, $\Gamma_e(g_{t'})(W)\le3\mathsf M^2\Pi^\flat_e\int\Gamma_K(G)\,d\mathsf P_W$ for the law $\mathsf P_W$ of (B3) and every $W\in\mathsf{Cut}_{t'}$.
\end{enumerate}
\end{proposition}

\begin{proof}[Proof sketch]
(a) Three facts: the monotonicity of $L^{j/2}[w]_j$ in $j$, with constant one; $\|D\mathcal M^j(U)w\|\le\sqrt6C_Q[w]_j$ when the plaquettes of $U$ lie within $a_QL^{-2j}$ of $\mathbf 1$ (equation~\eqref{eq:plain-average-l2-derivative}); and 
\[
[L^{-j}D\mathcal H_V(\zeta)X]_j\le\mathsf B\mathsf S\|X\|
\]
 (Section~\ref{sec:background-step-complements}). The key is that the gauge compensator is bounded pointwise: every bare velocity $w=L^{-j}D\mathcal H_U(0)X-\mathrm d_U\lambda$ is also $L^{-j}D\mathcal H_U(0)X^w+\mathrm d_U\nu$, $X^w=D\mathcal M^j(U)w$, with
\[
|\nu(x)|\le(k_6+\tfrac32)\,\mathsf B\,m^j_y(X^w)+\tfrac32L^j\max_{b_-\in B^j(y)}|w(b)|\qquad(x\in B^j(y)),
\]
at the block centres by equation~\YMMref{(996)}{446}, with $k_6=4\kappa_{\rm av}$ as fixed in Section~\ref{sec:background-step-complements}, and inside a block along a path, by $|\lambda(x)|\le|\lambda(y)|+\frac32L^j\max_{B^j(y)}|\mathrm d_U\lambda|$. Lift the velocity of level $k+1$ to a bare velocity $w$; the fluctuation $B$ moves the background of level $k$ by $O(\delta(x_k))$, so there is a curve of backgrounds of level $k$ over $\check V_{s\xi}(B)$, up to a gauge, whose bare velocity differs from $w$ by at most $C_\flat\delta(x_k)L^{1/2}[w]_{k+1}$ in $[\cdot]_k$; the monotonicity then costs nothing. The minimum in $N^\flat_W(\xi')$ is attained at some $w$; $\xi=X^w$ is the representative, and $\|\xi\|\le\sqrt6C_Q[w]_{k+1}$.

The argument of (a) is deterministic. The properties of $N^\flat_V$: a curve of backgrounds whose block field moves with velocity $X$ has a bare velocity $w$ with $L^{j/2}[w]_j\ge N^\flat_V(X)$; $\Gamma^\flat_K\le3\Gamma_K$ because $\langle X,\partial H\rangle^2=\langle X-\mathrm d_V\lambda,\partial H\rangle^2\le\|X-\mathrm d_V\lambda\|^2|\partial H|^2\le3[X-\mathrm d_V\lambda]_0^2|\partial H|^2$ for invariant $H$ and every $\lambda$; upper semicontinuity of $(V,X)\mapsto N^\flat_V(X)$ gives the Borel property.

(b) Differentiate Proposition~\ref{C-367-8}(c) at $\xi=0$; domination holds on the closed box, as the moving chart takes its values in a compact subset of $\mathsf{Cut}_k$. The left side depends on $\xi$ only modulo gauge: take the representative of (a); the first term is then at most $L^{-K/2}(1+C_\flat\delta(x_k))N^\flat_W(\xi')\int\Gamma^\flat_d(h_k)^{1/2}dp_W$. For the plain tower, the Poincar\'e inequality of $p_W$ used twice gives $|\mathrm{Cov}_p(f,\psi)|\le\frac4{\gamma_{\rm r}}\|\nabla_Bf\|_{L^2(p)}\sup|\nabla_B\psi|$. For the square roots, the log-Sobolev inequality of $p_W$ with $\kappa=\frac14\gamma_{\rm r}$ gives $\log\int e^{s(\psi-\int\psi)}dp\le s^2l^2/2\kappa$ when $|\nabla_B\psi|\le l$, and the variational formula of the entropy at $\phi^2$ and $s(\psi-\int\psi)$, optimised in $s$, gives
\[
|\mathrm{Cov}_p(\phi^2,\psi)|\le\frac{2l}\kappa\,\|\phi\|_{L^2(p)}\|\nabla_B\phi\|_{L^2(p)},
\]
with equality for a Gaussian $p$, $\phi=e^{ax/2}$, $\psi=x$. By Proposition~\ref{C-367-8}(d), $\sup|\nabla_B\psi_\xi|\le C_{\rm mx}\mathfrak b(x_k)\|\xi\|$; the chart factor $x_k^{-1/2}$ turns $\mathfrak b(x_k)$ into $\delta(x_k)$, and $L^{K/2}L^{-d/2}L^{-(k+1)/2}=L^{-1/2}$. Directions with $N^\flat_W(\xi')=0$ contribute zero by gauge invariance; for the other directions divide by $L^{-K/2}N^\flat_W(\xi')$ (for the square roots by $2\phi_{k+1}(W)=2\|\phi_k\circ V_W\|_{L^2(p_W)}$) and use Jensen's inequality for the probability $p_W$.

(c) Integrate the one-step bound against $\sigma^K_{d-1}$ (Lemma~\ref{C-367-7}(ii)) and iterate, using $\Gamma_d\le\mathsf M^2\Gamma^\flat_d$ at the top and $\Gamma^\flat_K\le3\Gamma_K$ at the bottom; for the pointwise form iterate at the fixed block field along the kernels. The constant does not depend on $K$ because $\sum_d\delta(L^d\mathsf b)<\infty$.
\end{proof}

\begin{proposition}[Transport in the $\ell^2$ metric]\label{C-367-19}
Assume (E). For $V\in\mathsf{Cut}_j$ with $\varpi(V)$ below the radius of Section~\ref{sec:background-step-complements}, $U$ in the minimal orbit over $V$ and $\mathsf Z=L^{-j}D\mathcal H_U(0)$, let $N_V(\xi)=\min_\lambda\|\mathsf Z\xi-\mathrm d_U\lambda\|$ ($N_V(\xi)=\min_\lambda\|\xi-\mathrm d_V\lambda\|$ at $j=0$). $N_V$ is a seminorm that does not depend on the choice of $U$ and vanishes exactly on the gauge directions; $N_V(\xi)^2\le\|\mathsf Z\xi\|^2\le C_{\rm rs}L^j\|\xi\|^2$, with $C_{\rm rs}$ depending only on $L$ and the constants of Section~\ref{sec:background-step-complements}. Every real $\xi$ is $\xi_*+\mathrm d_V\phi$ with exactly one horizontal datum $\xi_*$, that is, with $\mathrm d_U^{\mathsf T}\mathsf Z\xi_*=0$; on horizontal data $N_V(\xi_*)=\|\mathsf Z\xi_*\|$; every real $\xi$ has $\|\xi\|\le C_{\mathcal Q}L^{-j/2}\|\mathsf Z\xi\|$, so $\|\xi_*\|\le C_{\mathcal Q}L^{-j/2}N_V(\xi_*)$, with $C_{\mathcal Q}=34(L^3+2)$ as in Lemma~\ref{lem:constant-frame-row} (a number other than the $C_Q$ and $C'_Q$ of Lemma~\ref{lem:plain-average-remainder}); and the two metrics are equivalent: $3^{-1/2}N_V\le N^\flat_V\le\mathsf MC_{\mathcal Q}N_V$.

There are $C_\sharp$ and a depth $n_4\ge n_3$, depending only on $\mathsf b$ and the constants of the construction, such that for $K-k>n_4$, $W\in\mathsf{Cut}_{k+1}$ and $U_+$ a background over $W$ the following hold. All clauses of Proposition~\ref{C-367-8}(c) hold at $W$, with the frame $\check V_\xi(B)$ real-analytic in $(\xi,B)$ on $\{\sup|\xi|<r_0\}$ times the open box and with first derivatives in $\xi$ that are bounded there. For every horizontal datum $\xi$ at $W$, that is, every real $\xi$ on $\Lambda_{K-k-1,\ell}$ with $\mathrm d_{U_+}^{\mathsf T}\mathsf Z_{k+1,U_+}\xi=0$, $\mathsf Z_{k+1,U_+}=L^{-k-1}D\mathcal H^+(0)$, and all $B$ in the box, $s\mapsto\check V_{s\xi}(B)$ is defined and real-analytic near $0$ and
\[
N_{V_W(B)}\bigl(\check{\mathsf d}_k(W,B)\xi\bigr)\le\bigl(1+C_\sharp\,\delta(L^{K-k}\mathsf b)\bigr)N_W(\xi).
\]
Consequently Proposition~\ref{C-367-9}(b), (c) hold as well in the metric $N$ for $K>n\ge n_4$, (B3) for $K>e\ge n\ge n_4$, and (B2) in (a)--(c) for $n\ge\max\{n_4,n_{\rm c}\}$, in (d), under \textup{(Multiscale gauge)}, for $n\ge\max\{n_4,n_{\rm e}\}$ with $\mathsf R_n<1$, and in (e) for $n\ge n_4$, with $C_{\rm rs}\Pi_n$, $\Pi_n=\prod_{d>n}(1+\epsilon_d)$, $1+\epsilon_d=[1+(C_\sharp+\frac4{\gamma_{\rm r}}(C_{\rm ch}C_{\rm rs})^{1/2}C_{\mathcal Q}L^{-1/2}C_{\rm mx})\delta(L^d\mathsf b)]^2$, in the place of $3\mathsf M^2\Pi^\flat_n$ (in (B3): $C_{\rm rs}\Pi_e$ in the place of $3\mathsf M^2\Pi^\flat_e$).
\end{proposition}

\begin{proof}[Proof sketch]
\emph{The range of $\mathsf Z$.} The response is in the Landau gauge on the kernel of the site averages: by equation~\YMMref{(837)}{398}, $\mathsf R^{\rm La}D^{\mathsf T}\mathcal H=0$, $D=\mathrm d_U$, $\mathsf R^{\rm La}$ the orthogonal projection onto $D^{\mathsf T}D\mathcal N$ and $\mathcal N$ the site functions whose site averages $Q'_q$ vanish at the tops (Section~\ref{sec:response-step-conventions}); so $\mathsf R^{\rm La}\mathrm d_U^{\mathsf T}\mathsf Z\xi=0$ for every $\xi$. The response is covariant, as the averaging map is, so $\mathsf Z$ carries the gauge directions $\mathrm d_V\psi$ of the block field to gauge directions at $U$; these meet $\mathrm d_U\mathcal N$ only in $0$ (a $\mathrm d_U\mu$, $\mu\in\mathcal N$, in the range of $\mathsf Z$ has $\mathsf R^{\rm La}\mathrm d_U^{\mathsf T}\mathrm d_U\mu=0$ while $\mathrm d_U^{\mathsf T}\mathrm d_U\mu$ lies in the range of $\mathsf R^{\rm La}$, so $\mathrm d_U\mu=0$), and by the count of dimensions the two together span all the gauge directions at $U$.

A horizontal velocity of backgrounds at $U$, $h=\mathsf ZX-\mathrm d_U\lambda$ with $\mathrm d_U^{\mathsf T}h=0$, is then $\mathsf ZX'-\mathrm d_U\nu$ with $\nu\in\mathcal N$, and $\mathsf R^{\rm La}\mathrm d_U^{\mathsf T}\mathrm d_U\nu=\mathsf R^{\rm La}\mathrm d_U^{\mathsf T}\mathsf ZX'=0$: so $\mathrm d_U\nu=0$ and $h=\mathsf ZX'$. The range of $\mathsf Z$ holds every horizontal velocity of backgrounds, with no compensator.

\emph{The decomposition.} The part $h$ of $\mathsf Z\xi$ orthogonal to the range of $\mathrm d_U$ is a horizontal velocity of backgrounds, $h=\mathsf Z\xi_*$; $\mathsf Z(\xi-\xi_*)$ is a gauge direction at $U$, so $\xi-\xi_*$ is a gauge direction at $V$; $\xi_*$ is unique, as $\mathsf Z$ is injective (below) and a horizontal vector that is a gauge direction vanishes.

\emph{The bound.} Lemma~\ref{lem:constant-frame-row} applies in the constant frame, with bare spacing $L^{-j}$ and terminal spacing one. It gives $\mathcal Q(\mathcal H_U(\zeta))=\zeta$, and hence $L^jK_0\mathsf Z\xi=\xi$, together with $\|K_0\|\le2L^{-3j/2}\le C_{\mathcal Q}L^{-3j/2}$ on counting-measure $\ell^2$. Thus $\|\xi\|\le C_{\mathcal Q}L^{-j/2}\|\mathsf Z\xi\|$; on horizontal data $\|\mathsf Z\xi_*\|=N_V(\xi_*)$. The equivalence follows from $\|w\|^2\le3L^j[w]_j^2$ and $N^\flat_V(\xi)=N^\flat_V(\xi_*)\le L^{j/2}\mathsf M\|\xi_*\|$.

\emph{The estimate.} $U_+$ is a background of the level $k$ as well, and the same argument at the level $k$ shows that the velocity of a horizontal datum $\xi$ at $W$ is a response one level finer, with no compensator: $\mathsf Z_{k+1,U_+}\xi=\mathsf Z_{k,U_+}\zeta$, $\zeta=D\mathcal M^k(U_+)\mathsf Z_{k+1,U_+}\xi$; the fluctuation moves the background of level $k$ by $O(\delta(x_k))$, which costs the factor $1+C_\sharp\delta$.

\emph{The frame from the depth $n_4$ on.} As $n_4\ge n_3$, $\theta_d$ lies below the radius of Section~\ref{sec:background-step-complements}; $U_+$ is a reference of the levels $k+1$ and $k$; for small real $\xi$, $W'_\xi$ is a gauge transform of $e^{i\xi}W$ in $\mathsf{Cut}_{k+1}$, $\mathbf U_\xi$ lies in the minimal orbits over $W'_\xi$ and over $\mathcal M^k\mathbf U_\xi$, and $\xi\mapsto\mathcal M^k\mathbf U_\xi$, $\xi\mapsto W'_\xi$ are real-analytic. The value of $\mathcal M(V)$ at a bond $c$ of $\Lambda_{K-k-1,\ell}$ reads $V$ on the bonds of the double block of $c$ alone, of the pivot bonds the pivot bond of $c$ alone, and its derivative in the value at that pivot bond is invertible, with an inverse bounded by $2L^2$, on the compact closure of the set of the restrictions of the $V_W(B)$, $B$ in the box.

The implicit function theorem, with radii that are uniform over this compact set, gives about every base point $V_W(B)$ exactly one solution at each pivot near the value of the base point, real-analytic in the data, with bounded first derivatives. The frame $\check V_\xi(B)$ is $e^{ix_k^{-1/2}B(b)}(\mathcal M^k\mathbf U_\xi)(b)$ at the bonds that are not pivot bonds and this solution at the pivots; by the uniqueness at the pivots and continuation in $B$ along $\tau B$, $\tau\in[0,1]$, from $B=0$, it is $V_{W'_\xi}(\mathrm{Ad}(h_\xi)B)$ carried back by $h_\xi$. Hence, as in Proposition~\ref{C-367-8}(c): the tower identity; the density $e^{\mathcal E(P_\xi;B)}/\int e^{\mathcal E(P_\xi;B')}dB'$ of $\check p_\xi$, which reads the depth only through $n_3$; the values in a compact invariant subset of $\mathsf{Cut}_k$; analyticity on small $\xi$ times the open box, with $\xi$-derivatives bounded there, which carries the domination in the differentiation of Proposition~\ref{C-367-9}(b); and at the depths $\ge n_{\rm a}$ the frame is that of Proposition~\ref{C-367-8}(c).

In the differentiated step the representative is used only through $\|\xi\|\le c\,L^{-(k+1)/2}N_W(\xi)$, which the horizontal datum of the class satisfies with $c=C_{\mathcal Q}$; the rest of the proof of Proposition~\ref{C-367-9} is unchanged, with the square field $\Gamma^N$ built from $N$ as $\Gamma^\flat$ is built from $N^\flat$, $N_V(X)\le C_{\rm rs}^{1/2}L^{j/2}\|X\|$ at the top and $\Gamma^N_K\le\Gamma_K$ at the bottom.
\end{proof}

\begin{lemma}[Connected regions at the coarse scale]\label{C-367-11}
\begin{enumerate}
\item[(a)] Let $\mathsf n\ge\lambda_{{\rm BG},3}$, $0<\kappa'\le\kappa_0$, $\mathsf n\ge\kappa'^{-1/2}$, and $V$ a configuration on $(\mathbb Z/\mathsf n)^3$ with $\varpi(V)\le\kappa'$. Some gauge has $|V^u(b)-\mathbf 1|\le C_{\rm BG}\sqrt{\kappa'}$ at every bond; $\{\varpi<\kappa'\}$ lies in the path component of $\mathbf 1$ in $\{\varpi<C''\kappa'\}$; so, at $\kappa'=r/C''$, for $0<r\le C''\kappa_0$ and $\mathsf n\ge\max\{\lambda_{{\rm BG},3},(C''/r)^{1/2}\}$, every other component of $\{\varpi<r\}$ lies in the shell $\{r/C''\le\varpi<r\}$.

The estimate involving $J$ block levels used below is the additional hypothesis \textup{(Multiscale gauge)}. Its constants are finite constants in that hypothesis; no numerical values for them are asserted here.
\item[(b)] For a finite set $E$ of bonds, $0<\alpha<\pi$, $\mathsf K\subset\{A\in\mathfrak g^E:|A(b)|<\alpha\}$ open and convex and $V_0\in G^E$, $A\mapsto e^{iA}V_0$ maps $\mathsf K$ diffeomorphically onto its image $\mathsf B$; there Haar measure is $\prod_b(2\pi^2)^{-1}\mathsf j(A(b))\,dA(b)$, $\mathsf j(X)=(\sin|X|/|X|)^2$; $-\log\prod_b\mathsf j$ has Hessian $\ge\frac23$; and $\mathrm{Ent}_{dV|\mathsf B}(\psi^2)\le3\int_{\mathsf B}|\partial\psi|^2dV$ for bounded locally Lipschitz $\psi$, whatever the number of bonds.
\item[(c)] If $\nu|_{S_i}$ have log-Sobolev constants $c_i$ on open $S_1,S_2$ with $\nu(S_1\cap S_2)>0$, then $\nu|_{S_1\cup S_2}$ has $2\max_ic_i(1+16\nu(S_1\cup S_2)/\nu(S_1\cap S_2))$. If Haar measure on an open $S'$ has the constant $c$, $\rho(V)=\int\mathbf 1_{S'}(V^g)dg$, $\rho_0>0$ and $w=\min\{1,\rho/\rho_0\}$, then $\mathrm{Ent}_{w\mathsf g\,dV}(\phi^2)\le\frac{c\,\mathsf g_+}{\rho_0\mathsf g_-}\int|\partial\phi|^2w\,\mathsf g\,dV$ for invariant $\phi$ and Borel $0<\mathsf g_-\le\mathsf g\le\mathsf g_+$, whatever the boundary of $S'$. For an open invariant $\mathcal O$, a component $\mathsf C$ of it and a compact invariant $\mathcal K\subset\mathsf C$ there is a connected finite union $S'$ of sets $\{e^{iA}V_i:|A(b)|<\delta_i\}$ with $\mathcal K\subset S'\subset\mathsf C$.
\item[(d)] The set $\mathcal N\subset(0,2)$ of the $r=\varpi(V)$ at which $\sum_pc_p\partial a_p(V)=0$ for some $c\ge0$, $c\ne0$, carried by the plaquettes with $\varpi_p(V)=r$, is Lebesgue-null; for $r\notin\mathcal N$, $\{\varpi<r\}$ has finitely many components, each a connected open set with a Lipschitz boundary. The abelian fields whose $(1,2)$ plaquettes all equal $e^{if\sigma_3}$, $f=2\pi m/\mathsf n^2$, and whose other plaquettes are $\mathbf 1$, are such critical points: $2\sin(\pi m/\mathsf n^2)\in\mathcal N$ for integers $0<m<\mathsf n^2/2$.
\item[(e)] For $0<r<2$, every connected component of $\{\varpi<r\}$ is open, invariant, and contains a local minimiser of $\varpi$. For $\phi_1\in C^2[0,c)$, $c=\frac12r^2$, with $\phi_1'>0$ and $\phi_1\to\infty$ at $c$, the function $\Phi=\sum_p\phi_1(a_p)$ has compact sublevel sets in $\{\varpi<r\}$, every component contains a point where $\Phi$ restricted to it is least, and there $\sum_p\phi_1'(a_p)\,\partial a_p=0$ with weights $\phi_1'(a_p)>0$ at all plaquettes and $\nabla^2\Phi\ge0$; $\{\varpi<r\}$ is connected if and only if every local minimiser of $\Phi$ lies in the component of $\mathbf 1$; and $\{\varpi<r\}$ is the increasing union of the sets $\{\sum_pa_p^q<c^q\}$, $q\to\infty$.
\item[(f)] Let $0<\alpha\le\frac1{32}$, $r'=64\alpha^2$, $N$ the number of sites along an axis, and $\mathscr D=\{A:|A(b)|<\alpha\ \forall b,\ |Y_p(A)|<r'\ \forall p\}$, $e^{iY_p(A)}=\partial(e^{iA})(p)$, $|Y_p|\le\pi$. Then $\mathscr D$ is star-shaped, $\mathscr D=\{s\omega:|\omega|=1,\ 0\le s<\varrho(\omega)\}$ with $\frac12r'\le\varrho\le\rho_{\max}=\alpha\sqrt{3N^3}$ and $|\varrho(\omega)-\varrho(\omega')|\le\frac{4\pi\rho_{\max}^2}{r'}|\omega-\omega'|$; $A\mapsto e^{iA}$ maps $\mathscr D$ diffeomorphically onto $\exp\mathscr D=\{V:\mathsf d(V(b),\mathbf 1)<\alpha,\ \mathsf d(\partial V(p),\mathbf 1)<r'\}$, and $\mathrm{Ent}_{\mathbf 1_{\exp\mathscr D}dV}(\phi^2)\le C_{\mathscr D}\int_{\exp\mathscr D}|\partial\phi|^2dV$ for bounded locally Lipschitz $\phi$, with $C_{\mathscr D}$ of (B1)(e) at $\alpha_n=\alpha$.

With $w(V)=\int\mathbf 1\{\mathsf d(V^u(b),\mathbf 1)<\alpha\ \forall b\}\,du$ and $\mathsf{In}=\{V:\mathsf d(\partial V(p),\mathbf 1)<r'\ \forall p\}$: $\int_{\mathsf{In}}f\,w\,dV=\int_{\exp\mathscr D}f\,dV$ for invariant Borel $f\ge0$; $w(V)\ge[(\frac12\alpha-\sin\frac12\alpha)/2\pi]^{N^3}$ if some gauge brings all bonds of $V$ within $\frac14\alpha$ of $\mathbf 1$; and $\mathsf{In}\cap\{w>0\}$ is path connected.
\end{enumerate}
\end{lemma}

\begin{proof}[Proof sketch]
(a) The gauge is built cell by cell on a periodic grid of mesh $\kappa'^{-1/2}$, along the proof of Proposition~\ref{prop:uniform-box-gauge}; along the ray to $0$ in exponential coordinates, $\bigl||\partial e^{itA}(p)-\mathbf 1|-|t\,dA(p)|\bigr|\le10\alpha^2$, with $dA(p)=A(b_1)+A(b_2)-A(b_3)-A(b_4)$ the plain sum of the bond function $A$ around $p$, its bonds $b_1,\dots,b_4$ taken with the signs $(+,+,-,-)$ of Lemma~\ref{C-367-13}, when $|A|\le\alpha\le\frac18$, which keeps the ray in $\{\varpi<C''\kappa'\}$. The multilevel estimate in \textup{(Multiscale gauge)} is a separate input; this single-scale construction does not prove it.

For $\mathrm U(1)$ the flux sectors are distinct components. (b) The radial function $h(\rho)=-2\log(\sin\rho/\rho)$ on $\mathbb R^3$ has
\[
h''(\rho)=2\sum_{m\ge1}\bigl[(m\pi-\rho)^{-2}+(m\pi+\rho)^{-2}\bigr]\ge4\sum_{m\ge1}(m\pi)^{-2}=\tfrac23,\qquad\frac{h'(\rho)}\rho=4\sum_{m\ge1}(m^2\pi^2-\rho^2)^{-1}\ge\tfrac23 ;
\]
$|\nabla(\psi\circ\exp)|\le|\partial\psi|$ since $\frac d{ds}e^{i(X+sY)}=iZe^{iX}$ with $Z=\int_0^1\mathrm{Ad}(e^{iuX})Y\,du$, $|Z|\le|Y|$; then the log-Sobolev inequality on a convex set, and locally Lipschitz $\psi$ by the dilations $A_0+\vartheta(A-A_0)$, $\vartheta\to1$. (c) The first bound by comparing means on the overlap $S_1\cap S_2$; the second by applying the inequality on $S'$ in each orbit through the orbit average $\rho$, capped at $\rho_0$, and the bounded change of density $\mathsf g$; the cover of $\mathcal K$ by chart balls inside $\mathsf C$ is joined along paths in $\mathsf C$. (d) Sard's theorem on the strata indexed by the set of the maximal plaquettes; off $\mathcal N$ Gordan's alternative gives a common descent direction of all maximal plaquettes, which makes the boundary locally a Lipschitz graph.

At the abelian flux fields $\sum_p\partial a_p$ over the $(1,2)$ plaquettes vanishes. (e) Compactness of sublevel sets in the open set; a component is closed in $\{\varpi<r\}$, so a minimising sequence in it has its limit points in it. (f) Along rays the plaquette angle grows: $\frac d{dt}|Y(t)|\ge|Y|-32\alpha^2$, which gives the star shape and the Lipschitz radial function; the log-Sobolev inequality on a star-shaped domain with Lipschitz radial function follows by mapping it to a ball, at the cost of the Jacobian and of the Lipschitz constant, which is $C_{\mathscr D}$. The identity is Fubini over the gauge group; the lower bound counts, at each site, the gauge transformations within $\frac14\alpha$ of the given one, a ball of normalised volume $(\frac12\alpha-\sin\frac12\alpha)/2\pi$.
\end{proof}

\subsection{Gauge-fixed regularity and the admissibility hypothesis}

\begin{proposition}[Gauge-fixed fields and the gradient flow]\label{C-367-12}
Let $0<\kappa<\frac16$, and let $\mathsf N'(U)$ be the least $\|\mathbf B\|_{\kappa;a}$ over the $1$-forms $\mathbf B$ with $e^{ia_K\mathbf B}$ in the orbit of $U$. \emph{The gauge-fixed flow} is the lattice gradient flow of the action with a gauge term: $V_s=e^{ia_KA_s}$, $A_0=\mathbf B$, moves with the velocity of the Wilson flow $\Phi$ plus a gauge direction $\mathrm d_{V_s}\lambda_s$, $\lambda_s$ the lattice form of the divergence of $A_s$ that answers to the member $\mathrm d_{\mathfrak A}(\partial_\nu\mathfrak A_\nu)$ of the flow lines, fixed so that the linear part of the equation is the lattice heat operator. In the chart it reads $\frac d{ds}A_s(b)=a_K^{-3}\Theta^{\rm ch}(a_KA_s|_{\mathsf N(b)})$, $\mathsf N(b)$ the fifteen bonds near $b$, $\Theta^{\rm ch}$ holomorphic and bounded on $\{|a_KA|\le1\}$, and $\mathsf M_\Theta=\max\{1,\sup|\Theta^{\rm ch}|\}$, the supremum taken there, an absolute number; the members of $\Theta^{\rm ch}$ of the orders one, two and at least three in $A$ give
\[
\dot A=\Delta_aA+\mathcal N(A,A)+\mathcal R(A),\qquad\mathcal N(A,A')=\sum_jc_j\,(\tau_{h_j}\nabla_{\lambda_j}A_{\nu_j})\wedge(\tau_{h'_j}A'_{\nu'_j}),
\]
a finite sum, with absolute coefficients, over the directions and the fixed finite translation set used in the rough regularity quantity.

\emph{Fields in the chain gauge.} Assume (E); let $d_0\ge n_2$ be a top depth, $K>d_0$, $t=K-d_0$, $\mathbb E^{[d_0]}$ the normalised chain measure from the top depth $d_0$, and $V_k=\mathcal M^kU$. On a set of full measure: $U$ is a gauge transform of $\mathsf U_0$, where $\mathsf U_t=U_{\rm top}$ is a background over the top field $V_t$ and $\mathsf U_k=e^{ia_K\mathbf h_{K-k}}\mathsf U_{k+1}$ ($k<t$) is a gauge transform of the background over $V_k$, in a gauge that is a function of $V_{k+1}$ (the gauge of the chain); the layer $\mathbf h_d$ of the depth $d=K-k$ is a function of the fluctuation $B_k$ of the chart at $V_{k+1}$, with $|\mathbf h_d|\le c_{\mathsf h}L^d\delta(L^d\mathsf b)$, $c_{\mathsf h}$ a constant of the construction; its linear part in the fluctuation, carried by $\mathrm{Ad}(\gamma)$ below, is $\mathsf b^{-1/2}L^{d/2}\mathsf J_dB_k$, with the layer map $\mathsf J_d$, a function of $V_{k+1}$, of the norms $\|\mathsf J_dX\|_{L^2(T_\ell)}\le c_{\mathsf J}L^{-3d/2}\|X\|$ and $\|\mathsf J_dX\|_\infty\le c_{\mathsf J}\max|X|$, $c_{\mathsf J}$ a constant of the construction, and at every point of the box the derivative of $\mathbf h_d$ in $B_k$ is $\mathsf b^{-1/2}L^{d/2}$ times a map with the same two bounds; in its regular gauge $\gamma$ the top is $U_{\rm top}^\gamma=e^{ia_K\mathbf A^{\rm top}}$ with $|\mathbf A^{\rm top}|\le\mathsf a_{\rm top}\le CL^{10d_0}$, $\mathsf a_{\rm top}$ a number that reads neither $K$ nor the field; and $U$ is a gauge transform of $e^{ia_K\mathbf B_K}$,
\[
\mathbf B_K=\tfrac1{ia_K}\log\bigl(e^{ia_K\mathbf h'_K}\cdots e^{ia_K\mathbf h'_{d_0+1}}e^{ia_K\mathbf A^{\rm top}}\bigr),\qquad\mathbf h'_d=\mathrm{Ad}(\gamma)\mathbf h_d ,
\]
bond by bond, the logarithm being formed by its series if $\sum_da_K|\mathbf h_d|\le\frac14$ and $a_K\mathsf a_{\rm top}\le\frac18$, and otherwise by a fixed Borel $\mathfrak{su}(2)$ logarithm of norm at most $\pi$, principal off $-\mathbf1$, on positive bonds and extended by reversal. Admissibility and the conclusion \textup{(Adm)} have the meaning stated before Theorem~\ref{C-367-5}; the layer fields and their conditional estimates are specified in the proof below.

\begin{enumerate}
\item[(a)] There are $c_1,c_2\in(0,1]$, $c_{\rm Lip}\in(0,c_1]$ and $c_3,c_4,c_5$, depending only on $\kappa$, such that for all $K$, $m$ and $\beta>0$, every real $1$-form $\mathbf B$ on $\Lambda_{K,\ell}$ and $R\ge\max\{1,\|\mathbf B\|_{\kappa;a}\}$, with $s_0(R)=(c_1R^{-2})^{4/(1-2\kappa)}$, $s_{\rm Lip}(R)=(c_{\rm Lip}R^{-2})^{4/(1-2\kappa)}$ and $a_0(R)=(c_2R^{-1})^{2/(1-2\kappa)}$, the following hold. If $a_K\le a_0(R)$, the gauge-fixed flow $A_s$ from $A_0=\mathbf B$ exists on the closed interval $[0,s_0(R)]$, $\Phi_s(e^{ia_K\mathbf B})$ is gauge equivalent to $e^{ia_KA_s}$, and for $0<s\le s_0(R)$: $\|A_s-P_s\mathbf B\|_\infty\le c_3Rs^{-\kappa}$, $\|\nabla(A_s-P_s\mathbf B)\|_\infty\le c_3Rs^{-\frac12-\kappa}$, $\|A_s\|_\infty\le(1+c_3)Rs^{-\frac14-\frac\kappa2}$, $\|\nabla A_s\|_\infty\le(2+c_3)Rs^{-\frac34-\frac\kappa2}$.

If besides $\mathbf B'$ is a real $1$-form with $\|\mathbf B'\|_{\kappa;a}\le R$, $A'_s$ its gauge-fixed flow and $\mathsf D$ the rough distance of $\mathbf B$ and $\mathbf B'$, then for $0<s\le s_{\rm Lip}(R)$: $\|(A_s-A'_s)-P_s(\mathbf B-\mathbf B')\|_\infty\le c_5\mathsf Ds^{-\kappa}$, $\|\nabla[(A_s-A'_s)-P_s(\mathbf B-\mathbf B')]\|_\infty\le c_5\mathsf Ds^{-\frac12-\kappa}$, so $\|A_s-A'_s\|_\infty\le(1+c_5)\mathsf Ds^{-\frac14-\frac\kappa2}$. For every $a_K$: $A^a(\Phi_se^{ia_K\mathbf B})\le c_4\ell^3(1+R)^{24}s^{-\frac32-\kappa}$ for $0<s\le1$. These statements, the short-time theory of the gauge-fixed flow, are deterministic and rest on its equation alone; they hold on every lattice $\Lambda_{K,\ell}$, uniformly in $a_K$.

Assume (E): for every admissible $d_0$, $K>d_0$ and $q<\infty$, $\mathbb E^{[d_0]}[\|\mathbf B_K\|^q_{\kappa;a}]^{1/q}\le C(q,\kappa)L^{20d_0}$, with $C(q,\kappa)$ free of $K$ and $d_0$; and for every top depth $d_1\ge n_2$, $K>d_1$, invariant Borel $F\ge0$ and $r>1$,
\[
\langle F\rangle_{\mu^K}\le\mathbb E^{[d_1]}[F]+\sum_{d_1\le d\le K-2}\mathsf R_d^{1-1/r}\,\mathbb E^{[d+1]}[F^r]^{1/r}+\mathsf R_{K-1}\sup F ,
\]
where $\sum_d\mathsf R_d^{1/2}L^{pd}$ and $\sup_K\mathsf R_{K-1}L^{pK}$ are finite for every $p$. Under (E), the conclusion (Adm) proved below gives $\mathsf M_q=\mathsf M_q(\kappa)=\sup_K\langle\mathsf N'^q\rangle_{\mu^K}<\infty$ for every $q<\infty$, and: for every $\epsilon>0$ there are $\bar s\in(0,1]$, $\bar C$, $\bar K$ with $\mu^K(\mathsf O_K(\bar C,\bar s))\ge1-\epsilon$ for $K\ge\bar K$; $\sup_K\langle(\sup_{0<s\le1}s^{\frac32+\kappa}A^a(U_s))^q\rangle_{\mu^K}<\infty$ for every $q$; for $0<\sigma\le1$ and $1\le p<\infty$ there is $C_3$, bounded for $\sigma$ in compact subsets of $(0,1]$, with $\mu^K(\max_pa_p(U_\sigma)>10^{-4})\le C_3a_K^p$ and $\int[(a_K^{-4}\max_pa_p(U_\sigma))^p+(J,J)_a(U_\sigma)^p]\,d\mu^K\le C_3$ for all $K$; and for $0<\sigma\le1$ the laws of $\Lambda_\sigma$ under $\mu^K$ are tight.
\item[(b)] Assume (E) and \textup{(Flow approximation)}. For every closed path $C$ of bonds of some $\Lambda_{d,\ell}$ and $s>0$, the flowed Wilson loop $W_C\circ\Phi_s$ ($W_C$ the half trace of the holonomy) converges strongly along the full sequence of cutoffs; there are an invariant Borel set of threads of full measure and a Borel map $\mathfrak F$ from it into flow lines, $\mathfrak F(Y)_s$ represented by a $C^{1,1}$ connection for $s>0$, such that the limit is $W_C(\mathfrak F(Y)_s)$. For every flowed observable and $\sigma>0$, with $\Xi^\sigma_{K,D}(V)=H^\sigma_K(U)$ for $V\in\{\varpi<\theta_D\}\subset\Omega_D$, $U$ any element of the minimal orbit over $V$ in $\Omega_K$, and $\Xi^\sigma_{K,D}=0$ elsewhere:
\[
\lim_{D\to\infty}\sup_{K>D}\|H^\sigma_K-\Xi^\sigma_{K,D}\circ\pi_{K,D}\|_{L^2(\mu^K)}=0,\qquad\lim_{D\to\infty}\limsup_{K\to\infty}\|\Xi^\sigma_{K,D}(Y_D)-\mathsf G_\sigma\|_{L^2(\mu)}=0 :
\]
the flowed observable at the background of the block field of a large depth.
\item[(c)] Assume (E) and \textup{(Flow approximation)}. For $s>0$ and Lipschitz $f$ the joint laws of $\mathcal O(f)\circ\Phi_s$ and of the $W_C\circ\Phi_s$ under $\mu^K$ converge with all moments along the full sequence to the joint law of $\mathcal O^0(f)(\mathfrak F(Y)_s)$ and the $W_C(\mathfrak F(Y)_s)$ under $\mu$. Separately, under (E) alone, for $\epsilon,\eta>0$ there are Borel sets $\mathsf G_K$ with $\mu^K(\mathsf G_K)\ge1-\epsilon$, closed paths $C_1,\dots,C_q$ and a continuous $\Theta$ on $[-1,1]^q$ with $|\mathcal O(f)(U_s)-\Theta((W_{C_i}(U_s))_i)|\le\eta$ on $\mathsf G_K$ for all large $K$.
\end{enumerate}
\end{proposition}

\begin{proof}[Proof sketch]
(a)

\emph{The gauge-fixed flow.} Write $A_s=P_s\mathbf B+w_s$; the mild form $w_s=\int_0^sP_{s-r}[\mathcal N(A_r,A_r)+\mathcal R(A_r)]\,dr$ has a fixed point in the norm $\sup_{0<s\le s_0(R)}\max\{s^\kappa\|w_s\|_\infty,s^{\frac12+\kappa}\|\nabla w_s\|_\infty\}$. The first iterate $\int_0^sP_{s-r}\mathcal N(P_r\mathbf B,P_r\mathbf B)\,dr$ is bounded through the quadratic member of the rough norm, by $\|P_{s-r}f\|_\infty\le(s-r)^{-1/4}|f|_{-\frac12;a}$ and $\int_0^s(s-r)^{-\frac14}r^{-\frac34-\kappa}dr\le c\,s^{-\kappa}$; the cross terms $\mathcal N(P_r\mathbf B,w_r)$, $\mathcal N(w_r,P_r\mathbf B)$, $\mathcal N(w_r,w_r)$ give the unweighted factor $s^{\frac14-\frac32\kappa}$.

This is a positive power because $\kappa<\frac16$; multiplication by the weight $s^\kappa$ in the fixed-point norm gives the contraction condition $R^2s^{\frac14-\frac\kappa2}\le c_1$; the members of order at least three are bounded in the chart, $|\mathcal R(A)|\le2\sup|\Theta^{\rm ch}|\,\|A\|_\infty^3$ for $a_K\|A\|_\infty\le\frac12$ by Cauchy's estimate, and the flow does not leave the chart on $[0,s_0(R)]$ for $a_K\le a_0(R)$, by the choice of $c_2$, as $a_K\|\mathbf B\|_\infty\le e^{12}Ra_K^{\frac12-\kappa}$. The other two bounds follow with $s^{\frac14+\frac\kappa2}\|P_s\mathbf B\|_\infty\le R$, $s^{\frac34+\frac\kappa2}\|\nabla P_s\mathbf B\|_\infty\le2R$.

The velocities of $V_s$ and of $\Phi_s(V_0)$ differ by a gauge direction, so the two are gauge equivalent. For two data, $(A_s-A'_s)-P_s(\mathbf B-\mathbf B')$ solves the same kind of equation, into which the quadratic member enters through $\mathcal N(P_r\mathbf B,P_r\mathbf B)-\mathcal N(P_r\mathbf B',P_r\mathbf B')$, which the rough distance bounds; the fixed point closes on $(0,s_{\rm Lip}(R)]$. The action: for $s\le s_0(R)$ and $a_K\le a_0(R)$, $A^a\le\frac32\ell^3(a_K^{-2}\varpi)^2$ and $\varpi(e^{ia_KA})\le2a_K^2(\|A\|_\infty^2+\|\nabla A\|_\infty)$ (Lemma~\ref{C-367-13}(c), (e)) give with the four bounds $A^a(\Phi_se^{ia_K\mathbf B})\le6((1+c_3)^2+2+c_3)^2\ell^3R^4s^{-\frac32-\kappa}$; for $s_0(R)<s\le1$ the action does not increase (Lemma~\ref{C-367-13}(c)), and its value at $s_0(R)$ is at most $CR^4(R^2/c_1)^{4(\frac32+\kappa)/(1-2\kappa)}$, a power of $R$ below $24$; for $a_K>a_0(R)$, $A^a\le6\ell^3a_K^{-4}<6\ell^3(R/c_2)^{8/(1-2\kappa)}$, a power of $R$ below $12$.

So $c_4$ depends only on $\kappa$.

\emph{The field of the gauge of the chain.} Given $\mathcal F_{d-1}$, the $\sigma$-algebra of
$V_{K-d+1}$, the law of $B_{K-d}$ is $p_W$, $W=V_{K-d+1}$ (Lemma~\ref{C-367-7}(ii)). The layers are
made step by step, for $k=t-1,\dots,0$, starting from $\mathsf U_t=U_{\rm top}$, a background over
$V_t$, and ending at $k=0$, where $\mathcal M^0$ is the identity, so that $U$ is a gauge transform
of $\mathsf U_0$. At the step $k$, with $W=V_{k+1}$, let $U_+$ be a background over $W$ whose block
field $\mathcal M^kU_+$ is the base point $V_W(0)$ of the chart, a background of the level $k$ as
well, over $\mathcal M^kU_+$ (the background and chart identification in
Proposition~\YMMref{52.3}{433}), chosen as a Borel function of $W$, and let $\mathsf g_{k+1}$, a
function of $V_{k+1}$, be the gauge with $\mathsf U_{k+1}^{\mathsf g_{k+1}}=U_+$. The chart writes
$V_k=V_W(B_k)$, and $\zeta=x_k^{-1/2}B_k$, on the bonds that are not pivot bonds, has
$\sup|\zeta|<\delta(x_k)\le\frac16\varepsilon_{\mathrm A}$; by the global shift identification in
Proposition~\YMMref{51.5}{396}, and at $k=0$ by the pivot chart in equation~\YMMref{(800)}{391}, the
background over $V_k$ is a gauge transform of $e^{iL^{-k}\mathbf A[U_+;\zeta]}U_+$. As
$a_KL^d=L^{-k}$ and $(e^{iX}U)^g=e^{i\mathrm{Ad}(g)X}U^g$, the layer $\mathbf
h_d=L^d\mathrm{Ad}(\mathsf g_{k+1}^{-1})\mathbf A[U_+;\zeta]$ makes $\mathsf U_k=e^{ia_K\mathbf
h_d}\mathsf U_{k+1}$ the transform of $e^{iL^{-k}\mathbf A[U_+;\zeta]}U_+$ by $\mathsf
g_{k+1}^{-1}$, a function of $V_{k+1}$; $\mathsf U_k$ is a function of $V_k$, and so is the gauge
$\mathsf g_k$ of the next step. The bound $|\mathbf A[U_+;\zeta]|\le c_{\mathsf h}\sup|\zeta|$
(Proposition~\YMMref{51.5}{396} and the pivot estimate~\YMMref{(800)}{391}) gives $|\mathbf h_d|\le
c_{\mathsf h}L^d\delta(L^d\mathsf b)$. The product that defines $\mathbf B_K$ lies within $\frac38$
of $\mathbf 1$ when $\sum_da_K|\mathbf h_d|\le\frac14$ and $a_K\mathsf a_{\rm top}\le\frac18$.

\emph{The regular gauge of the top.} The regularity number of the level $k$ of a configuration $U'$
on the bare lattice is $L^{2k}\max_p|\partial U'(p)-\mathbf 1|$, the maximum over the bare
plaquettes $p$ (Definition~\ref{def:level-regularity}); at the level $0$ it is $\varpi(U')$, and at
the level $k$ it is $L^{2k}$ times its value at the level $0$ (by that definition in the two units).
For $k\ge1$, $\mathsf{Cut}_k$ consists of the block fields of the level $k$ whose backgrounds exist
and have regularity numbers of the level $k$ below $c_\rho\varepsilon(x_k)$ (the definition of the
field cut). As $V_t\in\mathsf I_{d_0}\subset\mathsf{Cut}_t$ and $x_t=L^{d_0}\mathsf b$, every
background $U'$ over $V_t$ has $L^{2t}\varpi(U')<c_\rho\varepsilon(L^{d_0}\mathsf b)$, so
$\varpi(U')<c_\rho\varepsilon(L^{d_0}\mathsf b)L^{-2t}=c_\rho\varepsilon(L^{d_0}\mathsf
b)L^{2d_0}a_K^2$. Lemma~\ref{C-367-11}(a) with $\mathsf n=N$, at
$\kappa'=\max\{c_\rho\varepsilon(L^{d_0}\mathsf b)L^{2d_0},\ell^{-2}\}a_K^2$, for which
$N\ge\kappa'^{-1/2}$, gives a gauge $\gamma$, which a measurable selection makes depend on $V_t$
alone in a Borel way, with $|U_{\rm top}^\gamma(b)-\mathbf 1|\le C_{\rm BG}\sqrt{\kappa'}\le C_{\rm
BG}\max\{(c_\rho\varepsilon(L^{d_0}\mathsf b))^{1/2},\ell^{-1}\}L^{d_0}a_K\le CL^{d_0}a_K$ at every
bond; at the finitely many $K$ with $\kappa'>\kappa_0$ or $N<\lambda_{{\rm BG},3}$ we take
$\gamma=\mathbf 1$, and there $\pi/a_K\le CL^{d_0}$. So $U_{\rm top}^\gamma=e^{ia_K\mathbf A^{\rm
top}}$, with $ia_K\mathbf A^{\rm top}$ a fixed Borel $\mathfrak{su}(2)$ logarithm of norm at most
$\pi$, equal to the principal logarithm off $-\mathbf1$, at each positive bond, and extended by
reversal, so that $\mathbf A^{\rm top}$ too depends on $V_t$ alone in a Borel way, and $|\mathbf
A^{\rm top}|\le\frac\pi2a_K^{-1}|U_{\rm top}^\gamma-\mathbf 1|$ gives $|\mathbf A^{\rm
top}|\le\mathsf a_{\rm top}=CL^{d_0}\le CL^{10d_0}$. The bounds of $\mathsf J_d$, and those of the
derivative of the layer on the box, come from the locality of the shift. Let $\mathfrak n$ be the
distance, in the maximum norm and in units of their side, of the cubes of $\pi^M_{k+1}$ that hold
$\beta_-$ and $b_-$, for a bare bond $\beta$ and a bond $b$ of the level $k$, and $e^{\bar\kappa}$
the radius of the polydisc $\mathbb D(Y)$ of the parameters of the shift systems
(Proposition~\YMMref{51.5}{396} and the hypotheses preceding it). For $\mathfrak n\ge3$ the value of
$\mathbf A[U_+;\zeta]$ at $\beta$ differs from that of the shift of the cube of $(2\mathfrak n-1)^3$
cubes about the cube of $\beta_-$, which does not read $\zeta(b)$, by at most $2c_{\mathsf
h}\sup|\zeta|e^{-\bar\kappa(\mathfrak n-2)}$: let every fine cube, that is every cube of $\pi^M_k$,
in the shell of the cubes of this cube at the distance $j$ from the cube of $\beta_-$, $2\le
j\le\mathfrak n-1$, one of its $\mathfrak n-2$ outer shells, carry a single variable $s_j$,
$|s_j|<e^{\bar\kappa}$, and every fine cube outside these shells the value $1$; by the holomorphy,
support and restriction properties of Proposition~\YMMref{51.5}{396} the difference is then the
value at $s_2=\dots=s_{\mathfrak n-1}=1$ of a holomorphic function of these $\mathfrak n-2$
variables, of modulus at most $2c_{\mathsf h}\sup|\zeta|$ on their polydisc, and zero as soon as
some $s_j$ is zero, since that shell then separates the cubes nearer than it to the cube of
$\beta_-$, among them those that hold the two ends of $\beta$, from the cubes beyond it; and
Schwarz's lemma in the $s_j$ gives the bound. So, by Cauchy's estimate, for
$\sup|\zeta|\le\frac12\varepsilon_{\mathrm A}$ the derivative of $\mathbf A[U_+;\zeta]$ at $\beta$
in the value of $\zeta$ at $b$ is at most $4c_{\mathsf h}e^{-\bar\kappa(\mathfrak n-2)_+}$, and for
$k=0$ it vanishes unless $\mathfrak n\le1$ (the double-block dependence of the pivot
chart~\YMMref{(800)}{391}). As at most $125$ cubes lie at a distance $\le2$ from a given one and at
most $24\mathfrak n^2+2$ at the distance $\mathfrak n$, Schur's test gives the two norms of $\mathsf
J_d$, and the same bounds for the derivative of the layer on the box, with $c_{\mathsf
J}=\sqrt{27}(LM)^3\cdot4c_{\mathsf h}[125+\sum_{\mathfrak n\ge3}(24\mathfrak
n^2+2)e^{-\bar\kappa(\mathfrak n-2)}]$, and $\|\mathsf J_d^*f\|\le c_{\mathsf
J}\min\{L^{-3d/2}\|f\|_2,\|f\|_1\}$.

Each layer splits as $\mathbf h'_d=\mathbf n_d+u_{(d)}+e_{(d)}$: $\mathbf n_d$ its conditional mean given $\mathcal F_{d-1}$; $u_{(d)}=\mathsf b^{-1/2}L^{d/2}\mathrm{Ad}(\gamma)\mathsf J_d\tilde B_d$, linear in the centred fluctuation $\tilde B_d=B_{K-d}-\int B\,dp_W$ with a predictable coefficient; $e_{(d)}$ the centred rest, which carries besides the factor $\delta(L^d\mathsf b)$. So $\mathbf B_K=u+\rho$, $u=\sum_du_{(d)}$, $\rho$ the sum of the $e_{(d)}$, of the $\mathbf n_d$, of $\mathbf A^{\rm top}$ and of the remainder of the ordered product beyond the sum of its factors.

\emph{The linear fluctuation.} By Lemma~\ref{C-367-7}(iv), given $\mathcal F_{d-1}$ every Lipschitz function of $B_{K-d}$ has Gaussian tails, with the variance proxy $\frac8{\gamma_{\rm r}}$ times the square of its Lipschitz constant; for $\langle u_{(d)},f\rangle$ this square is at most $c_{\mathsf J}^2\mathsf b^{-1}L^d\min\{\|f\|_1^2,L^{-3d}\|f\|_2^2\}$, so $(P_su_{(d)})(x)$, $f=k_s(x,\cdot)$ the heat kernel, has the proxy $C\mathsf b^{-1}L^d\min\{1,L^{-3d}s^{-3/2}\}$, whose sum over $d$ is $\le C\mathsf b^{-1}s^{-1/2}$. Exponential supermartingales over the depths, with test functions that are functions of $\mathcal F_{d-1}$, and a union over dyadic grids of times and sites, through the grid lemma for $|\cdot|_{\eta;a}$ and the bounds of the heat kernel, give Gaussian tails of $|u|_{-\frac12-\kappa;a}$.

\emph{The quadratic fluctuation and colour cancellation.} Sort the pairs in $\mathcal N_s[u,u]$ by the finer index. A pair of different depths is linear in the finer layer with a predictable coefficient: a martingale difference, with Gaussian tails as before. The diagonal $\mathcal N_s[u_{(d)},u_{(d)}]$, after $P_r$ at a point, is a quadratic form $(\tilde B_d,\mathcal Q_d\tilde B_d)$ of the centred fluctuation, $\|\mathcal Q_d\|_{\rm tr}\le C\mathsf b^{-1}L^ds^{-1/2}\min\{1,L^{-3d}s^{-3/2}\}$; its centred part has exponential moments by (Q), through the norm and the Hilbert--Schmidt norm of $\mathcal Q_d$. Its compensator, the conditional mean $\mathrm{tr}(\mathcal Q_d\mathrm{Cov}_W)$, would sum by the sizes to the order $\mathsf b^{-1}s^{-1}$, beyond the weight $s^{\frac34+\kappa}$; but it reads the conditional covariance of the layer on the bare lattice, in the gauge of the chain, only through its part that is antisymmetric in colour, as $X\wedge Y$ is antisymmetric.

By (Col) the scalar part drops out, and almost surely 
\[
|\mathrm{tr}(\mathcal Q_d\mathrm{Cov}_W)|\le Cs^{-1/2}\min\{1,L^{-3d}s^{-3/2}\}[L^{d/2}(\log x)^{\mathfrak p_0+6}+\mathsf a_{\rm top}],
\]
$x=L^d\mathsf b$, whose sum over $d$ is $\le C(1+\mathsf a_{\rm top})s^{-\frac34-\frac\kappa2}$, which the weight $s^{\frac34+\kappa}$ admits: the cancellation in colour.

\emph{Means and nonlinear remainders.} The $e_{(d)}$ by the bounds of the linear part, with the factor $\delta(L^d\mathsf b)$; the means by (M), summable in $|\cdot|_{-\kappa;a}$ through $\|\cdot\|_{L^p}$, $p=\lceil3/\kappa\rceil$; $\mathbf A^{\rm top}$ by $\mathsf a_{\rm top}$; the ordered product, at the second order linear in the finer factor with a predictable coefficient, at the third order small; each with $|\cdot|_{-\kappa;a}\le C(1+\mathsf a_{\rm top})$ in every $L^q$. The cross members of $u$ and $\rho$ by the sizes, as $(-\frac12-\kappa)+(-\kappa)\ge-\frac12-2\kappa$. Together $\mathbb E^{[d_0]}[\|\mathbf B_K\|_{\kappa;a}^q]^{1/q}\le C(q,\kappa)(1+\mathsf a_{\rm top})^2$, every constant reading $d_0$ only through $\mathsf a_{\rm top}$; with $\mathsf a_{\rm top}\le CL^{10d_0}$ this is the bound of the statement; for the finitely many $K$ with $a_K\mathsf a_{\rm top}>\frac18$ the principal logarithm has $\|\mathbf B_K\|_\infty<8\pi\mathsf a_{\rm top}$.

The statements on one fluctuation integral that enter, besides Lemma~\ref{C-367-7}(iv), the bounds of the response, the bounds of the heat kernel and the grid lemma, are, with $C$ free of $K$ and $d_0$: (M) $\|\mathbf n_d\|_\infty\le C\mathsf b^{-1}(\log L^d\mathsf b)^{2\mathfrak p_0}$ and $\|\mathbf n_d\|_{L^2}^2\le C\ell^3\mathsf b^{-3}L^{-d}(\log L^d\mathsf b)^{8\mathfrak p_0+12}$; (Col) in the gauge of the chain the kernel of $\mathrm{Ad}(\gamma)\mathsf J_d\mathsf A_W^{-1}\mathsf J_d^*\mathrm{Ad}(\gamma)^{\mathsf T}$ at two bare bonds whose partition cubes are at distance $N$ is a number times the identity in colour up to $\mathsf w_N$, $\sum_N(24N^2+2)\mathsf w_N\le C[x^{-1/2}(\log x)^{\mathfrak p_0+6}+\mathsf a_{\rm top}L^{-d}]$, the number being formed at the trivial background, where the Gaussian form and the response commute with the adjoint action; and $\|\mathrm{Cov}_W-\mathsf A_W^{-1}\|\le Cx^{-1/2}(\log x)^{\mathfrak p_0+3}$, $\mathsf A_W$ the Gaussian form of $p_W$.\par \textup{(Q)} There are universal constants $C_{\rm HW},c_{\rm HW}>0$ with the following property. Let the real symmetric $\mathcal Q_d$ be predictable, that is, functions of $\mathcal F_{d-1}$, and assume almost surely the deterministic bounds $\|\mathcal Q_d\|\le\mathsf q$ and $C_{\rm HW}\gamma_{\rm r}^{-2}\sum_d\|\mathcal Q_d\|_{\rm HS}^2\le\Sigma^2$. For the actual conditional covariance set $\mathsf X_d=(\tilde B_d,\mathcal Q_d\tilde B_d)-\mathrm{tr}(\mathcal Q_d\mathrm{Cov}_W)$. Then $\mathbb E^{[d_0]}e^{\epsilon\sum_d\mathsf X_d}\le e^{\epsilon^2\Sigma^2}$ for $|\epsilon|\mathsf q\le c_{\rm HW}\gamma_{\rm r}$.\par Indeed the conditional box law has Hessian at least $\gamma_{\rm r}/4$ and log-Sobolev constant $8/\gamma_{\rm r}$ (Lemma~\ref{C-367-7}(iv)); the convex-domain criterion \cite[Theorem~2.1]{KolesnikovMilman2016} applies. For nonzero $\mathcal Q$, Adamczak's dependent-vector inequality \cite[Theorem~2.3]{Adamczak2014} gives the centred quadratic form $Y$ the conditional tail $2\exp[-c\min\{\gamma_{\rm r}^2t^2/\|\mathcal Q\|_{\rm HS}^2,\gamma_{\rm r}t/\|\mathcal Q\|\}]$. For $\epsilon>0$, centering gives \[\mathbb Ee^{\epsilon Y}-1=\epsilon\int_0^\infty\bigl[(e^{\epsilon t}-1)\mathbb P(Y>t)+(1-e^{-\epsilon t})\mathbb P(Y<-t)\bigr]dt.\] Integrating the Gaussian and exponential tails bounds the conditional logarithmic moment by $C_{\rm HW}\epsilon^2\gamma_{\rm r}^{-2}\|\mathcal Q\|_{\rm HS}^2$ for $|\epsilon|\|\mathcal Q\|\le c_{\rm HW}\gamma_{\rm r}$; apply the same argument to $-Y$. Consequently $\exp[\epsilon\sum_{j\le d}\mathsf X_j-C_{\rm HW}\epsilon^2\gamma_{\rm r}^{-2}\sum_{j\le d}\|\mathcal Q_j\|_{\rm HS}^2]$ is a conditional supermartingale. Iteration and the pathwise budget prove \textup{(Q)}; independence of the layers is unnecessary.

\emph{The proof of \textup{(Adm)}.} Here $d_0<d\le K$, $k=K-d$, $x=L^d\mathsf b$, $g=x^{-1/2}$,
$N_k=(\ell L^d)^3$, $W\in\mathsf{Cut}_{k+1}$, $U_+$ is as above, $\bar B_W=\int B\,p_W(dB)$, and $v$
is the sum of the potentials, so that $\Psi_W=\frac12(B,\mathsf A_WB)-v$ up to a constant
(Proposition~\ref{C-367-8}(b)). (1)

\emph{Conditions on numbers.} Besides the chain, the bounds of (2) use only the following conditions
on numbers. At the top depth there are six,
\[
\begin{gathered}
1\le d_0<K,\qquad L^{d_0}\mathsf b\ge\mathsf t_\natural,\\
\mathsf R_{d_0}\le\tfrac12,\qquad
C_{\rm hs}\,\delta(L^{d_0}\mathsf b)\le\tfrac14\gamma_{\rm r},\\
4c_{\mathsf h}\delta(L^{d_0}\mathsf b)\le\tfrac14,\qquad
c_\rho\varepsilon(L^{d_0}\mathsf b)\le C_{\rm BG}^{-2},
\end{gathered}
\]
with $C_{\rm hs}=C_{\mathrm
E,1}[338K_P+\frac12K_\Lambda(1+\beta_{\rm s})N^{\rm s}_\star/\alpha_0]$, formed from $C_{\mathrm
E,1}$ and the numbers that make up $C_V$ in the setting; the fifth then holds at every depth $d>d_0$
too, $\delta$ being decreasing from the reciprocal of the squared coupling floor on, a number that
$\mathsf b$ does not fall below under (E) (the decreasing cutoff functions in
Section~\ref{sec:cutoff-calibration}). Together with them, (2) uses at the depths $d>d_0$ the bound
$C_{\rm Hs}\delta(L^d\mathsf b)\le\frac14\gamma_{\rm r}$, with the constant $C_{\rm
Hs}=144\,C_{\mathrm E,1}C_V$ of the setting, which exceeds $C_{\rm hs}$, as $3\cdot144>338$ and
$144>\frac12$, so that at $d_0$ this bound contains the fourth of the six; by the choice of $n_2$ it
holds at every depth $d\ge n_2$, and $n_2\le n_{\rm a}\le d_{\rm adm}$, so that it holds at $d_0$
and at every depth $d>d_0$ as soon as $d_0\ge d_{\rm adm}$ and asks nothing of $d_{\rm adm}$ beyond
its member $n_{\rm a}$. At each depth $d>d_0$, with $x=L^d\mathsf b$ as above, $\bar\kappa\ge1$
(Definition~\YMMref{47.2}{359}) and
\[
\begin{gathered}
\nu_0=\Bigl\lceil\frac{\log x}{2\bar\kappa}\Bigr\rceil,\qquad
\nu_1=\Bigl\lceil\frac{\lceil\log x\rceil+2}L\Bigr\rceil,\\
\nu_2=\max\Bigl\{\nu_1,\Bigl\lceil\frac{3\log x}{2\bar\kappa}\Bigr\rceil+1\Bigr\},\qquad
\nu_3=\max\Bigl\{\nu_1,2\Bigl\lceil\frac{2\log x}{\bar\kappa}\Bigr\rceil+19\Bigr\},
\end{gathered}
\]
there are three conditions on cubes, at each of the three
inner radii $\nu=\nu_1,\nu_2,\nu_3$:
\[
\begin{gathered}
2M(2\nu+2\nu_0+3)^3+4<L^{m+d-1},\\
4M(2\nu+2\nu_0+3)^3\,\theta_{d-1}\le\varepsilon_{\mathcal W}\,c_{*,\mathrm{rg}},\\
4B_*M(2\nu+2\nu_0+3)^3\,\theta_{d-1}
\le\tfrac12\varepsilon_{\mathcal W}\,a^{\rm us}(\alpha_0,\alpha_1),
\end{gathered}
\]
with $\theta_{d-1}=2c_\rho\varepsilon(x/L)$, the absolute
numbers $B_*\ge1$ and $c_{*,\mathrm{rg}}\in(0,1]$ of the local regular-gauge
Theorem~\YMMref{F}{199}, and the usable radius $a^{\rm us}(\alpha_0,\alpha_1)$ in
equation~\YMMref{(933)}{427} at the fixed radii. $(2\nu+2\nu_0+3)^3$ is the number of the cubes of
$\pi^M_{k+1}$ at distances up to $\nu+\nu_0+1$ from a given one, in units of their side; they form
the cube $Q_\nu$ about it, of side $(2\nu+2\nu_0+3)M$ in the units of the level $k+1$, whose cubes
at distances up to $\nu$ form its \emph{inner cube}, parted from the cubes outside $Q_\nu$ by its
\emph{collar}, the $\nu_0$ shells of the cubes at the distances $\nu+2,\dots,\nu+\nu_0+1$. The
\emph{trivial triple} on a polymer $Y$ is $P[0]=(\mathbf 1,0,0)$, the value at $0$ of the chart
$P[\omega]=(\mathbf 1,A,\mathbf J)$ at the trivial pair, $\omega=(A,\mathbf J)$, whose pair is
$\mathbf p(P[\omega])=(e^{iL^{-k-1}A},\mathbf J)$ (the triple and pair conventions of
Definition~\YMMref{49.1}{374}); a triple $P$ on $Y$ is \emph{within $r$ of the trivial triple in a
gauge $\mathfrak v$} if $P^{\mathfrak v}$ has the pair of $P[\omega]$ for a real $\omega$ with
$\|\omega\|_Y<r$, the triple norm, namely the largest of $|A|$, $|\nabla_{\mathbf 1}A|$ and
$|\mathbf J|$ on $\mathrm{Bd}(Y)$. Every term of the exponent of $p_W$ reads a triple through its
pair alone (proof of Proposition~\ref{C-367-8}(c)), and so do the structure of the step and the
shift (Theorem~\YMMref{U}{378}, Proposition~\YMMref{51.5}{396}, and the actual-core verification in
the proof of Theorem~\YMMref{W}{423}). Under the three conditions $Q_\nu$ is small against the side
of the torus of the level, and Theorem~\YMMref{F}{199} at the level $k+1$, on $Q_\nu$ with $\mathsf
l=(2\nu+2\nu_0+3)M$ and at $\alpha=2\theta_{d-1}/\varepsilon_{\mathcal W}$ (the first two conditions
give its conditions on the sides and on $\mathsf l\alpha$), gives a gauge $\mathfrak v$ and a real
$1$-form $A'$ with $U_+^{\mathfrak v}=e^{iL^{-k-1}A'}$ on the bare bonds of the concentric cube of
side $\mathsf l+4$ and $|A'|,|\nabla_{\mathbf 1}A'|<B_*\mathsf
l\alpha=2B_*(2\nu+2\nu_0+3)M\theta_{d-1}/\varepsilon_{\mathcal W}$ on $\mathrm{Bd}(Q_\nu)$ (the
potential and first-difference estimates of that theorem). Its hypotheses on $U_+$, at the spacing
$\eta=\eta_{k+1}$, are the bounds $|\partial U_+(p)-\mathbf 1|<\alpha\eta^2$ at the bare plaquettes
$p$ with a corner in the concentric cube of side $\mathsf l+4$ and $|(D^{\mathsf T}_{U_+}\partial
U_+)(b)|<\alpha\eta^2$ at the bare bonds $b$ with an end point in that cube; here, for a
configuration $U$, $D^{\mathsf T}_U\partial U$ is the $1$-form whose value at $b=\langle x,x+\eta
e_\mu\rangle$ is $\eta^{-1}\sum_{\nu\ne\mu}[X_{\mu\nu}(x)-\mathrm{Ad}(U(\langle x,x-\eta
e_\nu\rangle))X_{\mu\nu}(x-\eta e_\nu)]$, with $X_{\mu\nu}(y)$ the value of $\partial U$ at the
plaquette of the directions $\mu<\nu$ based at $y$ and $X_{\nu\mu}=-X_{\mu\nu}$ (the
covariant-difference convention of Section~\YMMref{22}{171}), that is, the transpose, at the spacing
$\eta$ and at $U$, of the covariant exterior derivative of $1$-forms into $2$-forms, applied to the
matrix-valued $2$-form $p\mapsto\partial U(p)$. Both bounds hold at all the bare plaquettes and
bonds of the torus: $U_+$, a background over $W\in\mathsf{Cut}_{k+1}$, has regularity number of the
level $k+1$ less than $c_\rho\varepsilon(x/L)=\frac12\theta_{d-1}$ (the definition of the field
cut), so $|\partial U_+-\mathbf 1|<\frac12\theta_{d-1}\eta_{k+1}^2<\alpha\eta_{k+1}^2$; its current
$\mathbf J_+$ is below $\theta_{d-1}/\varepsilon_{\mathcal W}$ at every bare bond, as
$\varpi(W)<\theta_{d-1}$ admits the plain-core verification in the proof of
Proposition~\YMMref{52.3}{433} at the free radii $\theta_{d-1}/\varepsilon_{\mathcal W}$, $\alpha_1$
(as in the proof of Proposition~\ref{C-367-8}(a)), which puts $P_W$ in the core with these radii,
and the condition (iii) of its simple set bounds the current by the first radius
(Definition~\YMMref{49.1}{374}); so $|D^{\mathsf T}_{U_+}\partial U_+|\le\eta_{k+1}^2|\mathbf
J_+|+\frac14\eta_{k+1}^3\theta_{d-1}^2<\alpha\eta_{k+1}^2$ (equation~\YMMref{(579)}{288}). As
$|\mathrm{Ad}(\mathfrak v)\mathbf J_+|=|\mathbf J_+|$, $P_W$ is within
$r_\nu:=4B_*M(2\nu+2\nu_0+3)^3\theta_{d-1}/\varepsilon_{\mathcal W}$ of the trivial triple in the
gauge $\mathfrak v$, with $\omega=(A',\mathrm{Ad}(\mathfrak v)\mathbf J_+)$, on $Q_\nu$ and on every
polymer in it; $r_\nu$ is below a multiple of $x^{-1/2}(\log x)^{\mathfrak p_0+3}$ and below half
the usable radius, so that the chart carries the ball of radius $2r_\nu$ into the analyticity cores.
The cubes $Q_{\nu_1}$, $Q_{\nu_2}$, $Q_{\nu_3}$ serve (i), (ii), (iii) of (2) in turn; in each the
collar is the room in which $\mathsf A_W^{-1}$ on the inner cube forgets the cubes outside, up to a
multiple of $x^{-1/2}$ (see (i)). The six conditions at the top depth and the conditions on cubes
hold for every integer $d_0\ge d_{\rm num}(\ell,\mathsf b)$, every $K>d_0$ and every depth $d>d_0$,
where
\[
\begin{aligned}
d_{\rm num}(\ell,\mathsf b):=\max\Bigl\{&1,\ \Bigl\lceil\frac{\log(\mathsf t_{\rm adm}/\mathsf b)}{\log L}\Bigr\rceil,\\
&\Bigl\lceil\frac{\mathsf y_{\mathsf R}(\ell,\mathsf b)-\log\mathsf b}{\log L}\Bigr\rceil,\\
&\Bigl\lceil\frac2{\log L}\log\frac{12\cdot10^5\,LM(\log\mathsf b)^3}\ell\Bigr\rceil\Bigr\},
\end{aligned}
\]
\[
\begin{aligned}
\mathsf t_{\rm adm}=\max\Bigl\{&\mathsf t_\natural,\\
&\Bigl(\frac{4\mathfrak p_0}e\Bigr)^{4\mathfrak p_0}
\max\Bigl\{\frac{4C_{\rm hs}A_1}{\gamma_{\rm r}},\,16c_{\mathsf h}A_1,\,C_{\rm BG}^2c_\rho A_0\Bigr\}^4,\\
&\Bigl(\frac{125000\,c_\rho A_0L^{1/2}M}{\varepsilon_{\mathcal W}}\Bigr)^4
\Bigl(\frac{4\mathfrak p_0+12}e\Bigr)^{4\mathfrak p_0+12}\\
&\qquad\max\Bigl\{\frac1{c_{*,\mathrm{rg}}},\,\frac{2B_*}{a^{\rm us}(\alpha_0,\alpha_1)}\Bigr\}^4\Bigr\},
\end{aligned}
\]
\[
\begin{gathered}
\mathsf y_{\mathsf R}(\ell,\mathsf b)=\max\Bigl\{
\Bigl(\frac6{c'A_0^2}\Bigr)^{\frac1{2\mathfrak p_0-1}},\
\Bigl(\frac{2\,\mathsf E(\ell,\mathsf b)}{c'A_0^2}\Bigr)^{\frac1{2\mathfrak p_0}}\Bigr\},\\
\mathsf E(\ell,\mathsf b)=\log4+3\log\frac\ell M+\max\{\mathsf k(\mathsf b),0\}\,\ell^3,
\end{gathered}
\]
with the rate $c'>0$ in the
bound defining $\mathsf R_d$ and $\mathsf k(\mathsf b)=\mathsf C^\Sigma_{\rm maj}+2+C_{\rm
low}(\mathsf b)$, the factor of $\ell^3$ in the last exponent of the definition of $\mathsf R_d$;
$\mathsf E(\ell,\mathsf b)>0$, as $\ell=2L^m\ge2MR(\mathsf b)\ge2M$ under (E). $\mathsf t_{\rm adm}$
is formed from the constants alone, and $d_{\rm num}(\ell,\mathsf b)$ from $\ell$, $\mathsf b$ and
the constants of the construction; it reads neither $K$ nor a level nor a field. For the proof: for
$\mathsf t>1$ and $r>0$, $\mathsf t^{-1/2}(\log\mathsf t)^r\le(4r/e)^r\mathsf t^{-1/4}$, as $r\log
z-\frac z4$ is largest at $z=4r$. Let $d_0\ge d_{\rm num}(\ell,\mathsf b)$ and $K>d_0$, which is the
first of the six. Then $L^{d_0}\mathsf b\ge\mathsf t_{\rm adm}$, which gives the second and, as
$\delta(\mathsf t)$ and $\varepsilon(\mathsf t)$ are $A_1$ and $A_0$ times $\mathsf
t^{-1/2}(\log\mathsf t)^{\mathfrak p_0}$, by that inequality at $r=\mathfrak p_0$ the last three;
and $z_0=\log L^{d_0}\mathsf b\ge\mathsf y_{\mathsf R}(\ell,\mathsf b)$, so that half of
$c'A_0^2z_0^{2\mathfrak p_0}$ is at least $3z_0$ and at least $\mathsf E(\ell,\mathsf b)$, while the
definition of $\mathsf R_{d_0}$, at the $(\ell L^{d_0})^3$ sites of the depth $d_0$ and with
$3d_0\log L<3z_0$, gives $\log2\mathsf R_{d_0}=\log4+3\log\frac\ell M+3d_0\log L+\mathsf k(\mathsf
b)\ell^3-c'A_0^2z_0^{2\mathfrak p_0}<\mathsf E(\ell,\mathsf b)+3z_0-c'A_0^2z_0^{2\mathfrak
p_0}\le0$: the third. At a depth $d>d_0$, $x\ge L\mathsf t_{\rm adm}$, and $\log x>3$, as
$\log\mathsf b>2$ (the label floor of Section~\ref{sec:cutoff-calibration}); as $\lceil
z\rceil<z+1$, $L\ge13$ and $\bar\kappa\ge1$, each of $\nu_1,\nu_2,\nu_3$ is less than $4\log x+21$
and $\nu_0<\frac12\log x+1$, so that $(2\nu+2\nu_0+3)^3<(9\log x+47)^3<15625(\log x)^3$. With
$\theta_{d-1}<2c_\rho A_0L^{1/2}x^{-1/2}(\log x)^{\mathfrak p_0}$ and that inequality at
$r=\mathfrak p_0+3$, the left member of the second condition on cubes is less than $125000\,c_\rho
A_0L^{1/2}M(\frac{4\mathfrak p_0+12}e)^{\mathfrak p_0+3}x^{-1/4}$ and that of the third less than
$B_*$ times this number; and $x\ge\mathsf t_{\rm adm}$ makes these two bounds at most
$\varepsilon_{\mathcal W}c_{*,\mathrm{rg}}$ and $\frac12\varepsilon_{\mathcal W}a^{\rm
us}(\alpha_0,\alpha_1)$. And as $\log x\le\log\mathsf b\,(1+d\log L)$ and $(1+z)^3<17.74\,e^{z/2}$
for $z\ge0$, the left member of the first is less than $31251M(\log x)^3<6\cdot10^5M(\log\mathsf
b)^3L^{d/2}$, which is below $L^{m+d-1}$, as $L^{d/2}>L^{d_0/2}\ge6\cdot10^5L^{1-m}M(\log\mathsf
b)^3$ by the last member of $d_{\rm num}(\ell,\mathsf b)$, $\ell$ being $2L^m$.

\emph{A uniform second moment.} For the fluctuation fields the bound holds under $p_W$ at every $W$.
Fix $W$ and a bond $b_0$. At a bond $b$, the divergence theorem in the variable $B(b)$ on the ball
$\{|B(b)|<\mathfrak b(x_k)\}$ of the box, applied to $B(b)e^{-\Psi_W}$, gives the flux $\mathfrak
b(x_k)\int_{|B(b)|=\mathfrak b(x_k)}e^{-\Psi_W}\ge0$, so $\int B(b)\cdot\nabla_b\Psi_W\,dp_W\le3$.
By Proposition~\ref{C-367-8}(b), $\Psi_W=\frac12(B,\mathsf A_WB)-v$ up to a constant, $v$ the sum of
the potentials; $\mathsf A_W\ge\frac12\gamma_{\rm r}$, and Cauchy's estimate gives $|\nabla_bv|\le
C\delta(x_k)\mathfrak b(x_k)\le\eta$ on the box, $\eta$ a number free of $x_k$ (the proof of
Lemma~\ref{C-367-7}(iv)). Take the weights $w_b=e^{-\mu\,\mathrm{dist}(b,b_0)}$, $\mathrm{dist}$ the
distance in the units of the level $k$. The entries of $\mathsf A_W$ fall off exponentially with the
distance of the bonds, as $\mathsf A_W$ is a sum of local terms with the bounds of the polymers
(Proposition~\ref{C-367-8}(b), Proposition~\YMMref{51.4}{393} and its proof), so Schur's test gives
$\|w^{1/2}\mathsf A_Ww^{-1/2}-\mathsf A_W\|\le\frac14\gamma_{\rm r}$ for a small $\mu>0$ that reads
neither $K$, $d$ nor $W$. Then $\frac14\gamma_{\rm
r}\sum_bw_b\int|B(b)|^2dp_W\le\sum_bw_b\int(3+\eta|B(b)|)\,dp_W$, and
$\eta|B(b)|\le2\eta^2/\gamma_{\rm r}+\frac18\gamma_{\rm r}|B(b)|^2$ gives $\int|B(b_0)|^2dp_W\le
m_2=\frac8{\gamma_{\rm r}}(3+\frac{2\eta^2}{\gamma_{\rm r}})\sum_bw_b$, which reads neither $K$,
$d$, $d_0$ nor $W$, as $\sum_bw_b\le C\mu^{-3}$.  The log-Sobolev bound of Lemma~\ref{C-367-7}(iv),
together with this second moment, gives Gaussian tails about a uniformly bounded mean. Consequently
the boundary layers used below have mass smaller than every power of $x$. (2)

\emph{The bounds on one integral}, for $d_0\ge d_{\rm adm}$; $C$ stands for numbers formed from
$\ell/\mathsf b$ and the constants of the construction. (i)

\emph{The mean of the fluctuation.} Integration by parts in each variable of the box, against a
cut-off that vanishes on its boundary, gives $\mathsf A_W\bar B_W$ as the mean of $\nabla v$ up to
terms carried by $\{|B(b)|\ge\frac12\mathfrak b(x)\}$, whose mass under $p_W$ is at most a multiple
of $e^{-c\,\mathfrak b(x)^2}$, $c>0$, and so smaller than every power of $x$; $\nabla v$ is at most
a multiple of $\delta(x)\mathfrak b(x)=A_1^2x^{-1/2}(\log x)^{2\mathfrak p_0}$ at each bond; and the
kernel of $\mathsf A_W^{-1}$ decays at the rate $\bar\kappa$: $\|\mathsf A_W^{-1}(b,b')\|\le
c_{\mathbf K}e^{-\bar\kappa(\mathfrak n-3)_+}$, $\mathfrak n$ the distance of the cubes of
$\pi^M_{k+1}$ that hold $b_-$ and $b'_-$, with $c_{\mathbf K}=\frac2{\gamma_{\rm
r}}+\frac4{\gamma_{\rm r}^2}(c_{\rm r}+\gamma_{\rm r})^2$, formed from $\gamma_{\rm r}$ and $c_{\rm
r}$ alone. Indeed $\mathsf A_W^{-1}$ is the whole covariance of the complete structure of the step
at $P_W$, a decoupled Gaussian structure with the constants $(\bar\kappa;\gamma_{\rm r},c_{\rm
r},\hat\varepsilon)$, $\hat\varepsilon=\frac12\mathsf u_{\rm G}(LM)^{-3}\le\frac12\gamma_{\rm r}$
(Definition~\YMMref{47.2}{359}, Theorem~\YMMref{U}{378}, and the actual-core verification in
Theorem~\YMMref{W}{423}); on the union of the two cubes it is $\mathbf C+\mathbf
C\Gamma\Gamma^{\mathsf T}\mathbf C$ at the decoupling parameters $1$ (the conditioning identity in
Theorem~\YMMref{U}{378}), where $\mathbf C=\mathbf C(Z_\circ,s)$ and
$\Gamma=\boldsymbol\Gamma(Z_\circ,s)$ are the conditioned covariance and the connecting operator,
the two families of operators that the structure assigns to a region $Z_\circ$, here the union of
the two cubes, and to the decoupling parameters $s$ (Definition~\YMMref{47.2}{359}); at all
parameters $s$ of the polydisc of radius $e^{\bar\kappa}$ this is an operator of norm at most
$c_{\mathbf K}$, as (St3) with $\hat\varepsilon\le\frac12\gamma_{\rm r}$ gives $\|\mathbf
C\|\le2/\gamma_{\rm r}$ and $\|\Gamma^{\mathsf T}\|=\|\Gamma\|\le c_{\rm r}+\hat\varepsilon\le
c_{\rm r}+\gamma_{\rm r}$; for $\mathfrak n\ge4$ its block between the two cubes is holomorphic in
$\mathfrak n-3$ variables, each the parameter of all the cubes of one of the shells at the distances
$2,\dots,\mathfrak n-2$ from the cube of $b_-$, the other parameters being $1$ ((St1)), and vanishes
as soon as one of them is $0$, since that shell then parts the two cubes into different
wall-components ((St2)); Schwarz's lemma in these variables gives the factor
$e^{-\bar\kappa(\mathfrak n-3)}$, and for $\mathfrak n\le3$ the bound is $\|\mathsf
A_W^{-1}\|\le\frac2{\gamma_{\rm r}}$. As at most $24\mathfrak n^2+2$ cubes lie at the distance
$\mathfrak n$ from a given one and $\bar\kappa\ge1$, the rows of $\mathsf A_W^{-1}$ have bounded
sums: so $\sup|\bar B_W|\le Cx^{-1/2}(\log x)^{2\mathfrak p_0}$. The same argument, with one
variable for each of the $\nu_0$ shells of the collar of a cube $Q_\nu$ of (1), compares the
parameters $1$ and $0$ on the cubes outside $Q_\nu$: as $\nu_0\bar\kappa\ge\frac12\log x$, on the
inner cube $\mathsf A_W^{-1}$ lies within $2c_{\mathbf K}x^{-1/2}$ of the operator at the parameters
$1$ on the collar and $0$ outside $Q_\nu$, which is that of the structure on $Q_\nu$ at the
restriction of $P_W$ (restriction consistency in Theorem~\YMMref{U}{378}); carried by the gauge
$\mathfrak v$ of $Q_\nu$ it is that of the structure at $P[\omega]$ (gauge covariance and dependence
on the pair), which at $P[z\omega]$ is holomorphic in $z$ for $|z|<a^{\rm
us}(\alpha_0,\alpha_1)/r_\nu$, by joint holomorphy in $(s,A',\mathbf J)$ at a fixed real factor,
taken at the real factor $\mathbf 1$, as the chart carries the ball of radius $a^{\rm
us}(\alpha_0,\alpha_1)$ into the cores (equation~\YMMref{(933)}{427} and its proof), and at most
$c_{\mathbf K}$ there, by (St3) as above, since Theorem~\YMMref{U}{378} gives at every triple of the
cores a structure with the same constants; so Schwarz's lemma in $z$ puts $\mathsf A_W^{-1}$ on the
inner cube, carried by $\mathfrak v$, within $2c_{\mathbf K}(x^{-1/2}+r_\nu/a^{\rm
us}(\alpha_0,\alpha_1))\le Cx^{-1/2}(\log x)^{\mathfrak p_0+3}$ of the operator formed at the
trivial triple on $Q_\nu$, which commutes with the constant gauge transformations, as these fix
$P[0]$ (gauge covariance). In the $\ell^2$ norm a further factor $x^{-1/2}$ comes, up to powers of
$\log x$: along the laws proportional to $e^{-\frac12(B,\mathsf A_WB)+\tau v}$ on the box,
$0\le\tau\le1$, whose exponents have Hessians $\ge\frac14\gamma_{\rm r}$, the Poincar\'e inequality
makes the means of $\nabla v$ at $\tau=0$ and at $\tau=1$ differ in $\ell^2$ by at most
$\frac4{\gamma_{\rm r}}C_{\rm Hs}\delta(x)\sup\|\nabla v\|$; and at $\tau=0$, the Gaussian law of
covariance $\mathsf A_W^{-1}$ conditioned on the box, the gradient at a bond $b$ of each potential
$v_Y$ with $b_-\in Y$ and $d(Y)<\lceil\log x\rceil$ has a mean smaller than its size by a factor
$Cx^{-1/2}(\log x)^{\mathfrak p_0+3}$: such a $Y$ lies within $(\lceil\log x\rceil+2)M\le\nu_1LM$ of
$b_-$ in the units of the level $k$, so in the inner cube of $Q_{\nu_1}$ about the cube of $b_-$; in
the gauge $\mathfrak v$ of $Q_{\nu_1}$ the triple is within $r_{\nu_1}$ of the trivial triple, at
which the potentials are invariant under the constant gauge transformations; carried by $\mathfrak
v$, the covariance lies on the inner cube within $Cx^{-1/2}(\log x)^{\mathfrak p_0+3}$ of an
operator that commutes with these (above); and $\mathfrak g$ holds no non-zero vector fixed by the
adjoint action. The potentials with $d(Y)\ge\lceil\log x\rceil$ add at most $C\delta(x)\mathfrak
b(x)x^{-1}$ at $b$, by the sum over polymers \eqref{eq:analytic-domain-tail} at $p=1$. So $\|\bar
B_W\|\le CN_k^{1/2}x^{-1}(\log x)^{3\mathfrak p_0+3}$. (ii)

\emph{The mean of the shift.} At a bare bond $\beta$ the parts of $\mathbf A[U_+;gB](\beta)$ of
degree two and of higher degree in $B$ are at most $c_{\mathsf h}\delta(x)^2/\varepsilon_{\mathrm
A}$ and $2c_{\mathsf h}\delta(x)^3/\varepsilon_{\mathrm A}^2$, by Cauchy's estimate; the mean of the
part of degree two is its value at $g\bar B_W$, at most a multiple of $g^2\sup|\bar B_W|^2$, plus
its contraction with $g^2\mathrm{Cov}_W$, and the latter is smaller than the size of that part by a
factor $x^{-1/2}$, up to powers of $\log x$: by the locality of the shift at $\mathfrak n=\nu_2+1$
(above), as $\bar\kappa(\nu_2-1)\ge\frac32\log x$, it may be replaced, at a cost below $4c_{\mathsf
h}\delta(x)x^{-3/2}$, by the shift of the inner cube of $Q_{\nu_2}$ about the cube of $\beta_-$, of
$(2\nu_2+1)^3$ cubes of $\pi^M_{k+1}$; by gauge covariance in Proposition~\YMMref{51.5}{396} it may
be formed in the gauge $\mathfrak v$ of $Q_{\nu_2}$, in which the triple is within $r_{\nu_2}\le
Cx^{-1/2}(\log x)^{\mathfrak p_0+3}$ of the trivial one, and by pair dependence and slice holomorphy
its part of degree two then differs from the one formed at the trivial triple by a multiple of
$r_{\nu_2}$ times its size; at the trivial triple the part of degree two is covariant under the
adjoint action of the constant gauge transformations, by covariance under constant gauges, so that
its contraction with an operator that commutes with them is an element of $\mathfrak g^c=\mathfrak
g\otimes\mathbb C$ fixed by the adjoint action, hence $0$; and on the inner cube the covariance of
$p_W$, carried by $\mathfrak v$, lies within $\epsilon^{\rm cv}_d+2c_{\mathbf
K}(x^{-1/2}+r_{\nu_2}/a^{\rm us}(\alpha_0,\alpha_1))\le Cx^{-1/2}(\log x)^{\mathfrak p_0+3}$ of such
an operator, by (i) and (iv). Hence $|\int\mathbf A[U_+;gB](\beta)\,p_W(dB)-(D\mathbf A[U_+;0]g\bar
B_W)(\beta)|\le\epsilon^{\rm s}_d\le Cx^{-3/2}(\log x)^{4\mathfrak p_0+6}$ at every bare bond; and
as $|D\mathbf A[U_+;0]X|\le c_{\mathsf h}\sup|X|$, and by the kernel bound of $D\mathbf A[U_+;0]$
above, which reads $\beta$ only through its cube, with Schur's test, $\sup_\beta|\int\mathbf
A[U_+;gB](\beta)\,p_W(dB)|\le\epsilon^{\rm m}_d$ and $\sum_y\max_{\beta_-\in B^k(y)}|\int\mathbf
A[U_+;gB](\beta)\,p_W(dB)|^2\le\epsilon^{\rm q}_d$, the sum over the sites $y$ of the level $k$,
where $\epsilon^{\rm m}_d$ is $c_{\mathsf h}g$ times the bound of $\sup|\bar B_W|$ in (i) plus
$\epsilon^{\rm s}_d$, and $\epsilon^{\rm q}_d$ is $2c_{\mathsf J}^2g^2$ times the square of the
bound of $\|\bar B_W\|$ in (i) plus $2N_k(\epsilon^{\rm s}_d)^2$. (iii)

\emph{The kernel.} The kernel of $D\mathbf A[U_+;0]\mathsf A_W^{-1}D\mathbf A[U_+;0]^*$ at two bare
bonds $\beta,\beta'$ is
\[
\sum_{b,b'}\partial_b\mathbf A[U_+;0](\beta)\,\mathsf
A_W^{-1}(b,b')\,\partial_{b'}\mathbf A[U_+;0](\beta')^{\mathsf T},
\]
with the kernel $\mathsf
A_W^{-1}(b,b')\in\mathrm{End}\,\mathfrak g$ of $\mathsf A_W^{-1}$, and its outer factors decay at
the rate $\bar\kappa$ in the distances of the cubes (above), and so does its middle one, by (i); so,
by the triangle inequality for these distances, it is at most $Ce^{-\bar\kappa\mathfrak n/2}$,
$\mathfrak n$ the distance of the cubes of $\beta_-$ and $\beta'_-$. Let $\mathfrak n\le
n_\flat:=\lceil2\log x/\bar\kappa\rceil+9$, so that $\nu_3=\max\{\nu_1,2n_\flat+1\}$, and let
$\mathfrak v$ be the gauge of $Q_{\nu_3}$ about the cube of $\beta_-$, in which $P_W$ is within
$r_{\nu_3}\le Cx^{-1/2}(\log x)^{\mathfrak p_0+3}$ of the trivial triple; the inner cube of
$Q_{\nu_3}$ holds every cube within $n_\flat$ of the cubes of $\beta_-$ and $\beta'_-$, and the
bonds $b$, $b'$ farther than $n_\flat$ from these add at most $Cx^{-2}(\log x)^2$, as
$\bar\kappa(n_\flat-2)\ge2\log x$. There both factors lie within $Cx^{-1/2}(\log x)^{\mathfrak
p_0+3}$ of the operators formed at the trivial triple, which commute with the adjoint action of the
constant gauge transformations, and a linear map of $\mathfrak g$ that commutes with it is a
multiple of $\mathbf 1$; so, up to that error, the factors are numbers formed at the trivial triple
times $\mathrm{Ad}(\mathfrak v(\beta_-)^{-1}\mathfrak v(b_-))$ and $\mathrm{Ad}(\mathfrak
v(b_-)^{-1}\mathfrak v(b'_-))$, the inner factors cancel, and the kernel is $\mathsf
s_0(\beta,\beta')\mathrm{Ad}(\mathfrak v(\beta_-)^{-1}\mathfrak v(\beta'_-))$ up to $Cx^{-1/2}(\log
x)^{\mathfrak p_0+3}$, with a number $\mathsf s_0(\beta,\beta')$ formed at the trivial triple. In
the gauge of the chain the kernel is multiplied on the left by $\mathrm{Ad}((\gamma\mathsf
g_{k+1}^{-1})(\beta_-))$ and on the right by the transpose of $\mathrm{Ad}((\gamma\mathsf
g_{k+1}^{-1})(\beta'_-))$, which turns $\mathrm{Ad}(\mathfrak v(\beta_-)^{-1}\mathfrak v(\beta'_-))$
into $\mathrm{Ad}(\mathfrak u(\beta_-)^{-1}\mathfrak u(\beta'_-))$, $\mathfrak u=\mathfrak v\mathsf
g_{k+1}\gamma^{-1}$. As $(\mathsf U^\gamma_{k+1})^{\mathfrak u}=U_+^{\mathfrak v}=e^{iL^{-k-1}A'}$
with $|A'|\le Cx^{-1/2}(\log x)^{\mathfrak p_0+3}$ on the cube, the ordered product of the values of
$\mathsf U^\gamma_{k+1}$ along a path of bare bonds from $\beta_-$ to $\beta'_-$ of length at most
$3(\mathfrak n+1)LML^{-d}$ lies within $3(\mathfrak n+1)M\sup|A'|$ of $\mathfrak
u(\beta_-)^{-1}\mathfrak u(\beta'_-)$, and within $3(\mathfrak n+1)LM[c_{\mathsf
h}\delta(x)+L^{-d}\mathsf a_{\rm top}]$ of $\mathbf 1$, as $\mathsf U^\gamma_{k+1}$ is the ordered
product of the layers of the depths below $d$ and of the top field, which in its regular gauge has
$|\mathbf A^{\rm top}|\le\mathsf a_{\rm top}$ at each bare bond. So in the gauge of the chain the
kernel at $\beta,\beta'$ is a number times $\mathbf 1$ up to a $\mathsf w_{\mathfrak n}$ that is
affine in $\mathsf a_{\rm top}L^{-d}$; for the pairs at distances $\mathfrak n>n_\flat$ the number
is taken $0$ and $\mathsf w_{\mathfrak n}$ is the bound $Ce^{-\bar\kappa\mathfrak n/2}\le
Cx^{-1}e^{-\bar\kappa(\mathfrak n-n_\flat)/2}$ of the kernel; and $\sum_{\mathfrak n}(24\mathfrak
n^2+2)\mathsf w_{\mathfrak n}\le\epsilon^\Sigma_d$, a number formed from $x$, the constants and
$\mathsf a_{\rm top}L^{-d}$. For $k=0$, where no shift system enters, (ii) and (iii) use the local
pivot construction and its holomorphy and gauge covariance in
equations~\YMMref{(800)}{391}--\YMMref{(803)}{391} in place of the corresponding shift properties:
off the pivot bonds $\mathbf A[U_+;\zeta]=\zeta$ is linear in $\zeta$ and free of the triple; at the
pivot bond $b_0(c)$ of a cube $c$ the pivot value $\beta_c(P,\zeta)$ depends on the block field and
on $\zeta$ through their values on $\mathsf P_c$ alone, so on the triple through its pair alone, the
block field being formed from the configuration of the pair; for $\sup|\zeta|<r^\circ_P$ and a fixed
real factor it is holomorphic in $(A',\mathbf J,\zeta)$ on the simple chart set, which contains the
cores by the common-domain verification in the proof of Theorem~\YMMref{W}{423}, with
$|\beta_c|\le2C_P\sup|\zeta|$; and
$\beta_c(P^u,\mathrm{Ad}(u)\zeta)=\mathrm{Ad}(u(b_0(c)_-))\,\beta_c(P,\zeta)$ for every gauge
transformation $u$ of the bare sites, constant or not. Carried by the gauge $\mathfrak v$ of
$Q_\nu$, the pivot value is therefore the one formed at $P[\omega]$, and
$z\mapsto\beta_c(P[z\omega],\zeta)$ is holomorphic and at most $2C_P\sup|\zeta|$ for $|z|<a^{\rm
us}(\alpha_0,\alpha_1)/r_\nu$, where $P[z\omega]$ lies in the cores; Schwarz's lemma in $z$, with
Cauchy's estimate in $\zeta$, puts its parts of degree one and two in $\zeta$ within a multiple of
$r_\nu$ times their size of the parts formed at the trivial triple, as pair dependence and slice
holomorphy do for $k\ge1$. (iv)

\emph{The covariance.} On the convex box $\mathsf A_W-C_{\rm Hs}\delta(x)\le D^2\Psi_W\le\mathsf
A_W+C_{\rm Hs}\delta(x)$, with $C_{\rm Hs}\delta(x)\le\frac14\gamma_{\rm r}$ as $d>d_0\ge n_2$ (the
bound $C_{\rm Hs}\delta(x)$ of the Hessians of the potentials in the proof of
Lemma~\ref{C-367-7}(iv), which Cauchy's estimate draws from their bound on the whole product of the
analyticity cores and the polydisc of radius $C_A\mathfrak b(x)$, and which so holds at the triple
$P_W$ of every $W\in\mathsf{Cut}_{k+1}$ and at every point of the box); the Brascamp--Lieb
inequality gives $\mathrm{Cov}_W\le(\mathsf A_W-C_{\rm Hs}\delta(x))^{-1}$, and integration by parts
and the Cauchy--Schwarz inequality against linear functions give
$\mathrm{Cov}_W\ge(1-\iota_d)(\mathsf A_W+C_{\rm Hs}\delta(x))^{-1}$, with a share $\iota_d$ of the
boundary of the box that is smaller than every power of $x$; so $\|\mathrm{Cov}_W-\mathsf
A_W^{-1}\|\le\epsilon^{\rm cv}_d=\frac8{\gamma_{\rm r}^2}C_{\rm Hs}\delta(x)+\frac2{\gamma_{\rm
r}}\iota_d$. (3)

\emph{One constant.} The numbers of (2) are formed from $x=L^d\mathsf b$, $\ell L^d$ and the
constants of the construction, $\epsilon^\Sigma_d$ besides from $\mathsf a_{\rm top}L^{-d}$, in
which it is affine; save for the factor $N_k=(\ell L^d)^3$ that multiplies both members of
$\epsilon^{\rm q}_d$ and stays in its bound below, $\ell L^d=(\ell/\mathsf b)x$ enters only through
counts of bonds that multiply masses of boundary layers of the box, which fall faster than every
power of $x$. Collecting the multiples of (2) gives four numbers $C_6$, $C'_6$, $C_\Sigma$, $C_{\rm
cv}$, formed from $\ell$, $\mathsf b$ and the constants, which read neither $K$, $d_0$, $d$, $W$, a
gauge nor $\mathsf a_{\rm top}$, with $\epsilon^{\rm m}_d\le C_6x^{-1}(\log x)^{2\mathfrak p_0}$,
$\epsilon^{\rm q}_d\le C'_6N_kx^{-3}(\log x)^{8\mathfrak p_0+12}$, $\epsilon^\Sigma_d\le
C_\Sigma[x^{-1/2}(\log x)^{\mathfrak p_0+6}+\mathsf a_{\rm top}L^{-d}]$ and $\epsilon^{\rm cv}_d\le
C_{\rm cv}x^{-1/2}(\log x)^{\mathfrak p_0+3}$ at $x=L^d\mathsf b$. From these, (M) and (Col): given
$\mathcal F_{d-1}$ the law of $B_k$ is $p_W$, $W=V_{k+1}$, and $\gamma$ and $\mathsf g_{k+1}$ are
functions of $V_{k+1}$ that act at each bond by isometries, so $|\mathbf n_d(\beta)|=L^d|\int\mathbf
A[U_+;gB](\beta)\,p_W(dB)|$; the factor $L^d$, with $L^dx^{-1}=\mathsf b^{-1}$, gives $C_6\mathsf
b^{-1}(\log L^d\mathsf b)^{2\mathfrak p_0}$ as a bound of $|\mathbf n_d(\beta)|$ for every bare bond
$\beta$; three bare bonds start at every bare site, and the norm of (M) takes at a site the
euclidean norm of the three values, so $\|\mathbf n_d\|_\infty\le\sqrt3\,C_6\mathsf b^{-1}(\log
L^d\mathsf b)^{2\mathfrak p_0}$; and as $\|\mathbf n_d\|_{L^2}^2=a_K^3\sum_\beta|\mathbf
n_d(\beta)|^2$ and $B^k(y)$ holds $L^{3k}$ bare sites, each the initial point of three bare bonds,
the sum over the blocks gives $\|\mathbf n_d\|_{L^2}^2\le3a_K^3L^{3k}L^{2d}\epsilon^{\rm
q}_d\le3C'_6\ell^3\mathsf b^{-3}L^{-d}(\log L^d\mathsf b)^{8\mathfrak p_0+12}$, as
$a_K^3L^{3k}=L^{-3d}$. The kernel of $\mathrm{Ad}(\gamma)\mathsf J_d\mathsf A_W^{-1}\mathsf
J_d^*\mathrm{Ad}(\gamma)^{\mathsf T}$ is that of (iii) in the gauge of the chain, in which the top
field is at most $\mathsf a_{\rm top}$ at each bare bond; so the first two clauses of (Col) hold
with $C_\Sigma$, $\mathbb E^{[d_0]}$-almost surely and at all pairs of bare bonds, and the third
holds with $C_{\rm cv}$ at every $W$. With the four explicit conditions, of which $d_0\ge n_{\rm a}$
and $d_0\ge d_{\rm num}(\ell,\mathsf b)$ follow from $d_0\ge d_{\rm adm}=\max\{n_{\rm a},d_{\rm
num}(\ell,\mathsf b)\}$ and the bounds of $c_\rho\varepsilon(L^{d_0}\mathsf b)$ and of $4c_{\mathsf
h}\delta(L^{d_0}\mathsf b)$ from (1), every top depth $d_0\ge d_{\rm adm}$ is admissible with
$C_{\rm adm}=\max\{\sqrt3\,C_6,3C'_6,C_\Sigma,C_{\rm cv}\}$: this is (Adm).

\emph{Layer estimates.} At an admissible $d_0$, with $C$ free of $K$ and $d_0$, the chain and these statements give: (L1) for $d_0<d\le K$ the ordered product $e^{ia_K\mathbf h'_d}\cdots e^{ia_K\mathbf h'_{d_0+1}}e^{ia_K\mathbf A^{\rm top}}$ of the top field and of the layers of the depths $\le d$ is $\mathsf U_{K-d}^\gamma$, a gauge transform of a background over the block field $V_{K-d}$ at depth $d$; (L2) $|\mathbf h_d|\le c_{\mathsf h}L^d\delta(L^d\mathsf b)$ at every bond, and $\sum_da_K|\mathbf h_d|\le\frac14$; (L3) for a real bond function $f$, every component $\xi$ of $\langle f,\mathbf h'_d-\mathbf n_d\rangle$ has $\mathbb E^{[d_0]}[e^{\lambda\xi}|\mathcal F_{d-1}]\le e^{4\lambda^2G_d(f)^2/\gamma_{\rm r}}$ for real $\lambda$, $G_d(f)^2=c_{\mathsf J}^2\mathsf b^{-1}L^d\min\{\|f\|_1^2,L^{-3d}\|f\|_2^2\}$, by Lemma~\ref{C-367-7}(iv), as the derivative of $\mathbf h'_d$ in the fluctuation obeys on the box the bounds of $\mathsf b^{-1/2}L^{d/2}\mathsf J_d$; (L4) $\sup|\mathbf n_d|\le C\mathsf b^{-1}(\log L^d\mathsf b)^{2\mathfrak p_0}$ by (M), a bound whose sums over $d_0<d'\le d$, times $L^{-d/2}$, tend to $0$ as $d\to\infty$, and almost surely $\|\sum_{d_0<d\le K}\mathbf n_d\|_{L^{12}(T_\ell)}\le\mathsf c_{12}$, a number free of $K$ and $d_0$, as $\|\mathbf n_d\|_{L^{12}(T_\ell)}\le\|\mathbf n_d\|_\infty^{5/6}\|\mathbf n_d\|_{L^2}^{1/6}$ is summable over $d$ by (M); (L5) the moments of the fluctuation fields and of the layers: at every bond $b$, $\mathbb E^{[d_0]}|B_{K-d}(b)|^2\le m_2$, and for $q<\infty$, $\|a_K(\sum_d|\mathbf h_d(b)|)^2\|_{L^q(\mathbb E^{[d_0]})}\le\mathsf c_{(q)}\mathsf b^{-1}$, with $m_2$ and $\mathsf c_{(q)}$ free of $K$, $d$ and $d_0$.

The fluctuation-field second moment was proved above before the mean and covariance estimates. For the layers: (L3) at $f=a_K^{-3}\delta_b$, for which $\langle f,\cdot\rangle$ is the value at $b$ and $G_d(f)^2=c_{\mathsf J}^2\mathsf b^{-1}L^d$, bounds the centred layers at $b$ in $L^{2q}(\mathbb E^{[d_0]})$ by $C\sqrt q\,\mathsf b^{-1/2}L^{d/2}$, whose sum over $d\le K$ is at most $C\sqrt q\,\mathsf b^{-1/2}a_K^{-1/2}$, and the means add at most $2a_K(\sum_{d\le K}C\mathsf b^{-1}(\log L^d\mathsf b)^{2\mathfrak p_0})^2$, a number free of $K$ times $\mathsf b^{-1}$.

In its regular gauge the top field is bounded by $\mathsf a_{\rm top}\le CL^{10d_0}$.

\emph{Passage from the chain measure to the Wilson measure.} The inequality over the top depths: on invariant functions the chain measure from the top depth $d$ has a density $\bar{\mathfrak w}_d\in[0,1]$ with respect to $\mu^K$, with $\int(1-\bar{\mathfrak w}_d)\,d\mu^K\le\mathsf R_d$ (Lemma~\ref{C-367-7}(i)); split $1=\bar{\mathfrak w}_{d_1}+\sum_{d_1<d\le K-1}(1-\bar{\mathfrak w}_{d_1})\cdots(1-\bar{\mathfrak w}_{d-1})\bar{\mathfrak w}_d+(1-\bar{\mathfrak w}_{d_1})\cdots(1-\bar{\mathfrak w}_{K-1})$, bound each product of the sum by its last two factors and the last product by its last factor, and use H\"older's inequality; $\sum_d\mathsf R_d^{1/2}L^{pd}$ and $\sup_K\mathsf R_{K-1}L^{pK}$ are finite, as $\mathsf R_d$ decays faster than every exponential in $d$.

By (Adm) let $d_{\rm adm}$ be a depth, free of $K$, such that every top depth $\ge d_{\rm adm}$ is admissible. At $d_1=d_{\rm adm}$ the moment bound applies at every top depth $d+1$ of the inequality, and $\mathsf N'\le\|\mathbf B_K\|_{\kappa;a}$, $\mathsf N'\le23L^{3K}$: this gives $\mathsf M_q<\infty$. The bound of the action at $R=\max\{1,\mathsf N'\}$ and $\mathsf M_{24q}$ give the moments of $\sup_ss^{\frac32+\kappa}A^a(U_s)$. The good sets: at an admissible $d_0$ with $\mathsf R_{d_0}\le\frac12\epsilon$, on $\{\|\mathbf B_K\|_{\kappa;a}\le R\}$ and for $a_K\le a_0(R)$ the four bounds put $U$ in $\mathsf O_K(\bar C,\bar s)$, $\bar C=(2+c_3)R$, $\bar s=s_0(R)$, and Chebyshev's inequality with Lemma~\ref{C-367-7}(i) gives $\mu^K(\mathsf O_K(\bar C,\bar s)^{\rm c})\le\mathsf R_{d_0}+C(1,\kappa)L^{20d_0}/R$.

With $E=A^a(U_{\sigma/2})$, Lemma~\ref{C-367-13}(f) bounds $a_K^{-4}\max_pa_p(U_\sigma)$ and $(J,J)_a(U_\sigma)$ by polynomials in $E$ if $a_KE\le e_0$, and otherwise by $2a_K^{-4}$ and $48\ell^3a_K^{-6}$, against $\mu^K(a_KE>e_0)\le(a_K/e_0)^{6p}\langle E^{6p}\rangle_{\mu^K}$. Tightness of $\Lambda_\sigma$: on $\mathsf O_K(\bar C,\bar s)$, which has mass $\ge1-\epsilon$, $\Lambda_\sigma$ is bounded by a number free of $K$ (Step~6).

(b) Choose $D_*$ so that for $D\ge D_*$
\[
 \begin{gathered}
 L^{m+D}\ge3R_1\mathfrak M,\qquad L^{m+D}>6\rho_*,\qquad
 2\theta_D\le\bar\varepsilon,\qquad
 2B_3\theta_D\le\alpha_*,\\
 2B_3\theta_DL^{-4}<\min\{\omega_M,(\kappa^{\rm SO}_2)^{-1}\},
 \end{gathered}
\]
where $\alpha_*,\rho_*$ are from
Proposition~\ref{prop:dyn-fixed-coarse-comparison}.
All these constants are fixed, so the conditions eventually hold
because $\theta_D\to0$.
For $V\in\mathcal K_D$, set $\eta=L^{-(K-D)}$ and
$N_D=(2L^{m+D})^3$. Take the minimal backgrounds
$U^{\rm A}=U_{K,D}(V)$ and $U^{\rm B}=U_{K+1,D}(V)$.
Their regularity on the larger open ball
$\{\varpi<2\theta_D\}$ and the proposition give
\[
 (U^{\rm A})^g=e^{i\eta A}\mathcal M(U^{\rm B}),\qquad
 \sup|A|\le2CB_3\theta_DN_D^{1/5}\eta^{1/5},\qquad
 \varpi(U^{\rm B})\le2B_3\theta_D(\eta/L)^2.
\]
The one-step contour estimate, Lemma~\YMMref{14.6}{65}(iii), gives
$\mathcal M(U^{\rm B})(c)=e^{iY_c}U^{\rm B}[c]$,
$|Y_c|\le\kappa_3\varpi(U^{\rm B})$, where $[c]$ subdivides
$c$ into its $L$ fine bonds. A path $C$ of physical length at most
$R$ uses at most $RL^D/\eta$ bonds of the first lattice.
Telescoping the products therefore gives
\[
 |W_C(U^{\rm A})-W_C(U^{\rm B})|
 \le2RL^DB_3\theta_D
   \left(CN_D^{1/5}\eta^{1/5}+\frac{\kappa_3}{L^2}\eta\right).
\]
For each fixed $D$ this is summable in $K$, uniformly in
$V\in\mathcal K_D$ and in the closed paths of length at most $R$
defined at that cutoff. Its constant may depend on $D$ through
$N_D$; no uniformity in $D$ is used in this step. Thus the unflowed
loops have limits $w_C^D(V)$.
For the positive-time assertion we now use
\textup{(Flow approximation)}. Compactness of $\mathcal K_D$ and
that hypothesis give a convergent subsequence of any sequence of
top data and background flows; its limit solves the specified
equation with the specified initial holonomies. Uniqueness
identifies the limiting orbit curve. Uniform convergence of the
potential representatives gives convergence of flowed holonomies:
the difference of their path products is bounded by the path
length times their supremum discrepancy. The sequential assertion
and compactness of $\mathcal K_D$ make this convergence uniform over
$V$ and bounded-length paths on compact positive time intervals. The cutoff-uniform comparison
between the bare field and the background is proved separately
below.

\emph{The bare field and the background over $\pi_{K,D}U$.} By Lemma~\ref{C-367-7}(i) at an admissible $d_0<D$ with $\mathsf R_{d_0}$ small, which the proved conclusion (Adm) provides, the comparison is made under $\mathbb E^{[d_0]}$. There $U$ is a gauge transform of $e^{ia_K\mathbf B_K}$; and $e^{ia_K\mathbf B^{\le D}_K}$, with $\mathbf B^{\le D}_K$ the logarithm of the ordered product of the top and of the layers of the depths $\le D$, the field cut at the depth $D$, a function of $\pi_{K,D}U$, is a gauge transform of $\mathsf U_{K-D}$, that is of the background over $\pi_{K,D}U$. (i) $\mathbf B_K-\mathbf B^{\le D}_K$ is small in the rough distance, in every $L^q$ uniformly in $K$, at the rate $L^{-\kappa D/8}$: its linear part is the sum of the martingale differences of the depths beyond $D$, its bilinear members are those of (a) with a factor of a depth beyond $D$, and its means and the remainders of the product are small in $L^p$.

(ii) The rough norm of the cut field has moments free of $K$ and $D$: its martingale part $u_{\le D}=\sum_{d\le D}u_{(d)}$ is the conditional mean $\mathbb E[u|\mathcal F_D]$ of that of the whole field, so $|u_{\le D}|_{-\frac12-\kappa;a}\le\mathbb E[|u|_{-\frac12-\kappa;a}|\mathcal F_D]$; for a bilinear member $\Upsilon$ of the rough norm (one tuple of directions and translations, one colour, one time, one point) the pairs of different depths whose finer depth exceeds $D$ have vanishing conditional means given $\mathcal F_D$, so $\mathbb E[\Upsilon(u,u)|\mathcal F_D]=\Upsilon(u_{\le D},u_{\le D})+\sum_{d>D}\mathbb E[\mathsf c_d|\mathcal F_D]$, $\mathsf c_d=\mathbb E[\Upsilon(u_{(d)},u_{(d)})|\mathcal F_{d-1}]$ the compensators of (a), whose sum is bounded almost surely by $C(1+\mathsf a_{\rm top})s^{-\frac34-\frac\kappa2}$; with the rest as in (a) and Jensen's inequality, $\mathbb E[\|\mathbf B^{\le D}_K\|^q_{\kappa;a}]^{1/q}\le C(q,\kappa)(1+\mathsf a_{\rm top})^2$.

So for a large $R$ and a large $D$, off a probability that is small uniformly in $K$, the two fields have rough norms $\le R$ and the rough distance $\le\eta$: they are the two data of the Lipschitz bound of (a).

\emph{From the short time to the time $s$.} Let $s_*=\min\{s_{\rm Lip}(R),\frac12s\}$ and $a_K\le a_0(R)$. At $s_*$ the two flowed fields are gauge transforms of $e^{ia_KA}$, $e^{ia_KA'}$ with the four bounds and $\|A-A'\|_\infty\le(1+c_5)\eta s_*^{-\frac14-\frac\kappa2}$; by Lemma~\ref{C-367-13}(c), (e) the actions of the flowed fields at $\frac12s_*$, and those of the fields $e^{ia_K((1-r)A+rA')}$, $0\le r\le1$, are at most $E=6(2+c_3)^4\ell^3R^4(\frac12s_*)^{-\frac32-\kappa}$. For $a_KE\le e_0$, Lemma~\ref{C-367-13}(f), started at $\frac12s_*$, bounds the curvature $a_K^{-2}\varpi$ of the flow of the bare field on $[s_*,s]$ by a number $F_1$, and, started at the time $0$ of the fields of the segment, bounds the number $\Lambda_{s-s_*}$ of their flows by a number $\Lambda_1$; both read $E$ and $s$ alone.

For a flowed observable $H$ of the flow time $s$: by Lemma~\ref{C-367-13}(b), (d), the square field of $H\circ\Phi_{s-s_*}$ is at most $\mathsf L^2$, $\mathsf L$ free of $K$, at the fields of the segment, whose $\ell^2$ length is at most $\sqrt3\ell^{3/2}a_K^{-1/2}\|A'-A\|_\infty$; so the two values of $H$ differ by at most $\sqrt3\ell^{3/2}\mathsf L\|A'-A\|_\infty$. With (i), (ii) this gives the first display: the mean squared difference is at most $3\ell^3\mathsf L^2(1+c_5)^2\eta^2s_*^{-\frac12-\kappa}$ plus $4\sup|\varphi|^2$ times a small probability, uniformly in $K$. For the loops, whose square field is of the order of the length times $a_K^{-2}$, the comparison is made in the supremum distance modulo gauge, in the units of $a_K$, which is at most $\|A-A'\|_\infty$ at $s_*$: by Lemma~\ref{C-367-11}(a) at $\kappa'=F_1a_K^2$ the flowed bare field has at every time of $[s_*,s]$ a gauge with a potential bounded by a number $C_0(F_1,\ell)$; the short-time theory of (a), restarted from these bounded data in steps of a length $\sigma_1(F_1,\ell,\kappa)$, multiplies the distance by at most a factor $C_1$ in each step, and asks nothing of the second field; after $m=\lceil(s-s_*)/\sigma_1\rceil$ steps the factor is $C_1^m$, free of $K$; and $|W_C(V)-W_C(V')|$ is at most the length of $C$ times the distance.

Hence $\lim_D\limsup_K\|W_C\circ\Phi_s-\Psi_D\circ\pi_{K,D}\|_{L^1(\mu^K)}=0$ with $\Psi_D$ free of $K$, which is strong convergence. Two configurations with equal loops are gauge equivalent (characters of $\mathrm{SU}(2)$), and the orbits with $\|\mathfrak A\|_\infty+\mathrm{Lip}\,\mathfrak A\le C$ form a compact set on which the loops give the topology: this gives $\mathfrak F$. The second display follows from the first, from (C3), and from total variation at the depth $D$.

(c) The good sets at the time $s$: by the moments of the action in (a) and Lemma~\ref{C-367-13}(f), started at $\frac14s$, for $\epsilon>0$ there is $F_1$, free of $K$, such that the closed invariant set $\mathsf G_K=\{a_K^{-2}\varpi(\Phi_rU)\le F_1\text{ for }\frac12s\le r\le s\}$ has $\mu^K(\mathsf G_K)\ge1-\epsilon$ for all large $K$. On $\mathsf G_K$, Lemma~\ref{C-367-11}(a) at $\kappa'=F_1a_K^2$ gives at the time $s-\sigma_1$, $\sigma_1=\min\{s_0''(R_1),\frac12s\}$, where $s_0''(R)=(c_1''R^{-2})^{4/(1-2\kappa)}\le s_0(R)$ and $c_1''\in(0,c_1]$ depends only on $\kappa$, a gauge with a potential bounded by $C_0(F_1,\ell)$, of rough norm at most $R_1=\max\{1,C_0+C_0^2\}$, as $\|\mathbf B\|_{\kappa;a}\le\|\mathbf B\|_\infty+\|\mathbf B\|_\infty^2$; the short-time theory of (a) restarted from it gives at the time $s$ a gauge transform $e^{ia_KA}$ of the flowed field with $\|A\|_\infty+\|\nabla A\|_\infty+\|\nabla\nabla A\|_\infty\le C$, a number $C\ge1$ free of $K$: the first two members by the bounds of (a) at the time $\sigma_1\le s_0(R_1)$ after the restart, the second differences by the same fixed point with one more derivative, which closes up to the time $s_0''(R_1)$ and gives $\|\nabla\nabla A_\sigma\|_\infty\le c_8R_1\sigma^{-\frac54-\frac\kappa2}$ for $0<\sigma\le s_0''(R_1)$, $\sigma$ the time after the restart, $c_8$ depending only on $\kappa$; here $\sigma=\sigma_1\le s_0''(R_1)$.

For $a_KC\le\frac1{24}$ and Lipschitz $f$, $|\mathcal O(f)(e^{ia_KA})-\mathcal O^0(f)(\mathfrak A)|\le c\,\ell^3C^5(\|f\|_\infty+\mathrm{Lip}f)\,a_K$ for the interpolated connection $\mathfrak A$; $\mathcal O^0(f)$ is continuous on the compact set of these orbits, so the theorems of Stone and Weierstrass and of Tietze give $\Theta$; the moments come from (a).
\end{proof}

\subsection{Deterministic estimates along the Wilson flow}

\begin{lemma}[Deterministic bounds]\label{C-367-13}
For all $K$, $m$, $\beta>0$, every $U$ and real $X$; $\partial U(p)=c_p\mathbf 1+i\mathsf s_p$, $Y_1,\dots,Y_4$ the values of $X$ at the four bonds of $p$ carried to its base point by the parallel transport of $U$, with the signs $\epsilon=(+,+,-,-)$, and $Z_p(X)=\sum_i\epsilon_iY_i$:
\begin{enumerate}
\item[(a)] $\langle X,\partial a_p\rangle=\langle Z_p(X),\mathsf s_p\rangle$, so $|\partial_ba_p|=|\mathsf s_p|$ for $b\in p$; 
\[
\nabla^2a_p[X,X]=c_p|Z_p(X)|^2+\sum_{i<j}\epsilon_i\epsilon_j\langle i[Y_i,Y_j],\mathsf s_p\rangle,
\]

\[
-3\varpi_p\sum_{b\in p}|X(b)|^2\le\nabla^2a_p[X,X]\le(4+3\varpi_p)\sum_{b\in p}|X(b)|^2,
\]
and 
\[
-12\varpi(U)\|X\|^2\le\nabla^2A_K(U)[X,X]\le40\|X\|^2
\]
 (second derivatives along $e^{isX}U$);
\item[(b)] with $w_s=D\Phi_s(U)X$: $\frac d{ds}\|w_s\|^2=-2a_K^{-2}\nabla^2A_K(U_s)[w_s,w_s]$; so $\|D\Phi_\sigma(U)\|\le e^{12\Lambda_\sigma(U)}$ and $\Gamma_K(G\circ\Phi_\sigma)(U)\le e^{24\Lambda_\sigma(U)}\Gamma_K(G)(U_\sigma)$ for $G\in C^1(\Omega_K)$; $D\Phi_s(U)\mathrm d_U\lambda=\mathrm d_{U_s}\lambda$ and $D\Phi_s(U)J(U)=J(U_s)$, and $D\Phi_\sigma(U)^{\mathsf T}$ maps horizontal vectors at $U_\sigma$ to horizontal vectors at $U$;
\item[(c)] $\partial_\sigma A^a(U_\sigma)=-(J,J)_a(U_\sigma)$; $(J,J)_a(U_\sigma)\le e^{24(\Lambda_\sigma-\Lambda_s)(U)}(J,J)_a(U_s)$ for $s\le\sigma$; $A^a(V)\le\frac32\ell^3(a_K^{-2}\varpi(V))^2$; $A^a(U_{\sigma/2})\le24\ell^3\Lambda_\sigma(U)^2\sigma^{-2}$; and $(J,J)_a(U_\sigma)\le\frac2\sigma e^{24\Lambda_\sigma(U)}A^a(U_{\sigma/2})$;
\item[(d)] $\Gamma_K(\mathcal O(f))\le2\|f\|_\infty^2(J,J)_a+256\|\nabla f\|_\infty^2A^a$ and $|v\,\mathcal O(f)|\le((J,J)_a\Gamma_K(\mathcal O(f)))^{1/2}$;
\item[(e)] on $\mathsf O_K(\bar C,\bar s)$, $a_K^{-2}\varpi(U_s)\le2\bar Cs^{-\frac34-\frac\kappa2}+2\bar C^2s^{-\frac12-\kappa}$ for $0<s\le\bar s$, and so $\Lambda_\sigma\le(8\bar C\sigma^{\frac14-\frac\kappa2}+4\bar C^2\sigma^{\frac12-\kappa})/(1-2\kappa)$ for $\sigma\le\bar s$;
\item[(f)] let $e_0=10^{-4}$, $c_{10}=5.1\cdot10^4$, $x_0=(1440c_{10})^{-1}$, $E=A^a(U_{s_1})$ and $a_KE\le e_0$; at $E=0$ set $E^{-2}=+\infty$, the flat field being stationary. Then for $s\ge s_1$, $a_p(U_s)\le e_0$, and $z_p=a_K^{-4}a_p(U_s)$ satisfies $\partial_sz_p\le\Delta_az_p+c_{10}\bar z_pz_p^{1/2}$ ($\Delta_a$ the Laplacian on parallel plaquettes, $\bar z_p$ the largest $z$ on a box of plaquettes about $p$); $a_K^{-4}\max_pa_p(U_s)\le\frac52E\min\{s-s_1,x_0^4E^{-2},1\}^{-3/2}$ for $s>s_1$; and at $s_1=\frac\sigma2\le\frac12$, $(J,J)_a(U_\sigma)\le e^{48}[\frac{4E}\sigma+7E^{3/2}\sigma^{-3/4}+3x_0^{-3}E^3]$;
\item[(g)] (Bochner.) For $H\in C^\infty(\Omega_K)$, with $\nabla^2$ the Hessian of the bi-invariant metric of $\Omega_K$, a product of unit spheres $S^3$ of Ricci curvature $2$,
\[
\|\mathcal L_KH\|^2_{L^2(\mu^K)}=a_K^{-2}\!\int\|\nabla^2H\|_{\rm HS}^2d\mu^K+2a_K^{-1}\mathcal E_K(H)+a_K^{-2}\beta\!\int\nabla^2A_K(\partial H,\partial H)\,d\mu^K ;
\]
consequently $(\int|\partial H|^4d\mu^K)^{1/2}\le\sup|H|\,(3\mathsf m+2(28\beta\sup|H|\,\mathsf m)^{1/2})$, $\mathsf m=a_K\|\mathcal L_KH\|_{L^2(\mu^K)}$, and $H_K\in C^\infty(\Omega_K)$ with $\sup_K\|\mathcal L_KH_K\|_{L^2(\mu^K)}<\infty$ and $\liminf_K\mathcal E_K(H_K)>0$ have $a_K^{-2}\beta\int\nabla^2A_K(\partial H_K,\partial H_K)\,d\mu^K\to-\infty$; and $(\mathcal L_K\cos G)^2+(\mathcal L_K\sin G)^2=(\mathcal L_KG)^2+\Gamma_K(G)^2$ for real $G\in C^\infty(\Omega_K)$;
\item[(h)] (Weitzenb\"ock and Kato.) $\nabla^2A_K(U)[z,z]=\langle z,(\nabla_U^{\mathsf T}\nabla_U+R_U)z\rangle$ for horizontal $z$ at $U$, where $(\nabla_U^{\mathsf T}\nabla_Uz)(b)=\sum_{b'\sim b}(z(b)-\tau_{bb'}z(b'))$, the sum over the six parallel neighbours with the parallel transports $\tau_{bb'}$ of $U$, and $\|R_U(b,b')\|\le4\varpi(U)\mathsf n(b,b')$, $\mathsf n(b,b')=1$ if $b=b'$ or $b,b'$ lie in a common plaquette and $0$ otherwise. For invariant $G\in C^1(\Omega_K)$ and every bond $b$,
\[
|\partial_b(G\circ\Phi_\sigma)(U)|\le\bigl[e^{4\Lambda_\sigma(U)\mathsf n}\,e^{-\sigma a_K^{-2}\Delta_0}\,|\partial G|(U_\sigma)\bigr](b),
\]
$\Delta_0$ the free Laplacian on the bonds (six parallel neighbours) and $|\partial G|(V)$ the function $b\mapsto|\partial_bG(V)|$; so $\|\partial(G\circ\Phi_\sigma)(U)\|_{\ell^q}\le e^{52\Lambda_\sigma(U)}\|\partial G(U_\sigma)\|_{\ell^q}$ for $1\le q\le\infty$, and $\sup_b|J(U_\sigma)(b)|\le e^{52(\Lambda_\sigma-\Lambda_s)(U)}\sup_b|J(U_s)(b)|$ for $s\le\sigma$;
\item[(i)] (One-bond fibres.) For $N\ge3$, $p\ni b_0$, $s\ge0$, the map $g\mapsto\partial(\Phi_sU^{(g)})(p)$ of $G$ to itself, $U^{(g)}$ the field $U$ with the value $g$ at $b_0$, is onto; so $\sup_U\varpi(U_s)=2$, and $\sup_U\mathcal O(f)(U_s)\ge2f(x(p))/a_K$ for $f\ge0$, at every flow time. At a field with values in the centre $\{\pm\mathbf 1\}$, $\partial A_K=0$ and $\nabla^2A_K[X,X]=\sum_pc_p|Z_p(X)|^2$, $c_p=\pm1$; at the field that is $-\mathbf 1$ at one bond and $\mathbf 1$ elsewhere, $\|D\Phi_\sigma\|\ge e^{4\sigma/a_K^2}$.
\end{enumerate}
\end{lemma}

\begin{proof}[Proof sketch]
(a) Expand $\partial(e^{itX}U)(p)$ to second order in $t$, with the values of $X$ carried to the base point of $p$; the bounds follow from $|\mathsf s_p|\le\varpi_p$, $|c_p|\le1$, $|Z_p|^2\le4\sum_{b\in p}|X(b)|^2$, and four plaquettes at each bond. (b) $\frac d{ds}\|D\Phi_sX\|^2=-2a_K^{-2}\nabla^2A_K(U_s)[D\Phi_sX,D\Phi_sX]$ and Gronwall; the flow commutes with gauge transformations and moves along the current. (c) The current solves the linearised flow; $A^a$ is non-increasing along the flow, $\varpi$ bounds $a_p$, and averaging the bound of the current over $[\frac\sigma2,\sigma]$ gives the last inequality. (d) $\partial_b\mathcal O(f)=a_K^{-1}f(x_b)\partial_bA_K+a_K^{-1}\sum_{p\ni b}(f(x(p))-f(x_b))\partial_ba_p$, $x_b$ a point of $b$, so that $|x(p)-x_b|\le2a_K$ for $p\ni b$, and $\sum_b(\sum_{p\ni b}\varpi_p)^2\le32A_K$.

(e) $\varpi(e^{iaA})\le2a^2(\|A\|_\infty^2+\|\nabla A\|_\infty)$, integrated in $s$. (f) The lattice Bianchi identity in a gauge near $p$ gives the differential inequality where $a_K^4\bar z_p\le e_0$; Duhamel's formula with $\|P_tz\|_\infty\le t^{-3/2}E$ bounds $\sup_tt^{3/2}\max_pz_p$ as long as $t^{1/4}E^{1/2}\le x_0$, and the action does not increase; then (c) on an interval before $\sigma$ on which $\Lambda$ grows by at most $1$. (g) Bochner's formula with the drift; the limit follows from the formula, as its first member on the right is $\ge0$ and the second, $2a_K^{-1}\mathcal E_K(H_K)$, tends to $\infty$; the bound follows from $\int\langle\partial H,\partial F\rangle d\mu^K=-a_K\int F\,\mathcal L_KH\,d\mu^K$ at $F=H|\partial H|^2$ and using $\nabla^2A_K\ge-28$; the identity is the chain rule $\mathcal L_K\phi(G)=\phi'(G)\mathcal L_KG+\phi''(G)\Gamma_K(G)$.

(h) Expanding $\sum_p|Z_p(X)|^2+\sum_x|(\mathrm d_U^{\mathsf T}X)(x)|^2$, the products of parallel bonds give the covariant Laplacian and those of the corners of a plaquette cancel but for $2\sum_p\langle X_1,(1-\mathrm{Ad}\,\partial U(p))X_4\rangle$; with (a) this is the identity. The adjoint of the linearised flow keeps horizontal vectors horizontal, and $|z(b)|$ is then a subsolution of the parabolic system with the off-diagonal entries $\ge0$; the comparison principle and the commutation of $\Delta_0$ with $\mathsf n$ ($\Delta_0$ is the same scalar parallel-neighbour convolution in every bond direction) give the bound; $e^{-t\Delta_0}$ is doubly stochastic and $\mathsf n$ has row sums $13$. The current is a forward horizontal solution, whence its maximum principle.

(i) $s\mapsto(g\mapsto\partial(\Phi_sU^{(g)})(p))$ is a homotopy from a translation of $S^3$, or a translation composed with $g\mapsto g^{-1}$, and a map of $S^3$ to itself that misses a point is homotopic to a constant; at the centre fields $\mathsf s_p=0$, and $X=e\,\delta_{b_0}$ gives $\langle X,\nabla^2A_KX\rangle=-4|e|^2$ at the flipped field.
\end{proof}

\begin{lemma}[Exact identities for the Laplacian along the flow]\label{C-367-21}
For all $K$, $m$, $\beta>0$:
\begin{enumerate}
\item[(a)] for $f\in C^\infty(\Omega_K)$, $[\Delta_K,v]f=-8a_K^{-1}vf-2a_K^{-3}\langle\nabla^2A_K,\nabla^2f\rangle_{\rm HS}$, and the diagonal blocks of $\nabla^2A_K$ are $\delta_{\alpha\beta}\sum_{p\ni b}(1-a_p)$;
\item[(b)] $\Phi_\sigma$ is a diffeomorphism of $\Omega_K$, and $\log\det D\Phi_\sigma(U)=-12a_K^{-2}\int_0^\sigma(3N^3-A_K(U_s))\,ds$;
\item[(c)] at every $U$ at which $\mathrm d_U$ is injective on the site functions (all $U$ off a closed null set), with $\mathsf G_U=(\mathrm d_U^*\mathrm d_U)^{-1}$, $\mathsf l(U)=\frac12\log\det\mathrm d_U^*\mathrm d_U$ and $\mathsf j^{\rm gh}_U(b)=i\sum_{\alpha,\gamma}\mathsf G_U((b_-,\alpha),(b_+,\gamma))[\sigma_\alpha,\mathrm{Ad}(U(b))\sigma_\gamma]$: $\mathsf j^{\rm gh}_U=\partial\mathsf l(U)$, and for invariant $f\in C^2(\Omega_K)$
\[
\Delta_Kf(U)=a_K^{-1}\sum_e\frac{d^2}{dt^2}f(e^{ite}U)\Big|_{t=0}+a_K^{-1}\sum_b\langle\mathsf j^{\rm gh}_U(b),\partial_bf(U)\rangle=\Delta^{\rm hor}_Kf+\Gamma_K(\mathsf l,f),
\]
$(e)$ an orthonormal basis of the horizontal vectors at $U$; for a site function $\lambda$, 
\[
\frac{d^2}{dt^2}f(e^{it\,\mathrm d_U\lambda}U)|_0=\sum_b\langle i[\lambda(b_-),\mathrm{Ad}(U(b))\lambda(b_+)],\partial_bf(U)\rangle;
\]
and at $f=G\circ\Phi_\sigma$ the second member is $a_K^{-1}\langle D\Phi_\sigma(U)\mathsf j^{\rm gh}_U,\partial G(U_\sigma)\rangle$: the gauge term is evaluated once, at the bare field, and then carried by the linearised flow.
\end{enumerate}
\end{lemma}

\begin{proof}[Proof sketch]
(a) On a Riemannian manifold, $[\Delta,\langle\nabla\phi,\nabla\cdot\rangle]f=\langle\nabla\Delta\phi,\nabla f\rangle+2\langle\nabla^2\phi,\nabla^2f\rangle+2\,{\rm Ric}(\nabla\phi,\nabla f)$; here ${\rm Ric}=2$, and $a_K\Delta_KA_K=\sum_{b,\alpha}(\partial^\alpha_b)^2A_K=12(3N^3-A_K)$, because each bond variable enters a plaquette through a matrix coefficient of the fundamental representation, an eigenfunction of the Laplacian of $S^3$ of eigenvalue $-3$, and four bonds enter each of the $3N^3$ plaquettes. (b) Liouville's formula $\frac d{ds}\log\det D\Phi_s=(\mathrm{div}\,v)\circ\Phi_s=-a_K^{-1}(\Delta_KA_K)\circ\Phi_s$. (c) The trace of the Hessian over an orthonormal basis splits into the horizontal part and the part along the gauge directions; along $\mathrm d_U\lambda$ the curve $e^{it\mathrm d_U\lambda}U$ differs from the gauge orbit at second order, which gives the formula for one site function by invariance of $f$; summing over a basis orthonormal for $\mathrm d_U^*\mathrm d_U$ gives $\mathsf j^{\rm gh}$, and differentiating $\log\det$ identifies it with $\partial\mathsf l$.

The last statement is the chain rule with (b) of Lemma~\ref{C-367-13}.
\end{proof}

\subsection{The lattice defect in generator estimates}

\begin{lemma}[A lattice defect]\label{C-367-14}
On $(\mathbb Z/N)^3$ let $\varsigma_p=-1$ on the four plaquettes that contain a bond $b_0$ and $+1$ elsewhere.
For $N\ge N_0$ the operator $\mathrm d^{\mathsf T}\varsigma\,\mathrm d$ on real $1$-forms, per colour the
Hessian of $A_K$ at the field that is $-\mathbf 1$ on $b_0$ and $\mathbf 1$ elsewhere, has exactly three
negative eigenvalues, $-\lambda_1$ simple and $-\lambda_2$ double, $6.6<\lambda_1<6.62$, $0.79<\lambda_2<0.8$,
with orthonormal eigenvectors $e_1$ (of $-\lambda_1$) and $e_2,e_3$ (of $-\lambda_2$) that decay exponentially
away from $b_0$. Let $\hat e_1,\hat e_2$ be the lattice unit vectors of the two axes orthogonal to $b_0$, so
that $b_0+\hat e_1$, $b_0+\hat e_2$ are the two translates of $b_0$ parallel to it that share a plaquette with
it, and the three bonds $b_0$, $b_0+\hat e_1$, $b_0+\hat e_2$ lie in ten plaquettes; the matrix $[e_i(b_0+\hat
e_j)]_{i=1,2,3;\,j=0,1,2}$, $\hat e_0:=0$, is upper triangular and invertible, with an inverse bounded by an
absolute constant. And $\langle X,\mathrm d^{\mathsf T}\varsigma\,\mathrm dX\rangle\ge c_1|\mathrm
dX|^2-7\sum_i\langle e_i,X\rangle^2$ with an absolute $c_1>0$. For $\varepsilon\in(0,1]$ and $\Lambda_0\ge1$
there are $N_*,T_*<\infty$ and $\delta_*,\rho_0,\theta,c_\sharp>0$, $0<\delta\le\frac1{10}$, $C_\sharp<\infty$,
depending only on $\varepsilon$ and $\Lambda_0$, such that for $N=\ell/a_K\ge N_*$, $T\ge T_*$, every bond $b_0$
and every $U$ whose flow in lattice time, $U^\circ_t=\Phi_{a_K^2t}U$, has $\sup_{t\le
T}\varpi(U^\circ_t)\le\delta_*$, $\int_0^T\varpi(U^\circ_t)\,dt\le\Lambda_0$, $T\le\Lambda_0N^2$ and
$A_K(U^\circ_{T/2})\le\delta_*^2$, there is a Borel set $\mathcal S$ of
$\zeta=(\zeta_0,\zeta_1,\zeta_2)\in\mathfrak g^3$, $|\zeta|\le\rho_0\le1$, of Lebesgue measure $\ge c_\sharp
e^{-(3\lambda_1+6\lambda_2+\varepsilon)T}$, such that the field $W_\zeta$, equal to $U$ but for
$W_\zeta(b_0)=-e^{i\zeta_0}U(b_0)$ and $W_\zeta(b_0+\hat e_j)=e^{i\zeta_j}U(b_0+\hat e_j)$ ($j=1,2$), and the
fields $\tilde V_t$, equal to $\Phi_{a_K^2t}W_\zeta$ but for the sign of their value at $b_0$, have the
properties (a)--(c) below. Here $a_p(\Phi_{a_K^2t}W_\zeta)=\varsigma_pa_p(\tilde V_t)+1-\varsigma_p$ for every
plaquette $p$, so that $A_K(\Phi_{a_K^2t}W_\zeta)=8+\sum_p\varsigma_pa_p(\tilde V_t)$, and $t\mapsto\tilde V_t$
is the gradient flow, in lattice time, of $\tilde A=\sum_p\varsigma_pa_p+8$. (a) $\varpi(\tilde V_t)\le\delta$
for $0\le t\le T$: the four plaquettes of $\Phi_{a_K^2t}W_\zeta$ at $b_0$ lie within $\delta$ of $-\mathbf 1$
and all the others within $\delta$ of $\mathbf 1$; (b) for every real function $\mathsf f$ of the plaquettes and
real $\mathsf f_0$,
\[
\bigl|\partial_{b_0}\bigl(\textstyle\sum_p\mathsf
f_p\,a_p\circ\Phi_{a_K^2T}\bigr)(W_\zeta)\bigr|\ge\Bigl[c_\sharp|\mathsf f_0|-C_\sharp\Bigl(\sum_p|\mathsf
f_p-\mathsf f_0|^2e^{-2\theta|p|}\Bigr)^{1/2}\Bigr]e^{(\lambda_1-\varepsilon)T},
\]
$|p|$ the distance of $p$ from $b_0$; in particular $|\partial_{b_0}(A_K\circ\Phi_{a_K^2T})(W_\zeta)|\ge
c_\sharp e^{(\lambda_1-\varepsilon)T}$; (c) $\int_{T/2}^T\sum_pa_p(\tilde
V_t)\,dt\le\frac32T\,A_K(U^\circ_{T/2})+C_\sharp$.

(d) (Traversing the fibre of one bond.) Let $\lambda'_1=\lambda_1-0.027$ and
$\lambda''_2=\lambda_2+\frac1{700}(\lambda_1-\lambda_2)+10^{-3}$; then $\lambda'_1-6\lambda''_2>\frac25$. For
$\Lambda_0\ge1$ there are $\epsilon_{\rm d},\rho_2,c_{\rm d}>0$ and $N_{\rm d}<\infty$, depending only on
$\Lambda_0$, such that the following holds for $N\ge N_{\rm d}$, $T\ge N_{\rm d}$ with $N^3\le c_{\rm d}T^2$,
every bond $b_0$ and every $U$ with $\int_0^T\varpi(U^\circ_t)\,dt+T/N^2\le\Lambda_0$ and $\sup_{t\le
T}\varpi(U^\circ_t)\le\epsilon_{\rm d}$. There is a compact set $\mathsf Y\subset\mathfrak g^2$ of Lebesgue
measure $\ge c_{\rm d}e^{-6\lambda''_2T}$ such that for every $(y_1,y_2)\in\mathsf Y$, with $W$ the field equal
to $U$ but for $W(b_0+\hat e_j)=e^{iy_j}U(b_0+\hat e_j)$ ($j=1,2$) and $W^{(g)}$ the field $W$ with the value
$g$ at $b_0$, and for every unit vector $e\in\mathfrak g$, there are $x_0\in\mathfrak g$, with
$|(x_0,y_1,y_2)|\le30C_\triangle\rho_2\le1$, where $C_\triangle=4c_\triangle$ and $c_\triangle$ is the absolute
constant that bounds the inverse of the matrix $[e_i(b_0+\hat e_j)]$, and
$r_*\in(0,\frac43\rho_2e^{-\lambda'_1T}]$ for which $\mathsf v(r)=A_K(\Phi_{a_K^2T}W^{(g_r)})$,
$g_r=-e^{ire}e^{ix_0}U(b_0)$, has $\mathsf v(0)\in[7,9]$ and $\mathsf v(0)-\mathsf v(r_*)\ge2\rho_2^2$. If $f\in
C^\infty(T_\ell)$ is real, $f(x_*)=\max|f|>0$ and $b_0$ is based at a site nearest to $x_*$, the same holds,
with constants that depend besides on $f$ and for $a_K\le c_{\rm d}(f)$, for $\mathsf
v_f(r)=\sum_pf(x(p))\,a_p(\Phi_{a_K^2T}W^{(g_r)})$ in the place of $\mathsf v$, and with the single condition
$\mathsf v_f(0)-\mathsf v_f(r_*)\ge f(x_*)\rho_2^2$ in the place of the two conditions on $\mathsf v$.
\end{lemma}

\begin{proof}[Sketch]
Flipping the sign of $U(b_0)$ conjugates the flow of $A_K$ to that of $\tilde A=\sum_p\varsigma_pa_p+8$, whose
Hessian at small fields is close to $\mathrm d^{\mathsf T}\varsigma\,\mathrm d$. The number of negative
eigenvalues is found exactly. For $E>0$, with $S=(\mathrm d^{\mathsf T}\mathrm d+E)^{1/2}$ and $\Pi_4$ the
projection onto the four plaquettes, $\mathrm d^{\mathsf T}\varsigma\,\mathrm d+E=S(1-2R^{\mathsf T}R)S$,
$R=\Pi_4\mathrm dS^{-1}$; by Sylvester's law of inertia the number of eigenvalues below $-E$ is the number of
eigenvalues $>1$ of $2RR^{\mathsf T}=K(E)=2\Pi_4[\frac{-\Delta}{-\Delta+E}-\mathrm d_2^{\mathsf T}G_E\mathrm
d_2]\Pi_4$, $\mathrm d_2$ the coboundary of $2$-forms and $G_E=(-\Delta+E)^{-1}$, which decreases in $E$. On
$\mathbb Z^3$, $K(E)$ is diagonal in three patterns of signs on the four plaquettes, taken relative to the signs
of $\mathrm d\delta_{b_0}$: $(1,1,1,1)$, the doublet $(1,0,-1,0)$, $(0,1,0,-1)$, and $(1,-1,1,-1)$, with the
entries
\[
K_1=2-\tfrac E3\bigl(1-EG_E(0)\bigr),\qquad K_2=2-F_2,\qquad
F_2=2E\bigl(G_E(0)+G_E(e_1)\bigr)+4\bigl(G_E(0)-G_E(e_1+e_2)\bigr),
\]
\[
K_3=2\bigl[1-E\bigl(G_E(0)-G_E(e_1)\bigr)-4G_E(0)+8G_E(e_1)-4G_E(e_1+e_2)\bigr],
\]
where $G_E(x)=(2\pi)^{-3}\int_{[-\pi,\pi]^3}e^{ik\cdot x}\bigl(E+\sum_\mu2(1-\cos k_\mu)\bigr)^{-1}dk$ is the
Green function of $-\Delta+E$ on the sites of $\mathbb Z^3$. Its Neumann series
$G_E(x)=\sum_{m\ge0}c_m(x)(6+E)^{-m-1}$, with the number $c_m(x)=\sum_n\binom mn\binom n{(n+x_1+x_2)/2}\binom
n{(n+x_1-x_2)/2}\binom{m-n}{(m-n+x_3)/2}$ of the walks of $m$ steps from $0$ to $x$ (a binomial coefficient
whose lower entry is not an integer between $0$ and its upper entry read as $0$), encloses $G_E(x)$ between its
partial sum up to $m=M$ and that sum plus $E^{-1}(\frac6{6+E})^{M+1}$; at $M=240$ the width is below $10^{-12}$
at the four energies $6.6$, $6.62$, $0.79$, $0.8$. This gives, to the digits shown, $E(1-EG_E(0))=2.997958$ at
$E=6.6$ and $3.002351$ at $E=6.62$, and $F_2=0.99944880$ at $E=0.79$ and $1.00144415$ at $E=0.8$; so
$K_1(6.6)>1>K_1(6.62)$ and $K_2(0.79)>1>K_2(0.8)$. As $E\to0$, $K_1\to2$, $K_2\to2-4(g_0-g_2)$ and
$K_3\to8(g_0-g_2)-\frac23$, $g_0-g_2=G_0(0)-G_0(e_1+e_2)$; since $0\le1-\cos k_1\cos k_2\le(1-\cos k_1)+(1-\cos
k_2)$, the symmetry of the axes gives $0\le g_0-g_2-[G_E(0)-G_E(e_1+e_2)]\le\frac E3G_E(0)$, and the series at
$E=\frac1{10}$ with $M=700$, of width $\le9.3\cdot10^{-5}$, give $0.19195<g_0-g_2<0.19970$; so $K_2(0+)>1.2$ and
$K_3(0+)<0.94$. As the entries decrease in $E$, the third pattern gives no negative eigenvalue, the first gives
the simple $-\lambda_1\in(-6.62,-6.6)$, the doublet gives the double $-\lambda_2\in(-0.8,-0.79)$, and there is
no other. On $(\mathbb Z/N)^3$ the Green function is $\sum_{k\in\mathbb Z^3}G_E(x+Nk)$, which differs from
$G_E(x)$ by at most $E^{-1}\sum_{k\ne0}(\frac6{6+E})^{|x+Nk|_1}$, uniformly small for $E\ge0.7$ and $|x|_1\le2$,
and the entries at $E\to0$ are Riemann sums of those on $\mathbb Z^3$: the strict inequalities persist for $N\ge
N_0$. The eigenvectors: $e_i\propto\mathrm d^{\mathsf T}G_{\lambda_i}\eta_i$, $\eta_i$ the pattern of the mode
on the four plaquettes, so they decay at the rate of $G_E(x)\le E^{-1}(\frac6{6+E})^{|x|_1}$: as
$\frac6{6+\lambda_i}\le\frac6{6.79}<e^{-0.12}$, $|e_i(b)|\le C_ee^{-c_e|b|}$ with $c_e=0.12$, an absolute number
$C_e$ and $|b|$ the distance of the bond $b$ from $b_0$; $e_1(b_0)^2\ge2/\lambda_1$; the doublet patterns are
odd under the reflections across the two planes through $b_0$ that contain $\hat e_2$, resp.\ $\hat e_1$, which
fix $b_0$ and $b_0+\hat e_2$, resp.\ $b_0$ and $b_0+\hat e_1$, so $e_2(b_0)=e_2(b_0+\hat
e_2)=0=e_3(b_0)=e_3(b_0+\hat e_1)$, while $|e_2(b_0+\hat e_1)|=|e_3(b_0+\hat e_2)|\ge0.049$: the matrix is
triangular. \emph{The coercivity.} On the range of $\mathrm d$ the form $\langle\eta,\varsigma\eta'\rangle$ is
that of $B=1-2\mathsf P\Pi_4\mathsf P$, $\mathsf P$ the projection onto the range of $\mathrm d$; the non-zero
eigenvalues of $\mathsf P\Pi_4\mathsf P$ are those of $\frac12K(0+)$, so the spectrum of $B$ lies in
$\{-1,1-K_2(0+),1-K_3(0+),1\}$: $\|B\|\le1$, $|B|\ge c_B=0.06$ for $N\ge N_0$, and $B$ has three negative
eigenvalues. On the span of the $\mathrm de_i$, three-dimensional, $B\le-\frac{\lambda_2}{12}$, as
$\langle\mathrm de_i,\varsigma\,\mathrm de_j\rangle=-\lambda_j\delta_{ij}$ and $|\mathrm dY|^2\le12|Y|^2$. A
self-adjoint $B$ with $\|B\|\le1$, $|B|\ge c_B$ and an $n$-dimensional negative spectral subspace is $\ge
c_-c_B$ on the $B$-orthogonal complement of every $n$-dimensional subspace on which $B\le-c_-$ (compare the two
spectral parts in the norm of $|B|$). For $X=\sum_ia_ie_i+X'$, $a_i=\langle e_i,X\rangle$, one has
$\langle\mathrm dX',\varsigma\,\mathrm de_i\rangle=-\lambda_i\langle X',e_i\rangle=0$, so $\langle X,\mathrm
d^{\mathsf T}\varsigma\,\mathrm dX\rangle=-\sum_i\lambda_ia_i^2+\langle\mathrm dX',\varsigma\,\mathrm
dX'\rangle\ge-\lambda_1\sum_ia_i^2+0.0039(\frac12|\mathrm dX|^2-12\sum_ia_i^2)$: the bound with $c_1=0.0019$, as
$\lambda_1+0.05\le7$. \emph{The unstable coordinates.} Let $\tilde\Phi_t$ be the flow of $\tilde A$ in lattice
time, $\tilde V_t=\tilde\Phi_t\tilde V_0$, $\tilde V_0$ the field $W_\zeta$ with the sign of its value at $b_0$
reversed. The nine coordinates $\xi_i(V)=\sum_p\omega_i(p)\,\mathrm{Ad}(T_p(V))\,\mathsf s_p(V)\in\mathfrak g$,
$i=1,2,3$, are the curvature paired with $\omega_i=-\lambda_i^{-1}\varsigma\,\mathrm de_i$ (so $\mathrm
d^{\mathsf T}\omega_i=e_i$, and $|\omega_i(p)|\le C_ee^{-c_e|p|}$ once $C_e$ is enlarged by an absolute factor),
cut off at the distance $\frac14N$ from $b_0$ and carried to one root along a comb of paths: $T_x(V)$ is the
ordered product of the values of $V$ along the path of the comb from the root to the site $x$, and $T_p=T_x$,
$T_b=T_x$ for a plaquette $p$ and a bond $b$ based at $x$; along a flow line of $\tilde A$,
$\dot\xi_i=\lambda_i\xi_i+r_i$ with $|r_i|\le C(q(\tilde V_t)+N^2e^{-c_eN/4})$,
$q(V)=\sum_pe^{-c|p|}\varpi_p(V)^2$ with $c=\frac12c_e$, a weighted sum of the plaquettes, at most
$C_q\varpi(V)^2$ with an absolute $C_q$. The other directions are controlled by the coercivity bound in the
rough field, through a lattice Weitzenb\"ock identity and Kato's inequality, as long as the nine coordinates
stay in the box $\mathfrak B=\{\max_i|\xi_i|\le\delta_u\}$, with the number $\delta_u>0$ fixed below, say on
$[0,t_1]$, $t_1\le T$; the bound does not grow with $T$, as the unstable coordinates are integrated backwards in
time. Let $D(t)$ be the $\ell^2$ distance of the gauge orbits of $\tilde V_t$ and $U^\circ_t$, an absolutely
continuous function with $D(0)\le|\zeta|\le\rho_0$; at a minimising gauge $\tilde V_t=e^{iY_t}U^\circ_t$, $Y_t$
horizontal at $U^\circ_t$; let $\mathcal D_t=\sum_p|Z_p(Y_t)|^2$ be the energy of the difference and
$\varpi^\circ_t=\varpi(U^\circ_t)$. (i) While $D\le\bar D$, a radius $\le\frac12$, half the derivative of $D^2$
is at most $-\sum_p\varsigma_p|Z_p(Y_t)|^2$ plus the source $2\sum_{p\ni b_0}\langle Z_p(Y_t),\mathsf
s_p(U^\circ_t)\rangle$ of the four plaquettes at $b_0$, a member $28\varpi^\circ_tD^2$ and members of the orders
three and four in $Y_t$; the coercivity bound applies to the first member, the three $\mathfrak g$-valued linear
coordinates of $Y_t$ that it leaves free equal the $\xi_i(\tilde V_t)$ up to errors bounded by $C\varpi^\circ_t$
and by $C(\varpi^\circ_t+D)(\mathcal D_t+(\varpi^\circ_t+N^{-2})D^2)^{1/2}$, and the Weitzenb\"ock identity with
Kato's inequality bounds $\sum_b|Y_t(b)|^4$ by $CD^2(\mathcal D_t+(\varpi^\circ_t+N^{-2})D^2)$; for $\bar D$ and
$\delta_*$ small this gives, with absolute numbers $c'>0$ and $C$, $\frac12\frac d{dt}D^2\le-c'\mathcal
D_t+C|\xi(\tilde V_t)|^2+C(\varpi^\circ_t+N^{-2})D^2+C(\varpi^\circ_t)^2$. (ii) Here the hypotheses
$\int_0^T\varpi^\circ_t\,dt\le\Lambda_0$ and $T\le\Lambda_0N^2$ enter: they give
$\int_0^T(\varpi^\circ_t+N^{-2})\,dt\le2\Lambda_0$, and with $\varpi^\circ_t\le\delta_*$ also
$\int_0^T(\varpi^\circ_t)^2\,dt\le\delta_*\Lambda_0$; so Gronwall's lemma gives, with $t_2\le t_1$ the largest
time with $D\le\bar D$ on $[0,t_2]$, $D(t)^2+2c'\int_0^t\mathcal D_s\,ds\le
e^{4C\Lambda_0}[\rho_0^2+2C\int_0^{t_2}|\xi(\tilde V_s)|^2ds+2C\delta_*\Lambda_0]$ for $t\le t_2$. (iii) The
equation $\dot\xi_i=\lambda_i\xi_i+r_i$ is integrated backwards from $t_2$:
$\xi_i(t)=e^{-\lambda_i(t_2-t)}\xi_i(t_2)-\int_t^{t_2}e^{-\lambda_i(u-t)}r_i(u)\,du$ with
$|\xi_i(t_2)|\le\delta_u$, and Young's inequality gives
$\int_0^{t_2}|\xi_i|^2\le\delta_u^2/\lambda_i+2\lambda_i^{-2}\int_0^{t_2}|r_i|^2$, a bound that does not grow
with $t_2$. (iv) By the bound of $r_i$, $\int_0^{t_2}|r_i|^2\le C[\sup_{[0,t_2]}q\int_0^{t_2}q+N^4e^{-c_eN/2}T]$
with $q=q(\tilde V_t)$, where $q(\tilde V_t)\le C(\mathcal D_t+(\varpi^\circ_t)^2+(\varpi^\circ_t+N^{-2})D^2)$,
$\sup q\le C(\delta_*+\bar D)^2$ and $N^4e^{-c_eN/2}T\le\Lambda_0N^6e^{-c_eN/2}$. (v) Inserting (iii) and (iv)
into (ii), $\mathcal J=\sup_{[0,t_2]}D^2+2c'\int_0^{t_2}\mathcal D_t\,dt$ obeys $\mathcal J\le
e^{C\Lambda_0}[\rho_0^2+\delta_u^2+\delta_*+N^6e^{-c_eN/2}]+e^{C\Lambda_0}(\delta_*+\bar D)^2\mathcal J$; for
$\bar D$ and $\delta_*$ small with $\Lambda_0$, and $N$ large with $\delta_*$, this gives $\mathcal
J\le\frac14\rho^2$ with $\rho^2=C_*(\rho_0^2+\delta_u^2+\delta_*)$, $C_*$ depending only on $\Lambda_0$, and if
$\rho\le\bar D$ the continuity of $D$ gives $t_2=t_1$. So on $[0,t_1]$ the distance of the two flows modulo
gauge stays below $\rho$, small with $\rho_0$, $\delta_u$ and $\delta_*$, and the energy of the difference is
integrable in time, $\int_0^{t_1}\mathcal D_t\,dt\le\rho^2/c'$. \emph{The numbers.} Given $\varepsilon$ and
$\Lambda_0$, the numbers are fixed in this order. First $C_*$ and $\bar D$, which read $\Lambda_0$ alone, and
$C_{**}=(C_*(64C_{\mathsf E}^2+2))^{1/2}$, with $C_{\mathsf E}$ the absolute bound of the inverse of the
triangular matrix. Next $\theta$ and a ceiling $\bar\delta\le\frac1{10}$ for $\delta$, small with $\varepsilon$
and $\Lambda_0$, such that along flows with $\varpi(\tilde V_t)\le\bar\delta$, for $N$ and $T$ large,
$\|\Lambda^9D\tilde\Phi_T\|\le e^{(3\lambda_1+6\lambda_2+\varepsilon)T}$, and the cone condition of (b) holds
with the loss $\varepsilon$ in the rate and leaves outside the strongest mode a part of $w_T$ that is small
against $|x_T|/(1+C_{**})$. Next $\delta_u$, small with all these. Then $\delta_*=\delta_u^2$ and
$\rho_0=8C_{\mathsf E}\delta_u$, so that the radius of (v) is $\rho=C_{**}\delta_u\le\bar D$ and
$\delta=\delta_*+4\rho\le(1+4C_{**})\delta_u\le\bar\delta$. Last $N_*$ and $T_*$, large with all these, so that
for $N\ge N_*$ the number $N^2e^{-c_eN/4}$ is small against $\delta_u$ and $N^6e^{-c_eN/2}\le\delta_*$. These
choices serve three steps below: on the faces of $\mathfrak B$, $|r_i|\le
C(C_q\delta^2+N^2e^{-c_eN/4})\le\frac12\lambda_2\delta_u$, so that the exit is strict; at $t=0$ the map
$\zeta\mapsto\xi(\tilde V_0)$ differs from its linear part, on the ball $|\zeta|\le\rho_0$, by a map with a
Lipschitz constant at most $1/(2C_{\mathsf E})$, so that its image holds the ball of radius $\rho_0/(2C_{\mathsf
E})=4\delta_u$ about $\xi(U)$, which holds $\mathfrak B$ as $|\xi(U)|\le C\delta_*\le\delta_u$; and at the end
$|\partial\tilde A(\tilde V_T)|\le C(\delta_*+\rho)\le C(1+C_{**})\delta_u$, so that the errors in (b) are small
multiples of $\delta_u|x_T|$. \emph{The set.} Let $\mathcal K_T$ be the set of the $\zeta$,
$|\zeta|\le\rho_0=8C_{\mathsf E}\delta_u$, $C_{\mathsf E}$ the bound of the inverse of the triangular matrix,
whose nine coordinates stay in $\mathfrak B$ up to the time $T$. At $t=0$ the map $\zeta\mapsto\xi(\tilde V_0)$
is close to the linear map $\zeta\mapsto(\sum_je_i(b_0+\hat e_j)\mathrm{Ad}(T_{b_0+\hat e_j}(U))\zeta_j)_i$,
invertible by the triangular matrix; on the faces of $\mathfrak B$ the exit is strict, as $\frac
d{dt}\frac12|\xi_i|^2\ge\delta_u(\lambda_2\delta_u-|r_i|)>0$ there; so by a degree argument the end map
$\Xi_T(\zeta)=\xi(\tilde V_T)$ carries $\mathcal K_T$ onto the interior of $\mathfrak B$. Put $\mathcal
S=\mathcal K_T\cap\Xi_T^{-1}(\mathsf S)$, $\mathsf S=\{\frac12\delta_u<|\xi_1|<\delta_u,\
|\xi_2|,|\xi_3|<\delta_u\}$: its image under $\Xi_T$ is the shell $\mathsf S$, of the volume $c\,\delta_u^9$,
and the area formula with $\|\Lambda^9D\tilde\Phi_T\|\le e^{(3\lambda_1+6\lambda_2+\varepsilon)T}$ bounds the
measure of $\mathcal S$ from below. (a) holds on $\mathcal K_T$, with $\delta=\delta_*+4\rho$. \emph{(b) on
$\mathcal S$.} With $w_T=D\tilde\Phi_T(\tilde V_0)[e\,\delta_{b_0}]$, $e\in\mathfrak g$,
\[
\bigl\langle e,\partial_{b_0}\bigl(\textstyle\sum_p\mathsf
f_pa_p\circ\Phi_{a_K^2T}\bigr)(W_\zeta)\bigr\rangle=\mathsf f_0\langle\partial\tilde A(\tilde
V_T),w_T\rangle+\sum_p(\mathsf f_p-\mathsf f_0)\varsigma_p\langle Z_p(w_T),\mathsf s_p(\tilde V_T)\rangle .
\]
By a cone condition, also with the weight $e^{\theta|b|}$, $w_T$ follows the strongest mode:
$|w_T|^2\le3.4|x_T|^2$, $x_T$ its coordinate along $e_1$ carried to the root, the rest of $w_T$ small against
$|x_T|$, and $|x_T|\ge\frac12|e|e^{(\lambda_1-\varepsilon)T}$. At the end the gradient is small,
$|\partial\tilde A(\tilde V_T)|\le C(\delta_*+\rho)$, as $A_K(U^\circ_{T/2})\le\delta_*^2$, and, since
$\varsigma\,\mathrm de_1=-\lambda_1\omega_1$, $\langle\partial\tilde A(\tilde
V_T),w_T\rangle=-\lambda_1\langle\xi_1(\tilde V_T),x_T\rangle$ up to a small multiple of $\delta_u|x_T|$. The
map $e\mapsto x_T(e)$ is onto; for $e$ with $x_T(e)$ a positive multiple of $\xi_1(\tilde V_T)$, of modulus
$>\frac12\delta_u$ on $\mathcal S$: $|\langle\partial\tilde A(\tilde
V_T),w_T\rangle|\ge\frac14\lambda_1\delta_u|x_T|$, while the second sum is at most $21\,\delta(\sum_p|\mathsf
f_p-\mathsf f_0|^2e^{-2\theta|p|})^{1/2}|x_T|$; divide by $|e|$. (c): on $\mathcal S$, up to a gauge, $\tilde
V_t=e^{iY_t}U^\circ_t$ with $Y_t$ horizontal at $U^\circ_t$ and small, and the energy $\mathcal
D_t=\sum_p|Z_p(Y_t)|^2$ of the difference has $\int_0^T\mathcal D_t\,dt$ bounded by a number that depends only
on $\varepsilon$ and $\Lambda_0$; as $\sum_pa_p(e^{iY}U)\le3A_K(U)+\frac32\sum_p|Z_p(Y)|^2+6\sum_p(\sum_{b\in
p}|Y(b)|^2)^2$, with a quartic member that is integrable in $t$ by the same bound and the inequality of Sobolev
on the discrete torus, and as $A_K(U^\circ_t)$ does not increase, the integral over $[\frac T2,T]$ gives (c).
(d): here the unstable part is read through nine coordinates of its own, taken from the velocity of the flow and
cut off at a fixed radius. Let $R$ be a radius that reads $\varepsilon_c$ and $\Lambda_0$ alone, $\chi_R$ a
cutoff of the bonds with values in $[0,1]$, equal to $1$ within the distance $\frac R2$ of $b_0$ and to $0$
beyond $R$, which varies by at most $\pi/R$ across a plaquette, and $\mathsf t_x(V)$ the transport from the root
to the site $x$ along the comb of the coordinates $\xi_i$; with $(\sigma_\alpha)$ an orthonormal basis of
$\mathfrak g$, $\lambda_3=\lambda_2$ and $i,\alpha=1,2,3$, let
$e_{i\alpha}(V)(b)=\chi_R(b)\,e_i(b)\,\mathrm{Ad}(\mathsf t_{b_-}(V))^{-1}\sigma_\alpha$ and
$u_{i\alpha}(V)=\lambda_i^{-1}\langle e_{i\alpha}(V),-\partial\tilde A(V)\rangle$. These nine numbers read the
velocity of the flow along the cut eigenvectors, and they are the nine unstable coordinates of (d):
$u_1=(u_{1\alpha})_\alpha\in\mathbb R^3$ is the coordinate in the strongest mode, and $(u_2,u_3)\in\mathbb R^6$
are the six slow coordinates. Along a flow line of $\tilde A$ the velocity solves the linearised flow exactly,
as the current does in Lemma~\ref{C-367-13}(c), so that $\frac d{ds}u_{i\alpha}(\tilde
V_s)=\lambda_iu_{i\alpha}(\tilde V_s)+\mathrm{err}_{i\alpha}(\tilde V_s)$ with $|\mathrm{err}_{i\alpha}(V)|\le
C_Rq_R(V)+4\lambda_2^{-1}\epsilon_Rq_R(V)^{1/2}$; here $q_R(V)$ is the sum of $\varpi_p(V)^2$ over the
plaquettes within the distance $R+1$ of $b_0$, $\epsilon_R=\max_i\|(\mathrm d^{\mathsf T}\varsigma\,\mathrm
d+\lambda_i)\chi_Re_i\|+\max_i\|(1-\chi_R)e_i\|_{\ell^{6/5}}\le CR^{5/2}(\frac6{6+\lambda_2})^{R/2}$ is the
defect of the cut eigenvectors, and the numbers $C_R$ read $R$ and $\Lambda_0$ alone. For $V'=e^{i\eta}V$ with
$\varpi(V)$ and $\|\eta\|$ small with $R$,
$|\mathrm{err}_{i\alpha}(V')-\mathrm{err}_{i\alpha}(V)|\le[C_R(\varpi(V)+\|\eta\|)+44\lambda_2^{-1}\epsilon_R]\|\eta\|$,
and $|u_{i\alpha}(V')-u_{i\alpha}(V)-\langle e_{i\alpha}(V),\eta\rangle|\le\frac{\varepsilon_c}{100}\|\eta\|$ if
$\varpi(V),\|\eta\|\le\epsilon_\sharp$, a number that reads $R$ and $\varepsilon_c$. Let $\tilde V_s$ and
$\tilde V'_s=e^{i\eta_s}\tilde V_s$ be two flow lines of $\tilde A$ in lattice time along which
$\varpi\le\epsilon_\sharp$ and $n_s=\|\eta_s\|\le\epsilon_\sharp$, and $p_s=|u_1(\tilde V'_s)-u_1(\tilde V_s)|$
the difference of their coordinates in the strongest mode. They obey the cone inequalities $\frac
d{ds}\frac12n^2\le\lambda_2n^2+(\lambda_1-\lambda_2)p^2+\varepsilon_cn^2$ and $\frac
d{ds}\frac12p^2\ge\lambda_1p^2-\varepsilon_cpn$, $\varepsilon_c\le10^{-3}$; so where $n^2\le700p^2$ the
difference $p$ grows at the rate $\ge\lambda_1-\varepsilon_c\sqrt{700}\ge\lambda'_1$, and this cone is never
left, while where $n^2>700p^2$ or $p=0$ the distance $n$ grows at the rate
$\le\lambda_2+\frac1{700}(\lambda_1-\lambda_2)+\varepsilon_c\le\lambda''_2$. As long as the nine unstable
coordinates stay in the ball $\{|u|\le\rho_1\}$, the flow of $\tilde A$ from the three changed bonds stays near
the flow of $A_K$ from $U$, modulo gauge and in $\ell^2$ of all the bonds: by an energy inequality, which holds
while the distance $D(s)$ of the two flows modulo gauge is at most a number $D_0$, $D(s)\le
C'(\Lambda_0)\rho_1+C_R\epsilon_{\rm d}^{1/2}$, with $C'(\Lambda_0)$ free of $\varepsilon_c$ and of the radii;
so $\varpi(\tilde V_s)\le\epsilon_{\rm d}+4D(s)\le5C'(\Lambda_0)\rho_1\le\epsilon_\sharp$ and
$|\mathrm{err}(\tilde V_s)|\le\frac1{12}\lambda_2\rho_1$ in $\mathbb R^9$ there. The values
$(x,y_1,y_2)\in\mathfrak g^3$ at $b_0$ and at the two neighbouring bonds parametrise the nine unstable
coordinates at the time $0$ by a map close to an invertible linear one, triangular through the matrix
$[e_i(b_0+\hat e_j)]$, whose inverse has the norm at most $c_\triangle$. The data are taken in the ball
$\{|(x,y_1,y_2)|\le30C_\triangle\rho_2\}$, of radius $C_\triangle\rho_1=4c_\triangle\rho_1$: there $\varpi$ and
the changes stay below $\epsilon_\sharp$, and the bound on $u_{i\alpha}(V')-u_{i\alpha}(V)$ above, with the
change of the $e_{i\alpha}(V)$ over the ball, small with $\rho_1$, leaves to the difference of the map and its
linearisation at the centre a Lipschitz constant at most $\frac1{2c_\triangle}$ on the ball; so the map carries
the ball homeomorphically onto its image, the images of the points of the boundary sphere lie at the distance
$\ge\frac{C_\triangle\rho_1}{2c_\triangle}=2\rho_1$ from $u(U)$, the image of the centre, and as $|u(U)|\le
C_R\epsilon_{\rm d}\le\rho_1$ the image holds the ball $\{|u|\le\rho_1\}$, $\rho_1=30\rho_2$; on the sphere
$\{|u|=\rho_1\}$ the exit is strict, as $\frac
d{ds}\frac12|u|^2\ge\lambda_2\rho_1^2-\frac1{12}\lambda_2\rho_1^2>0$ there; so by a shooting argument with
Brouwer's theorem, for every $\mathsf w\in\mathbb R^6$ with $|\mathsf w|\le\frac14\rho_1$ there are data in the
ball $\{|(x,y_1,y_2)|\le30C_\triangle\rho_2\}$ whose nine coordinates stay in the ball $\{|u|\le\rho_1\}$ up to
$T$, whose coordinate in the strongest mode vanishes at $T$, and whose six slow coordinates at $T$ equal
$\mathsf w$. Here $e_1(b_0)^2\ge4/\lambda_1>0.604$: as $e_2(b_0)=e_3(b_0)=0$, $\delta_{b_0}=e_1(b_0)e_1+z$ with
$z$ in the orthogonal complement of $e_1,e_2,e_3$, on which the form of $\mathrm d^{\mathsf T}\varsigma\,\mathrm
d$ is $\ge0$, as it has no other negative eigenvalue, and the cross term vanishes, $e_1$ being an eigenvector;
so $-4=\langle\delta_{b_0},\mathrm d^{\mathsf T}\varsigma\,\mathrm
d\,\delta_{b_0}\rangle\ge-\lambda_1e_1(b_0)^2$, the value $-4$ coming from the four plaquettes at $b_0$; thus
$|e_1(b_0)|>0.777$. For two such data let $s_1\le T$ be the last time up to which $n\le\epsilon_\sharp$. On
$[0,s_1]$ the pair stays out of the cone: inside it $p$ would grow from then on, against $p_T=0$ if $s_1=T$, as
the coordinates in the strongest mode both vanish at $T$, and against
$p_{s_1}\ge\epsilon_\sharp/\sqrt{700}>2\rho_1\ge p_{s_1}$ if $s_1<T$. Hence $n_{s_1}\le e^{\lambda''_2s_1}n_0$,
and the slow coordinates at $T$ differ by at most $3e^{\lambda''_2T}n_0$, by $\sqrt6\cdot1.01\,n_T$ if $s_1=T$
and by $2\rho_1<\epsilon_\sharp=n_{s_1}$ if $s_1<T$. At the time $0$ the difference $\eta_0$ lives on the three
changed bonds; $p_0\le n_0/\sqrt{700}$, $n_0\le|\eta_0(b_0)|+|\eta_0(b_0+\hat e_1)|+|\eta_0(b_0+\hat e_2)|$ and
$p_0\ge|e_1(b_0)||\eta_0(b_0)|-|\eta_0(b_0+\hat e_1)|-|\eta_0(b_0+\hat e_2)|-\varepsilon_cn_0$, so
$0.73|\eta_0(b_0)|\le1.04(|\eta_0(b_0+\hat e_1)|+|\eta_0(b_0+\hat e_2)|)$ and $n_0\le2.5(|\eta_0(b_0+\hat
e_1)|+|\eta_0(b_0+\hat e_2)|)\le4|y-y'|$, $y$, $y'$ the values of the data at the two neighbouring bonds: the
slow coordinates at $T$ differ by at most $12e^{\lambda''_2T}|y-y'|$; a Lipschitz map changes Lebesgue measure
at most by the sixth power of its constant, which gives the measure of $\mathsf Y$. On the fibre of $b_0$
through such a datum, along the geodesic $g_r$, the two flows differ at the time $0$ by
$\eta_0=re\,\delta_{b_0}$, so $p_0\ge(|e_1(b_0)|-\varepsilon_c)r\ge0.75\,r\ge0.75\,n_0$: the difference in the
strongest mode starts at $\ge0.75\,r$, inside the cone, and grows at the rate $\ge\lambda'_1$, with
$n_s\le1.35p_s$ by the same two inequalities; let $r_*$ be the first $r$ at which it reaches $\rho_2$ at the
time $T$ (meanwhile $n_s\le1.35\rho_2<\epsilon_\sharp$, and $|u|\le\frac13\rho_1+3.1n_s<\frac12\rho_1$, where
$\frac13\rho_1$ bounds $|u|$ along the datum by the equation of $u$ read backwards from $T$), so that
$r_*\le\frac43\rho_2e^{-\lambda'_1T}$. At the time $T$ the velocity of the flow is orthogonal to the strongest
mode, $\langle\partial\tilde A(\tilde V_T),e_{1\alpha}(\tilde V_T)\rangle=-\lambda_1u_{1\alpha}(\tilde V_T)=0$,
as the coordinate of the datum in that mode vanishes there, and it is of the order $\rho_2$ with a constant of
$\Lambda_0$ alone: $\|\partial\tilde A(\tilde V_T)\|\le\|\partial A_K(U^\circ_T)\|+16\epsilon_{\rm d}+44D(T)$,
$D(T)$ the distance of the two flows modulo gauge, where $\|\partial
A_K(U^\circ_T)\|^2\le48e^{24\Lambda_0}\Lambda_0^2N^3T^{-3}$ (the condition $N^3\le c_{\rm d}T^2$ enters here as
well) and $D(T)\le C'(\Lambda_0)\rho_1+C_R\epsilon_{\rm d}^{1/2}$; so $\|\partial\tilde A(\tilde
V_T)\|\le45\,C'(\Lambda_0)\rho_1=1350\,C'(\Lambda_0)\rho_2$, $\rho_2=\rho_1/30$, by the last condition on
$\epsilon_{\rm d}$ below. The numbers are chosen in this order: $C'(\Lambda_0)$;
$\varepsilon_c=\min\{10^{-3},(4000\,C'(\Lambda_0))^{-2}\}$; the radius $R$, which reads $\varepsilon_c$ and
$\Lambda_0$; $\epsilon_\sharp$, which bounds $\varpi$ and $n$ wherever the cone inequalities are used, and
$D_0$, the distance modulo gauge up to which the energy inequality holds, both depending on $R$ and
$\varepsilon_c$; $\rho_1$, small with $R$, with $5C'(\Lambda_0)\rho_1\le\epsilon_\sharp$,
$2\sqrt{700}\,\rho_1<\epsilon_\sharp$ and $5C_\triangle\rho_1\le\min\{1,\epsilon_\sharp\}$, so that
$30C_\triangle\rho_2\le1$ and, as $\epsilon_{\rm d}\le\frac15\epsilon_\sharp$ by the conditions on
$\epsilon_{\rm d}$ below, $\varpi$ and the changes stay below $\epsilon_\sharp$ on the ball of the data;
$\rho_2=\rho_1/30$; and last $\epsilon_{\rm d}$, $N_{\rm d}$ and $c_{\rm d}$, which depend on everything chosen
before: small enough for the energy inequality, and such that $C_R\epsilon_{\rm d}\le\rho_1$, $\epsilon_{\rm
d}+4C_R\epsilon_{\rm d}^{1/2}\le C'(\Lambda_0)\rho_1$, $\|\partial A_K(U^\circ_T)\|+16\epsilon_{\rm
d}+44C_R\epsilon_{\rm d}^{1/2}\le C'(\Lambda_0)\rho_1$ and $0.8\,e^{-11N_{\rm d}}\le\varepsilon_c$. At $r=r_*$
the ratio $\theta_s=n_s^2/p_s^2$ is defined on $[0,T]$, as $p_s>0$ there, and the cone inequalities give
$\theta'\le-2(\lambda_1-\lambda_2)(\theta-1)+2\varepsilon_c\theta(1+\theta^{1/2})$; as
$\lambda_1-\lambda_2>5.8$, the right side is at most $-11(\theta-1)+10\varepsilon_c$ where $1\le\theta\le2$, and
$\theta_0\le0.75^{-2}<1.8$, so that $\theta$ stays below $1.8$ and
$\theta_T-1\le0.8\,e^{-11T}+\varepsilon_c\le2\varepsilon_c$, by $T\ge N_{\rm d}$ and the last condition on
$N_{\rm d}$. Let $\hat e_\alpha=e_{1\alpha}(\tilde V_T)/\|\chi_Re_1\|$, three orthonormal vectors, and
$\eta_T=\sum_\alpha c_\alpha\hat e_\alpha+\eta_T^\perp$ with $\eta_T^\perp$ orthogonal to them; the bound on
$u_{i\alpha}(V')-u_{i\alpha}(V)-\langle e_{i\alpha}(V),\eta\rangle$ above, with $\|\chi_Re_1\|\le1$, gives
$|c|\ge p_T-\varepsilon_cn_T$, whence
$\|\eta_T^\perp\|^2=n_T^2-|c|^2\le(\theta_T-1)p_T^2+2\varepsilon_cp_Tn_T\le5\varepsilon_cp_T^2$, and
$p_T=\rho_2$ at $r=r_*$. Thus the displacement $\eta_T$ at $r=r_*$ lies along the strongest mode up to the part
$\eta_T^\perp$, of norm $\le(5\varepsilon_c)^{1/2}\rho_2$, so the member of the first order of the Taylor
expansion of $\tilde A$ at the end of the flow, which sees only that part, is at most
$1350\,C'(\Lambda_0)\rho_2\,(5\varepsilon_c)^{1/2}\rho_2\le0.76\,\rho_2^2$; the Hessian of $\tilde A$ on that
mode is close to $-\lambda_1$, and its quadratic form at every $X$ is at most $40\|X\|^2$, by the bounds of
Lemma~\ref{C-367-13}(a) at each plaquette, so that $\eta_T^\perp$ moves the member of the second order by at
most $0.2\,\rho_2^2$, and this member is at most $-3\rho_2^2$: the fall is at least $2\rho_2^2$. The value
$\mathsf v(0)$ is $8$ plus a number of modulus at most $\sum_pa_p$ of the field with the reversed sign at $b_0$,
and this sum is $\le1$ by the bound used for (c), as the two flows are close and
$A_K(U^\circ_T)\le24\Lambda_0^2N^3T^{-2}\le24\Lambda_0^2c_{\rm d}$, the action being non-increasing and
$\varpi(U^\circ_s)\le4\Lambda_0/T$ at some $s\in[\frac T4,\frac T2]$. For a general $f$ the same steps hold,
with $8000$ in the place of $4000$ in $\varepsilon_c$: the plaquettes within a fixed distance of $b_0$ change
$f(x(p))$ by $O(a_K)$ from $f(x_*)$, and the others see only the part of $\eta_T$ outside the strongest mode,
which is small; this gives the fall $f(x_*)\rho_2^2$.
\end{proof}

\subsection{Regularisation cutoffs and lower energy bounds}

\begin{lemma}[Condenser cutoffs]\label{C-367-15}
For all $K$, $m$, $\beta>0$, let $\eta\in D(\mathcal E_K)$, $0\le\eta\le1$, $\mathsf A=\{\eta=0\}$, $\mathsf B=\{\eta=1\}$, $D_0=\{\phi\in D(\mathcal E_K):\phi=0\text{ a.e.\ on }\mathsf A\cup\mathsf B\}$, $\mathcal E_{K,1}(\phi,\psi)=\mathcal E_K(\phi,\psi)+\int\phi\psi\,d\mu^K$. There is exactly one $w\in\eta+D_0$ with $\mathcal E_{K,1}(w,\phi)=0$ for $\phi\in D_0$; with $\varkappa=\mathcal E_{K,1}(w,w)$: $0\le w\le1$, $w=0$ on $\mathsf A$, $w=1$ on $\mathsf B$, $\Gamma_K(w,\cdot)=0$ a.e.\ on $\mathsf A\cup\mathsf B$, $\varkappa\le\mathcal E_{K,1}(\eta,\eta)$, $\int w\,d\mu^K\le\varkappa$, and $|\mathcal E_K(w,\phi)|\le\mathrm{osc}(\phi)\,\varkappa$ for every $\phi\in D(\mathcal E_K)\cap L^\infty$; $w$ is invariant if $\eta$ is.

With $q(u)=3u^2-2u^3$: $q(w)\in D(\mathcal E_K)$, $q(w)=0$ on $\mathsf A$, $q(w)=1$ on $\mathsf B$, $\mathcal E_K(q(w))\le\frac94\varkappa$, and $\mathcal E_K(q(w),\phi)=-\int\mathsf l\,\phi\,d\mu^K$ for $\phi\in D(\mathcal E_K)\cap L^\infty$, with $\mathsf l=6w^2(1-w)+6(1-2w)\Gamma_K(w)$, $\|\mathsf l\|_{L^1(\mu^K)}\le\frac{15}2\varkappa$.
\end{lemma}

\begin{proof}[Proof sketch]
$w$ is the element of least norm of the closed convex set $\{v\ge1\text{ on }\mathsf B,\ v\le0\text{ on }\mathsf A\}$; truncation does not raise the norm, so $0\le w\le1$. The variational inequality $\mathcal E_{K,1}(w,\psi)\ge0$ for $\psi\ge0$ on $\mathsf B$, $\psi\le0$ on $\mathsf A$, applied to $\psi=\phi\eta$, $(M-\phi)\eta$, $-\phi(1-\eta)$, $-(M-\phi)(1-\eta)$ for $0\le\phi\le M=\mathrm{osc}\,\phi$, gives
\[
M\Bigl(\int w\,d\mu^K-\varkappa\Bigr)\le\mathcal E_{K,1}(w,\phi)\le M\varkappa,\qquad0\le\int w\phi\,d\mu^K\le M\int w\,d\mu^K,
\]
so $|\mathcal E_K(w,\phi)|=|\mathcal E_{K,1}(w,\phi)-\int w\phi\,d\mu^K|\le M\varkappa$. Gauge transformations preserve $\mu^K$, $\Gamma_K$ and the constraint set, so uniqueness gives invariance. For $q(w)$, apply the equation of $w$ to $6w(1-w)\phi\in D_0$: $\Gamma_K(w,6w(1-w)\phi)=6w(1-w)\Gamma_K(w,\phi)+6(1-2w)\phi\,\Gamma_K(w)$; and $6w^2(1-w)\le\frac32w$, $\int\Gamma_K(w)\le\varkappa$.
\end{proof}

\begin{proposition}[Gauge-invariant regularisation cutoffs]\label{C-367-22}
For all $K$, $m$: let $0<\bar s\le1$ with $64\bar s^{7/16}\le\ell$ and $4a_K^2\le\bar s$. For a site $y$ and an integer $k\ge1$ let $Q(y,k)$ be the cube of $k^3$ sites with lowest corner $y$, $\mathsf e_k(V;y)=a_K^{-1}\sum\varpi_p(V)^2$ over the plaquettes with base point in $Q(y,k)$, $\mathsf e^*_k=\max_y\mathsf e_k(\cdot;y)$, $\mathfrak s_k(V)=a_K^3\sum_y\mathsf e_k(V;y)^8$, $k(r)=\lceil r^{7/16}/a_K\rceil$, and
\[
\mathscr N(U)=\int_{a_K^2}^{\bar s}r\,\mathfrak s_{16k(r)}(U_r)\,dr ,
\]
a smooth invariant function defined without choosing a gauge.
\begin{enumerate}
\item[(a)] If $\varpi(V)\le1$, $1\le k<N$ and $X$ is horizontal at $V$, then $\nabla^2A_K(V)[X,X]\ge\frac14\sum_p|Z_p(X)|^2-\hat\lambda_k(V)\|X\|^2$ with $a_K^{-2}\hat\lambda_k(V)=16[16(144C)^3\mathsf e^*_{k+1}(V)^2+(144(a_Kk)^2)^{-1}]$, $C$ absolute: on horizontal vectors the Hessian is bounded below by the square of the energy in windows of side $a_Kk$, with no supremum of the curvature and no gauge.
\item[(b)] Let $R>0$ and $K$ so large that $e^{64}(C_1R)^{1/8}a_K^{11/64}\le\frac1{144}$, $C_1$ absolute. If $\mathscr N(U)\le R$, then $\varpi(U_r)\le\frac1{12}$ for $0\le r\le\bar s$; $\Gamma_K(G\circ\Phi_\sigma)(U)\le e^{2\bar\Lambda(R)}\Gamma_K(G)(U_\sigma)$ for invariant $G\in C^2(\Omega_K)$ and $\sigma\le\bar s$, $\bar\Lambda(R)=1+c_\Lambda\bar s^{11/64}R^{1/4}+\bar s^{1/8}$, $c_\Lambda$ absolute; for $4a_K^2\le\sigma\le\bar s$, $A^a(U_{\sigma/2})\le\mathsf A_\sigma(R)=\frac12\cdot16^{-3}(\frac\sigma4)^{-25/16}\ell^{21/8}R^{1/8}$ and $(J,J)_a(U_\sigma)\le\frac2\sigma e^{2\bar\Lambda(R)}\mathsf A_\sigma(R)$; so $\Gamma_K(H^\sigma_K)$ and $|\partial_\sigma H^\sigma_K|$ are bounded on $\{\mathscr N\le R\}$ by numbers free of $K$ for $4a_K^2\le\sigma\le\bar s$, and for all $\sigma\ge4a_K^2$ if besides $2a_K\mathsf A_{\bar s}(R)\le1$; and $\Gamma_K(\mathscr N)\le512\,e^{2\bar\Lambda(R)}\mathscr N'$, with $\mathscr N'$ the same integral with $r^2(16\rho(r))^3a_K^3\sum_y\mathsf e_{16k(r)}(\cdot;y)^{15}$, $\rho(r)=a_Kk(r)$, in the place of $r\,\mathfrak s_{16k(r)}$.
\item[(c)] Assume (E), so that (Adm) of Proposition~\ref{C-367-12} holds, and let $\epsilon>0$. There are $\bar s$, $R$, $\bar K$, and for every $q<\infty$ a number $\mathsf C_q$, none depending on $K$, such that for $K\ge\bar K$ the function $\chi_K=\chi_0(\mathscr N/R)$, $\chi_0\in C^\infty(\mathbb R)$, $\chi_0=1$ on $(-\infty,\frac12]$, $\chi_0=0$ on $[1,\infty)$, $0\le-\chi_0'\le3$, is smooth and invariant, vanishes off $\{\mathscr N\le R\}$, is $1$ on the set $\mathsf O_K(\bar C,\bar s)$ of Proposition~\ref{C-367-12}(a) formed at $\kappa=\frac1{64}$ and this $\epsilon$, so that $\mu^K(\chi_K<1)\le\epsilon$, and has $\int\Gamma_K(\chi_K)^qd\mu^K\le\mathsf C_q$.
\item[(d)] In the setting of (c), for a flowed observable $H=H^\sigma_K$ at a fixed $\sigma>0$, $K\ge\bar K(\epsilon,\sigma)$, $\mathsf h=\sup_{\{\mathscr N\le R\}}\Gamma_K(H)$, $\mathsf h'=\sup_{\{\mathscr N\le R\}}|\partial_\sigma H|$: the vector field $\zeta_K=\chi_K\partial H$ has $a_K^{-1}\int|\zeta_K|^2d\mu^K\le\mathsf h$ and the divergence $\chi_K\Delta_KH+\mathsf b\,\chi_K\partial_\sigma H+\Gamma_K(\chi_K,H)$; for $F\in C^1(\Omega_K)$, $|\int F\chi_K\Delta_KH\,d\mu^K|\le\mathsf h^{1/2}\mathcal E_K(F)^{1/2}+[(\mathsf h\mathsf C_1)^{1/2}+\mathsf b\mathsf h']\|F\|_{L^2(\mu^K)}$; $|\int\chi_K(\mathsf b\partial_\sigma+\Delta_K)H\,d\mu^K|\le(\mathsf h\mathsf C_1)^{1/2}\mu^K(\frac12R<\mathscr N<R)^{1/2}$; every sequence of cutoffs has a subsequence along which $\chi_{K}H^\sigma_{K}$ converges strongly to an element of $D(\mathcal E_\star)$; and $\|\chi_KH-g^K_n\circ\pi_{K,n}\|^2_{L^2(\mu^K)}\le C\,a_n^2/\mathsf b+4\sup|\varphi|^2\mathsf R_n$ uniformly in $K$, $g^K_n$ the top of the tower of $\chi_KH$, $n\ge n''$.
\end{enumerate}
\end{proposition}

\begin{proof}[Proof sketch]
(a) As $\varpi(V)\le1$, $c_p\ge\frac12$, and Lemma~\ref{C-367-13}(a) gives $\nabla^2A_K(V)[X,X]\ge\frac12\sum_p|Z_p(X)|^2-3\sum_p|\mathsf s_p|\,|X|_p^2$, $|X|_p^2=\sum_{b\in p}|X(b)|^2$. With $f(x)=(\sum_\mu|X_\mu(x)|^2)^{1/2}$ and $t(x)$ the largest $|\mathsf s_p|$ over the plaquettes with the corner $x$, the remainder is at most $12\mathcal T$, $\mathcal T=\sum_xt(x)f(x)^2$. For horizontal $X$ the lattice Weitzenb\"ock identity of Lemma~\ref{C-367-13}(h), its remainder bounded plaquette by plaquette, gives $\|\nabla X\|^2\le2\sum_p|Z_p(X)|^2+36\mathcal T$, and Kato's inequality $\|\nabla f\|\le\|\nabla X\|$. In the windows $Q(y,k)$, each site lying in $k^3$ of them: H\"older's inequality and the inequality of Sobolev on a discrete cube, $\|f\|^2_{\ell^6(Q)}\le C_{\rm Sob}(\|\nabla f\|^2_{2,Q}+k^{-2}\|f\|^2_{\ell^2(Q)})$, $C_{\rm Sob}=864$, give $\mathcal T\le E^{1/2}\|f\|^{1/2}\mathsf B^{3/4}$, $E=4a_K\mathsf e^*_{k+1}(V)$, $\mathsf B=C_{\rm Sob}(\|\nabla f\|^2+k^{-2}\|f\|^2)\le C_{\rm Sob}(2\sum_p|Z_p(X)|^2+36\mathcal T+k^{-2}\|X\|^2)$; the $\ell^2$ member of the inequality of Sobolev is the source of the term $(144(a_Kk)^2)^{-1}$.

Young's inequality at $(144C_{\rm Sob})^{-1}$ gives $\frac34\mathcal T\le\frac1{72}\sum_p|Z_p(X)|^2+[(144C_{\rm Sob})^3E^2+(144k^2)^{-1}]\|X\|^2$, so $12\mathcal T\le\frac29\sum_p|Z_p(X)|^2+\hat\lambda_k(V)\|X\|^2$, and $\frac12-\frac29>\frac14$; the absolute $C$ of the statement is $C_{\rm Sob}$. (b) Let $\mathscr N(U)\le R$.

\emph{The window energies.} At every field a local energy inequality holds, from the derivative of a cut-off window energy along the flow: $\mathsf e_k(U_r;y)\le\mathsf e_{k+2h}(U_{r'};y-h)+64(a_Kh)^{-2}\int_{r'}^r\mathsf e_{k+2h}(U_\tau;y-h)\,d\tau$ for $r'<r$, $y-h=y-a_Kh(1,1,1)$. At $k=16k(r)$, $h=8k(r)$, the mean over $r'\in[\frac r2,r]$, the bound $\mathsf e^*_k(V)^8\le8^9(a_Kk)^{-3}\mathfrak s_k(V)$ (a window is covered by eight windows of the same side whose corners run over a sub-cube of side $\frac k2$) and Jensen's inequality give $\mathsf e^*_{16k(r)}(U_r)^8\le C_1r^{-53/16}R$ for $2a_K^2\le r\le\bar s$, directly from $\mathscr N(U)\le R$. So $\varpi(U_r)^2\le a_K\mathsf e^*_{16k(r)}(U_r)\le(C_1R)^{1/8}a_K^{11/64}$ there, and for $r<2a_K^2$, as $\varpi(U_{r'})^2\le e^{32|r-r'|/a_K^2}\varpi(U_r)^2$, with the factor $e^{64}$: by the condition on $K$, $\varpi(U_r)\le\frac1{12}$ on $[0,\bar s]$.

\emph{The floor in two regimes.} For horizontal $X$ at $U_r$, $a_K^{-2}\nabla^2A_K(U_r)[X,X]\ge\frac14a_K^{-2}\sum_p|Z_p(X)|^2-\lambda(r)\|X\|^2$ with $\lambda(r)=12a_K^{-2}\varpi(U_r)$ for $r<a_K^2$ (Lemma~\ref{C-367-13}(a), with $c_p\ge\frac12$; its integral over $[0,a_K^2]$ is $\le1$) and $\lambda(r)=a_K^{-2}\hat\lambda_{k(r)}(U_r)$ for $a_K^2\le r\le\bar s$ (by (a)); H\"older's inequality with the exponents $4,\frac43$ in $r$ bounds the integral of the window member by $c_\Lambda\bar s^{11/64}R^{1/4}$, and the member $(144\rho(r)^2)^{-1}$, $\rho(r)=a_Kk(r)\ge r^{7/16}$, contributes $\frac89\bar s^{1/8}\le\bar s^{1/8}$: $\int_0^{\bar s}\lambda(r)\,dr\le\bar\Lambda(R)$.

\emph{The bounds.} Lemma~\ref{C-367-13}(b)--(d), read on horizontal vectors with $\lambda$ in the place of $12a_K^{-2}\varpi$, give the bounds of $\Gamma_K(G\circ\Phi_\sigma)$, of $(J,J)_a$, of $\Gamma_K(H^\sigma_K)$ and of $\partial_\sigma H^\sigma_K$; $\mathsf A_\sigma(R)$ comes from $A^a=\frac12(16\rho)^{-3}a_K^3\sum_y\mathsf e_{16k}$, H\"older's inequality in $y$, the monotonicity of $A^a$ along the flow and the mean over $[\frac\sigma4,\frac\sigma2]$. For continuation at later fixed times use the action-only horizontal Hessian bound from Proposition~\ref{C-367-22}(a), as in Step~6 below: $a_KE\le\frac12$ gives $\max_pa_p\le\frac12$ along the later flow and a lower Hessian bound by $-c(E^2+\ell^{-2})$. This is distinct from the smaller small-action threshold in Lemma~\ref{C-367-13}(f).

\emph{The square field of $\mathscr N$.} With $M_r(V)=\int_r^{\bar s}w(\tau)\,\mathfrak s_{16k(\tau)}(\Phi_{\tau-r}V)\,d\tau$, $w(\tau)=\tau\mathbf 1_{\tau\ge a_K^2}$, and $\Xi_r=\partial M_r(U_r)$, horizontal, $\Xi_0=\partial\mathscr N(U)$, $\Xi_{\bar s}=0$: $\frac d{dr}\|\Xi_r\|^2=2a_K^{-2}\nabla^2A_K(U_r)[\Xi_r,\Xi_r]-2w(r)\langle\Xi_r,\partial\mathfrak s_{16k(r)}(U_r)\rangle$. The gradient of $\mathfrak s_{\hat k}$ is a covariant divergence: $\langle\Xi,\partial\mathfrak s_{\hat k}\rangle=\sum_p\mathsf f(p)\langle\mathsf s_p,Z_p(\Xi)\rangle\le\mathsf G_r^{1/2}(\sum_p|Z_p(\Xi)|^2)^{1/2}$, $\mathsf f(p)=16a_K^2\sum_{y:\,{\rm base}(p)\in Q(y,\hat k)}\mathsf e_{\hat k}(\cdot;y)^7$, $\mathsf G_r\le256a_K^{-1}(16\rho)^3a_K^3\sum_y\mathsf e_{\hat k}(\cdot;y)^{15}$; it is paid by the member $\frac14\sum_p|Z_p|^2$ of the floor, $2w\mathsf G_r^{1/2}(\sum_p|Z_p(\Xi)|^2)^{1/2}\le\frac12a_K^{-2}\sum_p|Z_p(\Xi)|^2+2a_K^2w^2\mathsf G_r$.

Gronwall's inequality from $\bar s$ back to $0$ gives 
\[
\|\partial\mathscr N(U)\|^2\le2a_K^2e^{2\bar\Lambda(R)}\int w^2\mathsf G_r\,dr\le512\,a_Ke^{2\bar\Lambda(R)}\mathscr N'(U),
\]
that is $\Gamma_K(\mathscr N)\le512\,e^{2\bar\Lambda(R)}\mathscr N'$. (c) $\{\mathscr N\le R\}$ contains $\mathsf O_K$, on which the window energies are bounded by the sizes of the flowed field; at each dyadic scale $r_j$ the square field of the functional is bounded on the good set of that scale by a positive power of $r_j$, and off it crudely by $r_j^{-1}$ (by $a_K^{-12}$ at the last scale) against a mass $\le t_Mr_j^M$ for every $M$, from the moments of Proposition~\ref{C-367-12}(a) and the super-exponential decay of $\mathsf R_d$; the sum over the dyadic scales converges in every $L^q$.

(d) The divergence is the product rule and (C1) of Theorem~\ref{C-367-5}; the bounds are Cauchy--Schwarz; the convergence is (A2) applied to $\chi_KH$, whose energy is at most $2\mathsf h+2\sup|\varphi|^2\mathsf C_1$; the last bound is the key estimate of Step~1 applied to $\chi_KH$.
\end{proof}

\begin{lemma}[Lower bounds for the flowed action]\label{C-367-16}
Assume (E), so that (Adm) of Proposition~\ref{C-367-12} holds; let $\mathbb P_K$ be the normalised chain measure from the least top depth $d_0$ with $\mathsf R_{d_0}\le\frac12$ such that every top depth $\ge d_0$ is admissible, $\mathbb E_K$ its mean, $\mathbf B_K$ and $\mathsf a_{\rm top}$ as in Proposition~\ref{C-367-12}, and $\|\cdot\|_2$ the norm of $L^2(T_\ell)$. For a real function $f$ and a $\mathfrak g$-valued function $h$ of the bare bonds, $\langle f,h\rangle=a_K^3\sum_bf(b)h(b)\in\mathfrak g$, $\|f\|_1=a_K^3\sum_b|f(b)|$, $\|f\|_2^2=a_K^3\sum_b|f(b)|^2$. For a function $Y$ of the bonds, $\mathrm dY(p)=Y(b_1)+Y(b_2)-Y(b_3)-Y(b_4)$ is the plain sum around the plaquette $p$, with the signs $(+,+,-,-)$ of Lemma~\ref{C-367-13} and no factor of length; for a $1$-form $A$ per unit length, $\mathrm d_aA=a_K^{-2}\mathrm d(a_KA)$ is its lattice curl, $(\mathrm d_aA)(p)=\partial^+_\mu A_\nu(x)-\partial^+_\nu A_\mu(x)$ for $p=(x;\mu\nu)$, $\|\mathrm d_aA\|_2^2=a_K^3\sum_p|\mathrm d_aA(p)|^2$, and $\mathrm d_aP_s=P_s\mathrm d_a$.
\begin{enumerate}
\item[(a)] There are $k_4=100\,c_{\mathsf J}\gamma_{\rm r}^{-1/2}$, $c_{\mathsf J}$ the constant of the layer maps of Proposition~\ref{C-367-12}, $p_\ast=9/(64\Lambda_+^2k_4^4)$, $k_\Sigma=15/(16\Lambda_+)$ and $\theta_0>0$, depending only on the constants of the construction, and $d_L$, free of $K$, with $\mathbb P_K(\|\mathrm d_aP_s\mathbf B_K\|_2^2\ge k_\Sigma\mathsf b^{-1}\ell^3a_d^{-3})\ge p_\ast$ for $d\ge d_L$, $K>d$ with $a_K\mathsf a_{\rm top}\le\frac18$, and $0\le s\le\theta_0a_d^2$. Let $k=K-d$ and $\mathbf Q^{\rm lin}_k$ the $k$-fold linear average at $\mathbf 1$ from the bare bonds to the bonds of depth $d$, the composite of $k$ maps \YMMref{(42)}{22}.

For a plaquette $P$ of depth $d$ none of whose bonds is a pivot bond, 
\[
X_P=\mathrm d\,\mathbf Q^{\rm lin}_k(a_KP_s\mathbf B_K)(P)\in\mathfrak g,
\]
the plain sum around $P$, with no factor of length, is $\langle P_sf_P,\mathbf B_K\rangle$ with the real bond function $f_P=a_K^{-2}q_P$, $q_P$ the kernel of $\mathsf a\mapsto\mathrm d\,\mathbf Q^{\rm lin}_k\mathsf a(P)$, $\|f_P\|_1\le4a_d$, $\|f_P\|_2^2\le4a_d^{-1}$; and $X_P$ has $\mathbb E_K|X_P|^2\ge\frac3{4\Lambda_+}x_k^{-1}$ and $\mathbb E_K|X_P|^4\le k_4^4x_k^{-2}$; at least $3(1-4L^{-3})(\ell L^d)^3$ of the plaquettes of depth $d$ have no pivot bond. Conditionally: let $d_0$ be any admissible top depth with $\mathsf R_{d_0}\le\frac12$, $\mathbb P_K$ the normalised chain measure from $d_0$, and $\mathbb E^0_K$ its conditional mean given the block field at depth $d_0$.

There are $d'_L>d_0$ and numbers $\tau_d\to0$, which depend on $d_0$ but not on $K$, such that for $d\ge d'_L$, $K>d$ with $a_K\mathsf a_{\rm top}\le\frac18$ and $0\le s\le\theta_0a_d^2$ there is an event $\mathcal T_d$ of the block field at depth $d_0$ with $\mathbb P_K(\mathcal T_d)\ge1-\tau_d$ on which, $N_P$ being the number of the plaquettes of depth $d$ without pivot bond and the sums running over them, $N_P^{-1}\sum_P\mathbb E^0_K|X_P|^2\ge\frac3{4\Lambda_+}x_k^{-1}$, $N_P^{-1}\sum_P\mathbb E^0_K|X_P|^4\le k_4^4x_k^{-2}$, and 
\[
\mathbb P_K(\|\mathrm d_aP_s\mathbf B_K\|_2^2\ge k_\Sigma\mathsf b^{-1}\ell^3a_d^{-3}\,|\,\text{block field at depth }d_0)\ge p_\ast.
\]

\item[(b)] Let $R\ge\max\{1,\|\mathbf B\|_{\kappa;a}\}$, $\bar C=(2+c_3)R$, $a_K\le a_1(R)$, $0<\sigma\le\min\{s_0(R),s_{\rm lin}(R)\}$, $a_K\bar C\sigma^{-1/4-\kappa/2}\le\frac18$, $P_\sigma\mathbf B\ne0$, $U=e^{ia_K\mathbf B}$, where $c_3$, $s_0(R)$, $a_0(R)$ are those of Proposition~\ref{C-367-12}(a) and $\mathsf M_\Theta$ that of its head, $s_{\rm lin}(R)=(2c_6\bar C^2)^{-4/(1-2\kappa)}$, $a_1(R)=\min\{a_0(R),(8e^{12}R)^{-2/(1-2\kappa)},(8(4\mathsf M_\Theta)^{1/6}c_3R)^{-1/(1-2\kappa)}\}$, a positive number that reads $R$ and $\kappa$ alone, and $c_6$, $c'$, $c_g$ are absolute. For $a_K\le a_1(R)$ the gauge-fixed flow $A_s$ from $\mathbf B$ has $a_K\|A_s\|_\infty\le\frac12$ on the closed interval $[0,s_0(R)]$, so that the chart and the linearised equation serve from $s=0$ on: $a_K\|\mathbf B\|_\infty\le e^{12}Ra_K^{\frac12-\kappa}\le\frac18$, as $\|f\|_\infty\le e^{12}\|P_{a_K^2}f\|_\infty$; on $[0,a_K^2/(4\mathsf M_\Theta)]$ the chart variable $a_KA_s$ moves at most at the speed $a_K^{-2}\mathsf M_\Theta$, so that $a_K\|A_s\|_\infty\le\frac18+\frac14$ there; later the first bound of Proposition~\ref{C-367-12}(a) gives $a_K\|A_s-P_s\mathbf B\|_\infty\le(4\mathsf M_\Theta)^\kappa c_3Ra_K^{1-2\kappa}\le(4\mathsf M_\Theta)^{1/6}c_3Ra_K^{1-2\kappa}\le\frac18$, and with $\|P_s\mathbf B\|_\infty\le\|\mathbf B\|_\infty$, $a_K\|A_s\|_\infty\le\frac14$.

Then
\[
\Gamma_K(A^a\circ\Phi_\sigma)(U)^{1/2}\ge\frac{\|\mathrm d_aP_{3\sigma/2}\mathbf B\|_2^2}{\|P_\sigma\mathbf B\|_2}-c'\ell^{3/2}\bar C^3\sigma^{-1-\kappa},\qquad\|P_\sigma\mathbf B\|_2\le\sqrt3\ell^{3/2}R\sigma^{-\frac14-\frac\kappa2},
\]
$\Gamma_K(A^a\circ\Phi_\sigma)(U)^{1/2}\le\frac87[\sqrt2c_g\sigma^{-1/2}\|\mathrm d_aP_\sigma\mathbf B\|_2+c'\ell^{3/2}\bar C^3\sigma^{-1-\kappa}]$, and $|(2A^a(U_\sigma))^{1/2}-\|\mathrm d_aP_\sigma\mathbf B\|_2|\le\sqrt3\ell^{3/2}(2\bar C+8\bar C^2)\sigma^{-\frac12-\kappa}$.
\item[(c)] Let $\mathfrak b>0$, $p=e^{-\Psi}dB$ a probability measure on the box $\prod_b\{|B(b)|<\mathfrak b\}$ with $\Psi\in C^2$ and $0\le D^2\Psi\le\Lambda_+$, $\partial_e$ the derivative in one coordinate of one bond $b_1$, and $\eta(B)=\phi(|B(b_1)|/\mathfrak b)$ with $\phi=1$ on $[0,\frac12]$ and $\phi(t)=\cos^2(\pi(t-\frac12))$ on $[\frac12,1]$. Then $\mathrm{Var}_p(u)\ge(\int\eta\,\partial_eu\,dp)^2/(\Lambda_++\pi^2\mathfrak b^{-2})$ for every bounded Lipschitz $u$.
\end{enumerate}
\end{lemma}

\begin{proof}[Proof sketch]
(c) Integrate by parts against $\eta$, which vanishes at the faces in the direction $e$, so there is no boundary term: $\int\eta\,\partial_eu\,dp=\int(u-\bar u)(\eta\partial_e\Psi-\partial_e\eta)\,dp$, and bound the second factor in $L^2(p)$ through $D^2\Psi\le\Lambda_+$ (one more integration by parts) and $|\partial_e\eta|\le\pi/\mathfrak b$. (a)

\emph{Signs.} They are those of the head of the lemma: the pairing $\langle f,\mathbf h\rangle=a_K^3\sum_bf(b)\mathbf h(b)\in\mathfrak g$ of a real bond function $f$ with a $\mathfrak g$-valued $\mathbf h$ on $\Lambda_{K,\ell}$; for a bond function $Y$ without a factor of length, its plain sum $\mathrm dY(p)=Y(b_1)+Y(b_2)-Y(b_3)-Y(b_4)$ around the plaquette $p$, with the signs of Lemma~\ref{C-367-13}, the bonds taken as in $\partial U(p)$; and for a $1$-form $A$ per unit length, its lattice curl $\mathrm d_aA=a_K^{-2}\mathrm d(a_KA)$, with $\|\mathrm d_aA\|_2^2=a_K^3\sum_p|\mathrm d_aA(p)|^2$ and $\mathrm d_aP_s=P_s\mathrm d_a$. Besides, $\|f\|_r^r=a_K^3\sum_b|f(b)|^r$ for $1\le r<\infty$.

The roman $\mathrm d$ and $\mathrm d_a$ are these sums around a plaquette, and the italic $d$ is a depth. For a plaquette $P$ of depth $d$, $k=K-d$, $X_P$ and $f_P$ are those of (a): $X_P=\mathrm d\,\mathbf Q^{\rm lin}_k(a_KP_s\mathbf B_K)(P)$ is the sum around $P$ of the $k$-fold linear average $\mathbf Q^{\rm lin}_k$ of equation~\YMMref{(42)}{22}, \YMMref{(42)}{22}, applied to the bond function $a_KP_s\mathbf B_K$, which carries no factor of length, and $f_P$ is the real bond function with $\mathrm d\,\mathbf Q^{\rm lin}_k(a_KA)(P)=\langle f_P,A\rangle$ for every $1$-form $A$, so that $X_P=\langle P_sf_P,\mathbf B_K\rangle$. The numbers $k_4$, $p_\ast$, $k_\Sigma$ of (a), and the number $\theta_0$ with $k_1$, are
\[
\begin{multlined}
\theta_0=\frac3{16\Lambda_+k_1^2},\qquad k_4=100\,c_{\mathsf J}\gamma_{\rm r}^{-1/2},\qquad p_\ast=\frac9{64\Lambda_+^2k_4^4},\qquad k_\Sigma=\frac{15}{16\Lambda_+},\qquad k_1 \\
=48\,c_\nabla\Bigl(\frac{24\,c_{\mathsf J}^2}{\gamma_{\rm r}(1-L^{-1})}\Bigr)^{1/2},
\end{multlined}
\]
with the absolute constant $c_\nabla$ of the bound $\|\partial^\pm_\nu P_sf\|_1\le c_\nabla s^{-1/2}\|f\|_1$ of the heat kernel and the constant $c_{\mathsf J}$ of the layer maps of Proposition~\ref{C-367-12}, which enters (L3) below; they depend only on the constants of the construction.

\emph{The linear average.} In the sum of \YMMref{(42)}{22} around a coarse plaquette the contour averages cancel, so one step maps the curl of a fine bond function to $L^{-3}$ times the sum, over the sites $x$ of the block, of its curls at the $L^2$ fine plaquettes of the square of side $L$ with the corner $x$: a matrix with the row sums $L^2$ and the column sums $L^{-1}$. By the Cauchy--Schwarz inequality, $\sum_P|\mathrm d\,\mathbf Q^{\rm lin}_k\mathsf a(P)|^2\le L^k\sum_p|\mathrm d\mathsf a(p)|^2$, and at $\mathsf a=a_KA$, as $\sum_p|\mathrm d(a_KA)(p)|^2=a_K\|\mathrm d_aA\|_2^2$ and $L^ka_K=a_d$,
\begin{equation}\label{C-367-31}
\|\mathrm d_aA\|_2^2\ \ge\ a_d^{-1}\sum_P\bigl|\mathrm d\,\mathbf Q^{\rm lin}_k(a_KA)(P)\bigr|^2 :
\end{equation}
the curls of the linear averages bound $\|\mathrm d_aA\|_2^2$ from below, with the constant one between the $L^2$ norms at the spacings $a_K$ and $a_d$. Iterated, $\mathrm d\,\mathbf Q^{\rm lin}_k\mathsf a(P)$ is the mean, over the $L^{3k}$ sites $u$ of the block, of the sum of $\mathsf a$ around the square of side $L^k$ with the corner $u$; so $a_K^2f_P$ is a signed sum of four tents along the bonds of $P$, $\|f_P\|_1\le4a_d$, $\|f_P\|_2^2\le4a_d^{-1}$ and $\sum_\nu\|\partial^+_\nu f_P\|_1\le24$; these norms do not grow under $P_s$, $\|P_sf_P\|_{12/11}\le\|P_sf_P\|_1^{5/6}\|P_sf_P\|_2^{1/6}\le4a_d^{3/4}$, and $(1-P_s)f=-\int_0^s\sum_\nu\partial^-_\nu P_r\partial^+_\nu f\,dr$ gives $\|(1-P_s)f_P\|_1\le48c_\nabla\sqrt s$.

\emph{The plaquette of a block field, by (c).} Let $W\in\mathsf{Cut}_{k+1}$ and let $P$ be a plaquette of level $k$ none of whose bonds is a pivot bond, $b_1$ its first bond. Its four bonds are coordinates of the chart, so $\partial V_W(B)(P)=e^{ix_k^{-1/2}B(b_1)}C$ with $C$ free of $B(b_1)$. Write $\partial V_W(B)(P)=c_P\mathbf 1+i\mathsf s_P$, let $(t_\alpha)$ be an orthonormal basis of $\mathfrak g$ and $\mathsf E=\int_0^1\mathrm{Ad}(e^{iux_k^{-1/2}B(b_1)})\,du$; the derivative of $\langle\mathsf s_P,t_\alpha\rangle$ in the coordinate $B(b_1)^\alpha$ is $x_k^{-1/2}\,\frac12\mathrm{tr}\bigl(t_\alpha(\mathsf Et_\alpha)\,\partial V_W(B)(P)\bigr)\ge x_k^{-1/2}(1-\delta(x_k)-\theta_d)$, as $|\mathsf Et_\alpha-t_\alpha|\le x_k^{-1/2}|B(b_1)|<\delta(x_k)$ on the box and $|\partial V_W(B)(P)-\mathbf 1|<\theta_d$ on the cut: in the variable $B$ of $p_W$, in which $D^2\Psi_W\le\Lambda_+$ holds (Lemma~\ref{C-367-7}(iv)), the derivative carries the chart factor $x_k^{-1/2}$.

(c) at $p=p_W$, $\mathfrak b=\mathfrak b(x_k)$ and $u=\langle\mathsf s_P,t_\alpha\rangle$, with $\int\eta\,dp_W\ge p_W(|B(b_1)|\le\frac12\mathfrak b(x_k))$, summed over the three $\alpha$ and integrated over $W=V_{k+1}$ with Jensen's inequality, and $\mathbb P_K(|B_k(b_1)|>\frac12\mathfrak b(x_k))\le4m_2\mathfrak b(x_k)^{-2}$, $m_2$ the bound, free of $k$, of the second moments of one fluctuation field in (L5) of the proof of Proposition~\ref{C-367-12}(a), recalled below (its Gaussian tails, Lemma~\ref{C-367-7}(iv)), give, as $|\partial V_W(B)(P)-\mathbf 1|\ge|\mathsf s_P|$,
\[
\mathbb E_K|\partial V_k(P)-\mathbf 1|^2\ \ge\ \frac{3x_k^{-1}\varsigma_d^2}{\Lambda_++\pi^2\mathfrak b(x_k)^{-2}},\qquad\varsigma_d=\bigl(1-\delta(x_k)-\theta_d\bigr)_+\bigl(1-4m_2\mathfrak b(x_k)^{-2}\bigr)_+\to1\quad(d\to\infty):
\]
the upper bound of the Hessian of one fluctuation integral keeps the plaquette of the block field from degenerating. There are $3(\ell L^{d-1})^3$ pivot bonds, each in four plaquettes, so at most $12L^{-3}(\ell L^d)^3$ of the $3(\ell L^d)^3$ plaquettes of level $k$ hold one.

\emph{The layers.} By (Adm), as in the head of the lemma, $d_0$ and every larger top depth are admissible. By Proposition~\ref{C-367-12} (the field of the gauge of the chain), $\mathbf B_K=\frac1{ia_K}\log\bigl(e^{ia_K\mathbf h'_K}\cdots e^{ia_K\mathbf h'_{d_0+1}}e^{ia_K\mathbf A^{\rm top}}\bigr)$, the logarithm of the ordered product of the top field, $|\mathbf A^{\rm top}|\le\mathsf a_{\rm top}$, and of the layers $\mathbf h'_{d'}=\mathrm{Ad}(\gamma)\mathbf h_{d'}$, $d_0<d'\le K$, with the regular gauge $\gamma$ of the top, a function of the top field, and $\mathbf h_{d'}$ a function of the block field at depth $d'$; let $\mathcal F_{d'}$ be the $\sigma$-algebra of that block field, which increases with $d'$, and $\mathbf n_{d'}=\mathbb E_K[\mathbf h'_{d'}|\mathcal F_{d'-1}]$.

The facts on the layers in the proof of Proposition~\ref{C-367-12}(a) give: (L1) the ordered product of the top field and of the layers of the depths $\le d$ is a gauge transform of a background over the block field $V_k$, $k=K-d$, and the $k$-fold average of a background over $V_k$ is $V_k$ (Definitions~\YMMref{9.3}{40} and~\YMMref{9.5}{41}, Section~\ref{sec:analytic-core-complements}); (L2) $|\mathbf h'_{d'}|=|\mathbf h_{d'}|\le c_{\mathsf h}L^{d'}\delta(L^{d'}\mathsf b)$ at every bond, with the constant $c_{\mathsf h}$ of the construction, and $\sum_{d'}a_K|\mathbf h'_{d'}|\le\frac14$, as $d_0$ is admissible; (L3) for a real bond function $f$, every component $\xi$ of $\langle f,\mathbf h'_{d'}-\mathbf n_{d'}\rangle$ has $\mathbb E_K[e^{\lambda\xi}|\mathcal F_{d'-1}]\le e^{4\lambda^2G_{d'}(f)^2/\gamma_{\rm r}}$ for real $\lambda$, $G_{d'}(f)^2=c_{\mathsf J}^2\mathsf b^{-1}L^{d'}\min\{\|f\|_1^2,L^{-3d'}\|f\|_2^2\}$: the log-Sobolev inequality of one fluctuation integral (Lemma~\ref{C-367-7}(iv)), carried by the map from the fluctuation field to the layer, whose derivative on the box obeys the bounds of the layer map with the constant $c_{\mathsf J}$; on $L^2(T_\ell)$ this is the covariance bound $Ca_{d'}^2/\mathsf b$; (L4) $\sup|\mathbf n_{d'}|\le\mathsf n_{d'}=C\mathsf b^{-1}(\log L^{d'}\mathsf b)^{2\mathfrak p_0}$, by the bounds (M), with $L^{-d/2}\mathsf s_d\to0$, $\mathsf s_d=\sum_{d_0<d'\le d}\mathsf n_{d'}$, and almost surely $\|\sum_{d_0<d'\le K}\mathbf n_{d'}\|_{L^{12}(T_\ell)}\le\mathsf c_{12}$, a number free of $K$; (L5) at every bond $b$, $\mathbb E_K|B_{K-d'}(b)|^2\le m_2$, and the moments of the layers are bounded so that $\|a_K(\sum_{d'}|\mathbf h'_{d'}(b)|)^2\|_{L^q(\mathbb P_K)}\le\mathsf c_{(q)}\mathsf b^{-1}$, with $m_2$ and $\mathsf c_{(q)}$ free of $K$ and $d'$.

For $\sum_j|x_j|=s'\le\frac12$ the product $e^{ix_1}\cdots e^{ix_m}$ of exponentials of elements of $\mathfrak g$ is $e^{iy}$ with $|y|\le1.02s'$ and $|y-\sum_jx_j|\le1.02s'^2$; so, with $a_K\mathsf a_{\rm top}\le\frac18$ and (L2), the remainder of the product formula, $\mathbf B_K-\mathbf A^{\rm top}-\sum_{d'}\mathbf h'_{d'}$, is at most $10a_K[(\sum_{d'}|\mathbf h'_{d'}|)^2+\mathsf a_{\rm top}^2]$ at every bond.

\emph{The field of the depths $\le d$.} Let $d>d_0$ and let $\mathbf B^{\le d}$ be the logarithm of the ordered product of the top field and of the layers of the depths $\le d$. At every bare bond the arguments of the exponentials have the sum $s'=a_K\mathsf a_{\rm top}+a_K\sum_{d'\le d}|\mathbf h'_{d'}|\le L^{-k}\bar y_d$, $\bar y_d=a_d\mathsf a_{\rm top}+c_{\mathsf h}\delta(L^d\mathsf b)/(1-L^{-1/2})$, by (L2) and $L^{d'}\delta(L^{d'}\mathsf b)\le L^{-(d-d')/2}L^d\delta(L^d\mathsf b)$ for $d'\le d$; so $y=L^k\sup a_K|\mathbf B^{\le d}|\le2\bar y_d$, and the remainder $\mathbf B^{\le d}-\mathbf A^{\rm top}-\sum_{d'\le d}\mathbf h'_{d'}$ is at most $1.02a_d^{-1}\bar y_d^2$ at every bond.

By (L1) and the gauge covariance of the average, $\mathcal M(U^u)=\mathcal M(U)^{u|}$ (Proposition~\YMMref{4.5}{20}, \YMMref{(40)}{20}), $\mathcal M^k(e^{ia_K\mathbf B^{\le d}})$ is a gauge transform of $V_k$, so its plaquette at $P$ has the modulus $|\partial V_k(P)-\mathbf 1|$. For $L^k\sup|\mathsf a|<c_Q$, $\mathcal M^k(e^{i\mathsf a})(c)=e^{iQ_k(\mathsf a)(c)}$ with $|Q_k(\mathsf a)(c)|\le C_QL^k\sup|\mathsf a|$ and $|Q_k(\mathsf a)(c)-\mathbf Q^{\rm lin}_k\mathsf a(c)|\le4C_Qc_Q^{-1}(L^k\sup|\mathsf a|)^2$: the linearisation of the average, with constants that depend on $L$ alone (Proposition~\YMMref{51.1}{390} and Lemma~\ref{lem:plain-average-remainder}).

At $\mathsf a=a_K\mathbf B^{\le d}$, with the product formula for the four factors around $P$,
\[
\bigl|\langle f_P,\mathbf B^{\le d}\rangle\bigr|\ \ge\ |\partial V_k(P)-\mathbf 1|-4\bigl(8C_Q^2+16C_Qc_Q^{-1}\bigr)\,\bar y_d^2 ,
\]
a bound of the modulus alone, which is the quantity that (c) bounds from below; the error tends to $0$ as $d\to\infty$, uniformly in $K$.

\emph{The heat kernel at the depths $\le d$.} Let $f^1=(1-P_s)f_P$, $\|f^1\|_1\le48c_\nabla\sqrt s$. Of the parts of $\mathbf B^{\le d}$, the top field, the conditional means $\mathbf n_{d'}$ and the remainder are bounded at every bond by $\mathsf a_{\rm top}$, $\mathsf s_d$ and $1.02a_d^{-1}\bar y_d^2$, and pair with $f^1$ to at most $\|f^1\|_1$ times these; the centred layers $\mathbf h'_{d'}-\mathbf n_{d'}$, $d'\le d$, are martingale differences for $(\mathcal F_{d'})$, so orthogonal, and by (L3) at $f=f^1$ their pairing with $f^1$ has the $L^2(\mathbb P_K)$ norm at most
\[
\Bigl(3\cdot\frac8{\gamma_{\rm r}}\,c_{\mathsf J}^2\mathsf b^{-1}\|f^1\|_1^2\sum_{d'\le d}L^{d'}\Bigr)^{1/2}\ \le\ k_1\Bigl(\frac s{a_d^2}\Bigr)^{1/2}x_k^{-1/2},
\]
as $\mathsf b^{-1}L^d=x_k^{-1}a_d^{-2}$. This cost does not vanish as $d$ grows: at $s=\theta_0a_d^2$ it is $\frac14(3/\Lambda_+)^{1/2}x_k^{-1/2}$, and the range $0\le s\le\theta_0a_d^2$ of (a) is made for it.

\emph{The finer layers.} Let $f=P_sf_P$. $\mathbf B_K-\mathbf B^{\le d}$ is the sum of the centred layers of the depths $d'>d$, of their conditional means and of the difference of the two remainders. $Y_P=\langle f,\mathbf B^{\le d}\rangle$ is $\mathcal F_d$-measurable and bounded, and $Z_P=\sum_{d'>d}\langle f,\mathbf h'_{d'}-\mathbf n_{d'}\rangle$ has $\mathbb E_K[Z_P|\mathcal F_d]=0$, so $\mathbb E_K|Y_P+Z_P|^2=\mathbb E_K|Y_P|^2+\mathbb E_K|Z_P|^2\ge\mathbb E_K|Y_P|^2$: in the second moment the finer layers can only add. The conditional means of the depths $>d$ pair with $f$ to at most $4a_d^{3/4}\mathsf c_{12}+4a_d\mathsf s_d$, by H\"older's inequality and (L4), and the remainders, by (L5), to at most $40a_d(\mathsf c_{(4)}\mathsf b^{-1}+a_K\mathsf a_{\rm top}^2)+4.08\,\bar y_d^2$ in $L^4(\mathbb P_K)$.

\emph{(a).} All these errors are at most $\mathsf z_dx_k^{-1/2}$ with $\mathsf z_d\to0$ uniformly in $K$, as $\bar y_d^2x_k^{1/2}\to0$, $a_dx_k^{1/2}=(a_d\mathsf b)^{1/2}\to0$, $a_dx_k^{1/2}\mathsf s_d\to0$ by (L4) and $a_d^{3/4}x_k^{1/2}\to0$; $d_L$ is the least depth $>d_0$ beyond which $\bar y_d\le\min\{\frac12,\frac1{16C_Q},\frac14c_Q\}$, $[3/(\Lambda_++\pi^2\mathfrak b(L^d\mathsf b)^{-2})]^{1/2}\varsigma_d-\mathsf z_d\ge\frac34(3/\Lambda_+)^{1/2}$ and $\mathsf z_d\le15c_{\mathsf J}\gamma_{\rm r}^{-1/2}$, a depth free of $K$. For $P$ without pivot bond, $d\ge d_L$ and $0\le s\le\theta_0a_d^2$,
\[
\|X_P\|_{L^2(\mathbb P_K)}\ \ge\ x_k^{-1/2}\Bigl\{\Bigl[\frac3{\Lambda_++\pi^2\mathfrak b(x_k)^{-2}}\Bigr]^{1/2}\varsigma_d-k_1\Bigl(\frac s{a_d^2}\Bigr)^{1/2}-\mathsf z_d\Bigr\}\ \ge\ x_k^{-1/2}\Bigl(\frac34-\frac14\Bigr)\Bigl(\frac3{\Lambda_+}\Bigr)^{1/2},
\]
as $k_1\theta_0^{1/2}=\frac14(3/\Lambda_+)^{1/2}$; since $(\frac34-\frac14)^2\frac3{\Lambda_+}=\frac3{4\Lambda_+}$, this is $\mathbb E_K|X_P|^2\ge\frac3{4\Lambda_+}x_k^{-1}$. For the fourth moment, the centred layers of all depths form a martingale whose components have, by (L3) at $f=P_sf_P$ with $\|f\|_1\le4a_d$ at the depths $d'\le d$ and $\|f\|_2^2\le4a_d^{-1}$ at the depths $d'>d$, the total variance proxy $v=\frac4{\gamma_{\rm r}}\sum_{d'}G_{d'}(f)^2\le\frac{4c_{\mathsf J}^2}{\gamma_{\rm r}}x_k^{-1}[\frac{16}{1-L^{-1}}+\frac4{L^2-1}]\le\frac{100\,c_{\mathsf J}^2}{\gamma_{\rm r}}x_k^{-1}$ ($L\ge3$); so each component $M^\alpha$ has $\mathbb E_K|M^\alpha|^4\le64v^2$, the martingale has the $L^4(\mathbb P_K)$ norm at most $60\sqrt2\,c_{\mathsf J}\gamma_{\rm r}^{-1/2}x_k^{-1/2}$, the other members at most $\mathsf z_dx_k^{-1/2}$, and $\mathbb E_K|X_P|^4\le k_4^4x_k^{-2}$, for every $P$.

\emph{Paley and Zygmund.} Let $\Sigma=a_d^{-1}\sum_P|X_P|^2$, the sum over the $N_P\ge3(1-4L^{-3})(\ell L^d)^3\ge\frac52(\ell/a_d)^3$ plaquettes without pivot bond ($L\ge3$). By \eqref{C-367-31} at $A=P_s\mathbf B_K$, $\|\mathrm d_aP_s\mathbf B_K\|_2^2\ge\Sigma$; $\mathbb E_K\Sigma\ge a_d^{-1}N_P\frac3{4\Lambda_+}x_k^{-1}$, and by the Cauchy--Schwarz inequality for the pairs of plaquettes $\mathbb E_K\Sigma^2\le a_d^{-2}N_P^2k_4^4x_k^{-2}$; so $\mathbb P_K(\Sigma\ge\frac12\mathbb E_K\Sigma)\ge(\mathbb E_K\Sigma)^2/4\mathbb E_K\Sigma^2\ge\frac14(3/4\Lambda_+k_4^2)^2=p_\ast$, and $\frac12\mathbb E_K\Sigma\ge\frac12\cdot\frac52\cdot\frac3{4\Lambda_+}\,\ell^3a_d^{-3}\,x_k^{-1}a_d^{-1}=k_\Sigma\mathsf b^{-1}\ell^3a_d^{-3}$, as $x_k^{-1}a_d^{-1}=\mathsf b^{-1}$.

This is (a). The conditional form: the bounds that hold at every bond or almost surely, and the martingale structure of the layers, are conditional on the block field at depth $d_0$ as they stand, so the argument goes through under $\mathbb E^0_K$ with the norms $(N_P^{-1}\sum_P\mathbb E^0_K|\cdot|^r)^{1/r}$, $r=2,4$, but for the two places where unconditional moments of one layer enter, the bound $m_2$ and (L5); there Markov's inequality bounds their conditional means off an event of probability at most $\tau_d$, and the inequality of Paley and Zygmund under $\mathbb E^0_K$, with the Cauchy--Schwarz inequality over $P$ for the second moment of $\Sigma$, gives the last bound.

(b) By Duhamel's formula the linearised gauge-fixed flow from an $L^2$ datum $z$ is a small perturbation of the heat flow, $\|w_s-P_sz\|_2\le2\epsilon(s)\|z\|_2$ with $\epsilon(s)=c_6\bar C^2s^{\frac14-\frac\kappa2}$; differentiating $A^a\circ\Phi_\sigma$ along $e^{ia_K(\mathbf B+rz)}$ at $z=P_\sigma\mathbf B$ and using $(\mathrm d_aP_\sigma\mathbf B,\mathrm d_aP_\sigma P_\sigma\mathbf B)=\|\mathrm d_aP_{3\sigma/2}\mathbf B\|_2^2$ and $\Gamma_K(F)^{1/2}\ge|\frac d{dr}F(e^{ia_K(\mathbf B+rz)})|/\|z\|_2$ gives the lower bound; the upper bound and the last inequality are the same expansion.
\end{proof}

The lower bound requires control of the frequency of the terminal
gradient. For the action this frequency is controlled by the change in
the squared current between two positive times. The following estimate
makes this control uniform in the lattice spacing.

\begin{lemma}[Lower propagation on a regular time interval]
\label{lem:short-time-lower-propagation}
Let $a=a_K$, $U_s=\Phi_sU$, $0<\kappa<\frac12$, and put
\[
 r=\frac34+\frac\kappa2,\qquad \theta=1-r=\frac14-\frac\kappa2.
\]
Suppose that $C_0>0$, $\sigma>0$, and
\begin{equation}\label{eq:propagation-regular-interval}
 a^{-2}\varpi(U_s)\le C_0s^{-r}\quad(0<s\le\sigma).
\end{equation}
Write $g(s)=(J,J)_a(U_s)$ and assume $g(\sigma)>0$. For
$x\ge0$ and $j>0$ define
\begin{align*}
 c_\kappa(x,j)
 &=e^{8x/\theta}\left(\log\max\{1,j\}
       +\frac{64x}{\theta}+1072x^2\right)+32x,\\
 E_\kappa(x,c)
 &=64x(1+\theta^{-1})
   +16e^{8x/\theta}x(c+32x)
   +\frac{672}{\theta}x^2
   +\frac{7168}{\theta}e^{8x/\theta}x^3.
\end{align*}
If $j\ge g(\sigma/2)/g(\sigma)$ and $x=C_0\sigma^\theta$, then
\begin{equation}\label{eq:short-time-action-propagation}
 \Gamma_K(A^a\circ\Phi_\sigma)(U)
 \ge g(\sigma)\exp\left[-2c_\kappa(x,j)
              -E_\kappa\bigl(x,c_\kappa(x,j)\bigr)\right].
\end{equation}
The factor on the right is positive, decreases as $j$ increases, and
depends only on $\kappa,C_0,\sigma,j$.
\end{lemma}

\begin{proof}
All the operators in the proof act on the finite periodic lattice at
spacing $a$. Use the scalar product
$(X,Y)_a=a^3\sum_b\langle X(b),Y(b)\rangle$ on positive bonds. Thus
$\mathrm{grad}_aF=a^{-2}\partial F$ satisfies
$\|\mathrm{grad}_aF\|_a^2=\Gamma_K(F)$, and
$\mathrm{grad}_aA^a=J$. At a field $V$, let
\[
 T_\mu h(x)=\mathrm{Ad}(V_\mu(x))h(x+ae_\mu),\qquad
 \nabla_\mu=a^{-1}(T_\mu-1),\qquad
 (DX)_{\mu\nu}=\nabla_\mu X_\nu-\nabla_\nu X_\mu.
\]
Here $D_0^*X=\sum_\mu\nabla_\mu^*X_\mu$. In this proof
$\Delta=\sum_\mu\nabla_\mu^*\nabla_\mu$ is positive and acts on each
component; it is distinct from the heat operator $\Delta_a$ of
Proposition~\ref{C-367-12}. Set
\[
 \phi(V)=a^{-2}\varpi(V),\qquad
 \mathcal F(p)=\frac{\partial V(p)-\partial V(p)^*}{2ia^2}.
\]
The orientations and base points of $\mathcal F$ agree with those of
$D$. Differentiating the plaquette action gives
$J=D^*\mathcal F$ and $a|J|_\infty\le4\phi$.

We first compare the action Hessian with $\Delta$. Define the symmetric
operator $H$ by
\[
 (X,HX)_a=\left.\frac{d^2}{dt^2}A^a(e^{iatX}V)\right|_{t=0}.
\]
For $p=p_{\mu\nu}(x)$ put
$\alpha_p=(-T_\nu X_\mu,-X_\nu,X_\mu,T_\mu X_\nu)$ and
$S_p=\sum_{i<j}[\alpha_{p,i},\alpha_{p,j}]$.
The four perturbed factors in the plaquette can be transported to $x$;
their sum is $a(DX)(p)$. For $P_t=\prod_i e^{it\alpha_{p,i}}$,
\[
 P_0''=-\left(\sum_i\alpha_{p,i}\right)^2-S_p.
\]
Every traceless Hermitian $2\times2$ matrix $Z$ has
$Z^2=|Z|^2\mathbf1$. Since
$\partial V(p)=(1-a_p)\mathbf1+ia^2\mathcal F(p)$, taking the trace
in the second variation gives
\begin{equation}\label{eq:propagation-hessian}
 (X,HX)_a
 =a^3\sum_p\left((1-a_p)|(DX)(p)|^2
                 +\langle\mathcal F(p),iS_p\rangle\right).
\end{equation}
Expansion of the squares also gives
\[
 \|DX\|_a^2+\|D_0^*X\|_a^2=(X,\Delta X)_a+w(X),\qquad
 w(X)=\sum_{\mu\ne\nu}(X_\mu,[\nabla_\mu,\nabla_\nu^*]X_\nu)_a.
\]
The commutator compares the two transports around a plaquette, so its
norm is at most $2\phi$. Thus $|w(X)|\le4\phi\|X\|_a^2$.
Moreover, $|\mathcal F(p)|\le\phi$,
$|S_p|\le3\sum_i|\alpha_{p,i}|^2$, and every bond belongs to four
plaquettes. Consequently the second term of
\eqref{eq:propagation-hessian} has modulus at most
$12\phi\|X\|_a^2$. Finally,
$a_p=\frac12|\partial V(p)-\mathbf1|^2$,
$|\partial V(p)-\mathbf1|\le2$, and
$|DX(p)|^2\le4a^{-2}\sum_{b\in p}|X(b)|^2$ imply
\[
 0\le a^3\sum_pa_p|DX(p)|^2\le16\phi\|X\|_a^2.
\]
It follows that
\begin{equation}\label{eq:propagation-hessian-error}
 H+D_0D_0^*=\Delta+\mathcal U,\qquad
 \|\mathcal U\|\le32\phi.
\end{equation}
In particular, $H=\Delta+\mathcal U$ on divergence-free forms.

We next record the two evolution equations used below. If $Y_s$ is the
differential of the flow, in the frame
$\delta U_s=iaY_sU_s$, then
\[
 Y_s'=-H_sY_s-\frac a2i[J_s,Y_s].
\]
Indeed, for $R_XF=\left.\frac d{dt}F(e^{iatX}V)\right|_0$ one has
$[R_X,R_Y]=-aR_{i[X,Y]}$. The symmetric part of the second derivative
is the Hessian, and hence
$R_YJ=HY-\frac a2i[J,Y]$. Differentiating
$\dot U=-iaJU$ now gives the displayed equation.
For a smooth gauge-invariant $G$ with $\Gamma_K(G)(U_\sigma)>0$, put
\[
 \xi(s)=\mathrm{grad}_a(G\circ\Phi_{\sigma-s})(U_s),\qquad
 n(s)=\frac{(\xi,H_s\xi)_a}{\|\xi\|_a^2}.
\]
The pairing $(\xi(s),Y_s)_a$ is constant. Since $D\Phi_s$ is
invertible, this gives
\begin{equation}\label{eq:propagation-adjoint}
 \xi'=H_s\xi-\frac a2i[J_s,\xi],\qquad
 \Gamma_K(G\circ\Phi_\sigma)(U)
 =\Gamma_K(G)(U_\sigma)\exp\left[-2\int_0^\sigma n(s)\,ds\right].
\end{equation}
Here $\xi$ never vanishes, and $D_0^*\xi=0$ because the flow is
gauge covariant. The current $j(s)=J(U_s)$ satisfies
\begin{equation}\label{eq:propagation-current-evolution}
 j'=-H_sj,\qquad D_0^*j=0,\qquad
 g'=-2(j,H_sj)_a.
\end{equation}
Thus $g$ vanishes at one time only if it vanishes identically.

The derivative of $\Delta$ has a useful cancellation. For fixed $X$,
$\partial_s\nabla_\mu X_\nu=-i[J_\mu,T_\mu X_\nu]$.
The part containing $a\nabla_\mu X_\nu$ pairs to zero with
$\nabla_\mu X_\nu$. Setting
$Y_\mu=\sum_\nu i[X_\nu,\nabla_\mu X_\nu]$ therefore gives
\[
 (X,\Delta'X)_a=-2(J,Y)_a=-2(\mathcal F,DY)_a.
\]
The discrete product rule
$\nabla_\mu[A,B]=[\nabla_\mu A,T_\mu B]+[A,\nabla_\mu B]$
shows explicitly that
\begin{align*}
 (DY)_{\mu\lambda}=\sum_\nu i\bigl(&
 [\nabla_\mu X_\nu,T_\mu\nabla_\lambda X_\nu]
 -[\nabla_\lambda X_\nu,T_\lambda\nabla_\mu X_\nu]\\
 &+[X_\nu,[\nabla_\mu,\nabla_\lambda]X_\nu]\bigr).
\end{align*}
The first two terms contain only first differences. Translations are
isometries, $|[A,B]|\le2|A||B|$, and the sum over the three pairs of
directions is bounded by twice the sum of the squared first differences.
The last term uses $\|[\nabla_\mu,\nabla_\lambda]\|\le2\phi$.
Multiplying by $2|\mathcal F|$ and summing gives
\begin{equation}\label{eq:propagation-laplacian-variation}
 |(X,\Delta'X)_a|
 \le8\phi(X,\Delta X)_a+24\phi^2\|X\|_a^2.
\end{equation}

For either $\eta=\xi$ or $\eta=j$, let
$h=(\eta,\Delta\eta)_a/\|\eta\|_a^2$ and
$n_\eta=(\eta,H\eta)_a/\|\eta\|_a^2$. Then
$h\ge0$, $|h-n_\eta|\le32\phi$, and
\[
 h'=\frac{2(\eta',(\Delta-h)\eta)_a
                  +(\eta,\Delta'\eta)_a}{\|\eta\|_a^2}.
\]
For $\xi$ the leading term is
$2\|(\Delta-h)\xi\|_a^2$, whereas for $j$ it is its negative.
The remaining part of $\xi'$ has norm at most
$(32+4)\phi\|\xi\|_a$, using $a|J|_\infty\le4\phi$.
The inequality
$72\phi\|\xi\|_a\|(\Delta-h)\xi\|_a
 \le2\|(\Delta-h)\xi\|_a^2+648\phi^2\|\xi\|_a^2$
and the corresponding bound with $64$ and $512$ for $j$ give
\begin{equation}\label{eq:propagation-frequency}
 h'\ge-8\phi h-672\phi^2\quad(\eta=\xi),\qquad
 h'\le8\phi h+536\phi^2\quad(\eta=j).
\end{equation}

Put $\psi(s)=C_0s^{-r}$, so
$\int_0^\sigma\psi=x/\theta$. For the adjoint solution, backward
integration of \eqref{eq:propagation-frequency} yields
\[
 h(s)\le e^{8x/\theta}
       \left(h(\sigma)+1344C_0^2s^{1-2r}\right),
 \qquad 0<s\le\sigma,
\]
because $2r-1\ge\frac12$. The square $\psi^2$ is not integrable at
zero, but $s\psi^2$ is. Accordingly, integrate against the weight $s$:
\[
 \int_0^\sigma h
 =\sigma h(\sigma)-\int_0^\sigma sh'
 \le\sigma h(\sigma)+\int_0^\sigma s(8\psi h+672\psi^2).
\]
There is no boundary term at zero: for each fixed lattice the flow,
the operators, and the nonzero adjoint solution are smooth there.
If $c\ge\sigma n(\sigma)$, then
$0\le\sigma h(\sigma)\le c+32x$. Combining the last two inequalities
with $|n-h|\le32\psi$ gives
\begin{equation}\label{eq:propagation-integrated-frequency}
 2\int_0^\sigma n(s)\,ds
 \le2\sigma n(\sigma)+E_\kappa(x,c).
\end{equation}
Indeed, the endpoint and $n-h$ terms give
$64x(1+\theta^{-1})$; the constant part of $h(s)$ gives at most
$16e^{8x/\theta}x(c+32x)$, since $2-r\ge1$; and the two remaining
integrals are
\[
 1344C_0^2\int_0^\sigma s^{1-2r}\,ds
       =\frac{672}{\theta}x^2,\qquad
 21504C_0^3e^{8x/\theta}\int_0^\sigma s^{2-3r}\,ds
       =\frac{7168}{\theta}e^{8x/\theta}x^3.
\]

It remains to control the terminal frequency when $G=A^a$.
Apply the second inequality of \eqref{eq:propagation-frequency} to
the current and integrate forward from $s$ to $\sigma$. Averaging
over $\sigma/2\le s\le\sigma$ gives
\[
 \frac\sigma2 h(\sigma)
 \le e^{8x/\theta}\left(
     \int_{\sigma/2}^\sigma h(s)\,ds
     +\frac\sigma2\,536\int_{\sigma/2}^\sigma\psi(s)^2\,ds\right).
\]
By \eqref{eq:propagation-current-evolution},
\[
 \int_{\sigma/2}^\sigma h(s)\,ds
 \le\frac12\log\frac{g(\sigma/2)}{g(\sigma)}
       +32\int_{\sigma/2}^\sigma\psi(s)\,ds,
 \qquad
 \int_{\sigma/2}^\sigma\psi(s)^2\,ds
       \le2C_0^2\sigma^{1-2r}.
\]
Consequently
\[
 0\le\sigma h(\sigma)
 \le e^{8x/\theta}\left(
       \log\frac{g(\sigma/2)}{g(\sigma)}
       +\frac{64x}{\theta}+1072x^2\right).
\]
The terminal adjoint vector for $G=A^a$ is the current itself.
Thus its terminal frequency satisfies
$\sigma n(\sigma)\le c_\kappa(x,j)$.
Equations~\eqref{eq:propagation-adjoint} and
\eqref{eq:propagation-integrated-frequency} prove
\eqref{eq:short-time-action-propagation}.
\end{proof}


\begin{proposition}[Lower energy bounds for flowed observables]\label{C-367-23}
For all $K$, $m$, $\beta>0$, with $a=a_K$, the lattice time $\tau=s/a^2$, $T=\sigma/a^2$ and $U^\circ_\tau=\Phi_{a^2\tau}U$: let $|X|_\infty=\max_b|X(b)|$, let $\nabla^3A_K(V)$ be the symmetric trilinear form with $\nabla^3A_K(V)[X,X,X]=\frac{d^3}{dt^3}A_K(e^{itX}V)|_{t=0}$, and let $Z_p$ be as in Lemma~\ref{C-367-13}, taken at the field $V$.
\begin{enumerate}
\item[(a)] (Log-convexity with a defect.) Let $G\in C^2(\Omega_K)$ be invariant, $U\in\Omega_K$, $\sigma>0$ and $\partial G(U_\sigma)\ne0$; for $0\le\tau\le T$ let $\gamma_\tau=\partial(G\circ\Phi_{a^2(T-\tau)})(U^\circ_\tau)$, $N(\tau)=\|\gamma_\tau\|^2$, $\mathsf q(\tau)=\nabla^2A_K(U^\circ_\tau)[\gamma_\tau,\gamma_\tau]$ and $\mathsf f_\tau=\partial A_K(U^\circ_\tau)$. Then $N>0$, $\dot N=2\mathsf q$, $(\log N)''\ge-\eta$ with $\eta(\tau)=2|\nabla^3A_K(U^\circ_\tau)[\mathsf f_\tau,\gamma_\tau,\gamma_\tau]|/N(\tau)$, and
\[
\Gamma_K(G\circ\Phi_\sigma)(U)\ \ge\ \Gamma_K(G)(U_\sigma)\,\exp\Bigl[-2\nu-\int_0^T\tau\,\eta(\tau)\,d\tau\Bigr],\qquad\nu=\frac{T\,\mathsf q(T)}{N(T)},
\]
$\nu$ the frequency of the gradient at the flowed field. For all $V$ and real $\beta,\gamma$: 
\[
\begin{multlined}
|\nabla^3A_K(V)[\beta,\gamma,\gamma]|\le|\gamma|_\infty\{4(\sum_p|Z_p(\beta)|^2)^{1/2}\|\gamma\|+8(\sum_p|Z_p(\gamma)|^2)^{1/2}\|\beta\| \\
+64\,\varpi(V)\|\gamma\|\|\beta\|\}.
\end{multlined}
\]

\item[(b)] (An effective form.) In (a) let besides $0<\kappa<\frac16$, $C_F\ge1$, $0<\sigma\le1$ and $a^{-2}\varpi(U_s)\le C_Fs^{-\frac34-\frac\kappa2}$ for $0<s\le\sigma$. Put $\lambda_F=6C_F$, $\mathsf e_1=(1+56\lambda_Fe^{56\lambda_F})^{1/2}$, 
\[
\mathsf r=(\ell/a)^{3/2}|\partial G(U_\sigma)|_\infty/\|\partial G(U_\sigma)\|,
\]
$\nu_+=\max\{\nu,0\}$, $\varrho_F=3.5\,e^{80\lambda_F}C_F\mathsf r$ and 
\[
\mathsf k_A=28e^{56\lambda_F}C_F+e^{28\lambda_F}\varrho_F(7\mathsf e_1+256C_F).
\]
If
\[
4e^{2\nu_++1}\varrho_F\bigl[48\mathsf e_1+384C_F+12\nu_+^{1/2}+32\mathsf k_A^{1/2}+768\varrho_F\bigr]\sigma^{\frac14-\frac\kappa2}\le1,
\]
then $\Gamma_K(G\circ\Phi_\sigma)(U)\ge e^{-2\nu_+-\frac12}\,\Gamma_K(G)(U_\sigma)$. At $G=A^a$: $\Gamma_K(A^a)(U_\sigma)=(J,J)_a(U_\sigma)$ and $\nu=-\frac\sigma2\,\partial_\sigma\log(J,J)_a(U_\sigma)$.
\item[(c)] (Short-time propagation and its extension.) Let $0<\kappa<\frac16$, $\mathbf B$ a real $1$-form, $R\ge\max\{1,\|\mathbf B\|_{\kappa;a}\}$ and
\[
a\le a_2(R):=\min\bigl\{a_1(R),\ (16e^{12}R)^{-2/(1-2\kappa)},\ \bigl(8(16\mathsf M_\Theta)^{1/6}c_3R\bigr)^{-1/(1-2\kappa)}\bigr\},
\]
a positive number that depends on $R$ and $\kappa$ alone, with $c_3$, $a_0(R)$, $s_0(R)$ of the short-time theory of the gauge-fixed flow in Proposition~\ref{C-367-12}(a), $a_1(R)\le a_0(R)$ the threshold of Lemma~\ref{C-367-16}(b), and $\mathsf M_\Theta=\max\{1,\sup|\Theta^{\rm ch}|\}$, the absolute number of the head of Proposition~\ref{C-367-12}, with the chart equation $\frac d{ds}A_s(b)=a^{-3}\Theta^{\rm ch}(aA_s|_{\mathsf N(b)})$ of the gauge-fixed flow set there and the supremum taken where $|aA|\le1$; so $a_2(R)\le a_1(R)$, and $a\le a_2(R)$ serves both Lemma~\ref{C-367-16}(b) and (c); let $U$ be gauge equivalent to $e^{ia\mathbf B}$, $\sigma>0$, $\sigma'=\min\{\sigma,s_0(R)\}$, $(J,J)_a(U_\sigma)>0$, $\Lambda\ge\Lambda_\sigma(U)$ and $\mathsf j_\sharp\ge(J,J)_a(U_{\sigma'/2})/(J,J)_a(U_\sigma)$.

For $0<\sigma\le s_0(R)$, put
\[
 C_0(R)=4\bigl((2+c_3)R\bigr)^2,\qquad
 x=C_0(R)\sigma^{1/4-\kappa/2},
\]
\[
 \mathsf G^{\rm short}_\kappa(R,\sigma,\mathsf j_\sharp)
 =\exp\left[-2c_\kappa(x,\mathsf j_\sharp)
 -E_\kappa\bigl(x,c_\kappa(x,\mathsf j_\sharp)\bigr)\right],
\]
where $c_\kappa$ and $E_\kappa$ are defined in
Lemma~\ref{lem:short-time-lower-propagation}. Then
\begin{equation}\label{eq:rough-data-short-propagation}
 \Gamma_K(A^a\circ\Phi_\sigma)(U)
 \ge(J,J)_a(U_\sigma)
       \mathsf G^{\rm short}_\kappa(R,\sigma,\mathsf j_\sharp).
\end{equation}
This positive factor is independent of $K$, $\ell$, and $\mathsf b$,
and decreases as $\mathsf j_\sharp$ increases.

By \emph{(Lower propagation)} we mean the additional assertion that
the lower estimate extends to every $\sigma>0$ under the hypotheses
of (c), in the form
$\Gamma_K(A^a\circ\Phi_\sigma)(U)\ge(J,J)_a(U_\sigma)
\mathsf G_\kappa(R,\sigma',\Lambda,\mathsf j_\sharp)$,
with a positive function of the four displayed arguments alone,
decreasing in $\Lambda$ and $\mathsf j_\sharp$, and independent of
$K$, $\ell$, and $\mathsf b$. On $0<\sigma\le s_0(R)$ the proved
factor $\mathsf G^{\rm short}_\kappa$ supplies this estimate without
any additional input. For arbitrary $\sigma$, the upper bound
$\Gamma_K(A^a\circ\Phi_\sigma)(U)\le
e^{24\Lambda_\sigma(U)}(J,J)_a(U_\sigma)$ also holds without the
additional assertion.
\item[(d)] (Positive energies and the flowed current.) Assume (E) and $0<\kappa<\frac16$; let $\varphi\in C^\infty(\mathbb R)$ be bounded with all its derivatives and $\varphi'>0$ on $[0,\infty)$, and $0<\sigma<\sigma_2\le1$. Say that (J$_\sigma$) holds if there is $c_0>0$ with $p_0:=\liminf_K\mu^K((J,J)_a(U_\sigma)\ge c_0)>0$, and (D$_{\sigma,\sigma_2}$) if $\liminf_K\langle A^a\circ\Phi_\sigma-A^a\circ\Phi_{\sigma_2}\rangle_{\mu^K}>0$. In (i) and in the first two implications of (iii) use the consequence (Adm) of (E) proved in Proposition~\ref{C-367-12}. (i) Assume also \textup{(Lower propagation)}. If (J$_\sigma$) holds, then $\liminf_K\int\chi_K\Gamma_K(\varphi(A^a)\circ\Phi_\sigma)\,d\mu^K>0$ for every sequence of Borel $\chi_K:\Omega_K\to[0,1]$ with $\limsup_K\mu^K(\chi_K<1)<p_0$.

(ii) Conversely, if $R\ge R_\flat$, the Borel $\chi_K:\Omega_K\to[0,1]$ vanish off $\mathsf O_K(8R,\bar s(R))$ and $\liminf_K\int\chi_K\Gamma_K(\varphi(A^a)\circ\Phi_\sigma)\,d\mu^K>0$, then (J$_\sigma$) holds. (iii) (J$_{\sigma_2}$) implies (D$_{\sigma,\sigma_2}$), and (D$_{\sigma,\sigma_2}$) implies (J$_\sigma$); and $\liminf_K\langle A^a\circ\Phi_\sigma\rangle_{\mu^K}>\sup_K\langle A^a\circ\Phi_1\rangle_{\mu^K}$ implies (D$_{\sigma,1}$).
\end{enumerate}
\end{proposition}

\begin{proof}[Proof sketch]
(a) $N>0$ because $D\Phi$ is invertible. As $G$ is invariant, $\gamma_\tau$ is horizontal, and along the trajectory $\dot\gamma_\tau=\nabla^2A_K(U^\circ_\tau)\gamma_\tau-\frac i2[\mathsf f_\tau,\gamma_\tau]$, the Hessian taken as a symmetric operator; since $\langle\gamma,i[\mathsf f,\gamma]\rangle=0$, $\dot N=2\mathsf q$, and a second differentiation, in which the Hessians cancel, gives $\dot{\mathsf q}=2\|\nabla^2A_K\gamma\|^2-\nabla^3A_K[\mathsf f,\gamma,\gamma]+\frac16\|i[\mathsf f,\gamma]\|^2$. So
\[
(\log N)''=\frac4{N^2}\bigl(\|\nabla^2A_K\gamma\|^2\|\gamma\|^2-\langle\gamma,\nabla^2A_K\gamma\rangle^2\bigr)+\frac2N\Bigl(\tfrac16\|i[\mathsf f,\gamma]\|^2-\nabla^3A_K[\mathsf f,\gamma,\gamma]\Bigr)\ge-\eta
\]
by the Cauchy--Schwarz inequality. Integrating twice from $T$, $\log\frac{N(\tau_1)}{N(T)}\ge-2(T-\tau_1)\frac{\mathsf q(T)}{N(T)}-\int_{\tau_1}^T(\tau-\tau_1)\eta(\tau)\,d\tau$ for $0\le\tau_1\le T$; at $\tau_1=0$, $N(0)=a\Gamma_K(G\circ\Phi_\sigma)(U)$ and $N(T)=a\Gamma_K(G)(U_\sigma)$. The weight $\tau$ under the integral keeps the defect finite at rough fields. The bound of $\nabla^3A_K$ follows from the expansion of the third derivative plaquette by plaquette, as in Lemma~\ref{C-367-13}(a), and the Cauchy--Schwarz inequality over the plaquettes. (b) The hypothesis gives $\Lambda_\sigma\le\lambda_F$; by Lemma~\ref{C-367-13}(b), (c), (h), $(J,J)_a(U_s)^{1/2}\le\mathsf j_0s^{-\frac54-\frac\kappa2}$ with $\mathsf j_0=3.5\,e^{28\lambda_F}\ell^{3/2}C_F$, and $|\gamma_\tau|_\infty\le e^{52\lambda_F}|\gamma_T|_\infty$, $\|\gamma_\tau\|\le e^{28\lambda_F}\|\gamma_T\|$.

Lemma~\ref{C-367-13}(a) gives $\sum_p|Z_p(X)|^2\le\nabla^2A_K(V)[X,X]+28\varpi(V)\|X\|^2$ for all $X$; with $\frac d{d\tau}\|\mathsf f_\tau\|^2=-2\nabla^2A_K[\mathsf f_\tau,\mathsf f_\tau]$ this makes the energy $\sum_p|Z_p(\mathsf f_\tau)|^2$ of the current integrable in $\tau$, $\int_0^T\sum_p|Z_p(\mathsf f_\tau)|^2d\tau\le\frac12\|\mathsf f_0\|^2+28\lambda_F\sup_\tau\|\mathsf f_\tau\|^2$, and with (a) it bounds the energy $\sum_p|Z_p(\gamma_\tau)|^2\le\mathsf q+28\varpi N$ of $\gamma_\tau$. On the largest interval $[\tau_1,T]$ on which $N\ge e^{-2\nu_+-1}N(T)$, the bound of $\nabla^3A_K$ and these estimates give $\int_{\tau_1}^T(\tau-\tau_1)\eta\le\frac12$ under the displayed condition; so (a) gives $\log\frac{N(\tau_1)}{N(T)}\ge-2\nu_+-\frac12$, and $\tau_1=0$ by continuity.

At $G=A^a$, $\partial A^a=a^{-1}\mathsf f$, and $\frac d{d\tau}\|\mathsf f_\tau\|^2=-2\langle\mathsf f_\tau,\nabla^2A_K\mathsf f_\tau\rangle$. (c) Let $0<\sigma\le s_0(R)$. Since $a\le a_2(R)\le a_0(R)$,
Proposition~\ref{C-367-12}(a) gives a gauge-fixed representative
$e^{iaA_s}$ of $U_s$ for $0\le s\le s_0(R)$, with
\[
 \|A_s\|_\infty\le(1+c_3)Rs^{-1/4-\kappa/2},\qquad
 \|\nabla A_s\|_\infty\le(2+c_3)Rs^{-3/4-\kappa/2}.
\]
For a plaquette, split its holonomy into the two differences between
opposite bond potentials and the group commutator of two adjacent
bond exponentials. The exponential is Lipschitz on Hermitian matrices,
and the group commutator has norm at most twice the product of their
norms. This gives
\[
 a^{-2}\varpi(e^{iaA_s})
 \le2\left(\|\nabla A_s\|_\infty+\|A_s\|_\infty^2\right),
\]
as in Lemma~\ref{C-367-13}(e). Put $\bar C=(2+c_3)R\ge1$.
Since $s_0(R)\le1$ and $1/2+\kappa\le3/4+\kappa/2$,
the right side is at most $4\bar C^2s^{-3/4-\kappa/2}$.
Curvature is gauge invariant, so
\eqref{eq:propagation-regular-interval} holds with $C_0=C_0(R)$.
Here $\sigma'=\sigma$, and the ratio assumed in (c) is exactly the
ratio required in Lemma~\ref{lem:short-time-lower-propagation}.
That lemma proves \eqref{eq:rough-data-short-propagation}.
The assertion for $\sigma>s_0(R)$ remains the stated additional
input. The upper bound is Lemma~\ref{C-367-13}(b) at $G=A^a$,
using $\Gamma_K(A^a)=(J,J)_a$.

(d) (i): under \textup{(Lower propagation)}, fix $R$, $R'$ and $s_+$ such that, for large $K$, the regularised fields $\mathsf O_K(8R,\bar s(R))$ of (C4), the event that the gauge orbit holds a $1$-form of rough norm $\le R'$ at the index $\frac\kappa2$, and $\{A^a\circ\Phi_\sigma\le s_+\}$ all have probability close to $1$ (Proposition~\ref{C-367-12}(a), by (Adm), and Chebyshev's inequality).

On their intersection with $\{\chi_K=1\}\cap\{(J,J)_a(U_\sigma)\ge c_0\}$, whose mass has lower limit $>0$, the proof of (C2)(a) bounds $\Lambda_\sigma$ and $(J,J)_a(U_{\sigma'/2})$, \textup{(Lower propagation)} at the index $\frac\kappa2$, for $a_K\le a_2(R')$, bounds $\Gamma_K(A^a\circ\Phi_\sigma)$ below by $c_0$ times a positive number free of $K$, and $\Gamma_K(\varphi(A^a)\circ\Phi_\sigma)=\varphi'(A^a\circ\Phi_\sigma)^2\Gamma_K(A^a\circ\Phi_\sigma)\ge(\min_{[0,s_+]}\varphi')^2\Gamma_K(A^a\circ\Phi_\sigma)$. (ii): on $\mathsf O_K(8R,\bar s(R))$, for large $K$, the proof of (C2)(a) gives $\Lambda_\sigma\le\lambda_\sharp$ and $(J,J)_a(U_\sigma)\le\mathsf j_\sigma$ with numbers free of $K$, and $\Gamma_K(\varphi(A^a)\circ\Phi_\sigma)\le\|\varphi'\|_\infty^2e^{24\lambda_\sharp}(J,J)_a(U_\sigma)$ (Lemma~\ref{C-367-13}(b)); splitting at a level $c$, $\int\chi_K\Gamma_K\le\|\varphi'\|_\infty^2e^{24\lambda_\sharp}[c+\mathsf j_\sigma\mu^K((J,J)_a(U_\sigma)\ge c)]$.

(iii): $A^a\circ\Phi_\sigma-A^a\circ\Phi_{\sigma_2}=\int_\sigma^{\sigma_2}(J,J)_a(U_s)\,ds$ (Lemma~\ref{C-367-13}(c)), and $(J,J)_a(U_{s'})\le e^{24\lambda_\sharp}(J,J)_a(U_s)$ for $s\le s'$ on $\mathsf O_K(8R,\bar s(R))$, whose complement has small mass for large $R$ by (Adm) (Step~8); in the second arrow the moments of $(J,J)_a(U_s)$ of Proposition~\ref{C-367-12}(a), under (Adm), and the Cauchy--Schwarz inequality give $\langle(J,J)_a(U_\sigma)\rangle_{\mu^K}\ge2c_0$, hence $\mu^K((J,J)_a(U_\sigma)\ge c_0)\ge c_0^2/\sup_K\langle(J,J)_a(U_\sigma)^2\rangle_{\mu^K}$. The last arrow is $\liminf_K(x_K-y_K)\ge\liminf_Kx_K-\limsup_Ky_K$.
\end{proof}

\subsection{Completion of the limit and flow arguments}

\emph{Step 1: compactness, motion and the Markov property, (A1)--(A3).} Let $K>n\ge n''$ and let $H\in C^1(\Omega_K)$ be real and invariant, with tower $(h_k)$. Under $\sigma^K$, $h_k(V_k)$ is a martingale from the top to the bare level (Lemma~\ref{C-367-7}(ii)), so $\int|H-h_t(V_t)|^2d\sigma^K=\sum_{k<t}\int\mathrm{Var}_{p_W}(h_k\circ V_W)\,\sigma^K_{d-1}(dW)$; Lemma~\ref{C-367-7}(iv) at $g=h_k$, \eqref{C-367-10}, $\sigma^K\le\mu^K$ on invariant functions, $\sum_{d>n}a_d^2=\frac{a_n^2}{L^2-1}$, and Lemma~\ref{C-367-7}(i) for the passage to $\mu^K$ give
\[
\|H-(\mathbf 1_{\mathsf I_n}h_t)\circ\pi_{K,n}\|^2_{L^2(\mu^K)}\le\frac{C^\sigma C_{\rm P}}{L^2-1}\frac{a_n^2}{\mathsf b}\,\mathcal E_K(H)+4\mathsf R_n\sup|H|^2,\qquad\int\Gamma_n(h_t)\,d\sigma^K_n\le C^\sigma\mathcal E_K(H).
\]
(A2): after replacing $H_j$ by smooth invariant functions with the same bounds (closure, a smooth contraction, the gauge average), let $c=\sup_j\sup|H_j|$, $e=\sup_j\mathcal E_{K_j}(H_j)$, and for $n\ge n''$ and $K_j>n$ let $h^j_n$ be the top of the tower of $H_j$ at $t=K_j-n$ and $g^j_n=\mathbf 1_{\mathsf I_n}h^j_n$. The tops are bounded by $c$ and, by the display and Lemma~\ref{C-367-7}(iii), have Dirichlet integrals on $\mathsf I_n$ for Haar measure bounded uniformly in $j$. Rellich's theorem on compact subsets of the open set $\mathsf I_n$, which needs no regularity of its boundary, the bound $c$, and diagonal sequences over an exhaustion of $\mathsf I_n$ by compact sets and over $n$, give one subsequence along which $h^j_n\to\phi_n$ in $L^2(\mathsf I_n,dV)$ for every $n\ge n''$.

The limit $\phi_n$ is a Haar class; as a function under $\mu_n$ we take a representative: let $\Psi_n$ be $\mathbf 1_{\mathsf I_n}$ times the gauge average of a Borel function in the class of $\phi_n$ with values in $[-c,c]$; it is Borel, invariant at every point, $|\Psi_n|\le c$, and $\Psi_n=\phi_n$ Haar-almost everywhere on $\mathsf I_n$, as $\phi_n$ is invariant almost everywhere. Proposition~\ref{C-367-17}(c) for $\mu^{K_j}_n$ at $f=|\Psi_n-g^j_n|^2\le4c^2$, which vanishes off $\mathsf I_n$, gives $\int f\,d\mu^{K_j}_n\le e^{\mathsf o_n}|\mathsf I_n|^{-1}\int_{\mathsf I_n}|\phi_n-h^j_n|^2dV+4c^2\mathsf R_n$, so with the display
\[
\limsup_j\|H_j-\Psi_n\circ\pi_{K_j,n}\|_{L^2(\mu^{K_j})}\ \le\ 2c\,\mathsf R_n^{1/2}+\Bigl(\frac{C^\sigma C_{\rm P}}{L^2-1}\frac{a_n^2}{\mathsf b}\,e+4\mathsf R_nc^2\Bigr)^{1/2}=:\epsilon_n,\qquad\epsilon_n\to0 .
\]
Total variation enters only here: $\|\Psi_n(Y_n)-\Psi_{n'}(Y_{n'})\|_{L^2(\mu)}=\lim_j\|\Psi_n\circ\pi_{K_j,n}-\Psi_{n'}\circ\pi_{K_j,n'}\|_{L^2(\mu^{K_j})}\le\epsilon_n+\epsilon_{n'}$, so the invariant cylinder functions $\Psi_n(Y_n)$ converge in $L^2(\mu)$ to some $G$, $|G|\le c$, and $\lim_n\limsup_j\|H_j-\Psi_n\circ\pi_{K_j,n}\|_{L^2(\mu^{K_j})}=0$: $H_j\to G$ strongly. This holds along every sequence of cutoffs. (A1): let $H_j\to\Phi(Y_n)$ be a recovery sequence of Lemma~\ref{C-367-6}(b), smooth, invariant, $\sup|H_j|\le\sup|\Phi|$ (truncation and the gauge average do not raise the energy). The conditional mean under $\sigma^K$ given $V_t$ is an $L^2(\sigma^K)$-contraction that fixes $\Phi(V_t)$, so the tops converge to $\Phi$ in $L^2(\sigma^K_n)$, hence in $L^2(\mathsf I_n,dV)$ by the lower bound of the density; the bound on $\int\Gamma_n(h_t)d\sigma^K_n$, $\inf\mathsf g\ge(1-\mathsf R_n)e^{-\mathsf o_n}/|\mathsf I_n|\ge e^{-\mathsf o_n}/2|\mathsf I_n|$ and the lower semicontinuity of the Dirichlet integral give \eqref{C-367-3}.

As $\Gamma_n(a_p)=4a_n^{-1}|\mathsf s_p|^2$, by the computation of Lemma~\ref{C-367-13}(a) on $\Omega_n$, is positive at the fields of $\mathsf I_n$ with $\varpi_p>0$, $\mathcal E_\star(a_p(Y_n))>0$; since $\nu_\Psi(\{\infty\})=0$ for invariant $\Psi$, $(\Psi,(1-Q_t)\Psi)=\int(1-e^{-t\lambda})d\nu_\Psi>0$. (A3): for an invariant cylinder function $\Psi$ and $t>0$, $H_j=P^{K_j}_t\Psi_{K_j}$ is invariant, $|H_j|\le\sup|\Psi|$, $\mathcal E_{K_j}(H_j)\le\sup|\Psi|^2/2et$ (Lemma~\ref{C-367-6}(a)), and $H_j\to Q_t\Psi$ weakly; by (A2) every subsequence has a strongly convergent subsequence, whose limit is the weak limit; so $H_j\to Q_t\Psi$ strongly, which is one of the equivalent properties of Lemma~\ref{C-367-6}(b), and with it all of them, (M1), and the convergence of $P^{K_j}_tH_j$ for all $H_j\to F$.

With (M2) of Lemma~\ref{C-367-6}(b) the forms converge in the strong and weak sense, which gives the strong convergence of the resolvents by the theorem on the convergence of resolvents under convergence of forms in the strong and weak sense; Lemma~\ref{C-367-6}(c) gives a dense domain and strong continuity. Compactness: fix a countable dense set $(\Phi_i)$ of invariant cylinder functions; for $G_m$ in the set choose a cylinder function $\Psi_m$ with $\|\Psi_m-G_m\|\le\frac1m$ and, by (M2), truncation at $\pm1$ and the gauge average, smooth invariant $H_m$ on $\Omega_{K_{j(m)}}$ with $|H_m|\le1$, $\mathcal E_{K_{j(m)}}(H_m)\le1+\frac1m$ and $\|H_m-\Psi_{m,K_{j(m)}}\|\le\frac1m$, where $j(m)$ increases so fast that $\|\Psi_{m,K}-\Phi_{i,K}\|$ is within $\frac1m$ of $\|\Psi_m-\Phi_i\|$ for $i\le m$ and $K\ge K_{j(m)}$; by (A2) a subsequence of $(H_m)$ converges strongly to some $G$, then $G_m\to G$ in $L^2(\mu)$ along it, and $|G|\le1$, $\mathcal E_\star(G)\le1$ by (M1).

The criterion without a cutoff: $W_j=(1-\mathcal L_{K_j})H_j$ is bounded in $L^2$, so a subsequence converges weakly to some $W$ (a diagonal sequence over a countable dense set of cylinder functions); for invariant cylinder $\Psi$, $\int H_j\Psi_{K_j}d\mu^{K_j}=\int W_j(1-\mathcal L_{K_j})^{-1}\Psi_{K_j}d\mu^{K_j}\to(W,(1+A_\star)^{-1}\Psi)$, the resolvents being symmetric and strongly convergent, so $F=(1+A_\star)^{-1}W\in D(A_\star)$; the weak limit $F-W=-A_\star F$ of $\mathcal L_{K_j}H_j$ does not depend on the subsequence, and $\mathcal E_{K_j}(H_j)=\langle H_j,W_j\rangle-\|H_j\|^2\to(F,W)-\|F\|^2=\mathcal E_\star(F)$. The bilinear forms: with $T^\alpha_j=\alpha(\alpha-\mathcal L_{K_j})^{-1}$ and $T^\alpha=\alpha(\alpha+A_\star)^{-1}$, $\mathcal E_{K_j}(T^\alpha_ju_j,v_j)=\alpha\langle u_j-T^\alpha_ju_j,v_j\rangle\to\alpha(u-T^\alpha u,v)=\mathcal E_\star(T^\alpha u,v)$, as $u_j-T^\alpha_ju_j\to u-T^\alpha u$ strongly; by the spectral theorem $\mathcal E_{K_j}(u_j-T^\alpha_ju_j)\le\mathcal E_{K_j}(u_j)-\alpha\langle u_j,u_j-T^\alpha_ju_j\rangle$, whose limit $\mathcal E_\star(u)-\int\frac{\alpha\lambda}{\alpha+\lambda}d\nu_u$ tends to $0$ as $\alpha\to\infty$; with $|\mathcal E_{K_j}(u_j-T^\alpha_ju_j,v_j)|\le\mathcal E_{K_j}(u_j-T^\alpha_ju_j)^{1/2}\mathcal E_{K_j}(v_j)^{1/2}$, $\mathcal E_\star(v)\le\liminf_j\mathcal E_{K_j}(v_j)$ by (M1) and $\mathcal E_\star(T^\alpha u,v)\to\mathcal E_\star(u,v)$ this gives the convergence.

\emph{Step 2: the connected coarse-scale law, (B1).} The regularity number of a background $U$ of level $k$ is $L^{2k}\max_p|\partial U(p)-\mathbf 1|$, the maximum over the bare plaquettes (Definition~\ref{def:sf-chart-data}), and $\mathsf{Cut}_k$ is the set of the block fields over which it is $<c_\rho\varepsilon(x_k)$ (the definitions preceding Lemma~\YMMref{113.10}{796}). By the curvature comparison in the proof of Lemma~\YMMref{113.10}{796}, these numbers contract from level to level, by the factor $L^{-2}$ up to $C_H\delta(x_k)$: for $W\in\mathsf{Cut}_{k+1}$ and $B$ in the box, the regularity number of every background over $V_W(B)$ is at most $L^{-2}$ times that of a background over $W$ plus $C_H\delta(x_k)$; starting the chain $j_1$ levels higher, at every $U$ in the support of its small-field weight,
\[
\varpi(V_t)\le2\vartheta\varepsilon_n,\qquad\vartheta=\vartheta_\infty+L^{-3j_1/2}c_\rho<1 .
\]
The set $\mathcal K_0=\{\varpi\le2\vartheta\varepsilon_n\}$ is compact and lies in $\mathsf I_n$; as $\mathsf D_n$ and $\mathsf I_n\setminus\mathsf D_n$ are open, it is the disjoint union of the compact invariant sets $\mathcal K_1=\mathcal K_0\cap\mathsf D_n$ and $\mathcal K_2=\mathcal K_0\setminus\mathsf D_n$, which are free of $K$. The convexity of one box joins $V_W(B)$ to $V_W(0)$ at a fixed $W$ only; the $j_1$ levels are passed with the tower of a continuous function. Let $\mu_{t+j_1}=c_\rho\varepsilon(x_{t+j_1})$ and $\mu_k=L^{-2}\mu_{k+1}+C_H\delta(x_k)$ for $t\le k<t+j_1$, so that $\mu_t\le\vartheta\varepsilon_n$, and let $\mathsf D'_k$ be the set of the $W\in\mathsf{Cut}_k$ over which the regularity number is $<\mu_k$: $\mathsf D'_{t+j_1}=\mathsf{Cut}_{t+j_1}$, $V_W(B)\in\mathsf D'_k$ for $W\in\mathsf D'_{k+1}$ and every $B$ in the box (Lemma~\YMMref{113.10}{796}, the curvature comparison in the proof of Lemma~\YMMref{113.10}{796}), and $\mathsf D'_t\subset\mathcal K_0$ (the small-curvature averaging bound of Proposition~\YMMref{14.1}{61}).

Let $h$ be continuous on $\Omega_n$ with values in $[0,1]$, $h=1$ on $\mathcal K_1$ and $h=0$ on $\mathcal K_2$ (Urysohn), and let $h_t=h$, $h_{k+1}(W)=\int h_k(V_W(B))\,p_W(dB)$, be its tower upward from the level $t$ in the chain from the depth $n-j_1$. The kernels carry bounded continuous functions on $\mathsf{Cut}_k$ to continuous functions on $\mathsf{Cut}_{k+1}$ (as in the proof of Lemma~\ref{C-367-7}(ii), where the integrands $F_k$ are shown continuous on the cuts): the kernel at $W$ is the normalised chart integral \eqref{eq:sf-chart-integral} of the integrand $F_k$ of the chain at the level $k$; on $\mathsf{Cut}_{k+1}$, which lies in the set $\mathsf O^{\rm im}$ of the chart (Lemma~\YMMref{113.10}{796}), the chart, its gauge-fixing form and the exponent of its Jacobian are continuous in $(W,B)$, and the chart integral of a continuous function is continuous (Lemma~\ref{lem:sf-chart-integral}); $F_k$ is continuous on $\mathsf{Cut}_k$, into which the chart carries the box (Lemma~\YMMref{113.10}{796}), and dominated convergence on the box applies.

By induction on $k$, $h_k$ is continuous on $\mathsf{Cut}_k$ and takes on $\mathsf D'_k$ the values $0$ and $1$ alone: for $W\in\mathsf D'_{k+1}$ the function $B\mapsto h_k(V_W(B))$ is continuous on the convex box with values in $\{0,1\}$, hence constant, and as $p_W$ is a probability, $h_{k+1}(W)=h_k(V_W(B))$ for every $B$ in the box. Take all the fluctuations $0$: for $W\in\mathsf{Cut}_{t+j_1}$ let $W^{(t+j_1)}=W$ and $W^{(k)}=V_{W^{(k+1)}}(0)=\mathcal M^kU^{(k)}_+$, $U^{(k)}_+$ a background over $W^{(k+1)}$. The backgrounds are nested: $U^{(k)}_+$ is a background over $W^{(k)}$ as well (Lemma~\ref{lem:sf-class-window}(2)), so the regularity number over $W^{(k)}$ is exactly $L^{-2}$ times that over $W^{(k+1)}$ (the identity $\mathsf m_k=L^{-2}\mathsf m_{k+1}$), with no $C_H\delta(x_k)$; hence $\varpi(W^{(t)})<2L^{-2j_1}c_\rho\varepsilon(x_{t+j_1})\le2L^{-3j_1/2}c_\rho\varepsilon_n\le2\varepsilon_n/C''$, $W^{(t)}\in\mathsf D_n$ by Lemma~\ref{C-367-11}(a), so $W^{(t)}\in\mathcal K_1$ and $h_{t+j_1}(W)=h(W^{(t)})=1$.

On the support of the small-field weight of the chain from the depth $n-j_1$, $V_{t+j_1}\in\mathsf{Cut}_{t+j_1}$ and $V_k=V_{V_{k+1}}(B_k)\in\mathsf D'_k$, so $h(V_t)=h_{t+1}(V_{t+1})=\dots=h_{t+j_1}(V_{t+j_1})=1$; as $V_t\in\mathcal K_0$ and $h=0$ on $\mathcal K_2$, $V_t\in\mathcal K_1$. Of the dependence on $W$ this argument uses the continuity of the kernels alone. So $\sigma^K_n(\Omega_n\setminus\mathcal K_1)\le\mathsf R_{n-j_1}$ by Lemma~\ref{C-367-7}(i) for the chain from depth $n-j_1$. Glue finitely many chart balls in $\mathsf D_n$ to a connected $S'\supset\mathcal K_1$ (Lemma~\ref{C-367-11}(b), (c)), put $w_n=\min\{1,\rho/\rho_0\}$, $\rho_0=\min_{\mathcal K_1}\rho>0$, so that $w_n=1$ on $\mathcal K_1$ and $w_n=0$ off $\mathsf D_n$: this is (a).

Lemma~\ref{C-367-11}(c) with the density of $\sigma^K_n$, $\sup\mathsf g\le e^{\mathsf o_n}\inf\mathsf g$ (Lemma~\ref{C-367-7}(iii)), and $\Gamma_n=a_n^{-1}|\partial\cdot|^2$ give (b) with $C^{\rm t}_n=a_nc(S')e^{\mathsf o_n}/\rho_0$. (c): the linearisation of (b) gives $\int(\phi-c)^2w_n\,d\sigma^K_n\le\frac12C^{\rm t}_n\int\Gamma_n(\phi)w_n\,d\sigma^K_n$, and the part where $w_n<1$ costs at most $\mathrm{osc}(\phi)^2\mathsf R_{n-j_1}$. (d): $\mathrm{Ent}_{\sigma^K_n}(\phi^2)\le\int\Xi(\phi^2,c')\,d\sigma^K_n=\mathrm{Ent}_{w_n\sigma^K_n}(\phi^2)+\int\Xi(\phi^2,c')(1-w_n)\,d\sigma^K_n$. (e): assume \textup{(Multiscale gauge)}. The conditions of $n_{\rm e}$ are those of \textup{(Multiscale gauge)} at $J=J_n$; a field of $\mathsf I_n$ whose block fields $\mathcal M^jV$, $j\le J_n$, lie in the cuts has a gauge with all bonds within $44\pi L^2\sum_{j<J_n}\theta_{n-j}+C_{\rm S}L^{-J_n}\le\frac14\alpha_n$ of $\mathbf 1$, so $w^{\rm e}_n\ge c_n$ there by Lemma~\ref{C-367-11}(f), $\chi_n=1$, and $\sigma^K_n(\chi_n<1)\le\mathsf R_{n-J_n}$ by Lemma~\ref{C-367-7}(i) for the chain from depth $n-J_n$.

Since $r'_n=2\arcsin\varepsilon_n=64\alpha_n^2$ is the angle of the radius $2\varepsilon_n$, $\mathsf I_n=\{\mathsf d(\partial V(p),\mathbf 1)<r'_n\ \forall p\}$, and the identity of Lemma~\ref{C-367-11}(f) carries the log-Sobolev inequality of Haar measure on $\exp\mathscr D$ to invariant functions under $w^{\rm e}_ndV$, then to $\chi_ndV$ at the cost $c_n^{-1}$, and to $\chi_n\sigma^K_n$ at the cost $e^{2\mathsf o_n}/(1-\mathsf R_n)$ of the density: this is $C^\flat_n$. The set $\{\chi_n>0\}\cap\mathsf I_n$ is the path connected set of Lemma~\ref{C-367-11}(f); it contains $\mathbf 1$, so it lies in $\mathsf D_n$. For the block-field laws, $\mathrm{Var}_{\mu^K_n}(\Phi)\le\int(\Phi-c)^2d\mu^K_n\le\int(\Phi-c)^2d\sigma^K_n+\mathsf R_n\mathrm{osc}(\Phi)^2$ by Lemma~\ref{C-367-7}(i), then (c) with $\chi_n$ and $\int\Gamma_n(\Phi)\,d\sigma^K_n\le\int_{\mathsf I_n}\Gamma_n(\Phi)\,d\mu^K_n$; for $\mu_n$ by total variation.

\emph{Step 3: weak Poincar\'e bounds and mixing, (A4).} Let $c$ be the mean of $H$ under $\varsigma^K$. The density $w_n(V_t)$ is a function of the first member of the chain, so it changes no kernel and no tower, and
\[
\int(H-c)^2d\varsigma^K=\int(H-h_t(V_t))^2d\varsigma^K+\int(h_t-c)^2w_n\,d\sigma^K_n\le\mathsf C_n\mathcal E_K(H),
\]
the first term by the martingale sum of Step~1, the second by (B1)(b) linearised and \eqref{C-367-10} at $k=t$; passing to $\mu^K$ costs $\mathsf r_n\mathrm{osc}(H)^2$ by (B2)(a). This is (a), first for smooth $H$ and then by closure. (b): $\frac d{dt}\mathrm{Var}(P^K_tH)=-2\mathcal E_K(P^K_tH)\le-\frac2{\mathsf C_n}[\mathrm{Var}(P^K_tH)-\mathsf r_n\mathrm{osc}(H)^2]$, as $\mathrm{osc}(P^K_tH)\le\mathrm{osc}(H)$, and $\mathrm{Var}(H)\le\frac14\mathrm{osc}(H)^2$. (c): truncated and gauge-averaged recovery sequences and (a); if $\mathcal E_\star(F)=0$, then $\mathrm{Var}_\mu(F)\le\mathsf r_n\mathrm{osc}(F)^2$ for all admitted $n$. For non-constant bounded $F\in L^2_{\rm inv}$ the spectral measure $\nu_F$ is not carried by $\{0\}$, as else $\mathcal E_\star(F)=0$, and $\nu_F(\{\infty\})=0$ (Lemma~\ref{C-367-6}(c)); so $(F,(1-Q_t)F)=\int(1-e^{-t\lambda})\,d\nu_F>0$.

(d): $\varphi(t)=\|Q_{t/2}(F-\int F)\|^2$ has $\varphi'(t)=-\mathcal E_\star(Q_{t/2}(F-\int F))\le-\mathsf C_n^{-1}[\varphi(t)-\mathsf r_n\mathrm{osc}(F)^2]$ by (c) and spectral calculus, and $\varphi(0)\le\frac14\mathrm{osc}(F)^2$. If $Q_t\rho=\rho$, then $(\rho,(1-Q_t)\rho)=0$, $\nu_\rho$ sits at $0$, $\mathcal E_\star(\rho)=0$, so $\mathcal E_\star(\rho\wedge M)=0$ for every $M$, and $\rho\wedge M$ is constant by (c): $\rho=1$. (e): apply (a) to $H_K$ truncated at $\pm M_F$, which does not raise the energy and keeps $\|H_K-F_K\|\to0$, so that $\mathrm{Var}_{\mu^K}$ of the truncation tends to $\mathrm{Var}_\mu F$; the second form is (a) itself.

(f): the same lines with the explicit top of (B1)(e).

\emph{Step 4: entropy along the chain and the conditional law, (B2), (B3).} (a): Lemma~\ref{C-367-7}(i) and $\int f(1-w_n(V_t))\,d\sigma^K\le\sup f\,\sigma^K_n(w_n<1)$. (b): with $f_k$ the tower of $G^2+\epsilon$ and $\phi_k=f_k^{1/2}$, since the $p_W$ are probabilities the entropy splits along the chain,
\[
\mathrm{Ent}_{\varsigma^K}(G^2+\epsilon)=\sum_{k<t}\int\mathrm{Ent}_{p_W}(f_k\circ V_W)\,\varsigma^K_{d-1}(dW)+\mathrm{Ent}_{w_n\sigma^K_n}(f_t);
\]
the fibres by Lemma~\ref{C-367-7}(iv), $\mathrm{Ent}_{p_W}(\phi_k^2\circ V_W)\le2C_{\rm P}\frac{a_d^2}{\mathsf b}\int\Gamma_d(\phi_k)\circ V_W\,dp_W$; the top by (B1)(b); and $\int\Gamma_d(\phi_k)\,d\varsigma^K_d\le3\mathsf M^2\Pi^\flat_n\int\Gamma_K(G)\,d\varsigma^K$ (Proposition~\ref{C-367-9}(b), (c), as $\Gamma_K((G^2+\epsilon)^{1/2})\le\Gamma_K(G)$); sum over $d>n$ and let $\epsilon\to0$. (c): $\mathrm{Ent}_\nu(G^2)\le\int\Xi(G^2,c')\,d\nu$ and (a). (d): under \textup{(Multiscale gauge)}, the function $\mathsf b\mapsto2\varepsilon(L^n\mathsf b)$ is smooth with a non-vanishing derivative where $\log L^n\mathsf b>2\mathfrak p_0$, so $2\varepsilon_n\notin\mathcal N$ for almost every $\mathsf b$, for all $m$ and $n$ at once; by Lemma~\ref{C-367-11}(d) $\mathsf D_n$ is then a connected Lipschitz domain, which carries a log-Sobolev inequality for the volume (the extension theorem, with the inequalities of Sobolev and Rellich), and the density bounds give $C^\sharp_n$; $\sigma^K_n(\mathsf I_n\setminus\mathsf D_n)\le\sigma^K_n(\chi_n<1)\le\mathsf R_{n-J_n}$ by Step~2(e).

For every $\mathsf b$, a radius $\tilde r\notin\mathcal N$ close to $2\varepsilon_n$ does the same at the cost $\eta$: the fields with $\chi_n=1$ and $\varpi<\tilde r$ lie in the component of $\mathbf 1$ in $\{\varpi<\tilde r\}$, since along the rays of Lemma~\ref{C-367-11}(f) the plaquette angles stay below the larger of $32\alpha_n^2$ and their final values, and the shells $\{\tilde r\le\varpi<2\varepsilon_n\}$ have $\sigma^K_n$-masses at most $e^{\mathsf o_n}$ times their relative Haar measure, uniformly in $K$. With $\chi_n$ the same lines as in (b). (e) is the splitting of (b) with the top law of $\sigma^K$ itself. (B3): the splitting of (b), stopped at depth $e$ and taken at a fixed block field $W$, needs no top: with $\Pi_{(e)}=\prod_{e<d\le K}(1+\epsilon^\flat_d)\le\Pi^\flat_e$,
\[
\mathrm{Ent}_{\mathsf P_W}(G^2)\le\sum_{e<d\le K}2C_{\rm P}\frac{a_d^2}{\mathsf b}\,3\mathsf M^2\Pi_{(e)}\int\Gamma_K(G)\,d\mathsf P_W\le C^{\rm fib}_e\int\Gamma_K(G)\,d\mathsf P_W ,
\]
and the one-step bound of Proposition~\ref{C-367-9}(b), iterated pointwise, gives the bound on $\Gamma_e(g_{t'})(W)$; no factor depends on $\ell$. $\mathsf P_W$ is a regular conditional law by Lemma~\ref{C-367-7}(ii), also under $\varsigma^K$, whose weight is a function of the top. Concentration: for real invariant $\psi$ the log-Sobolev inequality, applied to $e^{s\psi/2}$, gives $\int e^{s(\psi-\psi_{t'}(W))}d\mathsf P_W\le e^{s^2C^{\rm fib}_el^2/4}$, as $\int\psi\,d\mathsf P_W=\psi_{t'}(W)$, then Chebyshev's inequality; for $\mu^K$, the weight $\mathfrak w^{\rm sf}_t$ of $\sigma^K$ vanishes unless $\mathcal M^{t'}U\in\mathsf{Cut}_{t'}$, and Lemma~\ref{C-367-7}(i) bounds what $\mu^K$ adds to $\sigma^K$ by $\mathsf R_n$; both tails.

\emph{Step 5: entropy relaxation, (A5).} (a) is (B2)(c) at $\nu=\mu^K$, with $\Xi(x,c)\le\bar M[1+\log(\bar M/c)]$ for $0\le x\le\bar M$, $0<c\le\bar M$, and $c'\ge m_2-\mathsf r_nM_G^2$, as $\int G^2d\varsigma^K\ge\int G^2d\mu^K-\mathsf r_nM_G^2$ and $\varsigma^K(\Omega_K)\le1$. (b): (B2)(b) for the normalised measure $\varsigma^K/\varsigma^K(\Omega_K)$ gives Gaussian concentration about $\bar\psi$, and (B2)(a) adds $\mathsf r_n$. (c): $P^K_\tau f_0$ is smooth and strictly positive for $\tau>0$, the kernel of the elliptic diffusion being positive, and
\[
\begin{multlined}
\tfrac d{d\tau}\mathrm{Ent}_{\mu^K}(P^K_\tau f_0)=-\mathcal E_K(P^K_\tau f_0,\log P^K_\tau f_0)\le-4\mathcal E_K((P^K_\tau f_0)^{1/2})\le-\tfrac4{\mathsf C'_n}\bigl[\mathrm{Ent}_{\mu^K}(P^K_\tau f_0) \\
-\mathsf r_n\bar M(1+\log\tfrac{2\bar M}{m_1})\bigr]
\end{multlined}
\]
by (a) at $G=(P^K_\tau f_0)^{1/2}$, with $M_G^2\le\bar M$, $m_2=m_1$ and $m_1-\mathsf r_n\bar M\ge\frac12m_1$; Gronwall. For limits: (a) by recovery sequences, as integrals of continuous functions of bounded strongly convergent sequences converge; the decay first for invariant cylinder $f_0$ with values in $[0,\bar M]$, by (A3), and then by approximation in $L^2_{\rm inv}$ under the strict condition $\mathsf r_n\bar M<\frac12m_1$; $\tau\to\infty$ and then $n\to\infty$ give $\mathrm{Ent}_\mu(Q_\tau F)\to0$, and Pinsker's inequality gives the convergence in total variation.

\emph{The moments of the gauge-fixed field.} Steps~6--10 use Proposition~\ref{C-367-12}(a) through the following facts. Use the consequence (Adm) of (E) proved in Proposition~\ref{C-367-12}, and let $d_1$ be a depth, free of $K$, such that every top depth $\ge d_1$ is admissible (the depth $d_{\rm adm}$ of the proof of Proposition~\ref{C-367-12}(a)). With $\mathbb E^{[d]}$ the normalised chain measure from the top depth $d$ and $\mathbf B_K$ its gauge-fixed field: (i) for every admissible $d$, $K>d$, $q<\infty$ and every index $\kappa'\in(0,\frac16)$,
\begin{equation}\label{C-367-32}
\mathbb E^{[d]}\bigl[\|\mathbf B_K\|^q_{\kappa';a}\bigr]^{1/q}\ \le\ C(q,\kappa')\,L^{20d},
\end{equation}
with $C$ free of $K$ and $d$: given the coarser layers, the layer of depth $d'$ is a function of one fluctuation integral with conditional covariance at most $Ca_{d'}^2/\mathsf b$ and Gaussian tails (Lemma~\ref{C-367-7}(iv)), the products of layers of two depths in the bilinear part of the rough norm are martingale differences, those of one depth are controlled by the bound (Q) of centred quadratic forms and by the cancellation in colour (Col), the conditional means by (M), all three set in the proof of Proposition~\ref{C-367-12}(a), and the top field, of size at most a constant times $L^{10d}$ in its regular gauge, enters to the second power; (ii) the short-time theory of the gauge-fixed flow, which uses no measure: the numbers $c_1,c_2,c_3$, $s_0(R)$, $a_0(R)$ at every index $\kappa'$, the four bounds of $A_s$, $\nabla A_s$, $A_s-P_s\mathbf B$, $\nabla(A_s-P_s\mathbf B)$, monotone in $R$, the bound of the flowed action and the Lipschitz bound in the rough distance; (iii) for invariant Borel $F\ge0$ and $r>1$,
\begin{equation}\label{C-367-33}
\langle F\rangle_{\mu^K}\ \le\ \mathbb E^{[d_1]}[F]+\sum_{d_1\le d\le K-2}\mathsf R_d^{1-1/r}\,\mathbb E^{[d+1]}\bigl[F^r\bigr]^{1/r}+\mathsf R_{K-1}\sup F ,
\end{equation}
from the chains of all the top depths $d\ge d_1$, each compared with the bare density by Lemma~\ref{C-367-7}(i), and H\"older's inequality; and $\sum_d\mathsf R_d^{1/2}L^{pd}$ and $\sup_K\mathsf R_{K-1}L^{pK}$ are finite for every $p$, as $\mathsf R_d$ decays faster than every exponential in $d$; (iv) hence, with $\mathsf N'=\mathsf N'(U)$ of the head of Proposition~\ref{C-367-12} formed at the index $\kappa'$, the least $\|\mathbf B\|_{\kappa';a}$ over the $1$-forms $\mathbf B$ with $e^{ia_K\mathbf B}$ in the orbit of $U$, which is at most a fixed power of $L^K$ at every field, $\mathsf M_q(\kappa')=\sup_K\langle\mathsf N'^q\rangle_{\mu^K}<\infty$, and the moments under $\mu^K$, bounded uniformly in $K$, of $\sup_{0<s\le1}s^{\frac32+\kappa'}A^a(U_s)$, of $(J,J)_a(U_\sigma)$ and of $a_K^{-4}\max_pa_p(U_\sigma)$ in Proposition~\ref{C-367-12}(a), at every index $\kappa'\in(0,\frac16)$.

\emph{Step 6: flow identities and uniform estimates, (C1), (C2).} (C1) is Lemma~\ref{C-367-6}(a) with the integration by parts, and $\int\mathcal L_KH\,d\mu^K=0$. The formula for $P^K_tH^\sigma$: $\frac d{ds}P^K_sH^{\sigma+\mathsf b(t-s)}=P^K_s(\mathcal L_K-\mathsf b\,\partial_\sigma)H^{\sigma+\mathsf b(t-s)}=P^K_s\Delta_KH^{\sigma+\mathsf b(t-s)}$, integrated over $[0,t]$; the martingale is It\^o's formula for the time-dependent function $(r,U)\mapsto H^{\sigma+\mathsf b(t-r)}(U)$, whose drift is $(\partial_r+\mathcal L_K)H^{\sigma+\mathsf b(t-r)}=\Delta_KH^{\sigma+\mathsf b(t-r)}$. (C2): on $\mathsf O_K(\bar C,\bar s)$, for all $K$ and $\sigma\le\bar s$, $\Lambda_\sigma\le(8\bar C\sigma^{\frac14-\frac\kappa2}+4\bar C^2\sigma^{\frac12-\kappa})/(1-2\kappa)$ (Lemma~\ref{C-367-13}(e)) and $A^a(U_{\bar s/2})\le E_\sharp=6\ell^3(\bar C^2(\bar s/2)^{-\frac12-\kappa}+\bar C(\bar s/2)^{-\frac34-\frac\kappa2})^2$.

For $K\ge K_0$, the least $K$ with $a_KE_\sharp\le e_0$, Lemma~\ref{C-367-13}(f) at $s_1=\bar s/2$ gives $a_K^{-4}\max_pa_p(U_s)\le Z_\sharp=\frac52E_\sharp\min\{\bar s/2,x_0^4E_\sharp^{-2}\}^{-3/2}$ for $s\ge\bar s$, so $\Lambda_\sigma\le\Lambda_{\bar s}+(2Z_\sharp)^{1/2}$ for $\bar s<\sigma\le1$. With $G=\varphi(\mathcal O(f_1),\dots,\mathcal O(f_r))$,
\[
\Gamma_K(H^\sigma_K)\le e^{24\Lambda_\sigma}\|\nabla\varphi\|_\infty^2\sum_i\Gamma_K(\mathcal O(f_i))(U_\sigma),\qquad|\partial_\sigma H^\sigma_K|\le(\Gamma_K(G)(J,J)_a)^{1/2}(U_\sigma),
\]
with $\frac\sigma2(J,J)_a(U_\sigma)\le e^{24\Lambda_\sigma}A^a(U_{\sigma/2})$, $A^a(U_{\sigma/2})\le24\ell^3\Lambda_\sigma^2\sigma^{-2}$ for $\sigma\le\bar s$ and $\le E_\sharp$ beyond, and Lemma~\ref{C-367-13}(d). (b): the same pointwise inequalities without a set, with the moments (iv) of Proposition~\ref{C-367-12}(a) above, by (Adm); the Lipschitz bound in $\sigma$ is the integral of $\partial_\sigma$. (c): the second flow-time derivative involves $v^2\mathcal O(f)$ along the flow, which is bounded by $\mathsf h_s\Gamma_K(\mathcal O(f))^{1/2}$ plus $\|f\|_\infty$ times the square of the $L^2$ norm of the covariant gradient of the current and a term of the first level, $\mathsf h_s=a_K^{-7/2}\|\nabla^2A_K(U_s)\,\partial A_K(U_s)\|$; a second level of the energy estimates rests on a floor of the Hessian on horizontal vectors by the action alone: Proposition~\ref{C-367-22}(a) at $k=N-1$, where $\mathsf e^*_N=2A^a$ and $a_K(N-1)\ge\frac12\ell$, gives $a_K^{-2}\nabla^2A_K(V)[X,X]\ge-c(A^a(V)^2+\ell^{-2})\|X\|^2$ for horizontal $X$ when $\max_pa_p(V)\le\frac12$, with an absolute $c$.

With it, for every field, $0\le\sigma_0<\sigma$ and $E=A^a(U_{\sigma_0})$,
\[
(J,J)_a(U_\sigma)\le e\,E\,\bigl[(\sigma-\sigma_0)^{-1}+2c(E^2+\ell^{-2})\bigr]+3072\,\ell^3E^6 .
\]
Where $a_KE\le\frac12$, $\max_pa_p(U_s)\le a_KA^a(U_s)\le\frac12$ for $s\ge\sigma_0$, as the action does not increase along the flow, so the floor applies to the current, which is horizontal and is carried by the linearised flow: $\partial_s(J,J)_a(U_s)\le2c(E^2+\ell^{-2})(J,J)_a(U_s)$, and the mean of $(J,J)_a(U_s)$ over the last interval before $\sigma$ of length $\min\{\sigma-\sigma_0,(2c(E^2+\ell^{-2}))^{-1}\}$ is at most $E$ divided by that length, as $\partial_sA^a(U_s)=-(J,J)_a(U_s)$, while $(J,J)_a(U_\sigma)$ is at most $e$ times that mean; where $a_KE>\frac12$, $a_K^{-1}<2E$, and the crude bound $(J,J)_a\le48\ell^3a_K^{-6}$ gives the last member. The second level needs more than this floor: with $g=\partial A_K(U_s)$, which is horizontal, and $\mathsf Q_s=a_K^{-5}\nabla^2A_K(U_s)[g,g]$, $\partial_s(J,J)_a(U_s)=-2\mathsf Q_s$ and
\[
\partial_s\mathsf Q_s=-2\mathsf h_s^2-a_K^{-7}\nabla^3A_K(U_s)[g,g,g],
\]
and the cubic member has no sign. Plaquette by plaquette, with $|X|_p^2=\sum_{b\in p}|X(b)|^2$ and the covariant forward differences $\nabla$ of the field, $|\frac{d^3}{dt^3}a_p(e^{itX}V)|_{t=0}|\le6|X|_p^2(|\nabla_\nu X_\mu(x)|+|\nabla_\mu X_\nu(x)|)+25\varpi_p|X|_p^3$ for $p=p_{\mu\nu}(x)$, since the cubic term of the plaquette is a sum of triple products of the four values of $X$ at $p$, carried to its base point with their signs, each of which holds the sum of the values at two parallel bonds, a covariant difference up to $2\varpi_p|X|_p$; so $|\nabla^3A_K(V)[X,X,X]|\le24\sqrt2\,\|X\|_{\ell^4}^2\mathsf D(X)+200(2A_K(V))^{1/2}\|X\|_{\ell^6}^3$, $\mathsf D(X)$ the $\ell^2$ norm of $\nabla X$.

Kato's inequality, $||X_\mu(x+a_Ke_\nu)|-|X_\mu(x)||\le|\nabla_\nu X_\mu(x)|$, and the Sobolev inequality of the periodic lattice with constants free of $N$, $\|u\|_{\ell^6}\le N^{-1}\|u\|_{\ell^2}+\frac4{\sqrt3}\|\nabla u\|_{\ell^2}$ and $\|u\|_{\ell^4}^4\le\|u\|_{\ell^2}\|u\|_{\ell^6}^3$, bound these norms through $\|X\|$ and $\mathsf D(X)+N^{-1}\|X\|$. The same tools bound the remainder of the Weitzenb\"ock identity of Lemma~\ref{C-367-13}(h), at most a constant times $\sum_p\varpi_p|X|_p^2$, by $\frac18\mathsf D(X)^2+c(A_K(V)^2+N^{-2})\|X\|^2$, $c$ taken large enough: this is the floor in the form that keeps the covariant gradient, $\nabla^2A_K(V)[X,X]\ge\frac18\mathsf D(X)^2-c(A_K(V)^2+N^{-2})\|X\|^2$ for horizontal $X$ when $\max_pa_p(V)\le\frac12$, whence at $X=g$, $\mathsf D(g)^2+N^{-2}\|g\|^2\le8\|g\|\,\|\nabla^2A_K(U_s)g\|+9c(A_K(U_s)^2+N^{-2})\|g\|^2$.

Let $\sigma_1=\min I$, $\sigma_0=\frac14\sigma_1$, $\Delta=\sigma_1-\sigma_0$ and $E=A^a(U_{\sigma_0})$. Where $a_KE\le\frac12$, the display bounds $(J,J)_a(U_s)$ by $\bar{\mathsf y}=\mathrm e\,E[2\Delta^{-1}+2c(E^2+\ell^{-2})]$ for $s\ge\sigma_0+\frac12\Delta$, and Young's inequality absorbs the cubic member into $\mathsf h_s^2$ ($N^{-1}=a_K/\ell$, $A_K=a_KA^a$):
\[
\partial_s\mathsf Q_s\le-\mathsf h_s^2+\mathsf B,\qquad\mathsf B\le C\Bigl[\bar{\mathsf y}^{7/3}+E^2\bar{\mathsf y}^3+\bigl((E^2+\ell^{-2})^{5/4}+(E^2+\ell^{-2})^{3/2}E^{1/2}\bigr)\bar{\mathsf y}^{3/2}\Bigr],
\]
$C$ absolute. Above, $\mathsf Q_s$ is at most its mean over the interval of length $\frac12\Delta$ before $s$, that is half the fall of $(J,J)_a$ over that interval divided by its length, plus $\frac12\mathsf B\Delta$; below, the floor gives $\mathsf Q_s\ge-c(E^2+\ell^{-2})(J,J)_a(U_s)$. So for $\sigma_1\le r_1<r_2$, $\int_{r_1}^{r_2}\mathsf h_s^2ds\le\mathsf Q_{r_1}-\mathsf Q_{r_2}+\mathsf B(r_2-r_1)$, and $a_K^{-5}\mathsf D(g)^2\le8\mathsf Q_s+8c(E^2+\ell^{-2})(J,J)_a(U_s)$ bounds the square of the $L^2$ norm of the covariant gradient of the current. With $A_f=\sum_pf(x(p))a_p$, $v^2\mathcal O(f)=a_K^{-5}[\langle\nabla^2A_K\,g,\partial A_f\rangle+\nabla^2A_f[g,g]]$, and $|v^2\mathcal O(f)|(U_s)\le\mathsf h_s\Gamma_K(\mathcal O(f))(U_s)^{1/2}+\|f\|_\infty[3a_K^{-5}\mathsf D(g)^2+4c(E^2+\ell^{-2})(J,J)_a(U_s)]$; integrated over $s$ between two points of $I$, with the Cauchy--Schwarz inequality in $s$, these are polynomials in $E$ with powers at most $11$, times the square root of the length.

Where $a_KE>\frac12$, $a_K^{-1}<2E$, and the crude bounds $|v^2\mathcal O(f)|\le1536\|f\|_\infty\ell^3a_K^{-8}$ and $|v\mathcal O(f)|\le48\|f\|_\infty\ell^3a_K^{-6}$ become powers of $E$, the square of the second one the power $12$. As $E\le(4/\sigma_1)^{5/3}\sup_{0<s\le1}s^{\frac32+\kappa}A^a(U_s)$, this gives $|\partial_\sigma H^{\sigma''}_K-\partial_\sigma H^{\sigma'}_K|\le C_I|\sigma''-\sigma'|^{1/2}(1+\sup_{0<s\le1}s^{\frac32+\kappa}A^a(U_s))^{12}$ at every field and every cutoff, and the moments of Proposition~\ref{C-367-12}(a), by (Adm), give (c).

\emph{Step 7: convergence of flowed observables along the full sequence, (C3).} Bounded strongly convergent sequences are stable under continuous functions, so $\varphi(\Theta_1,\dots,\Theta_r)$, $\Theta_i$ of Proposition~\ref{C-367-12}(c), under (Adm) and \textup{(Flow approximation)}, for $f_i$ applied to the flowed loops, converges strongly, and $|H^\sigma_K-\varphi(\Theta_1,\dots,\Theta_r)|$ is at most $\mathrm{Lip}(\varphi)\,r\,\eta$ on $\bigcap_i\mathsf G_K^{(i)}$ and at most $2\sup|\varphi|$ elsewhere; let $\eta,\epsilon\to0$. The joint convergence in law identifies the limit as $\mathsf G_\sigma$, an expression independent of the cutoff; only the Lipschitz constants of the $f_i$ and of $\varphi$ were used.

By (C2)(b), $\|\mathsf G_\sigma-\mathsf G_{\sigma'}\|=\lim_K\|H^\sigma_K-H^{\sigma'}_K\|\le C|\sigma-\sigma'|$. For $\dot{\mathsf G}_\sigma$: by (C2)(b), (c) the functions $\sigma\mapsto\int\Psi_K\partial_\sigma H^\sigma_K\,d\mu^K$ are bounded and locally H\"older continuous of exponent $\frac12$, uniformly in $K$; their primitives $\int\Psi_KH^\sigma_K\,d\mu^K$ converge to $\int\Psi\mathsf G_\sigma\,d\mu$; so the derivatives converge, locally uniformly, to the derivative of the limit, which is then continuously differentiable with a H\"older continuous derivative; the bound in $L^2$ and the density of the cylinder functions represent the limit as $\int\Psi\dot{\mathsf G}_\sigma\,d\mu$ with $\dot{\mathsf G}_\sigma\in L^2_{\rm inv}$, along the full sequence.

\emph{Step 8: condenser cutoffs and the generator limit, (C4), (C5).} Let $c'_3=c_3(\kappa/2)$, $R_\flat=2(2+c'_3)$, $\rho=R/2(2+c'_3)$, $q=4p/\kappa$; $\bar s(R)=2^{-j_0(R)}$ is the largest dyadic number $\bar s$ with $\bar s\le(c'_1\rho^{-2})^8$, which lies below the time $s_0(\rho)$ of Proposition~\ref{C-367-12}(a) at the index $\frac\kappa2$, $c'_1=c_1(\frac\kappa2)$; $c_{\rm st}R^2\bar s^{\frac14-\frac\kappa2}\le1$; and $486R^2\bar s^{\frac12-\kappa}+108R\bar s^{\frac14-\frac\kappa2}\le1$, which by Lemma~\ref{C-367-13}(e) at $\bar C=9R$ gives $\Lambda_s\le1$ for $s\le\bar s$ on $\mathsf O_K(9R,\bar s)$. $\bar K(R)$ is the least integer above the depth $d_1$ of the conclusion (Adm) at which $a=a_K$ has $a\le(c'_2/\rho)^3$, which lies below $a_0(\rho)$ at the index $\frac\kappa2$, $c'_2=c_2(\frac\kappa2)$, $c_{\rm st}Ra^{\frac12-\kappa}\le1$, and $4a^2\le\bar s(R)$.

Here $c_{\rm st}$ is an absolute constant of the short-time theory of the gauge-fixed flow; all these conditions are upper bounds of $\bar s$ and of $a$ that decrease in $R$. Let $g_j=\sup\{\mathsf m_s(U_s):s\in(2^{-j-1},2^{-j}]\cap[a_K^2,\bar s(R)]\}$, $g_\flat=a_K^{1/2+\kappa}\min\|\mathbf B\|_\infty$ over $e^{ia_K\mathbf B}$ in the orbit of $U$, $\mathsf G^*=\max\{g_\flat,\max_jg_j\}$, and $\chi_{K,R}$ the average of $\min\{1,\max\{0,4-2\mathsf G^*/R\}\}$ against a product of smooth conjugation-invariant approximate identities supported on small rotations of the bond groups. On $\{\mathsf G^*\le3R\}$, $\mathsf G^*\le c$ implies $U\in\mathsf O_K(3c,\bar s)$ and $\Lambda_s\le1$ for $s\le\bar s$; there the pieces are locally Lipschitz for the $L^2$ metric, $\mathrm{lip}\,g_j\le C_L2^{(j+1)(1-\kappa)/2}$, $\mathrm{lip}\,g_\flat\le2a_K^{-1+\kappa}$, $C_L=24e^{24}$, because the linearised flow smooths from $L^2$ modulo gauge: $\mathrm{Ad}(u)D\Phi_s(U)Y=\zeta+a_K^{-1}\mathrm d\vartheta$ with $\|\zeta\|_\infty\le9e^{24}s^{-3/4}\|Y\|_2$ and $\|\nabla\zeta\|_\infty\le17e^{24}s^{-5/4}\|Y\|_2$.

The rotations are chosen so small that $\mathsf G^*$ changes by at most $\frac R2$ under them; so $\chi_{K,R}=1$ on $\{\mathsf G^*\le R\}$, $\chi_{K,R}=0$ off $\{\mathsf G^*\le\frac52R\}\subset\mathsf O_K(8R,\bar s(R))$, and the piece $j$ counts in $\Gamma_K(\chi_{K,R})$ only where $g_j\ge R$. The short-time theory at the index $\frac\kappa2$, which gains $s^{\kappa/4}$, and $\mathsf M_q=\mathsf M_q(\frac\kappa2)<\infty$, the number of Proposition~\ref{C-367-12}(a) at the index $\frac\kappa2$ under (Adm) ((iv) of the moments of the gauge-fixed field), give $\mu^K(g_j\ge R)\le\mathsf M_q\rho^{-q}2^{-pj}$ and $\mu^K(g_\flat\ge R)\le\mathsf M_q(2e^{12}/R)^qa_K^{2p}$; with $2^{(j+1)(1-\kappa)p-pj}\le2^p2^{-\kappa pj}$,
\[
\mu^K(\chi_{K,R}<1)\le\mathsf M_q[2\rho^{-q}+(2e^{12}/R)^q],\qquad\int\Gamma_K(\chi_{K,R})^pd\mu^K\le\bigl(\tfrac2R\bigr)^{2p}\mathsf M_q\Bigl[\frac{C_L^{2p}2^p\rho^{-q}}{1-2^{-\kappa p}}+4^p\bigl(\tfrac{2e^{12}}R\bigr)^q\Bigr],
\]
and we call the first of these bounds $\epsilon_p(R)$ and the second $C_\chi(p,R)$, $q=4p/\kappa$. As $\Gamma_K(\chi_{K,R})$ vanishes off the shell $\{0<\chi_{K,R}<1\}$, H\"older's inequality on the shell bounds the energy, $\mathcal E_K(\chi_{K,R})\le C_\chi(p,R)^{1/p}\epsilon_p(R)^{1-1/p}\le c(p)\mathsf M_qR^{-2-q}$, the exponents of $R$ adding up to $-(2+\frac qp)-(q-\frac qp)$, while the mass is only $\int(1-\chi_{K,R})^2d\mu^K\le\mu^K(\chi_{K,R}<1)\le\epsilon_p(R)\le c(p)\mathsf M_qR^{-q}$. The cutoff $\chi_{K,R}$ does not depend on $p$, so one cutoff serves all $p$, and
\[
\bar\varkappa(R):=\min_{p\in\mathbb N,\ p\le R}\Bigl[C_\chi(p,R)^{1/p}\epsilon_p(R)^{1-1/p}+\epsilon_p(R)\Bigr]>0
\]
bounds $\mathcal E_K(\chi_{K,R})+\int(1-\chi_{K,R})^2d\mu^K$ for $K\ge\bar K(R)$ and is at most $c(p)\mathsf M_{4p/\kappa}R^{-4p/\kappa}$ for every $p\le R$; so $R^{q'}\bar\varkappa(R)\to0$ for every $q'$. The numbers $\bar s(R)$, $\bar K(R)$ are monotone in $R$, and $\mathsf G^*$ does not increase when $\bar s$ decreases; so for $R'\ge\frac52R$, $\{\chi_{K,R}>0\}\subset\{\mathsf G^*_{(R)}\le\frac52R\}\subset\{\mathsf G^*_{(R')}\le R'\}$, on which $\chi_{K,R'}=1$. The condenser cutoff: Lemma~\ref{C-367-15} at $\eta=1-\chi_{K,R}$, with $\varkappa\le\mathcal E_{K,1}(\eta,\eta)\le\bar\varkappa(R)$. (C5): by (C1) and Lemma~\ref{C-367-15}, $\int\chi^{\rm c}_{K,R}\mathcal L_KH\,d\mu^K=\mathcal E_K(w_{K,R},H)$, so $|\int\chi^{\rm c}_{K,R}\mathcal L_KH^\sigma_K|\le\mathrm{osc}(\varphi)\bar\varkappa(R)$; moreover $|\int\chi^{\rm c}\mathsf b\,\partial_\sigma H-\mathsf b\,\partial_\sigma\langle H\rangle|=\mathsf b|\int w\,\partial_\sigma H|\le\mathsf bC_\partial\bar\varkappa(R)^{1/2}$, which gives the second line of (a).

(b) is the last part of Lemma~\ref{C-367-15}, as $q(1-u)=1-q(u)$. (c): $\int\chi^{\rm c}(H'\mathcal L_KH-H\mathcal L_KH')=\mathcal E_K(w,HH')-2\int H\,\Gamma_K(w,H')$, and $\Gamma_K(w,\cdot)$ vanishes off the shell. (d): by (C1) and Lemma~\ref{C-367-15},
\[
\Bigl|\int\chi^{\rm c}F\mathcal L_KH-\int H\mathcal L_KF\Bigr|\le\sup|H|\bigl[2\|F\|_\infty\varkappa+(2\mathcal E_K(F)^{1/2}+\|\mathcal L_KF\|)\varkappa^{1/2}\bigr];
\]
with a test sequence, $\int H^\sigma_{K_j}\mathcal L_{K_j}F_j\,d\mu^{K_j}\to-\int(A_\star F)\mathsf G_\sigma\,d\mu$ by (C3); the convergence of resolvents in (A3) gives $\alpha(\alpha-\mathcal L_{K_j})^{-1}W_j\to\alpha(\alpha+A_\star)^{-1}W$ strongly, bounded by $\sup_j\|W_j\|_\infty$, hence test sequences; $c_F$ uses $\bar\varkappa\le\bar\varkappa^{1/2}$. The flow-time half: $\int\chi^{\rm c}F_j\partial_\sigma H^\sigma_{K_j}$ differs from $\partial_\sigma\int F_jH^\sigma_{K_j}$ by at most $\sup_j\|F_j\|_\infty C_\partial\bar\varkappa(R)^{1/2}$, and $\int F_jH^\sigma_{K_j}\to(F,\mathsf G_\sigma)$, locally Lipschitz in $\sigma$ uniformly in $j$; this gives the identity of distributions.

Reversibility: take $F_j=\alpha(\alpha-\mathcal L_{K_j})^{-1}H'_{K_j}$, a test sequence for $F=\alpha(\alpha+A_\star)^{-1}\mathsf G'$, and note that $A_\star\alpha(\alpha+A_\star)^{-1}$ is self-adjoint. The criteria: for $F^\Lambda=\Lambda(\Lambda+A_\star)^{-1}\mathsf G_\sigma$ with its test sequence, $(A_\star F^\Lambda,\mathsf G_\sigma)=\int\frac{\Lambda\lambda}{\Lambda+\lambda}d\nu_{\mathsf G_\sigma}$ increases to $\mathcal E_\star(\mathsf G_\sigma)$ as $\Lambda\to\infty$, and is at most $\sup|\varphi|\,C_1$ after $R\to\infty$; in $L^2$ the same bound with Cauchy--Schwarz holds on a core of $A_\star$.

\emph{Step 9: lower energy bounds, (C6).} By (Adm), let $\mathbb P_K$ be the normalised chain measure from the least top depth $d_0$ with $\mathsf R_{d_0}\le\frac12$ such that every top depth $\ge d_0$ is admissible, as in the head of Lemma~\ref{C-367-16}. Take $R$ with $\mathbb P_K(\|\mathbf B_K\|_{\kappa;a}>R)\le\frac14p_\ast$ (by \eqref{C-367-32} at $q=1$ and Markov's inequality, $R=\max\{1,4C(1,\kappa)L^{20d_0}/p_\ast\}$), and $d$ with $\theta_0a_{d+1}^2<\frac32\sigma\le\theta_0a_d^2$. On the complement $\{\|\mathbf B_K\|_{\kappa;a}\le R\}$ of this event and on the event of Lemma~\ref{C-367-16}(a) at $s=\frac32\sigma$, which meet in a $\mathbb P_K$-probability at least $\frac34p_\ast$, Lemma~\ref{C-367-16}(b) and the flow bounds of Proposition~\ref{C-367-12}(a) give, for $\sigma\le\sigma_\ast$, where $\kappa<\frac16$ makes the error of lower order,
\[
\Gamma_K(A^a\circ\Phi_\sigma)\ge cR^{-2}\mathsf b^{-2}\ell^3\sigma^{-5/2+\kappa},\qquad A^a\circ\Phi_\sigma\le S_\ast\sigma^{-3/2-\kappa},
\]
on a closed invariant $\mathsf L_K$ with $\mathbb P_K(\mathsf L_K)\ge\frac34p_\ast$, so $\mu^K(\mathsf L_K)\ge\frac38p_\ast$ (Lemma~\ref{C-367-7}(i), $\mathsf R_{d_0}\le\frac12$); hence $\mu^K(\mathsf L_K\cap\{\chi_K=1\})\ge\frac14p_\ast$, and $\Gamma_K(\varphi(A^a)\circ\Phi_\sigma)=\varphi'(A^a\circ\Phi_\sigma)^2\,\Gamma_K(A^a\circ\Phi_\sigma)$. The mean flowed action from below is the same argument at $s=\sigma$ with the last inequality of Lemma~\ref{C-367-16}(b): $A^a\circ\Phi_\sigma\ge\frac18k''_\Sigma\mathsf b^{-1}\ell^3\sigma^{-3/2}$, $k''_\Sigma=k_\Sigma\theta_0^{3/2}L^{-3}$, with $\mu^K$-probability $\ge\frac38p_\ast$, so $c_\circ=\frac3{64}p_\ast k''_\Sigma$.

From above, under $\mu^K$ and not only under a chain measure: on invariant functions the chain measure $\sigma^K$ from a top depth $d$ has a density $\bar{\mathfrak w}\in[0,1]$ with respect to $\mu^K$, with $\int(1-\bar{\mathfrak w})\,d\mu^K\le\mathsf R_d$ (Lemma~\ref{C-367-7}(i)), so for invariant $F\ge0$, by the Cauchy--Schwarz inequality, $\langle F\rangle_{\mu^K}\le\mathbb E^{[d]}[F]+\mathsf R_d^{1/2}\langle F^2\rangle_{\mu^K}^{1/2}$. Take $F=A^a\circ\Phi_\sigma$, with $\langle F^2\rangle_{\mu^K}^{1/2}\le\mathsf C\sigma^{-3/2-\kappa}$ by the moments (iv) of Proposition~\ref{C-367-12}(a), under (Adm), and $d=d_\sigma$, the least depth $d\ge d_1$ with $\mathsf R_d\le\sigma^{2\kappa}$, an admissible top depth by (Adm); then the second member is at most $\mathsf C\sigma^{-3/2}$, and as $\mathsf R_d$ decays faster than every exponential in $d$, $L^{d_\sigma}=\sigma^{-o(1)}$ as $\sigma\to0$.

Under $\mathbb E^{[d]}$, $d$ admissible, the centred layers of all depths, with conditional covariances at most $Ca_{d'}^2/\mathsf b$, give $\mathbb E^{[d]}\|\mathrm d_aP_\sigma\mathbf B_K\|_2^2\le c\,\mathsf b^{-1}\ell^3\sigma^{-3/2}$ against the heat kernel, with $c$ free of $d$ and $K$; the top field, the conditional means of the layers and the remainder of the logarithm add a lower power of $\sigma$ times a number that grows at most like a power of $L^d$; on $\{\|\mathbf B_K\|_{\kappa;a}\le\sigma^{-\epsilon}\}$, $\epsilon>0$ small, the last inequality of Lemma~\ref{C-367-16}(b) passes from $\|\mathrm d_aP_\sigma\mathbf B_K\|_2^2$ to $2A^a\circ\Phi_\sigma$ at the cost of a lower power of $\sigma$; and off this event the bound $A^a\circ\Phi_\sigma\le c_4\ell^3(1+\|\mathbf B_K\|_{\kappa;a})^{24}\sigma^{-3/2-\kappa}$ of Proposition~\ref{C-367-12}(a) with \eqref{C-367-32} at a large $q$ costs at most $CL^{pd}\sigma^{-1}$.

So $\mathbb E^{[d]}[F]\le c\,\mathsf b^{-1}\ell^3\sigma^{-3/2}+CL^{pd}\sigma^{-3/2+\epsilon'}$ for large $K$, with some $\epsilon'>0$, and at $d=d_\sigma$ the second member is $o(\sigma^{-3/2})$: with the bound $\mathsf R_{d_\sigma}^{1/2}\langle F^2\rangle_{\mu^K}^{1/2}\le\mathsf C\sigma^{-3/2}$ this gives $C_\circ$, for $0<\sigma\le\sigma_{\ast\ast}$. For $\vartheta<1$ and for the probability that tends to $1$: let $\psi\in C^\infty(\mathbb R)$ take values in $[0,1]$. By Proposition~\ref{C-367-22}(b) the square field of $\psi(\sigma^{3/2}A^a)\circ\Phi_\sigma$ is bounded on $\{\mathscr N\le R\}$ uniformly in $K$ and in small $\sigma$, so by Proposition~\ref{C-367-22}(d), that is, by the key estimate of Step~1, which rests on \eqref{C-367-10}, the function $\chi'_K\,\psi(\sigma^{3/2}A^a\circ\Phi_\sigma)$, $\chi'_K$ the cutoff of Proposition~\ref{C-367-22}(c), is within a distance in $L^2(\mu^K)$ that is small for large $\hat d_0$, uniformly in $K$ and in small $\sigma$, of a function of the block field at depth $\hat d_0$; under (Adm), every $\hat d_0\ge d_1$ is an admissible top depth, and $\mathsf R_{\hat d_0}\le\frac12$ for large $\hat d_0$.

By the last clause of Lemma~\ref{C-367-16}(a), at the top depth $\hat d_0$ in the place of $d_0$, off an event of probability $\tau_d$ the lower bounds hold, given the block field at depth $\hat d_0$, with a conditional probability that is bounded below by a positive number free of $K$, $\sigma$ and $\hat d_0$; a function that is close to a function of the block field and exceeds a level with conditional probability bounded below exceeds a slightly smaller level with high probability. So the flowed action exceeds a fixed multiple of $\mathsf b^{-1}\ell^3\sigma^{-3/2}$ with a probability that tends to $1$ as $\sigma\to0$, uniformly in $K$; at the flow time $\frac32\sigma$ this bounds $\|\mathrm d_aP_{3\sigma/2}\mathbf B_K\|_2^2$ from below, by the last inequality of Lemma~\ref{C-367-16}(b), and the first inequality there, on $\{\|\mathbf B_K\|_{\kappa;a}\le R\}$ for a large $R$, gives the same for $\Gamma_K(A^a\circ\Phi_\sigma)$ and $\sigma^{-5/2+\kappa}$.

By Markov's inequality and the upper bound of the mean flowed action, $\mu^K(A^a\circ\Phi_\sigma>S_\ast\sigma^{-3/2-\kappa})\le2C_\circ\mathsf b^{-1}\ell^3\sigma^\kappa/S_\ast$ for $0<\sigma\le\sigma_{\ast\ast}$ and all large $K$, which tends to $0$ with $\sigma$. Given $\vartheta<1$, choose $\sigma_\vartheta$ so that the probability of these lower bounds, together with $A^a\circ\Phi_\sigma\le S_\ast\sigma^{-3/2-\kappa}$, exceeds $\frac12(1+\vartheta)$ for all large $K$; the set $\{\chi_K=1\}$, of mass $\ge1-\vartheta$, meets this event in a mass $\ge\frac12(1-\vartheta)$. The current: $\Gamma_K(A^a\circ\Phi_\sigma)(U)\le e^{24\Lambda_\sigma(U)}(J,J)_a(U_\sigma)$ by Lemma~\ref{C-367-13}(b), since $\Gamma_K(A^a)=(J,J)_a$; by the tightness of $\Lambda_\sigma$ (Proposition~\ref{C-367-12}(a), under (Adm)) there is $\lambda$ with $\mu^K(\Lambda_\sigma>\lambda)\le\frac18p_\ast$ for large $K$, and on $\mathsf L_K\cap\{\Lambda_\sigma\le\lambda\}$, of mass $\ge\frac14p_\ast$, $(J,J)_a(U_\sigma)\ge e^{-24\lambda}cR^{-2}\mathsf b^{-2}\ell^3\sigma^{-5/2+\kappa}$.

For small $\sigma$ the same property of the current follows by a second route from the lower bound of the mean flowed action, through Proposition~\ref{C-367-23}(d)(iii); and, under \textup{(Lower propagation)}, by Proposition~\ref{C-367-23}(d)(i) it gives back, conversely, the positivity in the limit of the energies $\int\chi_K\Gamma_K(\varphi(A^a)\circ\Phi_\sigma)\,d\mu^K$, $\varphi$ as there, for every weight with $\limsup_K\mu^K(\chi_K<1)<\frac14p_\ast$.

\emph{Step 10: growth of uncut generator images, (C7).} Let $U\in\mathsf O_K(\bar C,\bar s)$, of $\mu^K$-mass $\ge1-\epsilon$ (Proposition~\ref{C-367-12}(a), by (Adm)). For large $K$ its flow in lattice time satisfies the hypotheses of Lemma~\ref{C-367-14} with $T=\sigma/a_K^2$: by Lemma~\ref{C-367-13}(e) at the flow times $s\ge a_K^{9/4}$, where $\kappa<\frac16$ gives $\varpi(U_s)\le(2\bar C+2\bar C^2)a_K^{1/8}$, and, at smaller flow times, by the bound $16$ of the speed of the plaquette variables in lattice time, $\sup_{t\le T}\varpi(\Phi_{a_K^2t}U)\le(4\bar C^2+20)a_K^{1/8}$, $\int_0^T\varpi\,dt=\Lambda_\sigma$ is bounded, and the action at $T/2$ is small.

Given the other bonds, the law of the three bonds of Lemma~\ref{C-367-14} under $\mu^K$ has a density $\ge e^{-20\beta}$ with respect to Haar measure, as they lie in ten plaquettes, and $20\beta=20\mathsf b\,a_KT/\sigma=o(T)$. So, with $\varepsilon=0.002$,
\[
\int|\partial(A_K\circ\Phi_\sigma)|^q\,d\mu^K\ge e^{-20\beta}(1-\epsilon)\,c\,e^{[(q-3)\lambda_1-6\lambda_2-(q+1)\varepsilon]T},
\]
and the enclosures of $\lambda_1,\lambda_2$ give (a); at $q=4$, $\Gamma_K(A^a\circ\Phi_\sigma)^2=a_K^{-6}|\partial(A_K\circ\Phi_\sigma)|^4$ and $\lambda_1-6\lambda_2>1.8$. (b): $G=A^a\circ\Phi_\sigma$ takes its values in $[0,6\ell^3a_K^{-4}]$, as $A_K\le6N^3$, and $H=\arctan G$ has $|\partial H|=a_K^{-1}|\partial(A_K\circ\Phi_\sigma)|/(1+G^2)$; so (a) at $q=4$ gives $\int|\partial H|^4d\mu^K\ge c(\ell)\,a_K^{28}e^{1.76\,\sigma/a_K^2}$. Lemma~\ref{C-367-13}(g) with $\sup|H|\le\frac\pi2$ gives $\frac\pi2(3\mathsf m+14\beta^{1/2}\mathsf m^{1/2})\ge c(\ell)^{1/2}a_K^{14}e^{0.88\,\sigma/a_K^2}$; the right side tends to $\infty$ and $\beta^{1/2}=(\mathsf b/a_K)^{1/2}$, so $\mathsf m\ge1$ for large $K$ and then $\mathsf m\ge c\,a_K^{14.5}e^{0.88\,\sigma/a_K^2}$; and $\|\mathcal L_KH\|=\mathsf m/a_K\ge e^{0.8\,\sigma/a_K^2}$ for large $K$.

(c): $f=\mathfrak f^2$ with $\mathfrak f(x)=\sqrt2|\sin(\pi x_1/\ell)|$, of Lipschitz constant $\sqrt2\pi/\ell$. Let $b_0$ be the bond in the direction $3$ at the site $a_K(k,0,0)$, $k\ge1$ an integer fixed below, $\alpha=2\pi a_K/\ell$, $\mathsf f_p=f(x(p))$ and $\mathsf f_0=1-\cos(\alpha k)\ge\frac2{\pi^2}\alpha^2k^2$. For $p$ at the distance $j$ from $b_0$, $|\mathsf f_p-\mathsf f_0|\le\alpha^2(j+3)(k+j+3)$, so $(\sum_p|\mathsf f_p-\mathsf f_0|^2e^{-2\theta|p|})^{1/2}\le C(\theta)\alpha^2(k+1)$; fix $k$ with $C_\sharp C(\theta)(k+1)\le c_\sharp k^2/\pi^2$. Then Lemma~\ref{C-367-14}(b) gives $|\partial_{b_0}(\mathcal O(f)\circ\Phi_\sigma)(W_\zeta)|\ge a_K^{-1}\frac{c_\sharp}{\pi^2}\alpha^2k^2e^{(\lambda_1-\varepsilon)T}$ for $\zeta\in\mathcal S$.

The value: with $\tilde V_t$ of Lemma~\ref{C-367-14} and $\mathcal A_f(V)=\sum_p\mathsf f_p\varsigma_pa_p(V)$,
\[
\mathcal O(f)(\Phi_\sigma W_\zeta)=a_K^{-1}\Bigl[\mathcal A_f(\tilde V_T)+2\sum_{p\ni b_0}\mathsf f_p\Bigr]\ge0 ,
\]
and the second term is at most $4a_K^{-1}\alpha^2(k+3)^2$, which tends to $0$. For the first, a local energy inequality holds along the flow of $\tilde A$ in lattice time: with $\nu=-\partial\tilde A(\tilde V_t)$ its velocity, $f_b=f(a_Kb_-)$ and $a'=3a_K\sqrt2\pi/\ell$,
\[
\tfrac d{dt}\mathcal A_f(\tilde V_t)=-\sum_bf_b|\nu(b)|^2+\sum_b\Bigl\langle\nu(b),\sum_{p\ni b}(\mathsf f_p-f_b)\varsigma_p\,\partial_ba_p(\tilde V_t)\Bigr\rangle\le64a'^2\sum_pa_p(\tilde V_t),
\]
since $|\mathsf f_p-f_b|\le a'(2|\mathfrak f(a_Kb_-)|+a')$ for $b\in p$, $|\partial_ba_p|\le\varpi_p$, $|\nu(b)|\le\sum_{p\ni b}\varpi_p$, and $\sum_b(\sum_{p\ni b}\varpi_p)^2\le32\sum_pa_p$. So $\mathcal A_f(\tilde V_T)\le\mathcal A_f(\tilde V_t)+64a'^2\int_t^T\sum_pa_p(\tilde V_u)\,du$ for $t\le T$; the mean over $t\in[\frac T2,T]$, $|\mathcal A_f|\le2\sum_pa_p$, Lemma~\ref{C-367-14}(c) and $A_K(U^\circ_{T/2})\le a_KE^{\rm a}_\sigma$ (Lemma~\ref{C-367-13}(c), (e)) give
\[
\mathcal A_f(\tilde V_T)\le\Bigl(\frac4T+64a'^2\Bigr)\Bigl(\frac32T\,a_KE^{\rm a}_\sigma+C_\sharp\Bigr)\le a_K\Bigl(\frac4\sigma+\frac{11400}{\ell^2}\Bigr)\Bigl(\frac32\sigma E^{\rm a}_\sigma+C_\sharp a_K\Bigr),
\]
as $64a'^2=1152\pi^2a_K^2/\ell^2\le11400\,a_K^2/\ell^2$. So $0\le\mathcal O(f)\circ\Phi_\sigma\le M_\sigma$ on the fields $W_\zeta$, $\zeta\in\mathcal S$, for large $K$, and $\int|\partial\varphi(\mathcal O(f)\circ\Phi_\sigma)|^4d\mu^K\ge(\inf_{[0,M_\sigma]}|\varphi'|)^4e^{-20\beta}(1-\epsilon)\,c\,a_K^4e^{(\lambda_1-6\lambda_2-5\varepsilon)T}$; Lemma~\ref{C-367-13}(g) as in (b). (d): the last identity of Lemma~\ref{C-367-13}(g) gives $|\mathcal L_K\cos G|+|\mathcal L_K\sin G|\ge\Gamma_K(G)$, and (a) at $q=2p$, whose exponent $(2p-3)\lambda_1-6\lambda_2$ is positive for $p>\frac32+3\lambda_2/\lambda_1$, which is less than $1.8637$; for $p>4$ the $L^p$ norm dominates the $L^4$ norm.

(e): by Lemma~\ref{C-367-13}(g) and $\nabla^2A_K\ge-24$, $a_K^2M^2+24\mathsf b\,M\ge\int\|\nabla^2H\|^2_{\rm HS}\,d\mu^K$ for $M=\|\mathcal L_KH\|$ and $\sup|H|\le1$, since $\int|\partial H|^2d\mu^K=-a_K\int H\mathcal L_KH\,d\mu^K$. Let $H=\sin(\mathcal O(f)\circ\Phi_\sigma)$; as $\sin$ is odd we may replace $f$ by $-f$, so that $f(x_*)=\max|f|>0$ at some $x_*$, and we base $b_0$ at a site nearest to $x_*$. For large $K$ every $U\in\mathsf O_K(\bar C,\bar s)$ satisfies the hypotheses of Lemma~\ref{C-367-14}(d) with $\Lambda_0=\max\{1,\sigma/\ell^2+(8\bar C\sigma^{\frac14-\frac\kappa2}+4\bar C^2\sigma^{\frac12-\kappa})/(1-2\kappa)\}$ (Lemma~\ref{C-367-13}(e)): the supremum of $\varpi$ along the flow tends to $0$ as above, $T/N^2=\sigma/\ell^2$, and $N^3T^{-2}=\ell^3a_K\sigma^{-2}\to0$.

The fibre $g\mapsto W^{(g)}$ of the bond $b_0$ is totally geodesic in $\Omega_K$, so $\mathfrak h(g)=H(W^{(g)})$ has $\|\nabla^2_{S^3}\mathfrak h\|_{\rm HS}(g)\le\|\nabla^2H\|_{\rm HS}(W^{(g)})$. For $(y_1,y_2)\in\mathsf Y$ of Lemma~\ref{C-367-14}(d), the function $r\mapsto\mathcal O(f)(\Phi_\sigma W^{(g_r)})=a_K^{-1}\mathsf v_f(r)$ falls by at least $f(x_*)\rho_2^2/a_K\ge2\pi$ on $[0,r_*]$ for large $K$, so $\mathfrak h$ takes the values $1$ and $-1$ at two points $g',g''$ with $\mathsf d(g',g'')\le r_*\le\frac43\rho_2e^{-\lambda'_1T}$; the inequalities of Sobolev and Morrey on $S^3$ give $|\mathfrak h(g')-\mathfrak h(g'')|\le C\|\nabla^2_{S^3}\mathfrak h\|_{L^2(S^3)}\mathsf d(g',g'')^{1/2}$, so $\int\|\nabla^2H\|^2_{\rm HS}(W^{(g)})\,dg\ge c\,e^{\lambda'_1T}$ for the Haar probability measure $dg$.

The set $\mathsf Y$ has Lebesgue measure $\ge c_{\rm d}e^{-6\lambda''_2T}$, the Haar measure in the chart $e^{iy_j}U(b_0+\hat e_j)$ has a density $\ge(4\pi^2)^{-1}$ for $|y_j|\le1$, and given the other bonds the law of the three bonds has a density $\ge e^{-20\beta}$; so $\int\|\nabla^2H\|^2_{\rm HS}d\mu^K\ge c(1-\epsilon)\exp[(\lambda'_1-6\lambda''_2)T-20\beta]$, with $\lambda'_1-6\lambda''_2>\frac25$. This gives $a_K^2M^2+24\mathsf bM\ge c\exp[\frac25\sigma/a_K^2-20\mathsf b/a_K]$, hence $M\ge\exp[\frac\sigma{5a_K^2}-\frac{10\,\mathsf b}{a_K}-C]$ for $\sin$ and every $f\not\equiv0$. For a non-constant periodic $\varphi$ the same steps hold with $\varphi$ in the place of $\sin$ (and $\varphi(-\,\cdot\,)$ in the place of $\varphi$ if $f$ was replaced by $-f$): for large $K$ the fall $f(x_*)\rho_2^2/a_K$ exceeds a period of $\varphi$, so that $\varphi$ takes its maximum and its minimum on the fibre within the distance $r_*$, $\mathrm{osc}\,\varphi$ replaces $2$, and $\sup|\varphi|$ replaces $1$ in the first inequality; this gives the exponent $\sigma/20a_K^2$ for large $K$.

(f): on the fields $W_\zeta$ of Lemma~\ref{C-367-14} the four plaquettes at $b_0$ stay within $\delta\le\frac1{10}$ of $-\mathbf 1$ during the flow, so $\Lambda_\sigma(W_\zeta)\ge(2-\delta)\sigma/a_K^2$, while their mass is at least $(1-\epsilon)e^{-20\beta}c_\sharp e^{-(3\lambda_1+6\lambda_2+\varepsilon)\sigma/a_K^2}$; and $13(2-\delta)\ge24.7>3\lambda_1+6\lambda_2+\varepsilon$. \qed

\subsection{Comparison of backgrounds at consecutive cutoffs}
\label{sec:dyn-background-comparison}

The convergence argument in Proposition~\ref{C-367-12}(b) first fixes the
coarse lattice and then removes the ultraviolet cutoff. For this step a
comparison whose constant depends on the coarse volume is sufficient. We
prove that comparison here. The separate estimate between a rough field
and its background, used subsequently in that proposition, remains
uniform in the cutoff.

We measure lengths in units of the fixed coarse lattice. Its torus has
side $2L^n$ and volume $N=(2L^n)^3$. Let $\eta=L^{-j}$, $j\ge1$, and
$\eta'=\eta/L$. The two fine lattices have spacings $\eta$ and $\eta'$.
For a configuration $U$ at spacing $\xi$, write
\[
 F_U(p)=\frac{\partial U(p)-\partial U(p)^{-1}}{2i\xi^2},
 \qquad J(U)=D_U^*F_U,
 \qquad
 \mathcal A^\xi(U)=\xi^{-1}\sum_p
       \left(1-\tfrac12\operatorname{tr}\partial U(p)\right).
\]
For quaternionic matrices we use the real scalar product
$\operatorname{Re}\operatorname{tr}(X^*Y)/2$; its norm agrees with the
operator norm and restricts to the stated norm on the Lie algebra.
All sums over bonds and plaquettes below use positive orientations. The
pairing and $L^p$ norms at spacing $\xi$ include the measure $\xi^3$.
The notation $D_U$ denotes the covariant exterior derivative. We call
$U$ regular at radius $\alpha$ when
\[
 \sup_p|\partial U(p)-1|<\alpha\xi^2,
 \qquad \sup_b|J(U)(b)|<\alpha.
\]
These are the constant-scale regularity conditions used in the
background construction. All constants in this subsection depend only
on the fixed blocking factor $L$. We use the H\"older exponent $1/2$.

\begin{proposition}[Comparison at a fixed coarse lattice]
\label{prop:dyn-fixed-coarse-comparison}
There are $\alpha_*>0$, $C<\infty$ and a fixed room radius $\rho_*$ with
the following property. Suppose the constant-torus background setting
is admitted, $L^n\ge3L\mathfrak M$ and $L^n>6\rho_*$, where
$\mathfrak M$ is its fixed analytic block size. Suppose that $U^A$ and
$U^B$ are regular real critical backgrounds, at radius
$0<\alpha\le\alpha_*$, on the lattices of spacings $\eta$ and $\eta'$,
respectively, and that they have the same coarse datum:
\[
 \mathcal M^j U^A=\mathcal M^{j+1}U^B.
\]
Then there are a gauge transformation $g$ and a real one-form $A$ on
the $\eta$ lattice such that
\begin{equation}\label{eq:dyn-background-comparison}
 (U^A)^g=e^{i\eta A}\mathcal M(U^B),\qquad
 \|A\|_\infty\le C\alpha N^{1/5}\eta^{1/5}.
\end{equation}
The constants $\alpha_*,C,\rho_*$ are independent of $n,j$ and of the
coarse datum. The factor $N^{1/5}$ is retained in the estimate.
\end{proposition}

The proof compares the first variations of the actions on the two
lattices. The average of a critical field need not be critical. Its
failure of criticality is small when tested against a coarse variation
and a compatible fine variation. We first quantify this fact.

\begin{lemma}[One-step regularity and the action defect]
\label{lem:dyn-one-step-defect}
Suppose $U$ is regular at radius $\alpha$ on the $\eta'$ lattice, with
$\alpha$ sufficiently small and $L^n>6\rho_*$. Put $W=\mathcal M(U)$.
Then $W$ is regular at radius $C\alpha$ on the $\eta$ lattice. For every
real fine one-form $h$,
\begin{equation}\label{eq:dyn-action-defect}
 \left|\frac d{dt}
 \left[\mathcal A^{\eta'}(e^{it\eta'h}U)
       -\mathcal A^\eta\bigl(\mathcal M(e^{it\eta'h}U)\bigr)\right]_{t=0}
 \right|
 \le C\alpha\eta^{1/2}\|D_Uh\|_1
        +C\alpha\eta'\|h\|_{1,\sup}.
\end{equation}
Here $\|h\|_{1,\sup}=\eta^3\sum_{p'}\sup_{b\subset Y(p')}|h(b)|$;
$Y(p')$ is the union of the four fine blocks at the corners of the
coarse plaquette $p'$.
\end{lemma}

\begin{proof}
Choose $\rho_*$ to be a fixed multiple of the radius in
Theorem~\ref{I-56-65}, large enough to contain each averaging stencil
with a collar of width four. At every coarse site that theorem, with
constant level, supplies a gauge $U=e^{i\eta' P}$ on this cube with
\begin{equation}\label{eq:dyn-local-regularity}
 |P|+|\nabla P|+|D_1^*D_1P|+|\Delta P|\le C\alpha,
 \qquad
 |\nabla P(x)-\nabla P(y)|\le C\alpha|x-y|^{1/2}.
\end{equation}
The enlarged cube fits in the torus by $L^n>6\rho_*$. Its radius is
fixed, so the smallness condition in that theorem follows by reducing
$\alpha_*$. In particular the constants in
\eqref{eq:dyn-local-regularity} do not depend on $\eta$ or $N$.

We give the local estimates needed for averaging. On the fixed stencil
of a coarse bond write
\[
 \Phi(a)=\frac1i\log\mathcal M(e^{ia}),\qquad
 \Phi(0)=0,\qquad \Phi'(0)=Q_{\rm lin}.
\]
The exponential and logarithmic formulas defining $\mathcal M$ are
holomorphic on a ball of fixed radius, since each path contains a
bounded number, depending only on $L$, of factors. Their first three
derivatives are bounded on a smaller fixed ball by Cauchy's formula.
The differential $Q_{\rm lin}$ is the block mean of straight path
sums. Summing the oriented boundary of each path surface gives
\begin{equation}\label{eq:dyn-linear-stokes}
 dQ_{\rm lin}a(p')=L^2\langle da\rangle_{p'},\qquad
 \langle f\rangle_{p'}=
 L^{-5}\sum_{x\in B(p'_-)}\sum_{p\in S_x(p')}f(p).
\end{equation}
The interior edges cancel in each surface $S_x(p')$; its remaining
four edges are precisely the four path sums in the left member.
The mean has total mass one. Every fine plaquette occurs with total
weight $L^{-3}$ when $p'$ is summed, and hence
\begin{equation}\label{eq:dyn-mean-counting}
 \eta^3\sum_{p'}\langle f\rangle_{p'}
       =\eta'^3\sum_p f(p).
\end{equation}

For later use we also record the elementary nonlinear curvature
bounds. If $V=e^{i\xi P}$, $|P|\le a\le1$ and $\xi\le1/13$, then
\[
 F_V=D_1P+\mathsf n(P),\qquad |\mathsf n(P)|\le16a^2.
\]
Indeed, expand the four exponentials of a plaquette to first order;
the remainder is at most the square of the sum of the four angles.
Cauchy's formula applied to this quadratic remainder gives
$|\mathsf n(P')-\mathsf n(P)|\le64a\max|P'-P|$.
In the current, the difference of the two nonlinear remainders on
neighboring plaquettes contributes at most $64a^2$ per direction
when $|\nabla P|\le a$. The parallel transport in the other term
differs from the identity by at most $2\xi a$, while
$|F_V|\le2a+16a^2$. Thus
\begin{equation}\label{eq:dyn-current-linearization}
 |J(V)-D_1^*D_1P|\le200a^2.
\end{equation}

Let $P^M=\eta^{-1}\Phi(\eta'P)$ in the gauge
\eqref{eq:dyn-local-regularity}. The derivative bounds for $\Phi$
give $|P^M|+|\nabla P^M|\le C\alpha$. Split it into
$P^{\rm lin}=\eta^{-1}Q_{\rm lin}(\eta'P)$ and $P^{\rm nl}$.
For $\Psi=\Phi-\Phi'(0)$, $|\Psi'(a)|\le C|a|$ and
$|\Psi''(a)|\le C$. A mixed second difference of
$\Psi(\eta'P)$ is the double integral of
\[
 \Psi''(a(s,t))[\partial_sa,\partial_ta]
       +\Psi'(a(s,t))\partial_s\partial_ta,
\]
where $a(s,t)$ is the bilinear interpolation of its four translated
arguments. Here $|a|\le C\eta'\alpha$,
$|\partial_sa|+|\partial_ta|\le C\eta'\eta\alpha$ and
$|\partial_s\partial_ta|\le C\eta'\eta\alpha$.
After division by $\eta^3$ this bounds
$|D_1^*D_1P^{\rm nl}|$ by $C\alpha^2$.

There is a cancellation needed to bound the linear term without a
negative power of $\eta$. Fix distinct directions $\mu,\nu,\lambda$,
put $F_{\mu\sigma}=\nabla'_\mu P_\sigma-\nabla'_\sigma P_\mu$ and
$g_\sigma(w)=\eta'^{-1}[F_{\mu\sigma}(w)-
F_{\mu\sigma}(w-\eta'e_\sigma)]$. The linear Stokes identity gives
\[
 (D_1^*D_1P^{\rm lin})_\mu
   =K_\lambda(g_\nu+g_\lambda)+(K_\nu-K_\lambda)g_\nu,
\]
where
\[
 K_\sigma f(x)=L^{-6}\sum_{z\in B(x)}
    \sum_{0\le i,i',r<L}
       f(z+i\eta'e_\mu+(i'-r)\eta'e_\sigma).
\]
The first term is bounded by $C\alpha$ because
$g_\nu+g_\lambda=(D_1^*D_1P)_\mu$ and $K_\lambda$ is a positive
average. To treat the second, multiply $P$ by a product cutoff equal
to one on the stencil, supported in its collar, whose first and second
differences are bounded independently of the mesh, and extend it by
zero. The product rule and \eqref{eq:dyn-local-regularity} give
$\|\Delta'\widetilde P\|_\infty\le C\alpha$.
The two Fourier multipliers expressing
$(K_\nu-K_\lambda)g_\nu$ in terms of
$\Delta'\widetilde P_\nu,\Delta'\widetilde P_\mu$ are
\[
 -\kappa(\vartheta)\frac{\overline d_\nu d_\mu}{\sum_\sigma|d_\sigma|^2},
 \qquad
 \kappa(\vartheta)\frac{|d_\nu|^2}{\sum_\sigma|d_\sigma|^2},
\]
where $d_\sigma=e^{i\vartheta_\sigma}-1$,
$s(t)=L^{-1}\sum_{|q|\le(L-1)/2}e^{iqt}$,
$e_L(t)=L^{-1}\sum_{q=0}^{L-1}e^{iqt}$ and
\[
 \kappa=s(\vartheta_\mu)e_L(\vartheta_\mu)
 s(\vartheta_\nu)s(\vartheta_\lambda)
 [s(\vartheta_\nu)^2-s(\vartheta_\lambda)^2].
\]
The last bracket vanishes to second order at zero. Each multiplier
$m$ therefore satisfies $|\partial^r m(\vartheta)|\le
C|\vartheta|^{2-r}$ for $0\le r\le2$ near zero and is smooth elsewhere.
These are also weak derivative bounds: the boundary terms on a sphere
of radius $\delta$ tend to zero as $\delta\downarrow0$.
Twice integrating by parts and using Bessel's inequality gives
\[
 \sum_{z\in\mathbb Z^3}|\widehat m(z)|
 \le \left(\sum_z(1+|z|^2)^{-2}\right)^{1/2}
       \|m-\Delta_\vartheta m\|_{L^2}<\infty.
\]
The convolution kernels thus have $\ell^1$ norms bounded only in
terms of $L$. This bounds the second term by $C\alpha$. Together
with \eqref{eq:dyn-current-linearization} it proves
$|J(W)|\le C\alpha$.

We prove the curvature bound and the action estimate together. In
$Y(p')$ take the comb gauge based at $p'_-$. A comb comparison with
the elementary plaquettes gives $U=e^{ia_0}$ in this gauge with
$|a_0|\le C\varpi$, where $\varpi=\sup|\partial U-1|
\le\alpha\eta'^2$. Put $G_p=\partial U(p)-1$ in this fixed gauge and
$\bar G=\langle G\rangle_{p'}$. By
\eqref{eq:dyn-linear-stokes}, the map
\[
 R(a)=\partial\mathcal M(e^{ia})(p')-1
               -L^2\langle\partial e^{ia}-1\rangle_{p'}
\]
has zero value and derivative at zero. The holomorphic bounds above
imply
\begin{equation}\label{eq:dyn-plaquette-remainder}
 |R(a_0)|\le C\varpi^2,\qquad
 |R'(a_0)\dot a|\le C\varpi|\dot a|.
\end{equation}
Consequently $|\partial W-1|\le L^2\varpi+C\varpi^2
\le C\alpha\eta^2$, completing regularity of $W$.

In the gauge \eqref{eq:dyn-local-regularity}, curvature varies across
a stencil by $C\alpha\eta^{1/2}$ in physical units: this follows
from the H\"older bound for $\nabla P$ and the Lipschitz bound for
$\mathsf n(P)$. Passing to the comb gauge changes the comparison of
two curvatures by at most $C\alpha^2\eta$ in those units. Indeed,
each comb has length $C\eta$, each bond angle is at most
$C\eta'\alpha$, and $|\operatorname{Ad}(v)-1|\le2|v-1|$.
Thus
\begin{equation}\label{eq:dyn-curvature-oscillation}
 |G_p-\bar G|\le C\alpha\eta'^2\eta^{1/2}.
\end{equation}
The identity $1-\tfrac12\operatorname{tr}v=|v-1|^2/2$ for
$v\in SU(2)$ and \eqref{eq:dyn-mean-counting} give the exact defect
\[
 \mathcal A^{\eta'}(U)-\mathcal A^\eta(W)
   =\eta^{-1}\sum_{p'}\delta_{p'},\qquad
 \delta_{p'}=\frac{L^4}2\langle|G-\bar G|^2\rangle
                  -L^2(\bar G,R)-\frac12|R|^2.
\]
Differentiate at $U_t=e^{it\eta'h}U$, retaining the comb gauge chosen
at $t=0$. If $s_{p'}=\sup_{Y(p')}|h|$, the product rule on a
plaquette and the derivative of the local logarithm give
\[
 |\dot G_p|\le\eta'^2|D_Uh(p)|+4\varpi\eta's_{p'},
 \qquad |\dot a|\le2\eta's_{p'},
 \qquad |\dot R|\le C\varpi\eta's_{p'}.
\]
In the derivative of the variance, the term containing
$\dot{\bar G}$ vanishes because $\langle G-\bar G\rangle=0$.
Therefore
\[
 \dot\delta_{p'}=L^4\langle(G-\bar G,\dot G)\rangle
      -L^2(\bar G,\dot R)-L^2(R,\langle\dot G\rangle)
      -(R,\dot R).
\]
Substitute \eqref{eq:dyn-plaquette-remainder} and
\eqref{eq:dyn-curvature-oscillation}, use $\eta'=\eta/L$ and
$\alpha\le1$, and divide by $\eta$. This yields
\[
 \eta^{-1}|\dot\delta_{p'}|
 \le C\eta^3\bigl[\alpha\eta^{1/2}
       \langle|D_Uh|\rangle_{p'}+\alpha\eta's_{p'}\bigr].
\]
Summing and applying \eqref{eq:dyn-mean-counting} proves
\eqref{eq:dyn-action-defect}.
\end{proof}

\begin{proof}[Proof of Proposition~\ref{prop:dyn-fixed-coarse-comparison}]
Put $W=\mathcal M(U^B)$. The preceding lemma makes $W$ regular at
radius $C\alpha$, and it has the same reduced datum as $U^A$.
We check the setting of the companion's chart results. Take
$\mathfrak M=M_{\rm ps}$ and choose once and for all $R_1\ge L$ with
$R_1\mathfrak M\ge w_S$. At depths $0,\ldots,j$ let every region be
the whole torus, and let the next region be empty. The required
$\mathfrak M$-cubes exist since $\mathfrak M\le L^n$; the separation
distances are infinite. Every layer below $j$ is empty, so both its
full and its reduced determining rows are empty. At level $j$ all
bonds are interior determining rows. There are no inward or outward
rows, seam nodes or edge rows. Thus both determining maps are exactly
$\mathcal M^j$, and both row spaces are the one-forms on the unit
lattice. The local scale is one everywhere, so every weighted norm
is the corresponding supremum. The same conclusions hold at the
finer cutoff with depth $j+1$.

The companion's chart uses the matrix divergence regularity class.
The current bound implies that class with radius $r+r^2$ when the
curvature radius is $r$: the traceless part of the plaquette
divergence is the current, and its scalar part is a difference of
$1-\tfrac12|\partial U-1|^2$ at neighboring plaquettes.
This is the real $SU(2)$ conversion of Lemma~\YMMref{9.4}{40}.
It costs only a fixed multiple of $\alpha$ here. The terminal torus
has $2L^n\ge6LM_{\rm ps}$ sites per direction. All geometric and
regularity hypotheses of the common-chart argument in the proof of
Proposition~\YMMref{37.3}{284} are therefore satisfied. After reducing
$\alpha_*$ it gives
\begin{equation}\label{eq:dyn-critical-chart}
 (U^A)^g=e^{i\eta A}W,\qquad
 R D_W^*A=0,\qquad q(A)=0,\qquad \|A\|_X\le C\alpha,
\end{equation}
where $\|A\|_X$ controls both $\sup|A|$ and the covariant first
differences. Write $F(a)=\mathcal A^\eta(e^{i\eta a}W)$ and
$F_B(b)=\mathcal A^{\eta'}(e^{i\eta'b}U^B)$.
Here $q$ is the nonlinear coarse constraint and $K_a=q'(a)$.
Lemma~\YMMref{30.8}{247} converts criticality on the actual averaging
fiber into the multiplier equation. Applying
Proposition~\YMMref{30.9}{248} with zero constraint value bounds its
multiplier, first at the reference $W$ and then at $U^B$. Its gradient
at the reference is bounded by $C\alpha$ in both cases, by regularity.
It follows that on the common unit lattice
\begin{equation}\label{eq:dyn-multipliers}
 F'(A)=K_A^*\omega_A,\qquad
 F_B'(0)=(K_0^B)^*\omega_B,\qquad
 \|\omega_A\|_\infty+\|\omega_B\|_\infty\le C\alpha.
\end{equation}
Here the two constraints are the actual relative constraints, before
any orthogonal change of row frame. Their averaging recursions agree
after the first fine step, as checked below.

We recall the precise linear and nonlinear estimates needed from
Propositions~\YMMref{29.1}{231} and~\YMMref{30.5}{245}.
With $Q$ the corner linear constraint and $T$ its fixed orthogonal
row transport, $T^{-1}q(a)=Qa+\mathsf C(a)$, $\mathsf C(0)=0$, and
\begin{equation}\label{eq:dyn-constraint-bounds}
 \|\mathsf C'(a)\|_{2\to2}\le C\alpha,
 \|\mathsf C'(a)^*\chi-\mathsf C'(b)^*\chi\|_\infty
       \le C\|a-b\|_\infty\|\chi\|_\infty
\end{equation}
on a fixed multiple of the $\alpha$ ball containing the segment
$[0,A]$. The first inequality follows also directly from the stated
row and column bounds by applying them to the symmetric operator
$\mathsf C'(a)^*\mathsf C'(a)$.
Put $\Delta_R=F''(0)+D_WRD_W^*$ and
$G=(\Delta_R+Q^*Q)^{-1}$,
$M=QGQ^*$, $H=GQ^*M^{-1}$ and
$\widetilde G=G-GQ^*M^{-1}QG$. These operators obey
\begin{equation}\label{eq:dyn-linear-operators}
 QH=1,\quad \widetilde G\Delta_R=1-HQ,\quad
 \Delta_RH=Q^*(M^{-1}-1),\quad
 \|H\|_{2\to2}+\|\widetilde G\|_{2\to2}\le C.
\end{equation}
Theorem~\YMMref{I}{225} gives positivity of $G$ in this real regular
setting, and Proposition~\YMMref{29.1}{231} gives positivity of $M$.
Thus both $M$ and $G$ are positive definite, and $\widetilde G$ is
positive semidefinite. To see the last assertion directly, complete
the square in $(f,Gf)-(QGf,M^{-1}QGf)$.
The $L^2$ bounds follow from the sup bounds for $\widetilde G,H,H^*$:
a symmetric operator bounded by $C$ in any norm has all eigenvalues
bounded by $C$, and one applies this also to $H^*H$.

These facts imply coercivity along $A$. For $Qa=0$,
$a=\widetilde G\Delta_Ra$ and the spectral inequality
$\widetilde G^2\le C\widetilde G$ give
$(a,\Delta_Ra)\ge c\|a\|_2^2$. For arbitrary $h$, set
$a=h-HQh$. The mixed term in its quadratic form vanishes by
\eqref{eq:dyn-linear-operators}, whereas the remaining $HQh$ term
is at least $-\|Qh\|_2^2$. Hence, whenever
$C\|Qh\|_2\le\|h\|_2$,
\[
 (h,\Delta_Rh)\ge
 c(\|h\|_2-C\|Qh\|_2)^2-\|Qh\|_2^2.
\]
The constraint in \eqref{eq:dyn-critical-chart} gives
$QA=-\mathsf C(A)$, so \eqref{eq:dyn-constraint-bounds} implies
$\|QA\|_2\le C\alpha\|A\|_2$. The gauge term vanishes on $A$.
Finally $F''(0)=D_W^*D_W+\Delta'$ with
$\|\Delta'\|_{2\to2}\le C\alpha$ by the real Hessian bounds in
the same results. Combining the two resulting lower bounds and
reducing $\alpha_*$ proves
\begin{equation}\label{eq:dyn-coercivity}
 F''(0)[A,A]\ge c x^2,\qquad
 x^2=\|D_WA\|_2^2+\|A\|_2^2.
\end{equation}

We next construct the fine test form used to compare first
variations. For a fine bond $b=(x,\mu)$, let $Y$ be the lower corner
of its coarse cell and $r=x-Y$, $0\le r_\sigma<\eta$. Interpolate
the four parallel coarse bonds $c_s$, with initial points
$Y+\eta s$, $s\in\{0,1\}^{\{\nu:\nu\ne\mu\}}$, using the weights
\[
 w_s(x)=\prod_{\nu\ne\mu}
    (r_\nu/\eta)^{s_\nu}(1-r_\nu/\eta)^{1-s_\nu}.
\]
Transport their values to $x$ along the coordinate combs in $U^B$
and denote their weighted mean by $h^\uparrow(b)$. There is no
length factor in this definition: $h$ and $h^\uparrow$ are
potentials in the same physical units. Define $h^\downarrow$ by
\[
 \eta h^\downarrow=
 Q_{\mathsf m}(U^B)(\eta'h^\uparrow),
\]
where $Q_{\mathsf m}$ is the derivative of the relative one-step
average, including its straightening gauge. The elementary
interpolation estimates are
\begin{equation}\label{eq:dyn-lift-bounds}
 \begin{split}
 \|D_{U^B}h^\uparrow\|_1
     &\le\|D_Wh\|_1+C\alpha\eta N\|h\|_\infty,\\
 \|h^\uparrow\|_{1,\sup}&\le3N\|h\|_\infty,\\
 \|h-h^\downarrow\|_1&\le C\eta N\|h\|_X.
 \end{split}
\end{equation}
Here is the finite-stencil verification. In a comb gauge on the
box containing the relevant neighboring cells, the fine angles are
bounded by $C\alpha\eta\eta'$. Every comparison path has at most
$CL$ bonds. Replacing its transport by the identity changes the
interpolated value by $C\alpha\eta^2\sup|h|$ and its covariant curl
by $C\alpha\eta\sup|h|$. Without these errors, subtracting the
two linear interpolants in a plaquette gives the linear
interpolation of $D_Wh$ in the one transverse direction.
Its weights are nonnegative; their total fine mass over each
coarse plaquette is $L^3$. This proves the first estimate on
multiplying by $\eta'^3$ and summing. The second follows from
$|h^\uparrow|\le\|h\|_\infty$ and the $3N\eta^{-3}$ coarse
plaquettes. At zero fine angles the relative derivative
$Q_{\mathsf m}$ is straight path averaging; applied to the
interpolant it is a convex mean of neighboring values of $h$.
At the given fine angles its derivative differs by
$C\alpha\eta^2\sup|h|$ after division by $\eta$, by the bounded
second derivative of the relative average. Neighbors differ by
at most $3\eta\|\nabla^Wh\|_\infty+C\alpha\eta^2\|h\|_\infty$.
This gives the last inequality in \eqref{eq:dyn-lift-bounds}.

Let $a_1(t)$ be the relative logarithm of the first average of
$e^{it\eta'h^\uparrow}U^B$. It has derivative
$a_1'(0)=\eta h^\downarrow$. After this first step, the reference
configurations in the two averaging recursions are identical.
Their subsequent relative logarithms therefore agree, by induction
over the remaining $j$ steps. Differentiating the final constraint
gives
\begin{equation}\label{eq:dyn-constraint-chain}
 K_0^Bh^\uparrow=K_0h^\downarrow.
\end{equation}
Gauge invariance of the action removes the straightening gauge at
the first step. Thus its action defect is exactly
\[
 E(h)=F_B'(0)[h^\uparrow]-F'(0)[h^\downarrow].
\]
Put $d=F'(0)-K_0^*\omega_B$. From
\eqref{eq:dyn-multipliers} and \eqref{eq:dyn-constraint-chain},
$(d,h^\downarrow)=-E(h)$, and hence
\begin{equation}\label{eq:dyn-critical-defect-identity}
 F'(0)[h]=(d,h-h^\downarrow)+(\omega_B,K_0h)-E(h).
\end{equation}
Regularity of $W$, the transpose row bound for $K_0$ and
\eqref{eq:dyn-multipliers} give $\|d\|_\infty\le C\alpha$.
Apply \eqref{eq:dyn-action-defect} and
\eqref{eq:dyn-lift-bounds} with $h=A$, using
$\|A\|_X\le C\alpha$ and
$\|D_WA\|_1\le(3N)^{1/2}\|D_WA\|_2$. We obtain
\begin{equation}\label{eq:dyn-defect-bound}
 |(d,A-A^\downarrow)|+|E(A)|
 \le C\alpha\eta^{1/2}N^{1/2}x+C\alpha^2\eta N.
\end{equation}

The constraint multipliers contribute only quadratically. To see
this, for fixed $\chi$ the Hessian of $(\chi,\mathsf C(a))$ is
symmetric and has sup operator norm at most $C\|\chi\|_\infty$
by \eqref{eq:dyn-constraint-bounds}. Its $L^2$ operator norm has
the same bound, by the eigenvalue argument used above. Since
$q(0)=q(A)=0$, the fundamental theorem of calculus gives
\[
 T^{-1}K_0A=-\int_0^1[\mathsf C'(tA)-\mathsf C'(0)]A\,dt,
 \qquad
 T^{-1}K_AA=\int_0^1[\mathsf C'(A)-\mathsf C'(tA)]A\,dt.
\]
Pair these identities with $T^{-1}\omega_B$ and $T^{-1}\omega_A$,
whose suprema are unchanged by the orthogonal row transport.
Consequently
\begin{equation}\label{eq:dyn-multiplier-quadratic}
 |(\omega_B,K_0A)|+|F'(A)[A]|\le C\alpha\|A\|_2^2.
\end{equation}

For completeness we bound the third derivative of the action;
a bound on four separate bond derivatives would lose powers of the
mesh. Put $\phi(t)=F(tA)$. At a plaquette with reference holonomy
$P$, its varied holonomy has the form
$Z(t)=e^{tX_1}\cdots e^{tX_4}P$, where
$|X_i|\le\eta a_p$ and $a_p$ is the largest $|A|$ on its boundary.
Write $Z'=\Sigma Z$. Differentiating the four partial products gives
\[
 |\Sigma'|\le12\eta^2a_p^2,\qquad
 |\Sigma''|\le72\eta^3a_p^3,\qquad
 \Sigma(0)=i\eta^2Y,
 \quad |Y-D_WA(p)|\le C\alpha\eta a_p.
\]
The last error is precisely the transport around $P$ rather than
the identity; $|P-1|\le C\alpha\eta^2$. Set
$y_p=|Y|+12a_p^2$. Then
$|\Sigma(t)|\le\eta^2y_p$ and
$|Z(t)-1|\le\eta^2(C\alpha+y_p)$ for $0\le t\le1$.
Both $\Sigma''$ and $\Sigma^3$ are traceless, while
$\Sigma^2=-|\Sigma|^2 1$, since $\Sigma\in\mathfrak{su}(2)$.
In
\[
 \frac{d^3}{dt^3}\left(1-\tfrac12\operatorname{tr}Z\right)
 =-\tfrac12\operatorname{tr}
    [\Sigma''+2\Sigma'\Sigma+\Sigma\Sigma'+\Sigma^3]Z
\]
one may therefore replace $Z$ by $Z-1$ in the first and last
terms. With $a_p+|D_WA(p)|\le C\alpha$ the displayed bounds give
\[
 \left|\frac{d^3}{dt^3}\left(1-\tfrac12\operatorname{tr}Z\right)\right|
 \le C\eta^4\alpha\bigl(|D_WA(p)|^2+a_p^2\bigr).
\]
Each bond belongs to four plaquettes. Multiplication by $\eta^{-1}$
and summation therefore prove $|\phi'''(t)|\le C\alpha x^2$.

Taylor's formula, \eqref{eq:dyn-coercivity},
\eqref{eq:dyn-critical-defect-identity},
\eqref{eq:dyn-defect-bound} and
\eqref{eq:dyn-multiplier-quadratic} now imply
\[
 cx^2\le C\alpha x^2+
       C\alpha\eta^{1/2}N^{1/2}x+C\alpha^2\eta N.
\]
Choose $\alpha_*$ so the first term is absorbed. Young's inequality
absorbs half of the remaining left member and yields
\begin{equation}\label{eq:dyn-mean-square-comparison}
 \|A\|_2^2+\|D_WA\|_2^2\le C\alpha^2\eta N.
\end{equation}

Finally let $s=\|A\|_\infty$ and choose $r=C\alpha$ so that
$\|A\|_X\le r/2$. The norm of a component changes by at most
$\eta r/2$ across an adjacent site: parallel transport is an
isometry. If $s/r<\eta$, then $s<r\eta$. Otherwise set
$R=\lfloor s/(r\eta)\rfloor\ge s/(2r\eta)$.
On the lattice $\ell^1$ ball of radius $R$ about a maximizing site,
that component has norm at least $s/2$. Its physical radius is at
most $1/2$, so its sites are distinct on the torus. This ball has
more than $4R^3/3$ sites, giving
\[
 \|A\|_2^2\ge\eta^3\frac{4R^3}{3}\frac{s^2}{4}
            \ge\frac{s^5}{24r^3}.
\]
Thus $s\le\max\{r\eta,(24r^3\|A\|_2^2)^{1/5}\}$.
Substitution of \eqref{eq:dyn-mean-square-comparison}, with
$N\ge1$ and $\eta\le1$, proves
\eqref{eq:dyn-background-comparison}.
\end{proof}

\subsection{Uniformity and the order of limits}

Theorem~\ref{C-367-2} holds for every limit, with constants and rates that do not depend on the subsequence. The inequalities of (A4), (A5) and (B2)(c) carry a defect, and their constants depend on $\ell$; the constants of Theorem~\ref{C-367-4} that are free of the volume are those of the fine field given a block field: $C^{\rm fib}_e$ and $3\mathsf M^2\Pi^\flat_e$ in (B3), in the log-Sobolev inequality of $\mathsf P_W$, in the pointwise bound of $\Gamma_e(g_{t'})$ and in the tail under $\mathsf P_W$, and the same two numbers at $e=n$ in (B2)(e), whose first inequality is (B3) at $t'=K-n$ integrated over the top law by the chain rule of the entropy.

The tail under $\mu^K$ in (B3) adds $\mathsf R_n$; the first inequality of (B2)(e) keeps the entropy of the top law, which (B1) bounds at the cost $C^{\rm t}_n$; and the constants, defects and thresholds of (B1) and (B2)(a)--(c), $C^{\rm t}_n$, $C^\flat_n$, $\mathsf C'_n$, $\mathsf R_\cdot$, $\mathsf r_n$, $n_{\rm c}$, depend on $\ell$. The flowed observables are taken at a fixed positive flow time; under \textup{(Flow approximation)}, (C5) identifies the generator on them as a distribution, in the tested outer-$R$, inner-$\limsup_{K\to\infty}$ error sense of (C5)(d), and (C7) shows that without cutoffs the images under $\mathcal L_K$ of some of them are not bounded uniformly in $K$: for $p>1.8637$, $\sup_K[\|\mathcal L_K\cos(A^a\circ\Phi_\sigma)\|_{L^p(\mu^K)}+\|\mathcal L_K\sin(A^a\circ\Phi_\sigma)\|_{L^p(\mu^K)}]=\infty$, so the norm is unbounded at $\cos$ or at $\sin$, and at $\cos(\cdot+\theta)$ for all $\theta\in[0,\pi)$ but at most one, since $\cos$ and $\sin$ are linear combinations of any two of these functions.

\section{A finite-lattice criterion for confinement}\label{sec:conf}

Let $Z(n_1,n_2,n_3)$ be the periodic partition function at coupling $\beta>0$. We prove a fixed-lattice implication from
\[
 2\log Z(l,2l,2l)-\log Z(2l,2l,2l)\le x_0,
 \qquad l\ge3,\qquad 0<x_0\le0.069.
\]
The transfer spectrum controls the doubling defect. A vortex-tube recursion then gives area laws, electric-flux bounds and a gap over finite sections; excess-pressure estimates give uniqueness among periodic Gibbs states and convergence of periodic boxes of arbitrary shape.

All constants and parity restrictions are stated below. Physical conclusions uniform in the cutoff require the criterion at a fixed physical side for every sufficiently large cutoff. The stronger numerical tension and the simpler tension of Proposition~\ref{prop:flat} are kept separate.

\subsection{Notation and hypotheses}

\emph{Boxes.} A \emph{box} is a triple $\underline n=(n_1,n_2,n_3)$ of integers $\ge3$, $|\underline n|=n_1n_2n_3$; the cube of side $n$ is $(n,n,n)$. On the periodic lattice of a box a configuration $U$ gives each bond $b$ an element $U(b)$ of $G=\mathrm{SU}(2)$, with $U(\bar b)=U(b)^{-1}$ for the reversed bond; $dU$ is the product of the Haar probability measures, $\partial U(p)$ the ordered product along the boundary of a plaquette $p$, and $a_p=1-\frac12\operatorname{tr}\partial U(p)\in[0,2]$. For $\beta\ge0$ we write $Z(\underline n)=\int e^{-\beta\sum_pa_p}dU$, which lies in $[e^{-6\beta|\underline n|},1]$, $F=\log Z$, and $\langle\cdot\rangle_{\underline n}$ for the mean under $Z^{-1}e^{-\beta\sum_pa_p}dU$. The pressure is $P=\inf_{\underline n}F(\underline n)/|\underline n|\in[-6\beta,0]$, the \emph{excess} of a box is $\delta(\underline n)=F(\underline n)-|\underline n|P\ge0$, and $\delta_{\rm c}(n)=\delta(n,n,n)$. For integers $l\ge3$ we put
\[
\begin{multlined}
\mathsf e(l)=2F(l,2l,2l)-F(2l,2l,2l),\quad d_1(l)=2F(l,l,l)-F(2l,l,l),\quad d_2(l)=2F(l,2l,l) \\
-F(2l,2l,l),\quad X(l)=8F(l,l,l)-F(2l,2l,2l).
\end{multlined}
\]

\emph{Twists.} For directions $\mu<\nu$ with third direction $\rho$, a \emph{column} of the pair $\mu\nu$ is a set of plaquettes $p_{\mu\nu}(x+je_\rho)$, $0\le j<n_\rho$. For $z=(z_{12},z_{13},z_{23})\in\{0,1\}^3$, $Z^z$ is $Z$ with the sign of $\operatorname{tr}\partial U(p)$ reversed on one column of each pair $\mu\nu$ with $z_{\mu\nu}=1$, and the \emph{twisted box law} is the probability measure with this weight; $|z|=z_{12}+z_{13}+z_{23}$. We write $Z^-=Z^{(0,1,0)}$, $a_\pm=\frac12(Z\pm Z^-)$ and $z_\pm=a_\pm/Z$; $z_-$ is the \emph{flux weight}. On a cube, $r_j$ is the quotient $Z^z/Z$ for $|z|=j$ (Lemma~\ref{lem:tr}(c)); so $r_1=Z^-/Z=1-2z_-$.

\emph{Sections.} A direction $\mu_0$ taken as time has the side $t=n_{\mu_0}$; the two others, $\nu_1<\nu_2$, span the \emph{section} $\Sigma=(N_1,N_2)=(n_{\nu_1},n_{\nu_2})$, and $Z(t,\Sigma)=Z(t,N_1,N_2)$, likewise for $F$ and $\delta$; by Lemma~\ref{lem:tr} these numbers do not depend on the order of the sides. Over $\Sigma$, $\lambda_0>\lambda_1\ge\lambda_2\ge\dots\ge0$ are the eigenvalues of the transfer operator of Lemma~\ref{lem:tr}, with $Z(t,\Sigma)=\sum_i\lambda_i^t$ for $t\ge3$; $\lambda_0$ is the \emph{vacuum}, $\mathsf r_t(\Sigma)=\sum_{i\ge1}(\lambda_i/\lambda_0)^t$ the \emph{thermal number}, and a number $y>0$ \emph{lies at least $\eta$ above the vacuum} if $\log(\lambda_0/y)\ge\eta$.

The eigenvalues fall into four flux sectors $e\in\{0,1\}^2$ (Lemma~\ref{lem:tr}); $\Lambda_e$ is the top of the sector $e$, the sectors $e\ne00$ are those with electric flux, and over a section $(N_2,N)$, with the first direction as time, so that $\nu_1=2$ runs along the side $N_2$ and $\nu_2=3$ along the side $N$, $\Lambda_-=\max\{\Lambda_{01},\Lambda_{11}\}$ is the top of the sectors that are odd under $\mathsf J_2$, that is, under the \emph{flip of the side $N$}: the map $\iota_2$ of Lemma~\ref{lem:tr}, which changes the sign of the bonds along the side $N$ that start on one line along the side $N_2$. These sectors carry an electric flux along the side $N$, a flux line of the length $N$. With $\lambda_-(t,N)=\lim_sZ^-(t,s,N)^{1/s}$, the number $\varepsilon_{\rm v}(t,N)=\log(\lambda_0(t,N)/\lambda_-(t,N))\ge0$ is the free energy per unit length of the vortex line that runs along the second direction of the box $(t,s,N)$; $\varepsilon_{\rm v}(l')=\varepsilon_{\rm v}(l',l')$.

\emph{Functions.} $\varpi(x)=x/(1+\sqrt{1-x^2})$ on $[0,1]$; $\omega(D)=2\varpi(e^D-1)^2$ on $[0,\log2]$; $Y_1(x)=2x+\omega(x)$, $Y_2=2Y_1+\omega(Y_1)$, $\Psi=\omega(Y_2)$ and $\Phi=4\omega+2\omega(Y_1)+\omega(Y_2)$, where the arguments lie in $[0,\log2]$.

\emph{Cells.} For an integer $h\ge2$, $N=2h$, and a point $c\in(\mathbb Z+\frac12)^3$, the \emph{cell} $\mathsf C_h(c)$ is the cube of the $h$ layers below $c$, formed by the sites $x$ with $c_\mu-h<x_\mu<c_\mu$, $\mu=1,2,3$; a \emph{function of the cell} is a bounded Borel function of the bonds with both ends in it. On the cube $(N,N,N)$ the reflections $\mathfrak r_\mu\colon x_\mu\mapsto2c_\mu-x_\mu$ map the lattice onto itself, each fixing the two planes $x_\mu=c_\mu$ and $x_\mu=c_\mu+h$; for a function $A$ of the cell and $\epsilon\in\{0,1\}^3$ the \emph{copy} $A^{[\epsilon]}$ is $A$ composed with the $\mathfrak r_\mu$ with $\epsilon_\mu=1$ and complex conjugated if $|\epsilon|$ is odd, and
\[
|\!|\!|A|\!|\!|_N=\Bigl\langle\prod_{\epsilon\in\{0,1\}^3}A^{[\epsilon]}\Bigr\rangle_{(N,N,N)}^{1/8},
\]
the mean being a number $\ge0$ by Lemma~\ref{lem:mir}.

\emph{The infinite lattice.} $\Omega=G^{\mathbb B}$, $\mathbb B$ the bonds of $\mathbb Z^3$; local functions are bounded Borel functions of finitely many bonds. A \emph{cube of $h$ layers} is $Q(x_0,h)=\{x_0+(j_1,j_2,j_3):1\le j_\mu\le h\}$; a function of the bonds with both ends in $Q$ is a function on the configurations of every box with sides $\ge h+1$ through the projection of $\mathbb Z^3$, which is injective on $Q$, and box laws are laws on $\Omega$ by periodic extension. A \emph{Gibbs state} is a probability measure on $\Omega$ under which, given the bonds outside a finite set $\mathfrak A$, those of $\mathfrak A$ have the density proportional to $e^{-\beta\sum a_p}$, the sum over the plaquettes that meet $\mathfrak A$; it is \emph{periodic} if it is invariant under a subgroup of finite index of the translations.

The \emph{symmetries of the lattice} are the maps $x\mapsto Ox+v$, $O$ a signed permutation and $v\in\mathbb Z^3$. A \emph{lattice mirror} is a plane $\{x_\mu=c\}$, $c\in\frac12\mathbb Z$ (a \emph{site mirror} if $c\in\mathbb Z$, a \emph{bond mirror} otherwise), or a plane $\{x_\mu\pm x_\nu=c\}$, $\mu\ne\nu$, $c\in\mathbb Z$ (a \emph{diagonal mirror}); a state $s$ is \emph{reflection positive} at it if $s(\overline{\mathfrak rA}\,A)\ge0$ for every local $A$ of the bonds in one closed side of the mirror, $\mathfrak r$ the reflection in it (bonds in a site mirror or a diagonal mirror are admitted; at a bond mirror the bonds that cross it lie in neither closed side). For a closed lattice path $C$, $W_C=\frac12\operatorname{tr}U(C)$, $U(C)$ the ordered product along $C$; $W_C$ is real, $|W_C|\le1$, and it depends neither on the starting point nor on the sense; $|C|$ is the number of bonds of $C$, and $W_p$ is the loop along the boundary of one plaquette.

A \emph{rectangle} $R\times T$ of the $(\mu,\nu)$-plane runs $R$ steps along $e_\mu$, $T$ along $e_\nu$, and back.

\emph{Hypotheses.} The \emph{criterion at the side $l$} is the condition: $\beta>0$, $l\ge3$ is an integer of either parity, $0<x_0\le0.069$, and $\mathsf e(l)\le x_0$. Its numbers are $\theta=12.9\,x_0\le0.8901$, $\Lambda=\log\frac1\theta\ge0.11642$, $\varsigma=2\Lambda+1.67$, $\mathsf q=\varsigma/2\Lambda=1+0.835/\Lambda\le8.173$, and
\begin{equation}\label{eq:sigma}
\sigma^+=\frac{\varsigma}{4l^2}=\frac{\log\frac1\theta+0.835}{2l^2},\qquad\sigma^+l^2=\frac12\log\frac{e^{0.835}}{12.9x_0}\ge0.4757 .
\end{equation}
The \emph{cube condition at the side $l_0$} is $X(l_0)\le\frac1{20}$. For the physical scale let $L\ge2$ be an integer and $a_K=L^{-K}$. The \emph{uniform criterion} holds if there are $\mathsf b>0$, $l_*>0$, an integer $K_0$ and one $x_0\in(0,0.069]$ such that for every integer $K\ge K_0$ the number $l_*L^K$ is an integer $\ge3$ and the criterion holds at the side $l_*L^K$ and the coupling $L^K\mathsf b$ with this $x_0$. In the frame $L$ is the blocking factor, $a_K$ the spacing at the cutoff $K$, $\mathsf b=a_K\beta$ the label, and $\ell_0=2l_*$ the physical side of the cube of the criterion. In Lemma~\ref{lem:exc}, in the steps of proofs that use it, and in Proposition~\ref{prop:small}, $\mathsf L$ is the blocking factor, an odd integer $\ge13$; the $L$ of the uniform criterion is any integer $\ge2$, independent of it. We write $x/yz$ for $x/(yz)$.

\subsection{Principal results}

\begin{theorem}[One lattice]\label{thm:one}
Assume the criterion at the side $l$.
\begin{enumerate}
\item[(a)] \emph{The state.} For $n=2^kl$, $k\ge1$, and every bounded Borel function $A$ of the bonds of a cube of $h$ layers, $1\le h\le n/2-1$,
\[
|\langle A\rangle_{(n,n,n)}-\langle A\rangle_{(2n,2n,2n)}|\le1.29\,\sup|A|\,\theta^{n/2l}.
\]
So $\nu(A)=\lim_k\langle A\rangle_{(2^kl,2^kl,2^kl)}$ exists for every local $A$; and for $n=2^kl$, $k\ge1$, and $A$ a bounded Borel function of the bonds of a cube of $h$ layers with $1\le h\le n/2-1$, the $n$ and $A$ of the first bound, $|\nu(A)-\langle A\rangle_{(n,n,n)}|\le1.29\sup|A|\,\theta^{n/2l}/(1-\theta^{n/2l})$. The measure $\nu$ is a Gibbs state, invariant under the symmetries of the lattice, reflection positive at every lattice mirror and ergodic under the translations; it is the only periodic Gibbs state, and the laws of the periodic boxes converge weakly to $\nu$ as their three sides tend to infinity.
\item[(b)] \emph{Rectangles.} For integers $R,T\ge1$ and every rectangle $C$ of sides $R\times T$ in a coordinate plane,
\begin{equation}\label{eq:area}
0\le\nu(W_C)\le e^{-\sigma^+RT};
\end{equation}
on the cube of every integer side $n\ge2\mathsf ql$, for $R\le n/2$ and $T\le n-1$,
\[
0\le\langle W_C\rangle_{(n,n,n)}\le8^{T/n}\exp\Bigl[-\sigma^+RT+\frac{2\mathsf qRT}{ln}\Bigl(\frac12\log\frac{2\mathsf qn}l+0.86\Bigr)\Bigr];
\]
and, if $R,T\ge2$, also $\nu(W_C)\le(\frac{8\beta}{8\beta+3})^{2(R+T)-4}$, a bound that holds in every Gibbs state.
\item[(c)] \emph{Closed paths.} Let $C$ be a closed lattice path that runs a bond of the direction $\rho$, and $\rho'\ne\rho$; let $\mathsf A_{\rho'\rho}(C)$ be the number of unit squares of the $(\rho',\rho)$-plane about which the projection of $C$ to that plane winds an odd number of times, and $c_\rho(C)$ the number of changes from ``down'' to ``up'' in the cyclic sequence of the senses in which $C$ runs its bonds of the direction $\rho$. Then $|\nu(W_C)|\le2^{c_\rho(C)-1}e^{-\sigma^+\mathsf A_{\rho'\rho}(C)}$. More generally, for a gap $\tau$ between neighbouring layers of the direction $\rho$ let $\ell^{\rho'}_\tau(C)$ be the number of integers $\alpha$ such that an odd number of the bonds of the direction $\rho$ that $C$ runs in that gap have the $\rho'$-coordinate $\le\alpha$, so that $\sum_\tau\ell^{\rho'}_\tau(C)=\mathsf A_{\rho'\rho}(C)$; let $\mathsf A_\rho(C)=\sum_\tau\max_{\rho'\ne\rho}\ell^{\rho'}_\tau(C)$, the larger of the two counts taken gap by gap, so that $\mathsf A_\rho(C)\ge\max_{\rho'\ne\rho}\mathsf A_{\rho'\rho}(C)$; and let $D$ be a continuous unitary representation of $G$ on $\mathbb C^d$ with $D(-\mathbf 1)=-\mathbf 1$, and $W^D_C=d^{-1}\operatorname{tr}D(U(C))$.

Then $|\nu(W^D_C)|\le d^{\,c_\rho(C)-1}e^{-\sigma^+\mathsf A_\rho(C)}$; at $D(g)=g$, $d=2$, this is the first bound with $\mathsf A_\rho(C)$ in the place of $\mathsf A_{\rho'\rho}(C)$. In particular a closed path that in each such gap runs either no bond of the direction $\rho$ or exactly two, at the sites $o_\tau$, $r_\tau$ of the section, has $c_\rho(C)=1$, as its height in the direction $\rho$ has one maximum, and $|\nu(W_C)|\le\exp[-\sigma^+\sum_\tau|r_\tau-o_\tau|_\infty]$. So $|\nu(W_C)|\le e^{-\sigma^+\sum_j|r_j-o_j|}$ for every \emph{staircase loop}, a closed path in a coordinate plane that, in each gap between neighbouring layers of one direction $\rho$ of the plane, runs either no bond of the direction $\rho$ or exactly two, in opposite senses, at the $\rho'$-coordinates $o_j$ and $r_j$; for a simple loop $\sum_j|r_j-o_j|$ is the enclosed area.

On every periodic box, of any parity, the staircase loops of a $(1,3)$-plane described in Lemma~\ref{lem:two}(f), with $\rho=3$ and their bonds of the direction $3$ in at most $n_3-2$ gaps, have the bound stated there. For a simple closed curve in a coordinate plane $4c_\rho(C)\le|C|$, so that $|\nu(W_C)|\le\frac12\exp[-\sigma^+\mathsf A_{\rho'\rho}(C)+0.174\,|C|]$.
\item[(d)] \emph{Twisted cubes.} For every integer $n\ge2\mathsf ql$ and the seven $z\ne0$, on $(n,n,n)$:
\[
z_-\le\exp\Bigl[-\sigma^+n^2+\frac{2\mathsf qn}l\Bigl(\frac12\log\frac{2\mathsf qn}l+0.86\Bigr)\Bigr],\qquad0\le1-Z^z/Z\le6z_- .
\]
\item[(e)] \emph{Flux sectors and static charges.} For the same $n$, over the section $n\times n$ the top of every sector with electric flux lies at least $\sigma^+n-\frac{2\mathsf q}l(\frac12\log\frac{2\mathsf qn}l+0.86)-\frac{g_n}n$ above the vacuum, with $g_n=0.0017$ if $n=2^il$ and $g_n=0.105$ for every other $n$; and for $1\le R\le n/2$ the transfer operator $\hat{\mathsf T}_{o,r}$ with two static charges at the distance $R$ (Lemma~\ref{lem:two}) has
\[
\|\hat{\mathsf T}_{o,r}\|\le\lambda_0(n,n)\,\bigl[8(1+\mathsf r_n(n,n))\bigr]^{1/n}\exp\Bigl[-\sigma^+R+\frac{2\mathsf qR}{ln}\Bigl(\frac12\log\frac{2\mathsf qn}l+0.86\Bigr)\Bigr].
\]
The same bound holds for any two sites $o\ne r$ of the section, with $R=|r-o|_\infty$, the larger of the distances of their two coordinates on the circle of $n$ sites (Lemma~\ref{lem:two}(e)). So for any two sites $o\ne r$ of $\mathbb Z^2$, projected to the sections $n\times n$: $\limsup_{i\to\infty}\|\hat{\mathsf T}_{o,r}\|/\lambda_0(n,n)\le e^{-\sigma^+|r-o|_\infty}$ along $n=2^il$.
\item[(f)] \emph{Gap.} Over every section $N_1\times N_2$ with integers $N_1,N_2\ge2l$ the vacuum is simple, its eigenvector lies in the sector $00$, and every $\lambda_i$, $i\ge1$, the tops of the sectors with electric flux among them, lies at least $\frac\Lambda l+\frac{1.021}{n_0}$ above it, $n_0$ the largest $2^kl\le\min\{N_1,N_2\}$. Over $N_2\times n$, $n=2^kl$, $k\ge1$, and every integer $N_2\ge n$ the same holds with $\frac\Lambda l+\frac{3.5065}n$, and over $n\times n$ with $\frac\Lambda l+\frac{6.42}n$. For integers $N_1,N_2\ge n=2^kl$, $k\ge1$, with $0.01655\,N_1N_2\,\theta^{n/l}\le n^2\log2$: $\ \mathsf r_n(N_1,N_2)\le\frac{N_1N_2}{41.8\,n^2}\,\theta^{n/l}$.
\item[(g)] \emph{Flux lines.} Over a section $(N_2,N)$, with the first direction as time, $\Lambda_-=\max\{\Lambda_{01},\Lambda_{11}\}$ is the top of the sectors that are odd under the flip of the side $N$, the map $\iota_2$ of Lemma~\ref{lem:tr} that changes the sign of the bonds along the side $N$ starting on one line along the side $N_2$; these sectors carry a flux line of the length $N$.
\begin{enumerate}
\item[(g1)] Let $N=2wl$, $w\ge2$ a power of two, $\chi=0.03\,\theta^{N/l}$, $\hat N=\lceil2\sigma^+N^2/\log\frac1\chi\rceil$ and $\mathsf g=\log\hat N+1.71$. Then $4\le\hat N\le\mathsf qN/l+1$. For every integer $N_2\ge\hat NN/2$, $\Lambda_-$ lies at least $\sigma^+N-(2\log\hat N+1.86)/N$ above the vacuum, and the lower limit of $\log(\lambda_0/\Lambda_-)$ as $N_2\to\infty$ is at least $\sigma^+N-\mathsf g/N$. For every integer $N_2$ with $\max\{N,2\mathsf ql\}\le N_2<\hat NN/2$, $\Lambda_-$ lies at least $\sigma^+N-\frac{2\mathsf q}l(\log\frac{2\mathsf qN}l+0.86)$ above the vacuum. For integers $N'$ with $N\le N'<2N$ and $N_2\ge\hat NN/2$, over $(N_2,N')$ the top of the sectors odd under the flip of the side $N'$ lies at least $\sigma^+N'-(4\log\hat N+3.72)/N-(2N_2/N^2)\,0.01655\,\theta^{N/l}$ above the vacuum; the last term is at most $0.0331/N$ for $N_2\le N\theta^{-N/l}$.
\item[(g2)] Let $N\ge4l$ be any integer, $n_0=2wl$, $w\ge2$ a power of two, the dyadic multiple of $l$ with $n_0\le N<2n_0$, $\kappa=N/n_0$, $\tau=\theta^{n_0/l}$, $\chi_0=0.3602\,\tau$, $\mathsf E_0=\sigma^+n_0N$, $\hat N'$ the least even integer $\ge2\mathsf E_0/\log\frac1{\chi_0}$, $\mathsf h=0.01655\,\hat N'\tau$ and $\mathsf g'=\kappa\log\hat N'+3.22+(3\kappa+0.68)\mathsf h$. Then $6\le\hat N'<\mathsf qN/l+2$; for every integer $N_2\ge\hat N'n_0/2$, $\Lambda_-$ lies at least $\sigma^+N-\kappa\mathsf g'/N\ge\sigma^+N-(4\log\hat N'+18.2)/N$ above the vacuum; and for every integer $N_2$ with $\max\{N,2\mathsf ql\}\le N_2<\hat N'n_0/2$ it lies at least $\sigma^+N-\frac{2\mathsf q}l(\log\frac{2\mathsf qN}l+0.86)$ above it. Consequently, with
\[
C(\theta)=\max\Bigl\{\frac{\mathsf q}2+\frac14,\ 4\log\Bigl(\mathsf q+\frac12\Bigr)+18.2\Bigr\}\le27,
\]
for all integers $N\ge4l$ and $N_2\ge C(\theta)N^2/l$ the top $\Lambda_-$ lies at least $\sigma^+N-C(\theta)(1+\log\frac Nl)/N$ above the vacuum. If moreover $w\ge4$ and $\tau\le0.087$, then with $\hat N''=\lceil4\sigma^+n_0^2/\log\frac1{0.2728\,\tau}\rceil$ it lies at least $\sigma^+N-(4\log\hat N''+10.8)/n_0$ above the vacuum for every integer $N_2\ge\hat N''n_0/2$.
\item[(g3)] For every integer $N\ge4l$, every integer $N_2\ge\hat N'n_0/2$ and every rectangle $C$ of sides $R\times T$, $T\le N-1$, in a plane spanned by the first and the third direction,
\[
\limsup_{L_1\to\infty}\langle W_C\rangle_{(L_1,N_2,N)}\le4^{T/N}\exp\Bigl[-\sigma^+RT+\frac{RT}{N^2}\bigl(4\log\hat N'+18.2\bigr)\Bigr].
\]
\end{enumerate}
\item[(h)] \emph{Boxes of every shape, cylinders and twisted boxes.} For every even $N\ge2l$, every bounded real Borel function $A$ of the bonds of a cube of $N/2$ layers and every periodic box with all sides $\ge N$,
\[
|\langle A\rangle_{\underline n}-\nu(A)|\le5.42\,\theta^{N/4l}\sup|A|;
\]
for the events of these bonds the bound is $2.71\,\theta^{N/4l}$. In particular $|\langle A\rangle_{(m,m,m)}-\nu(A)|\le5.42\,\theta^{\lfloor m/2\rfloor/2l}\sup|A|$ for every integer $m\ge2l$ and every such $A$ of a cube of $\lfloor m/2\rfloor$ layers. Over every section $\Sigma$ the laws of the boxes $(t,\Sigma)$ converge as $t\to\infty$, on the bounded Borel functions of finitely many layers, to the law of the infinite cylinder over $\Sigma$; for every even $N'\ge2l$, every $N\ge N'$ and every $A$ as above of a cube of $N'/2$ layers, the cylinder over $N\times N$ satisfies the bound with $N'$, so that these laws converge to $\nu$ as $N\to\infty$; and on the cylinder over $n\times n$, $n\ge2l$, every rectangle with the side $R\le n/2$ across the axis and the side $T$ along it has a mean of modulus at most $[(1+\mathsf r_n(n,n))\,z_-(n,n,n)]^{RT/n^2}$. The laws of twisted boxes whose three sides tend to infinity and whose columns recede, in the sense that every finite set of plaquettes meets the columns of only finitely many members, converge weakly to $\nu$, and so do the mixtures, with any weights, of finitely many such sequences; if a column passes through a fixed plaquette $p_0$ in every member, the law of $\operatorname{tr}\partial U(p_0)$ tends to the law of $-\operatorname{tr}\partial U(p_0)$ under $\nu$, which is another law.
\item[(i)] \emph{The pressure.} There is an open interval $\mathsf U$ containing $\beta$ on which the criterion holds at the side $2l$ with $x_0=0.069$; on $\mathsf U$ the pressure is continuously differentiable, with $P'(\beta')=-3\,\nu_{\beta'}(a_p)$, $\nu_{\beta'}$ the state of (a) at $\beta'$; and $\nu(a_p)\le1-\frac\beta{4(\beta+1)}<1$.
\end{enumerate}
All constants are absolute; $\sigma^+$ and $\mathsf q$ depend only on $l$ and $x_0$. Decimal upper and lower bounds used in the finite recursions are understood in their indicated directions. The interval bounds recorded in the proof of Proposition~\ref{prop:tub} verify the numerical part of the stronger tension estimate.
\end{theorem}

\begin{theorem}[The criterion and the cube condition]\label{thm:crit}
Let $\beta\ge0$ and let $l\ge3$ be an integer of either parity.
\begin{enumerate}
\item[(a)] $X(l)=\mathsf e(l)+2d_2(l)+4d_1(l)$, and $\log(1+\mathsf r_l(\Sigma_j))\le d_j(l)\le2\log(1+\mathsf r_l(\Sigma_j))$ for $\Sigma_1=(l,l)$, $\Sigma_2=(2l,l)$, $\Sigma_3=(2l,2l)$ and $(d_1,d_2,d_3)=(d_1,d_2,\mathsf e)$; so the three terms are $\ge0$, and the cube condition at $l_0$ is the inequality $\mathsf e(l_0)\le\frac1{20}-2d_2(l_0)-4d_1(l_0)$. If $\beta>0$, the cube condition at $l_0$ gives the criterion at the same side $l_0$, with every $x_0\in[X(l_0),\frac1{20}]$, $x_0>0$, and $\theta\le0.645$.
\item[(b)] Under the criterion at $l$: $\mathsf e(2^kl)\le\theta^{2^k}/12.9$ for all $k\ge0$, and $X(2^kl)\le0.10643\,\theta^{2^k}\le0.04194$ for all $k\ge3$; so the cube condition holds at every side $2^kl$ with $k\ge3$, with $k\ge2$ if $x_0\le0.0650$, and with $k\ge1$ if $x_0\le0.0527$.
\item[(c)] Under the criterion at $l$, along $n=2^kl$ the rates $\mathsf m_{\mathsf e}(n)=\frac1n\log\frac1{12.9\,\mathsf e(n)}$ $(k\ge0)$ and $\mathsf m_X(n)=\frac1n\log\frac1{18X(n)}$ $(k\ge3)$, with values in $(0,\infty]$, do not decrease; for $k\ge3$, $\mathsf m_{\mathsf e}(n/2)-\frac{0.6501}n\le\mathsf m_X(n)\le\mathsf m_{\mathsf e}(n)-\frac{0.3331}n$; so the two rates have one common limit, which is $\ge\Lambda/l$.
\item[(d)] If $0<\mathsf e(l)<0.0740538$ at one $l$, then $\mathsf e(2^kl)\to0$ and $X(2^kl)\to0$ as $k\to\infty$; so the criterion holds at some side $2^kl$ with any prescribed $x_0\in(0,0.069]$, and the cube condition holds at $2^kl$ for all large $k$.
\end{enumerate}
\end{theorem}

\begin{theorem}[Physical scale, uniformly in the cutoff]\label{thm:phys}
Assume the uniform criterion, and let $\nu_K$ be the state of Theorem~\ref{thm:one} at $l=l_*L^K$ and $\beta=L^K\mathsf b$, read on the lattice $a_K\mathbb Z^3$. Put
\[
\sigma_*=\frac{\frac12\log\frac1\theta+0.4175}{l_*^2}\ \ge\ \frac{0.4757}{l_*^2}=\frac{1.9028}{\ell_0^2}.
\]
For every $K\ge K_0$, with physical lengths in $a_K\mathbb N$ and energies $a_K^{-1}\log(\lambda_0/\cdot)$:
\begin{enumerate}
\item[(a)] $0\le\nu_K(W_C)\le e^{-\sigma_*RT}$ for every rectangle $C$ of a coordinate plane with physical sides $R,T$; and if $R,T\ge2a_K$, also $\nu_K(W_C)\le\exp[-\frac12\sigma_*RT-\frac{3(|C|_{\rm ph}-4a_K)}{2(8\mathsf b+3a_K)}]$, $|C|_{\rm ph}=2(R+T)$ the physical length of $C$;
\item[(b)] on the periodic cube of physical side $\ell\ge2\mathsf ql_*$, $0\le1-Z^z/Z\le6\exp[-\sigma_*\ell^2+\frac{2\mathsf q\ell}{l_*}(\frac12\log\frac{2\mathsf q\ell}{l_*}+0.86)]$ for the seven $z\ne0$, and every sector with electric flux over $\ell\times\ell$ lies at least $\sigma_*\ell-\frac{2\mathsf q}{l_*}(\frac12\log\frac{2\mathsf q\ell}{l_*}+0.86)-\frac{0.105}\ell$ above the vacuum, with $\frac{0.0017}\ell$ in the place of $\frac{0.105}\ell$ if $\ell/l_*$ is a power of two;
\item[(c)] over every section whose physical sides are $\ge2l_*$ the gap is at least $\frac1{l_*}\log\frac1\theta$;
\item[(d)] a flux line of physical length $\ell=2wl_*$, $w\ge2$ a power of two, over a section $\ell_2\times\ell$ with $\ell_2\ge\hat N\ell/2$ lies at least $\sigma_*\ell-(2\log(\mathsf q\ell/l_*+1)+1.86)/\ell$ above the vacuum, and one of every physical length $\ell\ge4l_*$ over $\ell_2\times\ell$ with $\ell_2\ge C(\theta)\ell^2/l_*$ lies at least $\sigma_*\ell-C(\theta)(1+\log\frac\ell{l_*})/\ell$ above it, with the number $C(\theta)\le27$ of Theorem~\ref{thm:one}(g2);
\item[(e)] the cube condition holds at the side $8l_*L^K$, with $18X\le1.9157\,\theta^8$; conversely, the cube condition at the sides $l_*L^K$, $\beta=L^K\mathsf b$, for all $K\ge K_0$ gives the uniform criterion with $x_0=\frac1{20}$.
\end{enumerate}
The numbers $\sigma_*$, $\mathsf q$ and $C(\theta)$ depend on $l_*$ and $x_0$ alone, and $\hat N$ on $\ell/l_*$ and $x_0$ alone; none of them depends on $K$, $\mathsf b$ or $L$. In the first bound of (a) the factor in front of $e^{-\sigma_*RT}$ is $1$ and there is no perimeter term.
\end{theorem}

\begin{corollary}[The static potential]\label{cor:pot}
Assume the criterion at the side $l$, and let $m(R,T)=\nu(W_C)$ for a rectangle $C$ of sides $R\times T$ in a coordinate plane, $m(R,0)=1$.
\begin{enumerate}
\item[(a)] $m(R,T)$ depends on $R,T$ alone, $m(R,T)=m(T,R)$, and $w^{RT}\le m(R,T)\le e^{-\sigma^+RT}$, where $w=m(1,1)\ge\beta/(4(\beta+1))$.
\item[(b)] For every $n\ge0$ and real $c_0,\dots,c_n$: $\sum_{h,h'}c_hc_{h'}\,m(R,h+h')\ge0$ and $\sum_{h,h'}c_hc_{h'}\,m(R,h+h'+1)\ge0$. Hence $m(R,T)^2\le m(R,T-1)\,m(R,T+1)$ for $T\ge1$, and there is exactly one Borel probability measure $\varrho_R$ on $[0,1]$ with $m(R,T)=\int\lambda^T\,d\varrho_R(\lambda)$ for all $T\ge0$.
\item[(c)] $d(R;T)=\log(m(R,T)/m(R,T+1))$ does not increase in $T\ge0$; the static potential
\[
V(R)=\lim_{T\to\infty}d(R;T)=\inf_Td(R;T)=\lim_{T\to\infty}-\frac1T\log m(R,T)
\]
exists in $[0,\infty)$, $\max\operatorname{supp}\varrho_R=e^{-V(R)}$, and $m(R,T_2)\le e^{-(T_2-T_1)V(R)}m(R,T_1)$ for $T_2\ge T_1\ge0$, $V(R)$ being the largest number with this property. For $0\le T_1<T_2\le T_3<T_4$: $[m(R,T_2)/m(R,T_1)]^{1/(T_2-T_1)}\le[m(R,T_4)/m(R,T_3)]^{1/(T_4-T_3)}$. For $T\ge1$, $m(R,T)=\int_{[V(R),\infty)}e^{-TE}d\mu_R(E)$, $\mu_R$ the image of the restriction of $\varrho_R$ to $(0,1]$ under $\lambda\mapsto-\log\lambda$, a measure of mass in $(0,1]$.
\item[(d)] $\sigma^+R\le V(R)\le R\log\frac{4(\beta+1)}\beta$.
\item[(e)] Under the uniform criterion, for every $K\ge K_0$ and physical $R$: the potential $V_K(R)=a_K^{-1}V(R/a_K)$ of $\nu_K$ in physical units is $\ge\sigma_*R$; and for every real $\kappa$ the means $M_K(R,T)=e^{\kappa|C|_{\rm ph}}\nu_K(W_C)$ satisfy $M_K(R,T_2)\le e^{-(T_2-T_1)(V_K(R)-2\kappa)}M_K(R,T_1)$ for physical $T_2\ge T_1\ge0$, while in the inequality of four rectangles of (c) the factors $e^{\kappa|C|_{\rm ph}}$ cancel. Along a sequence of cutoffs: let $K_j\to\infty$, $K_j\ge K_0$, let $\kappa_j$ be arbitrary real numbers, and for $R,T$ in $\bigcup_ia_i\mathbb N$ and all large $j$ put $M_j(R,T)=e^{2\kappa_j(R+T)}\nu_{K_j}(W_C)$, $C$ a rectangle of the physical sides $R\times T$, and $v_j(R)=V_{K_j}(R)-2\kappa_j$, so that $v_j(R)\ge\sigma_*R-2\kappa_j$. If $T_1<T_2$ and $M_j(R,T_i)\to\mathsf W(R,T_i)\in(0,\infty)$ for $i=1,2$, then
\[
\limsup_jv_j(R)\le\frac1{T_2-T_1}\log\frac{\mathsf W(R,T_1)}{\mathsf W(R,T_2)},
\]
and if $v_j(R)\ge B$ for infinitely many $j$, then $\mathsf W(R,T_2)\le e^{-(T_2-T_1)B}\,\mathsf W(R,T_1)$; in particular $\mathsf W(R,T_2)\le e^{-(T_2-T_1)(\sigma_*R-2\kappa)}\,\mathsf W(R,T_1)$ if $\kappa_j\le\kappa$ for infinitely many $j$. If $M_j(R,T)\to\mathsf W(R,T)\in(0,\infty)$ for all $R,T$, then $\mathsf W(R,T)=\mathsf W(T,R)$, $\mathsf W(R,T_2)^{T_3-T_1}\le\mathsf W(R,T_1)^{T_3-T_2}\,\mathsf W(R,T_3)^{T_2-T_1}$ for $T_1<T_2<T_3$, the inequality of four rectangles of (c) holds for $\mathsf W$, and $V_{\rm ren}(R)=\lim_{T\to\infty}-\frac1T\log\mathsf W(R,T)$ exists in $[-\infty,\infty)$ and is the infimum of the quotients $\frac1{T_2-T_1}\log\frac{\mathsf W(R,T_1)}{\mathsf W(R,T_2)}$, $T_1<T_2$. The convergence of the $M_j$ is a hypothesis of these statements.
\item[(f)] With no criterion, in every periodic Gibbs state $s$ at $\beta\ge0$, $s(W_C)\ge0$ for every rectangle $C$ of a coordinate plane with one even side and for every square.
\item[(g)] With no criterion: let $\beta>0$ and let $s$ be a periodic Gibbs state that is positive at the bond mirrors on the functions of two crossing bonds, in the sense that
\[
s\Bigl(\sum_{i,j,k,l}\overline{(\mathsf z_p)_{ji}}\,(\mathsf z_q)_{kl}\,F_{jk}\,\overline{\mathfrak rF_{il}}\Bigr)\ge0
\]
for every bond mirror, $\mathfrak r$ the reflection in it, $\mathsf z_p$, $\mathsf z_q$ the values of $U$ at the bonds from two sites $x_p$, $x_q$ next to the mirror, equal or not, to their mirror images, and all continuous functions $F_{jk}$, $j,k\in\{1,2\}$, of finitely many bonds with both ends in the closed side of the mirror that holds $x_p$ and $x_q$, as in Lemma~\ref{lem:mir}(c). Let $m_s(R,T)=s(W_C)$ for a rectangle $C$ of sides $R\times T$ in a coordinate plane, $m_s(R,0)=1$. Then $m_s(R,T)$ depends on $R,T$ alone, $m_s(R,T)=m_s(T,R)$, and $[\beta/4(\beta+1)]^{RT}\le m_s(R,T)\le1$ for all integers $R,T\ge1$; both Hankel forms of (b) hold for $m_s$, and with them the log-convexity in $T$ and exactly one Borel probability measure $\varrho^s_R$ on $[0,1]$ with $m_s(R,T)=\int\lambda^T\,d\varrho^s_R(\lambda)$; $d_s(R;T)=\log(m_s(R,T)/m_s(R,T+1))$ does not increase in $T$, $V_s(R)=\lim_Td_s(R;T)=\inf_Td_s(R;T)=\lim_T-\frac1T\log m_s(R,T)$ exists in $[0,R\log\frac{4(\beta+1)}\beta]$, $\max\operatorname{supp}\varrho^s_R=e^{-V_s(R)}$, $m_s(R,T_2)\le e^{-(T_2-T_1)V_s(R)}m_s(R,T_1)$ for $T_2\ge T_1\ge0$, and the inequality of four rectangles of (c) holds for $m_s$.

On site mirrors alone, with no hypothesis at the bond mirrors: if $m_s(R,2)>0$, as is the case for every even $R$, then $V_{\rm e}(R)=\lim_h-\frac1{2h}\log m_s(R,2h)$ exists in $[0,\infty)$, and $m_s(R,T_2)\le e^{-(T_2-T_1)V_{\rm e}(R)}m_s(R,T_1)$ for every even $T_1\ge0$ and every integer $T_2\ge T_1$. The state $\nu$ of Theorem~\ref{thm:one} is positive at the bond mirrors in this sense, by Lemma~\ref{lem:mir}(c) on the cubes of sides $2^kl$ and Theorem~\ref{thm:one}(a), as in Step~2 of the proof; so (g) contains (a)--(c), the upper bound $e^{-\sigma^+RT}$ excepted.\qedhere
\end{enumerate}
\end{corollary}

\begin{proposition}[Where the criterion holds]\label{prop:where}
\begin{enumerate}
\item[(a)] For every integer $l\ge3$, every $x_0\in(0,0.069]$ and every $\beta\in(0,x_0/48l^3]$ the criterion holds at $l$.
\item[(b)] Let $0<\beta<\frac16$ and $l\ge4$ even. Then $X(l)\le64\,l^2(6\beta)^{l/2}$; so the cube condition, and the criterion with every $x_0\in[X(l),\frac1{20}]$, $x_0>0$, hold at every even $l$ with $64\,l^2(6\beta)^{l/2}\le\frac1{20}$, for instance at all even $l\ge10$ if $\beta=0.01$, $l\ge24$ if $\beta=0.05$, $l\ge62$ if $\beta=0.1$, $l\ge360$ if $\beta=0.15$.
\item[(c)] For $0<\beta<\frac16$ there is exactly one Gibbs state, periodic or not.
\end{enumerate}
\end{proposition}

\begin{proposition}[What the criterion says]\label{prop:says}
Let $\beta\ge0$ and let $l\ge3$ be an integer.
\begin{enumerate}
\item[(a)] With $\mathsf r=\mathsf r_l(2l,2l)$: $\log(1+\mathsf r)\le\mathsf e(l)\le2\log(1+\mathsf r)$. So $\mathsf r\le0.0351$ implies $\mathsf e(l)\le0.069$, and $\mathsf e(l)\le0.069$ implies $\mathsf r\le0.0357638$, in which case every $\lambda_i$, $i\ge1$, over $2l\times2l$, the tops of the sectors with electric flux among them, lies at least $3.33/l$ above the vacuum. Besides, $\mathsf e(l)\ge\log(1+\sum_{e\ne00}(\Lambda_e/\lambda_0)^l)$ over $2l\times2l$; so $\mathsf e(l)\ge\log(1+3y)$ if each of the three sectors with electric flux has $(\Lambda_e/\lambda_0)^l\ge y$, a floor that tends to $\log4$ as the three sectors approach the vacuum.
\item[(b)] Let $r_1=Z^-/Z$ on the cube of side $2l$, and $\mathsf f(x)=2\log(\sqrt{1-x}+\sqrt{2x})$ for $0\le x\le1$. Then $0<r_1\le1$ and $\mathsf e(l)\ge\mathsf f(1-r_1)$. Hence $\mathsf e(l)\le0.069$ forces $r_1>0.99937$, and $1-r_1\ge0.0039$ forces $\mathsf e(l)\ge0.16567$.
\item[(c)] If $\mathsf e(l)\le0.069$, then $1-Z^-/Z<0.0039$ on the cube of every integer side $N\ge16l$, and $\mathsf e(l')<0.0234$ for every integer $l'\ge16l$. So, if $1-Z^-/Z\ge0.0039$ on one cube of side $N$, then $\mathsf e(l)>0.069$ for every integer $3\le l\le N/16$.
\item[(d)] Under the criterion at $l$, the loop along the boundary of one plaquette has $l^2\ge0.4757/\log(1/\nu(W_p))$.
\end{enumerate}
\end{proposition}

\begin{remark}[Free fields]\label{rem:free}
For comparison, take three free massless fields in three dimensions, one for each generator of $\mathrm{SU}(2)$, each with one mode for every momentum, and leave out the modes of momentum $0$. Over the section $2l\times2l$ the energies of the modes are $\pi|m|/l$, $m\in\mathbb Z^2\setminus\{0\}$, so that the excited levels at the time $l$ give
\[
1+\mathsf r=\prod_{m\in\mathbb Z^2\setminus\{0\}}(1-e^{-\pi|m|})^{-3},\qquad1+\mathsf r'=\prod_{m\in\mathbb Z^2\setminus\{0\}}(1-e^{-2\pi|m|})^{-3},
\]
numbers free of $l$, with $\log(1+\mathsf r)\approx0.7201$ and $\log(1+\mathsf r')\approx0.0242$. The number $D=2\log(1+\mathsf r)-\log(1+\mathsf r')$ that stands for $\mathsf e(l)$ (Lemma~\ref{lem:seq}) is approximately $1.416$ at every scale, more than twenty times the bound $0.069$ of the criterion. This is a computation for free fields.
\end{remark}
For clause~(d) below we assume the following additional bound, denoted by \textbf{(Toron floor)}. For every integer $N'\ge4$ and every $\beta'\ge1.1\cdot10^9N'^{16}$, the partition function of the periodic cube at coupling $\beta'$, with normalized product Haar measure, satisfies \[ Z(N',N',N')\ge\frac{10^{-6}N'^2}{\pi^2\beta'^2}\,\mathsf X_0,\qquad \mathsf X_0=(2\pi^2)^{-(2\mathsf V-2)}\left(\frac{2\pi}{\beta'}\right)^{3\mathsf V-3}(\mathfrak d')^{-3/2},\qquad\mathsf V=N'^3, \] where $\mathfrak d'=\prod_{k\ne0}\sum_{\mu=1}^3(2-2\cos k_\mu)$, with $k\in(\frac{2\pi}{N'}\mathbb Z/2\pi\mathbb Z)^3$. This bound is separate from the construction theorem.
\begin{proposition}[Failure on small tori, given Theorem~\YMMref{D}{8}]\label{prop:small}
Given Theorem~\YMMref{D}{8} among the construction results, there are numbers $t_\circ>0$ and $\beta_\flat$, formed from the constants of Theorem~\YMMref{D}{8} and of the setting (Th) of Section~\ref{sec:cutoff-calibration}, the fixed calibration $\mathsf t_\natural$ of Section~\ref{sec:cutoff-calibration} among them, with the following properties.
\begin{enumerate}
\item[(a)] For every real $t\in(0,t_\circ]$ there is an integer $m_t$ such that on the periodic cube of $2\mathsf L^j$ sites a side at the coupling $\beta=2\mathsf L^j/t$ the quotient $r_1=Z^-/Z$ satisfies $r_1\le0.99596$ for every $j\ge m_t+2$; there $\mathsf e(\mathsf L^j)\ge0.16567$, and $\mathsf e(l)>0.069$ for every integer $3\le l\le\mathsf L^j/8$.
\item[(b)] For the sides $\ell_0=2\mathsf L^{m_0}$ and every label $\mathsf b$ with $\ell_0/\mathsf b\le t_\circ$: $\mathsf e(\ell_0/2a_K)\ge0.16567$ at $\beta=\mathsf L^K\mathsf b$ for all large $K$; so the uniform criterion with $L=\mathsf L$ and $l_*=\mathsf L^{m_0}$ fails at this label for every $K_0$ and every $x_0\le0.069$.
\item[(c)] For every real $\beta\ge\beta_\flat$: $\mathsf e(l)>0.069$ for every integer $3\le l\le\frac{t_\circ}{16\mathsf L}\beta$. So the uniform criterion, for any integer $L\ge2$, at any label $\mathsf b>0$, forces $\ell_0>\frac{t_\circ}{8\mathsf L}\mathsf b$.
\item[(d)] Assume in addition \textbf{(Toron floor)}. There are numbers $t_\star>0$, $\mathsf C_\star\ge1$ and $\mathsf p_\star\ge1$, formed from the same constants, such that for every integer $j\ge1$ and every real $\beta>0$ with $t=2\mathsf L^j/\beta\le t_\star$, on the periodic cube of $N=2\mathsf L^j$ sites a side,
\[
\frac{\lambda_0(N,N)^N}{Z(N,N,N)}\le\mathsf C_\star\,t\,\Bigl(\log\frac1t\Bigr)^{\mathsf p_\star},\qquad\text{so that}\qquad\mathsf e(\mathsf L^j)\ge\log\frac1t-\mathsf p_\star\log\log\frac1t-\log\mathsf C_\star .
\]
Under \textbf{(Toron floor)}, for a side $\ell_0=2\mathsf L^{m_0}$, $m_0\in\mathbb Z$, and a label $\mathsf b$ with $\ell_0/\mathsf b\le t_\star$ this says that at every cutoff $K\ge1-m_0$, at $\beta=\mathsf L^K\mathsf b$ and with $a_K=\mathsf L^{-K}$, $\mathsf e(\ell_0/2a_K)\ge\log\frac{\mathsf b}{\ell_0}-\mathsf p_\star\log\log\frac{\mathsf b}{\ell_0}-\log\mathsf C_\star$. The proof of (d) rests, besides Theorem~\YMMref{D}{8}, on \textbf{(Toron floor)} and the two estimates \eqref{eq:wind} and \eqref{eq:lowcut} stated in the next subsection: through them it bounds $\lambda_0(N,N)^N/Z(N,N,N)$ from above, and the floor of $\mathsf e$ follows from that ceiling.
\end{enumerate}
\end{proposition}

\begin{proposition}[What the trace identities alone allow]\label{prop:sheets}
Call \emph{data} every family of numbers $Z^z(\underline n)>0$, $\underline n$ a box and $z\in\{0,1\}^3$, with the properties (iii), (T1), (T2) and (T3) of Lemma~\ref{lem:tr}, and form $F$, $\mathsf e$, $d_1$, $d_2$, $X$, $z_-$ from them as above. Let $0<q<1$.
\begin{enumerate}
\item[(a)] $Z^z(\underline n)=\prod_\mu(1+q^{n_\mu})^{|\underline n|/n_\mu}$, the same for all $z$, are data, with $d_1:d_2:\mathsf e=1:2:4$ and $X=3\mathsf e$ at every $l$.
\item[(b)] The \emph{sheets} $Z^z(\underline n)=\prod_\rho\bigl(1+(-1)^{z_{\mu\nu}}q^{n_\mu n_\nu}\bigr)^{n_\rho}$, $\{\mu,\nu,\rho\}=\{1,2,3\}$, are data. With $s=q^{l^2}$ they have $\mathsf e(l)=8l\log(1+s^2)-4l\log(1+s^4)$ and $X(l)=24l\log(1+s)-6l\log(1+s^4)$, and $X/\mathsf e\ge2.95/s$ for $s\le0.03$ (at $l=4$, $s=0.03$: $\mathsf e\approx0.028774$, $X\approx2.837626$). So for no number $C$ does $X(l)\le C\,\mathsf e(l)$ hold for all data: at the same side, the criterion does not give the cube condition from these properties alone.
\item[(c)] For the sheets, $\frac12q^{n^2}\le z_-(n,n,n)\le nq^{n^2}$, so that $-n^{-2}\log z_-(n,n,n)\to\log\frac1q$, while with $x_0=\mathsf e(l)$ one has $\sigma^+=\log\frac1q-\frac{\log(103.2e^{-0.835}l)}{2l^2}+o(1)$ and $\Lambda=2l^2\log\frac1q-\log(103.2\,l)+o(1)$ as $q\to0$. So an inequality $\liminf_n-n^{-2}\log z_-(n,n,n)\ge c'\Lambda/l^2-C(l)$ with a number $c'>\frac12$ is false for the sheets: from the properties of data alone, the coefficient $\frac12$ of $\Lambda/l^2$ in $\sigma^+$ cannot be raised.
\end{enumerate}
\end{proposition}

\begin{proposition}[An area law without the numerical recursion of Proposition~\ref{prop:tub}]\label{prop:flat}
Assume the criterion at the side $l$, and let $l_j=2^{j+1}l$. Then $l_j\varepsilon_{\rm v}(l_j)\le0.0165\cdot2^{-j}\,\theta^{2\cdot4^j}$ for all $j\ge0$; $z_-(n,n,n)\le\max\{1,0.0096\sqrt{n/l}\}^{n/l}\,\theta^{n^2/2l^2}$ for $n=2^il$, $i\ge1$; and $0\le\nu(W_C)\le\theta^{RT/2l^2}$ for every rectangle $C$ of sides $R\times T$ in a coordinate plane. The tension $\log\frac1\theta/2l^2\ge0.0582/l^2$ so obtained does not use the values $E_1,\dots,E_6$ of Proposition~\ref{prop:tub}; the difference $0.4175/l^2$ to $\sigma^+$ is what that proposition adds.
\end{proposition}

\subsection{Construction estimates}

We use: (i) the periodic lattice of a box, its configurations and Haar measures, and Wilson's action; (ii) for a direction taken as time, the layers and the functions $A_\Sigma$ and $A_\times(V,S,V')=\sum_b(1-\frac12\operatorname{tr}V_bS_{b_+}V_b'^{-1}S_{b_-}^{-1})$ (Section~\YMMref{85}{677}, Section~\YMMref{85}{677}): the map $U\mapsto(V^\tau,S^\tau)_{\tau\in\mathbb Z/t}$, with $V^\tau\in\mathcal X=G^{\mathsf B(\Sigma)}$ the bonds of the layer $\tau$ and $S^\tau\in G^\Sigma$ those from the layer $\tau$ to the layer $\tau+1$, carries $dU$ to the product of the Haar measures, $\sum_pa_p=\sum_\tau[A_\Sigma(V^\tau)+A_\times(V^\tau,S^\tau,V^{\tau+1})]$ with continuous summands in $[0,2]$, and $A_\times(V',S^{-1},V)=A_\times(V,S,V')$; (iii) the crossing operator (Section~\YMMref{85}{677}, Proposition~\YMMref{85.1}{678}): the kernel $\mathsf K(V',V)=\int e^{-\beta A_\times(V,S,V')}dS$ is symmetric, continuous and has values in $(0,1]$, and the operator $\mathsf K$ on $\mathcal H=L^2(\mathcal X,dV)$ is Hilbert--Schmidt and $\ge0$; (iv) for Lemma~\ref{lem:sev} and Proposition~\ref{prop:small}, H\"older's inequality for traces of products of an even number of Hilbert--Schmidt operators (Lemma~\YMMref{86.1}{680}) and $\|A\|_{2k}^{2k}=\operatorname{Tr}(A^*A)^k$ (the norm definition preceding Lemma~\YMMref{86.1}{680}); (v) for the uniqueness, the boxes and the pressure, the chessboard estimate on the tori of side $2\mathsf L^j$ (Proposition~\YMMref{86.3}{681}), its form for two tori on one board (Corollary~\YMMref{86.6}{683}), and its bound on every periodic Gibbs state (Proposition~\YMMref{91.9}{705}); (vi) the invariance of every periodic Gibbs state under the symmetries of the lattice, its reflection positivity at every lattice mirror, the diagonal ones included, and its decomposition into ergodic components that are periodic Gibbs states (Propositions~\YMMref{103.4}{758} and~\YMMref{106.2}{764}); (vii) for Proposition~\ref{prop:small} alone, Theorem~\YMMref{D}{8} and the statements among the construction results through which it bears on the block fields at a fixed depth on tori that are small against the label: the setting (Th), what it gives, the weight of the small-field chain and the chain with its bound on the discarded mass (Section~\ref{sec:cutoff-calibration}, Lemma~\YMMref{118.1}{812}, the definitions preceding Proposition~\YMMref{115.2}{806}, Proposition~\textup{\YMMref{115.2}{806}}); the fixed calibration and the inner balls (Definitions~\ref{sec:cutoff-calibration}, \YMMref{113.10}{796}, Lemma~\YMMref{113.10}{796}); the setting (E) and the laws of the block fields (Definitions~\ref{sec:cutoff-calibration}, \textup{\ref{IV-62-62}}); the numerical conditions at the chosen constants (Proposition~\YMMref{69.1}{595}); the constants and the thresholds (Section~\ref{sec:cutoff-calibration} and equation~\YMMref{(634)}{339}); the block-averaging map with its locality, gauge covariance, measurability and covariance under the lattice isometries (equations~\YMMref{(37)--(39)}{20}, Proposition~\YMMref{4.5}{20}, \YMMref{4.5}{20}, \ref{lem:setup-lattice-covariance}); the bound on the holonomies along rearranged paths (Lemma~\YMMref{4.1}{17}); and, for its clause (d), the averages on small fields (Lemma~\YMMref{4.4}{19}, Lemma~\ref{lem:background-iterated-average}) and the comparison of two consecutive cutoffs (Section~\ref{sec:cutoff-calibration}, Section~\YMMref{116}{807}, Propositions~\YMMref{116.1}{808} and~\YMMref{116.3}{809} with their proofs, and Section~\ref{sec:block-field-laws}), on which Corollary~\textup{\ref{IV-62-66}} rests, read from one cutoff to the next and summed into the estimate \eqref{eq:lowcut} below, together with the signs of the constants of that comparison, among them $\kappa_{\rm b}>0$ and $c'>0$ (Proposition~\YMMref{114.4}{803} and Section~\ref{sec:cutoff-calibration}).

From the textbooks we use the spectral theorem for compact self-adjoint operators, Jentzsch's theorem, the min--max principle, the inequalities of Young, Jensen and Schwarz, Kolmogorov's extension theorem, the weak compactness of the probability measures on a compact metrisable space, and Hausdorff's moment theorem.

Clause (d) of Proposition~\ref{prop:small} rests, besides Theorem~\YMMref{D}{8} and the statements of (vii), on two estimates of its own, which its proof takes up (Lemma~\ref{lem:anti}, and Lemma~\ref{lem:cut} with the chain from a low cutoff). For a configuration $U$ of a periodic cube and a site $\xi$ let $\mathsf a(U,\xi)\in\mathbb R^5$ be the five invariants of the pair of the holonomies of $U$ along the closed straight paths through $\xi$ in the directions $1$ and $2$: their two class angles, $\alpha(g)=\arccos\frac12\operatorname{tr}g\in[0,\pi]$ for $g\in G$, and the three entries of $(\mathsf v^{\rm t}\mathsf v)^{1/2}$, $\mathsf v$ the real $3\times2$ matrix of their vector parts, numbers that do not change under a simultaneous conjugation of the pair. The first estimate concerns one lattice and assumes \textbf{(Toron floor)}: there are absolute numbers $C_{\rm a}$, $c_{\rm a}$, $c'_{\rm a}$, $c''_{\rm a}$ such that on the periodic cube of $N'\ge4$ sites a side, at every coupling $\beta'\ge c'_{\rm a}N'^{c''_{\rm a}}$, for every closed cube $\mathsf Q$ of the half side $w'$ in $\mathbb R^5$ with $(N'^5/2\beta')^{1/2}\le w'\le\frac1{20}$ and every $\omega>0$ below a fixed negative power of $N'$,
\begin{equation}\label{eq:wind}
\mu_{N',\beta'}\Bigl(\mathsf a(U,\xi)\in\mathsf Q,\ \max_p|\partial U(p)-\mathbf 1|<\omega\Bigr)\le C_{\rm a}N'^{c_{\rm a}}\,\beta'^{1/2}\,w'^3 ,
\end{equation}
$\mu_{N',\beta'}$ the law of that cube. Its proof combines a Gaussian ceiling with \textbf{(Toron floor)} for Wilson's measure on the cube in the gauge of a maximal tree, with constants that are powers of $N'$: the ceiling integrates out all bonds but the three winding holonomies of $\xi$ against the weight $\exp[-2\beta'N'^{-5}\sum_{\mu<\nu}|\mathsf v_\mu\times\mathsf v_\nu|^2]$, $\mathsf v_\mu$ the vector part of the holonomy in the direction $\mu$; the assumed floor of $Z(N',N',N')$ supplies the normalizing denominator; and the integral of that weight over the triples whose first two members have their five invariants in $\mathsf Q$ is at most a fixed multiple of $w'^3(\beta'/N'^5)^{-3/2}$. The second estimate is drawn from the lemmas of (vii) on two consecutive cutoffs.

Let $\mathsf b=\ell/t$ be a label and $\ell=2\mathsf L^m$ a side at which the setting (E) of Section~\ref{sec:cutoff-calibration} holds for it, with $\ell$ at most a fixed power of $y=\log\mathsf b$, and let $\mu^K_0$ be the law of the block field of the depth $0$ at the cutoff $K$ (Definition~\textup{\ref{IV-62-62}}(2)). For every $y$ above a number formed from the constants there is a lowest cutoff $K_1$, at which the bare cube has at most a fixed power of $y$ sites a side, such that for every integer $K\ge K_1$ and every Borel set $E$ of block fields of the depth $0$ that the gauge transformations carry onto itself (Definition~\textup{\ref{IV-62-62}}(3))
\begin{equation}\label{eq:lowcut}
\mu^K_0(E)\le e^{\mathsf S}\bigl(\mu^{K_1}_0(E)+\mathsf R_\Sigma\bigr),
\end{equation}
with $\mathsf S$ below an absolute number and $\mathsf R_\Sigma\le e^{-3y^2}$. Corollary~\textup{\ref{IV-62-66}} compares two consecutive cutoffs additively; here, from one cutoff to the next, the errors $\epsilon_{\rm fg}$ and $\epsilon_{\rm cl}$ of Section~\ref{sec:cutoff-calibration}, which decrease like powers of $\mathsf L^{-K}$ up to factors that are powers of $K$, $y$ and $\ell$, enter as factors on the probability, and the errors $\epsilon_{\rm sh}$ and $\mathsf R_{t_1}$, which are exponentially small in a power of $y$, as additive remainders and through the factors $(1-\epsilon_{\rm sh})^{-1}$ and $(1-\mathsf R_{t_1})^{-1}$ (Lemma~\ref{lem:cut}); $\mathsf S$ and $\mathsf R_\Sigma$ are the sums over $K\ge K_1$ of the exponents of the factors and of the remainders. So an event whose probability at the cutoff $K_1$ is at most of the order of $t$ has a probability at most of the order of $t$ at every cutoff $K\ge K_1$, since $y\ge\log\frac1t$ makes $e^{-3y^2}$ small against $t$.

Beyond (i)--(vii) the proofs use no result among the construction results; no statement of the section other than Proposition~\ref{prop:small} rests on Theorem~\YMMref{D}{8} or on Theorems~\YMMref{AJ}{673} and~\YMMref{AD}{607}, and none other than its clause (d) on \textbf{(Toron floor)} and the estimates \eqref{eq:wind} and \eqref{eq:lowcut}.

\subsection{Estimates and proofs}

\begin{lemma}[Two identities for the crossing weight]\label{lem:wt}
Let $(V^S)_b=S_{b_-}V_bS_{b_+}^{-1}$ and $\mathsf K_0(V,W)=\prod_b\exp[-\beta(1-\frac12\operatorname{tr}V_bW_b^{-1})]$ for $V,W\in\mathcal X$ and $S\in G^\Sigma$.
\begin{enumerate}
\item[(a)] $e^{-\beta A_\times(V,S,V')}=\mathsf K_0(V^g,V'^{gS})$ for all $V,V'\in\mathcal X$ and all $g,S\in G^\Sigma$.
\item[(b)] $\mathsf K_0(V,W)=\sum_\alpha c_\alpha\phi_\alpha(V)\overline{\phi_\alpha(W)}$, with $\alpha$ in a countable set, $c_\alpha\ge0$, $\sum_\alpha c_\alpha=e^{\beta\#\mathsf B(\Sigma)}$, $\phi_\alpha$ continuous and $|\phi_\alpha|\le1$.
\end{enumerate}
\end{lemma}

\begin{proof}[Proof sketch]
(a) $V_b((V'^S)_b)^{-1}=V_bS_{b_+}V_b'^{-1}S_{b_-}^{-1}$ is the matrix of the bond $b$ in $A_\times$, so $e^{-\beta A_\times(V,S,V')}=\mathsf K_0(V,V'^S)$; $(V^g)_b((W^g)_b)^{-1}=g_{b_-}V_bW_b^{-1}g_{b_-}^{-1}$, so $\mathsf K_0(V^g,W^g)=\mathsf K_0(V,W)$ by the cyclicity of the trace; and $(V^S)^g=V^{gS}$. (b) For $v,w\in G$, $\operatorname{tr}(vw^{-1})=\sum_{i,j}v_{ij}\overline{w_{ij}}$, with entries of modulus $\le1$. Expand each factor $e^{-\beta}\exp[\frac\beta2\operatorname{tr}V_bW_b^{-1}]$ in powers, write out the powers and multiply over $b$: $\alpha=(n_b,(i_{b,s},j_{b,s})_{s\le n_b})_b$, $c_\alpha=e^{-\beta\#\mathsf B(\Sigma)}\prod_b(\beta/2)^{n_b}/n_b!$, $\phi_\alpha(V)=\prod_b\prod_s(V_b)_{i_{b,s}j_{b,s}}$; as there are $4^{n_b}$ index families for each $n_b$, $\sum_\alpha c_\alpha=e^{-\beta\#\mathsf B(\Sigma)}\prod_be^{2\beta}$, and this bounds the series of the moduli.
\end{proof}

\begin{lemma}[Transfer operator and flux sectors]\label{lem:tr}
Fix a box and a direction $\mu_0$ as time, with the other directions $\nu_1<\nu_2$. For $i=1,2$ let $\iota_i$ change the sign of $V_b$ on the bonds of the direction $\nu_i$ that start on a fixed line of the other direction of the section, $(\mathsf J_i\psi)(V)=\psi(\iota_iV)$, $\mathsf J^k=\mathsf J_1^{k_1}\mathsf J_2^{k_2}$ and $P_e=\frac14(1+(-1)^{e_1}\mathsf J_1)(1+(-1)^{e_2}\mathsf J_2)$ for $e\in\{0,1\}^2$. Let $A^{(0)}_\Sigma=A_\Sigma$, let $A^{(1)}_\Sigma$ be $A_\Sigma$ with the summand $1-\frac12\operatorname{tr}$ of one plaquette replaced by $1+\frac12\operatorname{tr}$, $M_m$ the multiplication by $e^{-\beta A^{(m)}_\Sigma/2}$ and $\mathsf T_m=M_m\mathsf KM_m$; write $z\leftrightarrow(k;m)$ with $k=(z_{\mu_0\nu_1},z_{\mu_0\nu_2})$, each pair written in its order, and $m=z_{\nu_1\nu_2}$. Then:
\begin{enumerate}
\item[(i)] $\mathsf T_m\ge0$ is Hilbert--Schmidt with a continuous symmetric kernel with values in $(0,1]$; $\mathsf J_1$, $\mathsf J_2$ are commuting self-adjoint unitary involutions that commute with $\mathsf T_m$; none of these objects depends on $t$.
\item[(ii)] $Z^{(k;m)}=\operatorname{Tr}(\mathsf T_m^t\mathsf J^k)$ for every integer $t\ge3$.
\item[(iii)] $Z^z$ does not depend on the places of the columns, and $Z^{\pi z}(\pi\underline n)=Z^z(\underline n)$ for every permutation $\pi$ of the directions.
\end{enumerate}
Consequently: (T1) $Z$ is symmetric in the sides; (T2) $Z(t,\Sigma)=\sum_i\lambda_i^t$ for $t\ge3$, the $\lambda_i\ge0$ being the eigenvalues of $\mathsf T_0$, $\sum_i\lambda_i^2<\infty$, and $\lambda_0$ is simple with an eigenvector in the range of $P_{00}$; (T3) with the eigenvalues $\mu_i(e,m)\ge0$ of $\mathsf T_m$ on the range of $P_e$, which depend neither on $t$ nor on $k$,
\[
Z^{(k;m)}=\sum_e(-1)^{e\cdot k}W_t(e;m),\qquad W_t(e;m)=\operatorname{Tr}(\mathsf T_m^tP_e)=\sum_i\mu_i(e,m)^t\ge0;
\]
$\Lambda_e=\sup_i\mu_i(e,0)$ is the top of the sector $e$. Moreover:
\begin{enumerate}
\item[(a)] $W_t(00;0)\ge\lambda_0^t$;
\item[(b)] $0\le1-Z^{(k;0)}/Z\le2\mathsf r_t/(1+\mathsf r_t)$, so that $z_-\le\mathsf r_t/(1+\mathsf r_t)$ when the time is the direction $1$ or $3$;
\item[(c)] on a cube $Z^z/Z$ depends on $|z|$ alone, $r_2\le r_1$, $r_3\le r_1$, and $0\le1-Z^z/Z\le|z|(1-r_1)\le6z_-$; with any direction as time, the weights $w_e=W_n(e;0)/Z$ of the cube of side $n$ are $w_{00}=\frac14(1+2r_1+r_2)$, $w_{10}=w_{01}=\frac14(1-r_2)$, $w_{11}=\frac14(1-2r_1+r_2)$, and $1-r_1\le\frac43\,\mathsf r_n(n,n)/(1+\mathsf r_n(n,n))$;
\item[(d)] $a_-=\sum_i\mu_i^{n_1}=\sum_i\mu_i'^{\,n_3}$ and $Z^-=\sum_i\mu_i''^{\,n_2}$, with numbers $\mu_i,\mu_i',\mu_i''\ge0$ that do not depend on $n_1$, on $n_3$ and on $n_2$ respectively, and $\sup_i\mu''_i=\lambda_-(n_1,n_3)\le\lambda_0(n_1,n_3)$.
\end{enumerate}
\end{lemma}

\begin{proof}[Proof sketch]
Every plaquette is a summand of $A_\Sigma(V^\tau)$ or of $A_\times(V^\tau,S^\tau,V^{\tau+1})$. Put the columns at the origin of the section and at the crossing $\tau=0$. A column of the pair $\nu_1\nu_2$ is one summand in every layer and turns $A_\Sigma$ into $A^{(1)}_\Sigma$; a column of the pair $\mu_0\nu_i$ is the set of the crossing summands of the bonds flipped by $\iota_i$, and since $1-\frac12\operatorname{tr}(-g)=1+\frac12\operatorname{tr}g$, reversing it is the substitution $V^0\mapsto\iota_iV^0$ in that crossing. Integrating each $S^\tau$ first,
\[
Z^{(k;m)}=\int\mathsf T_m(V^1,\iota^kV^0)\prod_{\tau\ne0}\mathsf T_m(V^{\tau+1},V^\tau)\prod_\tau dV^\tau=\operatorname{Tr}(\mathsf T_m^t\mathsf J^k),
\]
the trace of a chain of $t\ge3$ bounded kernels: (ii). The functions $A^{(m)}_\Sigma$ and $A_\times$ are invariant under $\iota_i$, because the two bonds of the direction $\nu_i$ of a plaquette are both flipped or neither and $-\mathbf 1$ is central; so $\mathsf J_i$ commutes with $\mathsf T_m$, and $\langle\psi,\mathsf T_m\psi\rangle=\langle M_m\psi,\mathsf KM_m\psi\rangle\ge0$: (i). (iii): changing $U(b)$ into $-U(b)$ on a set $\mathfrak F$ of bonds keeps $dU$ and multiplies $\operatorname{tr}\partial U(p)$ by $(-1)^{\#(p\cap\mathfrak F)}$; for $\mathfrak F$ the bonds of the direction $\nu$ that start on the line of the direction $\rho$ through $x+e_\mu$ this sign is $-1$ exactly on the columns of the pair $\mu\nu$ through $x$ and through $x+e_\mu$ (the sides are $\ge3$), so a column moves by one step; a permutation of the axes renumbers the bonds and carries a plaquette to a plaquette or to its inverse, which has the same trace. (T2) and (T3) follow from the spectral theorem on the four ranges of the $P_e$.

By Jentzsch's theorem $\lambda_0$ is simple with an eigenvector $\varphi>0$, and $P_{00}\varphi=\frac14\sum_k\varphi\circ\iota^k\ge\frac14\varphi\ne0$: (a). (b): $Z-Z^{(k;0)}=2\sum_{e\cdot k\ \rm odd}W_t(e;0)\le2(Z-W_t(00;0))\le2(Z-\lambda_0^t)$. (c): on a cube (iii) makes $Z^z$ depend on $|z|$. With two twisted pairs, the time a direction of one of them but not of the other, $(k;m)=(10;1)$ and $r_2=\sum_e(-1)^{e_1}W(e;1)/Z\le\sum_eW(e;1)/Z=r^{(00;1)}=r_1$; in the same way $r_3=r^{(11;1)}\le r_1$; and with the time common to two twisted pairs, $(k;m)=(11;0)$, $4W(11;0)/Z=1-2r_1+r_2\ge0$ gives $1-r_2\le2(1-r_1)$. With all three pairs twisted, $(k;m)=(11;1)$ for every direction as time, and (T3) at $m=1$ gives $4W(11;1)/Z=r^{(00;1)}-r^{(10;1)}-r^{(01;1)}+r^{(11;1)}=r_1-2r_2+r_3\ge0$; added to $1-2r_1+r_2\ge0$ this is $1-r_3\le(1-r_1)+(1-r_2)$, and so $1-r_3\le3(1-r_1)$.

The formula for the weights is obtained by inverting $Z^{(k;0)}/Z=\sum_e(-1)^{e\cdot k}w_e$, with $Z^{(10;0)}=Z^{(01;0)}=Zr_1$ and $Z^{(11;0)}=Zr_2$. With $u=w_{10}$, $v=w_{11}$ one has $1-r_1=2(u+v)$, $1-r_2=4u$, so $r_2\le r_1$ means $v\le u$, which gives $2(u+v)\le\frac43(2u+v)=\frac43(1-w_{00})$, and $1-w_{00}\le1-\lambda_0^n/Z=\mathsf r_n/(1+\mathsf r_n)$ by (a). (d) is (T3) with the time $1$, $3$ and $2$; $\lambda_-\le\lambda_0$ because $Z^-\le Z$.
\end{proof}

\begin{lemma}[One positive sequence; the excess]\label{lem:seq}
Let $\kappa_0\ge\kappa_1\ge\dots\ge0$, $\kappa_0>0$, $\sum\kappa_i<\infty$; $p_M=\sum_i\kappa_i^M$ for real $M\ge1$, $\eta=\sum_{i\ge1}\kappa_i/\kappa_0$, $s_M=\sum_{i\ge1}(\kappa_i/\kappa_0)^M$, $D=2\log p_1-\log p_2$ and $D'=2\log p_2-\log p_4$. Then
\begin{enumerate}
\item[(a)] $\log(1+\eta)\le D\le2\log(1+\eta)$;
\item[(b)] if $D<\log2$, then $\eta\le\varpi(e^D-1)$;
\item[(c)] $D'\le2\log(1+s_2)\le2s_2\le2\eta^2$, so that $D'\le\omega(D)$ if $D<\log2$;
\item[(d)] $p_{MM'}\le p_M^{M'}$ for real $M,M'\ge1$;
\item[(e)] $s_M\le s_{M_0}^{M/M_0}$ for real $M\ge M_0\ge1$.
\end{enumerate}
For $\kappa_i=\lambda_i(\Sigma)^t$ one has $p_M=Z(Mt,\Sigma)$, $\eta=\mathsf r_t(\Sigma)$ and $s_M=\mathsf r_{Mt}(\Sigma)$ whenever $Mt$ is an integer $\ge3$. Hence:
\begin{enumerate}
\item[(f)] $F/|\underline n|$ and $\delta/|\underline n|$ do not increase in any side, $F/|\underline n|\to P$ as the three sides tend to infinity, $\log(1+\mathsf r_t(\Sigma))\le\delta(t,\Sigma)$, and $\log\lambda_0(N_1,N_2)\ge N_1N_2P$;
\item[(g)] for integers $A,B\ge N\ge3$, the Casimir term $\varepsilon(A,B)=\log\lambda_0(A,B)-ABP$ satisfies $0\le\varepsilon(A,B)\le(A+B)\,\delta_{\rm c}(N)/N^2$.
\end{enumerate}
\end{lemma}

\begin{proof}[Proof sketch]
$p_1=\kappa_0(1+\eta)$ and $p_M=\kappa_0^M(1+s_M)$, so $e^D=(1+\eta)^2/(1+s_2)$. For $u_i\ge0$ with the sum $u$ and real $m\ge1$, $\sum u_i^m\le u^m$; this gives (d), (e) and $s_2\le\min\{\eta,\eta^2\}$, hence (a), and (c) as $D'=2\log(1+s_2)-\log(1+s_4)$. (b): $e^D\ge(1+\eta)^2/(1+\eta^2)=1+2\eta/(1+\eta^2)$, and $\eta<1$ by (a); solve the quadratic inequality. (f): (d) with $M'=t'/t$ gives $Z(t',\Sigma)\le Z(t,\Sigma)^{t'/t}$ for $t'\ge t$, and (T1) the other sides; $\delta(t,\Sigma)=t(\log\lambda_0-N_1N_2P)+\log(1+\mathsf r_t)$, with $\log\lambda_0=\lim_sF(s,\Sigma)/s\ge N_1N_2P$. (g): put $\varkappa=\delta_{\rm c}(N)/N^2$. For $t,A,B\ge N$, $\delta(t,A,B)=t\,\varepsilon(A,B)+\log(1+\mathsf r_t(A,B))$, and by (f) $\log(1+\mathsf r_t(A,B))\le\log(1+\mathsf r_N(A,B))\le\delta(N,A,B)\le AB\varkappa$. By (f), $\varepsilon(A,B)/AB\le\delta(t,A,B)/tAB$ for every $t$, so $\varepsilon(A,B)/AB$ does not increase in $A$ or $B$ and tends to $0$ as $A,B\to\infty$. With $B$ as time, $\delta(s,t,B)=B\varepsilon(s,t)+\log(1+\mathsf r_B(s,t))\le B\varepsilon(s,t)+st\varkappa$; divide by $s$ and let $s\to\infty$: $\varepsilon(t,B)/t\le B\lim_s\varepsilon(s,t)/st+\varkappa$, and as $t\to\infty$, $\limsup_t\varepsilon(t,B)/t\le\varkappa$. With $A$ as time, $\delta(t,A,B)\le A\varepsilon(t,B)+tB\varkappa$; divide by $t$ and let $t\to\infty$: $\varepsilon(A,B)\le(A+B)\varkappa$.
\end{proof}

\begin{lemma}[Two static charges]\label{lem:two}
Fix a box, the time $3$, $o\in\Sigma$ and $r=o+Re_1$, $1\le R\le n_1-1$; let $\gamma$ be the straight path of $R$ bonds from $o$ to $r$ in the section, $V(\gamma)$ its ordered product, $\hat{\mathcal H}=L^2(\mathcal X;M_2(\mathbb C))$ with $\langle\Psi,\Psi'\rangle=\int\operatorname{tr}(\Psi^*\Psi')\,dV$,
\[
(\hat{\mathsf K}\Psi)(V)=\iint e^{-\beta A_\times(V,S,V')}\,S_o\Psi(V')S_r^{-1}\,dS\,dV',\qquad\hat{\mathsf T}=M_0\hat{\mathsf K}M_0,\qquad(\Upsilon\psi)(V)=\psi(V)V(\gamma);
\]
let $C$ be the rectangle $R\times T$, $1\le T\le n_3-1$, with the side $\gamma$ in the layer $0$, and $p_\xi$ half the trace of the line that winds once along the direction $3$ through $\xi\in\Sigma$. Then:
\begin{enumerate}
\item[(a)] $\hat{\mathsf K}$ and $\hat{\mathsf T}$ are Hilbert--Schmidt, self-adjoint and $\ge0$, they do not depend on $n_3$, and $\Upsilon^*\Upsilon=2$;
\item[(b)] $2Z\langle W_C\rangle=\operatorname{Tr}(\mathsf T_0^{n_3-T}\Upsilon^*\hat{\mathsf T}^T\Upsilon)$ and $\operatorname{Tr}\hat{\mathsf T}^{n_3}=4Z\langle p_op_r\rangle$;
\item[(c)] $0\le\langle W_C\rangle\le(4\langle p_op_r\rangle)^{T/n_3}$ and $\|\hat{\mathsf T}\|^{n_3}\le4Z\langle p_op_r\rangle$;
\item[(d)] $0\le\langle p_op_r\rangle\le z_+^{1-\varrho}z_-^\varrho+z_-^{1-\varrho}z_+^\varrho$, $\varrho=R/n_1$, and for $R\le n_1/2$
\[
0\le\langle W_C\rangle\le\bigl(8z_-^{R/n_1}\bigr)^{T/n_3},\qquad\frac{\|\hat{\mathsf T}\|}{\lambda_0(n_1,n_2)}\le\Bigl[8\bigl(1+\mathsf r_{n_3}(n_1,n_2)\bigr)z_-^{R/n_1}\Bigr]^{1/n_3}.
\]
\item[(e)] \emph{Any two sites.} For any two sites $o\ne r$ of the section, $\hat{\mathsf K}$ and $\hat{\mathsf T}=\hat{\mathsf T}_{o,r}$ are defined by the same formulas, (a) holds, $\operatorname{Tr}\hat{\mathsf T}^{n_3}=4Z\langle p_op_r\rangle$ and $\|\hat{\mathsf T}\|^{n_3}\le4Z\langle p_op_r\rangle$, and (d) holds for $\langle p_op_r\rangle$ and $\|\hat{\mathsf T}\|$ with $R$ the distance of $o_1$ and $r_1$ on the circle of $n_1$ sites, $1\le R\le n_1/2$; on a cube it holds as well with $R$ the distance of $o_2$ and $r_2$ on the circle of $n_2$ sites.
\item[(f)] \emph{Staircase loops in a box.} Let $1\le T\le n_3-2$, let $o$ be a site of the section, $\alpha_j$ and $R_j\ge1$ integers, $o_j=o+\alpha_je_1$ and $r_j=o+(\alpha_j+R_j)e_1$ for $0\le j<T$, and let $C$ be the closed path that runs in the layer $0$ straight along the line of the direction $1$ from $o_0$ to $r_0$, then for $j=0,\dots,T-1$ in turn one bond up at $r_j$ and, if $j<T-1$, straight from $r_j$ to $r_{j+1}$, in the layer $T$ straight from $r_{T-1}$ to $o_{T-1}$, and then for $j=T-1,\dots,0$ in turn one bond down at $o_j$ and, if $j>0$, straight from $o_j$ to $o_{j-1}$ (``straight'' from the site of the integer $\alpha$ to that of $\alpha'$: along the $|\alpha'-\alpha|$ bonds between them). Then, with the thermal numbers over $(n_1,n_2)$,
\[
\begin{multlined}
|\langle W_C\rangle|\le\bigl(1+\mathsf r_{n_3-T}\bigr)\prod_{j=0}^{T-1}\frac{\|\hat{\mathsf T}_{o_j,r_j}\|}{\lambda_0(n_1,n_2)},\qquad\text{and if all }R_j\le\tfrac{n_1}2:\quad|\langle W_C\rangle|\le\bigl(1 \\
+\mathsf r_{n_3-T}\bigr)\bigl[8(1+\mathsf r_{n_3})\bigr]^{T/n_3}z_-^{\sum_jR_j/n_1n_3}.
\end{multlined}
\]
\end{enumerate}
No side need be even.
\end{lemma}

\begin{proof}[Proof sketch]
The inequality that replaces H\"older's inequality for traces is this. Let $A,C\ge0$ be compact with eigenvalues $a_i$, $c_j$ and eigenbases $(e_i)$, $(f_j)$, $B$ bounded, $p,q\ge1$ integers and $s=p+q$. Then $\langle e_i,A^pB^*C^qBe_i\rangle=\sum_jw_{ij}a_i^pc_j^q$ with $w_{ij}=|\langle f_j,Be_i\rangle|^2\ge0$, the row and column sums of $(w_{ij})$ are $\le\|B\|^2$. A zero normaliser below makes the weighted product sum vanish, and infinite normalisers give only a trivial bound. For finite positive normalisers, Young's inequality $a^pc^q\le\mathsf a^{p/s}\mathsf c^{q/s}(\frac ps\frac{a^s}{\mathsf a}+\frac qs\frac{c^s}{\mathsf c})$, with $\mathsf a=\sum w_{ij}a_i^s$ and $\mathsf c=\sum w_{ij}c_j^s$, gives
\begin{equation}\label{eq:hold}
0\le\sum_i\langle e_i,A^pB^*C^qBe_i\rangle\le\|B\|^2\Bigl(\sum_ia_i^s\Bigr)^{p/s}\Bigl(\sum_jc_j^s\Bigr)^{q/s}.
\end{equation}
Its form with two operators bounds the modulus. For bounded $B_1$, $B_2$ the same eigenbases give $\sum_i\langle e_i,A^pB_1^*C^qB_2e_i\rangle=\sum_{i,j}a_i^pc_j^q\,\overline{\langle f_j,B_1e_i\rangle}\,\langle f_j,B_2e_i\rangle$, and $\sum_{i,j}a_i^pc_j^q|\langle f_j,B_ke_i\rangle|^2=\sum_i\langle e_i,A^pB_k^*C^qB_ke_i\rangle$ for $k=1,2$; Schwarz's inequality with the weights $a_i^pc_j^q\ge0$, and then \eqref{eq:hold} for $B_1$ and for $B_2$, give
\begin{equation}\label{eq:hold2}
\Bigl|\sum_i\langle e_i,A^pB_1^*C^qB_2e_i\rangle\Bigr|\le\|B_1\|\,\|B_2\|\Bigl(\sum_ia_i^s\Bigr)^{p/s}\Bigl(\sum_jc_j^s\Bigr)^{q/s}.
\end{equation}
(a) Symmetry follows from $A_\times(V',S^{-1},V)=A_\times(V,S,V')$. For positivity substitute $S=g^{-1}Q$, use $e^{-\beta A_\times(V,g^{-1}Q,V')}=\mathsf K_0(V^g,V'^Q)$ (Lemma~\ref{lem:wt}(a)), average over $g\in G^\Sigma$ (the left side does not depend on $g$) and expand $\mathsf K_0$ by Lemma~\ref{lem:wt}(b):
\[
\langle\Psi,\hat{\mathsf K}\Psi\rangle=\sum_\alpha c_\alpha\operatorname{tr}(X_\alpha^*X_\alpha)\ge0,\qquad X_\alpha=\iint\overline{\phi_\alpha(V^Q)}\,Q_o\Psi(V)Q_r^{-1}\,dQ\,dV.
\]
The argument uses of $o$ and $r$ only that they are sites of the section. (b) With $\mathsf P_\xi=S^0_\xi\cdots S^{T-1}_\xi$, $2W_C=\operatorname{tr}[V^0(\gamma)^*\mathsf P_oV^T(\gamma)\mathsf P_r^{-1}]$, and the layers write $2Z\langle W_C\rangle$ as the trace of the chain of the kernels of $\mathsf T_0^{n_3-T}$ and of $\Upsilon^*\hat{\mathsf T}^T\Upsilon$; at $T=n_3$ the trace over $M_2(\mathbb C)$ gives $\operatorname{tr}\mathsf P_o\operatorname{tr}\mathsf P_r^{-1}$. (c) is \eqref{eq:hold} with $A=\mathsf T_0$, $C=\hat{\mathsf T}$, $B=\Upsilon$, $p=n_3-T$, $q=T$ and $\sum a_i^{n_3}=Z$. (d) Take the time $1$. The two winding lines become the multiplication $\hat{\mathsf p}$ by one function of the variables of a layer, taken at the layers $o_1$ and $o_1+R$; each line holds exactly one bond flipped by $\iota_2$, so $\hat{\mathsf p}$ exchanges the ranges of $P_\pm=\frac12(1\pm\mathsf J_2)$, on which $\operatorname{Tr}(\mathsf T_0^{n_1}P_\pm)=a_\pm$ by Lemma~\ref{lem:tr}(ii). By \eqref{eq:hold}
\[
Z\langle p_op_r\rangle=\sum_\pm\operatorname{Tr}\bigl[(\mathsf T_0P_\pm)^{n_1-R}\,\hat{\mathsf p}\,(\mathsf T_0P_\mp)^R\,\hat{\mathsf p}\bigr]\le\sum_\pm a_\pm^{1-\varrho}a_\mp^\varrho ,
\]
and for $\varrho\le\frac12$ and $z_-\le z_+\le1$ each term is $\le z_-^\varrho Z$; finally $Z=\lambda_0^{n_3}(1+\mathsf r_{n_3})$. (e) The proof of (a) uses of $o$ and $r$ only that they are sites of the section, and the trace of $\hat{\mathsf T}^{n_3}$ in (b) and the second bound of (c) do not use $\gamma$; so, for any two sites, $\hat{\mathsf T}\ge0$, $4Z\langle p_op_r\rangle=\operatorname{Tr}\hat{\mathsf T}^{n_3}\ge0$ and $\|\hat{\mathsf T}\|^{n_3}\le4Z\langle p_op_r\rangle$: the sign of $\langle p_op_r\rangle$ comes from the time $3$. For the bound take the time $1$. The line through $\xi\in\Sigma$ is then the function $\mathsf p_{\xi_2}$ of the variables $V^{\xi_1}$ of the layer $\xi_1$, half the trace of the ordered product of the variables of its bonds of the direction $3$ with the second coordinate $\xi_2$; whatever $\xi_2$ is, the line holds exactly one bond flipped by $\iota_2$. So the multiplications $\hat{\mathsf p}_{o_2}$ and $\hat{\mathsf p}_{r_2}$, two different ones if $o_2\ne r_2$, are real, of norm $\le1$, and each exchanges the ranges of $P_\pm$. With $R$ the distance of $o_1$ and $r_1$ on the circle, $1\le R\le n_1/2$, the trace has $n_1-R$ factors $\mathsf T_0$ on one side of the two lines and $R$ on the other, and, as in (d), with the names $o$ and $r$ exchanged if need be,
\[
Z\langle p_op_r\rangle=\sum_\pm\operatorname{Tr}\bigl[(\mathsf T_0P_\pm)^{n_1-R}\,\hat{\mathsf p}_{o_2}\,(\mathsf T_0P_\mp)^R\,\hat{\mathsf p}_{r_2}\bigr];
\]
\eqref{eq:hold2} with $A=\mathsf T_0P_\pm$, $C=\mathsf T_0P_\mp$, $B_1=\hat{\mathsf p}_{o_2}$, $B_2=\hat{\mathsf p}_{r_2}$, $p=n_1-R$, $q=R$ bounds the term of $\pm$ in modulus by $a_\pm^{1-\varrho}a_\mp^\varrho$. With the sign from the time $3$ this is the first statement of (d) for $o$ and $r$, at every parity of $n_1$, $n_2$, $n_3$, and the bound of $\|\hat{\mathsf T}\|$ in (d) follows from it and from $\|\hat{\mathsf T}\|^{n_3}\le4Z\langle p_op_r\rangle$ as there. On a cube the time $2$ does the same with $o_2$, $r_2$, the twisted pair being then the pair $23$, whose quotient on the cube equals that of the pair $13$ (Lemma~\ref{lem:tr}(c)). (f) With the time $3$, $U(C)$ is the ordered product of the runs in the layers and of the bonds $S^j_{r_j}$, $(S^j_{o_j})^{-1}$; put $(\Upsilon_j\psi)(V)=\psi(V)V(o_j\to r_j)$, with $V(o_j\to r_j)$ the product along the straight run from $o_j$ to $r_j$, $(\Gamma_j\Psi)(V)=V(o_{j-1}\to o_j)\Psi(V)V(r_j\to r_{j-1})$, a unitary operator of $\hat{\mathcal H}$ that commutes with $M_0$, and $\mathsf Q=\Upsilon_0^*\hat{\mathsf T}_{o_0,r_0}\Gamma_1\hat{\mathsf T}_{o_1,r_1}\Gamma_2\cdots\Gamma_{T-1}\hat{\mathsf T}_{o_{T-1},r_{T-1}}\Upsilon_{T-1}$.

As in (b), $2Z\langle W_C\rangle=\operatorname{Tr}(\mathsf T_0^{n_3-T}\mathsf Q)$; in an eigenbasis of $\mathsf T_0$, $|\operatorname{Tr}(\mathsf T_0^{n_3-T}\mathsf Q)|\le\|\mathsf Q\|\lambda_0^{n_3-T}(1+\mathsf r_{n_3-T})$, with $n_3-T\ge2$; and $\|\mathsf Q\|\le2\prod_j\|\hat{\mathsf T}_{o_j,r_j}\|$, $Z\ge\lambda_0^{n_3}$. The second bound is (d) for each $j$.
\end{proof}

\begin{lemma}[Several static charges; closed paths]\label{lem:sev}
Fix a box with $n_1$ even and the time $3$.
\begin{enumerate}
\item[(a)] For a finite set $\mathcal V$ of runs, each with a site $u(v)\in\Sigma$ and a sense, let $D(S)=\bigotimes_vD_v(S_{u(v)})$ on $(\mathbb C^2)^{\otimes\mathcal V}$, with $D_v(g)=g$ for the sense up and $D_v(g)=\bar g$ for the sense down, $(\hat{\mathsf K}_D\Psi)(V)=\iint e^{-\beta A_\times(V,S,V')}D(S)\Psi(V')\,dS\,dV'$ and $\hat{\mathsf T}_D=M_0\hat{\mathsf K}_DM_0$. Then $\hat{\mathsf T}_D\ge0$ and $\|\hat{\mathsf T}_D\|^{n_3}\le\operatorname{Tr}\hat{\mathsf T}_D^{n_3}=2^{\#\mathcal V}Z\langle\prod_vp_{u(v)}\rangle$.
\item[(b)] Let $\#\mathcal V$ be even, $m_\tau$ the number of the $u(v)$ in the layer $\tau\in\{1,\dots,n_1\}$ of the direction $1$, and $\ell$ the number of $\tau$ with $m_1+\dots+m_\tau$ odd. If $\ell\le n_1/2$, then $|\langle\prod_vp_{u(v)}\rangle|\le2z_-^{\ell/n_1}$.
\item[(c)] Let $C$ be a closed lattice path in a cube of $h$ layers, $h+1\le\min_\mu n_\mu$, that runs bonds of the direction $3$ in $T$ gaps, $n_3\ge T+2$; in the gap $j$ let $\mathsf v_j$ be the number of these bonds and $\ell_j$ the number of integers $\alpha$ such that an odd number of them have first coordinate $\le\alpha$, and assume $\ell_j\le n_1/2$. Then, with $\mathsf r$ taken over the section $(n_1,n_2)$,
\[
|\langle W_C\rangle|\le2^{c_3(C)-1}\bigl(1+\mathsf r_{n_3-T}\bigr)\prod_j\Bigl[2^{\mathsf v_j+1}\bigl(1+\mathsf r_{n_3}\bigr)z_-^{\ell_j/n_1}\Bigr]^{1/n_3}.
\]
\item[(d)] Let $D$ be a continuous unitary representation of $G$ on $\mathbb C^d$ with $D(-\mathbf 1)=-\mathbf 1$, and $W^D_C=d^{-1}\operatorname{tr}D(U(C))$. Then (a)--(c) hold with $D_v=D$ for the sense up and $D_v=\overline D$ for the sense down, with $d^{-1}\operatorname{tr}D$ of the line that winds once along the direction $3$ through $\xi$ in the place of $p_\xi$, $d^{\#\mathcal V}$ in the place of $2^{\#\mathcal V}$, and $W^D_C$, $d^{\,c_3(C)-1}$, $2d^{\,\mathsf v_j}$ in the place of $W_C$, $2^{c_3(C)-1}$, $2^{\mathsf v_j+1}$.
\end{enumerate}
\end{lemma}

\begin{proof}[Proof sketch]
(a) As in Lemma~\ref{lem:two}(a), with $D(g^{-1}Q)=D(g)^{-1}D(Q)$; the trace of the $n_3$-th power gives $\prod_v\operatorname{tr}D_v(\mathsf P_{u(v)})$, and $\operatorname{tr}\bar g=\operatorname{tr}g$. (b) With the time $1$, $Z\langle\prod_vp_{u(v)}\rangle=\operatorname{Tr}\prod_\tau(\mathsf T_0\hat{\mathsf p}^\tau)$, $\hat{\mathsf p}^\tau$ the multiplication by the product of the line functions of the layer $\tau$. Each line holds one flipped bond, so $\hat{\mathsf p}^\tau$ maps the range of $P_s$ to that of $P_{s(-1)^{m_\tau}}$; splitting the trace over $s=\pm$, the factor of the layer $\tau$ becomes $\mathsf T_0P_{s_\tau}\hat{\mathsf p}^\tau$, with $s_\tau=-s$ for exactly $\ell$ layers. H\"older's inequality for the trace of the $n_1$ factors (Lemma~\YMMref{86.1}{680}), and $\|\mathsf C\hat{\mathsf p}\|_{n_1}^{n_1}=\operatorname{Tr}(\mathsf C\hat{\mathsf p}^2\mathsf C)^{n_1/2}\le\operatorname{Tr}\mathsf C^{n_1}=a_\pm$ for $\mathsf C=\mathsf T_0P_\pm$ (by min--max, as $0\le\mathsf C\hat{\mathsf p}^2\mathsf C\le\mathsf C^2$), bound the term of $s$ by $a_s^{1-\ell/n_1}a_{-s}^{\ell/n_1}$. (c) Cut $C$ into pieces in single layers and runs of the direction $3$: $2W_C=\Gamma_0D_0(S^0)\Gamma_1\cdots D_{T-1}(S^{T-1})\Gamma_T$, where $D_j$ is the tensor product over the runs of the gap $j$ and $\Gamma_j$, a function of $V^j$, is the tensor product over the pieces of the layer $j$ of unitary $2\times2$ matrices, of cups (linear forms on $\mathbb C^2\otimes\mathbb C^2$ with the entries of a matrix of $G$) and of caps (vectors with such entries).

Cups and caps have the norm $\sqrt2$ and occur at the changes from down to up and from up to down, $c_3(C)$ times each, so $\prod_j\sup_V\|\Gamma_j(V)\|=2^{c_3(C)}$. As in Lemma~\ref{lem:two}(b), $2Z\langle W_C\rangle=\operatorname{Tr}(\mathsf T_0^{n_3-T}\mathsf Q)$ with $\mathsf Q=\Gamma_0\hat{\mathsf T}_{D_0}\Gamma_1\cdots\hat{\mathsf T}_{D_{T-1}}\Gamma_T$; in an eigenbasis of $\mathsf T_0$, $|\operatorname{Tr}(\mathsf T_0^{n_3-T}\mathsf Q)|\le\|\mathsf Q\|\lambda_0^{n_3-T}(1+\mathsf r_{n_3-T})$, and (a), (b) bound $\|\hat{\mathsf T}_{D_j}\|^{n_3}\le2^{\mathsf v_j+1}\lambda_0^{n_3}(1+\mathsf r_{n_3})z_-^{\ell_j/n_1}$; numbering the layers so that $C$ lies in the layers $1,\dots,h$ makes the $\ell$ of (b) equal to $\ell_j$. (d) The proofs of (a)--(c) use of the representation $g\mapsto g$ only that it is continuous and unitary, which gives the positivity of $\hat{\mathsf T}_D$ and pointwise unitary pieces in the layers, and that it changes its sign at $-\mathbf 1$, which makes every line function odd under the flip and of modulus $\le1$; a cup or a cap of $D$, a linear form or a vector with the entries of a matrix of $D(G)$, has the norm $\sqrt d$, and the trace over $(\mathbb C^d)^{\otimes\mathcal V}$ of the $n_3$-th power gives $d^{\#\mathcal V}Z$ times the mean of the product of the line functions.
\end{proof}

\begin{lemma}[A floor at one length bounds the loops]\label{lem:road}
Let $\beta\ge0$ and $L_1,L_2,L_3\ge3$ integers; take the first direction as time over the section $(L_2,L_3)$, with its $\lambda_0$, $\mathsf r_s$ and $\Lambda_-$, and put $\phi=\log(\lambda_0/\Lambda_-)\in[0,\infty]$. Let $C$ be a rectangle $R\times T$ of a $(1,3)$-plane with $1\le R\le L_1-3$ and $1\le T\le L_3-1$, and $p_o$, $p_r$ the lines of Lemma~\ref{lem:two} through its corners in one layer of the direction $3$. Then
\[
0\le\langle W_C\rangle\le(4\langle p_op_r\rangle)^{T/L_3},\qquad\langle p_op_r\rangle\le\frac{1+\mathsf r_{L_1-R}}{1+\mathsf r_{L_1}}\Bigl[e^{-\phi R}+z_-(L_1-R,L_2,L_3)\Bigr];
\]
so, if $\phi>0$, $\limsup_{L_1\to\infty}\langle W_C\rangle_{(L_1,L_2,L_3)}\le4^{T/L_3}e^{-\phi RT/L_3}$.
\end{lemma}

\begin{proof}[Proof sketch]
The first inequality is Lemma~\ref{lem:two}(c). With the time $1$, as in the proof of Lemma~\ref{lem:two}(d), $Z\langle p_op_r\rangle=\sum_\pm\operatorname{Tr}(\mathsf T_\pm^{L_1-R}\hat{\mathsf p}\,\mathsf T_\mp^{R}\hat{\mathsf p})$ with $\mathsf T_\pm=\mathsf T_0P_\pm$ and $\|\hat{\mathsf p}\|\le1$. The nonzero eigenvalues of $\mathsf T_-$ are those of the sectors $01$ and $11$, so $\|\mathsf T_-\|=\Lambda_-$, and $\|\mathsf T_+\|\le\lambda_0$. For $B\ge0$ of trace class and $Y\ge0$ bounded, $\operatorname{Tr}(BY)=\operatorname{Tr}(B^{1/2}YB^{1/2})\le\|Y\|\operatorname{Tr}B$. With $B=\mathsf T_+^{L_1-R}$ and $Y=\hat{\mathsf p}\mathsf T_-^R\hat{\mathsf p}$ the term of $+$ is $\le\Lambda_-^R\lambda_0^{L_1-R}(1+\mathsf r_{L_1-R})$; with $B=\mathsf T_-^{L_1-R}$ and $Y=\hat{\mathsf p}\mathsf T_+^R\hat{\mathsf p}$ the term of $-$ is $\le\lambda_0^R\operatorname{Tr}\mathsf T_-^{L_1-R}=\lambda_0^Ra_-(L_1-R,L_2,L_3)=\lambda_0^{L_1}(1+\mathsf r_{L_1-R})z_-(L_1-R,L_2,L_3)$. Divide by $Z=\lambda_0^{L_1}(1+\mathsf r_{L_1})$. As $L_1\to\infty$ the quotient of the thermal factors tends to $1$, $\lambda_0$ being simple, and $z_-(s,L_2,L_3)\le\lambda_0^{-s}\operatorname{Tr}\mathsf T_-^s\le(\Lambda_-/\lambda_0)^{s-2}\lambda_0^{-2}\operatorname{Tr}\mathsf T_-^2\to0$ if $\phi>0$.
\end{proof}

\begin{lemma}[Flux weight and vortex line]\label{lem:vort}
\begin{enumerate}
\item[(a)] For integers $n_1'\ge n_1\ge3$, $n_3'\ge n_3\ge3$, $n_2\ge3$ and $\kappa=n_1'n_3'/(n_1n_3)$, a real number,
\[
a_-(n_1',n_2,n_3')\le a_-(\underline n)^\kappa,\qquad z_-(n_1',n_2,n_3')\le\bigl[z_-(\underline n)\,e^{\delta(\underline n)}\bigr]^\kappa .
\]
\item[(b)] For all integers $t,N,n\ge3$, with $\varepsilon=\varepsilon_{\rm v}(t,N)$,
\[
2z_-(t,n,N)\le n\varepsilon+\mathsf r_n(t,N),\qquad n\varepsilon\le\mathsf r_n(t,N)-\log\bigl(1-2z_-(t,2n,N)\bigr).
\]
\end{enumerate}
\end{lemma}

\begin{proof}[Proof sketch]
(a) By Lemma~\ref{lem:tr}(d), $a_-=\sum_i\mu_i^{n_1}$ with $\mu_i$ independent of $n_1$, so Lemma~\ref{lem:seq}(d) with a real exponent gives $a_-(n_1',n_2,n_3)\le a_-(n_1,n_2,n_3)^{n_1'/n_1}$; the same holds in the third side, and $Z(n_1',n_2,n_3')\ge e^{\kappa|\underline n|P}$. (b) With the time along the column, $Z(t,n,N)=\lambda_0^n(1+\mathsf r_n)$ and $Z^-(t,n,N)=\sum_i\mu_i''^{\,n}$ with $\sup_i\mu''_i=\lambda_0e^{-\varepsilon}$; so $1-2z_-=Z^-/Z\ge e^{-n\varepsilon}/(1+\mathsf r_n)\ge1-n\varepsilon-\mathsf r_n$. For the second inequality, which uses that the twisted spectrum is $\ge0$, $\sum\mu_i''^{\,2n}\le\lambda_-^n\sum\mu_i''^{\,n}$ gives $\lambda_-^n\ge Z^-(t,2n,N)/Z^-(t,n,N)\ge(1-2z_-(t,2n,N))\lambda_0^{2n}/(\lambda_0^n(1+\mathsf r_n))$.
\end{proof}

\begin{proposition}[Recursion; excess and thermal numbers]\label{prop:rec}
Let $\beta\ge0$ and let $l\ge3$ be an integer.
\begin{enumerate}
\item[(a)] With $\Sigma_1=(l,l)$, $\Sigma_2=(2l,l)$, $\Sigma_3=(2l,2l)$ and $(d_1,d_2,d_3)=(d_1,d_2,\mathsf e)$: $\log(1+\mathsf r_l(\Sigma_j))\le d_j(l)\le2\log(1+\mathsf r_l(\Sigma_j))$, and $X=4d_1+2d_2+\mathsf e$.
\item[(b)] If $Y_2(\mathsf e(l))<\log2$, then $d_1(2l)\le\omega(\mathsf e(l))$, $d_2(2l)\le\omega(Y_1(\mathsf e(l)))$, $\mathsf e(2l)\le\Psi(\mathsf e(l))$ and $X(2l)\le\Phi(\mathsf e(l))$.
\item[(c)] $\Psi(x)\le12.9x^2$ and $\Phi(x)\le19.89x^2$ for $0<x\le0.069$; $\Psi(x)\le11.03x^2$ and $\Phi(x)\le17.71x^2$ for $0<x\le\frac1{20}$; $\Psi(x)<x$ for $0<x<0.0740538$.
\item[(d)] Under the criterion, for $n=2^kl$, $k\ge1$, and $x=\mathsf e(n/2)\le\theta^{n/2l}/12.9$:
\[
\mathsf r_{n/2}(n,n)\le0.51832\,x,\quad\mathsf r_{n/2}(2n,n)\le1.09990\,x,\quad\mathsf r_{n/2}(2n,2n)\le2.53905\,x;
\]

\[
\mathsf r_n(n,n)\le0.001615\,\theta^{n/l},
\]

\[
\mathsf r_n(2n,n)\le0.00732\,\theta^{n/l},
\]
$\delta_{\rm c}(n)\le0.01655\,\theta^{n/l}$, $d_1(n)\le0.0032288\,\theta^{n/l}$, $d_2(n)\le0.0145396\,\theta^{n/l}$; and $\delta_{\rm c}(N)<0.1324\,\theta^{N/2l}$ for every integer $N\ge2l$.
\item[(e)] Let $\tilde\omega(D)=2\log(1+\varpi(e^D-1)^2)\le\omega(D)$, and let $\tilde Y_1,\tilde Y_2,\tilde\Psi,\tilde\Phi$ be $Y_1,Y_2,\Psi,\Phi$ with $\tilde\omega$ in the place of $\omega$. Then (b) holds with $\tilde\omega$, $\tilde Y_1$, $\tilde\Psi$, $\tilde\Phi$. With $X_0=0.069$ and $X_{k+1}=\tilde\Psi(X_k)$, the iterates $X_1,\dots,X_5$ are $0.0604446$, $0.0432902$, $0.0195722$, $0.0034309$, $0.0000960$ to the digits shown, the $X_k$ decrease to $0$, and $\bar\delta_j=\sum_{s\ge0}8^{-s-1}\tilde\Phi(X_{j-1+s})$ decreases in $j$, with $\bar\delta_4\le0.00073051$. If $\mathsf e(l)\le0.069$, with no other hypothesis, then $\mathsf e(2^kl)\le X_k$ and $\delta_{\rm c}(2^jl)\le\bar\delta_j$ for all $k\ge0$ and $j\ge1$.
\end{enumerate}
\end{proposition}

\begin{proof}[Proof sketch]
(a) Double the sides of the cube one after the other: by (T1), $X$ is the sum of $4[2F(l,l,l)-F(2l,l,l)]$, $2[2F(2l,l,l)-F(2l,2l,l)]$ and $2F(2l,2l,l)-F(2l,2l,2l)$, and each bracket is the $D$ of Lemma~\ref{lem:seq} for $\kappa_i=\lambda_i(\Sigma_j)^l$. (b) Over $(2l,2l)$ with $\kappa_i=\lambda_i^l$: $D=\mathsf e(l)$ and $D'=d_1(2l)\le\omega(\mathsf e)$. Doubling one side of the section at most doubles a defect, up to defects already bounded: by Lemma~\ref{lem:seq}(d) with the time in the second direction, $F(l,4l,2l)\le2F(l,2l,2l)$, so
\[
2F(l,4l,2l)-F(2l,4l,2l)\le2\mathsf e(l)+d_1(2l)\le Y_1(\mathsf e(l)),
\]
and Lemma~\ref{lem:seq}(c) over $(4l,2l)$ gives $d_2(2l)\le\omega(Y_1)$; in the same way $2F(l,4l,4l)-F(2l,4l,4l)\le Y_2$, and over $(4l,4l)$ one doubling of the time squares the defect: $\mathsf e(2l)\le\omega(Y_2)=\Psi$. Then (a) at $2l$. (c) $\varpi(x)/x$ and $(e^D-1)/D$ increase, so $\Psi(x)/x^2$ and $\Phi(x)/x^2$ do not decrease on $(0,0.0741]$; their values at $0.069$ are $12.893528$ and $19.881793$, at $\frac1{20}$ they are $11.029686$ and $17.709146$, to the digits shown; $\Psi(x)/x$ increases; its positive fixed point lies in $(0.0740538,0.0740539)$, and $Y_2<\log2$ throughout this interval. (d) Induction with (b), (c): $12.9\,\mathsf e(2^{k+1}l)\le(12.9\,\mathsf e(2^kl))^2$, so $\mathsf e(2^kl)\le\theta^{2^k}/12.9\le0.069$ and $X(2^{k+1}l)\le\frac{19.89}{166.41}\theta^{2^{k+1}}\le0.11953\,\theta^{2^{k+1}}$. At $l'=n/2$ the sections $(2l',2l')$, $(4l',2l')$, $(4l',4l')$ have at the time $l'$ the defects $x$, $\le Y_1(x)$, $\le Y_2(x)$, and Lemma~\ref{lem:seq}(b) with the monotone quotients at $0.069$ gives the first three bounds. Then $\mathsf r_n=s_2\le\eta^2$ gives $\mathsf r_n(n,n)\le\frac12\omega(x)\le0.268652\,x^2\le0.001615\,\theta^{n/l}$, likewise $\mathsf r_n(2n,n)$, and $d_1(n)\le\omega(x)$, $d_2(n)\le\omega(Y_1(x))$. For the excess, the pressure cancels in $X(n)=8\delta_{\rm c}(n)-\delta_{\rm c}(2n)$, and Lemma~\ref{lem:seq}(f) gives
\[
\delta_{\rm c}(n)=\sum_{i\ge0}8^{-i-1}X(2^in)\le0.11953\,\theta^{n/l}\sum_{i\ge0}8^{-i-1}\theta^{(2^i-1)n/l}\le0.01655\,\theta^{n/l}.
\]
For an integer $N$ with $2^jl\le N<2^{j+1}l$, $j\ge1$, Lemma~\ref{lem:seq}(f) gives $\delta_{\rm c}(N)\le8\,\delta_{\rm c}(2^jl)\le0.1324\,\theta^{2^j}<0.1324\,\theta^{N/2l}$. (e) By Lemma~\ref{lem:seq}(b), (c), $D'\le2\log(1+\eta^2)\le\tilde\omega(D)$, and the proof of (b) uses of $\omega$ only that it increases. The numbers $X_1,\dots,X_5$ and $\bar\delta_4$ are the values of the recursion $X_{k+1}=\tilde\Psi(X_k)$ from $X_0=0.069$ and of the series that defines $\bar\delta_4$; the $X_k$ decrease to $0$ because $\tilde\Psi\le\Psi$ and $\Psi(x)<x$ on $(0,0.0740538)$, and $\bar\delta_j$ decreases because $\tilde\Phi$ increases. If $\mathsf e(l)\le X_0$, the induction gives $\mathsf e(2^kl)\le X_k$, and $\delta_{\rm c}(2^jl)=\sum_{s\ge0}8^{-s-1}X(2^{j+s}l)\le\sum_s8^{-s-1}\tilde\Phi(X_{j+s-1})=\bar\delta_j$.
\end{proof}

\begin{proposition}[Vortex lines on tubes]\label{prop:tub}
Assume the criterion; put $l_j=2^{j+1}l$ and $w=2^j$.
\begin{enumerate}
\item[(a)] $l'\varepsilon_{\rm v}(l')\le0.01643\,\theta^{l'/l}$ for $l'=2^kl$, $k\ge1$.
\item[(b)] Let $l'=2^kl$, $k\ge1$, $l'\varepsilon_{\rm v}(l')\le e^{-E}$, $\hat N\ge2$ an integer, $y=\hat Ne^{-E}$, $\mathsf t=\theta^{l'/l}$, $\bar z_1=y+\frac12(0.001615\,\mathsf t)^{2\hat N}$ and $\bar z_2=\bar z_1^4e^{8\hat N\cdot0.01655\,\mathsf t}<\frac12$. Then
\[
2l'\varepsilon_{\rm v}(2l')\le\frac2{\hat N}\Bigl\{(0.001615\,\mathsf t^2)^{\hat N/2}+\frac{2\bar z_2}{1-2\bar z_2}\Bigr\}.
\]
\item[(c)] $l_j\varepsilon_{\rm v}(l_j)\le e^{-\sigma^+l_j^2}$ for all $j\ge0$.
\item[(d)] For real $M\ge w\mathsf q$ with $Ml_j$ an integer: $z_-(l_j,Ml_j,l_j)\le Me^{-\varsigma w^2}$ and $\delta(l_j,Ml_j,l_j)\le0.01655\,M\theta^{l_j/l}$.
\item[(e)] Let $l'=l_j$ and let $N_3\ge l'$ and $n_2\ge\mathsf ql'^2/2l$ be integers. Over the section $(n_2,N_3)$, with the first direction as time, the top $\Lambda_-$ of the sectors odd under the flip of the side $N_3$ lies at least $\sigma^+N_3-\frac{N_3}{l'^2}\bigl[\log\frac{n_2}{l'}+0.01655\,\frac{n_2}{l'}\theta^{l'/l}\bigr]$ above the vacuum.
\end{enumerate}
\end{proposition}

\begin{proof}[Proof sketch]
(a) The second inequality of Lemma~\ref{lem:vort}(b) at $t=N=n=l'$, with $z_-(l',2l',l')\le\mathsf r_{l'}(2l',l')\le0.00732\,\mathsf t$ (Lemma~\ref{lem:tr}(b), the time $1$) and $\mathsf r_{l'}(l',l')\le0.001615\,\mathsf t$ (Proposition~\ref{prop:rec}(d)). (b) Three moves. The first inequality of Lemma~\ref{lem:vort}(b) on the box $(l',2\hat Nl',l')$, with $\mathsf r_{2\hat Nl'}(l',l')\le\mathsf r_{l'}(l',l')^{2\hat N}$ (Lemma~\ref{lem:seq}(e)), gives $z_-\le\bar z_1$. Lemma~\ref{lem:vort}(a) with $\kappa=4$ passes to the box $(2l',2\hat Nl',2l')$, a fourth power at the cost $e^{4\delta}$, with $\delta(l',2\hat Nl',l')\le2\hat N\delta_{\rm c}(l')$. The second inequality of Lemma~\ref{lem:vort}(b) for the tube $2l'\times2l'$ at the time $\hat Nl'$, with $\mathsf r_{\hat Nl'}(2l',2l')\le\mathsf r_{2l'}(2l',2l')^{\hat N/2}$, returns to the vortex line. (c) Take $\hat N_j=4\cdot3^j$ and write $\mathsf t_j=\theta^{l_j/l}=e^{-2^j\cdot2\Lambda}$. For real $E$ and $\mathsf t\in(0,1]$ let $\mathsf H_j(E,\mathsf t)$ be the right member of (b) at $l'=l_j$ and $\hat N=\hat N_j$, formed with $y=\hat N_je^{-E}$, $\bar z_1=y+\frac12(0.001615\,\mathsf t)^{2\hat N_j}$ and $\bar z_2=\bar z_1^4e^{8\hat N_j\cdot0.01655\,\mathsf t}$, where $\bar z_2<\frac12$. Identically $\mathsf H_j=\frac2{\hat N_j}y^4e^{8\hat N_j\cdot0.01655\,\mathsf t}\mathsf B_j$, that is,
\[
-\log\mathsf H_j(E,\mathsf t)=4E-3\log\hat N_j-8\hat N_j\cdot0.01655\,\mathsf t-\log(2\mathsf B_j),\qquad\mathsf B_j=\frac{(0.001615\,\mathsf t^2)^{\hat N_j/2}}{y^4\,e^{8\hat N_j\cdot0.01655\,\mathsf t}}+\frac{2(\bar z_1/y)^4}{1-2\bar z_2}.
\]
\emph{One value of $\theta$ serves all.} Let $\Delta=2\Lambda-0.23284$, which is $\ge0$ as $2\log\frac1{0.8901}\ge0.23284$, and $\mathring{\mathsf t}_j=e^{-2^j\cdot0.23284}$, the value of $\mathsf t_j$ at $2\Lambda=0.23284$, so that $\mathsf t_j=\mathring{\mathsf t}_je^{-2^j\Delta}$. Let $E_0=4.3414$, and let $E_1,E_2,\dots$ be numbers with $E_{j+1}\le-\log\mathsf H_j(E_j,\mathring{\mathsf t}_j)$, $\bar z_2<\frac12$ at $(E_j,\mathring{\mathsf t}_j)$. Then $l_j\varepsilon_{\rm v}(l_j)\le e^{-E_j-4^j\Delta}$ for all $j\ge0$ and every $\theta\le0.8901$. At $j=0$ this is (a) at $l'=2l$, as $0.01643\,\theta^2=\exp[-\log\frac1{0.01643}-0.23284-\Delta]$ and $\log\frac1{0.01643}+0.23284\ge4.3414$. From $j$ to $j+1$ apply (b) with $E=E_j+4^j\Delta$ and $\mathsf t=\mathsf t_j$, and compare each number with its value at $(E_j,\mathring{\mathsf t}_j)$: $y$ is that value times $e^{-4^j\Delta}$; $(0.001615\,\mathsf t)^{2\hat N_j}$ is it times $e^{-2^{j+1}\hat N_j\Delta}\le e^{-4^j\Delta}$; so $\bar z_1$ is at most it times $e^{-4^j\Delta}$; as $e^{8\hat N_j\cdot0.01655\,\mathsf t}$ increases in $\mathsf t$, $\bar z_2$ is at most it times $e^{-4^{j+1}\Delta}$, and so is $\frac{2\bar z_2}{1-2\bar z_2}$, as $u\mapsto1/(1-2u)$ increases; and $(0.001615\,\mathsf t^2)^{\hat N_j/2}$ is it times $e^{-2^j\hat N_j\Delta}\le e^{-4^{j+1}\Delta}$, as $\hat N_j\ge4\cdot2^j$.

So each of the two terms of the right member of (b) is at most its value at $(E_j,\mathring{\mathsf t}_j)$ times $e^{-4^{j+1}\Delta}$, and
\[
l_{j+1}\varepsilon_{\rm v}(l_{j+1})\le\mathsf H_j(E_j+4^j\Delta,\mathsf t_j)\le\mathsf H_j(E_j,\mathring{\mathsf t}_j)\,e^{-4^{j+1}\Delta}\le e^{-E_{j+1}-4^{j+1}\Delta}.
\]
\emph{Numerical bounds.} The functions in the recursion are explicit and monotone on the intervals used above. Outward interval evaluation gives
\[
\begin{array}{c|cc}
x&\Psi(x)/x^2&\Phi(x)/x^2\\ \hline
0.069&(12.89352831,12.89352832)&(19.88179283,19.88179284)\\
0.05&(11.02968611,11.02968612)&(17.70914610,17.70914611).
\end{array}
\]
For the successive lower bounds $E_0,\ldots,E_6$ below, the differences
$-\log\mathsf H_j(E_j,\mathring{\mathsf t}_j)-E_{j+1}$ exceed, respectively,
\[
 9.7\cdot10^{-5},\quad5.1\cdot10^{-5},\quad3.8\cdot10^{-6},
 \quad6.4\cdot10^{-5},\quad3.5\cdot10^{-5},\quad7.6\cdot10^{-5}.
\]
In each evaluation $\bar z_2<1/2$. The geometric tail displayed below gives a lower bound greater than $1.909266$. The sharper defect recursion likewise gives $\bar\delta_4<0.00073051$. These bounds justify the directions of rounding used in the argument.

\emph{The numbers at $2\Lambda=0.23284$.} There $\mathsf B_j$ is a number close to $2$, and one may take $E_1,\dots,E_6=11.3558$, $35.5848$, $128.3243$, $495.6445$, $1962.8155$, $7829.1628$, each at most $-\log\mathsf H_j(E_j,\mathring{\mathsf t}_j)$ at the one before. For $j\ge6$ the losses $3\log\hat N_j+8\hat N_j\cdot0.01655\,\mathring{\mathsf t}_j+\log(2\mathsf B_j)$ are at most $3\log\hat N_j+\log4.01$, so that $E_{j+1}=4E_j-3\log\hat N_j-\log4.01$ may be taken, and this closed recursion gives
\[
4^{-j}E_j\ge4^{-6}E_6-\sum_{i\ge6}4^{-i-1}\bigl(3\log4+3i\log3+\log4.01\bigr)\ge1.909266 .
\]
So $4^{-j}E_j\ge0.23284+1.67$ for every $j\ge0$, hence $4^{-j}(E_j+4^j\Delta)\ge2\Lambda+1.67$, and $l_j\varepsilon_{\rm v}(l_j)\le\exp[-4^j(2\Lambda+1.67)]=e^{-\sigma^+l_j^2}$. (d) The first inequality of Lemma~\ref{lem:vort}(b) with $n=Ml_j$, (c), and $\mathsf r_{Ml_j}(l_j,l_j)\le(0.001615\,\theta^{2w})^M\le e^{-2\Lambda Mw}\le e^{-\varsigma w^2}$, as $2\Lambda\mathsf q=\varsigma$, give $2z_-\le(M+1)e^{-\varsigma w^2}$; Lemma~\ref{lem:seq}(f) gives $\delta\le M\delta_{\rm c}(l_j)$. (e) For every integer $t\ge l'$, by Lemma~\ref{lem:tr}(d) and Lemma~\ref{lem:vort}(a) at the real factors $t/l'$ and $N_3/l'$, $\Lambda_-^t\le a_-(t,n_2,N_3)\le a_-(l',n_2,l')^{tN_3/l'^2}$; so $\log\Lambda_-\le\frac{N_3}{l'^2}[F+\log z_-](l',n_2,l')$, while $\log\lambda_0(n_2,N_3)\ge n_2N_3P$. Subtracting, $\log(\lambda_0/\Lambda_-)\ge\frac{N_3}{l'^2}[-\log z_--\delta](l',n_2,l')$, and (d) with $M=n_2/l'\ge w\mathsf q$ bounds the bracket below by $\varsigma w^2-\log M-0.01655\,M\theta^{l'/l}$, with $\varsigma w^2=\sigma^+l'^2$.
\end{proof}

\begin{proposition}[Slabs]\label{prop:slab}
For integers $N,a,b,t\ge3$ let $\mathsf e^{(N)}(t)=2F(t,2t,N)-F(2t,2t,N)$, $d^{(N)}(t)=2F(t,t,N)-F(2t,t,N)$, $P^{(N)}=\inf_{a,b}F(a,b,N)/ab$ and $\delta^{(N)}(a,b)=F(a,b,N)-abP^{(N)}$.
\begin{enumerate}
\item[(a)] For every $\beta\ge0$: $0\le\delta^{(N)}(a,b)\le\delta(a,b,N)$, $\delta^{(N)}(a,b)/ab$ does not increase in $a$ or $b$, $\delta^{(N)}(a,a)/a^2\to0$, and $\log(1+\mathsf r_a(b,N))\le\delta^{(N)}(a,b)$; if $Y_1(\mathsf e^{(N)}(t))<\log2$, then $d^{(N)}(2t)\le\omega(\mathsf e^{(N)}(t))$ and $\mathsf e^{(N)}(2t)\le\omega(Y_1(\mathsf e^{(N)}(t)))$.
\end{enumerate}
Assume now the criterion.
\begin{enumerate}
\item[(b)] Let $n_0=2^kl\le N<2n_0$, $k\ge1$, $\chi'=0.3602\,\theta^{n_0/l}$ and $t_i=2^in_0$. Then $2.72\,\mathsf e^{(N)}(t_i)\le\chi'^{2^i}$, $\delta^{(N)}(t_i,t_i)\le0.133\,\chi'^{2^i}$, $\delta^{(N)}(n_0,n_0)\le0.01655\,\frac N{n_0}\theta^{n_0/l}$, $\mathsf r_{t_i}(t_i,N)\le0.14\,\chi'^{2^i}$, and $\mathsf r_{t_i}(b,N)\le0.28\,\chi'^{2^i}$ for $t_i\le b<2t_i$.
\item[(c)] Let $N=2^kl$, $k\ge1$, $\chi=0.03\,\theta^{N/l}$ and $t_i=2^iN$. Then $2.06\,\mathsf e^{(N)}(t_i)\le\chi^{2^i}$, $d^{(N)}(t_i)\le0.1202\,\chi^{2^i}$, $\mathsf r_{t_i}(t_i,N)\le0.0613\,\chi^{2^i}$, $\mathsf r_{t_i}(2t_i,N)\le\chi^{2^i}$ and $\delta^{(N)}(t_i,t_i)\le0.1826\,\chi^{2^i}$. If moreover $N=2wl$ with $w\ge2$, and $\hat N$, $\mathsf g$ are as in Theorem~\ref{thm:one}(g1), then $t_i\varepsilon_{\rm v}(t_i,N)\le\exp[-2^i(\sigma^+N^2-\mathsf g)]$ for all $i\ge0$.
\end{enumerate}
\end{proposition}

\begin{proof}[Proof sketch]
(a) The slab of thickness $N$ is a datum in two dimensions: $(a,b)\mapsto F(a,b,N)$ is symmetric and has the trace formula in each variable, so Lemma~\ref{lem:seq} and the proof of Proposition~\ref{prop:rec}(b) go over with one side of the section frozen; $P^{(N)}\ge NP$ gives $\delta^{(N)}\le\delta(\cdot,\cdot,N)$. (b) The pressures cancel at the start: $\mathsf e^{(N)}(n_0)\le2\delta(n_0,2n_0,N)\le4\frac N{n_0}\delta_{\rm c}(n_0)<0.1324\,\theta^{n_0/l}\le0.1049$. On $(0,0.1049]$, $\omega(Y_1(x))\le2.72x^2$ and $\omega(x)\le0.5593x^2$, so $2.72\,\mathsf e^{(N)}(t_{i+1})\le(2.72\,\mathsf e^{(N)}(t_i))^2$ and $2.72\,\mathsf e^{(N)}(n_0)\le\chi'$. Since $4F(t,t,N)-F(2t,2t,N)=2d^{(N)}(t)+\mathsf e^{(N)}(t)=4\delta^{(N)}(t,t)-\delta^{(N)}(2t,2t)$, and for $i\ge1$, $2d^{(N)}(t_i)+\mathsf e^{(N)}(t_i)\le(2\cdot0.5593+2.72)\,\mathsf e^{(N)}(t_{i-1})^2\le0.518842\,\chi'^{2^i}$, summation with the weights $4^{-j-1}$ gives $\delta^{(N)}(t_i,t_i)\le0.13236\,\chi'^{2^i}$ for $i\ge1$; at $i=0$, $\delta^{(N)}(n_0,n_0)\le\delta(n_0,n_0,N)\le\frac N{n_0}\delta_{\rm c}(n_0)<0.0331\,\theta^{n_0/l}\le0.0919\,\chi'$.

Finally $\log(1+\mathsf r_{t_i}(b,N))\le\delta^{(N)}(t_i,b)\le\frac b{t_i}\delta^{(N)}(t_i,t_i)$. (c) At the dyadic thickness the start is $\mathsf e^{(N)}(N)=d_2(N)\le0.0145396\,\theta^{N/l}\le0.01152$ and $d^{(N)}(N)=d_1(N)\le0.0032288\,\theta^{N/l}$; for $x\le0.0116$, $\omega(Y_1(x))\le2.06x^2$, $\omega(x)\le0.5059x^2$ and $e^x-1\le1.006x$, and for $D\le0.00286$, $\varpi(e^D-1)\le0.5008D$. So $2.06\,\mathsf e^{(N)}(t_i)\le\chi^{2^i}$; $d^{(N)}(t_i)\le0.5059\,\mathsf e^{(N)}(t_{i-1})^2\le0.1202\,\chi^{2^i}$; Lemma~\ref{lem:seq}(b) over $(t_i,N)$ gives $\mathsf r_{t_i}(t_i,N)\le0.5008\,d^{(N)}(t_i)$, and Lemma~\ref{lem:seq}(a) over $(2t_i,N)$ gives $\mathsf r_{t_i}(2t_i,N)\le e^{\mathsf e^{(N)}(t_i)}-1$; the sum with the weights $4^{-j-1}$ gives the bound on $\delta^{(N)}$.

In the flat tubes $t\times N$ one doubling of the width squares $t\varepsilon_{\rm v}(t,N)$. At $i=0$, Proposition~\ref{prop:tub}(c) gives $N\varepsilon_{\rm v}(N,N)\le e^{-E_0}$, $E_0=\sigma^+N^2$. From $t=t_i$ to $2t$, with $n=\hat Nt$ and $y_i=\hat Ne^{-E_i}$: the first inequality of Lemma~\ref{lem:vort}(b) at the time $2n$, with $\mathsf r_{2n}(t,N)\le\mathsf r_t(t,N)^{2\hat N}\le e^{-4E_i}$, gives $z_-(t,2n,N)\le(1+10^{-7})y_i$; Lemma~\ref{lem:vort}(a) with $\kappa=2$ in the width squares it at the cost $2F(t,2n,N)-F(2t,2n,N)\le4\hat N\delta^{(N)}(t,t)\le\xi_i:=0.732\,\hat N\chi^{2^i}$; and the second inequality of Lemma~\ref{lem:vort}(b) for the tube $2t\times N$ at the time $n$ gives
\[
2t\,\varepsilon_{\rm v}(2t,N)\le4.01\,\hat Ne^{-2E_i+\xi_i}=:e^{-E_{i+1}} .
\]
As $\hat N\chi\le0.7936$, $\sum_i2^{-i-1}\xi_i\le0.294$, and $2^{-i}E_i\ge\sigma^+N^2-\log\hat N-\log4.01-0.294\ge\sigma^+N^2-\mathsf g$.
\end{proof}

\begin{proposition}[Flux lines of every integer length]\label{prop:lines}
Assume the criterion; let $N\ge4l$ be an integer, $n_0$, $\kappa$, $\tau$, $\chi_0$, $\mathsf E_0$, $\hat N'$, $\mathsf h$ and $\mathsf g'$ as in Theorem~\ref{thm:one}(g2), $\mathsf E=\sigma^+n_0^2$ and $t_i=2^in_0$. Then $6\le\hat N'<\mathsf qN/l+2$, $\hat N'\le1.3453\,\mathsf E_0+2$, $\mathsf h\le0.4274\,\kappa+0.0208$, and $t_i\varepsilon_{\rm v}(t_i,N)\le\exp[-2^i(\mathsf E_0-\mathsf g')]$ for all $i\ge0$.
\end{proposition}

\begin{proof}[Proof sketch]
\emph{Numbers.} $\log\frac1{\chi_0}=-\log0.3602+2w\Lambda\ge1.48677$; the bounds on $\hat N'$ follow from $2\mathsf E_0/\log\frac1{\chi_0}<2\kappa\mathsf qw=\mathsf qN/l$, from the fact that the quotient exceeds $4$ at $w=2$ and increases in $w$, and from $2/1.48677\le1.3453$. As $\tau=\theta^{2w}\le0.8901^4$, the number $c_\tau=0.01655\,\tau$ satisfies $c_\tau\le0.010389$, and $\mathsf h=c_\tau\hat N'$. With $4\mathsf q\cdot0.01655\,we^{-2w\Lambda}\le0.8548$ (as in the proof of Theorem~\ref{thm:one}(d)), $\mathsf h<0.01655(2\kappa\mathsf qw+2)\tau\le0.4274\kappa+0.0208$. By the choice of $\hat N'$, $\chi_0^{\hat N'/2}\le e^{-\mathsf E_0}$. Here $\mathsf E_0=\kappa\mathsf E$ and $\mathsf E=\sigma^+n_0^2=\varsigma w^2\ge4\cdot1.90284=7.61136$. For $\kappa\in[1,2]$, $E\ge7.6$ and $0<c\le0.052$ the functions
\[
\gamma_c(\kappa,E)=\kappa\bigl[\log(1.3453\,\kappa E+2)+c\,(1.3453\,\kappa E+2)-E\bigr]
\]
increase in $c$, are convex in $\kappa$ and decrease in $E$. As $\log\hat N'\le\log(1.3453\,\kappa\mathsf E+2)$ and $\mathsf h\le c_\tau(1.3453\,\kappa\mathsf E+2)$,
\[
\log\bigl[\hat N'e^{-\mathsf E+\mathsf h}\bigr]^\kappa\le\gamma_{0.010389}(\kappa,\mathsf E),\qquad\log\bigl[\hat N'e^{-\mathsf E+5\mathsf h}\bigr]^\kappa\le\gamma_{0.051945}(\kappa,\mathsf E),\qquad0.051945=5\cdot0.010389 .
\]
On $\kappa\in[1,2]$, $E\ge7.61136$ each of the two functions is largest at $E=7.61136$ and at $\kappa=1$ or $\kappa=2$, where $\gamma_{0.010389}$ has approximate values $-4.97953$ and $-8.53047$, and $\gamma_{0.051945}$ approximate values $-4.47090$ and $-6.66219$. So $\bar z_0:=[\hat N'e^{-\mathsf E+\mathsf h}]^\kappa\le e^{-4.97952}\le0.00688$ and $6.1\,[\hat N'e^{-\mathsf E+5\mathsf h}]^\kappa\le6.1\,e^{-4.47090}\le0.0698$. \emph{The start.} The first inequality of Lemma~\ref{lem:vort}(b) for the tube $n_0\times n_0$ at the time $\hat N'n_0$, Proposition~\ref{prop:tub}(c) and $\mathsf r_{\hat N'n_0}(n_0,n_0)\le\chi_0^{\hat N'}\le e^{-\mathsf E}$ give $z_-(n_0,\hat N'n_0,n_0)\le\hat N'e^{-\mathsf E}$; with $\delta(n_0,\hat N'n_0,n_0)\le\hat N'\delta_{\rm c}(n_0)\le\mathsf h$, Lemma~\ref{lem:vort}(a) at the real factor $\kappa$ in the third side gives $z_-(n_0,\hat N'n_0,N)\le\bar z_0$; and the second inequality of Lemma~\ref{lem:vort}(b) for the tube $n_0\times N$ at the time $\hat N'n_0/2$, with $\mathsf r_{\hat N'n_0/2}(n_0,N)\le(0.14\chi_0)^{\hat N'/2}\le\bar z_0$ (Proposition~\ref{prop:slab}(b)), gives $n_0\varepsilon_{\rm v}(n_0,N)\le\frac{6.1}{\hat N'}\bar z_0=e^{-E'_0}$, $E'_0=\mathsf E_0-(\kappa-1)\log\hat N'-\kappa\mathsf h-\log6.1$.

\emph{The recursion.} As in Proposition~\ref{prop:slab}(c), with the slab of the integer thickness $N$ of Proposition~\ref{prop:slab}(b), one doubling of the width squares $t\varepsilon_{\rm v}(t,N)$: $E'_{i+1}=2E'_i-\log\hat N'-\log4.1-\xi'_i$, with $\xi'_0=4\kappa\mathsf h$ and $\xi'_i=0.532\,\hat N'\chi_0^{2^i}$ for $i\ge1$; the bound $6.1[\hat N'e^{-\mathsf E+5\mathsf h}]^\kappa\le0.0698$ keeps $y_i=\hat N'e^{-E'_i}$ so small that the step closes at every $i$, and $\sum_i2^{-i-1}\xi'_i\le(2\kappa+0.68)\mathsf h$. Hence $2^{-i}E'_i\ge E'_0-\log\hat N'-\log4.1-(2\kappa+0.68)\mathsf h\ge\mathsf E_0-\mathsf g'$, as $\log6.1+\log4.1\le3.22$. More precisely, the unused positive tail of this same recurrence gives \[\begin{multlined}E'_i=2^i\Bigl(E'_0-\log\hat N'-\log4.1-\sum_{j\ge0}2^{-j-1}\xi'_j\Bigr)+\log\hat N'+\log4.1+\sum_{j\ge i}2^{i-j-1}\xi'_j\\ \ge2^i(\mathsf E_0-\mathsf g')+\log\hat N'+\log4.1+\tfrac12\xi'_i.\end{multlined}\] Thus subtracting $\log\hat N'+\xi'_i/4$ still leaves the floor $2^i(\mathsf E_0-\mathsf g')$, also at $i=0$.
\end{proof}

\begin{proposition}[A local function in a long box]\label{prop:long}
Let $\beta\ge0$, $\Sigma$ a section, $h\ge1$ and $t\ge h+2$ integers, $s=t-h-1$, and $A$ a bounded Borel function of the variables of $h$ consecutive layers and of the crossings between them. There is a number $a_\infty(\Sigma,A)$, $|a_\infty|\le\sup|A|$, that depends neither on $t$ nor on the position of the layers, with $|\langle A\rangle_{(t,\Sigma)}-a_\infty(\Sigma,A)|\le2\sup|A|\,\mathsf r_s(\Sigma)$. As $\mathsf r_s(\Sigma)\to0$, the laws of the boxes $(t,\Sigma)$ converge, as $t\to\infty$, on the bounded Borel functions of finitely many layers: to the law of the infinite cylinder over $\Sigma$.
\end{proposition}

\begin{proof}[Proof sketch]
Integrating the weights of the $h$ layers and the $h+1$ crossings with $A$ inserted gives a kernel $\hat A$ dominated by $\sup|A|$ times that of $\mathsf T_0^{h+1}$, so $\|\hat A\|\le\sup|A|\lambda_0^{h+1}$, and $Z\langle A\rangle=\operatorname{Tr}(\mathsf T_0^s\hat A)=\sum_i\lambda_i^s\langle e_i,\hat Ae_i\rangle$ in an eigenbasis. With $a_\infty=\lambda_0^{-h-1}\langle e_0,\hat Ae_0\rangle$, which is determined because $\lambda_0$ is simple: $Z\langle A\rangle=\lambda_0^t(a_\infty+\vartheta\sup|A|\,\mathsf r_s)$ with $|\vartheta|\le1$, and $Z=\lambda_0^t(1+\mathsf r_t)$. The thermal number tends to $0$ by dominated convergence, $\lambda_0$ being simple and $\sum\lambda_i^2<\infty$; the limits are consistent, and Kolmogorov's theorem gives the law of the cylinder.
\end{proof}

\begin{lemma}[Mirrors of a box]\label{lem:mir}
Fix a box and a direction $\mu$ as time; let $c\in\mathbb Z+\frac12$, and $\mathfrak r$ the reflection $x_\mu\mapsto2c-x_\mu$ of the box.
\begin{enumerate}
\item[(a)] Let $\vartheta\ge2$ be an integer with $2\vartheta\le n_\mu$, and $A$ a bounded Borel function of the bonds with both ends in the $\vartheta$ layers below $c$. Then $\langle A\,\overline{\mathfrak rA}\rangle\ge0$ and $|\langle A\rangle|^2\le\langle A\,\overline{\mathfrak rA}\rangle$; $n_\mu$ may be odd. If $n_\mu=2\vartheta$ and $A'$ is another such function, $|\langle A\,\overline{\mathfrak rA'}\rangle|^2\le\langle A\,\overline{\mathfrak rA}\rangle\langle A'\,\overline{\mathfrak rA'}\rangle$.
\item[(b)] Let $N=2h\ge4$, $c\in(\mathbb Z+\frac12)^3$, and $H_\epsilon$, $\epsilon\in\{0,1\}^3$, functions of the cell $\mathsf C_h(c)$. On the cube $(N,N,N)$, $|\langle\prod_\epsilon H_\epsilon^{[\epsilon]}\rangle|\le\prod_\epsilon|\!|\!|H_\epsilon|\!|\!|_N$; so $|\!|\!|\cdot|\!|\!|_N$ is a seminorm, $|\langle A\rangle_{(N,N,N)}|\le|\!|\!|A|\!|\!|_N\le\sup|A|$, the mean of the product of the copies of $A$ over $\epsilon\in E$ is at most $|\!|\!|A|\!|\!|_N^{\#E}$ in modulus for every $E\subset\{0,1\}^3$, and $|\!|\!|A|\!|\!|_N^8\le\langle|A|^8\rangle_{(N,N,N)}$.
\item[(c)] Let $n_\mu$ be even and the mirror lie between the layers $0$ and $1$ of the direction $\mu$; let $x_p$, $x_q$ be two sites of the layer $1$, equal or not, $\mathsf z_p$, $\mathsf z_q$ the values of $U$ at the bonds from $x_p$ to $\mathfrak rx_p$ and from $x_q$ to $\mathfrak rx_q$, and $A_{jk}$, $j,k\in\{1,2\}$, bounded Borel functions of the bonds with both ends in the layers $1,\dots,n_\mu/2$. Then
\[
\Bigl\langle\sum_{i,j,k,l}\overline{(\mathsf z_p)_{ji}}\,(\mathsf z_q)_{kl}\,A_{jk}\,\overline{\mathfrak rA_{il}}\Bigr\rangle\ge0 .
\]
\end{enumerate}
\end{lemma}

\begin{proof}[Proof sketch]
Let $\mathsf T'=\mathsf K^{1/2}M_0^2\mathsf K^{1/2}$, which has the eigenvalues of $\mathsf T_0$. For a function $A$ of $\vartheta$ consecutive layers and the $\vartheta-1$ crossings between them let $\mathsf S_A$ be the operator whose kernel, a function of the first and the last of these layers, is the integral of $A$ against the weights of the $\vartheta$ layers and of the inner crossings over all other variables, and $\mathsf Y_A=\mathsf K^{1/2}\mathsf S_A\mathsf K^{1/2}$, a Hilbert--Schmidt operator; $\mathsf Y_1=\mathsf T'^{\,\vartheta}$. Integrating the layers in turn as in Lemma~\ref{lem:tr}, the mean of a product of such functions on consecutive slabs, with free layers between them, is the trace of the product of their operators and of powers of $\mathsf T'$; and, the weights being invariant under $\mathfrak r$ because $A_\times(V',S^{-1},V)=A_\times(V,S,V')$, the operator of $\overline{\mathfrak rA}$ on the mirrored slab is $\mathsf Y_A^*$.

So $Z\langle A\rangle=\operatorname{Tr}(\mathsf T'^{\,n_\mu-\vartheta}\mathsf Y_A)$ and $Z\langle A\,\overline{\mathfrak rA'}\rangle=\operatorname{Tr}(\mathsf T'^{\,n_\mu-2\vartheta}\mathsf Y_{A'}^*\mathsf Y_A)$. (a) Write $\operatorname{Tr}(\mathsf T'^{\,n_\mu-\vartheta}\mathsf Y_A)=\operatorname{Tr}(\mathsf T'^{\,n_\mu/2}\cdot\mathsf Y_A\mathsf T'^{\,v})$ with $v=\frac12(n_\mu-2\vartheta)\ge0$, a real power of a positive operator; Schwarz's inequality for Hilbert--Schmidt operators gives $|Z\langle A\rangle|^2\le\operatorname{Tr}\mathsf T'^{\,n_\mu}\,\|\mathsf Y_A\mathsf T'^{\,v}\|_2^2=Z\cdot Z\langle A\,\overline{\mathfrak rA}\rangle$, with $\|\mathsf Y_A\mathsf T'^{\,v}\|_2^2=\operatorname{Tr}(\mathsf T'^{\,2v}\mathsf Y_A^*\mathsf Y_A)\ge0$; at $n_\mu=2\vartheta$, Schwarz's inequality for $\operatorname{Tr}(\mathsf Y_{A'}^*\mathsf Y_A)$ is the second statement. (b) The reflection $\mathfrak r_\mu$ of the cube fixes the two bond mirrors $x_\mu=c_\mu$ and $x_\mu=c_\mu+h$, and the half below the first of them is the slab of the cell.

Group the eight factors into the four with $\epsilon_\mu=0$ and their four mirror images: the second statement of (a) gives $|\Lambda(H)|^2\le\Lambda(H^0)\Lambda(H^1)$ for $\Lambda(H)=\langle\prod_\epsilon H_\epsilon^{[\epsilon]}\rangle$, where $H^j$ is the family with the index $\epsilon_\mu$ set to $j$; three directions give the first bound, and the mean of the eight copies of one function is $\ge0$. Expanding $|\!|\!|H+H'|\!|\!|_N^8$ into $2^8$ such means gives the triangle inequality. With $H_\epsilon=A$ for $\epsilon\in E$ and $H_\epsilon=1$ otherwise one gets the bound for subsets; H\"older's inequality with eight exponents $8$ and the invariance of the cube law under the $\mathfrak r_\mu$ give the last bound. (c) Take $\mu$ as time; the mirror also fixes the plane between the layers $n_\mu/2$ and $n_\mu/2+1$.

Integrating in each half the layers and crossings between its two outer layers produces from $A_{jk}$ a function $\Phi_{jk}$ of $(V^1,V^{n_\mu/2})$, and from $\overline{\mathfrak rA_{il}}$, by the symmetry of the weights, $\overline{\Phi_{il}}$ at the mirrored layers $0$ and $n_\mu/2+1$. The crossing between the layers $0$ and $1$ carries the two cut bonds, which enter exactly as in the operator $\hat{\mathsf K}$ of Lemma~\ref{lem:two} with the charges at the sites of $x_p$ and $x_q$: the function under the mean is $\operatorname{tr}[(\overline{\mathfrak rA})^{\rm t}\mathsf z_p^{-1}A\,\mathsf z_q]$ for the matrices $A=(A_{jk})$, $\overline{\mathfrak rA}=(\overline{\mathfrak rA_{il}})$. The crossing between $n_\mu/2$ and $n_\mu/2+1$ enters through $\mathsf K$.

Hence $Z$ times the mean equals $\langle\Phi,(\hat{\mathsf K}\otimes\mathsf K)\Phi\rangle\ge0$, by Lemma~\ref{lem:two}(a), (e), whose proof uses of the two sites only that they lie in the section, and the positivity of $\mathsf K$.
\end{proof}

\begin{lemma}[Comparison of a functional with its dominant state]\label{lem:hinge}
Let $\mathcal A$ be a real vector space with an element $1$, $\mathsf N$ a seminorm and $\mathsf g$ a linear functional on $\mathcal A$ such that, with $\mathsf R(x)=(1+x)^8-1-8x$,
\[
\mathsf N(1+tG)^8\le1+8t\,\mathsf g(G)+\mathsf R(|t|\mathsf N(G))\quad\text{and}\quad|\mathsf g(G)|\le\mathsf N(G)
\]
for all $G\in\mathcal A$ and all real $t$. Let $\mathfrak s$ be a linear functional with $\mathfrak s(1)=1$ and $\mathfrak s(G)\le Q\,\mathsf N(G)$ for all $G$, where $1\le Q\le1.11$. Then $|\mathfrak s(G)-\mathsf g(G)|\le4.84\,[Q(Q-1)]^{1/2}\,\mathsf N(G)$.
\end{lemma}

\begin{proof}[Proof sketch]
From $(1+u)^{1/8}\le1+u/8$, $\mathsf N(1+tG)\le1+t\mathsf g(G)+\frac18\mathsf R(x)$, $x=|t|\mathsf N(G)$; so $1+t\mathfrak s(G)=\mathfrak s(1+tG)\le Q[1+t\mathsf g(G)+\frac18\mathsf R(x)]$, that is $t[\mathfrak s(G)-\mathsf g(G)]\le(Q-1)(1+x)+\frac Q8\mathsf R(x)$. If $\mathsf N(G)=0$ the bracket vanishes, and if $Q=1$ it vanishes as $x\to0$. Otherwise take $t=\pm x/\mathsf N(G)$ with $x^2=(Q-1)/4.74Q$, so that $x\le0.1446$ and $\mathsf R(x)/8x^2\le4.718$:
\[
|\mathfrak s(G)-\mathsf g(G)|\le\mathsf N(G)\Bigl[\frac{Q-1}x+(Q-1)+4.74\,Qx\Bigr]=\mathsf N(G)\,[Q(Q-1)]^{1/2}\Bigl[2\sqrt{4.74}+\Bigl(\frac{Q-1}Q\Bigr)^{1/2}\Bigr],
\]
and $2\sqrt{4.74}+(0.11/1.11)^{1/2}\le4.67$.
\end{proof}

\begin{proposition}[Boxes of every shape]\label{prop:shape}
Let $\beta\ge0$, $N=2h\ge4$, $c\in(\mathbb Z+\frac12)^3$, $\underline n$ a box with all sides $\ge N$, and $M_\mu=2^{j_\mu}h$, $j_\mu\ge1$, the largest number of this form that is $\le n_\mu$. For every function $A$ of the cell $\mathsf C_h(c)$:
\begin{enumerate}
\item[(a)] $|\langle A\rangle_{\underline n}|\le\exp\bigl\{\frac{\delta_{\rm c}(N)}8\bigl[1+3N\bigl(\frac1{M_1}+\frac1{M_2}+\frac1{M_3}\bigr)\bigr]\bigr\}\,|\!|\!|A|\!|\!|_N\le e^{\frac54\delta_{\rm c}(N)}\,|\!|\!|A|\!|\!|_N$;
\item[(b)] if $A$ is real and $\delta_{\rm c}(N)\le\frac12$, then $|\langle A\rangle_{\underline n}-\langle A\rangle_{(N,N,N)}|\le7.31\,\delta_{\rm c}(N)^{1/2}\,|\!|\!|A|\!|\!|_N$.
\end{enumerate}
Under the criterion at the side $l$ and for $N\ge2l$: $|\langle A\rangle_{\underline n}-\langle A\rangle_{(N,N,N)}|\le2.71\,\theta^{N/4l}\,|\!|\!|A|\!|\!|_N$ for every real $A$. None of these constants depends on the shape of the box.
\end{proposition}

\begin{proof}[Proof sketch]
(a) \emph{Doubling on the one box.} In the direction $1$ put $A_0=A$ and $A_{i+1}=A_i\,\overline{\mathfrak r^{(i)}A_i}$, $\mathfrak r^{(i)}$ the reflection in the plane $x_1=c_1+(2^i-1)h$; $A_i$ is a function of the bonds of the $2^ih$ layers below that plane, and by Lemma~\ref{lem:mir}(a), as long as $2^{i+1}h\le n_1$, $|\langle A_i\rangle_{\underline n}|^2\le\langle A_{i+1}\rangle_{\underline n}$. So $|\langle A\rangle_{\underline n}|\le\langle A_{j_1}\rangle_{\underline n}^{2^{-j_1}}$, and $A_{j_1}$ is a row of $M_1/h$ alternating copies of the cell, which lies below $c_2$ and $c_3$; the same in the directions $2$ and $3$, always on the box $\underline n$, gives $|\langle A\rangle_{\underline n}|\le\langle\mathfrak BA\rangle_{\underline n}^{1/\mathcal N}$, with $\mathfrak BA$ an array of $\mathcal N=M_1M_2M_3/h^3$ copies on a block of $M_1\times M_2\times M_3$ sites.

\emph{The free layers cost their top eigenvalue.} With the time $1$ and $\mathsf Y$ the operator of the slab of $h$ layers of $\mathfrak B_3\mathfrak B_2A$, $Z(\underline m)\langle\mathfrak BA\rangle_{\underline m}=\operatorname{Tr}(\mathsf T'^{\,m_1-M_1}(\mathsf Y^*\mathsf Y)^{M_1/N})$, and $\operatorname{Tr}(\mathsf O_1\mathsf O_2)\le\|\mathsf O_1\|\operatorname{Tr}\mathsf O_2$ for positive operators $\mathsf O_1,\mathsf O_2$, the second of trace class, gives 
\[
Z(\underline m)\langle\mathfrak BA\rangle_{\underline m}\le\lambda_0(m_2,m_3)^{m_1-M_1}Z(M_1,m_2,m_3)\langle\mathfrak BA\rangle_{(M_1,m_2,m_3)};
\]
with the times $2$ and $3$ in turn,
\[
Z(\underline n)\langle\mathfrak BA\rangle_{\underline n}\le\lambda_0(n_2,n_3)^{n_1-M_1}\lambda_0(M_1,n_3)^{n_2-M_2}\lambda_0(M_1,M_2)^{n_3-M_3}\,Z(\underline M)\langle\mathfrak BA\rangle_{\underline M}.
\]
On the block $\underline M$, whose sides lie in $N\mathbb N$, $\operatorname{Tr}(\mathsf Y^*\mathsf Y)^q\le(\operatorname{Tr}\mathsf Y^*\mathsf Y)^q$ in each direction gives $Z(\underline M)\langle\mathfrak BA\rangle_{\underline M}\le[Z(N,N,N)\,|\!|\!|A|\!|\!|_N^8]^{M_1M_2M_3/N^3}$. \emph{The count.} Write $\log\lambda_0(\mathsf A,\mathsf B)=\mathsf A\mathsf BP+\varepsilon(\mathsf A,\mathsf B)$ and $F=|\cdot|P+\delta$: the terms with $P$ cancel, and with $\varepsilon(\mathsf A,\mathsf B)\le(\mathsf A+\mathsf B)\delta_{\rm c}(N)/N^2$ (Lemma~\ref{lem:seq}(g)), $n_\mu<2M_\mu$ and $\delta(\underline n)\ge0$, the remaining exponent is at most $3\frac{\delta_{\rm c}(N)}{N^2}(M_1M_2+M_1M_3+M_2M_3)+\frac{M_1M_2M_3}{N^3}\delta_{\rm c}(N)$; divided by $\mathcal N=8M_1M_2M_3/N^3$ this is the exponent of (a), and $M_\mu\ge N$. (b) Lemma~\ref{lem:hinge} with $\mathcal A$ the real functions of the cell, $\mathsf N=|\!|\!|\cdot|\!|\!|_N$, $\mathsf g=\langle\cdot\rangle_{(N,N,N)}$, $\mathfrak s=\langle\cdot\rangle_{\underline n}$ and $Q=e^{\frac54\delta_{\rm c}(N)}$: by Lemma~\ref{lem:mir}(b), $|\!|\!|1+tG|\!|\!|_N^8$ is the sum over $E\subset\{0,1\}^3$ of $t^{\#E}$ times the mean of the copies of $G$ over $E$, which is $1$ for $E=\emptyset$, $\langle G\rangle_{(N,N,N)}$ for the eight one-element $E$, and at most $|\!|\!|G|\!|\!|_N^{\#E}$ else; and $\mathfrak s(G)\le Q\mathsf N(G)$ is (a).

So, if $\delta_{\rm c}(N)\le0.08348$, which suffices for $Q\le1.11$, the difference is at most $4.84\,[Q(Q-1)]^{1/2}|\!|\!|A|\!|\!|_N\le5.86\,\delta_{\rm c}(N)^{1/2}|\!|\!|A|\!|\!|_N$. For $0.08348<\delta_{\rm c}(N)\le\frac12$, (a) and Lemma~\ref{lem:mir}(b) bound it by $(e^{\frac54\delta_{\rm c}(N)}+1)|\!|\!|A|\!|\!|_N$, and $(e^{5\delta/4}+1)\delta^{-1/2}$ decreases on $(0,\frac12]$ and is $\le7.31$ at $0.08348$. Under the criterion $\delta_{\rm c}(N)<0.1324\,\theta^{N/2l}\le\frac12$ (Proposition~\ref{prop:rec}(d)), and $7.31\cdot0.1324^{1/2}\le2.71$.
\end{proof}

\begin{lemma}[The excess controls every periodic state]\label{lem:exc}
Let $\beta>0$, let $N=2\mathsf L^j$ be the side of one of the tori of the chessboard estimate (Proposition~\YMMref{86.3}{681}), and let $B$ be the set of the bonds with both ends in one of the eight closed cubes of side $N/2$, with faces through midpoints of bonds, that tile the torus $(N,N,N)$. For every periodic Gibbs state $s$ the law of $s$ on $B$ has a density $h$ with respect to the law of the torus $(N,N,N)$ on $B$, with $\int h^{8/7}\le e^{\delta_{\rm c}(N)/7}$, and
\[
\sup_A\bigl|s(A)-\langle\mathbf 1_A\rangle_{(N,N,N)}\bigr|\le\tfrac12(\sqrt\epsilon+\epsilon),\qquad\epsilon=22.3\bigl(e^{\delta_{\rm c}(N)/7}-1\bigr),
\]
the supremum over the Borel events $A$ of the bonds of $B$.
\end{lemma}

\begin{proof}[Proof sketch]
\emph{One board for all sides.} Let $N'=2\mathsf L^{j'}$, $j'\ge j$, be a side of the family. The cubes of side $N/2$, with faces through midpoints of bonds, that are the translates of the closed cube of $B$ by the multiples of $N/2$ in the three directions tile the torus $(N',N',N')$, $2N'/N$ of them a side, an even number, as $N'/N$ is a power of the odd number $\mathsf L$: these \emph{cells} form one board, and the planes between neighbouring cells are its walls. Number the cells by $k\in\{0,\dots,2N'/N-1\}^3$, the cell of $B$ being $k=0$. For a bounded Borel function $f$ of the bonds of $B$ the copy $f^{[k]}$ is $f$ composed, for $\mu=1,2,3$, with the reflections in the $k_\mu$ walls of the direction $\mu$ that carry the cell $0$ step by step to the cell $k$, and complex conjugated if $k_1+k_2+k_3$ is odd; it is a function of the bonds of the cell $k$. Put
\[
\|f\|_{N'}=\Bigl\langle\prod_kf^{[k]}\Bigr\rangle_{(N',N',N')}^{(N/2N')^3},
\]
the product over all $(2N'/N)^3$ cells. At $N'=N$ the cell of $B$ is a cell $\mathsf C_{N/2}(c)$, the eight copies are the copies $f^{[\epsilon]}$, $\epsilon\in\{0,1\}^3$, by which $|\!|\!|\cdot|\!|\!|_N$ is defined, and $\|f\|_N=|\!|\!|f|\!|\!|_N$. With the time $1$ the torus $(N',N',N')$ is cut by the walls into $2N'/N$ slabs of $N/2$ layers, and, as in the proof of Proposition~\ref{prop:shape}, with $\mathsf Y$ the operator of the first slab of the product (that of the cells with $k_1=0$), the operator of the next slab is $\mathsf Y^*$, and $Z(N',N',N')\langle\prod_kf^{[k]}\rangle_{(N',N',N')}=\operatorname{Tr}(\mathsf Y^*\mathsf Y)^{N'/N}$, a number $\ge0$ and at most $(\operatorname{Tr}\mathsf Y^*\mathsf Y)^{N'/N}$, where $\operatorname{Tr}\mathsf Y^*\mathsf Y$ is the same expression on the box $(N,N',N')$. The same in the directions $2$ and $3$ gives the chessboard estimate on this board for the two tori (Proposition~\YMMref{86.3}{681}, in the form of Corollary~\YMMref{86.6}{683}):
\[
Z(N',N',N')\,\|f\|_{N'}^{(2N'/N)^3}\le\bigl[Z(N,N,N)\,\|f\|_N^8\bigr]^{(N'/N)^3}.
\]
Its root of the order $(2N'/N)^3$ is
\[
\|f\|_{N'}\le\|f\|_N\exp\Bigl[\Bigl(\frac N2\Bigr)^3\Bigl(\frac{F(N,N,N)}{N^3}-\frac{F(N',N',N')}{N'^3}\Bigr)\Bigr]\le\|f\|_N\,e^{\delta_{\rm c}(N)/8},
\]
since $(N/2N')^3(N'/N)^3=\frac18$, $(N/2N')^3=(N/2)^3/N'^3$ and $F(N',N',N')/N'^3\ge P$. The domination of periodic Gibbs states (Proposition~\YMMref{91.9}{705}), read with this same board and seminorm, gives $|s(f)|\le\liminf_{j'\to\infty}\|f\|_{N'}$. H\"older's inequality and the reflection invariance of the torus law give $\|f\|_N\le\langle|f|^8\rangle_{(N,N,N)}^{1/8}$. So $|s(f)|\le e^{\delta_{\rm c}(N)/8}\|f\|_{L^8}$, and by duality $s$ has on $B$ a density $h$ with $\|h\|_{8/7}\le e^{\delta_{\rm c}(N)/8}$. With $p=\frac87$, $u^p-1-p(u-1)\ge0.045\min\{(u-1)^2,|u-1|^p\}$ for $u\ge0$; integrating at $u=h$ gives $\int_{|h-1|\le1}(h-1)^2+\int_{|h-1|>1}|h-1|^p\le\epsilon$, so $\int|h-1|\le\sqrt\epsilon+\epsilon$, and the total variation is half of it.
\end{proof}

\begin{lemma}[One bond; a bound by the perimeter]\label{lem:bond}
\begin{enumerate}
\item[(a)] For $\kappa\ge0$ let $\langle\cdot\rangle_\kappa$ be the probability measure on $G$ with the density proportional to $e^{\kappa x_0(V)}$ with respect to the Haar measure, $x_0(V)=\frac12\operatorname{tr}V$, and $r(\kappa)=\langle x_0\rangle_\kappa$. Then $\langle V\rangle_\kappa=r(\kappa)\mathbf 1$ and
\[
\frac\kappa{\kappa+4}\le r(\kappa)\le\frac{2\kappa}{2\kappa+3}.
\]
\item[(b)] Let $\beta\ge0$, $s$ a Gibbs state, periodic or not, $C$ a closed lattice path and $B$ a finite set of bonds which $C$ runs through exactly once each, as $b$ or as $\bar b$, no plaquette containing two of them. Then $|s(W_C)|\le(\frac{8\beta}{8\beta+3})^{\#B}\le\exp[-3\#B/(8\beta+3)]$. For a rectangle with $R,T\ge2$ one may take $\#B=2(R+T)-4$.
\end{enumerate}
\end{lemma}

\begin{proof}[Proof sketch]
(a) Conjugation invariance makes $\langle V\rangle_\kappa$ commute with $G$, so it is a multiple of $\mathbf 1$, and the multiple is $r(\kappa)$. Write $V=x_0\mathbf 1+\mathrm i\sum_jx_j\sigma_j$ with $\sum_{j=0}^3x_j^2=1$. By the invariance of the Haar measure under $V\mapsto e^{\mathrm is\sigma_j}V$, the derivatives at $s=0$ of $\int(x_je^{\kappa x_0})(e^{\mathrm is\sigma_j}V)\,dV$ and of $\int(x_jx_0e^{\kappa x_0})(e^{\mathrm is\sigma_j}V)\,dV$ vanish; as the derivatives of $x_0$ and $x_j$ along this path are $-x_j$ and $x_0$, summing over $j=1,2,3$ gives
\[
3\,r(\kappa)=\kappa\langle1-x_0^2\rangle_\kappa,\qquad4\langle x_0^2\rangle_\kappa-1=\kappa\langle x_0(1-x_0^2)\rangle_\kappa .
\]
The right member of the second identity lies in $[0,\kappa\langle1-x_0^2\rangle_\kappa]$, its sign by $V\mapsto-V$; so $(4+\kappa)\langle x_0^2\rangle_\kappa\le1+\kappa$, and the first identity gives $r(\kappa)\ge\frac\kappa3\cdot\frac3{4+\kappa}$; and $3r=\kappa\langle(1-x_0)(1+x_0)\rangle\le2\kappa(1-r)$. (b) As no plaquette contains two bonds of $B$, given the other bonds the weight of the bonds of $B$ factorises: $V(b)$ has the density proportional to $\exp[\beta\mathsf k_b\frac12\operatorname{tr}(V(b)g_b)]$, where the sum of the four staples of $b$ is $\mathsf k_bg_b$, $g_b\in G$, $0\le\mathsf k_b\le4$, as a sum of four unit quaternions. $W_C$ is a sum of products that hold one entry of each $V(b)^{\pm1}$, $b\in B$; by (a) and $V\mapsto Vg^{-1}$ the conditional mean of $V(b)^{\pm1}$ is $r(\beta\mathsf k_b)g_b^{\mp1}$, so the conditional mean of $W_C$ is $\prod_br(\beta\mathsf k_b)$ times half the trace of an element of $G$, of modulus at most $(\frac{8\beta}{8\beta+3})^{\#B}$; the Gibbs equation of $B$ carries this to $s$. In a rectangle with $R,T\ge2$ take all bonds but the first of each side: two bonds of one plaquette are either perpendicular with a common end, which happens on $C$ only at a corner, or parallel at the distance $1$, which does not happen for opposite sides at the distance $\ge2$. Finally $\log(1+y)\ge y/(1+y)$ at $y=3/8\beta$.
\end{proof}

\subsection{Proofs of the results}

\begin{proof}[Proof of Theorem~\ref{thm:one}]
\emph{Step 1 (the cubes converge).} In a box $(t,\Sigma)$ with $t\in\{n,2n\}$ a function $A$ of a cube of $h\le n/2-1$ layers is a function of $h$ consecutive layers, and Proposition~\ref{prop:long} applies with $s\ge n/2$ and one and the same $a_\infty(\Sigma,A)$ for both $t$. Doubling the three sides one after the other, each time with the time in the doubled direction,
\[
\begin{multlined}
|\langle A\rangle_{(n,n,n)}-\langle A\rangle_{(2n,2n,2n)}|\le4\sup|A|\bigl[\mathsf r_{n/2}(n,n)+\mathsf r_{n/2}(2n,n)+\mathsf r_{n/2}(2n,2n)\bigr]\le4 \\
\cdot4.15727\sup|A|\,\mathsf e(n/2)\le1.29\sup|A|\,\theta^{n/2l}
\end{multlined}
\]
by Proposition~\ref{prop:rec}(d). Summing over $n,2n,4n,\dots$ with $\theta^{2^in/2l}\le\theta^{n/2l}\theta^{in/2l}$ gives the second bound of (a). On each finite set of bonds the cube laws form a Cauchy sequence in total variation; the limits are consistent, and Kolmogorov's theorem gives $\nu$.

\emph{Step 2 (Gibbs state, symmetries).} For a finite set $\mathfrak A$ of bonds and a large cube, the action is the sum over the plaquettes that meet $\mathfrak A$ plus a remainder that does not depend on the bonds of $\mathfrak A$; integrating these bonds first shows that every cube law, and so $\nu$, satisfies the Gibbs equation of $\mathfrak A$. A symmetry of the lattice induces a map of each cube of side $2^kl$ that keeps $dU$ and the action, hence each cube law and $\nu$.

\emph{Step 3 (one periodic state).} By Step 2, $\nu$ is periodic. Let $s$ be a periodic Gibbs state and $A$ a local event. For all large $N=2\mathsf L^j$, $A$ is an event of the bonds inside one cell of side $N/2$, and Lemma~\ref{lem:exc}, applied to $s$ and to $\nu$, gives $|s(A)-\nu(A)|\le\sqrt\epsilon+\epsilon$ with $\epsilon=22.3(e^{\delta_{\rm c}(N)/7}-1)$, which tends to $0$ as $\delta_{\rm c}(N)<0.1324\,\theta^{N/2l}$ (Proposition~\ref{prop:rec}(d)). So $s=\nu$. Proposition~\YMMref{103.4}{758} gives reflection positivity at every lattice mirror. By Proposition~\YMMref{106.2}{764}, the ergodic components of $\nu$ are periodic Gibbs states; uniqueness makes them equal to $\nu$, so $\nu$ is ergodic.

\emph{Step 4 (all boxes).} Let the three sides of boxes $\underline n^{(k)}$ tend to infinity. Their laws on $\Omega$ lie in a compact set; every limit point satisfies the Gibbs equations, as in Step 2, and is invariant under all translations, as every box law is; so it is a periodic Gibbs state, that is $\nu$, and the whole sequence converges.

\emph{Step 5 (twisted cubes).} Let $n\ge2\mathsf ql$ be an integer, $j$ the largest integer with $4^j\le n/2\mathsf ql$, $w=2^j$, $l'=l_j=2wl$ and $M=n/l'$, a real number with $w\mathsf q\le M<4w\mathsf q$ and $M^2<2\mathsf qn/l$. Lemma~\ref{lem:vort}(a) from the box $(l',n,l')$ to the cube, with $\kappa=M^2$, and Proposition~\ref{prop:tub}(d) give
\[
z_-(n,n,n)\le\bigl[Me^{-\varsigma w^2+0.01655\,M\theta^{2w}}\bigr]^{M^2},\qquad M^2\varsigma w^2=\sigma^+n^2 ,
\]
and $0.01655\,M\theta^{2w}<4\mathsf q\cdot0.01655\,we^{-2w\Lambda}\le\frac{4\cdot0.01655}e(\frac1{2\Lambda}+\frac{1.67}{4\Lambda^2})\le0.855$ and $\log M<\frac12\log\frac{2\mathsf qn}l$: the first bound of (d). The bound for all seven $z\ne0$ then follows from Lemma~\ref{lem:tr}(c).

\emph{Step 6 (flux sectors, static charges).} With the time $1$ over $n\times n$ and $e\in\{01,11\}$, $\Lambda_e^n\le W_n(01;0)+W_n(11;0)=a_-(n,n,n)=z_-\lambda_0^n(1+\mathsf r_n(n,n))$; for $e=10$ the same holds because $Z^{(1,0,0)}=Z^{(0,1,0)}$ on the cube (Lemma~\ref{lem:tr}(iii)). So $\log(\lambda_0/\Lambda_e)\ge\frac1n[-\log z_-(n,n,n)-\log(1+\mathsf r_n(n,n))]$, and $\log(1+\mathsf r_n(n,n))$ is at most $\mathsf r_n(n,n)\le0.001615\,\theta^{n/l}\le0.0017$ at $n=2^il$, and in general, with $n_0=2^il\le n<2n_0$, at most $\delta_{\rm c}(n)\le8\delta_{\rm c}(n_0)\le8\cdot0.01655\cdot0.792279\le0.105$ (Lemma~\ref{lem:seq}(f), Proposition~\ref{prop:rec}(d)). With Step 5 this is the first statement of (e); the second is Lemma~\ref{lem:two}(d) on the cube with $n_3=n$, and Step 5. For any two sites $o\ne r$ it is Lemma~\ref{lem:two}(e) on the cube, with the time $1$ or $2$, whichever direction of the section realises $|r-o|_\infty$, and Step 5; for fixed $o$, $r$ and $n=2^il$, $i\to\infty$, the factor $[8(1+\mathsf r_n(n,n))]^{1/n}$ tends to $1$ and the correction $\frac{2\mathsf qR}{ln}(\frac12\log\frac{2\mathsf qn}l+0.86)$ to $0$, while $R=|r-o|_\infty$ once $n>2|r-o|_\infty$.

\emph{Step 7 (rectangles).} By the symmetries of $\nu$ and of the cube laws it suffices to take $C$ in a $(1,3)$-plane with the side $R$ along $e_1$. On the cube of an integer side $n\ge2\mathsf ql$ with $R\le n/2$, $T\le n-1$, Lemma~\ref{lem:two}(d) gives $0\le\langle W_C\rangle_{(n,n,n)}\le8^{T/n}z_-(n,n,n)^{RT/n^2}$, and Step 5 gives the bound on the cube. Along $n=2^kl$ the left side tends to $\nu(W_C)$ by Step 1, while $8^{T/n}\to1$ and $(\log n)/n\to0$: this is \eqref{eq:area}. The bound by the perimeter is Lemma~\ref{lem:bond}(b).

\emph{Step 8 (closed paths).} Let $C$ and $(\rho',\rho)$ be as in (c); by symmetry $(\rho',\rho)=(1,3)$. On the cubes $n=2^kl$, which are even, Lemma~\ref{lem:sev}(c) applies for large $k$: the sum of the $\ell_j$ is $\mathsf A_{13}(C)$, $\sum_j(\mathsf v_j+1)/n\to0$, the thermal factors tend to $1$, and $z_-(n,n,n)^{\mathsf A/n^2}\le e^{-\sigma^+\mathsf A+o(1)}$ by Step 5. For the general form let $\rho=3$: by the symmetry of the cube that exchanges the directions $1$ and $2$, Lemma~\ref{lem:sev}(b) holds with the direction $2$ in the place of the direction $1$, the two twisted pairs having one quotient on the cube (Lemma~\ref{lem:tr}(c)), and in the proof of Lemma~\ref{lem:sev}(c) each norm $\|\hat{\mathsf T}_{D_j}\|$ is bounded by (b) on its own; so in each gap $\ell_j$ may be taken as the larger of the two counts $\ell^1_j(C)$, $\ell^2_j(C)$, whose sum over the gaps is $\mathsf A_3(C)$, and Lemma~\ref{lem:sev}(d) carries all this to $D$, with $d^{\,c_3(C)-1}$ in front and $\sum_j\mathsf v_j\log d/n\to0$.

For a staircase loop the senses of its runs of the direction $\rho$ are up in $T$ consecutive gaps and then down in the same gaps, so $c_\rho=1$, and $\mathsf A_{\rho'\rho}(C)=\sum_j|r_j-o_j|$. For a simple closed curve of a coordinate plane each passage from down to up and each from up to down holds at least one bond of the other direction, and between two such passages there is at least one bond of the direction $\rho$; so $|C|\ge4c_\rho(C)$ and $2^{c_\rho-1}\le\frac12e^{(\log2/4)|C|}$.

\emph{Step 9 (gap).} Let $N=N_1\le N_2$ (by (T1)), $n_0\le N<2n_0$, and $i$ with $t_i\le N_2<2t_i$ in Proposition~\ref{prop:slab}(b): for $a\ge1$, $(\lambda_a/\lambda_0)^{t_i}\le\mathsf r_{t_i}(N_2,N)\le0.28\,\chi'^{2^i}\le\chi'^{t_i/n_0}$, so $\log(\lambda_0/\lambda_a)\ge\frac1{n_0}\log\frac1{\chi'}\ge\frac\Lambda l+\frac{1.021}{n_0}$. The vacuum is simple and its eigenvector lies in the sector $00$ (Lemma~\ref{lem:tr}); the tops of the sectors with electric flux are among the $\lambda_a$, $a\ge1$. For $N=n=2^kl$ and an integer $N_2\ge n$, with $t_i=2^in\le N_2<2t_i$, Proposition~\ref{prop:slab}(a), (c) give $\log(1+\mathsf r_{t_i}(N_2,n))\le\frac{N_2}{t_i}\delta^{(n)}(t_i,t_i)\le0.3652\,\chi^{2^i}$, so $(\lambda_a/\lambda_0)^{t_i}\le0.367\,\chi^{2^i}\le\chi^{t_i/n}$ and $\log(\lambda_0/\lambda_a)\ge\frac1n\log\frac1\chi\ge\frac\Lambda l+\frac{3.5065}n$. Over $n\times n$, $(\lambda_a/\lambda_0)^n\le\mathsf r_n(n,n)\le0.001615\,\theta^{n/l}$, and $\log\frac1{0.001615}\ge6.42$. Finally, for $N_1,N_2\ge n$, Lemma~\ref{lem:seq}(f) gives $\log(1+\mathsf r_n(N_1,N_2))\le\delta(n,N_1,N_2)\le\frac{N_1N_2}{n^2}\delta_{\rm c}(n)\le0.01655\,\frac{N_1N_2}{n^2}\theta^{n/l}=:u\le\log2$, and $e^u-1\le u/\log2$, $\log2/0.01655\ge41.8$.

\emph{Step 10 (flux lines of dyadic length).} $2\sigma^+N^2/\log\frac1\chi=2\varsigma w^2/(-\log0.03+2w\Lambda)$ increases in $w$, is $\ge3.83$ at $w=2$ and $<2\mathsf qw=\mathsf qN/l$: the bounds on $\hat N$. With the time $1$ over $(N_2,N)$, Lemma~\ref{lem:tr}(d) gives $\Lambda_-^t\le a_-(t,N_2,N)=z_-(t,N_2,N)\lambda_0^t(1+\mathsf r_t(N_2,N))$ for every $t$. Take $t=t_i=2^iN$ with $i$ the largest integer such that $\mathsf M=N_2/t\ge\hat N/2$, so that $\hat N/2\le\mathsf M<\hat N$. The first inequality of Lemma~\ref{lem:vort}(b) at the time $N_2$, Proposition~\ref{prop:slab}(c), $\mathsf r_{N_2}(t,N)\le\mathsf r_t(t,N)^{\mathsf M}\le e^{-2^i\sigma^+N^2}$ and $\log(1+\mathsf r_t(N_2,N))\le\mathsf M\delta^{(N)}(t,t)\le0.1826\,\mathsf M\chi^{2^i}$ give
\[
\log\frac{\lambda_0}{\Lambda_-}\ge\frac1t\Bigl[2^i(\sigma^+N^2-\mathsf g)-\log\mathsf M-0.1826\,\mathsf M\chi^{2^i}\Bigr]\ge\sigma^+N-\frac{2\log\hat N+1.86}N,
\]
as $\hat N\chi\le0.7936$. As $N_2\to\infty$, $i\to\infty$, and the bounded last two terms are divided by $2^iN$: the lower limit. For $\max\{N,2\mathsf ql\}\le N_2<\hat NN/2$ let $w'$ be the largest power of two with $w'\le w$ and $2\mathsf qw'^2l\le N_2$, and $l'=2w'l\le N$, so that $N_2\ge\mathsf ql'^2/2l$; Proposition~\ref{prop:tub}(e) with $N_3=N$, $n_2=N_2$ and $M=N_2/l'$ applies. If $w'<w$, then $8\mathsf qw'^2l>N_2$, so $N/l'^2<2\mathsf qN/N_2l\le2\mathsf q/l$ and $M<4\mathsf qw'$; if $w'=w$, then $N/l'^2=1/N\le2\mathsf q/l$ and $M<\hat N/2<4\mathsf qw$. In both cases $\log M<\log\frac{2\mathsf qN}l$ and $0.01655\,M\theta^{l'/l}<4\mathsf q\cdot0.01655\,w'e^{-2w'\Lambda}\le0.855$, which gives the second floor. For $N\le N'<2N$: with the time $1$, Lemma~\ref{lem:vort}(a) in the third side gives $\Lambda_-(N_2,N')^t\le a_-(t,N_2,N')\le a_-(t,N_2,N)^{N'/N}$, and $a_-(t,N_2,N)^{1/t}\to\Lambda_-(N_2,N)$; so $\log\Lambda_-(N_2,N')\le\frac{N'}N\log\Lambda_-(N_2,N)$. Also $\log\lambda_0(N_2,N')\ge N_2N'P$ and $\log\lambda_0(N_2,N)\le\frac1NF(N,N_2,N)\le N_2NP+\frac{N_2}{N^2}\delta_{\rm c}(N)$ (Lemma~\ref{lem:seq}(f)). Hence $\log\frac{\lambda_0}{\Lambda_-}(N_2,N')\ge\frac{N'}N[\log\frac{\lambda_0}{\Lambda_-}(N_2,N)-\frac{N_2}{N^2}\delta_{\rm c}(N)]$, and the first floor, $\delta_{\rm c}(N)\le0.01655\,\theta^{N/l}$ and $N'/N<2$ give the third.

\emph{Step 11 (flux lines of every integer length; the slabs).} Let $N\ge4l$. Proposition~\ref{prop:lines} gives $t_i\varepsilon_{\rm v}(t_i,N)\le e^{-E'_i}$ with $2^{-i}E'_i\ge\mathsf E_0-\mathsf g'$. As in Step 10, with $i$ the largest integer such that $\mathsf M=N_2/t_i\ge\hat N'/2$, the first inequality of Lemma~\ref{lem:vort}(b) at the time $N_2$, with $\mathsf r_{N_2}(t_i,N)\le\chi_0^{2^i\hat N'/2}\le e^{-2^i\mathsf E_0}$, gives $z_-(t_i,N_2,N)\le\hat N'e^{-E'_i}$, and Proposition~\ref{prop:slab}(b) gives $\log(1+\mathsf r_{t_i}(N_2,N))\le\mathsf M\delta^{(N)}(t_i,t_i)\le\frac14\xi'_i$; so $\log(\lambda_0/\Lambda_-)\ge\frac1{t_i}[E'_i-\log\hat N'-\frac14\xi'_i]\ge(\mathsf E_0-\mathsf g')/n_0=\sigma^+N-\kappa\mathsf g'/N$. The penultimate inequality uses the unused positive recurrence tail displayed in the proof of Proposition~\ref{prop:lines}.

As $\kappa\mathsf g'=\kappa^2\log\hat N'+\kappa[3.22+(3\kappa+0.68)\mathsf h]$ and $\mathsf h\le0.4274\kappa+0.0208$, $\kappa\mathsf g'$ increases in $\kappa$, and at $\kappa=2$ it is at most $4\log\hat N'+18.14$. The window $\max\{N,2\mathsf ql\}\le N_2<\hat N'n_0/2$ is treated as in Step 10, with $l'=2w'l\le n_0$. For the closed form, $\hat N'n_0/2<(\mathsf qN/l+2)N/2\le(\frac{\mathsf q}2+\frac14)N^2/l$ and $\hat N'<(\mathsf q+\frac12)N/l$, as $N\ge4l$; so $C(\theta)N^2/l\ge\hat N'n_0/2$ and $4\log\hat N'+18.2\le4\log\frac Nl+4\log(\mathsf q+\frac12)+18.2\le C(\theta)(1+\log\frac Nl)$; and $4\log8.673+18.2\le26.85$ at $\mathsf q\le8.173$. If $w\ge4$ and $\tau\le0.087$, the same road, with the slab of the thickness $N$ started from $\mathsf e^{(N)}(n_0)\le2\delta(n_0,2n_0,N)\le8\cdot0.01655\,\tau\le0.0116$, where $\omega(Y_1(x))\le2.06x^2$, so that $2.06\,\mathsf e^{(N)}(t_i)\le(0.2728\,\tau)^{2^i}$, and with the start from Proposition~\ref{prop:tub}(d) at $l'=n_0$ carried to the thickness $N$ by Lemma~\ref{lem:vort}(a) at the real factor $N/n_0$, gives the floor with $\hat N''$ and $10.8$. (g3) is Lemma~\ref{lem:road} with $L_2=N_2$, $L_3=N$ and $\phi\ge\sigma^+N-(4\log\hat N'+18.2)/N$.

\emph{Step 12 (boxes of every shape, cylinders, twisted boxes).} A cube of $N/2$ layers is a cell $\mathsf C_{N/2}(c)$, and all box laws are invariant under translations. For a real function $A$ of the cell, Proposition~\ref{prop:shape} gives $|\langle A\rangle_{\underline n}-\langle A\rangle_{(N,N,N)}|\le2.71\,\theta^{N/4l}|\!|\!|A|\!|\!|_N$ for every box with sides $\ge N$. If $A$ is continuous, let the sides of a second box tend to infinity: by Step 4, $|\nu(A)-\langle A\rangle_{(N,N,N)}|\le2.71\,\theta^{N/4l}|\!|\!|A|\!|\!|_N$, so $|\langle A\rangle_{\underline n}-\nu(A)|\le5.42\,\theta^{N/4l}|\!|\!|A|\!|\!|_N\le5.42\,\theta^{N/4l}\sup|A|$ by Lemma~\ref{lem:mir}(b). For a bounded real Borel $A$, take continuous $A_k$ with $|A_k|\le\sup|A|$ and $A_k\to A$ in $L^8$ of the mean of the images of $\langle\cdot\rangle_{(N,N,N)}$, $\langle\cdot\rangle_{\underline n}$ and $\nu$ on the bonds of the cell; both sides of the bound pass to the limit, the seminorm because $|\!|\!|A_k-A|\!|\!|_N\le\langle|A_k-A|^8\rangle_{(N,N,N)}^{1/8}$ (Lemma~\ref{lem:mir}(b)).

For an event take $A=\mathbf 1_E-\frac12$. For the cube of side $m\ge2l$ take $N=2\lfloor m/2\rfloor\le m$. The cylinder over $\Sigma$ is the limit of Proposition~\ref{prop:long}, and the bound passes to it; with $\Sigma=(n,n)$ and $M\to\infty$, Lemma~\ref{lem:vort}(a) in the third side gives $z_-(n,n,Mn)\le[z_-(n,n,n)\lambda_0^n(1+\mathsf r_n(n,n))]^M\lambda_0^{-Mn}=[(1+\mathsf r_n(n,n))z_-(n,n,n)]^M$, and Lemma~\ref{lem:two}(d) on the box $(n,n,Mn)$ gives the area bound on the cylinder, as $8^{T/Mn}\to1$. Twisted boxes: let the columns recede. For a finite set $\mathfrak A$ of bonds the plaquettes that meet $\mathfrak A$ carry no twist in all but finitely many members, so every limit point satisfies the Gibbs equations, as in Step 2.

The translate by $e_\mu$ of a twisted box law is the twisted law with the columns moved by one step, and by the proof of Lemma~\ref{lem:tr}(iii) this is the image of the law under $U(b)\mapsto-U(b)$ on a set of bonds next to the columns; a fixed local function depends on none of these bonds in all but finitely many members. So every limit point is a periodic Gibbs state, that is $\nu$; the two equations are linear in the law, so they hold for finite mixtures. If a column passes through $p_0$ for all members, one flip moves it off $p_0$ and changes the sign of $\operatorname{tr}\partial U(p_0)$, and the column so moved recedes after further flips that do not touch $p_0$; so the law of $\operatorname{tr}\partial U(p_0)$ tends to that of $-\operatorname{tr}\partial U(p_0)$ under $\nu$, which differs because $\nu(\operatorname{tr}\partial U(p_0))=2\nu(W_p)>0$ by Corollary~\ref{cor:pot}(a).

\emph{Step 13 (the pressure).} By Theorem~\ref{thm:crit}(b), $\mathsf e(2l)\le\theta^2/12.9\le0.0615<0.069$ at $\beta$. Each $F(\underline n;\beta')$ is a smooth function of $\beta'$, so there is an open interval $\mathsf U$ containing $\beta$ on which $\mathsf e(2l;\beta')\le0.069$: the criterion holds on $\mathsf U$ at the side $2l$ with $x_0=0.069$, so that what is proved above holds at each $\beta'\in\mathsf U$, with the state $\nu_{\beta'}$ and with $\delta_{\rm c}(N;\beta')<0.1324\cdot0.8901^{N/4l}$ for $N\ge4l$, uniformly on $\mathsf U$. Let $N_m=2\mathsf L^m$ and $f_m(\beta')=F(N_m,N_m,N_m;\beta')/N_m^3$: these are smooth, tend to $P$ pointwise, and $f_m'(\beta')=-3\langle a_p\rangle_{(N_m,N_m,N_m)}$ by the invariance of the cube under translations and permutations of the axes. As $0\le a_p\le2$ is a function of one cell of the torus, Lemma~\ref{lem:exc} gives $|\langle a_p\rangle_{(N_m,N_m,N_m)}-\nu_{\beta'}(a_p)|\le\sqrt{\epsilon_m}+\epsilon_m$, with $\epsilon_m\to0$ uniformly on $\mathsf U$. So the $f_m'$ converge uniformly on $\mathsf U$ to $-3\nu_{\beta'}(a_p)$, which is therefore continuous, and $P(\beta_2)-P(\beta_1)=\lim_m\int_{\beta_1}^{\beta_2}f_m'=\int_{\beta_1}^{\beta_2}(-3\nu_{\beta'}(a_p))\,d\beta'$. Finally $\nu(a_p)=1-\nu(W_p)\le1-\frac\beta{4(\beta+1)}$ by Corollary~\ref{cor:pot}(a).
\end{proof}

\begin{proof}[Proof of Theorem~\ref{thm:crit}]
(a) is Proposition~\ref{prop:rec}(a); if $\beta>0$, $\mathsf e(l_0)\le X(l_0)\le\frac1{20}\le0.069$, and $12.9/20=0.645$. (b) By Proposition~\ref{prop:rec}(b), (c), $12.9\,\mathsf e(2n)\le(12.9\,\mathsf e(n))^2$ as long as $\mathsf e(n)\le0.069$, and the induction keeps $\mathsf e(2^kl)\le\theta^{2^k}/12.9$. For $k\ge3$, $\mathsf e(2^{k-1}l)\le\theta^4/12.9\le0.04866\le\frac1{20}$, so
\[
X(2^kl)\le17.71\,\mathsf e(2^{k-1}l)^2\le\frac{17.71}{166.41}\,\theta^{2^k}\le0.10643\,\theta^{2^k}\le0.10643\cdot0.394013 .
\]
$\Phi(\Psi(x))=\frac1{20}$ at $x=0.065062\ldots$ and $\Phi(x)=\frac1{20}$ at $x=0.052727\ldots$; both functions increase, $X(2l)\le\Phi(x_0)$, $X(4l)\le\Phi(\Psi(x_0))$, and $\Phi(\Psi^k(x_0))$ does not increase in $k$ because $\Psi(x)\le x$ there. (c) $\mathsf m_{\mathsf e}(2n)\ge\mathsf m_{\mathsf e}(n)$ is $12.9\,\mathsf e(2n)\le(12.9\,\mathsf e(n))^2$. For $k\ge3$, $\mathsf e(n)\le X(n)\le\frac1{20}$ gives $X(2n)\le17.71\,\mathsf e(n)^2\le18X(n)^2$, so $\mathsf m_X$ does not decrease. $X(n)\ge\mathsf e(n)$ gives $\mathsf m_X(n)\le\mathsf m_{\mathsf e}(n)-\frac1n\log\frac{18}{12.9}$, with $\log\frac{18}{12.9}\ge0.3331$; $X(n)\le17.71\,\mathsf e(n/2)^2$ gives $18X(n)\le1.91563\,(12.9\,\mathsf e(n/2))^2$, so $\mathsf m_X(n)\ge\mathsf m_{\mathsf e}(n/2)-\frac1n\log1.91563$, with $\log1.91563\le0.6501$. Both rates being monotone and squeezed, they have one limit, and $\mathsf m_{\mathsf e}(l)\ge\Lambda/l$ is the criterion. (d) $Y_2(0.0740538)\approx0.3156<\log2$; for $0<x<0.0740538$ the iterates $\Psi^k(x)$ decrease, and to $0$, since a positive limit would be a fixed point below $0.0740538$; so $\mathsf e(2^kl)\le\Psi^k(\mathsf e(l))\to0$ and $X(2^{k+1}l)\le\Phi(\Psi^k(\mathsf e(l)))\to0$; and $\beta>0$ because $\mathsf e(l)>0$.
\end{proof}

\begin{proof}[Proof of Theorem~\ref{thm:phys}]
Apply Theorems~\ref{thm:one} and~\ref{thm:crit} at $l=l_*L^K$ and $\beta=L^K\mathsf b$, with the one $x_0$ of the uniform criterion; then $\theta$, $\Lambda$, $\mathsf q$, $C(\theta)$ do not depend on $K$. By \eqref{eq:sigma}, $\sigma^+l^2=\frac12\log\frac1\theta+0.4175$ depends on $\theta$ alone. A rectangle of physical sides $R,T$ has $RL^K\times TL^K$ bonds, and $\sigma^+\cdot RL^K\cdot TL^K=\frac{\sigma^+l^2}{l_*^2}RT=\sigma_*RT$. For $R,T\ge2a_K$, Lemma~\ref{lem:bond}(b) with $\#B=2(R+T)/a_K-4$ and $\beta=\mathsf b/a_K$ gives $\nu_K(W_C)\le\exp[-3(|C|_{\rm ph}-4a_K)/(8\mathsf b+3a_K)]$, and the smaller of two bounds is at most their geometric mean: (a). For $n=\ell/a_K$: $\sigma^+n^2=\sigma_*\ell^2$, $n/l=\ell/l_*$, $a_K^{-1}\sigma^+n=\sigma_*\ell$, $a_K^{-1}/l=1/l_*$ and $a_K^{-1}/n=1/\ell$; this turns Theorem~\ref{thm:one}(d), (e), (f) and (g) into (b), (c) and (d), with $\hat N\le\mathsf q\ell/l_*+1$ and $C(\theta)\le27$. (e) is Theorem~\ref{thm:crit}(b) at $k=3$, with $18\cdot\frac{17.71}{166.41}\le1.91563\le1.9157$, and Theorem~\ref{thm:crit}(a); finally $4\cdot0.4757=1.9028$.
\end{proof}

\begin{proof}[Proof of Corollary~\ref{cor:pot}]
\emph{Step 1 (symmetry).} The symmetries of $\nu$ carry a rectangle to every other one with the same sides, or with the sides exchanged and run backwards.

\emph{Step 2 (two Hankel forms).} At a site mirror orthogonal to $e_\mu$ that contains a site $x$, let $\gamma_h$ run from $x$ up $h$ steps along $e_\mu$, $R$ steps along $e_\lambda$, and down to $x+Re_\lambda$ ($\gamma_0$ lies in the mirror). Then $\sum_{j,k}U(\gamma_{h'})_{jk}\overline{U(\mathfrak r\gamma_h)_{jk}}=2W$ of a rectangle $R\times(h+h')$, and the reflection positivity of $\nu$ at this mirror (Proposition~\YMMref{103.4}{758}; $\nu$ is a periodic Gibbs state), applied to $A_{jk}=\sum_hc_hU(\gamma_h)_{jk}$, gives $\sum_{h,h'}c_hc_{h'}m(R,h+h')\ge0$. At a bond mirror, take the paths $\gamma'_h$ from a site $x_p$ next to the mirror up $h$ steps, along $R$ steps and down to $x_q=x_p+Re_\lambda$; then $\sum\overline{(\mathsf z_p)_{ji}}(\mathsf z_q)_{kl}A_{jk}\overline{\mathfrak rA_{il}}$ with $A_{jk}=\sum_hc_hU(\gamma'_h)_{jk}$ is $2\sum_{h,h'}c_hc_{h'}W$ of rectangles $R\times(h+h'+1)$. Its mean is $\ge0$ on the cubes of sides $2^kl$, which are even, by Lemma~\ref{lem:mir}(c) after a translation that brings the mirror between the layers $0$ and $1$, and it passes to $\nu$ by Theorem~\ref{thm:one}(a): $\sum_{h,h'}c_hc_{h'}m(R,h+h'+1)\ge0$.

\emph{Step 3 (log-convexity).} The $2\times2$ minors of the two forms give $m(R,T)^2\le m(R,T-1)m(R,T+1)$ for all $T\ge1$.

\emph{Step 4 ($w>0$).} Given the other bonds, $U(b)$ has the density proportional to $\exp[\beta\mathsf k\frac12\operatorname{tr}(U(b)g)]$, $\mathsf kg$ the sum of the four staples of $b$, $0\le\mathsf k\le4$; so, by the Gibbs equation and Lemma~\ref{lem:bond}(a), $4w=\nu(\mathsf k\,r(\beta\mathsf k))\ge\nu(\beta\mathsf k^2/(\beta\mathsf k+4))\ge\beta\nu(\mathsf k^2)/(4\beta+4)$. Two different staples of $b$ are mirror images at a mirror that contains $b$, a site mirror if they lie in one plane and a diagonal mirror if they lie in two; so $\nu(\mathsf k^2)=4+\sum_{s\ne s'}\nu(\frac12\operatorname{tr}ss'^{-1})\ge4$ by reflection positivity, and $w\ge\beta/(4(\beta+1))$.

\emph{Step 5 (positivity).} From $m(1,0)=1$, $m(1,1)=w>0$ and Step 3, $m(1,T)>0$ for all $T$; then $m(R,1)=m(1,R)>0$ and $m(R,T)>0$. As $T\mapsto-\log m(R,T)$ is concave and vanishes at $0$, $m(R,T)\ge m(R,1)^T\ge w^{RT}$; the upper bound of (a) is Theorem~\ref{thm:one}(b).

\emph{Step 6 (the potential).} The increments of a concave sequence do not increase; their limit is their infimum and also the limit of $-\frac1T\log m(R,T)$; summing the increments from $T_1$ to $T_2$ gives $m(R,T_2)\le e^{-(T_2-T_1)V(R)}m(R,T_1)$; if a number $v$ has this property for infinitely many $T_2$, dividing by $T_2$ gives $v\le V(R)$; and the increasing slopes of the convex function $\log m(R,\cdot)$ give the inequality of four rectangles.

\emph{Step 7 (the measure).} Let $\Lambda$ be the linear functional on the real polynomials with $\Lambda(\lambda^T)=m(R,T)$. Step 2 says $\Lambda(p^2)\ge0$ and $\Lambda(\lambda p^2)\ge0$. With $s_r=\Lambda(\lambda^rp^2)\ge0$, Schwarz's inequality gives $s_r^2\le s_0s_{2r}$, and, $|m|\le1$ bounding every $s_r$, iteration gives $s_2\le s_0$ and then $s_1\le s_0$: $\Lambda((1-\lambda)p^2)\ge0$; and $\Lambda(\lambda(1-\lambda)p^2)=\Lambda((1-\lambda)(\lambda p)^2)+\Lambda(\lambda((1-\lambda)p)^2)\ge0$. Writing $\lambda^T(1-\lambda)^k$ as $\lambda^\epsilon(1-\lambda)^{\epsilon'}q^2$ with $\epsilon,\epsilon'\in\{0,1\}$ gives $\sum_i(-1)^i\binom kim(R,T+i)\ge0$ for all $k,T\ge0$: the sequence is completely monotone, and Hausdorff's moment theorem gives the one measure $\varrho_R$. Since $m(R,T)\le(\max\operatorname{supp}\varrho_R)^T$ and $m(R,T)\ge(\lambda_*-\varepsilon)^T\varrho_R([\lambda_*-\varepsilon,1])$ for $\lambda_*=\max\operatorname{supp}\varrho_R>0$, $m(R,T)^{1/T}\to\lambda_*$, so $\lambda_*=e^{-V(R)}$; for $T\ge1$ the point $0$ gives nothing, and the substitution $\lambda=e^{-E}$ gives the representation, with the mass $\varrho_R((0,1])\ge m(R,1)>0$.

\emph{Step 8.} (d) follows from (a): $V(R)\ge\sigma^+R$ from the upper bound, $V(R)\le-R\log w$ from the lower one. (e) At $l=l_*L^K$, $\sigma^+=\sigma_*a_K^2$, so $a_K^{-1}V(R/a_K)\ge a_K^{-1}\sigma^+R/a_K=\sigma_*R$; and $M_K(R,T_2)/M_K(R,T_1)=e^{2\kappa(T_2-T_1)}m(R/a_K,T_2/a_K)/m(R/a_K,T_1/a_K)$, to which (c) applies; in the inequality of four rectangles the factors $e^{2\kappa(T'-T)}$, raised to the powers $1/(T'-T)$, give one and the same factor on both sides. Along a sequence, $M_j(R,T_2)\le e^{-(T_2-T_1)v_j(R)}M_j(R,T_1)$ at each $j$, as just shown with $\kappa=\kappa_j$; so $v_j(R)\le\frac1{T_2-T_1}\log\frac{M_j(R,T_1)}{M_j(R,T_2)}$, whose right member converges, and $v_j(R)\ge B$ for infinitely many $j$ gives $B\le\limsup_jv_j(R)$; and $V_{K_j}(R)\ge\sigma_*R$ gives $v_j(R)\ge\sigma_*R-2\kappa_j$.

The symmetry, the inequality of three rectangles, which is the log-convexity of (b) in the form of chords, and that of four rectangles hold for each $M_j$, the factors $e^{2\kappa_j(R+T)}$ cancelling in both inequalities, and pass to the limits. With $\mathsf F(T)=-\log\mathsf W(R,T)$, the inequality of three rectangles says that the slope $(\mathsf F(T')-\mathsf F(T_1))/(T'-T_1)$ does not increase in $T'>T_1$; so it has a limit in $[-\infty,\infty)$ as $T'\to\infty$, $\mathsf F(T')/T'$ tends to this limit for every $T_1$, and $V_{\rm ren}(R)$ is the infimum of all these slopes. (f) A rectangle with an even side $2h$ is $\gamma(\mathfrak r\gamma)^{-1}$ for the site mirror halfway along that side, and a square is $\gamma(\mathfrak r\gamma)^{-1}$ for a diagonal mirror through two of its corners; as in Step 2, Proposition~\YMMref{103.4}{758} gives $s(W_C)\ge0$.

\emph{Step 9 (no criterion).} (g) Steps 1--7 use of $\nu$ only that it is a periodic Gibbs state, invariant under the symmetries of the lattice and reflection positive at every lattice mirror (Proposition~\YMMref{103.4}{758}), and, for the second Hankel form, that it is positive at the bond mirrors on the functions of two crossing bonds. With the hypothesis of (g) they hold for $s$, but for the upper bound of (a), in whose place $m_s\le1$ stands; and $m_s(R,T)\ge w^{RT}$ gives $V_s(R)\le-R\log w\le R\log\frac{4(\beta+1)}\beta$. On site mirrors alone the first Hankel form, with all $c$ but two equal to $0$, gives $m_s(R,h+h')^2\le m_s(R,2h)\,m_s(R,2h')$. At $(h-1,h+1)$, and from $m_s(R,0)=1$ and $m_s(R,2)>0$, induction gives $m_s(R,2h)>0$ for all $h$, and $\psi(h)=-\log m_s(R,2h)$ is concave, $\ge0$, with $\psi(0)=0$; so $\psi(h)/h$ tends to $2V_{\rm e}(R)=\inf_h(\psi(h+1)-\psi(h))\ge0$, and $\psi(h_2)-\psi(h_1)\ge2(h_2-h_1)V_{\rm e}(R)$ is the inequality for even $T_1$ and $T_2$. For an odd $T_2=2h_2+1>T_1$, the pair $(h_2,h_2+1)$ gives $m_s(R,2h_2+1)\le[m_s(R,2h_2)\,m_s(R,2h_2+2)]^{1/2}$, and the even case twice gives the inequality. For even $R$: the pair $(0,1)$ gives $m_s(1,2)\ge m_s(1,1)^2=w^2>0$, by Step 4, and $m_s(2,2)\ge m_s(2,1)^2=m_s(1,2)^2>0$; the induction at the width $2$ gives $m_s(2,2h)>0$, and $m_s(2h,2)=m_s(2,2h)$.
\end{proof}

\begin{proof}[Proof of Proposition~\ref{prop:where}]
(a) $F(l,2l,2l)\le0$ and $-F(2l,2l,2l)\le6\beta(2l)^3$, so $0\le\mathsf e(l)\le48\beta l^3\le x_0$. (b) For $\mathsf k\in\mathbb R^4$ let $\pi_{\mathsf k}$ be the law on $G\simeq S^3$ with the density proportional to $e^{\langle U,\mathsf k\rangle}$. Every linear component $\langle U,\mathsf e\rangle$, $|\mathsf e|=1$, has variance $\le\frac14$ under $\pi_{\mathsf k}$: with $\kappa=|\mathsf k|$ and $x=\langle U,\mathsf k\rangle/\kappa$, the mean $A(\kappa)$ of $x$ satisfies $A'=\operatorname{Var}x$ and, by an integration by parts, $\langle1-x^2\rangle=3A/\kappa$, so $A'=1-3A/\kappa-A^2$; a comparison with the positive root $\mathsf B(\kappa)$ of $1-3a/\kappa-a^2=\frac14$, whose derivative is $<\frac14$, together with $A\le\kappa/4$, which follows from $3A=\kappa\langle1-x^2\rangle$ and $\langle x^2\rangle\ge\frac14$, the two identities in the proof of Lemma~\ref{lem:bond}(a), $x$ having the law of $x_0$ under $\langle\cdot\rangle_\kappa$, gives $A\ge\mathsf B$ and so $\operatorname{Var}x\le\frac14$, while the components orthogonal to $\mathsf k$ have the variance $A/\kappa\le\frac14$.

So two such laws differ in total variation by at most $\frac14|\mathsf k-\mathsf k'|$. The conditional law of one bond given the others is $\pi_{\beta\mathsf h}$, $\mathsf h$ the sum of the four unit vectors of its staples, and changing one of its $12$ neighbouring bonds moves $\mathsf h$ by at most $2$. For a bounded Borel function $f$ of finitely many bonds let $\delta_b(f)$ be its oscillation in the single bond $b$, and $\Delta(f)=\sum_b\delta_b(f)$. The conditional expectation $E_b$ in the bond $b$ gives $\delta_b(E_bf)=0$, $\delta_{b'}(E_bf)\le\delta_{b'}(f)+\frac\beta2\delta_b(f)$ for the $12$ neighbours $b'$ of $b$, and $\delta_{b'}(E_bf)\le\delta_{b'}(f)$ for the other bonds; so after any sequence of updates the oscillations are, bond by bond, at most the entries of the vector $v$ that arises from $(\delta_b(f))_b$ by the linear rules $v_b\to0$, $v_{b'}\to v_{b'}+\frac\beta2v_b$ for the neighbours $b'$ of the updated bond $b$.

Let $S_0$ be the set of the bonds of $f$, $S\supset S_0$ a finite set of bonds, $d(b)$ the distance of $b$ from $S_0$ in the graph of neighbouring bonds, $r=\min\{d(b):b\notin S\}$, $\alpha=6\beta<1$, and $f_k$ the function that $k$ sweeps of updates through $S$, each bond of $S$ once in a sweep, make of $f$. An update at $b$ takes $v_b\alpha^{-\min\{d(b),r\}}$ from the weighted sum $\sum_{b'}v_{b'}\alpha^{-\min\{d(b'),r\}}$ and adds at most $\frac\beta2\cdot12\,v_b\,\alpha^{-\min\{d(b),r\}-1}=v_b\,\alpha^{-\min\{d(b),r\}}$, as $d(b')\le d(b)+1$; so the weighted sum does not increase, it equals $\Delta(f)$ at the start, when $v$ lives on $S_0$, where $d=0$, and the bonds outside $S$ carry the weight $\alpha^{-r}$: $\sum_{b\notin S}v_b\le\alpha^r\Delta(f)$ at all times.

An update at $b\in S$ changes $\sum_{b'\in S}v_{b'}$ by at most $-(1-\alpha)v_b$, and in a sweep the entries of the bonds of $S$ at the times of their updates add up to at least the sum at the beginning of the sweep; so $\sum_{b\in S}v_b\le\alpha^k\Delta(f)$ after $k$ sweeps. As $f_k$ is a function of the bonds of $S$ and of their neighbours, and two configurations of these bonds are joined by changing one bond at a time, $\sup f_k-\inf f_k\le(\alpha^k+\alpha^r)\Delta(f)$. Now let $l\ge4$ be even, $S$ the set of the bonds with both ends in the cube $\{0,\dots,l-3\}^3$, and $f=a_p$ for the plaquette $p=p_{12}(c_0,c_0,c_0)$, $c_0=\frac{l-4}2$. The ends of the bonds at distance $\le j$ from the four bonds of $p$ lie in $[c_0-j,c_0+1+j]^3$, which lies in $[0,l-3]^3$ for $j\le c_0$; so every bond not in $S$ has $d\ge c_0+1=\frac l2-1$, that is $r\ge\frac l2-1$; and $\Delta(a_p)\le8$.

The bonds of $S$ and their neighbours lie in a cube of $l$ sites a side, which the projections modulo $l$ and modulo $2l$ map one to one into the tori of sides $l$ and $2l$, the plaquettes at the bonds of $S$ onto plaquettes; so on both tori the conditional laws of the bonds of $S$ are those used above, the functions $f_k$ are the same, and both means of $a_p$ equal means of $f_k$, which lie in $[\inf f_k,\sup f_k]$. Letting $k\to\infty$: $|\bar a_{2l}(b)-\bar a_l(b)|\le8(6b)^{l/2-1}$, $\bar a_n(b)$ the mean of $a_p$ on the cube of side $n$ at the coupling $b$. As $\frac d{d\beta}F(n,n,n)=-3n^3\bar a_n(\beta)$ and $X=0$ at $\beta=0$,
\[
X(l)=24\,l^3\int_0^\beta\bigl(\bar a_{2l}(b)-\bar a_l(b)\bigr)\,db\le64\,l^2(6\beta)^{l/2}.
\]
At each of the four pairs of a coupling and a side named in (b), $64\,l^2(6\beta)^{l/2}\le\frac1{20}$; and $l\mapsto l^2(6\beta)^{l/2}$ decreases for $l>4/\log\frac1{6\beta}$, which lies below each of the four sides. Theorem~\ref{thm:crit}(a) gives the criterion. (c) The same contraction, with $S$ the set of the bonds at the distance $<r$ from the bonds of $f$, applied to two Gibbs states $s$, $s'$, for which $s(E_bF)=s(F)$ by the Gibbs equation of one bond, gives $|s(f)-s'(f)|=|s(f_k)-s'(f_k)|\le(\alpha^k+\alpha^r)\Delta(f)$; let $k\to\infty$ and then $r\to\infty$.
\end{proof}

\begin{proof}[Proof of Proposition~\ref{prop:says}]
(a) Lemma~\ref{lem:seq}(a), (b) with $\kappa_i=\lambda_i(2l,2l)^l$, so that $D=\mathsf e(l)$ and $\eta=\mathsf r$: $2\log1.0351\le0.069$, and $\varpi(e^{0.069}-1)\le0.0357638$. The tops $\Lambda_e$, $e\ne00$, are among the $\lambda_i$, $i\ge1$ (Lemma~\ref{lem:tr}), so $(\Lambda_e/\lambda_0)^l\le\mathsf r$, and $\log\frac1{0.0357638}\ge3.33$; and $\mathsf r\ge\sum_{e\ne00}(\Lambda_e/\lambda_0)^l$ gives the last statement. (b) With the time $1$ over $(2l,2l)$ let $w_f=W_{2l}(f;0)/Z$ on the cube of side $2l$. Inside each sector $(\sum_i\mu_i^l)^2\ge\sum_i\mu_i^{2l}$, so
\[
\mathsf e(l)=2\log Z(l,2l,2l)-\log Z(2l,2l,2l)\ge2\log\sum_f\sqrt{w_f}.
\]
By Lemma~\ref{lem:tr}(c), $w_{00}=\frac14(1+2r_1+r_2)$, $w_{10}=w_{01}=\frac14(1-r_2)$, $w_{11}=\frac14(1-2r_1+r_2)$, and $w_{11}\ge0$, $r_2\le r_1$ give $2r_1-1\le r_2\le r_1$; here $r_1\le1$ by Lemma~\ref{lem:tr}(c), and $r_1>0$ because, with the time the direction that is not in the twisted pair, $Z^-=\operatorname{Tr}\mathsf T_1^{2l}$ is the integral over the diagonal of a kernel with values $>0$ (Lemma~\ref{lem:tr}(i), (ii)). The sum of the square roots is concave in $r_2$, so on this interval it is at least the smaller of its two end values; with $x=1-r_1\in[0,1)$ these are $\sqrt{1-x}+\sqrt{2x}$ at $r_2=2r_1-1$ and $\frac12\sqrt{4-3x}+\frac32\sqrt x$ at $r_2=r_1$, and the first is the smaller, as $\frac14(4-3x)\ge1-x$ and $\frac32>\sqrt2$.

The function $\mathsf f$ increases on $[0,\frac23]$ and decreases on $[\frac23,1]$, with $\mathsf f(1)=\log2$; and $\mathsf f(0.00063)>0.069$, $\mathsf f(0.0039)\ge0.16567$. (c) Let $\mathsf e(l)\le0.069$ and $N\ge16l$, and let $\bar N=2^jl\le N<2\bar N$, $j\ge4$. By Lemma~\ref{lem:seq}(f), $\log(1+\mathsf r_N(N,N))\le\log(1+\mathsf r_{\bar N}(N,N))\le\delta(\bar N,N,N)\le4\delta_{\rm c}(\bar N)\le4\bar\delta_4\le0.0029221$ (Proposition~\ref{prop:rec}(e)); and by Lemma~\ref{lem:tr}(c), $1-Z^-/Z=1-r_1\le\frac43(1-e^{-0.0029221})\le0.0038905<0.0039$. For an integer $l'\ge16l$ with $\bar N=2^jl\le l'<2\bar N$, $j\ge4$: $\mathsf e(l')\le2\log(1+\mathsf r_{l'}(2l',2l'))\le2\delta(\bar N,2l',2l')\le2(2l'/\bar N)^2\delta_{\rm c}(\bar N)<32\,\bar\delta_4\le0.02338$.

The last sentence is the contrapositive of the first, at $l\le N/16$. (d) Theorem~\ref{thm:one}(b) at $R=T=1$ gives $\sigma^+\le\log(1/\nu(W_p))$, and $\sigma^+l^2\ge0.4757$.
\end{proof}

\noindent\emph{The numbers of Remark~\ref{rem:free}.} Over a torus of side $2l$ the momenta are $\frac\pi lm$, $m\in\mathbb Z^2$; three fields with one mode for each momentum $m\ne0$ give the Fock levels whose sum at the time $l$ is the product displayed, and the same at the time $2l$; in the notation of Lemma~\ref{lem:seq}, $1+\eta=1+\mathsf r$ and $1+s_2=1+\mathsf r'$. Summing $-3\log(1-e^{-\pi|m|})$ over the shells $|m|^2=1,2,4,5,8,9,10,13,\dots$ gives approximately $0.72010$, and $-3\log(1-e^{-2\pi|m|})$ gives approximately $0.02415$; so $D\approx1.416$.

\begin{proof}[Proof of Proposition~\ref{prop:small}]
\emph{The exact part.} On every periodic cube of an even side $N$, at every $\beta$, take as time a direction of the twisted pair, let $\nu$ be the other direction of that pair, let $\vartheta\ge2$ and $k\ge2$ be integers with $2k\vartheta=N$, and cut the cube into $2k$ slabs of $\vartheta$ layers. The transfer operator over $\vartheta$ layers, $\mathsf Y=\mathsf T'^{\,\vartheta}$ in the notation of the proof of Lemma~\ref{lem:mir}, commutes with the flip $\mathsf J$ of one layer, the operator $\mathsf J_i$ of Lemma~\ref{lem:tr} that changes the sign of the bonds of the direction $\nu$ that cross one plane $\{x_\nu=c\}$, $c\in\mathbb Z+\frac12$; $Z=\operatorname{Tr}\mathsf Y^{2k}$ and $Z^-=\operatorname{Tr}(\mathsf J\mathsf Y^{2k})$ (Lemma~\ref{lem:tr}(ii), (c)); so $P_\pm=\frac12(1\pm\mathsf J)$ have $\operatorname{Tr}(P_\pm\mathsf Y^{2k})=\frac12(1\pm r_1)Z$: the weights of the two sectors of the flip are $\frac12(1\pm r_1)$.

A function of the bonds is \emph{odd under the flip} if it changes its sign when the bonds of the direction $\nu$ that cross the plane $\{x_\nu=c\}$ change their signs in all layers. Let $D$, $D'$ be bounded Borel functions of the bonds with both ends in one and the same slab of $\vartheta$ layers, $D'$ odd under the flip; of $D$ no parity is asked. For $1\le j\le2k-1$ let $\mathsf t_j$ be the translation by $j\vartheta$ along the time, and $\mathsf t=\mathsf t_k$ the translation by $N/2$. With the slab operators $\mathsf Y_D$, $\mathsf Y_{D'}$ of the proof of Lemma~\ref{lem:mir}, $Z\langle D\cdot\mathsf t_jD'\rangle=\operatorname{Tr}(\mathsf Y^{2k-j-1}\mathsf Y_{D'}\mathsf Y^{j-1}\mathsf Y_D)$.

The weights of the layers and of the crossings do not change when all layers are flipped (proof of Lemma~\ref{lem:tr}); so $P_\sigma\mathsf Y=\mathsf YP_\sigma$ and, $D'$ being odd, $\mathsf J\mathsf Y_{D'}\mathsf J=-\mathsf Y_{D'}$, that is $P_\sigma\mathsf Y_{D'}=\mathsf Y_{D'}P_{-\sigma}$. Splitting the trace over the two sectors, the sector changes between $D$ and $D'$:
\[
Z\langle D\cdot\mathsf t_jD'\rangle=\sum_{\sigma=\pm1}\operatorname{Tr}\bigl((P_\sigma\mathsf Y)^{2k-j-1}\,\mathsf Y_{D'}\,(P_{-\sigma}\mathsf Y)^{j-1}\,\mathsf Y_D\bigr),
\]
a trace of $2k$ Hilbert--Schmidt operators. H\"older's inequality for it (Lemma~\YMMref{86.1}{680}), with $\|P_\sigma\mathsf Y\|_{2k}^{2k}=\operatorname{Tr}(P_\sigma\mathsf Y^{2k})$ and $\|\mathsf Y_D\|_{2k}^{2k}=\operatorname{Tr}(\mathsf Y_D^*\mathsf Y_D)^k\le\sup|D|^{2k}Z$ (the norm definition preceding Lemma~\YMMref{86.1}{680}; the trace is $Z$ times the mean of a product of $2k$ copies of $D$ and of its mirror image), gives
\[
|\langle D\cdot\mathsf t_jD'\rangle|\le\sup|D|\,\sup|D'|\sum_{\sigma=\pm1}\Bigl(\frac{1+\sigma r_1}2\Bigr)^{\frac{2k-j-1}{2k}}\Bigl(\frac{1-\sigma r_1}2\Bigr)^{\frac{j-1}{2k}}
\]
(with $0^0=1$), and at $j=k$, where the two summands are equal,
\[
|\langle D\cdot\mathsf tD'\rangle|\le2\sup|D|\,\sup|D'|\,\Bigl[\frac{1-r_1^2}4\Bigr]^{\frac12-\frac1{2k}} .
\]

\emph{The calibrated construction.} The constants are chosen as in (Th1), (Th2) of Section~\ref{sec:cutoff-calibration}, the setting in which Theorem~\YMMref{D}{8} is used; among them are the cube size $M=\mathsf L^{m'}$, with $m'\ge5$ (the setting preceding equation~\YMMref{(634)}{339} has $M\ge L^4M_1$ and $M_1\ge L$), the exponents $r\ge2$ and $\mathfrak p_0$, the constant $A_0>0$ of the thresholds (equation~\YMMref{(634)}{339}), $\gamma\in(0,\frac12]$, and the numbers $c'>0$, $\mathsf C^\Sigma_{\rm maj}$, $c_{\rm low}$ of (Th2). The setting (E) of Section~\ref{sec:cutoff-calibration} gives besides a label $\mathsf b\ge\gamma^{-2}$ and an integer $m\ge1$ with $M\,R(\mathsf b)\le\mathsf L^m$, $R(x)=\min\{\mathsf L^i:i\in\mathbb Z_{\ge0},\ \mathsf L^i\ge(\log x)^r\}\ge1$, and $\ell=2\mathsf L^m$. At the cutoff $K$ the coupling is $\beta=\mathsf L^K\mathsf b$, and the cube has $N=\ell\mathsf L^K=2\mathsf L^{m+K}$ sites a side. For an integer $0\le d<K$ the block field of the depth $d$ is $V=\mathcal M^{K-d}U$, $\mathcal M$ the block-averaging map of equations~\YMMref{(37)--(39)}{20}; it lives on a periodic cube of
\[
\mathsf n=\ell\mathsf L^d=2\mathsf L^{m+d}
\]
blocks a side, each block having $\vartheta=\mathsf L^{K-d}$ sites a side, so that $N=\mathsf n\vartheta$; its law is $\mu^K_d$ (Definition~\textup{\ref{IV-62-62}}). As $R\ge1$, $\mathsf L^m\ge M=\mathsf L^{m'}\ge\mathsf L^5$, and with $d\ge1$ and $\mathsf L\ge13$:
\[
\mathsf n\ge2\cdot13^6 .
\]
Let $|\partial V-\mathbf 1|_\infty$ be the largest $|\partial V(p)-\mathbf 1|$ over the plaquettes of the block lattice, $x_d=\mathsf L^d\mathsf b$, $\varepsilon=A_0x_d^{-1/2}(\log x_d)^{\mathfrak p_0}$ the threshold of equation~\YMMref{(634)}{339} at the coupling $x_d$, $\delta=\exp[7\mathsf n^2\varepsilon]-1$, and $\mathsf R_d$ the number $\mathsf R_t$ of the definitions preceding Proposition~\YMMref{115.2}{806} at the level $t=K-d$; as there $N_t=\ell^3\mathsf L^{3d}$ and $x_t=x_d$, $\mathsf R_d$ does not depend on $K$. The hypotheses (H) are: $d\ge1$; $x_d\ge\mathsf t_\natural$, the fixed calibration of Section~\ref{sec:cutoff-calibration}; $\delta\le\frac1{20}$; and $\mathsf R_d\le\frac1{40}$.

\emph{Step 1: under (H) the block field is nearly flat with probability at least $\frac{39}{40}$, at every cutoff $K>d$.} The data $K$, $\mathsf b$, $m$ and the top level $t=K-d$ make the setting (Th) of Section~\ref{sec:cutoff-calibration}, as $1\le t\le K-1$ and $x_t=x_d\ge\mathsf t_\natural$; Lemma~\YMMref{118.1}{812} gives what Proposition~\textup{\YMMref{115.2}{806}} asks of it. Let $F=\mathbf 1_{\mathsf{In}_t}(\mathcal M^tU)$, $\mathsf{In}_t=\{|\partial V-\mathbf 1|_\infty<2\varepsilon\}$ the inner ball of the definitions preceding Lemma~\YMMref{113.10}{796}: a bounded gauge-invariant Borel function, by Proposition~\YMMref{4.5}{20}, \YMMref{4.5}{20} and~\YMMref{113.10}{796}. The weight $\mathfrak w^{\rm sf}_t$ of the definitions preceding Proposition~\YMMref{115.2}{806} carries the factor $\mathbf 1_{\mathsf{In}_t}(\mathcal M^tU)$, so $\mathfrak w^{\rm sf}_t=F\,\mathfrak w^{\rm sf}_t$, and the right inequality of Proposition~\textup{\YMMref{115.2}{806}}, domination and lost-mass bounds gives $Z-\int F\rho_0\,dU\le Z-\int\mathfrak w^{\rm sf}_t\rho_0\,dU\le\mathsf R_dZ$, $\rho_0=e^{-\beta A}$ at $\beta=\mathsf L^K\mathsf b$. That is
\[
\mu^K_d\bigl(|\partial V-\mathbf 1|_\infty\ge2\varepsilon\bigr)\le\mathsf R_d\le\tfrac1{40}.
\]
Theorem~\YMMref{D}{8} enters here, through Proposition~\textup{\YMMref{115.2}{806}}.

\emph{Step 2: a winding loop of the block field, its parity and its translate.} Fix a block site $\xi$; let $\Gamma_1$, $\Gamma_2$, $\Gamma_{12}$ be the closed paths of the block lattice from $\xi$ with the step sequences $(e_1^{\mathsf n})$, $(e_2^{\mathsf n})$, $(e_1^{\mathsf n},e_2^{\mathsf n})$, $P_i=V(\Gamma_i)$, and for each of them $D_\Gamma(U)=W_\Gamma(V)=\frac12\operatorname{tr}V(\Gamma)$. Take the direction $3$ as time, the flip of the bonds of the direction $1$ that cross the plane $\{x_1=c\}$, $c=\xi_1+\frac12\vartheta$, and $\mathsf t$ the translation by $N/2$ along $e_3$. Then: (i) $D_\Gamma$ is a real Borel function with $|D_\Gamma|\le1$ (Proposition~\YMMref{4.5}{20}); (ii) it is a function of the bonds with both ends in one slab of $\vartheta$ layers, the slab of the blocks of the sites of $\Gamma$, by the locality of $\mathcal M^{K-d}$ (Proposition~\YMMref{4.5}{20}); (iii) $D_{\Gamma_1}$ and $D_{\Gamma_{12}}$ are odd under the flip: as $\mathsf L$ is odd, the plane $\{x_1=c\}$ separates blocks at every level from $d$ to $K$, on each double block the flip is the gauge transformation that is $-\mathbf 1$ on one side of the plane, so by Proposition~\YMMref{4.5}{20} the averaging map carries the flip of the bonds that cross the plane at one level to the flip at the next level; and $\Gamma_1$, $\Gamma_{12}$ run exactly one block bond that crosses the plane, $\Gamma_2$ none; (iv) $\mathsf tD_\Gamma=W_{\mathsf t\Gamma}(V)$, as $\mathsf t$ is the translation by $\frac{\mathsf n}2$ blocks and $\mathcal M^{K-d}$ commutes with it (Lemma~\ref{lem:setup-lattice-covariance}).

And $\mathsf m_\Gamma:=\langle D_\Gamma\cdot\mathsf tD_\Gamma\rangle=\int W_\Gamma W_{\mathsf t\Gamma}\,d\mu^K_d$ has $\mathsf m_{\Gamma_2}=\mathsf m_{\Gamma_1}$: the exchange of $e_1$ and $e_2$ about $\xi$ keeps the law of the cube, commutes with $\mathsf t$ and with $\mathcal M^{K-d}$ (Lemma~\ref{lem:setup-lattice-covariance}), and carries $\Gamma_1$ to $\Gamma_2$.

On $\mathsf B=\{|\partial V-\mathbf 1|_\infty<2\varepsilon\}$ the two winding holonomies nearly commute: the step sequences $(e_1^{\mathsf n},e_2^{\mathsf n})$ and $(e_2^{\mathsf n},e_1^{\mathsf n})$ are rearrangements of one another, of the length $2\mathsf n$, and Lemma~\YMMref{4.1}{17} gives
\[
|P_2P_1P_2^{-1}P_1^{-1}-\mathbf 1|\le(1+|\partial V-\mathbf 1|_\infty)^{\mathsf n(2\mathsf n-1)}-1<e^{4\mathsf n^2\varepsilon}-1\le\delta .
\]
As $|g-\mathbf 1|^2=2-\operatorname{tr}g$ for $g\in G$, $\frac12\operatorname{tr}(P_2P_1P_2^{-1}P_1^{-1})\ge1-\frac12\delta^2$. The trace identity of Fricke, $\operatorname{tr}(ABA^{-1}B^{-1})=(\operatorname{tr}A)^2+(\operatorname{tr}B)^2+(\operatorname{tr}AB)^2-\operatorname{tr}A\,\operatorname{tr}B\,\operatorname{tr}AB-2$ for $A,B\in\mathrm{SL}(2,\mathbb C)$, at $A=P_2$, $B=P_1$, gives with $x_i=W_{\Gamma_i}(V)$ and $x_{12}=W_{\Gamma_{12}}(V)$
\[
x_1^2+x_2^2+x_{12}^2-2x_1x_2x_{12}\ge1-\tfrac14\delta^2,\qquad\text{so}\qquad2x_1^2+2x_2^2+x_{12}^2\ge1-\tfrac14\delta^2,
\]
as $|x_{12}|\le1$ gives $-2x_1x_2x_{12}\le x_1^2+x_2^2$. On $\mathsf B$ a winding loop is also close to its translate: with $\epsilon$ the path of $\frac{\mathsf n}2$ steps $e_3$ from $\xi$ to $\mathsf t\xi$ and $g=V(\epsilon)$, the step sequences of $\Gamma$ followed by $\epsilon$ and of $\epsilon$ followed by $\mathsf t\Gamma$ are rearrangements of one another, of the length $\le\frac52\mathsf n$, and Lemma~\YMMref{4.1}{17} gives $|gV(\mathsf t\Gamma)g^{-1}-V(\Gamma)|\le(1+|\partial V-\mathbf 1|_\infty)^{\frac{25}8\mathsf n^2}-1<e^{6.25\,\mathsf n^2\varepsilon}-1\le\delta$, so that $|W_{\mathsf t\Gamma}-W_\Gamma|\le\delta$. As $W_\Gamma W_{\mathsf t\Gamma}\ge W_\Gamma^2-|W_\Gamma-W_{\mathsf t\Gamma}|$ and $|W_\Gamma-W_{\mathsf t\Gamma}|\le2$ off $\mathsf B$, Step~1 gives $\mathsf m_\Gamma\ge\mu^K_d(W_\Gamma^2)-\delta-2\mathsf R_d$; with the weights $2,2,1$ and $\mathsf m_{\Gamma_2}=\mathsf m_{\Gamma_1}$,
\[
4\mathsf m_{\Gamma_1}+\mathsf m_{\Gamma_{12}}\ge\bigl(1-\tfrac14\delta^2\bigr)\bigl(1-\mathsf R_d\bigr)-5\bigl(\delta+2\mathsf R_d\bigr)\ge\tfrac{1599}{1600}\cdot\tfrac{39}{40}-\tfrac12=\tfrac{30361}{64000}>0.4743,
\]
the middle member decreasing in $\delta$ and in $\mathsf R_d$ on $[0,1]$. So $\mathsf m_\Gamma\ge0.0948>0.09$ for $\Gamma=\Gamma_1$ or for $\Gamma=\Gamma_{12}$; both loops give real functions of modulus $\le1$ that are odd under the flip and are functions of one slab of the thickness of a block.

\emph{Step 3: the ceiling.} The exact part at $D=D'=D_\Gamma$, with that $\Gamma$, $\vartheta=\mathsf L^{K-d}\ge2$ and $2k=\mathsf n$, so that $2k\vartheta=N$ and $\frac1{2k}=\frac1{\mathsf n}=\frac{a_d}\ell$, the ratio of the block size $a_d=\mathsf L^{-d}$ at the depth $d$ to the side $\ell$, gives $0.09\le\mathsf m_\Gamma\le2[\frac14(1-r_1^2)]^{\frac12(1-2/\mathsf n)}$, that is, with $\mathsf x=(0.045)^{2/(1-2/\mathsf n)}$,
\[
r_1^2\le1-4\mathsf x,\qquad r_1\le1-2\mathsf x=1-2\,(0.045)^{2/(1-2/\mathsf n)},
\]
the second because $(1-2\mathsf x)^2\ge1-4\mathsf x$, $1-2\mathsf x>0$ and $r_1>0$. The right member of the second inequality decreases in $\mathsf n$ to the limit $0.99595$, and it is $\le0.99596$ if and only if $\mathsf n\ge5020$; at $\mathsf n=2\cdot13^6$ it is $\le0.99595002$. So under (H), $r_1\le0.99596$ on the cube of $2\mathsf L^{m+K}$ sites a side at $\beta=\mathsf L^K\mathsf b$, for every $K>d$.

\emph{Step 4: the range of (H); the numbers $t_\circ$ and $m_t$; (a).} Let $\mathsf C_+=\max\{0,\mathsf C^\Sigma_{\rm maj}+2+c_{\rm low}\}$, and let $y_\circ\ge2$ be the least real number such that for all $y\ge y_\circ$
\[
\begin{multlined}
28\mathsf L^{11/2}M^2A_0\,y^{2r}(y+\log\mathsf L)^{\mathfrak p_0}e^{-y/2}\le\log\tfrac{21}{20},\qquad\log\bigl(16\mathsf L^9y^{3r}\bigr)+\bigl(\mathsf C_++\tfrac92y\bigr)(2\mathsf L^2M)^3y^{3r} \\
-c'A_0^2y^{2\mathfrak p_0}\le-\log40 ;
\end{multlined}
\]
it exists, as the first left member tends to $0$ and, in the second, $2\mathfrak p_0>3r+2$ at the chosen constants (the choice $\mathfrak p_0\ge6r+1$ in Proposition~\YMMref{69.1}{595}) and $c'A_0^2>0$. Put $t_\circ=\min\{2\mathsf L\gamma^2,\ 2\mathsf L^2/\mathsf t_\natural,\ 2\mathsf Le^{-y_\circ}\}$. For $0<t\le t_\circ$ let $m_t$ be the least integer $m\ge1$ with $M\,R(2\mathsf L^m/t)\le\mathsf L^m$, which exists as $R(x)\le\max\{1,\mathsf L(\log x)^r\}$; let $\ell_t=2\mathsf L^{m_t}$, $\mathsf b_t=\ell_t/t$ and $y=\log\mathsf b_t$. Then (E) holds for the label $\mathsf b_t$ and the integer $m_t$, as $\mathsf b_t\ge2\mathsf L/t_\circ\ge\gamma^{-2}$; $y\ge\log(2\mathsf L/t_\circ)\ge y_\circ$; and $\ell_t\le2\mathsf L^2My^r$, as $m_t=1$ or $\mathsf L^{m_t-1}<M\,R(\mathsf b_t/\mathsf L)\le M\mathsf Ly^r$.

At $d=1$: $x_1=\mathsf L\mathsf b_t\ge2\mathsf L^2/t_\circ\ge\mathsf t_\natural$; $\mathsf n=\mathsf L\ell_t\le2\mathsf L^3My^r$ and $x_1=\mathsf Le^y$ give $7\mathsf n^2\varepsilon\le\log\frac{21}{20}$ by the first inequality, that is $\delta\le\frac1{20}$; and $\log x_1\ge y$, $2M^{-3}\ell_t^3\mathsf L^3\le16\mathsf L^9y^{3r}$ and $\ell_t^3\le(2\mathsf L^2M)^3y^{3r}$ give $\mathsf R_1\le\frac1{40}$ by the second. So (H) holds for $(\mathsf b_t,m_t)$ at $d=1$. At the cutoff $K\ge2$ the cube has $N=2\mathsf L^{m_t+K}$ sites a side and $\beta=\mathsf L^K\mathsf b_t=N/t$; by Step~3, $r_1\le0.99596$ on the cube of $2\mathsf L^j$ sites a side at $\beta=2\mathsf L^j/t$, for every $j\ge m_t+2$. Then $1-r_1\ge0.00404\ge0.0039$, and Proposition~\ref{prop:says}(b), (c) give $\mathsf e(\mathsf L^j)\ge\mathsf f(0.0039)\ge0.16567$ and $\mathsf e(l)>0.069$ for $3\le l\le\mathsf L^j/8$: (a).

(b) The cube of the criterion at the cutoff $K$ has $2l_*\mathsf L^K=2\mathsf L^{m_0+K}$ sites a side and $\beta=\mathsf L^K\mathsf b=2\mathsf L^{m_0+K}/t$ with $t=\ell_0/\mathsf b\le t_\circ$; apply (a) with $j=m_0+K$, for all $K\ge m_t+2-m_0$. (c) Let $\bar y=\frac43[\log(2\mathsf L^2M)+r\log(4r)+\log(\mathsf L/t_\circ)]$ and $\beta_\flat=2\mathsf L^5M\bar y^r/t_\circ$. For $\beta\ge\beta_\flat$ choose $j$ with $2\mathsf L^j\le t_\circ\beta<2\mathsf L^{j+1}$ and put $t=2\mathsf L^j/\beta\in(t_\circ/\mathsf L,t_\circ]$. With the numbers of Step~4 at this $t$, $y\le\log\ell_t+\log(\mathsf L/t_\circ)$ and $\ell_t\le2\mathsf L^2My^r$; with $r\log y\le\frac y4+r\log(4r)-r$ this gives $y\le\bar y$, so $2\mathsf L^{m_t+2}=\mathsf L^2\ell_t\le2\mathsf L^4M\bar y^r=\frac{t_\circ}{\mathsf L}\beta_\flat\le\frac{t_\circ}{\mathsf L}\beta<t\beta=2\mathsf L^j$, that is $j\ge m_t+2$; and (a) gives $\mathsf e(l)>0.069$ for $3\le l\le\mathsf L^j/8$, while $\mathsf L^j/8>t_\circ\beta/16\mathsf L$. For the uniform criterion at a label $\mathsf b$ take $K\ge K_0$ with $L^K\mathsf b\ge\beta_\flat$: the criterion at $l=l_*L^K$ forces $l_*L^K>\frac{t_\circ}{16\mathsf L}L^K\mathsf b$.

(d) Assume \textbf{(Toron floor)} and put $\mathsf G_N=\log(Z(N,N,N)/\lambda_0(N,N)^N)=\log(1+\mathsf r_N(N,N))$. By Proposition~\ref{prop:says}(a), and as $\mathsf r_t$ decreases in $t$, $\mathsf e(N/2)\ge\log(1+\mathsf r_{N/2}(N,N))\ge\mathsf G_N$; and for every integer $q\ge2$, $Z(N,N,qN)\ge\lambda_0(N,N)^{qN}$ gives $\mathsf G_N\ge\mathsf d_q(N)/q$, $\mathsf d_q(N)=qF(N,N,N)-F(N,N,qN)$. Three facts bound $e^{-\mathsf d_q(N)}=Z(N,N,qN)/Z(N,N,N)^q$ from above: a swap on the long box, which bounds it through the largest probability on the cube $(N,N,N)$ that the five invariants of the pair of the winding holonomies at one site lie in one small cell; a bound of that probability on a bare cube (Lemma~\ref{lem:anti}); and a comparison of the laws of the block fields at two consecutive cutoffs in which the errors that bound quotients of densities enter as factors (Lemma~\ref{lem:cut}), which carries the bound from one low cutoff $K_1$ to all higher ones. Besides Theorem~\YMMref{D}{8} and the statements among the construction results named for this clause, (d) rests on \textbf{(Toron floor)} and these two lemmas, which give the estimate \eqref{eq:wind} and, with the chain from a low cutoff, the estimate \eqref{eq:lowcut}, and are stated below with their proofs.

\emph{A swap on the long box.} For $g\in G$ write $g=\cos\alpha(g)\,\mathbf 1+i\sum_j\mathsf v_j(g)\tau_j$, $\tau_1,\tau_2,\tau_3$ the Pauli matrices, with the class angle $\alpha(g)=\arccos\frac12\operatorname{tr}g\in[0,\pi]$ and the vector part $\mathsf v(g)\in\mathbb R^3$, $|\mathsf v(g)|=\sin\alpha(g)$. To a pair $(g_1,g_2)$ of elements of $G$ attach the real $3\times2$ matrix $\mathsf v$ with the columns $\mathsf v(g_1)$, $\mathsf v(g_2)$, the matrix $\mathsf s=(\mathsf v^{\rm t}\mathsf v)^{1/2}$, and the five invariants $\mathsf a(g_1,g_2)=(\alpha(g_1),\alpha(g_2),\mathsf s_{11},\mathsf s_{12},\mathsf s_{22})\in\mathbb R^5$: they do not change under a simultaneous conjugation of the pair, and if $|g_i-hg'_ih^{-1}|\le\delta\le1$ for $i=1,2$ and one $h\in G$, then $|\mathsf a(g_1,g_2)-\mathsf a(g'_1,g'_2)|_\infty\le2\delta$, as $|\mathsf v(g)-\mathsf v(g')|\le|g-g'|$, $|g-g'|\ge2\sin\frac12|\alpha(g)-\alpha(g')|$, and $\mathsf v\mapsto(\mathsf v^{\rm t}\mathsf v)^{1/2}$ has the Lipschitz constant $\sqrt2$ for the Hilbert--Schmidt norm.

Let $V$ be either the bare field, with $\mathsf n=N$, or the block field of the depth $0$, on a lattice of $\mathsf n=\ell$ sites a side in the section, formed on the long box by the same local formulas; let $\omega\in(0,1]$ with $\delta_q=(1+\omega)^{q\mathsf n^2}-1\le1$. To a site $\xi$ attach $\mathsf a(V,\xi)=\mathsf a(P_1(\xi),P_2(\xi))$, the five invariants of the pair of its winding holonomies in the directions $1$ and $2$. Cut $\mathbb R^5$ into half-open cubes $c_k$ of a side $w\ge2\delta_q$, and let $c^+_k$ be the half-open cube of side $3w$ with the same centre and the same choice of included faces, so that every point lies in exactly $3^5$ of them. On the event that every plaquette of $V$ on the box $(N,N,qN)$ is within $\omega$ of $\mathbf 1$, let $\xi'=\xi+ke_3$, $0\le k<q\mathsf n$, and $\mathsf t$ the transport of $V$ from $\xi$ to $\xi'$ along the line of the direction $3$: for $\mu=1,2$ the transports $\mathsf tP_\mu(\xi')$ and $P_\mu(\xi)\mathsf t$ run along the step sequences $(e_3^k,e_\mu^{\mathsf n})$ and $(e_\mu^{\mathsf n},e_3^k)$, which pass into one another by $k\mathsf n\le q\mathsf n^2$ exchanges of adjacent entries, and each exchange multiplies the transport by a conjugate of $\partial V(p)^{\pm1}$ for one plaquette $p$, as in the proof of Lemma~\YMMref{4.1}{17}; so $|\mathsf tP_\mu(\xi')\mathsf t^{-1}-P_\mu(\xi)|\le\delta_q$, with one $\mathsf t$ for both $\mu$, and $\mathsf a(V,\cdot)$ stays within $2\delta_q$ of $\mathsf a(V,\xi)$ at all the sites of that line.

Cut the box by walls into $q$ slabs of $N$ layers, with the mirrors $\theta_i$ in their middle planes and sites $\xi^-_i$ of one line of the direction $3$ just below these planes; then the indicator of the event is at most $\sum_k\prod_ig^k_i\,(g^k_i\circ\theta_i)$, $g^k_i=\mathbf 1\{\mathsf a(V,\xi^-_i)\in c^+_k\}$, a function of the bonds of the lower half of the slab $i$ (for the block field, whose walls of the blocks contain these planes, as $N=\mathsf n\mathsf L^K$ with $\mathsf n$ even, by the locality of the averaging map and its covariance under the lattice isometries, Proposition~\YMMref{4.5}{20}, \ref{lem:setup-lattice-covariance}). The slab operator of $g^k_i\,(g^k_i\circ\theta_i)$ is $\mathsf Y^*\mathsf Y\ge0$, $\mathsf Y$ that of the lower half with $g^k_i$, as in the proof of Lemma~\ref{lem:mir}; and $|\operatorname{Tr}(\mathsf O_1\cdots\mathsf O_q)|\le\prod_i\operatorname{Tr}\mathsf O_i$ for positive operators of trace class.

So every slab closes into a cube: $Z(N,N,qN)\langle\prod_ig^k_i\,(g^k_i\circ\theta_i)\rangle\le\prod_iZ(N,N,N)\langle g^k_i\rangle_{(N,N,N)}$. The plaquettes off the event are bounded by Schwarz's inequality at a mirror (Lemma~\ref{lem:mir}(a)) and the same closing. With $P_\ast$ the largest probability on the cube $(N,N,N)$ that $\mathsf a(V,\xi)$ lies in one $c^+_k$, and $\rho$ the largest probability on it that one plaquette of $V$ is not within $\omega$ of $\mathbf 1$, this gives $x\le243\,P_\ast^{q-1}+3q\mathsf n^3(\rho x)^{1/2}$ for $x=e^{-\mathsf d_q(N)}$, whence
\[
e^{-\mathsf d_q(N)}\le486\,P_\ast^{q-1}+9q^2\mathsf n^6\rho .
\]

\emph{Anticoncentration on a bare cube.} For an integer $N'\ge4$ and a coupling $\beta'>0$ let $\mu_{N',\beta'}$ be the law $Z^{-1}e^{-\beta'\sum_pa_p}dU$ of the periodic cube of $N'$ sites a side, $Z=Z(N',N',N')$, and $|\partial U-\mathbf 1|_\infty$ the largest $|\partial U(p)-\mathbf 1|$ over its plaquettes.

\begin{lemma}[Anticoncentration on a bare cube]\label{lem:anti}
Assume \textbf{(Toron floor)}. Let $N'\ge4$ be an integer, $\beta'\ge1.1\cdot10^9N'^{16}$ and $0<\omega\le1.6\cdot10^{-7}N'^{-13}$. For every site $\xi$, every $w'$ with $(N'^5/2\beta')^{1/2}\le w'\le\frac1{20}$ and every closed cube $\mathsf Q$ of the half side $w'$ in $\mathbb R^5$,
\[
\mu_{N',\beta'}\bigl(\mathsf a(U,\xi)\in\mathsf Q,\ \max_p|\partial U(p)-\mathbf 1|<\omega\bigr)\le10^{21}\,N'^{29.5}\,\beta'^{1/2}\,w'^3 .\qedhere
\]
\end{lemma}

Under \textbf{(Toron floor)}, the bound $C_{\rm a}N'^{c_{\rm a}}\beta'^{1/2}w'^3$ of \eqref{eq:wind} holds at the couplings $\beta'\ge c'_{\rm a}N'^{c''_{\rm a}}$ with the absolute numbers $C_{\rm a}=10^{21}$, $c_{\rm a}=29.5$, $c'_{\rm a}=1.1\cdot10^9$, $c''_{\rm a}=16$, for every $\omega$ below $1.6\cdot10^{-7}N'^{-13}$.

\emph{Sketch of the proof of Lemma~\ref{lem:anti}.} Write $\mathsf V=N'^3$, take $\xi=0$, let $|\cdot|_2$ be the Euclidean norm of a family indexed by bonds or by plaquettes, and $Z[F]=\int F\,e^{-\beta'\sum_pa_p}dU$. (i) \emph{The gauge of a maximal tree.} Let $T$ be the tree of the bonds $\langle x,x+e_3\rangle$ with $x_3\le N'-2$, $\langle x,x+e_2\rangle$ with $x_3=0$, $x_2\le N'-2$, and $\langle x,x+e_1\rangle$ with $x_2=x_3=0$, $x_1\le N'-2$. A bounded Borel function that the gauge transformations do not change has the same integral against $dU$ as against the Haar measures of the bonds off $T$ at $U=\mathbf 1$ on $T$; there the winding holonomies at $0$ are the values $P_\mu$ of $U$ at the three bonds $b^*_\mu$ from $(N'-1)e_\mu$ to $0$.

On the set $S_0$ of the other $2\mathsf V-2$ bonds write $U(b)=V^P(b)\,e^{a(b)}$, $a(b)\in\mathfrak{su}(2)$, $|a(b)|<\pi$, where $V^P(b)=P_\mu$ if $b$ is a bond of the direction $\mu$ that starts in the plane $x_\mu=N'-1$, and $V^P(b)=\mathbf 1$ else; in $a(b)$ the Haar measure has a density at most $(2\pi^2)^{-1}$. (ii) \emph{Small fields.} Along the tree every bond $b\in S_0$ is rebuilt from at most $2N'^2$ plaquettes, none of them a corner plaquette $p_{\mu\nu}(x)$ with $x_\mu=x_\nu=N'-1$: $|U(b)-V^P(b)|\le2N'^2|\partial U-\mathbf 1|_\infty$. The same rebuilding, linearised, gives $|a|_2\le2N'^{5/2}|d_Pa|_2$ for every triple $P$ and every family $a$ on $S_0$, $d_P$ being the linearisation at $V^P$ of $a\mapsto(\partial U(p))_p$, a covariant lattice curl.

So on $\{|\partial U-\mathbf 1|_\infty<\omega\}$ every $|a(b)|\le\pi N'^2\omega$, and at the corner plaquettes, where $\partial U(p)$ is the commutator $C_{\mu\nu}=P_\mu P_\nu P_\mu^{-1}P_\nu^{-1}$ up to four factors $e^{\pm a(b)}$, $|C_{\mu\nu}-\mathbf 1|\le9N'^2\omega$. Stokes' formula over the $N'^2$ plaquettes $p_{\mu\nu}(x)$ with $x_\rho=0$, $\rho$ the third direction, gives $2|\mathsf v(P_\mu)\times\mathsf v(P_\nu)|=|C_{\mu\nu}-\mathbf 1|\le\sum_p|\partial U(p)-\mathbf 1|$, and Schwarz's inequality then $\sum_pa_p\ge2N'^{-2}\sum_{\mu<\nu}|\mathsf v(P_\mu)\times\mathsf v(P_\nu)|^2$. (iii) \emph{The ceiling at a fixed triple.} Let $\kappa$ be the family of the logarithms of the $C_{\mu\nu}$, placed at the corner plaquettes of the planes $\mu\nu$.

With $|e^{X_1}\cdots e^{X_k}-e^{X_1+\dots+X_k}|\le\frac12(\sum_i|X_i|)^2$ for $X_i\in\mathfrak{su}(2)$, (ii) gives on $\{|\partial U-\mathbf 1|_\infty<\omega\}$ that $2\sum_pa_p\ge(1-\eta_1)^2|d_Pa+\kappa|_2^2$, $\eta_1=39\pi N'^{9/2}\omega\le2\cdot10^{-5}N'^{-8.5}$; keeping the fraction $N'^{-3}$ of the action for the holonomies,
\[
\sum_pa_p\ge\tfrac12(1-N'^{-3})(1-\eta_1)^2|d_Pa+\kappa|_2^2+\frac2{N'^5}\sum_{\mu<\nu}|\mathsf v(P_\mu)\times\mathsf v(P_\nu)|^2 .
\]
At a fixed triple the Gaussian integral over the $6\mathsf V-6$ real coordinates of $a$ is at most $(2\pi/\beta'')^{3\mathsf V-3}D(P)^{-1/2}$, with $\beta''=\beta'(1-N'^{-3})(1-\eta_1)^2$ and $D(P)$ the determinant of $d_P^*d_P$ on the families on $S_0$, the shift by $\kappa$ only lowering it; and $(\beta'/\beta'')^{3\mathsf V-3}\le(1-N'^{-3})^{-3N'^3}(1-\eta_1)^{-6N'^3}\le e^{3.2}$. (iv) \emph{Determinants.} A triple whose three commutators lie within $\epsilon$ of $\mathbf 1$ lies within $(2\epsilon)^{1/2}$ of a triple $P'$ with one axis, $P'_\mu=\cos\alpha_\mu\,\mathbf 1+i\sin\alpha_\mu\sum_j\hat n_j\tau_j$, and $D(P)\ge(1-48N'^{5/2}\delta)^{12\mathsf V}D(P')$ for $\max_\mu|P_\mu-P'_\mu|\le\delta$; at $\epsilon=9N'^2\omega$, $\omega\le1.6\cdot10^{-7}N'^{-13}$, this gives $D(P)^{-1/2}\le e^{0.5}D(P')^{-1/2}$.

At $P'$ the families split into the colour along $\hat n$, on which $d_{P'}$ is the lattice curl, and the complex colour across $\hat n$, on which it is the curl twisted by the phases $e^{-2i\alpha_\mu}$ at the bonds where $V^{P'}=P'_\mu$: $D(P')=D(0)\,D(2\alpha)^2$, where $D(\theta)$, for $\theta\in\mathbb R^3$, is the determinant of the curl twisted by the phases $e^{i\theta_\mu}$ on the families on $S_0$, an even function, $2\pi$-periodic in each $\theta_\mu$. For every momentum $k\in(\frac{2\pi}{N'}\mathbb Z/2\pi\mathbb Z)^3$ the twisted curl has the eigenvalue $\sum_\mu(2-2\cos(k_\mu+\theta_\mu/N'))$ twice, besides the gradients; let $\mathfrak d(\theta)$ be the product of these numbers over all $k$, and $\mathfrak d'$ their product over $k\ne0$ at $\theta=0$.

Then $D(0)=\mathfrak d'$, the change of coordinates along the tree being unimodular, and $D(\theta)\ge\mathfrak d'/(1.5\cdot10^{10}N'^{24})$ for every $\theta$: by $2\pi$-periodicity first take $\theta\in[-\pi,\pi]^3$; for $|\theta|_\infty\le\frac1{64}N'^{-11/2}$ by a comparison of the forms that uses the bound $|a|_2\le2N'^{5/2}|d_Pa|_2$ of (ii), and else by Cauchy's interlacing from the determinant $|1-e^{i\theta_\mu}|^2\mathfrak d(\theta)$ on the families that vanish on $T$ and at $b^*_\mu$ alone, $|\theta_\mu|=|\theta|_\infty$, together with $\mathfrak d(\theta)\ge4\sin^2\frac{\theta_\mu}2\,N'^{-2}\mathfrak d'$, which the product formula $\prod_{j=0}^{n-1}(2\cosh m-2\cos\frac{2\pi j+\vartheta}n)=2\cosh nm-2\cos\vartheta$ gives. This covers the triples near $P_\mu=\pm\mathbf 1$ as well.

Together, for every Borel set $\mathcal P$ of triples that the simultaneous conjugations carry onto itself,
\[
\begin{multlined}
Z\bigl[\mathbf 1_{\mathcal P}(P)\,\mathbf 1\{|\partial U-\mathbf 1|_\infty<\omega\}\bigr]\le10^{12}N'^{24}\,\mathsf X_0\int_{\mathcal P}\exp\Bigl[-\frac{2\beta'}{N'^5}\sum_{\mu<\nu}|\mathsf v(P_\mu)\times\mathsf v(P_\nu)|^2\Bigr]dP,\qquad\mathsf X_0 \\
=(2\pi^2)^{-(2\mathsf V-2)}\Bigl(\frac{2\pi}{\beta'}\Bigr)^{3\mathsf V-3}\mathfrak d'^{-3/2},
\end{multlined}
\]
$dP$ is Haar measure on $G^3$. (v) \emph{The normalizing lower bound.} Here we use \textbf{(Toron floor)}, namely $Z\ge10^{-6}N'^2\mathsf X_0/(\pi^2\beta'^2)$. The following coordinate identities describe its normalization. Write $P_1=\cos\alpha_1\mathbf1+i\sin\alpha_1\hat n\cdot\tau$, with $dP_1=(2/\pi)\sin^2\alpha_1\,d\alpha_1\,d\sigma(\hat n)$ and $\sigma$ probability area on $S^2$. Simultaneous conjugation of all three holonomies and all ordinary coordinates leaves the integral invariant, so the axis integral has mass one and can be represented by $\tau_3$. For $P_2,P_3$ use $P=e^{i\alpha\tau_3}e^{i(y_1\tau_1+y_2\tau_2)}$, $0\le\alpha<2\pi$, $|y|<\pi/2$, with $dP=(2\pi^2)^{-1}\{\sin(2|y|)/(2|y|)\}\,d\alpha\,d^2y$. At the axial triple $P^0=(e^{i\alpha_\mu\tau_3})_\mu$, the ordinary coordinates on $S_0$ and the four transverse coordinates $y,y'$ form the augmented variable $\xi=(a,y,y')$ of real dimension $6\mathsf V-2$. The linearized curl on these augmented coordinates, denoted by $D_{P^0}$, has Gaussian determinant $\mathfrak d'[4\sin^2\alpha_1\,\mathfrak d(2\alpha)]^2$: the parallel color contributes $\mathfrak d'$ and the two real transverse colors give the square of the determinant for the complex coordinates vanishing on $T$ and at $b_1^*$ alone. On $|a(b)|,|y|,|y'|\le r$, $r^2=120N'^5/\beta'$, the relative Haar factors obey \[ (1-r^2/3)^{2\mathsf V-2}(1-2r^2/3)^2\ge1-(2\mathsf V+2)r^2/3>0.99. \] Cancelling $\sin^2\alpha_1$ against its Gaussian determinant factor and integrating over $\alpha_1\in[\pi/4,3\pi/4]$, $\alpha_2,\alpha_3\in[0,2\pi)$ leaves the angular factor \[ (2\pi^2)^{-2}(2\pi)^2\int_{\pi/4}^{3\pi/4}\frac{2\sin^2\alpha_1}{\pi}\frac{d\alpha_1}{4\sin^2\alpha_1}=\frac1{4\pi^2}. \] These identities fix the Gaussian normalization; the quantitative lower bound on $Z$ remains the assumed \textbf{(Toron floor)}.

(vi) \emph{The toron integral.} Let $B=2\beta'/N'^5$, so that $B^{-1/2}\le w'\le\frac1{20}$. The triples with $\mathsf a(P_1,P_2)\in\mathsf Q$ carry at most $261\,w'^3B^{-3/2}$ of the integral of $\exp[-B\sum_{\mu<\nu}|\mathsf v(P_\mu)\times\mathsf v(P_\nu)|^2]$. Integrate $P_3$ first: with the singular values $s_1\ge s_2$ of the matrix $\mathsf v$ of the pair $(P_1,P_2)$, in the eigenframe of $\mathsf v\mathsf v^{\rm t}$, $\sum_{i=1,2}|\mathsf v(P_i)\times\mathsf v(P_3)|^2\ge s_2^2x_1^2+s_1^2(x_2^2+x_3^2)$, $x=\mathsf v(P_3)$, so that this integral is at most $\min\{1,(Bs_1^2)^{-1}\}\min\{1,3.81(Bs_2^2)^{-1/2}\}$; and $|\mathsf v(P_1)\times\mathsf v(P_2)|=s_1s_2$.

If $|\mathsf v(P_i)|<0.9$ for $i=1,2$ on the whole set, the Haar measure of the pair has in the coordinates $\mathsf s$ a density at most $\frac2{0.19\pi^2}\,4s_1s_2(s_1+s_2)$, the product of $4s_1s_2(s_1+s_2)$ with the bound of the $P_3$-integral is at most $30.5\,B^{-3/2}$, and the projection of $\mathsf Q$ onto the three coordinates $(\mathsf s_{11},\mathsf s_{12},\mathsf s_{22})$ has volume $8w'^3$. Else $\sin\alpha(P_i)\ge0.8$ on the set for one $i$, and the two angles alone give $12.6\,w'^2B^{-2}\le12.6\,w'^3B^{-3/2}$: in general the triples whose two angles lie in a set $\mathcal A$ carry at most $2.52\,B^{-2}\int_{\mathcal A}\min\{B^{1/2}\mathsf m,\mathsf m^{-1}\}\,d\alpha_1d\alpha_2$, $\mathsf m=\max_i\sin\alpha(P_i)$. (vii) By the ceiling of (iv) and \textbf{(Toron floor)}, $\mu_{N',\beta'}(P\in\mathcal P,\ |\partial U-\mathbf 1|_\infty<\omega)\le10^{18}\pi^2N'^{22}\beta'^2\int_{\mathcal P}\exp[-B\sum_{\mu<\nu}|\mathsf v(P_\mu)\times\mathsf v(P_\nu)|^2]\,dP$, and with (vi) the lemma follows, as $10^{18}\pi^2N'^{22}\beta'^2\cdot261\,w'^3(N'^5/2\beta')^{3/2}\le10^{21}N'^{29.5}\beta'^{1/2}w'^3$.

\emph{From one cutoff to the next.} In the setting (E) of Section~\ref{sec:cutoff-calibration}, with the side $\ell=2\mathsf L^m$ and the label $\mathsf b$, take at every integer $K\ge2$ the numbers of Section~\ref{sec:cutoff-calibration} at $n=1$: the depths $u=\lceil K/1260\rceil$ and $s=\lceil K/9\rceil$, $t_1=K-u$, and, with $N_{[d]}=\ell^3\mathsf L^{3d}$ and $\mathsf y_{[d]}=(\log(\mathsf L^d\mathsf b))^{2\mathfrak p_0}$, the four errors
\[
\begin{gathered}
\epsilon_{\rm fg}=C_{\rm old}C_{\rm fg}\tfrac{\mathsf L}{\mathsf L-1}E_0\,N_{[u]}\mathsf L^{-(s-u)},\qquad\\
\epsilon_{\rm cl}=4\mathsf K_{\rm w}c_\rho^2A_0^2\bigl(1+\tfrac{2\varrho}{\mathsf b}\bigr)\mathsf y_{[K]}N_{[s]}K^{1/2}\mathsf L^{-(K-s)/2},\\
\epsilon_{\rm sh}=3(s-u)N_{[s]}\exp\bigl[C_{\rm sh}-\tfrac14\kappa_{\rm b}A_1^2\mathsf y_{[u]}\bigr],\qquad\\
\mathsf R_{t_1}=2M^{-3}N_{[u]}\exp\bigl[-c'A_0^2\mathsf y_{[u]}+\bigl(\mathsf C^\Sigma_{\rm maj}+2+c_{\rm low}+\tfrac92\log\mathsf b\bigr)\ell^3\bigr],
\end{gathered}
\]
with the constants named there. The first two fall like powers of $\mathsf L^{-K}$, up to factors that are powers of $K$, $\ell$ and $\log\mathsf b$; the last two are exponentially small in $(\log(\mathsf L^u\mathsf b))^{2\mathfrak p_0}$. Corollary~\textup{\ref{IV-62-66}}(a) bounds the difference of the means of a gauge-invariant function at two consecutive cutoffs by $\mathfrak o_K=8\mathsf R_{t_1}+12\epsilon_{\rm fg}+8\epsilon_{\rm cl}+12\epsilon_{\rm sh}$ times its supremum; here its comparison is run with the errors that bound quotients of densities kept as factors, so that an event of a probability of the order of $t$ keeps this order.

\begin{lemma}[One step between consecutive cutoffs]\label{lem:cut}
In the setting (E), let $K\ge2$ be an integer at which the conditions (k0)--(k4) and (k6) of Section~\ref{sec:cutoff-calibration} hold at $n=1$, and put
\[
\epsilon^\times_K=4\epsilon_{\rm fg}+2\epsilon_{\rm cl}+1.06\,\epsilon_{\rm sh}+1.39\,\mathsf R_{t_1},\qquad\epsilon^+_K=\mathsf R_{t_1}+1.29\,\epsilon_{\rm sh}.
\]
Then for every integer $0\le d\le u$ and every Borel set $E$ of block fields of the depth $d$ that the gauge transformations carry onto itself (Definition~\textup{\ref{IV-62-62}}(3)),
\[
\mu^{K+1}_d(E)\le e^{\epsilon^\times_K}\mu^K_d(E)+\epsilon^+_K,\qquad\mu^K_d(E)\le e^{\epsilon^\times_K}\mu^{K+1}_d(E)+\epsilon^+_K .
\]
\end{lemma}

\emph{Sketch of the proof of Lemma~\ref{lem:cut}.} It suffices to take $d=u$: for $d<u$ replace $E$ by its preimage under the averaging maps from the depth $u$ to the depth $d$, which are the same on the two lattices (Section~\YMMref{116}{807}) and keep it Borel and gauge invariant (Proposition~\YMMref{4.5}{20}, \YMMref{4.5}{20}). Take the lattices $\mathrm A=(\ell,K)$ and $\mathrm B=(\ell,K+1)$ with the top levels $t_1$ and $t_1+1$, the objects of Section~\YMMref{116}{807} with $n=1$ and $f_1=1$, at the source $0$, and $\mathsf R=\mathsf R_{t_1}$, which is the same number on both (Section~\ref{sec:cutoff-calibration}). For a bounded Borel function $F\ge0$ of the block fields of the depth $u$ that the gauge transformations do not change, and for each lattice $\bullet$, of the cutoff $K^\bullet$ and the top level $t^\bullet$, let $\langle F\rangle^\bullet_0$ be the mean of $F(\mathcal M^{K^\bullet-u}U)$ at the cutoff $K^\bullet$, $\langle F\rangle^\bullet_1$ the normalised integral of $F$ against the measure of the chain on the inner ball $\mathsf{In}_{[u]}$, and $\langle F\rangle^\bullet_2$ that against the classical start on it. (1) \emph{The discarded mass, a remainder.} The measure of the chain is the image of $\mathfrak w^{\rm sf}_{t^\bullet}\rho_0\,dU$ under the averaging, with $\mathfrak w^{\rm sf}_{t^\bullet}\ge0$ (Proposition~\textup{\YMMref{115.2}{806}}, image identity); the right inequality of Proposition~\textup{\YMMref{115.2}{806}}, domination bound, for $F$ and for $\sup F-F$ composed with the averaging, which are gauge invariant (Proposition~\YMMref{4.5}{20}), and Proposition~\textup{\YMMref{115.2}{806}}, lost-mass bound, by which the discarded mass is at most $\mathsf R$ times the partition function, give $\langle F\rangle^\bullet_0\le\langle F\rangle^\bullet_1+\mathsf R\sup F$ and $\langle F\rangle^\bullet_1\le(1-\mathsf R)^{-1}\langle F\rangle^\bullet_0$. (2) \emph{Forgetting, a factor.} On $\mathsf{In}_{[u]}$ the two measures have the densities $e^{c^\bullet(0)}e^{\Phi^{\bullet,0}}$ and $e^{\hat c^\bullet(0)}e^{\hat\Phi^{\bullet,0}}$ (the proof of Proposition~\YMMref{116.1}{808}), both positive almost everywhere, as both measures are positive (Proposition~\textup{\YMMref{115.2}{806}}, image identity, Proposition~\YMMref{116.1}{808}); and the real parts of the two exponents differ by at most $\epsilon_{\rm fg}$ at every point of $\mathsf{In}_{[u]}$ (equation~\YMMref{(2180)}{803}).

So the quotient of the densities lies in $\varkappa^\bullet[e^{-\epsilon_{\rm fg}},e^{\epsilon_{\rm fg}}]$ with one number $\varkappa^\bullet>0$, which cancels between the numerator and the normaliser: $e^{-2\epsilon_{\rm fg}}\langle F\rangle^\bullet_2\le\langle F\rangle^\bullet_1\le e^{2\epsilon_{\rm fg}}\langle F\rangle^\bullet_2$. (3) \emph{The two classical starts, a factor and a remainder.} Both are integrals of $F(\mathcal M^{s-u}W)$ over the block fields $W$ of the depth $s$ against $\mathsf k^\bullet(W)\hat F^\bullet_0(W)\,dW$, $\mathsf k^\bullet\ge0$ (Section~\ref{sec:block-field-laws}); where $\mathsf k^{\rm A}$ and $\mathsf k^{\rm B}$ are both positive they are equal and $\hat F^{\rm B}_0=\hat F^{\rm A}_0e^{\mathsf q_0}$, $|\mathsf q_0|\le\epsilon_{\rm cl}$ (Proposition~\textup{\YMMref{116.1}{808}}, density comparison in the proof); and for $(\bullet,\circ)=(\mathrm A,\mathrm B)$ and $(\mathrm B,\mathrm A)$ the fields at which $\mathsf k^\circ$ vanishes carry at most the fraction $\epsilon_{\rm sh}$ of the mass of the start of $\bullet$ (Section~\ref{sec:block-field-laws}).

So $\langle F\rangle^\bullet_2\le e^{2\epsilon_{\rm cl}}(1-\epsilon_{\rm sh})^{-1}\langle F\rangle^\circ_2+\epsilon_{\rm sh}\sup F$. Chaining (1), (2), (3), (2), (1),
\[
\langle F\rangle^\bullet_0\le\frac{e^{4\epsilon_{\rm fg}+2\epsilon_{\rm cl}}}{(1-\epsilon_{\rm sh})(1-\mathsf R)}\langle F\rangle^\circ_0+\bigl(\mathsf R+e^{2\epsilon_{\rm fg}}\epsilon_{\rm sh}\bigr)\sup F,
\]
and (k2), (k3), (k6) give $e^{2\epsilon_{\rm fg}}\le e^{1/4}<1.29$, $-\log(1-\epsilon_{\rm sh})\le1.06\,\epsilon_{\rm sh}$ and $-\log(1-\mathsf R)\le2\log2\cdot\mathsf R\le1.39\,\mathsf R$. At $F=\mathbf 1_E$ this is the lemma.

\emph{The chain from a low cutoff.} Let $p=2\mathfrak p_0$, $y=\log\mathsf b$, $\bar\ell(y)=2\mathsf L^2My^r$, $\mathsf k(y)=\bar\ell(y)^3y^{p+1}$, $\mathsf N(y)=\mathsf L\,\bar\ell(y)^{31}y^{10(p+1)}$, and let $K_1=K_1(y)$ be the least integer $K\ge10$ with $\mathsf L^{K/10}\ge\mathsf k(y)$. With the constants of Section~\ref{sec:cutoff-calibration} and the number $\mathsf C_+$ of Step~4 put
\[
\begin{gathered}
a_1=C_{\rm old}C_{\rm fg}\tfrac{\mathsf L}{\mathsf L-1}E_0\mathsf L^4,\qquad\\
 a_2=4\mathsf K_{\rm w}c_\rho^2A_0^2(1+2\varrho)\mathsf L^{7/2},\qquad\\
 a_3=\sum_{K\ge1}(1+K\log\mathsf L)^pK^{1/2}\mathsf L^{-K/90},\\
a_4=529200\,e^{C_{\rm sh}}\sum_{i\ge1}i\,\mathsf L^{420i}e^{-\frac14\kappa_{\rm b}A_1^2(i\log\mathsf L)^p},\qquad\\
 a_5=2520\,M^{-3}\sum_{i\ge1}\mathsf L^{3i}e^{-c'A_0^2(i\log\mathsf L)^p},\qquad\\
 a_6=\frac{256\,\mathsf K_uc_\rho^2A_0^2\mathsf L^4}{\gamma_{\rm r}A_1^2},\\
\mathsf R(y)=1.29\,\bar\ell(y)^3\Bigl[\\
a_4e^{-\frac14\kappa_{\rm b}A_1^2y^p}\\
+a_5e^{(\mathsf C_++\frac92y)\bar\ell(y)^3-c'A_0^2y^p}\Bigr];
\end{gathered}
\]
these are formed from the constants alone, and the three series converge, as $p\ge2$, $\kappa_{\rm b},c'>0$ (Proposition~\YMMref{114.4}{803} and Section~\ref{sec:cutoff-calibration}) and $A_0,A_1>0$ (equation~\YMMref{(634)}{339}). Suppose that (E) holds, $\ell\le\bar\ell(y)$, $\mathsf b\ge\Theta_\star$, the number of the condition (k1), $y\ge\max\{8a_1,20a_2a_3,a_6\}$ and $\mathsf R(y)\le e^{-3y^2}$. Then (k0)--(k4) and (k6) hold at every integer $K\ge K_1$, and the sums $\mathsf S=\sum_{K\ge K_1}\epsilon^\times_K$ and $\mathsf R_\Sigma=\sum_{K\ge K_1}\epsilon^+_K$ satisfy
\[
\mathsf S\le\frac{17.72\,a_1+2a_2a_3}y+1.39\,\mathsf R(y),\qquad\mathsf R_\Sigma\le\mathsf R(y)\le e^{-3y^2},\qquad\ell\mathsf L^{K_1}<\mathsf N(y).
\]
Indeed, for $K\ge K_1$ one has $u<\frac K{1260}+1$ and $\frac K9\le s<\frac K9+1$, so $4u-s\le4-\frac K{10}$, $\frac{7s-K}2\le\frac72-\frac K{10}-\frac K{90}$ and $4s-K\le4-\frac K{10}$; and $\ell^3\mathsf L^{-K/10}\le\ell^3/\mathsf k(y)\le y^{-p-1}$. Hence $\epsilon_{\rm fg}\le a_1\ell^3\mathsf L^{-K/10}$, $\epsilon_{\rm cl}\le a_2y^p(1+K\log\mathsf L)^pK^{1/2}\mathsf L^{-K/90}\ell^3\mathsf L^{-K/10}$, and the left member of (k4) is at most $a_6\gamma_{\rm r}A_1^2\ell^3\mathsf L^{-K/10}$. Summed over $K\ge K_1$, as $(1-13^{-1/10})^{-1}\le4.43$, these give $\sum_{K\ge K_1}(4\epsilon_{\rm fg}+2\epsilon_{\rm cl})\le(17.72\,a_1+2a_2a_3)/y$. (k0) holds as $K\ge10$, (k1) as $\mathsf b\ge\Theta_\star$, and $y\ge\max\{8a_1,20a_2a_3,a_6\}$ gives (k2), (k4) and the first half of (k3).

For the two other errors, $(y+u\log\mathsf L)^p\ge y^p+(u\log\mathsf L)^p$, $s-u\le140u$, $N_{[s]}\le\ell^3\mathsf L^{420u}$, and at most $1260$ cutoffs have the same $u$: $\sum_{K\ge K_1}\epsilon_{\rm sh}\le a_4\ell^3e^{-\frac14\kappa_{\rm b}A_1^2y^p}$ and $\sum_{K\ge K_1}\mathsf R_{t_1}\le a_5\ell^3e^{(\mathsf C_++\frac92y)\ell^3-c'A_0^2y^p}$. So $\mathsf R_\Sigma\le\mathsf R(y)$, which gives the second half of (k3) and (k6), and $\sum_{K\ge K_1}(1.06\,\epsilon_{\rm sh}+1.39\,\mathsf R_{t_1})\le1.39\,\mathsf R(y)$. Finally $\mathsf k(y)>\mathsf L$, so $K_1\ge11$ and $\mathsf L^{(K_1-1)/10}<\mathsf k(y)$: $\mathsf L^{K_1}<\mathsf L\,\mathsf k(y)^{10}$ and $\ell\mathsf L^{K_1}<\mathsf N(y)$. Lemma~\ref{lem:cut} at $d=0$ and an induction on $K$ give $\mu^K_0(E)\le e^{\mathsf S_K}(\mu^{K_1}_0(E)+\mathsf R_K)$, with $\mathsf S_K$ and $\mathsf R_K$ the sums of the $\epsilon^\times_{K'}$ and of the $\epsilon^+_{K'}$ over $K_1\le K'<K$; so \eqref{eq:lowcut} holds: for every integer $K\ge K_1$ and every Borel set $E$ of block fields of the depth $0$ that the gauge transformations carry onto itself
\[
\mu^K_0(E)\le e^{\mathsf S}\bigl(\mu^{K_1}_0(E)+\mathsf R_\Sigma\bigr).
\]

\emph{The parameters.} Let $\mathsf y_\star$ be the least real number 
\[
y_0\ge\max\{122,\log\gamma^{-2},\log\Theta_\star,8a_1,20a_2a_3,a_6,2(17.72\,a_1+2a_2a_3)\}
\]
 such that at every $y\ge y_0$
\[
\begin{gathered}
\textup{(y1)}\ \mathsf R(y)\le e^{-3y^2};\qquad\\
\textup{(y2)}\ \log\mathsf L+30\log\bar\ell(y)+10(p+1)\log y\le y;\qquad\\
\textup{(y3)}\ 26\,\mathsf N(y)^{28}\,\bar\ell(y)\,y^4e^{-y}\le2.5\cdot10^{-14};\\
\textup{(y4)}\ 26\,\mathsf N(y)^6\,\bar\ell(y)\,y^2e^{-y}\le(c^{\rm SO}_2)^2;\qquad\\
\textup{(y5)}\ 31\,(y+1)^2\,\mathsf N(y)^6\le e^{2y^2},
\end{gathered}
\]
$c^{\rm SO}_2$ the number of Lemma~\ref{lem:background-iterated-average}. It exists: $\bar\ell$ and $\mathsf N$ are powers of $y$, so the left member of (y2) is $O(\log y)$, those of (y3) and (y4) tend to $0$, that of (y5) is a power of $y$, and the two exponents in $\mathsf R(y)e^{3y^2}$ are $-\frac14\kappa_{\rm b}A_1^2y^p+3y^2+O(\log y)$ and 
\[
-c'A_0^2y^p+(\mathsf C_++\frac92y)\bar\ell(y)^3+3y^2+O(\log y),
\]
which tend to $-\infty$, as $p=2\mathfrak p_0>3r+2$ at the chosen constants (the choice $\mathfrak p_0\ge6r+1$ in Proposition~\YMMref{69.1}{595}) and $r\ge2$; each condition compares continuous functions, so the numbers $y_0$ at which all hold on $[y_0,\infty)$ form a closed half line. Put $t_\star=2\mathsf Le^{-\mathsf y_\star}$. Now let $j\ge1$ and $\beta>0$ with $t=2\mathsf L^j/\beta\le t_\star$, and $N=2\mathsf L^j$; let $m$ be the least integer $m\ge1$ with $M\,R(2\mathsf L^m/t)\le\mathsf L^m$, $\ell=2\mathsf L^m$, $\mathsf b=\ell/t$, $y=\log\mathsf b$ and $K=j-m\in\mathbb Z$, so that $N=\ell\mathsf L^K$ and $\beta=\mathsf L^K\mathsf b$; and $q=\lceil y\rceil$.

Then $y\ge\log(2\mathsf L/t)\ge\mathsf y_\star$; as in Step~4, (E) holds for $(\mathsf b,m)$ and $\ell\le\bar\ell(y)$; by (y2), $\log\ell\le y/30$, so that $\log\frac1t\le y\le2\log\frac1t$, and $y\le q<y+1$, $q\ge122$; $e^{-y}\le t=\ell e^{-y}\le\bar\ell(y)e^{-y}$; the chain holds from $K_1=K_1(y)$, with $e^{\mathsf S}\le e^{0.53}<1.7$ and $\mathsf R_\Sigma\le e^{-3y^2}$; and $\beta_1=\mathsf L^{K_1}\mathsf b=\ell\mathsf L^{K_1}/t$ has $\log\beta_1=y+K_1\log\mathsf L<y+\log\mathsf L+10\log\mathsf k(y)\le2y$, by (y2).

\emph{Bare cubes.} Let $(N',\beta')$ be $(\ell\mathsf L^{K_1},\beta_1)$, the bare cube at the cutoff $K_1$, or $(N,\beta)$ if $N<\ell\mathsf L^{K_1}$. Then $N'$ is of the form $2\mathsf L^{j'}$, $j'\ge1$, so even, $26\le N'<\mathsf N(y)$, $\beta'=N'/t$, and $\log\beta'\le2y$. Put $\Lambda'=3N'^3\log(8e^8\beta'^{3/2})+\log(3N'^3)+3y^2$, $\omega_0=(2\Lambda'/\beta')^{1/2}$ and $\upsilon=N'^2\omega_0$. The Haar measure of $\{g\in G:|g-\mathbf 1|\le\beta'^{-1/2}\}$ is at least $\frac18\beta'^{-3/2}$, and $\beta'\sum_pa_p\le24N'^3$ when all bonds lie in this set; so $Z(N',N',N')\ge(8e^8\beta'^{3/2})^{-3N'^3}$, every plaquette has $|\partial U(p)-\mathbf 1|\ge\omega_0$ with a probability at most $Z^{-1}e^{-\beta'\omega_0^2/2}$, and
\[
\mu_{N',\beta'}\bigl(|\partial U-\mathbf 1|_\infty\ge\omega_0\bigr)\le3N'^3\,(8e^8\beta'^{3/2})^{3N'^3}e^{-\Lambda'}=e^{-3y^2}.
\]
Further $N'^3\le\Lambda'\le13N'^3y^2$, so $\upsilon^2\le26N'^6y^2t\le26\,\mathsf N(y)^6\bar\ell(y)y^2e^{-y}$; and (y3), (y4) give $\omega_0\le1.6\cdot10^{-7}N'^{-13}$, $q\upsilon\le1.7\cdot10^{-7}$, $\beta'\ge N'e^y/\bar\ell(y)\ge1.1\cdot10^9N'^{16}$ and $\upsilon\le c^{\rm SO}_2$. With $w=2.03\,q\upsilon$ and $w'=\frac32w+1.03\,\upsilon$ one has $(N'^5/2\beta')^{1/2}\le3\upsilon\le w'\le4q\upsilon\le\frac1{20}$, and, under \textbf{(Toron floor)}, Lemma~\ref{lem:anti} at $\omega=\omega_0$ gives, for every site $\xi'$ and every closed cube $\mathsf Q$ of the half side $w'$,
\[
\mu_{N',\beta'}\bigl(\mathsf a(U,\xi')\in\mathsf Q,\ |\partial U-\mathbf 1|_\infty<\omega_0\bigr)\le10^{21}N'^{29.5}\beta'^{1/2}(4q\upsilon)^3\le P^{\rm b}:=8.5\cdot10^{24}\,q^3y^3\,\mathsf N(y)^{39}\,t,
\]
as $(2\Lambda')^{3/2}\le132.6\,N'^{4.5}y^3$ and $N'^{40}/\beta'=N'^{39}t$.

\emph{The cutoffs $K\ge K_1$.} Take the swap with the block field $V=\mathcal M^KU$ of the depth $0$, $\mathsf n=\ell$, and $\omega=\theta_{\mathsf L}\mathsf L^{2K_1}\omega_0$, with $\omega_0$ and $\upsilon$ those of $(\ell\mathsf L^{K_1},\beta_1)$ and the number $\theta_{\mathsf L}=\exp\frac1{2(\mathsf L^2-1)}<1.003$ of Lemma~\ref{lem:background-iterated-average}: $\mathsf n^2\omega=\theta_{\mathsf L}\upsilon$, so $\delta_q\le e^{q\theta_{\mathsf L}\upsilon}-1\le1.015\,q\upsilon$, $2\delta_q\le w$ and $\omega\le1$. The sets $E_p=\{|\partial V(p)-\mathbf 1|\ge\omega\}$ and $\{\mathsf a(V,\xi')\in c^+_k\}$ are Borel sets of block fields of the depth $0$ that the gauge transformations carry onto themselves, as these conjugate $\partial V(p)$, and the two winding holonomies at $\xi'$ by one element; so the chain gives $\mu^K_0(E)\le1.7\,(\mu^{K_1}_0(E)+e^{-3y^2})$ for each of them.

At the cutoff $K_1$, where $|\partial U-\mathbf 1|_\infty<\omega_0$ on the bare cube: $\mathsf L^{2K_1}\omega_0\le\upsilon\le c^{\rm SO}_2$, so $|\partial V-\mathbf 1|_\infty<\omega$, and the winding holonomies of $V=\mathcal M^{K_1}U$ differ from the bare ones along the same lines by at most $13.3\,\ell\mathsf L^{2K_1}\omega_0=13.3\,\upsilon/\ell\le0.512\,\upsilon$ (Lemma~\ref{lem:averaged-line-holonomy}), so that $|\mathsf a(V,\xi')-\mathsf a(U,\xi')|_\infty\le1.03\,\upsilon$; and $c^+_k$ lies in a closed cube of the half side $\frac32w$. With the bounds on the bare cube $(\ell\mathsf L^{K_1},\beta_1)$ this gives, for the swap, $\rho\le3.4\,e^{-3y^2}$ and
\[
P_\ast\le1.7\bigl(P^{\rm b}+2e^{-3y^2}\bigr)\le\Pi(y)\,t,\qquad\Pi(y)=1.6\cdot10^{25}y^6\,\mathsf N(y)^{39},
\]
as $q^3\le1.025\,y^3$ and $e^{-3y^2}\le e^{-y}\le t$; and $9q^2\mathsf n^6\rho\le30.6\,(y+1)^2\bar\ell(y)^6e^{-3y^2}\le e^{-y^2}$, by (y5).

\emph{The cutoffs $K<K_1$}, among them those with $K\le0$. Take the swap with the bare field on the cube of $N<\ell\mathsf L^{K_1}$ sites a side, $\mathsf n=N$, and $\omega=\omega_0$ of $(N,\beta)$: $\mathsf n^2\omega=\upsilon$, $2\delta_q\le w$, $\rho\le e^{-3y^2}$, $P_\ast\le P^{\rm b}+e^{-3y^2}$, and $9q^2N^6\rho\le e^{-y^2}$; the same bounds hold.

\emph{The end.} As $\Pi(y)\ge1$, $t\ge e^{-y}$ and $q-1\le y$, $(\Pi(y)t)^{q-1}\ge e^{-y^2}$; so in all cases $e^{-\mathsf d_q(N)}\le486\,P_\ast^{q-1}+9q^2\mathsf n^6\rho\le487\,(\Pi(y)t)^{q-1}$. Where $\Pi(y)t<1$, as $q\ge y\ge\log\frac1t$ and $q\ge122$,
\[
\mathsf G_N\ge\frac{\mathsf d_q(N)}q\ge\frac{q-1}q\Bigl(\log\frac1t-\log\Pi(y)\Bigr)-\frac{\log487}q\ge\log\frac1t-\log\Pi(y)-1.06,
\]
that is $\lambda_0(N,N)^N/Z(N,N,N)=e^{-\mathsf G_N}\le2.89\,\Pi(y)\,t$. As $\mathsf N(y)^{39}=\mathsf L^{39}(2\mathsf L^2M)^{1209}y^{1209r+390p+390}$ and $y\le2\log\frac1t$, this is the bound of (d), with
\[
\mathsf p_\star=1209\,r+780\,\mathfrak p_0+396,\qquad\mathsf C_\star=4.9\cdot10^{25}\,\mathsf L^{39}(2\mathsf L^2M)^{1209}\,2^{\mathsf p_\star},
\]
formed from $\mathsf L$, $M$, $r$ and $\mathfrak p_0$, and with $t_\star$, formed from the constants. Where $\Pi(y)t\ge1$ the bound holds as well, as $\mathsf C_\star(\log\frac1t)^{\mathsf p_\star}\ge2.89\,\Pi(y)$ and $\lambda_0(N,N)^N\le Z(N,N,N)$. With $\mathsf e(\mathsf L^j)=\mathsf e(N/2)\ge\mathsf G_N$ this gives the floor of $\mathsf e$. The last sentence of (d) is the first at $j=m_0+K$, as $N=\ell_0/a_K$ and $t=\ell_0/\mathsf b$.
\end{proof}

\begin{proof}[Proof of Proposition~\ref{prop:sheets}]
(a) With the time $\mu_0$ and the section $(N_1,N_2)$, $Z(t,N_1,N_2)=(1+q^t)^{N_1N_2}\mathsf c^t=\sum_k\binom{N_1N_2}k(\mathsf cq^k)^t$, $\mathsf c=(1+q^{N_1})^{N_2}(1+q^{N_2})^{N_1}$: (T2), and (T3) with all the weight in the sector $00$; (T1) and (iii) are read off. With $f(u)=\log(1+q^u)$: $F(l,l,l)=3l^2f(l)$, $F(2l,l,l)=l^2f(2l)+4l^2f(l)$, $F(2l,2l,l)=4l^2f(2l)+4l^2f(l)$, $F(2l,2l,2l)=12l^2f(2l)$; so $d_1=l^2\mathsf h$, $d_2=2l^2\mathsf h$, $\mathsf e=4l^2\mathsf h$ with $\mathsf h=2f(l)-f(2l)$, and $X=12l^2\mathsf h$. (b) Expanding the two products of the factors that contain $t$,
\[
Z^{(k;m)}=\sum_{\mathsf a,\mathsf b}\binom{N_2}{\mathsf a}\binom{N_1}{\mathsf b}(-1)^{k_1\mathsf a+k_2\mathsf b}\bigl(\mathsf c_mq^{N_1\mathsf a+N_2\mathsf b}\bigr)^t,\qquad\mathsf c_m=1+(-1)^mq^{N_1N_2}>0,
\]
which is (T3) with the sector $e=(\mathsf a,\mathsf b)\bmod2$; (T1), (T2) and (iii) are read off. The formulas for $\mathsf e$ and $X$ follow from $F(\underline n)=\sum_\rho n_\rho\log(1+q^{n_\mu n_\nu})$; then $X\ge24l(s-\frac{s^2}2)-6ls^4$ and $\mathsf e\le8ls^2$, so $X/\mathsf e\ge\frac3s(1-\frac s2-\frac{s^3}4)\ge2.95/s$ for $s\le0.03$. (c) On the cube, $Z^-/Z=(\frac{1-u}{1+u})^n$ with $u=q^{n^2}$, so $\frac12q^{n^2}\le z_-(n,n,n)\le nq^{n^2}$, while $x_0=\mathsf e(l)=8ls^2(1+O(s^2))$ gives $\theta=103.2\,ls^2(1+O(s^2))$ and the two expansions, with the exact constant $103.2e^{-0.835}$. If the inequality held with $c'>\frac12$, then $\log\frac1q\ge2c'\log\frac1q-C'(l)$ for all small $q$, which is false.
\end{proof}

\begin{proof}[Proof of Proposition~\ref{prop:flat}]
Put $u_j=l_j\varepsilon_{\rm v}(l_j)$ and $\mathsf t_j=\theta^{2^{j+1}}$. By Proposition~\ref{prop:tub}(a), $u_0\le0.01643\,\theta^2$. Let $u_j\le0.0165\cdot2^{-j}\theta^{2\cdot4^j}$, and take $\hat N=2^{j+2}$ in Proposition~\ref{prop:tub}(b). Then the powers of $\theta$ match: $y=\hat Nu_j\le0.066\,\theta^{2\cdot4^j}$, $(0.001615\,\mathsf t_j)^{2\hat N}\le0.001615^8\theta^{2\cdot4^j}$, $(0.001615\,\mathsf t_j^2)^{\hat N/2}\le0.001615^2\theta^{2\cdot4^{j+1}}$, and $8\hat N\cdot0.01655\,\mathsf t_j=16\cdot0.01655\,\mathsf m\theta^{\mathsf m}\le0.835$ with $\mathsf m=2^{j+1}$, as $\mathsf m\,0.8901^{\mathsf m}\le3.153$. So $\bar z_1\le0.06601\,\theta^{2\cdot4^j}$, $\bar z_2\le4.38\cdot10^{-5}\theta^{2\cdot4^{j+1}}$, and $u_{j+1}\le2^{-j-1}\cdot9.04\cdot10^{-5}\,\theta^{2\cdot4^{j+1}}$, which closes the induction.

For the cubes let $n=2^il$, $i\ge1$, $j=\lfloor(i-1)/2\rfloor$ and $M=n/l_j$, an integer with $M^2\le n/l\le2M^2$. The first inequality of Lemma~\ref{lem:vort}(b) on $(l_j,n,l_j)$, with $\mathsf r_n(l_j,l_j)\le(0.001615\,\mathsf t_j)^M\le0.001615\,\theta^{l_j^2/2l^2}$, gives $z_-(l_j,n,l_j)\le0.00906\,M\theta^{l_j^2/2l^2}$; and $\delta(l_j,n,l_j)\le M\delta_{\rm c}(l_j)\le0.01655\,M\theta^{l_j/l}\le0.01655\,y\theta^y\le0.0523$ with $y=\sqrt{n/l}$, as $M\le y\le l_j/l$ and $y\theta^y\le1/e\Lambda$. Lemma~\ref{lem:vort}(a) with $\kappa=M^2$ gives $z_-(n,n,n)\le(0.0096\,M)^{M^2}\theta^{n^2/2l^2}$, and the first factor is $\le1$ if $0.0096\,M\le1$, else at most $(0.0096\sqrt{n/l})^{n/l}$. Lemma~\ref{lem:two}(d) on these cubes and Theorem~\ref{thm:one}(a) give the area law, as $\frac{RT}{nl}\log(0.0096\sqrt{n/l})\to0$; and $\log\frac1{0.8901}/2\ge0.0582$.
\end{proof}

\subsection{Scope of the conclusions}

Every statement that uses the criterion is an implication from it on one lattice. Proposition~\ref{prop:where} gives the criterion in two ranges of the coupling, both below $\frac16$: at every integer side $l\ge3$ for $0<\beta\le x_0/48l^3$ (its clause (a)), and at the even sides $l\ge4$ with $64\,l^2(6\beta)^{l/2}\le\frac1{20}$ for $0<\beta<\frac16$ (its clause (b)); from a side at which it holds the criterion passes, at the same coupling, to every integer side $l'\ge16l$, of either parity (Proposition~\ref{prop:says}(c)), and to the sides $2^kl$ (Theorem~\ref{thm:crit}(b)). With $\beta=L^K\mathsf b$ and $l=l_*L^K$ the first range reads $\mathsf b\le x_0L^{-4K}/48l_*^3$ and the second $\mathsf b<L^{-K}/6$, while $L^K\mathsf b\ge\frac16$ for all large $K$ at every label $\mathsf b>0$; so neither range, nor its passage to other sides at the same coupling, contains a fixed label as $K$ grows.

The uniform criterion is the hypothesis of Theorem~\ref{thm:phys} and of Corollary~\ref{cor:pot}(e). The number $\sigma_*$ is a floor valid at each cutoff. On the centre twists of the cubes of sides $n\ge2\mathsf ql$ the floors are those of the pure numbers $-\log z_-(n,n,n)\ge\sigma^+n^2-\frac{2\mathsf qn}l(\frac12\log\frac{2\mathsf qn}l+0.86)$ and $-\log(1-Z^z/Z)$, which is at least the same number less $\log6$ (Theorem~\ref{thm:one}(d); in physical units Theorem~\ref{thm:phys}(b)); the free energy $-\log(Z^z/Z)$ of a twist itself lies between $0$ and $-\log(1-6z_-)$ where $6z_-<1$, and the free energy $\varepsilon_{\rm v}$ of a vortex line per unit length is bounded from above (Propositions~\ref{prop:tub}, \ref{prop:flat}).

The energies of the flux sectors and of the flux lines, $\log(\lambda_0/\Lambda_e)$ and $\log(\lambda_0/\Lambda_-)$, have the floors of Theorem~\ref{thm:one}(e), (g) and Theorem~\ref{thm:phys}(b), (d), and the gap has those of Theorem~\ref{thm:one}(f) and Theorem~\ref{thm:phys}(c). For a path that lies in no coordinate plane, the area in Theorem~\ref{thm:one}(c) is that of a projection counted modulo $2$, which may vanish although every surface bounded by the path is large: two parallel squares of side $s$ at the distance $s$, joined by a segment run twice, form a path $C$ with $\mathsf A_{\rho'\rho}(C)=0$ for all $\rho'\ne\rho$, while every set of plaquettes whose boundary modulo $2$ is the path has at least $2s^2$ members.

For two static charges at sites $o$, $r$ of a section that lie on no common line of the lattice, Theorem~\ref{thm:one}(e) gives the floor $\sigma^+|r-o|_\infty$, with the maximum norm, which is at least $2^{-1/2}\sigma^+|r-o|$.

\section{From the box criterion to continuum confinement}\label{C-361-1}

A box inequality at one fixed physical size, imposed at every sufficiently large cutoff along a fixed label, controls the infrared limit. Under this hypothesis the periodic Gibbs state is unique at each such cutoff, the continuum limit is independent of the order of the volume and cutoff limits, and its reconstructed vacuum is unique with a positive Hamiltonian gap.

The fixed-lattice transfer and tube arguments are given in Section~\ref{sec:conf}. Here we prove the comparisons uniform in the cutoff, passage of the gap to the continuum, transport between labels and the single-cutoff certificates with their required thresholds and margins. The box inequality itself is an additional input; the elementary small-coupling examples do not verify it along a fixed positive label.

\subsection{Notation and hypotheses}

\emph{One lattice.} Let $\beta\ge0$. A \emph{box} $\underline n=(n_1,n_2,n_3)$ is the periodic lattice with $n_\mu\ge3$ sites along the side $\mu$; $|\underline n|=n_1n_2n_3$, and $N^3$ is the cube. Every bond carries the Haar probability measure of $\mathrm{SU}(2)$, every plaquette $p$ the variable $a_p=1-\frac12\operatorname{tr}\partial U(p)\in[0,2]$, and the box the weight $\exp(-\beta\sum_pa_p)$; $Z(\underline n)$ is its partition function, $\mu_{\underline n}$ its measure, $F=\log Z\le0$, $P=\inf_{\underline n}F(\underline n)/|\underline n|\in[-6\beta,0]$, and
\[
\delta(\underline n)=F(\underline n)-|\underline n|P\ge0,\qquad D_q(N)=q^3F(N^3)-F((qN)^3),\qquad\mathsf e(l)=2F(l,2l,2l)-F((2l)^3).
\]
We call $\delta$ the \emph{excess}. For a box and a direction $i$, $\mathsf d^{(i)}(\underline n)=2F(\underline n)-F(\underline n')$, $\underline n'$ the box with the side $i$ doubled. None of these numbers depends on a spacing. For a bounded function $\Psi$ of the bonds of $a\mathbb Z^3$ and a box, $\langle\Psi\rangle_{\underline n}$ is the mean of $\Psi$ under the measure of the box extended periodically to $a\mathbb Z^3$.

\emph{The transfer operator.} Fix a coordinate direction $\mu_0$ as the time, the two others $\nu_1<\nu_2$, and a cross-section $\Sigma=(A,B)$, $A,B\ge3$; the box with $t$ layers over $\Sigma$ is $(t,\Sigma)$. By Section~\YMMref{85}{677} and Proposition~\YMMref{85.1}{678} the configurations of a layer form a compact space $\mathcal X$ with the product Haar measure $dV$, a layer carries the continuous action $A_\Sigma\ge0$ of its plaquettes, and the crossing operator $\mathsf K$, which integrates the time-like bonds between two layers, is positive, Hilbert--Schmidt, with a continuous kernel with values in $(0,1]$. With $\mathsf m=e^{-\beta A_\Sigma}$ (the function, and the operator of multiplication) we put
\[
\mathsf D=\mathsf D_\Sigma=\mathsf m^{1/2}\mathsf K\,\mathsf m^{1/2},\qquad\mathsf T'=\mathsf K^{1/2}\mathsf m\,\mathsf K^{1/2}.
\]
$\mathsf D$ is positive, Hilbert--Schmidt, with the continuous kernel $\mathsf m(V')^{1/2}\mathsf K(V',V)\mathsf m(V)^{1/2}\in(0,1]$, and $Z(t,\Sigma)=\operatorname{Tr}\mathsf D^t$ for every integer $t\ge3$ (Proposition~\ref{C-361-7}). $\mathsf T'$, which stands with the slab operators of the trace formula (M1) below, is positive as well, as $\mathsf K\ge0$ and $\mathsf m>0$, and has the same nonzero eigenvalues $\lambda_0\ge\lambda_1\ge\dots>0$, with multiplicities; they depend neither on $t$ nor on the spacing. We put $\mathsf r_t(\Sigma)=\sum_{i\ge1}(\lambda_i/\lambda_0)^t$, $\mathsf G(t;\Sigma)=\log(1+\mathsf r_t)$, $\mathsf R_t=1-e^{-\mathsf G}$, $\varepsilon(\Sigma)=\log\lambda_0-|\Sigma|P$, and on a cube $\mathsf r_N=\mathsf r_N(N,N)$, $\mathsf G_N=\mathsf G(N;N,N)$.

\emph{The transfer operator of a state.} A periodic Gibbs state $\nu$ on $a\mathbb Z^3$ is invariant under $a\mathbb Z^3$ and reflection positive at every lattice mirror (M3), so it has the picture $(e_{\mu_0},\frac a2,ae_{\mu_0})$ defined immediately before Lemma~\YMMref{99.1}{737}. On the bounded Borel functions $F,F'$ of finitely many bonds with both end points in $\{x_{\mu_0}\ge a\}$ the form $\langle F,F'\rangle=\nu(\overline{\Theta_0F}\,F')$, $\Theta_0$ the reflection $x_{\mu_0}\mapsto a-x_{\mu_0}$, is positive; $\mathcal H(\nu)$ is the completion of its quotient by the null vectors, $[F]$ the class of $F$, $\Omega=[1]$, and $P_\Omega$ the projection on $\mathbb C\Omega$. The \emph{transfer operator} of $\nu$ is the operator $T_\nu$ with $T_\nu[F]=[\tau_{ae_{\mu_0}}F]$ of Lemma~\YMMref{99.1}{737}: it is self-adjoint, $0\le T_\nu\le1$, $T_\nu\Omega=\Omega$, and $\langle[F],T_\nu^n[F']\rangle=\nu(\overline{\Theta_0F}\,\tau_{nae_{\mu_0}}F')$, which is the number $h_{AB}(n)$ of Proposition~\ref{C-361-10}(f) for $A=\overline{\Theta_0F'}$ and $B=\overline{\Theta_0F}$, functions of the slab of the layers $x_{\mu_0}\in\{(1-\vartheta)a,\dots,0\}$ when $F,F'$ are functions of the bonds with both end points in the layers $a,\dots,\vartheta a$, $\vartheta\ge2$; there $h_{AA}(0)-|\nu(A)|^2=\|(1-P_\Omega)[F']\|^2$.

So, if $\mathsf e(l)\le x_0\le0.069$ at $\beta$ for one $l\ge3$, that proposition at $n_0=0$ gives, with $\theta=12.9\,x_0$, $|\langle[F],T_\nu^n[F']\rangle-\langle[F],\Omega\rangle\langle\Omega,[F']\rangle|\le\theta^{n/l}\|(1-P_\Omega)[F]\|\,\|(1-P_\Omega)[F']\|$; the classes being dense, $\|T_\nu^n-P_\Omega\|\le\theta^{n/l}$, hence $\ker(1-T_\nu)=\mathbb C\Omega$ and $\|T_\nu\|\le\theta^{1/l}$ on $\Omega^\perp$.

\emph{Twists and loops.} A \emph{column} of the plane $(\mu\nu)$ is the set of the plaquettes $p_{\mu\nu}(y+je_\rho)$, $j\in\mathbb Z$, $\rho$ the third direction, in lattice units. For $z\in\{0,1\}^3$, indexed by the three coordinate planes, $Z^z$ is the partition function with the sign of $\operatorname{tr}\partial U(p)$ reversed on one column of each plane with $z=1$; it does not depend on the columns chosen, and $r^z=Z^z/Z$. On a cube $r^z$ depends only on the number of twisted planes; we write $r_1,r_2,r_3$. With a time $\mu_0$ we write $z=(k;m)$, $k\in\{0,1\}^2$ recording the planes $(\mu_0\nu_i)$ and $m$ the plane $(\nu_1\nu_2)$; the weights of the four classes of electric flux $f\in\{0,1\}^2$ are $w_f=\frac14\sum_k(-1)^{f\cdot k}r^{(k;0)}$.

For a closed lattice path $C$, $W_C=\frac12\operatorname{tr}U(C)$, $U(C)$ the ordered product along $C$; $W_{R\times T}$ is the loop of a coordinate rectangle with $R\times T$ bonds. A \emph{staircase loop} of the plane $(13)$ is given by a site $o$, an integer $T\ge1$ and integers $\alpha_j$ and $R_j\ge1$, $0\le j<T$: with $o^k_j=o+\alpha_jae_1+kae_3$ and $r^k_j=o+(\alpha_j+R_j)ae_1+kae_3$, the loop $C(\alpha,R)$ runs from $o^0_0$ straight to $r^0_0$; then, for $j=0,\dots,T-1$ in turn, one step up from $r^j_j$ to $r^{j+1}_j$ and, if $j<T-1$, straight to $r^{j+1}_{j+1}$; from $r^T_{T-1}$ straight to $o^T_{T-1}$; and then, for $j=T-1,\dots,0$ in turn, one step down from $o^{j+1}_j$ to $o^j_j$ and, if $j>0$, straight to $o^j_{j-1}$.

It has one upward and one downward bond between any two consecutive layers that it meets, and consecutive rows need not overlap. Its physical area is $|A(C)|=a^2\sum_jR_j$; it has $c_3(C)=1$ and $A_{13}(C)=|A(C)|$ for the numbers $c_3$, $A_{13}$ of Theorem~\ref{C-361-3}(4). A staircase loop of another coordinate plane is the image of such a loop under a lattice isometry.

\emph{The construction.} Fix the common construction constants of Proposition~\YMMref{69.1}{595} and Section~\YMMref{82}{665}, and use the first-passage data of Proposition~\YMMref{74.1}{615}. The notation is: the blocking factor $L$, an odd integer; the spacings $a_K=L^{-K}$; the sides $\mathfrak L=\{2L^m:m\in\mathbb N\}$; the constants $\gamma$ and $\gamma_W\le\gamma$; the reference couplings $0<g\le\gamma$, the sides $\mathfrak L(g)\subset\mathfrak L$ at $g$ and the first-passage data at $g$; the numbers $M$, $r\ge2$, $\mathfrak p_0\in\mathbb N$ of the constants, $R(x)=\min\{L^j:j\in\mathbb Z_{\ge0},\ L^j\ge(\log x)^r\}$, and $\mathsf p_0=2\mathfrak p_0+\frac12$; and the interval of labels $I(g)=(g^{-2}+\bar B,\,X_g-\bar B]$, $X_g=L(g^{-2}+B_\sharp)$, $\bar B=LB_\sharp/(L-1)$, such that all pairs $(K,L^K\mathsf b)$ with $\mathsf b\in I(g)$ are first-passage data at $g$.

A \emph{label} is a real $\mathsf b>0$; its \emph{line} is the set of the pairs $(K,L^K\mathsf b)$. The torus $\ell$ at the cutoff $K$ is the cube $(\ell/a_K)^3$ of $a_K\mathbb Z^3$; $\mu_{K,\ell,\beta}$ is its measure, extended periodically to $a_K\mathbb Z^3$, and $Z_{K,\ell}$ its partition function at $\beta=L^K\mathsf b$; $p_\ell=\ell^{-3}\log Z_{K,\ell}$, $p_\infty=a_K^{-3}P$. On $a\mathbb Z^3$, for a real Lipschitz function $f$ of compact support, $\mathcal O_a(f)=a^{-1}\sum_pf(x(p))a_p$, $x(p)$ the centre of $p$, and $\mathcal O_K=\mathcal O_{a_K}$. For $X=(f_1,\dots,f_n)$, $S^T_{n;K,\ell,\beta}(X)$ and $S^T_n[\nu](X)$ are the joint cumulants of $\mathcal O_K(f_1),\dots,\mathcal O_K(f_n)$ on the torus and in a state $\nu$, and a periodic Gibbs state is a Gibbs state on $a_K\mathbb Z^3$ invariant under a subgroup of finite index of $a_K\mathbb Z^3$ (equation~\YMMref{(1819)}{698} and the definition following it).

Gibbs sequences and their components are defined in Definition~\ref{C-365-5}(2), and the volume-first and growing-torus limits in Section~\ref{sec:state-limit-conventions}. We write $S_n^\sigma$ for the connected functions, as in Proposition~\YMMref{92.1}{710}, and use the reflection-space notation $\mathcal H(\sigma)$, $H_\sigma$, $\Omega_\sigma$ introduced before and in Proposition~\YMMref{94.1}{720}. On the line of a label we write, at $\beta=L^K\mathsf b$,
\[
\begin{multlined}
\delta_K(\ell)=\delta((\ell/a_K)^3),\qquad\mathsf e_K(\mathsf b,\ell_0)=\mathsf e(\ell_0/2a_K),\qquad X_K(\mathsf b,\ell)=D_L(\ell/La_K)=L^3\log Z_{K,\ell/L} \\
-\log Z_{K,\ell},
\end{multlined}
\]
and call $X_K$ the \emph{cube number}; then $\ell^3(p_\ell-p_\infty)=\delta_K(\ell)$. A \emph{number of the constants} is formed from this fixed common construction assignment; it depends on no cutoff, coupling, side or label.

\emph{The chessboard seminorm.} For $\ell''\in\mathfrak L$ with $\ell''/2a_K\ge3$, a \emph{board} of side $\ell''/2$ is a board of Section~\ref{sec:periodic-chessboard} of the side $\ell''/2$ with an offset $o\in a_K\mathbb Z^3$: its cells are the closed cubes $o+\frac{\ell''}2\mathbf k+[-\frac{\ell''}4,\frac{\ell''}4]^3$, $\mathbf k\in\mathbb Z^3$, and its walls are the planes $x_\mu=o_\mu+\frac{\ell''}2(t+\frac12)$, $t\in\mathbb Z$, which lie in $a_K(\mathbb Z+\frac12)$, as $\ell''/2a_K$ is odd: planes of lattice reflections through no site. The \emph{bonds of a cell}, its \emph{inner bonds}, are the bonds with both end points in the closed cube. A function $F$ of the bonds of one cell has a mirror copy in every cell, its image under the reflections in the walls between, complex conjugated at an odd number of reflections.

For a torus $\ell\in\mathfrak L$ with $\ell\ge\ell''$, a period of the torus holds $(2\ell/\ell'')^3$ cells of the board, and $\|F\|^{(\ell'')}_{\mathrm{ch},\ell}$, the upper index in parentheses naming the side $\ell''$ of the board, is the root of order $(2\ell/\ell'')^3$ of the mean under $\mu_{K,\ell,\beta}$ of the product of the copies of $F$ in these cells: the chessboard seminorm of $F$ at the side $\ell$, for the board of $F$ (Section~\ref{sec:periodic-chessboard}, equation~\eqref{eq:period-box-seminorm}). At $\ell=\ell''$ there are eight copies, and we write $\|F\|_{\mathrm{ch},\ell''}=\|F\|^{(\ell'')}_{\mathrm{ch},\ell''}$. So $\|\cdot\|_{\mathrm{ch},\ell}$ without an upper index is always the seminorm of eight copies on the torus $\ell$, for a board of side $\ell/2$: a function of a cell of side $\ell''/2<\ell/2$ is then read under the board of side $\ell/2$ whose cell is centred at its own cell, whose walls are walls of its own board, $\ell/\ell''$ being a power of the odd $L$, and whose eight copies of it are among its copies under its own board; (M2) below compares the two seminorms.

On a box: for an even $N=2h$ and a corner $c\in a(\mathbb Z+\frac12)^3$, the cell below $c$ is $\mathsf C_h(c)=c+[-ha,0]^3$, and for a function $\Psi$ of its bonds $\|\Psi\|_{\mathrm{ch};N}$ is the eighth root of the mean on the cube $N^3$ of the product of its eight mirror copies in the planes $x_\mu=c_\mu$; at $N=\ell/a_K$, $\ell\in\mathfrak L$, it is the seminorm $\|\Psi\|_{\mathrm{ch},\ell}$ of the torus $\ell$ for the board of side $\ell/2$ whose cell is $\mathsf C_h(c)$.

Throughout, $0<x_0\le0.069$, $\theta=12.9\,x_0\le0.8901$ and $\Lambda=\log(1/\theta)\ge0.1164$.

\begin{definition}[The box inequality BI]\label{C-361-2}
Let $\mathsf b>0$ be real, $\ell_0=2cL^{m_0}$ with $c\in\mathbb N$ and $m_0\in\mathbb Z$, and let $K_0$ be an integer with $m_0+K_0\ge0$ and $cL^{m_0+K_0}\ge3$. BI$[\mathsf b,\ell_0,K_0;x_0]$ is the statement: for every integer $K\ge K_0$, with $l_K=\ell_0/2a_K=cL^{m_0+K}$,
\[
\mathsf e(l_K)=2\log Z(l_K,2l_K,2l_K)-\log Z(2l_K,2l_K,2l_K)\le x_0\qquad\text{at }\beta=L^K\mathsf b .
\]
Its rate is $m_*=2\log(1/\theta)/\ell_0\ge0.2328/\ell_0$; then $\theta^{1/l_K}=e^{-m_*a_K}$. For every $z\in\mathbb Z$ the data $(L^z\mathsf b,L^z\ell_0,K_0-z)$ give the same statement, since they name the same pairs $(l_K,L^K\mathsf b)$. In the proofs below only the fact that $l_K$ is an integer $\ge3$ is used, except where $\ell_0=2L^{m_0}$ is assumed.
\end{definition}

\begin{remark}[The spectral reading]\label{C-361-18}
At every $\beta\ge0$ and $l\ge3$, with $\mathsf r=\mathsf r_l(2l,2l)$ over the cross-section $(2l,2l)$ of the box $(l,2l,2l)$, Proposition~\ref{C-361-7}(b) at $q=2$ gives $\log(1+\mathsf r)\le\mathsf e(l)\le2\log(1+\mathsf r)$; so $\mathsf r\le e^{x_0/2}-1$ implies $\mathsf e(l)\le x_0$, and $\mathsf e(l)\le x_0$ implies $\mathsf r\le\varphi(e^{x_0}-1)$, $\varphi(\eta)=\eta/(1+\sqrt{1-\eta^2})$; at $x_0=0.069$ the two numbers are $0.035102$ and $0.035764$. Up to these constants BI is the statement that the thermal sum $\sum_{i\ge1}(\lambda_i/\lambda_0)^{l_K}$ over the cross-section $\ell_0\times\ell_0$ is small uniformly in the cutoff.
\end{remark}

\subsection{Principal results}

\begin{theorem}[One label: one state, one limit, one vacuum, a gap, an area law]\label{C-361-3}
Assume BI$[\mathsf b,\ell_0,K_0;x_0]$. Let $K\ge\max\{K_0,1\}$ and $\beta=L^K\mathsf b$. No constant below depends on $K$.
\begin{enumerate}
\item[(1)] \emph{The lattice.}
(a) For every $\ell\ge\ell_0$ with $\ell/a_K\in\mathbb N$: $\delta_K(\ell)<0.13736\,\theta^{\ell/\ell_0}$ and $\delta_K(\ell)\le0.5892\,\theta^{2\ell/\ell_0}=0.5892\,e^{-m_*\ell}$; at $\ell=2^j\ell_0$, $\delta_K(\ell)\le0.017170\,\theta^{2\ell/\ell_0}$; if $\ell/a_K$ is even, $\mathsf e(\ell/2a_K)<0.3702\,\theta^{\ell/2\ell_0}$ and $\mathsf e(\ell/2a_K)\le0.4079\,\theta^{\ell/\ell_0}$; on the cube $\ell/a_K$ all seven $0\le1-r^z<0.4121\,\theta^{\ell/\ell_0}$; for $\ell\ge L\ell_0$, $X_K(\mathsf b,\ell)<0.13736\,L^3\theta^{\ell/L\ell_0}$; and for every box whose least physical side $\ell_-$ is $\ge\ell_0$, $\delta(\underline n)<|\underline n|(a_K/\ell_-)^3\,0.13736\,\theta^{\ell_-/\ell_0}$.
(b) The lattice $a_K\mathbb Z^3$ has exactly one periodic Gibbs state $\nu_K$ at $\beta$; it is ergodic under $a_K\mathbb Z^3$, invariant under all isometries of the lattice, reflection positive at every lattice mirror, trivial on the tail $\sigma$-field, and an extreme point of the set of all Gibbs states at $\beta$; the pressure $P$ is continuously differentiable near $\beta$ with $P'(\beta)=-3\nu_K(a_p)$.
(c) For every $\ell\in\mathfrak L$ with $\ell\ge\ell_0$, on the inner bonds of a cell of side $\ell/2$ the laws of $\nu_K$ and of the torus $\ell$ are at a distance in total variation at most $0.541\,\theta^{\ell/2\ell_0}$, and $|\nu_K(G)-\mu_{K,\ell',\beta}(G)|\le0.66\,\theta^{\ell'/2\ell_0}\|G\|_{\mathrm{ch},\ell''}$ for real $G$ of a cell of side $\ell''/2$, $\ell_0\le\ell''\le\ell'$ in $\mathfrak L$. For every $\ell\ge\ell_0$ with $N=\ell/a_K$ even, every corner $c$, every bounded real Borel function $\Psi$ of the bonds of $\mathsf C_{N/2}(c)$ and every periodic box $\underline n$ with $n_1,n_2,n_3\ge N$, of any integer sides:
\[
|\langle\Psi\rangle_{\underline n}-\langle\Psi\rangle_{N^3}|\le2.71\,\theta^{\ell/2\ell_0}\|\Psi\|_{\mathrm{ch};N},\qquad|\langle\Psi\rangle_{\underline n}-\nu_K(\Psi)|\le5.42\,\theta^{\ell/2\ell_0}\|\Psi\|_{\mathrm{ch};N},
\]
and $|\mu_{\underline n}(A)-\nu_K(A)|\le2.71\,\theta^{\ell/2\ell_0}$ for the events $A$ of the bonds of the cell. Periodic boxes with receding closed twist sets, and finite mixtures of them, converge to $\nu_K$ as in Proposition~\ref{C-361-10}(d).
(d) Over every cross-section with physical sides $l_1\le l_2$, $l_1\ge\ell_0$, $l_i/a_K\in\mathbb N$, the top eigenvalue is simple, its eigenvector carries no electric flux, and
\[
-a_K^{-1}\log(\lambda_1/\lambda_0)>m_*+\frac{1.42}{l_1};\qquad\text{at }l_1=\ell_0:\ -a_K^{-1}\log(\lambda_1/\lambda_0)>\frac2{\ell_0}\log\frac1{2.42\,x_0}\ge\frac{3.579}{\ell_0};
\]
so the lowest level carrying electric flux, around either side, lies more than $m_*+1.42/l_1$ above the top; by contraposition, if over a cross-section with $l_1=l_h$ that level lies at most $1.42/l_h$ above the top at infinitely many cutoffs, then BI fails at the label $\mathsf b$ at every side $\ell_0\le l_h$. Over the square $(\ell,\ell)$, $\ell\ge\ell_0$, with $s_\ell=2^{k-1}\ell_0\in(\frac\ell4,\frac\ell2]$, $2^k\ell_0\le\ell<2^{k+1}\ell_0$, the levels $E_i=-a_K^{-1}\log(\lambda_i/\lambda_0)$ satisfy $E_1\ge m_*+1.5898/s_\ell\ge m_*+3.179/\ell$ and $\#\{i\ge1:E_i\le E\}\le0.20395\,e^{(E-m_*)s_\ell}$. For $t\in a_K\mathbb N$:
\[
\sum_{i\ge1}(\lambda_i/\lambda_0)^{t/a_K}\le\exp\Bigl[0.37\Bigl\langle\frac{l_1}{2t}\Bigr\rangle\Bigl\langle\frac{l_2}{2t}\Bigr\rangle e^{-m_*t}\Bigr]-1\le\exp\Bigl[\frac{0.37}{\ell_0^2}\,l_1l_2\,e^{-m_*t}\Bigr]-1\quad(t\ge\tfrac{\ell_0}2),
\]
$\langle u\rangle=\max\{1,u\}$, and the sum is $\le0.2314\,e^{-m_*t}$ for $t\ge l_2/2$. No such floor holds for magnetic flux: over the square $(\ell,\ell)$ the free energy of a vortex line, $E_{\mathrm v}=a_K^{-1}\log(\lambda_0/\|\mathsf T^{(1)}\|)$, satisfies $\ell E_{\mathrm v}<1.0292\,\theta^{\ell/\ell_0}$, and $\ell E_{\mathrm v}\le e^{-\sigma_*\ell^2}$ at $\ell=2^j\ell_0$, $j\ge0$, $\sigma_*$ as in (4).
(e) If $A,B$ are bounded functions of the bonds of two slabs between which lie $n$ layers, then $|\nu_K(AB)-\nu_K(A)\nu_K(B)|\le e^{-m_*na_K}\sup|A|\sup|B|$. For $F,G$ of two cells of side $\ell''/2$ of a board, $\ell''\in\mathfrak L$, $\ell''\ge\ell_0$, $j\ge1$ cells apart in a direction, $|\nu_K(FG)-\nu_K(F)\nu_K(G)|\le1.0350\,e^{-m_*(j-1)\ell''/2}\|F\|_{\mathrm{ch},\ell''}\|G\|_{\mathrm{ch},\ell''}$. For every $n$ with $na_K\ge\frac34\ell_0+\frac12a_K$ and every bounded real Borel $F$ of the bonds of the cube $\{|x|_\infty\le\lfloor n/3\rfloor a_K\}$, $\|\nu_K(F\,|\,\text{bonds outside }\{|x_3|<na_K\})-\nu_K(F)\|_{L^2(\nu_K)}\le0.7057\operatorname{osc}(F)\,e^{-m_*(2n-1)a_K/6}$.
(f) For $\ell\in\mathfrak L$, $\ell\ge\ell_0$, the free energy $E=\log Z_{\Gamma}-\log Z_{K,\ell}$ of the torus sheared by half a period, $\Gamma=\mathbb Z\ell(e_1-e_3)\oplus\mathbb Z\ell e_2\oplus\mathbb Z\frac\ell2(e_1+e_3)$, satisfies $0\le E\le\mathsf e(\ell/2a_K)\le0.4079\,\theta^{\ell/\ell_0}$; the gap alone gives $E\ge-1.24\,\theta^{\ell/\ell_0}$.
\item[(2)] \emph{One limit.} Let $n\ge1$, and let $X=(f_1,\dots,f_n)$ consist of real Lipschitz functions, none $0$, with pairwise disjoint compact supports in $(-\ell''/4,\ell''/4)^3$ for some $\ell''\in\mathfrak L$, $\ell''\ge\ell_0$. Put $X_i=\mathcal O_K(f_i)-\mu_{K,\ell'',\beta}(\mathcal O_K(f_i))$, $G_I=\prod_{i\in I}X_i$, $M_I=\sup_{K\in\mathbb N}\|G_I\|_{\mathrm{ch},\ell''}$. Then $M_I<\infty$, and with $C(X)=0.65\sum_\pi|\pi|!\prod_{I\in\pi}(1.0174\,M_I)$, $\pi$ running over the partitions of $\{1,\dots,n\}$, for every $\ell'\in\mathfrak L$ with $\ell'\ge\ell''$:
\[
\begin{multlined}
\sup_{K\ge\max\{K_0,1\}}\bigl|S^T_n[\nu_K](X)-S^T_{n;K,\ell',L^K\mathsf b}(X)\bigr|\le C(X)\,\theta^{\ell'/2\ell_0}, \\
\sup_{K\ge\max\{K_0,1\}}\ \relax \sup_{\min_\mu n_\mu\ge\ell'/a_K}\bigl|S^T_n[\mu_{\underline n}](X)-S^T_n[\nu_K](X)\bigr|\le C^{\rm sh}(X)\,\theta^{\ell'/2\ell_0},
\end{multlined}
\]
$C^{\rm sh}(X)=3.36\sum_\pi|\pi|!\prod_{I\in\pi}(1.208\,M_I)$. Hence, for $n\ge2$, $S_n(X)=\lim_KS^T_n[\nu_K](X)$ exists along all cutoffs; it equals $\lim_{\ell'}\lim_KS^T_{n;K,\ell',L^K\mathsf b}(X)$, the limit of the torus cumulants along every pair of sequences $K_j\to\infty$, $\mathfrak L\ni\ell_j\to\infty$, and the limit of the cumulants of periodic boxes $\underline n^{(j)}$ at cutoffs $K_j\to\infty$ along every sequence of boxes of any shape with $a_{K_j}\min_\mu n^{(j)}_\mu\to\infty$. If $\mathsf b\in I(g)$, $0<g\le\gamma$, and the supports lie in one closed ball of radius $R$ at pairwise distances $\ge\mathsf d$, then for $\ell'\in\mathfrak L(g)$, $\ell'\ge\max\{\ell_0,128n(R+1)\}$, the first bound holds with $C^\flat(X)=0.65\sum_\pi|\pi|!\prod_{I\in\pi}(1.0174\,\mathsf m^\flat_I)$, $\mathsf m^\flat_I=\mathsf M_{8|I|}(\mathsf d';R)^{1/8}\prod_{i\in I}\|f_i\|_\sharp$, which depends only on $g$, the constants (with the radius fixed in Theorems~\YMMref{AJ}{673} and~\YMMref{AD}{607}), $n$, $\mathsf d'=\min\{\mathsf d,1\}$, $R$ and the norms $\|f_i\|_\sharp=\max\{\sup|f_i|,40M\operatorname{Lip}f_i\}$; here $\mathsf M_k(\mathsf d';R)$ is the sum over the partitions of $\{1,\dots,k\}$ into blocks of at least two members of the products of the bounds $C^g_{|B|}(\mathsf d';R)$ of (M5).

With the same $C^\flat(X)$ it holds with a torus $\ell^+\ge\ell'$ of $\mathfrak L$ in the place of $\nu_K$, and with $0.69$ for $0.65$ with the box $(q_1N,q_2N,q_3N)$, $N=\ell'/a_K$, $q\in\mathbb N^3$, in the place of $\nu_K$; for every periodic box $\underline n$ with all sides $\ge\ell'/a_K$, $|S^T_n[\mu_{\underline n}](X)-S^T_n[\nu_K](X)|\le3.36\sum_\pi|\pi|!\prod_{I\in\pi}(1.19\,\mathsf m^\flat_I)\,\theta^{\ell'/2\ell_0}$; and for $n\ge2$, $|S_n(X)|\le C^g_n(\mathsf d;R)\prod_i\|f_i\|_\sharp$.
\item[(2$'$)] \emph{One far field.} Let $n\ge1$, $\rho\ge1$, $\mathsf d>0$, $f_1,\dots,f_n$ real Lipschitz, none $0$, with supports in $[-\rho,\rho]^3$ at pairwise sup distances $\ge\mathsf d$, $h$ real Lipschitz, not $0$, with support in $v+[-\rho,\rho]^3$, $\Pi=\|h\|_\sharp\prod_i\|f_i\|_\sharp$, and $\ell''\in\mathfrak L$, $\ell''\ge\max\{\ell_0,8\rho+4\mathsf d+4\}$. There is $\mathsf C_n$, formed from $n,\rho,\mathsf d,\ell''$, $\mathsf b$ and the constants, such that $|S^T_{n+1}[\nu_K](h,X)|\le\mathsf C_n\,e^{-m_*(|v|_\infty-\ell'')}\,\Pi$ for $|v|_\infty\ge\ell''$, and on every torus $\ell\ge L\ell''$ of $\mathfrak L$, $|S^T_{n+1;K,\ell,\beta}(h,X)|\le3.44\,\mathsf C_n\,e^{-m_*(|v|_{\infty,\ell}-\ell'')}\,\Pi$ for $|v|_{\infty,\ell}\ge\frac32\ell''$, $|\cdot|_{\infty,\ell}$ the distance on the torus.

The same two bounds hold with $h$ replaced by any bounded real Borel function $W$ of the bonds with both ends in $v+[-\rho,\rho]^3$, with $\mathsf C'_n\sup_\square\|W-\mu_{K,\ell'',\beta}(W)\|_{\mathrm{ch},\ell''}\prod_i\|f_i\|_\sharp$ in the place of $\mathsf C_n\Pi$, $\mathsf C'_n$ formed as $\mathsf C_n$ and the supremum over the cells $\square$ of side $\ell''/2$ of all boards that hold $v+[-\rho,\rho]^3$ at a sup distance $\ge\mathsf d/2$ from their walls; so the decay of a far insertion $W_K$, for instance a local density smeared at a fixed scale, is uniform in the cutoff whenever these seminorms are bounded in $K$. In particular, let $u$ be a real bounded Lipschitz function on $\mathbb R^3$ that vanishes on $\{|x|_\infty<r\}$, $r\ge\ell''+\rho$, and let $\chi_v=\chi_0(\cdot-v)$, $v\in\rho\mathbb Z^3$, be the translates of the tent $\chi_0(x)=\prod_\mu\max\{0,1-|x_\mu|/\rho\}$: a partition of unity by Lipschitz functions with $\operatorname{supp}\chi_v\subset v+[-\rho,\rho]^3$.

Then $u\chi_v=0$ unless $|v|_\infty\ge r-\rho\ge\ell''$, $\sup_v\|u\chi_v\|_\sharp<\infty$, and for every $K\ge\max\{K_0,1\}$
\[
\sum_{v\in\rho\mathbb Z^3}\bigl|S^T_{n+1}[\nu_K](u\chi_v,X)\bigr|\le\mathsf C_n\,\sup_v\|u\chi_v\|_\sharp\prod_i\|f_i\|_\sharp\sum_{v\in\rho\mathbb Z^3,\ |v|_\infty\ge r-\rho}e^{-m_*(|v|_\infty-\ell'')},
\]
a bound free of $K$ that tends to $0$ as $r\to\infty$ at a fixed bound of $\sup_v\|u\chi_v\|_\sharp$. Thus the connected-cumulant series $\sum_vS^T_{n+1}[\nu_K](u\chi_v,X)$ converges absolutely, uniformly in $K$, and is the limit of the cumulants with finite partial sums of the far insertions; it tends to $0$ as the region of $u$ recedes. On a torus $\ell\ge L\ell''$ of $\mathfrak L$, take $\rho=\ell/q\ge1$ with $q$ an integer and $\ell''\ge8\rho+4\mathsf d+4$, use the finite set $V=\rho\mathbb Z^3/\ell\mathbb Z^3$ and the periodized tents $\chi_v^{\rm per}(x)=\sum_{k\in\mathbb Z^3}\chi_0(x-v+\ell k)$. For a bounded Lipschitz function $u$ on the torus vanishing where $|x|_{\infty,\ell}<r$, $r\ge\frac32\ell''+\rho$, the finite sum over $V$ has the same bound with $3.44\,\mathsf C_n$ and periodic distances. Its tail is bounded by the displayed infinite-grid tail, so the decay as $r\to\infty$ is uniform over these tori.
\item[(3)] \emph{One vacuum, the Euclidean group, a gap.} Let $\mathsf b\in I(g)$, $0<g\le\gamma$. Every Gibbs sequence at $g$ with the pairs $(K_j,L^{K_j}\mathsf b)$, $K_j\ge\max\{K_0,1\}$, is $(\nu_{K_j})_j$; the full sequence $(\nu_K)$ converges to a functional $\sigma$ with $S^\sigma_n=S_n$ for $n\ge2$; $\sigma$ is its own and only ergodic component, invariant under $\mathbb R^3\rtimes\mathrm O(3)$ and reflection positive at every plane; $\ker H_\sigma=\mathbb C\Omega_\sigma$, and $\mathcal H(\sigma)$ carries a Wightman theory of one Hermitian scalar field on $\mathbb R^{1+2}$ with a unique vacuum and the Schwinger functions of $\sigma$. Every continuum limit at $g$ along growing tori with the pairs $(K_j,L^{K_j}\mathsf b)$, those with the spacing removed first among them, is $(S^\sigma_n)_{n\ge2}$. Finally $\operatorname{spec}H_\sigma\subset\{0\}\cup[m_*,\infty)$: for all $F,F'$ in the span of separated tuples in the upper half space and $t\ge0$,
\[
\bigl|\sigma(\Theta\bar F\cdot\tau_{te_3}F')-\overline{\sigma(F)}\,\sigma(F')\bigr|\le e^{-m_*t}\,\|(1-P_\Omega)[F]\|\,\|(1-P_\Omega)[F']\|,
\]
$\Theta$ the reflection $x_3\mapsto-x_3$, $P_\Omega$ the projection on $\Omega_\sigma$; the same in every direction.
\item[(4)] \emph{An area law.} For every rectangle $C$ of a coordinate plane with physical sides $R,T\in a_K\mathbb N$,
\[
0\le\nu_K(W_C)\le e^{-\sigma_*RT},\qquad\sigma_*=\frac{2\log(1/\theta)+1.67}{\ell_0^2}\ge\frac{1.9028}{\ell_0^2};
\]
$|\nu_K(W_C)|\le e^{-\sigma_*|A(C)|}$ for every staircase loop of a coordinate plane, $|A(C)|$ its physical area; and for every closed lattice path $C$ with a bond of the direction $3$, $|\nu_K(W_C)|\le2^{c_3(C)-1}e^{-\sigma_*A_{13}(C)}$, where $A_{13}(C)$ is the area of the union of the squares of side $a_K$ of the plane $(13)$ about which the projection of $C$ winds an odd number of times, and $c_3(C)$ is the number of passages from a downward to an upward bond of the direction $3$ in the cyclic order along $C$; the same holds for every ordered pair of directions. The static potential $V_K(r)=\lim_{T\to\infty}-T^{-1}\log\nu_K(W_{r\times T})$, $r,T\in a_K\mathbb N$ physical, exists in $[0,\infty]$, is concave and non-decreasing, and $V_K(r)\ge V_K(r')+\sigma_*(r-r')$ for $r\ge r'$. For every $\ell$ with $\ell/a_K$ an integer $\ge3$: $0\le1-r_1(\ell/a_K)\le2\exp(-8.8\cdot10^{-4}\,\ell^2/\ell_0^2)$, with $5.707\cdot10^{-3}$ in the place of $8.8\cdot10^{-4}$ for $\ell\ge128\,\ell_0$; and with $\mathsf q=1+0.835/\log(1/\theta)$, for $\ell\ge\mathsf q\ell_0$: $\frac12(1-r_1(\ell/a_K))\le\exp[-\sigma_*\ell^2+(4\mathsf q\ell/\ell_0)(\frac12\log(4\mathsf q\ell/\ell_0)+0.86)]$. None of these bounds depends on $K$, $\mathsf b$ or $L$.
\end{enumerate}
\end{theorem}

\begin{theorem}[All labels: one theory up to dilation]\label{C-361-4}
Assume BI$[\mathsf b,\ell_0,K_0;x_0]$ at one label $\mathsf b>0$.
\begin{enumerate}
\item[(1)] There are $\beta_0>0$ and $C_+<\infty$ such that for every real $\beta\ge\beta_0$ there is an integer $l$ with $3\le l\le C_+\beta$ and $\delta(l^3)\le0.008625$, hence $\mathsf e(l)\le0.069$, at $\beta$.
\item[(2)] For every $\mathsf b'>0$ and every integer $m$ with $L^m\ge48C_+\mathsf b'$, BI$[\mathsf b',2L^m,K_0';0.069]$ holds, $K_0'$ the least integer with $L^{K_0'}\mathsf b'\ge\beta_0$ and $K_0'+m\ge1$; on every compact set of labels one side and one threshold serve. Hence Theorem~\ref{C-361-3} holds at every label, its clause (3) at every label of every $I(g)$; and the set of the labels at which BI holds at some side and some threshold is empty or all of $(0,\infty)$.
\item[(2$'$)] \emph{Near the label, with the same threshold.} Let $\ell_1\in\mathfrak L$, $\ell_1\ge75\,\ell_0$, $f(x)=\max\{0,1-8|x|_\infty/\ell_1\}$ and $\mathsf M=\sup\{\|\mathcal O_K(f)-\mu_{K,\ell_1,L^K\mathsf b'}(\mathcal O_K(f))\|_{\mathrm{ch},\ell_1}:K\in\mathbb N,\ \mathsf b'\in[\frac{\mathsf b}2,2\mathsf b]\}$. Then $\mathsf M<\infty$, and with $\eta=\min\{\frac{\mathsf b}2,\ 0.0086/(0.17516\cdot1458\,\mathsf M)\}$: for all $|\mathsf b'-\mathsf b|\le\eta$ and all $K\ge\max\{K_0,1\}$, $\delta_K(\ell_1)\le0.008625$ at $\beta=L^K\mathsf b'$, so BI$[\mathsf b',2\ell_1,\max\{K_0,1\};0.069]$; $\mathsf b'\mapsto S_n^{(\mathsf b')}(X)$ is continuous there, and $S^T_n[\nu_{K_j,\mathsf b_j}](X)\to S_n^{(\mathsf b)}(X)$ along all $K_j\to\infty$, $\mathsf b_j\to\mathsf b$.

Moreover, let $\ell_0=2L^{m_0}$, let $\ell_1^\flat=2L^{m_1}$ with $m_1$ the least $m\in\mathbb N$ with $2L^m\ge\ell_0$ and $L^3\,0.13736\,\theta^{2L^m/\ell_0}\le0.03$, let $f^\flat$, $\mathsf M^\flat$ be formed as $f$, $\mathsf M$ with $\ell_1^\flat$ for $\ell_1$, and $r=\min\{\frac{\mathsf b}2,0.03/(16L^3\mathsf M^\flat)\}$. Then for all $|\mathsf b'-\mathsf b|\le r$ and all $K\ge\max\{K_0,1\}$, at $\beta=L^K\mathsf b'$: $0\le D_L(\ell_1^\flat/a_K)<0.06$, $\mathsf e(\ell_1^\flat/a_K)<2.2\cdot10^{-4}$, and BI$[\mathsf b',L^z\ell_1^\flat,\max\{K_0,1\};6.1\cdot10^{-9}]$ for every integer $z\ge1$. Neither $\eta$ nor $r$ depends on the cutoff.
\item[(3)] For every $K\in\mathbb N$ and every real $\beta\ge\beta_0$ the lattice $a_K\mathbb Z^3$ has exactly one periodic Gibbs state $\nu$ at $\beta$, which is also the limit of the tori; $|\nu(AB)-\nu(A)\nu(B)|\le e^{-c_*n/\beta}\sup|A|\sup|B|$ for bounded functions of two slabs between which lie $n$ layers, $c_*=\log(1/0.8901)/C_+$: the correlation length in lattice units is at most $\beta/c_*$, and the transfer operator $T_\nu$ of $\nu$ (the setting, in every coordinate direction as the time) has $\ker(1-T_\nu)=\mathbb C\Omega$ and norm at most $e^{-c_*/\beta}$ on $\Omega^\perp$. Moreover $0\le1-r_1(N)\le2\exp[-2.2\cdot10^{-4}(N/C_+\beta)^2]$ at $\beta$ for every integer $N\ge3$, and in lattice units every slope of the static potential of $\nu$ is at least $0.4757\,(C_+\beta)^{-2}$. For $0<g\le\gamma$ and $\ell_0^g=2C_+L(g^{-2}+\bar B)$: $|S^T_{n;K,\ell,\beta}(X)-S^T_{n;K,\ell',\beta}(X)|\le C_g(X)\,0.8901^{\ell/2\ell_0^g}$ for $\ell\le\ell'$ in $\mathfrak L(g)$ above a side that $X$ fixes, uniformly in the first-passage data $(K,\beta)$ at $g$ with $\beta\ge\beta_0$, with $C_g(X)$ free of $K,\beta,\ell,\ell'$; the same with $\nu$ for the torus $\ell'$.
\item[(4)] For every $c>0$ there is one functional $\sigma_c$ such that for all real spacings $a_j\to0$ and couplings $\beta_j$ with $a_j\beta_j\to c$, and every choice of periodic Gibbs states on $a_j\mathbb Z^3$ at $\beta_j$, the sequence converges to $\sigma_c$. With $X_s=(f_1(\cdot/s),\dots,f_n(\cdot/s))$,
\[
S^{\sigma_{tc}}_n(X)=t^{-n}\,S^{\sigma_c}_n(X_{1/t})\qquad(t,c>0,\ n\ge2),
\]
and $\sigma_c$ is the same for every admitted blocking factor. In lattice units, $H_n(Y)=\lim_{\beta\to\infty}\kappa_n(\beta;Y)$ exists along the real $\beta$, where $\kappa_n(\beta;Y)$ is the joint cumulant of the $\mathcal O_{1/\beta}(y_i)$ in the periodic Gibbs state on $\beta^{-1}\mathbb Z^3$ at $\beta$, and $S^{\sigma_c}_n(X)=c^{-n}H_n(X_{1/c})$. For $c\in I(g)$, $\sigma_c$ is the functional $\sigma$ of Theorem~\ref{C-361-3}(3) at the label $c$; its gap scales like $1/c$.
\item[(5)] \emph{The rotational Ward identity.} Let $\mathsf b\in I(g)$, $0<g\le\gamma$, $n\ge2$, $\mathsf d>0$, $R\ge1$, $0<\kappa<\frac1{16}$. There are $C$ and $p'$, depending on $g$, $n$, $\mathsf d$, $R$, $\kappa$, the data of BI and the constants, but on no side and on no place, such that for all large $K$, all $\ell\in\mathfrak L(g)$ and $v\in a_K\mathbb Z^3$ with $|v|_\infty+R+2\le\ell/4$, all $\omega\in\mathfrak{so}(3)$, and all real $\chi,f_1,\dots,f_n\in C^{1,1}(\mathbb R^3)$ that vanish outside $v+[-R,R]^3$, the $f_i$ at pairwise sup distances $\ge\mathsf d$, $\chi$ equal to $1$ at the centres and at the bonds of the plaquettes at which some $f_i$ does not vanish, the bulk term $\mathfrak U^{(K,\ell,L^K\mathsf b)}_{\omega,\chi}(f_1,\dots,f_n)$ of the rotational Ward identity~\YMMref{(2108)}{782} on the torus $\ell$ (the joint cumulant of the $\mathcal O_K(f_i)$ with one insertion of the local density of the rotation $\omega$ weighted by $\chi$, together with the like cumulants in which the insertion is weighted by one $f_i$) satisfies
\[
\bigl|\mathfrak U^{(K,\ell,L^K\mathsf b)}_{\omega,\chi}(f_1,\dots,f_n)\bigr|\le(1+2|v|_1)\,C\,|\omega|\,K^{p'}L^{-\kappa K}\,\|\chi\|_{C^{1,1}}\prod_i\|f_i\|_{C^{1,1}},
\]
and the same bound holds for the bulk term in the state $\nu_K$. The bound rests, besides on BI and on Theorem~\YMMref{D}{8}, on the bound (M7) of the member $\mathfrak U^\sharp_{\omega,\chi}$ of the bulk term of third order in the curvature.
\end{enumerate}
\end{theorem}

\begin{theorem}[One cutoff decides; a form without cutoff]\label{C-361-5}
Let $\mathsf b>0$ and $\ell=2L^m\in2L^{\mathbb Z}$. Let $\varsigma=\sum_{j\ge0}(1+j)^{\mathsf p_0}L^{-j/10}$, and let $z=z(\mathsf b,\ell)$ be the least integer $z\ge0$ with $L^z\mathsf b\ge\gamma^{-2}$, $m+z\ge2$ and $LM\,R(L^z\mathsf b)\le L^{m+z-1}$.
\begin{enumerate}
\item[(1)] There are an explicit threshold $K_\star(\mathsf b,\ell)$, formed from the constants, $\mathsf b$ and $\ell$, with $K_\star(\mathsf b,\ell)-z$ at or above the thresholds $K_{\rm bf}$ of (M4)(ii) on the two tori $L^{z-1}\ell$ and $L^z\ell$ at the label $L^z\mathsf b$, to which the dilation of Proposition~\ref{C-361-13}(c) by $z=z(\mathsf b,\ell)$ brings the data, and the margins
\[
\eta_K(\mathsf b,\ell)=3\varsigma\,\mathsf C_{\rm bf}(L^z\mathsf b)\,(L^z\ell)^3(K-z)^{\mathsf p_0}L^{-(K-z)/10}\le C(\mathsf b,\ell)\,K^{\mathsf p_0}L^{-K/10},
\]
$\mathsf C_{\rm bf}(\mathsf b')$ a number of the constants and of $\mathsf b'$, bounded on compact sets of labels, so that $C(\cdot,\ell)$ is bounded on compact sets of labels (where $z=0$, $\eta_K=3\varsigma\,\mathsf C_{\rm bf}(\mathsf b)\,\ell^3K^{\mathsf p_0}L^{-K/10}$), such that $|X_{K'}(\mathsf b,\ell)-X_K(\mathsf b,\ell)|\le\eta_K(\mathsf b,\ell)$ for $K'\ge K\ge K_\star(\mathsf b,\ell)$. So $X_\infty(\mathsf b,\ell)=\lim_KX_K(\mathsf b,\ell)\ge0$ exists, uniformly on compact sets of labels. For every $y>0$: if $X_{K_1}\le y-\eta_{K_1}$ at one $K_1\ge K_\star$, then $X_K\le y$ for all $K\ge K_1$; $X_\infty<y$ if and only if $X_{K_1}<y-\eta_{K_1}$ at one $K_1\ge K_\star$; $X_\infty>y$ if and only if $X_{K_1}>y+\eta_{K_1}$ at one $K_1\ge K_\star$, and then $X_K>y$ for all $K\ge K_1$.
\item[(2)] \emph{The cube criterion.} Let $2^k\le L/2<2^{k+1}$, $\Phi_L(x)=(L/2^{k+1})^3[\epsilon+\frac{20}7\epsilon^2]$ with $\epsilon=(12.9x)^{2^k}/12.9$, $x_1=\max\{x\in(0,0.069]:\Phi_L(x)\le x_0\}$, $y_1=\frac{L^3}8x_1$ and $y_0=\frac{L^3}8x_0$. If $X_{K_1}(\mathsf b,\ell)\le y_1-\eta_{K_1}(\mathsf b,\ell)$ at one $K_1\ge K_\star(\mathsf b,\ell)$, then BI$[\mathsf b,\ell,K_1;x_0]$ holds; if $X_{K_1}(\mathsf b,\ell)\le y_0-\eta_{K_1}(\mathsf b,\ell)$ there, then BI$[\mathsf b,2\ell/L,K_1;x_0]$ holds (with $c=2$), and with it Theorem~\ref{C-361-3} at the side $2\ell/L$, of rate $L\log(1/\theta)/\ell$. Conversely BI$[\mathsf b,2\ell/L,K_0;x_0]$ gives $X_K(\mathsf b,L\ell)\le L^3\bar\delta_{x_0}(L)$ for all $K\ge K_0$, with $\bar\delta_x(\nu)=(\nu/2^j)^3\mathsf c_0(12.9x)^{2^j}$ for $2^j\le\nu<2^{j+1}$, $\mathsf c_0=\frac{20}{7\cdot12.9^2}\approx0.0171693$; so $X_{K_1}(\mathsf b,L\ell)>L^3\bar\delta_{x_0}(L)+\eta_{K_1}(\mathsf b,L\ell)$ at one $K_1\ge K_\star(\mathsf b,L\ell)$ refutes it. For every $K$ with $K+m\ge2$,
\[
\min\Bigl\{0.069,\Bigl[\tfrac7{20}L^{-3}X_K(\mathsf b,L\ell)\Bigr]^{1/2}\Bigr\}\le\mathsf e_K(\mathsf b,\ell)\le\min\Bigl\{\tfrac2LX_K(\mathsf b,\ell),\ \Phi_L\bigl(\tfrac8{L^3}X_K(\mathsf b,\ell)\bigr)\Bigr\},
\]
the last member where its argument is $\le0.069$; so BI$[\mathsf b,\ell,K_0;x_0]$ gives $X_K(\mathsf b,L\ell)\le\frac{20}7L^3x_0^2$ for $K\ge K_0$, and if $X_{K_1}(\mathsf b,L\ell)>\frac{20}7L^3x_0^2+\eta_{K_1}(\mathsf b,L\ell)$ at one $K_1\ge K_\star(\mathsf b,L\ell)$, then $\mathsf e_K(\mathsf b,\ell)>x_0$ for every $K\ge K_1$. The same bracket holds for $\liminf_K\mathsf e_K$ and $\limsup_K\mathsf e_K$ with $X_\infty$.
\item[(3)] \emph{The inequality of BI itself at one cutoff.} Let $\ell_0=2L^{m_0}$, $i\ge1$ and $\ell=L^i\ell_0$. If $\mathsf e_{K_1}(\mathsf b,\ell_0)\le x_0$ at one $K_1\ge K_\star(\mathsf b,\ell)$ with $L^3\bar\delta_{x_0}(2L^{i-1})+\eta_{K_1}(\mathsf b,\ell)\le y_1$, then $X_K(\mathsf b,\ell)\le y_1$ for all $K\ge K_1$ and BI$[\mathsf b,L^i\ell_0,K_1;x_0]$ holds, of rate $L^{-i}m_*$; with $y_0$ in the place of $y_1$ it gives BI$[\mathsf b,2L^{i-1}\ell_0,K_1;x_0]$. The margin condition reads $\bar\delta_{x_0}(2L^{i-1})\le\frac18x_1-L^{-3}\eta_{K_1}$; at $x_0=0.069$ and $i=2$ it holds, once $\eta_{K_1}$ is small enough, at every odd $L\ge33$, so at every admitted blocking factor, and at $i=1$ it never holds. So, wherever $L^3\bar\delta_{x_0}(2L)<y_1$, with $K_{\star\star}(\mathsf b,\ell_0)$ the least $K\ge K_\star(\mathsf b,L^2\ell_0)$ such that $\eta_{K'}(\mathsf b,L^2\ell_0)\le y_1-L^3\bar\delta_{x_0}(2L)$ for all $K'\ge K$, an explicit expression, the inequality $\mathsf e_{K_1}(\mathsf b,\ell_0)\le x_0$ at one $K_1\ge K_{\star\star}(\mathsf b,\ell_0)$ gives BI$[\mathsf b,L^2\ell_0,K_1;x_0]$.
\item[(4)] The following are equivalent: (i) BI$[\mathsf b,\ell_0,K_0;x_0]$ holds for some $\ell_0\in2\mathbb NL^{\mathbb Z}$ and some $K_0$; (ii) $X_\infty(\mathsf b,\ell)<y_0$ for some $\ell\in2L^{\mathbb Z}$; (iii) $X_{K_1}(\mathsf b,\ell)<y_0-\eta_{K_1}(\mathsf b,\ell)$ for some $\ell\in2L^{\mathbb Z}$ and one $K_1\ge K_\star(\mathsf b,\ell)$; (iv) $X_\infty(\mathsf b,\ell)\to0$ as $2L^{\mathbb Z}\ni\ell\to\infty$. More precisely, (iii) at $(\ell,K_1)$ gives (i) at $\ell_0=2\ell/L$, $K_0=K_1$; and (i) at $(\ell_0,K_0)$ gives $X_K(\mathsf b,\ell)\le L^3\bar\delta_{x_0}(2\ell/L\ell_0)<0.13736\,L^3\theta^{\ell/L\ell_0}$ for all $K\ge K_0$ and all $\ell\in2L^{\mathbb Z}$ with $\ell\ge L\ell_0$.
\item[(5)] \emph{The half box.} Assume \textbf{(Hess)} from the setting of Theorem~\ref{C-366-4}. Let $\ell_0=2L^{m_0}$ and $J$ a compact set of labels. There are $C$ and $K_*$ with $|\mathsf e_{K+1}(\mathsf b',\ell_0)-\mathsf e_K(\mathsf b',\ell_0)|\le C\,L^{-K/10}$ for $\mathsf b'\in J$ and $K\ge K_*$. Hence $\mathsf e_\infty=\lim_K\mathsf e_K(\mathsf b,\ell_0)$ exists, and with $Q_K=\sum_{k\ge K}CL^{-k/10}$, $|\mathsf e_K-\mathsf e_\infty|\le Q_K$: if $\mathsf e_{K_1}(\mathsf b,\ell_0)\le x_0-Q_{K_1}$ at one $K_1\ge K_*$, then BI$[\mathsf b,\ell_0,K_1;x_0]$ holds, at the same side; BI$[\mathsf b,\ell_0,K_0;x_0']$ with $x_0'<x_0$ gives $\mathsf e_{K_1}\le x_0-Q_{K_1}$ at every large $K_1$; $\mathsf e_\infty<x_0$ if and only if $\mathsf e_{K_1}(\mathsf b,\ell_0)<x_0-Q_{K_1}$ at one $K_1\ge K_*$; and if $\mathsf e_{K_1}(\mathsf b,\ell_0)>x_0+Q_{K_1}$ at one $K_1\ge K_*$, then BI$[\mathsf b,\ell_0,K_0;x_0]$ fails for every $K_0$.
\item[(6)] \emph{A form without cutoff.} For every $\ell\in2L^{\mathbb Z}$ and every real $\mathsf b>0$ the limit $D_\infty(t)=\lim_KD_L(\ell/a_K)|_{\beta=L^K\mathsf b}$ exists, depends on $t=\ell/\mathsf b$ alone, is continuous on $(0,\infty)$, and $\sup_{t\in[t_-,t_+]}|D_L(\ell/a_K)-D_\infty(t)|\le C\,K^{\mathsf p_0}L^{-K/10}$ on compact sets for large $K$. The function $\delta_\infty(t)=\sum_{r\ge0}L^{-3(r+1)}D_\infty(L^rt)$ is continuous and positive, $\delta_K(\ell)\to\delta_\infty(\ell/\mathsf b)$ uniformly on compact sets of labels, $D_\infty(t)=L^3\delta_\infty(t)-\delta_\infty(Lt)$ and $(L-1)^3\delta_\infty\le D_\infty\le L^3\delta_\infty$. $\delta_\infty$ is the same for every admitted blocking factor; $t^{-3}\delta_\infty(t)$ does not increase, and $\delta_\infty(\rho t)\le(3\rho^2-3\rho+1)\delta_\infty(t)\le3\rho^2\delta_\infty(t)$ for real $\rho\ge1$. Either $\delta_\infty\ge0.008625$ on $(0,\infty)$, or $\delta_\infty(t_0)<0.008625$ at some $t_0$, and then $\delta_\infty(t)\le0.1374\,(103.2\,\delta_\infty(t_0))^{t/2t_0}$ for all real $t\ge2t_0$. The second case holds if and only if BI holds at some side of some label, and it holds as soon as $D_{q,\infty}(t_0)<(q-1)^3x_0/8$ at one $t_0$ for one admitted $q$, $D_{q,\infty}(t)=q^3\delta_\infty(t)-\delta_\infty(qt)$.
\item[(7)] \emph{The size of the thresholds.} Read as inequalities, the conditions of the list of constants admit exactly the odd blocking factors $L>48\cdot4096^2e^{24}=2.1331\ldots\cdot10^{19}$, the least being $21\,331\,858\,733\,892\,671\,083$. At every admissible choice of the constants, at a cutoff $K_1\ge K_\star(\mathsf b,\ell)$ the small cube of the criterion of (2) has more than $10^{1.42\cdot10^{11}}$ sites a side, and at a cutoff at which the criterion of (3) with $i=2$ applies each of the two boxes of BI has more than $10^{1.89\cdot10^{12}}$ sites. Independently of all constants but $L$, $\eta_K(\mathsf b,\ell)\ge48\,\varsigma_\flat(L)\,L^6(K-z)^{5/2}L^{-(K-z)/10}$ with $\varsigma_\flat(L)=\sum_{j\ge0}(1+j)^{5/2}L^{-j/10}$; as $X_K\ge0$, the margin alone forces more than $10^{152}$ sites a side on the small cube of any inequality $X_{K_1}\le y-\eta_{K_1}$ with $y\le y_0$.
\end{enumerate}
\end{theorem}
Clause~(5)(c) of the next theorem and Proposition~\ref{C-361-16}(b) assume the following additional estimate, denoted by \textbf{(Heaviest class)}. For every integer $j\ge1$, put $N=2L^j$. For every $\beta>0$ and every $\omega>0$ with $\Xi=\beta\omega^2\ge1$ and $\Xi N^5\omega\le3/640$, every $w$ with $N^2\omega\le w\le1/10$, every pair $Q\in G^2$ and every site $y$ of the periodic cube, its Wilson probability law satisfies \[ \mu_{N^3,\beta}\bigl(\mathsf D((P_1,P_2)(y),Q)\le w,\ \max_p|\partial U(p)-\mathbf1|<\omega\bigr)\le10^{11}\Xi^5N^7w^6\omega^{-4}. \] Here $P_1,P_2$ are the straight winding holonomies through $y$, and $\mathsf D(P,Q)=\min_{g\in G}\max_{i=1,2}|P_i-gQ_ig^{-1}|$ is the distance of pairs modulo simultaneous conjugation. This probability estimate is an additional input; the pointwise density maximum proved below does not alone give it.
\begin{theorem}[Small tori without the box criterion]\label{C-361-6}
\begin{enumerate}
\item[(1)] \emph{Comparison.} For $\beta\ge0$, let $N,l\ge3$ and $q\ge2$ be integers. Then $0\le\mathsf e(N)\le D_2(N)\le8\delta(N^3)$ and $D_q(N)\ge\frac{q^3}8D_2(N)$; $(q-1)^3\delta(N^3)\le D_q(N)\le q^3\delta(N^3)$, and $\delta(M^3)\le(3\rho^2-3\rho+1)\delta(N^3)\le3\rho^2\delta(N^3)$ for $M=\rho N\ge N$; $\delta(\underline n)\le(n_1n_2+n_1n_3+n_2n_3)N^{-2}\delta(N^3)$ for every box with all sides $\ge N$, so $\mathsf e(l)\le16(l/N)^2\delta(N^3)$ for $l\ge N$; $\mathsf e(l)\le\frac2LD_L(2l/L)$ if $l=Ll'$, $l'\ge2$; if $D_L(N)\le\frac{L^3}8x$ with $x\le0.069$, then $D_2(N)\le x$, $\mathsf e(N)\le x$ and, for even $N$, $\mathsf e(LN/2)\le\Phi_L(x)$; $\mathsf e(l)\le x\le0.069$ gives $D_L(2l)\le\frac{20}7L^3x^2$; $\delta(N^3)\le x_0/8$ gives $\mathsf e(N)\le x_0$; $\mathsf e(l)\le x_0$ gives $\delta(N'^3)\le0.008625$ for $N'\ge48l$; and $\mathsf e(l)\le24\beta l^3$, so $\mathsf e(l)\le x_0$ whenever $0<\beta\le x_0/24l^3$.
\item[(2)] \emph{Dichotomy.} At each $\beta\ge0$ the following are equivalent: $\mathsf e(l)\le0.069$ for one $l\ge3$; $\mathsf e(l)\to0$; $\delta(N^3)\to0$; $\delta(N^3)\le0.008625$ for one $N\ge3$. Then $\delta((2N)^3)\le183\,\delta(N^3)^2$ and $\delta(N'^3)\le0.13736\,(103.2\,\delta(N^3))^{N'/2N}$ for $N'\ge2N$. The numbers $0.07489$ and $0.00936$ may stand for $0.069$ and $0.008625$ in this equivalence. In the other case $\mathsf e(l)>0.069$ and $\delta(N^3)>0.008625$ for all $l,N\ge3$, and indeed $\mathsf e(2^kl)>0.07489$ for all $l$ and $k$; this case holds whenever there are two periodic Gibbs states at $\beta$, and then $\liminf_m\delta((2L^{m+K})^3)\ge\log2$. Every phase is reached by the tori: if $\liminf_m\delta((2L^{m+K})^3)\le D<\infty$ at $\beta$, there are sides along which the tori converge to a periodic Gibbs state $\mu^*$ with $\mu^*\ge e^{-D}\nu$ for every ergodic periodic Gibbs state $\nu$, and every periodic Gibbs state is at most $e^D\mu^*$.
\item[(3)] \emph{A floor.} Let $0<g\le\gamma_W$, let $\mathsf d_0(g,\vartheta)$ and $s(\mathsf d)$ be the threshold and the depth of the law of the two-point function at short distance of (M5), $\vartheta_1=3\sqrt3/8$, let $0<\mathsf d_\circ\le\mathsf d_0(g,\vartheta_1)$, let $(K,\beta)$ be a first-passage datum at $g$ with $K>s(\mathsf d_\circ)$, and $N_4=\lfloor4\mathsf d_\circ/a_K\rfloor$. Then for every integer $l\ge3$
\[
\mathsf e(l)\ge\min\bigl\{0.069,\ \tfrac1{12.9}\,31250^{-l/N_4}\bigr\}\ge\min\bigl\{0.069,\ 0.0775\,e^{-2.6\,la_K/\mathsf d_\circ}\bigr\};
\]
in particular $\mathsf e(\ell/2a_K)\ge\min\{0.069,0.0775\,e^{-1.3\ell/\mathsf d_0(g,\vartheta_1)}\}$ whenever $\ell/a_K\ge6$ is even, the same floor lies under $D_2$, $8\delta$ and $\frac8{L^3}D_L$ at the side $\ell/2$, and $\mathsf e(l)\le0.069$ fails for every $l$ with $la_K<0.0449\,\mathsf d_\circ$.
\item[(4)] \emph{What BI forces.} If BI$[\mathsf b,\ell_0,K_0;x_0]$ holds at some label, then $\ell_0\ge0.0895\,(\mathsf d_0^\star/\kappa_0)\,\mathsf b$, where $\kappa_0=L(\gamma_W^{-2}+\bar B)$, $\mathsf d_0^\star=\min\{\mathsf d_0(g_1,\vartheta_1),\mathsf d_0(g_2,\vartheta_1)\}$, $g_1=\gamma_W$ and $g_2=(\gamma_W^{-2}+2\bar B/L)^{-1/2}$. On one lattice, every box and $t\ge3$: $2F(t,\Sigma)-F(2t,\Sigma)\ge2\log\sum_fw_f^{1/2}$, the $w_f$ the weights of the box $(2t,\Sigma)$; on the cube $2l$, $\mathsf e(l)\ge2\log(\sqrt{1-x}+\sqrt{2x})$ for $x=1-r_1\le\frac12$, $\mathsf e(l)\ge\log\frac4{1+3r_1}$, and $\mathsf e(l)<\log2$ forces $r_1>\frac89$; so $\mathsf e(l)\le0.069$ forces $r_1\ge0.99937$. At every $K\ge K_0$, on the cube $2l_K$ at $\beta=L^K\mathsf b$, with $x_0=0.069$: $1-r_1\le0.000628$, $1-r_2\le0.001255$, $1-r_3\le0.001882$, $\delta_K(\ell_0)\le0.013603$, and $\sum_{i\ge1}(\lambda_i/\lambda_0)^{l_K}\le0.035764$ over the cross-section $(2l_K,2l_K)$, so $\log(\lambda_0/\lambda_i)\ge6.66\,a_K/\ell_0$ for $i\ge1$. BI$[\cdot;0.069]$ is false as soon as $r_1<0.99937$ on the cube $\ell_0$ at one cutoff $K\ge K_0$.
\item[(5)] \emph{Small tori.} Let $\vartheta(\tau,u)=\sum_{k\in\mathbb Z}e^{-\tau(2\pi k+u)^2}$,
\[
C_\Delta=\int_0^\infty\Bigl[1-e^{-\tau}-\vartheta(\tau,0)^3+(4\pi\tau)^{-3/2}\Bigr]\frac{d\tau}\tau,\qquad\Upsilon(u)=\exp\int_0^\infty\Bigl[e^{-\tau}-1+\vartheta(\tau,0)^3-\prod_{\mu=1}^3\vartheta(\tau,u_\mu)\Bigr]\frac{d\tau}\tau,
\]
$\mathscr J=\int_{[-\pi,\pi]^3}du/\Upsilon(u)$ and $c'=\log(\mathscr J/8\pi^2)-\frac32C_\Delta$; numerically $c'\approx2.0919$. There are numbers $t_1,t_2,t_4>0$ and $C,C_\flat$ of the constants, with $\mathsf a=\frac92r+\mathfrak p_0$ and $\mathsf q=1116\,r+720\,\mathfrak p_0+366$, such that for every side $\ell=2L^m$, $m\in\mathbb Z$, every real $\mathsf b>0$ and every cutoff $K\ge1-m$, with $t=\ell/\mathsf b$:
(a) if $t\le t_1$, then $|\delta_K(\ell)-\log(\mathsf b/\ell)-c'|\le C\,t^{1/1500}+C\,(a_K/\ell)^2\log^{5/2}(\ell/a_K)$, and $|\delta_\infty(t)-\log(1/t)-c'|\le C\,t^{1/1500}$;
(b) if $t\le t_2$, then $r_1(\ell/a_K)\le\eta_\flat(t)=C_\flat(\log\frac1t)^{\mathsf a}t^{1/500}$ at $\beta=L^K\mathsf b$, the weights of the four classes of electric flux on the cube satisfy $|w_f-\frac14|\le\frac34\eta_\flat(t)$, and $\delta_K(\ell)$ and $\mathsf e(\ell/2a_K)$ are both $\ge\mathsf G_{\ell/a_K}\ge\log\frac4{1+3\eta_\flat(t)}$; so $\liminf_{t\to0}\inf_K\delta_K(\ell)\ge\log4$, the same for $\mathsf e(\ell/2a_K)$, and if $\ell_0=2L^{m_0}$, $\ell_0/\mathsf b\le t_2$ and $\eta_\flat(\ell_0/\mathsf b)<0.99937$, the inequality of BI$[\mathsf b,\ell_0,K_0;x_0]$ fails at every cutoff $K\ge1-m_0$, for every $x_0\le0.069$;
(c) assuming \textbf{(Heaviest class)}, if $t\le t_4$, then over the cross-section $(\ell/a_K,\ell/a_K)$ the top level carries the weight $\lambda_0^{N}/Z(N^3)\le C\,t\,(\log\frac1t)^{\mathsf q}$, $N=\ell/a_K$; so $\delta_K(\ell)$ and $\mathsf e(\ell/2a_K)$ are $\ge\mathsf G_N\ge\log\frac{\mathsf b}\ell-\mathsf q\log\log\frac{\mathsf b}\ell-\log C$: the coefficient of $\log(\mathsf b/\ell)$ in (a) is reached from below at every cutoff;
(d) there are numbers $t_\circ>0$ and $\beta_\flat$ of the constants such that for every real $\beta\ge\beta_\flat$: $\mathsf e(l)>0.069$ for every integer $3\le l\le t_\circ\beta$; so BI at any label $\mathsf b$ forces $\ell_0>2t_\circ\mathsf b$. On one lattice, exactly: if $1-r_1(N)\ge0.0039$ on one cube, then $\mathsf e(l)>0.069$ for every integer $3\le l\le N/16$; and $\mathsf e(l)\le0.069$ gives $\mathsf e(l')<0.0234$ for every integer $l'\ge16l$;
(e) the pressure of one lattice: $|P(\beta)+3\log\beta-z_\infty|\le C\beta^{-1/1500}$ for all large real $\beta$, with $z_\infty=\log\frac2\pi-\frac3{16\pi^3}\int_{[-\pi,\pi]^3}\log\sum_\mu(2-2\cos p_\mu)\,dp\approx-2.96167$.
\item[(6)] \emph{An area from above.} If $t=\ell/\mathsf b>t_1$ and $j$ is the least integer with $L^{-j}t\le t_1$, then $\delta_K(\ell)\le3L^2(t/t_1)^2\,[\log(L/t_1)+c'+C]$ at every cutoff $K\ge1-m+j$. So $\sup_K\delta_K(\ell)<\infty$ on every torus at every label, and $\mathsf e(\ell/a_K)$, $D_L(\ell/a_K)$ are $O((\ell/\mathsf b)^2)$ uniformly in the cutoff.
\end{enumerate}
\end{theorem}

\subsection{Construction estimates}

Beyond the exact statements on one lattice proved below, the argument uses the following statements, (M1)--(M6) and (M8), one statement of the section on plane fields and the stress tensor (Section~\ref{C-360-1}), Lemma~\ref{C-360-19}, entered as (M7), and the comparison of two cutoffs with a factor, Proposition~\ref{C-361-23}, which is derived at the end of this subsection from the statements named in (M4)(ii).

(M1) The layers and the crossing operator of one direction (Section~\YMMref{85}{677}, Proposition~\YMMref{85.1}{678}): the objects of the setting; a bounded Borel function of the bonds with both ends in a slab of $\vartheta\ge2$ consecutive layers is $\Phi_{\tau_0}$ for a bounded Borel function $\Phi$ on the configurations $\mathcal Y_\vartheta$ of a slab, which depends neither on the length of the box in the time nor on $\tau_0$; its slab operator $\mathfrak Y[\Phi]=\mathsf K^{1/2}\mathfrak S[\Phi]\mathsf K^{1/2}$ is Hilbert--Schmidt, $\mathfrak Y[1_\vartheta]=\mathsf T'^{\,\vartheta}$, $\mathfrak Y[\Phi^{+j}]=\mathsf T'^{\,j}\mathfrak Y[\Phi]$ for $\Phi$ regarded as a function $\Phi^{+j}$ on a slab thicker by $j$ layers, read on the first $\vartheta$ of its $\vartheta+j$ layers, so that the $j$ added layers lie above, towards the later slabs of the trace formula below, in which the first slab stands rightmost, and $\mathfrak Y[\Phi^\Theta]=\mathfrak Y[\Phi]^*$ for the conjugated mirror image $\Phi^\Theta$, which, placed on a slab of $\vartheta$ layers, is the complex conjugate of $\Phi$ read on the configuration mirrored in the plane midway between that slab and the slab of $\Phi$; and the trace formula (Proposition~\YMMref{85.2}{679}): for $r\ge2$ functions $\Phi^1,\dots,\Phi^r$ on consecutive slabs filling the box $(t,\Sigma)$,
\[
Z(t,\Sigma)\,\mu_{(t,\Sigma)}\Bigl(\prod_s(\Phi^s)_{\tau_{s-1}}\Bigr)=\operatorname{Tr}\bigl(\mathfrak Y[\Phi^r]\cdots\mathfrak Y[\Phi^1]\bigr).
\]
(M2) The chessboard estimate on the tori $\ell\in\mathfrak L$, for one board of side $\ell''/2$, $\ell''\le\ell$ in $\mathfrak L$ (Corollary~\ref{IV-18-18}(b), (c), (e)): $\|\cdot\|^{(\ell'')}_{\mathrm{ch},\ell}$ is a seminorm on the functions of one cell, with $\|1\|^{(\ell'')}_{\mathrm{ch},\ell}=1$ and $|\mu_{K,\ell,\beta}(F)|\le\|F\|^{(\ell'')}_{\mathrm{ch},\ell}\le\sup|F|$; in its array form $|\mu_{K,\ell,\beta}(\prod_{\mathbf k}F_{\mathbf k})|\le\prod_{\mathbf k}\|F_{\mathbf k}\|^{(\ell'')}_{\mathrm{ch},\ell}$ for functions $F_{\mathbf k}$ of distinct cells $\mathbf k$ of one period, given in all of them or in some of them only; in its fine-to-coarse form, for one board and two tori, $\|F\|^{(\ell'')}_{\mathrm{ch},\ell}\le\|F\|^{(\ell'')}_{\mathrm{ch},\ell'}\exp[(\ell''/2)^3(p_{\ell'}-p_\ell)]$ and $p_\ell\le p_{\ell'}$ for $\ell''\le\ell'\le\ell$ in $\mathfrak L$; and the domination $|\nu(\prod_{\mathbf k}F_{\mathbf k})|\le\liminf_{\ell\to\infty}\prod_{\mathbf k}\|F_{\mathbf k}\|^{(\ell'')}_{\mathrm{ch},\ell}$ of a periodic Gibbs state $\nu$, for one board and growing tori (Proposition~\YMMref{91.9}{705}). Together: for a function $F$ of a cell of side $\ell''/2$ and its board of side $\ell/2$ of the setting, whose eight copies of $F$ are copies of $F$ under the board of $F$, the array form with the constant $1$ in the other cells and the fine-to-coarse form give
\[
\|F\|_{\mathrm{ch},\ell}\le\|F\|^{(\ell'')}_{\mathrm{ch},\ell}\le e^{(\ell''/2)^3(p_{\ell''}-p_\ell)}\,\|F\|_{\mathrm{ch},\ell''}:
\]
the seminorm of eight copies on the torus $\ell$ is at most that of all the copies under the board of $F$, which is at most a factor times the seminorm on the torus $\ell''$.
(M3) Limits of tori and of periodic boxes are periodic Gibbs states (Proposition~\YMMref{91.2}{699}); the DLR equations and the conditional laws of a Gibbs state; all periodic Gibbs states are invariant under the isometries of the lattice and reflection positive (Proposition~\textup{\YMMref{103.4}{758}}); a periodic Gibbs state is an integral of ergodic ones, and it is extreme if and only if it is ergodic (Proposition~\textup{\YMMref{106.1}{764}}).
(M4) Theorem~\YMMref{D}{8}, through the following consequences: (i) on one fixed torus the cumulants $S^T_{n;K,\ell,L^K\mathsf b}$ of fields with disjoint supports are bounded at all cutoffs, with bounds that depend on the side, a compact set of labels and the distance of the supports (Theorem~\textup{\YMMref{AS}{813}}, uniform bound~\YMMref{(2214)}{814}), and converge as $K\to\infty$, uniformly on compact sets of labels, to continuous functions of the label (Theorem~\textup{\YMMref{AS}{813}}, limit and local uniformity), with the Cauchy rate proved in Theorem~\YMMref{AS}{813}: on a compact interval of labels, $|S^T_{n;K',\ell,L^{K'}\mathsf b}(X)-S^T_{n;K,\ell,L^K\mathsf b}(X)|\le C\,K^{2\mathfrak p_0(n+1)+\frac12}L^{-K/10}\prod_i\|f_i\|_\sharp$ for all $K'\ge K$ above a threshold, where $C$ and the threshold depend on $n$, the side, the interval and a floor of the spread of the supports, and on the functions only through their norms, at every admitted choice of the constants; (ii) for a label $\mathsf b\ge\gamma^{-2}$, a side $\ell=2L^m$ with $MR(\mathsf b)\le L^m$ and a cutoff $K\ge2$ at which the conditions (k0)--(k4) and (k6) of Section~\ref{sec:cutoff-calibration} hold, with its numbers at $n=1$ ($u=\lceil K/1260\rceil$, $s=\lceil K/9\rceil$, $t_1=K-u$),
\[
\bigl|\log Z_{K+1,\ell}-\log Z_{K,\ell}-\mathfrak K_K(\ell,\mathsf b)\bigr|\le4\mathsf R_{t_1}+8\epsilon_{\rm fg}+8\epsilon_{\rm cl}+12\epsilon_{\rm sh}\le\mathfrak o_K(\ell,\mathsf b),
\]
where $\mathfrak K_K$, a sum of local vacuum constants, is the real part of the difference at the source $0$ of the contact terms of the two cutoffs, and $\mathfrak o_K=8\mathsf R_{t_1}+12\epsilon_{\rm fg}+8\epsilon_{\rm cl}+12\epsilon_{\rm sh}$ is the number~\textup{(\ref{IV-62-64})}. This is the comparison of the two cutoffs at the source $0$: the discarding of Proposition~\textup{\YMMref{115.2}{806}}, lost-mass bound gives $(1-\mathsf R_{t_1})Z\le T(0)\le Z$ at each cutoff for the chain $T$ of the top level; the form $T=e^cI$ and the contact terms of the proof of Proposition~\YMMref{116.1}{808} and Section~\YMMref{116}{807}, with the proof of Proposition~\YMMref{116.3}{809}, give $|I(0)/\hat I(0)-1|\le2\epsilon_{\rm fg}$ against the classical start $\hat T=e^{\hat c}\hat I$; Section~\ref{sec:block-field-laws} and Proposition~\YMMref{116.1}{808} give $|\hat T_{K+1}(0)/\hat T_K(0)-1|\le4\epsilon_{\rm cl}+6\epsilon_{\rm sh}$ for the classical starts of the two cutoffs; and $|\log(1+\zeta)|\le2|\zeta|$ for $|\zeta|\le\frac12$.

The rate $\mathfrak o_K\le\mathsf C_{\rm bf}\,\ell^3K^{\mathsf p_0}L^{-K/10}$ holds at every $K\ge K_{\rm bf}$, for the labels of an interval $J$ of labels admitted on the side $\ell$ (Definition~\textup{\ref{IV-62-62}}(5), Corollary~\textup{\ref{IV-62-66}}(a), \textup{(\ref{IV-62-68})}): $K_{\rm bf}=K_{\rm bf}(\ell,J)$ is the least cutoff from which the conditions hold and the two fast members, the large fields $8\bar{\mathsf R}$ and the shell term $12\bar\epsilon_{\rm sh}$, majorised on $J$, lie below half of $\mathsf C_{\rm bf}\ell^3K^{\mathsf p_0}L^{-K/10}$; it depends on the side, through the factor $\exp[(\mathsf C^\Sigma_{\rm maj}+2+C_{\rm low}(\mathsf b))\ell^3]$ of $\mathsf R_{t_1}$, and on $J$, while $\mathsf C_{\rm bf}$ depends on $J$ and on no side.

The same explicit numbers give a form uniform in the label, derived in the proof of Proposition~\ref{C-361-15}(b): there are numbers $y_\sharp$, $C_{\mathfrak o}$ of the constants such that at every label with $\log\mathsf b\ge y_\sharp$, on every side $2MR(\mathsf b)\le\ell\le2L^3M(\log\mathsf b)^r$ and at every $K\ge11(5+2\mathfrak p_0/r)\log_L\ell$, the conditions (k0)--(k4), (k6) hold and $\mathfrak o_K(\ell)\le C_{\mathfrak o}\,\ell^{3+2\mathfrak p_0/r}L^{-K/11}$, by the condition $2\mathfrak p_0>3r+2$ of the list; (iii) for a label $\mathsf b\ge\gamma^{-2}$, a side $\ell=2L^m$ with $MR(\mathsf b)\le L^m$, a cutoff $K\ge2$ at which the conditions (k0)--(k4) and (k6) of Section~\ref{sec:cutoff-calibration} hold, every integer $0\le d\le\lceil K/1260\rceil$ and every bounded Borel function $F$, invariant under the gauge transformations, of the configurations of the lattice of the spacing $L^{-d}$ on the torus $\ell$:
\[
\bigl|\langle F(\mathcal M^{K+1-d}U)\rangle_{K+1}-\langle F(\mathcal M^{K-d}U)\rangle_K\bigr|\le\mathfrak o_K\sup|F|
\]
(Corollary~\textup{\ref{IV-62-66}}(a)), with the number $\mathfrak o_K$ of (ii), which is $\le\mathsf C_{\rm bf}\,\ell^3K^{\mathsf p_0}L^{-K/10}$ at every $K\ge K_{\rm bf}$ by~\textup{(\ref{IV-62-68})}, so that $\sum_K\mathfrak o_K<\infty$; the block field is taken at the fixed spacing $L^{-d}$, and the number $K-d$ of its averaging steps grows with the cutoff; (iv) the exact identification of the zero-source contact terms with the vacuum constants of their first $K^\bullet-u$ steps, proved in Step~3(iv) of Proposition~\ref{C-365-11}, and their two-torus comparison in Lemma~\ref{lem:vacuum-prefix}; Proposition~\ref{C-361-13} below retains this stopping depth and proves the resulting summable estimate; (v) the averaging map among the construction results (equations~\YMMref{(37)--(39)}{20}), which sends a configuration $U$ to its block field $V_j=\mathcal M^jU$ of level $j$, whose blocks are the cubes of $L^j$ bare sites a side centred at the block sites, the points of $L^j\mathbb Z^3$, so that the block layer $1$ of level $q$ lies in the bare layers $\frac12(L^q+1),\dots,\frac12(3L^q-1)$: it is Borel (Proposition~\YMMref{4.5}{20}) and gauge covariant, it commutes with the isometries of the lattice that carry $L^j\mathbb Z^3$ onto itself (Lemma~\ref{lem:setup-lattice-covariance}), a block variable depends only on the bare bonds of the two blocks that its bond joins (Proposition~\YMMref{4.5}{20}), on bare fields whose plaquettes all lie within $\omega_0$ of $\mathbf 1$, with $L^{2q}\omega_0$ below a number of $L$, the block plaquettes of level $j\le q$ lie within $1.003\,L^{2j}\omega_0$ of $\mathbf 1$ (Lemma~\ref{lem:background-iterated-average}), and on a cube of $N$ bare sites a side the holonomies of the straight closed lines of the block field of level $q$ lie within $13.3\,NL^q\omega_0$ of those of the bare field (Lemma~\ref{lem:averaged-line-holonomy}); (vi) the comparison of (iii) with a factor, Proposition~\ref{C-361-23} at the end of this subsection, derived there from Proposition~\textup{\YMMref{115.2}{806}} and the statements of Section~\YMMref{116}{807} named in (ii): in the setting of (iii), for every gauge-invariant Borel function $F$ with values in $[0,1]$ of the block field of depth $d$, $\mu^{K+1}_d(F)\le e^{m_K}\mu^K_d(F)+\epsilon_K$ and $\mu^K_d(F)\le e^{m_K}\mu^{K+1}_d(F)+\epsilon_K$, where $\mu^K_d$ is the law of $\mathcal M^{K-d}U$ at the cutoff $K$ and $m_K$, $\epsilon_K$ are the explicit numbers of that proposition: the errors that are powers of $L^{-K}$ stand in the exponent $m_K$, the large fields of Proposition~\textup{\YMMref{115.2}{806}}, lost-mass bound and the shell term in the additive remainder $\epsilon_K$.

Multiplying these inequalities: there are a number $y_\times$ and a power $\mathsf N(y)$ of $y$, both formed from the constants, such that at every label with $y=\log\mathsf b\ge y_\times$, on the least side $\ell$ at which (iii) applies, and from the cutoff $K_1=K_1(y)$ of that proposition, at which the bare cube has $\ell L^{K_1}<\mathsf N(y)$ sites a side,
\[
\mu^K_0(E)\le1.7\,\bigl(\mu^{K_1}_0(E)+e^{-4y}\bigr)\qquad\text{for every }K\ge K_1
\]
and every Borel set $E$ of block fields of depth $0$ that the gauge transformations carry onto itself: a factor and a remainder, not an additive error summed over the cutoffs.
(M5) Theorems~\YMMref{AJ}{673} and~\YMMref{AD}{607}, through: (i) the bounds on cumulants at first-passage data: for $k\ge2$, every first-passage datum $(K,\beta)$ at $g$, every torus $\ell\in\mathfrak L(g)$ with $\ell\ge16k(R+1)$ and $k$ real Lipschitz functions with supports in $(-\frac\ell4,\frac\ell4)^3$, each in a ball of radius $R$, any two at a distance $\ge\mathsf d$, $|S^T_{k;K,\ell,\beta}(X)|\le C^g_k(\mathsf d;R)\prod_i\|f_i\|_\sharp$; the bounds depend, besides on $g$, $k$, $\mathsf d$, $R$ and the constants, only on a radius fixed once in that theorem, and neither on the cutoff nor on the coupling nor on the side nor on the place; (ii) the theorem on the canonical decomposition, the ergodic components and the reconstruction, Theorem~\textup{\YMMref{AP}{780}}: each ergodic component of the limit of a convergent Gibbs sequence is a limit along the same pairs whose approximants are periodic Gibbs states at these pairs (Theorem~\textup{\YMMref{AP}{780}}, component approximation in the proof), and it carries the reconstruction of Theorem~\textup{\YMMref{AP}{780}}; (iii) the law of the two-point function at short distance (Theorem~\textup{\YMMref{AV}{851}}, short-distance bound and torus version), with the ceiling $\gamma_W$ and the threshold $\mathsf d_0(g,\vartheta)$ of equations~\YMMref{(2346)}{850} and~\YMMref{(2349)--(2350)}{851}, the depth $s(\mathsf d)$ of equation~\YMMref{(2321)}{843} and the test pairs of Lemma~\ref{lem:ward-test-pairs}: for $0<g\le\gamma_W$, $\mathsf d\le\mathsf d_0(g,\vartheta)$, every first-passage datum $(K,\beta)$ at $g$ with $K>s(\mathsf d)$, every periodic Gibbs state $\nu$ at it and every test pair $(f,h)$ at $(\mathsf d;\vartheta)$,
\[
S^T_2[\nu](f,h)=\mathsf c_2\,\mathsf b^{-2}\,\mathcal K_{\rm free}(f,h)\,(1+\varepsilon'),\qquad|\varepsilon'|\le\tfrac13,\qquad\mathcal K_{\rm free}(f,h)=\iint f(x)h(y)|x-y|^{-6}\,dx\,dy,
\]
with $\mathsf c_2>0$ and $\mathcal K_{\rm free}$ as in Theorem~\YMMref{C}{7} and equation~\YMMref{(2240)}{822}; for $f(y)=(1-|y-2\mathsf de_3|^2/\mathsf d^2)_+^2$ and $0\le t\le4\mathsf d$ the pair $(f(\Theta\,\cdot),f(\cdot-te_3))$, $\Theta$ the reflection $x_3\mapsto-x_3$, is a test pair at $(\mathsf d;\vartheta_1)$ (Lemma~\ref{lem:ward-test-pairs}(c)), and $K>s(\mathsf d)$ gives $\mathsf d/a_K>4L^3M$. Besides, $I(g)=(u,Lu-2\bar B]$ with $u=g^{-2}+\bar B$, $\bar B\ge0$, and $g^{-2}>L^2\bar B$ for $g\le\gamma_W$ (the coupling ceiling preceding Theorem~\YMMref{AV}{851}, equation~\YMMref{(2346)}{850}). The labels $a_K\beta$ of all first-passage data $(K,\beta)$ at $g$ lie in the window $\bigl(g^{-2}+B_\sharp\frac{L-2}{L-1},\,L(g^{-2}+\bar B)\bigr]$ (\YMMref{(1531)}{615} in Proposition~\YMMref{74.1}{615}), whose upper end is $X_g+\bar B$.
(M6) The ordered choice leaves the blocking factor free upward (the choice of $L$ in the proof of Proposition~\YMMref{69.1}{595}, followed by the source choices in Section~\YMMref{82}{665}): every odd integer above a threshold is the value of $L$ at an admissible choice, and Theorems~\YMMref{D}{8}, \YMMref{AJ}{673} and~\YMMref{AD}{607} hold at every admissible choice.
(M7) The member of third order in the curvature of the bulk term of the rotational Ward identity, on growing tori: Lemma~\ref{C-360-19} of the section on plane fields and the stress tensor (Section~\ref{C-360-1}), which carries the bound of the cubic remainder on one torus, Lemma~\ref{C-360-10}, to all sides at once. For $\chi,f_1,\dots,f_n$ and $\omega$ as in Theorem~\ref{C-361-4}(5), this member of the bulk term of~\YMMref{(2108)}{782} on the torus $\ell$ is
\[
\mathfrak U^\sharp_{\omega,\chi}(f_1,\dots,f_n)=-\beta\bigl\langle\mathfrak K^\sharp[\chi\,\omega y];\mathcal O_K(f_1);\dots;\mathcal O_K(f_n)\bigr\rangle_{K,\ell,\beta}+a_K^{-1}\sum_{i=1}^n\bigl\langle\mathfrak K^\sharp[f_i\,\omega y];(\mathcal O_K(f_l))_{l\ne i}\bigr\rangle_{K,\ell,\beta},
\]
the joint cumulants in which the local density of the rotation is replaced by its part $\mathfrak K^\sharp[\boldsymbol\eta]=\sum_pa_K^{-1}\boldsymbol\eta^{\rho(p)}(x(p))\,\mathfrak r^\sharp_{\rho(p)}(p)$ of Section~\ref{sec:ward-orbital-complement}, built with the remainders $\mathfrak r^\sharp$ of third order in the field strength (Proposition~\ref{prop:ward-orbital-cubic}), with the weight $\beta/a_K$; here $\boldsymbol\eta^\rho(y)=\chi(y)(\omega y)_\rho$, resp.\ $f_i(y)(\omega y)_\rho$. As $\mathfrak K^\sharp$ depends on its field only through its values at the centres of the plaquettes, at which $\chi f_i=f_i$, this member is the cubic remainder $\mathfrak U^\sharp[\chi\,\omega y](f_1,\dots,f_n)$ of that section at the field $\chi\,\omega y$. Let $0<g\le\gamma$, $n\ge2$, $\mathsf d>0$ and $\rho_0\ge1$. There are $C_\sharp$ and $K_\sharp$, formed from $g$, $n$, $\mathsf d$, $\rho_0$ and the constants and depending on no side, such that for all first-passage data $(K,\beta)$ at $g$ with $K\ge K_\sharp$, all sides $\ell\ge2\rho_0+2$ of $\mathfrak L(g)$, all $\omega\in\mathfrak{so}(3)$ and all real $\chi,f_1,\dots,f_n\in C^{1,1}(\mathbb R^3)$ that vanish outside $[-\rho_0,\rho_0]^3$, $\chi$ equal to $1$ at the centres and at the bonds of the plaquettes at which some $f_i$ does not vanish, the supports of the $f_i$ at pairwise sup distances $\ge\mathsf d$,
\[
\bigl|\mathfrak U^\sharp_{\omega,\chi}(f_1,\dots,f_n)\bigr|\le C_\sharp\,|\omega|\,L^{-K/16}\,\|\chi\|_{C^{1,1}}\prod_i\|f_i\|_{C^{1,1}} .
\]
This is Lemma~\ref{C-360-19} at the field $\chi\,\omega y$, for the cube of side $\frac14\ell'$ centred at $0$, $\ell'$ the least side of $\mathfrak L(g)$ with $\ell'\ge8\rho_0$, and for a rate $c\in(\frac1{16},\frac1{14})$: there $K^pL^{-cK}\le C'L^{-K/16}$, the norm $\mathsf H(\eta)=\max_\rho\sup|\eta^\rho|+\max_\rho{\rm Lip}_\infty\eta^\rho$ of a field in that lemma, ${\rm Lip}_\infty$ the Lipschitz constant for the sup distance, is $\mathsf H(\chi\,\omega y)\le C\rho_0|\omega|\,\|\chi\|_{C^{1,1}}$ with a pure number $C$, and $\|f_i\|_\sharp\le\max\{1,40M\}\|f_i\|_{C^{1,1}}$. On each of the finitely many sides of $\mathfrak L(g)$ from $2\rho_0+2$ to $\ell'$ it is Lemma~\ref{C-360-10}, whose setting holds on that side at all first-passage data at $g$: their labels lie in $[\gamma^{-2},X_g+\bar B]$, and $MR(X_g+\bar B)\le L^m$ on every side $2L^m$ of $\mathfrak L(g)$.
The idea of the proof. On one torus the entry $\mathfrak K^\sharp$ enters each renormalisation step as one more local term, inside the radii of the one-step theorem with sources (Proposition~\ref{prop:source-entry}). Each of its terms vanishes to the third order at $\mathbf 1$ and is odd under the mirror of its plaquette, while the steps are covariant under the mirrors of their lattices; so only the differences of the backgrounds and their mirror images and the asymmetric part of the weight feed the families that the steps make of it, and what is fed at a level is forgotten at the rate $L^{-1/2}$ a level. The conditional mean of the entry, given the block field at a level $K-\lceil\phi K\rceil$ with a small $\phi>0$, is therefore a number up to an oscillation of the order $K^pL^{-cK}$ (Lemma~\ref{C-360-9}(b)), the large fields costing less than every power of $L^{-K}$; the cumulant is a mixed derivative at $0$ of the logarithm of the partition function with the entry and the sources, and Cauchy's estimates in the sources give $|\mathfrak U^\sharp|\le CK^pL^{-cK}$ times the norms at every rate $c<\frac1{14}$ (Lemma~\ref{C-360-10}), which is at most $C'L^{-K/16}$.

On all the sides at once, the carrier is put in one cell of a board of side $\frac12\ell'$, $\ell'$ a fixed side of $\mathfrak L(g)$; the chessboard estimate controls the generating function of the entry and of the sources on the torus $\ell$ by the mean of their mirror copies on the torus $\ell'$, with a factor $e^{(\ell'/2)^3(p_{\ell'}-p_\ell)}$ that is bounded uniformly at the first-passage data at $g$; on the torus $\ell'$ the bound of one torus applies to the copies, and the part of the cumulants that links different copies carries the oscillation of the conditional mean of Lemma~\ref{C-360-9}(b); the finitely many sides between $2\rho_0+2$ and $\ell'$ are tori of their own.
(M8) The rotational Ward identity among the construction results, which Theorem~\ref{C-361-4}(5) and its proof read: the bulk term~\YMMref{(2107)}{782} of equation~\YMMref{(2107)}{782} and the source calculation in the proof of Proposition~\YMMref{112.1}{782}, the bulk insertion of equation~\YMMref{(2107)}{782} and the source calculation in the proof of Proposition~\YMMref{112.1}{782}, whose mean at the sources gives the bulk term by differentiation in the source variables, and the identity~\YMMref{(2108)}{782}; the orbital potential $\Psi_\rho(p)$ and the remainder of third order $\mathfrak r^\sharp_\rho(p)$ of Section~\ref{sec:ward-orbital-complement}, with the bound of third order~\eqref{eq:ward-orbital-real-cubic}, and the summation by parts of Proposition~\ref{prop:ward-orbital-cubic}, which writes the bulk insertion as the sum of the spin part, a member of order $2$ in the curvature and the cubic member~\eqref{eq:ward-orbital-summation}; the small-field density with a local insertion and its derivative as a conditional mean in Lemma~\ref{C-360-9}(a), under its stated entry and cutoff hypotheses, obtained from the holomorphic fourth-family map of Proposition~\ref{prop:source-entry}; and the two-constants inequality of Lemma~\YMMref{116.2}{809}, used in the source variables.

\begin{proposition}[Two cutoffs with a factor]\label{C-361-23}
Let $\mathsf b\ge\gamma^{-2}$ and $\ell=2L^m$ with $MR(\mathsf b)\le L^m$; for a cutoff $K\ge2$ take the numbers $\epsilon_{\rm fg},\epsilon_{\rm cl},\epsilon_{\rm sh},\mathsf R_{t_1}$ of Section~\ref{sec:cutoff-calibration} at $n=1$ ($u=\lceil K/1260\rceil$, $s=\lceil K/9\rceil$, $t_1=K-u$), and put
\[
m_K=4\epsilon_{\rm fg}+2\epsilon_{\rm cl}+1.06\,\epsilon_{\rm sh}+1.39\,\mathsf R_{t_1},\qquad\epsilon_K=\mathsf R_{t_1}+1.29\,\epsilon_{\rm sh}.
\]
\begin{enumerate}
\item[(a)] At every cutoff $K\ge2$ at which (k0)--(k4) and (k6) hold, for every integer $0\le d\le\lceil K/1260\rceil$ and every Borel function $F$ with values in $[0,1]$ of the block fields of depth $d$ on the torus $\ell$ that is invariant under their gauge transformations,
\[
\mu^{K+1}_d(F)\le e^{m_K}\mu^K_d(F)+\epsilon_K,\qquad\mu^K_d(F)\le e^{m_K}\mu^{K+1}_d(F)+\epsilon_K,
\]
$\mu^K_d$ the law of $\mathcal M^{K-d}U$ at the cutoff $K$ (Definition~\textup{\ref{IV-62-62}}(2)).
\item[(b)] If the conditions hold at every $K\ge K_1$, and $\mathsf S=\sum_{K\ge K_1}m_K$ and $\mathsf R_\Sigma=\sum_{K\ge K_1}\epsilon_K$ are finite, then $e^{-\mathsf S}\mu^{K_1}_d(E)-\mathsf R_\Sigma\le\mu^K_d(E)\le e^{\mathsf S}\bigl(\mu^{K_1}_d(E)+\mathsf R_\Sigma\bigr)$ for every $K\ge K_1$, every integer $0\le d\le\lceil K_1/1260\rceil$ and every Borel set $E$ of block fields of depth $d$ that the gauge transformations carry onto itself.
\item[(c)] Let $p=2\mathfrak p_0$, $y=\log\mathsf b$, $\bar\ell(y)=2L^2My^r$, $\mathsf N(y)=L\,\bar\ell(y)^{31}y^{10(p+1)}$, and let $K_1(y)$ be the least integer $K\ge10$ with $L^{K/10}\ge\bar\ell(y)^3y^{p+1}$. There is a number $y_\times$, formed from the constants, such that at every label $\mathsf b$ with $y\ge y_\times$, on the least side $\ell=2L^m$, $m\in\mathbb N$, with $MR(\mathsf b)\le L^m$, which has $\ell\le\bar\ell(y)$: the conditions (k0)--(k4) and (k6) hold at every $K\ge K_1(y)$; $\ell L^{K_1(y)}<\mathsf N(y)$; from $K_1=K_1(y)$, $\mathsf S\le0.53$ and $\mathsf R_\Sigma\le e^{-4y}$; and so $\mu^K_0(E)\le1.7\,\bigl(\mu^{K_1(y)}_0(E)+e^{-4y}\bigr)$ for every $K\ge K_1(y)$ and every Borel set $E$ of block fields of depth $0$ that the gauge transformations carry onto itself.
\end{enumerate}
\end{proposition}

\begin{proof}
(a) Write $\langle F\rangle_K=\mu^K_d(F)$. The averaging maps from the depth $u$ to the depth $d$ are the same at the two cutoffs and commute with the gauge transformations, so $F\circ\mathcal M^{u-d}$ is a function of the same kind at the depth $u$, and we may take $d=u$. Three comparisons follow, each with the same function $F\ge0$ inserted on both sides. (1) \emph{Large fields.} By Proposition~\textup{\YMMref{115.2}{806}}, image identity and domination bound, the chain of the top level $t_1$, on its inner ball, is the image under $\mathcal M^{t_1}$ of $\mathfrak w^{\rm sf}_{t_1}\rho_0\,dU$, and its integral of a gauge-invariant $F\ge0$ is at most $\int F(\mathcal M^{t_1}U)\rho_0\,dU$; by the discarding of Proposition~\textup{\YMMref{115.2}{806}}, lost-mass bound its mass is at least $(1-\mathsf R_{t_1})Z$.

So the mean of $F$ under the law exceeds its mean under the normalised chain by at most $\mathsf R_{t_1}\sup F$, and the latter is at most $(1-\mathsf R_{t_1})^{-1}$ times the former. (2) \emph{The chain and the classical start.} On the small-field set the two measures have the positive densities $e^{c(0)}e^{\Phi}$ and $e^{\hat c(0)}e^{\hat\Phi}$ (the proof of Proposition~\YMMref{116.1}{808}), and $|\Phi-\hat\Phi|\le\epsilon_{\rm fg}$ at every block field of the cut set (equation~\YMMref{(2180)}{803}); so the two normalised means of $F$ differ by a factor within $e^{\pm2\epsilon_{\rm fg}}$, the constants $e^{c(0)}$, $e^{\hat c(0)}$ cancelling against $F=1$. (3) \emph{The classical starts of the two cutoffs.} They are integrals over the configurations of the depth $s$ of the weights $\mathsf k\ge0$ of the two cutoffs times densities that are positive on the cut sets (Section~\YMMref{116}{807}, Section~\ref{sec:block-field-laws}); where both weights are positive they are equal and the densities differ by a factor within $e^{\pm\epsilon_{\rm cl}}$ (Proposition~\textup{\YMMref{116.1}{808}}, density comparison in the proof), and where one of them vanishes the other carries at most the part $\epsilon_{\rm sh}$ of the mass (Section~\ref{sec:block-field-laws}); so the normalised mean at one cutoff is at most $e^{2\epsilon_{\rm cl}}(1-\epsilon_{\rm sh})^{-1}$ times that at the other plus $\epsilon_{\rm sh}\sup F$.

Chaining (1), (2), (3), (2), (1), for $\{K',K''\}=\{K,K+1\}$,
\[
\langle F\rangle_{K'}\le\frac{e^{4\epsilon_{\rm fg}+2\epsilon_{\rm cl}}}{(1-\epsilon_{\rm sh})(1-\mathsf R_{t_1})}\,\langle F\rangle_{K''}+\bigl(\mathsf R_{t_1}+e^{2\epsilon_{\rm fg}}\epsilon_{\rm sh}\bigr)\sup F,
\]
and (k2), (k3), (k6) give $e^{2\epsilon_{\rm fg}}\le e^{1/4}<1.29$ and, by convexity, $-\log(1-\epsilon_{\rm sh})\le1.06\,\epsilon_{\rm sh}$ and $-\log(1-\mathsf R_{t_1})\le1.39\,\mathsf R_{t_1}$. (b) By induction on $K$, $\langle F\rangle_K\le e^{\sigma_K}(\langle F\rangle_{K_1}+\rho_K)$ and $\langle F\rangle_K\ge e^{-\sigma_K}\langle F\rangle_{K_1}-\rho_K$, with the sums $\sigma_K$, $\rho_K$ of the $m_k$, $\epsilon_k$ over $K_1\le k<K$; $d\le\lceil K_1/1260\rceil\le\lceil K/1260\rceil$ keeps every step in the range of (a), and $F=\mathbf 1_E$. (c) Let $K\ge K_1(y)$; with $\ell\le\bar\ell(y)$, $\ell^3L^{-K/10}\le y^{-p-1}$. As $u\ge K/1260$ and $s\ge K/9$, the definitions~\YMMref{(2192)}{808} and~\YMMref{(2193)}{808} give $\epsilon_{\rm fg}\le a_1\ell^3L^{-K/10}$ and $\epsilon_{\rm cl}\le a_2\,y^p(1+K\log L)^pK^{1/2}L^{-K/90}\,\ell^3L^{-K/10}$ with numbers $a_1,a_2$ of the constants, so that $\sum_{K\ge K_1(y)}(4\epsilon_{\rm fg}+2\epsilon_{\rm cl})\le(17.72\,a_1+2a_2a_3)/y$, $a_3=\sum_{K\ge1}(1+K\log L)^pK^{1/2}L^{-K/90}$.

The numbers $\mathsf R_{t_1}$ and $\epsilon_{\rm sh}$ carry the factors $e^{-c'A_0^2\mathsf y_{[u]}}$ and $e^{-\frac14\kappa_{\rm b}A_1^2\mathsf y_{[u]}}$, where $\mathsf y_{[u]}=(y+u\log L)^p\ge y^p+(u\log L)^p$, against the growth $\ell^3L^{3s}$, $s\le140u$, and, in $\mathsf R_{t_1}$, $\exp[(\mathsf C^\Sigma_{\rm maj}+2+c_{\rm low}+\frac92y)\ell^3]\le\exp[O(y^{3r+1})]$; summed over the cutoffs, of which at most $1260$ have the same $u$, and by $2\mathfrak p_0>3r+2$, which follows from the exponent choice $r\ge2$, $\mathfrak p_0\ge6r+1$ in Theorem~\YMMref{D}{8}, $\sum_{K\ge K_1(y)}(\mathsf R_{t_1}+1.29\,\epsilon_{\rm sh})\le e^{-4y}$ for all large $y$. The same bounds give (k0)--(k4) and (k6) at every $K\ge K_1(y)$ for large $y$; $y_\times$ is the least number from which all this holds, with $\mathsf b\ge\max\{\gamma^{-2},\Theta_\star\}$ and $y\ge2(17.72\,a_1+2a_2a_3)$.

Then $\mathsf S\le\frac12+1.39\,e^{-4y}\le0.53$, $e^{0.53}<1.7$, and (b) at $d=0$ gives the last inequality. As $L<\bar\ell(y)^3$, $K_1(y)\ge11$, and by its least choice $L^{(K_1(y)-1)/10}<\bar\ell(y)^3y^{p+1}$, so $\ell L^{K_1(y)}<\bar\ell(y)\,L\,(\bar\ell(y)^3y^{p+1})^{10}=\mathsf N(y)$.
\end{proof}

\subsection{Statements on one lattice}

\begin{proposition}[The exact layer]\label{C-361-7}
Let $\beta\ge0$, $\Sigma=(A,B)$ a cross-section with sides $\ge3$, any coordinate direction being the time, and $t,t'\ge3$, $q\ge2$ integers.
\begin{enumerate}
\item[(a)] $Z(t,\Sigma)=\operatorname{Tr}\mathsf D^t=\sum_i\lambda_i^t$, $\sum_i\lambda_i^2<\infty$, $\lambda_0$ is simple and its eigenvector can be taken strictly positive. $F$ is symmetric in the three sides, and $e^{-6\beta|\underline n|}\le Z\le1$.
\item[(b)] $\mathsf r_{t+1}\le\mathsf r_t$, $\mathsf r_{t+t'}\le\mathsf r_t\mathsf r_{t'}$, $\mathsf r_{qt}\le\mathsf r_t^{\,q}$; $F(t,\Sigma)=t\log\lambda_0+\mathsf G(t;\Sigma)$, $F(t,\Sigma)/t$ does not increase, and $F(t+t',\Sigma)\le F(t,\Sigma)+F(t',\Sigma)$;
\[
qF(t,\Sigma)-F(qt,\Sigma)=q\,\mathsf G(t;\Sigma)-\mathsf G(qt;\Sigma)\in\bigl[(q-1)\mathsf G(t;\Sigma),\,q\,\mathsf G(t;\Sigma)\bigr].
\]
If $2F(t,\Sigma)-F(2t,\Sigma)\le x<\log2$, then $\mathsf r_t\le\varphi(e^x-1)$ and $2F(2t,\Sigma)-F(4t,\Sigma)\le\psi(x)$, where $\varphi(\eta)=\eta/(1+\sqrt{1-\eta^2})$ and $\psi(x)=2\log(1+\varphi(e^x-1)^2)$.
\item[(c)] $F/|\underline n|$ and $\delta/|\underline n|$ do not increase when one side grows, and $F/|\underline n|\to P$ along boxes all of whose sides tend to $\infty$; $\delta(t,\Sigma)=t\,\varepsilon(\Sigma)+\mathsf G(t;\Sigma)$ with $\varepsilon(\Sigma)\ge0$, so $\mathsf G\le\delta$; $D_q(N)=q^3\delta(N^3)-\delta((qN)^3)\in[0,q^3\delta(N^3)]$, $\mathsf e(l)=2\delta(l,2l,2l)-\delta((2l)^3)$, and
\[
\delta(N^3)=\sum_{s\ge0}q^{-3(s+1)}D_q(q^sN).
\]
\emph{From above, an area:} for all integers $A,B,t\ge N$, $\mathsf G(t;A,B)\le AB\,\varkappa$ and $\varepsilon(A,B)\le(A+B-N)\varkappa$, $\varkappa=\delta(N^3)/N^2$; so for every box with all sides $\ge N$
\[
\delta(\underline n)\le\bigl[n_1n_2+n_1n_3+n_2n_3-N(n_1+n_2+n_3)+N^2\bigr]\varkappa\le(n_1n_2+n_1n_3+n_2n_3)\,\varkappa,
\]
and $\delta(M^3)\le(3\rho^2-3\rho+1)\,\delta(N^3)$ for $M=\rho N\ge N$.
\item[(d)] There are commuting unitary involutions $\mathsf J_1,\mathsf J_2$ that commute with $\mathsf T^{(0)}=\mathsf D$ and with an operator $\mathsf T^{(1)}\ge0$ of the same kind ($\mathsf D$ with the sign of one plaquette of the cross-section reversed in $\mathsf m$), such that, with $\mathsf J^k=\mathsf J_1^{k_1}\mathsf J_2^{k_2}$ and $\mathsf P_e=\frac14\sum_k(-1)^{e\cdot k}\mathsf J^k$,
\[
Z^{(k;m)}(t,\Sigma)=\operatorname{Tr}\bigl(\mathsf J^k(\mathsf T^{(m)})^t\bigr)=\sum_e(-1)^{e\cdot k}W(e;m),\qquad W(e;m)=\operatorname{Tr}\bigl(\mathsf P_e(\mathsf T^{(m)})^t\bigr)\ge0 .
\]
Each $W(e;m)$ is a sum of $t$-th powers of positive numbers that do not depend on $t$; the eigenvector of $\lambda_0$ lies in the range of $\mathsf P_{00}$, and $W(00;0)\ge\lambda_0^t$; $Z^z>0$ does not depend on the columns chosen. On every box: $0\le1-r^{(k;0)}\le2\mathsf R_t$, $0<r^{(k;m)}\le r^{(00;m)}$, $r^{(11;m)}\ge r^{(10;m)}+r^{(01;m)}-r^{(00;m)}$, and
\[
\delta(t,\Sigma)=t\,\varepsilon(\Sigma)+\log(1+\mathsf r^{00}_t)+\log(1/\zeta_t),\qquad\zeta_t=\frac{W(00;0)}Z=\frac14\Bigl(1+\sum_{k\ne0}r^{(k;0)}\Bigr),
\]
three terms $\ge0$, $1+\mathsf r^{00}_t=W(00;0)/\lambda_0^t$. On a cube $N^3$: $r_2,r_3\le r_1\le1$; all seven $0\le1-r^z\le3\mathsf R_N\le3\min\{\mathsf G_N,\mathsf r_N\}$; $1-r_1\le\frac43\mathsf R_N$, $1-r_2\le2(1-r_1)$ and $1-r_3\le3(1-r_1)$.
\item[(e)] $2F(t,\Sigma)-F(2t,\Sigma)\ge2\log\sum_{f\in\{0,1\}^2}w_f^{1/2}$, the $w_f$ being the weights of the box $(2t,\Sigma)$; and $\mathsf G(t;\Sigma)\ge\log(1/w_{00})$, $w_{00}$ the weight of the box $(t,\Sigma)$.
\item[(f)] Let $\Gamma=\mathbb Zn_1e_1\oplus\mathbb Zn_2e_2\oplus\mathbb Z(v_1e_1+v_2e_2+n_3e_3)$, $v\in\mathbb Z^2$, and $Z_\Gamma$ the partition function of the lattice with these periods. With the time $e_3$ over $\Sigma=(n_1,n_2)$: $Z_\Gamma=\operatorname{Tr}(\mathsf U_v\mathsf D^{n_3})$, $\mathsf U_v$ the unitary of the translation of a layer by $v$, and $\lambda_0^{n_3}(1-\mathsf r_{n_3})\le Z_\Gamma\le Z(\underline n)$.
\item[(g)] \emph{The sheared cube.} Let $N=2h$ be even, $\Gamma^\diamond=\mathbb ZN(e_1-e_3)\oplus\mathbb ZNe_2\oplus\mathbb Zh(e_1+e_3)$, a period lattice of the index $N^3$, and $Z_{\Gamma^\diamond}$ the partition function of the lattice with these periods. Then $Z(N^3)\le Z_{\Gamma^\diamond}$.
\end{enumerate}
\end{proposition}

\begin{proof}
(a) The weight of the box is a product over the layers of $\mathsf m(V^\tau)$ and of the crossing weights $e^{-\beta A_\times(V^\tau,S^\tau,V^{\tau+1})}$, each $S^\tau$ standing in one factor; integrating the $S^\tau$ (Fubini) gives 
\[
\int\prod_\tau dV^\tau\prod_\tau\mathsf m(V^\tau)\mathsf K(V^{\tau+1},V^\tau),
\]
and splitting each $\mathsf m(V^\tau)$ into two roots this is 
\[
\int\prod_\tau\mathsf D(V^{\tau+1},V^\tau)\,dV^\tau=\operatorname{Tr}\mathsf D^t,
\]
for every $t\ge3$. With $\mathsf A=\mathsf K^{1/2}\mathsf m^{1/2}$, $\mathsf A^*\mathsf A=\mathsf D$ and $\mathsf A\mathsf A^*=\mathsf T'$ have the same nonzero spectrum. The kernel of $\mathsf D$ is continuous and bounded below by a positive number on a compact space; if $e_0$ is a real eigenvector of $\lambda_0$, so is $|e_0|$, hence so are $|e_0|\pm e_0\ge0$, one of which vanishes on a set of positive measure and is therefore $0$.

So every eigenvector of $\lambda_0$ has a strict sign, and no two of them are orthogonal. A permutation of the coordinates carries plaquettes to plaquettes and $\operatorname{tr}\partial U$ to itself; $0\le a_p\le2$ on $3|\underline n|$ plaquettes. (b) With $x_i=(\lambda_i/\lambda_0)^t\le1$: $x_i^{(t+1)/t}\le x_i$, $\sum x_i^q\le(\sum x_i)^q$, $\sum\lambda_i^{t+t'}\le\sum\lambda_i^t\sum\lambda_i^{t'}$; $F(t)/t=\log\|\lambda\|_{\ell^t}$ does not increase in $t$. The identity and its range follow from $F(t)=t\log\lambda_0+\mathsf G(t)$ and $0\le\mathsf G(qt)\le\mathsf G(t)$. If $2F(t)-F(2t)\le x$, then $\mathsf G(t)\le x<\log2$, $\mathsf r_t<1$, and
\[
e^x\ge\frac{(1+\mathsf r_t)^2}{1+\mathsf r_{2t}}\ge\frac{(1+\mathsf r_t)^2}{1+\mathsf r_t^2}=1+\frac{2\mathsf r_t}{1+\mathsf r_t^2},
\]
which gives $\mathsf r_t\le\varphi(e^x-1)$ as $r\mapsto2r/(1+r^2)$ increases on $[0,1]$; and $2F(2t)-F(4t)\le2\mathsf G(2t)\le2\log(1+\mathsf r_t^2)$. (c) is (b) in each direction, divided by the product of the two other sides; $\varepsilon(\Sigma)=\lim_t(F(t,\Sigma)-t|\Sigma|P)/t\ge0$; in $D_q$ and $\mathsf e$ the terms with $P$ cancel, and the series telescopes, its tail $q^{-3S}\delta((q^SN)^3)=N^3[F/|\cdot|-P]$ tending to $0$. For the area: $\mathsf G(t;A,B)\le\mathsf G(N;A,B)\le\delta(N,A,B)\le AB\varkappa$ by the monotonicity of $\mathsf G$ and of $\delta/|\underline n|$. With the side $B$ as the time, $\delta(s,t,B)=B\varepsilon(s,t)+\mathsf G(B;s,t)\le B\varepsilon(s,t)+st\varkappa$; dividing by $s$ and then $t$ and letting both tend to $\infty$, as $\varepsilon(s,t)/st\to0$, gives $\limsup_t\varepsilon(t,B)/t\le\varkappa$.

Since $\delta(a,b,c)=a\varepsilon(b,c)+\mathsf G(a;b,c)\le\delta(N,b,c)+(a-N)\varepsilon(b,c)$ for $a\ge N$, $\varepsilon(A,B)\le\varepsilon(N,B)+(A-N)\varkappa\le\varepsilon(N,N)+(A+B-2N)\varkappa$, and $\varepsilon(N,N)\le\delta(N^3)/N=N\varkappa$. The same inequality in the three sides in turn gives $\delta(\underline n)\le\delta(N^3)+(n_3-N)\varepsilon(N,N)+(n_2-N)\varepsilon(N,n_3)+(n_1-N)\varepsilon(n_2,n_3)$, which is the bracket. (d) As $-1$ is central in $\mathrm{SU}(2)$, reversing the sign of $\operatorname{tr}\partial U(p)$ on a column that crosses one time step is the same as multiplying by $-1$ the bonds of one line of the next layer; this change of variables keeps the Haar measure, appears under the trace as an involution $\mathsf J_i$ of the layer variables, and commutes with $\mathsf m$ and $\mathsf K$.

A reversed column in the plane of the cross-section lies in every layer and replaces $\mathsf m$ by $\mathsf m_1$. Moving a column by one step is a further flip of the bonds of a line, which does not change the trace. The $\mathsf P_e$ are the isotypical projections of $k\mapsto\mathsf J^k$; they commute with $\mathsf T^{(m)}$, so every eigenspace splits along them and $W(e;m)=\sum_{i:f(i)=e}\mu_i^t$. The positive top eigenvector is invariant under the flips, so it lies in the range of $\mathsf P_{00}$. Then $Z-Z^{(k;0)}=2\sum_{e\cdot k\text{ odd}}W(e;0)\le2(Z-\lambda_0^t)$, $Z^{(k;m)}\le\sum_eW(e;m)=Z^{(00;m)}$, $4W(11;m)=Z^{(00;m)}-Z^{(10;m)}-Z^{(01;m)}+Z^{(11;m)}\ge0$, and $\delta=t\varepsilon+\log(Z/\lambda_0^t)$ with $Z/\lambda_0^t=[W(00;0)/\lambda_0^t][Z/W(00;0)]$.

On a cube, by the cubic symmetry, with $u=w_{10}=w_{01}$, $v=w_{11}$: $1-r_1=2(u+v)$, $1-r_2=4u$, $2u+v=1-w_{00}\le\mathsf R_N$; two planes with a common direction as the time give $r_2\le r^{(00;1)}=r_1$, hence $v\le u$ and $1-r_1\le\frac43\mathsf R_N$; three planes give $r_3\ge2r_2-r_1$; and the sum over $e$ with a vanishing component gives $1-r^{(1,1,1)}\le(1-r^{z'})+(1-r^{z''})$, so $1-r_3\le(1-r_1)+(1-r_2)\le3(1-r_1)$. (e) Within each class $(\sum\mu_i^t)^2\ge\sum\mu_i^{2t}$, so $Z(t,\Sigma)\ge\sum_fW_{2t}(f)^{1/2}=Z(2t,\Sigma)^{1/2}\sum_fw_f^{1/2}$; and $Z/\lambda_0^t\ge Z/W(00;0)$. (f) Integrating the layers as in (a), the last layer is glued to the first through the translation by $v$, which produces $\mathsf U_v$ under the trace; $\mathsf U_v$ commutes with $\mathsf D$, so in a common eigenbasis $Z_\Gamma=\sum_i\omega_i\lambda_i^{n_3}$ with $|\omega_i|=1$ and $\omega_0=1$ (the top eigenvector is positive and unique, hence invariant); $Z_\Gamma$ is real and positive, so $\lambda_0^{n_3}-\sum_{i\ge1}\lambda_i^{n_3}\le Z_\Gamma\le\sum_i\lambda_i^{n_3}$. (g) Put $\mathsf v=e_1+e_3$.

The planes $x_1+x_3=c$, $c\in\mathbb Z$, are planes of lattice reflections, $(x_1,x_2,x_3)\mapsto(c-x_3,x_2,c-x_1)$; $N\mathbb Z^3=\Gamma_\perp\oplus\mathbb ZNe_1$ and $\Gamma^\diamond=\Gamma_\perp\oplus\mathbb Zh\mathsf v$ with $\Gamma_\perp=\mathbb ZN(e_1-e_3)\oplus\mathbb ZNe_2\perp\mathsf v$, and both generators $Ne_1$, $h\mathsf v$ raise the height $x_1+x_3$ by $N$. Cut either torus along the planes of the heights $-h,0,h$ into two slabs of the height $h$, periodic under $\Gamma_\perp$. No plaquette lies in such a plane, and no bond crosses it; a plaquette cut by it has two corners $z,z'$ in it and one on either side, and its weight is $e^{-\beta(1-\frac12\operatorname{tr}vv'^{-1})}$, $v$ and $v'$ the transports from $z$ to $z'$ through the upper and through the lower corner.

As $\frac12\operatorname{tr}vv'^{-1}=\frac12\sum_{i,j}v_{ij}\overline{v'_{ij}}$ on $\mathrm{SU}(2)$, the exponential series writes this weight as a sum with coefficients $\ge0$ of products $\Phi(v)\overline{\Phi(v')}$ of matrix elements. With the data $\phi$ of a plane, its bonds and the indices of this expansion at its cut plaquettes, under the product of Haar measure and of the coefficients, the partition function of the torus $\Gamma_\perp\oplus\mathbb Zg$ is $\iint k_2(\tau_g\phi_0,\phi_1)\,k_1(\phi_1,\phi_0)$, where $k_1$, $k_2$ are the kernels of the two slabs between the data of their walls and $\tau_g$, the translation by $g$, carries the data of the height $-h$ to those of the height $h$ and keeps their measure.

Schwarz' inequality gives $Z(N^3)^2\le\|k_1\|_2^2\,\|k_2\|_2^2$. The reflection in the middle plane exchanges the two sides of the cut plaquettes, conjugates the kernels, and acts on the plane of the height $-h$ as the translation by $h\mathsf v$; so $\|k_i\|_2^2$ is the partition function of the slab $i$ glued to its mirror image, that is of the torus $\Gamma_\perp\oplus\mathbb Zh\mathsf v=\Gamma^\diamond$, and $Z(N^3)\le Z_{\Gamma^\diamond}$.
\end{proof}

\begin{proposition}[The recursion under doubling]\label{C-361-8}
Let $\beta\ge0$, $l\ge3$, and $\mathsf e(l)\le x_0$ at $\beta$; put $\mathsf s(x)=\psi(4x+2\psi(x)+\psi(2x+\psi(x)))$.
\begin{enumerate}
\item[(a)] $\mathsf e(2l)\le\mathsf s(\mathsf e(l))\le12.696\,\mathsf e(l)^2\le12.9\,\mathsf e(l)^2$, $D_2(2l)\le19.669\,\mathsf e(l)^2\le20\,\mathsf e(l)^2$, and $\mathsf e(2^kl)\le\theta^{2^k}/12.9$ for $k\ge0$.
\item[(b)] $\delta((2^jl)^3)\le\frac{20}7\mathsf e(2^{j-1}l)^2\le0.017170\,\theta^{2^j}$, and $\delta((2^jl)^3)\le0.01655\,\theta^{2^j}$, for $j\ge1$; $\delta(l,2l,2l)\le\frac{x_0}2+\frac{10}7x_0^2$; for every integer $N\ge2l$: $\delta(N^3)<0.13736\,\theta^{N/2l}$; $\mathsf e(l')<0.3702\,\theta^{l'/2l}$ for every integer $l'\ge l$; $D_q(N)<q^3\,0.13736\,\theta^{N/2l}$; and $\delta(\underline n)<|\underline n|\,n_-^{-3}\,0.13736\,\theta^{n_-/2l}$ for every box whose least side $n_-$ is $\ge2l$.
\item[(c)] On every cube $N^3$ with $N\ge2l$: $\mathsf G_N\le0.06754\,\theta^{N/2l}$, $\mathsf r_N<0.1473\,\theta^{N/2l}$, all seven $0\le1-r^z<0.4121\,\theta^{N/2l}$, and $1-r_1\le0.09006\,\theta^{N/2l}$, $1-r_2\le0.1351\,\theta^{N/2l}$, $1-r_3\le0.2252\,\theta^{N/2l}$; at $N=2^{k+1}l$, $\mathsf r_N\le0.001615\,\theta^{N/l}$.
\item[(d)] If $x_0=0.069$: $1-r_1(N)<0.0039$ for every $N\ge16l$, and $\mathsf e(l')<0.0234$ for every integer $l'\ge16l$.
\item[(e)] \emph{The threshold.} $\mathsf s(x)\le0.99991\,x$ on $(0,0.07489]$, and the least positive fixed point of $\mathsf s$ lies in $(0.0748941,0.0748942)$. No threshold $\ge\log2$ can serve for a recursion of this kind: the function $F=|\underline n|P+\log2$, which in each direction is a sum of $t$-th powers of two equal positive numbers and has the properties of Proposition~\ref{C-361-7}(b), (c), has $\mathsf e\equiv\delta\equiv\log2$; and the positivity of the $\lambda_i$ is used: $e^F=e^{|\underline n|P}\prod_i(1+(-1)^{n_i})$ has $\mathsf e=-\infty$ at odd $l$ and $\mathsf e=\log8$ at even $l$.
\end{enumerate}
\end{proposition}

\begin{proof}
Let $\epsilon=\mathsf e(l)$, the doubling defect at $t=l$ over the cross-section $(2l,2l)$. We double the three sides in turn, each time with Proposition~\ref{C-361-7}(b). First $\mathsf a=2F((2l)^3)-F(2l,2l,4l)\le\psi(\epsilon)$. Next $y_1=2F(l,2l,4l)-F(2l,2l,4l)\le2\epsilon+\mathsf a$, because $F(l,2l,4l)\le2F(l,2l,2l)$, and $\mathsf b'=2F(2l,2l,4l)-F(2l,4l,4l)\le\psi(y_1)$. Last $2F(l,4l,4l)-F(2l,4l,4l)\le4\epsilon+2\mathsf a+\mathsf b'$, and
\[
\mathsf e(2l)\le\psi(4\epsilon+2\mathsf a+\mathsf b')\le\mathsf s(\epsilon),\qquad D_2(2l)=4\mathsf a+2\mathsf b'+\mathsf e(2l).
\]
$\psi$ and $\mathsf s$ are power series in $x$ with coefficients $\ge0$ beginning with $\frac12x^2$ and $8x^2$, so $\mathsf s(x)/x^2$ does not decrease, and $\epsilon\le0.069$ gives $12.696\,\epsilon^2$ and $19.669\,\epsilon^2$, all arguments staying below $0.293<\log2$; since $12.9\cdot0.069<1$, iteration gives (a). Put $\mathsf s_D(x)=4\psi(x)+2\psi(2x+\psi(x))+\mathsf s(x)$, a power series in $x$ with coefficients $\ge0$ beginning with $14x^2$, so that $\mathsf s_D(x)/x^2$ does not decrease; it is at most $19.668665$ at $0.069$, the number $19.669$ of (a) before rounding. As $\psi$ increases, the display, with $m=2^il$, $i\ge0$, in the place of $l$, gives $D_2(2m)\le\mathsf s_D(\mathsf e(m))\le19.668665\,\mathsf e(m)^2$, as $\mathsf e(m)\le\theta^{2^i}/12.9\le0.069$ by (a); and $\mathsf e$ decreases along these sides, as $\mathsf e(2m)\le12.9\cdot0.069\,\mathsf e(m)$. The series of Proposition~\ref{C-361-7}(c) at $q=2$, $\delta((2^jl)^3)=\sum_{s\ge0}8^{-(s+1)}D_2(2^{s+j}l)$, gives for $j\ge1$
\[
\delta((2^jl)^3)\le19.668665\sum_{s\ge0}8^{-(s+1)}\mathsf e(2^{s+j-1}l)^2\le\frac{19.668665}7\,\mathsf e(2^{j-1}l)^2\le\frac{19.668665}{7\cdot12.9^2}\,\theta^{2^j}\le0.016885\,\theta^{2^j},
\]
which is below $\frac{20}7\mathsf e(2^{j-1}l)^2$ and below $\mathsf c_0\theta^{2^j}$, $\mathsf c_0=\frac{20}{7\cdot12.9^2}\approx0.0171693$. Keeping the decay of the terms, $\mathsf e(2^{s+j-1}l)^2\le\theta^{2^{s+j}}/12.9^2$, the same series gives $\delta((2^jl)^3)\le\frac{19.668665}{8\cdot12.9^2}\,\theta^{2^j}\sum_{s\ge0}8^{-s}\theta^{2^{s+j}-2^j}\le0.0147743\cdot1.10720\,\theta^{2^j}\le0.01636\,\theta^{2^j}$, the sum being greatest at $j=1$ and $\theta=0.8901$. For $2^jl\le N<2^{j+1}l$, $j\ge1$, the monotone density gives $\delta(N^3)\le(N/2^jl)^3\delta((2^jl)^3)<8\cdot0.016885\,\theta^{2^j}<0.13736\,\theta^{N/2l}$, and boxes follow against the cube of their least side. For $\mathsf e(l')$ with $n=2^kl\le l'<2n$, $k\ge0$: $\mathsf e(l')\le2\mathsf G(n;2l',2l')\le2\delta(n,2l',2l')\le2(l'/n)^2\delta(n,2n,2n)<8\delta(n,2n,2n)=4[\mathsf e(n)+\delta((2n)^3)]$ by Proposition~\ref{C-361-7}(b), (c), and the bound above at $j=k+1$ with $\mathsf e(n)\le0.069$ gives $4[\mathsf e(n)+\delta((2n)^3)]\le4\,\mathsf e(n)\bigl(1+\frac{19.668665}7\,0.069\bigr)\le\frac{4\cdot1.193877}{12.9}\,\theta^{2^k}\le0.370195\,\theta^{2^k}<0.3702\,\theta^{l'/2l}$. (c) For $\bar N=2^jl\le N<2\bar N$, $j\ge1$, with the dyadic side as the time, $\mathsf G_N\le\mathsf G(\bar N;N,N)\le\delta(\bar N,N,N)\le4\delta(\bar N^3)\le\frac{4\cdot19.668665}{7\cdot12.9^2}\,\theta^{2^j}\le0.0675395\,\theta^{2^j}\le0.06754\,\theta^{N/2l}$.

Then Proposition~\ref{C-361-7}(d): in the signs of its proof, $1-r_1\le\frac43\mathsf R_N$, $1-r_2=4u\le2(2u+v)\le2\mathsf R_N$ and $1-r_3\le(1-r_1)+(1-r_2)\le\frac{10}3\mathsf R_N$, with $\mathsf R_N\le\mathsf G_N$, which give $0.0900527$, $0.135079$, $0.225132$ times $\theta^{N/2l}$; all seven $1-r^z\le3\mathsf R_N\le3\mathsf G_N$, and $\mathsf r_N=e^{\mathsf G_N}-1$ with $e^{cy}-1\le y(e^c-1)$ for $y\le1$. At $N=2^{k+1}l$, with $n=2^kl$: $\mathsf r_N\le\mathsf r_n(N,N)^2\le\varphi(e^{\mathsf e(n)}-1)^2$ by Proposition~\ref{C-361-7}(b), and $\varphi(e^x-1)^2/x^2$ does not decrease and is at most $0.268652$ at $0.069$; with $\mathsf e(n)\le\theta^{2^k}/12.9$ this is $\mathsf r_N\le0.268652\,\theta^{N/l}/12.9^2\le0.001615\,\theta^{N/l}$. (d) Iterate $X_{k+1}=\mathsf s(X_k)$ from $X_0=0.069$: $X_1,\dots,X_5\approx0.0604446,\ 0.0432902,\ 0.0195722,\ 0.0034309,\ 0.0000960$; the series gives $\delta((2^jl)^3)\le\sum_s8^{-(s+1)}\mathsf s_D(X_{j-1+s})$, with the $\mathsf s_D$ above and $\mathsf e(2^il)\le X_i$ ($\mathsf s$ increases), which at $j=4$ is at most $0.00073051$; so for $N\ge16l$, $\mathsf G_N\le0.0029221$ and $1-r_1\le\frac43(1-e^{-\mathsf G_N})<0.0039$, and $\mathsf e(l')\le32\cdot0.00073051<0.0234$. (e) $\mathsf s$ is star-shaped, and $\mathsf s(x)/x$ is approximately $0.99999914$ at $0.0748941$ and approximately $1.00000135$ at $0.0748942$.

The two examples satisfy the identities of Proposition~\ref{C-361-7}(a)--(c) with two equal top eigenvalues, respectively with an eigenvalue $-\lambda$.
\end{proof}

\begin{proposition}[Every cross-section: a simple top, a gap, the thermal law]\label{C-361-19}
Let $\beta\ge0$, $l\ge3$, $\mathsf e(l)\le x_0$ at $\beta$, and $A\le B$ integers $\ge2l$; let $n=2^kl$ with $2n\le A<4n$ and $n'=2^jn$ with $2n'\le B<4n'$; any direction is the time.
\begin{enumerate}
\item[(a)] Over $(A,B)$ the top $\lambda_0$ is simple and lies in the sector $00$, and $\lambda_1/\lambda_0<0.49067^{1/n}\theta^{1/l}<\theta^{1/l}$; so every level of a sector $e\ne00$, and every level of the sector $00$ but the top, lies more than $\frac1l\log\frac1\theta+\frac{1.4239}A$ a step above the top. If $A=2l$: $\lambda_1/\lambda_0<(2.42\,x_0)^{1/l}$, and $\mathsf r_t(2l,B)\le0.5628\,(2.42\,x_0)^{t/l}$ for $t\ge n'$. Over the square $(N,N)$, $N\ge2l$, with $2n\le N<4n$: $\mathsf r_n(N,N)\le0.20395\,\theta^{n/l}$.
\item[(b)] $\mathsf r_t(A,B)\le0.2314\,\theta^{t/l}$ for every integer $t\ge n'$, and for every integer $t\ge l$
\[
\mathsf G(t;A,B)\le0.37\,\Bigl\langle\frac A{2t}\Bigr\rangle\Bigl\langle\frac B{2t}\Bigr\rangle\,\theta^{t/l}\le0.37\,\frac{AB}{4l^2}\,\theta^{t/l},
\]
in detail $0.27263\frac{AB}{4t^2}\theta^{t/l}$ if $2t\le A$, $0.3695\frac B{2t}\theta^{t/l}$ if $A<2t\le B$, and $\mathsf r_t(A,B)\le0.26261\,\theta^{t/l}$ if $B<2t$.
\item[(c)] $\mathsf e(l')\le0.4079\,\theta^{l'/l}$ for every integer $l'\ge l$; $\delta(l',2l',2l')\le0.27263\,\theta^{l'/l}$; $D_2(N)\le4.1242\,\theta^{N/l}$ and $\delta(N^3)\le0.5892\,\theta^{N/l}$ for every integer $N\ge2l$.
\item[(d)] Over the square $(N,N)$, $N\ge2l$, the vortex free energy $\varepsilon_{\mathrm v}(N,N)=\log(\lambda_0/\|\mathsf T^{(1)}\|)$ satisfies $N\varepsilon_{\mathrm v}(N,N)<1.0292\,\theta^{N/2l}$.
\item[(e)] Conversely, on one lattice: if $\lambda_1/\lambda_0\le e^{-\mu}$ over the single cross-section $(2l,2l)$ and $3\le s<l$, then $\mathsf e(l)\le2e^{-\mu(l-s)}(e^{\delta(s,2l,2l)}-1)$; so $\mathsf e(l)\le x_0$ as soon as $\mu(l-s)\ge\delta(s,2l,2l)+\log(2/x_0)$.
\end{enumerate}
\end{proposition}

\begin{proof}
Freeze the side $A$: the boxes $(A,s,u)$ form a two-dimensional system, with the pressure $P_A=\inf F(A,s,u)/su$ and the excess $\delta_A(s,u)=F(A,s,u)-suP_A$, which has the properties of Proposition~\ref{C-361-7}(b), (c) in its two sides and bounds $\mathsf G(t;A,u)$ from above. The first two doublings of the proof of Proposition~\ref{C-361-8}, carried out in this system, give for its doubling defect $\eta=\mathsf e_A(m)=2F(A,m,2m)-F(A,2m,2m)$: if $\eta\le0.16521$, then $\mathsf e_A(2m)\le3.4522\,\eta^2\le0.5704\,\eta$ and $\delta_A(2m,2m)\le1.5514\,\eta^2$; the map $x\mapsto\psi(2x+\psi(x))$ has its least positive fixed point in $(0.2236,0.2237)$, the threshold of this recursion. It is started at the scale $A/2$: by Proposition~\ref{C-361-7}(b), (c) and Proposition~\ref{C-361-8}(b),
\[
\mathsf e_A(n)\le2\,\mathsf G(n;A,2n)\le2\,\delta(n,A,2n)\le\frac An\,\delta(n,2n,2n)\le0.185604\,\theta^{n/l}\le0.16521 .
\]
Iterating, $\mathsf e_A(2^jn)\le0.185604\,\theta^{2^jn/l}$ and $\delta_A(m,2m)\le0.11659\,\theta^{m/l}$ along $m=2^jn$; so for $2m\le u<4m$, $\mathsf r_m(A,u)\le(e^{0.23318}-1)\theta^{m/l}\le0.26261\,\theta^{m/l}<1$. At $m=n'$, $u=B$ this excludes a double top and gives $(\lambda_1/\lambda_0)^{n'}<\theta^{n'/l}$, and $\mathsf r_t\le\mathsf r_{n'}^{t/n'}$; the sector is that of Proposition~\ref{C-361-7}(d), as $\mathsf A=\mathsf K^{1/2}\mathsf m^{1/2}$ commutes with the flips and carries the eigenvectors of $\mathsf D$ to those of $\mathsf T'$. \emph{The constants of (a) and (b).} They come from a finer start and from finer constants of the same recursion. For $2n\le A<4n$, Proposition~\ref{C-361-7}(b) in the thickness and at the time $n$ over the cross-section $(2n,2n)$ give $\mathsf e_A(n)\le\frac A{2n}\bigl[\mathsf e(n)+\frac12\psi(\mathsf e(n))\bigr]\le2.037074\,\mathsf e(n)\le0.140559$, as $\mathsf e(n)\le0.069$.

On $[0,0.140559]$ the doubling of the slab, $\mathsf e_A(2m)\le\psi\bigl(2\mathsf e_A(m)+\psi(\mathsf e_A(m))\bigr)$, gives $\mathsf e_A(2m)\le3.1072\,\mathsf e_A(m)^2$, the inner argument staying below $0.292638<\log2$. So from a start $\mathsf e_A(n)\le x_1\le0.140559$, with $\vartheta=3.1072\,x_1<1$: $\mathsf e_A(2^jn)\le\vartheta^{2^j}/3.1072$; the series of the slab gives $\delta_A(n',2n')\le0.19314\,\vartheta^{n'/n}$; over $(B,A)$ with $2n'\le B<4n'$, $\mathsf r_{n'}\le e^{2\delta_A(n',2n')}-1\le0.4715\,\vartheta^{n'/n}<1$, so the top is simple, $\lambda_1/\lambda_0<\vartheta^{1/n}$, and $\mathsf r_t\le0.4715\,\vartheta^{t/n}$ for every $t\ge n'$. With $x_1=2.037074\,\mathsf e(n)$ and $\mathsf e(n)\le\theta^{n/l}/12.9$ (Proposition~\ref{C-361-8}(a)), $\vartheta\le6.3296\,\mathsf e(n)\le0.49067\,\theta^{n/l}$: this is the factor $0.49067^{1/n}$ in (a), and $0.4715\cdot0.49067\le0.2314$ is the number of (b).

At $A=2l$ the hypothesis itself is the start, $\mathsf e_{2l}(l)=\mathsf e(l)\le x_0\le0.069$, and on $[0,0.069]$ the same recursion has the constants $2.42$ and $0.5628$ in the places of $3.1072$ and $0.4715$: $\vartheta=2.42\,x_0$, $\lambda_1/\lambda_0<(2.42\,x_0)^{1/l}$ and $\mathsf r_t(2l,B)\le0.5628\,(2.42\,x_0)^{t/l}$ for $t\ge n'$. For the square, the box $(n,N,N)$ with the short side as the time gives $\mathsf G(n;N,N)\le\frac{N^2}{4n^2}\delta(n,2n,2n)\le0.185604\,\theta^{n/l}$ and $\mathsf r_n\le(e^{0.185604}-1)\theta^{n/l}$. The thermal law (b) follows in three zones: if $2t\le A$, over the cross-section $(2t,2t)$, with $m=2^il\le t<2m$, $\mathsf r_t(2t,2t)\le0.20395\,\theta^{t/l}$, so $\mathsf e(t)\le2\mathsf r_t\le0.4079\,\theta^{t/l}$ and $\delta(t,2t,2t)=\frac12[\mathsf e(t)+\delta((2t)^3)]\le0.27263\,\theta^{t/l}$, and the monotone density carries this to $(A,B)$; if $A<2t\le B$, the slab gives $\delta_A(t,2t)\le0.3695\,\theta^{t/l}$ and $\mathsf G(t;A,B)\le\frac B{2t}\delta_A(t,2t)$; if $B<2t$, $\mathsf r_t\le0.26261\,\theta^{t/l}$.

This also gives the first two clauses of (c); since $D_2(N)\le8\mathsf G(N;N,N)+4\mathsf G(N;N,2N)+2\mathsf G(N;2N,2N)$, the three zones and the series give the last two. (d) In the box $(N,N,h)$ with one twisted column of length $h$ and $\zeta=\frac12(1-r)$, the second inequality of step (ii) in the proof of Proposition~\ref{C-361-11} below, which is exact on one lattice, gives $N\varepsilon_{\mathrm v}(N,N)\le\mathsf d(N;N,N)-\log(1-2\zeta(N,N;2N))$, with $\zeta(N,N;2N)\le\delta(N,N,2N)\le2\delta(N^3)$ and $\mathsf d(N;N,N)\le2\mathsf G_N\le2\delta(N^3)$; with $\delta(N^3)<0.13736\,y$, $y=\theta^{N/2l}\le0.8901$, the first term is $<0.27472\,y$, and $-\log(1-0.54944\,y)\le0.75442\,y$ by convexity on $[0,0.8901]$; the sum is $<1.0292\,y$. (e) $\mathsf e(l)\le2\mathsf r_l(2l,2l)\le2e^{-\mu(l-s)}\mathsf r_s(2l,2l)$ and $\mathsf r_s=e^{\mathsf G(s)}-1\le e^{\delta(s,2l,2l)}-1$.
\end{proof}

\emph{The two seminorms on a larger torus.} By the setting and (M2), for $K\in\mathbb N$, $\beta>0$, $\ell''\le\ell$ in $\mathfrak L$ and a function $F$ of the bonds of a cell of side $\ell''/2$, with the seminorm $\|F\|^{(\ell'')}_{\mathrm{ch},\ell}$ of all its mirror copies on the torus $\ell$ under its own board and the seminorm $\|F\|_{\mathrm{ch},\ell}$ of its eight mirror copies under the board of side $\ell/2$ whose cell is centred at the cell of $F$, the cells of the two boards being nested (Lemma~\YMMref{104.1}{759}):
\[
\|F\|_{\mathrm{ch},\ell}\le\|F\|^{(\ell'')}_{\mathrm{ch},\ell}\le e^{(\ell''/2)^3(p_{\ell''}-p_\ell)}\,\|F\|_{\mathrm{ch},\ell''}\le e^{\delta((\ell''/a_K)^3)/8}\,\|F\|_{\mathrm{ch},\ell''},
\]
the last as $p_\ell\ge p_\infty=a_K^{-3}P$ and $(\ell''/2)^3(p_{\ell''}-p_\infty)=\frac18\delta((\ell''/a_K)^3)$.

\begin{proposition}[The excess bounds the distance from the torus]\label{C-361-9}
Let $K\in\mathbb N$, $\beta>0$, $\ell'\in\mathfrak L$, $\delta'=\delta((\ell'/a_K)^3)$ and $Q'=e^{\delta'/8}$; let $\mathfrak s$ be either $\mu_{K,\ell,\beta}$ with $\ell\in\mathfrak L$, $\ell\ge\ell'$, or a periodic Gibbs state at $\beta$.
\begin{enumerate}
\item[(a)] For every bounded function $F$ of the bonds of a cell of side $\ell'/2$, of any board: $|\mathfrak s(F)|\le Q'\|F\|_{\mathrm{ch},\ell'}\le Q'\,\mu_{K,\ell',\beta}(|F|^8)^{1/8}$.
\item[(b)] On the inner bonds of such a cell $\mathfrak s=h\,\mu_{K,\ell',\beta}$ with $\int h^{8/7}d\mu_{K,\ell',\beta}\le e^{\delta'/7}$, and the two laws are at a distance in total variation at most $\frac12(\sqrt\varepsilon+\varepsilon)$, $\varepsilon=22.3\,(e^{\delta'/7}-1)$.
\item[(c)] If $Q'\le1.11$, then for every real $G$ of the bonds of a cell of side $\ell''/2$, $\ell''\in\mathfrak L$, $\ell''\le\ell'$:
\[
\begin{multlined}
|\mathfrak s(G)-\mu_{K,\ell',\beta}(G)|\le4.84\,[Q'(Q'-1)]^{1/2}\,\|G\|_{\mathrm{ch},\ell'}\le4.84\,[Q'(Q'-1)]^{1/2}\,Q''\,\|G\|_{\mathrm{ch},\ell''}, \\
Q''=e^{\delta((\ell''/a_K)^3)/8};
\end{multlined}
\]
for complex $G$ the same with twice the constant.
\item[(d)] Let $\ell_i\in\mathfrak L$, $\ell_i\to\infty$, with $\limsup_i\delta((\ell_i/a_K)^3)\le D<\infty$ and $\mu_{K,\ell_i,\beta}\to\mu^*$ weakly. Then $\mu^*$ is a periodic Gibbs state; $\mu^*(A)\ge e^{-D}\mathfrak s(A)^8$ for every periodic Gibbs state $\mathfrak s$ and every Borel set $A$; $\mu^*\ge e^{-D}\nu$ for every ergodic periodic Gibbs state $\nu$; and $\mathfrak s\le e^D\mu^*$ for every periodic Gibbs state $\mathfrak s$. So $n$ mutually singular periodic Gibbs states at $\beta$ have $n\le e^D$, and indeed $n\le\liminf_{\mathfrak L\ni\ell'\to\infty}e^{\delta'}$.
\end{enumerate}
\end{proposition}

\begin{proof}
(a) By the fine-to-coarse form of (M2), $\|F\|^{(\ell')}_{\mathrm{ch},\ell}\le\|F\|_{\mathrm{ch},\ell'}\exp[(\ell'/2)^3(p_{\ell'}-p_\ell)]$, and since $p_\ell\ge p_\infty=a_K^{-3}P$ the exponent is at most $\frac18\ell'^3(p_{\ell'}-p_\infty)=\delta'/8$. For a torus $|\mu_{K,\ell,\beta}(F)|\le\|F\|^{(\ell')}_{\mathrm{ch},\ell}$, by the array form with the factor $1$ in the other cells; and a periodic Gibbs state is dominated by the tori, $|\mathfrak s(F)|\le\liminf_\ell\|F\|^{(\ell')}_{\mathrm{ch},\ell}$ (M2). The second inequality is Hölder's in the array form, the mirror copies of $|F|^8$ all having the same mean. (b) By (a), $\mathfrak s(A)\le Q'\mu_{K,\ell',\beta}(A)^{1/8}$ on the inner bonds, so $\mathfrak s=h\mu$ with $\|h\|_{8/7}\le Q'$ by duality. With $p=\frac87$ and $\phi(u)=u^p-1-p(u-1)\ge0.045\min\{(u-1)^2,|u-1|^p\}$ for $u\ge0$ (Taylor's formula on $[0,2]$, monotonicity of $\phi(1+v)/v^p$ beyond), $\int\phi(h)\le e^{\delta'/7}-1$ gives $\int|h-1|\le\sqrt\varepsilon+\varepsilon$. (c) We use an abstract form. Let $\mathcal A$ be a real vector space with a member $1$, $\mathsf N$ a seminorm and $\mathsf g$ a linear functional on it with
\[
\mathsf N(1+tG)^8\le1+8t\,\mathsf g(G)+R(|t|\mathsf N(G)),\quad R(x)=(1+x)^8-1-8x,\qquad|\mathsf g(G)|\le\mathsf N(G),
\]
and $\mathfrak s$ a linear functional with $\mathfrak s(1)=1$, $\mathfrak s(F)\le Q\,\mathsf N(F)$, $1\le Q\le1.11$. As $(1+u)^{1/8}\le1+\frac u8$: $t[\mathfrak s(G)-\mathsf g(G)]\le(Q-1)(1+x)+\frac Q8R(x)$, $x=|t|\mathsf N(G)$; if $\mathsf N(G)=0$, the bounds on $\mathsf g$ and the linearity bound on $\mathfrak s$, applied to $G$ and $-G$, give $\mathfrak s(G)=\mathsf g(G)=0$. For $Q=1$, use both signs of $t$ and let $t\to0$, since $R(x)=O(x^2)$. For $\mathsf N(G)>0$ and $Q>1$, take $t=\pm x/\mathsf N(G)$, $x^2=\frac{Q-1}{4.74Q}$ (so $x\le0.1446$, where $R(x)\le8\cdot4.7179\,x^2$); this gives $|\mathfrak s(G)-\mathsf g(G)|\le[2\sqrt{4.74}+(0.11/1.11)^{1/2}][Q(Q-1)]^{1/2}\mathsf N(G)\le4.6692\,[Q(Q-1)]^{1/2}\mathsf N(G)$. For $\mathsf N=\|\cdot\|_{\mathrm{ch},\ell'}$, $\mathsf g=\mu_{K,\ell',\beta}$ on the real functions of a cell, the product of the eight copies of $1+tG$ expands into $\sum_St^{|S|}m_S$ with $m_\emptyset=1$, $m_{\{\mathbf k\}}=\mu(G)$ for the eight one-point sets and $|m_S|\le\mathsf N(G)^{|S|}$ for $|S|\ge2$ (array form), and $Q=Q'$ by (a).

A function of a cell of side $\ell''/2$ is a function of the cell of side $\ell'/2$ of the setting, centred at its cell, as $\ell'/\ell''$ is a power of the odd $L$, and $\|G\|_{\mathrm{ch},\ell'}\le\|G\|^{(\ell'')}_{\mathrm{ch},\ell'}\le Q''\|G\|_{\mathrm{ch},\ell''}$ by the chain of (M2) recalled before the Proposition. (d) $\mu^*$ is a periodic Gibbs state invariant under $a_K\mathbb Z^3$ by (M3). For an event $A$ of the bonds of a cell of side $\ell_i/2$, (a) with $F=1_A$ gives $\mathfrak s(A)\le e^{\delta_i/8}\mu_{K,\ell_i,\beta}(A)^{1/8}$, $\delta_i=\delta((\ell_i/a_K)^3)$; the conditional laws of finitely many bonds given the others are continuous in the others, so the means of local events converge, and $\mu^*(A)\ge e^{-D}\mathfrak s(A)^8$ for local events, hence for all Borel sets by approximation.

An ergodic $\nu$ has an invariant carrier $N_\nu$ that carries no other ergodic state, so the ergodic decomposition of $\mu^*$ gives $\mu^*(\cdot\cap N_\nu)=\mu^*(N_\nu)\,\nu$ with $\mu^*(N_\nu)\ge e^{-D}\nu(N_\nu)^8=e^{-D}$; every $\mathfrak s$ is an integral of ergodic states (M3), each $\le e^D\mu^*$. Mutually singular states have disjoint carriers $C_j$, and $\sum_j\mu^*(C_j)\le1$ with $\mu^*(C_j)\ge e^{-D}$; for the last clause take sides along which $\delta'$ tends to its lower limit and a weakly convergent subsequence.
\end{proof}

\begin{lemma}[Boxes of every shape]\label{C-361-21}
Let $\beta\ge0$, $N=2h\ge4$, $c$ a corner, $\mathsf C=\mathsf C_h(c)$, and $\underline n$ any box with $n_\mu\ge N$; let $M_\mu=2^{j_\mu}h$, $j_\mu\in\mathbb N$, be the greatest number of this form that is $\le n_\mu$, so $N\le M_\mu\le n_\mu<2M_\mu$.
\begin{enumerate}
\item[(a)] Let $\mu$ be a direction, $\vartheta\ge2$ with $2\vartheta\le n_\mu$, $\theta'$ the reflection in a plane $x_\mu=c'_\mu$, $c'_\mu\in a(\mathbb Z+\frac12)$, and $F$ a bounded Borel function of the bonds with both ends in the $\vartheta$ layers below that plane. Then $\langle F\,\overline{\theta'F}\rangle_{\underline n}\ge0$ and $|\langle F\rangle_{\underline n}|^2\le\langle F\,\overline{\theta'F}\rangle_{\underline n}$, also for odd $n_\mu$. On the cube $N^3$, $|\langle\prod_\epsilon H_\epsilon^{[\epsilon]}\rangle_{N^3}|\le\prod_\epsilon\|H_\epsilon\|_{\mathrm{ch};N}$, so $\|\cdot\|_{\mathrm{ch};N}$ is a seminorm, and $|\langle\Psi\rangle_{N^3}|\le\|\Psi\|_{\mathrm{ch};N}\le\sup|\Psi|$.
\item[(b)] For every function $\Psi$ of the bonds of $\mathsf C$:
\[
|\langle\Psi\rangle_{\underline n}|\le\exp\Bigl\{\frac{\delta(N^3)}8\Bigl[1+3N\Bigl(\frac1{M_1}+\frac1{M_2}+\frac1{M_3}\Bigr)\Bigr]\Bigr\}\|\Psi\|_{\mathrm{ch};N}\le e^{\frac54\delta(N^3)}\|\Psi\|_{\mathrm{ch};N};
\]
and for the boxes $(q_1N,q_2N,q_3N)$, $q\in\mathbb N^3$: $|\langle\Psi\rangle_{\underline n}|\le\exp[\frac18(\delta(N^3)-\delta(\underline n)/q_1q_2q_3)]\|\Psi\|_{\mathrm{ch};N}$.
\item[(c)] Let $\Psi$ be real. If $\delta(N^3)\le\frac12$, then $|\langle\Psi\rangle_{\underline n}-\langle\Psi\rangle_{N^3}|\le7.31\,\delta(N^3)^{1/2}\|\Psi\|_{\mathrm{ch};N}$; for the boxes $(q_\mu N)$ and $\delta(N^3)\le0.8348$, the same with $1.86$ for $7.31$.
\end{enumerate}
None of these constants depends on the shape of the box.
\end{lemma}

\begin{proof}
(a) With the time $\mu$ and $\mathsf v=\frac12(n_\mu-2\vartheta)$, the trace formula (M1) and the cyclicity of the trace give $Z(\underline n)\langle F\rangle_{\underline n}=\operatorname{Tr}(\mathsf T'^{\,n_\mu/2}\cdot Y\mathsf T'^{\,\mathsf v})$, $Y=\mathfrak Y[\Phi_F]$; Schwarz's inequality in the Hilbert--Schmidt class, with $\|\mathsf T'^{\,n_\mu/2}\|_2^2=Z(\underline n)$ and $\|Y\mathsf T'^{\,\mathsf v}\|_2^2=\operatorname{Tr}(\mathsf T'^{\,2\mathsf v}Y^*Y)=Z(\underline n)\langle F\overline{\theta'F}\rangle_{\underline n}$, is the claim. Applied in the three directions in turn this is the array bound on the cube. (b) \emph{Doubling on a single box.} In the direction $1$ let $F_0=\Psi$ and $F_{i+1}=F_i\,\overline{\theta^{(i)}F_i}$, $\theta^{(i)}$ the reflection in $x_1=c_1+(2^i-1)ha$; by (a), as long as the doubled function fits, $|\langle F_i\rangle_{\underline n}|^2\le\langle F_{i+1}\rangle_{\underline n}$; so $|\langle\Psi\rangle_{\underline n}|\le\langle\mathfrak B_1\Psi\rangle_{\underline n}^{2^{-j_1}}$ with an array $\mathfrak B_1\Psi$ of $M_1/h$ cells.

The same in the directions $2$, $3$ gives an array $\mathfrak B\Psi$ on a block of $M_1\times M_2\times M_3$ sites with $\mathcal N=M_1M_2M_3/h^3$ cells and $|\langle\Psi\rangle_{\underline n}|\le\langle\mathfrak B\Psi\rangle_{\underline n}^{1/\mathcal N}$. \emph{The free layers cost their top eigenvalue.} Let $\underline m$ be a box with $m_\mu\ge M_\mu$ for each $\mu$; it is taken in turn as $\underline n$, $(M_1,n_2,n_3)$ and $(M_1,M_2,n_3)$. Take the direction $1$ as the time over the cross-section $(m_2,m_3)$, let $\Phi$ be the slab function of the cell of $\mathfrak B_3\mathfrak B_2\Psi$, which lies in the $h$ layers below $c_1$, and let $Y_\Phi=\mathfrak Y[\Phi]$ be its slab operator (another operator than the $Y$ of (a)).

In the direction $1$ the array $\mathfrak B\Psi$ is the row $\Phi,\Phi^\Theta,\Phi,\Phi^\Theta,\dots$ on $M_1/h$ consecutive slabs of $h$ layers, so its operators close up into $(Y_\Phi^*Y_\Phi)^{M_1/N}$, $M_1/N=2^{j_1-1}$, and the $m_1-M_1$ free layers carry $\mathsf T'^{\,m_1-M_1}$ (a free rest of one layer is $\mathfrak Y[\Phi^{+1}]=\mathsf T'\mathfrak Y[\Phi]$ of (M1)). By the trace formula (M1), $Z(\underline m)\langle\mathfrak B\Psi\rangle_{\underline m}=\operatorname{Tr}(\mathsf T'^{\,m_1-M_1}(Y_\Phi^*Y_\Phi)^{M_1/N})$ and $\operatorname{Tr}(Y_\Phi^*Y_\Phi)^{M_1/N}=Z(M_1,m_2,m_3)\langle\mathfrak B\Psi\rangle_{(M_1,m_2,m_3)}$; and $\operatorname{Tr}(\mathsf A\mathsf B)=\operatorname{Tr}(\mathsf B^{1/2}\mathsf A\mathsf B^{1/2})\le\|\mathsf A\|\operatorname{Tr}\mathsf B$ for $\mathsf A,\mathsf B\ge0$, $\mathsf B$ of trace class.

So, direction by direction,
\[
Z(\underline n)\langle\mathfrak B\Psi\rangle_{\underline n}\le\lambda_0(n_2,n_3)^{n_1-M_1}\lambda_0(M_1,n_3)^{n_2-M_2}\lambda_0(M_1,M_2)^{n_3-M_3}\,Z(M_1,M_2,M_3)\langle\mathfrak B\Psi\rangle_{(M_1,M_2,M_3)}.
\]
On the box $(M_1,M_2,M_3)$, whose sides lie in $N\mathbb N$, the inequality $\operatorname{Tr}(Y_\Phi^*Y_\Phi)^p\le(\operatorname{Tr}Y_\Phi^*Y_\Phi)^p$ at $p=M_\mu/N$ in each direction $\mu$, with $\Phi$ the cell of the array in that direction, reduces the array to the cube $N^3$, where it is the product of the eight mirror copies of $\Psi$: $Z(M_1,M_2,M_3)\langle\mathfrak B\Psi\rangle_{(M_1,M_2,M_3)}\le[Z(N^3)\|\Psi\|^8_{\mathrm{ch};N}]^{M_1M_2M_3/N^3}$. \emph{The count.} With $\log\lambda_0(\Sigma)=|\Sigma|P+\varepsilon(\Sigma)$ and $F=|\cdot|P+\delta$ the terms with $P$ cancel, and by the area bound $\varepsilon(\mathsf A,\mathsf B)\le(\mathsf A+\mathsf B)\delta(N^3)/N^2$ for $\mathsf A,\mathsf B\ge N$ (Proposition~\ref{C-361-7}(c)) the free layers cost at most $3\delta(N^3)N^{-2}(M_1M_2+M_1M_3+M_2M_3)$ against $\mathcal N=8M_1M_2M_3/N^3$ cells: an area against a volume.

\emph{The boxes $(q_1N,q_2N,q_3N)$, $q\in\mathbb N^3$,} are treated with another array, which fills the box for every $q$, so that no layer is free. For a function $F$ of the bonds in the $h$ layers below $c_\mu$, with slab function $\Phi_F$ and slab operator $Y_F=\mathfrak Y[\Phi_F]$, the time being $\mu$, let $\mathfrak A_\mu F$ be the product of the $2q_\mu$ cells $\Phi_F,\Phi_F^\Theta,\Phi_F,\Phi_F^\Theta,\dots$ on all the $2q_\mu$ consecutive slabs of $h$ layers: $F$ and its successive conjugated mirror images in the planes $x_\mu=c_\mu+jha$. By (M1), $Z(\underline n)\langle\mathfrak A_\mu F\rangle_{\underline n}=\operatorname{Tr}(Y_F^*Y_F)^{q_\mu}=\|Y_F\|_{2q_\mu}^{2q_\mu}$ and $Z(\underline n)=\operatorname{Tr}\mathsf T'^{\,2q_\mu h}=\|\mathsf T'^{\,h}\|_{2q_\mu}^{2q_\mu}$, the norms of the Schatten class of order $2q_\mu$; Hölder's inequality in these classes for the $2q_\mu$ factors $Y_F,\mathsf T'^{\,h},\dots,\mathsf T'^{\,h}$ gives $Z(\underline n)|\langle F\rangle_{\underline n}|=|\operatorname{Tr}(\mathsf T'^{\,h}\cdots\mathsf T'^{\,h}Y_F)|\le\|Y_F\|_{2q_\mu}\|\mathsf T'^{\,h}\|_{2q_\mu}^{2q_\mu-1}$, that is $|\langle F\rangle_{\underline n}|\le\langle\mathfrak A_\mu F\rangle_{\underline n}^{1/2q_\mu}$.

Applied on the one box $\underline n$ to $\Psi$ in the direction $1$, to $\mathfrak A_1\Psi$, which lies below $c_2$, in the direction $2$, and to $\mathfrak A_2\mathfrak A_1\Psi$, which lies below $c_3$, in the direction $3$, this gives $|\langle\Psi\rangle_{\underline n}|\le\langle\mathfrak A\Psi\rangle_{\underline n}^{1/8q_1q_2q_3}$, $\mathfrak A=\mathfrak A_3\mathfrak A_2\mathfrak A_1$, the product of the mirror copies of $\Psi$ in all the $8q_1q_2q_3$ cells of the box. Only then is the box made smaller: with the time $1$ and $F=\mathfrak A_3\mathfrak A_2\Psi$, $Z(q_1N,n_2,n_3)\langle\mathfrak A\Psi\rangle_{(q_1N,n_2,n_3)}=\operatorname{Tr}(Y_F^*Y_F)^{q_1}\le(\operatorname{Tr}Y_F^*Y_F)^{q_1}=[Z(N,n_2,n_3)\langle\mathfrak A\Psi\rangle_{(N,n_2,n_3)}]^{q_1}$, and the same in the directions $2$, $3$ give $Z(\underline n)\langle\mathfrak A\Psi\rangle_{\underline n}\le[Z(N^3)\|\Psi\|^8_{\mathrm{ch};N}]^{q_1q_2q_3}$.

As $|\underline n|=q_1q_2q_3N^3$, the terms with $P$ cancel, and $|\langle\Psi\rangle_{\underline n}|\le\exp\frac18[\delta(N^3)-\delta(\underline n)/q_1q_2q_3]\,\|\Psi\|_{\mathrm{ch};N}$. (c) The abstract form of the proof of Proposition~\ref{C-361-9}(c) with $\mathsf N=\|\cdot\|_{\mathrm{ch};N}$, $\mathsf g=\langle\cdot\rangle_{N^3}$, $\mathfrak s=\langle\cdot\rangle_{\underline n}$: its first condition is the array bound of (a) on the cube, the eight one-point terms being $\langle G\rangle_{N^3}$ by the mirror symmetry of the cube, and $\mathfrak s(F)\le Q\,\mathsf N(F)$ is (b). With $Q=e^{\frac54\delta(N^3)}$, which is $\le1.11$ if and only if $\delta(N^3)\le\frac45\log1.11$: as $u\mapsto e^u(e^u-1)/u$ increases, $Q(Q-1)\le\frac{1.11\cdot0.11}{\log1.11}\cdot\frac54\delta(N^3)$, so $[Q(Q-1)]^{1/2}\le1.209334\,\delta(N^3)^{1/2}$ and $4.84\cdot1.209334\le5.86\le7.31$; for $\frac45\log1.11<\delta(N^3)\le\frac12$, by (b) and (a) directly, $|\langle\Psi\rangle_{\underline n}-\langle\Psi\rangle_{N^3}|\le(e^{\frac54\delta}+1)\|\Psi\|_{\mathrm{ch};N}$, and $(e^{\frac54\delta}+1)\delta^{-1/2}$ decreases on $(0,\frac12]$, so it is at most its value at $\frac45\log1.11$, which is less than $7.3025$.

For the boxes $(q_\mu N)$ take $Q=e^{\delta(N^3)/8}$ by the second clause of (b): $Q\le1.11$ if and only if $\delta(N^3)\le8\log1.11$, and then $Q(Q-1)\le\frac{1.11\cdot0.11}{\log1.11}\cdot\frac{\delta(N^3)}8$, $[Q(Q-1)]^{1/2}\le0.382425\,\delta(N^3)^{1/2}$, and $4.84\cdot0.382425\le1.85094\le1.86$.
\end{proof}

\begin{lemma}[Cells are nearly independent under every boundary condition]\label{C-361-22}
\begin{enumerate}
\item[(a)] Let $\mathcal K$ be a finite set, $Y_\kappa$ measurable maps of a probability space $(E,R)$ to spaces $E_\kappa$ with probability measures $\mathrm{ref}_\kappa$, and $\Delta\ge0$ with $\log R(\exp\sum_\kappa\Phi_\kappa\circ Y_\kappa)\le\sum_\kappa[\Delta+\frac18\log\mathrm{ref}_\kappa(e^{8\Phi_\kappa})]$ for all bounded real $\Phi_\kappa$. Then for every probability $P=fR$ with $\mathsf H=R(f\log f)<\infty$ and all bounded real non-constant $G_\kappa$: $\sum_\kappa[P(G_\kappa\circ Y_\kappa)-\mathrm{ref}_\kappa(G_\kappa)]^2/\operatorname{osc}(G_\kappa)^2\le4\mathsf H+4\#\mathcal K\,\Delta$.
\item[(b)] Let $K\in\mathbb N$, $\beta>0$, $\ell'\le\ell$ in $\mathfrak L$, $n=\ell/a_K$, $n'=\ell'/a_K$ and $\delta'=\delta(n'^3)$; these letters serve in (c) as well. The $(2n/n')^3=8(n/n')^3$ cells of a period box of a board of side $\ell'/2$ on the torus $\ell$ satisfy the hypothesis of (a) with $R=\mu_{K,\ell,\beta}$, $\mathrm{ref}$ the law of a cell under $\mu_{K,\ell',\beta}$ and $\Delta=\delta'/8$.
\item[(c)] Let $D$ be the set of the sites of a period cube of the torus $\ell$, $n$ sites a side, that is the union of the $8(n/n')^3$ cells of side $\ell'/2$ of a period box of the board, let $W$ be the set of the positively oriented bonds with initial point in $D$, which holds the bonds of these cells, and let $\gamma_W(\cdot)(\xi)$ be the law of the bonds of $W$ given the configuration $\xi$ outside $W$. For every $\xi$ and all real non-constant $G_{\mathbf k}$ of the cells:
\[
\sum_{\mathbf k}\frac{[(\gamma_WG_{\mathbf k})(\xi)-\mu_{K,\ell',\beta}(G_{\mathbf k})]^2}{\operatorname{osc}(G_{\mathbf k})^2}\le144\,\beta n^2+4\,(n/n')^3\delta';
\]
the mean over the cells is at most $\varepsilon_1(n,n')^2=18\beta n'^3/n+\delta'/2$.
\item[(d)] Let $\lambda\ge3$ be odd and $\underline n=(\mathsf n_1\lambda,\mathsf n_2\lambda,\mathsf n_3\lambda)$ a periodic box with even $\mathsf n_1,\mathsf n_2,\mathsf n_3$, cut into the $\mathsf n_1\mathsf n_2\mathsf n_3$ cells of $\lambda$ sites a side of a board whose walls are planes of lattice reflections. The cells satisfy the hypothesis of (a) with $R=\mu_{\underline n}$, $\mathrm{ref}$ the law of one cell of the cube $(2\lambda)^3$ under $\mu_{(2\lambda)^3}$, and $\Delta=\delta((2\lambda)^3)/8$.
\end{enumerate}
\end{lemma}

\begin{proof}
(a) By the variational formula for the relative entropy, $P(\Phi)\le\mathsf H+\log R(e^\Phi)$ for $\Phi=\sum_\kappa t_\kappa(G_\kappa-\mathrm{ref}_\kappa G_\kappa)\circ Y_\kappa$; the hypothesis and Hoeffding's bound $\log\mathrm{ref}(e^{s(G-\mathrm{ref}G)})\le s^2\operatorname{osc}(G)^2/8$ give $\sum_\kappa t_\kappa[P(G_\kappa)-\mathrm{ref}_\kappa(G_\kappa)]\le\mathsf H+\#\mathcal K\Delta+\sum_\kappa t_\kappa^2\operatorname{osc}(G_\kappa)^2$; optimise over the $t_\kappa$. (b) This is the proof of Proposition~\ref{C-361-9}(a) for a product: for real $\Phi_{\mathbf k}$ of the cells, the array form of the chessboard estimate on the torus $\ell$ for the board of side $\ell'/2$ gives $\mu_{K,\ell,\beta}(\prod_{\mathbf k}e^{\Phi_{\mathbf k}})\le\prod_{\mathbf k}\|e^{\Phi_{\mathbf k}}\|^{(\ell')}_{\mathrm{ch},\ell}$, the fine-to-coarse form gives the factor $e^{\delta'/8}$ per cell, $\|e^{\Phi}\|^{(\ell')}_{\mathrm{ch},\ell}\le e^{\delta'/8}\|e^{\Phi}\|_{\mathrm{ch},\ell'}$, and Hölder's inequality on the torus $\ell'$ gives $\|e^{\Phi}\|_{\mathrm{ch},\ell'}\le\mu_{K,\ell',\beta}(e^{8\Phi})^{1/8}$, the eight mirror copies of $e^{8\Phi}$ having the mean of one cell. (c) Identify the bonds of $W$ with those of the torus $\ell$.

The conditional law of the bonds of $W$ and the law of the torus have the same weight but on the plaquettes that are not common to both, the $12n^2-3n$ plaquettes with a bond in $W$ that meet the outside and the $6n^2-3n$ plaquettes of the torus that wrap around it, a surface; as $a_p\in[0,2]$, the density of the first law against the second lies between $e^{-36\beta n^2}$ and $e^{36\beta n^2}$ (Proposition~\YMMref{91.8}{704}), so their relative entropy is $\mathsf H\le36\beta n^2$. Then (a), with (b) on the torus $\ell$ and $\#\mathcal K=8(n/n')^3$: the sum is at most $4\cdot36\beta n^2+4\cdot8(n/n')^3\delta'/8$, and the mean over the cells is at most $18\beta n'^3/n+\delta'/2$. (d) The reflection in a wall plane is a symmetry of the box, as every $\mathsf n_\mu$ is even, and it is reflection positive by the trace formula (Lemma~\ref{C-361-21}(a)).

So the chessboard estimate on the box holds (Proposition~\YMMref{86.3}{681}): the mean of a product over the cells is at most the product of the seminorms of the factors on the box (equation~\YMMref{(1756)}{681}), the $(\mathsf n_1\mathsf n_2\mathsf n_3)$-th roots of the means of the products of all their mirror copies. By the subadditivity in the volume (Proposition~\YMMref{86.5}{682}) at $\underline{\mathsf n}'=(2,2,2)$, which is the trace formula with the inequality $\operatorname{Tr}(Y^*Y)^q\le(\operatorname{Tr}Y^*Y)^q$ in each direction as in the proof of Lemma~\ref{C-361-21}(b), the seminorm on the box of a function $G$ of one cell is at most $\|G\|_{\mathrm{ch};2\lambda}$ times $\exp[\frac18F((2\lambda)^3)-\lambda^3F(\underline n)/|\underline n|]\le e^{\delta((2\lambda)^3)/8}$, as $F(\underline n)/|\underline n|\ge P$.

For $G=e^\Phi$, Hölder's inequality on the cube gives $\|e^\Phi\|_{\mathrm{ch};2\lambda}\le\langle e^{8\Phi}\rangle_{(2\lambda)^3}^{1/8}$, the eight mirror copies of $e^{8\Phi}$ having the mean of one cell; this is the hypothesis of (a) with the reference law of one cell of the cube $(2\lambda)^3$.
\end{proof}

\begin{lemma}[The operator of a cell; two cells cluster]\label{C-361-20}
Let $K\in\mathbb N$, $\beta>0$, a board of side $\mathsf l=\ell''/2$, $\ell''\in\mathfrak L$, and $\ell\in\mathfrak L$, $\ell\ge2\mathsf l$, $N=\ell/a_K=\mathsf n\mathsf w$ with $\mathsf w=\mathsf l/a_K$ layers to a cell; a direction as the time of the cube $N^3$, $\Sigma=(N,N)$.
\begin{enumerate}
\item[(a)] For $F$ of the bonds of one cell, $\Phi_F$ its slab function: $\|\mathfrak Y[\Phi_F]\|\le Z(N^3)^{1/\mathsf n}\|F\|^{(\ell'')}_{\mathrm{ch},\ell}=\lambda_0^{\mathsf w}(1+\mathsf r_N)^{1/\mathsf n}\|F\|^{(\ell'')}_{\mathrm{ch},\ell}$, with the seminorm of the setting of all the mirror copies of $F$ on the torus $\ell$ under the board of side $\mathsf l$.
\item[(b)] Let $\mathsf e(l)\le x_0$ at $\beta$, $\nu$ the periodic Gibbs state of Proposition~\ref{C-361-10}, $\ell''\ge2la_K$, $Q''=e^{\delta((\ell''/a_K)^3)/8}$. For $F,G$ of the bonds of two cells whose indices differ by $j\ge1$ in a direction: $|\nu(FG)-\nu(F)\nu(G)|\le\theta^{(j-1)\mathsf w/l}Q''^2\|F\|_{\mathrm{ch},\ell''}\|G\|_{\mathrm{ch},\ell''}$; on every torus $\ell\ge L\ell''$ of $\mathfrak L$, with $j_-=\min\{j,\mathsf n-j\}\ge2$ the distance on the torus, the same with $3.44\,\theta^{(j_--1)\mathsf w/l}$.
\end{enumerate}
The constants do not depend on $K$, $\beta$, $\ell$, $F$, $G$ or the offset.
\end{lemma}

\begin{proof}
(a) The operator norm is at most the Schatten norm of order $\mathsf n$ ($\mathsf n=2\ell/\ell''$ is even), and by the trace formula (M1) $\operatorname{Tr}(\mathfrak Y^*\mathfrak Y)^{\mathsf n/2}$ is $Z(N^3)$ times the mean of the product of the $\mathsf n$ mirror copies of $F$ along the time, cells of the board of side $\mathsf l$ in one column, which the array form of (M2) bounds by $(\|F\|^{(\ell'')}_{\mathrm{ch},\ell})^{\mathsf n}$. (b) \emph{The state $\nu$.} A function of the second cell is the reflected image of a function of the slab of the first, so the spectral formula of Proposition~\ref{C-361-10}(f) applies; over the cross-section $(N',N')$, $N'=\ell'/a_K$, $\ell'\ge\ell''$ in $\mathfrak L$, the covariance is $\lambda_0^{-2\mathsf w}\sum_{i\ge1}(\lambda_i/\lambda_0)^n\overline{\mathsf a_i}\mathsf b_i$ with $n=(j-1)\mathsf w$ layers between the cells (at $n=0$ with the term of the kernel of $\mathsf T'$ besides), and Schwarz's inequality bounds it by $(\lambda_1/\lambda_0)^n\lambda_0^{-2\mathsf w}\|\mathfrak Y[\Phi_F]\|\|\mathfrak Y[\Phi_G]\|$.

The gap of Proposition~\ref{C-361-19}(a) gives $(\lambda_1/\lambda_0)^n<\theta^{n/l}$, and (a) on the cube $N'$ with the fine-to-coarse form of (M2), $\|F\|^{(\ell'')}_{\mathrm{ch},\ell'}\le Q''\|F\|_{\mathrm{ch},\ell''}$ as $p_{\ell'}\ge p_\infty$, gives $\lambda_0^{-\mathsf w}\|\mathfrak Y[\Phi_F]\|\le(1+\mathsf r_{N'})^{1/\mathsf n'}Q''\|F\|_{\mathrm{ch},\ell''}$, $\mathsf n'=N'/\mathsf w$, the same for $G$; as $\ell'\to\infty$ the first factor tends to $1$, and the formula passes to $\nu$ (Proposition~\ref{C-361-10}(c), (f)). \emph{The torus $\ell\ge L\ell''$.} The cube $N=\mathsf n\mathsf w$ has $\mathsf n=2\ell/\ell''\ge2L\ge6$ cells along the time, an even number. The two cells, $j_-\ge2$ apart the shorter way, leave $n=(j_--1)\mathsf w$ layers between them on one side and $E=(\mathsf n-j_--1)\mathsf w\ge n$ on the other, and $n,E\ge\mathsf w\ge l$, as $\ell''\ge2la_K$.

Write, as in the proof of Lemma~\ref{C-361-29}, $\mathsf T'^m=\lambda_0^m(|e_0\rangle\langle e_0|+\mathsf R_{(m)})$, $e_0$ the top eigenvector, so that $\|\mathsf R_{(m)}\|=\varrho^m$, $\varrho=\lambda_1/\lambda_0$, and $\operatorname{Tr}|\mathsf R_{(m)}|=\mathsf r_m$, and put $y_C=\lambda_0^{-\mathsf w}\mathfrak Y[\Phi_C]$. The trace formula (M1) with the four slabs $F$, $1_n$, $G$, $1_E$ gives $\mu_{K,\ell,\beta}(FG)=\operatorname{Tr}\bigl((|e_0\rangle\langle e_0|+\mathsf R_{(E)})\,y_G\,(|e_0\rangle\langle e_0|+\mathsf R_{(n)})\,y_F\bigr)/(1+\mathsf r_N)$, and the mean of one cell is $\mu_{K,\ell,\beta}(C)=[\langle e_0,y_Ce_0\rangle+\operatorname{Tr}(\mathsf R_{(N-\mathsf w)}y_C)]/(1+\mathsf r_N)$.

Expanding, the products $\langle e_0,y_Fe_0\rangle\langle e_0,y_Ge_0\rangle$ cancel up to the factors $(1+\mathsf r_N)^{-1}$ and $(1+\mathsf r_N)^{-2}$, and $|\operatorname{Tr}(XW)|\le\|X\|\operatorname{Tr}|W|$ bounds the rest:
\[
|\mu_{K,\ell,\beta}(FG)-\mu_{K,\ell,\beta}(F)\mu_{K,\ell,\beta}(G)|\le\|y_F\|\,\|y_G\|\,\bigl[\mathsf r_N+\varrho^n+\varrho^E+\varrho^n\mathsf r_E+2\mathsf r_{N-\mathsf w}+\mathsf r_{N-\mathsf w}^2\bigr],
\]
the fourth term as the trace norm of $\mathsf R_{(E)}$ times the operator norm of $y_G\mathsf R_{(n)}y_F$. By (a) and the fine-to-coarse form, $\|y_C\|\le(1+\mathsf r_N)^{1/\mathsf n}Q''\|C\|_{\mathrm{ch},\ell''}$. Proposition~\ref{C-361-19} bounds $\varrho^n,\varrho^E\le\theta^{n/l}$ by (a), and $\mathsf r_N$, $\mathsf r_{N-\mathsf w}$, $\mathsf r_E$ by its thermal law (b) over the cross-section $(N,N)$: in the zone $2t>N$ by $0.26261\,\theta^{t/l}$, which serves for $t=N$, $N-\mathsf w$ and for $t=E$ if $2E>N$, and in the zone $2t\le N$ by $\mathsf r_t\le e^{0.27263(N/2t)^2\theta^{t/l}}-1$. Every term of the bracket, divided by $\theta^{n/l}$, and the factor $(1+\mathsf r_N)^{2/\mathsf n}$ increase with $u=\theta^{\mathsf w/l}\le\theta\le0.8901$; at $u=0.8901$ the product of the bracket over $\theta^{n/l}$ and $(1+\mathsf r_N)^{2/\mathsf n}$ is at most $3.33$ for all even $\mathsf n\ge6$ and $2\le j_-\le\mathsf n/2$, the greatest value $3.3214$ being taken at $\mathsf n=6$, $j_-=3$, where $n=E=2\mathsf w$ and the bracket is at most $u^2[2+0.6258+0.26261(u^4+2u^3+0.26261\,u^8)]$. This is the torus clause, with $3.33\le3.44$.
\end{proof}

\begin{proposition}[One lattice under the inequality]\label{C-361-10}
Let $K\in\mathbb N$, $\beta>0$, and $\mathsf e(l)\le x_0$ at $\beta$ for one integer $l\ge3$.
\begin{enumerate}
\item[(a)] The lattice $a_K\mathbb Z^3$ has exactly one periodic Gibbs state $\nu$ at $\beta$; it is ergodic under $a_K\mathbb Z^3$, invariant under all isometries of the lattice and reflection positive at every lattice mirror; there is an open interval $U\ni\beta$ on which $\mathsf e(2l)\le0.069$, $P$ is continuously differentiable on $U$, and $P'(\beta')=-3\nu_{\beta'}(a_p)$.
\item[(b)] For every bounded real Borel $F$ of the bonds of the cell $[-\frac{\ell'}4,\frac{\ell'}4]^3$, $\ell'\in\mathfrak L$, $n'=\ell'/a_K\ge2l$, every $n=\ell/a_K$ with $\ell\ge\ell'$ in $\mathfrak L$, and the cube $Q_M=\{|x|_\infty\le Ma_K\}$, $M\ge n$, of $2M+1$ sites a side: $\|\nu(F\,|\,\text{bonds outside }Q_M)-\nu(F)\|_{L^2(\nu)}\le\operatorname{osc}(F)\,\varepsilon_1(n,n')$, $\varepsilon_1(n,n')^2<18\beta n'^3/n+0.06868\,\theta^{n'/2l}$. So $\nu$ is trivial on the tail $\sigma$-field and an extreme point of the set of all Gibbs states at $\beta$, and for $\nu$-almost every boundary datum the laws in growing cubes converge weakly to $\nu$. The generator $\mathsf H_\nu\ge0$ of the reversible lattice Langevin dynamics of $\nu$, the operator of the closure of the form $\sum_b\nu(\nabla_bF\cdot\nabla_bF')$ on smooth local functions, has $\ker\mathsf H_\nu=\mathbb C1$, and $e^{-s\mathsf H_\nu}F\to\nu(F)$ in $L^2(\nu)$ for every $F$.
\item[(b$'$)] With a coordinate direction as the time, $\lambda=2h+1\ge\max\{3,l\}$, $t\ge\max\{2,l\}$ and $n=h+t$: for every bounded real Borel $F$ of the bonds of $\{|x|_\infty\le ha_K\}$, $\|\nu(F\,|\,\text{bonds outside }\{|x_3|<na_K\})-\nu(F)\|_{L^2(\nu)}\le\operatorname{osc}(F)\,\varepsilon_2$, $\varepsilon_2^2=0.27472\,(\lambda/t)^2\theta^{t/l}+0.06868\,\theta^{\lambda/l}$; with $h=\lfloor n/3\rfloor$, for $n\ge\frac12(3l+1)$, this is $\le0.7057\operatorname{osc}(F)\,\theta^{(2n-1)/6l}$.
\item[(c)] For every even $N\ge2l$, every corner $c$, every bounded real Borel $\Psi$ of the bonds of $\mathsf C_{N/2}(c)$ and every periodic box $\underline n$ with all sides $\ge N$: $|\langle\Psi\rangle_{\underline n}-\langle\Psi\rangle_{N^3}|\le2.71\,\theta^{N/4l}\|\Psi\|_{\mathrm{ch};N}$ and $|\langle\Psi\rangle_{\underline n}-\nu(\Psi)|\le5.42\,\theta^{N/4l}\|\Psi\|_{\mathrm{ch};N}\le5.42\,\theta^{N/4l}\sup|\Psi|$; for events $A$ of the cell, $|\mu_{\underline n}(A)-\nu(A)|\le2.71\,\theta^{N/4l}$. As all sides tend to $\infty$ the measures of the periodic boxes converge weakly to $\nu$, uniformly over all boxes with three sides $\ge N_0$.
\item[(d)] \emph{Twisted boxes.} Let $\Gamma_n$ be period lattices and $\mathsf T_n$ closed twist sets ($\Gamma_n$-periodic sets of plaquettes such that every elementary cube has an even number of faces in $\mathsf T_n$; the sign of $\operatorname{tr}\partial U$ is reversed on them) that recede: the shortest period of $\Gamma_n$ tends to $\infty$ and every finite set of plaquettes misses $\mathsf T_n$ at last. Then the laws $\mu^{\mathsf T_n}_{\Gamma_n,\beta}$, and all mixtures of a bounded number of such sequences with any weights (the sum over the eight sectors among them), converge weakly to $\nu$; and if $\mathsf T_n'=\mathsf T_n\triangle d\mathsf S_n$ with any $\Gamma_n$-periodic set $\mathsf S_n$ of bonds, $d\mathsf S$ the plaquettes with an odd number of bonds in $\mathsf S$, the joint cumulants of the twisted fields $\mathcal O^{\mathsf T_n'}_{a}(f)=a^{-1}\sum_pf(x(p))a^{\mathsf T_n'}_p$, $a^{\mathsf T}_p=1+\frac12\operatorname{tr}\partial U(p)$ on $\mathsf T$ and $a_p$ elsewhere, under $\mu^{\mathsf T_n'}_{\Gamma_n,\beta}$ converge to those of the $\mathcal O_a(f)$ under $\nu$. The word \emph{recede} cannot be dropped: $P(\beta)>-3\beta$, so $\nu(\operatorname{tr}\partial U(p))>0$, and a column through a fixed plaquette $p_0$ gives in the limit the law of $-\operatorname{tr}\partial U(p_0)$ under $\nu$, which differs from that of $\operatorname{tr}\partial U(p_0)$.
\item[(e)] Over every cross-section $(A,B)$ with $A,B\ge2l$ the top is simple, lies in the range of $\mathsf P_{00}$, and $\lambda_1/\lambda_0<\theta^{1/l}$ (Proposition~\ref{C-361-19}(a)).
\item[(f)] Let $\mathrm{Sl}$ be a slab of $\vartheta\ge2$ consecutive layers, $\Theta_n$ the lattice reflection such that exactly $n$ layers lie between $\mathrm{Sl}$ and $\Theta_n\mathrm{Sl}$, $A,B\in L^2(\nu)$ functions of the bonds of $\mathrm{Sl}$, and $h_{AB}(n)=\nu(\overline{\Theta_nA}\,B)$. For all integers $n\ge n_0\ge0$,
\[
\bigl|h_{AB}(n)-\overline{\nu(A)}\,\nu(B)\bigr|\le\theta^{(n-n_0)/l}\prod_{C=A,B}\bigl[h_{CC}(n_0)-|\nu(C)|^2\bigr]^{1/2},
\]
and $n\mapsto h_{CC}(n)-|\nu(C)|^2$ is nonnegative and log-convex; at $n_0=0$, for bounded $A,B$: $|\operatorname{Cov}_\nu|\le\theta^{n/l}\sup|A|\sup|B|$.
\end{enumerate}
\end{proposition}

\begin{proof}
(a) A periodic Gibbs state exists as a limit of tori (M3). If $\rho$ and $\nu$ are two of them and $A$ is a local event inside the central cell of side $\ell'/2$, Proposition~\ref{C-361-9}(b) gives $|\rho(A)-\nu(A)|\le\sqrt\varepsilon+\varepsilon$, $\varepsilon=22.3(e^{\delta'/7}-1)$, and $\delta'\to0$ as $\ell'\to\infty$ by Proposition~\ref{C-361-8}(b); so $\rho=\nu$. The symmetries and reflection positivity come from (M3); almost every ergodic component of $\nu$ is a periodic Gibbs state, hence $\nu$. By Proposition~\ref{C-361-8}(a), $\mathsf e(2l)\le12.9\cdot0.069^2<0.069$ at $\beta$, hence on an open interval $U$, the box functions being smooth in $\beta$; there the derivatives $-3\mu_{K,\ell_m,\beta'}(a_p)$ of $(\ell_m/a_K)^{-3}\log Z_{K,\ell_m}$, $\ell_m=2L^m$, converge uniformly to $-3\nu_{\beta'}(a_p)$ by Proposition~\ref{C-361-9}(b) and the bound of Proposition~\ref{C-361-8}(b), uniform on $U$. (b) Let $D$ be the cube of $n$ sites a side that is the union of the $8(n/n')^3$ cells $\square_{\mathbf k}=[-\frac{\ell'}4,\frac{\ell'}4]^3+\frac{\ell'}2\mathbf k$, $0\le\mathbf k_\mu<2n/n'$, $W$ the set of the positively oriented bonds with initial point in $D$, as in Lemma~\ref{C-361-22}(c), and $F_{\mathbf k}$ the translate of $F$ into $\square_{\mathbf k}$.

Lemma~\ref{C-361-22}(c), integrated against $\nu$, bounds the mean over $\mathbf k$ of $\|\nu(F_{\mathbf k}|\text{outside }W)-\mu_{K,\ell',\beta}(F)\|^2$ by $\operatorname{osc}(F)^2\varepsilon_1^2$, and by translation invariance the term of $\mathbf k$ is $\|\nu(F|\text{outside }W-\frac{\ell'}2\mathbf k)-\mu_{K,\ell',\beta}(F)\|^2$. Each translate $W-\frac{\ell'}2\mathbf k$ holds the cell and has its bonds in $[-\ell,\ell]^3$, so it lies in $Q_n\subset Q_M$; conditioning on the smaller field outside $Q_M$ is a contraction in $L^2(\nu)$, so every term, and with them the mean, is at least $\|\nu(F|\text{outside }Q_M)-\mu_{K,\ell',\beta}(F)\|^2$; $\delta'<0.13736\,\theta^{n'/2l}$; and the mean $\nu(F)$ is the best constant in $L^2(\nu)$.

Let $n\to\infty$, then $n'\to\infty$: tail triviality, which for a Gibbs state is extremality. A function in $\ker\mathsf H_\nu$ has $\nabla_bF=0$ for every bond, so (the conditional law of one bond being equivalent to Haar measure) it is measurable outside every finite set of bonds, hence constant; the rest is the spectral theorem. (b$'$) Over the cross-section the top eigenvector $e_0$ is positive and continuous; in the cylinder the joint density of the two layers $V^{\pm n}$ against $\pi\otimes\pi$, $\pi=e_0^2dV$, is $p_{2n}=\mathsf D^{2n}/\lambda_0^{2n}e_0\otimes e_0$, and given them the slab of the cells has the density $\varrho=p_t(x,y)p_t(y',x')/p_{2n}(x,x')$ against its own law; the mean relative entropy is $2\iint p_t\log p_t-\iint p_{2n}\log p_{2n}\le2\log\iint p_t^2=2\log(1+\mathsf r_{2t})$.

Lemma~\ref{C-361-22}(a) for the $\mathsf n_2^2$ cells of the slab (their sub-independence is Lemma~\ref{C-361-22}(d) on the boxes $(\lambda\mathsf n_1,N',N')$, $N'=\lambda\mathsf n_2$, with even $\mathsf n_1,\mathsf n_2$, in the limit $\mathsf n_1\to\infty$, the reference law being that of one cell of the cube $(2\lambda)^3$ and $\Delta=\delta((2\lambda)^3)/8$), translation invariance of the cylinder along the cross-section, and $\log(1+\mathsf r_{2t})\le\delta(2t,N',N')<0.13736(N'/2t)^2\theta^{t/l}$ give the bound; then $N'\to\infty$ by (c) and density. (c) Lemma~\ref{C-361-21}(c) with $\delta(N^3)<0.13736\,\theta^{N/2l}\le\frac12$, $7.31\cdot0.370622\le2.71$; a second box with all sides tending to $\infty$ gives $\nu$; from continuous to Borel $\Psi$ by approximation in $L^8$ of a measure that dominates both laws; for events take $\Psi=1_A-\frac12$.

The uniform convergence: limit points of box measures are periodic Gibbs states (M3), hence $\nu$. (d) For a $\Gamma$-periodic set $\mathsf S$ of bonds, the flip $U_b\mapsto-U_b$, $b\in\mathsf S$, keeps the Haar measure and carries $\mu^{\mathsf T}_\Gamma$ to $\mu^{\mathsf T\triangle d\mathsf S}_\Gamma$ ($-1$ is central). For a closed $\mathsf T$ and the bonds $\mathsf S$ whose plaquette swept in the direction $\mu$ lies in $\mathsf T$, $\tau_{-ae_\mu}\mathsf T=\mathsf T\triangle d\mathsf S$, by counting the faces of the elementary cubes; so the translates of the twisted measure are flips of it, which a local function does not see once $\mathsf T_n$ has receded. Limit points satisfy the DLR equations away from the twist sets, hence are periodic Gibbs states, hence $\nu$; both identities are linear in the law.

For $P(\beta)>-3\beta$: on the $\prod\lfloor n_i/3\rfloor\ge|\underline n|/125$ plaquettes $p_{12}(x)$, $x_i\in3\mathbb Z$, which share no site, Jensen's inequality gives $F/|\underline n|\ge-3\beta+\frac1{125}\log\int e^{\frac\beta2\operatorname{tr}u}du>-3\beta$; $P$ is convex, $P(0)=0$, so $P'(\beta)>-3$, and $P'=-3\nu(a_p)$ by (a). (e) is Proposition~\ref{C-361-19}(a). (f) In the box with $N=2\vartheta+n+E'$ layers over $(N',N')$, (M1) gives $\mu(\overline{\Theta_nA}\,B)=\operatorname{Tr}(\mathsf T'^{E'}Y_A^*\mathsf T'^nY_B)/\operatorname{Tr}\mathsf T'^N$, $Y_C=\mathfrak Y[\Phi_C]$ the slab operator of the function $C$. As $E'\to\infty$ the gap projects onto the top eigenvector $e_0$:
\[
\begin{multlined}
\mu\bigl(\overline{\Theta_nA}\,B\bigr)\longrightarrow\lambda_0^{-2\vartheta}\Bigl[\sum_{j\ge0}(\lambda_j/\lambda_0)^n\,\overline{\mathsf a_j}\,\mathsf b_j+1_{n=0}\,\langle\mathsf P_{\ker}Y_Ae_0,\mathsf P_{\ker}Y_Be_0\rangle\Bigr],\qquad\mathsf a_j=\langle e_j,Y_Ae_0\rangle,\ \relax \mathsf b_j \\
=\langle e_j,Y_Be_0\rangle,
\end{multlined}
\]
$e_j$ the eigenvectors of the nonzero eigenvalues and $\mathsf P_{\ker}$ the projection on $\ker\mathsf T'$; the term $j=0$ is the product of the means, which converge in the same way. Schwarz's inequality and $\lambda_j/\lambda_0<\theta^{1/l}$ for $j\ge1$ give the bound (for $n>n_0$ the kernel term is absent on the left); $N'\to\infty$ by (c), and $L^2(\nu)$ by density. The same formula exhibits $h_{CC}(n)-|\nu(C)|^2$ as a limit of $\int x^n\varrho(dx)$, $0^0=1$, with $\varrho=\lambda_0^{-2\vartheta}[\sum_{j\ge1}|\mathsf c_j|^2\delta_{\lambda_j/\lambda_0}+\|\mathsf P_{\ker}Y_Ce_0\|^2\delta_0]\ge0$ on $[0,1)$; such functions are $\ge0$ and log-convex.
\end{proof}

\begin{proposition}[Flux and loops on one lattice]\label{C-361-11}
Let $\beta>0$, $l\ge3$, $\mathsf e(l)\le x_0$, $\Lambda=\log(1/\theta)$, $\sigma^+=(2\Lambda+1.67)/4l^2$, so that $\sigma^+l^2=\frac\Lambda2+0.4175\ge0.4757$, $\mathsf q=1+0.835/\Lambda$, and $\nu$ as in Proposition~\ref{C-361-10}; lengths in lattice units.
\begin{enumerate}
\item[(a)] $0\le1-r_1(N)\le2\exp[-2.2\cdot10^{-4}(N/l)^2]$ for every integer $N\ge3$; $\frac12(1-r_1(N))\le\exp[-1.4268\cdot10^{-3}(N/l)^2]$ for $N\ge256\,l$; and for every $N\ge2\mathsf ql$,
\[
\tfrac12(1-r_1(N))\le\exp\Bigl[-\sigma^+N^2+\frac{2\mathsf qN}l\Bigl(\frac12\log\frac{2\mathsf qN}l+0.86\Bigr)\Bigr];
\]
and $1-r^z\le6\exp[\cdots]$ for all seven classes $z\ne0$, with each of the three exponents.
\item[(b)] $0\le\nu(W_C)\le e^{-\sigma^+RT}$ for every rectangle $C$ with $R\times T$ bonds in a coordinate plane; in the finite cube, for $n\ge2\mathsf ql$, $1\le R\le n/2$, $1\le T\le n-1$: $0\le\langle W_C\rangle_{(n,n,n)}\le8^{T/n}\exp[-\sigma^+RT+\frac{2\mathsf qRT}{ln}(\frac12\log\frac{2\mathsf qn}l+0.86)]$; and on every box, with no hypothesis on $\mathsf e$, for $1\le R\le n_1/2$, $1\le T\le n_3-1$ and $\zeta=\frac12(1-r)$ of one twisted column of the plane $(13)$: $0\le\langle W_{R\times T}\rangle_{\underline n}\le(8\,\zeta^{R/n_1})^{T/n_3}$.
\item[(c)] $|\nu(W_C)|\le e^{-\sigma^+|A(C)|}$ for every staircase loop of a coordinate plane, and $|\nu(W_C)|\le2^{c_3(C)-1}e^{-\sigma^+A_{13}(C)}$ for every closed lattice path $C$ with a bond of the direction $3$, with $A_{13},c_3$ as in Theorem~\ref{C-361-3}(4) at unit spacing; the same for every ordered pair of directions. Here a \emph{staircase loop} $C=C(\alpha,R)$ of the plane $(13)$ is the staircase loop of the setting (Twists and loops), read in lattice units: it is given by a site $o$, an integer $T\ge1$ and integers $\alpha_j$ and $R_j\ge1$, $0\le j<T$; with $o_j=o+\alpha_je_1$ and $r_j=o+(\alpha_j+R_j)e_1$ it runs at the height $0$ (in the direction $3$, counted from $o$) straight from $o_0$ to $r_0$; then, for $j=0,\dots,T-1$ in turn, one step up above $r_j$ from the height $j$ to the height $j+1$ and, if $j<T-1$, straight at that height to above $r_{j+1}$; at the height $T$ straight back from above $r_{T-1}$ to above $o_{T-1}$; then, for $j=T-1,\dots,0$ in turn, one step down above $o_j$ from the height $j+1$ to the height $j$ and, if $j>0$, straight at that height to above $o_{j-1}$.

Its area is $|A(C)|=\sum_jR_j$; the staircase loops of the other coordinate planes are the images of these under the symmetries of the lattice. A staircase loop winds once about each unit square of the rows $[\alpha_j,\alpha_j+R_j]\times[j,j+1]$ and about no other, and its steps in the direction $3$ are $T$ upward ones followed by $T$ downward ones; so $A_{13}(C)=|A(C)|$, $c_3(C)=1$, and the first clause is a case of the second.
\item[(d)] For every weak limit point of the measures of periodic boxes all of whose sides tend to infinity, the limits of the tori among them, $w(R,T)=\langle W_{R\times T}\rangle$ satisfies $0\le w\le1$, $w(R,T)^2\le w(R,T-1)w(R,T+1)$ and $w(R,T)^2\le w(R-1,T)w(R+1,T)$, with $w(R,0)=w(0,T)=1$; so $V(R)=\lim_T-T^{-1}\log w(R,T)=\inf_T(-T^{-1}\log w(R,T))$ exists in $[0,\infty]$, $V(0)=0$, $V$ is either finite everywhere or infinite at every $R\ge1$, $V$ is concave and non-decreasing, and, in the finite case, $V(R)\ge V(R')+\sigma_\infty(R-R')$ for $R\ge R'\ge0$, with $\sigma_\infty=\lim_RV(R)/R\in[0,\infty)$. Under the hypothesis, in $\nu$, which is the limit of the tori: $V(R)\ge V(R')+\sigma^+(R-R')$ for $R\ge R'\ge0$.
\item[(e)] With $\varepsilon_{\mathrm v}(l')=\varepsilon_{\mathrm v}(l',l')$ the vortex free energy of Proposition~\ref{C-361-19}(d) and $l_j=2^{j+1}l$: $l_j\varepsilon_{\mathrm v}(l_j)\le e^{-(2\Lambda+1.67)4^j}=e^{-\sigma^+l_j^2}$ for all $j\ge0$.
\end{enumerate}
\end{proposition}

\begin{proof}
For a box $(t,u;h)$, $t,u$ the sides of a plane with one twisted column of length $h$, let $\zeta(t,u;h)=\frac12(1-r)$ for that twist, and let $\varepsilon_{\mathrm v}(t,u)=\log(\|\mathsf T^{(0)}\|/\|\mathsf T^{(1)}\|)$ over the cross-section $(t,u)$, the time along $h$. (i) By Proposition~\ref{C-361-7}(d), in each of the two sides $t$, $u$ of the twisted plane $\zeta Z$ and $(1-\zeta)Z$ are sums of powers of numbers $\ge0$; Hölder's inequality for such sums and the excess give
\[
\zeta(t',u';h)\le[\zeta(t,u;h)\,e^{\delta(t,u,h)}]^{t'u'/tu}e^{-\delta(t',u',h)}\qquad(t'\ge t,\ u'\ge u),
\]
and $\zeta(t,u;h)\le1-e^{-\mathsf G(t;u,h)}\le\delta(t,u,h)$. (ii) In the third side $Z=\operatorname{Tr}(\mathsf T^{(0)})^h=\lambda_0^h(1+\mathsf r_h(t,u))$ and $Z^z=\operatorname{Tr}(\mathsf T^{(1)})^h$, with $\mathsf T^{(1)}\ge0$ of norm $\lambda_0e^{-\varepsilon_{\mathrm v}(t,u)}$. So $1-2\zeta(t,u;h)=Z^z/Z\ge e^{-h\varepsilon_{\mathrm v}}/(1+\mathsf r_h)\ge1-h\varepsilon_{\mathrm v}-\mathsf r_h$ if $h\varepsilon_{\mathrm v}\le1$, and $2\zeta<1$ always: the first inequality below. Further $\operatorname{Tr}(\mathsf T^{(1)})^{2h}\le\|\mathsf T^{(1)}\|^h\operatorname{Tr}(\mathsf T^{(1)})^h$, $Z^z\le Z$ at the length $h$ (Proposition~\ref{C-361-7}(d), with the time in the twisted plane) and $Z\ge\lambda_0^{2h}$ at the length $2h$ give $e^{-h\varepsilon_{\mathrm v}}\ge(1-2\zeta(t,u;2h))/(1+\mathsf r_h)$, the second:
\[
\begin{multlined}
2\zeta(t,u;h)\le h\,\varepsilon_{\mathrm v}(t,u)+\mathsf r_h(t,u),\qquad h\,\varepsilon_{\mathrm v}(t,u)\le\mathsf G(h;t,u)-\log\bigl(1-2\zeta(t,u;2h)\bigr)\le\mathsf d(h;t,u) \\
-\log\bigl(1-2\zeta(t,u;2h)\bigr),
\end{multlined}
\]
where $\mathsf d(h;t,u)=2F(t,u,h)-F(t,u,2h)=2\mathsf G(h;t,u)-\mathsf G(2h;t,u)\ge\mathsf G(h;t,u)$ is taken for the box $(t,u,h)$ with the side $h$ doubled; also $\mathsf r_h(t,u)\ge\mathsf G(h;t,u)$ may stand for $\mathsf G(h;t,u)$. (iii) By (i) and (ii), doubling the cross-section of a long tube, on which the thermal sums lie below the fourth power, replaces $l'\varepsilon_{\mathrm v}(l')$ by a bounded multiple of its fourth power. Let $\mathsf t=\theta^{l'/l}$ at the dyadic sides $l'=l_j=2^{j+1}l$, and let $\mathsf k_1=0.01655$ and $\mathsf k_2=0.001615$ be the constants of the bounds $\delta((2^jl)^3)\le\mathsf k_1\theta^{2^j}$, $j\ge1$, of Proposition~\ref{C-361-8}(b) and $\mathsf r_N\le\mathsf k_2\theta^{N/l}$ at $N=2^{k+1}l$ of Proposition~\ref{C-361-8}(c); so $\delta(l',l',h)\le(h/l')\,\mathsf k_1\mathsf t$ by the monotone density, $\mathsf r_{l'}(l',l')\le\mathsf k_2\mathsf t$ and $\mathsf r_{2l'}(2l',2l')\le\mathsf k_2\mathsf t^2$; and $\mathsf r_{Ms}\le\mathsf r_s^M$ for real $M\ge1$ with $Ms$ an integer, as in Proposition~\ref{C-361-7}(b).

\emph{The start.} The second inequality of (ii) at $t=u=h=2l$, with $\mathsf G\le\mathsf r$ and $\zeta(2l,2l;4l)\le\mathsf R_{2l}(2l,4l)=\mathsf r/(1+\mathsf r)$, $\mathsf r=\mathsf r_{2l}(2l,4l)$, by (i), gives $l_0\varepsilon_{\mathrm v}(l_0)\le\mathsf r_{2l}(2l,2l)+\log\frac{1+\mathsf r}{1-\mathsf r}$. By Proposition~\ref{C-361-7}(b), $\mathsf r_{2l}(2l,2l)\le\mathsf r_l(2l,2l)^2\le\varphi(e^{x_0}-1)^2$, and $\mathsf r\le\mathsf r_l(2l,4l)^2\le\varphi(e^{y_1}-1)^2$ with $y_1=2F(l,2l,4l)-F(2l,2l,4l)\le2x_0+\psi(x_0)$, as in the proof of Proposition~\ref{C-361-8}. Both bounds are power series in $x_0$ with coefficients $\ge0$ beginning with multiples of $x_0^2$, and so is $\log\frac{1+\mathsf r}{1-\mathsf r}$ in $\mathsf r$; the sum divided by $\theta^2=(12.9\,x_0)^2$ is therefore greatest at $x_0=0.069$, where it is at most $0.016154$: $l_0\varepsilon_{\mathrm v}(l_0)\le0.016154\,\theta^2$.

\emph{One doubling.} Let $\hat N\ge2$ be an integer. The first inequality of (ii) in the tube $(l',l';2\hat Nl')$, with $\mathsf r_{2\hat Nl'}(l',l')\le(\mathsf k_2\mathsf t)^{2\hat N}$, gives $\zeta(l',l';2\hat Nl')\le\bar z_1=\hat Nl'\varepsilon_{\mathrm v}(l')+\frac12(\mathsf k_2\mathsf t)^{2\hat N}$; (i) with $t'u'/tu=4$ and $\delta(l',l',2\hat Nl')\le2\hat N\mathsf k_1\mathsf t$ gives $\zeta(2l',2l';2\hat Nl')\le\bar z_2=\bar z_1^4e^{8\hat N\mathsf k_1\mathsf t}$; and the second inequality of (ii) in the tube $(2l',2l';\hat Nl')$, with $\mathsf G\le\mathsf r$, $\mathsf r_{\hat Nl'}(2l',2l')\le\mathsf r_{2l'}(2l',2l')^{\hat N/2}\le(\mathsf k_2\mathsf t^2)^{\hat N/2}$ and $-\log(1-x)\le x/(1-x)$, gives, if $\bar z_2<\frac12$,
\[
2l'\varepsilon_{\mathrm v}(2l')\le\mathsf U:=\frac2{\hat N}\Bigl\{(\mathsf k_2\mathsf t^2)^{\hat N/2}+\frac{2\bar z_2}{1-2\bar z_2}\Bigr\}.
\]
\emph{The iteration.} Take $\hat N_j=4\cdot3^j$ at $l'=l_j$, and let $\Lambda_0=\log(1/0.8901)\approx0.116421$. As $\hat N_j\ge4\cdot2^j$ and $2^{j+1}\hat N_j\ge4^j$, every term of $\mathsf U$ at $\theta\le0.8901$ is at most the term at $\theta=0.8901$ times $e^{-4^{j+1}(2\Lambda-2\Lambda_0)}$, if the bounds on $l_j\varepsilon_{\mathrm v}(l_j)$ carry the factor $e^{-4^j(2\Lambda-2\Lambda_0)}$; so it is enough to follow numbers $E_j$ with $l_j\varepsilon_{\mathrm v}(l_j)\le e^{-E_j-4^j(2\Lambda-2\Lambda_0)}$ at $\theta=0.8901$, $\mathsf t_j=0.8901^{2^{j+1}}$. The start gives such an $E_0\ge\log(1/0.016154)+2\Lambda_0\ge4.3414$; as $\mathsf U$ increases with $l'\varepsilon_{\mathrm v}(l')$, it is enough to take $E_0=4.3414$ and $E_{j+1}=-\log\mathsf U$ at $E_j$.

With $y_j=\hat N_je^{-E_j}$ this reads $E_{j+1}=4E_j-3\log\hat N_j-8\hat N_j\mathsf k_1\mathsf t_j-\log(2B_j)$, where $B_j=(\mathsf k_2\mathsf t_j^2)^{\hat N_j/2}y_j^{-4}e^{-8\hat N_j\mathsf k_1\mathsf t_j}+2(\bar z_1/y_j)^4/(1-2\bar z_2)$; along this sequence $B_j$ is at most $2.0920$ at $j=0$, $2.0001$ at $j=1$ and $2$ to ten digits from $j=2$ on, $\bar z_2\le1.2\cdot10^{-5}$, and the numbers $4^{-j}E_j-2\Lambda_0$ are at least $4.1085$, $2.6061$, $1.9912$, $1.7722$, $1.7032$, $1.6839$, $1.6785$ at $j=0,\dots,6$. For $j\ge6$, $8\hat N_j\mathsf k_1\mathsf t_j\le0.000131$ and $B_j\le2.0023$, so $E_{j+1}\ge4E_j-3\log\hat N_j-\log4.01$, and summing the losses, $4^{-j}E_j\ge4^{-6}E_6-\sum_{i\ge6}4^{-i-1}(3\log4+3i\log3+\log4.01)\ge1.909266$; so $4^{-j}E_j-2\Lambda_0\ge1.6764\ge1.67$ for every $j\ge0$, and $l_j\varepsilon_{\mathrm v}(l_j)\le e^{-(2\Lambda_0+1.67)4^j-(2\Lambda-2\Lambda_0)4^j}=e^{-(2\Lambda+1.67)4^j}$: this is (e).

\emph{From the tubes to the cubes.} For an integer $N\ge2\mathsf ql$ let $j\ge0$ be the largest integer with $4^j\le N/2\mathsf ql$, $w=2^j$, $l'=l_j=2wl$ and $M=N/l'$, a real number; then $w\mathsf q\le M<4w\mathsf q$ and $M^2<2\mathsf qN/l$. The first inequality of (ii) in the tube $(l',l';N)$, with $N\varepsilon_{\mathrm v}(l')=M\,l'\varepsilon_{\mathrm v}(l')\le Me^{-(2\Lambda+1.67)w^2}$ by (e) and $\mathsf r_N(l',l')\le(\mathsf k_2\mathsf t)^M\le e^{-2\Lambda wM}\le e^{-(2\Lambda+1.67)w^2}$, as $M\ge w\mathsf q$ and $2\Lambda\mathsf q=2\Lambda+1.67$, gives $\zeta(l',l';N)\le\frac{M+1}2e^{-(2\Lambda+1.67)w^2}\le Me^{-(2\Lambda+1.67)w^2}$; and (i) with $t'u'/tu=M^2$ and $\delta(l',l',N)\le M\mathsf k_1\mathsf t$ gives $\frac12(1-r_1(N))=\zeta(N,N;N)\le[Me^{-(2\Lambda+1.67)w^2}e^{M\mathsf k_1\mathsf t}]^{M^2}$.

Here $M^2(2\Lambda+1.67)w^2=\sigma^+N^2$; $M\mathsf k_1\mathsf t<4\mathsf q\mathsf k_1we^{-2\Lambda w}\le\frac{4\mathsf q\mathsf k_1}{2e\Lambda}=\frac{4\mathsf k_1}e\bigl(\frac1{2\Lambda}+\frac{1.67}{4\Lambda^2}\bigr)\le0.85478$, the last member decreasing in $\Lambda\ge\Lambda_0$; and $0\le\log M<\frac12\log\frac{2\mathsf qN}l$. So the exponent is at most $-\sigma^+N^2+M^2(\log M+0.855)\le-\sigma^+N^2+\frac{2\mathsf qN}l\bigl(\frac12\log\frac{2\mathsf qN}l+0.86\bigr)$: (a), the last display; the other two bounds of (a) are the same argument with the start at $64\,l$, where $D_2(64l)\le0.12019\,\theta^{64}$, and Proposition~\ref{C-361-7}(d) gives the seven classes. (iv) Two static charges at the ends $o,r$ of the side of length $R$ evolve by an operator $\hat{\mathsf T}=\mathsf m^{1/2}\hat{\mathsf K}\mathsf m^{1/2}$ on functions with values in the $2\times2$ matrices, $(\hat{\mathsf K}\Psi)(V)=\iint e^{-\beta A_\times(V,S,V')}S_o\Psi(V')S_r^{-1}dS\,dV'$, and
\[
2Z[W_C]=\operatorname{Tr}\bigl((\mathsf T^{(0)})^{n_3-T}\Phi_\gamma^*\hat{\mathsf T}^T\Phi_\gamma\bigr)\ge0,\qquad\operatorname{Tr}\hat{\mathsf T}^{n_3}=4Z[\mathsf P\mathsf P'],
\]
$\Phi_\gamma\psi=\psi\,V(\gamma)$, $\mathsf P,\mathsf P'$ the half-traces of the two straight closed lines through $o$ and $r$; the positivity of $\hat{\mathsf T}$ is proved as that of $\mathsf K$, with the crossing variables acting by a unitary representation. Hölder's inequality for the two spectral families gives $0\le\langle W_C\rangle\le(4\langle\mathsf P\mathsf P'\rangle)^{T/n_3}$, and (i) in the direction of $R$ gives $\langle\mathsf P\mathsf P'\rangle\le(1-\zeta)^{1-\varrho}\zeta^\varrho+\zeta^{1-\varrho}(1-\zeta)^\varrho\le2\zeta^{\varrho}$, $\varrho=R/n_1\le\frac12$: the bound on every box. With (a) on the cube $n$ and $n\to\infty$ (Proposition~\ref{C-361-10}(c)) the factor in front tends to $1$: (b). (c) is the same argument with one charge for each passage of $C$ through a layer: several static charges evolve by a positive operator on functions with values in $(\mathbb C^2)^{\otimes\mathcal V}$, and the projection of the path to the plane $(13)$ is counted modulo $2$ layer by layer; a passage from a downward to an upward bond costs a factor $2$. (d) On every periodic box, for $T\le n_3-2$ and $R\le n_1-2$, the positivity of the operators of (iv) gives the log-convexity in $T$, and the reflection positivity across a plane between two charges gives it in $R$; these are inequalities between finitely many means of continuous local functions, and they pass to every weak limit of box measures whose sides tend to infinity.

If $w(R,T_0)=0$ for one $T_0\ge1$, the log-convexity gives $w(R,T)=0$ for all $T\ge1$, and $V(R)=\infty$; else $T\mapsto w(R,T)^{1/T}$ does not decrease, and $V$ is the infimum. The concavity holds in $[0,\infty]$, so $V$ is infinite at every $R\ge1$ as soon as it is infinite at one; with $V(R)/R\to\sigma_\infty$ it gives the slopes. Under the hypothesis $w(R,T)\le e^{-\sigma^+RT}$, so $V(R)\ge\sigma^+R$ for all $R$, and a concave function with $V(R)\ge\sigma^+R$ has every slope $\ge\sigma^+$.
\end{proof}

\begin{lemma}[One step on a periodic box]\label{C-361-29}
Let $\beta\ge0$, $\Sigma$ a cross-section, integers $\vartheta\ge2$, $d\ge3$, $N=2d\vartheta$, $n=(d-1)\vartheta$, $\tau_s=\tau_0+s\vartheta$; let $\Phi$ be a bounded complex Borel function on $\mathcal Y_\vartheta$, $\Phi^s=\Phi$ for odd $s$ and $\Phi^s=\Phi^\Theta$ for even $s$, $\mu_N$ the measure of the box $(N,\Sigma)$, $\varrho=\lambda_1/\lambda_0$, and
\[
\mathsf p=\mu_N\Bigl(\prod_{s=1}^{2d}(\Phi^s)_{\tau_{s-1}}\Bigr),\qquad\mathsf y=[(1+\mathsf r_N)\,\mathsf p]^{1/2d}.
\]
Then $\mathsf p\ge0$, $\lambda_0^{-\vartheta}\|\mathfrak Y[\Phi]\|\le\mathsf y$, and
\[
0\le\mu_N\bigl(\Phi_{\tau_0}(\Phi^\Theta)_{\tau_d}\bigr)\le\bigl[(1+\mathsf r_N)\,|\mu_N(\Phi_{\tau_0})|+\mathsf r_{N-\vartheta}\,\mathsf y\bigr]^2+(2+\mathsf r_n)\,\varrho^{\,n}\,\mathsf y^2 .
\]
\end{lemma}

\begin{proof}
By (M1), $Z(N,\Sigma)\mathsf p=\operatorname{Tr}(Y^*Y)^d=\|Y\|^{2d}_{2d}\ge\|Y\|^{2d}$, $Y=\mathfrak Y[\Phi]$, and $Z(N,\Sigma)=\lambda_0^N(1+\mathsf r_N)$: the bound on $\|Y\|$. The mean of the two cells $d$ slabs apart is $\operatorname{Tr}(\mathsf T'^nY^*\mathsf T'^nY)/Z$; write $\mathsf T'^n=\lambda_0^n(|e_0\rangle\langle e_0|+\mathsf R_{(n)})$ with $\|\mathsf R_{(n)}\|=\varrho^n$ and $\operatorname{Tr}|\mathsf R_{(n)}|=\mathsf r_n$, and $\lambda_0^{-\vartheta}\langle e_0,Ye_0\rangle=(1+\mathsf r_N)\mu_N(\Phi_{\tau_0})-\lambda_0^{-\vartheta}\operatorname{Tr}(\mathsf R_{(N-\vartheta)}Y)$ from the mean of one cell; expanding, the term $|\langle e_0,Ye_0\rangle|^2$ gives the square, and the others are bounded by $\varrho^n$ and $\mathsf r_n$ times $\|Y\|^2$.
\end{proof}

\begin{lemma}[Contraction of the seminorm]\label{C-361-12}
Let $L'\ge13$ be an odd integer, $k\in\mathbb N$, $\beta>0$, on the lattice $L'^{-k}\mathbb Z^3$ with the sides $2L'^{\mathbb N}$. Let $\ell_f,\ell_c\in2L'^{\mathbb N}$ with $d=\ell_c/\ell_f\ge L'$, $\mathsf h=\ell_fL'^k/2$ the number of layers of a fine cell, and let $\mathsf e(l)\le x_0$ at $\beta$ with $24l\le\mathsf h$. Then for every real $G$ of the bonds of the cell $[-\ell_f/4,\ell_f/4]^3$,
\[
\bigl\|G-\mu_{k,\ell_c,\beta}(G)\bigr\|_{\mathrm{ch},\ell_c}\le1.343\,\theta^{(d-1)\mathsf h/8l}\,\bigl\|G-\mu_{k,\ell_f,\beta}(G)\bigr\|_{\mathrm{ch},\ell_f},
\]
where on the left $G$ is regarded as a function of the cell $[-\ell_c/4,\ell_c/4]^3$, which holds the cell $[-\ell_f/4,\ell_f/4]^3$. The constant $1.343$ depends on none of $G$, $k$, $\beta$, $d$, $L'$.
\end{lemma}

\begin{proof}
Let $X=G-\mu_{k,\ell_c,\beta}(G)$, $B_1=X$, $B_{2p}=B_p\cdot\Theta_iB_p$ with $\Theta_i$ the reflection in the coarse wall plane of the direction $i$, and $m_p=\mu_{k,\ell_c,\beta}(B_p)$; then $m_1=0$, and $m_8$ is the eighth power of the left member. In the direction $i$, $B_p$ is a function of one slab of $\mathsf h$ layers, and $\Theta_iB_p$ lies $n=(d-1)\mathsf h$ layers further, on the cube of $N=2d\mathsf h$ sites. Lemma~\ref{C-361-29} with $Y$ the slab operator of $B_p$, the gap of Proposition~\ref{C-361-19}(a), and $\lambda_0^{-\mathsf h}\|Y\|\le\mathsf c\,u^p$, $\mathsf c=(1+\mathsf r_N)^{1/2d}$, $u=\|X\|^{(\ell_f)}_{\mathrm{ch},\ell_c}$, the seminorm of $X$ on the torus $\ell_c$ for the fine board, of side $\ell_f/2$ (the array form of this board, applied to the pile of $2d$ copies of $B_p$, which are $2dp$ copies of $X$ in cells of the fine board, with the factor $1$ in the other cells, bounds the $\mathsf p$ of Lemma~\ref{C-361-29} by $u^{2dp}$), give
\[
0\le m_{2p}\le\bigl[(1+\mathsf r_N)\,m_p+\mathsf r_{N-\mathsf h}\,\mathsf c\,u^p\bigr]^2+(2+\mathsf r_n)\,\theta^{n/l}\,\mathsf c^2u^{2p}.
\]
With the bounds of Proposition~\ref{C-361-8} on the $\mathsf r$, three steps give $m_8\le3.68\,u^8\theta^{n/l}$. To bound $u$, put $G_0=G-\mu_{k,\ell_f,\beta}(G)$ and $Q_f=e^{\delta((2\mathsf h)^3)/8}$, $2\mathsf h=\ell_fL'^k$ being the number of sites of a side of the torus $\ell_f$. The fine-to-coarse form gives $u\le Q_f\|X\|_{\mathrm{ch},\ell_f}$; as $X=G_0-\mu_{k,\ell_c,\beta}(G_0)$ and $\|1\|_{\mathrm{ch},\ell_f}=1$, $\|X\|_{\mathrm{ch},\ell_f}\le\|G_0\|_{\mathrm{ch},\ell_f}+|\mu_{k,\ell_c,\beta}(G_0)|$; and Proposition~\ref{C-361-9}(c) at $\ell'=\ell''=\ell_f$, with $\mathfrak s=\mu_{k,\ell_c,\beta}$ and $\mu_{k,\ell_f,\beta}(G_0)=0$, gives $|\mu_{k,\ell_c,\beta}(G_0)|\le4.84\,[Q_f(Q_f-1)]^{1/2}\|G_0\|_{\mathrm{ch},\ell_f}$. By Proposition~\ref{C-361-19}(c) at $N=2\mathsf h\ge48l$, $\delta((2\mathsf h)^3)\le0.5892\,\theta^{48}\le0.0022046$, so $Q_f\le1.000276$ and $u\le Q_f(1+4.84\,[Q_f(Q_f-1)]^{1/2})\|G_0\|_{\mathrm{ch},\ell_f}\le1.0807\,\|G_0\|_{\mathrm{ch},\ell_f}$. As $m_8$ is the eighth power of the left member and $n=(d-1)\mathsf h$, the left member is at most $3.68^{1/8}\cdot1.0807\,\theta^{(d-1)\mathsf h/8l}\|G_0\|_{\mathrm{ch},\ell_f}\le1.272\,\theta^{(d-1)\mathsf h/8l}\|G_0\|_{\mathrm{ch},\ell_f}$, and $1.272\le1.343$.
\end{proof}

\begin{lemma}[A second blocking factor; dense labels]\label{C-361-28}
(a) For every integer $L\ge2$ there are an odd prime $L'\ge13$ with $L'>L$ and an admissible choice of the constants with the blocking factor $L'$; for such $L'$, $\alpha=\log L/\log L'\in(0,1)$ is irrational. (b) Let $\mathsf b>0$, $k(K)=\lfloor K\alpha\rfloor$ and $\mathsf c_K=\mathsf b\,L'^{K\alpha-k(K)}\in[\mathsf b,L'\mathsf b)$. Then $L^K\mathsf b=L'^{k(K)}\mathsf c_K$; $k$ does not decrease, $k(K+1)-k(K)\in\{0,1\}$ and $k(K)\to\infty$, so every integer $k>k(K_1)$ is $k(K)$ for some $K>K_1$; and $\{\mathsf c_K:K\ge K_1\}$ is dense in $[\mathsf b,L'\mathsf b]$ for every $K_1$.
\end{lemma}

\begin{proof}
(a) is (M6) with the infinitude of primes; $L^q=L'^p$ is impossible for a prime $L'\nmid L$. (b) is arithmetic; the density is that of $\{K\alpha\}$ modulo $1$ for irrational $\alpha$.
\end{proof}

\subsection{Statements uniform in the cutoff}

\begin{proposition}[Products of action densities on one torus]\label{C-361-30}
Let $\ell''\in\mathfrak L$, $J'=[\mathsf b_-,\mathsf b_+]\subset(0,\infty)$ and $\mathsf d>0$. For $k\ge2$ let $\mathsf B_k(\ell'',J',\mathsf d)$ be the bound of (M4)(i) on the $k$-th cumulants on the torus $\ell''$ at the labels of $J'$, for supports at pairwise sup distances $\ge\mathsf d$ and unit norms, and $\mathsf M_k=\sum_\pi\prod_{B\in\pi}\mathsf B_{|B|}$ over the partitions of $\{1,\dots,k\}$ into blocks of at least two members. Let $K\in\mathbb N$, $\mathsf b'\in J'$, $\beta=L^K\mathsf b'$, $\square$ a cell of a board of side $\ell''/2$ with any offset, $f_1,\dots,f_r$ real Lipschitz functions, none $0$, with supports in $\square$ at sup distance $\ge\mathsf d/2$ from its walls and pairwise $\ge\mathsf d$ apart, and $X_i=\mathcal O_K(f_i)-\mu_{K,\ell'',\beta}(\mathcal O_K(f_i))$. Then for $\emptyset\ne I\subset\{1,\dots,r\}$, $X_I=\prod_{i\in I}X_i$ is a real function of the bonds of $\square$ and
\[
\|X_I\|_{\mathrm{ch},\ell''}\le\mathsf M_{8|I|}^{1/8}\prod_{i\in I}\|f_i\|_\sharp .
\]
\end{proposition}

\begin{proof}
The walls of the board lie between lattice sites. Thus a plaquette whose centre lies strictly inside $\square$ has all its corners in the closed cell; the supports lie inside the cell, so $X_I$ depends only on its bonds. By the array form of (M2), $\|X_I\|^8_{\mathrm{ch},\ell''}$ is the mean on the torus $\ell''$ of the product of the eight mirror copies of $X_I$; the mirror copy of $\mathcal O_K(f)$ is $\mathcal O_K(f\circ\hat{\mathsf r})$ for the mirror map $\hat{\mathsf r}$, with the same mean. So it is a moment of $8|I|$ centred fields with pairwise disjoint supports on one torus, at distances $\ge\mathsf d$ inside one cell and across the walls, hence a sum over the partitions into blocks of at least two members of products of joint cumulants; (M4)(i) bounds each by $\mathsf B_{|B|}$ times the norms.
\end{proof}

\begin{proposition}[The vacuum constants of two tori; the cube number]\label{C-361-13}
\begin{enumerate}
\item[(a)] Let $\mathsf b\ge\gamma^{-2}$, $\ell\in\mathfrak L$ with $LM\,R(\mathsf b)\le\ell/2$, and $K\ge10$ with $L^{\lceil K/1260\rceil}\mathsf b\ge\Theta_\star$, a number of the constants, so that the auxiliary sequence of renormalization steps is formed. At every level of the auxiliary sequence started from zero families with zero counterterms, the local numbers on the tori $\ell$ and $L\ell$ are equal on every region that lifts to both. Consequently
\[
|L^3\mathfrak K_K(\ell,\mathsf b)-\mathfrak K_K(L\ell,\mathsf b)|\le C\,(L\ell)^3\exp\bigl[-\kappa_1\ell L^{\lceil K/1260\rceil}\bigr]
\]
with numbers $C,\kappa_1>0$ of the constants.
\item[(b)] For every $\mathsf b>0$, $\ell\in2L^{\mathbb Z}$ and $K\ge K_\star(\mathsf b,\ell)$: $|X_{K+1}(\mathsf b,\ell)-X_K(\mathsf b,\ell)|\le L^3\mathfrak o_K(\ell/L)+\mathfrak o_K(\ell)+C\ell^3e^{-\kappa_1\ell L^{\lceil K/1260\rceil-1}}$ after the dilation of (c), and $\sum_{k\ge K}|X_{k+1}-X_k|\le\eta_K(\mathsf b,\ell)$, with the $\eta_K$ of Theorem~\ref{C-361-5}(1).
\item[(c)] The dilation $(\mathsf b,\ell,K)\mapsto(L^z\mathsf b,L^z\ell,K-z)$ names the same lattices at the same couplings, so $X_K(\mathsf b,\ell)=X_{K-z}(L^z\mathsf b,L^z\ell)$, and it brings every label and side into the range of (a).
\end{enumerate}
\end{proposition}

\begin{proof}
(a) The zero initial families and zero counterterms have a common
infinite-lattice companion and satisfy the input size bound. The
induction in Theorem~\YMMref{Y}{438}, applied to the auxiliary
sequence, therefore constructs the same family on each common
lift. Its proof transports the entire local input through the
chart at every step, including its domain. The floor
$L^{\lceil K/1260\rceil}\mathsf b\ge\Theta_\star$ is the
coupling floor needed for this auxiliary run. The Gaussian
numbers also agree under these charts, by the local normalizer
construction in equations~\YMMref{(1399)--(1400)}{573}.

To estimate the contact terms, we keep the stopping depth in
their vacuum constants. Put $u=\lceil K/1260\rceil$, and mark
the cutoffs $K^{\mathrm A}=K$ and $K^{\mathrm B}=K+1$.
The zero-source chain identity proved in Step~3(iv) of
Proposition~\ref{C-365-11}, or directly
\eqref{eq:vacuum-step-normalization} summed over its steps, gives
\[
 c^\bullet_\lambda(0)=
 E^{\bullet}_{(K^\bullet-u)}[T_\lambda],\qquad
 \mathfrak K_K(\lambda,\mathsf b)
       =c^{\mathrm B}_\lambda(0)-c^{\mathrm A}_\lambda(0).
\]
Here $E^\bullet$ is built from the line
$x_j^\bullet=L^{K^\bullet-j}\mathsf b$ on $T_\lambda$.
These constants are real. The identity uses the actual chains;
it does not require the error estimate comparing their masses.
The line has zero increments and is admissible by the proof of
Lemma~\YMMref{118.1}{812}. The assumed room
$LMR(\mathsf b)\le\ell/2$ admits it on both tori. It also
implies the side floor $\ell\ge C_LM$: the retained room condition
\eqref{eq:running-source-history-room} gives
$C_L\le2LR(\gamma^{-2})\le2LR(\mathsf b)$.

Apply \eqref{eq:vacuum-prefix-comparison} at each cutoff, with
$n=K^\bullet-u$, to the sides $\ell$ and $L\ell$. Both lower
summation indices are $u$. The triangle inequality gives
\[
 \left|\frac{\mathfrak K_K(\ell,\mathsf b)}{\ell^3}
       -\frac{\mathfrak K_K(L\ell,\mathsf b)}{(L\ell)^3}\right|
 \le4C_{\rm tt}E_0\sum_{s\ge u}
                   L^{3s}e^{-\kappa\ell L^s/(5M)}
 \le8C_{\rm tt}E_0 L^{3u}e^{-\kappa\ell L^u/(5M)}.
\]
The last inequality is \eqref{eq:vacuum-prefix-tail}. Since
$\kappa\ell/(5M)>4$, the factor $L^{3u}$ can be absorbed in
half the exponent. For example, take
\[
 \kappa_1=\frac{\kappa}{10M},\qquad
 C=8C_{\rm tt}E_0\sup_{v\ge1}v^3e^{-2v}<\infty.
\]
Multiplication by $(L\ell)^3$ proves (a).

(b) The zero-source read-out (M4)(ii) gives
\[
 |\log Z_{K+1,\ell}-\log Z_{K,\ell}
                   -\mathfrak K_K(\ell,\mathsf b)|
       \le\mathfrak o_K(\ell,\mathsf b).
\]
Apply this on the sides $\ell/L$ and $\ell$, multiply the first
identity by $L^3$, and subtract. The contact terms cancel up to
the error in (a), giving the asserted bound for $X_{K+1}-X_K$.
The two read-out errors have the summable rate in (M4)(ii).
With $\varsigma=\sum_{j\ge0}(1+j)^{\mathsf p_0}L^{-j/10}$,
their tail is bounded by $\eta_K$ once $K$ is above the
thresholds where (M4)(ii) and (a) apply and where the exponential
term is at most
$\mathsf C_{\rm bf}\ell^3K^{\mathsf p_0}L^{-K/10}$.
After the dilation in (c), these requirements are imposed on
$K-z$ and on both sides $L^{z-1}\ell$, $L^z\ell$, at label
$L^z\mathsf b$, including their thresholds $K_{\rm bf}$.
This is the choice of $K_\star$ in
Theorem~\ref{C-361-5}(1).

(c) The dilation leaves both the bare lattice and its coupling
unchanged. This proves the identity for $X_K$ and permits the
fixed dilation used in (b).
\end{proof}

\begin{proposition}[Two cutoffs on boxes with mixed sides]\label{C-361-14}
Assume \textbf{(Hess)} from the setting of Theorem~\ref{C-366-4}. For every compact set $J$ of labels, every class $z\in\{0,1\}^3$, all sides $s_\mu\in L^{\mathbb Z}\cup2L^{\mathbb Z}$ such that each twisted plane has a side in $L^{\mathbb Z}$, and every $l\in2L^{\mathbb Z}$, there are $C$ and $K_*$ such that
\[
\Bigl|\log Z^z_{K+1}(\underline s;\mathsf b)-\log Z^z_K(\underline s;\mathsf b)-\frac{|\underline s|}{l^3}\bigl(\log Z_{K+1,l}-\log Z_{K,l}\bigr)\Bigr|\le C\,L^{-K/10}\qquad(\mathsf b\in J,\ K\ge K_*),
\]
where $Z^z_K(\underline s;\mathsf b)$ is the partition function of the box with $s_\mu/a_K$ sites along the side $\mu$ and the twist $z$, at $\beta=L^K\mathsf b$, and $|\underline s|=s_1s_2s_3$.
\end{proposition}

\begin{proof}
The box $\underline s$ is mixed for $z$ in the sense of Lemma~\ref{C-366-7}. After a common dilation by a power of $L$, choose a periodic covering cube $\ell_C\in2L^{\mathbb Z}$ for both this box and the reference cube $l^3$. The latter is mixed for class $0$. Theorem~\ref{C-366-4}(4), under \textbf{(Hess)}, gives the two comparisons with the same extensive term $\Delta_K[\ell_C,\mathsf b]$.

Subtract $|\underline s|/l^3$ times the comparison for $l^3$ from the comparison for $(\underline s,z)$. The coefficients $|\underline s|/\ell_C^3$ of $\Delta_K$ cancel. The two errors are bounded by $C_1L^{-K/10}$ and $(|\underline s|/l^3)C_2L^{-K/10}$, uniformly on $J$, after increasing the common threshold. This proves the stated estimate.





















The errors are $\le CL^{-K/10}$.
\end{proof}

\begin{lemma}[A summable bound closes the criterion]\label{C-361-26}
Let $(\mathsf e_K)_{K\ge K_*}$ be real numbers with $|\mathsf e_{K+1}-\mathsf e_K|\le\mathsf q_K$ and $Q_K=\sum_{k\ge K}\mathsf q_k<\infty$. Then (a) $\mathsf e_\infty=\lim\mathsf e_K$ exists and $|\mathsf e_K-\mathsf e_\infty|\le Q_K$; (b) if $\mathsf e_{K_1}\le x_0-Q_{K_1}$ at one $K_1\ge K_*$, then $\mathsf e_K\le x_0$ for all $K\ge K_1$; (c) if $\mathsf e_K\le x_0'<x_0$ for all large $K$, then $\mathsf e_{K_1}\le x_0-Q_{K_1}$ at every large $K_1$; (d) $\mathsf e_\infty<x_0$ if and only if $\mathsf e_{K_1}<x_0-Q_{K_1}$ at one $K_1$, and $\mathsf e_\infty>x_0$ if and only if $\mathsf e_{K_1}>x_0+Q_{K_1}$ at one $K_1$, and then $\mathsf e_K>x_0$ for all $K\ge K_1$; at $\mathsf e_\infty=x_0$ neither holds.
\end{lemma}

\begin{proof}
$|\mathsf e_K-\mathsf e_{K'}|\le Q_K$ for $K'\ge K$, and $Q_K\to0$.
\end{proof}

\begin{proposition}[Small tori]\label{C-361-15}
For an integer $N\ge3$ let $\lambda_N(k,u)=\sum_\mu(2-2\cos\frac{2\pi k_\mu+u_\mu}N)$ for $k\in(\mathbb Z/N)^3$, $u\in\mathbb R^3$; $\mathrm{Det}'(N)=\prod_{k\ne0}\lambda_N(k,0)$; $\Upsilon_N(u)=N^2\prod_k\lambda_N(k,u)/\mathrm{Det}'(N)$; $\mathscr J_N=\int_{[-\pi,\pi]^3}du/\Upsilon_N(u)$; $f_\Delta=(2\pi)^{-3}\int_{[-\pi,\pi]^3}\log\sum_\mu(2-2\cos p_\mu)\,dp$; and $c'_N=\log(\mathscr J_N/8\pi^2)-\frac32(\log\mathrm{Det}'(N)-f_\Delta N^3-2\log N)$.
\begin{enumerate}
\item[(a)] There are absolute numbers $c_3>0$ and $x_*\in(0,\frac12]$ with the following two properties. On one bare periodic cube: for every integer $N\ge3$ and every real $\beta\ge3$ with $\mathsf{er}_N=c_3(N^{66}\log\beta/\beta)^{1/8}\le x_*$,
\[
\begin{multlined}
\Bigl|\frac{\beta^{3N^3-1}Z(N^3)}{\mathcal Z(N)}-1\Bigr|\le\mathsf{er}_N,\qquad\mathcal Z(N)=\frac{(2\pi)^{3N^3-1}\,N^2\mathscr J_N}{(2\pi^2)^{2N^3}\,4\pi\,\mathrm{Det}'(N)^{3/2}},\qquad\log\mathcal Z(N) \\
=N^3z_\infty-\log N+c'_N,
\end{multlined}
\]
$Z(N^3)$ the partition function at $\beta$ and $z_\infty$ the number of (d). Consequently, for all integers $N\ge3$, $q\ge2$ and real $\beta\ge3$ with $\mathsf{er}=c_3((qN)^{66}\log\beta/\beta)^{1/8}\le x_*$:
\[
\bigl|D_q(N)-(q^3-1)\log(\beta/N)-q^3c'_N+c'_{qN}-\log q\bigr|\le2(q^3+1)\,\mathsf{er}.
\]
\item[(b)] There are numbers $t_0\in(0,\frac1{36}]$ and $C_B$ of the constants such that for every integer $i\ge1$, $N=2L^i$, and every real $\beta$ with $t=N/\beta\le t_0$: $|D_L(N)-(L^3-1)\log(1/t)-L^3c'_N+c'_{LN}-\log L|\le C_Bt^{1/1500}$.
\item[(c)] With $t_1=t_0^2$ and a number $C_C$ of the constants: $|\delta(N^3)-\log(1/t)-c'_N|\le C_Ct^{1/1500}$ for $N=2L^i$, $t=N/\beta\le t_1$. Moreover $c'_N\to c'$ and $|c'_N-c'|\le C_AN^{-2}(\log N)^{5/2}$ with an absolute $C_A$.
\item[(d)] With $z_\infty=\log\frac2\pi-\frac32f_\Delta\approx-2.96167$ there are numbers $\beta_1$, $C$ of the constants with $|P(\beta)+3\log\beta-z_\infty|\le C\beta^{-1/1500}$ for every real $\beta\ge\beta_1$.
\end{enumerate}
\end{proposition}

\begin{proof}
(a) \emph{The value.} As $\beta\to\infty$ the measure of the bare cube concentrates on the flat fields. Up to gauge these are the commuting triples of holonomies; at a generic triple the orbit is that of a maximal torus, parametrised by $u\in[0,2\pi]^3$, and the transverse fluctuations are Gaussian with the quadratic form of the lattice Laplacian twisted by $u$, whose lowest eigenvalue $\lambda_N(0,u)$ is of the order $|u|^2/N^2$. The volume of the orbits, the Haar measure near them and the Gaussian integral give $\beta^{3N^3-1}Z(N^3)\to\mathcal Z(N)$ and $\log\mathcal Z(N)=(3N^3-1)\log2\pi-2N^3\log2\pi^2-\log4\pi-\frac32\log\mathrm{Det}'(N)+2\log N+\log\mathscr J_N=N^3z_\infty-\log N+c'_N$, by $3\log2\pi-2\log2\pi^2=\log\frac2\pi$, $\log2\pi+\log4\pi=\log8\pi^2$ and the definition of $c'_N$.

\emph{The error, with all its powers of $N$.} (1) \emph{The cut.} In the gauge in which the bonds of a maximal tree carry $\mathbf 1$, which removes the gauge group exactly, with no gauge copies and no Jacobian, the field is bounded by its curvature: with $A=\sum_pa_p$, every other bond variable lies within $2NA^{1/2}$ of $\mathbf 1$ or, at a bond that crosses the boundary of the fundamental cube, of the holonomy $h_\mu$ of its direction, and $|h_\mu h_\nu h_\mu^{-1}h_\nu^{-1}-\mathbf 1|\le N(2A)^{1/2}$. The fields with $A\ge\alpha_\beta=4N^3\log\beta/\beta$ carry at most $\beta^{-4N^3}$, which is $(C/\beta)^{N^3}$ times $\mathcal Z(N)\beta^{1-3N^3}$: this cut is the source of $\log\beta$. (2) \emph{The generic stratum.} On $\{A<\alpha_\beta\}$ the integral is squeezed between two Gaussian integrals: in $d_\nu\le6N^3$ transverse variables $\nu$ with $\|\nu\|^2\le\Lambda_\delta A$ and $|A-Q|\le43N^2|\nu|_\infty\|\nu\|^2$, $Q$ the quadratic part of $A$ and $43N^2$ the bound of the cubic remainder of the plaquettes, the integral of $e^{-\beta A}$ over $\{A<\alpha_\beta\}$ lies between $e^{\mp(d_\nu\epsilon+C\rho_\delta)}$ times the Gaussian integral of $Q$, $\rho_\delta=(\Lambda_\delta\alpha_\beta)^{1/2}$, $\epsilon=43N^2\Lambda_\delta\rho_\delta$; the Gaussian determinants need not be computed, as the limit $\mathcal Z(N)$ names them.

Where some holonomy has $1-(\frac12\operatorname{tr}h_\mu)^2\ge\delta^2$, the base is the commuting triple, there are $6N^3-2$ transverse variables, and the softest of them costs of the order $\delta^2/N^2$: $\Lambda_\delta\le2\pi^2N^5+\frac{\pi^2}4N^2\delta^{-2}$. (3) \emph{The eight points} at which every $h_\mu$ is $\pm\mathbf 1$: there the stabiliser is all of $\mathrm{SU}(2)$, the nine constant modes $c=(c_\mu)$ have the quartic action $N^3|c\wedge c|^2$, $|c\wedge c|^2=\sum_{\mu<\nu}|c_\mu\times c_\nu|^2$, and the Gaussian picture fails for $|u|\lesssim(N/\beta)^{1/4}$. An upper bound suffices there. In Landau's gauge the action is at least $\frac12(1-\epsilon')\|d_1Y'\|^2+N^3|c\wedge c|^2$ for the non-constant part $Y'$ of the field, $d_1$ the lattice exterior derivative of $1$-forms, the Jacobian is at most $(\mathrm{Det}'(N)/N^3)^{3/2}(1+7N^2r_\delta)^{3N^3}$ (Kirchhoff's theorem), $r_\delta^2=6\pi^2N^2\delta^2+24N^5\alpha_\beta$, and the integral over the constant modes is at most $48\pi^3r_\delta N^{-3/2}(\beta N^3)^{-2}$; measured against $\mathcal Z(N)\beta^{1-3N^3}$, in which $\mathscr J_N=\int du/\Upsilon_N(u)$ lies between two absolute numbers, as $c_0|u|^2\le\Upsilon_N(u)\le c_0^{-1}$ with an absolute $c_0>0$ and $|u|^{-2}$ is integrable in three dimensions, the weight of these fields is at most $C\,N^{-1/2}r_\delta\,e^{CN^5r_\delta}$. (4) \emph{The balance.} The error of (2) is of the order of $N^8\delta^{-3}(N^3\log\beta/\beta)^{1/2}$, that of (3) of the order of $N^{1/2}\delta$; with $\delta=N^{-6}(N^{66}\log\beta/\beta)^{1/8}$ both, and the cut, are at most an absolute multiple of $(N^{66}\log\beta/\beta)^{1/8}$: the exponent $\frac18$ is the balance of $\delta^{-3}\beta^{-1/2}$ against $\delta$, all the dependence on $N$ sits in $N^{66}$, and $c_3$, $x_*$ are absolute.

\emph{The consequence.} With $|\log(1+x)|\le2|x|$ for $|x|\le\frac12$, the bound at $N$ and at $qN$ (the condition at $qN$ implies the one at $N$) gives $\log Z(N^3)=-(3N^3-1)\log\beta+\log\mathcal Z(N)+O(2\mathsf{er})$ and the same at $qN$; as $-q^3(3N^3-1)+(3q^3N^3-1)=q^3-1$, the terms in $N^3$ cancel in $D_q(N)=q^3F(N^3)-F((qN)^3)$, and $-q^3\log N+\log(qN)=-(q^3-1)\log N+\log q$. (b) Let $t\le t_0$. \emph{The frame.} Let $j_t$ be the least integer $j\ge1$ with $L^j\ge LM\,R(2L^j/t)$. As $R(Lx)\le L\,R(x)$ once $\log x\ge r$, the exponent of $R(2L^j/t)$ grows by at most one when $j$ grows by one, so equality holds at $j_t$: $\ell'=2L^{j_t}$, $\mathsf b'=\ell'/t$ and $y'=\log\mathsf b'$ satisfy $\ell'=2LM\,R(\mathsf b')$, which is $L$ times the least side at which (M4)(ii), (iii) apply at the label $\mathsf b'$, and $2LMy'^r\le\ell'<2L^2My'^r$, $\log\frac1t<y'<2\log\frac1t$ for small $t$.

Every pair $(N,\beta)$ with $N=2L^i$ and $N/\beta=t$ is the pair $(\ell'/a_K,L^K\mathsf b')$ of the cutoff $K=i-j_t$, of either sign, on the line of $\mathsf b'$ (Definition~\ref{C-361-2}). Write $D_K=D_L(\ell'/a_K)$ at $\beta=L^K\mathsf b'$, which is $X_K(\mathsf b',L\ell')$. \emph{Two cutoffs, uniformly in the label.} The number $\mathfrak o_K$ of (M4)(ii) is $8\mathsf R_{t_1}+12\epsilon_{\rm fg}+8\epsilon_{\rm cl}+12\epsilon_{\rm sh}$ (\textup{(\ref{IV-62-64})}), with the explicit numbers of Section~\ref{sec:cutoff-calibration} at $n=1$: with $y=\log\mathsf b$, $u=\lceil K/1260\rceil$, $s=\lceil K/9\rceil$ and the top level $t_1=K-u$ of (M4)(ii) in $\mathsf R_{t_1}$, $\epsilon_{\rm fg}\le c\,\ell^3L^{-34K/315}$ and $\epsilon_{\rm cl}\le c\,(y+K\log L)^{2\mathfrak p_0}K^{1/2}\ell^3L^{-K/9}$, while $\mathsf R_{t_1}$ carries the factor $e^{-c'A_0^2(y+u\log L)^{2\mathfrak p_0}}$ against $e^{(C+\frac92y)\ell^3}$, and $\epsilon_{\rm sh}$ the factor $e^{-\frac14\kappa_{\rm b}A_1^2(y+u\log L)^{2\mathfrak p_0}}$.

On the sides $2MR(\mathsf b)\le\ell\le2L^3My^r$ one has $\ell^3\le8L^9M^3y^{3r}$ and $y^{2\mathfrak p_0}\le(\ell/2M)^{2\mathfrak p_0/r}$, and the condition $2\mathfrak p_0>3r+2$ of the list of (H3) lets $e^{-c'A_0^2y^{2\mathfrak p_0}}$ beat $e^{(C+\frac92y)\ell^3}$. So there are numbers $y_\sharp$, $C_{\mathfrak o}$ of the constants such that at every label with $y\ge y_\sharp$, on every such side $\ell$ and at every $K\ge11(5+2\mathfrak p_0/r)\log_L\ell$, the conditions (k0)--(k4), (k6) of Section~\ref{sec:cutoff-calibration} hold and $\mathfrak o_K(\ell)\le C_{\mathfrak o}\,\ell^{3+2\mathfrak p_0/r}L^{-K/11}$, uniformly in the label; this is the form needed here, as $\mathsf b'\to\infty$ when $t\to0$, while the constant $\mathsf C_{\rm bf}$ of (M4)(ii) is formed on bounded sets of labels.

The two sides $\ell'$, $L\ell'$ lie in this range, and (M4)(ii) on both, with Proposition~\ref{C-361-13}(a), gives
\[
\begin{multlined}
|D_{K+1}-D_K|\le L^3\mathfrak o_K(\ell')+\mathfrak o_K(L\ell')+C\ell'^3L^{-K}\le C\,\ell'^{\,3+2\mathfrak p_0/r}L^{-K/11}\qquad\bigl(K\ge11(5 \\
+2\mathfrak p_0/r)\log_L(L\ell')\bigr).
\end{multlined}
\]
\emph{The base.} Let $K_\circ$ be the least integer with $N_\circ=\ell'L^{K_\circ}\ge t^{-1/130}$, so that $t^{-1/130}\le N_\circ<Lt^{-1/130}$; as $\ell'$ is a power of $\log\frac1t$, $K_\circ$ lies above the bound on $K$ of the display once $t\le t_0$. Let $N=2L^i$ and $\beta=N/t$. (1) If $N\le N_\circ$, then $(LN)^{66}\log\beta/\beta=L^{66}N^{65}t\log\frac Nt\le L^{131}t^{1/2}(\log L+\frac{131}{130}\log\frac1t)$, and (a) at $q=L$ applies for $t\le t_0$ and bounds the left member of (b) by $2(L^3+1)c_3L^{17}t^{1/16}(\log L+\frac{131}{130}\log\frac1t)^{1/8}$. (2) If $N>N_\circ$, then $K=i-j_t>K_\circ$, and with $\mathsf k_M=L^3c'_M-c'_{LM}+\log L$,
\[
D_L(N)-(L^3-1)\log\tfrac1t-\mathsf k_N=\bigl[D_K-D_{K_\circ}\bigr]+\bigl[D_L(N_\circ)-(L^3-1)\log\tfrac1t-\mathsf k_{N_\circ}\bigr]+\bigl[\mathsf k_{N_\circ}-\mathsf k_N\bigr]:
\]
the first bracket is at most $C(1-L^{-1/11})^{-1}\ell'^{\,3+2\mathfrak p_0/r+1/11}t^{1/1430}$ by the display, as $L^{-K_\circ}\le\ell't^{1/130}$; the second is case (1) at $N_\circ$; the third is at most $2(L^3+1)C_AN_\circ^{-2}(\log N_\circ)^{5/2}\le2(L^3+1)C_At^{1/65}(\log L+\frac1{130}\log\frac1t)^{5/2}$ by the rate of (c), whose proof does not use (b). The exponents $\frac1{16}$, $\frac1{1430}$, $\frac1{65}$ exceed $\frac1{1500}$, and the powers of $\log\frac1t$ are absorbed into $C_B$; $t_0$ is fixed by $\log\frac1{t_0}\ge y_\sharp$, by the place of $K_\circ$ and by the condition of (a) in case (1). (c) By Proposition~\ref{C-361-7}(c) and $\delta(M^3)\le3(M/N)^2\delta(N^3)$, $(1-3L^{-S-1})\delta(N^3)\le\sum_{s=0}^SL^{-3(s+1)}D_L(L^sN)\le\delta(N^3)$, and with $Y(M)=\log(\beta/M)+c'_M$ the main terms of (b) telescope:
\[
\sum_{s=0}^{S}L^{-3(s+1)}\bigl[(L^3-1)\log(\beta/L^sN)+L^3c'_{L^sN}-c'_{L^{s+1}N}+\log L\bigr]=Y(N)-L^{-3(S+1)}Y(L^{S+1}N);
\]
choosing $L^S$ of the order $t^{-1/2}$ gives (c). The limit $c'$ and the rate are those of Riemann sums for lattice determinants, written through heat kernels: $\log\mathrm{Det}'(N)$, $f_\Delta N^3$ and $2\log N$ are integrals $\int_0^\infty[\cdots]\frac{d\tau}\tau$ of one-dimensional heat kernels, which converge to those of $\vartheta(\tau,\cdot)$. (d) At $N=2L$ and $\beta\ge2L/t_1$, (a) gives $F(N^3)=-(3N^3-1)\log\beta+N^3z_\infty-\log N+c'_N+O((\log\beta/\beta)^{1/8})$ and (c) gives $\delta(N^3)=\log(\beta/N)+c'_N+O(\beta^{-1/1500})$; the difference is $N^3P=-3N^3\log\beta+N^3z_\infty+O(\beta^{-1/1500})$.
\end{proof}

\begin{proposition}[The flux quotient of small cubes at every cutoff]\label{C-361-16}
\begin{enumerate}
\item[(a)] There are numbers $t_2>0$, $C_\flat$ of the constants, and $\mathsf a=\frac92r+\mathfrak p_0$, such that for every integer $j\ge1$ and every real $\beta$ with $t=2L^j/\beta\le t_2$: $r_1(2L^j)\le\eta_\flat(t)=C_\flat(\log\frac1t)^{\mathsf a}t^{1/500}$ at $\beta$. Equivalently: for every side $\ell=2L^m$, $m\in\mathbb Z$, every real label $\mathsf b$ with $\ell/\mathsf b\le t_2$, and every cutoff $K\ge1-m$.
\item[(b)] Assume \textbf{(Heaviest class)}. There are numbers $t_4>0$, $C$ of the constants, and $\mathsf q=1116\,r+720\,\mathfrak p_0+366$, such that for every integer $j\ge1$ and every real $\beta$ with $t=2L^j/\beta\le t_4$, over the cross-section $(N,N)$, $N=2L^j$: $\lambda_0^N/Z(N^3)\le C\,t\,(\log\frac1t)^{\mathsf q}$.
\end{enumerate}
\end{proposition}

\begin{proof}
(a) \emph{On one lattice.} (i) Multiplying by $-1$ the bonds of a sheet (a half-plane of bonds bounded by two columns) moves a twisted column; so $Z^z$ does not depend on the place of its column, and reflection positivity across two layers with the column in the cut gives $0\le Z^z[\chi\chi']\le Z[\chi]^{1/2}Z[\chi']^{1/2}$ for $[0,1]$-valued $\chi,\chi'$ of the two halves. On configurations whose block plaquettes at a level $q$ (or the bare plaquettes) lie within $\omega$ of $\mathbf 1$, a twist forces the holonomies $P_1,P_3$ of the block field around the twisted plane to anticommute up to $\varpi=(1+\omega)^{\mathsf n^2}-1\le\frac1{20}$, $\mathsf n$ the number of block sites a side; hence $|W_1|\le\varpi/2$ and $W_2^2\ge1-0.52\,\varpi^2$ for the half-traces $W_\mu(y)$ of the straight closed lines of the block field through a site, and
\[
r_1(N)\le\max_y\mu_{N^3}\bigl(\mathcal C^\varpi(y)\bigr)+3\mathsf n^3\mathsf R^{1/2},\qquad\mathcal C^\varpi(y)=\bigl\{|W_1(y)|\le\tfrac\varpi2,\ W_2(y)^2\ge1-0.52\,\varpi^2\bigr\},
\]
$\mathsf R$ the largest probability of a rough block plaquette. (ii) On a bare cube small against $\beta$: with $\Xi_\beta=3N^3\log(8e^8\beta^{3/2})+2\log(3N^3\beta)$ and $\omega_0=(2\Xi_\beta/\beta)^{1/2}$, $Z\ge(8e^8\beta^{3/2})^{-3N^3}$, so a plaquette is rough, $|\partial U(p)-\mathbf 1|\ge\omega_0$, with probability $\le(3N^3\beta)^{-2}$. Then the angle of $P_1$ is shifted along the seam: with $P_1(y_0)=e^{i\alpha_0\hat n\cdot\sigma}$ at the site $y_0$ of the event, $0<\alpha_0<\pi$, $\sigma$ the Pauli matrices, $\mathsf N(x)$ the matrix $\hat n\cdot\sigma$ carried to $x$ along a path from $y_0$ that uses no seam bond $\langle x,x+e_k\rangle$, $x_k=0$, $k=1,2$, and $V(b)\mapsto V(b)e^{-is_k\mathsf N(b_+)}$ on the seam bonds of the direction $k$, this map, defined for $|s|_\infty\le\frac1{10}$, commutes with gauge transformations, turns $P_1(y_0)$ by $-s_1$ and $P_2(y_0)$ by $-s_2$ about $\hat n$, changes the action by at most $131N^4\beta\omega_0^2|s|_\infty$ on the fields whose plaquettes all lie within $\omega_0$ of $\mathbf 1$, and changes the Haar measure by the factor $\sin^2(\alpha_0-s_1)/\sin^2\alpha_0$.

Let $D$ be the event $\mathcal C^\varpi(y_0)$ intersected with that of no rough plaquette. For $s\in2\varpi\mathbb Z^2$ with $|s|_\infty\le s_0=(131N^4\beta\omega_0^2)^{-1}$, which is $<\frac1{10}$, the action rises by at most $1$, so the image of $D$ has the measure $\ge e^{-1}\cos^2(0.13)\,\mu_{N^3}(D)\ge0.36\,\mu_{N^3}(D)$; on $D$ the angles of $P_1(y_0)$ and $P_2(y_0)$ about $\hat n$ fix $s$ up to $1.02\varpi$ and $1.46\varpi$, so the images of distinct points of the grid are disjoint, and there are at least $(s_0/2\varpi)^2$ of them. Hence $\mu_{N^3}(D)\le\frac4{0.36}\varpi^2(131N^4\beta\omega_0^2)^2$, which is of the order $N^{12}\Xi_\beta^3/\beta$, as $\beta\omega_0^2=2\Xi_\beta$ and $\varpi\le2.1N^2\omega_0$; with the rough plaquettes, $\mu_{N^3}(\mathcal C^\varpi)\le[6.8\cdot10^6N^{12}\Xi_\beta^3+1]\beta^{-1}$, for the block field of a bare lattice that is not too large against $\beta$ as well, by (M4)(v). (iii) \emph{Through the cutoffs.} Write $2L^j=\ell L^K$ with $\ell=2L^m$ the least side at which (M4) applies at the label $\mathsf b=\ell/t$, $y=\log\mathsf b$, $\ell\le2L^2My^r$.

At the cutoffs $K\le K_1=\lceil y/(25\log L)\rceil$ the lattice is a bare cube small against $\beta$, and (ii) applies. Above $K_1$ we use (M4)(iii) at the depth $d=0$: at the cutoff $K$ the block field is $\mathcal M^KU$, on the unit lattice of the torus $\ell$, with $\ell$ sites a side at every cutoff; the roughness $\omega$ and the width $\varpi$ are chosen at $K_1$ and do not depend on $K$. The events $\mathcal C^\varpi(y')$ of this block field and the events that one of its plaquettes lies at a distance $\ge\omega$ from $\mathbf 1$ are closed and gauge invariant, and the same at all cutoffs; so by (M4)(iii) their probabilities move by at most $\mathfrak o_K$ from one cutoff to the next, and $\sum_{K\ge K_1}\mathfrak o_K\le Q(y)e^{-y/250}$, $Q$ a polynomial formed from the constants; (i) at the cutoff $K$ then gives $r_1\le4\bar\ell(y)^3Q(y)^{1/2}e^{-y/500}$, $\bar\ell(y)=2L^2My^r$, which is $\le C_\flat(\log\frac1t)^{\mathsf a}t^{1/500}$. (b) Under \textbf{(Heaviest class)}, use another label of a layer and a multiplicative comparison through the cutoffs.

\emph{The top level against the heaviest cell of a label.} Let $\Phi$ map the configurations on the bonds of the directions $1,2$ of one layer into a set $\mathsf X$, let $\mathcal P$ be a finite partition of $\mathsf X$ into cells whose preimages are Borel, $\sim$ a reflexive and symmetric relation on $\mathcal P$ with at most $\nu$ neighbours to each cell, and $\iota_k$ the cell of the label of the layer $x_3=k$. Then over the cross-section $(N,N)$ of an even cube $N=2h$
\[
e^{-\mathsf G_N}=\frac{\lambda_0^N}{Z(N^3)}\le\bigl(\nu^2\sqrt{p^\sharp}+\nu\sqrt{\varepsilon_2}+\sqrt{\varepsilon_1}\bigr)^2\le2\nu^4\bigl(p^\sharp+\max\{\varepsilon_1,\varepsilon_2\}\bigr)\quad(\nu\ge2),
\]
$p^\sharp=\max_I\mu_{N^3}(\iota_1=I)$, $\varepsilon_1=\mu_{N^3}(\iota_1\not\sim\iota_{\mathsf n/2-1})$, $\varepsilon_2=\mu_{N^3}(\iota_{\mathsf n/2-1}\not\sim\iota_{\mathsf n/2+1})$, in two cases: ($\alpha$) the label of a bare layer, $\mathsf n=N$; ($\beta$) the label of a layer of the block field $\mathcal M^qU$ of the bare cube, with $\mathsf n\ge8$ block sites a side. Indeed, let the first slab be the bare layers $1,\dots,s$, with $s=2$ in ($\alpha$), $s=\frac12(3L^q-1)$ in ($\beta$), and $2s<h$, which in ($\beta$) follows from $\mathsf n\ge8$. The label $\iota_1$ is a Borel function of the bonds with both ends in the first slab: in ($\alpha$) it reads the layer $x_3=1$ alone, and in ($\beta$), by the locality of the averaging map (Proposition~\YMMref{4.5}{20}), only the bare bonds of the blocks of the block layer $1$, which, centred at the height $L^q$, lie in the bare layers $\frac12(L^q+1),\dots,\frac12(3L^q-1)$.

The reflection in the plane $x_3=\frac h2$, the translation by $he_3$ and the reflection in the plane $x_3=0$ carry the first slab onto the layers $h-s,\dots,h-1$, the layers $h+1,\dots,h+s$ and the layers $N-s,\dots,N-1$; the layers $0$ and $h$ lie in none of these four slabs. The three maps are isometries of the bare lattice and, in ($\beta$), as $h=\frac{\mathsf n}2L^q$ is a multiple of $L^q$, of the block lattices, and the averaging maps commute with them (Lemma~\ref{lem:setup-lattice-covariance}); so the label $\iota_1$, read on the configuration moved by them, is $\iota_{\mathsf n/2-1}$, $\iota_{\mathsf n/2+1}$, $\iota_{\mathsf n-1}$ respectively. Let $\chi_I$ be the indicator function of the label $I$ on the first slab, $\mathfrak Y[\chi_I]$ its slab operator of (M1), and $Y_I=\mathfrak Y[\chi_I]\,\mathsf T'^{\,1/2}$, the operator $\mathsf T'=\mathsf K^{1/2}\mathsf m\,\mathsf K^{1/2}$ being $\ge0$; then $\sum_IY_I=\mathsf T'^{\,s+1/2}$, $R=\mathsf T'^{\,h-2s-1}\ge0$ and $\mathsf T'^{\,h}=\sum_{I,J}Y_J^*RY_I$.

The conjugated mirror image $\chi_J^\Theta$ of (M1), placed on the layers $h-s,\dots,h-1$, is $\chi_J$ read on the configuration mirrored in the plane $x_3=\frac h2$ midway between these layers and the first slab, that is the indicator of $\iota_{\mathsf n/2-1}=J$; placed on the layers $N-s,\dots,N-1$, it is $\chi_J$ read on the configuration mirrored in the plane $x_3=h$, which on the box of $N=2h$ layers is the reflection in the plane $x_3=0$, that is the indicator of $\iota_{\mathsf n-1}=J$; and $\chi_{I'}$ placed on the layers $h+1,\dots,h+s$ is the indicator of $\iota_{\mathsf n/2+1}=I'$. So the trace formula of (M1), in which the first slab stands rightmost, for the functions $\chi_I^{+(h-2s-1)}$, $(\chi_J^\Theta)^{+1}$, $\chi_{I'}^{+(h-2s-1)}$, $(\chi_{J'}^\Theta)^{+1}$ on the consecutive slabs of $h-s-1$, $s+1$, $h-s-1$, $s+1$ layers from the layer $1$ on, with $\mathfrak Y[\Phi^{+j}]=\mathsf T'^{\,j}\mathfrak Y[\Phi]$ (the $j$ added layers lying above those of $\Phi$), $\mathfrak Y[\chi_J^\Theta]=\mathfrak Y[\chi_J]^*$ and the cyclicity of the trace, gives $\operatorname{Tr}(Y_{J'}^*RY_{I'}\,Y_J^*RY_I)=\operatorname{Tr}(\mathsf T'\mathfrak Y[\chi_{J'}^\Theta]\,R\,\mathfrak Y[\chi_{I'}]\,\mathsf T'\mathfrak Y[\chi_J^\Theta]\,R\,\mathfrak Y[\chi_I])=Z(N^3)\,\mu_{N^3}(\iota_1=I,\iota_{\mathsf n/2-1}=J,\iota_{\mathsf n/2+1}=I',\iota_{\mathsf n-1}=J')$: probabilities of one and the same cube.

The terms with $I\not\sim J$ have together the Hilbert--Schmidt norm $\le(Z\varepsilon_1)^{1/2}$. The terms with $I\sim J$ are $\mathbb X\mathbb B\mathbb X^*$, with $\mathbb X(\varphi_I)_I=\sum_IY_I^*R^{1/2}\varphi_I$ and the neighbour matrix $\mathbb B$ of norm $\le\nu$ (Schur's test); $\mathbb X^*\mathbb X$ is the matrix of the $R^{1/2}Y_IY_J^*R^{1/2}$, whose Hilbert--Schmidt norms are $(Z\mu_{N^3}(\iota_1=\iota_{\mathsf n/2-1}=I,\ \iota_{\mathsf n/2+1}=\iota_{\mathsf n-1}=J))^{1/2}$, so Schur's test on the neighbouring entries and the Hilbert--Schmidt norm of the others give $\|\mathbb X^*\mathbb X\|\le\nu(Zp^\sharp)^{1/2}+(Z\varepsilon_2)^{1/2}$. Hence $\lambda_0^h=\|\mathsf T'^{\,h}\|\le Z^{1/2}[\nu(\nu\sqrt{p^\sharp}+\sqrt{\varepsilon_2})+\sqrt{\varepsilon_1}]$.

\emph{The label: the class of the pair of holonomies.} For pairs $P=(P_1,P_2)$, $Q$ in $\mathrm{SU}(2)^2$ let $\mathsf D(P,Q)=\min_g\max_{i=1,2}|P_i-gQ_ig^{-1}|$, a distance between the classes under simultaneous conjugation, $\mathcal O_w(Q)=\{P:\mathsf D(P,Q)\le w\}$, and $\mathsf v(r)$ the volume of a ball of radius $r$ of $\mathrm{SU}(2)$. For $0<w\le\frac2{35}$ a maximal family of points $Q^{(k)}$ at mutual class distances $>w$ gives a finite partition of $\mathrm{SU}(2)^2$ into Borel cells $I_k\subset\mathcal O_w(Q^{(k)})$, invariant under simultaneous conjugation; with $I_k\sim I_l$ if and only if $\mathsf D(Q^{(k)},Q^{(l)})\le3w$, two pairs at the class distance $\le w$ lie in neighbouring cells, and $\nu\le\mathsf v(\frac72w)^2/\mathsf v(\frac12w)^2\le1.21\cdot10^5$.

The label of a layer is the cell of the pair $(P_1,P_2)$ of the holonomies of its straight closed lines through one site, and $E_k$ is the event that this pair lies in $I_k$. The label is rigid across the layers: if all plaquettes of a lattice of $\mathsf n$ sites a side lie within $\varpi$ of $\mathbf 1$, one conjugation carries both holonomies at $a+ke_3$ to within $k\mathsf n\varpi$ of those at $a$, so that $\mathsf D\le k\mathsf n\varpi$. \emph{No class is heavy on a bare even cube.} Let $\rho_\omega$ be the density, against the Haar measure of $\mathrm{SU}(2)^3$, of the three holonomies $(P_1,P_2,P_3)$ through a site on the event that all plaquettes lie within $\omega$ of $\mathbf 1$; it vanishes unless $|[P_\mu,P_\nu]|\le N^2\omega$, and it is greatest at $(\mathbf 1,\mathbf 1,\mathbf 1)$, as follows.
Choose a pointwise version of $\rho_\omega$ by solving one seam bond on each of the three straight winding loops for $P_\mu$ and integrating over all other bonds. The loops have no bonds in common, and these substitutions preserve product Haar measure. Fix $\mu$ and cut at the site planes $x_\mu=0,h$, where $N=2h$. Let $B$ be their tangential bonds and $X_+,X_-$ the private bonds in the two halves, identified by reflection. No plaquette crosses a site plane, so the cut Wilson weight factors as $w_0(B)w_+(B,X_+)w_+(B,X_-)$, with $dU=dB\,dX_+\,dX_-$ and nonnegative factors bounded by one. The plaquette indicators factor in the same way. At fixed $B$, solve one private bond on the half-path from $0$ to $he_\mu$ for its transport. Its weighted Haar density $k_B$ is bounded and hence belongs to $L^2(G)$. If the two half-transports from $0$ to $he_\mu$ are $g_+,g_-$, then $P_\mu=g_+g_-^{-1}$ and its unnormalised conditional density is $q_B(p)=\int k_B(pg)k_B(g)\,dg\le\int k_B(g)^2dg=q_B(\mathbf1)$ by Schwarz and Haar invariance. The other two loops lie in $x_\mu=0$, so their holonomies and their solved seam bonds are functions of $B$. Integrating this inequality over the remaining boundary bonds with those holonomies fixed proves $\rho_\omega(P)\le\rho_\omega(P_1,\ldots,P_{\mu-1},\mathbf1,P_{\mu+1},\ldots,P_3)$ in the pointwise version just specified. Apply it for $\mu=1,2,3$. The factorization is also the site-mirror case of the proof of Proposition~\YMMref{91.4}{700}, which keeps the boundary bonds as separate Haar variables.
For the quantitative class probability we use \textbf{(Heaviest class)}. Take $\Xi'=\beta\omega_0^2\ge1$ with $\Xi' N^5\omega_0\le3/640$ and $N^2\omega_0\le w\le1/10$.

The input then gives, for every pair $Q$ and every site $y$,
\[
\mu_{N^3}\bigl(\mathsf D((P_1,P_2)(y),Q)\le w,\ \text{no rough plaquette}\bigr)\le10^{11}\,\Xi'^5N^7w^6\omega_0^{-4}.
\]
On the bare cube, with $\omega_0$ and $\Xi_\beta$ of (a)(ii), $\Xi'=2\Xi_\beta$ and $w=N^2\omega_0$, as long as $2\Xi_\beta N^5\omega_0\le\frac3{640}$: where no plaquette is rough, $\mathsf D$ between the labels of the layers $1$ and $\frac N2-1$ is $<w$ and between those of the layers $\frac N2\mp1$ it is $\le2N\omega_0\le w$, so $\varepsilon_1,\varepsilon_2$ are at most the probability, $\le\beta^{-2}$, that some plaquette is rough, $p^\sharp\le10^{11}(2\Xi_\beta)^6N^{19}\beta^{-1}+\beta^{-2}$, and, the bound increasing with $\nu$, with $1.21\cdot10^5$ for $\nu$ and $2\nu^4=4.2872\cdot10^{20}$: $e^{-\mathsf G_N}\le2.8\cdot10^{33}\Xi_\beta^6N^{18}\,t$ on the bare cube. \emph{Through the cutoffs.} With $y=\log\mathsf b$ and the side $\ell$ of (a)(iii), this serves at the cutoffs below the cutoff $K_1(y)$ of (M4)(vi), Proposition~\ref{C-361-23}(c), which in (b) we call $K_1$: the least $K\ge10$ with $L^{K/10}\ge\bar\ell(y)^3y^{2\mathfrak p_0+1}$, so that $L^{K_1}<L\,(\bar\ell(y)^3y^{2\mathfrak p_0+1})^{10}$ and the bare cube has $N_1=\ell L^{K_1}<\mathsf N(y)=L\,\bar\ell(y)^{31}y^{10(2\mathfrak p_0+1)}$ sites a side, a fixed power of $y$; at the bare cubes $N\le N_1$ one has $\Xi_\beta\le10N^3y$, and the condition $2\Xi_\beta N^5\omega_0\le\frac3{640}$ holds once $y$ is large, as $t\le\bar\ell(y)e^{-y}$.

Above $K_1$ the additive transport of (a) does not serve: even summed from the higher cutoff $\lceil y/(25\log L)\rceil$ of (a)(iii), its error $\sum_K\mathfrak o_K\le Q(y)e^{-y/250}$ is of the order of $t^{1/250}$ up to powers of $\log\frac1t$, and swamps the first power of $t$; summed from $K_1$ it is only of the order $1/y$. We use instead the comparison with a factor of (M4)(vi), Proposition~\ref{C-361-23}(c): from $K_1$ on its factors multiply to at most $e^{0.53}<1.7$, the sums from $K_1$ on of the numbers of Section~\ref{sec:cutoff-calibration} being of the order $1/y$, and its remainders add up to at most $e^{-4y}$. At $K_1$ take $\omega_0$ of (a)(ii) at $\beta_1=N_1/t$, $w=N_1^2\omega_0$, the cells $I_k$ of this $w$, and for the block field the roughness $\omega=1.003\,L^{2K_1}\omega_0$; none of these reads $K$.

By (M4)(v), where all bare plaquettes lie within $\omega_0$ of $\mathbf 1$ (for large $y$, $L^{2K_1}\omega_0\le N_1^2\omega_0$ lies below the number of $L$ that (M4)(v) asks), the block plaquettes of $\mathcal M^{K_1}U$ lie within $\omega$ of $\mathbf 1$ and its holonomies within $13.3\,N_1L^{K_1}\omega_0\le0.512\,w$ of the bare ones; so the bound on the classes of the bare cube $N_1$, taken at $1.52\,w$, gives for each of the finitely many events $E_k$ of the block field of depth $0$ at a site $\mu^{K_1}_0(E_k)\le\beta_1^{-2}+7.9\cdot10^{19}N_1^{36}y^6\,t$, and for the event $\mathfrak R$ of a rough block plaquette $\mu^{K_1}_0(\mathfrak R)\le\beta_1^{-2}$. By (M4)(vi), in the form of Proposition~\ref{C-361-23}(c), used for each of these gauge-invariant events apart (nothing is summed over the cells: $p^\sharp$ is a maximum), at every $K\ge K_1$ each has a probability at most $1.7$ times its probability at $K_1$ plus $1.7\,e^{-4y}$.

At the cutoff $K$ the label ($\beta$) of the block field of depth $0$, on the torus $\ell$ with $\mathsf n=\ell$ block sites a side, is rigid with $\omega$ for $\varpi$: between the layers $1$ and $\frac{\mathsf n}2-1$ the class distance is $<\frac12\ell^2\omega<w$, between $\frac{\mathsf n}2\mp1$ it is $\le2\ell\omega<w$, so $\varepsilon_1,\varepsilon_2\le\mu^K_0(\mathfrak R)$. As $\beta_1^{-2}\le t$ and $e^{-4y}\le t$, the bound of the top level by the heaviest cell gives $e^{-\mathsf G_N}\le6\cdot10^{40}N_1^{36}y^6\,t$, which is at most $C\,t\,(\log\frac1t)^{\mathsf q}$, as $N_1<\mathsf N(y)$ and $\log\frac1t\le y\le2\log\frac1t$; the bound $2.8\cdot10^{33}\Xi_\beta^6N^{18}t\le2.8\cdot10^{39}\mathsf N(y)^{36}y^6t$ of the bare cubes below $K_1$ is smaller.
\end{proof}

\subsection{Proof of the theorems}

\emph{Step 1: Theorem~\ref{C-361-3}(1).} BI is the hypothesis $\mathsf e(l)\le x_0$ of Propositions~\ref{C-361-8}, \ref{C-361-19}, \ref{C-361-10}, \ref{C-361-11} and of Lemma~\ref{C-361-20} at $l=l_K$, $\beta=L^K\mathsf b$, with one $x_0$ for all $K$. For $N=\ell/a_K$,
\[
\frac N{2l_K}=\frac\ell{\ell_0},\qquad\theta^{1/l_K}=e^{-m_*a_K},\qquad N\ge2l_K\iff\ell\ge\ell_0 .
\]
(a) is Proposition~\ref{C-361-8}(b), (c) and Proposition~\ref{C-361-19}(c), with $X_K(\mathsf b,\ell)\le L^3\delta((\ell/La_K)^3)$ (Proposition~\ref{C-361-7}(c)). (b) is Proposition~\ref{C-361-10}(a), (b). (c): with $y=\theta^{\ell/\ell_0}\le0.8901$, convexity gives $\varepsilon\le22.3(e^{0.13736/7}-1)\,y\le0.441912\,y$ in Proposition~\ref{C-361-9}(b), and $\frac12(\sqrt\varepsilon+\varepsilon)\le0.540843\sqrt y$; with $y=\theta^{\ell'/\ell_0}$, $Q',Q''\le e^{0.13736/8}\le1.017319$ and $Q'(Q'-1)\le0.0176182\,y$, so Proposition~\ref{C-361-9}(c) gives $4.84\cdot0.132734\cdot1.017319\le0.66$. The boxes are Proposition~\ref{C-361-10}(c) with $\theta^{N/4l_K}=\theta^{\ell/2\ell_0}$, the twisted boxes Proposition~\ref{C-361-10}(d). (d) is Proposition~\ref{C-361-19}(a), (b) with $A=l_1/a_K$, $B=l_2/a_K$ and $t/a_K$ for $t$: $\frac A{2(t/a_K)}=\frac{l_1}{2t}$, $\frac{AB}{4l_K^2}=\frac{l_1l_2}{\ell_0^2}$, $\theta^{(t/a_K)/l_K}=e^{-m_*t}$, $\frac{1.4239}{A}=a_K\frac{1.4239}{l_1}$, and $\sum_{i\ge1}(\lambda_i/\lambda_0)^{t/a_K}=e^{\mathsf G}-1$; at $l_1=\ell_0$, $(2.42x_0)^{1/l_K}=\exp[-a_K\frac2{\ell_0}\log\frac1{2.42x_0}]$, and $2\log(1/(2.42\cdot0.069))\approx3.5798$.

For the square, $n=s_\ell/a_K=2^kl_K$ and $(\lambda_i/\lambda_0)^n\le i^{-1}\mathsf r_n\le0.20395\,i^{-1}\theta^{n/l_K}$, which is the bound on $E_i$ and on their number; $\log(1/0.20395)\ge1.5898$. The vortex line: Proposition~\ref{C-361-19}(d) and Proposition~\ref{C-361-11}(e), at $\ell=2^j\ell_0$, $\ell/a_K=2^{j+1}l_K$, $\sigma^+a_K^{-2}=\sigma_*$. (e) is Proposition~\ref{C-361-10}(f) at $n_0=0$; Lemma~\ref{C-361-20}(b) with $\theta^{(j-1)\mathsf w/l_K}=e^{-m_*(j-1)\ell''/2}$ and $Q''^2\le1.0350$; and Proposition~\ref{C-361-10}(b$'$), $\theta^{(2n-1)/6l_K}=e^{-m_*(2n-1)a_K/6}$, $n\ge\frac12(3l_K+1)$ being $na_K\ge\frac34\ell_0+\frac12a_K$. (f) With the time $e_3$, the lattice $\Gamma$ is the box $(\ell/2a_K;2\ell/a_K,\ell/a_K)$ closed up with the translation by $\frac\ell2e_1$ (its periods are also $2\ell e_1$, $\ell e_2$, $\frac\ell2(e_1+e_3)$).

With $l'=\ell/2a_K$, Proposition~\ref{C-361-7}(f) over the cross-section $(4l',2l')$ gives $\lambda_0^{l'}(1-\mathsf r)\le Z_\Gamma\le Z(l',4l',2l')$, $\mathsf r=\mathsf r_{l'}(2l',4l')$. Above: $F(l',4l',2l')\le2F(l',2l',2l')$ (Proposition~\ref{C-361-7}(b)), so $E\le\mathsf e(l')$. Below: $\log\lambda_0=8l'^2P+\varepsilon(2l',4l')\ge8l'^2P$ and $\log Z((2l')^3)=8l'^3P+\delta_K(\ell)$, so $E\ge\log(1-\mathsf r)-\delta_K(\ell)$; by Proposition~\ref{C-361-19}(b) in the zone $2t\le A$, $\mathsf r\le e^c-1$ with $c=0.54526\,y$, $y=\theta^{\ell/\ell_0}$; as $c\mapsto\log(2-e^c)$ is concave and vanishes at $0$, $\log(2-e^c)\ge-2.0196\,c$ for $0\le c\le0.54526\cdot0.8901$; with $\delta_K(\ell)<0.13736\,y$ the sum is $\ge-1.24\,y$.

The sign $E\ge0$ is Proposition~\ref{C-361-7}(g) at $N=\ell/a_K$, the lattice $\Gamma$ being $\Gamma^\diamond$ in physical units.

\emph{Step 2: Theorem~\ref{C-361-3}(2).} Fix $K\ge\max\{K_0,1\}$, $\beta=L^K\mathsf b$, boards of the offset $0$. The $X_i$ and $G_I$ are real functions of the cell $[-\ell''/4,\ell''/4]^3$ at every cutoff. \emph{(i)} $M_I<\infty$: $\|G_I\|^8_{\mathrm{ch},\ell''}$ is a moment of $8|I|$ centred fields with disjoint supports on the single torus $\ell''$, a polynomial in their cumulants of orders $\ge2$, which converge as $K\to\infty$ by (M4)(i) (Proposition~\ref{C-361-30}). \emph{(ii) One factor.} With $y=\theta^{\ell'/\ell_0}$, by Step~1, $[Q'(Q'-1)]^{1/2}\le0.132734\sqrt y$ and $Q''\le1.0174$; Proposition~\ref{C-361-9}(a), (c) give $|\nu_K(G_I)|,|\mu_{K,\ell',\beta}(G_I)|\le1.0174\,M_I$ and $|\nu_K(G_I)-\mu_{K,\ell',\beta}(G_I)|\le0.65\,\sqrt y\,1.0174\,M_I$. \emph{(iii) The cumulant.} A cumulant of order $n\ge2$ does not see the centring, nor, for $n=1$, does the difference of the means of two laws; and
\[
\kappa_n=\sum_\pi(-1)^{|\pi|-1}(|\pi|-1)!\prod_{I\in\pi}\mathbb E\Bigl(\prod_{i\in I}X_i\Bigr),\qquad\Bigl|\prod_I\mathsf a_I-\prod_I\mathsf b_I\Bigr|\le\sum_I|\mathsf a_I-\mathsf b_I|\prod_{I'\ne I}\max\{|\mathsf a_{I'}|,|\mathsf b_{I'}|\},
\]
so the bound holds with $C(X)$. \emph{(iv) Limits.} As the torus cumulants converge at each fixed $\ell'$ by (M4)(i), $\limsup_K-\liminf_K$ of $S^T_n[\nu_K](X)$ is at most $2C(X)\theta^{\ell'/2\ell_0}$ for every $\ell'$, hence $0$; letting $K\to\infty$ in the bound gives the iterated limit, and along $(K_j,\ell_j)$ the difference is $\le C(X)\theta^{\ell_j/2\ell_0}$. \emph{(v) Every shape.} At $N'=\ell'/a_K$ the corner $c'=\frac{\ell'}4(1,1,1)$ has $\mathsf C_{N'/2}(c')=[-\frac{\ell'}4,\frac{\ell'}4]^3$, and $\|\cdot\|_{\mathrm{ch};N'}=\|\cdot\|_{\mathrm{ch},\ell'}$ there; by Lemma~\ref{C-361-21}(b), (c) with $\delta(N'^3)<0.13736\,y\le\frac12$: $|\langle G_I\rangle_{\underline n}|\le e^{0.1717}\,1.0174\,M_I\le1.208\,M_I$ and $|\langle G_I\rangle_{\underline n}-\mu_{K,\ell',\beta}(G_I)|\le7.31\cdot0.370622\sqrt y\,1.0174\,M_I$; with (ii), $|\langle G_I\rangle_{\underline n}-\nu_K(G_I)|\le3.36\sqrt y\,1.0174\,M_I$, and (iii) gives $C^{\rm sh}(X)$.

For boxes of any shape with $a_{K_j}\min_\mu n^{(j)}_\mu\to\infty$ take $\ell_j$ the greatest member of $\mathfrak L$ below that length. \emph{(vi) A constant free of label, side and place.} For $\mathsf b\in I(g)$ the pairs are first-passage data at $g$. Put the supports by a corner of the cell, $c_\mu=\lfloor z_\mu+R\rfloor+\frac32$, $z$ the centre of the ball: one corner serves all cutoffs because $L$ is odd; the cell $c+[-\frac{\ell'}2,0]^3$ holds the supports farther than $\frac12$ from its walls; and the $8n$ mirror images of the $f_i$ in the planes $x_\mu=c_\mu$ have the same norms, supports each in a ball of radius $R$, pairwise $\ge\mathsf d'$ apart (two images on the two sides of a wall are more than $1$ apart), all in $c+[-(2R+\frac32),2R+\frac32]^3$, which a lattice translation carries into $(-\frac{\ell'}4,\frac{\ell'}4)^3$.

So the expansion of Proposition~\ref{C-361-30}, with the bounds (M5)(i) in the place of (M4)(i), gives $\|G_I\|_{\mathrm{ch},\ell'}\le\mathsf m^\flat_I$, centring by the mean on the torus $\ell'$ itself (the floor $128n(R+1)$ is the floor $16k(R+1)$ of (M5)(i) at $k=8n$). Then (ii), (iii) at $\ell''=\ell'$, where $4.84\cdot0.132734=0.642433$ and no factor $Q''$ enters, give $C^\flat(X)$, for $\nu_K$ and for a torus $\ell^+\ge\ell'$ alike, both being states $\mathfrak s$ of Proposition~\ref{C-361-9}; for the boxes $(q_\mu N)$, Lemma~\ref{C-361-21}(b), (c) give $1.86\cdot0.370622=0.68936$ in the place of $0.642433$; for every box, Lemma~\ref{C-361-21}(b), (c) and (ii) give $|\langle G_I\rangle_{\underline n}-\nu_K(G_I)|\le(2.70925+0.642433)\sqrt y\,\mathsf m^\flat_I$ and $e^{\frac54\cdot0.13736}\le1.19$.

For $n\ge2$, $|S_n(X)|\le C^g_n\prod\|f_i\|_\sharp$ is the limit of (M5)(i) along the tori.

\emph{Step 3: Theorem~\ref{C-361-3}(2$'$).} For bounded $Y,X_1,\dots,X_n$, with $X_I=\prod_{i\in I}X_i$,
\[
\begin{multlined}
\kappa(Y,X_1,\dots,X_n)=\sum_{\emptyset\ne I\subset[n]}\operatorname{Cov}(Y,X_I)\,\mathsf c([n]\setminus I),\qquad\mathsf c(S)=\sum_{\pi\in\mathrm{Part}(S)}(-1)^{|\pi|}|\pi|!\prod_{B\in\pi}\mathbb E[X_B], \\
\mathsf c(\emptyset)=1 .
\end{multlined}
\]
Take a board of side $\ell''/2$ with a cell $\square$ about the tuple and a cell $\square'$ about the support of $h$; centring by the means on the torus $\ell''$, Proposition~\ref{C-361-30} bounds the seminorms of the products of the centred fields in each cell, Proposition~\ref{C-361-9}(a) bounds the means $\mathbb E[X_B]$, and Lemma~\ref{C-361-20}(b) bounds $\operatorname{Cov}(Y,X_I)$ by $e^{-m_*(j-1)\ell''/2}$ times the seminorms, the indices of $\square$, $\square'$ differing by $j\ge2(|v|_\infty-\ell'')/\ell''+1$ in some direction; on a torus both ways around enter. The constants depend on the label only through the compact set of Proposition~\ref{C-361-30}. For $W$ in the place of $h$ the same argument runs with $Y=W$, whose centred seminorm in its cell takes the place of the bound of Proposition~\ref{C-361-30} for $h$. Summing over the plaquettes of a far region with a partition of unity by translates of a fixed bump gives the absolute convergence and the decay: with the tents $\chi_v$, $v\in\rho\mathbb Z^3$, of (2$'$) and $u$ vanishing on $\{|x|_\infty<r\}$, $r\ge\ell''+\rho$, the bound for $h=u\chi_v$, which vanishes unless $|v|_\infty\ge r-\rho\ge\ell''$, summed over $v$, gives
\[
\sum_{v\in\rho\mathbb Z^3}\bigl|S^T_{n+1}[\nu_K](u\chi_v,X)\bigr|\le\mathsf C_n\,\sup_v\|u\chi_v\|_\sharp\prod_i\|f_i\|_\sharp\sum_{v\in\rho\mathbb Z^3,\ |v|_\infty\ge r-\rho}e^{-m_*(|v|_\infty-\ell'')},
\]
On the torus use $V$ and $\chi_v^{\rm per}$ from (2$'$). The conditions $\ell\ge L\ell''$ and $\ell''\ge8\rho+4\mathsf d+4$ give $2\rho<\ell$, so each periodized tent is a single local tent in its torus chart and $\sum_{v\in V}\chi_v^{\rm per}=1$. Choose half-open representatives of $V$ with minimal sup norm. They form a subset of $\rho\mathbb Z^3$ and satisfy $|v|_\infty=|v|_{\infty,\ell}$. The one-far-insertion bound summed over this finite set, for $r\ge\frac32\ell''+\rho$, is therefore bounded by the displayed infinite-grid tail, with the same $3.44\,\mathsf C_n$.

\emph{Step 4: Theorem~\ref{C-361-3}(3), the vacuum and the Euclidean group.} For $\mathsf b\in I(g)$ all pairs $(K,L^K\mathsf b)$ are first-passage data, and by Step~1 the only periodic Gibbs state there is $\nu_K$. The moments of the centred fields under $\nu_K$ are polynomials in the cumulants $S^T_{|I|}[\nu_K]$, $|I|\ge2$, which converge by Step~2; so $(\nu_K)$ converges, and $S^\sigma_n=S_n$. By (M5)(ii) each ergodic component of $\sigma$ is a limit along the full sequence whose approximants are periodic Gibbs states at the same pairs, that is $\nu_K$; so $\sigma$ is its own component, and Theorem~\YMMref{AP}{780} gives the invariance under $\mathbb R^3\rtimes\mathrm O(3)$, reflection positivity at every plane, $\ker H_\sigma=\mathbb C\Omega_\sigma$ and the Wightman theory with a unique vacuum. Step~2 (iv), with its iterated limit and its bound along every pair of sequences $K_j\to\infty$, $\mathfrak L\ni\ell_j\to\infty$, identifies every continuum limit at $g$ along growing tori with the pairs $(K_j,L^{K_j}\mathsf b)$ with $(S^\sigma_n)_{n\ge2}$, the limits with the spacing removed first among them; so each of them has the Schwinger functions of $\sigma$, and with them its invariance under $\mathbb R^3\rtimes\mathrm O(3)$, its reflection positivity, its reconstruction with a unique vacuum and, by Step~5, its gap.

\emph{Step 5: Theorem~\ref{C-361-3}(3), the gap.} Let $F=\sum_i\zeta_iX_i$ be a finite combination of separated tuples of fields supported in $\{\epsilon<x_3<\mathsf h'\}$, $\Theta$ the reflection $x_3\mapsto-x_3$, $\tau_t$ the translation by $te_3$ with $t\in a_K\mathbb N$, and $F_K$ the same combination of the lattice fields $\mathcal O_K(f)$, centred by their means in $\nu_K$, a function of one slab above the plane $x_3=0$, with $n_0$ layers between the slab and its mirror image. By translation invariance $\nu_K(\overline{\Theta F_K}\,\tau_tF_K)$ is $h_{F_KF_K}(n_0+t/a_K)$, the mirror being moved to the plane $x_3=-t/2$; so Proposition~\ref{C-361-10}(f) with $A=B=F_K$, $n=n_0+t/a_K$ and $l=l_K$ gives
\[
\nu_K\bigl(\overline{\Theta F_K}\,\tau_tF_K\bigr)-|\nu_K(F_K)|^2\le e^{-m_*t}\bigl[\nu_K\bigl(\overline{\Theta F_K}\,F_K\bigr)-|\nu_K(F_K)|^2\bigr],
\]
with a rate that does not depend on $K$ and a prefactor that converges. Fix $t\in a_{K_*}\mathbb N$. It belongs to every finer time grid, so we may let $K\to\infty$ at this fixed time. With $\psi=(1-P_\Omega)[F]\in\mathcal H(\sigma)$: $\langle\psi,e^{-tH_\sigma}\psi\rangle\le e^{-m_*t}\|\psi\|^2$ on the unbounded set $\bigcup_Ka_K\mathbb N$ of times; so the spectral measure of $\psi$ gives mass at most $e^{-\eta t}\|\psi\|^2$ to $[0,m_*-\eta]$ for arbitrarily large $t$, that is none, and such $\psi$ are dense in $\Omega_\sigma^\perp$. The bilinear bound is the spectral theorem on $\Omega_\sigma^\perp$; the other directions follow by the invariance.

\emph{Step 6: Theorem~\ref{C-361-3}(4).} Proposition~\ref{C-361-11} at $l=l_K$, lengths in units of $a_K$: $\sigma^+a_K^{-2}=(2\Lambda+1.67)/4l_K^2a_K^2=\sigma_*$, and $\sigma_*\ge(2\cdot0.116421+1.67)/\ell_0^2$ at $\theta\le0.8901$; the areas become physical areas. For the flux quotient, $N/l_K=2\ell/\ell_0$ turns $2.2\cdot10^{-4}(N/l_K)^2$ into $8.8\cdot10^{-4}\ell^2/\ell_0^2$, $1.4268\cdot10^{-3}$ into $5.707\cdot10^{-3}$ at $N\ge256l_K$, that is $\ell\ge128\ell_0$, and $2\mathsf qN/l_K$ into $4\mathsf q\ell/\ell_0$. For the potential, $V_K(r)=a_K^{-1}V(r/a_K)$ with the $V$ of Proposition~\ref{C-361-11}(d), whose slopes are $\ge\sigma^+=\sigma_*a_K^2$.

\emph{Step 7: Theorem~\ref{C-361-4}(1).} By Lemma~\ref{C-361-28} take a second admitted blocking factor $L'>L$, a prime with $\alpha=\log L/\log L'$ irrational; Theorem~\YMMref{D}{8} holds in the scheme of $L'$ (spacings $L'^{-k}$, sides $2L'^{\mathbb N}$), and $L^K\mathsf b=L'^{k(K)}\mathsf c_K$ with the labels $\mathsf c_K\in[\mathsf b,L'\mathsf b)$ dense. In the scheme of $L'$ let $\ell_f\in2L'^{\mathbb N}$, $\ell_f\ge24L'\ell_0$, $w_0=\ell_f/2L'$, $f(x)=\prod_i\max\{0,1-|x_i|/w_0\}$ (the tent), and $\mathsf M_k(\ell;c)=\|\mathcal O_k(f)-\mu_{k,\ell,L'^kc}(\mathcal O_k(f))\|_{\mathrm{ch},\ell}$. (i) $\mathsf M_k(\ell;\cdot)^8$ is a polynomial in torus cumulants of orders $2$ to $8$, which by (M4)(i) at $L'$ converges, uniformly on $[\mathsf b,L'\mathsf b]$, to a continuous function; so $\mathsf M_f=\sup_{k\in\mathbb N,\ c\in[\mathsf b,L'\mathsf b]}\mathsf M_k(\ell_f;c)<\infty$. (ii) At the labels $\mathsf c_K$, BI gives $\mathsf e(l_K)\le x_0$ at $\beta=L'^{k(K)}\mathsf c_K$ with $l_K=\ell_0L^K/2$, and a fine cell has $\mathsf h=\ell_fL'^{k(K)}/2$ layers with $\mathsf h/l_K\ge\ell_f/L'\ell_0\ge24$; Lemma~\ref{C-361-12} with $\ell_c=d\ell_f$, $d=L'^j$, gives
\[
\mathsf M_{k(K)}(\ell_c;\mathsf c_K)\le\epsilon_\star(d)=1.343\,\mathsf M_f\,\theta^{(d-1)\ell_f/8L'\ell_0}.
\]
(iii) Let $x(c)=L'^{-3}D_{L'}(N)$, $N=\ell_cL'^k$, at $\beta=L'^kc$. Since $\partial_\beta\log Z(n^3)$ is $-3n^3$ times the mean plaquette variable, which is the same for all plaquettes of a torus, and the translates of the tent are a partition of unity,
\[
x'(c)=-(\ell_c/w_0)^3\bigl[\mu_{k,\ell_c,\beta}(\mathcal O_k(f))-\mu_{k,L'\ell_c,\beta}(\mathcal O_k(f))\bigr],\qquad\ell_c/w_0=2L'd .
\]
Wherever $x\le0.0064$, $\delta(N^3)\le(L'/(L'-1))^3x\le0.0086$ by Theorem~\ref{C-361-6}(1), and Proposition~\ref{C-361-9}(c) bounds the bracket by $0.16\,\mathsf M_k(\ell_c;c)$; so $|x'|\le0.0058/((L'-1)\mathsf b)$ there as soon as $\mathsf M_k(\ell_c;c)\le\epsilon_{\rm req}(d)=0.036/[(2L'd)^3(L'-1)\mathsf b]$. The required smallness shrinks with $d$ like $d^{-3}$, while $\epsilon_\star(d)/\epsilon_{\rm req}(d)$ is a constant times $d^3\theta^{(d-1)\ell_f/8L'\ell_0}$, which tends to $0$ as $d=L'^j\to\infty$: fix $d$ with $\epsilon_\star(d)\le\frac12\epsilon_{\rm req}(d)$. By (i) the functions $\mathsf M_k(\ell_c;\cdot)$ converge uniformly on $[\mathsf b,L'\mathsf b]$ to a continuous limit; at the labels $\mathsf c_K$, $K\ge K_0$, which are dense, with $k(K)\to\infty$, they are at most $\frac12\epsilon_{\rm req}(d)$; so the limit is at most $\frac12\epsilon_{\rm req}(d)$ everywhere, and there is $\bar k\ge k(K_0)$ with $\mathsf M_k(\ell_c;c)\le\epsilon_{\rm req}(d)$ for all $c\in[\mathsf b,L'\mathsf b]$ and all $k\ge\bar k$. (iv) At each $k\ge\bar k$ there is $K$ with $k(K)=k$, and Proposition~\ref{C-361-8}(b) gives $x(\mathsf c_K)<0.000514$; by the mean value theorem the component of $\{x\le0.0064\}$ that contains $\mathsf c_K$ carries $x\le0.006314$, so it is all of $[\mathsf b,L'\mathsf b]$.

Hence $\delta(N^3)\le0.008625$ for all $c$ and all $k\ge\bar k$, and every $\beta\ge\beta_0=L'^{\bar k}\mathsf b$ equals $L'^kc$ with $l=N\le C_+\beta$, $C_+=\ell_c/\mathsf b$. Of $L$ the argument uses only that $\log L/\log L'$ is irrational.

\emph{Step 8: Theorem~\ref{C-361-4}(2), (2$'$), (3).} At $\beta=L^K\mathsf b'\ge\beta_0$ there is $l\le C_+\beta$ with $\mathsf e(l)\le0.069$, and Proposition~\ref{C-361-8}(b) gives $\delta(N'^3)\le0.13736\cdot0.8901^{24}<0.008625$, so $\mathsf e(N')\le0.069$, for all $N'\ge48l$. Since $L^{m+K}\ge48C_+\mathsf b'L^K\ge48l$, the cube of BI at the side $2L^m$ qualifies; on $[\mathsf b_-,\mathsf b_+]$ take $L^m\ge48C_+\mathsf b_+$ and the threshold of $\mathsf b_-$. If BI holds at one label, it holds at all: the set of the labels is empty or $(0,\infty)$. (2$'$): $\mathsf M<\infty$ by Proposition~\ref{C-361-30} and (M4)(i). Fix $K\ge\max\{K_0,1\}$ and write $\delta(s)$ for $\delta_K(\ell_1)$ at the label $s$. The derivative of the excess in the label is $L^K$ times the sum over the plaquettes of the torus of the differences between the mean plaquette variable on the torus and in the periodic Gibbs state at the label $s$ (Proposition~\ref{C-361-10}(a)); all plaquettes having one such difference, by the symmetries of both laws, it is a multiple $R_K$, $|R_K|\le1458$, of the difference of the means of $\mathcal O_K(f)$ on the torus and in that state:
\[
\delta(\mathsf b_2)-\delta(\mathsf b_1)=-\int_{\mathsf b_1}^{\mathsf b_2}R_K\bigl[\mu_{K,\ell_1,L^Ks}(\mathcal O_K(f))-\nu_{K,s}(\mathcal O_K(f))\bigr]\,ds .
\]
While $\delta(s)\le0.008625$, Proposition~\ref{C-361-9}(c) bounds the bracket by $0.17516\,\mathsf M$, with constants free of $K$. At $s=\mathsf b$, $\delta(\mathsf b)<0.13736\,\theta^{\ell_1/\ell_0}\le0.13736\cdot0.8901^{75}<0.000025$; so $\delta$ cannot reach $0.008625$ within $\eta$ of $\mathsf b$, at every cutoff. The continuity and the limits along labels follow from Step~2, made with a supremum over the labels of $[\mathsf b-\eta,\mathsf b+\eta]$ (one side $2\ell_1$, one threshold, $\theta=0.8901$), and the continuity in the label of (M4)(i). For the second form put $N=\ell_1^\flat/a_K$. At $\mathsf b$, Proposition~\ref{C-361-8}(b) gives $D_L(N)<L^3\,0.13736\,\theta^{\ell_1^\flat/\ell_0}\le0.03$.

At a label $s$ with $D_L(N)\le0.069$, $\delta(N^3)\le0.069/(L-1)^3$ by Theorem~\ref{C-361-6}(1), and the derivative of $D_L(N)=L^3F(N^3)-F((LN)^3)$ in the label is $-L^3R_K$ times the difference of the means of $\mathcal O_K(f^\flat)$ on the tori $\ell_1^\flat$ and $L\ell_1^\flat$, with the number $R_K$ of the first form, now for $f^\flat$, $|R_K|\le1458$; Proposition~\ref{C-361-9}(c) bounds that difference by $4.84\,[Q'(Q'-1)]^{1/2}\mathsf M^\flat\le0.010813\,\mathsf M^\flat$. So $D_L(N)$ rises by at most $0.03$ within $r$ of $\mathsf b$ and stays below $0.06$; then $\mathsf e(N)\le\frac8{L^3}D_L(N)<0.48/L^3<2.2\cdot10^{-4}$ by Theorem~\ref{C-361-6}(1), and Proposition~\ref{C-361-8}(b) at $l=N$, with $\theta_1=12.9\cdot2.2\cdot10^{-4}$, gives $\mathsf e(L^zN/2)<0.3702\,\theta_1^{L^z/4}\le6.1\cdot10^{-9}$, which is the inequality of BI at the side $L^z\ell_1^\flat$. (3): Proposition~\ref{C-361-10} at this $l$: the rate $\theta^{n/l}$ with $\theta=0.8901$ and $l\le C_+\beta$ is at most $e^{-c_*n/\beta}$; the transfer operator is the operator $T_\nu$ of the setting, whose matrix elements are the numbers $h_{AB}(n)$ of Proposition~\ref{C-361-10}(f): at $\theta=0.8901$ and $l\le C_+\beta$, $\ker(1-T_\nu)=\mathbb C\Omega$ and $\|T_\nu\|\le0.8901^{1/l}\le e^{-c_*/\beta}$ on $\Omega^\perp$; and Proposition~\ref{C-361-11}(a) with $l\le C_+\beta$ gives the flux quotient; Proposition~\ref{C-361-11}(d) gives the slopes, $\sigma^+l^2=\frac12\log\frac1{0.8901}+0.4175\ge0.4757$.

The labels $a_K\beta$ of all first-passage data $(K,\beta)$ at $g$ lie in the window $\bigl(g^{-2}+B_\sharp\frac{L-2}{L-1},\,L(g^{-2}+\bar B)\bigr]$ (\YMMref{(1531)}{615} in Proposition~\YMMref{74.1}{615}), so at a datum with $\beta\ge\beta_0$ the cube of the side $2la_K\le2C_+\beta a_K\le2C_+L(g^{-2}+\bar B)=\ell_0^g$ satisfies the inequality: Theorem~\ref{C-361-3}(1) and Step~2 hold there with $\ell_0^g$ for $\ell_0$ and $0.8901$ for $\theta$, uniformly over these data, the seminorms of Proposition~\ref{C-361-30} being bounded over the window.

\emph{Step 9: Theorem~\ref{C-361-4}(4).} On $a\mathbb Z^3$ the weight does not contain $a$, and $\mathcal O_a(f(\cdot/t))=t^{-1}\mathcal O_{a/t}(f)$; hence
\[
\kappa_n\bigl(\beta,a/t;X\bigr)=t^n\,\kappa_n\bigl(\beta,a;X_t\bigr)
\]
for the cumulants in the unique periodic Gibbs state at $\beta\ge\beta_0$. By Step~8, BI holds near every label with one side and one threshold, and Step~2, taken with a supremum over a compact set of labels and of dilations of $X$, gives one limit along every sequence of pairs whose labels converge, continuous in the label and the dilation. The supremum over the dilations is had through the rate of (M4)(i): by~the Cauchy estimate in the proof of Theorem~\YMMref{AS}{813} in Theorem~\textup{\YMMref{AS}{813}}, on one torus and on a compact interval of labels, $|S^T_{n;K',\ell,L^{K'}\mathsf b}(Y)-S^T_{n;K,\ell,L^K\mathsf b}(Y)|\le C\,K^{2\mathfrak p_0(n+1)+\frac12}L^{-K/10}\prod_i\|y_i\|_\sharp$ for all $K'\ge K$ above a threshold, where $C$ and the threshold depend on $n$, the side, the interval and a floor of the spread of the supports of $Y=(y_1,\dots,y_n)$, and on $Y$ otherwise only through these norms.

For $Y=X_{t'}$ with $t'$ in a compact interval $J''\subset(0,\infty)$, 
\[
\|f_i(\cdot/t')\|_\sharp\le\max\{\sup|f_i|,40M(\inf J'')^{-1}\operatorname{Lip}f_i\}
\]
 and the spread is at least $\inf J''$ times that of $X$; so the torus cumulants converge uniformly in the label and the dilation. At a fixed cutoff they are continuous in both, being finite sums with the coefficients $f_i(x(p)/t')$; so the limit is jointly continuous, and the seminorms $M_I$ of Step~2 are bounded uniformly in $t'\in J''$ in the same way. In the scheme of $L'$ the lattices of the line of $\mathsf b$ are the lattices of the dense labels $\mathsf c_K$, dilated by $\mathsf b/\mathsf c_K$; so the limit of the scheme of $L'$ at every label $c$ is the dilation of $\sigma_{\mathsf b}$ by $c/\mathsf b$.

The dilations by $L$ and by $L'$, which are trivial in their own schemes, generate a dense subgroup of $(0,\infty)$, and by continuity $\sigma_{tc}$ is the dilation of $\sigma_c$ for every real $t$; the identity above reduces arbitrary real spacings to the scheme of $L$, and $H_n$ is $S^{\sigma_1}_n$, so $S^{\sigma_c}_n(X)=c^{-n}H_n(X_{1/c})$. At $c\in I(g)$, Step~4 identifies $\sigma_c$; the gap of $\sigma_{tc}$ is that of $\sigma_c$ divided by $t$.

\emph{Step 10: Theorem~\ref{C-361-4}(5).} For $\varphi$ vanishing outside $v+[-\rho,\rho]^3$ the insertion of the rotation weighted by $\varphi$ depends only on the bonds within $2a_K$ of that cube and is bounded by $3402(1+|v|_\infty)\rho^4L^{4K}|\omega|\sup|\varphi|$, and $\sup|\mathcal O_K(f)|\le162\rho^3L^{4K}\sup|f|$ for $f$ vanishing outside $v+[-\rho,\rho]^3$; so in every state, for $\chi$ vanishing outside $v+[-\rho,\rho]^3$ and the $f_i$ outside $v+[-R,R]^3$, $|\mathfrak U|\le\mathsf c_n(1+\mathsf b)(1+|v|_\infty)\rho^4R^{3n}L^{(4n+5)K}|\omega|\sup|\chi|\prod\sup|f_i|$, with $\mathsf c_n=3402\cdot162^n\,\mathsf F_{n+1}$ and $\mathsf F_m=\sum_\pi|\pi|!$, the sum over the partitions $\pi$ of $\{1,\dots,m\}$; in the theorem $\rho=R$. (i) \emph{Torus against state.} The bulk term is a polynomial in means of products of local functions of one cube; by Theorem~\ref{C-361-3}(1)(c), in total variation, on all tori $\ell\ge\ell''$ and in $\nu_K$ it agrees up to $2.164\,\mathsf c_n(1+\mathsf b)(1+|v|_\infty)\rho^4R^{3n}L^{(4n+5)K}e^{-m_*\ell''/4}$ times the sup norms, which is $\le$ the same expression with $L^{-K}$ for $L^{(4n+5)K}e^{-m_*\ell''/4}$ once $\ell''\ge\Lambda_0K$, $\Lambda_0=4(4n+6)\log L/m_*$. (ii) \emph{The place.} For a state invariant under $a_K\mathbb Z^3$, translating the tuple by $v$ adds the translation $\omega v=\sum_\lambda v_\lambda\,\omega e_\lambda$ to the generator of the insertion, and the insertion of the translation by the vector $\omega e_\lambda$ is the difference of the bulk terms of the tuple moved by $e_\lambda$ and of the tuple, both carried by $[-R-1,R+1]^3$; summed over the components of $v$ this gives the factor $1+2|v|_1$. (iii) So it suffices to bound the bulk term of tuples carried by a fixed cube about the origin on one side in $[\Lambda_0K,L\Lambda_0K]$ of $\mathfrak L(g)$, which by (i) serves for $\nu_K$ and all larger sides, and on the sides $\ell\le L\Lambda_0K$ a power of $\ell$ is a power of $K$.

There $\mathfrak U$ of~\YMMref{(2108)}{782} is split as follows, with the objects of (M8). A joint cumulant with the $\mathcal O_K(f_i)$ is the derivative in the $z_i$ of a mean at the source $w_z=\sum_iz_if_i$, so the bulk term of~\YMMref{(2107)}{782} in equation~\YMMref{(2107)}{782} and the source calculation in the proof of Proposition~\YMMref{112.1}{782} is $-\partial^n_z$ at $z=0$ of the mean at the source $w_z$ of the bulk insertion $(\mathfrak S_{\omega,\chi}+\mathfrak K_{\omega,\chi})[\beta-a_K^{-1}w_z]$ (equation~\YMMref{(2107)}{782} and the source calculation in the proof of Proposition~\YMMref{112.1}{782}). With the orbital potential $\Psi_\rho(p)$ and the remainder of third order $\mathfrak r^\sharp_\rho(p)=\mathfrak r_\rho(p)-\Psi_\rho(p)+\Psi_\rho(p_-)$ of Section~\ref{sec:ward-orbital-complement}, the summation by parts of Proposition~\ref{prop:ward-orbital-cubic} writes this insertion as the sum of three members, affine in $z$: the spin part $\mathfrak S_\omega[\phi_z]$, $\phi_z=\beta\chi-a_K^{-1}w_z$; the member $-\sum_p\sum_\rho\tilde c_\rho(x(p))\Psi_\rho(p)$, $\tilde c_\rho(y)=(\omega y)_\rho\,a_K^{-1}[\phi_z(y+a_Ke_\rho)-\phi_z(y)]$, of order $2$ in the curvature with the weight $\beta$; and the cubic member $\mathfrak K^\sharp[\phi_z\,\omega y]$ of~\eqref{eq:ward-orbital-summation}, of order $3$ with the weight $\beta/a_K=\mathsf bL^{2K}$.

Accordingly $\mathfrak U=\mathfrak U^\sharp+(\mathfrak U-\mathfrak U^\sharp)$ with
\[
\mathfrak U^\sharp_{\omega,\chi}(f)=-\beta\bigl\langle\mathfrak K^\sharp[\chi\,\omega y];\mathcal O_K(f_1);\dots;\mathcal O_K(f_n)\bigr\rangle+a_K^{-1}\sum_{i=1}^n\bigl\langle\mathfrak K^\sharp[f_i\,\omega y];(\mathcal O_K(f_l))_{l\ne i}\bigr\rangle .
\]
\emph{All but the cubic member.} Use the exact lattice decomposition in Lemma~\ref{C-360-8}(b,c,e,f), as in Step~8 of the proof of Theorem~\ref{C-360-4}. For $\varphi\in\{\chi,f_1,\dots,f_n\}$, antisymmetry of $\omega$ rewrites $a_K^{-1}\mathfrak S_\omega[\varphi]$ as an order-$2$ entry with weights $k^\varphi_{b,p,\rho}=a_K^{-1}(\varphi(x(p))-\varphi(x_b))\sum_\lambda\omega_{\rho\lambda}a_K^{-1}(x_b-x(p))_\lambda$, bounded by $|\omega|\operatorname{Lip}\varphi/4$. The bilinear orbital entry has the vector weight $\mathsf k^{\varphi,\rho}(y)=(\omega y)_\rho\nabla^+_\rho\varphi(y)$, whose Lipschitz norm is at most $5|\omega|(1+\ell)\|\varphi\|_{C^{1,1}}$. Lemma~\ref{C-360-8}(c) gives $\mathfrak J^{\rm bl}[\mathsf k]=-\mathfrak N[\mathsf k]/4+\sum_{v=1}^5\mathfrak J_v[\mathsf k]-c_{\mathsf k}/4$: the five entries have orders $2,2,2,3,4$, and only the last two carry $\beta$. The null term is evaluated by the contact formulas (N3), (N4) of that lemma. They leave numbers, action densities and order-$2$ entries without a factor $\beta$; the numbers drop from the cumulants. Terms absorbing all the sources vanish, and merging the absorbed supports leaves spread at least $\mathsf d-3a_K\ge\mathsf d/2$. Thus Lemma~\ref{C-360-18}, applied in its spread form, bounds every remaining term.

For the dependence on $\ell$, use the proof of Lemma~\ref{C-360-18} before its constants are collected. A local insertion of order $m$ is carried to depth $u=\lceil\phi K\rceil$, $\phi=(\frac12-c_1)/(3n+4)$. The entry contributes at most a factor $\ell^3$ from its sites, and the source radii are $(2nN_t\mathsf C_t\|f_q\|_\sharp)^{-1}$, with $N_t=\ell^3L^{3u}$ and $\mathsf C_t\le CK^{2\mathfrak p_0}$. Their Cauchy cost contributes $\ell^{3n}$; the powers of $L$ are $L^{(3n+4)u-\mathsf g_mK}$, where $\mathsf g_2=1$ and $\mathsf g_m=3/2$ for $m\ge3$. The factors $\beta$ in the orders $3,4$ leave the decay $L^{-c_1K}$. The discarded mass of the definitions preceding Proposition~\YMMref{115.2}{806} has the form $2M^{-3}\ell^3L^{3u}\exp[-c\{\log(L^u\mathsf b')\}^{2\mathfrak p_0}+C_g\ell^3]$. Since $\mathfrak p_0\ge13$, the labels lie in a compact first-passage interval and $\ell\le L\Lambda_0K$, this mass and its fixed-order moment factors decay faster than every $L^{-AK}$, uniformly in these sides. The locality remainders have the same property. Collecting the weights from the preceding paragraph gives, times the norms, $|\mathfrak U-\mathfrak U^\sharp|\le C\ell^{3(n+1)}(R+5)|\omega|K^{(2n+1)\mathfrak p_0}L^{-c_1K}$ for every $0<c_1<\frac12$ on $2R+4\le\ell\le L\Lambda_0K$, uniformly over first-passage data at $g$, with $C$ free of the side.

\emph{The cubic member} obeys an estimate on all the sides at once, the bound (M7), which is Lemma~\ref{C-360-19} of the section on plane fields and the stress tensor (Section~\ref{C-360-1}) taken at the field $\chi\,\omega y$: for $n\ge2$, $\mathsf d>0$ and $\rho_0\ge1$ there are $C_\sharp$ and a threshold, depending on $g$, $n$, $\mathsf d$, $\rho_0$ and the constants and on no side, such that for all real $\chi,f_i\in C^{1,1}$ vanishing outside $[-\rho_0,\rho_0]^3$, $\chi$ covering the $f_i$, the supports of the $f_i$ at pairwise sup distances $\ge\mathsf d$,
\[
|\mathfrak U^\sharp_{\omega,\chi}(f)|\le C_\sharp|\omega|L^{-K/16}\|\chi\|_{C^{1,1}}\prod_i\|f_i\|_{C^{1,1}}
\]
at all first-passage data at $g$ with $K$ above the threshold, on all sides $\ell\ge2\rho_0+2$ of $\mathfrak L(g)$. It is proved in three steps, as in the idea of the proof of (M7). On one torus, the conditional mean of the cubic insertion given the block field of the level $K-\lceil\phi'K\rceil$, $\phi'>0$ small, under the small-field chain of Theorem~\YMMref{D}{8}, differs from a constant by at most a power of $K$ times $L^{-cK}$, for every $c<\frac1{14}$ (Lemma~\ref{C-360-9}(b)): the density is of third order, $|\mathfrak r^\sharp_{\rho(p)}(p)|\le C_{\rm orb}\,\bar a(p)^{3/2}$ (\eqref{eq:ward-orbital-real-cubic}), plain power counting against the weight $\mathsf bL^{2K}$ falls short by the factor $a_K^{-1/2}$, and the parity under the lattice mirrors, the covariance of the renormalisation step and a Hölder bound of a background against its mirror image win this factor back.

The joint cumulants of the insertion with the sources then obey the same bound, for every $c<\frac1{14}$, by the two-constants inequality of Lemma~\YMMref{116.2}{809} in the source variables; this is the bound of one torus, Lemma~\ref{C-360-10}. Last, the bound on one torus serves on the finitely many sides of $\mathfrak L(g)$ from $2\rho_0+2$ up to a fixed side $\ell'$, where (M7) is Lemma~\ref{C-360-10}, and it passes from $\ell'$ to every larger side (Lemma~\ref{C-360-19}) by the chessboard estimate (M2) on the board of side $\ell'/2$, which writes the difference of a cumulant on the torus $\ell$ and on the torus $\ell'$ through cumulants on the torus $\ell'$ of mirror copies of the entries that join cells, every block that holds a copy of the cubic insertion carrying the small factor, at the cost of the factor $e^{(\ell'/2)^3(p_{\ell'}-p_\ell)}\le e^{\delta_K(\ell')/8}$ of the fine-to-coarse form, which Theorem~\ref{C-361-6}(5)(a), (6) bound at the fixed side $\ell'$, uniformly in the cutoff and over the window of the labels at $g$.

With $\rho_0=R+1$ the sum is $\le C|\omega|K^{p'}L^{-\kappa K}$ times the norms for every $\kappa<\frac1{16}$. The condition $|v|_\infty+R+2\le\ell/4$ keeps the moved tuples and the insertion away from the seam of the torus, the planes $x_\mu=\pm\ell/2$ along which the cube $[-\ell/2,\ell/2]^3$ is closed up.

\emph{Step 11: Theorem~\ref{C-361-5}(1)--(4).} (1) is Proposition~\ref{C-361-13}(b), (c); the statements on a height $y$ follow from $|X_{K'}-X_K|\le\eta_K$ for $K'\ge K\ge K_\star$ and $\eta_K\to0$. (2) If $D_L(N)\le\frac{L^3}8x$ with $x\le0.069$, then $\mathsf e(N)\le D_2(N)\le\frac8{L^3}D_L(N)\le x$, because $D_L(N)-\frac{L^3}8D_2(N)=\frac{L^3}8\delta((2N)^3)-\delta((LN)^3)\ge0$ by the monotone density; Proposition~\ref{C-361-8} at $l=N$ gives $\mathsf e(2^kN)\le(12.9x)^{2^k}/12.9=\epsilon$, and since $2^kN\le LN/2<2^{k+1}N$,
\[
\mathsf e(LN/2)\le2\,\delta(\tfrac{LN}2,LN,LN)\le\Bigl(\frac L{2^{k+1}}\Bigr)^3\bigl[\epsilon+\tfrac{20}7\epsilon^2\bigr]=\Phi_L(x),\qquad D_L(2l)\le L^3\delta((2l)^3)\le\tfrac{20}7L^3\mathsf e(l)^2 .
\]
At $N=\ell/La_K$ we have $LN/2=\ell/2a_K$, and by (1) $X_K\le X_{K_1}+\eta_{K_1}\le y_1$ for all $K\ge K_1$; as $\Phi_L(x_1)\le x_0$, this gives BI$[\mathsf b,\ell,K_1;x_0]$; with $y_0$, $\mathsf e(\ell/La_K)\le x_0$, which is BI at the side $2\ell/L$ with $c=2$, and Theorem~\ref{C-361-3} holds at every side $2cL^{m_0}$. Conversely BI at the side $2\ell/L$ has $l_K=\ell/La_K$, and $X_K(\mathsf b,L\ell)=D_L(Ll_K)\le L^3\delta((Ll_K)^3)\le L^3\bar\delta_{x_0}(L)$ by the proof of Proposition~\ref{C-361-8}(b). The bracket: the upper members are Theorem~\ref{C-361-6}(1) at $l=\ell/2a_K$, $2l/L=\ell/La_K$; the lower member is the contrapositive of $D_L(2l)\le\frac{20}7L^3\mathsf e(l)^2$ for $\mathsf e(l)\le0.069$, at $2l=\ell/a_K$.

So BI gives $X_K(\mathsf b,L\ell)\le\frac{20}7L^3x_0^2$, and (1) excludes this for $K\ge K_1$ when $X_{K_1}(\mathsf b,L\ell)>\frac{20}7L^3x_0^2+\eta_{K_1}$; the members of the bracket are continuous and monotone in $X_K$, and $X_K\to X_\infty$, which gives the bracket for $\liminf_K\mathsf e_K$ and $\limsup_K\mathsf e_K$. (3) At the single cutoff $K_1$, Proposition~\ref{C-361-8}(b) at $l=l_{K_1}$ gives $X_{K_1}(\mathsf b,L^i\ell_0)=D_L(2L^{i-1}l_{K_1})\le L^3\bar\delta_{x_0}(2L^{i-1})$; with the margin, (1) gives $X_K\le y_1$ for all $K\ge K_1$, and (2) applies. At $i=2$, $\bar\delta_{x_0}(2L)<8\mathsf c_0\theta^{L}\le0.13736\,\theta^L$, which at $x_0=0.069$ and $L\ge33$ is at most $0.0030$; and for $L\ge33$, $2^k\ge16$, so $\Phi_L(0.05)\le8\cdot1.01\,(0.645)^{16}/12.9<0.069$ and $x_1\ge0.05$, $\frac18x_1\ge0.00625$: the margin condition holds once $\eta_{K_1}\le0.003L^3$.

At $i=1$ it never holds at $x_0=0.069$, as $\bar\delta_{0.069}(2)=\frac{20}7\cdot0.069^2>\frac18\cdot0.069$. (4) (iii) and (1) give $X_K\le y_0$ for $K\ge K_1$, hence $\mathsf e(\ell/La_K)\le x_0$, which is BI at the side $\ell_0=2\ell/L\in2\cdot2L^{\mathbb Z}$; (i) gives $X_K(\mathsf b,\ell)\le L^3\bar\delta_{x_0}(2\ell/L\ell_0)<0.13736L^3\theta^{\ell/L\ell_0}$ for $\ell\ge L\ell_0$ by Proposition~\ref{C-361-8}(b), hence (iv); (iv) implies (ii); and (ii) implies (iii) by (1).

\emph{Step 12: Theorem~\ref{C-361-5}(5), (6).} (5): under \textbf{(Hess)}, Proposition~\ref{C-361-14} at $\underline s=(\ell_0/2,\ell_0,\ell_0)$, $z=0$, $l=\ell_0$, where $|\underline s|/l^3=\frac12$, gives
\[
\mathsf e_{K+1}-\mathsf e_K=2\Bigl[\log Z_{K+1}(\underline s;\mathsf b)-\log Z_K(\underline s;\mathsf b)-\tfrac12\bigl(\log Z_{K+1,\ell_0}-\log Z_{K,\ell_0}\bigr)\Bigr],
\]
so $|\mathsf e_{K+1}-\mathsf e_K|\le2CL^{-K/10}$ uniformly on $J$, and Lemma~\ref{C-361-26} gives the rest. (6): $D_L(\ell/a_K)=X_K(\mathsf b,L\ell)$ converges uniformly on compact sets of labels by (1), with the rate of Proposition~\ref{C-361-13}(b); the pairs $(\ell/a_K,L^K\mathsf b)$ depend on $(\ell,\mathsf b)$ only through $\ell/\mathsf b$ up to the index of the cutoff, so $D_\infty$ is a continuous function of $t$. The series for $\delta_K(\ell)$ of Proposition~\ref{C-361-7}(c) at $q=L$, $\delta_K(\ell)=\sum_sL^{-3(s+1)}D_L(L^s\ell/a_K)$, has the majorant $L^{-3(s+1)}D_L(L^s\ell/a_K)\le L^{-3s}\delta_K(L^s\ell)\le C(t)\,L^{-s}$, uniformly in the large $K$, with $C(t)$ bounded on compact sets, by the bounds at a fixed cutoff of Theorem~\ref{C-361-6}(5)(a) and (6), whose proofs do not use the present clause; dominated convergence gives $\delta_K(\ell)\to\delta_\infty(t)$, locally uniformly, and $L^3\delta_\infty(t)-\delta_\infty(Lt)=D_\infty(t)$; the inequalities $(L-1)^3\delta\le D_L\le L^3\delta$ of Theorem~\ref{C-361-6}(1) pass to the limit.

\emph{One function for all block sizes.} Let $L'$ be a second admitted blocking factor with $\log L/\log L'$ irrational. At one coupling, cubes $N\le M$ with $M/N\le\rho$ have $\delta(M^3)\le\rho^3\delta(N^3)$ by the monotone density. Pass to a subsequence of the lattices $(N,\beta)$ that converge to $\delta_\infty^{(L)}(t)$, and choose along it cubes $M=2L'^{j}$ with $M/N\to\rho_0\in[1,\rho]$, which the density of $\{L'^j/L^i\}$ allows; the lattices $(M,\beta)$ are lattices of the scheme of $L'$ whose labels converge, so $\delta^{(L')}_\infty(\rho_0t)\le\rho^3\delta^{(L)}_\infty(t)$. Letting $\rho\to1$ and using the continuity, $\delta^{(L')}_\infty\le\delta^{(L)}_\infty$, and by symmetry the two are equal. The same comparison of nearly equal cubes, now with the lattice inequalities of Proposition~\ref{C-361-7}(c) and Theorem~\ref{C-361-6}(1),(2) between the cubes $N$ and $M=\lfloor\rho N\rfloor$, gives that $t^{-3}\delta_\infty(t)$ does not increase, the bound $(3\rho^2-3\rho+1)$ for real $\rho$, and the alternative for real $t\ge2t_0$.

\emph{Positivity.} $\delta_\infty\ge0$ from the series. At every $t$, choosing by a dilation a label in $I(g_1)\cup I(g_2)=(\gamma_W^{-2}+\bar B,\kappa_0]$ and the side $\ell=t\mathsf b$, Theorem~\ref{C-361-6}(1), (3) give $8\delta_K(\ell)\ge\mathsf e(\ell/a_K)\ge\min\{0.069,0.0775\,e^{-2.6\ell/\mathsf d_\circ}\}$ at all large $K$, so $\delta_\infty(t)>0$; this uses Theorems~\YMMref{AJ}{673} and~\YMMref{AD}{607} through (M5)(iii). \emph{The case of BI.} If BI holds at $(\mathsf b,\ell_0)$, then $\delta_\infty(\ell/\mathsf b)\le0.13736\,\theta^{\ell/\ell_0}$ for large $\ell$ by Theorem~\ref{C-361-3}(1); conversely, $\delta_\infty(t_0)<0.008625$ and $\ell/\mathsf b=t_0$ give $\delta_K(\ell)\le0.008625$ for all large $K$, hence $\mathsf e(\ell/a_K)\le0.069$: BI at the side $2\ell$.

If $D_{q,\infty}(t_0)<(q-1)^3x_0/8$ for an admitted $q$, then $\delta_\infty(t_0)<x_0/8\le0.008625$, as $(q-1)^3\delta_\infty\le D_{q,\infty}$. In the second case $\delta_\infty<0.008625$ on a whole interval $[t_2',Lt_2']$, and the uniform convergence on the compact set of labels $[\mathsf b,L\mathsf b]$ at one side gives, for every real $\beta$ above a threshold, a cube $N\le C\beta$ with $\delta(N^3)\le0.008625$: Theorem~\ref{C-361-4}(1) a second time.

\emph{Step 13: Theorem~\ref{C-361-5}(7).} The conditions of the list (H3), each attached to its constant, are inequalities between numbers; the one that binds $L$ alone, with its constant $4096$, is $L>48\cdot4096^2e^{24}\approx21\,331\,858\,733\,892\,671\,081.7\ldots$, and every odd integer above it is admitted. The list also forces floors on the other constants that enter $K_\star$ and $\eta_K$: $r\ge2$, $\mathfrak p_0\ge13$ (so $\mathsf p_0\ge26.5$), $M\ge L^7$, $\gamma^{-2}\ge(2L)^{1602}$, $A_0\ge256L^5A_1$, and three constants: $\mathsf c_{\rm a}=\mathsf C_{\rm ex}$, the expansion constant preceding equation~\YMMref{(705)}{360}, with $\mathsf C_{\rm ex}\ge8e^8$; $\mathsf c_{\rm c}=\kappa>0$, the constant of the list that enters $\lambda_{\rm out}=(1+4\beta_{\rm s})\kappa$ of~\YMMref{(889)}{415}; and $\mathsf c_{\rm b}=\lambda_1=\lambda_{\rm out}+2^3c_{\rm w}+\mathsf C_{\rm ex}+\lambda_{\rm s}$ of~\YMMref{(705)}{360}, where $c_{\rm w}=c_{\rm v}+3$, $c_{\rm v}=2+\mathsf C_{\rm ex}c_G(LM)^3u$ with $c_G=\mathsf C_{\rm ex}(\gamma_{\rm r}^{-1}+c_{\rm r}^2)\ge2\mathsf C_{\rm ex}$ and $(LM)^3u=\mathsf u^*_{\rm S}=3+c_{\rm r}+\mathsf u_{\rm G}>4$ by~\YMMref{(889)}{415}, so that $\mathsf c_{\rm b}>\mathsf c_{\rm c}+64\,\mathsf c_{\rm a}^2\ge\mathsf c_{\rm c}+4096\,e^{16}$, where $4096\,e^{16}\approx3.6397\cdot10^{10}$.

The threshold $K_\star$ is the least cutoff at which the conditions of (M4)(ii)--(iii) on the construction constants hold for both tori, with the thresholds $K_{\rm bf}$ of (M4)(ii) on both after the dilation by $z$, as in Theorem~\ref{C-361-5}(1), together with the comparison of Proposition~\ref{C-361-13}. At their floors these conditions give $(K-z)\log L>9\,\mathsf c_{\rm b}>36864\,e^{16}\approx3.2757\cdot10^{11}$ for every $K\ge K_\star(\mathsf b,\ell)$. Indeed, at the dilated data $(L^z\mathsf b,L^z\ell)$, whose cutoff $K-z$ is written $K$ in this sentence: the number $\kappa_{\rm b}$ of~\YMMref{(2181)}{803} is $w_\star e^{-\lambda_V}/32M^3$ with $w_\star=w_0e^{-4\lambda_1}/(\mathsf C_{\rm ex}(L+2)^3S_{\rm s})$, $w_0=(4S_{\rm s}L^3)^{-1}$ (equation~\YMMref{(705)}{360}), so $\kappa_{\rm b}^{-1}>L^6M^3e^{4\mathsf c_{\rm b}}$; the condition $\epsilon_{\rm sh}\le\frac1{20}$ of (k3) forces $\exp[\frac14\kappa_{\rm b}A_1^2\mathsf y_{[u]}]\ge60$, so $A_1^2\mathsf y_{[u]}>16/\kappa_{\rm b}$ and $A_0^2\mathsf y_{[u]}>L^{16}M^3e^{4\mathsf c_{\rm b}}$; then $\epsilon_{\rm cl}\le\frac1{20}$ forces $L^{(K-s)/2}\ge80A_0^2\mathsf y_{[u]}>e^{4\mathsf c_{\rm b}}$, $s=\lceil K/9\rceil$, and as $K-s\le\frac89K$, $K\log L>9\,\mathsf c_{\rm b}$.

So the small cube of the criterion of (2) has $2L^{m+K-1}\ge2L^{K-z+1}$ sites a side, as $m+z\ge2$, hence more than $e^{9\mathsf c_{\rm b}}>10^{1.4226\cdot10^{11}}$. Where the margin $\eta_K$ is below the height of the criterion, as the margin condition of (3) requires of $\eta_{K_1}(\mathsf b,L^2\ell_0)$, the same conditions, with $\mathsf C_{\rm bf}\ge64\mathsf K_{\rm w}c_\rho^2A_0^2\mathsf y_{[1]}L^{7/2}$ from~\textup{(\ref{IV-62-65})} and $\mathsf y_{[1]}(K-z)^{2\mathfrak p_0}\ge\mathsf y_{[u]}$, make $L^{(K-z)/10}$ larger than $e^{4\mathsf c_{\rm b}}\ge e^{4(\mathsf c_{\rm b}-\mathsf c_{\rm c})}$ times a power of $L$ and of $M$, and $e^{4(\mathsf c_{\rm b}-\mathsf c_{\rm c})}>e^{256\,\mathsf c_{\rm a}^2}\ge e^{16384\,e^{16}}$; so $(K-z)\log L>163840\,e^{16}\approx1.4559\cdot10^{12}$.

At the cutoff of (3) the half side $l=\ell_0/2a_{K_1}$ of the cube of BI is then $>10^{6.3228\cdot10^{11}}$, and the two boxes $(l,2l,2l)$ and $(2l)^3$, of $4l^3$ and $8l^3$ sites, have more than $10^{1.8968\cdot10^{12}}$ sites each. For the floor of $\eta_K$: $\varsigma\ge\varsigma_\flat(L)$ as $\mathsf p_0\ge\frac52$, $\mathsf C_{\rm bf}\ge2$, and $L^z\ell=2L^{m+z}\ge2L^2$ as $m+z\ge2$; so $\eta_K\ge3\varsigma_\flat\cdot2\cdot8L^6(K-z)^{5/2}L^{-(K-z)/10}$. An inequality $X_{K_1}\le y-\eta_{K_1}$ with $X_{K_1}\ge0$, $y\le y_0$, forces $\eta_{K_1}\le\frac{L^3}8\,0.069$, hence $K_1-z\ge K_\flat(L)$ for the least $K_\flat$ at which the floor drops below that number, and the small cube then has $2L^{m+K_1-1}\ge2L^{1+K_\flat}$ sites a side; at every admitted $L$ this exceeds $10^{152}$ (already at $L=13$, $2L^{1+K_\flat}>10^{152.9}$).

\emph{Step 14: Theorem~\ref{C-361-6}(1)--(4).} (1): $D_2(N)=4\mathsf d^{(3)}(N,N,N)+2\mathsf d^{(2)}(N,N,2N)+\mathsf d^{(1)}(N,2N,2N)$, three doubling defects $\ge0$ by Proposition~\ref{C-361-7}(b), the last being $\mathsf e(N)$; $D_2\le8\delta$ and $D_q-\frac{q^3}8D_2=\frac{q^3}8\delta((2N)^3)-\delta((qN)^3)\ge0$ by Proposition~\ref{C-361-7}(c); $(q-1)^3\delta(N^3)\le D_q(N)$ is $\delta((qN)^3)\le(3q^2-3q+1)\delta(N^3)$, and the box and cube bounds are Proposition~\ref{C-361-7}(c); $\mathsf e(l)\le2\delta(l,2l,2l)\le16(l/N)^2\delta(N^3)$. For $l=Ll'$, $n=2l/L$, with the time along a side of the cube $Ln=2l$: $D_L(n)\ge LF(n,2l,2l)-F(2l,2l,2l)=L\mathsf G(n)-\mathsf G(2l)\ge L\mathsf G(l)-\mathsf G(2l)\ge\frac L2\mathsf e(l)$ over the cross-section $(2l,2l)$, as $n\le l$ and $\mathsf e(l)=2\mathsf G(l)-\mathsf G(2l)$, the two other thermal defects of $D_L$ being $\ge0$; this is $\mathsf e(l)\le\frac2LD_L(2l/L)$.

The cube criterion and $D_L(2l)\le\frac{20}7L^3x^2$ are Step~11; $\mathsf e\le8\delta$; and $0.13736\cdot0.8901^{24}\le0.0084022$. Finally $F\le0$ and, by Jensen's inequality for the Haar measure, $F\ge-3\beta|\underline n|$, so $\mathsf e(l)\le-F((2l)^3)\le24\beta l^3$. (2) is Proposition~\ref{C-361-8} with $x_0=8\delta(N^3)$, $\theta=103.2\,\delta(N^3)$ (if $\delta(N^3)=0$ the excess vanishes for all larger cubes), $\delta((2N)^3)\le\frac{20}7\mathsf e(N)^2\le\frac{20}7\cdot64\,\delta(N^3)^2$, and Proposition~\ref{C-361-8}(e) for $0.07489$. Two periodic Gibbs states exclude the first alternative by Proposition~\ref{C-361-10}(a); by (M3) each is an integral of ergodic ones, two different ergodic ones are mutually singular, and Proposition~\ref{C-361-9}(d) gives $2\le\liminf e^{\delta}$.

The last sentence is Proposition~\ref{C-361-9}(d), the sides being chosen with $\delta$ tending to its lower limit and the laws converging, by the compactness of the laws on a compact space. (3): if $\mathsf e(l)>0.069$ there is nothing to prove; otherwise let $\nu$ be the state of Proposition~\ref{C-361-10}, $f$ the bump of (M5)(iii) at radius $\mathsf d_\circ$ and height $2\mathsf d_\circ$, and $C(n)$ the covariance under $\nu$ of $\mathcal O_K(f\circ\Theta)$ and $\mathcal O_K(f(\cdot-na_Ke_3))$, $\Theta$ the reflection $x_3\mapsto-x_3$. For $0\le n\le N_4$ the pair is a test pair at $(\mathsf d_\circ;\vartheta_1)$ and two points of the supports are between $2\mathsf d_\circ+na_K$ and $6\mathsf d_\circ+na_K$ apart; so (M5)(iii) gives $C(n)>0$ and
\[
\frac{C(0)}{C(N_4)}\le\frac{4/3}{2/3}\Bigl(\frac{10\,\mathsf d_\circ}{2\,\mathsf d_\circ}\Bigr)^6=2\cdot5^6=31250 .
\]
By translation invariance $C(n)$ is the covariance of $\mathcal O_K(f)$ and its mirror image in the plane $x_3=-\frac n2a_K$, so Proposition~\ref{C-361-10}(f) with $A=B=\mathcal O_K(f)$, as in Step~5, gives $C(n)\le\theta^{n/l}C(0)$, $C\ge0$, and $C$ is log-convex, so all $C(n)>0$ and the quotients $C(n+1)/C(n)$ do not decrease: $C(n)\ge C(N_4)\,31250^{-(n-N_4)/N_4}$. Taking $n$-th roots, $\theta^{1/l}\ge31250^{-1/N_4}$, with $\theta=12.9\,\mathsf e(l)$; and $\log31250<10.35\le2.6\,N_4a_K/\mathsf d_\circ$ as $a_K/\mathsf d_\circ\le\frac1{52}$. With $\mathsf d_\circ=\mathsf d_0$, $l=\ell/2a_K$; (1) carries the floor to $D_2$, $8\delta$, $\frac8{L^3}D_L$; and $\frac1{12.9}31250^{-l/N_4}>0.069$ if and only if $l/N_4<\log(1/(12.9\cdot0.069))/\log31250$, which $la_K<0.0449\,\mathsf d_\circ$ gives, as $N_4>4\mathsf d_\circ/a_K-1$ and $\mathsf d_\circ/a_K>4L^3M$. (4): dilating by $L^z$, the label enters the period $(\gamma_W^{-2}+\bar B,\kappa_0]$, which $I(g_1)\cup I(g_2)$ covers by the form of the intervals $I(g)$ in (M5); there (3) at the side $L^z\ell_0$ gives $L^{-z}m_*\le2.6/\mathsf d_0^\star$, and $L^z\le\kappa_0/\mathsf b$, so $m_*\le2.6\kappa_0/(\mathsf d_0^\star\mathsf b)$, which is the bound since $\log(1/\theta)/1.3\ge0.0895$.

On the cube $(2l)^3$, Proposition~\ref{C-361-7}(d), (e) give $w_{00}=\frac14(1+2r_1+r_2)$, $w_{10}=w_{01}=\frac14(1-r_2)=u$, $w_{11}=\frac14(1-2r_1+r_2)=v$; with $x=1-r_1=2(u+v)$ the sum $\sum_fw_f^{1/2}=(1-u-\frac x2)^{1/2}+2u^{1/2}+(\frac x2-u)^{1/2}$ is concave in $u$ on the range $\frac x4\le u\le\frac x2$ allowed by $r_2\le r_1$ and $v\ge0$, so least at an end of it, which is $u=\frac x2$ for $x\le\frac12$; this gives $\mathsf e(l)\ge2\log(\sqrt{1-x}+\sqrt{2x})$, which exceeds $0.069$ for $x>0.00062714$ and is at least $\log2$ for $\frac19\le x\le\frac12$; for $x>\frac12$ both end values of the sum are at least $\sqrt2$, so $\mathsf e(l)\ge\log2$ there; and $\mathsf e(l)\ge\mathsf G_{2l}\ge\log\frac1{w_{00}}\ge\log\frac4{1+3r_1}$.

Then Proposition~\ref{C-361-7}(d) for $r_2,r_3$, Proposition~\ref{C-361-8}(b) for $\delta((2l)^3)\le\frac{20}7\cdot0.069^2\le0.0136029$, and Proposition~\ref{C-361-7}(b) for $\mathsf r_{l}(2l,2l)\le\varphi(e^{0.069}-1)\le0.0357638$, $\log(1/0.0357638)\approx3.33082$, $l_Ka_K=\ell_0/2$.

\emph{Step 15: Theorem~\ref{C-361-6}(5), (6).} (5)(a) is Proposition~\ref{C-361-15}(c) at $N=\ell/a_K$, $\beta=L^K\mathsf b$, where $N/\beta=\ell/\mathsf b$, with $|c'_N-c'|\le C_AN^{-2}(\log N)^{5/2}$; and the limit $K\to\infty$ by Theorem~\ref{C-361-5}(6). (b) Proposition~\ref{C-361-16}(a) and Proposition~\ref{C-361-7}(d), (e): as $r_2\le r_1$, $w_{00}$ is the largest weight and $\frac14\le w_{00}\le\frac14(1+3r_1)$, $|4w_f-1|\le3r_1$; $\delta_K(\ell)\ge\mathsf G_N\ge\log\frac1{w_{00}}\ge\log\frac4{1+3r_1}$, and $\mathsf e(\ell/2a_K)\ge\mathsf G(\ell/2a_K;N,N)\ge\mathsf G_N$; the last assertion follows from (4): $\eta_\flat(\ell_0/\mathsf b)<0.99937$ gives $1-r_1>0.000628$ on the cube $\ell_0/a_K$ at every cutoff. (c) Under \textbf{(Heaviest class)}, Proposition~\ref{C-361-16}(b) gives $\mathsf G_N=\log(Z/\lambda_0^N)\ge\log\frac1t-\mathsf q\log\log\frac1t-\log C$. (d) Choose $t_\circ>0$ with $16Lt_\circ\le t_2$ and $\eta_\flat(t)\le0.9961$ for $t\le16Lt_\circ$, and $\beta_\flat=1/8t_\circ$.

For $\beta\ge\beta_\flat$ there is $j\ge1$ with $2L^j/\beta\in(16t_\circ,16Lt_\circ]$; then $1-r_1(2L^j)\ge0.0039$ by (b), and Proposition~\ref{C-361-8}(d), by contraposition, excludes $\mathsf e(l)\le0.069$ for every $l\le2L^j/16$, which includes $3\le l\le t_\circ\beta$. For BI take $K$ with $L^K\mathsf b\ge\beta_\flat$: $l_K=\ell_0L^K/2>t_\circ L^K\mathsf b$. The exact statement on one lattice is the contrapositive of Proposition~\ref{C-361-8}(d). (e) is Proposition~\ref{C-361-15}(d). (6): with $N=L^{-j}\ell/a_K$ at the same $\beta$, (a) at the side $L^{-j}\ell$, where $t_1/L<L^{-j}t\le t_1$, bounds $\delta(N^3)$ by $\log(L/t_1)+c'+C$, and Theorem~\ref{C-361-6}(1) gives $\delta_K(\ell)\le3L^{2j}\delta(N^3)$ with $L^j\le Lt/t_1$; then $\mathsf e\le8\delta$ and $D_L\le L^3\delta$.

\begin{remark}[Scope]\label{C-361-17}
The theorems do not assert BI at any label; by Theorem~\ref{C-361-6} it fails at every cutoff on cubes small against the label, and it can hold only on cubes of side at least $0.0895\,(\mathsf d_0^\star/\kappa_0)\,\mathsf b$. Theorem~\ref{C-361-4}(2) carries BI from one label to every other label at other sides and from other thresholds, with numbers $\beta_0$, $C_+$ that are not explicit. Uniqueness in Theorem~\ref{C-361-3}(1)(b) is uniqueness among the periodic Gibbs states; of Gibbs states without periods it says only that $\nu_K$ is extreme and trivial at infinity. The area law bounds the loops in the state $\nu_K$ at each cutoff. Theorem~\ref{C-361-4}(5) concerns tuples carried by a cube that keeps away from the seam of the torus.
\end{remark}

\appendix

\section{The lattice two-loop coefficient}\label{IV-73-0}

This section concerns the truncated two-point function $\langle a_p;a_q\rangle_{n,\beta}$ of the actions of two plaquettes $p,q$ under the Wilson measure of $SU(2)$ lattice gauge theory in a box $\Lambda_n\subset\mathbb Z^3$ with Dirichlet boundary data, at large $\beta$. We expand it up to order $\beta^{-3}$, pass to the infinite lattice in the coefficients of the expansion, determine the behaviour of the coefficient $\mathsf B_1^{\mathrm{lat}}(p,q)$ of order $\beta^{-3}$ at large separation of $p$ and $q$, and compute the constant $\kappa^{\mathrm{lat}}$ that governs this behaviour. The argument has four stages.

First, in the box $\Lambda_n$ the change of variables of Faddeev and Popov and Laplace's method give a constant $C$, depending on $n$, $\tau_{\mathrm{gf}}$, $p$ and $q$, such that
\[
\bigl|\langle a_p;a_q\rangle_{n,\beta}-\beta^{-2}\,\mathsf B_0^{[n]}(p,q)-\beta^{-3}\,\mathsf B_1^{[n]}(p,q;\tau_{\mathrm{gf}})\bigr|\le C\beta^{-4}\qquad(\beta\ge1),
\]
and neither coefficient depends on the gauge parameter $\tau_{\mathrm{gf}}>0$ (Proposition~\ref{IV-73-27}).

Second, the coefficients $\mathsf B_0^{\mathrm{lat}}$ and $\mathsf B_1^{\mathrm{lat}}$ of the lattice $\mathbb Z^3$ are defined by the same cumulants of a Gaussian family as those of a box (Definition~\ref{IV-73-22}). Here $\mathsf B_0^{\mathrm{lat}}(p,q)=\frac{\mathsf n}2\,\Pi^{\mathrm{lat}}(p,q)^2$ by \eqref{IV-73-23}, while $\mathsf B_1^{\mathrm{lat}}(p,q)$ is an absolutely convergent sum of lattice graphs with two loops (Lemma~\ref{IV-73-55}). By the method of images, the coefficients of the plaquettes $p+o_n$, $q+o_n$, translated to the centre $o_n$ of the box, converge to $\mathsf B_j^{\mathrm{lat}}(p,q)$ as $n\to\infty$ (Proposition~\ref{IV-73-76}). Consequently
\[
\mathsf B_1^{\mathrm{lat}}(p,q)=\lim_{n\to\infty}\ \lim_{\beta\to\infty}\ \beta^3\bigl[\langle a_{p+o_n};a_{q+o_n}\rangle_{n,\beta}-\beta^{-2}\,\mathsf B_0^{[n]}(p+o_n,q+o_n)\bigr],
\]
the two limits being taken one after the other, first in $\beta$ and then in $n$ (Theorem~\ref{IV-73-100}(c)).

Third, a calculus of lattice kernels, combined with identities at zero momentum for the self-energy and for the vertex function, yields functions $\mathsf U_{ij}\in C^1(\mathbb R^3\setminus\{0\})$, $i,j\in\{1,2,3\}$, homogeneous of degree $-5$, and pure constants $C,r_0$ such that
\[
\bigl|\mathsf B_1^{\mathrm{lat}}(p,q)-\mathsf U_{\lambda(p)\lambda(q)}(u)\bigr|\le C\,r^{-6}\lg r\quad(r\ge r_0),\qquad\sum_{i,j}\mathsf U_{ij}(u)=\kappa^{\mathrm{lat}}|u|^{-5},
\]
where $u=x(q)-x(p)$ joins the barycentres of the plaquettes, $r=|u|$, $\lambda(p)$ and $\lambda(q)$ are their dual directions, and the sum is taken over the nine pairs of dual directions (Theorem~\ref{IV-73-71}). From this we obtain a rate for the convergence of the smeared coefficient $\mathsf P_1^{\mathrm{lat}}(f,h;a)$ to $\kappa^{\mathrm{lat}}\mathcal K_5(f,h)$ on test pairs $(f,h)$ (Corollary~\ref{IV-73-73}).

Fourth, we compute $\kappa^{\mathrm{lat}}$. Six two-loop integrands of a continuum model carry the terms of degree $-5$ of the lattice classes of graphs; the constant $\kappa^{\mathrm{lat}}$ is tied to their coefficients of $\mathsf p^2\log|\mathsf p|$, and these are evaluated by means of finite parts, an expansion by regions and a functional $\mathcal R$. The outcome is $\kappa^{\mathrm{lat}}=\frac3{2\pi^3}$ (Theorem~\ref{IV-73-98}), and hence $c_1:=\kappa^{\mathrm{lat}}/\mathsf c_2=\frac8{3\pi}$ (Theorem~\ref{IV-73-100}). The relation of this coefficient calculation to the interacting continuum two-point function is stated separately in Section~\ref{C-350-1}.

The proof below includes the gauge-fixed finite-box expansion, convergence
of the lattice graph sums, comparison with the six continuum integrands,
and the residue calculation giving the coefficient. All finite-box
remainder constants retain their stated volume dependence. The final
scope statement records the order of the limits.

\subsection{Notation and hypotheses}

We collect the conventions, the standing data and the facts on $SU(2)$ that are used throughout the section, and then summarise its notation.

\begin{definition}[Standing setting]\label{IV-73-1}
The following conventions hold throughout this section.
\begin{enumerate}
\item[(1)]

\emph{The group.} We take $G=SU(2)$ and $N=2$, and we put $\mathsf n:=N^2-1=3$. We write $\operatorname{tr}$ for the trace of $2\times2$ matrices, not divided by $2$, and for a complex $2\times2$ matrix $W$ we denote by $|W|$ its operator norm. The group $G$ is a compact subset of $\mathbb C^{2\times2}$ for the distance of this norm. Its normalised Haar measure $du$ is the Borel probability measure on $G$ that is invariant under left translations; it is unique, and it is also invariant under right translations and under $g\mapsto g^{-1}$. We denote by $\mathfrak g$ the real vector space of Hermitian traceless $2\times2$ matrices, with the inner product $\langle X,Y\rangle:=\frac12\operatorname{tr}XY$ and the norm $|X|:=\langle X,X\rangle^{1/2}$, which on $\mathfrak g$ coincides with the operator norm; $dY$ is the Lebesgue measure of the Euclidean space $(\mathfrak g,\langle\cdot,\cdot\rangle)$.

The Pauli matrices
\[
\mathsf e_1=\begin{pmatrix}0&1\\1&0\end{pmatrix},\qquad\mathsf e_2=\begin{pmatrix}0&-i\\i&0\end{pmatrix},\qquad\mathsf e_3=\begin{pmatrix}1&0\\0&-1\end{pmatrix}
\]
form an orthonormal basis of $\mathfrak g$, and $\mathsf e_a\mathsf e_b=\delta_{ab}\mathbf 1+i\sum_c\epsilon_{abc}\mathsf e_c$, where $\epsilon_{abc}$ is totally antisymmetric in $a,b,c$ with $\epsilon_{123}=1$: it is the sign of the permutation $(a,b,c)$ of $(1,2,3)$, and it vanishes when two of its indices agree. Writing $X^a$ for the coordinates of $X\in\mathfrak g$ in this basis, we define the cross product $X\times Y:=\sum_{a,b,c}\epsilon_{abc}X^aY^b\mathsf e_c$. For $X,Y\in\mathfrak g$ we then have
\[
\begin{multlined}
XY=\langle X,Y\rangle\mathbf 1+i\,X\times Y,\qquad[X,Y]=2i\,X\times Y,\qquad e^{iX}=\cos|X|\,\mathbf 1+i\,\frac{\sin|X|}{|X|}\,X, \\
\tfrac12\operatorname{tr}e^{iX}=\cos|X| ;
\end{multlined}
\]
the third identity follows from $X^2=|X|^2\mathbf 1$, its right side being $\mathbf 1$ at $X=0$. Colour indices $a,a',c,c',\dots$, as well as the index $b$ in $\epsilon_{abc}$, take their values in $\{1,2,3\}$.

\item[(2)]

\emph{Facts on $G$; smooth maps.} Put $\sigma_{\mathrm H}=\frac1{2\pi^2}$; the entire function $\sigma_1(w)=\bigl(\sum_{n\ge0}(-w)^n/(2n+1)!\bigr)^2$, which satisfies $\sigma_1(|Y|^2)=(\sin|Y|/|Y|)^2\in(0,1]$ for $|Y|<\pi$; and, for $Y\in\mathfrak g$, the linear map $\mathsf{dE}(Y)H=\int_0^1e^{isY}He^{-isY}\,ds$ of $\mathfrak g$. For $\varrho>0$ we write $\mathbb B_\varrho=\{Y\in\mathfrak g:|Y|<\varrho\}$. The following formulas will be used in the finite-box change of variables. Their proofs follow this definition.
\begin{itemize}
\item For $X_1,\dots,X_n\in i\mathfrak g$ one has $|e^{X_1}\cdots e^{X_n}-\mathbf 1|\le e^{|X_1|+\dots+|X_n|}-1$, all norms being operator norms as in (1). The map $Y\mapsto e^{iY}$ is a homeomorphism of $\mathbb B_\pi$ onto $G\setminus\{-\mathbf 1\}$, whose inverse is denoted by $\frac1i\log$.
\item For $Y,H\in\mathfrak g$, $\frac d{ds}e^{i(Y+sH)}\big|_{s=0}=i\bigl(\mathsf{dE}(Y)H\bigr)e^{iY}$; more generally, $\frac d{ds}e^{iZ(s)}\big|_{s=0}=i\bigl(\mathsf{dE}(Z(0))Z'(0)\bigr)e^{iZ(0)}$ for every curve $s\mapsto Z(s)$ in $\mathfrak g$ that is differentiable at $s=0$. Moreover $\det\mathsf{dE}(Y)=\sigma_1(|Y|^2)$.
\item For every bounded Borel function $F$ on $G$, $\int_GF\,du=\sigma_{\mathrm H}\int_{\mathbb B_\pi}F(e^{iY})\,\sigma_1(|Y|^2)\,dY$. Further, $\int_{\mathfrak g}e^{-\frac12|Y|^2}\,dY=(2\pi)^{3/2}$.
\end{itemize}
Smooth means $C^\infty$. A map $f$ between subsets of finite-dimensional real vector spaces is smooth if every point of its domain has an open neighbourhood $O$ in the ambient space and a $C^\infty$ map $\tilde f$ on $O$ such that $f=\tilde f$ on the intersection of $O$ with the domain of $f$; it is a diffeomorphism if it is bijective and both $f$ and $f^{-1}$ are smooth. Composites of smooth maps are smooth, and for maps between open sets these are the usual notions. The group $G$ and its powers $G^m$ lie in powers of $\mathbb C^{2\times2}$; the product of $G$ and the inversion $g\mapsto g^{-1}=g^*$ are smooth, since they are restrictions of real polynomial maps. Finally, $Y\mapsto e^{iY}$ is a diffeomorphism of $\mathbb B_\pi$ onto $G\setminus\{-\mathbf 1\}$; the proof is given after the present definition.

\item[(3)]

\emph{Standard facts.} The following standard facts are used without further reference.
\begin{itemize}
\item

\emph{Measure and integration:} the existence of Haar measure on $G$; the theorems of Fubini and Tonelli and the monotone and dominated convergence theorems, for integrals and for sums; the Borel measurability of partial integrals; the change of variables formula for a diffeomorphism between open subsets of $\mathbb R^m$, in particular for translations, reflections, rotations, dilations and polar coordinates in $\mathbb R^3$ and in $\mathbb R^6$; the Cauchy--Schwarz inequality; the continuity of translations on $L^2(\mathbb T)$; Gauss' divergence theorem on a ball and on a cube less a ball. An integer matrix of determinant $\pm1$ acting on $(\mathbb R^3)^{\mathsf m}$ induces a measure-preserving bijection of $\mathbb T^{\mathsf m}$.

For a function that is continuous on an interval $[a,b]$ or on a compact cube, the Riemann sums over equal subintervals or equal subcubes converge to the integral as the mesh tends to $0$. Integration by parts holds in one variable; on $\mathbb R^3$, when one of the two factors has a compact support or lies in the Schwartz space $\mathcal S(\mathbb R^3)$ of rapidly decreasing smooth functions; and over a period, for smooth periodic functions, where no boundary terms arise.
\item

\emph{Calculus and topology:} the Leibniz rule, the chain rule, the mean value theorem, the intermediate value theorem and the fundamental theorem of calculus; Taylor's formula with integral remainder, in one and in several variables; the inverse function theorem for $C^\infty$ maps; the Heine--Borel theorem; the density of $\mathbb Q^3$ in $\mathbb R^3$. A function with continuous partial derivatives is of class $C^1$. Differentiation under the integral sign with respect to a real or a complex parameter is allowed when the integrand and its derivative have integrable majorants that are locally uniform in the parameter, and a series of $C^j$ functions may be differentiated termwise up to the order $j$ when the suprema of the derivatives of orders $\le j$ are summable. A continuous function on a compact set is bounded and uniformly continuous and attains its minimum, and a finite product of compact metric spaces and a closed subset of a compact metric space are compact.
\item

\emph{Series:} the power series of $\exp$, $\cos$ and $\sin$, the series of $\log(1+w)$ and the binomial series of $(1+w)^{-t}$ for $|w|<1$, and the Cauchy product; the inequalities $\frac2\pi t\le\sin t\le t$ for $t\in[0,\frac\pi2]$. A function holomorphic on a disc about $0$ is the sum of its Taylor series on that disc, and such a series may be evaluated at a square matrix of norm smaller than its radius of convergence, sums and products of series then giving sums and products of the matrices. For $\|\mathsf A\|<1$ we have the Neumann series $(\mathbf 1+\mathsf A)^{-1}=\sum_{n\ge0}(-\mathsf A)^n$ and the expansion $\log\det(\mathbf 1+\mathsf A)=\sum_{n\ge1}\frac{(-1)^{n-1}}n\operatorname{tr}\mathsf A^n$.

Formal power series in finitely many variables form a ring, in which division by a series with constant term $1$ is possible, and $\exp$, applied to series without constant term, and $\log$, applied to series with constant term $1$, are mutually inverse. Real analytic functions on open subsets of $\mathbb R^m$, with real, complex or matrix values, are given locally by absolutely convergent power series, are smooth, are locally the sums of their homogeneous Taylor polynomials, and are stable under sums, products, composites and quotients with non-vanishing denominator; examples are the polynomials, the restrictions of holomorphic functions to real arguments, $\log$ on $(0,\infty)$, and a power series as above evaluated at a matrix with real analytic entries and norm smaller than the radius of convergence.
\item

\emph{Linear algebra:} trace, determinant, adjoint, orthonormal bases and the polarisation of a quadratic form. A positive definite operator is invertible; the determinant is an affine function of a rank-one summand; a polynomial vanishing on a non-empty open set is zero; and Lagrange's identity $\sum_{\mu<\nu}|v_\mu w_\nu-v_\nu w_\mu|^2=|v|^2|w|^2-|v^*w|^2$ holds for $v,w\in\mathbb C^3$. The Cayley map $\mathsf A\mapsto(\mathbf 1-\mathsf A)(\mathbf 1+\mathsf A)^{-1}$ is continuous from the real antisymmetric $3\times3$ matrices onto the rotations without eigenvalue $-1$, and these rotations are dense in $SO(3)$.
\item

\emph{Fourier analysis:} Plancherel's theorem on $\mathbb R^3$ and for the map $L^2(\mathbb T,d^\circ k)\to\ell^2(\mathbb Z^3)$; the uniqueness of the Fourier coefficients of an integrable function on a torus; the Poisson summation formula for functions in $C_c^\infty(\mathbb R^D)$, the smooth functions on $\mathbb R^D$ with a compact support; and the finite inversion formula
\[
f(z)=(2n)^{-3}\sum_{k\in\frac\pi n\mathbb Z^3\cap(-\pi,\pi]^3}e^{ik\cdot z}\sum_{z'}e^{-ik\cdot z'}f(z')
\]
for functions $f$ on $(\mathbb Z/2n)^3$, where $z'$ ranges over $(\mathbb Z/2n)^3$. The inverse Fourier transform $\mathcal F^{-1}$ of (5) maps $\mathcal S(\mathbb R^3)$ into itself, and $\mathcal F^{-1}(x_j\phi)=-i\partial_{p_j}\mathcal F^{-1}\phi$ when $\phi$ is written as a function of $x$ and $\mathcal F^{-1}\phi$ as a function of $p$; further, $\widehat{\partial_jw}=ip_j\hat w$ for every continuously differentiable function $w$ with a compact support, and $\int_{\mathbb R^3}e^{-tp^2+ip\cdot x}\,dp=(\pi/t)^{3/2}e^{-x^2/4t}$ for $t>0$.
\item

\emph{Probability:} for every positive semidefinite covariance $\mathsf C$ there is a centred Gaussian family $(\phi_i)$ with covariance $\mathsf C$; its law is determined by $\mathsf C$, and $\langle e^{\sum_is_i\phi_i}\rangle=e^{\frac12\sum_{i,j}s_is_j\mathsf C(i,j)}$ for real $s_i$, finitely many of them non-zero. If $H$ is a real symmetric positive definite $m\times m$ matrix, then under the probability measure on $\mathbb R^m$ with a density proportional to $e^{-\frac12(z,Hz)}$ the coordinates of $\mathbb R^m$ form a centred Gaussian family with the covariance $H^{-1}$.
\item

\emph{Complex analysis:} the identity theorem on connected open sets, and Euler's Gamma function, which we write $\Gamma_{\mathrm E}$: it is given by $\Gamma_{\mathrm E}(z)=\int_0^\infty t^{z-1}e^{-t}\,dt$ and is holomorphic for $\operatorname{Re}z>0$, it satisfies $\Gamma_{\mathrm E}(z+1)=z\,\Gamma_{\mathrm E}(z)$ and $\Gamma_{\mathrm E}(\frac12)=\sqrt\pi$, and $1/\Gamma_{\mathrm E}$ extends to an entire function.
\end{itemize}

\item[(4)]

\emph{Three functions.} For $x\in\mathbb R^D$ we put $\langle x\rangle:=\max\{1,|x|\}$. For $t>0$ we put $\lg t:=1+|\log t|$; then $\lg t\ge1$, $\lg(1/t)=\lg t$, and $\lg(st)\le\lg s\,\lg t$ for $s,t>0$. Finally $\log_+t:=\max\{0,\log t\}$, $\log_+\infty:=\infty$, and $t/0:=\infty$ for $t>0$.

\item[(5)]

\emph{Fourier conventions.} We put $\mathbb T:=\mathbb R^3/2\pi\mathbb Z^3$ and identify functions on $\mathbb T$ with the functions on $\mathbb R^3$ that are periodic with the periods $2\pi e_\mu$, $\mu=1,2,3$; integrals over $\mathbb T$ are integrals over $[-\pi,\pi]^3$. We write $|k|_{\mathbb T}:=\operatorname{dist}(k,2\pi\mathbb Z^3)$, which equals $|k|$ on $[-\pi,\pi]^3$, and $d^\circ k:=(2\pi)^{-3}dk$. We put $\mathsf k_\mu(k):=e^{ik_\mu}-1$ and $\mathsf k(k):=(\mathsf k_\mu(k))_\mu\in\mathbb C^3$. Then
\[
|\mathsf k(k)|^2=\sum_\mu(2-2\cos k_\mu),\qquad\frac4{\pi^2}|k|^2\le|\mathsf k(k)|^2\le|k|^2\quad\text{on }[-\pi,\pi]^3,
\]
the inequalities following from $2-2\cos k_\mu=4\sin^2\frac{k_\mu}2$ and the bounds on $\sin t$ in (3). The \emph{kernel with the symbol} $\sigma\in L^1(\mathbb T)$ is the function $K(x):=\int_{\mathbb T}e^{ik\cdot x}\sigma(k)\,d^\circ k$ of $x\in\mathbb Z^3$. A bound on a symbol or on an integrand that is said to hold ``on $[-\pi,\pi]^3$'' or ``for $|\kappa|\le c_0$'' holds almost everywhere if the function is an $L^1$-class, and at every point $\ne0$ if the function is given by a formula. For a function $\phi$ on $\mathbb R^D$ that lies in the Schwartz space $\mathcal S(\mathbb R^D)$, or is continuous with a compact support, we put $\hat\phi(\kappa):=\int e^{-i\kappa\cdot\zeta}\phi(\zeta)\,d\zeta$, and on $\mathbb R^3$ we put $(\mathcal F^{-1}\phi)(x):=\int e^{ip\cdot x}\phi(p)\,d^\circ p$. On $\mathbb Z^D$ we use the differences $(\nabla_jf)(\zeta):=f(\zeta+e_j)-f(\zeta)$ and, for a multi-index $\gamma$, $\nabla^\gamma:=\prod_j\nabla_j^{\gamma_j}$.

\item[(6)]

\emph{Data and constants.} The data of a statement are the real numbers and integers quantified in it, and each of them may take every value in its range. The ranges are $\tau_{\mathrm{gf}}>0$ for the gauge parameter (abbreviated $\tau$ where this is said), $\beta>0$, $a>0$, $\mathsf b>0$, $\mathsf d>0$, $\vartheta\in(0,1)$, $\Lambda>1$, $L>1$ for the real number $L$ of Lemma~\ref{IV-73-102}, and $n\ge2$ for the side of the box $\Lambda_n$. The \emph{tools} are the eight statements marked (K), (S), (P), (L), (D), (E), (V) and (T), namely Lemmas~\ref{IV-73-39}--\ref{IV-73-46}, Proposition~\ref{IV-73-47} and Lemma~\ref{IV-73-75}; the ranges of their data are those stated in their hypotheses.

A relation between data holds only where a statement imposes it; thus $a\le\mathsf d$ is assumed in Corollary~\ref{IV-73-73} and in Theorem~\ref{IV-73-100}(b). For each constant whose existence a statement asserts, such as $C$, $C_\gamma$, $c_0$, $r_0$, $\varrho_1$, $\varrho_2$, $\varrho_3$, the statement specifies the quantities on which the constant depends. A \emph{pure constant} is a constant that depends on nothing. In proofs, $C,C',c,\dots$ denote positive constants whose values may change from one occurrence to the next and which depend only on the quantities on which the statement in question allows its constants to depend.

\item[(7)]

\emph{Length scale.} We work on the lattice $\mathbb Z^3$ of spacing $1$. For a function $f$ on $\mathbb R^3$, $\operatorname{supp}f$ is the closure of $\{f\ne0\}$, $\|f\|_\infty=\sup|f|$, $\operatorname{Lip}f=\sup_{x\ne y}|f(x)-f(y)|/|x-y|$, $\int f$ is the Lebesgue integral, and $\mathbb B(x,\rho)=\{y:|y-x|\le\rho\}$.

\item[(8)]

\emph{Brackets, indices and notations with several meanings.} On complex forms, the bracket $(\cdot,\cdot)$ of Definition~\ref{IV-73-2}(2) is Hermitian and antilinear in its first argument: it is the sum of the sesquilinear extension of $\langle\cdot,\cdot\rangle$, evaluated at the values of the two forms. On $\mathfrak g\otimes\mathbb C$, on the other hand, $\langle X,Y\rangle$ denotes the complex bilinear extension. For two sites $x,x'$, $\langle x,x'\rangle$ is a bond. The brackets $\langle\cdot\rangle_{n,\beta}$, $\langle\cdot\rangle_{\mathcal W,\beta}$ and $\langle\cdot\rangle_0$ (Definition~\ref{IV-73-4}, Lemma~\ref{IV-73-26} and Definition~\ref{IV-73-19}) are means; the same brackets with entries separated by semicolons are cumulants; $\langle x\rangle$ is the function of (4).

When $a,a',c,c'$ occur as upper indices of forms or of legs, as indices of $\mathsf e$, $\delta$, $\epsilon$ or $\partial$, or as the second entry of a pair $(b,a)$, they are colour indices; $b,b'$ denote bonds. Forms and site functions are $\mathfrak g$-valued unless they are said to be $\mathbb R$-valued or complex. Finally, $D$ denotes a dimension, equal to $3$ or to $6$.

Several notations have more than one meaning in this section, and the following rules tell the meanings apart. Without index, $\Gamma$ ranges over $\{1,\dots,6\}$ and numbers the integrands $F_\Gamma$ of Definition~\ref{IV-73-64}; $\Gamma_q$ is the kernel of Definition~\ref{IV-73-29}(5), and $\Gamma_{\mathrm E}$ is Euler's function of (3). With the index $n$, $\Lambda_n$ is the box; in Lemma~\ref{IV-73-42} and wherever that lemma is applied, $\Lambda$ is the function $\Lambda(\kappa)$ defined there, also when its argument is not written; otherwise $\Lambda$ is a real number $>1$. The momenta are $\kappa,\kappa_1,\kappa_2,\dots$, the $\kappa_J$ are the cumulants of Definition~\ref{IV-73-19}(2), $\kappa^{\mathrm{lat}}$ and $\kappa_\Gamma^{\mathrm{lat}}$ are numbers, and in the statement of Lemma~\ref{IV-73-70} $\kappa$ alone is a constant.

Before Definition~\ref{IV-73-38}, $\sigma_1$ is the entire function of (2); from that Definition on, $\sigma,\sigma_1,\sigma_2,\dots$ denote symbols, and an $s$ or $s_i$ that accompanies a symbol or a line is its degree, whereas $(s_1,s_2,s_3,s_4)=(+,+,-,-)$ are the signs of the four bonds of a plaquette (Lemma~\ref{IV-73-9}). Further, $\mathsf n$ is the number $3$ of (1), except in Lemma~\ref{IV-73-53}(a), (b) and in its proof, where it is a matrix $(n_{vv'})$ of integers. As an index, $\beta$ is a direction or a multi-index; in Proposition~\ref{IV-73-87} and in the proof of Proposition~\ref{IV-73-94} it is a real coefficient; otherwise it is the coupling $\beta>0$.

As a datum, $\vartheta\in(0,1)$, whereas the argument $\vartheta$ of the function $j$ of Lemma~\ref{IV-73-31} is a real variable. As a datum, the integer $n\ge2$ is the side of the box $\Lambda_n$; other integers $n$, such as the degree $n=s+3$ of a line, the index of the class $\mathfrak K(n)$, the exponent $n\in\{4,5\}$ of the sums $\mathsf s_n$ and the index $n\in\mathbb Z$ of the spaces $\mathcal L_n$, are introduced where they occur. Besides the kernel $\mathsf q$ of Lemma~\ref{IV-73-46} and the kernels derived from it, $\mathsf q$ is a vector of $\mathbb R^3$ where it is the argument of $\rho$ or of a symbol, and $\mathsf q(t)$ is the polynomial of Lemma~\ref{IV-73-90}.

In Definition~\ref{IV-73-88}(2), Lemma~\ref{IV-73-90} and Lemma~\ref{IV-73-92}, $\mathsf a,\mathsf b,\mathsf c$ are integers, while the datum $\mathsf b>0$ of (6) is a real number; in $\mathsf d_{b,a}$ (Lemma~\ref{IV-73-57}(I)) the $\mathsf d$ belongs to the name of a polynomial; and $\tau$ is a variable of integration in the proofs of Lemmas~\ref{IV-73-92} and~\ref{IV-73-102}. With an argument, $\mathcal S(\mathbb R^D)$ is the Schwartz space; with an upper index, $\mathcal S^s$ is a class of symbols (Definition~\ref{IV-73-38}(2)).

\item[(9)]

\emph{Further definitions.} The lattice and its forms are defined in Definition~\ref{IV-73-2}, the finite-box measure in Definition~\ref{IV-73-4}, and Gaussian cumulants in Definition~\ref{IV-73-19}. The test pairs and kernels are defined in Section~\ref{C-350-1}; their bounds are Lemma~\ref{lem:test-pair-kernels}.
\end{enumerate}
\end{definition}

\begin{proof}[Proof of the group formulas in \textup{(2)}]
The product bound follows by writing $e^{X_j}=1+Z_j$,
using $|Z_j|\le e^{|X_j|}-1$, and expanding the product:
the sum of the norms of its nonconstant terms is at most
$\prod_j(1+|Z_j|)-1\le e^{\sum_j|X_j|}-1$.
Every element of $SU(2)$ has the unique form
$a_0\mathbf1+i\sum_aa_a\mathsf e_a$ with real coefficients and
$a_0^2+\sum_aa_a^2=1$. The exponential formula in (1) therefore
gives a bijection from $\mathbb B_\pi$ to $G\setminus\{-\mathbf1\}$.
Its smooth inverse is verified below.

For matrices $A,B$, differentiate
$s\mapsto e^{s(A+rB)}e^{(1-s)A}$ and integrate to obtain
\[
 e^{A+rB}-e^A
 =r\int_0^1e^{s(A+rB)}Be^{(1-s)A}\,ds.
\]
Dividing by $r$ and passing to zero proves the stated
differential formula. If $\vartheta=|Y|>0$, conjugation by
$e^{isY}$ fixes the line spanned by $Y$ and acts on its
orthogonal plane as a rotation through $2s\vartheta$.
Indeed, for $H\perp Y$ the Pauli identities give
$HY=-YH$ and
$e^{isY}He^{-isY}
=\cos(2s\vartheta)H+\sin(2s\vartheta)iYH/\vartheta$.
Consequently the determinant of its average is
\[
 \det\mathsf{dE}(Y)
 =\left(\frac{\sin(2\vartheta)}{2\vartheta}\right)^2
  +\left(\frac{1-\cos(2\vartheta)}{2\vartheta}\right)^2
 =\left(\frac{\sin\vartheta}{\vartheta}\right)^2.
\]
The value at zero follows by continuity.

For completeness, this determinant also identifies the Haar
density. Let $\mu$ be the pushforward of
$\sigma_1(|Y|^2)dY$ on $\mathbb B_\pi$ under the exponential.
Right translation by $u\in G$ induces
$\tau_u(Y)=i^{-1}\log(e^{iY}u)$ wherever its argument is not
$-\mathbf1$. It is a diffeomorphism after removing at most one
point from each chart, with inverse $\tau_{u^{-1}}$.
Differentiating $e^{i\tau_u(Y)}=e^{iY}u$ gives
\[
 \mathsf{dE}(\tau_u(Y))D\tau_u(Y)=\mathsf{dE}(Y),\qquad
 \sigma_1(|\tau_u(Y)|^2)|\det D\tau_u(Y)|
 =\sigma_1(|Y|^2).
\]
The change of variables formula, with those null points omitted,
shows that $\mu$ is right invariant. Its mass is
$4\pi\int_0^\pi\sin^2\vartheta\,d\vartheta=2\pi^2$.
Uniqueness of normalized Haar measure proves the integration
formula. The Gaussian integral is the product of the three
one-dimensional Gaussian integrals in the orthonormal Pauli basis.
\end{proof}

\begin{proof}[Smoothness of the inverse exponential in \textup{(2)}]
Let $\mathrm{cs}(w):=\sum_{n\ge0}(-w)^n/(2n)!$ and $\mathrm{sn}(w):=\sum_{n\ge0}(-w)^n/(2n+1)!$. These series define entire functions which are real on $\mathbb R$, and comparing them term by term with the series of $\cos$ and $\sin$, and differentiating term by term, we obtain
\[
\mathrm{cs}(\varphi^2)=\cos\varphi,\qquad\mathrm{sn}(\varphi^2)=\frac{\sin\varphi}\varphi,\qquad\sigma_1=\mathrm{sn}^2,\qquad\mathrm{cs}'=-\tfrac12\,\mathrm{sn},
\]
for real $\varphi$, with $\varphi\ne0$ in the second identity; moreover $\mathrm{sn}(0)=1$. By (1), $e^{iY}=\mathrm{cs}(|Y|^2)\mathbf 1+i\,\mathrm{sn}(|Y|^2)\,Y$ for $Y\in\mathfrak g$, so the map $Y\mapsto e^{iY}$ is smooth on $\mathfrak g$. On the interval $(-\infty,\pi^2)$ the function $\mathrm{sn}$ is positive: for $w<0$ all terms of its series are positive, and for $0\le w<\pi^2$ it equals $\sin\sqrt w/\sqrt w>0$, with the value $1$ at $w=0$. Hence $\mathrm{cs}$ is strictly decreasing on $(-\infty,\pi^2)$ and its derivative does not vanish there. Moreover $\mathrm{cs}(w)\to+\infty$ as $w\to-\infty$, and $\mathrm{cs}(\pi^2)=\cos\pi=-1$. By the intermediate value theorem $\mathrm{cs}$ maps $(-\infty,\pi^2)$ onto $(-1,\infty)$, and by the inverse function theorem the inverse map $\mathrm{cs}^{-1}:(-1,\infty)\to(-\infty,\pi^2)$ is smooth.

Now let $W=e^{iY}$ with $Y\in\mathbb B_\pi$. Since $\operatorname{tr}Y=0$, $\operatorname{tr}\mathbf 1=2$ and $Y^*=Y$, the formula above gives $\frac12\operatorname{tr}W=\mathrm{cs}(|Y|^2)$ and $\frac1{2i}(W-W^*)=\mathrm{sn}(|Y|^2)\,Y$. As $|Y|^2\in[0,\pi^2)$, the first identity yields $|Y|^2=\mathrm{cs}^{-1}(\frac12\operatorname{Re}\operatorname{tr}W)$, and the second then yields
\[
\tfrac1i\log W=\Bigl[\mathrm{sn}\bigl(\mathrm{cs}^{-1}(\tfrac12\operatorname{Re}\operatorname{tr}W)\bigr)\Bigr]^{-1}\,\tfrac1{2i}(W-W^*)\qquad(W\in G\setminus\{-\mathbf 1\}).
\]
The right side is defined and smooth on the open set $\{W\in\mathbb C^{2\times2}:\frac12\operatorname{Re}\operatorname{tr}W>-1\}$: there $\mathrm{cs}^{-1}$ is smooth and takes its values in $(-\infty,\pi^2)$, where $\mathrm{sn}$ is positive. This open set contains $G\setminus\{-\mathbf 1\}$, because $\frac12\operatorname{tr}e^{iY}=\cos|Y|>-1$ for $Y\in\mathbb B_\pi$. Hence the inverse map $\frac1i\log$ is smooth, and the bijection $Y\mapsto e^{iY}$ of $\mathbb B_\pi$ onto $G\setminus\{-\mathbf 1\}$ given by (2) is a diffeomorphism.
\end{proof}

\medskip
For reference we summarise the notation introduced in the rest of the section, together with the places where it is defined.

\emph{Lattice and boxes} (Definition~\ref{IV-73-2}). We write $(x,\mu):=\langle x,x+e_\mu\rangle$ for the positively oriented bonds and, for a $1$-form $A$, $A_\mu(x):=A((x,\mu))$, with components $A^a_\mu(x)$. Bonds $b$ and plaquettes $p$ have the base points $b_-$ and $x_p$; the plaquette $p=p_{\mu\nu}(x)$ has the barycentre $x(p)=x+\frac12(e_\mu+e_\nu)$ and the dual direction $\lambda(p)\in\{1,2,3\}\setminus\{\mu,\nu\}$; $p_{[i]}(x)$ is the plaquette with base point $x$ and dual direction $i$; and $b+y$, $p+y$ are translates. The differences are
\[
\begin{multlined}
(\mathrm d\lambda)(b):=\lambda(b_+)-\lambda(b_-),\qquad(\mathrm dA)(p):=A(b_1)+A(b_2)-A(b_3)-A(b_4),\qquad(\mathrm d^*A)(x): \\
=\sum_\mu\bigl[A_\mu(x-e_\mu)-A_\mu(x)\bigr],
\end{multlined}
\]
and $(\cdot,\cdot)$, $\|\cdot\|$ are the sums of $\langle\cdot,\cdot\rangle$ over sites, over positively oriented bonds or over plaquettes; further $\bar\lambda(b):=\frac12(\lambda(b_-)+\lambda(b_+))$ and $(A\times\bar\lambda)(b):=A(b)\times\bar\lambda(b)$. The box is $\Lambda_n:=\{0,\dots,n\}^3$, with $\Lambda_n^\circ:=\{1,\dots,n-1\}^3$ and $\partial\Lambda_n:=\Lambda_n\setminus\Lambda_n^\circ$. Attached to it are the set $\mathsf{Bd}_n$ of \emph{free bonds}, which are the positively oriented bonds with both end points in $\Lambda_n$ and at least one of them in the interior $\Lambda_n^\circ$, the plaquettes $\mathsf{Pl}_n$, the configurations $\mathcal U_n\cong G^{\mathsf{Bd}_n}$, the forms $\mathsf{Form}_n$, the site functions $\mathsf X_n$, the group $\mathcal G_n:=\{u:\Lambda_n\to G,\ u=\mathbf 1\text{ on }\partial\Lambda_n\}$, and the $\mathbb R$-valued analogues $\mathsf{Form}_n^{\mathbb R}$, $\mathsf X_n^{\mathbb R}$ (Definition~\ref{IV-73-2}(3)).

Finally $|A|_\infty:=\max_b|A(b)|$, $|\lambda|_\infty:=\max_x|\lambda(x)|$, and the centre of the box is $o_n:=(\lfloor n/2\rfloor,\lfloor n/2\rfloor,\lfloor n/2\rfloor)$.

\emph{Measure} (Definition~\ref{IV-73-4}). The Wilson measure of the box is
\[
d\mu_{n,\beta}(U):=Z_{n,\beta}^{-1}\exp\Bigl[-\beta\sum_{p\in\mathsf{Pl}_n}a_p(U)\Bigr]\prod_{b\in\mathsf{Bd}_n}dU(b),
\]
with the mean $\langle\cdot\rangle_{n,\beta}$ and the truncated two-point function $\langle X;Y\rangle_{n,\beta}=\langle XY\rangle_{n,\beta}-\langle X\rangle_{n,\beta}\langle Y\rangle_{n,\beta}$.

\emph{Chart} (Definition~\ref{IV-73-6}, Lemma~\ref{IV-73-9}). The coefficients $\mathsf a_p^{(m)}(A)$ are given by $a_p(e^{itA})=\sum_{m\ge0}t^m\,\mathsf a_p^{(m)}(A)$. We put $A^\lambda(b):=\frac1i\log\bigl(e^{i\lambda(b_-)}e^{iA(b)}e^{-i\lambda(b_+)}\bigr)$; the gauge function is $A\mapsto\mathrm d^*A$, and the Faddeev--Popov operator is $\mathsf M(A)\lambda:=\frac d{ds}\mathrm d^*(A^{s\lambda})\big|_{s=0}$, with $\mathsf M(tA)=\sum_{j\ge0}t^j\mathsf M^{(j)}(A)$. The signs of the four bonds of a plaquette are $(s_1,s_2,s_3,s_4):=(+,+,-,-)$. With $Z:=\frac1i\log\partial(e^{itA})(p)$ we put $\mathfrak z_j:=\frac1{(j+1)!}\partial_t^{j+1}Z\big|_{t=0}$, and $\mathfrak z_j(p)$ denotes its value at the values of $A$ on the bonds of $p$. Finally $\mathsf a_Y:=-2\,Y\times\cdot$, and $\mathsf r(Y)$ is the remainder of fourth order in \eqref{IV-73-11}.

\emph{Slice} (Lemma~\ref{IV-73-12}). We write $S(U):=\sum_{p\in\mathsf{Pl}_n}a_p(U)$, $\mathsf T_n:=\{\langle x-e_1,x\rangle:x\in\Lambda_n^\circ\}$, $\mathcal U_n^{\mathsf T}:=\{V\in\mathcal U_n:V=\mathbf 1\text{ on }\mathsf T_n\}$ and $\mathsf{Form}_n^{\mathsf T}:=\{A\in\mathsf{Form}_n:A=0\text{ on }\mathsf T_n\}$, and $\mathsf j$ for the map $(u,V)\mapsto V^u$. The lemma provides the radii $\varrho_1,\varrho_2,\varrho_3$, the set $\mathfrak D$, the map $\Xi:(\lambda,A)\mapsto A^\lambda$ and the set $\mathcal W:=\Xi(\mathfrak D)$.

\emph{Action, propagators and vertices} (Definition~\ref{IV-73-15}). The upper index $\mathrm{lat}$ refers to $\mathbb Z^3$ and the upper index $[n]$ to $\Lambda_n$. We put $\mathsf S_2(A):=\frac12\|\mathrm dA\|^2+\frac1{2\tau_{\mathrm{gf}}}\|\mathrm d^*A\|^2$, $\mathsf C_{\tau_{\mathrm{gf}}}=(\mathrm d^*\mathrm d+\tau_{\mathrm{gf}}^{-1}\mathrm d\mathrm d^*)^{-1}$ and $\mathsf G=(\mathrm d^*\mathrm d)^{-1}$, and
\[
\begin{multlined}
\mathsf E_p(b):=\sum_is_i\,\mathsf C_{\tau_{\mathrm{gf}}}(b_i(p),b),\qquad\Pi(p,q):=\sum_js_j\,\mathsf E_p(b_j(q)),\qquad\mathsf g(b):=\mathsf G(b_+,b_+)+\mathsf G(b_-, \\
b_-)-2\mathsf G(b_+,b_-),
\end{multlined}
\]
\[
\mathsf D(b;b'):=\tfrac12\sum_{x\in\{b'_-,b'_+\}}\bigl[\mathsf G(x,b_+)-\mathsf G(x,b_-)\bigr].
\]
The vertices are $\mathsf S_3:=\sum_p\mathsf a_p^{(3)}$, the quartic Wilson vertex $\mathsf S_4:=\sum_p\mathsf a_p^{(4)}$, and
\[
\mathsf V_2:=\mathsf S_4+\frac13\sum_b|A(b)|^2+\frac23\sum_b\mathsf g(b)|A(b)|^2-4\sum_{b,b'}\mathsf D(b;b')\mathsf D(b';b)\langle A(b),A(b')\rangle,
\]
whose last three terms are the \emph{Haar term}, which is the quadratic term contributed by the Haar density in the exponential chart, the ghost tadpole and the ghost loop.

\emph{Gaussian calculus} (Definition~\ref{IV-73-19}, Lemma~\ref{IV-73-53}). The mean $\langle\cdot\rangle_0$ belongs to the centred Gaussian family with $\langle A^a(b)A^{a'}(b')\rangle_0=\delta_{aa'}\mathsf C^\Omega_{\tau_{\mathrm{gf}}}(b,b')$, and $\langle X_1;\dots;X_r\rangle_0$ is the joint cumulant of polynomials $X_1,\dots,X_r$ in this family. A \emph{leg} is a linear form that occurs as a factor of a monomial of a vertex: a \emph{plain leg} is a variable $A^a(b)$, and a \emph{curl leg} is $(\mathrm dA)^a(p)=\sum_is_iA^a(b_i(p))$. The \emph{standard writing} of $\mathsf a_p^{(2)}$, $\mathsf a_p^{(3)}$ and $\mathsf a_p^{(4)}$ is a fixed way of writing them as sums of monomials in plain and curl legs, chosen so that the counting of powers takes the curl legs into account.

On $\mathbb Z^3$, $\mathsf S_3$ and $\mathsf V_2$ are \emph{formal sums of local terms}: the local terms of $\mathsf S_3$ are the $\mathsf a_p^{(3)}$, and those of $\mathsf V_2$ are the $\mathsf a_p^{(4)}$, the terms of the sums over bonds $b$ and the terms of the sum over pairs $(b,b')$; the cumulants of formal sums are defined term by term under a condition of absolute convergence (Definition~\ref{IV-73-19}(3)). Further $\mathsf F[\![X_1,\dots,X_m]\!]:=\partial_{s_1}\cdots\partial_{s_m}\mathsf F(\sum_is_iX_i)$; for $\mathsf i=(b,a)$, $\delta_{\mathsf i}$ is the form equal to $\mathsf e_a$ at $b$ and to $0$ at the other positively oriented bonds; $\mathsf C(\mathsf i,\mathsf i'):=\delta_{aa'}\mathsf C(b,b')$ and $\mathsf E_p^a:=\mathsf E_p(\cdot)\mathsf e_a$; the values $\mathsf{val}(\mathsf n)$ and the \emph{classes} $\mathsf T_{\mathrm I},\mathsf T_1,\mathsf T_4,\mathsf T_{4'},\mathsf T_5,\mathsf T_6$, which are the topologies of the connected graphs of $\mathsf B_1$, are those of Lemma~\ref{IV-73-53}; and $\mathsf t_P(X,Y,Z):=-\langle\mathrm dX(P),\sum_{i<j}s_is_j\,Y(b_i)\times Z(b_j)\rangle$.

\emph{Coefficients} (Definition~\ref{IV-73-22}). We put $\mathsf B_0^\Omega(p,q):=\langle\mathsf a_p^{(2)};\mathsf a_q^{(2)}\rangle_0$; $\mathsf B_1^\Omega(p,q;\tau_{\mathrm{gf}})$ is the combination of seven cumulants in \eqref{IV-73-24}; $\mathsf B_j^{[n]}$ and $\mathsf B_j^{\mathrm{lat}}$ are these coefficients for $\Omega=\Lambda_n$ and for $\Omega=\mathbb Z^3$; and
\[
\mathsf B_1^{\mathrm{lat}}(p,q):=\mathsf B_1^{\mathrm{lat}}(p,q;1),\qquad\mathsf P_1^{\mathrm{lat}}(f,h;a):=a\sum_{p,q}f(a\,x(p))\,h(a\,x(q))\,\mathsf B_1^{\mathrm{lat}}(p,q).
\]

\emph{Laplace's method} (Lemma~\ref{IV-73-26}). Here $\langle\cdot\rangle_{\mathcal W,\beta}$ is the mean for the density proportional to $e^{-\beta\Phi-\Theta}$ on $\mathcal W$; we put $W_j:=\Phi_{j+2}+\Theta_j$, $X_j:=X'_{j+2}$ and $Y_j:=Y'_{j+2}$, and $\gamma_0,\gamma_2$ are the two coefficients defined in the lemma.

\emph{Zero momentum} (Definition~\ref{IV-73-29}, Lemmas~\ref{IV-73-31} and~\ref{IV-73-37}). We write $\partial_{b,a}$ for the derivative with respect to $A^a(b)$. The matrices $\mathsf N(B)$ and $\mathsf h_q(B)$ are the Hessians at $A=B$ of $\sum_pa_p(e^{iA})+\frac1{2\tau_{\mathrm{gf}}}\sum_x|\mathrm d^*A(x)|^2$ and of $a_q(e^{iA})$, respectively. The constant abelian field is $Y_\gamma((x,\mu)):=\gamma_\mu\mathsf e_3$. We put $\mathsf e_\pm:=2^{-1/2}(\mathsf e_1\pm i\mathsf e_2)$ and $\mathsf e_0:=\mathsf e_3$, call the index $e\in\{+1,-1,0\}$ of $\mathsf e_e$ its \emph{charge}, which is the eigenvalue of $\frac12\operatorname{ad}\mathsf e_3$ on $\mathsf e_e$, and set $\upsilon:=2e\gamma$.

The symbols at $Y_\gamma$ are $\hat{\mathsf N}_e$, $\hat{\mathsf h}_{q,e}$ and $\hat{\mathsf M}_e$; an additional index $j$, as in $\hat{\mathsf N}_{e,j}$, $\hat{\mathsf h}_{q,e,j}$ and $\hat{\mathsf M}_{e,j}$, marks the part that is homogeneous of degree $j$ in $\gamma$; and $\hat{\mathsf C}:=\hat{\mathsf N}_{e,0}^{-1}$. The integrands $\mathfrak s^{(i)}$ and the numbers $c^{(i)}_{\mu\nu}(\tau_{\mathrm{gf}})$, $i\le5$, describe the self-energy, and $Z_0:=\int_{\mathbb T}d^\circ k/|\mathsf k(k)|^2$; the integrands $\mathfrak v_q^{(j)}$ and the numbers $\mathsf g^{(j)}_{q;\mu\nu}(\tau_{\mathrm{gf}})$, $j\le4$, describe the vertex function; and $\Sigma_{\mathrm{tot}}$, $\Gamma_q$ are the kernels of Definition~\ref{IV-73-29}(5).

Moreover $j(\vartheta):=(e^{i\vartheta}-1)/i\vartheta$ for real $\vartheta\ne0$ and $j(0):=1$, $J:=\operatorname{diag}(j(\upsilon_\mu))$, $P(k):=|\mathsf k(k)|^2\mathbf 1-\mathsf k(k)\mathsf k(k)^*$, and, for $q=p_{\mu\nu}(x)$, $D_q(k)\alpha=\mathsf k_\mu(k)\alpha_\nu-\mathsf k_\nu(k)\alpha_\mu$ and $\varpi_q(k):=(|\mathsf k_\mu(k)|^2+|\mathsf k_\nu(k)|^2)/|\mathsf k(k)|^2$. Finally $R_c:x\mapsto2c-x$ is the point reflection at $c$, and $(x,\mu)^R=(2c-x-e_\mu,\mu)$.

\emph{Kernel calculus} (Definition~\ref{IV-73-38} and the tools). The cut-offs are $\theta_0(t):=e^{-1/t}$ for $t>0$, $\theta_0(t):=0$ for $t\le0$,
\[
\theta(s):=\frac{\theta_0(1-s)}{\theta_0(1-s)+\theta_0(s-\frac12)},\qquad\chi(\kappa):=\theta(|\kappa|/2).
\]
The symbol classes are the $\mathcal S^s$. A \emph{line symbol} of degree $s$, an integer $\ge-2$, is a symbol in $\mathcal S^s$ with a smooth leading part $\sigma^h$ that is homogeneous of degree $s$ and of parity $(-1)^s$, the remainder $\sigma-\chi\sigma^h$ lying in $\mathcal S^{s+1}$; its kernel is an \emph{elementary line} of degree $n=s+3$, which decays like $|x|^{-n}$. The \emph{line class} $\mathfrak K(n)$ is defined by (K1)--(K3): its elements are kernels that have a homogeneous leading symbol and a homogeneous leading kernel, up to one order and a logarithm, and an element $K$ has its \emph{leading data} $(\sigma^h,K^h)$ and its \emph{constants}.

A \emph{slowly varying function} is a lattice function whose differences of orders up to ten each gain a factor $R^{-1}$, as measured by $\|g\|_{(R)}:=\max_{|\gamma|\le10}R^{|\gamma|}\sup_\zeta|\nabla^\gamma g(\zeta)|$; with it we use $\check g(\kappa):=\sum_\zeta g(\zeta)e^{-i\kappa\cdot\zeta}$ and $\|G\|_{C^{10}}$. The \emph{majorants} of degree $s>-D$ are the Borel functions on $\mathbb R^D$ with values in $[0,\infty]$ that are homogeneous of degree $s$ and integrable over the ball $\{|\kappa|\le1\}$; they form the class $\mathcal H_s(D)$. The kernel $\mathcal F^{-1}_{\ne0}\sigma^h$ of a homogeneous symbol $\sigma^h$ is defined in Lemma~\ref{IV-73-39}(b).

For the one-loop functions of Lemma~\ref{IV-73-42}, which are built from $\mathsf m\in\{2,3\}$ line symbols $\sigma_1,\dots,\sigma_{\mathsf m}$ of the degrees $s_1,\dots,s_{\mathsf m}$, from signs $\epsilon_i\in\{\pm1\}$ and from linear maps $\lambda_i:\mathbb R^D\to\mathbb R^3$ with $D=3(\mathsf m-1)$, we use
\[
\begin{multlined}
\mathsf I(\kappa):=\int_{\mathbb T}\prod_i\sigma_i\bigl(\epsilon_il+\lambda_i(\kappa)\bigr)\,d^\circ l,\qquad l_i(\kappa):=-\epsilon_i\lambda_i(\kappa),\qquad\omega:=3+\sum_is_i,\qquad\Lambda(\kappa):=1 \\
+\sum_{i<j}\log_+\frac{|\kappa|}{|l_i(\kappa)-l_j(\kappa)|},
\end{multlined}
\]
together with \emph{the polynomial} $\mathsf Q$ and \emph{the homogeneous part} $\mathsf I^h$ of $\mathsf I$, which are its Taylor part and its non-analytic homogeneous part at small $\kappa$. In Lemma~\ref{IV-73-46}, a \emph{chain} is a convolution of lines, $(X*Y)(u):=\sum_zX(z)Y(u-z)$, and a \emph{bubble} is a product $K_{\mathrm a}K_{\mathrm b}$ of two lines, that is, a one-loop two-point graph; with the polynomial $\mathsf Q$ of the symbol of the bubble we put $\tilde{\mathsf Q}(k):=\mathsf Q(0)+\sum_j\partial_j\mathsf Q(0)\sin k_j$, and $\mathsf q$ is the kernel with the symbol $\tilde{\mathsf Q}$. The \emph{kite} of Proposition~\ref{IV-73-47} is the two-loop two-point graph with five elementary lines $K_1,\dots,K_5$ of the degrees $n_1,\dots,n_5$,
\[
\mathsf T(u):=\sum_{z,w}K_1(z)K_2(w)K_3(w-z)K_4(u-z)K_5(u-w),\qquad\varDelta:=\sum_in_i-6,
\]
and with it we use the functions $\hat{\mathsf T}_{\mathsf p},\hat{\mathsf T}_{\mathsf q}$, the numbers $\omega_{\mathsf p},\omega_{\mathsf q}$, the polynomials $\mathsf Q_\bullet$ with the corresponding $\tilde{\mathsf Q}_\bullet$ and kernels $\mathsf q_\bullet$, the function $\mathsf U$, and the \emph{short} and the \emph{long lines}, which are the lines of the kite cut off to $|y|\lesssim\varepsilon r$ and to $|y|\gtrsim\varepsilon r$, respectively. On discrete tori (Lemma~\ref{IV-73-75}) we use $\mathbb T_n:=\frac\pi n\mathbb Z^3\cap(-\pi,\pi]^3$, $|z|_n$, $\langle z\rangle_n$ and $K^{(n)}(z):=(2n)^{-3}\sum_{k\in\mathbb T_n\setminus\{0\}}e^{ik\cdot z}\sigma(k)$.

\emph{Graphs on $\mathbb Z^3$.} For $p_{[\iota]}=p_{\alpha\beta}$ we put $\mathsf c_{\iota,\nu}(k):=\mathsf k_\alpha(k)\delta_{\nu\beta}-\mathsf k_\beta(k)\delta_{\nu\alpha}$, and $\hat{\mathsf C}_{\mu\nu}(k)$ denotes the bracket in the Fourier representation of $\mathsf C^{\mathrm{lat}}_{\tau_{\mathrm{gf}}}$ in Definition~\ref{IV-73-15}(1) (Lemma~\ref{IV-73-51}). For $p=p_{\mu\nu}(x)$, $\epsilon_p:=-1$ if $(\mu,\nu)=(1,3)$ and $\epsilon_p:=+1$ otherwise, and $p^\flat:=(x,\lambda(p))$ (Lemma~\ref{IV-73-52}). Lemma~\ref{IV-73-57} uses $u':=x_q-x_p$, the index sets $\mathfrak G_\bullet(\iota,\iota')$, the coefficients $c_G$ of the \emph{elementary terms} $G$, each of them the contribution of one monomial and one pairing to one of the classes $\mathsf T_\bullet$, the lines $K^\to_{\iota\mu}(y):=\mathsf E_{p_{[\iota]}(0)}((y,\mu))$ and $K^\leftarrow_{\nu\iota'}(y):=\mathsf E_{p_{[\iota']}(y)}((0,\nu))$, the kernel $\Sigma_{\mu\nu}(y):=\Sigma_{\mathrm{tot}}((x,\mu),(x+y,\nu))$, the signs $\mathsf s_P(b)$, the decomposition $\partial_{b,a}\mathsf S_3=\mathsf d_{b,a}+\mathsf u_{b,a}$, the kernel $\mathsf q^0_{\mu\nu}$ of finite support, and the difference operators $\mathsf f_G(\nabla)$ of degree $d_G$.

Lemmas~\ref{IV-73-60} and~\ref{IV-73-69} use $\mathsf q_X,\mathsf q_Y$, $X':=X-\mathsf q_X$, $Y':=Y-\mathsf q_Y$ and the functions $\mathsf U_G$. Lemma~\ref{IV-73-62} uses the kernel $\mathsf q_G$ of the bubble $K_{\mathrm a}K_{\mathrm b}$ of each term $G$ of $\Sigma_{\mu\nu}$, and it splits $\Sigma_{\mu\nu}=\mathsf q^\Sigma_{\mu\nu}+\mathsf W_{\mu\nu}$ into a kernel $\mathsf q^\Sigma_{\mu\nu}$ of finite support and a remainder $\mathsf W_{\mu\nu}\in\mathfrak K(4)$,
\[
\mathsf q^\Sigma_{\mu\nu}:=\mathsf q^0_{\mu\nu}+\sum_Gc_G\,\mathsf f_G(\nabla)\mathsf q_G,\qquad\mathsf W_{\mu\nu}:=\sum_Gc_G\,\mathsf f_G(\nabla)\bigl[K_{\mathrm a}K_{\mathrm b}-\mathsf q_G\bigr];
\]
it also introduces a symmetric kernel $\gamma_q$ of finite range and the kernels $\gamma_{q;\mu\nu}$ derived from it. The \emph{local parts} are kernels of finite support that carry the Taylor polynomials, of degree at most $1$, of one-loop subgraphs, such as $\mathsf q$, $\mathsf q_\bullet$, $\mathsf q^\Sigma_{\mu\nu}$ and $\gamma_q$; by Lemma~\ref{IV-73-62}, the symbols of $\mathsf q^\Sigma_{\mu\nu}$ and of $\gamma_{q;\mu\nu}$ vanish to second order at zero momentum. In a second sense of the words, distinct from the local terms of the formal sums $\mathsf S_3$ and $\mathsf V_2$, the proofs of Lemma~\ref{IV-73-69} and Theorem~\ref{IV-73-71} call local terms those terms of the expansions of the elementary terms that contain one of the kernels $\mathsf q_X$, $\mathsf q_Y$, $\mathsf q_{\mathsf p}$, $\mathsf q_{\mathsf q}$, together with the two halves of $\mathsf T_{4'}$; they are grouped into the local terms at $q$ and those at $p$.

Two integrals over a loop momentum $l$ agree \emph{integrand by integrand} if their integrands coincide before integration, up to substitutions $l\to\pm l+\lambda(\kappa)$ with $\lambda$ linear.

\emph{Model} (Definition~\ref{IV-73-64}, Lemmas~\ref{IV-73-68}, \ref{IV-73-69} and~\ref{IV-73-70}, Theorem~\ref{IV-73-71}). The tensors are
\[
\begin{multlined}
V_2(k_1,k_2)_{\mu\nu}:=-(k_1\cdot k_2)\delta_{\mu\nu}+k_{1\nu}k_{2\mu},\qquad V_3(k_1,k_2,k_3)_{\alpha\beta\gamma}:=\delta_{\alpha\gamma}(k_1-k_3)_\beta+\delta_{\alpha\beta}(k_2 \\
-k_1)_\gamma+\delta_{\beta\gamma}(k_3-k_2)_\alpha,
\end{multlined}
\]
and the variables are $\mathsf x_0:=\mathsf p^2$, $\mathsf x_1:=k^2$, $\mathsf x_1':=(k-\mathsf p)^2$, $\mathsf x_2:=l^2$, $\mathsf x_2':=(l-\mathsf p)^2$, $\mathsf x_3:=(k-l)^2$. The six integrands $F_1,\dots,F_6$ are those of \eqref{IV-73-65}, each with its \emph{lines}, which are the vectors whose squares occur in its denominator, counted with multiplicity. With $\rho_0:=\frac\pi{16}$ and $\rho(\mathsf q):=\theta(|\mathsf q|/2\rho_0)$, $R_\Gamma$ is the product of $\rho$ over the lines of $\Gamma$, and
\[
\Psi_\Gamma(\mathsf p):=\iint R_\Gamma F_\Gamma\,d^\circ k\,d^\circ l,\qquad\mathsf T^{\mathrm{mod}}_\Gamma(u):=24\int_{\mathbb R^3}e^{i\mathsf p\cdot u}\Psi_\Gamma(\mathsf p)\,d^\circ\mathsf p .
\]
The \emph{model} consists of the continuum two-loop integrals $\Psi_\Gamma$, with the smooth ultraviolet cut-off $\rho$ on every line, and of the functions $\mathsf T^{\mathrm{mod}}_\Gamma$. Lemma~\ref{IV-73-68} introduces the \emph{leading integrand} $\mathsf F_G$ of an elementary term $G$, which is the product of the leading parts of the symbols of its lines, the sets $\mathfrak G_2,\mathfrak G_3$ and the symbols $\sigma_\Pi,\sigma_\to,\sigma_\leftarrow$; Lemma~\ref{IV-73-69} the functions $\mathsf U_\Gamma$, $\mathsf l_\Gamma$ and $\Pi^h:=\mathcal F^{-1}_{\ne0}\sigma_\Pi^h$; Lemma~\ref{IV-73-70} the numbers $c_{4,\Gamma}$ and $\kappa_\Gamma^{\mathrm{lat}}$, given by $\mathsf l_\Gamma=c_{4,\Gamma}|u|^{-4}$ and $\mathsf U_\Gamma(u)=\kappa_\Gamma^{\mathrm{lat}}|u|^{-5}$; and Theorem~\ref{IV-73-71} the functions $\mathsf U_{ij}$ and the constant $\kappa^{\mathrm{lat}}:=\sum_{\Gamma=1}^6\kappa_\Gamma^{\mathrm{lat}}$.

\emph{Value.} Lemma~\ref{IV-73-80} introduces $c(s):=\dfrac{2^s\,\Gamma_{\mathrm E}(\frac{3+s}2)}{\pi^{3/2}\,\Gamma_{\mathrm E}(-\frac s2)}$, and Lemma~\ref{IV-73-83} the sums $\mathsf s_n(\mathsf p):=\sum_{u\in\mathbb Z^3\setminus\{0\}}|u|^{-n}e^{-i\mathsf p\cdot u}$ and the functions $\mathsf e_4,\mathsf e_5$. Lemma~\ref{IV-73-84} uses the numbers $\mathsf a_0,\mathsf a_1,\mathsf a_2$, the function $\mathsf o$, and $R_\Gamma^\Lambda(k,l,\mathsf p):=R_\Gamma(k/\Lambda,l/\Lambda,\mathsf p/\Lambda)$. Definition~\ref{IV-73-85} introduces $\mathbb D:=\{\operatorname{Re}\mathsf z>-\frac12\}\setminus\{\frac12,\frac32,\dots\}$, the homogeneous factors of degree $d$, the class $\mathfrak P$, the finite part $\operatorname{f.p.}\!\int\mathsf g\,dk$, that is, the value at $\mathsf z=0$ of the continuation of $\int|k|^{2\mathsf z}\mathsf g\,dk$, the cut-off $\eta_c(k):=\theta(4|k|/c)$, the class $\mathfrak P_2(c)$ defined by ($\Phi0$)--($\Phi3$), the finite part $\operatorname{f.p.}\!\iint$, the class $\mathcal Y(c)$ defined by (Y1), (Y2), the scaled cut-offs $\Upsilon^\Lambda$, and
\[
\overline{G_2}(k,l):=\tfrac16\,\Delta_{\mathsf p}G(k,l,\mathsf p)\big|_{\mathsf p=0},\qquad\Xi_\Gamma[\Upsilon](\mathsf p):=\iint\Upsilon F_\Gamma\,d^\circ k\,d^\circ l .
\]
The homogeneous factors and the classes $\mathfrak P$, $\mathfrak P_2(c)$ and $\mathcal Y(c)$ describe the domain of the finite part in one and in two variables and the admissible cut-offs, which vanish near the origin. Lemma~\ref{IV-73-86}(d) treats a \emph{soft line}, that is, a factor $S(k,\mathsf p)$ of degree $-2$ that is singular at $k=0$ and at $k=\mathsf p$, and it uses $[SJ]_i(k;\mathsf p):=\sum_{|\alpha|=i}\frac{\mathsf p^\alpha}{\alpha!}\partial_{\mathsf p}^\alpha(SJ)(k,0)$ and the function $\mathsf h_S$; Proposition~\ref{IV-73-87} defines $\mathfrak f_\Gamma[\Upsilon]:=(2\pi)^{-6}\operatorname{f.p.}\!\iint\overline{(\Upsilon F_\Gamma)_2}\,dk\,dl$.

Definition~\ref{IV-73-88} introduces the spaces $\mathcal L_n$, the coefficient $r_{111}$ of $(\mathsf x_1\mathsf x_2\mathsf x_3)^{-1}$, the space $\mathfrak F$ and the functional $\mathcal R[F]:=r_{111}(\overline{F_2})$, which takes the coefficient of $(\mathsf x_1\mathsf x_2\mathsf x_3)^{-1}$ in the angular mean of the Taylor coefficient of order $\mathsf p^2$ of $F$; Lemma~\ref{IV-73-90} uses $\mathsf q(t):=\frac13t(2t-1)$; Proposition~\ref{IV-73-94}, the \emph{bridge}, expresses $\kappa_\Gamma^{\mathrm{lat}}$ through $\mathcal R[F_\Gamma]$ and defines $d_{2,\Gamma}:=\mathcal R[F_\Gamma]/16\pi^2$; Definition~\ref{IV-73-99} defines $c_1:=\kappa^{\mathrm{lat}}/\mathsf c_2$; and Lemma~\ref{IV-73-102} uses the \emph{shell}, that is, the annulus $\mathsf{Ann}_r:=\{u\in\mathbb R^3:r\le|u|<Lr\}$ with a real number $L>1$.

The lattice conventions below specify the oriented cells, gauge transformations and differences. A subscript $D$ on a bond or plaquette set denotes restriction to the domain $D$.

The joint cumulants are given by the partition formulas in Definition~\ref{IV-73-19}. The Gibbs kernel and the dilation of the smeared observable are written explicitly where they enter the finite-box argument.

The continuum test pairs $\mathfrak T_{\mathbb R^3}(\mathsf d;\vartheta)$ and the kernels $\mathcal K_q(f,h)$ are those of Section~\ref{C-350-1}. The leading constant is $\mathsf c_2=9/(16\pi^2)$. The curl projection and its Fourier kernel are specified below.



\subsection{Lattice, boxes, chart}

\begin{definition}[Lattice, forms, boxes with Dirichlet data]\label{IV-73-2}
\textup{(1)}

\emph{The lattice.} We work on the lattice $\mathbb Z^3$ of spacing $1$. A bond $b=\langle b_-,b_+\rangle$ is an ordered pair of neighbouring sites, $\bar b:=\langle b_+,b_-\rangle$ is its reverse, and the positively oriented bonds are
\[
(x,\mu):=\langle x,x+e_\mu\rangle\qquad\bigl(x\in\mathbb Z^3,\ \mu\in\{1,2,3\}\bigr);
\]
the base point of $(x,\mu)$ is $x=b_-$. For $x\in\mathbb Z^3$ and $1\le\mu<\nu\le3$ the plaquette $p=p_{\mu\nu}(x)$ has the base point $x_p:=x$, the four corners $x$, $x+e_\mu$, $x+e_\mu+e_\nu$, $x+e_\nu$, and the four bonds
\[
b_1=(x,\mu),\qquad b_2=(x+e_\mu,\nu),\qquad b_3=(x+e_\nu,\mu),\qquad b_4=(x,\nu).
\]
A plaquette is determined by the triple $(x,\mu,\nu)$.

Let $\mathcal B$ be a set of bonds which contains the reverse of each of its elements. A $G$-valued configuration on $\mathcal B$ is a map $U$ on $\mathcal B$ with $U(\bar b)=U(b)^{-1}$; the configuration with the constant value $\mathbf 1$ is again denoted by $\mathbf 1$. If the bonds of a plaquette $p$ belong to $\mathcal B$, we set
\[
\partial U(p):=U(b_1)\,U(b_2)\,U(b_3)^{-1}\,U(b_4)^{-1},
\]
and a $G$-valued function $u$ on a set of sites containing the end points of the bonds of $\mathcal B$ acts on $U$ by the gauge transformation $U^u(b):=u(b_-)\,U(b)\,u(b_+)^{-1}$. A $\mathfrak g$-valued $1$-form is a map $A$ on bonds with $A(\bar b)=-A(b)$, and $e^{iA}$ denotes the configuration $b\mapsto e^{iA(b)}$. The configurations on the set of all bonds form the space $\mathcal V(\mathbb L_1)$; a configuration on a set $B$ of positively oriented bonds, extended to the reverses by $U(\bar b)=U(b)^{-1}$, is an element of $G^B$. We write $A_\mu(x):=A((x,\mu))$, and $A^a_\mu(x)$ for the components of $A_\mu(x)$ in the basis $\mathsf e_1,\mathsf e_2,\mathsf e_3$.

The barycentre of $p=p_{\mu\nu}(x)$ is $x(p)=x+\frac12(e_\mu+e_\nu)$, and its \emph{dual direction} is the index $\lambda(p)\in\{1,2,3\}\setminus\{\mu,\nu\}$. For $x\in\mathbb Z^3$ and $i\in\{1,2,3\}$ we denote by $p_{[i]}(x)$ the plaquette with base point $x$ and dual direction $i$. The plaquettes correspond to the pairs $(z,\lambda)$, the barycentre of $(z,\lambda)$ being $z+\frac12(e_1+e_2+e_3-e_\lambda)$: the map $(z,\lambda)\mapsto p_{\mu\nu}(z)$, where $\{\mu,\nu\}=\{1,2,3\}\setminus\{\lambda\}$, is a bijection which preserves base point, barycentre and dual direction, and it sends $(x,i)$ to $p_{[i]}(x)$. For $y\in\mathbb Z^3$ the translates $b+y$ and $p+y$ are the images of $b$ and $p$ under the translation $\tau_y$.

For $U\in\mathcal V(\mathbb L_1)$ the plaquette action is $a_p(U)=1-\frac12\operatorname{tr}\partial U(p)$. It depends on the values of $U$ at the four bonds of $p$ only, and we use the same formula for every configuration $U$ on a set of bonds which is stable under reversal and contains the bonds of $p$. For every such $U$, and every $G$-valued $u$ on a set of sites containing the end points of the bonds on which $U$ is given,
\[
a_p(U)\in[0,2],\qquad a_p(U^u)=a_p(U).
\]
For the first relation we recall that every $g\in G$ can be written as $e^{iY}$ with $Y\in\mathfrak g$, $|Y|\le\pi$: for $g\ne-\mathbf 1$ this is the exponential chart in Definition~\ref{IV-73-1}(2), and $-\mathbf 1=e^{i\pi\mathsf e_3}$. Since $\frac12\operatorname{tr}e^{iY}=\cos|Y|$, the choice $e^{iY}=\partial U(p)$ gives $a_p(U)=1-\cos|Y|\in[0,2]$. For the second relation let $x$ be the base point of $p$; in the product
\begin{align*}
\partial U^u(p)={}&u(x)\,U(b_1)\,u(x+e_\mu)^{-1}\cdot u(x+e_\mu)\,U(b_2)\,u(x+e_\mu+e_\nu)^{-1}\\
&\cdot u(x+e_\mu+e_\nu)\,U(b_3)^{-1}\,u(x+e_\nu)^{-1}\cdot u(x+e_\nu)\,U(b_4)^{-1}\,u(x)^{-1}
\end{align*}
the inner factors cancel, so that $\partial U^u(p)=u(x)\,\partial U(p)\,u(x)^{-1}$, and the trace is invariant under conjugation. 

\textup{(2)}

\emph{Differences.} For a site function $\lambda$, $\mathfrak g$-valued or complex, and a $1$-form $A$ we set
\[
\begin{gathered}
(\mathrm d\lambda)(b):=\lambda(b_+)-\lambda(b_-),\qquad(\mathrm dA)(p):=A(b_1)+A(b_2)-A(b_3)-A(b_4),\\
(\mathrm d^*A)(x):=\sum_{\mu}\bigl[A_\mu(x-e_\mu)-A_\mu(x)\bigr].
\end{gathered}
\]
The relation $\mathrm d\mathrm d\lambda=0$ holds: at $p=p_{\mu\nu}(x)$,
\begin{align*}
(\mathrm d\mathrm d\lambda)(p)={}&\bigl[\lambda(x+e_\mu)-\lambda(x)\bigr]+\bigl[\lambda(x+e_\mu+e_\nu)-\lambda(x+e_\mu)\bigr]\\
&-\bigl[\lambda(x+e_\mu+e_\nu)-\lambda(x+e_\nu)\bigr]-\bigl[\lambda(x+e_\nu)-\lambda(x)\bigr]=0 .
\end{align*}
On $\mathfrak g$-valued $\lambda$ and $A$ the three operators act componentwise in the basis $\mathsf e_a$. By $(\cdot,\cdot)$ and $\|\cdot\|$ we denote the sums of $\langle\cdot,\cdot\rangle$ over the sites, over the positively oriented bonds or over the plaquettes, according to the arguments; on complex forms $(\cdot,\cdot)$ is the Hermitian sum, antilinear in its first argument. If $\lambda$ and $A$ have finite support, a shift of the summation variable gives
\begin{align*}
(\mathrm d\lambda,A)&=\sum_{x,\mu}\bigl\langle\lambda(x+e_\mu)-\lambda(x),A_\mu(x)\bigr\rangle\\
&=\sum_x\Bigl\langle\lambda(x),\sum_\mu\bigl[A_\mu(x-e_\mu)-A_\mu(x)\bigr]\Bigr\rangle=(\lambda,\mathrm d^*A).
\end{align*}
On plaquette functions, $\mathrm d^*$ denotes the adjoint of $\mathrm d$ on $1$-forms, which is again a difference operator. To write it down, we put $F_{\mu\nu}(x):=F(p_{\mu\nu}(x))$ for a plaquette function $F$ and $\mu<\nu$, and $F_{\nu\mu}(x):=-F_{\mu\nu}(x)$, $F_{\mu\mu}(x):=0$. Since $(\mathrm dA)(p_{\mu\nu}(x))=[A_\nu(x+e_\mu)-A_\nu(x)]-[A_\mu(x+e_\nu)-A_\mu(x)]$, the antisymmetry of $F_{\mu\nu}$ in $(\mu,\nu)$ and a shift of the summation variable give, for $A$ and $F$ of finite support,
\begin{align*}
(\mathrm dA,F)&=\sum_x\sum_{\mu,\nu}\bigl\langle A_\nu(x+e_\mu)-A_\nu(x),F_{\mu\nu}(x)\bigr\rangle\\
&=\sum_{x,\rho}\Bigl\langle A_\rho(x),\sum_\nu\bigl[F_{\nu\rho}(x-e_\nu)-F_{\nu\rho}(x)\bigr]\Bigr\rangle ,
\end{align*}
so that
\[
(\mathrm d^*F)((x,\rho))=-\sum_\nu\bigl[F_{\rho\nu}(x-e_\nu)-F_{\rho\nu}(x)\bigr];
\]
by this formula we apply $\mathrm d^*$ to plaquette functions of arbitrary support. Finally we set $\bar\lambda(b):=\frac12\bigl(\lambda(b_-)+\lambda(b_+)\bigr)$ and $(A\times\bar\lambda)(b):=A(b)\times\bar\lambda(b)$.

\textup{(3)}

\emph{Boxes.} Let $n\ge2$ be an integer, and let
\begin{equation}\label{IV-73-3}
\Lambda_n:=\{0,\dots,n\}^3,\qquad\Lambda_n^\circ:=\{1,\dots,n-1\}^3,\qquad\partial\Lambda_n:=\Lambda_n\setminus\Lambda_n^\circ .
\end{equation}
A positively oriented bond is \emph{free} if both its end points lie in $\Lambda_n$ and at least one of them lies in $\Lambda_n^\circ$; $\mathsf{Bd}_n$ denotes the set of free bonds, and $\mathsf{Pl}_n$ the set of plaquettes all of whose corners lie in $\Lambda_n$. A bond with both end points in $\Lambda_n$ which is neither free nor the reverse of a free bond has both end points in $\partial\Lambda_n$.
\begin{itemize}
\item $\mathcal U_n$ is the set of configurations on the bonds with both end points in $\Lambda_n$ which take the value $\mathbf 1$ at every bond that is neither free nor the reverse of a free bond; restriction to $\mathsf{Bd}_n$ identifies $\mathcal U_n$ with $G^{\mathsf{Bd}_n}$.
\item $\mathsf{Form}_n$ is the space of $1$-forms on the same bonds which vanish at every bond that is neither free nor the reverse of a free bond. $\mathsf X_n$ is the space of $\mathfrak g$-valued functions on $\Lambda_n$ which vanish on $\partial\Lambda_n$, and
\[
\mathcal G_n:=\{u:\Lambda_n\to G:\ u=\mathbf 1\text{ on }\partial\Lambda_n\}.
\]
\item For $A\in\mathsf{Form}_n$, $\mathrm dA$ is considered on $\mathsf{Pl}_n$, and $\mathrm d^*A$ on $\Lambda_n^\circ$, where it involves only bonds with both end points in $\Lambda_n$; we extend $\mathrm d^*A$ by $0$ to $\partial\Lambda_n$, so that $\mathrm d^*A\in\mathsf X_n$. For $\lambda\in\mathsf X_n$ we have $\mathrm d\lambda\in\mathsf{Form}_n$, because $\lambda$ vanishes at both end points of a bond in $\Lambda_n$ which is not free. The relation $(\mathrm d\lambda,A)=(\lambda,\mathrm d^*A)$ for $\lambda\in\mathsf X_n$, $A\in\mathsf{Form}_n$ follows from \textup{(2)}, applied to the extensions of $\lambda$ and $A$ by zero. For $u\in\mathcal G_n$ and $U\in\mathcal U_n$ we have $U^u\in\mathcal U_n$, since $U^u(b)=u(b_-)u(b_+)^{-1}=\mathbf 1$ at a bond in $\Lambda_n$ which is not free.
\item $\mathsf{Form}_n^{\mathbb R}$ and $\mathsf X_n^{\mathbb R}$ denote the real-valued analogues of $\mathsf{Form}_n$ and $\mathsf X_n$; an element of $\mathsf{Form}_n$, resp.\ of $\mathsf X_n$, has three components in $\mathsf{Form}_n^{\mathbb R}$, resp.\ in $\mathsf X_n^{\mathbb R}$. We write $|A|_\infty:=\max_b|A(b)|$ and $|\lambda|_\infty:=\max_x|\lambda(x)|$, and $o_n:=(\lfloor n/2\rfloor,\lfloor n/2\rfloor,\lfloor n/2\rfloor)$ is the centre of the box.
\item Geometrically, the plaquette $p_{\mu\nu}(x)$ is the square $\{x+se_\mu+te_\nu:\ s,t\in[0,1]\}$ with the four corners listed in \textup{(1)}, and the set of a bond $b$ is $[b_-,b_+]:=\{(1-s)b_-+sb_+:\ s\in[0,1]\}$. Write $\mathsf{Pl}(D)$ and $\mathsf{Bd}(D)$ for the plaquettes and positively oriented bonds contained in a set $D$. Then $\mathsf{Pl}_n=\mathsf{Pl}([0,n]^3)$, and $\mathsf{Bd}_n$ is the set of the bonds of $\mathsf{Bd}([0,n]^3)$ with an end point in $\Lambda_n^\circ$. Indeed, the square of a plaquette is the convex hull of its four corners, and $[b_-,b_+]$ is the convex hull of $b_-$ and $b_+$; since the cube $[0,n]^3$ is convex, it contains the square, resp.\ the set $[b_-,b_+]$, if and only if it contains these sites, and the sites of $\mathbb Z^3$ in $[0,n]^3$ are the points of $\Lambda_n$.
\end{itemize}
\end{definition}

\begin{definition}[Wilson measure of a box; truncated two-point function]\label{IV-73-4}
Let $n\ge2$ and $\beta>0$. The \emph{Wilson measure} of the box $\Lambda_n$ is the probability measure
\begin{equation}\label{IV-73-5}
d\mu_{n,\beta}(U):=Z_{n,\beta}^{-1}\exp\Bigl[-\beta\sum_{p\in\mathsf{Pl}_n}a_p(U)\Bigr]\prod_{b\in\mathsf{Bd}_n}dU(b)
\end{equation}
on $\mathcal U_n\cong G^{\mathsf{Bd}_n}$, where $dU(b)$ is the normalised Haar measure in the variable $U(b)$ and $Z_{n,\beta}>0$ is the integral of the exponential with respect to the product measure, so that $\mu_{n,\beta}$ has total mass $1$; its mean is denoted by $\langle\cdot\rangle_{n,\beta}$. For bounded Borel functions $X,Y$ on $\mathcal U_n$ define their joint cumulant by
\[
 \langle X;Y\rangle_{n,\beta}
 =\left.\partial_s\partial_t
       \log\langle e^{sX+tY}\rangle_{n,\beta}\right|_{s=t=0}.
\]
The exponential mean is holomorphic and equals $1$ at the origin, so its logarithm is defined near the origin. Differentiating the logarithm gives
\[
\langle X;Y\rangle_{n,\beta}=\langle XY\rangle_{n,\beta}-\langle X\rangle_{n,\beta}\langle Y\rangle_{n,\beta}.
\]

The measure $\mu_{n,\beta}$ is the Gibbs kernel of the set of free bonds, evaluated at the configuration $\mathbf 1$, in the following sense. For a finite set $B$ of positively oriented bonds define its Gibbs kernel by
\[
 (\gamma_BF)(U)=
 \frac{\int F(V\vee U_{B^c})e^{H_B(V\vee U_{B^c})}\,dV_B}
      {\int e^{H_B(V\vee U_{B^c})}\,dV_B},\qquad
 H_B=-\beta\sum_{p\in\mathsf{Pl}[B]}a_p,
\]
where $V\vee U_{B^c}$ equals $V$ on $B$ and $U$ off $B$, $\mathsf{Pl}[B]$ consists of the plaquettes containing a bond of $B$, and $dV_B$ is product Haar measure on $G^B$. Take $B=\mathsf{Bd}_n$, and for $V\in G^B$ let $V\vee\mathbf 1_{B^c}$ denote the configuration which equals $V$ on $B$ and $\mathbf 1$ at all other positively oriented bonds. Then, for every bounded Borel function $F$ on $\mathcal V(\mathbb L_1)$,
\[
(\gamma_BF)(\mathbf 1)=\int F(V\vee\mathbf 1_{B^c})\,d\mu_{n,\beta}(V).
\]
Indeed, by definition the left side is the quotient of $\int Fe^{H_B}\,dV_B$ by $\int e^{H_B}\,dV_B$, both integrands being evaluated at $V\vee\mathbf 1_{B^c}$. Here $dV_B$ is the product, over the bonds $b\in B$, of the normalised Haar measures in the variables $V(b)$, and $H_B=-\beta\sum_{p\in\mathsf{Pl}[B]}a_p$, the sum running over the set $\mathsf{Pl}[B]$ of the plaquettes which contain a bond of $B$. Every plaquette of $\mathsf{Pl}[B]$, that is, every plaquette containing a free bond $(x,\mu)$, lies in $\mathsf{Pl}_n$: one end point of $(x,\mu)$ lies in $\Lambda_n^\circ$, so both end points have the coordinates $x_\nu\in\{1,\dots,n-1\}$ for $\nu\ne\mu$, and the corners of such a plaquette are the end points of $(x,\mu)$ and their translates by one of the vectors $\pm e_\nu$, $\nu\ne\mu$; all of them lie in $\Lambda_n$.

Conversely, the bonds of a plaquette of $\mathsf{Pl}_n\setminus\mathsf{Pl}[B]$, which contains no free bond, have both end points in $\Lambda_n$ and are not free, so every configuration of $\mathcal U_n$ equals $\mathbf 1$ on them, and $a_p=0$ on $\mathcal U_n$. Hence $H_B(V\vee\mathbf 1_{B^c})=-\beta\sum_{p\in\mathsf{Pl}_n}a_p(V)$, and the quotient is the integral of $F(V\vee\mathbf 1_{B^c})$ with respect to $\mu_{n,\beta}$.
\end{definition}

\begin{definition}[Exponential chart, gauge function, Faddeev--Popov operator]\label{IV-73-6}
Let $\Omega$ be $\mathbb Z^3$ or a box $\Lambda_n$; in the second case forms are taken in $\mathsf{Form}_n$, site functions in $\mathsf X_n$ and plaquettes in $\mathsf{Pl}_n$.

\textup{(1)} For a plaquette $p$, a $1$-form $A$ and $m\ge0$, the number $\mathsf a^{(m)}_p(A)$ is the $m$-th Taylor coefficient of the plaquette action along the curve $t\mapsto e^{itA}$:
\begin{equation}\label{IV-73-7}
a_p(e^{itA})=\sum_{m\ge0}t^m\,\mathsf a^{(m)}_p(A)\qquad(t\in\mathbb R);
\end{equation}
the series converges for every real $t$ by Lemma~\ref{IV-73-9}\textup{(a)}.

\textup{(2)} If a $1$-form $A$ and a site function $\lambda$ satisfy $|A(b)|+|\lambda(b_-)|+|\lambda(b_+)|<1$ for every bond $b$, we set
\[
A^\lambda(b):=\tfrac1i\log\bigl(e^{i\lambda(b_-)}e^{iA(b)}e^{-i\lambda(b_+)}\bigr),\qquad\text{that is,}\qquad e^{iA^\lambda}=(e^{iA})^u,\quad u=e^{i\lambda}.
\]
This is well defined: by the product estimate in Definition~\ref{IV-73-1}(2),
\[
\bigl|e^{i\lambda(b_-)}e^{iA(b)}e^{-i\lambda(b_+)}-\mathbf 1\bigr|\le e^{|\lambda(b_-)|+|A(b)|+|\lambda(b_+)|}-1<e-1<2=|-\mathbf 1-\mathbf 1| ,
\]
so the argument of the logarithm lies in $G\setminus\{-\mathbf 1\}$, where $\frac1i\log$, the inverse of $Y\mapsto e^{iY}$ on $\mathbb B_\pi$, is defined (the exponential chart in Definition~\ref{IV-73-1}(2)).

\textup{(3)} The \emph{gauge function} is the map $A\mapsto\mathrm d^*A$. For a $1$-form $A$ with $|A(b)|<1$ at every bond $b$, the Faddeev--Popov operator $\mathsf M(A)$ acts on site functions $\lambda$ by
\begin{equation}\label{IV-73-8}
\mathsf M(A)\lambda:=\frac d{ds}\,\mathrm d^*\bigl(A^{s\lambda}\bigr)\Big|_{s=0},\qquad\mathsf M(tA)=\sum_{j\ge0}t^j\,\mathsf M^{(j)}(A);
\end{equation}
the second formula is the expansion at $t=0$ of the real-analytic function $t\mapsto\mathsf M(tA)$ (Lemma~\ref{IV-73-9}\textup{(d)}), and it defines the operators $\mathsf M^{(j)}(A)$; since $\mathsf M(t(cA))=\mathsf M((tc)A)$, the operator $\mathsf M^{(j)}(A)$ is homogeneous of degree $j$ in $A$. The derivative is taken site by site: $\mathrm d^*(A^{s\lambda})(x)$ involves the six bonds with end point $x$, and for each of them $A^{s\lambda}(b)$ is defined and smooth in $s$ for $|s|$ small. In particular the definition applies to every site function $\lambda$, bounded or not.
\end{definition}

\begin{lemma}[Coefficients of the chart]\label{IV-73-9}
Let $p$ be a plaquette with bonds $b_1,\dots,b_4$, let $A$ be a $1$-form, $A_i:=A(b_i)$, and $(s_1,s_2,s_3,s_4):=(+,+,-,-)$, so that $\mathrm dA(p)=\sum_is_iA_i$.
\begin{enumerate}
\item[\textup{(a)}] The series \eqref{IV-73-7} converges for every $t\in\mathbb R$. Each $\mathsf a^{(m)}_p$ is a homogeneous polynomial of degree $m$ in the twelve components of $A_1,\dots,A_4$, with real coefficients which do not depend on $p$. Moreover $\mathsf a^{(0)}_p=\mathsf a^{(1)}_p=0$, and $\mathsf a^{(m)}_p(gAg^{-1})=\mathsf a^{(m)}_p(A)$ for every $g\in G$, where $gAg^{-1}$ denotes the form $b\mapsto gA(b)g^{-1}$.
\item[\textup{(b)}] On the set $\{(t,A_1,\dots,A_4)\in\mathbb R\times\mathfrak g^4:\ |t|\sum_i|A_i|<\log3\}$ the function $Z:=\frac1i\log\partial(e^{itA})(p)\in\mathbb B_\pi$ is defined and smooth. For $j=0,1,2$ the maps $\mathfrak z_j:=\frac1{(j+1)!}\partial_t^{j+1}Z|_{t=0}$ are homogeneous polynomial maps $\mathfrak g^4\to\mathfrak g$ of degree $j+1$, which do not depend on $p$; $\mathfrak z_0=\mathrm dA(p)$, and
\begin{equation}\label{IV-73-10}
\begin{gathered}
\mathsf a^{(2)}_p=\tfrac12|\mathrm dA(p)|^2,\qquad\mathsf a^{(3)}_p=\langle\mathrm dA(p),\mathfrak z_1\rangle,\\
\mathsf a^{(4)}_p=\tfrac12|\mathfrak z_1|^2+\langle\mathrm dA(p),\mathfrak z_2\rangle-\tfrac1{24}|\mathrm dA(p)|^4,\qquad\mathfrak z_1=-\sum_{i<j}s_is_j\,A_i\times A_j .
\end{gathered}
\end{equation}
We write $\mathfrak z_j(p)$ for the value of $\mathfrak z_j$ at $\bigl(A(b_1(p)),\dots,A(b_4(p))\bigr)$.
\item[\textup{(c)}] The function $Y\mapsto-\log\sigma_1(|Y|^2)$ is real analytic and even on $\mathbb B_\pi$, and its Taylor polynomial of degree $3$ at $0$ is $\frac13|Y|^2$.
\item[\textup{(d)}] For $Y\in\mathfrak g$ let $\mathsf a_Y:=-2\,Y\times\cdot$ on $\mathfrak g$. Then $e^{isY}Xe^{-isY}=e^{s\mathsf a_Y}X$ for $s\in\mathbb R$ and $X\in\mathfrak g$. If $|A(b)|<1$ and $\lambda$ is a site function, then
\begin{equation}\label{IV-73-11}
\begin{aligned}
\frac d{ds}A^{s\lambda}(b)\Big|_{s=0}&=\mathsf{dE}(A(b))^{-1}\bigl(\lambda(b_-)-e^{\mathsf a_{A(b)}}\lambda(b_+)\bigr)\\
&=-\mathrm d\lambda(b)+2A(b)\times\bar\lambda(b)-\tfrac13A(b)\times\bigl(A(b)\times\mathrm d\lambda(b)\bigr)+\mathsf r(A(b))\,\mathrm d\lambda(b),
\end{aligned}
\end{equation}
where $\mathsf r:\mathbb B_\pi\to\operatorname{End}\mathfrak g$ is real analytic and vanishes at $0$ to fourth order. Consequently, if $|A(b)|<1$ at every bond, $\mathsf M(tA)$ is real analytic in $t$, and
\[
\begin{gathered}
\mathsf M^{(0)}=-\mathrm d^*\mathrm d,\qquad\mathsf M^{(1)}(A)\lambda=2\,\mathrm d^*(A\times\bar\lambda),\\
\mathsf M^{(2)}(A)\lambda=-\tfrac13\,\mathrm d^*\bigl(A\times(A\times\mathrm d\lambda)\bigr),\qquad\mathsf M^{(3)}(A)=0 .
\end{gathered}
\]
\end{enumerate}
\end{lemma}

\begin{proof}
\emph{Step 1: claim \textup{(a)}.} Since $U(b_3)^{-1}=e^{-itA_3}$ and $U(b_4)^{-1}=e^{-itA_4}$ for $U=e^{itA}$, we have $\partial(e^{itA})(p)=\prod_{i=1}^4e^{its_iA_i}$, a product of four exponential series in $t$ which converge for every $t$; their Cauchy product converges for every $t$, and hence so does \eqref{IV-73-7}. For $m\ge1$ the coefficient of $t^m$ is
\[
\mathsf a^{(m)}_p(A)=-\frac12\sum_{m_1+\dots+m_4=m}\operatorname{tr}\Bigl[\prod_{i=1}^4\frac{(is_iA_i)^{m_i}}{m_i!}\Bigr],
\]
that is, $-\frac12\operatorname{tr}$ of a sum of products of $m$ factors $is_iA_i$. It is therefore a homogeneous polynomial of degree $m$ in the components of $A_1,\dots,A_4$, given by the same formula for every plaquette. It takes real values on real arguments, since $a_p$ is real; hence, if we write it as $\sum_\alpha c_\alpha y^\alpha$ in the twelve real components $y$ of $A_1,\dots,A_4$, the polynomial $\sum_\alpha(\operatorname{Im}c_\alpha)\,y^\alpha$ vanishes on $\mathbb R^{12}$, so that $\operatorname{Im}c_\alpha=0$ and the coefficients are real. Replacing every $A_i$ by $gA_ig^{-1}$ conjugates each product of the factors $is_iA_i$ by $g$, which does not change the trace. The relations $\mathsf a^{(0)}_p=\mathsf a^{(1)}_p=0$ follow from \textup{(b)}, see Step~5.

\emph{Step 2: existence and smoothness of $Z$.} Let $X_i:=its_iA_i$ and $\mathsf x:=\partial(e^{itA})(p)=e^{X_1}e^{X_2}e^{X_3}e^{X_4}$. Since $|X_i|=|t|\,|A_i|$, the product estimate in Definition~\ref{IV-73-1}(2) gives on the set of \textup{(b)}
\[
|\mathsf x-\mathbf 1|\le e^{|t|\sum_i|A_i|}-1<e^{\log3}-1=2=|-\mathbf 1-\mathbf 1| ;
\]
hence $\mathsf x\ne-\mathbf 1$, and $Z=\frac1i\log\mathsf x\in\mathbb B_\pi$ is defined. The matrix $\mathsf x$ depends smoothly on $(t,A_1,\dots,A_4)$, and $W\mapsto\frac1i\log W$ is smooth on $G\setminus\{-\mathbf 1\}$, the inverse of the diffeomorphism $Y\mapsto e^{iY}$ of $\mathbb B_\pi$ onto $G\setminus\{-\mathbf 1\}$ (Definition~\ref{IV-73-1}); thus $Z$ is smooth. At $t=0$ we have $\mathsf x=\mathbf 1$ and $Z=0$. The maps $\mathfrak z_j$ are $t$-derivatives of $Z$ at $t=0$, hence smooth maps $\mathfrak g^4\to\mathfrak g$, given by the same formula for all plaquettes.

\emph{Step 3: homogeneity.} For $c>0$ the substitution $(t,A)\mapsto(t/c,cA)$ preserves the domain of $Z$ and leaves $\mathsf x$ unchanged. Hence $Z(t/c,cA)=Z(t,A)$, that is, $Z(s,cA)=Z(cs,A)$, and differentiating $j+1$ times at $s=0$ we find $\mathfrak z_j(cA)=c^{j+1}\mathfrak z_j(A)$. It remains to use the following fact:

\emph{a smooth function $f$ on a finite-dimensional real vector space with $f(cA)=c^mf(A)$ for all $c>0$ is a homogeneous polynomial of degree $m$.} To see it, let $f_i$ denote the homogeneous Taylor polynomials of $f$ at $0$; Taylor's formula at the points $cA$ gives
\[
c^mf(A)=f(cA)=\sum_{i\le m}c^if_i(A)+O(c^{m+1})\qquad(c\to0),
\]
and comparing the powers of $c$ as $c\to0$ we obtain $f_i(A)=0$ for $i<m$ and $f(A)=f_m(A)$. Applied to the components of $\mathfrak z_j$ in a basis of $\mathfrak g$, this shows that $\mathfrak z_j$ is a homogeneous polynomial map of degree $j+1$.

\emph{Step 4: $\mathfrak z_0$ and $\mathfrak z_1$.} At fixed $A$, Taylor's formula in $t$ gives $Z=t(\mathfrak z_0+t\mathfrak z_1+t^2\mathfrak z_2)+O(t^4)$. We compare the two expansions of $e^{iZ}=\mathsf x$,
\[
e^{iZ}=\mathbf 1+iZ-\tfrac12Z^2+O(t^3),\qquad\mathsf x=\mathbf 1+\sum_iX_i+\tfrac12\sum_iX_i^2+\sum_{i<j}X_iX_j+O(t^3).
\]
The terms of order $t$ give $it\mathfrak z_0=\sum_iX_i$, hence $\mathfrak z_0=\sum_is_iA_i=\mathrm dA(p)$. Since $-\frac12Z^2=\frac12(it\mathfrak z_0)^2+O(t^3)=\frac12\bigl(\sum_iX_i\bigr)^2+O(t^3)$, the terms of order $t^2$ give
\begin{align*}
it^2\mathfrak z_1&=\tfrac12\sum_iX_i^2+\sum_{i<j}X_iX_j-\tfrac12\Bigl(\sum_iX_i\Bigr)^2=\tfrac12\sum_{i<j}[X_i,X_j]\\
&=-\tfrac12t^2\sum_{i<j}s_is_j[A_i,A_j]=-it^2\sum_{i<j}s_is_j\,A_i\times A_j ,
\end{align*}
where $[X,Y]=2i\,X\times Y$ was used in the last step. This is the formula for $\mathfrak z_1$ in \eqref{IV-73-10}.

\emph{Step 5: the coefficients $\mathsf a^{(m)}_p$, $m\le4$.} Since $\frac12\operatorname{tr}e^{iZ}=\cos|Z|$, we have $a_p(e^{itA})=1-\cos|Z|=\frac12|Z|^2-\frac1{24}|Z|^4+O(t^6)$. By Step~4,
\[
|Z|^2=t^2\Bigl[|\mathfrak z_0|^2+2t\langle\mathfrak z_0,\mathfrak z_1\rangle+t^2\bigl(|\mathfrak z_1|^2+2\langle\mathfrak z_0,\mathfrak z_2\rangle\bigr)+O(t^3)\Bigr],
\]
and therefore
\[
a_p(e^{itA})=\tfrac12t^2|\mathfrak z_0|^2+t^3\langle\mathfrak z_0,\mathfrak z_1\rangle+t^4\Bigl(\tfrac12|\mathfrak z_1|^2+\langle\mathfrak z_0,\mathfrak z_2\rangle-\tfrac1{24}|\mathfrak z_0|^4\Bigr)+O(t^5).
\]
Comparing with the Taylor series \eqref{IV-73-7} we obtain $\mathsf a^{(0)}_p=\mathsf a^{(1)}_p=0$ and \eqref{IV-73-10}.

\emph{Step 6: claim \textup{(c)}.} From its power series, $\sigma_1(w)=\bigl(1-\frac w6+O(w^2)\bigr)^2=1-\frac w3+O(w^2)$, and $\sigma_1(|Y|^2)=(\sin|Y|/|Y|)^2>0$ on $\mathbb B_\pi$. Thus $Y\mapsto-\log\sigma_1(|Y|^2)$ is the composite of the real analytic maps $Y\mapsto|Y|^2$, $\sigma_1$ at real arguments, and $\log$ on $(0,\infty)$, hence real analytic. It depends on $|Y|^2$ only and is therefore even, and $-\log\bigl(1-\frac13|Y|^2+O(|Y|^4)\bigr)=\frac13|Y|^2+O(|Y|^4)$ shows that its Taylor polynomial of degree $3$ at $0$ is $\frac13|Y|^2$.

\emph{Step 7: conjugation.} By $[Y,X]=2i\,Y\times X$ we have $i[Y,X]=-2\,Y\times X=\mathsf a_YX$. The function $x(s):=e^{isY}Xe^{-isY}$ takes values in $\mathfrak g$ and satisfies $x'(s)=i[Y,x(s)]=\mathsf a_Yx(s)$, so that $\frac d{ds}\bigl(e^{-s\mathsf a_Y}x(s)\bigr)=0$; together with $x(0)=X$ this gives $x(s)=e^{s\mathsf a_Y}X$. Consequently
\[
\mathsf{dE}(Y)=\int_0^1e^{s\mathsf a_Y}\,ds=\mathsf f(\mathsf a_Y),\qquad\mathsf f(w):=\frac{e^w-1}w=\sum_{n\ge0}\frac{w^n}{(n+1)!}.
\]

\emph{Step 8: first equality in \eqref{IV-73-11}.} Let $Y:=A(b)$ and $\delta:=\frac d{ds}A^{s\lambda}(b)|_{s=0}$. For $|s|$ small, $A^{s\lambda}(b)$ is defined, depends smoothly on $s$, equals $Y$ at $s=0$, and satisfies $e^{iA^{s\lambda}(b)}=e^{is\lambda(b_-)}e^{iY}e^{-is\lambda(b_+)}$. We differentiate this identity at $s=0$, using the differential formula in Definition~\ref{IV-73-1}(2) on the left and Step~7 on the right:
\[
i\bigl(\mathsf{dE}(Y)\delta\bigr)e^{iY}=i\lambda(b_-)e^{iY}-ie^{iY}\lambda(b_+)=i\bigl(\lambda(b_-)-e^{\mathsf a_Y}\lambda(b_+)\bigr)e^{iY}.
\]
By the determinant formula in Definition~\ref{IV-73-1}(2), $\det\mathsf{dE}(Y)=\sigma_1(|Y|^2)>0$, so $\mathsf{dE}(Y)$ is invertible, and the first equality follows.

\emph{Step 9: expansion.} The double cross product formula $X\times(Y\times Z)=\langle X,Z\rangle Y-\langle X,Y\rangle Z$ gives $\mathsf a_YY=0$ and $\mathsf a_Y^2=-4|Y|^2$ on the orthogonal complement of $Y$; so $\mathsf a_Y$ has the eigenvalues $0$ and $\pm2i|Y|$, and Lagrange's identity $|Y\times X|^2=|Y|^2|X|^2-\langle Y,X\rangle^2$ gives $\|\mathsf a_Y\|\le2|Y|<2\pi$ on $\mathbb B_\pi$. The function $1/\mathsf f(w)=w/(e^w-1)$ is holomorphic for $|w|<2\pi$, and since $w/(e^w-1)+w/2$ is even,
\[
\frac1{\mathsf f(w)}=1-\frac w2+\frac{w^2}{12}+w^4\mathsf f_4(w),\qquad\frac{we^w}{e^w-1}=w+\frac w{e^w-1}\qquad(|w|<2\pi),
\]
with $\mathsf f_4$ holomorphic for $|w|<2\pi$. Evaluating these power series at $\mathsf a_Y$, Steps~7 and~8 give
\[
\begin{aligned}
\delta&=\tfrac1{\mathsf f}(\mathsf a_Y)\lambda(b_-)-\bigl(\mathsf a_Y+\tfrac1{\mathsf f}(\mathsf a_Y)\bigr)\lambda(b_+)=-\tfrac1{\mathsf f}(\mathsf a_Y)\,\mathrm d\lambda(b)-\mathsf a_Y\lambda(b_+)\\
&=-\mathrm d\lambda(b)-\mathsf a_Y\bigl(\lambda(b_+)-\tfrac12\mathrm d\lambda(b)\bigr)-\tfrac1{12}\mathsf a_Y^2\,\mathrm d\lambda(b)-\mathsf a_Y^4\mathsf f_4(\mathsf a_Y)\,\mathrm d\lambda(b).
\end{aligned}
\]
Now $\lambda(b_+)-\frac12\mathrm d\lambda(b)=\bar\lambda(b)$, $-\mathsf a_Y\bar\lambda(b)=2Y\times\bar\lambda(b)$ and $\mathsf a_Y^2=4\,Y\times(Y\times\cdot)$; with $\mathsf r(Y):=-\mathsf a_Y^4\mathsf f_4(\mathsf a_Y)$ this is the second equality in \eqref{IV-73-11}. The map $\mathsf r$ is real analytic on $\mathbb B_\pi$, being the sum of a power series at an operator which depends linearly on $Y$ and has norm $<2\pi$, and it vanishes at $0$ to fourth order because of the factor $\mathsf a_Y^4$. Since $\mathrm d^*$ involves at each site finitely many bonds, $\mathsf M(A)\lambda=\mathrm d^*\delta$, where $\delta$ is the form given at every bond by \eqref{IV-73-11}. If $A$ is replaced by $tA$, the first three terms of \eqref{IV-73-11} are of degree $0$, $1$, $2$ in $t$, and the last one is real analytic in $t$ and of order $t^4$. This proves that $t\mapsto\mathsf M(tA)$ is real analytic, and it gives the formulas for $\mathsf M^{(0)},\dots,\mathsf M^{(3)}$.
\end{proof}

\begin{lemma}[Free gauge action, absence of zero modes, a slice, and the Faddeev--Popov change of variables in a box]\label{IV-73-12}
Let $n\ge2$. On $\mathcal U_n$ let $S(U):=\sum_{p\in\mathsf{Pl}_n}a_p(U)$, and let
\[
\begin{gathered}
\mathsf T_n:=\{\langle x-e_1,x\rangle:\ x\in\Lambda_n^\circ\}\subset\mathsf{Bd}_n,\\
\mathcal U_n^{\mathsf T}:=\{V\in\mathcal U_n:\ V=\mathbf 1\text{ on }\mathsf T_n\},\qquad\mathsf{Form}_n^{\mathsf T}:=\{A\in\mathsf{Form}_n:\ A=0\text{ on }\mathsf T_n\}.
\end{gathered}
\]
\begin{enumerate}
\item[\textup{(a)}] A positively oriented bond $(x,\mu)$ with both end points in $\Lambda_n$ is free if and only if $x_\nu\in\{1,\dots,n-1\}$ for both $\nu\ne\mu$; otherwise it lies in a face $\{x_\nu=0\}$ or $\{x_\nu=n\}$ of the box, for some $\nu\ne\mu$.
\item[\textup{(b)}] The map $\mathsf j:\mathcal G_n\times\mathcal U_n^{\mathsf T}\to\mathcal U_n$, $(u,V)\mapsto V^u$, is a diffeomorphism, and for every bounded Borel function $F$ on $\mathcal U_n$
\begin{equation}\label{IV-73-13}
\int F(U)\prod_{b\in\mathsf{Bd}_n}dU(b)=\int\prod_{x\in\Lambda_n^\circ}du(x)\int\prod_{b\in\mathsf{Bd}_n\setminus\mathsf T_n}dV(b)\;F(V^u).
\end{equation}
The group $\mathcal G_n$ acts freely on $\mathcal U_n$. One has $S(U)=0$ if and only if $U=\mathbf 1^u$ for some $u\in\mathcal G_n$; the only zero of $S$ in $\mathcal U_n^{\mathsf T}$ is $\mathbf 1$.
\item[\textup{(c)}] If $A\in\mathsf{Form}_n$ satisfies $\mathrm dA=0$ on $\mathsf{Pl}_n$ and $\mathrm d^*A=0$, then $A=0$. Hence, for every $\tau_{\mathrm{gf}}>0$, the quadratic form $\frac12\|\mathrm dA\|^2+\frac1{2\tau_{\mathrm{gf}}}\|\mathrm d^*A\|^2$ is positive definite on $\mathsf{Form}_n$, and $\mathrm d^*\mathrm d$ is positive definite on $\mathsf X_n$.
\item[\textup{(d)}] There are numbers $0<\varrho_2,\varrho_3\le\varrho_1\le\frac1{20}$, depending on $n$, with the following properties, where
\[
\mathfrak D:=\{(\lambda,A)\in\mathsf X_n\times\mathsf{Form}_n^{\mathsf T}:\ |\lambda|_\infty<\varrho_1,\ |A|_\infty<\varrho_3\}.
\]
\begin{enumerate}
\item[\textup{(i)}] The map $\Xi:(\lambda,A)\mapsto A^\lambda$ extends to a diffeomorphism of a neighbourhood of $\overline{\mathfrak D}$ onto an open subset of $\{A'\in\mathsf{Form}_n:\ |A'|_\infty<1\}$.
\item[\textup{(ii)}] $\det\mathsf M(A')\ne0$ for every $A'\in\overline{\mathcal W}$, where $\mathcal W:=\Xi(\mathfrak D)$.
\item[\textup{(iii)}] If $A\in\mathsf{Form}_n^{\mathsf T}$ and $|A|_\infty\le\varrho_3$, then $\lambda\mapsto\mathrm d^*(A^\lambda)$ is injective on $\{\lambda\in\mathsf X_n:\ |\lambda|_\infty\le\varrho_1\}$, and the image of $\{\lambda\in\mathsf X_n:\ |\lambda|_\infty<\varrho_1\}$ contains $\{c\in\mathsf X_n:\ |c|_\infty<\varrho_2\}$.
\item[\textup{(iv)}] For every bounded Borel function $\Phi$ on $\mathcal W$ and every bounded Borel function $w\ge0$ on $\mathsf X_n$,
\begin{equation}\label{IV-73-14}
\begin{aligned}
&\int_{\mathcal W}\Phi(A')\,w(\mathrm d^*A')\,|\det\mathsf M(A')|\prod_{b\in\mathsf{Bd}_n}\sigma_1(|A'(b)|^2)\,dA'\\
&\qquad=\int_{|A|_\infty<\varrho_3}\prod_{b\in\mathsf{Bd}_n\setminus\mathsf T_n}\sigma_1(|A(b)|^2)\,dA\int_{\{\mathrm d^*(A^\lambda):\,|\lambda|_\infty<\varrho_1\}}w(c)\,\Phi\bigl(A^{\lambda(c)}\bigr)\,dc ,
\end{aligned}
\end{equation}
where $\lambda(c)$ is the solution of $\mathrm d^*(A^\lambda)=c$ with $|\lambda|_\infty<\varrho_1$, unique by \textup{(iii)}, and $dA'$, $dA$, $dc$ are the Lebesgue measures of the Euclidean spaces $\mathsf{Form}_n$, $\mathsf{Form}_n^{\mathsf T}$, $\mathsf X_n$, with the scalar product $(\cdot,\cdot)$ of Definition~\ref{IV-73-2}\textup{(2)}. Thus $dA'=\prod_{b\in\mathsf{Bd}_n}dA'(b)$, where $dA'(b)$ is the Lebesgue measure $dY$ of $(\mathfrak g,\langle\cdot,\cdot\rangle)$ in the variable $A'(b)$; likewise $dA$ is the product of these measures over the bonds of $\mathsf{Bd}_n\setminus\mathsf T_n$, and $dc$ the product over the sites of $\Lambda_n^\circ$.
\end{enumerate}
\end{enumerate}
\end{lemma}

\begin{proof}
\emph{Step 1: claim \textup{(a)}.} Let $(x,\mu)$ be a positively oriented bond with both end points in $\Lambda_n$; then $x_\mu\in\{0,\dots,n-1\}$, and the end points $x$, $x+e_\mu$ have the same coordinates $x_\nu$, $\nu\ne\mu$. If both these coordinates lie in $\{1,\dots,n-1\}$, then, as $n\ge2$, one of the numbers $x_\mu$, $x_\mu+1$ lies in $\{1,\dots,n-1\}$, and the corresponding end point lies in $\Lambda_n^\circ$: the bond is free. If $x_\nu\in\{0,n\}$ for some $\nu\ne\mu$, then both end points have this coordinate, neither lies in $\Lambda_n^\circ$, and the bond lies in the face $\{x_\nu=0\}$ or $\{x_\nu=n\}$.

\emph{Step 2: $\mathsf j$ is bijective and smooth.} Let $U\in\mathcal U_n$. For $u'\in\mathcal G_n$ the condition $U^{u'}=\mathbf 1$ on $\mathsf T_n$ is equivalent to
\[
u'(x)=u'(x-e_1)\,U(\langle x-e_1,x\rangle)\qquad(x\in\Lambda_n^\circ).
\]
Since $u'=\mathbf 1$ at the sites with first coordinate $0$, which lie in $\partial\Lambda_n$, this recursion in $x_1=1,\dots,n-1$ determines $u'$ uniquely: $u'(x)$ is the ordered product of the values of $U$ at the bonds $\langle y-e_1,y\rangle$ on the line from $(0,x_2,x_3)$ to $x$. Let $V:=U^{u'}\in\mathcal U_n^{\mathsf T}$ and $u:=u'^{-1}\in\mathcal G_n$; since $(U^v)^w=U^{wv}$, we have $V^u=U$. Conversely, if $U=V^u$ with $(u,V)\in\mathcal G_n\times\mathcal U_n^{\mathsf T}$, then $U^{u^{-1}}=V$ equals $\mathbf 1$ on $\mathsf T_n$, so $u^{-1}=u'$ and $V=U^{u'}$. Hence $\mathsf j$ is bijective. Both $\mathsf j$ and $\mathsf j^{-1}$ are composed of products and inverses of the values of $u$, $V$, $U$, and are therefore smooth (Definition~\ref{IV-73-1}).

\emph{Step 3: the action is free.} If $U^u=U$ for some $u\in\mathcal G_n$, then $u(b_+)=U(b)^{-1}u(b_-)U(b)$ at every bond $b\in\mathsf T_n$, i.e.\ $u(x)=U(\langle x-e_1,x\rangle)^{-1}u(x-e_1)\,U(\langle x-e_1,x\rangle)$ for $x\in\Lambda_n^\circ$; the recursion of Step~2, started from $u=\mathbf 1$ at $x_1=0$, gives $u=\mathbf 1$.

\emph{Step 4: the identity \eqref{IV-73-13}.} We fix $u\in\mathcal G_n$ on the right side of \eqref{IV-73-13} and substitute $U(b):=u(b_-)V(b)u(b_+)^{-1}$ for $b\in\mathsf{Bd}_n\setminus\mathsf T_n$. By Fubini's theorem and the invariance of the Haar measure under the maps $g\mapsto hgh'^{-1}$, $h,h'\in G$ (left and right invariance of Haar measure), the inner integral becomes $\int\prod_{b\in\mathsf{Bd}_n\setminus\mathsf T_n}dU(b)\,F(U)$, where $U$ takes the values $U(\langle x-e_1,x\rangle)=u(x-e_1)u(x)^{-1}$ on $\mathsf T_n$. In the integral over $u$ we now replace, on each line parallel to $e_1$ and for $x_1=n-1,n-2,\dots,1$ in this order, the variable $u(x)$ by the variable $U(\langle x-e_1,x\rangle)=u(x-e_1)u(x)^{-1}$, at fixed $u(x-e_1)$.

When this is done at $x$, the variable $u(x)$ enters no other value of $U$ on $\mathsf T_n$ (the value at $\langle x,x+e_1\rangle$, if $x+e_1\in\Lambda_n^\circ$, has already been replaced), and the map $u(x)\mapsto u(x-e_1)u(x)^{-1}$, an inversion followed by a left translation, preserves the Haar measure (left and right invariance of Haar measure). After the last substitution we arrive at the left side of \eqref{IV-73-13}.

\emph{Step 5: zeros of $S$.} If $\partial U(p)=e^{iY}$ with $|Y|\le\pi$, then $a_p(U)=1-\cos|Y|$; hence $a_p(U)=0$ if and only if $\partial U(p)=\mathbf 1$, and since $a_p\ge0$, $S(U)=0$ if and only if $\partial U(p)=\mathbf 1$ for every $p\in\mathsf{Pl}_n$. By the gauge invariance of $a_p$ (Definition~\ref{IV-73-2}\textup{(1)}), $S(\mathbf 1^u)=S(\mathbf 1)=0$ for $u\in\mathcal G_n$. Now let $V\in\mathcal U_n^{\mathsf T}$ with $S(V)=0$. We show that $V$ equals $\mathbf 1$ at every positively oriented bond with both end points in $\Lambda_n$, direction by direction.
\begin{itemize}
\item Bonds $(y,1)$ with $y_1\le n-2$: if $y+e_1\in\Lambda_n^\circ$, then $(y,1)\in\mathsf T_n$; otherwise one of $y_2$, $y_3$ lies in $\{0,n\}$, and $(y,1)$ is not free by \textup{(a)}. In both cases $V((y,1))=\mathbf 1$.
\item Bonds $(y,\nu)$, $\nu\in\{2,3\}$: for $y_1\le n-2$ the plaquette $p_{1\nu}(y)$ lies in $\mathsf{Pl}_n$, and $\partial V(p_{1\nu}(y))=\mathbf 1$ together with $V((y,1))=V((y+e_\nu,1))=\mathbf 1$ gives $V((y+e_1,\nu))=V((y,\nu))$. The bonds $(y,\nu)$ with $y_1=0$ are not free, so $V((y,\nu))=\mathbf 1$ for $y_1=0$, and by induction for $y_1\le n-1$; the bonds with $y_1=n$ are not free.
\item Bonds $(y,1)$ with $y_1=n-1$: for $y_2\le n-1$ the plaquette $p_{12}(y)$ gives $V((y+e_2,1))=V((y,1))$, because $V=\mathbf 1$ at the bonds in direction $2$. The bonds with $y_2=0$ are not free, and induction in $y_2$ gives $V((y,1))=\mathbf 1$ for $y_2\le n-1$; the bonds with $y_2=n$ are not free.
\end{itemize}
Hence $V=\mathbf 1$. Finally let $U\in\mathcal U_n$ with $S(U)=0$. By \textup{(b)}, $U=V^u$ with $(u,V)\in\mathcal G_n\times\mathcal U_n^{\mathsf T}$, and $S(V)=S(U^{u^{-1}})=S(U)=0$; thus $V=\mathbf 1$ and $U=\mathbf 1^u$.

\emph{Step 6: claim \textup{(c)}.} Let $A\in\mathsf{Form}_n$ with $\mathrm dA=0$ on $\mathsf{Pl}_n$ and $\mathrm d^*A=0$. We define $\lambda\in\mathsf X_n$ by $\lambda=0$ on $\partial\Lambda_n$ and $\lambda(x)=\lambda(x-e_1)+A(\langle x-e_1,x\rangle)$ for $x\in\Lambda_n^\circ$, recursively in $x_1$. Then $A':=A-\mathrm d\lambda\in\mathsf{Form}_n$ vanishes on $\mathsf T_n$, and $\mathrm dA'=\mathrm dA-\mathrm d\mathrm d\lambda=0$ on $\mathsf{Pl}_n$. The walk of Step~5, in additive form, shows $A'=0$: $A'$ vanishes at the bonds $(y,1)$ with $y_1\le n-2$, which lie in $\mathsf T_n$ or are not free; for $\nu\in\{2,3\}$ and $y_1\le n-2$,
\[
0=\mathrm dA'(p_{1\nu}(y))=A'((y+e_1,\nu))-A'((y,\nu)),
\]
so $A'=0$ at the bonds in the directions $2$ and $3$; and for $y_1=n-1$,
\[
0=\mathrm dA'(p_{12}(y))=A'((y,1))-A'((y+e_2,1)),
\]
so $A'=0$ at the remaining bonds. Hence $A=\mathrm d\lambda$, and $0=(\lambda,\mathrm d^*A)=(\lambda,\mathrm d^*\mathrm d\lambda)=\|\mathrm d\lambda\|^2$, so that $A=\mathrm d\lambda=0$. Consequently, for $\tau_{\mathrm{gf}}>0$ the quadratic form $\frac12\|\mathrm dA\|^2+\frac1{2\tau_{\mathrm{gf}}}\|\mathrm d^*A\|^2$ is non-negative and vanishes only at $A=0$; it is therefore positive definite. Similarly, $(\lambda,\mathrm d^*\mathrm d\lambda)=\|\mathrm d\lambda\|^2\ge0$ for $\lambda\in\mathsf X_n$, and $\mathrm d\lambda=0$ implies $\lambda(x)=\lambda(x-e_1)$ for $x\in\Lambda_n^\circ$, hence $\lambda=0$, since $\lambda=0$ at $x_1=0$; so $\mathrm d^*\mathrm d$ is positive definite on $\mathsf X_n$.

\emph{Step 7: size.} If $|\lambda|_\infty<\frac13$ and $|A|_\infty<\frac13$, then $|A(b)|+|\lambda(b_-)|+|\lambda(b_+)|<1$ at every bond, and $A^\lambda$ is defined. If $|\lambda|_\infty\le\frac1{10}$ and $|A|_\infty\le\frac1{10}$, then
\[
|A^\lambda(b)|\le\tfrac\pi2\,\bigl|e^{iA^\lambda(b)}-\mathbf 1\bigr|\le\tfrac\pi2\bigl(e^{3/10}-1\bigr)<1 .
\]
The second inequality is the product estimate in Definition~\ref{IV-73-1}(2). For the first one, let $Y\in\mathbb B_\pi$; then $(e^{iY}-\mathbf 1)^*(e^{iY}-\mathbf 1)=(2-2\cos|Y|)\mathbf 1$, so that $|e^{iY}-\mathbf 1|=2\sin\frac{|Y|}2\ge\frac2\pi|Y|$, by the inequality $\sin t\ge\frac2\pi t$ on $[0,\frac\pi2]$.

\emph{Step 8: $\Xi$ in the exponential charts.} The maps $\lambda\mapsto e^{i\lambda}$, $A\mapsto e^{iA}$, $A'\mapsto e^{iA'}$, taken site by site and bond by bond, are diffeomorphisms of the sets $\{|\lambda|_\infty<\pi\}\subset\mathsf X_n$, $\{|A|_\infty<\pi\}\subset\mathsf{Form}_n^{\mathsf T}$, $\{|A'|_\infty<\pi\}\subset\mathsf{Form}_n$ onto relatively open subsets of $\mathcal G_n$, $\mathcal U_n^{\mathsf T}$, $\mathcal U_n$ (Definition~\ref{IV-73-1}). On $\{|\lambda|_\infty<\frac13,\ |A|_\infty<\frac13\}$ we have $e^{iA^\lambda}=\mathsf j(e^{i\lambda},e^{iA})$ and $|A^\lambda|_\infty<\pi$; thus, in these charts, $\Xi$ is the map $\mathsf j$, and \textup{(b)} shows that $\Xi$ is a diffeomorphism of $\{|\lambda|_\infty<\frac13,\ |A|_\infty<\frac13\}$ onto an open subset of $\mathsf{Form}_n$, with $\Xi(0,0)=0$. We note that $\dim\mathsf X_n+\dim\mathsf{Form}_n^{\mathsf T}=\dim\mathsf{Form}_n$, since $\mathsf T_n$ has as many elements as $\Lambda_n^\circ$.

Let $\mathfrak D'\subset\{|\lambda|_\infty<\frac13,\ |A|_\infty<\frac13\}$ be open, and let $\Phi'$ be a bounded Borel function on $\Xi(\mathfrak D')$. We apply \eqref{IV-73-13} to the function $F$ on $\mathcal U_n$ which equals $\Phi'(A')$ at $U=e^{iA'}$, $A'\in\Xi(\mathfrak D')$, and $0$ elsewhere, and we express every Haar integral by the Haar formula in Definition~\ref{IV-73-1}(2) at $u_0=\mathbf 1$, bond by bond and site by site. Each of these integrals carries the measure $dY$ of $(\mathfrak g,\langle\cdot,\cdot\rangle)$ in the variable of its bond or site; since $(\cdot,\cdot)$ is the sum of $\langle\cdot,\cdot\rangle$ over the bonds, resp.\ the sites, the products of these measures are the Lebesgue measures $dA'$, $dA$, $d\lambda$ of the Euclidean spaces $\mathsf{Form}_n$, $\mathsf{Form}_n^{\mathsf T}$, $\mathsf X_n$, as in \textup{(iv)}.

On the left side this produces the factor $\sigma_{\mathrm H}^{\#\mathsf{Bd}_n}$ and an integral over $\Xi(\mathfrak D')$. On the right side, the injectivity of $\mathsf j$ and of the exponential charts shows that, for $(\lambda,A)$ in the domain of the charts, $(e^{iA})^{e^{i\lambda}}=e^{iA'}$ with $A'\in\Xi(\mathfrak D')$ holds if and only if $(\lambda,A)\in\mathfrak D'$, and then $A'=A^\lambda$; this produces the factor $\sigma_{\mathrm H}^{\#\Lambda_n^\circ+\#(\mathsf{Bd}_n\setminus\mathsf T_n)}=\sigma_{\mathrm H}^{\#\mathsf{Bd}_n}$ and an integral over $\mathfrak D'$. The powers of $\sigma_{\mathrm H}$ cancel, and
\begin{align*}
&\int_{\Xi(\mathfrak D')}\Phi'(A')\prod_{b\in\mathsf{Bd}_n}\sigma_1(|A'(b)|^2)\,dA'\\
&\qquad=\int_{\mathfrak D'}\Phi'(A^\lambda)\prod_{x\in\Lambda_n^\circ}\sigma_1(|\lambda(x)|^2)\prod_{b\in\mathsf{Bd}_n\setminus\mathsf T_n}\sigma_1(|A(b)|^2)\,d\lambda\,dA .
\end{align*}

\emph{Step 9: the derivative of the gauge function in $\lambda$.} Let $\mathsf c(\lambda,A):=\mathrm d^*(A^\lambda)$ for $\lambda\in\mathsf X_n$ and $A\in\mathsf{Form}_n^{\mathsf T}$ with $|\lambda|_\infty<\frac1{10}$, $|A|_\infty<\frac1{10}$. For $H\in\mathsf X_n$ and $|s|$ small let $\mu(s)\in\mathsf X_n$ be defined, site by site and with small values, by $e^{i\mu(s)}:=e^{i(\lambda+sH)}e^{-i\lambda}$. Then $\mu(0)=0$, and $\mu$ is smooth in $s$. Differentiating $e^{i\mu(s)}=e^{i(\lambda+sH)}e^{-i\lambda}$ at $s=0$ with the differential formula in Definition~\ref{IV-73-1}(2) on both sides, and using $\mathsf{dE}(0)=\mathbf 1$, we find at every site
\[
i\mu'(0)=i\bigl(\mathsf{dE}(\lambda)H\bigr)e^{i\lambda}e^{-i\lambda},\qquad\text{that is,}\qquad\mu'(0)=\vartheta_\lambda H,\quad\vartheta_\lambda:=\bigoplus_x\mathsf{dE}(\lambda(x)).
\]
Since $e^{i(\lambda+sH)}=e^{i\mu(s)}e^{i\lambda}$, we have $(e^{iA})^{e^{i(\lambda+sH)}}=\bigl((e^{iA})^{e^{i\lambda}}\bigr)^{e^{i\mu(s)}}$, that is, $A^{\lambda+sH}=(A^\lambda)^{\mu(s)}$ for $|s|$ small, both sides lying in $\mathbb B_\pi$ at every bond and having the same exponential. By Step~7, $|A^\lambda|_\infty<1$, and the definition \eqref{IV-73-8} of $\mathsf M$ with the chain rule gives
\[
\partial_\lambda\mathsf c(\lambda,A)H=\frac d{ds}\,\mathrm d^*\bigl((A^\lambda)^{\mu(s)}\bigr)\Big|_{s=0}=\mathsf M(A^\lambda)\,\vartheta_\lambda H .
\]
By the determinant formula in Definition~\ref{IV-73-1}(2), $\det\vartheta_\lambda=\prod_x\sigma_1(|\lambda(x)|^2)>0$. At $(\lambda,A)=(0,0)$ we have $\vartheta_0=\mathbf 1$ and $\partial_\lambda\mathsf c(0,0)=\mathsf M(0)=\mathsf M^{(0)}=-\mathrm d^*\mathrm d$, which is invertible by \textup{(c)}.

\emph{Step 10: choice of the radii; claims \textup{(i)}--\textup{(iii)}.} The map $\Theta_*:(\lambda,A)\mapsto(\mathsf c(\lambda,A),A)$, defined on the set of the pairs $(\lambda,A)\in\mathsf X_n\times\mathsf{Form}_n^{\mathsf T}$ with $|\lambda|_\infty<\frac1{10}$, $|A|_\infty<\frac1{10}$ and taking values in $\mathsf X_n\times\mathsf{Form}_n^{\mathsf T}$, is smooth, and its differential at $(0,0)$ is block triangular with the invertible diagonal blocks $-\mathrm d^*\mathrm d$ and $\mathbf 1$. By the inverse function theorem there is an open neighbourhood $\mathcal N$ of $(0,0)$ on which $\Theta_*$ is a diffeomorphism onto an open set; on $\mathcal N$ the differential of $\Theta_*$ is invertible, and therefore so is $\partial_\lambda\mathsf c$. We first choose $\varrho_1\le\frac1{20}$ with $\{|\lambda|_\infty\le2\varrho_1,\ |A|_\infty\le2\varrho_1\}\subset\mathcal N$, and then $\varrho_2,\varrho_3\le\varrho_1$ with
\[
\{c\in\mathsf X_n:\ |c|_\infty<\varrho_2\}\times\{A\in\mathsf{Form}_n^{\mathsf T}:\ |A|_\infty\le\varrho_3\}\subset\Theta_*\bigl(\{|\lambda|_\infty<\varrho_1,\ |A|_\infty<2\varrho_1\}\bigr);
\]
this is possible because the right side is an open set containing $\Theta_*(0,0)=(0,0)$. These numbers depend on $n$ only.

For \textup{(iii)}, let $A\in\mathsf{Form}_n^{\mathsf T}$ with $|A|_\infty\le\varrho_3$. Since $\Theta_*$ is injective on $\mathcal N$ and keeps the second component, $\lambda\mapsto\mathsf c(\lambda,A)$ is injective on $\{|\lambda|_\infty\le\varrho_1\}$. If $|c|_\infty<\varrho_2$, then $(c,A)=\Theta_*(\lambda,A'')$ with $|\lambda|_\infty<\varrho_1$, and necessarily $A''=A$; thus $c=\mathrm d^*(A^\lambda)$ with $|\lambda|_\infty<\varrho_1$. For \textup{(i)}: since $\varrho_1,\varrho_3\le\frac1{20}$, the closure $\overline{\mathfrak D}=\{|\lambda|_\infty\le\varrho_1,\ |A|_\infty\le\varrho_3\}$ lies in the open set $\{|\lambda|_\infty<\frac1{10},\ |A|_\infty<\frac1{10}\}$, on which $\Xi$ is a diffeomorphism onto an open set (Step~8) with values in $\{|A'|_\infty<1\}$ (Step~7).

For \textup{(ii)}: $\overline{\mathfrak D}$ is compact, so $\Xi(\overline{\mathfrak D})$ is compact; it contains $\mathcal W$ and is contained in $\overline{\mathcal W}$, hence equals $\overline{\mathcal W}$. For $(\lambda,A)\in\overline{\mathfrak D}\subset\mathcal N$ the operator $\partial_\lambda\mathsf c(\lambda,A)=\mathsf M(A^\lambda)\vartheta_\lambda$ is invertible and $\det\vartheta_\lambda\ne0$; hence $\det\mathsf M(A^\lambda)\ne0$.

\emph{Step 11: claim \textup{(iv)}.} By \eqref{IV-73-11}, the entries of $\mathsf M(A')$ are continuous functions of $A'$ for $|A'|_\infty<1$; hence $|\det\mathsf M|$ is bounded on the compact set $\overline{\mathcal W}$. We apply the identity of Step~8 with $\mathfrak D':=\mathfrak D$ and the bounded Borel function $\Phi':=\Phi\cdot(w\circ\mathrm d^*)\cdot|\det\mathsf M|$ on $\mathcal W$. Its left side is the left side of \eqref{IV-73-14}. On its right side we integrate first over $\lambda$, at fixed $A$ with $|A|_\infty<\varrho_3$ (Fubini's theorem). By \textup{(iii)} and Step~9, $\lambda\mapsto c=\mathsf c(\lambda,A)$ is an injective smooth map on $\{|\lambda|_\infty<\varrho_1\}$ with the invertible differential $\mathsf M(A^\lambda)\vartheta_\lambda$; hence it is a diffeomorphism onto its open image $\{\mathrm d^*(A^\lambda):|\lambda|_\infty<\varrho_1\}$, and
\[
dc=|\det\mathsf M(A^\lambda)|\prod_{x\in\Lambda_n^\circ}\sigma_1(|\lambda(x)|^2)\,d\lambda .
\]
Since $\Phi'(A^\lambda)=\Phi(A^\lambda)\,w(\mathsf c(\lambda,A))\,|\det\mathsf M(A^\lambda)|$, the substitution $c=\mathsf c(\lambda,A)$ turns the inner integral into $\int w(c)\,\Phi(A^{\lambda(c)})\,dc$ over this image, and \eqref{IV-73-14} follows.
\end{proof}

\subsection{The coefficient}

\begin{definition}[Quadratic form, covariances, vertices]\label{IV-73-15}
Let $\Omega$ be either $\mathbb Z^3$, indicated by the upper index $\mathrm{lat}$, or a box $\Lambda_n$, $n\ge2$, indicated by the upper index $[n]$, and let $\tau_{\mathrm{gf}}>0$; the index is omitted when $\Omega$ is fixed.

\textup{(1)} We set
\begin{equation}\label{IV-73-16}
\begin{gathered}
\mathsf S_2(A):=\tfrac12\|\mathrm dA\|^2+\tfrac1{2\tau_{\mathrm{gf}}}\|\mathrm d^*A\|^2,\\
\mathsf C_{\tau_{\mathrm{gf}}}=\bigl(\mathrm d^*\mathrm d+\tau_{\mathrm{gf}}^{-1}\mathrm d\mathrm d^*\bigr)^{-1},\qquad\mathsf G=(\mathrm d^*\mathrm d)^{-1}\ \text{on site functions.}
\end{gathered}
\end{equation}
In the box, $\mathrm d^*\mathrm d+\tau_{\mathrm{gf}}^{-1}\mathrm d\mathrm d^*$ is positive definite on $\mathsf{Form}_n^{\mathbb R}$ and $\mathrm d^*\mathrm d$ is positive definite on $\mathsf X_n^{\mathbb R}$, by Lemma~\ref{IV-73-12}\textup{(c)} applied to forms and functions with a single non-vanishing component. $\mathsf C^{[n]}_{\tau_{\mathrm{gf}}}$ and $\mathsf G^{[n]}$ are the inverses of these operators, with the matrices $\mathsf C^{[n]}_{\tau_{\mathrm{gf}}}(b,b')$, $b,b'\in\mathsf{Bd}_n$, and $\mathsf G^{[n]}(x,y)$, $x,y\in\Lambda_n^\circ$; we set $\mathsf C^{[n]}_{\tau_{\mathrm{gf}}}(b,b'):=0$ if $b$ or $b'$ is not a free bond, and $\mathsf G^{[n]}(x,y):=0$ if $x$ or $y$ does not lie in $\Lambda_n^\circ$. On $\mathbb Z^3$ we define, with $\mathsf k=\mathsf k(k)$,
\[
\begin{gathered}
\mathsf C^{\mathrm{lat}}_{\tau_{\mathrm{gf}}}\bigl((x,\mu),(y,\nu)\bigr):=\int_{\mathbb T}e^{ik\cdot(x-y)}\Bigl[\frac{\delta_{\mu\nu}}{|\mathsf k|^2}+(\tau_{\mathrm{gf}}-1)\frac{\mathsf k_\mu\overline{\mathsf k_\nu}}{|\mathsf k|^4}\Bigr]d^\circ k,\\
\mathsf G^{\mathrm{lat}}(x,y):=\int_{\mathbb T}\frac{e^{ik\cdot(x-y)}}{|\mathsf k|^2}\,d^\circ k .
\end{gathered}
\]
The matrix in brackets is Hermitian, with the eigenvalues $|\mathsf k|^{-2}$ (twice) and $\tau_{\mathrm{gf}}|\mathsf k|^{-2}$; as $|\mathsf k|^2\ge\frac4{\pi^2}|k|^2$ on $[-\pi,\pi]^3$, the integrands are bounded by $\frac{\pi^2}4\max\{1,\tau_{\mathrm{gf}}\}|k|^{-2}$, which is integrable. The bracket is positive, and it is replaced by its complex conjugate under $k\to-k$, because $\mathsf k_\mu(-k)=\overline{\mathsf k_\mu(k)}$. Hence the kernels are real and symmetric. They are also positive semidefinite: for finitely supported complex functions $f$ on the positively oriented bonds, resp.\ on the sites,
\[
\begin{gathered}
\sum_{b,b'}\overline{f(b)}\,\mathsf C^{\mathrm{lat}}_{\tau_{\mathrm{gf}}}(b,b')f(b')=\int_{\mathbb T}v(k)^*[\cdots]\,v(k)\,d^\circ k\ge0,\\
\sum_{x,y}\overline{f(x)}\,\mathsf G^{\mathrm{lat}}(x,y)f(y)=\int_{\mathbb T}\Bigl|\sum_ye^{-ik\cdot y}f(y)\Bigr|^2|\mathsf k|^{-2}\,d^\circ k\ge0,
\end{gathered}
\]
where $v_\nu(k):=\sum_ye^{-ik\cdot y}f((y,\nu))$ and $[\cdots]$ is the bracket.

The definition on $\mathbb Z^3$ is explained by plane waves. Let $A_\mu(x)=e^{ik\cdot x}\alpha_\mu$ with $\alpha\in\mathbb C^3$, let $\lambda(x)=e^{ik\cdot x}c$ with $c\in\mathbb C$, and let $F$ be the plaquette function with $F_{\rho\nu}(x)=e^{ik\cdot x}f_{\rho\nu}$ for all $\rho,\nu$, where $f_{\nu\rho}=-f_{\rho\nu}$; we write $\mathsf k^*\alpha:=\sum_\mu\overline{\mathsf k_\mu}\,\alpha_\mu$. The formulas for $\mathrm d$ and $\mathrm d^*$ in Definition~\ref{IV-73-2}\textup{(2)}, together with $e^{ik_\mu}-1=\mathsf k_\mu$ and $e^{-ik_\mu}-1=\overline{\mathsf k_\mu}$, give
\[
\begin{gathered}
\mathrm dA(p_{\mu\nu}(x))=e^{ik\cdot x}\bigl(\mathsf k_\mu\alpha_\nu-\mathsf k_\nu\alpha_\mu\bigr),\qquad\mathrm d^*A(x)=e^{ik\cdot x}\,\mathsf k^*\alpha,\\
(\mathrm d\lambda)((x,\rho))=e^{ik\cdot x}\,\mathsf k_\rho\,c,\qquad(\mathrm d^*F)((x,\rho))=-e^{ik\cdot x}\sum_\nu\overline{\mathsf k_\nu}\,f_{\rho\nu}.
\end{gathered}
\]
For $F=\mathrm dA$ we have $f_{\rho\nu}=\mathsf k_\rho\alpha_\nu-\mathsf k_\nu\alpha_\rho$, and for $\lambda=\mathrm d^*A$ we have $c=\mathsf k^*\alpha$; therefore
\[
(\mathrm d^*\mathrm dA)((x,\rho))=e^{ik\cdot x}\bigl(|\mathsf k|^2\alpha_\rho-\mathsf k_\rho\,\mathsf k^*\alpha\bigr),\qquad(\mathrm d\mathrm d^*A)((x,\rho))=e^{ik\cdot x}\,\mathsf k_\rho\,\mathsf k^*\alpha .
\]
Thus $\mathrm d^*\mathrm d+\tau_{\mathrm{gf}}^{-1}\mathrm d\mathrm d^*$ has the symbol $|\mathsf k|^2\mathbf 1-(1-\tau_{\mathrm{gf}}^{-1})\mathsf k\mathsf k^*$: it maps the plane wave with the amplitude $\alpha$ to the plane wave with the amplitude $\bigl(|\mathsf k|^2\mathbf 1-(1-\tau_{\mathrm{gf}}^{-1})\mathsf k\mathsf k^*\bigr)\alpha$. The inverse of this matrix is the bracket: as $\mathsf k^*\mathsf k=|\mathsf k|^2$,
\[
\begin{aligned}
&\bigl(|\mathsf k|^2\mathbf 1-(1-\tau_{\mathrm{gf}}^{-1})\mathsf k\mathsf k^*\bigr)\Bigl[\frac{\mathbf 1}{|\mathsf k|^2}+(\tau_{\mathrm{gf}}-1)\frac{\mathsf k\mathsf k^*}{|\mathsf k|^4}\Bigr]\\
&\qquad=\mathbf 1+\bigl[(\tau_{\mathrm{gf}}-1)-(1-\tau_{\mathrm{gf}}^{-1})-(1-\tau_{\mathrm{gf}}^{-1})(\tau_{\mathrm{gf}}-1)\bigr]\frac{\mathsf k\mathsf k^*}{|\mathsf k|^2}=\mathbf 1 ,
\end{aligned}
\]
since $(1-\tau_{\mathrm{gf}}^{-1})\bigl(1+(\tau_{\mathrm{gf}}-1)\bigr)=\tau_{\mathrm{gf}}-1$.

\textup{(2)}

\emph{Derived kernels.} With $b_i(p)$ and $s_i$ as in Lemma~\ref{IV-73-9} we set
\[
\begin{gathered}
\mathsf E_p(b):=\sum_is_i\,\mathsf C_{\tau_{\mathrm{gf}}}(b_i(p),b),\qquad\Pi(p,q):=\sum_js_j\,\mathsf E_p(b_j(q)),\\
\mathsf g(b):=\mathsf G(b_+,b_+)+\mathsf G(b_-,b_-)-2\mathsf G(b_+,b_-),\qquad\mathsf D(b;b'):=\tfrac12\sum_{x\in\{b'_-,b'_+\}}\bigl[\mathsf G(x,b_+)-\mathsf G(x,b_-)\bigr].
\end{gathered}
\]
The kernels $\mathsf E_p$ and $\Pi$ do not depend on $\tau_{\mathrm{gf}}$. On $\mathbb Z^3$, for $p=p_{\mu\nu}(x)$ and a bond $(y,\rho)$, the four terms of $\mathsf E_p$ combine to
\[
\sum_is_i\,\mathsf C^{\mathrm{lat}}_{\tau_{\mathrm{gf}}}\bigl(b_i(p),(y,\rho)\bigr)=\int_{\mathbb T}e^{ik\cdot(x-y)}\bigl(\mathsf k_\mu\,[\,\cdots]_{\nu\rho}-\mathsf k_\nu\,[\,\cdots]_{\mu\rho}\bigr)\,d^\circ k ,
\]
where $[\,\cdots]_{\alpha\rho}$ denotes the entry $(\alpha,\rho)$ of the bracket in \textup{(1)}; the $\tau_{\mathrm{gf}}$-dependent term of the bracket contributes $(\tau_{\mathrm{gf}}-1)(\mathsf k_\mu\mathsf k_\nu-\mathsf k_\nu\mathsf k_\mu)\overline{\mathsf k_\rho}|\mathsf k|^{-4}=0$. In the box,
\[
\mathsf C_{\tau_{\mathrm{gf}}}=\mathsf C_1+(\tau_{\mathrm{gf}}-1)\,\mathrm d\,\mathsf G^2\mathrm d^* ,\qquad\mathrm d\mathrm d=0 ;
\]
since $\mathsf E_p=\mathsf C_{\tau_{\mathrm{gf}}}\mathrm d^*\delta_p$ by the symmetry of $\mathsf C_{\tau_{\mathrm{gf}}}$, where $\delta_p$ is the plaquette function with the value $1$ at $p$ and $0$ elsewhere, and $\mathrm d^*\mathrm d^*=(\mathrm d\mathrm d)^*=0$, the second term does not contribute to $\mathsf E_p$ or to $\Pi$. To prove the identity for $\mathsf C_{\tau_{\mathrm{gf}}}$, let $\tau:=\tau_{\mathrm{gf}}$ and $\mathsf L_\tau:=\mathrm d^*\mathrm d+\tau^{-1}\mathrm d\mathrm d^*$. From $\mathrm d\mathrm d=0$ we get $\mathsf L_\tau\mathrm d=\tau^{-1}\mathrm d(\mathrm d^*\mathrm d)$, hence $\mathsf L_\tau\,\mathrm d\mathsf G^2\mathrm d^*=\tau^{-1}\mathrm d\mathsf G\mathrm d^*$.

Moreover $\mathrm d^*\mathsf L_1=\mathrm d^*\mathrm d\,\mathrm d^*$ gives $\mathrm d^*\mathsf C_1=\mathsf G\mathrm d^*$, and therefore $\mathsf L_\tau\mathsf C_1=\mathbf 1+(\tau^{-1}-1)\mathrm d\mathrm d^*\mathsf C_1=\mathbf 1+(\tau^{-1}-1)\mathrm d\mathsf G\mathrm d^*$. Adding the first identity multiplied by $\tau-1$ to the second we obtain $\mathsf L_\tau\bigl[\mathsf C_1+(\tau-1)\mathrm d\mathsf G^2\mathrm d^*\bigr]=\mathbf 1$.

On $\mathbb Z^3$, for $b=(x,\mu)$, the reality and symmetry of $\mathsf G^{\mathrm{lat}}$ give
\[
\mathsf g(b)=\int_{\mathbb T}\frac{2-e^{ik_\mu}-e^{-ik_\mu}}{|\mathsf k|^2}\,d^\circ k=\int_{\mathbb T}\frac{|\mathsf k_\mu|^2}{|\mathsf k|^2}\,d^\circ k=\frac13 ,
\]
since $|\mathsf k|^2$ is invariant under permutations of the coordinates of $k$, so that the three integrals $\int|\mathsf k_\mu|^2|\mathsf k|^{-2}d^\circ k$ are equal, with sum $1$.

\textup{(3)}

\emph{Vertices.} With sums over the plaquettes and the positively oriented bonds of $\Omega$, that is, over $\mathsf{Pl}_n$ and $\mathsf{Bd}_n$ in the box, we set
\begin{equation}\label{IV-73-17}
\begin{aligned}
\mathsf S_3&:=\sum_p\mathsf a_p^{(3)},\\
\mathsf V_2&:=\sum_p\mathsf a_p^{(4)}+\frac13\sum_b|A(b)|^2+\frac23\sum_b\mathsf g(b)|A(b)|^2-4\sum_{b,b'}\mathsf D(b;b')\mathsf D(b';b)\langle A(b),A(b')\rangle .
\end{aligned}
\end{equation}
The four terms of $\mathsf V_2$ are called the \emph{quartic Wilson vertex} $\mathsf S_4:=\sum_p\mathsf a_p^{(4)}$, the \emph{Haar term}, the \emph{ghost tadpole} and the \emph{ghost loop}. On $\mathbb Z^3$, $\mathsf S_3$ and $\mathsf V_2$ are formal sums of local terms in the sense of Definition~\ref{IV-73-19}\textup{(3)}. In the box,
\begin{equation}\label{IV-73-18}
\begin{gathered}
\tfrac23\sum_b\mathsf g(b)|A(b)|^2=-\operatorname{tr}\bigl[(\mathsf M^{(0)})^{-1}\mathsf M^{(2)}(A)\bigr],\\
-4\sum_{b,b'}\mathsf D(b;b')\mathsf D(b';b)\langle A(b),A(b')\rangle=\tfrac12\operatorname{tr}\bigl[\bigl((\mathsf M^{(0)})^{-1}\mathsf M^{(1)}(A)\bigr)^2\bigr],
\end{gathered}
\end{equation}
the traces being taken on $\mathsf X_n$, and $\operatorname{tr}\bigl[(\mathsf M^{(0)})^{-1}\mathsf M^{(1)}(A)\bigr]=0$.

To prove these identities we use $(\mathsf M^{(0)})^{-1}=-\mathsf G$ (Lemma~\ref{IV-73-9}\textup{(d)}). For a family $T=(T_b)$ of endomorphisms $T_b\in\operatorname{End}\mathfrak g$, acting on forms bond by bond, the cyclicity of the trace gives
\[
\operatorname{tr}\bigl[\mathsf G\,\mathrm d^*T\mathrm d\bigr]=\sum_b\operatorname{tr}T_b\cdot(\mathrm d\mathsf G\mathrm d^*)(b,b)=\sum_b\mathsf g(b)\operatorname{tr}T_b ,
\]
because the diagonal element of the kernel of $\mathrm d\mathsf G\mathrm d^*$ at $b$ is $\mathsf G(b_+,b_+)-\mathsf G(b_+,b_-)-\mathsf G(b_-,b_+)+\mathsf G(b_-,b_-)=\mathsf g(b)$. The formula $X\times(Y\times Z)=\langle X,Z\rangle Y-\langle X,Y\rangle Z$ gives $\operatorname{tr}[X\times(Y\times\cdot)]=\langle X,Y\rangle-3\langle X,Y\rangle=-2\langle X,Y\rangle$. By Lemma~\ref{IV-73-9}\textup{(d)}, $\mathsf M^{(2)}(A)=\mathrm d^*T\mathrm d$ with $T_b=-\frac13A(b)\times(A(b)\times\cdot)$, so $\operatorname{tr}T_b=\frac23|A(b)|^2$, and the first identity follows. Next let $\mathsf R:=\mathsf G\,\mathrm d^*(A\times\bar{\,\cdot\,})$, so that $(\mathsf M^{(0)})^{-1}\mathsf M^{(1)}(A)=-2\mathsf R$. Its kernel is
\[
\mathsf R(x,y)=\sum_b\bigl[\mathsf G(x,b_+)-\mathsf G(x,b_-)\bigr]\,\tfrac12\bigl[\delta_{y,b_-}+\delta_{y,b_+}\bigr]\,A(b)\times\cdot\ \in\operatorname{End}\mathfrak g .
\]
Since $\operatorname{tr}[X\times\cdot]=0$, we get $\operatorname{tr}\mathsf R=0$, which is the last assertion. Summing first over $x$ and $y$ in $\operatorname{tr}\mathsf R^2=\sum_{x,y}\operatorname{tr}[\mathsf R(x,y)\mathsf R(y,x)]$ we find
\[
\operatorname{tr}\mathsf R^2=\sum_{b,b'}\mathsf D(b;b')\mathsf D(b';b)\operatorname{tr}\bigl[A(b)\times(A(b')\times\cdot)\bigr]=-2\sum_{b,b'}\mathsf D(b;b')\mathsf D(b';b)\langle A(b),A(b')\rangle ,
\]
and $\frac12\operatorname{tr}[(-2\mathsf R)^2]=2\operatorname{tr}\mathsf R^2$ gives the second identity.
\end{definition}

\begin{definition}[Gaussian family; cumulants of polynomials; formal sums]\label{IV-73-19}
Let $\Omega$ and $\tau_{\mathrm{gf}}$ be as in Definition~\ref{IV-73-15}.

\textup{(1)} The variables $A^a(b)$, where $b$ ranges over the positively oriented bonds of $\Omega$ (over the free bonds if $\Omega=\Lambda_n$) and $a\in\{1,2,3\}$, form a centred Gaussian family with
\[
\langle A^a(b)A^{a'}(b')\rangle_0=\delta_{aa'}\mathsf C^\Omega_{\tau_{\mathrm{gf}}}(b,b');
\]
such a family exists because this covariance is positive semidefinite (Definition~\ref{IV-73-15}\textup{(1)}), and its mean is denoted by $\langle\cdot\rangle_0$. In the box we put $A^a(b):=0$ at the bonds which are not free, in accordance with the definition of $\mathsf{Form}_n$. A \emph{leg} is a variable $A^a(b)$ (\emph{plain leg}) or a variable $(\mathrm dA)^a(p)=\sum_is_iA^a(b_i(p))$ (\emph{curl leg}). By the definitions of $\mathsf E_p$ and $\Pi$,
\begin{equation}\label{IV-73-20}
\begin{gathered}
\langle A^a(b)A^{a'}(b')\rangle_0=\delta_{aa'}\mathsf C(b,b'),\qquad\langle(\mathrm dA)^a(p)A^{a'}(b)\rangle_0=\delta_{aa'}\mathsf E_p(b),\\
\langle(\mathrm dA)^a(p)(\mathrm dA)^{a'}(q)\rangle_0=\delta_{aa'}\Pi(p,q).
\end{gathered}
\end{equation}

\textup{(2)} Let $X_1,\dots,X_r$ be polynomials in the variables of the family, and $m_J:=\langle\prod_{i\in J}X_i\rangle_0$ for $\emptyset\ne J\subset\{1,\dots,r\}$. Their \emph{joint cumulants} $\langle X_i;i\in J\rangle_0:=\kappa_J$ are defined by the partition formula
\begin{equation}\label{eq:two-loop-cumulant-partition}
\kappa_J=\sum_{\pi\in\operatorname{Part}(J)}(-1)^{|\pi|-1}(|\pi|-1)!\prod_{I\in\pi}m_I .
\end{equation}
Here $\operatorname{Part}(J)$ denotes the partitions $\pi$ of $J$ into nonempty blocks $I$, and $|\pi|$ is the number of blocks. Equivalently, $(\kappa_J)$ is the unique family satisfying
\begin{equation}\label{eq:two-loop-moment-partition}
 m_J=\sum_{\pi\in\operatorname{Part}(J)}\prod_{I\in\pi}\kappa_I.
\end{equation}
The equivalence is proved below. The cumulants are symmetric and multilinear in the $X_i$: each index belongs to exactly one block of each partition, every moment is linear in each factor, and the formula involves $J$ only as a set. Formula~\eqref{eq:two-loop-cumulant-partition} requires only the indicated moments; no exponential moment of a polynomial in the Gaussian family is assumed.

When exponential moments exist, this definition agrees with differentiation of their logarithm. For bounded complex variables, the exponential mean is holomorphic near zero and can be differentiated under the integral. For real variables $Y_i=X_i$, $i\in J=\{1,\dots,n\}$, assume instead that $\nu(e^{t\sum_i|Y_i|})<\infty$ for some $t>0$, where $\nu(\cdot)=\langle\cdot\rangle_0$. Then $h(z):=\nu(e^{\sum_i z_iY_i})$ is holomorphic on $\{|\operatorname{Re}z_i|<t\}$ and $h(0)=1$. Write $\partial_I=\prod_{i\in I}\partial_{z_i}$. For $|\operatorname{Re}z_i|\le t/2$ and $I\subset J$,
\[
\prod_{i\in I}|Y_i|\;\bigl|e^{\sum_iz_iY_i}\bigr|\le\Bigl(\frac2t\Bigr)^{|I|}e^{t\sum_i|Y_i|},
\]
by $s\le e^s$ at $s=t|Y_i|/2$. This integrable majorant gives $\partial_Ih(0)=m_I$. The mixed derivative formula for the logarithm, proved below, is
\begin{equation}\label{eq:two-loop-log-partition}
 \partial_J\log h(0)
 =\sum_{\pi\in\operatorname{Part}(J)}
       (-1)^{|\pi|-1}(|\pi|-1)!
       \prod_{I\in\pi}\partial_Ih(0).
\end{equation}
It identifies this derivative with $\kappa_J$.

\textup{(3)} A polynomial in the Gaussian family can be written in several ways as a sum of monomials in legs, since a curl leg is a signed sum of four plain legs. The \emph{standard writing} of $\mathsf a_p^{(m)}$, $m=2,3,4$, is obtained from \eqref{IV-73-10} by multiplying out in components, writing each displayed factor $\mathrm dA(p)$ in curl legs $(\mathrm dA)^a(p)$ and the polynomials $\mathfrak z_1$, $\mathfrak z_2$ in plain legs $A^a(b_i(p))$. Thus $\mathsf a_p^{(2)}=\frac12\sum_a\bigl((\mathrm dA)^a(p)\bigr)^2$; every monomial of $\mathsf a_p^{(3)}$ consists of one curl leg and two plain legs; every monomial of $\mathsf a_p^{(4)}$ consists of four plain legs (those coming from $\frac12|\mathfrak z_1|^2$), or of one curl leg and three plain legs, or of four curl legs. The other local terms of $\mathsf V_2$ are written in plain legs.

\emph{Formal sums.} Let $X_i=\sum_{z\in\mathsf Z_i}X_{i,z}$, $i=1,\dots,r$, be formal sums of polynomials $X_{i,z}$ over countable index sets $\mathsf Z_i$. We set
\[
\langle X_1;\dots;X_r\rangle_0:=\sum_{z=(z_1,\dots,z_r)}\langle X_{1,z_1};\dots;X_{r,z_r}\rangle_0 ,
\]
provided the following absolute convergence holds: when every entry is in its standard writing and each cumulant of monomials is expanded as the sum over pairings of Lemma~\ref{IV-73-21}, the sum of the moduli of all the resulting terms is finite. Under this proviso $\sum_z|\langle X_{1,z_1};\dots;X_{r,z_r}\rangle_0|<\infty$, because at fixed $z$ the cumulant is a number, the same in every writing, and it is bounded by the sum of the moduli of its terms; hence the value of the sum does not depend on the writing or on the order of summation. The \emph{local terms} are the $\mathsf a_p^{(3)}$ for $\mathsf S_3$; for $\mathsf V_2$ they are the $\mathsf a_p^{(4)}$, the terms of the two sums over $b$, and the terms of the sum over pairs $(b,b')$.
\end{definition}

\begin{proof}[Proof of the cumulant identities in Definition~\ref{IV-73-19}]
Fix a finite nonempty index set $J_0$, and write
$z^I=\prod_{i\in I}z_i$ for $I\subset J_0$.
For a holomorphic function $h$ of finitely many variables and a
one-variable holomorphic function $\phi$ near $h(0)$, repeated
differentiation gives
\[
 \partial_J(\phi\circ h)
 =\sum_{\pi\in\operatorname{Part}(J)}
       \phi^{(|\pi|)}(h)\prod_{I\in\pi}\partial_Ih.
\]
Indeed, when a new variable is differentiated, it either
differentiates $\phi^{(|\pi|)}(h)$ and forms a new singleton
block, or differentiates one factor $\partial_Ih$ and joins
that block. Every partition of the enlarged set arises
exactly once. This proves the formula by induction from
the one-variable chain rule.
At $h(0)=1$, the derivatives
$(\log)^{(k)}(1)=(-1)^{k-1}(k-1)!$ give
\eqref{eq:two-loop-log-partition}. With $\phi=\exp$ and
$h(0)=0$, the same formula gives
\[
 \partial_Je^{h}(0)
 =\sum_{\pi\in\operatorname{Part}(J)}
       \prod_{I\in\pi}\partial_Ih(0).
\]

For arbitrary moments $m_I$, put
$H(z)=1+\sum_{\emptyset\ne I\subset J_0}m_Iz^I$ and
$\varphi=\log H$ near zero. The logarithm formula gives
$\partial_I\varphi(0)=\kappa_I$ as defined by
\eqref{eq:two-loop-cumulant-partition}. Since $e^\varphi=H$,
the exponential formula proves
\eqref{eq:two-loop-moment-partition}.
Conversely, suppose \eqref{eq:two-loop-moment-partition}
holds and put $\varphi(z)=\sum_{\emptyset\ne I\subset J_0}
\kappa_Iz^I$. The exponential formula gives
$\partial_Ie^\varphi(0)=m_I$, and applying the logarithm
formula to $e^\varphi$ recovers
\eqref{eq:two-loop-cumulant-partition}.
This also proves uniqueness.
\end{proof}

\begin{lemma}[Cumulants of monomials: connected pairings]\label{IV-73-21}
Let $(\phi_i)_{i\in I}$ be a centred Gaussian family with covariance $\mathsf C(i,j)$ and mean $\langle\cdot\rangle_0$, the cumulants of polynomials being those of Definition~\ref{IV-73-19}\textup{(2)}, and let $\mathsf m_v=\prod_{e\in E_v}\phi_{i(e)}$, $v=1,\dots,r$, where the sets $E_v$ of legs are finite and pairwise disjoint and $i(e)\in I$. Then
\[
\langle\mathsf m_1;\dots;\mathsf m_r\rangle_0=\sum_{\mathsf P}\ \prod_{\{e,e'\}\in\mathsf P}\mathsf C\bigl(i(e),i(e')\bigr),
\]
where the sum ranges over the pairings $\mathsf P$ of $E_1\cup\dots\cup E_r$ whose graph on $\{1,\dots,r\}$ is connected; this graph has an edge $vv'$ for each pair $\{e,e'\}\in\mathsf P$ with $e\in E_v$, $e'\in E_{v'}$, $v\ne v'$.
\end{lemma}

\begin{proof}
\emph{Step 1: moments.} Let $E$ be a finite set of legs, with $i(e)\in I$ for $e\in E$, and let $s=(s_e)_{e\in E}$ be real variables. The variables $\phi_{i(e)}$, $e\in E$, form a centred Gaussian family with covariance $\mathsf C(i(e),i(e'))$, so that
\[
\Bigl\langle\exp\sum_{e\in E}s_e\phi_{i(e)}\Bigr\rangle_0=\exp\Bigl(\tfrac12\sum_{e,e'\in E}s_es_{e'}\,\mathsf C\bigl(i(e),i(e')\bigr)\Bigr).
\]
Both sides are real analytic in $s$, and the left side may be differentiated under the mean. We compare the coefficients of the monomial $\prod_{e\in E}s_e$. On the left it is $\langle\prod_{e\in E}\phi_{i(e)}\rangle_0$. On the right, this monomial occurs only in the term $\frac1{m!}\bigl(\frac12\sum_{e,e'}s_es_{e'}\mathsf C(i(e),i(e'))\bigr)^m$ with $2m=\#E$; when the power is expanded, each pairing of $E$ arises from $m!\,2^m$ choices, namely the order of its $m$ pairs and the order within each pair. Hence
\[
\Bigl\langle\prod_{e\in E}\phi_{i(e)}\Bigr\rangle_0=\sum_{\mathsf P}\ \prod_{\{e,e'\}\in\mathsf P}\mathsf C\bigl(i(e),i(e')\bigr),
\]
the sum ranging over all pairings $\mathsf P$ of $E$; it vanishes if $\#E$ is odd.

\emph{Step 2: connected pairings.} For $\emptyset\ne J\subset\{1,\dots,r\}$ let $\tilde\kappa_J$ be the sum over the pairings of $\bigcup_{v\in J}E_v$ whose graph on $J$ is connected. A pairing of $\bigcup_{v\in J}E_v$ is the same as a partition $\pi$ of $J$, namely into the vertex sets of the connected components of its graph, together with a pairing of $\bigcup_{v\in I}E_v$ with connected graph on $I$ for each block $I\in\pi$; the product of the covariances factorises accordingly. Step~1, applied to $E=\bigcup_{v\in J}E_v$, therefore gives
\[
m_J:=\Bigl\langle\prod_{v\in J}\mathsf m_v\Bigr\rangle_0=\sum_{\pi\in\operatorname{Part}(J)}\ \prod_{I\in\pi}\tilde\kappa_I ,
\]
which is \eqref{eq:two-loop-moment-partition}. By Definition~\ref{IV-73-19}\textup{(2)}, $\tilde\kappa_J$ is the joint cumulant $\kappa_J$ of the $\mathsf m_v$, $v\in J$; for $J=\{1,\dots,r\}$ this is the claim.
\end{proof}

\begin{definition}[The lattice coefficients $\mathsf B_0$, $\mathsf B_1$, $\mathsf P_1^{\mathrm{lat}}$]\label{IV-73-22}
Let $\Omega$ and $\tau_{\mathrm{gf}}$ be as in Definition~\ref{IV-73-15}, let $p$, $q$ be plaquettes of $\Omega$, and let $\langle\cdot\rangle_0$ refer to the Gaussian family of $\Omega$ and $\tau_{\mathrm{gf}}$ (Definition~\ref{IV-73-19}). We set
\begin{equation}\label{IV-73-23}
\mathsf B_0^\Omega(p,q):=\langle\mathsf a_p^{(2)};\mathsf a_q^{(2)}\rangle_0=\tfrac{\mathsf n}2\,\Pi^\Omega(p,q)^2 .
\end{equation}
The second equality follows from Lemma~\ref{IV-73-21} and \eqref{IV-73-20}, applied to the legs, which form a centred Gaussian family as linear combinations of the variables $A^a(b)$. Indeed, $\mathsf a_p^{(2)}=\frac12\sum_a((\mathrm dA)^a(p))^2$ and multilinearity give $\mathsf B_0^\Omega(p,q)=\frac14\sum_{a,a'}\langle((\mathrm dA)^a(p))^2;((\mathrm dA)^{a'}(q))^2\rangle_0$; in each of these cumulants exactly two pairings are connected, namely those which pair each leg at $p$ with a leg at $q$, and each of them contributes $\delta_{aa'}\Pi^\Omega(p,q)^2$. Hence $\mathsf B_0^\Omega(p,q)=\frac14\cdot2\cdot3\,\Pi^\Omega(p,q)^2=\frac{\mathsf n}2\Pi^\Omega(p,q)^2$. Next we set
\begin{equation}\label{IV-73-24}
\begin{aligned}
\mathsf B_1^\Omega(p,q;\tau_{\mathrm{gf}}):={}&\langle\mathsf a_p^{(3)};\mathsf a_q^{(3)}\rangle_0+\langle\mathsf a_p^{(2)};\mathsf a_q^{(4)}\rangle_0+\langle\mathsf a_p^{(4)};\mathsf a_q^{(2)}\rangle_0\\
&-\langle\mathsf a_p^{(3)};\mathsf a_q^{(2)};\mathsf S_3\rangle_0-\langle\mathsf a_p^{(2)};\mathsf a_q^{(3)};\mathsf S_3\rangle_0\\
&-\langle\mathsf a_p^{(2)};\mathsf a_q^{(2)};\mathsf V_2\rangle_0+\tfrac12\langle\mathsf a_p^{(2)};\mathsf a_q^{(2)};\mathsf S_3;\mathsf S_3\rangle_0 .
\end{aligned}
\end{equation}
We write $\mathsf B_j^{[n]}$ for $\Omega=\Lambda_n$ and $\mathsf B_j^{\mathrm{lat}}$ for $\Omega=\mathbb Z^3$, $j=0,1$. On $\mathbb Z^3$ the four cumulants containing $\mathsf S_3$ or $\mathsf V_2$ are formal sums in the sense of Definition~\ref{IV-73-19}\textup{(3)}; the proviso of absolute convergence is satisfied for every $\tau_{\mathrm{gf}}>0$ by Lemma~\ref{IV-73-55}. Finally, for bounded Borel functions $f$, $h$ on $\mathbb R^3$ with compact supports and $a>0$,
\begin{equation}\label{IV-73-25}
\mathsf B_1^{\mathrm{lat}}(p,q):=\mathsf B_1^{\mathrm{lat}}(p,q;1),\qquad\mathsf P_1^{\mathrm{lat}}(f,h;a):=a\sum_{p,q}f(a\,x(p))\,h(a\,x(q))\,\mathsf B_1^{\mathrm{lat}}(p,q);
\end{equation}
the sum over pairs of plaquettes of $\mathbb Z^3$ is finite, because $f(a\,x(p))\ne0$ only for the finitely many plaquettes with $a\,x(p)$ in the compact support of $f$, and likewise for $h$.

The normalisation in \eqref{IV-73-25} follows from dilation. On the spacing-$a$ lattice define
\[
 \mathcal O_a(f)(V)=a^{-1}\sum_{P}f(x(P))a_P(V),
\]
with the sum over its plaquettes. The dilation sends a unit bond $b$ to $ab$ and defines $(\mathsf{dil}_aU)(ab)=U(b)$. It therefore preserves the ordered plaquette product and sends its barycentre $x(p)$ to $ax(p)$. Consequently
\[
\mathcal O_a(f)(\mathsf{dil}_aU)=a^{-1}\sum_pf(a\,x(p))\,a_p(U)
\]
for every configuration $U\in\mathcal V(\mathbb L_1)$, the sum running over the plaquettes of $\mathbb L_1=\mathbb Z^3$. If $\beta=\mathsf b/a$ with $\mathsf b>0$, then $a^{-2}\beta^{-3}=a\,\mathsf b^{-3}$, and \eqref{IV-73-25} gives
\[
a^{-2}\sum_{p,q}f(a\,x(p))\,h(a\,x(q))\,\beta^{-3}\,\mathsf B_1^{\mathrm{lat}}(p,q)=\mathsf b^{-3}\,\mathsf P_1^{\mathrm{lat}}(f,h;a).
\]
\end{definition}

\subsection{The asymptotic expansion in a box}

\begin{lemma}[Laplace's method for a truncated two-point function]\label{IV-73-26}
Let $m\ge1$ and $\rho>0$, let $\mathcal W\subset\mathbb R^m$ be a bounded Borel set containing $\{|y|<\rho\}$, and let $\Phi$, $\Theta$, $X$, $Y$ be real $C^\infty$ functions on an open set containing $\overline{\mathcal W}$ with the following properties: $\Phi(0)=0$, $\nabla\Phi(0)=0$, the Hessian $H$ of $\Phi$ at $0$ is positive definite, and $\inf\{\Phi(y):\ y\in\mathcal W,\ |y|\ge\rho'\}>0$ for every $\rho'\in(0,\rho]$; $X$ and $Y$ vanish at $0$ together with their gradients. For $\beta\ge1$ let $\langle\cdot\rangle_{\mathcal W,\beta}$ be the mean with respect to the probability density on $\mathcal W$, relative to the Lebesgue measure, proportional to $e^{-\beta\Phi-\Theta}$, and $\langle X;Y\rangle_{\mathcal W,\beta}:=\langle XY\rangle_{\mathcal W,\beta}-\langle X\rangle_{\mathcal W,\beta}\langle Y\rangle_{\mathcal W,\beta}$. Let $\Phi_j$, $\Theta_j$, $X'_j$, $Y'_j$ be the homogeneous Taylor polynomials of degree $j$ at $0$ of $\Phi$, $\Theta$, $X$, $Y$, and
\[
W_j:=\Phi_{j+2}+\Theta_j\quad(j\ge1),\qquad X_j:=X'_{j+2},\quad Y_j:=Y'_{j+2}\quad(j\ge0).
\]
Let $\langle\cdot\rangle_0$ be the centred Gaussian measure on $\mathbb R^m$ with covariance $H^{-1}$, the cumulants of polynomials being taken as in Definition~\ref{IV-73-19}\textup{(2)}. Then there is $C<\infty$, depending on the data $m,\rho,\mathcal W,\Phi,\Theta,X,Y$, such that for all $\beta\ge1$
\[
\bigl|\langle X;Y\rangle_{\mathcal W,\beta}-\beta^{-2}\gamma_0-\beta^{-3}\gamma_2\bigr|\le C\beta^{-4},
\]
where $\gamma_0:=\langle X_0;Y_0\rangle_0$ and
\[
\begin{aligned}
\gamma_2:={}&\langle X_1;Y_1\rangle_0+\langle X_0;Y_2\rangle_0+\langle X_2;Y_0\rangle_0-\langle X_1;Y_0;W_1\rangle_0-\langle X_0;Y_1;W_1\rangle_0\\
&-\langle X_0;Y_0;W_2\rangle_0+\tfrac12\langle X_0;Y_0;W_1;W_1\rangle_0 .
\end{aligned}
\]
\end{lemma}

\begin{proof}
\emph{Step 1: localisation and scaling.} By Taylor's formula at $0$ and the positivity of $H$, there is $\rho_0\in(0,\rho]$ with $\Phi(y)\ge\frac14(y,Hy)$ for $|y|\le\rho_0$. Let $t:=\beta^{-1/2}\in(0,1]$, and for bounded Borel functions $f$ on $\mathcal W$ let
\[
\mathsf N(f):=\int_{\mathcal W}f\,e^{-\beta\Phi-\Theta}\,dy .
\]
The contribution of $\{y\in\mathcal W:|y|\ge\rho_0\}$ to $\mathsf N(f)$ is bounded by $C\sup|f|\,e^{-c_*\beta}$, where $c_*:=\inf\{\Phi(y):\ y\in\mathcal W,\ |y|\ge\rho_0\}>0$. On $\{|y|<\rho_0\}\subset\mathcal W$ we substitute $y=tz$. Taylor's formula with integral remainder provides smooth functions $\Phi_\alpha$ ($|\alpha|=3$) and $X_\alpha$, $Y_\alpha$ ($|\alpha|=2$) on $\{|y|\le\rho_0\}$ such that
\[
\begin{gathered}
\Phi(y)=\tfrac12(y,Hy)+\sum_{|\alpha|=3}y^\alpha\Phi_\alpha(y),\\
X(y)=\sum_{|\alpha|=2}y^\alpha X_\alpha(y),\qquad Y(y)=\sum_{|\alpha|=2}y^\alpha Y_\alpha(y).
\end{gathered}
\]
We set
\begin{align*}
\varphi(t,z)&:=t^{-2}\Phi(tz)-\tfrac12(z,Hz)+\Theta(tz)-\Theta(0)=t\sum_\alpha z^\alpha\Phi_\alpha(tz)+\Theta(tz)-\Theta(0),\\
\mathsf x(t,z)&:=t^{-2}X(tz)=\sum_\alpha z^\alpha X_\alpha(tz),\qquad\mathsf y(t,z):=t^{-2}Y(tz)=\sum_\alpha z^\alpha Y_\alpha(tz).
\end{align*}
These functions are smooth on $\{(t,z):\ |tz|\le\rho_0\}$, the value $t=0$ included, and they satisfy $|\partial_t^j\varphi|,|\partial_t^j\mathsf x|,|\partial_t^j\mathsf y|\le C_j\langle z\rangle^{j+3}$ there. Since the Taylor coefficients at $t=0$ of $t\mapsto F(tz)$ are the homogeneous Taylor polynomials $F_j(z)$, their formal Taylor series in $t$ are $\sum_{j\ge1}t^jW_j(z)$, $\sum_{j\ge0}t^jX_j(z)$ and $\sum_{j\ge0}t^jY_j(z)$. Moreover
\[
e^{-\frac12(z,Hz)-\varphi(t,z)}=e^{-t^{-2}\Phi(tz)-\Theta(tz)+\Theta(0)}\le Ce^{-\frac14(z,Hz)}\qquad(|tz|\le\rho_0).
\]

\emph{Step 2: expansion of the integrals.} For $\mathsf f\in\{1,\mathsf x,\mathsf y,\mathsf x\mathsf y\}$ let $\mathsf F:=\mathsf fe^{-\varphi}$. By Step~1, $|\partial_t^j\mathsf F|\le C_j\langle z\rangle^{4j+6}e^{-\varphi}$ on $\{|tz|\le\rho_0\}$; if $|tz|\le\rho_0$, the points $(s,z)$ with $0\le s\le t$ belong to this set as well; and the coefficients $\mathsf F_j(z):=\frac1{j!}\partial_t^j\mathsf F(0,z)$ are polynomials in $z$, since $\varphi(0,z)=0$ and the $t$-derivatives of $\varphi$, $\mathsf x$, $\mathsf y$ at $t=0$ are polynomials. Taylor's formula in $t$ with a remainder of order $4$, together with the bound of Step~1 at the points $(s,z)$, gives
\[
\Bigl|\mathsf F(t,z)-\sum_{j<4}t^j\mathsf F_j(z)\Bigr|\,e^{-\frac12(z,Hz)}\le Ct^4\langle z\rangle^{22}e^{-\frac14(z,Hz)}\qquad(|tz|\le\rho_0).
\]
Let $\mathsf Z:=t^me^{-\Theta(0)}\int e^{-\frac12(z,Hz)}\,dz$, and let $\nu\in\{0,1,1,2\}$ be the number of factors $X$, $Y$ in $f\in\{1,X,Y,XY\}$, with the corresponding $\mathsf f\in\{1,\mathsf x,\mathsf y,\mathsf x\mathsf y\}$. We substitute $y=tz$ in $\mathsf N(f)$, use $X(tz)=t^2\mathsf x(t,z)$, $Y(tz)=t^2\mathsf y(t,z)$ and $e^{-\beta\Phi(tz)-\Theta(tz)}=e^{-\Theta(0)}e^{-\frac12(z,Hz)-\varphi(t,z)}$, integrate the last bound over $|z|<\rho_0/t$, and add the Gaussian tails $\int_{|z|\ge\rho_0/t}|\mathsf F_j|e^{-\frac12(z,Hz)}\,dz=O(e^{-c/t^2})$. Since the Gaussian measure with covariance $H^{-1}$ has a density proportional to $e^{-\frac12(z,Hz)}$, we obtain
\[
\mathsf N(f)=\mathsf Z\,t^{2\nu}\Bigl[\sum_{j<4}t^j\langle\mathsf F_j\rangle_0+O(t^4)\Bigr]\qquad(f=1,X,Y,XY),
\]
where the contribution of order $e^{-c_*/t^2}$ of $\{|y|\ge\rho_0\}$ is also absorbed in $O(t^4)$, because $\mathsf Z\,t^{2\nu+4}\ge c\,t^{m+8}$. For $f=1$ we have $\mathsf F_0=1$; hence there is $t_0\in(0,1]$ with $\mathsf N(1)\ge\frac12\mathsf Z$ for $t\le t_0$.

In the range $t_0<t\le1$, that is $1\le\beta<t_0^{-2}$, the claim holds with $C=t_0^{-8}\bigl(2\sup|X|\sup|Y|+|\gamma_0|+|\gamma_2|\bigr)$: the functions $X$ and $Y$ are continuous, hence bounded on the compact set $\overline{\mathcal W}$ (Heine--Borel), so that $|\langle X;Y\rangle_{\mathcal W,\beta}|\le2\sup|X|\sup|Y|$, and $\beta^{-4}>t_0^8$.

In the range $t\le t_0$ we have
\[
\langle X;Y\rangle_{\mathcal W,\beta}=\frac{\mathsf N(XY)}{\mathsf N(1)}-\frac{\mathsf N(X)\,\mathsf N(Y)}{\mathsf N(1)^2}=t^4\Bigl[\sum_{j<4}t^j\gamma_j+O(t^4)\Bigr],
\]
where $\sum_jt^j\gamma_j$ is the formal power series
\[
\mathsf e:=\frac{\langle\mathsf x\mathsf ye^{-\varphi}\rangle_0}{\langle e^{-\varphi}\rangle_0}-\frac{\langle\mathsf xe^{-\varphi}\rangle_0\,\langle\mathsf ye^{-\varphi}\rangle_0}{\langle e^{-\varphi}\rangle_0^2}\in\mathbb R[[t]],
\]
in which $\mathsf x=\sum_jt^jX_j$, $\mathsf y=\sum_jt^jY_j$ and $\varphi=\sum_{j\ge1}t^jW_j$ are formal series and $\langle\cdot\rangle_0$ acts coefficientwise. Indeed, the formal series $\sum_jt^j\langle\mathsf F_j\rangle_0$ are the series $\langle\mathsf fe^{-\varphi}\rangle_0$; the series $\langle e^{-\varphi}\rangle_0$ has the constant term $1$, so that the quotients are defined in $\mathbb R[[t]]$; and, the denominators being at least $\frac12$, the quotients of the truncated expansions differ from the truncations of the formal quotients by $O(t^4)$. Since $t^4=\beta^{-2}$, $t^6=\beta^{-3}$ and $t^8=\beta^{-4}$, it remains to identify $\gamma_0$ and $\gamma_2$ and to show $\gamma_1=\gamma_3=0$.

\emph{Step 3: the coefficients.} Let $\alpha_1$, $\alpha_2$ be formal variables, and let $Z:=\alpha_1\mathsf x+\alpha_2\mathsf y-\varphi$, a formal series in $(\alpha_1,\alpha_2,t)$ without constant term whose coefficients are polynomials in $z$. By \eqref{eq:two-loop-moment-partition} and multilinearity,
\[
\langle Z^r\rangle_0=\sum_{\pi\in\operatorname{Part}(\{1,\dots,r\})}\ \prod_{I\in\pi}\varkappa_{\#I},\qquad\varkappa_j:=\langle Z;\dots;Z\rangle_0\quad(j\text{ entries}).
\]
The number of partitions of $\{1,\dots,r\}$ with $m_j$ blocks of size $j$ ($\sum_jjm_j=r$) is $r!/\prod_j(j!)^{m_j}m_j!$; hence
\[
\langle e^Z\rangle_0=\sum_{r\ge0}\frac1{r!}\langle Z^r\rangle_0=\sum_{(m_j)}\ \prod_j\frac1{m_j!}\Bigl(\frac{\varkappa_j}{j!}\Bigr)^{m_j}=\exp\sum_{j\ge1}\frac1{j!}\varkappa_j .
\]
We compare the coefficients of $\alpha_1\alpha_2$ in $\log\langle e^Z\rangle_0$. Since $\langle e^Z\rangle_0=\langle e^{-\varphi}\rangle_0+\alpha_1\langle\mathsf xe^{-\varphi}\rangle_0+\alpha_2\langle\mathsf ye^{-\varphi}\rangle_0+\alpha_1\alpha_2\langle\mathsf x\mathsf ye^{-\varphi}\rangle_0$ up to terms containing $\alpha_1^2$ or $\alpha_2^2$, this coefficient is $\mathsf e$. On the other hand, by multilinearity and symmetry of the cumulants, the coefficient of $\alpha_1\alpha_2$ in $\varkappa_j$ is $j(j-1)\langle\mathsf x;\mathsf y;-\varphi;\dots;-\varphi\rangle_0$, with $j-2$ entries $-\varphi$. Therefore
\[
\mathsf e=\sum_{j\ge2}\frac{j(j-1)}{j!}\langle\mathsf x;\mathsf y;-\varphi;\dots;-\varphi\rangle_0=\sum_{r\ge0}\frac{(-1)^r}{r!}\langle\mathsf x;\mathsf y;\underbrace{\varphi;\dots;\varphi}_r\rangle_0 .
\]
Inserting the series for $\mathsf x$, $\mathsf y$, $\varphi$, the coefficient of $t^0$ is $\langle X_0;Y_0\rangle_0=\gamma_0$. In the coefficient of $t^2$ the term $r=0$ contributes $\langle X_1;Y_1\rangle_0+\langle X_0;Y_2\rangle_0+\langle X_2;Y_0\rangle_0$, the term $r=1$ contributes $-\langle X_1;Y_0;W_1\rangle_0-\langle X_0;Y_1;W_1\rangle_0-\langle X_0;Y_0;W_2\rangle_0$, and the term $r=2$ contributes $\frac12\langle X_0;Y_0;W_1;W_1\rangle_0$; the sum is $\gamma_2$. The coefficient of an odd power $t^k$ is a sum of cumulants of polynomials $X_{j_1}$, $Y_{j_2}$, $W_{i_1},\dots,W_{i_r}$ with $j_1+j_2+i_1+\dots+i_r=k$. Each of $X_j$, $Y_j$, $\Phi_{j+2}$, $\Theta_j$ is homogeneous of a degree with the parity of $j$; hence, in each term of \eqref{eq:two-loop-cumulant-partition} for such a cumulant, at least one moment is the mean of a polynomial which is odd under $z\mapsto-z$, and it vanishes, because the Gaussian measure is invariant under $z\mapsto-z$. Thus $\gamma_1=\gamma_3=0$.
\end{proof}

\begin{proposition}[The asymptotic expansion in a box]\label{IV-73-27}
Let $n\ge2$ be an integer, $\tau_{\mathrm{gf}}>0$, and $p,q\in\mathsf{Pl}_n$. There is $C<\infty$, depending on $n$, $\tau_{\mathrm{gf}}$, $p$, $q$, such that for all $\beta\ge1$
\begin{equation}\label{IV-73-28}
\bigl|\langle a_p;a_q\rangle_{n,\beta}-\beta^{-2}\,\mathsf B_0^{[n]}(p,q)-\beta^{-3}\,\mathsf B_1^{[n]}(p,q;\tau_{\mathrm{gf}})\bigr|\le C\beta^{-4}.
\end{equation}
In particular, $\mathsf B_1^{[n]}(p,q;\tau_{\mathrm{gf}})$ does not depend on $\tau_{\mathrm{gf}}$.
\end{proposition}

\begin{proof}
We write $\tau:=\tau_{\mathrm{gf}}$, and we use $S$, $\mathsf T_n$, $\varrho_1$, $\varrho_2$, $\varrho_3$, $\mathfrak D$, $\Xi$, $\mathcal W$ of Lemma~\ref{IV-73-12}. Throughout, $F$ denotes one of the functions $1$, $a_p$, $a_q$, $a_pa_q$ on $\mathcal U_n$; each of them satisfies $|F|\le4$ and, like $S$, is invariant under $\mathcal G_n$ (Definition~\ref{IV-73-2}\textup{(1)}). Let $d_0:=\dim\mathsf X_n$ and $d_1:=\dim\mathsf{Form}_n^{\mathsf T}$. The constants $C,c>0$ below may depend on $n$, $\tau$, $p$, $q$.

\emph{Step 1: axial gauge.} By \eqref{IV-73-13} and the $\mathcal G_n$-invariance of $Fe^{-\beta S}$, and then by the Haar formula in Definition~\ref{IV-73-1}(2) at $u_0=\mathbf 1$, bond by bond,
\begin{align*}
\int Fe^{-\beta S}\prod_{b\in\mathsf{Bd}_n}dU(b)&=\int F(V)\,e^{-\beta S(V)}\prod_{b\in\mathsf{Bd}_n\setminus\mathsf T_n}dV(b)\\
&=\sigma_{\mathrm H}^{d_1/3}\int_{\{A\in\mathsf{Form}_n^{\mathsf T}:\,|A(b)|<\pi\}}\bigl(Fe^{-\beta S}\bigr)(e^{iA})\prod_{b\in\mathsf{Bd}_n\setminus\mathsf T_n}\sigma_1(|A(b)|^2)\,dA ,
\end{align*}
as $\mathsf{Bd}_n\setminus\mathsf T_n$ has $d_1/3$ elements. Let $\mathsf I(F)$ denote the last integral restricted to $|A|_\infty<\varrho_3$. The space $\mathcal U_n^{\mathsf T}$ is compact, being a finite power of the compact group $G$, and the image of $\{|A|_\infty<\varrho_3\}$ under $A\mapsto e^{iA}$ is relatively open in it (Definition~\ref{IV-73-1}); hence its complement $K$ in $\mathcal U_n^{\mathsf T}$ is compact, and by the injectivity of the exponential on $\mathbb B_\pi$, $e^{iA}\in K$ for the remaining $A$ in the integral. The continuous function $S$ has no zero on $K$, because the only zero of $S$ in $\mathcal U_n^{\mathsf T}$ is $\mathbf 1=e^{i0}$ (Lemma~\ref{IV-73-12}\textup{(b)}); so its minimum $c_*$ on $K$ is positive, and the integral over the remaining $A$ is bounded by $Ce^{-c_*\beta}$.

By Lemma~\ref{IV-73-9}\textup{(a)},\textup{(b)}, $A\mapsto S(e^{iA})$ is smooth and vanishes at $0$ together with its gradient; hence $S(e^{iA})\le C|A|_\infty^2$ for $|A|_\infty\le\varrho_3$, and since $\sigma_1(|A(b)|^2)\ge c$ there,
\[
\mathsf I(1)\ge c\int_{|A|_\infty<\min\{\varrho_3,\beta^{-1/2}\}}e^{-C\beta|A|_\infty^2}\,dA\ge c\,\beta^{-d_1/2}.
\]
As $|\mathsf I(F)|\le4\,\mathsf I(1)$, it follows that $\langle F\rangle_{n,\beta}=\mathsf I(F)/\mathsf I(1)+O(e^{-c\beta})$.

\emph{Step 2: the gauge fixing term.} We apply \eqref{IV-73-14} with the bounded Borel functions $\Phi(A'):=(Fe^{-\beta S})(e^{iA'})$ and $w(c):=e^{-\beta\|c\|^2/2\tau}$. Since $e^{iA^\lambda}=(e^{iA})^{e^{i\lambda}}$ and $Fe^{-\beta S}$ is $\mathcal G_n$-invariant, $\Phi(A^\lambda)=\Phi(A)$; thus the inner integral on the right side of \eqref{IV-73-14} is $\Phi(A)$ times the integral of $w$ over $\{\mathrm d^*(A^\lambda):|\lambda|_\infty<\varrho_1\}$, a set which contains $\{|c|_\infty<\varrho_2\}$ by Lemma~\ref{IV-73-12}\textup{(d)(iii)}. Since $\|c\|^2=\sum_{x\in\Lambda_n^\circ}|c(x)|^2$ and $dc$ is the product of the measures $dY$ at the $d_0/3$ sites of $\Lambda_n^\circ$ (Lemma~\ref{IV-73-12}\textup{(d)(iv)}), the substitution $c\to(\tau/\beta)^{1/2}c$ and the Gaussian integral in Definition~\ref{IV-73-1}(2), site by site, give
\[
\int_{\mathsf X_n}w\,dc=\Bigl(\frac\tau\beta\Bigr)^{d_0/2}\prod_{x\in\Lambda_n^\circ}\int_{\mathfrak g}e^{-\frac12|Y|^2}\,dY=\Bigl(\frac\tau\beta\Bigr)^{d_0/2}(2\pi)^{d_0/2}=\Bigl(\frac{2\pi\tau}\beta\Bigr)^{d_0/2},
\]
and the integral of $w$ over $\{|c|_\infty\ge\varrho_2\}$ is at most $Ce^{-c\beta}(2\pi\tau/\beta)^{d_0/2}$. Hence the inner integral equals $\Phi(A)(2\pi\tau/\beta)^{d_0/2}(1-\epsilon(A))$ with $0\le\epsilon(A)\le Ce^{-c\beta}$. The left side of \eqref{IV-73-14} is
\begin{gather*}
\mathsf J(F):=\int_{\mathcal W}F(e^{iA'})\,e^{-\beta\Phi_\tau(A')-\Theta(A')}\,dA',\\
\Phi_\tau(A'):=S(e^{iA'})+\tfrac1{2\tau}\|\mathrm d^*A'\|^2,\qquad\Theta(A'):=-\sum_{b\in\mathsf{Bd}_n}\log\sigma_1(|A'(b)|^2)-\log|\det\mathsf M(A')|;
\end{gather*}
here $\Theta$ is defined on $\overline{\mathcal W}$, since $\sigma_1>0$ on $\mathbb B_\pi$ and $\det\mathsf M\ne0$ on $\overline{\mathcal W}$ (Lemma~\ref{IV-73-12}\textup{(d)(ii)}). Comparing with the definition of $\mathsf I(F)$ we obtain
\[
\bigl|\mathsf J(F)-(2\pi\tau/\beta)^{d_0/2}\,\mathsf I(F)\bigr|\le Ce^{-c\beta}(2\pi\tau/\beta)^{d_0/2}\,\mathsf I(1).
\]
Together with Step~1, and since $\mathsf J(F)/\mathsf J(1)$ and $\mathsf I(F)/\mathsf I(1)$ are bounded by $4$, this gives $\langle F\rangle_{n,\beta}=\mathsf J(F)/\mathsf J(1)+O(e^{-c\beta})$. Consequently, up to $O(e^{-c\beta})$, $\langle a_p;a_q\rangle_{n,\beta}$ is the truncated two-point function $\langle X;Y\rangle_{\mathcal W,\beta}$ of the functions $X(A'):=a_p(e^{iA'})$ and $Y(A'):=a_q(e^{iA'})$ for the probability density on $\mathcal W$ proportional to $e^{-\beta\Phi_\tau-\Theta}$.

\emph{Step 3: the hypotheses of Lemma~\ref{IV-73-26}.} We identify $\mathbb R^m$ with the Euclidean space $\mathsf{Form}_n$, so that $m=\dim\mathsf{Form}_n$, by means of an orthonormal basis for the scalar product $(\cdot,\cdot)$ of Definition~\ref{IV-73-2}\textup{(2)}. Then the Lebesgue measure of $\mathbb R^m$ becomes the measure $dA'$ of Lemma~\ref{IV-73-12}\textup{(d)(iv)}, the norm $|y|$ becomes $|A'|:=\|A'\|$, and $(y,Hy)$ becomes $(A',HA')$. We take the set $\mathcal W$ and the functions $\Phi_\tau$, $\Theta$, $X$, $Y$ of Step~2. The set $\mathcal W$ is open and bounded and contains $0=\Xi(0,0)$, hence a ball $\{|y|<\rho\}$; its closure is compact (Heine--Borel), contained in $\{|A'|_\infty<1\}$, and $\det\mathsf M\ne0$ on it (Lemma~\ref{IV-73-12}\textup{(d)}).

The four functions are smooth on an open set containing $\overline{\mathcal W}$: the maps $A'\mapsto a_P(e^{iA'})$ are smooth for every plaquette $P$ (Definition~\ref{IV-73-1}); $Y\mapsto\log\sigma_1(|Y|^2)$ is smooth on $\mathbb B_\pi$ (Lemma~\ref{IV-73-9}\textup{(c)}); and by \eqref{IV-73-11} the entries of $\mathsf M(A')$ are smooth on $\{|A'|_\infty<1\}$, being polynomials in the components of the $A'(b)$ plus entries of the $\mathsf r(A'(b))$; hence $\det\mathsf M$ is continuous, it does not vanish on an open set containing $\overline{\mathcal W}$, and $\log|\det\mathsf M|$ is smooth there. By Lemma~\ref{IV-73-9}\textup{(a)},\textup{(b)}, $X$, $Y$ and $\Phi_\tau$ vanish at $0$ together with their gradients, and the homogeneous Taylor polynomials of $A'\mapsto a_P(e^{iA'})$ at $0$ are the $\mathsf a^{(j)}_P(A')$, $j\ge0$, by \eqref{IV-73-7} and the homogeneity of the $\mathsf a^{(j)}_P$.

In particular the Taylor polynomial of degree $2$ of $\Phi_\tau$ is $\mathsf S_2$ (by \eqref{IV-73-10} and \eqref{IV-73-16}), which is positive definite by Lemma~\ref{IV-73-12}\textup{(c)}. Since $\frac12(A',HA')=\mathsf S_2(A')=\frac12\bigl(A',(\mathrm d^*\mathrm d+\tau^{-1}\mathrm d\mathrm d^*)A'\bigr)$, the Hessian $H$ acts in each colour as $\mathrm d^*\mathrm d+\tau^{-1}\mathrm d\mathrm d^*$, so that $H^{-1}=\mathsf C^{[n]}_\tau$ in each colour, and the Gaussian measure of Lemma~\ref{IV-73-26} is the Gaussian family of $\Lambda_n$ and $\tau$ (Definition~\ref{IV-73-19}\textup{(1)}). Finally let $\rho'\in(0,\rho]$. The function $\Phi_\tau\circ\Xi$ is continuous on the compact set $\overline{\mathfrak D}$, and by the gauge invariance of $S$,
\[
\Phi_\tau(\Xi(\lambda,A))=S(e^{iA})+\tfrac1{2\tau}\|\mathrm d^*(A^\lambda)\|^2 .
\]
If this vanishes, then $S(e^{iA})=0$, hence $e^{iA}=\mathbf 1$ and $A=0$ (Lemma~\ref{IV-73-12}\textup{(b)}), and $\mathrm d^*(0^\lambda)=0=\mathrm d^*(0^0)$, hence $\lambda=0$ by the injectivity in Lemma~\ref{IV-73-12}\textup{(d)(iii)}. Therefore $\Phi_\tau\circ\Xi$ has a positive minimum on the compact set $\overline{\mathfrak D}\setminus\Xi^{-1}(\{|A'|<\rho'\})$, which does not contain $(0,0)$, and $\inf\{\Phi_\tau(A'):\ A'\in\mathcal W,\ |A'|\ge\rho'\}>0$.

\emph{Step 4: the Taylor polynomials.} By Step~3, $X_j=\mathsf a_p^{(j+2)}$, $Y_j=\mathsf a_q^{(j+2)}$, $\Phi_{\tau,3}=\mathsf S_3$ and $\Phi_{\tau,4}=\sum_P\mathsf a_P^{(4)}$, the sums being over $\mathsf{Pl}_n$. By Lemma~\ref{IV-73-9}\textup{(c)}, as $t\to0$,
\[
-\sum_b\log\sigma_1\bigl(t^2|A(b)|^2\bigr)=\tfrac13t^2\sum_b|A(b)|^2+O(t^4).
\]
By Lemma~\ref{IV-73-9}\textup{(d)}, $\mathsf M(tA)=\mathsf M^{(0)}+t\mathsf M^{(1)}(A)+t^2\mathsf M^{(2)}(A)+O(t^4)$, where $\mathsf M^{(0)}=-\mathrm d^*\mathrm d$ is invertible. With $\mathsf A:=(\mathsf M^{(0)})^{-1}\bigl(\mathsf M(tA)-\mathsf M^{(0)}\bigr)=O(t)$ we have $\det(\mathbf 1+\mathsf A)>0$ for small $t$, and the series $\log\det(\mathbf 1+\mathsf A)=\sum_{k\ge1}\frac{(-1)^{k-1}}k\operatorname{tr}\mathsf A^k$ gives
\begin{align*}
\log|\det\mathsf M(tA)|-\log|\det\mathsf M^{(0)}|={}&t\operatorname{tr}\bigl[(\mathsf M^{(0)})^{-1}\mathsf M^{(1)}(A)\bigr]\\
&+t^2\Bigl\{\operatorname{tr}\bigl[(\mathsf M^{(0)})^{-1}\mathsf M^{(2)}(A)\bigr]-\tfrac12\operatorname{tr}\bigl[\bigl((\mathsf M^{(0)})^{-1}\mathsf M^{(1)}(A)\bigr)^2\bigr]\Bigr\}+O(t^3).
\end{align*}
By \eqref{IV-73-18} we conclude that $\Theta_1=0$ and
\[
\Theta_2=\tfrac13\sum_b|A(b)|^2+\tfrac23\sum_b\mathsf g(b)|A(b)|^2-4\sum_{b,b'}\mathsf D(b;b')\mathsf D(b';b)\langle A(b),A(b')\rangle ,
\]
so that $W_1=\mathsf S_3$ and $W_2=\mathsf V_2$, by \eqref{IV-73-17}. Lemma~\ref{IV-73-26} now gives \eqref{IV-73-28}, with $\gamma_0=\langle\mathsf a_p^{(2)};\mathsf a_q^{(2)}\rangle_0=\mathsf B_0^{[n]}(p,q)$ by \eqref{IV-73-23} and $\gamma_2=\mathsf B_1^{[n]}(p,q;\tau)$ by \eqref{IV-73-24}; the error $O(e^{-c\beta})$ of Step~2 is bounded by $C\beta^{-4}$.

\emph{Step 5: independence of the gauge parameter.} Neither $\langle a_p;a_q\rangle_{n,\beta}$ nor $\mathsf B_0^{[n]}(p,q)=\frac{\mathsf n}2\Pi^{[n]}(p,q)^2$ depends on $\tau$ (Definition~\ref{IV-73-15}\textup{(2)}). Multiplying \eqref{IV-73-28} by $\beta^3$ we obtain
\[
\beta^3\bigl[\langle a_p;a_q\rangle_{n,\beta}-\beta^{-2}\mathsf B_0^{[n]}(p,q)\bigr]=\mathsf B_1^{[n]}(p,q;\tau)+O(\beta^{-1}),
\]
and letting $\beta\to\infty$ we see that $\mathsf B_1^{[n]}(p,q;\tau)$ does not depend on $\tau$.
\end{proof}

\subsection{Zero momentum}

In this subsection we compute the values at zero momentum of the self-energy and of the vertex function of a plaquette action. Both are expressed through Hessians of the action at constant abelian fields, whose symbols can be computed explicitly.

\begin{definition}[Hessians; self-energy and vertex function; their values at zero momentum]
\label{IV-73-29}
We work on $\mathbb Z^3$ and fix a gauge parameter $\tau_{\mathrm{gf}}>0$. For a positively oriented bond $b$ and a colour index $a$ we write $\partial_{b,a}$ for the derivative with respect to the variable $A^a(b)$. On a formal sum $\mathsf F=\sum_z\mathsf F_z$ of local terms (Definition~\ref{IV-73-19}(3)) the derivative acts term by term, and the terms on which it vanishes are omitted. The derivatives $\partial_{b,a}\mathsf S_3$, $\partial_{b,a}\partial_{b',a'}\mathsf S_4$ and $\partial_{b,a}\partial_{b',a'}\mathsf V_2$ which occur below are then finite sums: a bond $(x,\mu)$ belongs to the two plaquettes on its two sides in each of the other two coordinate directions, hence to four plaquettes; of a sum over bonds $b''$ only the term with $b''=b$ survives; and of the sum over pairs of bonds in $\mathsf V_2$ only the terms of the pairs $(b,b')$ and $(b',b)$ survive.

\smallskip\noindent\emph{(1) Hessians.} Let $B$ be a $1$-form on $\mathbb Z^3$ with $\sup_b|B(b)|<1$, and let $q$ be a plaquette. We define real symmetric matrices $\mathsf N(B)$ and $\mathsf h_q(B)$, indexed by the pairs $(b,a)$ of a positively oriented bond $b$ and a colour index $a$, by
\[
\begin{multlined}
\mathsf N(B)_{(b,a),(b',a')}:=\partial_{b,a}\partial_{b',a'}\Bigl[\sum_pa_p(e^{iA})+\frac1{2\tau_{\mathrm{gf}}}\sum_x|\mathrm d^*A(x)|^2\Bigr]_{A=B},\qquad
\mathsf h_q(B)_{(b,a),(b',a')}: \\
=\partial_{b,a}\partial_{b',a'}\,a_q(e^{iA})\big|_{A=B}.
\end{multlined}
\]
The infinite sums over all plaquettes $p$ and over all sites $x$ are understood as follows: under $\partial_{b,a}\partial_{b',a'}$ only the finitely many terms which depend on $A(b)$ contribute, namely those of the plaquettes containing $b$ and those of the two end points of $b$. Consequently an entry of either matrix vanishes unless $b$ and $b'$ lie in a common plaquette or have a common end point. For fixed $(b,a)$ only finitely many entries $\mathsf N(B)_{(b,a),(b',a')}$ and $\mathsf h_q(B)_{(b,a),(b',a')}$ are different from $0$, and both matrices act on bounded complex $1$-forms $A$ by $(\mathsf N(B)A)^a(b):=\sum_{(b',a')}\mathsf N(B)_{(b,a),(b',a')}A^{a'}(b')$, and likewise for $\mathsf h_q(B)$.

\smallskip\noindent\emph{(2) Constant abelian fields.} For $\gamma\in\mathbb R^3$ with $|\gamma|<1$ let $Y_\gamma$ be the constant $1$-form $Y_\gamma((x,\mu)):=\gamma_\mu\mathsf e_3$. We put $\mathsf e_\pm:=2^{-1/2}(\mathsf e_1\pm i\mathsf e_2)$ and $\mathsf e_0:=\mathsf e_3$. These three vectors form a Hermitian-orthonormal basis of $\mathfrak g\otimes\mathbb C$, and
\[
[\mathsf e_3,\mathsf e_e]=2e\,\mathsf e_e\qquad(e\in\{+1,-1,0\});
\]
we call $e$ the \emph{charge}, and we write $\upsilon:=2e\gamma\in\mathbb R^3$. For $k\in\mathbb R^3$ we define two $3\times3$ matrices $\hat{\mathsf N}_e(k;\gamma)$, $\hat{\mathsf h}_{q,e}(k;\gamma)$ and a number $\hat{\mathsf M}_e(k;\gamma)$ by
\[
\begin{multlined}
\mathsf N(Y_\gamma)\bigl(e^{ik\cdot x}\alpha_\mu\mathsf e_e\bigr)=e^{ik\cdot x}(\hat{\mathsf N}_e\alpha)_\mu\,\mathsf e_e,\qquad
\bigl(e^{ik\cdot x}\alpha'\mathsf e_e,\,\mathsf h_q(Y_\gamma)\,e^{ik\cdot x}\alpha\,\mathsf e_e\bigr)=\alpha'^*\hat{\mathsf h}_{q,e}\alpha, \\
\mathsf M(Y_\gamma)\bigl(e^{ik\cdot x}\mathsf e_e\bigr)=\hat{\mathsf M}_e\,e^{ik\cdot x}\mathsf e_e
\end{multlined}
\]
for all $\alpha,\alpha'\in\mathbb C^3$. Here $e^{ik\cdot x}\alpha_\mu\mathsf e_e$ stands for the complex $1$-form $(x,\mu)\mapsto e^{ik\cdot x}\alpha_\mu\mathsf e_e$ and $e^{ik\cdot x}\mathsf e_e$ for the complex site function $x\mapsto e^{ik\cdot x}\mathsf e_e$; the bracket $(\cdot,\cdot)$ is Hermitian, and the sum which defines it is finite, since $\mathsf h_q(Y_\gamma)$ has non-zero entries only at bonds of $q$; the operator $\mathsf M(Y_\gamma)$ of Definition~\ref{IV-73-6}(3) is extended complex linearly. That the left sides have these forms is shown in Lemma~\ref{IV-73-31}, which also gives explicit formulas; by these formulas, for fixed $k$, the three objects are real analytic functions of $\gamma$ in $\{|\gamma|<1\}$.

A second index $j$ denotes the homogeneous polynomial of degree $j$ in $\gamma$ in the Taylor expansion at $\gamma=0$, so that $\hat{\mathsf N}_e=\sum_j\hat{\mathsf N}_{e,j}$, $\hat{\mathsf h}_{q,e}=\sum_j\hat{\mathsf h}_{q,e,j}$ and $\hat{\mathsf M}_e=\sum_j\hat{\mathsf M}_{e,j}$. For $k\notin2\pi\mathbb Z^3$ we put $\hat{\mathsf C}:=\hat{\mathsf N}_{e,0}^{-1}$. By Lemma~\ref{IV-73-31} at $\gamma=0$ we have $\hat{\mathsf N}_{e,0}(k)=|\mathsf k|^2\mathbf 1-(1-\tau_{\mathrm{gf}}^{-1})\mathsf k\mathsf k^*$ with $\mathsf k=\mathsf k(k)$, so that $\hat{\mathsf C}$ does not depend on $e$ and is the matrix with the entries
\[
\hat{\mathsf C}_{\mu\nu}(k)=\frac{\delta_{\mu\nu}}{|\mathsf k|^2}+(\tau_{\mathrm{gf}}-1)\frac{\mathsf k_\mu\overline{\mathsf k_\nu}}{|\mathsf k|^4}
\]
which appears in the formula for $\mathsf C^{\mathrm{lat}}_{\tau_{\mathrm{gf}}}$ in Definition~\ref{IV-73-15}(1); moreover $\hat{\mathsf M}_{e,0}=-|\mathsf k|^2$.

\smallskip\noindent\emph{(3) Five numbers of the self-energy.} For $k\in[-\pi,\pi]^3\setminus\{0\}$ and $\gamma\in\mathbb R^3$ we put, with all functions taken at $(k;\gamma)$,
\[
\begin{aligned}
\mathfrak s^{(1)}&:=-\tfrac14\sum_e\operatorname{tr}\bigl[(\hat{\mathsf C}\hat{\mathsf N}_{e,1})^2\bigr]&&\text{(gluon loop)},\\
\mathfrak s^{(2)}&:=\tfrac12\sum_e\operatorname{tr}\bigl[\hat{\mathsf C}\hat{\mathsf N}_{e,2}\bigr]&&\text{(gluon tadpole)},\\
\mathfrak s^{(3)}&:=\tfrac12\sum_e\operatorname{Re}\bigl[(\hat{\mathsf M}_{e,1}/\hat{\mathsf M}_{e,0})^2\bigr]&&\text{(ghost loop)},\\
\mathfrak s^{(4)}&:=-\sum_e\operatorname{Re}\bigl[\hat{\mathsf M}_{e,2}/\hat{\mathsf M}_{e,0}\bigr]&&\text{(ghost tadpole)},
\end{aligned}
\]
and $\mathfrak s^{(5)}(k;\gamma):=\frac13|\gamma|^2$ (Haar term). For fixed $k$ each $\mathfrak s^{(i)}(k;\cdot)$ is a homogeneous polynomial of degree $2$ in $\gamma$, bounded by $C|\gamma|^2|k|^{-2}$ (for $i\le4$ by Theorem~\ref{IV-73-34}(i), for $i=5$ because $|k|$ is bounded on the cube); hence it is integrable over $\mathbb T$, and its integral is a quadratic form in $\gamma$. For $i=1,\dots,5$ we let $c^{(i)}_{\mu\nu}(\tau_{\mathrm{gf}})$, $1\le\mu,\nu\le3$, be the entries of the symmetric matrix with
\[
\frac12\sum_{\mu\nu}c^{(i)}_{\mu\nu}(\tau_{\mathrm{gf}})\,\gamma_\mu\gamma_\nu=\int_{\mathbb T}\mathfrak s^{(i)}(k;\gamma)\,d^\circ k\qquad(\gamma\in\mathbb R^3).
\]
We also need the number
\begin{equation}\label{IV-73-30}
Z_0:=\int_{\mathbb T}\frac{d^\circ k}{|\mathsf k(k)|^2}=\int_{[-\pi,\pi]^3}\frac{dk}{(2\pi)^3}\Bigl(\sum_\mu4\sin^2\frac{k_\mu}2\Bigr)^{-1}.
\end{equation}

\smallskip\noindent\emph{(4) Four numbers of the vertex function of $a_q$.} For a plaquette $q$, $k\in[-\pi,\pi]^3\setminus\{0\}$ and $\gamma\in\mathbb R^3$ we put, again with all functions taken at $(k;\gamma)$,
\[
\begin{aligned}
\mathfrak v_q^{(1)}&:=\tfrac12\sum_e\operatorname{tr}\bigl[\hat{\mathsf h}_{q,e,2}\hat{\mathsf C}\bigr],&
\mathfrak v_q^{(2)}&:=-\tfrac12\sum_e\operatorname{tr}\bigl[\hat{\mathsf h}_{q,e,1}\hat{\mathsf C}\hat{\mathsf N}_{e,1}\hat{\mathsf C}\bigr],\\
\mathfrak v_q^{(3)}&:=-\tfrac12\sum_e\operatorname{tr}\bigl[\hat{\mathsf h}_{q,e,0}\hat{\mathsf C}\hat{\mathsf N}_{e,2}\hat{\mathsf C}\bigr],&
\mathfrak v_q^{(4)}&:=\tfrac12\sum_e\operatorname{tr}\bigl[\hat{\mathsf h}_{q,e,0}\hat{\mathsf C}\hat{\mathsf N}_{e,1}\hat{\mathsf C}\hat{\mathsf N}_{e,1}\hat{\mathsf C}\bigr].
\end{aligned}
\]
These are homogeneous polynomials of degree $2$ in $\gamma$, bounded by $C|\gamma|^2|k|^{-2}$ (Theorem~\ref{IV-73-36}). For $j=1,\dots,4$ we let $\mathsf g^{(j)}_{q;\mu\nu}(\tau_{\mathrm{gf}})$ be the entries of the symmetric matrix with
\[
\frac12\sum_{\mu\nu}\mathsf g^{(j)}_{q;\mu\nu}(\tau_{\mathrm{gf}})\,\gamma_\mu\gamma_\nu=\int_{\mathbb T}\mathfrak v_q^{(j)}(k;\gamma)\,d^\circ k\qquad(\gamma\in\mathbb R^3).
\]

\smallskip\noindent\emph{(5) Kernels.} Let $\langle\cdot\rangle_0$ refer to the Gaussian family of $\mathbb Z^3$ and $\tau_{\mathrm{gf}}$ (Definition~\ref{IV-73-19}(1)), and write $\partial=\partial_{b,a}$, $\partial'=\partial_{b',a'}$. We define
\[
\begin{aligned}
\Sigma_{\mathrm{tot}}(b,a;b',a')&:=\langle\partial\partial'\mathsf V_2\rangle_0-\langle\partial\mathsf S_3;\partial'\mathsf S_3\rangle_0,\\
\Gamma_q(b,a;b',a')&:=\langle\partial\partial'\mathsf a_q^{(4)}\rangle_0-\langle\partial\mathsf a_q^{(3)};\partial'\mathsf S_3\rangle_0-\langle\partial'\mathsf a_q^{(3)};\partial\mathsf S_3\rangle_0-\langle\mathsf a_q^{(2)};\partial\partial'\mathsf S_4\rangle_0+\langle\mathsf a_q^{(2)};\partial\mathsf S_3;\partial'\mathsf S_3\rangle_0 .
\end{aligned}
\]
These are means and cumulants (Definition~\ref{IV-73-19}(2)) of polynomials in the Gaussian family, and both kernels are symmetric under the exchange $(b,a)\leftrightarrow(b',a')$.

We determine the colour structure of the two kernels. For $g\in G$ the map $\operatorname{Ad}g:X\mapsto gXg^{-1}$ is orthogonal on $(\mathfrak g,\langle\cdot,\cdot\rangle)$; let $(\operatorname{Ad}g)_{aa'}$ be its matrix in the basis $\mathsf e_1,\mathsf e_2,\mathsf e_3$. The components $(gA(b)g^{-1})^a=\sum_{a'}(\operatorname{Ad}g)_{aa'}A^{a'}(b)$ of the $\mathfrak g$-valued family $(gA(b)g^{-1})_b$ are real linear combinations of the variables of the Gaussian family, so that this family is again centred Gaussian. By Definition~\ref{IV-73-19}(1) and the orthogonality of $\operatorname{Ad}g$, its covariance is
\[
\bigl\langle(gA(b)g^{-1})^a\,(gA(b')g^{-1})^{c}\bigr\rangle_0=\sum_{a'}(\operatorname{Ad}g)_{aa'}(\operatorname{Ad}g)_{ca'}\,\mathsf C^{\mathrm{lat}}_{\tau_{\mathrm{gf}}}(b,b')=\delta_{ac}\,\mathsf C^{\mathrm{lat}}_{\tau_{\mathrm{gf}}}(b,b'),
\]
which is the covariance of $(A(b))_b$; as a centred Gaussian family is determined in law by its covariance, $(gA(b)g^{-1})_b$ has the law of $(A(b))_b$. By Lemma~\ref{IV-73-9}(a) the polynomials $\mathsf a^{(m)}_p$ are invariant under $A\mapsto gAg^{-1}$, and so are $\mathsf S_3$, $\mathsf S_4$ and, term by term, $\mathsf V_2$, whose remaining terms involve $A$ only through $|A(b)|^2$ and $\langle A(b),A(b')\rangle$. By the chain rule, each of the two kernels, regarded at fixed $b,b'$ as a matrix $K$ in $(a,a')$, satisfies $K=(\operatorname{Ad}g)^{\mathsf T}K\operatorname{Ad}g$, that is, it commutes with every $\operatorname{Ad}g$. For $g=e^{isY}$, $Y\in\mathfrak g$, we have $\operatorname{Ad}g=e^{s\mathsf a_Y}$ (Lemma~\ref{IV-73-9}(d)), and differentiation at $s=0$ shows that $K$ commutes with $\mathsf e_c\times\cdot$ for $c=1,2,3$.

Such a matrix is a multiple of $\mathbf 1$: it preserves $\ker(\mathsf e_c\times\cdot)=\mathbb R\mathsf e_c$, so it is diagonal in the basis $\mathsf e_1,\mathsf e_2,\mathsf e_3$, with entries $m_1,m_2,m_3$ say; applying it to $\mathsf e_1\times\mathsf e_2=\mathsf e_3$ and to $\mathsf e_3\times\mathsf e_1=\mathsf e_2$ gives $m_3=m_2$ and $m_2=m_1$. Hence
\[
\Sigma_{\mathrm{tot}}(b,a;b',a')=\delta_{aa'}\,\Sigma_{\mathrm{tot}}(b,b'),\qquad\Gamma_q(b,a;b',a')=\delta_{aa'}\,\Gamma_q(b,b'),
\]
and these identities define the kernels $\Sigma_{\mathrm{tot}}(b,b')$ and $\Gamma_q(b,b')$. Finally, the kernels are covariant under the translations of $\mathbb Z^3$. Indeed, $\mathsf C^{\mathrm{lat}}_{\tau_{\mathrm{gf}}}((x,\mu),(y,\nu))$ and $\mathsf G^{\mathrm{lat}}(x,y)$ depend on the sites $x,y$ only through $x-y$ (Definition~\ref{IV-73-15}(1)), so that the law of the Gaussian family is invariant under translations; the coefficients of the polynomials $\mathsf a_p^{(m)}$ do not depend on $p$ (Lemma~\ref{IV-73-9}(a)); and $\mathsf g(b)=\frac13$ as well as $\mathsf D(b;b')$, which are built from $\mathsf G^{\mathrm{lat}}$ (Definition~\ref{IV-73-15}(2)), do not change when $b$ and $b'$ are translated together. Hence the formal sums $\mathsf S_3$, $\mathsf S_4$, $\mathsf V_2$ are invariant under translations as well, and $\Sigma_{\mathrm{tot}}(b+y,b'+y)=\Sigma_{\mathrm{tot}}(b,b')$ and $\Gamma_{q+y}(b+y,b'+y)=\Gamma_q(b,b')$ for all $y\in\mathbb Z^3$.
\end{definition}

The symbols of Definition~\ref{IV-73-29}(2) can be computed explicitly.

\begin{lemma}[The Hessian at a constant abelian field]
\label{IV-73-31}
Let $\gamma\in\mathbb R^3$ with $|\gamma|<1$, let $e\in\{\pm1,0\}$, $\upsilon=2e\gamma$ and $k\in\mathbb R^3$. Let $j(\vartheta):=(e^{i\vartheta}-1)/i\vartheta$ for real $\vartheta\ne0$ and $j(0):=1$, and put
\[
J:=\operatorname{diag}\bigl(j(\upsilon_1),j(\upsilon_2),j(\upsilon_3)\bigr),\qquad P(k):=|\mathsf k(k)|^2\mathbf 1-\mathsf k(k)\mathsf k(k)^*;
\]
the matrix $J$ is invertible, since $|\upsilon_\mu|<2$ and $j(\vartheta)\ne0$ for $|\vartheta|<2\pi$. For a plaquette $q=p_{\mu\nu}(x)$ let $D_q(k)$ be the $1\times3$ matrix with $D_q(k)\alpha=\mathsf k_\mu(k)\alpha_\nu-\mathsf k_\nu(k)\alpha_\mu$, $\alpha\in\mathbb C^3$. Then the left sides of the three identities in Definition~\ref{IV-73-29}(2) have the forms stated there, with
\[
\begin{multlined}
\hat{\mathsf N}_e(k;\gamma)=J^*P(k+\upsilon)J+\tau_{\mathrm{gf}}^{-1}\mathsf k(k)\mathsf k(k)^*,\qquad\hat{\mathsf h}_{q,e}(k;\gamma)=J^*D_q(k+\upsilon)^*D_q(k+\upsilon)J, \\
\hat{\mathsf M}_e(k;\gamma)=-\mathsf k(k)^*J^{-1}\mathsf k(k+\upsilon).
\end{multlined}
\]
Moreover the following hold.
\begin{enumerate}
\item[(a)] We have
\begin{equation}\label{IV-73-32}
\det\hat{\mathsf N}_e(k;\gamma)=\tau_{\mathrm{gf}}^{-1}\,|\det J|^2\,|\mathsf k(k+\upsilon)|^2\,|\hat{\mathsf M}_e(k;\gamma)|^2,\qquad\prod_{e}|\det J|=\prod_\mu\sigma_1(\gamma_\mu^2),
\end{equation}
where in the second identity $J$ is taken at $\upsilon=2e\gamma$ for each $e\in\{+1,-1,0\}$.
\item[(b)] If $\mathsf k(k+\upsilon)\ne0$ and $\hat{\mathsf M}_e(k;\gamma)\ne0$, then $\hat{\mathsf N}_e(k;\gamma)$ is invertible, and for $q=p_{\mu\nu}(x)$
\begin{equation}\label{IV-73-33}
\operatorname{tr}\bigl[\hat{\mathsf h}_{q,e}(k;\gamma)\hat{\mathsf N}_e(k;\gamma)^{-1}\bigr]=\varpi_q(k+\upsilon),\qquad\varpi_q(k):=\frac{|\mathsf k_\mu(k)|^2+|\mathsf k_\nu(k)|^2}{|\mathsf k(k)|^2}.
\end{equation}
\item[(c)] With $\mathsf k=\mathsf k(k)$,
\[
\hat{\mathsf M}_{e,1}=-\sum_\mu\upsilon_\mu\sin k_\mu,\qquad\hat{\mathsf M}_{e,2}=\frac1{12}\sum_\mu\upsilon_\mu^2|\mathsf k_\mu|^2,\qquad\operatorname{tr}\hat{\mathsf N}_{e,2}=\sum_\mu\upsilon_\mu^2\Bigl[2\cos k_\mu-\frac1{12}\bigl(|\mathsf k|^2-|\mathsf k_\mu|^2\bigr)\Bigr].
\]
\end{enumerate}
\end{lemma}

\begin{proof}
\emph{Step 1: second variation.} Let $a$ be a $1$-form of finite support, and let $u_\mu:=e^{i\gamma_\mu\mathsf e_3}$; these three matrices commute with each other. For a bond $b$ and small real $s$ the product $e^{i(Y_\gamma(b)+sa(b))}e^{-iY_\gamma(b)}$ lies close to $\mathbf 1$, in particular in $G\setminus\{-\mathbf 1\}$. By Definition~\ref{IV-73-1}(2) the map $Y\mapsto e^{iY}$ is a diffeomorphism of $\mathbb B_\pi$ onto $G\setminus\{-\mathbf 1\}$, with the smooth inverse $\frac1i\log$; we may therefore write the product as $e^{i\eta_s(b)}$, where $\eta_s(b):=\frac1i\log\bigl(e^{i(Y_\gamma(b)+sa(b))}e^{-iY_\gamma(b)}\bigr)\in\mathfrak g$ is a smooth function of $s$ with $\eta_0(b)=0$. The differential formula in Definition~\ref{IV-73-1}(2), applied with $Y=Y_\gamma(b)$ and $H=a(b)$, shows that the derivative of the product at $s=0$ equals $i\,\mathsf{dE}(Y_\gamma(b))a(b)$; hence
\[
\eta_s(b)=s\,\mathsf{dE}(Y_\gamma(b))\,a(b)+O(s^2).
\]
Let $p=p_{\mu\nu}(x)$ have the bonds $b_1,\dots,b_4$, and put $\eta_i:=\eta_s(b_i)$. Since $Y_\gamma(b_1)=Y_\gamma(b_3)=\gamma_\mu\mathsf e_3$ and $Y_\gamma(b_2)=Y_\gamma(b_4)=\gamma_\nu\mathsf e_3$, since $u_\mu$ and $u_\nu$ commute, and since $ue^{iX}u^{-1}=e^{iuXu^{-1}}$, we find
\[
\partial\bigl(e^{i(Y_\gamma+sa)}\bigr)(p)=e^{i\eta_1}u_\mu\,e^{i\eta_2}u_\nu\,u_\mu^{-1}e^{-i\eta_3}\,u_\nu^{-1}e^{-i\eta_4}=e^{i\eta_1}\,e^{iu_\mu\eta_2u_\mu^{-1}}\,e^{-iu_\nu\eta_3u_\nu^{-1}}\,e^{-i\eta_4}=e^{iZ},
\]
where $Z\in\mathfrak g$ is the image of this product under $\frac1i\log$; for small $s$ the product is close to $\mathbf 1$, and by the diffeomorphism property recalled above $Z$ is a smooth function of $s$, with $Z=0$ at $s=0$. To first order
\[
Z=\eta_1+u_\mu\eta_2u_\mu^{-1}-u_\nu\eta_3u_\nu^{-1}-\eta_4+O(s^2).
\]
Since $a_p(U)=1-\frac12\operatorname{tr}\partial U(p)$ and $\frac12\operatorname{tr}e^{iZ}=\cos|Z|$, we have $a_p(e^{i(Y_\gamma+sa)})=1-\cos|Z|=\frac12|Z|^2+O(|Z|^4)$. We put $(\mathsf J_\gamma a)(b):=\mathsf{dE}(Y_\gamma(b))a(b)$ and, for a $1$-form $a'$, $(\mathsf D_\gamma a')(p):=a'(b_1)+u_\mu a'(b_2)u_\mu^{-1}-u_\nu a'(b_3)u_\nu^{-1}-a'(b_4)$; by the expansion of $\eta_s(b)$, $Z=s\,(\mathsf D_\gamma\mathsf J_\gamma a)(p)+O(s^2)$, and therefore
\[
a_p\bigl(e^{i(Y_\gamma+sa)}\bigr)=\tfrac12s^2\bigl|(\mathsf D_\gamma\mathsf J_\gamma a)(p)\bigr|^2+O(s^3).
\]
The constant form $Y_\gamma$ satisfies $\mathrm d^*Y_\gamma=0$, so that $\frac1{2\tau_{\mathrm{gf}}}\sum_x|\mathrm d^*(Y_\gamma+sa)(x)|^2=\frac{s^2}{2\tau_{\mathrm{gf}}}\|\mathrm d^*a\|^2$. Taking second derivatives in $s$ at $s=0$ we obtain the quadratic forms of the two Hessians:
\[
\bigl(a,\mathsf h_p(Y_\gamma)a\bigr)=\bigl|(\mathsf D_\gamma\mathsf J_\gamma a)(p)\bigr|^2,\qquad\bigl(a,\mathsf N(Y_\gamma)a\bigr)=\sum_p\bigl|(\mathsf D_\gamma\mathsf J_\gamma a)(p)\bigr|^2+\tau_{\mathrm{gf}}^{-1}\|\mathrm d^*a\|^2 ,
\]
the sum over $p$ being finite. A real symmetric matrix of the kind considered in Definition~\ref{IV-73-29}(1) is determined by its quadratic form on the forms of finite support; by polarisation, therefore,
\[
\mathsf h_p(Y_\gamma)=(\delta_p\mathsf D_\gamma\mathsf J_\gamma)^*(\delta_p\mathsf D_\gamma\mathsf J_\gamma),\qquad\mathsf N(Y_\gamma)=(\mathsf D_\gamma\mathsf J_\gamma)^*\mathsf D_\gamma\mathsf J_\gamma+\tau_{\mathrm{gf}}^{-1}\mathrm d\mathrm d^*,
\]
where $\delta_p$ is the evaluation at $p$ and the adjoints are taken with respect to the brackets $(\cdot,\cdot)$.

\emph{Step 2: symbols.} We extend all operators complex linearly. From $[\mathsf e_3,\mathsf e_e]=2e\,\mathsf e_e$ we get $u_\mu\mathsf e_eu_\mu^{-1}=e^{2ie\gamma_\mu}\mathsf e_e=e^{i\upsilon_\mu}\mathsf e_e$, and hence
\[
\mathsf{dE}(\gamma_\mu\mathsf e_3)\,\mathsf e_e=\int_0^1e^{is\upsilon_\mu}\,ds\;\mathsf e_e=j(\upsilon_\mu)\,\mathsf e_e .
\]
So $\mathsf J_\gamma$ maps the form $e^{ik\cdot x}\alpha\,\mathsf e_e$ to $e^{ik\cdot x}J\alpha\,\mathsf e_e$, and for $\alpha'\in\mathbb C^3$
\[
\mathsf D_\gamma\bigl(e^{ik\cdot x}\alpha'\mathsf e_e\bigr)(p_{\mu\nu}(x))=e^{ik\cdot x}\bigl[\alpha'_\mu+e^{i(\upsilon_\mu+k_\mu)}\alpha'_\nu-e^{i(\upsilon_\nu+k_\nu)}\alpha'_\mu-\alpha'_\nu\bigr]\mathsf e_e=e^{ik\cdot x}D_{p_{\mu\nu}(x)}(k+\upsilon)\,\alpha'\,\mathsf e_e .
\]
Thus $\mathsf J_\gamma$, $\mathsf D_\gamma$ and $\mathrm d\mathrm d^*$ commute with the translations of $\mathbb Z^3$ and map a plane wave of momentum $k$ and charge $e$ to a plane wave of the same momentum and charge, the coefficient vector being multiplied by $J$, by the $1\times3$ matrices $D_{p_{\mu\nu}(x)}(k+\upsilon)$ ($\mu<\nu$), and by $\mathsf k(k)\mathsf k(k)^*$, respectively; the last symbol follows from the plane-wave formulas of Definition~\ref{IV-73-15}(1). The adjoint of such an operator acts on plane waves through the adjoint matrix. By Lagrange's identity,
\[
\sum_{\mu<\nu}|v_\mu\alpha_\nu-v_\nu\alpha_\mu|^2=|v|^2|\alpha|^2-|v^*\alpha|^2=\alpha^*\bigl(|v|^2\mathbf 1-vv^*\bigr)\alpha\qquad(v,\alpha\in\mathbb C^3),
\]
so that $\sum_{\mu<\nu}D_{p_{\mu\nu}(x)}(k')^*D_{p_{\mu\nu}(x)}(k')=P(k')$ for $k'\in\mathbb R^3$. Together with Step~1 this gives the formula for $\hat{\mathsf N}_e$; for $\hat{\mathsf h}_{q,e}$ only the plaquette $q$ enters, and the factors $e^{ik\cdot x}$ have modulus $1$.

For $\hat{\mathsf M}_e$ we use the middle expression in \eqref{IV-73-11}, which is real linear in $\lambda$ and is extended complex linearly, together with $e^{\mathsf a_Y}\mathsf e_e=u_\mu\mathsf e_eu_\mu^{-1}$ for $Y=\gamma_\mu\mathsf e_3$ (Lemma~\ref{IV-73-9}(d)). At $A=Y_\gamma$, $\lambda=e^{ik\cdot x}\mathsf e_e$ and the bond $(x,\mu)$ it gives
\[
\begin{multlined}
\frac d{ds}A^{s\lambda}((x,\mu))\Big|_{s=0}=\mathsf{dE}(\gamma_\mu\mathsf e_3)^{-1}\bigl(e^{ik\cdot x}\mathsf e_e-e^{i\upsilon_\mu}e^{ik\cdot(x+e_\mu)}\mathsf e_e\bigr)=j(\upsilon_\mu)^{-1}\bigl(1-e^{i(\upsilon_\mu+k_\mu)}\bigr)e^{ik\cdot x}\mathsf e_e= \\
-j(\upsilon_\mu)^{-1}\mathsf k_\mu(k+\upsilon)\,e^{ik\cdot x}\mathsf e_e .
\end{multlined}
\]
Since $\mathsf M(A)\lambda$ is $\mathrm d^*$ of this derivative and $\mathrm d^*$ multiplies plane waves by $\mathsf k(k)^*$ (Definition~\ref{IV-73-15}(1)), we obtain $\hat{\mathsf M}_e=-\mathsf k(k)^*J^{-1}\mathsf k(k+\upsilon)$. Here we have used that $j(\vartheta)\ne0$ for $|\vartheta|<2\pi$.

\emph{Step 3: proof of (a).} Let $v:=\mathsf k(k+\upsilon)$ and $\tilde u:=(J^*)^{-1}\mathsf k(k)$. Then
\[
\hat{\mathsf N}_e=J^*\bigl(P(k+\upsilon)+\tau_{\mathrm{gf}}^{-1}\tilde u\tilde u^*\bigr)J,\qquad v^*\tilde u=(J^{-1}v)^*\mathsf k(k)=-\overline{\hat{\mathsf M}_e},
\]
and $P(k+\upsilon)=|v|^2\mathbf 1-vv^*$. For $v,w\in\mathbb C^3$ we have $\det\bigl(|v|^2\mathbf 1-vv^*+ww^*\bigr)=|v|^2|v^*w|^2$: for $v=0$ both sides vanish; for $v\ne0$ we use an orthonormal basis whose first vector is $v/|v|$; in it the matrix is $|v|^2\operatorname{diag}(0,1,1)+ww^*$, the determinant is an affine function of the rank-one summand, and its value is $0+|v|^4|w_1|^2$ with $w_1=v^*w/|v|$. With $w=\tau_{\mathrm{gf}}^{-1/2}\tilde u$ this gives $\det\hat{\mathsf N}_e=|\det J|^2\,|v|^2\,\tau_{\mathrm{gf}}^{-1}|v^*\tilde u|^2$, the first identity in \eqref{IV-73-32}. For the second, $|j(\vartheta)|^2=(2-2\cos\vartheta)/\vartheta^2=\bigl(\sin\frac\vartheta2/\frac\vartheta2\bigr)^2$; since $\det J=\prod_\mu j(2e\gamma_\mu)$, the product over the three charges contributes, for each $\mu$, the factor $|j(2\gamma_\mu)|^2=(\sin\gamma_\mu/\gamma_\mu)^2=\sigma_1(\gamma_\mu^2)$.

\emph{Step 4: proof of (b).} By (a), $\hat{\mathsf N}_e$ is invertible. By the cyclicity of the trace,
\[
\operatorname{tr}\bigl[\hat{\mathsf h}_{q,e}\hat{\mathsf N}_e^{-1}\bigr]=D_q\bigl(P+\tau_{\mathrm{gf}}^{-1}\tilde u\tilde u^*\bigr)^{-1}D_q^*,
\]
with $D_q$ and $P$ taken at $k+\upsilon$. Now $D_qv=v_\mu v_\nu-v_\nu v_\mu=0$, hence $v^*D_q^*=0$ and $PD_q^*=|v|^2D_q^*$; moreover $Pv=0$, and $\tilde u^*v\ne0$ because $v^*\tilde u=-\overline{\hat{\mathsf M}_e}\ne0$. The vector
\[
\mathsf X:=|v|^{-2}D_q^*-\frac{\tilde u^*D_q^*}{|v|^2\,\tilde u^*v}\,v
\]
satisfies $P\mathsf X=D_q^*$ and $\tilde u^*\mathsf X=0$, hence $(P+\tau_{\mathrm{gf}}^{-1}\tilde u\tilde u^*)\mathsf X=D_q^*$. Therefore the trace equals $D_q\mathsf X=|v|^{-2}D_qD_q^*=|v|^{-2}\bigl(|v_\mu|^2+|v_\nu|^2\bigr)=\varpi_q(k+\upsilon)$.

\emph{Step 5: proof of (c).} We expand in powers of $\vartheta$ and of $\upsilon$; the dots stand for terms of degree at least $3$:
\[
j(\vartheta)=1+\frac{i\vartheta}2-\frac{\vartheta^2}6+\dots,\qquad j(\vartheta)^{-1}=1-\frac{i\vartheta}2-\frac{\vartheta^2}{12}+\dots,\qquad|j(\vartheta)|^2=1-\frac{\vartheta^2}{12}+\dots,
\]
\[
\begin{multlined}
\mathsf k_\mu(k+\upsilon)=\mathsf k_\mu+e^{ik_\mu}\bigl(i\upsilon_\mu-\tfrac12\upsilon_\mu^2+\dots\bigr),\qquad|\mathsf k_\mu(k+\upsilon)|^2=2-2\cos(k_\mu+\upsilon_\mu)=|\mathsf k_\mu|^2 \\
+2\upsilon_\mu\sin k_\mu+\upsilon_\mu^2\cos k_\mu+\dots .
\end{multlined}
\]
In $\hat{\mathsf M}_e=-\sum_\mu\overline{\mathsf k_\mu}\,j(\upsilon_\mu)^{-1}\mathsf k_\mu(k+\upsilon)$ the terms of degree $1$ are
\[
-i\sum_\mu\upsilon_\mu\overline{\mathsf k_\mu}\bigl[e^{ik_\mu}-\tfrac12\mathsf k_\mu\bigr]=-i\sum_\mu\upsilon_\mu\overline{\mathsf k_\mu}\,\tfrac12\bigl(e^{ik_\mu}+1\bigr)=-\sum_\mu\upsilon_\mu\sin k_\mu ,
\]
because $\overline{\mathsf k_\mu}(e^{ik_\mu}+1)=-2i\sin k_\mu$, and the terms of degree $2$ are
\[
-\sum_\mu\upsilon_\mu^2\,\overline{\mathsf k_\mu}\bigl[-\tfrac1{12}\mathsf k_\mu+\tfrac12e^{ik_\mu}-\tfrac12e^{ik_\mu}\bigr]=\frac1{12}\sum_\mu\upsilon_\mu^2|\mathsf k_\mu|^2 .
\]
The term $\tau_{\mathrm{gf}}^{-1}\mathsf k\mathsf k^*$ of $\hat{\mathsf N}_e$ does not depend on $\gamma$. Hence $\operatorname{tr}\hat{\mathsf N}_{e,2}$ is the homogeneous polynomial of degree $2$ in the expansion of
\[
\operatorname{tr}\bigl[J^*P(k+\upsilon)J\bigr]=\sum_\mu|j(\upsilon_\mu)|^2\sum_{\nu\ne\mu}|\mathsf k_\nu(k+\upsilon)|^2,
\]
namely
\[
\sum_\mu\Bigl[-\frac1{12}\upsilon_\mu^2\sum_{\nu\ne\mu}|\mathsf k_\nu|^2+\sum_{\nu\ne\mu}\upsilon_\nu^2\cos k_\nu\Bigr]=\sum_\mu\upsilon_\mu^2\Bigl[2\cos k_\mu-\frac1{12}\bigl(|\mathsf k|^2-|\mathsf k_\mu|^2\bigr)\Bigr].
\]
This is the asserted formula for $\operatorname{tr}\hat{\mathsf N}_{e,2}$, and (c) is proved.
\end{proof}

With these formulas we evaluate the self-energy at zero momentum.

\begin{theorem}[The five contributions to the self-energy at zero momentum]
\label{IV-73-34}
Let $\tau_{\mathrm{gf}}>0$.
\begin{enumerate}
\item[(i)] For $i=1,\dots,4$, all $k\in[-\pi,\pi]^3\setminus\{0\}$ and all $\gamma\in\mathbb R^3$ we have $|\mathfrak s^{(i)}(k;\gamma)|\le C|\gamma|^2|k|^{-2}$, with $C=C(\tau_{\mathrm{gf}})$.
\item[(ii)] $\sum_{i=1}^5c^{(i)}_{\mu\nu}(\tau_{\mathrm{gf}})=0$ for all $\mu,\nu$.
\item[(iii)] At $\tau_{\mathrm{gf}}=1$,
\begin{equation}\label{IV-73-35}
\bigl(c^{(1)}_{\mu\nu},\dots,c^{(5)}_{\mu\nu}\bigr)(1)=\delta_{\mu\nu}\cdot\Bigl(-20Z_0+\tfrac83,\ \ 16Z_0-\tfrac{28}9,\ \ 4Z_0-\tfrac23,\ \ \tfrac49,\ \ \tfrac23\Bigr).
\end{equation}
\end{enumerate}
\end{theorem}

\begin{proof}
\emph{Step 1: bounds.} By the formulas of Lemma~\ref{IV-73-31}, the entries of $J^{\pm1}$ and of $\mathsf k(k+\upsilon)$ are real analytic functions of $(k,\gamma)$. On $[-\pi,\pi]^3$ we have $P(k)=O(|k|^2)$ and $\partial_kP(k)=O(|k|)$. The terms of degree $1$ in $\gamma$ of $J^*P(k+\upsilon)J$ are terms containing the factor $P(k)$ and the term $\sum_\mu\upsilon_\mu\partial_\mu P(k)$; hence $|\hat{\mathsf N}_{e,1}|\le C|\gamma||k|$, while $|\hat{\mathsf N}_{e,2}|\le C|\gamma|^2$. Every term of $\hat{\mathsf M}_e$ contains the factor $\mathsf k(k)^*$, so that $|\hat{\mathsf M}_{e,j}|\le C|\gamma|^j|k|$. Finally $|\hat{\mathsf C}|\le C|k|^{-2}$ by the formula for $\hat{\mathsf C}$ in Definition~\ref{IV-73-29}(2), and $|\hat{\mathsf M}_{e,0}|^{-1}=|\mathsf k(k)|^{-2}\le C|k|^{-2}$, as $|\mathsf k(k)|^2\ge\frac4{\pi^2}|k|^2$.

Inserting these bounds into the definitions of $\mathfrak s^{(1)},\dots,\mathfrak s^{(4)}$ we obtain $|\hat{\mathsf C}\hat{\mathsf N}_{e,1}|^2\le C|\gamma|^2|k|^{-2}$, $|\hat{\mathsf C}\hat{\mathsf N}_{e,2}|\le C|\gamma|^2|k|^{-2}$, $|\hat{\mathsf M}_{e,1}/\hat{\mathsf M}_{e,0}|^2\le C|\gamma|^2|k|^{-2}$ and $|\hat{\mathsf M}_{e,2}/\hat{\mathsf M}_{e,0}|\le C|\gamma|^2|k|^{-1}\le C'|\gamma|^2|k|^{-2}$, which proves (i).

\emph{Step 2: identity of the integrands.} Fix $k\in[-\pi,\pi]^3\setminus\{0\}$, and let $\gamma$ be so close to $0$ that $\mathsf k(k+\upsilon)\ne0$ and $\hat{\mathsf M}_e(k;\gamma)\ne0$ for all $e$. We divide the first identity in \eqref{IV-73-32} by its value at $\gamma=0$, take one half of its logarithm, sum over $e$ and use the second identity in \eqref{IV-73-32}:
\[
\sum_e\Bigl\{\tfrac12\log\det\bigl[\hat{\mathsf C}\hat{\mathsf N}_e(k;\gamma)\bigr]-\log\bigl|\hat{\mathsf M}_e(k;\gamma)/\hat{\mathsf M}_{e,0}\bigr|\Bigr\}-\sum_\mu\log\sigma_1(\gamma_\mu^2)=\tfrac12\sum_e\log\frac{|\mathsf k(k+2e\gamma)|^2}{|\mathsf k(k)|^2}.
\]
Both sides are real analytic in $\gamma$ near $0$, and we compare their terms of degree $2$. On the left, the expansion $\log\det(\mathbf 1+\mathsf A)=\operatorname{tr}\mathsf A-\frac12\operatorname{tr}\mathsf A^2+\dots$ with $\mathsf A=\hat{\mathsf C}(\hat{\mathsf N}_{e,1}+\hat{\mathsf N}_{e,2}+\dots)$ yields $\mathfrak s^{(2)}+\mathfrak s^{(1)}$; the expansion $\log|1+z|=\operatorname{Re}(z-\frac12z^2+\dots)$ with $z=(\hat{\mathsf M}_{e,1}+\hat{\mathsf M}_{e,2}+\dots)/\hat{\mathsf M}_{e,0}$ yields, with the sign in front, $\mathfrak s^{(4)}+\mathfrak s^{(3)}$; and by Lemma~\ref{IV-73-9}(c), applied to $Y=\gamma_\mu\mathsf e_3$, the term $-\sum_\mu\log\sigma_1(\gamma_\mu^2)$ yields $\frac13|\gamma|^2=\mathfrak s^{(5)}$. On the right, with $\mathsf l(k):=\log|\mathsf k(k)|^2$, the terms of degree $2$ are $\frac12\sum_{e=\pm1}\frac12\sum_{\mu\nu}4\gamma_\mu\gamma_\nu\,\partial_\mu\partial_\nu\mathsf l(k)$. Hence
\[
\sum_{i=1}^5\mathfrak s^{(i)}(k;\gamma)=2\sum_{\mu\nu}\gamma_\mu\gamma_\nu\,\partial_\mu\partial_\nu\mathsf l(k)\qquad(k\ne0),
\]
an identity between homogeneous polynomials of degree $2$ in $\gamma$, valid for all $\gamma\in\mathbb R^3$.

\emph{Step 3: flux argument.} The function $\mathsf l$ is smooth on $\mathbb R^3\setminus2\pi\mathbb Z^3$ and periodic, and on the cube $|\partial\mathsf l|\le C|k|^{-1}$ and $|\partial\partial\mathsf l|\le C|k|^{-2}$. For $0<\varepsilon<1$ we apply Gauss' divergence theorem on the cube $[-\pi,\pi]^3$ less the ball $\{|k|\le\varepsilon\}$; the contributions of opposite faces of the cube cancel by periodicity, and
\[
\int_{[-\pi,\pi]^3\cap\{|k|>\varepsilon\}}\partial_\mu\partial_\nu\mathsf l\,dk=-\int_{|k|=\varepsilon}\frac{k_\mu}\varepsilon\,\partial_\nu\mathsf l\,dS=O(\varepsilon).
\]
Letting $\varepsilon\to0$ we get $\int_{\mathbb T}\partial_\mu\partial_\nu\mathsf l\,d^\circ k=0$. By Step~2 and the definition of the numbers $c^{(i)}_{\mu\nu}$, $\frac12\sum_{\mu\nu}\bigl(\sum_ic^{(i)}_{\mu\nu}\bigr)\gamma_\mu\gamma_\nu=0$ for all $\gamma$; as the matrices are symmetric, (ii) follows.

\emph{Step 4: elementary integrals.} We write $\mathsf k=\mathsf k(k)$. By the symmetry under permutations of the coordinates, $\int_{\mathbb T}|\mathsf k_\mu|^2|\mathsf k|^{-2}d^\circ k=\frac13$. As $\cos k_\mu=1-\frac12|\mathsf k_\mu|^2$, we get $\int_{\mathbb T}\cos k_\mu|\mathsf k|^{-2}d^\circ k=Z_0-\frac16$. Since $\partial_\mu|\mathsf k|^2=2\sin k_\mu$,
\[
\partial_\mu\bigl(\sin k_\mu|\mathsf k|^{-2}\bigr)=\cos k_\mu|\mathsf k|^{-2}-2\sin^2k_\mu|\mathsf k|^{-4};
\]
the function $\sin k_\mu|\mathsf k|^{-2}$ is periodic and $O(|k|^{-1})$, so the flux argument of Step~3 shows that the integral of the left side over $\mathbb T$ vanishes, and $\int_{\mathbb T}\sin^2k_\mu|\mathsf k|^{-4}d^\circ k=\frac12(Z_0-\frac16)$. Finally $\int_{\mathbb T}\sin k_\mu\sin k_\nu|\mathsf k|^{-4}d^\circ k=0$ for $\mu\ne\nu$, the integrand being odd under $k_\mu\to-k_\mu$.

\emph{Step 5: values.} We use Lemma~\ref{IV-73-31}(c) with $\upsilon=2e\gamma$, and $\hat{\mathsf M}_{e,0}=-|\mathsf k|^2$. For $e=0$ we have $\upsilon=0$, so that $\hat{\mathsf N}_{0,j}=0$ and $\hat{\mathsf M}_{0,j}=0$ for $j\ge1$; the charge $0$ contributes nothing, and for $e=\pm1$ we have $\upsilon_\mu^2=4\gamma_\mu^2$. We write $\delta$ for the matrix $(\delta_{\mu\nu})$.
\begin{itemize}
\item $\mathfrak s^{(5)}=\frac13|\gamma|^2$, whose integral is $\frac13|\gamma|^2$; hence $c^{(5)}=\frac23\delta$.
\item $\mathfrak s^{(4)}=\sum_{e=\pm1}\frac4{12}\sum_\mu\gamma_\mu^2|\mathsf k_\mu|^2|\mathsf k|^{-2}$, whose integral is $\frac29|\gamma|^2$ by Step~4; hence $c^{(4)}=\frac49\delta$.
\item $\mathfrak s^{(3)}=4\bigl(\sum_\mu\gamma_\mu\sin k_\mu\bigr)^2|\mathsf k|^{-4}$, whose integral is $4|\gamma|^2\cdot\frac12(Z_0-\frac16)$; hence $c^{(3)}=(4Z_0-\frac23)\delta$.
\item At $\tau_{\mathrm{gf}}=1$ we have $\hat{\mathsf C}=|\mathsf k|^{-2}\mathbf 1$ and
\[
\mathfrak s^{(2)}=\tfrac12\sum_{e=\pm1}|\mathsf k|^{-2}\operatorname{tr}\hat{\mathsf N}_{e,2}=4\sum_\mu\gamma_\mu^2\Bigl[2\cos k_\mu-\frac1{12}\bigl(|\mathsf k|^2-|\mathsf k_\mu|^2\bigr)\Bigr]|\mathsf k|^{-2},
\]
whose integral is $4|\gamma|^2\bigl[2Z_0-\frac13-\frac1{12}\cdot\frac23\bigr]=|\gamma|^2\bigl(8Z_0-\frac{14}9\bigr)$; hence $c^{(2)}=(16Z_0-\frac{28}9)\delta$.
\item By (ii), $c^{(1)}=-\sum_{i\ge2}c^{(i)}=-\bigl(20Z_0-\frac{28}9-\frac23+\frac49+\frac23\bigr)\delta=\bigl(-20Z_0+\frac83\bigr)\delta$.
\end{itemize}
This proves (iii).
\end{proof}

The same method applies to the vertex function of a plaquette action.

\begin{theorem}[The vertex function at zero momenta]
\label{IV-73-36}
Let $\tau_{\mathrm{gf}}>0$ and let $q$ be a plaquette. Then $|\mathfrak v_q^{(j)}(k;\gamma)|\le C|\gamma|^2|k|^{-2}$ for $j=1,\dots,4$, all $k\in[-\pi,\pi]^3\setminus\{0\}$ and all $\gamma\in\mathbb R^3$, where $C$ depends on $\tau_{\mathrm{gf}}$ only; and
\[
\sum_{j=1}^4\mathsf g^{(j)}_{q;\mu\nu}(\tau_{\mathrm{gf}})=0\qquad\text{for all }\mu,\nu .
\]
\end{theorem}

\begin{proof}
\emph{Step 1: bound.} The $1\times3$ matrix $D_q(k)$ is smooth in $k$, satisfies $D_q(k)=O(|k|)$, and depends on $q$ only through the two directions of $q$. In the expansion of $\hat{\mathsf h}_{q,e}=J^*D_q(k+\upsilon)^*D_q(k+\upsilon)J$ in powers of $\gamma$, each term of degree $j\le1$ retains at least $2-j$ factors $D_q(k)$ or $D_q(k)^*$; hence $|\hat{\mathsf h}_{q,e,j}|\le C|\gamma|^j|k|^{2-j}$ for $j=0,1,2$, with $C$ independent of $q$. As in Step~1 of the proof of Theorem~\ref{IV-73-34}, Lemma~\ref{IV-73-31} gives $|\hat{\mathsf N}_{e,1}|\le C|\gamma||k|$ and $|\hat{\mathsf N}_{e,2}|\le C|\gamma|^2$, and Definition~\ref{IV-73-15}(1) gives $|\hat{\mathsf C}|\le C|k|^{-2}$. Inserted into the definitions of $\mathfrak v_q^{(1)},\dots,\mathfrak v_q^{(4)}$, these bounds give the asserted bound.

\emph{Step 2: identity of the integrands.} Fix $k\in[-\pi,\pi]^3\setminus\{0\}$. For small $\gamma$ we write $\hat{\mathsf N}_e=\hat{\mathsf N}_{e,0}\bigl(\mathbf 1+\hat{\mathsf C}(\hat{\mathsf N}_e-\hat{\mathsf N}_{e,0})\bigr)$, and the Neumann series gives
\[
\hat{\mathsf N}_e^{-1}=\hat{\mathsf C}-\hat{\mathsf C}\hat{\mathsf N}_{e,1}\hat{\mathsf C}-\hat{\mathsf C}\hat{\mathsf N}_{e,2}\hat{\mathsf C}+\hat{\mathsf C}\hat{\mathsf N}_{e,1}\hat{\mathsf C}\hat{\mathsf N}_{e,1}\hat{\mathsf C}+O(|\gamma|^3).
\]
Hence the terms of degree $2$ in $\gamma$ of $\frac12\sum_e\operatorname{tr}\bigl[\hat{\mathsf h}_{q,e}\hat{\mathsf N}_e^{-1}\bigr]$ are $\sum_{j=1}^4\mathfrak v_q^{(j)}(k;\gamma)$. On the other hand, for small $\gamma$ the hypotheses of Lemma~\ref{IV-73-31}(b) hold, and by \eqref{IV-73-33} the same expression equals $\frac12\sum_e\varpi_q(k+2e\gamma)$, whose terms of degree $2$ are $2\sum_{\mu\nu}\gamma_\mu\gamma_\nu\partial_\mu\partial_\nu\varpi_q(k)$. Therefore
\[
\sum_{j=1}^4\mathfrak v_q^{(j)}(k;\gamma)=2\sum_{\mu\nu}\gamma_\mu\gamma_\nu\,\partial_\mu\partial_\nu\varpi_q(k)\qquad(k\ne0,\ \gamma\in\mathbb R^3).
\]

\emph{Step 3: flux argument.} The function $\varpi_q$ is smooth on $\mathbb R^3\setminus2\pi\mathbb Z^3$ and periodic, and on the cube $|\partial\varpi_q|\le C|k|^{-1}$ and $|\partial\partial\varpi_q|\le C|k|^{-2}$, since it is a quotient of smooth functions vanishing to second order at $0$ whose denominator is at least $\frac4{\pi^2}|k|^2$. As in Step~3 of the proof of Theorem~\ref{IV-73-34},
\[
\int_{[-\pi,\pi]^3\cap\{|k|>\varepsilon\}}\partial_\mu\partial_\nu\varpi_q\,dk=-\int_{|k|=\varepsilon}\frac{k_\mu}\varepsilon\,\partial_\nu\varpi_q\,dS=O(\varepsilon),
\]
hence $\int_{\mathbb T}\partial_\mu\partial_\nu\varpi_q\,d^\circ k=0$. By Step~2, $\frac12\sum_{\mu\nu}\bigl(\sum_j\mathsf g^{(j)}_{q;\mu\nu}\bigr)\gamma_\mu\gamma_\nu=0$ for all $\gamma$, and the symmetry of the matrices gives the claim.
\end{proof}

We next record the behaviour of these objects under point reflections.

\begin{lemma}[Point reflections]
\label{IV-73-37}
Let $c\in\frac12\mathbb Z^3$, and let $R_c:x\mapsto2c-x$ on $\mathbb Z^3$ (note that $2c\in\mathbb Z^3$). We let $R_c$ act on bonds by $R_c\langle x,x'\rangle:=\langle R_cx,R_cx'\rangle$, on plaquettes through their corners (the corners of $R_cp$ are the images of the corners of $p$), and on configurations $U$, $1$-forms $A$ and site functions $\lambda$ by
\[
(R_cU)(b):=U(R_cb),\qquad(R_cA)(b):=A(R_cb),\qquad(R_c\lambda)(x):=\lambda(R_cx);
\]
thus $(R_cA)_\mu(x)=-A_\mu(2c-x-e_\mu)$. Since $2c\in\mathbb Z^3$, $R_c$ preserves the lattice. The following hold.
\begin{enumerate}
\item[(a)] For all plaquettes $p$, configurations $U$, $1$-forms $A$ on $\mathbb Z^3$ and $m\ge0$: $a_p(R_cU)=a_{R_cp}(U)$ and $\mathsf a_p^{(m)}(R_cA)=\mathsf a^{(m)}_{R_cp}(A)$. Moreover $\mathrm d^*R_c=R_c\mathrm d^*$ on $1$-forms, and $\mathrm dR_c=R_c\mathrm d$ on site functions.
\item[(b)] For every $\tau_{\mathrm{gf}}>0$ the family $\bigl((R_cA)^a(b)\bigr)_{b,a}$ has the law of the Gaussian family of $\mathbb Z^3$ and $\tau_{\mathrm{gf}}$. Moreover $\mathsf D(R_cb;R_cb')=\mathsf D(b;b')$ for all bonds $b,b'$, where $\mathsf D$ is taken with oriented arguments, that is, it is given by the formula of Definition~\ref{IV-73-15}(2) for arbitrary bonds; in particular $\mathsf D(\bar b;b')=-\mathsf D(b;b')$ and $\mathsf D(b;\bar b')=\mathsf D(b;b')$.
\item[(c)] For a positively oriented bond $b$ let $b^R$ be the positively oriented one of the bonds $R_cb$ and $\overline{R_cb}$, that is, $(x,\mu)^R=(2c-x-e_\mu,\mu)$. Then for all positively oriented bonds $b,b'$ and all plaquettes $q$
\[
\Sigma_{\mathrm{tot}}(b^R,b'^R)=\Sigma_{\mathrm{tot}}(b,b'),\qquad\Gamma_{R_cq}(b^R,b'^R)=\Gamma_q(b,b').
\]
\end{enumerate}
\end{lemma}

\begin{proof}
The actions in the statement are defined by the same involution on sites; in particular $R_c^2$ is the identity on sites, bonds, plaquettes and configurations.

\emph{Step 1: proof of (a).} Let $p=p_{\mu\nu}(x)$. The loop $\partial p=(b_1,b_2,\bar b_3,\bar b_4)$ passes through $x$, $x+e_\mu$, $x+e_\mu+e_\nu$, $x+e_\nu$. The map $R_c$ carries it to the loop through $y+e_\mu+e_\nu$, $y+e_\nu$, $y$, $y+e_\mu$, where $y:=2c-x-e_\mu-e_\nu$ is the base point of $R_cp=p_{\mu\nu}(y)$; this image loop is $\partial(R_cp)$, traversed in the same sense but started at the opposite corner. Since $\partial(R_cU)(p)$ is the product of the values of $U$ along the image loop, it is conjugate in $G$ to $\partial U(R_cp)$, and the two matrices have the same trace; hence $a_p(R_cU)=a_{R_cp}(U)$. As $R_c(e^{itA})=e^{itR_cA}$, comparison of the Taylor coefficients in \eqref{IV-73-7} gives $\mathsf a^{(m)}_p(R_cA)=\mathsf a^{(m)}_{R_cp}(A)$.

Next, $(\mathrm d^*A)(x)=-\sum_{b:\,b_-=x}A(b)$, the sum over the six bonds starting at $x$, and $R_c$ maps the bonds starting at $x$ bijectively onto those starting at $R_cx$; therefore $\mathrm d^*R_cA=R_c\mathrm d^*A$. Finally $(\mathrm dR_c\lambda)(b)=\lambda(R_cb_+)-\lambda(R_cb_-)=(\mathrm d\lambda)(R_cb)$.

\emph{Step 2: proof of (b).} The variables $(R_cA)_\mu(x)=-A_\mu(2c-x-e_\mu)$ form a centred Gaussian family, with the same colour structure as $A$. By Definition~\ref{IV-73-15}(1) the covariance of $(R_cA)_\mu(x)$ and $(R_cA)_\nu(y)$, for equal colours, is
\[
\mathsf C\bigl((2c-x-e_\mu,\mu),(2c-y-e_\nu,\nu)\bigr)=\int_{\mathbb T}e^{-ik\cdot(x-y)}\,e^{-i(k_\mu-k_\nu)}\,\hat{\mathsf C}_{\mu\nu}(k)\,d^\circ k .
\]
Since $e^{-ik_\mu}\mathsf k_\mu(k)=-\mathsf k_\mu(-k)$, and correspondingly $e^{ik_\nu}\overline{\mathsf k_\nu(k)}=-\overline{\mathsf k_\nu(-k)}$, we have $e^{-i(k_\mu-k_\nu)}\hat{\mathsf C}_{\mu\nu}(k)=\hat{\mathsf C}_{\mu\nu}(-k)$, also for $\mu=\nu$. The substitution $k\to-k$ shows that the covariance equals $\mathsf C((x,\mu),(y,\nu))$, and a centred Gaussian family is determined in law by its covariance. Moreover $\mathsf G(R_cx,R_cy)=\mathsf G(x,y)$, again by the substitution $k\to-k$ in the formula for $\mathsf G^{\mathrm{lat}}$; with $(R_cb)_\pm=R_c(b_\pm)$ the formula for $\mathsf D$ in Definition~\ref{IV-73-15}(2) gives $\mathsf D(R_cb;R_cb')=\mathsf D(b;b')$. The same formula shows that reversing the first argument changes the sign of $\mathsf D$, while reversing the second does not change it.

\emph{Step 3: proof of (c).} By (a) and (b), $\mathsf S_3$ and $\mathsf V_2$ are invariant under $A\to R_cA$, term by term: the term of a plaquette $p$ is carried to that of $R_cp$, and the terms of a bond $b$ and of a pair $(b,b')$ in $\mathsf V_2$ to those of $b^R$ and of $(b^R,b'^R)$; here $\mathsf g(b)=\frac13$ is constant on $\mathbb Z^3$, $|A(R_cb)|=|A(b^R)|$, and the signs coming from reversed arguments of $\mathsf D$ cancel in pairs. Likewise $\mathsf a_q^{(m)}(A)=\mathsf a^{(m)}_{R_cq}(R_cA)$. Since $(R_cA)(b^R)=A(\bar b)=-A(b)$, the chain rule gives
\[
\partial_{b,a}(\mathsf F\circ R_c)=-(\partial_{b^R,a}\mathsf F)\circ R_c
\]
for every polynomial $\mathsf F$. Each term of the kernels in Definition~\ref{IV-73-29}(5) contains exactly two derivatives, so that the signs cancel, and by (b) the means and cumulants do not change under the substitution $A\to R_cA$. Hence $\Sigma_{\mathrm{tot}}(b,a;b',a')=\Sigma_{\mathrm{tot}}(b^R,a;b'^R,a')$ and $\Gamma_q(b,a;b',a')=\Gamma_{R_cq}(b^R,a;b'^R,a')$, which is (c).
\end{proof}

\subsection{Tools on kernels}

We now develop a calculus of kernels on $\mathbb Z^3$ given by symbols on $\mathbb T$; it serves to determine the behaviour of lattice graphs at large distances. Its basic objects are the lines of Definition~\ref{IV-73-38}(4): kernels whose symbols near $0$ and whose values at large distances are given, up to one order and a logarithm, by homogeneous functions. The lemmas after the definition relate these two homogeneous functions to each other; they compare lattice sums with integrals, and a kernel tested against a slowly varying function with its leading part; they expand one-loop integrals of line symbols at small external momenta; and they show that, under conditions on the degrees, convolutions of lines and, up to local terms, products of two lines are again lines.

\begin{definition}[Symbols, lines, slowly varying functions, majorants]
\label{IV-73-38}
\emph{(1) Two cut-offs.} Let $\theta_0(t)=e^{-1/t}$ for $t>0$ and $\theta_0(t)=0$ for $t\le0$. This function is smooth. To check the only issue, at zero, induction gives $\theta_0^{(n)}(t)=P_n(1/t)e^{-1/t}$ for $t>0$, where $P_0=1$ and $P_{n+1}(u)=u^2(P_n(u)-P_n'(u))$. Every polynomial times $e^{-u}$ tends to zero as $u\to\infty$, as does $uP_n(u)e^{-u}$. Thus each derivative extends continuously by zero and has derivative zero at the origin. For $s\ge0$ put
\[
\theta(s):=\frac{\theta_0(1-s)}{\theta_0(1-s)+\theta_0(s-\frac12)} ;
\]
the denominator is positive, since one of the numbers $1-s$, $s-\frac12$ is positive. Hence $\theta\in C^\infty([0,\infty))$, $0\le\theta\le1$, $\theta=1$ on $[0,\frac12]$ and $\theta=0$ on $[1,\infty)$. For $\kappa\in\mathbb R^D$, $D\in\{3,6\}$, let $\chi(\kappa):=\theta(|\kappa|/2)$. Then $\chi$ is smooth, $\chi=1$ for $|\kappa|\le1$ and $\chi=0$ for $|\kappa|\ge2$; since $2<\pi$, we also regard $\chi$ as a function on $\mathbb T^{D/3}$. The functions $\theta$ and $\chi$ are fixed by formulas, so constants which depend on $\theta$ and $\chi$ alone are pure constants.

\smallskip\noindent\emph{(2) Symbol classes.} For real $s$ let
\[
\mathcal S^s:=\Bigl\{\sigma\in C^\infty(\mathbb T\setminus\{0\}):\ \sup_{k\in[-\pi,\pi]^3\setminus\{0\}}|k|^{|\alpha|-s}|\partial^\alpha\sigma(k)|<\infty\ \text{for every multi-index }\alpha\Bigr\}.
\]
By Leibniz' rule $\mathcal S^s\mathcal S^{s'}\subset\mathcal S^{s+s'}$, and $\mathcal S^{s+1}\subset\mathcal S^s$ because $|k|$ is bounded on the cube.

\smallskip\noindent\emph{(3) Line symbols.} Let $s\ge-2$ be an integer. A \emph{line symbol of degree $s$} is a $\sigma\in\mathcal S^s$ for which there is a \emph{leading part} $\sigma^h\in C^\infty(\mathbb R^3\setminus\{0\})$ with
\[
\sigma^h(tk)=t^s\sigma^h(k)\quad(t>0),\qquad\sigma^h(-k)=(-1)^s\sigma^h(k),\qquad\sigma-\chi\sigma^h\in\mathcal S^{s+1};
\]
here $\chi\sigma^h$ is regarded as a function on $\mathbb T\setminus\{0\}$, which is possible since $\chi$ vanishes near the boundary of the cube. The kernel $K$ with the symbol $\sigma$ is then called an \emph{elementary line of degree $n=s+3$}. Sums of line symbols of the same degree are line symbols, and their leading parts add; the product of two line symbols whose degrees have a sum $\ge-2$ is a line symbol, and the leading parts multiply. For products this follows from
\[
\sigma_1\sigma_2-\chi\sigma_1^h\sigma_2^h=(\sigma_1-\chi\sigma_1^h)\sigma_2+\chi\sigma_1^h(\sigma_2-\chi\sigma_2^h)+(\chi^2-\chi)\sigma_1^h\sigma_2^h ,
\]
whose last term vanishes for $|k|\le1$ and for $|k|\ge2$. We list some examples.
\begin{itemize}
\item $e^{ik\cdot y}$, $y\in\mathbb Z^3$: degree $0$, leading part $1$.
\item $\mathsf k_\mu$ and $\overline{\mathsf k_\mu}$: degree $1$, leading parts $ik_\mu$ and $-ik_\mu$.
\item $|\mathsf k|^{-2}$: degree $-2$, leading part $|k|^{-2}$. Indeed, on the cube less $0$ we write $|\mathsf k|^2=|k|^2(1+\mathsf w)$ with $\mathsf w:=|\mathsf k|^2/|k|^2-1$; the numerator of $\mathsf w$, $|\mathsf k|^2-|k|^2$, vanishes to fourth order at $0$, so that $\mathsf w$ is smooth with the bounds of $\mathcal S^2$, and $1+\mathsf w\ge\frac4{\pi^2}$; then
\[
|\mathsf k|^{-2}-\chi|k|^{-2}=\chi|k|^{-2}\bigl((1+\mathsf w)^{-1}-1\bigr)+(1-\chi)|\mathsf k|^{-2}\in\mathcal S^0 .
\]
\item The periodic continuation of $\rho\sigma^h$, where $\sigma^h$ is as above and $\rho\in C_c^\infty(\mathbb R^3)$ satisfies $\operatorname{supp}\rho\subset\{|k|<\pi\}$ and $\rho=1$ near $0$: its leading part is $\sigma^h$.
\end{itemize}

\smallskip\noindent\emph{(4) Lines.} For an integer $n\ge2$ let $\mathfrak K(n)$ be the set of functions $K:\mathbb Z^3\to\mathbb C$ with the following properties.
\begin{itemize}
\item[(K1)] $K$ is the kernel with a symbol $\sigma\in L^1(\mathbb T)$.
\item[(K2)] There are a continuous function $\sigma^h:\mathbb R^3\setminus\{0\}\to\mathbb C$, homogeneous of degree $n-3$, and a number $c$ with $|\sigma(k)-\chi(k)\sigma^h(k)|\le c\,|k|^{n-2}\lg|k|$ on $[-\pi,\pi]^3$.
\item[(K3)] There are a function $K^h\in C^\infty(\mathbb R^3\setminus\{0\})$, homogeneous of degree $-n$, and numbers $c_\gamma$ with $|\nabla^\gamma(K-K^h)(x)|\le c_\gamma|x|^{-n-1-|\gamma|}\lg|x|$ for all multi-indices $\gamma$ and all $x\in\mathbb Z^3$ with $|x|\ge|\gamma|+1$.
\end{itemize}
The function $K$ determines the pair $(\sigma^h,K^h)$, which we call the \emph{leading data} of $K$. Indeed, $\sigma$ is determined almost everywhere by $K$ (uniqueness of Fourier coefficients). If $\sigma^h$ and $\sigma'^h$ both satisfy (K2), then $|\sigma^h-\sigma'^h|(k)\le C|k|^{n-2}\lg|k|$ for $0<|k|\le1$, first almost everywhere and then everywhere by continuity; at the points $tk$ with $t\to0$ this gives $t^{n-3}|\sigma^h(k)-\sigma'^h(k)|\le Ct^{n-2}|k|^{n-2}\lg(t|k|)$, whence $\sigma^h=\sigma'^h$. If $K^h$ and $K'^h$ both satisfy (K3), we take $\gamma=0$ and $x=j\mathsf m$ with $\mathsf m\in\mathbb Z^3\setminus\{0\}$ and $j\to\infty$: from $j^{-n}|K^h(\mathsf m)-K'^h(\mathsf m)|\le Cj^{-n-1}\lg(j|\mathsf m|)$ we get $K^h=K'^h$ on $\mathbb Z^3\setminus\{0\}$, hence on $\mathbb Q^3\setminus\{0\}$ by homogeneity, and everywhere by density and continuity.

The \emph{constants of $K$} are $c$, the numbers $c_\gamma$, $\|\sigma\|_{L^1}$, and the suprema of $|\sigma^h|$ and of the derivatives of $K^h$ on $S^2$.

\smallskip\noindent\emph{(5) Slowly varying functions.} For $D\in\{3,6\}$, $R\ge1$ and $g:\mathbb Z^D\to\mathbb C$ we put
\[
\|g\|_{(R)}:=\max_{|\gamma|\le10}R^{|\gamma|}\sup_\zeta|\nabla^\gamma g(\zeta)| ,
\]
and if $g$ has finite support, $\check g(\kappa):=\sum_\zeta g(\zeta)e^{-i\kappa\cdot\zeta}$. For $G\in C_c^\infty(\mathbb R^D)$ we put $\|G\|_{C^{10}}:=\max_{|\gamma|\le10}\sup|\partial^\gamma G|$.

\smallskip\noindent\emph{(6) Majorants.} For real $s>-D$ let $\mathcal H_s(D)$ be the set of Borel functions $\mathsf H:\mathbb R^D\to[0,\infty]$ which are homogeneous of degree $s$ and integrable over $\{|\kappa|\le1\}$. If $s<9-D$ and $\mathsf H\in\mathcal H_s(D)$, then $\int_{\mathbb R^D}\mathsf H(\kappa)\lg|\kappa|\,\langle\kappa\rangle^{-10}d\kappa<\infty$, as one sees in polar coordinates.
\end{definition}

\begin{lemma}[(K) From a symbol to its kernel]
\label{IV-73-39}
\begin{enumerate}
\item[(a)] Let $s>-3$ be real, $\sigma\in\mathcal S^s$, and let $K$ be the kernel with the symbol $\sigma$. Then for every multi-index $\gamma$ and all $x\in\mathbb Z^3$
\[
|\nabla^\gamma K(x)|\le C_\gamma\langle x\rangle^{-3-s-|\gamma|},
\]
where $C_\gamma$ depends on finitely many of the suprema in Definition~\ref{IV-73-38}(2).
\item[(b)] Let $s>-3$ be real and let $\sigma^h\in C^\infty(\mathbb R^3\setminus\{0\})$ be homogeneous of degree $s$. There is exactly one continuous function $K^h$ on $\mathbb R^3\setminus\{0\}$, homogeneous of degree $-3-s$, such that for every $\rho\in C_c^\infty(\mathbb R^3)$ with $\rho=1$ near $0$ and every $N$
\[
\sup_{|x|\ge1}|x|^N\Bigl|\int_{\mathbb R^3}e^{ik\cdot x}\rho(k)\sigma^h(k)\,d^\circ k-K^h(x)\Bigr|<\infty .
\]
The function $K^h$ is smooth. If $\sigma^h(-k)=\pm\sigma^h(k)$, then $K^h(-x)=\pm K^h(x)$; replacing $\sigma^h$ by $(ik)^\gamma\sigma^h$ replaces $K^h$ by $\partial^\gamma K^h$; and replacing $\sigma^h$ by $\sigma^h\circ O^{-1}$, $O$ a rotation, replaces $K^h$ by $K^h\circ O^{-1}$. We write $K^h=:\mathcal F^{-1}_{\ne0}\sigma^h$.
\item[(c)] Let $K$ be an elementary line of degree $n$ with the symbol $\sigma$, and $K^h:=\mathcal F^{-1}_{\ne0}\sigma^h$. Then
\[
|\nabla^\gamma(K-K^h)(x)|\le C_\gamma|x|^{-n-1-|\gamma|}\qquad(|x|\ge|\gamma|+1),
\]
with $C_\gamma=C_\gamma(\gamma,\sigma)$. In particular, if $n\ge2$, then $K\in\mathfrak K(n)$ with the leading data $(\sigma^h,\mathcal F^{-1}_{\ne0}\sigma^h)$.
\end{enumerate}
\end{lemma}

\begin{proof}
\emph{Step 1: proof of (a).} The kernel $\nabla^\gamma K$ has the symbol $\mathsf k^\gamma\sigma:=\prod_j\mathsf k_j^{\gamma_j}\sigma$, which lies in $\mathcal S^{s+|\gamma|}$ since $\mathsf k_j\in\mathcal S^1$; as $s+|\gamma|>-3$, it suffices to treat $\gamma=0$. Always $|K(x)|\le\|\sigma\|_{L^1}$, which settles $x=0$. Let $|x|\ge1$, and split $\sigma=\sigma'+\sigma''$ with $\sigma':=\chi(|x|k)\sigma$, where $\chi(|x|k)$ vanishes for $|k|\ge2/|x|$ and is regarded as a function on $\mathbb T$. Then
\[
\int_{\mathbb T}|\sigma'|\,d^\circ k\le C\int_{|k|\le2/|x|}|k|^s\,dk\le C|x|^{-3-s}.
\]
The function $\sigma''$ is smooth and periodic, vanishes for $|k|\le1/|x|$, and satisfies $|\partial^\alpha\sigma''(k)|\le C|k|^{s-|\alpha|}$: a derivative falling on $\chi(|x|k)$ produces a factor $|x|$, and $|x|\le2/|k|$ on the support of the derivatives of $\chi(|x|k)$. We choose $j$ with $|x_j|\ge|x|/\sqrt3$ and an integer $N>s+3$. After $N$ integrations by parts in $k_j$, which produce no boundary terms by periodicity,
\[
\Bigl|\int_{\mathbb T}e^{ik\cdot x}\sigma''(k)\,d^\circ k\Bigr|\le|x_j|^{-N}\int_{\mathbb T}|\partial_j^N\sigma''|\,d^\circ k\le C|x|^{-N}\int_{|k|\ge1/|x|}|k|^{s-N}\,dk\le C|x|^{-3-s}.
\]
These bounds involve only the finitely many suprema in Definition~\ref{IV-73-38}(2) with $|\alpha|\le N$.

\emph{Step 2: uniqueness in (b).} If $K^h$ and $K'^h$ both have the property in (b), then, with one and the same $\rho$, their difference is homogeneous of degree $-3-s$ and $O(|x|^{-N})$ on $\{|x|\ge1\}$ for every $N$; taking $N>3+s$ and letting $x=t\omega$, $|\omega|=1$, $t\to\infty$, we see that it vanishes.

\emph{Step 3: existence in (b).} Let $\rho$ be as in (b), $x\ne0$, and let $N$ be an integer with $N>\frac12(s+3)$. We put
\[
K^h_\rho(x):=\int e^{ik\cdot x}\rho\sigma^h\,d^\circ k+\mathsf E_\rho(x),\qquad\mathsf E_\rho(x):=|x|^{-2N}\int e^{ik\cdot x}(-\Delta_k)^N\bigl[(1-\rho)\sigma^h\bigr]\,d^\circ k ,
\]
with integrals over $\mathbb R^3$. The function $(1-\rho)\sigma^h$ is smooth, and its derivatives of order $j$ are $O(|k|^{s-j})$; as $2N>s+3$, the second integral converges absolutely. Since $|x|^2e^{ik\cdot x}=-\Delta_ke^{ik\cdot x}$, two integrations by parts over balls (Gauss' theorem) lead from $N$ to $N+1$; the boundary terms on the sphere of radius $\mathsf r$ are $O\bigl((1+|x|)\mathsf r^2\cdot\mathsf r^{s-2N}\bigr)$, which tends to $0$ as $\mathsf r\to\infty$ because $2N>s+3$. Hence $\mathsf E_\rho$ does not depend on $N$, so that $|\mathsf E_\rho(x)|\le C_N|x|^{-2N}$ for every $N$. For the same reason $\mathsf E_\rho$ is smooth on $\mathbb R^3\setminus\{0\}$: given an order $j$, we take $N$ with $2N>s+3+j$; since $(-\Delta_k)^N[(1-\rho)\sigma^h]=O(|k|^{s-2N})$ and each derivative of $e^{ik\cdot x}$ in $x$ brings a factor of modulus at most $|k|$, the derivatives of order $\le j$ in $x$ may be taken under the integral sign, with an integrable majorant independent of $x$, and $|x|^{-2N}$ is smooth on $\mathbb R^3\setminus\{0\}$.

If $\rho'$ is a second function as in (b), then $(\rho-\rho')\sigma^h\in C_c^\infty(\mathbb R^3\setminus\{0\})$, and integration by parts gives $\int e^{ik\cdot x}(-\Delta_k)^N[(\rho-\rho')\sigma^h]\,d^\circ k=|x|^{2N}\int e^{ik\cdot x}(\rho-\rho')\sigma^h\,d^\circ k$; therefore $K^h_\rho=K^h_{\rho'}=:K^h$. This function is smooth on $\mathbb R^3\setminus\{0\}$, and it has the property in (b), since $\int e^{ik\cdot x}\rho\sigma^h\,d^\circ k-K^h(x)=-\mathsf E_\rho(x)$. The substitution $k\to k/t$, $t>0$, gives $K^h_\rho(tx)=t^{-3-s}K^h_{\rho(\cdot/t)}(x)$, so that $K^h$ is homogeneous of degree $-3-s$. The parity follows from the substitution $k\to-k$ with an even $\rho$, the behaviour under rotations from $k\to Ok$ with a radial $\rho$.

Finally, $\mathsf E_\rho(x)=\lim_{t\to\infty}\int e^{ik\cdot x}(1-\rho)\sigma^h\rho(k/t)\,d^\circ k$, together with all derivatives in $x$ and locally uniformly in $x$, as one sees by integrating by parts as above; so we may differentiate under the integrals, and this gives the rule for $(ik)^\gamma\sigma^h$, a smooth function homogeneous of degree $s+|\gamma|>-3$.

\emph{Step 4: proof of (c).} Let $s:=n-3$. By (b) with $\rho=\chi$ and by (a) applied to $\sigma-\chi\sigma^h\in\mathcal S^{s+1}$,
\[
K(x)-K^h(x)=\int_{\mathbb T}e^{ik\cdot x}(\sigma-\chi\sigma^h)\,d^\circ k+\Bigl[\int_{\mathbb R^3}e^{ik\cdot x}\chi\sigma^h\,d^\circ k-K^h(x)\Bigr]=O(|x|^{-4-s})\qquad(x\in\mathbb Z^3\setminus\{0\}).
\]
The same argument applies to the line symbol $\mathsf k^\gamma\sigma$ of degree $s+|\gamma|$: its kernel is $\nabla^\gamma K$, its leading part is $(ik)^\gamma\sigma^h$, and $\mathcal F^{-1}_{\ne0}[(ik)^\gamma\sigma^h]=\partial^\gamma K^h$ by (b); hence $|\nabla^\gamma K-\partial^\gamma K^h|\le C|x|^{-4-s-|\gamma|}$. By Taylor's formula and the homogeneity of the derivatives of $K^h$, $|\nabla^\gamma K^h(x)-\partial^\gamma K^h(x)|\le C|x|^{-4-s-|\gamma|}$ for $|x|\ge|\gamma|+1$, because the points at which $K^h$ and its derivatives are evaluated have modulus at least $|x|/(|\gamma|+1)$. As $-4-s=-n-1$, this is the asserted bound. For $n\ge2$, (K1) holds as $\sigma\in\mathcal S^{n-3}\subset L^1(\mathbb T)$, (K2) holds as $|\sigma-\chi\sigma^h|\le C|k|^{n-2}$ on the cube, and (K3) holds since $\lg\ge1$.
\end{proof}

\begin{lemma}[(S) Sums]
\label{IV-73-40}
Let $D\in\{3,6\}$ and $R\ge1$.
\begin{enumerate}
\item[(a)] For real $a>D$, $\sum_{\zeta\in\mathbb Z^D}(R+|\zeta|)^{-a}\le C_aR^{D-a}$. If $y,\zeta\in\mathbb R^D$ and $|y|\le R$, then $R+|\zeta-y|\ge\frac12(R+|\zeta|)$.
\item[(b)] If $F\in C^1(\mathbb R^D)$ and $\sum_\zeta\sup_{\zeta+[0,1]^D}(|F|+|\nabla F|)<\infty$, then
\[
\Bigl|\sum_{\zeta\in\mathbb Z^D}F(\zeta)-\int_{\mathbb R^D}F\Bigr|\le\sqrt D\sum_\zeta\sup_{\zeta+[0,1]^D}|\nabla F| .
\]
\item[(c)] For real $a,a'\ge0$ with $a+a'>3$ and for $x\in\mathbb Z^3$,
\[
\sum_{z\in\mathbb Z^3}\langle z\rangle^{-a}\langle x-z\rangle^{-a'}\le C\langle x\rangle^{-\min\{a,a',a+a'-3\}}\lg\langle x\rangle,\qquad C=C(a,a').
\]
\end{enumerate}
\end{lemma}

\begin{proof}
(a) For $\zeta'\in\zeta+[0,1]^D$ we have $R+|\zeta'|\le R+|\zeta|+\sqrt D\le(1+\sqrt D)(R+|\zeta|)$, as $R\ge1$; so the factor $R+|\zeta|$ changes at most by the factor $1+\sqrt D$ over a cube of side $1$, and the sum is at most
\[
(1+\sqrt D)^a\int_{\mathbb R^D}(R+|\zeta|)^{-a}\,d\zeta=(1+\sqrt D)^aR^{D-a}\int_{\mathbb R^D}(1+|\zeta|)^{-a}\,d\zeta ,
\]
where the last integral is finite since $a>D$. If $|y|\le R$, then $R+|\zeta-y|\ge R$ and $R+|\zeta-y|\ge R+|\zeta|-|y|\ge|\zeta|$, so that $R+|\zeta-y|\ge\max\{R,|\zeta|\}\ge\frac12(R+|\zeta|)$.

(b) The cube $\zeta+[0,1]^D$ has volume $1$ and diameter $\sqrt D$; by the mean value theorem $\bigl|F(\zeta)-\int_{\zeta+[0,1]^D}F\bigr|\le\sqrt D\sup_{\zeta+[0,1]^D}|\nabla F|$. Summing over $\zeta$, with absolutely convergent series, we obtain the claim.

(c) Let first $|x|<2$. Then $|z|\le|x-z|+|x|<|x-z|+2\le3\langle x-z\rangle$, so that $\langle x-z\rangle\ge\frac13\langle z\rangle$ for all $z$, and the sum is at most $3^{a'}\sum_z\langle z\rangle^{-a-a'}$, which is finite as $a+a'>3$. On the other hand, the exponent $m:=\min\{a,a',a+a'-3\}$ is not negative, and $1\le\langle x\rangle<2$, $\lg\langle x\rangle\ge1$; so the right side of the claim is at least $2^{-m}C$, and the claim holds for $|x|<2$ once $C\ge2^m3^{a'}\sum_z\langle z\rangle^{-a-a'}$. Let now $|x|\ge2$. On $\{|z|\le|x|/2\}$ we have $\langle x-z\rangle\ge|x|/2$ and $\sum_{|z|\le|x|/2}\langle z\rangle^{-a}\le C\max\{1,|x|^{3-a}\}\lg|x|$; the sum over this set is therefore at most $C|x|^{-a'}\max\{1,|x|^{3-a}\}\lg|x|\le C|x|^{-\min\{a',a+a'-3\}}\lg|x|$.

The sum over $\{|x-z|\le|x|/2\}$ is treated in the same way, with $a$ and $a'$ exchanged. On the remaining set we have $|z|\ge|x|/2$ and $|x-z|\ge|x|/2$, hence $|z|\le|x|+|x-z|\le3|x-z|$ and $\langle x-z\rangle\ge\frac13\langle z\rangle$; the sum over it is at most $3^{a'}\sum_{|z|\ge|x|/2}\langle z\rangle^{-a-a'}\le C|x|^{3-a-a'}$.
\end{proof}

\begin{lemma}[(P) A kernel against a slowly varying function]
\label{IV-73-41}
Let $D\in\{3,6\}$, $R\ge1$, and let $g:\mathbb Z^D\to\mathbb C$ vanish outside $\{|\zeta|\le R\}$. We write $d^\circ\kappa:=(2\pi)^{-D}d\kappa$.
\begin{enumerate}
\item[(a)] $|\check g(\kappa)|\le C\,R^D\langle R\kappa\rangle^{-10}\|g\|_{(R)}$ for $\kappa\in[-\pi,\pi]^D$, with a pure constant $C$.
\item[(b)] Let $\hat{\mathsf T}\in L^1(\mathbb T^{D/3})$, $0<c_0\le1$, $-D<s<9-D$, $\mathsf H\in\mathcal H_s(D)$ and $\iota\in\{0,1\}$, and assume $|\hat{\mathsf T}(\kappa)|\le\mathsf H(\kappa)(\lg|\kappa|)^\iota$ for $|\kappa|\le c_0$. Then
\[
\Bigl|\int_{\mathbb T^{D/3}}\hat{\mathsf T}(\kappa)\,\check g(-\kappa)\,d^\circ\kappa\Bigr|\le C\,R^{-s}(\lg R)^\iota\|g\|_{(R)},\qquad C=C(\mathsf H,c_0,\|\hat{\mathsf T}\|_{L^1}).
\]
\item[(c)] Let, besides the hypotheses of (b), $s<8-D$. Let $\hat{\mathsf T}^h$ be a Borel function on $\mathbb R^D$, homogeneous of degree $s$, with $|\hat{\mathsf T}^h|\le\mathsf H$, and let $|\hat{\mathsf T}-\hat{\mathsf T}^h|\le\mathsf H'\lg|\kappa|$ on $\{|\kappa|\le c_0\}$ for some $\mathsf H'\in\mathcal H_{s+1}(D)$. Let $a\in\mathbb R$, $A\ge0$ and $G\in C_c^\infty(\mathbb R^D)$ with $\operatorname{supp}G\subset\{|\zeta|\le1\}$, and assume $\|g-R^{-a}G(\cdot/R)\|_{(R)}\le AR^{-a-1}\lg R$. Then
\[
\Bigl|\int_{\mathbb T^{D/3}}\hat{\mathsf T}(\kappa)\,\check g(-\kappa)\,d^\circ\kappa-R^{-a-s}\int_{\mathbb R^D}\hat{\mathsf T}^h(\kappa)\,\hat G(-\kappa)\,d^\circ\kappa\Bigr|\le C\bigl(A+\|G\|_{C^{10}}\bigr)R^{-a-s-1}\lg R
\]
and $\bigl|\int_{\mathbb R^D}\hat{\mathsf T}^h(\kappa)\hat G(-\kappa)\,d^\circ\kappa\bigr|\le C\|G\|_{C^{10}}$, where $C=C(\mathsf H,\mathsf H',c_0,\|\hat{\mathsf T}\|_{L^1})$ only, independent of $g$, $G$, $R$, $a$ and $A$.
\end{enumerate}
\end{lemma}

\begin{proof}
\emph{Step 1: proof of (a).} For $j\in\{1,\dots,D\}$ we have $\sum_\zeta g(\zeta+e_j)e^{-i\kappa\cdot\zeta}=e^{i\kappa_j}\check g(\kappa)$; hence the function $\nabla_j^Ng$, which again has finite support, satisfies
\[
\sum_\zeta(\nabla_j^Ng)(\zeta)\,e^{-i\kappa\cdot\zeta}=(e^{i\kappa_j}-1)^N\check g(\kappa).
\]
For $N\le10$ the function $\nabla_j^Ng$ vanishes outside $\{|\zeta|\le11R\}$, a set of at most $CR^D$ points, and $|\nabla_j^Ng|\le R^{-N}\|g\|_{(R)}$. For $N=0$ this gives $|\check g|\le CR^D\|g\|_{(R)}$. Choosing $j$ with $|\kappa_j|\ge|\kappa|/\sqrt D$ and using $|e^{i\kappa_j}-1|\ge\frac2\pi|\kappa_j|$ for $|\kappa_j|\le\pi$, we obtain from $N=10$ that $|\check g(\kappa)|\le C|\kappa|^{-10}R^{D-10}\|g\|_{(R)}$. The two bounds together give (a).

\emph{Step 2: proof of (b).} We split the integral at $|\kappa|=c_0$. For $|\kappa|\ge c_0$, (a) gives $|\check g(-\kappa)|\le CR^{D-10}c_0^{-10}\|g\|_{(R)}$; since $D-10<-s$ and $R\ge1$, the integral over $\{|\kappa|\ge c_0\}$ is at most $C\|\hat{\mathsf T}\|_{L^1}c_0^{-10}R^{-s}\|g\|_{(R)}$. For $|\kappa|\le c_0$ we use the hypothesis and (a), and substitute $\kappa=\kappa'/R$; as $\mathsf H(\kappa'/R)=R^{-s}\mathsf H(\kappa')$ and $\lg(|\kappa'|/R)\le\lg|\kappa'|\,\lg R$,
\[
\int_{|\kappa|\le c_0}\mathsf H(\kappa)(\lg|\kappa|)^\iota R^D\langle R\kappa\rangle^{-10}\,d\kappa\le R^{-s}(\lg R)^\iota\int_{\mathbb R^D}\mathsf H(\kappa')(\lg|\kappa'|)^\iota\langle\kappa'\rangle^{-10}\,d\kappa' ,
\]
and the last integral is finite by Definition~\ref{IV-73-38}(6), because $s<9-D$.

\emph{Step 3: proof of (c).} Let $g^h:=R^{-a}G(\cdot/R)$ on $\mathbb Z^D$. It vanishes outside $\{|\zeta|\le R\}$, and $\|g^h\|_{(R)}\le R^{-a}\|G\|_{C^{10}}$, since a difference of order $\gamma$ of $G(\cdot/R)$ is an integral of a derivative of order $|\gamma|$ of $G(\cdot/R)$, which is bounded by $R^{-|\gamma|}\|G\|_{C^{10}}$. On $\{|\kappa|\le c_0\}$ we have $|\hat{\mathsf T}|\le\mathsf H+\mathsf H'\lg|\kappa|\le\mathsf H+\mathsf H'/|\kappa|$, because $t\lg t\le1$ for $0<t\le1$; and $\mathsf H+\mathsf H'/|\kappa|\in\mathcal H_s(D)$, as one sees in polar coordinates. The function $h:=g-g^h$ vanishes outside $\{|\zeta|\le R\}$, and assertion (b), applied to $h$ in place of $g$, with this majorant and with $\iota=0$, gives
\[
\Bigl|\int\hat{\mathsf T}(\kappa)\,\check h(-\kappa)\,d^\circ\kappa\Bigr|\le CR^{-s}\|h\|_{(R)}\le CA\,R^{-s-a-1}\lg R .
\]
Next, the function $\hat{\mathsf T}-\chi(\cdot/c_0)\hat{\mathsf T}^h$ is integrable on the cube and bounded by $\mathsf H'\lg|\kappa|$ on $\{|\kappa|\le c_0\}$, where $\chi(\kappa/c_0)=1$. Assertion (b), applied with $s+1$ in place of $s$ (admissible, since $s+1<9-D$), with the majorant $\mathsf H'$, with $\iota=1$ and with $g^h$ in place of $g$, gives
\[
\Bigl|\int\bigl(\hat{\mathsf T}-\chi(\cdot/c_0)\hat{\mathsf T}^h\bigr)(\kappa)\,\check g^h(-\kappa)\,d^\circ\kappa\Bigr|\le CR^{-s-1}\lg R\,\|g^h\|_{(R)}\le CR^{-s-1-a}\lg R\,\|G\|_{C^{10}} .
\]
It remains to treat $\int\chi(\kappa/c_0)\hat{\mathsf T}^h(\kappa)\check g^h(-\kappa)\,d^\circ\kappa$. By the Poisson summation formula,
\[
\check g^h(\kappa)=R^{D-a}\sum_{\mathsf m\in\mathbb Z^D}\hat G\bigl(R(\kappa+2\pi\mathsf m)\bigr),\qquad\text{and}\qquad|\hat G(\xi)|\le C\langle\xi\rangle^{-10}\|G\|_{C^{10}} .
\]
For $|\kappa|\le2c_0\le2$ and $\mathsf m\ne0$ we have $|\kappa+2\pi\mathsf m|\ge\pi|\mathsf m|$; hence the terms with $\mathsf m\ne0$ contribute at most $CR^{D-a-10}\|G\|_{C^{10}}\le CR^{-a-s-1}\|G\|_{C^{10}}$, as $D-10<-s-1$. The term with $\mathsf m=0$ gives, after the substitution $\kappa=\kappa'/R$,
\[
R^{-a-s}\int_{\mathbb R^D}\chi\bigl(\kappa'/c_0R\bigr)\,\hat{\mathsf T}^h(\kappa')\,\hat G(-\kappa')\,d^\circ\kappa' .
\]
Since $\int_{|\kappa'|\ge c_0R}\mathsf H(\kappa')\langle\kappa'\rangle^{-10}d\kappa'\le CR^{s+D-10}\le CR^{-1}$, replacing $\chi(\kappa'/c_0R)$ by $1$ changes this expression by at most $CR^{-a-s-1}\|G\|_{C^{10}}$. The bound $|\hat{\mathsf T}^h(\kappa')\hat G(-\kappa')|\le C\mathsf H(\kappa')\langle\kappa'\rangle^{-10}\|G\|_{C^{10}}$ also gives the last assertion of (c). Collecting the estimates we obtain (c).
\end{proof}

\begin{lemma}[(L) One-loop functions]
\label{IV-73-42}
Let $\mathsf m\in\{2,3\}$ and $D:=3(\mathsf m-1)$. Let $\sigma_1,\dots,\sigma_{\mathsf m}$ be line symbols of degrees $s_1,\dots,s_{\mathsf m}$, let $\epsilon_i\in\{\pm1\}$, let $\lambda_i:\mathbb R^D\to\mathbb R^3$ be linear, and put $l_i(\kappa):=-\epsilon_i\lambda_i(\kappa)$. We assume:
\begin{itemize}
\item there is a number $c_\lambda\ge1$ with $\max_i|l_i(\kappa)|\le c_\lambda|\kappa|\le c_\lambda^2\max_{i<j}|l_i(\kappa)-l_j(\kappa)|$ for all $\kappa\in\mathbb R^D$;
\item $l_i\ne l_j$ as maps, for $i\ne j$;
\item the linear map $(l,\kappa)\mapsto(\epsilon_il+\lambda_i(\kappa))_{i=1,\dots,\mathsf m}$ of $(\mathbb R^3)^{\mathsf m}=\mathbb R^3\times\mathbb R^D$ has an integer matrix of determinant $\pm1$;
\item if $\mathsf m=3$, then $s_i+s_j\ge-3$ for $i\ne j$;
\item $\omega:=3+\sum_is_i\ge-1$.
\end{itemize}
We put
\[
\mathsf I(\kappa):=\int_{\mathbb T}\prod_{i=1}^{\mathsf m}\sigma_i\bigl(\epsilon_il+\lambda_i(\kappa)\bigr)\,d^\circ l,\qquad\Lambda(\kappa):=1+\sum_{i<j}\log_+\frac{|\kappa|}{|l_i(\kappa)-l_j(\kappa)|}\quad(\kappa\ne0).
\]
The function $\Lambda$ is homogeneous of degree $0$ and even, and $\Lambda(\kappa)=\infty$ exactly where two of the points $l_i(\kappa)$ agree.
\begin{enumerate}
\item[(0)] The integral defining $\mathsf I(\kappa)$ converges absolutely for almost every $\kappa$, and $\|\mathsf I\|_{L^1(\mathbb T^{D/3})}\le\prod_i\|\sigma_i\|_{L^1(\mathbb T)}$. It converges absolutely for every $\kappa$ with $0<c_\lambda|\kappa|\le\frac12$ at which the points $l_i(\kappa)$ are pairwise different.
\end{enumerate}
Moreover, there are numbers $c_0\in(0,1]$ and $C$, depending on the $\sigma_i$, $\epsilon_i$, $\lambda_i$ and $c_\lambda$, a polynomial $\mathsf Q$ in $\kappa$ of degree $\le\omega$, with $\mathsf Q=0$ if $\omega=-1$, and a Borel function $\mathsf I^h$ on $\mathbb R^D$, vanishing at $\kappa=0$ and wherever two of the points $l_i(\kappa)$ agree, with the following properties.
\begin{enumerate}
\item[(a)] $\mathsf I^h$ is homogeneous of degree $\omega$, $\mathsf I^h(-\kappa)=(-1)^{\omega+1}\mathsf I^h(\kappa)$, and $|\mathsf I^h(\kappa)|\le C|\kappa|^\omega\Lambda(\kappa)$. It depends only on the $\sigma_i^h$, $\epsilon_i$ and $\lambda_i$, and multilinearly on the $\sigma_i^h$. If $O$ is a rotation acting on each of the $\mathsf m-1$ arguments in $\mathbb R^3$ of $\kappa$, and the $\lambda_i$ commute with this action, then replacing every $\sigma_i^h$ by $\sigma_i^h\circ O^{-1}$ replaces $\mathsf I^h$ by $\mathsf I^h\circ O^{-1}$. If $\mathsf m=2$, then $\mathsf I^h$ is continuous on $\mathbb R^D\setminus\{0\}$.
\item[(b)] $|\mathsf I(\kappa)-\mathsf Q(\kappa)-\mathsf I^h(\kappa)|\le C|\kappa|^{\omega+1}\lg|\kappa|\,\Lambda(\kappa)$ for $0<|\kappa|\le c_0$ whenever the points $l_i(\kappa)$ are pairwise different.
\item[(c)] If $\omega\ge1$, then $\prod_i\sigma_i(\epsilon_il)$ is integrable over $\mathbb T$, and $\mathsf Q(0)=\int_{\mathbb T}\prod_i\sigma_i(\epsilon_il)\,d^\circ l$.
\item[(d)] (Uniqueness.) Let $\mathsf K^\sharp\subset\mathbb R^D\setminus\{0\}$ be open and non-empty, with $t\mathsf K^\sharp=\mathsf K^\sharp$ for every real $t\ne0$, such that the $l_i$ are pairwise different at every point of $\mathsf K^\sharp$. Let $\mathsf Q^\sharp$ be a polynomial of degree $\le\omega$, with $\mathsf Q^\sharp=0$ if $\omega=-1$, and let $\mathsf I^\sharp$ be a function on $\mathsf K^\sharp$, homogeneous of degree $\omega$, with $\mathsf I^\sharp(-\kappa)=(-1)^{\omega+1}\mathsf I^\sharp(\kappa)$. If $t^{-\omega}|\mathsf I(t\kappa)-\mathsf Q^\sharp(t\kappa)-\mathsf I^\sharp(t\kappa)|\to0$ as $t\downarrow0$ for every $\kappa\in\mathsf K^\sharp$, then $\mathsf Q^\sharp=\mathsf Q$, and $\mathsf I^\sharp=\mathsf I^h$ on $\mathsf K^\sharp$.

Thus $\mathsf Q$ and $\mathsf I^h$ are determined by $\mathsf I$ near $0$, and they are independent of $c_\lambda$; we call them \emph{the polynomial} and \emph{the homogeneous part} of $\mathsf I$. If $\tilde{\mathsf I}$ is a function of the same kind, with data $\tilde\sigma_i$, $\tilde\epsilon_i$, $\tilde\lambda_i$ and the same $\mathsf m$ and $\omega$, and if $\tilde{\mathsf I}(\kappa)=\mathsf I(-\kappa)$ for all small $\kappa\ne0$ at which both families of points, the $l_i(\kappa)$ and the $\tilde l_i(\kappa)$, are pairwise different, then the polynomial of $\tilde{\mathsf I}$ is the polynomial $\kappa\mapsto\mathsf Q(-\kappa)$, and the homogeneous part of $\tilde{\mathsf I}$ equals $\mathsf I^h(-\kappa)$ at every $\kappa\ne0$ at which both families of points are pairwise different.
\item[(e)] $|\kappa|^s\Lambda(\kappa)\in\mathcal H_s(D)$ for every real $s>-D$.
\end{enumerate}
\end{lemma}

\begin{proof}
\emph{Step 0: two bounds.} Let $t_1,\dots,t_{\mathsf m}>-3$ be real, with $t_i+t_j\ge-3$ for $i\ne j$ if $\mathsf m=3$; let $\kappa\ne0$ and $\varepsilon:=c_\lambda|\kappa|$. We first show
\begin{equation}\label{IV-73-43}
\int_{|l|\le2\varepsilon}\prod_i|l-l_i(\kappa)|^{t_i}\,dl\le C\varepsilon^{3+\sum_it_i}\Lambda(\kappa).
\end{equation}
By the scaling $l\to\varepsilon l$, $l_i(\kappa)\to\varepsilon l_i(\kappa)$ it suffices to treat $\varepsilon=1$; then $|l_i|\le1$ and $\max_{i<j}|l_i-l_j|\ge c_\lambda^{-2}=:2c$, and, the scaled distances being $|l_i-l_j|=|l_i(\kappa)-l_j(\kappa)|/(c_\lambda|\kappa|)$, we have $1+\log\frac1{|l_i-l_j|}\le(1+\log c_\lambda)\Lambda(\kappa)$ for $i\ne j$. If $\mathsf m=2$, the two singularities lie at distance at least $2c$ from each other, each of them is integrable, and the integral is at most $C$. Let $\mathsf m=3$, and number the points so that $d:=|l_1-l_2|$ is the least of the three distances. If $d\ge c/2$, the three singularities are apart from each other, and again the integral is at most $C$.

If $d<c/2$, then $|l_3-l_1|,|l_3-l_2|\ge c$, since one of them is at least $2c$ and they differ at most by $d$. On $\{|l-l_3|\le c/2\}$ the integrand is at most $C|l-l_3|^{t_3}$; off this set it is at most $C|l-l_1|^{t_1}|l-l_2|^{t_2}$. The integral of the latter over $\{|l-l_1|\le2d\}$ is at most $Cd^{3+t_1+t_2}\le C$, as $t_1+t_2\ge-3$ and $d<1$; over $\{2d\le|l-l_1|\le3\}$, where $\frac12|l-l_1|\le|l-l_2|\le\frac32|l-l_1|$, it is at most $C\int_{2d}^3\varrho^{2+t_1+t_2}\,d\varrho\le C(1+\log\frac1d)\le C\Lambda(\kappa)$. This proves \eqref{IV-73-43}.

We also prove (e) here. The function $|\kappa|^s\Lambda(\kappa)$ is non-negative, Borel and homogeneous of degree $s$, and, as $s>-D$, it is integrable over $\{|\kappa|\le1\}$ if $\Lambda$ is integrable over the sphere $\{|\kappa|=1\}$. For $i<j$ the linear map $l_i-l_j$ is not $0$, so there are a vector $v\in\mathbb R^D$ with $|v|=1$ and a number $c'>0$ with $|l_i(\kappa)-l_j(\kappa)|\ge c'|v\cdot\kappa|$; and $\log_+\bigl(1/(c'|v\cdot\kappa|)\bigr)$ is integrable over the sphere.

Next let $\tau_i\in\mathcal S^{t_i}$, $\mathsf A(l,\kappa):=\prod_i\tau_i(\epsilon_il+\lambda_i(\kappa))$, and $0<\varepsilon\le\frac12$. For $l\in[-\pi,\pi]^3$ with $|l|\ge2\varepsilon$ and for $|\kappa'|\le|\kappa|$ we have $\epsilon_il+\lambda_i(\kappa')=\epsilon_i(l-l_i(\kappa'))$ and $|l_i(\kappa')|\le\varepsilon\le\frac12|l|$, so that $\frac14|l|\le|\epsilon_il+\lambda_i(\kappa')|_{\mathbb T}\le\frac32|l|$. With the chain and Leibniz rules this gives
\begin{equation}\label{IV-73-44}
|\partial_\kappa^\alpha\mathsf A(l,\kappa')|\le C_\alpha|l|^{\sum_it_i-|\alpha|}\qquad(|\kappa'|\le|\kappa|).
\end{equation}
The same bound holds on $\mathbb R^3$, for all $l$ with $|l|\ge2\varepsilon$, if the $\tau_i$ are smooth on $\mathbb R^3\setminus\{0\}$ and homogeneous of degree $t_i$.

We now prove (0). By Fubini's theorem and the substitution $(l,\kappa)\to(\epsilon_il+\lambda_i(\kappa))_i$, which induces a measure-preserving bijection of $\mathbb T^{\mathsf m}$ since its matrix is an integer matrix of determinant $\pm1$, we get $\int_{\mathbb T^{D/3}}\int_{\mathbb T}\prod_i|\sigma_i(\epsilon_il+\lambda_i(\kappa))|\,d^\circ l\,d^\circ\kappa=\prod_i\|\sigma_i\|_{L^1(\mathbb T)}$; this gives the first two assertions. For $0<c_\lambda|\kappa|\le\frac12$ with pairwise different $l_i(\kappa)$ we have $\Lambda(\kappa)<\infty$; the integral over $\{|l|\le2\varepsilon\}$ converges by \eqref{IV-73-43} with $t_i=s_i$, since $|\sigma_i(\epsilon_il+\lambda_i(\kappa))|\le C|l-l_i(\kappa)|^{s_i}$ there, and the integral over $\{|l|\ge2\varepsilon\}$ by \eqref{IV-73-44} with $\alpha=0$.

\emph{Step 1: crude expansion.} Let $\tau_i\in\mathcal S^{t_i}$ with integers $t_i\ge-2$ as in Step~0, $\omega':=3+\sum_it_i\ge0$, $\mathsf A$ as in Step~0, and $\mathsf J(\kappa):=\int_{\mathbb T}\mathsf A(l,\kappa)\,d^\circ l$. We claim: there is a polynomial $\mathsf Q'$ of degree $<\omega'$ with
\[
|\mathsf J(\kappa)-\mathsf Q'(\kappa)|\le C|\kappa|^{\omega'}\lg|\kappa|\,\Lambda(\kappa)\qquad(0<c_\lambda|\kappa|\le\tfrac12),
\]
and $\mathsf Q'(0)=\int_{\mathbb T}\mathsf A(l,0)\,d^\circ l$ if $\omega'\ge1$. Indeed, the integral over $\{|l|\le2\varepsilon\}$ is $O(\varepsilon^{\omega'}\Lambda(\kappa))$ by \eqref{IV-73-43}. For $|l|\ge2\varepsilon$ we expand by Taylor's formula in $\kappa$ at $0$,
\[
\mathsf A(l,\kappa)=\sum_{|\alpha|<\omega'}\kappa^\alpha\mathsf A_\alpha(l)+\mathsf A_{\mathrm{rem}}(l,\kappa),\qquad|\mathsf A_{\mathrm{rem}}(l,\kappa)|\le C|\kappa|^{\omega'}|l|^{-3}
\]
by \eqref{IV-73-44}, so that $\int_{|l|\ge2\varepsilon}|\mathsf A_{\mathrm{rem}}|\,d^\circ l=O(|\kappa|^{\omega'}\lg|\kappa|)$. Moreover $|\mathsf A_\alpha(l)|\le C|l|^{\omega'-3-|\alpha|}$, which is integrable over $\mathbb T$ for $|\alpha|<\omega'$, and $|\kappa|^{|\alpha|}\int_{|l|\le2\varepsilon}|\mathsf A_\alpha|\,d^\circ l\le C|\kappa|^{\omega'}$. Hence the claim holds with $\mathsf Q'(\kappa):=\sum_{|\alpha|<\omega'}\kappa^\alpha\int_{\mathbb T}\mathsf A_\alpha\,d^\circ l$.

\emph{Step 2: leading parts.} Let $\mathsf A^h(l,\kappa):=\prod_i\sigma_i^h(\epsilon_il+\lambda_i(\kappa))$ on $\mathbb R^3\times\mathbb R^D$. It is homogeneous of degree $\omega-3$, and $\mathsf A^h(-l,-\kappa)=(-1)^{\omega+1}\mathsf A^h(l,\kappa)$. For $|l|\ge2\varepsilon$ we expand by Taylor's formula in $\kappa$ at $0$,
\[
\mathsf A^h(l,\kappa)=\sum_{j\le\omega}\mathsf A^h_j(l;\kappa)+\mathsf A^h_{\mathrm{rem}}(l;\kappa);
\]
here $\mathsf A^h_j$ is defined for $l\ne0$, homogeneous of degree $j$ in $\kappa$ and of degree $\omega-3-j$ in $l$, with $\mathsf A^h_j(-l;\kappa)=(-1)^{\omega+1+j}\mathsf A^h_j(l;\kappa)$, and $|\mathsf A^h_{\mathrm{rem}}(l;\kappa)|\le C|\kappa|^{\omega+1}|l|^{-4}$ by the bound \eqref{IV-73-44} for homogeneous factors. With $\varepsilon=c_\lambda|\kappa|$ we define, where $\kappa\ne0$ and the points $l_i(\kappa)$ are pairwise different,
\[
\mathsf I^h(\kappa):=\int_{|l|\le2\varepsilon}\mathsf A^h\,d^\circ l+\int_{|l|>2\varepsilon}\mathsf A^h_{\mathrm{rem}}\,d^\circ l-\sum_{j<\omega}\int_{|l|\le2\varepsilon}\mathsf A^h_j\,d^\circ l ,
\]
and $\mathsf I^h:=0$ elsewhere, that is, at $0$ and on a finite union of linear subspaces. The substitution $l\to tl$ gives the homogeneity, $l\to-l$ the parity, and $l\to Ol$ the behaviour under rotations. Each of the three terms is $O(|\kappa|^\omega\Lambda(\kappa))$, by \eqref{IV-73-43} and \eqref{IV-73-44}. The function $\mathsf I^h$ is Borel, as an absolutely convergent integral of a Borel function of $(l,\kappa)$. At fixed $c_\lambda$ it is built from the $\sigma_i^h$, $\epsilon_i$ and $\lambda_i$ alone, and it is multilinear in the $\sigma_i^h$.

\emph{Step 2$'$: continuity for $\mathsf m=2$.} Here $l_1(\kappa)\ne l_2(\kappa)$ for all $\kappa\ne0$, and by homogeneity it suffices to prove continuity on $\{|\kappa|=1\}$, where $\varepsilon=c_\lambda$ is fixed. The third term in the definition of $\mathsf I^h$ is a polynomial in $\kappa$: $\mathsf A^h_j(l;\kappa)=\sum_{|\alpha|=j}\kappa^\alpha\mathsf A^h_\alpha(l)$ with $|\mathsf A^h_\alpha(l)|\le C|l|^{\omega-3-j}$, a function integrable near $0$ for $j<\omega$, so that the coefficients $\int_{|l|\le2\varepsilon}\mathsf A^h_\alpha\,d^\circ l$ are finite. In the second term the integrand is continuous in $\kappa$ for every $l$ with $|l|>2\varepsilon$ and has the integrable majorant $C|l|^{-4}$; dominated convergence applies.

In the first term, $\mathsf A^h=\sigma_1^h(\epsilon_1(l-l_1(\kappa)))\,\sigma_2^h(\epsilon_2(l-l_2(\kappa)))$, and by the first hypothesis of the lemma $|l_1(\kappa)-l_2(\kappa)|\ge c_\lambda^{-1}=:4\delta$. We fix $\kappa_0$ with $|\kappa_0|=1$, put $\psi_1(l):=\theta(|l-l_1(\kappa_0)|/2\delta)$ and $\psi_2:=1-\psi_1$, and consider $\kappa$ with $|\kappa|=1$ so close to $\kappa_0$ that $|l_i(\kappa)-l_i(\kappa_0)|\le\frac12\delta$ for $i=1,2$. As $\theta=1$ on $[0,\frac12]$ and $\theta=0$ on $[1,\infty)$, we have $\psi_1=1$ where $|l-l_1(\kappa_0)|\le\delta$ and $\psi_1=0$ where $|l-l_1(\kappa_0)|\ge2\delta$. Hence, where $\psi_1\ne0$,
\[
|l-l_2(\kappa)|\ge|l_1(\kappa_0)-l_2(\kappa_0)|-|l-l_1(\kappa_0)|-|l_2(\kappa_0)-l_2(\kappa)|\ge4\delta-2\delta-\tfrac12\delta=\tfrac32\delta,
\]
and where $\psi_2\ne0$, that is, where $|l-l_1(\kappa_0)|>\delta$, we have $|l-l_1(\kappa)|\ge\delta-\frac12\delta=\frac12\delta$. We split the first term by the factors $\psi_1$ and $\psi_2$. In the integral with the factor $\psi_1$ we substitute $l\to l+l_1(\kappa)$; since $|l_1(\kappa)|\le c_\lambda|\kappa|=\varepsilon$, it becomes the integral over $\{|l+l_1(\kappa)|\le2\varepsilon\}\subset\{|l|\le3\varepsilon\}$ of
\[
\psi_1(l+l_1(\kappa))\,\sigma_1^h(\epsilon_1l)\,\sigma_2^h\bigl(\epsilon_2(l+l_1(\kappa)-l_2(\kappa))\bigr).
\]
The domain of integration moves with $\kappa$; we extend the integrand by $0$ off this domain and integrate over $\mathbb R^3$. Where $\psi_1(l+l_1(\kappa))\ne0$, the last factor is evaluated at points of modulus in $[\frac32\delta,4\varepsilon]$, a compact set not containing $0$, on which $\sigma_2^h$ is bounded and continuous. The singular factor $\sigma_1^h(\epsilon_1l)$ does not depend on $\kappa$, and $|\sigma_1^h(\epsilon_1l)|\le C|l|^{s_1}$ by homogeneity. Hence the extended integrand has the majorant $C|l|^{s_1}$, which does not depend on $\kappa$ and is integrable over $\{|l|\le3\varepsilon\}$ as $s_1\ge-2$; and as $\kappa\to\kappa_0$ it converges at every $l\ne0$ with $|l+l_1(\kappa_0)|\ne2\varepsilon$, that is, outside a null set. So dominated convergence applies. The integral with the factor $\psi_2$ is treated in the same way, with the substitution $l\to l+l_2(\kappa)$ and the majorant $C|l|^{s_2}$.

\emph{Step 2$''$: comparison.} Let $\mathsf X(l,\kappa):=\prod_i\chi(\epsilon_il+\lambda_i(\kappa))$, let $\chi_*(l):=\chi(l)^{\mathsf m}$, which is radial, and let $\mathsf I^{(0)}(\kappa):=\int_{\mathbb R^3}\mathsf X\mathsf A^h\,d^\circ l$ for $\varepsilon\le\frac14$. The function $\mathsf X-\chi_*$ is smooth and vanishes for $|l|\le\frac12$ and for $|l|\ge3$; therefore $\int(\mathsf X-\chi_*)\mathsf A^h\,d^\circ l$ is smooth in $\kappa$ and vanishes at $\kappa=0$, so it equals a polynomial of degree $\le\omega$ without constant term up to $O(|\kappa|^{\omega+1})$. Moreover
\[
\int\chi_*\mathsf A^h\,d^\circ l=\int_{|l|\le2\varepsilon}\mathsf A^h\,d^\circ l+\int_{|l|>2\varepsilon}\mathsf A^h_{\mathrm{rem}}\,d^\circ l-\int(1-\chi_*)\mathsf A^h_{\mathrm{rem}}\,d^\circ l+\sum_{j\le\omega}\int_{|l|>2\varepsilon}\chi_*\mathsf A^h_j\,d^\circ l .
\]
The third term is $O(|\kappa|^{\omega+1})$. For $j<\omega$ we have $\int_{|l|>2\varepsilon}\chi_*\mathsf A^h_j=\int\chi_*\mathsf A^h_j-\int_{|l|\le2\varepsilon}\mathsf A^h_j$, and $\int\chi_*\mathsf A^h_j\,d^\circ l$ is a homogeneous polynomial of degree $j$ in $\kappa$. For $j=\omega$ the function $\mathsf A^h_\omega(\cdot;\kappa)$ is odd, while $\chi_*$ and the domain $\{|l|>2\varepsilon\}$ are even, so that this term vanishes. Consequently
\[
\mathsf I^{(0)}=\mathsf Q^{(0)}+\mathsf I^h+O(|\kappa|^{\omega+1}),
\]
with a polynomial $\mathsf Q^{(0)}$ of degree $\le\omega$; if $\omega\ge1$, then $\mathsf Q^{(0)}(0)=\int\chi_*\mathsf A^h_0\,d^\circ l=\int\mathsf X(l,0)\mathsf A^h(l,0)\,d^\circ l$.

\emph{Step 3: proof of (b) and (c).} We write $\sigma_i=\chi\sigma_i^h+\sigma_i^\flat$ with $\sigma_i^\flat\in\mathcal S^{s_i+1}$ and multiply out the product in the definition of $\mathsf I$. For $\varepsilon\le\frac14$ the term in which all factors are leading is $\mathsf I^{(0)}$. Every other term is a function $\mathsf J$ as in Step~1, with $t_i\in\{s_i,s_i+1\}$ and $\omega'\ge\omega+1\ge0$, since $\chi\sigma_i^h\in\mathcal S^{s_i}$; the terms of degree $>\omega$ of its polynomial $\mathsf Q'$ are $O(|\kappa|^{\omega+1})$. This proves (b), with $\mathsf Q$ the sum of $\mathsf Q^{(0)}$ and of the polynomials $\mathsf Q'$ cut at degree $\omega$, and with $c_0:=\frac1{4c_\lambda}$, which lies in $(0,\frac14]$ as $c_\lambda\ge1$; for $0<|\kappa|\le c_0$ we have $\varepsilon\le\frac14$, as Step~$2''$ requires, and $c_\lambda|\kappa|\le\frac12$, as Step~1 requires.

If $\omega\ge1$, then $|\prod_i\sigma_i(\epsilon_il)|\le C|l|^{\omega-3}$ is integrable, and adding the constant terms we obtain $\mathsf Q(0)=\int_{\mathbb T}\prod_i(\chi\sigma_i^h+\sigma_i^\flat)(\epsilon_il)\,d^\circ l=\int_{\mathbb T}\prod_i\sigma_i(\epsilon_il)\,d^\circ l$, which is (c).

\emph{Step 4: uniqueness.} Let $\kappa\in\mathsf K^\sharp$. By (b), and since $\Lambda(t\kappa)=\Lambda(\kappa)<\infty$,
\[
t^{-\omega}\bigl|\mathsf I(t\kappa)-\mathsf Q(t\kappa)-\mathsf I^h(t\kappa)\bigr|\le Ct|\kappa|^{\omega+1}\lg(t|\kappa|)\,\Lambda(\kappa)\to0\qquad(t\downarrow0).
\]
Subtracting the hypothesis of (d) and writing $\mathsf Q_j$, $\mathsf Q^\sharp_j$ for the homogeneous terms of degree $j$, we get
\[
\sum_{j<\omega}t^{j-\omega}(\mathsf Q_j-\mathsf Q^\sharp_j)(\kappa)+(\mathsf Q_\omega-\mathsf Q^\sharp_\omega)(\kappa)+(\mathsf I^h-\mathsf I^\sharp)(\kappa)\to0 .
\]
Beginning with the least $j$, we conclude $(\mathsf Q_j-\mathsf Q^\sharp_j)(\kappa)=0$ for $j<\omega$, and then $(\mathsf Q_\omega-\mathsf Q^\sharp_\omega)(\kappa)+(\mathsf I^h-\mathsf I^\sharp)(\kappa)=0$. The same holds at $-\kappa\in\mathsf K^\sharp$, where the two terms acquire the factors $(-1)^\omega$ and $(-1)^{\omega+1}$; hence also $(\mathsf Q_\omega-\mathsf Q^\sharp_\omega)(\kappa)-(\mathsf I^h-\mathsf I^\sharp)(\kappa)=0$, and both terms vanish. (If $\omega=-1$, only the term $\mathsf I^h-\mathsf I^\sharp$ is present.) Polynomials which vanish on a non-empty open set vanish identically; so $\mathsf Q^\sharp=\mathsf Q$, and $\mathsf I^\sharp=\mathsf I^h$ on $\mathsf K^\sharp$.

The set of all $\kappa\ne0$ at which the $l_i(\kappa)$ are pairwise different is such a set $\mathsf K^\sharp$, being the complement of finitely many proper linear subspaces, and $\mathsf I^h=0$ outside it; therefore the pair $(\mathsf Q,\mathsf I^h)$ with (a) and (b) is unique and independent of $c_\lambda$. As, by Step~2, at fixed $c_\lambda$ the function $\mathsf I^h$ is built from the $\sigma_i^h$, $\epsilon_i$ and $\lambda_i$ alone, it depends only on these data.

For the last assertion of (d) let $\tilde{\mathsf Q}$ and $\tilde{\mathsf I}^h$ be the polynomial and the homogeneous part of $\tilde{\mathsf I}$, and let $\mathsf K^\sharp$ be the set of all $\kappa\ne0$ at which the $l_i(\kappa)$ are pairwise different and so are the $\tilde l_i(\kappa)$. It is again the complement of finitely many proper linear subspaces, hence open, non-empty and invariant under $\kappa\to t\kappa$ for real $t\ne0$. We put $\mathsf Q^\sharp(\kappa):=\tilde{\mathsf Q}(-\kappa)$ and $\mathsf I^\sharp(\kappa):=\tilde{\mathsf I}^h(-\kappa)$ on $\mathsf K^\sharp$; they have the degree, the homogeneity and the parity required in (d). For $\kappa\in\mathsf K^\sharp$ and small $t>0$ the hypothesis gives $\mathsf I(t\kappa)=\tilde{\mathsf I}(-t\kappa)$, and (b) for the data of $\tilde{\mathsf I}$, at the point $-t\kappa$, gives
\[
t^{-\omega}\bigl|\mathsf I(t\kappa)-\mathsf Q^\sharp(t\kappa)-\mathsf I^\sharp(t\kappa)\bigr|=t^{-\omega}\bigl|\tilde{\mathsf I}(-t\kappa)-\tilde{\mathsf Q}(-t\kappa)-\tilde{\mathsf I}^h(-t\kappa)\bigr|\to0\qquad(t\downarrow0).
\]
By the first part of (d), $\mathsf Q^\sharp=\mathsf Q$ and $\mathsf I^\sharp=\mathsf I^h$ on $\mathsf K^\sharp$. Thus $\tilde{\mathsf Q}(\kappa)=\mathsf Q(-\kappa)$ for all $\kappa\in\mathbb R^D$, as an identity of polynomials, and, $\mathsf K^\sharp$ being symmetric, $\tilde{\mathsf I}^h(\kappa)=\mathsf I^h(-\kappa)$ for all $\kappa\in\mathsf K^\sharp$.
\end{proof}

\begin{lemma}[(D) Differences of an asymptotic behaviour]
\label{IV-73-45}
Let $n>0$ be real and $f:\mathbb Z^3\to\mathbb C$. Assume that for every multi-index $\gamma$ there are a continuous function $\mathsf H_\gamma$ on $\mathbb R^3\setminus\{0\}$, homogeneous of degree $-n-|\gamma|$, and numbers $c_\gamma$, $r_\gamma$ with
\[
|\nabla^\gamma f(x)-\mathsf H_\gamma(x)|\le c_\gamma|x|^{-n-|\gamma|-1}\lg|x|\qquad(x\in\mathbb Z^3\setminus\{0\},\ |x|\ge r_\gamma).
\]
Then $\mathsf H_0\in C^\infty(\mathbb R^3\setminus\{0\})$, $\partial^\gamma\mathsf H_0=\mathsf H_\gamma$ for every $\gamma$, and
\[
|\nabla^\gamma(f-\mathsf H_0)(x)|\le c'_\gamma|x|^{-n-|\gamma|-1}\lg|x|\qquad\bigl(|x|\ge\max\{r_\gamma,|\gamma|+1\}\bigr),
\]
with $c'_\gamma=c'_\gamma(n,\gamma,c_\gamma,\mathsf H_0)$.
\end{lemma}

\begin{proof}
Let $y\in\mathbb R^3\setminus\{0\}$, $\mu\in\{1,2,3\}$, and let $t>0$ be so small that $y+se_\mu\ne0$ for $0\le s\le t$. For $R>0$ let $x_R:=(\lfloor Ry_\nu\rfloor)_\nu\in\mathbb Z^3$ and $N_R:=\lfloor Rt\rfloor$. Then
\[
\nabla^\gamma f(x_R+N_Re_\mu)-\nabla^\gamma f(x_R)=\sum_{j<N_R}\nabla^{\gamma+e_\mu}f(x_R+je_\mu).
\]
We multiply by $R^{n+|\gamma|}$ and let $R\to\infty$. By the hypothesis, the homogeneity and the continuity of $\mathsf H_\gamma$, the left side tends to $\mathsf H_\gamma(y+te_\mu)-\mathsf H_\gamma(y)$. The points $x_R+je_\mu$, $j<N_R$, have modulus at least $cR$ for some $c>0$; by the hypothesis for $\gamma+e_\mu$ and the homogeneity of $\mathsf H_{\gamma+e_\mu}$, the right side equals
\[
\frac1R\sum_{j<N_R}\mathsf H_{\gamma+e_\mu}\bigl((x_R+je_\mu)/R\bigr)+O(R^{-1}\lg R),
\]
and this tends to $\int_0^t\mathsf H_{\gamma+e_\mu}(y+se_\mu)\,ds$: it is a Riemann sum, up to the shift $x_R/R-y=O(R^{-1})$ and the uniform continuity of $\mathsf H_{\gamma+e_\mu}$ on a compact neighbourhood of $\{y+se_\mu:0\le s\le t\}$. Hence $\mathsf H_\gamma(y+te_\mu)-\mathsf H_\gamma(y)=\int_0^t\mathsf H_{\gamma+e_\mu}(y+se_\mu)\,ds$ for small $t>0$, and, applied at the point $y-te_\mu$, the corresponding identity holds for small $t<0$. Therefore $\partial_\mu\mathsf H_\gamma$ exists and equals the continuous function $\mathsf H_{\gamma+e_\mu}$. By induction on $|\gamma|$, all partial derivatives of $\mathsf H_0$ exist, are continuous, and $\partial^\gamma\mathsf H_0=\mathsf H_\gamma$; thus $\mathsf H_0\in C^\infty(\mathbb R^3\setminus\{0\})$.

Finally, Taylor's formula and the homogeneity of the derivatives of order $|\gamma|+1$ of $\mathsf H_0$ give $|\nabla^\gamma\mathsf H_0(x)-\partial^\gamma\mathsf H_0(x)|\le C|x|^{-n-|\gamma|-1}$ for $|x|\ge|\gamma|+1$, with $C$ depending on $n$, $\gamma$ and $\mathsf H_0$; together with the hypothesis this proves the bound for $\nabla^\gamma(f-\mathsf H_0)$.
\end{proof}

\begin{lemma}[(E) Chains (convolutions of lines), bubbles]
\label{IV-73-46}
For $X,Y:\mathbb Z^3\to\mathbb C$ we write $(X*Y)(u):=\sum_zX(z)Y(u-z)$, and $r:=|u|$.
\begin{enumerate}
\item[(a)] (Bound.) Let $X$, $Y$ be kernels with symbols $\hat X,\hat Y\in L^1(\mathbb T)$; let $a_X,a_Y\in(-3,3)$ and $n_X,n_Y\ge0$ with $n_X+n_Y>3$. Assume $|\hat X(k)|\le c|k|^{a_X}\lg|k|$ on $[-\pi,\pi]^3$ and $|\nabla^\gamma X(x)|\le c\langle x\rangle^{-n_X-|\gamma|}$ for $|\gamma|\le10$ and all $x$, and the same for $Y$ with $a_Y$, $n_Y$. Then the sum defining $(X*Y)(u)$ converges absolutely, and for $r\ge1$
\[
|(X*Y)(u)|\le C\bigl(r^{-a_X-n_Y}+r^{-a_Y-n_X}+r^{3-n_X-n_Y}\bigr)\lg r,\qquad C=C(c,a_X,a_Y,n_X,n_Y).
\]
\item[(b)] (Class.) Let $X\in\mathfrak K(n_X)$ and $Y\in\mathfrak K(n_Y)$ with $n_X,n_Y\le4$ and $n_X+n_Y\ge5$; we write $\hat X$ and $\hat Y$ for their symbols, as in (K1), and $(\hat X^h,X^h)$ and $(\hat Y^h,Y^h)$ for their leading data in the sense of Definition~\ref{IV-73-38}(4). Then $X*Y\in\mathfrak K(n_X+n_Y-3)$, with the symbol $\hat X\hat Y$. Its leading data consist of $\hat X^h\hat Y^h$ and of a function which depends only on the leading data of $X$ and $Y$, bilinearly. The constants of $X*Y$ depend only on those of $X$ and $Y$. The leading data are rotation covariant: if the leading data of $X$ and of $Y$ are composed with $O^{-1}$, $O$ a rotation, then so are those of $X*Y$.
\item[(c)] (Bubbles.) Let $K_{\mathrm a}$, $K_{\mathrm b}$ be elementary lines of degrees $n_{\mathrm a},n_{\mathrm b}\in\{1,2,3\}$, with symbols $\sigma_{\mathrm a}$, $\sigma_{\mathrm b}$; let $X:=K_{\mathrm a}K_{\mathrm b}$ (pointwise product) and $n:=n_{\mathrm a}+n_{\mathrm b}$. Then $X$ is the kernel with the symbol
\[
\mathsf I(k)=\int_{\mathbb T}\sigma_{\mathrm a}(k-l)\,\sigma_{\mathrm b}(l)\,d^\circ l ,
\]
a function as in Lemma~\ref{IV-73-42} with $\mathsf m=2$ and $\omega=n-3$. With its polynomial $\mathsf Q$ and its homogeneous part $\mathsf I^h$ we put
\[
\tilde{\mathsf Q}(k):=\mathsf Q(0)+\sum_j\partial_j\mathsf Q(0)\sin k_j ,
\]
and we let $\mathsf q$ be the kernel with the symbol $\tilde{\mathsf Q}$; it vanishes outside $\{0,\pm e_1,\pm e_2,\pm e_3\}$. Then $|\nabla^\gamma X(x)|\le C_\gamma\langle x\rangle^{-n-|\gamma|}$ for all $x$ and $\gamma$, and $|\mathsf I(k)-\tilde{\mathsf Q}(k)|\le C|k|^{\min\{n-3,2\}}$ on $[-\pi,\pi]^3$, where $C_\gamma$ and $C$ depend on $\gamma$, $\sigma_{\mathrm a}$, $\sigma_{\mathrm b}$. If $n\le4$, then $X-\mathsf q\in\mathfrak K(n)$, with the leading data $(\mathsf I^h,K_{\mathrm a}^hK_{\mathrm b}^h)$, where $K_{\mathrm a}^h=\mathcal F^{-1}_{\ne0}\sigma_{\mathrm a}^h$ and $K_{\mathrm b}^h=\mathcal F^{-1}_{\ne0}\sigma_{\mathrm b}^h$.
\item[(d)] (Uniqueness.) Let $K\in\mathfrak K(n)$ with $n\le6$ and leading data $(\sigma^h,K^h)$. If $\sigma^h=0$, then $K^h=0$. Hence, within $\mathfrak K(n)$, $K^h$ is determined by $\sigma^h$, and depends linearly on it.
\end{enumerate}
\end{lemma}

\begin{proof}
\emph{Step 0: splitting.} Let $r\ge8$; for $1\le r<8$ the bound in (a) follows from Lemma~\ref{IV-73-40}(c). For $z\in\mathbb R^3$ we put $\theta_1(z):=\theta(4|z|/r)$ and $\theta_2(z):=\theta(4|u-z|/r)$. These are smooth functions on $\mathbb R^3$, as $\theta\in C^\infty([0,\infty))$ is constant on $[0,\frac12]$; they vanish for $|z|\ge r/4$, resp. for $|u-z|\ge r/4$, so that their supports are disjoint, as $|u|=r$. We write $X*Y=\mathsf T_1+\mathsf T_2+\mathsf T_0$, where the three sums carry the factors $\theta_1$, $\theta_2$ and $1-\theta_1-\theta_2$. The function $g_1:=\theta_1\,Y(u-\cdot)$ on $\mathbb Z^3$ vanishes outside $\{|z|\le r/4\}$, a set on which $|u-z|\ge\frac34r$. Since $g_1$ has finite support and $X(z)=\int_{\mathbb T}e^{ik\cdot z}\hat X(k)\,d^\circ k$, we may interchange the sum and the integral:
\[
\mathsf T_1=\sum_zX(z)g_1(z)=\int_{\mathbb T}\hat X(k)\sum_zg_1(z)\,e^{ik\cdot z}\,d^\circ k=\int_{\mathbb T}\hat X(k)\,\check g_1(-k)\,d^\circ k .
\]
After the substitution $z\to u-z$ the same argument gives $\mathsf T_2=\int_{\mathbb T}\hat Y(k)\,\check g_2(-k)\,d^\circ k$ with $g_2:=\theta(4|\cdot|/r)\,X(u-\cdot)$.

\emph{Step 1: proof of (a).} The sum converges absolutely by Lemma~\ref{IV-73-40}(c) with $a=n_X$, $a'=n_Y$. By the discrete Leibniz rule $\nabla_j(fg)=(\nabla_jf)\,g(\cdot+e_j)+f\,\nabla_jg$, a difference of order $\le10$ of $g_1$ is a sum of products of differences of $\theta_1$, bounded by $|\nabla^\beta\theta_1|\le Cr^{-|\beta|}$, with differences of $Y$ of order $\le10$ taken at points $u-z'$ with $|z'|\le r/4+10$. At these points $\langle u-z'\rangle\ge r/40$: for $r\le40$ because $\langle u-z'\rangle\ge1$, and for $r>40$ because $|u-z'|\ge\frac34r-10\ge r/40$. With the bounds on the differences of $Y$ this gives $\|g_1\|_{(r)}\le Cr^{-n_Y}$; in the same way $\|g_2\|_{(r)}\le Cr^{-n_X}$.

Lemma~\ref{IV-73-41}(b), applied with $D=3$, $R=r$, $g=g_1$, $\hat{\mathsf T}=\hat X$, the majorant $\mathsf H=c|k|^{a_X}\in\mathcal H_{a_X}(3)$ (admissible since $-3<a_X<3<6$), $\iota=1$ and $c_0=1$, gives $|\mathsf T_1|\le Cr^{-a_X-n_Y}\lg r$; likewise $|\mathsf T_2|\le Cr^{-a_Y-n_X}\lg r$. On the support of $1-\theta_1-\theta_2$ we have $|z|\ge r/8$ and $|u-z|\ge r/8$, hence $\langle z\rangle\ge\frac19(r+|z|)$ and $\langle u-z\rangle\ge\frac19(r+|u-z|)$; together with $r+|u-z|\ge\frac12(r+|z|)$ (Lemma~\ref{IV-73-40}(a)) this gives
\[
|X(z)Y(u-z)|\le C(r+|z|)^{-n_X}(r+|u-z|)^{-n_Y}\le C(r+|z|)^{-n_X-n_Y},
\]
and Lemma~\ref{IV-73-40}(a) with the exponent $n_X+n_Y>3$ gives $|\mathsf T_0|\le Cr^{3-n_X-n_Y}$.

\emph{Step 2: proof of (b), conditions (K1) and (K2).} As $n_X,n_Y\ge2$, conditions (K3) and the homogeneity of $X^h$, $Y^h$ give $|X(x)|\le C\langle x\rangle^{-n_X}$ and $|Y(x)|\le C\langle x\rangle^{-n_Y}$, so that $X,Y\in\ell^2(\mathbb Z^3)$; by the uniqueness of Fourier coefficients $\hat X,\hat Y\in L^2(\mathbb T)$, and Plancherel's theorem gives $(X*Y)(u)=\int_{\mathbb T}e^{ik\cdot u}\hat X\hat Y\,d^\circ k$, which is (K1) with the symbol $\hat X\hat Y\in L^1(\mathbb T)$. For (K2) let $s_X:=n_X-3$ and $s_Y:=n_Y-3$, both $\ge-1$; then
\[
\hat X\hat Y-\chi\hat X^h\hat Y^h=(\hat X-\chi\hat X^h)\hat Y+\chi\hat X^h(\hat Y-\chi\hat Y^h)+(\chi^2-\chi)\hat X^h\hat Y^h=O\bigl(|k|^{s_X+s_Y+1}\lg|k|\bigr),
\]
and $\hat X^h\hat Y^h$ is continuous and homogeneous of degree $s_X+s_Y=(n_X+n_Y-3)-3$.

\emph{Step 3: proof of (b), condition (K3) for $\gamma=0$.} Let $n:=n_X+n_Y-3$, and let $\gamma$ be a multi-index; in this step $\gamma=0$, in Step~4 it is arbitrary. We take $r\ge200+8|\gamma|$. In $\mathsf T_1$ the line $Y$ enters through $g_1$, and in $\mathsf T_2$ the line $X$ through $g_2$; in the estimates below these lines are evaluated, together with their differences of the orders $\le|\gamma|+10$ (the order $\gamma$ of Step~4 and up to ten further differences from the norm $\|\cdot\|_{(r)}$), at points of modulus at least $r/2\ge|\gamma|+11$. In $\mathsf T_0$ both lines, with their differences of the orders $\le|\gamma|$, are evaluated at points of modulus at least $r/8\ge|\gamma|+1$. Thus all these points lie in the range $|x|\ge|\gamma'|+1$ in which (K3) for $X$ and $Y$ bounds the differences of order $\gamma'$. Let $\hat u:=u/r$ and
\[
G_1(\zeta):=\theta(4|\zeta|)\,Y^h(\hat u-\zeta),\qquad G_2(\zeta):=\theta(4|\zeta|)\,X^h(\hat u-\zeta);
\]
these functions lie in $C_c^\infty(\mathbb R^3)$, with support in $\{|\zeta|\le\frac14\}$, and $r^{-n_Y}G_1(z/r)=\theta_1(z)Y^h(u-z)$ by homogeneity. By (K3) for $Y$, $\|g_1-r^{-n_Y}G_1(\cdot/r)\|_{(r)}\le Cr^{-n_Y-1}\lg r$. We use Lemma~\ref{IV-73-41}(c) with $D=3$, $R=r$, $g=g_1$, $\hat{\mathsf T}=\hat X$, $\hat{\mathsf T}^h=\hat X^h$, $s=s_X\le1$, $\iota=1$, $c_0=1$, $a=n_Y$, $G=G_1$, the constant of the preceding bound as $A$, and the majorants $\mathsf H=2c|k|^{s_X}$, $\mathsf H'=c|k|^{s_X+1}$, where $c$ bounds both $|\hat X^h|$ on the unit sphere and the number in (K2) for $X$. The hypotheses of (b) and (c) of that lemma are satisfied: on $\{|k|\le1\}$ we have $\chi=1$, so that (K2) for $X$ gives $|\hat X-\hat X^h|\le\mathsf H'\lg|k|$ there, and hence $|\hat X|\le\mathsf H\lg|k|$, while $|\hat X^h|\le\mathsf H$ everywhere. We obtain
\[
\mathsf T_1=r^{-s_X-n_Y}\int\hat X^h(k)\,\hat G_1(-k)\,d^\circ k+O\bigl(r^{-s_X-n_Y-1}\lg r\bigr),
\]
and in the same way $\mathsf T_2=r^{-s_Y-n_X}\int\hat Y^h(k)\,\hat G_2(-k)\,d^\circ k+O(r^{-s_Y-n_X-1}\lg r)$; note $s_X+n_Y=s_Y+n_X=n$. In $\mathsf T_0$ we replace $X$, $Y$ by $X^h$, $Y^h$; by (K3) the error is at most
\[
C\sum_z(r+|z|)^{-n_X-n_Y-1}\lg(r+|z|)\le Cr^{2-n_X-n_Y}\lg r ,
\]
since $\lg(r+|z|)\le\lg r\,\lg(1+|z|/r)\le\lg r\,(r+|z|)/r$, and by Lemma~\ref{IV-73-40}(a) with the exponent $n_X+n_Y>3$. Next we replace the sum by the integral of $\mathsf F(z):=(1-\theta_1-\theta_2)(z)X^h(z)Y^h(u-z)$, $z\in\mathbb R^3$. The factor $1-\theta_1-\theta_2$ vanishes for $|z|\le r/8$ and for $|u-z|\le r/8$, that is, near the points $0$ and $u$ at which $X^h$ and $Y^h(u-\cdot)$ are singular; so $\mathsf F$, extended by $0$ at these two points, is smooth on $\mathbb R^3$, in particular $\mathsf F\in C^1(\mathbb R^3)$. Since moreover $|\mathsf F(z)|\le C(r+|z|)^{-n_X-n_Y}$ and $|\nabla\mathsf F(z)|\le C(r+|z|)^{-n_X-n_Y-1}$, Lemma~\ref{IV-73-40}(b) applies, and together with (a) of that lemma it bounds the error by $Cr^{2-n_X-n_Y}$. The substitution $z=r\zeta$ gives
\[
\int_{\mathbb R^3}\mathsf F=r^{3-n_X-n_Y}\int_{\mathbb R^3}\bigl(1-\theta(4|\zeta|)-\theta(4|\hat u-\zeta|)\bigr)X^h(\zeta)\,Y^h(\hat u-\zeta)\,d\zeta .
\]
Altogether $(X*Y)(u)=\mathsf H_0(u)+O(r^{-n-1}\lg r)$, where
\[
\begin{multlined}
\mathsf H_0(u):=r^{-n}\Bigl[\int\hat X^h(k)\,\hat G_1(-k)\,d^\circ k+\int\hat Y^h(k)\,\hat G_2(-k)\,d^\circ k+\int\bigl(1-\theta(4|\zeta|)-\theta(4|\hat u \\
-\zeta|)\bigr)X^h(\zeta)\,Y^h(\hat u-\zeta)\,d\zeta\Bigr].
\end{multlined}
\]
The function $\mathsf H_0$ is homogeneous of degree $-n$, the bracket depending on $\hat u$ only. It is continuous: $\hat u\mapsto G_i$ is continuous with respect to $\|\cdot\|_{C^{10}}$, so that the last bound of Lemma~\ref{IV-73-41}(c) applies to the first two integrals; and the third integral depends continuously on $\hat u$ by dominated convergence, since on the support of its integrand $|\zeta|\ge\frac18$ and $|\hat u-\zeta|\ge\frac18$, so that the integrand is bounded by $C\langle\zeta\rangle^{-n_X-n_Y}$ uniformly in $\hat u$, and $n_X+n_Y>3$. It is bilinear in the leading data of $X$ and $Y$, and it is rotation covariant, by the substitutions $k\to Ok$, $\zeta\to O\zeta$.

\emph{Step 4: proof of (b), condition (K3) for all $\gamma$.} Let $\gamma$ be a multi-index and $r\ge200+8|\gamma|$. We have $\nabla^\gamma(X*Y)(u)=\sum_zX(z)(\nabla^\gamma Y)(u-z)$, which we split by the factors $\theta_1$, $\theta_2$ and $1-\theta_1-\theta_2$ as before.

In the first sum the function $g_{1,\gamma}:=\theta_1(\nabla^\gamma Y)(u-\cdot)$ takes the place of $g_1$; it vanishes outside $\{|z|\le r/4\}$. The function $G_{1,\gamma}(\zeta):=\theta(4|\zeta|)(\partial^\gamma Y^h)(\hat u-\zeta)$ lies in $C_c^\infty(\mathbb R^3)$, with support in $\{|\zeta|\le\frac14\}$, and $r^{-n_Y-|\gamma|}G_{1,\gamma}(z/r)=\theta_1(z)(\partial^\gamma Y^h)(u-z)$ by homogeneity. By (K3) for $Y$, and since $|(\nabla^\gamma-\partial^\gamma)Y^h(x)|\le C|x|^{-n_Y-|\gamma|-1}$ at the points concerned, the function $g_{1,\gamma}$ differs from $r^{-n_Y-|\gamma|}G_{1,\gamma}(\cdot/r)$ by at most $Cr^{-n_Y-|\gamma|-1}\lg r$ in $\|\cdot\|_{(r)}$.

As in Step~0 the first sum equals $\int_{\mathbb T}\hat X(k)\,\check g_{1,\gamma}(-k)\,d^\circ k$, and Lemma~\ref{IV-73-41}(c), with the data of Step~3 except that now $g=g_{1,\gamma}$, $a=n_Y+|\gamma|$, $G=G_{1,\gamma}$ and $A$ is the constant of the bound just stated, shows that it equals $r^{-n-|\gamma|}\int\hat X^h(k)\,\hat G_{1,\gamma}(-k)\,d^\circ k+O(r^{-n-|\gamma|-1}\lg r)$; here $s_X+n_Y=n$.

In the second sum we sum by parts $|\gamma|$ times, by the rule $\sum_zf(z)(\nabla_jY)(u-z)=\sum_z(\nabla_jf)(z)Y(u-z)$ for $f$ of finite support, and obtain $\sum_{z'}Y(z')g_{2,\gamma}(z')$ with $g_{2,\gamma}(z'):=\nabla^\gamma(\theta_2X)(u-z')$. This function vanishes outside $\{|z'|\le r/4+|\gamma|\}$, a subset of $\{|z'|\le r\}$. The function
\[
G_{2,\gamma}(\zeta):=\partial^\gamma_y\bigl[\theta(4|\hat u-y|)X^h(y)\bigr]_{y=\hat u-\zeta}
\]
lies in $C_c^\infty(\mathbb R^3)$, with support in $\{|\zeta|\le\frac14\}$, and by (K3) for $X$, the discrete Leibniz rule and Taylor's formula, as for the first sum, $g_{2,\gamma}$ differs from $r^{-n_X-|\gamma|}G_{2,\gamma}(\cdot/r)$ by at most $Cr^{-n_X-|\gamma|-1}\lg r$ in $\|\cdot\|_{(r)}$. The second sum equals $\int_{\mathbb T}\hat Y(k)\,\check g_{2,\gamma}(-k)\,d^\circ k$, and Lemma~\ref{IV-73-41}(c), now taken with $D=3$, $R=r$, $g=g_{2,\gamma}$, $\hat{\mathsf T}=\hat Y$, $\hat{\mathsf T}^h=\hat Y^h$, $s=s_Y$, $\iota=1$, $c_0=1$, $\mathsf H=2c|k|^{s_Y}$, $\mathsf H'=c|k|^{s_Y+1}$ (with $c$ chosen for $Y$ as it was chosen for $X$ in Step~3), $a=n_X+|\gamma|$ and $G=G_{2,\gamma}$, yields the value $r^{-n-|\gamma|}\int\hat Y^h(k)\,\hat G_{2,\gamma}(-k)\,d^\circ k+O(r^{-n-|\gamma|-1}\lg r)$, as $s_Y+n_X=n$.

In the third sum we replace $X$ and $\nabla^\gamma Y$ by $X^h$ and $\partial^\gamma Y^h$, and then the sum by the integral over $\mathbb R^3$, exactly as in Step~3; the errors are $O(r^{-n-|\gamma|-1}\lg r)$, and the substitution $z=r\zeta$ produces the factor $r^{-n-|\gamma|}$. Altogether $\nabla^\gamma(X*Y)(u)=\mathsf H_\gamma(u)+O(r^{-n-|\gamma|-1}\lg r)$, where
\[
\begin{multlined}
\mathsf H_\gamma(u):=r^{-n-|\gamma|}\Bigl[\int\hat X^h(k)\,\hat G_{1,\gamma}(-k)\,d^\circ k+\int\hat Y^h(k)\,\hat G_{2,\gamma}(-k)\,d^\circ k+\int\bigl(1-\theta(4|\zeta|)-\theta(4|\hat u \\
-\zeta|)\bigr)X^h(\zeta)\,(\partial^\gamma Y^h)(\hat u-\zeta)\,d\zeta\Bigr].
\end{multlined}
\]
For the reasons given for $\mathsf H_0$ in Step~3, the function $\mathsf H_\gamma$ is homogeneous of degree $-n-|\gamma|$ and continuous on $\mathbb R^3\setminus\{0\}$: the bracket depends on $\hat u$ only; $\hat u\mapsto G_{1,\gamma}$ and $\hat u\mapsto G_{2,\gamma}$ are continuous with respect to $\|\cdot\|_{C^{10}}$, so that the last bound of Lemma~\ref{IV-73-41}(c) applies; and the third integrand is bounded by $C\langle\zeta\rangle^{-n_X-n_Y}$ uniformly in $\hat u$, so that dominated convergence applies. Lemma~\ref{IV-73-45}, applied to $f=X*Y$, with the present $n>0$, the functions $\mathsf H_\gamma$ and $r_\gamma=200+8|\gamma|$, gives (K3) with $K^h=\mathsf H_0$ for $|u|\ge200+8|\gamma|$.

For the finitely many other $u$ with $|u|\ge|\gamma|+1$, the left side of (K3) is a finite number, bounded through the constants of $X$ and $Y$, as $\mathsf H_0$ is evaluated only at points different from $0$; so (K3) holds with a larger $c_\gamma$, which still depends only on the constants of $X$ and $Y$. This proves (b).

\emph{Step 5: proof of (c).} We have $X(x)=\iint e^{i(k'+l)\cdot x}\sigma_{\mathrm a}(k')\sigma_{\mathrm b}(l)\,d^\circ k'\,d^\circ l$, and the substitution $k=k'+l$ shows that $X$ is the kernel with the symbol $\mathsf I$. The data of Lemma~\ref{IV-73-42} are $\lambda_1(k)=k$, $\epsilon_1=-1$, $\lambda_2=0$, $\epsilon_2=1$, so that $l_1(k)=k$, $l_2=0$; the hypotheses hold with $c_\lambda=1$, the map $(l,k)\mapsto(k-l,l)$ having an integer matrix of determinant $-1$, and $\omega=3+(n_{\mathrm a}-3)+(n_{\mathrm b}-3)=n-3\ge-1$; moreover $\Lambda=1$. The bound on $\nabla^\gamma X$ follows from Lemma~\ref{IV-73-39}(a) and the discrete Leibniz rule. For $|k|\ge c_0$ the function $\mathsf I$ is bounded, as an integral with two integrable singularities at distance at least $c_0$, and so is $\tilde{\mathsf Q}$. For $0<|k|\le c_0$, Lemma~\ref{IV-73-42}(b) gives
\[
\mathsf I-\tilde{\mathsf Q}=(\mathsf Q-\tilde{\mathsf Q})+\mathsf I^h+O(|k|^{\omega+1}\lg|k|),
\]
where $|\mathsf I^h(k)|\le C|k|^\omega$, and $\mathsf Q-\tilde{\mathsf Q}=O(|k|^2)$, even $O(|k|^3)$ if $\deg\mathsf Q\le1$. This gives the bound on $\mathsf I-\tilde{\mathsf Q}$, and, for $n\le4$ (where $\omega\le1$), conditions (K1) and (K2) for $X-\mathsf q$ with $\sigma^h=\mathsf I^h$, which is continuous on $\mathbb R^3\setminus\{0\}$ by Lemma~\ref{IV-73-42}(a). For (K3): at points where no difference involves a point of $\operatorname{supp}\mathsf q$,
\[
X-K_{\mathrm a}^hK_{\mathrm b}^h=(K_{\mathrm a}-K_{\mathrm a}^h)K_{\mathrm b}+K_{\mathrm a}^h(K_{\mathrm b}-K_{\mathrm b}^h),
\]
and Lemma~\ref{IV-73-39}(a),(c) gives the required bounds; at the finitely many $x$ with $|x|\ge|\gamma|+1$ at which a difference of order $\gamma$ involves a point of $\operatorname{supp}\mathsf q$, the left side of (K3) is a finite number, which we absorb into $c_\gamma$.

\emph{Step 6: proof of (d).} Let $G\in C_c^\infty(\mathbb R^3)$ with support in $\{\frac12<|\zeta|<1\}$, let $R\ge2$ and $g:=G(\cdot/R)$ on $\mathbb Z^3$. By (K2) with $\sigma^h=0$ we have $|\sigma(k)|\le c|k|^{n-2}\lg|k|$ on the cube, and Lemma~\ref{IV-73-41}(b), with $D=3$, $s=n-2<6$, $\mathsf H=c|k|^{n-2}$, $\iota=1$ and $c_0=1$, gives, since $\|g\|_{(R)}\le\|G\|_{C^{10}}$,
\[
\sum_xK(x)g(x)=\int_{\mathbb T}\sigma(k)\,\check g(-k)\,d^\circ k=O(R^{2-n}\lg R).
\]
On the other hand, $g$ vanishes outside $\{R/2<|x|<R\}$, a set of at most $CR^3$ points; by (K3) and by Lemma~\ref{IV-73-40}(b), applied to the function $K^hG(\cdot/R)$, whose gradient is bounded by $CR^{-n-1}$,
\[
\sum_xK(x)g(x)=\sum_xK^h(x)g(x)+O(R^{2-n}\lg R)=R^{3-n}\int_{\mathbb R^3}K^hG\,d\zeta+O(R^{2-n}\lg R).
\]
Letting $R\to\infty$ we obtain $\int K^hG\,d\zeta=0$ for all such $G$; as $K^h$ is continuous and homogeneous, $K^h=0$.
\end{proof}

\begin{proposition}[(V) The kite]\label{IV-73-47}
Let $K_1,\dots,K_5$ be elementary lines with symbols $\sigma_1,\dots,\sigma_5$ and degrees $n_1,\dots,n_5$, where
\[
n_1,n_2,n_4,n_5\in\{2,3\},\quad n_3\in\{1,2,3\},\quad n_1+n_2+n_3\le7,\quad n_3+n_4+n_5\le7 ,
\]
and put $\varDelta:=\sum_{i=1}^5n_i-6$. For $u\in\mathbb Z^3$ let $r:=|u|$ and
\[
\mathsf T(u):=\sum_{z,w\in\mathbb Z^3}K_1(z)\,K_2(w)\,K_3(w-z)\,K_4(u-z)\,K_5(u-w).
\]
\begin{enumerate}
\item[(i)] There is a constant $C$, depending on $n_1,\dots,n_5$ only, such that for every $u\in\mathbb Z^3$
\[
\sum_{z,w\in\mathbb Z^3}\langle z\rangle^{-n_1}\langle w\rangle^{-n_2}\langle w-z\rangle^{-n_3}\langle u-z\rangle^{-n_4}\langle u-w\rangle^{-n_5}\le C .
\]
In particular the series defining $\mathsf T(u)$ converges absolutely.
\item[(ii)] For $\kappa=(\kappa_1,\kappa_2)\in\mathbb R^3\times\mathbb R^3$ put
\[
\begin{aligned}
\hat{\mathsf T}_{\mathsf p}(\kappa_1,\kappa_2)&:=\int_{\mathbb T}\sigma_1(\kappa_1+l)\,\sigma_2(\kappa_2-l)\,\sigma_3(l)\,d^\circ l,\\
\hat{\mathsf T}_{\mathsf q}(\kappa_1,\kappa_2)&:=\int_{\mathbb T}\sigma_4(\kappa_1-l)\,\sigma_5(\kappa_2+l)\,\sigma_3(l)\,d^\circ l .
\end{aligned}
\]
Both are functions $\mathsf I$ of the kind treated in Lemma~\ref{IV-73-42}, with $\mathsf m=3$ and with data satisfying the hypotheses of that lemma; their numbers $\omega$ are
\[
\omega_{\mathsf p}:=n_1+n_2+n_3-6\qquad\text{and}\qquad\omega_{\mathsf q}:=n_3+n_4+n_5-6 ,
\]
respectively, and both lie in $\{-1,0,1\}$. For every function $g$ of finite support on $\mathbb Z^6$
\begin{equation}\label{IV-73-48}
\begin{aligned}
\sum_{z,w}K_1(z)\,K_2(w)\,K_3(w-z)\,g(z,w)&=\int_{\mathbb T^2}\hat{\mathsf T}_{\mathsf p}(\kappa)\,\check g(-\kappa)\,d^\circ\kappa,\\
\sum_{z,w}K_4(z)\,K_5(w)\,K_3(z-w)\,g(z,w)&=\int_{\mathbb T^2}\hat{\mathsf T}_{\mathsf q}(\kappa)\,\check g(-\kappa)\,d^\circ\kappa .
\end{aligned}
\end{equation}
\item[(iii)] Let $\mathsf Q_{\mathsf p}$ and $\mathsf Q_{\mathsf q}$ be the polynomials of $\hat{\mathsf T}_{\mathsf p}$ and of $\hat{\mathsf T}_{\mathsf q}$ in the sense of Lemma~\ref{IV-73-42}(d), let
\[
\tilde{\mathsf Q}_\bullet(\kappa):=\mathsf Q_\bullet(0)+\sum_{j=1}^6\partial_j\mathsf Q_\bullet(0)\,\sin\kappa_j\qquad(\bullet\in\{\mathsf p,\mathsf q\}),
\]
and let $\mathsf q_{\mathsf p},\mathsf q_{\mathsf q}$ be the kernels on $\mathbb Z^6$ with the symbols $\tilde{\mathsf Q}_{\mathsf p},\tilde{\mathsf Q}_{\mathsf q}$, that is, $\mathsf q_\bullet(\zeta)=\int_{\mathbb T^2}e^{i\kappa\cdot\zeta}\,\tilde{\mathsf Q}_\bullet(\kappa)\,d^\circ\kappa$; they vanish off $\{0,\pm e_1,\dots,\pm e_6\}$. There are constants $r_0$ and $C$, depending on $\sigma_1,\dots,\sigma_5$ only, and a function $\mathsf U\in C^1(\mathbb R^3\setminus\{0\})$, homogeneous of degree $-\varDelta$, such that for every $u$ with $r\ge r_0$
\[
\begin{aligned}
\Bigl|\mathsf T(u)&-\sum_{z,w}\mathsf q_{\mathsf p}(z,w)\,K_4(u-z)\,K_5(u-w)\\
&-\sum_{z,w}\mathsf q_{\mathsf q}(z,w)\,K_1(u-z)\,K_2(u-w)-\mathsf U(u)\Bigr|\le C\,r^{-\varDelta-1}\lg r .
\end{aligned}
\]
\item[(iv)] The function $\mathsf U$ depends only on the leading parts $\sigma_1^h,\dots,\sigma_5^h$, and it is linear in each of them. If every $\sigma_i^h$ is replaced by $\sigma_i^h\circ O^{-1}$, with one and the same rotation $O$, then $\mathsf U$ is replaced by $\mathsf U\circ O^{-1}$.
\end{enumerate}
\end{proposition}

\begin{proof}
\emph{Step 1: convergence.} The symbol $\sigma_i$ of $K_i$ is a line symbol of degree $s_i:=n_i-3\ge-2>-3$. Lemma~\ref{IV-73-39}(a), applied with $s=s_i$ and $\gamma=0$, therefore gives $|K_i(y)|\le C\langle y\rangle^{-n_i}$ for all $y\in\mathbb Z^3$, and the absolute convergence of the series for $\mathsf T(u)$ follows from the bound in (i). To prove that bound we estimate $\langle w-z\rangle^{-n_3}\le1$; what remains factorises into
\[
\Bigl(\sum_{z}\langle z\rangle^{-n_1}\langle u-z\rangle^{-n_4}\Bigr)\Bigl(\sum_{w}\langle w\rangle^{-n_2}\langle u-w\rangle^{-n_5}\Bigr).
\]
Lemma~\ref{IV-73-40}(c) applies to the first factor with $(a,a',x)=(n_1,n_4,u)$, since $a,a'\ge2$ and $a+a'\ge4>3$; it bounds that factor by $C\langle u\rangle^{-\min\{n_1,n_4,n_1+n_4-3\}}\lg\langle u\rangle\le C\langle u\rangle^{-1}\lg\langle u\rangle\le C$, because $t^{-1}\lg t=t^{-1}(1+\log t)\le1$ for $t\ge1$. The second factor is bounded in the same way, with $(n_2,n_5)$ in place of $(n_1,n_4)$.

\emph{Step 2: the hypotheses of Lemma~\ref{IV-73-42}.} For $\hat{\mathsf T}_{\mathsf p}$ we take $\mathsf m=3$, $D=6$, the line symbols $\sigma_1,\sigma_2,\sigma_3$ of degrees $s_1,s_2,s_3$, the signs $(\epsilon_1,\epsilon_2,\epsilon_3)=(1,-1,1)$ and the linear maps $\lambda_1(\kappa)=\kappa_1$, $\lambda_2(\kappa)=\kappa_2$, $\lambda_3(\kappa)=0$. The integrand of $\hat{\mathsf T}_{\mathsf p}$ is then $\prod_i\sigma_i(\epsilon_il+\lambda_i(\kappa))$, and $(l_1,l_2,l_3)(\kappa)=(-\kappa_1,\kappa_2,0)$. Up to signs, the pairwise differences of the $l_i(\kappa)$ are $\kappa_1+\kappa_2$, $\kappa_1$ and $\kappa_2$. Hence $\max_i|l_i(\kappa)|\le|\kappa|\le2|\kappa|$ and $|\kappa|\le\sqrt2\max\{|\kappa_1|,|\kappa_2|\}\le2\max_{i<j}|l_i(\kappa)-l_j(\kappa)|$, so that the condition on the $l_i$ holds with $c_\lambda=2$; the three maps $l_i$ are pairwise different.

The map $(l,\kappa_1,\kappa_2)\mapsto(\kappa_1+l,\kappa_2-l,l)$ of $(\mathbb R^3)^3$ acts on the three components in $\mathbb R^3$ by the integer matrix
\[
\begin{pmatrix}1&1&0\\-1&0&1\\1&0&0\end{pmatrix},
\]
whose determinant is $1$. Since $n_1,n_2\ge2$ and $n_3\ge1$, we have $s_1,s_2\ge-1$ and $s_3\ge-2$, hence $s_i+s_j\ge-3$ for $i\ne j$. Finally $\omega=3+s_1+s_2+s_3=n_1+n_2+n_3-6=\omega_{\mathsf p}$, and this number lies between $2+2+1-6=-1$ and $7-6=1$. For $\hat{\mathsf T}_{\mathsf q}$ the line symbols are $\sigma_4,\sigma_5,\sigma_3$, the signs are $(-1,1,1)$ and the maps are $\kappa\mapsto\kappa_1$, $\kappa\mapsto\kappa_2$, $\kappa\mapsto0$, so that $(l_1,l_2,l_3)(\kappa)=(\kappa_1,-\kappa_2,0)$; the matrix of $(l,\kappa_1,\kappa_2)\mapsto(\kappa_1-l,\kappa_2+l,l)$ is $\bigl(\begin{smallmatrix}-1&1&0\\1&0&1\\1&0&0\end{smallmatrix}\bigr)$, again of determinant $1$, and the remaining conditions follow as before, now with $n_3+n_4+n_5\le7$.

\emph{Step 3: proof of \eqref{IV-73-48}.} We insert $K_1(z)=\int_{\mathbb T}e^{ik_1\cdot z}\sigma_1(k_1)\,d^\circ k_1$, $K_2(w)=\int_{\mathbb T}e^{ik_2\cdot w}\sigma_2(k_2)\,d^\circ k_2$ and $K_3(w-z)=\int_{\mathbb T}e^{il\cdot(w-z)}\sigma_3(l)\,d^\circ l$. The sum over $(z,w)$ is finite and the three symbols are integrable, so Fubini's theorem gives
\[
\begin{aligned}
&\sum_{z,w}K_1(z)\,K_2(w)\,K_3(w-z)\,g(z,w)\\
&\qquad=\iiint\sigma_1(k_1)\,\sigma_2(k_2)\,\sigma_3(l)\,\check g\bigl(-(k_1-l),-(k_2+l)\bigr)\,d^\circ k_1\,d^\circ k_2\,d^\circ l ,
\end{aligned}
\]
since $\sum_{z,w}g(z,w)\,e^{i(k_1-l)\cdot z+i(k_2+l)\cdot w}=\check g\bigl(-(k_1-l),-(k_2+l)\bigr)$. For fixed $l$ we substitute $k_1=\kappa_1+l$ and $k_2=\kappa_2-l$, which are translations of $\mathbb T$, and we then carry out the integration over $l$ first; this is the first identity. The function $y\mapsto K_3(-y)$ is the kernel with the symbol $l\mapsto\sigma_3(-l)$, because $K_3(-y)=\int e^{-il\cdot y}\sigma_3(l)\,d^\circ l=\int e^{il\cdot y}\sigma_3(-l)\,d^\circ l$. The derivation of the first identity used only the integrability of the three symbols. Applied to $K_4$, $K_5$ and $K_3(-\cdot)$ in place of $K_1$, $K_2$ and $K_3$, it represents $\sum_{z,w}K_4(z)K_5(w)K_3(z-w)g(z,w)$ by the function $\kappa\mapsto\int_{\mathbb T}\sigma_4(\kappa_1+l)\,\sigma_5(\kappa_2-l)\,\sigma_3(-l)\,d^\circ l$, which becomes $\hat{\mathsf T}_{\mathsf q}(\kappa)$ after the substitution $l\to-l$.

\emph{Step 4: short and long lines.} Let $\varepsilon:=\frac18$ and, for $r>0$, $\theta^r(y):=\theta(|y|/\varepsilon r)$ and $\bar\theta^r:=1-\theta^r$, as functions on $\mathbb R^3$ and, by restriction, on $\mathbb Z^3$. By Definition~\ref{IV-73-38}(1), $\theta^r(y)=1$ for $|y|\le\frac12\varepsilon r$ and $\theta^r(y)=0$ for $|y|\ge\varepsilon r$. For the arguments of the five lines we write
\[
(y_1,y_2,y_3,y_4,y_5):=(z,\,w,\,w-z,\,u-z,\,u-w).
\]
For $\mathsf S\subset\{1,\dots,5\}$ let $\mathsf T_{\mathsf S}(u)$ be the series defining $\mathsf T(u)$, with the additional factor 
\[
\prod_{i\in\mathsf S}\theta^r(y_i)\prod_{i\notin\mathsf S}\bar\theta^r(y_i)
\]
 in each term; the lines with an index in $\mathsf S$ are called short, the others long. These series converge absolutely by (i), the additional factors taking values in $[0,1]$, and multiplying out $\prod_{i=1}^5\bigl(\theta^r(y_i)+\bar\theta^r(y_i)\bigr)=1$ we obtain $\mathsf T=\sum_{\mathsf S}\mathsf T_{\mathsf S}$.

The arguments satisfy
\[
u=y_1+y_4=y_2+y_5=y_1+y_3+y_5=y_2-y_3+y_4 .
\]
If $\mathsf S$ contains one of the sets $\{1,4\}$, $\{2,5\}$, $\{1,3,5\}$, $\{2,3,4\}$, then a non-vanishing term of $\mathsf T_{\mathsf S}(u)$ requires $|y_i|<\varepsilon r$ for the indices $i$ of that set, hence $r=|u|<3\varepsilon r<r$; thus $\mathsf T_{\mathsf S}=0$. There are three possibilities for $\mathsf S\cap\{1,4\}$, three for $\mathsf S\cap\{2,5\}$ and two for $\mathsf S\cap\{3\}$; of the resulting $3\cdot3\cdot2=18$ sets, two, namely $\{1,3,5\}$ and $\{2,3,4\}$, contain one of the forbidden triples, and $16$ sets remain.

The substitutions $(z,w)\to(w,z)$ and $(z,w)\to(u-z,u-w)$ are bijections of $\mathbb Z^6$; they carry $(y_1,\dots,y_5)$ into $(y_2,y_1,-y_3,y_5,y_4)$ and into $(y_4,y_5,-y_3,y_1,y_2)$, respectively. As $\theta^r$ is even, they show that $\mathsf T_{\mathsf S}$, computed for the lines $(K_1,K_2,K_3,K_4,K_5)$, equals $\mathsf T_{\mathsf S'}$ computed for $(K_2,K_1,K_3(-\cdot),K_5,K_4)$, where $\mathsf S'$ is the image of $\mathsf S$ under the permutation $(12)(45)$, and that it equals $\mathsf T_{\mathsf S''}$ computed for $(K_4,K_5,K_3(-\cdot),K_1,K_2)$, where $\mathsf S''$ is the image of $\mathsf S$ under $(14)(25)$. The new families satisfy the hypotheses of the proposition: $K_3(-\cdot)$ is an elementary line of degree $n_3$, its symbol $\sigma_3(-k)$ being a line symbol with leading part $\sigma_3^h(-k)$ (the cut-off $\chi$ is even), and the conditions on the degrees are invariant under both permutations.

Under the group generated by $(12)(45)$ and $(14)(25)$ the sixteen sets fall into seven classes: the set $\emptyset$ alone; the sets $\{1\},\{2\},\{4\},\{5\}$; the set $\{3\}$ alone; the sets $\{1,2\},\{4,5\}$; the sets $\{1,5\},\{2,4\}$; the sets $\{1,3\},\{2,3\},\{3,4\},\{3,5\}$; and the sets $\{1,2,3\},\{3,4,5\}$. We choose the representatives
\[
\emptyset,\quad\{1\},\quad\{3\},\quad\{1,2\},\quad\{1,5\},\quad\{1,3\},\quad\{1,2,3\}.
\]
It therefore suffices to expand $\mathsf T_{\mathsf S}$ for these seven sets, for an arbitrary family of lines obeying the hypotheses; for each of the remaining nine sets we use the expansion of its representative, computed for the transformed family.

\emph{Step 5: long factors.} Let $L_i:=\bar\theta^rK_i$ on $\mathbb Z^3$, let $K_i^h:=\mathcal F^{-1}_{\ne0}\sigma_i^h$, which by Lemma~\ref{IV-73-39}(b) is smooth on $\mathbb R^3\setminus\{0\}$ and homogeneous of degree $-3-s_i=-n_i$, and let $L_i^h:=\bar\theta^rK_i^h$, extended by $0$ at the origin. Since $\bar\theta^r$ vanishes on $\{|y|\le\frac12\varepsilon r\}$, the function $L_i^h$ is smooth on $\mathbb R^3$, and by homogeneity $L_i^h(y)=r^{-n_i}\,\bar\theta^1(y/r)\,K_i^h(y/r)$. We claim that there are $r_0$ and $C$ such that for $r\ge r_0$, $i=1,\dots,5$ and $|\gamma|\le10$
\begin{equation}\label{IV-73-49}
\begin{gathered}
|\nabla^\gamma L_i(y)|+|\partial^\gamma L_i^h(y)|\le C\,(r+|y|)^{-n_i-|\gamma|},\\
|\nabla^\gamma(L_i-L_i^h)(y)|\le C\,(r+|y|)^{-n_i-1-|\gamma|};
\end{gathered}
\end{equation}
here $y\in\mathbb Z^3$, except in the term $\partial^\gamma L_i^h(y)$, where $y\in\mathbb R^3$. Indeed, Lemma~\ref{IV-73-39}(a) gives $|\nabla^\gamma K_i(x)|\le C\langle x\rangle^{-n_i-|\gamma|}$, and Lemma~\ref{IV-73-39}(c) gives $|\nabla^\gamma(K_i-K_i^h)(x)|\le C|x|^{-n_i-1-|\gamma|}$ for $|x|\ge|\gamma|+1$. The derivatives of order $\gamma$ of $K_i^h$ are homogeneous of degree $-n_i-|\gamma|$, hence bounded by $C|x|^{-n_i-|\gamma|}$; the derivatives of order $j\ge1$ of $\bar\theta^r$ are bounded by $Cr^{-j}$ and vanish outside $\{\frac12\varepsilon r\le|y|\le\varepsilon r\}$. The discrete Leibniz rule $\nabla_j(fg)=(\nabla_jf)\,g(\cdot+e_j)+f\,\nabla_jg$ distributes $\nabla^\gamma$ over the two factors, the arguments being shifted by at most $|\gamma|\le10$.

Only points $y'$ with $|y'|\ge\frac12\varepsilon r$ contribute, and at such points, for $r\ge r_0$ with $r_0$ large enough, we have $|y'|\ge|\gamma|+1$ and $r+|y|\le C|y'|$ whenever $|y-y'|\le10$. These facts together give \eqref{IV-73-49}.

\emph{Step 6: the scheme for $\mathsf S\ne\emptyset$.} For each of the six non-empty representatives we write
\[
\mathsf T_{\mathsf S}(u)=\sum_{\zeta\in\mathbb Z^D}\mathsf T^{\mathsf S}(\zeta)\,g_{\mathsf S}(\zeta),
\]
where the kernel $\mathsf T^{\mathsf S}$ on $\mathbb Z^D$ collects the short lines, and we let $s$ be the degree of the leading part $\hat{\mathsf T}^{\mathsf S,h}$ of the symbol of $\mathsf T^{\mathsf S}$. The data are as follows.
\begin{itemize}
\item $\mathsf S=\{1\}$, $D=3$, $\zeta=z$: $\mathsf T^{\mathsf S}(z)=K_1(z)$, with symbol $\sigma_1(\kappa)$ and $s=n_1-3$; and
$g_{\mathsf S}(z)=\theta^r(z)\,L_4(u-z)\sum_wL_2(w)\,L_3(w-z)\,L_5(u-w)$.
\item $\mathsf S=\{3\}$, $D=3$, $\zeta=y$: $\mathsf T^{\mathsf S}(y)=K_3(y)$, with symbol $\sigma_3(\kappa)$ and $s=n_3-3$; and
$g_{\mathsf S}(y)=\theta^r(y)\sum_zL_1(z)\,L_2(z+y)\,L_4(u-z)\,L_5(u-z-y)$.
\item $\mathsf S=\{1,2\}$, $D=6$, $\zeta=(z,w)$: $\mathsf T^{\mathsf S}(z,w)=K_1(z)K_2(w)$, with symbol $\sigma_1(\kappa_1)\sigma_2(\kappa_2)$ and $s=n_1+n_2-6$; and
$g_{\mathsf S}(z,w)=\theta^r(z)\,\theta^r(w)\,L_3(w-z)\,L_4(u-z)\,L_5(u-w)$.
\item $\mathsf S=\{1,5\}$, $D=6$, $\zeta=(z,w')$: $\mathsf T^{\mathsf S}(z,w')=K_1(z)K_5(w')$, with symbol $\sigma_1(\kappa_1)\sigma_5(\kappa_2)$ and $s=n_1+n_5-6$; and
$g_{\mathsf S}(z,w')=\theta^r(z)\,\theta^r(w')\,L_2(u-w')\,L_3(u-w'-z)\,L_4(u-z)$.
\item $\mathsf S=\{1,3\}$, $D=6$, $\zeta=(z,y)$: $\mathsf T^{\mathsf S}(z,y)=K_1(z)K_3(y)$, with symbol $\sigma_1(\kappa_1)\sigma_3(\kappa_2)$ and $s=n_1+n_3-6$; and
$g_{\mathsf S}(z,y)=\theta^r(z)\,\theta^r(y)\,L_2(z+y)\,L_4(u-z)\,L_5(u-z-y)$.
\item $\mathsf S=\{1,2,3\}$, $D=6$, $\zeta=(z,w)$: $\mathsf T^{\mathsf S}(z,w)=K_1(z)K_2(w)K_3(w-z)$, with symbol $\hat{\mathsf T}_{\mathsf p}(\kappa)$ (by \eqref{IV-73-48}, applied to $g=\delta_\zeta$) and $s=\omega_{\mathsf p}$; and
$g_{\mathsf S}(z,w)=\theta^r(z)\,\theta^r(w)\,\theta^r(w-z)\,L_4(u-z)\,L_5(u-w)$.
\end{itemize}
In the cases $\{1,5\}$ and $\{1,3\}$ we have substituted $w=u-w'$, resp.\ $w=z+y$, in the sum over $w$; in the cases $\{1\}$ and $\{3\}$ the inner sums converge absolutely by Step~1. Each $g_{\mathsf S}$ vanishes off $\{|\zeta|\le r\}$, since the short factors confine every component of $\zeta$ to $\{|\cdot|<\varepsilon r\}$ and $\sqrt2\,\varepsilon<1$.

Let $g^h_{\mathsf S}$ be the function on $\mathbb R^D$ obtained from $g_{\mathsf S}$ by replacing every $L_i$ by $L_i^h$ and, in the cases $\{1\}$ and $\{3\}$, the inner sum by the integral over $\mathbb R^3$. Since $\theta^r(y)=\theta^1(y/r)$ and $L_i^h(y)=r^{-n_i}(\bar\theta^1K_i^h)(y/r)$, the substitution of $r$ times the variable in the inner integrals gives
\[
g^h_{\mathsf S}(\zeta)=r^{-a}\,G_{\mathsf S}(\zeta/r),\qquad\text{where}\quad a:=\begin{cases}\sum_{i\notin\mathsf S}n_i&\text{if }D=6,\\[2pt]\sum_{i\notin\mathsf S}n_i-3&\text{if }\mathsf S=\{1\},\{3\},\end{cases}
\]
so that $a+s=\varDelta$ in all six cases. Here $G_{\mathsf S}$ is the function $g^h_{\mathsf S}$ computed with $r=1$ and with $\hat u:=u/r$ in place of $u$; it belongs to $C_c^\infty(\mathbb R^D)$ and has its support in $\{|\zeta|\le1\}$. We claim
\begin{equation}\label{IV-73-50}
\|g_{\mathsf S}-g^h_{\mathsf S}\|_{(r)}\le C\,r^{-a-1}\qquad(r\ge r_0),
\end{equation}
where $g^h_{\mathsf S}$ is restricted to $\mathbb Z^D$ and $\|\cdot\|_{(r)}$ is the norm of Definition~\ref{IV-73-38}(5).

In the cases $\{1,2\}$, $\{1,5\}$, $\{1,3\}$, $\{1,2,3\}$ the functions $g_{\mathsf S}$ and $g^h_{\mathsf S}$ are products of the same short factors with the long factors $L_i$, resp.\ $L_i^h$, $i\notin\mathsf S$, evaluated at linear functions of $\zeta$. We write the difference as a sum of terms in each of which one factor $L_i$ has been replaced by $L_i-L_i^h$, and we distribute the differences $\nabla^\gamma$, $|\gamma|\le10$, by the Leibniz rule. By \eqref{IV-73-49} a difference falling on a long factor gains a factor $C(r+|y_i|)^{-1}\le Cr^{-1}$, a difference falling on a short factor gains $Cr^{-1}$, and the factor $L_i-L_i^h$ gains $(r+|y_i|)^{-1}\le r^{-1}$ in comparison with $L_i$; moreover $\prod_{i\notin\mathsf S}(r+|y_i|)^{-n_i}\le r^{-a}$. This gives $r^{|\gamma|}|\nabla^\gamma(g_{\mathsf S}-g^h_{\mathsf S})|\le Cr^{-a-1}$.

In the case $\mathsf S=\{1\}$ let
\[
\begin{aligned}
\mathsf w(z)&:=\sum_wL_2(w)\,L_3(w-z)\,L_5(u-w),\\
\mathsf w^h(z)&:=\int_{\mathbb R^3}L_2^h(w)\,L_3^h(w-z)\,L_5^h(u-w)\,dw .
\end{aligned}
\]
For $|z|\le r$ the differences in $z$ fall on $L_3(w-z)$ only. By \eqref{IV-73-49}, by $r+|u-w|\ge\frac12(r+|w|)$ (second assertion of Lemma~\ref{IV-73-40}(a), as $|u|=r$) and by the first assertion of Lemma~\ref{IV-73-40}(a) with $D=3$, $R=r$ and exponent $n_2+n_5\ge4>3$, we obtain
\[
|\nabla^\gamma\mathsf w(z)|\le Cr^{-n_3-|\gamma|}\sum_w(r+|w|)^{-n_2-n_5}\le Cr^{3-n_2-n_3-n_5-|\gamma|}.
\]
Each replacement of an $L_i$ by $L_i^h$ in $\mathsf w$ costs a factor $Cr^{-1}$, by the second inequality of \eqref{IV-73-49}. The replacement of the sum over $w$ by the integral costs a factor $Cr^{-1}$ as well: we apply Lemma~\ref{IV-73-40}(b), with $D=3$, to $F(w):=L_2^h(w)\,(\nabla_z^\gamma L_3^h)(w-z)\,L_5^h(u-w)$, whose gradient is bounded by $Cr^{-n_3-|\gamma|-1}(r+|w|)^{-n_2-n_5}$ on each cube $w+[0,1]^3$. Since $g_{\{1\}}(z)=\theta^r(z)L_4(u-z)\mathsf w(z)$, since $|\nabla^\gamma L_4(u-\cdot)|\le Cr^{-n_4-|\gamma|}$ on $\{|z|\le r\}$ by \eqref{IV-73-49}, and since the differences of order $\gamma$ of $\theta^r$ are bounded by $Cr^{-|\gamma|}$, the Leibniz rule gives \eqref{IV-73-50} with $a=n_2+n_3+n_4+n_5-3$.

In the case $\mathsf S=\{3\}$ let $\mathsf w(y):=\sum_zL_1(z)\,L_2(z+y)\,L_4(u-z)\,L_5(u-z-y)$, and let $\mathsf w^h$ be the integral over $\mathbb R^3$ of the corresponding product of the $L_i^h$. For $|y|\le r$ the differences in $y$ fall on $L_2$ and on $L_5$ (and, in $g_{\{3\}}$, on the factor $\theta^r(y)$), each difference gaining a factor $Cr^{-1}$. By Lemma~\ref{IV-73-40}(a) we have $r+|z+y|\ge\frac12(r+|z|)$, $r+|u-z|\ge\frac12(r+|z|)$, and, with $R=2r\ge|u-y|$, also $r+|u-y-z|\ge\frac14(r+|z|)$; the first assertion of that lemma, with exponent $n_1+n_2+n_4+n_5\ge8$, gives $\sum_z(r+|z|)^{-n_1-n_2-n_4-n_5}\le Cr^{3-n_1-n_2-n_4-n_5}$. Hence $|\nabla^\gamma\mathsf w(y)|\le Cr^{3-n_1-n_2-n_4-n_5-|\gamma|}$, and, as in the case $\{1\}$, the replacements $L_i\to L_i^h$ and the replacement of the sum by the integral (Lemma~\ref{IV-73-40}(b)) cost a factor $Cr^{-1}$ each. This proves \eqref{IV-73-50} in all six cases.

\emph{Step 7: symbols, and the pairing with slowly varying functions.} We verify the hypotheses of Lemma~\ref{IV-73-41}(c) for $\hat{\mathsf T}$, the symbol of $\mathsf T^{\mathsf S}$, and $\hat{\mathsf T}^h:=\hat{\mathsf T}^{\mathsf S,h}$.
In the cases $\{1\}$ and $\{3\}$ ($D=3$, $i=1$, resp.\ $i=3$, $s=s_i$) we take $c_0=1$. On $\{|\kappa|\le1\}$ we have $\chi=1$, so $\sigma_i-\chi\sigma_i^h\in\mathcal S^{s+1}$ gives $|\sigma_i-\sigma_i^h|\le C|\kappa|^{s+1}$ there; the majorants are $\mathsf H:=C|\kappa|^s\in\mathcal H_s(3)$ for $\sigma_i^h$ and $\mathsf H':=C|\kappa|^{s+1}\in\mathcal H_{s+1}(3)$, which belong to these classes since $s>-3$; the required bound $|\sigma_i-\sigma_i^h|\le\mathsf H'\lg|\kappa|$ holds on $\{|\kappa|\le1\}$, as $\lg|\kappa|\ge1$.
In the cases $\{1,2\}$, $\{1,5\}$, $\{1,3\}$ ($D=6$) the symbol is $\sigma_i(\kappa_1)\sigma_j(\kappa_2)$ with $(i,j)=(1,2),(1,5),(1,3)$, its leading part is $\sigma_i^h(\kappa_1)\sigma_j^h(\kappa_2)$, and $s=s_i+s_j$ with $s_i,s_j>-3$. Then $|\sigma_i^h(\kappa_1)\sigma_j^h(\kappa_2)|\le C|\kappa_1|^{s_i}|\kappa_2|^{s_j}$, a function in $\mathcal H_{s_i+s_j}(6)$, and on $\{|\kappa|\le1\}$
\[
\bigl|\sigma_i(\kappa_1)\sigma_j(\kappa_2)-\sigma_i^h(\kappa_1)\sigma_j^h(\kappa_2)\bigr|\le C\bigl(|\kappa_1|^{s_i+1}|\kappa_2|^{s_j}+|\kappa_1|^{s_i}|\kappa_2|^{s_j+1}\bigr),
\]
a function in $\mathcal H_{s_i+s_j+1}(6)$; here we wrote the difference as $(\sigma_i-\sigma_i^h)\sigma_j+\sigma_i^h(\sigma_j-\sigma_j^h)$ and used $|\sigma_j(\kappa_2)|\le C|\kappa_2|^{s_j}$.
In the case $\{1,2,3\}$ we use Lemma~\ref{IV-73-42}, whose hypotheses hold by Step~2, with the number $c_0$ given there. By part~(b) of that lemma, $|\hat{\mathsf T}_{\mathsf p}-\mathsf Q_{\mathsf p}-\mathsf I^h_{\mathsf p}|\le C|\kappa|^{\omega_{\mathsf p}+1}\lg|\kappa|\,\Lambda(\kappa)$ for almost every $\kappa$ with $|\kappa|\le c_0$, where $\mathsf I^h_{\mathsf p}$ is the homogeneous part of $\hat{\mathsf T}_{\mathsf p}$. The polynomial $\mathsf Q_{\mathsf p}$ has degree at most $\omega_{\mathsf p}\le1$, so $\mathsf Q_{\mathsf p}-\tilde{\mathsf Q}_{\mathsf p}=\sum_j\partial_j\mathsf Q_{\mathsf p}(0)(\kappa_j-\sin\kappa_j)=O(|\kappa|^3)$, and $|\kappa|^3\le|\kappa|^{\omega_{\mathsf p}+1}\lg|\kappa|\,\Lambda(\kappa)$ for $|\kappa|\le1$.

We therefore take $\hat{\mathsf T}:=\hat{\mathsf T}_{\mathsf p}-\tilde{\mathsf Q}_{\mathsf p}$, which is integrable on $\mathbb T^2$: Lemma~\ref{IV-73-42} bounds the $L^1$ norm of $\hat{\mathsf T}_{\mathsf p}$ by the product of the $L^1$ norms of $\sigma_1,\sigma_2,\sigma_3$, and $\tilde{\mathsf Q}_{\mathsf p}$ is a trigonometric polynomial. We further take $\hat{\mathsf T}^h:=\mathsf I^h_{\mathsf p}$ and the majorants $\mathsf H:=C|\kappa|^{\omega_{\mathsf p}}\Lambda$ (Lemma~\ref{IV-73-42}(a)) and $\mathsf H':=C|\kappa|^{\omega_{\mathsf p}+1}\Lambda$. They belong to $\mathcal H_{\omega_{\mathsf p}}(6)$ and to $\mathcal H_{\omega_{\mathsf p}+1}(6)$, respectively, since the same lemma asserts that $|\kappa|^{s}\Lambda(\kappa)$ belongs to $\mathcal H_s(6)$ for every real $s>-6$, and since $\omega_{\mathsf p}\ge-1$.
In all cases $-D<s<8-D$, as Lemma~\ref{IV-73-41}(c) requires: for $D=3$ we have $s\in\{-2,-1,0\}\subset(-3,2)$, and for $D=6$ we have $s\in\{-3,\dots,1\}\subset(-6,2)$.

Inserting the Fourier integrals of the short lines, we have $\mathsf T_{\mathsf S}(u)=\int_{\mathbb T^{D/3}}\hat{\mathsf T}^{\mathsf S}(\kappa)\,\check g_{\mathsf S}(-\kappa)\,d^\circ\kappa$, where $\hat{\mathsf T}^{\mathsf S}$ is the symbol of $\mathsf T^{\mathsf S}$: for the product kernels this follows from Fubini's theorem as in Step~3, and for $\{1,2,3\}$ it is \eqref{IV-73-48}. We apply Lemma~\ref{IV-73-41}(c) with $R=r$, $g=g_{\mathsf S}$, which vanishes off $\{|\zeta|\le r\}$, $G=G_{\mathsf S}$, the number $a$ of Step~6, and $A=C$, which is admissible by \eqref{IV-73-50} because $\lg r\ge1$. The norm $\|G_{\mathsf S}\|_{C^{10}}$ is bounded uniformly in $\hat u\in S^2$, by the smoothness in $(\zeta,\hat u)$ established in Step~9 and the compactness of $S^2$. In the first five cases, $\{1\}$, $\{3\}$, $\{1,2\}$, $\{1,5\}$, $\{1,3\}$, we obtain
\[
\mathsf T_{\mathsf S}(u)=\mathsf U_{\mathsf S}(u)+O\bigl(r^{-\varDelta-1}\lg r\bigr),\qquad\mathsf U_{\mathsf S}(u):=r^{-\varDelta}\int_{\mathbb R^D}\hat{\mathsf T}^{\mathsf S,h}(\kappa)\,\hat G_{\mathsf S}(-\kappa)\,d^\circ\kappa ,
\]
because $a+s=\varDelta$. In the case $\{1,2,3\}$, where $\hat{\mathsf T}^{\mathsf S,h}:=\mathsf I^h_{\mathsf p}$, the lemma is applied to $\hat{\mathsf T}_{\mathsf p}-\tilde{\mathsf Q}_{\mathsf p}$, and the remaining term is $\int_{\mathbb T^2}\tilde{\mathsf Q}_{\mathsf p}(\kappa)\check g_{\mathsf S}(-\kappa)\,d^\circ\kappa=\sum_{z,w}\mathsf q_{\mathsf p}(z,w)\,g_{\mathsf S}(z,w)$; thus
\[
\mathsf T_{\{1,2,3\}}(u)=\sum_{z,w}\mathsf q_{\mathsf p}(z,w)\,g_{\{1,2,3\}}(z,w)+\mathsf U_{\{1,2,3\}}(u)+O\bigl(r^{-\varDelta-1}\lg r\bigr).
\]
On the support of $\mathsf q_{\mathsf p}$, that is, for $(z,w)\in\{0,\pm e_1,\dots,\pm e_6\}$, we have $|z|,|w|,|w-z|\le1\le\frac12\varepsilon r$ and $|u-z|,|u-w|\ge r-1\ge\varepsilon r$ as soon as $r\ge r_0\ge16$; hence $g_{\{1,2,3\}}(z,w)=K_4(u-z)K_5(u-w)$ there, which produces the first local term in (iii). The image of $\{1,2,3\}$ under $(z,w)\to(u-z,u-w)$ is $\{3,4,5\}$, and the transformed family $(K_4,K_5,K_3(-\cdot),K_1,K_2)$ has, in place of $\hat{\mathsf T}_{\mathsf p}$, the function $\int_{\mathbb T}\sigma_4(\kappa_1+l)\sigma_5(\kappa_2-l)\sigma_3(-l)\,d^\circ l=\hat{\mathsf T}_{\mathsf q}(\kappa)$; its local kernel is therefore $\mathsf q_{\mathsf q}$, and $\mathsf T_{\{3,4,5\}}$ produces the second local term in (iii).

\emph{Step 8: the case $\mathsf S=\emptyset$.} Here $\mathsf T_\emptyset(u)=\sum_{z,w}\prod_{i=1}^5L_i(y_i)$. By \eqref{IV-73-49} with $\gamma=0$, by $(r+|y_3|)^{-n_3}\le r^{-n_3}$ and by Lemma~\ref{IV-73-40}(a) (as in Step~6, with $r+|u-z|\ge\frac12(r+|z|)$ and $r+|u-w|\ge\frac12(r+|w|)$),
\[
\sum_{z,w}\prod_{i=1}^5(r+|y_i|)^{-n_i}\le C\,r^{-n_3}\,r^{3-n_1-n_4}\,r^{3-n_2-n_5}=C\,r^{-\varDelta}.
\]
Replacing each $L_i$ by $L_i^h$ costs a factor $Cr^{-1}$, by the second inequality of \eqref{IV-73-49}; replacing the sum over $(z,w)\in\mathbb Z^6$ by the integral over $\mathbb R^6$ costs a factor $Cr^{-1}$ by Lemma~\ref{IV-73-40}(b) with $D=6$, applied to $F(z,w):=\prod_iL_i^h(y_i)$, whose gradient is bounded by $Cr^{-1}\prod_i(r+|y_i|)^{-n_i}$ on each cube of side $1$. Hence $\mathsf T_\emptyset(u)=\mathsf U_\emptyset(u)+O(r^{-\varDelta-1})$ with
\[
\mathsf U_\emptyset(u):=\iint_{\mathbb R^3\times\mathbb R^3}\prod_{i=1}^5L_i^h(y_i)\,dz\,dw ,
\]
which, after the substitution $(z,w)=r(z',w')$, is seen to be homogeneous of degree $6-\sum_in_i=-\varDelta$.

\emph{Step 9: the function $\mathsf U$.} Let $\mathsf U$ be the sum of the sixteen functions $\mathsf U_{\mathsf S}$, where for a set $\mathsf S$ which is not a representative, $\mathsf U_{\mathsf S}$ is the function of its representative computed for the transformed family (Step~4). By Steps~4, 7 and~8, the inequality in (iii) holds for $r\ge r_0$. Every $\mathsf U_{\mathsf S}$ is homogeneous of degree $-\varDelta$, and so is $\mathsf U$. We let a vector $v\in\mathbb R^3$ take the place of $\hat u$ in the definition of $G_{\mathsf S}$, and we write $G_{\mathsf S}(\zeta;v)$. In the factors $\bar\theta^1K_i^h$ the argument of $K_i^h$ has modulus at least $\frac12\varepsilon$ wherever $\bar\theta^1$ does not vanish, so these factors are smooth on $\mathbb R^3$; the inner integrals in the cases $\{1\}$ and $\{3\}$ may be differentiated under the integral sign, with integrable majorants that are uniform for $(\zeta,v)$ in compact sets.

Hence $(\zeta,v)\mapsto G_{\mathsf S}(\zeta;v)$ is smooth, with support in $\{|\zeta|\le1\}$ for every $v$, and $v\mapsto G_{\mathsf S}(\cdot\,;v)$ is of class $C^1$ with respect to the norm $\|\cdot\|_{C^{10}}$. The last bound of Lemma~\ref{IV-73-41}(c), $|\int\hat{\mathsf T}^h\hat G(-\kappa)\,d^\circ\kappa|\le C\|G\|_{C^{10}}$, which is linear in $G$, then shows that $\Phi_{\mathsf S}(v):=\int\hat{\mathsf T}^{\mathsf S,h}(\kappa)\,\hat G_{\mathsf S}(-\kappa;v)\,d^\circ\kappa$ is of class $C^1$ in $v$. Since $\mathsf U_{\mathsf S}(u)=|u|^{-\varDelta}\,\Phi_{\mathsf S}(u/|u|)$, the function $\mathsf U_{\mathsf S}$ is of class $C^1$ on $\mathbb R^3\setminus\{0\}$.

For $\mathsf U_\emptyset$ we differentiate under the integral sign.

The functions $\mathsf U_{\mathsf S}$ are built from the leading parts $\sigma_i^h$ of the short lines, from the functions $K_i^h=\mathcal F^{-1}_{\ne0}\sigma_i^h$ of the long lines, which depend linearly on $\sigma_i^h$ (Lemma~\ref{IV-73-39}(b)), from the homogeneous parts $\mathsf I^h_{\mathsf p}$ and $\mathsf I^h_{\mathsf q}$, which depend only on the leading parts and multilinearly on them (Lemma~\ref{IV-73-42}(a)), and from radial cut-offs; every line enters each $\mathsf U_{\mathsf S}$ exactly once. Hence $\mathsf U$ depends only on $\sigma_1^h,\dots,\sigma_5^h$ and is linear in each of them. If every $\sigma_i^h$ is replaced by $\sigma_i^h\circ O^{-1}$, then $K_i^h$ is replaced by $K_i^h\circ O^{-1}$ (Lemma~\ref{IV-73-39}(b)), and $\mathsf I^h_{\mathsf p}$, $\mathsf I^h_{\mathsf q}$ are replaced by their compositions with $O^{-1}$ acting on both arguments, since the linear maps $\lambda_i$ of Step~2 commute with this action (Lemma~\ref{IV-73-42}(a)).

The substitutions $\zeta\to O\zeta$ in the integrals defining $G_{\mathsf S}$ and $\mathsf U_\emptyset$, and $\kappa\to O\kappa$ in the integral defining $\mathsf U_{\mathsf S}$, together with the invariance of the radial cut-offs, show that $\mathsf U$ is replaced by $\mathsf U\circ O^{-1}$. This proves (iii) and (iv).
\end{proof}

\subsection{Classes of pairings; convergence}

\begin{lemma}[The lines of the lattice]\label{IV-73-51}
Let $\tau_{\mathrm{gf}}>0$; we work on $\mathbb Z^3$. For $\iota\in\{1,2,3\}$ write $p_{[\iota]}=p_{\alpha\beta}$ with $\alpha<\beta$, and put
\[
\mathsf c_{\iota,\nu}(k):=\mathsf k_\alpha(k)\,\delta_{\nu\beta}-\mathsf k_\beta(k)\,\delta_{\nu\alpha}\qquad(\nu\in\{1,2,3\},\ k\in\mathbb R^3),
\]
so that for $v\in\mathbb C^3$ the form $(x,\nu)\mapsto e^{ik\cdot x}v_\nu$ satisfies $\mathrm d(e^{ik\cdot x}v_\nu)(p_{[\iota]}(x))=e^{ik\cdot x}\sum_\nu\mathsf c_{\iota,\nu}(k)\,v_\nu$. Let $\hat{\mathsf C}_{\mu\nu}(k)$ be the bracket in the formula for $\mathsf C^{\mathrm{lat}}_{\tau_{\mathrm{gf}}}$ in Definition~\ref{IV-73-15}(1),
\[
\hat{\mathsf C}_{\mu\nu}(k)=\frac{\delta_{\mu\nu}}{|\mathsf k|^2}+(\tau_{\mathrm{gf}}-1)\,\frac{\mathsf k_\mu\overline{\mathsf k_\nu}}{|\mathsf k|^4}.
\]
Then for every $x\in\mathbb Z^3$, all directions $\mu,\nu$ and all $\iota,\iota'$, with $p_{[\iota']}=p_{\alpha'\beta'}$, $\alpha'<\beta'$, the following functions of $y\in\mathbb Z^3$ are elementary lines, with the symbols, the degrees $n$ and the leading parts of the symbols listed:
\begin{itemize}
\item $y\mapsto\mathsf C\bigl((x,\mu),(x+y,\nu)\bigr)$: symbol $\hat{\mathsf C}_{\nu\mu}(k)$; $n=1$; leading part $\delta_{\mu\nu}|k|^{-2}+(\tau_{\mathrm{gf}}-1)\,k_\mu k_\nu|k|^{-4}$;
\item $y\mapsto\mathsf E_{p_{[\iota]}(x)}\bigl((x+y,\nu)\bigr)$: symbol $\sum_\mu\overline{\mathsf c_{\iota,\mu}(k)}\,\hat{\mathsf C}_{\nu\mu}(k)=\overline{\mathsf c_{\iota,\nu}(k)}\,|\mathsf k|^{-2}$; $n=2$; leading part $-i(k_\alpha\delta_{\nu\beta}-k_\beta\delta_{\nu\alpha})|k|^{-2}$;
\item $y\mapsto\Pi\bigl(p_{[\iota]}(x),p_{[\iota']}(x+y)\bigr)$: symbol $\sum_\nu\overline{\mathsf c_{\iota,\nu}(k)}\,\mathsf c_{\iota',\nu}(k)\,|\mathsf k|^{-2}$; $n=3$; leading part $\sum_\nu(k_\alpha\delta_{\nu\beta}-k_\beta\delta_{\nu\alpha})(k_{\alpha'}\delta_{\nu\beta'}-k_{\beta'}\delta_{\nu\alpha'})|k|^{-2}$;
\item $y\mapsto\mathsf D\bigl((x,\mu);(x+y,\nu)\bigr)$: symbol $\frac12(1+e^{ik_\nu})\,\overline{\mathsf k_\mu(k)}\,|\mathsf k|^{-2}$; $n=2$; leading part $-ik_\mu|k|^{-2}$.
\end{itemize}
If $K(y)$ is an elementary line with symbol $\sigma(k)$ and $y_0\in\mathbb Z^3$, then $K(-y)$ and $K(y+y_0)$ are elementary lines of the same degree, with the symbols $\sigma(-k)$ and $e^{ik\cdot y_0}\sigma(k)$. Consequently, the contraction \eqref{IV-73-20} of a leg at a bond or a plaquette with base point $x+\xi$ with a leg at a bond or a plaquette with base point $x+y+\xi'$ equals $\delta_{aa'}$ times an elementary line in $y$ of degree $1+\#\{\text{curl legs among the two}\}$.
\end{lemma}

\begin{proof}
The matrix $\hat{\mathsf C}(k)$ is Hermitian, and $\hat{\mathsf C}(-k)=\overline{\hat{\mathsf C}(k)}$ because $\mathsf k_\mu(-k)=\overline{\mathsf k_\mu(k)}$. By Definition~\ref{IV-73-15}(1) and the substitution $k\to-k$,
\[
\begin{aligned}
\mathsf C\bigl((x,\mu),(x+y,\nu)\bigr)&=\int_{\mathbb T}e^{-ik\cdot y}\,\hat{\mathsf C}_{\mu\nu}(k)\,d^\circ k\\
&=\int_{\mathbb T}e^{ik\cdot y}\,\overline{\hat{\mathsf C}_{\mu\nu}(k)}\,d^\circ k=\int_{\mathbb T}e^{ik\cdot y}\,\hat{\mathsf C}_{\nu\mu}(k)\,d^\circ k .
\end{aligned}
\]
The plaquette $p=p_{[\iota]}(x)=p_{\alpha\beta}(x)$ has the bonds $(x,\alpha),(x+e_\alpha,\beta),(x+e_\beta,\alpha),(x,\beta)$, with the signs $+,+,-,-$. By the formula just proved, $\mathsf C\bigl((x+\xi,\mu'),(x+y,\nu)\bigr)=\int e^{ik\cdot(y-\xi)}\hat{\mathsf C}_{\nu\mu'}(k)\,d^\circ k$, and therefore
\[
\begin{aligned}
\mathsf E_p\bigl((x+y,\nu)\bigr)&=\sum_is_i\,\mathsf C\bigl(b_i(p),(x+y,\nu)\bigr)\\
&=\int_{\mathbb T}e^{ik\cdot y}\Bigl[\hat{\mathsf C}_{\nu\alpha}(k)\bigl(1-e^{-ik_\beta}\bigr)+\hat{\mathsf C}_{\nu\beta}(k)\bigl(e^{-ik_\alpha}-1\bigr)\Bigr]\,d^\circ k .
\end{aligned}
\]
As $\overline{\mathsf k_\mu(k)}=e^{-ik_\mu}-1$, the bracket equals $\sum_\mu\overline{\mathsf c_{\iota,\mu}(k)}\,\hat{\mathsf C}_{\nu\mu}(k)$. In this sum the part of $\hat{\mathsf C}_{\nu\mu}$ containing $\mathsf k_\nu\overline{\mathsf k_\mu}$ drops out, since $\sum_\mu\overline{\mathsf c_{\iota,\mu}}\,\mathsf k_\nu\overline{\mathsf k_\mu}=\mathsf k_\nu\,\overline{\sum_\mu\mathsf c_{\iota,\mu}\mathsf k_\mu}$ and $\sum_\mu\mathsf c_{\iota,\mu}\mathsf k_\mu=\mathsf k_\alpha\mathsf k_\beta-\mathsf k_\beta\mathsf k_\alpha=0$; what remains is $\overline{\mathsf c_{\iota,\nu}}\,|\mathsf k|^{-2}$. The kernel $\Pi$ arises from $\mathsf E_p$ by the operation $\mathrm d$ in the second argument (Definition~\ref{IV-73-15}(2)).

With $q=p_{[\iota']}(x+y)=p_{\alpha'\beta'}(x+y)$, the four bonds $(x+y,\alpha')$, $(x+y+e_{\alpha'},\beta')$, $(x+y+e_{\beta'},\alpha')$, $(x+y,\beta')$ of $q$ contribute, with the signs $+,+,-,-$, the symbols $\overline{\mathsf c_{\iota,\alpha'}}|\mathsf k|^{-2}$, $e^{ik_{\alpha'}}\overline{\mathsf c_{\iota,\beta'}}|\mathsf k|^{-2}$, $e^{ik_{\beta'}}\overline{\mathsf c_{\iota,\alpha'}}|\mathsf k|^{-2}$, $\overline{\mathsf c_{\iota,\beta'}}|\mathsf k|^{-2}$, and their sum is
\[
\bigl(\mathsf k_{\alpha'}\,\overline{\mathsf c_{\iota,\beta'}}-\mathsf k_{\beta'}\,\overline{\mathsf c_{\iota,\alpha'}}\bigr)|\mathsf k|^{-2}=\sum_\nu\overline{\mathsf c_{\iota,\nu}}\,\mathsf c_{\iota',\nu}\,|\mathsf k|^{-2};
\]
that is, $\mathrm d$ in the second argument multiplies the symbol by $\mathsf c_{\iota',\nu}(k)$ and sums over $\nu$. Finally, with $\mathsf G(x',x'')=\int e^{ik\cdot(x'-x'')}|\mathsf k|^{-2}\,d^\circ k$,
\[
\begin{aligned}
\mathsf D\bigl((x,\mu);(x+y,\nu)\bigr)&=\tfrac12\sum_{x'\in\{x+y,\,x+y+e_\nu\}}\bigl[\mathsf G(x',x+e_\mu)-\mathsf G(x',x)\bigr]\\
&=\int_{\mathbb T}e^{ik\cdot y}\,\tfrac12\bigl(1+e^{ik_\nu}\bigr)\bigl(e^{-ik_\mu}-1\bigr)\,|\mathsf k|^{-2}\,d^\circ k .
\end{aligned}
\]

All four symbols are sums of products of the examples of Definition~\ref{IV-73-38}(3), namely $e^{ik\cdot y}$ (degree $0$, leading part $1$), $\mathsf k_\mu$ and $\overline{\mathsf k_\mu}$ (degree $1$, leading parts $ik_\mu$ and $-ik_\mu$) and $|\mathsf k|^{-2}$ (degree $-2$, leading part $|k|^{-2}$), and no product of degree below $-2$ occurs: the term $\mathsf k_\nu\overline{\mathsf k_\mu}|\mathsf k|^{-4}$ is the product $(\mathsf k_\nu|\mathsf k|^{-2})(\overline{\mathsf k_\mu}|\mathsf k|^{-2})$ of two line symbols of degree $-1$. By the rules for sums and products in Definition~\ref{IV-73-38}(3), the four symbols are line symbols of the degrees $-2,-1,0,-1$, so the kernels are elementary lines of the degrees $1,2,3,2$; their leading parts arise by replacing $\mathsf k_\mu$, $\overline{\mathsf k_\mu}$, $|\mathsf k|^{-2}$, $e^{ik_\nu}$ by $ik_\mu$, $-ik_\mu$, $|k|^{-2}$, $1$, which gives the leading parts listed.

If $K$ has the symbol $\sigma$, then $K(-y)=\int e^{-ik\cdot y}\sigma(k)\,d^\circ k=\int e^{ik\cdot y}\sigma(-k)\,d^\circ k$ and $K(y+y_0)=\int e^{ik\cdot y}e^{ik\cdot y_0}\sigma(k)\,d^\circ k$. The function $\sigma(-k)$ is a line symbol of the same degree, with leading part $\sigma^h(-k)$, as $\chi$ is even; and $e^{ik\cdot y_0}\sigma(k)$ is the product of $\sigma$ with a line symbol of degree $0$ and leading part $1$. For the consequence we recall from \eqref{IV-73-20} that the contraction of two plain legs is $\delta_{aa'}\mathsf C$, that of a curl leg with a plain leg is $\delta_{aa'}\mathsf E$, and that of two curl legs is $\delta_{aa'}\Pi$. These kernels are invariant under simultaneous translations of their arguments, by their Fourier representations. Hence, by the list above and by the two rules just proved (which take care of the order of the legs and of the shifts $\xi,\xi'$), each contraction is an elementary line in $y$ of degree $1$, $2$, $3$, respectively.
\end{proof}

\begin{lemma}[$\Pi^{\mathrm{lat}}$ is the curvature projection $\Pi_\infty$ up to signs]\label{IV-73-52}
For a plaquette $p=p_{\mu\nu}(x)$, $\mu<\nu$, of $\mathbb Z^3$, let $\epsilon_p=-1$ for $(\mu,\nu)=(1,3)$ and $\epsilon_p=1$ otherwise. This converts the positive orientation to the cyclic orientation indexed by the dual direction. Write $p^\flat=(x,\lambda(p))$ in these cyclic coordinates, and let $\Pi_\infty$ be the orthogonal projection onto the closed range of curl in $\ell^2$ of cyclic plaquette cochains. Its Fourier formula is equation~\YMMref{(2242)}{822}. Then for all plaquettes $p,q$ of $\mathbb Z^3$
\[
\Pi^{\mathrm{lat}}(p,q)=\epsilon_p\,\epsilon_q\,\Pi_\infty(p^\flat,q^\flat),\qquad\text{and hence}\qquad\mathsf B_0^{\mathrm{lat}}(p,q)=\tfrac{\mathsf n}2\,\Pi_\infty(p^\flat,q^\flat)^2 .
\]
\end{lemma}

\begin{proof}
\emph{Step 1.} Write $p=p_{[\iota]}(x)$ and $q=p_{[\iota']}(x+y)$, and let $\epsilon_\iota:=-1$ for $\iota=2$ and $\epsilon_\iota:=1$ for $\iota\in\{1,3\}$. Since $p_{[2]}=p_{13}$, we have $\epsilon_p=\epsilon_\iota$ and $\epsilon_q=\epsilon_{\iota'}$. By the entry for $\Pi$ in Lemma~\ref{IV-73-51},
\[
\Pi^{\mathrm{lat}}(p,q)=\int_{\mathbb T}e^{ik\cdot y}\sum_\nu\overline{\mathsf c_{\iota,\nu}(k)}\,\mathsf c_{\iota',\nu}(k)\,|\mathsf k|^{-2}\,d^\circ k .
\]

\emph{Step 2.} We show that
\[
\sum_\nu\overline{\mathsf c_{\iota,\nu}}\,\mathsf c_{\iota',\nu}=\epsilon_\iota\,\epsilon_{\iota'}\bigl(\delta_{\iota\iota'}\,|\mathsf k|^2-\mathsf k_\iota\,\overline{\mathsf k_{\iota'}}\bigr).
\]
For $\iota=\iota'$, with $p_{[\iota]}=p_{\alpha\beta}$, the left side is $|\mathsf k_\alpha|^2+|\mathsf k_\beta|^2=|\mathsf k|^2-|\mathsf k_\iota|^2$, and $\epsilon_\iota^2=1$. Let $\iota\ne\iota'$, and let $\iota''$ be the third direction. As $\mathsf c_{\iota,\nu}=0$ for $\nu=\iota$ and $\mathsf c_{\iota',\nu}=0$ for $\nu=\iota'$, only $\nu=\iota''$ contributes. By definition $\mathsf c_{\iota,\iota''}=\pm\mathsf k_{\iota'}$, with the sign $+$ if and only if $\iota'<\iota''$, and $\mathsf c_{\iota',\iota''}=\pm\mathsf k_\iota$, with the sign $+$ if and only if $\iota<\iota''$. So the left side is $\pm\overline{\mathsf k_{\iota'}}\,\mathsf k_\iota$, with the sign $-$ if and only if $\iota''$ lies between $\iota$ and $\iota'$, that is, if and only if $\iota''=2$. On the right side the factor of $\mathsf k_\iota\overline{\mathsf k_{\iota'}}$ is $-\epsilon_\iota\epsilon_{\iota'}$, which equals $-1$ if and only if neither $\iota$ nor $\iota'$ is $2$, that is, if and only if $\iota''=2$.

\emph{Step 3.} By Fubini's theorem, $\int_{\mathbb T}e^{ik\cdot y}\,d^\circ k=\prod_\mu\frac1{2\pi}\int_{-\pi}^{\pi}e^{ity_\mu}\,dt=\delta_{y,0}$, because $\int_{-\pi}^\pi e^{itm}\,dt=0$ for integers $m\ne0$. The substitution $k\to-k$, together with $\mathsf k_\mu(-k)=\overline{\mathsf k_\mu(k)}$, gives
\[
\int_{\mathbb T}e^{ik\cdot y}\,\frac{\mathsf k_\iota\,\overline{\mathsf k_{\iota'}}}{|\mathsf k|^2}\,d^\circ k=\int_{[-\pi,\pi]^3}e^{-ik\cdot y}\,\frac{\overline{\mathsf k_\iota(k)}\,\mathsf k_{\iota'}(k)}{|\mathsf k(k)|^2}\,\frac{dk}{(2\pi)^3}=T_{\iota\iota'}(-y),
\]
where $T_{\lambda\lambda'}(y)=\int_{[-\pi,\pi]^3}e^{ik\cdot y}\,\overline{\mathsf q_\lambda}\,\mathsf q_{\lambda'}\,|\mathsf q|^{-2}\,\frac{dk}{(2\pi)^3}$ and $\mathsf q_\mu(k):=\mathsf k_\mu(k)=e^{ik_\mu}-1$. Together with Steps~1 and~2 we obtain
\[
\Pi^{\mathrm{lat}}(p,q)=\epsilon_\iota\,\epsilon_{\iota'}\bigl[\delta_{\iota\iota'}\,\delta_{y,0}-T_{\iota\iota'}(-y)\bigr].
\]
The Fourier formula in equation~\YMMref{(2242)}{822} states that $\Pi_\infty((z,\lambda),(z',\lambda'))=\delta_{pq}-T_{\lambda\lambda'}(z-z')$ for the plaquettes $p=(z,\lambda)$, $q=(z',\lambda')$ of $\mathbb Z^3$. We apply it to $p^\flat=(x,\iota)$ and $q^\flat=(x+y,\iota')$, for which $z-z'=-y$ and $\delta_{p^\flat q^\flat}=\delta_{\iota\iota'}\delta_{y,0}$; this is the first formula. The second follows from \eqref{IV-73-23}, as $\epsilon_p^2=\epsilon_q^2=1$.
\end{proof}

\begin{lemma}[The list of the pairings]\label{IV-73-53}
Let $\Omega$ and $\tau_{\mathrm{gf}}$ be as in Definition~\ref{IV-73-15}, and let $\langle\cdot\rangle_0$ refer to the Gaussian family of $\Omega$ and $\tau_{\mathrm{gf}}$. For a homogeneous polynomial $\mathsf F$ of degree $m$ in the variables $A^a(b)$ and for $\mathfrak g\otimes\mathbb C$-valued forms $X_1,\dots,X_m$ we put
\[
\mathsf F[\![X_1,\dots,X_m]\!]:=\partial_{s_1}\cdots\partial_{s_m}\,\mathsf F\Bigl(\sum_{i=1}^ms_iX_i\Bigr);
\]
this expression is symmetric and multilinear in $X_1,\dots,X_m$, and $\mathsf F[\![A,\dots,A]\!]=m!\,\mathsf F(A)$. For an index $\mathsf i=(b,a)$, where $b$ is a positively oriented bond of $\Omega$ (a free bond if $\Omega=\Lambda_n$) and $a\in\{1,2,3\}$, let $\delta_{\mathsf i}$ be the form which equals $\mathsf e_a$ at $b$ and $0$ at all other positively oriented bonds. We put $\mathsf C(\mathsf i,\mathsf i'):=\delta_{aa'}\,\mathsf C(b,b')$ for $\mathsf i'=(b',a')$, and $\mathsf E_p^a:=\mathsf E_p(\cdot)\,\mathsf e_a$.
\begin{enumerate}
\item[(a)] (Symmetry factors.) Let $\mathsf F_1,\dots,\mathsf F_r$ be homogeneous polynomials of degrees $m_1,\dots,m_r$. Then $\langle\mathsf F_1;\dots;\mathsf F_r\rangle_0=\sum_{\mathsf n}\mathsf{val}(\mathsf n)$. The sum extends over the symmetric matrices $\mathsf n=(n_{vv'})_{v,v'=1}^r$ of non-negative integers such that $2n_{vv}+\sum_{v'\ne v}n_{vv'}=m_v$ for every $v$ and such that the graph on $\{1,\dots,r\}$ with $n_{vv'}$ edges between $v$ and $v'\ne v$ is connected; and
\[
\begin{aligned}
\mathsf{val}(\mathsf n)=\Bigl[\prod_{v<v'}n_{vv'}!\prod_v2^{n_{vv}}\,n_{vv}!\Bigr]^{-1}&\sum\ \prod_{v=1}^r\mathsf F_v[\![\delta_{\mathsf i(v,1)},\dots,\delta_{\mathsf i(v,m_v)}]\!]\\
&\times\prod_{\text{lines }\{e,e'\}}\mathsf C\bigl(\mathsf i(e),\mathsf i(e')\bigr).
\end{aligned}
\]
Here, for every such $\mathsf n$, we fix one pairing of the slots $(v,1),\dots,(v,m_v)$, $v=1,\dots,r$, with exactly $n_{vv'}$ pairs between the slots of $v$ and those of $v'$, and $n_{vv}$ pairs among the slots of $v$; its pairs $\{e,e'\}$ are called lines, and the sum extends over all indices $\mathsf i(e)$ attached to the slots $e$. In (a), in (b) and in the proof of the lemma the symbol $\mathsf n$ stands for such a matrix; in the formula for $\mathsf T_{\mathrm I}$ and in \eqref{IV-73-54} below it is the number $\mathsf n=3$ of the setting.
\item[(b)] (Colour.) For every plaquette $P$, every real kernel $\mathsf c$ on pairs of bonds, symmetric or not, and every form $X$,
\[
\sum_{\mathsf i,\mathsf i'}\mathsf a_P^{(3)}[\![\delta_{\mathsf i},\delta_{\mathsf i'},X]\!]\,\mathsf c(b,b')\,\delta_{aa'}=0\qquad(\mathsf i=(b,a),\ \mathsf i'=(b',a')).
\]
Consequently $\mathsf{val}(\mathsf n)=0$ if $n_{vv}\ge1$ for some $v$ with $\mathsf F_v=\mathsf a_P^{(3)}$, and also if $n_{vv'}=2$ for some $v,v'$ with $\mathsf F_v=\mathsf a_p^{(2)}$ and $\mathsf F_{v'}=\mathsf a_P^{(3)}$.
\item[(c)] (Classes.) Let $\Sigma_{\mathrm{tot}}$ and $\Gamma_q$ be defined by the formulas of Definition~\ref{IV-73-29}(5), with the Gaussian family and the vertices of $\Omega$. Then
\[
\mathsf B_1^\Omega(p,q;\tau_{\mathrm{gf}})=\mathsf T_{\mathrm I}+\mathsf T_1+\mathsf T_4+\mathsf T_{4'}+\mathsf T_5+\mathsf T_6 ,
\]
where, with $P,P'$ denoting plaquettes and $\mathsf i,\mathsf i',\mathsf j,\mathsf j',\dots$ indices of $\Omega$,
\begin{align*}
\mathsf T_{\mathrm I}&:=-\mathsf n\,\Pi(p,q)\sum_{b,b'}\mathsf E_p(b)\,\Sigma_{\mathrm{tot}}(b,b')\,\mathsf E_q(b'),\\
\mathsf T_1&:=\tfrac12\sum_{P,P'}\sum_{a,a'}\sum_{\mathsf i,\mathsf j}\mathsf a_P^{(3)}[\![\mathsf E_p^a,\mathsf E_q^{a'},\delta_{\mathsf i}]\!]\,\mathsf C(\mathsf i,\mathsf j)\,\mathsf a_{P'}^{(3)}[\![\mathsf E_p^a,\mathsf E_q^{a'},\delta_{\mathsf j}]\!],\\
\mathsf T_4&:=-\tfrac14\sum_P\sum_{a,a'}\mathsf a_P^{(4)}[\![\mathsf E_p^a,\mathsf E_p^a,\mathsf E_q^{a'},\mathsf E_q^{a'}]\!],\\
\mathsf T_{4'}&:=\tfrac14\sum_a\sum_{\mathsf i,\mathsf j}\mathsf C(\mathsf i,\mathsf j)\bigl\{\mathsf a_q^{(4)}[\![\mathsf E_p^a,\mathsf E_p^a,\delta_{\mathsf i},\delta_{\mathsf j}]\!]+\mathsf a_p^{(4)}[\![\mathsf E_q^a,\mathsf E_q^a,\delta_{\mathsf i},\delta_{\mathsf j}]\!]\bigr\},\\
\mathsf T_5&:=-\tfrac12\sum_P\sum_a\sum_{\mathsf i,\mathsf i',\mathsf j,\mathsf j'}\mathsf C(\mathsf i,\mathsf j)\,\mathsf C(\mathsf i',\mathsf j')\bigl\{\mathsf a_q^{(3)}[\![\mathsf E_p^a,\delta_{\mathsf i},\delta_{\mathsf i'}]\!]\,\mathsf a_P^{(3)}[\![\mathsf E_p^a,\delta_{\mathsf j},\delta_{\mathsf j'}]\!]\\
&\qquad\qquad+\mathsf a_p^{(3)}[\![\mathsf E_q^a,\delta_{\mathsf i},\delta_{\mathsf i'}]\!]\,\mathsf a_P^{(3)}[\![\mathsf E_q^a,\delta_{\mathsf j},\delta_{\mathsf j'}]\!]\bigr\},\\
\mathsf T_6&:=\tfrac16\sum\mathsf a_p^{(3)}[\![\delta_{\mathsf i_1},\delta_{\mathsf i_2},\delta_{\mathsf i_3}]\!]\,\mathsf a_q^{(3)}[\![\delta_{\mathsf j_1},\delta_{\mathsf j_2},\delta_{\mathsf j_3}]\!]\prod_{\varsigma=1}^3\mathsf C(\mathsf i_\varsigma,\mathsf j_\varsigma),
\end{align*}
the last sum extending over all indices $\mathsf i_1,\mathsf i_2,\mathsf i_3,\mathsf j_1,\mathsf j_2,\mathsf j_3$. By the first half of $\mathsf T_{4'}$, resp.\ of $\mathsf T_5$, we mean the term with the first summand in braces, by the second half the term with the second summand. On $\Omega=\mathbb Z^3$ the formulas hold provided the sums converge absolutely, which is shown in Lemma~\ref{IV-73-55}. Moreover, if the summands of $\Gamma_q$ are taken in the order of Definition~\ref{IV-73-29}(5), the second and the third being grouped together, then
\begin{equation}\label{IV-73-54}
\begin{aligned}
\tfrac{\mathsf n}2\sum_{b,b'}\Gamma_q(b,b')\,\mathsf E_p(b)\,\mathsf E_p(b')&=\bigl[\text{first half of }\mathsf T_{4'}\bigr]\\
&\quad+\bigl[\text{first half of }\mathsf T_5\bigr]+\mathsf T_4+\mathsf T_1
\end{aligned}
\end{equation}
term by term, the four groups of summands of $\Gamma_q$ producing the four terms on the right in this order; the same holds with $p$ and $q$ exchanged and with the second halves of $\mathsf T_{4'}$ and $\mathsf T_5$.
\item[(d)] (Vertices.) For forms $X,Y,Z$ we have $\mathsf a_P^{(3)}[\![X,Y,Z]\!]=\sum\mathsf t_P(X',Y',Z')$, the sum extending over the six orderings $(X',Y',Z')$ of $(X,Y,Z)$, where
\[
\mathsf t_P(X,Y,Z):=-\Bigl\langle\mathrm dX(P),\sum_{i<j}s_is_j\,Y(b_i)\times Z(b_j)\Bigr\rangle ,
\]
with $b_i=b_i(P)$ and with the complex bilinear extensions of $\langle\cdot,\cdot\rangle$ and of $\times$. The standard writing of $\mathsf a_P^{(4)}$ consists of monomials of three kinds: with four plain legs (these come from $\frac12|\mathfrak z_1|^2$), with one curl leg and three plain legs, and with four curl legs.
\end{enumerate}
\end{lemma}

\begin{proof}
\emph{Step 1: proof of (a).} Writing $A_{\mathsf i}:=A^a(b)$ for $\mathsf i=(b,a)$, we have $A=\sum_{\mathsf i}A_{\mathsf i}\,\delta_{\mathsf i}$, and multilinearity together with $\mathsf F_v(A)=\frac1{m_v!}\mathsf F_v[\![A,\dots,A]\!]$ gives the finite sum
\[
\mathsf F_v=\frac1{m_v!}\sum_{\mathsf i_1,\dots,\mathsf i_{m_v}}\mathsf F_v[\![\delta_{\mathsf i_1},\dots,\delta_{\mathsf i_{m_v}}]\!]\,A_{\mathsf i_1}\cdots A_{\mathsf i_{m_v}} .
\]
We use the multilinearity of cumulants and Lemma~\ref{IV-73-21}, applied to the monomials $A_{\mathsf i(v,1)}\cdots A_{\mathsf i(v,m_v)}$, with the slots $(v,1),\dots,(v,m_v)$ as legs and with the covariance $\langle A_{\mathsf i}A_{\mathsf i'}\rangle_0=\mathsf C(\mathsf i,\mathsf i')$ of Definition~\ref{IV-73-19}(1). We find that $\langle\mathsf F_1;\dots;\mathsf F_r\rangle_0$ is $\prod_v\frac1{m_v!}$ times the sum, over all indices and over all pairings of the slots with connected graph on $\{1,\dots,r\}$, of $\prod_v\mathsf F_v[\![\cdots]\!]\prod_{\text{pairs}}\mathsf C$. As the forms $\mathsf F_v[\![\cdot]\!]$ are symmetric, the contribution of a pairing, summed over the indices, depends only on its matrix $(n_{vv'})$; and the graph of a pairing is connected if and only if that of its matrix is.

It remains to count the pairings with a given matrix $\mathsf n$. At each vertex $v$ we decide which slots are joined to which $v'\ne v$ and which are paired among themselves: there are $m_v!/[\prod_{v'\ne v}n_{vv'}!\,(2n_{vv})!]$ possibilities. For each pair $v<v'$ we match the chosen slots of $v$ with those of $v'$: $n_{vv'}!$ possibilities. At each $v$ we pair the $2n_{vv}$ slots reserved for this purpose: $(2n_{vv})!/(2^{n_{vv}}n_{vv}!)$ possibilities. Since $\prod_v\prod_{v'\ne v}n_{vv'}!=\prod_{v<v'}(n_{vv'}!)^2$, the number of pairings is
\[
\prod_vm_v!\Big/\Bigl[\prod_{v<v'}n_{vv'}!\prod_v2^{n_{vv}}\,n_{vv}!\Bigr],
\]
and multiplication by $\prod_v\frac1{m_v!}$ gives the factor in $\mathsf{val}(\mathsf n)$.

\emph{Step 2: proof of (d) and (b).} By \eqref{IV-73-10}, $\mathsf a^{(3)}_P(A)=\langle\mathrm dA(P),\mathfrak z_1(P)\rangle$ with $\mathfrak z_1=-\sum_{i<j}s_is_j\,A(b_i)\times A(b_j)$; thus $\mathsf a_P^{(3)}(A)=\mathsf t_P(A,A,A)$ with the trilinear form $\mathsf t_P$, and $\partial_{s_1}\partial_{s_2}\partial_{s_3}\,\mathsf t_P\bigl(\sum s_iX_i,\sum s_iX_i,\sum s_iX_i\bigr)$ is the sum of $\mathsf t_P$ over the six orderings of $X_1,X_2,X_3$. The three kinds of monomials of $\mathsf a_P^{(4)}$ are those of the three terms $\frac12|\mathfrak z_1|^2$, $\langle\mathrm dA(P),\mathfrak z_2\rangle$, $-\frac1{24}|\mathrm dA(P)|^4$ of \eqref{IV-73-10}, written as in Definition~\ref{IV-73-19}(3).

For (b) we insert (d): $\mathsf a_P^{(3)}[\![\delta_{(b,a)},\delta_{(b',a)},X]\!]$ is a sum of terms, each the product of a real number which depends on $b$, $b'$ and on the ordering, but not on $a$, with one of $\langle\mathsf e_a,\mathsf e_a\times X'\rangle$, $\langle\mathsf e_a,X'\times\mathsf e_a\rangle$, $\langle X',\mathsf e_a\times\mathsf e_a\rangle$, where $X'$ is a value of $X$ or of $\mathrm dX$. Summed over $a$, each of these three expressions vanishes, since $Y\times X'$ is orthogonal to $Y$ and $Y\times Y=0$; multiplying by $\mathsf c(b,b')$ and summing over $b,b'$ we obtain (b). A line between two slots of a vertex $\mathsf a_P^{(3)}$ produces exactly the sum in (b) with $\mathsf c=\mathsf C$.

By the rule for $\mathsf a_p^{(2)}$ in Step~3 below, two lines from $\mathsf a_p^{(2)}=\frac12\sum_a((\mathrm dA)^a(p))^2$ to the same vertex $\mathsf a_P^{(3)}$ produce the sum in (b) with $\mathsf c(b,b')=\mathsf E_p(b)\,\mathsf E_p(b')$. In both cases $\mathsf{val}(\mathsf n)=0$.

\emph{Step 3: rules for (c).} From $\mathsf a_p^{(2)}=\frac12\sum_a((\mathrm dA)^a(p))^2$ we obtain $\mathsf a_p^{(2)}[\![\delta_{\mathsf i},\delta_{\mathsf i'}]\!]=\delta_{aa'}\,\mathsf s_p(b)\,\mathsf s_p(b')$, where $\mathsf s_p(b_i(p))=s_i$ and $\mathsf s_p(b)=0$ for the bonds $b$ which do not belong to $p$. Since $\sum_b\mathsf s_p(b)\,\mathsf C(b,b'')=\mathsf E_p(b'')$, a line from a slot of $\mathsf a_p^{(2)}$ to a slot of another vertex inserts $\mathsf E_p^a$ into that slot, the colour $a$ being summed together with the other slot of $\mathsf a_p^{(2)}$. A line from $\mathsf a_p^{(2)}$ to $\mathsf a_q^{(2)}$ produces $\Pi(p,q)$ and identifies the colours of the two remaining slots. A vertex $\mathsf a^{(2)}$ carrying a line between its own two slots is not connected to the other vertices.

\emph{Step 4: enumeration for (c).} We go through the seven cumulants of \eqref{IV-73-24} and list the matrices $\mathsf n$ with connected graph whose value does not vanish by (b) or by Step~3; in particular there are no lines between two slots of a vertex $\mathsf a^{(2)}$ or $\mathsf a^{(3)}$. We index these matrices by names of the vertices: the vertices $\mathsf a_p^{(m)}$ and $\mathsf a_q^{(m)}$ are called $p$ and $q$, the vertices $\mathsf a_P^{(m)}$ and $\mathsf a_{P'}^{(3)}$ coming from $\mathsf S_3$ or from $\mathsf V_2$ are called $P$ and $P'$, and a quadratic local term $\mathsf F$ of $\mathsf V_2$, taken as a vertex, is called $\mathsf F$. Thus $n_{pq}$ is the number of lines between the slots of $p$ and those of $q$, $n_{qq}$ the number of lines joining two slots of $q$, and similarly for the other entries.
\begin{itemize}
\item $\langle\mathsf a_p^{(3)};\mathsf a_q^{(3)}\rangle_0$: here $n_{pq}=3$, with the factor $\frac1{3!}=\frac16$; this is $\mathsf T_6$.
\item $\langle\mathsf a_p^{(2)};\mathsf a_q^{(4)}\rangle_0$: here $n_{pq}=2$ and $n_{qq}=1$, with the factor $\frac1{2\cdot2}$; this is the first half of $\mathsf T_{4'}$. In the same way $\langle\mathsf a_p^{(4)};\mathsf a_q^{(2)}\rangle_0$ gives the second half.
\item $\langle\mathsf a_p^{(2)};\mathsf a_q^{(3)};\mathsf a_P^{(3)}\rangle_0$, from the term $-\langle\mathsf a_p^{(2)};\mathsf a_q^{(3)};\mathsf S_3\rangle_0$: the conditions $n_{pq}+n_{pP}=2$, $n_{pq}+n_{qP}=3$, $n_{pP}+n_{qP}=3$ force $n_{pq}=n_{pP}=1$ and $n_{qP}=2$, with the factor $\frac12$; with the sign of \eqref{IV-73-24} and the sum over $P$ this is the first half of $\mathsf T_5$. In the same way $-\langle\mathsf a_p^{(3)};\mathsf a_q^{(2)};\mathsf S_3\rangle_0$ gives the second half.
\item $\langle\mathsf a_p^{(2)};\mathsf a_q^{(2)};\mathsf F\rangle_0$ with a quadratic local term $\mathsf F$ of $\mathsf V_2$: here $n_{pq}=n_{p\mathsf F}=n_{q\mathsf F}=1$, with the factor $1$, and the value is $\Pi(p,q)\sum_a\mathsf F[\![\mathsf E_p^a,\mathsf E_q^a]\!]=\Pi(p,q)\sum_a\sum_{b,b'}\mathsf E_p(b)\,\mathsf E_q(b')\,\partial_{b,a}\partial_{b',a}\mathsf F$.
\item $\langle\mathsf a_p^{(2)};\mathsf a_q^{(2)};\mathsf a_P^{(4)}\rangle_0$: either $n_{pq}=n_{pP}=n_{qP}=n_{PP}=1$, with the factor $\frac12$ and the value
\[
\begin{aligned}
&\Pi(p,q)\sum_a\tfrac12\sum_{\mathsf i,\mathsf j}\mathsf C(\mathsf i,\mathsf j)\,\mathsf a_P^{(4)}[\![\mathsf E_p^a,\mathsf E_q^a,\delta_{\mathsf i},\delta_{\mathsf j}]\!]\\
&\qquad=\Pi(p,q)\sum_a\sum_{b,b'}\mathsf E_p(b)\,\mathsf E_q(b')\,\langle\partial_{b,a}\partial_{b',a}\mathsf a_P^{(4)}\rangle_0
\end{aligned}
\]
(by (a) for the quadratic polynomial $\partial_{b,a}\partial_{b',a}\mathsf a_P^{(4)}$, whose polarisation is $\mathsf a_P^{(4)}[\![\delta_{(b,a)},\delta_{(b',a)},\cdot,\cdot]\!]$); or $n_{pP}=n_{qP}=2$, with the factor $\frac14$, which with the sign of \eqref{IV-73-24} and the sum over $P$ gives $\mathsf T_4$. The matrix with $n_{pq}=2$ has a disconnected graph.
\item $\langle\mathsf a_p^{(2)};\mathsf a_q^{(2)};\mathsf a_P^{(3)};\mathsf a_{P'}^{(3)}\rangle_0$: subtracting the conditions at $p$ and $q$ from those at $P$ and $P'$ we get $n_{PP'}=1+n_{pq}$. For $n_{pq}=2$ the graph is disconnected. For $n_{pq}=1$ we get $n_{PP'}=2$ and either $n_{pP}=n_{qP'}=1$ or $n_{pP'}=n_{qP}=1$, each with the factor $\frac12$; summed over $P,P'$ and multiplied by the $\frac12$ of \eqref{IV-73-24}, they give
\[
\begin{aligned}
&\Pi(p,q)\sum_a\tfrac12\sum\mathsf a_P^{(3)}[\![\mathsf E_p^a,\delta_{\mathsf i},\delta_{\mathsf i'}]\!]\,\mathsf a_{P'}^{(3)}[\![\mathsf E_q^a,\delta_{\mathsf j},\delta_{\mathsf j'}]\!]\,\mathsf C(\mathsf i,\mathsf j)\,\mathsf C(\mathsf i',\mathsf j')\\
&\qquad=\Pi(p,q)\sum_a\sum_{b,b'}\mathsf E_p(b)\,\mathsf E_q(b')\,\langle\partial_{b,a}\mathsf S_3;\partial_{b',a}\mathsf S_3\rangle_0 ,
\end{aligned}
\]
where the sum on the left extends over $P,P'$ and all indices, and the equality is (a) for two quadratic polynomials. For $n_{pq}=0$ we get $n_{PP'}=1$ and $n_{pP}+n_{qP}=2=n_{pP'}+n_{qP'}$; the possibilities $(n_{pP},n_{qP})=(2,0)$ and $(0,2)$ vanish by (b), and there remains $n_{pP}=n_{qP}=n_{pP'}=n_{qP'}=1$, with the factor $1$; with the $\frac12$ of \eqref{IV-73-24} this is $\mathsf T_1$.
\item We collect the terms with $n_{pq}=1$ found above, with the signs of \eqref{IV-73-24}. By the definition of $\Sigma_{\mathrm{tot}}$ in Definition~\ref{IV-73-29}(5) they add up to $-\Pi(p,q)\sum_a\sum_{b,b'}\mathsf E_p(b)\,\mathsf E_q(b')\,\Sigma_{\mathrm{tot}}(b,a;b',a)$. The argument on the colour structure given there applies in the box as well and shows $\Sigma_{\mathrm{tot}}(b,a;b',a)=\Sigma_{\mathrm{tot}}(b,b')$ for every colour $a$; summing over the three colours we obtain $\mathsf T_{\mathrm I}$.
\end{itemize}

\emph{Step 5: proof of \eqref{IV-73-54}.} Let $\partial=\partial_{b,a}$, $\partial'=\partial_{b',a'}$, $\delta=\delta_{(b,a)}$, $\delta'=\delta_{(b',a')}$. By (a),
\begin{align*}
\langle\partial\partial'\mathsf a_q^{(4)}\rangle_0&=\tfrac12\sum_{\mathsf i,\mathsf j}\mathsf C(\mathsf i,\mathsf j)\,\mathsf a_q^{(4)}[\![\delta,\delta',\delta_{\mathsf i},\delta_{\mathsf j}]\!],\\
\langle\partial\mathsf a_q^{(3)};\partial'\mathsf S_3\rangle_0&=\tfrac12\sum_P\sum\mathsf C(\mathsf i,\mathsf j)\,\mathsf C(\mathsf i',\mathsf j')\,\mathsf a_q^{(3)}[\![\delta,\delta_{\mathsf i},\delta_{\mathsf i'}]\!]\,\mathsf a_P^{(3)}[\![\delta',\delta_{\mathsf j},\delta_{\mathsf j'}]\!],\\
\langle\mathsf a_q^{(2)};\partial\partial'\mathsf S_4\rangle_0&=\tfrac12\sum_P\sum_{a''}\mathsf a_P^{(4)}[\![\delta,\delta',\mathsf E_q^{a''},\mathsf E_q^{a''}]\!],\\
\langle\mathsf a_q^{(2)};\partial\mathsf S_3;\partial'\mathsf S_3\rangle_0&=\sum_{P,P'}\sum_{a''}\sum_{\mathsf i,\mathsf j}\mathsf C(\mathsf i,\mathsf j)\,\mathsf a_P^{(3)}[\![\delta,\mathsf E_q^{a''},\delta_{\mathsf i}]\!]\\
&\qquad\times\mathsf a_{P'}^{(3)}[\![\delta',\mathsf E_q^{a''},\delta_{\mathsf j}]\!];
\end{align*}
in the last line the three polynomials are quadratic, the only connected matrix has one line between any two of them, and its factor is $1$. We contract these expressions, and the one obtained from the second by exchanging $(b,a)$ and $(b',a')$, with $\frac12\sum_a\mathsf E_p^a\otimes\mathsf E_p^a$; that is, we set $a'=a$, multiply by $\frac12\,\mathsf E_p(b)\,\mathsf E_p(b')$ and sum over $b,b',a$. Since $\Gamma_q(b,a;b',a)=\Gamma_q(b,b')$, this turns $\Gamma_q$ into the left side of \eqref{IV-73-54}. By multilinearity, and with the signs of Definition~\ref{IV-73-29}(5), the first summand of $\Gamma_q$ produces the first half of $\mathsf T_{4'}$; the second and the third, whose contributions coincide after the symmetric contraction, together produce the first half of $\mathsf T_5$; the fourth produces $\mathsf T_4$; and the fifth produces $\mathsf T_1$. Exchanging $p$ and $q$ gives the statement for the second halves.
\end{proof}

\begin{lemma}[Convergence of the graph sums on $\mathbb Z^3$]\label{IV-73-55}
Let $\tau_{\mathrm{gf}}>0$.
\begin{enumerate}
\item[(i)] On $\mathbb Z^3$, for all plaquettes $p,q$, the sums over the positions of the local terms of $\mathsf S_3$ and $\mathsf V_2$ in \eqref{IV-73-24}, taken monomial by monomial in the standard writing and pairing by pairing, converge absolutely. Consequently $\mathsf B_1^{\mathrm{lat}}(p,q;\tau_{\mathrm{gf}})$ is defined, the formulas of Lemma~\ref{IV-73-53}(c) hold on $\mathbb Z^3$, and $\sup_{p,q}|\mathsf B_1^{\mathrm{lat}}(p,q;\tau_{\mathrm{gf}})|<\infty$.
\item[(ii)] (Uniform form.) Let $\mathsf c_0\ge1$, let $\Omega=\mathbb Z^3$ or $\Omega=\Lambda_n$, and substitute for $\mathsf C,\mathsf E,\Pi,\mathsf D,\mathsf g$ arbitrary real kernels which satisfy, with the base points $b_-$ of the bonds and $x_p$ of the plaquettes,
\begin{equation}\label{IV-73-56}
\begin{gathered}
|\mathsf C(b,b')|\le\frac{\mathsf c_0}{\langle b_--b'_-\rangle},\qquad|\mathsf E_p(b)|\le\frac{\mathsf c_0}{\langle x_p-b_-\rangle^{2}},\\
|\mathsf D(b;b')|\le\frac{\mathsf c_0}{\langle b_--b'_-\rangle^{2}},\qquad|\Pi(p,q)|\le\frac{\mathsf c_0}{\langle x_p-x_q\rangle^{3}},\qquad|\mathsf g(b)|\le\mathsf c_0 .
\end{gathered}
\end{equation}
Then the sum of the absolute values of all terms of the seven cumulants of \eqref{IV-73-24}, over monomials, pairings and positions, is at most $C\,\mathsf c_0^5$, with a pure constant $C$.
\end{enumerate}
\end{lemma}

\begin{proof}
\emph{Step 1.} By Lemma~\ref{IV-73-51}, the kernels $\mathsf C$, $\mathsf E$, $\Pi$, $\mathsf D$ of $\mathbb Z^3$ are, as functions of the difference of the base points, elementary lines of the degrees $1,2,3,2$, with finitely many different symbols. Lemma~\ref{IV-73-39}(a), applied with $s=-2,-1,0,-1$ and $\gamma=0$, bounds them by constants times $\langle\cdot\rangle^{-1}$, $\langle\cdot\rangle^{-2}$, $\langle\cdot\rangle^{-3}$, $\langle\cdot\rangle^{-2}$, and $\mathsf g=\frac13$ on $\mathbb Z^3$ by Definition~\ref{IV-73-15}(2). Hence these kernels satisfy \eqref{IV-73-56} with some $\mathsf c_0=\mathsf c_0(\tau_{\mathrm{gf}})$, and (i) follows from (ii): the bound $C\mathsf c_0^5$, which does not depend on $p,q$, is the absolute convergence required in Definition~\ref{IV-73-19}(3), so that $\mathsf B_1^{\mathrm{lat}}$ is defined and those rearrangements made in the proof of Lemma~\ref{IV-73-53}(c) are allowed whose terms are bounded by the majorants of Step~2 below, such as regroupings of whole terms of positions and regroupings which split one curl leg or, as in Step~5 of that proof and for $n_{pq}=n_{pP}=n_{qP}=n_{PP}=1$ in its Step~4, curl legs at $P$ or $P'$ joined to $p$ or $q$, whose lines then have the degree $2$ which these majorants have there.

For the kernels of $\mathbb Z^3$ we obtain the equality of the display for $n_{pq}=1$ in Step~4 of that proof, whose right side splits the curl legs at $P$ and at $P'$ joined to $p$ and to $q$, as in $\Pi(p,P)=\sum_b\mathsf s_P(b)\,\mathsf E_p(b)$, by regrouping one vertex at a time, with $u,z,w$ as in Step~2. At $P$ we write $\mathsf a_P^{(3)}[\![\mathsf E_p^a,\cdot,\cdot]\!]=\sum_b\mathsf E_p(b)\,\mathsf a_P^{(3)}[\![\delta_{(b,a)},\cdot,\cdot]\!]$, split the quadratic polynomial $\partial_{b,a}\mathsf a_P^{(3)}$, whose polarisation is the last factor, into its two terms as in \eqref{IV-73-58}, which follows from \eqref{IV-73-10} by the Leibniz rule, and keep the vertex at $P'$ in the standard writing; the resulting terms, indexed by $(P,b,P')$, the monomials and the pairings, are bounded by $C\langle z\rangle^{-2}\langle w-z\rangle^{-n_3}\langle u-w\rangle^{-n_5}$, with a constant $C$ which depends only on $\tau_{\mathrm{gf}}$, as the kernels of $\mathbb Z^3$ satisfy \eqref{IV-73-56} with $\mathsf c_0=\mathsf c_0(\tau_{\mathrm{gf}})$, and with the four pairs $(n_3,n_5)$ of Step~2, for which the computation there holds for every $n_1\ge2$; so they may be grouped by $(b,P')$, and the sum over $P\ni b$ turns the first entry into $\partial_{b,a}\mathsf S_3$.

Next, with $z,w$ the positions of $b,P'$ relative to the base point of $p$ and with $C=C(\tau_{\mathrm{gf}})$, we have $|\langle\partial_{b,a}\mathsf S_3;\partial_{b',a}\mathsf a_{P'}^{(3)}\rangle_0|\le C\langle w-z\rangle^{-3}$: by \eqref{IV-73-58}, a monomial of $\mathsf u_{b,a}$ has a curl leg and a plain leg, so that the two lines joining it to a monomial of $\partial_{b',a}\mathsf a_{P'}^{(3)}$ have degrees of sum at least $3$, and \eqref{IV-73-56} applies; and $\mathsf d_{b,a}$ contributes $\sum_{P\ni b}\mathsf s_P(b)\,\mathsf F(P)$ with $\mathsf F(P):=\langle\langle\mathsf e_a,\mathfrak z_1(P)\rangle;\partial_{b',a}\mathsf a_{P'}^{(3)}\rangle_0$, by Lemma~\ref{IV-73-53}(a) and Lemma~\ref{IV-73-51} a finite sum of products $K_{\mathrm a}K_{\mathrm b}$ of elementary lines in the base point of $P$ with $n_{\mathrm a}+n_{\mathrm b}\ge2$; since the four plaquettes containing $b$ form two pairs, each of one dual direction, with base points differing by one of the vectors $e_1,e_2,e_3$ and with opposite signs, this is a sum of two first differences in the base point of $P$, which Lemma~\ref{IV-73-46}(c) with $|\gamma|=1$ bounds by $C\langle w-z\rangle^{-3}$.

At $P'$, by Lemma~\ref{IV-73-53}(a) for two quadratic polynomials, the term of $(b,P')$ so obtained is the sum over $b'$ of $\mathsf E_p(b)\,\mathsf E_q(b')\,\langle\partial_{b,a}\mathsf S_3;\partial_{b',a}\mathsf a_{P'}^{(3)}\rangle_0$; these terms, indexed by $(b,P',b')$, are bounded by $C\langle z\rangle^{-2}\langle w-z\rangle^{-3}\langle u-w\rangle^{-2}$, the pair $(3,2)$ of Step~2; so they may be grouped by $(b,b')$, and the sum over $P'\ni b'$ gives the right side of that display. Thus this equality holds on $\mathbb Z^3$, and the sum over $b,b'$ in the formula for $\mathsf T_{\mathrm I}$ converges absolutely: for the summand $-\langle\partial_{b,a}\mathsf S_3;\partial_{b',a}\mathsf S_3\rangle_0$ of $\Sigma_{\mathrm{tot}}$ by these majorants, and for the summands of $\langle\partial_{b,a}\partial_{b',a}\mathsf V_2\rangle_0$, which by \eqref{IV-73-17} and \eqref{IV-73-56} are bounded and of finite range, or bounded by $8\mathsf c_0^2\langle b_--b'_-\rangle^{-4}$, by the sums of Step~2.

\emph{Step 2.} We prove (ii). We write every vertex in the standard writing and count a curl leg as one leg (Lemma~\ref{IV-73-53}(d)); at every position there are finitely many monomials, with coefficients independent of the position. We follow the enumeration in the proof of Lemma~\ref{IV-73-53}(c); at fixed positions all sums there are finite. By \eqref{IV-73-56} a term is bounded, up to a bounded factor, by $\prod_\ell\mathsf c_0\langle y_\ell\rangle^{-n_\ell}$, where $\ell$ ranges over its lines, $y_\ell$ is the difference of the positions of the two vertices joined by $\ell$, and $n_\ell=1+\#$(curl legs at $\ell$); the bounded factor accounts for shifts of length at most $3$ between the base points of the legs and the positions of the vertices, since $\langle y\rangle\le4\langle y+v\rangle$ for $|v|\le3$.

Every term has at most five lines, $\mathsf c_0\ge1$, and the numbers of monomials and pairings are pure constants; it therefore suffices to bound the sums over the positions of the products $\prod_\ell\langle y_\ell\rangle^{-n_\ell}$ by pure constants. Let $u$ be the base point of $q$ minus that of $p$, and let $z,w$ be the positions of $P,P'$ (or of the bonds of the local terms of $\mathsf V_2$) relative to the base point of $p$.
\begin{itemize}
\item $\mathsf T_1$: the lines starting at $p$ or ending at $q$ carry a curl leg of $p$ or of $q$, so their degrees are at least $2$, and the line between $P$ and $P'$ has degree at least $1$. As $\langle\cdot\rangle^{-n}$ decreases in $n$, the sum is at most the sum in Proposition~\ref{IV-73-47}(i) with $(n_1,\dots,n_5)=(2,2,1,2,2)$, which satisfies the hypotheses of that proposition.
\item The terms of $\mathsf T_{\mathrm I}$ with two cubic vertices, or with the ghost loop, have sums of the form
\[
\begin{gathered}
\langle u\rangle^{-3}\sum_{z,w}\langle z\rangle^{-n_1}\langle w-z\rangle^{-n_3}\langle u-w\rangle^{-n_5},\\
n_1,n_5\in\{2,3\},\qquad n_3\ge2,\qquad n_1+n_3+n_5=8 ;
\end{gathered}
\]
the factor $\langle u\rangle^{-3}$ comes from $\Pi(p,q)$, and $\langle w-z\rangle^{-n_3}$ from the two lines between the cubic vertices, resp.\ from $\mathsf D(b;b')\mathsf D(b';b)$. The possible pairs are $(n_3,n_5)=(2,3),(3,2),(3,3),(4,2)$, and for each of them Lemma~\ref{IV-73-40}(c) gives $\sum_w\langle w-z\rangle^{-n_3}\langle u-w\rangle^{-n_5}\le C\langle u-z\rangle^{-2}\lg\langle u-z\rangle$. Next, $\sum_z\langle z\rangle^{-2}\langle u-z\rangle^{-2}\lg\langle u-z\rangle\le C$: by $\lg t\le2t^{1/2}$ for $t\ge1$ and by Lemma~\ref{IV-73-40}(c) with $(a,a')=(2,\frac32)$ this sum is at most $C\langle u\rangle^{-1/2}\lg\langle u\rangle\le C$. Finally $\langle u\rangle^{-3}\le1$.
\item The other terms of $\mathsf T_{\mathrm I}$ (with a line between two slots of the quartic vertex at $P$, with the Haar term, with the ghost tadpole) have sums bounded by $\sum_z\langle z\rangle^{-2}\langle u-z\rangle^{-2}\le C$, by Lemma~\ref{IV-73-40}(c).
\item $\mathsf T_4$ leads to $\sum_z\langle z\rangle^{-4}\langle u-z\rangle^{-4}$, and $\mathsf T_5$ to $\sum_z\langle z\rangle^{-n_X}\langle u-z\rangle^{-n_Y}$ with $n_X,n_Y\ge2$, the line between $p$ and $q$ being bounded by $\mathsf c_0$; both sums are bounded by Lemma~\ref{IV-73-40}(c). The sums of $\mathsf T_{4'}$ and $\mathsf T_6$ are finite, as these classes have no vertex besides $p$ and $q$.
\item The graphs which vanish by colour in Lemma~\ref{IV-73-53}(b) vanish pairing by pairing, so they contribute no terms. Indeed $\langle\mathrm dA,A_i\times A_j\rangle=\sum\epsilon_{abc}(\mathrm dA)^aA_i^bA_j^c$, so every monomial of $\mathsf a_P^{(3)}$ has three legs of three different colours, while $((\mathrm dA)^a(p))^2$ has two legs of one colour; since the covariance is diagonal in the colours, a line between two legs of one monomial of $\mathsf a_P^{(3)}$, or two lines from $((\mathrm dA)^a(p))^2$ to one monomial of $\mathsf a_P^{(3)}$, give the factor $0$.
\end{itemize}
In the box the positions range over subsets of $\mathbb Z^3$, and the same bounds apply.
\end{proof}

\subsection{Elementary terms, their expansions, the local parts}

\begin{lemma}[Elementary terms and their degrees]\label{IV-73-57}
Let $\tau_{\mathrm{gf}}>0$, let us work on $\mathbb Z^3$, let $p=p_{[\iota]}(x_p)$ and $q=p_{[\iota']}(x_q)$ be plaquettes, and put $u':=x_q-x_p$. We write every vertex occurring in the formulas of Lemma~\ref{IV-73-53}(c) in the standard writing (Definition~\ref{IV-73-19}(3), Lemma~\ref{IV-73-53}(d)), we carry out the sums over the colours, and we write every contraction as an elementary line by Lemma~\ref{IV-73-51}. Then there are finite index sets $\mathfrak G_\bullet=\mathfrak G_\bullet(\iota,\iota')$, real numbers $c_G$ and elementary lines $K^G_\bullet$, all independent of $x_p$ and $x_q$, such that the following representations hold; every $G$ is called an elementary term. We omit the index $G$ on the lines, and $*$ denotes the convolution of Lemma~\ref{IV-73-46}.
\begin{enumerate}
\item[(1)] $\mathsf T_1=\sum_{G\in\mathfrak G_1}c_G\,\mathsf T^G(u')$, where every $\mathsf T^G$ is a kite as in Proposition~\ref{IV-73-47}, with degrees satisfying the hypotheses of that proposition and with $\varDelta=5$. Here $K_1$ and $K_2$ are contractions of $(\mathrm dA)(p)$ with legs at $P$ and at $P'$, $K_4$ and $K_5$ are contractions of legs at $P$ and at $P'$ with $(\mathrm dA)(q)$, and $K_3$ is the contraction of a leg at $P$ with a leg at $P'$.
\item[(4)] $\mathsf T_4=\sum_{G\in\mathfrak G_4}c_G\,(X*Y)(u')$, where $X=K_1K_2$ and $Y=K_3K_4$ with $n_i\in\{2,3\}$, and $(n_X,n_Y):=(n_1+n_2,n_3+n_4)\in\{(4,4),(5,4),(4,5),(6,6)\}$, with $(n_X,n_Y)=(4,4)$ for the monomials of $\frac12|\mathfrak z_1|^2$.
\item[(4$'$)] The first half of $\mathsf T_{4'}$ equals $\sum_{G\in\mathfrak G_{4'}}c_G\,K_1(u')\,K_2(u')$, where $n_i\in\{2,3\}$ and $K_1,K_2$ are contractions of $(\mathrm dA)(p)$ with legs at $q$.
\item[(5)] The first half of $\mathsf T_5$ equals $\sum_{G\in\mathfrak G_5}c_G\,K_0(u')\,(X*Y)(u')$ with $Y=K_{\mathrm a}K_{\mathrm b}$. Here $K_0$ and $X$ are contractions of $(\mathrm dA)(p)$ with a leg at $q$, resp.\ at $P$, and $n_0,n_X\in\{2,3\}$; $K_{\mathrm a},K_{\mathrm b}$ are contractions of legs at $P$ with legs at $q$, and $n_{\mathrm a},n_{\mathrm b}\in\{1,2,3\}$; finally $n_0+n_X+n_{\mathrm a}+n_{\mathrm b}=8$, so that $n_Y:=n_{\mathrm a}+n_{\mathrm b}\in\{2,3,4\}$.
\item[(6)] $\mathsf T_6=\sum_{G\in\mathfrak G_6}c_G\,K_1(u')\,K_2(u')\,K_3(u')$, with $n_1+n_2+n_3=5$.
\item[(I)] $\mathsf T_{\mathrm I}=-\mathsf n\,\Pi(p,q)\sum_{\mu,\nu}\bigl(K^\to_{\iota\mu}*\Sigma_{\mu\nu}*K^\leftarrow_{\nu\iota'}\bigr)(u')$, where $K^\to_{\iota\mu}(y):=\mathsf E_{p_{[\iota]}(0)}((y,\mu))$ and $K^\leftarrow_{\nu\iota'}(y):=\mathsf E_{p_{[\iota']}(y)}((0,\nu))$ are elementary lines of degree $2$, and $\Sigma_{\mu\nu}(y):=\Sigma_{\mathrm{tot}}((x,\mu),(x+y,\nu))$, which does not depend on $x$. With $\mathsf s_P(b):=s_i$ if $b=b_i(P)$ and $\mathsf s_P(b):=0$ otherwise,
\begin{equation}\label{IV-73-58}
\begin{gathered}
\partial_{b,a}\mathsf S_3=\mathsf d_{b,a}+\mathsf u_{b,a},\\
\mathsf d_{b,a}:=\sum_{P\ni b}\mathsf s_P(b)\,\langle\mathsf e_a,\mathfrak z_1(P)\rangle,\qquad\mathsf u_{b,a}:=\sum_{P\ni b}\langle\mathrm dA(P),\partial_{b,a}\mathfrak z_1(P)\rangle ;
\end{gathered}
\end{equation}
the monomials of $\mathsf d_{b,a}$ have two plain legs, those of $\mathsf u_{b,a}$ a curl leg and a plain leg. Moreover
\begin{equation}\label{IV-73-59}
\Sigma_{\mu\nu}=\mathsf q^0_{\mu\nu}+\sum_{G\in\mathfrak G_{\mathrm I}(\mu,\nu)}c_G\,\mathsf f_G(\nabla)\bigl[K_{\mathrm a}K_{\mathrm b}\bigr],
\end{equation}
where $\mathsf q^0_{\mu\nu}$ has finite support (it comes from the line between two slots of the quartic Wilson vertex, from the Haar term and from the ghost tadpole); $\mathsf f_G(k)$ is the product of two trigonometric polynomials, each a line symbol of degree $0$ or $1$, of total degree $d_G\in\{0,1,2\}$; $\mathsf f_G(\nabla)$ is the finite difference operator which multiplies symbols by $\mathsf f_G$; $n_{\mathrm a},n_{\mathrm b}\in\{1,2,3\}$; and $n_{\mathrm a}+n_{\mathrm b}+d_G=4$. The terms $G$ are the ghost loop $-8\,\mathsf D(b;b')\,\mathsf D(b';b)$, with $d_G=0$, and the terms of $-\langle\partial\mathsf S_3;\partial'\mathsf S_3\rangle_0$ multiplied out by \eqref{IV-73-58}, with $d_G$ the number of entries $\mathsf d$.
\end{enumerate}
The second halves of $\mathsf T_{4'}$ and $\mathsf T_5$ have the same forms, with $p$ and $q$ exchanged and $u'$ replaced by $-u'$.
\end{lemma}

\begin{proof}
We use two counting rules: every monomial of $\mathsf a^{(3)}$ has exactly one curl leg and $\mathsf a^{(2)}$ has two (Definition~\ref{IV-73-19}(3)); and the degree of a line is $1+\#$(its curl legs) (Lemma~\ref{IV-73-51}). By the consequence stated in Lemma~\ref{IV-73-51}, every contraction is an elementary line in the difference of the base points of the two vertices which it joins, its symbol depending only on the two legs; the shifts between the base points of the legs and those of the vertices are absorbed into the lines.

(1) A kite has five lines. There are four curl legs at $p$ and $q$ and one at each of $P$ and $P'$, hence $\sum_in_i=5+6=11$ and $\varDelta=5$. For $i\in\{1,2,4,5\}$ one end of the line $i$ is a curl leg of $p$ or of $q$, so that $n_i\ge2$, and $n_i\le3$. Besides the two curl legs of $p$, the lines $1,2,3$ carry at most two further curl legs, one of $P$ and one of $P'$; hence $n_1+n_2+n_3\le2+2+1+2=7$, and the same count applies to the lines $3,4,5$. With the base points $x_p+z$ and $x_p+w$ of $P$ and $P'$, the sum over $P,P'$ becomes the sum over $z,w$ and over the dual directions of $P,P'$, and the five lines become functions of $z$, $w$, $w-z$, $u'-z$, $u'-w$.

(4) Each of the four lines carries a curl leg at $p$ or at $q$; it has degree $2$, or degree $3$ if its leg at $P$ is a curl leg. The three kinds of monomials of $\mathsf a_P^{(4)}$ (Lemma~\ref{IV-73-53}(d)) give $(n_X,n_Y)=(4,4)$; $(5,4)$ or $(4,5)$; and $(6,6)$. With the base point $x_p+z$ of $P$, the two lines from $p$ make up $X(z)$, the two lines to $q$ make up $Y(u'-z)$, and the sum over $z$ is $(X*Y)(u')$.

(4$'$) There are two lines from $p$ to $q$, each with its curl leg at $p$; the line between two slots of the vertex at $q$ is a number, independent of the positions, and it is absorbed into $c_G$.

(5) There are four lines, and there are $2+1+1$ curl legs, at $p$, $q$ and $P$; hence $\sum n=4+4=8$. With the base point $x_p+z$ of $P$, the line from $p$ to $P$ is $X(z)$, the two lines from $P$ to $q$ make up $Y(u'-z)$, and the line from $p$ to $q$ is $K_0(u')$.

(6) There are three lines, and there are $1+1$ curl legs, at $p$ and $q$; hence $\sum n=3+2=5$.

(I) In the formula for $\mathsf T_{\mathrm I}$ we write $b=(x,\mu)$ and $b'=(x',\nu)$ with $x=x_p+z$, $x'=x_p+w$. By translation invariance $\mathsf E_p(b)=K^\to_{\iota\mu}(z)$ and $\mathsf E_q(b')=K^\leftarrow_{\nu\iota'}(u'-w)$, and by the translation covariance in Definition~\ref{IV-73-29}(5) $\Sigma_{\mathrm{tot}}(b,b')=\Sigma_{\mu\nu}(w-z)$; this gives the formula for $\mathsf T_{\mathrm I}$. The functions $K^\to_{\iota\mu}$ and $K^\leftarrow_{\nu\iota'}$ are elementary lines of degree $2$ by Lemma~\ref{IV-73-51}, the second because $K^\leftarrow_{\nu\iota'}(y)=\mathsf E_{p_{[\iota']}(0)}((-y,\nu))$. By Definition~\ref{IV-73-29}(5) and \eqref{IV-73-17}, with $\partial=\partial_{b,a}$ and $\partial'=\partial_{b',a}$,
\[
\begin{aligned}
\Sigma_{\mathrm{tot}}(b,b')={}&\Bigl[\langle\partial\partial'\mathsf S_4\rangle_0+\bigl(\tfrac23+\tfrac43\mathsf g(b)\bigr)\delta_{bb'}\Bigr]\\
&-4\bigl[\mathsf D(b;b')\mathsf D(b';b)+\mathsf D(b';b)\mathsf D(b;b')\bigr]-\langle\partial\mathsf S_3;\partial'\mathsf S_3\rangle_0 .
\end{aligned}
\]
The bracket has finite range and is translation invariant; it gives $\mathsf q^0_{\mu\nu}$. The middle term is the ghost loop $-8\,\mathsf D(b;b')\,\mathsf D(b';b)$, a product of two elementary lines of degree $2$ in $y$ (Lemma~\ref{IV-73-51}), with $\mathsf f_G=1$. From \eqref{IV-73-10} and the Leibniz rule we have
\[
\partial_{b,a}\langle\mathrm dA(P),\mathfrak z_1(P)\rangle=\mathsf s_P(b)\,\langle\mathsf e_a,\mathfrak z_1(P)\rangle+\langle\mathrm dA(P),\partial_{b,a}\mathfrak z_1(P)\rangle ,
\]
which gives \eqref{IV-73-58}. For a function $\mathsf F$ on the plaquettes of one dual direction, the map $b\mapsto\sum_{P\ni b}\mathsf s_P(b)\,\mathsf F(P)$ is a difference operator whose symbol is $\pm\overline{\mathsf k_\alpha}$ up to a factor $e^{ik\cdot y_0}$, a line symbol of degree $1$; this operator yields the factor of $\mathsf f_G$ which belongs to an entry $\mathsf d$. An entry $\mathsf u$ we split into its monomials and treat each of them separately. By the consequence stated in Lemma~\ref{IV-73-51}, the contraction of a leg with base point $x+\xi$ with a leg with base point $x+y+\xi'$ is an elementary line in $y$ in which the shifts $\xi$ and $\xi'$ are already contained; hence the factor of $\mathsf f_G$ which belongs to an entry $\mathsf u$ is $1$, a line symbol of degree $0$ with leading part $1$, just as for the ghost loop.

By Lemma~\ref{IV-73-53}(a), the cumulant of two quadratic monomials is a sum of two products of two lines. Hence every term of $-\langle\partial\mathsf S_3;\partial'\mathsf S_3\rangle_0$, multiplied out by \eqref{IV-73-58}, has the form $c_G\,\mathsf f_G(\nabla)[K_{\mathrm a}K_{\mathrm b}]$, with $d_G$ the number of entries $\mathsf d$ and with $n_{\mathrm a}+n_{\mathrm b}=2+\#$(curl legs on the two lines). As $\mathsf d$ has no curl leg and $\mathsf u$ has one, $n_{\mathrm a}+n_{\mathrm b}=2+(2-d_G)$. The second halves of $\mathsf T_{4'}$ and $\mathsf T_5$ are treated in the same way, the roles of $p$ and $q$ being exchanged.
\end{proof}

\begin{lemma}[Expansions of the elementary terms, with errors]\label{IV-73-60}
Let the setting be that of Lemma~\ref{IV-73-57}, let $r':=|u'|$ be large, and let $K^h:=\mathcal F^{-1}_{\ne0}\sigma^h$ for every line $K$ with symbol $\sigma$. For the bubbles $X$ and $Y$ let $\mathsf q_X$ and $\mathsf q_Y$ be the kernels $\mathsf q$ of Lemma~\ref{IV-73-46}(c), and put $X':=X-\mathsf q_X$, $Y':=Y-\mathsf q_Y$. Then the elementary terms of Lemma~\ref{IV-73-57} have the following expansions.
\begin{enumerate}
\item[(6)] $K_1K_2K_3(u')=K_1^hK_2^hK_3^h(u')+O(r'^{-6})$.
\item[(5)] $K_0(u')\,(X*Y)(u')=K_0(u')\,(\mathsf q_Y*X)(u')+K_0^h(u')\,(X*Y')^h(u')+O(r'^{-6}\lg r')$, where $(X*Y')^h$ is the second component of the leading data of $X*Y'\in\mathfrak K(n_X+n_Y-3)$.
\item[(4)] $(X*Y)(u')=(\mathsf q_X*Y)(u')+(\mathsf q_Y*X)(u')+\mathsf U_G(u')+O(r'^{-6}\lg r')$, where $\mathsf U_G:=(X'*Y')^h$ if $(n_X,n_Y)=(4,4)$ and $\mathsf U_G:=0$ otherwise.
\item[(1)] The expansion of Proposition~\ref{IV-73-47}(iii) holds, with the error $O(r'^{-6}\lg r')$.
\item[($*$)] In all four cases the function of degree $-5$ in the expansion (that is, $K_1^hK_2^hK_3^h$, $K_0^h(X*Y')^h$, $\mathsf U_G$, resp.\ the function $\mathsf U$ of Proposition~\ref{IV-73-47}) is determined by the leading parts $\sigma^h$ of the lines alone; it is multilinear in them, it belongs to $C^1(\mathbb R^3\setminus\{0\})$, and it is rotation covariant as in Proposition~\ref{IV-73-47}(iv).
\end{enumerate}
\end{lemma}

\begin{proof}
(6) By Lemma~\ref{IV-73-39}(a),(c), $K_i(u')=O(r'^{-n_i})$, $K_i^h(u')=O(r'^{-n_i})$ and $K_i(u')=K_i^h(u')+O(r'^{-n_i-1})$. As $n_1+n_2+n_3=5$, the product $K_1K_2K_3(u')$ differs from $K_1^hK_2^hK_3^h(u')$ by $O(r'^{-6})$, and $K_1^hK_2^hK_3^h$ is homogeneous of degree $-5$.

(5) The line $X$ is an elementary line of degree $n_X\in\{2,3\}$, so $X\in\mathfrak K(n_X)$ by Lemma~\ref{IV-73-39}(c). The bubble $Y=K_{\mathrm a}K_{\mathrm b}$ has $n_{\mathrm a},n_{\mathrm b}\in\{1,2,3\}$ and $n_Y\le4$, so $Y'\in\mathfrak K(n_Y)$ by Lemma~\ref{IV-73-46}(c). We write $X*Y=\mathsf q_Y*X+X*Y'$. Lemma~\ref{IV-73-46}(b) applies to $X$ and $Y'$, as $n_X,n_Y\le4$ and $n_X+n_Y=8-n_0\ge5$; it gives $X*Y'\in\mathfrak K(n_X+n_Y-3)$, and (K3) with $\gamma=0$ gives $X*Y'=(X*Y')^h+O(r'^{-(n_X+n_Y-3)-1}\lg r')$. Moreover $K_0=K_0^h+O(r'^{-n_0-1})$ by Lemma~\ref{IV-73-39}(c), and $K_0$, $K_0^h$, $X*Y'$, $(X*Y')^h$ are bounded by the corresponding powers of $r'$. As $n_0+(n_X+n_Y-3)=8-3=5$, the claim follows.

(4) Multiplying out $X*Y=(\mathsf q_X+X')*(\mathsf q_Y+Y')$ we obtain
\[
X*Y=\mathsf q_X*Y+\mathsf q_Y*X-\mathsf q_X*\mathsf q_Y+X'*Y',
\]
and $\mathsf q_X*\mathsf q_Y$ vanishes off $\{|u'|\le2\}$, because $\mathsf q_X$ and $\mathsf q_Y$ vanish off $\{0,\pm e_j\}$. If $(n_X,n_Y)=(4,4)$, then $X',Y'\in\mathfrak K(4)$ by Lemma~\ref{IV-73-46}(c), and Lemma~\ref{IV-73-46}(b) gives $X'*Y'\in\mathfrak K(5)$ and $X'*Y'=(X'*Y')^h+O(r'^{-6}\lg r')$. Otherwise we apply Lemma~\ref{IV-73-46}(a) to $X'$ and $Y'$, with $a_X=\min\{n_X-3,2\}$ and $a_Y=\min\{n_Y-3,2\}$. Its hypotheses hold by Lemma~\ref{IV-73-46}(c), which bounds the differences of a bubble of degree $n$ by $C_\gamma\langle x\rangle^{-n-|\gamma|}$ (the kernels $X',Y'$ differ from $X,Y$ at finitely many points only) and the symbol of $X'$, resp.\ $Y'$, by $|\mathsf I-\tilde{\mathsf Q}|\le C|k|^{\min\{n-3,2\}}$. The three exponents in the bound of Lemma~\ref{IV-73-46}(a) are $-2-4$, $-1-5$, $3-9$ for $(n_X,n_Y)=(5,4)$, the same by symmetry for $(4,5)$, and at most $-8$ for $(6,6)$; so $X'*Y'=O(r'^{-6}\lg r')$ in these cases.

(1) This is Proposition~\ref{IV-73-47}(iii) with $\varDelta=5$ (Lemma~\ref{IV-73-57}(1)), so that $r^{-\varDelta-1}\lg r=r'^{-6}\lg r'$.

($*$) In the case (6) this follows from Lemma~\ref{IV-73-39}(b). In the cases (5) and (4), Lemma~\ref{IV-73-46}(b) shows that the second component of the leading data of a convolution depends only on the leading data of the two factors, bilinearly and rotation covariantly. By Lemma~\ref{IV-73-46}(c) the leading data of $Y'$, and of $X'$ in the case (4), are $(\mathsf I^h,K_{\mathrm a}^hK_{\mathrm b}^h)$, where $\mathsf I^h$ depends only on $\sigma_{\mathrm a}^h,\sigma_{\mathrm b}^h$, bilinearly and rotation covariantly, by Lemma~\ref{IV-73-42}(a); the leading data of the lines $X$ and $K_0$ are $(\sigma^h,K^h)$ with $K^h=\mathcal F^{-1}_{\ne0}\sigma^h$, by Lemma~\ref{IV-73-39}(b),(c). The functions $K^h$ and the second components of leading data are smooth off the origin. In the case (1) the assertions are Proposition~\ref{IV-73-47}(iii),(iv).
\end{proof}

\begin{lemma}[A local insertion with vanishing moments]\label{IV-73-61}
\begin{enumerate}
\item[(a)] Let $\gamma_{\mu\nu}$, $1\le\mu,\nu\le3$, be functions of finite support on $\mathbb Z^3\times\mathbb Z^3$ whose symbols
\[
\hat\gamma_{\mu\nu}(\kappa_1,\kappa_2):=\sum_{y,y'}\gamma_{\mu\nu}(y,y')\,e^{-i(\kappa_1\cdot y+\kappa_2\cdot y')}
\]
vanish at $\kappa=0$ together with their gradients, and let $K_1,K_2,K_3$ be elementary lines of degree $2$. Then
\[
\Bigl|\sum_{\mu,\nu}\sum_{y,y'}\gamma_{\mu\nu}(y,y')\,K_\mu(u+y)\,K_\nu(u+y')\Bigr|\le C\,|u|^{-6}\qquad(u\in\mathbb Z^3\setminus\{0\}),
\]
where $C$ depends on the $\gamma_{\mu\nu}$ and on the symbols of the $K_\mu$.
\item[(b)] Let $\mathsf q$ be a function of finite support on $\mathbb Z^3$ with $\hat{\mathsf q}(k):=\sum_y\mathsf q(y)\,e^{-ik\cdot y}=O(|k|^2)$ at $k=0$, and let $K_1,K_2$ be elementary lines of degree $2$. Then $|(K_1*\mathsf q*K_2)(u)|\le C\langle u\rangle^{-3}$ for all $u\in\mathbb Z^3$, where $C$ depends on $\mathsf q$ and on the symbols of $K_1,K_2$.
\end{enumerate}
\end{lemma}

\begin{proof}
(a) As the $\gamma_{\mu\nu}$ have finite support, only points $y,y'$ in a fixed ball occur; let $|u|$ be large. By Lemma~\ref{IV-73-39}(a) with $s=-1$, $K_\mu(u)=O(|u|^{-2})$, $\nabla K_\mu(u)=O(|u|^{-3})$, and the second differences of $K_\mu$ near $u$ are $O(|u|^{-4})$. We write $K_\mu(u+y)-K_\mu(u)$ as a sum of first differences along a nearest-neighbour path from $u$ to $u+y$ and compare each of them with $\nabla_jK_\mu(u)$; this gives
\[
K_\mu(u+y)=K_\mu(u)+\sum_jy_j\,\nabla_jK_\mu(u)+O(|u|^{-4}).
\]
We insert this expansion for $K_\mu(u+y)$ and for $K_\nu(u+y')$. The terms of the orders $|u|^{-4}$ and $|u|^{-5}$ then carry the coefficients
\[
\sum_{y,y'}\gamma_{\mu\nu}(y,y')=\hat\gamma_{\mu\nu}(0),\qquad\sum_{y,y'}\gamma_{\mu\nu}(y,y')\,y_j,\qquad\sum_{y,y'}\gamma_{\mu\nu}(y,y')\,y'_j ,
\]
and the last two are $i$ times the derivatives of $\hat\gamma_{\mu\nu}$ with respect to $\kappa_{1,j}$, resp.\ $\kappa_{2,j}$, at $0$. All three vanish by hypothesis, and the remaining terms are $O(|u|^{-6})$. For $u\ne0$ in a bounded set the left side is bounded and $|u|^{-6}$ is bounded below, so the bound holds for all $u\ne0$ after enlarging $C$.

(b) The trigonometric polynomial $\hat{\mathsf q}$ is smooth and vanishes at $0$ together with its gradient, hence $\hat{\mathsf q}\in\mathcal S^2$. The kernel $\mathsf q*K_2$, a finite combination of translates of $K_2$, has the symbol $\hat{\mathsf q}\sigma_2\in\mathcal S^1$, where $\sigma_2\in\mathcal S^{-1}$ is the symbol of $K_2$; by Lemma~\ref{IV-73-39}(a) it is $O(\langle x\rangle^{-4})$ and lies in $\ell^2(\mathbb Z^3)$, as does $K_1=O(\langle x\rangle^{-2})$. By Plancherel's theorem $K_1*(\mathsf q*K_2)$ is the kernel with the symbol $\sigma_1\hat{\mathsf q}\sigma_2\in\mathcal S^0$, and Lemma~\ref{IV-73-39}(a) with $s=0$ gives the bound $C\langle u\rangle^{-3}$.
\end{proof}

\begin{lemma}[The local parts vanish to the second order]\label{IV-73-62}
Let $\tau_{\mathrm{gf}}>0$; we work on $\mathbb Z^3$.
\begin{enumerate}
\item[(a)] (Self-energy.) In \eqref{IV-73-59} let $\mathsf q_G$ be the kernel $\mathsf q$ of the bubble $K_{\mathrm a}K_{\mathrm b}$ in the sense of Lemma~\ref{IV-73-46}(c), and put
\[
\mathsf q^\Sigma_{\mu\nu}:=\mathsf q^0_{\mu\nu}+\sum_Gc_G\,\mathsf f_G(\nabla)\mathsf q_G,\qquad\mathsf W_{\mu\nu}:=\sum_Gc_G\,\mathsf f_G(\nabla)\bigl[K_{\mathrm a}K_{\mathrm b}-\mathsf q_G\bigr],
\]
so that $\Sigma_{\mu\nu}=\mathsf q^\Sigma_{\mu\nu}+\mathsf W_{\mu\nu}$. Then $\mathsf q^\Sigma_{\mu\nu}$ has finite support. Each $\mathsf f_G(\nabla)[K_{\mathrm a}K_{\mathrm b}-\mathsf q_G]$ belongs to $\mathfrak K(4)$, and the first component of its leading data is $\mathsf f_G^h\,\mathsf I_G^h$, where $\mathsf f_G^h$ is the leading part of $\mathsf f_G$ and $\mathsf I_G^h$ is the homogeneous part of the symbol $\mathsf I_G$ of the bubble (Lemma~\ref{IV-73-46}(c), Lemma~\ref{IV-73-42}(d)). Hence $\mathsf W_{\mu\nu}\in\mathfrak K(4)$, and its leading data are functions of the leading parts $\sigma^h$ and of the $\mathsf f_G^h$ occurring in \eqref{IV-73-59} alone. Moreover the symbol of $\mathsf q^\Sigma_{\mu\nu}$ satisfies $\hat{\mathsf q}^\Sigma_{\mu\nu}(k)=O(|k|^2)$ at $k=0$.
\item[(b)] (Vertex.) Let $q=p_{[\iota']}(x_q)$. Let $\gamma_q(b,b')$ be the symmetric kernel of finite range defined below; then, for every plaquette $p$ with $|u'|$ large, $\frac{\mathsf n}2\sum_{b,b'}\gamma_q(b,b')\,\mathsf E_p(b)\,\mathsf E_p(b')$ is the sum of the following terms of Lemma~\ref{IV-73-60}: the first half of $\mathsf T_{4'}$; the terms $c_G\,K_0(u')\,(\mathsf q_Y*X)(u')$ of the first half of $\mathsf T_5$; the terms $c_G\,(\mathsf q_Y*X)(u')$ of $\mathsf T_4$; and the terms $c_G\sum_{z,w}\mathsf q_{\mathsf q}(z,w)\,K_1(u'-z)\,K_2(u'-w)$ of $\mathsf T_1$. The kernel $\gamma_q$ is defined as follows. Each of $K_0$, $X$, $K_1$, $K_2$, at the points where it is evaluated in these terms, is either a value $\mathsf E_p(b)$ with $b$ a bond near $q$, or a value $\Pi(p,P)=\sum_is_i\,\mathsf E_p(b_i(P))$ with $P$ a plaquette near $q$.

Substituting this term by term in the four sums we obtain a quadratic expression in the values $\mathsf E_p(b)$, and $\frac{\mathsf n}2\gamma_q$ is its coefficient kernel, symmetrised. This kernel does not depend on $p$. Indeed, by the rule for $\mathsf a_p^{(2)}$ in the proof of Lemma~\ref{IV-73-53}(c), the plaquette $p$, with its base point and its dual direction, enters the four sums only through the forms $\mathsf E_p^a$ inserted into the two slots joined to $p$, that is, only through the values of $K_0$, $X$, $K_1$, $K_2$ just described. An elementary term $G$ of these sums is given by the monomials at the vertices at $q$, $P$, $P'$ and by the choice of their legs which are contracted with $p$, with $q$ and with one another, so that the numbers $c_G$ do not involve $p$; and the kernels $\mathsf q_Y$, $\mathsf q_{\mathsf q}$ and the line between two slots of the vertex at $q$ are built from lines among $q$, $P$, $P'$ alone.

We define $\gamma_q$ through this way of writing, since the numbers $\sum_{b,b'}\gamma_q(b,b')\mathsf E_p(b)\mathsf E_p(b')$ alone would not fix a kernel. Let $\gamma_{q;\mu\nu}(y,y'):=\gamma_q((x_q+y,\mu),(x_q+y',\nu))$. Then the symbols $\hat\gamma_{q;\mu\nu}$, defined as in Lemma~\ref{IV-73-61}(a), vanish at $\kappa=0$ together with their gradients.
\end{enumerate}
\end{lemma}

\begin{proof}
\emph{Step 1: a trace formula.} Let $\mathsf m^1,\dots,\mathsf m^{\mathsf r}$, $\mathsf r\le3$, be real symmetric matrices indexed by $\mathsf i=(x,\mu,a)$, that is, by a bond $(x,\mu)$ and a colour $a$, each with finitely many non-zero entries. For $y\in\mathbb Z^3$ let $\mathsf m_y$ be the translate of $\mathsf m$ by $y$, $(\mathsf m_y)_{(x,\mu,a),(x',\mu',a')}:=\mathsf m_{(x-y,\mu,a),(x'-y,\mu',a')}$, and let $\hat{\mathsf m}(l',l)$ be the $9\times9$ matrix, indexed by $(\mu,a)$ and $(\mu',a')$, with the entries
\[
\sum_{x,x'}e^{-il'\cdot x}\,\mathsf m_{(x,\mu,a),(x',\mu',a')}\,e^{il\cdot x'} .
\]
Then $\hat{\mathsf m}_y(l',l)=e^{-i(l'-l)\cdot y}\,\hat{\mathsf m}(l',l)$, and $\hat{\mathsf m}(l,l)$ is the symbol of the translation invariant matrix $\sum_y\mathsf m_y$, in the sense that $\sum_y\mathsf m_y$ maps $e^{il\cdot x}\alpha$ to $e^{il\cdot x}\,\hat{\mathsf m}(l,l)\,\alpha$. We insert $\mathsf C(\mathsf i,\mathsf i')=\delta_{aa'}\int_{\mathbb T}e^{il\cdot(x-x')}\,\hat{\mathsf C}_{\mu\mu'}(l)\,d^\circ l$ and let $\hat{\mathsf C}(l)$ act on the nine-dimensional space as $\hat{\mathsf C}(l)$ in the directions and as the identity in the colours. For $\mathsf r=3$ we obtain
\begin{equation}\label{IV-73-63}
\begin{aligned}
&\operatorname{Tr}\bigl[\mathsf m^1\mathsf C\,\mathsf m^2_{y_2}\mathsf C\,\mathsf m^3_{y_3}\mathsf C\bigr]\\
&\qquad=\iint e^{i(\kappa_2\cdot y_2+\kappa_3\cdot y_3)}\Bigl\{\int\operatorname{tr}\bigl[\hat{\mathsf m}^1(l_3,l)\,\hat{\mathsf C}(l)\,\hat{\mathsf m}^2(l,l_2)\,\hat{\mathsf C}(l_2)\\
&\qquad\qquad\qquad\hat{\mathsf m}^3(l_2,l_3)\,\hat{\mathsf C}(l_3)\bigr]\,d^\circ l\Bigr\}\,d^\circ\kappa_2\,d^\circ\kappa_3 ,
\end{aligned}
\end{equation}
with $l_2:=l+\kappa_2$, $l_3:=l+\kappa_2+\kappa_3$, all integrals being over $\mathbb T$. Indeed, the sums over the indices are finite and the integrals over the variables $l,l_2,l_3$ of the three covariances converge absolutely; Fubini's theorem allows us to carry out the sums first and then to pass to the variables $\kappa_2=l_2-l$, $\kappa_3=l_3-l_2$. For $\mathsf r=2$ one omits $\mathsf m^3$, the last factor $\hat{\mathsf C}$ and $\kappa_3$; for $\mathsf r=1$ one has $\operatorname{Tr}[\mathsf m^1\mathsf C]=\int\operatorname{tr}[\hat{\mathsf m}^1(l,l)\,\hat{\mathsf C}(l)]\,d^\circ l$. For the quadratic forms $\mathsf F_j:=\frac12(A,\mathsf m^jA)$, Lemma~\ref{IV-73-53}(a) gives
\[
\begin{gathered}
\langle\mathsf F_1\rangle_0=\tfrac12\operatorname{Tr}[\mathsf m^1\mathsf C],\qquad\langle\mathsf F_1;\mathsf F_2\rangle_0=\tfrac12\operatorname{Tr}[\mathsf m^1\mathsf C\,\mathsf m^2\mathsf C],\\
\langle\mathsf F_1;\mathsf F_2;\mathsf F_3\rangle_0=\operatorname{Tr}[\mathsf m^1\mathsf C\,\mathsf m^2\mathsf C\,\mathsf m^3\mathsf C].
\end{gathered}
\]

The formula \eqref{IV-73-63} also identifies integrands. For a leg $\ell$ of colour $a$ we define $\mathsf v_\ell(k)\in\mathbb C^3$ by $\ell(e^{ik\cdot x}\alpha\,\mathsf e_a)=\mathsf v_\ell(k)\cdot\alpha$, where $e^{ik\cdot x}\alpha\,\mathsf e_a$ denotes the form $(x,\mu)\mapsto e^{ik\cdot x}\alpha_\mu\mathsf e_a$; thus $\mathsf v_\ell(k)=e^{ik\cdot\xi}e_\nu$ for the plain leg at the bond $(\xi,\nu)$, and $\mathsf v_\ell(k)=e^{ik\cdot\xi}\bigl(\mathsf c_{\iota,\nu}(k)\bigr)_{\nu=1,2,3}$ for the curl leg at $p_{[\iota]}(\xi)$. The contraction of $\ell$ with the translate by $y$ of a leg $\ell'$ of the same colour is $\int e^{-ik\cdot y}\,\mathsf v_\ell(k)^{\mathsf T}\,\hat{\mathsf C}(k)\,\mathsf v_{\ell'}(-k)\,d^\circ k$, and the Hessian $\mathsf m$ of the product $\ell_1\ell_2$ of two legs has
\[
\hat{\mathsf m}(l',l)=\mathsf v_{\ell_1}(-l')\,\mathsf v_{\ell_2}(l)^{\mathsf T}+\mathsf v_{\ell_2}(-l')\,\mathsf v_{\ell_1}(l)^{\mathsf T},
\]
the colour indices being those of the legs. The matrices $\mathsf m^j$ used below are sums of Hessians of monomials in legs. Multiplying out, we find that for all $(l,\kappa)$ the integrand in braces in \eqref{IV-73-63} is a sum of products of the symbols of $\mathsf r$ lines evaluated at $\pm l,\pm l_2,\pm l_3$, and that these products are, term by term, the integrands of the one-loop functions (Lemma~\ref{IV-73-42}, Lemma~\ref{IV-73-46}(c), Proposition~\ref{IV-73-47}) of the elementary terms of Lemma~\ref{IV-73-57}, up to substitutions $l\to\pm l+\lambda(\kappa)$ with $\lambda$ linear. At $\kappa=0$ the integrand in braces is the sum of the integrands of these one-loop functions at $\kappa=0$, up to the substitution $l\to-l$ in single terms. We express this by saying that the two sides agree integrand by integrand.

\emph{Step 2: the Hessians at a constant abelian field.} By \eqref{IV-73-7} with $t=1$, $a_p(e^{iA})=\sum_m\mathsf a_p^{(m)}(A)$. Let $\mathsf N_j(B)$ and $\mathsf h_q^{(j)}(B)$ be the parts of degree $j$ in $B$ of the Hessians $\mathsf N(B)$ and $\mathsf h_q(B)$ of Definition~\ref{IV-73-29}(1). Since $\mathsf S_3$ is cubic and $\mathsf S_4$ quartic, $\mathsf N_1(B)$ is the Hessian of the quadratic polynomial $\sum_{b,a}B^a(b)\,\partial_{b,a}\mathsf S_3$, and $\mathsf N_2(B)$ is that of $\frac12\sum B^a(b)\,B^{a'}(b')\,\partial_{b,a}\partial_{b',a'}\mathsf S_4$; the matrix $\mathsf h_q^{(j)}(B)$ is obtained in the same way from $\mathsf a_q^{(2+j)}$.

Let $\mathsf H_b$ be the Hessian of $\partial_{b,3}\mathsf S_3$, and let $\mathsf H_b^{\mathsf u}$ and $\mathsf H_b^{\mathsf d}$ be those of $\mathsf u_{b,3}$ and $\mathsf d_{b,3}$ in \eqref{IV-73-58}; these matrices are translation covariant, $\mathsf H_{b+y}=(\mathsf H_b)_y$. For the constant field $Y_\gamma$ of Definition~\ref{IV-73-29}(2), with $\mu(b)$ the direction of $b$, we have
\[
\sum_b\gamma_{\mu(b)}\,\mathsf d_{b,3}=\sum_P\langle\mathrm dY_\gamma(P),\mathfrak z_1(P)\rangle=0
\]
(as a quadratic form, every entry of its Hessian being a finite sum), because $\sum_b\gamma_{\mu(b)}\,\mathsf s_P(b)\,\mathsf e_3=\mathrm dY_\gamma(P)$ and $\mathrm dY_\gamma=0$. Hence $\mathsf N_1(Y_\gamma)=\sum_b\gamma_{\mu(b)}\,\mathsf H_b^{\mathsf u}=\sum_\mu\gamma_\mu\sum_y(\mathsf H^{\mathsf u}_{(0,\mu)})_y$, and by Step~1 and Definition~\ref{IV-73-29}(2) its symbol is
\[
\sum_\mu\gamma_\mu\,\hat{\mathsf H}^{\mathsf u}_{(0,\mu)}(l,l)=\bigoplus_e\hat{\mathsf N}_{e,1}(l;\gamma)
\]
in the basis $\mathsf e_e$ of the colours.

\emph{Step 3: the structure in (a).} Let $K\in\mathfrak K(n)$ have the symbol $\sigma$ and the leading data $(\sigma^h,K^h)$, and let $y_0\in\mathbb Z^3$. The kernels with the symbols $e^{ik\cdot y_0}\sigma$, $\mathsf k_\mu\sigma$ and $\overline{\mathsf k_\mu}\sigma$ lie in $\mathfrak K(n)$, $\mathfrak K(n+1)$ and $\mathfrak K(n+1)$, respectively, with the leading data $(\sigma^h,K^h)$, $(ik_\mu\sigma^h,\partial_\mu K^h)$ and $(-ik_\mu\sigma^h,-\partial_\mu K^h)$. Indeed, the symbol $\mathsf k_\mu\sigma$ belongs to the kernel $\nabla_\mu K$. For (K2) we write
\[
\mathsf k_\mu\sigma-\chi\,ik_\mu\sigma^h=\mathsf k_\mu(\sigma-\chi\sigma^h)+(\mathsf k_\mu-ik_\mu)\,\chi\sigma^h=O(|k|^{n-1}\lg|k|),
\]
and for (K3), by Taylor's formula and the homogeneity of the derivatives of $K^h$,
\[
\nabla^\gamma(\nabla_\mu K-\partial_\mu K^h)=\nabla^{\gamma+e_\mu}(K-K^h)+\nabla^\gamma(\nabla_\mu-\partial_\mu)K^h=O(|x|^{-n-2-|\gamma|}\lg|x|);
\]
at the finitely many points with $|\gamma|+1\le|x|<|\gamma|+2$ the left side is a finite number, which is absorbed into the constant. Next, $\overline{\mathsf k_\mu}=-e^{-ik_\mu}\mathsf k_\mu$. For the shift, $(e^{ik\cdot y_0}-1)\,\chi\sigma^h=O(|k|^{n-2})$ and $\nabla^\gamma[K^h(\cdot+y_0)-K^h]=O(|x|^{-n-1-|\gamma|})$. By Lemma~\ref{IV-73-46}(c), $K_{\mathrm a}K_{\mathrm b}-\mathsf q_G\in\mathfrak K(n_{\mathrm a}+n_{\mathrm b})$, with the leading data $(\mathsf I_G^h,K_{\mathrm a}^hK_{\mathrm b}^h)$; that lemma applies, since $n_{\mathrm a},n_{\mathrm b}\in\{1,2,3\}$ and $n_{\mathrm a}+n_{\mathrm b}=4-d_G\le4$. A trigonometric polynomial $\mathsf f$ is a finite sum of terms $c\,e^{ik\cdot y_0}$.

As a line symbol of degree $0$ it has the leading part $\mathsf f(0)$. As a line symbol of degree $1$ it has $\mathsf f(0)=0$, its leading part is its linear part at $0$, and $\mathsf f$ is a finite sum of terms $c\,e^{ik\cdot y_0}\mathsf k_\mu$: indeed $\mathsf f$ is a linear combination of the functions $e^{ik\cdot y}-1$, and each of these is a telescoping sum, along a nearest-neighbour path from $0$ to $y$, of terms $\pm e^{ik\cdot y_0}\mathsf k_\mu$ (with $\overline{\mathsf k_\mu}=-e^{-ik_\mu}\mathsf k_\mu$ for the steps in the negative direction). Since $\mathfrak K(n)$ is a linear space and the leading data depend linearly on the kernel, multiplication of the symbol by $\mathsf f$ maps $\mathfrak K(n)$ into $\mathfrak K(n)$, resp.\ $\mathfrak K(n+1)$, and multiplies the first component of the leading data by $\mathsf f^h$.

Applying this once for each of the two factors of $\mathsf f_G$, we obtain $\mathsf f_G(\nabla)[K_{\mathrm a}K_{\mathrm b}-\mathsf q_G]\in\mathfrak K(n_{\mathrm a}+n_{\mathrm b}+d_G)=\mathfrak K(4)$, with $\mathsf f_G^h\,\mathsf I_G^h$ as the first component of its leading data. Hence $\mathsf W_{\mu\nu}\in\mathfrak K(4)$. The first component of its leading data, $\sum_Gc_G\,\mathsf f_G^h\,\mathsf I_G^h$, depends only on the leading parts $\sigma^h$ (Lemma~\ref{IV-73-42}(a)) and on the $\mathsf f_G^h$; by Lemma~\ref{IV-73-46}(d) the second component is determined by the first, linearly. Finally $\mathsf q^\Sigma_{\mu\nu}$ has finite support, since $\mathsf q^0_{\mu\nu}$ and the $\mathsf q_G$ have and $\mathsf f_G(\nabla)$ is a finite difference operator.

\emph{Step 4: the value of $\hat{\mathsf q}^\Sigma_{\mu\nu}$ at $0$.} The symbol of $\mathsf q_G$ is $\tilde{\mathsf Q}_G$, whose value at $0$ is $\mathsf Q_G(0)$, and $\mathsf f_G(0)=0$ if $d_G\ge1$. Hence
\[
\hat{\mathsf q}^\Sigma_{\mu\nu}(0)=\hat{\mathsf q}^0_{\mu\nu}(0)+\sum_{G:\,d_G=0}c_G\,\mathsf f_G(0)\,\mathsf Q_G(0).
\]
For $d_G=0$ we have $n_{\mathrm a}+n_{\mathrm b}=4$, so that the bubble symbol has $\omega=1$, and Lemma~\ref{IV-73-42}(c), with the data of the bubble given in the proof of Lemma~\ref{IV-73-46}(c), yields $\mathsf Q_G(0)=\int_{\mathbb T}\sigma_{\mathrm a}(-l)\,\sigma_{\mathrm b}(l)\,d^\circ l$. The terms with $d_G=0$ are the ghost loop and the terms of $-\langle\mathsf u_{b,3};\mathsf u_{b',3}\rangle_0=-\frac12\operatorname{Tr}[\mathsf H_b^{\mathsf u}\mathsf C\,\mathsf H_{b'}^{\mathsf u}\mathsf C]$ (Step~1). We compute $\frac12\sum_{\mu,\nu}\gamma_\mu\gamma_\nu$ times each contribution to $\hat{\mathsf q}^\Sigma_{\mu\nu}(0)$.
\begin{itemize}
\item Gluon loop. By \eqref{IV-73-63}, $-\frac12\operatorname{Tr}[\mathsf H^{\mathsf u}_{(0,\mu)}\mathsf C\,\mathsf H^{\mathsf u}_{(y,\nu)}\mathsf C]$ is the kernel in $y=b'_--b_-$ with the symbol $\kappa_2\mapsto-\frac12\int\operatorname{tr}[\hat{\mathsf H}^{\mathsf u}_{(0,\mu)}(l_2,l)\,\hat{\mathsf C}(l)\,\hat{\mathsf H}^{\mathsf u}_{(0,\nu)}(l,l_2)\,\hat{\mathsf C}(l_2)]\,d^\circ l$, and this symbol is $\sum c_G\,\mathsf f_G\,\mathsf I_G$, summed over the corresponding $G$, integrand by integrand. At zero momentum the integrand is $-\frac12\operatorname{tr}[\hat{\mathsf H}^{\mathsf u}_{(0,\mu)}(l,l)\,\hat{\mathsf C}(l)\,\hat{\mathsf H}^{\mathsf u}_{(0,\nu)}(l,l)\,\hat{\mathsf C}(l)]$, and by Step~2
\[
\begin{aligned}
&\tfrac12\sum_{\mu,\nu}\gamma_\mu\gamma_\nu\Bigl(-\tfrac12\int\operatorname{tr}\bigl[\hat{\mathsf H}^{\mathsf u}_{(0,\mu)}(l,l)\,\hat{\mathsf C}\,\hat{\mathsf H}^{\mathsf u}_{(0,\nu)}(l,l)\,\hat{\mathsf C}\bigr]\,d^\circ l\Bigr)\\
&\qquad=-\tfrac14\sum_e\int\operatorname{tr}\bigl[(\hat{\mathsf C}\hat{\mathsf N}_{e,1})^2\bigr]\,d^\circ l=\int_{\mathbb T}\mathfrak s^{(1)}(l;\gamma)\,d^\circ l .
\end{aligned}
\]
\item Gluon tadpole. Let $\mathsf H_{bb'}$ be the Hessian of $\partial_{b,3}\partial_{b',3}\mathsf S_4$, and let $\mathsf m:=\frac12\sum_{\mu,\nu}\sum_y\gamma_\mu\gamma_\nu\,\mathsf H_{(0,\mu),(y,\nu)}$, a finite sum. By Step~2 and translation covariance $\mathsf N_2(Y_\gamma)=\sum_x\mathsf m_x$, hence $\hat{\mathsf m}(l,l)=\bigoplus_e\hat{\mathsf N}_{e,2}(l;\gamma)$, and Step~1 with $\mathsf r=1$ gives
\[
\begin{aligned}
\tfrac12\sum_{\mu,\nu}\gamma_\mu\gamma_\nu\sum_y\langle\partial_{(0,\mu),3}\partial_{(y,\nu),3}\mathsf S_4\rangle_0&=\tfrac12\operatorname{Tr}[\mathsf m\,\mathsf C]\\
&=\tfrac12\int\operatorname{tr}\bigl[\hat{\mathsf m}(l,l)\,\hat{\mathsf C}(l)\bigr]\,d^\circ l=\int_{\mathbb T}\mathfrak s^{(2)}(l;\gamma)\,d^\circ l .
\end{aligned}
\]
\item Haar term: $\frac12\cdot\frac23|\gamma|^2=\frac13|\gamma|^2=\int\mathfrak s^{(5)}$. Ghost tadpole: $\frac12\cdot\frac43\cdot\frac13|\gamma|^2=\frac29|\gamma|^2=\int\mathfrak s^{(4)}$, since by Lemma~\ref{IV-73-31}(c) $\int\mathfrak s^{(4)}=\sum_{e=\pm1}\frac4{12}\sum_\mu\gamma_\mu^2\int|\mathsf k_\mu|^2|\mathsf k|^{-2}\,d^\circ k$, and $\int|\mathsf k_\mu|^2|\mathsf k|^{-2}\,d^\circ k=\frac13$.
\item Ghost loop. Here $\mathfrak s^{(3)}=4\bigl(\sum_\mu\gamma_\mu\sin k_\mu\bigr)^2|\mathsf k|^{-4}$. By Lemma~\ref{IV-73-51} the symbols of $y\mapsto\mathsf D((x,\mu);(x+y,\nu))$ and of $y\mapsto\mathsf D((x+y,\nu);(x,\mu))$ are $\frac12(1+e^{ik_\nu})\,\overline{\mathsf k_\mu}\,|\mathsf k|^{-2}$ and its image under $\mu\leftrightarrow\nu$, $k\to-k$. Since $(1+e^{ik_\nu})\,\overline{\mathsf k_\nu}=-2i\sin k_\nu$, the ghost loop $G$, for which $c_G=-8$ and $\mathsf f_G=1$, has $\mathsf Q_G(0)=\int\frac14(-2i\sin l_\mu)(-2i\sin l_\nu)\,|\mathsf k(l)|^{-4}\,d^\circ l$, and $\frac12\sum_{\mu,\nu}\gamma_\mu\gamma_\nu\,(-8)\,\mathsf Q_G(0)=\int\mathfrak s^{(3)}$.
\end{itemize}
Summing up, we find $\frac12\sum_{\mu,\nu}\gamma_\mu\gamma_\nu\,\hat{\mathsf q}^\Sigma_{\mu\nu}(0)=\sum_{i=1}^5\int\mathfrak s^{(i)}=\frac12\sum_{\mu,\nu}\bigl(\sum_ic^{(i)}_{\mu\nu}\bigr)\gamma_\mu\gamma_\nu$ for small $\gamma$ (Definition~\ref{IV-73-29}(3)); hence $\hat{\mathsf q}^\Sigma_{\mu\nu}(0)+\hat{\mathsf q}^\Sigma_{\nu\mu}(0)=2\sum_ic^{(i)}_{\mu\nu}$, and this number vanishes, because by Theorem~\ref{IV-73-34} the five matrices $c^{(i)}$ add up to zero, $\sum_{i=1}^5c^{(i)}_{\mu\nu}=0$ for all $\mu,\nu$. It remains to show that $\hat{\mathsf q}^\Sigma_{\mu\nu}(0)=\hat{\mathsf q}^\Sigma_{\nu\mu}(0)$. Let $\hat\Sigma_{\mu\nu}:=\hat{\mathsf q}^\Sigma_{\mu\nu}+\hat{\mathsf W}_{\mu\nu}$ be the symbol of $\Sigma_{\mu\nu}$.

Since $\Sigma_{\mathrm{tot}}$ is symmetric and translation covariant, $\Sigma_{\mu\nu}(y)=\Sigma_{\nu\mu}(-y)$, and the uniqueness of Fourier coefficients gives $\hat\Sigma_{\mu\nu}(k)=\hat\Sigma_{\nu\mu}(-k)$ for almost every $k$. As $\hat{\mathsf q}^\Sigma$ is continuous and $\hat{\mathsf W}(k)\to0$ for $k\to0$, by (K2) for $\mathsf W\in\mathfrak K(4)$, letting $k\to0$ gives the symmetry, and therefore $\hat{\mathsf q}^\Sigma_{\mu\nu}(0)=0$.

\emph{Step 5: the gradient of $\hat{\mathsf q}^\Sigma_{\mu\nu}$ at $0$.} The function $\hat\Sigma_{\mu\nu}$ is integrable on $\mathbb T$, with the Fourier coefficients $\Sigma_{\mu\nu}(y)$. Lemma~\ref{IV-73-37}(c) with $c=0$, for which $(x,\mu)^R=(-x-e_\mu,\mu)$, gives
\[
\Sigma_{\mu\nu}(y)=\Sigma_{\mathrm{tot}}\bigl((-e_\mu,\mu),(-y-e_\nu,\nu)\bigr)=\Sigma_{\mu\nu}(-y+e_\mu-e_\nu),
\]
hence $\hat\Sigma_{\mu\nu}(k)=e^{-ik\cdot(e_\mu-e_\nu)}\,\hat\Sigma_{\mu\nu}(-k)$ for almost every $k$, by the uniqueness of Fourier coefficients. With $\varphi(k):=\frac12(k_\mu-k_\nu)$ the function $\tilde\Sigma(k):=e^{i\varphi(k)}\hat\Sigma_{\mu\nu}(k)$ is therefore even. By (K2) for $\mathsf W_{\mu\nu}\in\mathfrak K(4)$, $\hat{\mathsf W}_{\mu\nu}=\hat{\mathsf W}^h+O(|k|^2\lg|k|)$ near $0$, where $\hat{\mathsf W}^h=\sum_Gc_G\,\mathsf f_G^h\,\mathsf I_G^h$ is homogeneous of degree $1$ and even: by Lemma~\ref{IV-73-42}(a) the parity of $\mathsf f_G^h\mathsf I_G^h$ is $(-1)^{d_G}(-1)^{\omega_G+1}$, where $\omega_G=n_{\mathrm a}+n_{\mathrm b}-3$, and $\omega_G+d_G=1$.

Consequently the odd part of the smooth function $e^{i\varphi}\hat{\mathsf q}^\Sigma_{\mu\nu}=\tilde\Sigma-e^{i\varphi}\hat{\mathsf W}_{\mu\nu}$ equals the odd part of $-e^{i\varphi}\hat{\mathsf W}^h+O(|k|^2\lg|k|)$, which is $O(|k|^2\lg|k|)$, because $\hat{\mathsf W}^h$ is even and $(e^{i\varphi}-1)\hat{\mathsf W}^h=O(|k|^2)$. So the gradient of $e^{i\varphi}\hat{\mathsf q}^\Sigma_{\mu\nu}$ at $0$ vanishes; as $\hat{\mathsf q}^\Sigma_{\mu\nu}(0)=0$ by Step~4, the gradient of $\hat{\mathsf q}^\Sigma_{\mu\nu}$ at $0$ vanishes as well. Since $\hat{\mathsf q}^\Sigma_{\mu\nu}$ is a trigonometric polynomial, this proves $\hat{\mathsf q}^\Sigma_{\mu\nu}(k)=O(|k|^2)$ and completes the proof of (a).

\emph{Step 6: set-up for (b).} Let $\Gamma_{q;\mu\nu}(y,y'):=\Gamma_q((x_q+y,\mu),(x_q+y',\nu))$. In the rest of the proof $c_G$ denotes $\frac2{\mathsf n}$ times the number $c_G$ of Lemma~\ref{IV-73-57}; the reason is that the sums with the numbers $c_G$ of that lemma equal $\frac{\mathsf n}2\Gamma_q$, by \eqref{IV-73-54}, resp.\ $\frac{\mathsf n}2\gamma_q$, by the definition of $\gamma_q$ in (b). By the first, second, third and fourth summand we mean the four groups of summands of $\Gamma_q$ as in \eqref{IV-73-54}. In each summand the derivatives $\partial_{b,a}$ and $\partial_{b',a'}$ of Definition~\ref{IV-73-29}(5) act on two legs of the vertices at $q$, $P$, $P'$; we call these two legs the outer legs.

In the terms on the right of \eqref{IV-73-54} the outer legs are the legs which are contracted with the two curl legs $(\mathrm dA)(p)$, namely the legs at which the lines $K_1$, $K_2$ of Lemma~\ref{IV-73-57}(4$'$), (4) and (1), resp.\ the lines $K_0$, $X$ of Lemma~\ref{IV-73-57}(5), end. An outer leg is either plain or a curl leg; in the second case we speak of an outer curl leg. The symbol of a kernel on $\mathbb Z^6$, such as $\Gamma_{q;\mu\nu}$, is defined as in Lemma~\ref{IV-73-61}(a), so that $\kappa_1$ is the variable dual to $y$ and $\kappa_2$ the variable dual to $y'$; accordingly we say that the outer leg on which $\partial_{b,a}$ acts, $b=(x_q+y,\mu)$, carries the momentum $\kappa_1$, and that the outer leg on which $\partial_{b',a'}$ acts, $b'=(x_q+y',\nu)$, carries the momentum $\kappa_2$.

By \eqref{IV-73-54} and Lemma~\ref{IV-73-57}, $\Gamma_{q;\mu\nu}$ is a finite sum of kernels on $\mathbb Z^6$ with the symbols
\[
c_G\,\mathsf f_{G,1}(\kappa_1)\,\mathsf f_{G,2}(\kappa_2)\,\mathsf J_G(\kappa),
\]
symmetrised under $(\mu,\kappa_1)\leftrightarrow(\nu,\kappa_2)$. Here the $\mathsf f_{G,i}$ are trigonometric polynomials, line symbols of degree $0$ if the corresponding outer leg is plain and of degree $1$ if it is a curl leg. Further, $\mathsf J_G=1$ for the first summand; $\mathsf J_G$ is the bubble symbol $\mathsf I_G(\pm\kappa_2)$ of $Y$ for the second summand, the outer leg at $q$ carrying the momentum $\kappa_1$; $\mathsf J_G=\mathsf I_G(\pm(\kappa_1+\kappa_2))$ for the third summand, both outer legs sitting at $P$; and $\mathsf J_G=\hat{\mathsf T}_{\mathsf q}(\pm\kappa)$ for the fourth summand, the Fourier representation being \eqref{IV-73-48}.

By the definition of $\gamma_q$ in (b), the symbol $\hat\gamma_{q;\mu\nu}$ is the same sum with $\mathsf J_G$ replaced by $\tilde{\mathsf Q}_G$, where $\tilde{\mathsf Q}_G:=1$ for the first summand, which is entirely local (the line between two slots of the vertex at $q$ is a number, a factor of $c_G$), and otherwise $\tilde{\mathsf Q}_G$ is the function $\tilde{\mathsf Q}$ of Lemma~\ref{IV-73-46}(c), resp.\ $\tilde{\mathsf Q}_{\mathsf q}$ of Proposition~\ref{IV-73-47}, taken at the argument of $\mathsf J_G$. Indeed, each of the local terms listed in (b) arises from the corresponding full term of Lemma~\ref{IV-73-57} by replacing a single kernel by its local kernel, while the lines which end at the outer legs are left as they are.

In the first half of $\mathsf T_5$ the full term is $K_0(u')\sum_zY(z)\,X(u'-z)$, and the local term has $\mathsf q_Y$ in place of $Y$, in accordance with $X*Y=\mathsf q_Y*X+X*Y'$; in $\mathsf T_4$ the full term is $\sum_zY(z)\,(K_1K_2)(u'-z)$, and again $Y$ is replaced by $\mathsf q_Y$. In $\mathsf T_1$ the substitution $(z,w)\to(u'-z,u'-w)$ turns the full term into
\[
\sum_{z,w}\bigl[K_4(z)\,K_5(w)\,K_3(z-w)\bigr]\,K_1(u'-z)\,K_2(u'-w),
\]
where the bracket, a kernel in $(z,w)$, has the symbol $\hat{\mathsf T}_{\mathsf q}$ by the second identity of \eqref{IV-73-48}; the local term has $\mathsf q_{\mathsf q}(z,w)$, with the symbol $\tilde{\mathsf Q}_{\mathsf q}$, in place of the bracket. The first half of $\mathsf T_{4'}$ is entirely local and is its own local term. In each case the symbol $\mathsf J_G$ of the replaced kernel turns into $\tilde{\mathsf Q}_G$, while the number $c_G$, the factors $\mathsf f_{G,1}$, $\mathsf f_{G,2}$ and the argument of $\mathsf J_G$ remain the same. Let $\hat\Gamma:=\hat\Gamma_{q;\mu\nu}$ be the sum of the symbols of $\Gamma_{q;\mu\nu}$; in $\hat\Gamma-\hat\gamma_{q;\mu\nu}$ the first summand cancels.

\emph{Step 7: counting degrees.} Let $d_{G,i}\in\{0,1\}$ be the degree of $\mathsf f_{G,i}$ and $d_G:=d_{G,1}+d_{G,2}$, the number of outer curl legs; for $G$ in the second, third and fourth summand let $\omega_G$ be the number $\omega$ of $\mathsf J_G$. We use again that the degree of a line is $1+\#$(its curl legs) (Lemma~\ref{IV-73-51}), and that $\mathsf a^{(2)}$ has two curl legs and every monomial of $\mathsf a^{(3)}$ one (Definition~\ref{IV-73-19}(3)). In the second summand the vertices at $q$ and $P$ have two curl legs, of which $2-d_G$ lie on the two bubble lines, so that $\omega_G=2+(2-d_G)-3=1-d_G$. In the fourth summand the vertices at $q$, $P$, $P'$ have four curl legs, of which $4-d_G$ lie on the three lines of $\hat{\mathsf T}_{\mathsf q}$, so that $\omega_G=3+(4-d_G)-6=1-d_G$.

In the third summand the bubble lines carry the two curl legs at $q$: for the monomials of $\frac12|\mathfrak z_1|^2$ we get $(d_G,\omega_G)=(0,1)$; for the monomials of $\mathsf a_P^{(4)}$ with one curl leg we get $(1,1)$ or $(0,2)$, according as this leg is outer or not; for those with four curl legs we get $(2,3)$. Thus $d_G+\omega_G=1$, except for these last three kinds, for which $d_G+\min\{\omega_G,2\}\ge2$ and $\omega_G\ge1$.

\emph{Step 8: expansion of $\hat\Gamma-\hat\gamma_{q;\mu\nu}$.} We show that for almost every $\kappa$ with $|\kappa|\le c$, $c$ small,
\[
\begin{gathered}
\hat\Gamma-\hat\gamma_{q;\mu\nu}=\hat\Gamma^h+O\bigl(|\kappa|^2\lg|\kappa|\,\Lambda(\kappa)\bigr),\\
\hat\Gamma^h:=\sum_{G:\,d_G+\omega_G=1}c_G\,\mathsf f_{G,1}^h\,\mathsf f_{G,2}^h\,\mathsf J_G^h,\qquad|\hat\Gamma^h|\le C|\kappa|\,\Lambda(\kappa).
\end{gathered}
\]
Here $\mathsf J_G^h$ is the homogeneous part of $\mathsf J_G$ (Lemma~\ref{IV-73-42}(d)), taken at the argument of $\mathsf J_G$, and $\Lambda\ge1$ is the function $\Lambda$ of Lemma~\ref{IV-73-42} for $\hat{\mathsf T}_{\mathsf q}$. The function $\hat\Gamma^h$ is homogeneous of degree $1$ and even, each of its terms having the parity $(-1)^{d_G}(-1)^{\omega_G+1}=1$. Let $\kappa'$ denote the argument of $\mathsf J_G$; then $|\kappa'|\le\sqrt2|\kappa|$. For $d_G+\omega_G\ge2$, Lemma~\ref{IV-73-46}(c) gives
\[
\bigl|\mathsf f_{G,1}\mathsf f_{G,2}\,(\mathsf J_G-\tilde{\mathsf Q}_G)\bigr|\le C|\kappa|^{d_G}\,|\kappa'|^{\min\{\omega_G,2\}}\le C|\kappa|^2 .
\]
For the other $G$ we have $\omega_G\le1$, hence $\deg\mathsf Q_G\le1$ and $\mathsf Q_G-\tilde{\mathsf Q}_G=O(|\kappa'|^3)$, and Lemma~\ref{IV-73-42}(b) gives
\[
\mathsf J_G-\tilde{\mathsf Q}_G=\mathsf J_G^h+O\bigl(|\kappa'|^{\omega_G+1}\lg|\kappa'|\,\Lambda_G(\kappa')\bigr)+O(|\kappa'|^3),
\]
where $\Lambda_G$ is bounded for the bubbles and equals $\Lambda$ for $\hat{\mathsf T}_{\mathsf q}$. Moreover $\mathsf f_{G,i}(\kappa_i)=\mathsf f_{G,i}^h(\kappa_i)+O(|\kappa_i|^{d_{G,i}+1})$, and the functions $t\lg t$ and $t^2\lg t$ are increasing on $(0,1]$, their derivatives being $-\log t$ and $t(1-2\log t)$. We go through the summands.
\begin{itemize}
\item Third summand, $\omega_G=1$, $d_G=0$: $|\mathsf J_G^h|\le C|\kappa'|\le C|\kappa|$, and the errors are $O(|\kappa|^2\lg|\kappa|)$.
\item Fourth summand: $\kappa'=\pm\kappa$, $|\mathsf J_G^h|\le C|\kappa|^{\omega_G}\Lambda$ and $|\mathsf f_{G,1}\mathsf f_{G,2}|\le C|\kappa|^{d_G}$; the leading term is $O(|\kappa|\Lambda)$, and the errors are $O(|\kappa|^2\lg|\kappa|\,\Lambda)$.
\item Second summand: $\kappa'=\pm\kappa_2$, and $\omega_G=-1$ occurs. We take the factor of the outer leg at $P$, which carries the momentum $\kappa_2$, together with $\mathsf J_G$. As $d_{G,2}+\omega_G=1-d_{G,1}\ge0$, we have $|\mathsf f^h_{G,2}(\kappa_2)\,\mathsf J_G^h|\le C|\kappa_2|^{1-d_{G,1}}$ and $|\mathsf f^h_{G,1}\mathsf f^h_{G,2}\,\mathsf J_G^h|\le C|\kappa_1|^{d_{G,1}}|\kappa_2|^{1-d_{G,1}}\le C|\kappa|$. Replacing $\mathsf f_{G,i}$ by its leading part costs a factor $|\kappa_i|$, and
\[
|\mathsf f_{G,1}\mathsf f_{G,2}|\,|\kappa_2|^{\omega_G+1}\lg|\kappa_2|\le C|\kappa_1|^{d_{G,1}}\,|\kappa_2|^{2-d_{G,1}}\lg|\kappa_2|\le C|\kappa|^2\lg|\kappa| .
\]
\end{itemize}
The symmetrised terms are treated in the same way, with the indices $1$ and $2$ exchanged.

\emph{Step 9: the value of $\hat\gamma_{q;\mu\nu}$ at $0$.} We have $\hat\gamma_{q;\mu\nu}(0)=\sum c_G\,\mathsf f_{G,1}(0)\,\mathsf f_{G,2}(0)\,\tilde{\mathsf Q}_G(0)$, the sum extending over the $G$ with two plain outer legs, as an outer curl leg gives $\mathsf f_{G,i}(0)=0$. For such $G$ in the first summand, $\tilde{\mathsf Q}_G=1$, and the line between two slots of the vertex at $q$, contained in $c_G$, is the integral of its symbol (the case $\mathsf r=1$ of Step~1). For such $G$ in the other summands, $d_G=0$, hence $\omega_G\ge1$ by Step~7, and $\tilde{\mathsf Q}_G(0)=\mathsf Q_G(0)$ is the integral of the integrand at $\kappa=0$ (Lemma~\ref{IV-73-42}(c)). By Step~1 the four summands of $\Gamma_q(b,3;b',3)$ have the trace forms
\[
\begin{gathered}
\tfrac12\operatorname{Tr}[\mathsf m_{bb'}\mathsf C],\qquad-\tfrac12\operatorname{Tr}[\mathsf m_b\mathsf C\,\mathsf H_{b'}\mathsf C]-(b\leftrightarrow b'),\\
-\tfrac12\operatorname{Tr}[\mathsf h_q^{(0)}\mathsf C\,\mathsf H_{bb'}\mathsf C],\qquad\operatorname{Tr}[\mathsf h_q^{(0)}\mathsf C\,\mathsf H_b\mathsf C\,\mathsf H_{b'}\mathsf C],
\end{gathered}
\]
where $\mathsf m_{bb'}$, $\mathsf m_b$ and $\mathsf H_{bb'}$ are the Hessians of $\partial\partial'\mathsf a_q^{(4)}$, $\partial\mathsf a_q^{(3)}$ and $\partial\partial'\mathsf S_4$ ($\partial=\partial_{b,3}$, $\partial'=\partial_{b',3}$), $\mathsf H_b$ is the Hessian of $\partial_{b,3}\mathsf S_3$ (Step~2), and $\mathsf h_q^{(0)}$ is the Hessian of $\mathsf a_q^{(2)}$. By Step~2,
\[
\begin{gathered}
\sum_{b,b'}\tfrac12\gamma_{\mu(b)}\gamma_{\mu(b')}\,\mathsf m_{bb'}=\mathsf h_q^{(2)}(Y_\gamma),\qquad\sum_b\gamma_{\mu(b)}\,\mathsf m_b=\mathsf h_q^{(1)}(Y_\gamma),\\
\sum_\mu\gamma_\mu\,\hat{\mathsf H}_{(0,\mu)}(l,l)=\bigoplus_e\hat{\mathsf N}_{e,1}(l;\gamma),\\
\sum_{\mu,\nu}\sum_y\tfrac12\gamma_\mu\gamma_\nu\,\hat{\mathsf H}_{(0,\mu),(y,\nu)}(l,l)=\bigoplus_e\hat{\mathsf N}_{e,2}(l;\gamma),
\end{gathered}
\]
the sums in the first line being finite; and $\hat{\mathsf h}(l,l)$ is the matrix of the form $(e^{il\cdot x}\alpha'\mathsf e_e,\mathsf h\,e^{il\cdot x}\alpha\,\mathsf e_e)$, which for $\mathsf h=\mathsf h_q^{(j)}(Y_\gamma)$ is $\bigoplus_e\hat{\mathsf h}_{q,e,j}(l;\gamma)$ by Definition~\ref{IV-73-29}(2). The parts of the Hessians which come from an outer curl leg drop out here, since $\mathrm dY_\gamma=0$; they belong to the $G$ with $\mathsf f_{G,1}(0)\,\mathsf f_{G,2}(0)=0$. We use \eqref{IV-73-63} at $\kappa=0$, integrand by integrand, and write $\hat{\mathsf h}_j:=\bigoplus_e\hat{\mathsf h}_{q,e,j}$ and $\hat{\mathsf N}_j:=\bigoplus_e\hat{\mathsf N}_{e,j}$. The four summands contribute to $\frac12\sum_{\mu,\nu}\gamma_\mu\gamma_\nu\,\hat\gamma_{q;\mu\nu}(0)$ the amounts
\[
\begin{gathered}
\tfrac12\!\int\!\operatorname{tr}\bigl[\hat{\mathsf h}_2\hat{\mathsf C}\bigr],\qquad\tfrac12\cdot2\cdot\bigl(-\tfrac12\bigr)\!\int\!\operatorname{tr}\bigl[\hat{\mathsf h}_1\hat{\mathsf C}\hat{\mathsf N}_1\hat{\mathsf C}\bigr],\\
-\tfrac12\!\int\!\operatorname{tr}\bigl[\hat{\mathsf h}_0\hat{\mathsf C}\hat{\mathsf N}_2\hat{\mathsf C}\bigr],\qquad\tfrac12\!\int\!\operatorname{tr}\bigl[\hat{\mathsf h}_0\hat{\mathsf C}\hat{\mathsf N}_1\hat{\mathsf C}\hat{\mathsf N}_1\hat{\mathsf C}\bigr]
\end{gathered}
\]
(integrals over $\mathbb T$ with respect to $d^\circ l$), and these are $\int\mathfrak v_q^{(j)}(l;\gamma)\,d^\circ l$, $j=1,\dots,4$, by Definition~\ref{IV-73-29}(4). Hence, by Definition~\ref{IV-73-29}(4) and Theorem~\ref{IV-73-36},
\[
\tfrac12\sum_{\mu,\nu}\gamma_\mu\gamma_\nu\,\hat\gamma_{q;\mu\nu}(0)=\tfrac12\sum_{\mu,\nu}\gamma_\mu\gamma_\nu\sum_{j=1}^4\mathsf g^{(j)}_{q;\mu\nu}=0
\]
for all small $\gamma$. Since $\gamma_q$ is a symmetric kernel, $\hat\gamma_{q;\mu\nu}(0)$ is symmetric in $\mu,\nu$, and therefore $\hat\gamma_{q;\mu\nu}(0)=0$.

\emph{Step 10: the gradient of $\hat\gamma_{q;\mu\nu}$ at $0$.} Lemma~\ref{IV-73-37}(c) with $c=x(q)$, for which $R_cq=q$, gives
\[
\Gamma_{q;\mu\nu}(y,y')=\Gamma_{q;\mu\nu}(\mathsf o-e_\mu-y,\ \mathsf o-e_\nu-y'),\qquad\mathsf o:=2\bigl(x(q)-x_q\bigr),
\]
and the uniqueness of Fourier coefficients gives $\hat\Gamma(\kappa)=e^{-2i\varphi(\kappa)}\,\hat\Gamma(-\kappa)$ for almost every $\kappa$, with the real linear form $\varphi(\kappa):=\frac12[\kappa_1\cdot(\mathsf o-e_\mu)+\kappa_2\cdot(\mathsf o-e_\nu)]$; hence $e^{i\varphi}\hat\Gamma$ is even almost everywhere. Since $\hat\Gamma^h$ is even, the odd part of $e^{i\varphi}\hat\Gamma^h$ is $O(|\kappa|^2\Lambda)$. By Step~8 the odd part of the smooth function $e^{i\varphi}\hat\gamma_{q;\mu\nu}$ is therefore $O(|\kappa|^2\lg|\kappa|\,\Lambda)$ almost everywhere near $0$. The function $\Lambda$ is bounded on an open cone. Indeed, for $\hat{\mathsf T}_{\mathsf q}$ we have $(l_1,l_2,l_3)(\kappa)=(\kappa_1,-\kappa_2,0)$ (Step~2 of the proof of Proposition~\ref{IV-73-47}), so that the differences $l_i(\kappa)-l_j(\kappa)$, $i<j$, are $\kappa_1+\kappa_2$, $\kappa_1$ and $-\kappa_2$. Let
\[
\mathsf K:=\bigl\{\kappa\in\mathbb R^6:\ \min\{|\kappa_1|,|\kappa_2|,|\kappa_1+\kappa_2|\}>\tfrac12|\kappa|\bigr\}.
\]
This set is open, it is invariant under $\kappa\to t\kappa$ for real $t\ne0$, and it is not empty, since it contains $(v,v)$ for every $v\in\mathbb R^3\setminus\{0\}$: there $|\kappa|=\sqrt2\,|v|$, while $|\kappa_1|=|\kappa_2|=|v|$ and $|\kappa_1+\kappa_2|=2|v|$. On $\mathsf K$ each of the three terms $\log_+\bigl(|\kappa|/|l_i(\kappa)-l_j(\kappa)|\bigr)$ in the definition of $\Lambda$ is at most $\log2$, so that $\Lambda\le1+3\log2$ there. On $\mathsf K$ the odd part is therefore $O(|\kappa|^2\lg|\kappa|)$ almost everywhere near $0$. The odd part of the smooth function $e^{i\varphi}\hat\gamma_{q;\mu\nu}$ differs from its linear part $\kappa\mapsto\kappa\cdot\nabla(e^{i\varphi}\hat\gamma_{q;\mu\nu})(0)$ by $O(|\kappa|^3)$.

Hence $|\kappa\cdot\nabla(e^{i\varphi}\hat\gamma_{q;\mu\nu})(0)|\le C|\kappa|^2\lg|\kappa|$ for almost every, and by continuity for every, $\kappa\in\mathsf K$ with $|\kappa|\le c$. Replacing $\kappa$ by $t\kappa$ and letting $t\downarrow0$, we see that this linear function vanishes on $\mathsf K$, a non-empty open set, and hence everywhere: the gradient of $e^{i\varphi}\hat\gamma_{q;\mu\nu}$ at $0$ vanishes. As $\hat\gamma_{q;\mu\nu}(0)=0$ by Step~9, the gradient of $\hat\gamma_{q;\mu\nu}$ at $0$ vanishes as well. This proves (b).
\end{proof}

\subsection{The six terms, the continuum model, the kernel of degree $-5$}

In this subsection we introduce six continuum two-loop integrands and a model built from them, with a smooth ultraviolet cut-off on every line. We show that the model has the same parts of degree $-5$ as the lattice classes of Lemma~\ref{IV-73-57}, that these parts are multiples of $|u|^{-5}$, and we deduce the large-distance asymptotics of $\mathsf B_1^{\mathrm{lat}}$ and its consequence for smeared plaquette observables.

\begin{definition}[The six two-loop integrands; the model]\label{IV-73-64}
\emph{(1) Tensors and variables.} For $k_1,k_2,k_3\in\mathbb R^3$ we put
\begin{gather*}
V_2(k_1,k_2)_{\mu\nu}:=-(k_1\cdot k_2)\,\delta_{\mu\nu}+k_{1\nu}k_{2\mu},\\
V_3(k_1,k_2,k_3)_{\alpha\beta\gamma}:=\delta_{\alpha\gamma}(k_1-k_3)_\beta+\delta_{\alpha\beta}(k_2-k_1)_\gamma+\delta_{\beta\gamma}(k_3-k_2)_\alpha .
\end{gather*}
The tensor $V_2$ is even and the tensor $V_3$ is odd under the simultaneous change of sign $k_i\to-k_i$ of all arguments, and $|V_2(k_1,k_2)_{\mu\nu}|\le2|k_1||k_2|$ for all $\mu,\nu$. For $k,l,\mathsf p\in\mathbb R^3$ we write
\[
\mathsf x_0:=\mathsf p^2,\quad\mathsf x_1:=k^2,\quad\mathsf x_1':=(k-\mathsf p)^2,\quad\mathsf x_2:=l^2,\quad\mathsf x_2':=(l-\mathsf p)^2,\quad\mathsf x_3:=(k-l)^2 .
\]

\emph{(2) The six integrands.} With summation over repeated tensor indices we define
\begin{equation}\label{IV-73-65}
\begin{aligned}
F_1&:=-\tfrac12\,V_2(k,\mathsf p-k)_{\mu\nu}\,V_2(-l,l-\mathsf p)_{\rho\sigma}\\
&\qquad\qquad\times V_3(-k,l,k-l)_{\mu\rho\alpha}\,V_3(k-\mathsf p,\mathsf p-l,l-k)_{\nu\sigma\alpha}\,/\,\mathsf x_1\mathsf x_1'\mathsf x_2\mathsf x_2'\mathsf x_3,\\
F_2&:=-\tfrac12\,V_2(k,\mathsf p-k)_{\mu\nu}\,V_2(-k,k-\mathsf p)_{\rho\nu}\\
&\qquad\qquad\times V_3(-k,l,k-l)_{\mu\beta\gamma}\,V_3(k,-l,l-k)_{\rho\beta\gamma}\,/\,\mathsf x_1^2\mathsf x_1'\mathsf x_2\mathsf x_3,\\
F_3&:=-\,V_2(k,\mathsf p-k)_{\mu\nu}\,V_2(-k,k-\mathsf p)_{\rho\nu}\,l_\rho(l-k)_\mu\,/\,\mathsf x_1^2\mathsf x_1'\mathsf x_2\mathsf x_3,\\
F_4&:=-\tfrac12\bigl[V_2(k,\mathsf p-k)_{\mu\mu}V_2(l,\mathsf p-l)_{\nu\nu}\\
&\qquad\qquad-\tfrac12V_2(k,\mathsf p-k)_{\mu\nu}\bigl(V_2(l,\mathsf p-l)_{\mu\nu}+V_2(l,\mathsf p-l)_{\nu\mu}\bigr)\bigr]/\,\mathsf x_1\mathsf x_1'\mathsf x_2\mathsf x_2',\\
F_5&:=V_3(\mathsf p-k,l,k-l)_{\alpha\beta\gamma}\,V_3(k,-l,l-k)_{\sigma\beta\gamma}\,V_2(k-\mathsf p,-k)_{\alpha\sigma}\,/\,\mathsf x_1'\mathsf x_1\mathsf x_2\mathsf x_3,\\
F_6&:=-\tfrac16\,V_3(l,k-l,\mathsf p-k)_{\alpha\beta\gamma}\,V_3(-l,l-k,k-\mathsf p)_{\alpha\beta\gamma}\,/\,\mathsf x_2\mathsf x_3\mathsf x_1' .
\end{aligned}
\end{equation}
We call $F_1$ the kite, $F_2$ the gluon loop on a line, $F_3$ the ghost loop on a line, $F_4$ the quartic vertex, $F_5$ the term of $\mathsf a^{(3)}$ with $\mathsf S_3$, and $F_6$ the term of $\mathsf a^{(3)}$ with $\mathsf a^{(3)}$. Each $F_\Gamma$, $\Gamma\in\{1,\dots,6\}$, is a rational function of $(k,l,\mathsf p)$, homogeneous of degree $-4$, even under $(k,l,\mathsf p)\to-(k,l,\mathsf p)$, and invariant under simultaneous rotations of $k$, $l$ and $\mathsf p$. The \emph{lines of $\Gamma$} are the vectors whose squares occur in the denominator of $F_\Gamma$, counted with multiplicity: for $\Gamma=1$ the vectors $k,\mathsf p-k,l,\mathsf p-l,k-l$; for $\Gamma=2,3$ the vector $k$ twice and the vectors $\mathsf p-k,l,k-l$; for $\Gamma=4$ the vectors $k,\mathsf p-k,l,\mathsf p-l$; for $\Gamma=5$ the vectors $k,\mathsf p-k,l,k-l$; for $\Gamma=6$ the vectors $\mathsf p-k,l,k-l$.

\emph{(3) The model.} Let $\rho_0:=\frac\pi{16}$, and let $\rho(\mathsf q):=\theta(|\mathsf q|/2\rho_0)$ for $\mathsf q\in\mathbb R^3$, with the function $\theta$ of Definition~\ref{IV-73-38}(1); thus $\rho(\mathsf q)=1$ for $|\mathsf q|\le\rho_0$ and $\rho(\mathsf q)=0$ for $|\mathsf q|\ge2\rho_0$. Let $R_\Gamma(k,l,\mathsf p)$ be the product of the values of $\rho$ at the lines of $\Gamma$. We define
\[
\Psi_\Gamma(\mathsf p):=\iint R_\Gamma F_\Gamma\,d^\circ k\,d^\circ l\quad(\mathsf p\in\mathbb R^3\setminus\{0\}),\qquad
\mathsf T^{\mathrm{mod}}_\Gamma(u):=24\int_{\mathbb R^3}e^{i\mathsf p\cdot u}\,\Psi_\Gamma(\mathsf p)\,d^\circ\mathsf p\quad(u\in\mathbb Z^3),
\]
the first integral extending over $\mathbb R^3\times\mathbb R^3$.

These integrals converge absolutely. Indeed, every entry of $V_3(k_1,k_2,k_3)$ is bounded by $C\max_i|k_i|$; together with the bound on the entries of $V_2$ and with the denominators this gives, on every bounded set of points $(k,l,\mathsf p)$ and with a constant $C$ depending on that set,
\begin{equation}\label{IV-73-66}
\begin{gathered}
|F_1|\le\frac C{|k|\,|\mathsf p-k|\,|l|\,|\mathsf p-l|\,|k-l|^2},\qquad|F_2|,\ |F_3|\le\frac{C(|k|+|l|)}{|k|^2|l|^2|k-l|^2},\\
|F_4|\le\frac C{|k|\,|\mathsf p-k|\,|l|\,|\mathsf p-l|},\qquad|F_5|\le\frac C{|k|\,|\mathsf p-k|\,|l|^2|k-l|^2},\qquad
|F_6|\le\frac C{|\mathsf p-k|^2|l|^2|k-l|^2}.
\end{gathered}
\end{equation}
We integrate these majorants first over $l$ and then over $k$. For $v,w$ in a bounded subset of $\mathbb R^3$, and with integrals over a bounded set of $l$, we have the elementary bounds
\[
\int|l-v|^{-2}|l-w|^{-2}\,dl\le C|v-w|^{-1},\qquad\int|l-v|^{-1}|l-w|^{-2}\,dl\le C\lg|v-w| ,
\]
with constants depending on these sets; moreover $|l|^{-1}|\mathsf p-l|^{-1}\le|\mathsf p|^{-1}\bigl(|l|^{-1}+|\mathsf p-l|^{-1}\bigr)$, since $|\mathsf p|\le|l|+|\mathsf p-l|$, and $|k|+|l|\le2|k|+|k-l|$. For $\Gamma=1$ the majorant is at most $C|\mathsf p|^{-2}\bigl(|k|^{-1}+|\mathsf p-k|^{-1}\bigr)\bigl(|l|^{-1}+|\mathsf p-l|^{-1}\bigr)|k-l|^{-2}$; its integral over $l$ is at most $C|\mathsf p|^{-2}\bigl(|k|^{-1}+|\mathsf p-k|^{-1}\bigr)\bigl(\lg|k|+\lg|k-\mathsf p|\bigr)$, which is integrable over bounded sets of $k$, with a bound uniform for $\mathsf p$ in a bounded set after the factor $|\mathsf p|^{-2}$ is removed. For $\Gamma=4$ the same argument applies without the factor $|k-l|^{-2}$.

For $\Gamma=2,3$ the majorant is at most $2C|k|^{-1}|l|^{-2}|k-l|^{-2}+C|k|^{-2}|l|^{-2}|k-l|^{-1}$, whose integral over $l$ is at most $C|k|^{-2}\lg|k|$, a function integrable over bounded sets of $k$. For $\Gamma=5$ the integral over $l$ is at most $C|k|^{-2}|\mathsf p-k|^{-1}$, and the subsequent integral over $k$ is at most $C\lg|\mathsf p|$. For $\Gamma=6$ the integral over $l$ is at most $C|k|^{-1}|\mathsf p-k|^{-2}$, and again the integral over $k$ is at most $C\lg|\mathsf p|$. Hence, for $\mathsf p$ in a bounded set, the integral of each majorant over a bounded set of $(k,l)$ is at most $C|\mathsf p|^{-2}$, since $1$ and $\lg|\mathsf p|$ are bounded by $C|\mathsf p|^{-2}$ there.

Where $R_\Gamma\ne0$, every line of $\Gamma$ has modulus less than $2\rho_0$. Hence $|l|<2\rho_0$ and $|k|<4\rho_0$ there (by the lines $l$ and $k$, and for $\Gamma=6$ by the lines $l$ and $k-l$), and $|\mathsf p|\le|\mathsf p-k|+|k|<6\rho_0$, by the line $\mathsf p-k$, which occurs for every $\Gamma$. Consequently the integral defining $\Psi_\Gamma(\mathsf p)$ converges absolutely, $|\Psi_\Gamma(\mathsf p)|\le C|\mathsf p|^{-2}$, and $\Psi_\Gamma$ is integrable, radial (by the rotation invariance of $F_\Gamma$ and of $R_\Gamma$, and the substitution $(k,l)\to(Ok,Ol)$), and equal to $0$ for $|\mathsf p|\ge6\rho_0=\frac{3\pi}8$; in particular $\mathsf T^{\mathrm{mod}}_\Gamma(u)$ is well defined.
\end{definition}

In order to compare the model with the elementary terms of Lemma~\ref{IV-73-57}, we write these terms, too, as integrals over momenta.

\begin{lemma}[Momentum-space forms]\label{IV-73-67}
\textup{(a)} Let $K_1,\dots,K_5$ be elementary lines as in Proposition~\ref{IV-73-47}, with symbols $\sigma_1,\dots,\sigma_5$, and let $\mathsf T$ be their kite. Then, for every $u\in\mathbb Z^3$,
\[
\mathsf T(u)=\int_{\mathbb T}e^{i\mathsf p\cdot u}\Bigl[\iint_{\mathbb T^2}\sigma_1(k)\,\sigma_2(\mathsf p-k)\,\sigma_3(k-l)\,\sigma_4(l)\,\sigma_5(\mathsf p-l)\,d^\circ k\,d^\circ l\Bigr]d^\circ\mathsf p ,
\]
and the triple integral converges absolutely.

\textup{(b)} If two kernels have symbols $\sigma_{\mathrm a},\sigma_{\mathrm b}\in L^1(\mathbb T)$, then their product is the kernel whose symbol is the convolution $k\mapsto\int_{\mathbb T}\sigma_{\mathrm a}(k-l)\,\sigma_{\mathrm b}(l)\,d^\circ l$ of the two symbols. If $X,Y\in\ell^2(\mathbb Z^3)$ are kernels with symbols $\hat X,\hat Y\in L^1(\mathbb T)$, then $X*Y$ \textup{(}Lemma~\ref{IV-73-46}\textup{)} is the kernel with the symbol $\hat X\hat Y$.
\end{lemma}

\begin{proof}
We begin with (b). Let $K_{\mathrm a},K_{\mathrm b}$ be the two kernels. The function $(k,l)\mapsto\sigma_{\mathrm a}(k)\sigma_{\mathrm b}(l)$ is integrable over $\mathbb T^2$; by Fubini's theorem and the substitution $k\to k-l$ on $\mathbb T$,
\[
K_{\mathrm a}(x)\,K_{\mathrm b}(x)=\iint e^{i(k+l)\cdot x}\,\sigma_{\mathrm a}(k)\,\sigma_{\mathrm b}(l)\,d^\circ k\,d^\circ l=\int_{\mathbb T}e^{ik\cdot x}\Bigl[\int_{\mathbb T}\sigma_{\mathrm a}(k-l)\,\sigma_{\mathrm b}(l)\,d^\circ l\Bigr]d^\circ k ,
\]
and the bracket is an integrable function of $k$. For the second assertion we note that $X\in\ell^2(\mathbb Z^3)$ is, by Plancherel's theorem, the sequence of Fourier coefficients of a function in $L^2(\mathbb T)$, which coincides with $\hat X$ by the uniqueness of Fourier coefficients; thus $\hat X\in L^2(\mathbb T)$, and likewise $\hat Y\in L^2(\mathbb T)$. The kernel $z\mapsto Y(u-z)$ has the symbol $k\mapsto e^{-ik\cdot u}\hat Y(-k)$, and Plancherel's theorem in the form $\sum_za(z)b(z)=\int_{\mathbb T}\hat a(-k)\,\hat b(k)\,d^\circ k$, for square summable kernels $a,b$ with symbols $\hat a,\hat b$, gives $(X*Y)(u)=\int_{\mathbb T}e^{ik\cdot u}\hat X(k)\hat Y(k)\,d^\circ k$.

We turn to (a).

\emph{Step 1.} For $i\ne3$ the degree $n_i$ is at least $2$, so that $|\sigma_i(\mathsf q)|\le C|\mathsf q|^{-1}$ on $[-\pi,\pi]^3$ (Definition~\ref{IV-73-38}(2),(3)); hence $\sigma_i\in L^2(\mathbb T)$ for $i\ne3$, while $\sigma_3\in L^1(\mathbb T)$. The functions
\[
\mathsf A(v):=\int_{\mathbb T}\sigma_1(v+\mathsf q)\,\sigma_4(\mathsf q)\,e^{iu\cdot\mathsf q}\,d^\circ\mathsf q,\qquad
\mathsf B(v):=\int_{\mathbb T}\sigma_2(v+\mathsf q)\,\sigma_5(\mathsf q)\,e^{iu\cdot\mathsf q}\,d^\circ\mathsf q
\]
are bounded, by the Cauchy--Schwarz inequality, and continuous, since translations act continuously on $L^2(\mathbb T)$. Consequently
\[
\Phi(\eta_1,\eta_2):=\int_{\mathbb T}\sigma_3(\mathsf q)\,\mathsf A(\eta_1+\mathsf q)\,\mathsf B(\eta_2-\mathsf q)\,d^\circ\mathsf q
\]
is bounded and continuous on $\mathbb T^2$, by dominated convergence, and the integral over $(\mathsf q_3,\mathsf q_4,\mathsf q_5)\in\mathbb T^3$ which results from inserting the definitions of $\mathsf A$ and $\mathsf B$ (with $\mathsf q_3$ in place of $\mathsf q$ and $\mathsf q_4,\mathsf q_5$ the variables of $\mathsf A,\mathsf B$) converges absolutely.

\emph{Step 2.} Let $\phi\in C_c^\infty(\mathbb R^3)$ with support in $\{|x|\le1\}$ and $\phi(0)=1$, and let $\phi_R:=\phi(\cdot/R)$. By Proposition~\ref{IV-73-47}, the double sum defining $\mathsf T(u)$ converges absolutely; by dominated convergence, $\mathsf T(u)=\lim_{R\to\infty}\mathsf T_R(u)$, where $\mathsf T_R(u)$ denotes the same sum with the additional factor $\phi_R(z)\phi_R(w)$. For fixed $R$ we insert $K_i(y)=\int_{\mathbb T}e^{i\mathsf q_i\cdot y}\sigma_i(\mathsf q_i)\,d^\circ\mathsf q_i$. The sums over $z$ and $w$ are then finite, and by Fubini's theorem they produce the factor
\[
\check\phi_R\bigl(-(\mathsf q_1-\mathsf q_3-\mathsf q_4)\bigr)\,\check\phi_R\bigl(-(\mathsf q_2+\mathsf q_3-\mathsf q_5)\bigr)\,e^{iu\cdot(\mathsf q_4+\mathsf q_5)} .
\]
We pass to the variables $\eta_1:=\mathsf q_1-\mathsf q_3-\mathsf q_4$, $\eta_2:=\mathsf q_2+\mathsf q_3-\mathsf q_5$, $\mathsf q_3$, $\mathsf q_4$, $\mathsf q_5$; this is a linear change of variables with integer matrix of determinant $\pm1$, hence a measure preserving bijection of $\mathbb T^5$, and it gives
\[
\mathsf T_R(u)=\iint_{\mathbb T^2}\check\phi_R(-\eta_1)\,\check\phi_R(-\eta_2)\,\Phi(\eta_1,\eta_2)\,d^\circ\eta_1\,d^\circ\eta_2 .
\]
Here $\int_{\mathbb T}\check\phi_R(-\eta)\,d^\circ\eta=\phi_R(0)=1$, and $|\check\phi_R(\eta)|\le CR^3\langle R\eta\rangle^{-10}$ by Lemma~\ref{IV-73-41}(a), applied with $D=3$ to the restriction of $\phi_R$ to $\mathbb Z^3$; this function vanishes off $\{|\zeta|\le R\}$, and its norm $\|\phi_R\|_{(R)}$ is at most $\|\phi\|_{C^{10}}$, a difference being an integral of a derivative. Thus the functions $\eta\mapsto\check\phi_R(-\eta)$ have integral $1$, bounded $L^1$ norms, and integrals over $\{|\eta|\ge\delta\}$ tending to $0$ as $R\to\infty$, for every $\delta>0$. As $\Phi$ is bounded and continuous at $(0,0)$, $\mathsf T_R(u)\to\Phi(0,0)$.

\emph{Step 3.} In the triple integral $\Phi(0,0)$ we change variables to $l=\mathsf q_4$, $k=\mathsf q_3+\mathsf q_4$, $\mathsf p=\mathsf q_4+\mathsf q_5$, again a measure preserving bijection of $\mathbb T^3$. Then $\mathsf q_3=k-l$, $\mathsf q_5=\mathsf p-l$ and $\mathsf q_5-\mathsf q_3=\mathsf p-k$, so that the integrand becomes $\sigma_1(k)\,\sigma_2(\mathsf p-k)\,\sigma_3(k-l)\,\sigma_4(l)\,\sigma_5(\mathsf p-l)\,e^{i\mathsf p\cdot u}$. The triple integral converges absolutely by Step~1, and Fubini's theorem yields the formula of (a).
\end{proof}

The next lemma computes the weights with which the leading parts of the elementary terms add up to the six integrands of Definition~\ref{IV-73-64}.

\begin{lemma}[The weights of the six terms]\label{IV-73-68}
Let $\tau_{\mathrm{gf}}=1$, and let $p=p_{[\iota]}(x_p)$, $q=p_{[\iota']}(x_q)$. To an elementary term $G$ of Lemma~\ref{IV-73-57}, with its lines named as there, we attach its \emph{leading integrand} $\mathsf F_G(k,l,\mathsf p)$, where $\sigma^h$ denotes the leading part of the symbol of the line in question:
\[
\begin{array}{ll}
G\in\mathfrak G_1: & \sigma_1^h(k)\,\sigma_2^h(\mathsf p-k)\,\sigma_3^h(k-l)\,\sigma_4^h(l)\,\sigma_5^h(\mathsf p-l);\\[3pt]
G\in\mathfrak G_4,\ (n_X,n_Y)=(4,4): & \sigma_1^h(k)\,\sigma_2^h(\mathsf p-k)\,\sigma_3^h(l)\,\sigma_4^h(\mathsf p-l);\\[3pt]
G\in\mathfrak G_5: & \sigma_0^h(\mathsf p-k)\,\sigma_X^h(k)\,\sigma_{\mathrm a}^h(k-l)\,\sigma_{\mathrm b}^h(l);\\[3pt]
G\in\mathfrak G_6: & \sigma_1^h(\mathsf p-k)\,\sigma_2^h(k-l)\,\sigma_3^h(l);\\[3pt]
G\in\mathfrak G_{\mathrm I}(\mu,\nu): & \sigma_\Pi^h(\mathsf p-k)\,\sigma^h_{\to,\iota\mu}(k)\,\mathsf f_G^h(k)\,\sigma_{\mathrm a}^h(k-l)\,\sigma_{\mathrm b}^h(l)\,\sigma^h_{\leftarrow,\nu\iota'}(k);
\end{array}
\]
here $\sigma_\Pi$, $\sigma_\to$ and $\sigma_\leftarrow$ are the symbols of $y\mapsto\Pi\bigl(p_{[\iota]}(0),p_{[\iota']}(y)\bigr)$, of $K^\to_{\iota\mu}$ and of $K^\leftarrow_{\nu\iota'}$. Let $\mathfrak G_2=\mathfrak G_2(\iota,\iota')$, resp.\ $\mathfrak G_3=\mathfrak G_3(\iota,\iota')$, be the set of the triples $(\mu,\nu,G)$ such that $G\in\mathfrak G_{\mathrm I}(\mu,\nu)$ is a term of $-\langle\partial\mathsf S_3;\partial'\mathsf S_3\rangle_0$, resp.\ of the ghost loop, in \eqref{IV-73-59}; such a triple carries the coefficient $-\mathsf n\,c_G$, which takes the place of $c_G$ in the sums below. Then, with sums over $\iota,\iota'\in\{1,2,3\}$ and over $G\in\mathfrak G_\Gamma(\iota,\iota')$, where for $\Gamma=4$ only the terms with $(n_X,n_Y)=(4,4)$ are taken,
\[
\sum_{\iota,\iota'}\sum_G c_G\,\mathsf F_G=24\,F_\Gamma\quad(\Gamma=1,2,3,4,6),\qquad\sum_{\iota,\iota'}\sum_{G\in\mathfrak G_5}c_G\,\mathsf F_G=12\,F_5 .
\]
\end{lemma}

\begin{proof}
\emph{Step 1 (Feynman rules).} Let $K(y)=\int e^{i\mathsf q\cdot y}\sigma(\mathsf q)\,d^\circ\mathsf q$ be a line, where $y$ is the base point of its second leg minus that of its first leg; we say that the line carries the momentum $\mathsf q$ from the first vertex to the second. The momentum entering a vertex through this line is then $\mathsf q$ at the second vertex and $-\mathsf q$ at the first. Leading parts of line symbols are additive and multiplicative (Definition~\ref{IV-73-38}(3)); hence the sum over the monomials of a vertex may be carried out on the leading parts. By Lemma~\ref{IV-73-51} at $\tau_{\mathrm{gf}}=1$, every line contributes a factor $\delta_{\mu\nu}|\mathsf q|^{-2}$, where at a plain leg the index of the line is the direction of the bond of that leg; a curl leg at a plaquette $p_{\alpha\beta}$ contributes the factor $i(\mathsf q'_\alpha\delta_{\nu\beta}-\mathsf q'_\beta\delta_{\nu\alpha})$, where $\mathsf q'$ is the momentum entering the vertex; and shifts of base points contribute the factor $1$, the leading part of $e^{ik\cdot y_0}$.

\emph{Step 2 (the vertex $\mathsf a^{(2)}$).} Two curl legs at $p_{[\iota]}$, with entering momenta $k_1,k_2$ and line indices $\mu,\nu$, give after summation over $\iota$, that is, over $\alpha<\beta$,
\[
\begin{aligned}
&i^2\sum_{\alpha<\beta}(k_{1\alpha}\delta_{\mu\beta}-k_{1\beta}\delta_{\mu\alpha})(k_{2\alpha}\delta_{\nu\beta}-k_{2\beta}\delta_{\nu\alpha})\\
&\qquad=-\bigl[(k_1\cdot k_2)\,\delta_{\mu\nu}-k_{1\nu}k_{2\mu}\bigr]=V_2(k_1,k_2)_{\mu\nu} .
\end{aligned}
\]

\emph{Step 3 (the vertex $\mathsf a^{(3)}$).} By Lemma~\ref{IV-73-53}(d), $\mathsf a^{(3)}_P[\![X,Y,Z]\!]$ is the sum of $\mathsf t_P$ over the six orders of $X,Y,Z$. Let $P=p_{\mu\nu}$. To leading order, $(Y(b_1),\dots,Y(b_4))$ is replaced by $(Y_\mu,Y_\nu,Y_\mu,Y_\nu)$; hence
\[
\sum_{i<j}s_is_j\,Y(b_i)\times Z(b_j)\to Y_\mu\times Z_\nu-Y_\nu\times Z_\mu-Y_\mu\times Z_\mu-Y_\nu\times Z_\nu ,
\]
and, by the antisymmetry of the cross product,
\[
\mathsf t_P(X,Y,Z)+\mathsf t_P(X,Z,Y)\to-2\bigl\langle(\mathrm dX)_{\mu\nu},\,Y_\mu\times Z_\nu-Y_\nu\times Z_\mu\bigr\rangle .
\]
Both factors are antisymmetric in $\mu\nu$. Summing over $\mu<\nu$ and inserting $(\mathrm dX)_{\mu\nu}\to i(k_{1\mu}X_\nu-k_{1\nu}X_\mu)$, we obtain
\[
-2i\sum_{\mu,\nu}k_{1\mu}\bigl\langle X_\nu,\,Y_\mu\times Z_\nu-Y_\nu\times Z_\mu\bigr\rangle .
\]
For legs of colours $a,b,c$, with indices $\alpha,\beta,\gamma$ and entering momenta $k_1,k_2,k_3$, this equals $-2i\,\epsilon_{abc}\bigl[\delta_{\alpha\gamma}k_{1\beta}-\delta_{\alpha\beta}k_{1\gamma}\bigr]$. We add the terms in which the second, resp.\ the third, leg is the curl leg; they arise by cyclic permutation, under which $\epsilon_{abc}$ is invariant. The leading part of $\sum_\iota\mathsf a^{(3)}_{p_{[\iota]}}[\![\cdot,\cdot,\cdot]\!]$ is therefore
\[
-2i\,\epsilon_{abc}\,V_3(k_1,k_2,k_3)_{\alpha\beta\gamma} .
\]

\emph{Step 4 ($\Gamma=1$).} By Lemma~\ref{IV-73-67}(a) the five lines $1,\dots,5$ of the kite carry $k$ (from $p$ to $P$), $\mathsf p-k$ (from $p$ to $P'$), $k-l$ (from $P$ to $P'$), $l$ (from $P$ to $q$) and $\mathsf p-l$ (from $P'$ to $q$). The entering momenta are $(-k,k-\mathsf p)$ at $p$; $(k,-l,l-k)$ at $P$, for the legs towards $p$, $q$, $P'$; $(\mathsf p-k,l-\mathsf p,k-l)$ at $P'$; and $(l,\mathsf p-l)$ at $q$. At both cubic vertices the colours are $(a,a',c)$, and $\sum_{a,a',c}\epsilon_{aa'c}^2=6$. With the factor $\frac12$ of $\mathsf T_1$ in Lemma~\ref{IV-73-53}(c) we obtain
\[
\begin{aligned}
&\tfrac12(-2i)^2\cdot6\cdot V_2(-k,k-\mathsf p)_{\mu\nu}\,V_2(l,\mathsf p-l)_{\rho\sigma}\\
&\qquad\times V_3(k,-l,l-k)_{\mu\rho\alpha}\,V_3(\mathsf p-k,l-\mathsf p,k-l)_{\nu\sigma\alpha}\,/\,\mathsf x_1\mathsf x_1'\mathsf x_2\mathsf x_2'\mathsf x_3 ,
\end{aligned}
\]
and $\frac12(-4)\cdot6=24\cdot(-\frac12)$. Since $V_2$ is even, and the product of the two factors $V_3$ is even under the change of sign of all their arguments, this is $24F_1$.

\emph{Step 5 ($\Gamma=6$).} The factor is $\frac16(-2i)^2\cdot6=24\cdot(-\frac16)$. The entering momenta are $-(\mathsf p-k)$, $-(k-l)$, $-l$ at $p$ and the opposite ones at $q$. The tensor $V_3$ is invariant under the cyclic permutations of its arguments together with its indices, and it changes sign under a transposition, like $\epsilon_{abc}$; this gives $24F_6$.

\emph{Step 6 ($\Gamma=5$).} The line $K_0$ carries $\mathsf p-k$ from $p$ to $q$, the line $X$ carries $k$ from $p$ to $P$, and $K_{\mathrm a}$, $K_{\mathrm b}$ carry $k-l$ and $l$ from $P$ to $q$. The factor is $-\frac12(-2i)^2\cdot6=12$, with the factor $-\frac12$ of $\mathsf T_5$, and the tensors are
\[
V_2(k-\mathsf p,-k)_{\alpha\sigma}\,V_3(\mathsf p-k,l,k-l)_{\alpha\beta\gamma}\,V_3(k,-l,l-k)_{\sigma\beta\gamma};
\]
with the factors $|\mathsf q|^{-2}$ of the four lines this is $12F_5$.

\emph{Step 7 ($\Gamma=2$).} By Lemma~\ref{IV-73-53}(a), applied to two quadratic polynomials,
\[
-\langle\partial_{b,3}\mathsf S_3;\partial_{b',3}\mathsf S_3\rangle_0=-\tfrac12\sum\mathsf a_P^{(3)}[\![\delta_{b,3},\delta_{\mathsf i},\delta_{\mathsf i'}]\!]\,\mathsf a_{P'}^{(3)}[\![\delta_{b',3},\delta_{\mathsf j},\delta_{\mathsf j'}]\!]\,\mathsf C(\mathsf i,\mathsf j)\,\mathsf C(\mathsf i',\mathsf j'),
\]
the sum extending over $P,P'$ and over the indices $\mathsf i,\mathsf i',\mathsf j,\mathsf j'$. Here $\sum_{c,c'}\epsilon_{3cc'}^2=2$, and together with the factor $-\mathsf n$ of $\mathsf T_{\mathrm I}$ the total factor is $-\mathsf n\cdot(-\frac12)(-2i)^2\cdot2=-12=24\cdot(-\frac12)$. The line $\Pi$ carries $\mathsf p-k$ from $p$ to $q$, and the chain carries $k$. The curl legs give $V_2(-k,k-\mathsf p)_{\mu\nu}$ at $p$ and $V_2(k,\mathsf p-k)_{\rho\nu}$ at $q$, where $\nu$ is the index of the line $\Pi$, and the cubic vertices give $V_3(k,-l,l-k)_{\mu\beta\gamma}\,V_3(-k,l,k-l)_{\rho\beta\gamma}$. With the factors of the five lines, whose squares give the denominator $\mathsf x_1^2\mathsf x_1'\mathsf x_2\mathsf x_3$, and with the parities of $V_2$ and $V_3$, this is $24F_2$.

\emph{Step 8 ($\Gamma=3$).} The ghost loop enters $\Sigma_{\mathrm{tot}}$ through the kernel $-8\,\mathsf D(b;b')\,\mathsf D(b';b)$. By Lemma~\ref{IV-73-51}, the leading parts of the symbols of $y\mapsto\mathsf D((x,\mu);(x+y,\rho))$ and of $y\mapsto\mathsf D((x+y,\rho);(x,\mu))$ are $-i\mathsf q_\mu|\mathsf q|^{-2}$ and $+i\mathsf q_\rho|\mathsf q|^{-2}$. With $\mathsf q=k-l$, resp.\ $\mathsf q=l$, and with the factor $-\mathsf n$ of $\mathsf T_{\mathrm I}$ we obtain
\[
-\mathsf n\,(-8)\,(k-l)_\mu\,l_\rho\,/\,\mathsf x_3\mathsf x_2=24\cdot(-1)\,l_\rho(l-k)_\mu\,/\,\mathsf x_2\mathsf x_3 ;
\]
together with the factors at $p$ and $q$ and on the lines $\Pi$, $K^\to$, $K^\leftarrow$, as in Step~7, this is $24F_3$.

\emph{Step 9 ($\Gamma=4$).} By Lemma~\ref{IV-73-57}(4), the terms $G\in\mathfrak G_4$ with $(n_X,n_Y)=(4,4)$ are exactly those which arise from the monomials of $\frac12|\mathfrak z_1|^2$ in the standard writing of $\mathsf a_P^{(4)}$ (Lemma~\ref{IV-73-53}(d)); here, by \eqref{IV-73-10} in Lemma~\ref{IV-73-9}(b), with $A_i:=A(b_i(P))$,
\[
\mathsf a_P^{(4)}=\tfrac12|\mathfrak z_1|^2+\langle\mathrm dA(P),\mathfrak z_2\rangle-\tfrac1{24}|\mathrm dA(P)|^4,\qquad\mathfrak z_1=-\sum_{i<j}s_is_j\,A_i\times A_j .
\]
Thus the vertex of the class $\Gamma=4$ is $\frac12|\mathfrak z_1|^2$. At $P=p_{\alpha\beta}$ the replacement of Step~3, $(A_1,\dots,A_4)\to(A_\alpha,A_\beta,A_\alpha,A_\beta)$, gives to leading order
\[
\mathfrak z_1\to-\bigl(A_\alpha\times A_\beta-A_\beta\times A_\alpha-A_\alpha\times A_\alpha-A_\beta\times A_\beta\bigr)=-2A_\alpha\times A_\beta ;
\]
since $\frac12|2A_\alpha\times A_\beta|^2=2|A_\alpha\times A_\beta|^2$ and the terms with $\alpha=\beta$ vanish, the sum over $\iota$, that is, over $\alpha<\beta$, gives $\sum_\iota\frac12|\mathfrak z_1|^2\to\mathsf Q(A):=\sum_{\alpha,\beta}|A_\alpha\times A_\beta|^2$, where $|Z|^2:=\langle Z,Z\rangle$ with the complex bilinear extension of $\langle\cdot,\cdot\rangle$. For $X=\xi\mathsf e_a$ and $Y=\eta\mathsf e_{a'}$ with $\xi,\eta\in\mathbb C^3$, that is, $X_\alpha=\xi_\alpha\mathsf e_a$ and $Y_\alpha=\eta_\alpha\mathsf e_{a'}$, we have, with the bilinear dot product on $\mathbb C^3$,
\begin{gather*}
\mathsf Q(sX+tY)=s^2t^2(1-\delta_{aa'})\sum_{\alpha,\beta}(\xi_\alpha\eta_\beta-\xi_\beta\eta_\alpha)^2,\\
\sum_{\alpha,\beta}(\xi_\alpha\eta_\beta-\xi_\beta\eta_\alpha)^2=2\bigl[(\xi\cdot\xi)(\eta\cdot\eta)-(\xi\cdot\eta)^2\bigr].
\end{gather*}
Polarising, with $X_i=\xi_i\mathsf e_a$ and $Y_j=\eta_j\mathsf e_{a'}$ ($i,j\in\{1,2\}$), we find
\[
\mathsf Q[\![X_1,X_2,Y_1,Y_2]\!]=4(1-\delta_{aa'})\cdot2\bigl[(\xi_1\cdot\xi_2)(\eta_1\cdot\eta_2)-\tfrac12(\xi_1\cdot\eta_1)(\xi_2\cdot\eta_2)-\tfrac12(\xi_1\cdot\eta_2)(\xi_2\cdot\eta_1)\bigr].
\]
Since $\sum_{a,a'}(1-\delta_{aa'})=6$, the factor is $-\frac14\cdot4\cdot2\cdot6=-12=24\cdot(-\frac12)$, with the factor $-\frac14$ of $\mathsf T_4$. The pairs of legs contracted with $p$ and with $q$ give $\xi_{1\mu}\xi_{2\nu}\to V_2(k,\mathsf p-k)_{\mu\nu}/\mathsf x_1\mathsf x_1'$ and $\eta_{1\mu}\eta_{2\nu}\to V_2(l,\mathsf p-l)_{\mu\nu}/\mathsf x_2\mathsf x_2'$, and the result is $24F_4$.
\end{proof}

We now show that the model $\mathsf T^{\mathrm{mod}}_\Gamma$ of Definition~\ref{IV-73-64} has, up to a part of degree $-4$, the same part of degree $-5$ as the corresponding lattice terms. For a kernel $K$ in a class $\mathfrak K(n)$ we write $K^h$, or $(K)^h$ when $K$ is given by a longer expression, for the second component of its leading data (Definition~\ref{IV-73-38}(4)), that is, for the function $K^h$ of condition (K3), and we call it the \emph{leading kernel} of $K$; it is smooth off $0$ and homogeneous of degree $-n$.

\begin{lemma}[The model has the degree $-5$ parts of the lattice classes]\label{IV-73-69}
Let $\tau_{\mathrm{gf}}=1$. For $G\in\mathfrak G_1\cup\mathfrak G_4\cup\mathfrak G_5\cup\mathfrak G_6$ let $\mathsf U_G$ be the function of degree $-5$ attached to $G$ in Lemma~\ref{IV-73-60}. For $(\mu,\nu,G)\in\mathfrak G_2\cup\mathfrak G_3$ let
\[
\mathsf U_G:=\Pi^h\cdot\bigl(K^\to_{\iota\mu}*\mathsf f_G(\nabla)[K_{\mathrm a}K_{\mathrm b}-\mathsf q_G]*K^\leftarrow_{\nu\iota'}\bigr)^h,\qquad\Pi^h:=\mathcal F^{-1}_{\ne0}\sigma_\Pi^h ,
\]
where the chain lies in $\mathfrak K(2)$ by Lemma~\ref{IV-73-62}(a) and two applications of Lemma~\ref{IV-73-46}(b). Let
\begin{gather*}
\mathsf U_\Gamma:=\sum_{\iota,\iota'}\sum_{G\in\mathfrak G_\Gamma}c_G\,\mathsf U_G\quad(\Gamma\ne5),\\
\mathsf U_5(u):=\sum_{\iota,\iota'}\sum_{G\in\mathfrak G_5}c_G\bigl[\mathsf U_G(u)+\mathsf U_G(-u)\bigr],
\end{gather*}
with the coefficients $-\mathsf n\,c_G$ of Lemma~\ref{IV-73-68} for $\Gamma=2,3$. Then for each $\Gamma=1,\dots,6$ there are a continuous function $\mathsf l_\Gamma$ on $\mathbb R^3\setminus\{0\}$, homogeneous of degree $-4$, and a pure constant $C$ such that
\[
\bigl|\mathsf T_\Gamma^{\mathrm{mod}}(u)-\mathsf l_\Gamma(u)-\mathsf U_\Gamma(u)\bigr|\le C|u|^{-6}\lg|u|\qquad(u\in\mathbb Z^3\setminus\{0\}).
\]
\end{lemma}

\begin{proof}
\emph{Step 0 (model lines).} For a line whose symbol has the leading part $\sigma^h$, of degree $s$, let $\sigma^{\mathrm{mod}}$ be the periodic continuation of $\rho\sigma^h$, with $\rho$ of Definition~\ref{IV-73-64}(3). As $\rho\in C_c^\infty(\mathbb R^3)$, $\operatorname{supp}\rho\subset\{|k|\le2\rho_0\}\subset\{|k|<\pi\}$ and $\rho=1$ near $0$, the function $\sigma^{\mathrm{mod}}$ is a line symbol with the same leading part (Definition~\ref{IV-73-38}(3)), and its parity is exact: $\sigma^{\mathrm{mod}}(-k)=(-1)^s\sigma^{\mathrm{mod}}(k)$. Its kernel $K^{\mathrm{mod}}(x)=\int_{\mathbb R^3}e^{ik\cdot x}\rho\,\sigma^h\,d^\circ k$ satisfies $K^{\mathrm{mod}}=K^h+O(|x|^{-N})$ for every $N$, with $K^h=\mathcal F^{-1}_{\ne0}\sigma^h$, by Lemma~\ref{IV-73-39}(b). For $(\mu,\nu,G)\in\mathfrak G_2\cup\mathfrak G_3$ the elementary term is, by Lemma~\ref{IV-73-57}(I), the value at $u=u'$ of the function
\[
u\mapsto K_\Pi(u)\,\bigl(K^\to_{\iota\mu}*\mathsf f_G(\nabla)[K_{\mathrm a}K_{\mathrm b}]*K^\leftarrow_{\nu\iota'}\bigr)(u),\qquad K_\Pi(y):=\Pi\bigl(p_{[\iota]}(0),p_{[\iota']}(y)\bigr)\quad(y\in\mathbb Z^3);
\]
here $K_\Pi$ is the elementary line of degree $3$ of Lemma~\ref{IV-73-51}, with the symbol $\sigma_\Pi$ of Lemma~\ref{IV-73-68}, and $\Pi(p,q)=K_\Pi(u')$ by translation invariance. For an elementary term $G$ let $G^{\mathrm{mod}}(u)$ be the same expression with every line $K$ replaced by $K^{\mathrm{mod}}$, in particular $K_\Pi$ by $K_\Pi^{\mathrm{mod}}$; for $G$ of type (I) we moreover drop the operator $\mathsf f_G(\nabla)$ and replace $K^\to$ by the elementary line $K_1'$ whose symbol is the periodic continuation of $\rho\,\sigma_\to^h\mathsf f_G^h$; it has degree $2+d_G$.

\emph{Step 1 (the model in momentum space).} We show that $\mathsf T^{\mathrm{mod}}_\Gamma=\sum c_G\,G^{\mathrm{mod}}$, the sum extending over $\iota,\iota'$ and $G$ as in Lemma~\ref{IV-73-68}, where for $\Gamma=5$ the term $G^{\mathrm{mod}}(u)$ is replaced by $G^{\mathrm{mod}}(u)+G^{\mathrm{mod}}(-u)$. By Lemma~\ref{IV-73-67},
\[
G^{\mathrm{mod}}(u)=\int_{\mathbb T}e^{i\mathsf p\cdot u}\Bigl[\iint_{\mathbb T^2}\prod\sigma^{\mathrm{mod}}\,d^\circ k\,d^\circ l\Bigr]d^\circ\mathsf p ,
\]
where the product extends over the lines of $G$ and the arguments of the symbols are those in $\mathsf F_G$. For the kites this is assertion (a) of that lemma, the model lines having the degrees of the lattice lines. For the other types we apply assertion (b) to the products and convolutions of which $G^{\mathrm{mod}}$ is built, one at a time, and we verify its hypotheses. A model line of degree $n$ is bounded by $C\langle x\rangle^{-n}$ (Lemma~\ref{IV-73-39}(a)), and its symbol is integrable. A kernel in $\ell^2(\mathbb Z^3)$ with an integrable symbol has its symbol in $L^2(\mathbb T)$, as in the proof of Lemma~\ref{IV-73-67}(b); hence the symbol of the convolution of two such kernels, a product of two functions in $L^2(\mathbb T)$, is integrable, and so is the symbol of the product of two kernels with integrable symbols, a convolution of two integrable functions.
\begin{itemize}
\item For $G\in\mathfrak G_6$, $G^{\mathrm{mod}}$ is a product of three model lines, and the rule for products is used twice.
\item For $G\in\mathfrak G_4$ with $(n_X,n_Y)=(4,4)$, $G^{\mathrm{mod}}=X*Y$, where $X$ and $Y$ are products of two model lines of degree $2$. Their symbols are given by the rule for products, and $|X(x)|,|Y(x)|\le C\langle x\rangle^{-4}$, so that $X,Y\in\ell^2(\mathbb Z^3)$ and the rule for $X*Y$ applies.
\item For $G\in\mathfrak G_5$, $G^{\mathrm{mod}}(u)=K_0(u)\,(X*Y)(u)$ with model lines $K_0$, $X$ and a product $Y$ of two model lines, where $n_X\in\{2,3\}$ and $n_Y\in\{2,3,4\}$. Thus $|X(x)|\le C\langle x\rangle^{-n_X}$ and $|Y(x)|\le C\langle x\rangle^{-n_Y}$ with exponents at least $2$, so that $X,Y\in\ell^2(\mathbb Z^3)$; the kernel $X*Y$ has the symbol $\hat X\hat Y$, and the product with $K_0$ is treated by the rule for products.
\item For $G$ of type (I), $G^{\mathrm{mod}}=K_\Pi^{\mathrm{mod}}\cdot(K_1'*X*K^{\leftarrow,\mathrm{mod}})$ with $X:=K_{\mathrm a}^{\mathrm{mod}}K_{\mathrm b}^{\mathrm{mod}}$; the degrees of $K_1'$, $X$ and $K^{\leftarrow,\mathrm{mod}}$ are $2+d_G$, $4-d_G$ and $2$, all at least $2$. Hence $K_1'*X$ has the symbol $\hat K_1'\hat X$, and Lemma~\ref{IV-73-40}(c), with $(a,a')=(2+d_G,4-d_G)$ and $a+a'-3=3$, gives
\[
|(K_1'*X)(x)|\le C\langle x\rangle^{-\min\{2+d_G,\,4-d_G,\,3\}}\lg\langle x\rangle ;
\]
the exponent is at least $2$, so that $K_1'*X\in\ell^2(\mathbb Z^3)$, the logarithm included. The rule for convolutions applies once more to $(K_1'*X)*K^{\leftarrow,\mathrm{mod}}$, and the product with $K_\Pi^{\mathrm{mod}}$ is treated by the rule for products.
\end{itemize}

For $k,l,\mathsf p\in[-\pi,\pi]^3$ the periodic continuations may be replaced by $\rho\sigma^h$. We use the following observation: if $v\in[-\pi,\pi]^3$, $|w|<4\rho_0$ and $|v-w|_{\mathbb T}<2\rho_0$, then $|v-w|<2\rho_0$, because the distance of $v-w$ from any point $2\pi\mathsf m$ with $\mathsf m\in\mathbb Z^3\setminus\{0\}$ is at least $\pi-4\rho_0>2\rho_0$. Let $k,l,\mathsf p\in[-\pi,\pi]^3$, and assume first that every line $\mathsf q$ of $\Gamma$ satisfies $|\mathsf q|_{\mathbb T}<2\rho_0$; this is the case wherever all the factors $\sigma^{\mathrm{mod}}$ are non-zero. Then $|l|=|l|_{\mathbb T}<2\rho_0$, since $l\in[-\pi,\pi]^3$. For $\Gamma\ne6$ we have likewise $|k|<2\rho_0$; then $|k-l|<4\rho_0<\pi$, so that $|k-l|=|k-l|_{\mathbb T}<2\rho_0$ wherever $k-l$ is a line ($\Gamma=1,2,3,5$), and the observation with $v=\mathsf p$ and $w=k$, resp.\ $w=l$, gives $|\mathsf p-k|<2\rho_0$, resp.\ $|\mathsf p-l|<2\rho_0$, wherever these vectors are lines.

For $\Gamma=6$ the observation with $(v,w)=(k,l)$ gives $|k-l|<2\rho_0$, hence $|k|<4\rho_0$, and with $(v,w)=(\mathsf p,k)$ it gives $|\mathsf p-k|<2\rho_0$. Thus every line satisfies $|\mathsf q|<2\rho_0$, and $\sigma^{\mathrm{mod}}(\mathsf q)=\rho(\mathsf q)\sigma^h(\mathsf q)$. If, on the other hand, some line satisfies $|\mathsf q|_{\mathbb T}\ge2\rho_0$, then $\sigma^{\mathrm{mod}}(\mathsf q)=0$ and also $\rho(\mathsf q)=0$, since $|\mathsf q|\ge|\mathsf q|_{\mathbb T}$. In either case the product of the factors $\sigma^{\mathrm{mod}}$ equals $R_\Gamma\mathsf F_G$ at $(k,l,\mathsf p)$, the lines of $G$ being the lines of $\Gamma$ (for type (I) the factor $\rho\,\sigma^h_\to\mathsf f^h_G$ of $K_1'$ provides the factors $\sigma^h_\to(k)\,\mathsf f^h_G(k)$ of $\mathsf F_G$).

Since $R_\Gamma$ vanishes unless $k,l,\mathsf p\in[-\pi,\pi]^3$ (Definition~\ref{IV-73-64}(3)), we obtain
\[
G^{\mathrm{mod}}(u)=\int_{\mathbb R^3}e^{i\mathsf p\cdot u}\Bigl[\iint R_\Gamma\,\mathsf F_G\,d^\circ k\,d^\circ l\Bigr]d^\circ\mathsf p .
\]
By Lemma~\ref{IV-73-68}, $\sum c_G\iint R_\Gamma\mathsf F_G\,d^\circ k\,d^\circ l=24\Psi_\Gamma$ for $\Gamma\ne5$, whence $\sum c_GG^{\mathrm{mod}}=\mathsf T^{\mathrm{mod}}_\Gamma$. For $\Gamma=5$ the term $G^{\mathrm{mod}}(-u)$ carries the factor $e^{-i\mathsf p\cdot u}$, and $\sum c_G\iint R_5\mathsf F_G\,d^\circ k\,d^\circ l=12\Psi_5$ is even; hence $\sum c_G[G^{\mathrm{mod}}(u)+G^{\mathrm{mod}}(-u)]=12\int(e^{i\mathsf p\cdot u}+e^{-i\mathsf p\cdot u})\Psi_5\,d^\circ\mathsf p=\mathsf T^{\mathrm{mod}}_5(u)$.

\emph{Step 2 (expansions).} Lemma~\ref{IV-73-60} applies to $G^{\mathrm{mod}}$, the model lines being elementary lines of the same degrees as the lattice lines; the functions of degree $-5$ which it provides are determined by the leading parts of the symbols of the lines alone (Lemma~\ref{IV-73-60}), and therefore coincide with the $\mathsf U_G$. For $G$ of type (I) let $X:=K^{\mathrm{mod}}_{\mathrm a}K^{\mathrm{mod}}_{\mathrm b}$, of degree $n_X=4-d_G$, and write $X=\mathsf q_X+X'$ with $X'\in\mathfrak K(n_X)$ (Lemma~\ref{IV-73-46}(c)). Since $K_1'$ is a line of degree $2+d_G\le4$, two applications of Lemma~\ref{IV-73-46}(b) give $K_1'*X'*K^{\leftarrow,\mathrm{mod}}\in\mathfrak K(2)$, and the first component of its leading data is $\sigma_\to^h\mathsf f_G^h\mathsf I^h\sigma_\leftarrow^h$, where the homogeneous part $\mathsf I^h$ of the bubble depends only on $\sigma^h_{\mathrm a}$, $\sigma^h_{\mathrm b}$ (Lemma~\ref{IV-73-42}(a)).

The lattice chain $K^\to_{\iota\mu}*\mathsf f_G(\nabla)[K_{\mathrm a}K_{\mathrm b}-\mathsf q_G]*K^\leftarrow_{\nu\iota'}$ has the same first component of its leading data, by Lemma~\ref{IV-73-62}(a) and by Lemma~\ref{IV-73-46}(b), used for each of its two convolutions. Within $\mathfrak K(2)$ the leading kernel is determined by the first component of the leading data (Lemma~\ref{IV-73-46}(d)); hence the two chains have the same leading kernel. Now $G^{\mathrm{mod}}=K^{\mathrm{mod}}_\Pi\cdot(K_1'*\mathsf q_X*K^{\leftarrow,\mathrm{mod}})+K^{\mathrm{mod}}_\Pi\cdot(K_1'*X'*K^{\leftarrow,\mathrm{mod}})$, and the first term is local. In the second term, $K_\Pi^{\mathrm{mod}}=\Pi^h+O(|u|^{-N})$ with $\Pi^h=O(|u|^{-3})$, and $K_1'*X'*K^{\leftarrow,\mathrm{mod}}$ differs from its leading kernel, which is $O(|u|^{-2})$, by $O(|u|^{-3}\lg|u|)$, by condition (K3).

We obtain in all cases
\[
G^{\mathrm{mod}}=[\text{local terms}]+\mathsf U_G+O(|u|^{-6}\lg|u|).
\]

\emph{Step 3 (local terms).} In this step all lines are model lines, and we omit the superscript $\mathrm{mod}$. Because of the exact parity of the model symbols, every one-loop function $\mathsf I$ built from model lines satisfies $\mathsf I(-\kappa)=(-1)^{\omega+1}\mathsf I(\kappa)$, as one sees by the substitution $l\to-l$; its homogeneous part $\mathsf I^h$ has the same parity (Lemma~\ref{IV-73-42}(a)). By Lemma~\ref{IV-73-42}(b) the polynomial $\mathsf Q(\kappa)-(-1)^{\omega+1}\mathsf Q(-\kappa)$, of degree at most $\omega$, is therefore $O(|\kappa|^{\omega+1}\lg|\kappa|\,\Lambda(\kappa))$, hence equal to $0$. Consequently $\mathsf Q=0$ if $\omega\le0$, and $\mathsf Q=\mathsf Q(0)$ if $\omega=1$; in both cases $\mathsf q=\mathsf Q(0)\,\delta_0$.

The local terms are thus: $\mathsf Q_{\mathsf p}(0)K_4K_5(u)+\mathsf Q_{\mathsf q}(0)K_1K_2(u)$ for the kites, where $n_4+n_5=4$ if $\omega_{\mathsf p}=1$ and $n_1+n_2=4$ if $\omega_{\mathsf q}=1$; $\mathsf Q_Y(0)K_0X(u)$, with $n_0+n_X=4$, for $\mathfrak G_5$; $\mathsf Q_X(0)Y(u)+\mathsf Q_Y(0)X(u)$, with $n_X=n_Y=4$, for $\mathfrak G_4$; and, for type (I) with $d_G=0$, $\mathsf Q_X(0)K_\Pi(u)\,(K_1'*K^\leftarrow)(u)$, where $K_1'*K^\leftarrow$ is, by Lemma~\ref{IV-73-67}(b), the elementary line of degree $1$ whose symbol is the periodic continuation of $\rho^2\sigma_\to^h\mathsf f^h_G\sigma_\leftarrow^h$ (a line symbol by Definition~\ref{IV-73-38}(3), with $\rho^2$ in place of $\rho$), $\mathsf f^h_G=\mathsf f_G(0)$ being a number.

By Step~0, and by Lemma~\ref{IV-73-39}(b) with $\rho^2$ in place of $\rho$ for the line $K_1'*K^\leftarrow$, each of these local terms is, up to $O(|u|^{-N})$, its factor $\mathsf Q_\bullet(0)$ times the product of the functions $K^h=\mathcal F^{-1}_{\ne0}\sigma^h$ of its lines, a function homogeneous of degree $-4$ and smooth off $0$. We let $\mathsf l_\Gamma$ be the sum of these functions, each taken with the coefficient with which its term $G^{\mathrm{mod}}$ enters $\mathsf T^{\mathrm{mod}}_\Gamma$ in Step~1. For $\Gamma=5$, where $\mathsf T^{\mathrm{mod}}_5(u)=\sum c_G\bigl[G^{\mathrm{mod}}(u)+G^{\mathrm{mod}}(-u)\bigr]$, the expansion of Step~2 is used at $u$ and at $-u$; accordingly, every such function $\mathsf h$ of a term $G\in\mathfrak G_5$ enters $\mathsf l_5$ together with its reflection $u\mapsto\mathsf h(-u)$, just as $\mathsf U_5$ contains $\mathsf U_G(u)$ and $\mathsf U_G(-u)$.

Adding the expansions of Step~2 with these coefficients, we obtain the claimed bound for large $|u|$; for small $|u|$ all terms are bounded. All data being fixed, the constant $C$ is pure.
\end{proof}

The functions $\mathsf l_\Gamma$ and $\mathsf U_\Gamma$ are in fact multiples of $|u|^{-4}$ and $|u|^{-5}$, since $\mathsf T^{\mathrm{mod}}_\Gamma$ is radial.

\begin{lemma}[Rotation invariance of the homogeneous parts; definition of $\kappa_\Gamma^{\mathrm{lat}}$]\label{IV-73-70}
Let $\mathsf t\colon[0,\infty)\to\mathbb C$, let $\mathsf l,\mathsf U$ be continuous functions on $\mathbb R^3\setminus\{0\}$, homogeneous of degrees $-4$ and $-5$ respectively, and assume that for some $C$
\[
\bigl|\mathsf t(|u|)-\mathsf l(u)-\mathsf U(u)\bigr|\le C|u|^{-6}\lg|u|\qquad(u\in\mathbb Z^3\setminus\{0\}).
\]
Then $\mathsf l(u)=c_4|u|^{-4}$ and $\mathsf U(u)=\kappa|u|^{-5}$ with constants $c_4,\kappa$. In particular the functions of Lemma~\ref{IV-73-69} satisfy $\mathsf l_\Gamma(u)=c_{4,\Gamma}|u|^{-4}$ and
\[
\mathsf U_\Gamma(u)=\kappa_\Gamma^{\mathrm{lat}}\,|u|^{-5}\qquad(\Gamma=1,\dots,6);
\]
this defines the numbers $c_{4,\Gamma}$ and $\kappa_\Gamma^{\mathrm{lat}}$.
\end{lemma}

\begin{proof}
Let $O$ be a rotation with rational entries, $\mathsf N$ a common denominator of these entries, and $\mathsf m\in\mathbb Z^3\setminus\{0\}$. For integers $j\ge1$ the points $j\mathsf N\mathsf m$ and $j\mathsf NO\mathsf m$ lie in $\mathbb Z^3$ and have the same norm $jr_1$, $r_1:=\mathsf N|\mathsf m|$. Let $\hat{\mathsf m}:=\mathsf m/|\mathsf m|$. Subtracting the hypothesis at these two points, at which $\mathsf t$ takes the same value, we obtain by homogeneity
\[
\Bigl|(jr_1)^{-4}\bigl[\mathsf l(\hat{\mathsf m})-\mathsf l(O\hat{\mathsf m})\bigr]+(jr_1)^{-5}\bigl[\mathsf U(\hat{\mathsf m})-\mathsf U(O\hat{\mathsf m})\bigr]\Bigr|\le2C(jr_1)^{-6}\lg(jr_1) .
\]
Multiplying by $(jr_1)^4$ and letting $j\to\infty$ we obtain $\mathsf l(\hat{\mathsf m})=\mathsf l(O\hat{\mathsf m})$; then, multiplying the remaining difference by $(jr_1)^5$, we obtain $\mathsf U(\hat{\mathsf m})=\mathsf U(O\hat{\mathsf m})$. The directions $\mathsf m/|\mathsf m|$ are dense in $S^2$, since $\mathbb Q^3$ is dense in $\mathbb R^3$. The rational rotations are dense in $SO(3)$: the Cayley map $\mathsf A\mapsto(\mathbf 1-\mathsf A)(\mathbf 1+\mathsf A)^{-1}$ is continuous from the real antisymmetric $3\times3$ matrices onto the rotations without eigenvalue $-1$, which are dense in $SO(3)$; the antisymmetric matrices with rational entries are dense among the real ones, and by the formula they are mapped to rotations with rational entries, since the inverse of an invertible matrix with rational entries has rational entries.

By continuity, $\mathsf l$ and $\mathsf U$ are invariant under all rotations; as $SO(3)$ acts transitively on $S^2$, they are constant on $S^2$, and homogeneity gives the claim. For the last assertion we note that $\mathsf T^{\mathrm{mod}}_\Gamma$ is the restriction to $\mathbb Z^3$ of a radial function, since $\Psi_\Gamma$ is radial (Definition~\ref{IV-73-64}(3)); thus $\mathsf T^{\mathrm{mod}}_\Gamma(u)=\mathsf t(|u|)$ for a function $\mathsf t$, and Lemma~\ref{IV-73-69} provides the hypothesis.
\end{proof}

We can now describe $\mathsf B_1^{\mathrm{lat}}(p,q)$ at large distances of the two plaquettes. The local parts of the elementary terms drop out by the zero-momentum identities of Lemma~\ref{IV-73-62}, and the parts of degree $-5$ add up, by Lemmas~\ref{IV-73-69} and~\ref{IV-73-70}, to a multiple of $|u|^{-5}$.

\begin{theorem}[The kernel of degree $-5$: large-distance asymptotics of $\mathsf B_1^{\mathrm{lat}}$]\label{IV-73-71}
There are functions $\mathsf U_{ij}\in C^1(\mathbb R^3\setminus\{0\})$, $i,j\in\{1,2,3\}$, homogeneous of degree $-5$, and pure constants $C,r_0$ such that for all plaquettes $p,q$ of $\mathbb Z^3$, with $u:=x(q)-x(p)$ and $r:=|u|\ge r_0$,
\begin{equation}\label{IV-73-72}
\begin{gathered}
\bigl|\mathsf B_1^{\mathrm{lat}}(p,q)-\mathsf U_{\lambda(p)\lambda(q)}(u)\bigr|\le C\,r^{-6}\lg r,\\
\sum_{i,j}\mathsf U_{ij}(u)=\kappa^{\mathrm{lat}}|u|^{-5},\qquad\kappa^{\mathrm{lat}}:=\sum_{\Gamma=1}^6\kappa_\Gamma^{\mathrm{lat}} ,
\end{gathered}
\end{equation}
with the numbers $\kappa_\Gamma^{\mathrm{lat}}$ of Lemma~\ref{IV-73-70}. Moreover $\mathsf B_1^{\mathrm{lat}}$ is bounded.
\end{theorem}

\begin{proof}
We work at $\tau_{\mathrm{gf}}=1$; recall that $\mathsf B_1^{\mathrm{lat}}(p,q)=\mathsf B_1^{\mathrm{lat}}(p,q;1)$ by \eqref{IV-73-25}. Boundedness is contained in Lemma~\ref{IV-73-55}, which gives $\sup_{p,q}|\mathsf B_1^{\mathrm{lat}}(p,q;\tau_{\mathrm{gf}})|<\infty$ for every $\tau_{\mathrm{gf}}>0$. Let $p=p_{[\iota]}(x_p)$, $q=p_{[\iota']}(x_q)$ and $u':=x_q-x_p$; then $|u-u'|\le\sqrt2$, and we take $r$ large. By Lemma~\ref{IV-73-53}(c) and Lemma~\ref{IV-73-57}, $\mathsf B_1^{\mathrm{lat}}(p,q)$ is the sum of $\mathsf T_{\mathrm I}$ and of finitely many elementary terms, to which Lemma~\ref{IV-73-60} applies.

\emph{Step 1 (local terms at $q$).} These are the first half of $\mathsf T_{4'}$ and the terms with $\mathsf q_Y$ in the first half of $\mathsf T_5$, with $\mathsf q_Y$ in $\mathsf T_4$ and with $\mathsf q_{\mathsf q}$ in $\mathsf T_1$, in the expansions of Lemma~\ref{IV-73-60}. By Lemma~\ref{IV-73-62}(b), and since $\mathsf E_p((x_q+y,\mu))=K^\to_{\iota\mu}(u'+y)$ by translation invariance, their sum equals
\[
\tfrac{\mathsf n}2\sum_{b,b'}\gamma_q(b,b')\,\mathsf E_p(b)\,\mathsf E_p(b')=\tfrac{\mathsf n}2\sum_{\mu,\nu}\sum_{y,y'}\gamma_{q;\mu\nu}(y,y')\,K^\to_{\iota\mu}(u'+y)\,K^\to_{\iota\nu}(u'+y') .
\]
The kernels $\gamma_{q;\mu\nu}$ have finite support, and their symbols vanish at $\kappa=0$ together with their gradients (Lemma~\ref{IV-73-62}(b)); the $K^\to_{\iota\mu}$ are elementary lines of degree $2$. Lemma~\ref{IV-73-61}(a), applied with $\gamma_{\mu\nu}:=\gamma_{q;\mu\nu}$ and $K_\mu:=K^\to_{\iota\mu}$, shows that the sum is $O(|u'|^{-6})=O(r^{-6})$. By the translation invariance in Lemma~\ref{IV-73-57}, the kernels $\gamma_{q;\mu\nu}$ depend on $q$ only through $\iota'$, so that only finitely many constants occur.

\emph{Step 2 (local terms at $p$).} These are the second half of $\mathsf T_{4'}$ and the terms with $\mathsf q_Y$ in the second half of $\mathsf T_5$, with $\mathsf q_X$ in $\mathsf T_4$ and with $\mathsf q_{\mathsf p}$ in $\mathsf T_1$. We show that they are the local terms at the second plaquette of the ordered pair $(q,p)$; by Step~1, applied to this pair, their sum is then $\frac{\mathsf n}2\sum_{b,b'}\gamma_p(b,b')\,\mathsf E_q(b)\,\mathsf E_q(b')=O(r^{-6})$. For the second halves of $\mathsf T_{4'}$ and $\mathsf T_5$ this holds because they are the first halves for $(q,p)$ (Lemma~\ref{IV-73-57}). The formulas for $\mathsf T_1$ and $\mathsf T_4$ in Lemma~\ref{IV-73-53}(c) are symmetric under $p\leftrightarrow q$.

An elementary term $G$ of $\mathsf T_1$ or $\mathsf T_4$ is given by the monomials at $P,P'$ and by the choice of the legs contracted with $p$, with $q$, and with one another; taken for the ordered pair $(q,p)$, it becomes a term $\tilde G$ with $c_{\tilde G}=c_G$, and $G\mapsto\tilde G$ is a bijection $\mathfrak G_\bullet(\iota,\iota')\to\mathfrak G_\bullet(\iota',\iota)$. In $\tilde G$ the base points are counted from $x_q$, $u'$ is replaced by $-u'$, and the contractions, which are symmetric in their two legs, are taken from the other end: for the kite
\[
\tilde K_1=K_4(-\cdot),\quad\tilde K_2=K_5(-\cdot),\quad\tilde K_3=K_3,\quad\tilde K_4=K_1(-\cdot),\quad\tilde K_5=K_2(-\cdot),
\]
and for $\mathsf T_4$ we have $\tilde X=Y(-\cdot)$, $\tilde Y=X(-\cdot)$. Since $K(-\cdot)$ has the symbol $\sigma(-k)$, the function $\hat{\mathsf T}_{\mathsf q}$ of $\tilde G$ is
\[
\int_{\mathbb T}\sigma_1(-\kappa_1+l)\,\sigma_2(-\kappa_2-l)\,\sigma_3(l)\,d^\circ l=\hat{\mathsf T}_{\mathsf p}(-\kappa),
\]
and the symbol of $\tilde Y$ is $\mathsf I_X(-k)$, after the substitution $l\to-l$. By the last assertion of Lemma~\ref{IV-73-42}(d), whose hypothesis holds here for all $\kappa$, the corresponding polynomials are $\mathsf Q_{\mathsf p}(-\kappa)$ and $\mathsf Q_X(-k)$. The passage from a polynomial $\mathsf Q$ to $\tilde{\mathsf Q}(\kappa)=\mathsf Q(0)+\sum_j\partial_j\mathsf Q(0)\sin\kappa_j$ in Proposition~\ref{IV-73-47} and Lemma~\ref{IV-73-46}(c) commutes with the change of sign of the argument, the sine being odd. Hence the kernel $\mathsf q_{\mathsf q}$ of $\tilde G$ is $(z,w)\mapsto\mathsf q_{\mathsf p}(-z,-w)$, and $\mathsf q_{\tilde Y}=\mathsf q_X(-\cdot)$, so that
\begin{gather*}
\sum_{z,w}\mathsf q_{\mathsf p}(-z,-w)\,\tilde K_1(-u'-z)\,\tilde K_2(-u'-w)=\sum_{z,w}\mathsf q_{\mathsf p}(z,w)\,K_4(u'-z)\,K_5(u'-w),\\
(\mathsf q_{\tilde Y}*\tilde X)(-u')=(\mathsf q_X*Y)(u') .
\end{gather*}
These are the local terms at $p$ in $\mathsf T_1$ and in $\mathsf T_4$.

\emph{Step 3 (the class $\mathsf T_{\mathrm I}$).} By Lemma~\ref{IV-73-62}(a),
\[
K^\to*\Sigma_{\mu\nu}*K^\leftarrow=K^\to*\mathsf q^\Sigma_{\mu\nu}*K^\leftarrow+K^\to*\mathsf W_{\mu\nu}*K^\leftarrow .
\]
The first term is $O(r^{-3})$ by Lemma~\ref{IV-73-61}(b), since $\hat{\mathsf q}^\Sigma_{\mu\nu}(k)=O(|k|^2)$. The second lies in $\mathfrak K(2)$, by Lemma~\ref{IV-73-46}(b) applied to each of its two convolutions, and it equals its leading kernel, of degree $-2$, up to $O(r^{-3}\lg r)$; by the linearity of the leading data, this leading kernel is the sum over $G$, with the coefficients $c_G$, of the leading kernels of the chains in the definition of $\mathsf U_G$ in Lemma~\ref{IV-73-69}. Moreover $\Pi(p,q)=\Pi^h(u')+O(r^{-4})$, with $\Pi^h$ homogeneous of degree $-3$ (Lemma~\ref{IV-73-51}, Lemma~\ref{IV-73-39}(c)). Therefore
\[
\mathsf T_{\mathrm I}=\sum_{(\mu,\nu,G)\in\mathfrak G_2\cup\mathfrak G_3}(-\mathsf n\,c_G)\,\mathsf U_G(u')+O(r^{-6}\lg r).
\]

\emph{Step 4 (assembly).} Let
\[
\begin{aligned}
\mathsf U_{\iota\iota'}(u):={}&\sum_{\Gamma\ne5}\ \sum_{G\in\mathfrak G_\Gamma(\iota,\iota')}c_G\,\mathsf U_G(u)\\
&+\sum_{G\in\mathfrak G_5(\iota,\iota')}c_G\,\mathsf U_G(u)+\sum_{G\in\mathfrak G_5(\iota',\iota)}c_G\,\mathsf U_G(-u),
\end{aligned}
\]
with the coefficients $-\mathsf n\,c_G$ for $\Gamma=2,3$; the last sum comes from the second half of $\mathsf T_5$, which is the first half for $(q,p)$ with $u'$ replaced by $-u'$. The function $\mathsf U_{\iota\iota'}$ is homogeneous of degree $-5$ and of class $C^1$ off $0$, by Lemma~\ref{IV-73-60} and since the leading kernels of the classes $\mathfrak K(n)$ are smooth off $0$. By Steps 1--3 and Lemma~\ref{IV-73-60}, $\mathsf B_1^{\mathrm{lat}}(p,q)=\mathsf U_{\iota\iota'}(u')+O(r^{-6}\lg r)$. The gradient of $\mathsf U_{\iota\iota'}$ is homogeneous of degree $-6$, and all points $u'+t(u-u')$, $0\le t\le1$, have modulus at least $r/2$; hence $|\mathsf U_{\iota\iota'}(u')-\mathsf U_{\iota\iota'}(u)|\le Cr^{-6}$.

With $\mathsf U_{ij}:=\mathsf U_{\iota\iota'}$ for $(i,j)=(\iota,\iota')=(\lambda(p),\lambda(q))$ we obtain the bound in \eqref{IV-73-72}; as $\tau_{\mathrm{gf}}=1$, the finite sets $\mathfrak G_\bullet(\iota,\iota')$, the coefficients $c_G$ and the lines are fixed once and for all, the constants $C,r_0$ are pure. Finally, exchanging $\iota$ and $\iota'$ in the last sum and using Lemmas~\ref{IV-73-69} and~\ref{IV-73-70},
\[
\sum_{\iota,\iota'}\mathsf U_{\iota\iota'}=\sum_{\Gamma=1}^6\mathsf U_\Gamma=\sum_{\Gamma=1}^6\kappa_\Gamma^{\mathrm{lat}}\,|u|^{-5} .
\]
\end{proof}

For the smeared observables of \eqref{IV-73-25} the theorem gives the following rate of convergence on test pairs. Here $\mathfrak T_{\mathbb R^3}(\mathsf d;\vartheta)$ is the set of test pairs at the scale $\mathsf d$ with the shape parameter $\vartheta$ defined in Section~\ref{C-350-1}, with $\mathbb E=\mathbb R^3$.

\begin{corollary}[Rate of convergence of $\mathsf P_1^{\mathrm{lat}}$ to $\kappa^{\mathrm{lat}}\mathcal K_5$]\label{IV-73-73}
There is a pure constant $C$ such that the following holds for all $\mathsf d>0$, $\vartheta\in(0,1)$, $(f,h)\in\mathfrak T_{\mathbb R^3}(\mathsf d;\vartheta)$ and $0<a\le\mathsf d$. For the kernel of Section~\ref{C-350-1}, at $q=5$,
\[
\mathcal K_5(f,h)=\iint f(x)\,h(y)\,|x-y|^{-5}\,dx\,dy ,
\]
the integrand being $0$ where $f(x)h(y)=0$. Then
\begin{equation}\label{IV-73-74}
\bigl|\mathsf P_1^{\mathrm{lat}}(f,h;a)-\kappa^{\mathrm{lat}}\,\mathcal K_5(f,h)\bigr|\le C\,\vartheta^{-1}\,a\,\lg(\mathsf d/a)\,\|f\|_\infty\|h\|_\infty .
\end{equation}
\end{corollary}

We recall three properties of a test pair $(f,h)\in\mathfrak T_{\mathbb R^3}(\mathsf d;\vartheta)$: the supports of $f$ and $h$ lie in closed balls $\mathbb B(x_f,\mathsf d)$ and $\mathbb B(x_h,\mathsf d)$ of radius $\mathsf d$; $\mathsf d\le|v-w|\le10\,\mathsf d$ for $v\in\operatorname{supp}f$ and $w\in\operatorname{supp}h$ (the support separation in Section~\ref{C-350-1}); and $\operatorname{Lip}f\le\|f\|_\infty/\vartheta\mathsf d$, $\operatorname{Lip}h\le\|h\|_\infty/\vartheta\mathsf d$ (by the definition of that test class).

\begin{proof}
Let $\varrho:=\mathsf d/a\ge1$, $\mathsf A:=\|f\|_\infty\|h\|_\infty$ and $\varrho_0:=\max\{r_0,8\}$, with $r_0$ of Theorem~\ref{IV-73-71}. A ball of radius $\varrho+2$ contains at most $C\varrho^3$ points of $\frac12\mathbb Z^3$. The barycentres $x(p)$ lie in $\frac12\mathbb Z^3$, distinct plaquettes have distinct barycentres, and $f(ax(p))\ne0$ implies $x(p)\in\mathbb B(x_f/a,\varrho)$; likewise for $h$.

\emph{Step 1 ($\varrho\le\varrho_0$).} Then the sum defining $\mathsf P_1^{\mathrm{lat}}(f,h;a)$ has at most $C$ non-zero terms, and since $\mathsf B_1^{\mathrm{lat}}$ is bounded (Theorem~\ref{IV-73-71}), $|\mathsf P_1^{\mathrm{lat}}(f,h;a)|\le Ca\mathsf A$. Moreover $|x-y|\ge\mathsf d$ where $f(x)h(y)\ne0$, whence $|\mathcal K_5(f,h)|\le(\frac{4\pi}3\mathsf d^3)^2\mathsf d^{-5}\mathsf A\le C\varrho_0a\mathsf A$. As $\vartheta^{-1}>1$ and $\lg(\mathsf d/a)\ge1$, \eqref{IV-73-74} holds in this case.

\emph{Step 2 ($\varrho>\varrho_0$, first error).} If $f(ax(p))h(ax(q))\ne0$, then $r=|x(q)-x(p)|\in[\varrho,10\varrho]$ by the separation of the supports; in particular $r\ge r_0$. There are at most $C\varrho^6$ such pairs, and $r^{-6}\lg r\le\varrho^{-6}\lg(10\varrho)$ for these $r$. By \eqref{IV-73-72}, the replacement of $\mathsf B_1^{\mathrm{lat}}(p,q)$ by $\mathsf U_{\lambda(p)\lambda(q)}(x(q)-x(p))$ in $\mathsf P_1^{\mathrm{lat}}(f,h;a)$ costs at most
\[
a\cdot C\varrho^6\cdot C\varrho^{-6}\lg(10\varrho)\,\mathsf A\le Ca\lg\varrho\,\mathsf A .
\]

\emph{Step 3 (second error).} We write $p=p_{[i]}(x)$, $q=p_{[j]}(y)$, so that $x(p)=x+\mathsf o_i$ with $\mathsf o_i:=x(p_{[i]}(0))$, $|\mathsf o_i|=2^{-1/2}$; all terms below are set equal to $0$ where the product of the values of $f$ and $h$ vanishes. After Step~2 we have to compare
\[
\mathfrak R_1:=a\sum_{i,j}\sum_{x,y\in\mathbb Z^3}f(a(x+\mathsf o_i))\,h(a(y+\mathsf o_j))\,\mathsf U_{ij}(y-x+\mathsf o_j-\mathsf o_i)
\]
with
\[
\mathfrak R_2:=a\sum_{x,y}f(ax)\,h(ay)\sum_{i,j}\mathsf U_{ij}(y-x)=a\,\kappa^{\mathrm{lat}}\sum_{x,y}\frac{f(ax)\,h(ay)}{|y-x|^5},
\]
the last equality by \eqref{IV-73-72}. We replace $f$, $h$ and $\mathsf U_{ij}$ one after the other. In every non-zero term, the arguments of $\mathsf U_{ij}$ which occur, and the points between them, have modulus at least $\varrho-2\ge\varrho/2$, and $x,y$ lie in balls of radius $\varrho+1$. By homogeneity, $|\mathsf U_{ij}|\le C\varrho^{-5}$ and $|\nabla\mathsf U_{ij}|\le C\varrho^{-6}$ at these points, and the Lipschitz bound gives $|f(a(x+\mathsf o_i))-f(ax)|\le a\|f\|_\infty/\vartheta\mathsf d$, and likewise for $h$. Therefore
\[
|\mathfrak R_1-\mathfrak R_2|\le a\,C\varrho^6\Bigl[\frac a{\vartheta\mathsf d}\,\varrho^{-5}+\varrho^{-6}\Bigr]\mathsf A\le Ca\vartheta^{-1}\mathsf A .
\]

\emph{Step 4 (third error).} Let $\mathsf F(v,w):=f(v)h(w)|v-w|^{-5}$, equal to $0$ where $f(v)h(w)=0$. Then $\mathfrak R_2=\kappa^{\mathrm{lat}}\sum_{x,y}a^6\mathsf F(ax,ay)$ and $\mathcal K_5(f,h)=\iint\mathsf F$. Let $(v,w)\in(ax,ay)+[0,a)^6$. Unless $\mathsf F$ vanishes at both points $(v,w)$ and $(ax,ay)$, we have $|v-w|\ge\mathsf d-2\sqrt3a\ge\mathsf d/2$ at both points and at all points $(1-t)(ax,ay)+t(v,w)$, $0\le t\le1$ (recall $\varrho>8$), and the Lipschitz bounds for $f,h$ together with the bound $5|z|^{-6}$ for the gradient of $|z|^{-5}$ give
\[
|\mathsf F(v,w)-\mathsf F(ax,ay)|\le\sqrt3a\Bigl[\frac2{\vartheta\mathsf d}\Bigl(\frac{\mathsf d}2\Bigr)^{-5}+10\Bigl(\frac{\mathsf d}2\Bigr)^{-6}\Bigr]\mathsf A .
\]
At most $C\varrho^6$ of the cubes $(ax,ay)+[0,a)^6$, of volume $a^6$, contain points where $\mathsf F\ne0$. Hence
\[
|\mathfrak R_2-\kappa^{\mathrm{lat}}\mathcal K_5(f,h)|\le C\varrho^6a^6\cdot a\vartheta^{-1}\mathsf d^{-6}\mathsf A=Ca\vartheta^{-1}\mathsf A .
\]
Adding the bounds of Steps 2--4 and using $\vartheta^{-1}>1$, $\lg\varrho\ge1$, we obtain \eqref{IV-73-74}.
\end{proof}

\subsection{The boxes exhaust the lattice}

We now show that the coefficients $\mathsf B_j^{[n]}$ of the boxes converge to the coefficients $\mathsf B_j^{\mathrm{lat}}$ of the lattice $\mathbb Z^3$; together with Proposition~\ref{IV-73-27}, this shows that $\mathsf B_1^{\mathrm{lat}}$ is independent of the gauge parameter. The kernels of a box are expressed by the method of images through kernels on a discrete torus, which we first bound uniformly in the size of the torus.

\begin{lemma}[Kernels on discrete tori, uniformly]\label{IV-73-75}
For an integer $n\ge2$ let $\mathbb T_n:=\frac\pi n\mathbb Z^3\cap(-\pi,\pi]^3$; for $z\in\mathbb Z^3$ let $|z|_n$ be the least norm of the points of $z+2n\mathbb Z^3$, and $\langle z\rangle_n:=\max\{1,|z|_n\}$. For real $s>-3$ and $\sigma\in\mathcal S^s$ let
\[
K^{(n)}(z):=(2n)^{-3}\sum_{k\in\mathbb T_n\setminus\{0\}}e^{ik\cdot z}\,\sigma(k)\qquad(z\in\mathbb Z^3).
\]
\textup{(a)} $|K^{(n)}(z)|\le C\langle z\rangle_n^{-3-s}$ for all $n\ge2$ and $z\in\mathbb Z^3$, where $C$ depends on $s$ and on finitely many of the suprema in Definition~\ref{IV-73-38}(2), but not on $n$ and $z$.

\textup{(b)} For every $z\in\mathbb Z^3$, $K^{(n)}(z)\to K(z)=\int_{\mathbb T}e^{ik\cdot z}\sigma(k)\,d^\circ k$ as $n\to\infty$.
\end{lemma}

\begin{proof}
Let $\mathsf h:=\pi/n$, so that $(2n)^{-3}=(2\pi)^{-3}\mathsf h^3$.

\emph{Step 0.} If $0\ne k\in\mathsf h\mathbb Z^3$ and $k'$ lies in the cube of side $\mathsf h$ centred at $k$, then $(1-\frac{\sqrt3}2)|k|\le|k'|\le(1+\frac{\sqrt3}2)|k|$, as $|k'-k|\le\frac{\sqrt3}2\mathsf h\le\frac{\sqrt3}2|k|$. Comparing every term with the integral over its cube, we obtain
\begin{gather*}
(2n)^{-3}\sum_{k\in\mathsf h\mathbb Z^3,\ 0<|k|\le\varrho}|k|^t\le C\varrho^{3+t}\quad(t>-3),\\
(2n)^{-3}\sum_{k\in\mathsf h\mathbb Z^3,\ |k|\ge\varrho}|k|^t\le C\varrho^{3+t}\quad(t<-3),
\end{gather*}
with $C$ depending on $t$ only.

\emph{Step 1 (proof of (b)).} Let $0<\varepsilon<1$, and split $\sigma=\chi(k/\varepsilon)\sigma+(1-\chi(k/\varepsilon))\sigma$ on $[-\pi,\pi]^3$. The first summand is bounded by $C|k|^s$ and vanishes for $|k|\ge2\varepsilon$; by Step~0 its contributions to $K^{(n)}(z)$ and to $K(z)$ are $O(\varepsilon^{3+s})$, uniformly in $n$. The second summand, set equal to $0$ at $k=0$, is continuous on $[-\pi,\pi]^3$; its contribution to $K^{(n)}(z)$ is a Riemann sum over the $(2n)^3$ subcubes of side $\mathsf h$, which converges to its contribution to $K(z)$. Letting first $n\to\infty$ and then $\varepsilon\to0$ we obtain (b).

\emph{Step 2 (proof of (a)).} Since $K^{(n)}$ is $2n$-periodic in each coordinate, we may take the representative of $z$ with $|z|=|z|_n$; then $|z_j|\le n$ for all $j$. For $z=0$ the bound follows from Step~0. Let $|z|\ge1$, and write $\sigma=\sigma'+\sigma''$ with $\sigma':=\chi(|z|k)\sigma$ on $[-\pi,\pi]^3$. As $\sigma'$ is bounded by $C|k|^s$ and vanishes for $|k|\ge2/|z|$, Step~0 shows that the contribution of $\sigma'$ is at most $C|z|^{-3-s}$. The function $\sigma''$ is smooth and periodic, it vanishes where $|k|_{\mathbb T}\le1/|z|$, and $|\partial^\alpha\sigma''(k)|\le C|k|_{\mathbb T}^{s-|\alpha|}$, because every derivative falling on $\chi(|z|k)$ produces a factor $|z|\le2/|k|$ on its support. Since $\mathbb T_n$ represents $\mathsf h\mathbb Z^3/2\pi\mathbb Z^3$, we have $\sum_ke^{ik\cdot z}\sigma''(k-\mathsf he_j)=e^{i\mathsf hz_j}\sum_ke^{ik\cdot z}\sigma''(k)$; with $(\delta_j\mathsf F)(k):=\mathsf F(k)-\mathsf F(k-\mathsf he_j)$ it follows that
\[
(1-e^{i\mathsf hz_j})^N\sum_{k\in\mathbb T_n}e^{ik\cdot z}\,\sigma''(k)=\sum_{k\in\mathbb T_n}e^{ik\cdot z}\,\delta_j^N\sigma''(k)
\]
for every integer $N\ge1$; the term $k=0$ vanishes, as $\sigma''(0)=0$. Here $|1-e^{i\mathsf hz_j}|=2|\sin\frac{\pi z_j}{2n}|\ge2|z_j|/n$, since $\sin t\ge\frac2\pi t$ on $[0,\frac\pi2]$ and $|z_j|\le n$; moreover $|\delta_j^N\sigma''(k)|\le\mathsf h^N\sup_{0\le t\le N\mathsf h}|\partial_j^N\sigma''(k-te_j)|$. We fix an integer $N>s+3$; as $|z|\le\sqrt3n$, we have $N\mathsf h\le C/|z|$. For $|k|\ge2N\mathsf h$ the points $k-te_j$, $0\le t\le N\mathsf h$, satisfy $|k-te_j|_{\mathbb T}\in[\frac12|k|,\frac32|k|]$; they may leave $[-\pi,\pi]^3$, but their distance from $2\pi\mathbb Z^3$ equals $|k|$ at $t=0$ and changes by at most $t\le\frac12|k|$. Hence for these $k$ the supremum is at most $C|k|^{s-N}$, and it vanishes unless $|k|\ge\frac2{3|z|}$. For $|k|<2N\mathsf h$ the supremum is at most $C|z|^{N-s}$, as $\sigma''$ vanishes where $|k|_{\mathbb T}\le1/|z|$ and $s-N<0$; there are at most $C$ such points $k$, and $(2n)^{-3}\le C|z|^{-3}$. Together with Step~0 we find
\[
(2n)^{-3}\sum_{k\in\mathbb T_n}\bigl|\delta_j^N\sigma''(k)\bigr|\le C\mathsf h^N|z|^{N-s-3} .
\]
Choosing $j$ with $|z_j|\ge|z|/\sqrt3$, we conclude that the contribution of $\sigma''$ is at most 
\[
(n/2|z_j|)^N(\pi/n)^N\,C|z|^{N-s-3}\le C|z|^{-3-s}.
\]
This proves (a).
\end{proof}

\begin{proposition}[The boxes exhaust the lattice]\label{IV-73-76}
Let $\tau_{\mathrm{gf}}>0$, let $p,q$ be plaquettes of $\mathbb Z^3$ and $j\in\{0,1\}$, and write $\mathsf B_0^\Omega(p,q;\tau_{\mathrm{gf}}):=\mathsf B_0^\Omega(p,q)$. The numbers $\mathsf B_j^{[n]}(p+o_n,q+o_n;\tau_{\mathrm{gf}})$ are defined for all $n\ge2$ with $p+o_n,q+o_n\in\mathsf{Pl}_n$, that is, for all but finitely many $n$, and
\begin{equation}\label{IV-73-77}
\lim_{n\to\infty}\mathsf B_j^{[n]}(p+o_n,q+o_n;\tau_{\mathrm{gf}})=\mathsf B_j^{\mathrm{lat}}(p,q;\tau_{\mathrm{gf}}).
\end{equation}
\end{proposition}

\begin{proof}
We use the method of images. Fix $n$ and write $\tau:=\tau_{\mathrm{gf}}$.

\emph{Step 1 (reflections).} We regard forms and functions on the torus $(\mathbb Z/2n)^3$ as $2ne_\nu$-periodic forms and functions on $\mathbb Z^3$, and we write $(\cdot,\cdot)_{\mathsf t}$ for the sums over one period. Let $\varrho_\nu$ be the reflection $z_\nu\to-z_\nu$, acting on bonds by $\varrho_\nu\langle z,z'\rangle=\langle\varrho_\nu z,\varrho_\nu z'\rangle$, and put $(\varrho_\nu^*\tilde A)(b):=\tilde A(\varrho_\nu b)$ and $\varrho_\nu^*\lambda:=\lambda\circ\varrho_\nu$; these maps commute with $\mathrm d$ and $\mathrm d^*$. Let $\mathsf F^{\mathrm{as}}$ and $\mathsf X^{\mathrm{as}}$ be the spaces of real periodic forms $\tilde A$ and functions $\lambda$ with $\varrho_\nu^*\tilde A=-\tilde A$ and $\varrho_\nu^*\lambda=-\lambda$ for all $\nu$; in components, $\tilde A_\mu(\varrho_\nu z)=-\tilde A_\mu(z)$ for $\nu\ne\mu$, and $\tilde A_\mu(\varrho_\mu z-e_\mu)=\tilde A_\mu(z)$.

We claim that restriction to the box is a bijection of $\mathsf F^{\mathrm{as}}$ onto $\mathsf{Form}_n^{\mathbb R}$ and of $\mathsf X^{\mathrm{as}}$ onto $\mathsf X_n^{\mathbb R}$. If $z_\nu\in\{0,n\}$, then $\varrho_\nu z\equiv z$ modulo $2n$, hence $\tilde A_\mu(z)=0$ for $\nu\ne\mu$, and $\lambda(z)=0$; by Lemma~\ref{IV-73-12}(a), $\tilde A$ vanishes on the non-free bonds of the box, and $\lambda$ vanishes on $\partial\Lambda_n$. For $\varepsilon\in\{0,1\}^3$ let $R_\varepsilon:=\prod_\nu\varrho_\nu^{\varepsilon_\nu}$, $|\varepsilon|:=\sum_\nu\varepsilon_\nu$, and $R_\varepsilon(z,\mu):=(R_\varepsilon z-\varepsilon_\mu e_\mu,\mu)$, the positively oriented one among the image bond and its reverse.

Every positively oriented bond of the torus has the form $R_\varepsilon(x,\mu)$ with $x_\mu\in\{0,\dots,n-1\}$ and $x_\nu\in\{0,\dots,n\}$ for $\nu\ne\mu$; here $(x,\mu)$ is unique, and $\varepsilon$ is unique if $(x,\mu)$ is free. Hence every $A\in\mathsf{Form}_n^{\mathbb R}$ has the unique continuation $\tilde A(R_\varepsilon b)=(-1)^{|\varepsilon|+\varepsilon_\mu}A(b)$, $b=(x,\mu)$, and every $\lambda\in\mathsf X_n^{\mathbb R}$ the unique continuation $\tilde\lambda(R_\varepsilon x)=(-1)^{|\varepsilon|}\lambda(x)$.

\emph{Step 2 (operators).} The eight boxes $R_\varepsilon[0,n]^3$ cover the torus, and a plaquette lying in two of them lies in a hyperplane $\{z_\nu=0\}$ or $\{z_\nu=n\}$, on whose bonds $\tilde A=0$. The functions $|\mathrm d\tilde A|$ and $|\mathrm d^*\tilde A|$ are invariant under the $\varrho_\nu$, and $\mathrm d^*\tilde A\in\mathsf X^{\mathrm{as}}$ vanishes on these hyperplanes. Hence
\[
(\tilde A,\tilde A')_{\mathsf t}=8(A,A'),\qquad\|\mathrm d\tilde A\|^2_{\mathsf t}=8\|\mathrm dA\|^2,\qquad\|\mathrm d^*\tilde A\|^2_{\mathsf t}=8\|\mathrm d^*A\|^2 .
\]
The operator $\mathsf L_\tau:=\mathrm d^*\mathrm d+\tau^{-1}\mathrm d\mathrm d^*$ on the torus commutes with the $\varrho_\nu^*$ and preserves $\mathsf F^{\mathrm{as}}$. By polarisation, $(\tilde A',\mathsf L_\tau\tilde A)_{\mathsf t}=8(A',\mathsf L_\tau^{[n]}A)$ for all $A,A'\in\mathsf{Form}_n^{\mathbb R}$, where $\mathsf L^{[n]}_\tau$ is the operator on $\mathsf{Form}_n^{\mathbb R}$ of the quadratic form in Definition~\ref{IV-73-15}(1); hence $\mathsf L_\tau\tilde A$ is the continuation of $\mathsf L_\tau^{[n]}A$. The forms in $\mathsf F^{\mathrm{as}}$ have mean $0$, and on periodic forms of mean $0$ the inverse of $\mathsf L_\tau$ is, by the finite Fourier inversion formula, the convolution with
\[
\mathsf C^{(n)}_{\mu\nu}(z):=(2n)^{-3}\sum_{k\in\mathbb T_n\setminus\{0\}}e^{ik\cdot z}\,\hat{\mathsf C}_{\mu\nu}(k) .
\]
The continuation of $\delta_{b'}$, for a free bond $b'=(x',\mu')$, is $\sum_\varepsilon(-1)^{|\varepsilon|+\varepsilon_{\mu'}}\delta_{R_\varepsilon b'}$. For functions we have in the same way $(\tilde\lambda,\tilde\lambda')_{\mathsf t}=8(\lambda,\lambda')$ and $\|\mathrm d\tilde\lambda\|^2_{\mathsf t}=8\|\mathrm d\lambda\|^2$, since $\mathrm d\tilde\lambda\in\mathsf F^{\mathrm{as}}$ is the continuation of $\mathrm d\lambda$; the elements of $\mathsf X^{\mathrm{as}}$ have mean $0$, and the continuation of $\delta_{x'}$, $x'\in\Lambda_n^\circ$, is $\sum_\varepsilon(-1)^{|\varepsilon|}\delta_{R_\varepsilon x'}$. By polarisation, as for forms, $\mathrm d^*\mathrm d\tilde\lambda$ is the continuation of $\mathrm d^*\mathrm d\lambda$, and on periodic functions of mean $0$ the inverse of $\mathrm d^*\mathrm d$ is the convolution with $\mathsf G^{(n)}(z):=(2n)^{-3}\sum_{k\in\mathbb T_n\setminus\{0\}}e^{ik\cdot z}|\mathsf k(k)|^{-2}$.

We apply this to $B:=(\mathsf L_\tau^{[n]})^{-1}\delta_{b'}$, whose value at a free bond $b$ is $\mathsf C^{[n]}_\tau(b,b')$; its continuation $\tilde B$ has mean $0$ and satisfies $\mathsf L_\tau\tilde B=\tilde\delta_{b'}$, the continuation of $\delta_{b'}$, so that $\tilde B$ is the convolution of $\mathsf C^{(n)}$ with $\tilde\delta_{b'}$. In the same way we treat $(\mathrm d^*\mathrm d)^{-1}\delta_{x'}$ on $\mathsf X_n^{\mathbb R}$, whose value at $x\in\Lambda_n^\circ$ is $\mathsf G^{[n]}(x,x')$. Consequently, for free bonds $(x,\mu),(x',\mu')$ and for $x,x'\in\Lambda_n^\circ$,
\begin{equation}\label{IV-73-78}
\begin{gathered}
\mathsf C^{[n]}_\tau\bigl((x,\mu),(x',\mu')\bigr)=\sum_\varepsilon(-1)^{|\varepsilon|+\varepsilon_{\mu'}}\,\mathsf C^{(n)}_{\mu\mu'}\bigl(x-R_\varepsilon x'+\varepsilon_{\mu'}e_{\mu'}\bigr),\\
\mathsf G^{[n]}(x,x')=\sum_\varepsilon(-1)^{|\varepsilon|}\,\mathsf G^{(n)}(x-R_\varepsilon x').
\end{gathered}
\end{equation}
In the first argument the right sides are values of elements of $\mathsf F^{\mathrm{as}}$, resp.\ $\mathsf X^{\mathrm{as}}$, and therefore vanish at non-free bonds, resp.\ on $\partial\Lambda_n$; in the second argument the same follows by symmetry. Since the box kernels vanish there by definition, \eqref{IV-73-78} holds for all bonds and sites of $\Lambda_n$.

\emph{Step 3 (uniform bounds).} The kernels $\mathsf E^{[n]},\Pi^{[n]},\mathsf D^{[n]},\mathsf g^{[n]}$ arise from \eqref{IV-73-78} by the differences and means of Definition~\ref{IV-73-15}(2). A difference in the first argument acts on $\mathsf C^{(n)}$ and $\mathsf G^{(n)}$. In the second argument, $R_\varepsilon$ maps the four bonds of $q$ onto those of $R_\varepsilon q$, and $\sum_js_j(-1)^{\varepsilon_{\mu_j}}\delta_{R_\varepsilon b_j(q)}=\pm\sum_js_j\delta_{b_j(R_\varepsilon q)}$, with $\mu_j$ the direction of $b_j(q)$, since $\mathrm d$ commutes with the $\varrho_\nu^*$; the same holds for the end points of a bond. Hence each of these kernels is a sum of eight terms $\pm K^{(n)}(x-R_\varepsilon x'+y_0)$ with $|y_0|\le3$, where $x$ and $x'$ are the base points of its two arguments, as in \eqref{IV-73-56} (the base point $b_-$ of a bond $b$, the base point $x_p$ of a plaquette $p$; for $\mathsf g^{[n]}(b)$ both are $b_-$), and where $K^{(n)}$ is built as in Lemma~\ref{IV-73-75} with the symbol of the corresponding line of Lemma~\ref{IV-73-51}, an element of $\mathcal S^{n_\ell-3}$, $n_\ell$ being the degree of the line.

For $x,x'\in\Lambda_n$ we have $|x-R_\varepsilon x'|_n\ge|x-x'|$: if $\varepsilon_\nu=0$, then $|x_\nu-x'_\nu|\le n$, and if $\varepsilon_\nu=1$, then $\min\{x_\nu+x'_\nu,2n-x_\nu-x'_\nu\}\ge|x_\nu-x'_\nu|$. By Lemma~\ref{IV-73-75}(a), the kernels of all boxes therefore satisfy the bounds \eqref{IV-73-56} with one constant $\mathsf c_0=\mathsf c_0(\tau_{\mathrm{gf}})$.

\emph{Step 4 (pointwise limits).} The arguments are translates by $o_n$ of fixed bonds, sites and plaquettes. For $\varepsilon=0$ the term $K^{(n)}(x-x'+y_0)$ has a fixed argument and converges to the corresponding kernel of $\mathbb Z^3$, by Lemma~\ref{IV-73-75}(b) and Lemma~\ref{IV-73-51}. If $\varepsilon_\nu=1$, the $\nu$-th coordinate of the argument is $\tilde x_\nu+\tilde x'_\nu+2\lfloor n/2\rfloor+O(1)$, where $\tilde x=x-o_n$ and $\tilde x'=x'-o_n$ are fixed; its distance from $2n\mathbb Z$ is therefore at least $n-C$, and the term tends to $0$ by Lemma~\ref{IV-73-75}(a).

\emph{Step 5 (conclusion).} For $j=0$, $\mathsf B_0^{[n]}=\frac{\mathsf n}2(\Pi^{[n]})^2$ by \eqref{IV-73-23}, which converges by Step~4 to $\frac{\mathsf n}2(\Pi^{\mathrm{lat}})^2=\mathsf B_0^{\mathrm{lat}}$. For $j=1$, the formulas of Lemma~\ref{IV-73-53}(c) in the box, which are finite sums, express $\mathsf B_1^{[n]}(p+o_n,q+o_n;\tau_{\mathrm{gf}})$ as a sum, over monomials, pairings and positions, of products of box kernels. We write the positions in the box as translates $P+o_n$, $P'+o_n$, $b+o_n,\dots$ of plaquettes $P,P'$ and bonds $b,\dots$ of $\mathbb Z^3$, and we index the terms by the monomial, the pairing and the positions $P,P',b,\dots$ in $\mathbb Z^3$; this set of indices does not depend on $n$, and a term whose positions do not all lie in the box is set equal to $0$.

For fixed $P,P',b,\dots$ all positions lie in the box for every large $n$, and by Step~4 the term converges, as $n\to\infty$, to the term of the formulas of Lemma~\ref{IV-73-53}(c) on $\mathbb Z^3$ with the same monomial, the same pairing and the positions $P,P',b,\dots$. By Step~3, the box kernels satisfy \eqref{IV-73-56} with one constant $\mathsf c_0$; these bounds depend only on differences of base points, which the translation by $o_n$ does not change, and the uniform form of Lemma~\ref{IV-73-55}, valid for any kernels with these bounds, therefore dominates the moduli of the terms, uniformly in $n$, by a family of majorants over the same set of indices whose sum is at most $C\mathsf c_0^5$.

By dominated convergence, $\mathsf B_1^{[n]}(p+o_n,q+o_n;\tau_{\mathrm{gf}})$ tends to the sum of the terms on $\mathbb Z^3$, which converges absolutely and equals $\mathsf B_1^{\mathrm{lat}}(p,q;\tau_{\mathrm{gf}})$ by Lemma~\ref{IV-73-55}. This is \eqref{IV-73-77}.
\end{proof}

\begin{corollary}[Independence of the gauge parameter]\label{IV-73-79}
For every $\tau_{\mathrm{gf}}>0$ and all plaquettes $p,q$ of $\mathbb Z^3$,
\[
\mathsf B_1^{\mathrm{lat}}(p,q;\tau_{\mathrm{gf}})=\mathsf B_1^{\mathrm{lat}}(p,q).
\]
\end{corollary}

\begin{proof}
By Proposition~\ref{IV-73-27}, $\mathsf B_1^{[n]}(p+o_n,q+o_n;\tau_{\mathrm{gf}})=\mathsf B_1^{[n]}(p+o_n,q+o_n;1)$ for every $n$ with $p+o_n,q+o_n\in\mathsf{Pl}_n$. Letting $n\to\infty$, Proposition~\ref{IV-73-76} gives $\mathsf B_1^{\mathrm{lat}}(p,q;\tau_{\mathrm{gf}})=\mathsf B_1^{\mathrm{lat}}(p,q;1)$, and the right side is $\mathsf B_1^{\mathrm{lat}}(p,q)$ by \eqref{IV-73-25}.
\end{proof}

\subsection{The value: Fourier constants, coefficient of $\mathsf p^2\log|\mathsf p|$}

To compute $\kappa^{\mathrm{lat}}$ we relate each $\kappa_\Gamma^{\mathrm{lat}}$ to the coefficient of $\mathsf p^2\log|\mathsf p|$ in the expansion of $\Psi_\Gamma$ at small $\mathsf p$. This requires two Fourier constants and the corresponding expansions of two lattice sums. In this subsection $(\mathcal F^{-1}\phi)(p)=\int e^{ip\cdot x}\phi(x)\,d^\circ x$ for a function $\phi$ of $x\in\mathbb R^3$.

\begin{lemma}[Two Fourier constants]\label{IV-73-80}
Let $\phi\in\mathcal S(\mathbb R^3)$ vanish near $0$.

\textup{(a)} For $s\in\mathbb C$ with $\operatorname{Re}s>-3$,
\begin{equation}\label{IV-73-81}
\begin{gathered}
\int_{\mathbb R^3}|p|^s\,(\mathcal F^{-1}\phi)(p)\,dp=c(s)\int_{\mathbb R^3}|x|^{-3-s}\,\phi(x)\,dx,\\
c(s):=\frac{2^s\,\Gamma_{\mathrm E}(\frac{3+s}2)}{\pi^{3/2}\,\Gamma_{\mathrm E}(-\frac s2)},
\end{gathered}
\end{equation}
where $1/\Gamma_{\mathrm E}$ is entire. Moreover $c(s)=0$ for $s=0,2,4,\dots$, and $c(1)=-\pi^{-2}$.

\textup{(b)} We have
\begin{equation}\label{IV-73-82}
\int_{\mathbb R^3}p^2\log|p|\,(\mathcal F^{-1}\phi)(p)\,dp=c'(2)\int_{\mathbb R^3}|x|^{-5}\,\phi(x)\,dx,\qquad c'(2)=\frac3{2\pi}.
\end{equation}
\end{lemma}

\begin{proof}
\emph{Step 1.} Let $\psi:=\mathcal F^{-1}\phi\in\mathcal S(\mathbb R^3)$, and assume first $-3<\operatorname{Re}s<0$. The substitution $t\to t/p^2$ in Euler's integral gives $\int_0^\infty t^{-s/2-1}e^{-tp^2}\,dt=\Gamma_{\mathrm E}(-\frac s2)\,|p|^s$ for $p\ne0$. Since
\[
\iint t^{-\operatorname{Re}s/2-1}e^{-tp^2}\,|\psi(p)|\,dt\,dp=\Gamma_{\mathrm E}\bigl(-\tfrac12\operatorname{Re}s\bigr)\int|p|^{\operatorname{Re}s}|\psi(p)|\,dp<\infty ,
\]
Fubini's theorem yields
\[
\int|p|^s\psi(p)\,dp=\Gamma_{\mathrm E}\bigl(-\tfrac s2\bigr)^{-1}\int_0^\infty t^{-s/2-1}\Bigl[\int e^{-tp^2}\psi(p)\,dp\Bigr]dt .
\]
By Fubini's theorem and the Gaussian integral $\int_{\mathbb R^3}e^{-tp^2+ip\cdot x}\,dp=(\pi/t)^{3/2}e^{-x^2/4t}$,
\[
\int e^{-tp^2}\psi(p)\,dp=\int\phi(x)\Bigl[\int e^{-tp^2+ip\cdot x}\,d^\circ p\Bigr]dx=\int(4\pi t)^{-3/2}e^{-x^2/4t}\,\phi(x)\,dx .
\]
For $x\ne0$ the substitution $t\to x^2/4t$ gives
\[
\int_0^\infty t^{-(s+3)/2-1}e^{-x^2/4t}\,dt=\Gamma_{\mathrm E}\bigl(\tfrac{s+3}2\bigr)\Bigl(\frac{x^2}4\Bigr)^{-(s+3)/2},
\]
as $\operatorname{Re}s+3>0$. The same identity with $\operatorname{Re}s$ in place of $s$,
\[
\int_0^\infty t^{-(\operatorname{Re}s+3)/2-1}e^{-x^2/4t}\,dt=\Gamma_{\mathrm E}\bigl(\tfrac{\operatorname{Re}s+3}2\bigr)\Bigl(\frac{x^2}4\Bigr)^{-(\operatorname{Re}s+3)/2},
\]
together with $\int|x|^{-\operatorname{Re}s-3}|\phi(x)|\,dx<\infty$, which holds as $\phi$ vanishes near $0$ and lies in $\mathcal S(\mathbb R^3)$, shows that the modulus $(4\pi)^{-3/2}t^{-(\operatorname{Re}s+3)/2-1}e^{-x^2/4t}|\phi(x)|$ of the integrand $t^{-s/2-1}(4\pi t)^{-3/2}e^{-x^2/4t}\phi(x)$ is integrable over $(0,\infty)\times\mathbb R^3$. Hence Fubini's theorem applies once more, and
\[
\int|p|^s\psi(p)\,dp=\Gamma_{\mathrm E}\bigl(-\tfrac s2\bigr)^{-1}(4\pi)^{-3/2}\,2^{s+3}\,\Gamma_{\mathrm E}\bigl(\tfrac{s+3}2\bigr)\int|x|^{-s-3}\phi(x)\,dx ,
\]
with $(4\pi)^{-3/2}2^{s+3}=2^s\pi^{-3/2}$. This is \eqref{IV-73-81} for $-3<\operatorname{Re}s<0$.

\emph{Step 2.} The left side of \eqref{IV-73-81} is holomorphic in $s$ on $\{\operatorname{Re}s>-3\}$, since $|p|^{\operatorname{Re}s}|\psi|$ provides integrable majorants locally uniformly in $s$; the function $s\mapsto\int|x|^{-3-s}\phi(x)\,dx$ is entire, since $\phi$ vanishes near $0$ and lies in $\mathcal S(\mathbb R^3)$; and $c$ is holomorphic on $\{\operatorname{Re}s>-3\}$, since $\Gamma_{\mathrm E}$ is holomorphic on the right half plane and $1/\Gamma_{\mathrm E}$ is entire. Both sides agree for $-3<\operatorname{Re}s<0$, hence on the connected set $\{\operatorname{Re}s>-3\}$ by the identity theorem. The entire functions $1/\Gamma_{\mathrm E}(z)$ and $z/\Gamma_{\mathrm E}(z+1)$ agree for $\operatorname{Re}z>0$, hence for all $z$; therefore $1/\Gamma_{\mathrm E}$ vanishes at $0$, and then successively at $-1,-2,\dots$, which gives $c(s)=0$ for $s=0,2,4,\dots$.

Moreover $1/\Gamma_{\mathrm E}(-\frac12)=-\frac12/\Gamma_{\mathrm E}(\frac12)=-\frac1{2\sqrt\pi}$, whence $c(1)=\frac{2\,\Gamma_{\mathrm E}(2)}{\pi^{3/2}}\cdot\bigl(-\frac1{2\sqrt\pi}\bigr)=-\pi^{-2}$.

\emph{Step 3.} For (b) we differentiate \eqref{IV-73-81} with respect to $s$ at $s=2$. On the left we obtain $\int p^2\log|p|\,\psi(p)\,dp$, by differentiation under the integral; on the right $c'(2)\int|x|^{-5}\phi-c(2)\int|x|^{-5}\log|x|\,\phi$, where $c(2)=0$. From $1/\Gamma_{\mathrm E}(z)=z(z+1)/\Gamma_{\mathrm E}(z+2)$ the derivative of $1/\Gamma_{\mathrm E}$ at $z=-1$ is $(2z+1)/\Gamma_{\mathrm E}(z+2)|_{z=-1}=-1$; with $z=-\frac s2$ this gives $\partial_s\bigl[1/\Gamma_{\mathrm E}(-\frac s2)\bigr]_{s=2}=\frac12$. As $1/\Gamma_{\mathrm E}(-1)=0$, only this derivative contributes, and $c'(2)=2^2\,\Gamma_{\mathrm E}(\frac52)\,\pi^{-3/2}\cdot\frac12$ with $\Gamma_{\mathrm E}(\frac52)=\frac32\cdot\frac12\sqrt\pi=\frac34\sqrt\pi$; thus $c'(2)=\frac3{2\pi}$.
\end{proof}

\begin{lemma}[Two Fourier series: Epstein-type sums near $\mathsf p=0$]\label{IV-73-83}
For $n\in\{4,5\}$ and $\mathsf p\in\mathbb R^3$ let $\mathsf s_n(\mathsf p):=\sum_{u\in\mathbb Z^3\setminus\{0\}}|u|^{-n}e^{-i\mathsf p\cdot u}$. There are functions $\mathsf e_4,\mathsf e_5\in C^\infty(\{|\mathsf p|<1\})$, invariant under the $48$ signed permutations of the coordinates, such that
\[
\mathsf s_4(\mathsf p)=-\pi^2|\mathsf p|+\mathsf e_4(\mathsf p),\qquad\mathsf s_5(\mathsf p)=\tfrac{2\pi}3\,\mathsf p^2\log|\mathsf p|+\mathsf e_5(\mathsf p)\qquad(0<|\mathsf p|<1).\qedhere
\]
\end{lemma}

\begin{proof}
We treat the two cases $(\varpi,n,\mathsf c)=(|p|,4,-\pi^{-2})$ and $(\varpi,n,\mathsf c)=(p^2\log|p|,5,\frac3{2\pi})$ together. The function $\varpi$ is continuous, radial and smooth off $0$, and its derivatives of order $j$ are $O(|p|^{2-j}\lg|p|)$ at infinity. For an integer $\mathsf N\ge3$ let
\[
\mathsf K(x):=\int e^{ip\cdot x}\chi(p)\,\varpi(p)\,d^\circ p,\qquad\mathsf E(x):=|x|^{-2\mathsf N}\int e^{ip\cdot x}(-\Delta_p)^{\mathsf N}\bigl[(1-\chi)\varpi\bigr](p)\,d^\circ p\quad(x\ne0).
\]

\emph{Step 1.} The integrand of $\mathsf E$ vanishes for $|p|\le1$ and is $O(|p|^{2-2\mathsf N}\lg|p|)$, with $2\mathsf N-2>3$; hence the integral converges absolutely. The value of $\mathsf E(x)$ does not depend on $\mathsf N$: we write $|x|^2e^{ip\cdot x}=-\Delta_pe^{ip\cdot x}$ and integrate by parts twice over balls (Gauss' theorem); the boundary terms on the sphere of radius $\mathsf r$ are $O\bigl((1+|x|)\mathsf r^2\cdot\mathsf r^{2-2\mathsf N}\lg\mathsf r\bigr)$, which tends to $0$ as $\mathsf r\to\infty$ because $\mathsf N\ge3$; hence $\mathsf N$ and $\mathsf N+1$ give the same value. Consequently $\mathsf E$ is continuous off $0$ and radial, and $|\mathsf E(x)|\le C_{\mathsf N}|x|^{-2\mathsf N}$ for every $\mathsf N\ge3$.

\emph{Step 2.} We show $\mathsf K+\mathsf E=\mathsf c\,|x|^{-n}$ on $\mathbb R^3\setminus\{0\}$. Let $\phi\in C_c^\infty(\mathbb R^3\setminus\{0\})$. By Fubini's theorem, $\int\mathsf K\phi\,dx=\int\chi\varpi\,\mathcal F^{-1}\phi\,dp$. Let $\phi_{\mathsf N}:=|x|^{-2\mathsf N}\phi\in C_c^\infty(\mathbb R^3\setminus\{0\})$. By Fubini's theorem and $2\mathsf N$ integrations by parts, which are permitted as $\mathcal F^{-1}\phi_{\mathsf N}\in\mathcal S(\mathbb R^3)$, and by $(-\Delta_p)^{\mathsf N}\mathcal F^{-1}\phi_{\mathsf N}=\mathcal F^{-1}(|x|^{2\mathsf N}\phi_{\mathsf N})=\mathcal F^{-1}\phi$,
\[
\int\mathsf E\phi\,dx=\int(-\Delta_p)^{\mathsf N}\bigl[(1-\chi)\varpi\bigr]\,\mathcal F^{-1}\phi_{\mathsf N}\,dp=\int(1-\chi)\varpi\,\mathcal F^{-1}\phi\,dp .
\]
The sum is $\int\varpi\,\mathcal F^{-1}\phi\,dp=\mathsf c\int|x|^{-n}\phi\,dx$, by Lemma~\ref{IV-73-80}(a) at $s=1$ in the first case and by Lemma~\ref{IV-73-80}(b) in the second. As $\mathsf K+\mathsf E-\mathsf c|x|^{-n}$ is continuous on $\mathbb R^3\setminus\{0\}$, it vanishes there.

\emph{Step 3.} The periodic continuation of $\chi\varpi$, whose support lies in $\{|p|\le2\}$ and thus in the interior of $[-\pi,\pi]^3$, is continuous, and its Fourier coefficients are the numbers $\mathsf K(u)$, $u\in\mathbb Z^3$. By Steps 1 and 2, $\sum_u|\mathsf K(u)|<\infty$; hence the continuous function $\sum_u\mathsf K(u)e^{-i\mathsf p\cdot u}$ has the same Fourier coefficients and, both functions being continuous, coincides with the periodic continuation of $\chi\varpi$. For $|\mathsf p|<1$ we have $\chi(\mathsf p)=1$, and therefore
\[
\varpi(\mathsf p)=\mathsf K(0)+\mathsf c\,\mathsf s_n(\mathsf p)-\sum_{u\ne0}\mathsf E(u)\,e^{-i\mathsf p\cdot u} .
\]
We put $\mathsf e_n(\mathsf p):=\mathsf c^{-1}\bigl[\sum_{u\ne0}\mathsf E(u)e^{-i\mathsf p\cdot u}-\mathsf K(0)\bigr]$. This function is smooth, by the rapid decay of $\mathsf E$, and invariant under the signed permutations of the coordinates, since $\mathsf E$ is radial. Finally $\mathsf c^{-1}=-\pi^2$, resp.\ $\mathsf c^{-1}=\frac{2\pi}3$, which gives the claim.
\end{proof}

\begin{lemma}[Coefficient of $\mathsf p^2\log|\mathsf p|$ in $\Psi_\Gamma$; scaling of the cut-off]\label{IV-73-84}
Let $\Gamma\in\{1,\dots,6\}$.

\textup{(a)} There are numbers $\mathsf a_0,\mathsf a_1,\mathsf a_2$ and a function $\mathsf o$ on $\{|\mathsf p|<1\}$ with $\mathsf o(\mathsf p)/\mathsf p^2\to0$ as $\mathsf p\to0$, depending on $\Gamma$ only, such that for almost every $\mathsf p$ with $|\mathsf p|<1$
\[
\Psi_\Gamma(\mathsf p)=\mathsf a_0+\mathsf a_1|\mathsf p|+\mathsf a_2\mathsf p^2+\tfrac\pi{36}\,\kappa_\Gamma^{\mathrm{lat}}\,\mathsf p^2\log|\mathsf p|+\mathsf o(\mathsf p).
\]

\textup{(b)} For real $\Lambda>1$ let $R_\Gamma^\Lambda(k,l,\mathsf p):=R_\Gamma(k/\Lambda,l/\Lambda,\mathsf p/\Lambda)$. Then $\iint R_\Gamma^\Lambda F_\Gamma\,d^\circ k\,d^\circ l=\Lambda^2\Psi_\Gamma(\mathsf p/\Lambda)$ for every $\mathsf p\ne0$, and for almost every $\mathsf p$ with $|\mathsf p|<1$
\[
\begin{aligned}
\iint(R_\Gamma^\Lambda-R_\Gamma)F_\Gamma\,d^\circ k\,d^\circ l={}&(\Lambda^2-1)\mathsf a_0+(\Lambda-1)\mathsf a_1|\mathsf p|-\tfrac\pi{36}\,\kappa_\Gamma^{\mathrm{lat}}\log\Lambda\;\mathsf p^2\\
&+\Lambda^2\mathsf o(\mathsf p/\Lambda)-\mathsf o(\mathsf p).
\end{aligned}
\]
\end{lemma}

\begin{proof}
\emph{Step 1.} By Lemmas~\ref{IV-73-69} and~\ref{IV-73-70},
\begin{gather*}
\mathsf T_\Gamma^{\mathrm{mod}}(u)=c_{4,\Gamma}|u|^{-4}+\kappa_\Gamma^{\mathrm{lat}}|u|^{-5}+\mathsf r(u),\\
|\mathsf r(u)|\le C|u|^{-6}\lg|u|\qquad(u\in\mathbb Z^3\setminus\{0\});
\end{gather*}
we put $\mathsf r(0):=\mathsf T^{\mathrm{mod}}_\Gamma(0)$. Hence $\sum_u|\mathsf T_\Gamma^{\mathrm{mod}}(u)|<\infty$. The function $\Psi_\Gamma$ is integrable and vanishes for $|\mathsf p|\ge\frac{3\pi}8$ (Definition~\ref{IV-73-64}(3)), that is, outside a ball contained in $[-\pi,\pi]^3$. Therefore the periodic function which equals $24\Psi_\Gamma$ on $[-\pi,\pi]^3$ has the Fourier coefficients $24\int_{[-\pi,\pi]^3}e^{i\mathsf p\cdot u}\Psi_\Gamma(\mathsf p)\,d^\circ\mathsf p=\mathsf T_\Gamma^{\mathrm{mod}}(u)$, which are also the Fourier coefficients of the continuous function $\sum_u\mathsf T_\Gamma^{\mathrm{mod}}(u)e^{-i\mathsf p\cdot u}$. By the uniqueness of Fourier coefficients the two periodic functions agree almost everywhere; as the first one coincides with $24\Psi_\Gamma$ on $[-\pi,\pi]^3$, we obtain, for almost every $\mathsf p\in[-\pi,\pi]^3$,
\begin{gather*}
24\,\Psi_\Gamma(\mathsf p)=c_{4,\Gamma}\,\mathsf s_4(\mathsf p)+\kappa_\Gamma^{\mathrm{lat}}\,\mathsf s_5(\mathsf p)+\hat{\mathsf r}(\mathsf p),\\
\hat{\mathsf r}(\mathsf p):=\sum_u\mathsf r(u)\,e^{-i\mathsf p\cdot u} .
\end{gather*}

\emph{Step 2.} Since $\sum_u|u|^2|\mathsf r(u)|<\infty$, the function $\hat{\mathsf r}$ is of class $C^2$. As $\mathsf T_\Gamma^{\mathrm{mod}}$ is the restriction to $\mathbb Z^3$ of a radial function, $\mathsf r$ and $\hat{\mathsf r}$ are invariant under the signed permutations of the coordinates. A function $\mathsf f$ of class $C^2$ near $0$ with this invariance has vanishing gradient at $0$ (invariance under $-\mathbf 1$), a Hessian at $0$ with vanishing off-diagonal entries (change of sign of one coordinate) and equal diagonal entries (permutations); hence $\mathsf f(\mathsf p)=\mathsf f(0)+\mathsf a\,\mathsf p^2+o(\mathsf p^2)$. We apply this to $\hat{\mathsf r}$, $\mathsf e_4$ and $\mathsf e_5$ and insert Lemma~\ref{IV-73-83}; this gives (a), the coefficient of $\mathsf p^2\log|\mathsf p|$ being $\frac1{24}\cdot\frac{2\pi}3\,\kappa_\Gamma^{\mathrm{lat}}=\frac\pi{36}\,\kappa_\Gamma^{\mathrm{lat}}$, and $\mathsf o$ being $\frac1{24}$ times the sum of the remainders $o(\mathsf p^2)$ of $c_{4,\Gamma}\mathsf e_4$, $\kappa_\Gamma^{\mathrm{lat}}\mathsf e_5$ and $\hat{\mathsf r}$.

\emph{Step 3.} For (b) we substitute $k\to\Lambda k$, $l\to\Lambda l$ and use $F_\Gamma(\Lambda k,\Lambda l,\Lambda\mathsf q)=\Lambda^{-4}F_\Gamma(k,l,\mathsf q)$ with $\mathsf q=\mathsf p/\Lambda$; this produces the factor $\Lambda^{6-4}=\Lambda^2$ and gives the first assertion. For the second we insert (a) at $\mathsf p/\Lambda$ and at $\mathsf p$, both of modulus less than $1$. The terms with $\mathsf a_2$ cancel, the terms with $\mathsf a_0$, $\mathsf a_1$ give $(\Lambda^2-1)\mathsf a_0+(\Lambda-1)\mathsf a_1|\mathsf p|$, and
\[
\Lambda^2\,\tfrac\pi{36}\,\kappa_\Gamma^{\mathrm{lat}}\Lambda^{-2}\mathsf p^2\bigl(\log|\mathsf p|-\log\Lambda\bigr)-\tfrac\pi{36}\,\kappa_\Gamma^{\mathrm{lat}}\,\mathsf p^2\log|\mathsf p|=-\tfrac\pi{36}\,\kappa_\Gamma^{\mathrm{lat}}\log\Lambda\;\mathsf p^2 .
\]
\end{proof}

\subsection{The value: finite parts and regions}

The coefficient $\frac\pi{36}\kappa_\Gamma^{\mathrm{lat}}$ of Lemma~\ref{IV-73-84} will be computed by an expansion by regions of two-loop integrals with cut-offs vanishing near the origin. We first define the finite parts which occur in this expansion, by analytic regularisation, and the classes of cut-offs.

\begin{definition}[Finite parts (analytic regularisation); cut-offs vanishing near the origin]\label{IV-73-85}
Let $\mathbb D:=\{\mathsf z\in\mathbb C:\operatorname{Re}\mathsf z>-\frac12\}\setminus\{\frac12,\frac32,\frac52,\dots\}$; this set is open and connected, and it contains $0$.

\emph{(1) One variable.} A \emph{homogeneous factor of degree $d\in\mathbb Z$} is a function $\mathsf h\in C^\infty(\mathbb R^3\setminus\{0\})$ with $\mathsf h(tk)=t^d\mathsf h(k)$ for $t>0$ and $\mathsf h(-k)=(-1)^d\mathsf h(k)$. Let $\mathfrak P$ be the set of functions on $\mathbb R^3\setminus\{0\}$ of the form $\mathsf g=\sum_\iota\mathsf h_\iota\mathsf g_\iota$, a finite sum, with homogeneous factors $\mathsf h_\iota$ of degrees $d_\iota$ and with $\mathsf g_\iota\in C_c^\infty(\mathbb R^3)$. For such $\mathsf g$ the integral $\int|k|^{2\mathsf z}\mathsf g(k)\,dk$ converges absolutely and is holomorphic in $\mathsf z$ for $2\operatorname{Re}\mathsf z+\min_\iota d_\iota>-3$; it has a holomorphic continuation to $\mathbb D$ (Lemma~\ref{IV-73-86}(a)), which is unique. The \emph{finite part} $\operatorname{f.p.}\!\int\mathsf g\,dk$ is the value of this continuation at $\mathsf z=0$; it depends on $\mathsf g$ only, and it is linear in $\mathsf g$.

\emph{(2) Two variables.} Let $c>0$ and $\eta_c(k):=\theta(4|k|/c)$, so that $\eta_c(k)=1$ for $|k|\le\frac c8$ and $\eta_c(k)=0$ for $|k|\ge\frac c4$. Let $\mathfrak P_2(c)$ be the set of Borel functions $\Phi$ on $\{(k,l):k\ne0,\ l\ne0,\ k\ne l\}$ for which there is a constant $C_\Phi$ with the following properties.
\begin{itemize}
\item[($\Phi$0)] $\Phi(k,l)=0$ if $|k|+|l|<c$ or $|k|+|l|>C_\Phi$.
\item[($\Phi$1)] For $0<|k|<\frac c3$, $\Phi(k,l)=\sum_\iota\mathsf h_\iota(k)\Phi_\iota(k,l)$, a finite sum with homogeneous factors $\mathsf h_\iota$ and with functions $\Phi_\iota\in C^\infty(\{|k|<\frac c3\}\times\mathbb R^3)$ which vanish where $|k|+|l|<c$ or $|k|+|l|>C_\Phi$.
\item[($\Phi$2)] For $0<|l|<\frac c3$, $\Phi(k,l)=\sum_\iota\mathsf h'_\iota(l)\Phi'_\iota(k,l)$, with the same properties after the exchange of $k$ and $l$.
\item[($\Phi$3)] $|\Phi(k,l)|\le C_\Phi|k-l|^{-2}$ for $|k|,|l|\ge\frac c8$.
\end{itemize}
For $\Phi\in\mathfrak P_2(c)$ the function
\[
k\mapsto\eta_c(k)\int\Phi(k,l)\,dl=\sum_\iota\mathsf h_\iota(k)\,\eta_c(k)\int\Phi_\iota(k,l)\,dl
\]
lies in $\mathfrak P$, since $\eta_c$ vanishes for $|k|\ge\frac c4$ and the $\Phi_\iota$ are smooth and vanish for $|l|>C_\Phi$; likewise $l\mapsto\eta_c(l)\int\Phi(k,l)\,dk$ lies in $\mathfrak P$. We define
\[
\begin{aligned}
\operatorname{f.p.}\!\iint\Phi\,dk\,dl:={}&\iint\bigl(1-\eta_c(k)-\eta_c(l)\bigr)\Phi\,dk\,dl\\
&+\operatorname{f.p.}\!\int\Bigl[\eta_c(k)\!\int\Phi(k,l)\,dl\Bigr]dk+\operatorname{f.p.}\!\int\Bigl[\eta_c(l)\!\int\Phi(k,l)\,dk\Bigr]dl .
\end{aligned}
\]
On the support of $\Phi$, $|k|<\frac c8$ implies $|l|>\frac78c$, hence $\eta_c(k)=1$ and $\eta_c(l)=0$; the same holds with $k$ and $l$ exchanged. Therefore the first integrand vanishes unless $|k|,|l|\ge\frac c8$, and the first integral converges absolutely, by ($\Phi$0) and ($\Phi$3): its integrand is bounded by $C_\Phi|k-l|^{-2}$, which is integrable over bounded sets of $(k,l)$, and it vanishes for $|k|+|l|>C_\Phi$. This finite part is linear on $\mathfrak P_2(c)$ and symmetric under the exchange of $k$ and $l$.

\emph{(3) Cut-offs.} Let $\mathcal Y(c)$ be the set of functions $\Upsilon\in C^\infty(\mathbb R^3\times\mathbb R^3\times\mathbb R^3)$ for which there is a constant $C_\Upsilon$ with the following properties.
\begin{itemize}
\item[(Y1)] $\Upsilon(k,l,\mathsf p)=0$ if $|\mathsf p|<c$ and ($|k|+|l|<c$ or $|k|+|l|>C_\Upsilon$).
\item[(Y2)] $\Upsilon(-k,-l,-\mathsf p)=\Upsilon(k,l,\mathsf p)=\Upsilon(Ok,Ol,O\mathsf p)$ for every rotation $O$.
\end{itemize}
Then $\mathcal Y(c')\subset\mathcal Y(c)$ for $c<c'$. For $\Lambda>1$ and a function $\Upsilon$ of $(k,l,\mathsf p)$ or of $(k,l)$, let $\Upsilon^\Lambda$ denote $\Upsilon$ evaluated at the arguments divided by $\Lambda$; if $\Upsilon\in\mathcal Y(c)$, then $\Upsilon^\Lambda\in\mathcal Y(\Lambda c)\subset\mathcal Y(c)$. For a function $G(k,l,\mathsf p)$ which is of class $C^2$ in $\mathsf p$ near $0$ we put
\[
\overline{G_2}(k,l):=\tfrac16\,\Delta_{\mathsf p}G(k,l,\mathsf p)\big|_{\mathsf p=0} ;
\]
this is the mean over the directions of $\mathsf p$ of the term of second order in the Taylor expansion of $G$ in $\mathsf p$, divided by $\mathsf p^2$, since the mean of $\mathsf p_i\mathsf p_j$ over $\{|\mathsf p|=t\}$ is $\frac13t^2\delta_{ij}$. Finally, for $\Upsilon\in\mathcal Y(c)$ and $0<|\mathsf p|<c$ we put
\[
\Xi_\Gamma[\Upsilon](\mathsf p):=\iint\Upsilon F_\Gamma\,d^\circ k\,d^\circ l ;
\]
this integral converges absolutely by \eqref{IV-73-66} and the estimates following it in Definition~\ref{IV-73-64}(3), since by (Y1) the integrand vanishes for $|k|+|l|>C_\Upsilon$.
\end{definition}

\begin{lemma}[Properties of finite parts]\label{IV-73-86}
\begin{enumerate}
\item[\textup{(a)}] \textup{(}Continuation with a parameter.\textup{)} Let $\mathsf h$ be a homogeneous factor of degree $d$, and let $\mathsf g\colon\mathbb R^3\times\mathbb C\to\mathbb C$ have the following properties: for every $\mathsf z\in\mathbb C$ the function $\mathsf g(\cdot;\mathsf z)$ is smooth, and its support lies in a compact set which does not depend on $\mathsf z$; for every multi-index $\alpha$ the derivative $\partial_k^\alpha\mathsf g$ is continuous on $\mathbb R^3\times\mathbb C$ and, for each fixed $k$, an entire function of $\mathsf z$. Then the function
\[
\mathsf z\longmapsto\int|k|^{2\mathsf z}\,\mathsf h(k)\,\mathsf g(k;\mathsf z)\,dk ,
\]
which is defined for $2\operatorname{Re}\mathsf z+d>-3$, has a holomorphic continuation to $\mathbb D$, and the value of this continuation at $\mathsf z=0$ is $\operatorname{f.p.}\!\int\mathsf h\,\mathsf g(\cdot;0)\,dk$.
\item[\textup{(b)}] If $\mathsf g\in\mathfrak P$ is integrable, then $\operatorname{f.p.}\!\int\mathsf g\,dk=\int\mathsf g\,dk$.
\item[\textup{(c)}] Let $\mathsf h$ be a homogeneous factor of degree $d\le-4$, and let $\eta(k)=\theta_1(|k|)$, where $\theta_1\in C^\infty([0,\infty))$ has compact support and is equal to $1$ near $0$. Then
\[
\operatorname{f.p.}\!\int\eta\,\mathsf h\,dk=-\int(1-\eta)\,\mathsf h\,dk .
\]
\item[\textup{(d)}] \textup{(}A soft line.\textup{)} Let $\mu_0>0$ and $\eta(k):=\theta(|k|/2\mu_0)$. Let
\[
S(k,\mathsf p)=\mathsf P(k,\mathsf p)\,|k|^{-2\mathsf a}\,|k-\mathsf p|^{-2\mathsf a'} ,
\]
where $\mathsf a,\mathsf a'\ge0$ are integers and $\mathsf P$ is a polynomial, homogeneous of degree $2\mathsf a+2\mathsf a'-2$, and assume that $|S(k,\mathsf p)|\le C_S\bigl(|k|^{-2}+|k-\mathsf p|^{-2}\bigr)$. Let $J\in C^\infty\bigl(\{|k|<3\mu_0\}\times\{|\mathsf p|<\mu_0\}\bigr)$ satisfy $J(-k,-\mathsf p)=J(k,\mathsf p)$, and for $k\ne0$ and $i\in\{0,1,2\}$ put
\[
[SJ]_i(k;\mathsf p):=\sum_{|\alpha|=i}\frac{\mathsf p^\alpha}{\alpha!}\,\partial_{\mathsf p}^\alpha(SJ)(k,0).
\]
Then $\eta\,[SJ]_i(\cdot;\mathsf p)\in\mathfrak P$. Moreover, there are a Borel function $\mathsf h_S$ on $\mathbb R^3$, which depends on $S$ only, is bounded on the unit sphere and satisfies $\mathsf h_S(t\mathsf p)=t\,\mathsf h_S(\mathsf p)$ for $t>0$, and a constant $C=C(\mu_0,S,J)$, such that for $0<|\mathsf p|\le\frac12\mu_0$
\[
\Bigl|\int\eta SJ\,dk-\operatorname{f.p.}\!\int\eta\,[SJ]_0\,dk-\operatorname{f.p.}\!\int\eta\,[SJ]_2(k;\mathsf p)\,dk-J(0,0)\,\mathsf h_S(\mathsf p)\Bigr|\le C\,|\mathsf p|^3\lg|\mathsf p| .
\]
\item[\textup{(e)}] \textup{(}Two variables.\textup{)} Let $\Phi\in\mathfrak P_2(c)$ and $(\epsilon,\epsilon')\in\{(1,1),(1,0),(0,1)\}$. If $\epsilon=0$, assume that all degrees of the homogeneous factors in \textup{($\Phi1$)} are $\ge-2$; if $\epsilon'=0$, assume that all degrees of the homogeneous factors in \textup{($\Phi2$)} are $\ge-2$. Then the integral
\[
\iint|k|^{2\epsilon\mathsf z}\,|l|^{2\epsilon'\mathsf z}\,\Phi\,dk\,dl
\]
converges absolutely for large $\operatorname{Re}\mathsf z$, and it has a holomorphic continuation to $\mathbb D$ whose value at $\mathsf z=0$ is $\operatorname{f.p.}\!\iint\Phi\,dk\,dl$. If $0<c'<c$, then $\mathfrak P_2(c)\subset\mathfrak P_2(c')$, and every $\Phi\in\mathfrak P_2(c)$ has the same finite part whether it is regarded as an element of $\mathfrak P_2(c)$ or of $\mathfrak P_2(c')$. If $\Phi\in\mathfrak P_2(c)$ is integrable, then $\operatorname{f.p.}\!\iint\Phi\,dk\,dl=\iint\Phi\,dk\,dl$.
\end{enumerate}
\end{lemma}

\begin{proof}
\emph{Proof of \textup{(a)}.} Put $m:=\max\{0,-d-3\}$; then $d+m+1\ge-2$. For the proof of (a) let $\eta(k):=\theta(|k|)$, with the function $\theta$ of Definition~\ref{IV-73-38}(1), so that $\eta=1$ for $|k|\le\frac12$ and $\eta=0$ for $|k|\ge1$, and let
\[
T\mathsf g(k;\mathsf z):=\sum_{|\alpha|\le m}\frac{k^\alpha}{\alpha!}\,\partial_k^\alpha\mathsf g(0;\mathsf z)
\]
be the Taylor polynomial of degree $m$ of $\mathsf g(\cdot;\mathsf z)$ at $k=0$. For $2\operatorname{Re}\mathsf z+d>-3$ we decompose
\begin{equation*}
\int|k|^{2\mathsf z}\mathsf h\,\mathsf g\,dk=\int|k|^{2\mathsf z}\mathsf h\,\bigl[\mathsf g-\eta\,T\mathsf g\bigr]\,dk+\sum_{|\alpha|\le m}\frac{\partial_k^\alpha\mathsf g(0;\mathsf z)}{\alpha!}\int|k|^{2\mathsf z}\,\eta(k)\,\mathsf h(k)\,k^\alpha\,dk .
\end{equation*}
The integrand of the first integral on the right vanishes outside a compact set which does not depend on $\mathsf z$. For $|k|\le\frac12$ we have $\eta=1$, and Taylor's formula, together with the continuity of the derivatives of order $m+1$ of $\mathsf g$ on $\mathbb R^3\times\mathbb C$, bounds the modulus of the integrand by $C|k|^{2\operatorname{Re}\mathsf z+d+m+1}$, with a constant $C$ that stays bounded as long as $\mathsf z$ ranges over a compact set. Since $d+m+1\ge-2$, the exponent is larger than $-3$ as soon as $\operatorname{Re}\mathsf z>-\frac12$. For each $k\ne0$ the integrand is an entire function of $\mathsf z$; differentiating under the integral sign, we see that the first integral is holomorphic on $\{\operatorname{Re}\mathsf z>-\frac12\}$.

In the integrals of the sum we pass to polar coordinates and integrate by parts once. With
\[
\mathsf m_\alpha:=\int_{S^2}\mathsf h(\omega)\,\omega^\alpha\,d\omega
\]
we obtain, for $2\operatorname{Re}\mathsf z+d+|\alpha|>-3$,
\[
\int|k|^{2\mathsf z}\,\eta\,\mathsf h\,k^\alpha\,dk=\mathsf m_\alpha\int_0^\infty\theta(s)\,s^{2\mathsf z+d+|\alpha|+2}\,ds=-\frac{\mathsf m_\alpha}{2\mathsf z+d+|\alpha|+3}\int_0^\infty\theta'(s)\,s^{2\mathsf z+d+|\alpha|+3}\,ds .
\]
As $\theta'$ is supported in $[\frac12,1]$, the last integral is an entire function of $\mathsf z$. If $d+|\alpha|$ is odd, then $\omega\mapsto\mathsf h(\omega)\omega^\alpha$ is odd and $\mathsf m_\alpha=0$. If $d+|\alpha|$ is even, the denominator vanishes only at $\mathsf z=-\frac12(d+|\alpha|+3)$, which is half an odd integer and hence not a point of $\mathbb D$. Consequently the right side of the decomposition is holomorphic on $\mathbb D$; it is the asserted continuation, which is unique because $\mathbb D$ is connected.

The value of this continuation at $\mathsf z=0$ involves $\mathsf g$ only through $\mathsf g(\cdot;0)$: it contains the function $\mathsf g(\cdot;0)-\eta\,T\mathsf g(\cdot;0)$ and the numbers $\partial_k^\alpha\mathsf g(0;0)$. It is therefore equal to the value at $0$ of the same expression built with the function $\mathsf g(\cdot;0)$, which does not depend on $\mathsf z$. By the argument just given, applied to this function, that expression continues $\mathsf z\mapsto\int|k|^{2\mathsf z}\mathsf h\,\mathsf g(\cdot;0)\,dk$ to $\mathbb D$, and by Definition~\ref{IV-73-85}(1) its value at $0$ is $\operatorname{f.p.}\!\int\mathsf h\,\mathsf g(\cdot;0)\,dk$.

\emph{Proof of \textup{(b)}.} Let $\mathsf g\in\mathfrak P$ be integrable. Like every element of $\mathfrak P$, the function $\mathsf g$ vanishes outside a bounded set; hence for every $\mathsf N>0$ there is $C_{\mathsf N}$ with $|k|^{2\operatorname{Re}\mathsf z}|\mathsf g(k)|\le C_{\mathsf N}|\mathsf g(k)|$ whenever $0\le\operatorname{Re}\mathsf z\le\mathsf N$. By differentiation under the integral sign and by dominated convergence, $\int|k|^{2\mathsf z}\mathsf g\,dk$ is holomorphic on $\{\operatorname{Re}\mathsf z>0\}$ and continuous on $\{\operatorname{Re}\mathsf z\ge0\}$. The open set $\mathbb D\cap\{\operatorname{Re}\mathsf z>0\}$ is connected, and on it this function coincides with the continuation of Definition~\ref{IV-73-85}(1), because the two functions agree for large $\operatorname{Re}\mathsf z$ (identity theorem). Letting $\mathsf z\downarrow0$ along the real axis, we obtain $\operatorname{f.p.}\!\int\mathsf g\,dk=\int\mathsf g\,dk$.

\emph{Proof of \textup{(c)}.} Let $\mathsf m_0:=\int_{S^2}\mathsf h\,d\omega$. For $2\operatorname{Re}\mathsf z+d>-3$, polar coordinates and one integration by parts give
\[
\int|k|^{2\mathsf z}\,\eta\,\mathsf h\,dk=\mathsf m_0\int_0^\infty\theta_1(s)\,s^{2\mathsf z+d+2}\,ds=-\frac{\mathsf m_0}{2\mathsf z+d+3}\int_0^\infty\theta_1'(s)\,s^{2\mathsf z+d+3}\,ds .
\]
The support of $\theta_1'$ is a compact subset of $(0,\infty)$, so the last integral is an entire function of $\mathsf z$. As in the proof of (a), either $d$ is odd and $\mathsf m_0=0$, or the zero of the denominator is half an odd integer and lies outside $\mathbb D$. Hence the right side is the continuation of the left side to $\mathbb D$, and its value at $\mathsf z=0$ is
\[
\operatorname{f.p.}\!\int\eta\,\mathsf h\,dk=-\frac{\mathsf m_0}{d+3}\int_0^\infty\theta_1'(s)\,s^{d+3}\,ds .
\]
On the other hand, $1-\theta_1$ vanishes near $0$ and $d+3<0$; an integration by parts without boundary terms therefore yields
\[
\int(1-\eta)\,\mathsf h\,dk=\mathsf m_0\int_0^\infty\bigl(1-\theta_1(s)\bigr)\,s^{d+2}\,ds=\frac{\mathsf m_0}{d+3}\int_0^\infty\theta_1'(s)\,s^{d+3}\,ds .
\]
Comparing the last two displays, we obtain (c).

\emph{Proof of \textup{(d)}: preparation.} The function $S$ is smooth on $\{k\ne0,\ k\ne\mathsf p\}$, homogeneous of degree $-2$ in $(k,\mathsf p)$ jointly, and even, since the polynomial $\mathsf P$ is homogeneous of even degree. The set $\{(k,\mathsf q):|k|=1,\ |\mathsf q|\le\frac12\}$ is compact, and $S$ is smooth in a neighbourhood of it; together with homogeneity this gives, for every multi-index $\alpha$,
\[
|\partial_{\mathsf p}^\alpha S(k,\mathsf q)|\le C_\alpha\,|k|^{-2-|\alpha|}\qquad\bigl(|\mathsf q|\le\tfrac12|k|\bigr).
\]
For every multi-index $\beta$ the function $\partial_{\mathsf p}^\beta S(\cdot,0)$ is a homogeneous factor of degree $-2-|\beta|$: it is smooth off $0$, homogeneous of this degree, and of parity $(-1)^{|\beta|}$ because $S$ is even. By the Leibniz rule, $\eta\,[SJ]_i(\cdot;\mathsf p)$ is a finite sum of products of such factors with the functions $\eta\,\partial_{\mathsf p}^\gamma J(\cdot,0)$; these belong to $C_c^\infty(\mathbb R^3)$, since $\eta$ vanishes for $|k|\ge2\mu_0$. Hence $\eta\,[SJ]_i(\cdot;\mathsf p)\in\mathfrak P$.

Since $J$ is even, the terms of odd degree in its Taylor expansion at $(0,0)$ vanish, and its Taylor polynomial of degree $3$ is $J(0,0)+J_2$ with a quadratic form $J_2$ in $(k,\mathsf p)$. Thus $J=J(0,0)+J_2+J_4$, and Taylor's formula with integral remainder, applied to $J$ and to its derivatives, gives
\[
|\partial_{\mathsf p}^\alpha J_4(k,\mathsf p)|\le C\,(|k|+|\mathsf p|)^{4-|\alpha|}\qquad(|\alpha|\le3)\quad\text{on }\{|k|\le2\mu_0,\ |\mathsf p|\le\tfrac12\mu_0\}.
\]
Each of the four expressions inside the modulus in the asserted estimate, namely $\int\eta SJ\,dk$, $\operatorname{f.p.}\!\int\eta\,[SJ]_0\,dk$, $\operatorname{f.p.}\!\int\eta\,[SJ]_2(k;\mathsf p)\,dk$ and $J(0,0)\,\mathsf h_S(\mathsf p)$, depends linearly on $J$, because $J\mapsto[SJ]_i$ and the finite part are linear. We insert $J=J(0,0)+J_2+J_4$. The quadratic form $J_2$ and the function $J_4$ are even and vanish at $(0,0)$, so that for them the last of the four expressions is $0$. It therefore suffices to consider, in place of $SJ$, the three functions
\[
N:=S,\qquad N:=SJ_2,\qquad N:=SJ_4 ,
\]
and to show that the difference $\int\eta N\,dk-\operatorname{f.p.}\!\int\eta\,[N]_0\,dk-\operatorname{f.p.}\!\int\eta\,[N]_2\,dk$ is equal to $\mathsf h_S(\mathsf p)+O(|\mathsf p|^3)$ for $N=S$, the case which enters with the factor $J(0,0)$, and that it is $O(|\mathsf p|^3\lg|\mathsf p|)$ for $N=SJ_2$ and for $N=SJ_4$. For each of the three functions we write $N_i:=[N]_i$ ($i=0,1,2$), and for $|k|\ge2|\mathsf p|$ we put $N_{\mathrm r}:=N-N_0-N_1-N_2$; moreover, $\mathsf B:=\{|k|\le2|\mathsf p|\}$. Since $|\mathsf p|\le\frac12\mu_0$, we have $\eta=1$ on $\mathsf B$. In all three cases $N(-k,-\mathsf p)=N(k,\mathsf p)$, as $S$, $J_2$ and $J_4$ are even; differentiating with respect to $\mathsf p$ at $\mathsf p=0$, we see that $N_1$ is odd in $k$ while $N_0$ and $N_2$ are even.

\emph{The case $N=S$.} By the bound on the derivatives of $S$, $|N_i|\le C|\mathsf p|^i|k|^{-2-i}$. On the complement $\mathsf B^c$ of $\mathsf B$, Taylor's formula in $\mathsf p$, applied at the points $t\mathsf p$ with $0\le t\le1$, all of which satisfy $|t\mathsf p|\le\frac12|k|$, gives $|N_{\mathrm r}|\le C|\mathsf p|^3|k|^{-5}$. We write
\[
\int\eta N\,dk=\int_{\mathsf B}N\,dk+\int_{\mathsf B^c}\eta\,(N_0+N_1+N_2+N_{\mathrm r})\,dk
\]
and consider the four terms of the second integral one after the other. First, $\int_{\mathsf B^c}\eta N_0=\int\eta N_0-\int_{\mathsf B}N_0$. Second, $\int_{\mathsf B^c}\eta N_1=0$, because $N_1$ is odd and integrable over $\mathsf B^c\cap\operatorname{supp}\eta$, while $\eta$ and $\mathsf B^c$ are invariant under $k\to-k$. Third, $N_2(\cdot;\mathsf p)$ is a homogeneous factor of degree $-4$; assertion (c), applied with $\theta_1(s):=\theta(s/2\mu_0)$, and the vanishing of $1-\eta$ on $\mathsf B$ give
\[
\int_{\mathsf B^c}\eta N_2=\int_{\mathsf B^c}N_2-\int(1-\eta)N_2=\int_{\mathsf B^c}N_2+\operatorname{f.p.}\!\int\eta N_2 .
\]
Fourth, $\int_{\mathsf B^c}\eta N_{\mathrm r}=\int_{\mathsf B^c}N_{\mathrm r}-\int(1-\eta)N_{\mathrm r}$, and
\[
\Bigl|\int(1-\eta)N_{\mathrm r}\Bigr|\le C|\mathsf p|^3\int_{|k|\ge\mu_0}|k|^{-5}\,dk\le C|\mathsf p|^3 .
\]
Collecting the terms, we arrive at
\[
\begin{multlined}
\int\eta S\,dk=\int\eta N_0\,dk+\operatorname{f.p.}\!\int\eta N_2\,dk+\mathsf h_S(\mathsf p)+O(|\mathsf p|^3),\qquad
\mathsf h_S(\mathsf p):=\int_{|k|\le2|\mathsf p|}(S-N_0)\,dk \\
+\int_{|k|>2|\mathsf p|}(N_2+N_{\mathrm r})\,dk .
\end{multlined}
\]
Both integrals in the definition of $\mathsf h_S$ converge absolutely, because $|S|+|N_0|\le C(|k|^{-2}+|k-\mathsf p|^{-2})$, and $|N_2+N_{\mathrm r}|\le C|\mathsf p|^2|k|^{-4}$ on $\mathsf B^c$. The substitution $k\to tk$, together with $N_i(tk;t\mathsf p)=t^{-2}N_i(k;\mathsf p)$, yields $\mathsf h_S(t\mathsf p)=t\,\mathsf h_S(\mathsf p)$ for $t>0$; the same bounds, taken at $|\mathsf p|=1$, show that $\mathsf h_S$ is bounded on the unit sphere $\{|\mathsf p|=1\}$, so that $|\mathsf h_S(\mathsf p)|\le C|\mathsf p|$; and $\mathsf h_S$ is a Borel function by Fubini's theorem. By its definition $\mathsf h_S$ depends on $S$ only. Finally, $\eta N_0\in\mathfrak P$ is integrable, so that $\int\eta N_0\,dk$ is its finite part by (b).

\emph{The case $N=SJ_2$.} Now $N$ is homogeneous of degree $0$ in $(k,\mathsf p)$, and as above $|\partial_{\mathsf p}^\alpha N(k,\mathsf q)|\le C|k|^{-|\alpha|}$ for $|\mathsf q|\le\frac12|k|$. Hence $|N_i|\le C|\mathsf p|^i|k|^{-i}$, which is integrable near $k=0$ for $i\le2$; further $|N_{\mathrm r}|\le C|\mathsf p|^3|k|^{-3}$ on $\mathsf B^c$, and $|N|\le C|\mathsf p|^2(|k|^{-2}+|k-\mathsf p|^{-2})$ on $\mathsf B$. Therefore
\[
\int\eta N\,dk=\sum_{i\le2}\int\eta N_i\,dk+\Bigl[\int_{\mathsf B}(N-N_0-N_1-N_2)\,dk+\int_{\mathsf B^c}\eta N_{\mathrm r}\,dk\Bigr],
\]
and the bracket is bounded in modulus by
\[
C|\mathsf p|^3+C|\mathsf p|^3\int_{2|\mathsf p|\le|k|\le2\mu_0}|k|^{-3}\,dk\le C|\mathsf p|^3\lg|\mathsf p| .
\]
The term with $i=1$ vanishes because $N_1$ is odd, and the terms with $i=0$ and $i=2$ are finite parts by (b).

\emph{The case $N=SJ_4$.} By the Leibniz rule and the bounds on the derivatives of $S$ and of $J_4$,
\[
|\partial_{\mathsf p}^\alpha N(k,\mathsf q)|\le C|k|^{2-|\alpha|}\qquad\bigl(|\alpha|\le3,\ |\mathsf q|\le\tfrac12|k|,\ |\mathsf q|\le\tfrac12\mu_0,\ |k|\le2\mu_0\bigr).
\]
Hence $|N_i|\le C|\mathsf p|^i|k|^{2-i}$, $|N_{\mathrm r}|\le C|\mathsf p|^3|k|^{-1}$ on $\mathsf B^c\cap\operatorname{supp}\eta$, and $|N|\le C|\mathsf p|^4(|k|^{-2}+|k-\mathsf p|^{-2})$ on $\mathsf B$. Arguing as in the previous case, we find $\int\eta N\,dk=\int\eta N_0\,dk+\int\eta N_2\,dk+O(|\mathsf p|^3)$, and the two integrals on the right are finite parts by (b).

Adding the three cases, the first one multiplied by $J(0,0)$, and using $|\mathsf p|^3\le|\mathsf p|^3\lg|\mathsf p|$, we obtain the estimate of (d), with a constant depending on $\mu_0$, $S$ and $J$.

\emph{Proof of \textup{(e)}.} We split $1=\bigl(1-\eta_c(k)-\eta_c(l)\bigr)+\eta_c(k)+\eta_c(l)$ and treat the three resulting integrals separately; each of them converges absolutely for large $\operatorname{Re}\mathsf z$, as the following discussion shows.

On the support of $\Phi$, $|k|<\frac c8$ implies $|l|>\frac78c$ by \textup{($\Phi0$)}, hence $\eta_c(k)=1$ and $\eta_c(l)=0$; the same holds with $k$ and $l$ exchanged. Therefore the first integrand, which contains the factor $1-\eta_c(k)-\eta_c(l)$, vanishes unless $|k|,|l|\ge\frac c8$, and by \textup{($\Phi0$)} it vanishes unless $|k|+|l|\le C_\Phi$. On this bounded set the factor $|k|^{2\epsilon\mathsf z}|l|^{2\epsilon'\mathsf z}$ is an entire function of $\mathsf z$, bounded for $\mathsf z$ in compact sets, while $|\Phi|\le C_\Phi|k-l|^{-2}$ by \textup{($\Phi3$)}. Hence the first integral is an entire function of $\mathsf z$, whose value at $\mathsf z=0$ is $\iint\bigl(1-\eta_c(k)-\eta_c(l)\bigr)\Phi\,dk\,dl$.

In the second integral, $\eta_c(k)\ne0$ forces $|k|\le\frac c4<\frac c3$, so that \textup{($\Phi1$)} applies, and Fubini's theorem turns this integral into
\[
\sum_\iota\int|k|^{2\epsilon\mathsf z}\,\mathsf h_\iota(k)\,\mathsf g_\iota(k;\mathsf z)\,dk,\qquad\mathsf g_\iota(k;\mathsf z):=\eta_c(k)\int|l|^{2\epsilon'\mathsf z}\,\Phi_\iota(k,l)\,dl .
\]
Since $\Phi_\iota$ vanishes where $|k|+|l|<c$ or $|k|+|l|>C_\Phi$, the integrand defining $\mathsf g_\iota$ vanishes unless $\frac34c\le|l|\le C_\Phi$; consequently $\mathsf g_\iota$ has the properties required in (a). If $\epsilon=1$, assertion (a) provides a holomorphic continuation to $\mathbb D$ whose value at $\mathsf z=0$ is $\operatorname{f.p.}\!\int\mathsf h_\iota\,\mathsf g_\iota(\cdot;0)\,dk$. If $\epsilon=0$, all degrees of the $\mathsf h_\iota$ are $\ge-2$; the $k$-integral then converges absolutely for every $\mathsf z$ and defines an entire function, whose value at $0$ is again the finite part, by (b). Summing over $\iota$, we obtain at $\mathsf z=0$ the number $\operatorname{f.p.}\!\int\bigl[\eta_c(k)\int\Phi(k,l)\,dl\bigr]dk$.

The hypotheses of (e) are symmetric under the exchange of $k$ with $l$, of $\epsilon$ with $\epsilon'$ and of \textup{($\Phi1$)} with \textup{($\Phi2$)}: the set $\{(1,1),(1,0),(0,1)\}$ is invariant under $(\epsilon,\epsilon')\mapsto(\epsilon',\epsilon)$, conditions \textup{($\Phi0$)} and \textup{($\Phi3$)} do not change when $k$ and $l$ are interchanged, and the assumptions on the degrees in \textup{($\Phi1$)} and \textup{($\Phi2$)} are interchanged. Under this exchange the third integral, which contains the factor $\eta_c(l)$, becomes the second one, and the argument just given applies to it. In detail: $\eta_c(l)\ne0$ forces $|l|\le\frac c4<\frac c3$, so that \textup{($\Phi2$)} applies, and Fubini's theorem turns the third integral into
\[
\sum_\iota\int|l|^{2\epsilon'\mathsf z}\,\mathsf h'_\iota(l)\,\mathsf g'_\iota(l;\mathsf z)\,dl,\qquad\mathsf g'_\iota(l;\mathsf z):=\eta_c(l)\int|k|^{2\epsilon\mathsf z}\,\Phi'_\iota(k,l)\,dk .
\]
Since $\Phi'_\iota$ vanishes where $|k|+|l|<c$ or $|k|+|l|>C_\Phi$, the integrand defining $\mathsf g'_\iota$ vanishes unless $\frac34c\le|k|\le C_\Phi$, and $\mathsf g'_\iota$ has the properties required in (a). If $\epsilon'=1$, assertion (a) provides a holomorphic continuation to $\mathbb D$ whose value at $\mathsf z=0$ is $\operatorname{f.p.}\!\int\mathsf h'_\iota\,\mathsf g'_\iota(\cdot;0)\,dl$. If $\epsilon'=0$, all degrees of the $\mathsf h'_\iota$ are $\ge-2$; the $l$-integral then converges absolutely for every $\mathsf z$ and defines an entire function, whose value at $0$ is the finite part, by (b). Summing over $\iota$, we obtain at $\mathsf z=0$ the number $\operatorname{f.p.}\!\int\bigl[\eta_c(l)\int\Phi(k,l)\,dk\bigr]dl$. The sum of the three values at $\mathsf z=0$ is the right side of the definition of $\operatorname{f.p.}\!\iint\Phi\,dk\,dl$ in Definition~\ref{IV-73-85}(2).

\emph{Change of the parameter $c$.} Let $0<c'<c$ and $\Phi\in\mathfrak P_2(c)$. Conditions \textup{($\Phi0$)}, \textup{($\Phi1$)} and \textup{($\Phi2$)} at $c'$ follow from the same conditions at $c$ by restriction. For \textup{($\Phi3$)} at $c'$, let $|k|,|l|\ge\frac{c'}8$ and $\Phi(k,l)\ne0$. If $|k|,|l|\ge\frac c8$, \textup{($\Phi3$)} at $c$ applies. Both $|k|$ and $|l|$ smaller than $\frac c8$ is impossible by \textup{($\Phi0$)}. If $\frac{c'}8\le|k|<\frac c8$, we use $\Phi=\sum_\iota\mathsf h_\iota\Phi_\iota$ from \textup{($\Phi1$)}: the $\mathsf h_\iota$ are bounded on $\{\frac{c'}8\le|k|\le\frac c8\}$, the $\Phi_\iota$ are bounded there as well (they are continuous on $\{|k|\le\frac c8\}\times\mathbb R^3$ and vanish for $|l|>C_\Phi$), and $|k-l|\le|k|+|l|\le C_\Phi$; hence $|\Phi|\le C\le C\,C_\Phi^2|k-l|^{-2}$.

The case $\frac{c'}8\le|l|<\frac c8$ is symmetric, with \textup{($\Phi2$)}. Thus $\mathfrak P_2(c)\subset\mathfrak P_2(c')$. The integral with $(\epsilon,\epsilon')=(1,1)$ does not involve $c$, so the value at $\mathsf z=0$ of its continuation, which is the finite part with respect to $c$ and also the finite part with respect to $c'$, is the same.

Finally, let $\Phi$ be integrable. By Fubini's theorem the functions $k\mapsto\eta_c(k)\int\Phi(k,l)\,dl$ and $l\mapsto\eta_c(l)\int\Phi(k,l)\,dk$ are integrable elements of $\mathfrak P$, so by (b) their finite parts are their integrals, and the definition of $\operatorname{f.p.}\!\iint\Phi$ reduces to $\iint\Phi\,dk\,dl$.
\end{proof}

\begin{proposition}[Expansion by regions of $\Xi_\Gamma{[\Upsilon]}$ at small $\mathsf p$]\label{IV-73-87}
Let $\Gamma\in\{1,\dots,6\}$, $c>0$ and $\Upsilon\in\mathcal Y(c)$. Then $\overline{(\Upsilon F_\Gamma)_2}\in\mathfrak P_2(c)$, and there are numbers $\alpha$, $\beta$ and $C$, depending only on $\Gamma$, $c$ and $\Upsilon$, such that
\[
\bigl|\Xi_\Gamma[\Upsilon](\mathsf p)-\alpha-\beta\,|\mathsf p|-\mathfrak f_\Gamma[\Upsilon]\,\mathsf p^2\bigr|\le C\,|\mathsf p|^3\lg|\mathsf p|\qquad\bigl(0<|\mathsf p|\le\tfrac c{16}\bigr),
\]
where
\[
\mathfrak f_\Gamma[\Upsilon]:=(2\pi)^{-6}\operatorname{f.p.}\!\iint\overline{(\Upsilon F_\Gamma)_2}\,dk\,dl .
\]
The map $\Upsilon\mapsto\mathfrak f_\Gamma[\Upsilon]$ is linear on $\mathcal Y(c)$, and $\mathfrak f_\Gamma[\Upsilon^\Lambda]=\mathfrak f_\Gamma[\Upsilon]$ for every $\Lambda>1$.
\end{proposition}

\begin{proof}
\emph{Step 1: two factorisations of $F_\Gamma$.} We write the integrand $F_\Gamma$ of \eqref{IV-73-65} in two ways,
\[
F_\Gamma=\sum S\cdot H=\sum S'\cdot H' ,
\]
where the sums extend over the tensor indices carried by $S$, resp.\ by $S'$, the factor $S$ depends on $(k,\mathsf p)$ and $S'$ on $(l,\mathsf p)$ as in the following table, and $H$, resp.\ $H'$, denotes the product of all the other factors of $F_\Gamma$:
\[
\begin{array}{c|c|c}
\Gamma&S(k,\mathsf p)&S'(l,\mathsf p)\\\hline
1&V_2(k,\mathsf p-k)_{\mu\nu}/\mathsf x_1\mathsf x_1'&V_2(-l,l-\mathsf p)_{\rho\sigma}/\mathsf x_2\mathsf x_2'\\
2,\ 3&V_2(k,\mathsf p-k)_{\mu\nu}V_2(-k,k-\mathsf p)_{\rho\nu}/\mathsf x_1^2\mathsf x_1'&1/\mathsf x_2\\
4&V_2(k,\mathsf p-k)_{\mu\nu}/\mathsf x_1\mathsf x_1'&V_2(l,\mathsf p-l)_{\mu\nu}/\mathsf x_2\mathsf x_2'\\
5&V_2(k-\mathsf p,-k)_{\alpha\sigma}/\mathsf x_1\mathsf x_1'&1/\mathsf x_2\\
6&1/\mathsf x_1'&1/\mathsf x_2
\end{array}
\]
For $\Gamma=4$ this means
\[
H_{\mu\nu}=-\tfrac12\bigl[\delta_{\mu\nu}V_2(l,\mathsf p-l)_{\lambda\lambda}-\tfrac12V_2(l,\mathsf p-l)_{\mu\nu}-\tfrac12V_2(l,\mathsf p-l)_{\nu\mu}\bigr]/\mathsf x_2\mathsf x_2',
\]
and $H'$ has the same form, with $V_2(k,\mathsf p-k)$ and $\mathsf x_1\mathsf x_1'$ in place of $V_2(l,\mathsf p-l)$ and $\mathsf x_2\mathsf x_2'$. Every component of $S$ has the form required in Lemma~\ref{IV-73-86}(d): it is a homogeneous polynomial divided by $|k|^{2\mathsf a}|k-\mathsf p|^{2\mathsf a'}$, with $(\mathsf a,\mathsf a')=(1,1)$ for $\Gamma\in\{1,4,5\}$, $(2,1)$ for $\Gamma\in\{2,3\}$ and $(0,1)$ for $\Gamma=6$, and since $|V_2(k_1,k_2)|\le2|k_1||k_2|$ its modulus is at most $2/|k||k-\mathsf p|\le|k|^{-2}+|k-\mathsf p|^{-2}$ for $\Gamma\in\{1,4,5\}$, at most $4|k|^{-2}$ for $\Gamma\in\{2,3\}$, and equal to $|k-\mathsf p|^{-2}$ for $\Gamma=6$.

The same holds for the components of $S'$, as functions of $l$. The factor $S$ contains all factors $\mathsf x_1,\mathsf x_1'$ of the denominator of $F_\Gamma$, and $S'$ all factors $\mathsf x_2,\mathsf x_2'$. Hence $H$ is a polynomial divided by powers of $\mathsf x_2,\mathsf x_2',\mathsf x_3$, and $H'$ a polynomial divided by powers of $\mathsf x_1,\mathsf x_1',\mathsf x_3$; in both, $\mathsf x_3$ occurs at most to the first power. Moreover, $H$ and $H'$ are even under the reflection of all their arguments: they are quotients of homogeneous polynomials of even degree, because besides factors $V_2$ each of them contains two factors $V_3$, or the factor $l_\rho(l-k)_\mu$, or neither.

\emph{Step 2: where the lines are long.} Put $\mu_0:=\frac c8$ and $\eta:=\eta_c=\theta(|\cdot|/2\mu_0)$. Let $|\mathsf p|<\mu_0$. By \textup{(Y1)}, $\Upsilon(k,l,\mathsf p)\ne0$ implies $|k|+|l|\ge8\mu_0$. If, in addition, $|k|<3\mu_0$, then
\[
|l|>5\mu_0,\qquad|l-\mathsf p|>4\mu_0,\qquad|k-l|>2\mu_0 ;
\]
these are the lengths whose squares $\mathsf x_2,\mathsf x_2',\mathsf x_3$ occur in the denominator of $H$. Consequently the product $\Upsilon H$, extended by $0$ where $\Upsilon=0$, is smooth on $\{|k|<3\mu_0\}\times\mathbb R^3\times\{|\mathsf p|<\mu_0\}$: this set is covered by the open set $\{|k|+|l|<8\mu_0,\ |\mathsf p|<c\}$, on which $\Upsilon$ vanishes identically, and by the open set $\{|k|<3\mu_0,\ |l|>5\mu_0,\ |\mathsf p|<\mu_0\}$, on which $H$ is smooth. Symmetrically, if $\Upsilon(k,l,\mathsf p)\ne0$ and $|l|<3\mu_0$, then $|k|\ge8\mu_0-|l|$, and
\[
|k|>5\mu_0,\qquad|k-\mathsf p|>4\mu_0,\qquad|k-l|>2\mu_0 ;
\]
these are the lengths whose squares $\mathsf x_1,\mathsf x_1',\mathsf x_3$ occur in the denominator of $H'$. The set $\mathbb R^3\times\{|l|<3\mu_0\}\times\{|\mathsf p|<\mu_0\}$ is covered by the open set $\{|k|+|l|<8\mu_0,\ |\mathsf p|<c\}$, on which $\Upsilon$ vanishes identically, and by the open set $\{|k|>5\mu_0,\ |l|<3\mu_0,\ |\mathsf p|<\mu_0\}$, on which $H'$ is smooth; hence $\Upsilon H'$, extended by $0$ where $\Upsilon=0$, is smooth on it. Finally, if $|k|,|l|\ge\mu_0$ and $|\mathsf p|\le\frac12\mu_0$, then $\mathsf x_1,\mathsf x_1',\mathsf x_2,\mathsf x_2'\ge\frac14\mu_0^2$, and therefore $F_\Gamma=F^\sharp/\mathsf x_3$ with a function $F^\sharp$ which is smooth on this set.

\emph{Step 3: $\Phi:=\overline{(\Upsilon F_\Gamma)_2}$ belongs to $\mathfrak P_2(c)$.} Condition \textup{($\Phi0$)} follows from \textup{(Y1)}: where $|k|+|l|<c$ or $|k|+|l|>C_\Upsilon$, the function $\Upsilon F_\Gamma$ vanishes for all $|\mathsf p|<c$, and so does $\Phi$. For \textup{($\Phi1$)} let $0<|k|<\frac c3$; then $|k|<3\mu_0$, and
\[
\Phi=\frac16\sum_j\sum\Bigl[\partial_{\mathsf p_j}^2S\cdot\Upsilon H+2\,\partial_{\mathsf p_j}S\cdot\partial_{\mathsf p_j}(\Upsilon H)+S\cdot\partial_{\mathsf p_j}^2(\Upsilon H)\Bigr]_{\mathsf p=0}
\]
by the Leibniz rule. Here $\partial_{\mathsf p}^\beta S(\cdot,0)$ is a homogeneous factor of degree $-2-|\beta|$, and $\partial_{\mathsf p}^\beta(\Upsilon H)(k,l,0)$ is smooth by Step~2. The latter function vanishes wherever $\Upsilon(\cdot,\cdot,\mathsf p)=0$ for all small $\mathsf p$, that is, by \textup{(Y1)}, where $|k|+|l|<c$ or $|k|+|l|>C_\Upsilon$; this is the support condition of \textup{($\Phi1$)}, with any $C_\Phi\ge C_\Upsilon$. Condition \textup{($\Phi2$)} is obtained from the factorisation $F_\Gamma=\sum S'\cdot H'$ by the same argument with the roles of $k$ and $l$ exchanged: for $0<|l|<\frac c3$ the Leibniz rule writes $\Phi$ as the corresponding sum of products of the functions $\partial_{\mathsf p}^\beta S'(\cdot,0)$, which are homogeneous factors of degree $-2-|\beta|$ in $l$, with the functions $\partial_{\mathsf p}^\beta(\Upsilon H')(k,l,0)$, which are smooth on $\mathbb R^3\times\{|l|<3\mu_0\}$ by Step~2 and vanish where $|k|+|l|<c$ or $|k|+|l|>C_\Upsilon$.

For \textup{($\Phi3$)} let $|k|,|l|\ge\frac c8=\mu_0$; by Step~2,
\[
\Phi=\tfrac16\,\Delta_{\mathsf p}\bigl(\Upsilon F^\sharp\bigr)\big|_{\mathsf p=0}\big/\mathsf x_3 ,
\]
and the numerator is bounded, being smooth with compact support in $(k,l)$.

\emph{Step 4: three regions.} For $0<|\mathsf p|\le\frac12\mu_0$ we split $\Xi_\Gamma[\Upsilon]=\Xi_{\mathrm H}+\Xi_{\mathrm A}+\Xi_{\mathrm B}$, inserting into the integral the factors $1-\eta(k)-\eta(l)$, $\eta(k)$ and $\eta(l)$, respectively. If $\Upsilon\ne0$ and $|k|<\mu_0$, then $|l|>7\mu_0$, hence $\eta(k)=1$ and $\eta(l)=0$; thus on the support of $\Upsilon$ the first factor vanishes unless $|k|,|l|\ge\mu_0$, and the same holds with $k$ and $l$ exchanged. Each of the three terms is invariant under rotations of $\mathsf p$, by \textup{(Y2)}, by the rotation invariance of $F_\Gamma$, and because $\eta$ is radial.

\emph{Step 5: the hard region.} The integrand of $\Xi_{\mathrm H}$ is $(1-\eta(k)-\eta(l))\,\Upsilon F^\sharp/\mathsf x_3$; it is supported in $\{\mu_0\le|k|,|l|\le C_\Upsilon\}$, and all its derivatives with respect to $\mathsf p$ are bounded by $C|k-l|^{-2}$. Hence $\Xi_{\mathrm H}$ is smooth on $\{|\mathsf p|<\frac12\mu_0\}$; it is even, by \textup{(Y2)} and the evenness of $F_\Gamma$, and rotation invariant. At $\mathsf p=0$ its gradient and its third derivatives therefore vanish, and its Hessian is a multiple of the identity, so that
\[
\Xi_{\mathrm H}(\mathsf p)=\Xi_{\mathrm H}(0)+\tfrac16\,\Delta\Xi_{\mathrm H}(0)\,\mathsf p^2+O(|\mathsf p|^4),\qquad\tfrac16\,\Delta\Xi_{\mathrm H}(0)=(2\pi)^{-6}\iint\bigl(1-\eta(k)-\eta(l)\bigr)\,\Phi\,dk\,dl .
\]

\emph{Step 6: $k$ soft.} Let $J(k,\mathsf p):=\int\Upsilon H\,d^\circ l$, one function for each value of the tensor indices of $S$. By Step~2, and since $\Upsilon H$ vanishes for $|l|>C_\Upsilon$, the function $J$ is smooth on $\{|k|<3\mu_0\}\times\{|\mathsf p|<\mu_0\}$, and it is even, as the substitution $l\to-l$ shows. By Fubini's theorem $\Xi_{\mathrm A}=\sum\int\eta SJ\,d^\circ k$. We apply Lemma~\ref{IV-73-86}(d), with the present $\mu_0$, to each component of $S$ and the corresponding component of $J$; for $0<|\mathsf p|\le\frac12\mu_0$ this gives
\[
\begin{multlined}
(2\pi)^3\,\Xi_{\mathrm A}(\mathsf p)=\sum\Bigl\{\operatorname{f.p.}\!\int\eta\,[SJ]_0\,dk+\sum_{i,j}\mathsf p_i\mathsf p_j\,\tfrac12\operatorname{f.p.}\!\int\eta\,\partial_{\mathsf p_i}\partial_{\mathsf p_j}(SJ)(k,0)\,dk+J(0,0)\,\mathsf h_S(\mathsf p)\Bigr\} \\
+O(|\mathsf p|^3\lg|\mathsf p|).
\end{multlined}
\]
We take the mean of both sides over the sphere $|\mathsf p|=t$. The left side is constant there, by Step~4. The mean of $\mathsf h_S$ is $t$ times its mean over the unit sphere, by homogeneity, and the mean of $\mathsf p_i\mathsf p_j$ is $\frac13t^2\delta_{ij}$. For $k\ne0$, differentiation under the integral sign gives $\sum\frac16\Delta_{\mathsf p}(SJ)(k,0)=\int\Phi(k,l)\,d^\circ l$. Since the finite part is linear, we conclude that
\[
\Xi_{\mathrm A}(\mathsf p)=\alpha_{\mathrm A}+\beta_{\mathrm A}|\mathsf p|+(2\pi)^{-6}\,\mathsf p^2\operatorname{f.p.}\!\int\Bigl[\eta(k)\int\Phi(k,l)\,dl\Bigr]dk+O(|\mathsf p|^3\lg|\mathsf p|)
\]
with numbers $\alpha_{\mathrm A}$, $\beta_{\mathrm A}$.

\emph{Step 7: $l$ soft, and the sum.} For $\Xi_{\mathrm B}$ we exchange the roles of $k$ and $l$ in Step~6. We use the factorisation $F_\Gamma=\sum S'\cdot H'$ of Step~1 and put $J'(l,\mathsf p):=\int\Upsilon H'\,d^\circ k$, one function for each value of the tensor indices of $S'$. By Step~2, and since $\Upsilon H'$ vanishes for $|k|>C_\Upsilon$, the function $J'$ is smooth on $\{|l|<3\mu_0\}\times\{|\mathsf p|<\mu_0\}$, and it is even, as the substitution $k\to-k$ shows; by Fubini's theorem $\Xi_{\mathrm B}=\sum\int\eta S'J'\,d^\circ l$. We apply Lemma~\ref{IV-73-86}(d), with $l$ in place of $k$, to each component of $S'$ and the corresponding component of $J'$, and we take the mean over the sphere $|\mathsf p|=t$ as in Step~6, now with $\sum\frac16\Delta_{\mathsf p}(S'J')(l,0)=\int\Phi(k,l)\,d^\circ k$ for $l\ne0$. This gives
\[
\Xi_{\mathrm B}(\mathsf p)=\alpha_{\mathrm B}+\beta_{\mathrm B}|\mathsf p|+(2\pi)^{-6}\,\mathsf p^2\operatorname{f.p.}\!\int\Bigl[\eta(l)\int\Phi(k,l)\,dk\Bigr]dl+O(|\mathsf p|^3\lg|\mathsf p|)\qquad\bigl(0<|\mathsf p|\le\tfrac12\mu_0\bigr)
\]
with numbers $\alpha_{\mathrm B}$, $\beta_{\mathrm B}$. Since $\eta=\eta_c$ and $\Phi\in\mathfrak P_2(c)$ by Step~3, Definition~\ref{IV-73-85}(2) shows that the sum of the three coefficients of $\mathsf p^2$ obtained in Steps~5--7 is $\mathfrak f_\Gamma[\Upsilon]$. Adding the three expansions, and noting that $\frac12\mu_0=\frac c{16}$, we obtain the asserted estimate, with $\alpha:=\Xi_{\mathrm H}(0)+\alpha_{\mathrm A}+\alpha_{\mathrm B}$ and $\beta:=\beta_{\mathrm A}+\beta_{\mathrm B}$.

\emph{Step 8: linearity and scaling.} The map $\Upsilon\mapsto\overline{(\Upsilon F_\Gamma)_2}$ and the finite part are linear; hence so is $\mathfrak f_\Gamma$. Let $\Lambda>1$. The substitution $k\to\Lambda k$, $l\to\Lambda l$, together with the homogeneity of degree $-4$ of $F_\Gamma$, gives $\Xi_\Gamma[\Upsilon^\Lambda](\mathsf p)=\Lambda^2\,\Xi_\Gamma[\Upsilon](\mathsf p/\Lambda)$, so that by the estimate already proved
\[
\Xi_\Gamma[\Upsilon^\Lambda](\mathsf p)=\Lambda^2\alpha+\Lambda\beta\,|\mathsf p|+\mathfrak f_\Gamma[\Upsilon]\,\mathsf p^2+O(|\mathsf p|^3\lg|\mathsf p|).
\]
On the other hand $\Upsilon^\Lambda\in\mathcal Y(c)$, and the estimate applied to $\Upsilon^\Lambda$ gives $\Xi_\Gamma[\Upsilon^\Lambda](\mathsf p)=\alpha'+\beta'|\mathsf p|+\mathfrak f_\Gamma[\Upsilon^\Lambda]\,\mathsf p^2+O(|\mathsf p|^3\lg|\mathsf p|)$ with numbers $\alpha',\beta'$. If $\mathsf a+\mathsf bt+\mathsf ct^2=o(t^2)$ as $t\downarrow0$, then $\mathsf a=\mathsf b=\mathsf c=0$; comparing the two expansions, we obtain in particular $\mathfrak f_\Gamma[\Upsilon^\Lambda]=\mathfrak f_\Gamma[\Upsilon]$.
\end{proof}

\subsection{The value: the functional $\mathcal R$, the residue, the bridge}

By Proposition~\ref{IV-73-87}, the coefficient of $\mathsf p^2$ in the expansion of $\Xi_\Gamma[\Upsilon]$ at small $\mathsf p$ is the number $\mathfrak f_\Gamma[\Upsilon]$, which is given by a finite part. We now compute it for the differences $R_\Gamma^\Lambda-R_\Gamma$ of the cut-offs of Lemma~\ref{IV-73-84}(b). The result is $d_{2,\Gamma}\log\Lambda$ with $d_{2,\Gamma}=\mathcal R[F_\Gamma]/16\pi^2$, where $\mathcal R$ is a linear functional defined below; the comparison with Lemma~\ref{IV-73-84} then expresses $\kappa_\Gamma^{\mathrm{lat}}$ through $\mathcal R[F_\Gamma]$.

\begin{definition}[The functional $\mathcal R$]\label{IV-73-88}
We use the variables $\mathsf x_0,\mathsf x_1,\mathsf x_1',\mathsf x_2,\mathsf x_2',\mathsf x_3$ of Definition~\ref{IV-73-64}(1).
\begin{enumerate}
\item[(1)] The scalar products of $k$, $l$, $\mathsf p$ are linear in these variables:
\[
k\cdot l=\tfrac12(\mathsf x_1+\mathsf x_2-\mathsf x_3),\qquad k\cdot\mathsf p=\tfrac12(\mathsf x_1-\mathsf x_1'+\mathsf x_0),\qquad l\cdot\mathsf p=\tfrac12(\mathsf x_2-\mathsf x_2'+\mathsf x_0).
\]
\item[(2)] For $n\in\mathbb Z$ let $\mathcal L_n$ be the real linear span of the functions $\mathsf x_1^{-\mathsf a}\mathsf x_2^{-\mathsf b}\mathsf x_3^{-\mathsf c}$, where $\mathsf a,\mathsf b,\mathsf c\in\mathbb Z$ and $\mathsf a+\mathsf b+\mathsf c=-n$, on the set $\{(k,l):k\ne0,\ l\ne0,\ k\ne l\}$. These functions are linearly independent. Indeed, the image of the map $(k,l)\mapsto(\mathsf x_1,\mathsf x_2,\mathsf x_3)$ contains a non-empty open set: for $k=(t_1,t_2,0)$ and $l=(0,t_3,0)$ the map becomes
\[
(t_1,t_2,t_3)\longmapsto\bigl(t_1^2+t_2^2,\ t_3^2,\ t_1^2+(t_2-t_3)^2\bigr),
\]
whose differential at $(1,0,1)$ has determinant $8$, so that the inverse function theorem applies. A vanishing finite linear combination of the functions, multiplied by a sufficiently high power of $\mathsf x_1\mathsf x_2\mathsf x_3$, becomes a polynomial in $(\mathsf x_1,\mathsf x_2,\mathsf x_3)$ which vanishes on a non-empty open set, and such a polynomial is zero. We denote by $r_{111}\colon\mathcal L_{-3}\to\mathbb R$ the linear map which assigns to an element of $\mathcal L_{-3}$ its coefficient of $(\mathsf x_1\mathsf x_2\mathsf x_3)^{-1}$.
\item[(3)] Let $\mathfrak F$ be the real linear span of the monomials
\[
\mathsf x_0^n\,\mathsf x_1^{-a_1}\mathsf x_1'^{-a_2}\mathsf x_2^{-a_3}\mathsf x_2'^{-a_4}\mathsf x_3^{-a_5},\qquad\text{integers }n\ge0,\ a_1,\dots,a_5,\quad\textstyle\sum_ia_i=2+n ;
\]
these monomials are homogeneous of degree $-4$. The numerators of the integrands in \eqref{IV-73-65} are complete contractions of tensors built from $k$, $l$, $\mathsf p$ and Kronecker symbols; hence they are polynomials in the scalar products of $k$, $l$, $\mathsf p$, and by (1) polynomials in the six variables. Since moreover $\mathsf x_3$ occurs in the denominators of \eqref{IV-73-65} at most to the first power, and not at all in $F_4$, each $F_\Gamma$ belongs to $\mathfrak F$: it is a linear combination of monomials with $a_5\le1$, and $F_4$ is a linear combination of monomials with $a_5\le0$.
\item[(4)] For $F\in\mathfrak F$ we put
\begin{equation}\label{IV-73-89}
\mathcal R[F]:=r_{111}\bigl(\overline{F_2}\bigr),
\end{equation}
with $\overline{F_2}$ as in Definition~\ref{IV-73-85}(3); by Lemma~\ref{IV-73-90} below, $\overline{F_2}\in\mathcal L_{-3}$.
\end{enumerate}
\end{definition}

\begin{lemma}[Values of $\mathcal R$ on monomials]\label{IV-73-90}
For every $F\in\mathfrak F$ the function $\overline{F_2}$ belongs to $\mathcal L_{-3}$; hence $\mathcal R$ is well defined on $\mathfrak F$, and it is linear. If $F$ is a linear combination of monomials with $a_5\le1$, then $\overline{F_2}$ is a linear combination of functions $\mathsf x_1^{-\mathsf a}\mathsf x_2^{-\mathsf b}\mathsf x_3^{-\mathsf c}$ with $\mathsf c\le1$. For integers $n\ge0$ and $a_1,\dots,a_5$ with $\sum_ia_i=2+n$, put
\[
(\alpha_1,\alpha_2,\alpha_3):=(a_1+a_2,\ a_3+a_4,\ a_5),\qquad\mathsf q(t):=\tfrac13t(2t-1).
\]
Then
\begin{equation}\label{IV-73-91}
\mathcal R\bigl[\mathsf x_0^n\,\mathsf x_1^{-a_1}\mathsf x_1'^{-a_2}\mathsf x_2^{-a_3}\mathsf x_2'^{-a_4}\mathsf x_3^{-a_5}\bigr]=
\begin{cases}
1&n=1,\ (\alpha_1,\alpha_2,\alpha_3)=(1,1,1),\\
\mathsf q(a_2)+\frac23a_2a_4&n=0,\ (\alpha_1,\alpha_2,\alpha_3)=(0,1,1),\\
\mathsf q(a_4)+\frac23a_2a_4&n=0,\ (\alpha_1,\alpha_2,\alpha_3)=(1,0,1),\\
-\frac23a_2a_4&n=0,\ (\alpha_1,\alpha_2,\alpha_3)=(0,0,2),\\
0&\text{otherwise.}
\end{cases}
\end{equation}
\end{lemma}

\begin{proof}
The map $F\mapsto\overline{F_2}$ is linear; it therefore suffices to consider a monomial $F$ as on the left of \eqref{IV-73-91}. We write $x:=k\cdot\mathsf p$ and $y:=l\cdot\mathsf p$; then $\mathsf x_1'=\mathsf x_1-2x+\mathsf x_0$ and $\mathsf x_2'=\mathsf x_2-2y+\mathsf x_0$. For an integer $t$, the binomial series $(1+w)^{-t}=1-tw+\frac12t(t+1)w^2+O(w^3)$, taken at $w=(\mathsf x_0-2x)/\mathsf x_1$, gives for fixed $(k,l)$ and small $\mathsf p$
\[
\mathsf x_1'^{-t}=\mathsf x_1^{-t}\Bigl[1+2t\,\frac x{\mathsf x_1}-t\,\frac{\mathsf x_0}{\mathsf x_1}+2t(t+1)\,\frac{x^2}{\mathsf x_1^2}+O(|\mathsf p|^3)\Bigr],
\]
and $\mathsf x_2'^{-t}$ has the same expansion, with $y$ and $\mathsf x_2$ in place of $x$ and $\mathsf x_1$. We use these expansions at $t=a_2$ and at $t=a_4$. On the terms of second order the operator $\frac16\Delta_{\mathsf p}$ acts by
\[
\tfrac16\Delta_{\mathsf p}\mathsf x_0=1,\qquad\tfrac16\Delta_{\mathsf p}x^2=\tfrac13\mathsf x_1,\qquad\tfrac16\Delta_{\mathsf p}y^2=\tfrac13\mathsf x_2,\qquad\tfrac16\Delta_{\mathsf p}(xy)=\tfrac13\,k\cdot l=\tfrac16(\mathsf x_1+\mathsf x_2-\mathsf x_3),
\]
the last equality by Definition~\ref{IV-73-88}(1).

If $n\ge2$, then $F=O(|\mathsf p|^4)$ and $\overline{F_2}=0$. If $n=1$, the term of second order in $\mathsf p$ of $F$ is $\mathsf x_0\,\mathsf x_1^{-\alpha_1}\mathsf x_2^{-\alpha_2}\mathsf x_3^{-\alpha_3}$, and $\overline{F_2}=\mathsf x_1^{-\alpha_1}\mathsf x_2^{-\alpha_2}\mathsf x_3^{-\alpha_3}$. If $n=0$, the terms of second order are
\[
\mathsf x_1^{-\alpha_1}\mathsf x_2^{-\alpha_2}\mathsf x_3^{-\alpha_3}\Bigl\{\Bigl[-a_2\frac{\mathsf x_0}{\mathsf x_1}+2a_2(a_2+1)\frac{x^2}{\mathsf x_1^2}\Bigr]+\Bigl[-a_4\frac{\mathsf x_0}{\mathsf x_2}+2a_4(a_4+1)\frac{y^2}{\mathsf x_2^2}\Bigr]+4a_2a_4\,\frac{xy}{\mathsf x_1\mathsf x_2}\Bigr\},
\]
and since $-t+\frac23t(t+1)=\mathsf q(t)$, the rules above give
\[
\overline{F_2}=\mathsf x_1^{-\alpha_1}\mathsf x_2^{-\alpha_2}\mathsf x_3^{-\alpha_3}\Bigl\{\mathsf q(a_2)\,\mathsf x_1^{-1}+\mathsf q(a_4)\,\mathsf x_2^{-1}+\tfrac23a_2a_4\bigl[\mathsf x_2^{-1}+\mathsf x_1^{-1}-\mathsf x_3\mathsf x_1^{-1}\mathsf x_2^{-1}\bigr]\Bigr\}.
\]
In each case $\overline{F_2}\in\mathcal L_{-3}$, because $\alpha_1+\alpha_2+\alpha_3=2+n$; and the exponent $\mathsf c$ of $\mathsf x_3$ in the terms obtained is $a_5$ or $a_5-1$, which proves the assertion on $\mathsf c\le1$. Finally, the monomial $(\mathsf x_1\mathsf x_2\mathsf x_3)^{-1}$ occurs for $n=1$ exactly when $(\alpha_1,\alpha_2,\alpha_3)=(1,1,1)$, with coefficient $1$. For $n=0$ it arises from the terms containing $\mathsf x_1^{-1}$ exactly when $(\alpha_1,\alpha_2,\alpha_3)=(0,1,1)$, from the terms containing $\mathsf x_2^{-1}$ exactly when $(\alpha_1,\alpha_2,\alpha_3)=(1,0,1)$, and from the last term exactly when $(\alpha_1,\alpha_2,\alpha_3)=(0,0,2)$. Collecting the coefficients, we obtain \eqref{IV-73-91}.
\end{proof}

\begin{lemma}[The residue]\label{IV-73-92}
Let $\mathsf M=\mathsf x_1^{-\mathsf a}\mathsf x_2^{-\mathsf b}\mathsf x_3^{-\mathsf c}$ with integers $\mathsf a,\mathsf b,\mathsf c$ such that $\mathsf a+\mathsf b+\mathsf c=3$ and $\mathsf c\le1$. Let $\phi\in C_c^\infty(\mathbb R^3\times\mathbb R^3)$ be equal to $1$ near $0$, let $\Lambda>1$, and let $\phi^\Lambda$ be as in Definition~\ref{IV-73-85}(3). Then $(\phi^\Lambda-\phi)\mathsf M\in\mathfrak P_2(c)$ for all sufficiently small $c>0$, and
\begin{equation}\label{IV-73-93}
\begin{gathered}
\operatorname{f.p.}\!\iint(\phi^\Lambda-\phi)\,\mathsf M\,dk\,dl=\begin{cases}4\pi^4\log\Lambda&(\mathsf a,\mathsf b,\mathsf c)=(1,1,1),\\0&\text{otherwise;}\end{cases}\qquad\\
\oint_{S^5}\frac{d\omega}{\mathsf x_1\mathsf x_2\mathsf x_3}=4\pi^4,\qquad\\
\frac{4\pi^4}{(2\pi)^6}=\frac1{16\pi^2},
\end{gathered}
\end{equation}
where $d\omega$ denotes the surface measure of the sphere $S^5\subset\mathbb R^3\times\mathbb R^3$.
\end{lemma}

\begin{proof}
\emph{Step 0: the class.} Let $\Upsilon_0\in C_c^\infty(\mathbb R^6)$ vanish for $|k|+|l|<c$. We show that $\Upsilon_0\mathsf M\in\mathfrak P_2(c)$. Condition \textup{($\Phi0$)} holds, since $\Upsilon_0$ vanishes for $|k|+|l|<c$ and has compact support. If $|k|<\frac c3$ and $\Upsilon_0(k,l)\ne0$, then $|l|\ge c-|k|>\frac23c$ and $|k-l|\ge|l|-|k|>\frac c3$. Hence $\Upsilon_0\,\mathsf x_2^{-\mathsf b}\mathsf x_3^{-\mathsf c}$, extended by $0$ where $\Upsilon_0$ vanishes, is smooth on $\{|k|<\frac c3\}\times\mathbb R^3$ (it vanishes on the open set $\{|k|+|l|<c\}$, and it is smooth on the open set $\{|k|<\frac c3,\ |l|>\frac23c\}$), and it vanishes where $|k|+|l|<c$ or where $(k,l)$ lies outside the support of $\Upsilon_0$.

Thus \textup{($\Phi1$)} holds with the single homogeneous factor $|k|^{-2\mathsf a}$, of degree $-2\mathsf a$. Symmetrically, if $|l|<\frac c3$ and $\Upsilon_0(k,l)\ne0$, then $|k|\ge c-|l|>\frac23c$ and $|k-l|\ge|k|-|l|>\frac c3$. Hence $\Upsilon_0\,\mathsf x_1^{-\mathsf a}\mathsf x_3^{-\mathsf c}$, extended by $0$ where $\Upsilon_0$ vanishes, is smooth on $\mathbb R^3\times\{|l|<\frac c3\}$ (it vanishes on the open set $\{|k|+|l|<c\}$, and it is smooth on the open set $\{|k|>\frac23c,\ |l|<\frac c3\}$), and it vanishes where $|k|+|l|<c$ or where $(k,l)$ lies outside the support of $\Upsilon_0$. Thus \textup{($\Phi2$)} holds with the single homogeneous factor $|l|^{-2\mathsf b}$, of degree $-2\mathsf b$.

Condition \textup{($\Phi3$)} holds because $\mathsf c\le1$ and $\Upsilon_0$ has compact support. We apply this to $\Upsilon_0=\phi^\Lambda-\phi$, which vanishes near $0$, since $\phi$ and $\phi^\Lambda$ are both equal to $1$ there.

\emph{Step 1: independence of the cut-off.} Let $\Upsilon_0$ be as in Step~0. For large $\operatorname{Re}\mathsf z$ the substitution $(k,l)\to(\Lambda k,\Lambda l)$ and the homogeneity of degree $-6$ of $\mathsf M$ give
\[
\iint(|k||l|)^{2\mathsf z}\,\Upsilon_0^\Lambda\,\mathsf M\,dk\,dl=\Lambda^{4\mathsf z+6-6}\iint(|k||l|)^{2\mathsf z}\,\Upsilon_0\,\mathsf M\,dk\,dl .
\]
Since $\Upsilon_0^\Lambda$ also vanishes for $|k|+|l|<c$, Lemma~\ref{IV-73-86}(e), with $(\epsilon,\epsilon')=(1,1)$, continues both sides holomorphically to $\mathbb D$, and at $\mathsf z=0$ we obtain $\operatorname{f.p.}\!\iint\Upsilon_0^\Lambda\mathsf M=\operatorname{f.p.}\!\iint\Upsilon_0\mathsf M$. If $\tilde\phi$ is a second function with the properties of $\phi$, we take $\Upsilon_0:=\phi-\tilde\phi$; then
\[
(\phi^\Lambda-\phi)-(\tilde\phi^\Lambda-\tilde\phi)=\Upsilon_0^\Lambda-\Upsilon_0 ,
\]
and the finite part of $(\phi^\Lambda-\phi)\mathsf M$ does not depend on the choice of $\phi$. We may therefore choose $\phi$ conveniently in each case. In the following steps we use the function $\varphi(s):=\theta(s)$ of one variable $s\ge0$; it is equal to $1$ on $[0,\frac12]$ and to $0$ on $[1,\infty)$.

\emph{Step 2: the case $\mathsf c\le0$.} We choose $\phi(k,l):=\varphi(|k|)\varphi(|l|)$. Since $\mathsf x_3^{-\mathsf c}=(k^2-2k\cdot l+l^2)^{|\mathsf c|}$ is a polynomial, $\mathsf M$ is a finite sum
\[
\mathsf M=\sum c_{\alpha\beta}\,|k|^{-2\mathsf a_\alpha}k^\alpha\cdot|l|^{-2\mathsf b_\beta}l^\beta
\]
with numbers $c_{\alpha\beta}$, multi-indices $\alpha,\beta$ and integers $\mathsf a_\alpha,\mathsf b_\beta$. For large $\operatorname{Re}\mathsf z$ put
\[
\mathsf Z(\mathsf z):=\iint(|k||l|)^{2\mathsf z}\,\phi\,\mathsf M\,dk\,dl=\sum c_{\alpha\beta}\int|k|^{2\mathsf z-2\mathsf a_\alpha}k^\alpha\varphi(|k|)\,dk\int|l|^{2\mathsf z-2\mathsf b_\beta}l^\beta\varphi(|l|)\,dl .
\]
As $\mathsf M$ is homogeneous of degree $-6$, we have $\iint(|k||l|)^{2\mathsf z}\phi^\Lambda\mathsf M\,dk\,dl=\Lambda^{4\mathsf z}\mathsf Z(\mathsf z)$. The functions $|k|^{-2\mathsf a_\alpha}k^\alpha$ and $|l|^{-2\mathsf b_\beta}l^\beta$ are homogeneous factors, so by Lemma~\ref{IV-73-86}(a) each factor in the expression for $\mathsf Z$ continues holomorphically to $\mathbb D$. Hence the continuation of $\iint(|k||l|)^{2\mathsf z}(\phi^\Lambda-\phi)\mathsf M\,dk\,dl$ is $(\Lambda^{4\mathsf z}-1)\mathsf Z(\mathsf z)$, which vanishes at $\mathsf z=0$; by Lemma~\ref{IV-73-86}(e) this value is the finite part.

\emph{Step 3: the case $\mathsf c=1$, $(\mathsf a,\mathsf b)\ne(1,1)$.} Now $\mathsf a+\mathsf b=2$. The finite part is symmetric under $k\leftrightarrow l$ (Definition~\ref{IV-73-85}(2)), and $\phi$ may be chosen freely; we may therefore assume $\mathsf b\le0$ and $\mathsf a\ge2$. The degree of the homogeneous factor in \textup{($\Phi2$)} is then $-2\mathsf b\ge0$, so that Lemma~\ref{IV-73-86}(e) applies with $(\epsilon,\epsilon')=(1,0)$. We choose $\phi(k,l):=\varphi(|k|)\varphi(|l-k|)$ and write $\mathsf m:=l-k$. Then
\[
\mathsf M=|k|^{-2\mathsf a}\,|k+\mathsf m|^{2|\mathsf b|}\,|\mathsf m|^{-2}=\sum c_{\alpha\beta}\,|k|^{-2\mathsf a}k^\alpha\cdot\mathsf m^\beta|\mathsf m|^{-2},
\]
and for large $\operatorname{Re}\mathsf z$
\[
\mathsf Z(\mathsf z):=\iint|k|^{2\mathsf z}\,\phi\,\mathsf M\,dk\,dl=\sum c_{\alpha\beta}\int|k|^{2\mathsf z-2\mathsf a}k^\alpha\varphi(|k|)\,dk\int\mathsf m^\beta|\mathsf m|^{-2}\varphi(|\mathsf m|)\,d\mathsf m ,
\]
while $\iint|k|^{2\mathsf z}\phi^\Lambda\mathsf M\,dk\,dl=\Lambda^{2\mathsf z}\mathsf Z(\mathsf z)$. The integrals over $\mathsf m$ converge absolutely, and the integrals over $k$ continue to $\mathbb D$ by Lemma~\ref{IV-73-86}(a). The continuation $(\Lambda^{2\mathsf z}-1)\mathsf Z(\mathsf z)$ vanishes at $\mathsf z=0$.

\emph{Step 4: the case $(\mathsf a,\mathsf b,\mathsf c)=(1,1,1)$.} We choose $\phi(k,l):=\varphi\bigl((k^2+l^2)^{1/2}\bigr)$. All degrees of the homogeneous factors in \textup{($\Phi1$)} and \textup{($\Phi2$)} are now $-2$; hence $(\phi^\Lambda-\phi)\mathsf M$ is integrable, and by Lemma~\ref{IV-73-86}(e) its finite part is its integral. In polar coordinates of $\mathbb R^6$,
\[
\iint(\phi^\Lambda-\phi)\,\mathsf M\,dk\,dl=\oint_{S^5}\mathsf M\,d\omega\cdot\int_0^\infty\bigl[\varphi(s/\Lambda)-\varphi(s)\bigr]\frac{ds}s .
\]
The integrand on the right vanishes for $s\le\frac12$, and for $0<\varepsilon\le\frac12$
\[
\int_\varepsilon^\infty\bigl[\varphi(s/\Lambda)-\varphi(s)\bigr]\frac{ds}s=\int_{\varepsilon/\Lambda}^\varepsilon\varphi(s)\,\frac{ds}s=\log\Lambda .
\]
Together with the value of the integral over $S^5$, computed in Step~5, this gives $4\pi^4\log\Lambda$.

\emph{Step 5: the integral over the sphere.} We parametrise $S^5$ by $k=\cos\psi\,\omega_1$, $l=\sin\psi\,\omega_2$ with $\psi\in(0,\frac\pi2)$ and $\omega_1,\omega_2\in S^2$. In $\mathbb R^6$ we have $dk\,dl=\mathsf r_1^2d\mathsf r_1\,d\omega_1\,\mathsf r_2^2d\mathsf r_2\,d\omega_2$, with $\mathsf r_1=s\cos\psi$, $\mathsf r_2=s\sin\psi$ and $d\mathsf r_1\,d\mathsf r_2=s\,ds\,d\psi$; hence
\[
d\omega=\cos^2\psi\,\sin^2\psi\,d\psi\,d\omega_1\,d\omega_2 .
\]
On $S^5$ we have $\mathsf x_1\mathsf x_2=\cos^2\psi\sin^2\psi$ and $\mathsf x_3=1-\sin2\psi\;\omega_1\cdot\omega_2$. For Borel functions $g\ge0$ on $[-1,1]$ we have $\iint_{S^2\times S^2}g(\omega_1\cdot\omega_2)\,d\omega_1\,d\omega_2=4\pi\cdot2\pi\int_{-1}^1g(\tau)\,d\tau$, and for $0\le z<1$
\[
\int_{-1}^1\frac{d\tau}{1-z\tau}=\sum_{j\ge0}z^j\int_{-1}^1\tau^j\,d\tau=2\sum_{n\ge0}\frac{z^{2n}}{2n+1}.
\]
Applying this with $z=\sin2\psi$ for $\psi\ne\frac\pi4$, and interchanging sum and integral, which is allowed since all terms are non-negative, we find
\[
\begin{multlined}
\oint_{S^5}\frac{d\omega}{\mathsf x_1\mathsf x_2\mathsf x_3}=8\pi^2\int_0^{\pi/2}2\sum_{n\ge0}\frac{\sin^{2n}2\psi}{2n+1}\,d\psi=16\pi^2\sum_{n\ge0}\frac{\mathsf w_n}{2n+1},\qquad\mathsf w_n:=\int_0^{\pi/2}\sin^{2n}2\psi\,d\psi \\
=\int_0^{\pi/2}\sin^{2n}\psi\,d\psi .
\end{multlined}
\]
An integration by parts gives $2n\,\mathsf w_n=(2n-1)\,\mathsf w_{n-1}$, and $\mathsf w_0=\frac\pi2$. Hence $\mathsf w_n=\frac\pi2\,\mathsf b_n$ with
\[
\mathsf b_n:=\prod_{j=1}^n\frac{2j-1}{2j},
\]
the coefficients of the binomial series $(1-t^2)^{-1/2}=\sum_n\mathsf b_nt^{2n}$, $|t|<1$. Integrating this series over $[0,x]$ and letting $x\uparrow1$ (monotone convergence), we obtain
\[
\sum_n\frac{\mathsf b_n}{2n+1}=\int_0^1(1-t^2)^{-1/2}\,dt=\frac\pi2 ,
\]
where we have substituted $t=\sin\psi$, $0\le\psi<\frac\pi2$, so that $(1-t^2)^{-1/2}dt=d\psi$. Therefore the integral over $S^5$ equals $16\pi^2\cdot\frac\pi2\cdot\frac\pi2=4\pi^4$, and $4\pi^4/(2\pi)^6=1/16\pi^2$.
\end{proof}

\begin{proposition}[The bridge: $\kappa_\Gamma^{\mathrm{lat}}$ from $\mathcal R{[F_\Gamma]}$]\label{IV-73-94}
Let $\Gamma\in\{1,\dots,6\}$ and $\Lambda>1$; let $\mathfrak f_\Gamma$ be as in Proposition~\ref{IV-73-87}, $R_\Gamma^\Lambda$ as in Lemma~\ref{IV-73-84}(b), and $\kappa_\Gamma^{\mathrm{lat}}$ as in Lemma~\ref{IV-73-70}. Then
\begin{equation}\label{IV-73-95}
\mathfrak f_\Gamma\bigl[R_\Gamma^\Lambda-R_\Gamma\bigr]=d_{2,\Gamma}\log\Lambda,\quad d_{2,\Gamma}:=\frac{\mathcal R[F_\Gamma]}{16\pi^2};\qquad\kappa_\Gamma^{\mathrm{lat}}=-\frac3{2\pi}\cdot24\cdot d_{2,\Gamma}=-\frac3{32\pi^3}\cdot24\cdot\mathcal R[F_\Gamma].
\end{equation}
\end{proposition}

\begin{proof}
\emph{Step 1: cut-offs.} The function $R_\Gamma$ is smooth, even and rotation invariant. If $|k|+|l|<\frac12\rho_0$ and $|\mathsf p|<\frac12\rho_0$, then each of the vectors $k$, $l$, $k-l$, $\mathsf p-k$, $\mathsf p-l$, among which are the lines of $\Gamma$, has length less than $\rho_0$, and therefore $R_\Gamma=1=R_\Gamma^\Lambda$ there. If $R_\Gamma\ne0$, then $|l|<2\rho_0$ and $|k|<4\rho_0$: for $\Gamma\ne6$ because $l$ and $k$ are lines of $\Gamma$, and for $\Gamma=6$ because $l$ and $k-l$ are. Put $\phi(k,l):=\theta\bigl((k^2+l^2)^{1/2}\bigr)$; then $\phi=1$ for $|k|+|l|\le\frac12$. Consequently the functions
\[
\Theta:=R_\Gamma^\Lambda-R_\Gamma,\qquad\Upsilon_0:=R_\Gamma-\phi,\qquad\phi^\Lambda-\phi ,
\]
regarded as functions of $(k,l,\mathsf p)$, belong to $\mathcal Y(c)$ with $c:=\frac12\rho_0$. Indeed, they are smooth and satisfy \textup{(Y2)}; for $|\mathsf p|<c$ and $|k|+|l|<c$ they vanish, because $R_\Gamma$, $R_\Gamma^\Lambda$, $\phi$ and $\phi^\Lambda$ are all equal to $1$ there ($c<\frac12$); and they vanish when $|k|+|l|$ is large, because $R_\Gamma$, $R_\Gamma^\Lambda$, $\phi$ and $\phi^\Lambda$ do. We note that, by Lemma~\ref{IV-73-86}(e), the finite part, and with it $\mathfrak f_\Gamma$, does not depend on the parameter $c$ with respect to which an element of $\mathcal Y(c)$ is considered.

\emph{Step 2: $\mathfrak f_\Gamma[\Theta]=-\frac\pi{36}\kappa_\Gamma^{\mathrm{lat}}\log\Lambda$.} The function $\Xi_\Gamma[\Theta](\mathsf p)$ is the left side of the identity in Lemma~\ref{IV-73-84}(b). We thus have two expansions of it: that of Lemma~\ref{IV-73-84}(b), valid for almost every $\mathsf p$ with $|\mathsf p|<1$, and that of Proposition~\ref{IV-73-87}, valid for every $\mathsf p$ with $0<|\mathsf p|\le\frac c{16}$. Since the remainder $\Lambda^2\mathsf o(\mathsf p/\Lambda)-\mathsf o(\mathsf p)$ of the former is $o(\mathsf p^2)$, they give
\[
(\Lambda^2-1)\,\mathsf a_0+(\Lambda-1)\,\mathsf a_1|\mathsf p|-\tfrac\pi{36}\,\kappa_\Gamma^{\mathrm{lat}}\log\Lambda\;\mathsf p^2+o(\mathsf p^2)=\alpha+\beta\,|\mathsf p|+\mathfrak f_\Gamma[\Theta]\,\mathsf p^2+o(\mathsf p^2).
\]
Along a sequence $\mathsf p\to0$ of points at which both expansions hold, we compare first the constant terms, then the coefficients of $|\mathsf p|$, and then those of $\mathsf p^2$.

\emph{Step 3: $\mathfrak f_\Gamma[\Theta]=\mathfrak f_\Gamma[\phi^\Lambda-\phi]$.} We have $\Theta=(\phi^\Lambda-\phi)+(\Upsilon_0^\Lambda-\Upsilon_0)$. By the linearity of $\mathfrak f_\Gamma$ and the identity $\mathfrak f_\Gamma[\Upsilon_0^\Lambda]=\mathfrak f_\Gamma[\Upsilon_0]$ of Proposition~\ref{IV-73-87}, the second difference contributes nothing.

\emph{Step 4: the value.} Since $\phi$ does not depend on $\mathsf p$, we have $\overline{((\phi^\Lambda-\phi)F_\Gamma)_2}=(\phi^\Lambda-\phi)\,\overline{(F_\Gamma)_2}$. By Definition~\ref{IV-73-88}(3) and Lemma~\ref{IV-73-90}, $\overline{(F_\Gamma)_2}=\sum_{\mathsf M}r_{\mathsf M}\mathsf M$ is a finite linear combination of monomials $\mathsf M$ as in Lemma~\ref{IV-73-92}, and the coefficient of $\mathsf M=(\mathsf x_1\mathsf x_2\mathsf x_3)^{-1}$ is $r_{\mathsf M}=\mathcal R[F_\Gamma]$. The finite part is linear and independent of the parameter $c$ (Lemma~\ref{IV-73-86}(e)); hence \eqref{IV-73-93} yields
\[
\mathfrak f_\Gamma[\phi^\Lambda-\phi]=(2\pi)^{-6}\cdot4\pi^4\log\Lambda\cdot\mathcal R[F_\Gamma]=d_{2,\Gamma}\log\Lambda .
\]

\emph{Step 5: conclusion.} By Steps~2--4, $\mathfrak f_\Gamma[R_\Gamma^\Lambda-R_\Gamma]=\mathfrak f_\Gamma[\Theta]=d_{2,\Gamma}\log\Lambda$ and $-\frac\pi{36}\kappa_\Gamma^{\mathrm{lat}}=d_{2,\Gamma}$. Hence
\[
\kappa_\Gamma^{\mathrm{lat}}=-\frac{36}\pi\,d_{2,\Gamma}=-\frac3{2\pi}\cdot24\cdot d_{2,\Gamma},
\]
and $\frac3{2\pi}\cdot\frac1{16\pi^2}=\frac3{32\pi^3}$ gives the last expression in \eqref{IV-73-95}.
\end{proof}

\begin{proposition}[The six numbers $\mathcal R{[F_\Gamma]}$]\label{IV-73-96}
For the integrands $F_1,\dots,F_6$ of \eqref{IV-73-65},
\begin{equation}\label{IV-73-97}
\begin{multlined}
\mathcal R[F_1]=-\tfrac43,\quad\mathcal R[F_2]=-\tfrac{13}{12},\quad\mathcal R[F_3]=-\tfrac1{12},\quad\mathcal R[F_4]=0,\quad\mathcal R[F_5]=+\tfrac{13}6,\quad\mathcal R[F_6]= \\
-\tfrac13;\qquad\sum_{\Gamma=1}^6\mathcal R[F_\Gamma]=-\tfrac23 .
\end{multlined}
\end{equation}
\end{proposition}

\begin{proof}
\emph{Step 0: notation.} Throughout this proof we use the abbreviations
\[
A,\ A',\ B,\ B',\ C,\ P:=\mathsf x_1,\ \mathsf x_1',\ \mathsf x_2,\ \mathsf x_2',\ \mathsf x_3,\ \mathsf x_0 ,
\]
and we put $k':=k-\mathsf p$, $l':=l-\mathsf p$ and $q:=k-l=k'-l'$, so that $A'=k'^2$, $B'=l'^2$ and $C=q^2$. The relations $P=(k-k')^2=(l-l')^2$, $C=(k-l)^2=(k'-l')^2$, $k\cdot\mathsf p=A-k\cdot k'$ and $l\cdot\mathsf p=B-l\cdot l'$ express the remaining scalar products through these variables:
\[
\begin{gathered}
a:=k\cdot k'=\tfrac12(A+A'-P),\quad\\
 b:=l\cdot l'=\tfrac12(B+B'-P),\quad\\
 e:=k\cdot l=\tfrac12(A+B-C),\quad\\
 e':=k'\cdot l'=\tfrac12(A'+B'-C),\\
f:=k\cdot l'=e-k\cdot\mathsf p=\tfrac12(A'+B-C-P),\qquad\\
 f':=k'\cdot l=e-l\cdot\mathsf p=\tfrac12(A+B'-C-P).
\end{gathered}
\]
We shall use two symmetries. The exchange of primes, $(k,l)\leftrightarrow(k',l')$, interchanges $A\leftrightarrow A'$, $B\leftrightarrow B'$, $e\leftrightarrow e'$, $f\leftrightarrow f'$ and fixes $a$ and $b$. The exchange $(k,k')\leftrightarrow(l,l')$ interchanges $(A,A')\leftrightarrow(B,B')$, $a\leftrightarrow b$, $f\leftrightarrow f'$ and fixes $e$ and $e'$. Finally, $\mathsf q(1)=\frac13$, $\mathsf q(0)=0$ and $\mathsf q(-1)=1$.

\emph{Step 1: two identities for $V_3$.} Let
\[
V:=V_3(k_1,k_2,k_3)_{\alpha\beta\gamma}=\delta_{\alpha\gamma}\mathsf a_\beta+\delta_{\alpha\beta}\mathsf b_\gamma+\delta_{\beta\gamma}\mathsf c_\alpha,\qquad\mathsf a:=k_1-k_3,\quad\mathsf b:=k_2-k_1,\quad\mathsf c:=k_3-k_2 .
\]
(i) Since $\mathsf a+\mathsf b+\mathsf c=0$,
\[
V_{\alpha\beta\gamma}V_{\alpha\beta\gamma}=3(\mathsf a^2+\mathsf b^2+\mathsf c^2)+2(\mathsf a\cdot\mathsf b+\mathsf a\cdot\mathsf c+\mathsf b\cdot\mathsf c)=2(\mathsf a^2+\mathsf b^2+\mathsf c^2),\qquad\mathsf a^2+\mathsf b^2+\mathsf c^2=3\sum_ik_i^2-\Bigl(\sum_ik_i\Bigr)^2 .
\]
(ii) If $k_1+k_2+k_3=0$, then
\[
k_{1\alpha}V_{\alpha\beta\gamma}=(\mathsf P_2-\mathsf P_3)_{\beta\gamma},\qquad\mathsf P_i:=k_i^2\,\delta_{\beta\gamma}-k_{i\beta}k_{i\gamma} .
\]
Indeed, $k_{1\alpha}V_{\alpha\beta\gamma}=k_{1\gamma}(k_1-k_3)_\beta+k_{1\beta}(k_2-k_1)_\gamma+\delta_{\beta\gamma}\,k_1\cdot(k_3-k_2)$; here $k_{1\gamma}(k_1-k_3)_\beta+k_{1\beta}(k_2-k_1)_\gamma=k_{1\beta}k_{2\gamma}-k_{1\gamma}k_{3\beta}$, which equals $k_{3\beta}k_{3\gamma}-k_{2\beta}k_{2\gamma}$ because $k_1=-k_2-k_3$, and for the same reason $k_1\cdot(k_3-k_2)=k_2^2-k_3^2$.

\emph{Step 2: $\Gamma=4$.} By Definition~\ref{IV-73-88}(3), $F_4$ is a linear combination of monomials with $a_5\le0$, whereas every non-zero case of \eqref{IV-73-91} has $\alpha_3=a_5\ge1$. Hence $\mathcal R[F_4]=0$.

\emph{Step 3: $\Gamma=6$.} Since $V_3$ is odd, $V_3(-l,l-k,k-\mathsf p)=-V_3(l,k-l,\mathsf p-k)$, and therefore
\[
F_6=+\tfrac16\,V_{\alpha\beta\gamma}V_{\alpha\beta\gamma}/A'BC,\qquad V:=V_3(l,q,-k'),
\]
where the three arguments of $V_3$ add up to $l+q-k'=\mathsf p$. By Step~1(i), $V_{\alpha\beta\gamma}V_{\alpha\beta\gamma}=2\bigl[3(B+C+A')-P\bigr]$, hence
\[
F_6=\frac1{A'C}+\frac1{A'B}+\frac1{BC}-\frac13\,\frac P{A'BC}.
\]
We evaluate $\mathcal R$ by \eqref{IV-73-91}. The monomial $1/A'C$ has $n=0$, $(\alpha_1,\alpha_2,\alpha_3)=(1,0,1)$ and $a_2=1$, $a_4=0$, so its value is $\mathsf q(0)=0$; the monomial $1/A'B$ has $(\alpha_1,\alpha_2,\alpha_3)=(1,1,0)$ and value $0$; the monomial $1/BC$ has $n=0$, $(\alpha_1,\alpha_2,\alpha_3)=(0,1,1)$ and $a_2=0$, so its value is $\mathsf q(0)=0$; and $P/A'BC$ has $n=1$ and $(\alpha_1,\alpha_2,\alpha_3)=(1,1,1)$, so its value is $1$. Therefore $\mathcal R[F_6]=-\frac13$.

\emph{Step 4: $\Gamma=2,3$, a common factor.} Put
\[
T_{\mu\nu}:=V_2(k,\mathsf p-k)_{\mu\nu}=V_2(k,-k')_{\mu\nu}=a\,\delta_{\mu\nu}-k_\nu k'_\mu ;
\]
as $V_2$ is even, also $V_2(-k,k-\mathsf p)_{\rho\nu}=T_{\rho\nu}$. Let
\[
H_{\mu\rho}:=T_{\mu\nu}T_{\rho\nu}/A'=\bigl[a^2\delta_{\mu\rho}-a\,(k_\rho k'_\mu+k_\mu k'_\rho)+A\,k'_\mu k'_\rho\bigr]/A',
\]
and $V:=V_3(-k,l,q)_{\mu\beta\gamma}$, so that $V_3(k,-l,-q)=-V$. In this notation
\[
F_2=\tfrac12\,H_{\mu\rho}V_{\mu\beta\gamma}V_{\rho\beta\gamma}/A^2BC,\qquad F_3=-H_{\mu\rho}\,l_\rho(l-k)_\mu/A^2BC ,
\]
and $\mathsf p$ enters these expressions through $H$ only. We compute $\overline{(H_{\mu\rho})_2}$. We insert $k'=k-\mathsf p$, $a=A-x$ and $A'=A-2x+P$, with $x:=k\cdot\mathsf p$, into $H$ and expand: the term of second order in $\mathsf p$ is a sum of products of $\mathsf p_i\mathsf p_j$ with tensors built from $\delta$ and $k$, the coefficients being powers of $A$. Applying $\frac16\Delta_{\mathsf p}$, we obtain
\[
\overline{(H_{\mu\rho})_2}=\mathsf a_1\,\delta_{\mu\rho}+\mathsf a_2\,k_\mu k_\rho/A
\]
with numbers $\mathsf a_1,\mathsf a_2$, since $H$ is homogeneous of degree $2$ in $(k,\mathsf p)$. Next, $k_\mu T_{\mu\nu}=a\,k_\nu-k_\nu(k\cdot k')=0$ for all $\mathsf p$; hence $k_\mu\overline{(H_{\mu\rho})_2}=0$, which gives $\mathsf a_2=-\mathsf a_1$. To determine $\mathsf a_1$ we take the trace. From $T_{\mu\nu}T_{\mu\nu}=3a^2-2a^2+AA'$ we get $H_{\mu\mu}=a^2/A'+A$, and with $a=A-x$, $A'=A-2x+P$,
\[
\frac{a^2}{A'}=\Bigl(A-2x+\frac{x^2}A\Bigr)\Bigl(1+\frac{2x}A-\frac PA+\frac{4x^2}{A^2}+O(|\mathsf p|^3)\Bigr),
\]
whose term of second order is $-P+4x^2/A-4x^2/A+x^2/A$; its image under $\frac16\Delta_{\mathsf p}$ is $-1+\frac13$. Since the trace of $\overline{(H_{\mu\rho})_2}$ is $3\mathsf a_1+\mathsf a_2=2\mathsf a_1$, we obtain $2\mathsf a_1=-\frac23$, that is,
\[
\overline{(H_{\mu\rho})_2}=-\tfrac13\,\bigl(\delta_{\mu\rho}-k_\mu k_\rho/A\bigr).
\]

\emph{Step 5: $\Gamma=2$.} By Step~4,
\[
\overline{(F_2)_2}=-\tfrac16\,\mathsf W/A^2BC,\qquad\mathsf W:=V_{\mu\beta\gamma}V_{\mu\beta\gamma}-|k_\mu V_{\mu\beta\gamma}|^2/A .
\]
The function $A\mathsf W$ is a polynomial in $A,B,C$, homogeneous of degree $2$, and $\overline{(F_2)_2}=-\frac16A\mathsf W/A^3BC$. A monomial $A^iB^jC^m$ of $A\mathsf W$ thus contributes a multiple of $A^{i-3}B^{j-1}C^{m-1}$, and only $A^2$ yields $(ABC)^{-1}$; we may therefore set $B=C=0$. By Step~1(i), applied to the arguments $-k,l,q$, whose sum is $0$, $V_{\mu\beta\gamma}V_{\mu\beta\gamma}=6(A+B+C)$, which becomes $6A$. By Step~1(ii), applied with $(k_1,k_2,k_3)=(-k,l,q)$, we have $|k_\mu V_{\mu\beta\gamma}|^2=\operatorname{tr}(\mathsf P_l-\mathsf P_q)^2$ with $\mathsf P_l=B\delta-ll^{\mathsf T}$ and $\mathsf P_q=C\delta-qq^{\mathsf T}$. Here $\operatorname{tr}\mathsf P_l^2=2B^2$, $\operatorname{tr}\mathsf P_q^2=2C^2$, $\operatorname{tr}\mathsf P_l\mathsf P_q=BC+(l\cdot q)^2$ and $l\cdot q=e-B=\frac12(A-B-C)$; thus $|k_\mu V_{\mu\beta\gamma}|^2=2B^2+2C^2-2BC-2(l\cdot q)^2$ is a polynomial in $A,B,C$, which becomes $-2(\frac12A)^2=-\frac12A^2$ at $B=C=0$.

Hence $A\mathsf W$ becomes $6A^2+\frac12A^2$, and
\[
\mathcal R[F_2]=-\tfrac16\cdot\tfrac{13}2=-\tfrac{13}{12}.
\]

\emph{Step 6: $\Gamma=3$.} By Step~4,
\[
\begin{multlined}
\overline{(F_3)_2}=+\tfrac13\,\mathsf W'/A^2BC,\qquad\mathsf W':=\bigl(\delta_{\mu\rho}-k_\mu k_\rho/A\bigr)\,l_\rho(l-k)_\mu=l\cdot(l-k)-\frac{(k\cdot l)\,(k\cdot(l-k))}A \\
=B-e-\frac{e(e-A)}A .
\end{multlined}
\]
Again $A\mathsf W'$ is a polynomial in $A,B,C$, homogeneous of degree $2$, and as in Step~5 only its coefficient of $A^2$ matters. At $B=C=0$ we have $e=\frac12A$, and $A\mathsf W'$ becomes $-\frac12A^2+\frac14A^2$. Hence
\[
\mathcal R[F_3]=\tfrac13\cdot\bigl(-\tfrac14\bigr)=-\tfrac1{12}.
\]

\emph{Step 7: values for the denominator $AA'BC$.} For a monomial $\mathsf m=P^nA^iA'^{i'}B^jB'^{j'}C^m$ of degree $2$ we put $\mathsf L_5(\mathsf m):=\mathcal R[\mathsf m/AA'BC]$. In the notation of Lemma~\ref{IV-73-90}, the monomial $\mathsf m/AA'BC$ has the exponents $(a_1,\dots,a_5)=(1-i,1-i',1-j,-j',1-m)$, and $(\alpha_1,\alpha_2,\alpha_3)=(2-i-i',\,1-j-j',\,1-m)$. By \eqref{IV-73-91}, $\mathsf L_5(\mathsf m)\ne0$ requires $m=0$, and the values are as follows.
\begin{itemize}
\item $n=1$ (then $i+i'=1$, $j=j'=0$): $\mathsf L_5(PA)=\mathsf L_5(PA')=1$.
\item $n=0$, $(\alpha_1,\alpha_2,\alpha_3)=(0,1,1)$ (then $i+i'=2$, $a_4=0$, value $\mathsf q(1-i')$): $\mathsf L_5(A^2)=\frac13$, $\mathsf L_5(AA')=0$, $\mathsf L_5(A'^2)=1$.
\item $n=0$, $(\alpha_1,\alpha_2,\alpha_3)=(1,0,1)$ (then $i+i'=1$, $j+j'=1$, value $\mathsf q(-j')-\frac23(1-i')j'$): $\mathsf L_5(AB)=\mathsf L_5(A'B)=0$, $\mathsf L_5(AB')=1-\frac23=\frac13$, $\mathsf L_5(A'B')=1$.
\end{itemize}
All other monomials of degree $2$ have the value $0$; in particular, every monomial containing $B$ or $C$ has the value $0$.

\emph{Step 8: $\Gamma=5$.} Put
\[
\tilde T_{\alpha\sigma}:=V_2(k-\mathsf p,-k)_{\alpha\sigma}=V_2(k',-k)_{\alpha\sigma}=a\,\delta_{\alpha\sigma}-k_\alpha k'_\sigma .
\]
Then $k'_\alpha\tilde T_{\alpha\sigma}=0=\tilde T_{\alpha\sigma}k_\sigma$; we write $\tilde T(X,Y):=X_\alpha\tilde T_{\alpha\sigma}Y_\sigma=a\,X\cdot Y-(k\cdot X)(k'\cdot Y)$, and we note $\tilde T_{\alpha\alpha}=2a$. Further,
\[
\begin{aligned}
V_3(\mathsf p-k,l,k-l)&=V_3(-k',l,q)\\
&=\delta_{\alpha\gamma}u_{1\beta}+\delta_{\alpha\beta}v_{1\gamma}+\delta_{\beta\gamma}w_{1\alpha},\\
u_1&=-k'-q,\qquad v_1=l+k',\qquad w_1=q-l,\\
V_3(k,-l,l-k)&=\delta_{\sigma\gamma}u_{2\beta}+\delta_{\sigma\beta}v_{2\gamma}+\delta_{\beta\gamma}w_{2\sigma},\\
u_2&=k+q,\qquad v_2=-(k+l),\qquad w_2=l-q .
\end{aligned}
\]
Hence $F_5=N_5/AA'BC$, where, after multiplying out the nine products and summing over $\beta,\gamma$,
\[
\begin{multlined}
N_5=2a\,(u_1\cdot u_2+v_1\cdot v_2)+\tilde T(v_2,u_1)+\tilde T(u_1,w_2)+\tilde T(u_2,v_1)+\tilde T(v_1,w_2)+\tilde T(w_1,u_2) \\
+\tilde T(w_1,v_2)+3\,\tilde T(w_1,w_2).
\end{multlined}
\]
The first argument of $\tilde T$ matters only modulo multiples of $k'$, the second only modulo multiples of $k$. Using $q=k-l$ we find, in the first argument,
\[
u_1\equiv l-k,\qquad v_1\equiv l,\qquad w_1=k-2l,\qquad u_2=2k-l,\qquad v_2=-(k+l),
\]
and in the second argument
\[
u_1\equiv l-k',\qquad v_1=l+k',\qquad u_2\equiv-l,\qquad v_2\equiv-l,\qquad w_2\equiv2l .
\]
Therefore
\[
\begin{gathered}
\tilde T(u_1,w_2)=-2\tilde T(k,l)+2\tilde T(l,l),\qquad\\
\tilde T(v_1,w_2)=2\tilde T(l,l),\qquad\\
\tilde T(w_1,u_2)=\tilde T(w_1,v_2)=-\tilde T(k,l)+2\tilde T(l,l),\\
3\,\tilde T(w_1,w_2)=6\tilde T(k,l)-12\tilde T(l,l),\qquad\\
\tilde T(v_2,u_1)+\tilde T(u_2,v_1)=-\tilde T(k+l,l-k')+\tilde T(2k-l,l+k')=\tilde T(k,l)-2\tilde T(l,l)+3\tilde T(k,k'),
\end{gathered}
\]
and the sum of these seven terms is $3\tilde T(k,l)-6\tilde T(l,l)+3\tilde T(k,k')$. We regard $N_5$ as a polynomial in $A,A',B,B',C,P$; by the last assertion of Step~7 we may set $B=C=0$ in it. Then $e=\frac12A$, $f'=\frac12(A+B'-P)$, $k\cdot q=A-e=\frac12A$ and $k'\cdot q=a-f'$, and with the values of Step~7 we find:
\[
\begin{aligned}
\tilde T(k,l)&=ae-Af'=\tfrac14A\,(-A+A'-2B'+P):\\
\mathsf L_5&=\tfrac14\bigl(-\tfrac13+0-\tfrac23+1\bigr)=0;\\
\tilde T(l,l)&=aB-ef'=-\tfrac14A\,(A+B'-P):\\
\mathsf L_5&=-\tfrac14\bigl(\tfrac13+\tfrac13-1\bigr)=\tfrac1{12};\\
\tilde T(k,k')&=a^2-AA'=\tfrac14(A+A'-P)^2-AA':\\
\mathsf L_5&=\tfrac14\bigl(\tfrac13+1-2-2\bigr)=-\tfrac23;\\
u_1\cdot u_2+v_1\cdot v_2&\\
=-(a+k'\cdot q+k\cdot q+C)-(e+B+a+f')\\
=-(3a+A);\\
a\,(3a+A)&=\tfrac34(A+A'-P)^2+\tfrac12A\,(A+A'-P):\\
\mathsf L_5&=\tfrac34\bigl(-\tfrac83\bigr)+\tfrac12\bigl(\tfrac13-1\bigr)=-\tfrac73 .
\end{aligned}
\]
In the third line the four numbers in the bracket come from the monomials $A^2$, $A'^2$, $-2PA$, $-2PA'$ of $(A+A'-P)^2$, the others having the value $0$; thus $\mathsf L_5\bigl((A+A'-P)^2\bigr)=\frac13+1-2-2=-\frac83$, the value used in the last line. Consequently
\[
\mathcal R[F_5]=-2\bigl(-\tfrac73\bigr)+3\cdot0-6\cdot\tfrac1{12}+3\bigl(-\tfrac23\bigr)=\tfrac{14}3-\tfrac52=\tfrac{13}6 .
\]

\emph{Step 9: values for the denominator $AA'BB'C$.} For a monomial $\mathsf m=P^nA^iA'^{i'}B^jB'^{j'}C^m$ of degree $3$ we put $\mathsf L(\mathsf m):=\mathcal R[\mathsf m/AA'BB'C]$. Now $(a_1,\dots,a_5)=(1-i,1-i',1-j,1-j',1-m)$ and $(\alpha_1,\alpha_2,\alpha_3)=(2-i-i',\,2-j-j',\,1-m)$, and again $\mathsf L(\mathsf m)\ne0$ requires $m=0$. Let $X\in\{A,A'\}$ and $Y\in\{B,B'\}$. The values are as follows.
\begin{itemize}
\item $n=1$: $\mathsf L(PXY)=1$.
\item $n=0$, $(\alpha_1,\alpha_2,\alpha_3)=(0,1,1)$ (then $i+i'=2$, $j+j'=1$, value $\mathsf q(1-i')+\frac23(1-i')(1-j')$): $\mathsf L(A^2B)=\frac13+\frac23=1$, $\mathsf L(A^2B')=\frac13$, $\mathsf L(A'^2B)=1-\frac23=\frac13$, $\mathsf L(A'^2B')=1$, $\mathsf L(AA'Y)=0$.
\item $n=0$, $(\alpha_1,\alpha_2,\alpha_3)=(1,0,1)$ (then $i+i'=1$, $j+j'=2$, value $\mathsf q(1-j')+\frac23(1-i')(1-j')$): $\mathsf L(AB^2)=1$, $\mathsf L(A'B^2)=\frac13$, $\mathsf L(AB'^2)=\frac13$, $\mathsf L(A'B'^2)=1$, $\mathsf L(XBB')=0$.
\end{itemize}
All other monomials of degree $3$ have the value $0$: those containing $C$ or $P^2$, and those of the forms $PX_1X_2$, $PY_1Y_2$, $X_1X_2X_3$, $Y_1Y_2Y_3$ with $X_i\in\{A,A'\}$, $Y_i\in\{B,B'\}$. These values are invariant under the exchange of primes and under $(A,A')\leftrightarrow(B,B')$.

\emph{Step 10: $\Gamma=1$, reduction.} As in Step~4, $T_{\mu\nu}=V_2(k,\mathsf p-k)_{\mu\nu}=a\,\delta_{\mu\nu}-k_\nu k'_\mu$, and $k_\mu T_{\mu\nu}=0=T_{\mu\nu}k'_\nu$. Similarly,
\[
S_{\rho\sigma}:=V_2(-l,l-\mathsf p)_{\rho\sigma}=V_2(-l,l')_{\rho\sigma}=b\,\delta_{\rho\sigma}-l_\sigma l'_\rho,\qquad l_\rho S_{\rho\sigma}=0=S_{\rho\sigma}l'_\sigma .
\]
The tensor $V_3(-k,l,q)_{\mu\rho\alpha}=\delta_{\mu\alpha}(l-2k)_\rho+\delta_{\mu\rho}(k+l)_\alpha+\delta_{\rho\alpha}(k-2l)_\mu$ is contracted with $T_{\mu\nu}S_{\rho\sigma}$, which annihilates the terms with $l_\rho$ and with $k_\mu$; in this contraction it may therefore be replaced by
\[
\mathsf W_{\mu\rho\alpha}:=-2k_\rho\,\delta_{\mu\alpha}+(k+l)_\alpha\,\delta_{\mu\rho}-2l_\mu\,\delta_{\rho\alpha}.
\]
Likewise, $V_3(k-\mathsf p,\mathsf p-l,l-k)_{\nu\sigma\alpha}=V_3(k',-l',l'-k')_{\nu\sigma\alpha}=\delta_{\nu\alpha}(2k'-l')_\sigma-\delta_{\nu\sigma}(k'+l')_\alpha+\delta_{\sigma\alpha}(2l'-k')_\nu$, contracted with $T_{\mu\nu}S_{\rho\sigma}$, may be replaced by $-\mathsf W'_{\nu\sigma\alpha}$, where $\mathsf W'$ denotes $\mathsf W$ with $(k',l')$ in place of $(k,l)$. Hence
\[
F_1=+\tfrac12\,N'/AA'BB'C,\qquad N':=T_{\mu\nu}S_{\rho\sigma}\,\mathsf W_{\mu\rho\alpha}\mathsf W'_{\nu\sigma\alpha}.
\]
We use the bilinear forms
\[
T(X,Y):=X_\mu T_{\mu\nu}Y_\nu=a\,X\cdot Y-(k'\cdot X)(k\cdot Y),\qquad S(X,Y):=X_\rho S_{\rho\sigma}Y_\sigma=b\,X\cdot Y-(l'\cdot X)(l\cdot Y),
\]
the vectors $\mathsf s_\sigma:=k_\rho S_{\rho\sigma}=(bk-fl)_\sigma$ and $\mathsf t_\rho:=S_{\rho\sigma}k'_\sigma=(bk'-f'l')_\rho$, and the identities $T_{\mu\mu}=2a$, $S_{\rho\rho}=2b$, $T_{\mu\nu}S_{\mu\nu}=3ab-ab-ab+ee'=ab+ee'$. Let $N'_{ij}$ be the contribution to $N'$ of the $i$-th term of $\mathsf W$ and the $j$-th term of $\mathsf W'$. Then
\[
\begin{gathered}
N'_{11}=4T_{\mu\mu}S(k,k')=8a(ab-ff'),\qquad\\
 N'_{33}=4T(l,l')S_{\rho\rho}=8b(ab-ff'),\qquad\\
 N'_{22}=(k+l)\cdot(k'+l')\,(ab+ee'),\\
N'_{13}=4T(\mathsf s,l'),\quad\\
 N'_{31}=4T(l,\mathsf t),\quad\\
 N'_{12}=-2T(k'+l',\mathsf s),\quad\\
 N'_{21}=-2T(\mathsf t,k+l),\\
N'_{23}=-2\,T_{\mu\nu}l'_\nu\,S_{\mu\sigma}(k+l)_\sigma,\qquad\\
 N'_{32}=-2\,l_\mu T_{\mu\nu}\,(k'+l')_\rho S_{\rho\nu}.
\end{gathered}
\]
Under the exchange of primes, $T_{\mu\nu}$ goes over into $T_{\nu\mu}$, so that $T(X,Y)$ becomes the value of $T$ at the images of $Y$ and $X$, in this order, and $\mathsf s$ goes over into $\mathsf t$; hence $N'_{13}$ goes over into $N'_{31}$, and $N'_{12}$ into $N'_{21}$. Under the exchange $(k,k')\leftrightarrow(l,l')$, the tensors $T_{\mu\nu}$ and $S_{\mu\nu}$ are interchanged; hence $N'_{12}=-2(k'+l')_\mu T_{\mu\nu}k_\rho S_{\rho\nu}$ goes over into $N'_{32}$, $N'_{21}=-2S_{\mu\sigma}k'_\sigma T_{\mu\nu}(k+l)_\nu$ into $N'_{23}$, and $N'_{11}$ into $N'_{33}$. By the invariance of the values in Step~9,
\[
\mathsf L(N'_{33})=\mathsf L(N'_{11}),\qquad\mathsf L(N'_{31})=\mathsf L(N'_{13}),\qquad\mathsf L(N'_{21})=\mathsf L(N'_{32})=\mathsf L(N'_{23})=\mathsf L(N'_{12}).
\]

\emph{Step 11: $\Gamma=1$, values.} By Step~9 we may set $C=0$. Then
\[
4(ab-ff')=(A+A'-P)(B+B'-P)-(A'+B-P)(A+B'-P)=AB'+A'B-AA'-BB'=:\mathsf D .
\]
\begin{itemize}
\item $N'_{11}=(A+A'-P)\,\mathsf D$. Of its monomials only $A^2B'$, $A'^2B$, $-PAB'$ and $-PA'B$ have non-zero values, and $\mathsf L(N'_{11})=\frac13+\frac13-1-1=-\frac43$.
\item $N'_{13}$: since $\mathsf s\cdot l'=bf-fb=0$ and $k'\cdot\mathsf s=ab-ff'$, we have $N'_{13}=-4f(ab-ff')=-\frac12(A'+B-P)\,\mathsf D$. In $(A'+B-P)\,\mathsf D$ the monomials with non-zero values are $A'^2B$, $A'B^2$, $-PAB'$ and $-PA'B$; they give $\frac13+\frac13-2=-\frac43$, and therefore $\mathsf L(N'_{13})=\frac23$.
\item $N'_{22}$: here $(k+l)\cdot(k'+l')=a+f+f'+b=A+A'+B+B'-2P$, and
\[
\begin{multlined}
4(ab+ee')=(A+A'-P)(B+B'-P)+(A+B)(A'+B')=\mathsf E+AA'+BB'-P(A \\
+A'+B+B')+P^2,\qquad\mathsf E:=AB+2AB'+2A'B+A'B' .
\end{multlined}
\]
We have $\mathsf L\bigl[(A+A'+B+B')\,\mathsf E\bigr]=(1+\frac23)+(\frac23+1)+(1+\frac23)+(\frac23+1)=\frac{20}3$, the four contributions coming from $A(AB+2AB')$, $A'(2A'B+A'B')$, $B(AB+2A'B)$ and $B'(2AB'+A'B')$. Further, $\mathsf L\bigl[(A+A'+B+B')(AA'+BB')\bigr]=0$, $\mathsf L\bigl[-P(A+A'+B+B')^2\bigr]=-2\cdot4=-8$ and $\mathsf L[-2P\,\mathsf E]=-2\cdot6=-12$, while the remaining monomials, which contain $P^2$ or have one of the forms $PX_1X_2$, $PY_1Y_2$, have the value $0$. Hence $\mathsf L(N'_{22})=\frac14\bigl(\frac{20}3-20\bigr)=-\frac{10}3$.
\item $N'_{12}=-2\bigl[a\,(k'+l')\cdot\mathsf s-\bigl(k'\cdot(k'+l')\bigr)(k\cdot\mathsf s)\bigr]$. Since $(k'+l')\cdot\mathsf s=ab-ff'$, the first term is $a(ab-ff')=\frac18(A+A'-P)\,\mathsf D$, with the value $\frac18\cdot(-\frac43)=-\frac16$. Next, $k'\cdot(k'+l')=A'+e'=\frac12(3A'+B')$ and
\[
k\cdot\mathsf s=bA-fe=\tfrac14\bigl[2A(B+B'-P)-(A'+B-P)(A+B)\bigr]=\tfrac14\bigl[AB+2AB'-AA'-A'B-B^2-PA+PB\bigr].
\]
In $(3A'+B')\bigl[AB+2AB'-AA'-A'B-B^2-PA+PB\bigr]$ the monomials with non-zero values are $-3A'^2B$, $-3A'B^2$, $+3PA'B$, $+2AB'^2$ and $-PAB'$; they give $-1-1+3+\frac23-1=\frac23$. The second term, $\frac18(3A'+B')\bigl[\cdots\bigr]$, therefore has the value $\frac18\cdot\frac23=\frac1{12}$, and $\mathsf L(N'_{12})=-2\bigl(-\frac16-\frac1{12}\bigr)=\frac12$.
\item Summing, with the identities at the end of Step~10,
\[
\mathsf L(N')=2\bigl(-\tfrac43\bigr)+2\cdot\tfrac23-\tfrac{10}3+4\cdot\tfrac12=-\tfrac83,\qquad\text{hence}\qquad\mathcal R[F_1]=\tfrac12\bigl(-\tfrac83\bigr)=-\tfrac43 .
\]
\end{itemize}

\emph{Step 12: the sum.} Adding the six values obtained in Steps~2, 3, 5, 6, 8 and~11, we find
\[
\sum_{\Gamma=1}^6\mathcal R[F_\Gamma]=-\tfrac43-\tfrac{13}{12}-\tfrac1{12}+0+\tfrac{13}6-\tfrac13=\frac{-16-13-1+26-4}{12}=-\tfrac23 ,
\]
which completes the proof of \eqref{IV-73-97}.
\end{proof}

\subsection{The value: $\kappa^{\mathrm{lat}}$ and $c_1$}

Adding the six contributions of Proposition~\ref{IV-73-94} and Proposition~\ref{IV-73-96}, we now obtain the value of the constant $\kappa^{\mathrm{lat}}$ of Theorem~\ref{IV-73-71}, and with it the number $c_1$.

\begin{theorem}[The value of $\kappa^{\mathrm{lat}}$]\label{IV-73-98}
The number $\kappa^{\mathrm{lat}}$ of Theorem~\ref{IV-73-71} is
\[
\kappa^{\mathrm{lat}}=\frac3{2\pi^3}.
\]
\end{theorem}

\begin{proof}
In \eqref{IV-73-72}, $\kappa^{\mathrm{lat}}$ is the sum of the numbers $\kappa_\Gamma^{\mathrm{lat}}$ of Lemma~\ref{IV-73-70}. Expressing each of them by \eqref{IV-73-95} and inserting the sum from \eqref{IV-73-97}, we obtain
\[
\kappa^{\mathrm{lat}}=\sum_{\Gamma=1}^6\kappa_\Gamma^{\mathrm{lat}}=-\frac3{32\pi^3}\cdot24\cdot\sum_{\Gamma=1}^6\mathcal R[F_\Gamma]=-\frac3{32\pi^3}\cdot24\cdot\Bigl(-\frac23\Bigr)=\frac{48}{32\pi^3}=\frac3{2\pi^3},
\]
as asserted.
\end{proof}

\begin{definition}[The number $c_1$]\label{IV-73-99}
We put
\[
c_1:=\frac{\kappa^{\mathrm{lat}}}{\mathsf c_2},
\]
where $\kappa^{\mathrm{lat}}$ is the number of Theorem~\ref{IV-73-71} and $\mathsf c_2=\frac{3\dim\mathfrak g}{16\pi^2}=\frac9{16\pi^2}$ is the leading constant fixed in Section~\ref{C-350-1}; here $\dim\mathfrak g=3=\mathsf n$.

This fixes the coefficient used in Section~\ref{C-350-1}; it does not by itself give a cutoff-uniform remainder for the interacting measure.
\end{definition}

\begin{theorem}[The number $c_1$]\label{IV-73-100}
\begin{enumerate}
\item[\textup{(a)}] $c_1=\dfrac8{3\pi}$, and $\mathsf c_2c_1=\dfrac3{2\pi^3}$.
\item[\textup{(b)}] There is a pure constant $C$ with the following property. Let $\mathsf d>0$, $\vartheta\in(0,1)$, $(f,h)\in\mathfrak T_{\mathbb R^3}(\mathsf d;\vartheta)$ and $0<a\le\mathsf d$, the class of test pairs and the functional $\mathcal K_5$ being those of Corollary~\ref{IV-73-73}. Then
\begin{equation}\label{IV-73-101}
\Bigl|\frac{\mathsf P_1^{\mathrm{lat}}(f,h;a)}{\mathsf c_2\,\mathcal K_5(f,h)}-c_1\Bigr|\le C\,\vartheta^{-3}\,\frac a{\mathsf d}\,\lg\frac{\mathsf d}a .
\end{equation}
Consequently, $\lim_{a\to0}\mathsf P_1^{\mathrm{lat}}(f,h;a)/\mathsf c_2\mathcal K_5(f,h)=c_1$ for every such pair.
\item[\textup{(c)}] For all plaquettes $p,q$ of $\mathbb Z^3$, with $n$ ranging over the integers for which $p+o_n,q+o_n\in\mathsf{Pl}_n$,
\[
\mathsf B_1^{\mathrm{lat}}(p,q)=\lim_{n\to\infty}\ \lim_{\beta\to\infty}\ \beta^3\Bigl[\langle a_{p+o_n};a_{q+o_n}\rangle_{n,\beta}-\beta^{-2}\,\mathsf B_0^{[n]}(p+o_n,q+o_n)\Bigr].
\]
The right side is gauge independent: no gauge fixing enters it.
\end{enumerate}
\end{theorem}

\begin{proof}
(a) By Theorem~\ref{IV-73-98},
\[
c_1=\frac3{2\pi^3}\cdot\frac{16\pi^2}9=\frac{48}{18\pi}=\frac8{3\pi},
\]
and $\mathsf c_2c_1=\kappa^{\mathrm{lat}}=\frac3{2\pi^3}$.

(b) Lemma~\ref{lem:test-pair-kernels} on $\mathbb E=\mathbb R^3$, with $q=5$, gives
\[
\mathcal K_5(f,h)\ge10^{-5}\,\vartheta^2\,\mathsf d\,\|f\|_\infty\|h\|_\infty ,
\]
and the right side is positive, since neither $f$ nor $h$ vanishes identically. We divide \eqref{IV-73-74} by $\mathsf c_2\mathcal K_5(f,h)$ and use $\kappa^{\mathrm{lat}}=\mathsf c_2c_1$. The left side of \eqref{IV-73-101} is then at most
\[
C\,\vartheta^{-1}\,a\,\lg(\mathsf d/a)\cdot10^5\,\mathsf c_2^{-1}\,\vartheta^{-2}\,\mathsf d^{-1},
\]
which is the asserted bound, with a pure constant. Since $\frac a{\mathsf d}\lg\frac{\mathsf d}a\to0$ as $a\to0$, the limit follows.

(c) We apply \eqref{IV-73-28} at $\tau_{\mathrm{gf}}=1$ to the plaquettes $p+o_n,q+o_n\in\mathsf{Pl}_n$: for $\beta\ge1$ the expression $\beta^3[\,\cdots]$ differs from $\mathsf B_1^{[n]}(p+o_n,q+o_n;1)$ by at most $C\beta^{-1}$, with a constant $C$ depending on $n$, $p$ and $q$. This gives the inner limit. By \eqref{IV-73-77} with $j=1$ and $\tau_{\mathrm{gf}}=1$, the outer limit equals $\mathsf B_1^{\mathrm{lat}}(p,q;1)$, which is $\mathsf B_1^{\mathrm{lat}}(p,q)$ by \eqref{IV-73-25}. The mean $\langle\cdot\rangle_{n,\beta}$ is taken in the Wilson measure of the box, and $\mathsf B_0^{[n]}=\frac{\mathsf n}2(\Pi^{\Lambda_n})^2$ by \eqref{IV-73-23}, where $\Pi^{\Lambda_n}$ does not depend on $\tau_{\mathrm{gf}}$ (Definition~\ref{IV-73-15}(2)); thus no gauge fixing enters the right side.
\end{proof}

\subsection{What one shell adds; scope}

We close the section with two items. The first, Lemma~\ref{IV-73-102}, concerns a kernel $\mathsf U(u)=\mathsf c\,|u|^{-5}\varphi(u)$ with $\varphi(Lu)=\varphi(u)$ for a real number $L>1$, for instance $\kappa^{\mathrm{lat}}|u|^{-5}$, the case $\mathsf c=\kappa^{\mathrm{lat}}$, $\varphi\equiv1$; it computes the part of the quadratic form $\frac12\iint w(x)\,w(y)\,\mathsf U(x-y)\,dx\,dy$ that comes from the shell $r\le|x-y|<Lr$, which covers the scales from $r$ to $Lr$. The second, Remark~\ref{IV-73-105}, delimits the results of the section.

\begin{lemma}[What one shell adds]\label{IV-73-102}
Let $L>1$ and $\mathsf c\ge0$ be real numbers, and let $\varphi\colon\mathbb R^3\setminus\{0\}\to\mathbb R$ be bounded, Borel and even, with $\varphi(Lu)=\varphi(u)$. Put $\mathsf U(u):=\mathsf c\,|u|^{-5}\varphi(u)$ and, for $r>0$,
\[
\mathsf{Ann}_r:=\{u\in\mathbb R^3:r\le|u|<Lr\}.
\]
\begin{enumerate}
\item[\textup{(a)}] For every $r>0$,
\begin{equation}\label{IV-73-103}
\int_{\mathsf{Ann}_r}\mathsf U(u)\,|u|^2\,du=\mathsf c\int_{S^2}\int_1^L\varphi(t\omega)\,\frac{dt}t\,d\omega ;
\end{equation}
for $\varphi\equiv1$ this equals $4\pi\,\mathsf c\log L$.
\item[\textup{(b)}] Assume moreover that $\int_{\mathsf{Ann}_r}\mathsf U(u)\,u_iu_j\,du$ is a multiple of $\delta_{ij}$; this is the case if $\varphi$ is invariant under the signed permutations of the coordinates. Then for every real $w\in C_c^2(\mathbb R^3)$
\begin{equation}\label{IV-73-104}
\begin{gathered}
\tfrac12\iint_{x-y\in\mathsf{Ann}_r}w(x)\,w(y)\,\mathsf U(x-y)\,dx\,dy=\\
\tfrac12\Bigl(\int_{\mathsf{Ann}_r}\mathsf U\Bigr)\int w^2\\
-\tfrac1{12}\Bigl(\int_{\mathsf{Ann}_r}\mathsf U(u)\,|u|^2\,du\Bigr)\int|\nabla w|^2+\mathrm{Rem},\\
|\mathrm{Rem}|\le\tfrac\pi{24}\,\mathsf c\,\|\varphi\|_\infty\,(Lr)^2\,\|\Delta w\|_2^2 .
\end{gathered}
\end{equation}
\item[\textup{(c)}] If $\mathsf c=\mathsf c_2c_1\mathsf b^{-3}$ with $\mathsf b>0$, and $\varphi\equiv1$, then the middle term on the right of \eqref{IV-73-104} equals
\[
-\frac1{2\pi^2}\,\mathsf b^{-1}\log L\int\bigl|\nabla(w/\mathsf b)\bigr|^2 .
\]
\end{enumerate}
\end{lemma}

\begin{proof}
(a) In polar coordinates the left side of \eqref{IV-73-103} equals
\[
\mathsf c\int_{S^2}\int_r^{Lr}t^{-5}\varphi(t\omega)\,t^2\,t^2\,dt\,d\omega=\mathsf c\int_{S^2}\int_r^{Lr}\varphi(t\omega)\,\frac{dt}t\,d\omega .
\]
For fixed $\omega$ the function $\tau\mapsto\varphi(e^\tau\omega)$ is bounded, Borel and periodic with period $\log L$; its integral over an interval of length $\log L$ is therefore the same for all such intervals, and
\[
\int_r^{Lr}\varphi(t\omega)\,\frac{dt}t=\int_{\log r}^{\log r+\log L}\varphi(e^\tau\omega)\,d\tau=\int_1^L\varphi(t\omega)\,\frac{dt}t .
\]
For $\varphi\equiv1$ the right side of \eqref{IV-73-103} is $\mathsf c\cdot4\pi\cdot\log L$.

(b) If $\varphi$ is invariant under the signed permutations of the coordinates, then so are $\mathsf{Ann}_r$ and $\mathsf U$. Changing the sign of $u_i$ shows that the off-diagonal entries of the matrix $\bigl(\int_{\mathsf{Ann}_r}\mathsf U\,u_iu_j\bigr)_{ij}$ vanish, and permuting the coordinates shows that its diagonal entries are equal.

We turn to the expansion. The substitution $x=y+u$ (Fubini's theorem and translation invariance) turns the left side of \eqref{IV-73-104} into
\[
\tfrac12\int_{\mathsf{Ann}_r}\mathsf U(u)\,\mathsf a(u)\,du,\qquad\mathsf a(u):=\int w(y+u)\,w(y)\,dy .
\]
Since $w(\cdot+u)$ has the transform $e^{ip\cdot u}\hat w(p)$, Plancherel's theorem gives
\[
\mathsf a(u)=\int|\hat w(p)|^2\,e^{ip\cdot u}\,d^\circ p=\int|\hat w(p)|^2\cos(p\cdot u)\,d^\circ p ,
\]
the second equality because $\mathsf a$ is real. For real $\vartheta'$, Taylor's formula with integral remainder and $|\cos^{(4)}|\le1$ give $|\cos\vartheta'-1+\frac12\vartheta'^2|\le\frac1{24}\vartheta'^4$. The functions $w$, $\partial_iw$ and $\Delta w$ belong to $L^2$, with the transforms $\hat w$, $ip_i\hat w$ and $-|p|^2\hat w$, by the rule for the transform of a derivative, applied to $w$ and to $\partial_iw\in C_c^1(\mathbb R^3)$. Plancherel's theorem therefore yields
\[
\int|\hat w|^2\,d^\circ p=\int w^2,\qquad\int|\hat w|^2\,p_ip_j\,d^\circ p=\int\partial_iw\,\partial_jw,\qquad\int|\hat w|^2\,|p|^4\,d^\circ p=\|\Delta w\|_2^2 ,
\]
and, with $\vartheta'=p\cdot u$ and $|p\cdot u|\le|p||u|$,
\[
\Bigl|\mathsf a(u)-\int w^2+\tfrac12\sum_{i,j}u_iu_j\int\partial_iw\,\partial_jw\Bigr|\le\tfrac1{24}\,|u|^4\,\|\Delta w\|_2^2 .
\]
By the hypothesis of (b), taking the trace, $\int_{\mathsf{Ann}_r}\mathsf U\,u_iu_j=\frac13\delta_{ij}\int_{\mathsf{Ann}_r}\mathsf U\,|u|^2$. Inserting the expansion of $\mathsf a$, we obtain the right side of \eqref{IV-73-104}: its first term is $\frac12(\int_{\mathsf{Ann}_r}\mathsf U)\int w^2$, its second term has the coefficient $\frac12\cdot\frac12\cdot\frac13=\frac1{12}$, and the remainder satisfies
\[
|\mathrm{Rem}|\le\tfrac12\cdot\tfrac1{24}\,\|\Delta w\|_2^2\int_{\mathsf{Ann}_r}|\mathsf U|\,|u|^4\,du\le\tfrac1{48}\,\|\Delta w\|_2^2\,\mathsf c\,\|\varphi\|_\infty\,4\pi\int_r^{Lr}t\,dt\le\tfrac1{48}\,\|\Delta w\|_2^2\,\mathsf c\,\|\varphi\|_\infty\,4\pi\,\frac{(Lr)^2}2 ,
\]
which is the bound in \eqref{IV-73-104}.

(c) By (a) and Theorem~\ref{IV-73-100}(a), the middle term equals
\[
-\tfrac1{12}\cdot4\pi\cdot\frac3{2\pi^3}\,\mathsf b^{-3}\log L\int|\nabla w|^2=-\frac1{2\pi^2}\,\mathsf b^{-3}\log L\int|\nabla w|^2 ,
\]
and $\mathsf b^{-3}\int|\nabla w|^2=\mathsf b^{-1}\int|\nabla(w/\mathsf b)|^2$.
\end{proof}

\begin{remark}[Scope of the results]\label{IV-73-105}
The results of this section are delimited as follows.
\begin{enumerate}
\item[(i)] We make no statement about the cumulants of the smeared plaquette observables $\mathcal O_a(f)$ in Wilson measures on tori. In particular, it is not asserted that $\mathsf B_1^{\mathrm{lat}}$ is an expansion coefficient of these cumulants when $\beta=\mathsf b/a$ and the side of the torus, measured in lattice spacings, grow jointly as $a\to0$.
\item[(ii)] The constant $C$ of \eqref{IV-73-28} is not controlled as a function of $n$. Accordingly, the limits in Theorem~\ref{IV-73-100}(c) are iterated limits, to be taken in the order written there: first $\beta\to\infty$, then $n\to\infty$.
\item[(iii)] Corollary~\ref{IV-73-73} concerns test pairs, whose supports are at distance at least $\mathsf d$ from each other; it gives no information on the terms with $p=q$ of the double sum \eqref{IV-73-25}, on constant functions $f,h$, or on functions $f,h$ with intersecting supports.
\item[(iv)] The group is $SU(2)$ throughout; nothing is asserted for $N>2$.
\item[(v)] Of the term of degree $-6$ in the large-distance behaviour of $\mathsf B_1^{\mathrm{lat}}$, we prove only the bound \eqref{IV-73-72}.
\end{enumerate}
\end{remark}

\section{Construction estimates used in the limit arguments}\label{sec:upstream-interfaces}

We use one assignment of the construction constants throughout each
argument. The finite renormalisation theorem is \cite[Theorem~D]{YMM1};
the sourced effective measures are those of \cite[Theorem~AJ]{YMM1}.
The factorisation into local source energies and a remaining partition
function, with the corresponding stability bounds, is
\cite[Theorem~AD]{YMM1}. In particular, the latter bound applies after
extracting the local energies and the vacuum constant. It is not a
uniform bound on the complete sum of local energies.

\subsection{Periodic lifts and chessboard estimates}\label{sec:periodic-chessboard}

The estimates for a periodic lift follow from the finite-torus
estimates by pulling back each observable. We record this passage
because the same cell functions will be used at several volumes.
Fix a cutoff $K$, a coupling $\beta\ge0$, and a board $(i,o)$ as
defined in Section~\ref{C-355-1}; write $\mathsf l=a_i$.
Let $\mathcal A(\square_{\mathbf k})$ be the bounded Borel functions
of the bonds whose endpoints lie in $\square_{\mathbf k}$.
For an allowed side $\ell=2L^m\ge2\mathsf l$, let
$\mathsf{per}_\ell$ send a torus configuration to its periodic lift,
and set
\[
 \tilde\mu_{K,\ell,\beta}(F)
   =\langle F\circ\mathsf{per}_\ell\rangle_{K,\ell,\beta}.
\]
A period box of cell indices is
$\mathcal K_\ell(\mathbf k_0)
=\mathbf k_0+\{0,\ldots,\ell/\mathsf l-1\}^3$.
Its projection onto the torus cells is a bijection. Since
$\ell/\mathsf l$ is even, projection also preserves the parity
used to conjugate mirror copies. Thus the chessboard seminorm
defined earlier is independent of this period box, and
\begin{equation}\label{eq:period-box-seminorm}
 \|F\|_{{\rm ch},\ell}^{(\ell/\mathsf l)^3}
 =\tilde\mu_{K,\ell,\beta}
       \left(\prod_{\mathbf k\in\mathcal K_\ell(\mathbf k_0)}
                    F^{[\mathbf k]}\right)\ge0.
\end{equation}
Here an action-density observable with its weight supported in a cell
interior is a function of that cell's inner bonds. Indeed, in each
plaquette direction its barycentre and the parallel walls have
coordinates in \(a_K(\mathbb Z+\tfrac12)\). A barycentre strictly
between two walls is at least \(a_K\) from either wall; moving by
\(\pm a_K/2\) to a corner stays inside. In the remaining direction
all four corners have the barycentre's coordinate. Thus every
plaquette contributing to the observable has all four bonds in
the cell.

\begin{corollary}[Chessboard estimates for periodic lifts]\label{IV-18-18}
Write $\tilde\mu=\tilde\mu_{K,\ell,\beta}$.
\begin{enumerate}
\item[(a)] If $\theta$ is reflection in the bond mirror $x_\mu=c$,
$c\in a_K(\mathbb Z+\frac12)$, then
\[
 (F,F')_\theta=\tilde\mu(\overline{\theta F}F')
\]
is a positive semidefinite Hermitian form on bounded local Borel
functions of bonds in $c<x_\mu<c+\ell/2$. In particular it satisfies
the Cauchy--Schwarz inequality.
\item[(b)] For a subset $J$ of a period box and
$F_{\mathbf k}\in\mathcal A(\square_{\mathbf k})$,
\[
 \left|\tilde\mu\left(\prod_{\mathbf k\in J}F_{\mathbf k}\right)\right|
 \le\prod_{\mathbf k\in J}\|F_{\mathbf k}\|_{{\rm ch},\ell}.
\]
\item[(c)] On each cell, $\|\cdot\|_{{\rm ch},\ell}$ is a seminorm,
is unchanged by conjugation or a mirror copy, and obeys
\[
 \|1\|_{{\rm ch},\ell}=1,\qquad
 |\tilde\mu(F)|\le\|F\|_{{\rm ch},\ell}\le\sup|F|.
\]
\item[(d)] Let $f_1,\ldots,f_n$ be bounded real sources supported
in the interior of one cell, $z\in\mathbb C^n$, $v=\sum_qz_qf_q$, and let
$v^{\langle\ell\rangle}$ denote periodisation. With the usual
action-density observable $\mathcal O_K(v)$ and chessboard
continuation, including conjugation in odd cells,
\[
 \bigl\|e^{\mathcal O_K(v)}\bigr\|_{{\rm ch},\ell}^{(\ell/\mathsf l)^3}
 =\frac{Z_{K,\ell,\beta}[(v^{\langle\ell\rangle})^{\rm ch}]}
        {Z_{K,\ell,\beta}}\ge0.
\]
\item[(e)] If $\ell\ge\ell'\ge2\mathsf l$ are allowed sides and
$p_\ell=\ell^{-3}\log Z_{K,\ell,\beta}$, then
\[
 p_\ell\le p_{\ell'},\qquad
 \|F\|_{{\rm ch},\ell}
 \le e^{\mathsf l^3(p_{\ell'}-p_\ell)}
          \|F\|_{{\rm ch},\ell'}.
\]
\end{enumerate}
\end{corollary}

\begin{proof}
Pullback by $\mathsf{per}_\ell$ preserves products, conjugation
and reflection. A function supported in the open slab in (a)
becomes a function of bonds in the corresponding torus half.
Bond reflection positivity on the torus therefore gives (a).
For (b), pull back each factor and apply
Proposition~\YMMref{86.3}{681}; fill the unused cells with $1$.
Proposition~\YMMref{86.4}{682} gives (c), since pullback preserves
the supremum of a function of the bonds inside a cell.

Periodisation is linear and the source vanishes outside the
cell interior. Hence
$\mathcal O_K(v)\circ\mathsf{per}_\ell
=\mathcal O(v^{\langle\ell\rangle})$.
Multiplication of all mirror copies of its exponential adds the
reflected, conjugated sources in the exponent. Integration gives
the partition-function ratio in (d); its nonnegativity is the
chessboard positivity already used to define the seminorm.
Finally the cell observable is represented by the same finite
set of bond variables at both sides. Corollary~\YMMref{86.6}{683}
applied to this common function gives (e), including pressure
monotonicity.
\end{proof}


\subsection{The chart integral of a small-field step}
\label{sec:sf-chart-interface}

The fluctuation variables describe a fibre of the averaging map. We
record their domain and the resulting measure identity, including the
normalising factors. This will also specify what is meant by the chart
integrals used in the chains below.

Work at one step $k\to k+1$, with $K,m\ge1$ and $0\le k<K$.
Let $\Gamma$ be the positively oriented coarse bonds, $\mathfrak X_0$
the fine bonds and $\mathfrak X$ the nonpivot fine bonds. Write $N_j$
for the number of sites at level $j$. The chart and its pivot coefficients
are those of Definition~\YMMref{51.3}{391} and
Proposition~\YMMref{53.2}{441}. Fix admissible constants
$0<\varpi_\chi<\omega_M$, $0<r_\chi\le1$ and $C_\chi\ge1$
as in that proposition. These are chosen with the local bounds
$|(L\mathsf m_c)^{-1}|\le2L^2$ and
$|\mathsf{Jb}_c|\le\tfrac12$ used in its proof.
Their values depend only on $L$. We retain the compatible choices
\[
 \varpi_1=2\varpi_\chi,\qquad
 \varpi_\chi+8r_\chi\le\varpi_1\le\omega_{\rm p},
\]
with $\varpi_1$ below the small-curvature ceiling in that proof. These are obtained by decreasing $\varpi_\chi$ and then $r_\chi$; they ensure that the chart field has plaquette defect below $\varpi_1$ when $C_\chi\sup|\mathcal A|<r_\chi$.

\begin{definition}[Chart domain and thresholds]\label{def:sf-chart-data}
Let $p(W)=V_k^{\rm bf}(W)$ be the level-$k$ average of a real regular
critical background over $W$, chosen in the block axial gauge. It is
independent of the axial representative. We use it on its stable domain,
that is, on the open set where it is defined continuously, and put
\[
 O^\diamond=\mathrm{dom}^{\circ}p,\qquad
 O^\diamond(v)=\{W\in O^\diamond:
          \max_q|\partial p(W)(q)-1|<v\},\qquad
 O^{\rm im}=O^\diamond(\varpi_\chi).
\]
Here the maximum is over the fine plaquettes. For a configuration $U$
on the bare lattice, its regularity number at level $j$ is
$\mathsf m_j(U)=L^{2j}\max_q|\partial U(q)-1|$, where this
maximum is over bare plaquettes. If the minimal background orbit
over a level-$j$ field $W$ exists, put $\mu_j(W)=\mathsf m_j(U)$
for any $U$ in that orbit. This number is independent of the
representative: a real gauge conjugates each plaquette holonomy
by a unitary matrix and hence preserves its defect norm.

For the double-block thresholds put
\[
 r_{\mathsf y}=\min\{r_*/20,(1600L^2)^{-1}\},\quad
 \varpi_{\rm db}=\frac{r_{\mathsf y}^2}{8\pi L^3},\quad
 r_{\rm db}=\frac{r_{\mathsf y}^2}{64L^2},\quad C_{\rm db}=8L^3,
\]
where $r_*$ is the radius of Lemma~\YMMref{14.7}{66}. With
$c_2^{\rm SO}$ as in Proposition~\YMMref{14.1}{61} and
$\omega_{\rm p}$ as in Proposition~\YMMref{14.11}{75}, set
\[
 a_{\rm sf}=\min\{c_2^{\rm SO},\tfrac12L^2\varpi_{\rm db},
                       \tfrac14L^2\omega_{\rm p}\},\quad
 \delta_{\rm sf}=\min\{r_{\rm db}/C_{\rm db},\omega_{\rm p}/2\},
\]
and
\[
 \varpi_{\rm im}=\varpi_\chi,\quad
 \delta_{\rm im}=\min\{\delta_{\rm sf},r_\chi/C_\chi\},\quad
 a_{\rm im}=\min\{a_{\rm sf},L^2\varpi_\chi/4\},\quad
 \mathsf c_{\rm im}=C_\chi.
\]
At $\mathsf t>1$ write
\[
 g=\mathsf t^{-1/2},\quad
 \varepsilon(\mathsf t)=A_0\mathsf t^{-1/2}(\log\mathsf t)^{\mathfrak p_0},
 \quad\delta(\mathsf t)=\frac{A_1}{A_0}\varepsilon(\mathsf t),\quad
 \mathfrak b(\mathsf t)=g^{-1}\delta(\mathsf t).
\]
The \emph{data under (Im)} consist of one such lattice and the construction
constants with
\begin{equation}\label{eq:sf-Im-conditions}
 A_0/A_1\ge C_\chi,\qquad
 \varepsilon(\mathsf t)\le1,\qquad
 \delta(\mathsf t)\le\delta_{\rm im}.
\end{equation}
The additional condition $b_0\le a_{\rm im}$, when imposed on a
background radius $b_0$, is separate from (Im).
\end{definition}

These chart constants can be used in the common construction choice
of Proposition~\YMMref{69.1}{595}. Fix the chosen admissible
$(\varpi_\chi,r_\chi,C_\chi)$ among the constants depending only on
$L$. The subsequent choice of $A_0$ retains $A_0\ge C_\chi A_1$.
After all these constants have been fixed, the final choice of
$\gamma$ makes $\varepsilon(x)\le1$ and
$\delta(x)\le\min\{\delta_{\rm sf},r_\chi/C_\chi\}$ for every
$x\ge\gamma^{-2}$, since both functions tend to zero as
$x\to\infty$. Thus the same chart constants satisfy (Im) at
every step of the admitted construction.

For $W\in O^{\rm im}$, the chart and its scaled form are
\[
 V_W[\mathcal A]=e^{i\psi_{p(W)}^{\rm lin}(\mathsf C_{p(W)}\mathcal A)}p(W),
 \qquad V_W(B)=V_W[gB],\qquad
 \mathsf{Rg}_{\mathsf t}=\{B:|B(b)|<\mathfrak b(\mathsf t)
                              \text{ for every }b\in\mathfrak X^+\}.
\]
The set $\mathsf{Rg}_{\mathsf t}$ is a product of Euclidean balls in
$\mathfrak g\cong\mathbb R^3$. The small-field test
$\chi^{\rm sf}_{\mathsf t}(V)$ equals one when $p(\mathcal MV)$ is
stably defined and all nonpivot intrinsic logarithms
$i^{-1}\log(V(b)p(\mathcal MV)(b)^{-1})$ have norm less than
$\delta(\mathsf t)$; otherwise it equals zero. There is no pivot
cutoff. Let $\mathfrak G_k[\mathsf t]$ be the block gauge factor of
Section~\YMMref{52.3}{418}, and $G_W(B)$ its nonnegative gauge action
at $V_W(B)$. With $\mathsf J_W(B)$ the logarithmic Jacobian in
equation~\YMMref{(985)}{442}, put
\begin{equation}\label{eq:sf-chart-integral}
 \mathsf I[F](W)=\mathsf{pf}_k(\mathsf t)\,
 z^{(k)}_{p(W)}(\Gamma)
 \int_{\mathsf{Rg}_{\mathsf t}}
      e^{-\mathsf tG_W(B)+\mathsf J_W(B)}F(V_W(B))\,dB,
\end{equation}
where
\[
 z^{(k)}_p(\Gamma)=\prod_{c\in\Gamma}|\det\mathsf m_c(p)|^{-1},
 \qquad
 \mathsf{pf}_k(\mathsf t)=
 \mathsf z(\mathsf t,\varepsilon(\mathsf t))^{-(N_k-N_{k+1})}
 (\sigma_{\rm H}\mathsf t^{-3/2})^{\#\mathfrak X^+}L^{-3\#\Gamma},
 \qquad \sigma_{\rm H}=(2\pi^2)^{-1}.
\]
For a real symmetric positive definite form $\mathsf A^\circ$ on the
nonpivot fields we also use the Gaussian normaliser
\[
 Z^{(k)}(W;\mathsf A^\circ)
 =z^{(k)}_{p(W)}(\Gamma)
       \int e^{-\frac12(B,\mathsf A^\circ B)_1}\,dB,
\]
where this integral is over the full real field space. The chart
integral, by contrast, retains the field box.

\begin{lemma}[The chart and its measure]\label{lem:sf-chart-integral}
Assume (Im). For $W\in O^{\rm im}$ and
$B\in\mathsf{Rg}_{\mathsf t}$ the chart is real, belongs to the fibre
$\mathcal MV_W(B)=W$, and has nonpivot intrinsic logarithms $gB$.
It passes the small-field test and all the gauge cutoffs. Its factors obey
\begin{equation}\label{eq:sf-chart-factor-bounds}
 0<z^{(k)}_{p(W)}(\Gamma)\le(8L^9)^{\#\Gamma},\qquad
 0<e^{\mathsf J_W(B)}\le(27/8)^{\#\Gamma}.
\end{equation}
The chart, its Jacobian, its gauge action and its elimination factor
are continuous in their real arguments on this domain.
If $B\mapsto F(V_W(B))$ is bounded and Borel, the integral
\eqref{eq:sf-chart-integral} is absolutely convergent. For bounded Borel
$F$ on the configuration space it is a bounded Borel density, on
$O^{\rm im}$, of
\[
 \mathcal M_*\bigl(\mathfrak G_k[\mathsf t]\,
                 \chi^{\rm sf}_{\mathsf t}F\,dV\bigr).
\]
It is continuous on any subset on which $F(V_W(B))$ is continuous in
$W$ for each $B$, with a bound uniform in $W$ and $B$.
\end{lemma}

\begin{proof}
The definitions give
\[
 C_\chi g\mathfrak b(\mathsf t)=C_\chi\delta(\mathsf t)\le r_\chi,
 \qquad C_\chi\delta(\mathsf t)\le\varepsilon(\mathsf t).
\]
The strict field inequalities therefore put the chart in
Proposition~\YMMref{53.2}{441}. Its fibre identity and converse give
precisely the configurations selected by the nonpivot test. The contour
bound following that proposition is
$|V_W(B)(y,x)-1|\le C_\chi g\max_{b\subset B(y)}|B(b)|
 <\varepsilon(\mathsf t)$, so every gauge indicator is one.

The pivot space has real dimension three. Thus
$|(L\mathsf m_c)^{-1}|\le2L^2$ gives
$0<|\det\mathsf m_c|^{-1}\le L^3(2L^2)^3=8L^9$.
In the Jacobian formula each positive determinant
$\det(1-\mathsf{Jb}_c)$ is at most $(3/2)^3=27/8$;
the remaining factors are $(\sin|B'(b)|/|B'(b)|)^2\le1$.
This proves \eqref{eq:sf-chart-factor-bounds}.
Since $G_W\ge0$ and the field box has finite volume, these bounds
give absolute convergence and a uniform integrable majorant.

Apply equation~\YMMref{(986)}{442} with the exterior empty and
reference $p(W)$, first to a nonnegative integrand with $W$ restricted
to an arbitrary Borel subset of $O^{\rm im}$. The converse ensures
that this change of variables accounts for the entire tested fibre.
The pivot factors are
$\exp(\sum_c\mathsf z_c)=L^{-3\#\Gamma}z^{(k)}_{p(W)}(\Gamma)$,
and the gauge factor is
$\mathsf z(\mathsf t,\varepsilon)^{-(N_k-N_{k+1})}
 e^{-\mathsf tG_W(B)}$.
Their product gives exactly \eqref{eq:sf-chart-integral}.
Tonelli, followed by linearity, gives the bounded complex Borel
version. The same fixed field box and the majorant just proved give
the stated continuity by dominated convergence. This is also the
whole-lattice density calculation in equation~\YMMref{(912)}{421}.
\end{proof}

\subsection{Regularity classes and descent of a background}
\label{sec:sf-class-window}

We shall use three classes of block fields. The first records only the
curvature of the critical background. The other two require that the
background can be used in the local analytic functions of a step. The
small-curvature window identifies these classes without imposing a new
condition on the configurations.

Retain the geometric and numerical setting of
Section~\YMMref{52.7}{432}, in particular the admitted lattice sizes,
the choice of $B^{(\mathrm{cg})}$, and the uniform number
$\varepsilon_{\mathcal W}$ of Proposition~\YMMref{52.3}{433}.
For $j\ge1$ let $\mathcal C_j^\flat(e)$ consist of the block fields
$W$ whose minimal regular orbit exists and satisfies
$\mathsf m_j(U)<e$. Let $\mathcal C_j(b_0,b_1;e)$ be the subclass
whose pairs on every domain belong to the cores with radii $(b_0,b_1)$,
for every representative of the orbit, as in that section. Write
$\mathcal C_j^\circ(b_0,b_1;e)$ when the actual plain triples
$(U,0,J^{(j)}(U))$ belong to those cores. Thus
\[
 \mathcal C_j^\circ(b_0,b_1;e)
 \subset\mathcal C_j(b_0,b_1;e)\subset\mathcal C_j^\flat(e).
\]

\begin{lemma}[The window and descent]\label{lem:sf-class-window}
On lattices admitted at the largest level under consideration the
following assertions hold.
\begin{enumerate}
\item For $0<b_0,b_1\le1$, $0<e\le c_2^{\rm SO}$ and
$2e\le\varepsilon_{\mathcal W}b_0$,
\begin{equation}\label{eq:sf-class-window}
 \mathcal C_j^\flat(e)=\mathcal C_j(b_0,b_1;e)
                      =\mathcal C_j^\circ(b_0,b_1;e).
\end{equation}
Every background in this class has scaled curvature and current less
than $b_0$. If also $\mathsf t>1$,
$\varepsilon(\mathsf t)\le1$ and
$\delta(\mathsf t)\le\varepsilon_{\mathcal W}$, this class is
contained in the actual step domain
$\mathcal W_j(\mathsf t;b_0,b_1)$.
\item Let $1\le j\le j'$, $0<e,e'\le1$ and
$L^{-2(j'-j)}e'\le e$. If $W'\in\mathcal C_{j'}^\flat(e')$
and $U$ is in its critical orbit, then the critical orbit over
$\mathcal M^jU$ exists, contains $U$, and
$\mathcal M^jU\in\mathcal C_j^\flat(e)$.
\item Every isometry of the level-$j$ lattice carries
$\mathcal C_j^\flat(e)$ into itself and carries its critical orbits
to the corresponding critical orbits. The same holds for the two
classes in the window \eqref{eq:sf-class-window}.
\end{enumerate}
\end{lemma}

\begin{proof}
For the first assertion take $W\in\mathcal C_j^\flat(e)$ and a
representative $U$ of its critical orbit. The small-curvature averaging
bound of Proposition~\YMMref{14.1}{61} gives
\[
 \varpi(W)=\varpi(\mathcal M^jU)<2e
                                  \le\varepsilon_{\mathcal W}b_0.
\]
The first assertion of Proposition~\YMMref{52.3}{433}, and the
plain-core verification in its proof, put $(U,0,J^{(j)}(U))$ in
every required core. This argument uses the small-plaquette assertion
of that proposition; the separate characterization under
$2b_0\le\varepsilon_{\mathcal W}$ is not needed. The two inclusions
preceding the lemma give equality. On the whole torus the core
conditions bound $\mathsf m_j(U)$ and $|J^{(j)}(U)|$ by $b_0$.
The same small-plaquette assertion gives the step domain when the
conditions on $\mathsf t$ hold.

For descent, restrict the constant determining sequence to depth $j$.
Lemma~\YMMref{51.13}{408} preserves criticality and regularity.
Uniqueness of the regular critical orbit, Theorem~\YMMref{E}{10},
therefore identifies $U$ with a member of the orbit over
$\mathcal M^jU$. Its regularity number is
\[
 \mathsf m_j(U)=L^{-2(j'-j)}\mathsf m_{j'}(U)<e.
\]
All representatives have the same number, since a real gauge
conjugates each plaquette holonomy by a unitary matrix. This proves
the second assertion.

Finally an isometry transports the constant determining sequence,
the paths defining its averages, and the action to their counterparts.
It preserves regularity and criticality. Applying it and its inverse
therefore carries the full critical orbit to the full critical orbit.
It permutes plaquettes, possibly reversing them; conjugation and
inversion preserve $|\partial U-1|$. Thus it preserves
$\mathsf m_j$ and $\mathcal C_j^\flat(e)$. The equality in the
first assertion gives the last claim.
\end{proof}

\subsection{Measurability of the orbit-defined array density}
\label{sec:array-borel-density}

The chart formula gives a density once its input is a bounded
Borel function on the whole configuration space. For an array
action, this is not immediate from continuity on each analytic
core: the minimal background may fail to exist at other block
fields. We verify the assertion with the function set equal to
zero wherever that background is undefined.

Retain the geometric smallness condition
\begin{equation}\label{eq:array-borel-radius}
             \varepsilon_{\rm v}\le\frac1{4L^2}.
\end{equation}
All configuration spaces in this argument are finite products
of $SU(2)$, with their compact metric topology. Averages are
the real centered averages. Write
$\varpi(V)=\max_p|\partial V(p)-\mathbf1|$ and
$\omega_M=(200L^2)^{-1}$.

\begin{lemma}[Regular averaging is a submersion]
\label{lem:regular-average-submersion}
For $0\le j<k<K$, the one-step average is $C^1$ with
surjective differential on
$\mathcal S_j=\{V:\varpi(V)<\omega_M\}$.
If a bare configuration is regular at radius
$\varepsilon_{\rm v}$ in level-$k$ units, then all inputs to
its first $k$ averaging steps lie in these domains. Their composite is $C^1$
with surjective differential on an open neighbourhood of that
configuration.
\end{lemma}

\begin{proof}
We first record the estimates which keep averaging in its
small-curvature branch. Put $h=(L-1)/2$ and
$\lambda=L^2\varpi(V)<1/200$. Exchanging two consecutive
path segments across a plaquette changes their transports by
at most $\varpi(V)$; reversing a consecutive bond pair changes
nothing. Thus the six coordinate paths inside one block have
diameter at most $3h^2\varpi(V)\le3\lambda/4$.
If a family of unitary matrices has diameter $d<1/40$, its
mean $\bar V$ satisfies
\[
 |\log(V_i\bar V^{-1})|\le7d/6.
\]
Indeed Proposition~\YMMref{4.3}{18} supplies the unique mean
and an entry within $1/10$ of it. The diameter bound puts
all entries within $1/8$, and their logarithms within $1/7$.
For antihermitian $X,Y$,
$|e^Xe^{-Y}-\mathbf1-(X-Y)|\le
(|X|+|Y|)|X-Y|/2$; average this inequality with
$\sum_i\log(V_i\bar V^{-1})=0$. The exponential inequality
follows by integrating the derivative of $e^{sX}e^{-sY}$.
The contour mean therefore differs from any coordinate path
by at most $7\lambda/8$.

For a coarse bond $c$ and a site $x$ in its initial block,
let $H_x(c)$ be the path from $c_-$ to $x$, straight across
to $x+Le_\nu$, and back to $c_+$. Put
$W_x=H_x(c)V[c]^{-1}$. Moving the straight path past the
two transverse coordinate pieces costs at most
$L(L-1)\varpi(V)\le\lambda$. The two contour-mean errors give
\[
 |W_x-\mathbf1|\le11\lambda/4<1/70,
 \qquad \mathcal M(V)(c)=e^{\Sigma_c}V[c],
 \qquad \Sigma_c=L^{-3}\sum_x\log W_x.
\]
In particular $|\log W_x|,|\Sigma_c|\le2.8\lambda$.
Set $l_x(c)=\log(H_x(c)\mathcal M(V)(c)^{-1})$.
For antihermitian $Z_i$ with $s=\sum_i|Z_i|\le1/4$,
the exponential and logarithm series give
\[
 \left|\log\prod_i e^{Z_i}-\sum_i Z_i\right|
       \le\tfrac32s^2.
\]
To see the constant, the terms of total degree at least two
in the product have norm at most $e^s-1-s\le0.65s^2$;
the remainder in its logarithm is at most
$s^2/[2(1-s)]\le2s^2/3$. Applied to
$e^{\log W_x}e^{-\Sigma_c}$ this gives
\[
 |l_x(c)|\le2.1(2.8\lambda),\qquad
 \left|L^{-3}\sum_xl_x(c)\right|\le6(2.8\lambda)^2.
\]

For a coarse plaquette, take the four paths $H_x$ with a
common offset in their four blocks. Their contour factors
cancel, leaving a conjugate of an $L$-by-$L$ plaquette loop.
Its distance from the identity is at most $\lambda$ and its
logarithm has norm at most $\lambda+0.51\lambda^2$.
The coarse plaquette is the product of this loop and four
transported factors $e^{\pm l_x}$. Their logarithmic product
has total input size at most $24.6\lambda\le1/4$ and
remainder at most $908\lambda^2$. Average over the common
offset. The four linear $l_x$ averages each cost at most
$47.1\lambda^2$, so
\begin{equation}\label{eq:array-borel-average-bound}
 \varpi(\mathcal MV)\le\lambda(1+1100\lambda).
\end{equation}

Let $U$ be regular and put
$\theta=L^{2k}\varpi(U)<\varepsilon_{\rm v}$.
Induction using \eqref{eq:array-borel-average-bound} gives
\[
 \varpi(\mathcal M^jU)\le2\theta L^{-2(k-j)}
      <\tfrac12L^{-4}<\omega_M\quad(0\le j<k).
\]
In the induction the accumulated multiplicative error is at
most
$\exp(550\sum_{n\ge2}L^{-2n})<e^{0.02}<2$.
This proves that every participating average uses the displayed
small-curvature formula.

That formula is $C^1$. For completeness, the equation defining
a contour mean in exponential coordinates is
$\sum_i\log(e^{iA_i}V_i\bar V^{-1}e^{-Z})=0$.
At $A=Z=0$, its derivative in $Z$ differs from $-nI$ by
at most $0.07n$: the logarithms have norm less than $0.03$,
and the derivative remainder of $\log T$ is bounded by
$r|\dot T|/(1-r)$ when $|T-I|\le r$.
The implicit function theorem makes the mean $C^1$.
Products, logarithms and exponentials then give the same
property for $\mathcal M$.

To prove surjectivity, vary only the $L^2$ corridor bonds
between the blocks of one coarse bond $c$. A corridor bond
belongs to no other positively oriented coarse-bond corridor.
Transport a fixed $X\in\mathfrak g$ to each such bond along
one coordinate path and use these values as the variation.
The other coarse bonds do not change. For the selected bond,
the transported variation $X_x$ satisfies
\[
 |X_x-X|\le7\lambda|X|/4,
 \quad W_x(t)=e^{itX_x}W_xe^{-itX},
 \quad \mathcal M(V_t)(c)\mathcal M(V)(c)^{-1}
             =e^{\Sigma(t)}e^{itX}e^{-\Sigma(0)}.
\]
Here $|\dot W_x(0)|\le7.25\lambda|X|$,
$|\dot\Sigma(0)|\le0.0369|X|$ and $|\Sigma(0)|\le0.014$.
Differentiation of the last product therefore differs from
$iX$ by at most
$(e^{0.014}0.0369+2(0.014))|X|\le0.066|X|$.
This is an invertible map on the three-dimensional Lie algebra.
Doing it separately at each coarse bond proves surjectivity.
The preimages of the open sets $\mathcal S_j$ are open;
composition and the chain rule finish the proof.
\end{proof}

\begin{lemma}[Bounded Borel array densities]
\label{lem:array-borel-density}
Under \eqref{eq:array-borel-radius}, let $S$ be the action
of a bare array at an admitted place of height $k$ and let
$\operatorname{Re}x>0$. The function
\[
 F_{x,S}(V)=
 \begin{cases}
 e^{-xA^{\eta_k}(U_k(V))+S(U_k(V))},&
 \text{if the minimal orbit exists and $S$ is defined on it},\\
 0,&\text{otherwise},
 \end{cases}
\]
is bounded and Borel on $\mathcal V_k$. So is its product
with the small-field cutoff.
\end{lemma}

\begin{proof}
The action domain $D\subset\mathcal V_0$ is open.
Indeed the pair domain of each local member is open by
Lemma~\ref{lem:analytic-open-pairs}, and the pair of a real
configuration is a continuous polynomial in its entries and
their conjugates. Only finitely many members occur. If a
required trivial pair is absent, $D$ is empty; otherwise
it is the intersection of these finitely many open preimages.
The action is continuous and gauge invariant there.
It is bounded: local terms have bounded family norms, and
each real plaquette action is in $[0,2]$.

For $k=0$, the array is empty and the background is the
configuration itself precisely on the open regularity set.
The claimed function is $e^{-xA}$ there and zero elsewhere,
with modulus at most one. Suppose henceforth that $k\ge1$.
Let $R$ be the open set of regular bare configurations, and
$C\subset R$ its configurations critical on their averaging
fibre. We show that $C$ is relatively closed in $R$.
By Lemma~\ref{lem:regular-average-submersion}, in local
exponential coordinates the average is a $C^1$ submersion
$f$, and the action is a smooth function $a$.
The implicit function theorem identifies tangent vectors to
a fibre with $\ker Df$. Thus criticality is equivalent to
$Da\,\xi=0$ for every $\xi\in\ker Df$.
If critical points $A_n$ tend to $0$, write $T_n=Df(A_n)$.
For $\xi\in\ker Df(0)$, the vectors
\[
 \xi_n=\xi-T_n^{\mathsf T}(T_nT_n^{\mathsf T})^{-1}T_n\xi
\]
lie in $\ker T_n$ and tend to $\xi$. Hence
$Da(0)\xi=\lim_nDa(A_n)\xi_n=0$. This proves relative
closedness. Residual gauge transformations preserve $C$ and
each fibre, by gauge covariance of the averages and invariance
of the action.

A relatively closed subset of an open subset of a compact
metric space is a countable union of compact sets.
Continuous images of such sets are again countable unions
of compact sets, hence Borel. Apply this first to $C$ and
the average $M=\mathcal M^k|_R$. Its image is the set of
block fields with a nonempty regular critical fibre.
The residual gauge relation
$\{(U,U^u):U\in\mathcal V_0,\ u\in\mathcal G^{\rm res}\}$
is closed, being the continuous image of a compact space.
Consequently
\[
 T=\{(U,U')\in C\times C:MU=MU',\quad
                  U'\ne U^u\text{ for every residual }u\}
\]
is also a countable union of compact sets: equality of
averages is relatively closed and the complement of the
gauge relation is open. Its image under $(U,U')\mapsto MU$
is the set of fibres containing more than one residual orbit.
Subtracting this image from $M(C)$ gives the Borel set
$\mathrm{Ex}$ where the minimal background exists.

Since $D$ is gauge invariant, the orbit over
$V\in\mathrm{Ex}$ lies in $D$ exactly when it meets $D$.
Thus the required base domain is
$\mathrm{Ex}\cap M(C\cap D)$, again Borel.
For any continuous orbit-invariant function $q$ on $D$ and
closed set $B\subset\mathbb C$, the set where its orbit
value lies in $B$ is
$\mathrm{Ex}\cap M(C\cap q^{-1}(B))$.
The same compact-union argument makes this Borel, since
$q^{-1}(B)$ is relatively closed in $D$.
Take $q=-xA^{\eta_k}+S$ and then exponentiate.
Extension by zero gives the stated Borel function, and
\[
 |F_{x,S}|\le e^{\sup_D|S|},
\]
because the real action is nonnegative and
$\operatorname{Re}x>0$. The cutoff is a Borel indicator,
so multiplication by it preserves both assertions.
\end{proof}


\subsection{Weighted progressions and support sums}
\label{sec:weighted-progression-interface}

The weighted progression keeps the counterterms fixed and inserts
the source into the action. Its parameter domain is determined by
the room left in the bare-array norm. We recall that domain, so that
the use of several sources does not require an unstated extension
of a two-source result.

\begin{definition}\label{def:weighted-progression-data}
The data of a weighted progression through height $k$ are as follows.
\begin{itemize}
\item[(Wp1)] Use the canonical datum, auxiliary maps, response maps
and constants of Section~\YMMref{122}{825}, with all its
small-coupling hypotheses. In particular use the weighted
old-potential constant $K_\Lambda$, the weights
$\omega_n=(n+1)L^{-(n-1)}$, its ceiling $s^\star_\omega$, and
$K_{\mathbf M}$ evaluated at that weighted constant. The
counterterms are fixed independently of the sources.
\item[(Wp2)] Choose calibrations $t_j\ge\gamma^{-2}$ and
$x_j\in D(t_j,\sigma_\star t_j)$ for $0\le j\le k$, with
the ordinary canonical families $E^{(i)}$ and counterterms $e_i$.
No relation among the calibrations is required at this point.
\item[(Wp3)] Choose one lattice $\mathfrak w=(\ell,K)\in\mathfrak W_k$,
$\ell=2L^m$, so that all these places and their domains exist.
\item[(Wp4)] At each positive height use the bracket localization
and its full hypotheses from Section~\YMMref{122}{825}. At
height at least two its octet property is required on every
lattice in the defining class, not only on $\mathfrak w$.
\item[(Wp5)] At height zero retain
$\alpha_\iota\le a_P^\circ$ for $\iota=0,1$ and
$C_{\mathcal A}\delta(t_0)\le r_P^\circ/2$.
\item[(Wp6)] Choose any finite number $p\ge1$ of plaquette
weights $w_a$, and put $w_\lambda=\sum_a\lambda_aw_a$ for
$\lambda\in\mathbb C^p$.
\end{itemize}
The source weights may be complex. Reality conclusions are made
only for real weights and real running coefficients.
\end{definition}

Here is the parameter room in this notation. Put
\begin{equation}\label{eq:weighted-progression-numbers}
\begin{gathered}
 s_j=(1+\beta_{\rm s})\delta(t_j)/\alpha_0,
 \qquad\theta_j=K_{\mathbf M}s_j,
 \qquad N_{\star,j}=w_\star/(4K_\Lambda s_j),\\
 m_{j,a}=12M^3\alpha_0^2L^{K-j}|w_a|_\infty,
 \qquad c^\flat=81C_P^2C_{\mathcal A}^2/K_\Lambda,\\
 d_{1,a}=c^\flat\theta_0m_{0,a},\qquad
 n_{j,a}=m_{j,a}+\sum_{i=1}^j\omega_{j+1-i}d_{i,a},
 \qquad d_{j+1,a}=\theta_jn_{j,a},\\
 N_j^{\rm rm}=N_{\star,j}-
       \Big(\sum_{n=1}^j\omega_n\Big)
                         (E_0+\tfrac32\alpha_0^2B_\sharp).
\end{gathered}
\end{equation}
The recursion in the last two lines is used for $1\le j\le k$;
the admitted ceiling makes $N_j^{\rm rm}>0$. Define
\begin{equation}\label{eq:weighted-progression-domain}
\begin{split}
 \Lambda^{\rm w}_0
 &=\{\lambda:c^\flat\sum_a|\lambda_a|m_{0,a}<N_{\star,0}\},\\
 \Lambda^{\rm w}_j
 &=\{\lambda\in\Lambda^{\rm w}_{j-1}:
                 \sum_a|\lambda_a|n_{j,a}<N_j^{\rm rm}\}.
\end{split}
\end{equation}
For $j\ge1$ write $\mathsf{Arr}^{[j]}$ for the admitted bare
arrays with slot $\{j\}$ and weighted norm at most $N_{\star,j}$.
The localized step map is $\mathsf{St}_j$. Its Lipschitz and
derivative bounds are equations~\YMMref{(2251)--(2252)}{825}.

For real running coefficients put $\mu_{j,a}=\theta_jm_{j,a}$,
$C_0^\flat=\max(2,c^\flat)$ and
$C_1^\flat=\max(2,(c^\flat)^2/8)$. The convolution constants are
\begin{equation}\label{eq:weighted-running-numbers}
 \mathsf S_A=(4/3)^{1/2}\sum_{n\ge1}
                    (n+1)^{\mathfrak p_0+1}L^{1-n/2},\qquad
 \mathsf S_B=(4/3)^{1/2}\sum_{n\ge0}
                    (n+1)^{\mathfrak p_0}L^{-n/2}.
\end{equation}
For two sources let $D_j$ be the least sup distance between
corners of their supporting plaquettes in level-$j$ units, and put
\begin{equation}\label{eq:weighted-mixed-numbers}
 \nu_j=\frac{32K_{\mathbf M}K_\Lambda}{w_\star}
                    s_j^2m_{j,1}m_{j,2},\qquad
 \mathsf T(u)=\sum_{n\ge1}(n+1)^{2\mathfrak p_0+1}
              e^{-4\beta_{\rm s}\kappa[L^{n-1}u-2]_+}\quad(u>0).
\end{equation}
Set $\mathsf T(0)=+\infty$. These are the constants in the
running and mixed-derivative estimates of
Section~\YMMref{122}{825}. Their series converge: the first two
have geometric decay, and for fixed $u>0$ the last summand has
exponential decay in $L^{n-1}u$. Each summand of $\mathsf T$
decreases with $u$, so $\mathsf T$ does too.

\begin{proposition}[A finite source system]
\label{prop:weighted-finite-system}
At the data of Definition~\ref{def:weighted-progression-data},
the recursion
\begin{equation}\label{eq:weighted-progression-recursion}
\begin{split}
 f_{\lambda,j}
 &=((E^{(i)}_\lambda,e_i)_{i\le j},\emptyset)
                 \sqcup\mathsf W_{L^{K-j}w_\lambda},\\
 E^{(j+1)}_\lambda
 &=E^{(j+1)}+\mathsf{St}_j(f_{\lambda,j})
                         -\mathsf{St}_j(f_{0,j})
\end{split}
\end{equation}
exists on $\Lambda^{\rm w}_j$. At height zero use the bracketed
plaquette action in the fourth place of the same step.
The conclusions of Proposition~\YMMref{122.1}{826} hold for
these $p$ sources: joint holomorphy, the Lipschitz bound
\[
 \|E^{(j)}_\lambda-E^{(j)}_{\lambda'}\|_{j,+}
       \le\sum_a|\lambda_a-\lambda'_a|d_{j,a},
\]
local dependence on the weight, gauge invariance, and the exact
sourced sum rule. A mixed source derivative vanishes on a domain
which misses any one of its differentiated source supports.
The parameter sets are open convex neighbourhoods of zero and
are stable under coordinatewise contraction. If all weights and
running coefficients are real, conjugation and the reality of
source derivatives at zero hold as well.
\end{proposition}

\begin{proof}
The proof of Proposition~\YMMref{122.1}{826} uses the sources
only through the sums in \eqref{eq:weighted-progression-numbers}.
At height zero the bracketed action is linear in $w_\lambda$.
For a single plaquette its displacement satisfies
$\rho\le2C_PC_{\mathcal A}s_0\alpha_0$, and the bracket is
bounded by $19\rho^2+6\rho\alpha_0L^{-2}$, hence by
$81C_P^2C_{\mathcal A}^2s_0\alpha_0^2$.
There are at most $12M^3$ plaquettes in one elementary domain.
After division by $K_\Lambda s_0$, its normalized size is
at most $c^\flat\sum_a|\lambda_a|m_{0,a}$.
The height-zero hypotheses in (Wp5) put this family inside
the first-step domain. The step estimate gives $d_{1,a}$.

At positive height a weight member has norm at most $m_{j,a}/2$;
its array weight is $\omega_1=2$. Thus the induction yields
\[
 N_j(f_{\lambda,j}-f_{\lambda',j})
 \le\sum_a|\lambda_a-\lambda'_a|
       \left(m_{j,a}+\sum_{i\le j}\omega_{j+1-i}d_{i,a}\right).
\]
The ordinary array uses at most the amount subtracted in
$N_j^{\rm rm}$. Equation~\eqref{eq:weighted-progression-domain}
therefore leaves the input strictly inside the step ball.
Its Lipschitz estimate proves the next bound and its holomorphy
passes through the recursion. This proves the induction for
any finite $p$; no step singles out two coordinates.

If two weights agree on a domain, every earlier member used
there agrees, by induction, and bracket locality identifies the
new members. In particular a member is independent of a parameter
whose source support misses its domain. Differentiation proves
the support assertion. Uniform Cauchy estimates in the weighted
family norm justify these derivatives. The parameter sets are
finite intersections of strict inequalities with nonnegative
coefficients and positive right sides, so they have the stated
geometric properties.

Finally the weight-member action is
$\eta_j^{-1}L^{K-j}\sum_pw_\lambda(p)a_p
=a_K^{-1}\sum_pw_\lambda(p)a_p$, independently of $j$.
It therefore enters the exact step with precisely the desired
source. At height zero its exponential is a continuous bounded
function on the finite product of groups. At positive height
it is the action of an admitted bare array, so the bounded Borel
hypothesis of the sum rule holds. This is the actual-fibre
argument in Proposition~\YMMref{122.1}{826}. Gauge invariance
and, for real data, conjugation pass through the same induction.
\end{proof}

\begin{corollary}[Canonical data for the weighted progression]
\label{cor:weighted-canonical-data}
Fix the admitted canonical construction, including
\eqref{eq:array-borel-radius}, the weighted
old-potential constant and the reduced ceiling in
Section~\YMMref{122}{825}. For arbitrary calibrations
$t_j\ge\gamma^{-2}$, coefficients
$|x_j/t_j-1|<\sigma_\star$ and one lattice
$\mathfrak w\in\mathfrak W_k$, the canonical absolute places
give \textup{(Wp1)--(Wp5)} of
Definition~\ref{def:weighted-progression-data}. They use
\[
 K_\Lambda=K^\omega_\Lambda
       =(1+\beta_{\rm s})^4C_\Lambda C_\Phi,
 \qquad \omega_n=(n+1)L^{-(n-1)}.
\]
The canonical provision is the provision of the actual step at
each height, and its output radii are $\alpha_0,\alpha_1$.
On the global background class of that step, the canonical
interaction part and the array rule have their exact sum rules.
Every admitted array gives a bounded complex Borel density.
At height zero the added action is $\mathcal O(w_\lambda)$.
The constants of these assertions are independent of the height,
lattice, calibrations and sources. Each estimate for a running
sequence retains its additional running and separation hypotheses.
\end{corollary}

\begin{proof}
The common choice in Proposition~\YMMref{69.1}{595}, with
the smaller weighted ceiling specified in
Section~\YMMref{122}{825}, fixes the geometry, auxiliary maps
and response maps before the calibrations. The admissible-radius
verification gives, at every $t_j\ge\gamma^{-2}$, the absolute
place with output radii $\alpha_0,\alpha_1$. Its bracket
localization has the octet property on all the lattices of its
defining class. Thus the bracket hypotheses at positive heights
hold, including the substitution and outer-substitution bounds.
At height zero the same verification gives
\[
 C_{\mathcal A}\delta(t_0)\le\varepsilon_P\le r_P^\circ/2,
 \qquad \alpha_\iota\le a_P\le a_P^\circ\quad(\iota=0,1).
\]
These are precisely the two first-step conditions in (Wp5).
The absolute-place construction and the regular progression use
the same canonical auxiliary maps, restricted to the arguments
of the step. Their three groups of data therefore agree. In
particular the old-potential rule uses $K_\Lambda^\omega$ and
the displayed lag weights, as required in (Wp1).

The global-domain verification of Proposition~\YMMref{52.3}{433}
and the actual sum rule of Theorem~\YMMref{W}{423} identify
the minimal orbit, the shift, its Hessian and the gauge quadratic
form on that class. They also identify the interaction totals
and the total of the array rule. These assertions concern the
place and its operators; changing the permitted array ceiling
does not change them. Theorem~\YMMref{V}{416} consequently
applies with the height-dependent ceiling $N_{\star,j}$.

Here $\operatorname{Re}(x_j/t_j)>0$, since
$\sigma_\star\le1/2$. Lemma~\ref{lem:array-borel-density}
gives the bounded Borel assertion for an array.
At height zero, multiplication by
$\exp\mathcal O(w_\lambda)$ preserves this assertion:
the plaquette action is a finite polynomial in the matrix
entries on the compact configuration space, so this multiplier
is continuous and bounded. These facts give the Borel hypotheses
of the sum rule. The numerical constants were fixed before the
calibrations, the lattice and the weights, which proves the last
assertion. No running comparison is used here.
\end{proof}


\begin{lemma}[Summing domains which meet a support]
\label{lem:source-support-sum}
Let $L\ge13$ be odd, $M=L^{m'}$ with $m'\ge0$, and
$1\le j\le t\le K$, with $m'<m+K-t$. Let $\mathcal S$
be a set of bare plaquettes, and let $N_t(\mathcal S)$ count
the level-$t$ $M$-cubes containing a corner of one of them.
For every integer $q\ge0$, $D\ge0$ and
$\kappa'>\log9216+2$,
\begin{equation}\label{eq:source-support-sum}
 \sum_{\substack{X\in\mathcal D_j:
              \mathcal P(X)\cap\mathcal S\ne\emptyset\\d_j(X)\ge D}}
 \#(X\cap\mathbb T^{(j)})(d_j(X)+1)^q e^{-\kappa'd_j(X)}
 \le8M^3C_{\mathrm E,q+1}L^{3(t-j)}N_t(\mathcal S)e^{-D},
\end{equation}
where $C_{\mathrm E,r}=2048\sum_{n\ge0}(n+2)^re^{-n}$.
\end{lemma}

\begin{proof}
Each bare site belongs to exactly one cube of a partition.
A level-$t$ cube contains $L^{3(t-j)}$ level-$j$ cubes.
There are therefore at most $L^{3(t-j)}N_t(\mathcal S)$
level-$j$ cubes containing a corner of the support, and every
domain in the sum contains one of these cubes.
Also
\[
 \#(X\cap\mathbb T^{(j)})=M^3\#X
                         \le8M^3(d_j(X)+1),
\]
by Proposition~\YMMref{45.1}{353}. For each chosen cube the
same proposition bounds the number of domains of size $n$
containing it by $2048\,9216^n$. Consequently their contribution
to the left side of \eqref{eq:source-support-sum} is at most
\[
 8M^3\,2048\sum_{n\ge D}(n+2)^{q+1}
                  e^{-(\kappa'-\log9216)n}
 \le8M^3C_{\mathrm E,q+1}e^{-D}.
\]
Here $\kappa'-\log9216>2$ and
$e^{-2n}\le e^{-D}e^{-n}$ for $n\ge D$.
Summing over the possible cubes proves the result.
\end{proof}

\subsection{Source families along a finite run}
\label{sec:running-source-chain}

The one-step estimates must be iterated with the increments of the
chosen coupling sequence. These increments remain in the old array
when a source is inserted. We give the recursion and its bounds
before identifying the measure it represents.

Fix $K\ge1$ and a finite real sequence $x_0,\ldots,x_K$ with
\begin{equation}\label{eq:running-source-couplings}
 b_i=L^{-1}x_{i-1}-x_i,\qquad |b_i|\le B_\sharp,
 \qquad x_K\ge\gamma^{-2},\qquad
 \bar B=\frac{L}{L-1}B_\sharp.
\end{equation}
Use one lattice $(\ell,K)$, $\ell=2L^m$, with
$2L^{m+1}\ge C_LM$, and the absolute radius systems with fixed
output radii $(\alpha_0,\alpha_1)$ at all heights. The step data
are the canonical local data of Sections~\YMMref{51}{389}
and~\YMMref{52}{415}, with their all-height auxiliary maps and
octet property. More generally the argument applies to these same
local constructions whenever all the hypotheses below hold.

\begin{definition}\label{def:running-source-setting}
The source recursion uses the admitted local-step hypotheses of
Proposition~\ref{prop:source-entry} at every height $j<K$, with
one common choice of constants and auxiliary maps. In particular,
retain the core and Gaussian hypotheses of
Definition~\YMMref{52.1}{415} and Theorem~\YMMref{V}{416},
the absolute-system and octet hypotheses for the old-potential
bound in Proposition~\YMMref{51.16}{412}, and the source-domain
hypotheses of Lemma~\YMMref{114.1}{800}. The numerical restrictions
on the fixed radii and the step factor are
\begin{equation}\label{eq:running-source-radii}
 \begin{gathered}
 4\beta_{\rm s}\kappa\ge1,\qquad
 \gamma^{-2}\ge2B_\sharp,\qquad \log\gamma^{-2}>2\mathfrak p_0,\\
 \alpha_0+\alpha_1\le a^\star_{\rm S},\qquad
 c_{\rm sh}M\alpha_0<\alpha_1,\qquad
 \alpha_1^2\le c_\natural\alpha_0,\\
 s_\gamma=\frac{(1+\beta_{\rm s})\delta(\gamma^{-2})}{\alpha_0}
                  \le s^\star_{\rm S}.
 \end{gathered}
\end{equation}
The constants in these inequalities are those of the admitted
local-step estimates. The source size $E_0^{\rm s}>0$ and old-array
bound $N_\star^{\rm s}\ge0$ satisfy
\begin{equation}\label{eq:running-source-room}
 \begin{gathered}
 E_0^{\rm s}\ge2E_{\mathsf N},\qquad
 N_\star^{\rm s}\ge\frac{L}{L-1}
       (E_0^{\rm s}+\tfrac32\alpha_0^2B_\sharp),\\
 K_{\mathbf M}\left(\delta(\gamma^{-2})+
 s_\gamma N_\star^{\rm s}
       +e^{-\gamma_2\mathfrak b(\gamma^{-2})^2/2}\right)
                    \le\tfrac14 E_0^{\rm s},\\
 8K_\Lambda s_\gamma N_\star^{\rm s}\le w_\star,
 \qquad 8K_Ps_\gamma\le w_\star.
 \end{gathered}
\end{equation}
Here $\gamma_2=2(LM)^{-3}$ and $K_{\mathbf M}=K_{\mathbf M}(K_\Lambda)$
are the constants used in the recursion preceding
Proposition~\YMMref{114.2}{801}. The constants
$K_{\rm Q}^\vee,K_{\rm L}^\vee$ are chosen by
\YMMref{(2166)}{800}; the common source-weight radii and
$K^{\rm d},K_{\mathbf M}^{\rm f},\mu_\star$ are those of
\eqref{eq:source-weight-room}, with $E_0=E_0^{\rm s}$.
\end{definition}

For a complex function $w$ on the bare plaquettes put
\begin{equation}\label{eq:running-source-recursion}
 \begin{gathered}
 v_j^w=-L^{K-j}w/x_j,\qquad
 \mathfrak f_j[w]=\bigl((E^{(i)}[w],b_i)_{1\le i\le j},\emptyset\bigr),\\
 E^{(j+1)}[w]=\mathbf N^{(j+1)}+
 \mathbf F[x_j;x_j;\mathfrak a_j^{\rm abs};
            V'[\mathfrak f_j[w]]-\mathbb W^{\rm src}[v_j^w]]
                       \qquad(0\le j<K).
 \end{gathered}
\end{equation}
The array at height zero is empty. All objects in a step use the
same lattice and the same auxiliary data. The last increment $b_K$
enters the final array but is not used to form another local family.
The norm $\|E\|_{j,\lambda}$ is the supremum of
$e^{\lambda d_j(X)}|E(X,P,z)|$ over pointed domains and their full
level-$j$ cores.

The largest useful parameter domain here is
\begin{equation}\label{eq:running-source-level-domain}
 \begin{split}
 \|w\|_{\rm lw}
 &=\max_{0\le j<K}\left\{
       \frac{|v_j^w|_\infty}{\sigma_{\rm src}},\quad
       \frac{\operatorname{sw}_j(v_j^w)}{o_{\rm src}(x_j)}\right\},\\
 \mathsf{Src}^{\rm lw}&=\{w:\|w\|_{\rm lw}<1\}.
 \end{split}
\end{equation}
It is an open convex neighbourhood of zero in the finite-dimensional
space of plaquette functions. To express a bound directly in terms
of a physical source, fix $\varrho>0$, let $\mathsf{Bc}$ be the set of plaquette
barycentres, with the torus sup distance, and put
\[
 \mathsf{Src}_{\mathsf{Bc}}(2\varrho,2\varrho)
 =\{w:|w|_\infty<2\varrho,\quad
        (40/\sqrt3)M\operatorname{Lip}_\infty(w;\mathsf{Bc})
                                                   <2\varrho\}.
\]
Require
\begin{equation}\label{eq:running-source-radius}
 2\varrho\le\sigma_{\rm src}(x_K-\bar B),\qquad
 2K_{\rm L}^\vee\varrho\gamma\le(40/\sqrt3)\mu_\star,
\end{equation}
and define
\begin{equation}\label{eq:running-source-difference}
 \mathsf D(w)=2K_{\mathbf M}^{\rm f}
 \left(\frac{2K_{\rm Q}^\vee|w|_\infty}{x_K-\bar B}
       +K_{\rm L}^\vee\gamma M
                           \operatorname{Lip}_\infty(w;\mathsf{Bc})\right).
\end{equation}

\begin{proposition}\label{prop:running-source-chain}
Under Definition~\ref{def:running-source-setting}, the recursion
\eqref{eq:running-source-recursion} is defined through level $K$
for every $w\in\mathsf{Src}^{\rm lw}$. Its families are local
terms on their full cores, jointly holomorphic in the source and
the core variables on each real-background slice, and
\begin{equation}\label{eq:running-source-bounds}
 \|E^{(j)}[w]\|_{j,\kappa+1}<\tfrac{15}{16}E_0^{\rm s},
 \qquad \|E^{(j)}[0]\|_{j,\kappa+1}\le\tfrac34E_0^{\rm s},
 \qquad N_{\mathfrak a_j^{\rm abs}}(\mathfrak f_j[w])
                                                 \le N_\star^{\rm s}.
\end{equation}
A term on $X$ reads $w$ only on its plaquettes. The families commute
with translations by $M\mathbb Z^3$ in their output scale. If the
auxiliary maps commute with conjugation, the families at real
sources are conjugation real. At zero source the recursion agrees
with the finite-run progression formed from the same increments
and auxiliary data.

If \eqref{eq:running-source-radius} holds, its source class lies
in $\mathsf{Src}^{\rm lw}$, and on that class
\begin{equation}\label{eq:running-source-lipschitz}
 \|E^{(j)}[w]-E^{(j)}[0]\|_{j,\kappa+1}
       \le\tfrac23\mathsf D(w),\qquad
 \mathsf D(w)<\tfrac38E_0^{\rm s}.
\end{equation}
\end{proposition}

\begin{proof}
The proof of Proposition~\YMMref{114.2}{801} uses the same local
estimates. We retain the increments in the arrays and verify the
bounds needed to repeat that induction.

First, $x_{j-1}=L(x_j+b_j)>x_j$ whenever $x_j\ge\gamma^{-2}$:
$|b_j|\le x_j/2$ and $L>2$. Thus all $x_j\ge\gamma^{-2}$.
Unrolling the recursion gives
\[
 x_K-\bar B\le L^{-(K-j)}x_j\le x_K+\bar B,
 \qquad x_K>\bar B.
\]
The lattice admitted at height $K-1$ is admitted at every earlier
height since its side in level-$j$ units is
$2L^{m+K-j}\ge2L^{m+1}$. The absolute systems have common output
radii and step factors $s_j\le s_\gamma$. Their restrictions to
older levels are exactly the domains required by the old-array
map. These are the same place and domain checks as in the
absolute-system construction of Section~\YMMref{51}{389}.

Suppose the families through level $j$ have norm at most
$E_0^{\rm s}$. Their bare-array norm is bounded by
\[
 \sum_{i=1}^jL^{-(j-i)}
       (E_0^{\rm s}+\tfrac32\alpha_0^2|b_i|)
 \le\frac{L}{L-1}(E_0^{\rm s}+\tfrac32\alpha_0^2B_\sharp)
 \le N_\star^{\rm s}.
\]
Hence $8K_\Lambda s_jN(\mathfrak f_j[w])\le w_\star$.
Every $w\in\mathsf{Src}^{\rm lw}$ has an admitted relative weight
at this step. Proposition~\ref{prop:source-entry} therefore forms
$E^{(j+1)}[w]$ on the full output core. The normalizer and
local-output estimates used in the proof of
Proposition~\YMMref{114.2}{801} give
\[
 \begin{split}
 \|E^{(j+1)}[w]\|_{j+1,\kappa+1}
 &\le\tfrac12E_0^{\rm s}
  +K_{\mathbf M}\bigl(\delta(x_j)+s_jN(\mathfrak f_j[w])
                   +e^{-\gamma_2\mathfrak b(x_j)^2/2}\bigr)
  +K^{\rm d}\mathsf{sz}_{T_\ell}(v_j^w)\\
 &<\left(\tfrac12+\tfrac14+\tfrac3{16}\right)E_0^{\rm s}.
 \end{split}
\]
Here $\delta(x_j)\le\delta(\gamma^{-2})$,
$\mathfrak b(x_j)\ge\mathfrak b(\gamma^{-2})$, and the last
term is less than $3E_0^{\rm s}/16$ by
\eqref{eq:source-weight-room}. At zero weight the last term
vanishes. This closes the norm induction, starting with the empty
array; $4\beta_{\rm s}\kappa\ge1$ allows its output rate to be
replaced by $\kappa+1$.

For the smaller physical source class,
\[
 |v_j^w|_\infty\le\frac{|w|_\infty}{x_K-\bar B},\qquad
 \operatorname{sw}_j(v_j^w)
       \le x_j^{-1}M\operatorname{Lip}_\infty(w;\mathsf{Bc}).
\]
Together with $x_j^{-1/2}\le\gamma$ and
\eqref{eq:running-source-radius}, these inequalities place every
relative weight in its open weight class. They also give
$K_{\mathbf M}^{\rm f}\mathsf{sz}_{T_\ell}(v_j^w)
\le\mathsf D(w)/2$ and
$\mathsf D(w)<6K_{\mathbf M}^{\rm f}\mu_\star
\le3E_0^{\rm s}/8$.

Write $D_i=\|E^{(i)}[w]-E^{(i)}[0]\|_{i,\kappa+1}$.
The counterterms cancel in the difference of the old arrays.
Their norm is thus at most $L(L-1)^{-1}\max_{i\le j}D_i$.
Linearity of the old-potential map and
Proposition~\ref{prop:source-entry} imply
\[
 D_{j+1}\le K^{\rm d}K_\Lambda s_\gamma
                 \frac{L}{L-1}\max_{i\le j}D_i
                      +\tfrac12\mathsf D(w).
\]
The smallness conditions imply
\[
 K^{\rm d}K_\Lambda s_\gamma\frac{L}{L-1}
 \le\frac{(351/2)}{243}\,\frac14
 =\frac{13}{72}<\frac14.
\]
Indeed $K_{\mathbf M}\ge243L^3K_\Pi K_\Lambda$ and
$N_\star^{\rm s}\ge L(L-1)^{-1}E_0^{\rm s}$ in
\eqref{eq:running-source-room}. Induction, with maximum zero
at height zero, gives $D_i\le2\mathsf D(w)/3$.

The old-array parameters are holomorphic by induction, and
$w\mapsto v_j^w$ is linear; the parameter assertion in
Proposition~\ref{prop:source-entry} proves joint holomorphy.
If two sources agree on $X$, their older terms agree on every
subdomain of $X$, and their direct source potentials agree on $X$.
Locality of the same map proves the asserted source dependence.
Conjugation follows inductively when the auxiliary maps have that
covariance. An output translation in $M\mathbb Z^3$ is an input
translation in $LM\mathbb Z^3$ and is an admitted translation at
every older level, so the local translation covariance also
passes through the induction. Finally at $w=0$ the source family
vanishes; both recursions have identical arrays, increments and
step maps at every height. This proves their equality.
\end{proof}

At the canonical source choice we retain the following numerical
part of the choice explicitly. Put
\begin{equation}\label{eq:canonical-source-quarter-room}
 \begin{gathered}
 N_1=\frac{L}{L-1}(E_0+\tfrac32\alpha_0^2B_\sharp),\qquad
 N_\star\ge\frac{2L-1}{L-1}N_1,\qquad
 K_\Lambda^\omega=(1+\beta_{\rm s})^4K_\Lambda,\\
 K_{\mathbf M}(K_\Lambda^\omega)
 \left(\delta(\gamma^{-2})+s_\gamma N_1
                 +e^{-\gamma_2\mathfrak b(\gamma^{-2})^2/2}\right)
                                      \le E_0/4,\\
 s_\gamma\le\frac{w_\star}{4K_\Lambda^\omega N_\star},
 \qquad 8K_Ps_\gamma\le w_\star.
 \end{gathered}
\end{equation}
The allowance $N_\star$ is chosen after $E_0,B_\sharp$ and before
$\gamma$. The last three requirements are upper bounds tending
to zero with $\gamma$, with all other quantities already fixed.
They can therefore be retained in the ordered common choice of
Proposition~\YMMref{69.1}{595}. They are the source-room
conditions also imposed in Section~\ref{sec:cutoff-calibration}
and Lemma~\ref{C-363-23}; a half-size bound alone is not used in
the source induction.

\begin{corollary}\label{cor:canonical-running-source}
Use the canonical construction data of
Propositions~\YMMref{69.1}{595} and~\YMMref{69.3}{603}, with
the numerical source choice \eqref{eq:canonical-source-quarter-room}.
Take $E_0^{\rm s}=E_0$ and $N_\star^{\rm s}=N_1$.
Every admitted finite run and lattice satisfy the setting of
Proposition~\ref{prop:running-source-chain}. A positive source
radius can be chosen independently of the run and the lattice by
requiring \eqref{eq:running-source-radius} with $\gamma^{-2}$
in place of $x_K$. The zero-source families are the canonical
finite-run families restricted to this lattice.
\end{corollary}

\begin{proof}
The common construction supplies the all-height auxiliary maps,
their octet and conjugation properties, the absolute-radius
inequalities and the local-step hypotheses. Its size choice gives
$E_0\ge2E_{\mathsf N}$. Moreover
$N_\star\ge(2L-1)(L-1)^{-1}N_1\ge2N_1$.
Since $K_{\mathbf M}(K_\Lambda)\le K_{\mathbf M}(K_\Lambda^\omega)$
the quarter-room inequality gives
\eqref{eq:running-source-room} at $N_\star^{\rm s}=N_1$.
For its remaining array allowance,
\[
 8K_\Lambda s_\gamma N_1
 \le 8K_\Lambda\frac{w_\star}{4K_\Lambda^\omega N_\star}
                      \frac{N_\star}{2}
 =\frac{K_\Lambda}{K_\Lambda^\omega}w_\star\le w_\star.
\]
The other source ceiling is explicitly included in the choice.
Since $\gamma^{-2}>\bar B$, both upper bounds for $\varrho$ are
positive. The final assertion follows from the zero-source
identification in the proposition.
\end{proof}

\subsection{Source energies and the exact density}
\label{sec:running-source-density}

The values of the local families at the trivial triple are
configuration-independent. Their differences from zero source
are the source energies:
\begin{equation}\label{eq:running-source-energies}
 \begin{gathered}
 e_j(X,z;w)=E^{(j)}[w](X,\mathbf1,z)
                         -E^{(j)}[0](X,\mathbf1,z),\\
 e_j(X;w)=\sum_{z\in X\cap\mathbb T^{(j)}}e_j(X,z;w),\qquad
 \mathcal E_n(w)=\sum_{j=1}^n\sum_{X\in\mathcal D_j}e_j(X;w).
 \end{gathered}
\end{equation}
The trivial triple belongs to every core here. Each energy is
holomorphic, vanishes at zero, and reads the source only on $X$.
If the family difference has norm at most $D$, then
\begin{equation}\label{eq:running-source-energy-bound}
 |e_j(X,z;w)|\le De^{-(\kappa+1)d_j(X)},\qquad
 |e_j(X;w)|\le8M^3D e^{-\kappa d_j(X)}.
\end{equation}
Indeed $\#(X\cap\mathbb T^{(j)})=M^3\#_jX
\le8M^3(d_j(X)+1)$ and $(d+1)e^{-d}\le1$.
The polymer count of Proposition~\YMMref{45.1}{353} then gives,
for $\kappa>\kappa_{\mathcal D}+2$,
\[
 \sum_{X\ni z}|e_j(X;w)|\le8M^3C_{\rm E,0}D,
 \qquad \sum_{X,z}|e_j(X,z;w)|
                       \le C_{\rm E,0}D\,\#\mathbb T^{(j)}.
\]
One may use $D=\mathsf D(w)$ on the physical source class.
On the full level-wise ball use $D=2E_0^{\rm s}$ and apply
Schwarz's lemma to the radial complex line to obtain
\begin{equation}\label{eq:running-source-energy-linear}
 |e_j(X;w)|\le16M^3E_0^{\rm s}\|w\|_{\rm lw}
                                             e^{-\kappa d_j(X)}.
\end{equation}
The locality and conjugation statements follow by evaluation
of the corresponding family statements.

For the measure identity we require the actual history tests,
not merely the bounds on the local families.
\begin{definition}\label{def:running-source-history}
In addition to Definition~\ref{def:running-source-setting}, use
an admitted history frame for the same finite sequence through
level $K$, and its restrictions to earlier levels. Retain all
numerical, width, collar, response and curvature hypotheses in
item~\textup{(J.4)} of Definition~\ref{def:weighted-history-setting}:
in particular the actual test constructions of
Propositions~\YMMref{53.7}{450}, \YMMref{53.13}{459} and
\YMMref{53.14}{461}, with \YMMref{(1030)}{459}, apply at each
step. The bare-array action and local sum rule have the hypotheses
of Theorem~\YMMref{V}{416}. Require also
\begin{equation}\label{eq:running-source-history-room}
 \begin{gathered}
 0<\alpha_\iota\le1,\quad \alpha_0\le a_{\rm H},\quad
 m'\ge6,\quad C_L\le2LR(\gamma^{-2}),\\
 2\varepsilon(\gamma^{-2})\le\varepsilon_{\mathcal W}\alpha_0,
 \qquad \delta(\gamma^{-2})\le\varepsilon_{\mathcal W}.
 \end{gathered}
\end{equation}
The constants have the history ranges, including
$\max\{A_0,A_1\}\gamma(\log\gamma^{-2})^{\mathfrak p_0}\le1$
and $(L^{1/r}-1)\log\gamma^{-2}\ge\log(2L)$.
The canonical Gaussian families, prefactors and vacuum numbers
are used in the local steps and the history. The coefficients of
the bare arrays are the increments $b_i$ of
\eqref{eq:running-source-couplings}. This definition does not
require these increments to be the counterterms of a regular
progression.
\end{definition}

Write $\mathcal V_n^\chi=\{W:\chi_n(T_\ell)(W)=1\}$ and
$h_n$ for the history with no large-field region through level $n$.
Put
\begin{equation}\label{eq:running-weighted-action}
 A^{\eta_n}(c,U)=L^n\sum_pc(p)a_p(U),\qquad
 x_n^w=x_n-L^{K-n}w,\qquad
 \rho_0^w(U)=e^{-A^{\eta_0}(x_0^w,U)}.
\end{equation}

\begin{proposition}\label{prop:running-source-density}
Under Definition~\ref{def:running-source-history}, every
$W\in\mathcal V_n^\chi$, $1\le n\le K$, has a minimal background
orbit. The sourced bare-array action is defined on that orbit
and is constant there. There are source-independent complex numbers
$c_j$, real when the auxiliary maps commute with conjugation, such that, for $w\in\mathsf{Src}^{\rm lw}$,
\begin{equation}\label{eq:running-source-density}
 \begin{split}
 \tau_{h_n}[\rho_0^w]&=f_n^w\,dW,\\
 f_n^w(W)&=\mathbf1_{\mathcal V_n^\chi}(W)
 \exp\left[\sum_{j=1}^nc_j+\mathcal E_n(w)
       -A^{\eta_n}(x_n^w,U_n(W))
       +S[\mathfrak f_n[w]](U_n(W))\right].
 \end{split}
\end{equation}
In the vacuum notation for the history,
$E_{(n)}[\mathsf{vd}^{0}]=\sum_{j\le n}c_j$ and
$E_{(n)}[\mathsf{vd}^{w}]=E_{(n)}[\mathsf{vd}^{0}]+\mathcal E_n(w)$.
The density is bounded and Borel. Any representative of the
minimal orbit can be used in this formula. At $n=0$ the measure
is $\rho_0^w\,dU$.
\end{proposition}

\begin{proof}
A passing density test bounds every level-$n$ plaquette by
$2\varepsilon(x_n)$. The room conditions therefore put
$\mathcal V_n^\chi$ inside the global class of the sum rule and
its chart image, by Proposition~\YMMref{52.3}{433}.
The proof of Lemma~\ref{lem:weighted-passing-fields} applies
verbatim to these tests: it uses their geometric hypotheses and
the coupling bounds, not the rule defining the increments.
Thus every cutoff-chart field over a passing $W$ passes all
preceding density tests and all the $Q$ tests of the step.

The global class supplies the minimal orbit and its plain core
representatives at every earlier level. Hence each pair read by
$\mathfrak f_n[w]$ is in the required domain. At the top level
use the plain-core verification in the proof of
Proposition~\YMMref{52.3}{433}. Gauge invariance of the local
terms makes the action constant on the orbit. The density
domain and Borel argument in Theorem~\YMMref{W}{423} applies
to these bounded complex local families as well: the domains
are independent of their values, and composition on the
background charts is Borel.

We first compute one chart integral. Let $F_k^w$ equal
\[
 \exp[-A^{\eta_k}(x_k^w,U_k(V))
                       +S[\mathfrak f_k[w]](U_k(V))]
\]
on the orbit domain where all its array pairs are admitted,
and zero otherwise. For $k=0$ its value on every chart over the global class equals
$\rho_0^w$. The source sum rule of
Proposition~\ref{prop:source-entry}, with the local and
whole-torus identities of Theorem~\YMMref{V}{416}, gives
\begin{equation}\label{eq:running-source-one-step}
 \mathsf I[F_k^w](W)=
 \exp\left[c_{k+1}+\sum_Xe_{k+1}(X;w)
       -A^{\eta_{k+1}}(x_{k+1}^w,U_+)
       +S[\mathfrak f_{k+1}[w]](U_+)\right]
\end{equation}
on the global class, where $U_+=U_{k+1}(W)$ and
\[
 c_{k+1}=\log(\mathsf{pf}_k(x_k)\mathsf Z^{(k)}[\mathbf1])
                  +\sum_{X,z}E^{(k+1)}[0](X,\mathbf1,z).
\]
To check the normalization, separate the new local sum into its
value minus its value at $\mathbf1$, its zero-source value at
$\mathbf1$, and the source energy. The first part enters the
new action; the latter two give the displayed constant and
energy. More explicitly, the empty-slot action has the increment identity
\[
 S[\mathfrak f_{k+1}[w]](U)-S[\mathfrak f_k[w]](U)
 =\sum_{X,z}\bigl(E^{(k+1)}[w](X,U,z)
               -E^{(k+1)}[w](X,\mathbf1,z)\bigr)
                      -b_{k+1}A^{\eta_{k+1}}(1,U).
\]
Since $A^{\eta_k}(c,U)=A^{\eta_{k+1}}(L^{-1}c,U)$,
its coupling part uses
\[
 L^{-1}x_k^w-b_{k+1}
 =L^{-1}x_k-b_{k+1}-L^{K-k-1}w=x_{k+1}^w.
\]
This proves the exponent identity with the actual increment.
The vacuum system keeps the same Gaussian and gauge-fixing
normalizers at every source. Only its base-point values change,
by $e_j(X,z;w)$. Summing those changes over $X,z$ and then over
the levels gives the asserted vacuum identity.

We now induct on the history length. The empty-outcome factor
is the product displayed in
\eqref{eq:weighted-history-factor}, with the present couplings.
Its derivation in Proposition~\ref{prop:weighted-history-step}
uses only the definition of the history and the actual tests,
and is unchanged for these increments. It is a bounded
nonnegative Borel factor, so its pushforward identity holds
also for a bounded complex initial density. At a passing
$W$ every additional fibre test is one, as checked above.
If the formula is known at level $k$, its chart integrand is
therefore $F_k^w$ times the constant
$\exp[\sum_{j\le k}c_j+\mathcal E_k(w)]$.
Equation~\eqref{eq:running-source-one-step} proves the formula
at $k+1$; for the first step the chart integrand is $\rho_0^w$.

The chart integral of a bounded Borel function is bounded and
Borel on its open image domain, by
Lemma~\ref{lem:sf-chart-integral}. It agrees with the new density
on the open set $\mathcal V_{k+1}^\chi$ and is set to zero off
that set. Thus each displayed density is bounded and Borel,
and the pushforward identity holds on every Borel set.
The argument applies on the level-wise source ball because
all its steps have the admitted relative weights.
\end{proof}

\subsection{The weighted progression on an actual history}
\label{sec:weighted-history-interface}

The sum rule computes an integral over a fluctuation chart. To use
it for a history measure, one must check that the extra tests in
the history factor equal one throughout that chart. We give this
check before identifying the sourced density.

\begin{definition}\label{def:weighted-history-setting}
The setting \textup{(J)} is the following specialization of
Definition~\ref{def:weighted-progression-data}.
\begin{itemize}
\item[(J.1)] The calibrations and running coefficients agree:
$t_j=x_j\ge\gamma^{-2}$ for $0\le j\le k$, where $k\ge1$.
They are real and satisfy
\[
 x_j=L^{-1}x_{j-1}-e_j\quad(1\le j\le k+1),
 \qquad x_{k+1}:=L^{-1}x_k-e_{k+1}.
\]
The counterterms are those of this same regular progression.
\item[(J.2)] At every height $j\le k$ use the actual canonical
place and the absolute radius system with output radii
$(\alpha_0,\alpha_1)$, the auxiliary maps and response maps
of Sections~\YMMref{51}{389}, \YMMref{52}{415} and
\YMMref{69}{595}. All their hypotheses hold; the classes of
the sum rule are built with the same fixed shift radii. In
particular $0<\alpha_\iota\le1$,
$\alpha_0\le a_{\rm H}$, and
$\delta(x_j)\le\varepsilon_{\mathcal W}$.
Require also the following identification, denoted \textup{(Pv)}:
the three groups defining the regular step, restricted to this
place, are its actual canonical step data, and the summand
$\mathbf N^{(j+1)}$ in its localized output is the localized
Gaussian family used in that sum rule. All common symbols
denote the same objects. This identification is part of the
setting; norm bounds for an abstract progression do not imply it.
\item[(J.3)] The next coefficient is greater than one, and
\begin{equation}\label{eq:weighted-history-smallness}
 2\varepsilon(x_t)\le\varepsilon_{\mathcal W}\alpha_0
 \qquad(1\le t\le k+1).
\end{equation}
\item[(J.4)] The lattice has an admitted history frame through
level $k+1$, with these same coefficients. Its restrictions to
earlier levels are the corresponding frames. Retain all the
numerical, radius and collar hypotheses of the actual test
construction, including Propositions~\YMMref{53.7}{450},
\YMMref{53.13}{459} and~\YMMref{53.14}{461} and
equation~\YMMref{(1030)}{459}. Thus the smallness of $A_1/A_0$,
the exponential collar bound, the response and curvature ceilings,
and $\check R_j\le L\check R_{j+1}$ are required. The density
thresholds also satisfy $\varepsilon_t\le\varepsilon_\flat$.
These are hypotheses on the same constants and frame, not a
change of constants at a later step.
\end{itemize}
We refer to the first two items alone as \textup{(J.1), (J.2)}.
\end{definition}

Write $h_t$ for the history with no large-field region through
level $t$ and put
\[
 \mathcal V_t^\chi=\{W:\chi_t(T_\ell)(W)=1\}.
\]
The tests use strict inequalities, so this set is open. At the
step $j\to t=j+1$, let $q_t$ be the product of all the $Q$
tests, and let $q_j^0=1$ for $j>0$ and $q_0^0=\chi_0(T_\ell)$.
The symbols $\mathfrak G_j[x_j]$, $\chi^{\rm sf}_{x_j}$,
$O^{\rm im}$ and $\mathsf I$ have the meanings of
Section~\ref{sec:sf-chart-interface}.

\begin{lemma}\label{lem:weighted-passing-fields}
In setting \textup{(J)}, for $1\le t\le k+1$,
\begin{equation}\label{eq:weighted-passing-fields}
 \mathcal V_t^\chi
 \subset\{W:\varpi(W)<2\varepsilon_t\}
 \subset\mathcal W_t(x_{t-1};\alpha_0,\alpha_1)
 \subset O^{\rm im}.
\end{equation}
Every member of the minimal orbit over a field in the middle
set is regular at level $t$ with radius $2B_3\varepsilon_t$.
Consequently its bare plaquettes satisfy
\begin{equation}\label{eq:weighted-background-plaquettes}
 0\le a_p(U_t(W))<2B_3^2\varepsilon_t^2L^{-4t}.
\end{equation}
For every $W\in\mathcal V_t^\chi$ and every point $V_W(B)$
of its cutoff chart, the tests satisfy
\begin{equation}\label{eq:weighted-fibre-tests}
 q_t(V_W(B))=q_j^0(V_W(B))
             =\chi_j(T_\ell)(V_W(B))=1.
\end{equation}
\end{lemma}

\begin{proof}
At every level-$t$ plaquette, one of the passing density tests
gives a defect less than $2\varepsilon_t$, by the density
comparison in equation~\YMMref{(650)}{344}. Taking the maximum
over the finite set of plaquettes proves the first inclusion.
The next follows from \eqref{eq:weighted-history-smallness}
and Proposition~\YMMref{52.3}{433}, with
$\mathsf t=x_{t-1}$ and $b_\iota=\alpha_\iota$.
That proposition includes the existence of the whole cutoff
chart and its image domain. The background theorem used in
its proof, Theorem~\YMMref{E}{10}, applied at curvature
$2\varepsilon_t$, gives the stated regularity for every
representative. Since for $u\in\mathrm{SU}(2)$
$1-\tfrac12\operatorname{tr}u=\tfrac12|u-1|^2$, its
bare-plaquette estimate gives
\eqref{eq:weighted-background-plaquettes}.

For the last assertion, put $V=V_W(B)$. The chart identity gives
$\mathcal MV=W$ and $\chi^{\rm sf}_{x_j}(V)=1$.
Here all regions of the history are the whole torus, so there
are no external zones or stored large-field factors. The
comparison of the background of a test cube with the whole
background, in Proposition~\YMMref{53.7}{450}, bounds their
bond difference by $\delta_j$. The cutoff chart bounds the
fluctuation bond by another $\delta_j$. These are precisely
the two contributions to the $Q$ test. The verification in
the proof of Proposition~\YMMref{53.14}{461} therefore makes
every $Q$ test pass. The new density tests supply the $P$
tests, so $(V,W)\in\mathsf{PQ}_t(T_\ell)$.
Proposition~\YMMref{53.13}{459}, with $Y=T_\ell$, now makes
every level-$j$ density test pass. Its hypotheses include the
collar and response bounds retained in (J.4). At $j=0$ this
also proves $q_0^0=\chi_0=1$; at positive height $q_j^0=1$
by definition. This proves \eqref{eq:weighted-fibre-tests}.
\end{proof}

\begin{proposition}\label{prop:weighted-history-step}
In setting \textup{(J)}, the step factor of $h_t$ is
\begin{equation}\label{eq:weighted-history-factor}
 g_t(V)=\chi_t(T_\ell)(\mathcal MV)\,
        q_j^0(V)q_t(V)\mathfrak G_j[x_j](V)
                                      \chi^{\rm sf}_{x_j}(V).
\end{equation}
It is bounded, nonnegative and Borel. For every bounded
complex Borel bare density $\varsigma$,
\[
 \tau_{h_t}[\varsigma]
       =\mathcal M_*\bigl(g_t\tau_{h_j}[\varsigma]\bigr).
\]
If bounded Borel functions $F,F'$ on $\mathcal V_j$ satisfy
$F'=e^cF$ on $\mathcal V_j^\chi$, then
\begin{equation}\label{eq:weighted-history-chart}
 \mathcal M_*(g_tF'\,dV)
       =\chi_t(T_\ell)e^c\mathsf I[F]\,dW.
\end{equation}
The right-hand density is zero off $\mathcal V_t^\chi$.
\end{proposition}

\begin{proof}
In the history with no large-field region, every previous
proper coordinate set is empty. The current field is $V$,
its new field is $\mathcal MV$, and the earlier unused
readings are the identity. The density-test part of the
empty outcome is $\chi_t(\mathcal MV)$ at positive height,
with the additional factor $\chi_0(V)$ in the first step.
Its $Q$ part is $q_t$, and its bond-fluctuation part is the
product defining $\chi^{\rm sf}_{x_j}$. The gauge factor is
taken over all coarse sites and equals $\mathfrak G_j[x_j]$.
This gives \eqref{eq:weighted-history-factor}. The bounded
Borel assertions follow from the corresponding assertions
for these factors. The identity of history weights is
\[
 \mathfrak w_{h_t}(U)=\mathfrak w_{h_j}(U)
                                  g_t(\mathcal M^jU).
\]
Integrating a bounded Borel test function of $\mathcal M^tU$
and then pushing forward by $\mathcal M^j$ proves the measure
recursion.

Apply the chart formula of Lemma~\ref{lem:sf-chart-integral}
to $F_1=q_j^0q_tF'$. By
\eqref{eq:weighted-history-factor}, the resulting density
on the passing set is $\mathsf I[F_1]$ and is zero elsewhere.
For $W$ in that set and $B$ in the integration domain,
\eqref{eq:weighted-fibre-tests} gives
$F_1(V_W(B))=F'(V_W(B))=e^cF(V_W(B))$.
The chart integral is linear and depends on its argument
only through these values. Hence
$\mathsf I[F_1](W)=e^c\mathsf I[F](W)$, which proves
\eqref{eq:weighted-history-chart} as an identity of complex
Borel measures.
\end{proof}

\subsection{Exact array steps and covariance quotients}
\label{sec:weighted-covariance-interface}

The history identity must also be available when just one array
is varied. This is what allows a derivative created in one step
to be distinguished from derivatives carried by the earlier
families. The input below is an arbitrary admitted array; it
need not be an array already produced by the source recursion.

For an array $g$ of height $j$, let $\mathsf{Df}_j$ be the
common domain of its action $S[g]$: every pair needed by its
members belongs to the prescribed analytic core. The domain
is gauge invariant and the action is constant on each gauge
orbit. Define $F_g$ on $\mathcal V_j$ by
\[
 F_g(V)=\exp\{-x_jA^{\eta_j}(U_j(V))+S[g](U_j(V))\}
\]
when the minimal orbit exists and lies in $\mathsf{Df}_j$,
and by zero otherwise. The notation $U_j(V)$ in this formula
means any representative of that orbit. Write
$\mathsf Q_j=\{(X,z):X\in\mathcal D_j,\ z\in X\cap\mathbb T^{(j)}\}$,
and let $\mathsf{Sum}_j$ denote the sum over this finite set.

\begin{lemma}[An array in one exact step]
\label{lem:weighted-array-step}
Assume \eqref{eq:array-borel-radius} and (J.1), (J.2), including (Pv), of
Definition~\ref{def:weighted-history-setting}. For $1\le j\le k$
and $g\in\mathsf{Arr}^{[j]}$, both $F_g$ and
$\chi^{\rm sf}_{x_j}F_g$ are bounded complex Borel functions.
Their small-field pushforward has chart density $\mathsf I[F_g]$.
On the global background class $\mathcal W_{j+1}(x_j;\alpha_0,\alpha_1)$,
\begin{equation}\label{eq:weighted-array-step}
 \mathsf I[F_g](W)=\exp\left\{
 c_j+\mathsf{Sum}_{j+1}
       [\mathbf N^{(j+1)}+\mathsf{St}_j(g)](U_{j+1}(W))
 -x_jA^{\eta_j}(U_{j+1}(W))+S[g](U_{j+1}(W))\right\}.
\end{equation}
Here
$c_j=\log\mathsf{pf}_j(x_j)+\log\mathsf Z^{(j)}[(\mathbf1,0,0)]$
is independent of $g$. All terms are defined and single-valued
on the minimal orbit.

At height zero, for $\lambda\in\Lambda^{\rm w}_0$, put
$\rho_0^\lambda=\exp\{-x_0A^{\eta_0}+\mathcal O(w_\lambda)\}$.
It is continuous and bounded, and on
$\mathcal W_1(x_0;\alpha_0,\alpha_1)$ its chart density is
\begin{equation}\label{eq:weighted-first-step}
 \mathsf I[\rho_0^\lambda](W)=
 \exp\left\{c_0+\mathsf{Sum}_1[E^{(1)}_\lambda](U_1(W))
 -x_0A^{\eta_0}(U_1(W))+\mathcal O(w_\lambda)(U_1(W))\right\}.
\end{equation}
\end{lemma}

\begin{proof}
The orbit domain and gauge invariance give the stated definition
of $F_g$. Lemma~\ref{lem:array-borel-density}, together with the chart formula of
Lemma~\ref{lem:sf-chart-integral}, give the bounded Borel
functions and their pushforward density. The hypotheses on the
actual place include $\alpha_0\le a_{\rm H}$, the chart
smallness and the required ratio $A_0/A_1$; these are part of
(J.2) and its canonical provision, not extra conclusions from
the bound on $N_j(g)$.

At $W$ in the stated class, choose one representative of the
minimal orbit to define all its operators. The global-domain
proof of Proposition~\YMMref{52.3}{433} and
Theorem~\YMMref{W}{423} identify the shift and Hessian,
and make the interaction and array totals sum to their physical
actions. Apply the whole-torus identity of
Theorem~\YMMref{V}{416}, equation~\YMMref{(902)}{417},
at $\tau=1$ and the fourth potential formed from $g$.
The logarithm of the integral is
\[
 \log\mathsf{pf}_j(x_j)+\log Z^{(j)}(W)
 +\mathsf{Sum}_{j+1}[\mathsf{St}_j(g)](U_{j+1}(W))
 -x_jA^{\eta_j}(U_{j+1}(W))+S[g](U_{j+1}(W)).
\]
The Gaussian identification in (Pv) gives
\[
 \log Z^{(j)}(W)=\log\mathsf Z^{(j)}[(\mathbf1,0,0)]
           +\mathsf{Sum}_{j+1}[\mathbf N^{(j+1)}](U_{j+1}(W)).
\]
The logarithms here are of positive Gaussian normalizations.
Every restriction of the background triple is its plain
triple on that domain, so the pointed terms can indeed be
evaluated at $U_{j+1}(W)$. This proves
\eqref{eq:weighted-array-step}.

For the first step the bare sourced density is continuous
on the compact configuration space. Its orbit-defined version
equals it wherever the height-zero regular background exists,
and vanishes elsewhere. That regularity set is open: it is
given by finitely many strict curvature and current inequalities.
The two functions agree everywhere on every chart fibre in
the global background class, since all those fibre backgrounds
exist. They therefore have the same chart integral. Applying
the same whole-torus identity, now to the height-zero
fourth-place source family, gives
\eqref{eq:weighted-first-step}, because
$E^{(1)}_\lambda=\mathbf N^{(1)}+\mathbf G_1+
\mathbf F[\mathbb V^{\rm w}_\lambda]$. The first two families
are independent of $\lambda$. The same restriction and
gauge-invariance argument proves single-valuedness.
\end{proof}

For two weights in setting (J), write
$X_{j,a}(V)=\mathcal O(w_a)(U_j(V))$ on the passing set,
continued by zero elsewhere; at height zero put
$X_{0,a}=\mathcal O(w_a)$. Let
$m_j=\tau_{h_j}[\rho_0]$ and
$\sigma_j[F]=\mathcal M_*(Fg_jm_{j-1})$. The step quotient
is the Radon--Nikodym derivative
$\mathsf{Qt}_j[F]=d\sigma_j[F]/dm_j$; its covariance is
$\mathsf{Qt}_j[F;G]=\mathsf{Qt}_j[FG]-
\mathsf{Qt}_j[F]\mathsf{Qt}_j[G]$.
These quotients exist for bounded Borel functions, since
$|\sigma_j[F]|\le\|F\|_\infty m_j$.

\begin{lemma}[The covariance created in one step]
\label{lem:weighted-step-covariance}
In setting (J), with \eqref{eq:array-borel-radius}, at height $i\ge1$, leave every earlier
family and counterterm equal to its zero-source value and
insert $s_1\mathsf W_{L^{K-i}w_1}+s_2\mathsf W_{L^{K-i}w_2}$
in the slot. Let $H_s$ be the output of that one step.
For $i=0$, let $H_s=\mathbf F[\mathbb V^{\rm w}_s]$.
Then $K^{(i+1)}=\partial_1\partial_2H_s|_{s=0}$ is defined,
vanishes on domains which miss either support, and
\begin{equation}\label{eq:weighted-step-covariance}
 \mathsf{Qt}_{i+1}[X_{i,1};X_{i,2}](W)
       =\mathsf{Sum}_{i+1}[K^{(i+1)}](U_{i+1}(W))
       \quad\hbox{$m_{i+1}$-almost everywhere}.
\end{equation}
For $i=0$, $K^{(1)}=\ddot E^{(1)}$. For $i\ge1$ these
are exactly the family and identity in
equations~\YMMref{(2278)--(2279)}{831}.
\end{lemma}

\begin{proof}
At positive height let $u_a$ be the array with just the
indicated slot member. Its norm is at most $m_{i,a}$,
including the lag weight $\omega_1=2$. Thus
$g_s=f_{0,i}+s_1u_1+s_2u_2$ is admitted whenever
$\sum_a|s_a|m_{i,a}<N_i^{\rm rm}$, an open neighbourhood
of zero. Its action differs from the ordinary one by
$\mathcal O(w_s)$, since
$\eta_i^{-1}L^{K-i}=L^K$. Holomorphy of the step gives
the mixed derivative, and locality makes it vanish when
a source support misses the domain. At height zero the
same facts follow from Proposition~\ref{prop:weighted-finite-system};
the Gaussian and normalizer terms have zero source derivative,
so $K^{(1)}=\ddot E^{(1)}$.

Apply Lemma~\ref{lem:weighted-array-step} at $g_s$ and
$g_0$, and use Proposition~\ref{prop:weighted-history-step}.
The normalizer and all source-independent terms cancel.
Consequently
\[
 \sigma_{i+1}[e^{s_1X_{i,1}+s_2X_{i,2}}]
       =e^{b_s}m_{i+1},\qquad
 b_s(W)=\mathsf{Sum}_{i+1}[H_s-H_0](U_{i+1}(W))
                    +\sum_as_aX_{i+1,a}(W).
\]
For $i=0$ the same identity follows from
\eqref{eq:weighted-first-step}. The quotient is Borel on
the passing set because it is a quotient of Borel chart
densities with nonzero denominator; extend it by zero outside.
The finite sum and the uniform Lipschitz bounds give a bound
on $b_s$ independent of $W$ on every smaller parameter polydisc.

Test this measure identity against a bounded Borel function.
On both sides the integrands are uniformly bounded on a
smaller polydisc, so differentiation under the integrals is
valid. The first derivative of $e^{b_s}$ at zero is
$\partial_ab_s$, and its mixed derivative is
$\partial_1\partial_2b_s+\partial_1b_s\,\partial_2b_s$.
The product term cancels in the covariance. The last term of
$b_s$ is linear in $s$ and contributes no mixed derivative.
The remaining derivative is the right side of
\eqref{eq:weighted-step-covariance}. This also proves the
assertion at the first step.
\end{proof}


We also record exactly how large the running coefficients
must be in applications. This keeps the threshold used in
the source estimates independent of the number of sources.
Put
\[
 \theta(x)=K_{\mathbf M}(1+\beta_{\rm s})\delta(x)/\alpha_0.
\]
Let $x_{\rm J1}$ be the least
$x\ge\max(e,e^{2\mathfrak p_0})$ for which
$2\varepsilon(x)\le\varepsilon_{\mathcal W}\alpha_0$.
Let $x_{\rm J}$ be the least $x\ge x_{\rm J1}$ such that
\begin{equation}\label{eq:weighted-history-threshold}
\begin{gathered}
 x\ge\max(4\bar B,eL),\qquad
 C_0^\flat\mathsf S_A\theta(x)\le1,\qquad
 \tfrac43C_1^\flat L\mathsf T(3)\theta(x)\le1,\\
 (1+\beta_{\rm s})\delta(x)/\alpha_0\le s^\star/2,
 \qquad\varepsilon(x)\le c_A\alpha_0/(8MB_3).
\end{gathered}
\end{equation}
Here the constants and $\mathsf T$ are those of
Section~\ref{sec:weighted-progression-interface},
and $c_A$ is the regular-background radius constant of
equation~\YMMref{(2262)}{828}.
Both least values exist. Indeed, $\varepsilon$, $\delta$
and $\theta$ are positive constant multiples of
$x^{-1/2}(\log x)^{\mathfrak p_0}$, continuous and
nonincreasing for $x\ge e^{2\mathfrak p_0}$, and tending
to zero. The series defining $\mathsf S_A$ converges by
the ratio test. Since $3L^{n-1}-2\ge n$, the same test
gives $\mathsf T(3)<\infty$. Each inequality therefore
defines a closed terminal half-line; their finite
intersection has a least point. All the inequalities
remain true above the resulting threshold.

\begin{proposition}\label{prop:weighted-canonical-history}
Fix one admitted assignment of the canonical constants,
with the weighted old-potential constant and auxiliary data
of Definition~\ref{def:weighted-progression-data}. Let
$x_j=x_j(\beta)$ be its real canonical run, and choose
$2\le K\le j_{\rm fp}(\beta)-1$ and an admitted frame
$(K,(x_0,\ldots,x_K),\bar{\mathfrak w})$ on
$\mathfrak w=(2L^m,K)$. For $1\le k\le K-1$ and every
finite list of sources, (J.1), (J.2), including (Pv), and
(J.4) hold. If $x_{k+1}\ge x_{\rm J1}$, then all of (J)
holds. If $x_{k+1}\ge x_{\rm J}$, the running comparison,
half-ceiling and regular-background radius hypotheses of
Propositions~\YMMref{122.1}{826}, \YMMref{122.2}{827}
and equation~\YMMref{(2264)}{828} hold at every earlier
height; the mixed estimate still requires the displayed
positive separation of the two supports.
\end{proposition}

\begin{proof}
The common choice and the actual-step verification in
Propositions~\YMMref{69.1}{595} and~\YMMref{69.3}{603}
give the canonical auxiliary data, weighted radius bounds,
octet condition and all the test hypotheses in (J.4).
Theorem~\YMMref{D}{8} supplies the regular progression
at every initial segment of the run. Restricting a longer
progression to such a segment does not change its earlier
families or counterterms, since the recursion uses only
the preceding members. Thus its $e_j$ are the $e_j$ in
the running recursion, which proves (J.1).

At height $j$, the canonical shift and fluctuation maps
are the same maps used to define the actual place, merely
restricted to its core. This proves the first part of
(Pv). The Gaussian summand in the regular output and the
one in the sum rule are both the localized Gaussian
family at that height, on the same lattice and with the
same auxiliary data. They agree term by term. The actual
place verification also gives
$\alpha_0<a_{\rm pg}^\star\le a_{\rm H}$ and
$(1+\beta_{\rm s})\delta(x_j)/\alpha_0\le s_\omega^\star$.
With $s_\omega^\star\le1$ and
$a_{\rm pg}^\star\le\varepsilon_{\mathcal W}/2$, this
implies $\delta(x_j)\le\varepsilon_{\mathcal W}$.
These are the remaining requirements of (J.2).

The frame restricts to every earlier height. Its run
decreases and satisfies
$|x_j-L^{-(j-i)}x_i|\le\bar B$. Hence every earlier
$x_j$ is at least $x_{k+1}$. The definition of $x_{\rm J1}$
now gives (J.3). If $x_{k+1}\ge x_{\rm J}$, all the
inequalities in \eqref{eq:weighted-history-threshold}
hold at each earlier coefficient. The first gives
$x_k\ge4\bar B$ and $\log x_k\ge1+\log L$; the next
three give the stated running and half-ceiling bounds.
The last gives
$c_A\alpha_0/(2MB_3\varepsilon_t)\ge4$. Together with
the lattice room $L^{m+K-t}\ge LM\ge8M+4$, this makes
the radius in equation~\YMMref{(2262)}{828} at least three.
For two supports with $D_{k^\flat+1}\ge3M$,
$D_j=L^{k^\flat+1-j}D_{k^\flat+1}\ge3M$ at earlier
heights. Monotonicity of $\mathsf T$ therefore gives
\[
 \tfrac43C_1^\flat L\theta(x_j)\mathsf T(D_j/M)
 \le\tfrac43C_1^\flat L\theta(x_j)\mathsf T(3)\le1,
\]
which is the remaining mixed-estimate hypothesis.
\end{proof}

\subsection{Regular backgrounds and the two-source tower}
\label{sec:weighted-calculus-interface}

We need two consequences of the local step. First, a family evaluated
at a regular background differs little from its value at the identity.
Second, integrating one fluctuation carries the full source exponential
to the next level. The second statement concerns the normalized integral;
it is stronger than the corresponding identity for its first derivative.

\begin{lemma}[A pointed family at a regular background]
\label{lem:weighted-regular-family}
Use the constants and absolute radius systems of
Definition~\ref{def:weighted-history-setting}. Let
$1\le j\le T\le k+1$, let $\varepsilon>0$, and let $U$ be a
real bare configuration regular of level $T$ at radius $\varepsilon$:
\[
 |\partial U-\mathbf1|<\varepsilon\eta_T^2,
 \qquad |J^{(T)}(U)|<\varepsilon.
\]
Suppose that the actions of the bare arrays through level $T$ are
defined at $U$. Let $f$ be a pointed family of level $j$, vanishing
on domains disjoint from a set $\mathsf{Cr}$ of bare sites. Put
\begin{equation}\label{eq:weighted-regular-radius}
 d_*(j;T,\varepsilon)=
 \min\left\{\frac{c_A\alpha_0L^{2(T-j)}}{M\varepsilon},
                  \frac{L^{m+K-j}-4}{2M}\right\}-1.
\end{equation}
Then
\begin{equation}\label{eq:weighted-regular-family}
\begin{split}
 &\sum_{X\in\mathcal D_j}\sum_{z\in X\cap\mathbb T^{(j)}}
                       |f(X,U,z)-f(X,\mathbf1,z)|\\
 &\quad\le8M^3\mathsf n^{\rm cb}_j(\mathsf{Cr})
          \|f\|_{j,\lambda_{\rm out}}
 \left[C_AC_{\mathrm E,3}M^2
             \frac{\varepsilon^2}{\alpha_0^2}L^{-4(T-j)}
       +2C_{\mathrm E,1}e^{-[d_*(j;T,\varepsilon)]_+}\right].
\end{split}
\end{equation}
Here $\mathsf n^{\rm cb}_j(\mathsf{Cr})$ counts the level-$j$
cubes meeting $\mathsf{Cr}$, and $C_A,c_A$ are the constants
in the local regular-background estimate of
Section~\YMMref{122}{825}.
\end{lemma}

\begin{proof}
Write $d=d_j(X)$ and $\vartheta=L^{-(T-j)}$.
The regular-background estimate in equation~\YMMref{(2263)}{828}
applies to the member $f(X,\cdot,z)$ when
\[
 2(1+d)M+4<L^{m+K-j},\qquad
 (1+d)M\varepsilon L^{-2(T-j)}\le c_A\alpha_0.
\]
Its contribution is at most
\[
 C_A\|f\|_{j,\lambda_{\rm out}}e^{-\lambda_{\rm out}d}
       M^2(1+d)^2\frac{\varepsilon^2}{\alpha_0^2}L^{-4(T-j)}.
\]
The radius hypotheses required there, including
$M\alpha_0\le\alpha_1$ and $\alpha_0+\alpha_1\le a^\sharp$,
are part of the fixed setting. If either geometric inequality
fails, then $d\ge[d_*(j;T,\varepsilon)]_+$. Both values still
belong to their cores, by the hypothesis on the array actions,
so their difference is at most
$2\|f\|_{j,\lambda_{\rm out}}e^{-\lambda_{\rm out}d}$.

There are $M^3\#X\le8M^3(d+1)$ base points in $X$.
Every contributing domain contains a cube meeting $\mathsf{Cr}$;
choose one lattice site in each such cube. The cube count in the
proof of Lemma~\ref{lem:source-support-sum} is
\[
 \sum_{X\ni z,\ d_j(X)\ge D}
      e^{-\lambda_{\rm out}d_j(X)}(d_j(X)+2)^p
           \le C_{\mathrm E,p}e^{-D}.
\]
The setting includes $\lambda_{\rm out}>\kappa_{\mathcal D}+2$.
For the first kind of domain use $(d+1)^3\le(d+2)^3$
and $D=0$. For the second use $d+1\le d+2$ and
$D=[d_*]_+$. Summing over the cubes meeting $\mathsf{Cr}$
gives \eqref{eq:weighted-regular-family}.
\end{proof}

\begin{lemma}[Normalized step integrals]
\label{lem:weighted-step-integrals}
At an actual step of height $i$ in
Definition~\ref{def:weighted-history-setting}, let
$P\in\mathcal U^\circ_{i+1}(T_\ell;\alpha_0,\alpha_1)$.
The integral defining $\Xi(x_i;P)$ is absolutely convergent and
nonzero. Its normalized measure $\nu_i[P]$ has finite total
variation and mass one. For every bare array $g$ through level
$i$, the function $B\mapsto\mathbb V'[g]^{\rm tot}(P;B)$ is
bounded and continuous on $\mathsf{Rg}_{x_i}$, with
\begin{equation}\label{eq:weighted-total-bound}
 |\mathbb V'[g]^{\rm tot}(P;B)|
       \le\#\mathcal D_i\,K_\Lambda s_iN_i(g).
\end{equation}
The first and mixed insertion identities are
equation~\YMMref{(2298)}{837}. For a real chart point of norm
less than $r_\alpha$, the normalized measure is a probability
measure by Proposition~\YMMref{124.5}{836}.
\end{lemma}

\begin{proof}
The rule sending an array to the fourth potential bounds each
of its local members by $K_\Lambda s_iN_i(g)$ on the step
domain. Summing over the finitely many domains gives
\eqref{eq:weighted-total-bound}. Each member is holomorphic
in the fluctuation field there, hence continuous on the real box.
To evaluate the normalization, insert $\mu g$ in the source
slot and restrict $\mu$ so that its array norm is smaller
than the remaining room. The whole-torus identity of
Theorem~\YMMref{V}{416}, equation~\YMMref{(902)}{417},
expresses the integral as the exponential
of the sum of the localized output, so it is nonzero and
absolutely convergent. Here $\tau=1$ removes the additional
quadratic potential. At $\mu=0$, division by this nonzero
integral gives mass one and finite total variation. The bound
\eqref{eq:weighted-total-bound} justifies differentiation
of the normalized integral; the first derivative of its
logarithm is the mean of the inserted total, and the mixed
derivative is its covariance. This is precisely the argument
for equation~\YMMref{(2298)}{837}. Positivity is asserted only
at the real chart points covered by
Proposition~\YMMref{124.5}{836}.
\end{proof}

For the tower identity we use the real two-source frame of
Lemma~\ref{C-352-11}(a), retaining its actual canonical places,
flat-top hypotheses and weighted progression. The chart, lowering
map and chain classes are those of Section~\YMMref{124}{832}.
In particular, a member of $\mathsf{Cc}_n(r)$ is a real bare
potential $A$ satisfying
\begin{equation}\label{eq:weighted-chain-class}
 \|v_A^{[n]}\|\le r,\qquad
 v_A^{[n]}=(A,J^{(n)}(e^{i\eta_nA})),\qquad
 e^{i\eta_nA}=U^u,\quad
 U\in\mathcal C(\mathbf B_n,W),\quad
 \varpi(W)<\varepsilon_{\mathcal W}\alpha_0.
\end{equation}
Thus the norm condition is accompanied by an actual background
condition. The two radii in equation~\YMMref{(2285)}{833} are
denoted here by $\varrho_{\rm ch}$ and $\varrho_{\rm fl}$.

\begin{lemma}[The domains of a chain background]
\label{lem:weighted-chain-background}
In this frame, let $1\le n\le K$, $r\ge0$ and
$A\in\mathsf{Cc}_n(r)$. The array actions through level $n$
are defined at $\widetilde U=e^{i\eta_nA}$. If, in addition,
$r<r'\le\varrho_{\rm ch}$, then $\widetilde U$ is regular
of level $n$ at radius $3r'$, and
\[
 P[v_A^{[n]}]\in\mathcal U^\circ_n(T_\ell;\alpha_0,\alpha_1),
 \qquad 0\le a_p(\widetilde U)<\tfrac92(r')^2\eta_n^4.
\]
If $r=0$, then $A=0$ and $\widetilde U=\mathbf1$.
\end{lemma}

\begin{proof}
Under the additional radius condition the chart norm is at most
$r<r'$. Its image lies in the simple domain with radii
$(3r',r',r')$, and in the core because
$r'\le r_\alpha$. This is the usable-chart construction in
equation~\YMMref{(933)}{427} and the domain verification
preceding equation~\YMMref{(2287)}{834}. Hence
\[
 |\partial\widetilde U-\mathbf1|<3r'\eta_n^2,
 \qquad |J^{(n)}(\widetilde U)|<r',
\]
which gives regularity. To check the earlier array actions, use
the background $U$ and block field $W$ in
\eqref{eq:weighted-chain-class}. The small-plaquette clause of
Proposition~\YMMref{52.3}{433} puts $W$ in the actual step
domain. At height $n-1$ the canonical sum rule has its array
actions defined at every representative of this orbit. This
gives all levels through $n-1$; at level $n$ the plain-core
clause in the proof of the same proposition supplies the
remaining pairs. At $n=1$ there are no earlier levels to check.
The trivial pairs lie in their cores as well. These are exactly
the domain conditions for the absolute array actions. Their
domains are gauge invariant, so they hold also at $U^u=
\widetilde U$.

For an $SU(2)$ plaquette,
$a_p(\widetilde U)=\tfrac12|\partial\widetilde U(p)-\mathbf1|^2$;
the regularity bound gives the asserted estimate. Finally the
chart norm dominates $|A|$, so $r=0$ forces $A=0$.
\end{proof}

Write $\mathsf{Sum}[f](U)$ for the finite sum over the domains
and base points of the family $f$, evaluated at the pairs of $U$.
For $w_\lambda=\lambda_1w_1+\lambda_2w_2$ and
$\lambda\in\Lambda^{\rm w}=\Lambda^{\rm w}_{K-1}$, set
\begin{equation}\label{eq:weighted-relative-functions}
\begin{gathered}
 D_i^\lambda(U)=\mathcal O(w_\lambda)(U)
       +\sum_{j=1}^i\mathsf{Sum}[E^{(j)}_\lambda-E^{(j)}](U),
 \qquad \mathcal O(w)(U)=L^K\sum_pw(p)a_p(U),\\
 \Phi_{a,i}=\partial_aD_i^\lambda|_0,
 \qquad \mathsf F_i=\partial_1\partial_2D_i^\lambda|_0,\\
 \Phi^{\rm ir}_{a,i}(U)=\sum_{j=1}^i
       \bigl(\mathsf{Sum}[\dot E^{(j)}_a](U)
                       -\mathsf{Sum}[\dot E^{(j)}_a](\mathbf1)\bigr).
\end{gathered}
\end{equation}
These functions are defined when the array actions through level
$i$ are defined at $U$. Then the relative exponent is holomorphic
in $\lambda$ and $D_i^0=0$, by
Proposition~\ref{prop:weighted-finite-system}.

\begin{proposition}[The two-source tower]
\label{prop:weighted-pair-tower}
In the frame just specified, let $1\le i<K$ and $r\ge0$ satisfy
\begin{equation}\label{eq:weighted-tower-smallness}
 r<\varrho_{\rm ch},\qquad
 \delta_i<\varrho_{\rm fl},\qquad
 6L^{-2}r+8C_{\rm db}\delta_i<\varepsilon_{\mathcal W}\alpha_0.
\end{equation}
For $A\in\mathsf{Cc}_{i+1}(r)$ put
$\widetilde U=e^{i\eta_{i+1}A}$, $v=v_A^{[i+1]}$,
$P=P[v]$ and $\nu=\nu_i[P]$. For
$B\in\mathsf{Rg}_{x_i}$ put
\[
 \widetilde U^{\rm lw}(B)
        =e^{i\eta_iA^{\rm lw}(v,g_iB)}.
\]
Then $\nu$ is a probability measure. The relative exponents and
their derivatives in \eqref{eq:weighted-relative-functions} are
defined at $\widetilde U$ through level $i+1$ and at
$\widetilde U^{\rm lw}(B)$ through level $i$. Their latter
values, as well as $\mathcal O(w_a)(\widetilde U^{\rm lw}(B))$,
are bounded continuous functions of $B$. For every
$\lambda\in\Lambda^{\rm w}$,
\begin{equation}\label{eq:weighted-pair-tower}
 e^{D_{i+1}^\lambda(\widetilde U)}
     =\int e^{D_i^\lambda(\widetilde U^{\rm lw}(B))}\,d\nu(B).
\end{equation}
Consequently,
\begin{equation}\label{eq:weighted-pair-tower-derivatives}
\begin{split}
 \Phi_{a,i+1}(\widetilde U)
    &=\int\Phi_{a,i}(\widetilde U^{\rm lw})\,d\nu,\\
 \mathsf F_{i+1}(\widetilde U)
    &=\int\mathsf F_i(\widetilde U^{\rm lw})\,d\nu
       +\big\langle\Phi_{1,i}(\widetilde U^{\rm lw});
                         \Phi_{2,i}(\widetilde U^{\rm lw})\big\rangle_\nu.
\end{split}
\end{equation}
\end{proposition}

\begin{proof}
Proposition~\YMMref{124.2}{834} applies under
\eqref{eq:weighted-tower-smallness}. It gives
\[
 A^{\rm lw}(v,g_iB)\in
  \mathsf{Cc}_i(L^{-1}r+\tfrac{11}{10}c_{\mathfrak h}\delta_i),
\]
and Lemma~\ref{lem:weighted-chain-background} puts $P$ in
the level-$(i+1)$ core and verifies the array actions at both
backgrounds. Restriction to earlier
levels preserves this definedness: it retains a subset of the
required core conditions. The chart point is real and has norm
less than $r_\alpha$, so Proposition~\YMMref{124.5}{836} gives
the asserted positivity and normalization of $\nu$.

The source enters only through the bare array. At height $i$,
the difference $g_\lambda=f_{\lambda,i}-f_{0,i}$ has zero
counterterms and equals
\[
 g_\lambda=((E^{(j)}_\lambda-E^{(j)},0)_{j\le i},\emptyset)
                 \sqcup\mathsf W_{L^{K-i}w_\lambda}.
\]
Indeed $E^{(j)}_0=E^{(j)}$ and the counterterms were kept fixed.
The action of a zero-counterterm array is the sum of its pointed
totals minus their values at $\mathbf1$. Its weight member adds
$\eta_i^{-1}\sum_pL^{K-i}w_\lambda(p)a_p(U)=\mathcal O(w_\lambda)(U)$.
Since the plaquette action vanishes at $\mathbf1$, this gives
\[
 S[g_\lambda](U)=D_i^\lambda(U)-D_i^\lambda(\mathbf1).
\]
Apply the array sum rule, equation~\YMMref{(2289)}{834}, to
$g_\lambda$. The constant at $\mathbf1$ cancels, yielding
\begin{equation}\label{eq:weighted-tower-action}
 \mathbb V'[g_\lambda]^{\rm tot}(P;B)
      =D_i^\lambda(\widetilde U^{\rm lw}(B))
                            -D_i^\lambda(\widetilde U).
\end{equation}
The same calculation for the first and mixed derivatives of
$g_\lambda$, and for its weight member alone, gives the analogous
identities for $\Phi_{a,i}$, $\mathsf F_i$ and $\mathcal O(w_a)$.
Lemma~\ref{lem:weighted-step-integrals} makes these potential
totals bounded and continuous on the fluctuation box; hence
so are the listed functions. Finally
$\Phi^{\rm ir}_{a,i}=\Phi_{a,i}-\mathcal O(w_a)
-\sum_{j\le i}\mathsf{Sum}[\dot E^{(j)}_a](\mathbf1)$,
which gives the remaining continuity assertion.

Let $\Xi_\lambda(P)$ be the step integral with array
$f_{\lambda,i}$ and with all other data unchanged. The finite-source
estimate puts this array in the admitted norm ball. The actual
one-step theorem therefore gives an absolutely convergent integral
and the identity
\begin{equation}\label{eq:weighted-tower-integral}
 \Xi_\lambda(P)
   =\exp\mathsf{Sum}[\mathsf{St}_i(f_{\lambda,i})](P)\ne0.
\end{equation}
Here the calibration equals the running coefficient, so $\tau=1$
and the additional quadratic potential, bounded by a constant times
$|1-\tau|$, is zero. The whole torus is an admitted localization
domain, and the output radii are $(\alpha_0,\alpha_1)$. These are
the hypotheses for the whole-torus identity of Theorem~\YMMref{V}{416};
the same specialization is used in
Proposition~\ref{prop:weighted-finite-system}.

At zero source this integral is precisely the normalization of
$\nu_i[P]$, equation~\YMMref{(2293)}{836}. The provision is linear
in its bare array, so division of the two integrals and
\eqref{eq:weighted-tower-action} give
\[
 \frac{\Xi_\lambda(P)}{\Xi_0(P)}
  =e^{-D_i^\lambda(\widetilde U)}
       \int e^{D_i^\lambda(\widetilde U^{\rm lw}(B))}\,d\nu(B).
\]
On the other hand the recursion
\eqref{eq:weighted-progression-recursion} and
\eqref{eq:weighted-tower-integral} give
\[
 \frac{\Xi_\lambda(P)}{\Xi_0(P)}
   =\exp\mathsf{Sum}[E^{(i+1)}_\lambda-E^{(i+1)}](P)
   =\exp\bigl(D_{i+1}^\lambda(\widetilde U)
                                  -D_i^\lambda(\widetilde U)\bigr).
\]
The sum at $P$ equals the sum at its pair, which is the pair of
$\widetilde U$. Comparing the two quotients proves
\eqref{eq:weighted-pair-tower}.

To differentiate this identity, the pointwise holomorphy in the
sources needs a bound uniform in the fluctuation field. Define the
finite numbers
\[
 K_{i,a}=2L^K\sum_p|w_a(p)|
       +\sum_{j\le i}\sum_{X\in\mathcal D_j}
                         \#(X\cap\mathbb T^{(j)})d_{j,a},
\]
where $d_{j,a}$ are the Lipschitz bounds in
\eqref{eq:weighted-progression-numbers}. All pairs entering the
lowered array actions lie in their cores. The family norm bounds
each of their values, and $0\le a_p(U)\le2$ on real configurations.
Consequently
\[
 |D_i^\lambda(\widetilde U^{\rm lw}(B))|
                  \le\sum_{a=1}^2|\lambda_a|K_{i,a}
 \quad(\lambda\in\Lambda^{\rm w},\ B\in\mathsf{Rg}_{x_i}).
\]
Thus the exponential is uniformly bounded on each compact source
set, independently of $B$. Cauchy's formula and integration against
the probability measure justify differentiation. Since $D_i^0=0$,
the first derivatives give the first identity in
\eqref{eq:weighted-pair-tower-derivatives}. The mixed derivative
gives the integral of
$\mathsf F_i+\Phi_{1,i}\Phi_{2,i}$. Subtracting the product
of the two first-derivative identities leaves the covariance
in its second identity.
\end{proof}

\subsection{The weighted step and its derivatives}

The derivative estimates used below concern the actual local step.
In the notation of \cite[equations~(2249)--(2250)]{YMM1}, put
\[
 s_j=\frac{(1+\beta_s)\delta(t_j)}{\alpha_0},\qquad
 \theta_j=K_{\mathbf M}s_j,\qquad
 N_{\star,j}=\frac{w_\star}{4K_\Lambda s_j}.
\]
Here $t_j\ge\gamma^{-2}$ is the calibration, the coefficients satisfy
$|x_j/t_j-1|<\sigma_\star$, $N_j$ is the weighted bare-array norm at
height $j$, and $\|\cdot\|_{j,+}$ is the pointed-family norm at rate
$\kappa^+=(1+4\beta_s)\kappa$. With the same canonical provision and
counterterms at every parameter value, the local step obeys
\begin{equation}\label{upstream-weighted-step}
 \|\mathsf{St}_j(f)-\mathsf{St}_j(f')\|_{j+1,+}
     \le\theta_jN_j(f-f')
 \qquad\bigl(N_j(f),N_j(f')\le N_{\star,j}\bigr).
\end{equation}
The map is jointly holomorphic with the field variables and is real on
real data. Its member on a domain $X$ depends only on the restriction
of the input array to $X$; the step uses bracket localisation before
applying the cluster map. At $N_j(f)\le N_{\star,j}/2$,
\begin{equation}\label{upstream-weighted-derivatives}
 \begin{split}
 \|D\mathsf{St}_j(f)g\|_{j+1,+}
       &\le\theta_jN_j(g),\\
 \|D^2\mathsf{St}_j(f)[g_1,g_2]\|_{j+1,+}
       &\le\frac{2\theta_j}{N_{\star,j}}N_j(g_1)N_j(g_2).
 \end{split}
\end{equation}
For completeness, the fourth-potential and cluster difference bounds
give the coefficient
\[
162L^3K_\Pi K_\Lambda s_j
\]
in
\eqref{upstream-weighted-step}; the choice
$K_{\mathbf M}\ge243L^3K_\Pi K_\Lambda$ gives the displayed bound.
The first derivative follows by a difference quotient. For the second, if $N_j(g_2)=0$, the first derivative bound vanishes in this direction at every admissible point, so the mixed second derivative is zero. Suppose $N_j(g_2)>0$;
apply the first estimate at $f+zg_2$ and use Cauchy's formula on the
circle 
\[
|z|=N_{\star,j}/(2N_j(g_2)),
\]
with a limit from smaller circles
at the boundary. The Cauchy formula for each local member, together
with the uniform weighted bounds, also gives the ordinary chain rule
for holomorphic arrays. These estimates retain their stated array
domains; a different source or coupling must be put into these domains
before they can be applied.

\subsection{The lattice five-point inequality}

The symmetries and reflection positivity of a periodic Gibbs state
are \cite[Proposition~103.4]{YMM1}. Its coordinate reflection pictures
and the lattice spectral inequality are
\cite[Lemma~99.1 and Proposition~99.3]{YMM1}. The following consequence
records the precise lattice inequality used in the continuum cone
argument. We include the passage from one lattice step to every integer power.

\begin{proposition}\label{upstream-lattice-five-point}
Let $a>0$, $\beta\ge0$ and $\nu\in\mathcal G^{\rm per}_{a,\beta}$.
Choose distinct directions $\mu,\lambda$, a coordinate mirror
$x_\mu=c_0$ with $c_0\in\tfrac a2\mathbb Z$, and a bounded Borel
local function $F$ of the bonds in its positive half-space. Set
\[
 u(w)=\nu\bigl(\overline{\theta_{e_\mu,c_0}F}\,\tau_wF\bigr),
 \qquad w\in a\mathbb Z^3.
\]
Then, for every $\mathsf t\in a\mathbb N$,
\begin{equation}\label{IV-54-29}
 u(0)-4u(\mathsf te_\mu)+u(2\mathsf te_\mu)
   +u(\mathsf te_\mu+\mathsf te_\lambda)
   +u(\mathsf te_\mu-\mathsf te_\lambda)\ge0.
\end{equation}
\end{proposition}

\begin{proof}
Let $T$ and $V$ be the normal transfer contraction and tangential
translation unitary in this picture. Their joint spectral measure is
supported where
\[
 (1-t)^2\ge t|1-\zeta_\lambda|^2.
\]
For $t=e^{-x}>0$ and $\zeta_\lambda=e^{iy}$ this is
$\sinh(x/2)\ge|\sin(y/2)|$. For every integer $k\ge1$,
\[
 \sinh(kx/2)\ge k\sinh(x/2)
       \ge k|\sin(y/2)|\ge|\sin(ky/2)|.
\]
The first and last inequalities follow by induction from the addition
formulas. Thus the spectral support also satisfies
\[
 (1-t^k)^2\ge t^k|1-\zeta_\lambda^k|^2;
\]
at $t=0$ this remains true. Functional calculus gives
\[
 (1-T^k)^2\ge T^k\bigl(2-V(kae_\lambda)-V(kae_\lambda)^*\bigr).
\]
Take its quadratic form on $[F]$. The picture identity
$\langle[F],T^jV(v)[F]\rangle=u(jae_\mu+v)$ identifies its five
terms with \eqref{IV-54-29}, where $\mathsf t=ka$.
\end{proof}

\subsection{The calibrated cutoff comparison}\label{sec:cutoff-calibration}

We record the choices used in the comparison of two successive cutoffs.
They must be made before fixing the torus or its physical coupling.
The setting has three parts, denoted by \textup{(Th1)}--\textup{(Th3)}.

In \textup{(Th1)}, choose the real construction constants and canonical
auxiliary data as in Lemma~\YMMref{118.1}{812}. Impose the global
background, response, localization and actual-domain hypotheses of
Section~\YMMref{113}{786}, and the source and recursion hypotheses of
Section~\YMMref{114}{799}. In particular the fixed pair radii are
$(\alpha_0,\alpha_1)$, $\alpha_0\le a_{\rm H}$, the lag weights are
$L^{-(j-i)}$, and
\[
 E_0^{\rm s}=E_0,\qquad
 N_\star^{\rm s}=\frac L{L-1}
                 (E_0+\tfrac32\alpha_0^2B_\sharp).
\]
The constants $K_{\rm Q}^\vee,K_{\rm L}^\vee,\mu_\star,
\sigma_{\rm src}$ and the three smallness inequalities are those
in equations~\YMMref{(2170)}{801} and the surrounding source construction.
Take
\[
 \varrho=\min\left\{
 \frac{20\mu_\star}{\sqrt3K_{\rm L}^\vee\gamma},
 \frac{\sigma_{\rm src}}2(\gamma^{-2}-\bar B)\right\}>0,
 \qquad
 \mathsf{Src}=\mathsf{Src}_{\mathsf{Bc}}(2\varrho,2\varrho).
\]
Thus the coupling sequence, Gaussian structures, auxiliary data and
pair cores are independent of the source.

In \textup{(Th2)}, fix $0<\eth<1$ before decreasing $\gamma$, and use
the tail calibration of Section~\YMMref{115}{805}:
\[
 c'=(1-\eth)c_\star,\qquad
 C_{\rm maj}^{\Sigma}=C_{\rm maj}+(6L^3+1)\log2,\qquad
 C_{\rm low}(x)=c_{\rm low}+\tfrac92\log x.
\]
The ceiling $\gamma$ is chosen so that $\log\gamma^{-2}$ exceeds
both the case-II history threshold and the summability threshold
specified there. These thresholds, the source smallness conditions,
and the background and response conditions are part of the
calibration. The proof of Lemma~\YMMref{118.1}{812} verifies that they
can be imposed simultaneously. None depends on the finite lattice.

For later use put $c_{\log}=\log L/\log\gamma^{-2}$ and
\[
 c_\rho=\max\left\{2B_3,
       \frac{C_HA_1}{(1-L^{-3/2})A_0}\right\}.
\]
We write $\mathsf t_\natural$ for the fixed calibration denoted by
$\tau_\natural$ in equation~\YMMref{(2134)}{790}. It is the least
$\tau\ge\gamma^{-2}$ at which
\[
 \begin{gathered}
 2c_\rho\varepsilon(\tau)
 \le\min\{\varepsilon_{\mathcal W}\alpha_0,c_2^{\rm SO},
                                  \varepsilon_{\rm v}/4\},\\
 \delta(\tau)\le\varepsilon_{\mathcal W},\qquad
 4C_{\rm db}\delta(\tau)
       \le\min\{1,r_{\rm B},r_{\rm B}^{\rm s}/2\}.
 \end{gathered}
\]
The thresholds are continuous, decreasing and tend to zero;
their positive fixed ceilings therefore make this minimum finite.
All three inequalities persist at larger $\tau$.

In \textup{(Th3)}, fix
\[
 K\ge2,\quad \mathsf b'\ge\gamma^{-2},\quad m\in\mathbb N,\quad
 MR(\mathsf b')\le L^m,\quad \ell=2L^m,\quad
 x_j=L^{K-j}\mathsf b'\quad(0\le j\le K).
\]
Choose an integer top level $1\le t_*\le K-1$ such that
$x_{t_*}\ge\mathsf t_\natural$. These are precisely the
finite-lattice data in Section~\YMMref{115}{805}. The setting
\textup{(E)} retains \textup{(Th1)}, \textup{(Th2)}, the label and
the side, allowing $K$ to vary. A compact label interval
$J=[b_-,b_+]$ is admissible on this side when
$b_-\ge\gamma^{-2}$ and $MR(b_+)\le L^m$, as in
equation~\YMMref{(2201)}{810}.

For a comparison involving $n$ insertions, the depths are
$u=\lceil K/[180(3n+4)]\rceil$ and $s=\lceil K/9\rceil$.
The errors and smallness conditions are those in
Section~\YMMref{116}{807}; the additional support condition is
$d_\infty(\underline f)L^u\ge6M$ when $n\ge2$.
For bounded functions of block fields one takes $n=1$ and omits
this support condition. Section~\ref{sec:block-field-laws} displays
all the resulting errors and conditions \textup{(k0)}--\textup{(k4)},
\textup{(k6)}.

\begin{lemma}\label{lem:cutoff-admissible-labels}
Let $0<g\le\gamma$ and $\ell=2L^m\in\mathfrak L(g)$. Then
\[
 J_g=[g^{-2}+B_\sharp-L^{-1}\bar B,\ X_g+\bar B]
\]
is an admissible label interval on $T_\ell$, and it contains $I(g)$
and every first-passage label at $g$. The larger interval
$\mathsf J_g=[\gamma^{-2},L(X_g+\bar B)]$ is also admissible
on $T_\ell$.
\end{lemma}
\begin{proof}
The lower endpoint is at least $\gamma^{-2}$, since
$B_\sharp-L^{-1}\bar B=B_\sharp(L-2)/(L-1)\ge0$.
The side floor in equation~\YMMref{(1528)}{615} gives
$L^m\ge LM R(X_g+\bar B)$, and hence the required room
$MR(\max J_g)\le L^m$. The containment of $I(g)$ follows from
its endpoints. The first-passage label bound is
equation~\YMMref{(1531)}{615}.
For the larger interval, the radius scaling inequality in the proof
of Proposition~\YMMref{74.1}{615} gives
$R(LX)\le LR(X)$ for $X\ge\gamma^{-2}$. Apply it at
$X=X_g+\bar B$ and use the same side floor to obtain
$MR(\max\mathsf J_g)\le LMR(X_g+\bar B)\le L^m$.
\end{proof}

\subsection{The laws of the block fields}\label{sec:block-field-laws}

The cutoff comparison also controls bounded functions of a fixed block
field. We prove the resulting total-variation convergence here. We use
the calibrated small-field setting of Section~\YMMref{115}{805}: all constants,
auxiliary data, source radii and tail thresholds are fixed before the
cutoff, with $E_0^{\rm s}=E_0$,
$N_\star^{\rm s}=L(E_0+3\alpha_0^2B_\sharp/2)/(L-1)$,
$\mathsf b\ge\gamma^{-2}$, $MR(\mathsf b)\le L^m$ and $\ell=2L^m$.
This is the setting denoted by \textup{(E)}. The precise calibration,
including the source and tail smallness conditions, is that of
Proposition~\YMMref{115.2}{806}, not merely the displayed lower bound on
$\mathsf b$.

At cutoffs $K$ and $K+1$ use the same physical depth, and put
\[
 u=\lceil K/1260\rceil,\qquad s=\lceil K/9\rceil,
 \qquad t_1=K-u,\qquad t_0=K-s.
\]
\[
 x_{[d]}=L^d\mathsf b,\qquad N_{[d]}=\ell^3L^{3d},
 \qquad y_{[d]}=(\log x_{[d]})^{2\mathfrak p_0}.
\]
The errors in Proposition~\YMMref{116.1}{808} and the discarded mass in
Proposition~\YMMref{115.2}{806} are
\begin{align*}
 \epsilon_{\rm fg}&=C_{\rm old}C_{\rm fg}\frac L{L-1}E_0
                           N_{[u]}L^{-(s-u)},\\
 \epsilon_{\rm cl}&=4\mathsf K_{\rm w}c_\rho^2A_0^2
    (1+2\varrho/\mathsf b)y_{[K]}N_{[s]}\sqrt K L^{-(K-s)/2},\\
 \epsilon_{\rm sh}&=3(s-u)N_{[s]}
    \exp[C_{\rm sh}-\kappa_{\rm b}A_1^2y_{[u]}/4],\\
 \mathsf R_{t_1}&=2M^{-3}N_{[u]}
    \exp[-c'A_0^2y_{[u]}+(\mathsf C^\Sigma_{\rm maj}+2
                                      +C_{\rm low}(\mathsf b))\ell^3].
\end{align*}
Here $c'=(1-\eth)c_\star>0$, with the fixed $0<\eth<1$ of that
calibration. We retain the following names for the comparison conditions:
\begin{gather*}
 \textup{(k0)}\quad1\le u<s<K,\qquad
 \textup{(k1)}\quad L^u\mathsf b\ge\Theta_\star,\\
 \textup{(k2)}\quad\epsilon_{\rm fg}\le1/8,\qquad
 \textup{(k3)}\quad\epsilon_{\rm cl},\epsilon_{\rm sh}\le1/20,\\
 \textup{(k4)}\quad256\mathsf K_uc_\rho^2A_0^2N_{[s]}L^{-(K-s)}
                         \le\gamma_{\rm r}A_1^2,\qquad
 \textup{(k6)}\quad\mathsf R_{t_1}\le1/2.
\end{gather*}
There is no support-separation condition in this measure statement.
The constants $\Theta_\star$, $\mathsf K_u$, $\mathsf K_{\rm w}$
and the other comparison constants are those in
Sections~\YMMref{116}{807} and~\YMMref{115}{805}, with identical level
units. In particular the block lattice at depth $d$ is
$\Lambda_{d,\ell}$ on both cutoffs.

For uniformity in $\mathsf b$, let
$\mathsf J=[\mathsf b_-,\mathsf b_+]$ satisfy
$\mathsf b_-\ge\gamma^{-2}$ and $MR(\mathsf b_+)\le L^m$.
Define $\bar\epsilon_{\rm sh}$ by replacing $\mathsf b$ by
$\mathsf b_-$ in $y_{[u]}$, and define $\bar{\mathsf R}$ by that
same replacement in its negative exponent and by replacing
$C_{\rm low}(\mathsf b)$ with $C_{\rm low}(\mathsf b_+)$.
They dominate their unbarred counterparts. The uniform threshold
$K_0$ in Theorem~\YMMref{AR}{810}, taken at $n=1$, makes (k0)--(k4) and
(k6) valid for all $K\ge K_0$. Its proof shows that
$K^ae^{a'K}(\bar{\mathsf R}+\bar\epsilon_{\rm sh})\to0$
for every fixed real $a,a'$. The reason is
$u\ge K/1260$ and $2\mathfrak p_0\ge2$ in the negative exponents.

We first record the measure-theoretic passage to the limit.

\begin{lemma}[A summable sequence of probability measures]\label{IV-62-61}
Let $(\Omega,\mathfrak B)$ be a measurable space, let $K_*\in\mathbb N$, and let $\mu^K$, $K\ge K_*$, be probability measures on $\mathfrak B$. Let $q_K\ge0$ be numbers with $\sum_{K\ge K_*}q_K<\infty$ and
\[
\sup_{E\in\mathfrak B}\bigl|\mu^{K+1}(E)-\mu^K(E)\bigr|\le q_K\qquad(K\ge K_*).
\]
Then $\mu(E):=\lim_{K\to\infty}\mu^K(E)$ exists for every $E\in\mathfrak B$,
\[
\sup_{E\in\mathfrak B}\bigl|\mu^K(E)-\mu(E)\bigr|\le\sum_{k\ge K}q_k\qquad(K\ge K_*),
\]
and $\mu$ is a probability measure on $\mathfrak B$.
\end{lemma}

\begin{proof}
\emph{Step 1: the limit and the bound.} For $K'>K\ge K_*$ and $E\in\mathfrak B$,
\[
\bigl|\mu^{K'}(E)-\mu^K(E)\bigr|\le\sum_{k=K}^{K'-1}q_k\le\sum_{k\ge K}q_k ,
\]
and the right-hand side tends to $0$ as $K\to\infty$. Hence $(\mu^K(E))_K$ is a Cauchy sequence, and its limit $\mu(E)$ exists. Letting $K'\to\infty$ in the display gives the asserted bound, uniformly in $E$.

\emph{Step 2: elementary properties.} As limits of the corresponding properties of the $\mu^K$, we have $\mu(E)\ge0$ for all $E\in\mathfrak B$, $\mu(\emptyset)=0$, $\mu(\Omega)=1$, and $\mu$ is finitely additive.

\emph{Step 3: $\sigma$-additivity.} Let $E_1,E_2,\dots\in\mathfrak B$ be pairwise disjoint and put $R_N:=\bigcup_{i>N}E_i$. By finite additivity,
\[
\mu\Bigl(\bigcup_iE_i\Bigr)=\sum_{i\le N}\mu(E_i)+\mu(R_N).
\]
Let $\eta>0$, and choose $K\ge K_*$ with $\sum_{k\ge K}q_k\le\eta$. By Step~1, $\mu(R_N)\le\mu^K(R_N)+\eta$ for every $N$. Since $\mu^K$ is a measure and $R_N$ decreases to $\emptyset$, we have $\mu^K(R_N)\to0$ as $N\to\infty$; hence $\limsup_N\mu(R_N)\le\eta$. As $\eta>0$ is arbitrary and $\mu(R_N)\ge0$, $\mu(R_N)\to0$, and $\mu(\bigcup_iE_i)=\sum_i\mu(E_i)$.
\end{proof}

\begin{definition}[The laws of the block fields]\label{IV-62-62}
Assume (E), with label $\mathsf b$ and side $\ell=2L^m$.
\begin{enumerate}
\item[(1)] At cutoff $K$ set $\beta=L^K\mathsf b$ and
\[
 \rho_0(U)=e^{-A^{a_K}(\mathsf b,U)},\qquad
 Z_K=\int\rho_0\,dU>0,\qquad
 \langle H\rangle_K=Z_K^{-1}\int H\rho_0\,dU.
\]
Here $dU$ is product Haar probability measure and
$A^{a_K}(\mathsf b,U)=\beta\sum_pa_p(U)$ is the Wilson action.
\item[(2)] For $K\ge\max\{d,1\}$ the law at depth $d$ is
\begin{equation}\label{IV-62-63}
 \mu_d^K(E)=\langle\mathbf1_E(\mathcal M^{K-d}U)\rangle_K
 \qquad(E\subset\mathcal V(\Lambda_{d,\ell})\text{ Borel}).
\end{equation}
The averaging map is Borel and gauge covariant. Its map from
$\Lambda_{d+1,\ell}$ to $\Lambda_{d,\ell}$ is independent of the
bare cutoff: its definition uses precisely these two lattices.
These properties are proved in Proposition~\YMMref{4.5}{20}.
Thus $\mu_d^K$ is a Borel probability measure.
\item[(3)] Write $\mathfrak B_d$ for the sigma-algebra of invariant
Borel sets in $\mathcal V(\Lambda_{d,\ell})$. Gauge transformations
are bijections, so complements and countable unions of invariant
sets remain invariant. The map $\mathcal M$ in the consistency
identity below is the fixed map from depth $d+1$ to depth $d$.
\item[(4)] For $K\ge2$, use the errors above and set
\begin{equation}\label{IV-62-64}
 \mathfrak o_K=8\mathsf R_{t_1}+12\epsilon_{\rm fg}
                          +8\epsilon_{\rm cl}+12\epsilon_{\rm sh}.
\end{equation}
\item[(5)] For an interval $\mathsf J=[\mathsf b_-,\mathsf b_+]$
as above, put
\begin{equation}\label{IV-62-65}
\mathsf C_{\mathrm{bf}}:=2\Bigl[1+12\,C_{\mathrm{old}}C_{\mathrm{fg}}\tfrac L{L-1}E_0\,L^4+32\,\mathsf K_{\mathrm w}c_\rho^2A_0^2\bigl(1+\tfrac{2\varrho}{\mathsf b_-}\bigr)\bigl(\log(L\mathsf b_+)\bigr)^{2\mathfrak p_0}L^{7/2}\Bigr]
\end{equation}
and
\[
K_{\mathrm{bf}}:=\min\Bigl\{k\in\mathbb Z:\ k\ge K_0,\ \text{and }\ 8\,\bar{\mathsf R}+12\,\bar\epsilon_{\mathrm{sh}}\le\tfrac12\,\mathsf C_{\mathrm{bf}}\,\ell^3K^{\mathsf p_0}L^{-K/10}\ \text{ at all integers }K\ge k\Bigr\},
\]
where $\mathsf p_0=2\mathfrak p_0+1/2$ and $K_0$ is the uniform
threshold specified above. This minimum exists: the bracket defining
$\mathsf C_{\rm bf}$ is at least $1$, and the barred errors decay
faster than every exponential in $K$. Its dependencies are only
$\ell$, the endpoints of $\mathsf J$ and the fixed constants.
\end{enumerate}
\end{definition}

The comparison of the two lattices made above also controls the laws of Definition~\ref{IV-62-62}, and Lemma~\ref{IV-62-61} turns this control into their convergence as the cutoff tends to infinity.

\begin{corollary}[Convergence of the laws of the block fields]\label{IV-62-66}
Assume the setting (E), and take $u,s,t_1,t_0$ and the comparison errors defined above.
\begin{enumerate}
\item[\textup{(a)}] Let $K\ge2$ be an integer at which the conditions (k0), (k1), (k2), (k3), (k4) and (k6) hold. Then, for every integer $0\le d\le u$ and every bounded Borel function $F:\mathcal V(\Lambda_{d,\ell})\to\mathbb C$ invariant under the gauge transformations of $\Lambda_{d,\ell}$,
\begin{equation}\label{IV-62-67}
\Bigl|\bigl\langle F(\mathcal M^{K+1-d}U)\bigr\rangle_{K+1}-\bigl\langle F(\mathcal M^{K-d}U)\bigr\rangle_K\Bigr|\ \le\ \mathfrak o_K\,\sup|F| .
\end{equation}
If moreover $\mathsf J$ is an interval of labels under (Tor) on the side $\ell$ with $\mathsf b\in\mathsf J$, then at every integer $K\ge K_{\mathrm{bf}}$ the conditions (k0)--(k4) and (k6) hold, and
\begin{equation}\label{IV-62-68}
\mathfrak o_K\ \le\ \mathsf C_{\mathrm{bf}}\,\ell^3\,K^{\mathsf p_0}\,L^{-K/10};
\end{equation}
here $\mathsf C_{\mathrm{bf}}$ and $K_{\mathrm{bf}}$ are the numbers of Definition~\ref{IV-62-62}(5), which depend on the label $\mathsf b$ only through $\mathsf J$.
\item[\textup{(b)}] Let $K_{\mathrm{bf}}$ be taken at $\mathsf J=\{\mathsf b\}$. For every integer $d\ge0$ and every $E\in\mathfrak B_d$ the limit $\mu_d(E):=\lim_{K\to\infty}\mu^K_d(E)$ exists, and
\begin{equation}\label{IV-62-69}
\sup_{E\in\mathfrak B_d}\bigl|\mu^K_d(E)-\mu_d(E)\bigr|\ \le\ \sum_{K'\ge K}\mathfrak o_{K'}<\infty\qquad\text{for every integer }K\ge K_{\mathrm{bf}}\text{ with }\lceil K/1260\rceil\ge d .
\end{equation}
The set function $\mu_d$ is a probability measure on $\mathfrak B_d$. The family $(\mu_d)_{d\ge0}$ is consistent: if $E\in\mathfrak B_d$, then $\mathcal M^{-1}(E)\in\mathfrak B_{d+1}$ and
\begin{equation}\label{IV-62-70}
\mu_{d+1}\bigl(\mathcal M^{-1}E\bigr)=\mu_d(E),
\end{equation}
where $\mathcal M:\mathcal V(\Lambda_{d+1,\ell})\to\mathcal V(\Lambda_{d,\ell})$ is the averaging map of Definition~\ref{IV-62-62}(3).
\item[\textup{(c)}] For every integer $d\ge0$, every integer $K\ge\max\{d,1\}$ and every Borel set $E\subset\mathcal V(\Lambda_{d,\ell})$,
\begin{equation}\label{IV-62-71}
\mu^K_d(E)=\int\bar{\mathbf 1}_E\,d\mu^K_d,\qquad\bar{\mathbf 1}_E(V):=\int\mathbf 1_E(V^v)\,dv,
\end{equation}
where $dv$ is the product of one Haar probability measure of $G$ for each site of $\Lambda_{d,\ell}$, a measure on the group of the gauge transformations $v$ of $\Lambda_{d,\ell}$.
The function $\bar{\mathbf 1}_E$ is Borel, takes its values in $[0,1]$ and is gauge invariant.
Consequently the limit in (b), the bound \eqref{IV-62-69} and the identity \eqref{IV-62-70} hold for all Borel sets: $\mu^K_d\to\mu_d$ in total variation, that is, uniformly on the Borel subsets of $\mathcal V(\Lambda_{d,\ell})$; $\mu_d$ is a Borel probability measure on $\mathcal V(\Lambda_{d,\ell})$ that is gauge invariant; and the image of $\mu_{d+1}$ under $\mathcal M$ is $\mu_d$.
\end{enumerate}
\end{corollary}

\begin{proof}
We first compare the normalized means at depth $u$. On either lattice
$\bullet\in\{\mathrm A,\mathrm B\}$, let $\Lambda_0^\bullet$
be the full unnormalized law at that depth, $\Lambda_1^\bullet$ its
small-field chain restricted to $\mathsf{In}_{[u]}$, and
$\Lambda_2^\bullet$ the chain restarted from the classical action at
depth $s$, again restricted to $\mathsf{In}_{[u]}$. We regard them as
functionals on bounded gauge-invariant Borel functions $F$ of the
depth-$u$ field. These are precisely the chain and classical-start
measures used in Propositions~\YMMref{115.2}{806} and~\YMMref{116.1}{808}, with source
zero. All three functionals are positive and have strictly positive
mass. For the full law this is immediate. The chain is the image of
the nonnegative small-field weight times the bare density, by
Proposition~\YMMref{115.2}{806}. For the classical start, the initial density
is positive, the chart kernels preserve positivity, and the final
exponential density does not vanish on the open set
$\mathsf{In}_{[u]}$ containing the trivial configuration; these are
the positivity statements in the proof of Proposition~\YMMref{116.1}{808}.

We use the following elementary normalization bound. If positive
functionals $P,Q$ satisfy $P(1)>0$ and
\[
 |Q(F)-P(F)|\le\epsilon\|F\|_\infty P(1),
 \qquad0\le\epsilon\le1/2,
\]
then $Q(1)\ge(1-\epsilon)P(1)>0$, and subtraction of the two normalized
fractions gives
\[
 \left|\frac{Q(F)}{Q(1)}-\frac{P(F)}{P(1)}\right|
 \le\frac{2\epsilon}{1-\epsilon}\|F\|_\infty
 \le4\epsilon\|F\|_\infty.
\]

The discarded-mass estimate in Proposition~\YMMref{115.2}{806} gives
$0\le\Lambda_0^\bullet-\Lambda_1^\bullet$ on this algebra and
\[
 (\Lambda_0^\bullet-\Lambda_1^\bullet)(1)
            \le\mathsf R_{t_1}\Lambda_0^\bullet(1).
\]
Consequently their normalized means differ by at most
$4\mathsf R_{t_1}\|F\|_\infty$.

Forgetting the early array changes the final exponent by at most
$\epsilon_{\rm fg}$, by equation~\YMMref{2180}{803} and the definition of
$\epsilon_{\rm fg}$ above. Write the two final densities as
$p_1=e^{c^\bullet+\Phi^\bullet}$ and
$p_2=e^{\hat c^\bullet+\hat\Phi^\bullet}$ on
$\mathsf{In}_{[u]}$. Positivity gives $p_1=|p_1|$ almost everywhere
and $p_2=|p_2|>0$ there. Thus, with
$\varkappa=|e^{c^\bullet-\hat c^\bullet}|$,
\[
 \frac{p_1}{p_2}\in
 \varkappa[e^{-\epsilon_{\rm fg}},e^{\epsilon_{\rm fg}}]
 \quad\hbox{almost everywhere}.
\]
Since $\epsilon_{\rm fg}\le1/8$,
$e^{\epsilon_{\rm fg}}-1\le(5/4)\epsilon_{\rm fg}\le1/2$.
Apply the normalization bound to
$P=\varkappa\Lambda_2^\bullet$, $Q=\Lambda_1^\bullet$.
The factor $\varkappa$ cancels, and these two normalized means differ
by at most $5\epsilon_{\rm fg}\|F\|_\infty$.

The two classical starts can be written on the common depth-$s$
configuration space as $k^\bullet\hat F_0^\bullet\,dW$, followed
by the same block averaging down to depth $u$. The final
total-variation estimate in the proof of Proposition~\YMMref{116.1}{808} is
\[
 \int|k^{\mathrm A}\hat F_0^{\mathrm A}
             -k^{\mathrm B}\hat F_0^{\mathrm B}|\,dW
 \le(2\epsilon_{\rm cl}+3\epsilon_{\rm sh})
                           \Lambda_2^{\mathrm A}(1).
\]
We will also use the relative version of the mismatch bound. For every
bounded nonnegative gauge-invariant function $H$ of the final block field,
the proof of Proposition~\YMMref{116.1}{808} gives
\[
 \int H(\mathcal M^{s-u}W)\,k^\bullet
           \mathbf1_{\{k^\circ=0\}}\hat F_0^\bullet\,dW
       \le\epsilon_{\rm sh}\Lambda_2^\bullet(H).
\]
Indeed, at a mismatching test some fluctuation of lattice $\bullet$
exceeds half its threshold. Proposition~\YMMref{114.4}{803} bounds the
marked next measure by
$\exp(C_{\rm sh}-\kappa_{\rm b}\mathfrak b(x_{(k)})^2/4)$ times the
ordinary next measure. Each later positive cut and image operation
preserves this order. Sum over at most $s-u$ steps and $3N_{[s]}$
bonds at each step, and use
$\mathfrak b(x_{(k)})^2\ge A_1^2y_{[u]}$ to obtain the displayed bound.
This proves an inequality of measures before testing against $H$.

The total-variation proof bounds separately the mass where a fluctuation
test differs and the difference of the classical action on the common cut. Inserting
$F$ in this inequality costs only $\|F\|_\infty$.
By (k3), $2\epsilon_{\rm cl}+3\epsilon_{\rm sh}\le1/4$,
so normalization gives an error
$(8\epsilon_{\rm cl}+12\epsilon_{\rm sh})\|F\|_\infty$.
Adding the five comparisons
\[
 \Lambda_0^{\mathrm A},\ \Lambda_1^{\mathrm A},\
 \Lambda_2^{\mathrm A},\ \Lambda_2^{\mathrm B},\
 \Lambda_1^{\mathrm B},\ \Lambda_0^{\mathrm B}
\]
bounds the difference of the normalized full laws by
\[
 [8\mathsf R_{t_1}+10\epsilon_{\rm fg}
                   +8\epsilon_{\rm cl}+12\epsilon_{\rm sh}]
               \|F\|_\infty\le\mathfrak o_K\|F\|_\infty.
\]
For $d\le u$, apply this to $F\circ\mathcal M^{u-d}$.
Block averaging is Borel and gauge covariant, and its maps at equal
physical depth coincide. Hence this composition remains bounded and
gauge invariant, with no increase of its supremum norm. This proves
\eqref{IV-62-67}.

For the rate, the ceiling inequalities in Theorem~\YMMref{AR}{810} give
$4u-s\le-K/10+4$ and $(7s-K)/2\le-K/9+7/2$.
Therefore
\[
 12\epsilon_{\rm fg}\le
 12C_{\rm old}C_{\rm fg}\frac L{L-1}E_0L^4\ell^3L^{-K/10},
\]
and, since $\log(L^K\mathsf b_+)\le K\log(L\mathsf b_+)$,
\[
 8\epsilon_{\rm cl}\le
 32\mathsf K_{\rm w}c_\rho^2A_0^2
 (1+2\varrho/\mathsf b_-)(\log(L\mathsf b_+))^{2\mathfrak p_0}
 L^{7/2}\ell^3K^{2\mathfrak p_0+1/2}L^{-K/9}.
\]
Their sum is at most half the right side of \eqref{IV-62-68};
the other half bounds $8\mathsf R_{t_1}+12\epsilon_{\rm sh}$
by the choice of $K_{\rm bf}$. This proves the uniform rate.

For (b), fix $d$ and let $K_*$ be the least $K\ge K_{\rm bf}$
with $\lceil K/1260\rceil\ge d$. For every invariant Borel set $E$,
apply (a) to $F=\mathbf1_E$. The resulting differences of the
probability measures are bounded by the summable sequence
$\mathfrak o_K$, since a fixed power of $K$ times $L^{-K/10}$ is
summable. Lemma~\ref{IV-62-61} gives the limit, the rate and countable
additivity. Gauge covariance makes $\mathcal M^{-1}E$ invariant,
and the exact finite-cutoff identity
\[
 \mu^K_{d+1}(\mathcal M^{-1}E)=\mu^K_d(E)
\]
passes to the limit. This proves consistency.

For (c), invariance of the Wilson density and Haar measure makes the
bare law gauge invariant. Gauge covariance of block averaging and
extension of a coarse site gauge to the fine sites give the same
invariance for $\mu^K_d$. The action $(V,v)\mapsto V^v$ is continuous,
so Fubini's theorem shows that $\bar{\mathbf1}_E$ is Borel.
Right invariance of Haar measure gives
\[
 \bar{\mathbf1}_E(V^{v'})
 =\int\mathbf1_E(V^{vv'})\,dv=\bar{\mathbf1}_E(V).
\]
Integrating the invariance of $\mu^K_d$ over $v$ proves
\eqref{IV-62-71}. Now apply (a) to this function, whose values lie
in $[0,1]$. The same summable bound holds for every Borel set $E$,
so Lemma~\ref{IV-62-61} applies on the full Borel sigma-algebra.
The finite-cutoff consistency and gauge invariance pass to these
limits. For invariant $E$, $\bar{\mathbf1}_E=\mathbf1_E$, so this
extension agrees with the measure in (b).
\end{proof}

\subsection{Local H\"older gauges}

The compactness argument needs a gauge with a uniformly H\"older
continuous first derivative. The regional gauge theorem controls the
potential and its first differences. Its construction also controls
the componentwise Laplacian; after a cutoff, the scalar inverse then
gives the required H\"older estimate. We give these two steps, including
the uniformity for every exponent below one.

We use the ordinary positive Laplacian
$\Delta=\sum_\nu\nabla_\nu^*\nabla_\nu$. Thus the ordinary
curl--divergence $\partial^*\partial$ on one-forms and the
componentwise operator $\Delta$ are different operators. A bound
``on a cube'' for a one-form or its derivatives means a bound at
every positively oriented bond with an endpoint in that cube.

\begin{proposition}[Local regular gauges on a cube]\label{I-56-4}
There are constants $B^{\mathrm p}_1\ge1$ and
$c^{\mathrm p}_1\in(0,1]$, depending only on $L$, with the following
property; fix the choices made in the proof and put
$\mathfrak r_{\mathrm p}=L$.
Let $1\le k\le K$, and measure lengths at level $k$.
Let $\mathsf l$ be a positive integer multiple of $\mathfrak r_{\mathrm p}$;
let $\square$ be a compact cube of side $\mathsf l$ that is a union
of cubes of $\pi[\mathfrak r_{\mathrm p}a_{K-k}]$; and let
$\tilde\square$ and $\square^\natural$ be the concentric cubes of
sides $\mathsf l+4\mathfrak r_{\mathrm p}$ and
$\mathsf l+4L\mathfrak r_{\mathrm p}$. Assume the last side is less
than half the torus period. Let $\alpha>0$ satisfy
$21L^2\mathsf l\alpha\le c^{\mathrm p}_1$. Suppose
$U\in\mathcal V_0$ satisfies
$|\partial U(p)-\mathbf1|<\alpha\eta^2$ at every bare plaquette
with a corner in $\square^\natural$, and
$|(D_U^{\mathsf T}\partial U)(b)|<\alpha\eta^2$ at every positively
oriented bare bond with an endpoint there. Then there are a real gauge
$u$ on the bare sites of $\tilde\square$ and a real one-form $A_0$
on its inner bonds such that $U^{u^{-1}}=e^{i\eta A_0}$ and
\begin{equation}\label{I-56-5}
 |A_0|,\quad|\nabla^\eta A_0|,\quad
 |\partial^{\eta*}\partial^\eta A_0|,\quad|\Delta^\eta A_0|
 <21L^2B^{\mathrm p}_1\mathsf l\alpha
 \qquad\text{on }\square.
\end{equation}
All derivatives in this display are taken at the trivial configuration.
\end{proposition}

\begin{proof}
We use the construction in the proof of \cite[Theorem~F]{YMM1},
keeping one estimate furnished by that construction in addition to
the three estimates in its statement. In the box of side
$s=\mathsf l+4$, put $\epsilon=s\alpha$. The normalized covariant
derivative $\omega$ of the smoothed comb section satisfies, away
from the boundary,
\[
 |\omega|\le20\epsilon,\qquad
 |\nabla^\eta\omega|\le C_M\epsilon,\qquad
 |\Delta^\eta\omega|\le55\epsilon+\sigma|\nabla^\eta\omega|,
 \quad \sigma=120\epsilon.
\]
These are respectively the normalization bound, the absorption
estimate for ordinary differences, and the equation for $\omega$
in that proof. Its smallness choice gives $\sigma\le1$.
The constructed potential is
$A_0=\eta^{-1}\Theta(\mathbf1+\eta\omega)$, where $\Theta$ is
the smooth quaternion logarithm on $|q-\mathbf1|\le1/2$.
Take $c_\Theta$ to bound its first two derivatives there.
Taylor's formula in each of the three positive and three negative
coordinate directions gives
\[
 |\Delta^\eta A_0|
 \le c_\Theta(55+C_M)\epsilon
       +6c_\Theta\eta C_M^2\epsilon^2
 \le c_\Theta(55+11C_M)\epsilon.
\]
For the last inequality use $C_M=2560c_I$ and
$\epsilon\le(120c_I)^{-1}$, as in the absorption step, and
$\eta\le1/13$. Thus the same gauge has all four estimates.
Increase the absolute box constant $B_\square$ to include this
last coefficient. With the constants of that proof put
$B_*=5B_\square$ and $c_*=c_\square/5$. Then all four estimates
are bounded by $B_*\mathsf l\alpha$ on $\square$, provided
$\mathsf l\alpha\le c_*$.

To specify the constants in \eqref{I-56-5}, set
\[
 B^{\mathrm p}_1=\max\{1,LB_*/21\},\qquad
 c^{\mathrm p}_1=\min\{1,21c_*/L\}.
\]
The assumed smallness implies
$\mathsf l\alpha\le c^{\mathrm p}_1/(21L^2)\le c_*/L^3$.
The unit-partition cubes refine the $L$-partition cubes because
$L$ is odd. Hence the box construction applies to $\square$
and its concentric enlargement of side $\mathsf l+4$, both
inside $\tilde\square$.

The gauge is required on the larger cube $\tilde\square$,
although its estimates are required only on $\square$.
Extend the constructed gauge by $\mathbf1$ at the remaining
sites of $\tilde\square$. Keep $A_0$ on the bonds of the smaller
enlargement; on each other positive bond choose an
$\mathrm{SU}(2)$ logarithm of the transformed bond variable and
extend by reversal. Such a logarithm exists by the elementary
formula $g=\cos t\,\mathbf1+i\sin t\,n$ with $n^2=\mathbf1$.
Every bond entering the four quantities in \eqref{I-56-5} at a
bond touching $\square$ lies within distance $2\eta<2$ of
$\square$. These quantities are therefore unchanged by the
extension. Finally
$B_*\mathsf l\alpha\le21L^2B^{\mathrm p}_1\mathsf l\alpha$.
\end{proof}

We next treat the scalar block inverse. In level-$k$ units put
$\eta=L^{-k}$. The blocks $B^k(y)$ have volume one and
$L^{3k}$ bare sites. Define
\[
 (\mathsf Q'f)(y)=\eta^3\sum_{x\in B^k(y)}f(x),\qquad
 (\mathsf Q'^*v)(x)=v(y)\quad(x\in B^k(y)),
 \qquad
 \mathsf G'=(\Delta+\varkappa_k\mathsf Q'^*\mathsf Q')^{-1}.
\]
Here $\varkappa_1=1$ and
$\varkappa_{k+1}=\varkappa_k/(L^{-2}+\varkappa_k)$; induction
gives $1-L^{-2}<\varkappa_k\le1$.
The adjoint uses the pairing $\eta^3\sum_x$ on bare sites and
the counting pairing on the unit-spaced sites. At the trivial
configuration all transports in the scalar average are the
identity, so these are precisely the construction's scalar
average and inverse.

\begin{proposition}[H\"older bound for the gradient of $\mathsf G'$ at $\mathbf1$]\label{I-56-57}
There is a number $\delta_{\mathrm H}\in(0,1]$, depending only on
$L$, with the following property. For every
$\delta_{\mathrm p}\in(0,\delta_{\mathrm H}]$ and every
$\beta_0\in[0,1)$ there is a number $B_{\mathrm p}(\beta_0)\ge1$,
depending only on $L$ and $\beta_0$, for which the assertions below hold.
Let $1\le k\le K$, use level-$k$ lengths, and allow any number
of sites of $\mathbb T^{(k)}$ per direction.
Let $y,y'\in\mathbb T^{(k)}$. Put
$\mathsf e_k=\frac12(1-\eta)(1,1,1)$, and let $\zeta$ be the
restriction of a smooth function on the torus supported in
$\{x:|x-(y-\mathsf e_k)|_\infty<1\}$, with $|\zeta|\le1$ and
first derivatives bounded by $c_\zeta$.
Let $h$ be a real bare site function supported in $B^k(y')$.
Then $\mathsf G'$ exists and, for $0\le\beta\le\beta_0$, every
direction $\nu$, and all bare sites with $0<|x-x'|_1\le1$,
\begin{equation}\label{I-56-58}
 \bigl|(\zeta\nabla_\nu\mathsf G'h)(x')
             -(\zeta\nabla_\nu\mathsf G'h)(x)\bigr|
 \le B_{\mathrm p}(\beta_0)(1+c_\zeta)
 e^{-\delta_{\mathrm p}|y-y'|_\infty}
 \sup|h|\,|x'-x|_1^\beta.
\end{equation}
We fix the constants $\delta_{\mathrm H}$ and
$B_{\mathrm p}(\beta_0)=2B_{\mathrm H}(\beta_0)$ chosen in the proof.
\end{proposition}

\begin{proof}
The argument is the scalar weighted-energy and local Poisson
argument of \cite[Propositions~18.1 and~21.3]{YMM1}.
We give the gradient H\"older step, the choice uniform in
$\beta\le\beta_0$, and the passage to small tori.

\emph{Coercivity and exponential weight.}
Set $\gamma=1-L^{-2}$ and
$\delta_0=\gamma/[24(L^2+1)]$.
On each unit block the path-product Poincar\'e inequality gives
\[
 \|\nabla v\|_{B^k(y)}^2\ge4\|v-\mathsf Q'v(y)\|_{B^k(y)}^2.
\]
Adding the block-average term, and then the nonnegative bonds
between blocks, shows that
$\mathsf A:=\Delta+\varkappa_k\mathsf Q'^*\mathsf Q'\ge\gamma$;
also $\langle v,\mathsf Av\rangle\ge\|\nabla v\|^2$.
This proves invertibility on every finite torus.

Let $u=\mathsf G'h$, $\varphi(x)=|x-y'|_1$, and
$v=e^{\delta_0\varphi}u$. The conjugated quadratic-form
calculation in the proof of \cite[Proposition~21.3]{YMM1}
uses $|\varphi(x)-\varphi(x')|\le|x-x'|_1$ and
$|\varphi-\varphi(y)|\le3/2$ on each block. The loss in the
Laplacian form is at most $6L^2\delta_0^2\|v\|^2$,
and that in the block-average form is at most
$(e^{3\delta_0}-1)\|v\|^2$. Since their sum is at most
$\gamma/4$, pairing the equation with $e^{\delta_0\varphi}v$
gives
\[
 \tfrac12\|\nabla v\|^2+\tfrac14\gamma\|v\|^2
 \le\|v\|\,\|e^{\delta_0\varphi}h\|,
 \qquad
 \|e^{\delta_0\varphi}u\|\le C(L)\sup|h|.
\]
The last inequality uses the unit volume and bounded diameter
of the input block. If a fixed-size cube $Q$ is within a
fixed distance of $y$, it follows that
\begin{equation}\label{eq:local-gauge-scalar-energy}
 \|u\|_Q\le C(L)e^{-\delta_0|y-y'|_1}\sup|h|.
\end{equation}
The constant may depend on the fixed size of $Q$, but not on
the spacing or the torus period.

\emph{Local gradient differences on a sufficiently large torus.}
Choose a power $M_\sharp=L^{\mu_\sharp}$ so large that the
cubes of radius $L+10$ used below have a single torus chart
when the period is at least $2M_\sharp$. This choice depends
only on $L$. Let $o=y-\mathsf e_k$, and let $Q^*$ be the
cube of radius $L+8$ about $o$. On the cube of radius $8$
about $o$ the equation is
\[
 \Delta u=g,\qquad g=h-\varkappa_k\mathsf Q'^*\mathsf Q'u.
\]
Every block meeting this smaller cube lies in $Q^*$.
Its volume is one, so Cauchy--Schwarz bounds its average
by $\|u\|_{Q^*}$. Thus
\[
 |g|\le |h|_{Q^*}+\|u\|_{Q^*}.
\]
If $h$ meets $Q^*$, the distance $|y-y'|_1$ is bounded by
a constant depending only on $L$; if it does not, its local
supremum is zero. Consequently \eqref{eq:local-gauge-scalar-energy}
bounds both terms by
$C(L)e^{-\delta_0|y-y'|_1}\sup|h|$.

Apply \cite[Proposition~18.1(b)]{YMM1} with radius $8$ and
$F=0$. Its hypothesis $8\ge16\eta$ holds since
$\eta\le1/13$. It gives, for points in the cube of radius
$2$ about $o$, a bound for $D=\nabla_\nu u$ and, for each
$0<\alpha<1$, a bound for its difference:
\[
 |D(x)|\le C(L)e^{-\delta_0|y-y'|_1}\sup|h|,
 \quad
 |D(x')-D(x)|\le C(L,\alpha)e^{-\delta_0|y-y'|_1}
                    \sup|h|\,|x'-x|_1^\alpha.
\]
Choose $\alpha=\beta_0$ when $\beta_0>0$. Since
$0<|x-x'|_1\le1$, this difference bound holds with
$|x-x'|_1^\beta$ for every $0\le\beta\le\beta_0$.
For $\beta_0=0$ use twice the supremum bound.
If $\zeta(x)=\zeta(x')=0$ there is nothing to prove.
Otherwise interchange the points so that $\zeta(x)\ne0$;
then both points lie within distance $2$ of $o$ in the
supremum norm. The product identity
\[
 (\zeta D)(x')-(\zeta D)(x)
 =\zeta(x)\bigl(D(x')-D(x)\bigr)
       +\bigl(\zeta(x')-\zeta(x)\bigr)D(x')
\]
and
$|\zeta(x')-\zeta(x)|\le c_\zeta|x'-x|_1
\le c_\zeta|x'-x|_1^\beta$
give \eqref{I-56-58} on these tori, with a constant
$B^\circ_{\mathrm H}(\beta_0)$ and rate
$\delta_{\mathrm H}:=\delta_0/2$.
All bounds hold componentwise for $\mathfrak g$-valued $h$:
test the difference against a unit vector of $\mathfrak g$
and use $|\langle E,h\rangle|\le|h|$.

\emph{Small tori.}
Write the period in level-$k$ units as $S=2L^{m+K-k}$ and put
$j=\max\{0,\mu_\sharp-(m+K-k)\}$. On the cover with period
$\widehat S=L^jS$, pullback commutes with $\nabla$, $\Delta$
and $\mathsf Q'^*\mathsf Q'$, because the lifted blocks are
the disjoint translates of the original blocks. Invertibility
therefore gives
\[
 (\mathsf G'h)\circ p=\widehat{\mathsf G}'(h\circ p),
 \qquad h\circ p=\sum_{p(\hat y')=y'}h_{\hat y'}.
\]
Each $h_{\hat y'}$ is supported in a single lifted block
and has supremum at most $\sup|h|$. There are $L^{3j}$ terms.
Choose a lift $\hat o$ of $o$ and retain just the copy of
$\zeta$ supported in $|\hat x-\hat o|_\infty<1$.
This cutoff is smooth: the original cutoff vanishes in a
neighbourhood of the boundary of that cube.
Its supremum and first-derivative bound are unchanged.

For the tested pair with $\zeta(x)\ne0$, lift $x$ within
distance one of $\hat o$, and lift $x'-x$ by a shortest vector.
Then the lifted pair has the same distance, and its second
point is within distance two of $\hat o$. Since $S\ge2L>4$,
the retained cutoff agrees at both points with the original
one. Apply the large-torus estimate to each $h_{\hat y'}$.
Projection decreases torus distance, so
$|\hat y-\hat y'|_1\ge|y-y'|_1$.
Summing costs at most $L^{3j}\le M_\sharp^3$.
Thus the bound holds on every torus with
\[
 B_{\mathrm H}(\beta_0)
   =\max\{1,M_\sharp^3B^\circ_{\mathrm H}(\beta_0)\},
 \qquad B_{\mathrm p}(\beta_0)=2B_{\mathrm H}(\beta_0).
\]
For a complex cutoff apply the real-cutoff argument to its
real and imaginary parts; the last factor two covers their
sum. Finally
$\delta_{\mathrm p}|y-y'|_\infty
\le\delta_{\mathrm H}|y-y'|_1$,
which gives the stated rate.
\end{proof}

We localise with cut-off functions of cubes and use a discrete Leibniz rule.

\begin{lemma}[Cut-off functions; a discrete Leibniz rule]\label{I-56-59}
Let $\varphi_\chi:\mathbb R\to[0,1]$ be the function of class $C^2$, equal to $1$ on $(-\infty,1]$ and to $0$ on $[2,\infty)$, fixed once and for all; let $s'_\chi:=\sup|\varphi_\chi'|$, $s''_\chi:=\sup|\varphi_\chi''|$, and $c_{\chi,\mathrm{triv}}=\max\{1,s'_\chi,(s'_\chi)^2,s''_\chi\}$.
Let $0\le k\le K$, use level-$k$ lengths, $\eta=L^{-k}$, and let $\mathsf K=\prod_\mu I_\mu$ be a product of three compact intervals of one length, the side of $\mathsf K$; the length $0$ is admitted.
Let $\xi\ge\eta$, assume that the side of $\mathsf K$ and $\xi$ are less than $\frac18$ of the side of the torus, and let $\chi:=\chi_{\mathsf K,\xi}$,
\[
\chi(x)=\prod_\mu\varphi_\chi\bigl(\operatorname{dist}(x_\mu,I_\mu)/\xi\bigr);
\]
The same formula applies when the side is zero. Let $\nabla=\nabla^{\mathbf 1}$, $\partial=D_{\mathbf 1}$, and $\Delta=\partial^*\partial$ on site functions, on the bare lattice.

\textup{(a)} $s'_\chi$ and $s''_\chi$ are finite; $0\le\chi\le1$; $\chi=1$ at sup distance $\le\xi$ from $\mathsf K$, and $\chi=0$ at sup distance $\ge2\xi$ from $\mathsf K$.

\textup{(b)} For all bare sites $x$ and all directions $\mu$, $\nu$,
\[
|(\nabla_\mu\chi)(x)|\le c_{\chi,\mathrm{triv}}\,\xi^{-1},\qquad|(\nabla_\mu\nabla_\nu\chi)(x)|\le c_{\chi,\mathrm{triv}}\,\xi^{-2}.
\]

\textup{(c)} Let $A$ be a matrix-valued bare $1$-form, $(\chi A)(b):=\chi(b_-)A(b)$, and let $\lambda$ be a matrix-valued bare site function. Then, for all bare sites $x$ and all directions $\mu$, $\nu$,
\[
\begin{gathered}
|(\nabla_\mu(\chi A)_\nu)(x)|\le|(\nabla_\mu A_\nu)(x)|+c_{\chi,\mathrm{triv}}\,\xi^{-1}|A_\nu(x+\eta e_\mu)|,\qquad\\
|(\nabla_\mu(\chi\lambda))(x)|\le|(\nabla_\mu\lambda)(x)|+c_{\chi,\mathrm{triv}}\,\xi^{-1}|\lambda(x+\eta e_\mu)|,\\
|(\Delta(\chi\lambda))(x)|\le|(\Delta\lambda)(x)|+6c_{\chi,\mathrm{triv}}\,\xi^{-1}\max_b|(\partial\lambda)(b)|\\
+3c_{\chi,\mathrm{triv}}\,\xi^{-2}|\lambda(x)|,
\end{gathered}
\]
where the maximum is taken over the six bare bonds with end point $x$.
\end{lemma}

\begin{proof}
(a), (b).
The derivatives $\varphi'_\chi$ and $\varphi''_\chi$ are continuous and vanish outside $[1,2]$; hence their suprema are finite.
Let $I=[c_-,c_+]$ with $c_-\le c_+$, and $g(t):=\varphi_\chi(\operatorname{dist}(t,I)/\xi)$ for real $t$.
Then $g=1$ on $\{t:\operatorname{dist}(t,I)<\xi\}$, $g(t)=\varphi_\chi((t-c_+)/\xi)$ for $t>c_+$, and $g(t)=\varphi_\chi((c_--t)/\xi)$ for $t<c_-$.
These three open sets cover the line, also when $c_-=c_+$; hence $g\in C^2$, $0\le g\le1$, $|g'|\le s'_\chi\xi^{-1}$ and $|g''|\le s''_\chi\xi^{-2}$.
On a circle of length $\ell'$, with $|I|,\xi<\frac18\ell'$, the point at distance $\frac12\ell'$ from the midpoint of $I$, which is the point of the circle farthest from $I$, lies at distance $\frac12(\ell'-|I|)>\frac7{16}\ell'>2\xi$ from $I$, so that $g=0$ near it; at all other points the three formulas above hold in a chart.

The function $\chi$ is the product $g_1g_2g_3$, where $g_\mu$ denotes the function $g$ of the interval $I_\mu$, evaluated at the coordinate $x_\mu$; hence $0\le\chi\le1$.
The sup distance of $x$ from $\mathsf K$ is $\max_\mu\operatorname{dist}(x_\mu,I_\mu)$: if it is at most $\xi$, all three factors equal $1$; if it is at least $2\xi$, one factor vanishes.
A first difference quotient of $\chi$ in a direction $\mu$ is the first difference quotient of $g_\mu$, multiplied by the two other factors, which take values in $[0,1]$; by the mean value theorem it is at most $s'_\chi\xi^{-1}$ in modulus.
A second difference quotient in two different directions $\mu\ne\nu$ is the product of the first difference quotients of $g_\mu$ and of $g_\nu$ with the third factor, which takes values in $[0,1]$; hence it is at most $(s'_\chi)^2\xi^{-2}$ in modulus.
A second difference quotient in one direction $\mu$ is the second difference quotient of $g_\mu$, multiplied by the two other factors, with values in $[0,1]$; the mean value theorem, applied twice, bounds it by $s''_\chi\xi^{-2}$.
Since $c_{\chi,\mathrm{triv}}\ge\max\{s'_\chi,(s'_\chi)^2,s''_\chi\}$, this proves (b).

(c) For functions $f$, $h$ of the bare sites, $f$ scalar and real-valued,
\[
(\nabla_\mu(fh))(x)=\eta^{-1}\bigl[f(x+\eta e_\mu)h(x+\eta e_\mu)-f(x)h(x)\bigr]=f(x)(\nabla_\mu h)(x)+(\nabla_\mu f)(x)\,h(x+\eta e_\mu).
\]
With $f=\chi$ and $h=A_\nu$, resp.\ $h=\lambda$, and with (a), (b), this gives the first two inequalities.

By the definition of the positive scalar Laplacian,
\[
(\Delta\lambda)(x)=\sum_\mu\eta^{-1}\bigl[(\nabla_\mu\lambda)(x-\eta e_\mu)-(\nabla_\mu\lambda)(x)\bigr]=\eta^{-2}\sum_\mu\bigl[2\lambda(x)-\lambda(x+\eta e_\mu)-\lambda(x-\eta e_\mu)\bigr].
\]
Fix $\mu$, and write $c:=\chi(x)$, $c_\pm:=\chi(x\pm\eta e_\mu)$, $l:=\lambda(x)$, $l_\pm:=\lambda(x\pm\eta e_\mu)$. Then
\[
2cl-c_+l_+-c_-l_-=c\,(2l-l_+-l_-)+(c-c_+)(l_+-l)+(c-c_-)(l_--l)+(2c-c_+-c_-)\,l .
\]
We divide by $\eta^2$ and sum over $\mu$.
The first terms give $\chi(x)(\Delta\lambda)(x)$, of modulus at most $|(\Delta\lambda)(x)|$.
In the second and the third terms, $\eta^{-1}|c-c_\pm|\le c_{\chi,\mathrm{triv}}\xi^{-1}$ by (b), and $\eta^{-1}(l_\pm-l)$ are the values of $\pm\partial\lambda$ at the bonds $\langle x,x+\eta e_\mu\rangle$ and $\langle x-\eta e_\mu,x\rangle$; there are six such terms.
In the fourth terms, $\eta^{-2}(2c-c_+-c_-)=-(\nabla_\mu\nabla_\mu\chi)(x-\eta e_\mu)$, of modulus at most $c_{\chi,\mathrm{triv}}\xi^{-2}$ by (b); there are three such terms.
This gives the third inequality.
\end{proof}

Finally we compute the operator $\mathsf Q'^*\mathsf Q'$ at $\mathbf 1$.

\begin{lemma}[The operator $\mathsf Q'^*\mathsf Q'$ at $\mathbf 1$]\label{I-56-60}
Let $1\le k\le K$, use level-$k$ lengths, and let $\mathsf Q'$, $\mathsf Q'^*$ and $\varkappa_k$ be as defined before Proposition~\ref{I-56-57}; the adjoint is taken for $(\cdot,\cdot)_\eta$ on real bare site functions and for $\sum_y\langle\mu(y),\mu'(y)\rangle$ on real site functions on $\mathbb T^{(k)}$, whose cells have volume $1$.
For a real bare site function $\lambda$, a bare site $x$ and the site $y\in\mathbb T^{(k)}$ with $x\in B^k(y)$,
\[
(\mathsf Q'^*\mathsf Q'\lambda)(x)=L^{-3k}\sum_{x''\in B^k(y)}\lambda(x'') .
\]
Moreover $\sup|\varkappa_k\mathsf Q'^*\mathsf Q'\lambda|\le\sup|\lambda|$, and $(\mathsf Q'^*\mathsf Q'\lambda)(x)=0$ unless $\lambda(x'')\ne0$ at some bare site $x''$ with $|x''-x|_\infty<1$.
\end{lemma}

\begin{proof}
The adjoint formula follows by grouping the fine-site pairing by blocks: $\eta^3\sum_x\langle\lambda(x),\mu(y(x))\rangle=\sum_y\langle\mathsf Q'\lambda(y),\mu(y)\rangle$. Hence $(\mathsf Q'^*\mu)(x)=\mu(y)$ on $B^k(y)$, and the asserted formula is the block mean.
A mean has modulus at most $\sup|\lambda|$, and $0<\varkappa_k\le1$; this gives the bound.
The sites of $B^k(y)$ lie within sup distance $<\frac12$ of $y$, by their centered cube coordinates; hence any two of them are at sup distance $<1$, which gives the last assertion.
\end{proof}

\begin{definition}[The numbers of the local gauges]\label{I-56-61}
Use the constants $B^{\mathrm p}_1,c^{\mathrm p}_1$ of Proposition~\ref{I-56-4}, $c_{\chi,\mathrm{triv}}$ of Lemma~\ref{I-56-59}, and $B_{\mathrm p}(\beta_0)$ of Proposition~\ref{I-56-57}.
We define $\phi_0\colon\mathbb R\to\mathbb R$ by $\phi_0(s):=e^{-1/s}$ for $s>0$ and $\phi_0(s):=0$ for $s\le0$, and we set
\[
\begin{multlined}
\psi_0(t):=\phi_0\bigl(\tfrac78-|t|\bigr)\,\Bigl[\phi_0\bigl(\tfrac78-|t|\bigr)+\phi_0\bigl(|t|-\tfrac34\bigr)\Bigr]^{-1}\quad(t\in\mathbb R),\qquad
\zeta_0(z):=\psi_0(z_1)\,\psi_0(z_2)\,\psi_0(z_3) \\
(z\in\mathbb R^3).
\end{multlined}
\]
Lemma~\ref{I-56-64} below shows that $\psi_0$ is well defined and collects the properties of $\psi_0$ and $\zeta_0$ on which we rely. With $\partial_\mu$ the partial derivative in the direction $\mu$ and the supremum taken over $\mathbb R^3$, we put
\begin{equation}\label{I-56-62}
\begin{gathered}
c_\zeta:=\max_\mu\sup|\partial_\mu\zeta_0|,\qquad a_{\mathrm g}:=\min\Bigl\{\frac15,\ \frac{c^{\mathrm p}_1}{126L^2}\Bigr\},\qquad\mathsf c_{\mathrm g}:=\max\{160,\ 126L^2B^{\mathrm p}_1\},\\
\mathsf c_{\mathrm g}(\beta_0):=8^3B_{\mathrm p}(\beta_0)(1+c_\zeta)(2+9c_{\chi,\mathrm{triv}})\,\mathsf c_{\mathrm g}\qquad(\beta_0\in[0,1)).
\end{gathered}
\end{equation}
The number $c_\zeta$ is finite by Lemma~\ref{I-56-64}.
\end{definition}

The next two lemmas record the elementary properties of these numbers and functions.

\begin{lemma}\label{I-56-63}
The numbers $a_{\mathrm g}$ and $\mathsf c_{\mathrm g}$ depend only on $L$. Moreover $a_{\mathrm g}\in(0,\frac15]$ and $\mathsf c_{\mathrm g}\ge160$.
\end{lemma}

\begin{proof}
Proposition~\ref{I-56-4} fixes $B^{\mathrm p}_1,c^{\mathrm p}_1$ as functions of $L$ alone.
Besides these two numbers, the formulas \eqref{I-56-62} for $a_{\mathrm g}$ and $\mathsf c_{\mathrm g}$ contain only $L$ and explicit numbers, so $a_{\mathrm g}$ and $\mathsf c_{\mathrm g}$ depend only on $L$.
Since $c^{\mathrm p}_1>0$ by Proposition~\ref{I-56-4}, both entries of the minimum that defines $a_{\mathrm g}$ are positive, and $a_{\mathrm g}>0$.
Finally, $\frac15$ is an entry of that minimum and $160$ an entry of the maximum that defines $\mathsf c_{\mathrm g}$; hence $a_{\mathrm g}\le\frac15$ and $\mathsf c_{\mathrm g}\ge160$.
\end{proof}

\begin{lemma}\label{I-56-64}
The function $\psi_0$ is defined on all of $\mathbb R$ and of class $C^\infty$; it satisfies $0\le\psi_0\le1$, $\psi_0(t)=1$ for $|t|\le\frac34$ and $\psi_0(t)=0$ for $|t|\ge\frac78$.
The function $\zeta_0$ belongs to $C^\infty(\mathbb R^3)$ and satisfies $0\le\zeta_0\le1$, $\zeta_0(z)=1$ for $|z|_\infty\le\frac34$ and $\zeta_0(z)=0$ for $|z|_\infty\ge\frac78$.
The numbers $c_{\chi,\mathrm{triv}}\ge1$ and $c_\zeta\ge0$ are finite absolute constants.
For every $\beta_0\in[0,1)$ the number $\mathsf c_{\mathrm g}(\beta_0)$ depends only on $L$ and $\beta_0$, and $\mathsf c_{\mathrm g}(\beta_0)\ge8^3\mathsf c_{\mathrm g}\ge\mathsf c_{\mathrm g}$.
\end{lemma}

\begin{proof}
\emph{Step 1: $\phi_0$ is smooth.} Define polynomials $P_q$ by $P_0:=1$ and $P_{q+1}(u):=u^2\bigl(P_q(u)-P_q'(u)\bigr)$.
On $(0,\infty)$ we have $\phi_0^{(q)}(s)=P_q(1/s)\,e^{-1/s}$ for every $q\ge0$: this holds for $q=0$, and the derivative of $P_q(1/s)e^{-1/s}$ equals $s^{-2}\bigl(P_q(1/s)-P_q'(1/s)\bigr)e^{-1/s}=P_{q+1}(1/s)e^{-1/s}$.
For $u>0$ and every integer $l\ge0$, the power series of the exponential gives $e^u\ge u^{l+1}/(l+1)!$, that is, $u^le^{-u}\le(l+1)!/u$.
Since $s^{-1}P_q(1/s)$ is a polynomial in $u=1/s$, and since each term $u^le^{-u}$ is at most $(l+1)!/u$, which tends to $0$ as $u\to\infty$, it follows that $s^{-1}P_q(1/s)e^{-1/s}\to0$ as $s\downarrow0$.
We now show by induction on $q$ that $\phi_0^{(q)}$ exists on all of $\mathbb R$, vanishes on $(-\infty,0]$ and equals $P_q(1/s)e^{-1/s}$ on $(0,\infty)$.
For $q=0$ this is the definition.
If it holds for $q$, then $\phi_0^{(q)}$ vanishes identically on $(-\infty,0)$, so that its derivative vanishes there, and on $(0,\infty)$ its derivative is $P_{q+1}(1/s)e^{-1/s}$.
At $s=0$ the difference quotient $\phi_0^{(q)}(s)/s$ is $0$ for $s<0$ and equals $s^{-1}P_q(1/s)e^{-1/s}$ for $s>0$, which tends to $0$ as $s\downarrow0$; hence $\phi_0^{(q+1)}(0)$ exists and equals $0$.
Thus $\phi_0$ has derivatives of all orders on $\mathbb R$, and $\phi_0\in C^\infty(\mathbb R)$.

\emph{Step 2: $\psi_0$.} We have $\phi_0\ge0$ everywhere and $\phi_0(s)>0$ for $s>0$.
The two arguments $\frac78-|t|$ and $|t|-\frac34$ add up to $\frac18>0$, so at least one of them is positive, and the bracket in the definition of $\psi_0$ is positive; thus $\psi_0$ is defined on $\mathbb R$.
The numerator is one of the two non-negative terms of the bracket, whence $0\le\psi_0\le1$.
For $|t|\le\frac34$ the second term of the bracket vanishes and $\psi_0(t)=1$; for $|t|\ge\frac78$ the numerator vanishes and $\psi_0(t)=0$.
On the open set $\{t\ne0\}$ the function $t\mapsto|t|$ is smooth, so $\psi_0$ is there a quotient of smooth functions with positive denominator, hence smooth; on the open set $\{|t|<\frac34\}$ the function $\psi_0$ is constant.
These two open sets cover $\mathbb R$, and therefore $\psi_0\in C^\infty(\mathbb R)$.

\emph{Step 3: $\zeta_0$.} As a product of three smooth functions with values in $[0,1]$, the function $\zeta_0$ is smooth with values in $[0,1]$. If $|z|_\infty\le\frac34$, each factor equals $1$; if $|z|_\infty\ge\frac78$, some coordinate satisfies $|z_\mu|\ge\frac78$ and the corresponding factor vanishes. In particular $\zeta_0$ is supported in $\{|z|_\infty\le\frac78\}$.

\emph{Step 4: the constants.} The number $c_{\chi,\mathrm{triv}}=\max\{1,s'_\chi,(s'_\chi)^2,s''_\chi\}$, with $s'_\chi=\sup|\varphi_\chi'|$ and $s''_\chi=\sup|\varphi_\chi''|$, is computed from the single function $\varphi_\chi$ fixed once and for all; it is finite by Lemma~\ref{I-56-59}(a), it is at least $1$ because $1$ is an entry of the maximum, and it is an absolute constant.
Next, $\partial_\mu\zeta_0(z)=\psi_0'(z_\mu)\prod_{\nu\ne\mu}\psi_0(z_\nu)$; since $0\le\psi_0\le1$ and $\psi_0(0)=1$, we get $\sup|\partial_\mu\zeta_0|=\sup|\psi_0'|$ for each $\mu$, so that $c_\zeta=\sup|\psi_0'|\ge0$.
The derivative $\psi_0'$ is continuous and vanishes outside the compact set $\{\frac34\le|t|\le\frac78\}$, so $c_\zeta$ is finite; it is determined by the explicit formula for $\psi_0$ and is therefore absolute.
By Lemma~\ref{I-56-63}, $\mathsf c_{\mathrm g}$ depends only on $L$, and by Proposition~\ref{I-56-57}, $B_{\mathrm p}(\beta_0)\ge1$ depends only on $L$ and $\beta_0$.
Consequently $\mathsf c_{\mathrm g}(\beta_0)$ depends only on $L$ and $\beta_0$.
Each of the factors $B_{\mathrm p}(\beta_0)$, $1+c_\zeta$ and $2+9c_{\chi,\mathrm{triv}}$ is at least $1$, which gives $\mathsf c_{\mathrm g}(\beta_0)\ge8^3\mathsf c_{\mathrm g}\ge\mathsf c_{\mathrm g}$.
\end{proof}

We now show that a configuration whose local radii are small on a neighbourhood of a cube admits, near this cube, a gauge in which it is the exponential of a small real $1$-form with a H\"older continuous gradient.

The second assertion of the local-gauge theorem concerns a critical
configuration. We make explicit the local inequality in its admission
hypothesis before using it. The subsequent profile argument applies to
every admitted triple; it does not replace an arbitrary cut set by a
reduced determining set.

\begin{definition}[The regional data and the local radii]\label{def:local-gauge-data}
Data \textup{(D)} consist of a depth $0\le n\le K$, powers
$\mathfrak M$ and $R_1\ge L$ of $L$, $\rho=R_1\mathfrak M$,
and a decreasing sequence $\Omega_0,\ldots,\Omega_n$.
Each $\Omega_j$ is a union of $\mathfrak M$-cubes at level $j$,
and in bare units
\[
 \operatorname{dist}_\infty(\Omega_j^c,\Omega_{j+1})>\rho L^j.
\]
The level-$n$ torus has at least $6\rho$ sites per direction.
Empty regions and repeated regions are allowed; put
$\Omega_{n+1}=\varnothing$. The index of a lattice site is the
largest $j$ for which it lies in $\Omega_j$. The determining
set $\mathbf B$ consists, at level $j$, of the bonds with both
endpoints in $\Omega_0$ and at least one endpoint in
$\Omega_j\setminus\Omega_{j+1}$; it includes inward bonds.
In level-$n$ units put $\xi_j=L^{j-n}$. For a connected finite
union $\Gamma$ of segments define its weighted length by
\[
 |\Gamma|_{\mathbf B}=
 \sum_{j=1}^n\xi_j^{-1}\operatorname{length}
       (\Gamma\cap(\Omega_j\setminus\Omega_{j+1}))
       +\xi_0^{-1}\operatorname{length}(\Gamma\setminus\Omega_1).
\]
The distance $d_{\mathbf B}(x,x')$ is the infimum of these
lengths over such sets containing $x,x'$. It is unchanged by
rescaling all lengths.

An element is a bare plaquette in $\Omega_0$, or a positively
oriented bare bond at which all four incident plaquettes lie
in $\Omega_0$, so that the current $J(U)$ is defined.
For such an element let $q$ be the largest index of its corners
and put $\mathsf L_{\mathbf B}=\xi_q$. Its base site is the
base of the plaquette or the initial site of the bond. With
$\eta=L^{-n}$ define
\[
 \mathsf r_U(p)=\eta^{-2}\xi_q^2|\partial U(p)-\mathbf1|,
 \qquad
 \mathsf r_U(b)=\xi_q^3|J(U)(b)|.
\]
Thus $\mathsf r_U(p)=L^{2q}|\partial U(p)-\mathbf1|$ and,
if the current is written in level-$k$ units,
$|J(U)(b)|=L^{3(k-q)}\mathsf r_U(b)$.
Regularity at radius $a$ means $\mathsf r_U(\mathfrak e)<a$
for all elements.

A cut set $\mathfrak F$ is a reversal-invariant set of inward
bonds, and $\mathbf B_{\mathfrak F}=\mathbf B\setminus\mathfrak F$.
A real datum assigns group elements, with reversal by inversion,
to these bonds. The constraint space $C(\mathbf B_{\mathfrak F},V)$
fixes the iterated averages at these values; criticality means
stationarity of the Wilson action under variations preserving them.
Its tests include the tested plaquettes and the pinched pairs.
At level $j$, a plaquette is tested if at least one of its bonds
is retained and each bond is either retained or, when $j\ge1$,
has both endpoints in $\Omega_{j-1}\setminus\Omega_j$.
Use the value of $V_j$ on a retained bond and the average of
$V_{j-1}$ on a bond of the second kind. The latter is defined
only when all bonds needed for that average are retained.
The tested holonomy is the oriented product of these four values.
For completeness, at a retained inward bond $b$ ending at a seam node
$x$, let $y_b$ be the parent of $b_-$ and set
\[
 \mathsf f_V(b)=V_{j+1}(\langle[x],y_b\rangle)
                       V_j(y_b,b_-)V_j(b).
\]
Here $V_j(y_b,b_-)$ is the averaged contour transport inside
the parent block. The frame is defined only when the coarse bond
$\langle[x],y_b\rangle$ and all bonds internal to that parent
block are retained, and the averaged contour variable is defined.
Two distinct retained inward bonds $b,\tilde b$
ending at $x$ form a pinched pair when
$b_-+\tilde b_--x\in\Omega_{j+1}$; its holonomy is
$\mathsf f_V(b)\mathsf f_V(\tilde b)^{-1}$.
The sites of a test are its plaquette corners or its seam node.
Distances to tests mean the minimum over those sites.
\end{definition}

\begin{definition}[Admission for the local inequality]\label{def:local-gauge-admission}
Fix the localization powers $\mathfrak M_{\mathrm c}$ and
$R_{\mathrm c}$, and put
$\rho_{\mathrm c}=R_{\mathrm c}\mathfrak M_{\mathrm c}\ge L^3$.
A triple $(\sigma,K_{\mathrm{loc}},a_{\mathrm c})$ is admitted
if $\sigma\ge L^2\rho_{\mathrm c}$ is a power of $L$,
$K_{\mathrm{loc}}\ge4$, $0<a_{\mathrm c}\le1$, and the following
local inequality holds for all the data in its scope:
data \textup{(D)} with $\mathfrak M\ge\mathfrak M_{\mathrm c}$
and $\rho\ge L^2\sigma$, every cut set $\mathfrak F$, every real
datum $V$ on $\mathbf B_{\mathfrak F}$, and every critical
configuration $U$ regular at radius $a_{\mathrm c}$.
All its test holonomies and frames must be defined. Put
\[
 R_{\mathrm{loc}}=15L\sigma,\qquad
 \mathsf r(x)=\max_{x(\mathfrak e)=x}\mathsf r_U(\mathfrak e),
\]
with empty maximum zero, and
\[
 \mathsf r_R(x)=\max_{d_{\mathbf B}(x,z)\le R_{\mathrm{loc}}}
                       \mathsf r(z),\qquad
 \mathsf v_R(x)=\max_{d_{\mathbf B}(x,\tau)\le R_{\mathrm{loc}}}
                       |\partial V(\tau)-\mathbf1|.
\]
The required inequality is
\begin{equation}\label{eq:local-gauge-admitted-inequality}
 \mathsf r(x)\le K_{\mathrm{loc}}\mathsf v_R(x)
                          +\tfrac12\mathsf r_R(x).
\end{equation}
The five associated numbers are $\mathfrak M_{\mathrm c}$,
$\sigma$, $a_{\mathrm c}$ and
\[
 \delta^{\mathrm c}_0=\frac{\log2}{30L\sigma},\qquad
 \mathsf c^{\mathrm c}_{\mathrm{loc}}=10K_{\mathrm{loc}}.
\]
Admission is always with respect to the torus and all the data
quantified above. On reduced data, \cite[Lemma~53.9]{YMM1}
proves this local inequality. For arbitrary cut sets it is part
of the admission hypothesis of the theorem below.
\end{definition}

\begin{lemma}[From the local inequality to a local radius]\label{lem:local-gauge-profile}
Let the triple be admitted, and let $V,U$ be in its scope.
For every element $\mathfrak e$ and each of its corners $x'$
one has
\begin{equation}\label{eq:local-gauge-profile}
 \mathsf r_U(\mathfrak e)\le
 \mathsf c^{\mathrm c}_{\mathrm{loc}}P(x'),\qquad
 P(x')=\max_\tau e^{-\delta^{\mathrm c}_0d_{\mathbf B}(x',\tau)}
                         |\partial V(\tau)-\mathbf1|.
\end{equation}
If $W\subset\Omega_0$, $0<\varepsilon'\le\varepsilon_1$ and
$D\ge0$, and the test holonomies have size less than
$\varepsilon_1$ everywhere and less than $\varepsilon'$
at distance at most $D$ from $W$, then at every element with
a corner in $W$,
\[
 \mathsf r_U(\mathfrak e)<
 \mathsf c^{\mathrm c}_{\mathrm{loc}}
       (\varepsilon'+\varepsilon_1e^{-\delta^{\mathrm c}_0D}).
\]
\end{lemma}

\begin{proof}
The triangle inequality gives
$P(z)\le e^{\delta^{\mathrm c}_0d_{\mathbf B}(x,z)}P(x)$.
Since $e^{\delta^{\mathrm c}_0R_{\mathrm{loc}}}=\sqrt2$,
\[
 \mathsf v_R(x)\le\sqrt2P(x),\qquad
 P(z)\le\sqrt2P(x)
 \quad\text{if }d_{\mathbf B}(x,z)\le R_{\mathrm{loc}}.
\]
For $t>0$ maximize $\mathsf r(x)/(P(x)+t)$ over the finite
set of bare sites and call the maximum $S$. At a maximizing
site, \eqref{eq:local-gauge-admitted-inequality} gives
\[
 S(P(x)+t)\le\sqrt2K_{\mathrm{loc}}P(x)
                         +\frac{S}{\sqrt2}(P(x)+t).
\]
Hence $S\le2(\sqrt2+1)K_{\mathrm{loc}}\le5K_{\mathrm{loc}}$.
Letting $t$ decrease to zero proves
$\mathsf r(x)\le5K_{\mathrm{loc}}P(x)$.
The base and every corner of an element have weighted distance
at most two: join them by at most two bare edges, on which the
weight is at most the inverse bare spacing.
Thus moving to a corner costs at most
$e^{2\delta^{\mathrm c}_0}\le2$, proving
\eqref{eq:local-gauge-profile}. If there are no bare sites,
there are no elements and the assertion is empty.

Now let $x'\in W$ be a corner of $\mathfrak e$.
A test at distance at most $D$ from $W$ contributes less than
$\varepsilon'$ to $P(x')$. Every other test has distance
greater than $D$ from $x'$, and contributes less than
$\varepsilon_1e^{-\delta^{\mathrm c}_0D}$.
There are finitely many tests. Their maximum is therefore
strictly smaller than the larger of these two positive
numbers, which is at most their sum. If there are no tests,
$P=0$ and \eqref{eq:local-gauge-profile} gives
$\mathsf r_U=0$. This proves the last assertion also in that case.
\end{proof}

\begin{theorem}[Local gauges on cubes]\label{I-56-65}
There exist numbers $\mathsf c_{\mathrm g}\ge1$ and $a_{\mathrm g}\in(0,1]$ depending only on $L$ and, for every $\beta_0\in[0,1)$, a number $\mathsf c_{\mathrm g}(\beta_0)\ge1$ depending only on $L$ and $\beta_0$, such that the two statements below hold for every $\beta_0\in[0,1)$.
The numbers $\mathsf c_{\mathrm g}$, $a_{\mathrm g}$ and $\mathsf c_{\mathrm g}(\beta_0)$ of Definition~\ref{I-56-61} have this property, and they are the numbers we use whenever the theorem is applied.

\emph{First statement.} Let data \textup{(D)} of Definition~\ref{def:local-gauge-data} be given with $\rho\ge L^7$, and let $0\le j\le n$.
We use level-$n$ lengths, so that $\xi_j=L^{j-n}$ and $\eta=\eta_n$.
Let $m'$ be a positive integer multiple of $\rho$, let $\square$ be a compact cube of side $2m'\xi_j$, and let $\tilde\square$ be the compact cube of side $(2m'+4\rho)\xi_j$ concentric with $\square$; we assume that this side is smaller than half the side of $T_\ell$ and that $\tilde\square\subset\Omega_j$.
Let $U$ be any configuration on the inner bare bonds of $\Omega_0$.
Let $\varepsilon_\square>0$ be such that $\mathsf r_U(\mathfrak e)<\varepsilon_\square$ for every element $\mathfrak e$ with a corner in $\tilde\square$, and such that $m'\varepsilon_\square\le a_{\mathrm g}$.
Then there exist a cube $\square'\supset\{x:\operatorname{dist}_\infty(x,\square)\le4\xi_j\}$, a map $u$ from the bare sites of $\square'$ to $G$ and a real $1$-form $A$ on the bare bonds inside $\square'$ such that $U^{u^{-1}}=e^{i\eta A}$ on these bonds and such that, on $\square$,
\begin{equation}\label{I-56-66}
\begin{gathered}
\xi_j|A|,\quad\xi_j^2|\nabla^\eta A|,\quad\xi_j^3|\partial^{\eta*}\partial^\eta A|,\quad\xi_j^3|\Delta^\eta A|<\mathsf c_{\mathrm g}\,m'\varepsilon_\square,\\
\bigl|(\nabla^\eta_\nu A_\mu)(x')-(\nabla^\eta_\nu A_\mu)(x)\bigr|\le\mathsf c_{\mathrm g}(\beta_0)\,m'\varepsilon_\square\,\xi_j^{-2-\beta}\,|x'-x|_1^{\beta}\qquad(0\le\beta\le\beta_0).
\end{gathered}
\end{equation}
The second line holds for all bare sites $x\ne x'$ of $\square$ and all directions $\mu,\nu$.
All derivatives are taken at the configuration $\mathbf 1$, the operator $\Delta^\eta=\sum_\nu\nabla^{\eta\mathsf T}_\nu\nabla^\eta_\nu$ acts on $A$ componentwise, and a bound on $\square$ holds at every positively oriented bare bond $b$ with $\{b_-,b_+\}\cap\square\ne\emptyset$.

\emph{Second statement.} Let $(\sigma,K_{\mathrm{loc}},a_{\mathrm c})$ be an admitted triple, with its five numbers $\mathfrak M_{\mathrm c},\sigma,\delta^{\mathrm c}_0,\mathsf c^{\mathrm c}_{\mathrm{loc}},a_{\mathrm c}$ (Definition~\ref{def:local-gauge-admission}). Let the data (D) with $\rho\ge L^2\sigma$, the index $j$, the number $m'$ and the cubes $\square$, $\tilde\square$ be as in the first statement, and assume moreover $\mathfrak M\ge\mathfrak M_{\mathrm c}$.
Let $\mathfrak F$, $V$ and $U$ be in the scope of the admitted triple: $\mathfrak F$ is a reversal-invariant set of inward bonds, $V$ a real datum on $\mathbf B_{\mathfrak F}$, and $U$ a critical configuration of $C(\mathbf B_{\mathfrak F},V)$ regular at the radius $a_{\mathrm c}$.
Let $\varepsilon'$, $\varepsilon_1$ and $D$ satisfy $0<\varepsilon'\le\varepsilon_1$, $D\ge0$, $|\partial V(\tau)-\mathbf 1|<\varepsilon_1$ at every test $\tau$, and $|\partial V(\tau)-\mathbf 1|<\varepsilon'$ at every test $\tau$ with $d_{\mathbf B}(\tau,\tilde\square)\le D$.
Set
\begin{equation}\label{I-56-67}
\varepsilon_\square:=\mathsf c^{\mathrm c}_{\mathrm{loc}}\,\varepsilon^\sharp,\qquad\varepsilon^\sharp:=\varepsilon'+\varepsilon_1e^{-\delta^{\mathrm c}_0D}.
\end{equation}
If $\mathsf c^{\mathrm c}_{\mathrm{loc}}\,m'\varepsilon^\sharp\le a_{\mathrm g}$, then there exist $\square'$, $u$ and $A$ as in the first statement, with $U^{u^{-1}}=e^{i\eta A}$ on the bare bonds inside $\square'$, such that the bounds \eqref{I-56-66} hold with this $\varepsilon_\square$.
\end{theorem}

\begin{proof}
By Lemmas~\ref{I-56-63} and~\ref{I-56-64}, the numbers $a_{\mathrm g}\in(0,\frac15]\subset(0,1]$ and $\mathsf c_{\mathrm g}\ge160\ge1$ depend only on $L$, and $\mathsf c_{\mathrm g}(\beta_0)\ge\mathsf c_{\mathrm g}\ge1$ depends only on $L$ and $\beta_0$; these are the ranges and dependences claimed.
Since $a_{\mathrm g}\le\frac15$ and $m'\ge1$, the condition $m'\varepsilon_\square\le a_{\mathrm g}$ implies $\varepsilon_\square\le\frac15$.
The numerical bound $\rho\ge L^7$ is part of the hypotheses of the first statement.

\emph{Step 1: scaling.} A $1$-form expressed in level-$j$ lengths is $\xi_j$ times the same $1$-form expressed in level-$n$ lengths, each difference quotient contributes one further factor $\xi_j$, and distances in level-$j$ lengths are the distances in level-$n$ lengths divided by $\xi_j$.
Consequently the quantities on the left of the first line of \eqref{I-56-66}, and the quotient of the left-hand side of the second line by $\xi_j^{-2-\beta}|x'-x|_1^\beta$, are scale-free, and it suffices to prove \eqref{I-56-66} in level-$j$ lengths, in which $\xi_j=1$.
From now on we use level-$j$ lengths: $\eta=L^{-j}$, the cube $\square$ has side $2m'$, and $\tilde\square$ has side $2m'+4\rho$.
Let $\mathfrak e$ be an element with a corner in $\tilde\square$.
Since $\tilde\square\subset\Omega_j$, this corner has index at least $j$, so the largest index $j'$ of a corner of $\mathfrak e$ satisfies $j'\ge j$.
The scale-free formulas in Definition~\ref{def:local-gauge-data}, written in level-$j$ lengths, give $|\partial U(p)-\mathbf 1|=\mathsf r_U(p)L^{-2j'}\le\mathsf r_U(p)\eta^2$ if $\mathfrak e=p$ is a plaquette, and $|J(U)(b)|=\mathsf r_U(b)L^{3(j-j')}\le\mathsf r_U(b)$ if $\mathfrak e=b$ is a bond.
Hence
\begin{equation}\label{I-56-68}
|\partial U(p)-\mathbf 1|<\varepsilon_\square\eta^2,\qquad|J(U)(b)|<\varepsilon_\square
\end{equation}
at every element $p$, resp.\ $b$, with a corner in $\tilde\square$. We extend $U$ by $\mathbf 1$ to a configuration in $\mathcal V_0$.

\emph{Step 2: the case $j=0$.} Here $\eta=1$.
Let $\square':=\{x:\operatorname{dist}_\infty(x,\square)\le4\}$ and let $\mathsf Y$ be the set of bare sites of $\square'$.
In a chart of $\tilde\square$, the set $\square'$ is a product of three compact intervals of length $2m'+8$, and $\mathsf Y$ is the set of integer points in it; $\mathsf Y$ is consequently a box of bare sites with at most $s+1$ sites in each direction, where $s:=2m'+8\le3m'$ (recall $m'\ge\rho\ge L^7$).
Since $4\le2\rho$, we have $\square'\subset\tilde\square$.
Every plaquette inside $\mathsf Y$ is an element with a corner in $\tilde\square$; let $\varpi$ be the largest of the numbers $|\partial U(p)-\mathbf 1|$ over these plaquettes.
By \eqref{I-56-68}, $\varpi<\varepsilon_\square$, and thus $2s\varpi<6m'\varepsilon_\square\le\frac65<2$.
The comb-gauge estimate \cite[Proposition~14.4]{YMM1}, applied to the box $\mathsf Y$ and to $U$ on the bonds with both end points in $\mathsf Y$, provides $h\colon\mathsf Y\to G$ with $U^h=e^{ia^0}$ for a real bond function $a^0$ with
\[
\sup|a^0|\le\pi s\varpi<3\pi m'\varepsilon_\square<10m'\varepsilon_\square .
\]
We set $u:=h^{-1}$ and $A:=a^0$; then $U^{u^{-1}}=e^{i\eta A}$ on the bonds inside $\square'$, and $|A|<10m'\varepsilon_\square$.
The bounds on $\square$ concern the positively oriented bare bonds $b$ with $\{b_-,b_+\}\cap\square\ne\emptyset$.
At such a bond $b$, each of the values of $\nabla A$, $\partial^*\partial A$ and $\Delta A$ is a sum of at most $16$ values of $\pm A$ at bonds joining sites within sup distance $2$ of $\square$, all of which lie in $\mathsf Y$.
Indeed, a difference quotient involves two values; by the discrete curl--divergence formula, the value $(\partial^*\partial A)(b)$ is made of the values of $\pm\partial A$ at the four plaquettes containing $b$, each of which is a sum of four values of $\pm A$; and $(\Delta A)_\mu(x)=\sum_\nu[2A_\mu(x)-A_\mu(x+e_\nu)-A_\mu(x-e_\nu)]$ involves twelve values.
Therefore $|\nabla A|$, $|\partial^*\partial A|$ and $|\Delta A|$ are smaller than $160\,m'\varepsilon_\square$ on $\square$.
For bare sites $x\ne x'$ of $\square$ we have $|x'-x|_1\ge1$, and so
\[
|(\nabla_\nu A_\mu)(x')-(\nabla_\nu A_\mu)(x)|<2\cdot2\cdot10\,m'\varepsilon_\square\le160\,m'\varepsilon_\square\,|x'-x|_1^\beta .
\]
Thus \eqref{I-56-66} holds with $160$ in place of both constants, and hence as stated, because $\mathsf c_{\mathrm g}(\beta_0)\ge\mathsf c_{\mathrm g}\ge160$.

\emph{Step 3: the case $j\ge1$; the cube.} We choose a chart in which $\square=\prod_\mu[t_\mu,t_\mu+2m']$.
Put $\mathfrak r_{\mathrm p}=L$; thus $13\le\mathfrak r_{\mathrm p}\le\rho$.
In the level-$j$ chart, the cubes of $\pi[\mathfrak r_{\mathrm p}a_{K-j}]$ are the images of the products $\prod_\mu[g_\mu,g_\mu+\mathfrak r_{\mathrm p}]$ with $g_\mu\in\mathfrak r_{\mathrm p}(\mathbb Z+\frac12)$.
Let $g_\mu$ be the largest number in $\mathfrak r_{\mathrm p}(\mathbb Z+\frac12)$ with $g_\mu\le t_\mu-4$, and put
\[
\mathsf l:=2m'+2\mathfrak r_{\mathrm p},\qquad\square_6:=\prod_\mu[g_\mu,g_\mu+\mathsf l].
\]
Since $\rho$ and $\mathfrak r_{\mathrm p}$ are powers of $L$ with $\rho\ge\mathfrak r_{\mathrm p}$, the number $m'$ is a multiple of $\mathfrak r_{\mathrm p}$, and hence $\mathsf l$ is a positive integer multiple of $\mathfrak r_{\mathrm p}$.
Moreover $\mathsf l\le2m'+4\rho$, and $2m'+4\rho$, the side of $\tilde\square$, is smaller than half the side of the torus by hypothesis.
Subdividing each coordinate interval into intervals of length $\mathfrak r_{\mathrm p}$ shows that $\square_6$ is a union of cubes of $\pi[\mathfrak r_{\mathrm p}a_{K-j}]$.
Its side satisfies $\mathsf l\le3m'$, because $2\mathfrak r_{\mathrm p}=2L\le\rho\le m'$.
From $t_\mu-4-\mathfrak r_{\mathrm p}<g_\mu\le t_\mu-4$ and $\mathfrak r_{\mathrm p}\ge13$ we obtain $g_\mu\ge t_\mu-2\mathfrak r_{\mathrm p}$ and
\[
g_\mu+\mathsf l>t_\mu+2m'+\mathfrak r_{\mathrm p}-4\ge t_\mu+2m'+4,\qquad g_\mu+\mathsf l\le t_\mu+2m'+2\mathfrak r_{\mathrm p},
\]
and therefore
\[
\{x:\operatorname{dist}_\infty(x,\square)\le4\}\subset\square_6\subset\{x:\operatorname{dist}_\infty(x,\square)\le2\mathfrak r_{\mathrm p}\}.
\]
Let $\square_6^\natural$ be the compact cube of side $\mathsf l+4L\mathfrak r_{\mathrm p}$ concentric with $\square_6$.
Its points lie at sup distance at most $(2L+2)\mathfrak r_{\mathrm p}\le2\rho-1$ from $\square$; hence $\square_6^\natural\subset\tilde\square$, and the side of $\square_6^\natural$ is smaller than half the side of the torus.
A bare plaquette that has a corner $x\in\square_6^\natural$, or that contains a bare bond with an end point $x\in\square_6^\natural$, has all its corners within sup distance $\eta\le1$ of $x$, hence within $2\rho$ of $\square$, that is, in $\tilde\square\subset\Omega_0$.
In particular the current is defined at the bonds with an end point in $\square_6^\natural$, and these plaquettes and positively oriented bonds are elements with a corner in $\tilde\square$.
We convert the current bound to the full matrix divergence required by
Proposition~\ref{I-56-4}. For $X\in\mathrm{SU}(2)$ write
$X=c\mathbf1+iH$, where $c=\frac12\operatorname{tr}X$ and
$1-c=\frac12|X-\mathbf1|^2$. Thus
$\partial U(p)=c(p)\mathbf1+i\eta^2F_U(p)$ and
\[
 (D_U^{\mathsf T}\partial U)(b)
       =a(b)\mathbf1+i\eta^2J(U)(b).
\]
The scalar term is the sum of two transported differences
$\pm\eta^{-1}(c(p)-c(p'))$. With
$r=\varepsilon_\square\le1/5$, \eqref{I-56-68} implies
$0\le1-c(p)\le r^2\eta^4/2$, so that
$|a(b)|\le r^2\eta^3\le r^2\eta^2$.
Consequently every positively oriented bare bond $b$ with
an endpoint in $\square_6^\natural$ satisfies
\[
|(D^{\mathsf T}_U\partial U)(b)|\le(r+r^2)\eta^2\le\tfrac65\,\varepsilon_\square\eta^2 .
\]
We apply Proposition~\ref{I-56-4} at the level $j$, to the cube $\square_6$ of side $\mathsf l$ and with $\alpha:=2\varepsilon_\square$.
Its hypotheses are satisfied: $\mathsf l$ is a positive integer multiple of $\mathfrak r_{\mathrm p}$, the cube $\square_6$ is a union of cubes of $\pi[\mathfrak r_{\mathrm p}a_{K-j}]$, the concentric cube $\square_6^\natural$ of side $\mathsf l+4L\mathfrak r_{\mathrm p}$ has side smaller than half the side of the torus, the plaquette variables and the values of $D^{\mathsf T}_U\partial U$ just estimated are bounded by $\frac65\varepsilon_\square\eta^2<\alpha\eta^2$, and
\[
21L^2\mathsf l\alpha\le21L^2\cdot3m'\cdot2\varepsilon_\square=126L^2m'\varepsilon_\square\le126L^2a_{\mathrm g}\le c^{\mathrm p}_1 .
\]
The proposition yields a map $u$ on the bare sites of the cube $\square'$ of side $\mathsf l+4\mathfrak r_{\mathrm p}$ concentric with $\square_6$, which contains $\square_6$ and hence $\{x:\operatorname{dist}_\infty(x,\square)\le4\}$, and a real $1$-form $A:=A_0$ on the bare bonds inside $\square'$ with $U^{u^{-1}}=e^{i\eta A}$ there; since $21L^2B^{\mathrm p}_1\mathsf l\alpha\le126L^2B^{\mathrm p}_1m'\varepsilon_\square$, its bound \eqref{I-56-5} gives
\begin{equation}\label{I-56-69}
|A|,\ |\nabla A|,\ |\partial^*\partial A|,\ |\Delta A|<\bar{\mathsf a}:=126L^2B^{\mathrm p}_1m'\varepsilon_\square\qquad\text{on }\square_6 .
\end{equation}
As $\square\subset\square_6$ and $\bar{\mathsf a}\le\mathsf c_{\mathrm g}m'\varepsilon_\square$, the first line of \eqref{I-56-66} follows.

\emph{Step 4: the case $j\ge1$; differences of the gradient.} Let $x\ne x'$ be bare sites of $\square$, and fix $\mu$, $\nu$ and $\beta\in[0,\beta_0]$.

If $|x-x'|_\infty>\frac14$, then $4|x'-x|_1>1$, and \eqref{I-56-69} together with $4^\beta\le4$ gives
\[
|(\nabla_\nu A_\mu)(x')-(\nabla_\nu A_\mu)(x)|<2\bar{\mathsf a}\le2\bar{\mathsf a}\,(4|x'-x|_1)^\beta\le8\bar{\mathsf a}\,|x'-x|_1^\beta .
\]

Let now $0<|x-x'|_\infty\le\frac14$, so that $|x-x'|_1\le\frac34\le1$.
Let $\chi:=\chi_{\{x\},1}$ be the cut-off function of Lemma~\ref{I-56-59} for the cube $\mathsf K=\{x\}$ of side $0$ and for $\xi=1$.
The two conditions of that lemma are met: $\xi=1\ge\eta=L^{-j}$, and the side $0$ of $\mathsf K$ and the number $\xi=1$ are smaller than one eighth of the side of the torus, which is at least $6\rho>8$ by data \textup{(D)}.
By clause (a) of Lemma~\ref{I-56-59}, $\chi$ takes values in $[0,1]$, $\chi=1$ at sup distance at most $1$ from $x$, and $\chi=0$ at sup distance at least $2$ from $x$.
Let $f:=\chi A_\mu$, extended by $0$; since $A$ is real and $\chi$ is real-valued, $f$ is a real site function.
The support of $f$ and the neighbours of its points lie within sup distance $3$ of $x$, hence in $\square_6$.
Let $\Delta$ and $\varkappa_j\mathsf Q'^*\mathsf Q'$ be the operators of Proposition~\ref{I-56-57} at the level $j$, and put $h:=(\Delta+\varkappa_j\mathsf Q'^*\mathsf Q')f$.
That proposition provides the inverse $\mathsf G'$ of $\Delta+\varkappa_j\mathsf Q'^*\mathsf Q'$, so that $f=\mathsf G'h$.
The function $h$ is real: the values of $\Delta f$ are real linear combinations of values of $f$, and by Lemma~\ref{I-56-60} the values of $\mathsf Q'^*\mathsf Q'f$ are means of values of $f$.
Lemma~\ref{I-56-59}(c) with $\xi=1$, in which $\Delta$ acts on the component $A_\mu$ as on a site function, Lemma~\ref{I-56-60} and \eqref{I-56-69} give
\[
\sup|h|\le\bigl(1+6c_{\chi,\mathrm{triv}}+3c_{\chi,\mathrm{triv}}\bigr)\bar{\mathsf a}+\bar{\mathsf a}=(2+9c_{\chi,\mathrm{triv}})\,\bar{\mathsf a}.
\]
Moreover $h$ vanishes at sup distance at least $3$ from $x$: the term $\Delta f$ because $f$ vanishes at such a point and at its neighbours, the term $\varkappa_j\mathsf Q'^*\mathsf Q'f$ by the last assertion of Lemma~\ref{I-56-60}.
As the sites of a block $B^j(y')$ lie within sup distance less than $\frac12$ of $y'$, the function $h$ is supported over the sites $y'\in\mathbb T^{(j)}$ with $|y'-x|_\infty<\frac72$, and there are at most $7^3\le8^3$ of them.

Let $\mathsf e_j=\frac12(1-\eta)(1,1,1)$ be the corner vector used in Proposition~\ref{I-56-57}, and let $y\in\mathbb T^{(j)}$ satisfy $|x-(y-\mathsf e_j)|_\infty\le\frac12$; such a site exists because the points $y-\mathsf e_j$, $y\in\mathbb T^{(j)}$, form a lattice of spacing $1$.
Let $\zeta:=\zeta_0\bigl(\cdot-(y-\mathsf e_j)\bigr)$ on the torus; this is meaningful since the side of the torus exceeds $2$ and $\zeta_0$ vanishes outside $\{|z|_\infty<\frac78\}$.
By Lemma~\ref{I-56-64}, $\zeta$ is the restriction of a smooth function on $T_\ell$ supported in the open cube $\{w:|w-(y-\mathsf e_j)|_\infty<1\}$, with $|\zeta|\le1$ and first derivatives bounded in modulus by $c_\zeta$.
The sites $x$ and $x'$ lie within sup distance $\frac12+\frac14=\frac34$ of $y-\mathsf e_j$, so $\zeta=1$ at $x$ and at $x'$; and $\chi=1$ at $x$, $x'$ and at their neighbours.
Hence $\zeta\nabla_\nu\mathsf G'h=\nabla_\nu A_\mu$ at $x$ and at $x'$.
We write $h$ as the sum of its restrictions to the blocks $B^j(y')$; each of them is a real bare site function supported over $y'$ and bounded by $\sup|h|$.
Applying \eqref{I-56-58} with $k=j$ and $\delta_{\mathrm p}=\delta_{\mathrm H}$ to each of them, using $e^{-\delta_{\mathrm H}|y-y'|_\infty}\le1$ and summing over the at most $8^3$ sites $y'$, we obtain
\[
\bigl|(\nabla_\nu A_\mu)(x')-(\nabla_\nu A_\mu)(x)\bigr|\le8^3B_{\mathrm p}(\beta_0)(1+c_\zeta)(2+9c_{\chi,\mathrm{triv}})\,\bar{\mathsf a}\,|x'-x|_1^\beta .
\]

In both cases the second line of \eqref{I-56-66} follows, since $\bar{\mathsf a}\le\mathsf c_{\mathrm g}m'\varepsilon_\square$ and $8\le8^3B_{\mathrm p}(\beta_0)(1+c_\zeta)(2+9c_{\chi,\mathrm{triv}})$.

\emph{Step 5: the second statement.} The admitted-triple bounds give $\rho\ge L^2\sigma\ge L^4\rho_{\mathrm c}\ge L^7$. Lemma~\ref{lem:local-gauge-profile}, applied with $W=\tilde\square\subset\Omega_0$, shows that every element $\mathfrak e$ with a corner in $\tilde\square$ satisfies $\mathsf r_U(\mathfrak e)<\mathsf c^{\mathrm c}_{\mathrm{loc}}\varepsilon^\sharp=\varepsilon_\square$.
Since $\varepsilon_\square>0$ and $m'\varepsilon_\square=\mathsf c^{\mathrm c}_{\mathrm{loc}}m'\varepsilon^\sharp\le a_{\mathrm g}$, the first statement applies with this $\varepsilon_\square$.
\end{proof}

To apply the first statement on a small-field background at level $n$,
use data \textup{(D)} with $\rho\ge L^7$. Fix a power
$\rho_{\rm H}\ge\rho$ once and take $m'=\rho_{\rm H}$; since both
numbers are powers of $L$, $m'$ is an integer multiple of $\rho$.
A cube of side $2\rho_{\rm H}$ in level-$n$ units has the prescribed
concentric enlargement of side $2\rho_{\rm H}+4\rho\le6\rho_{\rm H}$.
On the full small-field background this enlargement lies in $\Omega_n$.
Its side is smaller than half the period when
\[
 6\rho_{\rm H}<L^{m+K-n}.
\]
If $\mathsf r_U<\varepsilon_n$ on that enlargement and
$\rho_{\rm H}\varepsilon_n\le a_{\rm g}$, the theorem, with $j=n$,
gives
\[
 |A|,\ |\nabla A|\le\mathsf c_{\rm g}\rho_{\rm H}\varepsilon_n,
 \qquad
 |\nabla_\nu A_\mu(x')-\nabla_\nu A_\mu(x)|
 \le\mathsf c_{\rm g}(\beta_0)\rho_{\rm H}\varepsilon_n
                                      |x'-x|_1^\beta.
\]
Here $U^{u^{-1}}=e^{iL^{-n}A}$ on the larger gauge cube.
If the available radius is $C_{\rm reg}\varepsilon_n$, that radius
must instead be used in both the smallness condition and these
bounds. In particular, the exponents $3/4$ and $7/8$ are both
admitted. At a level $n=K-\lceil\psi K\rceil$, with fixed $\psi>0$,
the period condition and the smallness condition hold for all
sufficiently large $K$ whenever the small-field radius has the
stated construction bound
$\varepsilon_n\le CK^pL^{-\lceil\psi K\rceil/2}$; the threshold is
uniform on each compact admissible interval of labels.

\section{Quantitative formulations of the tensor results}\label{sf-sec}

We collect the quantitative tensor statements used below, including equivalence of lattice representatives with the same symbol, continuity in the label and the radial conservation identity. The proofs below refer to the detailed estimates in the preceding sections. The assertions about overlapping tensor sources carry the additional qualification stated in Theorem~\ref{sf-coin}(b).

\subsection{Notation and hypotheses}

\emph{Notation.} All constants depend only on the basic constants (a common real choice as in Proposition~\YMMref{69.1}{595}, among them $\mathfrak p_0$, $A_0$, $M=L^{m'}$, $C_{\rm lat}$ and the last, $\gamma$) and on the arguments shown; a constant is called basic if it depends on the basic constants alone. The blocking factor $L$ is odd; $a_K=L^{-K}$, $\mathbb L_K=a_K\mathbb Z^3$; $T_\ell$ is the torus of side $\ell=2L^m$ and $\Lambda_{K,\ell}$ its lattice of spacing $a_K$; $\langle\cdot\rangle_{K,\ell,\beta}$ is Wilson's measure on $\Lambda_{K,\ell}$, $\langle\mathcal A_1;\dots;\mathcal A_n\rangle$ a cumulant, $\mathcal G_{K,\beta}$ the set of the periodic Gibbs states on $\mathbb L_K$ at $\beta$, that is, of the Gibbs states of Wilson's weight $e^{-\beta\sum_pa_p}$ that are invariant under the translations of a subgroup of finite index of $a_K\mathbb Z^3$ (Section~\ref{sec:state-limit-conventions}); by Proposition~\textup{\YMMref{103.4}{758}} every $\nu\in\mathcal G_{K,\beta}$ is invariant under $a_K\mathbb Z^3\rtimes W_3$ and reflection positive at the lattice mirrors $\{n\cdot x=c\}$, $n\in\{e_\mu,e_\mu\pm e_\lambda\}$, of the nine families.

On $\mathbb R^3$ and on $T_\ell$, $|x|$ is the Euclidean norm and $|x|_\infty=\max_\mu|x_\mu|$ the sup norm; the balls $B(x,r)$ (open) and $\bar B(x,r)$ (closed), the distances $|x-y|$ of points and the distance $\mathrm{dist}$ of sets are Euclidean. Where sets or supports are said to be pairwise $\ge\delta$ apart, to have a diameter $\le\varrho$, or where a weight vanishes within $\delta$ of a set, the sup distance $\mathrm{dist}_\infty$ is meant, unless the Euclidean distance is named, as in Theorem~\ref{sf-coin}(a), (c). In Lemmas~\ref{sf-jet} and~\ref{sf-tr}, in Theorem~\ref{sf-gibbs}(e) and in Theorem~\ref{sf-coin}(c), (d) the separations of sets are read with the Euclidean distance $\mathrm{dist}(A,B)=\inf_{x\in A,y\in B}|x-y|$, which these statements name; the factor $1/\sqrt3$ between it and the sup distance is paid inside their proofs.

The depth at the level $k$ of the renormalisation is $s=K-k$ (\ref{sec:setup-first-passage-arithmetic}); the depth of a distance $d$ is the least $s=s(d)\ge0$ with $L^sd\ge C_W(s)$, where $C_W(s)=c_W(\log\bar x_s)^{2\mathfrak p_0}$, $\bar x_s=L^{s+1}(g^{-2}+\bar B)$ and $c_W$ is basic; it depends on $g$, it does not decrease as $d$ falls or as $g^{-2}$ grows, and $s(d)\to\infty$ as $d\to0$; if $s(d)\ge1$, then $C_W(s)\le L^sd<LC_W(s)$ at $s=s(d)$, so that $d$ is about $C_W(s)$ units of the level of that depth, and, for $0<d<1$, $Y_{s(d)}:=\log\bar x_{s(d)}\le2\log\frac1d+C(g)$. The reference coupling $g>0$ is that of~\YMMref{74}{614}, $\mathfrak d(g)$ is the set of its first-passage data $(K,\beta)$ (\YMMref{74}{614}, \YMMref{74}{614}), $g_t$ is the running coupling at the level $t$, and $\mathsf b=a_K\beta$ is the label (\ref{sec:setup-first-passage-arithmetic}). $B_\sharp$ bounds $|L^{-1}x_{j-1}-x_j|$ as in Section~\ref{sec:setup-first-passage-arithmetic}, $\bar B=B_\sharp L/(L-1)$, $X_g=L(g^{-2}+B_\sharp)$ as in \YMMref{(1526)}{615}, $I(g)=(g^{-2}+\bar B,X_g-\bar B]$ as in \YMMref{(1527)}{615}, and $\mathfrak L(g)=\{2L^m:m\ge m_{\mathfrak L}(g)\}$ (\YMMref{(1528)}{615}, \YMMref{(1528)}{615}) is the set of the sides at $g$, $\ell_g$ its least member.

Every $(K,\beta)\in\mathfrak d(g)$ has $g^{-2}<a_K\beta\le L(g^{-2}+\bar B)$, and $(K,L^K\mathsf b)\in\mathfrak d(g)$ for every integer $K\ge1$ if $\mathsf b\in I(g)$. (H) is the condition $0<g\le\gamma^{\rm ref}_W$, where $\gamma^{\rm sd}_W\le\gamma^{\rm ref}_W\le\gamma$ are the thresholds of the choices in equations~\YMMref{(2326)--(2327)}{845}; (T) is (H) together with $\ell\in\mathfrak L(g)$ and $\mathsf b\in I(g)$, and then $\langle\cdot\rangle_K:=\langle\cdot\rangle_{K,\ell,L^K\mathsf b}$.

\emph{The field.} $V=\mathrm{Sym}_3(\mathbb R)=V_A\oplus V_E\oplus V_T$, the multiples of $\mathbf 1$, the traceless diagonal matrices and the matrices with zero diagonal, is the splitting under the group $W_3$ of the $48$ signed permutation matrices; $V_{\mathbb C}=\mathrm{Sym}_3(\mathbb C)$; $E_{\lambda\lambda'}$ are the matrix units and $e_{ij}=E_{ij}+E_{ji}$. As $F_{\mu\nu}=\varepsilon_{\mu\nu\lambda}B_\lambda$ in three dimensions, the six fields are the entries of $T_{\lambda\lambda'}=\langle B_\lambda,B_{\lambda'}\rangle$. A weight is a bounded Borel function $H$ on $\mathbb R^3$ (or on $T_\ell$) with values in $V$ or $V_{\mathbb C}$ and compact support, $|H|=\sup_x\max_{\lambda\lambda'}|H_{\lambda\lambda'}(x)|$; an isometry $rx=Rx+v$ acts by $(rH)(x)=RH(r^{-1}x)R^t$, so that a reflection acts on points and on indices.

At a site $x$ the $24$ plaquette loops are $W_{x,l}$, $l=(\alpha,\alpha')$ an ordered pair of signed directions with $\alpha'\ne\pm\alpha$; $\lambda(l)$ is the normal of the plane of $l$ and $s(l)=\mathrm{sgn}\,\alpha\,\mathrm{sgn}\,\alpha'\,\varepsilon_{\lambda(l)|\alpha||\alpha'|}$ its orientation. With $w(U)=\frac{U-U^{-1}}{2i}$, $Q_x(l,l')=\frac12\mathrm{tr}\,w(W_{x,l})w(W_{x,l'})$, $a_p=1-\frac12\mathrm{tr}\,U(\partial p)$, and, for a bond $b=\langle x,x+a_Ke_\mu\rangle$ with $\lambda_b<\lambda'_b$ the two other directions, $q_b=\frac14\sum_{\epsilon,\epsilon'=\pm}\epsilon\epsilon'Q_x((\mu,\epsilon\lambda_b),(\mu,\epsilon'\lambda'_b))$, the lattice form of $\frac12\int\mathrm{tr}(HT)$ is
\eqref{sf-lat}
($x(p)$, $x_b$ the midpoints, $\lambda(p)$ the normal of $p$). $O^T_K(f\mathbf 1)=\mathcal O(f)$ is the action density, and $\mathcal O_\lambda(\phi)=O^T_K(\phi E_{\lambda\lambda})=a_K^{-1}\sum_{\lambda(p)=\lambda}\phi(x(p))a_p$ are the three plane fields. $\bar O^T_K(H)$ is \eqref{sf-lat} with $H(x(p))$, $H(x_b)$ replaced by the mean of $H$ over the cell $x+a_K[-\frac12,\frac12)^3$, $x$ the least corner of $p$, resp.\ of $b$. A corner field is $\mathcal Q_{l,l';\xi}(\phi)=a_K^{-1}\sum_x\phi(x+a_K\xi)Q_x(l,l')$, $|\xi|_\infty\le\frac12$; a site form is $a_K^{-1}\sum_x\sum_{l,l'}v(x;l,l')Q_x(l,l')$, of norm $\|v\|=\sup_x\sum_{l,l'}|v(x;l,l')|$; a mixed source is a pair $\mathfrak s=(w,\mathsf w)$ of real weights of the plaquettes and of the bonds, with the lattice form $\mathcal O^T(\mathfrak s)=a_K^{-1}[\sum_pw(p)a_p-\sum_b\mathsf w(b)q_b]$ and $|\mathfrak s|=\sup|w|+8\sup|\mathsf w|$. $\mathfrak S(O)$ is the set of the tuples $X=(H_1,\dots,H_n)$, $n\ge0$, of continuous weights with pairwise disjoint supports in an open set $O$, and $\mathfrak F(O)$ its complex span; $rX$, $\theta F$, $\tau_vF$ and $\bar F$ act entry by entry, $F'\cdot F$ concatenates, and $A^{\nu'}_K(X)=\prod_a[O^T_K(H_a)-\nu'(O^T_K(H_a))]$, $A^\nu[\mathcal A]=\mathcal A-\nu(\mathcal A)$. A pair at scale $d$ is a pair of real weights vanishing off closed balls $\bar B(x_0,d_-)$, $\bar B(y_0,d_-)$, $0<d_-<d$.

\emph{Free values and the stress tensor.} Let
\begin{equation}\label{sf-free}
\mathcal W(H_1,H_2)=\tfrac32\iint\mathrm{tr}[H_1(x)K_3(x-y)H_2(y)K_3(x-y)]\,dx\,dy,\qquad K_3(u)=\frac{3\hat u\hat u^t-\mathbf 1}{4\pi|u|^3}.
\end{equation}
The free values are the separated connected Wick contractions, with no contraction within one insertion, as in Theorem~\ref{C-371-4}(f)(i). Use three independent colors with off-diagonal covariance $\langle B^c_\lambda(x)B^d_\mu(y)\rangle=\delta^{cd}\mathsf b^{-1}K_3(x-y)_{\lambda\mu}$ for $x\ne y$. The cumulants of the smeared bilinears formally denoted by $\frac12\int\mathrm{tr}(H_aT^{\rm free})$, $T^{\rm free}=\sum_cB^c(B^c)^t$, at separated supports are: $\mathsf b^{-2}\mathcal W(H_1,H_2)$ at two points; the kernel $3\mathsf b^{-3}k_3$, $k_3(x,y,z)=\mathrm{tr}[K_3(x-y)K_3(y-z)K_3(z-x)]$, for three action densities and $3\mathsf b^{-4}k_4$, $k_4$ the sum over the three rings $(x_{\varsigma(1)},\dots,x_{\varsigma(4)})$ of $\mathrm{tr}\prod_lK_3(x_{\varsigma(l)}-x_{\varsigma(l+1)})$, for four; and $3\mathsf b^{-n}K^{(n)}_{\underline\lambda}(\phi)$ for $n$ plane fields $\mathcal O_{\lambda_a}(\phi_a)$.

On the lattice, $3\mathsf b^{-n}K^{(n)}_K(w_1,\dots,w_n)$ is the cumulant of the plaquette forms $a_K^{-1}\sum_pw_a(p)a_p$ for three free lattice photon fields, that is with $a_p$ replaced by $\frac12\sum_c(F^c_p)^2$, the $F^c$ independent Gaussian plaquette variables of covariance $\beta^{-1}C^{\rm free}$, $C^{\rm free}$ the free lattice covariance; likewise, for mixed sources the free lattice value of the cumulant of $\mathcal O^T(\mathfrak s_1)$, $\mathcal O^T(\mathfrak s_2)$ is the cumulant obtained when every $a_p$ and every $Q_x(l,l')$ is replaced by its quadratic part in three such independent fields; it is $\frac32\mathsf b^{-2}\mathrm{Tr}(\mathsf W_1C^{\rm free}\mathsf W_2C^{\rm free})$, $\mathsf W_a$ the symmetric matrix with $\mathcal O^T(\mathfrak s_a)=(2a_K)^{-1}\sum_cF^{c\,t}\mathsf W_aF^c$ after this replacement. $\mathsf w_0=\frac1{12}388^{-6}$.

For $\eta\in C^2_c(\mathbb R^3;\mathbb R^3)$ put $H^{\rm cons}[\eta]=(\nabla\cdot\eta)\mathbf 1-\nabla\eta-(\nabla\eta)^t$, $(\nabla\eta)_{ij}=\partial_i\eta_j$, and $\|\eta\|_{C^2}=\sup|\eta|+\sup|\nabla\eta|+\sup|\nabla^2\eta|$; then $\frac12\mathrm{tr}(H^{\rm cons}[\eta]T)=\sum_{ij}\partial_i\eta_j\Theta_{ij}$ with the stress tensor $\Theta_{ij}=\frac12\delta_{ij}\mathrm{tr}\,T-T_{ij}$: insertions of $H^{\rm cons}[\eta]$ express the conservation of $\Theta$, and $\frac12\mathrm{tr}(s_{ij}T)=\Theta_{ij}$ for $s_{ij}=\delta_{ij}\mathbf 1-e_{ij}$. For a unit vector $n$, $\Theta_n=\mathbf 1-2nn^t$.

\emph{Chessboards and limits.} A board of side $\mathsf l=L^{-i}$, $i\in\mathbb Z$, divides space into cubes, its cells, by bond mirrors, its walls; for $\ell'\in\mathfrak L(g)$, $\ell'\ge2\mathsf l$, and $D$ a function of the bonds of one cell, $\|D\|_{{\rm ch},\ell'}$ is the $N$-th root of $|\langle\prod_\kappa D_\kappa\rangle_{K,\ell',\beta}|$, the $D_\kappa$ being the $N=(\ell'/\mathsf l)^3$ mirror images of $D$ in the cells of $T_{\ell'}$, conjugated after an odd number of reflections as in the main chessboard convention. A Gibbs sequence at $g$ is $(K_j,\beta_j,\nu_j)$ with $(K_j,\beta_j)\in\mathfrak d(g)$, $K_j\to\infty$, $\nu_j\in\mathcal G_{K_j,\beta_j}$; a limit at $g$ is $\sigma(X)=\lim_j\nu_j(A^{\nu_j}_{K_j}(X))$, $X\in\mathfrak S(\mathbb R^3)$, along a Gibbs sequence with $a_{K_j}\beta_j\to\mathsf b$, its label.

Its trace part $\sigma_0(f_1,\dots,f_n)=\sigma(f_1\mathbf 1,\dots,f_n\mathbf 1)$ is a scalar limit, to which Theorem~\YMMref{AP}{780} (the scalar action density on $\mathbb R^3$), with Theorem~\YMMref{AL}{746} applies, with a canonical decomposition $(\mathsf X_0,\lambda_0,(\sigma_{0,\zeta})_{\zeta\in\mathsf X_0})$. With $\theta$ the reflection in $\{x_3=0\}$, $\mathcal H(\sigma)$ is the Hilbert space of the form $\langle[F'],[F]\rangle=\sigma(\theta\bar F'\cdot F)$ on $\mathfrak F(\{x_3>0\})$, and $\Omega=[\emptyset]$. More generally, let $\sigma$ have the properties of Theorem~\ref{sf-gibbs}(c), as every limit at $g$ has, and let $u$ be a unit vector at which $\sigma$ is reflection positive: $\sigma(\theta_u\bar F\cdot F)\ge0$ for $F\in\mathfrak F(\{u\cdot x>0\})$, $\theta_u$ the reflection in $\{u\cdot x=0\}$, acting on points and on indices by $E\mapsto\Theta_uE\Theta_u$.

The \emph{picture of $u$} is formed in the same way: the Hilbert space $\mathcal H_u$ of the form $\langle[F'],[F]\rangle_u=\sigma(\theta_u\bar F'\cdot F)$ on $\mathfrak F(\{u\cdot x>0\})$, $\Omega_u=[\emptyset]_u$, the self-adjoint $H_u\ge0$ with $e^{-tH_u}[F]_u=[\tau_{tu}F]_u$, $t\ge0$, the unitary group $\mathcal U^u(v)[F]_u=[\tau_vF]_u$, $v\in u^\perp$, which commutes with it, and their joint spectral measure $\mathcal P_u$ on $[0,\infty)\times u^\perp$, with $e^{-tH_u}\mathcal U^u(v)=\int e^{-tE+iv\cdot q}\,d\mathcal P_u(E,q)$; the picture of $e_3$ is $(\mathcal H(\sigma),H,\mathcal U,\mathcal P)$ of Theorem~\ref{sf-wight}(a). On a torus, $\mathcal S_{n;K}(X)=\langle O^T_K(H_1);\dots;O^T_K(H_n)\rangle_K$.

\emph{Sources.} Theorem~\YMMref{D}{8} follows an insertion through the levels as a source under a majorant. A \emph{majorant} on $\Lambda_{K,\ell}$ is a pair $(\mathsf s,S)$ of a number $\mathsf s\ge0$, its \emph{size}, and a set $S$ of plaquettes; its weight is the plaquette source $\mathsf s\mathbf 1_S$, and $\mathsf m_j=12M^3\alpha_0^2L^{K-j}\mathsf s$, $j\ge0$, are the numbers that the calculus of Theorem~\YMMref{D}{8} attaches to that weight at the level $j$, $\alpha_0\in(0,1]$ a basic constant of the calculus. A \emph{source under $(\mathsf s,S)$} is a triple of: a \emph{first member} $F$, a continuous function of the configurations of $\Lambda_{K,\ell}$, invariant under the gauge transformations, with $|F(U)|\le a_K^{-1}\mathsf s\sum_{p\in S}a_p(U)$; a local family at the level $0$ of norm at most $\mathsf c^\flat\mathsf m_0$, $\mathsf c^\flat$ a basic number of the calculus; and local families at the levels $j\ge1$ of norms at most $\frac12\mathsf m_j$; each family vanishes at the regions of its level that hold no plaquette of $S$, and each sums at its level to $F$, up to the value at the trivial field.

A \emph{source system under majorants} is a tuple of sources under majorants $(\mathsf s_a,S_a)$; all the numbers of the calculus are then formed at the weights $\mathsf s_a\mathbf 1_{S_a}$. A source under $(\mathsf s,S)$ is a source under every $(\mathsf s',S')$ with $\mathsf s'\ge\mathsf s$ and $S'\supset S$, the majorant entering only the right members of three inequalities and the two clauses of vanishing; the sum of two sources, member by member, is a source under the sum of the sizes and the union of the sets, and $c$ times a source is one under $|c|$ times the size. For the value at short distance the calculus reads besides, of a real source, the kernel $\mathsf t(p,p')=\mathsf t(p',p)$ of its first member to the second order at a small field, vanishing unless $p,p'\in S$, with $2\sum_{p'}|\mathsf t(p,p')|\le\mathsf s$.

Lemma~\ref{sf-src} says which insertions are sources.

\subsection{Principal results}

\begin{theorem}[The tori]\label{sf-tori}
Assume \textup{(T)}, $n\ge2$, $d>0$, $\theta_n=1/(1260n(8n+3))$, and let the supports of $X=(H_1,\dots,H_n)$ be pairwise $\ge d$ apart.
\begin{enumerate}
\item[(a)] There are an exponent $\mathsf p$ depending only on the basic constants (one may take $\mathsf p=2\mathfrak p_0+2$), a number $C^{\rm od}$ depending only on $n$ and on the constant $C_{\rm bf}$ of the comparison of two cutoffs at the label $\mathsf b$, and a threshold $K_1$ depending besides on $g$, $\mathsf b$, $\ell$, $d$, on the threshold $K_{\rm bf}$ of that comparison and on the number $\Theta_\star$ ($C_{\rm bf}$, $K_{\rm bf}$ of Definition~\textup{\ref{IV-62-62}}(5), formed at a compact interval $J\subset I(g)$ that holds $\mathsf b$, so that they read the label only through $J$ (Corollary~\textup{\ref{IV-62-66}}(a)), $\Theta_\star$ of Section~\YMMref{113}{786}), such that for $K\ge K_1$ and Lipschitz weights
\begin{equation}\label{sf-rate}
|\mathcal S_{n;K+1}(X)-\mathcal S_{n;K}(X)|\le6^nC^{\rm od}\ell^{3(n+1)}K^{\mathsf p}L^{-\theta_nK}\prod_a(|H_a|+\mathrm{Lip}\,H_a).
\end{equation}
So $\mathcal S_n=\mathcal S^{\ell,\mathsf b}_n=\lim_K\mathcal S_{n;K}$ exists along the full sequence of cutoffs, with $|\mathcal S_{n;K}(X)-\mathcal S_n(X)|\le6^nC^{\rm od}\ell^{3(n+1)}\sum_{K'\ge K}K'^{\mathsf p}L^{-\theta_nK'}\prod_a(|H_a|+\mathrm{Lip}\,H_a)$; it is $n$-linear and symmetric, $|\mathcal S_n(X)|\le60^nK^{\rm ub}_n(g,\ell,\frac d2)\prod_a|H_a|$ with the number of Lemma~\ref{sf-sym}(b), and the centred moments converge to $\sigma^{\ell,\mathsf b}(X)=\sum_\pi\prod_{B\in\pi}\mathcal S_{|B|}(X_B)$, the sum over the partitions $\pi$ of the entries without singletons; $\sigma^{\ell,\mathsf b}(\emptyset)=1$, $\sigma^{\ell,\mathsf b}((H))=0$. With $\bar O^T_K$, \eqref{sf-rate} holds for bounded Borel weights vanishing off closed sets $\Sigma_a$ pairwise $\ge d$ apart, with $\prod_a|H_a|$ on the right and a number $C^{\rm av}$ of the same kind as $C^{\rm od}$; the limit extends $\mathcal S_n$, and in each argument $A\mapsto\mathcal S_n(\dots,1_AE,\dots)$, $E\in V$, is a countably additive measure on $\Sigma_a$ vanishing on Lebesgue null sets, so it is given by a density in $L^1(\Sigma_a;V)$.
\item[(b)] $\mathcal S_n(rX)=\mathcal S_n(X)$ for $rx=Rx+v$, $R\in W_3$, $v\in\mathbb R^3$; $\sigma^{\ell,\mathsf b}(\theta\bar F\cdot F)\ge0$ for the reflection $\theta$ in a plane $\{x_\mu=c\}$, acting on points and indices, and $F\in\mathfrak F(\{c<x_\mu<c+\frac\ell2\})$; at $H_a=f_a\mathbf 1$, $\mathcal S_n$ is the limit of the truncated functions of the action density.
\item[(c)] A class insertion over a closed set $\Sigma$ is a finite sum $\mathcal E$ of terms $\mathcal O(\phi1_\lambda)=a_K^{-1}\sum_{\lambda(p)=\lambda}\phi(x(p))a_p$ and $\mathcal Q_{l,l';\xi}(\phi)$ with real Lipschitz $\phi$ supported in $\Sigma$; its symbol $\mathrm{Sym}\,\mathcal E$ is the sum of the weights $\phi E_{\lambda\lambda}$, resp.\ $\phi\,s(l)s(l')(E_{\lambda(l)\lambda(l')}+E_{\lambda(l')\lambda(l)})$, of its terms, and $\|\mathcal E\|$ is the sum of their numbers $\sup|\phi|+\mathrm{Lip}\,\phi$. There is $C^{\rm u}$ depending only on $n,g,\ell,d$ such that for $K\ge K_1+1$ and class insertions $\mathcal E_a$, $\mathcal E'_a$ over closed sets pairwise $\ge d$ apart with $\mathrm{Sym}\,\mathcal E_a=\mathrm{Sym}\,\mathcal E'_a$,
\[|\langle\mathcal E_1;\dots;\mathcal E_n\rangle_K-\langle\mathcal E'_1;\dots;\mathcal E'_n\rangle_K|\le C^{\rm u}K^{\mathsf p}L^{-\theta_nK}\prod_a(\|\mathcal E_a\|+\|\mathcal E'_a\|),\]
and the Lipschitz numbers drop out of the bound if $\mathcal E_a$, $\mathcal E'_a$ carry the same weights at the same sites. So $\lim_K\langle\mathcal E_1;\dots;\mathcal E_n\rangle_K=\mathcal S_n(\mathrm{Sym}\,\mathcal E_1,\dots,\mathrm{Sym}\,\mathcal E_n)$: the bond form, the clover form and the single corners give one field, and $\mathrm{Sym}\,O^T_K(H)=H$.
\item[(d)] Let $(H_1,H_2)$ be a pair at scale $d<1$ with $3d\le|x_0-y_0|\le\frac\ell4$. There are a threshold $s_1(g,\ell)$ on the depth and a number $\mathsf e^T_{g,\ell}(d)$, tending to $0$ like $(\log\frac1d)^{-2\mathfrak p_0}$, such that if $s(d)\ge s_1(g,\ell)$, then for Lipschitz pairs, and for bounded ones with $\bar O^T_K$,
\begin{equation}\label{sf-two}
|\mathcal S_2(H_1,H_2)-\mathsf b^{-2}\mathcal W(H_1,H_2)|\le\mathsf e^T_{g,\ell}(d)\,\mathsf b^{-2}|H_1||H_2|,
\end{equation}
and the same holds for class insertions over the two balls, with $\mathrm{Sym}\,\mathcal E_a$ for $H_a$. At finite cutoff, for $(K,\beta)\in\mathfrak d(g)$ with $K\ge s(d)+2$ and $4a_K\le d-d_-$, and for continuous pairs, \eqref{sf-two} holds for $\langle O^T_K(H_1);O^T_K(H_2)\rangle_{K,\ell,\beta}$ with $a_K\beta$ for $\mathsf b$ and the further error $\mathsf b^{-2}[c_K|H_1||H_2|+\frac98\mathrm{rm}_K(H_1,H_2)]$, where $c_K=28000\,C_{\rm lat}(1+C_{\mathsf k})(a_K/d)^2$ with $C_{\mathsf k}$ basic, and $\mathrm{rm}_K\ge0$ depends only on $K$, $d$, the two weights and the choice of $O^T_K$ or $\bar O^T_K$, and tends to $0$ as $K\to\infty$. One may take $s_1(g,\ell)\le\hat s(g,\ell):=\max\{s_\flat(g),\lceil2\Lambda_\ell/\log L\rceil\}$, with a depth $s_\flat(g)$ depending only on $g$ and $\Lambda_\ell=(12\log\ell/(c'A_0^2))^{1/(2\mathfrak p_0)}$, $c'$ basic; from this depth on, the error of the finite-cutoff law against the lattice Wick value of Lemma~\ref{sf-depth} is independent of $\ell$. Further, under the same conditions on $d$:
(i) $|\mathcal S_2(rH_1,rH_2)-\mathcal S_2(H_1,H_2)|\le10\,\mathsf e^T_{g,\ell}(d)\mathsf b^{-2}|H_1||H_2|$ for $rx=x_0+R(x-x_0)$, $R\in O(3)$;
(ii) for the reflection $\theta$ in a plane $\{x_\mu=c\}$, a point $z$ at distance $\frac32d$ from it, a Lipschitz function $\phi$ with $1_{B(z,d/128)}\le\phi\le1_{B(z,d/64)}$ and $E\in V$: $\mathcal S_2(\theta(\phi E),\phi E)\ge(\mathsf w_0-\mathsf e^T_{g,\ell}(d))\mathsf b^{-2}\mathrm{tr}\,E^2$; so no non-zero real combination of the six fields vanishes.
And at every $d$: (iii) for $0<d_-<d<1$, $\frac32d\le c-x_{0,\mu}\le\frac\ell2-\frac32d$ and real Lipschitz $H$ vanishing off $\bar B(x_0,d_-)$, $0\le\mathcal S_2(\theta H,H)\le3600\,C^*_3(g,\ell)\mathsf b^{-2}|H|^2$, with $C^*_3(g,\ell)\le c_3(g)(1+\log\ell)^8e^{16\Lambda_\ell}$ independent of $d$.
\item[(e)] For $q\ge2$, $\eta\in C^2(T_\ell;\mathbb R^3)$ and Lipschitz $H_2,\dots,H_q$, the $q$ supports pairwise $\ge d$ apart: $\mathcal S_q(H^{\rm cons}[\eta],H_2,\dots,H_q)=0$, and $\sigma^{\ell,\mathsf b}((H^{\rm cons}[\eta])\cdot X')=0$ for $X'\in\mathfrak S(T_\ell)$ off $\mathrm{supp}\,\eta$. There are $C,\mathsf p,K_0$ depending only on $g,\ell,q,d$ such that for all $(K,\beta)\in\mathfrak d(g)$ with $K\ge K_0$ and all bounded $H_a$,
\[|\langle O^T_K(H^{\rm cons}[\eta]);O^T_K(H_2);\dots;O^T_K(H_q)\rangle_{K,\ell,\beta}|\le CK^{\mathsf p}L^{-K/56}\|\eta\|_{C^2}\prod_{a\ge2}|H_a|,\]
also with $\bar O^T_K$; the bound is uniform in the label, and it stays true, with another $C$, if $O^T_K(H^{\rm cons}[\eta])$ is changed by a site form plus a plaquette form, carried within $2a_K$ of $\mathrm{supp}\,\eta$, whose weights are at most $36a_K\|\eta\|_{C^2}$. At $H_a=h_a\mathbf 1$ this is the conservation of the stress tensor among action densities.
\item[(f)] The numbers $\mathsf p$, $C^{\rm od}$, $C^{\rm av}$ and $K_1$ of (a) may be taken the same at all the labels of a compact interval $J\subset I(g)$ (Lemma~\ref{sf-cut}(b)), and $\mathsf b\mapsto\mathcal S^{\ell,\mathsf b}_n(X)$ is continuous on $I(g)$. For Lipschitz weights, and for bounded ones with $\bar O^T_K$: if $\mathsf b$ is an inner point of $I(g)$ and $a_K\beta_K\to\mathsf b$, then $\langle O^T_K(H_1);\dots;O^T_K(H_n)\rangle_{K,\ell,\beta_K}\to\mathcal S^{\ell,\mathsf b}_n(X)$ along the full sequence of cutoffs.
\end{enumerate}
By Lemma~\ref{sf-dil}, (a)--(f) hold on every torus of side $2L^m$, $m\ge0$, and at every label $\mathsf b>0$, with the constants of a triple $(g',L^z\mathsf b,L^z\ell)$ with \textup{(T)}, taken at the distance $L^zd$, and powers of $L^z$; for large $z$ one may take $g'=(L^z\mathsf b/2)^{-1/2}$, and then $L^z\mathsf b$ is an inner point of $I(g')$.
\end{theorem}

\begin{theorem}[Periodic Gibbs states and limits on $\mathbb R^3$]\label{sf-gibbs}
Assume \textup{(H)}.
\begin{enumerate}
\item[(a)] Let $n\ge2$ and $\delta_0,\varrho_0>0$. For $(K,\beta)\in\mathfrak d(g)$ beyond a threshold depending only on $g,n,\delta_0,\varrho_0$, all $\nu,\nu'\in\mathcal G_{K,\beta}$ and all real weights vanishing off closed sets of diameter $\le\varrho_0$ pairwise $\ge2\delta_0$ apart,
\[|\nu(A^{\nu'}_K(X))|\le C^{\rm lc}_n(g,\delta_0,\varrho_0)\prod_a|H_a|,\]
a number independent of the cutoff, the states, the positions and the distances of the supports; the same holds for the measures on the tori $T_\ell$, $\ell\in\mathfrak L(g)$, if the supports lie in a cube of side $\ell-2$. With a count in $n$, through the following thresholds, which depend only on the basic constants, $g$, $\delta$ and $i$. For $0<\delta\le\frac14$ let $z_\sharp:=0$ and $g_-:=g$ if $g\le\gamma^{\rm sd}_W$; for $\gamma^{\rm sd}_W<g\le\gamma^{\rm ref}_W$ let $z_\sharp$ be the least integer $z\ge1$ with $g_-^{-2}:=L^{z-1}(g^{-2}-\bar B)-B_\sharp\ge(\gamma^{\rm sd}_W)^{-2}$. Let $i_E(g,\delta):=z_\sharp+i_1$, $i_1$ the least integer $\ge1$ with $L^{i_1-1}(g_-^{-2}-\bar B)-B_\sharp\ge\mathsf x_E(\delta)$, $\mathsf x_E(\delta)\le\mathsf x_*/\delta$ the number of the proof of Lemma~\ref{sf-stop}, $\mathsf x_*$ basic; so $L^{i_E(g,\delta)}\le c_\sharp(g)/\delta$ with $c_\sharp(g)$ depending only on $g$.

Let $K_E(i)=K_E(i;g,\delta):=i+s_\sharp(i)$, $s_\sharp(i)$ the depth of the distance $\delta$ formed at the coupling of inverse square $L^{i-1}(X_g+2\bar B)-B_\sharp$. For $g\le\gamma^{\rm sd}_W$ these are thresholds of Lemma~\ref{sf-stop}, the depth not decreasing as the inverse square of the coupling grows; for $\gamma^{\rm sd}_W<g$ they are those of Lemma~\ref{sf-stop} read at the stopped coupling $g_{(z_\sharp)}\le\gamma^{\rm sd}_W$ of the datum. Then for $i\ge i_E(g,\frac1{8n})$, $(K,\beta)\in\mathfrak d(g)$ with $K\ge\max\{K_E(i;g,\frac1{8n}),i+\log(4n)/\log L\}$, and weights vanishing off sets of diameter $\le a_i/(4n)$ pairwise more than $a_i$ apart, in every periodic Gibbs state and on every torus of $\mathfrak L(g)$, $|\nu(A^\nu_K(X))|\le(2c_gn^4Y_i^{q_E}L^{11i/2})^n\prod_a|H_a|$, where $c_g$ depends only on $g$, $q_E=(4+\mathfrak p_0)(2\mathfrak p_0-1)+10\mathfrak p_0$ and $Y_i=2[\log(8nL^2c_W)+2\mathfrak p_0\log(4\mathfrak p_0)+i\log L+\log(X_g+2\bar B)]$.
\item[(b)] For $p_T>\frac{11}2$ and $0<\delta\le\frac14$ there are $C_T$ depending only on $g,\delta,p_T$ and thresholds $K_E(i)$, those of (a) at the margin $\delta$ for $i\ge i_E(g,\delta)$ and depending only on $g,\delta,i$ for the finitely many smaller $i$, such that for $i\ge0$, $(K,\beta)\in\mathfrak d(g)$ with $K\ge K_E(i)$, every board of side $a_i$ and every cell of it, every $\ell\in\mathfrak L(g)$, every bounded $H:\mathbb R^3\to V_{\mathbb C}$ vanishing within $\delta a_i$ of the complement of the cell, and $c$ the mean of $O^T_K(H)$ on any torus of $\mathfrak L(g)$ or in any periodic Gibbs state:
\[\|O^T_K(H)-c\|_{{\rm ch},\ell}\le C_TL^{ip_T}\mathsf b^{-1}|H|.\]
\item[(c)] Every Gibbs sequence at $g$ has a subsequence along which $a_{K_j}\beta_j\to\mathsf b\in[g^{-2},L(g^{-2}+\bar B)]$ and the limit exists on $\mathfrak S(\mathbb R^3)$; a Gibbs sequence whose centred correlations converge at all tuples of $C^1$ weights converges on all of $\mathfrak S(\mathbb R^3)$. A limit at $g$ is real, symmetric and multilinear, $\sigma(\emptyset)=1$, $\sigma((H))=0$; it has the bounds of (a), so that for fixed separated tuples $X,Y$ there is $R(X,Y)$ with $\sup_{|v|\ge R(X,Y)}|\sigma(X\cdot\tau_vY)|<\infty$; choose $R(X,Y)$ so the translated groups stay a fixed positive distance apart; its truncated functions are the limits of the cumulants; $\sigma(rX)=\sigma(X)$ for $r\in\mathbb R^3\rtimes W_3$; and $\sigma(\theta\bar F\cdot F)\ge0$ for the reflection $\theta$ in every plane $\{n\cdot x=c\}$, $n\in\{e_\mu,e_\mu\pm e_\lambda\}$, and $F\in\mathfrak F(\{n\cdot x>c\})$. Its trace part is the scalar limit of its sequence. Along tori $T_{\ell_j}$, $\ell_j\in\mathfrak L(g)$, $\ell_j\to\infty$, at data $(K_j,\beta_j)\in\mathfrak d(g)$ with $K_j\to\infty$, the centred correlations have limit points with the same bounds and invariance, positive at the coordinate mirrors.
\item[(d)] There are $d_0(g)>0$ and $\mathsf e^\infty_g(d)$ depending only on $g$, with $\sup_{d<d_0(g)}(\log\frac1d)^{\mathfrak p_0/4}\mathsf e^\infty_g(d)<\infty$, such that for Lipschitz pairs at scale $d<d_0(g)$ with $3d\le|x_0-y_0|\le4d$:
\[|\mathcal S(H_1,H_2)-\mathsf b^{-2}\mathcal W(H_1,H_2)|\le\mathsf e^\infty_g(d)\mathsf b^{-2}|H_1||H_2|\]
for $\mathcal S=\mathcal S^{\ell,\mathsf b}_2$, every $\ell\in\mathfrak L(g)$ and $\mathsf b\in I(g)$, and for $\mathcal S=\sigma$, every limit at $g$; at finite cutoff the same holds in every periodic Gibbs state and on every torus of $\mathfrak L(g)$, under the conditions and with the further error of Theorem~\ref{sf-tori}(d), which are independent of $\beta$, $\ell$ and the state, the condition $K\ge s(d)+2$ being replaced by $K\ge s(d)+3$ if $\gamma^{\rm sd}_W<g$ (the depth of the dilated distance $L^zd$ at the coupling of Lemma~\ref{sf-dil}(b), shifted by $z$, is at most $s(d)+1$). So the label of a limit is determined by the limit; $\sigma(\theta(\phi E),\phi E)\ge(\mathsf w_0-\mathsf e^\infty_g(d))\mathsf b^{-2}\mathrm{tr}\,E^2$ with $\theta,\phi$ as in Theorem~\ref{sf-tori}(d)(ii), so that no non-zero real combination of the six fields vanishes in any limit; and $|\sigma(rH_1,rH_2)-\sigma(H_1,H_2)|\le10\,\mathsf e^\infty_g(d)\mathsf b^{-2}|H_1||H_2|$ for $rx=x_0+R(x-x_0)$, $R\in O(3)$.
\item[(e)] Let $0<g\le\gamma$, $\ell\in\mathfrak L(g)$, $(K,\beta)\in\mathfrak d(g)$ with $K\ge K_\circ(g)$, a threshold depending on $g$, $0<16r\le d'\le\frac\ell8$, $s(r)\le K$, $|x_1-x_2|=d'$, and let $\phi_a$ be real Lipschitz functions supported in $B(x_a,r)$ with $\mathrm{Lip}\,\phi_a\le|\phi_a|/(\vartheta r)$, $0<\vartheta\le1$. There are a finite function $\varepsilon^{\rm cl}(g,\ell,\cdot)$ on $(0,\frac\ell8]$, with $\varepsilon^{\rm cl}(g,\ell,d')\le C(g,\ell)(\log\frac1{d'})^{19\mathfrak p_0}d'^{1/2}$ for small $d'$, and basic $C_\flat$, $C'_\flat$ such that
\[\big|\langle\mathcal O_\lambda(\phi_1);\mathcal O_\mu(\phi_2)\rangle_{K,\ell,\beta}-\mathsf b^{-2}\mathcal W(\phi_1E_{\lambda\lambda},\phi_2E_{\mu\mu})\big|\le\Big[\varepsilon^{\rm cl}+\tfrac32C_\flat\Big(\frac{a_K}{\vartheta r}+\frac{d'^3}{\ell^3}\Big)\Big]\Big(\frac r{d'}\Big)^6\frac{|\phi_1||\phi_2|}{\mathsf b^2};\]
for $O^T_K(\phi_1e_{\lambda\lambda'})$, $\lambda\ne\lambda'$, against $\mathcal O_\gamma(\phi_2)$ the same holds with $\mathcal W(\phi_1e_{\lambda\lambda'},\phi_2E_{\gamma\gamma})$ and $8\varepsilon^{\rm cl}+3C'_\flat(\frac{a_K}{\vartheta r}+\frac{d'^3}{\ell^3})$ in the bracket; and for all real mixed sources $\mathfrak s_a$ carried by $B(x_a,r)$, against the free lattice value: $|\langle\mathcal O^T(\mathfrak s_1);\mathcal O^T(\mathfrak s_2)\rangle_{K,\ell,\beta}-\frac32\mathsf b^{-2}\mathrm{Tr}(\mathsf W_1C^{\rm free}\mathsf W_2C^{\rm free})|\le\varepsilon^{\rm cl}(g,\ell,d')(\frac r{d'})^6\mathsf b^{-2}|\mathfrak s_1||\mathfrak s_2|$. For $g$ below a basic threshold, the first display holds in every periodic Gibbs state at finite cutoff, for bumps of radius $r$ in the balls of radius $\frac18d'$ about two points at a distance $d'$ below a threshold depending on $g$, with $E_2(g,\frac12d')\le C(g)(\log\frac1{d'})^{19\mathfrak p_0}d'^{1/2}$ for $\varepsilon^{\rm cl}$, where $E_2(g,t)=\varepsilon_2(g,t)+C(g)t^{3/2}+\frac32C_\flat(t/\ell_g)^3$, $\varepsilon_2$ of Lemma~\ref{sf-jet}(a), with $(2r/d')^6$ for $(r/d')^6$ and without the term $d'^3/\ell^3$, the bracket being independent of the state; and every limit at $g$ has, for $\lambda,\mu\le3$ and $0<|x-y|<d_\sigma(g)$, a threshold depending on $g$,
\[
\begin{multlined}
\big|s_2(x,y;E_{\lambda\lambda},E_{\mu\mu})-\tfrac32\mathsf b^{-2}[K_3(x-y)_{\lambda\mu}]^2\big|\le(\tfrac3\pi)^2\varepsilon_\sigma(g,|x-y|)\mathsf b^{-2}|x-y|^{-6},\qquad\varepsilon_\sigma(g,t) \\
\le C(g)(\log\tfrac1t)^{19\mathfrak p_0}t^{1/2},
\end{multlined}
\]
$s_2$ the kernel of Theorem~\ref{sf-wight}(d); the sum over $\lambda,\mu$ of these kernels is the kernel of two action densities.
\item[(f)] Let $0<g\le\gamma$, $n\ge2$, $d,R>0$. There are $C,\mathsf p$ depending only on $g,n,d,R$ and a threshold $K_0$ such that for $(K,\beta)\in\mathfrak d(g)$ with $K\ge K_0$, all $\nu\in\mathcal G_{K,\beta}$, $\eta\in C^2_c(\mathbb R^3;\mathbb R^3)$ and bounded $H_2,\dots,H_n$, the $n$ supports in a ball of radius $R$ and pairwise $\ge d$ apart,
\[\Big|\nu\Big(A^\nu[O^T_K(H^{\rm cons}[\eta])]\prod_{a\ge2}A^\nu[O^T_K(H_a)]\Big)\Big|\le CK^{\mathsf p}L^{-K/36}\|\eta\|_{C^2}\prod_{a\ge2}|H_a|,\]
and the same holds for the joint cumulant; so $\sigma((H^{\rm cons}[\eta])\cdot X)=0$ for every limit at $g$ and every $X$ off $\mathrm{supp}\,\eta$. Every slot that holds a conservation or stress insertion brings its own rate: let $0<c<\frac1{28}$ and $1\le n_1\le n$, let the slot $a\le n_1$ hold any of $O^T_K(H^{\rm cons}[\eta_a])$, $\mathrm{St}_K[\eta_a]$, $\Delta_K[\eta_a]$, $\mathrm{Rm}_K[\eta_a]$ of Proposition~\ref{sf-stress}, $\eta_a\in C^2_c(\mathbb R^3;\mathbb R^3)$, and the slot $a>n_1$ a tensor entry $O^T_K(H_a)$ with a continuous weight, the $n$ supports lying in a ball of radius $R$ and pairwise $\ge d$ apart; then, with $C,\mathsf p$ depending only on $g,n,R,c$ and for $K$ beyond a threshold depending on $g,n,c,d$, in every $\nu\in\mathcal G_{K,\beta}$
\[\Big|\nu\Big(\prod_{a=1}^nA^\nu[X_a]\Big)\Big|\le CK^{\mathsf p}L^{-cn_1K}\prod_{a\le n_1}\|\eta_a\|_{C^2}\prod_{a>n_1}|H_a|,\]
$X_a$ the entry of the slot $a$, and the same holds for the joint cumulant; so the stress insertion has its tensor among any number of tensor entries. At two points a shorter road, one small seminorm against the crude bound of a tensor entry, gives: if $\eta,\eta'\in C^2_c(\mathbb R^3;\mathbb R^3)$ and a bounded $H'$ lie in a ball of radius $R$ and $\mathrm{supp}\,\eta$ is at distance $\ge d$ from $\mathrm{supp}\,\eta'$ and from $\mathrm{supp}\,H'$, then $|\nu(A^\nu[O^T_K(H^{\rm cons}[\eta])]A^\nu[O^T_K(H')])|\le CK^{\mathsf p}L^{-K/60}\|\eta\|_{C^2}|H'|$, and the stress insertion has its tensor in every state: $|\nu(A^\nu[\Delta_K[\eta]]A^\nu[Z])|\le CK^{\mathsf p}L^{-K/60}\|\eta\|_{C^2}\mathsf z$ for $Z=O^T_K(H')$, $\mathsf z=|H'|$, and for $Z=\mathrm{St}_K[\eta']$, $\Delta_K[\eta']$, $\mathsf z=\|\eta'\|_{C^2}$. With two small seminorms the rates are better than those of the display above at $n=2$: two stress insertions of Proposition~\ref{sf-stress} have $|\nu(A^\nu[\mathrm{St}_K[\eta']]A^\nu[\mathrm{St}_K[\eta]])|\le CK^{\mathsf p}L^{-K/14}\|\eta'\|_{C^2}\|\eta\|_{C^2}$, and two conservation insertions have $|\nu(A^\nu[O^T_K(H^{\rm cons}[\eta'])]A^\nu[O^T_K(H^{\rm cons}[\eta])])|\le CK^{\mathsf p}L^{-K/20}\|\eta'\|_{C^2}\|\eta\|_{C^2}$.
\end{enumerate}
\end{theorem}

\begin{theorem}[Invariance, the Wightman field, the decomposition]\label{sf-wight}
Assume \textup{(H)} and let $\sigma$ be a limit at $g$ with label $\mathsf b$.
\begin{enumerate}
\item[(a)] $\mathcal H(\sigma)$ is separable; $e^{-tH}[F]=[\tau_{te_3}F]$, $t\ge0$, and $\mathcal U(v)[F]=[\tau_{(v,0)}F]$ define a self-adjoint $H\ge0$ with $H\Omega=0$ and a unitary group $\mathcal U(v)=e^{iv\cdot P}$, $v\in\mathbb R^2$, commuting with it; the joint spectral measure $\mathcal P$ of $(H,P_1,P_2)$ has $\mathcal P(\{E<|q|\})=0$; and in the picture of each unit normal $u=n/|n|$, $n\in\{e_\mu,e_\mu\pm e_\lambda\}$, of the nine families (Setting) the joint spectral measure has $\mathcal P_u(\{(E,q):E<|q|\})=0$, $|q|$ the Euclidean length of $q\in u^\perp$. The eight $w\in W_3$ with $we_3=e_3$ act by unitary operators $[F]\mapsto[wF]$ that fix $\Omega$ and commute with $H$. The map $[F]_0\mapsto[F\mathbf 1]$, every entry $f$ replaced by $f\mathbf 1$, carries $\mathcal H(\sigma_0)$ isometrically onto a closed subspace $\mathcal H_0$ that reduces $H$, $\mathcal U$ and $\mathcal P$; on $\mathcal H_0$ these are the operators of the scalar theory of Theorem~\YMMref{AP}{780}.
\item[(b)] For a cell $Q=(v,0)+(-\frac{\mathsf l}2,\frac{\mathsf l}2)^2\times(0,\mathsf l)$, $\mathsf l=L^{-i}$, and $X=(h_1,\dots,h_n)\in\mathfrak S(Q)$ with supports $\ge\delta$ from $\partial Q$ and $\ge2\delta$ apart, the slab operator $\mathsf B_{\mathsf l}(X)[F]=[X\cdot\tau_{\mathsf le_3}F]$ is bounded, $\|\mathsf B_{\mathsf l}(X)\|\le C^{\rm sl}_n(g;\mathsf l,\delta)\prod_a|h_a|$, with $\mathsf B_{\mathsf l}(X)\Omega=[X]$ and $\mathsf B_{\mathsf l}(X)^*=\mathsf B_{\mathsf l}(\theta_{e_3,\mathsf l/2}X)$. There is exactly one norm-continuous family of bounded operators $\mathsf B(s,s';e)$, $s,s'>0$, linear in $e\in V$, with
\[\sigma(\theta\tau_{se_3}\bar G'\cdot(h)\cdot\tau_{s'e_3}G)=\int\langle[G'],\mathcal U(\vec x)\mathsf B(s+x_3,s'-x_3;h(x))\mathcal U(-\vec x)[G]\rangle\,dx\]
for $G,G'\in\mathfrak F(\{x_3>0\})$ and $h$ supported in $\{-s<x_3<s'\}$; $e^{-tH}\mathsf B(s,s';e)e^{-t'H}=\mathsf B(s+t,s'+t';e)$, $\mathsf B(s,s';e)^*=\mathsf B(s',s;\Theta_{e_3}e\Theta_{e_3})$, and $\|\mathsf B(s,s';e)\|\le C_{\rm pt}(\min\{s,s'\})|e|$ with $C_{\rm pt}(\eta)\le C\eta^{-3-p_T}$ for $0<\eta\le1$ and every $p_T>\frac{11}2$, $C$ depending only on $g$ and $p_T$. For every $p_V>\frac{17}2$ there is $C_V\le2^{p_V}C_{p_V}g^2$, $C_{p_V}$ depending only on $g$ and $p_V$, the same for all limits at $g$, such that for $0<\eta\le1$, $F,F'\in\mathfrak F(\{x_3>\eta\})$ and $h$ supported in $\{|x_3|<\frac\eta2\}$:
\begin{equation}\label{sf-slab}
|\sigma(\theta\bar F\cdot(h)\cdot F')|\le C_V\eta^{-p_V}\|h\|_{L^1}\|[\tau_{-\eta e_3}F]\|\,\|[\tau_{-\eta e_3}F']\|.
\end{equation}
\item[(c)] $\sigma(RX)=\sigma(X)$ for $R\in O(3)$; $\sigma(\theta_{n,c}\bar F\cdot F)\ge0$ for the reflection $\theta_{n,c}$ in every plane $\{n\cdot x=c\}$, $n\in S^2$, acting on indices by $E\mapsto\Theta_nE\Theta_n$, and $F\in\mathfrak F(\{n\cdot x>c\})$.
\item[(d)] There is exactly one continuous function $s_n(x;e)$ at pairwise different points $x_1,\dots,x_n$, $n$-linear in $e\in V^n$, with $\sigma(X)=\int s_n(x;H_1(x_1),\dots,H_n(x_n))\,dx$ for all $X\in\mathfrak S(\mathbb R^3)$. It is real-analytic, symmetric, invariant under all isometries acting on points and indices; at points ordered in $x_3$ it is the boundary value of the vacuum expectations of the ordered products of the operators $\mathsf B$; and $|s_n(x;e)|\le(cC_V)^n\max\{1,2n^2/\delta(x)\}^{p_Vn}\prod_k|e_k|$, $\delta(x)=\min_{a\ne b}|x_a-x_b|$. On Schwartz functions $\varphi$ vanishing with all derivatives at coinciding points, $|\int s_n(x;\varphi(x))\,dx|\le C_E(g)^n(n!)^{18}|\varphi|_{9n}$, $|\cdot|_k$ the norms of Section~\YMMref{93}{713}.
\item[(e)] With $U(x)=e^{ix^0H}\mathcal U(\vec x)$, $x\in\mathbb R^{1+2}$, and $\mathcal D_\infty=\bigcap_j\mathrm{dom}\,H^j$: for $f\in\mathcal S(\mathbb R^{1+2};V_{\mathbb C})$, $\langle\psi,T(f)\psi'\rangle=\int\lim_{s\downarrow0}\langle U(x)^{-1}\psi,\mathsf B(s,s;f(x))U(x)^{-1}\psi'\rangle\,dx$ defines $T(f):\mathcal D_\infty\to\mathcal D_\infty$, with $\|(1+H)^jT(f)\psi\|\le c_j(f)\|(1+H)^{2l+j}\psi\|$ for $l>p_V$, $c_j$ a Schwartz norm of $f$, and $T(f)^*\supset T(\Theta_{e_3}\bar f\Theta_{e_3})$: a component with $r$ indices $3$ is Hermitian after the factor $i^r$. The span $\mathcal D_0$ of $\Omega$ and of the $T(f_1)\cdots T(f_k)\Omega$ is dense. There is a strongly continuous unitary representation $U$ of $\mathbb R^{1+2}\rtimes\mathcal L$, $\mathcal L$ the identity component of $O(1,2)$, with $U(x,1)=U(x)$, $U(a,\Lambda)\Omega=\Omega$, $U(a,\Lambda)\mathcal D_0=\mathcal D_0$ and $U(a,\Lambda)T(f)U(a,\Lambda)^{-1}=T(\hat\Lambda f(\Lambda^{-1}(\cdot-a))\hat\Lambda^t)$ on $\mathcal D_0$, where $\hat\Lambda=\iota\Lambda\iota^{-1}$, $\iota(t,\vec x)=(\vec x,-it)$; its spectrum lies in the closed forward cone; and $T(f)T(f')=T(f')T(f)$ on $\mathcal D_0$ if $f,f'$ have spacelike separated supports, all six components with all six. So $\mathrm{tr}\,F_{ij}F_{kl}$ is a tempered, local, Poincar\'e covariant tensor field with cyclic vacuum and Euclidean kernels $s_n$.
\item[(f)] There are a compact metrisable space $\mathsf X$, a Borel probability measure $\lambda$ with support $\mathsf X$ and real symmetric functionals $\sigma_\xi$, $\xi\in\mathsf X$, with $\xi\mapsto\sigma_\xi(X)$ continuous and $\xi\mapsto\sigma_\xi$ injective, such that $\sigma$ has iterated plateaux at every $(X^1,\dots,X^q)$, $q\ge1$, $X^i\in\mathfrak S(\mathbb R^3)$, in the sense of Definition~\ref{C-355-4}(5) with tuples of weights for tuples of functions, and
\[\lim_{|v_q|\to\infty}\cdots\lim_{|v_2|\to\infty}\sigma(X^1\cdot\tau_{v_2}X^2\cdots\tau_{v_q}X^q)=\int\prod_i\sigma_\xi(X^i)\,d\lambda(\xi):\]
with $G^{(1)}(v_2,\dots,v_q)=\sigma(X^1\cdot\tau_{v_2}X^2\cdots\tau_{v_q}X^q)$ at the $(v_2,\dots,v_q)$ at which the supports of $X^1,\tau_{v_2}X^2,\dots,\tau_{v_q}X^q$ are pairwise disjoint, the limits 
\[
G^{(p)}(v_{p+1},\dots,v_q)=\lim_{|v_p|\to\infty}G^{(p-1)}(v_p,v_{p+1},\dots,v_q),
\]
$v_p\in\mathbb R^3$, the others fixed, exist successively for $p=2,\dots,q$, and $G^{(q)}$, the iterated plateau, is the right member (at $q=1$: $\sigma(X^1)=\int\sigma_\xi(X^1)\,d\lambda$); and for every $\xi$: (i) $\sigma_\xi(X)=\lim_j\rho_j(A^{\nu_j}_{K_j}(X))$ along the full sequence of $\sigma$, $\rho_j$ the state $\nu_j$ conditioned on a translation-invariant event of positive mass, a periodic Gibbs state at $(K_j,\beta_j)$ (Proposition~\textup{\YMMref{106.1}{764}}), so that $(K_j,\beta_j,\rho_j)$ is a Gibbs sequence at $g$; (ii) $\sigma_\xi((H))=\frac{c_\xi}3\int\mathrm{tr}\,H$ with a real $c_\xi$, $\sigma_\xi$ has the bounds of Theorem~\ref{sf-gibbs}(a) with the same numbers, and the re-centred functional $\sigma^c_\xi(X)=\sum_S\prod_{a\notin S}(-\frac{c_\xi}3\int\mathrm{tr}\,H_a)\,\sigma_\xi(X_S)$ is the limit at $g$ of $(K_j,\beta_j,\rho_j)$ along the full sequence, so that Theorem~\ref{sf-gibbs}(c)--(f) and (a)--(e) hold for it, and $\mathcal H(\sigma^c_\xi)$ is unitarily $\mathcal H(\sigma_\xi)$; (iii) $\ker H_\xi=\mathbb C\Omega_\xi$, and 
\[
|\sigma_\xi(X\cdot\tau_vY)-\sigma_\xi(X)\sigma_\xi(Y)|\le C(X,Y)/|v|
\]
 for $|v|\ge R(X,Y)$, where $R(X,Y)$ keeps the two translated groups a fixed positive distance apart, with $C(X,Y)$ independent of $\xi$; (iv) there is a continuous map $\pi$ of $\mathsf X$ onto $\mathsf X_0$ with $\pi_*\lambda=\lambda_0$ and $\sigma_\xi(f_1\mathbf 1,\dots,f_n\mathbf 1)=\sigma_{0,\pi(\xi)}(f_1,\dots,f_n)$, so that $c_\xi=c_{\pi(\xi)}$, $c_\zeta$ the number with $\sigma_{0,\zeta}((f))=c_\zeta\int f$, and every component of the scalar decomposition is the trace part of a component $\sigma_\xi$.

The map $P_0[F]\mapsto(\xi\mapsto\sigma_\xi(F))$, $P_0$ the projection onto $\ker H$, extends to a unitary map of $\ker H$ onto $L^2(\mathsf X,\lambda)$, which carries $\ker H\cap\mathcal H_0$ onto $\{u\circ\pi:u\in L^2(\mathsf X_0,\lambda_0)\}$; so $\ker H=\mathbb C\Omega$ exactly when $\mathsf X$ is a point. Every re-centred component $\sigma^c_\xi$ satisfies all Wightman axioms, the uniqueness of the vacuum among them.
\item[(g)] With $\Theta^E_{ij}(f)=T(fs_{ij})$, $\Theta_{00}=-\Theta^E_{33}$, $\Theta_{0k}=-i\Theta^E_{3k}$ and $\Theta_{kk'}=\Theta^E_{kk'}$ for $k,k'\le2$: $\Theta_{0\nu}(\partial_0f)\psi=\sum_{k\le2}\Theta_{k\nu}(\partial_kf)\psi$ for $f\in\mathcal S(\mathbb R^{1+2})$, $\nu\le2$ and $\psi\in\mathcal D_0$.
\item[(h)] For real $z\ne0$ let $Z_z=P_A+P_E+zP_T$, $P_\bullet$ the projections of $V$ onto $V_A$, $V_E$, $V_T$, and $\sigma^z(X)=\sigma(Z_zX)$. Then $\sigma^z$ has the properties of Theorem~\ref{sf-gibbs}(c) with $\max\{1,|z|\}^n$ times the bounds, $[F]_z\mapsto[Z_zF]$ maps $\mathcal H(\sigma^z)$ unitarily onto $\mathcal H(\sigma)$ and intertwines $e^{-tH}$ and $\mathcal U$, so that (a), the circular cone included, and (b) hold for $\sigma^z$. If $z^2\ne1$, then $\sigma^z$ is not invariant under $O(3)$, it violates the law of Theorem~\ref{sf-gibbs}(d), and it is reflection positive exactly at the planes of the nine families. A real symmetric functional with the properties of Theorem~\ref{sf-gibbs}(c) and the square cone $\mathcal P(\{E<\max_k|q_k|\})=0$ that is reflection positive at one plane whose normal is not in the nine families is invariant under $O(3)$ and reflection positive at every plane.
\end{enumerate}
\end{theorem}

\begin{definition}[Overlapping-source input]\label{sf-overlap-input}
In the setting \textup{(T)}, consider the source derivatives of the step maps on their admitted analytic domains, with the weighted local-family norms and source sizes of the construction. For two bond forms $Y_a,Y_b$ whose weights overlap, assume that the derivative of order two created at depth $u$ decomposes into a local part $\mathfrak l$ of norm at most $cL^u$ and a remainder of norm at most $cK^{\mathsf p}L^{u/2}$, in the units of the two source sizes. The local part vanishes at the trivial background. At the real background $U_{u'}$ of a small block field of a coarser level, of depth $u'\le u$, assume the evaluated sum estimate
\[
 \left|\sum_{X,z}\mathfrak l(X,U_{u'},z)\right|
 \le K^{\mathsf p}\mathsf b^{-2}N L^{-2(u-u')},
\]
where the sum runs over every polymer and base point and $N$ is the number of sites of that coarser level. The family-norm bounds on the decomposition hold on the admitted analytic domains; the last estimate is required at these real backgrounds. The constants may depend on the number of derivatives, the fixed torus and a compact label interval, but not on the cutoff or the depths. The sources in the overlapping cases are the off-diagonal bond forms in clause \textup{(b)} below.
\end{definition}

The cancellation behind the evaluated sum estimate concerns the linear parts of the two insertions in the fluctuation field. If these parts have coefficient vectors $r_a,r_b$, their pairing through the fluctuation covariance $C_u$ satisfies
\[
 |\langle r_a,C_ur_b\rangle|
 \le \|C_u\|_{2\to2}\|r_a\|_2\|r_b\|_2.
\]
Thus a uniform covariance bound and the estimates
\[
 \|r_a\|_2^2\le CK^{\mathsf p}\mathsf b^{-2}N
 L^{-2(u-u')}\|\phi_a\|_\sharp^2
\]
give the required estimate, after absorbing fixed constants and the two source sizes. The proposed gain comes from writing the linear curvature response as a discrete curl and summing by parts: a difference then falls on the product of the weight and the background curvature. Over one lattice spacing at depth $u$, the local gauge bounds at depth $u'$ give the factor $L^{-(u-u')}$ in the $L^2$ norm, whose square gives the displayed gain. To complete this argument for the actual off-diagonal step insertions, one still needs the quantitative decomposition above and the identification and bounds of its linear coefficient vectors, including the remaining nonlinear terms. These are the additional input here; the local gauge estimate and the step-map Cauchy estimate alone do not supply that identification.

\begin{theorem}[Coinciding points, overlapping insertions, more points]\label{sf-coin}
\begin{enumerate}
\item[(a)] Assume \textup{(H)}, $\ell\in\mathfrak L(g)$, $(K,\beta)\in\mathfrak d(g)$, $0<d<1$ and $s_{\rm th}(n,g)\le s(d)\le K$, $s_{\rm th}$ a threshold depth depending on $n$ and $g$. For complex weights whose supports lie within the Euclidean distance $0.87d$ of points $x_a$ and are at pairwise Euclidean distances $\ge1.26d$, with no bound on their mutual distances,
\[|\langle O^T_K(H_1);\dots;O^T_K(H_n)\rangle_{K,\ell,\beta}|\le60^nC_n(g,\ell)\big[2\log\tfrac1d+C(g)\big]^{8n\mathfrak p_0}\mathsf b^{-n}\prod_a|H_a|:\]
the power counting $d^{-3}$ of an entry, the dimension of $\mathrm{tr}\,FF$, up to logarithms. On a side that is small over the canonical sequence of $(K,\beta)$ (equation~\YMMref{(2312)}{841}), for $g\le\gamma^{\rm sd}_W$ the least side $\ell_g$ (Lemma~\YMMref{126.3}{846}), the number $C_n(g,\ell)$ is free of $\ell$; on every other $\ell\in\mathfrak L(g)$ it carries the factor $\max\{1,\mathsf c(\ell)\}^n$, $\mathsf c(\ell)=2\ell^3e^{\mathsf B(\bar x_0)\ell^3}/(c'A_0^2)$, with the function $\mathsf B$ and the number $c'$ of equation~\YMMref{(2311)}{841}. Read at the reference coupling $\gamma^{\rm sd}_W$ after a dilation, the bound takes a scale-covariant form: let $\delta_\star=[L((\gamma^{\rm sd}_W)^{-2}+\bar B)]^{-1}$ and $\beta>X_{\gamma^{\rm sd}_W}$, let $z\in\mathbb Z$ be fixed by the first-passage window of $\beta$ at $\gamma^{\rm sd}_W$, so that $\delta_\star\mathsf b\le L^{-z}<(\gamma^{\rm sd}_W)^2\mathsf b$, and let $L^z\ell$ be the least side at $\gamma^{\rm sd}_W$; then for every $d<\delta_\star\mathsf b$ with $s_{\rm th}(n,\gamma^{\rm sd}_W)\le s(L^zd)\le K-z$, the depth and the threshold being formed at $\gamma^{\rm sd}_W$, the bound holds with $[2\log\frac{\mathsf b}d+C]^{8n\mathfrak p_0}$ in the place of $[2\log\frac1d+C(g)]^{8n\mathfrak p_0}$ and numbers depending only on $n$ and the basic constants, and this statement is invariant under $(K,\ell,\mathsf b,d)\mapsto(K-1,L\ell,L\mathsf b,Ld)$.

For supports only at pairwise Euclidean distances $\ge d$ the bound holds at $d''=\min\{\frac34d,L^{-s_{\rm th}(n,g)}\}$, for $K\ge s(d'')$, with $\sum_Q\sup_Q|H_a|$ in the place of $|H_a|$ and $\mathsf b^{-n}$ kept, $Q$ running through the boxes of a partition of $T_\ell$ into at most $(2\ell/d'')^3$ boxes of sides $\le\frac{99}{100}d''$: each box lies within the Euclidean distance $0.87d''$ of its centre and pieces of two weights are at Euclidean distances $\ge d\ge1.26d''$, so every product of the pieces $H_a\mathbf 1_Q$ has the geometry of the first display at $d''$. In the amalgam norm $d''^3\sum_Q\sup_Q|H_a|$ this is the bound with $\mathsf bd''^3$ in the place of $\mathsf b$: the power $d^{-3}$ of an entry.

At $n=2$ and for $g\le\gamma^{\rm sd}_W$ the off-diagonal components need no logarithm and no number of the side: for every $\ell\in\mathfrak L(g)$, $(K,\beta)\in\mathfrak d(g)$, $h>0$ with $1.31h<1$ and $s(1.31h)\le K$, and complex weights $\mathsf w_1,\mathsf w_2$ of the bonds with $x_b$ in closed cubes $Q_1,Q_2$ of side $\frac32h$ at a sup distance $\ge\frac52h$ wherever $\mathsf w_i(b)\ne0$, no bound being asked on the distance of the cubes, the bond forms $Y_i=a_K^{-1}\sum_b\mathsf w_i(b)q_b$ have $|\langle Y_1;Y_2\rangle_{K,\ell,\beta}|\le36\,C_\infty(g)\,\mathsf b^{-2}\sup|\mathsf w_1|\sup|\mathsf w_2|$, with the number $C_\infty(g)$ of Lemma~\ref{sf-far}, which depends on $g$ alone; so the off-diagonal part of $\mathcal S^{\ell,\mathsf b}_2$ has the power counting of Proposition~\ref{sf-ext} with $q=0$ and a number $A$ free of $\ell$.

The bound passes to $\mathcal S^{\ell,\mathsf b}_n$ for weights vanishing off closed sets in balls of radius $0.86d$ at pairwise Euclidean distances $\ge1.27d$, with a number depending on $\ell$. So $\mathcal S^{\ell,\mathsf b}_n$ has the power counting of Proposition~\ref{sf-ext} with $q=8n\mathfrak p_0$, a number $A=A_n(g,\ell)\mathsf b^{-n}$ and $h_0\in\ell2^{-\mathbb N}$ with $h_0\le\min\{\frac\ell{64n},\frac12L^{-s_{\rm th}(n,g)}\}$: it extends over all the diagonals to a distribution $\tilde{\mathcal S}_n$ of order $\le3n-2$ with the same power counting and one more logarithm. The lattice functionals $\tilde u_K$ formed from the cumulants at the cutoff $K$ with the subtractions of Proposition~\ref{sf-ext}, the sum running over the products of Whitney cubes of the sides $h$ with $s(1.52h)\le K$ alone, obey $|\tilde u_K(F)|\le C_EA\,\mathfrak l(\lambda)^{8n\mathfrak p_0+1}\|F\|_{(\lambda),3n-2}$, one number for all $K$, for $8h_K\le\lambda\le h_0$, $h_K\le2C_W(K)a_K$ the least side admitted at the cutoff $K$, and every $F$ supported in a product of closed cubes of side $2\lambda$, wherever their centres lie; and $\tilde u_K(F)\to\tilde{\mathcal S}_n(F)$ for every smooth $F$.

Within this power counting the extension is unique exactly up to the sums $\sum_{\pi,\alpha}\partial^\alpha\delta_\pi\otimes T_{\pi,\alpha}$ over the partitions $\pi$ into at most $n-1$ blocks and $|\alpha|\le3|\pi|$. At two points the real, exchange-symmetric, translation-invariant contact terms in $\mathfrak E_2(0)$ form a twelve-dimensional space under $W_3$ on $\mathbb R^3$ or $T_\ell$, and a six-dimensional space under $O(3)$ on $\mathbb R^3$, by Lemma~\ref{C-355-34}(f). $\mathcal S_2$ is no measure across the diagonal.
\item[(b)] The convergence rates in this clause assume the overlapping-source input of Definition~\ref{sf-overlap-input}. Under \textup{(T)} let $r\ge2$ tensor entries $O^T_K(\phi_aE_a)$, $E_a\in V$, and $n$ action densities $\mathcal O(f_i)$ be given, all with Lipschitz weights, and let $d>0$. Their joint cumulants $\kappa_K$ converge as $K\to\infty$ at a rate: $|\kappa_{K+1}-\kappa_K|\le CK^{\mathsf p}L^{-cK}\prod_a\|\phi_a\|_\sharp|E_a|\prod_i\|f_i\|_\sharp$ in the norms $\|\phi\|_\sharp=\max\{\sup|\phi|,40M\,\mathrm{Lip}\,\phi\}$, $|E|$ the largest modulus of the entries of $E$, for $K$ beyond a threshold, with $C,\mathsf p$ depending on $r$, $n$, $\ell$ and a compact interval of labels about $\mathsf b$ and the threshold depending besides on the geometry of the supports (on $d$ in the cases where $d$ is named), in the following cases:
\begin{enumerate}[label=\textup{(\roman*)}]
\item $c=\frac1{120r}$ if the tensor supports are pairwise $\ge d$ apart, nothing being asked of the supports of the $f_i$, which may meet one another and the tensor supports
\item $c=\frac1{120r}$ also if $\phi_a\phi_b=0$ for $a\ne b$, the tensor weights not overlapping (their supports may touch), provided that of any $r+n$ points, one from each support, the supports of the $f_i$ included, two are $\ge d$ apart, so that no point is common to all the supports (at $n=0$ the numbers $C,\mathsf p$ depend on $d$ as well)
\item $c=\frac1{360r}$ for off-diagonal $E_a$, $r\ge3$, if no three tensor supports share a point, nothing being asked of the $f_i$
\item $c=\frac1{864}$ for $r=2$, off-diagonal $E_a$ and $n\ge1$, the two tensor supports overlapping at will, if of any $n+2$ points, one from each support, two are $\ge d$ apart.
\end{enumerate}
\item[(c)] Assume $\mathrm{(Q_n)}$ of
Definition~\ref{hyp:quantitative-torus}, with its real plaquette
weights, first-creation bound, numerical prerequisites and
error $\mathsf e_n$, and suppose $36d<\ell$.
For three free lattice photon fields the cyclic formula in
that hypothesis is exact. For the interacting fields the
following are consequences of the hypothesis.
For the plane fields, with real Lipschitz $\phi_a$ in the open balls $B(x_a,d)$, these in a ball of radius $\le\frac\ell4$, $\mathrm{Lip}\,\phi_a\le|\phi_a|/(\vartheta_0d)$ and $|x'-y'|\ge d$ for $x'\in\mathrm{supp}\,\phi_a$, $y'\in\mathrm{supp}\,\phi_b$, $a\ne b$ (Euclidean distances, as for the $w_a$): $|\langle\mathcal O_{\lambda_1}(\phi_1);\dots;\mathcal O_{\lambda_n}(\phi_n)\rangle_{K,\ell,\beta}-3\mathsf b^{-n}K^{(n)}_{\underline\lambda}(\phi)|\le[\mathsf e_n+C_n(\frac{a_K}{\vartheta_0d}+\frac{d^3}{\ell^3})]\mathsf b^{-n}\prod_a|\phi_a|$, $C_n=\frac32n![6^nC_{\rm lat}(1+3C_{\rm lat})^{n-1}+43\cdot11^n]$; the same with $3^nC_n$ for real combinations $\sum_\lambda k_{a,\lambda}\mathcal O_\lambda(\phi_a)$, $k_a\in[-1,1]^3$, among them the diagonal entries $\mathcal O(\phi)-2\mathcal O_m(\phi)=O^T_K(\phi s_{mm})$ of the stress tensor.

For $2\le n\le7$ and $g\le\min\{\gamma_n,\gamma^{\rm sd}_W\}$, with $d$ below a threshold depending on $n$ and $g$ and the balls $B(x_a,d)$ in one ball of radius $9d$, the lattice inequality in Definition~\ref{hyp:quantitative-torus}, with $\mathcal K^{(n)}_K$ formed on $T_{\ell_g}$ and with $\varepsilon^{\rm rt}_n(g,d)$ of Lemma~\ref{sf-tr}\textup{(a)} added to $\mathsf e_n$, holds also in every periodic Gibbs state and on every torus of $\mathfrak L(g)$; $\varepsilon^{\rm rt}_n(g,d)$ tends to $0$ with $d$ like a positive power of $d$ times a power of $\log\frac1d$.
\item[(d)] Let $g$ be below a basic threshold. On the least torus, for $n=2,3,4$ and a label $\mathsf b\in I(g)$, the limit $\mathcal S^{\ell_g,\mathsf b}_n$ of Theorem~\ref{sf-tori}(a) has on each set of points named below a Borel kernel: $\mathcal S^{\ell_g,\mathsf b}_n(f_1\mathbf 1,\dots,f_n\mathbf 1)=\int s_n\prod_af_a$ for continuous $f_a$ whose supports form a product inside that set, $s_n$ the kernel of the truncated $n$-point function of the action density, and $\mathcal S^{\ell_g,\mathsf b}_2(f_1E_{\lambda\lambda},f_2E_{\mu\mu})=\iint\tau_{\lambda\mu}f_1f_2$ at $0<|x-y|<d_\sigma(g)$. For a limit $\sigma$ at $g$ the kernel of three action densities is $s_3(\cdot;\mathbf 1,\mathbf 1,\mathbf 1)$ of Theorem~\ref{sf-wight}(d), the third centred moment being the third cumulant, and $\tau_{\lambda\mu}=s_2(\cdot;E_{\lambda\lambda},E_{\mu\mu})$; for a component $\sigma_\xi$ they are the kernels of its re-centred $\sigma^c_\xi$. On $T_{\ell_g}$ almost everywhere, and for every limit at $g$ and every re-centred component at every such triple of points, the kernel $s_3$ of three action densities obeys
\[
\begin{multlined}
|s_3-3\mathsf b^{-3}k_3|\le\varepsilon_A(g,d)\,\mathsf b^{-3}r^{-3}d^{-6}\quad\text{for }0<|x-y|=r\le d=\mathrm{dist}(z,\{x,y\})\le d_A(g), \\
\varepsilon_A\le C(g)\Big[\big(\log\tfrac1d\big)^{39\mathfrak p_0+2}d^{1/2}+\Big(\frac d{\ell_g}\Big)^3\Big],
\end{multlined}
\]
uniformly in $\frac rd\in(0,1]$, $d_A(g)>0$ depending on $g$. At a cutoff the same law holds for narrow weights, in every periodic Gibbs state and on $T_{\ell_g}$: for real plaquette weights $w_1,w_2,w_3$ carried by sets $S_a\subset B(x_a,\rho)$ with $\mathrm{dist}(S_1,S_2)\ge r$, $|x_1-x_2|\le8r$, $\mathrm{dist}(S_a,S_3)\ge d$ and $|x_a-x_3|\le18d$ for $a=1,2$, where $0<\rho\le r\le d$, $d$ below a threshold depending on $g$, and $s(\rho)\le K$ (the geometry of Lemma~\ref{sf-jet}(c)), the third cumulant of the forms $a_K^{-1}\sum_pw_a(p)a_p$ is $3\mathsf b^{-3}K^{(3)}_K(w_1,w_2,w_3)$ up to $[C(g)Y^{39\mathfrak p_0+2}d^{1/2}+C(g)d^{3/2}]\rho^9r^{-3}d^{-6}\mathsf b^{-3}\prod_a|w_a|$, $Y=\log\bar x_{s(d)}$; for Lipschitz bumps $w_a=f_a\circ x$ with $\mathrm{Lip}\,f_a\le|f_a|/(\vartheta\rho)$, $0<\vartheta\le1$, the same holds against $3\mathsf b^{-3}\int k_3\prod_af_a$ with $C(g)(\frac{a_K}{\vartheta\rho}+(\frac d{\ell_g})^3)$ more in the bracket. For $d=|c-z|\le\frac12d_A(g)$ and $0<r\le\frac d2$, almost everywhere on the torus and at every $(c,z)$ where the kernels are continuous, in particular for limits and re-centred components,
\[
\begin{multlined}
\Big|\tfrac{2\pi}3\mathsf br^3s_3\big(c+\tfrac r2e_\lambda,c-\tfrac r2e_\lambda,z\big)-t_\lambda(c,z)\Big|\le\Big[\varepsilon'(g,d)+2.7\frac{r^2}{d^2}\Big]\mathsf b^{-2}d^{-6},\qquad\varepsilon'(g,d) \\
\le C(g)\big(\log\tfrac1d\big)^{39\mathfrak p_0+2}d^{1/2},
\end{multlined}
\]
$t_\lambda=\sum_\mu(\tau_{\lambda\mu}-\frac13\sum_{\lambda'}\tau_{\lambda'\mu})$, the kernel of the traceless diagonal component $E_{\lambda\lambda}-\frac13\mathbf 1$ against the action density, formed with the kernels $\tau_{\lambda\mu}$ of the same torus, limit or component: the tensor is recovered, to leading order, from two action densities that approach one another along $e_\lambda$, the right member being small against the natural size $\mathsf b^{-2}d^{-6}$ as $d\to0$ and $r/d\to0$. Four action densities at two scales are free in the same sense, on the torus and at a cutoff: on $T_{\ell_g}$, for $r\le|x_1-x_2|,|x_3-x_4|\le2r$ and $r\le d=\mathrm{dist}(\{x_1,x_2\},\{x_3,x_4\})\le d'_A(g)$, the kernel $s_4$ of the truncated four-point function of the action density has $|s_4-3\mathsf b^{-4}k_4|\le\varepsilon_4(g,d)\mathsf b^{-4}r^{-6}d^{-6}$ almost everywhere, $\varepsilon_4\le C(g)[(\log\frac1d)^{47\mathfrak p_0+2}d^{1/2}+(\frac d{\ell_g})^3]$, uniformly in $\frac rd\in(0,1]$; and at a cutoff, in every periodic Gibbs state and on $T_{\ell_g}$, for real plaquette weights $w_1,\dots,w_4$ in two such narrow pairs $\{1,2\}$, $\{3,4\}$, with $\mathrm{dist}(S_3,S_4)\ge r$, $|x_3-x_4|\le8r$, $\mathrm{dist}(S_a,S_b)\ge d$ and $|x_a-x_b|\le18d$ for $a\le2<b$, $d$ below a threshold depending on $g$, the fourth cumulant of the plaquette forms is $3\mathsf b^{-4}K^{(4)}_K(w_1,\dots,w_4)$ up to $[C(g)Y^{47\mathfrak p_0+2}d^{1/2}+C(g)d^{3/2}]\rho^{12}r^{-6}d^{-6}\mathsf b^{-4}\prod_a|w_a|$.
\end{enumerate}
\end{theorem}

\begin{corollary}[A tensor against the action density]\label{sf-ode}
Let $\sigma$ be a limit at $g$ under \textup{(H)}, or a re-centred component $\sigma^c_\xi$. There are real-analytic functions $A,B$ on $(0,\infty)$ with $s_2(x,y;e,\mathbf 1)=\mathrm{tr}\big(e\,[A(r)\mathbf 1+B(r)\hat u\hat u^t]\big)$, $u=x-y$, $r=|u|$, for all $e\in V$; $F=3A+B$ is the kernel of two action densities, and
\[B'(r)+\frac3rB(r)=\frac14F'(r),\qquad\text{so}\qquad B(r)=r^{-3}\Big(c+\frac14\int_1^rt^3F'(t)\,dt\Big)\]
for one real number $c$.
\end{corollary}

\subsection{Construction estimates}

The proof uses Theorem~\YMMref{D}{8}, the inductive description of the effective densities, on every torus, read through its calculus of sources under majorants (Setting), into which the insertions enter by Lemma~\ref{sf-src}; on one torus nothing else of the renormalisation enters. Theorem~\YMMref{D}{8} is read at two choices of the basic constants, the second with $\frac{93}{100}A_1$ in the place of $A_1$ and all the other constants that the deposit constant reads kept, $A_0$ among them (proof of Lemma~\ref{sf-depth}(b)). Theorems~\YMMref{AJ}{673} and~\YMMref{AD}{607} are used here through their estimates at source $0$, that is, through the pressure gap of Lemma~\ref{sf-gap}.

Besides, the proof uses the displays \YMMref{(1526)}{615}, \YMMref{(1527)}{615}, \YMMref{(1528)}{615}, \YMMref{(1528)}{615}, the canonical sequence stopped at a level (Corollary~\YMMref{74.3}{616} and Lemma~\ref{lem:setup-stopped-arithmetic}), and the following statements among the construction results: the sourced bare density (equation~\YMMref{(10)}{14}); the reflection positivity of Wilson's measure at bond mirrors (Corollary~\YMMref{85.3}{679}, Proposition~\YMMref{91.4}{700}) and of periodic Gibbs states at the lattice mirrors of the nine families $e_\mu$, $e_\mu\pm e_\lambda$ (Proposition~\textup{\YMMref{103.4}{758}}), the invariance of periodic Gibbs states under $a_K\mathbb Z^3\rtimes W_3$ (Proposition~\textup{\YMMref{103.4}{758}}) and their definition (Section~\ref{sec:state-limit-conventions}), and that a periodic Gibbs state stays one under a bounded density measurable for the translation-invariant events (Proposition~\textup{\YMMref{106.1}{764}}); the chessboard estimate (Proposition~\YMMref{86.3}{681}, Corollary~\ref{IV-18-18}), and the domination of periodic Gibbs states by chessboard seminorms along growing tori (Proposition~\YMMref{91.9}{705}); the comparison of the block fields of two successive cutoffs (Corollary~\textup{\ref{IV-62-66}}(a)), with a constant and a threshold that depend on the label only through an interval of labels under (Tor) (Definitions~\ref{sec:cutoff-calibration} and~\textup{\ref{IV-62-62}}(5)), formed at a compact interval of labels, with the backgrounds, the bound on their radius (equation~\YMMref{(2148)}{794}) and their minimality (the setting preceding Lemma~\YMMref{113.8}{794}, Proposition~\YMMref{113.9}{795}, Propositions~\YMMref{113.5}{792} and~\YMMref{113.7}{793}); the comparison of the lattice kernel with the free kernel (equations~\YMMref{(2245)--(2246)}{824}) and the decay of the curvature slices (Propositions~\YMMref{120.5}{821} and~\YMMref{120.6}{821}); the rarity of large-field regions on one torus (equation~\eqref{eq:ward-volume-constants}, Proposition~\ref{prop:ward-fixed-volume-rarity}), and the sides small over the canonical sequence (equation~\YMMref{(2312)}{841}, Lemma~\YMMref{126.3}{846}); integration by parts in the Haar measure along a field of directions; the scalar action-density theorem (Theorem~\YMMref{AP}{780}), its canonical decomposition (Theorem~\YMMref{AL}{746}), and the lattice clustering and conditional estimates of Propositions~\YMMref{99.7}{742} and~\YMMref{100.1}{744}; the lemma on a measure on a half plane seen along two directions at the angles $\pm\frac\pi4$, with the theorem on two mirrors at the angle $\frac\pi4$ for scalar functionals (Lemma~\YMMref{94.3}{722}, Proposition~\YMMref{94.4}{722}); and the reconstruction theorem for scalar functionals (Theorem~\textup{\YMMref{AO}{778}}).

\subsection{Estimates and proofs}

\begin{lemma}[The pressure gap]\label{sf-gap}
Let $0<g\le\gamma$ and $p_\ell=\ell^{-3}\log Z_{K,\ell,\beta}$. There is $G_g$ depending only on $g$ such that $0\le p_{\ell'}-p_\ell\le G_g$ for all $(K,\beta)\in\mathfrak d(g)$ and all $\ell'\le\ell$ in $\mathfrak L(g)$.
\end{lemma}
\begin{proof}[Proof sketch]
The lower bound, the monotony of the pressure in the side, is Corollary~\ref{IV-18-18}(e). For the upper bound, Theorems~\YMMref{AJ}{673} and~\YMMref{AD}{607} at the source $0$ factor the partition function as $Z_{K,\ell,\beta}=e^{E_{K,\ell,\beta}}\tilde Z$, where $E_{K,\ell,\beta}$ is the sum of the constants of the renormalisation, the canonical source energies vanish at the source $0$ (Lemma~\ref{lem:canonical-source-system}), and the normalised partition function obeys $e^{-C_g\ell^3}\le\tilde Z\le e^{C_S\ell^3}$; the energies per unit volume $E_{K,\ell,\beta}/\ell^3$ of two sides differ by at most a basic number (Lemma~\ref{lem:canonical-source-system}, \eqref{eq:source-system-pressure}). Taking logarithms and dividing by the volumes, $p_{\ell'}-p_\ell\le C_S+C_g+C$, which is the bound in \eqref{eq:source-system-pressure}.
\end{proof}

\begin{lemma}[Insertions are sources]\label{sf-src}
Let the data of the calculus of Theorem~\YMMref{D}{8} be given on $\Lambda_{K,\ell}$. A \emph{cell} is a positively oriented bond $b=\langle x,x+a_Ke_\mu\rangle$, with its base $x$, its four loops $(\mu,\alpha')$, $|\alpha'|\ne\mu$, and its four plaquettes, those that contain $b$; or a site $x$, with all the $24$ loops at $x$ and its twelve plaquettes, those with the corner $x$. A weight $v$ of the bonds, resp.\ of the sites, assigns a complex number $v(\mathsf c;l,l')$ to every cell $\mathsf c$ and pair of its loops ($16$ pairs at a bond, $576$ at a site); $\|v\|=\sup_{\mathsf c}\sum_{l,l'}|v(\mathsf c;l,l')|$, and $S[v]$ is the set of the plaquettes of the cells at which $v\ne0$.
\textup{(a)} For a complex weight $w$ of the plaquettes, $\mathcal O(w)=a_K^{-1}\sum_pw(p)a_p$ is the first member of a source under $(\sup|w|,\mathrm{supp}\,w)$.
\textup{(b)} For a weight $v$ of the bonds, $a_K^{-1}\sum_b\sum_{l,l'}v(b;l,l')Q_x(l,l')$, $x$ the base of $b$, is the first member of a source under $(59\|v\|,S[v])$; for a weight $v$ of the sites, the site form $a_K^{-1}\sum_x\sum_{l,l'}v(x;l,l')Q_x(l,l')$ is one under $(23\|v\|,S[v])$. In particular the bond form $a_K^{-1}\sum_b\mathsf w(b)q_b$ is one under $(59\sup|\mathsf w|,S[\mathsf w])$, $S[\mathsf w]$ the plaquettes that contain a bond $b$ with $\mathsf w(b)\ne0$; and a corner field $\mathcal Q_{l,l';\xi}(\phi)$ is one under $(23\sup|\phi|,S)$, $S$ the plaquettes with a corner $x$ at which $\phi(x+a_K\xi)\ne0$.
\textup{(c)} So $O^T_K(H)$ and $\bar O^T_K(H)$ are first members of sources under $(60|H|,S_H)$, $S_H$ the plaquettes of the plaquette part together with the plaquettes that contain a bond of the bond part at which the weight does not vanish. For real weights the sources are real, and the kernels of the second order of their first members have the row sums $2\sum_{p'}|\mathsf t(p,p')|\le\sup|w|$ in (a), $\le8\|v\|$ in (b), $\le2\sup|\mathsf w|$ for the bond form and $\le3|H|$ for $O^T_K(H)$.
\end{lemma}
\begin{proof}[Proof sketch]
(a) is the plaquette source of Theorem~\YMMref{D}{8}. (b) The chain of Theorem~\YMMref{D}{8} reads a source only through the clauses of its definition, and these follow from four properties of a corner form $Q=Q_x(l,l')$, $p,p'$ the plaquettes of $l,l'$: it is a polynomial in the entries of the values of the configuration at the bonds of $p$, $p'$ and at their reverses, and vanishes at the trivial field, since $W_{x,l}$ is the plaquette variable of $p$ or its inverse, conjugated by a product of at most one value at a bond of $p$ for a cell that is a bond, of at most two for a site; it is gauge invariant; at $SU(2)$-valued fields it is real with $|Q|\le a_p+a_{p'}$; and at a small point of the chart its variation is of the second order, $|Q(e^{iB'}U)-Q(U)|\le784\rho^2+112e^{2t/13}\rho\zeta_0$ for a $1$-form $B'$ with $\rho=\max|B'|\le\frac1{128}$, $\zeta_0$ a bound of the distances of the plaquette variables of $p$, $p'$ from $\mathbf 1$, $t=1$ at a bond and $t=2$ at a site.

So $|F(U)|\le a_K^{-1}\sum_{p\in S[v]}8\|v\|a_p(U)$, a plaquette lying in the cells of its four bonds, resp.\ of its four corners. The four, resp.\ twelve, plaquettes of a cell have their corners in a region of at most six, resp.\ seven, cubes of the level $j$ with a common corner, and at most $18M^3\eta_j^{-3}$, resp.\ $7M^3\eta_j^{-3}$, cells have one region of the level $j$, $\eta_j=L^{-j}$. The family at the level $0$ collects, region by region, the variations of the forms of its cells; by the bound of the second order its norm is at most $3138\cdot18$, resp.\ $3138\cdot7$, units of the calculus times $\|v\|$, while $\mathsf c^\flat\mathsf m_0$ is $972$ units for each unit of the size: $3138\cdot18\le972\cdot59$ and $3138\cdot7\le972\cdot23$.

At the levels $j\ge1$ the forms of the cells, read at the configurations of the level and collected by regions, are families of norm at most $4\cdot18\,e^{4/13}$, resp.\ $4\cdot7\,e^{8/13}$, times $M^3\alpha_0^2L^{K-j}\|v\|$, below $\frac12\mathsf m_j=6\cdot59$, resp.\ $6\cdot23$, times the same; and the sums over the regions give the form, cell by cell. The weight $v(b;(\mu,\epsilon\lambda_b),(\mu,\epsilon'\lambda'_b))=\frac14\epsilon\epsilon'\mathsf w(b)$, $0$ at the other twelve pairs, gives $\mathsf w(b)q_b$ with $\|v\|=\sup|\mathsf w|$; the weight $v(x;l,l')=\phi(x+a_K\xi)$ at the one pair gives $\mathcal Q_{l,l';\xi}(\phi)$. To the second order $Q_x(l,l')$ is $\pm$ the inner product of the angles of the plaquette variables of $p$ and $p'$, the signs telling whether the loops run along the boundaries of their plaquettes or against them; this gives the kernel of the second order, and its row sums are counted as $|F|$ was. (c) $O^T_K(H)$ is $\mathcal O(w_H)$ less the bond form of $\mathsf w_H(b)=H_{\lambda_b\lambda'_b}(x_b)$, resp.\ of the means of $H$ over the lattice cells $x+a_K[-\frac12,\frac12)^3$ of the Setting, with $\sup|w_H|,\sup|\mathsf w_H|\le|H|$, and sizes add: $1+59=60$.
\end{proof}

\begin{lemma}[Exact symmetries; separated insertions]\label{sf-sym}
\textup{(a)} Let $rx=Rx+v$, with $R\in W_3$, $Re_\mu=s_\mu e_{\pi(\mu)}$, and $v\in a_K\mathbb Z^3$, act on configurations by moving the bonds. Then $rO^T_K(H)=O^T_K(rH)$ for every complex weight, on $\Lambda_{K,\ell}$ and on $\mathbb L_K$; in components $ra_p=a_{rp}$ and $rq_b=s_{\lambda_b}s_{\lambda'_b}q_{rb}$. Under a law $\nu$ invariant under these maps, as the torus measures and the periodic Gibbs states are (Proposition~\textup{\YMMref{103.4}{758}}), $\nu(q_b)=0$ and $\nu(a_p)=e_\nu$ does not depend on $p$, so that $\nu(O^T_K(H))=e_\nu a_K^{-1}\sum_pH_{\lambda(p)\lambda(p)}(x(p))$: the off-diagonal components need no subtraction; moments and cumulants are invariant. $O^T_K(H)$ depends only on the bonds within $a_K$ of a point where $H\ne0$.

\textup{(b)} Assume \textup{(H)}, $\ell\in\mathfrak L(g)$, $q\ge2$, $\delta>0$. In the calculus of sources of Theorem~\YMMref{D}{8} let $C_{(m)}$, $m\ge1$, be the numbers, depending only on $m$, with which the mixed derivative families of order $m$ of a source system under majorants are bounded at the level $t$ by $C_{(m)}L^{4m(T-t)}$ times the product of the sizes of the sources, $T$ the working level; put $K_r=\sum_{\pi}\prod_{B\in\pi}C_{(|B|)}$, the sum over the partitions of $\{1,\dots,r\}$; and let $h_q(s)$ be the largest left member, at the depth $s$, of the smallness conditions that this calculus asks of the families of the orders $2,\dots,q$ for majorants at least $3M$ units of the working level apart, a function of the depth depending only on $q$ and $g$ and tending to $0$ as $s\to\infty$.

Let $s_q(\delta)$ be the least $s\ge0$ with $L^s\delta\ge3M$ and $h_q(s)\le1$; at $s=s_q(\delta)$ put $M_r=K_r[1+L^{4r}V_\ell(4q)L^{-(4q+1)s}]$, where $V_\ell(b)=4M^{-3}\ell^3e^{\mathsf B(X_g)\ell^3}\mathsf C_{\rm ag}(b+1,\eth)$, with the function $\mathsf B$ and the number $\mathsf C_{\rm ag}$ of equation~\eqref{eq:ward-volume-constants} at a fixed $\eth\in(0,1)$, so that $K_rL^{4r}V_\ell(b)L^{-(b+1)s}$, $b\ge4r$, bounds, in units of the sizes, the part of a centred mean of $r$ first members that is carried by the histories in which a large field is created below the working level (Proposition~\ref{prop:ward-fixed-volume-rarity}); and
\[K^{\rm ub}_q(g,\ell,\delta)=(4\ell^3L^{4s})^q\sum_\pi(|\pi|-1)!\prod_{B\in\pi}M_{|B|},\]
the sum over all partitions of $\{1,\dots,q\}$. Let $(K,\beta)\in\mathfrak d(g)$ with $K\ge s_q(\delta)+2$ and $3a_K\le\delta$, and let each $\mathcal F_a$ be one of: $O^T_K(H_a)$ or $\bar O^T_K(H_a)$ with $H_a$ vanishing off a closed set $\Sigma_a$ (size $\mathsf s_a=60|H_a|$); a site form with weights at the sites within $\frac12a_K$ of $\Sigma_a$ (size $23\|v_a\|$); $\mathcal O(w_a)$ with $w_a(p)=0$ unless $x(p)$ is within $a_K$ of $\Sigma_a$ (size $|w_a|$); or a sum of two of these (the sum of the sizes); the $\Sigma_a$ being pairwise $\ge2\delta$ apart. Then
\[|\langle\mathcal F_1;\dots;\mathcal F_q\rangle_{K,\ell,\beta}|\le K^{\rm ub}_q(g,\ell,\delta)\prod_a\mathsf s_a.\]
\end{lemma}
\begin{proof}[Proof sketch]
(a) $\mathrm{tr}\,U^{-1}=\mathrm{tr}\,U$ on $SU(2)$ gives $ra_p=a_{rp}$, and $w(U^{-1})=-w(U)$ gives the signs in $rq_b$; inserted into \eqref{sf-lat} they give $rO^T_K(H)=O^T_K(rH)$, as $(rH)_{\pi\lambda,\pi\lambda'}(rx)=s_\lambda s_{\lambda'}H_{\lambda\lambda'}(x)$. The reflection through $b$ normal to $\lambda_b$ lies in the group, fixes $b$ and reverses $q_b$; the group is transitive on the plaquettes. (b) By Lemma~\ref{sf-src} each $\mathcal F_a$ is the first member of a source, in the calculus of Theorem~\YMMref{D}{8}, under the majorant made of its size and of the plaquettes of its form (for $O^T_K(H)$ the size is $1+59$, for its plaquette and its bond part), and so, the notion being monotone in the majorant, under the majorant made of its size and of the plaquettes whose corners lie within $\frac32a_K$ of $\Sigma_a$; as $3a_K\le\delta$, corners of plaquettes of different majorants are $\ge\delta$ apart. At the depth $s=s_q(\delta)$, $L^s\delta$ exceeds the range of the block terms, and the mean of a centred product of $r$ first members is at most $M_r\prod_a4\ell^3L^{4s}\mathsf s_a$, up to constants common to all the subproducts, which cancel in cumulants; the M\"obius formula for cumulants gives $K^{\rm ub}_q$; the number grows with $\ell$ and as $\delta\to0$, since only the sizes of the sources enter.
\end{proof}

\begin{lemma}[Two successive cutoffs]\label{sf-cut}
Assume \textup{(T)}. An entry is one of $\mathcal O(w)$, $\mathcal O(w1_\lambda)$, 
\[
a_K^{-1}\sum_{b\parallel\mu}w(x_b)q_b
\]
 and $\mathcal Q_{l,l';\xi}(w)$, with a real function $w$ read at every cutoff at the barycentres of the plaquettes, at the midpoints of the bonds, resp.\ at the shifted sites.
\textup{(a)} \eqref{sf-rate} holds for $n$ entries with supports pairwise $\ge d$ apart, without the factor $6^n$ and with $|w_a|$ for $|H_a|$; and for entries whose weights are the cell means of bounded Borel functions, with $\prod_a|w_a|$ on the right.
\textup{(b)} Every compact interval $J\subset I(g)$ is an interval of labels under (Tor) on every side $\ell\in\mathfrak L(g)$ (Section~\ref{sec:cutoff-calibration}): its labels exceed $g^{-2}\ge\gamma^{-2}$, and $MR(x)\le L^{m_{\mathfrak L}(g)-1}$ for $x\le L(g^{-2}+\bar B)$ by \YMMref{(1528)}{615}. With $C_{\rm bf}$ and $K_{\rm bf}$ formed at $J$ (Definition~\textup{\ref{IV-62-62}}(5)), which serve the comparison of Corollary~\textup{\ref{IV-62-66}}(a) at every label of $J$, the numbers $\mathsf p$, $C^{\rm od}$ and $C^{\rm av}$ of Theorem~\ref{sf-tori}(a) serve all labels of $J$, and so does $K_1$, formed with $K_{\rm bf}$ of $J$ and at the least label of $J$.
\textup{(c)} Let $(t,\varepsilon_r,V)$ be data of the comparison of the cutoffs $K$ and $K+1$ (the setting preceding Lemma~\YMMref{113.8}{794}): a level $1\le t\le K$ of the cutoff $K$ with $L^{K-t}\mathsf b\ge\Theta_\star$, the radius 
\[
\varepsilon_r=2c_\rho\varepsilon(L^{K-t}\mathsf b)
\]
 of that Definition, and a block field $V$ of the depth $K-t$ with $\varpi(V)<\varepsilon_r$; let $U^A$ and $U^B$ be its backgrounds at the level $t$ of the cutoff $K$ and at the level $t+1$ of the cutoff $K+1$, and $\eta_t=L^{-t}$, $N_t=\ell^3L^{3(K-t)}$. The \emph{corner form} of a weight $v$ on the sites and of loops $l,l'$ at a configuration $U$ of a lattice is $\eta^{-1}\sum_xv(x)Q_x(l,l')(U)$, with $\eta=\eta_{t+1}$ on $\Lambda_{K+1,\ell}$ at $U^B$ and $\eta=\eta_t$ on $\Lambda_{K,\ell}$ at $\mathcal MU^B$ and at $U^A$; as $a_K^{-1}=L^{K-t}\eta_t^{-1}$ and 
\[
a_{K+1}^{-1}=L^{K-t}\eta_{t+1}^{-1},
\]
an entry is $L^{K-t}$ times a corner form, and $L^{K-t}N_t=\ell^3L^{4(K-t)}$ gives the factor $L^{4\sigma}$ below, $\sigma=K-t$. $B_3$ is the fixed background regularity constant of Theorem~\YMMref{E}{10}; its bound at these data is $B_3\varepsilon_r$, as used in the comparison preceding Lemma~\YMMref{113.8}{794}. For loops $l,l'$ and $l_1,l'_1$ in the same planes and complex $v$ on the sites,
\[\Big|\eta_{t+1}^{-1}\sum_xv(x)\big[s(l)s(l')Q_x(l,l')-s(l_1)s(l'_1)Q_x(l_1,l'_1)\big](U^B)\Big|\le2K^{\rm cf}_0\sup|v|\,\varepsilon_r^2N_tL^{-t/2},\]
$K^{\rm cf}_0$ basic: at a background the corner forms depend, in the mean, on their planes and signs alone.
\textup{(d)} Let $U^A$ be the background of the coarser lattice and $\mathcal MU^B$ the block average of $U^B$. For weights $v_f$, $v_c$ on the sites of the fine and of the coarse lattice, with $v_{\max}=\max\{\sup|v_f|,\sup|v_c|\}$ and the mismatch $m=\max_y|L^{-3}\sum_{x\in B(y)}v_f(x)-v_c(y)|$, $B(y)$ the block of $y$:
\[\Big|\eta_{t+1}^{-1}\sum_xv_f(x)Q_x(l,l')(U^B)-\eta_t^{-1}\sum_yv_c(y)Q_y(l,l')(\mathcal MU^B)\Big|\le\big[4B_3^2m+K^{\rm cf}_0v_{\max}L^{-t/2}\big]\varepsilon_r^2N_t;\]
and $\mathcal MU^B$, a nearly minimal competitor at the level $t$, has $|\eta_t^{-1}\sum_xv(x)[Q_x(l,l')(\mathcal MU^B)-Q_x(l,l')(U^A)]|\le12\sup|v|\,(C\,t\,A\,\Delta A)^{1/2}$, $A$ its action and $\Delta A$ its excess. So an entry with a Lipschitz weight, read at the two cutoffs, changes over one block field of depth $\sigma$ by at most $C\ell^3L^{4\sigma}K^{1/2}L^{-(K-\sigma)/2}(\sup|w|+L^{-\sigma}\mathrm{Lip}\,w)$, and an entry whose weights at the two cutoffs are tied by the block mean by the same with $\sup|w|$ alone.
\end{lemma}
\begin{proof}[Proof sketch]
(a) The entries are first members of sources under majorants (Lemma~\ref{sf-src}), of the sizes $\sup|w|$, $59\sup|w|$, resp.\ $23\sup|w|$, and the corners of the plaquettes of two majorants are $\ge d-3a_K\ge d'=\frac d2$ apart. \emph{First link, at one cutoff.} By Theorem~\YMMref{D}{8} a source system of $n$ members is carried by its first members read at the backgrounds of the block field of a depth $\sigma$ up to the error
\[2\mathfrak K'_n\big[2L^{4n}V_\ell(b)L^{-(b+1)s}+n\,\epsilon_{\rm fg}(\sigma-s)+2e^{3-L^sd'/M}\big](4\ell^3L^{4s})^n\prod_a\mathsf s_a\]
at a second depth $s<\sigma$, for every real $b\ge4n$, with $V_\ell$ and $M_r$ of Lemma~\ref{sf-sym}(b) formed at the depth $s$ and $\mathfrak K'_n=\sum_\pi|\pi|!\prod_{B\in\pi}M_{|B|}$: the first term is carried by the histories in which a large field is created below the working level of the depth $s$ (Proposition~\ref{prop:ward-fixed-volume-rarity}); the second is the forgetting of what the members of the sources make below the level $K-\sigma$, $\epsilon_{\rm fg}(\nu)=C_{\rm fg}(\nu+1)^{2\mathfrak p_0+2}L^{-\nu/2}$ with $C_{\rm fg}$ basic; the third is the decay over the distance $d'$ in the units of the working level; no number of it reads $K$ or $\beta$. \emph{Second link, between the cutoffs.} The laws of the block field of the depth $\sigma$ at the cutoffs $K$ and $K+1$ differ in total variation by at most $C_{\rm bf}\ell^3K^{\mathsf p_0}L^{-K/10}$, $\mathsf p_0=2\mathfrak p_0+\frac12$ the exponent of Definition~\textup{\ref{IV-62-62}}(5) (Corollary~\textup{\ref{IV-62-66}}(a)); over one block field an entry read at the two backgrounds changes by the defect of (d); and a first member read at a background is at most $\frac32B_3^2\ell^3L^{4\sigma}$ times its size.

With $\sigma=\lfloor K/(1260n)\rfloor$, $s=\lfloor\sigma/(8n+3)\rfloor$ and $b=4n+1$ the three errors are at most $L^2L^{-\theta_nK}$, $L^{-K/3}$ and $L^{-K/11}$ times the sizes, up to powers of $K$ and of $\ell$: the first as $\sigma-s\ge(8n+2)s$ makes each term of the bracket at most a power of $K$ times $L^{-(4n+1)s}$, and $L^{-s}\le L^2L^{-\theta_nK}$; the second as $L^{4n\sigma}L^{-(K-\sigma)/2}\le L^{-K/3}$; the third as $L^{4n\sigma}L^{-K/10}\le L^{-K/11}$: this is $\theta_n$. A Lipschitz weight read at the points of two lattices has the mismatch $m\le\frac32a_K\mathrm{Lip}\,w$, whence the term $L^{-\sigma}\mathrm{Lip}\,w$; cell means have $m=0$, as the cell of a coarse site is the union of the cells of the sites of its block, whence sup norms. (b) The first link reads the label only through the data of the depth $s$ of $(K,L^K\mathsf b)$, whose numbers are uniform on the compact $J$, the numbers of its error being free of the label.

Otherwise the label enters the proof of (a) only through $C_{\rm bf}$ and $K_{\rm bf}$, which depend on it only through $J$ (Corollary~\textup{\ref{IV-62-66}}(a)), through the condition $L^\sigma\mathsf b\ge\Theta_\star$ of the data of the comparison at the depth $\sigma$, which holds at every label of $J$ once it holds at the least one, and through the radius $\varepsilon_r$ of the backgrounds, which the calibration conditions in equation~\YMMref{(2148)}{794} bound at the levels where $L^\sigma\mathsf b\ge\Theta_\star$, and of which the defects of (c), (d) read only $\varepsilon_r\le1$. (c) Over a block, in a suitable gauge, each loop at the background is $\pm$ one plaquette variable of the block, depending only on the plane of the loop, up to $O(\varpi^2)$, $\varpi$ the size of the background; the variance of the plaquette variables of a block is bounded, plaquette by plaquette, by the defect of Jensen's inequality, whose sum is the action that one averaging step loses, and this is controlled by the weighted actions of the backgrounds (Proposition~\YMMref{113.9}{795}). (d) The first display is (c) carried across one averaging step, the block mean tying the weights; for the second, $\mathcal MU^B$ is nearly minimal at the level $t$ (Propositions~\YMMref{113.5}{792}, \YMMref{113.7}{793}), and a nearly minimal competitor has nearly the corner forms of the background, by its energy; this gives the factor $t^{1/2}$ and, with $t\le K$, the defect.
\end{proof}

\begin{lemma}[Two bond forms far apart]\label{sf-far}
Let $0<g\le\gamma^{\rm sd}_W$. There is a number $C_\infty(g)$ depending only on $g$ such that the following holds for every $\ell\in\mathfrak L(g)$, every $(K,\beta)\in\mathfrak d(g)$ and every $0<d<1$ with $s(d)\le K$. Let $\mathsf w_1,\mathsf w_2$ be real weights of the bonds, each carried by the bonds of one direction, and $x_1,x_2$ points with $|x_b-x_i|\le d-2a_K$ wherever $\mathsf w_i(b)\ne0$ and $|x_1-x_2|_\infty\ge3d+a_K$, no upper bound being asked on the distance of $x_1$ and $x_2$; let $Y_i=a_K^{-1}\sum_b\mathsf w_i(b)q_b$. Then
\[|\langle Y_1;Y_2\rangle_{K,\ell,\beta}|\le C_\infty(g)\,\mathsf b^{-2}\sup|\mathsf w_1|\sup|\mathsf w_2|.\]
For complex weights carried by the bonds of all directions the same holds with $36C_\infty(g)$, each weight being the sum of six real weights carried by the bonds of one direction.
\end{lemma}
\begin{proof}[Proof sketch]
The means vanish (Lemma~\ref{sf-sym}(a)), so the cumulant is the mean of $Y_1Y_2$; a lattice isometry carries a bond form of one direction to one, with the same sup norm of the weight. \emph{Far pairs.} There are a direction $\mu$ in which $x_1$ and $x_2$ are $\ge3d+a_K$ apart on the torus and a bond mirror $\{x_\mu=c\}$ between the two supports whose two planes $\{x_\mu=c\}$, $\{x_\mu=c+\frac\ell2\}$ are at distances $\ge\frac32d$ from $x_1$ and from $x_2$. Reflection positivity of Wilson's measure at bond mirrors (Corollary~\YMMref{85.3}{679}, Proposition~\YMMref{91.4}{700}) and Schwarz's inequality for the positive form give $|\langle Y_1Y_2\rangle|\le\langle\theta Y_1\cdot Y_1\rangle^{1/2}\langle\theta Y_2\cdot Y_2\rangle^{1/2}$.

So it suffices to bound a mirror pair $\langle\theta Y\cdot Y\rangle$ of a form carried within $d-2a_K$ of a point at the distance $t\in[\frac32d,\frac\ell2-\frac32d]$ from the mirror, a bound that reads neither the side nor the distance of the two supports. \emph{The distance from the mirror.} Let $F_r$ be the form moved by $r$ lattice steps $a_K$ away from the mirror and $\varphi(r)=\langle\theta F_r\cdot F_r\rangle$. As the translation by one lattice step normal to the mirror, away from it on the side of $F$, carries $\theta F_r\cdot F_r$ to $\theta F_{r-1}\cdot F_{r+1}$, the reflection turning it into the opposite translation, the invariance of the measure and Schwarz's inequality give $\varphi(r)^2\le\varphi(r-1)\varphi(r+1)$; so on an interval of $r$ the maximum of $\varphi$ is at an end, and the interval may be cut so that each end lies at a distance between $\frac32d$ and $2d$ from one of the two planes of the mirror.

\emph{A close mirror pair, $d<d_\flat:=L^{-s_8(g)}$.} Moved by a lattice translation into the board of side $\frac12\ell_g$, with the mirror a wall, the pair is $Y$ in one cell and $\theta Y$ in the next; the chessboard estimate and the pressure gap (Lemma~\ref{sf-gap}) give $\langle\theta Y\cdot Y\rangle\le e^{2G_g(\ell_g/2)^3}M_8^{1/4}$ on every torus of $\mathfrak L(g)$, $M_8$ the mean on $T_{\ell_g}$ of the eight mirror images of $Y$ in the cells of the board. They form four close pairs far from one another, and at the depth $s(d)\ge s_8(g)$, the calculus being read at the working level $\max\{K-s,2\}$, of the depth $\min\{s,K-2\}$, $s=s(d)$, so that every $K\ge s(d)$ is admitted, the cutoffs $K=s(d)$ and $K=s(d)+1$ among them, the conditional cumulants at the working level make their mean the product of the Wick values of the four pairs, up to an error that is small in the units of the sizes: blocks that link two pairs are small, as the pairs are far apart in the units of that level; a close pair has its Wick value; a single source is small; the levels with large fields are summed against their weights.

A Wick value of a close pair is at most a basic multiple of $\mathsf b^{-2}\sup|\mathsf w|^2$, the covariance of the background being bounded by a basic multiple of $(a_K/d)^3$ at the distances of the pair. So $M_8\le C_8\mathsf b^{-8}\sup|\mathsf w|^8$, with $C_8$ and $s_8(g)$ depending on $g$ alone. \emph{$d\ge d_\flat$.} The carriers of the eight images are $\ge d\ge d_\flat$ apart; their moments are bounded on the least torus by the bounds of at most eight separated sources at the scale $d_\flat$ (Theorem~\ref{sf-coin}(a), the bound for supports only at pairwise Euclidean distances $\ge d$, whose numbers read on $T_{\ell_g}$ the coupling alone) for $K\ge s(d''_\flat)$, $d''_\flat=\min\{\frac34d_\flat,L^{-s_{\rm th}(8,g)}\}$, and at the finitely many smaller cutoffs by the crude bound $|Y|\le c\,a_K^{-4}\sup|\mathsf w|\le c\,L^{4s(d''_\flat)}\sup|\mathsf w|$, $c$ basic, with $1\le\bar x_0^2\mathsf b^{-2}$; this gives a number depending on $g$ alone. $C_\infty(g)$ is the larger of the two numbers.
\end{proof}

\begin{lemma}[Dilation]\label{sf-dil}
For an integer $z\ge0$, with $K\ge z$ in the finite-lattice identity, and $H^{(z)}=H(L^{-z}\cdot)$: $O^T_K(H)=L^zO^T_{K-z}(H^{(z)})$ at corresponding configurations of $\Lambda_{K,\ell}$ and $\Lambda_{K-z,L^z\ell}$, at the same $\beta$, the label becoming $L^z\mathsf b$, and the two Wilson measures correspond; $X\mapsto X^{(z)}$ commutes with $O(3)$, carries $\tau_v$ to $\tau_{L^zv}$, the plane $\{n\cdot x=c\}$ to $\{n\cdot x=L^zc\}$ and $H^{\rm cons}[\eta]$ to $H^{\rm cons}[L^z\eta(L^{-z}\cdot)]$, and keeps $\mathcal W$ and the sup norms. \textup{(a)} For $\mathsf b_*>0$ and $\ell_*=2L^{m_*}$ there are $z$ and $g'\le\gamma^{\rm ref}_W$ with \textup{(T)} at $(g',L^z\mathsf b_*,L^z\ell_*)$; for large $z$ the coupling $g'=(L^z\mathsf b_*/2)^{-1/2}$ serves, and $L^z\mathsf b_*$ is an inner point of $I(g')$.

\textup{(b)} Let $0<g\le\gamma$. For $L^zg^{-2}$ large, $g'_1=(L^zg^{-2}-\bar B)^{-1/2}$ and $g'_2=(L^{z+1}g^{-2}-3\bar B)^{-1/2}$ are $\le\gamma^{\rm sd}_W$, every $(K,\beta)\in\mathfrak d(g)$ with $K>z$ has $(K-z,\beta)\in\mathfrak d(g'_1)\cup\mathfrak d(g'_2)$, and $L^z\mathfrak L(g)\subset\mathfrak L(g'_i)$; this arithmetic reads $g^{-2}>\bar B$ alone. Under \textup{(H)}, a limit at $g$ with label $\mathsf b$ is $\sigma(X)=L^{nz}\sigma_z(X^{(z)})$, $\sigma_z$ a limit at $(L^z\mathsf b/2)^{-1/2}\le\gamma^{\rm sd}_W$; $[X]\mapsto L^{nz}[X^{(z)}]$ is a unitary map of $\mathcal H(\sigma)$ onto $\mathcal H(\sigma_z)$ with $e^{-tH}\mapsto e^{-L^ztH_z}$; and if $\sigma_z$ has \eqref{sf-slab} with $(C_V,p_V)$, then $\sigma$ has it with $(L^{4z}2^{p_V}C_V,p_V)$.
\end{lemma}
\begin{proof}[Proof sketch]
$a_K^{-1}=L^za_{K-z}^{-1}$, while $a_p$ and $q_b$ do not depend on the spacing; $x\mapsto L^zx$ carries $\Lambda_{K,\ell}$ onto $\Lambda_{K-z,L^z\ell}$, bonds to bonds and plaquettes to plaquettes, and the Wilson weight at $\beta$ to the Wilson weight at $\beta$, and the periodic Gibbs states correspond, their definition depending on no spacing. $K_3$ is homogeneous of order $-3$, so $\mathcal W(H_1^{(z)},H_2^{(z)})=\mathcal W(H_1,H_2)$, and $\nabla(L^z\eta(L^{-z}\cdot))=(\nabla\eta)(L^{-z}\cdot)$. The rest is the arithmetic of first passage with \YMMref{(1526)}{615}, \YMMref{(1527)}{615}, \YMMref{(1528)}{615}, \YMMref{(1528)}{615}: with $\mathsf b'=L^z\mathsf b_*$, the coupling with $g'^{-2}=(\mathsf b'+2\bar B)/L-\bar B$ has $\mathsf b'$ as the upper end of $I(g')$, and $g'=(\mathsf b'/2)^{-1/2}$ makes $\mathsf b'$ an inner point; $m_{\mathfrak L}(g')$ grows like $\log z$, so $L^z\ell_*\in\mathfrak L(g')$ for large $z$; put $x=g^{-2}$ and $\Lambda=L^z$. Then $I(g'_1)=(\Lambda x,L\Lambda x-2\bar B]$ and $I(g'_2)=(L\Lambda x-2\bar B,L^2\Lambda x-(2L+2)\bar B]$. Their union contains $\Lambda(x,L(x+\bar B)]$ for $x>\bar B$, $L\ge13$ and $z\ge1$. Every first-passage label has $x<\mathsf b\le L(x+\bar B)$ by \YMMref{(1531)}{615}. Hence $\Lambda\mathsf b$ belongs to one of the intervals; since $\beta=L^{K-z}\Lambda\mathsf b$ and $K-z\ge1$, Section~\YMMref{74}{614} gives the first passage at $g'_1$ or $g'_2$. In \eqref{sf-slab} the dilation multiplies the distances by $L^z$ and the $L^1$ norm by $L^{3z}$, and the shift of the slab $\{|x_3|<\frac\eta2\}$ of \eqref{sf-slab} costs $2^{p_V}$.
\end{proof}

\begin{lemma}[Free values]\label{sf-freev}
\textup{(a)} $K_3(Ru)=RK_3(u)R^t$ for $R\in O(3)$, and $K_3$ is homogeneous of order $-3$; so $\mathcal W(rH_1,rH_2)=\mathcal W(H_1,H_2)$ for every isometry $r$, and $\mathcal W(H_1^{(z)},H_2^{(z)})=\mathcal W(H_1,H_2)$. $|K_3(u)_{\lambda\lambda'}|\le(2\pi|u|^3)^{-1}$, and for a pair at scale $d$ with $|x_0-y_0|\ge3d$, $|\mathcal W(H_1,H_2)|\le54|H_1||H_2|$.
\textup{(b)} In the setting of Theorem~\ref{sf-tori}\textup{(d)(ii)}, $\mathcal W(\theta(\phi E),\phi E)\ge\mathsf w_0\,\mathrm{tr}\,E^2$ for every $E\in V$.
\textup{(c)} $K_3(X)^2=(3\hat X\hat X^t+\mathbf 1)/(16\pi^2|X|^6)$. So the free kernel of two action densities is $\frac9{16\pi^2}\mathsf b^{-2}|X|^{-6}$, that of $\mathcal O_\lambda$ against $\mathcal O_\mu$ is $\frac32\mathsf b^{-2}[K_3(X)_{\lambda\mu}]^2$, and that of the traceless component along a unit vector $n$ against the action density is $t^\circ(n,X)=\frac9{32\pi^2}\mathsf b^{-2}|X|^{-6}\big((n\cdot\hat X)^2-\frac13\big)$.
\textup{(d)} For $0<r\le\frac d2$, $d=|c-z|$, and every unit vector $n$: $|\frac{2\pi}3\mathsf br^3\cdot3\mathsf b^{-3}k_3(c+\frac r2n,c-\frac r2n,z)-t^\circ(n,c-z)|\le2.66\,(r/d)^2\mathsf b^{-2}d^{-6}$.
\end{lemma}
\begin{proof}[Proof sketch]
(a) $\widehat{Ru}=R\hat u$. Each of the $81$ terms of $\mathrm{tr}(H_1K_3H_2K_3)$ is at most $|H_1||H_2|(4\pi^2|u|^6)^{-1}$, $|u|\ge d$ on the supports, and each support has volume at most $\frac{4\pi}3d^3$: $\frac32\cdot\frac{81}{4\pi^2}\cdot\frac{16\pi^2}9=54$. (b) By the translations and the symmetry of $\mathcal W$ one may take for $\theta$ the reflection in $\{x_3=0\}$ and $z=he_3$, $h=\frac32d$. With $\Theta=\mathrm{diag}(1,1,-1)$ and $M(\hat u)=3\hat u\hat u^t-\mathbf 1$, $\mathcal W(\theta(\phi E),\phi E)=\frac32\iint\phi(\theta x)\phi(y)\,\mathrm{tr}(\Theta E\Theta M(\hat u)EM(\hat u))\,(16\pi^2|u|^6)^{-1}dx\,dy$, $u=x-y$. At $\hat u=-e_3$, $\Theta M=\mathrm{diag}(-1,-1,-2)$ and the trace is $\sum_{\alpha\beta}c_\alpha c_\beta E_{\alpha\beta}^2\ge\mathrm{tr}\,E^2$, $c=(1,1,2)$; on the supports $|\hat u+e_3|\le2r/h$, $r=\frac d{64}$, and $\|M(\hat u)-M(-e_3)\|\le12r/h$, which changes the trace by at most $\frac12\mathrm{tr}\,E^2$; with $|u|\le\frac{97}{32}d$ and $\int\phi\ge\frac{4\pi}3(\frac r2)^3$ the integral is $\ge\frac34\cdot\frac19(\frac{32}{128\cdot97})^6\mathrm{tr}\,E^2=\mathsf w_0\,\mathrm{tr}\,E^2$. (c) $(3\hat X\hat X^t-\mathbf 1)^2=3\hat X\hat X^t+\mathbf 1$, and the free kernel of two entries with the labels $e$, $e'$ is $\frac32\mathsf b^{-2}\mathrm{tr}(eK_3(X)e'K_3(X))$. (d) $k_3(c+\frac r2n,c-\frac r2n,z)=\mathrm{tr}[K_3(rn)K_3(X-\frac r2n)K_3(X+\frac r2n)]$, $X=c-z$, $K_3(rn)=(3nn^t-\mathbf 1)/(4\pi r^3)$, and $K_3(X-\frac r2n)K_3(X+\frac r2n)=K_3(X)^2+[K_3(X),D]+O(r^2|X|^{-8})$ with $D=\frac r2(n\cdot\nabla)K_3(X)$; the commutator has trace zero against the symmetric $K_3(rn)$, the leading term gives $t^\circ(n,X)$ by (c), and the second-order terms are bounded by $\|(h\cdot\nabla)K_3(y)\|\le\frac6{4\pi}|h||y|^{-4}$ and $\|(h\cdot\nabla)^2K_3(y)\|\le\frac{24}{4\pi}|h|^2|y|^{-5}$.
\end{proof}

\begin{lemma}[The two-point law at a depth]\label{sf-depth}
\textup{(a)} In the setting of Theorem~\ref{sf-tori}\textup{(d)} at finite cutoff, with $(K,\beta)\in\mathfrak d(g)$, $\mathsf b=a_K\beta$, $s=s(d)\ge s_1(g,\ell)$, $K\ge s+2$ and $4a_K\le d-d_-$, let $F_i=O^T_K(H_i)$, with kernels $\mathfrak t_i$ as quadratic forms in the plaquette variables, let $C$ be the covariance of the background at the depth $s$, $v$ the coupling number of the working level, so that $v^2\le\frac98\mathsf b^{-2}$ and $|v^2-\mathsf b^{-2}|\le3\bar Bg^2L^{2-s}\mathsf b^{-2}$, and $\mathcal G_C=6v^2\sum\mathfrak t_1(p,p')\mathfrak t_2(q,q')C(p,q)C(p',q')$ the Wick value in three independent Gaussian fields. Then $|\langle F_1;F_2\rangle_{K,\ell,\beta}-\mathcal G_C|\le3600\,\varepsilon^e_{g,\ell}(s)v^2|H_1||H_2|$ with $Y_s^{2\mathfrak p_0}\varepsilon^e_{g,\ell}(s)$ bounded; and $\epsilon_p\epsilon_qC(p,q)=a_K^3[K_3(x(p)-x(q))_{\lambda(p)\lambda(q)}+\rho(p,q)]$, $\epsilon_p$ the orientations, with $|\rho|\le E_1d^{-3}$, $E_1=C_{\rm lat}[(\frac{a_K}d)^2+(\frac d\ell)^3]+C^{\rm fl}t^3e^{-cLt/\sqrt3}$, $C_{\rm lat},C^{\rm fl},c$ basic and $t=L^sd$; the term $(a_K/d)^2$ drops as $K\to\infty$.
\textup{(b)} The weight of the histories in which a large field is born at the level $t$ or at an earlier one is at most $C\,M^{-3}\ell^3L^{3(K-t)}\exp[-\frac34c_\star A_0^2(\log g_t^{-2})^{2\mathfrak p_0}]$, $1\le t\le K$, with $c_\star$ basic, $C$ formed with the basic constants and with those of the second choice of them that the proof fixes once, and $\frac34c_\star\ge\frac12c'$. So $s_1(g,\ell)\le\hat s(g,\ell)$, and for $s\ge\hat s(g,\ell)$ the number $\varepsilon^e_{g,\ell}(s)$ may be replaced by a function $\varepsilon^\circ_g(s)$ of $g$ and $s$ alone, with $Y_s^{2\mathfrak p_0}\varepsilon^\circ_g(s)$ bounded.
\textup{(c)} For the reflection $\theta$ in a bond mirror $\{x_\mu=c\}$, every $0<d_-<d<1$ with $s(d)\le K$ and $\frac32a_K\le d-d_-$, and real $H$ vanishing off $\bar B(x_0,d_-)$ with $\frac32d\le c-x_{0,\mu}\le\frac\ell2-\frac32d$: $0\le\langle\theta O^T_K(H);O^T_K(H)\rangle_{K,\ell,\beta}\le3600\,C^*_3(g,\ell)\mathsf b^{-2}|H|^2$, with $C^*_3(g,\ell)=\max\{(4/d'+2)^6C_\star(g),\,c_{\rm cr}L^{8s'}\bar x_0^2\}$, $d'=L^{-\hat s(g,\ell)-1}$, $s'=s(d')$, $C_\star(g)$ depending only on $g$ and $c_{\rm cr}$ basic, and $C^*_3(g,\ell)\le c_3(g)(1+\log\ell)^8e^{16\Lambda_\ell}$.
\end{lemma}
\begin{proof}[Proof sketch]
(a) The two insertions are sources of Theorem~\YMMref{D}{8} under majorants, carried to the working level of depth $s$. The plaquette part $a_K^{-1}\sum_pH_{\lambda(p)\lambda(p)}(x(p))a_p$ of $O^T_K(H)$ is a source under the majorant made of $|H|$ and of its plaquettes. Its bond part, the bond form $-a_K^{-1}\sum_b\mathsf w(b)q_b$ with $\mathsf w(b)=H_{\lambda_b\lambda'_b}(x_b)$, is a source under the majorant made of $59\sup|\mathsf w|$ and of the set $S[\mathsf w]$ of the plaquettes that contain a bond $b$ with $\mathsf w(b)\ne0$: it is a gauge-invariant polynomial in the bond variables, bounded by $a_K^{-1}\sum_{p\in S[\mathsf w]}2\sup|\mathsf w|\,a_p$, as $|q_b|\le\frac12\sum_{p\ni b}a_p$; at the first step of the renormalisation each term $q_b$ is a function of the at most six cubes of the level $0$ that hold a corner of one of the four plaquettes through $b$, at most $18M^3$ bonds having one such region, and it is there holomorphic in the fluctuation field, with a bound of the kind of a plaquette term; at every further level its members are families of that level bounded by half of the number that the size of the majorant gives at that level; the count of the bonds of one region against the bound of one plaquette term gives the number $59$.

A sum of sources is a source under the sum of the majorants, so that $O^T_K(H)$ is the first member of a source of the size $60|H|$ (Lemma~\ref{sf-sym}(b)); and to the second order at a small field this source is the quadratic form $\sum_{p,p'}\mathfrak t_H(p,p')\langle\Phi_p,\Phi_{p'}\rangle$ in the plaquette angles $\Phi_p$, with the diagonal kernel $\frac12H_{\lambda(p)\lambda(p)}(x(p))\delta_{pp'}$ for the plaquette part and, for the bond part, the form $-\sum_b\mathsf w(b)\langle u_{b,\lambda_b}(\Phi),u_{b,\lambda'_b}(\Phi)\rangle$, $u_{b,\lambda}(\Phi)$ half the signed sum of the angles of the two plaquettes through $b$ in the plane of $b$ and $e_\lambda$; for real $H$, $2\sum_{p'}|\mathfrak t_H(p,p')|\le3|H|$.

At the working level, on the small-field set, the conditional cumulant of the close pair is $\mathcal G_C$ up to $\varepsilon^e_{g,\ell}(s)$ times the sizes, since at a background a site form is a quadratic form in the fluctuation field, $Q_x(l,l')=\pm\langle\bar G_l,\bar G_{l'}\rangle$ to leading order with $\bar G_l$ depending only on the plane of $l$, and two quadratic forms in a Gaussian field have the Wick value as their cumulant. The contributions of the earlier levels, where large fields occur, are summed against the weights of the large-field sets, whose rarity (b) bounds the first admissible depth $s_1(g,\ell)$. The covariance $C$ is the free lattice covariance less the curvature slices of the levels above the working level; the comparison of the lattice kernel with the free kernel (equations~\YMMref{(2245)--(2246)}{824}), which gives the term $C_{\rm lat}[(a_K/d)^2+(d/\ell)^3]$, and the decay of the curvature slices (Propositions~\YMMref{120.5}{821} and~\YMMref{120.6}{821}), which gives the term $C^{\rm fl}t^3e^{-cLt/\sqrt3}$, give the form of $C$, $t$ being the distance in units of the working level. (b) A history with a large field at the level $t$ has a first creation at a level $t'\le t$, in one of the at most $CM^{-3}\ell^3L^{3(K-t')}$ cells of that level; the weight of the first creations at the level $t'$ is at most a sum, over the boundedly many test cubes of the step $t'$ in each of these cells (and, at $t'=1$, over the density cubes of the level $0$), of the means of test functions $F$: local gauge-invariant functions with $0\le F\le1$, the indicator that a test of the step fails at its cube, or, where the test reads a gauge variable of the step, the indicator that the field is farther from the background of the cube than $\frac{94}{100}$ of the threshold of the step, averaged over the gauge variables of the cube with the density of the gauge fixing of the step, an average that the sum over the histories undoes exactly.

The mean of one test function is bounded through the chessboard estimate, whose root removes the volume factor of a bound on the whole torus, and a second choice of the basic constants. First, $\langle F\rangle\le\|F\|_{\rm ch}$ on a board of $T_\ell$ centred at $F$ whose cells are $L^2$ times the cells of the level $t'$ (Proposition~\YMMref{86.3}{681}), and the $N$-th power of this seminorm is the mean of the $N$ mirror images of $F$, $N$ the number of the cells. Second, let the second choice have $\frac{93}{100}A_1$ in the place of $A_1$ and keep all the other constants that the deposit constant reads, $A_0$ among them, so that its deposit constant $c''_\star$ is at least $(\frac{93}{100})^2c_\star$ and its thresholds lie below $\frac{94}{100}$ of those of the first choice: a test function with the value $1$ then sits in the large-field region of the histories of the second choice, and the mean of the product of the $N$ images is at most the mass, at the second choice, of the histories that have all the $N$ centres in their large-field region; as one creation serves one cell, the cells of the board being larger than the reach of a creation, these histories have at least $N$ creations up to the level $t'$.

Third, Theorem~\YMMref{D}{8} at the second choice, with its large-field marks read at the dial $\frac9{10}$, and Markov's inequality bound this mass by $\exp[-\frac9{10}c''_\star p_0(g_{t'})^2N+C'(\log g_{t'}^{-2})\,n_{t'}N]$, $p_0(g)^2=A_0^2(\log g^{-2})^{2\mathfrak p_0}$, $n_{t'}$ the number of the sites of the level $t'$ in one cell of the board and $C'$ formed with the constants of the two choices; this price of a cell is at most $\frac1{40}c_\star p_0(g_{t'})^2$, by the floor $2\mathfrak p_0>3\mathsf r+2$ of the basic constants; and the $N$-th root gives $\langle F\rangle\le\exp[-\frac34c_\star p_0(g_{t'})^2]$, as $\frac9{10}(\frac{93}{100})^2-\frac1{40}\ge\frac34$. On the tori too small for such a board, Theorem~\YMMref{D}{8} at the first choice, its large-field marks read at the same dial, with one creation and Markov's inequality, gives the same bound directly, as $\frac9{10}-\frac1{40}\ge\frac34$: the side of such a torus is less than $2L^2$ cell sides of the level $t'$, the side of two cells of the board, and is twice a power of $L$ in these units, so that it is at most $2L$ of them; the torus then holds at most $8L^3\le L^6$ cells of that level, no more than one cell of the board holds, and the same fortieth pays their price; at the levels at which the run lies below the thresholds of these conditions, stopping the run at the last level above them and dilating gives the bound with the number $C$.

Summing over $t'\le t$, the couplings decreasing towards the cutoff, gives the bound. The two-point law needs the bound of (b) at the level $t=K-s$ of the depth $s$, $CM^{-3}\ell^3L^{3s}$ times the exponential, times $\bar x_s^2$, to be small; there $(\log g_t^{-2})^{2\mathfrak p_0}\ge(\frac12s\log L)^{2\mathfrak p_0}$, one half of the exponent pays the factor $\ell^3$ from the depth $\lceil2\Lambda_\ell/\log L\rceil$ on, and the other half pays $L^{3s}\bar x_s^2$ and a factor polynomial in $s$ from a depth $s_\flat(g)$ on; so from $\hat s(g,\ell)$ on the factor of the volume is absorbed into the exponential, which leaves an error function independent of $\ell$. (c) From the depth $\hat s(g,\ell)$ on, two entries of size $d$ whose centres are $\ge3d$ apart, with no upper limit on the distance, have a two-point function at most $C_\star(g)\mathsf b^{-2}$ times their sizes, for every $K\ge s(d)$ and independently of $\ell$.

For the argument of (a) and (b), read at the working level $\max\{K-s,2\}$, of the depth $\min\{s,K-2\}$, holds for every $K\ge s(d)$ and asks no upper limit on the distance of the two entries; the margin $\frac32a_K\le d-d_-$ puts the plaquettes of the source of $O^T_K(H)$, whose midpoints lie within $d_-+\frac12a_K$ of the centre, within $d-a_K$ of it, as the calculus asks; and it gives the Wick value up to $\varepsilon^\circ_g(s)v^2$ times the sizes. The comparison of the lattice kernel with the free kernel, the slices between the cutoff and the working level being of a range far below $d$, gives $|C(p,q)|\le C a_K^3d^{-3}$ at these distances, which, with $a_K^3\sum_{p,p'}|\mathfrak t_H(p,p')|\le45d^3|H|$ for an entry of size $d$, bounds the Wick value by a basic multiple of $v^2$ times the sizes.

If $s(d)\ge\hat s(g,\ell)$, this bounds the mirror pair, whose supports are $\ge d$ and whose centres are $\ge3d$ apart on the torus, by a number independent of $d$. Otherwise $d>3d'$, $d'=L^{-\hat s-1}$, $\hat s=\hat s(g,\ell)$, and two cases remain. If $K\ge s'$, $s'=s(d')\ge\hat s+2$, cut $H$ by a grid of mesh $\frac12d'$ into at most $(4/d'+2)^3$ pieces, each in a closed ball of radius $\frac{\sqrt3}4d'$, which has the margin at $d'$, as $a_K\le d'/C_W(s')$; a piece of $\theta H$ and a piece of $H$ lie $\ge d\ge3d'$ apart, and the bound above at $d'$ bounds their two-point function by $C_\star(g)\mathsf b^{-2}$ times their sizes, independently of $\ell$. Summing over the at most $(4/d'+2)^6$ pairs of pieces gives the first member.

If $s(d)\le K<s'$, the crude bound $|O^T_K(H)|\le c_1a_K^{-4}|H|\le c_1L^{4s'}|H|$, $c_1$ basic, from the number of the plaquettes and bonds in the ball and the factor $a_K^{-1}$, together with $1\le\bar x_0^2\mathsf b^{-2}$, gives the second member, with $c_{\rm cr}=c_1^2/1800$. Finally $s'\le\hat s+1+\lceil\log_LC_W(2\hat s+2)\rceil$, so $L^{8s'}\le L^{8\hat s+16}C_W(2\hat s+2)^8$; $L^{\hat s}\le L^{s_\flat(g)+1}e^{2\Lambda_\ell}$; and $C_W(2\hat s+2)^8\le C(g)(1+\Lambda_\ell)^{16\mathfrak p_0}\le C'(g)(1+\log\ell)^8$, as $\Lambda_\ell^{2\mathfrak p_0}$ is a multiple of $\log\ell$. So the second member is at most $C(g)(1+\log\ell)^8e^{16\Lambda_\ell}$, the first, $(4L^{\hat s+1}+2)^6C_\star(g)$, is at most $C(g)e^{12\Lambda_\ell}$, and $C^*_3(g,\ell)\le c_3(g)(1+\log\ell)^8e^{16\Lambda_\ell}$.

Positivity, for real $H$, for which $\theta O^T_K(H)=O^T_K(\theta H)$, is the reflection positivity of Wilson's measure at bond mirrors (Corollary~\YMMref{85.3}{679}, Proposition~\YMMref{91.4}{700}).
\end{proof}

\begin{lemma}[Chessboards]\label{sf-chess}
Assume \textup{(H)} and let $(K,\beta)\in\mathfrak d(g)$; the boards have the side $\mathsf l=L^{-i}$ with $i\le K-1$, and their walls are bond mirrors of $\mathbb L_K$.
\textup{(a)} For $\ell'\le\ell$ in $\mathfrak L(g)$, a board of side $\mathsf l\le\frac{\ell'}2$, functions $D_k$ of the bonds of distinct cells and $\nu$ a periodic Gibbs state, or Wilson's measure on $T_\ell$ with the cells distinct modulo $\ell$, in one period box: $|\nu(\prod_kD_k)|\le\prod_ke^{G_g\mathsf l^3}\|D_k\|_{{\rm ch},\ell'}$; no number depends on the positions of the cells.
\textup{(b)} For real weights in one cell, $\ge\delta$ from its walls and pairwise $\ge2\delta$ apart, $0<\delta\le\frac{\mathsf l}4$, with $c_a$ their means on $T_{\ell'}$ and $N=(\ell'/\mathsf l)^3$, and for $K\ge\max_{2\le q\le nN}s_q(\delta)+2$ and $3a_K\le\delta$, the depths $s_q(\delta)$ of Lemma~\ref{sf-sym}(b): $\|\prod_a(O^T_K(H_a)-c_a)\|_{{\rm ch},\ell'}\le C^\flat_n\prod_a|H_a|$, where $C^\flat_n=60^n(\mathrm{Mo}_{nN})^{1/N}$ and $\mathrm{Mo}_q=\sum_\pi\prod_{B\in\pi}K^{\rm ub}_{|B|}(g,\ell',\delta)$, the sum over the partitions without singletons. So in every periodic Gibbs state and on every torus of side $\ge\ell'$, $|\nu(\prod_a(O^T_K(H_a)-c_a))|\le e^{G_g\mathsf l^3}C^\flat_n\prod_a|H_a|$, and $|\nu(A^{\nu'}_K(X))|\le C^G_n\prod_a|H_a|$ with $C^G_n=\sum_{k=0}^n\binom nk(e^{G_g\mathsf l^3}C^\flat_1)^{n-k}e^{G_g\mathsf l^3}C^\flat_k$, $C^\flat_0=1$.
\textup{(c)} Let $h\in a_K\mathbb N$, with $\mathsf l/h$ an integer and $2nh<\mathsf l$. For $n$ sets of diameter $\le\frac{\mathsf l}{4n}$, some board of side $\mathsf l$ with its offset in $h\mathbb Z^3$ keeps each set at distance $\ge\frac{\mathsf l}{8n}$ from its walls. The offsets of the boards of $\mathbb L_K$ run over $a_K\mathbb Z^3$; at $h=a_K$ this asks $2na_K<\mathsf l$.
\textup{(d)} For a wall plane $\{x_3=c\}$, cells $\square_k$ of the slab $\{c\le x_3\le c+\mathsf l\}$, functions $D_k$ of the bonds of $\square_k$, a periodic Gibbs state $\nu$, and local functions $F,F'$ of the bonds in $\{x_3\ge c\}$:
\[\Big|\nu\Big(\overline{\theta_cF'}\prod_kD_k\,\tau_{\mathsf le_3}F\Big)\Big|\le\prod_ke^{G_g\mathsf l^3}\|D_k\|_{{\rm ch},\ell'}\;\nu(\overline{\theta_cF'}F')^{1/2}\,\nu(\overline{\theta_cF}F)^{1/2}.\]
\end{lemma}
\begin{proof}[Proof sketch]
(a) The chessboard estimate on $T_\ell$ (Proposition~\YMMref{86.3}{681}, Corollary~\ref{IV-18-18}(b)) bounds the mean by the product of the seminorms on $T_\ell$, and $\|D\|_{{\rm ch},\ell}\le e^{\mathsf l^3(p_{\ell'}-p_\ell)}\|D\|_{{\rm ch},\ell'}\le e^{G_g\mathsf l^3}\|D\|_{{\rm ch},\ell'}$ by Corollary~\ref{IV-18-18}(e) and Lemma~\ref{sf-gap}; a periodic Gibbs state is dominated by the seminorms along growing tori of $\mathfrak L(g)$ (Proposition~\YMMref{91.9}{705}). (b) $\|D\|^N_{{\rm ch},\ell'}$ is the mean on $T_{\ell'}$ of $nN$ centred insertions, the mirror images being insertions by Lemma~\ref{sf-sym}(a), pairwise $\ge2\delta$ apart, since a path from one cell to another crosses a wall; the mean of a product of centred variables is the sum over the partitions without singletons of products of cumulants, each bounded by Lemma~\ref{sf-sym}(b) at an order $q\le nN$, under its conditions $K\ge s_q(\delta)+2$ and $3a_K\le\delta$; the last also makes the product a function of the bonds of the cell, each $O^T_K(H_a)$ reading the bonds within $a_K$ of its support.

The rest is (a) and the expansion $\prod_a(E_a-\nu'(E_a))=\sum_J\prod_{a\notin J}(c_a-\nu'(E_a))\prod_{a\in J}(E_a-c_a)$. (c) In each direction the offsets that bring a wall within $\frac{\mathsf l}{8n}$ of one of the sets form, modulo $\mathsf l$, an open interval of length $\le\frac{\mathsf l}{4n}+\frac{\mathsf l}{4n}=\frac{\mathsf l}{2n}$, which holds fewer than $\frac{\mathsf l}{2nh}+1$ points of $h\mathbb Z$; the $n$ intervals hold fewer than $\frac{\mathsf l}{2h}+n$ of the $\frac{\mathsf l}h$ points of a period, and $\frac{\mathsf l}{2h}+n<\frac{\mathsf l}h$ as $2nh<\mathsf l$: a good offset exists in each direction. (d) Multiple reflections: Schwarz's inequality for the form $\nu(\overline{\theta_cA'}A)$, positive by the reflection positivity of $\nu$ at the bond mirror $\{x_3=c\}$ (Proposition~\textup{\YMMref{103.4}{758}}), is iterated over the walls of the board, which turns the middle factor into the products of the mirror images of the $D_k$ over growing tori; domination and (a) give the bound.
\end{proof}

\begin{lemma}[One torus serves all]\label{sf-one}
Let $0<g\le\gamma^{\rm sd}_W$. There are a depth $s_{16}(g)$, the threshold depth of the bounds for sixteen sources, $d_0(g)=L^{-s_{16}(g)}$, and functions $\omega'_g(s)$, $\delta_g(s)$, $\mu_g(s)$ tending to $0$ as $s\to\infty$, all depending only on $g$, with $Y_s^{2\mathfrak p_0}\omega'_g(s)$ bounded, $\delta_g(s)\le C(g)Y_s^{-\mathfrak p_0/4}$ and $\mu_g(s)\le C(g)Y_s^{-7\mathfrak p_0}$, such that the following holds. Let $(H_1,H_2)$ be a pair at scale $d<d_0(g)$ with $3d\le|x_0-y_0|\le4d$, $F_i=O^T_K(H_i)$, $(K,\beta)\in\mathfrak d(g)$ with $K\ge s(d)+2$ and $\frac32a_K\le d-d_-$, $s=s(d)$, $v$ the coupling number of the working level, $G=G_g(\ell_g/2)^3$, and $\mathfrak s$ a periodic Gibbs state or Wilson's measure on $T_\ell$, $\ell\in\mathfrak L(g)$. Then, the weights being read on $T_{\ell_g}$ as well,
\[
\begin{multlined}
|\mathfrak s(F_1;F_2)-\langle F_1;F_2\rangle_{K,\ell_g,\beta}|\le\big[\delta_g(s)+e^{2G}\mu_g(s)^2\big]3600v^2|H_1||H_2|,\qquad|\mathfrak s(F_i)-\langle F_i\rangle_{K,\ell_g,\beta}| \\
\le e^G\mu_g(s)\,60v|H_i|,
\end{multlined}
\]
and on $T_{\ell_g}$, $|\langle F_1;F_2\rangle_{K,\ell_g,\beta}-\mathcal G_C|\le\omega'_g(s)\,3600v^2|H_1||H_2|$ with the Wick value $\mathcal G_C$ of Lemma~\ref{sf-depth}.
\end{lemma}
\begin{proof}[Proof sketch]
Move the pair, by a lattice translation, to the middle of a cell of the board of eight cells of $T_{\ell_g}$; its mirror images are sixteen centred sources $\tilde P_\tau$, two in each cell. Their mean $\langle\prod_\tau\tilde P_\tau\rangle$ factorises, up to $\omega^\sharp_g(s)$ times the product of their sizes, by the conditional cumulants at the working level of depth $s$: blocks linking two pairs are small, as the pairs are far apart in units of that level; a close pair has its Gaussian value up to $\varepsilon_g(s)$, which also gives the last display; a single source obeys $|\Phi[F_a]-e_a|\le e_g(s)\hat m_a$ for its conditional mean $\Phi[F_a]$ given the field of the working level, $e_a$ its mean on $T_{\ell_g}$ and $\hat m_a=60v|H_a|$ its size, with no hypothesis on its mean, which admits the diagonal components; the levels with large fields are summed against their weights.

So $P=(F_1-c_1)(F_2-c_2)$, $c_i$ the means on $T_{\ell_g}$, has a small centred seminorm, its eighth power being the factorised mean, and Lemma~\ref{sf-chess}(a) carries it to every state and torus: $|\mathfrak s(P)-\langle P\rangle_{K,\ell_g,\beta}|\le e^G\|P-\langle P\rangle\|_{{\rm ch},\ell_g}$. One source alone gives the means: the mean of its eight centred mirror images is at most $e_g(s)^8$ times the product of their sizes, up to the blocks that link two copies and the tails, so that $\mu_g(s)\le e_g(s)+[\text{the linking blocks}+\text{the tails}]^{1/8}$; $e_g(s)\le C(g)Y_s^{-7\mathfrak p_0}$ up to terms that fall faster than every power of $s$, and the linking blocks and the tails are at most a power of $s$ times $L^{19s}e^{-cs^2}$, $c>0$ depending on $g$; so $\mu_g(s)\le C(g)Y_s^{-7\mathfrak p_0}$, $Y_s$ growing linearly in $s$, and $e^{2G}\mu_g(s)^2\le C(g)Y_s^{-14\mathfrak p_0}$.

The eighth root of $\omega^\sharp_g$, with $Y_s^{2\mathfrak p_0}\omega^\sharp_g(s)$ bounded, gives $\delta_g(s)\le C(g)Y_s^{-\mathfrak p_0/4}$.
\end{proof}

\begin{lemma}[A stopped coupling]\label{sf-stop}
Let $0<g\le\gamma^{\rm sd}_W$ and $0<\delta\le\frac14$. There are $i_E(g,\delta)$, thresholds $K_E(i)$ and a number $C_\delta$, depending only on $g$, $\delta$ and $i$, such that for $i\ge i_E(g,\delta)$ and $(K,\beta)\in\mathfrak d(g)$ with $K\ge K_E(i)$: after the dilation by $L^i$, $(K-i,\beta)$ is a first-passage datum at the running coupling $g_{(i)}$ of the scale $a_i$, with $L^{i-1}(g^{-2}-\bar B)-B_\sharp\le g_{(i)}^{-2}\le L^{i-1}(X_g+\bar B)-B_\sharp$ and the label $L^i\mathsf b$; $g_{(i)}<g$, so that \textup{(H)} holds at $g_{(i)}$; its least side $\ell''$, the least member of $\mathfrak L(g_{(i)})$, has $8\le\ell''^3\le CY_i^{2\mathfrak p_0-1}$, $Y_i=2(c_E+i\log L)$ with $c_E$ depending only on $g,\delta$; $L^i\ell\in\mathfrak L(g_{(i)})$, so that $L^i\ell\ge\ell''$, for every $\ell\in\mathfrak L(g)$; its pressure gap is $G_{g_{(i)}}\le G_g+\frac92i\log L+\frac18$; and for bounded $H$ vanishing within $\delta$ of the complement of a cell of side $1$, $c$ its mean on $T_{\ell''}$:
\[\|O^T_{K-i}(H)-c\|_{{\rm ch},\ell''}\le C_\delta Y_i^{q_E}L^i(L^i\mathsf b)^{-1}|H|.\]
One may take for $i_E(g,\delta)$ the least integer $i\ge1$ with $L^{i-1}(g^{-2}-\bar B)-B_\sharp\ge\mathsf x_E(\delta)$, $\mathsf x_E(\delta)\le\mathsf x_*/\delta$ the number of the proof, $\mathsf x_*$ basic, and $K_E(i)=i+s_\sharp(i)$, $s_\sharp(i)$ the depth of the distance $\delta$ formed at the coupling of inverse square $L^{i-1}(X_g+\bar B)-B_\sharp$, or of inverse square $L^{i-1}(X_g+2\bar B)-B_\sharp$, as in Theorem~\ref{sf-gibbs}\textup{(a)}, the depth not decreasing as the inverse square of the coupling grows; then $L^{i_E(g,\delta)}\le c_\sharp/\delta$, and one may take $C_\delta=\delta^{-4}C_1$, with $c_\sharp$ and $C_1$ depending only on $g$.
\end{lemma}
\begin{proof}[Proof sketch]
The first assertions are the arithmetic of first passage with \YMMref{(1526)}{615}, \YMMref{(1527)}{615}, \YMMref{(1528)}{615}, \YMMref{(1528)}{615}, as the canonical sequence stopped at a level states it (Corollary~\YMMref{74.3}{616} and Lemma~\ref{lem:setup-stopped-arithmetic}): the stopped datum is a first-passage datum at $g_{(i)}$ whose run is the cut of the run of $(K,\beta)$; for $i\ge1$ every coupling at which it is a first-passage datum lies below $g$, so that $g_{(i)}<g\le\gamma^{\rm sd}_W$; and $L^i\ell\in\mathfrak L(g_{(i)})$ for $\ell\in\mathfrak L(g)$: the least exponent of the sides at a coupling $g'$ reads $g'$ through $\log_LR(L(g'^{-2}+\bar B))$ alone, $R(x)$ the least power of $L$ that is $\ge(\log x)^{\mathsf r}$, whose argument grows by at most the factor $L^{i+1}$ from $g$ to $g_{(i)}$, as $g_{(i)}^{-2}+\bar B\le L^i(X_g+\bar B)=L^{i+1}(g^{-2}+\bar B)$, and $(\log L^{i+1}x)^{\mathsf r}\le(\log x)^{\mathsf r}L^{\mathsf r(i+1)/\log x}$, so that the exponent grows by at most $i$ at the arguments $\log x\ge2\mathsf r$ and the corollary gives the inclusion in general.

The bound on the pressure gap is that of Lemma~\ref{sf-gap} read at $g_{(i)}$, whose constant $C_{g_{(i)}}\le C_g+\frac92i\log L+\frac18$ grows by at most $\frac92\log L$ a step (Corollary~\YMMref{74.3}{616} and Lemma~\ref{lem:setup-stopped-arithmetic}). After the dilation the seminorm is formed on the least torus $T_{\ell''}$ at $g_{(i)}$, at cells of side $1$; let $H$ be real first, its real and imaginary parts costing a factor $2$. The points of the cell at a sup distance $\ge\delta$ from its walls form a cube of side $1-2\delta$, the union of at most $2\delta^{-3}$ boxes of sides $\le\frac{99}{100}\delta$, and $O^T_{K-i}(H)$ is the sum of the plaquette and the bond parts of these boxes, at most $4\delta^{-3}$ pieces; the seminorm is subadditive.

A piece with its $N=\ell''^3$ mirror images is a family of $N$ sources, the plaquette part under the majorant of size $\le|H|$, the bond part under that of size $\le59|H|$ (proof of Lemma~\ref{sf-depth}(a)), carried by balls of radius $<\frac78\delta$ that are $\ge\frac54\delta$ apart: the geometry of the calculus at the distance $\delta$, at the depth $s=s(\delta)$ formed at $g_{(i)}$. The $N$-th power of the seminorm of a centred piece is a moment of its $N$ centred images, and on the least torus the calculus of Theorem~\YMMref{D}{8} bounds the moments of $n\le\ell''^3$ sources: there are constants $c_a$ such that for every non-empty set $B$ of them $|\langle\prod_{a\in B}(F_a-c_a)\rangle|\le\prod_{a\in B}A_n(s)v\,\mathsf s_a$, $\mathsf s_a$ the sizes, with
\[A_n(s)=2^{2\mathfrak p_0+1}\,n\,\max\{C^I,C_2n^{3+\mathfrak p_0}\}\,(LC_W(s))^3\,\bar x_s\,Y_s^{2\mathfrak p_0},\]
$C^I$ and $C_2$ basic. Here the mixed derivative families of order $r$ have sizes $(c_*r)^r$, $c_*$ basic, whose root of order $r$, summed over the levels, is at most $C_2r^{3+\mathfrak p_0}$; the number of the partitions of $r$ sources is at most $n^r$; and the histories with a large field, summed by the level of their first creation against the sizes of the sources at that level, add at most as much as the histories without one, since on a small side the weights of the first creations at the level $t$ carry $e^{-\frac12c'p_0(g_t)^2}$ with $\frac12c'p_0(g_{K-i})^2\ge(\frac92\log x_{K-i}+2)\ell''^3$, $x_{K-i}$ the last member of the run, while $p_0(g_t)^2$ grows towards the cutoff fast enough to pay the growth of the sizes. This bound asks $2Y_s^{2\mathfrak p_0}\ge u_*(n)$, a number of the order of $n\log n$ formed with the basic constants, and $L^{3s}\le x_{K-i}^{20}$. One constant for all the images, their common mean $m$ on $T_{\ell''}$, costs a factor $2$: $|\langle\prod_\kappa(F_\kappa-m)\rangle|\le\sum_B(A_N(s)v\mathsf s)^N=(2A_N(s)v\mathsf s)^N$. So
\[\|O^T_{K-i}(H)-c\|_{{\rm ch},\ell''}\le2\delta^{-3}\cdot2A_N(s)v(1+59)|H|=240\,\delta^{-3}A_N(s)v|H|,\qquad v^2\le\tfrac98(L^i\mathsf b)^{-2}.\]
At $n=N\le CY_i^{2\mathfrak p_0-1}$ and the depth of $\delta$, the factor $N\max\{C^I,C_2N^{3+\mathfrak p_0}\}$ gives $Y_i^{(4+\mathfrak p_0)(2\mathfrak p_0-1)}$, $(LC_W(s))^3$ gives $Y_i^{6\mathfrak p_0}$, $\bar x_s\le C\delta^{-1}Y_i^{2\mathfrak p_0}L^i$ gives $Y_i^{2\mathfrak p_0}$, $L^i$ and $\delta^{-1}$, and $Y_s^{2\mathfrak p_0}\le Y_i^{2\mathfrak p_0}$: in all $Y_i^{q_E}$ with $q_E=(4+\mathfrak p_0)(2\mathfrak p_0-1)+10\mathfrak p_0$, and the factors $\delta^{-3}$ of the cut and $\delta^{-1}$ of $\bar x_s$ give $C_\delta=\delta^{-4}C_1$. The two conditions of the moment bound hold where $g_{(i)}^{-2}$ exceeds a number $\mathsf x_E(\delta)\le\mathsf x_*/\delta$, which fixes the threshold and gives $L^{i_E}\le c_\sharp/\delta$.
\end{proof}

\begin{lemma}[Plateaux and components]\label{sf-plat}
Assume \textup{(H)}; let $\sigma$ be a limit at $g$ along $(K_j,\beta_j,\nu_j)$, $A_j=A^{\nu_j}_{K_j}$, and $\mathscr I$ the $\sigma$-field of the translation-invariant events.
\textup{(a)} For $X,Y$ supported in $(-R,R)^3$, large $j$, $\rho\in\mathcal G_{K_j,\beta_j}$ and $v\in a_{K_j}\mathbb Z^3$ with $|v|\ge4\sqrt3R$: 
\[
|\rho(A_j(X\cdot\tau_vY))-\rho(E_\rho[A_j(X)|\mathscr I]E_\rho[A_j(Y)|\mathscr I])|\le C(X)C(Y)R/|v|,
\]
the numbers depending only on $g$ and the tuples.
\textup{(b)} The moments $\nu_j(\prod_iE_{\nu_j}[A_j(X^i)|\mathscr I])$ converge to the iterated plateaux of $\sigma$.
\textup{(c)} There are $(\mathsf X,\lambda,(\sigma_\xi))$ as in Theorem~\ref{sf-wight}\textup{(f)} with these plateaux $\int\prod_i\sigma_\xi(X^i)\,d\lambda$, and for every $\xi$ invariant events $E_j$ with $\nu_j(E_j)>0$ such that $\|E_{\rho_j}[A_j(X)|\mathscr I]-\sigma_\xi(X)\|_{L^\infty(\rho_j)}\to0$ for every $X$, $\rho_j=\nu_j(\cdot|E_j)$; each $\rho_j$ is a periodic Gibbs state at $(K_j,\beta_j)$ (Proposition~\textup{\YMMref{106.1}{764}}), invariant under $a_{K_j}\mathbb Z^3\rtimes W_3$ and reflection positive at the lattice mirrors of the nine families (Proposition~\textup{\YMMref{103.4}{758}}), and $(K_j,\beta_j,\rho_j)$ is a Gibbs sequence at $g$.
\textup{(d)} If $\rho_j\in\mathcal G_{K_j,\beta_j}$ are periodic Gibbs states, such as the conditioned states of \textup{(c)} (Proposition~\textup{\YMMref{106.1}{764}}), and $\sigma'(X)=\lim_j\rho_j(A_j(X))$ with $\|E_{\rho_j}[A_j(X)|\mathscr I]-\sigma'(X)\|_{L^\infty(\rho_j)}\to0$ for all $X$, then 
\[
|\sigma'(X\cdot\tau_vY)-\sigma'(X)\sigma'(Y)|\le C(X,Y)/|v|
\]
 for $|v|$ large, $C(X,Y)$ independent of $\sigma'$.
\end{lemma}
\begin{proof}[Proof sketch]
(a) Reflection positivity of $\rho$ at the lattice mirrors of the nine families (Proposition~\textup{\YMMref{103.4}{758}}) and the nested cubes give clustering on one lattice at the rate $|v|^{-1}$, by Proposition~\YMMref{99.7}{742} and the conditional-expectation argument in Proposition~\YMMref{100.1}{744}; applied to the conditional expectations on $\mathscr I$, with the bounds of Theorem~\ref{sf-gibbs}(a), which are uniform in $j$, $\rho$ and the positions, it gives the display, whose right member is the power $|v|^{-1}$ that positivity gives without a mass gap. (b) Iterate (a) and use the convergence of $\sigma$ at the translated tuples. (c) The $e_j(X)=E_{\nu_j}[A_j(X)|\mathscr I]$ are bounded uniformly in $j$; for a countable set of tuples, dense in the sup norms, the limits of their joint moments define, by the theorems of Stone--Weierstrass and Riesz, a probability measure on a countable product of compact intervals; $\mathsf X$ is its support, $\sigma_\xi$ the coordinates continued by continuity, and the events $E_j$ that put finitely many of the $e_j(X_k)$ near the coordinates of $\xi$ give the approximants. (d) Under $\rho_j$ the conditional expectations are nearly the constants $\sigma'(X)$, $\sigma'(Y)$, and (a) at $\rho=\rho_j$ gives the product at the rate $|v|^{-1}$.
\end{proof}

\begin{lemma}[Jets of narrow bumps]\label{sf-jet}
Let $g$ be below a basic threshold, $\ell\in\mathfrak L(g)$, $(K,\beta)\in\mathfrak d(g)$, and let $Y=\log\bar x_{s(d')}\le2\log\frac1{d'}+C(g)$.
\textup{(a)} For real plaquette weights $w_1$, $w_2$ carried by balls of radius $\rho$ about points at distance between $d'$ and $18d'$, the balls at the Euclidean distance $\ge d'$, $0<\rho\le d'\le d_4(g)$, a threshold depending on $g$, $64d'\le\ell$ and $s(\rho)\le K$: the cumulant of the two plaquette forms is the free lattice value $\frac32\mathsf b^{-2}\sum_{p,q}w_1(p)w_2(q)C^{\rm free}(p,q)^2$ up to $\varepsilon_2(g,d')(\rho/d')^6\mathsf b^{-2}|w_1||w_2|$, plus the part of the large fields, $\varepsilon_2(g,d')\le C(g)Y^{19\mathfrak p_0}d'^{1/2}$; on $T_\ell$ the large fields add at most $C(g,\ell)d'^{3/2}(\rho/d')^6\mathsf b^{-2}|w_1||w_2|$, $C(g,\ell)$ formed with the moments of the weights of the large-field histories on $T_\ell$ (Lemma~\ref{sf-depth}(b)), and on $T_{\ell_g}$ at most $C(g)d'^{3/2}(\rho/d')^6\mathsf b^{-2}|w_1||w_2|$. The same holds for mixed sources, against $\frac32\mathsf b^{-2}\mathrm{Tr}(\mathsf W_1C^{\rm free}\mathsf W_2C^{\rm free})$.
\textup{(b)} Let besides the two centres be at the Euclidean distance $\le8d'$. Let $\mathfrak b_n=A_1(\log g_n^{-2})^{\mathfrak p_0}$ and $\delta_n=g_n\mathfrak b_n$ at the running coupling $g_n$ of the level $n$, $A_1$ a basic constant. For a level $n$ from the level $k_{d'}=\max\{K-s(d'),2\}$ of the distance $d'$ to the last one, at which $\mathfrak b_n\le A_1Y^{\mathfrak p_0}$, the jet of the pair at the level $n$ is the mixed derivative $\mathsf F_n(U)$, in the two strengths at $0$, of the exponent of the effective density of the level $n$ in which the two plaquette forms are entered as sources, read at a field $U=e^{i\eta_nA}$, $\eta_n=L^{-n}$, $A$ a field of the level $n$ in its small-field class, of a radius of the order of $\delta_n$; its Gaussian part is the quadratic form $q_n(\Phi)=a_K^{-1}\sum_{p,q}w_1(p)\mathcal S_n(p,q)w_2(q)\langle\Phi_p,\Phi_q\rangle$ in the plaquette angles $\Phi_p$ of $U$, with $\mathcal S_n=\sum_{k_{d'}<j\le n}v_j\mathcal C_{(j)}$ the sum of the covariance slices $\mathcal C_{(j)}$ of these levels, each times the coupling number $v_j$ of its level; and $(d'/a_{K-n})^3\mathsf b^{-2}|w_1||w_2|$ is the natural size of this Gaussian part. Then
\[|\mathsf F_n(U)-\mathsf F_n(\mathbf 1)-q_n(\Phi)|\le\big[C\mathfrak b_n^2\delta_n+C'Y^{15\mathfrak p_0+2}\big(g\,d'^{1/2}+e^{-Y^{2\mathfrak p_0}/\sqrt3}\big)\big]\Big(\frac\rho{d'}\Big)^6\Big(\frac{d'}{a_{K-n}}\Big)^3\mathsf b^{-2}|w_1||w_2|,\]
with $C$, $C'$ basic; there is no gain as $\rho\to0$ at fixed $d'$.
\textup{(c)} Let $w_1,w_2,w_3$ be real plaquette weights carried by sets $S_a\subset B(x_a,\rho)$ with $\mathrm{dist}(S_1,S_2)\ge r$, $|x_1-x_2|\le8r$, and $\mathrm{dist}(S_a,S_3)\ge d'$, $|x_a-x_3|\le18d'$ for $a=1,2$, where $0<\rho\le r\le d'\le d_4(g)$, $64d'\le\ell$ and $s(\rho)\le K$: a narrow pair and a far bump. Then the third cumulant of the plaquette forms is $3\mathsf b^{-3}K^{(3)}_K$ up to $C(g)Y^{39\mathfrak p_0+2}d'^{1/2}\rho^9r^{-3}d'^{-6}\mathsf b^{-3}\prod_a|w_a|$. For four weights in two such pairs $\{1,2\}$, $\{3,4\}$, with $\mathrm{dist}(S_3,S_4)\ge r$, $|x_3-x_4|\le8r$ as well, $\mathrm{dist}(S_a,S_b)\ge d'$ and $|x_a-x_b|\le18d'$ for $a\le2<b$, and $d'$ below a threshold $d^{(4)}_4(g)\le d_4(g)$, the fourth cumulant is its free value up to $C(g)Y^{47\mathfrak p_0+2}d'^{1/2}\rho^{12}r^{-6}d'^{-6}\mathsf b^{-4}\prod_a|w_a|$. To both the parts of the large fields are to be added; neither bound depends on $\rho/r$ or $r/d'$.
\end{lemma}
\begin{proof}[Proof sketch]
A cumulant of plaquette forms is a mixed derivative at the strength $0$ of the logarithm of the partition function with the forms entered as sources; Theorem~\YMMref{D}{8} follows it through the levels as a sum of jets of the effective densities. Below the level of the distance the two sources meet only in polymers that reach across the distance, and these contribute at most $C(\rho/d')^6e^{-u/(2\sqrt3M)}$ times the natural sizes, $u$ the distance in units of the level. A single source of any width and any real weight has at every level its jet within $C\mathfrak b_n^2\delta_n$ of its Gaussian value. At the levels above the distance the kernel families of the pair are their Gaussian contraction up to errors in the running coupling and in the decay of the curvature slices (Propositions~\YMMref{120.5}{821} and~\YMMref{120.6}{821}), independently of the width; so each level adds the free size times the coupling at that level, which sums to (b); there the geometry $|x_1-x_2|\le8d'$ is read: every corner of a plaquette of the two supports lies within the sup distance $11d'$ of a bare site next to $x_1$, which bounds, in the units of each level, the regions of the kernel families of the pair; summed over the levels, the Gaussian parts give the free lattice value, and the large-field histories are summed against their weights, by Lemma~\ref{sf-depth}(b).

At three and four points the crossing families, which meet both clusters, are bounded at the level of the far distance by the decay of the kernels (Proposition~\YMMref{120.6}{821}), with sizes $(\rho/d')^6$ times the natural sizes of the clusters; a narrow pair at its own level is (b).
\end{proof}

\begin{lemma}[Transfer to every state]\label{sf-tr}
Let $0<g\le\gamma^{\rm sd}_W$, $(K,\beta)\in\mathfrak d(g)$, $\mathsf b=a_K\beta$, $\bar{\mathsf b}=L(g^{-2}+\bar B)$, $\mathsf Q_0(g)=e^{G_g\ell_g^3/8}$ with $G_g$ of Lemma~\ref{sf-gap}, and let $\mathfrak s$ be a periodic Gibbs state in $\mathcal G_{K,\beta}$ or Wilson's measure on a torus of $\mathfrak L(g)$ at $\beta$. For real weights $w_1,\dots,w_n$ of the plaquettes of $\mathbb L_K$ vanishing off bounded sets let $\mathsf S_n[\mathfrak s](w)$ be the joint cumulant in $\mathfrak s$ of the plaquette forms $a_K^{-1}\sum_pw_a(p)a_p$, and $\mathsf S_n[\ell_g](w)$ the same on $T_{\ell_g}$ at $\beta$, the weights being carried to a torus by the projection, which is one to one on the balls below; the support of $w_a$ is the closure of $\{x(p):w_a(p)\ne0\}$, and $|w_a|=\sup|w_a|$.
\begin{enumerate}
\item[(a)] \emph{One scale.} Let $2\le n\le7$ and $0<\sigma<1$. There are $d_2\in(0,1)$ and numbers $C,c$, depending only on $n,g,\sigma$, such that for $0<d\le d_2$ with $s(d)\le K$ and weights whose supports lie in closed balls of radius $d$ contained in one ball of radius $9d$ and are at pairwise Euclidean distances $\ge d$,
\[|\mathsf S_n[\mathfrak s](w)-\mathsf S_n[\ell_g](w)|\le\varepsilon^{\rm rt}_n(g,d)\,\mathsf b^{-n}\prod_a|w_a|,\qquad\varepsilon^{\rm rt}_n(g,d)\le C\Big(\log\frac1d\Big)^cd^{\,4-\frac n2-3\sigma}.\]
The exponent $4-\frac n2-3\sigma=3(1-\sigma)-(\frac n2-1)$ is positive for $\sigma<\frac13(4-\frac n2)$, and some $\sigma\in(0,1)$ makes it positive exactly when $n\le7$.
\item[(b)] \emph{Two scales.} Let besides $g$ be below a basic threshold, $n\in\{2,3,4\}$, $0<r\le d_\times\le1$, $r$ below a threshold depending on $g$, $s(r)\le K$, and let the supports of the weights lie in balls $B(x_a,r)$ contained in one ball of radius $\le\frac18\ell_g-1$, at the Euclidean distance $\ge d_\times$ at $n=2$, and at $n=3,4$ forming the clusters $\{1,2\}$, $\{3\}$, resp.\ $\{1,2\}$, $\{3,4\}$, the supports being at Euclidean distances $\ge r$ within a cluster and $\ge d_\times$ across. Then
\[|\mathsf S_n[\mathfrak s](w)-\mathsf S_n[\ell_g](w)|\le C^{\rm tr}_n(g)\Big(\frac r{d_\times}\Big)^6d_\times^{3/2}\prod_a|w_a|,\]
with $C^{\rm tr}_n(g)$ depending only on $n$ and $g$.
\item[(c)] \emph{Three lengths.} In (b) at $n\in\{3,4\}$ let the supports lie in balls $B(x_a,\rho)$, $0<\rho\le r$, $\rho$ below a threshold depending on $g$, $s(\rho)\le K$. Then the difference of (b) is at most $C(g)\,d_\times^{3/2}\rho^9r^{-3}d_\times^{-6}\prod_a|w_a|$ at $n=3$ and $C(g)\,d_\times^{3/2}\rho^{12}r^{-6}d_\times^{-6}\prod_a|w_a|$ at $n=4$: $d_\times^{3/2}$ times the natural size of the cumulant of the narrow clusters, $\mathsf b$ being bounded by $\bar{\mathsf b}$.
\end{enumerate}
\end{lemma}
\begin{proof}[Proof sketch]
In each case move the tuple, by a lattice translation, into the middle of one cell of the board of eight cells of $T_{\ell_g}$, of side $\frac12\ell_g$; it stays $\ge\frac18\ell_g$ from the walls. Its eight mirror images are $8n$ sources on $T_{\ell_g}$ (Lemma~\ref{sf-sym}(a)). By Lemma~\ref{sf-chess}(a) and the pressure gap, the mean of a product $P$ of centred sources of one cell in $\mathfrak s$ differs from that on $T_{\ell_g}$ by at most $\mathsf Q_0(g)\|P-\langle P\rangle_{K,\ell_g,\beta}\|_{{\rm ch},\ell_g}$, and the eighth power of this seminorm is the mean on $T_{\ell_g}$ of the product of the eight mirror images of $P-\langle P\rangle_{K,\ell_g,\beta}$; the expansion of cumulants into centred moments passes from moments to cumulants.

Theorem~\YMMref{D}{8} writes the mean of the mirror sources as a mean, over the histories, of the exponential of their joint jets, the histories with large fields being summed against their weights, whose moments are bounded on the least torus (Lemma~\ref{sf-depth}(b)). (a) The eight copies are compared at the level of spacing about $d^\sigma$, where each copy fills a part $d^{3(1-\sigma)}$ of a block of that level; each derivative in a source is paid by Cauchy's estimate on a polydisc whose height is the square root of the coupling number at the level of the supports, about $(\mathsf b/d)^{1/2}$, and one factor of that number comes back, which costs $(\frac43\bar{\mathsf b}/d)^{\frac n2-1}$; the eighth root, $\mathsf Q_0(g)$ and the sum over the partitions give $\varepsilon^{\rm rt}_n(g,d)\le C(\log\frac1d)^cd^{3(1-\sigma)-(\frac n2-1)}$, the logarithms coming from the sizes of the sources at the depth $s(d)$. (b), (c) For narrow clusters the joint jets of the $8n$ mirror sources, less constants carried by single cells, are bounded almost everywhere by a number depending on $n$ and $g$ times $\alpha^{-6}$, in units fixed by the sizes of the clusters, where the large-field variable of a history is $\alpha=a_{K-t'}$, $t'$ the last level before its first large field ($t'=t-1$ if the first large field is born at the level $t$), and $\alpha=1$ on the histories with no large field: the jets within one copy as in the proof of Lemma~\ref{sf-jet}, and the jets that link two copies factor into the same sizes; with the moments of $\alpha^{-48n}$ on the least torus, $\int\alpha^{-b}\le1+L^bV_{\ell_g}(b)$ for $b\ge0$, $V_\ell$ the number of Lemma~\ref{sf-sym}(b), which read neither $K$ nor $\beta$, by the rarity of the large fields on one torus (Proposition~\ref{prop:ward-fixed-volume-rarity}), the eighth root of the mean of the mirror images and $\mathsf Q_0(g)$ give the comparison with a constant $\eta(1+\eta)^{n-1}$, $\eta$ depending only on $n$ and $g$, times the product of the units, which is the right member of (b), resp.\ of (c), up to a number depending on $n$ and $g$.
\end{proof}

\begin{proposition}[The stress insertion]\label{sf-stress}
Let $\eta\in C^2_c(\mathbb R^3;\mathbb R^3)$ and $\mathsf b=a_K\beta$.
\textup{(a)} Let $\Theta_\beta[\eta]$ be the bilinear form in the loops at the bonds with the weights $\pm\frac12\beta(\eta^\rho(x_b)-\eta^\rho(x(p)))$, $p\ni b$, and $\mathrm{St}_K[\eta]:=\Theta_\beta[\eta]-\mathsf b\,\mathcal O(\nabla^{\rm s}\cdot\eta)$, $\nabla^{\rm s}\cdot$ a lattice divergence: the derivative $\beta\mathfrak v_\eta S_W$ of the action $S_W=\sum_pa_p$ along the field of directions $\mathfrak v_\eta$ that moves the bonds by $\eta$, less the orbital insertion. Identically in the configuration, $\mathrm{St}_K[\eta]=\mathcal N_\eta-\frac14\mathcal N'_\eta-\mathrm{Rm}_K[\eta]-c_\eta$, with a constant $c_\eta$, $\mathcal N_\eta=\beta\mathfrak v_\eta S_W-\mathrm{div}\,\mathfrak v_\eta$ and $\mathcal N'_\eta$ of the same form for another field of directions, both carried by the bonds within $a_K$ of $\mathrm{supp}\,\eta$, and $\mathrm{Rm}_K[\eta]=\mathcal S[\mathsf b\eta]-J_\eta+\frac32a_K\mathcal O(w_\eta)$, where $J_\eta=\sum_{v=1}^5J_v[\mathsf k_\eta]$ is a sum of five local forms in the forward differences $\mathsf k^\rho_\eta=\nabla^+_\rho\eta^\rho$, a local form of the order $m$ being a sum $\sum_p\varkappa(p)j(p)$, with weights $\varkappa(p)$, of a gauge-invariant density $j(p)$ of the bonds next to $p$ that vanishes to the order $m$ at the trivial configuration: three of them of the order $2$ in the plaquette variables, with weights at most $(3+5\mathsf b)\|\eta\|_{C^2}$, one of the order $3$, and one of the order $4$, with the density $a_p^2$, these two with weights at most $L^K\mathsf b\sup|\nabla\eta|$; $w_\eta$ is a lattice divergence of $\eta$, with $|w_\eta|\le2\sup|\nabla\eta|$; and the cubic remainder is $\mathcal S[\mathsf b\eta]=\sum_pL^{2K}\mathsf b\,\eta^{\rho(p)}(x(p))\,r^\sharp_{\rho(p)}(p)$, $\rho(p)$ the normal of $p$, with third-order remainders $|r^\sharp_\rho(p)|\le C\bar a(p)^{3/2}$, $\bar a(p)$ the largest $a_q$ over the plaquettes $q$ next to $p$. $\Delta_K[\eta]:=\mathrm{St}_K[\eta]-\mathsf b\,O^T_K(H^{\rm cons}[\eta])$ is a site form plus a plaquette form of vanishing symbol with weights at most $22\mathsf b\sup|\nabla\eta|$, plus a site form and a plaquette form with weights at most $60\mathsf ba_K\sup|\nabla^2\eta|$, all carried within $2a_K$ of $\mathrm{supp}\,\eta$; and
\begin{equation}\label{sf-ward}
\mathsf b\,O^T_K(H^{\rm cons}[\eta])=\mathrm{St}_K[\eta]-\Delta_K[\eta],\qquad\nu(A^\nu[\mathrm{St}_K[\eta]]\,Y)=-\nu(A^\nu[\mathrm{Rm}_K[\eta]]\,Y)
\end{equation}
for every Gibbs state or torus measure $\nu$ at $\beta$ and every bounded $Y$ depending on no bond within $2a_K$ of $\mathrm{supp}\,\eta$.
\textup{(b)} For $0<g\le\gamma$, $\ell'\in\mathfrak L(g)$ and $\eta$ in a cell of a board of eight cells of $T_{\ell'}$ with its offset in $L^{-\lceil K/400\rceil}\mathbb Z^3$, $\ge\delta$ from its walls, $0<\delta\le\frac{\ell'}{16}$, there are $C,\mathsf p,K_0$ depending only on $g,\ell',\delta$ such that for $(K,\beta)\in\mathfrak d(g)$ with $K\ge K_0$: $\|\mathrm{St}_K[\eta]-\langle\mathrm{St}_K[\eta]\rangle\|_{{\rm ch},\ell'}=\|\mathrm{Rm}_K[\eta]-\langle\mathrm{Rm}_K[\eta]\rangle\|_{{\rm ch},\ell'}\le CK^{\mathsf p}L^{-K/28}\|\eta\|_{C^2}$, and the same with $L^{-K/36}$ for $\Delta_K[\eta]$ and for $O^T_K(H^{\rm cons}[\eta])$; a tensor entry in a cell has the crude bound $\|O^T_K(H')-\langle O^T_K(H')\rangle\|_{{\rm ch},\ell'}\le CK^{2\mathfrak p_0}L^{K/100}|H'|$.
\textup{(c)} On $T_\ell$, $\ell=2L^m$, at the labels of a compact interval $[\mathsf b_-,\mathsf b_+]$ with $\gamma^{-2}\le\mathsf b_-$ and $MR(\mathsf b_+)\le L^m$, $R(x)$ the least power of $L$ that is $\ge(\log x)^{\mathsf r}$, $\mathsf r$ basic (the condition of Section~\ref{sec:cutoff-calibration} on the side $\ell$), for $0<\phi<\frac1{42}$, $t=K-\lceil\phi K\rceil$ and $K\ge K_\circ$: (i) for every finite sum $J$ of gauge-invariant polynomial densities on boxes of at most $3\times3\times3$ sites that vanish to second order at $U=\mathbf 1$, and $|s|$ small, the image of the measure tilted by $e^{sJ}$ and restricted to the small fields under $k\le t$ block averagings has the density $\exp[\sum_{i\le k}c_i(s)]1_{{\rm Cut}_k}e^{\Phi^s_k}$ of Theorem~\YMMref{D}{8} with $sJ$ entered as one more local family; (ii) on the small-field set the conditional mean of $\mathcal S[\mathsf b\eta]$ given the block field at the level $t$ is constant up to $C_\circ(\sup|\eta|+\mathrm{Lip}\,\eta)K^{p_\circ}L^{-c'_0K}$, $c'_0=\min\{\frac1{28},\frac1{14}-3\phi\}$, with $K_\circ$, $C_\circ$, $p_\circ$ depending only on $\ell$, the interval and $\phi$.
\end{proposition}
\begin{proof}[Proof sketch]
(a) Integration by parts in the Haar measure along a field of directions $\mathfrak v$ makes $\beta\mathfrak vS_W-\mathrm{div}\,\mathfrak v$ vanish against every $Y$ that depends on no bond moved by $\mathfrak v$. The derivative $\beta\mathfrak v_\eta S_W$ expands into $\Theta_\beta[\eta]$, $\mathsf b\mathcal O(\nabla^{\rm s}\cdot\eta)$ and the orbital insertion, whose cubic part is $\mathcal S[\mathsf b\eta]$; the divergence is $\mathrm{div}\,\mathfrak v_\eta=-\frac32a_K\mathcal O(w_\eta)$. This gives the identity and the second relation of \eqref{sf-ward}. A direct computation of the symbol shows that $\Theta_\beta[\eta]$ has the symbol $\mathsf b[2(\nabla\cdot\eta)\mathbf 1-\nabla\eta-(\nabla\eta)^t]$ and $\mathsf b\mathcal O(\nabla^{\rm s}\cdot\eta)$ the symbol $\mathsf b(\nabla\cdot\eta)\mathbf 1$, so that $\Delta_K[\eta]$ has vanishing symbol up to the terms with the factor $a_K$, which come from Taylor's formula for $\eta$ between $x_b$, $x(p)$ and the sites. (c) By Theorem~\YMMref{D}{8} with $sJ$, here $s\mathcal S[\mathsf b\eta]$, as one more local family of small size, the tilted measure has after $k$ block averagings the stated density, only the first step needing work; the conditional mean is the $s$-derivative of the exponent at $s=0$; the local terms born from the cubic entry are small by three gains: the third order, $\sup|s_{n_0}|\lesssim L^{\frac92y-K}$ at the depth $y$ of birth; the parity $B\mapsto-B$ of the fluctuation field; and the oddness of $r^\sharp$ under point reflections, read with a H\"older gauge of the background; the count at $y=\lceil K/7\rceil$ gives $c'_0$. (b) Take the depth $u=\lceil K/400\rceil$ and the level $t=K-u$; as the offset of the board lies in $L^{-u}\mathbb Z^3$, its mirror maps carry the lattices of the levels up to $t$ onto themselves, so that the copies of a local form are local forms of the same order and weights, with the same conditional means up to the map.

The eighth power of a seminorm is the mean on $T_{\ell'}$ of the eight centred mirror images, a sum over the partitions of the eight copies without singletons of products of their cumulants. In every such cumulant $\mathrm{St}_K$ becomes $-\mathrm{Rm}_K$ by \eqref{sf-ward}, slot by slot, as each copy depends on no bond near the other copies of $\mathrm{supp}\,\eta$, and the signs multiply to $1$ over a partition: the two seminorms are equal. A cumulant of $r+1\ge2$ local forms of orders $\ge2$ with carriers far apart is at most a pure number times the product of the oscillations $\omega_t$ of their conditional means given the block field of the level $t$ (the sup norm, over the small fields of that level, of the conditional mean less a constant), plus a term of the polymers that link two carriers, at most $\sum_iN_ie^{-\delta L^{K-i}/M}$ with $N_i$ the number of the sites of the level $i$, and a term of the histories with a large field born at a level up to $t$, at most the product of the suprema, $\le CL^{5K}$ each, times their weight $R_t\le CL^{3u}\exp[-c'A_0^2(\frac K{400}\log L)^{2\mathfrak p_0}]$; these two are below $L^{-8K}$ for large $K$.

The oscillations: of $\mathcal S[\mathsf b\eta]$, at most $CK^{\mathsf p}L^{-K/28}\|\eta\|_{C^2}$ by (c) at $\phi=\frac1{400}$, where $c'_0=\frac1{28}$; by power counting at the backgrounds of the level $t$, a form of the order $2$ with the weights $\varkappa$ has $\omega_t\le CK^{2\mathfrak p_0}L^{4u-K}|\varkappa|$, and one of an order $\ge3$ has $\omega_t\le CK^{3\mathfrak p_0}L^{4u-\frac32K}|\varkappa|$; so the four forms of the order $2$ of $\mathrm{Rm}_K$, the three of $J_\eta$ and $\frac32a_K\mathcal O(w_\eta)$, with weights $O(\mathsf b)\|\eta\|_{C^2}$, have $\omega_t\le CK^{2\mathfrak p_0}L^{4u-K}\|\eta\|_{C^2}$, and the forms of the orders $3$, $4$, with weights $L^K\mathsf b\sup|\nabla\eta|$, have $\omega_t\le CK^{3\mathfrak p_0}L^{4u-\frac K2}\|\eta\|_{C^2}$; $4u-\frac K2\le-\frac K{28}$ for large $K$.

Taking the eighth root gives $CK^{\mathsf p}L^{-K/28}\|\eta\|_{C^2}$ for $\mathrm{St}_K$ and $\mathrm{Rm}_K$. By counting, $\Delta_K[\eta]$ is a site form plus a plaquette form of vanishing symbol, with weights at most $22\mathsf b\sup|\nabla\eta|$, plus a site form and a plaquette form with weights at most $60\mathsf ba_K\sup|\nabla^2\eta|$; written as local forms $\sum_p\varkappa(p)j(p)$ of the order $2$, the latter have weights $\le60\mathsf b\sup|\nabla^2\eta|$, the factor $a_K$ cancelling the factor $a_K^{-1}$ of a site form, and $\omega_t\le CK^{2\mathfrak p_0}L^{4u-K}\|\eta\|_{C^2}$. The form of vanishing symbol, with weights $\le22\mathsf bL^K\sup|\nabla\eta|$ as a local form, is cut at the level $n_\Delta=K-y_\Delta$, $y_\Delta=\lceil K/14\rceil$: at the regular configurations of that level it equals its symbol, which is $0$, up to $CK^{2\mathfrak p_0}L^{\frac72y_\Delta-\frac K2}\mathsf b\sup|\nabla\eta|$; what the steps below $n_\Delta$ create is created at a small coupling and adds at most $CK^{\mathsf p}L^{3u-y_\Delta/2}\le CK^{\mathsf p}L^{3-K/36}$, as $3u-\frac12y_\Delta\le3-(\frac1{28}-\frac3{400})K\le3-\frac K{36}$; and above $n_\Delta$ nothing is created, a conditional mean being at most a supremum.

So $\|\Delta_K[\eta]-\langle\Delta_K[\eta]\rangle\|_{{\rm ch},\ell'}\le CK^{\mathsf p}L^{-K/36}\|\eta\|_{C^2}$, and, as $\mathsf bO^T_K(H^{\rm cons}[\eta])=\mathrm{St}_K[\eta]-\Delta_K[\eta]$, the seminorm is subadditive and $\mathsf b\ge\gamma^{-2}\ge1$, the same holds for $O^T_K(H^{\rm cons}[\eta])$. The crude bound: $O^T_K(H')$ is one form of the order $2$ with weights $\le3L^K|H'|$, whose oscillation at the level $t$ is at most $CK^{2\mathfrak p_0}L^{4u}|H'|$, and $4u\le\frac K{100}+4$.
\end{proof}

\begin{proposition}[Spectrum, point operators and totality]\label{sf-spec}
Let $\sigma$ be a real symmetric functional on $\mathfrak F(\mathbb R^3)$, multilinear at disjoint supports, with $\sigma(\emptyset)=1$, $\sigma((H))=0$, bounded by $C_n(d;R)\prod_a|H_a|$ at supports in balls of radius $R$ pairwise $\ge d$ apart, invariant under $\mathbb R^3\rtimes W_3$ and reflection positive at the planes of the nine families: the properties of Theorem~\ref{sf-gibbs}\textup{(c)}.
\textup{(a)} $\mathcal H(\sigma)$, $H$, $\mathcal U$, $\mathcal P$ and the eight $w$ are as in Theorem~\ref{sf-wight}\textup{(a)}, and $\mathcal P(\{E<\max_k|q_k|\})=0$. If $(u,u',u'')$ is an orthonormal basis, $n_\beta=\cos\beta\,u+\sin\beta\,u'$, and $\sigma$ is reflection positive at the four planes through the axis $\mathbb Ru''$ with the normals $n_0,n_{\pi/4},n_{\pi/2},n_{3\pi/4}$ and at one plane with the normal $(n_{\beta_1}+u'')/\sqrt2$, $\beta_1\in\frac\pi4\mathbb Z$, then in the picture of each $n_\beta$, $\beta\in\frac\pi4\mathbb Z$, the spectral measure is carried by $\{E\ge|q|\}$; the nine families give this in the picture of each of their mirrors. The hypotheses ask for no invariance beyond $\mathbb R^3\rtimes W_3$, and the conclusion is stable under $\sigma\mapsto\sigma^z$ of Theorem~\ref{sf-wight}\textup{(h)}.
\textup{(b)} $\Omega$ and the classes of the tuples whose entries lie in pairwise disjoint slabs $\{t<x_3<t'\}$ are total in $\mathcal H(\sigma)$: for every $X\in\mathfrak S(\{x_3>0\})$ and $\epsilon>0$ there are finitely many such ordered tuples $X_i$, with entries $\psi\cdot(H_a*\varphi')$, $\varphi'$ a mollifier and $\psi$ smooth, $0\le\psi\le1$, such that $\|[X]-\sum_i[X_i]\|<\epsilon$. So two such functionals that agree at the ordered tuples are equal.
\textup{(c)} If the slab operators $\mathsf B_{\mathsf l}(X)[F]=[X\cdot\tau_{\mathsf le_3}F]$, $X$ in a cell of $\{0<x_3<\mathsf l\}$, are bounded by $C_{\mathsf l}\prod_a|h_a|$, then the point operators of Theorem~\ref{sf-wight}\textup{(b)} exist and are unique, with the relations stated there and $\|\mathsf B(s,s';e)\|\le C_{\rm pt}(\min\{s,s'\})|e|$, $C_{\rm pt}(\eta)=6(80/\mathsf l_\eta)^3C_{\mathsf l_\eta}$, $\mathsf l_\eta$ the largest $L^{-i}$ with $\frac9{16}\mathsf l_\eta\le\eta$ and $C_{\mathsf l}$ the bound for one entry at the margin $\frac{\mathsf l}{16}$. The functions
\[
\begin{multlined}
\mathbf W_n(z;e)=\langle\Omega,\mathsf B(\epsilon,\epsilon;e_1)\mathsf E(z^1-2\epsilon e_3)\mathsf B(\epsilon,\epsilon;e_2)\cdots\mathsf E(z^{n-1}-2\epsilon e_3)\mathsf B(\epsilon,\epsilon;e_n)\Omega\rangle,\qquad\mathsf E(z) \\
=\int e^{-z_3E+i(z_1q_1+z_2q_2)}d\mathcal P,
\end{multlined}
\]
do not depend on $\epsilon$, are holomorphic where all $z^k$ lie in the square tube $\{\mathrm{Re}\,z_3>|\mathrm{Im}\,z_1|+|\mathrm{Im}\,z_2|\}$, and in the circular tube $\{\mathrm{Re}\,z_3>|(\mathrm{Im}\,z_1,\mathrm{Im}\,z_2)|\}$ if the circular cone holds, are bounded by $\prod_kC_{\rm pt}(\epsilon)|e_k|$ where all $z^k-2\epsilon e_3$ lie in the tube, and give the kernels at supports ordered in $x_3$: $\sigma(h_1,\dots,h_n)=\int\mathbf W_n(x^2-x^1,\dots,x^n-x^{n-1};h_1(x^1),\dots,h_n(x^n))\,dx$; by $W_3$ the same holds in the other coordinate directions. If $\|\mathsf B_{\mathsf l}((h))\|\le c\sup|h|$ for $h$ supported in a cube of half side $\rho$, then $\|e^{-\varepsilon H}\mathsf B_{\mathsf l}((h))e^{-\varepsilon H}\|\le4c\rho^{-3}\|h\|_{L^1}$ for $\varepsilon\ge12\rho$.
\end{proposition}
\begin{proof}[Proof sketch]
(a) Iterated Schwarz inequalities for the positive form, $|u(k)|\le u(0)^{1-2^{-m}}|u(2^mk)|^{2^{-m}}$ with $u(k)=\langle A,\tau_{ke_3}A\rangle$, together with the bounds of $\sigma$, which are uniform in the positions, give contractions; so $e^{-tH}$ is a self-adjoint contraction semigroup and the spectral theorem gives $\mathcal P$. The square cone comes from two pictures of one correlation. Let $e_\pm=(e_3\pm e_k)/\sqrt2$, $\theta_+$ the reflection in the diagonal plane $\{x_3+x_k=0\}$, and $F$ a combination of tuples in the wedge $\{x_3>|x_k|\}$; then $F$ and $\theta_+\theta F$ lie on the open side of the diagonal plane. For $s,s'\ge0$ one and the same number $\sigma(\theta\bar F\cdot\tau_{se_++s'e_-}F)$ is, in the picture of $e_3$, the transform $\int e^{-sw-s'\bar w}\,d\mu$, $w=(E-iq_k)/\sqrt2$, of the image $\mu$ of the spectral measure of $[F]$ under $(E,q)\mapsto(E,q_k)$, a positive measure, and, in the picture of the diagonal mirror, the matrix element $\langle[\theta_+\theta F],e^{-sH_+}e^{-is'Q}[F]\rangle$ with the energy $H_+\ge0$ and a momentum $Q$ of that picture; the reflections act on the labels too, and $\theta_+\theta_+=1$.

By the lemma on a measure on a half plane seen along two directions at the angles $\pm\frac\pi4$ (Lemma~\YMMref{94.3}{722}), $\mu(\{|q_k|>E\})=0$, as in Proposition~\YMMref{94.4}{722}, which treats the scalar functionals in the same way. For a combination $F'$ of tuples in $\{x_3>0\}$ with supports in $B(0,r)$, $\tau_{re_3}F'$ lies in the wedge, $[\tau_{re_3}F']=e^{-rH}[F']$, $\mathcal P$ commutes with $e^{-rH}$ and $\ker e^{-rH}=0$; the classes are dense: $\mathcal P(\{E<|q_k|\})=0$. The same at the orthonormal pairs $(n_\beta,n_{\beta+\pi/2})$ and $(n_{\beta_1},u'')$, with the tilted mirrors $n_{\beta+\pi/4}$ and $(n_{\beta_1}+u'')/\sqrt2$, gives the square cone in each of these pictures, and the nine families give it in the picture of each of their mirrors.

For the circular cone, let $F$ be a combination of tuples in $\{x\cdot u>0\}$ with supports in the open ball $B(0,R)$, $A=\theta_u\bar F$, $B=F$, and $\mathsf u(x)=\sigma(A\cdot\tau_xB)$ for $|x|\ge2R$. By the symmetry of $\sigma$ this one number is a matrix element in every picture whose mirror separates $A$ from $\tau_xB$: for $\beta\in\frac\pi4\mathbb Z$, $n'_\beta=n_{\beta+\pi/2}$ and $x\cdot n_\beta\ge2R$, $\mathsf u(x)=\int\exp(-(x\cdot n_\beta-2R)E+ix\cdot q)\,dm^\beta(E,q)$, $m^\beta$ the spectral measure of the picture of $n_\beta$ between $[\theta_{n_\beta}\tau_{-Rn_\beta}\theta_uF]$ and $[\tau_{Rn_\beta}F]$, and $m^0=e^{-2RE}$ times the spectral measure of $[F]$ in the picture of $u$, a positive measure.

By the square cone, $m^\beta$ is carried by $\{E\ge|q\cdot n'_\beta|\}$ for every $\beta$, and $m^{\beta_1}$ by $\{E\ge|q\cdot u''|\}$, the cone in the picture of $n_{\beta_1}$ that the tilted mirror $(n_{\beta_1}+u'')/\sqrt2$ gives. Let $\kappa>0$ and cut every $m^\beta$ to $\{|q\cdot u''|\ge\kappa\}$, the momentum along the axis common to the eight pictures: as $s\mapsto\mathsf u(\rho n_t+su'')$ is the Fourier transform of the images of the $m^\beta$, weighted by the exponentials, under $q\mapsto q\cdot u''$, the cut transforms agree too, two complex measures with one Fourier transform being equal. With $m_\beta$ the image of the cut $m^\beta$ under $(E,q)\mapsto(E,q\cdot n'_\beta)$, carried by $\{E\ge|q|\}$, $q$ now real, the functions $f_\beta(c;\rho)=\int\exp(2RE-\rho[E\cos(c-\beta)-iq\sin(c-\beta)])\,dm_\beta$, $c$ the angle $t$ continued to complex values, agree where both are defined; $\mu=e^{2RE}m_0$ is a positive measure, and $m_{\beta_1}$ is carried by $\{E\ge\kappa\}$.

Apply Lemma~\ref{C-355-26}(c) to the eight cut measures just constructed. Its band comparison gives $\mu(\{E^2-q^2<\kappa^2\})=0$. So the spectral measure of $[F]$ in the picture of $u$ gives no mass to $\{|q\cdot u''|\ge\kappa,\ E\ge|q\cdot u'|,\ E^2-(q\cdot u')^2<\kappa^2\}$; a point with $|q\cdot u'|\le E<|q|$ lies in one of these sets for a rational $\kappa$, and the square cone excludes $E<|q\cdot u'|$; so $\{E<|q|\}$ is a null set, and the classes $[F]$ are dense. The map $[F]\mapsto[Z_zF]$ is unitary and intertwines $e^{-tH}$ and $\mathcal U$, whence the stability. (b) For two groups $Y_L$, $Y_R$ of smooth entries parted by a vertical plane $\{x_\mu=c\}$, $\mu\le2$, at a distance $\ge r$, $s\mapsto[\tau_{s_1e_3}Y_L\cdot\tau_{s_2e_3}Y_R]$ is $C^2$ in norm, indeed real-analytic, with bounds independent of the entries: its Gram kernel $\langle v(x-y),v(x+y)\rangle$ moves each group rigidly with its mirror image, by $y$, and, after the element of $\mathbb R^3\rtimes W_3$ that turns the plane into $\{x_3=0\}$, is a matrix element of $e^{-rH}\mathcal U(\eta e_\mu)$ between two vectors, smooth in $\eta$ by the square cone; this uses the bounds of $\sigma$ at four times the numbers of entries.

With the norms $\|g\|=\max\{\sup|g|,\sup|\partial_3g|\}$ and $\|g\|_{\infty,1}$, the supremum over $(x_1,x_2)$ of the $L^1$ norm in $x_3$, writing an entry as $\partial_3$ of its primitive in $x_3$ plus a bump in $x_3$ times its integral over the height gives $\|[Y]\|\le C\|g_a\|_{\infty,1}\|g_b\|_{\infty,1}\prod_{e\ne a,b}\|g_e\|$ for $g_a\in Y_L$, $g_b\in Y_R$: $L^1$ in the height in two entries at once. Now mollify the entries, $H_a*\varphi'$, and cut all of them, by a smooth partition of unity of mesh $r$ in the three coordinates, into pieces carried by boxes of sides $\le2r$ in $\{x_3\ge2t_1\}$, $t_1>0$, pairwise $\ge4nr$ apart in the sup distance, $n$ the number of the entries; call a pair of boxes bad if their intervals of heights meet, and argue by induction on the number of the bad pairs, $n$, $r$, $t_1$ being fixed.

With no bad pair the tuple is ordered. For a bad pair $\{a,b\}$ the two boxes are $\ge4nr$ apart in $x_1$ or $x_2$, say in $x_\mu$; the planes $\{x_\mu=c\}$ that part them at a distance $\ge r$ from both fill an interval of length $\ge(4n-2)r$, and those that come within $r$ of one of the $n-2$ other boxes fill at most $(4n-8)r$: there is a plane that parts $a$ from $b$ and puts every box at a distance $\ge r$ on one side, a cut into two groups. Cut $g_a$ and $g_b$ alone into layers of height $\frac1N$ in $x_3$, the other entries being kept: the products of two layers at least two steps apart are tuples with one bad pair fewer, and the at most $3(2rN+3)$ products of near layers have, by the $L^1$ bound, norms at most $K_0N^{-2}\sup|g_a|\sup|g_b|\prod_{e\ne a,b}\|g_e\|$, $K_0$ independent of $N$, so that their sum tends to $0$ as $N\to\infty$; the number $N$ of each stage is chosen after the number of the terms of the stage is known.

Equally, $\varrho\mapsto[Y_L\cdot\tau_{\varrho e_3}Y_R]$ is real-analytic in norm and has one bad pair fewer for large $\varrho$, so that a vector orthogonal to the ordered classes is orthogonal to it for every $\varrho$, at $\varrho=0$ too. (c) For $\omega,\omega'\in\mathcal H(\sigma)$, $\mu(h)=\langle e^{-\delta H}\omega',\mathsf B_{\mathsf l}((h))e^{-\delta H}\omega\rangle$ has $|\mu(h)|\le C|h|$; by the covariance under translations $y\mapsto\mu(\tau_yh)$ continues holomorphically to $|\mathrm{Re}\,y_3|+|\mathrm{Im}\,y_1|+|\mathrm{Im}\,y_2|<\delta$, by the square cone, so Cauchy's estimate gives $|\mu(\partial_Sh)|\le(\frac3\delta)^{|S|}C|h|$; a compactly supported measure whose derivative $\partial_1\partial_2\partial_3$ is a measure of mass $m$ is a function bounded by $m$; this gives the continuous densities, which define the point operators.

The last bound has a road of its own. Let $\chi$ be the tent of mass $1$ on the cube, $\sup\chi=\rho^{-3}$, and write $h=\int\tau_u\big(\chi\,h(\cdot+u)\big)\,du$, the integral over $|u|_\infty\le2\rho$ of functions carried by the cube. The covariance of the slab operators under the semigroup and the translations, $e^{-tH}\mathsf B_{\mathsf l}((h))e^{-t'H}=\mathsf B_{\mathsf l+t+t'}((\tau_{te_3}h))$ and $\mathcal U(v)\mathsf B_{\mathsf l}((h))\mathcal U(-v)=\mathsf B_{\mathsf l}((\tau_{(v,0)}h))$, puts $u$ into the two exponentials, $e^{-(\varepsilon+u_3)H}\mathcal U(\vec u)$ on the left and $\mathcal U(-\vec u)e^{-(\varepsilon-u_3)H}$ on the right. Their power series in $u$ under $e^{-\varepsilon H}$ converge in norm, as $\|A^\alpha e^{-\varepsilon H}\|\le|\alpha|!\,\varepsilon^{-|\alpha|}$ for the monomials $A^\alpha$ in $H$, $P_1$, $P_2$, by the square cone, and $|u|_1\le6\rho\le\frac\varepsilon2$; integrated over $u$, the coefficients are slab operators of the functions $k_\gamma(y)=\chi(y)\int(x-y)^\gamma h(x)\,dx$, with $\sup|k_\gamma|\le\rho^{-3}(2\rho)^{|\gamma|}\|h\|_{L^1}$; and the double series is at most $c\rho^{-3}\|h\|_{L^1}(\sum_n(6\rho/\varepsilon)^n)^2\le4c\rho^{-3}\|h\|_{L^1}$.

For supports ordered in $x_3$ and cut by planes $x_3=c_k$, $[X]=\mathsf A_{c_1}(h_1)\mathsf A_{c_2-c_1}(\tau_{-c_1e_3}h_2)\cdots\Omega$ with $\mathsf A_c(h)=\int\mathcal U(\vec x)\mathsf B(x_3,c-x_3;h(x))\mathcal U(-\vec x)\,dx$; the semigroup relations give $\mathbf W_n$, by induction on $n$, and its holomorphy is that of $\mathsf E(z)$, a contraction on the tube by the cone.
\end{proof}

\begin{proposition}[Rotations from conservation]\label{sf-rot}
Let $\sigma$ be as in Proposition~\ref{sf-spec}, with bounded slab operators, and assume: \textup{(i)} $\sigma(\theta H^{\rm cons}[\eta],H^{\rm cons}[\eta])=0$ for $\eta\in C^2_c(\{x_3>0\};\mathbb R^3)$; \textup{(ii)} $\sigma(\theta(he_{\mu\nu}),he_{\mu\nu})\le C|h|^2$ for $\mu\ne\nu$ and real $h$ supported in $B(x_0,\mathsf d)$, $0<\mathsf d<1$, $\frac32\mathsf d\le x_{0,3}\le2\mathsf d$; \textup{(iii)} for some $a<4$, $\sigma(\theta(f\mathbf 1),f\mathbf 1)\le C_\varphi t^{-a}$ for $0<t\le1$ and $f=\alpha(x_3,x_1)\varphi(x_2)$ with $\mathrm{supp}\,\alpha\subset[t,2t]\times[-t,t]$, $|\alpha|\le3$, $\mathrm{Lip}\,\alpha\le\frac{12}t$; \textup{(iv)} for some $p<\infty$, $\|\mathsf B_{2d}((\tau_{de_3}h))\|\le C\min\{1,d-d'\}^{-p}\|h\|_{L^1}$ for $0\le d'<d$ and $h$ supported in $\{|x_3|\le d'\}$.

Then $\sigma(RX)=\sigma(X)$ for all $R\in O(3)$, and $\sigma((H^{\rm cons}[\eta])\cdot X)=0$ for $X$ off $\mathrm{supp}\,\eta$. Write $u_\pm=e_3\pm ie_1$ and use the basis $\varepsilon_{\mp2}=u_\pm u_\pm^t$, $\varepsilon_{\mp1}=u_\pm e_2^t+e_2u_\pm^t$, $\varepsilon_0=e_3e_3^t+e_1e_1^t$, $\varepsilon'_0=e_2e_2^t$ of $V_{\mathbb C}$, on which the rotations $R_t$ about $e_2$ act by the phases $e^{ik(a)t}$, $k(\varepsilon_j)=j$, $k(\varepsilon'_0)=0$. A helicity label is a pair $\lambda=(\varphi,a)$ of a real $\varphi\in C^1_c(\mathbb R)$, the profile in $x_2$, and a member $a$ of this basis or its negative, with $k(\lambda)=k(a)$; the plate of a continuous $\alpha$ with compact support in the plane $(x_3,x_1)$ is the entry $(\alpha\otimes\lambda)(x)=\alpha(x_3,x_1)\varphi(x_2)a$.

In this plane write $\xi=(\xi^0,\xi^1)$ for $(x_3,x_1)$, $\xi^\pm=\xi^0\pm i\xi^1$ for complex $\xi$, $(R_\zeta\xi)^\pm=e^{\pm i\zeta}\xi^\pm$ for complex $\zeta$, $\mathsf{Tb}=\{\xi\in\mathbb C^2:\mathrm{Re}\,\xi^\pm>0\}$, and $T(\xi)=\int e^{-\xi^0E+i\xi^1q_1}\,d\mathcal P$ on the closure of $\mathsf{Tb}$, contractions by the square cone. The point-plate operators $\mathsf B_\lambda(\xi,\xi')$, $(\xi,\xi')\in\mathsf{Tb}^2$, are the bounded operators, holomorphic in norm, with $T(w)\mathsf B_\lambda(\xi,\xi')T(w')=\mathsf B_\lambda(w+\xi,\xi'+w')$ for $w,w'$ in the closure of $\mathsf{Tb}$ and $\mathsf B_{\mathsf l}((\alpha\otimes\lambda))=\int\alpha(y)\,\mathsf B_\lambda(y,\mathsf le_3-y)\,dy$ for $\alpha$ carried by $\{0<y^0<\mathsf l\}$, $e_3$ standing for the point $(1,0)$ of the plane; the slab bound gives them, and these two relations fix them.

The vectors of $j$ point plates are $\Psi(x;\lambda_1,\dots,\lambda_j)=\mathsf B_{\lambda_1}(x_1,\frac12(x_2-x_1))\mathsf B_{\lambda_2}(\frac12(x_2-x_1),\frac12(x_3-x_2))\cdots\mathsf B_{\lambda_j}(\frac12(x_j-x_{j-1}),w)\Omega$ for $x_1$ and the $x_{a+1}-x_a$ in $\mathsf{Tb}$, which do not depend on $w\in\mathsf{Tb}$. The kernels $\mathcal K_n(Z;\underline\lambda)$ are the ordered kernels of $\sigma$ at plates, $\sigma(\alpha_1\otimes\lambda_1,\dots,\alpha_n\otimes\lambda_n)=\int\prod_a\alpha_a(y_a)\,\mathcal K_n(Y;\underline\lambda)\,dY$ at supports ordered in $x_3$, continued and glued over the permutations of the points and the frames $R_\upsilon$, $\upsilon\in\frac\pi2\mathbb Z$, to one holomorphic function on the union $\mathcal Z_n$ of the tubes of the points $Z\in(\mathbb C^2)^n$ whose successive differences, in some order, lie in $R_\upsilon\mathsf{Tb}$.

For helicity labels $\underline\lambda$ let $\kappa=\sum_ak(\lambda_a)$. Then also $\mathcal K_n(R_{is}Z;\underline\lambda)=e^{\kappa s}\mathcal K_n(Z;\underline\lambda)$ for $Z\in\mathcal Z_n$, and $e^{-\kappa s}\|\Psi(R_{is}z;\underline\lambda)\|=\|\Psi(z;\underline\lambda)\|$ for real $z$ with $z_1$ and the $z_{a+1}-z_a$ in $\mathsf{Tb}$, for real $s$. In \textup{(iv)} a weaker growth suffices: that the vectors of $j$ point plates grow along $R_{is}$ at most like $e^{A_j|s|}$ with $\liminf_jA_j/j^2<2$, that is $\sup_s\sup_ze^{-A_j|s|}\|\Psi(R_{is}z;\underline\lambda)\|<\infty$ for every $j$-tuple of labels and every compact set of real $z$ whose $z_1$ and $z_{a+1}-z_a$ have angles $<\frac{5\pi}{16}$ with $e_3$, nothing being asked of the suprema.

Given the slab bound and invariance at two points under the $R_t$, these are equivalent: invariance under $O(3)$; the exact norms along the complex rotations; the growth with $A_j=2j$; a bound $e^{\gamma|s|}$, for one $\gamma<\infty$, of the operators $\mathsf B_\lambda(R_{is}\xi,R_{is}\xi')$ of one point plate, on compact sets of real pairs $(\xi,\xi')$ of angles $<\frac{5\pi}{16}$ with $e_3$; and if invariance fails, then $A_a\ge a(a-a'-2)$ for some $a'$ and all $a\ge a'$. Invariance at two points needs of \textup{(iv)} only some exponential growth of one point plate; given it and the invariance of the trace part, invariance at tuples with at most two entries outside $V_A$ needs no growth bound.
\end{proposition}
\begin{proof}[Proof sketch]
The ordered kernels are holomorphic in the plane $(x_3,x_1)$, on the tube $\mathsf{Tb}=\{\mathrm{Re}(\xi^0\pm i\xi^1)>0\}$ of the statement, and the kernels of ordered tuples glue, by the frames $R_\upsilon$, $\upsilon\in\frac\pi2\mathbb Z$, of $W_3$, to the one function $\mathcal K_n$ on the union $\mathcal Z_n$ of the tubes. Write, as there, $\xi^\pm=\xi^0\pm i\xi^1$, and besides $d(\xi)=\min\mathrm{Re}\,\xi^\pm$ and $\partial_\pm=\partial/\partial\xi^\pm$, so that $(R_{is}\xi)^\pm=e^{\mp s}\xi^\pm$ and $d(R_{is}\xi)\ge e^{-|s|}d(\xi)$; and let $\varepsilon_{\mp2}=u_\pm u_\pm^t$, $\varepsilon_{\mp1}=u_\pm e_2^t+e_2u_\pm^t$, $\varepsilon_0=e_3e_3^t+e_1e_1^t$, $\varepsilon'_0=e_2e_2^t$ be the basis of the statement, on which the rotations $R_t$ act by the phases $e^{ikt}$, the helicity $k$ being the index ($k=0$ at $\varepsilon'_0$), and a helicity label $\lambda=(\varphi,a)$ having the helicity $k(\lambda)=k(a)$; $\mathbf 1=\varepsilon_0+\varepsilon'_0$ and $\varepsilon_{-2}-\varepsilon_2=2i(e_3e_1^t+e_1e_3^t)$.

By (iv), applied to plates that shrink to a point, the operators of one point plate obey $\|\mathsf B_\lambda(\xi,\xi')\|\le C\min\{1,d(\xi),d(\xi')\}^{-p}$, $C$ depending on the label, so the vectors of $j$ point plates grow along the complex rotations $R_{is}$ at most like $e^{A_j|s|}$ with $A_j=pj$; below, $A_j$ is any such growth of the vectors of $j$ point plates, as in the weaker hypothesis. At a tuple of $n$ points parted by a mirror the kernel is the scalar product of the vectors of $\lfloor n/2\rfloor$ and of $\lceil n/2\rceil$ point plates (the Gram formula), so Schwarz's inequality and $|\kappa|\le2n$ bound the twisted kernel by a multiple of $e^{(A_{\lfloor n/2\rfloor}+A_{\lceil n/2\rceil}+2n)|s|}$, and as $R_{\pi/2}\in W_3$,
\[e^{-\kappa s}\mathcal K_n(R_{is}Z;\underline\lambda)=\sum_{4|m|\le A_{\lfloor n/2\rfloor}+A_{\lceil n/2\rceil}+2n}\mathsf c_m(Z;\underline\lambda)e^{-4ms}:\]
finitely many harmonics, an entire function of period $\frac\pi2$ and exponential growth being a trigonometric polynomial. Let the top frequency $\bar m(n)$ be the largest $|m|$ with $\mathsf c_m\ne0$ at some labels; so $\bar m(2a)\le\frac12A_a+a$. Schwarz's inequality makes the top harmonic positive: at tuples parted by a mirror, the matrices of $\mathsf c_{-\bar m(2a)}$ over the configurations of $a$ points on either side are limits of $e^{-4\bar m(2a)s}$ times Gram matrices; so the top frequencies obey $2\bar m(a+b)\le\bar m(2a)+\bar m(2b)$, and regular polygons, kept by the rotation by $\frac\pi a$, under which a harmonic takes the phase $e^{i(4m-\kappa)\pi/a}$, give $a\mid2\bar m(2a)$. The diagonal mirror makes every $\bar m(2a)$ even: positivity at $\{x_1+x_3=0\}$, at tuples of $a$ plates, makes the matrices $\big(\sum_m(-1)^m\mathsf c_m(\theta S_i\cup S_{i'})\big)_{i,i'}$ positive at the real configurations $S_i$ of $a$ labelled points of the half plane; positivity continues from the real points to the tube, by the power series about one real point, so that along $R_{is}$ the matrices $\big(\sum_m(-1)^me^{-4ms}\mathsf c_m(\theta S_i\cup S_{i'})\big)$ are positive, and as $s\to\infty$ the form $(-1)^{\bar m(2a)}\mathsf c_{-\bar m(2a)}$ is positive; with the positivity of the top harmonic, which is not zero, an odd $\bar m(2a)$ is excluded.

A convex sequence of integers $\nu(a)=\bar m(2a)$, $\nu(0)=0$, with $a\mid2\nu(a)$ is linear, $\nu(a)=a\nu(1)$, or has $\nu(a)\ge a\nu(1)+\frac12a(a-a')$ for all $a\ge a'$, $a'$ the least $a$ with $\nu(a+1)/(a+1)>\nu(a)/a$; and if every $\nu(a)$ is even, then in the second case $\nu(a)\ge a\nu(1)+a(a-a'-\frac32)$ for $a>a'$, as $2\nu(a)/a$ grows by at least $2$ over two steps from $a'$ on and is then a multiple of $4$ at odd $a$, so that over two steps from an odd $a$ it grows by at least $4$. With $\bar m(2a)\le\frac12A_a+a$ the second case forces $A_a\ge2a^2-(2a'+5)a$, so $\liminf_aA_a/a^2\ge2$: a growth with $\liminf_jA_j/j^2<2$ leaves $\bar m(2a)=a\bar m(2)$, and then $\bar m(n)\le\frac n2\bar m(2)$ by convexity.

At two points, if $\bar m(2)=M\ge1$, $M$ is even, so $M\ge2$. The top harmonic is a positive Hermitian form $\mathsf G$ on the finite combinations of point plates of the half plane, the limit of $e^{-4Ms}$ times the Gram form of the twisted vectors $e^{-k(\lambda)s}\Psi(R_{is}x;\lambda)$; a point plate whose twisted vector grows more slowly than $e^{2Ms}$ is null for $\mathsf G$, and by Schwarz's inequality a null plate has the top harmonic zero against every plate. By (ii), applied to the pieces of size about $d(z)$ into which a point plate is cut in $x_2$, summed by the triangle inequality, an off-diagonal label $o\in V_T\otimes\mathbb C$ has $\|\Psi(z;(\varphi,o))\|\le C\min\{1,d(z)\}^{-3}$; so at the helicity $k$ the twisted square norm of such a label is at most $Ce^{(6-2k)|s|}$, below $e^{4Ms}$ at $k=+1$.

The helicity $+2$ follows through the off-diagonal label $e_3e_1^t+e_1e_3^t=\frac1{2i}(\varepsilon_{-2}-\varepsilon_2)$: at $z=R_{is}x$ the square norm of $\Psi(z;(\varphi,\varepsilon_{-2}))-\Psi(z;(\varphi,\varepsilon_2))$, at most $Ce^{6|s|}$, is $e^{4s}N_{++}-2\,\mathrm{Re}\,N_{+-}+e^{-4s}N_{--}$ in the Gram entries of the twisted vectors of $\varepsilon_2$ and of $\varepsilon_{-2}$, all $O(e^{4Ms})$, $N_{++}$ that of $\varepsilon_2$; so $\mathsf G$ at the plate of $\varepsilon_2$, $\lim e^{-4Ms}N_{++}$, is at most $\lim Ce^{(2-4M)s}=0$. By (iii) the square norm of the vector of the label $\mathbf 1$, of helicity $0$, grows at most like $e^{(4+a)|s|}$, $a<4$ the exponent of (iii), and $4+a<8\le4M$.

So the point plates of the labels $\varepsilon_1$, $\varepsilon_2$ and $\mathbf 1$ are null. The labels $\varepsilon'_0$, $\varepsilon_{-1}$, $\varepsilon_{-2}$ fall by (i), which gives $[(H^{\rm cons}[\eta])]=0$: for $\eta=\alpha(x_3,x_1)\varphi(x_2)w$, $w\in\{u_+,u_-,e_2\}$, $H^{\rm cons}[\eta]=\partial_+(\alpha\varphi)\,c(u_+,w)+\partial_-(\alpha\varphi)\,c(u_-,w)+\alpha\varphi'\,c(e_2,w)$ with $c(p,q)=(p\cdot q)\mathbf 1-pq^t-qp^t$, and integration by parts in the plane gives
\[\partial_+\Psi(\xi;(\varphi,c(u_+,w)))+\partial_-\Psi(\xi;(\varphi,c(u_-,w)))=\Psi(\xi;(\varphi',c(e_2,w)))\qquad(\xi\in\mathsf{Tb})\]
for smooth profiles $\varphi$; so conservation holds harmonic by harmonic. Let $\mathsf h^\varphi_b(\xi)$ be the top harmonic of the kernel of a fixed point plate against the point plate $(\varphi,b)$ at the relative position $\xi\in\mathsf{Tb}$. As $c(u_\pm,u_\pm)=-2\varepsilon_{\mp2}$, $c(u_+,u_-)=2\varepsilon'_0$, $c(u_\pm,e_2)=-\varepsilon_{\mp1}$, $c(e_2,e_2)=\varepsilon_0-\varepsilon'_0$, and as $\mathsf h^\varphi_{\varepsilon_1}=\mathsf h^\varphi_{\varepsilon_2}=0$, $\mathsf h^\varphi_{\varepsilon_0}=-\mathsf h^\varphi_{\varepsilon'_0}$ by the null plates found, the choices $w=u_-,e_2,u_+$ give three equations, which are solved one after the other:
\[\partial_+\mathsf h^\varphi_{\varepsilon'_0}=0,\qquad\partial_+\mathsf h^\varphi_{\varepsilon_{-1}}=2\mathsf h^{\varphi'}_{\varepsilon'_0},\qquad\partial_+\mathsf h^\varphi_{\varepsilon_{-2}}=\partial_-\mathsf h^\varphi_{\varepsilon'_0}+\tfrac12\mathsf h^{\varphi'}_{\varepsilon_{-1}}.\]
In turn, for $b=\varepsilon'_0,\varepsilon_{-1},\varepsilon_{-2}$, the function $\mathsf h^\varphi_b$, now with $\partial_+\mathsf h^\varphi_b=0$ and taken against the point plate of the same label, is a holomorphic function of $\xi^-=\xi^0-i\xi^1$ alone, homogeneous along $R_{is}$ with the exponent $2k(b)+4M$, that is $4M$, $4M-2$, $4M-4$, positive as $M\ge2$, and bounded where $\mathrm{Re}\,\xi^\pm\ge1$, the semigroup being a contraction; so it vanishes, the plate is null, $\mathsf h^\varphi_b=0$ against every plate, and the next equation becomes $\partial_+\mathsf h^\varphi_b=0$ for the next label $b$. With $\varepsilon_0=\mathbf 1-\varepsilon'_0$ every point plate with a smooth profile is null, and by approximation every point plate; so the top harmonic at two points vanishes, against $M\ge1$: $M=0$.

Then every $\bar m(n)=0$, the kernels are covariant under the $R_t$, and the exact norms follow; slides along coordinate planes, along which $s\mapsto\sigma(Y_I\cdot\tau_{su}Y_J)$ is real-analytic, pass to all tuples, and $W_3$ with the $R_t$ generates $O(3)$. By (i), $[(H^{\rm cons}[\eta])]=0$ for $\eta$ above a plane, and slides give conservation among any entries. The equivalences: invariance gives a unitary group of boosts on the closed span of the plate vectors, $U(s)\Psi(w;\underline\lambda)=e^{-\kappa s}\Psi(R_{is}w;\underline\lambda)$, so exact norms and $A_j=2j$; conversely each growth named gives finitely many harmonics and the argument above, a bound $e^{\gamma|s|}$ of the operators of one point plate giving $A_j=\gamma j$; as $\nu(1)=\bar m(2)=0$ by the invariance at two points, a failing invariance leaves the second case for $\nu(a)=\bar m(2a)$, so $\bar m(2a)\ge\frac12a(a-a')$ and, with $\bar m(2a)\le\frac12A_a+a$, $A_a\ge a(a-a'-2)$ for all $a\ge a'$.

With at most two entries outside $V_A$ the invariant scalar vectors have exact norms along the boosts, Schwarz's inequality bounds the harmonics by them, and Liouville's theorem leaves the constant one.
\end{proof}

\begin{proposition}[Reconstruction with values in $V$]\label{sf-rec}
Let $\sigma$ be a real symmetric functional on $\mathfrak S(\mathbb R^3)$, multilinear at disjoint supports, with $\sigma(\emptyset)=1$, invariant under all isometries acting on points and indices, positive at $\{x_3=0\}$, with separable $\mathcal H(\sigma)$, with $H$, $\mathcal U$, $\mathcal P$ as in Theorem~\ref{sf-wight}\textup{(a)}, the circular cone $\mathcal P(\{E<|q|\})=0$, and \eqref{sf-slab} for some $C_V$ and $p_V$. Then: \textup{(i)} there is exactly one continuous $s_n$ at non-coinciding points, $n$-linear in the labels, with $\sigma(X)=\int s_n(x;H_1(x_1),\dots,H_n(x_n))\,dx$ for all $X$; it is real-analytic, symmetric and invariant, and $|s_n(x;e)|\le(cC_V)^n\max\{1,2n^2/\delta(x)\}^{p_Vn}\prod_k|e_k|$; \textup{(ii)} the fields $T(f)$ of Theorem~\ref{sf-wight}\textup{(e)} are defined on $\mathcal D_\infty$, with the bound and the adjoint stated there, covariant under the translations, and their vacuum expectations are tempered, the boundary values of the kernels continued to the forward tube; \textup{(iii)} $\Omega$ and the classes of the tuples in disjoint slabs are total in $\mathcal H(\sigma)$, and $\mathcal D_0$ is dense; \textup{(iv)} there is a strongly continuous unitary representation of $\mathbb R^{1+2}\rtimes\mathcal L$ on $\mathcal H(\sigma)$ as in Theorem~\ref{sf-wight}\textup{(e)}, acting on the labels by $D(\Lambda)e=\hat\Lambda e\hat\Lambda^t$, with spectrum in the closed forward cone; \textup{(v)} the fields commute at spacelike separation.

Nothing is supposed or said of $\ker H$; if $\ker H=\mathbb C\Omega$, all Wightman axioms hold.
\end{proposition}
\begin{proof}[Proof sketch]
It is the proof of Theorem~\textup{\YMMref{AO}{778}} with values in $V$. The ordered kernels $\mathbf W_n$ of Proposition~\ref{sf-spec}(c) continue to the circular tube, and \eqref{sf-slab} gives $|\mathbf W_n|\le(C_V\eta^{-p_V})^n\prod|e_k|$ at distance $2\eta$ from its boundary, with no factorial. For $n$ distinct points more than half of the directions $u$ have $|u\cdot(x_a-x_b)|\ge n^{-2}|x_a-x_b|$ for all $a\ne b$ (zones of Archimedes); rotating such a $u$ to $e_3$, the frames in $SO(3)$ glue the ordered kernels into one continuous, real-analytic function $s_n$, and the bound at $\eta=\min\{1,\delta(x)/(2n^2)\}$ is (i). The limits $\lim_{s\downarrow0}\langle\psi,\mathsf B(s,s;e)\psi'\rangle$ exist for $\psi,\psi'\in\mathrm{dom}\,H^l$, $l>p_V$, the power being paid by $(1+H)^l$; this defines $T(f)$, tempered, with $\langle\Omega,T(f_1)\cdots T(f_n)\Omega\rangle$ the boundary values of the $\mathbf W_n$ at $\iota$ of the real points, from the directions of the forward cone.

The totality (iii) follows from (i): for $X\in\mathfrak S(\{x_3>0\})$ cut every entry by a smooth partition of unity in $x_3$ of mesh $\epsilon$; the terms in which two entries lie in one slab or in neighbouring slabs form a remainder $R_\epsilon$, and $\|[R_\epsilon]\|^2=\sigma(\theta\bar R_\epsilon\cdot R_\epsilon)$ is, by (i), an integral of $s_{2n}$ against the weights of $\theta\bar R_\epsilon$ and $R_\epsilon$, where $s_{2n}$ is continuous, so bounded, on the compact product of the supports of the entries of $\theta X$ and of $X$, which are pairwise disjoint, while the weights of $R_\epsilon$ are carried by a set of measure $O(\epsilon)$; so $[R_\epsilon]\to0$, and $[X]$ is a limit of sums of classes of tuples in disjoint slabs.

A vector orthogonal to $\mathcal D_0$ gives functions holomorphic in the forward tube whose boundary values vanish; their values at imaginary times are its scalar products with the classes of the tuples in disjoint slabs, so by the totality it is $0$, and $\mathcal D_0$ is dense. The complex rotations $R_{is}$ of the plane $(x_3,x_1)$ are the boosts: invariance gives exact norms along them, so they are unitary on the closed span of the point vectors, which is $\mathcal H(\sigma)$ by (iii), and with the translations and the rotations about $e_3$ they generate $U(a,\Lambda)$ on $\mathcal H(\sigma)$, with $U(a,\Lambda)\mathcal D_0=\mathcal D_0$; $D$ is a continuous homomorphism of $\mathcal L$ into $GL(V_{\mathbb C})$, $\hat\Lambda^t\hat\Lambda=\mathbf 1$, with $D(\Lambda)e^*=(D(\Lambda)e)^*$ for $e^*=\Theta_{e_3}\bar e\Theta_{e_3}$, as $\overline{\hat\Lambda}=\Theta_{e_3}\hat\Lambda\Theta_{e_3}$.

Locality: the kernels continue to tubes tilted by an angle of the order of $\eta$, and the symmetry of $s_n$ under permutations of points and labels gives the commutation at spacelike separation, all components with all.
\end{proof}

\begin{proposition}[A count in $n$, without the rotations]\label{sf-count}
Assume \textup{(H)}, $n\ge2$, $\delta_n=\frac1{8n}$, and let $c_g$, $Y_i$ be as in Theorem~\ref{sf-gibbs}\textup{(a)}, $\mathsf g_{n,i}=c_gn^4Y_i^{q_E}L^{11i/2}$. Let $i_E(g,\delta)$ and $K_E(i)=K_E(i;g,\delta)$ be the thresholds of Theorem~\ref{sf-gibbs}\textup{(a)}, written there for all of \textup{(H)}: for $g\le\gamma^{\rm sd}_W$ they are numbers of Lemma~\ref{sf-stop}; for $\gamma^{\rm sd}_W<g\le\gamma^{\rm ref}_W$, $z_\sharp$ is the least integer $z\ge1$ for which the coupling $g_-$ with $g_-^{-2}=L^{z-1}(g^{-2}-\bar B)-B_\sharp$ is $\le\gamma^{\rm sd}_W$, $i_E(g,\delta)=z_\sharp+i_E(g_-,\delta)$ and $K_E(i)=i+s_\sharp(i)$, $s_\sharp(i)$ the depth of the distance $\delta$ formed at the coupling with the inverse square $L^{i-1}(X_g+2\bar B)-B_\sharp$. Here $K_E(i)$ is formed with the margin $\delta=\delta_n$, so that it reads $n$.
\textup{(a)} For $i\ge i_E(g,\delta_n)$, $\mathsf l=L^{-i}$, $(K,\beta)\in\mathfrak d(g)$ with $K\ge\max\{K_E(i),i+\log(4n)/\log L\}$, every law $\mathfrak s$ that is a periodic Gibbs state or Wilson's measure on a torus of $\mathfrak L(g)$, and bounded $H_a$ whose supports have diameters $\le\frac{\mathsf l}{4n}$ and are pairwise more than $\mathsf l$ apart, with $c_a$ the means on $T_{\ell_g}$: $|\mathfrak s(\prod_{a\in B}(O^T_K(H_a)-c_a))|\le\prod_{a\in B}\mathsf g_{n,i}|H_a|$ for every $B\subset\{1,\dots,n\}$, and $|\mathfrak s(\prod_aA^{\mathfrak s}[O^T_K(H_a)])|\le(2\mathsf g_{n,i})^n\prod_a|H_a|$.
\textup{(b)} For a limit $\sigma$ at $g$ and $X\in\mathfrak S(\mathbb R^3)$ with entries pairwise $\ge2\mathsf l$ apart: $|\sigma(X)|\le\prod_a2\mathsf g_{n,i}(4n/\mathsf l)^3\|H_a\|_\natural$, with $\|H\|_\natural=\sum_{\mathsf c}|\mathsf c|\sup_{\mathsf c}|H|$ over the boxes $\mathsf c=\frac{\mathsf l}{4n}(k+[0,1)^3)$, $k\in\mathbb Z^3$, an upper Riemann sum of $\int|H|$.
\textup{(c)} There is $C_G$ depending only on $g$ such that for all $n\ge2$, $0<d\le1$ and $X\in\mathfrak S(\mathbb R^3)$ with $n$ entries pairwise $\ge2d$ apart, ordered or not,
\[|\sigma(H_1,\dots,H_n)|\le(C_Gn^{16}d^{-9})^n\prod_a\|H_a\|_\natural\le(e^{16}C_G)^n(n!)^{16}d^{-9n}\prod_a\|H_a\|_\natural,\]
the norms at the mesh $L^{-i}/(4n)$, $i$ the least integer with $L^{-i}\le d$ and $i\ge i_E(g,\delta_n)$. The same holds for every component $\sigma_\xi$, with the same numbers; and with the extension of bounds on small boxes to Schwartz functions it gives, without the rotations, $|\int s_n(x;\varphi(x))\,dx|\le C^n(n!)^{43/2}|\varphi|_{15n}$ for Schwartz $\varphi$ vanishing with all derivatives at coinciding points.
\end{proposition}
\begin{proof}[Proof sketch]
(a) With the margin $\delta_n=\frac1{8n}$ one board of side $\mathsf l$ can be shifted, by Lemma~\ref{sf-chess}(c) at $h=a_K$, whose condition $2na_K<\mathsf l$ the condition $K\ge i+\log(4n)/\log L$ gives, so that the $n$ small pieces lie well inside $n$ different cells; each piece, centred, has by the display of Step~7 of the proof of the principal theorems below at the margin $\delta_n$, for $i\ge i_E(g,\delta_n)$ and $K\ge K_E(i)$, a seminorm at most $C'_{\delta_n}Y_i^{q_E}L^{11i/2}\mathsf b^{-1}|H_a|$, the constant $C_\delta=\delta^{-4}C_1$ of Lemma~\ref{sf-stop} giving the factor $n^4$, and $\mathsf b^{-1}<g^2$. For $\gamma^{\rm sd}_W<g\le\gamma^{\rm ref}_W$ the display holds in the same form, with the numbers $i_E(g,\delta)$, $K_E(i)$ set in the statement and in Theorem~\ref{sf-gibbs}(a), with $Y_i$ formed with $\log(X_g+2\bar B)$ and a constant proportional to $\delta^{-4}$: by the arithmetic of first passage, \YMMref{(1526)}{615} and \YMMref{(1527)}{615}, the datum stopped at the scale $a_{z_\sharp}$ is a first-passage datum at a coupling $g_{(z_\sharp)}$ with $g_{(z_\sharp)}^{-2}\ge g_-^{-2}$, so $g_{(z_\sharp)}\le\gamma^{\rm sd}_W$, and Lemma~\ref{sf-stop} applies at $g_{(z_\sharp)}$; and $L^{i_E(g,\delta)}\le c_\sharp/\delta$ with $c_\sharp=L^{z_\sharp}c_\sharp(g_-)$, a number of $g$.

Lemma~\ref{sf-chess}(a) at $g$ multiplies the cells, with $e^{G_g\mathsf l^3}\le e^{G_g}$ for each, as $\mathsf l\le1$; the centring gives the $2^n$. (b) Cut each entry by the boxes and apply (a) to every product of pieces. (c) Take $\mathsf l=L^{-i}$ comparable with $d$: the threshold $L^{i_E(g,\delta_n)}\le8nc_\sharp(g)$ of Theorem~\ref{sf-gibbs}(a) at $\delta_n=\frac1{8n}$ and $(4n/\mathsf l)^3L^{11i/2}Y_i^{q_E}$ give $C_Gn^{16}d^{-9}$ an entry, and $n^{16n}\le e^{16n}(n!)^{16}$. The components: Lemma~\ref{sf-plat}(c), the bounds of (a) being uniform in the states, and the conditioned states of a component being periodic Gibbs states (Proposition~\textup{\YMMref{106.1}{764}}). The growth: a Whitney decomposition of the complement of the diagonals into products of small boxes, the bound on each product, and the Taylor subtraction at the nearest diagonal, as in Proposition~\ref{sf-ext}.

For $\gamma^{\rm sd}_W<g\le\gamma^{\rm ref}_W$ the limits, (b), (c) and the growth follow also from the couplings $g'\le\gamma^{\rm sd}_W$ by $\sigma=L^{nz}\sigma_z(X^{(z)})$ of Lemma~\ref{sf-dil}(b), the powers of $L^z$ entering $C_G$ and $C$. As the bounds of (a) are lattice bounds, uniform in the states and the positions, the count holds at tuples ordered in no direction, independently of the invariance under $O(3)$ and of the analyticity of the kernels.
\end{proof}

\begin{proposition}[Extension over the diagonals]\label{sf-ext}
Let $M$ be $T_\ell$ or $\mathbb R^3$ and $u$ a functional of $n\ge2$ entries, $n$-linear on smooth $V_{\mathbb C}$-valued functions with pairwise disjoint supports, with the power counting: there are $A,h_0>0$ and $q\ge0$, with $h_0\in\ell2^{-\mathbb N}$ and $h_0\le\ell/(64n)$ on $T_\ell$, such that for $0<h\le h_0$, all closed cubes $Q_1,\dots,Q_n\subset M$ of side $\frac32h$, wherever they lie, at pairwise sup distances $\ge\frac52h$, and all $f_a$ supported in the interior of $Q_a$: $|u(f_1,\dots,f_n)|\le A\,\mathfrak l(h)^q\prod_a\sup|f_a|$, $\mathfrak l(h)=1+\log\frac{h_0}h$. \textup{(a)} There is exactly one distribution $u^\circ$ on the complement of the diagonals in $M^n$ with $u^\circ(f_1\otimes\cdots\otimes f_n)=u(f_1,\dots,f_n)$. \textup{(b)} It extends to a distribution $\tilde u$ on $M^n$ of order $\le3n-2$ with $|\tilde u(F)|\le C_EA\,\mathfrak l(\lambda)^{q+1}\|F\|_{(\lambda),3n-2}$ for $0<\lambda\le h_0$ and every $F$ supported in a product of closed sup balls of radius $\lambda$, that is of closed cubes of side $2\lambda$, wherever their centres lie, $\|F\|_{(\lambda),m}=\max_{|\gamma|\le m}\lambda^{|\gamma|}\sup|\partial^\gamma F|$; it can be chosen invariant under the symmetries of $u$ among the permutations, the translations and a closed subgroup of $O(3)$ that maps $M$ onto itself. The constant $C_E$ depends only on $n$ and $q$, with the label norm and auxiliary smooth cutoffs fixed.

\textup{(c)} Two extensions with this power counting differ by a sum $\sum_{|\pi|\le n-1}\sum_{|\alpha|\le3|\pi|}\partial_\xi^\alpha\delta_\pi\otimes T_{\pi,\alpha}$, $\delta_\pi$ the delta of the diagonal of the partition $\pi$, $\xi$ the coordinates normal to it and $T_{\pi,\alpha}$ distributions in $|\pi|$ points with the power counting $|\alpha|$; and every such sum may be added.
\end{proposition}
\begin{proof}[Proof sketch]
This is Lemma~\ref{C-355-34}(a)--(d) at $\kappa=r=0$. The admissible products of dyadic cubes, at mutual distances at least three times their side and maximal among such, form a Whitney decomposition of the complement of the diagonals: below the coarsest side $h_0$, the distance to the diagonals is comparable with the side; neighbours have comparable sides and the overlaps are bounded; a smooth partition of unity $\varphi_{\mathbf Q}$ has $|\partial^\gamma\varphi_{\mathbf Q}|\le c_\gamma h(\mathbf Q)^{-|\gamma|}$. On a product of cubes a smooth function is an absolutely convergent sum of products of functions of one point, so $u$ acts on it with the bound $A\mathfrak l(h)^q$ times a scaled $C^N$ norm, which gives (a). For each product $\mathbf Q$, with its clusters, the points closer than $12h$, forming $k$ blocks, subtract from $F$, when $k\le n-1$, the Taylor polynomial of order $3k$ averaged over the piece of size $h$ of the corresponding stratum next to $\mathbf Q$; put $P_{\mathbf Q}=0$ for the other products: $\tilde u(F)=\sum_{\mathbf Q}u_{\mathbf Q}(\varphi_{\mathbf Q}(F-P_{\mathbf Q}F))$; on the products with $k\le n-1$ the subtraction gains $h^{3k+1}$, which beats the number of the products of side $h$ near a point of a stratum with $k$ blocks, and the sum converges, with one more logarithm from the borderline count; for the invariant choice, take the grid partition and Taylor averages covariantly under simultaneous translations by $h_0\mathbb Z^3$, as in that lemma. The resulting extension $T$ is invariant under these translations. Its normalized mean $h_0^{-3}\int_{[0,h_0)^3}T(\tau_tF)\,dt$ is invariant under every simultaneous translation, since the integrand is periodic in each coordinate. Average afterward over permutations and the compact orthogonal subgroup. These means preserve the extension, its order and its power counting by part~(d) of that lemma. General rotations may enlarge sup-support radii by $\sqrt3$; the fixed-scale localization in its part~(c) handles the endpoint scales, with constants depending on $n,q$.

A distribution supported on the diagonals with this power counting is a sum of derivatives of deltas of the strata, and a distribution in $k$ points with a power counting $\kappa'>3k$ vanishes; this gives (c).
\end{proof}

\begin{proposition}[Fields and kernels from the square cone]\label{sf-norot}
Let $0<g\le\gamma^{\rm sd}_W$ and let $\sigma$ be a limit at $g$. Theorem~\ref{sf-gibbs}\textup{(a)}--\textup{(c)} and \eqref{sf-slab}, with the square cone $\mathcal P(\{E<\max_k|q_k|\})=0$ in the place of the circular one, give the following, independently of the invariance under $O(3)$: the part of Theorem~\ref{sf-wight}\textup{(d)}, \textup{(e)} that rests on the square cone, the one step at which the rotations enter Lorentz covariance, and bounded slab operators for entries side by side.
\textup{(a)} There are bounded operators $\mathsf b(w|e|w')$, for $w,w'$ in the open square tube and $e\in V_{\mathbb C}$, linear in $e$ and holomorphic in norm, with $\|\mathsf b(w|e|w')\|\le\sqrt2C_Vs^{-p_V}|e|$ for $0<s\le\min\{1,d_\square(w),d_\square(w')\}$, $d_\square(z)=\mathrm{Re}\,z_3-|\mathrm{Im}\,z_1|-|\mathrm{Im}\,z_2|$, and $\mathsf B_{2\eta}((\tau_{\eta e_3}h))=\int\mathsf b(y+\eta e_3|h(y)|-y+\eta e_3)\,dy$; and $|\mathbf W_n(z;e)|\le(C_V\eta^{-p_V})^n\prod_k|e_k|$ where $d_\square(z^k)\ge2\eta$, $0<\eta\le1$.
\textup{(b)} The fields $T(f)$ of Theorem~\ref{sf-wight}\textup{(e)} are defined on $\mathcal D_\infty$, with the same bound and adjoint, covariant under the translations and the dihedral group of order eight that fixes $e_3$, with joint spectrum in $\{E\ge\max_k|q_k|\}$ and $\Omega$ cyclic in $\mathcal H(\sigma)$; their vacuum expectations are tempered and are the boundary values of the $\mathbf W_n$ from every direction of $\{\alpha^0>|\alpha^1|+|\alpha^2|\}$; and they commute in the square form: $T(f)T(f')=T(f')T(f)$ on $\mathcal D_0$ whenever $\max_k|x^k-x'^k|>|t-t'|$ on the supports. Lorentz covariance needs the rotations at one step only: for one one-parameter group $r_t$ of rotations, the identity $\frac d{dt}\mathbf W_n(r_t\xi;r_te_1r_t^t,\dots,r_te_nr_t^t)\big|_{t=0}=0$ at all points $\xi$ and labels $e$ is equivalent to invariance under $O(3)$.
\textup{(c)} For $m\ge0$ the kernel $s(x;y_1,\dots,y_m;e)=\lim_{\epsilon\downarrow0}\sigma(\delta^x_\epsilon e,\delta^{y_1}_\epsilon\mathbf 1,\dots,\delta^{y_m}_\epsilon\mathbf 1)$ of one tensor entry among $m$ action densities, $\delta^z_\epsilon$ a bump of width $\epsilon$ at $z$, exists at all non-coinciding points, is real-analytic there and represents $\sigma$ at such tuples; where for one direction $\mu$ the coordinate $x_\mu$ differs from all $y_{k,\mu}$, it is the matrix element of a point operator between two point vectors of the scalar theory, the points below and above $x$ in that direction. The trace part is data of Theorem~\textup{\YMMref{AO}{778}}, and its point vectors exist and are norm continuous at all non-coinciding points of the half space.
\textup{(d)} For $l>0$ and a tuple $Y=(h_1,\dots,h_n)\in\mathfrak S(\{0<x_3<l\})$, its entries side by side or not, let $d(Y)=\min\{1,\mathrm{dist}(\mathrm{supp}\,Y,\{x_3=0\}\cup\{x_3=l\}),\min_{a\ne b}\mathrm{dist}(\mathrm{supp}\,h_a,\mathrm{supp}\,h_b)\}$. The slab operator $\mathsf B_l(Y)[F]=[Y\cdot\tau_{le_3}F]$ is bounded,
\[\|\mathsf B_l(Y)\|\le\mathsf K_n\big(\sqrt2C_V(3/d(Y))^{p_V}\big)^n\prod_a\|h_a\|_{L^1},\qquad\mathsf K_n=(10n)^{\lceil p_Vn\rceil n(n-1)/2};\]
and $\sigma(X)=\int s_n(x;h_1(x_1),\dots,h_n(x_n))\,dx$ for every $X=(h_1,\dots,h_n)\in\mathfrak S(\mathbb R^3)$, with one function $s_n$ continuous at all non-coinciding points, given at points of pairwise different heights by $\mathbf W_n$ at the points ordered in $x_3$, and $|s_n(x;e)|\le\mathsf K_n(\sqrt2C_V(3/d_\infty(x))^{p_V})^n\prod_a|e_a|$, $d_\infty(x)=\min\{1,\min_{a\ne b}|x_a-x_b|_\infty\}$.
\end{proposition}
\begin{proof}[Proof sketch]
(a) As in Proposition~\ref{sf-spec}(c), with \eqref{sf-slab} in the place of the finite function $C_{\rm pt}$: covariance moves the point of the plate into the exponentials of the spectral measure, which are contractions on the square tube by the square cone, and the bound is \eqref{sf-slab} at the distance $d_\square$. (b) As in the proof of Proposition~\ref{sf-rec}, with the square tube for the circular one: boundary values as tempered distributions from the directions of the square cone, the translations and the dihedral group of $W_3$ fixing $e_3$ acting, and Proposition~\ref{sf-spec}(b) for the cyclicity. The commutation in the square form: by a partition of unity of $f$ and $f'$ into compact pieces it suffices to treat two pieces, and for them to fix a horizontal direction $k$ with $|x^k-x'^k|-|t-t'|\ge c>0$ on their supports.

First, at a configuration of $n$ points ordered in $x^k$ with $|\Delta x^k|>|\Delta t|$ for every pair, $\iota(t,\vec x)=(\vec x,-it)$, the points $\iota(x)+\eta he_3$, $h=(h_1,\dots,h_n)$ increasing along an order of the points, lie in the $e_3$-tube of that order for $\eta>0$, for either order of the two entries that are exchanged, and inside the $e_k$-tube with a margin free of $\eta$, as $\mathrm{Re}\,\Delta z_k=\Delta x^k>|\Delta t|=|\mathrm{Im}\,\Delta z_3|$; the tubes are convex and have a common real open set, on which $\sigma$ has one kernel, so the two boundary values as $\eta\downarrow0$ are the one function of the picture $e_k$ at $\iota(x)$. Second, $F(a,b)=\langle T(g_1(\cdot-b_1))\cdots T(g_m(\cdot-b_m))\Omega,[T(f),T(f')]T(f_1(\cdot-a_1))\cdots T(f_{m'}(\cdot-a_{m'}))\Omega\rangle$ is continuous where $\mathrm{Im}\,a_1$, the $\mathrm{Im}(a_{i+1}-a_i)$, $\mathrm{Im}\,b_1$ and the $\mathrm{Im}(b_{i+1}-b_i)$ lie in the closed cone $\{\alpha^0\ge|\alpha^1|+|\alpha^2|\}$, dual to the square cone, and holomorphic in $(a,\bar b)$ inside, by the spectrum and the bounds of $T$ on $\mathcal D_\infty$; it vanishes on the open set of the real shifts far along $e_k$, by the first step; the edge-of-the-wedge theorem against $0$ gives $F=0$, and Schwartz $f$, $f'$ follow by density.

The boosts are the complex rotations, and their covariance is the analytic continuation of the rotation covariance from one infinitesimal generator. (c) The trace part has the bounds, the invariance and the positivity required by Theorem~\textup{\YMMref{AO}{778}}, whose reconstruction gives the point vectors $\Phi_n(y)$ of the scalar theory. For $(x,y)$ with $x_\mu\ne y_{k,\mu}$ for all $k$ cut by two planes $x_\mu=c_\pm$ next to $x_\mu$, and write the kernel as $\langle\Phi(\text{reflected points below}),\mathcal U\,\mathsf B_e(x_\mu-c_-,c_+-x_\mu)\,\mathcal U^{-1}\Phi(\text{points above})\rangle$ in the picture $e_\mu$; the three pictures agree where they overlap, by the uniqueness of continuous kernels.

The remaining points form a null set of codimension three; there holomorphy in each coordinate separately with one bound gives joint holomorphy, by Bernstein's estimate of the Chebyshev coefficients, so the kernel is real-analytic at all non-coinciding points. (d) First the continuation and the bound, then the integral. By the symmetry of $\sigma$ the class $[Y\cdot\tau_{le_3}F]$ does not depend on the order of the entries. At points $x$ of pairwise different heights let $\mathsf G_l(x;e)$ be the product of the operators $\mathsf b(w|e|w')$ of (a), taken in the order of the heights, whose arguments are the half gaps of the heights and the distances to the faces of the slab; the spectators enter as two vectors $\mathcal U(v)[F']$, $\mathcal U(w)[F]$, and $f(x;v,w)=\langle\mathcal U(v)[F'],\mathsf G_l(x;e)\mathcal U(w)[F]\rangle$ is analytic in $(v,w)$.

Take a pair $a,b$ of entries with $|x^k_a-x^k_b|\ge d=d(Y)$ in a horizontal direction $k$, and let $\zeta$ be their relative height, the mean height and the other points being fixed: the two orders of the pair give functions $f_\pm(\zeta;v,w)$, holomorphic in $\pm\mathrm{Re}\,\zeta>0$, with $|f_\pm|\le\mathsf M(|\mathrm{Re}\,\zeta|/2)^{-\mathsf P}$ by (a), $\mathsf P$ the power of the bound of (a) in the half gap of the pair ($\mathsf P=2p_V$ at $n=2$). For spectators shifted far along $e_k$, planes $x_k={\rm const}$ part them from the entries; in the picture $e_k$ all is then ordered in $x_k$, $\zeta$ is a translation along the mirror, and the kernel is holomorphic in $|\mathrm{Im}\,\zeta|<|x^k_a-x^k_b|$, equal to $f_+$ and to $f_-$ at the real $\zeta$ of either sign, by the symmetry of $\sigma$ and the uniqueness of continuous kernels.

For such spectators the Cauchy integral over $|\zeta|=r<d$ of $(\zeta^2+r^2)^{\mathsf N}f(\zeta;v,w)$, $\mathsf N=\lceil\mathsf P\rceil$, reproduces $(\zeta_0^2+r^2)^{\mathsf N}f(\zeta_0;v,w)$ inside the circle; it is analytic in $(v,w)$, the factor $(\zeta^2+r^2)^{\mathsf N}$ taming $f$ on the circle, and the identity theorem in $(v,w)$ carries the continuation across $\mathrm{Re}\,\zeta=0$ to all spectators. On the circle $|\zeta^2+r^2|=2r|\mathrm{Re}\,\zeta|$, and at real $\zeta_0$ inside $|\zeta_0^2+r^2|\ge r^2$, so the maximum principle gives $|f|\le2^{\mathsf N}(2/r)^{\mathsf P}\mathsf M$ there, while where $|\mathrm{Re}\,\zeta|\ge r$ the bound (a) serves directly; at $n=2$, $r=\frac23d$ gives the factor $2^{\mathsf N}(3/d)^{2p_V}\mathsf M$.

Made once for each of the $\frac12n(n-1)$ pairs, with the total power at most $p_Vn$ and radii of the order $d(Y)/n$, this gives $\|\mathsf G_l(x;e)\|\le\mathsf K_n(\sqrt2C_V(3/d(Y))^{p_V})^n\prod_a|e_a|$ at all points of pairwise different heights, and the continuation of $\mathsf G_l$ with this bound across the sets of equal heights, where the entries are apart in a horizontal direction. Then the integrand is bounded, and $\sigma$ charges no set of points with two equal heights: by Proposition~\ref{sf-spec}(b) the unordered remainder of a cut in $x_3$ tends to $0$, and at ordered tuples the class is the product of the operators of (a) applied to $[F]$, by an induction on the number of the entries; so $[Y\cdot\tau_{le_3}F]$ is the integral, over the points of pairwise different heights, a set of full measure, of $\mathsf G_l(x;h_1(x_1),\dots,h_n(x_n))[F]$, which gives the $L^1$ bound.

So $s_n(x;e)=\langle\Omega,\mathsf G_l(x;e)\Omega\rangle$ is continuous at all non-coinciding points, and $\sigma(X)=\langle\Omega,[X]\rangle=\int s_n$ after a translation of $X$ into a slab.
\end{proof}

\subsection{Proof of Theorems~\ref{sf-tori}--\ref{sf-coin} and of Corollary~\ref{sf-ode}}

\emph{Step 1: the tori, \textup{(a)}--\textup{(c)} and \textup{(f)}.} By \eqref{sf-lat}, $O^T_K(H)=\sum_\lambda\mathcal O(H_{\lambda\lambda}1_\lambda)+\sum_{\lambda<\lambda'}O^{\rm od}_{\lambda\lambda'}(H_{\lambda\lambda'})$, $O^{\rm od}_{\lambda\lambda'}(\phi)=-a_K^{-1}\sum_{b\parallel\mu}\phi(x_b)q_b$ for $\mu$ the third direction, with weights supported in $\mathrm{supp}\,H$, of sup norm $\le|H|$ and Lipschitz number $\le\mathrm{Lip}\,H$; so a cumulant of the $O^T_K(H_a)$ is a sum of $6^n$ cumulants of entries. Lemma~\ref{sf-cut}(a) gives \eqref{sf-rate}, and summing over $K'\ge K$ the limit with its rate; the bound of $\mathcal S_n$ is Lemma~\ref{sf-sym}(b) at $\delta=\frac d2$; the formula for the centred moments holds at every $K$.

With $\bar O^T_K$ the cumulant at fixed $K$ is a polynomial in finitely many cell means of the $H_a$, each continuous under bounded pointwise convergence, and the rate is uniform on balls of the sup norm: so the limit is continuous under bounded pointwise convergence, which makes it countably additive in each argument, and it vanishes when all the cell means vanish, in particular on null sets; the theorem of Radon--Nikodym gives the density. For (b), Lemma~\ref{sf-sym}(a) gives the invariance under $a_K\mathbb Z^3\rtimes W_3$ at every cutoff; a real and a lattice translation of a weight differ by a weight of sup norm $\le a_K\mathrm{Lip}\,H$, and the bounds are in sup norms. As $L$ is odd, $\{x_\mu=c\}$ with $c\in a_{K_0}(\mathbb Z+\frac12)$ is a bond mirror of all the lattices with $K\ge K_0$, at which Wilson's measure is reflection positive (Corollary~\YMMref{85.3}{679}, Proposition~\YMMref{91.4}{700}) and $\theta O^T_K(H)=O^T_K(\theta H)$; such $c$ are dense, and the reflection at another $c$ is conjugate to one of these by a translation.

At $H_a=f_a\mathbf 1$, $O^T_K(f_a\mathbf 1)=\mathcal O(f_a)$. For (c): every class insertion is, at the rate, its standard form, built from the corner fields with the loops $l_\lambda=(\kappa,\kappa')$, $\kappa<\kappa'$ the two directions other than $\lambda$, at the sites. The shift $\xi$ costs a site form of norm $\le\frac12a_K\mathrm{Lip}\,\phi$, bounded by Lemma~\ref{sf-sym}(b); loops in one plane are exchanged, with the sign $s(l)s(l')$, inside the rate of Lemma~\ref{sf-cut}(a), the difference being tied at the depth $\sigma$ with the defect of Lemma~\ref{sf-cut}(c); and $a_p=\frac12Q_x(l_\lambda,l_\lambda)+\frac12a_p^2$, whose second term is of fourth order at the backgrounds. The standard form depends on the symbol alone; $O^T_K(\mathrm{Sym}\,\mathcal E)$ is a class insertion with the same symbol, and (a) gives the limit; with the same weights at the same sites no shift and no Lipschitz number occur.

For (f), take the numbers $C_{\rm bf}$, $K_{\rm bf}$ of Definition~\textup{\ref{IV-62-62}}(5) at a compact interval $J\subset I(g)$ that holds $\mathsf b$ in its interior: by Corollary~\textup{\ref{IV-62-66}}(a) they read the label only through $J$, and the comparison of the block fields of two successive cutoffs holds with them at every label of $J$. The rate of (a) reads the label at these places alone: $C_{\rm bf}$ and $K_{\rm bf}$; the condition $L^\sigma\mathsf b\ge\Theta_\star$, which is monotone in the label; the radius of the backgrounds, which the calibration conditions in equation~\YMMref{(2148)}{794} bound at the levels where $L^\sigma\mathsf b\ge\Theta_\star$; and the data of the depth $s$ of $(K,L^K\mathsf b)$, whose numbers are uniform on the compact $J$.

So, by Lemma~\ref{sf-cut}(b), one exponent $\mathsf p$, one $C^{\rm od}$, one $C^{\rm av}$ and one threshold $K_1$ serve all labels of $J$, for the entries and for the cell means. At a fixed $K$, $\mathcal S_{n;K}$ is continuous in the label, an integral over a compact group against a density continuous in $\beta=L^K\mathsf b$; so $\mathcal S^{\ell,\cdot}_n$ is a uniform limit on $J$ of functions continuous in the label, hence continuous on $J$; with $\mathsf b_K=a_K\beta_K\in J$ for large $K$,
\[|\mathcal S_{n;K}(\mathsf b_K)-\mathcal S_n(\mathsf b)|\le|\mathcal S_{n;K}(\mathsf b_K)-\mathcal S_n(\mathsf b_K)|+|\mathcal S_n(\mathsf b_K)-\mathcal S_n(\mathsf b)|\to0,\]
as in Corollary~\YMMref{117.1}{812}; with $\bar O^T_K$ the same.

\emph{Step 2: every label and side; all of \textup{(H)}.} Lemma~\ref{sf-dil}(a) gives
\[\langle O^T_K(H_1);\dots;O^T_K(H_n)\rangle_{K,\ell_*,L^K\mathsf b_*}=L^{zn}\langle O^T_{K-z}(H^{(z)}_1);\dots;O^T_{K-z}(H^{(z)}_n)\rangle_{K-z,L^z\ell_*,L^{K-z}(L^z\mathsf b_*)},\]
which carries Theorem~\ref{sf-tori} to every label and side, with $L^zd$ for $d$; with $g'=(L^z\mathsf b/2)^{-1/2}$ the label $L^z\mathsf b$ is an inner point of $I(g')$, which carries (f). Lemmas~\ref{sf-one} and~\ref{sf-stop}, Proposition~\ref{sf-count} for its limits and its growth, and Steps~6, 7 and 10--13 are first carried out at couplings $g\le\gamma^{\rm sd}_W$; Lemma~\ref{sf-dil}(b) carries the data, the states, the limits, the Hilbert spaces, the semigroups, \eqref{sf-slab} and the decomposition to all of (H), and $\sigma=L^{nz}\sigma_z(X^{(z)})$ carries the bounds of Proposition~\ref{sf-count}(b), (c) and its growth, the powers of $L^z$ entering $C_G$ and $C$; the count of Proposition~\ref{sf-count}(a) at a cutoff holds on all of (H) with the numbers set in its statement. In particular one may take $d_0(g)=L^{-z}\min_id_0(g'_i)$, with the error $\max_i\mathsf e^\infty_{g'_i}(L^zd)$ at $d$, for $z$ large enough at $g$.

\emph{Step 3: two points on one torus.} By Lemma~\ref{sf-depth}(a) and $|K_3(u)_{\lambda\lambda'}|\le(2\pi|u|^3)^{-1}$, $|u|\ge d$ on the supports, the Wick value is a Riemann sum of $\mathcal W$: with $a_K^3\sum_{p,p'}|\mathfrak t_H(p,p')|\le45d^3|H|$, by the row sums $2\sum_{p'}|\mathfrak t_H(p,p')|\le3|H|$ of the kernel of $O^T_K(H)$ (Step~14(a)) and the counts of the plaquettes and the bonds in a ball,
\[\Big|6\sum\mathfrak t_1\mathfrak t_2CC-\mathcal W(H_1,H_2)\Big|\le12150\,E_1\big(\tfrac1\pi+E_1\big)|H_1||H_2|+\mathrm{rm}_K(H_1,H_2),\]
where $\mathrm{rm}_K$ is the difference between a Riemann sum over the anchors, that is the barycentres, the midpoints or the sites of the cells, and the integral $\mathcal W$; it depends only on $K$, $d$ and the weights, and it tends to $0$, by the uniform continuity of $K_3$ on the two balls, for continuous weights and for cell means. With $|\mathcal W|\le54|H_1||H_2|$ (Lemma~\ref{sf-freev}(a)) and the bounds on $v^2$, \eqref{sf-two} follows at finite cutoff with
\[\mathsf e^T_{g,\ell}(d)=\tfrac98\big[3600\,\varepsilon^e_{g,\ell}(s)+12150\,E^\infty\big(\tfrac1\pi+E^\infty\big)\big]+162\bar Bg^2L^{2-s},\qquad E^\infty=E_1-C_{\rm lat}(a_K/d)^2,\]
and $\frac98\cdot12150\,[E_1(\frac1\pi+E_1)-E^\infty(\frac1\pi+E^\infty)]\le c_K$; the term $162\bar Bg^2L^{2-s}$ is $54$ times $3\bar Bg^2L^{2-s}$. As $K\to\infty$, (a) gives \eqref{sf-two}; $\mathsf e^T\to0$ like $(\log\frac1d)^{-2\mathfrak p_0}$ because $Y_s^{2\mathfrak p_0}\varepsilon^e_{g,\ell}(s)$ is bounded and $t=L^sd\ge C_W(s)\to\infty$. The class: (c) at $n=2$. The bound on $s_1(g,\ell)$ and the independence of $\ell$: Lemma~\ref{sf-depth}(b). (i): $\mathcal W$ is invariant (Lemma~\ref{sf-freev}(a)) and $|rH|\le3|H|$, so the law at $(H_1,H_2)$ and at $(rH_1,rH_2)$ gives $(1+9)\mathsf e^T$. (ii): Lemma~\ref{sf-freev}(b) and the law at the pair $(\theta(\phi E),\phi E)$, whose centres are $3d$ apart, with $|\phi E|^2\le\mathrm{tr}\,E^2$. (iii): Lemma~\ref{sf-depth}(c) at $c\in a_{K_0}(\mathbb Z+\frac12)$ and the limit; a general $c$ by (b).

\emph{Step 4: conservation on the torus.} Cut $\eta$ by a partition of unity into pieces of diameter $<1$; by multilinearity it suffices to treat one piece. In its slot, by \eqref{sf-ward}, $\mathsf bO^T_K(H^{\rm cons}[\eta])=\mathrm{St}_K[\eta]-\Delta_K[\eta]$, and $\mathrm{St}_K[\eta]$ may be replaced by $-\mathrm{Rm}_K[\eta]$, the null terms $\mathcal N_\eta$, $\mathcal N'_\eta$ vanishing against the far entries. This leaves the remainders of Proposition~\ref{sf-stress}(a): the cubic one $\mathcal S[\mathsf b\eta]$, the five local forms of $J_\eta$, the form $\frac32a_K\mathcal O(w_\eta)$ and the forms of $\Delta_K[\eta]$, each with a conditional mean, given the block field of depth $\phi K$, that is constant on $T_\ell$ up to $K^{\mathsf p}L^{-(\frac1{28}-3\phi)K}\|\eta\|_{C^2}$: by Proposition~\ref{sf-stress}(c) for the cubic one; by power counting at the backgrounds of the level $K-\phi K$, as in the proof of Proposition~\ref{sf-stress}(b), for the local forms, those of the order $2$ (the three of $J_\eta$, with weights at most $(3+5\mathsf b)\|\eta\|_{C^2}$, the form $\frac32a_K\mathcal O(w_\eta)$ and the two forms of $\Delta_K[\eta]$ with weights at most $60\mathsf ba_K\sup|\nabla^2\eta|$) having oscillations at most $CK^{2\mathfrak p_0}L^{4\phi K-K}\|\eta\|_{C^2}$, and the two of the orders $3$ and $4$ of $J_\eta$, with weights at most $L^K\mathsf b\sup|\nabla\eta|$, oscillations at most $CK^{3\mathfrak p_0}L^{4\phi K-\frac K2}\|\eta\|_{C^2}$; and by the cut at the level $K-\lceil K/14\rceil$ for the form of vanishing symbol of $\Delta_K[\eta]$ alone.

A joint cumulant of entries far apart is at most the product of the oscillations of their conditional means at that level, a spectator costing $L^{4\phi K}$; at $\phi=1/(56(4q-1))$ the exponent is $-\frac1{28}+(4q-1)\phi=-\frac1{56}$. (a) gives the limit, and the moments vanish, being sums of products of cumulants of which one holds the slot of $H^{\rm cons}[\eta]$. The estimate of Proposition~\ref{sf-stress}(c) is uniform on the interval of labels, which gives the uniformity in the label; $\bar O^T_K$ as in Step~1; a change by forms with weights $\le36a_K\|\eta\|_{C^2}$ costs, by Lemma~\ref{sf-sym}(b), $a_K$ times a number independent of $K$. At $H_a=h_a\mathbf 1$ the spectators are action densities.

The rate $L^{-K/36}$ also holds, by the argument of Step~8 on $T_\ell$.

\emph{Step 5: states and $\mathbb R^3$.} Lemma~\ref{sf-chess}(a), (b) bound products in one cell in every state, centred by any state. For Theorem~\ref{sf-gibbs}(a) cut the weights by a fine grid of cubes of side $\frac1{4n}$, into at most $(4n\varrho_0+2)^3$ pieces each; for each product of pieces Lemma~\ref{sf-chess}(c), at $h=a_K$, gives a board of side $1$ whose walls avoid the pieces, with the margin $\delta=\min\{\delta_0,\frac1{8n}\}$, the threshold on $K$ of Theorem~\ref{sf-gibbs}(a) including $2na_K<1$, so that
\[\Big|\nu\Big(\prod_a\big(O^T_K(H_a)-c(H_a)\big)\Big)\Big|\le\sum_{\rm pieces}\ \prod_{\rm cells}e^{G_g}\|D_{\rm cell}\|_{{\rm ch},\ell_g}\le C^{\rm loc}_n\prod_a|H_a|,\]

\[
C^{\rm loc}_m=(4n\varrho_0+2)^{3m}(\max_{r\le n}\max\{1,e^{G_g}C^\flat_r\})^m,
\]
with $c(H)$ the mean on $T_{\ell_g}$, which is $e_{\ell_g}a_K^{-1}\sum_pH_{\lambda(p)\lambda(p)}(x(p))$ by Lemma~\ref{sf-sym}(a); the centring by $\nu'$ gives $C^{\rm lc}_n=\sum_k\binom nk(C^{\rm loc}_1)^{n-k}C^{\rm loc}_k$. The count in $n$ is Proposition~\ref{sf-count}(a). The bounds are in sup norms and independent of the positions, so the $\nu_j(A^{\nu_j}_{K_j}(\cdot))$ are equicontinuous on continuous weights with supports in fixed open sets; a diagonal extraction on a countable family of tuples, dense in the sup norm, gives the limits of (c) with their labels, and the same equicontinuity shows that convergence at the tuples of $C^1$ weights gives convergence on all of $\mathfrak S(\mathbb R^3)$.

Lemma~\ref{sf-sym}(a), with the invariance of every periodic Gibbs state under $a_K\mathbb Z^3\rtimes W_3$ (Proposition~\textup{\YMMref{103.4}{758}}, periodic in the sense of Section~\ref{sec:state-limit-conventions}), gives the invariance under $W_3$ and the lattice translations at every cutoff, and the positivity of periodic Gibbs states at the lattice mirrors of the nine families (Proposition~\textup{\YMMref{103.4}{758}}), the linear parts of these reflections lying in $W_3$, so that $\theta A^\nu_K(X)=A^\nu_K(\theta X)$, gives the positivity there; both pass, by the equicontinuity of the bounds of Theorem~\ref{sf-gibbs}(a), to real translations and to all positions of the mirrors, a real $v$ and a real $c$ being limits $v_j\to v$, $c_j\to c$ of lattice ones; the labels lie in $[g^{-2},L(g^{-2}+\bar B)]$ by first passage; the trace part is the limit of the action density along the same sequence.

On growing tori the measures obey the same bounds, by Lemma~\ref{sf-chess}(a) for torus measures, are invariant under the lattice isometries and positive at the coordinate bond mirrors (Proposition~\YMMref{91.4}{700}).

\emph{Step 6: one error and one threshold.} Let first $g\le\gamma^{\rm sd}_W$ and $d_0(g)=L^{-s_{16}(g)}$. Lemma~\ref{sf-one} reduces every periodic Gibbs state and every torus of $\mathfrak L(g)$ to $T_{\ell_g}$, where Step~3 holds with $\omega'_g(s)$ for $\varepsilon^e_{g,\ell}(s)$:
\[\mathsf e^\infty_g(d)=\tfrac98\Big\{3600\big[\omega'_g(s)+\delta_g(s)+e^{2G_g(\ell_g/2)^3}\mu_g(s)^2\big]+12150\,E^\infty\big(\tfrac1\pi+E^\infty\big)\Big\}+162\bar Bg^2L^{2-s},\]
which depends on $g$ alone; $\delta_g(s)\le C(g)Y_s^{-\mathfrak p_0/4}$ gives the rate $(\log\frac1d)^{-\mathfrak p_0/4}$. The tori at the labels $\mathsf b\in I(g)$ are among the measures considered, since $(K,L^K\mathsf b)\in\mathfrak d(g)$; the limits follow along their sequences. The label is determined: by Lemma~\ref{sf-freev}(a), (b) the free value $\mathcal W(\theta(\phi E),\phi E)$ of a mirror pair is independent of $d$ and $\ge\mathsf w_0\mathrm{tr}\,E^2$, and the law gives $\mathsf b^{-2}=\lim_{d\to0}\sigma(\theta(\phi E),\phi E)/\mathcal W(\theta(\phi E),\phi E)$. The lower bound and the near invariance follow as in Step~3. All of (H): Step~2.

\emph{Step 7: the power.} Let first $g\le\gamma^{\rm sd}_W$. For $i\ge i_E(g,\delta)$ and $K\ge K_E(i;g,\delta)$, let $c_i$ be the mean of $O^T_{K-i}(H^{(i)})$ on the least torus at $g_{(i)}$ in Lemma~\ref{sf-stop}. The thresholds are those of Theorem~\ref{sf-gibbs}(a), provided in this coupling range by that lemma. Lemmas~\ref{sf-dil}, \ref{sf-stop} and~\ref{sf-chess}(a), at the coupling $g_{(i)}$, give
\[\|O^T_K(H)-L^ic_i\|_{{\rm ch},\ell}=L^i\|O^T_{K-i}(H^{(i)})-c_i\|_{{\rm ch},L^i\ell}\le L^ie^{G_{g_{(i)}}}C_\delta Y_i^{q_E}L^i(L^i\mathsf b)^{-1}|H|\le C'_\delta Y_i^{q_E}L^{\frac{11}2i}\mathsf b^{-1}|H|,\]
the cells having side $1$ after the dilation and $e^{G_{g_{(i)}}}\le e^{G_g+1/8}L^{9i/2}$; this is $\le C_TL^{ip_T}\mathsf b^{-1}|H|$ for $p_T>\frac{11}2$, $Y_i$ being linear in $i$. The mean on another torus or in a periodic Gibbs state differs from $L^ic_i$ by the same bound, by Lemma~\ref{sf-chess}(a) at one cell after the same dilation. For the finitely many $i<i_E(g,\delta)$: Lemma~\ref{sf-chess}(b) on the least torus, with $N=(\ell_gL^i)^3$ copies and Lemma~\ref{sf-sym}(b) at the orders up to $N$. All of (H): Lemma~\ref{sf-dil}(b), with the thresholds of Theorem~\ref{sf-gibbs}(a), which on the strip $\gamma^{\rm sd}_W<g\le\gamma^{\rm ref}_W$ are set through $z_\sharp$ and $g_-$ and at which the display holds in the same form by the stopped datum of the proof of Proposition~\ref{sf-count}(a); the small $i$ again by Lemma~\ref{sf-chess}(b).

\emph{Step 8: conservation in every state.} Let $0<g\le\gamma$. Dilate by $L^z$, with $O^T_K(H)=L^zO^T_{K-z}(H^{(z)})$ and the states and torus measures corresponding (Lemma~\ref{sf-dil}), to the two couplings $g'_1=(L^zg^{-2}-\bar B)^{-1/2}$ and $g'_2=(L^{z+1}g^{-2}-3\bar B)^{-1/2}$ of Lemma~\ref{sf-dil}(b). By $LB_\sharp=(L-1)\bar B$, $I(g'_1)=(L^zg^{-2},L^{z+1}g^{-2}-2\bar B]$ and $I(g'_2)=(L^{z+1}g^{-2}-2\bar B,L^{z+2}g^{-2}-(2L+2)\bar B]$, and these two intervals cover the labels $L^z\mathsf b$, $\mathsf b\in(g^{-2},L(g^{-2}+\bar B)]$, for every odd $L$ and $z\ge1$, as $g^{-2}\ge\gamma^{-2}>\bar B$; so every $(K,\beta)\in\mathfrak d(g)$ with $K>z$ has $(K-z,\beta)\in\mathfrak d(g'_\iota)$ for one $\iota$, by the label $L^z\mathsf b\in I(g'_\iota)$, and $L^z\ell\in\mathfrak L(g'_\iota)$ for $z$ large.

This arithmetic reads $g^{-2}>\bar B$ alone and serves every $0<g\le\gamma$. Let $z$ be the least integer with $g'_1,g'_2\le\gamma^{\rm sd}_W$ and $\ell_{g'_\iota}\le L^zd$, and write $g'=g'_\iota$; such a $z$ exists, as $\ell_{g'}$ grows only like a power of $z$, that is of $\log g'^{-2}$. Cut $\eta$ and the spectators into pieces of diameter $\le\ell_{g'}/(8n)$ and shift the board of side $\ell_{g'}/2$, whose seminorms are taken on $T_{\ell_{g'}}$ with its eight cells, on the grid $h\mathbb Z^3$, $h=L^{-\lceil K'/400\rceil}$, $K'=K-z$ the cutoff after the dilation, as Proposition~\ref{sf-stress}(b) asks of the offset of its board, so that each piece lies in a cell of its own, $\ge\ell_{g'}/(16n)$ from its walls (Lemma~\ref{sf-chess}(c): $h\in a_{K'}\mathbb N$, $\ell_{g'}/(2h)$ is an integer and $2nh<\ell_{g'}/2$ for large $K$).

The piece of $\eta$ has a seminorm $\le CK^{\mathsf p}L^{-K/36}\|\eta\|_{C^2}$ by Proposition~\ref{sf-stress}(b); the others are bounded by Lemma~\ref{sf-sym}(b) at the orders $2,\dots,8$, their eight mirror images being separated insertions on $T_{\ell_{g'}}$; Lemma~\ref{sf-chess}(a) transfers the product to every periodic Gibbs state, in particular $|\nu(A^\nu[F_1]A^\nu[F_2])|\le2e^{2G}\|F_1-c_1\|_{{\rm ch}}\|F_2-c_2\|_{{\rm ch}}$ for two cells, $G$ the gap of one cell; summing over the pieces and the $2^n$ terms of the centring and dilating back gives the bound, and the joint cumulant follows by the M\"obius formula. The limit follows along the sequence. The rest of this step is made after the same dilation, at the coupling $g'\le\gamma^{\rm sd}_W$ and the cutoff $K-z$, written $K$ below, where Lemma~\ref{sf-chess} applies; the powers of $L^z$ enter the constants.

For the rate $L^{-cn_1K}$, let $u=\lceil\phi K\rceil$, $\phi=(\frac1{28}-c)/(4n-1)$, and $K$ so large that $dL^u\ge K^2$; shift a board on a fine grid so that each piece lies inside one cell, several pieces being allowed in one cell. In one cell the chessboard seminorm of a product of centred slots is at most a number times the product of the oscillations of their conditional means given the block field of depth $u$: at most $CK^{\mathsf p}L^{3u-K/28}\|\eta_a\|_{C^2}$ for a slot that holds a conservation or stress insertion, as in Step~4 (Proposition~\ref{sf-stress}(c) and power counting), and at most $CK^{\mathsf p}L^{4u}|H_a|$ for a tensor entry; Lemma~\ref{sf-chess}(a) carries the product of the cells to every periodic Gibbs state, and $n_1(3u-\frac K{28})+(n-n_1)4u\le-cn_1K+4n$, since $(4n-n_1)\phi\le\frac1{28}-c$ and $(n_1-1)c\le\frac{n_1-1}{28}$.

At two points the same transfer gives the rates directly: $(L^{-K/28})^2$ for two stress insertions, as the mean of the copies turns $\mathrm{St}_K$ into $-\mathrm{Rm}_K$; $(L^{-K/36})^2\le L^{-K/20}$ for two conservation insertions, by Proposition~\ref{sf-stress}(b) for each; $L^{-K/36}$ times the crude $K^{2\mathfrak p_0}L^{K/100}$ of a tensor entry, which is $\le L^{-K/60}$ up to powers of $K$, for $H^{\rm cons}[\eta]$ or $\Delta_K[\eta]$ against $O^T_K(H')$; and the same for $\Delta_K[\eta]$ against $\mathrm{St}_K[\eta']$ or $\Delta_K[\eta']$.

\emph{Step 9: narrow bumps; the kernels of the plane fields.} In this step, as in Lemmas~\ref{sf-jet} and~\ref{sf-tr}, in Theorem~\ref{sf-gibbs}(e), in Theorem~\ref{sf-coin}(c), (d) and in Step~14(c), (d), the distance of two sets is the Euclidean one, $\mathrm{dist}(A,B)=\inf_{x\in A,y\in B}|x-y|$, and the separations of sets are read with it; the factor $1/\sqrt3$ between it and the sup distance is paid inside the proofs of these lemmas, and the numbers $\frac78d'$, $\frac67d'$, $\frac32t$, $\frac52t$, $3t$, $(2r/d')^6$ and $E_2(g,\frac12d')$ below are those of this reading. Let first $g$ be below the threshold of Lemma~\ref{sf-jet}, $\ell\in\mathfrak L(g)$, and let the two bumps be read as mixed sources carried by $B(x_a,r)$, $|x_1-x_2|=d'$, $16r\le d'\le\frac\ell8$; their balls are $\ge\frac78d'$ apart.

For $d'$ below a threshold depending on $g$, Lemma~\ref{sf-jet}(a), (b) at the separation $d=\frac67d'$, the distance $d'=\frac76d$ of the centres lying between $d$ and $18d$, as (a) asks, and below $8d$, as (b) asks: the cumulant of the two bumps is its free lattice value up to $\varepsilon_2(g,d)(r/d)^6\mathsf b^{-2}|\mathfrak s_1||\mathfrak s_2|$ plus the part of the large fields, which on $T_\ell$ is at most $C(g,\ell)(r/d)^6d^{3/2}\mathsf b^{-2}|\mathfrak s_1||\mathfrak s_2|$; and $(r/d)^6=(\frac76)^6(r/d')^6$. Distances $d'$ above that threshold, up to $\frac\ell8$: the argument of Lemma~\ref{sf-jet}(a), (b) at a fixed small separation $d=t(g)$, a threshold depending on $g$, asks no ceiling on the actual distance of the two balls but that they lie in a ball of radius $<\frac14\ell$, and gives a tolerance that is a finite function of that distance; where the bumps are wider than the widths that the law admits at the separation $d$, both members being bilinear, the sources are cut into pieces carried by the cubes of a grid of a mesh depending on $g$, the number of pairs of pieces being at most a basic multiple of the sixth power of the width over the mesh, and the sum of the laws of the pairs of pieces costs, against $(r/d')^6$, the finite factor $(d'/d)^6$ times a basic number.

So $\varepsilon^{\rm cl}(g,\ell,\cdot)$ is a finite function on $(0,\frac\ell8]$, and $\varepsilon^{\rm cl}(g,\ell,d')\le C(g,\ell)(\log\frac1{d'})^{19\mathfrak p_0}d'^{1/2}$ for small $d'$, the large fields adding $C(g,\ell)d'^{3/2}$. For Lipschitz bumps the lattice sums against the free kernel, $\sum_{\lambda(p)=\lambda,\lambda(q)=\mu}\phi_1(x(p))\phi_2(x(q))C^{\rm free}(p,q)^2$, differ from $\iint\phi_1\phi_2[K_3]_{\lambda\mu}^2$ by at most $C_\flat(\frac{a_K}{\vartheta r}+\frac{d'^3}{\ell^3})(\frac r{d'})^6|\phi_1||\phi_2|$, by the comparison of the lattice kernel with the free kernel (equations~\YMMref{(2245)--(2246)}{824}), the term $a_K/(\vartheta r)$ being the Lipschitz defect of a Riemann sum of a bump of width $r$ and $(d'/\ell)^3$ the images on the torus; the off-diagonal bond forms likewise with $C'_\flat$, the factor $8$ coming from the four corner forms of a bond and the two directions.

Couplings $0<g\le\gamma$ above the threshold of Lemma~\ref{sf-jet} are reached by the dilation of Step~8, whose arithmetic reads $g^{-2}>\bar B$ alone and whose two couplings lie below any basic threshold once $L^zg^{-2}$ is large, at $L^zd'$ and $L^z\ell$, both members of the law being multiplied by $L^{-2z}$; this defines $\varepsilon^{\rm cl}$ and $K_\circ(g)$. The tolerance depends on $d'$ alone: there is no gain as $r\to0$ at fixed $d'$. In every periodic Gibbs state, for $g$ below a basic threshold: put $t=\frac12d'$; bumps of radius $r$ in the balls of radius $\frac18d'$ about two points at the distance $d'$ have supports $\ge\frac32t$ apart, at centres $\le\frac52t$ apart, in one ball of radius $3t\le\frac18\ell_g-1$ when $d'$ is below a threshold depending on $g$.

On $T_{\ell_g}$, Lemma~\ref{sf-jet}(a) at $(\rho,d')=(r,t)$ gives the law of the plaquette forms with $\varepsilon_2(g,t)+C(g)t^{3/2}$ in units of $(r/t)^6\mathsf b^{-2}|\phi_1||\phi_2|$; Lemma~\ref{sf-tr}(b) at $n=2$, $d_\times=t$, carries it to every periodic Gibbs state with the further $C^{\rm tr}_2(g)(r/t)^6t^{3/2}|\phi_1||\phi_2|$, at most $C(g)t^{3/2}$ in these units, as $\mathsf b\le L(g^{-2}+\bar B)$; and the Riemann sums on $T_{\ell_g}$ add $\frac32C_\flat(\frac{a_K}{\vartheta r}+(\frac t{\ell_g})^3)$ in the same units. This is the law in every state with $E_2(g,t)=\varepsilon_2(g,t)+C(g)t^{3/2}+\frac32C_\flat(t/\ell_g)^3$ at $t=\frac12d'$ for $\varepsilon^{\rm cl}$, and $(r/t)^6=(2r/d')^6$; $E_2(g,t)\le C(g)(\log\frac1t)^{19\mathfrak p_0}t^{1/2}$.

First take the chosen continuum subsequential limit at each fixed positive width; then let the widths tend to $0$. The estimates are uniform in the underlying states. The limits have bounded Borel kernels within $(\frac3\pi)^2\varepsilon_\sigma(g,|x-y|)\mathsf b^{-2}|x-y|^{-6}$ of the free ones almost everywhere, by Lebesgue's differentiation theorem, functions of $x-y$ alone by the translations; for a limit at $g$ the kernels are continuous (Theorem~\ref{sf-wight}(d)), so the bound holds at every pair; the sum over $\lambda,\mu$ is the scalar kernel, the trace part being the scalar limit.

\emph{Step 10: operators.} By Theorem~\ref{sf-gibbs}(c), Proposition~\ref{sf-spec} applies to $\sigma$: (a) gives Theorem~\ref{sf-wight}(a), the circular cone from the nine families, and the scalar subspace through the trace part, the map $[F]_0\mapsto[F\mathbf 1]$ being isometric as $\Theta_{e_3}\mathbf 1\Theta_{e_3}=\mathbf 1$ and commuting with the translations. Multiple reflections with domination, Lemma~\ref{sf-chess}(d), bound the slab operators by the chessboard seminorms of the products of the centred insertions, which Lemma~\ref{sf-chess}(b) bounds by $C^{\rm sl}_n$, the number $C^G_n$ at the margin $\frac\delta2$ and the least side $\ge2\mathsf l$; the limit is taken with a moving mirror, by equicontinuity; the adjoint follows from $\theta_{e_3,\mathsf l/2}=\tau_{\mathsf le_3}\theta$. Proposition~\ref{sf-spec}(c) gives the point operators. By Theorem~\ref{sf-gibbs}(b) at the margin $\frac1{16}$ and Lemma~\ref{sf-chess}(d), a plate in a cell of side $a_i$ has $\|\mathsf B_{a_i}((h))\|\le2C_TL^{ip_T}\mathsf b^{-1}|h|$; Proposition~\ref{sf-spec}(c) at cubes of half side $\rho$ of the order of $a_i$, with $a_i$ of the order of $d-d'$, gives
\[\|\mathsf B_{2d}((\tau_{de_3}h))\|\le4\cdot32^3a_i^{-3}\,2C_TL^{ip_T}\mathsf b^{-1}\|h\|_{L^1}\le C_p\min\{1,d-d'\}^{-p}\mathsf b^{-1}\|h\|_{L^1},\qquad p=3+p_T>\tfrac{17}2,\]
and so \eqref{sf-slab} with $C_V=2^{p_V}C_{p_V}\mathsf b^{-1}\le2^{p_V}C_{p_V}g^2$, one number for all limits at $g$; and $C_{\rm pt}(\eta)\le12\cdot80^3e^{2G_g}C_T(\frac{9L}{16})^{3+p_T}\eta^{-3-p_T}$ for $0<\eta\le1$.

\emph{Step 11: invariance; \textup{(h)}.} A limit has the hypotheses of Proposition~\ref{sf-rot}: (i) by Theorem~\ref{sf-gibbs}(f); (ii) by Theorem~\ref{sf-gibbs}(d) for $\mathsf d<d_0(g)$, the free value of a mirror pair being scale invariant, and by Theorem~\ref{sf-gibbs}(a) for larger $\mathsf d$; (iii) with $a=2$, as follows. The free value of a plate $f=\alpha(x_3,x_1)\varphi(x_2)$ against its mirror image is of the order $t^{-1}$, a tube of cross-section about $t^2$ at the height about $t$, and so is $\sigma(\theta(f\mathbf 1),f\mathbf 1)$ wherever the kernel law of Theorem~\ref{sf-gibbs}(e) holds; so (iii) does not hold with $a=0$. Cut $f$ by a smooth partition of unity into at most $C_\varphi/t$ Lipschitz pieces $f_k$ of diameter $\le\frac t4$, $C_\varphi$ depending on the length of $\mathrm{supp}\,\varphi$; the centre $x_k$ of $f_k$ has $\frac78t\le x_{k,3}\le\frac{17}8t$, so $(\theta(f_k\mathbf 1),f_k\mathbf 1)$ is a mirror pair at every scale $d\in[\frac12x_{k,3},\frac23x_{k,3}]$, and for $t$ below a threshold depending on $g$ the law of Theorem~\ref{sf-gibbs}(d) with Lemma~\ref{sf-freev}(a) gives $\|[(f_k\mathbf 1)]\|^2=\sigma(\theta(f_k\mathbf 1),f_k\mathbf 1)\le(54+\mathsf e^\infty_g(d))\mathsf b^{-2}|f|^2$.

The triangle inequality in $\mathcal H(\sigma)$ gives $\sigma(\theta(f\mathbf 1),f\mathbf 1)=\|\sum_k[(f_k\mathbf 1)]\|^2\le(C_\varphi/t)^2(54+\mathsf e^\infty_g(d))\mathsf b^{-2}|f|^2$, and for the larger $t\le1$ Theorem~\ref{sf-gibbs}(a) bounds the left member by a number depending on $g$ and $\varphi$; so (iii) holds with $a=2$, and in the harmonic argument $4+a=6<8$. (iv) by Step~10. So $\sigma$ is invariant under $O(3)$ and conserves. Positivity at the plane with the normal $n$ is positivity at a coordinate plane with a normal $m$, conjugated by an $R$ with $Rm=n$. For (h): $Z_z$ commutes with $W_3$ and with the reflections of the nine families, $|Z_zH|\le\max\{1,|z|\}|H|$, and $[F]_z\mapsto[Z_zF]$ is unitary and intertwines the semigroup and the translations, so the cones and the slab and point operators carry over.

The rotation $R$ by $\frac\pi4$ about $e_3$ has $Re_{12}R^t=E_{22}-E_{11}\in V_E$, while $\sigma(\phi_1e_{12},\phi_2e_{12})\ne0$ for a suitable pair by the lower bound of Theorem~\ref{sf-gibbs}(d); so invariance of $\sigma^z$ would force $z^2\sigma(\phi_1e_{12},\phi_2e_{12})=\sigma^z(R(\phi_1e_{12}),R(\phi_2e_{12}))=\sigma(\phi_1e_{12},\phi_2e_{12})$. The law of Theorem~\ref{sf-gibbs}(d) for $\sigma^z$ at off-diagonal pairs would give $|z^2-1||\mathcal W|\le(1+z^2)\mathsf e^\infty_g(d)|H_1||H_2|$ for all small $d$, which forces $z^2=1$ by the lower bound of $\mathcal W$ at mirror pairs. A positive form is Hermitian, so positivity at a plane gives invariance under its reflection, labels turned; the invariance group of a functional with the properties of Theorem~\ref{sf-gibbs}(c) is a closed subgroup of $O(3)$ containing $W_3$, and a closed subgroup of $O(3)$ that contains $W_3$ and one further element is $O(3)$, since the rotations of $W_3$ form a maximal finite subgroup of $SO(3)$ and an infinite closed subgroup containing them contains the rotations about several axes; so positivity at one further plane gives invariance, and conjugation gives positivity at every plane.

So $\sigma^z$, $z^2\ne1$, positive at the nine families and not invariant, is positive exactly there.

\emph{Step 12: the field; real time.} By Steps~10 and~11 and Proposition~\ref{sf-spec}(b), the restriction of $\sigma$ to $C^1$ entries has the hypotheses of Proposition~\ref{sf-rec}, with $p_V$ and $C_V$, and the Hilbert spaces of $\sigma$ and of this restriction are one, the classes of $C^1$ tuples being dense. Proposition~\ref{sf-rec} gives Theorem~\ref{sf-wight}(e) and the kernels of (d), which represent $\sigma$ at all continuous entries by the convergence on $\mathfrak S(\mathbb R^3)$; at supports ordered in $x_3$ the kernel is the boundary value of $\mathbf W_n$ (Proposition~\ref{sf-spec}(c)). The growth: for $\varphi$ vanishing with all its derivatives at coinciding points, Taylor's formula at the nearest point with two equal entries, to the order $\lceil p_Vn\rceil\le9n$, absorbs $\delta(x)^{-p_Vn}$; then $(1+|x|^2)^{-k/2}$, $k=9n>3n$ the index of the norm below, is integrable, there are $6^n$ components, and $(2n^2)^{p_Vn}\le(2e^2)^{p_Vn}(n!)^{2p_V}$, which at $p_V=9$ gives $(n!)^{18}|\varphi|_{9n}$.

For (g): let $\eta\in C^2_c(\mathbb R^3;\mathbb R^3)$ be carried by the slab $\{|y_3|<a\}$, $a>0$, and let the spectators be $\tau_{ae_3}F$ and $\theta\tau_{ae_3}F'$, $F,F'\in\mathfrak F(\{x_3>0\})$, on either side of it. As $H^{\rm cons}[\eta]=\sum_{ij}\partial_i\eta_j\,s_{ij}$, Theorem~\ref{sf-gibbs}(f), with the point operators of Theorem~\ref{sf-wight}(b) and an integration by parts in $y$, gives $\sum_i\partial_{y_i}\langle[F'],\mathcal U(\vec y)\mathsf B(a+y_3,a-y_3;s_{ij})\mathcal U(-\vec y)[F]\rangle=0$ for $|y_3|<a$ and every $j$; both members are holomorphic in $y$, and at $y=\iota(x)$, where $y_3=-it$ and $\partial_{y_3}=i\partial_t$, this is, smeared with $f$, $-i\,\Theta^E_{3j}(\partial_0f)=\sum_{k\le2}\Theta^E_{kj}(\partial_kf)$ between the vectors $e^{-aH}[F']$ and $e^{-aH}[F]$: the identity (g) at $\nu=j\le2$, and, multiplied by $-i$, at $\nu=0$, $j=3$.

These vectors span a core of every power of $1+H$, as $(1+H)^me^{-aH}$ is bounded for every $m$, the $[F]$ are dense and the spectral truncations of a vector lie in the range of $e^{-aH}$; read on the core of $(1+H)^{2l}$, $l>p_V$, the bound $\|(1+H)^jT(f)\psi\|\le c_j(f)\|(1+H)^{2l+j}\psi\|$ of (e) at $j=0$ extends the identity to $\mathcal D_0$.

\emph{Step 13: the decomposition.} Lemma~\ref{sf-plat} gives $(\mathsf X,\lambda,(\sigma_\xi))$, the plateaux and (i). The conditioned states $\rho_j=\nu_j(\cdot\,|E_j)$, $E_j$ an event invariant under $a_{K_j}\mathbb Z^3$ with $\nu_j(E_j)>0$, are periodic Gibbs states at $(K_j,\beta_j)$: by Proposition~\textup{\YMMref{106.1}{764}} a periodic Gibbs state stays one under a bounded density $h\ge0$, measurable for the invariant events, with mean $1$, here $h=1_{E_j}/\nu_j(E_j)$. So each $\rho_j$ is invariant under $a_{K_j}\mathbb Z^3\rtimes W_3$ and reflection positive at the lattice mirrors of the nine families (Proposition~\textup{\YMMref{103.4}{758}}), and $(K_j,\beta_j,\rho_j)$ is a Gibbs sequence at $g$.

By Lemma~\ref{sf-sym}(a) at $\rho_j$ and at $\nu_j$, $\rho_j(q_b)=0$ and $\rho_j(a_p)=e_{\rho_j}$ does not depend on $p$, so that $\rho_j(O^T_K(H))-\nu_j(O^T_K(H))=(e_{\rho_j}-e_{\nu_j})a_K^{-1}\sum_pH_{\lambda(p)\lambda(p)}(x(p))$; hence the one-point function of a component depends on $\mathrm{tr}\,H$ alone and is $\frac{c_\xi}3\int\mathrm{tr}\,H$. The re-centring map $S_\xi(X)=\sum_J\prod_{a\notin J}(\frac{c_\xi}3\int\mathrm{tr}\,H_a)\,X_J$ is a bijection, multiplicative for the concatenation, commuting with the conjugation and with every isometry, and $\sigma_\xi=\sigma^c_\xi\circ S_\xi$; as $\int\mathrm{tr}\,H^{\rm cons}[\eta]=\int\nabla\cdot\eta=0$, the re-centred $\sigma^c_\xi$, the limit of $(K_j,\beta_j,\rho_j)$, inherits Theorems~\ref{sf-gibbs}(c)--(f) and~\ref{sf-wight}(a)--(e), and $[F]\mapsto[S_\xi F]$ is unitary: (ii).

By Lemma~\ref{sf-plat}(d) the plateau of a component is the product, at the rate $|v|^{-1}$ with numbers independent of $\xi$, and the mean ergodic theorem for the translations gives $\ker H_\xi=\mathbb C\Omega_\xi$: (iii). Restricted to scalar tuples the $\sigma_\xi$ decompose $\sigma_0$ into clustering functionals with the same plateaux; the uniqueness of the canonical decomposition in Proposition~\YMMref{101.1}{747} gives the map $\pi$, continuous, with $\pi_*\lambda=\lambda_0$, and $\pi(\mathsf X)=\mathsf X_0$ as $\mathsf X$ is compact and $\mathrm{supp}\,\lambda_0=\mathsf X_0$; the means agree, $c_\xi=c_{\pi(\xi)}$: (iv). Finally $\|P_0[F]\|^2=\lim_{|v|\to\infty}\sigma(\theta\bar F\cdot\tau_vF)=\int|\sigma_\xi(F)|^2d\lambda$; the map is onto by the theorem of Stone--Weierstrass, on $\mathcal H_0$ it is the map of the scalar decomposition composed with $\pi$, and $L^2(\mathsf X,\lambda)$ is one-dimensional exactly when $\lambda$ is a point mass, that is, when $\mathsf X=\mathrm{supp}\,\lambda$ is a point.

By (ii), (iii) and Proposition~\ref{sf-rec}, every $\sigma^c_\xi$ satisfies all the Wightman axioms.

\emph{Step 14: coinciding points, overlaps and more points.} (a): $O^T_K(H)$ is the first member of a source of Theorem~\YMMref{D}{8} under a majorant of size $60|H|$, carried by the plaquettes $p$ with $H_{\lambda(p)\lambda(p)}(x(p))\ne0$ and by the plaquettes through the bonds $b$ with $H_{\lambda_b\lambda'_b}(x_b)\ne0$ (Lemma~\ref{sf-sym}(b)): its plaquette part has the size $|H|$, and its bond part $-a_K^{-1}\sum_bH_{\lambda_b\lambda'_b}(x_b)q_b$ the size $59|H|$. Indeed $q_b$ is a bilinear form in the vector parts $w(W_{x,(\mu_b,\alpha)})$, $\alpha\in\{\pm\lambda_b,\pm\lambda'_b\}$, of the four plaquette loops at the bond $b=\langle x,x+a_Ke_{\mu_b}\rangle$; for a form $-a_K^{-1}\sum_b\mathsf w(b)q_b$ the members at the height $0$ and at every level are carried by the at most six cubes with a common corner that hold the four plaquettes through $b$, and they are bounded there as those of a plaquette source of the weight $59\sup|\mathsf w|$ on these plaquettes; to the second order at a small field the form is $a_K^{-1}$ times the quadratic form $-\sum_b\mathsf w(b)\langle\mathsf u_{b,\lambda_b},\mathsf u_{b,\lambda'_b}\rangle$ in the plaquette variables, $\mathsf u_{b,\lambda}$ linear forms with the coefficients $0,\pm\frac12$, whose symmetric kernel has the row sums $2\sum_{p'}|\mathfrak t(p,p')|\le2\sup|\mathsf w|$, so that the kernel $\mathfrak t_H$ of $O^T_K(H)$ has $2\sum_{p'}|\mathfrak t_H(p,p')|\le3|H|$.

So at the depth $s(d)$ an insertion of size $d$ is a source of size $v|H|$ in the units of the working level, and Theorem~\YMMref{D}{8} bounds the cumulant of $n$ such sources, wherever they lie, by $60^nk'_n(g,s)\mathsf b^{-n}\prod_a|H_a|$ on a side that is small over the canonical sequence of $(K,\beta)$ (equation~\YMMref{(2312)}{841}), for $g\le\gamma^{\rm sd}_W$ the least side $\ell_g$ (Lemma~\YMMref{126.3}{846}), the bound reading only the sizes of the sources; $k'_n(g,s)\le C_n(g)Y_s^{8n\mathfrak p_0}$, of the $8\mathfrak p_0$ of one entry $2\mathfrak p_0$ coming from the supremum of the plaquette action, in units of the working level, on the small fields, and $6\mathfrak p_0$ from the number $(LC_W(s))^3$ of the unit cubes of the working level in a ball of radius $d$; and $Y_{s(d)}\le2\log\frac1d+C(g)$.

On every other side $\ell\in\mathfrak L(g)$ the histories with a large field on $T_\ell$, summed against their weights (Proposition~\ref{prop:ward-fixed-volume-rarity}), add the factor $\max\{1,\mathsf c(\ell)\}^n$, $\mathsf c(\ell)=2\ell^3e^{\mathsf B(\bar x_0)\ell^3}/(c'A_0^2)$, with the function $\mathsf B$ and the number $c'$ of equation~\YMMref{(2311)}{841}, as $V_\ell$ does in Lemma~\ref{sf-sym}(b): this is the $\ell$ of $C_n(g,\ell)$. The scale-covariant form is the same bound read at the coupling $\gamma^{\rm sd}_W$ after the dilation of Lemma~\ref{sf-dil} by $L^z$: $(K-z,\beta)$ is a first-passage datum at $\gamma^{\rm sd}_W$ with the label $L^z\mathsf b\in((\gamma^{\rm sd}_W)^{-2},L((\gamma^{\rm sd}_W)^{-2}+\bar B)]$, which gives $\delta_\star\mathsf b\le L^{-z}<(\gamma^{\rm sd}_W)^2\mathsf b$, so that $\log\frac1{L^zd}$ differs from $\log\frac{\mathsf b}d$ by at most a basic number, and the factor $L^{zn}$ of the dilation turns $(L^z\mathsf b)^{-n}$ into $\mathsf b^{-n}$; the depth and the threshold are those of the dilated distance at $\gamma^{\rm sd}_W$, and the shift $(K,\ell,\mathsf b,d)\mapsto(K-1,L\ell,L\mathsf b,Ld)$ is one more step of the dilation.

Supports only at pairwise Euclidean distances $\ge d$ are cut into the boxes $Q$ of sides $\le\frac{99}{100}d''$ of a partition of $T_\ell$ into at most $(2\ell/d'')^3$ boxes, $d''=\min\{\frac34d,L^{-s_{\rm th}(n,g)}\}$: a piece $H_a1_Q$ lies within the Euclidean distance $\frac{\sqrt3}2\cdot\frac{99}{100}d''<0.8574\,d''\le0.87\,d''$ of the centre of $Q$, and pieces of different entries are at Euclidean distances $\ge d\ge\frac43d''\ge1.26\,d''$; the bound at the distance $d''$ for each product of pieces, summed by the $n$-linearity, the sum of the products being the product of the sums, gives $\prod_a\big(\sum_Q\sup_Q|H_a|\big)$ in the place of $\prod_a|H_a|$, at the same $\mathsf b^{-n}$: that is $(\mathsf bd''^3)^{-n}$ times the product of the norms $d''^3\sum_Q\sup_Q|H_a|$, a factor $d''^{-3}$ for each entry.

At $n=2$, for $g\le\gamma^{\rm sd}_W$, every $\ell\in\mathfrak L(g)$, $1.31h<1$ and $s(1.31h)\le K$, two off-diagonal bond forms $Y[\mathsf w_\iota]=-a_K^{-1}\sum_b\mathsf w_\iota(b)q_b$ whose bonds have their midpoints in cubes of side $\frac32h$ at the sup distance $\ge\frac52h$ have $|\langle Y[\mathsf w_1];Y[\mathsf w_2]\rangle_{K,\ell,\beta}|\le36\,C_\infty(g)\mathsf b^{-2}\sup|\mathsf w_1|\sup|\mathsf w_2|$, with the number $C_\infty(g)$ of Lemma~\ref{sf-far}, of $g$ alone. This rests on Lemma~\ref{sf-far}, applied at $d=1.31h$: at every $0<d<1$ with $s(d)\le K$, two real bond forms of one direction each, with $|x_b-x_\iota|\le d-2a_K$ on $\mathrm{supp}\,\mathsf w_\iota$ and $|x_1-x_2|_\infty\ge3d+a_K$, with no upper bound on that distance, have $|\langle Y[\mathsf w_1];Y[\mathsf w_2]\rangle_{K,\ell,\beta}|\le C_\infty(g)\mathsf b^{-2}\sup|\mathsf w_1|\sup|\mathsf w_2|$.

The means vanish, by Lemma~\ref{sf-sym}(a); a bond mirror parts the two forms at the distance $\ge\frac32d$ from both points, and reflection positivity there (Corollary~\YMMref{85.3}{679}, Proposition~\YMMref{91.4}{700}), that is Schwarz's inequality for the positive form, bounds the cumulant by the square roots of the two mirror pairs $\langle\theta Y[\mathsf w_\iota];Y[\mathsf w_\iota]\rangle$. By the transfer operator of the normal direction, positive by reflection positivity, a mirror pair is, as a function of the distance $t$ of the form from the mirror, a sum of exponentials with positive coefficients, symmetric under $t\mapsto\frac\ell2-t$, as the reflection fixes a second plane of the torus at the distance $\frac\ell2$; so on $[\frac32d,\frac\ell2-\frac32d]$ it is largest at $t=\frac32d$, where the form and its image are at the distance $3d$.

The proof of Lemma~\ref{sf-far} bounds such a mirror pair on every torus of $\mathfrak L(g)$ and at every $K\ge s(d)$ by a number of $g$ alone times $\mathsf b^{-2}$ and the sizes, the cutoffs $K=s(d)$ and $K=s(d)+1$ among them; for $d<d_0(g)$ and $K\ge s(d)+2$ this is also the range of Lemma~\ref{sf-one}: on every torus of $\mathfrak L(g)$ the pair has its Wick value of Lemma~\ref{sf-depth}(a) on $T_{\ell_g}$ up to a number of $g$ alone times $v^2$ and the sizes, and the Wick value of two entries of size $d$ at the distance $3d$ is at most a basic multiple of $v^2\le\frac98\mathsf b^{-2}$ times the sizes, by $|C(p,q)|\le Ca_K^3d^{-3}$ and $a_K^3\sum_{p,p'}|\mathfrak t_H(p,p')|\le45d^3|H|$, as in Lemma~\ref{sf-depth}(c); for $d\ge d_0(g)$ the forms are cut into pieces of a size below $d_0(g)$, whose pairs are bounded so, in a number that depends on $g$ alone; and where $K$ lies below the depth that these pieces ask, the crude bound $c\,a_K^{-4}$ of a form serves, as in Lemma~\ref{sf-depth}(c).

At $d=1.31h$, $x_\iota$ the centres of the cubes: $|x_b-x_\iota|\le1.2991h\le d-2a_K$ and $|x_1-x_2|_\infty\ge4h\ge3d+a_K$; and a complex weight of the bonds is the sum of six weights, each $i^\epsilon$ times a real weight carried by the bonds of one direction, of sup norm at most that of the weight: thirty-six pairs. The bound passes to $\mathcal S_n$ by Theorem~\ref{sf-tori}(a) for $\bar O^T_K$ and continuity, which gives the power counting of Proposition~\ref{sf-ext} with $q=8n\mathfrak p_0$ and $h_0\in\ell2^{-\mathbb N}$ with $h_0\le\min\{\frac\ell{64n},\frac12L^{-s_{\rm th}(n,g)}\}$; Proposition~\ref{sf-ext} gives the extension and the contact terms; the same subtractions made at each cutoff, with the products of cubes restricted to those of a side $h$ with $s(1.52h)\le K$, give lattice functionals that are bounded uniformly in $K$ and converge to the extension.

The counts at two points are those of Lemma~\ref{C-355-34}(f): take real, exchange-symmetric, translation-invariant contact terms in $\mathfrak E_2(0)$. They are $\sum_{|\alpha|\le3}\partial^\alpha\delta(x-y)\,c_\alpha(e,e')$ with bilinear $c_\alpha$, symmetric under the exchange of the two entries; under $O(3)$ on $\mathbb R^3$, which contains $-\mathbf 1$, the odd orders drop, and there remain $\mathrm{tr}\,e\,\mathrm{tr}\,e'$ and $\mathrm{tr}(ee')$ at the order $0$ and the contractions of $\partial_i\partial_j$ with $\delta_{ij}\mathrm{tr}\,e\,\mathrm{tr}\,e'$, $\delta_{ij}\mathrm{tr}(ee')$, $e_{ij}\mathrm{tr}\,e'+e'_{ij}\mathrm{tr}\,e$ and $(ee'+e'e)_{ij}$ at the order $2$: six; under $W_3$, on $\mathbb R^3$ or $T_\ell$, there are twelve. $\mathcal S_2$ is no measure across the diagonal: at the scale $d$ a unit cube holds $\gtrsim d^{-3}$ disjoint mirror pairs of balls, each of mass $\ge(\mathsf w_0-\mathsf e^T)\mathsf b^{-2}\mathrm{tr}\,E^2$ by Theorem~\ref{sf-tori}(d)(ii). (b): the action densities stay in the exponent: their weights form the source $w_z=\sum_iz_if_i$ of the sourced bare density $\rho^{w_z}_0=\exp[-\beta A_K+\sum_iz_i\mathcal O(f_i)]$ of equation~\YMMref{(10)}{14}, which Theorem~\YMMref{D}{8} carries through all the levels as a source at the plaquettes; the tensor entries, in the third and fourth cases the bond forms $Y_a=-a_K^{-1}\sum_{b\parallel\mu_a}\phi_a(x_b)q_b$ of the off-diagonal labels $E_a$, enter through their sources of Lemma~\ref{sf-sym}(b) and are derivatives at the strength $0$: $\kappa_K=\partial_{z_1}\cdots\partial_{z_n}\partial_{\tau_1}\cdots\partial_{\tau_r}\log\int\exp[\sum_a\tau_aY_a]\,\rho^{w_z}_0\,dU$ at $z=0$, $\tau=0$; and likewise, in the first two cases, where the tensor weights do not overlap, for diagonal labels with the plaquette forms $Y_a=a_K^{-1}\sum_p\phi_a(x(p))(E_a)_{\lambda(p)\lambda(p)}a_p$, and for every $E_a\in V$ with $Y_a=O^T_K(\phi_aE_a)$, by \eqref{sf-lat} the plaquette form of the diagonal part of $E_a$ plus the bond forms of its off-diagonal part.

Differentiated at the strength $0$ along the tensor entries, the step maps of Theorem~\YMMref{D}{8} give a recursion for the members of each order, each derivative of a step being bounded by Cauchy's estimate in the strengths. The input of Definition~\ref{sf-overlap-input} is used at order two: for two bond forms $Y_a$, $Y_b$ whose weights overlap, one renormalisation step at the depth $u$ creates a part $\mathfrak l$ bounded by $cL^u$ and a remainder bounded by $cK^{\mathsf p}L^{u/2}$ in the units of the sizes, and $\mathfrak l$, read at a background of a lower depth $u'$, is at most $K^{\mathsf p}\mathsf b^{-2}NL^{-2(u-u')}$, $N$ the number of the sites of the level of that depth: the arrays of order two born above a depth are forgotten below it, and this carries the case of two overlapping entries.

We assume this second-derivative estimate, including its stated forgetting bound, for these bond forms; and it is here that the third and fourth cases use that the labels $E_a$ are off-diagonal. Where the tensor supports are $\ge d$ apart, a member of order $\ge2$ vanishes unless its polymer reaches across the distance $d$; where only $\phi_a\phi_b=0$, the arrays of order $m\ge2$ are at most $K^{\mathsf p}L^{-mu/2}$ in the units of the sizes, the weights not overlapping, and with the condition on the $r+n$ supports this gives the same rate. For $r\ge3$ off-diagonal entries of which any two may overlap and no three supports share a point, each overlapping pair is carried by its arrays of order two, which are forgotten below their depth as for two entries; an array of order $m\ge3$ joins $m$ entries whose supports have no common point, so, the supports being compact, its polymer must cross a positive distance fixed by their geometry, and it is bounded as where the supports are apart; the depths proportional to $K$ then give $c=\frac1{360r}$.

Two cutoffs are then compared as in Lemma~\ref{sf-cut}, at depths proportional to $K$, which gives the rates. (c), with the distances of sets read as in Step~9, Euclidean, here and in (d): at the depth $s(d)$ the $n$-th cumulant of the quadratic insertions in the Gaussian part is the sum over the cyclic orders, $3\mathsf b^{-n}K^{(n)}_K$, the value for three free lattice photon fields, for which the left member vanishes identically; the uniform interacting remainder is the additional input $\mathrm{(Q_n)}$, with its explicit first-creation term bounded by the weights $\varpi_t$; for the plane fields the free lattice cumulants differ from $3\mathsf b^{-n}K^{(n)}_{\underline\lambda}$ by at most $\frac13C_n(\frac{a_K}{\vartheta_0d}+\frac{d^3}{\ell^3})\prod_a|\phi_a|$, by the comparison of the lattice kernel with the free kernel (equations~\YMMref{(2245)--(2246)}{824}).

Every periodic Gibbs state and every torus of $\mathfrak L(g)$, for $2\le n\le7$: hypothesis $\mathrm{(Q_n)}$ on $T_{\ell_g}$ and Lemma~\ref{sf-tr}(a), the weights being bounded functions of the plaquettes in balls of radius $d$ within one ball of radius $9d$, give the lattice inequality in Definition~\ref{hyp:quantitative-torus} with $\varepsilon^{\rm rt}_n(g,d)$ added to $\mathsf e_n$, and $\varepsilon^{\rm rt}_n(g,d)\to0$ with $d$ since some $\sigma\in(0,1)$ has $4-\frac n2-3\sigma>0$ exactly when $n\le7$. (d): Lemma~\ref{sf-jet}(c) gives the laws at a cutoff on $T_{\ell_g}$, the large fields being bounded as in Lemma~\ref{sf-jet}(a), and the comparison of the free lattice cycles with $k_3$, $k_4$, by equations~\YMMref{(2245)--(2246)}{824}, gives the terms $(d/\ell_g)^3$; Lemma~\ref{sf-tr}(c) carries the laws to every periodic Gibbs state, at three points with the further $C(g)d^{3/2}\rho^9r^{-3}d^{-6}\prod_a|w_a|$ and at four with $C(g)d^{3/2}\rho^{12}r^{-6}d^{-6}\prod_a|w_a|$, $d^{3/2}$ times the natural sizes; first passing to a continuum subsequential limit at fixed positive widths, then letting the widths tend to $0$ and using Lebesgue's differentiation theorem gives the kernels, with estimates uniform in the states, those of the truncated functions; the law of four points is that on the torus and at a cutoff; for limits and re-centred components $\sigma^c_\xi$, Lemma~\ref{sf-plat} and Theorem~\ref{sf-wight}(d), (f), the kernels being continuous, and for a limit the third centred moment being the third cumulant; the point-split law follows from the three-point law, Lemma~\ref{sf-freev}(d) and the kernels of Theorem~\ref{sf-gibbs}(e), the constant $2.7$ covering $2.66$ and the passage from $\mathrm{dist}(z,\{x,y\})$ to $|c-z|$.

\emph{Proof of Corollary~\ref{sf-ode}.} By Theorem~\ref{sf-wight}(c), (d) the kernel $s_2(x,y;e,\mathbf 1)$ is real-analytic, a function of $u=x-y$, and linear in $e$, so $s_2(x,y;e,\mathbf 1)=\mathrm{tr}(e\,\tau(u))$ with $\tau(u)\in V$; invariance under $O(3)$ gives $\tau(Ru)=R\tau(u)R^t$, so $\tau(u)$ commutes with the stabiliser of $\hat u$ and is $A(r)\mathbf 1+B(r)\hat u\hat u^t$, and $\mathrm{tr}\,\tau=3A+B=F$. Conservation against one action density, $\sigma((H^{\rm cons}[\eta]),(f\mathbf 1))=0$ for $\mathrm{supp}\,\eta$ off $\mathrm{supp}\,f$, reads $\iint[(\nabla\cdot\eta)(x)F(x-y)-2\sum_{ij}\partial_i\eta_j(x)\tau_{ij}(x-y)]f(y)\,dx\,dy=0$; integration by parts in $x$, for all such $\eta$ and $f$, gives $2\sum_i\partial_i\tau_{ij}=\partial_jF$ off $u=0$, that is $2(A'+B'+\frac2rB)=F'$, as $\sum_i\partial_i(B\hat u_i\hat u_j)=(B'+\frac2rB)\hat u_j$. With $A=\frac13(F-B)$ this is $B'+\frac3rB=\frac14F'$, and $(r^3B)'=\frac14r^3F'$. $\square$

\begin{remark}
The limits on $\mathbb R^3$ are taken along subsequences of Gibbs sequences, and each theorem speaks of the limit of the sequence chosen; the label, the bounds and the laws at short distance are the same for all of them. The uniqueness of the vacuum in Theorem~\ref{sf-wight}\textup{(f)} is that of each component $\sigma_\xi$, and the vacuum of $\mathcal H(\sigma)$ is unique exactly when the index space $\mathsf X$ is a point. Theorem~\ref{sf-wight} concerns limits through periodic Gibbs states; of the limits along growing tori Theorem~\ref{sf-gibbs}\textup{(c)} is asserted, with positivity at the coordinate mirrors alone. The laws at short distance are laws to leading order, with errors that tend to $0$ with the distance, and Theorem~\ref{sf-tori} and Theorem~\ref{sf-coin}\textup{(a)}, \textup{(b)} concern one torus at a time, with constants that depend on its side.
\end{remark}

\section{Uniform forms and distributional bounds}\label{sec:bil}

The estimates in this section apply to local forms beyond the fixed tensor representatives. We retain the explicit dependence on the cell size, cutoff and torus side, together with the distributional bounds at coinciding points. These formulations also record the distinction between uniqueness of a limiting scalar functional and uniqueness of the vacuum in its components.

\subsection{Notation and hypotheses}

\emph{Data.} The radius, the constants and the data are those of the main theorems (Proposition~\YMMref{69.1}{595} and Section~\YMMref{82}{665}; the constants chosen in the order of that proof; the uniform free-response geometry, radius and decay conditions in Section~\YMMref{52.6}{426}, retained in Section~\YMMref{120.3}{820}; $c_\star$ by equation~\YMMref{(1429)}{582}): $L$ odd, $M$ a power of $L$, $A_0$, $A_1$, $P=2\mathfrak p_0\ge18$, $\eth\in(0,1)$, $c'=(1-\eth)c_\star$, the last constant $\gamma$, and the ceilings $\gamma_1=\gamma^{\rm sd}_W\le\gamma_2=\gamma^{\rm ref}_W\le\gamma$ of the choices in equations~\YMMref{(2326)--(2327)}{845}. $B_\sharp\ge0$ is the flow bound: $|L^{-1}x_{j-1}-x_j|\le B_\sharp$ at every step (Section~\ref{sec:setup-first-passage-arithmetic}), and $\bar B=B_\sharp L/(L-1)$ (Section~\YMMref{74}{614}); both are numbers of the constants alone. $T_\ell$, $\ell=2L^m$, is the torus of side $\ell$, $\Lambda_{K,\ell}$ its lattice of spacing $a_K=L^{-K}$, $\mathbb L_K=a_K\mathbb Z^3$; $|\cdot|$ and $|\cdot|_\infty$ are the Euclidean and the maximum norm, on a torus minimised over lifts; $\hat u=u/|u|$.

A first-passage datum $(K,\beta)$ at the reference coupling $g$ (\YMMref{74.2}{616}, \YMMref{74.1}{615}) has the running couplings $x_0=\beta,\dots,x_K$, with $|L^{-1}x_{t-1}-x_t|\le B_\sharp$ for $1\le t\le K$ (Proposition~\YMMref{74.1}{615}), so that, the differences telescoping, $|L^{-z}x_{K-z}-x_K|\le B_\sharp\sum_{0\le i<z}L^{-i}\le\bar B$ for $0\le z\le K$; $y_t=\log x_t$, $g_t=x_t^{-1/2}$, the large-field threshold $p_0(g_t)^2=A_0^2y_t^P$; the first passage $x_K\le X_g<x_{K-1}$ \YMMref{74}{614} at $X_g=L(g^{-2}+B_\sharp)$, so that $X_g+\bar B=L(g^{-2}+\bar B)$ \YMMref{(1526)}{615}; and the label $b=a_K\beta\in I^{\rm lab}(g):=[g^{-2},\bar x_0]$, $\bar x_0=L(g^{-2}+\bar B)=X_g+\bar B$; for $\mathsf b$ in the interval $I(g)=(g^{-2}+\bar B,X_g-\bar B]\subset I^{\rm lab}(g)$ of Section~\YMMref{74}{614}, $X_g-\bar B=L(g^{-2}+\bar B)-2\bar B$, $(K,L^K\mathsf b)$ is a first-passage datum for every integer $K\ge1$. $\mathfrak L(g)=\{2L^m:m\ge m_{\mathfrak L}(g)\}$ is the family of the sides at $g$ (equation~\YMMref{(1528)}{615}), $\ell'=2\mathsf l$, $\mathsf l=L^{m_{\mathfrak L}(g)}$, its least member, $\mathsf l\ge LM$.

A side $\ell=2L^m$ is \emph{admissible} if $MR(x_K)\le L^m$, $R(x)$ the least power of $L$ that is at least $(\log x)^r$, $r$ a constant; it is \emph{small} if it is small over the running couplings in the sense of equation~\YMMref{(2312)}{841}: admissible, $LM\le L^m$, $\ell\ge N_{\mathsf H}$, the lower bound of condition (H), and $\mathsf B(x_K)\ell^3\le\frac12c'p_0(g_K)^2$, $\mathsf B$ the function equation~\YMMref{(2311)}{841} of the volume factor $e^{\mathsf B\ell^3}$ (equation~\eqref{eq:ward-volume-constants}); for $g\le\gamma_1$ the side $\ell'$ is small for every datum (Lemma~\YMMref{126.3}{846}). As $R$ does not decrease, $x_K\le X_g<\bar x_0$ and $L^{m_{\mathfrak L}(g)}\ge LMR(\bar x_0)$ \YMMref{(1528)}{615}, every side $2L^m\in\mathfrak L(g)$ has $LMR(x_K)\le L^m$ for every datum at $g$: the sides of $\mathfrak L(g)$ are admissible.

\emph{Depth.} $Y_s=(s+1)\log L+\log(g^{-2}+\bar B)=\log\bar x_s$, $\bar x_s=L^{s+1}(g^{-2}+\bar B)$ and $C_W(s)=C_{W,0}A_0^2Y_s^P$ as in equation~\YMMref{(2321)}{843}, with the room constant $C_{W,0}=\max\{1,4L^2M/A_0^2\}$ (equation~\YMMref{(2320)}{843}), so that $C_W(s)\ge4L^2MY_s^P$ and $L^{-2}C_W(s)\ge4MY_s^P\ge M$; $\mathsf{Lm}_s=L^2C_W(s)Y_s^{3\mathfrak p_0}$, $x^-_s=4L^{s-2}g^{-2}/(5\mathsf{Lm}_s)$. For $0<d<1$, $s(d)$ is the least $s$ with $L^sd\ge C_W(s)$; so $Y_{s(d)}\le2\log\frac1d+C(g)$. At the working level $\mathsf k=\max\{K-s(d),2\}$: $y=y_{\mathsf k}$, $\hat d=d/a_{K-\mathsf k}$, $\upsilon^{-1}=a_{K-\mathsf k}x_{\mathsf k}$, and $\upsilon^2\le\frac98b^{-2}$; if $s(d)\le K$, then $\hat d\ge L^{s(d)-2}d\ge L^{-2}C_W(s(d))\ge M$ and $d\ge C_W(s(d))a_{s(d)}\ge4L^2Ma_K\ge676a_K$.

The free template of Section~\YMMref{120}{817} has, for $0\le n\le K$, the orthogonal projection $\Pi_n=\Pi^{\rm tp}_n$ onto the background space of the level $n$ of the exact curvatures, with its kernel $\Pi_n(p,q)$ on the plaquettes (Section~\YMMref{120}{817}); the free kernel is $\Pi_0$, the background kernel of the depth $s'$ is $C^{\rm bg}_{K,s'}=\Pi_{K-s'}$, $0\le s'\le K$, and the slice of the step $j$ is $\mathcal C_{(j)}=\Pi_{j-1}-\Pi_j$, $1\le j\le K$ \YMMref{(2227)}{818}. So $\Pi_0=\Pi_n+\sum_{1\le j\le n}\mathcal C_{(j)}$ for every $n$, and $\Pi_k-\Pi_T=\sum_{k<j\le T}\mathcal C_{(j)}$ for $k\le T$ \YMMref{(2227)}{818}; in particular the free kernel is the sum of the background kernel $\Pi_K$ of the top level and of all the slices.

With $\eta_n=L^{-n}$ and $\mathsf r=|x(p)-x(q)|_\infty$, $|\Pi_n(p,q)|\le C^{\rm bg}_{\rm tp}\eta_n^3$ and $|\mathcal C_{(j)}(p,q)|\le C_{\rm tp}\eta_{j-1}^3e^{-\delta_{\rm tp}\mathsf r/a_{K-j+1}}$ (Propositions~\YMMref{120.5}{821} and~\YMMref{120.6}{821}). The background kernel of the working level is $C^{\rm bg}=C^{\rm bg}_{K,K-\mathsf k}=\Pi_{\mathsf k}$, and the free kernel is $C^{\rm bg}$ plus the slices $\mathcal C_{(j)}$, $j\le\mathsf k$, the slice $\mathcal C_{(j)}$ decaying exponentially beyond the scale $a_{K-j+1}$. $K_3(u)=(3\hat u\hat u^{\rm T}-\mathbf1)/(4\pi|u|^3)$.

\emph{Large fields; the chain.} The chain has the levels $0\le j\le K$, $\eta_j=L^{-j}$, local families with the sizes of \YMMref{(759)}{375}, the geometry of Proposition~\YMMref{45.1}{353} and Lemma~\ref{lem:analytic-domain-geometry}, and the contraction numbers $\theta^{[j]}$. A history is a sequence of large-field regions up to a step; $Z_t(h)$ is its region at the level $t$, $\tilde\nu_K(h)=Z^{-1}\int\mathfrak w_h\rho_0\,dU$ the law of the histories; $\mathsf w_t$ is the probability that the first large field is created at the step $t$, so that $\sum_{t'\le t}\mathsf w_{t'}=\tilde\nu_K(\{Z_t\ne\emptyset\})$; $\mathsf W_T=\sum_{t\le T}\mathsf w_t$. $\Phi^{(t)}[X]$ is the conditional mean of $X$ given the block field of the level $t$ on the history without large fields. For sources with majorants (Lemma~\ref{bil:forms}) $\mathsf D_j$ is the least sup distance, in units of $a_{K-j}$, of corners of plaquettes of different sets $S_a$. For two sources at the working level, $\mathsf k<j_1\le T$ are the middle and the upper level that go with $(\mathsf k,\hat d,y)$ (equation~\YMMref{(2316)}{842}), $n_T=T-\mathsf k$, $L^{n_T}\ge\hat dy^{3\mathfrak p_0}$; for $\mathsf L\ge0$, $c>0$,
\[
{\rm sm}^{\mathsf L,c}_T=\max\Bigl\{x_T^{-1/2},\ x_{\mathsf k}e^{-d_*(T)},\ x_{\mathsf k}e^{-(\mathsf D_{\mathsf k}/M-2)},\ \max_{\mathsf k\le i<T}(\delta_i+\varepsilon^{\rm co}_i),\ x_{\mathsf k}^2\mathsf LL^{3(K-\mathsf k)}e^{-cp_0(g_T)^2}\Bigr\},
\]
and ${\rm sm}_T$ is its value at $c=\frac12c'$ and $\mathsf L=\mathsf L_\ell:=2M^{-3}\ell^3$, the number of the weights of a small side: by the rarity of large fields on one torus (Proposition~\ref{prop:ward-fixed-volume-rarity}), $\mathsf w_t\le\tilde\nu_K(\{Z_t\ne\emptyset\})\le2M^{-3}\ell^3L^{3(K-t)}e^{-c'p_0(g_t)^2}e^{\mathsf B(x_K)\ell^3}$ for $1\le t\le K$, for $g\le\gamma_2$, $K\ge3$ and every admissible side, and on a small side $\mathsf B(x_K)\ell^3\le\frac12c'p_0(g_K)^2\le\frac12c'p_0(g_t)^2$, as $x_t\ge x_K$; so $\mathsf w_t\le\mathsf L_\ell L^{3(K-t)}e^{-\frac12c'p_0(g_t)^2}$ and $\mathsf W_T\le\mathsf L_\ell L^{3(K-T)}e^{-\frac12c'p_0(g_T)^2}$ there. The symbols $\alpha_0,\beta_s,\kappa,K_M,K_\Lambda,C_{\mathcal A}$ retain their meanings in the common choice of Proposition~\YMMref{69.1}{595} and Section~\YMMref{82}{665} and $C_P$ is the number of Definition~\YMMref{51.2}{390}; $\delta(t)=A_1t^{-1/2}(\log t)^{\mathfrak p_0}$ and $\delta_i=\delta(x_i)$; $\varepsilon^{\rm co}_i$ is the error $e_i^{\rm co}$ defined before equation~\YMMref{(2302)}{837}, $d_*$ that of equation~\YMMref{(2262)}{828}, $M_{\rm pl}$ that of equation~\YMMref{(2285)}{833}; $\mathsf c_X$, $\mathsf c_I$, $\mathsf C_{\mathsf M}$, $\mathsf C_i$, $\mathsf c_{\rm sb}$ and $\mathsf c_{\rm tl}$ are fixed in \eqref{eq:short-distance-constants}--\eqref{eq:short-distance-collected-constants} and the preceding paragraph; $\mathsf p_0$ is the exponent of Corollary~\textup{\ref{IV-62-66}}; $C_{\rm tp}$, $C^{\rm bg}_{\rm tp}$, $\delta_{\rm tp}$ are the numbers of Propositions~\YMMref{120.5}{821} and~\YMMref{120.6}{821}, $C_{\rm lat}$ that of equation~\YMMref{(2245)}{824}, and $C^{\rm fl}_\star$, $C^{\rm tm}_\flat$ those of equation~\eqref{eq:ward-template-constants}; $c_\theta$ and $\mathsf S_3$ are numbers of the constants alone.

The short-distance estimates use the constants of
Lemmas~\YMMref{125.2}{840} and~\YMMref{125.5}{842}.
We write these constants with the font used here:
$\mathsf C_{\rm sz}=C_{\rm sz}$,
$\mathsf C_{\mathsf M}=C_{\mathsf M}$,
$\mathsf c_{\rm up}=c_{\rm up}$, $\mathsf c_I=c_I$ and
$\mathsf C_i=C_i$, with the same convention for $C'_3,C''_3$.
In particular $\mathsf c_X=405B_3^2A_0^2/2$.
Fix these choices before the source, lattice and levels, and put
\begin{equation}\label{eq:short-distance-constants}
 \mathsf c_{\mathfrak p}=\sup_{a\in\mathbb Z_{\ge0}}
                  (a+2)^{4\mathfrak p_0}L^{-a},\qquad
 \mathsf c_{\rm sb}=\mathsf C_{\rm sz}+\mathsf c_{\rm up},\qquad
 \mathsf c_{\rm tl}=2L^8\mathsf c_{\mathfrak p}+1.
\end{equation}
The supremum is finite because exponential decay dominates the
fixed power of $a$. For $\vartheta\in\{1,2\}$ define
\begin{equation}\label{eq:short-distance-collected-constants}
\begin{split}
 \mathsf C_B^{(\vartheta)}={}&\mathsf C_{\mathsf M}
       +4\mathsf c_X\mathsf c_I+6\mathsf c_I^2\\
 &+\mathsf c_{\rm tl}\bigl[\mathsf C_{\mathsf M}
   +(2+\vartheta^2+\vartheta^2\mathsf c_{\rm tl})
       \mathsf c_{\rm sb}^2
   +2\vartheta\mathsf c_{\rm sb}(\mathsf c_X+\mathsf c_I)\bigr],\\
 \mathsf C_{\rm B}={}&\mathsf C_7+2(1+A_1)^4\mathsf C_5
       +\mathsf C'_3+\mathsf C''_3+\mathsf C_6
       +\mathsf C_B^{(1)},\\
 \mathsf C'_B={}&\mathsf C_B^{(2)},\qquad
 \mathsf C^e_B=\mathsf C_7+2(1+A_1)^4\mathsf C_5
       +\mathsf C'_3+\mathsf C''_3+\mathsf C_6+\mathsf C'_B.
\end{split}
\end{equation}
These are fixed numbers, and $\mathsf C_{\rm B}\le\mathsf C^e_B$.

\begin{lemma}[The coefficient of the collected error]
\label{lem:short-distance-collected-constant}
Under the hypotheses of Corollary~\YMMref{125.6}{843},
the coefficient $C_B$ in its estimate may be chosen to be
$\mathsf C_{\rm B}$ in
\eqref{eq:short-distance-collected-constants}.
\end{lemma}

\begin{proof}
Use the notation of that corollary:
$p=\mathfrak p_0$, $q=\widehat d_{\mathsf k}$,
$y=y_{\mathsf k}$, $n=T-\mathsf k\ge1$, and $s_T\le1$.
Divide the errors by $\widehat m_1\widehat m_2$, and write
$\mathfrak m=n^4q^6y^{15p}$.
The first four terms in equation~\YMMref{(2319)}{843}
are the direct terms of the seven-error estimate.
Its remaining terms are bounded by
\[
 [\mathsf C_7+2(1+A_1)^4\mathsf C_5
             +\mathsf C'_3+\mathsf C''_3+\mathsf C_6]
                  \mathfrak m s_T.
\]
Indeed $x_T^{-1}\le s_T^2\le s_T$, $W_T\le t_T\le s_T$,
and $1+A_1y^p\le(1+A_1)y^p$.

In the remainder of equation~\YMMref{(2315)}{842}, use
$i_T\le\mathsf c_Iq^3y^{3p}s_T$ and
$z_T\le\mathsf c_Xy^{-7p}+i_T$.
The separated-support term costs
$\mathsf C_{\mathsf M}\mathfrak m s_T$.
The cross terms cost
$(4\mathsf c_X\mathsf c_I+2\mathsf c_I^2)\mathfrak m s_T$,
and the top term is at most
\[
 2z_T^2\le4\mathsf c_X^2y^{-14p}
                   +4\mathsf c_I^2\mathfrak m s_T.
\]
For its last line, the size and tail bounds give
\[
 2L^8\mathsf c_{\mathfrak p}+n\le\mathsf c_{\rm tl}n,
 \qquad b_0\le\mathsf c_{\rm sb}nq^3x_{\mathsf k}y^{2p},
 \qquad x_{\mathsf k}^2t_T\le s_T.
\]
Consequently the bracket in that line is at most
\[
 [\mathsf C_{\mathsf M}+(3+\mathsf c_{\rm tl})\mathsf c_{\rm sb}^2
       +2\mathsf c_{\rm sb}(\mathsf c_X+\mathsf c_I)]
                  n^3q^6y^{5p}x_{\mathsf k}^2.
\]
Multiplication by its remaining factor, at most
$\mathsf c_{\rm tl}nt_T$, gives the last part of
$\mathsf C_B^{(1)}\mathfrak m s_T$.
Adding these coefficients proves the assertion.
\end{proof}

We use this admissible value of the collected coefficient in
the depth error of equation~\YMMref{(2330)}{846}.
We use the following names for its auxiliary functions:
\begin{equation}\label{eq:short-distance-envelope-names}
 \mathsf{Lm}_s=\mathsf L_s=L^2C_W(s)Y_s^{3\mathfrak p_0},\qquad
 \mathsf{sm}_g(s)=\sigma_g(s).
\end{equation}
Here $\mathsf L_s$ is defined in equation~\YMMref{(2322)}{844},
and $\sigma_g$, $\epsilon_g$, $Y_s^-$, $x_s^-$ and $s_1(g)$
are the functions and threshold of
equations~\YMMref{(2328)--(2332)}{846}.
The proof there gives $s^a\mathsf{sm}_g(s)\to0$ for every
$a\ge0$ and boundedness of $Y_s^{2\mathfrak p_0}\epsilon_g(s)$.
Changing the fixed collected coefficient to the admissible value
above preserves that proof: its only additional occurrence is
a fixed multiple of a polynomial times $\mathsf{sm}_g(s)$.

\emph{Laws.} $\langle\cdot\rangle_\ell=\langle\cdot\rangle_{K,\ell,\beta}$ is the Wilson measure on $T_\ell$, $\tilde\mu_\ell$ its periodic extension to $\mathbb L_K$, $\mathcal G^{\rm per}_{K,\beta}$ the set of the periodic Gibbs states on $\mathbb L_K$; a \emph{law} at $(K,\beta)$ is a member of $\mathcal G^{\rm per}_{K,\beta}$ or a $\tilde\mu_\ell$, $\ell\in\mathfrak L(g)$, and for the latter distances are those of $T_\ell$. $\langle X_1;\dots;X_n\rangle_\nu$ is the joint cumulant (complex multilinear), $\nu(X;Y)=\langle X;Y\rangle_\nu$, $A[F]=A^\nu[F]=F-\nu(F)$; for a tuple $X=(H_1,\dots,H_n)$, $A^\nu(X)=\prod_aA^\nu[O^T_K(H_a)]$, extended linearly to combinations of tuples; along a sequence of laws $\nu_j$, $A_j=A^{\nu_j}$, and $A(X)=A^\rho(X)$ for the law $\rho$ at hand. $p_\ell=\ell^{-3}\log Z_{K,\ell,\beta}$ is the pressure, and $G_g=C_S+C_g+2C_{\rm tt}C'_{\rm tt}E_0$ is formed with the constants of Section~\ref{sec:source-system-complements}, with which the canonical source energies form a system of localised source energies (Lemma~\ref{lem:canonical-source-system}).

A \emph{mirror} $\theta$ is the reflection in a coordinate plane through midpoints of bonds; it acts on points, on functions of the field by $(\theta F)(U)=F(\theta U)$, and on tensor labels. The \emph{nine families} are the planes whose normals are $e_\mu$ or $(e_\mu\pm e_\lambda)/\sqrt2$, $\mu\ne\lambda$ (eighteen unit normals); every periodic Gibbs state is invariant under the lattice isometries and reflection positive, in the sense of the reflection form below, at the lattice planes of the nine families (Proposition~\textup{\YMMref{103.4}{758}}). For a mirror $\theta_c$ in $\{x_{\mu_0}=c\}$ and a law $\nu$, $\|F\|_{\nu,c}=\nu(\overline{\theta_cF}F)^{1/2}$ on the functions of the field on $\{x_{\mu_0}>c\}$.

\emph{Boards.} A board of $\mathbb L_K$ has a side $\mathsf l_0=a_i$, $i\le K-1$ an integer (negative $i$ admitted, so that $\mathsf l=L^{m_{\mathfrak L}(g)}$ is a side of a board), and an offset $o\in a_K\mathbb Z^3$; it is the partition of $\mathbb R^3$ into the cubes $o+\mathsf l_0\mathbf n+[-\frac12\mathsf l_0,\frac12\mathsf l_0)^3$, $\mathbf n\in\mathbb Z^3$, its \emph{cells}, so that the offset is the centre of a cell, and its \emph{walls} are the planes $\{x_\mu=o_\mu+\mathsf l_0(t+\frac12)\}$, $t\in\mathbb Z$, $\mu=1,2,3$. As $\mathsf l_0/a_K=L^{K-i}$ is odd, the walls lie in $a_K(\mathbb Z+\frac12)$ and pass through midpoints of bonds: they are planes of mirrors. A \emph{function of a cell} is a function of the variables of the bonds with both ends in the open cell; its images under successive reflections in the walls are its \emph{copies}, with complex conjugation after an odd number of reflections. For a side $\ell_0\in2\mathsf l_0\mathbb N$ the chessboard seminorm $\|\Phi\|_{{\rm ch},\ell_0}\ge0$ of a bounded function $\Phi$ of a cell has as its $N_0$-th power, $N_0=(\ell_0/\mathsf l_0)^3$, the $\tilde\mu_{\ell_0}$-mean of the product of the $N_0$ copies of $\Phi$ in a period box.

\emph{Field and label conventions.} The field is $O_K^T$ of \eqref{sf-lat}, with the same $w(G)$, $a_p$, $q_b$ and midpoint convention. We use $\langle X,Y\rangle=\frac12\operatorname{tr}(XY)$ and the sup norm $|H|_\infty$. In this section $E_{\lambda\lambda'}$ denotes the symmetric label with entries $1$ at $(\lambda,\lambda')$ and $(\lambda',\lambda)$, rather than a matrix unit. A loop is a pair $l=(\alpha,\alpha')$, $\alpha,\alpha'\in\{\pm1,\pm2,\pm3\}$, $|\alpha|\ne|\alpha'|$; $W_x(l)$ is the holonomy from $x$ along $e_\alpha$, then $e_{\alpha'}$, around the plaquette $p(x,l)$, and $\sigma(l)=\mathrm{sgn}\,\alpha\,\mathrm{sgn}\,\alpha'\,\mathrm{sgn}(|\alpha'|-|\alpha|)$. 
We write $S_n^T[\nu](H_1,\ldots,H_n)$ for the joint cumulant of the $O_K^T(H_a)$ under $\nu$. The free covariance is $\mathcal W$ in \eqref{sf-free}. The test norm used here is
\[
 N(H)=\max_{\lambda\lambda'}\bigl(|H_{\lambda\lambda'}|_\infty+d\,\operatorname{Lip}H_{\lambda\lambda'}\bigr).
\]

\emph{Forms.} A \emph{cell} is a plaquette, a positively oriented bond with the four loops that start along it, or a site with all its $24$ loops. A \emph{form} is $F=a_K^{-1}\sum_ph(p)a_p$ or $F=a_K^{-1}\sum_{\mathsf c}\sum_{l,l'}v(\mathsf c,l,l')\langle w(W_l),w(W_{l'})\rangle$, the cells $\mathsf c$ being bonds, resp. sites, with real weights of finite support (complex weights where this is said); its \emph{size} $|F|_\star$ is $\sup|h|$, $59\|v\|$, resp. $23\|v\|$, $\|v\|=\sup_{\mathsf c}\sum_{l,l'}|v(\mathsf c,l,l')|$; its \emph{carrier} is the set of the centres of the cells of its support, $S[F]$ the set of the plaquettes of these cells; its \emph{kernel} is $\mathsf t_F(p,p')=\frac12h(p)\delta_{pp'}$, resp. $\mathsf t_F(p',p'')=\frac12\sum v(\mathsf c,l,l')\sigma(l)\sigma(l')[\delta_{p'p(l)}\delta_{p''p(l')}+\delta_{p'p(l')}\delta_{p''p(l)}]$. For a weight $\mathsf w$ of the bonds, $Y[\mathsf w]=-a_K^{-1}\sum\mathsf w(\mathsf e)q_{\mathsf e}$ is a bond form of size $59\sup|\mathsf w|$; $O^T_K(H)$ is a sum of six forms, each on the plaquettes of one normal or on the bonds of one direction, of total size at most $180|H|_\infty$. The \emph{Gaussian value} of two forms is
\[
G(F_1,F_2)=6\upsilon^2\sum\mathsf t_{F_1}(p,p')\,\mathsf t_{F_2}(q,q')\,C^{\rm bg}(p,q)\,C^{\rm bg}(p',q'),
\]
Wick's value for three free photon fields, $2(N^2-1)=6$ at $N=2$.

\emph{Limits.} $S$ denotes a volume-first limit, or a continuum limit along growing tori, of the action density, with its defining sequence (Section~\ref{sec:state-limit-conventions}). A tuple is \emph{separated} if its supports are pairwise disjoint. Along a sequence of laws at which all $S^T_n[\nu_j](X)$ converge, $\sigma^T(X)=\sum_\pi\prod_{I\in\pi}S^T_{\#I}(X_I)$, the sum over the partitions into blocks of at least two, is the limit of the means of the products of the insertions centred in the states, and $\sigma^T(\emptyset)=1$. A \emph{limit through periodic Gibbs states} is a limit along $\nu_j\in\mathcal G^{\rm per}_{K_j,\beta_j}$, $K_j\to\infty$, that converges at all separated tuples in $C^1_c(\mathbb R^3;V)$; $H,P_1,P_2,\Omega$ are the Hamiltonian, the momenta and the vacuum reconstructed from $\sigma^T$ at $\{x_3=0\}$ (Proposition~\YMMref{94.1}{720}; classes $[F]$ of combinations $F$ of tuples in $\{x_3>0\}$, $\langle[F'],[F]\rangle=\sigma^T(\theta\bar F'\cdot F)$), with $e^{-tH}[F]=[\tau_{te_3}F]$, the unitary group $\mathcal U(v)[F]=[\tau_{(v,0)}F]$, $v\in\mathbb R^2$, of the horizontal translations, and the joint spectral measure $\mathcal P$ of $(H,P_1,P_2)$; and alike at a plane of each of the nine families: the \emph{nine pictures}. $\tau_v$ is the translation by $v$.

\emph{Numbers.} Each symbol $C_\cdot$, $c_\cdot$, $k_\cdot$, $s_\cdot(g)$, and each error function $\varepsilon_\cdot,\omega_\cdot,o_\cdot,e_\cdot,{\rm Df}_\cdot$ of $s$, is a function of the constants, of $g$ and of the arguments shown alone; it does not depend on $K$, $\beta$, $\ell$, the law, a weight or a position unless these are shown, and error functions tend to $0$ as $s\to\infty$. $s_1(g)$ is the first depth of the envelopes of equations~\YMMref{(2328)--(2336)}{846}, with the notation in \eqref{eq:short-distance-envelope-names}, $s_p(g)$ ($p$ even) that of Lemma~\ref{bil:far}.

\subsection{Principal results}

\begin{theorem}\label{thm:bil}
Let $(K,\beta)$ be a first-passage datum at $g$ and $b=a_K\beta$.

\textbf{(A)} Let $0<g\le\gamma_2$ and let $\ell$ be small. Then the following hold.
\begin{itemize}
\item[(A1)] For $n\ge2$, $0<d<1$, $K\ge K_0'(n,d)$ and forms $F_1,\dots,F_n$ on $\Lambda_{K,\ell}$, with complex weights allowed at bonds and sites, whose carriers are pairwise at distance $\ge d$,
\[
|\langle F_1;\dots;F_n\rangle_\ell|\le k^\circ(n,d)\,b^{-n}\prod_a|F_a|_\star,
\]
where $d''=\min\{\frac34d,L^{-s_{\rm th}(n)}\}$, $K_0'(n,d)=s(d'')$, $k^\circ(n,d)=(2\ell_{\max}/d'')^{3n}k'_n(g,s(d''))$, $\ell_{\max}^3=\frac12c'A_0^2Y_{s(d'')}^P$, with $s_{\rm th}$, $k'_n$ of Lemma~\ref{bil:nsrc}: the power $d^{-3}$ for each entry, times a power of $\log\frac1d$. No upper bound on the distances is asked.
\item[(A2)] Let $s_1(g)\le s(d)\le K$, $3d\le|x_1-x_2|\le4d$, and let the carrier of $F_i$ lie within $d-2a_K$ of $x_i$. Then $|\langle F_1;F_2\rangle_\ell-G(F_1,F_2)|\le\varepsilon^\natural_g(s(d))\,\upsilon^2|F_1|_\star|F_2|_\star$, $\varepsilon^\natural_g$ of Lemma~\ref{bil:two}. For complex tensor weights $H_i$ supported within $d-2a_K$ of $x_i$ the same holds for $O^T_K(H_1)$, $O^T_K(H_2)$, the Gaussian value being formed with their kernels (Lemma~\ref{bil:forms}) and the error being $4\cdot60^2\varepsilon^\natural_g(s(d))\upsilon^2|H_1|_\infty|H_2|_\infty$. If $H_i=\varphi_iE_{\lambda_i\lambda_i'}$, $\varphi_i$ real and Lipschitz with support in $B(x_i,d-2a_K)$, and $\lambda_1\ne\lambda_1'$, then
\[
|\langle O^T_K(H_1);O^T_K(H_2)\rangle_\ell-b^{-2}\mathcal W(H_1,H_2)|\le e^{\rm od}\,b^{-2}N(H_1)N(H_2),
\]
with the number $e^{\rm od}$ of Step~2, at most $\bar e^{\rm od}_g(s(d))$, $\bar e^{\rm od}_g$ an error function.
\item[(A3)] Let $s_1(g)\le s(d)\le K$, and let the carrier of $Y=Y[\mathsf w]$, $\mathsf w$ real on the bonds of one direction, lie within $d-2a_K$ of a point at a distance between $\frac32d$ and $2d$ from the plane of a mirror $\theta$. Then
\[
0\le\langle\theta Y;Y\rangle_\ell\le C^\circ_{\rm od}(g)\,b^{-2}\sup|\mathsf w|^2,\qquad C^\circ_{\rm od}(g)=\tfrac98\bigl[600C_{\mathsf k}^2+59^2\bar\varepsilon^\natural_g\bigr],\quad\bar\varepsilon^\natural_g=\sup_{s\ge s_1(g)}\varepsilon^\natural_g(s),
\]
$C_{\mathsf k}$ of Lemma~\ref{bil:ker}. For a real form $F$ at bonds, resp. sites, with the weight $v$, placed in the same way, $0\le\langle\theta F;F\rangle_\ell\le\frac98[6N^2C_{\mathsf k}^2+c^2\bar\varepsilon^\natural_g]\,b^{-2}\|v\|^2$, $(N,c)=(30,59)$, resp. $(10,23)$. For real tensor weights $H_i$ as in (A2), $|\langle O^T_K(H_1);O^T_K(H_2)\rangle_\ell|\le\frac98[12150C_{\mathsf k}^2+60^2\bar\varepsilon^\natural_g]\,b^{-2}|H_1|_\infty|H_2|_\infty$.
\end{itemize}

\textbf{(B)} Let $g\le\gamma_1$, $s_8(g)\le s(d)\le K$, let $\theta'_1,\theta'_2,\theta'_3$ be mirrors of $T_{\ell'}$, $\theta'_\mu$ normal to $e_\mu$, and let $Y$ be as in (A3), carried within $d-2a_K$ of a point at a distance between $\frac32d$ and $2d$ from a plane of $\theta'_{\mu_0}$ and, up to $a_K/2$, midway between the planes of the two other mirrors. With $\theta'^\eta=\theta'^{\eta_1}_1\theta'^{\eta_2}_2\theta'^{\eta_3}_3$,
\[
\begin{multlined}
0\le M_8=\Bigl\langle\prod_{\eta\in\{0,1\}^3}\theta'^\eta Y\Bigr\rangle_{\ell'}\le C_8\,b^{-8}\sup|\mathsf w|^8,\qquad0\le M_8-\langle\theta'_{\mu_0}Y;Y\rangle_{\ell'}^4 \\
\le(\tfrac98)^459^8\omega^\sharp_{g,8}(s(d))\,b^{-8}\sup|\mathsf w|^8,
\end{multlined}
\]
$C_8=(\frac98)^4[(106C'^2_{\mathsf k})^4+59^8\bar\omega_{g,8}]$, $\bar\omega_{g,8}=\sup_{s\ge s_8(g)}\omega_{g,8}(s)$, with $C'_{\mathsf k}$ of Lemma~\ref{bil:ker} and $\omega_{g,8}$, $\omega^\sharp_{g,8}$ of Lemma~\ref{bil:far}.

\textbf{(C)} Let $g\le\gamma_1$ and let $\nu,\nu'$ be laws at $(K,\beta)$.
\begin{itemize}
\item[(C1)] For $n\ge2$, $d,R>0$, every $K$, and forms $F_1,\dots,F_n$ on $\mathbb L_K$ whose carriers lie in balls of radius $R$ and are pairwise at distance $\ge d$: $|\langle F_1;\dots;F_n\rangle_\nu|\le C^T_n(d;R)\,b^{-n}\prod_a|F_a|_\star$, with $C^T_n$ of Step~4; complex weights cost a factor $2^n$. So 
\[
|S^T_n[\nu](H_1,\dots,H_n)|\le60^nC^T_n(d;R)\,b^{-n}\prod_a|H_a|_\infty
\]
 for supports in such balls at mutual distance $\ge d$, and $|\nu(F)-\nu'(F)|\le2C^T_1(R)\,b^{-1}|F|_\star$.
\item[(C2)] If $s(d)\le K$, $\mathrm{supp}\,H_i\subset B(x_i,d-3a_K)$ and $|x_1-x_2|_\infty\ge3d+a_K$ (on a torus for all lifts): 
\[
|S^T_2[\nu](H_1,H_2)|\le180^2C^\diamond_\infty(g)\,b^{-2}|H_1|_\infty|H_2|_\infty,
\]
$C^\diamond_\infty$ of Step~5. For $Y[\mathsf w_i]$, $\mathsf w_i$ real on the bonds of one direction, carried within $d-2a_K$ of $x_i$: 
\[
|\nu(Y[\mathsf w_1];Y[\mathsf w_2])|\le C^{\rm per}_\infty(g)\,b^{-2}\sup|\mathsf w_1|\sup|\mathsf w_2|,
\]
with a number $C^{\rm per}_\infty\ge e^{2G_g\mathsf l^3}C_8^{1/4}$.
\item[(C3)] Let $s_{16}(g)\le s=s(d)\le K$, $3d\le|x_1-x_2|\le4d$, let the carriers of real forms $F_i$ lie within $d-2a_K$ of $x_i$, $\hat m=\upsilon^2|F_1|_\star|F_2|_\star$, and let $F_i'$ be the forms on $T_{\ell'}$ with the weights of $F_i$ on the board of side $\mathsf l$ whose cell holds the pair in its middle. Then
\[
\begin{multlined}
|\nu(F_i)-\tilde\mu_{\ell'}(F_i)|\le e^{G_g\mathsf l^3}o_{g,16}(s)^{1/8}\upsilon|F_i|_\star,\quad|\nu(F_1;F_2)-\tilde\mu_{\ell'}(F_1;F_2)|\le{\rm Df}_g(s)\,\hat m, \\
|\nu(F_1;F_2)-G(F_1',F_2')|\le[\omega'_{g,16}+{\rm Df}_g](s)\,\hat m,
\end{multlined}
\]
the Gaussian value being formed on $T_{\ell'}$ and ${\rm Df}_g$ given in Step~5; the error function and the first depth are the same for all the laws. The same holds for $F_i=O^T_K(H_i)$, $H_i$ real tensor weights supported within $d-2a_K$ of $x_i$, with $60|H_i|_\infty$ in the place of $|F_i|_\star$.
\end{itemize}

\textbf{(D)} Let $g\le\gamma_1$, let $(K_j,\beta_j)$ be first-passage data at $g$ and $\nu_j$ laws at $(K_j,\beta_j)$: all periodic Gibbs states, or all $\tilde\mu_{\ell_j}$ with $\ell_j\to\infty$. Tuples $X=(H_1,\dots,H_n)$ are in $C_c(\mathbb R^3;V)$ with pairwise disjoint supports.
\begin{itemize}
\item[(D1)] Along a subsequence ($K_j\to\infty$ is not asked) $a_{K_j}\beta_j\to\mathsf b\in I^{\rm lab}(g)$ and $S^T_n[\nu_j](X)\to S^T_n(X)$ for all $n\ge2$ and all $X$; $S^T_n$ is symmetric, multilinear, bounded as in (C1) with $\mathsf b$ for $b$, and in each entry a $V^*$-valued Radon measure off the other supports. For laws $\nu'_j$ at the same data, $\lim(\nu_j-\nu'_j)(O^T_{K_j}(H))$ exists along a further subsequence, of modulus $\le120C^T_1(R)\mathsf b^{-1}|H|_\infty$ if $\mathrm{supp}\,H$ lies in a ball of radius $R$. If each $\nu_j$ is a weak limit of measures $\tilde\mu_\ell$ at $(K_j,\beta_j)$, then for any numbers $\Lambda_j$ there are sides $\ell_j\ge\max\{\Lambda_j,j\}$ in $\mathfrak L(g)$ with $S^T_n[\tilde\mu_{\ell_j}](X)\to S^T_n(X)$ along the full sequence: a limit through limit states of the tori is a limit along growing tori.
\item[(D2)] Over every volume-first limit $S$ and every continuum limit along growing tori there is such an $S^T$, with $S^T_n(f_1\mathbf1,\dots,f_n\mathbf1)=S_n(f_1,\dots,f_n)$, $f_i\in C^1_c$.
\item[(D3)] If $K_j\to\infty$: $S^T_n(rX)=S^T_n(X)$ for $rx=Ax+y$, $A\in W_3$; and for every coordinate plane $\{x_\mu=c\}$, with its reflection $\theta$ acting on points and on indices, all finitely many tuples $X_i$ (the empty one included) with supports in $\{x_\mu>c\}$ and all complex $c_i$: $\sum_{i,k}\bar c_ic_k\sigma^T(\theta X_i\cdot X_k)\ge0$. For limits through periodic Gibbs states the same holds at the planes with the normals $e_\mu\pm e_\lambda$, the Hilbert space is separable, and $H\ge|P_1|$, $H\ge|P_2|$. The numbers $\mathfrak c_j=a_{K_j}^{-4}(\nu_j(a_p)-\tilde\mu_{\ell'}(a_p))$, the same at all plaquettes, obey $|\mathfrak c_j|\le C(g)\mathsf b^{-1}$ for all large $j$; along a further subsequence $\mathfrak c_j\to\mathfrak c$, and along it the moments centred at the means of $\tilde\mu_{\ell'}$ in the place of those of the $\nu_j$ converge too, to $\sum_B\prod_{i\in B}(\mathfrak c\int\mathrm{tr}\,H_i)\,\sigma^T(X_{B^c})$, the sum over the subsets $B$ of the entries, a functional with the same properties whose one-point function is $\mathfrak c\int\mathrm{tr}\,H$.
\item[(D4)] For $\mathsf b\in I(g)$, $\ell\in\mathfrak L(g)$ and continuous $H_a:T_\ell\to V$ with disjoint supports, $S^{T,(\ell,\mathsf b)}_n(H)=\lim_K\langle O^T_K(H_1);\dots;O^T_K(H_n)\rangle_{K,\ell,L^K\mathsf b}$ exists along the full sequence of cutoffs (for $g\le\gamma_2$) and obeys the bound of (D1). Along a subsequence of every $\ell_i\to\infty$ these limits converge, and the limit is also the limit of $S^T_n[\tilde\mu_{\ell_i}]$ at suitable cutoffs $K_i\to\infty$: a limit along growing tori, with the properties of (D3) for these limits.
\item[(D5)] Let $K_j\to\infty$, $s(d)\ge s_{16}(g)$, $3d\le|x_1-x_2|\le4d$, and $H_i$ Lipschitz with compact support in the open ball $B(x_i,d)$. Then
\[
|S^T_2(H_1,H_2)-\mathsf b^{-2}\mathcal W(H_1,H_2)|\le36\,\bar e_g(s(d))\,\mathsf b^{-2}N(H_1)N(H_2),
\]
the function $\bar e_g$ being the same for all sequences; the same holds for the limits of (D4), on every $\ell$. At $H_i=f_i\mathbf1$ the leading term is $\frac9{16\pi^2}\mathsf b^{-2}\iint f_1(x)f_2(y)|x-y|^{-6}dx\,dy$, the scalar leading coefficient in Theorem~\YMMref{AW}{852}.
\end{itemize}

\textbf{(E)} Let $g\le\gamma_2$, $K\ge3$, and let $\ell=2L^m$ be admissible.
\begin{itemize}
\item[(E1)] With $\tilde\nu_K$ the law of the histories and $Z_t$ the large-field region of the level $t$, for $1\le t\le K$:
\[
0\le\mathsf w_t\le\tilde\nu_K(\{Z_t\ne\emptyset\})\le C_\flat M^{-3}\ell^3L^{3(K-t)}\exp\bigl[-\tfrac34c_\star p_0(g_t)^2\bigr];
\]
$C_\flat$ is a function of the constants at two admissible sets of values alone, and does not depend on $g$, $K$, $\beta$, $\ell$, $t$; $\frac34c_\star\ge\frac12c'$, and $\frac34c_\star\ge c'$ if $\eth\ge\frac14$. The same holds with $(c^\natural_2,C_\natural)$ for $(\frac34c_\star,C_\flat)$, $c^\natural_2$ the constant of Definition~\ref{def:rarity-choices} with which Proposition~\ref{prop:local-large-field-rarity} bounds the rarity of a large field at a point, and $C_\natural$ a function of the constants.
\item[(E2)] Let $(c,C)$ be a pair for which (E1) holds with $(c,C)$ for $(\frac34c_\star,C_\flat)$, as $(\frac12c',C_\flat)$, $(\frac34c_\star,C_\flat)$ and $(c^\natural_2,C_\natural)$ do; $\sigma_c=\frac2{\log L}(\frac6{cA_0^2})^{1/P}$, $\hat s(g,\ell)=\max\{s_\flat(g),\lceil\sigma_c(\log\ell)^{1/P}\rceil\}$, the numbers below being formed at $(c,C)$ and $s_\flat(g)$ being the first depth of the envelopes at $(c,C)$, chosen so that also $2c\mathfrak p_0A_0^2y^{P-1}\ge13$ at the working level. Let $\hat s\le s(d)\le K$, $|x_1-x_2|\ge3d$ (no upper bound), and let $F_i$ be sums of forms within $d-2a_K$ of $x_i$, of size the sum of the sizes. Then
(i) $|\langle F_1;F_2\rangle_\ell-G(F_1,F_2)|\le\varepsilon^\circ_g(s(d))\upsilon^2|F_1|_\star|F_2|_\star$, with $Y_s^P\varepsilon^\circ_g(s)$ bounded;
(ii) $|\langle F_1;F_2\rangle_\ell|\le C^\circ_\star(g)b^{-2}|F_1|_\star|F_2|_\star$;
(iii) if $x_1$ is at a distance $\ge\frac32d$ from both planes of $\theta$, $0\le\langle\theta F_1;F_1\rangle_\ell\le C^\circ_\star(g)b^{-2}|F_1|_\star^2$;
(iv) for all $d$ with $s(d)\le K$ and $F$ placed as in (iii), $0\le\langle\theta F;F\rangle_\ell\le C^\circ_3(g,\ell)b^{-2}(|F|_\star/59)^2$, which is $C^\circ_3b^{-2}\sup|\mathsf w|^2$ at $F=Y[\mathsf w]$, $\mathsf w$ real on the bonds of one direction, and $C^\circ_3(g,\ell)\le c_3(g)(1+\log\ell)^8\exp[16(\frac{6\log\ell}{cA_0^2})^{1/P}]$, so that $\ell^{-\epsilon}C^\circ_3(g,\ell)$ is bounded for every $\epsilon>0$;
(v) if all points of the carriers of $F_1$ and $F_2$ are $\ge r$ apart and $K\ge s(\min\{\frac r8,L^{-\hat s}\})$: $|\langle F_1;F_2\rangle_\ell|\le(\ell\max\{\frac8r,L^{\hat s}\}+1)^6C^\circ_\star(g)b^{-2}|F_1|_\star|F_2|_\star$.
On the sides with $\log\ell\le(s_\flat(g)/\sigma_c)^P$ the first depth is $\hat s=s_\flat(g)$, free of $\ell$.
\item[(E3)] Let $\mathsf L<\infty$ and $c''\ge\frac12c'$ be numbers with $\mathsf w_t\le\mathsf LL^{3(K-t)}e^{-c''p_0(g_t)^2}$ for $1\le t\le K$. Then: (i) Lemma~\ref{bil:nsrc}, for sources under majorants, holds on $T_\ell$ with $\max\{1,M^3\mathsf L/(c'A_0^2)\}^nk'_n$ for $k'_n$, under its hypotheses $n\ge2$ and $s_{\rm th}(n)\le s(d)\le K$ and in its geometry: $|\langle F_1;\dots;F_n\rangle_\ell|\le\max\{1,M^3\mathsf L/(c'A_0^2)\}^nk'_n(g,s(d))\,b^{-n}\prod_a\mathsf s_a$, with no factor of a covering; (ii) from (i), by the covering of Step~1, (A1) holds on $T_\ell$ with $2\ell$ for $2\ell_{\max}$ and $\max\{1,M^3\mathsf L/(c'A_0^2)\}^nk'_n$ for $k'_n$. Two such pairs are had on every admissible side: $\mathsf L=2M^{-3}\ell^3e^{\mathsf B(\bar x_0)\ell^3}$, $c''=c'$, by the rarity of large fields on one torus (Proposition~\ref{prop:ward-fixed-volume-rarity}), as $\mathsf B$ increases and $x_K\le X_g<\bar x_0$; and the pair of (E1), $\mathsf L=C_\flat M^{-3}\ell^3$, $c''=\frac34c_\star\ge\frac12c'$. So (i) and (ii) hold on every admissible side with $\max\{1,c(\ell)\}^nk'_n$ for $k'_n$, (ii) with $2\ell$ for $2\ell_{\max}$, where $c(\ell)=2\ell^3e^{\mathsf B(\bar x_0)\ell^3}/(c'A_0^2)$, the number $M^3\mathsf L/(c'A_0^2)$ at the first pair; with the weights of (E1), $c(\ell)$ may be replaced by $C_\flat\ell^3/(c'A_0^2)$, a power of $\ell$.
\end{itemize}
\end{theorem}

\begin{theorem}\label{bil:lim}
Let $g\le\gamma_1$ unless another range is given, and let $\sigma=\sigma^T$, with cumulants $S^T_n$, be a limit through periodic Gibbs states.
\begin{itemize}
\item[(F1)] (\emph{The free law in sup norms.}) $\mathsf b=\lim_ja_{K_j}\beta_j$ exists and is fixed by $\sigma$. With $d_0(g)=L^{-s_{16}(g)}$ there is a function $e^\infty_g$ on $(0,d_0(g))$, given in Step~8, with $(\log\frac1d)^{\mathfrak p_0/4}e^\infty_g(d)$ bounded, such that for $0<d_-<d<d_0(g)$, $3d\le|x_1-x_2|\le4d$ and $H_i\in C^1_c$ supported in $\bar B(x_i,d_-)$:
\[
|S^T_2(H_1,H_2)-\mathsf b^{-2}\mathcal W(H_1,H_2)|\le e^\infty_g(d)\,\mathsf b^{-2}|H_1|_\infty|H_2|_\infty;
\]
the same holds for $S^{T,(\ell,\mathsf b)}_2$ of (D4), every $\ell\in\mathfrak L(g)$, $\mathsf b\in I(g)$, Lipschitz $H_i$; and both bounds hold with $C(g)(\log\frac1d)^{-2\mathfrak p_0}$ in the place of $e^\infty_g(d)$. At a cutoff, for $K\ge s(d)+2$, $4a_K\le d-d_-$, continuous $H_i$ and every law $\nu$ at a datum $(K,\beta)$: $|\nu(O^T_K(H_1);O^T_K(H_2))-b^{-2}\mathcal W(H_1,H_2)|\le b^{-2}\{[e^\infty_g(d)+C(g)(a_K/d)^2]|H_1|_\infty|H_2|_\infty+r_K(H_1,H_2)\}$, $r_K$ free of $\beta$, $\ell$, $\nu$, and $r_K\to0$ as $K\to\infty$. For $R\in O(3)$ and $rx=x_1+R(x-x_1)$, $|\sigma(rH_1,rH_2)-\sigma(H_1,H_2)|\le10e^\infty_g(d)\mathsf b^{-2}|H_1|_\infty|H_2|_\infty$. If $e^\infty_g(d)<w_0=\frac1{12}388^{-6}$, $\theta$ is the reflection in $\{x_{\mu_0}=c\}$, $|c-x_{0,\mu_0}|=\frac32d$, and $\varphi\in C^1_c$, $\mathbf1_{B(x_0,d/128)}\le\varphi\le\mathbf1_{B(x_0,d/64)}$, then $\sigma(\theta(\varphi E),\varphi E)\ge(w_0-e^\infty_g(d))\mathsf b^{-2}\mathrm{tr}\,E^2$ for every $E\in V$: no nonzero real combination of the six components vanishes.
\item[(F2)] (\emph{Powers of the distance; growth; one insertion.}) Let $g\le\gamma_2$, $n\ge2$, $q_E=(4+\mathfrak p_0)(P-1)+10\mathfrak p_0$, $\delta_n=\frac1{8n}$. There are $c_g$, $C_{\rm sn}$ and $i_E(g,n)$ such that, for integers $i\ge\max\{0,i_E(g,n)\}$, $Y_i=c_g(1+i+\log n)$, $K_E(i,n)=i+\max\{s_E(i,n),\lceil\log_L4n\rceil\}$, where $s_E(i,n)$, the depth $s_E(i)$ of Lemma~\ref{bil:semi} at $\delta=\delta_n$, is the depth $s(\delta_n)$ formed at the coupling $g^\sharp_i$, $(g^\sharp_i)^{-2}=L^{i-1}(X_g+\bar B)-B_\sharp$, so that $K_E(i,n)$ is a number of the constants, $g$, $i$ and $n$ (the coupling $g_{(i)}$ stopped at the level $i$ of a datum, $g_{(i)}^{-2}=L^{-1}x_{K-i}-B_\sharp$, has $g_{(i)}^{-2}\le(g^\sharp_i)^{-2}$, as $x_{K-i}\le L^i(x_K+\bar B)\le L^i(X_g+\bar B)$, and $C_W(s)$, so the depth, does not decrease as the inverse coupling grows; so, for $K\ge K_E(i,n)$, the depth $s(\delta_n)$ formed at $g_{(i)}$ is at most $s_E(i,n)\le K-i$), every datum with $K\ge K_E(i,n)$ and every law $\nu$:
(i) on every board of side $a_i$ and every torus of $\mathfrak L(g)$ the chessboard seminorm of $A[O^T_K(H)]$, $H$ bounded and $0$ within $\delta_na_i$ of the complement of a cell, is at most $C_{\rm sn}n^4Y_i^{q_E}L^{11i/2}b^{-1}|H|_\infty$;
(ii) for $H_1,\dots,H_n$ with supports of sup diameter $\le\frac{a_i}{4n}$ at pairwise sup distances $>a_i$: $|\nu(\prod_aA[O^T_K(H_a)])|\le(2c_gn^4Y_i^{q_E}L^{11i/2})^n\prod_a|H_a|_\infty$.
For $g\le\gamma_1$:
(iii) at separated tuples with pairwise sup distances $\ge2a_i$, $|\sigma(H_1,\dots,H_n)|\le\prod_a2c_gn^4Y_i^{q_E}L^{11i/2}(4n/a_i)^3\|H_a\|_\natural$, where $\|H\|_\natural=\sum_{\mathfrak b}|\mathfrak b|\sup_{\mathfrak b}|H|_\infty$ over the boxes $\mathfrak b=h(\mathbf n+[0,1)^3)$, $\mathbf n\in\mathbb Z^3$, $h=\frac{a_i}{4n}$;
(iv) there is $C_G$ such that for all $n\ge2$, $0<d\le1$ and tuples at pairwise sup distances $\ge2d$, ordered or not, $|\sigma(H_1,\dots,H_n)|\le(C_Gn^{16}d^{-9})^n\prod_a\|H_a\|_\natural\le(e^{16}C_G)^n(n!)^{16}d^{-9n}\prod_a\|H_a\|_\natural$, the boxes being those of $i=\max\{\lceil\log_L\frac1d\rceil,i_E(g,n)\}$;
(v) so $\sigma$ extends to exactly one complex linear $\mathsf S^\sigma_n$ on the $V_{\mathbb C}^{\otimes n}$-valued functions on $\mathbb R^{3n}$ whose components are Schwartz functions flat at coinciding points that is bounded for one of the norms $|\varphi|_k=\max_{\mathbf e}\sup_{|\alpha|\le k}(1+|x|^2)^{k/2}|\partial^\alpha\varphi_{\mathbf e}(x)|$, and $|\mathsf S^\sigma_n(\varphi)|\le C_{\rm gr}^n(n!)^{24}|\varphi|_{17n}$: the growth hypothesis in Proposition~\YMMref{93.5}{719} for each of the scalar functionals $\sigma(f_1E_{\mathbf e_1},\dots,f_nE_{\mathbf e_n})$, $\mathbf e_a$ among the six pairs $\lambda\le\lambda'$, which obey (iv) with the same numbers; a $V_{\mathbb C}^{\otimes n}$-valued function is the sum of its $6^n$ components times the tensors $E_{\mathbf e_1}\otimes\dots\otimes E_{\mathbf e_n}$, which costs the factor $6^n$ in $C_{\rm gr}^n$;
(vi) for $\eta\in(0,1]$, combinations $F,F'$ of tuples in $\{x_3>\eta\}$, $h\in C_c(\{|x_3|<\frac\eta2\};V_{\mathbb C})$ and every $p>\frac{17}2$: $|\sigma(\theta\bar F'\cdot(h)\cdot F)|\le C_p\mathsf b^{-1}\eta^{-p}\|h\|_1\|[\tau_{-\eta e_3}F']\|\|[\tau_{-\eta e_3}F]\|$, $\|h\|_1=\int\max_{\lambda\lambda'}|h_{\lambda\lambda'}|$; this holds for $g\le\gamma_2$;
(vii) for $g\le\gamma_2$, $\Omega$ and the classes of the tuples with entries in pairwise disjoint slabs $\{c<x_3<c'\}$ are total in the Hilbert space;
(viii) given (F3), $\sigma$ has one continuous kernel $\mathsf s_n$ on the configurations of pairwise different points, $\sigma(X)=\int\mathsf s_n(x;H_1(x_1),\dots,H_n(x_n))dx$ at all separated tuples, with $|\mathsf s_n(x;e)|\le(cC_p\mathsf b^{-1})^n\max\{1,2n^2/\varrho(x)\}^{9n}\prod_k|e_k|$, $\varrho(x)$ the least distance of two points, and $|\mathsf S^\sigma_n(\varphi)|\le C^n(n!)^{18}|\varphi|_{9n}$.
\item[(F3)] (\emph{Rotations.}) $\sigma(RX)=\sigma(X)$ and $S^T_n(RX)=S^T_n(X)$ for all $R\in O(3)$, $(RH)(x)=RH(R^{-1}x)R^{\rm T}$, at all separated tuples with continuous entries. For $0<g\le\gamma$ the same holds for every limit of a sequence of periodic Gibbs states that converges at all tuples; on the entries $f\mathbf1$ it follows also from the invariance in Theorem~\YMMref{AN}{763}, applied to the functional centred at the means of $\tilde\mu_{\ell'}$ (Step~11). For real $\zeta\ne1$ let $\Lambda_\zeta$ be the linear map of $V$ that is $1$ on the diagonal and $\zeta$ on the off-diagonal matrices: $\sigma(\Lambda_\zeta\,\cdot)$ has all the properties of (D1) and (D3), the cone of (F4), the insertion bound of (F2) and the plateau of (F6), and it is not invariant under $O(3)$ (Proposition~\ref{bil:zeta}).
\item[(F4)] (\emph{Mirrors, cone.}) For $g\le\gamma_2$ the joint spectral measure of $(H,P_1,P_2)$ is carried by $\{E\ge|q|\}$, $|q|$ the Euclidean length, in each of the nine pictures, and $\int e^{-\mathfrak z_3E+i(\mathfrak z_1q_1+\mathfrak z_2q_2)}d\mathcal P$ is a holomorphic semigroup of contractions on the tube $\{\mathrm{Re}\,\mathfrak z_3>|(\mathrm{Im}\,\mathfrak z_1,\mathrm{Im}\,\mathfrak z_2)|\}$. For $0<g\le\gamma$, $\sigma(\theta_{n,c}\bar F\cdot F)\ge0$ for every plane $\{n\cdot x=c\}$ and every combination $F$ of tuples in $\{n\cdot x>c\}$, the labels turned by $1-2nn^{\rm T}$.
\item[(F5)] (\emph{Conservation.}) Let $0<g\le\gamma$, $H^{\rm cons}[\eta]=(\nabla\cdot\eta)\mathbf1-\nabla\eta-(\nabla\eta)^{\rm T}$, $(\nabla\eta)_{\nu\rho}=\partial_\nu\eta^\rho$, for $\eta\in C^2_c(\mathbb R^3;\mathbb R^3)$, $\|\eta\|_{C^2}=\sup|\eta|+\sup|\nabla\eta|+\sup|\nabla^2\eta|$; let $q\ge2$ and $H_2,\dots,H_q$ be bounded, the $q$ supports in one cube of side $2R$, pairwise at sup distance $\ge d$. For $K\ge K_0$ and every periodic Gibbs state $\nu$,
\[
\Bigl|\nu\Bigl(A[O^T_K(H^{\rm cons}[\eta])]\prod_{a\ge2}A[O^T_K(H_a)]\Bigr)\Bigr|\le CK^pL^{-K/36}\|\eta\|_{C^2}\prod_{a\ge2}|H_a|_\infty,
\]
$C,p,K_0$ of the constants, $g,q,d,R$ (at every $K\ge z+1$ the dilation of Lemma~\ref{bil:dil}(d), $z=z(g,d)$, carries every datum at $g$ to one at one of two couplings $g'_1,g'_2\le\gamma_1$ fixed by $g$ and $d$, with all its periodic Gibbs states, boards and seminorms); the same for the joint cumulant; on $T_\ell$ with numbers depending on $\ell$. For $0<g\le\gamma$, $\ell\in\mathfrak L(g)$, every datum $(K,\beta)$ at $g$ with $K\ge K_0$, $\eta\in C^2$ on $T_\ell$ and bounded $H_a$, the supports pairwise at sup distance $\ge d$ on $T_\ell$, the joint cumulant on $T_\ell$ is at most $CK^pL^{-e_qK}\|\eta\|_{C^2}\prod_a|H_a|_\infty$ with $e_2=\frac{51}{2800}$, $e_3=\frac{23}{2800}$, $e_q=\frac1{56}$ for $q\ge4$, and $C,p,K_0$ of the constants, $g$, $\ell$, $q$, $d$. So $\sigma(H^{\rm cons}[\eta]\cdot X)=0$ for every limit of periodic Gibbs states converging at all tuples and every $X$ with continuous entries whose compact supports are pairwise disjoint and off $\mathrm{supp}\,\eta$; and $S^{T,(\ell,\mathsf b)}_q(H^{\rm cons}[\eta],H_2,\dots,H_q)=0$ for $g\le\gamma_2$, $\mathsf b\in I(g)$, $\ell\in\mathfrak L(g)$, $\eta\in C^2$ on $T_\ell$, Lipschitz $H_a$, pairwise disjoint supports.
\item[(F6)] (\emph{Plateau; components.}) For combinations $F,G$ of tuples in $(-\rho_0,\rho_0)^3$ the limit $m_2(\bar F,G)=\lim_{|v|\to\infty}\sigma(\bar F\cdot\tau_vG)$ exists, with $|\sigma(\bar F\cdot\tau_vG)-m_2(\bar F,G)|\le4C_{\rm dc}\mathcal N_{\rho_0}(F)^{1/2}\mathcal N_{\rho_0}(G)^{1/2}\rho_0/|v|_\infty$ for $|v|_\infty>4\rho_0$, where $\mathcal N_{\rho_0}(F)=\max_\mu\max_{\mathsf g}\sigma(\theta_\mu\mathsf g\bar F\cdot\tau_{2\rho_0e_\mu}\mathsf gF)$ over the eight sign changes $\mathsf g$ and the coordinate mirrors $\theta_\mu$ through $0$, $C_{\rm dc}$ the absolute number in equation~\YMMref{(1968)}{742}; and $m_2(\bar F,F)\ge0$. For $g\le\gamma_2$ there are a compact metrisable space $\mathsf X$, a Borel probability measure $\lambda$ on it of full support, and real symmetric functionals $\sigma_\xi$, with $\xi\mapsto\sigma_\xi(X)$ continuous and $\xi\mapsto\sigma_\xi$ injective, such that $\sigma=\int\sigma_\xi\,\lambda(d\xi)$ and $\int\prod_{i\le q}\sigma_\xi(X^i)\,\lambda(d\xi)=\lim\sigma(X^1\cdot\tau_{v_2}X^2\cdots\tau_{v_q}X^q)$, the $|v_k|\to\infty$ one after the other; every $\sigma_\xi$ is the limit, along the full sequence, of the $\nu_j$ conditioned on invariant events of positive measure, and re-centred in these conditioned states it is itself a limit through periodic Gibbs states; $\ker H_\xi=\mathbb C\Omega_\xi$; for a canonical decomposition $(\mathsf X_0,\lambda_0,\sigma_{0,\xi_0})$ of the limit functional of the action density along the same sequence (Theorem~\textup{\YMMref{AP}{780}}) there is exactly one map $\pi:\mathsf X\to\mathsf X_0$ with $\sigma_\xi(f_1\mathbf1,\dots,f_n\mathbf1)=\sigma_{0,\pi(\xi)}(f_1,\dots,f_n)$ for all $\xi$ and all scalar tuples; it is continuous and onto, $\pi_*\lambda=\lambda_0$, and $\sigma_\xi((H))=\frac13c_{\pi(\xi)}\int\mathrm{tr}\,H$, $c_{\xi_0}$ the one-point constant of the component $\xi_0$.
\item[(F7)] (\emph{One limit, one vacuum.}) For $0<g\le\gamma$ and $\mathsf b\in I(g)$: the cumulants of the action density converge as $K\to\infty$ for every choice of states in $\mathcal G^{\rm per}_{K,L^K\mathsf b}$ if and only if all convergent sequences of periodic Gibbs states at $g$ with $a_{K_j}\beta_j=\mathsf b$ for all $j$ have the same $S_n$, $n\ge2$; then, for the limit functional $\sigma_0$ of the action density and every canonical decomposition $(\mathsf X_0,\lambda_0,\sigma_{0,\xi})$ of it, $\sigma_{0,\xi}=\sigma_0$ for $\lambda_0$-almost every $\xi$, $\mathsf X_0$ is a point, and $\ker H_{\sigma_0}=\mathbb C\Omega$, $H_{\sigma_0}$ the Hamiltonian reconstructed from $\sigma_0$.
\item[(F8)] (\emph{Coinciding points.}) Let $g\le\gamma_2$, $\mathsf b\in I(g)$, $\ell\in\mathfrak L(g)$, $n\ge2$. $S^{T,(\ell,\mathsf b)}_n$ extends to a symmetric, translation invariant distribution $\tilde S_n$ on $T_\ell^n$ of order $\le3n-2$ with
\[
|\tilde S_n(F)|\le C\mathsf b^{-n}\bigl(1+\log\tfrac{h_0}\lambda\bigr)^{8n\mathfrak p_0+1}\max_{|\gamma|\le3n-2}\lambda^{|\gamma|}\sup|\partial^\gamma F|
\]
for $F$ supported in a product of closed sup-balls of radius $\lambda\le h_0$ about any points, coinciding or not ($C,h_0$ of $g,\ell,n$); the lattice cumulants with their Taylor parts at the diagonals subtracted (Lemma~\ref{bil:ext}) obey this bound for all $K$ and all $8h_K\le\lambda\le h_0$, $h_K\le2C_W(K)a_K$ the least scale at the cutoff $K$, with one number for all cutoffs, and they converge to $\tilde S_n$ at every smooth $F$; two extensions with this bound differ by a sum $\sum_{|\pi|\le n-1}\sum_{|\alpha|\le3|\pi|}\partial^\alpha_\xi\delta_\pi\otimes T_{\pi,\alpha}$ of contact terms, and every such sum may be added; at $n=2$ the translation invariant, real, symmetric contact terms covariant under the cubic group form a space of dimension twelve, under $O(3)$ on $\mathbb R^3$ one of dimension six.

The reconstruction (the Hilbert space, the semigroup, the momenta and the kernels of Lemma~\ref{bil:vec}) uses the values of $\sigma$ at separated tuples alone, so no choice of the extension and no contact term enters it; beyond the separated tuples it rests on the insertion bound of (F2)(vi). For $g\le\gamma_1$ every limit point of $\langle Y[\varphi_1];Y[\varphi_2]\rangle_{K,\ell,\beta}$ as $K,\ell\to\infty$, $\ell\in\mathfrak L(g)$, in any order, is bounded at separated cubes by $36C^{\rm od}_\infty(g)(\inf b)^{-2}\sup|\varphi_1|\sup|\varphi_2|$, with no power of the distance and no logarithm, and so extends to $\mathbb R^6$ as a distribution of order $\le4$.
\item[(F9)] (\emph{Mirror bound in every state; the power $6$.}) Let $s(d)\le K$. If $X=A[O^T_K(H)]$ is carried in $B(x_0,d-2a_K)$ and $x_0$ is at a distance $\ge\frac32d$ from the plane of $\theta$, then $0\le\nu(\theta X\cdot X)\le C_\infty(g)\,b^{-2}|H|_\infty^2$ in every periodic Gibbs state $\nu$, $C_\infty(g)$ of Step~9; for a bond form $Y=Y[\mathsf w]$ with a complex weight carried there, $\nu(Y)=0$ and $0\le\nu(\overline{\theta Y}Y)\le36C^{\rm od}_\infty(g)b^{-2}|\mathsf w|_\infty^2$, with $C^{\rm od}_\infty(g)$ alone for real weights of one direction. For $H_1,H_2\in C_c(\mathbb R^3;V)$ with disjoint supports and off-diagonal values,
\[
\begin{multlined}
|\sigma(H_1,H_2)|\le A_6C^{\rm od}_\infty(g)\,\mathsf b^{-2}\iint\|H_1(x)\|_1\|H_2(y)\|_1\min\{|x-y|_\infty,3\}^{-6}dx\,dy,\qquad\|E\|_1 \\
=\sum_{\lambda<\lambda'}|E_{\lambda\lambda'}|,
\end{multlined}
\]
$A_6=64\cdot36\cdot((8/\pi^3)32^3)^2$; for the whole tensor the same holds with a pure number $A$ in the place of $A_6$, with $C_\infty(g)$ in the place of $C^{\rm od}_\infty(g)$ and with $\|E\|_1=\sum_{\lambda\le\lambda'}|E_{\lambda\lambda'}|$; by (D5) the power $6$ is the true one.
\item[(F10)] (\emph{Larger couplings.}) If $\gamma_1<g\le\gamma_2$, there are $z=z(g)$ and $g'_1,g'_2\le\gamma_1$ such that for $K\ge z+1$ the dilation by $L^z$ carries every law at $(K,\beta)$ at $g$ to one at $(K-z,\beta)$ at $g'_1$ or $g'_2$; so (C1) for tensor entries, (D1)--(D5) and (F1) hold at $g$, with $d_0(g)=L^{-z}\min_sd_0(g'_s)$ and $e^\infty_g(d)=\max_se^\infty_{g'_s}(L^zd)$, and (C1) for all forms holds with numbers depending on $\ell'$. If $\gamma_2<g\le\gamma$, let $z$ be the least $z\ge1$ with $g'_1\le\gamma_1$ in Lemma~\ref{bil:dil}(d); then $g'_2<g'_1\le\gamma_1$, and for $K\ge z+1$ the dilation by $L^z$ carries every periodic Gibbs state at $(K,\beta)$ at $g$ to one at $(K-z,\beta)$ at $g'_1$ or $g'_2$, so (C1) for tensor entries holds at $g$ in every periodic Gibbs state at $(K,\beta)$, with $\max_\iota C^T_n(L^zd;L^zR)$, formed at $g'_\iota$, in the place of $C^T_n(d;R)$. If $0<g\le\gamma$, every limit $\sigma$ of a sequence of periodic Gibbs states at $g$ that converges at all tuples, with a limit point $\mathsf b$ of its labels, is $\sigma(X)=L^{nz}\sigma_z((H_a(L^{-z}\cdot))_a)$, $\sigma_z$ such a limit at $g_z=(L^z\mathsf b/2)^{-1/2}\le\gamma_1$ ($z$ large), and it has the properties (D1), (D3), (F1), (F3), (F4), with numbers depending on $\mathsf b$.
\end{itemize}
\end{theorem}

\subsection{Construction estimates}

We use Theorem~\YMMref{D}{8}, at the common numerical choice of Proposition~\YMMref{69.1}{595} and the admissible histories of Corollary~\YMMref{69.4}{604} and the stability bound, Theorems~\YMMref{AJ}{673} and~\YMMref{AD}{607}, with the results among the construction results on the step map and its sizes (\YMMref{(759)}{375}, Proposition~\YMMref{45.1}{353} and Lemma~\ref{lem:analytic-domain-geometry}), the height-zero source step (Proposition~\ref{prop:source-entry} and the first-step part of Lemma~\ref{lem:weighted-array-step}), the flow up to the first passage (\YMMref{74.1}{615}, \YMMref{(1523)}{614} and~\YMMref{(1525)}{615}), the barycentre interpolation by tents in Lemma~\ref{bil:sym}(e) and the running-coupling estimate in the proof of
Lemma~\YMMref{125.4}{841}, with its exponent replacement
explained in the proof of Lemma~\ref{bil:two}(b). Theorem~\YMMref{D}{8} enters every clause, and parts (A1), (B), (E1) and (E3) rest on it alone. Theorems~\YMMref{AJ}{673} and~\YMMref{AD}{607} enter only through the constant $G_g$ of Proposition~\ref{bil:press}(a), which may be replaced in every bound by the constant $G^F_g$ obtained there from Theorem~\YMMref{D}{8} alone, through statements on periodic Gibbs states, through the decomposition of Theorem~\textup{\YMMref{AP}{780}}, and, in (E2), through the chain for two sources on an admissible side. No stability bound with a tensor-valued source is needed.

\subsection{Estimates and proofs}

\begin{lemma}[Matrices and corner forms]\label{bil:mat}
(a) For complex $2\times2$ matrices, with the operator norm: $|\langle X,Y\rangle|\le|X||Y|$; if $\det V=\det V'=1$ and $T$ is invertible, then $w(TVT^{-1})=Tw(V)T^{-1}$, $w(V^{-1})=-w(V)$, $|V^{-1}-\mathbf1|=|V-\mathbf1|$ and $|V^{-1}-V'^{-1}|=|V-V'|$; so $|w(V)|\le|V-\mathbf1|$ and $|w(V)-w(V')|\le|V-V'|$.

(b) For loops $l,l'$ of one cell the corner form $Q=\langle w(W_l),w(W_{l'})\rangle$ is a polynomial in the bond variables, $0$ at the trivial configuration, gauge invariant, and real at $\mathrm{SU}(2)$ configurations, where $|Q|\le a_{p(l)}+a_{p(l')}$. Let $\eta,b_0,b_1>0$ be arbitrary. At complex configurations with $|U(\mathsf e)|\le e^{\eta b_1}$ on the bonds of the cell and $|\partial U(p)-\mathbf1|<b_0\eta^2$ at its two plaquettes, $|Q|\le4e^{4t\eta b_1}b_0^2\eta^4$, $t=1$ at bonds and $t=2$ at sites. If $|U(\mathsf e)|\le e^{1/13}$ at the bonds of $p=p(l)$, $p'=p(l')$ and at their reverses, $|\partial U(p)-\mathbf1|\le\zeta_0$ and $|\partial U(p')-\mathbf1|\le\zeta_0$ with $e^{4/13}\zeta_0\le1$, and $B'$ is a $1$-form with $\rho=\max|B'|\le\frac1{128}$, then $|Q(e^{iB'}U)-Q(U)|\le784\rho^2+112e^{2t/13}\rho\zeta_0$. At the small points of a chart of height $i$, $|Q-\sigma(l)\sigma(l')\langle\mathrm{ang}_{p(l)},\mathrm{ang}_{p(l')}\rangle|\le12M_{\rm pl}^2\eta_i^4\mathsf{rt}^{-3}$, where $\mathrm{ang}_p$ is the angle of $p$ and $\mathsf{rt}\ge2$ is the quotient of the radii of the chart over the point (Lemma~\YMMref{124.4}{835}).
\end{lemma}

\begin{proof}[Proof sketch]
(a) The trace of a $2\times2$ matrix is the sum of two entries, each at most the norm. For $\det V=1$, $V^{-1}=(\mathrm{tr}\,V)\mathbf1-V$; so $V^{-1}-\mathbf1=(\mathrm{tr}\,E)\mathbf1-E$ with $E=V-\mathbf1$, and $V^{-1}-V'^{-1}=(\mathrm{tr}\,E)\mathbf1-E$ with $E=V-V'$; for $2\times2$ matrices $(\mathrm{tr}\,E)\mathbf1-E=JE^{\rm T}J^{-1}$ with the unitary $J=\begin{pmatrix}0&1\\-1&0\end{pmatrix}$, and $|E^{\rm T}|=|E|$; finally $w(V)=[(V-\mathbf1)-(V^{-1}-\mathbf1)]/2i$. (b) For $V=\exp(i\phi\,n\cdot\sigma)\in\mathrm{SU}(2)$, $w(V)=\sin\phi\;n\cdot\sigma$ and $\langle w(V),w(V)\rangle=\sin^2\phi\le2(1-\cos\phi)$, which with the Schwarz inequality for $\langle\cdot,\cdot\rangle$ gives $|Q|\le a_{p(l)}+a_{p(l')}$. The complex bounds follow from (a) applied to the holonomies of the two plaquettes; the bound on a shift and the second order follow from the expansion of the holonomy in the $1$-form and in the chart, $w(W_l)=\sigma(l)\,\mathrm{ang}_{p(l)}$ to the first order; every bound of a vector part at a complex configuration goes through (a).
\end{proof}

\begin{lemma}[Forms are sources]\label{bil:forms}
(a) Let $\bar k\ge1$ be an integer, the \emph{top index}. Call $F$ a \emph{source of top index $\bar k$ under the majorant} $(\mathsf s,S)$, $\mathsf s\ge0$ a number and $S$ a set of plaquettes, if $F$ is continuous, gauge invariant, $|F|\le L^K\mathsf s\sum_{p\in S}a_p$, and $F$ is the sum, at height $0$, of a family in the fourth place over a point (Section~\ref{sec:step-input-conventions}) of size $\le c^\flat\mathsf m_0$, $c^\flat=81C_P^2C_{\mathcal A}^2/K_\Lambda$, and, at each height $1\le j\le\bar k$, of a family without base point (Definition~\YMMref{49.5}{375}) of size $\le\frac12\mathsf m_j$, $\mathsf m_j=12M^3\alpha_0^2L^{K-j}\mathsf s$, both vanishing on the regions that miss $S$. The top index is that of the data of the chain at which the source is used: the working level $\mathsf k$ for $n$ sources (Lemmas~\ref{bil:nsrc} and~\ref{bil:prod}, where the chain runs up to the level $\mathsf k+1$), and $T-1$, $T$ the upper level, for two sources and for far pairs (Lemmas~\ref{bil:two} and~\ref{bil:far}), where the one-step covariance of the level $j$, $\mathsf k<j\le T$, is formed with the members of the height $j-1$; since $n_T\ge2$, there $T-1>\mathsf k$.

A source is a source of the top index of the data at hand. Call it \emph{quadratic}, with the kernel $\mathsf t$, if besides $\mathsf t$ is a real symmetric function on $S\times S$ with $2\sum_{p'}|\mathsf t(p,p')|\le\mathsf s$ such that $|a_KF-\sum\mathsf t(p,p')\langle\mathrm{ang}_p,\mathrm{ang}_{p'}\rangle|\le M_{\rm pl}^2\eta_i^4\mathsf{rt}^{-3}(\#S)\mathsf s$ at the small points of the chart of height $i$. Then every form $F$ is a source under $(|F|_\star,S[F])$ of every top index, since it has its member at every height $j\ge1$ of the chain, linearly in its weights, and it is quadratic with the kernel $\mathsf t_F$ if it is real; a plaquette form with the weight $h$ is one under $(\sup|h|,\mathrm{supp}\,h)$, with $\mathsf t=\frac12h\delta$; for real weights at bonds or sites $2\sum_{p'}|\mathsf t_F(p,p')|\le8\|v\|$; $O^T_K(H)$ is one under $(60|H|_\infty,S_H)$, $S_H$ the plaquettes of its six forms, with $2\sum_{p'}|\mathsf t_H(p,p')|\le3|H|_\infty$.

Under sums the sizes add and the sets unite, and a source under $(\mathsf s,S)$ is one under $(\mathsf s',S)$ for every $\mathsf s'\ge\mathsf s$.

(b) The proof of the two-point law of the action density on a torus of small side (Theorem~\YMMref{AU}{847}), and the proofs of the statements on which it rests (Theorem~\YMMref{AT}{839}, Proposition~\YMMref{125.1}{838}, Lemmas~\YMMref{125.2}{840}, \YMMref{125.4}{841} and~\YMMref{125.5}{842}; the weighted progression of Proposition~\textup{\YMMref{122.1}{826}}; the first step and the term of the history without large-field regions, the first-step part of Lemma~\ref{lem:weighted-array-step} and Proposition~\YMMref{123.1}{829}; the quotients of the sources and of a step, their tower and the reduction of the covariance to one-step covariances, Proposition~\YMMref{123.1}{829} and equation~\YMMref{(2268)}{829}, the quotient definitions in Section~\ref{sec:weighted-covariance-interface}, the quotient identities preceding equation~\YMMref{(2273)}{830}, equation~\YMMref{(2273)}{830}, Lemma~\ref{lem:weighted-step-covariance} and Proposition~\YMMref{123.2}{830}; the carried family and the flat top, Proposition~\YMMref{123.3}{832}, equation~\YMMref{(2283)}{832}; and the sum over the level of the first creation, equation~\eqref{eq:ward-first-creation}), use the source only through the properties of (a) and its boundedness; they therefore apply to every source in the sense of (a) whose top index is that of the data of the statement applied, and to every quadratic source where the kernel is used.

(c) The step map and its Lipschitz bound hold for forms: for weights $v$, $v'$ at the cells the outputs of one step differ by at most $\theta^{[j]}8n_{\mathfrak c}e^{4t_{\mathfrak c}/13}M^3\alpha_0^2\|v-v'\|$ in the norms of the families of Definition~\YMMref{49.5}{375}, a pure number times the number for plaquette weights ($n_{\mathfrak c}=18$, $t_{\mathfrak c}=1$ at bonds; $n_{\mathfrak c}=7$, $t_{\mathfrak c}=2$ at sites); the step map depends only on its input array.
\end{lemma}

\begin{proof}[Proof sketch]
(a) The cubes of the partition of height $k$ that hold a corner of a plaquette of a cell form a connected region of at most $6$ (bond) or $7$ (site) cubes with a common corner, and a region of height $k$ serves at most $n_{\mathfrak c}M^3\eta_k^{-3}$ cells. At height $0$, with $\mathbf U$ the configuration of the bond variables and $B'$ the $1$-form in the fourth place, of which the family is a function, the family $\sum_{X_0(\mathsf c)=Y}[f_{\mathsf c}(e^{iB'}\mathbf U)-f_{\mathsf c}(\mathbf U)]$, $f_{\mathsf c}=\sum_{l,l'}v(\mathsf c,l,l')\langle w(W_l),w(W_{l'})\rangle$, $X_0(\mathsf c)$ the region of the cell, is holomorphic and gauge covariant, and Lemma~\ref{bil:mat}(b) bounds it at complex configurations; with the count, $n_{\mathfrak c}\cdot3138/(81\cdot12)\le59$, resp. $23$, it gives the size $c^\flat\mathsf m_0$ with the factor $59$, resp. $23$, of $|F|_\star$.

At height $j\ge1$ the member over a region $X$ is $\eta_j^{-1}\sum_{X_j(\mathsf c)=X}f_{\mathsf c}$, of size $\le\frac12\mathsf m_j$ by the same count. To second order $w(W_l)=\sigma(l)\,\mathrm{ang}_{p(l)}$, whence $\sum\mathsf t_F(p,p')\langle\Phi_p,\Phi_{p'}\rangle=\sum v(\mathsf c,l,l')\sigma(l)\sigma(l')\langle\Phi_{p(l)},\Phi_{p(l')}\rangle$, with the error of Lemma~\ref{bil:mat}(b). Finally $O^T_K(H)=a_K^{-1}\sum_pw_H(p)a_p+Y[\mathsf w_H]$ with the plaquette weight $w_H(p)=H_{\lambda(p)\lambda(p)}(x(p))$ and the bond weight $\mathsf w_H(\mathsf e)=H_{\lambda\lambda'}(x_{\mathsf e})$: a plaquette form of size $\le|H|_\infty$ and a bond form of size $\le59|H|_\infty$, whose sizes add to $60|H|_\infty$. (b) The source enters these proofs at four places.

First, through the clause on $F$ itself: its continuity and its gauge invariance in the first step, where the Wilson density is multiplied by $e^{F_\lambda}$, $F_\lambda=\sum_a\lambda_aF_a$ (the first-step part of Lemma~\ref{lem:weighted-array-step}), and in the sum over the level of the first creation (equation~\eqref{eq:ward-first-creation}); its pointwise majorant, which bounds $F$ at the backgrounds of the small fields of the level $t$ by $L^K\mathsf s\sum_{p\in S}\sup a_p$, the supremum over these backgrounds, that is by the size of the source at the level $t$, and by $2L^K\mathsf s\,\#S$ at the height $0$, in the quotients, their tower, the reduction and the identity of the quotients (Proposition~\YMMref{123.1}{829} and equation~\YMMref{(2268)}{829}, the quotient identities preceding equation~\YMMref{(2273)}{830}, equation~\YMMref{(2273)}{830} and Proposition~\YMMref{123.2}{830}); it is the source of the power $2\mathfrak p_0$ of Lemma~\ref{bil:nsrc}.

Second, at the height $0$, through its family in the fourth place over a point and the size of that family (the height-zero part of Proposition~\ref{prop:weighted-finite-system}; here the clause of the height $0$ of (a)). Third, at the heights $1\le j\le\bar k$, through its families without base point, their sizes and the action formula (Lemma~\ref{lem:plaquette-weight-member}), the one-step covariance of the level $j+1$ being formed with the member of the height $j$ (Lemma~\ref{lem:weighted-step-covariance}); here the clause of the heights of (a), up to the top index. Fourth, at the second order at the small points of a chart, through its kernel (Lemma~\YMMref{124.4}{835}), in the carried family and the flat top (Propositions~\YMMref{125.1}{838} and~\YMMref{123.3}{832}, equation~\YMMref{(2283)}{832}); here the kernel $\mathsf t$ of (a).

Every other step of these proofs is stated there for arbitrary bare arrays and local families, for bounded Borel functions of the field, or for numbers. (c) The step map is formed from its input array alone; its Lipschitz constant for the arrays of plaquette weights, with the count of (a) and Lemma~\ref{bil:mat}(b), gives the factor $8n_{\mathfrak c}e^{4t_{\mathfrak c}/13}$.
\end{proof}

\begin{lemma}[The kernel]\label{bil:ker}
Let $s'\le K$, $d\ge a_{s'}$, $\hat d=d/a_{s'}$, and $|u|=|x(p)-x(q)|\ge d$. Then
\[
|C^{\rm bg}_{K,s'}(p,q)-\epsilon_p\epsilon_qa_K^3K_3(u)_{\lambda(p)\lambda(q)}|\le(a_K/d)^3E_1,\qquad E_1=C_{\rm lat}[(a_K/d)^2+(d/\ell)^3]+C^{\rm fl}_\star\hat d^3e^{-\delta_{\rm tp}L\hat d/\sqrt3},
\]
$\epsilon_p=\pm1$ the orientation of $p$; so $|C^{\rm bg}(p,q)|\le C_{\mathsf k}(a_K/d)^3$, $C_{\mathsf k}=1+3C_{\rm lat}+C^{\rm fl}_\star C^{\rm tm}_\flat$, and in the form in which far pairs use it the same holds with $C'_{\mathsf k}=3^{3/2}C_{\rm tp}\mathsf S_3+C^{\rm bg}_{\rm tp}$. If $8a_K\le d$ and $4d<\ell$, a ball of radius $d$ holds at most $10(d/a_K)^3$ sites, at most as many centres of bonds of one direction, and as many centres of plaquettes of one normal.
\end{lemma}

\begin{proof}[Proof sketch]
The free kernel is the background kernel plus the fluctuation kernel $C^{\rm fl}_{K,s'}=\Pi_0-\Pi_{K-s'}$, the sum of the slices of the scales below $a_{s'}$ \YMMref{(2227)}{818}. The comparison of the lattice free kernel with $a_K^3K_3$, with the number $C_{\rm lat}$ of equation~\YMMref{(2245)}{824}, gives the first term of $E_1$, and the exponential decay of the fluctuation kernel beyond the scale $a_{s'}$, $|C^{\rm fl}_{K,s'}(p,q)|\le C^{\rm fl}_\star(a_K/a_{s'})^3e^{-\delta_{\rm tp}L\mathsf r/a_{s'}}$ for $\mathsf r=|x(p)-x(q)|_\infty\ge a_{s'}/\sqrt3$ (Proposition~\YMMref{120.6}{821} at $t_0=3^{-1/2}$), with $\mathsf r\ge d/\sqrt3$, gives the second; condition (H) fixes the constants, and $|K_3(u)_{\lambda\lambda'}|\le(2\pi)^{-1}|u|^{-3}$. The count is a count of lattice points in a ball.
\end{proof}

\begin{lemma}[$n$ sources on a small torus]\label{bil:nsrc}
Let $0<g\le\gamma_2$, $n\ge2$, let $\ell$ be small, and let $s_{\rm th}(n)\le s(d)\le K$, where $s_{\rm th}(n)$ is the least integer $s\ge3$ with $2T_n(3)\theta(L^{s-2}g^{-2})\le1$, $\theta(t)=K_M(1+\beta_s)\delta(t)/\alpha_0$, which does not increase on $[e^P,\infty)$, $\delta(t)=A_1t^{-1/2}(\log t)^{\mathfrak p_0}$, and $T_n$ is the series \eqref{eq:source-mixed-series}. Let $(\mathsf s_a,S_a)$ be majorants with points $x_a$, $|x(p)-x_a|\le\frac78d$ on $S_a$, $|x(p)-x(q)|\ge\frac54d$ for $p\in S_a$, $q\in S_{a'}$, $a\ne a'$, and let $F_a$ be (complex) sources of top index $\mathsf k$, the working level, under them. Then
\[
|\langle F_1;\dots;F_n\rangle_\ell|\le k'_n(g,s(d))\,b^{-n}\prod_a\mathsf s_a,\qquad k'_n(g,s)\le C_n(g)Y_s^{8n\mathfrak p_0},
\]
so $k'_n(g,s(d))\le C_n(g)[2\log\frac1d+C(g)]^{8n\mathfrak p_0}$. No upper bound on the distances is asked.
\end{lemma}

\begin{proof}[Proof sketch]
Insert $F_\lambda=\sum_a\lambda_aF_a$ as the source into the weighted progression of \eqref{eq:weighted-progression-recursion} (Proposition~\textup{\YMMref{122.1}{826}}). A chain rule for the step map in any number of parameters (Cauchy estimates on polydiscs) and locality give, for every set $\beta$ of $r=\#\beta\ge2$ sources: the mixed derivative $E^{(j)}_\beta=\partial_\beta E^{(j)}_\lambda|_0$ of the local family of level $j$ vanishes on the polymers that miss some $S_a$, $a\in\beta$, so that its norm is at most $e^{-4\beta_s\kappa[\mathsf D_j/M-2]_+}$ times its size $\mathsf d_{j,\beta}$; splitting into the top block and the rest, by induction on $r$ and on $j$, $\mathsf d_{j+1,\beta}\le c_\theta\hat G_r\prod_{a\in\beta}\hat g_{j,a}$, $\hat G_r$ given by a recursion in $r$ and $\hat g_{j,a}$ the sizes of the sources at the level $j$, as long as $\theta^{[j]}T_n(\mathsf D_j/M)\le\frac12$; the threshold $s_{\rm th}(n)$ gives this for $1\le j\le\mathsf k$, since $\mathsf D_j/M\ge3$ and $x_j>L^{s-2}g^{-2}$ there.

The chain runs up to the level $\mathsf k+1$, and its top index is the working level $\mathsf k$, so that the sources need their families without base point at the heights $1\le j\le\mathsf k$ alone. On the history without large fields, with $e^\Psi$ the density of the level $t\le\mathsf k$, $\Psi$ formed with the source $F_\lambda$, the identity $\partial_Be^\Psi=e^\Psi\sum_\pi\prod_{\beta\in\pi}\partial_\beta\Psi$ gives, almost everywhere and for any complex constants $c_a$,
\[
\Phi^{(t)}\Bigl[\prod_{a\in B}(F_a-c_a)\Bigr]=\sum_{\pi\in{\rm Part}(B)}\prod_{\beta\in\pi}\varkappa^{(t)}_\beta,\qquad\varkappa^{(t)}_{\{a\}}=\partial_a\Psi|_0-c_a,\qquad\varkappa^{(t)}_\beta=\partial_\beta\Psi|_0\quad(\#\beta\ge2),
\]
since the constants multiply the density by $e^{-\sum_a\lambda_ac_a}$; the blocks of two or more elements do not depend on the $c_a$, and all the blocks are bounded by the sizes \YMMref{(759)}{375}. The sum over the level $t$ of the first creation (equation~\eqref{eq:ward-first-creation}) adds the weights $\mathsf w_t\le2M^{-3}\ell^3L^{3(K-t)}e^{-\frac12c'p_0(g_t)^2}$ (Lemma~\YMMref{125.3}{841}), and on a small side $2M^{-3}\ell^3\le M^{-3}c'A_0^2Y_s^P$; the Gaussian decay in $K-t$ makes the sum converge. M\"obius inversion gives the cumulant from the moments of the $F_a-c_a$. These steps alone give a power of $\frac1d$: the sizes of the blocks hold the factor $x_{\mathsf k}^{r/2}\le\bar x_s^{r/2}$, and the sum over $t$ the factor $L^{3s}$.

Three estimates remove both factors. (i) $\hat d\ge4MY_s^P$ (Lemma~\YMMref{126.1}{844}, equation~\YMMref{(2323)}{844}), and corners of plaquettes of different sets $S_a$ are at the sup distance $\ge\frac5{4\sqrt3}d-a_K$; so $\mathsf D_{\mathsf k}/M\ge2Y_s^P$, as in Lemma~\YMMref{126.1}{844}, and a block of an order $r\ge2$ at the working level is at most a number times $\hat d^3(x_{\mathsf k}^{1/2}y^{\mathfrak p_0})^re^{-(\mathsf D_{\mathsf k}/M-2)}\prod_{a\in\beta}\upsilon\mathsf s_a$, in which $e^{2-2Y_s^P}\bar x_s^{r/2}=e^{2-2Y_s^P+rY_s/2}$ is bounded for $r\le n$ (at the levels $t<\mathsf k$ the blocks are larger by $L^{(3+r/2)(\mathsf k-t)}(\mathsf k-t+2)^{r\mathfrak p_0}$, and they enter the sum over $t$ alone). (ii) In the sum over $t$, $y_{\mathsf k-m}\ge y+m\log\frac{3L}4$ for $m\ge0$, $(y+u)^P\ge y^P+u^P$, and $u^P\ge u^2$ at $u=m\log\frac{3L}4$ give $e^{-\frac12c'p_0(g_{\mathsf k-m})^2}\le e^{-\frac12c'A_0^2y^P}e^{-\frac12c'A_0^2(m\log\frac{3L}4)^2}$; the first factor is kept, and as $y>Y^-_s=Y_s-\log(2L^3)$ (Lemma~\YMMref{126.1}{844}, equation~\YMMref{(2324)}{844}, equations~\YMMref{(2328)--(2336)}{846}, with the notation in \eqref{eq:short-distance-envelope-names}) it is at most $e^{-\frac12c'A_0^2(Y^-_s)^P}$, against which $L^{3s}\bar x_s^r$ is bounded. (iii) Take for $c_a$ the constant part of the first order of $F_a$ at the working level, $c_a=\sum_{i\le\mathsf k}\sum_{(X,z)}\dot E^{(i)}_a(X,\mathbf 1,z)$, $\dot E^{(i)}_a$ the derivative in $\lambda_a$ at $0$ of the local family of the level $i$.

Then the singleton $\varkappa^{(\mathsf k)}_{\{a\}}$ is the first-order part of $F_a$ at the block field of that level, at most its supremum over the small fields, the size $\mathsf X_{\mathsf k,a}\le\mathsf c_Xy^P\hat d^3\upsilon\mathsf s_a$ of the size definitions preceding and in Proposition~\YMMref{123.2}{830} (Lemma~\YMMref{125.2}{840}, equation~\YMMref{(2310)}{840}, at each pair of the majorants), plus an irrelevant part, smaller by a factor that tends to $0$ as $s\to\infty$; neither holds a factor $x_{\mathsf k}$. So the leading term of a moment is the product of the singletons, at most $\prod_aC\,Y_s^P(LC_W(s))^3\upsilon\mathsf s_a$; the partitions with a block of an order $\ge2$ and the sum over $t$ add at most a number of $n$ and $g$ times this, and with $\upsilon^2\le\frac98b^{-2}$ M\"obius inversion gives $k'_n(g,s)\le C_n(g)Y_s^{8n\mathfrak p_0}$.

Of the power $8\mathfrak p_0$ for each entry, $2\mathfrak p_0$ come from the supremum of $x_{\mathsf k}a_p$ over the small fields, of order $A_0^2y^{2\mathfrak p_0}$, and $6\mathfrak p_0$ from the number $\hat d^3<(LC_W(s))^3$ of the unit cubes of the working level in a ball of radius $d$. Of the side the proof reads the settings of the chain at the $n$ majorants, the conditions of the sum over the level of the first creation and the bound on the weights $\mathsf w_t$, which serves only to bound that sum. The weights enter only the tail term of each block of the M\"obius inversion, the sum over the level of the first creation in the moment of that block: linearly in their prefactor, and through the Gaussian sum of the estimate (ii) above, in which $e^{-c''p_0(g_t)^2}$ with $c''\ge\frac12c'$ may stand for $e^{-\frac12c'p_0(g_t)^2}$.

As the settings and the conditions hold on every admissible side for $K\ge3$ (Step~14), the lemma holds on every admissible side, for $K\ge3$ and weights $\mathsf w_t\le\mathsf LL^{3(K-t)}e^{-c''p_0(g_t)^2}$, $1\le t\le K$, $c''\ge\frac12c'$, with the factor $\max\{1,M^3\mathsf L/(c'A_0^2)\}^n$ on $k'_n$: as $Y_s\ge1$, each block gains at most the factor $\max\{1,M^3\mathsf L/(c'A_0^2)\}$ over the tail of a small side, and there are at most $n$ blocks (Theorem~\ref{thm:bil}(E3)).
\end{proof}

\begin{lemma}[As many sources as unit cubes]\label{bil:prod}
Let $0<g'\le\gamma_2$, $(K',\beta)$ a first-passage datum at $g'$, $\ell$ a side small over its running couplings, $0<d<1$ with $s=s(d)\le K'$, the depth, the working level and $\upsilon$ being formed at $g'$, and let $n$ be an integer with $2\le n\le\ell^3$. Let $(\mathsf s_a,S_a)$, $1\le a\le n$, be majorants placed as in Lemma~\ref{bil:nsrc} and $F_a$ sources of top index $\mathsf k$, the working level, under them, and suppose
\[
2Y_s^P\ge u_*(n),\qquad y_{K'}\ge2\log L+\mathfrak p_0\log3,\qquad L^{3s}\le x_{K'}^{20},
\]
where $u_*(n)=\max\{4,(2/A_\kappa)(n[c_1+\log(c^*n)]+\log3)\}\le Cn\log(en)$, $A_\kappa=4\beta_s\kappa$, $c^*=\max\{4,\frac12c^\flat\}$, $c_1$ a number of the constants. Then there are complex numbers $c_a$ such that for every nonempty $B\subset\{1,\dots,n\}$
\[
\Bigl|\Bigl\langle\prod_{a\in B}(F_a-c_a)\Bigr\rangle_{K',\ell,\beta}\Bigr|\le\prod_{a\in B}\mathsf A_n(s)\upsilon\mathsf s_a,\qquad\mathsf A_n(s)\le Cn^{4+\mathfrak p_0}(LC_W(s))^3\bar x_sY_s^P .
\]
\end{lemma}

\begin{proof}[Proof sketch]
It is the proof of Lemma~\ref{bil:nsrc} with two changes. First, the mixed derivative families of every order $r\le n$ are bounded with the numbers $(c^*r)^r$ in the place of $\hat G_r$, by Cauchy estimates on polydiscs of radii proportional to $1/r$; for this it suffices that $T_m(u)(c^*m)^m\le1$ for $2\le m\le n$ at $u=\mathsf D_j/M$, which holds as soon as $u\ge4$ and $\frac12A_\kappa u\ge m[c_1+\log(c^*m)]+\log3$, so for $u\ge u_*(n)$, and the first hypothesis gives $\mathsf D_j/M\ge u_*(n)$. Second, the moments are bounded directly, from the conditional cumulants at the upper level and the sum over the partitions, and not through the cumulants of the sources; the sum over the level of the first creation (equation~\eqref{eq:ward-first-creation}) is bounded by the count of the cubes of a small side, through the last two hypotheses. Hence a bound in product form, with no factorial in $n$ beyond the power in $\mathsf A_n(s)$.
\end{proof}

\begin{lemma}[Two sources at one scale]\label{bil:two}
(a) Let $g\le\gamma_2$, $\ell$ small, $s_1(g)\le s(d)\le K$, $3d\le|x_1-x_2|\le4d$, and let $(F_a,\mathsf t_a)$ be real quadratic sources of top index $T-1$ under $(\mathsf s_a,S_a)$ in the sense of Lemma~\ref{bil:forms}(a), with their families without base point at the heights $1\le j\le T-1$, $T$ the upper level (forms have a member at every height $j\ge1$), $S_a$ within $d-a_K$ of $x_a$. Then
\[
\begin{multlined}
\Bigl|\langle F_1;F_2\rangle_\ell-6\upsilon^2\sum\mathsf t_1(p,p')\mathsf t_2(q,q')C^{\rm bg}(p,q)C^{\rm bg}(p',q')\Bigr|\le\varepsilon^\natural_g(s(d))\upsilon^2\mathsf s_1\mathsf s_2,\qquad|\langle F_1;F_2\rangle_\ell| \\
\le C^\sharp_g(s(d))\upsilon^2\mathsf s_1\mathsf s_2,
\end{multlined}
\]
where $\varepsilon^\natural_g(s)=\varepsilon^\flat_g(s)+2(1+A_1)^4\mathsf C_5s^2(LC_W(s))^6Y_s^{2P}{\rm sm}_g(s)$,
\[
\begin{aligned}
\varepsilon^\flat_g(s)={}&(\mathsf C_2+\mathsf C_3)(Y^-_s)^{-2\mathfrak p_0}+\mathsf C_3\,s\,(Y^-_s)^{-5\mathfrak p_0}\\
&+\mathsf C_6(Y^-_s)^{-6\mathfrak p_0}+\mathsf C_8\bigl((LC_W(s))^3+1\bigr)(Y^-_s)^{-9\mathfrak p_0}\\
&+4\mathsf c_X^2(Y^-_s)^{-14\mathfrak p_0}\\
&+(\mathsf C'_B+\mathsf C_7+\mathsf C'_3+\mathsf C''_3+\mathsf C_6)\,s^4(LC_W(s))^6Y_s^{11\mathfrak p_0}{\rm sm}_g(s),
\end{aligned}
\]
$Y^-_s=Y_s-\log(2L^3)$ and ${\rm sm}_g$ the functions of equations~\YMMref{(2328)--(2336)}{846}, with the notation in \eqref{eq:short-distance-envelope-names}, $\mathsf C_i$ the constants of Lemma~\YMMref{125.5}{842}, with the convention in \eqref{eq:short-distance-constants}--\eqref{eq:short-distance-collected-constants}: $\varepsilon^\flat_g$ bounds all the errors of the proof but those of the Gaussian values at the flat top, which the second summand bounds, and both tend to $0$ as $s\to\infty$ (the decay argument preceding equation~\YMMref{(2332)}{846}); and 
\[
C^\sharp_g(s)=\mathsf C_2(LC_W(s))^6Y_s^{2P}+4\mathsf c_X^2+\mathsf C'_Bs^4(LC_W(s))^6Y_s^{6\mathfrak p_0}{\rm sm}_g(s)\le C(g)Y_s^{16\mathfrak p_0}.
\]


(b) Let $\ell$ be admissible, $K\ge3$, $s=s(d)\le K$, $\mathsf{Lm}_s\le L^{s-3}$, $x^-_s\ge e^P$; then the settings of the chain hold at the levels $\mathsf k<j_1\le T\le K-2$, with $n_T\le s$, $676a_K\le d$, $20d<\ell$, $M\le\hat d<LC_W(s)$. Let $|x_1-x_2|\ge3d$, $|L^{-1}x_{t-1}-x_t|\le B_\sharp$ for all $t$, $B_\sharp$ the flow bound of the running couplings and $\bar B=LB_\sharp/(L-1)$ (Section~\YMMref{74}{614}), $y_K\ge2\log(1+\bar B)$; let $\mathsf L\ge0$, $c>0$ with $\mathsf w_t\le\mathsf LL^{3(K-t)}e^{-cp_0(g_t)^2}$ ($t\le T$), $\mathsf W_T\le\mathsf LL^{3(K-T)}e^{-cp_0(g_T)^2}$, $2c\mathfrak p_0A_0^2y^{P-1}\ge13$, and ${\rm sm}^{\mathsf L,c}_T\le1$. Then on $T_\ell$ the error in (a) is at most
\[
\bigl\{(\mathsf C_2+\mathsf C_3)y^{-2\mathfrak p_0}+\mathsf C_3n_Ty^{-5\mathfrak p_0}+\mathsf C_6y^{-6\mathfrak p_0}+\mathsf C_8(\hat d^3+1)y^{-9\mathfrak p_0}+4\mathsf c_X^2y^{-14\mathfrak p_0}+\mathsf C^e_Bn_T^4\hat d^6y^{15\mathfrak p_0}{\rm sm}^{\mathsf L,c}_T\bigr\}\upsilon^2\mathsf s_1\mathsf s_2,
\]
and the remainder ${\rm Rest}'_T$ of the proof of (a) obeys $|{\rm Rest}'_T|\le\bigl[4\mathsf c_X^2y^{-14\mathfrak p_0}+\mathsf C'_Bn_T^4\hat d^6y^{6\mathfrak p_0}{\rm sm}^{\mathsf L,c}_T\bigr]\upsilon^2\mathsf s_1\mathsf s_2$; for two plaquette weights $\mathsf C^{(1)}_B$ may replace $\mathsf C'_B$.

(c) In the setting of (a), almost everywhere at the upper level $T$ and for all the quotients,
\[
\bigl|\Phi^{(T)}[F_1F_2]-\Phi^{(T)}[F_1]\,\Phi^{(T)}[F_2]-G\bigr|\le\varepsilon_g(s(d))\,\upsilon^2\mathsf s_1\mathsf s_2,
\]
$G=6\upsilon^2\sum\mathsf t_1(p,p')\mathsf t_2(q,q')C^{\rm bg}(p,q)C^{\rm bg}(p',q')$ the Gaussian value of (a) and $\varepsilon_g=\epsilon_g$ the error function in equation~\YMMref{(2330)}{846}, with the coefficient fixed in Lemma~\ref{lem:short-distance-collected-constant}; and almost everywhere at every level $j\le T$ the conditional covariance $\Phi^{(j)}[F_1F_2]-\Phi^{(j)}[F_1]\Phi^{(j)}[F_2]$ is at most the crude size $\mathsf M'_j$ of the size definitions preceding and in Proposition~\YMMref{123.2}{830}, formed at the majorants, which Lemma~\YMMref{125.2}{840} bounds by $[\mathsf m_\star L^{4n}(n+1)^{2\mathfrak p_0}+2n_T\mathsf{sb}_0^2]\upsilon^2\mathsf s_1\mathsf s_2$, $n=[\mathsf k-j]_+$.
\end{lemma}

\begin{proof}[Proof sketch]
(a) The covariance is a sum of one-step covariances plus a rest, $\langle F_1;F_2\rangle=\sum_{\mathsf k<j\le T}{\rm qm}_j+{\rm Rest}'_T$, ${\rm qm}_j=\int{\rm Qt}_j[X_{j-1,1};X_{j-1,2}]\,d\mathsf m_{j|T}$, with the step quotients ${\rm Qt}_j$ of the quotient definitions in Section~\ref{sec:weighted-covariance-interface}, their measures $\mathsf m_{j|T}$ and the first-order parts $X_{j-1,a}$ of the sources (Proposition~\YMMref{123.2}{830}, which holds for these sources by Lemma~\ref{bil:forms}(b)). The steps $j_1<j\le T$ are bounded crudely, by $2\mathsf X_{j-1,1}\mathsf X_{j-1,2}$ each: the error $\mathsf{Er}_2$ of equation~\YMMref{(2307)}{839}. For $\mathsf k<j\le j_1$ the one-step covariance is almost everywhere the total of the mixed derivative family $\mathsf K^{(j)}$ of the step $j$ (Lemma~\ref{lem:weighted-step-covariance}); carried to the upper level $T$ (Proposition~\YMMref{123.3}{832}) and evaluated at the trivial background it is the number $\mathsf{sum}^{(j)}_T$, up to the error $\mathsf{Ft}^{(j)}_T$ of the flat top, which holds at every point of the small-field set of the level $T$ (Proposition~\YMMref{123.3}{832}, equation~\YMMref{(2283)}{832}, in its sup-norm form); the mass $1-\mathsf W_T$ of $\mathsf m_{T|T}$ costs $\mathsf W_T(\bar{\mathsf n}^{\rm G}_j+\mathsf e_{j,T})$; and $\mathsf{sum}^{(j)}_T$ is the Gaussian number $\mathsf n^{\rm G}_{j,T}$ up to the error $\mathsf e_{j,T}$ of equations~\YMMref{(2303)--(2305)}{838} (Proposition~\YMMref{125.1}{838}).

With kernels $\mathsf t_a$ in the place of $\frac12w_a\delta$ this proof reads the source at three places alone: the covariance of one step at a triple is the covariance of the first-order variations of the two sources under the fluctuation measure of that step; the second order of a source at a small point of the chart is the quadratic form of its kernel; and the forms of the slices are built with the kernels. So it gives the Gaussian numbers $2(N^2-1)\upsilon_j^2\sum\mathsf t_1(p,p')\mathsf t_2(q,q')\mathcal C_{(j)}(p,q)\mathcal C_{(j)}(p',q')+4(N^2-1)\sum_{j<i'\le T}\upsilon_j\upsilon_{i'}\sum\mathsf t_1(p,p')\mathsf t_2(q,q')\mathcal C_{(j)}(p',q')\mathcal C_{(i')}(p,q)$, with the same errors $\mathsf e_{j,T}$ and the same constants, for every real quadratic source, the bond forms under $(59\sup|\mathsf w|,S[\mathsf w])$ among them (Lemma~\ref{bil:forms}(a)).

The Gaussian numbers sum to $6\upsilon^2\sum\mathsf t_1\mathsf t_2DD$, $D=\sum_{\mathsf k<j\le T}\mathcal C_{(j)}=\Pi_{\mathsf k}-\Pi_T$, since the slices telescope and the terms $(j,i')$, $(i',j)$ agree under the symmetry of the kernels, and $C^{\rm bg}=\Pi_{\mathsf k}$; the Gaussian numbers above $j_1$, the replacing of the couplings $\upsilon_j$ of the steps by $\upsilon$ and the dropping of $\Pi_T$, with $|\Pi_T(p,q)|\le C^{\rm bg}_{\rm tp}\eta_T^3$ (Proposition~\YMMref{120.5}{821}), cost the errors $\mathsf{Er}_6$, $\mathsf{Er}_7$, $\mathsf{Er}_8$ of equation~\YMMref{(2307)}{839}. With the deposit inequality $L^{3a}\mathsf e_{\mathsf k-a}\le L^{-10a}\mathsf e_{\mathsf k}$, $\mathsf e_k=e^{-\frac12c'p_0(g_k)^2}$ (the proof of Lemma~\YMMref{125.4}{841}; its hypotheses
hold by equation~\YMMref{(2326)}{845} and the logarithmic
floors immediately following that equation),
\[
|{\rm Rest}'_T|\le\bigl[4\mathsf c_X^2y^{-14\mathfrak p_0}+\mathsf C'_Bn_T^4\hat d^6y^{6\mathfrak p_0}{\rm sm}_T\bigr]\upsilon^2\mathsf s_1\mathsf s_2 .
\]
Collected, the left member of (a) is at most $|{\rm Rest}'_T|+\mathsf{Er}_2+\sum_{\mathsf k<j\le j_1}\mathsf{Ft}^{(j)}_T+\mathsf W_T\sum_{\mathsf k<j\le j_1}\bar{\mathsf n}^{\rm G}_j+2\sum_{\mathsf k<j\le j_1}\mathsf e_{j,T}+\mathsf{Er}_6+\mathsf{Er}_7+\mathsf{Er}_8$. By the bounds of the errors at the attached levels (Lemma~\YMMref{125.5}{842}), with $Y^-_s<y\le Y_s$ (Lemma~\YMMref{126.1}{844}, equation~\YMMref{(2324)}{844}), $1\le n_T\le s$, $\hat d<LC_W(s)$ and ${\rm sm}_T\le{\rm sm}_g(s)$, all these terms but the sum of the $\mathsf e_{j,T}$ are together at most $\varepsilon^\flat_g(s)\upsilon^2\mathsf s_1\mathsf s_2$, and $2\sum_{\mathsf k<j\le j_1}\mathsf e_{j,T}\le2(1+A_1)^4\mathsf C_5s^2(LC_W(s))^6Y_s^{2P}{\rm sm}_g(s)\upsilon^2\mathsf s_1\mathsf s_2$.

The crude bound is the same sum with the Gaussian numbers bounded by their sizes. (b) In the proof of (a) the side enters only through the weights $\mathsf w_t$ and the sum over the level of the first creation (equation~\eqref{eq:ward-first-creation}), and the distance only through a lower bound; the factor $e^{-cp_0(g_t)^2}$ is used at three places alone: in the deposit inequality at $c$, in its monotonicity in $t$,
and in the bound on $\mathsf w_t$. For the first, the proof of
Lemma~\YMMref{125.4}{841} applies with $c'/2$ replaced by $c$.
To make this substitution explicit, put $P=2\mathfrak p_0$,
$y=y_{\mathsf k}$ and
$\Delta=y_{\mathsf k-a}-y$ for $1\le a\le\mathsf k$.
The flow bound and $y_K\ge2\log(1+\bar B)$ give
\[
 \Delta\ge a\log L+\log(1-\bar B/x_{\mathsf k})>2a,
 \qquad
 y(a\log L-\Delta)
 \le\frac8{3e}\frac{\bar B}{1+\bar B}<1.
\]
Keeping the linear and quadratic terms of the power gives
\[
 (y+\Delta)^P-y^P
 \ge Py^{P-1}a\log L
   +Py^{P-2}\left[\frac{P-1}{2}\Delta^2
                         -y(a\log L-\Delta)\right]
 \ge Py^{P-1}a\log L.
\]
Thus, for $\mathsf e_j=e^{-cp_0(g_j)^2}$,
\[
 L^{3a}\mathsf e_{\mathsf k-a}
 \le L^{-(2c\mathfrak p_0A_0^2y^{P-1}-3)a}
               \mathsf e_{\mathsf k}
 \le L^{-10a}\mathsf e_{\mathsf k}.
\]
The case $a=0$ is immediate. This uses exactly the rate condition
$2c\mathfrak p_0A_0^2y^{P-1}\ge13$ in (b), at the working
level $\mathsf k$; no stronger rate condition at $K$ is needed. The settings follow from the arithmetic of the working level under the two depth conditions, the flow bound and the lower bound on $y_K$. (c) A block of two is the conditional covariance $\Phi^{(n)}_{12}=\Phi^{(n)}[F_1F_2]-\Phi^{(n)}[F_1]\Phi^{(n)}[F_2]$.

Almost everywhere, by the tower of the quotients (equation~\YMMref{(2273)}{830}), $\Phi^{(n)}_{12}={\rm Qt}_n[\Phi^{(n-1)}_{12}]+{\rm Qt}_n[\Phi^{(n-1)}[F_1];\Phi^{(n-1)}[F_2]]$, and $\Phi^{(n-1)}[F_a]$ is a constant plus the first-order part $X_{n-1,a}$ plus an irrelevant part (equation~\YMMref{(2267)}{829} and Proposition~\YMMref{123.1}{829}); this gives the crude bound at every level (the proof of Proposition~\YMMref{123.2}{830}, through Lemma~\ref{bil:forms}(b)). For $\mathsf k<j\le j_1$ the one-step covariance ${\rm Qt}_j[X_{j-1,1};X_{j-1,2}]$ is almost everywhere the total $\mathsf F^{(j)}_j$ of the family $\mathsf K^{(j)}$ (Lemma~\ref{lem:weighted-step-covariance}), and ${\rm Qt}_{i+1}[\mathsf F^{(j)}_i]=\mathsf F^{(j)}_{i+1}$ almost everywhere (the proof of Proposition~\YMMref{123.3}{832}).

So, by induction on $n$ from $\mathsf k$ to $T$, the remainder $R_n=\Phi^{(n)}_{12}-\sum_{\mathsf k<j\le\min\{n,j_1\}}\mathsf F^{(j)}_n$ is almost everywhere at most $\mathsf M_{\mathsf k}$, plus twice the products of the sizes of the first-order and of the irrelevant parts at the levels $\mathsf k\le j<n$, plus twice the products of the sizes of the first-order parts at the steps above $j_1$. At $n=T$ the measure $\mathsf m_T$ is carried by the small-field set of the level $T$, at every point of which $\mathsf F^{(j)}_T$ is within $\mathsf{Ft}^{(j)}_T$ of the number $\mathsf{sum}^{(j)}_T$ (Proposition~\YMMref{123.3}{832}, equation~\YMMref{(2283)}{832}, before any integration); $\mathsf{sum}^{(j)}_T$ is $\mathsf n^{\rm G}_{j,T}$ up to $\mathsf e_{j,T}$, and the Gaussian numbers give $G$ up to $\mathsf{Er}_6+\mathsf{Er}_7+\mathsf{Er}_8$, as in (a).

No sum over the levels of the first creation enters, so that there are no tails and no weight $\mathsf W_T$: the error is at most the bracket of Corollary~\YMMref{125.6}{843} times $\upsilon^2\mathsf s_1\mathsf s_2$, and with $Y^-_s<y\le Y_s$, $n_T\le s$, $\hat d<LC_W(s)$ and ${\rm sm}_T\le{\rm sm}_g(s)$ this bracket is at most $\varepsilon_g(s)$, term by term.
\end{proof}

\begin{lemma}[Symmetries and copies]\label{bil:sym}
(a) For a lattice isometry $r$ with $re_\mu=s_\mu e_{\pi(\mu)}$: $q_{(r\mathsf e)^+}(rU)=s_\lambda s_{\lambda'}q_{\mathsf e}(U)$, $\lambda<\lambda'$ the directions other than that of $\mathsf e$; so $rO^T_K(H)=O^T_K(rH)$, and there is a linear bijection $r^\sharp$ of the weights, with $\|r^\sharp v\|=\|v\|$, real weights kept real and carriers and sets $S[\cdot]$ carried by $r$, such that the image of a form is the form of the same kind with the weight $r^\sharp v$; in particular $Y[r^\sharp\mathsf w](rU)=Y[\mathsf w](U)$ and $\theta Y[\mathsf w]=Y[\theta^\sharp\mathsf w]$.

(b) A law that is invariant under the reflections in the site planes has $\nu(q_{\mathsf e})=0$; all laws are such, so $\nu(Y[\mathsf w])=0$.

(c) For a board of side $\mathsf l_0$ and $\ell_0=2\mathsf l_0$, a form whose carrier lies in a cell at a distance $\ge\rho>a_K$ from the walls is a function of the cell; its copies on $T_{\ell_0}$ are forms of the same kind, size and mean, with the mirror images as carriers; copies in different cells of a period box are $\ge2\rho$ apart, two in one cell at the original distance.

(d) If $3d\le|x_0-y_0|\le4d$, the midpoint is within $\frac{a_K}2$ of the offset of the board of side $\mathsf l$ in each coordinate, and the weights are carried within $d-2a_K$ of $x_0$, $y_0$, then the corners of their plaquettes lie within the sup distance $3d$ of the centre of the cell, and the sixteen copies on $T_{\ell'}$ form eight pairs, partners $3d$ to $4d$ apart, different pairs at sup distance $\ge\mathsf l-6d\ge11M$, the suprema of the weights kept; the same holds for tensor weights and arbitrary majorants.

(e) For majorants with plaquettes within $d-a_K$ of $x_0$, $y_0$, $3d\le|x_0-y_0|\le4d$, the tents $\tilde f_a=\mathsf s_a\sum_{p\in S_a}\max\{0,1-2a_K^{-1}|u-x(p)|_\infty\}$ are Lipschitz, equal $\mathsf s_a\mathbf1_{S_a}$ at the barycentres, have their supports in $B(x_a,d)$, and $d\le{\rm dist}(\mathrm{supp}\,\tilde f_1,\mathrm{supp}\,\tilde f_2)\le6d$. Let besides $0<g\le\gamma_2$, $(K,\beta)$ a first-passage datum at $g$, $\ell$ small, $0<d<1$ and $s_1(g)\le s=s(d)\le K$, so that $\mathsf{Lm}_s\le L^{s-3}$ and $x^-_s\ge e^P$ (equation~\YMMref{(2331)}{846}). Then $s\ge3$, $2\le\mathsf k\le K-1$, $M\le L^{-2}C_W(s)\le\hat d<LC_W(s)$, $\mathsf D_{\mathsf k}\ge3M$ for the pair, $20d\le\ell$, and the middle and the upper level have $\mathsf k<j_1\le T\le K-2$, $1\le n_T\le s$, $L^{n_T}<\mathsf{Lm}_s$, $x_T>x^-_s$.

Let $n\ge1$ pairs of such majorants be given, $p=2n$, the corners of plaquettes of different pairs at the sup distance $\ge11M$. Then $\mathsf D_{\mathsf k}$, the least over all $a\ne a'$, is $\ge3M$: partners give $3M$, members of different pairs $11ML^{K-\mathsf k}$. The settings of the chain for the $p$ sources hold with the top index $T-1$, with the smallness conditions at the heights $0,\dots,T-1$ and $d_*(t)\ge3$ for $1\le t\le T$; the levels, the lattices, the histories, the measures and the quotients of the chain are the same for all pairs, and for the chain of each pair alone, and do not depend on the sources; and $\theta^{[j]}T_p(\mathsf D_j/M)\le\frac12$ for $1\le j<T$ as soon as
\[
{\rm th}_{g,p}(s)=2s^{2+p\mathfrak p_0}\mathsf{Lm}_s^{p/2-1}T^+_p\,\theta(x^-_s)\le1,\qquad T^+_p=T_p(3)+(\tfrac43)^{p/2}L,
\]
$\theta$ the function of Lemma~\ref{bil:nsrc}; ${\rm th}_{g,p}(s)$ tends to $0$ as $s\to\infty$, being a power of $s$ times $L^{-s/2}$.
\end{lemma}

\begin{proof}[Proof sketch]
(a) $r$ carries the loops at $\mathsf e$ to those at $r\mathsf e$; reversal sends $w(V)$ to $w(V^{-1})=-w(V)$, conjugation keeps $\langle\cdot,\cdot\rangle$, and the sign is the product of the signs of the two vectors $u_\lambda$, $u_{\lambda'}$. (b) The reflection in the site plane through $x_{\mathsf e}$ normal to $\lambda$ flips $s_\lambda$ alone, so that it changes the sign of $q_{\mathsf e}$; torus measures are invariant under it, periodic Gibbs states by Proposition~\textup{\YMMref{103.4}{758}}. (c) Each coordinate of a copy lies between two consecutive walls, at a distance $\ge\rho$ from both, and the mean is kept by the invariance of $\tilde\mu_{\ell_0}$ under the lattice isometries. (d) The corners lie within the sup distance $3d$ of the centre of the cell, and $\mathsf l\ge LM$. (e) A tent is carried within $\frac{\sqrt3}2a_K$ of its barycentre and vanishes at the other barycentres; the supports need not come near the centres, which is why only $d\le{\rm dist}\le6d$ holds, and nothing uses an upper bound on the distance.

The settings and the numbers come from the arithmetic of the working level (Lemma~\YMMref{126.1}{844}) and from Lemma~\YMMref{126.2}{845}, by which the objects of the construction satisfy the settings of the chain; neither involves a source, and of the majorants the settings for the $p$ sources read the weights $\mathsf s_a\mathbf 1_{S_a}$ and the bound $\mathsf D_{\mathsf k}\ge3M$ alone. \emph{The threshold.} $T_p(u)=\sum_{n'\ge1}(\frac43)^{p/2}(n'+1)^{1+p\mathfrak p_0}L^{1+(p/2-1)n'}e^{-4\beta_s\kappa[L^{n'-1}u-2]_+}$ does not increase. If $L^qu\ge3$, $q\ge1$, the $q$ terms with $n'\le q$ are each at most $(\frac43)^{p/2}(q+1)^{1+p\mathfrak p_0}L^{1+(p/2-1)q}$, and the term of the index $n'=q+n''$ is at most $(q+1)^{1+p\mathfrak p_0}L^{(p/2-1)q}$ times the term of the index $n''$ of $T_p(3)$, as $n'+1\le(q+1)(n''+1)$ and $L^{n'-1}u\ge3L^{n''-1}$; so $T_p(u)\le(q+1)^{2+p\mathfrak p_0}L^{(p/2-1)q}T^+_p$.

For $1\le j\le\mathsf k$, $\mathsf D_j/M\ge3$ and $T_p(\mathsf D_j/M)\le T_p(3)\le T^+_p$; for $\mathsf k<j<T$, $q=j-\mathsf k\le n_T-1$ has $L^q\mathsf D_j/M=\mathsf D_{\mathsf k}/M\ge3$, $q+1\le n_T\le s$ and $L^{(p/2-1)q}<\mathsf{Lm}_s^{p/2-1}$. And $\theta^{[j]}=\theta(x_j)\le\theta(x^-_s)$, since $x_j\ge x_T>x^-_s\ge e^P$ and $\theta$ does not increase on $[e^P,\infty)$. So $\theta^{[j]}T_p(\mathsf D_j/M)\le\frac12{\rm th}_{g,p}(s)\le\frac12$ for $1\le j<T$.
\end{proof}

\begin{lemma}[Far pairs; centred sources]\label{bil:far}
Let $0<g\le\gamma_2$, $\ell$ small, $n\ge1$, $p=2n$, $s=s(d)\le K$. For pairs $\tau\le n$ let $3d\le|z_{(\tau,1)}-z_{(\tau,2)}|\le4d$, let the majorants have their plaquettes within $d-a_K$ of $z_a$, and let corners of plaquettes of different pairs be at sup distance $\ge11M$; let $(F_a,\mathsf t_a)$ be real quadratic sources of top index $T-1$ under the majorants, with their families without base point at the heights $1\le j\le T-1$ (as forms have them), $\tilde F_a=F_a-\langle F_a\rangle_\ell$, $\tilde P_J=\prod_{\tau\in J}\tilde F_{(\tau,1)}\tilde F_{(\tau,2)}$, $G_\tau$ the Gaussian value of the pair, $\hat m_a=\upsilon\mathsf s_a$, $\hat m_\tau=\hat m_{(\tau,1)}\hat m_{(\tau,2)}$, $\hat m_J=\prod_{\tau\in J}\hat m_\tau$. For $s\ge s_p(g)$, a first depth $\ge s_1(g)$, and all nonempty $J$ and $B$:
\[
|\langle\tilde P_J\rangle_\ell-\textstyle\prod_JG_\tau|\le\omega'_{g,p}(s)\hat m_J,\qquad|\langle\tilde P_J\rangle_\ell-\prod_J\langle\tilde P_\tau\rangle_\ell|\le\omega'^\sharp_{g,p}(s)\hat m_J,\qquad|G_\tau|\le g_0\hat m_\tau;
\]
if $N_2$ pairs lie in $B$ and $N_1$ meet it in one element,
\[
\Bigl|\Bigl\langle\prod_B\tilde F_a\Bigr\rangle_\ell\Bigr|\le\bigl[e_g(s)^{N_1}(g_0+\delta_g(s))^{N_2}+\mathrm{Bl}(p){\rm far}_{g,p}(s)(1+g_0+\varepsilon_g(s))^n+{\rm tl}_{g,p}(s)\bigr]\prod_B\hat m_a,
\]
so that this mean is $\le o_{g,p}(s)\prod_B\hat m_a$, $o_{g,p}=e_g+\mathrm{Bl}(p){\rm far}_{g,p}(1+g_0+\varepsilon_g)^n+{\rm tl}_{g,p}$, if $B$ holds no pair, and $\le{\rm Mo}_{g,p}(s)\prod_B\hat m_a$ in every case, $\sup_{s\ge s_p(g)}{\rm Mo}_{g,p}(s)<\infty$; $Y_s^P\omega'^\sharp_{g,p}(s)$ is bounded. Here
\[
\begin{multlined}
\omega'_{g,p}=n\delta_g(1+g_0+\delta_g)^{n-1}+\mathrm{Bl}(p){\rm far}_{g,p}(1+g_0+\varepsilon_g)^n+{\rm sm}_g(1+g_0)^n+{\rm tl}_{g,p},\qquad\omega'^\sharp_{g,p} \\
=\omega'_{g,p}\bigl[1+n(1+g_0+\omega'_{g,p})^{n-1}\bigr],
\end{multlined}
\]
with $g_0=\frac32 81^2C'^2_{\mathsf k}$ for arbitrary majorants and $g_0=422C'^2_{\mathsf k}$ for majorants of bond weights of one direction. For bond forms $Y[\mathsf w_a]$ with real weights on the bonds of one direction, whose means vanish, $|G_\tau|\le106C'^2_{\mathsf k}\upsilon^2\sup|\mathsf w_{(\tau,1)}|\sup|\mathsf w_{(\tau,2)}|$, and the same functions at $g_0=106C'^2_{\mathsf k}/59^2$, written $\omega_{g,p}$, $\omega^\sharp_{g,p}$, serve in units of $\hat m_a=59\upsilon\sup|\mathsf w_a|$. Here ${\rm far}_{g,p}(s)=C^{\rm far}_p(LC_W(s))^3\bar x_s^{p/2}Y_s^{p\mathfrak p_0}e^{-11L^{s-2}/\mathsf{Lm}_s}$, $\varepsilon_g$ is the error function of equations~\YMMref{(2328)--(2336)}{846}, with the notation in \eqref{eq:short-distance-envelope-names} for a close pair, $\delta_g=\varepsilon_g+e_g^2$, $e_g$ and ${\rm tl}_{g,p}$ are the error functions of the singletons and of the tails,
\[
\begin{gathered}
e_g(s)=2\mathsf c_X(Y^-_s)^{-7\mathfrak p_0}+2(\mathsf c_I+\mathsf c_{\rm tl}\mathsf c_{\rm sb})\,s^2(LC_W(s))^3Y_s^{3\mathfrak p_0}{\rm sm}_g(s),\\
{\rm tl}_{g,p}(s)=\mathrm{Bl}(p)\,\mathsf c_\sigma^p\mathsf c_{\rm sb}^p(1+S^{\rm tl}_p)\,s^{1+3p/2}(LC_W(s))^{3p}Y_s^{2p\mathfrak p_0}\,M^{-3}c'A_0^2Y_s^PL^{3s}\bar x_s^p\,e^{-\frac12c'A_0^2(\log x^-_s)^P},
\end{gathered}
\]
$\mathsf c_\sigma=4+\mathsf C_{\mathsf M}^{1/2}/\mathsf C_{\rm sz}+\sqrt2$, $\mathsf C_{\rm sz}$ of Lemma~\YMMref{125.2}{840} and the constant convention preceding \eqref{eq:short-distance-constants}, $S^{\rm tl}_p=\sum_{a\ge0}L^{(3+4p)a+4p}(a+2)^{2p\mathfrak p_0}e^{-\frac12c'A_0^2(a\log\frac{3L}4)^2}$, ${\rm Mo}_{g,p}=(1+g_0+\delta_g)^n+\mathrm{Bl}(p){\rm far}_{g,p}(1+g_0+\varepsilon_g)^n+{\rm tl}_{g,p}$, $\mathrm{Bl}(p)$ is the number of partitions of $p$ elements, and $s_p(g)$ is the least integer $\ge s_1(g)$ from which ${\rm th}_{g,p}$ of Lemma~\ref{bil:sym}(e), ${\rm far}_{g,p}$ and $e_g$ are at most $1$ and $2\mathsf c_{\rm tl}s\,{\rm sm}_g(s)\le1$; it exists, since these four functions tend to $0$.
\end{lemma}

\begin{proof}[Proof sketch]
By Lemma~\ref{bil:sym}(e) all pairs share one working level, one upper level $T$ and one tower of quotients, and, as ${\rm th}_{g,p}(s)\le1$ for $s\ge s_p(g)$, $\theta^{[j]}T_p(\mathsf D_j/M)\le\frac12$ for $1\le j<T$, so that the mixed derivative families of every order $\le p$ are bounded up to the level $T$ (the bound of the proof of Lemma~\ref{bil:nsrc}). At $T$, almost everywhere, $\Phi^{(T)}[\prod_B\tilde F_a]=\sum_\pi\prod_{\beta\in\pi}\varkappa^{(T)}_\beta$ (the identity of the proof of Lemma~\ref{bil:nsrc}, which holds at every level of the chain in use, here up to $T$, at $c_a=\langle F_a\rangle_\ell$); the conditional cumulants of order $\ge2$ do not depend on the centring constants.

A block that links two pairs is at most ${\rm far}_{g,p}(s)\prod\hat m_a$: by the vanishing of mixed derivatives on polymers that miss a support it is exponentially small in the distance $11M$ measured in units of the level $T$, that is in $L^{K-T}>L^{s-2}/\mathsf{Lm}_s$. A block of one pair is the conditional covariance $\Phi^{(T)}[F_{(\tau,1)}F_{(\tau,2)}]-\Phi^{(T)}[F_{(\tau,1)}]\Phi^{(T)}[F_{(\tau,2)}]$, the quotients being those of the chain of that pair alone (Lemma~\ref{bil:sym}(e)); by Lemma~\ref{bil:two}(c) it is almost everywhere within $\varepsilon_g(s)\hat m_\tau$ of $G_\tau$. A singleton obeys, by the sum over the level of the first creation (equation~\eqref{eq:ward-first-creation}) applied to $F_a$ minus its flat-top constant,
\[
|\Phi^{(T)}[F_a]-\langle F_a\rangle_\ell|\le2\mathsf Y_{T,a}+\rho'_a\le e_g(s)\hat m_a,
\]
$\mathsf Y_{T,a}$ the size of $F_a$ at the level $T$ and $\rho'_a$ its first-creation sum; for bond forms the mean is $0$, and the bound is used only in its last line, so that centring replaces the vanishing of the mean. So the partitions of $B$ without a linking block add up to the product, over the pairs that meet $B$, of $\varkappa^{(T)}_{\{(\tau,1),(\tau,2)\}}+\varkappa^{(T)}_{\{(\tau,1)\}}\varkappa^{(T)}_{\{(\tau,2)\}}=G_\tau+O(\delta_g(s))\hat m_\tau$ if both members lie in $B$, and of a singleton if one does. Crude bounds of the conditional moments at all the levels $t\le T$ (for a block of one pair, the crude clause of Lemma~\ref{bil:two}(c)), multiplied by $\mathsf w_t$ and summed, give the tails ${\rm tl}_{g,p}$.

Summing over the at most $\mathrm{Bl}(p)$ partitions gives the bounds of the lemma, and $\langle\tilde P_\tau\rangle=G_\tau+O(\delta_g)\hat m_\tau$ gives the second from the first. \emph{The size of $G_\tau$.} For plaquettes $p\in S_{(\tau,1)}$, $q\in S_{(\tau,2)}$, $|x(p)-x(q)|\ge3d-2(d-a_K)>d$; so $|C^{\rm bg}(p,q)|\le C'_{\mathsf k}(a_K/d)^3$ (Lemma~\ref{bil:ker}), and $|G_\tau|\le6C'^2_{\mathsf k}(a_K/d)^6(\sum|\mathsf t_{(\tau,1)}|)(\sum|\mathsf t_{(\tau,2)}|)\upsilon^2$. The two sums are counted by volume. The midpoints of the bonds of one direction form a translate of $a_K\mathbb Z^3$. Let $\mathsf w$ be a weight on the bonds of one direction with $S[\mathsf w]$ within $d-a_K$ of $z_a$.

A bond of its support has its four plaquettes in $S[\mathsf w]$, and its midpoint is the midpoint of the centres of two of them, opposite in one plane, so that it lies within $d-a_K$ of $z_a$; the open cubes of side $a_K$ about these midpoints are disjoint and lie within $d-a_K+\frac{\sqrt3}2a_K<d$ of $z_a$, in $B(z_a,d)$, a Euclidean ball as $20d\le\ell$; so the support of $\mathsf w$ has at most $\frac{4\pi}3(d/a_K)^3$ bonds. For a bond form $Y[\mathsf w]$ on the bonds of one direction, $q_b$ being a quarter of a signed sum of four corner forms, $\sum_{p,p'}|\mathsf t(p,p')|\le\sum_b|\mathsf w(b)|\le\frac{4\pi}3(d/a_K)^3\sup|\mathsf w|$, and $6(\frac{4\pi}3)^2<106$. For a majorant of such a form, $\#S\le4\cdot\frac{4\pi}3(d/a_K)^3$, each bond lying in four plaquettes, and $\sum|\mathsf t|\le\frac12(\#S)\mathsf s$, as $2\sum_{p'}|\mathsf t(p,p')|\le\mathsf s$; so $6(\frac{8\pi}3)^2<422$.

For an arbitrary majorant, $\#S\le81(d/a_K)^3$ (the plaquette count following equation~\YMMref{(2300)}{837}) and $\sum|\mathsf t|\le\frac12(\#S)\mathsf s$ give $6(\frac{81}2)^2=\frac32 81^2$.
\end{proof}

\begin{proposition}[Pressure, domination, boards, extraction]\label{bil:press}
(a) For $0<g\le\gamma$, $\ell_1\le\ell_2$ in $\mathfrak L(g)$ and every datum, $0\le p_{\ell_1}-p_{\ell_2}\le G_g$; from Theorem~\YMMref{D}{8} alone, $0\le p_{\ell_1}-p_{\ell_2}\le G^F_g=C_g+C_{\rm maj}+2+2C_{\rm tt}C'_{\rm tt}E_0\le2C_g+2C_{\rm tt}C'_{\rm tt}E_0$.

(b) Let $K$, $\beta>0$, a board of side $\mathsf l_0$, a side $\ell_0\ge2\mathsf l_0$ and $G$ with $p_{\ell_0}-p_\ell\le G$ on an unbounded set of sides be given. For bounded functions $\Phi_{\mathbf k}$ of distinct cells and every $\nu\in\mathcal G^{\rm per}_{K,\beta}$ (and for $\tilde\mu_\ell$ if the cells are distinct modulo the period): $|\nu(\prod\Phi_{\mathbf k})|\le\prod e^{G\mathsf l_0^3}\|\Phi_{\mathbf k}\|_{{\rm ch},\ell_0}$. No reflection invariance of $\nu$ is used.

(c) Let $h>0$ with $\mathsf l_0\in h\mathbb N$, as $h=a_K$, or $h=a_{K-t}$ for a level $t$. If $n(r+2\rho+h)<\mathsf l_0$, any $n$ cubes of side $r$ admit an offset in $h\mathbb Z^3$ with all their points at a distance $\ge\rho$ from the walls, also after translations by $\mathsf l_0\mathbb Z^3$; at $h=a_{K-t}$ the walls lie along faces of the blocks of the level $t$.

(d) Let functionals $\Lambda^{(j)}_n$ on tuples in $C_c(\mathbb R^3;V)$ with disjoint supports be linear in each entry on the tuples in $B(0,R')$ from an index $j_0(n,d,R,R')$ on, and bounded there by $c_n(d;R)\prod|H_a|_\infty$ at the tuples of type $(d;R)$ (supports in balls of radius $R$, at mutual distance $\ge d$). Then every subsequence has a subsequence that converges at all tuples, to multilinear limits with the same bounds.
\end{proposition}

\begin{proof}[Proof sketch]
(a) At the source $0$ the localised source energies of Lemma~\ref{lem:canonical-source-system} give $E/\ell^3-C_g\le p_\ell\le E/\ell^3+C_S$ and $|E_\ell/\ell^3-E_{\ell'}/\ell'^3|\le2C_{\rm tt}C'_{\rm tt}E_0$, and $p_\ell$ does not increase with $\ell$ by the chessboard estimate (Corollary~\YMMref{86.6}{683}). Without Theorems~\YMMref{AJ}{673} and~\YMMref{AD}{607}, $\tilde Z(0)=e^{-E}Z$ with $e^{-C_g\ell^3}\le\tilde Z(0)\le e^{(C_{\rm maj}+2)\ell^3}$ gives $G^F_g$. (b) The chessboard estimate on $T_\ell$ (Proposition~\YMMref{86.3}{681}) gives $|\langle\prod\Phi_{\mathbf k}\rangle_\ell|\le\prod\|\Phi_{\mathbf k}\|_{{\rm ch},\ell}$, and comparing $Z_\ell$ with $Z_{\ell_0}$ cell by cell gives $\|\Phi\|_{{\rm ch},\ell}\le e^{\mathsf l_0^3(p_{\ell_0}-p_\ell)}\|\Phi\|_{{\rm ch},\ell_0}$ (Corollary~\YMMref{86.6}{683}).

For every periodic Gibbs state $\nu$ and every unbounded set of sides, the domination theorem (Proposition~\YMMref{91.9}{705}) gives $|\nu(\prod\Phi_{\mathbf k})|\le\liminf_\ell\prod\|\Phi_{\mathbf k}\|_{{\rm ch},\ell}$, the lower limit taken along the sides of the set; along the sides with $p_{\ell_0}-p_\ell\le G$ the comparison bounds the right member by $\prod e^{G\mathsf l_0^3}\|\Phi_{\mathbf k}\|_{{\rm ch},\ell_0}$. (c) In each direction there are $\mathsf l_0/h$ offsets modulo $\mathsf l_0$, and fewer than $n((r+2\rho)/h+1)<\mathsf l_0/h$ of them are bad. (d) The bound gives the equicontinuity $|\Lambda(X)-\Lambda(X')|\le c_n\sum_a|H_a-H_a'|_\infty\prod_{a'\ne a}\max\{|H_{a'}|_\infty,|H'_{a'}|_\infty\}$; a countable set of Lipschitz functions is dense with control of the supports; a diagonal sequence converges on it, hence everywhere.
\end{proof}

\begin{lemma}[The rate on a fixed torus]\label{bil:rate}
Let $g\le\gamma_2$, $\mathsf b\in I(g)$, $\ell\in\mathfrak L(g)$, $q\ge2$, $d>0$, $\theta_q=1/(1260q(8q+3))$, $\mathsf p=\max\{2\mathfrak p_0+2,\mathsf p_0,1\}$, where $\mathsf C_{\rm bf}\ell^3K^{\mathsf p_0}L^{-K/10}$, $K\ge K_{\rm bf}$, is the comparison of the laws of the block fields at two cutoffs of Corollary~\textup{\ref{IV-62-66}}. There are $C$, of the constants, $q$ and $\mathsf C_{\rm bf}$, and $K_1$, besides of $g,\mathsf b,\ell,d$, such that for $K\ge K_1$ and $G_K=\langle\mathcal E_1;\dots;\mathcal E_q\rangle_{K,\ell,L^K\mathsf b}$, the $\mathcal E_a$ any of $\mathcal O_K(\varphi)=O^T_K(\varphi\mathbf 1)$, $O^T_K(\varphi E_{\lambda\lambda})$ (the plaquettes of the plane normal to $\lambda$ alone) and $O^T_K(\varphi E_{\lambda\lambda'})$, $\lambda\ne\lambda'$, read at the centres of the plaquettes and the bonds, with real $\varphi_a$, Lipschitz in $|\cdot|_\infty$, whose supports are pairwise at sup distance $\ge d$:
\[
|G_{K+1}-G_K|\le C\ell^{3(q+1)}K^{\mathsf p}L^{-\theta_qK}\prod_a(\sup|\varphi_a|+\mathrm{Lip}\,\varphi_a).
\]
So $\lim_KG_K$ exists, with the tail of the series as the error; any number of the entries may be off-diagonal, and one pair $(C,K_1)$ serves a compact interval of labels. The same holds with $\prod_a\sup|\psi_a|$ for bounded Borel functions $\psi_a$ that vanish off closed sets at pairwise sup distances $\ge d$, entered through their cell means: at a site $x$ of $\Lambda_{K,\ell}$ the mean $\psi_a^{(K)}(x)=a_K^{-3}\int_{x+a_K[-\frac12,\frac12)^3}\psi_a$ is the weight of the plaquette with the corner $x$ and the edges $a_Ke_\kappa$, $a_Ke_{\kappa'}$, for the component $\lambda\lambda$, $\{\kappa,\kappa',\lambda\}=\{1,2,3\}$, and of the bond $\langle x,x+a_Ke_\mu\rangle$, for the component $\lambda\lambda'$, $\lambda\ne\lambda'$, $\mu$ the third direction, for all six components.

As $L$ is odd, the cells of $\Lambda_{K+1,\ell}$ at the sites of the block $B(y)$ of a site $y$ of $\Lambda_{K,\ell}$ make up the cell of $y$, so that $\psi^{(K)}(y)=L^{-3}\sum_{x\in B(y)}\psi^{(K+1)}(x)$: the weights of two cutoffs are tied exactly. So the limit exists for test functions of no regularity and agrees with that for Lipschitz functions. For one function evaluated at the points of the lattices of two cutoffs no bound of this kind by the sup norms holds. For one off-diagonal entry among action densities the rate is free of their number: if $n\ge1$, $\lambda\ne\lambda'$, and $\varphi,f_1,\dots,f_n$ are real Lipschitz functions on $T_\ell$, none of them $0$, whose $n+1$ supports have no common point, then $\kappa_K=\langle O^T_K(\varphi E_{\lambda\lambda'});\mathcal O_K(f_1);\dots;\mathcal O_K(f_n)\rangle_{K,\ell,L^K\mathsf b}$ obeys
\[
|\kappa_{K+1}-\kappa_K|\le CK^pL^{-K/120}(\sup|\varphi|+\mathrm{Lip}\,\varphi)\prod_{i=1}^n(\sup|f_i|+\mathrm{Lip}\,f_i)\qquad(K\ge K_2),
\]
with $C,p,K_2$ of the constants, $n$, $\ell$, the end points of a compact interval of labels that holds $\mathsf b$ and the spread $\mathsf d_\infty$ of the $n+1$ supports (the definition preceding Lemma~\YMMref{114.6}{805}).
\end{lemma}

\begin{proof}[Proof sketch]
Let $s<\sigma'\le\lceil K/1260\rceil$ be depths, the second at most the depth up to which Corollary~\textup{\ref{IV-62-66}} compares the laws of the block fields, and $\sigma'\ge(8q+3)s$. At each of the cutoffs $K$, $K+1$ the cumulant is, up to three errors, the cumulant of the first members of the sources, that is of the functions $\mathcal E_a$ themselves and not of their families at the heights, evaluated at the background of the block field of the depth $\sigma'$: the progression with the sources switched on at the level $K-\sigma'$ forgets what lies below, with the error $\epsilon(\sigma'-s)$, $\epsilon(\nu)=C(\nu+1)^{2\mathfrak p_0+2}L^{-\nu/2}$; the large fields cost $L^{-(b'+1)s}$, for any $b'\ge4q$; and the distance of the supports costs $e^{3-L^sd/M}$.

Each of the three multiplies $(4\ell^3L^{4s})^q\prod_a(\sup|\varphi_a|+\mathrm{Lip}\,\varphi_a)$, the bound of the entries at the depth $s$; as $L^{4qs}L^{-(\sigma'-s)/2}\le L^{-s}$ for $\sigma'-s\ge(8q+2)s$ and $L^{4qs}L^{-(b'+1)s}\le L^{-s}$ for $b'\ge4q$, together they are of the order $K^{2\mathfrak p_0+2}L^{-s}$. Over one block field $V$ of the depth $\sigma'$, with $t=K-\sigma'$, the background $U^{\rm A}$ on $\Lambda_{K,\ell}$ and the background $U^{\rm B}$ on $\Lambda_{K+1,\ell}$ are compared in three links. (1) $U'=\mathcal M(U^{\rm B})$, $U^{\rm B}$ averaged once, is a configuration of the coarse lattice over the same block field, a competitor whose excess energy over $U^{\rm A}$ is at most $C\varepsilon_{\rm r}^2\eta_tN_t$, $\varepsilon_{\rm r}$ the radius of the small fields of the depth $\sigma'$ and $N_t$ the number of the cubes of the level $t$: its energy is at least that of $U^{\rm A}$, by the minimality of $U^{\rm A}$, and at most that of $U^{\rm B}$ up to a number of the same order, by the inequality of Jensen of (2); and the energy of $U^{\rm B}$ exceeds that of $U^{\rm A}$ by at most a number of the same order, since the prolongation of $U^{\rm A}$ to the fine lattice is a competitor of $U^{\rm B}$ whose energy exceeds that of $U^{\rm A}$ by at most such a number. (2) The fine background against its average: for a bilinear form of the loops at one site or at one bond, the defect of Jensen's inequality $\sum_p\mathsf W(p',p)|g_p|^2\ge L^{-2}|G'_{p'}|^2$, with the averaging kernel $\mathsf W$, $g_p$ the fine and $G'_{p'}$ the averaged curvature, is the variance of the fine curvature on the block of $p'$, and its sum over the blocks is the difference of the two energies; so the curvature is nearly constant in mean square under the shifts of a plaquette in its own plane within a block, with the error $L^{-t/2}$, and the weights enter only through the mismatch of their block means, which is $0$ for cell means and at most $\frac32a_K\mathrm{Lip}\,\varphi$ for one function read at the centres on the two lattices. (3) The competitor against $U^{\rm A}$: along the $t$ consecutive levels each pair of backgrounds differs by one gauge transformation and one small $1$-form $A_k$, with $\sum_kL^{t-k}\|A_k\|^2$ at most a number times the excess, which costs $(t\eta_t)^{1/2}$.

So the corner fields $a_K^{-1}\sum_x\varphi(x+a_K\xi)Q_x(l,l')$, $Q_x(l,l')$ the corner forms of Lemma~\ref{bil:mat}(b), of which $O^T_K(\varphi E_{\lambda\lambda'})$ is a sum of four with the coefficients $\pm\frac14$, agree at the two backgrounds up to $C\ell^3L^{4\sigma'}K^{1/2}L^{-(K-\sigma')/2}(\sup|\varphi|+L^{-\sigma'}\mathrm{Lip}\,\varphi)$ for Lipschitz weights, and up to the same number with $\sup|\psi|$ for weights tied by the cell mean; for one function evaluated at the points of two lattices the mismatch of the block means is not small in the sup norm, and there the Lipschitz number pays. The block-field laws at $K$ and $K+1$ differ by $\mathsf C_{\rm bf}\ell^3K^{\mathsf p_0}L^{-K/10}$ (Corollary~\textup{\ref{IV-62-66}}(a), whose numbers read the label only through an interval that holds it).

With $\sigma'=\lfloor K/(1260q)\rfloor$ and $s=\lfloor\sigma'/(8q+3)\rfloor$ the four errors give the rate $\theta_q$. For one off-diagonal entry among action densities whose supports have no common point another road keeps the action densities as sources in the exponent of the progression, as in the comparison of two cutoffs with sources (Theorem~\YMMref{AR}{810}), and passes the one off-diagonal entry alone to the block field as a bounded function: it is put into the chain as a member whose tower is split at the depth $\lceil K/36\rceil$; what it creates below the level $K-\lceil K/36\rceil$ is born at a small coupling, and above it the member is carried along the run and forgotten.

No factor $L^{4q\sigma'}$ enters, the rate is free of $n$, and it is $L^{-K/120}$.
\end{proof}

\begin{lemma}[The seminorm of one tensor insertion]\label{bil:semi}
Let $g\le\gamma_1$ and $0<\delta\le\frac14$, and let $s_E(i)$ be the depth $s(\delta)$ formed at the coupling $g^\sharp_i$, $(g^\sharp_i)^{-2}=L^{i-1}(X_g+\bar B)-B_\sharp$, with the flow bound $B_\sharp$ of Lemma~\ref{bil:two}(b). There are $i_E(g,\delta)$ and numbers $c_E$ and $C_E$ ($C_E$ proportional to $\delta^{-4}$), of the constants, $g$ and $\delta$, such that for $i\ge i_E(g,\delta)$, every datum $(K,\beta)$ at $g$ with $K\ge K_E(i)=i+s_E(i)$, every board of side $a_i$ of $\mathbb L_K$, every cell of it, every $\ell\in\mathfrak L(g)$ and every bounded $H:\mathbb R^3\to V_{\mathbb C}$ that vanishes at the points whose sup distance from the complement of the cell is $<\delta a_i$:
\[
\|O^T_K(H)-c\|_{{\rm ch},\ell}\le2\mathsf E_i\,b^{-1}|H|_\infty,\qquad\mathsf E_i=510\,\delta^{-3}e^{G_g+1/8}L^{9i/2}\mathsf A_N(s)\le C_E\mathsf Y_i^{q_E}L^{11i/2},
\]
for $c$ the mean of $O^T_K(H)$ in any $\tilde\mu_{\ell_1}$, $\ell_1\in\mathfrak L(g)$, or in any periodic Gibbs state at $(K,\beta)$; here $\mathsf Y_i=2[c_E+\log(1/a_i)]$, $q_E=(4+\mathfrak p_0)(P-1)+10\mathfrak p_0$, $8\le N\le\frac19c'A_0^2\mathsf Y_i^{P-1}$ is the number of cells of the least torus at the coupling stopped at the level $i$, and $s=s(\delta)$ is formed at that coupling. By the dilation of Lemma~\ref{bil:dil}(b) the same holds for $g\le\gamma_2$, with numbers depending on $g$.
\end{lemma}

\begin{proof}[Proof sketch]
Let $k=K-i$ and let $g_{(i)}$ be the coupling stopped at the level $i$, $g_{(i)}^{-2}=L^{-1}x_k(\beta)-B_\sharp$. By Corollary~\YMMref{74.3}{616} and Lemma~\ref{lem:setup-stopped-arithmetic}, $(k,\beta)$ is a first-passage datum at $g_{(i)}<g$ with the running couplings $x_0,\dots,x_k$ and the label $L^ib$, $L^i\ell\in\mathfrak L(g_{(i)})$, and $C_{g_{(i)}}\le C_g+\frac92i\log L+\frac18$. The dilation $x\mapsto L^ix$ (Lemma~\ref{lem:lattice-law-dilation}) carries $\mathbb L_K$ to $\mathbb L_k$, the board of side $a_i$ to a board of side $1$, $O^T_K(H)$ to $O^T_k(H')$, $H'=L^iH(L^{-i}\cdot)$, with $|H'|_\infty/(L^ib)=|H|_\infty/b$, and seminorms into seminorms. At $g_{(i)}$ a cell of side $1$ is a cell of the least torus $T_{\ell''}$, with $N=\ell''^3$ cells, and the $N$-th power of the seminorm on $T_{\ell''}$ is the mean of the product of the $N$ copies.

\emph{The threshold.} The hypotheses of Lemma~\ref{bil:prod} at $n=N$, $g'=g_{(i)}$, $K'=k$, $d=\delta$ are the depth condition $s(\delta)\le k$, $s(\delta)$ formed at $g_{(i)}$, and three inequalities. As $g^{-2}<x_K\le X_g$ and $L^i(x_K-\bar B)+\bar B\le x_k\le L^i(x_K+\bar B)-\bar B$ (equation~\YMMref{(1530)}{615}),
\[
L^{i-1}(g^{-2}-\bar B)-B_\sharp\le g_{(i)}^{-2}\le(g^\sharp_i)^{-2}=L^{i-1}(X_g+\bar B)-B_\sharp ;
\]
$g_{(i)}$ reads the datum, the two bounds do not. Since $C_W(s)$ increases with the inverse square of the coupling at which it is formed, the depth $s(\delta)$ formed at $g_{(i)}$ is at most $s_E(i)$, and $K\ge K_E(i)=i+s_E(i)$ gives $s(\delta)\le s_E(i)\le K-i=k$: the depth condition. The least torus $T_{\ell''}$ at $g_{(i)}<g\le\gamma_1$ is small over the running couplings $x_0,\dots,x_k$ (Lemma~\YMMref{126.3}{846}). The three inequalities hold once $g_{(i)}^{-2}\ge x_E(\delta)$, a number of the constants and $\delta$. The first holds as soon as $Y_s$ exceeds a number of the constants alone, since $Y_s\ge\log g_{(i)}^{-2}$ and $N\le\frac19c'A_0^2Y_s^{P-1}$, so that $u_*(N)\le CN\log(eN)$ grows more slowly than $Y_s^P$; the second is a lower bound on $y_k$, and $y_k>\log g_{(i)}^{-2}$; for the third, $L^s\delta<LC_W(s)$ gives $L^{3s}<(Lc_W/\delta)^3\bigl(2[c_Y(\delta)+\log2g_{(i)}^{-2}]\bigr)^{6\mathfrak p_0}$, $c_W=C_{W,0}A_0^2$, $c_Y(\delta)=\log(L^2c_W/\delta)+2\mathfrak p_0\log(4\mathfrak p_0)$, and this is $\le g_{(i)}^{-40}<x_k^{20}$ from a number $\le x_*/\delta$ on, $x_*$ a number of the constants. So $x_E(\delta)\le x_*/\delta$. Let $i_E(g,\delta)$ be the least $i\ge1$ with $L^{i-1}(g^{-2}-\bar B)-B_\sharp\ge x_*/\delta$; as $g_{(i)}^{-2}\ge L^{i-1}(g^{-2}-\bar B)-B_\sharp$, the hypotheses hold for $i\ge i_E(g,\delta)$, and $L^{i_E(g,\delta)}\le C(g)\delta^{-1}$.

\emph{One piece.} Let $H$ be real first. The points of the cell at sup distance $\ge\delta$ from its walls form a cube of side $1-2\delta$, the union of at most $2\delta^{-3}$ boxes of sides $\le0.99\delta$, each within $\frac78\delta$ of its centre; on each box $O^T_k(H')$ has a plaquette piece, of size $\le|H'|_\infty$, and a bond piece, of size $\le59|H'|_\infty$ (Lemma~\ref{bil:forms}(a)). Let $F$ be one piece, of size $\mathsf s$, and $F_\kappa$ its $N$ mirror copies, one in each cell: forms of the same kind and size (Lemma~\ref{bil:sym}(c)), placed as in Lemma~\ref{bil:nsrc} at the distance $\delta$, since copies in different cells are $\ge2\delta$ apart. Lemma~\ref{bil:prod} gives complex numbers $c_\kappa$ with $|\langle\prod_{\kappa\in B}(F_\kappa-c_\kappa)\rangle_{k,\ell'',\beta}|\le(\mathsf A_N(s)\upsilon\mathsf s)^{\#B}$ for every nonempty set $B$ of cells. The mean $m_F$ of $F$ on $T_{\ell''}$ is that of each copy, so $|m_F-c_\kappa|\le\mathsf A_N(s)\upsilon\mathsf s$, the case $B=\{\kappa\}$; expanding $\prod_\kappa(F_\kappa-m_F)=\prod_\kappa\bigl((F_\kappa-c_\kappa)+(c_\kappa-m_F)\bigr)$ into its $2^N$ terms gives $\|F-m_F\|_{{\rm ch},\ell''}\le2\mathsf A_N(s)\upsilon\mathsf s$.

\emph{The sum.} The seminorm is subadditive (Proposition~\YMMref{86.4}{682}); with $m$ the mean of $O^T_k(H')$ on $T_{\ell''}$, the sum of the $m_F$, and $\upsilon\le(\frac98)^{1/2}(L^ib)^{-1}$,
\[
\|O^T_k(H')-m\|_{{\rm ch},\ell''}\le2\delta^{-3}\cdot2\mathsf A_N(s)\upsilon\cdot60|H'|_\infty\le255\,\delta^{-3}\mathsf A_N(s)\,b^{-1}|H|_\infty ;
\]
for complex $H$ the real and imaginary parts give the number $510$. The domination of $T_\ell$ by $T_{\ell''}$ (Proposition~\ref{bil:press}(b)) costs $e^{G_{g_{(i)}}}\le e^{G_g+1/8}L^{9i/2}$, and the means in the other laws differ from the mean on $T_{\ell''}$ by the same bound. Since $\bar x_s$ at the stopped coupling is at most $C\delta^{-1}\mathsf Y_i^PL^i$ and $N$ is a power of $\mathsf Y_i$, $\mathsf A_N(s)$ is a power of $\mathsf Y_i$ times $L^i$: of the power $\frac{11}2$, $\frac92$ come from the pressure gap at the stopped coupling and $1$ from $\bar x_s$.
\end{proof}

\begin{lemma}[Vectors and kernels at ordered heights]\label{bil:vec}
Let $\sigma$ be a real symmetric functional on the separated tuples with continuous entries, multilinear and bounded as in (D1), invariant under the translations, reflection positive at $\{x_3=0\}$, with $H\ge|P_1|$, $H\ge|P_2|$. In (b) let it have besides the insertion bound of (F2)(vi) with a number $C_V$ for $C_p\mathsf b^{-1}$ and a power $p$. Let $\mathsf{Tb}=\{\mathfrak z\in\mathbb C^3:\mathrm{Re}\,\mathfrak z_3>|\mathrm{Im}\,\mathfrak z_1|+|\mathrm{Im}\,\mathfrak z_2|\}$ and $\mathsf E(\mathfrak z)=\int e^{-\mathfrak z_3E+i(\mathfrak z_1q_1+\mathfrak z_2q_2)}d\mathcal P$.

(a) $\mathsf E$ is a holomorphic semigroup of contractions on $\mathsf{Tb}$, with $\mathsf E((v,t))=e^{-tH}\mathcal U(v)$; if $\mathcal P$ is carried by $\{E\ge|q|\}$, the same holds on $\{\mathrm{Re}\,\mathfrak z_3>|(\mathrm{Im}\,\mathfrak z_1,\mathrm{Im}\,\mathfrak z_2)|\}$.

(b) For $n\ge2$ there are functions $\mathbf W_n(\mathfrak z;e)$ on $\mathsf{Tb}^{n-1}\times V_{\mathbb C}^n$, holomorphic, and $n$-linear in $e$, with $|\mathbf W_n(\mathfrak z;e)|\le(C_V\eta^{-p})^n\prod_k|e_k|$ whenever $\mathrm{Re}\,\mathfrak z_{k,3}-|\mathrm{Im}\,\mathfrak z_{k,1}|-|\mathrm{Im}\,\mathfrak z_{k,2}|\ge2\eta$ for all $k$, $0<\eta\le1$, such that
\[
\sigma(H_1,\dots,H_n)=\int\mathbf W_n(x_2-x_1,\dots,x_n-x_{n-1};H_1(x_1),\dots,H_n(x_n))\,dx
\]
for every tuple whose supports are ordered in the third coordinate, $\sup_{\mathrm{supp}\,H_k}x_3<\inf_{\mathrm{supp}\,H_{k+1}}x_3$; so at ordered heights $\sigma$ has real-analytic kernels.

(c) Suppose besides that $\sigma$ is invariant under $W_3$, or else, in the place of $W_3$ and of the bounds $H\ge|P_1|$, $H\ge|P_2|$, that $\sigma$ is reflection positive at the four planes through $0$ with the normals $(e_k\pm e_3)/\sqrt2$, $k=1,2$. Then $\Omega$ and the classes of the tuples with entries in pairwise disjoint slabs $\{c<x_3<c'\}$ are total in the Hilbert space: for every tuple $X$ in $\{x_3>0\}$ and every $\epsilon>0$ there are finitely many such tuples $X_i$ with $\|[X]-\sum_i[X_i]\|<\epsilon$. The insertion bound is not used in (c).
\end{lemma}

\begin{proof}[Proof sketch]
(a) On the support $\{E\ge\max(|q_1|,|q_2|)\}$ of $\mathcal P$ the integrand has modulus at most $e^{-(\mathrm{Re}\,\mathfrak z_3-|\mathrm{Im}\,\mathfrak z_1|-|\mathrm{Im}\,\mathfrak z_2|)E}\le1$, and its derivatives are bounded locally uniformly by $Ee^{-\delta E}\le1/(e\delta)$; on $\{E\ge|q|\}$ the Schwarz inequality gives the larger tube. (b) The slab operator of $h$ supported in $\{0<x_3<l\}$, $\langle[F'],\mathsf{Sl}_l(h)[F]\rangle=\sigma(\theta\bar F'\cdot(h)\cdot\tau_{le_3}F)$, has $\|\mathsf{Sl}_{2\eta}(\tau_{\eta e_3}h)\|\le C_V\eta^{-p}\|h\|_1$ by the insertion bound, and $e^{-tH}\mathsf{Sl}_l(h)e^{-t'H}=\mathsf{Sl}_{l+t+t'}(\tau_{te_3}h)$. For approximate $\delta$-functions $h=\delta_\epsilon(\cdot-y)e$ these operators converge in norm, by the covariance under the translations and the semigroup, to point operators of norm $\le C_Vs^{-p}|e|$, $|e|=\max_{\lambda\lambda'}|e_{\lambda\lambda'}|$, at the distance $s$ from the boundary of the tube (the insertion bound holds for complex $h$), holomorphic in their arguments.

The vectors of $n$ point insertions at ordered heights are products of these operators and of $\mathsf E(\mathfrak z)$ applied to $\Omega$; their scalar products with $\Omega$ are the functions $\mathbf W_n$, with one factor $C_V\eta^{-p}$ for each insertion and no factorial; at ordered tuples the reconstruction writes $\sigma$ as the integral of $\mathbf W_n$. (c) Under the second hypothesis the reflections in $\{x_3=0\}$ and in the four diagonal planes leave $\sigma$ invariant (Proposition~\ref{bil:cone}(a)) and generate $W_3$, and positivity at two planes at the angle $\frac\pi4$ gives $H\ge|P_1|$, $H\ge|P_2|$ (Proposition~\YMMref{94.4}{722}); so the first hypothesis holds, and we suppose it. By the symmetry of $\sigma$ a class does not depend on the order of the entries.

Let $X$ be a tuple of $n$ entries in $\{x_3>0\}$ and $\epsilon>0$, the supports of the entries at pairwise sup distances $\ge4\delta_0>0$. Cut the entries by a partition of unity subordinate to the boxes of a grid of a mesh $h$ with $(n-2)h\le\delta_0$ and $h\le\delta_0$, and mollify the pieces; by multilinearity and the bounds of (D1) the class changes by less than $\frac\epsilon2$, and it becomes a finite sum of classes of tuples of $n$ smooth pieces, one from each entry, each piece carried by one box of side about $h$, the boxes of different pieces at the sup distance $\ge2\delta_0$. Call two pieces of such a tuple a \emph{bad pair} if the ranges of their heights meet. The two boxes of a bad pair are then $\ge2\delta_0$ apart in a horizontal direction $\mu\in\{1,2\}$; the projections on the axis of $x_\mu$ of the $n-2$ other boxes cover at most $(n-2)h\le\delta_0$ of the gap between them, so that they leave at most $n-1$ free intervals of a total length $\ge\delta_0$, one of length $\ge\delta_0/(n-1)$.

The vertical plane $\{x_\mu=c\}$ through its midpoint is at the distance $\ge\delta'=\delta_0/(2(n-1))$ from all the boxes, and it parts the pieces into two groups $X_{\rm L}$ and $X_{\rm R}$, the bad pair on both sides. The quarter turn of $W_3$ that exchanges $e_\mu$ and $e_3$, with a translation, carries the picture of $\{x_3=0\}$ to that of $\{x_\mu=c\}$, with its positivity and its bounds $H\ge|P_1|$, $H\ge|P_2|$. Put $\psi(\varrho)=[X_{\rm L}\cdot\tau_{\varrho e_3}X_{\rm R}]$ for $\varrho$ above a negative number $\varrho_0$, so that the shifted tuple lies in $\{x_3>0\}$; vertical shifts leave the horizontal projections unmoved, so the shifted tuple stays separated and the plane parts it for every $\varrho$. $\psi$ is continuous in norm, since $\|\psi(\varrho)-\psi(\varrho')\|^2$ is a combination of values of $\sigma$ that are continuous in the translations by (D1).

Its Gram kernel $G(s,s')=\langle\psi(s),\psi(s')\rangle$ is the value of $\sigma$ at a tuple whose two groups, $\theta\bar X_{\rm L}\cdot X_{\rm L}$ and $B(s,s')=\tau_{-se_3}\theta\bar X_{\rm R}\cdot\tau_{s'e_3}X_{\rm R}$, the plane $\{x_\mu=c\}$ parts; as $\tau_{re_3}B(s,s')=B(s-r,s'+r)$, in the picture of that plane the translation along $e_3$, a direction of momentum there, moves $(s,s')$ along the antidiagonals alone:
\[
G(w-u,w+u)=\bigl\langle\Phi_{\rm L},\mathcal U_\mu(ue_3)\Phi_{\rm R}(w)\bigr\rangle_\mu,\qquad\Phi_{\rm L}=e^{-\delta'H_\mu}\Phi'_{\rm L},\quad\Phi_{\rm R}(w)=e^{-\delta'H_\mu}\Phi'_{\rm R}(w),
\]
with the classes $\Phi'$ of the two groups moved by $\delta'$ towards the plane. As $H_\mu\ge|P|$ in that picture, $u\mapsto G(w-u,w+u)$ continues holomorphically to the strip $|\mathrm{Im}\,u|<2\delta'$, bounded there by $\|\Phi'_{\rm L}\|\,\|\Phi'_{\rm R}(w)\|$, locally uniformly in real $w$ by (D1); nothing is said in $w$. This is enough, by a lemma on Gram kernels: if $\psi$ is continuous in norm on an interval and $u\mapsto G(w-u,w+u)$ continues to a strip $|\mathrm{Im}\,u|<\kappa$ with a bound uniform for $w$ in compact parts of the interval, then $\psi$ continues holomorphically in norm to a complex neighbourhood of the interval. Indeed, put $\psi_\lambda(z)=(\lambda/\pi)^{1/2}\int e^{-\lambda(z-s)^2}\psi(s)\,ds$, over a compact part of the interval; then
\[
\|\psi_\lambda(x+iy)\|^2=\frac{2\lambda}\pi\iint e^{-2\lambda(x-w)^2-2\lambda(u-iy)^2}G(w-u,w+u)\,dw\,du ,
\]
the contour shift $u\mapsto u+iy$ takes the factor $e^{2\lambda y^2}$ away, its ends cost $e^{-2\lambda(\kappa_1^2-y^2)}$ for a $\kappa_1<\kappa$, and so $\sup_\lambda\|\psi_\lambda\|$ is finite on a complex neighbourhood of every inner point; $\psi_\lambda\to\psi$ on the reals as $\lambda\to\infty$, and Vitali's theorem, applied to $\langle\chi,\psi_\lambda\rangle$ for every vector $\chi$, gives the continuation. So for every vector $\chi$ the function $\varrho\mapsto\langle\chi,\psi(\varrho)\rangle$ is real-analytic on $(\varrho_0,\infty)$. For large $\varrho$ the shifted tuple has at least one bad pair fewer, and its bad pairs are earlier bad pairs inside $X_{\rm L}$ or inside $X_{\rm R}$, in unchanged relative position, so that the same $h$, $\delta_0$ and $\delta'$ serve at every step. By induction on the number of bad pairs its class then lies in the closed span $\mathcal H^{\rm ord}$ of $\Omega$ and of the classes of the tuples with entries in pairwise disjoint slabs; for a vector $\chi$ orthogonal to $\mathcal H^{\rm ord}$ the function $\langle\chi,\psi(\varrho)\rangle$ vanishes for large $\varrho$, hence on $(\varrho_0,\infty)$ and at $\varrho=0$. So $[X]$ is within $\frac\epsilon2$ of $\mathcal H^{\rm ord}$, and finitely many tuples $X_i$ with entries in disjoint slabs give $\|[X]-\sum_i[X_i]\|<\epsilon$.
\end{proof}

\begin{lemma}[A Ward bound at the top levels; the conserved combination]\label{bil:ward}
Let $0<g\le\gamma$, $\ell'=2\mathsf l\in\mathfrak L(g)$, $0<\delta\le\frac18\mathsf l$, $\phi=\frac1{400}$, $t=K-\lceil\phi K\rceil$. There are $C$, $p$, $K_0$, of the constants, $g$, $\ell'$ and $\delta$, such that for every datum $(K,\beta)$ at $g$ with $K\ge K_0$, every board of side $\mathsf l$ of $\mathbb L_K$ with an offset in $a_{K-t}\mathbb Z^3$, every cell of it and every $\eta\in C^2_c(\mathbb R^3;\mathbb R^3)$ supported in the cell at a sup distance $\ge\delta$ from its walls, the means being those of the periodisations on $T_{\ell'}$:
\[
\bigl\|O^T_K(H^{\rm cons}[\eta])-\langle O^T_K(H^{\rm cons}[\eta])\rangle_{\ell'}\bigr\|_{{\rm ch},\ell'}\le CK^pL^{-K/36}\|\eta\|_{C^2}.
\]
Let $\mathsf{St}_K[\eta]=\Theta_\beta[\eta]-b\,\mathcal O_K(\nabla^{\rm s}\!\cdot\eta)$, with the stress insertion $\Theta_\beta[\eta]$ and the lattice divergence $\nabla^{\rm s}\!\cdot\eta$ of Section~\ref{sec:ward-general-calculus} and equations~\YMMref{(2104)--(2105)}{781}, taken at the plaquette weight $\beta$ and at the vector weight $p\mapsto\eta(x(p))$; it is the derivative of the action $\beta\sum_pa_p$ along the displacement derivation of $\eta$ (Section~\ref{sec:ward-general-calculus}) less the orbital insertion $\mathfrak K^{\rm v}[\beta\eta]$ of Section~\ref{sec:ward-general-calculus} and equations~\YMMref{(2104)--(2105)}{781}. It obeys the same bound with $L^{-K/28}$, and $\mathsf{St}_K[\eta]-b\,O^T_K(H^{\rm cons}[\eta])$ obeys it with $L^{-K/36}$.
\end{lemma}

\begin{proof}[Proof sketch]
\emph{Entries.} An entry of order two is a form of the Setting, at plaquettes, bonds or sites, with weights $v$, the factor $a_K^{-1}$ standing outside them, and $\|v\|$ the largest of $\sup|h|$ and $\sup_{\mathsf c}\sum_{l,l'}|v(\mathsf c,l,l')|$. Its symbol is the $V$-valued weight $S$ for which its second order in the curvature, at a background whose curvature is constant on the cells, is $\frac12\sum_{\lambda\lambda'}S_{\lambda\lambda'}\langle B_\lambda,B_{\lambda'}\rangle$ summed over the cells; for $O^T_K(H)$ it is $H$. A \emph{null entry} is an entry of order two whose symbol is $0$, such as the difference of two corner forms $Q_x(l,l')$ whose plaquettes have the same normals and the same signs $\sigma(l)\sigma(l')$; it is another thing than the null terms of (i). Entries of the orders three and four are sums over cells of polynomials in the bond variables of a cell that vanish to the third, resp.\ fourth, order at the trivial configuration, with weights $\varkappa$ that hold the factor $a_K^{-1}$. The regular configurations of a level are the backgrounds of the block fields of its small-field set. The oscillation of an entry at the level $t$ is $\omega_t=\inf_c\|\Phi_t-c\|_\infty$, $\Phi_t$ its conditional mean at the level $t$ on the small-field set of that level, at the source $0$ on $T_{\ell'}$.

(i) \emph{Null terms.} For a field of directions $\mathbf X$ that vanishes at the bonds outside a finite set $\Lambda$ (Section~\ref{sec:ward-general-calculus}), the function $\mathsf N=\beta\,\mathfrak v_{\mathbf X}A-\mathrm{div}\,\mathfrak v_{\mathbf X}$, the derivative of the action along the vector field $\mathfrak v_{\mathbf X}$ on $\mathrm{SU}(2)^\Lambda$ minus its divergence, has $\nu(\mathsf NY)=0$ for every bounded Borel $Y$ that does not depend on the variables of the bonds of $\Lambda$, in every Gibbs state and every torus measure: integration by parts on the Haar measure (Lemma~\ref{lem:ward-general-calculus}). (ii) \emph{The identity.} Identically in the field, by the derivative of a plaquette term (Lemma~\ref{lem:ward-general-calculus}) and a summation by parts, $\mathsf{St}_K[\eta]=\mathfrak v_\eta A(\beta,\cdot)-\mathfrak K^{\rm v}[\beta\eta]$; the divergence of the displacement derivation is $\mathrm{div}\,\mathfrak v_\eta=-\frac32a_K\mathcal O_K(w_\eta)$, $w_\eta(p)=[\nabla^{\rm c}_\mu\eta^\mu+\nabla^{\rm c}_\nu\eta^\nu](x(p))$, $\mu<\nu$ the directions of $p$ (Lemma~\ref{lem:ward-general-calculus}); and the orbital insertion is $\mathfrak K^{\rm v}[\beta\eta]=\mathcal S[b\eta]-\beta\sum_{p,\rho}(\nabla^+_\rho\eta^\rho)(p)\Psi_\rho(p)$ (Proposition~\ref{prop:ward-orbital-cubic}(c)), with the orbital potentials $\Psi_\rho$ and the cubic term $\mathcal S[b\eta]=\beta\mathfrak K^\sharp[\eta]=\sum_pL^{2K}b\,\eta^{\rho(p)}(x(p))\mathfrak r^\sharp_{\rho(p)}(p)$ of the remainders $\mathfrak r^\sharp$, which are of the third order in the field strength (Section~\ref{sec:ward-orbital-complement}, Proposition~\ref{prop:ward-orbital-cubic}(a)).

The term with the orbital potentials is, up to a null term of the kind of (i), whose field of directions moves each bond $b$ along $\sum_\rho(\nabla^+_\rho\eta^\rho)(x_b)$ times the sum of the vector parts of its two plaquettes of the direction $\rho$, and up to a constant, a sum of entries of order two with weights at most $Cb\,a_K\|\eta\|_{C^2}$ and of entries of the orders three and four with weights at most $CL^Kb\|\eta\|_{C^2}$. So $\mathsf{St}_K[\eta]$ is the sum of null terms of the kind of (i), of a constant, and of $-\mathsf{Rm}_K[\eta]$, where the remainder $\mathsf{Rm}_K[\eta]$ is $\mathcal S[b\eta]$ plus these entries of the orders two to four plus $\frac32a_K\mathcal O_K(w_\eta)$; in every joint cumulant whose other entries do not depend on the bonds within $2a_K$ of $\mathrm{supp}\,\eta$, $\mathsf{St}_K[\eta]$ may be replaced by $-\mathsf{Rm}_K[\eta]$. (iii) \emph{The collapse.} $\mathsf{St}_K[\eta]-b\,O^T_K(H^{\rm cons}[\eta])$ is the sum of a null entry of order two, with weights at most $22\,b\sup|\nabla\eta|$, and of an entry of order two with weights at most $60\,b\,a_K\sup|\nabla^2\eta|$, both carried within $2a_K$ of $\mathrm{supp}\,\eta$: to the first order in $\nabla\eta$ the symbol of $\mathsf{St}_K[\eta]$ is $b\,H^{\rm cons}[\eta]$, and this is the lattice form of the identity $\partial_\nu(2T_{\nu\rho}-\delta_{\nu\rho}\mathrm{tr}\,T)=0$ for $T_{\lambda\lambda'}=\langle B_\lambda,B_{\lambda'}\rangle$. (iv) \emph{Oscillations.} On $T_{\ell'}$ at the source $0$, the joint cumulants of entries whose carriers are at least $d_\times$ apart are bounded by the products of their oscillations at the level $t$, up to terms exponentially small in $d_\times L^{K-t}$ and the weight of the large fields: the chain with one entry put in as a source gives the cumulants of several entries by polarisation, its families being local in the entry.

An entry of order two with the weights $v$ has $\omega_t\le CK^{2\mathfrak p_0}L^{4\lceil\phi K\rceil}\|v\|$, and an entry of an order $\ge3$ with the weights $\varkappa$ has $\omega_t\le CK^{3\mathfrak p_0}L^{4\lceil\phi K\rceil-\frac32K}\sup|\varkappa|$. A null entry is much smaller at the regular configurations of the level $n_\Delta=K-\lceil K/14\rceil$, where the null entry of (iii) is at most $CK^{2\mathfrak p_0}L^{-K/4}b\sup|\nabla\eta|$, and it has $\omega_t\le CK^pL^{3\lceil\phi K\rceil-K/28}\|v\|$: what is born below the level $n_\Delta$ is born at the coupling $x_j^{-1/2}\le CL^{-K/28}$, $j\le n_\Delta$, and is counted over the $L^{3\lceil\phi K\rceil}$ cubes of the level $t$ in a cell; above $n_\Delta$ nothing is born, a conditional mean being not larger than a supremum over the regular configurations. (v) \emph{The root.} For all data at $g$, whose labels lie in the compact interval $[\gamma^{-2},X_g+\bar B]$, and $K\ge K_0$, the conditional mean at the level $t$ of the cubic term $\mathcal S[b\eta]$ is almost everywhere on the small-field set of that level a constant up to $CK^pL^{-c_0'K}\|\eta\|_{C^2}$, $c_0'=\min\{\frac1{14}-\frac\tau7,\frac1{14}-3\phi\}$, with constants allowed to depend on the fractions, for $0<\tau\le\frac14$ and $0<\phi<\frac1{42}$; at $\tau=\frac14$, $\phi=\frac1{400}$, $c_0'=\frac1{28}$.

The cubic term has no part of order two, its remainders being of the third order (Proposition~\ref{prop:ward-orbital-cubic}(a)). Its families are born along the chain and stopped at the level $K-\lceil K/7\rceil$: what is born below that level is born at a small coupling and is carried upward through chains of steps, which a stopping rule with the free fraction $\tau$ cuts, at the cost $L^{-(\frac12-\tau)\lceil K/7\rceil}$, while a chain that reaches the level $t$ costs $L^{3\lceil\phi K\rceil-\lceil K/7\rceil/2}$; at the stopping level the cubic term costs its supremum over the regular configurations; the other exponents that enter are larger, and $c_0'$ is the least. \emph{The passage.} The eighth power of the seminorm on $T_{\ell'}$ is the mean of the product of the eight copies in the cells of a period box; centred at the mean, it is a sum over the partitions of products of joint cumulants of the copies, which are $\ge2\delta$ apart.

The copies of a null term are null terms, and, the walls lying along faces of the blocks of the level $t$ by the offset in $a_{K-t}\mathbb Z^3$ (Proposition~\ref{bil:press}(c) at $h=a_{K-t}$), the copies have the mirrored conditional means and the same oscillations; so the seminorm is at most a number times the oscillation at the level $t$, plus the eighth root of the remaining terms. For $\mathsf{St}_K[\eta]$, (ii) removes the null terms and the constant; (v) bounds the oscillation of the cubic term by $CK^pL^{-K/28}\|\eta\|_{C^2}$; and by (iv) the entries of order two of $\mathsf{Rm}_K[\eta]$ have $\omega_t\le CK^{2\mathfrak p_0}L^{4\lceil\phi K\rceil-K}b\|\eta\|_{C^2}$, those of the orders three and four $\omega_t\le CK^{3\mathfrak p_0}L^{4\lceil\phi K\rceil-K/2}b\|\eta\|_{C^2}$, and $\frac32a_K\mathcal O_K(w_\eta)$ has $\omega_t\le CK^{2\mathfrak p_0}L^{4\lceil\phi K\rceil-K}\|\eta\|_{C^2}$: the rate $\frac1{28}$.

For the difference, (iii) and (iv) give $\omega_t\le CK^pL^{3\lceil\phi K\rceil-K/28}b\sup|\nabla\eta|+CK^{2\mathfrak p_0}L^{4\lceil\phi K\rceil-K}b\sup|\nabla^2\eta|$, and $\frac3{400}-\frac1{28}=-\frac{79}{2800}<-\frac1{36}$: the rate $\frac1{36}$. As $b^{-1}\le\gamma^2$, $O^T_K(H^{\rm cons}[\eta])=b^{-1}\bigl(\mathsf{St}_K[\eta]-(\mathsf{St}_K[\eta]-b\,O^T_K(H^{\rm cons}[\eta]))\bigr)$ has the rate $\frac1{36}$.
\end{proof}

\begin{proposition}[Rotations]\label{bil:rot}
Let $g\le\gamma_1$ and let $\sigma$ be a limit through periodic Gibbs states, with the label $\mathsf b$.

(a) On the pairs with disjoint supports, $\sigma(H,H')=\iint\mathcal G(y-x)(H(x),H'(y))\,dx\,dy$ with a real-analytic map $\mathcal G$ from $\mathbb R^3\setminus\{0\}$ to the bilinear forms on $V_{\mathbb C}$; $\mathcal G(u)(a,b)=\mathcal G(-u)(b,a)$, $\mathcal G(wu)(\rho(w)a,\rho(w)b)=\mathcal G(u)(a,b)$ for $w\in W_3$, $\rho(w)a=waw^{\rm T}$; $\mathcal G$ is real at real labels; and $\mathcal G(\cdot)(a,b)$ continues holomorphically to the polydisc of radius $|u|/(6\sqrt3)$ about $u$, bounded there by $\Gamma(|u|/2\sqrt3)$, $\Gamma(t)$ the largest of the numbers $\int e^{-tE}d\rho_{cc}$, $c=\rho(w)a,\rho(w)b$, $w\in W_3$, with the spectral measures $\rho_{cc}$ of the label $c$. And $\mathcal G(Ru)=(\rho(R)\otimes\rho(R))\mathcal G(u)$ for all $R\in O(3)$: $\sigma(RH,RH')=\sigma(H,H')$.

(b) $\sigma(RX)=\sigma(X)$ for all $R\in O(3)$ and all separated tuples $X$ with continuous entries. At the tuples with at most two entries outside $C_c(\mathbb R^3;\mathbb R\mathbf1)$ the proof uses neither the insertion bound nor Lemma~\ref{bil:vec}.
\end{proposition}

\begin{proof}[Proof sketch]
\emph{1. The kernel at two points.} A test function $\chi_d\ge0$ with support in a ball of radius $\sqrt3d$ about a point of height $3d$, whose Laplace--Fourier transform has no zero on the spectrum, defines through the spectral measure $\mathcal P$ a positive matrix of measures $\rho_{ab}$, antilinear in $a$ and linear in $b$, with $\langle[(ha)],\mathcal P(\cdot)[(h'b)]\rangle=\overline{\tilde h}\tilde h'\rho_{ab}$; $\Gamma_a(t)=\int e^{-tE}d\rho_{aa}$ is finite, $\mathcal G^+_{ab}(\mathfrak z)=\int e^{-\mathfrak z_3E+i(\mathfrak z_1q_1+\mathfrak z_2q_2)}d\rho_{ab}$ is holomorphic on $\mathsf{Tb}$ with $|\mathcal G^+_{ab}(\mathfrak z)|\le[\Gamma_a\Gamma_b]^{1/2}(\mathrm{Re}\,\mathfrak z_3-|\mathrm{Im}\,\mathfrak z_1|-|\mathrm{Im}\,\mathfrak z_2|)$, and $\mathcal G(u)(a,b)=\mathcal G^+_{a^\dagger b}(u)$ on $\{u_3>0\}$, $a^\dagger=\rho(\theta)\bar a$, extended by $W_3$ to the six half spaces that cover $\mathbb R^3\setminus\{0\}$; this gives (a) but for the last sentence.

\emph{2. Harmonics.} Let $R_t$ be the rotation by the angle $t$ about $e_2$, $R_te_3=\cos t\,e_3+\sin t\,e_1$, $R_te_1=-\sin t\,e_3+\cos t\,e_1$, and $u_\pm=e_3\pm ie_1$, so that $R_tu_\pm=e^{\mp it}u_\pm$; the labels $\varepsilon_{\mp2}=u_\pm u_\pm^{\rm T}$, $\varepsilon_{\mp1}=u_\pm e_2^{\rm T}+e_2u_\pm^{\rm T}$, $\varepsilon_0=e_3e_3^{\rm T}+e_1e_1^{\rm T}$, $\varepsilon'_0=e_2e_2^{\rm T}$ form a basis of $V_{\mathbb C}$ with $\rho(R_t)\varepsilon_k=e^{ikt}\varepsilon_k$. For a label $a$ of this basis $k(a)$ is its index, $k(\varepsilon'_0)=0$, and $-k(a)$ is its \emph{helicity}; as $\theta\bar u_\pm=-u_\pm$, $a^\dagger=(-1)^{k(a)}a$. For a label $a$ and a real $\varphi\in C_c(\mathbb R)$ on the axis of $e_2$, $\varphi a$ is the label smeared along the axis, and the function $\Gamma_a$ of step~1 at it is
\[
\begin{multlined}
\Gamma_{\varphi a}(t)=\int e^{-tE}|\hat\varphi(q_2)|^2d\rho_{aa}=\iint\varphi(s)\varphi(s')\,\mathcal G\bigl((0,s'-s,t)\bigr)(a^\dagger,a)\,ds\,ds'\quad(t>0),\qquad\hat\varphi(q_2) \\
=\int\varphi(s)e^{iq_2s}ds .
\end{multlined}
\]
For two labels $a$, $b$ of the basis, $\kappa=k(a)+k(b)$, real $\varphi,\psi\in C_c(\mathbb R)$ on the axis of $e_2$ and $r>0$, with $n_\zeta=(\cos\zeta,\sin\zeta)$ in the plane of $(x_3,x_1)$ and $(rn_\zeta,\tau)=r\cos\zeta\,e_3+r\sin\zeta\,e_1+\tau e_2$, the function $K(\zeta)=\iint\varphi(s)\psi(s')\,\mathcal G((rn_\zeta,s'-s))(a,b)\,ds\,ds'$ is holomorphic on $|\mathrm{Re}\,\zeta|<\frac\pi2$; the quarter turn of $W_3$ continues it to an entire function with $K(\zeta+\frac\pi2)=e^{-i\kappa\pi/2}K(\zeta)$, so that $K(\zeta)=\sum_m\ell_me^{i(4m-\kappa)\zeta}$, $\ell_m=\ell_m(r;\varphi a,\psi b)$, and $|K(\vartheta+is)|\le[\Gamma_{\varphi a}\Gamma_{\psi b}]^{1/2}(re^{-|s|}/\sqrt2)$; if $\Gamma_{\varphi a}(t)\le C\max\{1,t^{-1}\}^{A_a}$ and likewise for $\psi b$, only the $m$ with $|4m-\kappa|\le\frac12(A_a+A_b)$ occur. Covariance under the $R_t$ is $\ell_m=0$ for $m\ne0$.

\emph{3. Positivity.} Positivity at $\{x_3=0\}$, and at the diagonal plane with normal $(e_1+e_3)/\sqrt2$ through Bochner's theorem and the analytic continuation of characteristic functions, make the Laurent polynomials $Q(x)=(\sum_m\ell_m(r;\varphi_ia_i^\dagger,\varphi_ja_j)x^m)_{ij}$ Hermitian and positive semi-definite for every real $x\ne0$; so their extreme exponents are even and the off-diagonal entries are bounded by the diagonal ones.

\emph{4. Growth.} For an off-diagonal label $o$, the mirror bound (F9) applied to $\chi_do$ gives $\Gamma_o(t)\le\mathsf c_2C^{\rm od}_\infty(g)\mathsf b^{-2}\max\{1,t^{-1}\}^6$; with (a), $|\mathcal G(u)(o,o')|\le C(1+|u|^{-6})$ for the off-diagonal labels $o_{\mu\nu}=e_\mu e_\nu^{\rm T}+e_\nu e_\mu^{\rm T}$, $\mu\ne\nu$, which $W_3$ permutes up to signs. Conservation (F5) at two points gives, with $\mathcal G(u)(a,b)=\sum_{ij}a_{ij}G_{ij}(u;b)$,
\[
2\,\partial_iG_{ij}(u;b)=\partial_j\,\mathrm{tr}\,G(u;b)\qquad(u\ne0),
\]
and, with Cauchy's estimate on the polydiscs of (a) and the off-diagonal growth, $|\mathcal G(u)(\bar a_\mu,\bar a_\mu)|\le C(1+|u|^{-6})$ for the diagonal labels $\bar a_\mu=2e_\mu e_\mu^{\rm T}-\mathbf1$. Smeared along the axis of $e_2$, these pointwise bounds give one power less: for $c$ among the $o_{\mu\nu}$ and the $\bar a_\mu$, as $c^\dagger=\pm c$ and $\int_{\mathbb R}(t^2+\tau^2)^{-3}d\tau=\frac{3\pi}8t^{-5}$,
\[
\Gamma_{\varphi c}(t)\le\|\varphi\|_\infty^2\lambda_\varphi\int_{|\tau|\le\lambda_\varphi}C\bigl(1+(t^2+\tau^2)^{-3}\bigr)d\tau\le C_\varphi\max\{1,t^{-1}\}^5,
\]
$\lambda_\varphi$ the length of the convex hull of $\mathrm{supp}\,\varphi$; and $\rho_{c+c',c+c'}\le2\rho_{cc}+2\rho_{c'c'}$. The labels with $k=\pm1,\pm2$ are combinations of these, $\varepsilon_{\mp2}=\frac12(\bar a_3-\bar a_1)\pm io_{31}$, $\varepsilon_{\mp1}=o_{32}\pm io_{12}$; the labels with $k=0$ involve the trace label besides, $\varepsilon'_0=\frac12(\mathbf1+\bar a_2)$, $\varepsilon_0=\frac12(\mathbf1-\bar a_2)$, and $\Gamma_{\varphi\mathbf1}$ is the function $\Gamma_\varphi$ of the two-point function of the action density, which grows with the exponent $4+a_\Gamma$ (the proof of Proposition~\YMMref{96.3}{728}), $a_\Gamma<4$ the exponent of its growth (R7) of Definition~\YMMref{95.1}{723}. So $A_a=5$ for $k(a)=\pm1,\pm2$ and $A_a=\max\{5,4+a_\Gamma\}<8$ for $k(a)=0$; what the window of step~2 needs is $A_a<6$ at $|k(a)|=1$ and $A_a<8$ at $k(a)\in\{0,\pm2\}$.

At a diagonal entry of $Q$, $\kappa=2k(a)$: at $k(a)=0$ the window leaves $|m|\le1$, at $k(a)=\pm1$ it leaves $m\in\{0,\pm1\}$, and as the extreme exponents are even only $m=0$ remains; at the label $\varepsilon_{\mp2}$ it leaves $m\in\{0,\mp1,\mp2\}$. The off-diagonal entries being bounded by the diagonal ones near $x=0$ and $x=\infty$: $\ell_m=0$ for $m\ne0$ if $|k(a)|,|k(b)|\le1$; if $k(a)=\mp2$ and $|k(b)|\le1$ (or the two exchanged), $\ell_m=0$ unless $m\in\{0,\mp1\}$; if $k(a)=k(b)=\mp2$, unless $m\in\{0,\mp1,\mp2\}$; if $k(a)=-k(b)=\mp2$, unless $|m|\le1$. So everything is covariant but possibly the part of helicity $\pm2$ of $T_{33}-T_{11}\pm2iT_{31}$, the label $\varepsilon_{\mp2}$, through these harmonics.

\emph{5. Conservation decides at two points.} In the variable $z=u_3+iu_1$, $\partial_z=\frac12(\partial_3-i\partial_1)$, the identity of step~4, contracted with $u_+$ and with $u_-$, takes the form
\[
\partial_z\mathcal G(\varepsilon_{-2},b)=\partial_{\bar z}\mathcal G(\varepsilon'_0,b)-\tfrac12\partial_2\mathcal G(\varepsilon_{-1},b),\qquad\partial_{\bar z}\mathcal G(\varepsilon_{+2},b)=\partial_z\mathcal G(\varepsilon'_0,b)-\tfrac12\partial_2\mathcal G(\varepsilon_{+1},b).
\]
For labels $a$, $b$ of the basis, $\mathcal G((rn_\zeta,\tau))(a,b)=\sum_m\ell_m(r,\tau;a,b)e^{i(4m-\kappa)\zeta}$ is the Fourier series in $\zeta$ of a real-analytic function with the quasi-period of step~2; its coefficients are real-analytic in $(r,\tau)$, the series may be differentiated term by term, and $\ell_m(r;\varphi a,\psi b)=\iint\varphi(s)\psi(s')\ell_m(r,s'-s;a,b)\,ds\,ds'$, so that what step~4 says of the smeared harmonics for all $\varphi,\psi$ holds for the $\ell_m(r,\tau;a,b)$; by (a), $|\ell_m(r,\tau;a,b)|\le\Gamma(\frac1{2\sqrt3})$ for $r\ge1$: they stay bounded as $r\to\infty$. In polar coordinates $\partial_z(fe^{i\nu\zeta})=\frac12e^{i(\nu-1)\zeta}(f'+\frac\nu rf)$ and $\partial_{\bar z}(fe^{i\nu\zeta})=\frac12e^{i(\nu+1)\zeta}(f'-\frac\nu rf)$; so a frequency of the member with the label $\varepsilon_{\mp2}$ that the other member has not gives $\ell=c(\tau)\,r^n$, with $n=-\nu$ in the first identity and $n=\nu$ in the second, $\nu$ the frequency of the harmonic $\ell$ before the derivative, and $c=0$ if $n\ge1$, by the boundedness.

This decides in three stages. First, $|k(b)|\le1$: at $(\varepsilon_{-2},b)$ step~4 leaves the frequencies $2-k(b)$ ($m=0$) and $-2-k(b)$ ($m=-1$), which $\partial_z$ lowers to $1-k(b)$ and $-3-k(b)$, while the right member, formed with pairs of labels with $|k|\le1$, has the one frequency $1-k(b)$; so $\ell_{-1}=c\,r^{2+k(b)}$, $2+k(b)\ge1$, and $\ell_{-1}=0$. In the same way $\ell_1=c\,r^{2-k(b)}=0$ at $(\varepsilon_{+2},b)$, and by $\mathcal G(u)(a,b)=\mathcal G(-u)(b,a)$ the pairs $(\varepsilon_{\mp2},b)$ and $(b,\varepsilon_{\mp2})$ with $|k(b)|\le1$ are covariant. Second, at $(\varepsilon_{-2},\varepsilon_{-2})$, $\kappa=-4$, the left member has the frequencies $3$, $-1$, $-5$ ($m=0,-1,-2$), the right member, by the first stage, the one frequency $3$; so $\ell_{-2}=c\,r^4=0$, and alike $\ell_2=0$ at $(\varepsilon_{+2},\varepsilon_{+2})$.

Third, step~3: at the label $\varphi\varepsilon_{\mp2}$ the diagonal entry of $Q$ is now $\ell_0+\ell_{\mp1}x^{\mp1}$, non-negative at both signs of $x$, and its extreme exponent $\mp1$ is odd unless $\ell_{\mp1}=0$. So all the diagonal entries of $Q$ are constant, and the off-diagonal ones, bounded by them on $\mathbb R\setminus\{0\}$, are constant too: $\ell_m=0$ for $m\ne0$ at every pair, the pair $(\varepsilon_{-2},\varepsilon_{+2})$ and different smearing functions included. Hence $\mathcal G$ is covariant under the $R_t$, with $W_3$ under the rotations about the three axes, which generate $SO(3)$, and $-\mathbf1\in W_3$.

\emph{6. All tuples.} Write $x=(y,x_2)$, $y=(x_3,x_1)$ being the point of the plane of $(x_3,x_1)$, which is normal to the axis of $e_2$. A \emph{plate} is a function $(\alpha\otimes\varphi a)(x)=\alpha(y)\varphi(x_2)a$ with a \emph{profile} $\alpha\in C^1_c(\mathbb R^2;\mathbb C)$ in that plane and a \emph{helicity label} $\varphi a$, $\varphi\in C^1_c(\mathbb R)$ real and $a$ a label of the basis of step~2; the \emph{planar projection} of a function on $\mathbb R^3$ is the set of the points $y$ of its support . For $n$ helicity labels $\underline a=(\varphi_ja_j)_j$ let $S=-\sum_jk(a_j)$ be their total helicity; $R_t(\alpha\otimes\varphi a)=e^{ik(a)t}(\alpha\circ R_{-t})\otimes\varphi a$, and $R_\zeta$ acts on configurations $Y=(y_1,\dots,y_n)$ of points of the plane, continued to complex $\zeta$.

At plates whose profiles are ordered in $x_3$, Lemma~\ref{bil:vec}(b), smeared along the axis, writes $\sigma$ as the integral of the profiles against a kernel $\mathcal K_n(Y;\underline a)$, holomorphic on the tube of that order; the quarter turns of $W_3$ carry the order in $x_3$ to the orders along $\pm e_1$ and $-e_3$, the labels taking the phases of step~2, and the kernels of all the orders form one holomorphic function $\mathcal K_n(\cdot;\underline a)$ on the union $\mathcal Z_n$ of the tubes, with $\mathcal K_n(R_{\pi/2}Z;\underline a)=e^{iS\pi/2}\mathcal K_n(Z;\underline a)$ (Lemma~\YMMref{97.5}{733}, with this phase); its real points are the configurations whose $x_3$-coordinates, or whose $x_1$-coordinates, are pairwise different.

At a configuration $Y$ whose consecutive differences make angles $<\frac\pi8$ with $e_3$, $R_\zeta Y\in\mathcal Z_n$ for every complex $\zeta$, every real angle lying within $\frac\pi4$ of a multiple of $\frac\pi2$; so $h_Y(\zeta)=e^{-iS\zeta}\mathcal K_n(R_\zeta Y;\underline a)$ is entire, of period $\frac\pi2$, of growth $e^{(p+2)n|\mathrm{Im}\,\zeta|}$ by the bound of Lemma~\ref{bil:vec}(b) and $|S|\le2n$, and $h_Y(\zeta)=\sum_{4|m|\le(p+2)n}c_m(Y)e^{4im\zeta}$, that is $\mathcal K_n(R_{i\tau}Y;\underline a)=e^{-S\tau}\sum_mc_m(Y)e^{-4m\tau}$. As the frequencies are bounded uniformly, $c_m$ is one fixed linear combination of the values $e^{Sl}\mathcal K_n(R_{il}Z;\underline a)$, $l=0,1,\dots$, with the coefficients of an inverse Vandermonde matrix, so a holomorphic function on $\mathcal Z_n$, and the expansion holds on all of $\mathcal Z_n$ by the identity theorem on each tube (Lemma~\YMMref{98.1}{734}).

The phases are exact, $c_m(R_tY)=e^{i(S+4m)t}c_m(Y)$: first at the configurations without two orthogonal differences of points, a dense set that all the rotations keep in $\mathcal Z_n$, then by continuity; so $\tilde c_m(Y)=e^{-i(S+4m)t}c_m(R_tY)$, for any $t$ with $R_tY\in\mathcal Z_n$, is a real-analytic function on all the configurations of pairwise different points, with the same phases under the rotations of the plane about any point (Lemma~\YMMref{98.1}{734}); under these maps and the reflections of the plane the points move and the labels stay. Let $\bar m(n)$ be the largest $|m|$ with $c_m\ne0$ for some $n$ helicity labels. The argument of Proposition~\YMMref{98.3}{735} for scalar entries goes through, the labels entering at three places. (i) The mirror image of a helicity label is the label up to a sign, $\rho(\theta)\overline{\varphi\varepsilon_k}=(-1)^k\varphi\varepsilon_k$; so the Gram matrix of the vectors of point insertions at ordered heights, turned by $R_{i\tau}$, is the matrix of the kernels up to the signs $(-1)^S$ and the factors $e^{-S\tau}$, which cancel in the Schwarz inequality: $2\bar m(N_1+N_2)\le\bar m(2N_1)+\bar m(2N_2)$ for integers $N_1,N_2\ge0$ (Lemma~\YMMref{98.2}{735}). (ii) These Gram matrices, multiplied on both sides by the diagonal matrix of the factors $e^{S\tau}$ and by $e^{-4\bar m(2N)\tau}$, tend as $\tau\to\infty$ to the matrix of the top harmonic $\mathsf t_N=\tilde c_{-\bar m(2N)}$ of $2N$ points; so $\mathsf t_N$ is positive semi-definite, up to the signs $(-1)^S$, at the line $\{x_3=0\}$ of the plane, and, with its phases under the rotations, at every line $\mathsf L$ of the plane up to phases: $|\mathsf t_N(\theta_{\mathsf L}\mathsf P\cup\mathsf P')|^2\le|\mathsf t_N(\theta_{\mathsf L}\mathsf P\cup\mathsf P)|\,|\mathsf t_N(\theta_{\mathsf L}\mathsf P'\cup\mathsf P')|$ for labelled configurations $\mathsf P$, $\mathsf P'$ of $N$ points on one side of $\mathsf L$ (Lemma~\YMMref{98.2}{735}, equation~\YMMref{(1947)}{735}). (iii) A turn by $t$ about a point that maps a labelled configuration to itself gives $(1-e^{i(S+4m)t})\tilde c_m=0$ there; at a regular $2N$-gon all of whose points carry one helicity label, of the integer helicity $h$, the turn by $\frac\pi N$ gives the factor $e^{i(2Nh-4\bar m(2N))\pi/N}=e^{-4\pi i\bar m(2N)/N}$ at $m=-\bar m(2N)$.

If $N$ did not divide $2\bar m(2N)$, $\mathsf t_N$ would vanish at every such polygon; the inequality of (ii) at the perpendicular bisectors of its chords, with the identity theorem for the real-analytic $\mathsf t_N$ on the connected set of the configurations of pairwise different points, would carry the vanishing of $\mathsf t_N$ at all the labelled configurations that contain an arc of $j$ consecutive vertices to those that contain an arc of $\lceil\frac j2\rceil$ of them, from $j=2N$ down to $j=1$, so that $\mathsf t_N$ would vanish identically, against the definition of $\bar m(2N)$. So $N$ divides $2\bar m(2N)$ (the proof of Proposition~\YMMref{98.3}{735}). As $\bar m(2N)\le\frac{p+2}2N$, the convex-integer-sequence argument at the end of the proof of Proposition~\YMMref{98.3}{735} gives $\bar m(2N)=N\,\bar m(2)$, and $\bar m(n)\le\frac n2\bar m(2)$.

By step~5, $\bar m(2)=0$; so $\bar m(n)=0$, and the kernel $\tilde{\mathcal K}_n=\tilde c_0$ has $e^{-iSt}\tilde{\mathcal K}_n(R_tY;\underline a)=\tilde{\mathcal K}_n(Y;\underline a)$ at all the configurations of pairwise different points. At plates whose profiles have pairwise disjoint supports $\sigma$ is the integral of the profiles against $\tilde{\mathcal K}_n$, with no part at coinciding planar coordinates: by induction on the number of the pairs of profiles whose ranges in $x_1$ overlap, the group of profiles above a level in $x_3$ being slid along $e_1$, both members are real-analytic in the slide and agree once the overlap is gone, and small pieces of any profiles have such levels (the proof of Proposition~\YMMref{98.4}{736}).

So $\sigma(R_tX)=\sigma(X)$ at the tuples of plates whose profiles have pairwise disjoint supports, the phases $e^{ik(a_j)t}$ of the labels compensating $e^{-iSt}$; Riemann sums of translated small plates with constant matrix values, and mollification, with the bound (D1), give it at the tuples with $C^1$ entries whose planar projections are pairwise disjoint (Proposition~\YMMref{98.4}{736}), and the slides of Proposition~\YMMref{95.4}{725}, in the picture of $e_2$, give it at all separated tuples with $C^1$ entries; these proofs read the entries through multilinearity, the translations, the positivity, $H\ge|P_k|$ and the bound (D1) alone. $O(3)$ follows as in step~5, and continuous entries by mollification with radial functions, which commutes with the rotations.

\emph{7. At most two tensor entries.} At tuples with at most two entries outside the scalar part, the argument of Theorem~\YMMref{AK}{737} for the action density, with the covariance at two points of step~5 for the pair of tensor entries, gives the invariance without the insertion bound.
\end{proof}

\begin{proposition}[Mirrors and the cone]\label{bil:cone}
Let $\sigma$ be a real symmetric functional on the separated tuples with continuous entries, multilinear and bounded as in (D1), and invariant under the translations.

(a) If $\sigma$ is reflection positive at a plane $\{n\cdot x=c\}$, the labels turned by $\Theta_n=1-2nn^{\rm T}$, then $\sigma$ is invariant under the reflection in that plane, and under the reflections in all the parallel planes.

(b) A closed subgroup of $O(3)$ that contains $W_3$ and one element more is $O(3)$; and $\Theta_n\in W_3$ if and only if $n$ is one of the eighteen normals of the nine families. So, if $\sigma$ is besides invariant under $W_3$ and reflection positive at the nine families, the following are equivalent: $\sigma$ is invariant under $O(3)$; it is invariant under one $R\notin W_3$; it is reflection positive at one plane whose normal is not one of the eighteen; it is reflection positive at every plane.

(c) Let $(u,u',u'')$ be an orthonormal basis, $n_\beta=\cos\beta\,u+\sin\beta\,u'$, and let $\sigma$ be reflection positive at the four planes through the axis $\mathbb Ru''$ with the normals $n_0,n_{\pi/4},n_{\pi/2},n_{3\pi/4}$, and at one plane with the normal $(n_{\beta_1}+u'')/\sqrt2$, $\beta_1\in\frac\pi4\mathbb Z$. Then in each of the pictures $n_\beta$, $\beta\in\frac\pi4\mathbb Z$, the joint spectral measure of energy and momenta is carried by $\{E\ge|q|\}$, $|q|$ the Euclidean length. No invariance under $W_3$, under a reflection or under a rotation is supposed.

(d) If $\sigma$ is invariant under $O(3)$ and reflection positive at one plane, it is reflection positive at every plane $\{n\cdot x=c\}$, with $(\theta_{n,c}H)(x)=\Theta_nH(x-2(n\cdot x-c)n)\Theta_n$.
\end{proposition}

\begin{proof}[Proof sketch]
(a) Positivity makes the form $\langle F',F\rangle=\sigma(\theta\bar F'\cdot F)$ on the combinations of tuples on one side Hermitian; at $F'=\emptyset$ and a real tuple $X$ on that side this says $\sigma(X)=\sigma(\theta X)$. Every tuple lies on one side of a parallel plane far away, and the reflections in parallel planes differ by translations. (b) The identity component of such a subgroup is normalised by $W_3$; it is not trivial, since no finite subgroup of $O(3)$ contains the octahedral group $W_3$ properly, and it is no group of rotations about one axis, since $W_3$ fixes no axis; so it is $SO(3)$, and $-\mathbf1\in W_3$. A mirror whose normal is not one of the eighteen gives by (a) an element $\Theta_n\notin W_3$ of the group of invariance; conversely (d). (c) Positivity at two planes at the angle $\frac\pi4$ bounds, in the picture of the one, the energy by the momentum along the other normal (Proposition~\YMMref{94.4}{722}); so the five planes give $E\ge|q''|$, the momentum along $u''$, in the picture $n_{\beta_1}$, and the bounds along $u'$ in the eight pictures in the plane of $u$, $u'$.

In these eight pictures the Laplace transforms of the spectral measures, at fixed $q''$, are restrictions of one holomorphic function on a band, because by the symmetry of $\sigma$ one and the same number $\sigma(A\cdot\tau_xB)$ is a matrix element in every picture whose mirror parts $A$ from $\tau_xB$, also when $x$ has moved about $A$ to the other side; the comparison of Lemma~\ref{C-355-26}(b,c) carries the bound of one picture to $E^2\ge q'^2+q''^2$ in all of them. (d) $\theta_{n,c}=\mathsf g\theta_m\mathsf g^{-1}$ with the rigid motion $\mathsf gx=Rx+cn$, $Rm=n$, and $\Theta_n=R\Theta_mR^{\rm T}$.
\end{proof}

\begin{proposition}[What the properties of the limits do not decide]\label{bil:zeta}
Let $\sigma$ be a limit through periodic Gibbs states, $g\le\gamma_1$, and for real $\zeta$ let $\Lambda_\zeta$ be the linear map of $V$ that is $1$ on the diagonal and $\zeta$ on the off-diagonal matrices, $\sigma^\zeta(X)=\sigma(\Lambda_\zeta X)$.

(a) $\sigma^\zeta$ is a real symmetric functional, bounded as in (D1) with the factor $\max\{1,|\zeta|\}^n$, invariant under translations and $W_3$, reflection positive at the nine families; $[F]\mapsto[\Lambda_\zeta F]$ is an isometry of its Hilbert space into that of $\sigma$ that intertwines the semigroups, the translations and the vacua; so $\sigma^\zeta$ has the cone of (F4), the insertion bound of (F2) with the factor $\max\{1,|\zeta|\}$ for each entry, and the plateau of (F6).

(b) For $\zeta\ne1$, $\sigma^\zeta$ is not invariant under the rotations about $e_3$; it is reflection positive at no plane outside the nine families; and it does not obey the free law of (F1) at any label.
\end{proposition}

\begin{proof}[Proof sketch]
(a) $\Lambda_\zeta$ commutes with $\rho(w)$ for $w\in W_3$, since a signed permutation carries diagonal matrices to diagonal and off-diagonal to off-diagonal ones, and the mirrors of the nine families act on the labels by elements of $W_3$; the rest follows from the isometry. (b) Let $D=E_{11}-E_{22}$ and $O=E_{12}$, the symmetric matrix with the entries $1$ at $(1,2)$ and $(2,1)$, so that $\rho(R_t)D=\cos2t\,D+\sin2t\,O$, $\rho(R_t)O=-\sin2t\,D+\cos2t\,O$ for the rotations $R_t$ about $e_3$, and let $C=E_{33}$. If $\sigma$ and $\sigma^\zeta$ were both invariant under the $R_t$, the mixed kernel would obey $(\zeta-1)\sin2t\;\mathcal G(R_tu)(O,C)=0$, so $\mathcal G(\cdot)(O,C)\equiv0$; but by (F1) at short distance $\mathcal G(u)(O,C)$ is close to $\mathsf b^{-2}$ times the mixed free kernel $3K_3(u)_{13}K_3(u)_{23}$, which does not vanish identically.

By Proposition~\ref{bil:cone}(b) $\sigma^\zeta$ is then positive at no further plane. If $\sigma^\zeta$ obeyed the free law at a label $\mathsf b'$, the law at two diagonal entries, where $\sigma^\zeta=\sigma$, would give $\mathsf b'=\mathsf b$; and the mixed law depends on $\zeta$ with its sign, $\sigma^\zeta(H_1,H_2)=\zeta\sigma(H_1,H_2)$ for $H_1$ off-diagonal and $H_2$ diagonal, while the mixed free covariance does not vanish identically; so $\zeta=1$. In the proof of Proposition~\ref{bil:rot} every input but conservation holds for $\sigma^\zeta$: it is conservation that excludes it at all distances, and the mixed free law at short distance.
\end{proof}

\begin{lemma}[Extension over coinciding points]\label{bil:ext}
Let $\mathsf M=\mathbb R^3$ or $T_\ell$, $n\ge2$, $A,h_0>0$, $q\ge0$, and let $u$ be a functional of $n$ entries, $n$-linear on $\prod_aC^\infty_c(O_a;V_{\mathbb C})$ for open $O_a$ with pairwise disjoint closures, with the bound
\[
|u(f_1,\dots,f_n)|\le A\,\bigl(1+\log\tfrac{h_0}h\bigr)^q\prod_a\sup|f_a|
\]
for $0<h\le h_0$ and all $f_a$ supported in closed cubes of side $\frac32h$, wherever they lie, at pairwise sup distances $\ge\frac52h$; on $T_\ell$ the number $h_0$ lies in $\ell2^{-\mathbb N}$ and $h_0\le\ell/(64n)$.

(a) There is a distribution $\tilde u$ on $\mathsf M^n$ of order $\le3n-2$ that agrees with $u$ on separated tensor products and has
\[
|\tilde u(F)|\le C_{\rm ex}A\bigl(1+\log\tfrac{h_0}\lambda\bigr)^{q+1}\max_{|\gamma|\le3n-2}\lambda^{|\gamma|}\sup|\partial^\gamma F|
\]
for every $F$ supported in a product of closed sup-balls of radius $\lambda\le h_0$ about any points, $C_{\rm ex}$ depending only on $n$ and $q$, with the label norm and auxiliary smooth cutoffs fixed; if $u$ is invariant under permutations of the entries, translations and a closed subgroup of $O(3)$ that maps $\mathsf M$ to itself, $\tilde u$ can be chosen invariant.

(b) Two extensions with this bound differ by $\sum_{|\pi|\le n-1}\sum_{|\alpha|\le3|\pi|}\partial^\alpha_\xi\delta_\pi\otimes T_{\pi,\alpha}$, $\pi$ the partitions of the entries, $\delta_\pi$ the diagonal of $\pi$, $\xi$ the differences of the points in the blocks, and $T_{\pi,\alpha}$ distributions on $\mathsf M^{|\pi|}$ with the bound of (a) at the power $\lambda^{|\alpha|}$; and every such sum can be added. A distribution with the bound of (a) at a power $\lambda^{\kappa'}$, $\kappa'>3k$, on $\mathsf M^k$ is $0$.

(c) At $n=2$ the translation invariant terms of (b) are $\sum_{|\alpha|\le3}\int\mathsf c_\alpha(\partial^\alpha f_1(x),f_2(x))dx$ with bilinear forms $\mathsf c_\alpha$; the real ones that are symmetric in the two entries and covariant under $O(3)$ on $\mathbb R^3$ form a space of dimension six, spanned in Fourier space by $\mathrm{tr}\,H\,\mathrm{tr}\,H'$, $\mathrm{tr}\,HH'$, $p^2\mathrm{tr}\,H\,\mathrm{tr}\,H'$, $p^2\mathrm{tr}\,HH'$, $(p\cdot Hp)\mathrm{tr}\,H'+(p\cdot H'p)\mathrm{tr}\,H$, $(Hp)\cdot(H'p)$; covariant under the cubic group on $T_\ell$, they form a space of dimension twelve.
\end{lemma}

\begin{proof}[Proof sketch]
(a) Apply Lemma~\ref{C-355-34}(a)--(d) at $\kappa=r=0$, with the stated torus convention. Its construction reads as follows. $\mathsf M^n$ minus the big diagonal is the union of the admissible products of dyadic cubes of sides $h_j=2^{-j}h_0$, pairwise at sup distances $\ge3h_j$, that lie in no larger admissible product (on $T_\ell$ the dyadic cubes tile the torus, as $h_0\in\ell2^{-\mathbb N}$, and each cluster lies in a sup ball of radius $13nh_0<\frac\ell4$, which lifts to $\mathbb R^3$ and defines the differences within that cluster); there is a partition of unity $\varphi_{\mathbf Q}$ with $|\partial^\gamma\varphi_{\mathbf Q}|\le c_\gamma h(\mathbf Q)^{-|\gamma|}$, and on each product $u$ extends to the smooth functions with the bound, by an absolutely convergent expansion in products.

The cubes of a product at mutual distances $\le12h$ form its clusters, a partition $\pi(\mathbf Q)$ with $k$ blocks, where $k\le n-1$ below the coarsest side $h_0$; put $\tilde u(F)=\sum_{\mathbf Q}u_{\mathbf Q}(\varphi_{\mathbf Q}(F-P_{\mathbf Q}F))$, $P_{\mathbf Q}F$ the Taylor polynomial of order $3k$ averaged over the piece of size $h$ of the stratum $\Delta_{\pi(\mathbf Q)}$ next to $\mathbf Q$ when $k\le n-1$, and $P_{\mathbf Q}=0$ otherwise. On the products with $k\le n-1$ the remainder is of order $(h/\lambda)^{3k+1}$, the number of products of side $h$ near a given point is of order $(\lambda/h)^{3k}$, and the sum over the scales converges with one more logarithm; the order of the extension is $3(n-1)+1=3n-2$. Choose the grid partition and the Taylor averages covariantly under $h_0\mathbb Z^3$. Average the resulting extension over the parameters $t\in[0,h_0)^3$ with normalized Lebesgue measure; periodicity in $t$ makes this mean invariant under every translation. Then average over permutations and the compact orthogonal subgroup, as in Lemma~\ref{C-355-34}(d). Its support-radius comparison and the fixed-scale localization in part~(c) preserve the stated bound, with a constant depending on $n,q$. (b) A distribution supported on the big diagonal with the bound of (a) is, stratum by stratum, a sum of transversal derivatives of order $\le3|\pi|$ of $\delta_\pi$; the vanishing for $\kappa'>3k$ is the scaling of the bound. (c) A translation invariant contact term at $n=2$ is a polynomial of order $\le3$ in the relative momentum $p$; $-\mathbf1\in O(3)$ acts trivially on $V$ and changes the sign of $p$, so only the orders $0$ and $2$ remain, and the invariant theory of $O(3)$ on symmetric matrices gives two forms at the order $0$ and four at the order $2$; under the cubic group $3+9$.
\end{proof}

\begin{lemma}[Dilation]\label{bil:dil}
Let $z\ge1$ be an integer and $H^{(z)}(x)=H(L^{-z}x)$.

(a) For $K\ge z$, $O^T_K(H)(U)=L^zO^T_{K-z}(H^{(z)})(U^{(z)})$, $U^{(z)}$ the image of $U$ under $x\mapsto L^zx$; so 
\[
\langle O^T_K(H_1);\dots;O^T_K(H_n)\rangle_{K,\ell,\beta}=L^{nz}\langle O^T_{K-z}(H^{(z)}_1);\dots;O^T_{K-z}(H^{(z)}_n)\rangle_{K-z,L^z\ell,\beta};
\]
periodic Gibbs states at $(K,\beta)$ go to periodic Gibbs states at $(K-z,\beta)$, boards and seminorms are carried along, the dilation commutes with rotations about the origin, carries $\tau_v$ to $\tau_{L^zv}$, and carries the mirror $\{n\cdot x=c\}$ to $\{n\cdot x=L^zc\}$, and $\mathcal W(H^{(z)}_1,H^{(z)}_2)=\mathcal W(H_1,H_2)$.

(b) Let $\gamma_1<g\le\gamma_2$, $t_1=L^zg^{-2}-\bar B$, $t_2=L^{z+1}g^{-2}-3\bar B$, $g'_s=t_s^{-1/2}$. For every datum $(K,\beta)$ at $g$ with $K\ge z+1$, $(K-z,\beta)$ is a datum at $g'_1$ or at $g'_2$; $I(g'_1)\cup I(g'_2)$ contains $L^zb$ for every label $b$ at $g$; $L^z\ell\in\mathfrak L(g'_s)$ for $\ell\in\mathfrak L(g)$; and $g'_s\le\gamma_1$ as soon as $(L^z-1)g^{-2}$ exceeds a number of the constants; $z(g)$ is the least such $z$.

(c) Let $0<g\le\gamma$, $\sigma$ a limit of a sequence of periodic Gibbs states at $g$ that converges at all tuples, $\mathsf b$ a limit point of its labels, and $g_z=(L^z\mathsf b/2)^{-1/2}$. For every sufficiently large integer $z$, $L^z\mathsf b$ is an inner point of $I(g_z)$; along a subsequence the dilated states at $(K_j-z,\beta_j)$ form a convergent sequence at $g_z$ with the limit $\sigma_z$, $\sigma(X)=L^{nz}\sigma_z((H_a^{(z)})_a)$; and $g_z\le\gamma_1$ as soon as $L^z\mathsf b$ exceeds a number of the constants.

(d) Let $0<g\le\gamma$, $d>0$, $q\ge2$, $\Lambda=L^z$, $g'^{-2}_1=\Lambda g^{-2}/2$, $g'^{-2}_2=L\Lambda g^{-2}/4$, and let $2\mathsf l_\iota$ be the least side at $g'_\iota$, $\iota=1,2$. For every datum $(K,\beta)$ at $g$ with $K\ge z+1$, $(K-z,\beta)$ is a datum at $g'_1$ or at $g'_2$. Let $z=z(g,d)$ be the least $z\ge1$ with $g'_1\le\gamma_1$ and $\mathsf l_2\le\Lambda d/2$; then $g'_2<g'_1\le\gamma_1$ and $\mathsf l_1\le\mathsf l_2\le\Lambda d/2$, and we put $\delta_\iota=\mathsf l_\iota/(8q)$.
\end{lemma}

\begin{proof}[Proof sketch]
(a) $a_K^{-1}=L^za_{K-z}^{-1}$, and $a_p$, $q_{\mathsf e}$ do not depend on the spacing; the dilation carries bonds to bonds and plaquettes to plaquettes, the Haar measure and the action along, so $\mu_{K,\ell,\beta}$ to $\mu_{K-z,L^z\ell,\beta}$ (the exact dilation of Lemma~\ref{lem:lattice-law-dilation}); $K_3$ is homogeneous of order $-3$, so the factors $L^{6z}$ of the volume and $L^{-6z}$ of $K_3^2$ cancel in $\mathcal W$. (b) The running couplings of $(K-z,\beta)$ are $x_0,\dots,x_{K-z}$, and by the flow (\YMMref{74.1}{615}) $|L^{-z}x_{K-z}-x_K|\le\bar B$; so the first passage at $g'_1$ or $g'_2$ happens at $K-z$. By Section~\YMMref{74}{614}, $I(g'_1)=(L^zg^{-2},L^{z+1}g^{-2}-2\bar B]$ and $I(g'_2)=(L^{z+1}g^{-2}-2\bar B,L^{z+2}g^{-2}-(2L+2)\bar B]$; their union contains $L^z\mathsf b$ for every $\mathsf b\in(g^{-2},L(g^{-2}+\bar B)]$, so for every label of a datum at $g$ and every $\mathsf b\in I(g)$. (c) $I(g_z)=(g_z^{-2}+\bar B,L(g_z^{-2}+\bar B)-2\bar B]$ holds $L^z\mathsf b=2g_z^{-2}$ in its interior once $L^z\mathsf b$ is large, as $L\ge13$; the dilated data are first-passage data at $g_z$ along a subsequence whose labels converge, and (a) carries the cumulants. (d) Put $x=g^{-2}$.

By Section~\YMMref{74}{614}, $I(g'_1)=(\frac{\Lambda x}2+\bar B,\frac{L\Lambda x}2+(L-2)\bar B]$ and $I(g'_2)=(\frac{L\Lambda x}4+\bar B,\frac{L^2\Lambda x}4+(L-2)\bar B]$; as $\frac{L\Lambda x}4+\bar B\le\frac{L\Lambda x}2+(L-2)\bar B$, their union is the interval $(\frac{\Lambda x}2+\bar B,\frac{L^2\Lambda x}4+(L-2)\bar B]$. It contains $\Lambda I^{\rm lab}(g)=[\Lambda x,L\Lambda(x+\bar B)]$: $\frac{\Lambda x}2+\bar B<\Lambda x$ because $2\bar B<\Lambda x$, and $L\Lambda(x+\bar B)\le\frac{L^2\Lambda x}4+(L-2)\bar B$ because $\bar B\le(\frac L4-1)x$; both follow from $\bar B\le\frac{13}{24}g^{-2}$ (Lemma~\ref{lem:setup-stopped-arithmetic}) and $\Lambda\ge L\ge13$. The label $b=a_K\beta$ of a datum $(K,\beta)$ at $g$ lies in $I^{\rm lab}(g)$ (\YMMref{74.1}{615}); so $L^zb$ lies in $I(g'_1)$ or in $I(g'_2)$, and, as $\beta=L^{K-z}\cdot L^zb$ and $K-z\ge1$, $(K-z,\beta)$ is a first-passage datum at $g'_1$ or at $g'_2$ by the property of the intervals $I(\cdot)$ (Section~\YMMref{74}{614}).

As $z\to\infty$, $g'_1\to0$, while $\mathsf l_2$ grows like a power of $\log(\Lambda g^{-2})$ (equation~\YMMref{(1528)}{615}) and $\Lambda d$ like $\Lambda$; so $z(g,d)$ exists, and $\mathsf l_1\le\mathsf l_2$ because the least side does not decrease with the inverse coupling.
\end{proof}

\begin{proof}[Proof of Theorems~\ref{thm:bil} and~\ref{bil:lim}]
\emph{Step 1: (A1).} Cut $T_\ell$ into at most $(2\ell/d'')^3$ boxes of radius $\le0.86\,d''$, and each form into its pieces on these boxes. A choice of nonzero pieces, one of each form, has the geometry of Lemma~\ref{bil:nsrc} at $d''$, with the majorants of Lemma~\ref{bil:forms}: the pieces lie within $\frac78d''$ of points, and the plaquettes of pieces of different forms are at least $d-2a_K\ge\frac54d''$ apart, as $d\ge\frac43d''$; $s(d'')\ge s_{\rm th}(n)$ since $L^{s(d'')}\ge1/d''$, and $\ell\le\ell_{\max}$ on a small side. Multilinearity gives
\[
|\langle F_1;\dots;F_n\rangle_\ell|\le(2\ell_{\max}/d'')^{3n}k'_n(g,s(d''))\,b^{-n}\prod_a|F_a|_\star .
\]

\emph{Step 2: (A2), (A3).} (A2) is Lemma~\ref{bil:two}(a) with Lemma~\ref{bil:forms}(a): a real form is a real quadratic source under $(|F|_\star,S[F])$ with the kernel $\mathsf t_F$, and its plaquettes lie within $d-a_K$ of $x_i$ when its cells lie within $d-2a_K$; for complex tensor weights the real and imaginary parts give four real pairs, and $O^T_K(H)$ is a source under $(60|H|_\infty,S_H)$. For the continuum form, Lemma~\ref{bil:ker} gives $C^{\rm bg}(p,q)=\epsilon_p\epsilon_qa_K^3K_3(x(p)-x(q))_{\lambda(p)\lambda(q)}+O((a_K/d)^3E_1)$; the sums over plaquettes and bonds are Riemann sums of the integral, with an error $(6/\pi)(a_K/d)$ for each factor at Lipschitz weights; $\upsilon^2$ against $b^{-2}$ costs $8\bar Bg^2L^{2-s}$; and
\[
\begin{aligned}\tfrac32\mathrm{tr}[E_{\lambda_1\lambda_1'}KE_{\lambda_2\lambda_2'}K]&=3[K_{\lambda_1\lambda_2}K_{\lambda_1'\lambda_2'}+K_{\lambda_1\lambda_2'}K_{\lambda_1'\lambda_2}],\quad(\lambda_1\ne\lambda_1',\ \lambda_2\ne\lambda_2'),\\\tfrac32\mathrm{tr}[E_{\lambda\lambda'}KE_{\gamma\gamma}K]&=3K_{\lambda\gamma}K_{\lambda'\gamma}\quad(\lambda\ne\lambda')\end{aligned}
\]
identify the Gaussian value with $b^{-2}\mathcal W$. So, with $s=s(d)$,
\[
e^{\rm od}=\tfrac98\bigl[59^2\varepsilon^\natural_g(s)+600E_2(\tfrac1\pi+E_2)+636/C_W(s)\bigr]+8\bar Bg^2L^{2-s},\qquad E_2=E_1+\tfrac6\pi\tfrac{a_K}d,
\]
which is at most an error function of $s$, since $a_K/d\le C_W(s)^{-1}$, $d<LC_W(s)L^{-s}$ and $M\le L^{-2}C_W(s)\le\hat d<LC_W(s)$ (Lemma~\ref{bil:sym}(e)), so that also the term $C^{\rm fl}_\star\hat d^3e^{-\delta_{\rm tp}L\hat d/\sqrt3}$ of $E_1$ is at most an error function of $s$. For (A3) apply (A2) to $(\theta Y,Y)$, whose points are $3d$ to $4d$ apart. By Lemma~\ref{bil:ker} and its count, at most $10(d/a_K)^3$ centres of bonds of one direction in a ball of radius $d$,
\[
|G(\theta Y,Y)|\le6C_{\mathsf k}^2(a_K/d)^6\Bigl(\sum|\mathsf t|\Bigr)^2\upsilon^2\le600C_{\mathsf k}^2\upsilon^2\sup|\mathsf w|^2,
\]
and with $\upsilon^2\le\frac98b^{-2}$, $|Y|_\star=59\sup|\mathsf w|$ the bound follows. Since $\langle Y\rangle_\ell=0$ (Lemma~\ref{bil:sym}(b)), $\langle\theta Y;Y\rangle_\ell=\langle\theta Y\cdot Y\rangle_\ell\ge0$ by reflection positivity at a plane through midpoints of bonds (Corollary~\YMMref{85.3}{679}); for other forms the centred form $F-\langle F\rangle_\ell$ gives the sign. For forms at bonds and sites the count gives $\sum|\mathsf t|\le N(d/a_K)^3\|v\|$, $N=30$, resp. $10$, and for $O^T_K(H)$ the same count gives the number $12150=6\cdot45^2$, with $60^2\bar\varepsilon^\natural_g$ from the size $60|H|_\infty$.

\emph{Step 3: (B).} By Lemma~\ref{bil:sym}(a) the copies $\theta'^\eta Y$ are bond forms of the same size; by Lemma~\ref{bil:sym}(d) partners across the plane of $\theta'_{\mu_0}$ are $3d$ to $4d$ apart and other pairs $\ge11M$; their means vanish by Lemma~\ref{bil:sym}(b). Lemma~\ref{bil:far} with $n=4$, for bond forms, gives $M_8\le(g_0^4+\bar\omega_{g,8})\hat m_J$ with $\hat m_J=59^8\upsilon^8\sup|\mathsf w|^8$ and $59^2g_0=106C'^2_{\mathsf k}$, so $M_8\le C_8b^{-8}\sup|\mathsf w|^8$; and $|M_8-\prod_\tau\langle P_\tau\rangle|\le\omega^\sharp_{g,8}\hat m_J$ with $\langle P_\tau\rangle=\langle\theta'_{\mu_0}Y;Y\rangle_{\ell'}$ for all four pairs. The signs come from reflection positivity: $M_8$ is a mean of mirror copies, and the chessboard estimate gives $\langle\theta'_{\mu_0}Y\cdot Y\rangle_{\ell'}\le M_8^{1/8}M_8^{1/8}$ times the seminorms $1$ of the six other cells.

\emph{Step 4: (C1).} With $\delta=\min\{d,\frac12,\mathsf l/(4n)\}$ cut each form into its pieces on the cubes of side $\delta$ (at most $(2R/\delta+2)^{3n}$ tuples of pieces); only the boundedness of the weights is used. For each tuple of pieces Proposition~\ref{bil:press}(c) gives a board of side $\mathsf l$ with all pieces at distance $\frac\delta2$ from the walls; on a torus, pieces are first moved by periods into one period box. Centre by $c_a=\tilde\mu_{\ell'}(F_a)$, write the cumulant through moments, and dominate each moment by Proposition~\ref{bil:press}(b). The eighth power of the seminorm of a product of centred pieces in one cell is a centred moment of at most $8m$ forms on $T_{\ell'}$ at mutual distance $\ge\delta$ (Lemma~\ref{bil:sym}(c)), bounded by (A1):
\[
\Bigl\|\prod_a(F_a-c_a)\Bigr\|_{{\rm ch},\ell'}\le\mathrm{Bl}_{8m}^{1/8}\prod_ab^{-1}|F_a|_\star,\qquad\mathrm{Bl}_N=\sum_{\pi\in{\rm Part}_2(N)}\prod_{I\in\pi}k^\circ(\#I,\delta),
\]
${\rm Part}_2(N)$ the partitions into blocks of at least two. So 
\[
C^T_n(d;R)=\max\{(2R/\delta+2)^{3n}\mathsf f_ne^{nG_g\mathsf l^3}\mathsf{Bm}_n^n,\ \relax C^{\rm cr}_n(d;R)\},
\]
$\mathsf f_n=\sum_\pi(\#\pi-1)!$, $\mathsf{Bm}_n=\max_{m\le n}\max\{1,\mathrm{Bl}^{1/8}_{8m}\}$; below the cutoff $K_n=\max_{r\le8n}K'_0(r,\delta)$ bounds by suprema give $C^{\rm cr}_n(d;R)=\mathsf f_n[6\bar x_0L^{K_n}(2RL^{K_n}+1)^3]^n$. The one-point number $C^T_1(R)$ is the case $n=1$. Complex weights: real and imaginary parts.

\emph{Step 5: (C2), (C3).} For (C2) let $D=|x_1-x_2|_\infty$, on a torus minimised over the lifts, be attained in the coordinate $\mu_0$. If $D>3$, the supports lie in balls of radius $1$ more than $1$ apart, and (C1) at $(\frac12;1)$ bounds the covariance. Let $D\le3$. One board of side $\mathsf l$ serves both insertions: it has a wall normal to $e_{\mu_0}$ at the midpoint of $x_{1,\mu_0}$ and $x_{2,\mu_0}$, up to $\frac{a_K}2$, and its offset $o$ at $\frac12(x_{1,\nu}+x_{2,\nu})$, up to $\frac{a_K}2$, in the two coordinates $\nu\ne\mu_0$. The distance $t_i$ of $x_i$ to the wall lies in $[\frac D2-\frac{a_K}2,\frac D2+\frac{a_K}2]\subset[\frac32d,2)$, and $|x_{i,\nu}-o_\nu|<2$, so that the corners of the plaquettes of either insertion lie within $3$ of the middle of its cell in the coordinates $\nu\ne\mu_0$.

Of the eight copies of one insertion on $T_{\ell'}$, the partners across the wall are $2t_i\in[3d',4d']$ apart, $d'=\max\{d,\frac{t_i}2\}$, and copies that differ in a reflection normal to an $e_\nu$, $\nu\ne\mu_0$, are at least $\mathsf l-6\ge11M$ apart ($\mathsf l\ge LM$, $M\ge3$); the carriers lie within $d'-2a_K$ of their points, and $s(d')\le s(d)\le K$. So Lemma~\ref{bil:far} with $n=4$ at the distance $d'$, for bond forms, whose means vanish, bounds the eighth power of the seminorm of a bond insertion by $(g_0^4+\bar\omega_{g,8})\hat m_J=C_8b^{-8}\sup|\mathsf w|^8$, as in Step~3, on this board, which the two insertions fix together, and its centred one-direction case bounds the seminorm of each of the forms of a tensor insertion; domination bounds the covariance by the product of the two seminorms; if $s(d')<s_8(g)$, (A1) is used at the distance $d_\flat=L^{-s_8(g)}$.

This gives
\[
\begin{multlined}
C^\diamond_\infty=\max\{2(c^\diamond_\infty)^2,C^T_2(\tfrac12;1)\},\qquad c^\diamond_\infty=e^{G_g\mathsf l^3}\max\bigl\{(\tfrac98)^{1/2}((422C'^2_{\mathsf k})^4+\bar\omega^{\rm o}_{g,8})^{1/8},\ \relax \mathrm{Bl}_8(d_\flat)^{1/8}, \\
\ \relax 12\bar x_0L^{K_\flat}(2L^{K_\flat}+1)^3\bigr\},
\end{multlined}
\]
$\bar\omega^{\rm o}_{g,8}$ the supremum of the error of Lemma~\ref{bil:far} for majorants of one direction, $\mathrm{Bl}_8(d_\flat)$ the number of Step~4 at the distance $d_\flat$ and $K_\flat$ its cutoff. For bond forms the mean vanishes in every law, and $C^{\rm per}_\infty$ is formed in the same way with $C_8^{1/4}$; complex weights of all three directions cost the factor $36$. For (C3) place the pair in the middle of a cell (Lemma~\ref{bil:sym}(d)) and use Lemma~\ref{bil:far} with $n=8$ for the sixteen copies (in its proof the blocks of one pair rest on the clause of Lemma~\ref{bil:two}(c), almost everywhere at the upper level): for $X_i=F_i-c_i$, $c_i=\tilde\mu_{\ell'}(F_i)$, and $m=\tilde\mu_{\ell'}(X_1X_2)=\tilde\mu_{\ell'}(F_1;F_2)$,
\[
\|X_1X_2-m\|_{{\rm ch},\ell'}^8=\sum_J(-m)^{8-\#J}\langle\tilde P_J\rangle_{\ell'}\le\omega'^\sharp_{g,16}(s)(|m|+\hat m)^8,
\]
since $\langle\tilde P_J\rangle$ replaced by $m^{\#J}$ gives $(m-m)^8=0$ and $|\langle\tilde P_J\rangle-m^{\#J}|\le\omega'^\sharp_{g,16}\hat m^{\#J}$. For one point, $\|X_i\|^8_{{\rm ch},\ell'}$ is a centred moment of eight of the sixteen forms, one from each pair, a set that holds no pair, so it is at most $o_{g,16}(s)(\upsilon|F_i|_\star)^8$. Domination now gives the three bounds of (C3), with
\[
{\rm Df}_g(s)=e^{G_g\mathsf l^3}[\omega'^\sharp_{g,16}(s)]^{1/8}(1+g_0+\omega'_{g,16}(s))+e^{2G_g\mathsf l^3}o_{g,16}(s)^{1/4},
\]
since $|m|\le(g_0+\omega'_{g,16})\hat m$, and $|m-G(F_1',F_2')|\le\omega'_{g,16}\hat m$ by Lemma~\ref{bil:far} at one pair. For $O^T_K(H_i)$ the same, with the majorants $(60|H_i|_\infty,S_{H_i})$.

\emph{Step 6: (D1)--(D3).} (C1), Proposition~\ref{bil:press}(d), the compactness of $I^{\rm lab}(g)$ and the theorem of F.~Riesz give (D1); the second clause is (C1) at $n=1$; for the third, the local functions $O^T_K(H)$ and the weak convergence give some $\ell_j\ge\max\{\Lambda_j,j\}$ in $\mathfrak L(g)$ at which the first $j$ tuples of a dense countable family are matched within $\frac1j$. $O^T_K(f\mathbf1)=\mathcal O_K(f)$ gives (D2), along a subsequence of the defining sequence of $S$. Lemma~\ref{bil:sym}(a) gives $rO^T_K(H)=O^T_K(rH)$; the invariance and the reflection positivity of the laws at the lattice mirrors pass to the limits, since lattice isometries approximate the real ones and the bounds are equicontinuous; periodic Gibbs states are positive at the nine families.

The semigroup of symmetric contractions $[F]\mapsto[\tau_{te_3}F]$ is $e^{-tH}$, Stone's theorem gives the momenta and the joint spectral measure, and the lattice energy--momentum inequality (Proposition~\textup{\YMMref{99.3}{739}}) gives $H\ge|P_k|$. Centred at the means of $\tilde\mu_{\ell'}$, the moments in $\nu_j$ are $\sum_B\prod_{i\in B}(\nu_j-\tilde\mu_{\ell'})(O^T_{K_j}(H_i))\,\nu_j(A_j(X_{B^c}))$. The bond parts have mean $0$ in both laws, and all plaquettes have one mean in each law (Proposition~\textup{\YMMref{103.4}{758}}), so that $(\nu_j-\tilde\mu_{\ell'})(O^T_{K_j}(H))=\mathfrak c_j\,a_{K_j}^3\sum_pH_{\lambda(p)\lambda(p)}(x(p))$, a Riemann sum of $\mathfrak c_j\int\mathrm{tr}\,H$, with $\mathfrak c_j=a_{K_j}^{-4}(\nu_j(a_p)-\tilde\mu_{\ell'}(a_p))$; (C1) at $n=1$, applied to a plaquette form, bounds $|\mathfrak c_j|$ by $C(g)\mathsf b^{-1}$ for large $j$.

So the $\mathfrak c_j$ converge along a further subsequence, to $\mathfrak c$, and along it the recentred moments converge to the functional of (D3), whose one-point function is $\mathfrak c\int\mathrm{tr}\,H$; it is reflection positive, since the means of $\tilde\mu_{\ell'}$ are invariant under the lattice mirrors.

\emph{Step 7: (D4).} For Lipschitz entries Lemma~\ref{bil:rate} makes the cumulants Cauchy in $K$, and its form for cell means gives the same for bounded Borel weights; continuous entries are approximated by Lipschitz ones uniformly in $K$, by the torus form of (C1); Proposition~\ref{bil:press}(d) applies in $\ell$ to periodised tuples, which are defined and linear once $\ell$ is large against the supports; a diagonal choice of the cutoffs gives the limits along growing tori.

\emph{Step 8: (D5), (F1).} On $T_{\ell'}$, (C3) and the continuum form of Step~2, with $\frac32\mathrm{tr}[E_{\lambda\lambda}KE_{\gamma\gamma}K]=\frac32K_{\lambda\gamma}^2$, give the law for all $36$ pairs of components; then $j\to\infty$. For two off-diagonal components the error is $\bar e^{\rm od}_g+\frac98 59^2{\rm Df}_g$, for an off-diagonal and a diagonal one $\bar e^{\rm od}_g+\frac98 59\,{\rm Df}_g$, for two diagonal ones $\frac98[\omega'_{g,16}+{\rm Df}_g+150\bar E_1(\frac1\pi+\bar E_1)+159/C_W(s)]+2\bar Bg^2L^{2-s}$, $\bar E_1$ the bound of $E_1$ by a function of $s$ ((C3) being used with the majorants of the respective kinds); $\bar e_g$ is their maximum. At $H_i=f_i\mathbf1$, $\frac32\mathrm{tr}\,K_3(u)^2=\frac9{16\pi^2}|u|^{-6}$; the leading term is covariant under $O(3)$. For (F1), Lemma~\ref{bil:far} with arbitrary majorants, centred, at $n=8$ bounds the means, $|\nu(O^T_K(H))-\tilde\mu_{\ell'}(O^T_K(H))|\le60e^{G_g\mathsf l^3}{\rm mn}_g(s)\upsilon|H|_\infty$, where ${\rm mn}_g=[e_g^8+\mathrm{Bl}(16){\rm far}_{g,16}(1+g_0+\varepsilon_g)^8+{\rm tl}_{g,16}]^{1/8}\le o_{g,16}^{1/8}$ comes from its moment bound at $N_1=8$, so sup norms suffice, and the law holds at a finite cutoff for continuous weights, the Riemann sums of continuous weights giving $r_K\to0$ and the lattice term of $E_1$ the term $C(g)(a_K/d)^2$; letting $j\to\infty$,
\[
\begin{multlined}
e^\infty_g(d)=\tfrac98\Bigl\{3600\bigl[\omega'_{g,16}+e^{G_g\mathsf l^3}(\omega'^\sharp_{g,16})^{1/8}(1+g_0+\omega'_{g,16})+e^{2G_g\mathsf l^3}{\rm mn}_g^2\bigr](s)+12150E^\infty_1(\tfrac1\pi \\
+E^\infty_1)\Bigr\}+162\bar Bg^2L^{2-s},\qquad s=s(d),
\end{multlined}
\]
$E^\infty_1=E_1-C_{\rm lat}(a_K/d)^2$. The errors are powers of $Y_{s(d)}^{-1}$, the worst the root $[\omega'^\sharp_{g,16}]^{1/8}$, and $Y_{s(d)}\sim\log\frac1d$. This root can be avoided: in Step~5 the eighth power of the seminorm of $X_1X_2-m$ is an alternating sum of means of products of the sixteen copies, and at the upper level the parts of these means without a block that links two pairs recombine to $\prod_\tau(\Psi_\tau-m)$, $\Psi_\tau$ the conditional means of the pairs, with $|\Psi_\tau-m|\le(\delta_g+\omega'_{g,16})\hat m$ almost everywhere at the upper level (the blocks of one pair by Lemma~\ref{bil:two}(c), the singletons as in the proof of Lemma~\ref{bil:far}, and $|G(F_1',F_2')-m|\le\omega'_{g,16}\hat m$ as in Step~5); so the seminorm itself is at most a multiple of $(\delta_g+\omega'_{g,16})\hat m$ plus the eighth roots of the terms ${\rm far}$ and ${\rm tl}$, which decay faster than every power of $Y_s^{-1}$, and the means are bounded with ${\rm mn}_g$, so that the second term of ${\rm Df}_g$ may be taken $e^{2G_g\mathsf l^3}{\rm mn}_g^2$; as $\varepsilon_g$, $\omega'_{g,16}$ and $1/C_W(s)$ are of order $Y_s^{-2\mathfrak p_0}$ and ${\rm mn}_g^2$ is of smaller order, the error is at most $C(g)(\log\frac1d)^{-2\mathfrak p_0}$; by Lemma~\ref{bil:dil} the same holds for $g\le\gamma_2$, as $(\log\frac1{L^zd})^{-2\mathfrak p_0}\le2^{2\mathfrak p_0}(\log\frac1d)^{-2\mathfrak p_0}$ for $d\le L^{-2z}$.

On the tori of $\mathfrak L(g)$ the torus limits of (D4) are limits along cutoffs at the labels $L^K\mathsf b$, and the same bound passes to them. The law fixes $\mathsf b$, since $\mathcal W\ne0$. Since $\mathcal W$ is invariant under $O(3)$ and $|rH|_\infty\le3|H|_\infty$, the law gives $|\sigma(rH_1,rH_2)-\sigma(H_1,H_2)|\le(9+1)e^\infty_g\mathsf b^{-2}|H_1|_\infty|H_2|_\infty$. Non-degeneracy: $\mathcal W(\theta(\varphi E),\varphi E)\ge w_0\,\mathrm{tr}\,E^2$ for this bump.

\emph{Step 9: (F2), (F9).} (i) is Lemma~\ref{bil:semi} at the margin $\delta_n$; as the number $C_E$ there is proportional to $\delta^{-4}$, the margin costs $n^4$, $C_E$ at $\delta_n$ being $C_{\rm sn}n^4$, and $\mathsf Y_i$ at $\delta_n$ is linear in $i$ and $\log n$; for $\gamma_1<g\le\gamma_2$ by Lemma~\ref{bil:dil}(b). (ii) For $K\ge i+\lceil\log_L4n\rceil$, $n(\frac{a_i}{4n}+2\delta_na_i+a_K)<a_i$, and Proposition~\ref{bil:press}(c) gives a board of side $a_i$ with an offset in $a_K\mathbb Z^3$ on which each support lies in a cell of its own, at the distance $\ge\delta_na_i$ from the walls; domination by a torus of $\mathfrak L(g)$ bounds the moments centred at the means of $\tilde\mu_{\ell'}$ by $\prod_a2e^{G_g}\mathsf E_ib^{-1}|H_a|_\infty$, $b^{-1}<g^2$, and centring in $\nu$ costs $2^n$. (iii) Cut the entries into their pieces on the boxes of side $\frac{a_i}{4n}$; pieces of different entries are more than $a_i$ apart; (ii) at the cutoffs and multilinearity give the bound, which passes to the limit. (iv) Take for $i$ the least integer with $L^{-i}\le d$ and $i\ge i_E(g,n)$; by the threshold in the proof of Lemma~\ref{bil:semi}, $L^{i_E(g,\delta)}\le C(g)\delta^{-1}$, so that $L^{i_E(g,\delta_n)}\le Cn$; then $2c_gn^4Y_i^{q_E}L^{11i/2}(4n/a_i)^3\le C_Gn^{16}d^{-9}$, the powers of the logarithm being absorbed, and $n^{16n}\le e^{16n}(n!)^{16}$. (v) The extension theorem, Proposition~\YMMref{93.3}{717}: a multilinear functional $\mathsf T$ on separated tuples of smooth functions with $|\mathsf T(X)|\le\mathsf c_{\mathsf T}\max\{1,\mathsf d^{-\mathsf a}\}\prod\|f_i\|_\sharp$ for supports in unit balls at least $\mathsf d$ apart, $\|f\|_\sharp=\max\{\sup|f|,40M\,\mathrm{Lip}\,f\}$, has exactly one extension to the Schwartz functions flat at coinciding points that is bounded for one norm $|\cdot|_k$, and it is bounded by $C_{\ne}^{n+\mathsf a}M^n(n!)^8\mathsf c_{\mathsf T}|\varphi|_{8n+\mathsf a-1}$; at $\mathsf a=9n$ and $\mathsf c_{\mathsf T}=C^n(e^{16}C_G)^n(n!)^{16}$ from (iv) this is the bound of (v). (vi) Let $i\ge i_E(g,\frac1{16})$, $z_3=\frac12a_i$, and $h$ supported in $\{|x-z|_\infty\le\frac38a_i\}$.

At the cutoff $K_j$ a board of side $a_i$ with the wall plane $\{x_3=c_j\}$, $c_j=\frac12a_{K_j}$, puts $h$ in one cell with the margin $\frac1{16}a_i$; by Lemma~\ref{bil:semi} the centred insertion $D_j$ of $h$ has seminorm $\le\mathsf c_E(i)b_j^{-1}|h|_\infty$, $\mathsf c_E(i)=2C_E\mathsf Y_i^{q_E}a_i^{-11/2}$, on every torus of $\mathfrak L(g)$; the reflection inequality for an insertion between two half spaces (Proposition~\YMMref{104.2}{759}) gives $|\rho_j(\overline{\theta_{c_j}A(F')}\cdot D_j\cdot\tau_{a_ie_3}A(F))|\le\mathsf c_E(i)b_j^{-1}|h|_\infty\|A(F')\|_{\rho_j,c_j}\|A(F)\|_{\rho_j,c_j}$ in the states of the sequence, and in the limit the slab operator of $h$ has norm $\le\mathsf c_E(i)\mathsf b^{-1}|h|_\infty$.

The semigroup and the translations, which are contractions, carry this bound at the scale $a_i\approx\eta$ to every $\eta\in(0,1]$, with the power $\eta^{-11/2-\epsilon}$; and the slab operators being holomorphic in the translation of their argument, the mean value inequality on a polydisc turns the sup norm on a cube of side $2\rho$ into the $L^1$ norm at the cost $4\rho^{-3}$. So the bound holds for every $p>\frac{11}2+3=\frac{17}2$. For $\gamma_1<g\le\gamma_2$ by Lemma~\ref{bil:dil}(b). (vii) is Lemma~\ref{bil:vec}(c), its hypotheses being (D1) and (D3), which give the invariance under $W_3$ and $H\ge|P_1|$, $H\ge|P_2|$; the insertion bound is not used. (viii) For pairwise different points $x_1,\dots,x_n$, more than half of the unit vectors $u$ have $|u\cdot(x_i-x_j)|\ge n^{-2}|x_i-x_j|$ for all $i\ne j$: the band $\{|u\cdot w|<n^{-2}\}$ has the fraction $n^{-2}$ of the sphere (Archimedes), and there are $\frac12n(n-1)$ pairs.

A rotation that carries such a $u$ to $e_3$ orders the points in the third coordinate with gaps $\ge n^{-2}\varrho(x)$, and the kernels $\mathbf W_n$ of Lemma~\ref{bil:vec} with $p=9$ give the bound at $\eta=\min\{1,\varrho(x)/(2n^2)\}$; the kernels obtained from different $u$ agree, since by (F3) they all represent $\sigma$, and form one continuous kernel. The growth follows from Taylor's formula at a point with two equal entries, of the order $\lceil pn\rceil$, and $(2n^2)^{9n}\le2^{9n}e^{18n}(n!)^{18}$.

(F9). In a periodic Gibbs state $\nu$, $\langle F',F\rangle_{\nu,c}=\nu(\overline{\theta_cF'}F)$ is a positive semi-definite Hermitian form on the local functions on $\{x_{\mu_0}>c\}$, the translation by $a_Ke_{\mu_0}$ is a contraction for it, and so the mirror may be moved away: $\|F\|_{\nu,c}\le\|F\|_{\nu,c'}$ for $c\le c'$, $c'-c\in a_K\mathbb Z$. It therefore suffices that $x_0$ lies at a distance between $\frac32d$ and $2d$ from the plane. There the plane is a wall of a board of side $\mathsf l$, $\theta X$ is the copy of $X$ in the neighbouring cell, and, as $\nu(\theta X\cdot X)=\nu(\theta(F-c)\cdot(F-c))-(\nu(F)-c)^2$ for $F=O^T_K(H)$ and $c$ its mean on $T_{\ell'}$, domination by $T_{\ell'}$ gives $\nu(\theta X\cdot X)\le e^{2G_g\mathsf l^3}\|F-c\|^2_{{\rm ch},\ell'}$, the eighth power of the seminorm being the mean of eight copies. For bond forms (B) bounds it, and for the whole tensor Lemma~\ref{bil:far} with arbitrary majorants, centred, at $n=4$; if $s(d)<s_8(g)$, (C1) at the distance $L^{-s_8(g)}$ bounds the covariance directly. So
\[
\begin{multlined}
C^{\rm od}_\infty(g)=\max\{e^{2G_g\mathsf l^3}C_8^{1/4},\,59^2C^T_2(L^{-s_8(g)};1)\},\qquad C_\infty(g)=\max\{e^{2G_g\mathsf l^3}(C'_8)^{1/4}, \\
\,60^2C^T_2(L^{-s_8(g)};1)\},
\end{multlined}
\]
$C'_8=(\frac98)^460^8[g_0^4+\bar\omega'_{g,8}]$, $g_0=\frac32 81^2C'^2_{\mathsf k}$; complex weights of all directions cost $36$. For the kernel bound at off-diagonal entries, Proposition~\ref{bil:rot}(a) bounds the two-point kernel at $u$ by $\Gamma(|u|/2\sqrt3)$, and the mirror bound for bond forms applied to bumps at the height $t$ bounds $\Gamma(t)$ by a multiple of $C^{\rm od}_\infty(g)\mathsf b^{-2}\max\{1,t^{-1}\}^6$; the holomorphy of the two-point function in the relative position and the mean value inequality on a polydisc of radius proportional to $|x-y|$ turn this bound for bumps at the scale $|x-y|$ into the pointwise bound $A_6C^{\rm od}_\infty(g)\mathsf b^{-2}\min\{|x-y|_\infty,3\}^{-6}$ of the kernel. For the whole tensor the same steps, with the mirror bound for $X=A[O^T_K(H)]$ and its number $C_\infty(g)$, give the bound with a pure number $A$ in the place of $A_6$. By (D5) the kernel is of order $|x-y|^{-6}$ at short distance, so the power is the true one.

\emph{Step 10: (F5).} Let $z=z(g,d)$, $g'_1$, $g'_2$, $\mathsf l_\iota$ and $\delta_\iota$ be those of Lemma~\ref{bil:dil}(d), which holds for $0<g\le\gamma$, and $\Lambda=L^z$. With $z$ so fixed, cut $\eta$ and the $H_a$ by a partition of unity into pieces of sup diameter $\le\mathsf l_1/(2q\Lambda)$, a number of $g$, $d$, $q$; the number of the pieces enters the number $C$ of (F5). The dilation by $\Lambda$ (Lemma~\ref{bil:dil}(a)) carries the datum $(K,\beta)$, $K\ge z+1$, to the datum $(K-z,\beta)$ at $g'_1$ or at $g'_2$, the pieces to pieces of sup diameter $r\le\mathsf l_1/(2q)$, and the supports to supports at sup distance $\ge\Lambda d\ge2\mathsf l_\iota$, so that two pieces of different entries never lie in one cell of side $\mathsf l_\iota$.

Lemma~\ref{bil:ward} at the cutoff $K'=K-z$ asks a board with an offset in $a'\mathbb Z^3$, $a'=a_{\lceil\phi K'\rceil}=L^{-\lceil\phi K'\rceil}$. The pigeonhole of Proposition~\ref{bil:press}(c) at $h=a'$, on the coarser lattice $a'\mathbb Z^3$, $\mathsf l_\iota/a'$ being an integer, gives such an offset with all the points of $q$ cubes of side $r$ at a distance $\ge\rho$ from the walls as soon as $q(r+2\rho+a')<\mathsf l_\iota$: in each direction there are $\mathsf l_\iota/a'$ offsets, and fewer than $q((r+2\rho)/a'+1)$ of them are bad. At $\rho=\delta_\iota=\mathsf l_\iota/(8q)$ this is $r+a'<\frac{3\mathsf l_\iota}{4q}$, which holds for $r\le\mathsf l_1/(2q)$ once $a'<\mathsf l_1/(4q)$, that is for $K\ge K_0(g,d,q)$.

So every piece of a tuple of pieces gets a cell of its own on a board of side $\mathsf l_\iota$, with an offset in $a'\mathbb Z^3$, of the least torus $T_{2\mathsf l_\iota}$ at the smaller coupling, the piece of $\eta$ at a distance $\ge\delta_\iota$ from the walls; domination (Proposition~\ref{bil:press}(b)) bounds the mean of the product by the product of the seminorms. The seminorm of the piece of $\eta$ is bounded by Lemma~\ref{bil:ward}, at the cutoff $K-z$ and the margin $\delta_\iota$, and the seminorms of the pieces of the $H_a$ by (A1) at the orders $2,\dots,8$, the side $2\mathsf l_\iota$ being small as $g'_\iota\le\gamma_1$, and the eighth power of a seminorm being a moment of eight copies; the cumulant is a polynomial in such means.

On $T_\ell$ directly, the cumulant at the source $0$ is bounded through the oscillations at the level $K-\lceil\phi_qK\rceil$ (step (iv) of the proof of Lemma~\ref{bil:ward}). There the slot of $H^{\rm cons}[\eta]$ has an oscillation $\le CK^pL^{(3\phi_q-\frac1{28})K}\|\eta\|_{C^2}$: its cubic term contributes $L^{-K/28}$ by the root (v), and the null entry of order two of the collapse (iii), whose oscillation at that level is $\le CK^pL^{3\lceil\phi_qK\rceil-K/28}$ for each unit of its weights, which are of size $b\sup|\nabla\eta|$, the factor $a_K^{-1}$ standing outside them (steps (iii) and (iv) at the level $K-\lceil K/14\rceil$), contributes $L^{(3\phi_q-\frac1{28})K}$. Each of the $q-1$ spectators costs $CK^{2\mathfrak p_0}L^{4\phi_qK}$.

The exponent $3\phi_q-\frac1{28}+4(q-1)\phi_q=(4q-1)\phi_q-\frac1{28}$ is $-\frac{51}{2800}$, $-\frac{23}{2800}$ at $\phi_q=\frac1{400}$ for $q=2,3$, and $-\frac1{56}$ at $\phi_q=\frac1{56(4q-1)}$ for $q\ge4$. In the limits the moments and cumulants with the entry $H^{\rm cons}[\eta]$ tend to $0$; on the tori the limits of (D4) along the full sequence of cutoffs vanish.

\emph{Step 11: (F3), (F4).} (F3) at $g\le\gamma_1$ is Proposition~\ref{bil:rot}(b), whose inputs are (D1), (D3), the bound $H\ge|P_k|$, the insertion bound (F2)(vi), the mirror bound (F9), conservation (F5) and the growth (R7) of Definition~\YMMref{95.1}{723} of the two-point function of the action density, with an exponent $a_\Gamma<4$; the domination of periodic Gibbs states in these inputs uses that $\mathfrak L(g)$ is unbounded. For the cumulants, by induction on $n$,
\[
S^T_n(X)=\sigma^T(X)-\sum_{\pi\ne\{X\}}\prod_{I\in\pi}S^T_{\#I}(X_I),
\]
the sum over the partitions into at least two blocks of at least two entries, $R$ acting entry by entry. For $0<g\le\gamma$ Lemma~\ref{bil:dil}(c) gives $\sigma(RX)=L^{nz}\sigma_z((RX)^{(z)})=L^{nz}\sigma_z(R(X^{(z)}))=\sigma(X)$, the rotations about $0$ commuting with the dilation. On the entries $f\mathbf1$ there is a second road. Along a further subsequence with $\mathfrak c_j\to\mathfrak c$ (Step~6), the functional centred at the means of the least torus, restricted to the entries $f\mathbf1$, is the limit of the moments of the action density in the $\nu_j$ recentred in the invariant laws $\tilde\mu_{\ell'}$, the limit functional of the sequence of pairs of laws $(\nu_j,\tilde\mu_{\ell'})$; Theorem~\YMMref{AN}{763} gives its invariance under $O(3)$ for $0<g\le\gamma$.

By the formula of (D3), inverted, $\sigma$ at the entries $f\mathbf1$ is a polynomial in this functional and in the numbers $\mathfrak c\int f$, which are invariant; so $\sigma$ is invariant there. The last sentence is Proposition~\ref{bil:zeta}. (F4) at $g\le\gamma_1$: Proposition~\ref{bil:cone}(c) in each of the nine pictures, with positivity at the nine families (D3): for the picture with the normal $e_3$ take $u''=e_2$, whose four planes have the normals $e_3$, $(e_3+e_1)/\sqrt2$, $e_1$, $(e_1-e_3)/\sqrt2$, and the fifth plane the normal $(e_3+e_2)/\sqrt2$; for a diagonal normal such as $(e_1+e_2)/\sqrt2$ take $u''=e_3$, the normals $e_1$, $(e_1+e_2)/\sqrt2$, $e_2$, $(e_2-e_1)/\sqrt2$ and $(e_1+e_3)/\sqrt2$; all of them belong to the nine families.

The tube is Lemma~\ref{bil:vec}(a); for $\gamma_1<g\le\gamma_2$ by Lemma~\ref{bil:dil}(b). Every mirror, for $0<g\le\gamma$: Proposition~\ref{bil:cone}(d) with (F3) and (D3).

\emph{Step 12: (F6)--(F8).} (F6). Let $X$, $Y$ be tuples in $(-\rho_0,\rho_0)^3$, and let $\rho$ be a state of the sequence at the cutoff $K=K_j$, or its conditioning on an invariant event of positive measure, again a periodic Gibbs state, so invariant under $\mathbb L_K$ and reflection positive at the three site mirrors, the three bond mirrors and one diagonal mirror of each coordinate plane (Proposition~\textup{\YMMref{103.4}{758}}). Let $r=r_j$ be the least integer with $ra_K\ge\rho_0+2a_K$; then $A(X)$, $A(Y)$ are functions of the bonds of the cube $Q_r=[-ra_K,ra_K]^3$. The clustering estimate~\textup{(1968)} of Proposition~\textup{\YMMref{99.7}{742}} and its sup-distance consequence in the proof, which rests on reflection positivity at these nine mirrors of one lattice, gives for $v\in\mathbb L_K$ with $|v|_\infty\ge4ra_K$
\[
\bigl|\rho\bigl(\overline{A(X)}\,\tau_vA(Y)\bigr)-\rho\bigl(\overline{\mathbb E[A(X)|\mathcal I]}\,\mathbb E[A(Y)|\mathcal I]\bigr)\bigr|\le4\mathsf C_{\rm dc}\,\mathcal N_r(A(X))^{1/2}\mathcal N_r(A(Y))^{1/2}\,\frac{ra_K}{|v|_\infty},
\]
$\mathcal I$ the $\sigma$-field of the events invariant under $\mathbb L_K$, $\mathsf C_{\rm dc}$ the absolute number in equation~\YMMref{(1968)}{742}, and $\mathcal N_r(A)=\max_\mu\max_{\mathsf g}\rho(\theta_\mu\mathsf g\bar A\cdot\tau_{2ra_Ke_\mu}\mathsf gA)$ the reflection size in equation~\YMMref{(1959)}{740}, $\theta_\mu$ the site mirrors through $0$ and $\mathsf g$ the eight sign changes. The tuples $\theta_\mu\mathsf g\bar X\cdot\tau_{2ra_Ke_\mu}\mathsf gX$ are separated, and (C1) bounds their values uniformly in $j$; as $j\to\infty$, $r_ja_{K_j}\to\rho_0$, the reflection sizes tend to $\mathcal N_{\rho_0}(X)$ by the convergence at separated tuples and the equicontinuity of (D1), and lattice vectors $v_j\to v$ carry the bound to every $v$ with $|v|_\infty>4\rho_0$: $\sigma$ obeys it with $4\mathsf C_{\rm dc}\mathcal N_{\rho_0}(F)^{1/2}\mathcal N_{\rho_0}(G)^{1/2}\rho_0/|v|_\infty$, for combinations $F$, $G$ of such tuples.

This gives the plateau, its rate and $m_2(\bar F,G)=\lim_j\nu_j(\overline{\mathbb E[A_j(F)|\mathcal I]}\,\mathbb E[A_j(G)|\mathcal I])$, so $m_2(\bar F,F)\ge0$. With the uniform bound the iterated plateaux commute with $j\to\infty$: the moments $\nu_j(\prod_i\mathbb E[A_j(X^i)|\mathcal I])$ converge to the iterated plateaux of $\sigma$. \emph{The decomposition.} Let $(X^i)_{i\ge1}$ be a countable family of tuples, dense in the sense of the proof of Proposition~\ref{bil:press}(d). The conditional means $\mathbb E[A_j(X^i)|\mathcal I]$ are bounded uniformly in $j$, by (C1) in the conditioned states, which are periodic Gibbs states; so the laws of $(\mathbb E[A_j(X^i)|\mathcal I])_i$ under $\nu_j$ live on one countable product $Q$ of compact intervals.

Their moments converge, and the polynomials are dense in $C(Q)$ (Stone--Weierstrass); so these laws converge weakly to a law $\lambda$ on $Q$ (F.~Riesz). Let $\mathsf X$ be the support of $\lambda$, a compact metrisable space, and $\sigma_\xi(X^i)=\xi_i$ the coordinates, extended to all tuples by the equicontinuity of Proposition~\ref{bil:press}(d). Then $\lambda$ has full support on $\mathsf X$, $\xi\mapsto\sigma_\xi(X)$ is continuous, $\xi\mapsto\sigma_\xi$ is injective by construction, $\int\prod_i\sigma_\xi(X^i)\lambda(d\xi)$ is the iterated plateau and, at $q=1$, $\sigma=\int\sigma_\xi\lambda(d\xi)$. For every $\xi$ every neighbourhood $U$ of $\xi$ has $\lambda(U)>0$, so its preimage under the map $(\mathbb E[A_j(X^i)|\mathcal I])_i$ is an invariant event of positive $\nu_j$-measure for large $j$; shrinking the neighbourhoods along a diagonal sequence gives invariant events $E_j$ with $\nu_j(E_j)>0$ such that the conditioned states $\rho_j$ have $\|\mathbb E_{\rho_j}[A_j(X)|\mathcal I]-\sigma_\xi(X)\|_{L^\infty(\rho_j)}\to0$, so $\sigma_\xi(X)=\lim_j\rho_j(A_j(X))$ along the full sequence; conditioned on an invariant event of positive measure a periodic Gibbs state is again a periodic Gibbs state, so the $\sigma_\xi$ re-centred in the $\rho_j$ are limits through periodic Gibbs states.

The plateau of a component factorises, so the vectors of its Hilbert space that are invariant under the semigroup, which form $\ker H_\xi$, are the multiples of $\Omega_\xi$. \emph{The map $\pi$.} Theorem~\textup{\YMMref{AP}{780}} gives a canonical decomposition $(\mathsf X_0,\lambda_0,\sigma_{0,\xi_0})$ of the limit functional of the action density along the sequence, with $\xi_0\mapsto\sigma_{0,\xi_0}$ continuous and injective, and its one-point constants $c_{\xi_0}$. At scalar tuples the iterated plateaux of $\sigma$ are those of the action density; so, for a countable dense family of scalar tuples $X^i_0$, the laws of $(\sigma_\xi(X^i_0))_i$ under $\lambda$ and of $(\sigma_{0,\xi_0}(X^i_0))_i$ under $\lambda_0$ have the same moments and are equal.

Both measures have full support and both maps are continuous on compact spaces, so the two images coincide: for every $\xi$ there is a $\xi_0=\pi(\xi)$ with $\sigma_\xi(f_1\mathbf1,\dots,f_n\mathbf1)=\sigma_{0,\pi(\xi)}(f_1,\dots,f_n)$ at all scalar tuples, and it is unique, $\xi_0\mapsto\sigma_{0,\xi_0}$ being injective. The graph of $\pi$ is closed in a compact space, so $\pi$ is continuous; it is onto, as the images coincide; and $\pi_*\lambda=\lambda_0$, as the two laws are equal. The one-point function of $\sigma_\xi$ is $\lim_j(\rho_j-\nu_j)(O^T_{K_j}(H))$; the bond parts have mean $0$ and the plaquettes have one mean in each of the two states, so it is $c'_\xi\int\mathrm{tr}\,H$ for a number $c'_\xi$, and at $H=f\mathbf1$ it is $\sigma_{0,\pi(\xi)}((f))=c_{\pi(\xi)}\int f$; so $\sigma_\xi((H))=\frac13c_{\pi(\xi)}\int\mathrm{tr}\,H$. (F7).

Interlace two choices of states, even $K$ from the one and odd $K$ from the other; conversely every subsequence has a convergent subsequence with the one limit. If all the limits agree, each $\sigma_{0,\xi}$, a limit through conditioned states, has $S^{\sigma_{0,\xi}}_n=S_n$, since cumulants of order $\ge2$ ignore the centring, and $S_2(f,h)=\int[S_2(f,h)+c_\xi^2\int f\int h]\lambda_0(d\xi)$ forces $\int c_\xi^2\,d\lambda_0=0$; the image of the vacuum projection is dense in $L^2(\mathsf X_0,\lambda_0)$, hence of dimension one. (F8). Every side $\ell=2L^m\in\mathfrak L(g)$ is admissible over the run of every datum at $g$: $LM\,R(x_K)\le L^m$ (equation~\YMMref{(1533)}{615}, as in the paragraph \emph{Data} of the Setting).

On $T_\ell$ the clause (i) of (E3), on Lemma~\ref{bil:nsrc} (the bound of that Lemma on an admissible side with $\max\{1,c(\ell)\}^nk'_n$ for $k'_n$, without the factor of the covering), at the geometry of Lemma~\ref{bil:nsrc} at $d=1.52h$ (a cube of side $\frac32h$ lies within $1.3h\le\frac78d$ of its centre, and cubes at sup distance $\ge\frac52h$ are more than $\frac54d$ apart), bounds the cumulants of $n$ entries at the cubes of side $\frac32h$ at sup distances $\ge\frac52h$, for $0<h\le h_0$, where $h_0\in\ell2^{-\mathbb N}$ and $h_0\le\min\{\ell/(64n),\frac12L^{-s_{\rm th}(n)}\}$, so that $1.52h\le1.52h_0<L^{-s_{\rm th}(n)}$, by $60^n\max\{1,c(\ell)\}^nC_n(g)[2\log\frac1{1.52h}+C(g)]^{8n\mathfrak p_0}\,b^{-n}\prod|H_a|_\infty$ uniformly in the cutoff, which is at most $C(1+\log\frac{h_0}h)^{8n\mathfrak p_0}b^{-n}\prod|H_a|_\infty$; through the cell means of Lemma~\ref{bil:rate} the bound passes to the limits of (D4), and Lemma~\ref{bil:ext} gives the extension, its bound, the contact terms and the counts.

At a cutoff the same construction, with the products of the cubes restricted to the scales $h$ with $s(1.52h)\le K$, gives subtracted lattice cumulants with the bound for all $\lambda\ge8h_K$, $h_K\le2C_W(K)a_K$ the least scale, uniformly in $K$; by Lemma~\ref{bil:rate} they converge to $\tilde S_n$; their difference to the lattice cumulants depends on $F$ only through the derivatives of the orders $\le3k$ on the strata of $k$ blocks and the values at distances $<6h_K$. On $\mathbb R^3$: let two cubes of side $\frac32h$ lie at sup distance $\ge\frac52h$, $d=1.31h$, $s(d)\le K$. Their centres are at sup distance $\ge4h\ge3d+a_K$, the cubes lie within $d-2a_K$ of their centres, and a mirror plane $\{x_{\mu_0}=c\}$ between them, normal to a coordinate in which the centres are $\ge4h$ apart, has each centre at a distance $\ge2h-\frac{a_K}2\ge\frac32d$.

On a torus $T_\ell$, $\ell\in\mathfrak L(g)$, the mirror $\theta$ has the second plane $\{x_{\mu_0}=c+\frac\ell2\}$, at a distance $\ge\frac\ell4-\frac{a_K}2\ge2h-\frac{a_K}2\ge\frac32d$ from each centre, as the sup distance $\ge4h$ of the centres on $T_\ell$ is at most $\frac\ell2$. For a bond form $Y$ with a complex weight of all directions, carried within $d-2a_K$ of a point at a distance $\ge\frac32d$ from both planes, and $\tau$ the translation by $a_Ke_{\mu_0}$, the invariance of the torus measure gives $\psi(k):=\langle\overline{\theta\tau^kY}\,\tau^kY\rangle_\ell=\langle\theta\bar Y\,\tau^{2k}Y\rangle_\ell$, and the Schwarz inequality for the positive form of $\theta$ on the half torus between the two planes (Corollary~\YMMref{85.3}{679}) gives $\psi(k)^2\le\psi(k-1)\psi(k+1)$ as long as the translates keep the distance $\ge\frac32d$ from both planes.

A non-negative sequence with this property takes its largest value at an end of its range, and there the point lies at a distance between $\frac32d$ and $2d$ from one of the two planes. That plane is a wall of a board of side $\mathsf l$, $\frac\ell2$ being a multiple of $\mathsf l$, the cells are distinct modulo the period, and domination by $T_{\ell'}$ with (B), or (C1) at the distance $L^{-s_8(g)}$ if $s(d)<s_8(g)$, bounds $\psi$ there as in (F9); so the mirror bound of (F9) for bond forms with complex weights of all directions, $0\le\langle\overline{\theta Y}Y\rangle_\ell\le36C^{\rm od}_\infty(g)b^{-2}|\mathsf w|_\infty^2$, holds on the tori of $\mathfrak L(g)$. With the Schwarz inequality for the positive form of that mirror, the means of bond forms being $0$, this mirror bound gives $|\langle Y[\mathsf w_1];Y[\mathsf w_2]\rangle_{K,\ell,\beta}|\le36C^{\rm od}_\infty(g)b^{-2}|\mathsf w_1|_\infty|\mathsf w_2|_\infty$ on every $\ell\in\mathfrak L(g)$: the bound of Lemma~\ref{bil:ext} with $q=0$ and a number free of the side, which every limit point inherits.

\emph{Step 13: (F10).} Lemma~\ref{bil:dil}(a), (b): by the exact dilation (Lemma~\ref{lem:lattice-law-dilation}) the cumulants, states, boards and seminorms at $g$ are those at $g'_1$ or $g'_2$, the free covariance is invariant, and the windows of labels of $g'_1$, $g'_2$ cover $L^zb$; the clauses at $g'_s\le\gamma_1$ carry over with the stated $d_0(g)$ and $e^\infty_g$; for all forms (C1) follows from (A1) on the least side, with numbers depending on that side. For $\gamma_2<g\le\gamma$ let $z$ be the least $z\ge1$ with $g'_1\le\gamma_1$ in Lemma~\ref{bil:dil}(d); then $g'_2<g'_1\le\gamma_1$, as $g'^{-2}_2=\frac L2g'^{-2}_1$. By (d) and (a), for $K\ge z+1$ the dilation by $L^z$ carries every periodic Gibbs state $\nu$ at a datum $(K,\beta)$ at $g$ to a periodic Gibbs state $\nu^{(z)}$ at the datum $(K-z,\beta)$ at $g'_1$ or at $g'_2$, whose label is $a_{K-z}\beta=L^zb$; the cumulant in $\nu$ of tensor entries $H_1,\dots,H_n$ with supports in balls of radius $R$ at mutual distance $\ge d$ is $L^{nz}$ times the cumulant in $\nu^{(z)}$ of the $H^{(z)}_a$, and these have the same sup norms and supports in balls of radius $L^zR$ at mutual distance $\ge L^zd$.

As $L^{nz}(L^zb)^{-n}=b^{-n}$, (C1) for tensor entries at $g'_1$ and $g'_2$ gives it at $g$ in every periodic Gibbs state at $(K,\beta)$, $K\ge z+1$, with the numbers at $g'_\iota$ and the distances and radii multiplied by $L^z$: $|S^T_n[\nu](H_1,\dots,H_n)|\le60^n\max_{\iota=1,2}C^T_n(L^zd;L^zR)\,b^{-n}\prod_a|H_a|_\infty$, $C^T_n$ formed at $g'_\iota$. For $0<g\le\gamma$, Lemma~\ref{bil:dil}(c): $L^z\mathsf b$ is an inner point of $I(g_z)$, and along a subsequence the dilated sequence converges at $g_z$.

\emph{Step 14: (E1)--(E3).} (E1): $\tilde\nu_K(\{Z_t\ne\emptyset\})=\sum_{t'\le t}\mathsf w_{t'}$, the weights of the first creations, and the run may be stopped at the level $t$ without changing this number. A first creation forces the failure of a local test in some test cube of the step; averaging the test over the gauge variables of its cube with the density of the gauge fixing of the step gives Borel functions $0\le F\le1$ that are local, gauge invariant and covariant under the lattice isometries, and the average is undone exactly inside the sum over the histories of any admissible set of values with the same $A_0$; so $\mathsf w_t$ is at most a sum of means of such functions, two for each test cube, its plaquette test and its averaged witness test, and at $t=1$ one for each density cube too: at most $2N_\tau+N_\delta$ for each cell of the level $t$, $N_\tau$ the number of the test cubes and $N_\delta$ that of the density cubes of a cell ($N_\delta$ entering at $t=1$ alone), and at most $(2N_\tau+N_\delta)M^{-3}\ell^3L^{3(K-t)}$ in all.

These means pass to a second admissible set of values with $\frac{93}{100}A_1$ for $A_1$ and the other constants that enter $c_\star$ kept, whose deposit constant $c_\star''$ has $(\frac{93}{100})^2c_\star\le c_\star''\le c_\star$: a product of test functions with disjoint supports is dominated by the probability, at the second set of values, that all the centres lie in the large-field region. Let $s_t$ be the side of a cell of the level $t$. If the torus has at least $2L^2$ cells of the level $t$ along a side, take the board of side $L^2s_t$ whose offset is the centre of the test function $F$, its \emph{witness}; the chessboard estimate with the mirrors in the walls of this board bounds $\langle F\rangle$ by the $N$-th root of the mean of the product of the copies of $F$ in its $N$ cells on the torus, and this product asks a large field in every cell of the board.

A creation reaches less than $49s_t<L^2s_t$ and so serves one cell alone: such a history has at least $N$ creations. Theorem~\YMMref{D}{8}, which bounds the density of these histories from above by $\exp[-\frac9{10}c''_\star p_0(g_t)^2N]$, at the fraction $\frac9{10}$ of the deposit, and that of the empty history from below, with the volume term $Cy_t(L^2MR(x_t))^3N$, which $\frac1{40}c_\star p_0(g_t)^2N$ pays, gives $\langle F\rangle\le\exp[-\frac34c_\star p_0(g_t)^2]$, as $\frac9{10}(\frac{93}{100})^2-\frac1{40}\ge\frac34$; with fewer cells of the level $t$ Theorem~\YMMref{D}{8} is used without mirrors, at the first set of values; a threshold on $y_K$ is removed by stopping the couplings at the last level above it.

The union bound over the cells gives $C_\flat$. The pair $(c^\natural_2,C_\natural)$: if $Z_t\ne\emptyset$, some cell of the level $t$ has its centre in $Z_t$, and the rarity at a point of Proposition~\ref{prop:local-large-field-rarity} with the union bound over the at most $M^{-3}\ell^3L^{3(K-t)}$ cells gives the bound. (E2): Lemma~\ref{bil:two}(b) with $\mathsf L=CM^{-3}\ell^3$; half of the exponent $cp_0(g_t)^2$ pays $\ell^3$ once $\frac12cA_0^2(\frac12s\log L)^P\ge3\log\ell$, which is $s\ge\sigma_c(\log\ell)^{1/P}$; the other half pays $\log C+3s\log L+2Y_s$ and the envelopes from $s_\flat(g)$ on; the depth condition $2c\mathfrak p_0A_0^2y^{P-1}\ge13$ is the rate condition of the deposit lemma at $c$. For (iv) there are three cases.

If $s(d)\ge\hat s$, (iv) is (iii). If $s(d)<\hat s$, then $d\ge L^{-s(d)}\ge L^2d'$, $d'=L^{-\hat s-1}$; let $s'=s(d')\ge\hat s$. If besides $K\ge s'$, cover by balls of radius $d'$, at which (ii) and (iii) hold; this costs the factor $(4/d'+2)^6\le5^6L^{6\hat s+6}$. If $K<s'$, the suprema of the forms give the number $100L^{8s'}\bar x_0^2$. So $C^\circ_3$ is the larger of $(4/d'+2)^6$ times the number of (ii) and (iii), in the units of (iv), and of $100L^{8s'}\bar x_0^2$; as $s'\le\hat s+1+\lceil\log_LC_W(2\hat s+2)\rceil$, $L^{8s'}\le L^{8\hat s+16}C_W(2\hat s+2)^8$, and with $L^{\hat s}\le L^{s_\flat(g)+1}\exp[2(\frac{6\log\ell}{cA_0^2})^{1/P}]$ this gives the bound of (iv), whose factor $(1+\log\ell)^8\exp[16(\frac{6\log\ell}{cA_0^2})^{1/P}]$ comes from the third case, the covering costing less; for (v) cover by balls of radius $\min\{\frac r8,L^{-\hat s}\}$. (E3): let $\ell=2L^m$ be admissible, $K\ge3$, and $s_{\rm th}(n)\le s(d)\le K$.

The proof of Lemma~\ref{bil:nsrc} reads the side at the following places, and each of them holds on an admissible side. The arithmetic of the working level (Lemma~\YMMref{126.1}{844}) asks nothing of the side. Lemma~\YMMref{126.2}{845} asks (Sd.1) alone, $MR(x_K)\le L^m$, and gives the frame, the order $x_0>x_1>\dots>x_K$ of the run (so $y_K\le y$), (Sd.2) and (Sd.3) ($LM\le L^m$, $\ell\ge N_{\mathsf H}$), the data of the construction at every level $\le K$, and the settings of the chain at the levels $\mathsf k<j_1=T=\mathsf k+1$ at every pair of the majorants; the setting of the chain for $n$ sources (Proposition~\ref{prop:weighted-canonical-history} at $p=n$) asks these data and the frame alone, and the threshold $s_{\rm th}(n)$ reads no side.

The sum over the level of the first creation (equation~\eqref{eq:ward-first-creation}) holds at all the levels $\le K$ (Theorem~\YMMref{D}{8} and Corollary~\YMMref{69.4}{604}, as in the proof of Lemma~\YMMref{126.2}{845}), and the rarity on one torus, Proposition~\ref{prop:ward-fixed-volume-rarity} at the level $K$, whose hypotheses hold by that canonical construction, bounds the weights before any smallness of the side is used:
\[
\mathsf w_t\le\tilde\nu_K(\{Z_t\ne\emptyset\})\le2M^{-3}\ell^3e^{\mathsf B(x_K)\ell^3}L^{3(K-t)}e^{-c'p_0(g_t)^2},\qquad\mathsf B(x_K)\le\mathsf B(\bar x_0).
\]
From there on the side enters the proof only through this number $2M^{-3}\ell^3e^{\mathsf B(\bar x_0)\ell^3}$, in the tail of the sum over the level of the first creation, where a small side has $2M^{-3}\ell^3\le M^{-3}c'A_0^2Y_s^P$, and through the number of the boxes of the covering. As $Y_s\ge1$ the quotient of the two numbers is at most $c(\ell)$, each block of a partition takes at most the factor $\max\{1,c(\ell)\}$, and there are at most $n$ blocks: the bound of Lemma~\ref{bil:nsrc} holds on every admissible side with $\max\{1,c(\ell)\}^nk'_n(g,s(d))$ for $k'_n(g,s(d))$, which is (E3)(i) at this pair; the covering of Step~1, with $2\ell$ for $2\ell_{\max}$, gives (E3)(ii). The same argument serves every pair $(\mathsf L,c'')$ of (E3): the weights enter only the tail term of each block, linearly in $\mathsf L$ and through a Gaussian sum in which $c''\ge\frac12c'$ may stand for $\frac12c'$, and no smallness of $\mathsf w_t$ is used; so, as $Y_s\ge1$, each block takes at most the factor $\max\{1,M^3\mathsf L/(c'A_0^2)\}$, which is $\max\{1,c(\ell)\}$ at the pair above, and (E3)(i) and (ii) hold with it. The weights of (E1), $\mathsf w_t\le C_\flat M^{-3}\ell^3L^{3(K-t)}e^{-\frac12c'p_0(g_t)^2}$, give $C_\flat\ell^3/(c'A_0^2)$ in the place of $c(\ell)$.
\end{proof}

\begin{remark}
Limits are taken along subsequences of given sequences of laws; only (D4), on a fixed torus, and (F7) concern the agreement of different subsequences. Theorem~\ref{bil:lim} speaks of limits through periodic Gibbs states (and, in (F5) and (F8), of the torus limits of (D4)); for limits along growing tori Theorem~\ref{thm:bil} asserts the bounds, the trace, the translations, the cubic group, the positivity at the coordinate planes and the free law. In (E) the first depth $\hat s(g,\ell)$ and the number $C^\circ_3(g,\ell)$ depend on the side, and (E2) concerns two entries; (E3) concerns any number of entries on an admissible side, with numbers that depend on the side through $c(\ell)$. All bounds concern separated supports, except (F8).
\end{remark}


\section{A complex scale for one fluctuation step}
\label{sec:complex-field-scale}

The coefficient arguments in Sections~\ref{C-350-1} and
\ref{C-351-1} require a disc of fluctuation scales whose radius
increases as the running coupling increases. Analyticity near the
real step alone does not give this disc. We obtain it by constructing
the potentials at a fixed lower real coupling, where the unscaled
field has a larger admissible range, and then evaluating the same
functions at $\lambda\mathsf t^{-1/2}B$.

We use the local input and the localized expansion of
Definition~\YMMref{52.1}{415} and Theorem~\YMMref{V}{416}.
In particular $k+1$ is a single admitted step,
$0\le k<K$, $M=L^{m'}$ satisfies $m'<m+K-k-1$, and
$P=(U,A',\mathbf J)$ lies in the output core
$\mathcal U^\circ(X)$. On a fixed real-$U$ slice its variables
$(A',\mathbf J)$ lie in the open set $\mathcal O_U(X)$.
The Gaussian structure, its restriction consistency, the
gauge-fixing operators and their bounds are those in that
definition. We retain the constants and numerical assumptions of
Theorem~\YMMref{V}{416}. Write
\[
 \delta(t)=A_1t^{-1/2}(\log t)^{\mathfrak p_0},\qquad
 \mathfrak b(t)=t^{1/2}\delta(t),\qquad
 s_{\mathfrak a}(t)=\delta(t)/a_{0,k+1},
\]
and let $\mathbb D_t^{\rm f}(Y)$ be the complex field polydisc
$\sup_b|B(b)|<C_{\mathcal A}\mathfrak b(t)$ on the nonpivot bonds
in $Y$. The potential norm uses the weight
\[
 Q_Y(B)=1+M^{-3}\sum_{b\in\mathfrak X^{\circ+}\langle Y\rangle}|B(b)|^2.
\]
All parameter sets below are either points or open subsets of a
finite-dimensional complex vector space; a point factor is omitted
from a product domain.

\begin{definition}[Scaled input]\label{IV-71-26}
Fix $\mathsf t>1$ in the admitted range and a radius $R>0$.
Suppose that four potentials $V^{J,\lambda},V^{G,\lambda},
V^{A,\lambda},V^{\prime,\lambda}_\theta$ are given for $|\lambda|<R$,
$\theta\in\Theta$, on every local core and field polydisc.
They have the pair dependence and gauge covariance of
Definition~\YMMref{52.1}{415}, and are jointly holomorphic on
\[
 D(0,R)\times\Theta\times\mathcal O_U(Y)
                   \times\mathbb D_{\mathsf t}^{\rm f}(Y).
\]
The first three are independent of $\theta$. The $J$ and fourth
potentials vanish at $B=0$; the $G$ and $A$ potentials vanish
through degree two in $B$. Suppose also that
\begin{equation}\label{eq:complex-input-bounds}
 \begin{aligned}
 |V^{o,\lambda}(Y;P,B)|
    &\le |\lambda|k_{\rm int}e^{-\lambda_Vd_k(Y)}Q_Y(B)
            &&(o=J,G,A),\\
 |V^{\prime,\lambda}_\theta(Y;P,B)|
    &\le |\lambda|k_{\rm iv}e^{-\lambda_Vd_k(Y)}Q_Y(B),\\
 Rk_{\rm int}&\le w_*/8,\qquad Rk_{\rm iv}\le w_*/4,
 \end{aligned}
\end{equation}
and $V^{o,-\lambda}(Y;P,-B)=V^{o,\lambda}(Y;P,B)$ for all four
potentials. The Gaussian structure and gauge-fixing operators are
fixed as $\lambda,\theta$ vary. The coefficient $x$ is an independent
parameter in $D(\mathsf t,\sigma_*\mathsf t)$, and $\tau=x/\mathsf t$.
\end{definition}

Use exactly the localized interpolation and the three bond terms
in the proof of Theorem~\YMMref{V}{416}: $\mathbf G_\tau$ is the
Gaussian term, $E^\chi$ the cutoff term, and
$E^{\bullet,\lambda}_\theta$ the interpolated interaction term.
For the last term the interpolation is
\[
 V^{\rm int}_v(B)=V^{J,\lambda}(vB)
       +v^{-2}[V^{G,\lambda}+\tau V^{A,\lambda}](vB),
 \qquad V'_v(B)=V^{\prime,\lambda}_\theta(vB).
\]
The value at $v=0$ is its removable continuation. Distribute bond
terms to sites with the same nonnegative tents $h_z$:
\[
 \mathbf G_\tau(X,P,z)=\sum_bh_z(b)\mathbf G_\tau(X,b;P),
\]
\[
 F^\lambda_\theta(X,P,z)=\sum_bh_z(b)
            (E^\chi+E^{\bullet,\lambda}_\theta)(X,b;P),\qquad
 F^\chi(X,P,z)=\sum_bh_z(b)E^\chi(X,b;P).
\]
Thus $F^\lambda-F^\chi$ is the interaction-dependent part. It is
this difference that vanishes at $\lambda=0$.

\begin{proposition}[The step on the full complex disc]\label{IV-71-37}
For the data of Definition~\ref{IV-71-26}, put
$M_\bullet=K_\Pi R(32k_{\rm int}+9k_{\rm iv})
\le(25/4)K_\Pi w_*$ and $d=d_{k+1}(X)$.
The following assertions hold on the full product domain
\begin{equation}\label{eq:complex-step-domain}
 D(\mathsf t,\sigma_*\mathsf t)\times D(0,R)
                    \times\Theta\times\mathcal O_U(X).
\end{equation}
\begin{enumerate}
\item[(cx1)] All three bond terms are defined and vanish unless
$b^\sharp\subset X$. Their tent sums are pointed local families.
\item[(cx2)] For fixed $x,P$, $\mathbf G_\tau$ and $E^\chi$ are
independent of $\lambda,\theta$ and of all four potentials. They
have the Gaussian and cutoff bounds of Theorem~\YMMref{V}{416}.
\item[(cx3)] The interaction term obeys
\[
 |E^{\bullet,\lambda}_\theta(X,b;P)|
 \le K_\Pi|\lambda|(32k_{\rm int}+9k_{\rm iv})
                                      e^{-\lambda_{\rm out}d}.
\]
For $F^\lambda-F^\chi$ the same bound holds with the factor $3L^3$.
In particular the respective bounds on the disc are
$M_\bullet e^{-\lambda_{\rm out}d}$ and
$3L^3M_\bullet e^{-\lambda_{\rm out}d}$.
\item[(cx4)] $E^{\bullet,0}_\theta=0$ and $F^0_\theta=F^\chi$.
\item[(cx5)] All terms are jointly holomorphic on
\eqref{eq:complex-step-domain}; the terms in (cx2) are holomorphic
in $x$ and the slice variables alone.
\item[(cx6)] Each individual bond term $E^{\bullet,\lambda}_\theta$
and each site term $F^\lambda_\theta$ is even in $\lambda$.
\item[(cx7)] The terms depend only on the local pair and are gauge
invariant, with the same restriction and locality properties as
the original step. No sum over domains is needed for these properties.
\item[(cx8)] The Taylor series starts at order two:
\[
 F^\lambda_\theta=F^\chi+
        \sum_{m\ge1}\lambda^{2m}F_{[2m],\theta},\qquad
 |F_{[2m],\theta}(X,P,z)|
   \le3L^3M_\bullet R^{-2m}e^{-\lambda_{\rm out}d}.
\]
The coefficients are jointly holomorphic in all remaining
parameters and inherit locality and gauge invariance. The
corresponding bond bounds omit $3L^3$.
\item[(cx9)] If $R>1$, then
\[
 |F^1-F^\chi|\le3L^3M_\bullet R^{-2}e^{-\lambda_{\rm out}d},
 \qquad
 |F^1-F^\chi-F_{[2]}|
 \le\frac{3L^3M_\bullet R^{-4}}{1-R^{-2}}
                                      e^{-\lambda_{\rm out}d}.
\]
\item[(cx10)] The local sum rules and the whole-torus integral
identity of Theorem~\YMMref{V}{416} hold for these potentials.
More explicitly, put $V^{\rm int,\lambda}=V^{J,\lambda}
+V^{G,\lambda}+\tau V^{A,\lambda}$ and
$V^{\rm q}=-(1-\tau)\sum_\square
(B,\widehat G_\square B)/2$. Then
\begin{align*}
 \Xi^\lambda_\theta(x;P_*)
 :={}&\int d\mu_{\tau^{-1}C_*(P_*)}(B)
       \prod_{b\in\mathfrak X^+}\mathbf1_{\{|B(b)|<\mathfrak b(\mathsf t)\}}\,
       \exp\sum_Y(V^{\rm q}+V^{\rm int,\lambda}
                         +V^{\prime,\lambda}_\theta)(Y;P_*|_Y,B|_Y)\\
 ={}&\exp\sum_{X,z}(\mathbf G_\tau+F^\lambda_\theta)
                         (X,P_*|_X,z)\ne0.
\end{align*}
The inverse covariance has positive real part, the integral
converges absolutely, and its integrand is defined as zero where
the cutoff product vanishes. The Gaussian measure and integral
are complex when their data are complex.
\item[(cx11)] If $R>1$ and the four potentials at $\lambda=1$
are the admitted potentials of the original step, the terms
at $\lambda=1$ are exactly its terms, with the same localization.
\end{enumerate}
\end{proposition}

\begin{proof}
We check the quantitative bounds in the construction rather than
increasing the constant $K_P$ in its global smallness conditions.
The Gaussian scaling calculation in Theorem~\YMMref{V}{416}
depends only on $x$, the fixed structure and the numerical
constants. Thus on its stated $x$ disc the scaled structure
$(\tau^{-1}C,\tau^{1/2}\Gamma)$ is admitted, the whole covariance
has inverse with positive real part, and
$\|V^{\rm q}\|\le w_*/2$.

Put $\nu_{\rm int}=|\lambda|k_{\rm int}$ and
$\nu_{\rm iv}=|\lambda|k_{\rm iv}$. The interpolation estimate,
Proposition~\YMMref{48.4}{372}, uses only the potential norms and
their indicated orders of vanishing. It gives
\[
 \|V_v^{\rm int}\|\le\tfrac72\nu_{\rm int}\le7w_*/16,
 \quad \|V'_v\|\le\nu_{\rm iv}\le w_*/4,
 \quad
 \|\dot V^{\rm int}_{v,b}\|\le32\nu_{\rm int},
 \quad \|\dot V'_{v,b}\|\le9\nu_{\rm iv}.
\]
These are precisely the margins used in the proof of
Theorem~\YMMref{V}{416}. Its insertion expansion therefore defines
$E^{\bullet,\lambda}_\theta$ and gives (cx3), with its support.
The first two terms use no interaction potentials, proving (cx2).
The tent mass at one site is at most $3L^3$, giving the site bounds.
At $\lambda=0$ all four potentials vanish by
\eqref{eq:complex-input-bounds}; hence so do the interpolation
and its insertion, proving (cx4).

The same insertion majorant is uniform on the full product domain.
Each refined integrand is holomorphic in its parameters on the
slice and fixed field polydisc. The common majorant permits
parameter integration and locally uniform summation, so the terms
are jointly holomorphic. The cutoff terms use their fixed-density
formula, whose differentiated masses have the same summable
majorant, exactly as in that theorem. This proves (cx5).

For parity apply $B\mapsto-B$ to each refined Gaussian integral.
Its covariance, every quadratic gauge term and every cutoff are
unchanged. The hypothesis on the four potentials gives
$V_v^{o,-\lambda}(-B)=V_v^{o,\lambda}(B)$ for the localized
interpolations, including their bond velocities: multiplication
by the interpolation field commutes with $-1$, and $v^{-2}$ is
unchanged. In a refined derivative the signs from differentiating
the potentials are matched by those from the corresponding
Gaussian derivatives; equivalently both are conjugated by the
orthogonal map $-I$ on the bond-coordinate space. The insertion
expansion is covariant under this map. Thus equality holds for
each domain and each marked bond before summation. It proves
(cx6). Pair dependence, gauge covariance and restriction commute
with the same construction, proving (cx7).

For fixed remaining parameters, $E^{\bullet,\lambda}_\theta$ is
even, holomorphic, zero at zero and bounded by
$M_\bullet e^{-\lambda_{\rm out}d}$. Cauchy's formula on
$|\lambda|=r<R$ bounds its degree-$2m$ coefficient by
$M_\bullet r^{-2m}e^{-\lambda_{\rm out}d}$; let $r\uparrow R$.
The odd coefficients vanish. Coefficient extraction by this
contour integral preserves holomorphy in the remaining parameters
and all the established invariances. Finite tent summation gives
(cx8). Schwarz's lemma applied after the order-two zero gives
the first bound in (cx9); summing the coefficient bounds for
$m\ge2$ gives its second bound.

For (cx10), integrate the insertion along $v$: the sum of its bond
velocities is the derivative of the interpolated potential.
Thus its sum is the difference of the full cluster sum and the
one with only $V^{\rm q}$. The cutoff term supplies the difference
between the finite-threshold and infinite-threshold sums for
$V^{\rm q}$, and the Gaussian term supplies the latter sum.
The tents have total mass one when summed over their sites.
Restriction consistency identifies the domain terms with those
computed on the whole torus. The complete Gaussian expansion
therefore gives exactly the displayed exponential identity.
Its common majorant proves absolute convergence; its exponential
form proves nonvanishing. Finally the formulas defining every
term are unchanged at $\lambda=1$, which proves (cx11).
\end{proof}

We now construct this input on a large disc. The required
one-variable estimate is included to distinguish the gain from
the mere vanishing of a potential at the origin.

\begin{lemma}[Rescaling a holomorphic potential]\label{lem:complex-ray}
Let $E_r=\{B:\sup|B|<r\}$, $0<r_\flat\le r$, and
$|\mu|\le r_\flat/r$. Suppose $f$ is holomorphic on $E_{r_\flat}$.
If $f(0)=0$ and $|f(B')|\le aQ_Y(B')$, then
\[
 |f(\mu B)|\le a|\mu|(r/r_\flat)Q_Y(B),\qquad B\in E_r.
\]
If the bound is $a$ instead, omit $Q_Y$. If $f$ vanishes through
degree two, then the stronger bound is
\[
 |f(\mu B)|\le a|\mu|^3(r/r_\flat)
                   (1+M^3/r_\flat^2)(Q_Y(B)-1).
\]
\end{lemma}

\begin{proof}
For the first assertion apply Schwarz's lemma to
$z\mapsto f(z(r_\flat/r)B)$ on the unit disc. Its bound is
$aQ_Y(B)$, since $Q_Y(cB)=1+|c|^2(Q_Y(B)-1)$.
For the third-order assertion let $\beta=\sup|B|>0$;
the case $B=0$ is immediate. For $|\mu|\le s<r_\flat/\beta$,
the function $z\mapsto f(zsB)$ vanishes to order three and is
bounded on the unit disc by $a[1+s^2(Q_Y(B)-1)]$. Schwarz's
lemma of order three gives
\[
 |f(\mu B)|\le a|\mu|^3
             [s^{-3}+s^{-1}(Q_Y(B)-1)].
\]
Let $s\uparrow r_\flat/\beta$ and use
$\beta^2\le M^3(Q_Y(B)-1)$ and $\beta<r$.
This proves the displayed estimate.
\end{proof}

\begin{definition}[A lower calibration]\label{IV-71-12}
Retain one lattice, one system of radii, its open slices and one
Gaussian structure and gauge-fixing family. Let
$1<\mathsf t_\flat\le\mathsf t$ and
$\delta(\mathsf t)\le\delta(\mathsf t_\flat)$.
All numerical and radius conditions of Theorem~\YMMref{V}{416}
are required at $\mathsf t_\flat$, with
$\varepsilon(\mathsf t_\flat),\delta(\mathsf t_\flat)\le1$,
$\mathfrak b(\mathsf t_\flat)\ge\mathfrak b_*$ and
$s_{\mathfrak a}(\mathsf t_\flat)\le s^*$.
Suppose four functions of the unscaled field $\mathcal A$ are
defined for $\sup|\mathcal A|<C_{\mathcal A}\delta(\mathsf t_\flat)$.
For $s\in\{\mathsf t_\flat,\mathsf t\}$ set
\[
 V^J[s](B)=v^J(s^{-1/2}B),\quad
 V^o[s](B)=-s\,v^o(s^{-1/2}B)\ (o=G,A),\quad
 V'_\theta[s](B)=v'_\theta(s^{-1/2}B).
\]
The dependence on $Y,P$ is suppressed. At $s=\mathsf t_\flat$
require the full local input of Definition~\YMMref{52.1}{415},
jointly in $\theta\in\Theta$, with fourth-family size at most
$n_\flat\le N_*$. No bound at a larger unscaled field is assumed.
\end{definition}

\begin{proposition}[The disc supplied by a lower calibration]
\label{IV-71-46}
For Definition~\ref{IV-71-12}, put
\[
 R=\frac{\delta(\mathsf t_\flat)}{2\delta(\mathsf t)},\qquad
 k_{\rm int}=2K_P\delta(\mathsf t),\qquad
 k_{\rm iv}=K_\Lambda s_{\mathfrak a}(\mathsf t)n_\flat.
\]
On $|\lambda|<2R$, define
\begin{equation}\label{eq:complex-scaled-families}
 V^{J,\lambda}(B)=v^J(\lambda\mathsf t^{-1/2}B),\qquad
 V^{\prime,\lambda}_\theta(B)=v'_\theta(\lambda\mathsf t^{-1/2}B),
 \qquad
 V^{o,\lambda}(B)=-\mathsf t\lambda^{-2}
             v^o(\lambda\mathsf t^{-1/2}B)\quad(o=G,A),
\end{equation}
with the last two values at zero defined as zero.
Restricted to $|\lambda|<R$, these satisfy
Definition~\ref{IV-71-26}. In particular
Proposition~\ref{IV-71-37} holds on its full product domain, with
\[
 M_\bullet=K_\Pi
    [32K_P\delta(\mathsf t_\flat)
       +\tfrac92K_\Lambda s_{\mathfrak a}(\mathsf t_\flat)n_\flat]
       \le\tfrac{41}8K_\Pi w_*.
\]
If $R>1$ and the input at $\mathsf t$ is the original admitted
input, the value at $\lambda=1$ is its original step. The
order-two gain and fourth-order remainder carry respectively the
factors $R^{-2}=4\delta(\mathsf t)^2/\delta(\mathsf t_\flat)^2$
and $R^{-4}=16\delta(\mathsf t)^4/\delta(\mathsf t_\flat)^4$.
\end{proposition}

\begin{proof}
Let $r=C_{\mathcal A}\mathfrak b(\mathsf t)$,
$r_\flat=C_{\mathcal A}\mathfrak b(\mathsf t_\flat)$ and
$\mu=\lambda(\mathsf t_\flat/\mathsf t)^{1/2}$.
Then $r_\flat\le r$ and
$|\mu|r/r_\flat=|\lambda|/(2R)$.
Thus $\mu B$ lies in the lower field polydisc for $|\lambda|<2R$,
and \eqref{eq:complex-scaled-families} equals
\[
 V^J[\mathsf t_\flat](\mu B),\qquad
 V'_\theta[\mathsf t_\flat](\mu B),\qquad
 \mu^{-2}V^o[\mathsf t_\flat](\mu B).
\]
The last functions extend holomorphically through $\lambda=0$:
their numerators vanish through order two, so division by
$\lambda^2$ has a removable zero, also jointly in the other
parameters by Cauchy's coefficient formula.

Apply Lemma~\ref{lem:complex-ray} at the lower place. For $J$
the bound is $|\lambda|K_P\delta(\mathsf t)Q_Y(B)
e^{-\lambda_Vd_k(Y)}$. For $G,A$, division by $\mu^2$ leaves
the bound
\[
 |\lambda|K_P\delta(\mathsf t)
 (1+M^3/r_\flat^2)(Q_Y(B)-1)e^{-\lambda_Vd_k(Y)}.
\]
The threshold floor gives
$\mathfrak b(\mathsf t_\flat)^2>(LM)^3/2$, and
$C_{\mathcal A}\ge3$, so $1+M^3/r_\flat^2<11/9<2$.
The constant-bound version of the lemma gives for the fourth
potential $|\lambda|K_\Lambda s_{\mathfrak a}(\mathsf t)n_\flat
e^{-\lambda_Vd_k(Y)}$.
Moreover
\[
 Rk_{\rm int}=K_P\delta(\mathsf t_\flat)\le w_*/8,
 \qquad
 Rk_{\rm iv}=\tfrac12K_\Lambda
 s_{\mathfrak a}(\mathsf t_\flat)n_\flat\le w_*/8.
\]
These prove all the size requirements. Pair dependence and gauge
covariance are preserved by scalar multiplication of the field.
The products $(-\lambda)(-B)$ and $\lambda B$ agree, as do
$(-\lambda)^{-2}$ and $\lambda^{-2}$, proving parity.
The orders of vanishing in $B$ are unchanged. At $\lambda=1$
the defining formulas are precisely those at $\mathsf t$.
The coefficient and remainder claims now follow from
Proposition~\ref{IV-71-37}.
\end{proof}

\begin{proposition}[The actual construction at two calibrations]
\label{IV-71-50}
Fix the constants, the tower and shift systems, and one lag-weight
case $(\omega,\mathsf K)$ of the actual step
Theorem~\YMMref{W}{423}. Suppose its geometric, core, shift and
octet hypotheses hold at a lower real calibration
$\mathsf t_\flat>1$, with the same radii $\mathfrak a$, and choose
positive lag weights satisfying
\[
 \varrho_j^2L^{-(k-j)}\le c_\omega\omega_{k+1-j}
 \quad(1\le j\le k),\qquad \mathsf K=K_\Lambda c_\omega.
\]
Retain the rate bounds and all numerical hypotheses of
Theorem~\YMMref{V}{416} at $\mathsf t_\flat$, with
$K_\Lambda$ replaced by this $\mathsf K$. Choose
$\check s>0$ such that
\[
 \check s\le s^*,\qquad8K_P\check s\le w_*,\qquad
 4\mathsf K N_*\check s\le w_*,\qquad
 s_{\mathfrak a}(\mathsf t_\flat)\le\check s,
 \quad \mathfrak b(\mathsf t_\flat)\ge\mathfrak b_*.
\]
Let $\mathsf t\ge\mathsf t_\flat$ with
$\delta(\mathsf t)\le\delta(\mathsf t_\flat)$. For a bare array
$f_\circ$ on these same cores, assume
$N_{\mathfrak a,\omega}(f_\circ)\le N_*$; at height zero the
array is empty. Then the canonical Gaussian structure and
the actual four potentials give Definition~\ref{IV-71-12}, with
$K_\Lambda$ replaced by $\mathsf K$ and
$n_\flat\le N_{\mathfrak a,\omega}(f_\circ)$.
The step is consequently defined on
\[
 x\in D(\mathsf t,\sigma_*\mathsf t),\qquad
 |\lambda|<\delta(\mathsf t_\flat)/(2\delta(\mathsf t)).
\]
If the latter radius exceeds one, its terms at one coincide with
the actual canonical step terms. This statement concerns
$\mathbf G_\tau,F^\lambda$; the separate local normalizer is
unchanged and is not included in $F^\lambda$.

The same conclusion holds jointly for a holomorphic family of
arrays with this uniform size bound. If an array has a reserve
$N_{\mathfrak a,\omega}(f_\circ)+2q\le N_*$ with $q>0$, its first and mixed
second array derivatives, in directions $h_1,h_2$, satisfy
\begin{equation}\label{eq:complex-array-derivatives}
 \begin{split}
 |D_fF^\lambda[h_1](X,P,z)|
   &\le Cq^{-1}\|h_1\|\,\frac{(|\lambda|/R)^2}{1-(|\lambda|/R)^2}
                                  e^{-\lambda_{\rm out}d},\\
 |D_f^2F^\lambda[h_1,h_2](X,P,z)|
   &\le Cq^{-2}\|h_1\|\|h_2\|
       \frac{(|\lambda|/R)^2}{1-(|\lambda|/R)^2}
                                  e^{-\lambda_{\rm out}d},
 \end{split}
\end{equation}
where $R=\delta(\mathsf t_\flat)/(2\delta(\mathsf t))$,
$\|h\|=N_{\mathfrak a,\omega}(h)$, and $C$ depends only on
the fixed constants and $\check s$. The derivatives retain the
local dependence on the array and weights.
\end{proposition}

\begin{proof}
We verify both the common objects and their domains.
The radius and octet conditions of the actual step involve only
the chosen radii, geometric constants and shift systems.
Their coupling-dependent inequalities involve
$\delta(t)$, $s_{\mathfrak a}(t)$ and the lower bound on
$\mathfrak b(t)$. The first two decrease between the two specified
calibrations and the third increases. Thus every such hypothesis
valid at $\mathsf t_\flat$ is valid at $\mathsf t$.
The lag ratios and their weight inequalities do not involve $t$.

The actual covariance, its square root, the local backgrounds
and the gauge-fixing matrices are built from the pair, the lattice,
the fixed shifts and the radii. Changing $t$ with these data fixed
changes none of them. In the Jacobian formula the field occurs
as $\mathcal A=t^{-1/2}B$. In the gauge and action potentials the
same field occurs inside their cubic remainders, with the sole
additional factor $-t$. In the fourth-potential formula the
substitutions, finite vertex differences and octet evaluations
are fixed; the field again enters only as $\mathcal A$.
These facts follow directly from the actual formulas in the
proofs of Theorem~\YMMref{W}{423} and
Proposition~\YMMref{51.16}{412}. They supply the four unscaled
functions in Definition~\ref{IV-71-12} on
$\sup|\mathcal A|<C_{\mathcal A}\delta(\mathsf t_\flat)$:
this is exactly the image of the lower field polydisc under
$B\mapsto\mathsf t_\flat^{-1/2}B$.

At that lower calibration Theorem~\YMMref{W}{423} gives the
three canonical interaction bounds and the Gaussian input.
Proposition~\YMMref{51.16}{412} bounds the fourth potential by
\[
 \mathsf K s_{\mathfrak a}(\mathsf t_\flat)
 N_{\mathfrak a,\omega}(f_\circ)e^{-\lambda_Vd_k(Y)}.
\]
It is linear in the array, holomorphic, and zero at zero field.
Thus its lower size is at most $N_{\mathfrak a,\omega}(f_\circ)$.
All premises of Definition~\ref{IV-71-12} are verified, with
$\mathsf K$ as its fourth-potential constant.
Proposition~\ref{IV-71-46} proves the asserted domain and the
identification at one. It also gives the uniform bound
\[
 M_\bullet\le K_\Pi(32K_P+\tfrac92\mathsf K N_*)\check s.
\]
Uniform holomorphic array parameters are permitted throughout
the localized construction, so the parameter statement follows.

For the derivative estimates, vary the array by
$z_1h_1+z_2h_2$ on the two discs
$|z_i|<q/\|h_i\|$, omitting a zero direction. The reserve keeps
their sum within the admitted ball. The term $F^\chi$ is
independent of the array. For each even coefficient of
$F^\lambda-F^\chi$, (cx8) and Cauchy's formula in the $z_i$
give respectively the additional factors
$q^{-1}\|h_1\|$ and $q^{-2}\|h_1\|\|h_2\|$.
Sum the coefficient bounds for $m\ge1$ to obtain
\eqref{eq:complex-array-derivatives}. The finite local
substitutions use only input domains contained in the output
domain, and taking these derivatives does not change that
dependence.
\end{proof}

When the step is subsequently evaluated at an even holomorphic
$x=\xi(\lambda)$, this substitution is legitimate only while
$|\xi(\lambda)-\mathsf t|<\sigma_*\mathsf t$. It preserves
holomorphy and evenness. Although the Gaussian and cutoff terms
may then vary with $\lambda$ through $\xi$, they remain independent
of the bare array; the array-derivative argument is unchanged.


\subsection{Source families and the local step}
\label{sec:source-entry-calculus}

An insertion changes the potentials integrated in a step. It does
not change its Gaussian structure. We specify the norm in which this
change is small and the parameter dependence needed to differentiate
the resulting local terms.

Use one admitted step of Definition~\YMMref{52.1}{415} and
Theorem~\YMMref{V}{416}, including their lattice, core, Gaussian
and numerical hypotheses. At height $k$ write $P$ for a background
triple and let $B$ range over the complex nonpivot-field polydisc
$\sup_b|B(b)|<C_{\mathcal A}\mathfrak b(\mathsf t)$.
For a family $V$ define
\begin{equation}\label{eq:source-entry-norm}
 \mathsf n^{\rm f}(V)=
 \sup_{X,P}\ \sup_{Y^\sharp\subset X,\,B}
 \frac{e^{\lambda_Vd_k(Y)}|V(Y;P|_Y,B)|}{1+q_Y(B)},
 \qquad q_Y(B)=M^{-3}\sum_{b\in\mathfrak X^{\circ+}\langle Y\rangle}
                                      |B(b)|^2.
\end{equation}
Keeping $X,P$ fixed defines $\mathsf n^{\rm f}_{X,P}$;
keeping only $X$ fixed defines $\mathsf n^{\rm f}_X$.
A family in the fourth place is holomorphic on the open real-$U$
slices and the field polydiscs, depends on the background pair,
is gauge invariant, vanishes at $B=0$, and has finite norm
\eqref{eq:source-entry-norm}. Holomorphic families of such inputs
are required to be jointly holomorphic in their parameters and
these variables, with a uniform norm bound.

The three ordinary interactions $V^{\rm J},V^{\rm G},V^{\rm A}$
are also families in the fourth place. Indeed their analyticity,
pair dependence and gauge invariance are part of the local input;
$V^{\rm J}(0)=0$, and $V^{\rm G},V^{\rm A}$ vanish through
degree two. Their pointwise bounds give
\begin{equation}\label{eq:source-internal-size}
 \mathsf n^{\rm f}_X(V^\natural)
 \le\mathsf n^{\rm f}(V^\natural)
 \le K_P\delta(\mathsf t),\qquad \natural={\rm J},{\rm G},{\rm A}.
\end{equation}
The same bound holds for their even and odd parts, since $B\mapsto-B$
preserves the polydisc and $q_Y$.

For the canonical absolute-radius step, the source potentials are
defined on the larger unscaled-field sets
\begin{equation}\label{eq:source-entry-domain}
 \mathfrak D^{\rm src}(Y)
 =\{(P,\mathcal A):P\in\mathcal U^\circ(Y),\quad
                              \sup_b|\mathcal A(b)|<\varepsilon_{\rm src}\},
 \qquad
 \varepsilon_{\rm src}=
 \begin{cases}
 \min\{\varepsilon_{\mathfrak h},c_{\rm os}\alpha_0/(2c_{\mathfrak h})\},
                                                        &k\ge1,\\
 r_P^\circ,                                             &k=0.
 \end{cases}
\end{equation}
Retain $C_{\mathcal A}\delta(\mathsf t)\le\varepsilon_{\rm src}/2$
and $0<\lambda_V\le\bar\kappa-2$, as in
Lemma~\YMMref{114.1}{800}. Thus $\mathcal A=\mathsf t^{-1/2}B$
lies strictly inside this source domain. Its radius is fixed by the
cores and does not shrink with $\mathsf t$.

Let $\mathbb W^{\rm src}[v]$ be the direct weighted-action family
of that lemma, with the reference linear term subtracted bond by
bond. Choose $K_{\rm Q}^\vee,K_{\rm L}^\vee$ satisfying
equation~\YMMref{(2166)}{800}, and put
\begin{equation}\label{eq:source-weight-room}
 \begin{gathered}
 K^{\rm d}=\tfrac{351}{2}L^3K_\Pi,\qquad
 \mu_\star=\min\{w_\star/32,E_0/(16K_{\mathbf M}^{\rm f})\},
 \qquad K^{\rm d}\le K_{\mathbf M}^{\rm f},\\
 \sigma_{\rm src}=\mu_\star/K_{\rm Q}^\vee,\qquad
 o_{\rm src}(\mathsf t)=\mathsf t^{-1/2}\mu_\star/K_{\rm L}^\vee,\\
 \mathsf{Wgt}_k^{\rm src}(\mathsf t)
 =\{v:|v|_\infty<\sigma_{\rm src},\quad
                 \operatorname{sw}_k(v)<o_{\rm src}(\mathsf t)\}.
 \end{gathered}
\end{equation}
Here $E_0$ denotes the size $E^{\rm s}_0$ in the recursion
preceding Proposition~\YMMref{114.2}{801}, and
$K_{\mathbf M}^{\rm f}$ is its constant
$3L^3\max\{K_Ee^{4\lambda_1},32K_\Pi K_P,117K_\Pi/2\}$.
The oscillation compares plaquettes whose barycentres are at sup
distance at most $M$ in level-$k$ units. Write
\[
 \mathsf{sz}_X(v)=K_{\rm Q}^\vee
       (1+2C_{\mathcal A}\delta/\varepsilon_{\rm src})|v|_X
                 +K_{\rm L}^\vee\mathsf t^{1/2}\operatorname{sw}_k(v).
\]

\begin{proposition}[Changing the fourth family]\label{prop:source-entry}
Retain the preceding admitted step. At the real field scale
$x=\mathsf t$, let the three internal families have total norm
at most $3w_\star/8$. The pointed integration map $\mathbf F$
is holomorphic in a holomorphic fourth-family input on its
admitted norm ball and is local in that input. It preserves pair
dependence, gauge invariance, conjugation and the lattice
covariances of the data. If $V_0,V_1$ have fourth-family norms
at most $w_\star/4$ and
$\mathsf n^{\rm f}(V_1-V_0)\le w_\star/2$, then
\[
 |\mathbf F[V_1]-\mathbf F[V_0]|(X,P,z)
 \le K^{\rm d}\mathsf n_X^{\rm f}(V_1-V_0)
                         e^{-\lambda_{\rm out}d_{k+1}(X)}.
\]
In particular let $\mathfrak f$ be a canonical old array satisfying
$8K_\Lambda s_{\mathfrak a}N_{\mathfrak a}(\mathfrak f)\le w_\star$,
with the empty array at $k=0$. For $v\in\mathsf{Wgt}_k^{\rm src}$,
the fourth family
$V'[\mathfrak f]-\mathbb W^{\rm src}[v]$ has norm less than
$7w_\star/32$. Its output
$E_v=\mathbf N+\mathbf F[V'[\mathfrak f]-\mathbb W^{\rm src}[v]]$
is therefore defined on the full output core, jointly holomorphic
in the weight and any holomorphic old-array parameters, and
\[
 |E_v-E_{v=0}|(X,P,z)
 \le K^{\rm d}\mathsf{sz}_X(v)
                         e^{-\lambda_{\rm out}d_{k+1}(X)}.
\]
The maps obey the exact local and whole-torus sum rules of
Theorem~\YMMref{V}{416}. On an actual real fibre admitted by
Lemma~\ref{lem:sf-chart-integral}, the source term sums to the
weighted-action difference of equation~\YMMref{(2168)}{800}.
\end{proposition}

\begin{proof}
At a fixed background the four inputs are families on the field
polydiscs. Lemma~\YMMref{77.1}{631} constructs their bond insertion
terms, and Proposition~\YMMref{77.2}{632} compares two fourth
families. In its second estimate use
$n_{\rm int}\le3w_\star/8$, $n_{\rm iv}\le w_\star/4$ and
$n_W\le w_\star/2$. Its coefficient is at most
\[
 9(3L^3-3)K_\Pi\{1+8(3/8+1/4)\}
 <162L^3K_\Pi<K^{\rm d}.
\]
The same proof at fixed $X$ gives the local norm. Holomorphy follows
from the parameter-dependent insertion theorem with its common
summable majorant. Restriction identifies every Gaussian matrix,
potential and insertion on a smaller domain, and hence identifies
the localized terms. The covariance assertions follow from the
corresponding changes of variables in those same terms.

Proposition~\YMMref{51.16}{412} bounds the old fourth family by
$K_\Lambda s_{\mathfrak a}N_{\mathfrak a}(\mathfrak f)\le w_\star/8$.
Lemma~\YMMref{114.1}{800} gives
$\mathsf n_X^{\rm f}(\mathbb W^{\rm src}[v])\le\mathsf{sz}_X(v)$.
The half-radius condition and \eqref{eq:source-weight-room} imply
\[
 \mathsf{sz}_{T_\ell}(v)<2\mu_\star+\mu_\star
       \le3w_\star/32.
\]
Their sum is strictly below $7w_\star/32<w_\star/4$.
The old-potential map preserves holomorphic parameters and the
source family is linear in $v$, so the preceding parameter result
applies. Here $E_{v=0}$ is the output with the zero weight and
the same old array. The source-independent normalizer cancels
in $E_v-E_{v=0}$.
Finally the local sum rule is precisely the pointed insertion
identity. The whole-torus identity follows by restriction and
summation of the tents. For the source summation the reference
linear terms cancel bond by bond by criticality of the real
background; the remaining terms are its weighted-action
difference, as proved in Lemma~\YMMref{114.1}{800}.
\end{proof}

\begin{proposition}[Holomorphic parameters and parity]
\label{prop:source-entry-parity}
Use the same local Gaussian and numerical hypotheses, allowing
$x\in D(\mathsf t,\sigma_\star\mathsf t)$ and $\tau=x/\mathsf t$.
Let four families depend jointly holomorphically on
$(\mu,\zeta)\in D(0,\rho)\times\mathcal O$, where $\mathcal O$
is an open finite-dimensional parameter set or a point. The
Gaussian structure, $\tau$ and the gauge-fixing operators do not
depend on $\mu$. Suppose each internal family has norm at most
$\bar\nu_{\rm int}\le w_\star/8$, the fourth has norm at most
$\bar\nu'\le w_\star/4$, and all four retain the stated vanishing
orders at $B=0$. Put
\[
 \bar B=K_\Pi(32\bar\nu_{\rm int}+9\bar\nu'),\qquad
 \Phi(\mu,\zeta)=\sum_bh_z(b)E^\bullet(X,b;\mu,\zeta),\qquad
 d=d_{k+1}(X).
\]
Then the Gaussian and cutoff terms do not depend on $\mu$,
$\Phi$ is holomorphic, and
$|\Phi|\le3L^3\bar B e^{-\lambda_{\rm out}d}$.
If each of the four families satisfies
$V^\natural_{-\mu,\zeta}(Y;-B)=V^\natural_{\mu,\zeta}(Y;B)$,
then $\Phi(-\mu,\zeta)=\Phi(\mu,\zeta)$.
Its odd Taylor coefficients vanish. With
$\Phi_{(n)}=(n!)^{-1}\partial_\mu^n\Phi(0,\zeta)$ one has
\[
 |\Phi_{(n)}|\le3L^3\bar B\rho^{-n}e^{-\lambda_{\rm out}d},
 \qquad
 \left|\Phi-\sum_{n\le N}\Phi_{(n)}\mu^n\right|
 \le3L^3\bar B\frac{(|\mu|/\rho)^{N+1}}{1-|\mu|/\rho}
                                       e^{-\lambda_{\rm out}d}.
\]
For $\rho>1$, parity improves the remainder after degree two at
$\mu=1$ to
$3L^3\bar B\rho^{-4}(1-\rho^{-2})^{-1}e^{-\lambda_{\rm out}d}$.
\end{proposition}

\begin{proof}
Use the interpolation displayed in Section~\ref{sec:complex-field-scale}.
The bond insertions from the three internal families have total
norm at most $32\bar\nu_{\rm int}$, and that from the fourth
has norm at most $9\bar\nu'$. The insertion theorem,
Theorem~\YMMref{S}{366}, bounds their interpolated integral by
$\bar B e^{-\lambda_{\rm out}d}$. The nonnegative tents have
total at most $3L^3$, giving the asserted site bound. The common
majorant also permits integration of the holomorphic parameter
dependence. The Gaussian and cutoff terms use none of the four
families.

For parity take the orthogonal map $O_b=-I$ at every nonpivot
bond. It preserves the field polydiscs and cutoffs. The covariance,
coupling and gauge-fixing operators are unchanged, since
$(-I)T(-I)=T$ for each of these linear operators. Their quadratic
potentials and marked quadratic insertions are unchanged as well.
The interpolation commutes with $B\mapsto-B$, by
Proposition~\YMMref{48.4}{372}. Thus the families and bond
insertions at $\mu$ are carried to those at $-\mu$.
The orthogonal covariance of the insertion and cutoff expansions,
Theorems~\YMMref{S}{366} and~\YMMref{T}{370}, identifies their
individual localized terms. Integration of these identical terms
and then summation with the fixed tents proves that $\Phi$ is even.
This argument compares the localized terms, not only their total
partition function.

Cauchy's formula on circles of radius increasing to $\rho$ gives
the coefficient bounds. Summing their geometric majorant gives
the remainder estimate. Evenness removes the odd coefficients;
at $\mu=1$ the remaining tail is bounded by the geometric series
$\rho^{-4}+\rho^{-6}+\cdots$.
\end{proof}

\section{Sources and the first large-field region}
\label{sec:levelwise-first-creation}

The comparison in Proposition~\ref{C-350-8} needs a source domain
that does not shrink when an early large-field region is created.
We construct the source dependence with the real Gaussian structure
at every step. Its total constant term agrees near zero with that
of the weighted progression; the individual localized families
need not agree.

Fix the real canonical run $x_0,\ldots,x_K$, its counterterms
$e_i=L^{-1}x_{i-1}-x_i$, and its all-small-field histories.
We retain the global core, actual-sum-rule, source-radius and
normalizer hypotheses of Lemma~\YMMref{114.1}{800} and
Proposition~\YMMref{114.2}{801}, together with the canonical
history tests of Proposition~\YMMref{123.1}{829}.
In particular $|e_i|\le B_\sharp$, the source potential
$\mathbb W^{\rm src}[v]$ has its actual, rather than merely
formal, sum rule, and all steps use their prescribed gauge factors
and density tests. Write $E_{\rm s}=E^{\rm s}_0$ and
$N_{\rm s}=N^{\rm s}_\star$. The numerical conditions used below are
\begin{gather*}
 E_{\rm s}\ge2E_{\mathsf N},\qquad
 N_{\rm s}\ge\frac L{L-1}
                   (E_{\rm s}+\tfrac32\alpha_0^2B_\sharp),\\
 K_{\mathbf M}\bigl(\delta(\gamma^{-2})+
       \mathsf s_\gamma N_{\rm s}
       +e^{-\gamma_2\mathfrak b(\gamma^{-2})^2/2}\bigr)
       \le E_{\rm s}/4,\qquad
 8K_\Lambda\mathsf s_\gamma N_{\rm s}\le w_\star,
 \qquad8K_P\mathsf s_\gamma\le w_\star,
\end{gather*}
where $\mathsf s_\gamma=(1+\beta_{\rm s})
\delta(\gamma^{-2})/\alpha_0$. We also retain
$4\beta_{\rm s}\kappa\ge1$ and the normalizer decay bound
$\delta_{\mathsf N}M\ge\kappa+1$. Put
\[
 K^{\rm d}=\tfrac{351}{2}L^3K_\Pi\le K_{\mathbf M}^{\rm f},
 \qquad
 \mu_\star=\min\{w_\star/32,E_{\rm s}/(16K_{\mathbf M}^{\rm f})\},
 \qquad \sigma_{\rm src}=\mu_\star/K_{\rm Q}^\vee.
\]
Here $K_{\mathbf M}^{\rm f}$ and $K_{\rm Q}^\vee,K_{\rm L}^\vee$
are the constants in that source construction. The level-wise
source domain is the open unit ball of
\begin{equation}\label{eq:levelwise-source-norm}
 \|w\|_{\rm lw}=\max_{0\le j<K}
  \max\left\{\frac{|v_j^w|_\infty}{\sigma_{\rm src}},
      \frac{\operatorname{sw}_j(v_j^w)}
                  {x_j^{-1/2}\mu_\star/K_{\rm L}^\vee}\right\},
 \qquad v_j^w=-L^{K-j}w/x_j.
\end{equation}
The swing compares bare plaquettes whose barycentres have sup
distance at most $M$ in level-$j$ units. Since the amplitude
is included, this is a norm on the finite-dimensional complex
space of bare plaquette weights.

\begin{proposition}[The sourced history density]
\label{prop:levelwise-source-density}
On $\|w\|_{\rm lw}<1$ there are local pointed families
$F_j[w]$, $1\le j\le K$, jointly holomorphic in the source
and in every open core slice, with
\begin{equation}\label{eq:levelwise-family-bound}
 \|F_j[w]\|_{j,\kappa+1}\le E_{\rm s},\qquad
 \|\Delta F_j[w]\|_{j,\kappa+1}
       \le2E_{\rm s}\|w\|_{\rm lw},\qquad
 \Delta F_j[w]=F_j[w]-F_j[0].
\end{equation}
A member indexed by $X$ depends only on the source plaquettes
in $X$. At zero source these are the ordinary progression's
families. Define the finite sums
\begin{align}
 \mathcal E_t(w)&=\sum_{j\le t}\sum_{X,z}
                              \Delta F_j[w](X,\mathbf1,z),
 \qquad \mathcal E=\mathcal E_K,\label{eq:levelwise-energy}\\
 \mathsf R_t(w,W)&=\mathcal O(w)(U_t(W))
       +\sum_{j\le t}\sum_{X,z}
       [\Delta F_j[w](X,U_t(W),z)
                    -\Delta F_j[w](X,\mathbf1,z)].
                                      \label{eq:levelwise-remainder}
\end{align}
At $t=0$ set $\mathcal E_0=0$ and
$\mathsf R_0(w,U)=\mathcal O(w)(U)$.
If $m_t^w$ is the all-small-field history measure with bare
input $e^{\mathcal O(w)}\rho_0\,dU$, and $m_t=m_t^0$, then
\begin{equation}\label{eq:levelwise-density}
 m_t^w=e^{\mathcal E_t(w)+\mathsf R_t(w,\cdot)}m_t.
\end{equation}
For $t\ge1$ the right side is evaluated on the passing density
set and is zero elsewhere. Moreover, on the common neighborhood
of zero in any finite set of source parameters, $\mathcal E_t$
equals the total source energy of the weighted progression of
Proposition~\YMMref{122.1}{826}. In particular its mixed second
derivative at $t=K$ is $\mathsf J_2$.
\end{proposition}

\begin{proof}
Starting with the empty array, use the actual old-potential
rule and the canonical localized step at $x_j$, with $\tau=1$:
\[
 \mathfrak f_j[w]=((F_i[w],e_i)_{1\le i\le j},\emptyset),
 \qquad
 F_{j+1}[w]=\mathbf N^{(j+1)}+
   \mathbf F[\mathbb V'[\mathfrak f_j[w]]
                               -\mathbb W^{\rm src}[v_j^w]].
\]
This is the real-Gaussian source recursion of
Proposition~\YMMref{114.2}{801}, with the counterterms of
the running coupling retained. Their contribution to the
array size is explicit:
\[
 N(\mathfrak f_j[w])=
   \sum_{i=1}^j L^{-(j-i)}
      (\|F_i[w]\|_{i,\kappa}+\tfrac32\alpha_0^2|e_i|)
 \le\frac L{L-1}(E_{\rm s}+\tfrac32\alpha_0^2B_\sharp)
 \le N_{\rm s}.
\]
The induction therefore has the same ceilings as that
proposition. The old potential has norm at most $w_\star/8$.
Lemma~\YMMref{114.1}{800} and \eqref{eq:levelwise-source-norm}
give
\[
 \mathsf n^{\rm f}(\mathbb W^{\rm src}[v_j^w])
 \le2K_{\rm Q}^\vee|v_j^w|_\infty
           +K_{\rm L}^\vee x_j^{1/2}\operatorname{sw}_j(v_j^w)
 \le3\mu_\star.
\]
Thus the fourth potential has norm at most $7w_\star/32$;
the three internal potentials together have norm at most
$3w_\star/8$. Lemma~\YMMref{77.1}{631} applies. The
unperturbed output is bounded by $E_{\rm s}/4$, the normalizer
by $E_{\rm s}/2$, and the perturbation estimate (1601),
p.~631, bounds the source contribution by
$3K^{\rm d}\mu_\star\le3E_{\rm s}/16$.
The new norm is at most $15E_{\rm s}/16<E_{\rm s}$.
The inequalities on the decay rates give the norm with rate
$\kappa+1$ used here.

The same step is jointly holomorphic in its potentials and
open core variables. Since $v_j^w$ is linear in $w$, induction
proves the asserted holomorphy. The local rule for a member
on $X$ uses only earlier members and source plaquettes inside
$X$, proving locality. At $w=0$ the source potential vanishes,
so the recursion is exactly the ordinary one. Finally, on the
complex line through $w$ the holomorphic difference
$F_j[\zeta w/\|w\|_{\rm lw}]-F_j[0]$ vanishes at zero and
has norm at most $2E_{\rm s}$ for $|\zeta|<1$.
Schwarz's lemma proves the second bound in
\eqref{eq:levelwise-family-bound}; the assertion at $w=0$
is immediate.

We next verify that this recursion gives the same history
operator used for the weighted progression. With $j=t-1$,
its all-small-field step factor is
\[
 g_t=(\chi_t\circ\mathcal M)\,\mathsf q_j^0\mathsf q_t
                    \mathfrak G_j[x_j]\chi_j^{\rm sf},
 \qquad m_t^w=\mathcal M_*(g_t m_j^w).
\]
The factors are bounded, nonnegative and Borel, and independent
of $w$. Here $\mathsf q_j^0,\mathsf q_t$ are the compatibility
tests of this history, and $\chi_t$ is its full density test.
For a passing output $W$, every point $V_W(B)$ of the actual
small-field chart satisfies
\[
 \mathsf q_j^0(V_W(B))=\mathsf q_t(V_W(B))
                    =\chi_j(V_W(B))=1.
\]
These are the chart and test properties used in the proof of
Proposition~\YMMref{123.1}{829}. Thus, if $F'=e^cF$ on the
passing input set, the image of $g_tF'\,dV$ has the pointwise
chart density $\chi_t(W)e^c\mathsf I_j[F](W)$.
Indeed, write the image integral against a bounded Borel test
function, insert $\chi_t(\mathcal M V)$, and apply the physical
chart formula to $\mathsf q_j^0\mathsf q_tF'$. On the chart
the displayed identities replace this integrand by $e^cF$.
This proves the claim for every passing $W$.

The source-potential sum rule in Lemma~\YMMref{114.1}{800}
identifies the sum of the fourth potential with the change
along that fibre of
\[
 S[\mathfrak f_j[w]](U)-x_jA^{\eta_j}(v_j^w,U).
\]
Its subtracted linear terms sum to zero by criticality.
Theorem~\YMMref{W}{423} supplies the other three sum rules,
the same Gaussian covariance and the same positive flat
normalizing factor $\exp(\mathsf{cs}_j)$, independently of
$w$. The exact cluster identity then gives the next chart
integral, exactly as in Proposition~\YMMref{114.2}{801}.
Keeping the counterterms changes only the following action
identity:
\[
 L^t x_t+\sum_{i=1}^tL^ie_i=x_0.
\]
Consequently, after $t$ steps the pointwise exponent is
\[
 A_t^w(U)=-x_0A^{\eta_0}(U)+\mathcal O(w)(U)
       +\sum_{i\le t}\sum_{X,z}F_i[w](X,U,z)
       +\sum_{j<t}\mathsf{cs}_j.
\]
To see this directly, the counterterms combine with
$-x_tA^{\eta_t}$ by the displayed identity, and the flat
values subtracted in the bare-array action cancel the flat
values in the vacuum factor. Induction using the preceding
chart identity gives density $\exp A_t^w(U_t(W))$ on the
passing set. The backgrounds are in all required cores, and
the local functions are invariant on their critical gauge
orbits. The chart formula proves Borel measurability; the
finite number of bounded core terms proves boundedness.
At $w=0$ this density is strictly positive. Dividing it out
and adding and subtracting the flat values proves
\eqref{eq:levelwise-density}.

For completeness, the comparison of source energies uses
this pointwise chart identity, rather than a choice of a
version of a measure density. Both constructions start with
the identical bare density $e^{\mathcal O(w)}\rho_0$ and use
the identical factors $g_t$ just displayed. The exact
weighted sum rule, on its admitted parameter neighborhood,
therefore evaluates the same successive chart integrals.
At $W=\mathbf1$ all density and compatibility tests pass:
their tested curvatures and their deviations from the trivial
datum are zero. The canonical background is $U_t(\mathbf1)
=\mathbf1$; this also follows from the trivial-triple
assertion of Proposition~\YMMref{51.9}{402} and the uniqueness
of the canonical orbit. Hence $\mathcal O(w)(\mathbf1)=0$
and all centered terms vanish. The common, source-independent
flat factors cancel, giving
\[
 e^{\mathcal E_t(w)}=e^{\mathcal E_t^{\rm weighted}(w)}.
\]
On a connected common neighborhood of zero their difference
is holomorphic, takes values in $2\pi i\mathbb Z$, and is
zero at zero. It is therefore zero throughout that
neighborhood. This proves the claimed equality and the
identification of its mixed derivative.
\end{proof}

\begin{proposition}[Envelopes and first creation]
\label{prop:first-creation-envelopes}
Retain the preceding hypotheses and the calibrated real
history bounds of Theorem~\YMMref{D}{8}. Suppose that the
source is supported in one physical ball of radius $1$.
For $1\le t\le K$ put
\begin{equation}\label{eq:first-creation-radius}
 d_*(t)=\min\left\{\frac{c_A\alpha_0}{2MB_3\varepsilon_t},
                    \frac{L^{m+K-t}-4}{2M}\right\}-1.
\end{equation}
Assume $M\alpha_0\le\alpha_1$,
$\alpha_0+\alpha_1\le a^\sharp$, and $d_*(t)\ge3$.
There are numbers $\mathsf G_t,\mathsf H_t$ such that
\begin{equation}\label{eq:first-creation-envelopes}
 |\mathsf R_t(w,\cdot)|\le\|w\|_{\rm lw}\mathsf G_t,
 \qquad
 |\mathcal E(w)-\mathcal E_t(w)|
                       \le\|w\|_{\rm lw}\mathsf H_t,
\end{equation}
on the passing set, with
\begin{align*}
 \mathsf G_0&\le Cb_+L^{4K},\\
 \mathsf G_t&\le C\left[b_+\varepsilon_t^2L^{4(K-t)}
       +(\varepsilon_t^2/\alpha_0^2+e^{-d_*(t)})L^{3(K-t)}\right]
                                      &&(1\le t\le K),\\
 \mathsf H_t&\le CL^{3(K-t)} &&(0\le t\le K).
\end{align*}
Here $C$ depends only on the fixed construction constants,
and $b_+$ is the upper coupling bound used in
Proposition~\ref{C-350-8}. On the least torus and for fixed
$g$ these imply
\begin{equation}\label{eq:first-creation-count}
 \mathsf G_K\le C(g),\qquad
 \mathsf G_{t-1}+\mathsf H_{t-1}
       \le C_{\rm cnt}(g)L^{4(K-t+1)},\qquad C_{\rm cnt}(g)\ge1.
\end{equation}

Let $\mathsf{Cr}_t(w)$ be the contribution whose first
large-field region is created at step $t$, and set
$\mathsf w_t=\mathsf{Cr}_t(0)/Z[0]$. Then
\begin{align}
 Z[w]&=m_K^w(\mathcal V_K)+\sum_{t=1}^K\mathsf{Cr}_t(w),
                                           \label{eq:first-creation-split}\\
 \left|e^{-\mathcal E(w)}m_K^w(\mathcal V_K)/Z[0]\right|
        &\le e^{\|w\|_{\rm lw}\mathsf G_K},\nonumber\\
 \left|e^{-\mathcal E(w)}\mathsf{Cr}_t(w)/Z[0]\right|
        &\le e^{\|w\|_{\rm lw}(\mathsf G_{t-1}+\mathsf H_{t-1})}
                    \mathsf w_t,\label{eq:first-creation-complex}\\
 \sum_{t=1}^K L^{8(K-t)}\mathsf w_t&\le C_8(g)<\infty,
                                           \label{eq:first-creation-rarity}
\end{align}
uniformly in the datum at fixed $g$. All the sourced
quantities in these formulas are holomorphic on the
level-wise unit ball.
\end{proposition}

\begin{proof}
Put $r=\|w\|_{\rm lw}$. The amplitude part of the norm and
the coupling bounds give $|w|_\infty\le Cb_+r$.
There are at most $CL^{3K}$ bare plaquettes in the support
ball. At a passing output of level $t$, the background is
regular of that level at radius $2B_3\varepsilon_t$.
Consequently
\[
 0\le a_p(U_t)\le2B_3^2\varepsilon_t^2L^{-4t},\qquad
 |\mathcal O(w)(U_t)|
       \le Cb_+r\varepsilon_t^2L^{4(K-t)}.
\]
At $t=0$ use $a_p\le2$, giving $Cb_+rL^{4K}$.

We spell out the sum of centered local terms. The
regular-background estimate (2261), p.~827, applies to any
local pointed member $E$ at level $j\le t$:
\[
 |E(U_t)-E(\mathbf1)|\le\|E\|_X
 \left[C(\varepsilon_t/\alpha_0)^2L^{-4(t-j)}
             M^2(1+d_j(X))^2
         +2\mathbf1_{\{d_j(X)\ge d_*(j)\}}\right],
\]
where
\[
 d_*(j)=\min\left\{
   \frac{c_A\alpha_0L^{2(t-j)}}{2MB_3\varepsilon_t},
   \frac{L^{m+K-j}-4}{2M}\right\}-1.
\]
Apply it to $E=\Delta F_j[w]$ and use
\eqref{eq:levelwise-family-bound}. Locality makes this
member zero unless $X$ meets a source plaquette. Anchor
such a domain at a level-$j$ cube meeting the support ball;
there are at most $CL^{3(K-j)}$ choices. The rooted-domain
sum with rate $\kappa+1$ absorbs the base-point count and
the polynomial $(1+d)^2$, and its tail beyond $d_*(j)$
is at most $Ce^{-d_*(j)}$. These are the same rooted sums
used in (2262), p.~827. Since
\[
 d_*(j)\ge L^{t-j}(d_*(t)+1)-1,
 \qquad
 L^{3(t-j)}e^{-d_*(j)}\le L^{-(t-j)}e^{-d_*(t)}
\]
when $d_*(t)\ge3$, the two level sums are bounded by
\[
 CrL^{3(K-t)}
       (\varepsilon_t^2/\alpha_0^2+e^{-d_*(t)})
                    \sum_{j\le t}L^{-(t-j)}.
\]
This geometric sum is bounded independently of $t$ and $K$,
proving the stated $\mathsf G_t$ bound. For the energy tail,
the same anchoring and the uncentered family bound give
\[
 |\mathcal E-\mathcal E_t|
       \le Cr\sum_{j>t}L^{3(K-j)}\le CrL^{3(K-t)}.
\]
This includes $t=0$; at $t=K$ the tail is zero.

The radius condition is available at every datum used in
Proposition~\ref{C-350-8}. Decrease $\gamma_{\rm C}$ so that
\[
 \varepsilon(x)\le c_A\alpha_0/(8MB_3)
                      \quad(x\ge\gamma_{\rm C}^{-2}).
\]
The least-torus choice has $L^m\ge LM$, and $L\ge13$,
$M\ge1$ imply $L^{m+K-t}\ge LM\ge8M+4$. Both entries
of the minimum in \eqref{eq:first-creation-radius} are
therefore at least $4$, so $d_*(t)\ge3$. The first-passage
window bounds $x_K$ and $b_+$ in terms of $g$, and the
regularity radii are bounded by fixed constants throughout
the run. This proves \eqref{eq:first-creation-count}.
The dependence of $C_{\rm cnt}$ on $g$ is allowed here;
in particular it absorbs the factor $b_+$ in $\mathsf G_0$.

Let $g_t^{\rm w}$ be the sum of the nonnegative step factors
for the other outcomes, after the all-small-field prefix
of length $t-1$. These factors too are source independent.
Partition the histories according to their first such
outcome. Summing all later extensions leaves the prefix
mass, by the normalized test partition. This gives
\eqref{eq:first-creation-split} and the explicit formula
\[
 \mathsf{Cr}_t(w)=\int g_t^{\rm w}\,dm_{t-1}^w,
 \qquad \mathsf w_t=Z[0]^{-1}\int g_t^{\rm w}\,dm_{t-1}\ge0.
\]
This is the decomposition (2268), p.~829. Substitution of
\eqref{eq:levelwise-density}, followed by
\eqref{eq:first-creation-envelopes}, proves
\eqref{eq:first-creation-complex}. The all-small-field
zero-source mass is at most $Z[0]$, since it is one
nonnegative summand of the partition. On each finite
lattice the factors are bounded and the input source
exponential is entire; integration and the preceding
finite holomorphic sums prove the holomorphy assertions.

Finally the event of first creation at $t$ is contained
in $\{Z_t\ne\emptyset\}$. Lemma~\YMMref{115.1}{805}
and the partition-function lower bound used immediately
after it give, exactly as in (2185), p.~805,
\[
 \mathsf w_t\le
  2M^{-3}N_t\exp\{-c'A_0^2(\log x_t)^{2\mathfrak p_0}
        +(C_{\rm maj}^{\Sigma}+2+C_{\rm low}(b_+))\ell_g^3\}.
\]
The canonical least torus satisfies the side assumptions of
that history estimate. Here $N_t=\ell_g^3L^{3(K-t)}$,
$x_t\ge L^{K-t}b^\vee$, and $b^\vee=g^{-2}-\bar B>1$
after decreasing $\gamma_{\rm C}$.
With $n=K-t$, the sum in \eqref{eq:first-creation-rarity}
is consequently bounded by a fixed multiple, depending on
$g$, of
\[
 \sum_{n\ge0}\exp\{11n\log L
       -c'A_0^2(\log b^\vee+n\log L)^{2\mathfrak p_0}\}.
\]
It converges because $2\mathfrak p_0>1$: from some index
depending only on $g$ and the fixed constants the negative
term exceeds $12n\log L$, and the remaining tail is
bounded by $\sum_nL^{-n}$. This gives $C_8(g)$ uniformly
in $K$. No uniform bound as $g$ tends to zero is required.
\end{proof}

\begin{lemma}[Packed source amplitudes]
\label{lem:packed-levelwise-source}
Let $u_\iota$ be real Lipschitz functions with disjoint
nonzero sets, $\|u_\iota\|_\infty\le A_\iota$, and
$\operatorname{Lip}u_\iota\le A_\iota/(\vartheta d)$.
For $w=(\sum_\iota\zeta_\iota u_\iota)\circ x$ and
$B=\max_\iota|\zeta_\iota|A_\iota$, one has
\[
 \|w\|_{\rm lw}\le B/\mathsf A_\star(d),\qquad
 \mathsf A_\star(d)=\min\left\{\sigma_{\rm src}b^\vee,
  \frac{\mu_\star}{K_{\rm L}^\vee}
   \frac{b^\vee}{(b^\wedge)^{1/2}}
             \left(\frac{\vartheta d}{4\sqrt3M}\right)^{1/2}\right\},
\]
where $b^\vee=g^{-2}-\bar B$ and $b^\wedge=b_+$.
\end{lemma}

\begin{proof}
Disjointness gives $|w|_\infty\le B$. The sum of the
weighted functions has Lipschitz constant at most
$2B/(\vartheta d)$: for points in different nonzero
sets, bound each value by its Lipschitz constant times
the distance to the complement of that set. Thus, with
$n=K-j$,
\[
 \operatorname{sw}_j(w)
     \le2B\min\{1,\sqrt3ML^{-n}/(\vartheta d)\}.
\]
The running-coupling bounds
$L^nb^\vee\le x_j\le L^nb^\wedge$ now imply
\[
 |v_j^w|_\infty\le B/b^\vee,\qquad
 x_j^{1/2}\operatorname{sw}_j(v_j^w)
 \le\frac{2B(b^\wedge)^{1/2}}{b^\vee}
       \min\{L^{n/2},\sqrt3ML^{-n/2}/(\vartheta d)\}
 \le\frac{2B(b^\wedge)^{1/2}}{b^\vee}
                       \left(\frac{\sqrt3M}{\vartheta d}\right)^{1/2}.
\]
The last inequality uses $\min(a,b)\le\sqrt{ab}$.
Insertion in \eqref{eq:levelwise-source-norm} proves the claim.
\end{proof}



\begin{corollary}[Packed sources in finitely many balls]
\label{cor:packed-source-coefficients}
Retain the numerical, canonical-run and actual-history hypotheses of
Propositions~\ref{prop:levelwise-source-density}
and~\ref{prop:first-creation-envelopes}, and the calibrated rarity
hypotheses of Proposition~\ref{prop:ward-fixed-volume-rarity}.
Work on the least torus at a first-passage datum at fixed $g$.
Let $d>0$, $0<\vartheta\le1$ and $\vartheta d\le1$.
Let $I$ have $q\ge1$ elements. For $i\in I$, let
$F_i=\sum_\iota h_{i,\iota}$ be a finite sum of real Lipschitz
functions whose nonzero sets, for all pairs $(i,\iota)$, are
pairwise disjoint. Suppose
\[
 \|h_{i,\iota}\|_\infty\le A_i,
 \qquad \operatorname{Lip}h_{i,\iota}\le A_i/(\vartheta d),
 \qquad A_i>0,
\]
and that every bump vanishes off the union of $N\ge1$ physical
balls of radius $1$. Use the amplitude
$\mathsf A_\star(d)$ of Lemma~\ref{lem:packed-levelwise-source},
and put
\[
 u_i=A_i/\mathsf A_\star(d),\qquad
 w_z=\Bigl(\sum_{i\in I}z_iF_i\Bigr)\circ x,
 \qquad \mathbb P=\{z:|z_i|u_i<1\text{ for every }i\}.
\]
Then
\begin{equation}\label{eq:packed-source-coefficients}
\begin{gathered}
 \Omega(z)=Z[w_z]e^{-\mathcal E(w_z)}/Z[0]
 \quad\text{is holomorphic on }\mathbb P,\\
 \Omega(0)=1,\qquad
 |[z^B]\Omega|\le Q_{N,q}(g)\prod_{i\in B}u_i
 \quad(\emptyset\ne B\subset I).
\end{gathered}
\end{equation}
Here $[z^B]\Omega=(\prod_{i\in B}\partial_{z_i})\Omega(0)$.
The number $Q_{N,q}(g)\ge1$ depends only on $N,q,g$ and the
fixed construction data, not on the cutoff, label, balls, packing,
$d$, $\vartheta$ or functions.
\end{corollary}

\begin{proof}
Lemma~\ref{lem:packed-levelwise-source} gives
$\|w_z\|_{\rm lw}\le\max_i|z_i|u_i$.
The bump index may include several functions with the same
parameter $z_i$; the same maximum bounds each of them.
Thus $w_z$ is in the level-wise ball on $\mathbb P$.
The finite lattice partition function is entire in its source,
the source energy is holomorphic on this ball, and $Z[0]>0$.
This proves holomorphy and $\Omega(0)=1$.

In the proof of Proposition~\ref{prop:first-creation-envelopes},
the support enters only through the number of bare plaquettes
and the number of level-$j$ cubes meeting it. For a union of
$N$ balls each count is at most $N$ times its one-ball bound.
The local-family bound is applied to the whole source; no
decomposition of that source in the level-wise norm is needed.
Consequently the proof gives envelopes
\[
 \mathsf G_t^{[N]}=N\mathsf G_t,
 \qquad \mathsf H_t^{[N]}=N\mathsf H_t,
 \qquad
 \mathsf G_K^{[N]}\le NC(g),\qquad
 \mathsf G_{t-1}^{[N]}+\mathsf H_{t-1}^{[N]}
       \le NC_{\rm cnt}(g)L^{4(K-t+1)}.
\]
Split $\Omega=\Omega_\emptyset+\sum_{t=1}^K\Omega_t$
by \eqref{eq:first-creation-split}. The bounds
\eqref{eq:first-creation-complex}, with these envelopes, give
\[
 |\Omega_\emptyset|\le e^{NC(g)}\quad\text{on }\mathbb P,
 \qquad
 |\Omega_t|\le e\mathsf w_t
 \quad\text{on }\{|z_i|u_i\mathsf{Cnt}_t<1\},
 \qquad
 \mathsf{Cnt}_t=NC_{\rm cnt}(g)L^{4(K-t+1)}\ge1.
\]
Cauchy's formula, separately in the variables of $B$, yields
\[
 |[z^B]\Omega_\emptyset|\le e^{NC(g)}u^B,
 \qquad
 |[z^B]\Omega_t|\le e\mathsf w_t\mathsf{Cnt}_t^{|B|}u^B
 \le e\mathsf w_t[NC_{\rm cnt}(g)L^4]^q
                      L^{4q(K-t)}u^B,
 \qquad u^B=\prod_{i\in B}u_i.
\]
Proposition~\ref{prop:ward-fixed-volume-rarity}, with exponent
$4q$ and $t=K$, bounds the sum of the last weights by
\[
 \sum_{t=1}^K L^{4q(K-t)}\mathsf w_t\le C_{4q}(g),
 \qquad
 C_{4q}(g)=4M^{-3}\ell_g^3
       e^{\mathsf B(X_g)\ell_g^3}\mathsf C_{\rm ag}(4q+1).
\]
Indeed $x_K\le X_g$, the function $\mathsf B$ in
\eqref{eq:ward-volume-constants} increases, and the remaining
factor $e^{-c'p_K^2/2}$ is at most $1$.
The displayed bound is finite at fixed $g$. Therefore one may take
\[
 Q_{N,q}(g)=e^{NC(g)}
       +e[NC_{\rm cnt}(g)L^4]^qC_{4q}(g).
\]
All factors have the asserted dependence.
\end{proof}


\subsection{The lattice spectral inequality at several steps}
\label{sec:spectral-steps}

The spectral restriction must be used at a fixed physical displacement
while the lattice spacing tends to zero. Its one-step form therefore
needs the following elementary extension.

\begin{lemma}\label{lem:spectral-steps}
Let \(T\) be a positive self-adjoint contraction and \(V\) a commuting
unitary operator. If
\[
 (1-T)^2\ge T(2-V-V^*),
\]
then, for every integer \(k\ge1\),
\[
 (1-T^k)^2\ge T^k(2-V^k-V^{*k}).
\]
For the lattice transfer operators of Proposition~\YMMref{99.3}{739},
this implies the five-point inequality
\[
 u(0)-4u(te_\mu)+u(2te_\mu)
             +u(te_\mu+te_\lambda)+u(te_\mu-te_\lambda)\ge0,
 \qquad t=ka,
\]
where \(u(w)=\nu(\overline{\theta F}\,\tau_wF)\) in the corresponding
reflection Hilbert space.

Conversely, suppose strongly commuting self-adjoint operators \(H\ge0\)
and \(P\) have joint spectral measure \(\mathcal P\). If, on a dense
set of vectors, the bounded quadratic form
\[
 1-4e^{-tH}+e^{-2tH}+e^{-tH}(e^{itP}+e^{-itP})
\]
is nonnegative for every \(t\) in a subset of \((0,\infty)\) with zero as a limit point,
then \(\mathcal P\{(E,q):E<|q|\}=0\).
\end{lemma}

\begin{proof}
At a point of the joint spectrum of \(T,V\), write \(t_0=e^{-x}>0\)
and \(v=e^{iy}\), where \(x\ge0\). The assumed inequality is
\(\sinh(x/2)\ge|\sin(y/2)|\).
The addition formulas, by induction, give
\[
 \sinh(kx/2)\ge k\sinh(x/2)
       \ge k|\sin(y/2)|\ge|\sin(ky/2)|.
\]
This is the desired inequality at \(t_0^k,v^k\); at \(t_0=0\)
it is immediate. Joint spectral calculus proves the operator
inequality. Taking its matrix element at \([F]\) and using the
translation identity in Lemma~\YMMref{99.1}{737} gives the five
displayed correlations.

For the continuum assertion put
\[
 f_t(E,q)=4e^{-tE}\bigl(\sinh^2(tE/2)-\sin^2(tq/2)\bigr).
\]
This is the real bounded function
\(1-4e^{-tE}+e^{-2tE}+2e^{-tE}\cos(tq)\).
Density extends the assumed positivity to the whole Hilbert space,
so \(\mathcal P\{f_t<0\}=0\). Choose a sequence \(t_j\downarrow0\)
from the indicated set. Outside the union of these countably many
null sets,
\[
 \frac{2}{t_j}\sinh(t_jE/2)
       \ge\frac{2}{t_j}|\sin(t_jq/2)|
\]
for every \(j\). Taking the ordinary numerical limit gives \(E\ge|q|\).
The spectral projection of the complementary set is therefore zero.
\end{proof}

\section{Lattice conventions and first-passage arithmetic}\label{sec:setup-complements}

\subsection{Lattice symmetries}
\label{sec:setup-lattice-symmetries}

The block averages must commute with the lattice symmetries used in
reflection arguments. Gauge covariance alone does not imply this fact.
Write $a_j=L^{-j}$ and $\Lambda_j=a_j\mathbb Z^3/\ell\mathbb Z^3$.
For $i\ge-m$, let $\mathsf{Iso}_i$ consist of the maps of $T_\ell$
induced by $x\mapsto Rx+v$, where $R$ is a signed permutation matrix
and $v\in a_i\mathbb Z^3$. Such a map acts on oriented bonds and paths
by acting on their vertices, and on a configuration by
$(\mathsf rU)(b)=U(\mathsf r^{-1}b)$. In the following lemma every
lattice and every standard cube partition entering an object has index
at least $i$.

\begin{lemma}\label{lem:setup-lattice-covariance}
Let $\mathsf r\in\mathsf{Iso}_i$. It preserves the periodic Euclidean,
sup and $\ell^1$ distances. It carries blocks, corridors, pivots,
coordinate-order contour families and determining regions to those of
their images. For real configurations, it preserves plaquette defect
norms, the Wilson action and product Haar measure, and
\[
 \mathcal M^n(\mathsf rU)=\mathsf r\mathcal M^n(U),\qquad
 \mathbf Q^{\rm lin}(\mathsf rA)
       =\mathsf r(\mathbf Q^{\rm lin}A),\qquad
 Q(\mathsf rA)=\mathsf r(QA).
\]
The identities hold on the corresponding domains for partial
configurations and for iterated linear averages. The block lifts and
the determining-data maps commute with $\mathsf r$ as well.

The map carries constrained critical regular configurations onto those
of the transported datum, with the same regularity radius. Thus
\[
 \mathcal C(\mathsf r\mathbf B,\mathsf rV;r)
       =\mathsf r\mathcal C(\mathbf B,V;r).
\]
Whenever these sets consist of one residual gauge orbit, the background
orbits obey
\[
 U_{\mathsf r\mathbf B}(\mathsf rV)
          =\mathsf rU_{\mathbf B}(V).
\]
This includes constant-torus and localized background orbits. It does
not identify representatives chosen in a fixed axial frame: that frame
must also be transported, except for translations.
\end{lemma}

\begin{proof}
The matrix $R$ preserves each of the three norms and each lattice
$a_j\mathbb Z^3$. For $j\ge i$, $a_i=L^{j-i}a_j$, so translation by
$v$ preserves that lattice too. Both operations preserve the period
lattice. They therefore descend to $T_\ell$ and preserve its periodic
distances. The cubes are sup balls centered at lattice sites. The
definitions of blocks, corridors, determining regions, their
enlargements and their local block sides use these cubes, distances
and set operations; these definitions are consequently preserved.
The pivot is the unique bond of the center-to-center straight path
in the corridor. Coordinate permutations and changes of signs carry
coordinate-order paths bijectively onto coordinate-order paths.

For every path $\Gamma$,
$(\mathsf rU)(\mathsf r\Gamma)=U(\Gamma)$, since the ordered products
have identical factors. Returning a transported plaquette to the
chosen base point and orientation replaces its holonomy by a
conjugate of the original holonomy or its inverse. Conjugation and
inversion preserve the defect norm for unitary matrices; they
preserve the SU(2) trace too. This proves the assertions about the
action and the defects. On the independent positive-bond variables,
$U\mapsto\mathsf rU$ is a permutation followed by inversion in some
factors, so it preserves product Haar measure.

The contour families defining the nonabelian means are merely
renumbered. Their means, their diameters and their domains of
definition are therefore unchanged. The same is true of the summands
in the block average when the orientation of the coarse bond is
preserved. To check orientation reversal, write $T=U[c]$ for the
straight holonomy of a positive coarse bond $c$. The loop variables
$W_x$ of the small-field formula and $Y_{x,\Gamma,\Gamma'}$ of the
second real formula satisfy, at the reversed bond,
\[
 W^{\bar c}_{x'}=T^{-1}W_x^{-1}T,\qquad
 Y^{\bar c}_{x',\Gamma',\Gamma}
       =T^{-1}Y_{x,\Gamma,\Gamma'}^{-1}T,
 \quad x'=x+(c_+-c_-).
\]
Thus the same one of the three real branches in
equations~\YMMref{(37)--(38)}{20} applies. In the small-field branch,
$\log W^{\bar c}_{x'}=-T^{-1}(\log W_x)T$; the intrinsic logarithm
has the same identity in the second branch. Summing and exponentiating
gives the inverse of the average at $c$. The third branch is the
straight holonomy and has the same property. Hence the defining
formula works at either orientation, and
$\mathcal M(\mathsf rU)=\mathsf r\mathcal M(U)$. Iteration and
restriction prove the assertions for partial fields and determining
data. The additive computation with path sums proves covariance of
$\mathbf Q^{\rm lin}$ and $Q=L^{-1}\mathbf Q^{\rm lin}$. A block lift
is the identity on intra-block bonds and the coarse value on corridor
bonds, so it is transported by the same geometric bijections.

For completeness, transport of the current uses the tangent action,
not just the signed pullback of a one-form. At positive image bonds
$b'$ define
\[
 (\mathsf r_U A)(b')=
 \begin{cases}
 A(b),&\mathsf r^{-1}b'=b,\ b\text{ positive},\\
 -\operatorname{Ad}(U(b)^{-1})A(b),
       &\mathsf r^{-1}b'=\bar b,\ b\text{ positive}.
 \end{cases}
\]
Then
$\mathsf r(e^{it\xi A}U)
 =e^{it\xi\mathsf r_U A}\mathsf rU$.
The map $\mathsf r_U$ is an isometry for the trace pairing and for the
bondwise norms. Differentiate action invariance and use
Proposition~\YMMref{9.2}{38}; nondegeneracy of the pairing gives
$J(\mathsf rU)=\mathsf r_UJ(U)$. The plaquette and current bounds in
Definition~\YMMref{9.3}{40} are therefore preserved, including their
local block-side factors.

Covariance of the determining-data map carries each constraint fiber
onto the transported fiber. It carries differentiable curves in one
fiber onto those in the other, preserving their actions; hence it
preserves criticality. It also carries the nodes and residual gauge
group onto their images, since
$\mathsf r(U^u)=(\mathsf rU)^{\mathsf ru}$. This proves the assertion
about regular critical sets and then about their gauge orbits. The
localized background in equation~\YMMref{(77)}{42} uses precisely the
transported tower, lift and determining datum. Finally, coordinate
paths chosen in a frame are transported to paths in the transported
frame, which explains the qualification about representatives.
\end{proof}

\subsection{Elementary first-passage estimates}
\label{sec:setup-first-passage-arithmetic}

We use the canonical sequence and the constants of
Section~\YMMref{74}{614}. In particular,
\[
 x_j=L^{-1}x_{j-1}-b_j,\quad |b_j|\le B_\sharp,\quad
 \bar B=\frac L{L-1}B_\sharp,\quad
 X_g=L(g^{-2}+B_\sharp),\quad
 I(g)=(g^{-2}+\bar B,X_g-\bar B].
\]
The label is $\mathsf b=L^{-K}\beta$, the comparison line is $\theta_j=L^{K-j}\mathsf b$, and the depth at level $k$ is $s=K-k$. The side floor is the integer in equation~\YMMref{(1528)}{615}:
$\mathfrak L(g)=\{2L^m:m\ge m_{\mathfrak L}(g)\}$.
The following arithmetic also records the pressure constant used
when the sequence is stopped. Choose a fixed
$C_{\rm st}\ge\max\{c_{\rm low},0\}$, where $c_{\rm low}$ is the
constant of Proposition~\YMMref{74.4}{617}, and put
\begin{equation}\label{eq:setup-pressure-constants}
 C_g=C_{\rm st}+\frac92\log X_g,\qquad
 C_{\rm fp}^{\rm st}
   =\max\left\{\frac92,C_{\rm st}+\frac92\log\frac{3L}{2}\right\}.
\end{equation}
These constants precede the choice of the coupling, cutoff and side.
The upper-bound constant $C_{\rm S}$ of Theorem~\YMMref{AD}{607}
may and will be chosen at least $C_{\rm fp}^{\rm st}$.

\begin{lemma}\label{lem:setup-pressure-stability}
For every first-passage datum $(K,\beta)$ at $0<g\le\gamma$ and
every $\ell\in\mathfrak L(g)$, with the normalizing energy of
Theorem~\YMMref{AD}{607},
\[
 -C_g\le \ell^{-3}\bigl(\log Z_{K,\ell,\beta}
                      -E_{K,\ell,\beta}\bigr)\le C_{\rm S},
 \qquad
 C_{\rm low}(x_K)\le C_g
     \le C_{\rm fp}^{\rm st}(1+\log g^{-2}).
\]
\end{lemma}

\begin{proof}
The zero-source factorization in Theorem~\YMMref{AD}{607}
uses the accumulated vacuum constant $E_{(K)}$, which is
$E_{K,\ell,\beta}$ in the definition preceding
Proposition~\YMMref{74.4}{617}. Its local source energies vanish
at zero. Thus both results use exactly
$\widetilde Z(0)=e^{-E_{K,\ell,\beta}}Z_{K,\ell,\beta}$.
Proposition~\YMMref{74.1}{615} verifies the sequence and room
hypotheses of Proposition~\YMMref{74.4}{617}, which gives
\[
 \widetilde Z(0)\ge
 e^{-[c_{\rm low}+(9/2)\log x_K]\ell^3}.
\]
Since $x_K\le X_g$ and $c_{\rm low}\le C_{\rm st}$, this is the
claimed lower bound. The upper bound is the zero-source case of
Theorem~\YMMref{AD}{607}. Finally
$B_\sharp\le\frac12g^{-2}$ gives $X_g\le(3L/2)g^{-2}$, so
\[
 C_g\le C_{\rm st}+\frac92\log\frac{3L}{2}
                     +\frac92\log g^{-2}
       \le C_{\rm fp}^{\rm st}(1+\log g^{-2}).
\]
All constants used to choose $C_{\rm S}$ in this argument are fixed
before the coupling and the source radius, as in that theorem.
\end{proof}

\begin{lemma}\label{lem:setup-stopped-arithmetic}
Let $0<g\le\gamma$, let $(K,\beta)$ be a first-passage datum at $g$,
and assume the conditions in equation~\YMMref{(1523)}{614}.
For $0\le k\le K$,
\[
 K-k\le\frac{\log x_k}{\log L},\qquad
 \bar B\le\frac{13}{24}g^{-2}.
\]
For $1\le i\le K-1$, put $k=K-i$ and
$g_{(i)}^{-2}=L^{-1}x_k-B_\sharp$. Then $X_{g_{(i)}}=x_k$, and,
with $C_h=C_{\rm st}+\frac92\log X_h$,
\begin{equation}\label{eq:setup-stopped-pressure}
 C_{g_{(i)}}\le C_g+\frac92i\log L+\frac18
                    \le C_g+5i\log L.
\end{equation}
The prefix ending at $k$ is a first-passage sequence at a reference
coupling $g'\le\gamma$ exactly when
\begin{equation}\label{eq:setup-stopped-window}
 L^{-1}x_k-B_\sharp\le g'^{-2}<x_k+b_k-B_\sharp.
\end{equation}
For $i\ge1$ every such $g'$ is smaller than $g$.

If $p_0(g_j)^2=A_0^2y_j^{2p}$, where $y_j=\log x_j$ and $p\ge1$,
then for $0\le t\le k\le K$,
\begin{equation}\label{eq:setup-threshold-convexity}
 p_0(g_t)^2-p_0(g_k)^2
 \ge2pA_0^2y_k^{2p-1}(k-t)\log(L/2).
\end{equation}
\end{lemma}

\begin{proof}
Iterating the recurrence gives equation~\YMMref{(1530)}{615}.
In particular
\[
 L^{K-k}(x_K-\bar B)+\bar B
       \le x_k\le L^{K-k}(x_K+\bar B)-\bar B.
\]
Here $x_K>g^{-2}\ge\gamma^{-2}$, and
$L^{-1}\gamma^{-2}-B_\sharp>1$ implies
$\gamma^{-2}-\bar B>1$. Thus $x_k>L^{K-k}$, proving the depth
bound. Since $g^{-2}\ge2B_\sharp$ and $L\ge13$,
$\bar B\le\frac{L}{2(L-1)}g^{-2}\le\frac{13}{24}g^{-2}$.

The definition of $g_{(i)}$ gives $X_{g_{(i)}}=x_k$. Since
$x_K\le X_g$, the upper affine bound gives
$x_k\le L^i(X_g+\bar B)$. Moreover
$X_g=L g^{-2}+(L-1)\bar B\ge3(L-1)\bar B\ge36\bar B$.
Consequently
\[
 \frac92\log x_k
 \le\frac92\log X_g+\frac92i\log L
          +\frac92\log(1+\bar B/X_g)
 \le\frac92\log X_g+\frac92i\log L+\frac18.
\]
Adding $C_{\rm st}$ proves the first inequality in
\eqref{eq:setup-stopped-pressure}; the second follows from
$1/8\le\frac12i\log L$.

The sequence is strictly decreasing. Hence its prefix first crosses
$X_{g'}$ at $k$ exactly when $x_k\le X_{g'}<x_{k-1}$.
Using $x_{k-1}=L(x_k+b_k)$ gives
\eqref{eq:setup-stopped-window}. If $k<K$, then
$x_k\ge x_{K-1}>X_g$, so the lower endpoint
$L^{-1}x_k-B_\sharp$ exceeds $g^{-2}$.

Finally $x_j\ge\gamma^{-2}\ge2B_\sharp$ and the recurrence give
$x_{j-1}=L(x_j+b_j)\ge(L/2)x_j$. Hence
$y_t-y_k\ge(k-t)\log(L/2)$. Apply convexity of $y\mapsto y^{2p}$
on the positive real axis and multiply by $A_0^2$ to obtain
\eqref{eq:setup-threshold-convexity}.
\end{proof}

\begin{lemma}\label{lem:setup-label-coverage}
Assume the constants of Section~\YMMref{74}{614}. Every
$\lambda>\gamma^{-2}+\bar B$ belongs to $I(h)$ for
\[
 h^{-2}=\max\{\gamma^{-2},L^{-1}(\lambda+\bar B)-B_\sharp\}.
\]
If $\lambda\ge2\gamma^{-2}$, one may instead take
$h^{-2}=\lambda/2$. For either choice, $h\le\gamma$ and
$(K,L^K\lambda)$ is a first-passage datum at $h$ for every $K\ge1$.
Also, $0<h\le g\le\gamma$ implies
$\mathfrak L(h)\subset\mathfrak L(g)$.
\end{lemma}

\begin{proof}
The definition of the maximum gives $\lambda\le X_h-\bar B$.
If the maximum is $\gamma^{-2}$, then
$h^{-2}+\bar B<\lambda$ by hypothesis. In the other case,
using $\bar B-B_\sharp=L^{-1}\bar B$, we have
$h^{-2}+\bar B=L^{-1}(\lambda+2\bar B)<\lambda$, since
$L\ge3$ and $\lambda>\bar B$. Thus $\lambda\in I(h)$.
For the second choice,
$\bar B<\gamma^{-2}\le\lambda/2$ gives
$h^{-2}+\bar B<\lambda$, whereas
$X_h-\bar B=(L/2)\lambda+(L-2)\bar B\ge\lambda$.
The first-passage assertion is Corollary~\YMMref{74.2}{616}.
Finally $R$ is nondecreasing, so the explicit side floor in
equation~\YMMref{(1528)}{615} is nondecreasing as the inverse square
coupling increases. This proves the inclusion of side families.
\end{proof}

\subsection{Uniqueness from real values}

\begin{lemma}\label{lem:setup-real-identity}
A holomorphic function on a connected open subset of $\mathbb C^N$
that vanishes on a nonempty relatively open set of real points
vanishes throughout the domain.
\end{lemma}

\begin{proof}
At a point in that real open set, all real derivatives vanish.
They are the coefficients of the complex power series, so the
function vanishes on a complex open ball. The set of points having
a neighborhood on which the function vanishes is open and nonempty.
It is closed in the domain too: at a limit point every derivative
vanishes by continuity, so the power series is zero there as well.
Connectedness therefore proves the assertion.
\end{proof}

\section{The cubic orbital remainder}\label{sec:ward-orbital-complement}

The orbital insertion contains a discrete derivative of a quadratic
potential. Subtracting that derivative leaves a cubic local function. We
give the calculation, including the uniformity needed on complex fields.
The constants below are fixed numbers, independent of the lattice and of
its spacing.

Write $a$ for the lattice spacing, $p=p_{\mu\nu}(x)$, and
$p_\pm=p\pm ae_\rho$. Put $a_q=1-\frac12\operatorname{tr}\partial U(q)$
and $\mathbf w(V)=(V-V^{-1})/(2i)$. For $\rho$ in the plane of $p$, let
$R=p\cup p_+$ be the rectangle consisting of the two plaquettes, and set
\begin{equation}\label{eq:ward-orbital-potential}
 \Psi_\rho(p)=\tfrac12a_R-a_p-a_{p_+}+\tfrac12a_pa_{p_+}.
\end{equation}
For $\rho$ normal to $p$, let $W_p^y$ be the holonomy of $p$ from its
corner $y$, with the orientation of $p$, and set
\begin{equation}\label{eq:ward-orbital-normal-potential}
 \Psi_\rho(p)=-\tfrac1{16}\sum_{y\text{ a corner of }p}
 \left\langle \mathbf w(W_p^y)-\operatorname{Ad}(U(y,y+ae_\rho))
      \mathbf w(W_{p_+}^{y+ae_\rho}),
 \mathbf w(W_p^y)-\operatorname{Ad}(U(y,y+ae_\rho))
      \mathbf w(W_{p_+}^{y+ae_\rho})\right\rangle .
\end{equation}
The pairing is $\langle X,Y\rangle=\frac12\operatorname{tr}(XY)$; on
complex fields it remains bilinear. With the bond directions and
densities of the Ward identity, put
\begin{equation}\label{eq:ward-orbital-remainder}
 \mathfrak r_\rho(p)=\sum_{b\subset\partial p}
  \langle\mathbf u_{b,\rho},\mathbf w(W_{b,p})\rangle
  -\tfrac12(a_{p_+}-a_{p_-}),\qquad
 \mathfrak r^\sharp_\rho(p)=\mathfrak r_\rho(p)
             -\Psi_\rho(p)+\Psi_\rho(p_-).
\end{equation}
These definitions are local and therefore apply on the infinite lattice
as well as on a torus on which the two cubes adjacent to $p$ have distinct
vertices.

\begin{proposition}\label{prop:ward-orbital-cubic}
The following statements hold.
\begin{enumerate}
\item[(a)] For real configurations, $-(a_p+a_{p_+})\le\Psi_\rho(p)\le0$.
If $\rho$ is in the plane of $p$, then $\mathfrak r^\sharp_\rho(p)=0$.
For a normal $\rho$, let $\bar a(p)$ be the largest plaquette action on
the eleven faces of the two cubes adjacent to $p$. There is a fixed
$C_{\rm orb}<\infty$ such that
\begin{equation}\label{eq:ward-orbital-real-cubic}
 |\mathfrak r^\sharp_\rho(p)|\le C_{\rm orb}\bar a(p)^{3/2}.
\end{equation}
In particular, if their plaquette matrices are within $\mu$ of the
identity, then $|\mathfrak r^\sharp_\rho(p)|\le c_{\rm orb}\mu^3$,
where $c_{\rm orb}=2^{-3/2}C_{\rm orb}$.
\item[(b)] The remainder is a gauge invariant polynomial in the twenty
bond matrices of those two cubes. On $SL(2,\mathbb C)$ configurations
with bond norms at most $2$, it has a fixed finite height. There is a
fixed $c_\sharp$ such that, if all eleven plaquette matrices are within
$r\le1$ of the identity, then
\begin{equation}\label{eq:ward-orbital-complex-cubic}
 |\mathfrak r^\sharp_\rho(p)|\le c_\sharp r^3.
\end{equation}
\item[(c)] If $\mathsf r e_\rho=s_\rho e_{\rho'}$ is the linear part of a
lattice isometry, then
\[
 \mathfrak r^\sharp_{\rho'}(\mathsf rp)(\mathsf rU)
       =s_\rho\mathfrak r^\sharp_\rho(p)(U).
\]
Every real torus mean of $\mathfrak r^\sharp_\rho(p)$ is zero.
For every vector weight $\eta$,
\begin{equation}\label{eq:ward-orbital-summation}
 \mathfrak K^{\rm v}[\eta]
 =-\sum_{p,\rho}(\nabla^+_\rho\eta^\rho)(p)\Psi_\rho(p)
   +a^{-1}\sum_p\eta^{\rho(p)}(p)\mathfrak r^\sharp_{\rho(p)}(p).
\end{equation}
\end{enumerate}
\end{proposition}

\begin{proof}
For a common bond of two coplanar plaquettes, put
$V_1=W_y(\nu',\rho)$ and $V_2=W_y(\nu',-\rho)$, with the signs fixed
by their traversal from that bond. Cancellation of the common bond
gives $\frac12\operatorname{tr}(V_2^{-1}V_1)=1-a_R$. The $SU(2)$
product formula consequently gives
\[
 \langle\mathbf w(V_1),\mathbf w(V_2)\rangle
 =a_p+a_{p_+}-a_pa_{p_+}-a_R.
\]
Using $|\mathbf w(V)|^2=a_V(2-a_V)$ in
\eqref{eq:ward-orbital-potential}, we obtain
\begin{equation}\label{eq:ward-orbital-inplane}
 \Psi_\rho(p)=-\tfrac14|\mathbf w(V_1)+\mathbf w(V_2)|^2
              -\tfrac14(a_p^2+a_{p_+}^2).
\end{equation}
The two nonzero bond terms, after subtracting
$\tfrac12(a_{p_+}-a_{p_-})$ as in
\eqref{eq:ward-orbital-remainder}, give
$\Psi_\rho(p)-\Psi_\rho(p_-)$; the other two have
$\mathbf u_{b,\rho}=0$. This proves the in-plane identity. The inequality
$|u+v|^2\le2|u|^2+2|v|^2$ in the last display proves its two sign bounds.
For a normal direction, the same inequality in
\eqref{eq:ward-orbital-normal-potential} bounds each of its four terms
by $2[a_p(2-a_p)+a_{p_+}(2-a_{p_+})]$; this gives the asserted bounds
there too.

We next prove the cubic cancellation. All twenty bonds can be identified
with the positive edges of
$\{0,1\}^2\times\{-1,0,1\}$. Thus the normal remainder is one universal
function $\Phi$ of twenty matrices. Inverses are polynomial on
$SL(2,\mathbb C)$, since $V^{-1}=(\operatorname{tr}V)\mathbf1-V$.
The displayed formulas show that $\Phi$ is polynomial and gauge
invariant. Write $g(X)=\Phi((e^{iX_b})_b)$.

Orient the four edges of the central plaquette positively, and let
$\epsilon_b=1,1,-1,-1$ be their incidence signs around its boundary.
For $U_t(b)=e^{itX_b}$, write $f_0=X(\partial p)$ and
$f_\pm=X(\partial p_\pm)$ for the unnormalised oriented sums of $X$
around these plaquettes, with no factor of the lattice spacing.
The side plaquette sums $h^\pm(b)$ obey
\[
 \sum_{b\subset\partial p}\epsilon_bh^+(b)=f_0-f_+,
 \qquad
 \sum_{b\subset\partial p}\epsilon_bh^-(b)=f_0-f_-:
\]
in both sums the vertical bond contributions cancel pairwise, and the
remaining bonds are those of the two parallel faces. Expanding the
holonomies to first order gives
\[
 \mathbf u_{b,\rho}=-\tfrac t2(h^+(b)-h^-(b))+O(t^2),
 \quad \mathbf w(W_{b,p})=t\epsilon_bf_0+O(t^2),
 \quad a_{p_\pm}=\tfrac{t^2}{2}|f_\pm|^2+O(t^3).
\]
Consequently
\begin{align*}
 \mathfrak r_\rho(p)(U_t)
 &=t^2\left[\tfrac12\langle f_0,f_+-f_-\rangle
              -\tfrac14(|f_+|^2-|f_-|^2)\right]+O(t^3)\\
 &=t^2\left[-\tfrac14|f_+-f_0|^2
                   +\tfrac14|f_0-f_-|^2\right]+O(t^3),
\end{align*}
whereas the two potentials have quadratic terms
$-t^2|f_+-f_0|^2/4$ and $-t^2|f_0-f_-|^2/4$.
They cancel in $g(tX)$. Thus $g$ has no terms of degrees at most two.
This computation is polynomial in the coordinates of $X$ and therefore
holds for complex $X$ with the bilinear pairing as well.

For the real estimate, root the tree at the lowest corner $x$ and take
the central edges $b_1=\langle x,x+ae_\mu\rangle$,
$b_2=\langle x+ae_\mu,x+ae_\mu+ae_\nu\rangle$,
$b_4=\langle x,x+ae_\nu\rangle$, together with all eight vertical edges.
Its eleven edges join
the twelve vertices without a cycle. Gauge transformation along its
paths makes every tree bond equal to the identity. The fourth central
edge is then the inverse of the central plaquette holonomy. If $b^\pm$
is a translate of a central edge $b$, the corresponding side plaquette
gives
\[
 U(b^\pm)=W_y(\alpha,\pm\rho)^{-1}U(b).
\]
Unitarity and $|\partial U(q)-\mathbf1|^2=2a_q$ imply that all twenty
bond matrices are within $2(2\bar a(p))^{1/2}$ of the identity.
Choose their logarithms with
$|X_b|\le\min\{\pi,\frac\pi2|U(b)-\mathbf1|\}$.
On $\max_b|X_b|\le\pi$, the third derivatives of the fixed function $g$
are bounded by a fixed $C_3$. Taylor's formula and the cancellation
above now give
\[
 |\Phi(U)|\le\tfrac16C_3\max_b|X_b|^3
 \le\tfrac16C_3\pi^3(2\bar a(p))^{3/2}.
\]
This proves \eqref{eq:ward-orbital-real-cubic}, for example with
$C_{\rm orb}=C_3\pi^3 2^{3/2}/6$. The statement in terms of $\mu$
uses $\bar a(p)\le\mu^2/2$.

The same tree and the same Taylor cancellation prove the complex
estimate; only the constants in the gauge bounds change. Its paths
from the root have at most three edges. For an $SL(2,\mathbb C)$ matrix
the operator norms of the matrix and its inverse coincide, so both are
at most $2$ under the hypothesis. A plaquette holonomy based at another
corner, or conjugated by one of these tree paths, therefore differs
from the identity by at most a fixed multiple of $r$. After the tree
gauge, the fourth central edge is the inverse of such a holonomy; the
displayed side-face identity bounds each of the other eight edges by
$C r$ from the identity when $r\le r_0$, for fixed positive $C,r_0$.
Their principal logarithms are then bounded by $C' r$, after reducing
$r_0$ if necessary. Taylor's integral formula on a fixed complex ball
gives $|g(X)|\le C''\max_b|X_b|^3$, because its first three homogeneous
terms vanish. This proves \eqref{eq:ward-orbital-complex-cubic} for
$r\le r_0$. On the original bond domain $|U(b)|\le2$, the polynomial
$\Phi$ has a fixed height $H$ by compactness; for $r_0\le r\le1$ use
$H\le Hr_0^{-3}r^3$. This completes (b).

For covariance, reversing an orientation changes both vector parts in
each paired bond term by the same sign. Transport to another base point
conjugates both factors. Hence the bond calculation for the Ward
densities gives
$\mathfrak r_{\rho'}(\mathsf rp)(\mathsf rU)
=s_\rho\mathfrak r_\rho(p)(U)$.
In the potential the paired factors are unchanged if their order or
their common orientation is reversed. Its two faces are interchanged
when the normal direction is reversed. It follows directly from
\eqref{eq:ward-orbital-potential} and
\eqref{eq:ward-orbital-normal-potential} that
\[
 \Psi_{\rho'}(\mathsf rp)(\mathsf rU)=
 \begin{cases}\Psi_\rho(p)(U),&s_\rho=1,\\
               \Psi_\rho(p_-)(U),&s_\rho=-1.
 \end{cases}
\]
Subtracting the adjacent potentials proves (c). Reflection in the plane
through the centre of $p$ normal to $\rho$ fixes $p$, reverses the
remainder, and preserves the real torus measure; its mean is therefore
zero. Finally insert \eqref{eq:ward-orbital-remainder} in the orbital
sum and change $p$ to $p+ae_\rho$ in its translated potential term.
This gives \eqref{eq:ward-orbital-summation}.
\end{proof}

\section{Test pairs and the fluctuation template}\label{sec:ward-kernel-complements}

We record two consequences of the free-kernel comparison which are used
below in forms not included in its statement. Distances in this section
are Euclidean distances and their periodic counterparts.

\begin{lemma}\label{lem:ward-test-pairs}
Use the test pairs of scale $d$ and shape $\vartheta\in(0,1)$ defined
before Theorem~\YMMref{C}{7}: each function is nonnegative and nonzero,
is supported in a ball of radius $d$, has integral at least
$\vartheta d^3$ times its supremum, and has Lipschitz constant at most
its supremum divided by $\vartheta d$; the two supports have distance
between $d$ and $6d$. On a torus require $20d\le\ell$.
\begin{enumerate}
\item[(a)] Both supports lie in one ball of radius $5d$, and every
cross-distance lies in $[d,10d]$. Decreasing $\vartheta$ preserves the
test-pair conditions. If $c\ge\vartheta d$, then
$\max\{\|f\|_\infty,c\operatorname{Lip}f\}
\le c(\vartheta d)^{-1}\|f\|_\infty$, and similarly for $h$.
\item[(b)] A pair and its lift to a Euclidean ball of radius at most
$\ell/4$ have the same suprema, Lipschitz constants, integrals and
cross-distances. Translation and periodisation inside such a ball
preserve its test-pair properties. Dilation by $t>0$, on $\mathbb R^3$
or between the corresponding tori, takes $(f,h)$ to
$(f(\cdot/t),h(\cdot/t))$ at scale $td$. For
$\mathcal K_q(f,h)=\iint f(x)h(y)|x-y|^{-q}\,dx\,dy$,
\begin{equation}\label{eq:ward-test-scaling}
 \mathcal K_q(f(\cdot/t),h(\cdot/t))=t^{6-q}\mathcal K_q(f,h).
\end{equation}
On the family of admissible tori we use $t=L^z$, $z\ge0$.
\item[(c)] There are smooth test pairs for every $\vartheta\in(0,1)$,
including a pair $(f\circ\theta,f)$ mirrored through $x_3=0$, with $f$
supported in the ball of radius $d$ about $2de_3$. A concrete $C^1$
choice, when $\vartheta\le3\sqrt3/8$, is
$f_x(z)=(1-|z-x|^2/d^2)_+^2$ and $f_y$ with
$3d\le|x-y|\le8d$.
\end{enumerate}
\end{lemma}

\begin{proof}
The centres of the supporting balls are at most $8d$ apart: approximate
the minimum distance of the supports and add their two radii. Their
midpoint is within $4d$ of each centre. This proves the common-ball
claim and the upper bound $10d$; the lower bound is the separation
hypothesis. The assertions about $\vartheta$ and the norm follow
directly from its two inequalities.

If two points lie in a Euclidean ball of radius $\ell/4$, their
difference has each coordinate of absolute value at most $\ell/2$;
it is a least periodic lift. Projection is consequently an isometry
on that ball. For a function supported there, its extension by zero
has the same Lipschitz constant before and after projection. Indeed,
if one point is outside the ball, follow a shortest segment from the
inside point until its first boundary point, where the function is
zero; that segment bounds the difference by its Lipschitz constant
times the original distance. The integral and supremum are unchanged
by the injective projection. Dilation multiplies distances by $t$,
Lipschitz constants by $t^{-1}$ and integrals by $t^3$, proving all
the test conditions and, by changing both integration variables,
\eqref{eq:ward-test-scaling}.

For (c), let $x=2de_3$ and choose $r<1$ with
$r>\max\{\vartheta,(3\vartheta/\pi)^{1/3}\}$. Smooth the cone
$g(z)=(1-|z-x|/(rd))_+$ by a nonnegative mollifier of integral one
supported in the ball of radius $\epsilon d$, with
$0<\epsilon<\min\{1-r,r-\vartheta\}$. The resulting $f$ satisfies
\[
 \|f\|_\infty\le1,\quad f(x)\ge1-\epsilon/r,
 \quad\operatorname{Lip}f\le(rd)^{-1},\quad
 \int f=(\pi/3)r^3d^3.
\]
Thus its mass and Lipschitz inequalities hold at shape $\vartheta$.
Its support lies in the prescribed ball. The reflected support has
distance at least $2d$ and at most $4d$, since both centres belong to
the respective supports. Periodisation is covered by (b).
For the explicit polynomial bump, radial integration and differentiation
give $\|f_x\|_\infty=1$, $\int f_x=32\pi d^3/105$ and
$\operatorname{Lip}f_x=8/(3\sqrt3 d)$. The two supporting balls have
distance $|x-y|-2d$, proving its assertion.
\end{proof}

\begin{lemma}\label{lem:ward-fluctuation-template}
The free-kernel comparison of Proposition~\YMMref{121.2}{824} also holds
for $k=C^{\rm fl}_{K,s}$ when $a_s\ge d$, with the additional error
$e=(d/a_s)^3$. Its constant may be enlarged by a fixed number depending
only on the constants of the flat curvature kernels. In particular,
its statement for $k=C^{\rm fl}_{K,K-n}$ has
$e=(d/a_{K-n})^3$ whenever $a_{K-n}\ge d$.
\end{lemma}

\begin{proof}
The orthogonal decomposition~\YMMref{(2227)}{818} gives
$C^{\rm fl}_{K,s}=C^{\rm free}_K-C^{\rm bg}_{K,s}$.
The all-side background estimate of Proposition~\YMMref{120.5}{821}
is $|C^{\rm bg}_{K,s}(p,q)|\le C^{\rm bg}_{\rm tp}(a_K/a_s)^3$.
On the separated supports put $t=d/a_s\le1$. Thus the two kernels
differ by at most $C^{\rm bg}_{\rm tp}t^3a_K^3d^{-3}$, whereas
\YMMref{(2246)}{824} bounds the free kernel by
$(1+3C_{\rm lat})a_K^3d^{-3}$.
The squared-kernel estimate in the proof of
Proposition~\YMMref{121.2}{824}, obtained by applying
$|u^2-v^2|\le|u-v|(2|v|+|u-v|)$ to its plaquette sum, consequently gives
\[
 |\mathcal K[C^{\rm fl}_{K,s}]-\mathcal K[C^{\rm free}_K]|
 \le324C^{\rm bg}_{\rm tp}
       (2+6C_{\rm lat}+C^{\rm bg}_{\rm tp})t^3
       \|f\|_\infty\|h\|_\infty.
\]
Adding its free-kernel conclusion proves the assertion with the stated
fixed enlargement, under the same $8a_K\le d$ and $20d\le\ell$
hypotheses.
\end{proof}

When explicit kernel constants are needed, we use
\begin{equation}\label{eq:ward-template-constants}
 C^{\rm fl}_\star=C^{\rm fl}_{\rm tp}(1/\sqrt3),\qquad
 C^{\rm tm}_\flat=\left(\frac{3\sqrt3}{e\delta_{\rm tp}L}\right)^3.
\end{equation}
The first constant is the convergent slice sum of
Proposition~\YMMref{120.6}{821}; the second bounds
$\sup_{t\ge0}t^3e^{-\delta_{\rm tp}Lt/\sqrt3}$, whose maximum is at
$t=3\sqrt3/(\delta_{\rm tp}L)$. The comparison constant $C_{\rm tm}$
is chosen to cover both Proposition~\YMMref{121.2}{824} and
Lemma~\ref{lem:ward-fluctuation-template}.

\section{The Ward identity for a general displacement}
\label{sec:ward-general-calculus}

The finite rotational identity of Proposition~\YMMref{112.1}{782}
uses a displacement of the form $\eta(y)=\chi(y)\omega y$.
Its elementary bond calculation also applies to an arbitrary vector
weight. We record that form to fix the stress and null-term conventions.

For a positive bond $b$, let $\partial_b^{\mathsf a}$ differentiate
along $U(b)\mapsto e^{i\epsilon\sigma_{\mathsf a}}U(b)$, where the
$\sigma_{\mathsf a}$ are the Pauli matrices. For a field of directions
$\mathbf X_b$ whose matrix components are continuous functions of the
configuration and continuously differentiable along the bond $b$,
put
\[
 v_{\mathbf X}=\sum_{b,\mathsf a}
       \langle\sigma_{\mathsf a},\mathbf X_b\rangle\partial_b^{\mathsf a},
 \qquad \operatorname{div}v_{\mathbf X}
   =\sum_{b,\mathsf a}\partial_b^{\mathsf a}
       \langle\sigma_{\mathsf a},\mathbf X_b\rangle.
\]
For the displacement $v_\eta$ use
$\mathbf X_b=a^{-1}\sum_\rho\eta^\rho(x_b)\mathbf u_{b,\rho}$,
with $\mathbf u$ as in Definition~\ref{stc:fields} and the paragraph
preceding it. Write $A(v)=\sum_pv(p)a_p$ and
\[
 \Theta_v[\eta]=a^{-1}\sum_{p,b\subset\partial p,\rho}
 v(p)[\eta^\rho(x_b)-\eta^\rho(x(p))]
       \langle\mathbf u_{b,\rho},\mathbf w(W_{b,p})\rangle.
\]
Products $v\eta$ below are evaluated at plaquette centres, separately
for each orientation of a plaquette. Differences of such weights
preserve that orientation.

\begin{lemma}\label{lem:ward-general-calculus}
For every vector weight $\eta$ and plaquette weight $v$,
\begin{equation}\label{eq:ward-general-displacement}
 v_\eta A(v)=\Theta_v[\eta]+\mathfrak K^{\rm v}[v\eta]
                   -A(\nabla^{\rm s}\!\cdot(v\eta)).
\end{equation}
Moreover, with $P(p)=\{\mu,\nu\}$,
\begin{equation}\label{eq:ward-general-divergence}
 \operatorname{div}v_\eta
   =-\tfrac32\sum_pa_p
     [\nabla^{\rm c}_\mu\eta^\mu+
      \nabla^{\rm c}_\nu\eta^\nu](x(p)).
\end{equation}
For a field $\mathbf X$ with this regularity, supported on finitely many bonds,
$\beta v_{\mathbf X}A(1)-\operatorname{div}v_{\mathbf X}$ is a null
term in the sense of Lemma~\ref{stc:null}(b).
\end{lemma}

\begin{proof}
The bond calculation in the proof of
Proposition~\YMMref{112.1}{782} gives, for $b\subset\partial p$,
\[
 \partial_b^{\mathsf a}a_p
   =\langle\sigma_{\mathsf a},\mathbf w(W_{b,p})\rangle,
 \qquad
 \partial_b^{\mathsf a}
    \langle\sigma_{\mathsf a},\mathbf w(W_{b,p})\rangle=1-a_p.
\]
Insert the first identity in $v_\eta A(v)$ and subtract
$\eta^\rho(x(p))$ from $\eta^\rho(x_b)$. The differences give
$\Theta_v[\eta]$; the remaining sum is
\[
 a^{-1}\sum_{p,\rho}v(p)\eta^\rho(x(p))
 \left[\mathfrak r_\rho(p)+\tfrac12(a_{p_+}-a_{p_-})\right].
\]
Reindexing the translated plaquette actions gives
$-A(\nabla^{\rm s}\!\cdot(v\eta))$, proving
\eqref{eq:ward-general-displacement}. The second bond identity gives
\[
 \sum_{\mathsf a}\partial_b^{\mathsf a}
       \langle\sigma_{\mathsf a},\mathbf u_{b,\rho}\rangle
       =\tfrac32(a_{p(b,\rho)}-a_{p(b,-\rho)}).
\]
Collect the two parallel bonds of each plaquette. Their midpoints
differ by $ae_\rho$, so this sum is precisely
\eqref{eq:ward-general-divergence}. Finally, Haar invariance gives
$\int\partial_b^{\mathsf a}F\,dU(b)=0$. Apply it to
$\langle\sigma_{\mathsf a},\mathbf X_b\rangle e^{-\beta A(1)}Y$
and sum over the supported bonds, with $Y$ independent of those
bonds. This is the finite-volume null identity; conditioning on their
complement gives the Gibbs-state identity, as in
Lemma~\ref{stc:null}(b).
\end{proof}

\section{Dilation and normalization of the lattice laws}
\label{sec:lattice-law-complements}

Dilation changes the physical spacing while preserving the bond variables
and the Wilson action. Besides the torus states treated in
Proposition~\YMMref{119.3}{816}, we will need this observation for all
periodic Gibbs states and for chessboard seminorms.

\begin{lemma}\label{lem:lattice-law-dilation}
Let \(K,z\in\mathbb Z\), \(\beta>0\), and define
\(U^{(z)}(L^zb)=U(b)\), \(f^{[z]}(x)=f(L^{-z}x)\).
Dilation is an affine bijection between the Gibbs laws at
\((a_K,\beta)\) and \((a_{K-z},\beta)\). It preserves period numbers,
translation invariance, invariant events and conditioning on such events.
For every law \(\nu\) and compactly supported bounded \(f_i\),
\[
 \mathcal O_K(f)(U)=L^z\mathcal O_{K-z}(f^{[z]})(U^{(z)}),
 \qquad
 S_n^T[\nu](\underline f)
   =L^{zn}S_n^T[\nu^{(z)}](\underline f^{[z]}).
\]
It carries a board \((i,o)\) to \((i-z,L^zo)\), with the same cell
indices and relative parities. On corresponding finite tori,
\[
 \|F^{(z)}\|_{{\rm ch};L^z\ell}^{(i-z,L^zo)}
       =\|F\|_{{\rm ch};\ell}^{(i,o)}.
\]
Here each board and torus must satisfy its usual lattice divisibility
and size conditions; dilation preserves those conditions.
\end{lemma}

\begin{proof}
The configuration map is a permutation of the coordinates, with inverse
given by \(-z\). For a finite set \(B\) of bonds, the plaquettes meeting
\(L^zB\) are exactly the images of the plaquettes meeting \(B\).
Their holonomies and actions agree. The map also preserves the product
Haar measure on the integrated bonds and commutes with joining their
variables to a fixed exterior configuration. Thus the numerator and
denominator of the finite Gibbs kernel transform separately, giving
\[
 \gamma_{L^zB}F^{(z)}=(\gamma_BF)^{(z)}.
\]
Consequently \(\nu\gamma_B=\nu\) for every finite \(B\) if and only if
the analogous equations hold for \(\nu^{(z)}\). This proves the
Gibbs assertion in both directions. Moreover dilation conjugates
\(\tau_v\) to \(\tau_{L^zv}\). It therefore carries each period subgroup
to one of the same index and carries invariant events bijectively to
invariant events. The identity
\(\nu^{(z)}(E^{(z)})=\nu(E)\) also proves that pushforward commutes
with conditioning whenever the conditioning event has positive mass.

Relabel the finitely many plaquettes in each observable and use
\(a_{K-z}=L^za_K\). This proves the observable identity; the joint
generating functions then differ by replacing each source variable by
\(L^z\) times that variable, which proves the cumulant identity.
Finite Wilson measures transform by the same change of variables.
Finally the dilation of every cell is the cell with the same index
in the new board. The reflections conjugate to the new reflections,
and relative parity is unchanged. Hence the folded products have
equal integrals under the corresponding Wilson measures. The number
of cells is unchanged, so taking the same root proves the seminorm
identity.
\end{proof}

We also record the normalization estimate needed when a positive
part of a history integral is discarded. Its loss of mass must be
accounted for before comparing cumulants.

\begin{lemma}\label{lem:positive-functional-comparison}
Let \(\Lambda,\Lambda'\) be positive linear functionals on a unital
algebra of bounded functions closed under conjugation and absolute
values. Suppose that \(\Lambda(1)>0\), \(0\le\epsilon\le1/2\), and
\[
 |(\Lambda-\Lambda')(F)|\le\epsilon\|F\|_\infty\Lambda(1).
\]
Then
\[
 \Lambda'(1)\ge(1-\epsilon)\Lambda(1),\qquad
 \left|\frac{\Lambda'(F)}{\Lambda'(1)}
       -\frac{\Lambda(F)}{\Lambda(1)}\right|
 \le\frac{2\epsilon}{1-\epsilon}\|F\|_\infty
 \le4\epsilon\|F\|_\infty .
\]
The hypothesis holds in particular if \(\Lambda-\Lambda'\) is
positive and its mass is at most \(\epsilon\Lambda(1)\).
If \(|Y_i|\le s_i\), the cumulants of the normalized functionals obey
\[
 |\operatorname{cum}_{\Lambda'}(Y_1,\ldots,Y_n)
      -\operatorname{cum}_{\Lambda}(Y_1,\ldots,Y_n)|
 \le4\epsilon\left(\sum_{\pi}|\pi|!\right)\prod_i s_i,
\]
where the sum is over partitions of \(\{1,\ldots,n\}\).
\end{lemma}

\begin{proof}
Positivity gives
\(|\Lambda(F)|\le\Lambda(|F|)\le\|F\|_\infty\Lambda(1)\):
multiply \(F\) by a unit scalar making \(\Lambda(F)\) nonnegative
real and use \(\Re F\le|F|\). Apply the assumed estimate at \(F=1\).
Subtract the two normalized values by first replacing the numerator
and then the denominator; each resulting term is at most
\(\epsilon\|F\|_\infty/(1-\epsilon)\).
Applying the same positivity bound to \(\Lambda-\Lambda'\) proves
the assertion concerning discarded mass.

Every normalized moment indexed by a set \(I\) is bounded by
\(\prod_{i\in I}s_i\), and its change is at most
\(4\epsilon\prod_{i\in I}s_i\). For a partition with \(q\) blocks,
telescope the product of its \(q\) moments. The resulting bound is
\(4q\epsilon\prod_i s_i\). The coefficient of that product in the
cumulant formula has modulus \((q-1)!\). Summing proves the claim,
as in Lemma~\YMMref{87.2}{684}.
\end{proof}

\subsection{States, first-passage data and their limits}
\label{sec:state-limit-conventions}

A state at \(g\), at a fixed pair \((K,\beta)\), is either the
periodic lift of the Wilson law on a side \(\ell\in\mathfrak L(g)\),
or a translation-invariant periodic Gibbs state at that pair.
This definition alone does not impose first passage.
Write \(\mathcal G^{\rm inv}_{K,\beta}\) for the latter states.
Proposition~\YMMref{103.4}{758} shows that every periodic Gibbs
state is invariant under all lattice translations, so
\(\mathcal G^{\rm inv}_{K,\beta}=\mathcal G^{\rm per}_{K,\beta}\).
Torus-limit states in \(\mathcal G^{\rm tor}_{K,\beta}\) exist and belong to this set by
Proposition~\YMMref{91.2}{699}.
The finite Wilson law is translation invariant by relabelling its
bonds. Consequently every state just defined has symmetric
multilinear cumulants invariant under simultaneous lattice
translations. For a lifted torus law those cumulants are the torus
cumulants of the periodized test functions.

Fix the construction constants. The canonical family
\(\mathfrak d(g)\) consists of the first-passage pairs \((K,\beta)\)
at \(g\), as in Section~\YMMref{74}{614}. Their labels lie in
\[
 I^{\rm lab}(g)=
 [g^{-2}+B_\sharp(L-2)/(L-1),\,L(g^{-2}+\bar B)].
\]
This is Proposition~\YMMref{74.1}{615}; in particular every label
is at least \(g^{-2}\). For every \(b\in I(g)\) and \(K\ge1\),
\((K,L^Kb)\in\mathfrak d(g)\), by Corollary~\YMMref{74.2}{616}.

A separated tuple in an open set \(O\) is a tuple of real functions
in \(C_c^1(O)\) with pairwise disjoint supports.
A limit at \(g\) along growing tori is a family \((S_n)_{n\ge2}\)
obtained on every separated tuple from the cumulants along
\[
 (K_j,\ell_j,\beta_j),\qquad
 (K_j,\beta_j)\in\mathfrak d(g),\quad
 \ell_j\in\mathfrak L(g),\quad K_j,\ell_j\longrightarrow\infty,
\]
with the test functions periodized on each torus.
A volume-first limit at \(g\) is obtained instead from
\[
 S_n^T[\nu_j],\qquad
 (K_j,\beta_j)\in\mathfrak d(g),\quad K_j\longrightarrow\infty,\quad
 \nu_j\in\mathcal G^{\rm tor}_{K_j,\beta_j}.
\]
A label along a defining sequence is a subsequential limit of
\(a_{K_j}\beta_j\). Passing to that subsequence leaves the limiting
family unchanged. These are the two classes preceding
Proposition~\YMMref{119.4}{817}, restricted to first-passage pairs
at \(g\) and, in the first case, to \(\mathfrak L(g)\).

\section{Local large-field probabilities}\label{sec:local-large-field-rarity}

The history bounds of Section~\YMMref{70}{604} include an exponential
of the volume. To bound the probability that a given point belongs to a
large-field region, we first attach a gauge invariant test to each possible
creation. Reflecting that test across a board forces many creations. The
root in the chessboard estimate then removes the volume factor.

We use the real construction of Theorem~\YMMref{D}{8}, with the constants
constructed in Section~\YMMref{69}{595} and the final source ceilings
of Section~\YMMref{82}{665}. The couplings obey
$x_{j-1}=L(x_j+b_j)$, $|b_j|\le B_\sharp$, and the lattice is admissible
through level $K$. Write $y_j=\log x_j$, $g_j=x_j^{-1/2}$ and
\[
 p_j=p_0(g_j)=A_0y_j^{\mathfrak p_0},\quad
 \varepsilon_j=g_jp_j,\quad \delta_j=(A_1/A_0)\varepsilon_j,\quad
 s_j=M\check R_j a_{K-j},\quad w_j=M_2\check R_j a_{K-j}.
\]
Here $\check R_j$ is the least nonnegative integral power of $L$ at least
$y_j^r$, $M_2=L^2M_1$, $M\ge L^2M_2$, $L\ge13$, and
$2\mathfrak p_0>3r+2$. A cell of level $j$ has side $s_j$; a test cube
has side $Lw_j$, and its enlargement by $a$ layers has side $(L+2a)w_j$.
Distances in this section are periodic sup distances. The probability of
a history is $\widetilde\nu_k(h)=\mathsf t(h)/Z$, where
$\mathsf t(h)=\int\mathfrak w_h\rho_0\,dU$, $\mathfrak w_h=\mathbf1_h\mathfrak G_h$,
and $\rho_0$ is the Wilson density. The exact history sum gives
$\sum_h\mathsf t(h)=Z$.

\begin{definition}[The two choices and their threshold]\label{def:rarity-choices}
Put
\[
 \kappa_{\rm w}=\min\{(4\Theta_1)^{-1},\theta_c\},\qquad
 \lambda_\natural=\frac{\kappa_{\rm w}}{2\pi},\qquad
 \lambda'_\natural=\frac{\lambda_\natural}{2+30L}.
\]
A lowered choice keeps the constants before $A_1$ and the value of $A_0$,
replaces $A_1$ by $A'_1=\lambda'_\natural A_1$, and completes the
construction with $B'_\sharp\ge B_\sharp$. Primes denote the constants
of this second choice. Set
\begin{equation}\label{eq:rarity-constants}
\begin{split}
 c_1^\natural&=\frac{c'_\star}{4\cdot49^3},\qquad
 C_1^\natural=\mathsf C'_{\rm maj}+c'_{\rm low}+\frac{13}{2},\\
 c'_{\rm low}&=C'_5+3\log7+24(B'_3)^2,\\
 c_2^\natural&=\frac12\min\left\{\frac{c_1^\natural}{8},
                          \frac{\lambda_\natural^2A_1^2}{2A_0^2}\right\},\\
 N_\tau&=\left(\frac{M}{LM_2}\right)^3,\qquad N_\delta=\left(\frac M{M_2}\right)^3,\\
 C_3^\natural&=27(N_\tau+N_\delta)
       +540N_\tau\bigl((L+6)L^2M_2\bigr)^3(2L)^{3/2}.
\end{split}
\end{equation}
The positive constants $c'_{\star},\mathsf C'_{\rm maj},C'_5$ are those of
the deposit and history bounds in Section~\YMMref{70}{604} at the second
choice. Let $\bar B=LB_\sharp/(L-1)$,
$Y_0=\max\{\log\max(e,2B_\sharp),2\log(2L)/(L^{1/r}-1)\}$, and let
$Y'_{\rm mark}$ be the threshold below equation~\YMMref{(1455)}{592}
with $w=\eth=1/2$, $c=c'_\star$, $c_{\rm Dep}=c'_\star$ and $c_a=0$,
at the second choice. Explicitly it is
\[
 Y'_{\rm mark}=\max\left\{
 (Y^{\rm II})'\left(c'_\star/4,48^3(D'_0+\lambda'_L)\right),
 (Y^{\rm S})'(c'_\star/4)\right\},\qquad
 D'_0=\mathsf C'_{\rm maj}+27(2+\log(2L))C'_{\flat}.
\]
Here $C'_{\flat}$ is the vacuum surface constant used in
equation~\YMMref{(1487)}{606}; $Y^{\rm II}$ and $Y^{\rm S}$ are the
explicit functions in Section~\YMMref{67}{587}. Finally put
\[
 Y_{\rm pay}(c,C)=\max\left\{1,\log(1+\bar B),
 \left(\frac4{cA_0}\right)^{1/\mathfrak p_0},
 \left(\frac{4CM^3L^3}{cA_0^2}\right)^{1/(2\mathfrak p_0-3r-2)}\right\}.
\]
The threshold $\mathsf Y^\natural$ is the least number at least $1$ such
that every $y\ge\mathsf Y^\natural$ satisfies
\begin{equation}\label{eq:rarity-seven-conditions}
\begin{aligned}
 &\textup{(Y.1)}&&y\ge\log(\gamma'^{-2}),&
 &\textup{(Y.2)}&&y\ge Y_0,\\
 &\textup{(Y.3)}&&y\ge Y'_{\rm mark},&
 &\textup{(Y.4)}&&y\ge Y_{\rm pay}(c_1^\natural,C_1^\natural),\\
 &\textup{(Y.5)}&&
 C_{\rm cu}\sqrt{2L}(A_0/A_1)e^{-\delta_{\rm lc}M_2y^r/4}
                   \le\lambda_\natural,\\
 &\textup{(Y.6)}&&y\ge
       \frac{\log C_3^\natural+3r+3/2}{c_2^\natural A_0^2},\\
 &\textup{(Y.7)}&&y\ge2\log(1+\bar B),\qquad
       2c_2^\natural\mathfrak p_0A_0^2y^{2\mathfrak p_0-1}\ge4.
\end{aligned}
\end{equation}
The constants $C_{\rm cu}$ and $\delta_{\rm lc}$ are those of the local
background comparison~\YMMref{(1012)}{451} at the original choice.
All these numbers are independent of the lattice and the coupling sequence.
We also write $\mathsf Y_i^\natural$ for the lower bound on $y$ in
(Y.$i$), taking the maximum with $1$ in (Y.6). Thus
$\mathsf Y^\natural=\max_i\mathsf Y_i^\natural$; in particular the
four thresholds used for the witness estimate are exactly those in
(Y.1)--(Y.4).
\end{definition}

\begin{lemma}\label{lem:rarity-lowered-choice}
For a choice made by the ordered construction of
Section~\YMMref{69}{595}, a lowered choice exists by resuming that
construction after $A_0$. If a lattice is admissible at the original
choice and $x_K\ge\gamma'^{-2}$, it is admissible at the lowered choice.
The lattices, block averages, local background maps and their stable
domains, plaquette tests, gauge densities and Wilson measure are the same
for both choices. The fluctuation thresholds satisfy
$\delta'_j=\lambda'_\natural\delta_j$.
\end{lemma}

\begin{proof}
In the ordered parameter choice, $A_1$ is free in $(0,\infty)$ and has
no attached condition. The next constant is $A_0$. Its eight conditions
are
\[
\begin{gathered}
 A_0\ge c_{\rm im}A_1,\quad A_0>c_{\rm ct}A_1,\quad
 A_0\ge C_{\rm sf}L^3A_1,\quad A_1/A_0\le c_{\rm Q},\quad
 A_1/A_0\le c_{\rm no},\\
 A_0\ge4B_3C_HA_1,\quad A_0/A_1\ge64L^3,\quad
 A_0/A_1\ge384\pi L^3B_5\mathsf S(\delta_{0,\rm tw}).
\end{gathered}
\]
Every other constant in these conditions precedes $A_1$. Decreasing
$A_1$ while keeping $A_0$ preserves all eight inequalities. The admissible
partial choice can be completed by the following continuation of
Section~\YMMref{69}{595}. Rechoose $C_0,C_1$, then the radii
$\alpha_1,\alpha_0,\alpha_2,\alpha_3,\alpha_4,a_\flat$ in the stated
order, retaining the strict reserves for every difference of radii.
The lower bounds on $\theta_c$ involve only the earlier fixed
$\Theta_1,M$ and chart constants, so its original value may be retained.
Evaluate the normalizer constants, increase $E_0$ as required, and
choose $B'_\sharp$ above both its new lower bound and $B_\sharp$.
All these are finite choices before the coupling ceiling. Finally
reduce $\gamma'$ so that the real numerical inequalities hold on
$x\ge\gamma'^{-2}$. The last part of Section~\YMMref{69}{595}
expresses them as finitely many quantities tending to zero, or
finite lower thresholds, on that interval. The source conditions
of Section~\YMMref{82}{665} impose only further such restrictions
on this final ceiling. This constructs the two choices actually used
here; no extension assertion for arbitrary partially specified lists
is required.

The listed common objects depend only on the constants before $A_1$,
on $A_0$, and on $(K,m,x_\bullet)$. In particular their definitions do
not use the later fluctuation thresholds. The formula for $\delta_j$
gives the asserted scaling. The data conditions $\varepsilon_j\le1$
and $M\check R_j\le L^{m+K-j}$ are unchanged, $\delta'_j\le\delta_j$,
and the same increments satisfy $|b_j|\le B'_\sharp$. The remaining
coupling condition is precisely $x_K\ge\gamma'^{-2}$.
\end{proof}

\subsection{Witnesses for a creation}

Write $V_j=\mathcal M^jU$. For a test cube $Q$ of step $t$, set
$Y_Q=(Q^{\sim3})^{(t)}$ and dress a bond of level $t-1$ by
\[
 V^{\rm dr}_{t-1}(b)=
 V_{t-1}([b_-]_t,b_-)V_{t-1}(b)
                    V_{t-1}([b_+]_t,b_+)^{-1}.
\]
Define $F_{t,Q}(U)=0$ exactly when all the following conditions hold:
the plaquette test of $Q$ is one; the local background
$V^{\rm bg}_{t-1,Q}$ is stably defined; all contour averages
$V_{t-1}(y,x)$, $y\in Y_Q$, $x\in B(y)$, are defined; and
\begin{equation}\label{eq:rarity-witness-test}
 |V^{\rm dr}_{t-1}(b)(V^{\rm bg}_{t-1,Q}(b))^{-1}-1|
       <\lambda_\natural\delta_{t-1}
\end{equation}
for every non-pivot bond incident to $Q^{\sim2}$. Otherwise set
$F_{t,Q}=1$. At step one include the witnesses $F_{0,Q}=1-\chi_0(Q)$
for the bare density cubes. Their support cubes are respectively
$Q^{\sim4}$ and $Q^{\sim1}$.

\begin{lemma}\label{lem:rarity-witnesses}
Under (Y.2), these witnesses are Borel, gauge invariant, and covariant under the
lattice isometries that preserve the partitions. Their support cubes
have diameter at most $(L+8)w_t\le(L+8)L^{-2}s_k$ when $t\le k$.
If also (Y.5) holds, a cell $\mathfrak c$ is created at step $t$ and
$\mathfrak w_h(U)>0$, then either a witness whose cube meets
$\mathfrak c$ equals one, or there is a test cube $Q$ meeting
$\mathfrak c$ and a non-coarse site $x\in(Q^{\sim3})^{(t-1)}$ such that
$[x]_t\in P''_t(h)$ and the defined contour variable satisfies
\begin{equation}\label{eq:rarity-gauge-alternative}
 |V_{t-1}([x]_t,x)-1|>\lambda_\natural\delta_{t-1}.
\end{equation}
Under (Y.1), on every history at the lowered choice with
$\mathfrak w'_h(U)>0$, a witness with value one has its centre in $Z'_t$.
\end{lemma}

\begin{proof}
The stable domains of the local background maps are open, and the maps
are continuous there. Contour averages and iterated block averages are
Borel on their domains. Thus all four tests defining $F$ are Borel.
The local background reads only the inner bonds of $Q^{\sim4}$.
The blocks needed for every contour or block average are contained in
that cube, since all its faces are block faces. This proves the stated
support. The bare density witness reads only the plaquettes of
$Q^{\sim1}$. The side bounds follow from $M\ge L^2M_2$ and the
monotonicity of $s_t$ in Lemma~\YMMref{67.1}{587}.

Under a gauge transformation $u$, the dressed bond transforms with the
values $u([b_-]_t),u([b_+]_t)$ at the coarse endpoints. The local
background has exactly the same transformation law, so their quotient
is conjugated. The plaquette test and the domains of the contour and
background maps are gauge invariant. This proves gauge invariance of
$F$. An isometry preserving the block partitions carries each contour,
pivot bond, local minimizing orbit and axial representative to the
corresponding object of the image cube. It also carries the stable
domain onto the image stable domain, because it is a homeomorphism
intertwining the two partial maps. Thus it preserves all four tests.
In particular these conclusions apply to the isometries with
translations in the lattice of spacing $s_k/L$: every partition used
above has side at most $Lw_t\le s_k/L$, an integral power of $L$.

For the creation alternative, first observe that, with
$w_\pm=V_{t-1}([b_\pm]_t,b_\pm)$,
\[
 |V_{t-1}(b)-V^{\rm bg}_{t-1,Q}(b)|
 \le |w_--1|+|w_+-1|
                +|V^{\rm dr}_{t-1}(b)-V^{\rm bg}_{t-1,Q}(b)|.
\]
Consequently a bond discrepancy at least $3\lambda_\natural\delta_{t-1}$
either makes $F=1$ or forces one endpoint contour to satisfy
\eqref{eq:rarity-gauge-alternative}. If $Q^{\sim3}\subset B(P''_t)$,
that endpoint lies in the asserted gauge-fixing region. It cannot be a
coarse site, where the contour variable is one.

The created cells are the hulls of the failed plaquette cubes, failed
$Q$-test cubes, failed fluctuation cubes, selected intermediate bonds,
and failed low-field cells; at step one there are also failed bare
density cubes. A failed plaquette or bare density test directly gives
a witness. A failed $Q$-test gives either an undefined local background
or a non-pivot discrepancy at least $2\delta_{t-1}$. Its cube satisfies
$Q^{\sim3}\subset B(P''_t)$, so the preceding observation applies.

In each of the remaining three cases a non-pivot bond incident to a
test cube $Q\subset\Omega_t^{\sim-1}$ meeting $\mathfrak c$ satisfies
\[
 |\mathcal A_{t-1}(b)|\ge\kappa_{\rm w}\delta_{t-1}.
\]
Indeed the respective failed tests have thresholds $\delta_{t-1}$,
$\delta_{t-1}/(4\Theta_1)$ and $\theta_c\delta_{t-1}$. The geometry
gives $Q^{\sim4}\subset\Omega_t\subset B(P''_t)$. The regularity and
$P,Q$ tests on the support of the history meet the hypotheses of
Proposition~\YMMref{53.12}{457}. Its principal-log estimate gives
\[
 |V_{t-1}(b)-V^{\rm bg,step}_{t-1}(b)|
       \ge(2/\pi)|\mathcal A_{t-1}(b)|
       \ge4\lambda_\natural\delta_{t-1}.
\]
The local comparison~\YMMref{(1012)}{451} is applicable on this cube:
its four-layer enlargement is contained in $\Omega_t$ and both endpoints
of $b$ belong to $Q^{\sim3}$. It gives a difference of backgrounds at
most $C_{\rm cu}\varepsilon_te^{-\delta_{\rm lc}M_2\check R_t/4}$.
Since $\varepsilon_t\le\sqrt{2L}(A_0/A_1)\delta_{t-1}$ and
$\check R_t\ge y_K^r$, condition (Y.5) bounds this by
$\lambda_\natural\delta_{t-1}$. The discrepancy from the cube background
is therefore at least $3\lambda_\natural\delta_{t-1}$, as required.

Finally suppose that a witness centre belongs to $\Omega'_t$. Its
whole test cube lies there. The lowered plaquette and $Q$ tests are
one, the local background is stably defined, and its bond discrepancy
is less than $2\delta'_{t-1}$. Proposition~\YMMref{53.12}{457}, now
at the lowered choice, gives all contour averages and bounds them by
$15L\delta'_{t-1}$. Hence the dressed discrepancy is less than
$(2+30L)\delta'_{t-1}=\lambda_\natural\delta_{t-1}$. Thus $F=0$.
The bare density test is also one on $\Omega'_1$. A witness with value
one consequently has its centre outside $\Omega'_t$, hence in
$Z'_t=(\Lambda'_t)^c$.
\end{proof}

\subsection{Gauge integration and the chessboard estimate}

For a class function $0\le\phi\le1$, let $\mathsf{gv}_{t,x}$ equal
$\phi(V_{t-1}([x]_t,x))$ when the contour is defined and
$[x]_t\in P''_t(h)$, and zero otherwise. Only non-coarse sites are used.
The density of one gauge variable is
\begin{equation}\label{eq:rarity-gauge-density}
 f^{\rm gv}_j(u)=z_j^{-1}\mathbf1_{\{|u-1|<\varepsilon_j\}}
                      e^{-x_j(1-\frac12\operatorname{tr}u)}.
\end{equation}

\begin{lemma}\label{lem:rarity-gauge-integration}
For a finite set $J$ of distinct pairs $(t,x)$, class functions
$0\le\phi_{t,x}\le1$, and a bounded nonnegative gauge invariant Borel
function $u$ of the bare field,
\begin{equation}\label{eq:rarity-gauge-integration}
 \sum_h\int\mathfrak w_h u\prod_{(t,x)\in J}\mathsf{gv}_{t,x}\,dU
 \le\left(\prod_{(t,x)\in J}\int_G\phi_{t,x}f^{\rm gv}_{t-1}\,dv\right)
                                                     \int u\,dU.
\end{equation}
For the indicator of $|v-1|>\lambda_\natural\delta_j$, its integral $q_j$
satisfies
\begin{equation}\label{eq:rarity-gauge-tail}
 q_j\le20\exp\left[\tfrac32y_j
             -\tfrac12\lambda_\natural^2(A_1/A_0)^2p_j^2\right].
\end{equation}
\end{lemma}

\begin{proof}
Induct on the last step. Fix a prefix history and the plaquette outcome
of its final step. Summing the remaining outcome indicators gives at
most one. The remaining last-step gauge factor is a product of the
densities \eqref{eq:rarity-gauge-density} over
$B(P''_t)\setminus P''_t$. Average over independent gauge transformations
at these sites, extended constantly over each finer block and equal
to one at the coarse centres. This preserves the prefix, the plaquette
outcome, Haar measure and $u$. Earlier contour factors are only
conjugated, and their class functions are unchanged. A last-step
contour becomes $V_{t-1}([x]_t,x)v(x)^{-1}$. If the necessary contour
domains fail, its contribution is zero; otherwise Haar translation
turns each site factor into $\int\phi_{t,x}f^{\rm gv}_{t-1}$, or into
one at a site not in $J$. The plaquette indicators sum to the prefix
indicator. Applying the induction hypothesis proves the displayed
inequality, including $J=\emptyset$.

On $\mathrm{SU}(2)$, $1-\frac12\operatorname{tr}v=|v-1|^2/2$.
The Haar radial density for $r=|v-1|$ is
$(2/\pi)r^2\sqrt{1-r^2/4}\,dr$. Since
$x_j^{-1/2}\le\varepsilon_j\le1$,
\[
 z_j\ge\frac{\sqrt3}{3\pi\sqrt e}x_j^{-3/2}
                   >\frac1{20}x_j^{-3/2}.
\]
On the tail set the exponential is at most
$e^{-x_j\lambda_\natural^2\delta_j^2/2}$. Now
$x_j\delta_j^2=(A_1/A_0)^2p_j^2$, which proves the tail bound.
\end{proof}

\begin{lemma}\label{lem:rarity-chessboard}
Under (Y.1)--(Y.4), every witness $F$ of step $t$ satisfies
\begin{equation}\label{eq:rarity-one-witness}
 \langle F\rangle\le e^{-c_1^\natural p_t^2}.
\end{equation}
More generally, at a level $k\ge t$, the same bound holds for its
chessboard seminorm on any of the eight boards of side $s_k$ with
offset $(L-1)s_k\sigma/(2L)$, $\sigma\in\{0,1\}^3$, whose cell
contains its support cube. These boards preserve all construction
partitions.
For a family of witnesses with support cubes at pairwise distance
greater than $s_k$, their joint mean is at most
$\prod_F e^{-c_1^\natural p_{t(F)}^2/8}$.
\end{lemma}

\begin{proof}
Use the eight boards of cell side $s_k$ whose offsets are
$(L-1)s_k\sigma/(2L)$, $\sigma\in\{0,1\}^3$. In one coordinate the
walls of the two families have minimum separation $(L-1)s_k/(2L)$.
An interval shorter than this separation lies between two consecutive
walls of at least one family. The support diameter
$(L+8)s_k/L^2$ is shorter, since $L\ge13$. Choosing the family in
each coordinate puts every witness in a cell of one of these boards.
All board reflections and translations preserve its construction,
by Lemma~\ref{lem:rarity-witnesses}.

Let $N=(\ell/s_k)^3$ be the number of cells and let $\Pi$ be the product
of all reflected copies of $F$. Each copy is a witness of the same
step $t$. On a history at the lowered choice with $\mathfrak w'_h(U)>0$,
$\Pi=1$ forces all $N$
witness centres into $Z'_t$. By Proposition~\YMMref{67.2}{588}, such
a centre is within
$\rho_{q,t}=7s_q+11\sum_{j=q+1}^t s_j$ of a creation at some $q\le t$.
The centre of its board cell is at most
$s_k/2+\rho_{q,t}+s_q/2$ from the centre of that creation. Since
$2\rho_{q,t}+s_q<48s_t\le48s_k$, a creation can be assigned at most
$49^3$ board centres. There are therefore at least $N/49^3$ creations.

For the mark
$\Theta=(c'_\star/2)\sum_q p_q^2 n_q^{\rm cr}$, this gives
$\Theta\ge2c_1^\natural p_t^2N$. Equation~\YMMref{(1487)}{606}, at
$w=\eth=1/2$ and $c_a=0$, and Proposition~\YMMref{70.2}{606} yield
\[
 \langle\Pi\rangle\le
 \exp\left[-2c_1^\natural p_t^2N+
       C_1^\natural y_k(M\check R_k)^3N\right].
\]
The normalization is the same Wilson partition function for the two
choices. Condition (Y.3) is exactly the marked-mass hypothesis.
Condition (Y.4) and $\check R_k<Ly_k^r$, $y_k\ge1$, give
$C_1^\natural y_k(M\check R_k)^3\le c_1^\natural p_k^2
\le c_1^\natural p_t^2$. Thus
$\langle\Pi\rangle^{1/N}\le e^{-c_1^\natural p_t^2}$.
The chessboard estimate of Proposition~\YMMref{86.3}{681} gives
\eqref{eq:rarity-one-witness} and the seminorm assertion.

For several witnesses assign each to one of the eight boards and let
$\Phi_\sigma$ be the product in that class. Distinct witnesses in a
class occupy distinct cells, by the separation assumption. Since
$\Phi_\sigma^8=\Phi_\sigma$, H\"older's inequality and then the
chessboard estimate give
$\langle\prod\Phi_\sigma\rangle\le
\prod\langle\Phi_\sigma\rangle^{1/8}
\le\prod_F e^{-c_1^\natural p_{t(F)}^2/8}$.
\end{proof}

\subsection{From witnesses to points}

\begin{proposition}\label{prop:local-large-field-rarity}
Suppose $y_K\ge\mathsf Y^\natural$. For any prescribed cell
$\mathfrak c$ of level $t\le k$, the probability that it is created
at step $t$ is at most $e^{-c_2^\natural p_t^2}$. For every
$x\in T_\ell$ and $1\le n\le k\le K$,
\begin{equation}\label{eq:rarity-point}
 \widetilde\nu_k\{x\in Z_n\}
             \le2\cdot24^3 e^{-c_2^\natural p_n^2}.
\end{equation}
The same bound holds for a prescribed cell contained in $Z_n$.
For $0\le a<k$ the right side for $n=k-a$ is at most
\begin{equation}\label{eq:rarity-age}
 2\cdot24^3e^{-c_2^\natural p_k^2}
 L^{-2c_2^\natural\mathfrak p_0A_0^2y_k^{2\mathfrak p_0-1}a}.
\end{equation}
\end{proposition}

\begin{proof}
At most $27(N_\tau+N_\delta)$ witness cubes meet $\mathfrak c$.
Indeed a partition of side $a\le s_t$ has at most
$(s_t/a+2)^3\le27(s_t/a)^3$ cubes meeting it. Apply this to the test
cubes, and, at $t=1$, to the bare density cubes; their quotients are
bounded respectively by $M/(LM_2)$ and $M/M_2$. Each test cube has
at most $((L+6)LM_2\check R_t)^3$ sites in its three-layer enlargement.
There are therefore at most $27N_\tau$ times this many possible
gauge-tail sites.

Multiply the alternative of Lemma~\ref{lem:rarity-witnesses} by
$\mathfrak w_h\rho_0$ and sum and integrate over histories.
Lemma~\ref{lem:rarity-gauge-integration}, first with $u=F\rho_0$ and
then with $u=\rho_0$, and Lemma~\ref{lem:rarity-chessboard} bound the
creation probability by
\[
 e_t=27(N_\tau+N_\delta)e^{-c_1^\natural p_t^2/8}
 +27N_\tau((L+6)LM_2\check R_t)^3q_{t-1}.
\]
Here weakening the first exponent to $c_1^\natural/8$ keeps the same
constant as the bound for separated creations. The inequalities
$p_{t-1}\ge p_t$, $y_{t-1}\le y_t+\log(2L)$ and
$\check R_t<Ly_t^r$ imply
\[
 e_t\le e^{-2c_2^\natural p_t^2}C_3^\natural
             e^{(3r+3/2)y_t}\le e^{-c_2^\natural p_t^2}.
\]
The last inequality follows from (Y.6), $y_t\ge y_K\ge1$ and
$2\mathfrak p_0\ge2$.

If $x\in Z_n$, Proposition~\YMMref{67.2}{588} supplies a creation
of level $t\le n$ within distance $\rho_{t,n}$. There are at most
$24^3L^{3(n-t)}$ possible cells: their count is at most
$(2\rho_{t,n}/s_t+2)^3$, and
\[
 2\rho_{t,n}/s_t+2
 \le16+22\sum_{j=1}^{n-t}L^j\le24L^{n-t}.
\]
We record the coupling estimate needed to sum these bounds, including
the small defect from an exact geometric progression. Put $P=2\mathfrak p_0$,
$\Delta=y_{n-a}-y_n$, and $v=\bar B/x_n$. Condition (Y.7) gives
$x_n\ge(1+\bar B)^2\ge4\bar B$. The affine flow bound gives, for
$a\ge1$,
\[
 \Delta\ge a\log L+\log(1-v)>2a,\qquad
 a\log L-\Delta\le\tfrac43\bar B e^{-y_n},\qquad
 y_n(a\log L-\Delta)\le\frac8{3e}\frac{\bar B}{1+\bar B}<1.
\]
The last inequality uses $ze^{-z/2}\le2/e$. The first two terms of
the power expansion therefore give
\[
 y_{n-a}^P-y_n^P\ge Py_n^{P-1}a\log L
 +Py_n^{P-2}\left[\frac{P-1}{2}\Delta^2
                  -y_n(a\log L-\Delta)\right]
 \ge Py_n^{P-1}a\log L.
\]
By (Y.7), the rate $c_2^\natural PA_0^2y_n^{P-1}$ is at least four.
Thus the union bound over cells and levels is at most
\[
 24^3\sum_{a=0}^{n-1}L^{3a}e^{-c_2^\natural p_{n-a}^2}
 \le24^3e^{-c_2^\natural p_n^2}\sum_{a\ge0}L^{-a}
 \le2\cdot24^3e^{-c_2^\natural p_n^2}.
\]
A cell event implies the point event at its centre. The same power
estimate with $n=k$ proves \eqref{eq:rarity-age}.
\end{proof}


\section{Gaussian restrictions and history tails}
\label{sec:ward-history-complements}

\subsection{Removing the fluctuation boxes}

\begin{lemma}\label{lem:ward-box-cumulants}
Consider a complete Gaussian structure on a torus satisfying
Theorem~\YMMref{R}{361}, with at most three admitted homogeneous quadratic origins
of sort $\Sigma_{\rm q}$ and no origin of sort $\Sigma_{\rm v}$.
Let $\nu^{\rm r}$ be its normalized measure restricted to the
fluctuation box of radius $\mathfrak b$, and let $\nu^\circ$ be the
normalized unrestricted measure with the same quadratic origins.
For $n$ insertions $W_i$ of sort $\Sigma_{\rm q}$, with nonzero norms
$a_i=\|W_i\|_{\star,\lambda_V}$, set
\[
 r_i=\frac{w_\star}{na_i},\qquad
 \Pi_{\rm ins}=\prod_i r_i^{-1},\qquad
 N_{\rm ins}=\sum_{X:\,X\text{ meets every insertion}}
                         e^{-\lambda_{\rm out}d(X)}.
\]
Here ``meets'' means that $X$ contains the enlarged localization
domain of at least one nonzero member of that insertion. Then
\begin{equation}\label{eq:ward-box-cumulants}
 \left|\langle W_1^{\rm tot};\ldots;W_n^{\rm tot}\rangle_{\nu^{\rm r}}
       -\langle W_1^{\rm tot};\ldots;W_n^{\rm tot}\rangle_{\nu^\circ}\right|
 \le K_P^\Pi\mathsf p_\star\Pi_{\rm ins}N_{\rm ins},
\end{equation}
where, in the notation of equation~\YMMref{(705)}{360},
\[
 \mathsf p_\star=\mathsf C_{\rm ex}((L+2)M)^3e^{4\lambda_1}
                         e^{-\gamma_2\mathfrak b^2/2}.
\]
If $\lambda_{\rm out}\ge c_{\rm E}$, then for $n=2$ or $4$ and
single field components at arbitrary bonds,
the right side is at most $C_{\rm bx}e^{-\gamma_2\mathfrak b^2/2}$,
where $C_{\rm bx}$ is independent of the side, cutoff and bonds.
Both measures are even, so odd joint cumulants of field components vanish. In the unrestricted
measure the second cumulant is its covariance and the fourth is zero.
\end{lemma}

\begin{proof}
Add the single origin $\sum_i z_iW_i$ on the polydisc $|z_i|<r_i$.
Its norm is at most $\sum_i|z_i|a_i<w_\star$, so the system has at
most four admitted origins. Proposition~\YMMref{47.4}{365} makes its
cluster sums holomorphic. The cluster integral identity of
Theorem~\YMMref{R}{361} gives the generating function as
\[
 \int e^{\sum_i z_iW_i^{\rm tot}}d\nu^{\rm r}
       =\exp\left[\sum_X(E_z(X)-E_0(X))\right].
\]
The same identity holds at threshold infinity. All these integrals
are absolutely convergent. To justify differentiation, replace the
quadratic origins by their real parts and the inserted origin by
$\sum_i r_i|W_i|$. Monotonicity of the potential norm leaves it
admitted. Its absolutely integrable density majorizes the integrands
on the closed source polydisc; Cauchy's formula on a smaller polydisc
and Fubini's theorem justify every derivative.

The mixed derivative of $E_z(X)$ vanishes unless $X$ meets every
insertion: otherwise one source coordinate changes no potential in
$X$, and the locality clause of Proposition~\YMMref{47.4}{365}
makes $E_z(X)$ independent of that coordinate. Thus cumulants are
the sums of the mixed cluster derivatives on the set in $N_{\rm ins}$.
Equation~\YMMref{(724)}{364}, applied at every $z$, bounds
$|E_z(X)-E_z^\circ(X)|$ by
$K_P^\Pi\mathsf p_\star e^{-\lambda_{\rm out}d(X)}$; there are no
$\Sigma_{\rm v}$ origins, and the inserted origin remains of sort
$\Sigma_{\rm q}$. At infinite threshold the same structure and norms
apply, while $\mathsf p_\star$ becomes zero. Cauchy's estimate with
radii increasing to $r_i$ gives the factor $\Pi_{\rm ins}$.
Summing proves \eqref{eq:ward-box-cumulants}.

A component at one bond may be placed on a fixed elementary
localization domain $Y_i$. Its norm is bounded by
$\frac12M^{3/2}e^{\lambda_Vd(Y_i)}$, since
$s\le\frac12M^{3/2}(1+M^{-3}s^2)$ for $s\ge0$. The elementary
domains have uniformly bounded size. The polymer sum anchored at
one such domain is bounded independently of the torus, by the
exponential domain-counting estimate~\YMMref{(684)}{353} of
Proposition~\YMMref{45.1}{353}.
Thus $\Pi_{\rm ins}N_{\rm ins}$ is uniformly bounded for $n=2,4$.
The symmetry $B\mapsto-B$ proves evenness. Completing the Gaussian
square gives logarithmic generating function
$\frac12\sum_{i,j}z_iz_j(\varphi_i,C\varphi_j)$ in the unrestricted
measure, proving its stated cumulants.
\end{proof}

In the applications to the quadratic systems of
Proposition~\ref{C-350-6}, the Gaussian structure and the quadratic
origin are precisely the admitted one-step system at field scale
zero in Proposition~\ref{IV-71-37}. The complete structure is that
of Proposition~\YMMref{120.3}{819}; its covariance after the quadratic
origin is included is $\mathsf M_i^{-1}$. Taking
$\mathfrak b=\mathfrak b(x_i)=A_1(\log x_i)^{\mathfrak p_0}$ gives
both exponential error bounds used in the field-scale expansions.
Here $\lambda_{\rm out}=(1+4\beta_{\rm s})\kappa>c_{\rm E}$:
the entropy floor in the choice of $\kappa$ in
Section~\YMMref{69}{595} already implies $\kappa>c_{\rm E}$.

\subsection{First creation on a fixed torus}

We use the all-small-field history measures $m_t$, step factors $g_t$,
and conditional quotients $\Phi_t$ of Section~\YMMref{123}{828}.
The history law is $\widetilde\nu_k(h)=Z^{-1}\int\mathfrak w_h\rho_0\,dU$;
the measure on the all-small-field history is
$m_t[F]=(\mathcal M^t)_*(F\mathfrak w_{h_t}\rho_0\,dU)$, and
$m_t=m_t[1]$. Its quotient is the bounded density
$\Phi_t[F]=dm_t[F]/dm_t$, defined up to an $m_t$-null set.
The exact identities~\YMMref{(2269)--(2271)}{830} hold for every
bounded gauge invariant insertion $F$, provided the prescribed tests
and the ceilings $\varepsilon_j\le\varepsilon_\flat$ hold through
the levels used. In particular
\begin{equation}\label{eq:ward-first-creation}
 \langle F\rangle=Z^{-1}\int\Phi_T[F]dm_T
   +Z^{-1}\sum_{t=1}^T\int\Phi_{t-1}[F]g_t^{\rm w}\,dm_{t-1},
 \qquad
 \sum_{t\le T}w_t=\widetilde\nu_K\{Z_T\ne\emptyset\}.
\end{equation}
These formulas are the partition of histories according to their
first non-small-field outcome, followed by the exact sum over all
extensions of each prefix. They are valid independently of the
special source-dependent representation of $\Phi_t$.

\begin{proposition}\label{prop:ward-fixed-volume-rarity}
Use the actual real history masses of Theorem~\YMMref{D}{8} on an
admissible lattice. Fix $0<\eth<1$, put $c'=(1-\eth)c_\star$, and
assume the case-II absorption and summability thresholds of
Proposition~\YMMref{67.6}{591} at $(c_\star,\eth)$, as well as
$y_K\ge\max\{Y_1,2\log(1+\bar B)\}$, with $Y_1$ from Lemma~\YMMref{67.1}{587}. These are the hypotheses used for the
calibrated history tail in Lemma~\YMMref{115.1}{805}; the couplings
here may be any admissible running sequence. Define
\begin{equation}\label{eq:ward-volume-constants}
 \begin{split}
 \mathsf B(x)&=\mathsf C_{\rm maj}^\Sigma+2+c_{\rm low}+\tfrac92\log x,
       \qquad \mathsf C_{\rm maj}^\Sigma=\mathsf C_{\rm maj}+(6L^3+1)\log2,\\
 \mathsf C_{\rm ag}(q)&=\exp\left[\frac{(3+q)^2(\log L)^2}
                   {2c'A_0^2(\log(L/2))^2}\right]\qquad(q\ge0).
 \end{split}
\end{equation}
We also write $\mathsf C_{\rm ag}(q,\eth)$ when the dependence on $\eth$ needs to be displayed. For $1\le t\le k\le K$, with $N_t=\ell^3L^{3(K-t)}$,
\begin{equation}\label{eq:ward-fixed-volume-tail}
 0\le w_t\le\widetilde\nu_k\{Z_t\ne\emptyset\}
 \le2M^{-3}N_te^{-c'p_t^2}e^{\mathsf B(x_K)\ell^3}.
\end{equation}
If $\mathsf B(x_K)\ell^3\le c'p_K^2/2$, one may replace its last two
exponentials by $e^{-c'p_t^2/2}$.
Write $P_t=\widetilde\nu_K\{Z_t\ne\emptyset\}$. Then
$m_t(\mathcal V_t)/Z=1-P_t$ and, for every $q\ge0$,
\begin{equation}\label{eq:ward-weighted-tail}
 \sum_{t'=1}^t L^{q(K-t')}P_{t'}
 \le4M^{-3}\ell^3e^{\mathsf B(x_K)\ell^3}\mathsf C_{\rm ag}(q+1)
              L^{-(K-t)}e^{-c'p_K^2/2}.
\end{equation}
If the stronger condition $2\mathsf B(x_K)\ell^3\le c'p_K^2/2$ holds,
then the probability in \eqref{eq:ward-fixed-volume-tail} is also
at most $L^{3(K-t)}e^{-c'p_t^2/2}$; its weighted sum is at most
$2\mathsf C_{\rm ag}(q+1)L^{-(K-t)}e^{-c'p_K^2/2}$.
Both weighted-sum bounds also hold with $w_{t'}$ in place of $P_{t'}$.
All history probabilities of an event determined through level $t$
are unchanged when the terminal level $k\ge t$ is changed.
\end{proposition}

\begin{proof}
We give the dependence on the sequence explicitly. Proposition~\YMMref{70.1}{605}
and equation~\YMMref{(1485)}{605} bound the real masses by the
summation input of Proposition~\YMMref{67.6}{591}. Its proof and
the proof of Proposition~\YMMref{67.5}{590} bound the normalized
weight for an exact set $W_j$ of creations by
\[
 e^{\mathsf C_{\rm maj}N_K}
 \prod_j(2^{-n_j^{\rm bit}}\zeta_j)^{\#W_j},\qquad
 \zeta_j=e^{-c'p_j^2-2p_j}.
\]
The multiplicity estimate~\YMMref{(1447)}{589} cancels the binary
factors even if a creation cell is prescribed. Since $Z_t$ is
nonempty exactly when a creation occurs at some $j\le t$, the union
bound in this finite subset product gives
\[
 \sum_{h:Z_t\ne\emptyset}e^{-E_{(K)}}\mathsf t(h)
 \le\Pi_K(\mathsf C_{\rm maj},c')
                   \sum_{j\le t}\#\mathcal P_j e^{-c'p_j^2}
 \le2M^{-3}N_te^{-c'p_t^2}
                         \Pi_K(\mathsf C_{\rm maj}^\Sigma,c').
\]
The last sum is the stopped summability estimate~\YMMref{(1451)}{591},
under the threshold assumed in the statement. Dividing by the
partition lower bound in Proposition~\YMMref{70.2}{606}, and using
$\Pi_K(\mathsf C_{\rm maj}^\Sigma,c')\le
e^{(\mathsf C_{\rm maj}^\Sigma+2)N_K}$, proves
\eqref{eq:ward-fixed-volume-tail}. This proof uses the admissible
flow bound, not an exact geometric line. Consistency follows by
summing the extensions of a prefix; the small-field ceilings are
part of the construction hypotheses. Equation~\eqref{eq:ward-first-creation}
gives the assertion for $w_t$.

For $a=K-t$, the inequality $y_t-y_K\ge a\log(L/2)$ and the first
three terms of the power expansion, with $y_K\ge1$ and
$P=2\mathfrak p_0\ge2$, imply
$p_t^2\ge p_K^2+A_0^2a^2(\log(L/2))^2$. Completing the square gives
\begin{equation}\label{eq:ward-age-square}
 L^{(3+q)a}e^{-c'p_t^2/2}
 \le \mathsf C_{\rm ag}(q)e^{-c'p_K^2/2}.
\end{equation}
Use $q+1$ and sum $\sum_{a\ge K-t}L^{-a}\le2L^{-(K-t)}$ to
obtain \eqref{eq:ward-weighted-tail}. Monotonicity $p_t\ge p_K$
absorbs the volume exponential under the weaker small-side
condition. Under the stronger condition use also
$2M^{-3}\le1$ and $\ell^3\le e^{\mathsf B(x_K)\ell^3}$, since
$\mathsf B(x_K)\ge1$. This gives the final assertions.
\end{proof}

For a first-passage datum at $g$, $x_K\le X_g$, so the preceding
bounds remain valid with $\mathsf B(X_g)$ in place of $\mathsf B(x_K)$. The canonical
choice takes its coupling ceiling below the fixed thresholds in
this proposition. Its admissibility on all levels is proved in
Theorem~\YMMref{D}{8}, and the first-passage and stopped-data
conditions in Proposition~\YMMref{74.1}{615} and
Lemma~\ref{lem:setup-stopped-arithmetic}. On a small side the
specialization is Lemma~\YMMref{125.3}{841}; the least admissible
side is small below the ceiling of Lemma~\YMMref{126.3}{846}.


\section{Source systems, pressure comparison and dilation}
\label{sec:source-system-complements}

The source estimates used below require one choice of constants for all
cutoffs and admissible sides. Fix the constants of
Theorem~\YMMref{AD}{607} at source radius \(2\varrho\), and use
\(C_g\) from \eqref{eq:setup-pressure-constants}. Put
\[
 C_{\rm gap}=2C_{\rm tt}C'_{\rm tt}E_0,\qquad
 G_g=C_{\rm S}+C_g+C_{\rm gap},
\]
where \(C_{\rm tt},C'_{\rm tt}\) are the constants of
Proposition~\YMMref{74.5}{617}. As in the proof of
Theorem~\YMMref{AD}{607}, choose \(C_{\rm S}\ge C_{\rm gap}\).

\begin{lemma}\label{lem:canonical-source-system}
For every first-passage datum at \(0<g\le\gamma\), the canonical
source energies on every \(T_\ell\), \(\ell\in\mathfrak L(g)\),
form the system of Section~\YMMref{88}{685}, with constants
\[
 (\rho_*;C_e,\kappa_e;E_+,C_{\rm low})
       =(2\varrho;C_{\rm S},\kappa;C_{\rm S},C_g).
\]
The same constants apply on all these sides. Their vacuum constants
and pressures satisfy
\begin{equation}\label{eq:source-system-pressure}
 \left|\ell^{-3}E_{K,\ell,\beta}
       -(\ell')^{-3}E_{K,\ell',\beta}\right|\le C_{\rm gap},
 \qquad |p_\ell-p_{\ell'}|\le G_g.
\end{equation}
Here \(C_g\le C_{\rm S}(1+\log g^{-2})\). The constants
\(C_{\rm S},\kappa,E_0,C_{\rm tt},C'_{\rm tt}\) are fixed before
the source radius and the coupling; \(C_g\) and \(G_g\) also depend on \(g\).
\end{lemma}
\begin{proof}
The source ball of Theorem~\YMMref{AD}{607} contains every
\(\sum_i z_if_i\) with
\(\sum_i|z_i|\|f_i\|_\sharp<2\varrho\), and is preserved by
conjugation and lattice isometries. It also contains the source
obtained by putting two such cell sources on the even and odd
reflected cells. On each cell the supremum and Lipschitz bounds
are those of one source; a segment crossing cell walls splits into
segments on which the sources vanish at the intervening walls.
This proves the same bounds globally, as in Section~\YMMref{89}{688}.
Restriction of holomorphy to these finite-dimensional source
families gives the holomorphy required in Section~\YMMref{88}{685}.

Theorem~\YMMref{AD}{607} supplies the local energies, their
exponential decay, their vanishing at zero, and the factorization
of the sourced partition function. Set \(e_0=0\) if the system is
indexed from zero. Lemma~\ref{lem:setup-pressure-stability} gives
the stated vacuum floor with \(C_g\). Proposition~\YMMref{74.1}{615}
checks the common-sequence and room hypotheses of
Proposition~\YMMref{74.5}{617}, whose equation~\YMMref{(1541)}{618}
gives the first estimate in \eqref{eq:source-system-pressure}.
Write \(p_\ell=\ell^{-3}E_{K,\ell,\beta}+\vartheta_\ell\), where
\(-C_g\le\vartheta_\ell\le C_{\rm S}\), and subtract for the two
sides. This proves the second estimate. The assertions about the
constants follow from those same results.
\end{proof}

For clarity, the finite reflection positivity used here includes
both coordinate site mirrors and coordinate bond mirrors. Indeed,
\(\Gamma=\ell\mathbb Z^3\) is adapted to every coordinate normal
with height period \(\ell\). Proposition~\YMMref{91.4}{700} applies
at every \(c\in(a_K/2)\mathbb Z\). A function of bonds in either
closed half of the torus lifts to a function of bonds in the closed
slab \(c\le x_\mu\le c+\ell/2\): choose a lift of each bond in
that slab. Periodization commutes with the reflection, so the
positive form descends to the finite Wilson measure. Positivity
of this sesquilinear form gives its Cauchy--Schwarz inequality.
This also applies to polynomials in centred observables supported
in the open half, by their reflection covariance.

\begin{lemma}\label{lem:dilated-torus-cumulant-bound}
Let \(n\ge2\), let \((K_0,\beta)\) be a first-passage datum at
\(0<g\le\gamma\), let
\(q\in\mathbb Z\), and suppose that \(L^q\ell\in\mathfrak L(g)\).
Set \(K=K_0+q\). If a real Lipschitz tuple on \(T_\ell\) has
one pair of supports at distance at least \(d>0\), then its
cumulant is bounded by
\[
 |S^T_{n;K,\ell,\beta}(\underline f)|
 \le \max\{1,L^q\}^n B_n(L^qd,L^q\ell,L^qR)
                    \prod_i\|f_i\|_\sharp.
\]
Here \(B_n\) can be any applicable bound from
Propositions~\YMMref{88.3}{688}, \YMMref{89.4}{693} and
\YMMref{89.5}{694}, evaluated on the dilated torus and tuple.
For the last two bounds every support must lie in a sup-ball of
radius \(R\); the middle bound also requires
\(L^q\ell\ge16n(L^qR+1)\).
\end{lemma}
\begin{proof}
Proposition~\YMMref{119.1}{815} identifies the original cumulant
with \(L^{qn}\) times the cumulant at cutoff \(K_0\), side
\(L^q\ell\), and test functions \(f_i^{[q]}\). Separation and
support radii scale by \(L^q\), while
\(\|f_i^{[q]}\|_\sharp\le\max\{1,L^{-q}\}\|f_i\|_\sharp\).
Apply the indicated bound there and use
\(L^q\max\{1,L^{-q}\}=\max\{1,L^q\}\).
\end{proof}

\subsection{The conjugation specific to two by two matrices}
\label{sec:two-by-two-conjugation}

The complex conjugation used for \(\mathrm{SU}(2)\) preserves the
matrix norm. This fact is needed when comparing the domains of a
local term and its conjugate; it is not a general property of
\(X\mapsto(X^*)^{-1}\) in higher matrix dimension.

\begin{lemma}\label{lem:two-by-two-conjugation}
Put \(\epsilon=\left(\begin{smallmatrix}0&1\\-1&0\end{smallmatrix}\right)\)
and \(\Theta(X)=\epsilon\overline X\epsilon^{-1}\).
Then \(\Theta\) is a conjugate-linear multiplicative involution
on complex two by two matrices and
\[
 |\Theta(X)|=|X|,\qquad
 \Theta(X)=\overline{\operatorname{tr}X}\,I-X^*.
\]
For \(U\in\mathrm{SL}(2,\mathbb C)\) and traceless \(A\),
\[
 \Theta(U)=(U^*)^{-1},\qquad \Theta(iA)=iA^*.
\]
In particular \(|(U^*)^{-1}-I|=|U-I|\).
\end{lemma}
\begin{proof}
Entrywise conjugation and conjugation by \(\epsilon\) preserve
products. Since \(\epsilon\) is real and \(\epsilon^2=-I\),
their composite squares to the identity. Entrywise conjugation
preserves the operator norm, as does conjugation by the unitary
matrix \(\epsilon\). Direct multiplication gives
\(\epsilon Y^{\mathsf T}\epsilon^{-1}=(\operatorname{tr}Y)I-Y\).
Take \(Y=X^*\) to obtain the displayed formula for \(\Theta(X)\).
For determinant-one \(Y\), the same two by two identity gives
\(Y^{-1}=(\operatorname{tr}Y)I-Y\); apply it to \(Y=U^*\).
For traceless \(A\), the formula gives \(\Theta(iA)=-(iA)^*=iA^*\).
Finally apply norm preservation to \(U-I\).
\end{proof}

\section{Local domains and finite differences}
\label{sec:analytic-domain-complements}

The estimates below retain polynomial factors in a domain's size while
summing over domains that join separated supports. We use the domains
$\mathcal D_j$ and their linear sizes from Section~\YMMref{45}{351}.
Lengths in the next lemma are measured in level-$j$ units, so a cube
has side $M$. The same assertions apply to finite localization domains
of the infinite lattice.

\begin{lemma}\label{lem:analytic-domain-geometry}
Set
\[
 \kappa_{\mathcal D}=\log9216,\qquad c_{\rm E}=\log18432,
 \qquad C_{\rm E}=4096,\qquad
 C_{{\rm E},p}=2048\sum_{n\ge0}(n+2)^pe^{-n}\quad(p=0,1,\ldots).
\]
For every site $z$, $D\ge0$ and $\kappa>\kappa_{\mathcal D}+2$,
\begin{equation}\label{eq:analytic-domain-tail}
 \sum_{\substack{X\in\mathcal D_j\colon z\in X\\d_j(X)\ge D}}
 e^{-\kappa d_j(X)}(d_j(X)+2)^p\le C_{{\rm E},p}e^{-D}.
\end{equation}
For all $x,z\in X$,
\begin{equation}\label{eq:analytic-domain-diameter}
 |x-z|_\infty\le M(d_j(X)+2).
\end{equation}
If $X_0\subset X$ are domains, then
\begin{equation}\label{eq:analytic-domain-extension}
 d_j(X)\le d_j(X_0)+\#X-\#X_0.
\end{equation}
Here $\#X$ counts cubes. On the torus, the distance is periodic.
\end{lemma}

\begin{proof}
The cube side $M$ is odd, its centre is an integer site and its faces
lie at half-integer coordinates. Every site therefore belongs to
exactly one cube. Proposition~\YMMref{45.1}{353}, equation~\YMMref{(683)}{353},
bounds the number of domains containing that cube and having size $n$
by $2048\cdot9216^n$. For $n\ge D$,
$9216^ne^{-\kappa n}\le e^{-2n}\le e^{-D}e^{-n}$.
Multiplication by $(n+2)^p$ and summation proves
\eqref{eq:analytic-domain-tail}; the defining series converges by the
ratio test. The same rooted count on the infinite lattice follows from
the identical finite encoding: a walk of length $2n$, beginning at one
of the eight corners of the root cube, visits a spanning connected
edge set; its at most $n+1$ vertices meet at most $8(n+1)$ cubes.
There are at most $8\cdot6^{2n}2^{8(n+1)}$ such encodings.
The estimate for the infinite sum follows first for finite partial
sums, then by monotonicity.

Choose a connected edge set with $d_j(X)$ edges meeting every cube of
$X$. The cubes containing $x,z$ meet that set at vertices $v_x,v_z$.
A path in the set from $v_x$ to $v_z$ uses at most $d_j(X)$ edges;
$x,v_x$ and $z,v_z$ are each at distance at most $M$. The triangle
inequality proves \eqref{eq:analytic-domain-diameter}, also after
projection to the torus.

Enumerate the cubes of $X\setminus X_0$ so that each new cube is adjacent
to a previous one; a rooted spanning tree of the cube graph supplies
such an order. A spanning edge set of the previous domain meets the
adjacent old cube at a corner. That corner is either on the common
face or one cube edge from it. Adding that edge, if needed, makes the
set meet the new cube and costs at most one. Iteration proves
\eqref{eq:analytic-domain-extension}.
\end{proof}

\begin{lemma}\label{lem:near-domain-geometry}
Let $j$ be a level with $m'<m+K-j$, let $D_j=L^{j-K}$,
and put $W=MD_j$. Thus the cubes of level $j$ have physical
side $W$ and the torus has side $\ell=2L^m$. Let
$\operatorname{pr}:\mathbb R^3\to T_\ell$ be the covering map.
A finite localization domain $\widehat X$ in the infinite
level-$j$ cube lattice is called near if it lies in a cuboid
whose three side lengths are strictly less than $\ell/2$.
A torus domain is non-wrapping if it is the image of a near set.
Then the following hold.
\begin{enumerate}
\item The image of every infinite localization domain belongs to
$\mathcal D_j$. Every near lift of a non-wrapping domain belongs
to the infinite domain family. Its cubes correspond bijectively
to those of the torus domain, adjacency is preserved in both
directions, and $d_j(\widehat X)=d_j(X)$.
\item All near lifts of one domain are its translates by
$\ell\mathbb Z^3$. They are pairwise disjoint; specifying one
lift of a point of the domain specifies exactly one near lift
containing that point.
\item Every infinite domain $\widehat X$ lies in a cuboid with
side lengths at most $W(d_j(\widehat X)+2)$. A torus domain with
$W(d_j(X)+2)<\ell/2$ is non-wrapping. Consequently a wrapping
domain, or an infinite domain which is not near, has size at least
$\ell/(2W)-2$.
\end{enumerate}
\end{lemma}

\begin{proof}
The hypothesis gives $\ell/W=2L^{m+K-j-m'}\ge2L$.
Projection carries cubes to cubes and shared faces to shared
faces. It therefore carries a finite wall-connected union to
a wall-connected union. On a near set projection is injective.
The coordinate differences of two of its points have magnitude
less than $\ell/2$; a different lift of such a difference differs
by a period. Each cube of the image lifts to a whole cube:
the difference between the inverse chart on that cube and any
one of its cube lifts is continuous and period-valued, hence
constant. Thus the inverse chart preserves the faces, edges
and corners of its cubes, as well as their adjacency. Lifting
a wall chain shows that a near lift is a localization domain.
Connected edge sets meeting every cube are carried in both
directions with their number of edges unchanged. Taking the
minimum of that number proves the equality of the linear sizes.

To compare two near lifts, translate one by a period so that
one pair of points over the same torus point agrees. For any
other pair over the same point, their displacement is a period.
Each coordinate of this displacement is the difference of two
displacements from the matched points, each of magnitude less
than $\ell/2$. It has magnitude less than $\ell$, so is zero.
The lifts therefore agree after that translation. Distinct
period translates are disjoint, since two points in a near
cuboid cannot differ by a nonzero period. Translating a chosen
lift to contain the prescribed point proves existence and
uniqueness in the second assertion.

For an infinite domain choose a connected edge set with
$d=d_j(\widehat X)$ edges meeting every cube. If $d=0$ this
set is one corner. Otherwise enumerate its edges so that each
edge after the first meets an earlier edge. Adding an edge of
length $W$ increases each coordinate span by at most $W$.
The edge set is therefore contained in a cuboid with side
lengths at most $dW$. Every cube of the domain meets it and
lies within sup distance $W$ of it. Enlarging the cuboid by
$W$ on both sides proves the bound $W(d+2)$.

The same construction lifts a small torus domain. Enumerate a
minimal connected edge set as above, lift the first edge, and
translate each later edge by a period so that it meets the
chosen lift of an earlier edge. Their union projects onto the
whole edge set and lies in a cuboid of side lengths at most
$dW$. For every cube of the domain choose a lift containing
a lifted point where that cube meets the edge set. The union
of these cubes projects onto the domain and lies in the
$W$-enlargement of the cuboid. If $W(d+2)<\ell/2$, this is a
near lift. The last two lower bounds are the contrapositives
of the cuboid bounds.
\end{proof}


The brackets of equation~\YMMref{(806)}{392} are alternating finite
differences. Their useful support property has an elementary proof
which we record to make the localization of a shifted local term
explicit.

\begin{lemma}\label{lem:analytic-rooted-bracket}
Let $\mathsf P$ be a finite set of cubes and let $\mathsf N_0\subset
\mathsf P$ be nonempty and wall-connected. Suppose that, at every
vertex $\mathbf1_W$, $F(\mathbf1_W)$ depends only on the wall-components
of $W$ meeting $\mathsf N_0$: deleting all other components does not
change its value. For $\mathsf N_0\subset W\subset\mathsf P$,
\begin{equation}\label{eq:analytic-bracket-resummation}
 F(\mathbf1_W)-F(0)
   =\sum_{\mathcal K\subset W\setminus\mathsf N_0}
                              F^*[\mathcal K\mid\mathsf N_0].
\end{equation}
Every term whose union $\mathcal K\cup\mathsf N_0$ is not
wall-connected is zero.
\end{lemma}

\begin{proof}
The resummation is equation~\YMMref{(809)}{392}: expand the alternating
sums, and use $\sum_{A\subset B}(-1)^{\#A}=0$ for nonempty $B$.
For the second assertion choose a wall-component $c$ of
$\mathcal K\cup\mathsf N_0$ disjoint from $\mathsf N_0$ and a cube
$\Delta\in c$. For every $S\subset\mathcal K\setminus\{\Delta\}$,
all components within $c$ are unrooted, even after $\Delta$ is added.
Thus $F(\mathbf1_{S\cup\mathsf N_0})=
F(\mathbf1_{S\cup\{\Delta\}\cup\mathsf N_0})$.
The two terms have opposite signs in the bracket. Pairing them proves
its vanishing. The same cancellation applies to $F-F(0)$.
\end{proof}

\section{Local analytic cores and substitution}
\label{sec:analytic-core-complements}

We use the simple domains and hereditary analytic cores of
Definition~\YMMref{49.1}{374}, with the geometric and regularity
thresholds of Proposition~\YMMref{51.9}{402}. These thresholds, together
with the positivity of the local gauge constants, will be called the
tower hypotheses below. In particular, the four geometric parameters
are powers of $L$, $L\ge13$ is odd, $R_1\ge L$,
$M_1\ge R_1\mathfrak M$, and $M\ge L^4M_1$; the background and response
thresholds are those fixed at the beginning of Section~\YMMref{51.5}{398}.
Write $a^\sharp\le1/60$ and $K^\sharp_{\rm tw}\ge1$ for its tower
radius and bound. Every use below keeps the domain and subtower
restrictions in that proposition. A domain on the infinite lattice
is finite; the same statement there is read in a containing lattice
chart, with the chart compatibility established in
Section~\YMMref{52.5}{425}, equation~\YMMref{(926)}{425}.

The triple is $P=(U,A',\mathbf J)$, where $U$ is real and
$\mathbf p(P)=(e^{i\eta_jA'}U,\mathbf J)$ is its represented pair.
Its independent bond coordinates and gauge action are those preceding
Definition~\YMMref{49.1}{374}. For a real bare configuration $V$, its
pair at level $j$ is $(V,J^{(j)}(V))$ restricted to the incident bonds
of the domain. A local term is evaluated at this pair only when that
pair belongs to its core. This convention does not assert that an
arbitrary real configuration belongs to the core.

For a block field $W$, the background $U_j(W)$ denotes an element of
the unique regular critical orbit over its determining datum, when
that orbit exists (Definitions~\YMMref{9.3}{40} and~\YMMref{9.5}{41}).
On the full torus the reduced determining set consists exactly of
the terminal bonds. Gauge invariant expressions have the same value
at every representative. The total of a pointed family is
equation~\YMMref{(935)}{427}; it is defined only if the orbit exists and
every pair read by the sum belongs to the respective core.

\begin{lemma}[Open pair domains]\label{lem:analytic-open-pairs}
Under the tower hypotheses, let $X$ be a localization domain and
$b_0,b_1,c_0>0$, $b_0+b_1\le a^\sharp$. For fixed real $U$, the set
\[
 \mathcal O_U=\{(A',\mathbf J):(U,A',\mathbf J)
                   \in\mathcal U_j^\circ(X;b_0,b_1,c_0)\}
\]
is open in the independent bond coordinates. The map
$\Psi_U(A',\mathbf J)=(e^{i\eta_jA'}U,\mathbf J)$ is biholomorphic
from this set onto an open set of pairs. The union of these pair
domains over $U$ is open. A function on the core is analytic and
depends only on the represented pair if and only if it descends to
a holomorphic function on that union.
\end{lemma}

\begin{proof}
The conditions on the real factor and its real subtowers are fixed
when $U$ is fixed. If any fails, the slice is empty. Otherwise the
simple-domain conditions on $A'$ and $\mathbf J$ are finitely many
strict inequalities between continuous functions. On this open set,
Proposition~\YMMref{51.9}{402}(f) constructs every required complex
subtower holomorphically. There are finitely many subtowers, bonds
and plaquettes in $X$. Their strict curvature and current inequalities
therefore define an open subset. The holomorphic-continuation clause
of the core is preserved on this subset by the same tower construction.

On the slice, $\eta_j|A'(b)|<\eta_j b_1<1/2<\log2$. Thus the
logarithm at $I$ gives the holomorphic inverse
\[
 (\mathbf U,\mathbf J)\longmapsto
 \bigl((i\eta_j)^{-1}\log(\mathbf U U^{-1}),\mathbf J\bigr)
\]
of $\Psi_U$. The exponential and this logarithm identify open
neighbourhoods in the bond-coordinate spaces, so the image of the
open slice is open. Their union is open as well. Composition proves
one direction of the final assertion. Conversely a function of the
pair gives a single well-defined function on the union. On each open
image it is the analytic slice function composed with the displayed
inverse. These local holomorphic functions agree on intersections,
and hence define the required holomorphic function.
\end{proof}

The chart at the trivial pair is $P[v]=(I,A,\mathbf J)$, where
$\|v\|_X=\max\{\sup|A|,\sup|\nabla A|,\sup|\mathbf J|\}$ and
only differences defined on the incident bonds enter the maximum.
Equation~\YMMref{(933)}{427}, proved in Section~\YMMref{52.6}{426}, gives
\[
 r_{\rm us}=\min\{b_1,a^\sharp/4,b_0/(4K^\sharp_{\rm tw})\},
 \qquad P[\{\|v\|_X<r_{\rm us}\}]
                         \subset\mathcal U_j^\circ(X;b_0,b_1).
\]
In particular this is a genuine ball in the full core, not just in
the simple domain.

\begin{proposition}[Three pairs of radii]
\label{prop:analytic-three-radii}
Assume the tower hypotheses and let $X$ be a level-$(j-1)$ domain,
with lengths measured in level-$j$ units. Let positive radii satisfy
\[
\begin{gathered}
 \beta_0<b_0\le\gamma_0,\quad\beta_1<b_1\le\gamma_1,
 \quad\gamma_0+\gamma_1\le a^\sharp,\\
 \beta_0+\tfrac94s\le\gamma_0,\quad
 \beta_1+\tfrac{11}{10}s\le\gamma_1,\\
 \beta_0+\tfrac94b_2\le b_0,\quad
 \beta_1+\tfrac{11}{10}b_2\le b_1,\\
 2K^\sharp_{\rm tw}(\gamma_0+\gamma_1^2)b_2
                                      \le(b_0-\beta_0)s.
\end{gathered}
\]
If $P=(U,A',\mathbf J)\in\mathcal U_j^\circ(X;\beta_0,\beta_1,c_0)$,
$\mathbf U=e^{i\eta_jA'}U$, and
$|\mathbf A|,|\nabla_{\mathbf U}\mathbf A|<b_2$,
$|\mathbf J+\mathbf K|<c'_0$, then
\[
 P\dotplus(\mathbf A,\mathbf K)
 =\bigl(U,(i\eta_j)^{-1}\log(e^{i\eta_j\mathbf A}e^{i\eta_jA'}),
                  \mathbf J+\mathbf K\bigr)
 \in\mathcal U_j^\circ(X;b_0,b_1,c'_0).
\]
Its pair is $(e^{i\eta_j\mathbf A}\mathbf U,\mathbf J+\mathbf K)$.
\end{proposition}

\begin{proof}
The logarithm and the simple-domain bounds follow from the pointwise
increment calculation in the proof of Proposition~\YMMref{51.10}{403}:
the potential and its first derivative increase by less than
$11b_2/10$, and the plaquette bound by less than $9b_2/4$.
The real factor does not change. Put $R=s/b_2$. The last inequality
and $K^\sharp_{\rm tw}\ge1$ imply $R>2$. For $|z|<R$ the same
increment calculation, with $z\mathbf A$ in place of $\mathbf A$,
places the substituted field in the simple domain at outer radii
$(\gamma_0,\gamma_1)$. Proposition~\YMMref{51.9}{402}(f) therefore
constructs all its subtowers holomorphically on that disc.

Consider a level-$p$ subtower and set $\lambda_p=L^{-(j-p)}$.
After its curvature and current are multiplied by $L^{2p}$ and
$L^{-2(n-p)}$, respectively, each is bounded on the disc by
$\lambda_pK^\sharp_{\rm tw}(\gamma_0+\gamma_1^2)$.
At $z=0$ its norm is strictly below $\lambda_p\beta_0$.
The Schwarz estimate for the difference from its value at zero,
as used in the same proof of Proposition~\YMMref{51.10}{403}, bounds
its increment at $z=1$ by
\[
 \frac2R\lambda_pK^\sharp_{\rm tw}(\gamma_0+\gamma_1^2)
                         \le\lambda_p(b_0-\beta_0).
\]
This proves the strict target inequality for each complex subtower.
The real subtower inequalities are unchanged and their old radius
$\beta_0$ is smaller than $b_0$. All clauses of the target core follow.
\end{proof}

For the passage from half the fixed radii to three quarters, take
$\beta_i=\alpha_i/2$, $b_i=3\alpha_i/4$, $\gamma_i=\alpha_i$,
and $s=\min\{2\alpha_0/9,5\alpha_1/11\}$. A positive $b_2$ satisfying
\[
 b_2\le\min\left\{\alpha_0/9,5\alpha_1/22,
 \frac{\alpha_0s}{8K^\sharp_{\rm tw}(\alpha_0+\alpha_1^2)}\right\}
\]
satisfies all four displayed conditions. This choice depends only on
the fixed radii and the tower constants.

\subsection{Regular gauges, background response and weight members}
\label{sec:regular-background-complements}

We record three consequences used when evaluating an interaction at
an earlier level. The first puts a regular field in a local gauge.
The second controls the change in its background when the datum moves.
The third identifies a plaquette weight with an actual local family.

\begin{definition}\label{def:level-regularity}
Put $\eta_j=L^{-j}$. A real bare configuration $U$ is regular
at level $j$ with radius $\alpha$ on $S$ if
\[
 |\partial U(p)-1|<\alpha\eta_j^2,\qquad
 |J^{(j)}(U)(b)|<\alpha
\]
for every bare plaquette with a corner in $S$ and every positive
bare bond with an endpoint in $S$, respectively. The current is
computed with spacing $\eta_j$, as in equation~\YMMref{(69)}{38}.
For $j\le n$, changing units gives
$J^{(j)}=L^{-3(n-j)}J^{(n)}$. Thus a field regular at level $n$
with radius $\varepsilon$ is regular at level $j$ with radius
$\varepsilon L^{-2(n-j)}$ on the same physical set.
\end{definition}

Write $c_{*,\mathrm{rg}}=c_*$ and use $B_*$ for the constants
in Theorem~\YMMref{F}{199}. These are the constants of the local
regular-gauge theorem throughout the manuscript.

\begin{lemma}[A gauge on a regular cube]\label{lem:regular-cube-gauge}
Let $1\le j\le K$, and measure lengths at level $j$. Let $Q$ be
a closed cube of integer side $l$, made of unit cells, and let
$\widetilde Q$ be its concentric enlargement of side $l+4$.
Suppose
\[
 l+4<L^{m+K-j},\qquad 0<\alpha\le1,
 \qquad 2l\alpha\le c_{*,\mathrm{rg}}.
\]
If $U$ is regular with radius $\alpha$ on $\widetilde Q$, there
are a real gauge $u$ on its bare sites and a real one-form $A$ on
its inner bonds such that $U^u=e^{i\eta_j A}$ there and
\[
 |A_\mu(x)|,\ |\nabla^{\eta_j}_\nu A_\mu(x)|
       <2B_*l\alpha
\]
at every positive bare bond with an endpoint in $Q$.
\end{lemma}

\begin{proof}
The faces of $Q$ have half-integer coordinates, so the enlargement
by two units is again a region of unit cells. Set $\eta=\eta_j$.
Every plaquette in the star of a bond having an endpoint in
$\widetilde Q$ has that endpoint as a corner. The current identity
used in the proof of Lemma~\YMMref{113.12}{798} therefore gives
\[
 |D_U^{\mathsf T}\partial U(b)|
 \le\eta^{-1}(\alpha\eta^2)^2+\eta^2|J^{(j)}(U)(b)|
 <\alpha\eta^2(1+\alpha\eta)<2\alpha\eta^2,
\]
since $\eta\le1/13$. The plaquette bound also holds at radius
$2\alpha$. Apply Theorem~\YMMref{F}{199} at that radius.
Its smallness condition is $2l\alpha\le c_{*,\mathrm{rg}}$,
and its conclusion is precisely the displayed bound. If its gauge
is denoted by $u_0^{-1}$, take $u=u_0^{-1}$.
\end{proof}

Local terms depend on the represented pair and are gauge invariant
(Definition~\YMMref{49.4}{375}). Thus, when $U^u=e^{i\eta_jA}$,
their evaluation in this gauge uses
\[
 v=(A,\operatorname{Ad}(u)J^{(j)}(U)),
\]
not an untransformed current. Unitary conjugation preserves its norm.
On the usable ball of Section~\ref{sec:analytic-core-complements}
this is a holomorphic coordinate for the same value of the local term;
the trivial pair corresponds to $v=0$ and belongs to that ball.

For background response we use the same representative constants as
in Section~\YMMref{113}{786}, immediately before
equation~\YMMref{(2147)}{794}:
\[
 K_{\rm B}=K^{\rm s}_{\rm B}(\rho_{\rm s}),\qquad
 C_H=4C_{\rm db}\max\{2,K_{\rm B}\}.
\]
Here $r^{\rm s}_{\rm B}=r_{\rm G}(\rho_{\rm s})$ and
$K^{\rm s}_{\rm B}$ is the curvature-response bound obtained
in the proof of Theorem~\YMMref{E}{10}. The admitted datum radius
satisfies $r_{\rm B}\le r^{\rm s}_{\rm B}$.

\begin{lemma}[Change of the regularity number]\label{lem:background-curvature-response}
Use a constant determining sequence at level $j\ge1$, with the
geometric and numerical hypotheses of Theorem~\YMMref{E}{10}.
If $\varpi(V_0)<\bar\varepsilon$ and $B$ is real with
\[
 \sup|B|<\min\{r_{\rm B},r^{\rm s}_{\rm B}/2\},
\]
the backgrounds over $V_0$ and $e^{iB}V_0$ exist, depend
real-analytically on $B$, and their canonical representatives satisfy
\begin{equation}\label{eq:background-curvature-response}
 \eta_j^{-2}|\partial U_j(e^{iB}V_0)(p)
                   -\partial U_j(V_0)(p)|
 \le2K_{\rm B}\sup|B|.
\end{equation}
Consequently
$\mu_j(e^{iB}V_0)\le\mu_j(V_0)+2K_{\rm B}\sup|B|$.
\end{lemma}

\begin{proof}
Theorem~\YMMref{E}{10} supplies the actual holomorphic
representative $\Phi_{V_0}$ on the row ball of radius
$r^{\rm s}_{\rm B}$. Its proof Section~\YMMref{38}{288}
bounds the holomorphic map
\[
 F_p(B)=\eta_j^{-2}
       [\partial\Phi_{V_0}(B)(p)-\partial U_j(V_0)(p)]
\]
by $K_{\rm B}r^{\rm s}_{\rm B}$ on that ball. There is no
layer weight here because the determining sequence is constant.
Cauchy's estimate bounds its derivative on the half ball by
$2K_{\rm B}$ in the row sup norm. Integrating along $sB$,
$0\le s\le1$, proves \eqref{eq:background-curvature-response}.
At real $B$ these are the canonical real representatives.
Taking the maximum over plaquettes proves the last assertion;
the number is independent of the representative by unitary gauge
invariance.
\end{proof}

\begin{lemma}[The local family of a plaquette weight]\label{lem:plaquette-weight-member}
Let $j\ge1$ be an admitted height and use the cores with radii
$(b_0,b_1)$. For a bare plaquette $p$, let $X(p)$ be the union
of the level-$j$ cubes containing a corner of $p$. For a complex
weight $w$ define
\begin{equation}\label{eq:plaquette-weight-member}
 \mathsf W_w(X,P)=\eta_j^{-1}\sum_{p:X(p)=X}w(p)a_p(\mathbf U),
 \qquad a_p(\mathbf U)=1-\tfrac12\operatorname{tr}\partial\mathbf U(p),
\end{equation}
where $(\mathbf U,\mathbf J)$ is the pair represented by $P$.
This is a local family, linear in $w$, vanishing at the trivial
pair, and
\[
 \|\mathsf W_w\|_{j,\kappa}\le6M^3b_0^2\sup_p|w(p)|.
\]
Adding it to the level-$j$ unpointed slot of a bare array adds
$\eta_j^{-1}\sum_p w(p)a_p(U)$ to its action, wherever that
action is defined, and changes no counterterm.
\end{lemma}

\begin{proof}
The cubes of $X(p)$ are either one cube, two adjacent cubes, or
four cubes meeting along an edge. Thus they have a common corner,
$d_j(X(p))=0$, and all bonds of $p$ lie in the domain.
The formula is holomorphic on every core slice, depends only on
the represented pair, and is gauge invariant. For $V\in
\mathrm{SL}(2,\mathbb C)$,
$1-\frac12\operatorname{tr}V=\frac12\det(V-1)$. Hence the
core's plaquette bound gives
\[
 |a_p(\mathbf U)|\le\tfrac12|\partial\mathbf U(p)-1|^2
                    <\tfrac12 b_0^2\eta_j^4.
\]
Each $X(p)$ has at most four cubes; there are at most
$12M^3\eta_j^{-3}$ bare plaquettes with that domain. Multiplication
by $\eta_j^{-1}$ gives the asserted norm bound, since the domain
has size zero. At the trivial pair every plaquette term vanishes.
In the array action, equation~\YMMref{(874)}{412}, each plaquette
occurs at exactly one $X(p)$. Summing proves the action identity.
\end{proof}


\subsection{Normalizers for overlapping entries}
\label{sec:overlapping-normalizers}

In the setting of Lemma~\ref{lem:canonical-source-system}, let
\(A_g=C_{\rm S}+C_g+C_{\rm gap}=G_g\). We use the following
normalizer and cumulant constants:
\begin{equation}\label{eq:overlap-bound-constants}
\begin{split}
 \mathsf Q_1(L)&=\exp\{A_gL^3+\Phi_1(2,1;L)\},\\
 Q^\natural(n,R)&=\exp\{A_gh(n,R)^3+\Phi_1(R,R+1;M)\},\\
 C_n^{\rm per}(d;R)&=(n/\varrho)^n
 \bigl[\log2\,(2Q^\natural(n,R)^n+2)^n+\mathsf U(R,d)\bigr].
\end{split}
\end{equation}
Here \(h(n,R)\) is the least power of \(L\) at least \(8n(R+1)\);
\(\Phi_1\) and \(\mathsf U\) are the functions in Section~\YMMref{89}{688},
with source-energy constants \((C_{\rm S},\kappa)\).
These are equations~\YMMref{(1844)}{708} and~\YMMref{(1847)}{709}
with the floor and gap constants of Lemma~\ref{lem:canonical-source-system}.

We will use one part of the normalizer proof for tuples whose supports
may overlap. Fix a board of side \(l\), and suppose each entry is
supported in a ball of radius \(R\) whose centre is at distance at
least \(R+D\) from the walls, where \(D>0\). On the least admissible
torus of side at least \(2l\), form the cell normalizers and set
\(E(z)=\sum_c E_c(z)\). Each \(E_c\) reads only source coordinates
assigned to cell \(c\). The construction in the proof of
Proposition~\YMMref{91.12}{708} gives, on the source polydisc of
radius \(\varrho\),
\begin{equation}\label{eq:overlap-graded-normalizer}
 |e^{-E(z)}|\Pi_\ell(z)
 \le\exp\{A_gl^3+\Phi_1(R,D;l)\}^{\#B(z)}
 \qquad(\ell\in\mathfrak L(g),\ \ell\ge2l).
\end{equation}
Indeed, the cell bound gives one factor
\(\exp\{A_gl^3+\#B_c(z)\Phi_1(R,D;l)\}\) for each cell containing
an active coordinate. The number of such cells is at most \(\#B(z)\),
and \(\sum_c\#B_c(z)=\#B(z)\). Multiplying proves the inequality.
Neither this count nor the holomorphy of \(E\) uses separation of
the supports. If at least two cells are occupied, its full mixed
derivative vanishes since every summand omits a coordinate.
Separation is needed only for the additional bound on that derivative
when all entries occupy one cell. The same normalizer works for
periodic Gibbs states by Proposition~\YMMref{91.11}{707}.

\section{Background estimates and the normalized step}
\label{sec:background-step-complements}

The background estimates used here are those for reduced determining
sets. The region sequence may have an arbitrary first member. It is
separated, its terminal lattice has at least $6R_1\mathfrak M$ sites
in each direction, and the geometric, regular-reference and datum
thresholds are those of Section~\YMMref{51.5}{398} and the localization
construction in Section~\YMMref{53.4}{445}. In particular we retain
\[
 0<\varepsilon_{\rm r}\le\bar\varepsilon,
 \qquad C_{\mathfrak P}\mathfrak M B_3\varepsilon_{\rm r}
          \le a_{\rm R},\qquad
 \varrho=\min\{r_{\rm B},\varrho_{\rm R}^0\},\qquad B_3\varepsilon_{\rm r}\le L^{-2}a_{\mathsf m}.
\]
Here $a_{\mathsf m}$ is the radius of Proposition~\YMMref{30.2}{238}; the last inequality ensures its curvature hypothesis for the regular backgrounds, whose bare curvature is at most $B_3\varepsilon_{\rm r}\eta_n^2$. All depth thresholds chosen from this regularity radius use the decreased value.
Theorem~\YMMref{E}{10} gives a single regular critical orbit and its
holomorphic axial representative. Lemma~\YMMref{51.6}{398} writes it as
$(e^{i\eta_n\mathcal H(B)}U_0)^{u_B}$ on $\sup|B|<\varrho$, with
$\mathcal H(0)=0$, $u_0=I$ and the stated Landau condition.

For the response constants in this paper we take
$B_5^{\rm e}\ge\max\{8,B_5\}$ and
$\delta_0^{\rm e}=\delta_{\rm lc}$ from equations~\YMMref{(996)}{446}
and~\YMMref{(997)}{446}. Thus, with
$m_y(B)=\sum_{c>0}e^{-\delta_0^{\rm e}d(y,y_c)}|B(c)|$,
\[
 \xi_y|\mathcal H(B)(b)|\le B_5^{\rm e}m_y(B),\quad
 \xi_y^2|\nabla_{U_0}\mathcal H(B)(x)|\le B_5^{\rm e}m_y(B),\quad
 |u_B(y)-I|\le4\kappa_{\rm av}B_5^{\rm e}m_y(B).
\]
The first two bounds also hold for differences at two real data,
with $m_y(B-B')$ on the right. The derivative estimate
\YMMref{(838)}{399} supplies the corresponding linear response and its
second covariant difference, where all incident plaquettes are inner.
These estimates hold on the specified datum ball; they are not
assertions for arbitrary block configurations. Here $\kappa_{\rm av}$ denotes the
joint contour constant called $\kappa_6$ in the localization construction.
When a bound for the derivative of the gauge at a top is written with
a single coefficient $k_6$, we take $k_6=4\kappa_{\rm av}$.

The localization estimate itself is Proposition~\YMMref{53.5}{447}.
Its hypotheses include the one- or two-component complement,
the four-layer collars and their level assignments, the separation
of different components, the joined sequence and its regular datum,
and $\varepsilon_{\rm r}\le\epsilon_{\rm lc}$. The definitions and
all those conditions are equations~\YMMref{(998)}{446}--\YMMref{(1000)}{446}
and the paragraphs preceding the proposition. We keep its wide
three-layer collars and its weighted distance to the tower rows.

\begin{lemma}[Iterated plaquette bound]\label{lem:background-iterated-average}
Let $p'$ be a plaquette at level $q\ge1$. At level $j\le q$, take the
union of the four $(q-j)$-blocks below its corners and let $\varpi_j$
be the largest plaquette curvature of $\mathcal M^jU$ within that
union, in all three orientations. Set
\[
 \theta_L=\exp\frac1{2(L^2-1)},\qquad
 c_2^{\rm SO}=\min\{\varpi'_L/2,L^2/(4C_0^{\rm SO})\}.
\]
The constants in the latter formula are those of
Lemma~\YMMref{14.9}{68}. If $\beta=L^{2q}\varpi_0\le c_2^{\rm SO}$,
then every average in this cone is defined and
\[
 \varpi_j\le\theta_LL^{2(j-q)}\beta\quad(0\le j\le q),
 \qquad 1<\theta_L<1.003.
\]
\end{lemma}

\begin{proof}
The cone at level $j+1$ is a union of blocks contained in the cone
at level $j$. Consequently the one-step plaquette estimate in
Lemma~\YMMref{14.9}{68} applies there whenever
$\varpi_j\le\varpi'_L$, and gives
\[
 \varpi_{j+1}\le L^2\varpi_j
                          (1+C_0^{\rm SO}L^{-2}\varpi_j).
\]
Put $\theta_0=1$ and
$\theta_j=\prod_{i<j}(1+\frac12L^{2(i-q)})$. The geometric series
and $1+x\le e^x$ give $\theta_j\le\theta_L<2$.
If $\varpi_j\le\theta_jL^{2(j-q)}\beta$ and $j<q$, then
$\varpi_j\le2L^{2(j-q)}c_2^{\rm SO}\le\varpi'_L$, so the
small-curvature branch of the real average applies. Moreover
$C_0^{\rm SO}L^{-2}\varpi_j\le\frac12L^{2(j-q)}$ by the definition
of $c_2^{\rm SO}$. The recurrence proves the next induction step.
The case $j=0$ is the definition of $\beta$. Finally $L\ge13$ gives
$\theta_L\le e^{1/336}\le336/335<1.003$.
\end{proof}

\begin{lemma}[Averaging along a straight closed line]
\label{lem:averaged-line-holonomy}
Let $U$ be a real $\mathrm{SU}(2)$-valued configuration and $q\ge1$.
Suppose every bare plaquette has curvature at most $\omega_0$ and
$L^{2q}\omega_0\le c_2^{\rm SO}$, with the constant of
Lemma~\ref{lem:background-iterated-average}. Let $\Gamma$ be a
straight closed line of $N$ bare bonds through a level-$q$ site,
where $L^q$ divides $N$. Let $\Gamma_q$ be the same line in the
level-$q$ lattice. Then
\[
 \bigl|\mathcal M^qU(\Gamma_q)-U(\Gamma)\bigr|
       \le13.3\,NL^q\omega_0.
\]
The assertion also holds on a rectangular periodic lattice under
these same local hypotheses.
\end{lemma}

\begin{proof}
We first keep the curvature parameter in the proof of
Lemma~\YMMref{4.4}{19}. Write $t=L^2\omega\le1/200$ at one
step. The contour diameter is at most
$d=e^{9t/8}-1$, and the rearrangement bound for the middle
path is $e^{25t/8}-1$ (Lemma~\YMMref{4.1}{17}).
Proposition~\YMMref{4.3}{18} writes the contour mean relative
to any chosen contour as $e^Y$, with $Y$ skew-Hermitian and
$|Y|\le2d$. Since $|e^Y-1|\le|Y|$, the two contour replacements
and the middle-path replacement, telescoped between unitary
matrices, give
\[
 |W_x-1|\le4(e^{9t/8}-1)+(e^{25t/8}-1)<8t.
\]
Indeed $e^u-1\le u/(1-u)$ and $t\le1/200$ give the coefficient
$7200/1591+200/63<8$. Hence
$|\log W_x|\le8t/(1-8t)<12t$. The average in
\YMMref{(37)}{20} is the exponential of the mean of these
skew-Hermitian logarithms times the straight transport. It differs
from that transport by at most $12L^2\omega$.

At level $j$, Lemma~\ref{lem:background-iterated-average} gives
$\omega_j\le\theta_LL^{2j}\omega_0$. The small-curvature
admission in that lemma also gives
$L^2\omega_j\le L^2\varpi'_L<1/200$ for $j<q$.
The line has $N/L^{j+1}$ bonds after the next average. A product
of unitary matrices changes by at most the sum of the changes
of its factors, so summing the one-step errors gives
\[
 \sum_{j=0}^{q-1}12\frac{N}{L^{j+1}}L^2\omega_j
 \le12\theta_LN\omega_0\sum_{j=0}^{q-1}L^{j+1}
 \le\frac{12L\theta_L}{L-1}\,NL^q\omega_0
 <13.3\,NL^q\omega_0.
\]
Here $L\ge13$ and $\theta_L<1.003$. All paths and exchanges
are local to their double blocks. The proof is unchanged when
the periods in the three directions differ.
\end{proof}

\begin{lemma}[Derivatives and remainder of the plain average]
\label{lem:plain-average-remainder}
Take $a_Q=a_{\rm plain}$, $c_Q=c_{\rm plain}$ and $C_Q$
from Proposition~\YMMref{51.1}{390}, with the smallness choices
made in its proof, and put $C'_Q=4C_Q/c_Q$.
On its tower over a double block, let $U$ be real with
$\varpi(U)\le a_QL^{-2j}$. For a complex bond form $a$, write
\[
 Q_j(U,a)(c)=\frac1i\log\bigl[
       \mathcal M^j(e^{ia}U)(c)\mathcal M^j(U)(c)^{-1}\bigr],
 \qquad \tau_c=L^j\sup_{\mathsf D_c^j}|a|<c_Q.
\]
The form is defined on the positive bare bonds inside the endpoint
double-block tower; negative bonds are extended covariantly.
Then
\begin{equation}\label{eq:plain-average-remainder}
 |\dot Q_j(U)a(c)|\le C_Q\tau_c,
 \qquad
 |Q_j(U,a)(c)-\dot Q_j(U)a(c)|\le C'_Q\tau_c^2.
\end{equation}
At $U=1$, $\dot Q_j(1)=\mathbf Q^{\rm lin}_j$, the actual
iterated linear average of Lemma~\YMMref{4.6}{22}.
On a whole real periodic lattice satisfying the same curvature
bound, the left-trivialized derivative obeys
\begin{equation}\label{eq:plain-average-l2-derivative}
 \|D\mathcal M^j(U)w\|\le\sqrt6C_Q[w]_j,
 \qquad
 [w]_j^2=\sum_y\left(L^j\max_{b_-\in B^j(y)}|w(b)|\right)^2,
\end{equation}
where both sums defining the lattice $\ell^2$ norms use positive
bonds and counting measure.
\end{lemma}

\begin{proof}
Fix the row $c$. Proposition~\YMMref{51.1}{390} gives a
holomorphic function $F(a)=Q_j(U,a)(c)$ on the ball $\tau_c<c_Q$,
with $F(0)=0$ and $|F(a)|\le C_Q\tau_c$. On every complex
radial line the Cauchy bound gives $|F'(0)a|\le C_Q\tau_c$.
If $\tau_c\le c_Q/2$, its Taylor remainder is bounded by
\[
 \sum_{n\ge2}C_Qc_Q(\tau_c/c_Q)^n
      =\frac{C_Q\tau_c^2}{c_Q-\tau_c}
      \le\frac{2C_Q}{c_Q}\tau_c^2.
\]
The coefficients can first be estimated on circles of radius less
than $c_Q$ and then passed to their supremum. If
$c_Q/2<\tau_c<c_Q$, the bounds for $F$ and $F'(0)$ instead give
$|F(a)-F'(0)a|\le2C_Q\tau_c\le4C_Q\tau_c^2/c_Q$.
The derivative of the logarithm at $1$ is the identity, so the
linearization of the iterated average in
Lemma~\YMMref{4.6}{22} identifies the derivative when $U=1$.

For the final assertion put
$m_y=L^j\max_{b_-\in B^j(y)}|w(b)|$.
Every bond in the determining double block of $c$ has its initial
site in one of its two endpoint blocks. Differentiating the same
row estimate at $a=0$ therefore gives
\[
 |D\mathcal M^j(U)w(c)|^2
       \le C_Q^2\max\{m_{c_-},m_{c_+}\}^2
       \le C_Q^2(m_{c_-}^2+m_{c_+}^2).
\]
Each site is an endpoint of six positive bonds. Summing proves
\eqref{eq:plain-average-l2-derivative}. For $j=0$ the derivative
is the identity and the same bound follows by the same count.
\end{proof}

\subsection{The normalized fluctuation measure}
\label{sec:normalized-step-measure}

Consider a real canonical step of height $t$. Retain its actual
array-size bounds, the full-core radii and numerical ceilings of
Theorem~\YMMref{W}{423}, and the whole-chart hypotheses of
Proposition~\YMMref{52.3}{433}. For the tame class
$\mathsf O_{t+1}=\{\varpi(W)<2\varepsilon_{t+1}\}$ these hypotheses
follow under the coupling ceiling used in the chain; explicitly it
is chosen so that
\[
 2\varepsilon_{t+1}\le\bar\varepsilon,
 \quad2\varepsilon_{t+1}\le\varepsilon_{\mathcal W}\alpha_0,
 \quad\delta_t\le\varepsilon_{\mathcal W},
 \quad4\varepsilon_0<\varepsilon_{\rm v}.
\]
All geometric room and response requirements of that proposition
are retained. At a tame $W$ choose $U_+$ in its regular orbit so that
$\mathcal M^tU_+$ is in the axial gauge, and put
\[
 P_W=(U_+,0,J^{(t+1)}(U_+)),\qquad
 \mathsf{Rg}_{x_t}=\{B\text{ real on nonpivot bonds}:|B(b)|<\mathfrak b_t\}.
\]
Theorem~\YMMref{E}{10} gives regularity at radius
$2B_3\varepsilon_{t+1}$, including the current bound in level-$(t+1)$
units. Proposition~\YMMref{52.3}{433} places every restriction of
$P_W$ in its full core and supplies the whole chart $V_W(B)$ and its
lower backgrounds. The complete Gaussian structure has real positive
covariance $C_W=(\mathsf B_W+\widehat{\mathsf G}_W)^{-1}$, as in
Proposition~\YMMref{52.2}{419} and Theorem~\YMMref{W}{423}.

Let $S_t$ be the action of the canonical bare array, $S_0=0$, and set
\[
 F_t(V)=\exp[-x_tA^{\eta_t}(U_t(V))+S_t(U_t(V))]
\]
on its defined set. The chart only reads this set. Define
\begin{align}
 \mathsf E_W(B)&=\sum_{Y\in\mathcal D_t}
    (\mathbb V^{\rm J}+\mathbb V^{\rm G}+\mathbb V^{\rm A}
                   +\mathbb V'[S_t])(Y;P_W|_Y,B|_Y),
                   \label{eq:normalized-step-potential}\\
 \Xi_W&=\int_{\mathsf{Rg}_{x_t}}e^{\mathsf E_W(B)}\,d\mu_{C_W}(B),
 \qquad d\nu_W=\Xi_W^{-1}\mathbf1_{\mathsf{Rg}_{x_t}}
                           e^{\mathsf E_W}\,d\mu_{C_W}.
                   \label{eq:normalized-step-measure}
\end{align}

\begin{proposition}\label{prop:normalized-step-kernel}
The integral $\Xi_W$ converges absolutely and is nonzero. Thus
$\nu_W$ is a complex Borel measure of finite variation and mass one.
Its pushforward $\kappa_{t+1}(W,\cdot)=(V_W)_*\nu_W$ is supported on
the fibre over $W$ and, for bounded Borel $G$, satisfies
\begin{equation}\label{eq:normalized-step-quotient}
 \int G\,d\kappa_{t+1}(W,\cdot)
       =\frac{\mathsf I_t[GF_t](W)}{\mathsf I_t[F_t](W)}.
\end{equation}
The two potentials $\mathbb V^{\rm int}=\mathbb V^{\rm J}
+\mathbb V^{\rm G}+\mathbb V^{\rm A}$ and $\mathbb V'[S_t]$, read
on real fields, have quadratic-growth norms bounded respectively by
\[
 \tfrac72K_P\delta_t\le\tfrac12w_\star,
 \qquad K_\Lambda s_t^{\rm ab}N_\star\le\tfrac14w_\star,
 \qquad s_t^{\rm ab}=(1+\beta_{\rm s})\delta_t/\alpha_0.
\]
They are therefore each at most a constant times $\delta_t$.
On the actual all-small-field history measures and with their step
factor $\mathsf h_{t+1}$, the quotient has the disintegration identity
\[
 \int H(\mathcal MV)G(V)\mathsf h_{t+1}(V)\,dm_t(V)
 =\int H(W)\left[\int G\,d\kappa_{t+1}(W,\cdot)\right]dm_{t+1}(W).
\]
All statements use the prescribed tests and gauge factors of the
canonical history.
\end{proposition}

\begin{proof}
The actual construction in Theorem~\YMMref{W}{423} supplies the
Gaussian structure and all four local potentials on the full core,
with Definition~\YMMref{52.1}{415}'s bounds. Restricting the field
polydisc to the real cutoff box gives the quadratic-growth norm.
The triangle inequality and the calculation in the proof of
Theorem~\YMMref{V}{416} give the displayed two bounds (its
interaction bound is in fact $7w_\star/16$). The fourth norm uses
the admitted bare-array bound $N(\mathfrak f_t)\le N_\star$.
The exact integral identity~\YMMref{(902)}{417} at $x=\mathsf t=x_t$,
hence $\tau=1$, is precisely $\Xi_W$: the additional quadratic
potential is zero there. It proves absolute convergence and
nonvanishing. The sum defining $\mathsf E_W$ is finite and continuous
on the box. Division by $\Xi_W$ gives finite variation and mass one.

To identify the quotient, the action and gauge remainders in the
chart satisfy
\[
 e^{-x_tG_W(B)+\mathsf J_W(B)}F_t(V_W(B))
 =e^{v_W}e^{-(B,C_W^{-1}B)/2}e^{\mathsf E_W(B)},\qquad
 v_W=-x_tA^{\eta_t}(U_+)+S_t(U_+).
\]
This is the sum of the definitions of the Jacobian, gauge and action
remainders; their quadratic parts are $\widehat{\mathsf G}_W$ and
$\mathsf B_W$. The physical chart identity of
Lemma~\ref{lem:sf-chart-integral} therefore gives
\[
 \mathsf I_t[GF_t](W)=k_W\int_{\mathsf{Rg}_{x_t}}
                G(V_W(B))e^{\mathsf E_W(B)}\,d\mu_{C_W}(B)
\]
with $k_W\ne0$ independent of $G$. Divide by its value at $G=1$.
The chart is continuous and $\mathcal MV_W(B)=W$, proving the Borel
and fibre assertions. The averaged axial field and the complete
operators are independent of the representative $U_+$; the action
and the array are gauge invariant. Thus the measure is independent
of that choice as well.

For disintegration, use the actual density and passing-test identity
proved in Proposition~\ref{prop:levelwise-source-density} at zero
source. On the passing input set, $dm_t=e^{E_{(t)}}F_t\,dV$;
in the step factor the compatibility tests equal one on the whole
small-field chart. The physical chart formula consequently gives
\[
 \int H(\mathcal MV)G(V)\mathsf h_{t+1}(V)\,dm_t(V)
 =e^{E_{(t)}}\int\chi_{t+1}(W)H(W)
                                  \mathsf I_t[GF_t](W)\,dW.
\]
At $G=1$ this is the density formula for $m_{t+1}$. Insert
\eqref{eq:normalized-step-quotient} to obtain the assertion. Null
sets in the density identities do not affect these integrals.
Bounded Borel insertions are allowed by the absolute chart formula;
their measurability and integrability follow from that same formula.
At $t=0$ the bare density is $F_0=\rho_0$ on the passing set, using
$4\varepsilon_0<\varepsilon_{\rm v}$, so the argument is unchanged.
\end{proof}

\subsection{Recovering the datum from the response}
\label{sec:constant-frame-row}

The response estimates give an upper bound for the fine field in
terms of its datum. For the comparison of metrics we also need the
reverse bound. Only the constant frame is needed: the bare spacing
is $\eta=L^{-j}$, the terminal spacing is one, all regions through
level $j$ are the whole torus and the next region is empty. Thus
every constraint row has scale one and there are no interface rows.
All norms in the following statement use counting measure on positive
bonds. Retain the geometric and response hypotheses of
Section~\ref{sec:background-step-complements}.

\begin{lemma}\label{lem:constant-frame-row}
Let $U$ be a regular real background in this constant frame, with
$|\partial U(p)-1|\le\alpha_0\eta^2$ and
$\alpha_0\le L^{-2}a_{\mathsf m}$, where $a_{\mathsf m}$ is the
radius in Proposition~\YMMref{30.2}{238}.
Let $\mathcal Q(A)=a_j(A)$ be its centered straightened row
logarithm, and set $K_0=D\mathcal Q(0)$. Then
\begin{equation}\label{eq:constant-frame-row-kernel}
 \sup_c\sum_b|K_0(c,b)|\le2,\qquad
 \sup_b\sum_c|K_0(c,b)|\le2L^{-3j},\qquad
 \|K_0\|_{2\to2}\le2L^{-3j/2}.
\end{equation}
For the normalized response of equation~\YMMref{(837)}{398},
$\mathcal Q(\mathcal H_U(\zeta))=\zeta$ near zero. Consequently,
with $\mathsf Z=L^{-j}D\mathcal H_U(0)$,
\begin{equation}\label{eq:constant-frame-response-lower}
 \|\xi\|\le2L^{-j/2}\|\mathsf Z\xi\|
                \le C_{\mathcal Q}L^{-j/2}\|\mathsf Z\xi\|,
 \qquad C_{\mathcal Q}:=34(L^3+2).
\end{equation}
At $j=0$ the same conclusions hold with $K_0=\mathsf Z=1$.
\end{lemma}

\begin{proof}
The row estimate is the zero-field estimate for $Q_j^{\mathsf m}$
in Proposition~\YMMref{30.2}{238}, with $\xi_j=1$ and
$\vartheta_c=1$. Its level-$j$ column estimate is
equation~\YMMref{(468)}{238}, again with $\xi_j=1$ and
$\eta=L^{-j}$. These are bounds on sums of the norms of the
actual derivative kernels. For clarity, the counting-measure Schur inequality here is
\[
 \sum_c\left|\sum_bK_0(c,b)A(b)\right|^2
 \le2\sum_{c,b}|K_0(c,b)|\,|A(b)|^2
 \le4L^{-3j}\sum_b|A(b)|^2.
\]
It gives the asserted norm bound.

The straightening recursion gives
\[
 \mathcal M^j(e^{i\eta A}U)
                  =(e^{i\mathcal Q(A)}V)^{v_j[A]},
 \qquad V=\mathcal M^jU.
\]
This is the accumulated one-step identity used in
Proposition~\YMMref{30.2}{238} and in the gauge-removed row
formula before Lemma~\YMMref{30.6}{246}. For the response,
equation~\YMMref{(837)}{398} says that its gauge factor $u_\zeta$
satisfies $u_\zeta(y)v_j[\mathcal H_U(\zeta)](y)=1$ at every
terminal site. Its actual average is $e^{i\zeta}V$.
Gauge covariance and the displayed cancellation therefore give
$e^{i\mathcal Q(\mathcal H_U(\zeta))}V=e^{i\zeta}V$.
The two logarithms lie on the branch through zero, proving their
equality. Differentiating at zero yields
$K_0D\mathcal H_U(0)\xi=\xi$, or $L^jK_0\mathsf Z\xi=\xi$.
Equation~\eqref{eq:constant-frame-response-lower} follows from
\eqref{eq:constant-frame-row-kernel}; the retained constant
$34(L^3+2)$ is larger than two. At $j=0$ no averaging is performed,
so both maps are the identity.
\end{proof}

\begin{lemma}[Continuity on the full core]\label{lem:full-core-continuity}
Let $j\ge1$ and let $X$ be an admitted finite domain. On the
hereditary core of Definition~\YMMref{49.1}{374}, take
$0<b_1\le\pi L^j$. Suppose $F$ depends only on the represented
pair and is holomorphic on every slice with fixed real factor.
Then $F$ is continuous on the full core, with its topology of
matrix entries and bond-form coordinates.
\end{lemma}

\begin{proof}
Put $\xi=L^{-j}$ and let
$P_\nu=(U_\nu,A_\nu,J_\nu)\to P_0$ in the core. Write
$\mathbf U_\nu=e^{i\xi A_\nu}U_\nu$. We move the real factor
to $U_0$, keeping this represented pair fixed.
For each positive bond, $|i\xi A_0(b)|<\pi$. The differential
of the exponential at this matrix is invertible: after a left
translation it is the entire function
$(1-e^{-\mathrm{ad}Z})/(\mathrm{ad}Z)$, whose only possible
zeros have nonzero eigenvalues in $2\pi i\mathbb Z$; here
$\|\mathrm{ad}Z\|<2\pi$. A local inverse $\Lambda_b$ therefore
exists at $e^{i\xi A_0(b)}$. For all large $\nu$ set
\[
 A'_\nu(b)=(i\xi)^{-1}
       \Lambda_b(\mathbf U_\nu(b)U_0(b)^{-1}),\qquad
 \widetilde P_\nu=(U_0,A'_\nu,J_\nu).
\]
There are finitely many bonds, so these choices hold together,
$A'_\nu\to A_0$, and $\widetilde P_\nu$ represents the pair
$(\mathbf U_\nu,J_\nu)$ of $P_\nu$.

We check that $\widetilde P_\nu$ remains in the core. Its real
conditions hold at $U_0$, and its pair-only conditions hold because
they hold for $P_\nu$. The finitely many strict bounds on $A'_\nu$
and its covariant differences follow by continuity at fixed $U_0$.
For each tower, its complex averaged datum is a continuous function
of $\mathbf U_\nu$ on the admitted analytic branch. Its real datum
is now the fixed datum of $U_0$. Their logarithmic quotient thus
converges to the quotient for $P_0$. At $P_0$ that quotient is
strictly inside the response ball; the same ball contains the
quotients for all large $\nu$. Holomorphy of the tower response
gives convergence of its configuration, curvature and current.
The strict bounds defining that tower consequently persist.
There are only finitely many subdomains and tower tests on $X$,
so all these assertions hold simultaneously. This proves
$\widetilde P_\nu$ lies in the core for all large $\nu$.

Both $\widetilde P_\nu$ and $P_0$ are now on the slice with real
factor $U_0$. Slice holomorphy and pair dependence give
\[
 F(P_\nu)=F(\widetilde P_\nu)\longrightarrow F(P_0).
\]
Since the core is a subspace of a finite-dimensional metric space,
sequential continuity proves the assertion.
\end{proof}


\subsection{The half-H\"older estimate for the ordinary response}
\label{sec:ordinary-holder-response}

The ordinary response to a block datum has a half-H\"older first
covariant difference. This assertion concerns ordinary forcing in
the constrained system. It does not give the same-exponent estimate
for a variable anisotropy in Lemma~\ref{C-360-22}.
We retain the natural-pair geometry, actual averaging rows, local
operators and decoupling parameters of Theorem~\YMMref{M}{315}.
All constants below depend only on $L$. The common regularity radius
and block-number floor must meet the local inverse conditions used
below.

For $0<\beta<1$ use the transported H\"older seminorm of
Proposition~\YMMref{34.1}{263}, with exponent $\beta$:
\[
 [f]_{\mathbf U,\beta;y}
 =\sup_{i\ne i'}
  \left(\frac{\xi_y}{|i-i'|_1}\right)^\beta
  |\operatorname{Ad}(\mathbf U(\Gamma_{x,x'}))f(i')-f(i)|.
\]
The indices have the same kind, their base points lie in the
prescribed enlargement $\Delta^{3/2}(y)$, and the coordinate comb
$\Gamma_{x,x'}$ stays in the local domain. At a boundary only
admissible indices and combs are used. Write
$|f|_{s;y}^{\beta}=\xi_y^{-s}
 (\sup|f|+[f]_{\mathbf U,\beta;y})$ and
$|\Psi|_y^{+\beta}=\max\{|\Psi|_y,
 |\nabla_{\mathbf U}A|_{-2;y}^{\beta}\}$ for the enhanced solution
norm. The forcing norms and their divergence injection are those of
\eqref{eq:weighted-source-L2-norm}, with suprema instead of
mean squares; only the first-slot divergence tensor receives the
additional H\"older norm.

\begin{lemma}\label{lem:ordinary-holder-kernels}
There are a radius $a_H>0$, a block floor $M_H$, a decay rate
$\delta_H>0$ and constants $\mathsf B_\sharp(0),\mathsf B_\sharp(1/2)$ such
that the following holds. Take the local response setting above,
with the four local-system radii at most $a_H$ and block number
$\mathfrak M\ge M_H$. Choose the common decay rate
$\mathring\delta_0\le\delta_H$ in the decoupled construction.
Then its normalized operators $H_s,K_s,\mathfrak G_s$ satisfy
\begin{equation}\label{eq:ordinary-holder-kernel-tests}
 \begin{split}
 \mathfrak n_{y,\zeta}(H_sX),\quad
 \mathfrak n_{y,\zeta}(K_sX)
 &\le6L^3\mathsf B_\sharp(1/2)e^{8\varkappa}
       \sum_z e^{-\mu d(y,z)}|X|_z,\\
 \mathfrak n_{y,\zeta}(\mathfrak G_sf)
 &\le6L^3\mathsf B_\sharp(1/2)e^{8\varkappa}
       \sum_z e^{-\mu d(y,z)}|f|_z^\sharp,
 \qquad \mu=3\mathring\delta_0/16,
 \end{split}
\end{equation}
where $|f|_z^\sharp=\xi_z^3\sup_{\mathsf y(b)=z}|f(b)|$ and
\[
 \mathfrak n_{y,\zeta}(A)
  =\frac{\xi_y^{5/2}[\zeta\nabla_{\mathbf U}A]_{1/2;\xi_y}}
          {\xi_y^{1/2}[\zeta]_{1/2;\xi_y}+|\zeta|}.
\]
Here $\zeta\ne0$ is a smooth scalar cutoff supported in the local
comparison cube, whose intersection with the ambient region lies
in $W$. The subscript $1/2;\xi_y$ restricts pairs to physical
Manhattan distance at most $\xi_y$, and uses the same transports
and domain restriction. In this display $[\cdot]_{1/2;\xi_y}$
is the physical quotient, with denominator $|x-x'|_1^{1/2}$;
thus $\xi_y^{1/2}[\zeta]_{1/2;\xi_y}$ is dimensionless. The constants are uniform for all admitted
local domains and $|s(\Delta)|\le e^{\varkappa}$ and for arbitrary subsums
of the localized expansion.
\end{lemma}

\begin{proof}
We give the extension of the two local estimates in the proof of
Proposition~\YMMref{34.1}{263}. Only the fixed exponents $1/2$
and $3/4$ are needed.

At the identity, eliminate the scalar and row components as in
Section~\YMMref{23.4}{182}. The potential is the vector Green
operator applied to
\[
 f'_1+\nabla^{\mathsf T}X^{(p)}
       +\nabla^{\mathsf T}X_1+Q^{\mathsf T}\omega.
\]
The term $X_1$ is absent for ordinary forcing; $Q^{\mathsf T}\omega$
is absent in system $[4]$. To extend the cutoff proof there from $\beta=1/2$ to a fixed
$0<\beta<1$, take $\beta'=(1+\beta)/2$. The scalar elimination
bounds $X^{(p)}$ in $C^{\beta'}$ with weight $-2$, and $f'_1$
and $Q^{\mathsf T}\omega$ with weight $-3$. Apply the general
Green-kernel H\"older bounds
\YMMref{(225)}{140} and~\YMMref{(227)}{140}, with output exponent $\beta$ and,
for $X^{(p)}$, input exponent $\beta'$; their gap is
$(1-\beta)/2>0$. The same cutoff calculation therefore bounds
the first difference of the potential in $C^\beta$ for ordinary
forcing. We use this extension only at $\beta=3/4$; the case
$\beta=1/2$ is equation~\YMMref{349}{182} itself. For the other
estimate needed below, that same equation~\YMMref{349}{182}
already bounds the $C^{1/2}$ first difference for a $C^{3/4}$
divergence source. Its exponent gap is $1/4$ and is not set to zero.

For completeness, the passage to a small regular field does not
require absorption in the enhanced norm. In a local real gauge write
$\mathbb K_0-\mathbb K=\iota_0 V$. The already proved ordinary
inverse controls $|\Psi|_y$, and the local coefficient estimate gives
\[
 |V\Psi|_y\le C\epsilon
       \max_{d(y,z)\le\rho}|\Psi|_z.
\]
For $\Psi=\mathbb E\iota_0g$ the identity
$\Psi=\mathbb E_0\iota_0(g+V\Psi)$ puts only ordinary forcing
into the flat enhanced estimate. For a divergence datum $g^+$
one uses
$\Psi=\mathbb E_0\iota_0^+(g^++(V\Psi,0))$ instead; the ordinary
bound for that solution first uses the positive-exponent divergence
estimate. The extra term has no divergence component, so no
H\"older norm of $V\Psi$ is required. Converting the injection and
the transported norms between the flat and small gauges costs a
fixed local factor. These are the same coefficient, shift and
transport comparisons used to prove equation~\YMMref{519}{263}.
At layer zero, every comparison distance is a bounded number of
bare spacings and the massive boundary inverse gives the estimates
directly. This proves the three required local estimates, with one
radius and floor chosen for the two fixed exponents.

The cutoff commutator has ordinary forcing. Consequently in each
walk the enhanced estimate is needed only at its output factor;
the remaining factors use the ordinary strong inverse. More
explicitly, if $T_\kappa$ is the local inverse with output cutoff
and $R^\sharp_\kappa$ is its commutator forcing, the ordinary walk is
\[
 T_{\kappa_0}R^\sharp_{\kappa_1}\cdots R^\sharp_{\kappa_m}.
\]
The output factor is bounded in $|\cdot|^{+\beta}$, and each of
the other factors costs
$q=C\mathfrak M^{-1}e^{-\delta D_0}$. The enlarged comparison
cube and its two difference shifts lie in the common
$[5/6]$ region, so the rows, transports and seminorms used by the
local and global systems agree. This is the same identification
used in Proposition~\ref{prop:weighted-system-L2}. A finite number
of pieces can follow a given piece. Increasing the common block
floor until their number times $q$ is at most $1/2$ gives the
summed bound and the termwise walk bound.

For $H_s$ and $K_s$, apply this estimate to the last-row forcing.
For $\mathfrak G_s$ use the doubled inverse in the proof of
Proposition~\YMMref{34.1}{263}. Its coupling is
$D(D^{\mathsf T}A'-w')$, a divergence source plus the explicitly
retained layer-zero boundary correction. Apply the ordinary
$3/4$ estimate to the second copy. The transported product and
shift formulas in that proof bound the coupling tensor in
$C^{3/4}$, with weight $-2$. Apply the divergence estimate with
exponents $1/2,1/4$ to the first copy. This bounds its first
difference in $C^{1/2}$. The doubled cutoff commutator in
equation~\YMMref{526}{266} again has ordinary forcing, so the
same walk multiplication and summation apply. The local multiplier
insertions see only the weak norm, already bounded by these output
norms; their tree factors and summability are unchanged. This
identifies the estimate for the actual projected operator, rather
than for a different scalar Green function.

Finally insert the parameters. The carrier count in the decoupled
construction gives
\[
 \#\mathsf P(\omega)\le8+68D_\omega/\mathsf w,
 \qquad
 |s^\omega|e^{-aD_\omega}
  \le e^{8\varkappa}
          e^{-(a-68\varkappa/\mathsf w)D_\omega}.
\]
The standing lower bound on $\mathsf w$ leaves rate
$\mathring\delta_0/4$. Fix $\mathsf B_\sharp(1/2)$ large enough to include
the preceding local constants, the absolute tree masses and the
bounded number of output-window cells. The normalized column
bounds are therefore
\[
 \begin{split}
 [\zeta\nabla H_se_c]_{1/2;\xi_y},\quad
 [\zeta\nabla K_se_c]_{1/2;\xi_y}
 &\le \mathsf B_\sharp(1/2)e^{8\varkappa}
  (\xi_y^{1/2}[\zeta]_{1/2;\xi_y}+|\zeta|)
  \xi_y^{-3/2}\xi_z^{-1}e^{-\mathring\delta_0d(y,z)/4}|e|,\\
 [\zeta\nabla\mathfrak G_sf]_{1/2;\xi_y}
 &\le \mathsf B_\sharp(1/2)e^{8\varkappa}
  (\xi_y^{1/2}[\zeta]_{1/2;\xi_y}+|\zeta|)
  \xi_y^{1/2}e^{-\mathring\delta_0d(y,z)/4}\sup|f|,
 \end{split}
\]
when $c$ is at $z$ and $f$ is supported over $z$. The ordinary
product inequality $[\zeta v]_\beta\le
|\zeta|[v]_\beta+[\zeta]_\beta|v|$ explains the cutoff factor.
After multiplying by $\xi_y^{5/2}$, the scale ratios are
$\xi_y/\xi_z$ and $(\xi_y/\xi_z)^3$ respectively. The layer
comparison absorbs them with a factor $L^3$ and a loss
$\mathring\delta_0/16$ in the rate. There are at most six rows
at each top. Summing the columns or source cells proves
\eqref{eq:ordinary-holder-kernel-tests}. All steps use absolute
bounds on individual terms, so arbitrary subsums obey the same
estimate.
\end{proof}

\begin{proposition}\label{prop:ordinary-holder-response}
Retain the hypotheses of Lemma~\ref{lem:ordinary-holder-kernels}
and the datum radius of Theorem~\YMMref{M}{315}. Choose
$\mathsf B_\sharp(0)$ as the ordinary kernel constant in that theorem,
$\mathsf B_\sharp(1/2)\ge \mathsf B_\sharp(0)$ as above, and put
\[
 \mathsf b_0=6L^3\mathsf B_\sharp(0),\qquad
 \mathsf b_{\mathrm D}=\mathsf b_0c_S(\mathring\delta_0/16),\qquad
 \widehat{\mathsf b}=\mathsf b_{\mathrm D}e^{8\varkappa}.
\]
For its ordinary response $\mathbf H_s$ and any two admissible
block data,
\begin{equation}\label{eq:ordinary-holder-response}
 \sup_y w(y)\mathfrak n_{y,\zeta}
       (\mathbf H_s(B_1)-\mathbf H_s(B_2))
 \le3\widehat{\mathsf b}\frac{\mathsf B_\sharp(1/2)}{\mathsf B_\sharp(0)}
                           |B_1-B_2|_w .
\end{equation}
Here $w$ is any of the admissible exponential weights of
Theorem~\YMMref{M}{315}. In particular, taking $B_2=0$ and $w=1$
gives the same bound for the response itself.
\end{proposition}

\begin{proof}
Write $A'_i=K_sB_i+\mathbf A_0(B_i)$. The response equations
and constraint correction give the exact difference identity
\[
 \mathbf H_s(B_1)-\mathbf H_s(B_2)
 =K_s(B_1-B_2)-\mathfrak G_s(W_s(A'_1)-W_s(A'_2))
                   -H_s(D_s(A'_1)-D_s(A'_2)).
\]
The proof of Theorem~\YMMref{J}{256}, with the radius in
Theorem~\YMMref{M}{315}, gives
\[
 |W_s(A'_1)-W_s(A'_2)|_w^\sharp\le|B_1-B_2|_w,
 \qquad
 |D_s(A'_1)-D_s(A'_2)|_w\le\tfrac12|B_1-B_2|_w.
\]
Indeed $\|A'_i\|\le2\widehat{\mathsf b}\sup|B_i|$ and
$\|A'_1-A'_2\|_w\le2\widehat{\mathsf b}|B_1-B_2|_w$; the first bound
uses $8C_W\widehat{\mathsf b}^4\sup|B_i|\le1$, and the second uses
$32C_2\mathsf e\widehat{\mathsf b}\varepsilon_3\le1$, exactly the
radius inequalities in equation~\YMMref{494}{256}.
Multiplication by $w$ in
\eqref{eq:ordinary-holder-kernel-tests} leaves decay rate
$\mathring\delta_0/16$ for the top sum. Each operator is thus
bounded by $\widehat{\mathsf b}\mathsf B_\sharp(1/2)/\mathsf B_\sharp(0)$ in the indicated
seminorm. The triangle inequality gives the factor
$1+1+1/2\le3$. Finally $\mathbf H_s(0)=0$.
\end{proof}

\subsection{The local perturbation of the Hessian}
\label{sec:weighted-hessian-complement}

A plaquette weight changes the principal part of the Hessian. The
remainder is a local curvature term. This elementary fact identifies
the perturbation that a weighted inverse estimate must control.

\begin{lemma}\label{lem:weighted-local-hessian}
Let $\mathsf A$ be a finite ambient region of a three-dimensional lattice
of spacing $\xi$, and let $\mathbf U$ be an $\mathrm{SL}(2,\mathbb C)$
configuration on its inner bonds. Give one-forms the complex bilinear
pairing $(A,B)_\xi=\xi^3\sum_b\frac12\operatorname{tr}(A(b)B(b))$,
where the sum uses positive inner bonds. For a complex function $w$
on the positive inner plaquettes, define
\[
 F_w(A)=\xi^{-1}\sum_p w(p)
       \bigl(1-\tfrac12\operatorname{tr}\partial(e^{i\xi A}\mathbf U)(p)\bigr).
\]
There is a unique symmetric Hessian $\Delta_w$ with
$F_w''(0)[E,A]=(E,\Delta_w A)_\xi$. Its entries depend linearly on
$w$ and polynomially on the matrices $\mathbf U(b)^{\pm1}$.

Write $\nabla_\mu$ for the covariant forward difference, set
$D^{(1)}A=\nabla_\mu A_\nu-\nabla_\nu A_\mu$ on a $\mu\nu$ plaquette,
and use the $\xi^3$-weighted plaquette pairing to define its transpose.
Then
\begin{equation}\label{eq:weighted-local-hessian}
 \Delta_w=D^{(1)\mathsf T}wD^{(1)}+Z_{1,w}.
\end{equation}
Suppose that $|\partial\mathbf U(p)-\mathbf1|\le1$ and
$|\operatorname{Ad}(\mathbf U(b'))^{\pm1}|\le2$ on every plaquette
containing a given bond $b$ and on its four bonds. Put
\[
 w_b=\max_{p\ni b}|w(p)|,\qquad
 f_b=\xi^{-2}\max_{p\ni b}|\partial\mathbf U(p)-\mathbf1|,
 \qquad |A|_{(b)}=\max_{p\ni b}\max_{b'\in p}|A(b')|.
\]
With absolute constants, the kernel of $\Delta_w$ satisfies
$|\Delta_w(b,b')|\le Cw_b\xi^{-2}$, and
\begin{equation}\label{eq:weighted-local-remainder}
 |(Z_{1,w}A)(b)|\le Cw_b f_b|A|_{(b)}.
\end{equation}
Both kernels vanish unless $b,b'$ belong to a common plaquette.
The same assertions hold after restriction to a subregion, with fields
extended by zero and with the ambient plaquettes retained.
\end{lemma}

\begin{proof}
We give the weighted version of the plaquette calculation in
Lemma~\YMMref{25.3}{194}. At $p=p_{\mu\nu}(x)$ put
$Y_p=\partial\mathbf U(p)$, $y_0=\tfrac12\operatorname{tr}Y_p$,
$\mathsf U_\mu=\operatorname{Ad}(\mathbf U_\mu)$ and
\[
\begin{aligned}
 a_1&=A_\mu(x),& a_2&=\mathsf U_\mu(x)A_\nu(x+\xi e_\mu),\\
 a_3&=\operatorname{Ad}(Y_p)\mathsf U_\nu(x)A_\mu(x+\xi e_\nu),&
 a_4&=\operatorname{Ad}(Y_p)A_\nu(x).
\end{aligned}
\]
For $\epsilon=(1,1,-1,-1)$ and $S_p(A)=\sum_i\epsilon_i a_i$,
transporting the four factors to $x$ gives
\[
 \partial(e^{it\xi A}\mathbf U)(p)
   =e^{it\xi a_1}e^{it\xi a_2}e^{-it\xi a_3}e^{-it\xi a_4}Y_p.
\]
The quadratic term of an ordered product is
$\frac12(\sum_iX_i)^2+\frac12\sum_{i<j}[X_i,X_j]$.
Since a traceless two-by-two matrix $S$ has
$S^2=\frac12\operatorname{tr}(S^2)\mathbf1$, and
$Y_p=y_0\mathbf1+i\xi^2F_{\mathbf U}(p)$, differentiation yields
\begin{equation}\label{eq:weighted-plaquette-polarization}
 (A,\Delta_w A)_\xi
 =\xi^3\sum_p w(p)\left\{
 y_0\langle\xi^{-1}S_p(A),\xi^{-1}S_p(A)\rangle
 +i\sum_{i<j}\epsilon_i\epsilon_j
       \langle[a_i,a_j],F_{\mathbf U}(p)\rangle\right\}.
\end{equation}
The pairing is nondegenerate; polarizing this finite quadratic form
therefore defines exactly one symmetric operator. The formula also
proves its asserted dependence on $w$ and $\mathbf U$.

The exact identity
\[
 \xi^{-1}S_p(A)=D^{(1)}A(p)+\rho_p(A),\qquad
 \rho_p(A)=\xi^{-1}(I-\operatorname{Ad}(Y_p))
    \{\mathsf U_\nu(x)A_\mu(x+\xi e_\nu)+A_\nu(x)\}
\]
shows that the quadratic form of $Z_{1,w}$ is the sum, multiplied by
$\xi^3w(p)$, of
\[
\begin{split}
 &(y_0-1)\langle D^{(1)}A,D^{(1)}A\rangle
 +2y_0\langle D^{(1)}A,\rho_p(A)\rangle
 +y_0\langle\rho_p(A),\rho_p(A)\rangle\\
 &\hspace{20mm}
 +i\sum_{i<j}\epsilon_i\epsilon_j
       \langle[a_i,a_j],F_{\mathbf U}(p)\rangle.
\end{split}
\]
Here all values in one summand are at $p$. If
$\kappa_p=|Y_p-\mathbf1|\le1$, the two-by-two determinant identity gives
\[
 |y_0-1|\le\kappa_p,\quad
 \xi^2|F_{\mathbf U}(p)|\le\kappa_p,\quad
 |I-\operatorname{Ad}(Y_p^{\pm1})|\le3\kappa_p.
\]
Consequently $|D^{(1)}A(p)|\le C\xi^{-1}|A|_{(p)}$ and
$|\rho_p(A)|\le C\xi^{-1}\kappa_p|A|_{(p)}$.
Polarize against a one-form supported at $b$, of norm at most one.
Each of the at most four contributing plaquettes then costs
$C|w(p)|\xi^{-2}\kappa_p|A|_{(p)}$; the factor $\xi^3$ cancels
against that of the one-bond pairing. The nondegenerate finite
dimensional pairing gives \eqref{eq:weighted-local-remainder}.
The same polarization of \eqref{eq:weighted-plaquette-polarization},
with the second field supported on one bond, gives the kernel bound.
Every summand involves only the four bonds of its plaquette. This proves
the locality assertions, including restriction by zero extension.
\end{proof}

In particular the perturbation at $w=1+v$ is
\[
 \Delta_{1+v}-\Delta_1=D^{(1)\mathsf T}vD^{(1)}+Z_{1,v}.
\]
The principal insertion contains the divergence source $vD^{(1)}A$.
Multiplication by $v$ costs at most $|v|_\infty$ in $L^2$.
This calculation identifies the insertion; it does not assert the
uniform inverse or the endpoint H\"older estimate for that insertion.

\subsection{Mean-square bounds for the weighted system}
\label{sec:weighted-mean-square}

The weight in Lemma~\ref{lem:weighted-local-hessian} need not have
bounded differences. We keep its principal contribution in divergence
form. The relevant inverse estimate is therefore a mean-square estimate
with a divergence source. We give the extension of the proof of
Theorem~\YMMref{G}{204} that is needed for this purpose.

We use precisely the separated regions, inner positive bonds, reduced
coarse rows and four systems of Sections~\YMMref{23}{176} and
\YMMref{27}{203}. In particular the row space consists of the positive
level-$n$ bonds with neither endpoint in $\Omega_{n+1}$. The bilinear
pairings have weights $\xi^3$ on fine sites and bonds, $\xi_y^3$ on
block tops, and $\xi_c^3$ on a coarse row $c$. All transposes below use
these pairings. The averaging row is the derivative of the physical
constraint,
\[
 Q A(c)=\left.\frac{d}{ds}\right|_{s=0}
     \xi_c^{-1}\mathcal Q_n(\mathbf U;s\xi_c A)(c),
\]
with the convention of equation~\YMMref{370}{190}; $Q'$ is the
ordered scalar average. At $\mathbf U=\mathbf1$, these are the straight
bond average and the arithmetic scalar average. At nontrivial fields
we retain their actual first and second derivatives. In particular,
\[
 (A,\Delta^{(2)}[\mathsf j]B)_\xi
       =2(C^{(2)}(A,B),\mathsf j)_{\mathrm{rows}}.
\]
Thus the result below does not replace the physical constraint by its
linearization at the identity.

For a fine-indexed field $f$ define the local mean-square seminorm
\begin{equation}\label{eq:weighted-local-L2}
 \|f\|_{s;y}
 =\xi_y^{-s-3/2}
       \left(\xi^3\sum_{\mathsf y(i)=y}|f(i)|^2\right)^{1/2}.
\end{equation}
For a field on tops or rows, use
$\|f\|_{s;y}=\xi_y^{-s}\max_{\mathsf y(i)=y}|f(i)|$.
Values are zero outside the region. The system differences omit translated
indices outside their inner-bond or site index set. The ambient Hessian
retains boundary plaquettes of the zero-extended field; its additional
boundary terms are treated explicitly in the proof below.
For $\Psi=(A;v,w,\sigma;p,q,r;\omega)$ put
\begin{equation}\label{eq:weighted-system-L2-norm}
\begin{split}
 \|\Psi\|_y=\max\{&\|A\|_{-1;y},\|\nabla A\|_{-2;y},
 \|v\|_{0;y},\|Dv\|_{-1;y},\|w\|_{-2;y},\|Dw\|_{-3;y},
 \|\sigma\|_{-4;y},\\
 &\|p\|_{-2;y},\|q\|_{0;y},\|Dq\|_{-1;y},
 \|r\|_{-4;y},\|\omega\|_{-3;y}\}.
\end{split}
\end{equation}
Absent components are omitted for a compressed system. Write
$g^+=(f_1,\ldots,f_8;X_5,X_1)$ and set
\begin{equation}\label{eq:weighted-source-L2-norm}
\begin{split}
 \|g^+\|_y=\max\{&\|f_1\|_{-3;y},\|f_2\|_{-2;y},
 \|f_3\|_{-4;y},\|f_4\|_{0;y},\|f_5\|_{-4;y},
 \|f_6\|_{-2;y},\|f_7\|_{0;y},\\
 &\|f_8\|_{-1;y},\|X_5\|_{-3;y},\|X_1\|_{-2;y}\}.
\end{split}
\end{equation}
Here $X_1$ has one fine bond index and one difference direction. The
source map $\iota^+$ replaces $f_1$ by
$f_1+\nabla^{\mathsf T}X_1$ and $f_5$ by
$f_5+D^{\mathsf T}X_5$. The original map $\iota$ is the case $X_1=0$.

\begin{proposition}\label{prop:weighted-system-L2}
There are constants $a_2,r_2,B_2,\delta_2>0$ and a power $M_2$ of $L$,
depending only on $L$, with the following properties. Take the full
geometric hypotheses and the natural-pair class of Theorem~\YMMref{G}{204},
with block number $\mathfrak M\ge M_2$ and all four radii
$\alpha_0,\alpha_1,\gamma_0,\gamma'$ at most $a_2$.
For any complex plaquette function $v$ with $|v|_\infty<r_2$, replace
$\Delta_Y$ by $\Delta_{1+v}$ in the first slot of each applicable system;
retain the other coefficients and rows. The resulting four systems
$\mathbb K_v^{[i]}$, $i=1,2,4,5$, are invertible. For $i=4,5$,
\begin{equation}\label{eq:weighted-system-L2-decay}
 \| (\mathbb K_v^{[i]})^{-1}\iota^+g^+\|_y
 \le B_2\sup_z e^{-\delta_2d(y,z)}\|g^+\|_z.
\end{equation}
The two scalar systems are unchanged. Their bounds are the corresponding
compressed bounds. Constants are independent of the volume, the number
of layers and the differences of $v$.

The inverse has a sum over walks of the pieces of
Section~\YMMref{27.1}{204}. If
$\varpi=(\kappa_0,\ldots,\kappa_m)$, its term $E_{v,\varpi}\iota^+$
vanishes unless consecutive enlarged pieces meet. It is supported in
$\kappa_0^{[2/3]}$ at its output and depends only on the coefficients
and weights in the union of the local regions of the visited pieces.
With $y_0=y$ it satisfies
\begin{equation}\label{eq:weighted-L2-walk}
 \|E_{v,\varpi}\iota^+g^+\|_y
 \le B_2 q_2^m\sup
       \left\{e^{-\delta_2\sum_{a=0}^m d(y_a,y_{a+1})}
                          \|g^+\|_{y_{m+1}}\right\},
 \qquad q_2=\frac{C_2}{\mathfrak M}e^{-\delta_2D_0}.
\end{equation}
The intermediate tops belong to
$\kappa_{a-1}^{[5/6]}\cap\kappa_a^{[5/6]}$ and the final top to
$\kappa_m^{[5/6]}$. The number $D_0$ is the separation between the
input cutoff and the commutator support in the same piece construction.
An empty supremum is zero. The rate can be fixed first below the
unweighted local rate; increasing $M_2$ and decreasing $a_2,r_2$ then
preserves that rate for all weights under consideration.
\end{proposition}

\begin{proof}
We retain the proof's order: the flat system, a small field, and then
localization to regular pieces. This also checks that the mean-square
estimate uses the same rows as the nonlinear problem.

\emph{The flat inverse.}
Write
\[
 G'=(\Delta'_a)^{-1},\quad C'=(Q'G'^2Q'^{\mathsf T})^{-1},
 \quad N=G'Q'^{\mathsf T}C',\quad P=NQ'G',\quad R=I-P,
\]
\[
 G_0=(\Delta_Y(\mathbf1)+DRD^{\mathsf T}
                    +Q^{\mathsf T}\mathsf m^{-2}Q)^{-1},
 \qquad Z=QG_0Q^{\mathsf T}.
\]
These are the scalar, vector and row operators in
Section~\YMMref{23.2}{179}. In the notation $T:s\Rightarrow s'$,
the arrow means an estimate
$\|Tf\|_{s';y}\le C\sup_z e^{-\delta d(y,z)}\|f\|_{s;z}$.
At a common sufficiently small rate, their mean-square bounds are
\begin{equation}\label{eq:weighted-flat-L2-arrows}
\begin{gathered}
 O:s\Rightarrow s+2,\qquad
 \nabla O, O\nabla^{\mathsf T}:s\Rightarrow s+1,
 \qquad \nabla O\nabla^{\mathsf T}:s\Rightarrow s,\\
 \nabla\nabla G_0:s\Rightarrow s,\qquad
 C':0\Rightarrow-4,\qquad Z^{-1}:s\Rightarrow s-2,
 \qquad Q,Q^{\mathsf T},Q',Q'^{\mathsf T}:s\Rightarrow s.
\end{gathered}
\end{equation}
Here $O=G',G_0$, and only weights $-4\le s\le2$ are used.
For completeness, these bounds follow from the local $L^2$ estimates
in the proof of Proposition~\YMMref{21.3}{143}, the six bounds of
Proposition~\YMMref{21.1}{140}, and the top and row kernel estimates
in Sections~\YMMref{21}{138} and~\YMMref{22}{171}, as follows.
If $N_y(f)=(\xi^3\sum_{\mathsf y(i)=y}|f(i)|^2)^{1/2}$, the local
bounds have the form
\[
 N_y(T\mathbf1_z f)\le C\xi_y^a e^{-\delta_0d(y,z)}N_z(f),
 \qquad a=2,1,0.
\]
Normalization introduces
$(\xi_z/\xi_y)^{s+3/2}$. The layer-gap estimate of
Lemma~\YMMref{21.2}{141} bounds this factor by
$L^6e^{\delta_0d(y,z)/4}$ after increasing the separation threshold.
Summation over $z$ spends another quarter of the decay; the remaining
half gives the arrow. The same argument applies to the differentiated
scalar energy estimate: pairing
$\Delta'_a u=\nabla^{\mathsf T}X$ with an exponentially weighted $u$
bounds its weighted gradient by $C\|X\|_2$, using the coercive
conjugated form in that proof. Thus the mixed scalar arrow costs no
difference of $X$.

The top and row seminorms are suprema, so their arrows are their
counting-kernel bounds. For an average, Cauchy--Schwarz gives
\[
 |Qf(c)|^2\le\Bigl(\sum_b|Q(c,b)|\Bigr)
                       \sum_b|Q(c,b)|\,|f(b)|^2.
\]
The straight average has row sum one and
$|Q(c,b)|\le(\xi/\xi_c)^3$; $Q'$ is the arithmetic mean.
A double block meets only $C(L)$ cells at comparable scales. This
proves their mean-square arrows. The weighted transposes have bounded
row sums, and top or row input is already measured by its supremum;
there are at most nine fine components per site. This proves the two
transpose arrows with no mesh-dependent factor. These observations
also justify compositions of the arrows: the supremum estimate
composes by the triangle inequality for $d$.

The elimination formulas are unchanged by adding $X_1$. Set
\[
 p'=RG'f_5+Pf_6+Nf_7,\qquad p''=RG'D^{\mathsf T}X_5,
 \qquad \varphi=Pf_2+RG'f_3-Nf_4,
\]
\[
 f_0=f_1+D\varphi+Dp'+Dp'',\qquad
 f_*=f_0+\nabla^{\mathsf T}X_1.
\]
Direct substitution in the eight equations gives
\begin{equation}\label{eq:weighted-flat-elimination}
\begin{aligned}
 p&=p'+p'',&q&=G'(f_6-p),\\
 r&=C'\{Q'G'^2(f_5+D^{\mathsf T}X_5)-Q'G'f_6-f_7\},\\
 A&=G_0f_* &&(i=4),\\
 \omega&=Z^{-1}(f_8-QG_0f_*),& A&=G_0(f_*+Q^{\mathsf T}\omega)
                      \quad(i=5),\\
 w&=PD^{\mathsf T}A+\varphi,&v&=G'(f_2+D^{\mathsf T}A-w),\\
 \sigma&=C'\{Q'G'(f_2+D^{\mathsf T}A)-Q'G'^2f_3-f_4\}.
\end{aligned}
\end{equation}
For example, $p',p'',\varphi$ have weight $-2$ and their gradients
have weight $-3$; for $Dp''$ use
$DRG'D^{\mathsf T}=DG'D^{\mathsf T}-(DP)G'D^{\mathsf T}$.
Hence $f_0$ has weight $-3$. The mixed arrow in
\eqref{eq:weighted-flat-L2-arrows} gives $A$ of weight $-1$ and
$\nabla A$ of weight $-2$ even when $X_1$ is only in mean square.
The row expression gives $\omega$ of weight $-3$.

To finish the remaining slots, write
$D^{\mathsf T}A=T_1\nabla A+T_2A$.
The first matrix has at most three entries per row and one per column,
each of size one. The second is supported on the layer of bare cells
meeting the boundary, where $\xi_y=\xi$, and has at most three
entries per row and two per column, each of size $\xi^{-1}$.
Cauchy--Schwarz therefore bounds $D^{\mathsf T}A$ with weight $-2$.
The last line of \eqref{eq:weighted-flat-elimination} now gives exactly
the remaining entries of \eqref{eq:weighted-system-L2-norm}.
This proves the flat estimate. When $X_1=0$ the same calculation
also bounds $\nabla\nabla A$ with weight $-3$.

\emph{Small backgrounds and the actual rows.}
The defect factorization of Section~\YMMref{25.4}{196} is
$\mathbb K_0-\mathbb K=\iota_0 V$. Its coefficients satisfy the
same bound in the mean-square norms,
\begin{equation}\label{eq:weighted-small-defect-L2}
 \|V\Psi\|_y\le C\epsilon
                     \max_{d(y,z)\le\rho_0}\|\Psi\|_z,
 \qquad \rho_0=C(L).
\end{equation}
Here is the kernel check needed to pass from the displayed
pointwise defect to this assertion. The local difference and current
terms have uniformly bounded row and column multiplicities, scale
factors $\xi_y^{-a}$, and coefficients $O(\epsilon)$; their second-order
part is expressed using one difference of $A$. Cauchy--Schwarz gives
\eqref{eq:weighted-small-defect-L2} for these terms. For the averaging
terms, the proof of Lemma~\YMMref{25.2}{192} gives
\[
 \sum_b|\delta Q(c,b)|\le C\epsilon,
 \qquad |\delta Q(c,b)|\le
                     C\epsilon(\xi/\xi_c)^3,
\]
with double-block support; the same estimates hold for $Q'$.
The preceding Cauchy--Schwarz calculation proves their mean-square
bounds. Their transposes use the specified physical pairings and the
bounded number of incident rows. Products such as
$\delta(Q^{\mathsf T}\mathsf m^{-2}Q)$ are expanded with one changed
factor, so satisfy the required weighted bound.

Finally $\Delta^{(2)}[\mathsf j]$ has a symmetric bilinear kernel
supported in the same double blocks. The bound for $C^{(2)}$ in that
lemma, polarized against a field supported at one bond, bounds its
absolute row sum by $C\epsilon\xi_y^{-2}$. In Pauli coordinates the
bilinear transpose has norm at most twice the original matrix norm;
therefore the column sum obeys the same bound with a fixed extra
factor. The weighted Schur inequality proves its mean-square bound.
These are all the types of terms in $V$. In particular this argument
uses the Hessian of the actual nonlinear row, not merely a coincident
row label.

Choose $\epsilon$ so that the norm of
$g^+\mapsto(V\mathbb E_0\iota_0^+g^+,0)$ is at most $1/4$ at the
chosen decay rate. Then
\[
 \mathbb E\iota_0^+
 =\mathbb E_0\iota_0^+
       [I-(V\mathbb E_0\iota_0^+,0)]^{-1}
\]
converges in these norms. The difference
$\iota^+-\iota_0^+$ is a local source with coefficients
$O(\epsilon)$, since it contains only the changed transports in the
two divergences. Covariant and ordinary first differences differ by
local zeroth-order terms of the same size. This proves the small-field
estimate in the covariant seminorms. Gauge changes by $\mathrm{SU}(2)$
matrices preserve all the mean-square norms, so the conclusion also
holds for a triple small in a local real gauge.

\emph{Inserting the weight.}
Only the first equation changes. By
Lemma~\ref{lem:weighted-local-hessian},
\[
 (\mathbb K_v-\mathbb K)\Psi
       =e_1\{D^{(1)\mathsf T}vD^{(1)}A+Z_{1,v}A\}
       =\iota^+ V_v\Psi.
\]
To read the last expression in the system's source space, put the
antisymmetric components of $vD^{(1)}A$ whose indices belong to that
space into $X_1$. Put $Z_{1,v}A$ and the omitted boundary contributions
into $f_1$; the other entries are zero. Indeed, the ambient plaquette
form retains plaquettes meeting the boundary, whereas the system's
$\nabla$ omits a translated bond outside its inner-bond index set.
Their difference is supported on bonds a bounded number of bare steps
from the boundary. Every such cell is in the bare layer, where
$\xi_y=\xi$. Each difference has bounded row and column multiplicity
and size at most $C|v|_\infty\xi^{-2}|A|$ on neighboring bonds.
The weighted Schur estimate therefore gives it source weight $-3$
from the potential weight $-1$. This is the same boundary treatment
used above for $D^{\mathsf T}A$. In the interior the identification
with the transpose of
$D^{(1)}A=\nabla_\mu A_\nu-\nabla_\nu A_\mu$ is exact.
The curvature bound of the natural pair and the local plaquette
support consequently give
\begin{equation}\label{eq:weighted-vertex-L2}
 \|V_v\Psi\|_y\le C|v|_\infty
                       \max_{d(y,z)\le\rho_0}\|\Psi\|_z.
\end{equation}
No derivative of $v$ occurs. For a small-field local inverse,
$\|V_v\mathbb E\iota^+\|\le1/4$ follows by decreasing $r_2$.
Consequently
\[
 \mathbb E_v\iota^+
 =\mathbb E\iota^+[I+V_v\mathbb E\iota^+]^{-1}
\]
converges at the same fixed local rate. It is a right inverse on all
sources: $\iota^+$ is surjective, since the plain source components
alone span the equations. Finite dimensionality therefore gives
invertibility. The scalar compressions do not contain the changed
first equation. This proves the local weighted conclusion for all
four systems.

\emph{Localization.}
Use the pieces, $\widetilde h_\kappa$, $\theta_\kappa$ and local
sequences of Lemma~\YMMref{27.5}{207}. Lemma~\YMMref{27.6}{209}
identifies the global and local physical averaging germs, including
$C^{(2)}$, on $\kappa^{[5/6]}$. The pair on a regular piece has a
real gauge in which its size is at most
$8L^2(\alpha_0+\alpha_1+\gamma_0+\gamma')$.
Choose $a_2$ to place this below the small-field threshold above.
A boundary piece has only the bare layer; its mean-square norms are
finite-component Euclidean norms, uniformly equivalent to the local
supremum norms. The boundary-piece inverse in the proof of
Theorem~\YMMref{G}{204}, followed by the same bounded insertion,
gives its local estimate. Thus every local weighted inverse has a
uniform mean-square bound at a rate $2\delta_2$.

The source cutoff must respect divergence form. Define
$h\star g^+$ by multiplying all components by $h$ and subtracting
$[\nabla^{\mathsf T},h]X_1$ from the first component and
$[D^{\mathsf T},h]X_5$ from the fifth. Then
\begin{equation}\label{eq:weighted-source-cutoff}
 \iota^+(h\star g^+)=h\,\iota^+g^+.
\end{equation}
The additional commutators have coefficients bounded by the first
scaled difference of $h$, so are bounded in mean square. For the
system commutator write
\[
 [\mathbb K_v,h]=\iota^+ C_{v,h}^+.
\]
The unweighted terms have the same three kernel types checked in
\eqref{eq:weighted-small-defect-L2}, with $\epsilon$ replaced by the
scaled first or second difference of $h$. For example,
$[Q,h](c,b)=Q(c,b)(h(b)-h(c))$; the quadratic-row term is treated by
its symmetric Schur bounds. The weighted principal term is exactly
\[
 [D^{(1)\mathsf T}vD^{(1)},h]A
  =D^{(1)\mathsf T}\{v[D^{(1)},h]A\}
                +[D^{(1)\mathsf T},h]vD^{(1)}A.
\]
Keep the admitted indices of its first term as a divergence source.
Its omitted boundary indices are placed in the ordinary first slot,
by the same bare-layer calculation as for $V_v$. Commuting their local
kernel with $h$ introduces the difference of $h$ across at most a
fixed number of bare steps. Thus both the interior and boundary terms
cost a first scaled difference of $h$ and $|v|_\infty$; no difference
of $v$ is taken. The local remainder $[Z_{1,v},h]$ has the same bound
by its plaquette support. The piece cutoffs consequently give
\[
 \|C_{v,\widetilde h_\kappa}^+\Psi\|_y
      \le C\mathfrak M^{-1}
                    \max_{d(y,z)\le2\rho_0}\|\Psi\|_z.
\]
The commutator is supported where the cutoff is nonconstant, at
distance at least $D_0$ from the input cutoff. All these norms coincide
under the local/global identification: the values, differences and
one-step shifts that occur here lie in $\kappa^{[5/6]}$. This is the
same margin used in Lemma~\YMMref{27.6}{209}; raising $M_2$ to ensure
$\mathfrak M\ge32L^4$ also includes the two shifts in the differentiated
norms.

Define maps on divergence-form sources by
\[
 T_\kappa^+=\widetilde h_\kappa\mathbb E_{v,\kappa}\iota^+
                                      (\theta_\kappa\star\cdot),
 \qquad
 R_\kappa^+=-C_{v,\widetilde h_\kappa}^+
                   \mathbb E_{v,\kappa}\iota^+
                                      (\theta_\kappa\star\cdot).
\]
The identity $\sum_\kappa\theta_\kappa=1$ gives
\[
 \mathbb K_v\sum_\kappa T_\kappa^+
       =\iota^+\left(I-\sum_\kappa R_\kappa^+\right).
\]
Use half the local exponential decay for the input-to-commutator
distance. The remaining half yields $\|T_\kappa^+\|\le C$ and
$\|R_\kappa^+\|\le q_2$ with the supports in the statement.
The local geometry bounds the number of possible next pieces by a
number $N=N(L)$; at a fixed output top at most $N$ first pieces
contribute. Increase $M_2$ until $Nq_2\le1/2$. Then
\[
 \mathbb E_v\iota^+
 =\sum_{m\ge0}\ \sum_{\kappa_0,\ldots,\kappa_m}
       T_{\kappa_0}^+R_{\kappa_1}^+\cdots R_{\kappa_m}^+.
\]
Composition gives \eqref{eq:weighted-L2-walk}; summing the geometric
series $\sum_m N^{m+1}q_2^m$ gives
\eqref{eq:weighted-system-L2-decay}. The local construction also proves
the stated dependence on the visited regions. This completes the proof.
\end{proof}

A convenient precise meaning of a connected walk term is now available.
Make a graph whose vertices are the pieces and whose edges join pieces
with meeting $[5/6]$ enlargements. For a finite vertex set $S$, sum the
walks whose set of visited pieces is exactly $S$. There is no such walk
when $S$ is disconnected. If $|S|=s$, every contributing walk has at
least $s-1$ transitions. Given $0<\epsilon_*<1$, increase $M_2$ until
$Nq_2\le\min\{1/2,\epsilon_*\}$. At each fixed output top its sum is
then bounded, in \eqref{eq:weighted-system-L2-decay}, by
\[
 C\sum_{m\ge s-1}(Nq_2)^m\le2C\epsilon_*^{s-1}.
\]
Locality follows from the union of the corresponding local regions.
The same grouping applies to a difference of two local inverses:
$E_v-E_0=-E_0(\mathbb K_v-\mathbb K_0)E_v$ and the displayed local
vertex identify its support. This is the connected-set factor used
for the weighted fluctuation form; its cubes are these pieces, with
the indicated enlargements, rather than an arbitrary partition.

The proposition also permits an envelope $h$ with
$h(z)/h(y)\le C_h e^{\delta_2d(y,z)/2}$. Dividing
\eqref{eq:weighted-system-L2-decay} by $h(y)$ gives the corresponding
weighted mean-square estimate, with constant $B_2C_h$ and at least
half the decay rate. This quantitative ratio is what is meant here
when an envelope varies slowly.

\begin{corollary}\label{cor:weighted-energy-L2}
Under the hypotheses of Proposition~\ref{prop:weighted-system-L2}, let
the only forcing be the first-slot divergence $D^{(1)\mathsf T}Y$,
with homogeneous row data. The potential component of the solution
satisfies
\begin{equation}\label{eq:weighted-global-energy-L2}
 \xi^3\sum_b\xi_{y(b)}^{-2}|A(b)|^2
     +\xi^3\sum_{b,\mu}|\nabla_\mu A(b)|^2
       \le C\xi^3\sum_p|Y(p)|^2.
\end{equation}
Ambient boundary plaquettes are included on the right. The constant
depends only on $L$.
\end{corollary}

\begin{proof}
Decompose the divergence source by its cell positions and apply the
local estimate to one cell at a time. The antisymmetric embedding
and its bare-layer boundary correction have uniformly bounded
multiplicities, as in the proof of
\eqref{eq:weighted-vertex-L2}. Thus the block norm of the map from
$Y$ to $(A,\nabla A)$ is at most $Ce^{-\delta_2d(y,z)}$, with weights
$-2$ on $Y$ and $\nabla A$ and $-1$ on $A$.
The squares of these normalized cell norms acquire the common factor
$\xi_y^{-1}$ when expressed in the energy norm of
\eqref{eq:weighted-global-energy-L2}. On the resulting ordinary
sequence $\ell^2$ spaces the block bound is therefore
\[
 Ce^{-\delta_2d(y,z)}(\xi_z/\xi_y)^{1/2}
       \le C'L^{1/2}e^{-\delta_2d(y,z)/2},
\]
after the same increase of the layer separation used in the flat
estimate. The row and column sums of this nonnegative matrix are
bounded by the cell-summation estimate of
Lemma~\YMMref{21.2}{141}. Cauchy--Schwarz, first with its row sum and
then with its column sum, proves its boundedness on $\ell^2$.
This is \eqref{eq:weighted-global-energy-L2}.
\end{proof}

\section{A uniform gauge on a box}
\label{sec:uniform-box-gauge}

A comb gauge gives a bond bound proportional to the side of the
box. To obtain a bound independent of that side, divide the box
into cells of side comparable to the inverse square root of the
plaquette bound. First fix the gauge on their edges. Extend it over
each face and then each cell, keeping its already prescribed boundary
values. For $SU(2)$, the required extensions can be constructed
quantitatively on the unit sphere in the quaternions.

The following constants will also fix the numerical conventions in
Lemma~\ref{C-367-11}:
\begin{equation}\label{eq:uniform-box-gauge-constants}
\begin{gathered}
 \lambda_{{\rm BG},1}=\pi+2,\qquad \rho_{{\rm BG},1}=1/250,\qquad
 \lambda_{{\rm BG},2}=2+12\sqrt3\lambda_{{\rm BG},1}\rho_{{\rm BG},1}^{-2},
 \\
 \lambda'_{{\rm BG},2}=\lambda_{{\rm BG},2}+4,\qquad
 \rho_{{\rm BG},2}=(250000\lambda_{{\rm BG},2}'^2)^{-1},
 \\
 \lambda_{{\rm BG},3}=4+12\sqrt3\lambda'_{{\rm BG},2}\rho_{{\rm BG},2}^{-2},
 \qquad C_{\rm BG}=2\lambda_{{\rm BG},3}.
\end{gathered}
\end{equation}

\begin{proposition}\label{prop:uniform-box-gauge}
Let $U$ be an $SU(2)$ configuration on the bonds of
$\{0,\ldots,n\}^3$, where $n\ge0$. If
$0<\kappa\le C_{\rm BG}^{-2}$ and
$|\partial U(p)-\mathbf1|\le\kappa$ at every plaquette, then some
gauge $u$ satisfies
\begin{equation}\label{eq:uniform-box-gauge}
 |U^u(b)-\mathbf1|\le C_{\rm BG}\sqrt\kappa
 \quad\hbox{at every bond.}
\end{equation}
The same conclusion holds on $(\mathbb Z/n\mathbb Z)^3$ if
$n\ge\kappa^{-1/2}$.
\end{proposition}

\begin{proof}
We use the operator norm, which on quaternions is their Euclidean
norm, and identify $SU(2)$ with the unit sphere $S^3$. Unitary
multiplication preserves this norm. In particular, reversing a bond
does not change its distance from $\mathbf1$.

\emph{Avoiding a point of the sphere.}
If a subset $S\subset S^3$ is covered by $m$ closed balls of radius
$0<\rho<1/8$ and $m<(8\rho)^{-3}$, there is $q\in S^3$ such that
$|z+q|>\rho$ for every $z\in S$.
Indeed, put $k=\lceil(8\rho)^{-1}\rceil-1\ge1$. The points

\[
 p_v=\frac{(1,v)}{|(1,v)|},\qquad
 v\in\{-1+2a/k:0\le a\le k\}^3,
\]
number $(k+1)^3\ge(8\rho)^{-3}>m$. For distinct grid points $v,w$,
the wedge product of $(1,v)$ and $(1,w)$ has norm at least $|v-w|$,
whereas both vectors have norm at most $2$. Consequently

\[
 |p_v-p_w|\ge\sin\angle((1,v),(1,w))
 \ge\tfrac14|v-w|\ge\frac1{2k}>4\rho.
\]
Each of the $m$ ball centres is within distance $2\rho$ of at most
one $p_v$. Choose a remaining $p_v$ and put $q=-p_v$.

\emph{Extending a map which avoids that point.}
Let $Y$ be a nonempty subset of a finite metric space $(X,d)$.
Suppose $h:Y\to S^3$ is $a$-Lipschitz and $|h+q|\ge\rho$.
Then it extends to $H:X\to S^3$ with Lipschitz constant
$4\sqrt3\rho^{-2}a$. To see this, use the inversion

\[
 I_q(z)=-q+2\frac{z+q}{|z+q|^2},\qquad
 |I_q(z)-I_q(w)|=
 \frac{2|z-w|}{|z+q|\,|w+q|}.
\]
It is its own inverse, maps $S^3\setminus\{-q\}$ to $q^\perp$,
and maps $q^\perp$ to $S^3$. Its Lipschitz constants on the two
sets in use are $2\rho^{-2}$ and $2$, respectively. Each coordinate
$f$ of $I_qh$ in an orthonormal basis of $q^\perp$ extends by

\[
 F(x)=\min_{y\in Y}\{f(y)+2\rho^{-2}a\,d(x,y)\}.
\]
The Lipschitz inequality for $f$ shows that $F=f$ on $Y$; the
triangle inequality shows that $F$ has the same Lipschitz constant.
The three coordinate extensions together have Lipschitz constant
$2\sqrt3\rho^{-2}a$. Composition with $I_q$ gives $H$.

\emph{The cell geometry and the comb.}
Put $N=\lfloor\kappa^{-1/2}\rfloor$. Then

\[
 N^2\kappa\le1,\qquad N\kappa\le N^{-1},\qquad
 N\ge\lambda_{{\rm BG},3},\qquad N^{-1}\le2\sqrt\kappa.
\]
For a coordinate rectangle or box with side lengths between $N$
and $2N-1$, any two boundary sites are joined along its boundary
by a path of length at most three times their $\ell^1$ distance.
For sites on one face, or on intersecting faces, a monotone path
has exactly that distance. For opposite faces, use either one of
the two intervening faces. If their separating direction is $\mu$
and the chosen other direction is $\nu$, the shorter detour costs
at most $n_\nu+n_\mu$ in those directions; this is at most
$3n_\mu$, since $n_\nu<2n_\mu$. The remaining coordinate contributes
its ordinary difference. The rectangle proof is the same argument
with two coordinates.

In a cell with lowest corner $y$, let $r_y(x)$ be the transport
from $y$ to $x$, first in direction $1$, then $2$, then $3$.
For $b=\langle x,x+e_\mu\rangle$ put
$m_b=r_y(x)U(b)r_y(x+e_\mu)^{-1}$. Moving the last step through
the later directions in the comb exchanges one plaquette at a time.
The holonomy of a strip of $k$ plaquettes is a product of $k$
unitary conjugates of those plaquettes, so its distance from
$\mathbf1$ is at most $k\kappa$. Thus, writing $x-y=(k_1,k_2,k_3)$,

\[
 |m_b-\mathbf1|\le
 \begin{cases}
 (k_2+k_3)\kappa,&\mu=1,\\
 k_3\kappa,&\mu=2,\\
 0,&\mu=3.
 \end{cases}
\]
This is the comb calculation of Lemma~\YMMref{14.4}{64}, with
the coordinate order reversed. It proves the estimate on an
ordinary coordinate box and on the lift of a periodic cell,
without any condition on the ambient torus side. On a face the
bound is $2/N$, and on a cell it is $4/N$.
If $n<N$ in the nonperiodic case, the same comb on the whole box
gives $2n\kappa\le4\sqrt\kappa$, which proves the assertion.

Suppose now $n\ge N$. Divide each coordinate interval into lengths
between $N$ and $2N-1$: take $t_i=iN$ for
$0\le i<\lfloor n/N\rfloor$ and the last endpoint $n$.
Their products give nodes, edges, faces and cells. Each new
piece meets the union of lower-dimensional pieces exactly in its
boundary; different interiors are disjoint. Every bond is contained
in a cell. We can therefore define the gauge successively on nodes,
edges, faces and cells.

\emph{Edges.}
Set $u=\mathbf1$ at the nodes. On an edge from $y$ to
$y'=y+n_ce_\mu$, write $U([y,y'])=e^{iY}$ with $|Y|\le\pi$, and set

\[
 u(y+ke_\mu)=e^{-ikY/n_c}U([y,y+ke_\mu]),
 \qquad 0\le k\le n_c.
\]
Both endpoint values are $\mathbf1$. At an edge bond the transformed
field is $e^{iY/n_c}$, so its distance from $\mathbf1$ is at most
$\pi/N$.

For the next two steps, write $h=ur_y^{-1}$ on the prescribed
part of a piece and $w_b=U^u(b)$. Since
$h(b_-)m_bh(b_+)^{-1}=w_b$, we have
\begin{equation}\label{eq:box-gauge-boundary-comparison}
 \begin{split}
 |h(b_+)-h(b_-)|&\le |w_b-\mathbf1|+|m_b-\mathbf1|,\\
 |w_b-\mathbf1|&\le |m_b-\mathbf1|+|h(b_+)-h(b_-)|.
 \end{split}
\end{equation}

\emph{Faces.}
On the boundary of a face, the bond differences of $h$ are at most
$\lambda_{{\rm BG},1}/N$; the boundary-path bound makes $h$
$3\lambda_{{\rm BG},1}/N$-Lipschitz in $\ell^1$ distance.
The boundary cycle has at most $8N$ bonds. Sample it at intervals
$s_1=\lfloor\rho_{{\rm BG},1}N/\lambda_{{\rm BG},1}\rfloor$.
The constants ensure $s_1\ge\rho_{{\rm BG},1}N/(2\lambda_{{\rm BG},1})\ge1$.
The image is covered by balls of radius $\rho_{{\rm BG},1}$ whose
number is at most

\[
 16\lambda_{{\rm BG},1}/\rho_{{\rm BG},1}+1
 \le17\lambda_{{\rm BG},1}/\rho_{{\rm BG},1}
 <21853<(8\rho_{{\rm BG},1})^{-3}.
\]
The avoidance and extension arguments give an extension $H$ over
the face with bond differences bounded by
$12\sqrt3\lambda_{{\rm BG},1}\rho_{{\rm BG},1}^{-2}/N$.
Set $u=Hr_y$ there. It retains the prescribed edge values, and
\eqref{eq:box-gauge-boundary-comparison} gives
$|U^u(b)-\mathbf1|\le\lambda_{{\rm BG},2}/N$ on every face bond.

\emph{Cells.}
On the boundary of a cell, the bond differences of $h$ are at most
$\lambda'_{{\rm BG},2}/N$, and its $\ell^1$ Lipschitz constant is
$3\lambda'_{{\rm BG},2}/N$. On each face sample both coordinates
at multiples of
$s_2=\lfloor\rho_{{\rm BG},2}N/(2\lambda'_{{\rm BG},2})\rfloor$.
Here $s_2\ge\rho_{{\rm BG},2}N/(4\lambda'_{{\rm BG},2})\ge1$.
Lowering the two coordinates to these multiples uses fewer than
$2s_2$ steps, so the image lies within $\rho_{{\rm BG},2}$ of a
sampled value. The number of samples on all six faces is at most

\[
 m=486\lambda_{{\rm BG},2}'^2\rho_{{\rm BG},2}^{-2},\qquad
 m(8\rho_{{\rm BG},2})^3\le248832/250000<1.
\]
The same extension gives $H$ over the cell, with bond differences
at most $12\sqrt3\lambda'_{{\rm BG},2}\rho_{{\rm BG},2}^{-2}/N$.
Again set $u=Hr_y$. The boundary values agree, and
\eqref{eq:box-gauge-boundary-comparison} gives

\[
 |U^u(b)-\mathbf1|\le\lambda_{{\rm BG},3}/N
 \le C_{\rm BG}\sqrt\kappa.
\]
Every bond belongs to a cell, proving the box assertion.

\emph{The periodic lattice.}
Divide each coordinate circle of length $n$ into the same integer
lengths, now identifying its last endpoint with its first. Work
on each piece in its coordinate lift. Put $u=\mathbf1$ at the
periodic nodes and perform the edge, face and cell constructions
just given once for each periodic piece. The edge formula agrees
at both endpoints even when an edge is a full circle. At every
later stage the boundary values are the values of the already
defined periodic function $u$. The auxiliary function $h=ur_y^{-1}$
is defined on the lifted boundary; it need not have equal values
on opposite lifted faces, since the comb transports there may
differ. After extension, however, $Hr_y$ equals the prescribed
$u$ at every boundary point. Thus those values agree at all
identified points. No new interior point has two preimages,
even when a cell spans an entire coordinate circle. The extensions
therefore descend to a single gauge on the torus. Every periodic
bond has a lift in one cell, with the same plaquette and comb
bounds as above. This proves the periodic assertion.
\end{proof}

For completeness, this also gives the first component assertion
of Lemma~\ref{C-367-11}(a). Let
$0<\kappa'\le(8\pi C_{\rm BG})^{-2}$ and use the periodic gauge.
Write $V^u=e^{iA}$, choosing the principal logarithm at positive
bonds and extending by reversal. Then

\[
 |A|\le\alpha:=\tfrac\pi2 C_{\rm BG}\sqrt{\kappa'}\le1/16.
\]
For the product of four exponentials around a plaquette,
Taylor's formula in $t$ gives
$|\partial e^{itA}(p)-\mathbf1-it\,dA(p)|\le8\alpha^2$:
the norm of its second derivative is bounded by
$(\sum_{b\in\partial p}|A(b)|)^2\le16\alpha^2$.
Here $dA(p)$ is the unnormalised oriented sum around $p$.
In particular, the bound $10\alpha^2$ used in
Lemma~\ref{C-367-11} holds. Applying it at $t$ and at $1$ yields

\[
 |\partial e^{itA}(p)-\mathbf1|
 \le |\partial V^u(p)-\mathbf1|+20\alpha^2
 \le(1+5\pi^2C_{\rm BG}^2)\kappa'
 <C''\kappa',\qquad 0\le t\le1.
\]
The gauge group is a finite product of connected copies of $SU(2)$,
so a path of gauge transformations joins $V$ to $V^u$ without
changing its plaquette norms. Joining this path to $e^{itA}$
places $V$ in the component of $\mathbf1$ in
$\{\varpi<C''\kappa'\}$. Taking $\kappa'=r/C''$ gives the shell
assertion there. This is a single-scale construction; it does not
prove the separate hypothesis \textup{(Multiscale gauge)}.

\section{Response and one-step conventions}
\label{sec:response-step-conventions}

We record the response bound needed when differentiating the
normalised measure of a step. Work on a constant torus sequence,
at level $j\ge1$, with the real regular background, geometric
thresholds and datum ball of Section~\ref{sec:background-step-complements}.
Lengths in the following lemma are in level-$j$ units, so the bare
spacing is $\eta_j=L^{-j}$ and every terminal pillar has scale $1$.
Let $d_j$ be the periodic $\ell^1$ distance on the terminal lattice.
In this setting it is the weighted path distance in
equation~\YMMref{(838)}{399}. For a field $X$ on its positive bonds put
\[
 m_y(X)=\sum_{c>0}e^{-\delta_0d_j(y,c_-)}|X(c)|,
 \qquad
 S_0=\left(\frac{1+e^{-\delta_0}}{1-e^{-\delta_0}}\right)^3,
 \qquad \mathsf S=3S_0,
\]
where $\delta_0>0$ is no larger than the rate in that equation.

\begin{lemma}\label{lem:response-current-zero}
There is a constant $\mathsf B$, depending only on the fixed
construction constants, such that the response $\mathcal H_U$ obeys
\[
 \max\{|\mathcal H_U(X)(b)|,
                  |\nabla_U\mathcal H_U(X)(b)|\}
 \le\mathsf Bm_y(X)
\]
on its datum ball, whenever $b$ has an endpoint in the pillar of $y$.
At zero datum, put
\[
 h=D\mathcal H_U(0)X,\qquad
 j_X=\left.\frac d{ds}J^{(j)}
          (e^{i\eta_j\mathcal H_U(sX)}U)\right|_{s=0}.
\]
Then
\begin{equation}\label{eq:response-current-zero}
 \max\{|h(b)|,|\nabla_Uh(b)|,|j_X(b)|\}
 \le\mathsf Bm_y(X),\qquad
 \sum_ym_y(X)^2\le\mathsf S^2\sum_{c>0}|X(c)|^2.
\end{equation}
All constants are independent of $j$ and of the torus side.
\end{lemma}

\begin{proof}
For one input bond, equation~\YMMref{(838)}{399} bounds
$h$, $\nabla_Uh$ and $D_U^{\mathsf T}D_Uh$ with exponential decay.
Sum this bound over the input bonds. Integrating the same estimate
along $sX$, $0\le s\le1$, proves the first assertion, since
$\mathcal H_U(0)=0$ and the datum ball is convex.

The current requires one further identity. Let $F(A)$ be the local
Wilson action at $e^{i\eta_jA}U$ and set
$f(z)=(e^z-1)/z$, with its power-series value at zero.
Proposition~\YMMref{13.1}{53} gives, on a complete plaquette star,
\[
 J^{(j)}(e^{i\eta_jA}U)(b)
   =f(-i\eta_j\operatorname{ad}_{A(b)})^{-1}F'(A)(b),
 \qquad F''(0)=D_U^{\mathsf T}D_U+\Delta'.
\]
Since $f(z)^{-1}=1-z/2+O(z^2)$, differentiation at $A=0$ gives
\[
 j_X(b)=(D_U^{\mathsf T}D_Uh)(b)+(\Delta'h)(b)
                  +\tfrac{i\eta_j}{2}[h(b),J^{(j)}(U)(b)].
\]
All stars are complete on the torus. The same proposition bounds
$\Delta'h$ by $C\varepsilon$ times the supremum of $h$ on this
fixed star, where $\varepsilon$ is the background regularity
radius. Regularity also gives $|J^{(j)}(U)|\le\varepsilon$.
The star meets only pillars at a fixed distance from $y$; translating
the exponential weight over that distance costs a fixed factor.
Together with $\eta_j\le1$ and $|[h,J]|\le2|h||J|$, the three terms
therefore have the asserted bound, after increasing $\mathsf B$.

Finally the sum of $e^{-\delta_0d_j(y,z)}$ over a periodic lattice
is at most $S_0$: choose least coordinate representatives and bound
their sum by the sum over $\mathbb Z^3$. There are three positive
bonds starting at each site. The nonnegative kernel defining
$m_y$ has row sums at most $3S_0$ and column sums at most $S_0$.
Cauchy--Schwarz, first with this kernel and then after summing in
$y$, gives an operator norm at most $\sqrt3S_0\le\mathsf S$.
\end{proof}

The gauge equation used with this lemma is exactly
equation~\YMMref{(837)}{398}. Thus the projection denoted by
$\mathsf R^{\rm La}$ in the dynamics section is the companion's
$\mathsf R_{\mathsf m}$: the projection onto $D_U^{\mathsf T}D_U\mathcal N$,
where $\mathcal N$ consists of site fields with zero prescribed site
averages. This is the mean-Landau choice of Theorem~\YMMref{K}{269}.

\subsection{The local input and its global sum}
\label{sec:step-input-conventions}

At height $k$, the actual input is the one in
Section~\YMMref{52.4}{422} and Theorem~\YMMref{W}{423}.
Keep its geometric thresholds, admissible radius system, shift
hypotheses, octet condition when $k\ge2$, numerical ceilings and
array bound $N\le N_*$. In particular, $E_0>0$ is the fixed size
of the output family there; it is chosen before the height, lattice
and input array. For the density identity, keep in addition the
whole-data response neighbourhood, full-torus core membership and
bounded Borel integrand required in that theorem.

A family in the fourth place means the fourth potential of
Definition~\YMMref{52.1}{415}. Its members depend on the represented
pair, are gauge invariant, vanish at $B=0$, and are jointly
holomorphic on each core slice and the field polydisc
$|B|<C_{\mathcal A}\mathfrak b(\mathsf t)$. With additional parameters,
joint holomorphy is required in those parameters too. Its size is
equation~\YMMref{(888)}{415}. For a parameter space consisting of
one point this is the same definition without a parameter.

The phrase that this family \emph{sums to $S$ at $W$} has the
following precise meaning. The averaged background is real,
admissible, and has invertible pivots; the real cutoff chart is
defined throughout its box, every background over a point of
that chart exists, and $P_W$ belongs to the full-torus core.
The gauge-invariant function $S$ is defined on both critical
orbits. For every real cutoff field,
\begin{equation}\label{eq:fourth-family-global-sum}
 \sum_Y\mathbb V'(Y;P_W|_Y,B|_Y)
   =S(U_k(V_W(B)))-S(U_{k+1}(W)).
\end{equation}
These are the background and chart conditions in
Section~\ref{sec:normalized-step-measure}. For a bare-array
potential the identity is equation~\YMMref{(882)}{413}.
For a source it is the additional sum condition imposed at its
entry into the step.

The block field of a triple is
$V[P]=\mathcal M^k(e^{i\eta_{k+1}A'}U)$ on the level-$k$ bonds
of its domain, with the actual complex averaging domains retained.
The fine-cube gauge form $\widehat{\mathsf G}_\square(P)$ is evaluated
at this field by equation~\YMMref{(823)}{394}. Equivalently it is
the quadratic form obtained from $\mathsf t\Psi^{\rm G}(\square;P,B)$
in the proof of Proposition~\YMMref{51.4}{393}.
Definition~\YMMref{52.1}{415},
item~2, records its pair dependence, slice holomorphy, covariance
and compatibility with restrictions; its sum on the full torus
is equation~\YMMref{(828)}{395}. At real data the inverse whole
covariance differs from its real reference form, with floor
$\gamma_{\rm r}$ and ceiling $c_{\rm r}$, by at most
$\frac12u_{\rm G}(LM)^{-3}$. These are the actual constants of
Theorem~\YMMref{U}{378}, used with the choice $u_{\rm G}=4C_{\rm Str}$
and the radius inequalities of Section~\YMMref{52.4}{422}.
Completeness is Definition~\YMMref{47.2}{359}, condition (St4).

The expansion numbers $w_0,\mathsf C_{\rm ex},c_G,c_{\rm v},c_{\rm w},
\lambda_1,\lambda_V,w_\star$ are those defined immediately before
and in equation~\YMMref{(705)}{360}; in particular
$\mathsf C_{\rm ex}\ge8e^8$.
Their one-step substitution is equation~\YMMref{(889)}{415},
including $\lambda_{\rm out}=(1+4\beta_{\rm s})\kappa$ and
$(LM)^3u=3+c_{\rm r}+u_{\rm G}$.

\section{Geometric convergence of the Gaussian one-loop coefficient}
\label{sec:geometric-oneloop}

The bounded Gaussian coefficient has a limit as the number of
averaging steps increases. We prove this by comparing two successive
meshes in units in which the terminal lattice has spacing one.
Throughout this section the group is $\mathrm{SU}(2)$, $L\ge13$ is odd,
and the averages, Gaussian normalisers and coefficient functional
are those of Sections~\YMMref{42}{318}--\YMMref{43}{324}.
In particular, Theorem~\YMMref{O}{325} gives
\[
 c_k=\mathsf b[\Psi_k]=\sum_{l=1}^kL^{-(k-l)}b_l^0,
 \qquad c_0=0,\qquad
 c_2^{\rm loop}:=\sup_{k\ge0}|c_k|<\infty.
\]
We write $c_k$ here for the numbers denoted by $\mathsf c_k$ elsewhere.

\begin{theorem}\label{III-12-105}
There are real numbers $c_\infty$ and
$b_\infty^0=(1-L^{-1})c_\infty$, depending only on $L$, and a finite
constant $C_{\rm geo}=C_{\rm geo}(L)$ such that
\[
 |c_k-c_\infty|\le\tfrac12C_{\rm geo}L^{-k/6}\quad(k\ge0),
 \qquad
 |b_l^0-b_\infty^0|\le C_{\rm geo}L^{-l/6}\quad(l\ge1).
\]
Moreover
\[
 |b_{l+1}^0-b_l^0|
 \le(1-L^{-1/6})C_{\rm geo}L^{-l/6}\quad(l\ge1).
\]
\end{theorem}

The local determinant analysis already gives a uniformly holomorphic
remainder. The additional issue is convergence of that remainder
on changing Fourier spaces. We first remove the explicit massive
term, then compare the remaining finite expression. The comparison
will be made at real momentum; a separate finite-difference argument
will recover the derivative defining the coefficient.

\subsection{The explicit massive term}

Put $\eta=L^{-k}$ and
$\psi(q)=\sum_\mu(2-2\cos q_\mu)$, with normalised integration over
$[-\pi,\pi]^3$. The massive determinant calculation in
Proposition~\YMMref{43.3}{335} gives fixed real trigonometric
polynomials $g_1,g_2$ such that
\[
 \widetilde\beta(z)=I_2[g_2](z)+I_1[g_1](z),\qquad
 I_n[g](z)=\int(\psi(q)+z)^{-n}g(q)\,dq.
\]
More explicitly,
$g_2=4\operatorname{tr}\sigma_1^2-\psi_{11}\psi_{22}/3$ and
$g_1=-8\operatorname{tr}\sigma_2+(\psi^{(1)}+\psi^{(2)})/2$
in that calculation. Define
\[
 E(\eta)=\eta\bigl[\widetilde\beta(\eta^2)
       -\eta^2\widetilde\beta'(\eta^2)
       +\tfrac12\eta^4\widetilde\beta''(\eta^2)\bigr].
\]
With $\delta=\delta_{\rm loop}\le1/64$, let
\[
 F_k(p)=\widehat\pi^0_{12}(p)-\mathcal L_k^\sharp(p),\qquad
 r_k=c_k-E(L^{-k}).
\]
The determinant representation in Section~\YMMref{43.3}{329},
Proposition~\YMMref{43.2}{334} and the massive calculation give
\begin{equation}\label{eq:loop-remainder-reduction}
 r_k=-\partial_1\partial_2F_k(0),\qquad
 |F_k(p)|\le C_{\rm rem}\quad(|p_\mu|<\delta),
\end{equation}
with $C_{\rm rem}$ independent of $k$.

\begin{lemma}\label{lem:loop-explicit-limit}
There is an absolute real number $E_\infty$ such that
$C_E:=\sup_{0<\eta\le1}\eta^{-1}|E(\eta)-E_\infty|<\infty$.
\end{lemma}

\begin{proof}
Differentiation under the integral gives $I_n'=-nI_{n+1}$.
Writing $z=\eta^2$, we obtain
\[
 E(\eta)=z^{1/2}(I_2+2zI_3+3z^2I_4)[g_2](z)
            +z^{1/2}(I_1+zI_2+z^2I_3)[g_1](z).
\]
For a fixed trigonometric polynomial $g$, put
$i_n=(2\pi)^{-3}\int_{\mathbb R^3}(1+|y|^2)^{-n}\,dy$.
We claim
\begin{equation}\label{eq:loop-massive-integrals}
 |I_1[g](z)|\le C_g,\qquad
 |z^{n-3/2}I_n[g](z)-g(0)i_n|\le C_g\sqrt z
 \quad(n=2,3,4; 0<z\le1).
\end{equation}
Indeed, on the integration cube,
\[
 a|q|^2\le\psi(q)\le|q|^2,\qquad
 0\le|q|^2-\psi(q)\le|q|^4/12,\qquad a=4/\pi^2.
\]
The first bound follows by integrating $C|q|^{-2}$ over a fixed
ball. For the others replace $g$ by its even part, whose difference
from $g(0)$ is $O(|q|^2)$. The error is bounded by
$C\int_{|q|\le\sqrt3\pi}|q|^2(a|q|^2+z)^{-n}\,dq$.
Replacing $\psi$ by $|q|^2$ costs at most
$C\int_{|q|\le\sqrt3\pi}|q|^4(a|q|^2+z)^{-n-1}\,dq$.
For either integral, multiplication by $z^{n-3/2}$ gives
$O(\sqrt z)$ when $n=2$, by the integrable radial majorant
$C|q|^{-2}$, and $O(z)$ when $n=3,4$, by $q=\sqrt z\,y$.
Extending the cube to $\mathbb R^3$ costs
$O(z^{n-3/2})$, since its complement has $|q|\ge\pi$.
The same substitution then gives $i_n$ exactly. This proves
\eqref{eq:loop-massive-integrals}.

It follows that
$E_\infty=(i_2+2i_3+3i_4)g_2(0)$ and
$E(\eta)-E_\infty=O(\eta)$. The contribution of $g_1$ is
$O(\sqrt z)$ by the same estimates. Reality follows from that of
the determinant coefficient. Since $E$ is defined by
$\widetilde\beta$, its limit and $C_E$ do not depend on a choice
of the polynomial representation.
\end{proof}
\subsection{A finite expression for the remainder}

We use the physical pairings of weight $\ell_j^3$ on
$\ell_j\mathbb Z^3$, $\ell_j=L^{j-k}$, and the scaled exterior
derivative. The operators $D_U,Q_U,R_U,K_U,\Delta_U,\Xi_U,Y_U$
are those defined before Section~\YMMref{43.2}{327}. Thus
$Q_U=\eta\mathring{\mathcal Q}^{[k]}_U$ and
$R_U=\mathfrak a_U^{[k]}$. Unsubscripted $d,K,\Delta,Q,R$
denote their free values, and
$G=(K+1)^{-1}$, $g=(\Delta+1)^{-1}$.
The response is $h=\Xi^{-1}Q^{\mathsf T}
(Q\Xi^{-1}Q^{\mathsf T})^{-1}$.
Along
\[
 \zeta_\pm=\eta e^{\pm ip\cdot x}\alpha_\pm,\qquad
 \alpha_\pm=\widehat h(\pm p)e_{\mu_\pm}t_3,
\]
bold letters denote the jets with components $00,10,01,11$.
The two nonzero field directions and their curls are bounded
uniformly in $k$ on the fixed strip, by
Section~\YMMref{43.2}{327}. In this subsection either pair of
directions is allowed; we later take $(\mu_+,\mu_-)=(1,2)$.

For completeness, the real domains for these jets are
$K^{00}=K=[-\pi-\delta,\pi+\delta]^3$,
$K^{11}=[-\pi,\pi]^3$ and
$K^{10}=K\cap(K-p)$, $K^{01}=K\cap(K+p)$.
Their norm is the sum of the four supremum operator norms.
Products are the jet products: if $a,b\in\{0,1\}$, then
\begin{equation}\label{eq:loop-jet-product}
 (\mathbf T^1\cdots\mathbf T^s)_{ab}(\varkappa)
 =\sum_{\sum a_i=a,\,\sum b_i=b}
   \prod_{i=1}^s T^i_{a_ib_i}
       \left(\varkappa+\sum_{j>i}(a_j-b_j)p\right).
\end{equation}
Every argument lies in its stated domain. Consequently the norm
is submultiplicative, and a product containing three factors with
zero $00$ component vanishes. Transposes use the bilinear pairing;
their arguments are $-\varkappa$, $-\varkappa-p$,
$-\varkappa+p$, $-\varkappa$, respectively, and preserve the norm.
Put
\[
 \mathsf t(\mathbf T)=\int\operatorname{tr}T_{11}(\varkappa)
                      \,d\varkappa,
\]
and, when $C_{00}$ is invertible,
\[
 \ell(\mathbf C)=\int\operatorname{tr}\bigl[
 C_{00}^{-1}C_{11}
 -C_{00}(\varkappa)^{-1}C_{10}(\varkappa-p)
 C_{00}(\varkappa-p)^{-1}C_{01}(\varkappa)\bigr]\,d\varkappa.
\]
The integrals are normalised over $[-\pi,\pi]^3$.

The local Hessian formula of Section~\YMMref{43.4}{332} reads
$H_U=(d+E_U)^{\mathsf T}z_U(d+E_U)+f_U$.
Write $\mathbf e=\mathbf D-d$, and define
\[
 \kappa=G^{1/2}(\mathbf K-K)G^{1/2},\quad
 x=\mathbf QG^{1/2},\quad y=G^{1/2}\mathbf D\mathbf R^{\mathsf T},
 \qquad
 \kappa_g=g^{1/2}(\boldsymbol\Delta-\Delta)g^{1/2},\quad
 x_g=\mathbf Rg^{1/2}.
\]
All of these are jets. For either datum
$(\mathsf G,\kappa,x,y)=(G,\kappa,x,y)$ or
$(g,\kappa_g,x_g,0)$, put
\[
 \omega=x^{\mathsf T}x-yy^{\mathsf T},\quad
 B_t=(1+\kappa+t\omega)^{-1},\quad
 B^\flat_\tau=(1+\kappa+\omega-(1-\tau)\mathsf G)^{-1}.
\]
The strip estimate, Proposition~\YMMref{43.1}{327}, and the
bounded jets in Section~\YMMref{43.4}{332} bound all these jets
and inverses by one constant $A=A(L)$. Indeed the four free
$00$ inverses are precisely the sandwiched inverses in that
proposition; the inverse jet follows by multiplying out its
$10,01,11$ components, each involving at most three such inverses.
The $00$ inverse of $xB^\flat_0x^{\mathsf T}$ is bounded by a
constant $\iota$, also uniformly in $k$. These statements hold
for $0\le t,\tau\le1$, continuously in those parameters.

For these data define the following explicit expression:
\begin{equation}\label{eq:loop-Phi}
\begin{split}
 \Phi[\mathsf G,\kappa,x,y]={}&
 \int_0^1[\mathsf t(xB_tx^{\mathsf T})
                  -\mathsf t(y^{\mathsf T}B_ty)]\,dt
       +\ell(xB^\flat_0x^{\mathsf T})\\
 &-\int_0^1(1-\tau)^2
       \mathsf t(\mathsf GB_0\mathsf GB_0\mathsf GB^\flat_\tau)\,d\tau\\
 &+\int_0^1[\mathsf t(xB^\flat_\tau\mathsf GB_0x^{\mathsf T})
       -\mathsf t(y^{\mathsf T}B^\flat_\tau\mathsf GB_0y)]\,d\tau\\
 &+\int_0^1(1-\tau)[
       \mathsf t(xB^\flat_\tau\mathsf GB_0\mathsf GB_0x^{\mathsf T})
       -\mathsf t(y^{\mathsf T}B^\flat_\tau\mathsf GB_0\mathsf GB_0y)]\,d\tau.
\end{split}
\end{equation}
\begin{lemma}\label{lem:loop-finite-expression}
For real $|p_\mu|\le\delta$,
\[
 \widehat\pi^0_{\mu_+\mu_-}(p)-\mathcal L_k^\sharp(p)
       =-\tfrac12\Phi[G,\kappa,x,y]+\Phi[g,\kappa_g,x_g,0].
\]
\end{lemma}

\begin{proof}
Set $V=\mathbf Q^{\mathsf T}\mathbf Q
-\mathbf D\mathbf R^{\mathsf T}\mathbf R\mathbf D^{\mathsf T}$.
The Schur factorisation of the vector saddle matrix gives
$\ell(\boldsymbol\Xi)+
\ell(\mathbf Q\boldsymbol\Xi^{-1}\mathbf Q^{\mathsf T})$;
the scalar version gives
$\ell(\mathbf Y)+\ell(\mathbf R\mathbf Y^{-1}\mathbf R^{\mathsf T})$.
The triangular factors have determinant one and hence zero
polarisation. The coarse factors are invertible by the strip
estimate, so these identities take place inside the same jet algebra.

For either pair $(K_\circ,V_\circ)=(\mathbf K,V)$ or
$(\boldsymbol\Delta,\mathbf R^{\mathsf T}\mathbf R)$, let
$G_U=(K_\circ+1)^{-1}$ and
$P_\tau=(K_\circ+V_\circ+\tau)^{-1}$.
Integration of the logarithmic derivative along $tV_\circ$ and
then along $\tau$ gives
\[
 \ell(K_\circ+V_\circ)=\ell(K_\circ+1)
 +\int_0^1\mathsf t((K_\circ+1+tV_\circ)^{-1}V_\circ)\,dt
 -\int_0^1\mathsf t(P_\tau)\,d\tau.
\]
Twice using the resolvent identity yields exactly
\[
 P_\tau=G_U+(1-\tau)G_U^2
 +(1-\tau)^2G_U^2P_\tau-G_UV_\circ P_\tau
 -(1-\tau)G_U^2V_\circ P_\tau.
\]
The first two terms give
$[1-\partial_z+\tfrac12\partial_z^2]\ell(K_\circ+z)|_{z=1}$.
Substitute
$G_U=\mathsf G^{1/2}B_0\mathsf G^{1/2}$,
$P_\tau=\mathsf G^{1/2}B^\flat_\tau\mathsf G^{1/2}$ and
$\mathsf G^{1/2}V_\circ\mathsf G^{1/2}=\omega$ in the remaining
terms. Cyclicity of $\mathsf t$ gives the four integrals in
\eqref{eq:loop-Phi}, and the coarse Schur factor gives its last
logarithmic term. Finally take minus one half of the vector
identity plus the scalar identity. This also explains why the
only fine-space trace left has three propagator factors.
\end{proof}

\subsection{Comparing the two Fourier spaces}

We compare height $k$ with height $k+1$; primes indicate the latter.
Their bare spacings are $\eta$ and $\eta'=\eta/L$, and all later
lattices agree: $\Lambda'_{j+1}=\Lambda_j$.
Let $N=(\eta^{-1}-1)/2$ and
$\mathfrak m=\{m\in\mathbb Z^3:|m|_\infty\le N\}$, with
$\mathfrak m'$ defined at $\eta'$. Fourier coefficients identify
each fine site, bond or plaquette fibre isometrically with
$\ell^2(\mathfrak m;\mathsf V)$, where $\mathsf V$ is its colour
and component space. The bilinear pairing pairs $m$ with $-m$.
Let $J$ extend coefficients by zero from $\mathfrak m$ to
$\mathfrak m'$, and let $J=1$ on the common coarse fibres.
Then $J^{\mathsf T}=J^*$ and $J^*J=1$.
For $1\le\Lambda\le N$ let $P_\Lambda,P'_\Lambda$ retain
the modes $|m|_\infty\le\Lambda$. These projections commute
with $J$ in the corresponding sense and are their own transposes.

For two jets between corresponding fibres define
$\operatorname{df}(T',T)=\|T'J-JT\|$ in the real jet norm.
The product rule \eqref{eq:loop-jet-product} gives
\begin{equation}\label{eq:loop-defect-rules}
\begin{split}
 \operatorname{df}(T'S',TS)&\le
       \|T'\|\operatorname{df}(S',S)
                 +\operatorname{df}(T',T)\|S\|,\\
 \operatorname{df}(T'^{\mathsf T},T^{\mathsf T})&\le
       \operatorname{df}(T',T)+\|T'(1-JJ^*)\|,\\
 \operatorname{df}((1+S')^{-1},(1+S)^{-1})&\le
       \|(1+S')^{-1}\|\operatorname{df}(S',S)
                              \|(1+S)^{-1}\|.
\end{split}
\end{equation}
The first and third follow by inserting $J$ between the factors.
For the second, insert $JJ^*+(1-JJ^*)$ before transposing.
This extra tail will be retained throughout the comparison.

Here are the elementary multiplier estimates we need. For
$h=\eta,\eta'$ and $q=\varkappa+2\pi m$, write
\[
 d^h_\mu(q)=h^{-1}(e^{ihq_\mu}-1),\quad
 D_h(q)=\sum_\mu|d^h_\mu(q)|^2,\quad
 s_h(t)=h\sum_{|j|\le(h^{-1}-1)/2}e^{ihjt},
\]
\[
 v_h(q)=\prod_\nu s_h(q_\nu),\qquad
 u_{h,\mu}(q)=e^{i(1-h)q_\mu/2}s_h(q_\mu)v_h(q),\qquad
 g_h(q)=(1+D_h(q))^{-1/2}.
\]
Define $s_1(t)=L^{-1}\sum_{|j|\le(L-1)/2}e^{i\eta'jt}$,
$v_1=\prod s_1$ and
$u_{1,\mu}=e^{i(\eta-\eta')q_\mu/2}s_1(q_\mu)v_1$.
Division with centred remainder gives
\begin{equation}\label{eq:loop-finest-multipliers}
 s_{\eta'}=s_\eta s_1,\qquad
 v_{\eta'}=v_\eta v_1,\qquad
 u_{\eta',\mu}=u_{\eta,\mu}u_{1,\mu}.
\end{equation}
The free finest-step averages therefore act by
\begin{equation}\label{eq:loop-finest-aliases}
 (Q^{\rm f}\varphi')_{\mu,m}
 =\sum_{m'\equiv m}u_{1,\mu}(q_{m'})\varphi'_{\mu,m'},
 \qquad
 (R^{\rm f}\varphi')_m
 =\sum_{m'\equiv m}v_1(q_{m'})\varphi'_{m'},
\end{equation}
where congruence is modulo $\eta^{-1}\mathbb Z^3$.
They have norm at most one. The transpose of the second map
sends $\varphi_m$ to $v_1(q_{m'})\varphi_m$ at every such alias.
These formulas follow by averaging each plane wave over the
centred block and, for $Q^{\rm f}$, over its straight path.

Put $w(n)=\max(1,2|n|-1)^{-1}$,
$w^{(3)}(m)=\prod_\nu w(m_\nu)$ and
$w_3(m)=w^{(3)}(m)/\max(1,|m|_\infty)$.
For $\varkappa\in K$ and the respective mode sets,
\begin{equation}\label{eq:loop-elementary-bounds}
\begin{gathered}
 D_h(q)\ge\tfrac9{25}|q|^2,\quad
 |u_{h,\mu}(q)|,|v_h(q)|\le2w^{(3)}(m),\quad
 |q|\le11(1+|m|_\infty),\\
 |g_\eta-g_{\eta'}|\le6\eta,\quad
 |g_\eta^2-g_{\eta'}^2|\le3\eta^2,\quad
 |d^{\eta'}_\mu g_{\eta'}-d^\eta_\mu g_\eta|
                                      \le8\eta|q|,\\
 \sum_{|m|_\infty\ge\Lambda}w_3(m)\le252\Lambda^{-1/2},
 \quad\sum_mw_3(m)\le253,\quad
 \eta(3+\log\eta^{-1})^3\le49\eta^{1/2}.
\end{gathered}
\end{equation}
For the middle line the modes are common modes. To verify the
bounds, use
$h^{-1}|\sin(ht/2)|\ge3|t|/10$ for
$|t|\le\pi/h+1$, and
$h^{-1}\sin(ht/2)s_h(t)=\sin(t/2)$.
Taylor's formula gives
$|d^\eta-d^{\eta'}|\le\eta|q|^2$ and
$|D_\eta-D_{\eta'}|\le\eta^2|q|^4/12$;
the root estimates follow by subtracting the two positive
square roots. Also
\[
 |1-v_1|\le\eta^2|q|^2/8,\quad
 |1-u_{1,\mu}|\le\eta|q_\mu|/2+\eta^2|q|^2/4,
\]
and $1-|u_{1,\mu}|^2,1-v_1^2\le\eta^2|q|^2/2$.
For the sums, $\sum_nw(n)^2\le4$ and its tail from $n_0$
is at most $4/n_0$. On the shell $|m|_\infty=n$,
\[
 \sum w_3(m)\le6n^{-2}(3+\log 2n)^2\le84n^{-3/2}.
\]
Summing the shells proves the displayed tail. The final bound
follows by maximising $e^{-u/2}(3+u)^3$ at $u=3$.

\subsection{The backgrounds at the two heights}

The comparison must include the response fields themselves.
Write their components as
$\alpha_{\pm,\mu}=a^\pm_\mu t_3$ and
$f_{\pm,\mu\nu}=f^\pm_{\mu\nu}t_3$, where
$d(e^{\pm ipx}\alpha_\pm)=e^{\pm ipx}f_\pm$.
For any one of these scalar functions write $\widehat a_m$
for its Fourier coefficient. Constants below depend only on $L$.

\begin{lemma}\label{lem:loop-background-comparison}
For common modes,
\[
 |\widehat a_m|,|\widehat a'_m|\le Cw_3(m),\qquad
 |\widehat a'_m-\widehat a_m|\le C\eta w^{(3)}(m).
\]
The first bound also holds on all primed modes. Consequently
\begin{equation}\label{eq:loop-background-l1}
 \sum|\widehat a_m|\le C,\quad
 \sum_{|m|_\infty\ge\Lambda}|\widehat a_m|
                 \le C\Lambda^{-1/2},\quad
 \sum_{m\in\mathfrak m}|\widehat a'_m-\widehat a_m|
   +\sum_{m\in\mathfrak m'\setminus\mathfrak m}|\widehat a'_m|
                 \le C\eta^{1/2}.
\end{equation}
Moreover $Q^{\rm f}(e^{\pm ipx}\alpha'_\pm)$ differs from
$e^{\pm ipx}\alpha_\pm$ in supremum norm by $C\eta^{1/2}$,
and $e^{\pm ipx}a^{\prime\pm}_\mu(x)$ has oscillation at most
$C\eta^{1/2}$ between sites at distance at most $\eta$.
\end{lemma}

\begin{proof}
It suffices to use the plus sign. For
$b=\widehat h(p)\varepsilon$, $n=\widehat\nu(p)\varepsilon$
and $\gamma=\widehat R(p)\widehat{d^{\mathsf T}}(p)b$,
Section~\YMMref{43.2}{327} gives
\begin{equation}\label{eq:loop-mode-response}
 D_\eta(q_m)b_{\mu,m}
  =u_{\eta,\mu}(-q_m)n_\mu
                  +d^\eta_\mu(q_m)v_\eta(q_m)\gamma,
 \qquad\|b\|,|n|,|\gamma|\le C_h|\varepsilon|.
\end{equation}
For $m\ne0$, $|q_m|\ge\pi|m|_\infty$.
The multiplier bounds give $|b_m|\le Cw_3(m)$.
Taking a curl cancels the two $\gamma$ terms and gives
\[
 f_{\mu\nu,m}=
 \frac{d_\mu(q_m)u_\nu(-q_m)n_\nu
          -d_\nu(q_m)u_\mu(-q_m)n_\mu}{D_\eta(q_m)},
\]
which obeys the same bound. The zero modes are bounded by the
response norm.

To compare zero modes without dividing by $|p|$, use
$Qh=1$, $d^{\mathsf T}\nu=0$ and $\sum_mv(q_m)^2=1$.
They imply
\[
 \gamma=v(p)^{-1}\sum_\nu d_\nu(-p)b_{\nu,0},\qquad
 n_\mu=u_\mu(-p)^{-1}
       [D_\eta(p)b_{\mu,0}-d_\mu(p)\sum_\nu d_\nu(-p)b_{\nu,0}].
\]
Here $v(p)\ge(2/\pi)^3$ and $|u_\mu(p)|\ge(2/\pi)^4$.
Set
$A_\mu=\sum_{m\ne0}|u_\mu(q_m)|^2/D_\eta(q_m)$ and
$A_0=\sum_{m\ne0}v(q_m)^2/D_\eta(q_m)$.
Write $d_\mu=d^\eta_\mu$ and $d^{\rm c}_\mu(p)=e^{ip_\mu}-1$.
The constraint becomes $Hb_0=\varepsilon$, where
\[
 H_{\mu\nu}=\delta_{\mu\nu}
       [u_\mu(p)+A_\mu D_\eta(p)/u_\mu(-p)]
 +[d^{\rm c}_\mu(p)A_0/v(p)
                      -A_\mu d_\mu(p)/u_\mu(-p)]d_\nu(-p).
\]
This holds for every $\varepsilon$, so $H$ is invertible and
$\|H^{-1}\|\le C_h$ by the response bound.
Factorisation \eqref{eq:loop-finest-multipliers} gives
$|A'_\mu-A_\mu|+|A'_0-A_0|\le C\eta^2$:
on common nonzero modes the reciprocal-denominator and averaging
defects are $O(\eta^2)$ times $|u|^2$ or $v^2$, whose sums are
at most one; on added modes $|q_m|^{-2}\le C\eta^2$.
All other entries in $H$ change by $O(\eta)$, with denominators
bounded away from zero. Thus $H'-H=O(\eta)$ and
$b'_0-b_0,n'-n,\gamma'-\gamma=O(\eta)$.

Subtract \eqref{eq:loop-mode-response} and the curl formula
at a common nonzero mode, one factor at a time. The defects of
$u,v,d,D^{-1}$ are bounded respectively by
$C\eta|q|w^{(3)}$, $C\eta^2|q|^2w^{(3)}$,
$\eta|q|^2$ and $C\eta^2$.
Since $|q|\ge1$ and $\eta|q|\le6$, each resulting term is
$O(\eta w^{(3)}(m))$. This proves the coefficient comparison.
Summation uses \eqref{eq:loop-elementary-bounds}; in particular
$\eta\sum_{|m|_\infty\le N}w^{(3)}(m)
\le\eta(3+\log\eta^{-1})^3$.
The added-mode tail starts at $N+1\ge(2\eta)^{-1}$.

Finally use the alias formula \eqref{eq:loop-finest-aliases}.
The principal alias contributes the coefficient difference and
$(u_1-1)\widehat a'_m$, bounded by
$C\eta(1+|m|_\infty)w_3(m)$; other aliases are controlled by
the added-mode tail. For the oscillation, multiply the coefficients
by $\min(2,19\eta(1+|m|_\infty))$ and split at $N$.
The same two sums give $C\eta^{1/2}$.
\end{proof}

\subsection{The actual averaging maps}

We now compare the field-dependent averages. This requires their
actual local definitions: equality of the free multipliers alone
would not identify their derivatives. Let $\rho_*,C_*$ be the
uniform field radius and bounds from Section~\YMMref{43.4}{332},
and $C_h$ the response bound from Section~\YMMref{43.2}{327}.
Enlarge $C_*$ to be at least $16$, so it also bounds
$\ell_j^{-1}|\mathcal M^j(e^{i\eta a})-1|_\infty/|a|_\infty$.
Decrease $\rho_*$ if necessary so that $\rho_*\le1$, and put
\[
 r_1=\frac{\rho_*}{5C_*},\qquad
 r_2=\frac{\rho_*}{2C_*L},\qquad r_0=\frac{r_1}{2C_h}.
\]
For $|s_\pm|<r_0$ the fields
$a(s)=s_+e^{ipx}\alpha_++s_-e^{-ipx}\alpha_-$ and $a'(s)$
have supremum norm less than $r_1$.

A dilation taking one averaging step to another takes centred
blocks, coordinate contours and their orders to the corresponding
blocks and contours. It therefore preserves their products and
the equations defining the local mean. The holomorphic mean is
the branch through the identity, so uniqueness on its defining
polydisc preserves this branch as well. Differentiating these
identities preserves the corrected differential average and the
scalar transport average. In the operator norms the factor
$\lambda^3$ in the physical pairing occurs on both sides and
cancels. Thus the estimates for a step with terminal spacing one
apply to every step. Splitting the product after the first step
is consequently the exact identity
\begin{equation}\label{eq:loop-actual-split}
 \mathbf Q'=\mathbf Q^\sharp\mathbf Q^{\rm f},\qquad
 \mathbf R'=\mathbf R^\sharp\mathbf R^{\rm f}.
\end{equation}
Here the sharp maps are the remaining height-$k$ averages
evaluated at $\mathcal M(e^{i\eta'a'(s)})$.

We record explicitly how field bounds give bounds on these jets.
If a local operator function $T(s)$ is holomorphic on
$|s_\pm|<r_0$ and $\|T(s)-T(0)\|\le M$, pairing with two
finitely supported unit vectors and applying Cauchy's formula
bounds its $10,01,11$ derivatives by $M/r_0,M/r_0,M/r_0^2$.
For real $p$, the first derivative kernels become periodic after
multiplication at the output by $e^{\mp ipx}$, an isometry;
the mixed derivative kernel is already periodic. The Fourier
symbol norm is bounded by the corresponding operator norm.
This proves the same bounds in the real jet norm. Applying this
argument to $T(s)-\widetilde T(s)$ also proves the difference
version, including the $00$ component when it is nonzero.

For a single step and an absorbed field $A$, the uniform bound
and Schwarz's lemma on $|A|<\rho_*/L$ give
\[
 \|L^{-1}(\mathring{\mathcal Q}_{e^{iA}}
                 -\mathring{\mathsf Q})\|
 +\|\mathfrak a_{e^{iA}}-\mathfrak a\|
 \le4C_*L\rho_*^{-1}|A|_\infty.
\]
Indeed each difference vanishes at zero and is bounded by $2C_*$
on that ball; apply the scalar Schwarz lemma after pairing.
For the intermediate logarithm
$A_{j-1}=i^{-1}\log\mathcal M^{j-1}(e^{i\eta a})$,
the uniform mean bound gives
$|A_{j-1}|_\infty\le2C_*\ell_{j-1}|a|_\infty$.
The preceding estimate and the jet Cauchy estimate therefore give
\begin{equation}\label{eq:loop-step-jets}
 \mathbf A_j=\mathfrak a+\mathbf a_j,\quad
 (\mathbf a_j)_{00}=0,\quad
 \|\mathbf a_j\|\le c_a\ell_{j-1},\qquad
 \|\mathbf Q^{\rm f}-Q^{\rm f}\|
 +\|\mathbf R^{\rm f}-R^{\rm f}\|\le C\eta,
\end{equation}
where $\mathbf A_j$ is the jet of the $j$th scalar step and
$c_a=(8C_*^2L/\rho_*)r_1(2r_0^{-1}+r_0^{-2})$.
All these constants depend only on $L$.

The scalar average also suppresses high Fourier modes:
\begin{equation}\label{eq:loop-R-tail}
 \|\mathbf R(1-P_\Lambda)\|,
 \|(1-P_\Lambda)\mathbf R^{\mathsf T}\|
 \le28(1+c_a)^2\Lambda^{-1/2}.
\end{equation}
To prove this, the product $a^{[n]}$ of the first $n$ free steps
has multiplier
$v^{(n)}(q)=\prod_\nu L^{-n}\sum_{|j|\le(L^n-1)/2}e^{i\eta jq_\nu}$.
Its output modes have aliases $m=\bar m+\ell_n^{-1}e$.
The sine identity for this finite sum gives
$|v^{(n)}(q_m)|^2\le2w^{(3)}(e)^2$.
If $|m|_\infty>\Lambda$ and $\ell_n\Lambda\ge1$, some
$|e_\nu|>\ell_n\Lambda-1/2$; thus the sum of the squares in
each alias class is at most $768/(\ell_n\Lambda)$.
Cauchy--Schwarz gives
$\|a^{[n]}(1-P_\Lambda)\|\le28(\ell_n\Lambda)^{-1/2}$.
When $\ell_n\Lambda<1$ this follows instead from the contraction
bound. Expand $\mathbf R=\prod_{j=k}^1(\mathfrak a+\mathbf a_j)$.
There are only the free term, terms with one insertion, and
terms with two insertions: three zero-$00$ jets have zero product.
Applying the free tail bound to the rightmost free segment gives
\[
 \|\mathbf R(1-P_\Lambda)\|
 \le28\Lambda^{-1/2}
 \left(1+c_a\sum_j\ell_{j-1}^{1/2}
       +c_a^2\sum_{i<j}\ell_{j-1}\ell_{i-1}^{1/2}\right).
\]
Since $\sum_j\ell_{j-1}^{1/2}\le(\sqrt L-1)^{-1}\le1$,
this proves the first bound; transposition proves the second.

It remains to compare the sharp maps with the unprimed maps.
For the finest step define
\[
 \mathcal N(A)=i^{-1}\log\mathcal M(e^{iA})-\mathsf QA.
\]
On $|A|_\infty\le r_2$, the mean differs from the identity
by at most $\rho_*/2$, and hence the logarithm is defined.
The logarithm and the linear average bound $\mathcal N$ by $3$.
Its value and first derivative vanish at zero because
$\mathcal M'(1)=\mathsf Q$. Schwarz's lemma for a double zero gives
$|\mathcal N(A)|_\infty\le3r_2^{-2}|A|_\infty^2$.
The relation between the ordinary and corrected linear averages
is $\mathsf Q=\mathring{\mathsf Q}-d_0\mathsf n$, with $d_0$
the unscaled coarse difference. Therefore, for $A=\eta'a'(s)$,
\begin{equation}\label{eq:loop-mean-comparison}
 a^\sharp=Q^{\rm f}a'(s)
       -\eta^{-1}d_0\mathsf n(\eta'a'(s))
       +\eta^{-1}\mathcal N(\eta'a'(s)).
\end{equation}
The operator $\mathsf n$ annihilates constant forms; its paths
have at most $3L/2$ bonds within one block. Subtract a constant
value on each such block and use the oscillation bound in
Lemma~\ref{lem:loop-background-comparison}. The middle term in
\eqref{eq:loop-mean-comparison} is bounded by $C\eta^{1/2}$;
the last is bounded by $C\eta$. The first differs from $a(s)$
by $C\eta^{1/2}$. Thus $|a^\sharp-a(s)|_\infty\le C\eta^{1/2}$
on the fixed bidisc. Both fields remain within $\rho_*$,
and $|a(s)|<\rho_*/2$.

The uniformly bounded full averages are Lipschitz between such
fields, with constant $4C_*/\rho_*$. To see this, restrict to the
complex line from the first field to the second. If their distance
is at most $\rho_*/2$, Schwarz's lemma on the disc of radius
$\rho_*/(2|a_1-a_0|)$ gives the bound; for a larger distance the
bound $2C_*$ suffices. Applying the difference Cauchy estimate
on the bidisc now proves
\begin{equation}\label{eq:loop-sharp-comparison}
 \|\mathbf Q^\sharp-\mathbf Q\|
       +\|\mathbf R^\sharp-\mathbf R\|\le C\eta^{1/2}.
\end{equation}

\subsection{The local multiplication terms}

We next compare the curvature and covariant-derivative jets.
The comparison uses their actual first derivatives, including
the lattice shifts. For a periodic scalar function $a$, let
\[
 M_h[a;\varpi;\chi_1,\chi_2]_{m,\widetilde m}
 =\widehat a_{m-\widetilde m}
 e^{ih(\varpi+2\pi(m-\widetilde m))\cdot\chi_1}
 e^{ihq_{\widetilde m}\cdot\chi_2},
\]
where the coefficient index is read modulo $h^{-1}\mathbb Z^3$,
$\varpi=\pm p$ and $\chi_1,\chi_2\in\{-1,0,1\}^3$.
In position space this is multiplication by
$e^{i\varpi\cdot(x+h\chi_1)}a(x+h\chi_1)$ followed by
the shift $x\mapsto x+h\chi_2$. Schur's inequality gives
$\|M_h\|\le\sum_m|\widehat a_m|$.
If $3\Lambda\le N$, it also gives
\begin{equation}\label{eq:loop-multiplication-comparison}
\begin{split}
 \|(1-P_{2\Lambda})M_\eta P_\Lambda\|
       &\le\sum_{|m|_\infty\ge\Lambda}|\widehat a_m|,\\
 \|P'_{2\Lambda}(M_{\eta'}J-JM_\eta)P_\Lambda\|
       &\le\sum_{|m|_\infty\le3\Lambda}
                         |\widehat a'_m-\widehat a_m|
                    +90\eta\Lambda\sum_m|\widehat a_m|.
\end{split}
\end{equation}
For the first inequality an output above $2\Lambda$ and an input
below $\Lambda$ have coefficient index at least $\Lambda$,
even after periodic reduction, since $N+1-\Lambda\ge\Lambda$.
For the second the coefficient indices lie below $3\Lambda$,
so no reduction is needed. Subtract the two phase factors and
use $|e^{iu}-e^{iv}|\le|u-v|$; their difference is at most
$90\eta\Lambda$. The row and column sums prove the claim.

Here are the particular multiplication terms in the Hessian.
Write $\theta=\operatorname{ad}t_3$, so $\|\theta\|=2$ and
$\theta^{\mathsf T}=-\theta$. On the commuting background
$U=\exp(ih a t_3)$,
\[
 (e_U\lambda)(b)=h^{-1}(e^{ih a(b)\theta}-1)\lambda(b_+).
\]
Thus its first derivative is $ia_\mu\theta$ followed by the
shift $he_\mu$. Its transpose has shifts $(-e_\mu,-e_\mu)$
and coefficient $-i\theta$. The mixed derivative is
$-h a_\mu^+a_\mu^-\theta^2$ followed by that shift, hence $O(h)$.

On a positively oriented $\mu\nu$ plaquette enumerate the bonds
as bottom, right, top, left, with signs $(1,1,-1,-1)$ and origins
$o_1=o_4=0$, $o_2=e_\mu$, $o_3=e_\nu$.
For their field values $a_i$, put
$\vartheta_1=0$, $\vartheta_2=a_1$,
$\vartheta_3=a_1+a_2-a_3$, $\vartheta_4=h\,da$.
Then
\[
 E_Ub=h^{-1}\sum_i\varsigma_i
                  (e^{ih\vartheta_i\theta}-1)b_i,
 \qquad z_U=\cos(h^2da).
\]
The first derivative of $E$ in this plane is
\[
 i\theta\{a_\mu(x)b_\nu(x+he_\mu)
 -[a_\mu(x)+a_\nu(x+he_\mu)-a_\mu(x+he_\nu)]
                         b_\mu(x+he_\nu)
 -h\,da(p)b_\nu(x)\}.
\]
This is a list of five multiplication terms per plane.
Its mixed derivative is
$-h\sum_i\varsigma_i\vartheta_i^+\vartheta_i^-\theta^2b_i$,
which is $O(h)$. The first derivatives of $z$ vanish and
$z_{11}=-h^4f_+f_-$.

Finally the remaining curvature term has quadratic form
\[
 (b,f_Ub)=-ih^3\sum_p h^{-2}\sin(h^2da(p))
 \sum_{i<j}\varsigma_i\varsigma_j
       \langle b_i,\theta
               e^{ih(\vartheta_j-\vartheta_i)\theta}b_j\rangle.
\]
Its first derivative is the symmetrisation of the matrix with
entries $-i\varsigma_i\varsigma_j da\,\theta$ from $b_j$ to
$b_i$. Its multiplication terms have offsets
$(-o_i,o_j-o_i)$, together with the transposed offsets
$(-o_j,o_i-o_j)$, and coefficients of norm at most one.
There are $36$ terms over the three planes.
The mixed derivative comes from differentiating the exponential
once and $da$ once; it is a sum of
$h(da)(\vartheta_j-\vartheta_i)\theta^2$ and its polarisation,
so is $O(h)$. The sine factor has no quadratic term.

Consequently the first jets of $e,e^{\mathsf T},E,f$ are finite
matched lists of the operators in
\eqref{eq:loop-multiplication-comparison}, with coefficients
$a^\pm_\mu$, $f^\pm_{\mu\nu}$ or $h f^\pm_{\mu\nu}$.
The lists and shifts are identical at both heights.
Lemma~\ref{lem:loop-background-comparison} controls their
coefficient sums and differences. Their mixed jets are $O(\eta)$
and $O(\eta')$; this last observation is what permits the
comparison without introducing additional second-order
background coefficients.

\subsection{Defects of the determinant data}

Let $1\le\Lambda$ and $3\Lambda\le N$, and put
$\epsilon_1=\Lambda^{-1/2}+\eta\Lambda+\eta^{1/2}$.
All constants in the following comparison are independent of
$k,p,\Lambda$ in their stated ranges.

\begin{lemma}\label{lem:loop-data-comparison}
For each of the vector and scalar data in
\eqref{eq:loop-Phi},
\[
 \operatorname{df}(\kappa',\kappa)
 +\operatorname{df}(x',x)+\|x'(1-JJ^*)\|
 +\operatorname{df}(y',y)\le C\epsilon_1.
\]
\end{lemma}

\begin{proof}
The free roots $G^{1/2},g^{1/2}$ have diagonal entries
$g_\eta(q_m)$. Denote this diagonal matrix by $g_*$, and the
matrix on the same modes with entries $g_{\eta'}(q_m)$ by
$g_\flat$. The elementary bounds imply
\begin{equation}\label{eq:loop-root-tails}
 \|(1-P_\Lambda)g_*\|,\|(1-P_\Lambda)g_\flat\|
                    \le2/\Lambda,
 \quad \|G'^{1/2}(1-JJ^*)\|\le4\eta,
 \quad\|g_\flat-g_*\|\le6\eta;
\end{equation}
the corresponding scalar statements hold as well.
Indeed a nonzero mode has $|q_m|\ge\pi|m|_\infty/2$,
and an added mode has $|m|_\infty\ge N+1\ge(2\eta)^{-1}$.

Consider the free blocks
$D_1=dG^{1/2}$, $D_0=d^{\mathsf T}G^{1/2}$,
$D_g=dg^{1/2}$ and their transposes. They are diagonal in the
Fourier modes, with matrices of at most $3\times3$ entries
$0$ or $\pm d^\eta_\mu(\pm q)g_\eta(q)$ and the same positions
and signs at both heights. Consequently
\begin{equation}\label{eq:loop-free-blocks}
 \|D\|,\|D'\|\le3,\quad
 D'J-JD=J\delta_D,\quad
 \|\delta_D\|\le6,\quad
 \|\delta_DP_\Lambda\|\le528\eta\Lambda.
\end{equation}
The last estimate follows from
$|d^{\eta'}_\mu g_{\eta'}-d^\eta_\mu g_\eta|
\le8\eta|q|$ and $|q|\le22\Lambda$ on these modes.

For a multiplication jet $M=E,e^{\mathsf T},f$ or $e$,
let $X=M G^{1/2}$ or $M g^{1/2}$ with the appropriate spaces.
We claim
\begin{equation}\label{eq:loop-X-comparison}
 \|X'(1-JJ^*)\|\le C\eta,\quad
 \|(1-P_{2\Lambda})X\|\le C\epsilon_1,\quad
 \operatorname{df}(X',X)\le C\epsilon_1.
\end{equation}
The first bound is the added-mode root bound times the bounded
multiplication jet. For the second split the input at
$P_\Lambda$: the high input costs $C/\Lambda$, the low input
followed by the high output costs $C\Lambda^{-1/2}$ by
\eqref{eq:loop-multiplication-comparison} and the coefficient
tail, and the mixed jet costs $C\eta$.
For the third use
\[
 X'J-JX=(M'J-JM)g_\flat+JM(g_\flat-g_*).
\]
The second term is $O(\eta)$. Split the first at $P_\Lambda$
on the input and $P'_{2\Lambda}$ on the output. The high input
costs $C/\Lambda$; the low input and high output are bounded
by the preceding tail argument at the two heights; on the
remaining modes \eqref{eq:loop-multiplication-comparison}
costs $C(\eta^{1/2}+\eta\Lambda)$.
This proves \eqref{eq:loop-X-comparison}.

Now put $X_1=EG^{1/2}$, $X_0=e^{\mathsf T}G^{1/2}$,
$X_f=fG^{1/2}$ and $X_g=eg^{1/2}$. The Hessian formula gives
\begin{align*}
 \kappa={}&D_1^{\mathsf T}(z-1)D_1
 +D_1^{\mathsf T}zX_1+X_1^{\mathsf T}zD_1
 +X_1^{\mathsf T}zX_1+G^{1/2}X_f\\
 &+D_0^{\mathsf T}X_0+X_0^{\mathsf T}D_0
                                      +X_0^{\mathsf T}X_0,\\
 \kappa_g={}&D_g^{\mathsf T}X_g+X_g^{\mathsf T}D_g
                                      +X_g^{\mathsf T}X_g.
\end{align*}
The terms containing $z-1$ are $O(\eta^4)$ at either height,
since its only nonzero component is $11$.
For $D^{\mathsf T}X$, split $X$ at $P_{2\Lambda}$ and use
\eqref{eq:loop-free-blocks} on the low modes and the output
tail in \eqref{eq:loop-X-comparison} on the other modes.
This gives defect $C\epsilon_1$. For $X^{\mathsf T}D$ apply
the transpose rule to $D^{\mathsf T}X$, retaining its added-mode
tail, bounded by $3\|X'(1-JJ^*)\|$. Products
$X^{\mathsf T}X$ and $G^{1/2}X_f$ follow from the product
and transpose rules. These observations prove the claims for
$\kappa$ and $\kappa_g$.

For $x$ the exact split \eqref{eq:loop-actual-split} gives
\begin{align*}
 x'J-x={}&(\mathbf Q^\sharp-\mathbf Q)
                              \mathbf Q^{\rm f}G'^{1/2}J\\
 &+\mathbf Q(\mathbf Q^{\rm f}-Q^{\rm f})G'^{1/2}J
 +\mathbf Q(Q^{\rm f}G'^{1/2}J-G^{1/2}).
\end{align*}
The first two terms cost $C\eta^{1/2}$ and $C\eta$.
The last difference is diagonal, with entries
$u_{1,\mu}g_{\eta'}-g_\eta$, of modulus at most $11\eta$:
use the bounds for $u_1-1$, $g_{\eta'}$ and $g_{\eta'}-g_\eta$.
The scalar version uses $v_1g_{\eta'}-g_\eta$, bounded by
$8\eta$. The added-mode tails of $x'$ and $x'_g$ are bounded
by $4\eta$ times the uniform averaging bound. This proves
the assertions concerning $x$.

For $y$, set
$X'=G'^{1/2}\mathbf D'\mathbf R^{{\rm f}\mathsf T}$ and
$X=G^{1/2}\mathbf D$. Then
\begin{align*}
 y'-Jy={}&X'(\mathbf R^\sharp-\mathbf R)^{\mathsf T}
 +(X'-JX)P_\Lambda\mathbf R^{\mathsf T}\\
 &+(X'-JX)(1-P_\Lambda)\mathbf R^{\mathsf T}.
\end{align*}
The first term costs $C\eta^{1/2}$; the last costs
$C\Lambda^{-1/2}$ by \eqref{eq:loop-R-tail}. To estimate the
middle term, expand $X'-JX$ as
\begin{align*}
 &[G'^{1/2}d'R^{{\rm f}\mathsf T}-JG^{1/2}d]\\
 &\quad+G'^{1/2}\mathbf D'
                    (\mathbf R^{\rm f}-R^{\rm f})^{\mathsf T}\\
 &\quad+G'^{1/2}e'(R^{{\rm f}\mathsf T}-J)
       +[G'^{1/2}e'J-JG^{1/2}e].
\end{align*}
The second line is $O(\eta)$ by \eqref{eq:loop-step-jets}.
For the free aliases at $|m|_\infty\le\Lambda$, the principal
coefficient satisfies $|v_1(q_m)-1|\le11\eta\Lambda$.
For any other alias, some component of $q_{m'}$ has modulus
at least $\pi\eta^{-1}$; the sine quotient defining $s_1$
then gives $|v_1(q_{m'})|\le6\eta\Lambda$.
There are at most $L^3$ aliases. Hence
\[
 \|(R^{{\rm f}\mathsf T}-J)P_\Lambda\|
                              \le11L^{3/2}\eta\Lambda.
\]
The same calculation, with the normalized derivative, bounds
the first line by $187\sqrt3 L^{3/2}\eta\Lambda$;
on the principal alias use
$|d^{\eta'}g_{\eta'}-d^\eta g_\eta|\le8\eta|q|$,
and on the others use $|d^{\eta'}g_{\eta'}|\le1$.
The third line's first term is thus $O(\eta\Lambda)$.
For its last term, split the output at $P'_{2\Lambda}$.
On the low output, the multiplication comparison and the root
difference give $C(\eta^{1/2}+\eta\Lambda+\eta)$; on the
high output the left root gives $C/\Lambda$.
This proves $\operatorname{df}(y',y)\le C\epsilon_1$ and
completes the proof.
\end{proof}

\subsection{Stability of the finite expression}

The coarse traces in \eqref{eq:loop-Phi} have dimension at most
nine. Its one remaining fine trace is controlled instead by
the three propagators in it. Write $\|\cdot\|_3$ for the
Schatten norm of exponent three on a Fourier fibre.

\begin{lemma}\label{lem:loop-Phi-stability}
Suppose two data for \eqref{eq:loop-Phi} satisfy the following
bounds, with $A,C_3,\iota\ge1$ and $\epsilon>0$.
At every real $\varkappa\in K$ the diagonal propagators satisfy
\[
 \|\mathsf G\|_3,\|\mathsf G'\|_3\le C_3,\qquad
 \|\mathsf G'(1-JJ^*)\|_3,
 \|(\mathsf G'J-J\mathsf G)J^*\|_3\le\epsilon.
\]
The jets $\kappa,x,y,B_t,B^\flat_\tau$ have norms at most $A$
at both heights, for $t,\tau\in[0,1]$, and the $00$ inverses
of $xB^\flat_0x^{\mathsf T}$ have norms at most $\iota$.
Finally
\[
 \operatorname{df}(\kappa',\kappa),\quad
 \operatorname{df}(x',x),\quad\|x'(1-JJ^*)\|,\quad
 \operatorname{df}(y',y)\le\epsilon.
\]
Then, with $a=\max(A,C_3)$,
\[
 |\Phi[\mathsf G',\kappa',x',y']
          -\Phi[\mathsf G,\kappa,x,y]|
                   \le2600\iota^3a^{11}\epsilon.
\]
\end{lemma}

\begin{proof}
The propagator operator norms are at most $C_3$, and their
defect is at most $\epsilon$. The transpose defects of $x,y$
are at most $2\epsilon,\epsilon$, respectively; for $y$ its
input is coarse and there is no added input space. The product
and inverse rules therefore give
\[
 \operatorname{df}(\omega',\omega)\le5A\epsilon,
 \quad\operatorname{df}(B'_t,B_t)\le6A^3\epsilon,
 \quad\operatorname{df}(B^{\flat\prime}_\tau,B^\flat_\tau)
                                      \le7A^3\epsilon.
\]
A product of $n$ displayed factors has defect at most
$7na^{n+2}\epsilon$.

For an endomorphism of a coarse fibre, the trace of its $11$
component is bounded by nine times its jet norm. The six
coarse-trace products in \eqref{eq:loop-Phi} have respectively
$3,3,5,5,7,7$ factors. Remembering that
$\int_0^1(1-\tau)\,d\tau=1/2$, their total difference is at most
\[
 9(42a^5+70a^7+49a^9)\epsilon\le1449a^9\epsilon.
\]
For the logarithmic term put $C=xB^\flat_0x^{\mathsf T}$.
Then $\|C\|,\|C'\|\le a^3$ and
$e:=\|C'-C\|\le21a^5\epsilon$.
The inverse identity bounds the difference of $C_{00}^{-1}$
by $\iota^2e$. Subtracting the two products in the definition
of $\ell$ one factor at a time gives at most
$6\iota^3a^6e$ before the trace, hence
$1134\iota^3a^{11}\epsilon$ after integration.

It remains to compare
$W=\mathsf GB_0\mathsf GB_0\mathsf GB^\flat_\tau$.
Insert the projection onto the common modes in the trace:
\[
 \operatorname{tr}W'_{11}-\operatorname{tr}W_{11}
 =\operatorname{tr}[W'_{11}(1-JJ^*)]
       +\operatorname{tr}[(W'J-JW)_{11}J^*].
\]
In the first term cycle the projection next to a propagator.
H\"older's inequality for the three Schatten-three factors
bounds it by $\epsilon C_3^2A^3\le a^5\epsilon$.
In the second term telescope the six factors. Insert $J^*J=1$
between common-space factors so that each term is a product
on the primed space. The three propagator-difference terms
are bounded by $3a^5\epsilon$ in total. The three inverse-jet
differences, bounded by $6A^3\epsilon,6A^3\epsilon,
7A^3\epsilon$, contribute at most $19a^8\epsilon$.
In each case the other three propagator factors, including
a difference when present, are estimated in Schatten norm;
the jet product norm sums the finitely many $11$ terms.
The remaining integral has weight $(1-\tau)^2$, whose integral
is $1/3$. Its contribution is therefore at most $8a^8\epsilon$.
Adding the three bounds and using $a,\iota\ge1$ proves the lemma.
\end{proof}

\begin{proposition}\label{prop:loop-real-comparison}
There is a finite $C_{\rm re}=C_{\rm re}(L)$ such that
\[
 |F_{k+1}(p)-F_k(p)|\le C_{\rm re}L^{-k/3}
 \quad(k\ge1,\ p\in\mathbb R^3,\ |p_\mu|\le\delta).
\]
\end{proposition}

\begin{proof}
Choose $\Lambda=\lfloor\eta^{-2/3}/4\rfloor$.
For $\eta^{-1}\ge13$ this is at least one, is at least
$\eta^{-2/3}/8$, and satisfies $3\Lambda\le N$.
For the last assertion the function
$(y-1)/2-3y^{2/3}/4$ is positive at $y=13$ and increasing
thereafter. Thus
\[
 \epsilon_1\le(\sqrt8+1/4+1)\eta^{1/3}
                                      \le5\eta^{1/3}.
\]
Lemma~\ref{lem:loop-data-comparison} supplies the operator
defects in Lemma~\ref{lem:loop-Phi-stability} with
$\epsilon=C\eta^{1/3}$. The uniform inverse and jet bounds
were established before \eqref{eq:loop-Phi}.

We verify the required Schatten bounds explicitly. Both
propagators have eigenvalues $g_h(q_m)^2$, each repeated
at most nine times. Since $g_h\le\min(1,2/|m|_\infty)$
away from zero and the shell $|m|_\infty=n$ has
$24n^2+2\le26n^2$ points,
\[
 9\sum_mg_h(q_m)^6
       \le9\left(1+64\cdot26\sum_{n\ge1}n^{-4}\right).
\]
This gives a uniform $C_3$. The sum over added modes starts
beyond $N\ge6\eta^{-1}/13$, and its tail is $O(\eta^3)$.
On common modes \eqref{eq:loop-elementary-bounds} gives
$|g_{\eta'}^2-g_\eta^2|\le3\eta^2$; there are $\eta^{-3}$
such modes, so their cubed Schatten difference is at most
$9\eta^{-3}(3\eta^2)^3=243\eta^3$.
Both relevant Schatten norms are therefore $O(\eta)$,
and hence at most $C\eta^{1/3}$.

Apply Lemma~\ref{lem:loop-Phi-stability} to the vector and
scalar data and use Lemma~\ref{lem:loop-finite-expression}.
This proves the stated bound. In particular the choice
\[
 C_{\rm re}:=\sup_{k\ge1,\ |p_\mu|\le\delta,\ p\in\mathbb R^3}
       L^{k/3}|F_{k+1}(p)-F_k(p)|
\]
is finite and depends only on $L$.
\end{proof}

\subsection{Recovering the coefficient and summing the increments}

The real-momentum bound does not by itself bound derivatives.
We use the uniform complex bound from
\eqref{eq:loop-remainder-reduction} as follows. If $f$ is
holomorphic on $|p_\mu|<\delta$, bounded by $M$ there and
by $\epsilon$ on the real cube $|p_\mu|\le\delta/4$, then
\begin{equation}\label{eq:loop-finite-difference}
 |\partial_1\partial_2f(0)|
                  \le17\sqrt{\epsilon M}\,\delta^{-2}.
\end{equation}
Indeed expand $f(s,t,0)=\sum c_{ij}s^it^j$, with
$|c_{ij}|\le M\delta^{-i-j}$. The centred mixed difference
at the four points $(\pm h,\pm h,0)$ divided by $4h^2$
retains exactly the odd powers of $s,t$. For $h\le\delta/4$
the terms other than $c_{11}$ total at most
$3Mh^2\delta^{-4}$. Therefore
\[
 |c_{11}|\le\epsilon h^{-2}+3Mh^2\delta^{-4}.
\]
When $0<\epsilon<M$, choose
$h^2=\delta^2\sqrt{\epsilon/M}/16$.
When $\epsilon\ge M$, the ordinary Cauchy bound suffices;
when $\epsilon=0$, let $h$ tend to zero. These cases prove
\eqref{eq:loop-finite-difference}.

\begin{proof}[Proof of Theorem~\ref{III-12-105}]
Apply \eqref{eq:loop-finite-difference} to $F_{k+1}-F_k$,
with $M=2C_{\rm rem}$ and
$\epsilon=C_{\rm re}L^{-k/3}$. By
\eqref{eq:loop-remainder-reduction},
\[
 C_{\rm inc}:=\sup_{k\ge1}L^{k/6}|r_{k+1}-r_k|
 \le17(2C_{\rm rem}C_{\rm re})^{1/2}\delta^{-2}<\infty.
\]
Define
\begin{equation}\label{eq:loop-geometric-constant}
 C_\circ=\max\{C_{\rm inc}+2C_E,c_2^{\rm loop}\},
 \qquad C_{\rm geo}=\frac{2C_\circ}{1-L^{-1/6}}.
\end{equation}
These are functions of $L$ alone: each is defined from the
Gaussian coefficients, their fixed massive term and their
fixed determinant remainder.
Lemma~\ref{lem:loop-explicit-limit} gives
\[
 |c_{k+1}-c_k|\le C_\circ L^{-k/6}\qquad(k\ge1).
\]
The same bound holds at $k=0$ because $c_0=0$ and
$|c_1|\le c_2^{\rm loop}$. For $n>k\ge0$ it follows that
\[
 |c_n-c_k|\le\sum_{j\ge k}C_\circ L^{-j/6}
                         =\tfrac12C_{\rm geo}L^{-k/6}.
\]
Thus $c_k$ converges to a real $c_\infty$, with the asserted
error, and $|c_\infty|\le c_2^{\rm loop}$.
Also $r_k\to c_\infty-E_\infty$.

The defining identity for $c_k$ gives
$b_l^0=c_l-L^{-1}c_{l-1}$ for $l\ge1$. Therefore
\[
 |b_{l+1}^0-b_l^0|
 \le C_\circ(1+L^{-5/6})L^{-l/6}
 \le(1-L^{-1/6})C_{\rm geo}L^{-l/6}.
\]
This includes $l=1$ by the already verified $k=0$ bound.
The limit is $(1-L^{-1})c_\infty$; summing this last
increment estimate from $l$ to infinity gives
$|b_l^0-b_\infty^0|\le C_{\rm geo}L^{-l/6}$, as claimed.
\end{proof}


\subsection{The bound used in the coupling flow}

The coupling recursion needs a bound on the increments, not only
on their accumulated coefficients. The following elementary form
also applies to the coefficients of a general group once their
stated bound has been established.

\begin{lemma}\label{lem:gaussian-flow-bound}
Let $L\ge13$, let $c_0=0$ and $|c_k|\le c_2^{\rm loop}$ for all
$k\ge0$, and put $b_l^0=c_l-L^{-1}c_{l-1}$. If
$B_\sharp\ge(1+L^{-1})c_2^{\rm loop}+1$ and
$\bar B=LB_\sharp/(L-1)$, then
\[
 B_\sharp\ge1,\qquad |b_l^0|\le B_\sharp-1\quad(l\ge1),
 \qquad B_\sharp<\bar B\le\tfrac{13}{12}B_\sharp.
\]
\end{lemma}

\begin{proof}
The triangle inequality gives
$|b_l^0|\le(1+L^{-1})c_2^{\rm loop}\le B_\sharp-1$.
Since $c_2^{\rm loop}\ge0$, we have $B_\sharp\ge1$.
Finally $1<L/(L-1)\le13/12$, which proves the bounds on $\bar B$.
\end{proof}


\begin{thebibliography}{9}
\bibitem{YMM1} \emph{Scalar fields from three-dimensional lattice Yang--Mills theory}, companion manuscript \href{run:YM3_I.pdf}{YMM1}.
\bibitem{StoyanovLin2013} J. Stoyanov and G. D. Lin, \emph{Hardy's condition in the moment problem for probability distributions}, Theory Probab. Appl. \textbf{57} (2013), no.~4, 699--708. \href{https://doi.org/10.1137/S0040585X9798631X}{doi:10.1137/S0040585X9798631X}. \bibitem{Adamczak2014} R. Adamczak, \emph{A note on the Hanson--Wright inequality for random vectors with dependencies}, \href{https://arxiv.org/abs/1409.8457v1}{arXiv:1409.8457v1} (2014). \bibitem{KolesnikovMilman2016} A. V. Kolesnikov and E. Milman, \emph{Riemannian metrics on convex sets with applications to Poincar\'e and log-Sobolev inequalities}, \href{https://arxiv.org/abs/1510.02971v2}{arXiv:1510.02971v2} (2016).
\bibitem{LawlerLimic} G. F. Lawler and V. Limic,
\emph{Random Walk: A Modern Introduction}, Cambridge University Press,
2010. \href{https://www.math.uchicago.edu/~lawler/srwbook.pdf}{Author-hosted text},
Theorem~6.3.8(a).
\bibitem{Nussbaum1965} A. E. Nussbaum, \emph{Quasi-analytic vectors},
Ark. Mat. \textbf{6} (1965), 179--191.
\href{https://doi.org/10.1007/BF02591357}{doi:10.1007/BF02591357}.
\end{thebibliography}
\end{document}
